%% Éléments de géométrie algébrique : texte français cumulatif complet.
%% Ce fichier contient six documents LaTeX autonomes, couverture comprise.
%% Extraire avec reconstruction/split_complete.py ; compiler les composantes
%% et assembler leurs PDF avec BUILD_ALL.ps1 dans le projet éditable joint.
%% Les six préambules sont conservés pour préserver les mises en pages natives.
%% Ne pas compiler ce conteneur multi-document directement avec pdfLaTeX.

% EGA-COMPONENT-BEGIN COVER 8B049F2A6693F5DFEDC20B5DAC72844BE33D77D0A8C1AC2B01D1C0D8C3A955F6
\documentclass[11pt,a4paper]{article}
\usepackage[T1]{fontenc}
\usepackage[utf8]{inputenc}
\usepackage[french]{babel}
\usepackage{lmodern}
\usepackage[margin=27mm]{geometry}
\usepackage{xcolor}
\definecolor{navy}{RGB}{23,49,71}
\definecolor{muted}{RGB}{78,89,98}
\newcommand{\editionlink}[2]{\leavevmode\pdfstartlink user{/Subtype /Link /A << /S /URI /URI (#1) >>}\textcolor{navy}{#2}\pdfendlink}
\pagestyle{empty}
\begin{document}
\color{navy}
\noindent\textsc{Bibliothèque de géométrie algébrique}\par
\vspace{22mm}
\noindent{\fontsize{32}{38}\selectfont\bfseries Éléments de géométrie\par algébrique}\par
\vspace{9mm}
\noindent{\Large Alexandre Grothendieck et Jean Dieudonné}\par
\vspace{18mm}
\noindent{\LARGE Édition française}\par
\vspace{3mm}
\noindent{\large Texte diplomatique --- lecteur cumulatif EGA I--IV}\par
\vspace{10mm}
\color{muted}
\noindent Cette édition réunit les volumes I, II, III-1, III-2 et IV-1 à IV-4,
y compris les préliminaires (EGA~0) qui leur sont associés.
Les signets et les références internes permettent de parcourir le lecteur.\par
\vspace{5mm}
\noindent Édition diplomatique établie à partir des volumes imprimés
diffusés par NUMDAM.\par
\vspace{9mm}
\noindent\editionlink{https://doi.org/10.5281/zenodo.21921588}{Édition française historique : DOI 10.5281/zenodo.21921588}\par
\vspace{3mm}
\noindent\editionlink{https://doi.org/10.5281/zenodo.21921591}{Édition anglaise historique : DOI 10.5281/zenodo.21921591}\par
\vspace{3mm}
\noindent\editionlink{https://github.com/KokunoYumeto/ega-fr}{Sources, fichiers LaTeX et versions : github.com/KokunoYumeto/ega-fr}\par
\vfill
\small
\noindent Préparation des sources et assemblage : OpenAI Codex,
GPT-6 Astra, effort Ultra.\par
\vspace{3mm}
\noindent Le PDF est accompagné du texte LaTeX et du projet éditable,
avec les fichiers nécessaires à la compilation.\par

% BEGIN RIGHTS AND AI TRANSCRIPTION NOTICE
\clearpage
\begingroup
\small
\noindent\textbf{Droits, disponibilité et provenance}

Les droits sur les œuvres originales et sur toute contribution protégée à ces traductions et transcriptions dérivées restent ceux de leurs titulaires respectifs. Ce projet ne place pas ces œuvres ou contributions dans le domaine public et n'accorde aucune licence sur elles. Les mentions de droits de tiers et les autorisations déjà accordées restent applicables.

Ces éditions sont mises à disposition gratuitement, à des fins non commerciales et exclusivement pédagogiques, dans la présomption que les titulaires de droits ne s'opposeraient pas à la libre disponibilité de ces traductions et transcriptions aux côtés des originaux français sur NUMDAM (www.numdam.org/journals/PMIHES/). Cette présomption ne signifie pas qu'une autorisation a été obtenue ; le libre accès aux originaux n'autorise pas, à lui seul, leur traduction, leur adaptation ou leur republication.

Si vous détenez des droits sur un élément présenté ici et vous opposez à sa disponibilité, veuillez ouvrir une demande sur GitHub (github.com/KokunoYumeto/ega-fr/issues). Le responsable de la publication retirera l'élément concerné à la demande du titulaire des droits.

La transcription moderne, les corrections, la composition et les notes mathématiques sont produites par IA. Les travaux documentés comprennent ceux d'OpenAI Codex, GPT-6 Astra, effort Ultra. Les indications de provenance des contributions antérieures restent applicables. Les textes originaux restent l'œuvre d'Alexandre Grothendieck et Jean Dieudonné ; ils ne sont pas attribués à l'IA. Cette édition n'est approuvée ni par les auteurs, ni par l'IHÉS, ni par NUMDAM.
\par
\endgroup
\clearpage
% END RIGHTS AND AI TRANSCRIPTION NOTICE
\end{document}
% EGA-COMPONENT-END COVER

% EGA-COMPONENT-BEGIN BASE 0ADC35F981F1EC3DB0259067D5A6DCCE4C615D3B7ADAA24B1DBDB8D72E1A7EB4
\documentclass[10pt,oneside]{book}

\usepackage[T1]{fontenc}
\usepackage[utf8]{inputenc}
\usepackage[french]{babel}
\usepackage{amsmath,amssymb,mathrsfs}
\usepackage{longtable}
\usepackage[all]{xy}
\usepackage{tikz-cd}
\usepackage[a4paper,margin=22mm]{geometry}
\usepackage[
  colorlinks=true,
  linkcolor=blue,
  hypertexnames=false,
  unicode=true,
  pdfencoding=auto
]{hyperref}
%
% DEBUT SOURCE preamble-base.tex
\def\sh{\mathcal}                   % sheaf font
\def\bb{\mathbf}                    % bold font
\def\cat{\mathtt}                   % category font
\def\leq{\leqslant}                 % <=
\def\geq{\geqslant}                 % >=
\def\setmin{-}                      % set minus
\def\rad{\mathfrak{r}}              % radical
\def\nilrad{\mathfrak{N}}           % nilradical
\def\emp{\varnothing}               % empty set
\def\vphi{\phi}                     % for switching \phi and \varphi, change if needed
\def\HH{\mathrm{H}}                 % cohomology H
\def\CHH{\check{\HH}}               % Čech cohomology H
\def\RR{\mathrm{R}}                 % right derived R
\def\LL{\mathrm{L}}                 % left derived L
\def\dual#1{{#1}^\vee}              % dual
\def\kres{k}                        % residue field k
\def\C{\cat{C}}                     % category C
\def\op{^\cat{op}}                  % opposite category
\def\Set{\cat{Set}}                 % category of sets
\def\CHom{\cat{Hom}}                % functor category
\def\supertilde{{\,\widetilde{\,}\,}}   % use \supertilde instead of ^\sim
\def\GL{\bb{GL}}
\def\bbGamma{\boldsymbol\Gamma}
\def\red{\mathrm{red}}
\def\rg{\operatorname{rg}}
\def\gr{\operatorname{gr}}
\def\Gr{\operatorname{Gr}}
\def\Hom{\operatorname{Hom}}
\def\Proj{\operatorname{Proj}}
\def\Tor{\operatorname{Tor}}
\def\Ext{\operatorname{Ext}}
\def\Supp{\operatorname{Supp}}
\def\Ker{\operatorname{Ker}}
\def\Im{\operatorname{Im}}
\def\Coker{\operatorname{Coker}}
\def\Spec{\operatorname{Spec}}
\def\Spm{\operatorname{Spm}}
\def\Spf{\operatorname{Spf}}
\def\grad{\operatorname{grad}}
\def\dimc{\operatorname{dimc}}
\def\codim{\operatorname{codim}}
\def\id{\operatorname{id}}
\def\Der{\operatorname{Der}}
\def\Diff{\operatorname{Diff}}
\def\Hyp{\operatorname{Hyp}}
\def\Tr{\operatorname{Tr}}

\renewcommand{\to}{\mathchoice{\longrightarrow}{\rightarrow}{\rightarrow}{\rightarrow}}
\newcommand{\from}{\mathchoice{\longleftarrow}{\leftarrow}{\leftarrow}{\leftarrow}}
\let\mapstoo\mapsto
\renewcommand{\mapsto}{\mathchoice{\longmapsto}{\mapstoo}{\mapstoo}{\mapstoo}}
\def\isoto{\simeq} % isomorphism
\def\simto{\xrightarrow{\sim}}

\renewcommand{\hat}[1]{\widehat{#1}}
\renewcommand{\tilde}[1]{\widetilde{#1}}

% FIN SOURCE preamble-base.tex

\hypersetup{
  pdftitle={Éléments de géométrie algébrique -- édition française cumulative},
  pdfauthor={Alexander Grothendieck et Jean Dieudonné},
  pdfsubject={Texte diplomatique français cumulatif depuis l'autorité NUMDAM}
}

\renewcommand{\thesection}{\arabic{section}}
\renewcommand{\thesubsection}{\thesection.\arabic{subsection}}
\newcommand{\oldpage}[2][]{%
  \setcounter{footnote}{0}%
  \marginpar{\scriptsize\textbf{#1}~|~#2}\ignorespaces}
\newenvironment{env}[1][]{\par\medskip\noindent\textbf{(#1)}\ }{\par}
\newenvironment{definition}[1][]{%
  \par\medskip\noindent\itshape Définition (#1). ---\ }{\par}
\newenvironment{proposition}[1][]{%
  \par\medskip\noindent\itshape Proposition (#1). ---\ }{\par}
\newenvironment{theorem}[1][]{%
  \par\medskip\noindent\itshape Théorème (#1). ---\ }{\par}
\newenvironment{corollary}[1][]{%
  \par\medskip\noindent\itshape Corollaire (#1). ---\ }{\par}
\newenvironment{remark}[1][]{%
  \par\medskip\noindent\itshape Remarque (#1). ---\ }{\par}
\newenvironment{remarks}[1][]{%
  \par\medskip\noindent\itshape Remarques (#1). ---\ }{\par}
\newenvironment{scholium}[1][]{%
  \par\medskip\noindent\textit{Scholie (#1). ---\ }}{\par}
\newenvironment{lemma}[1][]{%
  \par\medskip\noindent\itshape Lemme (#1). ---\ }{\par}
\newenvironment{example}[1][]{%
  \par\medskip\noindent\textit{Exemple (#1). ---\ }}{\par}
\newenvironment{examples}[1][]{%
  \par\medskip\noindent\textit{Exemples (#1). ---\ }}{\par}
\newenvironment{notation}[1][]{%
  \par\medskip\noindent\textit{Notations (#1). ---\ }}{\par}
\newenvironment{proof}{%
  \par\medskip\noindent\textit{Démonstration. ---\ }}{\par}
\newenvironment{namedtheorem}[2][]{%
  \par\medskip\noindent\itshape #2 (#1). ---\ }{\par}
\newenvironment{namedlemma}[2][]{%
  \par\medskip\noindent\itshape #2 (#1). ---\ }{\par}

\makeatletter
\newcommand{\EGAReaderTarget}[2]{\hypertarget{#1}{}\@bsphack\begingroup\def\@currentHref{#1}\label{#2}\endgroup\@esphack}
\makeatother
\begin{document}
% BEGIN RIGHTS AND AI TRANSCRIPTION NOTICE
\clearpage
\begingroup
\small
\noindent\textbf{Droits, disponibilité et provenance}

Les droits sur les œuvres originales et sur toute contribution protégée à ces traductions et transcriptions dérivées restent ceux de leurs titulaires respectifs. Ce projet ne place pas ces œuvres ou contributions dans le domaine public et n'accorde aucune licence sur elles. Les mentions de droits de tiers et les autorisations déjà accordées restent applicables.

Ces éditions sont mises à disposition gratuitement, à des fins non commerciales et exclusivement pédagogiques, dans la présomption que les titulaires de droits ne s'opposeraient pas à la libre disponibilité de ces traductions et transcriptions aux côtés des originaux français sur NUMDAM (www.numdam.org/journals/PMIHES/). Cette présomption ne signifie pas qu'une autorisation a été obtenue ; le libre accès aux originaux n'autorise pas, à lui seul, leur traduction, leur adaptation ou leur republication.

Si vous détenez des droits sur un élément présenté ici et vous opposez à sa disponibilité, veuillez ouvrir une demande sur GitHub (github.com/KokunoYumeto/ega-fr/issues). Le responsable de la publication retirera l'élément concerné à la demande du titulaire des droits.

La transcription moderne, les corrections, la composition et les notes mathématiques sont produites par IA. Les travaux documentés comprennent ceux d'OpenAI Codex, GPT-6 Astra, effort Ultra. Les indications de provenance des contributions antérieures restent applicables. Les textes originaux restent l'œuvre d'Alexandre Grothendieck et Jean Dieudonné ; ils ne sont pas attribués à l'IA. Cette édition n'est approuvée ni par les auteurs, ni par l'IHÉS, ni par NUMDAM.
\par
\endgroup
\clearpage
% END RIGHTS AND AI TRANSCRIPTION NOTICE

%
% DEBUT SOURCE ega1/frontmatter-fr.tex
\begin{titlepage}
\centering

{\Large INSTITUT\par}
\vspace{0.4em}
{\Large DES HAUTES ÉTUDES\par}
\vspace{0.4em}
{\Large SCIENTIFIQUES\par}

% The printed title page bears the graphic device of the Institut here.
\vfill

{\Large\bfseries ÉLÉMENTS DE GÉOMÉTRIE ALGÉBRIQUE\par}
\vspace{1.5em}
{\large\itshape par A. GROTHENDIECK\par}
\vspace{0.7em}
{\itshape Rédigés avec la collaboration de J. DIEUDONNÉ\par}

\vspace{2em}
I

\vspace{1.5em}
{\Large\bfseries LE LANGAGE DES SCHÉMAS\par}

\vfill
{\Large 1960\par}
\vspace{1.2em}
{\Large PUBLICATIONS MATHÉMATIQUES, N\textsuperscript{o} 4\par}
\vspace{0.5em}
{\large 5, ROND-POINT BUGEAUD --- PARIS (XVI\textsuperscript{e})\par}
\end{titlepage}

\thispagestyle{empty}
\vspace*{\fill}
\begin{center}
\fbox{%
  \begin{minipage}{0.50\textwidth}
  \centering
  DÉPOT LÉGAL

  \smallskip
  1\textsuperscript{re} édition \dotfill 3\textsuperscript{e} trimestre 1960

  \smallskip
  TOUS DROITS

  réservés pour tous pays

  \smallskip
  \copyright\ 1960, \emph{Institut des Hautes Études Scientifiques}
  \end{minipage}%
}
\end{center}
\vspace*{\fill}
\clearpage

% FIN SOURCE ega1/frontmatter-fr.tex

%
% DEBUT SOURCE ega1/intro-fr.tex
\section*{Introduction}

\begin{center}
\emph{A Oscar Zariski et André Weil.}
\end{center}

\oldpage[I]{5}

Ce mémoire, et les nombreux autres qui doivent lui faire suite, sont destinés à former un traité sur les fondements de la Géométrie algébrique.
Ils ne présupposent en principe aucune connaissance particulière de cette discipline, et il s'est même avéré qu'une telle connaissance, malgré ses avantages évidents, pouvait parfois (par l'habitude trop exclusive du point de vue birationnel qu'elle implique) être nuisible à celui qui désire se familiariser avec le point de vue et les techniques exposés ici.
Par contre, nous supposerons que le lecteur a une bonne connaissance des sujets suivants~:

\begin{enumerate}
  \item[a)] L'\emph{Algèbre commutative}, telle qu'elle est exposée par exemple dans les volumes en cours de préparation des \emph{Éléments} de N.~Bourbaki (et, en attendant la parution de ces volumes, dans Samuel-Zariski~\cite{I-13} et Samuel~\cite{I-11},~\cite{I-12}).
  \item[b)] L'\emph{Algèbre homologique}, pour laquelle nous renvoyons à Cartan-Eilenberg~\cite{I-2} (cité~(M)) et Godement~\cite{I-4} (cité~(G)), ainsi qu'à l'article récent de A.~Grothendieck~\cite{I-6} (cité~(T)).
  \item[c)] La \emph{Théorie des faisceaux}, où nos principales références seront~(G) et~(T); cette dernière théorie fournit le langage indispensable pour interpréter en termes <<~géométriques~>> les notions essentielles de l'Algèbre commutative, et pour les <<~globaliser~>>.
  \item[d)] Enfin, il sera utile au lecteur d'avoir une certaine familiarité avec le \emph{langage fonctoriel}, qui sera constamment employé dans ce Traité, et pour lequel le lecteur pourra consulter~(M),~(G) et surtout~(T); les principes de ce langage et les principaux résultats de la théorie générale des foncteurs seront exposés plus en détail dans un ouvrage en cours de préparation par les auteurs de ce Traité.
\end{enumerate}

\bigskip
\begin{center}*\quad*\quad*\end{center}
\bigskip

Ce n'est pas le lieu, dans cette Introduction, de donner une description plus ou moins sommaire du point de vue des <<~schémas~>> en Géométrie algébrique, ni la longue liste des raisons qui ont rendu nécessaire son adoption, et en particulier l'acceptation systématique d'éléments nilpotents dans les anneaux locaux des <<~variétés~>> que nous considérons (ce qui, nécessairement, relègue au second plan la notion d'application rationnelle, au profit de celle d'application régulière ou <<~morphisme~>>).
Le présent Traité vise précisément à développer de façon systématique le langage des <<~schémas~>> et démontrera, nous l'espérons, sa nécessité.
Encore qu'il serait facile de le faire, nous

\oldpage[I]{6}
n'essayerons pas non plus de donner ici une introduction <<~intuitive~>> aux notions développées dans le chapitre premier.
Le lecteur qui désirerait avoir un aperçu préliminaire des matières de ce Traité pourra se reporter à la conférence faite par A.~Grothendieck au Congrès international des Mathématiciens à Edinburgh en 1958~\cite{I-7}, et à l'exposé~\cite{I-8} du même auteur.
Le travail~\cite{I-14} (cité~(FAC)) de J.-P.~Serre peut aussi être considéré comme un exposé intermédiaire entre le point de vue classique et le point de vue des schémas en Géométrie algébrique, et à ce titre, sa lecture peut constituer une excellente préparation à celle de nos \emph{Éléments}.

\bigskip
\begin{center}*\quad*\quad*\end{center}
\bigskip

A titre informatif, nous donnons ci-dessous le plan général prévu pour ce Traité, d'ailleurs sujet à modifications ultérieures, surtout en ce qui concerne les derniers chapitres~:

\medskip
\begin{tabular}{rl}
  Chapitre Premier. & --- Le langage des schémas.\\
  --- \textbf{II}. & --- Étude globale élémentaire de quelques classes de morphismes.\\
  --- \textbf{III}. & --- Cohomologie des faisceaux algébriques cohérents. Applications.\\
  --- \textbf{IV}. & --- Étude locale des morphismes.\\
  --- \textbf{V}. & --- Procédés élémentaires de construction de schémas.\\
  --- \textbf{VI}. & --- Technique de descente. Méthode générale de construction des schémas.\\
  --- \textbf{VII}. & --- Schémas de groupes, espaces fibrés principaux.\\
  --- \textbf{VIII}. & --- Étude différentielle des espaces fibrés.\\
  --- \textbf{IX}. & --- Le groupe fondamental.\\
  --- \textbf{X}. & --- Résidus et dualité.\\
  --- \textbf{XI}. & --- Théories d'intersection, classes de Chern, théorème de Riemann-Roch.\\
  --- \textbf{XII}. & --- Schémas abéliens et schémas de Picard.\\
  --- \textbf{XIII}. & --- Cohomologie de Weil.\\
\end{tabular}
\medskip

En principe, tous les chapitres sont considérés comme ouverts, et des paragraphes supplémentaires pourront toujours leur être ajoutés ultérieurement; de tels paragraphes paraîtront en fascicules séparés, pour diminuer les inconvénients du mode de publication adopté.
Lorsqu'un tel paragraphe est prévu ou en préparation au moment de la publication d'un chapitre, il sera mentionné dans le sommaire dudit chapitre, même si en raison de certains ordres d'urgence sa publication effective devait être nettement postérieure.
Pour la commodité du lecteur, nous donnons dans un <<~Chapitre~0~>> des compléments divers d'Algèbre commutative, d'Algèbre homologique, de Théorie des faisceaux, utilisés au cours des chapitres de ce Traité, qui sont plus ou moins bien connus, mais pour lesquels il n'a pas été possible de donner des références commodes.
Il est recommandé au lecteur de ne se reporter au chapitre~0 qu'en cours de lecture du Traité proprement dit, et dans la mesure où les résultats auxquels nous référons ne lui sont pas

\oldpage[I]{7}
suffisamment familiers.
Nous pensons d'ailleurs que de cette façon, la lecture de ce Traité pourra être pour le débutant une bonne méthode lui permettant de se familiariser avec l'Algèbre commutative et l'Algèbre homologique, dont l'étude, lorsqu'elle ne s'accompagne pas d'applications tangibles, est jugée fastidieuse, voire déprimante, par un assez grand nombre.

\bigskip
\begin{center}*\quad*\quad*\end{center}
\bigskip

Il est hors de notre compétence de donner dans cette Introduction un aperçu historique, même sommaire, des notions et résultats exposés.
Le texte ne contiendra que des références jugées particulièrement utiles pour sa compréhension, et nous n'indiquerons l'origine que des résultats les plus importants.
Formellement du moins, les sujets traités dans notre ouvrage sont assez neufs, ce qui expliquera la rareté des références faites aux Pères de la Géométrie algébrique du \textsc{xix}\textsuperscript{e} siècle et du début du \textsc{xx}\textsuperscript{e} siècle, dont nous ne connaissons les travaux que par ouï-dire.
Il convient cependant de dire quelques mots ici sur les ouvrages qui ont le plus directement influencé les auteurs et contribué au développement du point de vue des schémas.
Il faut en tout premier lieu citer le travail fondamental (FAC) de J.-P.~Serre, qui a servi d'introduction à la Géométrie algébrique pour plus d'un jeune adepte (dont l'un des auteurs du présent Traité), rebuté par l'aridité des classiques \emph{Foundations} de A.~Weil~\cite{I-18}.
C'est là qu'il est démontré pour la première fois que la <<~topologie de Zariski~>> d'une variété algébrique <<~abstraite~>> est parfaitement appropriée pour lui appliquer certaines techniques de la Topologie algébrique et donner lieu notamment à une théorie cohomologique.
De plus, la définition d'une variété algébrique qui y est donnée est celle qui se prête le plus naturellement à l'extension de cette notion que nous développons ici\footnote{Ainsi que J.-P.~Serre nous l'a signalé, il convient de noter que l'idée de définir la structure de variété par la donnée d'un faisceau d'anneaux est due à H.~Cartan, qui a pris cette idée comme point de départ de sa théorie des espaces analytiques.
Bien entendu, tout comme en Géométrie algébrique, il importerait, en <<~Géométrie analytique~>>, de donner droit de cité aux éléments nilpotents dans les anneaux locaux des espaces analytiques.
Cette extension de la définition de H.~Cartan et J.-P.~Serre a été récemment abordée par H.~Grauert~\cite{I-5}, et il y a lieu d'espérer qu'un exposé systématique de Géométrie analytique dans ce cadre général verra bientôt le jour.
Il est d'ailleurs évident que les notions et techniques développées dans ce Traité gardent un sens en Géométrie analytique, bien qu'il faille s'attendre à des difficultés techniques plus considérables dans cette dernière théorie.
On peut prévoir que la Géométrie algébrique, par la simplicité de ses méthodes, pourra servir comme une sorte de modèle formel pour de futurs développements dans la théorie des espaces analytiques.}.
Serre avait d'ailleurs remarqué lui-même que la théorie cohomologique des variétés algébriques affines pouvait se transcrire sans difficulté en remplaçant les algèbres affines sur un corps par des anneaux commutatifs quelconques.
Les chapitres I et II de ce Traité et les deux premiers paragraphes du chapitre III peuvent donc être considérés, pour l'essentiel, comme des transpositions faciles, dans ce cadre élargi, des résultats principaux de (FAC) et d'un article ultérieur du même auteur~\cite{I-15}.
Nous avons aussi retiré grand profit du \emph{Séminaire de Géométrie algébrique} de C.~Chevalley~\cite{I-1}; en particulier, l'usage systématique des <<~ensembles constructibles~>> introduits par lui, s'est révélé fort utile en théorie des schémas (cf. chap.~IV).
Nous lui avons aussi emprunté l'étude des morphismes du point

\oldpage[I]{8}
de vue de la dimension (chap.~IV), qui se transcrit sans changement notable dans le cadre des schémas.
Il convient de noter par ailleurs que la notion de <<~schémas d'anneaux locaux~>>, introduite par Chevalley, se prête naturellement à une extension de la Géométrie algébrique (n'ayant pas cependant toute la souplesse et la généralité que nous entendons lui donner ici); pour les rapports entre cette notion et notre théorie, voir chapitre premier, \S~8.
Une telle extension a été développée par M.~Nagata dans une série de mémoires~\cite{I-9} contenant de nombreux résultats spéciaux concernant la Géométrie algébrique sur les anneaux de Dedekind\footnote{Parmi les travaux qui se rapprochent de notre point de vue en Géométrie algébrique, signalons l'important travail de E.~Kähler~\cite{I-22}, et une Note récente de Chow et Igusa~\cite{I-3}, qui reprennent dans le cadre de la théorie de Nagata-Chevalley certains résultats de (FAC) et donnent aussi une formule de Künneth.}.

\bigskip
\begin{center}*\quad*\quad*\end{center}
\bigskip

Enfin, il va sans dire qu'un livre sur la Géométrie algébrique, et surtout un livre portant sur les fondements, est nécessairement influencé, ne serait-ce que par personnes interposées, par des mathématiciens tels que O.~Zariski et A.~Weil.
En particulier, la \emph{Théorie des fonctions holomorphes} de Zariski~\cite{I-20}, convenablement assouplie grâce aux méthodes cohomologiques et complétée par un théorème d'existence (chap.~III, \S\S~4 et 5) est (avec la technique de descente exposée au chap.~VI) un des principaux outils employés dans ce Traité, et nous semble un des plus puissants dont on dispose en Géométrie algébrique.

La technique générale dans laquelle elle s'insère peut être esquissée de la façon suivante (un exemple typique en sera fourni au chap.~IX, dans l'étude du groupe fondamental).
On a un morphisme propre (chap.~II) $f:X\to Y$ d'une variété algébrique dans une autre (plus généralement, d'un schéma dans un autre) qu'on veut étudier au voisinage d'un point $y\in Y$, en vue de résoudre un problème P relatif à un voisinage de $y$.
On opère par étapes successives~:
\begin{enumerate}
  \item[$1^\circ$] On peut supposer $Y$ affine, de sorte que $X$ devient un schéma défini sur l'anneau affine $A$ de $Y$, et on peut même remplacer $A$ par l'anneau local de $y$.
  Cette réduction est toujours facile en pratique (chap.~V) et nous ramène au cas où $A$ est un anneau \emph{local}.
  \item[$2^\circ$] On étudie le problème envisagé lorsque $A$ est un anneau local \emph{artinien}.
  Pour qu'il garde effectivement un sens lorsque $A$ n'est pas supposé intègre, il y a lieu parfois de reformuler le problème P, et il apparaît que l'on obtient souvent ainsi une meilleure compréhension du problème, de nature <<~infinitésimale~>> à ce stade.
  \item[$3^\circ$] La théorie des schémas formels (chap.~III, \S\S~3, 4 et 5) permet de passer du cas d'un anneau artinien au cas d'un \emph{anneau local complet}.
  \item[$4^\circ$] Enfin, si $A$ est un anneau local quelconque, la considération de <<~sections multiformes~>> sur des schémas convenables sur $X$, approchant une section <<~formelle~>> donnée (chap.~IV) permettra souvent de passer d'un résultat connu pour le schéma

  \oldpage[I]{9}
  déduit de $X$ par extension des scalaires au complété de $A$, à un résultat analogue pour une extension finie assez simple (par exemple non ramifiée) de $A$.
\end{enumerate}

Cette esquisse montre l'importance de l'étude systématique des schémas définis sur un anneau artinien $A$.
Le point de vue de Serre dans sa formulation de la théorie du corps de classes local, et des travaux récents de Greenberg, semblent suggérer qu'une telle étude pourrait être entreprise en attachant fonctoriellement à un tel schéma $X$ un schéma $X'$ sur le corps résiduel $k$ de $A$ (supposé parfait) de dimension égale (dans les cas favorables) à $n\dim X$, où $n$ est la longueur de $A$.

Quant à l'influence de A.~Weil, qu'il nous suffise de dire que c'est la nécessité de développer l'outillage nécessaire pour formuler avec toute la généralité voulue la définition de la <<~cohomologie de Weil~>> et pour aborder la démonstration\footnote{Pour éviter tout malentendu, précisons que cette tâche vient à peine d'être entreprise au moment où cette Introduction est écrite, et n'a donc pas encore abouti à la démonstration des conjectures de Weil.} de toutes les propriétés formelles nécessaires pour établir ses célèbres conjectures en Géométrie diophantienne~\cite{I-19}, qui a été une des principales motivations de la rédaction du présent Traité, au même titre que le désir de trouver le cadre naturel des notions et méthodes usuelles en Géométrie algébrique, et de donner aux auteurs l'occasion de comprendre lesdites notions et techniques.

\bigskip
\begin{center}*\quad*\quad*\end{center}
\bigskip

Pour terminer, nous croyons utile de prévenir les lecteurs que, tout comme les auteurs eux-mêmes, ils auront sans doute quelque difficulté avant de s'accoutumer au langage des schémas, et de se convaincre que les constructions habituelles que suggère l'intuition géométrique peuvent se transcrire, essentiellement d'une seule façon raisonnable, dans ce langage.
Comme dans beaucoup de parties de la Mathématique moderne, l'intuition première s'éloigne de plus en plus, en apparence, du langage propre à l'exprimer avec toute la précision et la généralité voulues.
En l'occurrence, la difficulté psychologique tient à la nécessité de transporter aux objets d'une catégorie déjà assez différente de la catégorie des ensembles (à savoir la catégorie des préschémas, ou la catégorie des préschémas sur un préschéma donné) des notions familières pour les ensembles~: produits cartésiens, lois de groupe, d'anneau, de module, fibrés, fibrés principaux homogènes, etc.
Il sera sans doute difficile au mathématicien, dans l'avenir, de se dérober à ce nouvel effort d'abstraction, peut-être assez minime, somme toute, en comparaison de celui fourni par nos pères, se familiarisant avec la Théorie des Ensembles.

\bigskip
\begin{center}*\quad*\quad*\end{center}
\bigskip

Les références seront données suivant le système décimal; par exemple, dans III, 4.9.3, le chiffre III indique le chapitre, le chiffre 4 le paragraphe, le chiffre 9 la section du paragraphe.
A l'intérieur du même chapitre, on supprimera la mention du chapitre.

% FIN SOURCE ega1/intro-fr.tex

%
% DEBUT SOURCE ega1/ega0-1-fr.tex
\section{Anneaux de fractions}
\label{section:0.1-fr}

\setcounter{subsection}{-1}
\subsection{Anneaux et algèbres}
\label{subsection:0.1.0-fr}

\oldpage[0\textsubscript{I}]{11}
\begin{env}[1.0.1]
\label{0.1.0.1-fr}
Tous les anneaux considérés dans ce Traité posséderont un \emph{élément unité};
tous les modules sur un tel anneau seront supposés \emph{unitaires}; les homomorphismes d'anneaux seront toujours supposés \emph{transformer l'élément unité en élément unité}; sauf mention expresse du contraire, un sous-anneau d'un anneau $A$ sera supposé \emph{contenir l'élément unité de $A$}.
Nous considérerons surtout des anneaux \emph{commutatifs}, et lorsque nous parlerons d'anneau sans préciser, il sera sous-entendu qu'il s'agit d'un anneau commutatif.
Si $A$ est un anneau non nécessairement commutatif, par $A$-module nous entendrons toujours un module \emph{à gauche}, sauf mention expresse du contraire.
\end{env}

\begin{env}[1.0.2]
\label{0.1.0.2-fr}
Soient $A$, $B$ deux anneaux non nécessairement commutatifs, $\varphi:A\to B$ un homomorphisme.
Tout $B$-module à gauche (resp. à droite) $M$ peut être muni d'une structure de $A$-module à gauche (resp. à droite) en posant $a.m=\varphi(a).m$ (resp. $m.a=m.\varphi(a)$); lorsqu'il sera nécessaire de distinguer sur $M$ les structures de $A$-module et de $B$-module, nous désignerons par $M_{[\varphi]}$ le $A$-module à gauche (resp. à droite) ainsi défini.
Si $L$ est un $A$-module, un homomorphisme $u:L\to M_{[\varphi]}$ est donc un homomorphisme de groupes commutatifs tel que $u(a.x)=\varphi(a).u(x)$ pour $a\in A$, $x\in L$; on dira aussi que c'est un \emph{$\varphi$-homomorphisme} $L\to M$, et que le couple $(\varphi,u)$ (ou, par abus de langage, $u$) est un \emph{di-homomorphisme} de $(A,L)$ dans $(B,M)$.
Les couples $(A,L)$ formés d'un anneau $A$ et d'un $A$-module $L$ forment donc une \emph{catégorie} pour laquelle les morphismes sont les di-homomorphismes.
\end{env}

\begin{env}[1.0.3]
\label{0.1.0.3-fr}
Sous les hypothèses de (1.0.2), si $\mathfrak{J}$ est un idéal à gauche (resp. à droite) de $A$, nous noterons $B\mathfrak{J}$ (resp. $\mathfrak{J}B$) l'idéal à gauche (resp. à droite) $B\varphi(\mathfrak{J})$ (resp. $\varphi(\mathfrak{J})B$) de $B$ engendré par $\varphi(\mathfrak{J})$; c'est aussi l'image de l'homomorphisme canonique $B\otimes_A\mathfrak{J}\to B$ (resp. $\mathfrak{J}\otimes_A B\to B$) de $B$-modules à gauche (resp. à droite).
\end{env}

\begin{env}[1.0.4]
\label{0.1.0.4-fr}
Si $A$ est un anneau (commutatif), $B$ un anneau non nécessairement commutatif, la donnée d'une structure de \emph{$A$-algèbre} sur $B$ équivaut à la donnée d'un homomorphisme d'anneaux $\varphi:A\to B$ tel que $\varphi(A)$ soit contenu dans le centre de $B$.
Pour tout idéal $\mathfrak{J}$ de $A$, $\mathfrak{J}B=B\mathfrak{J}$ est alors un idéal bilatère de $B$, et pour tout $B$-module $M$, $\mathfrak{J}M$ est alors un $B$-module égal à $(B\mathfrak{J})M$.
\end{env}

\begin{env}[1.0.5]
\label{0.1.0.5-fr}
Nous ne reviendrons pas sur les notions de \emph{module de type fini} et d'\emph{algèbre} (commutative) \emph{de type fini}; dire qu'un $A$-module $M$ est de type fini signifie qu'il existe

\oldpage[0\textsubscript{I}]{12}
une suite exacte $A^p\to M\to 0$.
On dit qu'un $A$-module $M$ admet une \emph{présentation finie} s'il est isomorphe au conoyau d'un homomorphisme $A^p\to A^q$, autrement dit s'il existe une suite exacte $A^p\to A^q\to M\to 0$.
On notera que sur un anneau \emph{noethérien} $A$, tout $A$-module de type fini admet une présentation finie.

Rappelons qu'une $A$-algèbre $B$ est dite \emph{entière sur $A$} si tout élément de $B$ est racine dans $B$ d'un polynôme unitaire à coefficients dans $A$; il revient au même de dire que tout élément de $B$ est contenu dans une sous-algèbre de $B$ qui est un \emph{$A$-module de type fini}.
Lorsqu'il en est ainsi, et que $B$ est commutative, la sous-algèbre de $B$ engendrée par une partie finie de $B$ est un $A$-module de type fini; pour que l'algèbre commutative $B$ soit entière et de type fini sur $A$, il faut et il suffit donc que $B$ soit un $A$-module de type fini; on dit alors aussi que $B$ est une $A$-algèbre \emph{entière finie} (ou simplement \emph{finie} si aucune confusion n'en résulte).
On observera que dans ces définitions, on ne suppose pas que l'homomorphisme $A\to B$ définissant la structure de $A$-algèbre soit injectif.
\end{env}

\begin{env}[1.0.6]
\label{0.1.0.6-fr}
Un anneau \emph{intègre} est un anneau dans lequel le produit d'une famille finie d'éléments $\neq 0$ est $\neq 0$; il revient au même de dire que dans un tel anneau on a $0\neq 1$ et le produit de deux éléments $\neq 0$ est non nul.
Un idéal \emph{premier} d'un anneau $A$ est un idéal $\mathfrak{p}$ tel que $A/\mathfrak{p}$ soit intègre; cela entraîne donc $\mathfrak{p}\neq A$.
Pour qu'un anneau $A$ ait au moins un idéal premier, il faut et il suffit que $A\neq\{0\}$.
\end{env}

\begin{env}[1.0.7]
\label{0.1.0.7-fr}
Un anneau \emph{local} est un anneau $A$ dans lequel il existe un seul idéal maximal, qui est alors le complémentaire des éléments inversibles et contient tous les idéaux $\neq A$.
Si $A$ et $B$ sont deux anneaux locaux, $\mathfrak{m}$ et $\mathfrak{n}$ leurs idéaux maximaux respectifs, on dit qu'un homomorphisme $\varphi:A\to B$ est \emph{local} si $\varphi(\mathfrak{m})\subset\mathfrak{n}$ (ou, ce qui revient au même, si $\varphi^{-1}(\mathfrak{n})=\mathfrak{m}$).
Par passage aux quotients, un tel homomorphisme définit alors un monomorphisme du corps résiduel $A/\mathfrak{m}$ dans le corps résiduel $B/\mathfrak{n}$.
Le composé de deux homomorphismes locaux est un homomorphisme local.
\end{env}

\subsection{Racine d'un idéal. Nilradical et radical d'un anneau}
\label{subsection:0.1.1-fr}

\begin{env}[1.1.1]
\label{0.1.1.1-fr}
Soit $\mathfrak{a}$ un idéal d'un anneau $A$; la \emph{racine} de $\mathfrak{a}$, notée $\mathfrak{r}(\mathfrak{a})$, est l'ensemble des $x\in A$ tels que $x^n\in\mathfrak{a}$ pour un entier $n>0$ au moins; c'est un idéal contenant $\mathfrak{a}$.
On a $\mathfrak{r}(\mathfrak{r}(\mathfrak{a}))=\mathfrak{r}(\mathfrak{a})$; la relation $\mathfrak{a}\subset\mathfrak{b}$ entraîne $\mathfrak{r}(\mathfrak{a})\subset\mathfrak{r}(\mathfrak{b})$; la racine d'une intersection finie d'idéaux est l'intersection de leurs racines.
Si $\varphi$ est un homomorphisme d'un anneau $A'$ dans $A$, on a $\mathfrak{r}(\varphi^{-1}(\mathfrak{a}))=\varphi^{-1}(\mathfrak{r}(\mathfrak{a}))$ pour tout idéal $\mathfrak{a}\subset A$.
Pour qu'un idéal soit racine d'un idéal, il faut et il suffit qu'il soit intersection d'idéaux premiers.
La racine d'un idéal $\mathfrak{a}$ est l'intersection des idéaux premiers \emph{minimaux} parmi ceux qui contiennent $\mathfrak{a}$; si $A$ est noethérien, ces idéaux premiers minimaux sont en nombre fini.

La racine de l'idéal $(0)$ est encore appelée le \emph{nilradical} de $A$; c'est l'ensemble $\mathfrak{N}$ des éléments nilpotents de $A$.
On dit que l'anneau $A$ est \emph{réduit} si $\mathfrak{N}=(0)$; pour tout anneau $A$, le quotient $A/\mathfrak{N}$ de $A$ par son nilradical est un anneau réduit.
\end{env}

\begin{env}[1.1.2]
\label{0.1.1.2-fr}
Rappelons que le \emph{radical} $\mathfrak{R}(A)$ d'un anneau $A$ (non nécessairement commutatif) est l'intersection des idéaux à gauche maximaux de $A$ (et aussi l'intersection des idéaux à droite maximaux).
Le radical de $A/\mathfrak{R}(A)$ est $(0)$.
\end{env}

\subsection{Modules et anneaux de fractions}
\label{subsection:0.1.2-fr}

\oldpage[0\textsubscript{I}]{13}
\begin{env}[1.2.1]
\label{0.1.2.1-fr}
On dit qu'une partie $S$ d'un anneau $A$ est \emph{multiplicative} si $1\in S$ et si le produit de deux éléments de $S$ est dans $S$.
Les exemples qui seront les plus importants pour la suite sont~: $1^\circ$ l'ensemble $S_f$ des puissances $f^n$ ($n\geq 0$) d'un élément $f\in A$; $2^\circ$ le complémentaire $A-\mathfrak{p}$ d'un idéal \emph{premier} $\mathfrak{p}$ de $A$.
\end{env}

\begin{env}[1.2.2]
\label{0.1.2.2-fr}
Soient $S$ une partie multiplicative d'un anneau $A$, $M$ un $A$-module; dans l'ensemble $M\times S$, la relation entre couples $(m_1,s_1)$, $(m_2,s_2)$~:
\[
  \text{<<~il existe $s\in S$ tel que $s(s_1m_2-s_2m_1)=0$~>>}
\]
est une relation d'équivalence.
On désigne par $S^{-1}M$ l'ensemble quotient de $M\times S$ par cette relation, par $m/s$ l'image canonique dans $S^{-1}M$ du couple $(m,s)$; on appelle \emph{application canonique} de $M$ dans $S^{-1}M$ l'application $i_M^S:m\mapsto m/1$ (aussi notée $i^S$).
Cette application n'est en général ni injective ni surjective; son noyau est l'ensemble des $m\in M$ tels qu'il existe un $s\in S$ pour lequel $sm=0$.

Dans $S^{-1}M$ on définit une loi de groupe additif en prenant
\[
  (m_1/s_1)+(m_2/s_2)=(s_2m_1+s_1m_2)/(s_1s_2)
\]
(on vérifie que c'est bien indépendant des expressions des éléments de $S^{-1}M$ considérés).
Sur $S^{-1}A$ on définit en outre une loi multiplicative en prenant $(a_1/s_1)(a_2/s_2)=(a_1a_2)/(s_1s_2)$, et enfin une loi externe sur $S^{-1}M$, ayant $S^{-1}A$ comme ensemble d'opérateurs, en posant $(a/s)(m/s')=(am)/(ss')$.
On vérifie ainsi que $S^{-1}A$ est muni d'une structure d'anneau (dit \emph{anneau de fractions de $A$ à dénominateurs dans $S$}) et $S^{-1}M$ d'une structure de $S^{-1}A$-module (dit \emph{module des fractions de $M$ à dénominateurs dans $S$}); pour tout $s\in S$, $s/1$ est inversible dans $S^{-1}A$, son inverse étant $1/s$.
L'application canonique $i_A^S$ (resp. $i_M^S$) est un homomorphisme d'anneaux (resp. un homomorphisme de $A$-modules, $S^{-1}M$ étant considéré comme $A$-module au moyen de l'homomorphisme $i_A^S:A\to S^{-1}A$).
\end{env}

\begin{env}[1.2.3]
\label{0.1.2.3-fr}
Si $S_f=\{f^n\}_{n\geq 0}$ pour un $f\in A$, on écrit $A_f$ et $M_f$ au lieu de $S_f^{-1}A$ et $S_f^{-1}M$; quand $A_f$ est considéré comme algèbre sur $A$, on peut écrire $A_f=A[1/f]$.
$A_f$ est isomorphe à l'algèbre quotient $A[T]/(fT-1)A[T]$.
Lorsque $f=1$, $A_f$ et $M_f$ s'identifient canoniquement à $A$ et $M$; si $f$ est nilpotent, $A_f$ et $M_f$ sont réduits à $0$.

Lorsque $S=A-\mathfrak{p}$, où $\mathfrak{p}$ est un idéal premier de $A$, on écrit $A_{\mathfrak{p}}$ et $M_{\mathfrak{p}}$ au lieu de $S^{-1}A$ et $S^{-1}M$; $A_{\mathfrak{p}}$ est un \emph{anneau local} dont l'idéal maximal $\mathfrak{q}$ est engendré par $i_A^S(\mathfrak{p})$, et on a $(i_A^S)^{-1}(\mathfrak{q})=\mathfrak{p}$; par passage aux quotients, $i_A^S$ donne un monomorphisme de l'anneau intègre $A/\mathfrak{p}$ dans le corps $A_{\mathfrak{p}}/\mathfrak{q}$, qui s'identifie au corps des fractions de $A/\mathfrak{p}$.
\end{env}

\begin{env}[1.2.4]
\label{0.1.2.4-fr}
L'anneau de fractions $S^{-1}A$ et l'homomorphisme canonique $i_A^S$ sont solution d'un \emph{problème d'application universelle}~: tout homomorphisme $u$ de $A$ dans un anneau $B$ tel que $u(S)$ se compose d'éléments \emph{inversibles} dans $B$ se factorise d'une seule manière
\[
  u:A\xrightarrow{i_A^S}S^{-1}A\xrightarrow{u^*}B
\]

\oldpage[0\textsubscript{I}]{14}
où $u^*$ est un homomorphisme d'anneaux.
Sous les mêmes hypothèses, soient $M$ un $A$-module, $N$ un $B$-module, $v:M\to N$ un homomorphisme de $A$-modules (pour la structure de $B$-module sur $N$ définie par $u:A\to B$); alors $v$ se factorise d'une seule manière
\[
  v:M\xrightarrow{i_M^S}S^{-1}M\xrightarrow{v^*}N
\]
où $v^*$ est un homomorphisme de $S^{-1}A$-modules (pour la structure de $S^{-1}A$-module sur $N$ définie par $u^*$).
\end{env}

\begin{env}[1.2.5]
\label{0.1.2.5-fr}
On définit un isomorphisme canonique $S^{-1}A\otimes_A M\xrightarrow{\sim}S^{-1}M$ de $S^{-1}A$-modules, en faisant correspondre à l'élément $(a/s)\otimes m$ l'élément $(am)/s$, l'isomorphisme réciproque appliquant $m/s$ sur $(1/s)\otimes m$.
\end{env}

\begin{env}[1.2.6]
\label{0.1.2.6-fr}
Pour tout idéal $\mathfrak{a}'$ de $S^{-1}A$, $\mathfrak{a}=(i_A^S)^{-1}(\mathfrak{a}')$ est un idéal de $A$, et $\mathfrak{a}'$ est l'idéal de $S^{-1}A$ engendré par $i_A^S(\mathfrak{a})$, qui s'identifie à $S^{-1}\mathfrak{a}$ (1.3.2).
L'application $\mathfrak{p}'\mapsto(i_A^S)^{-1}(\mathfrak{p}')$ est un isomorphisme, pour la structure d'ordre, de l'ensemble des idéaux \emph{premiers} de $S^{-1}A$ sur l'ensemble des idéaux premiers $\mathfrak{p}$ de $A$ tels que $\mathfrak{p}\cap S=\varnothing$.
En outre, les anneaux locaux $A_{\mathfrak{p}}$ et $(S^{-1}A)_{S^{-1}\mathfrak{p}}$ sont alors canoniquement (1.5.1) isomorphes.
\end{env}

\begin{env}[1.2.7]
\label{0.1.2.7-fr}
Lorsque $A$ est un anneau \emph{intègre}, dont on désigne par $K$ le corps des fractions, l'application canonique $i_A^S:A\to S^{-1}A$ est injective pour toute partie multiplicative $S$ ne contenant pas $0$, et $S^{-1}A$ s'identifie alors canoniquement à un sous-anneau de $K$ contenant $A$.
En particulier, pour tout idéal premier $\mathfrak{p}$ de $A$, $A_{\mathfrak{p}}$ est un anneau local contenant $A$, d'idéal maximal $\mathfrak{p}A_{\mathfrak{p}}$, et on a $\mathfrak{p}A_{\mathfrak{p}}\cap A=\mathfrak{p}$.
\end{env}

\begin{env}[1.2.8]
\label{0.1.2.8-fr}
Si $A$ est un anneau \emph{réduit} (1.1.1), il en est de même de $S^{-1}A$~: en effet, si $(x/s)^n=0$ pour $x\in A$, $s\in S$, cela signifie qu'il existe $s'\in S$ tel que $s'x^n=0$, d'où $(s'x)^n=0$, ce qui, par hypothèse, entraîne $s'x=0$, donc $x/s=0$.
\end{env}

\subsection{Propriétés fonctorielles}
\label{subsection:0.1.3-fr}

\begin{env}[1.3.1]
\label{0.1.3.1-fr}
Soient $M$, $N$ deux $A$-modules, $u$ un $A$-homomorphisme $M\to N$.
Si $S$ est une partie multiplicative de $A$, on définit un $S^{-1}A$-homomorphisme $S^{-1}M\to S^{-1}N$, noté $S^{-1}u$, en posant $(S^{-1}u)(m/s)=u(m)/s$; si $S^{-1}M$ et $S^{-1}N$ sont canoniquement identifiés à $S^{-1}A\otimes_A M$ et $S^{-1}A\otimes_A N$ (1.2.5), $S^{-1}u$ est identifié à $1\otimes u$.
Si $P$ est un troisième $A$-module, $v$ un $A$-homomorphisme $N\to P$, on a $S^{-1}(v\circ u)=(S^{-1}v)\circ(S^{-1}u)$; autrement dit $S^{-1}M$ est un \emph{foncteur covariant en $M$}, de la catégorie des $A$-modules dans celle des $S^{-1}A$-modules ($A$ et $S$ étant fixés).
\end{env}

\begin{env}[1.3.2]
\label{0.1.3.2-fr}
Le foncteur $S^{-1}M$ est \emph{exact}; autrement dit, si la suite
\[
  M\xrightarrow{u}N\xrightarrow{v}P
\]
est exacte, il en est de même de la suite
\[
  S^{-1}M\xrightarrow{S^{-1}u}S^{-1}N\xrightarrow{S^{-1}v}S^{-1}P.
\]
En particulier, si $u:M\to N$ est injectif (resp. surjectif), il en est de même de $S^{-1}u$;

\oldpage[0\textsubscript{I}]{15}
si $N$ et $P$ sont deux sous-modules de $M$, $S^{-1}N$ et $S^{-1}P$ s'identifient canoniquement à des sous-modules de $S^{-1}M$, et l'on a
\[
  S^{-1}(N+P)=S^{-1}N+S^{-1}P
  \quad\text{et}\quad
  S^{-1}(N\cap P)=(S^{-1}N)\cap(S^{-1}P).
\]
\end{env}

\begin{env}[1.3.3]
\label{0.1.3.3-fr}
Soit $(M_\alpha,\varphi_{\beta\alpha})$ un système inductif de $A$-modules; alors $(S^{-1}M_\alpha,S^{-1}\varphi_{\beta\alpha})$ est un système inductif de $S^{-1}A$-modules.
Exprimant les $S^{-1}M_\alpha$ et $S^{-1}\varphi_{\beta\alpha}$ comme des produits tensoriels (1.2.5 et 1.3.1), il résulte de la permutabilité des opérations de produit tensoriel et de limite inductive, que l'on a un isomorphisme canonique
\[
  S^{-1}\varinjlim M_\alpha\xrightarrow{\sim}\varinjlim S^{-1}M_\alpha,
\]
ce qu'on exprime encore en disant que le foncteur $S^{-1}M$ (en $M$) \emph{commute avec les limites inductives}.
\end{env}

\begin{env}[1.3.4]
\label{0.1.3.4-fr}
Soient $M$, $N$ deux $A$-modules; il existe un isomorphisme canonique \emph{fonctoriel} (en $M$ et $N$)
\[
  (S^{-1}M)\otimes_{S^{-1}A}(S^{-1}N)
  \xrightarrow{\sim}
  S^{-1}(M\otimes_A N)
\]
qui transforme $(m/s)\otimes(n/t)$ en $(m\otimes n)/st$.
\end{env}

\begin{env}[1.3.5]
\label{0.1.3.5-fr}
On a de même un homomorphisme \emph{fonctoriel} (en $M$ et $N$)
\[
  S^{-1}\operatorname{Hom}_A(M,N)
  \longrightarrow
  \operatorname{Hom}_{S^{-1}A}(S^{-1}M,S^{-1}N)
\]
qui, à $u/s$, fait correspondre l'homomorphisme $m/t\mapsto u(m)/st$.
Lorsque $M$ a une \emph{présentation finie}, l'homomorphisme précédent est un \emph{isomorphisme}~: c'est immédiat lorsque $M$ est de la forme $A^r$, et on passe de là au cas général en partant de la suite exacte $A^p\to A^q\to M\to 0$, et en utilisant l'exactitude du foncteur $S^{-1}M$ et l'exactitude à gauche du foncteur $\operatorname{Hom}_A(M,N)$ en $M$.
On notera que ce cas se présente toujours lorsque $A$ est \emph{noethérien} et le $A$-module $M$ \emph{de type fini}.
\end{env}

\subsection{Changement de partie multiplicative}
\label{subsection:0.1.4-fr}

\begin{env}[1.4.1]
\label{0.1.4.1-fr}
Soient $S$, $T$ deux parties multiplicatives d'un anneau $A$ telles que $S\subset T$; il existe un homomorphisme canonique $\rho_A^{T,S}$ (ou simplement $\rho^{T,S}$) de $S^{-1}A$ dans $T^{-1}A$, faisant correspondre à l'élément noté $a/s$ de $S^{-1}A$ l'élément noté $a/s$ dans $T^{-1}A$; on a
\[
  i_A^T=\rho_A^{T,S}\circ i_A^S.
\]
Pour tout $A$-module $M$, il existe de même une application $S^{-1}A$-linéaire de $S^{-1}M$ dans $T^{-1}M$ (ce dernier étant considéré comme $S^{-1}A$-module grâce à l'homomorphisme $\rho_A^{T,S}$), qui fait correspondre à l'élément $m/s$ de $S^{-1}M$ l'élément $m/s$ de $T^{-1}M$; on note cette application $\rho_M^{T,S}$, ou simplement $\rho^{T,S}$, et on a encore
\[
  i_M^T=\rho_M^{T,S}\circ i_M^S;
\]
dans l'identification canonique (1.2.5), $\rho_M^{T,S}$ s'identifie à $\rho_A^{T,S}\otimes 1$.
L'homomorphisme $\rho_M^{T,S}$ est un \emph{morphisme fonctoriel} (ou transformation naturelle) du foncteur $S^{-1}M$ dans le foncteur $T^{-1}M$, autrement dit, le diagramme
\[
\begin{tikzcd}[column sep=large,row sep=large]
S^{-1}M \arrow[r,"S^{-1}u"] \arrow[d,"\rho_M^{T,S}"'] & S^{-1}N \arrow[d,"\rho_N^{T,S}"] \\
T^{-1}M \arrow[r,"T^{-1}u"] & T^{-1}N
\end{tikzcd}
\]

\oldpage[0\textsubscript{I}]{16}
est commutatif, pour tout homomorphisme $u:M\to N$; on notera en outre que $T^{-1}u$ est entièrement déterminé par $S^{-1}u$, car pour $m\in M$ et $t\in T$, on a
\[
  (T^{-1}u)(m/t)=(t/1)^{-1}\rho^{T,S}\bigl((S^{-1}u)(m/1)\bigr).
\]
\end{env}

\begin{env}[1.4.2]
\label{0.1.4.2-fr}
Avec les mêmes notations, pour deux $A$-modules $M$, $N$, les diagrammes (cf. (1.3.4) et (1.3.5))
\[
\begin{tikzcd}[column sep=large,row sep=large]
(S^{-1}M)\otimes_{S^{-1}A}(S^{-1}N) \arrow[r,"\sim"] \arrow[d]
  & S^{-1}(M\otimes_A N) \arrow[d]
  \\
(T^{-1}M)\otimes_{T^{-1}A}(T^{-1}N) \arrow[r,"\sim"]
  & T^{-1}(M\otimes_A N)
\end{tikzcd}
\]
\[
\begin{tikzcd}[column sep=large,row sep=large]
S^{-1}\operatorname{Hom}_A(M,N) \arrow[r] \arrow[d]
  & \operatorname{Hom}_{S^{-1}A}(S^{-1}M,S^{-1}N) \arrow[d] \\
T^{-1}\operatorname{Hom}_A(M,N) \arrow[r]
  & \operatorname{Hom}_{T^{-1}A}(T^{-1}M,T^{-1}N)
\end{tikzcd}
\]
sont commutatifs.
\end{env}

\begin{env}[1.4.3]
\label{0.1.4.3-fr}
Il y a un cas important dans lequel l'homomorphisme $\rho^{T,S}$ est \emph{bijectif}, savoir lorsque tout élément de $T$ est diviseur d'un élément de $S$; on identifie alors par $\rho^{T,S}$ les modules $S^{-1}M$ et $T^{-1}M$.
On dit que $S$ est \emph{saturé} si tout diviseur dans $A$ d'un élément de $S$ est dans $S$; en remplaçant $S$ par l'ensemble $T$ de tous les diviseurs des éléments de $S$ (ensemble qui est multiplicatif et saturé), on voit qu'on peut toujours, si l'on veut, se limiter à la considération de modules de fractions $S^{-1}M$, où $S$ est saturé.
\end{env}

\begin{env}[1.4.4]
\label{0.1.4.4-fr}
Si $S$, $T$, $U$ sont trois parties multiplicatives de $A$ telles que $S\subset T\subset U$, on a
\[
  \rho^{U,S}=\rho^{U,T}\circ\rho^{T,S}.
\]
\end{env}

\begin{env}[1.4.5]
\label{0.1.4.5-fr}
Considérons une \emph{famille filtrante croissante} $(S_\alpha)$ de parties multiplicatives de $A$ (on écrira $\alpha\leq\beta$ pour $S_\alpha\subset S_\beta$), et soit $S$ la partie multiplicative $\bigcup_\alpha S_\alpha$; posons
\[
  \rho_{\beta\alpha}=\rho_A^{S_\beta,S_\alpha}\quad\text{pour }\alpha\leq\beta;
\]
en vertu de (1.4.4), les homomorphismes $\rho_{\beta\alpha}$ définissent un anneau $A'$ \emph{limite inductive} du système inductif d'anneaux $(S_\alpha^{-1}A,\rho_{\beta\alpha})$.
Soit $\rho_\alpha$ l'application canonique $S_\alpha^{-1}A\to A'$, et posons $\varphi_\alpha=\rho_A^{S,S_\alpha}$; comme $\varphi_\alpha=\varphi_\beta\circ\rho_{\beta\alpha}$ pour $\alpha\leq\beta$ d'après (1.4.4), on peut définir de façon unique un homomorphisme $\varphi:A'\to S^{-1}A$ tel que le diagramme
\[
\begin{tikzcd}[column sep=huge,row sep=large]
  & S_\alpha^{-1}A
      \arrow[ddl,"\rho_\alpha"']
      \arrow[d,"\rho_{\beta\alpha}"]
      \arrow[ddr,"\varphi_\alpha"] & \\
  & S_\beta^{-1}A
      \arrow[dl,"\rho_\beta"]
      \arrow[dr,"\varphi_\beta"'] & \\
A' \arrow[rr,"\varphi"'] & & S^{-1}A
\end{tikzcd}
\qquad(\alpha\leq\beta)
\]
soit commutatif.
En fait, $\varphi$ est un \emph{isomorphisme}~: il est en effet immédiat par construction que $\varphi$ est surjectif.
D'autre part, si $\rho_\alpha(a/s_\alpha)\in A'$ est tel que $\varphi(\rho_\alpha(a/s_\alpha))=0$, cela signifie que $a/s_\alpha=0$ dans $S^{-1}A$, c'est-à-dire qu'il existe $s\in S$ tel que $sa=0$; mais il y a un $\beta\geq\alpha$ tel que $s\in S_\beta$, et par suite, comme $\rho_\alpha(a/s_\alpha)=\rho_\beta(sa/ss_\alpha)=0$, on voit que $\varphi$ est injectif.
On traite de même le cas d'un $A$-module $M$, et on a ainsi défini des isomorphismes canoniques
\[
  \varinjlim S_\alpha^{-1}A\xrightarrow{\sim}
  (\varinjlim S_\alpha)^{-1}A,
  \qquad
  \varinjlim S_\alpha^{-1}M\xrightarrow{\sim}
  (\varinjlim S_\alpha)^{-1}M,
\]
le second étant \emph{fonctoriel en $M$}.
\end{env}

\oldpage[0\textsubscript{I}]{17}
\begin{env}[1.4.6]
\label{0.1.4.6-fr}
Soient $S_1$, $S_2$ deux parties multiplicatives de $A$; alors $S_1S_2$ est aussi une partie multiplicative de $A$.
Désignons par $S_2'$ l'image canonique de $S_2$ dans l'anneau $S_1^{-1}A$, qui est une partie multiplicative de cet anneau.
Pour tout $A$-module $M$, il existe alors un isomorphisme fonctoriel
\[
  S_2'^{-1}(S_1^{-1}M)\xrightarrow{\sim}(S_1S_2)^{-1}M
\]
qui fait correspondre à $(m/s_1)/(s_2/1)$ l'élément $m/(s_1s_2)$.
\end{env}

\subsection{Changement d'anneau}
\label{subsection:0.1.5-fr}

\begin{env}[1.5.1]
\label{0.1.5.1-fr}
Soient $A$, $A'$ deux anneaux, $\varphi$ un homomorphisme $A'\to A$, $S$ (resp. $S'$) une partie multiplicative de $A$ (resp. $A'$), telle que $\varphi(S')\subset S$; l'homomorphisme composé
\[
  A'\xrightarrow{\varphi}A\longrightarrow S^{-1}A
\]
se factorise en
\[
  A'\longrightarrow S'^{-1}A'\xrightarrow{\varphi^{S'}}S^{-1}A
\]
en vertu de (1.2.4); on a $\varphi^{S'}(a'/s')=\varphi(a')/\varphi(s')$.
Si $A=\varphi(A')$ et $S=\varphi(S')$, $\varphi^{S'}$ est \emph{surjective}.
Si $A'=A$ et si $\varphi$ est l'identité, $\varphi^{S'}$ n'est autre que l'homomorphisme $\rho_A^{S,S'}$ défini en (1.4.1).
\end{env}

\begin{env}[1.5.2]
\label{0.1.5.2-fr}
Sous les hypothèses de (1.5.1), soit $M$ un $A$-module.
Il existe un homomorphisme canonique fonctoriel
\[
  \sigma:S'^{-1}(M_{[\varphi]})\longrightarrow
  (S^{-1}M)_{[\varphi^{S'}]}
\]
de $S'^{-1}A'$-modules, faisant correspondre à tout élément $m/s'$ de $S'^{-1}(M_{[\varphi]})$ l'élément $m/\varphi(s')$ de $(S^{-1}M)_{[\varphi^{S'}]}$; on vérifie en effet immédiatement que cette définition ne dépend pas de l'expression $m/s'$ de l'élément considéré.
Lorsque $S=\varphi(S')$, l'homomorphisme $\sigma$ est \emph{bijectif}.
Lorsque $A'=A$ et que $\varphi$ est l'identité, $\sigma$ n'est autre que l'homomorphisme $\rho_M^{S,S'}$ défini en (1.4.1).

Lorsqu'on prend en particulier $M=A$, l'homomorphisme $\varphi$ définit sur $A$ une structure de $A'$-algèbre; $S'^{-1}(A_{[\varphi]})$ est alors muni d'une structure d'anneau, pour laquelle il s'identifie à $(\varphi(S'))^{-1}A$, et l'homomorphisme
\[
  \sigma:S'^{-1}(A_{[\varphi]})\longrightarrow S^{-1}A
\]
est un homomorphisme de $S'^{-1}A'$-algèbres.
\end{env}

\begin{env}[1.5.3]
\label{0.1.5.3-fr}
Soient $M$ et $N$ deux $A$-modules; en composant les homomorphismes définis dans (1.3.4) et (1.5.2), on obtient un homomorphisme
\[
  (S^{-1}M\otimes_{S^{-1}A}S^{-1}N)_{[\varphi^{S'}]}
  \longleftarrow
  S'^{-1}((M\otimes_A N)_{[\varphi]})
\]
qui est un isomorphisme lorsque $\varphi(S')=S$.
De même, en composant les homomorphismes (1.3.5) et (1.5.2), on obtient un homomorphisme
\[
  S'^{-1}((\operatorname{Hom}_A(M,N))_{[\varphi]})
  \longrightarrow
  (\operatorname{Hom}_{S^{-1}A}(S^{-1}M,S^{-1}N))_{[\varphi^{S'}]}
\]
qui est un isomorphisme lorsque $\varphi(S')=S$ et que $M$ admet une présentation finie.
\end{env}

\begin{env}[1.5.4]
\label{0.1.5.4-fr}
Considérons maintenant un $A'$-module $N'$, et formons le produit tensoriel $N'\otimes_{A'}A_{[\varphi]}$, qui peut être considéré comme un $A$-module en posant
\[
  a.(n'\otimes b)=n'\otimes(ab).
\]
Il existe un isomorphisme fonctoriel de $S^{-1}A$-modules
\[
  \tau:(S'^{-1}N')\otimes_{S'^{-1}A'}(S^{-1}A)_{[\varphi^{S'}]}
  \xrightarrow{\sim}
  S^{-1}(N'\otimes_{A'}A_{[\varphi]})
\]
\oldpage[0\textsubscript{I}]{18}
qui, à l'élément $(n'/s')\otimes(a/s)$, fait correspondre l'élément
$(n'\otimes a)/(\varphi(s')s)$; on vérifie en effet séparément que lorsqu'on
remplace $n'/s'$ (resp. $a/s$) par une autre expression du même élément,
$(n'\otimes a)/(\varphi(s')s)$ ne change pas; d'autre part, on peut définir un
homomorphisme réciproque de $\tau$ en faisant correspondre à $(n'\otimes a)/s$
l'élément $(n'/1)\otimes(a/s)$: on utilise le fait que
$S^{-1}(N'\otimes_{A'}A_{[\varphi]})$ est canoniquement isomorphe à
$(N'\otimes_{A'}A_{[\varphi]})\otimes_A S^{-1}A$ (1.2.5), donc aussi à
$N'\otimes_{A'}(S^{-1}A)_{[\psi]}$, en désignant par $\psi$ l'homomorphisme
composé $a'\mapsto\varphi(a')/1$ de $A'$ dans $S^{-1}A$.
\end{env}

\begin{env}[1.5.5]
\label{0.1.5.5-fr}
Si $M'$ et $N'$ sont deux $A'$-modules, en composant les isomorphismes (1.3.4)
et (1.5.4), on obtient un isomorphisme
\[
  S'^{-1}M'\otimes_{S'^{-1}A'}S'^{-1}N'
  \otimes_{S'^{-1}A'}S^{-1}A
  \xrightarrow{\sim}
  S^{-1}(M'\otimes_{A'}N'\otimes_{A'}A).
\]
De même, si $M'$ admet une présentation finie, on a en vertu de (1.3.5) et
(1.5.4), un isomorphisme
\[
  \operatorname{Hom}_{S'^{-1}A'}(S'^{-1}M',S'^{-1}N')
  \otimes_{S'^{-1}A'}S^{-1}A
  \xrightarrow{\sim}
  S^{-1}(\operatorname{Hom}_{A'}(M',N')\otimes_{A'}A).
\]
\end{env}

\begin{env}[1.5.6]
\label{0.1.5.6-fr}
Sous les hypothèses de (1.5.1), soit $T$ (resp. $T'$) une seconde partie
multiplicative de $A$ (resp. $A'$) telle que $S\subset T$ (resp. $S'\subset T'$)
et $\varphi(T')\subset T$. Alors le diagramme
\[
\begin{tikzcd}[column sep=huge,row sep=large]
S'^{-1}A' \arrow[r,"\varphi^{S'}"] \arrow[d,"\rho^{T',S'}"']
  & S^{-1}A \arrow[d,"\rho^{T,S}"] \\
T'^{-1}A' \arrow[r,"\varphi^{T'}"']
  & T^{-1}A
\end{tikzcd}
\]
est commutatif. Si $M$ est un $A$-module, le diagramme
\[
\begin{tikzcd}[column sep=huge,row sep=large]
S'^{-1}(M_{[\varphi]}) \arrow[r,"\sigma"] \arrow[d,"\rho^{T',S'}"']
  & (S^{-1}M)_{[\varphi^{S'}]} \arrow[d,"\rho^{T,S}"] \\
T'^{-1}(M_{[\varphi]}) \arrow[r,"\sigma"']
  & (T^{-1}M)_{[\varphi^{T'}]}
\end{tikzcd}
\]
est commutatif. Enfin, si $N'$ est un $A'$-module, le diagramme
\[
\begin{tikzcd}[column sep=huge,row sep=large]
(S'^{-1}N')\otimes_{S'^{-1}A'}(S^{-1}A)_{[\varphi^{S'}]}
  \arrow[r,"\tau","{\sim}"'] \arrow[d]
  & S^{-1}(N'\otimes_{A'}A_{[\varphi]}) \arrow[d,"\rho^{T,S}"] \\
(T'^{-1}N')\otimes_{T'^{-1}A'}(T^{-1}A)_{[\varphi^{T'}]}
  \arrow[r,"\tau"',"{\sim}"]
  & T^{-1}(N'\otimes_{A'}A_{[\varphi]})
\end{tikzcd}
\]
est commutatif, la flèche verticale de gauche étant obtenue en appliquant
$\rho_{N'}^{T',S'}$ à $S'^{-1}N'$ et $\rho_A^{T,S}$ à $S^{-1}A$.
\end{env}

\oldpage[0\textsubscript{I}]{19}

\begin{env}[1.5.7]
\label{0.1.5.7-fr}
Soient $A''$ un troisième anneau, $\varphi':A''\to A'$ un homomorphisme
d'anneaux, $S''$ une partie multiplicative de $A''$ telle que
$\varphi'(S'')\subset S'$. Posons $\varphi''=\varphi\circ\varphi'$; alors on a
\[
  \varphi''^{S''}=\varphi^{S'}\circ\varphi'^{S''}.
\]
Soit $M$ un $A$-module; on a évidemment
$M_{[\varphi'']}=(M_{[\varphi]})_{[\varphi']}$; si $\sigma'$ et $\sigma''$ sont
les homomorphismes définis à partir de $\varphi'$ et $\varphi''$ comme $\sigma$
est défini dans (1.5.2) à partir de $\varphi$, on a la formule de transitivité
\[
  \sigma''=\sigma\circ\sigma'.
\]
Enfin, soit $N''$ un $A''$-module; le $A$-module
$N''\otimes_{A''}A_{[\varphi'']}$ s'identifie canoniquement à
\[
  (N''\otimes_{A''}A'_{[\varphi']})\otimes_{A'}A_{[\varphi]},
\]
et de même, le $S^{-1}A$-module
\[
  (S''^{-1}N'')\otimes_{S''^{-1}A''}
  (S^{-1}A)_{[\varphi''^{S''}]}
\]
s'identifie canoniquement à
\[
  \bigl((S''^{-1}N'')\otimes_{S''^{-1}A''}
  (S'^{-1}A')_{[\varphi'^{S''}]}\bigr)
  \otimes_{S'^{-1}A'}(S^{-1}A)_{[\varphi^{S'}]}.
\]
Avec ces identifications, si $\tau'$ et $\tau''$ sont les isomorphismes définis
à partir de $\varphi'$ et $\varphi''$ comme $\tau$ est défini dans (1.5.4) à
partir de $\varphi$, on a la formule de transitivité
\[
  \tau''=\tau\circ(\tau'\otimes 1).
\]
\end{env}

\begin{env}[1.5.8]
\label{0.1.5.8-fr}
Soit $A$ un sous-anneau d'un anneau $B$; pour tout idéal premier
\emph{minimal} $\mathfrak p$ de $A$, il existe un idéal premier minimal
$\mathfrak q$ de $B$ tel que $\mathfrak p=A\cap\mathfrak q$. En effet,
$A_{\mathfrak p}$ est un sous-anneau de $B_{\mathfrak p}$ (1.3.2) et possède
un seul idéal premier $\mathfrak p'$ (1.2.6); comme $B_{\mathfrak p}$ n'est pas
réduit à $0$, il possède au moins un idéal premier $\mathfrak q'$ et on a
nécessairement $\mathfrak q'\cap A_{\mathfrak p}=\mathfrak p'$; l'idéal premier
$\mathfrak q_1$ de $B$, image réciproque de $\mathfrak q'$ est donc tel que
$\mathfrak q_1\cap A=\mathfrak p$, et \emph{a fortiori} on a
$\mathfrak q\cap A=\mathfrak p$ pour tout idéal premier minimal $\mathfrak q$
de $B$ contenu dans $\mathfrak q_1$.
\end{env}

\subsection{Identification du module $M_f$ à une limite inductive}
\label{subsection:0.1.6-fr}

\begin{env}[1.6.1]
\label{0.1.6.1-fr}
Soient $M$ un $A$-module, $f$ un élément de $A$. Considérons une suite
$(M_n)$ de $A$-modules, tous identiques à $M$, et pour tout couple d'entiers
$m\leq n$, soit $\varphi_{nm}$ l'homomorphisme $z\mapsto f^{n-m}z$ de $M_m$
dans $M_n$; il est immédiat que $((M_n),(\varphi_{nm}))$ est un \emph{système
inductif} de $A$-modules; soit $N=\varinjlim M_n$ la limite inductive de ce
système. Nous allons définir un $A$-isomorphisme canonique \emph{fonctoriel} de
$N$ sur $M_f$. Pour cela, remarquons que, pour tout $n$,
$\theta_n:z\mapsto z/f^n$ est un $A$-homomorphisme de $M=M_n$ dans $M_f$, et
il résulte des définitions que l'on a $\theta_n\circ\varphi_{nm}=\theta_m$ pour
$m\leq n$. Il existe donc un $A$-homomorphisme $\theta:N\to M_f$ tel que, si
$\varphi_n$ désigne l'homomorphisme canonique $M_n\to N$, on ait
$\theta_n=\theta\circ\varphi_n$ pour tout $n$. Comme par hypothèse tout élément
de $M_f$ est de la forme $z/f^n$ pour un $n$ au moins, il est clair que
$\theta$ est surjectif. D'autre part, si $\theta(\varphi_n(z))=0$, autrement dit
$z/f^n=0$, il existe un entier $k>0$ tel que $f^kz=0$, donc
$\varphi_{n+k,n}(z)=0$, ce qui entraîne $\varphi_n(z)=0$. On peut donc
identifier $M_f$ et $\varinjlim M_n$ au moyen de $\theta$.
\end{env}

\begin{env}[1.6.2]
\label{0.1.6.2-fr}
Écrivons maintenant $M_{f,n}$, $\varphi_{nm}^f$ et $\varphi_n^f$ au lieu de
$M_n$, $\varphi_{nm}$ et $\varphi_n$. Soit $g$ un second élément de $A$. Comme
$f^n$ divise $f^ng^n$, on a un homomorphisme fonctoriel
\[
  \rho_{fg,f}:M_f\longrightarrow M_{fg}\qquad\text{(1.4.1 et 1.4.3)};
\]

\oldpage[0\textsubscript{I}]{20}

si on identifie $M_f$ et $M_{fg}$ à $\varinjlim M_{f,n}$ et
$\varinjlim M_{fg,n}$ respectivement, $\rho_{fg,f}$ s'identifie à la
\emph{limite inductive} des applications
$\rho_{fg,f}^n:M_{f,n}\to M_{fg,n}$ définies par
$\rho_{fg,f}^n(z)=g^nz$. En effet, cela résulte immédiatement de la
commutativité du diagramme
\[
\begin{tikzcd}[column sep=huge,row sep=large]
M_{f,n} \arrow[r,"\rho_{fg,f}^n"] \arrow[d,"\varphi_n^f"']
  & M_{fg,n} \arrow[d,"\varphi_n^{fg}"] \\
M_f \arrow[r,"\rho_{fg,f}"]
  & M_{fg}.
\end{tikzcd}
\]
\end{env}

\subsection{Support d'un module}
\label{subsection:0.1.7-fr}

\begin{env}[1.7.1]
\label{0.1.7.1-fr}
Étant donné un $A$-module $M$, on appelle \emph{support} de $M$ et on note
$\operatorname{Supp}(M)$ l'ensemble des idéaux premiers $\mathfrak p$ de $A$
tels que $M_{\mathfrak p}\neq0$. Pour que $M=0$, il faut et il suffit que
$\operatorname{Supp}(M)=\varnothing$, car si $M_{\mathfrak p}=0$ pour tout
$\mathfrak p$, l'annulateur d'un élément $x\in M$ ne peut être contenu dans
aucun idéal premier de $A$, donc est $A$ tout entier.
\end{env}

\begin{env}[1.7.2]
\label{0.1.7.2-fr}
Si $0\to N\to M\to P\to0$ est une suite exacte de $A$-modules, on a
\[
  \operatorname{Supp}(M)=\operatorname{Supp}(N)\cup\operatorname{Supp}(P)
\]
car pour tout idéal premier $\mathfrak p$ de $A$, la suite
$0\to N_{\mathfrak p}\to M_{\mathfrak p}\to P_{\mathfrak p}\to0$ est exacte
(1.3.2) et pour que $M_{\mathfrak p}=0$, il faut et il suffit que
$N_{\mathfrak p}=P_{\mathfrak p}=0$.
\end{env}

\begin{env}[1.7.3]
\label{0.1.7.3-fr}
Si $M$ est somme d'une famille $(M_\lambda)$ de sous-modules,
$M_{\mathfrak p}$ est somme des $(M_\lambda)_{\mathfrak p}$ pour tout idéal
premier $\mathfrak p$ de $A$ (1.3.3 et 1.3.2), donc
\[
  \operatorname{Supp}(M)=\bigcup_\lambda\operatorname{Supp}(M_\lambda).
\]
\end{env}

\begin{env}[1.7.4]
\label{0.1.7.4-fr}
Si $M$ est un $A$-module \emph{de type fini}, $\operatorname{Supp}(M)$ est
l'ensemble des idéaux premiers \emph{contenant l'annulateur de $M$}. En effet,
si $M$ est monogène et engendré par $x$, dire que $M_{\mathfrak p}=0$ signifie
qu'il existe $s\notin\mathfrak p$ tel que $s.x=0$, donc que $\mathfrak p$ ne
contient pas l'annulateur de $x$. Si maintenant $M$ admet un système fini
$(x_i)_{1\leq i\leq n}$ de générateurs et si $\mathfrak a_i$ est l'annulateur
de $x_i$, il résulte de (1.7.3) que $\operatorname{Supp}(M)$ est l'ensemble des
$\mathfrak p$ contenant l'un des $\mathfrak a_i$, ou, ce qui revient au même,
l'ensemble des $\mathfrak p$ contenant $\mathfrak a=\bigcap_i\mathfrak a_i$,
qui est l'annulateur de $M$.
\end{env}

\begin{env}[1.7.5]
\label{0.1.7.5-fr}
Si $M$ et $N$ sont deux $A$-modules \emph{de type fini}, on a
\[
  \operatorname{Supp}(M\otimes_A N)
  =\operatorname{Supp}(M)\cap\operatorname{Supp}(N).
\]
Il s'agit de voir que si $\mathfrak p$ est un idéal premier de $A$, la
condition $M_{\mathfrak p}\otimes_{A_{\mathfrak p}}N_{\mathfrak p}\neq0$ est
équivalente à «~$M_{\mathfrak p}\neq0$ et $N_{\mathfrak p}\neq0$~» (compte
tenu de (1.3.4)). Autrement dit, il s'agit de voir que si $P$, $Q$ sont deux
modules de type fini sur un anneau \emph{local} $B$, non réduits à $0$, alors
$P\otimes_B Q\neq0$. Soit $\mathfrak m$ l'idéal maximal de $B$. En vertu du
lemme de Nakayama, les espaces vectoriels $P/\mathfrak mP$ et
$Q/\mathfrak mQ$ ne sont pas réduits à $0$, donc il en est de même de leur
produit tensoriel
\[
  (P/\mathfrak mP)\otimes_{B/\mathfrak m}(Q/\mathfrak mQ)
  =(P\otimes_B Q)\otimes_B(B/\mathfrak m),
\]
d'où la conclusion.

En particulier, si $M$ est un $A$-module de type fini, $\mathfrak a$ un idéal
de $A$, $\operatorname{Supp}(M/\mathfrak aM)$ est l'ensemble des idéaux
premiers contenant à la fois $\mathfrak a$ et l'annulateur $\mathfrak n$ de
$M$ (1.7.4), autrement dit l'ensemble des idéaux premiers contenant
$\mathfrak a+\mathfrak n$.
\end{env}

\oldpage[0\textsubscript{I}]{21}

\section{Espaces irréductibles. Espaces noethériens}
\label{section:0.2-fr}

\subsection{Espaces irréductibles}
\label{subsection:0.2.1-fr}

\begin{env}[2.1.1]
\label{0.2.1.1-fr}
On dit qu'un espace topologique $X$ est \emph{irréductible} s'il est non vide
et s'il n'est pas réunion de deux sous-espaces fermés distincts de $X$. Il
revient au même de dire que $X\neq\varnothing$ et que l'intersection de deux
ouverts (et par suite d'un nombre fini d'ouverts) non vides de $X$ est non
vide, ou que tout ouvert non vide est partout dense, ou que toute partie
fermée $\neq X$ est \emph{rare}, ou enfin que tout ouvert de $X$ est
\emph{connexe}.
\end{env}

\begin{env}[2.1.2]
\label{0.2.1.2-fr}
Pour qu'un sous-espace $Y$ d'un espace topologique $X$ soit irréductible, il
faut et il suffit que son adhérence $\overline Y$ soit irréductible. En
particulier, tout sous-espace qui est l'adhérence $\overline{\{x\}}$ d'un
sous-espace réduit à un point est irréductible; nous exprimerons la relation
$y\in\overline{\{x\}}$ (équivalente à
$\overline{\{y\}}\subset\overline{\{x\}}$) en disant que $y$ est
\emph{spécialisation de $x$} ou que $x$ est une \emph{générisation de $y$}.
Lorsqu'il existe dans un espace irréductible $X$ un point $x$ tel que
$X=\overline{\{x\}}$, nous dirons que $x$ est \emph{point générique} de $X$.
Tout ouvert non vide de $X$ contient alors $x$ et tout sous-espace contenant
$x$ admet $x$ pour point générique.
\end{env}

\begin{env}[2.1.3]
\label{0.2.1.3-fr}
Rappelons qu'on appelle \emph{espace de Kolmogoroff} un espace topologique $X$
vérifiant l'axiome de séparation~:

\smallskip
\noindent
\textup{(T$_0$)} Si $x\neq y$ sont deux points quelconques de $X$, il existe
un ensemble ouvert contenant l'un des points $x$, $y$ et non l'autre.

Si un espace de Kolmogoroff irréductible admet un point générique, il n'en
admet \emph{qu'un seul} puisqu'un ouvert non vide contient tout point
générique.

Rappelons qu'un espace topologique $X$ est dit \emph{quasi-compact} si, de tout
recouvrement ouvert de $X$, on peut extraire un recouvrement fini de $X$ (ou,
ce qui revient au même, si toute famille filtrante décroissante d'ensembles
fermés non vides a une intersection non vide). Si $X$ est un espace
quasi-compact, toute partie fermée non vide $A$ de $X$ contient un ensemble
fermé non vide \emph{minimal} $M$, car l'ensemble des parties fermées non vides
de $A$ est inductif pour la relation $\supset$; si en outre $X$ est un espace
de Kolmogoroff, $M$ est nécessairement réduite à un seul point (ou, comme on
dit par abus de langage, est un \emph{point fermé}).
\end{env}

\begin{env}[2.1.4]
\label{0.2.1.4-fr}
Dans un espace irréductible $X$, tout sous-espace ouvert non vide $U$ est
irréductible, et si $X$ admet un point générique $x$, $x$ est aussi point
générique de $U$.

Soit $(U_\alpha)$ un recouvrement (dont l'ensemble d'indices est non vide) d'un
espace topologique $X$, formé d'ouverts non vides; pour que $X$ soit
irréductible, il faut et il suffit que $U_\alpha$ soit irréductible pour tout
$\alpha$, et que $U_\alpha\cap U_\beta\neq\varnothing$ quels que soient
$\alpha$, $\beta$. La condition est évidemment nécessaire; pour voir qu'elle
est suffisante, il suffit de prouver que si $V$ est un ouvert non vide de $X$,
$V\cap U_\alpha$ est non vide pour tout $\alpha$, car alors $V\cap U_\alpha$
est dense dans $U_\alpha$ pour tout $\alpha$, et par suite $V$ est dense dans
$X$. Or, il y a au moins un indice $\gamma$ tel que
$V\cap U_\gamma\neq\varnothing$, donc $V\cap U_\gamma$ est dense dans
$U_\gamma$, et comme, pour tout $\alpha$,
$U_\alpha\cap U_\gamma\neq\varnothing$, on a aussi
$V\cap U_\alpha\cap U_\gamma\neq\varnothing$.
\end{env}

\oldpage[0\textsubscript{I}]{22}

\begin{env}[2.1.5]
\label{0.2.1.5-fr}
Soient $X$ un espace irréductible, $f$ une application continue de $X$ dans un
espace topologique $Y$. Alors $f(X)$ est irréductible, et si $x$ est point
générique de $X$, $f(x)$ est point générique de $f(X)$ et par suite aussi de
$\overline{f(X)}$. En particulier, si de plus $Y$ est irréductible et a un seul
point générique $y$, pour que $f(X)$ soit partout dense, il faut et il suffit
que $f(x)=y$.
\end{env}

\begin{env}[2.1.6]
\label{0.2.1.6-fr}
Tout sous-espace irréductible d'un espace topologique $X$ est contenu dans un
sous-espace irréductible \emph{maximal}, qui est nécessairement fermé. Les
sous-espaces irréductibles maximaux de $X$ sont appelés les \emph{composantes
irréductibles} de $X$. Si $Z_1$, $Z_2$ sont deux composantes irréductibles
distinctes de l'espace $X$, $Z_1\cap Z_2$ est un ensemble fermé \emph{rare}
dans chacun des sous-espaces $Z_1$, $Z_2$; en particulier, si une composante
irréductible de $X$ admet un point générique (2.1.2) un tel point ne peut
appartenir à aucune autre composante irréductible. Si $X$ n'a qu'un nombre
\emph{fini} de composantes irréductibles $Z_i$ ($1\leq i\leq n$), et si, pour
chaque $i$, on pose
\[
  U_i=Z_i\cap\complement\Bigl(\bigcup_{j\neq i}Z_j\Bigr),
\]
les $U_i$ sont ouverts, irréductibles, deux à deux sans point commun et leur
réunion est dense dans $X$.

Soit $U$ une partie ouverte d'un espace topologique $X$. Si $Z$ est une partie
irréductible de $X$ rencontrant $U$, $Z\cap U$ est ouvert et dense dans $Z$,
donc irréductible; inversement, pour toute partie fermée irréductible $Y$ de
$U$, l'adhérence $\overline Y$ de $Y$ dans $X$ est irréductible et
$\overline Y\cap U=Y$. On en conclut qu'il y a \emph{correspondance
biunivoque} entre les composantes irréductibles de $U$ et les composantes
irréductibles de $X$ qui rencontrent $U$.
\end{env}

\begin{env}[2.1.7]
\label{0.2.1.7-fr}
Si un espace topologique $X$ est réunion d'un nombre \emph{fini} de
sous-espaces irréductibles fermés $Y_i$, les composantes irréductibles de $X$
sont les éléments maximaux de l'ensemble des $Y_i$, car si $Z$ est une partie
fermée irréductible de $X$, $Z$ est réunion des $Z\cap Y_i$, d'où on tire
aussitôt que $Z$ doit être contenu dans un des $Y_i$. Soit $Y$ un sous-espace
d'un espace topologique $X$, et supposons que $Y$ n'ait qu'un nombre fini de
composantes irréductibles $Y_i$ ($1\leq i\leq n$); alors les adhérences
$\overline{Y_i}$ dans $X$ sont les composantes irréductibles de $\overline Y$.
\end{env}

\begin{env}[2.1.8]
\label{0.2.1.8-fr}
Soit $Y$ un espace irréductible admettant un seul point générique $y$. Soient
$X$ un espace topologique, $f$ une application continue de $X$ dans $Y$.
Alors, pour toute composante irréductible $Z$ de $X$ rencontrant $f^{-1}(y)$,
$f(Z)$ est dense dans $Y$. La réciproque n'est pas nécessairement vraie;
toutefois, si $Z$ admet un point générique $z$, et si $f(Z)$ est dense dans
$Y$, on a nécessairement $f(z)=y$ (2.1.5); en outre, $Z\cap f^{-1}(y)$ est
alors l'adhérence de $\{z\}$ dans $f^{-1}(y)$ et est donc irréductible, et
comme toute partie irréductible de $f^{-1}(y)$ contenant $z$ est
nécessairement contenue dans $Z$ (2.1.6), $z$ est point générique de
$Z\cap f^{-1}(y)$. Comme toute composante irréductible de $f^{-1}(y)$ est
contenue dans une composante irréductible de $X$, on voit que si toute
composante irréductible $Z$ de $X$ rencontrant $f^{-1}(y)$ admet un point
générique, alors il y a \emph{correspondance biunivoque} entre l'ensemble de
ces composantes et l'ensemble des composantes irréductibles
$Z\cap f^{-1}(y)$ de $f^{-1}(y)$, les points génériques de $Z$ étant
identiques à ceux de $Z\cap f^{-1}(y)$.
\end{env}

\oldpage[0\textsubscript{I}]{23}

\subsection{Espaces noethériens}
\label{subsection:0.2.2-fr}

\begin{env}[2.2.1]
\label{0.2.2.1-fr}
On dit qu'un espace topologique $X$ est \emph{noethérien} si l'ensemble des
ouverts de $X$ vérifie la condition \emph{maximale}, ou, ce qui revient au
même, si l'ensemble des fermés de $X$ vérifie la condition \emph{minimale}. On
dit que $X$ est \emph{localement noethérien} si tout $x\in X$ admet un
voisinage qui est un sous-espace noethérien.
\end{env}

\begin{env}[2.2.2]
\label{0.2.2.2-fr}
Soit $E$ un ensemble ordonné vérifiant la condition \emph{minimale}, et soit
$P$ une propriété des éléments de $E$ soumise à la condition suivante~: si
$a\in E$ est tel que pour tout $x<a$, $P(x)$ soit vraie, alors $P(a)$ est
vraie. Dans ces conditions, \emph{$P(x)$ est vraie pour tout $x\in E$}
(«~principe de récurrence noethérienne~»). En effet, soit $F$ l'ensemble des
$x\in E$ pour lesquels $P(x)$ est fausse~; si $F$ était non vide, il aurait un
élément minimal $a$, et comme alors $P(x)$ est vraie pour tout $x<a$, $P(a)$
serait aussi vraie, ce qui est contradictoire.

Nous appliquerons en particulier ce principe lorsque $E$ est un \emph{ensemble
de parties fermées d'un espace noethérien}.
\end{env}

\begin{env}[2.2.3]
\label{0.2.2.3-fr}
Tout sous-espace d'un espace noethérien est noethérien. Inversement, tout
espace topologique réunion finie de sous-espaces noethériens est noethérien.
\end{env}

\begin{env}[2.2.4]
\label{0.2.2.4-fr}
Tout espace noethérien est quasi-compact~; inversement, tout espace topologique
dans lequel tout ouvert est quasi-compact est noethérien.
\end{env}

\begin{env}[2.2.5]
\label{0.2.2.5-fr}
Un espace noethérien n'a qu'un nombre \emph{fini} de composantes
irréductibles, comme on le voit par récurrence noethérienne.
\end{env}

\section{Compléments sur les faisceaux}
\label{section:0.3-fr}

\subsection{Faisceaux à valeurs dans une catégorie}
\label{subsection:0.3.1-fr}

\begin{env}[3.1.1]
\label{0.3.1.1-fr}
Soient $\mathbf{K}$ une catégorie, $(A_\alpha)_{\alpha\in I}$,
$(A_{\alpha\beta})_{(\alpha,\beta)\in I\times I}$ deux familles d'objets de
$\mathbf{K}$ telles que $A_{\beta\alpha}=A_{\alpha\beta}$,
$(\rho_{\alpha\beta})_{(\alpha,\beta)\in I\times I}$ une famille de morphismes
$\rho_{\alpha\beta}:A_\alpha\to A_{\alpha\beta}$. Nous dirons qu'un couple
formé d'un objet $A$ de $\mathbf{K}$ et d'une famille de morphismes
$\rho_\alpha:A\to A_\alpha$ est \emph{solution du problème universel} défini
par la donnée des familles $(A_\alpha)$, $(A_{\alpha\beta})$ et
$(\rho_{\alpha\beta})$ si, pour tout objet $B$ de $\mathbf{K}$, l'application
qui, à tout $f\in\operatorname{Hom}(B,A)$ fait correspondre la famille
$(\rho_\alpha\circ f)\in\prod_\alpha\operatorname{Hom}(B,A_\alpha)$ est une
\emph{bijection} de $\operatorname{Hom}(B,A)$ sur l'ensemble des $(f_\alpha)$
telles que
$\rho_{\alpha\beta}\circ f_\alpha=\rho_{\beta\alpha}\circ f_\beta$ pour tout
couple d'indices $(\alpha,\beta)$. On voit aussitôt que s'il existe une telle
solution, elle est unique à un isomorphisme près.
\end{env}

\begin{env}[3.1.2]
\label{0.3.1.2-fr}
Nous ne rappellerons pas la définition d'un \emph{préfaisceau}
$U\to\mathcal{F}(U)$ sur un espace topologique $X$, à valeurs dans une
catégorie $\mathbf{K}$ (G, I, 1.9)~; nous dirons qu'un tel préfaisceau est un
\emph{faisceau à valeurs dans $\mathbf{K}$} s'il satisfait à l'axiome
suivant~:

\smallskip
\noindent
\textup{(F)} \emph{Pour tout recouvrement $(U_\alpha)$ d'un ouvert $U$ de $X$
par des ouverts $U_\alpha$ contenus dans $U$, si on désigne par $\rho_\alpha$
(resp. $\rho_{\alpha\beta}$) le morphisme de restriction}
\[
  \mathcal{F}(U)\to\mathcal{F}(U_\alpha)
  \qquad
  \bigl(\text{resp. }\mathcal{F}(U_\alpha)
  \to\mathcal{F}(U_\alpha\cap U_\beta)\bigr),
\]

\oldpage[0\textsubscript{I}]{24}

\emph{le couple formé de $\mathcal{F}(U)$ et de la famille $(\rho_\alpha)$ est
solution du problème universel pour $(\mathcal{F}(U_\alpha))$,
$(\mathcal{F}(U_\alpha\cap U_\beta))$ et $(\rho_{\alpha\beta})$ (3.1.1)}%
\footnote{C'est un cas particulier de la notion générale de \emph{limite
projective} (non filtrante) (\emph{voir} (T, I, 1.8) et le livre en préparation
annoncé dans l'Introduction).}.

Il revient au même de dire que, pour tout objet $T$ de $\mathbf{K}$, la
famille $U\to\operatorname{Hom}(T,\mathcal{F}(U))$ est un \emph{faisceau
d'ensembles}.
\end{env}

\begin{env}[3.1.3]
\label{0.3.1.3-fr}
Supposons que $\mathbf{K}$ soit la catégorie définie par une «~espèce de
structure avec morphismes~» $\Sigma$, les objets de $\mathbf{K}$ étant donc
les ensembles munis de structures d'espèce $\Sigma$ et les morphismes ceux de
$\Sigma$. Supposons que la catégorie $\mathbf{K}$ vérifie en outre la
condition suivante~:

\smallskip
\noindent
\textup{(E)} Si $(A,(\rho_\alpha))$ est solution d'un problème d'application
universelle \emph{dans la catégorie $\mathbf{K}$} pour des familles
$(A_\alpha)$, $(A_{\alpha\beta})$, $(\rho_{\alpha\beta})$, alors c'est aussi
une solution du problème d'application universelle pour les mêmes familles
\emph{dans la catégorie des ensembles} (c'est-à-dire quand on considère $A$,
les $A_\alpha$ et $A_{\alpha\beta}$ comme des ensembles, les $\rho_\alpha$ et
$\rho_{\alpha\beta}$ comme des applications)%
\footnote{On peut prouver que cela signifie aussi que le foncteur canonique
$\mathbf{K}\to(\mathrm{Ens})$ \emph{permute aux limites projectives} (non
nécessairement filtrantes).}.

Dans ces conditions, la condition (F) entraîne que, considéré comme
préfaisceau \emph{d'ensembles}, $U\to\mathcal{F}(U)$ est un \emph{faisceau}.
En outre, pour qu'une application $u:T\to\mathcal{F}(U)$ soit un morphisme de
$\mathbf{K}$, il faut et il suffit, en vertu de (F), que chaque application
$\rho_\alpha\circ u$ soit un morphisme $T\to\mathcal{F}(U_\alpha)$, ce qui
signifie que la structure d'espèce $\Sigma$ sur $\mathcal{F}(U)$ est
\emph{structure initiale} pour les morphismes $\rho_\alpha$. Réciproquement,
supposons qu'un préfaisceau $U\to\mathcal{F}(U)$ sur $X$, à valeurs dans
$\mathbf{K}$, soit un \emph{faisceau d'ensembles}, et vérifie la condition
précédente~; il est clair alors qu'il satisfait à (F), donc est un
\emph{faisceau à valeurs dans $\mathbf{K}$}.
\end{env}

\begin{env}[3.1.4]
\label{0.3.1.4-fr}
Lorsque $\Sigma$ est l'espèce de structure de groupe ou d'anneau, le fait que
le préfaisceau $U\to\mathcal{F}(U)$ à valeurs dans $\mathbf{K}$ est un
\emph{faisceau d'ensembles} entraîne \emph{ipso facto} que c'est un
\emph{faisceau à valeurs dans $\mathbf{K}$} (autrement dit, un faisceau de
groupes ou d'anneaux au sens de (G))%
\footnote{Cela tient à ce que dans la catégorie $\mathbf{K}$, tout morphisme
qui est une \emph{bijection} (en tant qu'application d'ensembles) est un
\emph{isomorphisme}. Cela n'est plus vrai lorsque $\mathbf{K}$ est la
catégorie des espaces topologiques, par exemple.}. Mais il n'en est plus de
même lorsque par exemple $\mathbf{K}$ est la catégorie des \emph{anneaux
topologiques} (avec pour morphismes les représentations continues)~: un
faisceau à valeurs dans $\mathbf{K}$ est un faisceau d'anneaux
$U\to\mathcal{F}(U)$ tel que, pour tout ouvert $U$ et tout recouvrement de $U$
par des ouverts $U_\alpha\subset U$, la topologie de l'anneau $\mathcal{F}(U)$
soit \emph{la moins fine} rendant continues les représentations
$\mathcal{F}(U)\to\mathcal{F}(U_\alpha)$. On dira dans ce cas que
$U\to\mathcal{F}(U)$ considéré comme faisceau d'anneaux (sans topologie) est
\emph{sous-jacent} au faisceau d'anneaux topologiques
$U\to\mathcal{F}(U)$. Les morphismes
$u_V:\mathcal{F}(V)\to\mathcal{G}(V)$ ($V$ ouvert arbitraire de $X$) de
faisceaux d'anneaux topologiques sont donc des homomorphismes des faisceaux
d'anneaux sous-jacents, tels que $u_V$ soit \emph{continu} pour tout ouvert
$V\subset X$~; pour les distinguer des homomorphismes quelconques des
faisceaux d'anneaux sous-jacents, on les appellera homomorphismes
\emph{continus} de faisceaux d'anneaux topologiques. On a des définitions et
conventions analogues pour les faisceaux d'espaces topologiques ou de groupes
topologiques.
\end{env}

\oldpage[0\textsubscript{I}]{25}

\begin{env}[3.1.5]
\label{0.3.1.5-fr}
Il est clair que pour toute catégorie $\mathbf{K}$, si $\mathcal{F}$ est un
préfaisceau (resp. un faisceau) sur $X$ à valeurs dans $\mathbf{K}$ et $U$ un
ouvert de $X$, les $\mathcal{F}(V)$ pour les ouverts $V\subset U$ constituent
un préfaisceau (resp. un faisceau) à valeurs dans $\mathbf{K}$, que l'on
appelle préfaisceau (resp. faisceau) \emph{induit} par $\mathcal{F}$ sur $U$ et
que l'on note $\mathcal{F}|U$.

Pour tout morphisme $u:\mathcal{F}\to\mathcal{G}$ de préfaisceaux sur $X$ à
valeurs dans $\mathbf{K}$, on désignera par $u|U$ le morphisme
$\mathcal{F}|U\to\mathcal{G}|U$ formé des $u_V$ pour $V\subset U$.
\end{env}

\begin{env}[3.1.6]
\label{0.3.1.6-fr}
Supposons maintenant que la catégorie $\mathbf{K}$ admette des \emph{limites
inductives} (T, 1.8)~; alors, pour tout préfaisceau (et en particulier tout
faisceau) $\mathcal{F}$ sur $X$ à valeurs dans $\mathbf{K}$ et tout $x\in X$,
on peut définir la \emph{fibre} $\mathcal{F}_x$ comme l'objet de $\mathbf{K}$
limite inductive des $\mathcal{F}(U)$ selon l'ensemble filtrant (pour
$\supset$) des voisinages ouverts $U$ de $x$ dans $X$, et pour les morphismes
$\rho_U^V:\mathcal{F}(V)\to\mathcal{F}(U)$. Si
$u:\mathcal{F}\to\mathcal{G}$ est un morphisme de préfaisceaux à valeurs dans
$\mathbf{K}$, on définit pour tout $x\in X$ le morphisme
$u_x:\mathcal{F}_x\to\mathcal{G}_x$ comme la limite inductive des
$u_U:\mathcal{F}(U)\to\mathcal{G}(U)$ selon l'ensemble des voisinages ouverts
de $x$~; on définit ainsi $\mathcal{F}_x$ comme foncteur covariant en
$\mathcal{F}$, à valeurs dans $\mathbf{K}$, pour tout $x\in X$.

Lorsque $\mathbf{K}$ est en outre définie par une espèce de structure avec
morphismes $\Sigma$, on appelle encore \emph{sections au-dessus de $U$} d'un
\emph{faisceau $\mathcal{F}$} à valeurs dans $\mathbf{K}$ les éléments de
$\mathcal{F}(U)$, et on écrit alors $\Gamma(U,\mathcal{F})$ au lieu de
$\mathcal{F}(U)$~; pour $s\in\Gamma(U,\mathcal{F})$, $V$ ouvert contenu dans
$U$, on écrit $s|V$ au lieu de $\rho_V^U(s)$~; pour tout $x\in U$, l'image
canonique de $s$ dans $\mathcal{F}_x$ est le \emph{germe} de $s$ au point $x$,
noté $s_x$ (\emph{nous n'emploierons jamais la notation $s(x)$ dans ce sens},
cette notation étant réservée pour une autre notion relative aux faisceaux
particuliers qui seront considérés dans ce Traité (5.5.1)).

Si alors $u:\mathcal{F}\to\mathcal{G}$ est un morphisme de faisceaux à valeurs
dans $\mathbf{K}$, on écrira $u(s)$ au lieu de $u_V(s)$ pour tout
$s\in\Gamma(U,\mathcal{F})$.

Si $\mathcal{F}$ est un faisceau de groupes commutatifs, ou d'anneaux, ou de
modules, on dit que l'ensemble des $x\in X$ tels que
$\mathcal{F}_x\neq\{0\}$ est le \emph{support} de $\mathcal{F}$, noté
$\operatorname{Supp}(\mathcal{F})$~; cet ensemble n'est pas nécessairement
fermé dans $X$.

Lorsque $\mathbf{K}$ est définie par une espèce de structure avec morphismes,
\emph{nous nous abstiendrons systématiquement de faire intervenir le point de
vue des «~espaces étalés~»} en ce qui concerne les faisceaux à valeurs dans
$\mathbf{K}$~; autrement dit, nous ne considérerons jamais un faisceau comme
un espace topologique (ni même comme l'ensemble réunion de ses fibres), et
nous ne considérerons pas davantage un morphisme
$u:\mathcal{F}\to\mathcal{G}$ de tels faisceaux sur $X$ comme une application
continue d'espaces topologiques.
\end{env}

\subsection{Préfaisceaux sur une base d'ouverts}
\label{subsection:0.3.2-fr}

\begin{env}[3.2.1]
\label{0.3.2.1-fr}
Nous nous restreindrons dans ce qui suit à des catégories $\mathbf{K}$
admettant des \emph{limites projectives} (généralisées, c'est-à-dire
correspondant à des ensembles préordonnés non nécessairement filtrants,
cf. (T, 1.8)). Soient $X$ un espace topologique, $\mathfrak{B}$ une base
d'ouverts pour la topologie de $X$. Nous appellerons \emph{préfaisceau sur
$\mathfrak{B}$, à valeurs dans $\mathbf{K}$}, une famille d'objets
$\mathcal{F}(U)\in\mathbf{K}$, attachés à chaque $U\in\mathfrak{B}$, et une
famille de morphismes
$\rho_U^V:\mathcal{F}(V)\to\mathcal{F}(U)$ définis pour tout couple $(U,V)$
d'éléments de $\mathfrak{B}$ tels que $U\subset V$,

\oldpage[0\textsubscript{I}]{26}

avec les conditions $\rho_U^U=\text{identité}$ et
$\rho_U^W=\rho_U^V\circ\rho_V^W$ si $U$, $V$, $W$ dans $\mathfrak{B}$ sont
tels que $U\subset V\subset W$. On peut lui associer un \emph{préfaisceau à
valeurs dans $\mathbf{K}$}~: $U\to\mathcal{F}'(U)$ au sens ordinaire, en
prenant pour tout ouvert $U$,
\[
  \mathcal{F}'(U)=\varprojlim_{V\subset U}\mathcal{F}(V),
\]
où $V$ parcourt l'ensemble ordonné (pour $\subset$, \emph{non filtrant en
général}) des ensembles $V\in\mathfrak{B}$ tels que $V\subset U$, car les
$\mathcal{F}(V)$ forment un système projectif pour les $\rho_V^W$
($V\subset W\subset U$, $V\in\mathfrak{B}$, $W\in\mathfrak{B}$). En effet,
si $U$, $U'$ sont deux ouverts de $X$ tels que $U\subset U'$, on définit
$\rho'_U{}^{U'}$ comme la limite projective (pour $V\subset U$) des morphismes
canoniques $\mathcal{F}'(U')\to\mathcal{F}(V)$, autrement dit l'unique
morphisme $\mathcal{F}'(U')\to\mathcal{F}'(U)$, qui, composé avec les
morphismes canoniques $\mathcal{F}'(U)\to\mathcal{F}(V)$, donne les morphismes
canoniques $\mathcal{F}'(U')\to\mathcal{F}(V)$~; la vérification de la
transitivité de $\rho'_U{}^{U'}$ est alors immédiate. De plus, si
$U\in\mathfrak{B}$, le morphisme canonique
$\mathcal{F}'(U)\to\mathcal{F}(U)$ est un isomorphisme, permettant d'identifier
ces deux objets%
\footnote{Si $X$ est un espace \emph{noethérien}, on peut encore définir
$\mathcal{F}'(U)$ et montrer que c'est un préfaisceau (au sens ordinaire)
lorsqu'on suppose seulement que $\mathbf{K}$ admet des limites projectives
pour les systèmes projectifs \emph{finis}. En effet, si $U$ est un ouvert
quelconque de $X$, il y a un recouvrement \emph{fini} $(V_i)$ de $U$ formé
d'ensembles de $\mathfrak{B}$~; pour tout couple $(i,j)$ d'indices, soit
$(V_{ijk})$ un recouvrement fini de $V_i\cap V_j$ formé d'ensembles de
$\mathfrak{B}$. Soit $I$ l'ensemble formé des $i$ et des triplets $(i,j,k)$,
ordonné par les seules relations $i>(i,j,k)$, $j>(i,j,k)$~; on prend alors
pour $\mathcal{F}'(U)$ la limite projective du système des $\mathcal{F}(V_i)$
et $\mathcal{F}(V_{ijk})$~; on vérifie aisément que cela ne dépend pas des
recouvrements $(V_i)$ et $(V_{ijk})$ et que $U\to\mathcal{F}'(U)$ est un
préfaisceau.}.
\end{env}

\begin{env}[3.2.2]
\label{0.3.2.2-fr}
Pour que le préfaisceau $\mathcal{F}'$ ainsi défini soit un \emph{faisceau},
il faut et il suffit que le préfaisceau $\mathcal{F}$ sur $\mathfrak{B}$
vérifie la condition~:

\smallskip
\noindent
\textup{(F\textsubscript{0})} \emph{Pour tout recouvrement $(U_\alpha)$ de
$U\in\mathfrak{B}$ par des ensembles $U_\alpha\in\mathfrak{B}$ contenus dans
$U$, et pour tout objet $T\in\mathbf{K}$, l'application qui, à tout
$f\in\operatorname{Hom}(T,\mathcal{F}(U))$ fait correspondre la famille
\[
  (\rho_{U_\alpha}^{U}\circ f)
  \in\prod_\alpha\operatorname{Hom}(T,\mathcal{F}(U_\alpha))
\]
est une bijection de $\operatorname{Hom}(T,\mathcal{F}(U))$ sur l'ensemble des
$(f_\alpha)$ tels que
$\rho_V^{U_\alpha}\circ f_\alpha=\rho_V^{U_\beta}\circ f_\beta$ pour tout
couple d'indices $(\alpha,\beta)$ et tout $V\in\mathfrak{B}$ tel que
$V\subset U_\alpha\cap U_\beta$}%
\footnote{Cela signifie encore que le couple formé par $\mathcal{F}(U)$ et les
$\rho_\alpha=\rho_{U_\alpha}^{U}$ est \emph{solution du problème universel}
défini (3.1.1) par la donnée de
$A_\alpha=\mathcal{F}(U_\alpha)$,
$A_{\alpha\beta}=\prod\mathcal{F}(V)$ (pour les $V\in\mathfrak{B}$ tels que
$V\subset U_\alpha\cap U_\beta$) et
$\rho_{\alpha\beta}=(\rho_V''):\mathcal{F}(U_\alpha)\to
\prod\mathcal{F}(V)$ défini par la condition que pour
$V\in\mathfrak{B}$, $V'\in\mathfrak{B}$, $W\in\mathfrak{B}$,
$V\cup V'\subset U_\alpha\cap U_\beta$, $W\subset V\cap V'$,
$\rho_W^V\circ\rho_V''=\rho_W^{V'}\circ\rho_{V'}''$.}.

La condition est évidemment nécessaire. Pour montrer qu'elle est suffisante,
considérons d'abord une seconde base $\mathfrak{B}'$ de la topologie de $X$,
\emph{contenue dans $\mathfrak{B}$}, et montrons que si $\mathcal{F}''$
désigne le préfaisceau déduit de la sous-famille
$(\mathcal{F}(V))_{V\in\mathfrak{B}'}$, $\mathcal{F}''$ est
\emph{canoniquement isomorphe} à $\mathcal{F}'$. En effet, tout d'abord la
limite projective (pour $V\in\mathfrak{B}'$, $V\subset U$) des morphismes
canoniques $\mathcal{F}'(U)\to\mathcal{F}(V)$ est un morphisme
$\mathcal{F}'(U)\to\mathcal{F}''(U)$ pour tout ouvert $U$. Si
$U\in\mathfrak{B}$, ce morphisme est un isomorphisme, car par hypothèse les
morphismes canoniques $\mathcal{F}''(U)\to\mathcal{F}(V)$ pour
$V\in\mathfrak{B}'$, $V\subset U$, se factorisent en
\[
  \mathcal{F}''(U)\to\mathcal{F}(U)\to\mathcal{F}(V),
\]
et il est immédiat de voir que les composés des morphismes
$\mathcal{F}(U)\to\mathcal{F}''(U)$ et
$\mathcal{F}''(U)\to\mathcal{F}(U)$ ainsi définis sont les identités. Ceci
étant, pour tout ouvert $U$, les morphismes
$\mathcal{F}''(U)\to\mathcal{F}''(W)=\mathcal{F}(W)$ pour
$W\in\mathfrak{B}$ et $W\subset U$ vérifient les conditions caractérisant la
limite projective des $\mathcal{F}(W)$ ($W\in\mathfrak{B}$, $W\subset U$), ce
qui démontre notre assertion compte tenu de l'unicité d'une limite projective
à un isomorphisme près.

Cela posé, soit $U$ un ouvert quelconque de $X$, $(U_\alpha)$ un recouvrement
de $U$ par des ouverts contenus dans $U$, et soit $\mathfrak{B}'$ la
sous-famille de $\mathfrak{B}$ constituée par les ensembles

\oldpage[0\textsubscript{I}]{27}

de $\mathfrak{B}$ contenus dans un $U_\alpha$ au moins~; il est clair que
$\mathfrak{B}'$ est encore une base de la topologie de $X$, donc
$\mathcal{F}'(U)$ (resp. $\mathcal{F}''(U_\alpha)$) est limite projective des
$\mathcal{F}(V)$ pour $V\in\mathfrak{B}'$ et $V\subset U$ (resp.
$V\subset U_\alpha$)~; l'axiome (F) se vérifie alors aussitôt en vertu de la
définition de la limite projective.

Lorsque (F\textsubscript{0}) est vérifié, nous dirons par abus de langage que
le préfaisceau $\mathcal{F}$ sur la base $\mathfrak{B}$ est un \emph{faisceau}.
\end{env}

\begin{env}[3.2.3]
\label{0.3.2.3-fr}
Soient $\mathcal{F}$, $\mathcal{G}$ deux préfaisceaux sur la base
$\mathfrak{B}$, à valeurs dans $\mathbf{K}$~; on définit un \emph{morphisme}
$u:\mathcal{F}\to\mathcal{G}$ comme une famille
$(u_V)_{V\in\mathfrak{B}}$ de morphismes
$u_V:\mathcal{F}(V)\to\mathcal{G}(V)$ satisfaisant aux conditions de
compatibilité usuelles avec les morphismes de restriction $\rho_V^W$. Avec
les notations de (3.2.1), on en déduit un morphisme
$u':\mathcal{F}'\to\mathcal{G}'$ des préfaisceaux (ordinaires) correspondants
en prenant pour $u_U'$ la limite projective des $u_V$ pour
$V\in\mathfrak{B}$ et $V\subset U$~; la vérification des conditions de
compatibilité avec $\rho'_U{}^{U'}$ découle du caractère fonctoriel de la
limite projective.
\end{env}

\begin{env}[3.2.4]
\label{0.3.2.4-fr}
Si la catégorie $\mathbf{K}$ admet des limites inductives, et si $\mathcal{F}$
est un préfaisceau sur la base $\mathfrak{B}$, à valeurs dans $\mathbf{K}$,
pour tout $x\in X$ les voisinages de $x$ appartenant à $\mathfrak{B}$ forment
un ensemble cofinal (pour $\supset$) dans l'ensemble des voisinages de $x$,
donc, si $\mathcal{F}'$ est le préfaisceau (ordinaire) correspondant à
$\mathcal{F}$, la fibre $\mathcal{F}_x'$ est égale à
\[
  \varinjlim_{\mathfrak{B}}\mathcal{F}(V)
\]
selon l'ensemble des $V\in\mathfrak{B}$ contenant $x$.
Si $u:\mathcal{F}\to\mathcal{G}$ est un morphisme de préfaisceaux sur
$\mathfrak{B}$ à valeurs dans $\mathbf{K}$,
$u':\mathcal{F}'\to\mathcal{G}'$ le morphisme correspondant de préfaisceaux
ordinaires, $u_x'$ est de même la limite inductive des morphismes
$u_V:\mathcal{F}(V)\to\mathcal{G}(V)$ pour $V\in\mathfrak{B}$, $x\in V$.
\end{env}

\begin{env}[3.2.5]
\label{0.3.2.5-fr}
Revenons aux conditions générales de (3.2.1). Si $\mathcal{F}$ est un
\emph{faisceau} ordinaire à valeurs dans $\mathbf{K}$, $\mathcal{F}_1$ le
faisceau \emph{sur $\mathfrak{B}$} obtenu par restriction de $\mathcal{F}$ à
$\mathfrak{B}$, le faisceau ordinaire $\mathcal{F}_1'$ obtenu à partir de
$\mathcal{F}_1$ par le procédé de (3.2.1) est canoniquement isomorphe à
$\mathcal{F}$, en vertu de la condition (F) et des propriétés d'unicité de la
limite projective. On identifiera d'ordinaire $\mathcal{F}$ et
$\mathcal{F}_1'$.

Si $\mathcal{G}$ est un second faisceau (ordinaire) sur $X$ à valeurs dans
$\mathbf{K}$, et $u:\mathcal{F}\to\mathcal{G}$ un morphisme, la remarque
précédente montre que la donnée des
$u_V:\mathcal{F}(V)\to\mathcal{G}(V)$ \emph{pour les $V\in\mathfrak{B}$
seulement} détermine complètement $u$~; inversement, il suffit, les $u_V$
étant donnés pour $V\in\mathfrak{B}$, de vérifier le diagramme de commutativité
avec les morphismes de restriction $\rho_V^W$ pour $V\in\mathfrak{B}$,
$W\in\mathfrak{B}$ et $V\subset W$, pour qu'il existe un morphisme $u'$ et un
seul de $\mathcal{F}$ dans $\mathcal{G}$ tel que $u_V'=u_V$ pour tout
$V\in\mathfrak{B}$ (3.2.3).
\end{env}

\begin{env}[3.2.6]
\label{0.3.2.6-fr}
Supposons toujours que $\mathbf{K}$ admette des limites projectives. Alors la
catégorie des \emph{faisceaux sur $X$ à valeurs dans $\mathbf{K}$} admet aussi
des \emph{limites projectives}~; si $(\mathcal{F}_\lambda)$ est un système
projectif de faisceaux sur $X$ à valeurs dans $\mathbf{K}$, les
\[
  \mathcal{F}(U)=\varprojlim_\lambda\mathcal{F}_\lambda(U)
\]
définissent en effet un préfaisceau à valeurs dans $\mathbf{K}$, et la
vérification de l'axiome (F) résulte de la transitivité des limites
projectives~; le fait que $\mathcal{F}$ est alors limite projective des
$\mathcal{F}_\lambda$ est immédiat.

Lorsque $\mathbf{K}$ est la catégorie des ensembles, pour tout système
projectif $(\mathcal{H}_\lambda)$ tel

\oldpage[0\textsubscript{I}]{28}

que $\mathcal{H}_\lambda$ soit un \emph{sous-faisceau} de
$\mathcal{F}_\lambda$ pour tout $\lambda$,
$\varprojlim_\lambda\mathcal{H}_\lambda$ s'identifie canoniquement à un
\emph{sous-faisceau} de $\varprojlim_\lambda\mathcal{F}_\lambda$. Si
$\mathbf{K}$ est la catégorie des groupes commutatifs, le foncteur covariant
$\varprojlim_\lambda\mathcal{F}_\lambda$ est \emph{additif et exact à gauche}.
\end{env}

\subsection{Recollement de faisceaux}
\label{subsection:0.3.3-fr}

\begin{env}[3.3.1]
\label{0.3.3.1-fr}
Supposons encore que la catégorie $\mathbf{K}$ admette des limites projectives
(généralisées). Soient $X$ un espace topologique,
$\mathfrak{U}=(U_\lambda)_{\lambda\in L}$ un recouvrement ouvert de $X$, et
pour chaque $\lambda\in L$, soit $\mathcal{F}_\lambda$ un faisceau sur
$U_\lambda$, à valeurs dans $\mathbf{K}$~; pour tout couple d'indices
$(\lambda,\mu)$, supposons donné un \emph{isomorphisme}
\[
  \theta_{\lambda\mu}:
  \mathcal{F}_\mu|(U_\lambda\cap U_\mu)
  \xrightarrow{\sim}
  \mathcal{F}_\lambda|(U_\lambda\cap U_\mu);
\]
en outre, supposons que pour tout triplet $(\lambda,\mu,\nu)$, en désignant par
$\theta'_{\lambda\mu}$, $\theta'_{\mu\nu}$, $\theta'_{\lambda\nu}$ les
restrictions de $\theta_{\lambda\mu}$, $\theta_{\mu\nu}$,
$\theta_{\lambda\nu}$ à $U_\lambda\cap U_\mu\cap U_\nu$, on ait
\[
  \theta'_{\lambda\nu}
  =\theta'_{\lambda\mu}\circ\theta'_{\mu\nu}
\]
(\emph{condition de recollement pour les $\theta_{\lambda\mu}$}). Alors, il
existe un faisceau $\mathcal{F}$ sur $X$, à valeurs dans $\mathbf{K}$, et pour
chaque $\lambda$ un isomorphisme
\[
  \eta_\lambda:\mathcal{F}|U_\lambda\xrightarrow{\sim}\mathcal{F}_\lambda
\]
tels que, pour tout couple $(\lambda,\mu)$, en désignant par
$\eta'_\lambda$ et $\eta'_\mu$ les restrictions de $\eta_\lambda$ et
$\eta_\mu$ à $U_\lambda\cap U_\mu$, on ait
\[
  \theta_{\lambda\mu}=\eta'_\lambda\circ{\eta'_\mu}^{-1};
\]
en outre, $\mathcal{F}$ et les $\eta_\lambda$ sont déterminés à un
isomorphisme unique près par ces conditions. L'unicité résulte en effet
aussitôt de (3.2.5). Pour établir l'existence de $\mathcal{F}$, désignons par
$\mathfrak{B}$ la base d'ouverts formée des ouverts contenus dans un
$U_\lambda$ au moins, et pour tout $U\in\mathfrak{B}$, choisissons (par la
fonction $\tau$ de Hilbert) un des $\mathcal{F}_\lambda(U)$ pour un des
$\lambda$ tels que $U\subset U_\lambda$~; si on désigne cet objet par
$\mathcal{F}(U)$, les $\rho_U^V$ pour $U\subset V$, $U\in\mathfrak{B}$,
$V\in\mathfrak{B}$ se définissent de façon évidente (au moyen des
$\theta_{\lambda\mu}$), et la condition de transitivité est conséquence de la
condition de recollement~; en outre, la vérification de
(F\textsubscript{0}) est immédiate, donc le préfaisceau sur $\mathfrak{B}$
ainsi défini est bien un faisceau, et on en déduit par le procédé général
(3.2.1) un faisceau (ordinaire) encore noté $\mathcal{F}$ et qui répond à la
question. On dit que $\mathcal{F}$ est obtenu par \emph{recollement des
$\mathcal{F}_\lambda$ au moyen des $\theta_{\lambda\mu}$} et on identifiera
d'ordinaire $\mathcal{F}_\lambda$ et $\mathcal{F}|U_\lambda$ au moyen de
$\eta_\lambda$.

Il est clair que tout faisceau $\mathcal{F}$ sur $X$ à valeurs dans
$\mathbf{K}$ peut être considéré comme obtenu par recollement des faisceaux
$\mathcal{F}_\lambda=\mathcal{F}|U_\lambda$ (où $(U_\lambda)$ est un
recouvrement ouvert arbitraire de $X$), au moyen des isomorphismes
$\theta_{\lambda\mu}$ réduits à l'identité.
\end{env}

\begin{env}[3.3.2]
\label{0.3.3.2-fr}
Avec les mêmes notations, soit $\mathcal{G}_\lambda$ un second faisceau sur
$U_\lambda$ (pour tout $\lambda\in L$) à valeurs dans $\mathbf{K}$, et soit
donné pour tout couple $(\lambda,\mu)$ un isomorphisme
\[
  \omega_{\lambda\mu}:
  \mathcal{G}_\mu|(U_\lambda\cap U_\mu)
  \xrightarrow{\sim}
  \mathcal{G}_\lambda|(U_\lambda\cap U_\mu),
\]
ces isomorphismes vérifiant la condition de recollement. Supposons enfin donné
pour tout $\lambda$ un morphisme
$u_\lambda:\mathcal{F}_\lambda\to\mathcal{G}_\lambda$, et que les diagrammes
\begin{equation}
\begin{tikzcd}[column sep=huge,row sep=large]
\mathcal{F}_\mu|(U_\lambda\cap U_\mu)
  \arrow[r,"u_\mu"] \arrow[d]
  & \mathcal{G}_\mu|(U_\lambda\cap U_\mu) \arrow[d] \\
\mathcal{F}_\lambda|(U_\lambda\cap U_\mu)
  \arrow[r,"u_\lambda"']
  & \mathcal{G}_\lambda|(U_\lambda\cap U_\mu)
\end{tikzcd}
\tag{3.3.2.1}\label{0.3.3.2.1-fr}
\end{equation}
soient commutatifs. Alors, si $\mathcal{G}$ est obtenu par recollement des
$\mathcal{G}_\lambda$ au moyen des $\omega_{\lambda\mu}$, il existe un
morphisme $u:\mathcal{F}\to\mathcal{G}$ et un seul tel que les diagrammes

\oldpage[0\textsubscript{I}]{29}

\[
\begin{tikzcd}[column sep=huge,row sep=large]
\mathcal{F}|U_\lambda
  \arrow[r,"u|U_\lambda"] \arrow[d]
  & \mathcal{G}|U_\lambda \arrow[d] \\
\mathcal{F}_\lambda
  \arrow[r,"u_\lambda"']
  & \mathcal{G}_\lambda
\end{tikzcd}
\]
soient commutatifs~; cela résulte aussitôt de (3.2.3). La correspondance entre
la famille $(u_\lambda)$ et $u$ est une bijection fonctorielle de la partie de
$\prod_\lambda\operatorname{Hom}(\mathcal{F}_\lambda,\mathcal{G}_\lambda)$
vérifiant les conditions (3.3.2.1) sur
$\operatorname{Hom}(\mathcal{F},\mathcal{G})$.
\end{env}

\begin{env}[3.3.3]
\label{0.3.3.3-fr}
Avec les notations de (3.3.1), soit $V$ un ouvert de $X$~; il est immédiat que
les restrictions à $V\cap U_\lambda\cap U_\mu$ des $\theta_{\lambda\mu}$
satisfont à la condition de recollement pour les faisceaux induits
$\mathcal{F}_\lambda|(V\cap U_\lambda)$ et que le faisceau sur $V$ obtenu par
recollement de ces derniers s'identifie canoniquement à $\mathcal{F}|V$.
\end{env}

\subsection{Images directes de préfaisceaux}
\label{subsection:0.3.4-fr}

\begin{env}[3.4.1]
\label{0.3.4.1-fr}
Soient $X$, $Y$ deux espaces topologiques, $\psi:X\to Y$ une application
continue. Soit $\mathcal{F}$ un préfaisceau sur $X$ à valeurs dans une
catégorie $\mathbf{K}$~; pour tout ouvert $U\subset Y$, soit
\[
  \mathcal{G}(U)=\mathcal{F}(\psi^{-1}(U)),
\]
et si $U$, $V$ sont deux parties ouvertes de $Y$ telles que $U\subset V$, soit
$\rho_U^V$ le morphisme
\[
  \mathcal{F}(\psi^{-1}(V))\to\mathcal{F}(\psi^{-1}(U));
\]
il est immédiat que les $\mathcal{G}(U)$ et les $\rho_U^V$ définissent un
\emph{préfaisceau} sur $Y$ à valeurs dans $\mathbf{K}$, que l'on appelle
l'\emph{image directe de $\mathcal{F}$ par $\psi$} et que l'on note
$\psi_*(\mathcal{F})$. Si $\mathcal{F}$ est un faisceau, on vérifie aussitôt
l'axiome (F) pour le préfaisceau $\mathcal{G}(U)$, donc
$\psi_*(\mathcal{F})$ est un faisceau.
\end{env}

\begin{env}[3.4.2]
\label{0.3.4.2-fr}
Soient $\mathcal{F}_1$, $\mathcal{F}_2$ deux préfaisceaux sur $X$ à valeurs
dans $\mathbf{K}$, et soit $u:\mathcal{F}_1\to\mathcal{F}_2$ un morphisme.
Lorsque $U$ parcourt l'ensemble des parties ouvertes de $Y$, la famille de
morphismes
\[
  u_{\psi^{-1}(U)}:
  \mathcal{F}_1(\psi^{-1}(U))\to\mathcal{F}_2(\psi^{-1}(U))
\]
satisfait aux conditions de compatibilité avec les morphismes de restriction,
et définit par suite un morphisme
\[
  \psi_*(u):\psi_*(\mathcal{F}_1)\to\psi_*(\mathcal{F}_2).
\]
Si $v:\mathcal{F}_2\to\mathcal{F}_3$ est un morphisme de $\mathcal{F}_2$ dans
un troisième préfaisceau sur $X$ à valeurs dans $\mathbf{K}$, on a
\[
  \psi_*(v\circ u)=\psi_*(v)\circ\psi_*(u);
\]
autrement dit, $\psi_*(\mathcal{F})$ est un \emph{foncteur covariant} en
$\mathcal{F}$, de la catégorie des préfaisceaux (resp. faisceaux) sur $X$ à
valeurs dans $\mathbf{K}$, dans celle des préfaisceaux (resp. faisceaux) sur
$Y$ à valeurs dans $\mathbf{K}$.
\end{env}

\begin{env}[3.4.3]
\label{0.3.4.3-fr}
Soient $Z$ un troisième espace topologique, $\psi':Y\to Z$ une application
continue, et soit $\psi''=\psi'\circ\psi$. Il est clair que l'on a
\[
  \psi_*''(\mathcal{F})=\psi_*'(\psi_*(\mathcal{F}))
\]
pour tout préfaisceau $\mathcal{F}$ sur $X$ à valeurs dans $\mathbf{K}$~; en
outre, pour tout morphisme $u:\mathcal{F}\to\mathcal{G}$ de tels préfaisceaux,
on a
\[
  \psi_*''(u)=\psi_*'(\psi_*(u)).
\]
En d'autres termes, $\psi_*''$ est le \emph{composé} des foncteurs $\psi_*'$
et $\psi_*$, ce qu'on peut écrire
\[
  (\psi'\circ\psi)_*=\psi_*'\circ\psi_*.
\]

En outre, pour tout ouvert $U$ de $Y$, l'image par la restriction
$\psi|\psi^{-1}(U)$ du préfaisceau induit
$\mathcal{F}|\psi^{-1}(U)$ n'est autre que le préfaisceau induit
$\psi_*(\mathcal{F})|U$.
\end{env}

\begin{env}[3.4.4]
\label{0.3.4.4-fr}
Supposons que la catégorie $\mathbf{K}$ admette des limites inductives, et soit
$\mathcal{F}$ un préfaisceau sur $X$ à valeurs dans $\mathbf{K}$~; pour tout
$x\in X$, les morphismes
\[
  \Gamma(\psi^{-1}(U),\mathcal{F})\to\mathcal{F}_x
\]
($U$ voisinage ouvert de $\psi(x)$ dans $Y$) forment un système inductif, qui
donne par passage

\oldpage[0\textsubscript{I}]{30}

à la limite un morphisme
\[
  \psi_x:(\psi_*(\mathcal{F}))_{\psi(x)}\to\mathcal{F}_x
\]
des fibres~; en général, ce morphisme n'est \emph{ni injectif ni surjectif}.
Il est fonctoriel~; en effet, si
$u:\mathcal{F}_1\to\mathcal{F}_2$ est un morphisme de préfaisceaux sur $X$ à
valeurs dans $\mathbf{K}$, le diagramme
\[
\begin{tikzcd}[column sep=huge,row sep=large]
(\psi_*(\mathcal{F}_1))_{\psi(x)}
  \arrow[r,"\psi_x"] \arrow[d,"{(\psi_*(u))_{\psi(x)}}"']
  & (\mathcal{F}_1)_x \arrow[d,"u_x"] \\
(\psi_*(\mathcal{F}_2))_{\psi(x)}
  \arrow[r,"\psi_x"']
  & (\mathcal{F}_2)_x
\end{tikzcd}
\]
est commutatif. Si $Z$ est un troisième espace topologique,
$\psi':Y\to Z$ une application continue, et $\psi''=\psi'\circ\psi$, on a
\[
  \psi_x''=\psi_x\circ\psi'_{\psi(x)}
\]
pour $x\in X$.
\end{env}

\begin{env}[3.4.5]
\label{0.3.4.5-fr}
Sous les hypothèses de (3.4.4), supposons en outre que $\psi$ soit un
\emph{homéomorphisme} de $X$ sur le sous-espace $\psi(X)$ de $Y$. Alors, pour
tout $x\in X$, $\psi_x$ est un \emph{isomorphisme}. Ceci s'applique en
particulier à l'injection canonique $j$ d'une partie $X$ de $Y$ dans $Y$.
\end{env}

\begin{env}[3.4.6]
\label{0.3.4.6-fr}
Supposons que $\mathbf{K}$ soit la catégorie des groupes, ou des anneaux, etc.
Si $\mathcal{F}$ est un faisceau sur $X$ à valeurs dans $\mathbf{K}$, de
support $S$, et si $y\notin\overline{\psi(S)}$, il résulte de la définition de
$\psi_*(\mathcal{F})$ que
\[
  (\psi_*(\mathcal{F}))_y=\{0\},
\]
autrement dit le support de $\psi_*(\mathcal{F})$ est contenu dans
$\overline{\psi(S)}$~; mais il n'est pas nécessairement contenu dans
$\psi(S)$. Sous les mêmes hypothèses, si $j$ est l'injection canonique d'une
partie $X$ de $Y$ dans $Y$, le faisceau $j_*(\mathcal{F})$ induit
$\mathcal{F}$ sur $X$~; si de plus $X$ est \emph{fermée} dans $Y$,
$j_*(\mathcal{F})$ est le faisceau sur $Y$ qui induit $\mathcal{F}$ sur $X$
et $0$ sur $Y-X$ (G, II, 2.9.2), mais il est en général distinct de ce dernier
lorsqu'on suppose $X$ localement fermé mais non fermé.
\end{env}
 
\subsection{Images réciproques de préfaisceaux}

\begin{env}[3.5.1]
\label{0.3.5.1-fr}
Sous les hypothèses de (3.4.1), si $\mathcal{F}$ (resp. $\mathcal{G}$) est un
préfaisceau sur $X$ (resp. $Y$) à valeurs dans $\mathbf{K}$, tout morphisme
$u:\mathcal{G}\to\psi_*(\mathcal{F})$ de préfaisceaux sur $Y$ s'appelle encore
un $\psi$-\emph{morphisme} de $\mathcal{G}$ dans $\mathcal{F}$, et se note
aussi $\mathcal{G}\to\mathcal{F}$. On désigne aussi par
$\operatorname{Hom}_\psi(\mathcal{G},\mathcal{F})$ l'ensemble
$\operatorname{Hom}_Y(\mathcal{G},\psi_*(\mathcal{F}))$ des
$\psi$-morphismes de $\mathcal{G}$ dans $\mathcal{F}$.

Pour tout couple $(U,V)$, où $U$ est un ouvert de $X$, $V$ un ouvert de $Y$
tel que $\psi(U)\subset V$, on a un morphisme
$u_{U,V}:\mathcal{G}(V)\to\mathcal{F}(U)$ en composant le morphisme de
restriction
$\mathcal{F}(\psi^{-1}(V))\to\mathcal{F}(U)$ et le morphisme
$u_V:\mathcal{G}(V)\to\psi_*(\mathcal{F})(V)=\mathcal{F}(\psi^{-1}(V))$~; il
est immédiat que ces morphismes rendent commutatifs les diagrammes
\begin{equation*}
\begin{tikzcd}[column sep=huge,row sep=large]
\mathcal{G}(V) \arrow[r,"u_{U,V}"] \arrow[d]
  & \mathcal{F}(U) \arrow[d] \\
\mathcal{G}(V') \arrow[r,"u_{U',V'}"']
  & \mathcal{F}(U')
\end{tikzcd}
\tag{3.5.1.1}\label{0.3.5.1.1-fr}
\end{equation*}
pour $U'\subset U$, $V'\subset V$, $\psi(U')\subset V'$. Inversement, la
donnée d'une famille $(u_{U,V})$ de morphismes rendant commutatifs les
diagrammes (3.5.1.1) définit un $\psi$-morphisme $u$, car il suffit de prendre
\[
  u_V=u_{\psi^{-1}(V),V}.
\]

\oldpage[0\textsubscript{I}]{31}

Si la catégorie $\mathbf{K}$ admet des limites projectives (généralisées), et
si $\mathfrak{B}$, $\mathfrak{B}'$ sont des bases des topologies de $X$ et
$Y$ respectivement, pour définir un $\psi$-morphisme $u$ de \emph{faisceaux},
on peut se borner à se donner les $u_{U,V}$ pour
$U\in\mathfrak{B}$, $V\in\mathfrak{B}'$ et $\psi(U)\subset V$, vérifiant les
conditions de compatibilité (3.5.1.1) pour $U$, $U'$ dans $\mathfrak{B}$ et
$V$, $V'$ dans $\mathfrak{B}'$~; il suffit en effet de définir $u_W$, pour
tout ouvert $W\subset Y$, comme limite projective des $u_{U,V}$ pour
$V\in\mathfrak{B}'$ et $V\subset W$, $U\in\mathfrak{B}$ et
$\psi(U)\subset V$.

Lorsque la catégorie $\mathbf{K}$ admet des limites inductives, on a, pour
tout $x\in X$, un morphisme
\[
  \mathcal{G}(V)\to\mathcal{F}(\psi^{-1}(V))\to\mathcal{F}_x
\]
pour tout voisinage ouvert $V$ de $\psi(x)$ dans $Y$, et ces morphismes
forment un système inductif qui donne par passage à la limite un morphisme
\[
  \mathcal{G}_{\psi(x)}\to\mathcal{F}_x.
\]
\end{env}

\begin{env}[3.5.2]
\label{0.3.5.2-fr}
Sous les hypothèses de (3.4.3), soient $\mathcal{F}$, $\mathcal{G}$,
$\mathcal{H}$ des préfaisceaux à valeurs dans $\mathbf{K}$ sur $X$, $Y$, $Z$
respectivement, et soient
$u:\mathcal{G}\to\psi_*(\mathcal{F})$,
$v:\mathcal{H}\to\psi'_*(\mathcal{G})$ un $\psi$-morphisme et un
$\psi'$-morphisme respectivement. On en déduit un $\psi''$-morphisme
\[
  w:\mathcal{H}\xrightarrow{v}\psi'_*(\mathcal{G})
  \xrightarrow{\psi'_*(u)}
  \psi'_*(\psi_*(\mathcal{F}))=\psi''_*(\mathcal{F}),
\]
que l'on appelle, par définition, le \emph{composé} de $u$ et de $v$. On peut
donc considérer les couples $(X,\mathcal{F})$ formés d'un espace topologique
$X$ et d'un préfaisceau $\mathcal{F}$ sur $X$ (à valeurs dans $\mathbf{K}$)
comme formant une \emph{catégorie}, les morphismes étant les couples
$(\psi,\theta):(X,\mathcal{F})\to(Y,\mathcal{G})$ formés d'une application
continue $\psi:X\to Y$ et d'un $\psi$-morphisme
$\theta:\mathcal{G}\to\mathcal{F}$.
\end{env}

\begin{env}[3.5.3]
\label{0.3.5.3-fr}
Soient $\psi:X\to Y$ une application continue, $\mathcal{G}$ un
\emph{préfaisceau} sur $Y$ à valeurs dans $\mathbf{K}$. Nous appellerons
\emph{image réciproque de $\mathcal{G}$ par $\psi$} un couple
$(\mathcal{G}',\rho)$, où $\mathcal{G}'$ est un \emph{faisceau} sur $X$ à
valeurs dans $\mathbf{K}$, et
$\rho:\mathcal{G}\to\mathcal{G}'$ un $\psi$-morphisme (autrement dit un
homomorphisme $\mathcal{G}\to\psi_*(\mathcal{G}'))$ tels que, pour tout
\emph{faisceau} $\mathcal{F}$ sur $X$ à valeurs dans $\mathbf{K}$,
l'application
\begin{equation*}
  \operatorname{Hom}_X(\mathcal{G}',\mathcal{F})
  \longrightarrow \operatorname{Hom}_\psi(\mathcal{G},\mathcal{F})
  =\operatorname{Hom}_Y(\mathcal{G},\psi_*(\mathcal{F}))
\tag{3.5.3.1}\label{0.3.5.3.1-fr}
\end{equation*}
transformant $v$ en $\psi_*(v)\circ\rho$, soit une \emph{bijection}~; cette
application, étant fonctorielle en $\mathcal{F}$, définira alors un
isomorphisme de foncteurs en $\mathcal{F}$. Le couple
$(\mathcal{G}',\rho)$ étant solution d'un problème universel, on sait qu'il
est \emph{déterminé à un isomorphisme unique près} lorsqu'il existe. On écrira
alors $\mathcal{G}'=\psi^*(\mathcal{G})$, $\rho=\rho_{\mathcal{G}}$, et par
abus de langage, on dira que $\psi^*(\mathcal{G})$ est \emph{le faisceau image
réciproque} de $\mathcal{G}$ par $\psi$, étant entendu que
$\psi^*(\mathcal{G})$ est considéré comme muni du $\psi$-morphisme
\emph{canonique}
$\rho_{\mathcal{G}}:\mathcal{G}\to\psi^*(\mathcal{G})$, c'est-à-dire de
l'\emph{homomorphisme canonique} de préfaisceaux sur $Y$~:
\begin{equation*}
  \rho_{\mathcal{G}}:\mathcal{G}\to\psi_*(\psi^*(\mathcal{G})).
\tag{3.5.3.2}\label{0.3.5.3.2-fr}
\end{equation*}
Pour tout homomorphisme
$v:\psi^*(\mathcal{G})\to\mathcal{F}$ (où $\mathcal{F}$ est un faisceau sur
$X$ à valeurs dans $\mathbf{K}$), on posera
$v^b=\psi_*(v)\circ\rho_{\mathcal{G}}:
\mathcal{G}\to\psi_*(\mathcal{F})$. Par définition, \emph{tout} morphisme de
préfaisceaux $u:\mathcal{G}\to\psi_*(\mathcal{F})$ est de la forme $v^b$ pour
un $v$ et un seul, que l'on notera $u^\sharp$. En d'autres termes, tout
morphisme $u:\mathcal{G}\to\psi_*(\mathcal{F})$ de préfaisceaux se factorise
de façon unique en
\begin{equation*}
  u:\mathcal{G}\xrightarrow{\rho_{\mathcal{G}}}
  \psi_*(\psi^*(\mathcal{G}))
  \xrightarrow{\psi_*(u^\sharp)}\psi_*(\mathcal{F}).
\tag{3.5.3.3}\label{0.3.5.3.3-fr}
\end{equation*}
\end{env}

\oldpage[0\textsubscript{I}]{32}

\begin{env}[3.5.4]
\label{0.3.5.4-fr}
Supposons maintenant que la catégorie $\mathbf{K}$ soit telle\footnote{Dans
le livre cité dans l'Introduction, nous donnerons des conditions très
générales sur la catégorie $\mathbf{K}$ assurant l'existence des images
réciproques de préfaisceaux à valeurs dans $\mathbf{K}$.} que \emph{tout}
préfaisceau $\mathcal{G}$ sur $Y$ à valeurs dans $\mathbf{K}$ admette une
image réciproque par $\psi$, que nous noterons $\psi^*(\mathcal{G})$.

Nous allons voir qu'on peut définir $\psi^*(\mathcal{G})$ comme \emph{foncteur
covariant} en $\mathcal{G}$, de la catégorie des préfaisceaux sur $Y$ à
valeurs dans $\mathbf{K}$, dans celle des faisceaux sur $X$ à valeurs dans
$\mathbf{K}$, de telle façon que l'isomorphisme $v\to v^b$ soit un
\emph{isomorphisme de bifoncteurs}
\begin{equation*}
  \operatorname{Hom}_X(\psi^*(\mathcal{G}),\mathcal{F})
  \xrightarrow{\sim}
  \operatorname{Hom}_Y(\mathcal{G},\psi_*(\mathcal{F}))
\tag{3.5.4.1}\label{0.3.5.4.1-fr}
\end{equation*}
en $\mathcal{G}$ et $\mathcal{F}$.

En effet, pour tout morphisme
$w:\mathcal{G}_1\to\mathcal{G}_2$ de préfaisceaux sur $Y$ à valeurs dans
$\mathbf{K}$, considérons le morphisme composé
\[
  \mathcal{G}_1\xrightarrow{w}\mathcal{G}_2
  \xrightarrow{\rho_{\mathcal{G}_2}}
  \psi_*(\psi^*(\mathcal{G}_2))~;
\]
il lui correspond un morphisme
$(\rho_{\mathcal{G}_2}\circ w)^\sharp:
\psi^*(\mathcal{G}_1)\to\psi^*(\mathcal{G}_2)$, que nous noterons
$\psi^*(w)$. On a donc, en vertu de (3.5.3.3)
\begin{equation*}
  \psi_*(\psi^*(w))\circ\rho_{\mathcal{G}_1}
  =\rho_{\mathcal{G}_2}\circ w.
\tag{3.5.4.2}\label{0.3.5.4.2-fr}
\end{equation*}
Pour tout morphisme $u:\mathcal{G}_2\to\psi_*(\mathcal{F})$, où
$\mathcal{F}$ est un faisceau sur $X$ à valeurs dans $\mathbf{K}$, on a,
d'après (3.5.3.3), (3.5.4.2) et la définition de $v^b$
\[
  (u^\sharp\circ\psi^*(w))^b
  =\psi_*(u^\sharp)\circ\psi_*(\psi^*(w))\circ\rho_{\mathcal{G}_1}
  =\psi_*(u^\sharp)\circ\rho_{\mathcal{G}_2}\circ w
  =u\circ w
\]
ou encore
\begin{equation*}
  (u\circ w)^\sharp=u^\sharp\circ\psi^*(w).
\tag{3.5.4.3}\label{0.3.5.4.3-fr}
\end{equation*}
Si on prend en particulier pour $u$ un morphisme
\[
  \mathcal{G}_2\xrightarrow{w'}\mathcal{G}_3
  \xrightarrow{\rho_{\mathcal{G}_3}}
  \psi_*(\psi^*(\mathcal{G}_3)),
\]
il vient
\[
  \psi^*(w'\circ w)
  =(\rho_{\mathcal{G}_3}\circ w'\circ w)^\sharp
  =(\rho_{\mathcal{G}_3}\circ w')^\sharp\circ\psi^*(w)
  =\psi^*(w')\circ\psi^*(w),
\]
d'où notre assertion.

Enfin, pour tout faisceau $\mathcal{F}$ sur $X$ à valeurs dans
$\mathbf{K}$, soit $i_{\mathcal{F}}$ le morphisme identique de
$\psi_*(\mathcal{F})$ et notons
\[
  \sigma_{\mathcal{F}}:
  \psi^*(\psi_*(\mathcal{F}))\to\mathcal{F}
\]
le morphisme $(i_{\mathcal{F}})^\sharp$~; la formule (3.5.4.3) donne en
particulier la factorisation
\begin{equation*}
  u^\sharp:\psi^*(\mathcal{G})
  \xrightarrow{\psi^*(u)}
  \psi^*(\psi_*(\mathcal{F}))
  \xrightarrow{\sigma_{\mathcal{F}}}\mathcal{F}
\tag{3.5.4.4}\label{0.3.5.4.4-fr}
\end{equation*}
pour tout morphisme $u:\mathcal{G}\to\psi_*(\mathcal{F})$. Nous dirons que le
morphisme $\sigma_{\mathcal{F}}$ est \emph{canonique}.
\end{env}

\begin{env}[3.5.5]
\label{0.3.5.5-fr}
Soit $\psi':Y\to Z$ une application continue, et supposons que tout
préfaisceau $\mathcal{H}$ sur $Z$ à valeurs dans $\mathbf{K}$ admette une
image réciproque $\psi'^*(\mathcal{H})$ par $\psi'$. Alors (avec les
hypothèses de (3.5.4)) tout préfaisceau $\mathcal{H}$ sur $Z$ à valeurs dans
$\mathbf{K}$ admet une image réciproque par $\psi''=\psi'\circ\psi$ et l'on a
un isomorphisme canonique fonctoriel
\begin{equation*}
  \psi''^*(\mathcal{H})\xrightarrow{\sim}
  \psi^*(\psi'^*(\mathcal{H}))
\tag{3.5.5.1}\label{0.3.5.5.1-fr}
\end{equation*}

\oldpage[0\textsubscript{I}]{33}

Cela résulte en effet aussitôt des définitions, tenant compte de ce que
$\psi''_*=\psi'_*\circ\psi_*$. En outre, si
$u:\mathcal{G}\to\psi_*(\mathcal{F})$ est un $\psi$-morphisme,
$v:\mathcal{H}\to\psi'_*(\mathcal{G})$ un $\psi'$-morphisme, et
$w=\psi'_*(u)\circ v$ leur composé (3.5.2), on constate aussitôt que
$w^\sharp$ est le morphisme composé
\[
  w^\sharp:\psi^*(\psi'^*(\mathcal{H}))
  \xrightarrow{\psi^*(v^\sharp)}
  \psi^*(\mathcal{G})
  \xrightarrow{u^\sharp}\mathcal{F}.
\]
\end{env}

\begin{env}[3.5.6]
\label{0.3.5.6-fr}
Prenons en particulier pour $\psi$ l'application identique
$1_X:X\to X$. Alors si l'image réciproque par $\psi$ d'un préfaisceau
$\mathcal{F}$ sur $X$ à valeurs dans $\mathbf{K}$ existe, on dit que cette
image réciproque est le \emph{faisceau associé au préfaisceau $\mathcal{F}$}.
Tout morphisme $u:\mathcal{F}\to\mathcal{F}'$ de $\mathcal{F}$ dans un
\emph{faisceau} $\mathcal{F}'$ à valeurs dans $\mathbf{K}$ se factorise donc
de façon unique en
\[
  \mathcal{F}\xrightarrow{\rho_{\mathcal{F}}}
  1_X^*(\mathcal{F})\xrightarrow{u^\sharp}\mathcal{F}'.
\]
\end{env}

\subsection{Faisceaux simples et faisceaux localement simples}

\begin{env}[3.6.1]
\label{0.3.6.1-fr}
Nous dirons qu'un \emph{préfaisceau} $\mathcal{F}$ sur $X$, à valeurs dans
$\mathbf{K}$, est \emph{constant} si les morphismes canoniques
$\mathcal{F}(X)\to\mathcal{F}(U)$ sont des \emph{isomorphismes} pour tout
ouvert non vide $U\subset X$~; on notera que $\mathcal{F}$ n'est pas
nécessairement un faisceau. On dit qu'un \emph{faisceau} est \emph{simple}
s'il est associé (3.5.6) à un préfaisceau constant. On dit qu'un faisceau
$\mathcal{F}$ est \emph{localement simple} si tout $x\in X$ admet un voisinage
ouvert $U$ tel que $\mathcal{F}|U$ soit simple.
\end{env}

\begin{env}[3.6.2]
\label{0.3.6.2-fr}
Supposons que $X$ soit \emph{irréductible} (2.1.1)~; alors les propriétés
suivantes sont équivalentes~:
\begin{enumerate}
  \item[a)] \emph{$\mathcal{F}$ est un préfaisceau constant sur $X$}~;
  \item[b)] \emph{$\mathcal{F}$ est un faisceau simple sur $X$}~;
  \item[c)] \emph{$\mathcal{F}$ est un faisceau localement simple sur $X$}.
\end{enumerate}

En effet, soit $\mathcal{F}$ un préfaisceau constant sur $X$~; si $U$, $V$
sont deux ouverts non vides dans $X$, $U\cap V$ est non vide, donc
$\mathcal{F}(X)\to\mathcal{F}(U)\to\mathcal{F}(U\cap V)$ et
$\mathcal{F}(X)\to\mathcal{F}(U)$ étant des isomorphismes, il en est de même
de $\mathcal{F}(U)\to\mathcal{F}(U\cap V)$ et de même
$\mathcal{F}(V)\to\mathcal{F}(U\cap V)$ est un isomorphisme. On en conclut
aussitôt que l'axiome (F) de (3.1.2) est bien vérifié, $\mathcal{F}$ est
isomorphe à son faisceau associé, et par suite a) entraîne b).

Soit maintenant $(U_\alpha)$ un recouvrement ouvert de $X$ par des ouverts
non vides et un faisceau $\mathcal{F}$ sur $X$ tel que
$\mathcal{F}|U_\alpha$ soit simple pour tout $\alpha$~; comme $U_\alpha$ est
irréductible, $\mathcal{F}|U_\alpha$ est un préfaisceau constant en vertu de
ce qui précède. Comme $U_\alpha\cap U_\beta$ n'est pas vide,
$\mathcal{F}(U_\alpha)\to\mathcal{F}(U_\alpha\cap U_\beta)$ et
$\mathcal{F}(U_\beta)\to\mathcal{F}(U_\alpha\cap U_\beta)$ sont des
isomorphismes, d'où on déduit un isomorphisme canonique
$\theta_{\alpha\beta}:\mathcal{F}(U_\alpha)\to\mathcal{F}(U_\beta)$ pour tout
couple d'indices. Mais alors, si on applique la condition (F) pour $U=X$, on
voit que pour tout indice $\alpha_0$, $\mathcal{F}(U_{\alpha_0})$ et les
$\theta_{\alpha_0\alpha}$ sont solutions du problème universel, ce qui (en
vertu de l'unicité) implique que
$\mathcal{F}(X)\to\mathcal{F}(U_{\alpha_0})$ est un isomorphisme, et démontre
donc que c) entraîne a).
\end{env}

\subsection{Images réciproques de préfaisceaux de groupes ou d'anneaux}

\oldpage[0\textsubscript{I}]{34}

\begin{env}[3.7.1]
\label{0.3.7.1-fr}
Nous allons montrer que lorsqu'on prend pour $\mathbf{K}$ la catégorie des
ensembles, l'image réciproque par $\psi$ de tout préfaisceau $\mathcal{G}$ à
valeurs dans $\mathbf{K}$ \emph{existe toujours} (les notations et hypothèses
sur $X$, $Y$, $\psi$ étant celles de (3.5.3)). En effet, pour tout ouvert
$U\subset X$, définissons $\mathcal{G}'(U)$ comme suit~: un élément $s'$ de
$\mathcal{G}'(U)$ est une famille $(s'_x)_{x\in U}$, où
$s'_x\in\mathcal{G}_{\psi(x)}$ pour tout $x\in U$, et où, pour tout $x\in U$,
la condition suivante est remplie~: il existe un voisinage ouvert $V$ de
$\psi(x)$ dans $Y$, un voisinage $W\subset\psi^{-1}(V)\cap U$ de $x$ et un
élément $s\in\mathcal{G}(V)$ tels que $s'_z=s_{\psi(z)}$ pour tout $z\in W$.
On vérifie immédiatement que $U\mapsto\mathcal{G}'(U)$ satisfait bien aux axiomes
des \emph{faisceaux}.

Soit maintenant $\mathcal{F}$ un faisceau d'ensembles sur $X$, et soient
$u:\mathcal{G}\to\psi_*(\mathcal{F})$,
$v:\mathcal{G}'\to\mathcal{F}$ des morphismes. On définit $u^\sharp$ et $v^b$
de la façon suivante~: si $s'$ est une section de $\mathcal{G}'$ au-dessus
d'un voisinage $U$ de $x\in X$ et si $V$ est un voisinage ouvert de $\psi(x)$
et $s\in\mathcal{G}(V)$ tel que l'on ait $s'_z=s_{\psi(z)}$ pour $z$ dans un
voisinage de $x$ contenu dans $\psi^{-1}(V)\cap U$, on prend
$u_x^\sharp(s'_x)=u_{\psi(x)}(s_{\psi(x)})$. De même, si
$s\in\mathcal{G}(V)$ ($V$ ouvert dans $Y$), $v^b(s)$ est la section de
$\mathcal{F}$ au-dessus de $\psi^{-1}(V)$, image par $v$ de la section $s'$ de
$\mathcal{G}'$ telle que $s'_x=s_{\psi(x)}$ pour tout
$x\in\psi^{-1}(V)$. En outre, l'homomorphisme canonique (3.5.3)
$\rho:\mathcal{G}\to\psi_*(\psi^*(\mathcal{G}))$ se définit de la façon
suivante~: pour tout ouvert $V\subset Y$ et toute section
$s\in\Gamma(V,\mathcal{G})$, $\rho(s)$ est la section
$(s_{\psi(x)})_{x\in\psi^{-1}(V)}$ de $\psi^*(\mathcal{G})$ au-dessus de
$\psi^{-1}(V)$. La vérification des relations
$(u^\sharp)^b=u$, $(v^b)^\sharp=v$ et
$v^b=\psi_*(v)\circ\rho$ est immédiate, et démontre notre assertion.

On vérifie que, si $w:\mathcal{G}_1\to\mathcal{G}_2$ est un homomorphisme de
préfaisceaux d'ensembles sur $Y$, $\psi^*(w)$ s'explicite de la façon
suivante~: si $s'=(s'_x)_{x\in U}$ est une section de
$\psi^*(\mathcal{G}_1)$ au-dessus d'un ouvert $U$ de $X$,
$(\psi^*(w))(s')$ est la famille
$(w_{\psi(x)}(s'_x))_{x\in U}$. Enfin, il est immédiat que pour tout ouvert
$V$ de $Y$, l'image réciproque de $\mathcal{G}|V$ par la restriction de
$\psi$ à $\psi^{-1}(V)$ est identique au faisceau induit
$\psi^*(\mathcal{G})|\psi^{-1}(V)$.

Lorsque $\psi$ est l'identité $1_X$, on retrouve la définition d'un faisceau
d'ensembles associé à un préfaisceau (G, II, 1.2). Les considérations
précédentes s'appliquent sans changement lorsque $\mathbf{K}$ est la
catégorie des groupes ou des anneaux (non nécessairement commutatifs).

Lorsque $X$ est une partie quelconque d'un espace topologique $Y$, et $j$
l'injection canonique $X\to Y$, pour tout faisceau $\mathcal{G}$ sur $Y$ à
valeurs dans une catégorie $\mathbf{K}$, on appelle faisceau \emph{induit} sur
$X$ par $\mathcal{G}$ l'image réciproque $j^*(\mathcal{G})$ (lorsqu'elle
existe)~; pour les faisceaux d'ensembles (ou de groupes, ou d'anneaux) on
retrouve la définition usuelle (G, II, 1.5).
\end{env}

\begin{env}[3.7.2]
\label{0.3.7.2-fr}
Conservant les notations et hypothèses de (3.5.3), supposons que
$\mathcal{G}$ soit un \emph{faisceau} de groupes (resp. d'anneaux) sur $Y$.
La définition des sections de $\psi^*(\mathcal{G})$ (3.7.1) montre (compte
tenu de (3.4.4)) que l'homomorphisme de fibres
$\psi_x\circ\rho_{\psi(x)}:\mathcal{G}_{\psi(x)}\to
(\psi^*(\mathcal{G}))_x$ est un \emph{isomorphisme fonctoriel en
$\mathcal{G}$}, qui permet d'identifier ces deux fibres~; avec cette
identification, $u_x^\sharp$ est identique à l'homomorphisme défini dans
(3.5.1), et en particulier, on a
$\operatorname{Supp}(\psi^*(\mathcal{G}))=
\psi^{-1}(\operatorname{Supp}(\mathcal{G}))$.

Une conséquence immédiate de ce résultat est que \emph{le foncteur
$\psi^*(\mathcal{G})$ est exact en $\mathcal{G}$} dans la catégorie abélienne
des faisceaux de groupes commutatifs.
\end{env}

\subsection{Faisceaux d'espaces pseudo-discrets}

\oldpage[0\textsubscript{I}]{35}

\begin{env}[3.8.1]
\label{0.3.8.1-fr}
Soit $X$ un espace topologique dont la topologie admet une base
$\mathfrak{B}$ formée d'ensembles ouverts \emph{quasi-compacts}. Soit
$\mathcal{F}$ un \emph{faisceau d'ensembles} sur $X$~; si on munit chacun des
$\mathcal{F}(U)$ de la topologie \emph{discrète},
$U\mapsto\mathcal{F}(U)$ est un \emph{préfaisceau d'espaces topologiques}.
Nous allons voir qu'il existe un \emph{faisceau d'espaces topologiques
$\mathcal{F}'$ associé à $\mathcal{F}$} (3.5.6) tel que
$\Gamma(U,\mathcal{F}')$ soit l'espace discret $\mathcal{F}(U)$ pour tout
ouvert \emph{quasi-compact} $U$. Il suffira pour cela de montrer que le
préfaisceau $U\mapsto\mathcal{F}(U)$ d'espaces topologiques discrets
\emph{sur $\mathfrak{B}$} vérifie la condition $(\mathrm{F}_0)$ de (3.2.2),
et plus généralement que si $U$ est un ouvert quasi-compact et si
$(U_\alpha)$ est un recouvrement de $U$ par des ensembles de $\mathfrak{B}$,
la topologie la moins fine $\mathcal{T}$ sur $\Gamma(U,\mathcal{F})$ rendant
continues les applications
$\Gamma(U,\mathcal{F})\to\Gamma(U_\alpha,\mathcal{F})$ est la topologie
\emph{discrète}. Or, il existe un nombre fini d'indices $\alpha_i$ tels que
$U=\bigcup_i U_{\alpha_i}$. Soit $s\in\Gamma(U,\mathcal{F})$ et soit $s_i$ son
image dans $\Gamma(U_{\alpha_i},\mathcal{F})$~; l'intersection des images
réciproques des ensembles $\{s_i\}$ est par définition un voisinage de $s$
pour $\mathcal{T}$~; mais puisque $\mathcal{F}$ est un faisceau d'ensembles et
que les $U_{\alpha_i}$ recouvrent $U$, cette intersection se réduit à $s$,
d'où notre assertion.

On notera que si $U$ est un ouvert non quasi-compact de $X$, l'espace
topologique $\Gamma(U,\mathcal{F}')$ a encore $\Gamma(U,\mathcal{F})$ comme
ensemble sous-jacent, mais sa topologie n'est pas discrète en général~:
c'est la moins fine rendant continues les applications
$\Gamma(U,\mathcal{F})\to\Gamma(V,\mathcal{F})$, pour $V\in\mathfrak{B}$ et
$V\subset U$ (les $\Gamma(V,\mathcal{F})$ étant discrets).

Les considérations précédentes s'appliquent sans modification aux faisceaux
de groupes ou d'anneaux (non nécessairement commutatifs), et leur associent
respectivement des faisceaux de \emph{groupes topologiques} ou d'\emph{anneaux
topologiques}. Pour abréger, nous dirons que le faisceau $\mathcal{F}'$ est le
faisceau d'\emph{espaces} (resp. \emph{groupes}, \emph{anneaux})
\emph{pseudo-discrets} associé au faisceau d'ensembles (resp. groupes,
anneaux) $\mathcal{F}$.
\end{env}

\begin{env}[3.8.2]
\label{0.3.8.2-fr}
Soient $\mathcal{F}$, $\mathcal{G}$ deux faisceaux d'ensembles (resp. groupes,
anneaux) sur $X$, $u:\mathcal{F}\to\mathcal{G}$ un homomorphisme. Alors $u$
est aussi un homomorphisme \emph{continu} $\mathcal{F}'\to\mathcal{G}'$, en
désignant par $\mathcal{F}'$ et $\mathcal{G}'$ les faisceaux pseudo-discrets
associés à $\mathcal{F}$ et $\mathcal{G}$~; cela résulte en effet de
(3.2.5).
\end{env}

\begin{env}[3.8.3]
\label{0.3.8.3-fr}
Soient $\mathcal{F}$ un faisceau d'ensembles, $\mathcal{H}$ un sous-faisceau de
$\mathcal{F}$, $\mathcal{F}'$ et $\mathcal{H}'$ les faisceaux pseudo-discrets
associés à $\mathcal{F}$ et $\mathcal{H}$ respectivement. Alors, pour tout
ouvert $U\subset X$, $\Gamma(U,\mathcal{H}')$ est \emph{fermé} dans
$\Gamma(U,\mathcal{F}')$~: en effet, c'est l'intersection des images
réciproques des $\Gamma(V,\mathcal{H})$ (pour $V\in\mathfrak{B}$,
$V\subset U$) par les applications continues
$\Gamma(U,\mathcal{F})\to\Gamma(V,\mathcal{F})$, et
$\Gamma(V,\mathcal{H})$ est fermé dans l'espace discret
$\Gamma(V,\mathcal{F})$.
\end{env}

\section{Espaces annelés}
\label{section:0.4-fr}

\subsection{Espaces annelés, $\mathcal{A}$-Modules, $\mathcal{A}$-Algèbres}
\label{subsection:0.4.1-fr}

\begin{env}[4.1.1]
\label{0.4.1.1-fr}
Un \emph{espace annelé} (resp. \emph{topologiquement annelé}) est un couple
$(X,\mathcal{A})$ formé d'un espace topologique $X$ et d'un faisceau d'anneaux
(non nécessairement commutatifs) (resp. d'un faisceau d'anneaux topologiques)
$\mathcal{A}$ sur $X$~; on dit que $X$ est l'espace
\oldpage[0\textsubscript{I}]{36}
topologique \emph{sous-jacent} à l'espace annelé $(X,\mathcal{A})$, et
$\mathcal{A}$ le \emph{faisceau structural}. Ce dernier se note
$\mathcal{O}_X$, et sa fibre en un point $x\in X$ se note $\mathcal{O}_{X,x}$
ou simplement $\mathcal{O}_x$ lorsqu'il n'en résulte pas de confusion.

On désignera par $1$ ou $e$ la \emph{section unité} de $\mathcal{O}_X$
au-dessus de $X$ (élément unité de $\Gamma(X,\mathcal{O}_X)$).

Comme dans ce Traité nous aurons surtout à considérer des faisceaux d'anneaux
\emph{commutatifs}, il sera sous-entendu, lorsque nous parlerons d'un espace
annelé $(X,\mathcal{A})$ sans préciser, que $\mathcal{A}$ est un faisceau
d'anneaux commutatifs.

Les espaces annelés à faisceau structural non nécessairement commutatif
(resp. les espaces topologiquement annelés) forment une \emph{catégorie},
lorsqu'on définit un \emph{morphisme}
$(X,\mathcal{A})\to(Y,\mathcal{B})$ comme un couple
$(\psi,\theta)=\Psi$ formé d'une application continue $\psi:X\to Y$ et d'un
$\psi$-\emph{morphisme} $\theta:\mathcal{G}\to\mathcal{F}$ (3.5.1) de
faisceaux d'anneaux (resp. de faisceaux d'anneaux topologiques)~; le
\emph{composé} d'un second morphisme
$\Psi'=(\psi',\theta'):(Y,\mathcal{B})\to(Z,\mathcal{C})$ et de $\Psi$, noté
$\Psi''=\Psi'\circ\Psi$, est le morphisme $(\psi'',\theta'')$ où
$\psi''=\psi'\circ\psi$, et $\theta''$ est le composé de $\theta$ et $\theta'$
(égal à $\psi'_*(\theta)\circ\theta'$, cf. 3.5.2). Pour les espaces
annelés, rappelons que l'on a alors
$\theta''^\sharp=\theta^\sharp\circ\psi^*(\theta'^\sharp)$ (3.5.5)~; donc si
$\theta'^\sharp$ et $\theta^\sharp$ sont des homomorphismes \emph{injectifs}
(resp. \emph{surjectifs}), il en est de même de $\theta''^\sharp$, compte tenu
du fait que $\psi_x\circ\rho_{\psi(x)}$ est un isomorphisme pour tout $x\in X$
(3.7.2). On vérifie aussitôt, grâce à ce qui précède, que lorsque $\psi$
est une application continue \emph{injective} et $\theta^\sharp$ un
homomorphisme \emph{surjectif} de faisceaux d'anneaux, le morphisme
$(\psi,\theta)$ est un \emph{monomorphisme} (T, 1.1) dans la catégorie des
espaces annelés.

Par abus de langage, on remplacera souvent $\psi$ par $\Psi$ dans les
notations, par exemple en écrivant $\Psi^{-1}(U)$ au lieu de $\psi^{-1}(U)$
pour une partie $U$ de $Y$, lorsque cela ne risquera pas d'entraîner confusion.
\end{env}

\begin{env}[4.1.2]
\label{0.4.1.2-fr}
Pour toute partie $M$ de $X$, le couple $(M,\mathcal{A}|M)$ est évidemment un
espace annelé, dit \emph{induit} sur $M$ par l'espace annelé
$(X,\mathcal{A})$ (et appelé encore la \emph{restriction} de
$(X,\mathcal{A})$ à $M$). Si $j$ est l'injection canonique $M\to X$ et
$\omega$ l'application identique de $\mathcal{A}|M$, $(j,\omega^b)$ est un
monomorphisme $(M,\mathcal{A}|M)\to(X,\mathcal{A})$ d'espaces annelés, appelé
l'\emph{injection canonique}. Le composé d'un morphisme
$\Psi:(X,\mathcal{A})\to(Y,\mathcal{B})$ et de cette injection est appelé la
\emph{restriction} de $\Psi$ à $M$.
\end{env}

\begin{env}[4.1.3]
\label{0.4.1.3-fr}
Nous ne reviendrons pas sur la définition des $\mathcal{A}$-Modules ou
faisceaux algébriques sur un espace annelé $(X,\mathcal{A})$
(G, II, 2.2)~; lorsque $\mathcal{A}$ est un faisceau d'anneaux non
nécessairement commutatifs, par $\mathcal{A}$-Module, il faudra toujours
sous-entendre \og $\mathcal{A}$-Module à gauche\fg{} sauf mention expresse du
contraire. Les sous-$\mathcal{A}$-Modules de $\mathcal{A}$ seront qualifiés de
faisceaux d'idéaux (à gauche, à droite ou bilatères) dans $\mathcal{A}$ ou de
$\mathcal{A}$-Idéaux.

Lorsque $\mathcal{A}$ est un faisceau d'anneaux commutatifs, et que, dans la
définition des $\mathcal{A}$-Modules, on remplace partout la structure de
module par celle d'algèbre, on obtient la définition d'une
$\mathcal{A}$-Algèbre sur $X$. Il revient au même de dire qu'une
$\mathcal{A}$-Algèbre (non nécessairement commutative) est un
$\mathcal{A}$-Module $\mathcal{C}$, muni d'un homomorphisme de
$\mathcal{A}$-Modules
$\varphi:\mathcal{C}\otimes_{\mathcal{A}}\mathcal{C}\to\mathcal{C}$ et
d'une section $e$ au-dessus de $X$, tels que~: 1\textsuperscript{o} Le diagramme
\oldpage[0\textsubscript{I}]{37}
\[
\begin{tikzcd}[column sep=large,row sep=large]
\mathcal{C}\otimes_{\mathcal{A}}\mathcal{C}\otimes_{\mathcal{A}}\mathcal{C}
  \arrow[r,"\varphi\otimes 1"]
  \arrow[d,"1\otimes\varphi"']
&
\mathcal{C}\otimes_{\mathcal{A}}\mathcal{C}
  \arrow[d,"\varphi"]
\\
\mathcal{C}\otimes_{\mathcal{A}}\mathcal{C}
  \arrow[r,"\varphi"']
&
\mathcal{C}
\end{tikzcd}
\]
soit commutatif~; 2\textsuperscript{o} pour tout ouvert $U\subset X$ et toute section
$s\in\Gamma(U,\mathcal{C})$, on ait
$\varphi((e|U)\otimes s)=\varphi(s\otimes(e|U))=s$. Dire que
$\mathcal{C}$ est une $\mathcal{A}$-Algèbre commutative revient en outre à
dire que le diagramme
\[
\begin{tikzcd}[column sep=large,row sep=large]
\mathcal{C}\otimes_{\mathcal{A}}\mathcal{C}
  \arrow[rr,"\sigma"]
  \arrow[dr,"\varphi"']
&&
\mathcal{C}\otimes_{\mathcal{A}}\mathcal{C}
  \arrow[dl,"\varphi"]
\\
& \mathcal{C} &
\end{tikzcd}
\]
est commutatif, $\sigma$ désignant la symétrie canonique du produit tensoriel
$\mathcal{C}\otimes_{\mathcal{A}}\mathcal{C}$.

Les homomorphismes de $\mathcal{A}$-Algèbres se définissent aussi comme les
homomorphismes de $\mathcal{A}$-Modules dans (G, II, 2.2), mais naturellement
ne forment plus un groupe abélien.

Si $\mathcal{M}$ est un sous-$\mathcal{A}$-Module d'une
$\mathcal{A}$-Algèbre $\mathcal{C}$, la sous-$\mathcal{A}$-Algèbre de
$\mathcal{C}$ engendrée par $\mathcal{M}$ est la somme des images des
homomorphismes $\bigotimes^n\mathcal{M}\to\mathcal{C}$ (pour les
$n\geqslant 0$). C'est aussi le faisceau associé au préfaisceau
$U\mapsto\mathcal{B}(U)$ d'algèbres, $\mathcal{B}(U)$ étant la sous-algèbre
de $\Gamma(U,\mathcal{C})$ engendrée par le sous-module
$\Gamma(U,\mathcal{M})$.
\end{env}

\begin{env}[4.1.4]
\label{0.4.1.4-fr}
On dit qu'un faisceau d'anneaux $\mathcal{A}$ sur un espace topologique $X$
est \emph{réduit en un point} $x$ de $X$ si la fibre $\mathcal{A}_x$ est un
anneau \emph{réduit} (1.1.1)~; on dit que $\mathcal{A}$ est \emph{réduit}
s'il est réduit en tout point de $X$. Rappelons qu'un anneau $A$ est dit
\emph{régulier} si chacun des anneaux locaux $A_{\mathfrak p}$ (où
$\mathfrak p$ parcourt l'ensemble des idéaux premiers de $A$) est un anneau
local régulier~; nous dirons qu'un faisceau d'anneaux $\mathcal{A}$ sur $X$
est \emph{régulier en un point} $x$ (resp. \emph{régulier}) si la fibre
$\mathcal{A}_x$ est un anneau régulier (resp. si $\mathcal{A}$ est régulier
en tout point). Enfin, nous dirons qu'un faisceau d'anneaux $\mathcal{A}$ sur
$X$ est \emph{normal en un point} $x$ (resp. \emph{normal}) si sa fibre
$\mathcal{A}_x$ est un anneau intègre et intégralement clos (resp. si
$\mathcal{A}$ est normal en tout point). Nous dirons qu'un espace annelé
$(X,\mathcal{A})$ a l'une des propriétés précédentes si le faisceau d'anneaux
$\mathcal{A}$ a cette propriété.

Un faisceau d'anneaux \emph{gradués} $\mathcal{A}$ est par définition un
faisceau d'anneaux qui est somme directe (G, II, 2.7) d'une famille
$(\mathcal{A}_n)_{n\in\mathbf Z}$ de faisceaux de groupes abéliens avec les
conditions $\mathcal{A}_m\mathcal{A}_n\subset\mathcal{A}_{m+n}$~; un
$\mathcal{A}$-\emph{Module gradué} est un $\mathcal{A}$-Module
$\mathcal{F}$ somme directe d'une famille
$(\mathcal{F}_n)_{n\in\mathbf Z}$ de faisceaux de groupes abéliens,
satisfaisant aux conditions
$\mathcal{A}_m\mathcal{F}_n\subset\mathcal{F}_{m+n}$. Il revient d'ailleurs
au même de dire que l'on a
$(\mathcal{A}_m)_x(\mathcal{A}_n)_x\subset(\mathcal{A}_{m+n})_x$
(resp.
$(\mathcal{A}_m)_x(\mathcal{F}_n)_x\subset(\mathcal{F}_{m+n})_x$)
en tout point $x$.
\end{env}

\begin{env}[4.1.5]
\label{0.4.1.5-fr}
Étant donné un espace annelé $(X,\mathcal{A})$ (non nécessairement
commutatif), nous ne rappellerons pas ici les définitions des bifoncteurs
$\mathcal{F}\otimes_{\mathcal{A}}\mathcal{G}$,
$\mathcal{H}\!om_{\mathcal{A}}(\mathcal{F},\mathcal{G})$ et
$\operatorname{Hom}_{\mathcal{A}}(\mathcal{F},\mathcal{G})$
(G, II, 2.8 et 2.2) dans les catégories des $\mathcal{A}$-Modules à gauche ou
à droite (selon les cas), à valeurs dans la catégorie des faisceaux de groupes
abéliens (ou plus

\oldpage[0\textsubscript{I}]{38}
généralement des $\mathcal{C}$-Modules, si $\mathcal{C}$ est le centre de
$\mathcal{A}$). La fibre
$(\mathcal{F}\otimes_{\mathcal{A}}\mathcal{G})_x$ en tout point $x\in X$
s'identifie canoniquement à
$\mathcal{F}_x\otimes_{\mathcal{A}_x}\mathcal{G}_x$ et on définit un
homomorphisme canonique fonctoriel
\[
  (\mathcal{H}\!om_{\mathcal{A}}(\mathcal{F},\mathcal{G}))_x
  \longrightarrow
  \operatorname{Hom}_{\mathcal{A}_x}(\mathcal{F}_x,\mathcal{G}_x)
\]
qui, en général, n'est ni injectif ni surjectif. Les bifoncteurs considérés
ci-dessus sont additifs et, en particulier, commutent aux sommes directes
finies~; $\mathcal{F}\otimes_{\mathcal{A}}\mathcal{G}$ est exact à droite en
$\mathcal{F}$ et en $\mathcal{G}$, commute aux limites inductives, et
$\mathcal{A}\otimes_{\mathcal{A}}\mathcal{G}$
(resp. $\mathcal{F}\otimes_{\mathcal{A}}\mathcal{A}$) s'identifie
canoniquement à $\mathcal{G}$ (resp. $\mathcal{F}$). Les foncteurs
$\mathcal{H}\!om_{\mathcal{A}}(\mathcal{F},\mathcal{G})$ et
$\operatorname{Hom}_{\mathcal{A}}(\mathcal{F},\mathcal{G})$ sont
\emph{exacts à gauche} en $\mathcal{F}$ et $\mathcal{G}$~; de façon précise,
si on a une suite exacte de la forme
$0\to\mathcal{G}'\to\mathcal{G}\to\mathcal{G}''$, la suite
\[
  0\longrightarrow
  \mathcal{H}\!om_{\mathcal{A}}(\mathcal{F},\mathcal{G}')
  \longrightarrow
  \mathcal{H}\!om_{\mathcal{A}}(\mathcal{F},\mathcal{G})
  \longrightarrow
  \mathcal{H}\!om_{\mathcal{A}}(\mathcal{F},\mathcal{G}'')
\]
est exacte, et si on a une suite exacte de la forme
$\mathcal{F}'\to\mathcal{F}\to\mathcal{F}''\to0$, la suite
\[
  0\longrightarrow
  \mathcal{H}\!om_{\mathcal{A}}(\mathcal{F}'',\mathcal{G})
  \longrightarrow
  \mathcal{H}\!om_{\mathcal{A}}(\mathcal{F},\mathcal{G})
  \longrightarrow
  \mathcal{H}\!om_{\mathcal{A}}(\mathcal{F}',\mathcal{G})
\]
est exacte, avec des propriétés analogues pour le foncteur
$\operatorname{Hom}$. En outre,
$\mathcal{H}\!om_{\mathcal{A}}(\mathcal{A},\mathcal{G})$ s'identifie
canoniquement à $\mathcal{G}$~; enfin, pour tout ouvert $U\subset X$, on a
\[
  \Gamma\bigl(U,\mathcal{H}\!om_{\mathcal{A}}(\mathcal{F},\mathcal{G})\bigr)
  =
  \operatorname{Hom}_{\mathcal{A}|U}(\mathcal{F}|U,\mathcal{G}|U).
\]

Pour tout $\mathcal{A}$-Module à gauche $\mathcal{F}$ (resp. à droite), on
appelle \emph{dual de $\mathcal{F}$} et on note $\check{\mathcal{F}}$ le
$\mathcal{A}$-Module à droite (resp. à gauche)
$\mathcal{H}\!om_{\mathcal{A}}(\mathcal{F},\mathcal{A})$.

Enfin, si $\mathcal{A}$ est un faisceau d'anneaux commutatifs,
$\mathcal{F}$ un $\mathcal{A}$-Module,
$U\mapsto\bigwedge^p\Gamma(U,\mathcal{F})$ est un préfaisceau dont le
faisceau associé est un $\mathcal{A}$-Module qui se note
$\bigwedge^p\mathcal{F}$ et s'appelle la \emph{puissance extérieure $p$-ème
de $\mathcal{F}$}~; on vérifie aisément que l'application canonique du
préfaisceau $U\mapsto\bigwedge^p\Gamma(U,\mathcal{F})$ dans le faisceau
associé $\bigwedge^p\mathcal{F}$ est \emph{injective}, et que pour tout
$x\in X$, on a
$(\bigwedge^p\mathcal{F})_x=\bigwedge^p(\mathcal{F}_x)$. Il est clair que
$\bigwedge^p\mathcal{F}$ est un foncteur covariant en $\mathcal{F}$.
\end{env}

\begin{env}[4.1.6]
\label{0.4.1.6-fr}
Supposons que $\mathcal{A}$ soit un faisceau d'anneaux non nécessairement
commutatifs, $\mathcal{J}$ un faisceau d'idéaux à gauche de $\mathcal{A}$,
$\mathcal{F}$ un $\mathcal{A}$-Module à gauche~; on note alors
$\mathcal{J}\mathcal{F}$ le sous-$\mathcal{A}$-Module de $\mathcal{F}$,
image de $\mathcal{J}\otimes_{\mathbf Z}\mathcal{F}$ (où $\mathbf Z$ est le
faisceau associé au préfaisceau constant $U\to\mathbf Z$) par l'application
canonique $\mathcal{J}\otimes_{\mathbf Z}\mathcal{F}\to\mathcal{F}$~; il est
clair que pour tout $x\in X$, on a
$(\mathcal{J}\mathcal{F})_x=\mathcal{J}_x\mathcal{F}_x$. Lorsque
$\mathcal{A}$ est commutatif, $\mathcal{J}\mathcal{F}$ est aussi l'image
canonique de $\mathcal{J}\otimes_{\mathcal{A}}\mathcal{F}\to\mathcal{F}$.
Il est immédiat que $\mathcal{J}\mathcal{F}$ est aussi le
$\mathcal{A}$-Module associé au préfaisceau
$U\to\Gamma(U,\mathcal{J})\Gamma(U,\mathcal{F})$. Si
$\mathcal{J}_1$, $\mathcal{J}_2$ sont deux faisceaux d'idéaux à gauche de
$\mathcal{A}$, on a
$\mathcal{J}_1(\mathcal{J}_2\mathcal{F})=(\mathcal{J}_1\mathcal{J}_2)\mathcal{F}$.
\end{env}

\begin{env}[4.1.7]
\label{0.4.1.7-fr}
Soit $(X_\lambda,\mathcal{A}_\lambda)_{\lambda\in L}$ une famille d'espaces
annelés~; pour tout couple $(\lambda,\mu)$, supposons donnée une partie ouverte
$V_{\lambda\mu}$ de $X_\lambda$, et un isomorphisme d'espaces annelés
\[
  \varphi_{\lambda\mu}:
  (V_{\mu\lambda},\mathcal{A}_\mu|V_{\mu\lambda})
  \xrightarrow{\sim}
  (V_{\lambda\mu},\mathcal{A}_\lambda|V_{\lambda\mu}),
\]
avec $V_{\lambda\lambda}=X_\lambda$, $\varphi_{\lambda\lambda}$ étant
l'identité. Supposons de plus que, pour tout triplet $(\lambda,\mu,\nu)$, si
on désigne par $\varphi'_{\mu\lambda}$ la restriction de
$\varphi_{\mu\lambda}$ à $V_{\lambda\mu}\cap V_{\lambda\nu}$,
$\varphi'_{\mu\lambda}$ soit un isomorphisme de
\[
  (V_{\lambda\mu}\cap V_{\lambda\nu},
   \mathcal{A}_\lambda|(V_{\lambda\mu}\cap V_{\lambda\nu}))
\]
sur
\[
  (V_{\mu\nu}\cap V_{\mu\lambda},
   \mathcal{A}_\mu|(V_{\mu\nu}\cap V_{\mu\lambda}))
\]
et que l'on ait
$\varphi'_{\lambda\nu}=\varphi'_{\lambda\mu}\circ\varphi'_{\mu\nu}$
(\emph{condition de recollement} pour les $\varphi_{\lambda\mu}$). On peut
tout d'abord considérer alors l'espace topologique obtenu par recollement (au
moyen des $\varphi_{\lambda\mu}$) des $X_\lambda$

\oldpage[0\textsubscript{I}]{39}
le long des $V_{\lambda\mu}$~; si on identifie $X_\lambda$ à la partie ouverte
$X'_\lambda$ correspondante dans $X$, les hypothèses entraînent que les trois
ensembles $V_{\lambda\mu}\cap V_{\lambda\nu}$,
$V_{\mu\nu}\cap V_{\mu\lambda}$,
$V_{\nu\lambda}\cap V_{\nu\mu}$ s'identifient à
$X'_\lambda\cap X'_\mu\cap X'_\nu$. On peut alors transporter à
$X'_\lambda$ la structure d'espace annelé de $X_\lambda$, et si
$\mathcal{A}'_\lambda$ est le faisceau d'anneaux transporté de
$\mathcal{A}_\lambda$, les $\mathcal{A}'_\lambda$ vérifient la condition de
recollement (3.3.1) et définissent donc un faisceau d'anneaux $\mathcal{A}$ sur
$X$~; on dit que $(X,\mathcal{A})$ est l'espace annelé obtenu par
\emph{recollement des $(X_\lambda,\mathcal{A}_\lambda)$ le long des
$V_{\lambda\mu}$}, au moyen des $\varphi_{\lambda\mu}$.
\end{env}

\subsection{Image directe d'un $\mathcal{A}$-Module}
\label{subsection:0.4.2-fr}

\begin{env}[4.2.1]
\label{0.4.2.1-fr}
Soient $(X,\mathcal{A})$, $(Y,\mathcal{B})$ deux espaces annelés,
$\Psi=(\psi,\theta)$ un morphisme
$(X,\mathcal{A})\to(Y,\mathcal{B})$~; $\psi_*(\mathcal{A})$ est donc un
faisceau d'anneaux sur $Y$, et $\theta$ un homomorphisme
$\mathcal{B}\to\psi_*(\mathcal{A})$ de faisceaux d'anneaux. Soit alors
$\mathcal{F}$ un $\mathcal{A}$-Module~; son image directe
$\psi_*(\mathcal{F})$ est un faisceau de groupes abéliens sur $Y$. En outre,
pour tout ouvert $U\subset Y$,
\[
  \Gamma(U,\psi_*(\mathcal{F}))
  =\Gamma(\psi^{-1}(U),\mathcal{F})
\]
est muni d'une structure de module par rapport à l'anneau
$\Gamma(U,\psi_*(\mathcal{A}))=\Gamma(\psi^{-1}(U),\mathcal{A})$~; les
applications bilinéaires qui définissent ces structures étant compatibles
avec les opérations de restriction, définissent sur $\psi_*(\mathcal{F})$ une
structure de $\psi_*(\mathcal{A})$-Module. L'homomorphisme
$\theta:\mathcal{B}\to\psi_*(\mathcal{A})$ permet alors de définir aussi sur
$\psi_*(\mathcal{F})$ une structure de $\mathcal{B}$-Module~; nous dirons que
ce $\mathcal{B}$-Module est l'\emph{image directe de $\mathcal{F}$ par le
morphisme $\Psi$}, et nous le noterons $\Psi_*(\mathcal{F})$. Si
$\mathcal{F}_1$, $\mathcal{F}_2$ sont deux $\mathcal{A}$-Modules sur $X$ et
$u$ un $\mathcal{A}$-homomorphisme $\mathcal{F}_1\to\mathcal{F}_2$, il est
immédiat (en considérant les sections au-dessus des ouverts de $Y$) que
$\psi_*(u)$ est un $\psi_*(\mathcal{A})$-homomorphisme
$\psi_*(\mathcal{F}_1)\to\psi_*(\mathcal{F}_2)$, et \emph{a fortiori} un
$\mathcal{B}$-homomorphisme
$\Psi_*(\mathcal{F}_1)\to\Psi_*(\mathcal{F}_2)$~; en tant que
$\mathcal{B}$-homomorphisme, nous le noterons $\Psi_*(u)$. On voit donc que
$\Psi_*$ est un \emph{foncteur covariant} de la catégorie des
$\mathcal{A}$-Modules dans celle des $\mathcal{B}$-Modules. En outre, il est
immédiat que ce foncteur est \emph{exact à gauche} (G, II, 2.12).

Sur $\psi_*(\mathcal{A})$, la structure de $\mathcal{B}$-Module et la
structure de faisceaux d'anneaux définissent une structure de
$\mathcal{B}$-Algèbre~; on notera $\Psi_*(\mathcal{A})$ cette
$\mathcal{B}$-Algèbre.
\end{env}

\begin{env}[4.2.2]
\label{0.4.2.2-fr}
Soient $\mathcal{M}$, $\mathcal{N}$ deux $\mathcal{A}$-Modules. Pour tout
ouvert $U$ de $Y$, on a une application canonique
\[
  \Gamma(\psi^{-1}(U),\mathcal{M})
  \times\Gamma(\psi^{-1}(U),\mathcal{N})
  \longrightarrow
  \Gamma(\psi^{-1}(U),\mathcal{M}\otimes_{\mathcal{A}}\mathcal{N})
\]
qui est bilinéaire sur l'anneau
$\Gamma(\psi^{-1}(U),\mathcal{A})=\Gamma(U,\psi_*(\mathcal{A}))$, et
\emph{a fortiori} sur $\Gamma(U,\mathcal{B})$~; elle définit donc un
homomorphisme
\[
  \Gamma(U,\Psi_*(\mathcal{M}))
  \otimes_{\Gamma(U,\mathcal{B})}
  \Gamma(U,\Psi_*(\mathcal{N}))
  \longrightarrow
  \Gamma(U,\Psi_*(\mathcal{M}\otimes_{\mathcal{A}}\mathcal{N}))
\]
et comme on vérifie aussitôt que ces homomorphismes sont compatibles avec les
opérations de restriction, ils donnent un homomorphisme canonique fonctoriel
de $\mathcal{B}$-Modules
\[
  \Psi_*(\mathcal{M})\otimes_{\mathcal{B}}\Psi_*(\mathcal{N})
  \longrightarrow
  \Psi_*(\mathcal{M}\otimes_{\mathcal{A}}\mathcal{N}).
  \tag{4.2.2.1}\label{0.4.2.2.1-fr}
\]

\oldpage[0\textsubscript{I}]{40}
qui ne sera en général ni injectif ni surjectif. Si $\mathcal{P}$ est un
troisième $\mathcal{A}$-Module, on vérifie aussitôt que le diagramme
\[
\begin{tikzcd}[column sep=large,row sep=large]
\Psi_*(\mathcal{M})\otimes_{\mathcal{B}}\Psi_*(\mathcal{N})
  \otimes_{\mathcal{B}}\Psi_*(\mathcal{P})
  \arrow[r]
  \arrow[d]
&
\Psi_*(\mathcal{M}\otimes_{\mathcal{A}}\mathcal{N})
  \otimes_{\mathcal{B}}\Psi_*(\mathcal{P})
  \arrow[d]
\\
\Psi_*(\mathcal{M})\otimes_{\mathcal{B}}
  \Psi_*(\mathcal{N}\otimes_{\mathcal{A}}\mathcal{P})
  \arrow[r]
&
\Psi_*(\mathcal{M}\otimes_{\mathcal{A}}\mathcal{N}
  \otimes_{\mathcal{A}}\mathcal{P})
\end{tikzcd}
\tag{4.2.2.2}\label{0.4.2.2.2-fr}
\]
est commutatif.
\end{env}

\begin{env}[4.2.3]
\label{0.4.2.3-fr}
Soient $\mathcal{M}$, $\mathcal{N}$ deux $\mathcal{A}$-Modules. Pour tout
ouvert $U\subset Y$, on a par définition
\[
  \Gamma\bigl(\psi^{-1}(U),
    \mathcal{H}\!om_{\mathcal{A}}(\mathcal{M},\mathcal{N})\bigr)
  =
  \operatorname{Hom}_{\mathcal{A}|V}
    (\mathcal{M}|V,\mathcal{N}|V),
\]
où on a posé pour alléger $V=\psi^{-1}(U)$~; l'application
$u\to\Psi_*(u)$ est un homomorphisme
\[
  \operatorname{Hom}_{\mathcal{A}|V}(\mathcal{M}|V,\mathcal{N}|V)
  \longrightarrow
  \operatorname{Hom}_{\mathcal{B}|U}
    (\Psi_*(\mathcal{M})|U,\Psi_*(\mathcal{N})|U)
\]
pour les structures de $\Gamma(U,\mathcal{B})$-modules~; ces homomorphismes
étant compatibles aux opérations de restriction, définissent donc un
homomorphisme canonique fonctoriel de $\mathcal{B}$-Modules
\[
  \Psi_*\bigl(\mathcal{H}\!om_{\mathcal{A}}(\mathcal{M},\mathcal{N})\bigr)
  \longrightarrow
  \mathcal{H}\!om_{\mathcal{B}}
    (\Psi_*(\mathcal{M}),\Psi_*(\mathcal{N})).
  \tag{4.2.3.1}\label{0.4.2.3.1-fr}
\]
\end{env}

\begin{env}[4.2.4]
\label{0.4.2.4-fr}
Si $\mathcal{C}$ est une $\mathcal{A}$-Algèbre, l'homomorphisme composé
\[
  \Psi_*(\mathcal{C})\otimes_{\mathcal{B}}\Psi_*(\mathcal{C})
  \longrightarrow
  \Psi_*(\mathcal{C}\otimes_{\mathcal{A}}\mathcal{C})
  \longrightarrow
  \Psi_*(\mathcal{C})
\]
définit sur $\Psi_*(\mathcal{C})$ une structure de $\mathcal{B}$-Algèbre,
comme il résulte de (4.2.2.2). On voit de même que si $\mathcal{M}$ est un
$\mathcal{C}$-Module, $\Psi_*(\mathcal{M})$ est muni canoniquement d'une
structure de $\Psi_*(\mathcal{C})$-Module.
\end{env}

\begin{env}[4.2.5]
\label{0.4.2.5-fr}
Considérons en particulier le cas où $X$ est un sous-espace \emph{fermé} de
$Y$ et où $\psi$ est l'injection canonique $j:X\to Y$. Si
$\mathcal{B}'=\mathcal{B}|X=j^*(\mathcal{B})$ est la restriction du faisceau
d'anneaux $\mathcal{B}$ à $X$, un $\mathcal{A}$-Module $\mathcal{M}$ peut
être considéré comme $\mathcal{B}'$-Module au moyen de l'homomorphisme
$\theta^\sharp:\mathcal{B}'\to\mathcal{A}$~; $\Psi_*(\mathcal{M})$ est alors
le $\mathcal{B}$-Module qui induit $\mathcal{M}$ sur $X$ et $0$ ailleurs. Si
$\mathcal{N}$ est un second $\mathcal{A}$-Module,
$\Psi_*(\mathcal{M})\otimes_{\mathcal{B}}\Psi_*(\mathcal{N})$ s'identifie
donc à $\Psi_*(\mathcal{M}\otimes_{\mathcal{B}'}\mathcal{N})$ et
$\mathcal{H}\!om_{\mathcal{B}}(\Psi_*(\mathcal{M}),\Psi_*(\mathcal{N}))$
à
$\Psi_*\bigl(\mathcal{H}\!om_{\mathcal{B}'}(\mathcal{M},\mathcal{N})\bigr)$.
\end{env}

\begin{env}[4.2.6]
\label{0.4.2.6-fr}
Soient $(Z,\mathcal{C})$ un troisième espace annelé,
$\Psi'=(\psi',\theta')$ un morphisme
$(Y,\mathcal{B})\to(Z,\mathcal{C})$~; si $\Psi''$ est le morphisme composé
$\Psi'\circ\Psi$, il est clair que l'on a
$\Psi''_*=\Psi'_*\circ\Psi_*$.
\end{env}

\oldpage[0\textsubscript{I}]{41}
\subsection{Image réciproque d'un $\mathcal{B}$-Module}
\label{subsection:0.4.3-fr}

\begin{env}[4.3.1]
\label{0.4.3.1-fr}
Les hypothèses et notations étant celles de (4.2.1), soit $\mathcal{G}$ un
$\mathcal{B}$-Module et $\psi^*(\mathcal{G})$ son image réciproque (3.7.1)
qui est donc un faisceau de groupes abéliens sur $X$. La définition des
sections de $\psi^*(\mathcal{G})$ et de $\psi^*(\mathcal{B})$ (3.7.1)
montre que $\psi^*(\mathcal{G})$ est canoniquement muni d'une structure de
$\psi^*(\mathcal{B})$-Module. D'autre part, l'homomorphisme
$\theta^\sharp:\psi^*(\mathcal{B})\to\mathcal{A}$ munit $\mathcal{A}$ d'une
structure de $\psi^*(\mathcal{B})$-Module, que nous noterons
$\mathcal{A}_{[\theta]}$ lorsqu'il y a lieu d'éviter des confusions~; le
produit tensoriel
$\psi^*(\mathcal{G})\otimes_{\psi^*(\mathcal{B})}\mathcal{A}_{[\theta]}$
est alors muni d'une structure de $\mathcal{A}$-Module. Nous dirons que ce
$\mathcal{A}$-Module est l'image réciproque de $\mathcal{G}$ par le morphisme
$\Psi$ et nous le noterons $\Psi^*(\mathcal{G})$. Si $\mathcal{G}_1$,
$\mathcal{G}_2$ sont deux $\mathcal{B}$-Modules sur $Y$, $v$ un
$\mathcal{B}$-homomorphisme $\mathcal{G}_1\to\mathcal{G}_2$,
$\psi^*(v)$, comme on le vérifie aussitôt, est un
$\psi^*(\mathcal{B})$-homomorphisme de $\psi^*(\mathcal{G}_1)$ dans
$\psi^*(\mathcal{G}_2)$~; par suite $\psi^*(v)\otimes1$ est un
$\mathcal{A}$-homomorphisme $\Psi^*(\mathcal{G}_1)\to\Psi^*(\mathcal{G}_2)$,
que nous noterons $\Psi^*(v)$. On a donc défini $\Psi^*$ comme un
\emph{foncteur covariant} de la catégorie des $\mathcal{B}$-Modules dans celle
des $\mathcal{A}$-Modules. Ici, ce foncteur (contrairement à $\psi^*$) n'est
plus exact en général, mais seulement \emph{exact à droite}, la tensorisation
par $\mathcal{A}$ étant un foncteur exact à droite dans la catégorie des
$\psi^*(\mathcal{B})$-Modules.

Pour tout $x\in X$, on a
\[
  (\Psi^*(\mathcal{G}))_x
  =\mathcal{G}_{\psi(x)}\otimes_{\mathcal{B}_{\psi(x)}}\mathcal{A}_x,
\]
en vertu de (3.7.2). Le support de $\Psi^*(\mathcal{G})$ est donc contenu dans
$\psi^{-1}(\operatorname{Supp}\mathcal{G})$.
\end{env}

\begin{env}[4.3.2]
\label{0.4.3.2-fr}
Soit $(\mathcal{G}_\lambda)$ un système inductif de $\mathcal{B}$-Modules, et
soit $\mathcal{G}=\varinjlim\mathcal{G}_\lambda$ sa limite inductive. Les
homomorphismes canoniques $\mathcal{G}_\lambda\to\mathcal{G}$ définissent des
$\psi^*(\mathcal{B})$-homomorphismes
$\psi^*(\mathcal{G}_\lambda)\to\psi^*(\mathcal{G})$, qui donnent un
homomorphisme canonique
\[
  \varinjlim\psi^*(\mathcal{G}_\lambda)
  \longrightarrow
  \psi^*(\mathcal{G}).
\]
Comme la fibre en un point d'une limite inductive de faisceaux est la limite
inductive des fibres au même point (G, II, 1.11), l'homomorphisme canonique
précédent est \emph{bijectif} (3.7.2). En outre, le produit tensoriel commute
aux limites inductives de faisceaux, et on a donc un \emph{isomorphisme
canonique fonctoriel}
\[
  \varinjlim\Psi^*(\mathcal{G}_\lambda)
  \xrightarrow{\sim}
  \Psi^*\bigl(\varinjlim\mathcal{G}_\lambda\bigr)
\]
de $\mathcal{A}$-Modules.

D'autre part, pour une somme directe finie $\bigoplus_i\mathcal{G}_i$ de
$\mathcal{B}$-Modules, il est clair que
\[
  \psi^*\bigl(\bigoplus_i\mathcal{G}_i\bigr)
  =\bigoplus_i\psi^*(\mathcal{G}_i)~;
\]
donc, par produit tensoriel avec $\mathcal{A}_{[\theta]}$,
\[
  \Psi^*\bigl(\bigoplus_i\mathcal{G}_i\bigr)
  =\bigoplus_i\Psi^*(\mathcal{G}_i).
  \tag{4.3.2.1}\label{0.4.3.2.1-fr}
\]
Par passage à la limite inductive, on en déduit, en vertu de ce qui précède,
que l'égalité précédente est encore vraie pour une somme directe
\emph{quelconque}.
\end{env}

\begin{env}[4.3.3]
\EGAReaderTarget{canon-ref.0.4.3.3-fr}{0.4.3.3-fr}
Soient $\mathcal{G}_1$, $\mathcal{G}_2$ deux $\mathcal{B}$-Modules~; de la
définition des images réciproques de faisceaux de groupes abéliens (3.7.1),
on déduit aussitôt un homomorphisme canonique
\[
  \psi^*(\mathcal{G}_1)\otimes_{\psi^*(\mathcal{B})}\psi^*(\mathcal{G}_2)
  \longrightarrow
  \psi^*(\mathcal{G}_1\otimes_{\mathcal{B}}\mathcal{G}_2)
\]
de $\psi^*(\mathcal{B})$-Modules, et la fibre en un point d'un produit
tensoriel de faisceaux étant le produit tensoriel des fibres en ce point
(G, II, 2.8), on déduit de (3.7.2) que l'homomorphisme précédent est en fait
un \emph{isomorphisme}. Par produit tensoriel avec $\mathcal{A}$, on en déduit
un \emph{isomorphisme canonique fonctoriel}
\[
  \Psi^*(\mathcal{G}_1)\otimes_{\mathcal{A}}\Psi^*(\mathcal{G}_2)
  \xrightarrow{\sim}
  \Psi^*(\mathcal{G}_1\otimes_{\mathcal{B}}\mathcal{G}_2).
  \tag{4.3.3.1}\label{0.4.3.3.1-fr}
\]
\end{env}

\begin{env}[4.3.4]
\label{0.4.3.4-fr}
Soit $\mathcal{C}$ une $\mathcal{B}$-Algèbre~; la donnée de la structure
d'Algèbre sur $\mathcal{C}$ revient à la donnée d'un $\mathcal{B}$-homomorphisme
\[
  \mathcal{C}\otimes_{\mathcal{B}}\mathcal{C}\longrightarrow\mathcal{C}
\]
satisfaisant aux conditions d'associativité et de commutativité (conditions
qui se vérifient sur chaque fibre)~; l'isomorphisme précédent permet de
transporter cet homomorphisme en un homomorphisme de $\mathcal{A}$-Modules
\[
  \Psi^*(\mathcal{C})\otimes_{\mathcal{A}}\Psi^*(\mathcal{C})
  \longrightarrow\Psi^*(\mathcal{C})
\]
satisfaisant aux mêmes conditions, donc $\Psi^*(\mathcal{C})$ est ainsi muni
d'une structure de $\mathcal{A}$-Algèbre. En particulier, il résulte aussitôt
des définitions que la $\mathcal{A}$-Algèbre $\Psi^*(\mathcal{B})$ est égale à
$\mathcal{A}$ (à un isomorphisme canonique près).

De même, si $\mathcal{M}$ est un $\mathcal{C}$-Module, la donnée de cette
structure de Module revient
\oldpage[0\textsubscript{I}]{42}
à celle d'un $\mathcal{B}$-homomorphisme
$\mathcal{C}\otimes_{\mathcal{B}}\mathcal{M}\to\mathcal{M}$ vérifiant la
condition d'associativité~; d'où par transport une structure de
$\Psi^*(\mathcal{C})$-Module sur $\Psi^*(\mathcal{M})$.
\end{env}

\begin{env}[4.3.5]
\label{0.4.3.5-fr}
Soit $\mathcal{J}$ un faisceau d'idéaux de $\mathcal{B}$~; comme le foncteur
$\psi^*$ est exact, le $\psi^*(\mathcal{B})$-Module $\psi^*(\mathcal{J})$
s'identifie canoniquement à un faisceau d'idéaux de $\psi^*(\mathcal{B})$~;
l'injection canonique $\psi^*(\mathcal{J})\to\psi^*(\mathcal{B})$ donne alors
un homomorphisme de $\mathcal{A}$-Modules
\[
  \Psi^*(\mathcal{J})
  =\psi^*(\mathcal{J})\otimes_{\psi^*(\mathcal{B})}\mathcal{A}_{[\theta]}
  \longrightarrow\mathcal{A}~;
\]
nous noterons $\Psi^*(\mathcal{J})\mathcal{A}$, ou
$\mathcal{J}\mathcal{A}$ si aucune confusion n'est à craindre, l'image de
$\Psi^*(\mathcal{J})$ par cet homomorphisme. On a donc par définition
\[
  \mathcal{J}\mathcal{A}
  =\theta^\sharp\bigl(\psi^*(\mathcal{J})\bigr)\mathcal{A}
\]
et en particulier, pour tout $x\in X$,
\[
  (\mathcal{J}\mathcal{A})_x
  =\theta_x^\sharp(\mathcal{J}_{\psi(x)})\mathcal{A}_x,
\]
en tenant compte de l'identification canonique des fibres de
$\psi^*(\mathcal{J})$ et de celles de $\mathcal{J}$ (3.7.2). Si
$\mathcal{J}_1$, $\mathcal{J}_2$ sont deux idéaux de $\mathcal{B}$, on a
\[
  (\mathcal{J}_1\mathcal{J}_2)\mathcal{A}
  =\mathcal{J}_1(\mathcal{J}_2\mathcal{A})
  =(\mathcal{J}_1\mathcal{A})(\mathcal{J}_2\mathcal{A}).
\]
Si $\mathcal{F}$ est un $\mathcal{A}$-Module, on posera
$\mathcal{J}\mathcal{F}=(\mathcal{J}\mathcal{A})\mathcal{F}$.
\end{env}

\begin{env}[4.3.6]
\label{0.4.3.6-fr}
Soient $(Z,\mathcal{C})$ un troisième espace annelé,
$\Psi'=(\psi',\theta')$ un morphisme $(Y,\mathcal{B})\to(Z,\mathcal{C})$~;
si $\Psi''$ est le morphisme composé $\Psi'\circ\Psi$, il résulte de la
définition (4.3.1) et de (4.3.3.1) que l'on a
${\Psi''}^*=\Psi^*\circ{\Psi'}^*$.
\end{env}

\subsection{Relations entre images directes et images réciproques}
\label{subsection:0.4.4-fr}

\begin{env}[4.4.1]
\label{0.4.4.1-fr}
Les hypothèses et notations étant celles de (4.2.1), soit $\mathcal{G}$ un
$\mathcal{B}$-Module. Par définition, un homomorphisme
$u:\mathcal{G}\to\Psi_*(\mathcal{F})$ de $\mathcal{B}$-Modules s'appelle
encore un \emph{$\Psi$-morphisme de $\mathcal{G}$ dans $\mathcal{F}$}, ou
simplement un \emph{homomorphisme de $\mathcal{G}$ dans $\mathcal{F}$} et on
l'écrit $u:\mathcal{G}\to\mathcal{F}$ quand aucune confusion n'en résulte.
Se donner un tel homomorphisme revient à se donner, pour tout couple $(U,V)$
où $U$ est un ouvert de $X$, $V$ un ouvert de $Y$ tels que
$\psi(U)\subset V$, un \emph{homomorphisme}
\[
  u_{U,V}:\Gamma(V,\mathcal{G})\longrightarrow\Gamma(U,\mathcal{F})
\]
\emph{de $\Gamma(V,\mathcal{B})$-modules}, $\Gamma(U,\mathcal{F})$ étant
considéré comme $\Gamma(V,\mathcal{B})$-module au moyen de l'homomorphisme
d'anneaux
$\theta_{U,V}:\Gamma(V,\mathcal{B})\to\Gamma(U,\mathcal{A})$~; les $u_{U,V}$
doivent en outre rendre commutatifs les diagrammes (3.5.1.1). Il suffit
d'ailleurs pour définir $u$ de se donner les $u_{U,V}$ lorsque $U$ (resp.
$V$) parcourt une base $\mathfrak{B}$ (resp. $\mathfrak{B}'$) de la topologie
de $X$ (resp. $Y$) et de vérifier la commutativité de (3.5.1.1) pour ces
restrictions.
\end{env}

\begin{env}[4.4.2]
\label{0.4.4.2-fr}
Sous les hypothèses de (4.2.1) et (4.2.6), soient $\mathcal{H}$ un
$\mathcal{C}$-Module, $v:\mathcal{H}\to\Psi'_*(\mathcal{G})$ un
$\Psi'$-morphisme~; alors
\[
  w:\mathcal{H}\xrightarrow{v}\Psi'_*(\mathcal{G})
  \xrightarrow{\Psi'_*(u)}\Psi'_*(\Psi_*(\mathcal{F}))
\]
est un $\Psi''$-morphisme que l'on appelle le \emph{composé de $u$ et $v$}.
\end{env}

\begin{env}[4.4.3]
\label{0.4.4.3-fr}
Nous allons maintenant voir que l'on peut définir un \emph{isomorphisme}
canonique de \emph{bifoncteurs} en $\mathcal{F}$ et $\mathcal{G}$
\[
  \operatorname{Hom}_{\mathcal{A}}(\Psi^*(\mathcal{G}),\mathcal{F})
  \xrightarrow{\sim}
  \operatorname{Hom}_{\mathcal{B}}(\mathcal{G},\Psi_*(\mathcal{F}))
  \tag{4.4.3.1}\label{0.4.4.3.1-fr}
\]
que nous désignerons par $v\mapsto v_\theta^\flat$ (ou simplement
$v\mapsto v^\flat$ si aucune confusion n'est possible)~; nous noterons
$u\mapsto u_\theta^\sharp$, ou $u\mapsto u^\sharp$ l'isomorphisme réciproque.
Cette définition est la suivante~: en composant
$v:\Psi^*(\mathcal{G})\to\mathcal{F}$ avec l'application canonique
$\psi^*(\mathcal{G})\to\Psi^*(\mathcal{G})$, on obtient un homomorphisme de
faisceaux de groupes $v':\psi^*(\mathcal{G})\to\mathcal{F}$, qui est aussi un
homomorphisme de $\psi^*(\mathcal{B})$-modules. On en déduit (3.7.1) un
homomorphisme
${v'}^\flat:\mathcal{G}\to\psi_*(\mathcal{F})=\Psi_*(\mathcal{F})$, qui est
aussi un homomorphisme de $\mathcal{B}$-Modules comme on
\oldpage[0\textsubscript{I}]{43}
le vérifie sans peine~; on prend $v_\theta^\flat={v'}^\flat$. De même, de
$u:\mathcal{G}\to\Psi_*(\mathcal{F})$, qui est un homomorphisme de
$\mathcal{B}$-Modules, on déduit (3.7.1) un homomorphisme
$u^\sharp:\psi^*(\mathcal{G})\to\mathcal{F}$ de
$\psi^*(\mathcal{B})$-Modules, d'où par tensorisation avec $\mathcal{A}$ un
homomorphisme de $\mathcal{A}$-Modules
$\Psi^*(\mathcal{G})\to\mathcal{F}$, que nous désignerons par
$u_\theta^\sharp$. Il est immédiat de vérifier que
$(u_\theta^\sharp)_\theta^\flat=u$ et
$(v_\theta^\flat)_\theta^\sharp=v$, ainsi que le caractère fonctoriel en
$\mathcal{F}$ de l'isomorphisme $v\mapsto v_\theta^\flat$. Le caractère
fonctoriel en $\mathcal{G}$ de $u\mapsto u_\theta^\sharp$ s'en déduit alors
formellement comme dans (3.5.4) (raisonnement qui démontrerait aussi le
caractère fonctoriel de $\Psi^*$ établi dans (4.3.1) directement).

Si on prend pour $v$ l'homomorphisme identique de $\Psi^*(\mathcal{G})$,
$v_\theta^\flat$ est un homomorphisme
\[
  \rho_{\mathcal{G}}:\mathcal{G}\longrightarrow
  \Psi_*(\Psi^*(\mathcal{G}))~;
  \tag{4.4.3.2}\label{0.4.4.3.2-fr}
\]
si on prend pour $u$ l'homomorphisme identique de $\Psi_*(\mathcal{F})$,
$u_\theta^\sharp$ est un homomorphisme
\[
  \sigma_{\mathcal{F}}:\Psi^*(\Psi_*(\mathcal{F}))
  \longrightarrow\mathcal{F}~;
  \tag{4.4.3.3}\label{0.4.4.3.3-fr}
\]
ces homomorphismes seront dits \emph{canoniques}. Ils ne sont en général ni
injectifs ni surjectifs. On a des factorisations canoniques analogues à
(3.5.3.3) et (3.5.4.4).

On notera que si $s$ est une section de $\mathcal{G}$ au-dessus d'un ouvert
$V$ de $Y$, $\rho_{\mathcal{G}}(s)$ est la section $s'\otimes1$ de
$\Psi^*(\mathcal{G})$ au-dessus de $\psi^{-1}(V)$, $s'$ étant telle que
$s'_x=s_{\psi(x)}$ pour tout $x\in\psi^{-1}(V)$. Notons aussi que si
$u:\mathcal{G}\to\psi_*(\mathcal{F})$ est un homomorphisme, il définit pour
tout $x\in X$ un homomorphisme
$u_x:\mathcal{G}_{\psi(x)}\to\mathcal{F}_x$ sur les fibres, obtenu en
composant
$(u^\sharp)_x:(\Psi^*(\mathcal{G}))_x\to\mathcal{F}_x$ et l'homomorphisme
canonique $s_x\mapsto s_x\otimes1$ de $\mathcal{G}_{\psi(x)}$ dans
$(\Psi^*(\mathcal{G}))_x
=\mathcal{G}_{\psi(x)}\otimes_{\mathcal{B}_{\psi(x)}}\mathcal{A}_x$.
L'homomorphisme $u_x$ s'obtient aussi par passage à la limite inductive à
partir des homomorphismes
\[
  \Gamma(V,\mathcal{G})\xrightarrow{u}
  \Gamma(\psi^{-1}(V),\mathcal{F})\longrightarrow\mathcal{F}_x,
\]
où $V$ parcourt les voisinages de $\psi(x)$.
\end{env}

\begin{env}[4.4.4]
\label{0.4.4.4-fr}
Soient $\mathcal{F}_1$, $\mathcal{F}_2$ des $\mathcal{A}$-Modules,
$\mathcal{G}_1$, $\mathcal{G}_2$ des $\mathcal{B}$-Modules, $u_i$ ($i=1,2$)
un homomorphisme de $\mathcal{G}_i$ dans $\mathcal{F}_i$. Nous désignerons par
$u_1\otimes u_2$ l'homomorphisme
\[
  u:\mathcal{G}_1\otimes_{\mathcal{B}}\mathcal{G}_2
  \longrightarrow
  \mathcal{F}_1\otimes_{\mathcal{A}}\mathcal{F}_2
\]
tel que $u^\sharp=(u_1)^\sharp\otimes(u_2)^\sharp$ (compte tenu de
(4.3.3.1))~; on vérifie que $u$ est aussi le composé
\[
  \mathcal{G}_1\otimes_{\mathcal{B}}\mathcal{G}_2
  \longrightarrow
  \Psi_*(\mathcal{F}_1)\otimes_{\mathcal{B}}\Psi_*(\mathcal{F}_2)
  \longrightarrow
  \Psi_*(\mathcal{F}_1\otimes_{\mathcal{A}}\mathcal{F}_2),
\]
où la première flèche est le produit tensoriel ordinaire
$u_1\otimes_{\mathcal{B}}u_2$ et la seconde l'homomorphisme canonique
(4.2.2.1).
\end{env}

\begin{env}[4.4.5]
\label{0.4.4.5-fr}
Soient $(\mathcal{G}_\lambda)_{\lambda\in L}$ un système inductif de
$\mathcal{B}$-Modules, et, pour tout $\lambda\in L$, soit $u_\lambda$ un
homomorphisme $\mathcal{G}_\lambda\to\Psi_*(\mathcal{F})$, formant un système
inductif~; posons
$\mathcal{G}=\varinjlim\mathcal{G}_\lambda$ et
$u=\varinjlim u_\lambda$~; alors les $(u_\lambda)^\sharp$ forment un système
inductif d'homomorphismes $\Psi^*(\mathcal{G}_\lambda)\to\mathcal{F}$, et la
limite inductive de ce système n'est autre que $u^\sharp$.
\end{env}

\begin{env}[4.4.6]
\label{0.4.4.6-fr}
Soient $\mathcal{M}$, $\mathcal{N}$ deux $\mathcal{B}$-Modules, $V$ un ouvert
de $Y$, $U=\psi^{-1}(V)$~; l'application $v\mapsto\Psi^*(v)$ est un
homomorphisme
\[
  \operatorname{Hom}_{\mathcal{B}|V}(\mathcal{M}|V,\mathcal{N}|V)
  \longrightarrow
  \operatorname{Hom}_{\mathcal{A}|U}
  (\Psi^*(\mathcal{M})|U,\Psi^*(\mathcal{N})|U)
\]
\begin{sloppypar}
pour les structures de $\Gamma(V,\mathcal{B})$-module
($\operatorname{Hom}_{\mathcal{A}|U}
(\Psi^*(\mathcal{M})|U,\allowbreak\Psi^*(\mathcal{N})|U)$ est normalement muni d'une
structure de $\Gamma(U,\psi^*(\mathcal{B}))$-module, et grâce à
l'homomorphisme
\oldpage[0\textsubscript{I}]{44}
canonique (3.7.2)
$\Gamma(V,\mathcal{B})\to\Gamma(U,\psi^*(\mathcal{B}))$, c'est donc aussi un
$\Gamma(V,\mathcal{B})$-module).
\end{sloppypar}
On voit aussitôt que ces homomorphismes sont
compatibles avec les restrictions, et par suite définissent un homomorphisme
canonique fonctoriel
\[
  \gamma:\mathcal{H}\!om_{\mathcal{B}}(\mathcal{M},\mathcal{N})
  \longrightarrow
  \Psi_*\bigl(\mathcal{H}\!om_{\mathcal{A}}
  (\Psi^*(\mathcal{M}),\Psi^*(\mathcal{N}))\bigr)~;
\]
il correspond par suite aussi à cet homomorphisme l'homomorphisme
\[
  \gamma^\sharp:
  \Psi^*\bigl(\mathcal{H}\!om_{\mathcal{B}}(\mathcal{M},\mathcal{N})\bigr)
  \longrightarrow
  \mathcal{H}\!om_{\mathcal{A}}
  (\Psi^*(\mathcal{M}),\Psi^*(\mathcal{N}))
\]
et ces homomorphismes canoniques sont fonctoriels en $\mathcal{M}$ et
$\mathcal{N}$.
\end{env}

\begin{env}[4.4.7]
\label{0.4.4.7-fr}
Supposons que $\mathcal{F}$ (resp. $\mathcal{G}$) soit une
$\mathcal{A}$-Algèbre (resp. une $\mathcal{B}$-Algèbre). Si
$u:\mathcal{G}\to\Psi_*(\mathcal{F})$ est un homomorphisme de
$\mathcal{B}$-Algèbres, $u^\sharp$ est un homomorphisme
$\Psi^*(\mathcal{G})\to\mathcal{F}$ de $\mathcal{A}$-Algèbres~; cela résulte
de la commutativité du diagramme
\[
\begin{tikzcd}[column sep=huge,row sep=large]
\mathcal{G}\otimes_{\mathcal{B}}\mathcal{G}
  \arrow[r] \arrow[d]
  & \mathcal{G} \arrow[d,"u"] \\
\Psi_*(\mathcal{F}\otimes_{\mathcal{A}}\mathcal{F})
  \arrow[r]
  & \Psi_*(\mathcal{F})
\end{tikzcd}
\]
et de (4.4.4). De même, si $v:\Psi^*(\mathcal{G})\to\mathcal{F}$ est un
homomorphisme de $\mathcal{A}$-Algèbres,
$v^\flat:\mathcal{G}\to\Psi_*(\mathcal{F})$ est un homomorphisme de
$\mathcal{B}$-Algèbres.
\end{env}

\begin{env}[4.4.8]
\label{0.4.4.8-fr}
Soient $(Z,\mathcal{C})$ un troisième espace annelé,
$\Psi'=(\psi',\theta')$ un morphisme $(Y,\mathcal{B})\to(Z,\mathcal{C})$, et
$\Psi'':(X,\mathcal{A})\to(Z,\mathcal{C})$ le morphisme composé
$\Psi'\circ\Psi$. Soit $\mathcal{H}$ un $\mathcal{C}$-Module, $u'$ un
homomorphisme de $\mathcal{H}$ dans $\mathcal{G}$~; le composé
$v''=v\circ v'$ est par définition l'homomorphisme de $\mathcal{H}$ dans
$\mathcal{F}$ défini par
\[
  \mathcal{H}\xrightarrow{v'}\Psi'_*(\mathcal{G})
  \xrightarrow{\Psi'_*(v)}\Psi'_*(\Psi_*(\mathcal{F}))~;
\]
on vérifie que ${v''}^\sharp$ est l'homomorphisme
\[
  \Psi^*({\Psi'}^*(\mathcal{H}))
  \xrightarrow{\Psi^*({v'}^\sharp)}
  \Psi^*(\mathcal{G})
  \xrightarrow{v^\sharp}\mathcal{F}.
\]
\end{env}

\section{Faisceaux quasi-cohérents et faisceaux cohérents}
\label{section:0.5-fr}

\subsection{Faisceaux quasi-cohérents}
\label{subsection:0.5.1-fr}

\begin{env}[5.1.1]
\label{0.5.1.1-fr}
Soient $(X,\mathcal{O}_X)$ un espace annelé, $\mathcal{F}$ un
$\mathcal{O}_X$-Module. La donnée d'un homomorphisme
$u:\mathcal{O}_X\to\mathcal{F}$ de $\mathcal{O}_X$-Modules équivaut à celle
de la section $s=u(1)\in\Gamma(X,\mathcal{F})$. En effet, lorsque $s$ est
donnée, pour toute section $t\in\Gamma(U,\mathcal{O}_X)$, on a nécessairement
$u(t)=t\cdot(s|U)$~; on dit que $u$ est \emph{défini par la section $s$}. Si
maintenant $I$ est un ensemble d'indices quelconque, considérons le faisceau
somme directe $\mathcal{O}_X^{(I)}$, et pour tout $i\in I$, soit $h_i$
l'injection canonique du $i$-ème facteur dans $\mathcal{O}_X^{(I)}$~; on sait
que $u\mapsto(u\circ h_i)$ est un isomorphisme de
$\operatorname{Hom}_{\mathcal{O}_X}(\mathcal{O}_X^{(I)},\mathcal{F})$ sur le
produit
$\bigl(\operatorname{Hom}_{\mathcal{O}_X}(\mathcal{O}_X,\mathcal{F})\bigr)^I$.
Il y a donc correspondance biunivoque canonique entre les homomorphismes
$u:\mathcal{O}_X^{(I)}\to\mathcal{F}$ et les \emph{familles de sections
$(s_i)_{i\in I}$ de $\mathcal{F}$ au-dessus de $X$}. L'homomorphisme $u$
correspondant à $(s_i)$ applique un élément
$(a_i)\in(\Gamma(U,\mathcal{O}_X))^{(I)}$ sur
$\sum_{i\in I}a_i\cdot(s_i|U)$.

On dit que $\mathcal{F}$ est \emph{engendré par la famille $(s_i)$} si
l'homomorphisme $\mathcal{O}_X^{(I)}\to\mathcal{F}$ défini
\oldpage[0\textsubscript{I}]{45}
par cette famille est \emph{surjectif} (autrement dit, si, pour tout
$x\in X$, $\mathcal{F}_x$ est un $\mathcal{O}_x$-module engendré par les
$(s_i)_x$). On dit que $\mathcal{F}$ est \emph{engendré par ses sections
au-dessus de $X$} s'il est engendré par la famille de toutes ces sections (ou
par une sous-famille)~; autrement dit, s'il existe un homomorphisme surjectif
$\mathcal{O}_X^{(I)}\to\mathcal{F}$ pour un $I$ convenable.

On notera qu'un $\mathcal{O}_X$-Module $\mathcal{F}$ peut être tel qu'il
existe un point $x_0\in X$ pour lequel $\mathcal{F}|U$ n'est pas engendré par
ses sections au-dessus de $U$, \emph{quel que soit le voisinage ouvert $U$ de
$x_0$}~: il suffit de prendre $X=\mathbb{R}$, pour $\mathcal{O}_X$ le faisceau
simple $\mathbb{Z}$, pour $\mathcal{F}$ le sous-faisceau algébrique de
$\mathcal{O}_X$ tel que $\mathcal{F}_0=\{0\}$, $\mathcal{F}_x=\mathbb{Z}$ pour
$x\ne0$, et enfin $x_0=0$~: la seule section de $\mathcal{F}|U$ au-dessus de
$U$ est $0$ pour un voisinage $U$ de $0$.
\end{env}

\begin{env}[5.1.2]
\label{0.5.1.2-fr}
Soit $f:X\to Y$ un morphisme d'espaces annelés. Si $\mathcal{F}$ est un
$\mathcal{O}_X$-Module engendré par ses sections au-dessus de $X$, alors
l'homomorphisme canonique
$f^*(f_*(\mathcal{F}))\to\mathcal{F}$ (4.4.3.3) est \emph{surjectif}~: en
effet, avec les notations de (5.1.1), $s_i\otimes1$ est une section de
$f^*(f_*(\mathcal{F}))$ au-dessus de $X$, et son image dans $\mathcal{F}$ est
$s_i$. L'exemple de (5.1.1) où $f$ est l'identité, montre que la réciproque de
cette proposition est inexacte en général.
\end{env}

\begin{env}[5.1.3]
\label{0.5.1.3-fr}
On dit qu'un $\mathcal{O}_X$-Module $\mathcal{F}$ est \emph{quasi-cohérent}
si, pour tout $x\in X$, il y a un voisinage ouvert $U$ de $x$ tel que
$\mathcal{F}|V$ soit isomorphe au \emph{conoyau} d'un homomorphisme de la
forme
$\mathcal{O}_X^{(I)}|V\to\mathcal{O}_X^{(J)}|V$, où $I$ et $J$ soient des
ensembles d'indices arbitraires. Il est clair que $\mathcal{O}_X$ lui-même est
un $\mathcal{O}_X$-Module quasi-cohérent, et que toute somme directe de
$\mathcal{O}_X$-Modules quasi-cohérents est un $\mathcal{O}_X$-Module
quasi-cohérent. On dit qu'une $\mathcal{O}_X$-Algèbre $\mathcal{A}$ est
\emph{quasi-cohérente} si elle l'est en tant que $\mathcal{O}_X$-Module.
\end{env}

\begin{env}[5.1.4]
\label{0.5.1.4-fr}
Soit $f:X\to Y$ un morphisme d'espaces annelés. Si $\mathcal{G}$ est un
$\mathcal{O}_Y$-Module quasi-cohérent, alors $f^*(\mathcal{G})$ est un
$\mathcal{O}_X$-Module quasi-cohérent. En effet, pour tout $x\in X$, il y a
un voisinage ouvert $V$ de $f(x)$ dans $Y$ tel que $\mathcal{G}|V$ soit le
conoyau d'un homomorphisme
$\mathcal{O}_Y^{(I)}|V\to\mathcal{O}_Y^{(J)}|V$. Si $U=f^{-1}(V)$, et si
$f_U$ est la restriction de $f$ à $U$, on a
$f^*(\mathcal{G})|U=f_U^*(\mathcal{G}|V)$~; comme $f_U^*$ est exact à droite
et commute aux sommes directes, $f_U^*(\mathcal{G}|V)$ est conoyau d'un
homomorphisme
$\mathcal{O}_X^{(I)}|U\to\mathcal{O}_X^{(J)}|U$.
\end{env}

\subsection{Faisceaux de type fini}
\label{subsection:0.5.2-fr}

\begin{env}[5.2.1]
\label{0.5.2.1-fr}
On dit qu'un $\mathcal{O}_X$-Module $\mathcal{F}$ est \emph{de type fini} si,
pour tout $x\in X$, il existe un voisinage ouvert $U$ de $x$ tel que
$\mathcal{F}|U$ soit engendré par une famille \emph{finie} de sections
au-dessus de $U$, ou encore soit isomorphe à un faisceau quotient d'un
faisceau de la forme $(\mathcal{O}_X|U)^p$ où $p$ est fini. Tout faisceau
quotient d'un faisceau de type fini est de type fini, ainsi que toute somme
directe finie et tout produit tensoriel fini de faisceaux de type fini. Un
$\mathcal{O}_X$-Module de type fini n'est pas nécessairement quasi-cohérent,
comme on peut le voir pour le $\mathcal{O}_X$-Module
$\mathcal{O}_X/\mathcal{F}$, où $\mathcal{F}$ est l'exemple de (5.1.1). Si
$\mathcal{F}$ est de type fini, $\mathcal{F}_x$ est un
$\mathcal{O}_x$-module de type fini pour tout $x\in X$, mais l'exemple
(5.1.1) montre que cette condition nécessaire n'est pas en général
suffisante.
\end{env}

\begin{env}[5.2.2]
\label{0.5.2.2-fr}
Soit $\mathcal{F}$ un $\mathcal{O}_X$-Module \emph{de type fini}. Si $s_i$
$(1\leq i\leq n)$ sont des sections de $\mathcal{F}$ au-dessus d'un voisinage
ouvert $U$ d'un point $x\in X$ et si les $(s_i)_x$ engendrent
$\mathcal{F}_x$, il existe un voisinage ouvert $V\subset U$ de $x$ tel que
les $(s_i)_y$ engendrent $\mathcal{F}_y$ pour tout $y\in V$ (FAC, I, 2, 12,
prop.~1). On en conclut en particulier que le support de $\mathcal{F}$ est
\emph{fermé}.

\oldpage[0\textsubscript{I}]{46}

De même, si $u:\mathcal{F}\to\mathcal{G}$ est un homomorphisme tel que
$u_x=0$, alors il existe un voisinage $U$ de $x$ tel que $u_y=0$ pour tout
$y\in U$.
\end{env}

\begin{env}[5.2.3]
\label{0.5.2.3-fr}
Supposons $X$ \emph{quasi-compact}, et soient $\mathcal{F}$, $\mathcal{G}$ deux
$\mathcal{O}_X$-Modules tels que $\mathcal{G}$ soit \emph{de type fini},
$u:\mathcal{F}\to\mathcal{G}$ un homomorphisme \emph{surjectif}. Supposons de
plus que $\mathcal{F}$ soit limite inductive d'un système inductif
$(\mathcal{F}_\lambda)$ de $\mathcal{O}_X$-Modules. Alors il existe un indice
$\mu$ tel que l'homomorphisme $\mathcal{F}_\mu\to\mathcal{G}$ soit
\emph{surjectif}. En effet, pour tout $x\in X$, il existe un système fini de
sections $s_i$ de $\mathcal{G}$ au-dessus d'un voisinage ouvert $U(x)$ de $x$
tel que les $(s_i)_y$ engendrent $\mathcal{G}_y$ pour tout $y\in U(x)$~; il y
a donc un voisinage ouvert $V(x)\subset U(x)$ de $x$ et $n$ sections $t_i$ de
$\mathcal{F}$ au-dessus de $V(x)$ telles que
$s_i|V(x)=u(t_i)$ pour tout $i$~; on peut en outre supposer que les $t_i$ sont
les images canoniques de sections d'un même faisceau
$\mathcal{F}_{\lambda(x)}$ au-dessus de $V(x)$. Couvrons alors $X$ avec un
nombre fini de voisinages $V(x_k)$, et soit $\mu$ un indice supérieur à tous
les $\lambda(x_k)$~; il est clair que cet indice répond à la question.

Supposant toujours $X$ quasi-compact, soit $\mathcal{F}$ un
$\mathcal{O}_X$-Module de type fini engendré par ses sections au-dessus de $X$
(5.1.1)~; alors $\mathcal{F}$ est engendré par une sous-famille \emph{finie}
de ces sections~: il suffit en effet de couvrir $X$ par un nombre fini de
voisinages ouverts $U_k$ tels que, pour chaque $k$, il y ait un nombre fini de
sections $s_{ik}$ de $\mathcal{F}$ au-dessus de $X$ dont les restrictions à
$U_k$ engendrent $\mathcal{F}|U_k$~; il est clair que les $s_{ik}$ engendrent
alors $\mathcal{F}$.
\end{env}

\begin{env}[5.2.4]
\label{0.5.2.4-fr}
Soit $f:X\to Y$ un morphisme d'espaces annelés. Si $\mathcal{G}$ est un
$\mathcal{O}_Y$-Module de type fini, alors $f^*(\mathcal{G})$ est un
$\mathcal{O}_X$-module de type fini. En effet, pour tout $x\in X$, il y a un
voisinage ouvert $V$ de $f(x)$ dans $Y$ et un homomorphisme surjectif
$v:\mathcal{O}_Y^p|V\to\mathcal{G}|V$. Si $U=f^{-1}(V)$ et si $f_U$ est la
restriction de $f$ à $U$, on a
$f^*(\mathcal{G})|U=f_U^*(\mathcal{G}|V)$~; comme $f_U$ est exact à droite
(4.3.1) et commute aux sommes directes (4.3.2), $f_U^*(v)$ est un
homomorphisme surjectif
$\mathcal{O}_X^p|U\to f^*(\mathcal{G})|U$.
\end{env}

\begin{env}[5.2.5]
\label{0.5.2.5-fr}
On dit qu'un $\mathcal{O}_X$-Module $\mathcal{F}$ \emph{admet une présentation
finie} si, pour tout $x\in X$, il existe un voisinage ouvert $U$ de $x$ tel
que $\mathcal{F}|U$ soit isomorphe à un \emph{conoyau d'un
$(\mathcal{O}_X|U)$-homomorphisme}
$\mathcal{O}_X^p|U\to\mathcal{O}_X^q|U$, $p$ et $q$ étant deux entiers $>0$.
Un tel $\mathcal{O}_X$-Module est donc de type fini et quasi-cohérent. Si
$f:X\to Y$ est un morphisme d'espaces annelés, et si $\mathcal{G}$ est un
$\mathcal{O}_Y$-Module admettant une présentation finie,
$f^*(\mathcal{G})$ admet une présentation finie, comme le montre le
raisonnement de (5.1.4).
\end{env}

\begin{env}[5.2.6]
\label{0.5.2.6-fr}
Soit $\mathcal{F}$ un $\mathcal{O}_X$-Module admettant une présentation finie
(5.2.5)~; alors, pour tout $\mathcal{O}_X$-Module $\mathcal{H}$,
l'homomorphisme canonique fonctoriel
\[
  \bigl(\mathcal{H}\!om_{\mathcal{O}_X}(\mathcal{F},\mathcal{H})\bigr)_x
  \longrightarrow
  \operatorname{Hom}_{\mathcal{O}_x}(\mathcal{F}_x,\mathcal{H}_x)
\]
est \emph{bijectif} (T, 4.1.1).
\end{env}

\begin{env}[5.2.7]
\label{0.5.2.7-fr}
Soient $\mathcal{F}$, $\mathcal{G}$ deux $\mathcal{O}_X$-Modules ayant une
présentation finie. Si, pour un $x\in X$, $\mathcal{F}_x$ et $\mathcal{G}_x$
sont deux $\mathcal{O}_x$-modules \emph{isomorphes}, alors il existe un
voisinage ouvert $U$ de $x$ tel que $\mathcal{F}|U$ et $\mathcal{G}|U$ soient
\emph{isomorphes}. En effet, si
$\varphi:\mathcal{F}_x\to\mathcal{G}_x$ et
$\psi:\mathcal{G}_x\to\mathcal{F}_x$ sont deux isomorphismes réciproques, il
existe, d'après (5.2.6), un voisinage ouvert $V$ de $x$ et une section $u$
(resp. $v$) de
$\mathcal{H}\!om_{\mathcal{O}_X}(\mathcal{F},\mathcal{G})$ (resp.
$\mathcal{H}\!om_{\mathcal{O}_X}(\mathcal{G},\mathcal{F})$) au-dessus de $V$
telle

\oldpage[0\textsubscript{I}]{47}

que $u_x=\varphi$ (resp. $v_x=\psi$). Comme $(u\circ v)_x$ et
$(v\circ u)_x$ sont les automorphismes identiques, il existe un voisinage
ouvert $U\subset V$ de $x$ tel que $(u\circ v)|U$ et $(v\circ u)|U$ soient les
automorphismes identiques, d'où la proposition.
\end{env}

\subsection{Faisceaux cohérents}
\label{subsection:0.5.3-fr}

\begin{env}[5.3.1]
\label{0.5.3.1-fr}
On dit qu'un $\mathcal{O}_X$-Module $\mathcal{F}$ est \emph{cohérent} s'il
vérifie les deux conditions suivantes~:
\begin{enumerate}
  \item[$a)$] $\mathcal{F}$ est de type fini.
  \item[$b)$] Pour tout ouvert $U\subset X$, tout entier $n>0$, et tout
    homomorphisme
    $u:\mathcal{O}_X^n|U\to\mathcal{F}|U$, le noyau de $u$ est de type fini.
\end{enumerate}
On notera que ces deux conditions sont de caractère \emph{local}.

Pour la plupart des démonstrations des propriétés des faisceaux cohérents
rappelées dans ce qui suit, cf. (FAC), I, 2.
\end{env}

\begin{env}[5.3.2]
\label{0.5.3.2-fr}
Tout $\mathcal{O}_X$-Module cohérent admet une présentation finie (5.2.5)~; la
réciproque n'est pas nécessairement exacte, car $\mathcal{O}_X$ lui-même n'est
pas nécessairement un $\mathcal{O}_X$-Module cohérent.

Tout sous-$\mathcal{O}_X$-Module \emph{de type fini} d'un
$\mathcal{O}_X$-Module cohérent est cohérent~; toute somme directe \emph{finie}
de $\mathcal{O}_X$-Modules cohérents est un $\mathcal{O}_X$-Module cohérent.
\end{env}

\begin{env}[5.3.3]
\label{0.5.3.3-fr}
Si $0\to\mathcal{F}\to\mathcal{G}\to\mathcal{H}\to0$ est une suite exacte de
$\mathcal{O}_X$-Modules et si deux de ces $\mathcal{O}_X$-Modules sont
cohérents, il en est de même du troisième.
\end{env}

\begin{env}[5.3.4]
\label{0.5.3.4-fr}
Si $\mathcal{F}$ et $\mathcal{G}$ sont deux $\mathcal{O}_X$-Modules cohérents,
$u:\mathcal{F}\to\mathcal{G}$ un homomorphisme, $\operatorname{Im}(u)$,
$\operatorname{Ker}(u)$ et $\operatorname{Coker}(u)$ sont des
$\mathcal{O}_X$-Modules cohérents. En particulier, si $\mathcal{F}$ et
$\mathcal{G}$ sont des sous-$\mathcal{O}_X$-Modules cohérents d'un
$\mathcal{O}_X$-Module cohérent, $\mathcal{F}+\mathcal{G}$ et
$\mathcal{F}\cap\mathcal{G}$ sont cohérents.
\end{env}

\begin{env}[5.3.5]
\label{0.5.3.5-fr}
Si $\mathcal{F}$ et $\mathcal{G}$ sont deux $\mathcal{O}_X$-Modules cohérents,
il en est de même de $\mathcal{F}\otimes_{\mathcal{O}_X}\mathcal{G}$ et de
$\mathcal{H}\!om_{\mathcal{O}_X}(\mathcal{F},\mathcal{G})$.
\end{env}

\begin{env}[5.3.6]
\label{0.5.3.6-fr}
Soient $\mathcal{F}$ un $\mathcal{O}_X$-Module cohérent, $\mathcal{J}$ un
faisceau cohérent d'idéaux de $\mathcal{O}_X$. Alors le
$\mathcal{O}_X$-Module $\mathcal{J}\mathcal{F}$ est cohérent, en tant qu'image
de $\mathcal{J}\otimes_{\mathcal{O}_X}\mathcal{F}$ par l'homomorphisme
canonique
$\mathcal{J}\otimes_{\mathcal{O}_X}\mathcal{F}\to\mathcal{F}$
(5.3.4 et 5.3.5).
\end{env}

\begin{env}[5.3.7]
\label{0.5.3.7-fr}
On dit qu'une $\mathcal{O}_X$-Algèbre $\mathcal{A}$ est \emph{cohérente} si
$\mathcal{A}$ est un $\mathcal{O}_X$-Module cohérent. En particulier,
$\mathcal{O}_X$ est un \emph{faisceau cohérent d'anneaux} si, et seulement si,
pour tout ouvert $U\subset X$ et tout homomorphisme de la forme
$u:\mathcal{O}_X^p|U\to\mathcal{O}_X|U$, le noyau de $u$ est un
$(\mathcal{O}_X|U)$-Module de type fini.

Si $\mathcal{O}_X$ est un faisceau cohérent d'anneaux, tout
$\mathcal{O}_X$-Module $\mathcal{F}$ admettant une présentation finie (5.2.5)
est cohérent, en vertu de (5.3.4).

L'\emph{annulateur} d'un $\mathcal{O}_X$-Module $\mathcal{F}$ est le noyau
$\mathcal{J}$ de l'homomorphisme canonique
$\mathcal{O}_X\to\mathcal{H}\!om_{\mathcal{O}_X}(\mathcal{F},\mathcal{F})$
qui, à toute section $s\in\Gamma(U,\mathcal{O}_X)$ fait correspondre la
multiplication par $s$ dans
$\operatorname{Hom}(\mathcal{F}|U,\mathcal{F}|U)$~; si $\mathcal{O}_X$ est
cohérent et si $\mathcal{F}$ est un $\mathcal{O}_X$-Module cohérent,
$\mathcal{J}$ est cohérent (5.3.4 et 5.3.5) et pour tout $x\in X$,
$\mathcal{J}_x$ est l'annulateur de $\mathcal{F}_x$ (5.2.6).
\end{env}

\oldpage[0\textsubscript{I}]{48}

\begin{env}[5.3.8]
\label{0.5.3.8-fr}
Supposons $\mathcal{O}_X$ cohérent~; soient $\mathcal{F}$ un
$\mathcal{O}_X$-Module cohérent, $x$ un point de $X$, $M$ un sous-module de
type fini de $\mathcal{F}_x$~; il existe alors un voisinage ouvert $U$ de $x$
et un sous-$(\mathcal{O}_X|U)$-Module cohérent $\mathcal{G}$ de
$\mathcal{F}|U$ tel que $\mathcal{G}_x=M$ (T, 4.1, lemme 1).

Ce résultat, joint aux propriétés des sous-$\mathcal{O}_X$-Modules d'un
$\mathcal{O}_X$-Module cohérent, impose des conditions nécessaires aux anneaux
$\mathcal{O}_x$ pour que $\mathcal{O}_X$ soit cohérent. Par exemple (5.3.4),
l'intersection de deux idéaux de type fini de $\mathcal{O}_x$ doit encore être
un idéal de type fini.
\end{env}

\begin{env}[5.3.9]
\label{0.5.3.9-fr}
Supposons $\mathcal{O}_X$ cohérent, et soit $M$ un $\mathcal{O}_x$-module
admettant une présentation finie, donc isomorphe à un conoyau d'un
homomorphisme $\varphi:\mathcal{O}_x^p\to\mathcal{O}_x^q$~; alors il existe un
voisinage ouvert $U$ de $X$ et un $(\mathcal{O}_X|U)$-Module cohérent
$\mathcal{F}$ tel que $\mathcal{F}_x$ soit isomorphe à $M$. En effet, en vertu
de (5.2.6), il existe une section $u$ de
$\mathcal{H}\!om_{\mathcal{O}_X}(\mathcal{O}_X^p,\mathcal{O}_X^q)$ au-dessus
d'un voisinage ouvert $U$ de $x$ telle que $u_x=\varphi$~; le conoyau
$\mathcal{F}$ de l'homomorphisme
$u:\mathcal{O}_X^p|U\to\mathcal{O}_X^q|U$ répond à la question (5.3.4).
\end{env}

\begin{env}[5.3.10]
\label{0.5.3.10-fr}
Supposons $\mathcal{O}_X$ cohérent, et soit $\mathcal{J}$ un faisceau cohérent
d'idéaux dans $\mathcal{O}_X$. Pour qu'un
$(\mathcal{O}_X/\mathcal{J})$-Module $\mathcal{F}$ soit cohérent, il faut et il
suffit qu'il soit cohérent en tant que $\mathcal{O}_X$-Module. En particulier
$\mathcal{O}_X/\mathcal{J}$ est un faisceau cohérent d'anneaux.
\end{env}

\begin{env}[5.3.11]
\label{0.5.3.11-fr}
Soit $f:X\to Y$ un morphisme d'espaces annelés, et supposons
$\mathcal{O}_X$ cohérent~; alors, pour tout $\mathcal{O}_Y$-Module cohérent
$\mathcal{G}$, $f^*(\mathcal{G})$ est un $\mathcal{O}_X$-Module cohérent. En
effet, avec les notations de (5.2.4), on peut supposer que $\mathcal{G}|V$ est
conoyau d'un homomorphisme
$v:\mathcal{O}_Y^q|V\to\mathcal{O}_Y^p|V$~; comme $f_U^*$ est exact à droite,
$f^*(\mathcal{G})|U=f_U^*(\mathcal{G}|V)$ est conoyau de l'homomorphisme
$f_U^*(v):\mathcal{O}_X^q|U\to\mathcal{O}_X^p|U$, d'où notre assertion.
\end{env}

\begin{env}[5.3.12]
\label{0.5.3.12-fr}
Soient $Y$ une partie fermée de $X$, $j:Y\to X$ l'injection canonique,
$\mathcal{O}_Y$ un faisceau d'anneaux sur $Y$, et posons
$\mathcal{O}_X=j_*(\mathcal{O}_Y)$. Pour qu'un $\mathcal{O}_Y$-Module
$\mathcal{G}$ soit de type fini (resp. quasi-cohérent, cohérent), il faut et il
suffit que $j_*(\mathcal{G})$ soit un $\mathcal{O}_X$-Module de type fini
(resp. quasi-cohérent, cohérent).
\end{env}

\subsection{Faisceaux localement libres}
\label{subsection:0.5.4-fr}

\begin{env}[5.4.1]
\label{0.5.4.1-fr}
Soit $X$ un espace annelé. On dit qu'un $\mathcal{O}_X$-Module $\mathcal{F}$
est \emph{localement libre} si, pour tout $x\in X$, il existe un voisinage
ouvert $U$ de $x$ tel que $\mathcal{F}|U$ soit isomorphe à un
$(\mathcal{O}_X|U)$-Module de la forme $\mathcal{O}_X^{(I)}|U$, où $I$ peut
dépendre de $U$. Si pour tout $U$, $I$ est fini, on dit que $\mathcal{F}$ est
\emph{de rang fini}~; si pour tout $U$, $I$ a le même nombre fini d'éléments
$n$, on dit que $\mathcal{F}$ est \emph{de rang $n$}. Un
$\mathcal{O}_X$-Module localement libre de rang 1 est encore appelé
\emph{inversible} (cf. (5.4.3)). Si $\mathcal{F}$ est un
$\mathcal{O}_X$-Module localement libre de rang fini, pour tout $x\in X$,
$\mathcal{F}_x$ est un $\mathcal{O}_x$-module libre de rang fini $n(x)$, et il
existe un voisinage $U$ de $x$ tel que $\mathcal{F}|U$ soit de rang $n(x)$~;
si $X$ est connexe, $n(x)$ est donc \emph{constant}.

Il est clair que tout faisceau localement libre est quasi-cohérent, et si
$\mathcal{O}_X$ est un faisceau cohérent d'anneaux, tout
$\mathcal{O}_X$-Module localement libre de rang fini est cohérent.

Si $\mathcal{L}$ est localement libre,
$\mathcal{L}\otimes_{\mathcal{O}_X}\mathcal{F}$ est un foncteur \emph{exact}
en $\mathcal{F}$ dans la catégorie des $\mathcal{O}_X$-Modules.

Nous aurons surtout à considérer des $\mathcal{O}_X$-Modules localement libres
de rang fini,

\oldpage[0\textsubscript{I}]{49}

et lorsque nous parlerons de faisceaux localement libres sans préciser, il
sera sous-entendu qu'ils sont de \emph{rang fini}.
\end{env}

\begin{env}[5.4.2]
\label{0.5.4.2-fr}
Si $\mathcal{L}$, $\mathcal{F}$ sont deux $\mathcal{O}_X$-Modules, on a un
homomorphisme canonique fonctoriel
\[
  \check{\mathcal{L}}\otimes_{\mathcal{O}_X}\mathcal{F}
  =\mathcal{H}\!om_{\mathcal{O}_X}(\mathcal{L},\mathcal{O}_X)
   \otimes_{\mathcal{O}_X}\mathcal{F}
  \longrightarrow
  \mathcal{H}\!om_{\mathcal{O}_X}(\mathcal{L},\mathcal{F})
  \tag{5.4.2.1}\label{0.5.4.2.1-fr}
\]
défini de la façon suivante~: pour tout ouvert $U$, à tout couple $(u,t)$, où
\[
  u\in\Gamma\bigl(U,
    \mathcal{H}\!om_{\mathcal{O}_X}(\mathcal{L},\mathcal{O}_X)\bigr)
  =\operatorname{Hom}(\mathcal{L}|U,\mathcal{O}_X|U)
  \quad\text{et}\quad
  t\in\Gamma(U,\mathcal{F}),
\]
on fait correspondre l'élément de
$\operatorname{Hom}(\mathcal{L}|U,\mathcal{F}|U)$ qui, pour tout $x\in U$,
fait correspondre à $s_x\in\mathcal{L}_x$ l'élément $u_x(s_x)t_x$ de
$\mathcal{F}_x$. Si $\mathcal{L}$ est \emph{localement libre de rang fini},
cet homomorphisme est \emph{bijectif}~; la propriété étant locale, on peut en
effet se borner au cas où $\mathcal{L}=\mathcal{O}_X^n$~; comme pour tout
$\mathcal{O}_X$-Module $\mathcal{G}$,
$\mathcal{H}\!om_{\mathcal{O}_X}(\mathcal{O}_X^n,\mathcal{G})$ est
canoniquement isomorphe à $\mathcal{G}^n$, on est ramené au cas
$\mathcal{L}=\mathcal{O}_X$, qui est immédiat.
\end{env}

\begin{env}[5.4.3]
\EGAReaderTarget{canon-ref.0.5.4.3-fr}{0.5.4.3-fr}
Si $\mathcal{L}$ est inversible, il en est de même de son dual
$\check{\mathcal{L}}=
\mathcal{H}\!om_{\mathcal{O}_X}(\mathcal{L},\mathcal{O}_X)$, car on se ramène
aussitôt (la question étant locale) au cas $\mathcal{L}=\mathcal{O}_X$. En
outre, on a un isomorphisme canonique
\[
  \mathcal{H}\!om_{\mathcal{O}_X}(\mathcal{L},\mathcal{O}_X)
  \otimes_{\mathcal{O}_X}\mathcal{L}
  \xrightarrow{\sim}
  \mathcal{O}_X
  \tag{5.4.3.1}\label{0.5.4.3.1-fr}
\]
en effet, d'après (5.4.2), il suffit de définir un isomorphisme canonique
$\mathcal{H}\!om_{\mathcal{O}_X}(\mathcal{L},\mathcal{L})
\xrightarrow{\sim}\mathcal{O}_X$. Or, pour \emph{tout}
$\mathcal{O}_X$-Module $\mathcal{F}$, on a un homomorphisme canonique
$\mathcal{O}_X\xrightarrow{\sim}
\mathcal{H}\!om_{\mathcal{O}_X}(\mathcal{F},\mathcal{F})$ (5.3.7). Il reste
à prouver que si $\mathcal{F}=\mathcal{L}$ est inversible, cet homomorphisme
est bijectif, et comme la question est locale, on est ramené au cas
$\mathcal{L}=\mathcal{O}_X$, qui est immédiat.

En raison de ce qui précède, on pose
$\mathcal{L}^{-1}=
\mathcal{H}\!om_{\mathcal{O}_X}(\mathcal{L},\mathcal{O}_X)$, et on dit que
$\mathcal{L}^{-1}$ est l'\emph{inverse} de $\mathcal{L}$. La terminologie de
\og faisceau inversible\fg{} peut se justifier de la façon suivante lorsque $X$ est
réduit à un point et $\mathcal{O}_X$ est un anneau \emph{local} $A$ d'idéal
maximal $\mathfrak m$~; si $M$ et $M'$ sont deux $A$-modules ($M$ étant de
type fini) tels que $M\otimes_A M'$ soit isomorphe à $A$, comme
$(A/\mathfrak m)\otimes_A(M\otimes_A M')$ s'identifie à
$(M/\mathfrak mM)\otimes_{A/\mathfrak m}(M'/\mathfrak mM')$, ce dernier
produit tensoriel d'espaces vectoriels sur le corps $A/\mathfrak m$ est
isomorphe à $A/\mathfrak m$, ce qui exige que $M/\mathfrak mM$ et
$M'/\mathfrak mM'$ soient de dimension 1. Pour tout élément $z\in M$
n'appartenant pas à $\mathfrak mM$, on a donc $M=Az+\mathfrak mM$, ce qui
entraîne $M=Az$ d'après le lemme de Nakayama, $M$ étant de type fini.
D'ailleurs, comme l'annulateur de $z$ annule $M\otimes_A M'$, isomorphe à
$A$, cet annulateur est $\{0\}$, et $M$ est par suite \emph{isomorphe à $A$}.
Dans le cas général, cela montre que si $\mathcal{L}$ est un
$\mathcal{O}_X$-Module de type fini, tel qu'il existe un
$\mathcal{O}_X$-Module $\mathcal{F}$ pour lequel
$\mathcal{L}\otimes_{\mathcal{O}_X}\mathcal{F}$ soit isomorphe à
$\mathcal{O}_X$, et si en outre les anneaux $\mathcal{O}_x$ sont des anneaux
locaux, alors $\mathcal{L}_x$ est un $\mathcal{O}_x$-module isomorphe à
$\mathcal{O}_x$ pour tout $x\in X$. Si $\mathcal{O}_X$ et $\mathcal{L}$ sont
supposés \emph{cohérents}, on en conclut donc que $\mathcal{L}$ est inversible
en vertu de (5.2.7).
\end{env}

\begin{env}[5.4.4]
\label{0.5.4.4-fr}
Si $\mathcal{L}$ et $\mathcal{L}'$ sont deux $\mathcal{O}_X$-Modules
inversibles, il en est de même de
$\mathcal{L}\otimes_{\mathcal{O}_X}\mathcal{L}'$, car la question étant
locale, on peut supposer que $\mathcal{L}=\mathcal{O}_X$, et le résultat est
alors trivial. Pour tout entier $n\geq1$, on désigne par
$\mathcal{L}^{\otimes n}$ le produit tensoriel de $n$ faisceaux

\oldpage[0\textsubscript{I}]{50}

identiques à $\mathcal{L}$~; on pose par convention
$\mathcal{L}^{\otimes0}=\mathcal{O}_X$ et pour $n\geq1$,
$\mathcal{L}^{\otimes(-n)}=(\mathcal{L}^{-1})^{\otimes n}$. Avec ces
notations, il existe alors un \emph{isomorphisme canonique fonctoriel}
\[
  \mathcal{L}^{\otimes m}\otimes_{\mathcal{O}_X}\mathcal{L}^{\otimes n}
  \xrightarrow{\sim}
  \mathcal{L}^{\otimes(n+m)}
  \tag{5.4.4.1}\label{0.5.4.4.1-fr}
\]
quels que soient les entiers rationnels $m,n$~: en effet, en vertu des
définitions, on se ramène aussitôt au cas où $m=-1$, $n=1$, et
l'isomorphisme en question a alors été défini en (5.4.3).
\end{env}

\begin{env}[5.4.5]
\EGAReaderTarget{canon-ref.0.5.4.5-fr}{0.5.4.5-fr}
Soit $f:Y\to X$ un morphisme d'espaces annelés. Si $\mathcal{L}$ est un
$\mathcal{O}_X$-Module localement libre (resp. inversible),
$f^*(\mathcal{L})$ est un $\mathcal{O}_Y$-Module localement libre (resp.
inversible)~: cela résulte aussitôt de ce que les images réciproques de deux
$\mathcal{O}_X$-Modules localement isomorphes sont localement isomorphes, de
ce que $f^*$ commute aux sommes directes finies et de ce que
$f^*(\mathcal{O}_X)=\mathcal{O}_Y$ (4.3.4). En outre, on sait qu'on a un
homomorphisme canonique fonctoriel
\[
  f^*(\check{\mathcal{L}})
  \longrightarrow
  \check{\bigl(f^*(\mathcal{L})\bigr)}
\]
(4.4.6), et lorsque $L$ est localement libre, cet homomorphisme est
\emph{bijectif}~: on est en effet encore ramené au cas où
$\mathcal{L}=\mathcal{O}_X$ qui est trivial. On en conclut que si
$\mathcal{L}$ est inversible, $f^*(\mathcal{L}^{\otimes n})$ s'identifie
canoniquement à $(f^*(\mathcal{L}))^{\otimes n}$ pour tout entier rationnel
$n$.
\end{env}

\begin{env}[5.4.6]
\label{0.5.4.6-fr}
Soit $\mathcal{L}$ un $\mathcal{O}_X$-Module inversible~; on désigne par
$\Gamma_*(X,\mathcal{L})$ ou simplement $\Gamma_*(\mathcal{L})$ le groupe
abélien somme directe
$\bigoplus_{n\in\mathbb{Z}}\Gamma(X,\mathcal{L}^{\otimes n})$~; on le munit
d'une structure d'\emph{anneau gradué}, en faisant correspondre au couple
$(s_n,s_m)$, où $s_n\in\Gamma(X,\mathcal{L}^{\otimes n})$,
$s_m\in\Gamma(X,\mathcal{L}^{\otimes m})$, la section de
$\mathcal{L}^{\otimes(n+m)}$ au-dessus de $X$ qui correspond canoniquement
(5.4.4.1) à la section $s_n\otimes s_m$ de
$\mathcal{L}^{\otimes n}\otimes_{\mathcal{O}_X}\mathcal{L}^{\otimes m}$~;
l'associativité de cette multiplication se vérifie de façon immédiate. Il est
clair que $\Gamma_*(X,\mathcal{L})$ est un foncteur covariant en
$\mathcal{L}$, à valeurs dans la catégorie des anneaux gradués.

Si maintenant $\mathcal{F}$ est un $\mathcal{O}_X$-Module quelconque, on pose
\[
  \Gamma_*(\mathcal{L},\mathcal{F})
  =\bigoplus_{n\in\mathbb{Z}}
   \Gamma\bigl(X,
     \mathcal{F}\otimes_{\mathcal{O}_X}\mathcal{L}^{\otimes n}\bigr).
\]
On munit ce groupe abélien d'une structure de \emph{module gradué} sur
l'anneau gradué $\Gamma_*(\mathcal{L})$ de la façon suivante~: au couple
$(s_n,u_m)$, où $s_n\in\Gamma(X,\mathcal{L}^{\otimes n})$ et
$u_m\in\Gamma(X,
\mathcal{F}\otimes_{\mathcal{O}_X}\mathcal{L}^{\otimes m})$, on fait
correspondre la section de
$\mathcal{F}\otimes_{\mathcal{O}_X}\mathcal{L}^{\otimes(m+n)}$ qui correspond
canoniquement (5.4.4.1) à $s_n\otimes u_m$~; la vérification des axiomes des
modules est immédiate. Pour $X$ et $\mathcal{L}$ fixés,
$\Gamma_*(\mathcal{L},\mathcal{F})$ est un foncteur covariant en
$\mathcal{F}$ à valeurs dans la catégorie des
$\Gamma_*(\mathcal{L})$-modules gradués~; pour $X$ et $\mathcal{F}$ fixés,
c'est un foncteur covariant en $\mathcal{L}$ à valeurs dans la catégorie des
groupes abéliens.

Si $f:Y\to X$ est un morphisme d'espaces annelés, l'homomorphisme canonique
(4.4.3.2)
\[
  \rho:\mathcal{L}^{\otimes n}
  \longrightarrow
  f_*\bigl(f^*(\mathcal{L}^{\otimes n})\bigr)
\]
définit un homomorphisme de groupes abéliens
\[
  \Gamma(X,\mathcal{L}^{\otimes n})
  \longrightarrow
  \Gamma\bigl(Y,f^*(\mathcal{L}^{\otimes n})\bigr),
\]
et comme
$f^*(\mathcal{L}^{\otimes n})=(f^*(\mathcal{L}))^{\otimes n}$, il résulte des
définitions des homomorphismes canoniques (4.4.3.2) et (5.4.4.1) que les
homomorphismes précédents définissent un \emph{homomorphisme fonctoriel
d'anneaux gradués}
$\Gamma_*(\mathcal{L})\to\Gamma_*(f^*(\mathcal{L}))$. Le même homomorphisme
canonique (4.4.3) définit de même un homomorphisme de groupes abéliens
\[
  \Gamma\bigl(X,
    \mathcal{F}\otimes_{\mathcal{O}_X}\mathcal{L}^{\otimes n}\bigr)
  \longrightarrow
  \Gamma\bigl(Y,
    f^*(\mathcal{F}\otimes_{\mathcal{O}_X}\mathcal{L}^{\otimes n})\bigr),
\]
et comme
\[
  f^*(\mathcal{F}\otimes_{\mathcal{O}_X}\mathcal{L}^{\otimes n})
  =f^*(\mathcal{F})\otimes_{\mathcal{O}_Y}
   (f^*(\mathcal{L}))^{\otimes n}
  \tag{4.3.3.1}
\]

\oldpage[0\textsubscript{I}]{51}

ces homomorphismes (pour $n$ variable) définissent un
\emph{di-homomorphisme de modules gradués}
\[
  \Gamma_*(\mathcal{L},\mathcal{F})
  \longrightarrow
  \Gamma_*(f^*(\mathcal{L}),f^*(\mathcal{F})).
\]
\end{env}

\begin{env}[5.4.7]
\EGAReaderTarget{canon-ref.0.5.4.7-fr}{0.5.4.7-fr}
On peut montrer qu'il existe un \emph{ensemble} $\mathfrak{M}$ (noté aussi
$\mathfrak{M}(X)$) de $\mathcal{O}_X$-Modules inversibles tel que tout
$\mathcal{O}_X$-Module inversible soit isomorphe à un élément de
$\mathfrak{M}$ et un seul\footnote{Voir le livre en préparation cité dans
l'Introduction.}~; on définit dans $\mathfrak{M}$ une loi de composition en
faisant correspondre à deux éléments $\mathcal{L}$, $\mathcal{L}'$ de
$\mathfrak{M}$ l'unique élément de $\mathfrak{M}$ isomorphe à
$\mathcal{L}\otimes_{\mathcal{O}_X}\mathcal{L}'$. Avec cette loi de
composition, $\mathfrak{M}$ est un \emph{groupe isomorphe au groupe de
cohomologie} $H^1(X,\mathcal{O}_X^*)$, où $\mathcal{O}_X^*$ est le
sous-faisceau de $\mathcal{O}_X$ tel que $\Gamma(U,\mathcal{O}_X^*)$ soit le
groupe des éléments inversibles de l'anneau $\Gamma(U,\mathcal{O}_X)$ pour
tout ouvert $U\subset X$ ($\mathcal{O}_X^*$ est donc un faisceau de groupes
abéliens \emph{multiplicatifs}).

On notera pour cela que pour tout ouvert $U\subset X$, le groupe des sections
$\Gamma(U,\mathcal{O}_X^*)$ s'identifie canoniquement au \emph{groupe des
automorphismes} du $(\mathcal{O}_X|U)$-Module $\mathcal{O}_X|U$,
l'identification faisant correspondre à une section $\varepsilon$ de
$\mathcal{O}_X^*$ au-dessus de $U$ l'automorphisme $u$ de $\mathcal{O}_X|U$
tel que $u_x(s_x)=\varepsilon_xs_x$ pour tout $x\in X$ et tout
$s_x\in\mathcal{O}_x$. Soit alors $\mathfrak{U}=(U_\lambda)$ un recouvrement
ouvert de $X$~; la donnée, pour tout couple d'indices $(\lambda,\mu)$ d'un
automorphisme $\theta_{\lambda\mu}$ de
$\mathcal{O}_X|(U_\lambda\cap U_\mu)$ revient à se donner une
$1$-\emph{cochaîne} du recouvrement $\mathfrak{U}$, à valeurs dans
$\mathcal{O}_X^*$, et dire que les $\theta_{\lambda\mu}$ vérifient la
condition de recollement (3.3.1) signifie que la cochaîne correspondante est
un \emph{cocycle}. De même, la donnée, pour tout $\lambda$, d'un automorphisme
$\omega_\lambda$ de $\mathcal{O}_X|U_\lambda$ revient à la donnée d'une
$0$-cochaîne du recouvrement $\mathfrak{U}$, à valeurs dans
$\mathcal{O}_X^*$, et son \emph{cobord} correspond à la famille des
automorphismes
$(\omega_\lambda|U_\lambda\cap U_\mu)\circ
 (\omega_\mu|U_\lambda\cap U_\mu)^{-1}$.

On peut faire correspondre à tout $1$-cocycle de $\mathfrak{U}$ à valeurs
dans $\mathcal{O}_X^*$ l'élément de $\mathfrak{M}$ isomorphe au
$\mathcal{O}_X$-Module inversible obtenu par recollement à partir de la
famille d'automorphismes $(\theta_{\lambda\mu})$ correspondant à ce cocycle,
et à deux cocycles cohomologues correspondent deux éléments égaux de
$\mathfrak{M}$ (3.3.2)~; autrement dit, on a ainsi défini une application
$\varphi_{\mathfrak{U}}:H^1(\mathfrak{U},\mathcal{O}_X^*)\to\mathfrak{M}$.
En outre, si $\mathfrak{B}$ est un second recouvrement ouvert de $X$, plus fin
que $\mathfrak{U}$, le diagramme
\[
\begin{tikzcd}[column sep=huge,row sep=large]
H^1(\mathfrak{U},\mathcal{O}_X^*)
  \arrow[dd]
  \arrow[rrd,"\varphi_{\mathfrak{U}}"] & & \\
& & \mathfrak{M} \\
H^1(\mathfrak{B},\mathcal{O}_X^*)
  \arrow[rru,swap,"\varphi_{\mathfrak{B}}"] & &
\end{tikzcd}
\]
où la flèche verticale est l'homomorphisme canonique (G, II, 5.7), est
commutatif, comme il résulte de (3.3.3). Par passage à la limite inductive, on
obtient donc bien une application
$H^1(X,\mathcal{O}_X^*)\to\mathfrak{M}$, le groupe de cohomologie de \v{C}ech
$\check{H}^1(X,\mathcal{O}_X^*)$ s'identifiant comme on sait au premier groupe
de cohomologie $H^1(X,\mathcal{O}_X^*)$ (G, II, 5.9, cor. du th.~5.9.1).
Cette application est \emph{surjective}~: en effet, par définition, pour tout
$\mathcal{O}_X$-Module inversible $\mathcal{L}$, il y a un recouvrement ouvert
$(U_\lambda)$ de $X$ tel que $\mathcal{L}$ s'obtienne par recollement des
faisceaux $\mathcal{O}_X|U_\lambda$ (3.3.1). Elle est aussi
\emph{injective}, car il suffit de le prouver pour les applications
$H^1(\mathfrak{U},\mathcal{O}_X)\to\mathfrak{M}$, et cela résulte alors de
(3.3.2). Il reste à montrer que

\oldpage[0\textsubscript{I}]{52}

la bijection ainsi définie est un homomorphisme de groupe. Or, étant donnés
deux $\mathcal{O}_X$-Modules inversibles $\mathcal{L}$, $\mathcal{L}'$, il y a
un recouvrement ouvert $(U_\lambda)$ tel que $\mathcal{L}|U_\lambda$ et
$\mathcal{L}'|U_\lambda$ soient isomorphes à $\mathcal{O}_X|U_\lambda$ pour
tout $\lambda$~; il y a donc pour chaque indice $\lambda$ un élément
$a_\lambda$ (resp. $a'_\lambda$) de $\Gamma(U_\lambda,\mathcal{L})$ (resp.
$\Gamma(U_\lambda,\mathcal{L}'))$ tel que les éléments de
$\Gamma(U_\lambda,\mathcal{L})$ (resp. $\Gamma(U_\lambda,\mathcal{L}'))$
soient les $s_\lambda\cdot a_\lambda$ (resp.
$s_\lambda\cdot a'_\lambda$), où $s_\lambda$ parcourt
$\Gamma(U_\lambda,\mathcal{O}_X)$. Les cocycles correspondants
$(\varepsilon_{\lambda\mu})$, $(\varepsilon'_{\lambda\mu})$ sont tels que
$s_\lambda\cdot a_\lambda=s_\mu\cdot a_\mu$ (resp.
$s_\lambda\cdot a'_\lambda=s_\mu\cdot a'_\mu$) au-dessus de
$U_\lambda\cap U_\mu$ soit équivalent à
$s_\lambda=\varepsilon_{\lambda\mu}s_\mu$ (resp.
$s_\lambda=\varepsilon'_{\lambda\mu}s_\mu$) au-dessus de
$U_\lambda\cap U_\mu$. Comme les sections de
$\mathcal{L}\otimes_{\mathcal{O}_X}\mathcal{L}'$ au-dessus de $U_\lambda$
sont les sommes finies des
$s_\lambda s'_\lambda\cdot(a_\lambda\otimes a'_\lambda)$ où $s_\lambda$ et
$s'_\lambda$ parcourent $\Gamma(U_\lambda,\mathcal{O}_X)$, il est clair que
le cocycle $(\varepsilon_{\lambda\mu}\varepsilon'_{\lambda\mu})$ correspond
à $\mathcal{L}\otimes_{\mathcal{O}_X}\mathcal{L}'$, ce qui achève la
démonstration\footnote{Pour une forme générale de ce résultat, voir le livre
cité dans la note de la p.~51.}.
\end{env}

\begin{env}[5.4.8]
\label{0.5.4.8-fr}
Soit $f=(\psi,\omega)$ un morphisme $Y\to X$ d'espaces annelés. Le foncteur
$f^*(\mathcal{L})$ dans la catégorie des $\mathcal{O}_X$-Modules libres
définit une application (encore notée $f^*$ par abus de langage) de l'ensemble
$\mathfrak{M}(X)$ dans l'ensemble $\mathfrak{M}(Y)$. D'autre part, on a un
homomorphisme canonique (T, 3.2.2)
\[
  H^1(X,\mathcal{O}_X^*)\longrightarrow H^1(Y,\mathcal{O}_Y^*).
  \tag{5.4.8.1}\label{0.5.4.8.1-fr}
\]
Lorsqu'on identifie canoniquement (5.4.7) $\mathfrak{M}(X)$ et
$H^1(X,\mathcal{O}_X^*)$ (resp. $\mathfrak{M}(Y)$ et
$H^1(Y,\mathcal{O}_Y^*)$), l'homomorphisme (5.4.8.1) \emph{s'identifie à
l'application $f^*$}. En effet, si $\mathcal{L}$ provient d'un cocycle
$(\varepsilon_{\lambda\mu})$ correspondant à un recouvrement ouvert
$(U_\lambda)$ de $X$, il suffit de montrer que $f^*(\mathcal{L})$ provient
d'un cocycle dont la classe de cohomologie est image par (5.4.8.1) de celle de
$(\varepsilon_{\lambda\mu})$. Or, si $\theta_{\lambda\mu}$ est
l'automorphisme de $\mathcal{O}_X|(U_\lambda\cap U_\mu)$ qui correspond à
$\varepsilon_{\lambda\mu}$, il est clair que $f^*(\mathcal{L})$ s'obtient par
recollement des $\mathcal{O}_Y|\psi^{-1}(U_\lambda)$ au moyen des
automorphismes $f^*(\theta_{\lambda\mu})$, et il suffit donc de vérifier que
ces derniers correspondent au cocycle
$(\omega^\#(\varepsilon_{\lambda\mu}))$, ce qui résulte aussitôt des
définitions (on a identifié $\varepsilon_{\lambda\mu}$ à son image canonique
par $\rho$ (3.7.2), section de $\psi^*(\mathcal{O}_X^*)$ au-dessus de
$\psi^{-1}(U_\lambda\cap U_\mu)$).
\end{env}

\begin{env}[5.4.9]
\label{0.5.4.9-fr}
Soient $\mathcal{E}$, $\mathcal{F}$ deux $\mathcal{O}_X$-Modules,
$\mathcal{F}$ étant supposé \emph{localement libre}, et soit $\mathcal{G}$ un
$\mathcal{O}_X$-Module \emph{extension de $\mathcal{F}$ par $\mathcal{E}$},
autrement dit tel qu'il existe une suite exacte
\[
  0\longrightarrow\mathcal{E}\xrightarrow{i}\mathcal{G}
  \xrightarrow{p}\mathcal{F}\longrightarrow 0.
\]
Alors, pour tout $x\in X$, il existe un voisinage ouvert $U$ de $x$ tel que
$\mathcal{G}|U$ soit isomorphe à la \emph{somme directe}
$\mathcal{E}|U\oplus\mathcal{F}|U$. On peut en effet se borner au cas où
$\mathcal{F}=\mathcal{O}_X^n$~; soient $e_i$ ($1\leq i\leq n$) les sections
canoniques (5.5.5) de $\mathcal{O}_X^n$~; il existe alors un voisinage ouvert
$U$ de $x$ et $n$ sections $s_i$ de $\mathcal{G}$ au-dessus de $U$ telles que
$p(s_i|U)=e_i|U$ pour $1\leq i\leq n$. Cela étant, soit $f$
l'homomorphisme $\mathcal{F}|U\to\mathcal{G}|U$ défini par les sections
$s_i|U$ (5.1.1). Il est immédiat que pour tout ouvert $V\subset U$, et toute
section $s\in\Gamma(V,\mathcal{G})$ on a
$s-f(p(s))\in\Gamma(V,\mathcal{E})$, d'où notre assertion.
\end{env}

\begin{env}[5.4.10]
\label{0.5.4.10-fr}
Soient $f:X\to Y$ un morphisme d'espaces annelés, $\mathcal{F}$ un
$\mathcal{O}_X$-Module, $\mathcal{L}$ un $\mathcal{O}_Y$-Module localement
libre de rang fini. Il existe alors un isomorphisme canonique
\[
  f_*(\mathcal{F})\otimes_{\mathcal{O}_Y}\mathcal{L}
  \xrightarrow{\sim}
  f_*\bigl(\mathcal{F}\otimes_{\mathcal{O}_X}f^*(\mathcal{L})\bigr).
  \tag{5.4.10.1}\label{0.5.4.10.1-fr}
\]

\oldpage[0\textsubscript{I}]{53}

En effet, pour tout $\mathcal{O}_Y$-Module $\mathcal{L}$, on a un
homomorphisme canonique
\[
  f_*(\mathcal{F})\otimes_{\mathcal{O}_Y}\mathcal{L}
  \xrightarrow{1\otimes\rho}
  f_*(\mathcal{F})\otimes_{\mathcal{O}_Y}f_*(f^*(\mathcal{L}))
  \xrightarrow{\alpha}
  f_*\bigl(\mathcal{F}\otimes_{\mathcal{O}_X}f^*(\mathcal{L})\bigr),
\]
$\rho$ étant l'homomorphisme (4.4.3.2) et $\alpha$ l'homomorphisme (4.2.2.1).
Pour montrer que lorsque $\mathcal{L}$ est localement libre, cet homomorphisme
est bijectif, il suffit, la question étant locale sur $Y$, de considérer le
cas où $\mathcal{L}=\mathcal{O}_Y^n$~; en outre, $f_*$ et $f^*$ commutant aux
sommes directes finies, on peut supposer $n=1$, et dans ce cas la proposition
résulte aussitôt des définitions et de la relation
$f^*(\mathcal{O}_Y)=\mathcal{O}_X$.
\end{env}

\subsection{Faisceaux sur un espace annelé en anneaux locaux}
\label{subsection:0.5.5-fr}

\begin{env}[5.5.1]
\label{0.5.5.1-fr}
Nous dirons qu'un espace annelé $(X,\mathcal{O}_X)$ est un \emph{espace
annelé en anneaux locaux} si, pour tout $x\in X$, $\mathcal{O}_x$ est un
anneau local~; ces espaces annelés seront de loin les plus fréquents de ceux
qui seront considérés dans ce travail. Nous désignerons alors par
$\mathfrak{m}_x$, l'\emph{idéal maximal} de $\mathcal{O}_x$, par $k(x)$ le
\emph{corps résiduel} $\mathcal{O}_x/\mathfrak{m}_x$~; pour tout
$\mathcal{O}_X$-Module $\mathcal{F}$, tout ouvert $U$ de $X$, tout point
$x\in U$ et toute section $f\in\Gamma(U,\mathcal{F})$, nous désignerons par
$f(x)$ la \emph{classe} du germe
$f_x\in\mathcal{F}_x\pmod{\mathfrak{m}_x\mathcal{F}_x}$, et nous dirons que
c'est la \emph{valeur} de $f$ au point $x$. La relation $f(x)=0$ signifie
donc que $f_x\in\mathfrak{m}_x\mathcal{F}_x$~; lorsqu'elle est remplie, on dira
(par abus de langage) que $f$ \emph{s'annule en $x$}. On aura soin de ne pas
confondre cette relation avec $f_x=0$.
\end{env}

\begin{env}[5.5.2]
\label{0.5.5.2-fr}
Soient $X$ un espace annelé en anneaux locaux, $\mathcal{L}$ un
$\mathcal{O}_X$-Module inversible, $f$ une section de $\mathcal{L}$ au-dessus
de $X$. Il y a alors \emph{équivalence} entre les trois propriétés suivantes
en un point $x\in X$~:
\begin{enumerate}
  \item[$a)$] \emph{$f_x$ est un générateur de $\mathcal{L}_x$}~;
  \item[$b)$] \emph{$f_x\notin\mathfrak{m}_x\mathcal{L}_x$} (autrement dit,
    $f(x)\neq0$)~;
  \item[$c)$] \emph{Il existe une section $g$ de $\mathcal{L}^{-1}$ au-dessus
    d'un voisinage ouvert $V$ de $x$ telle que l'image canonique de $f\otimes
    g$ dans $\Gamma(V,\mathcal{O}_X)$ (5.4.3) soit la section unité}.
\end{enumerate}

En effet, la question étant locale, on peut se borner au cas où
$\mathcal{L}=\mathcal{O}_X$~; l'équivalence de $a)$ et $b)$ est alors
évidente, et il est clair que $c)$ entraîne $b)$. D'autre part, si
$f_x\notin\mathfrak{m}_x$, $f_x$ est inversible dans $\mathcal{O}_x$, soit
$f_xg_x=1_x$. Par définition des germes de sections, cela veut dire qu'il
existe un voisinage $V$ de $x$ et une section $g$ de $\mathcal{O}_X$ au-dessus
de $V$ telle que $fg=1$ dans $V$, d'où $c)$.

Il résulte aussitôt de la condition $c)$ que l'ensemble $X_f$ des $x$
vérifiant les conditions équivalentes $a)$, $b)$, $c)$ est \emph{ouvert} dans
$X$~; suivant la terminologie introduite dans (5.5.1), c'est l'ensemble des
$x$ où $f$ \emph{ne s'annule pas}.
\end{env}

\begin{env}[5.5.3]
\label{0.5.5.3-fr}
Sous les hypothèses de (5.5.2), soit $\mathcal{L}'$ un second
$\mathcal{O}_X$-Module inversible~; alors, si
$f\in\Gamma(X,\mathcal{L})$, $g\in\Gamma(X,\mathcal{L}')$, on a
\[
  X_f\cap X_g=X_{f\otimes g}.
\]
On se ramène en effet aussitôt au cas où
$\mathcal{L}=\mathcal{L}'=\mathcal{O}_X$ (la question étant locale)~; comme
$f\otimes g$ s'identifie alors canoniquement au produit $fg$, la proposition
est évidente.
\end{env}

\oldpage[0\textsubscript{I}]{54}

\begin{env}[5.5.4]
\label{0.5.5.4-fr}
Soit $\mathcal{F}$ un $\mathcal{O}_X$-Module localement libre de rang $n$~;
il est immédiat que $\bigwedge^p\mathcal{F}$ est un $\mathcal{O}_X$-Module
localement libre de rang $\binom{n}{p}$ si $p\leq n$, réduit à $0$ si $p>n$,
car la question est locale et on est ramené au cas où
$\mathcal{F}=\mathcal{O}_X^n$~; en outre, pour tout $x\in X$,
\[
  (\bigwedge^p\mathcal{F})_x/
  \mathfrak{m}_x(\bigwedge^p\mathcal{F})_x
\]
est un espace vectoriel de dimension $\binom{n}{p}$ sur $k(x)$, qui
s'identifie canoniquement à
\[
  \bigwedge^p(\mathcal{F}_x/\mathfrak{m}_x\mathcal{F}_x).
\]
Soient
$s_1,\ldots,s_p$ des sections de $F$ au-dessus d'un ouvert $U$ de $X$, et
soit $s=s_1\wedge\cdots\wedge s_p$, qui est une section de
$\bigwedge^p\mathcal{F}$ au-dessus de $U$ (4.1.5)~; on a
$s(x)=s_1(x)\wedge\cdots\wedge s_p(x)$, et par suite, dire que
$s_1(x),\ldots,s_p(x)$ sont \emph{linéairement dépendants} signifie que
$s(x)=0$. On en conclut que \emph{l'ensemble des $x\in X$ tels que
$s_1(x),\ldots,s_p(x)$ soient linéairement indépendants est ouvert dans
$X$}~: il suffit en effet, en se ramenant au cas où
$\mathcal{F}=\mathcal{O}_X^n$, d'appliquer (5.5.2) à la section image de $s$
par une des projections de
$\bigwedge^p\mathcal{F}=\mathcal{O}_X^{\binom{n}{p}}$ sur les
$\binom{n}{p}$ facteurs.

En particulier, si $s_1,\ldots,s_n$ sont $n$ sections de $\mathcal{F}$
au-dessus de $U$ telles que $s_1(x),\ldots,s_n(x)$ soient linéairement
indépendantes en tout point $x\in U$, l'homomorphisme
$u:\mathcal{O}_X^n|U\to\mathcal{F}|U$ défini par les $s_i$ (5.1.1) est un
\emph{isomorphisme}~: on peut en effet se restreindre au cas où
$\mathcal{F}=\mathcal{O}_X^n$ et où on identifie canoniquement
$\bigwedge^n\mathcal{F}$ et $\mathcal{O}_X$~;
$s=s_1\wedge\cdots\wedge s_n$ est alors une section de $\mathcal{O}_X$
\emph{inversible} au-dessus de $U$, et on définit un homomorphisme réciproque
de $u$ au moyen des formules de Cramer.
\end{env}

\begin{env}[5.5.5]
\label{0.5.5.5-fr}
Soient $\mathcal{E}$, $\mathcal{F}$ deux $\mathcal{O}_X$-Modules localement
libres (de rang fini), et soit $u:\mathcal{E}\to\mathcal{F}$ un
homomorphisme. Pour qu'il existe un voisinage $u$ de $x\in X$ tel que $u|U$
soit \emph{injectif} et que $\mathcal{F}|U$ soit \emph{somme directe de
$u(\mathcal{E})|U$ et d'un sous-$(\mathcal{O}_X|U)$-Module localement libre
$\mathcal{G}$}, il faut et il suffit que
$u_x:\mathcal{E}_x\to\mathcal{F}_x$ donne, par passage aux quotients, un
homomorphisme \emph{injectif} d'espaces vectoriels
\[
  \mathcal{E}_x/\mathfrak{m}_x\mathcal{E}_x
  \longrightarrow
  \mathcal{F}_x/\mathfrak{m}_x\mathcal{F}_x.
\]
La condition est en effet \emph{nécessaire}, car $\mathcal{F}_x$ est alors
somme directe des $\mathcal{O}_x$-modules libres $u_x(\mathcal{E}_x)$ et
$\mathcal{G}_x$, donc $\mathcal{F}_x/\mathfrak{m}_x\mathcal{F}_x$ est somme
directe de
$u_x(\mathcal{E}_x)/\mathfrak{m}_x u_x(\mathcal{E}_x)$ et de
$\mathcal{G}_x/\mathfrak{m}_x\mathcal{G}_x$. La condition est
\emph{suffisante}, car on peut se borner au cas où
$\mathcal{E}=\mathcal{O}_X^m$~; soient $s_1,\ldots,s_m$ les images par $u$ des
sections $e_i$ de $\mathcal{O}_X^m$ telles que $(e_i)_y$ soit égal au $i$-ème
élément de la base canonique de $\mathcal{O}_y^m$ pour tout $y\in X$
(\emph{sections canoniques de $\mathcal{O}_X^m$})~; par hypothèse
$s_1(x),\ldots,s_m(x)$ sont linéairement indépendants, donc, si $\mathcal{F}$
est de rang $n$, il existe $n-m$ sections $s_{m+1},\ldots,s_n$ de
$\mathcal{F}$ au-dessus d'un voisinage $V$ de $x$ telles que les $s_i(x)$
$(1\leq i\leq n)$ forment une base de
$\mathcal{F}_x/\mathfrak{m}_x\mathcal{F}_x$. Il existe alors (5.5.4) un
voisinage $U\subset V$ de $x$ tel que les $s_i(y)$ $(1\leq i\leq n)$ forment
une base de $\mathcal{F}_y/\mathfrak{m}_y\mathcal{F}_y$ pour tout $y\in V$,
et on en conclut (5.5.4) qu'il y a un isomorphisme de $\mathcal{F}|U$ sur
$\mathcal{O}_X^n|U$, appliquant les $s_i|U$ $(1\leq i\leq m)$ sur les
$e_i|U$, ce qui achève la démonstration.
\end{env}

\section{Platitude}
\label{section:0.6-fr}

\begin{env}[6.0]
\label{0.6.0.0-fr}
La notion de platitude est due à J.-P.~Serre [16]~; dans ce qui suit, on omet
les démonstrations des résultats qui sont exposés dans l'\emph{Algèbre
commutative} de N.~Bourbaki, à laquelle nous renvoyons le lecteur. Nous
supposons tous les anneaux commutatifs.\footnote{Voir l'exposé cité de
N.~Bourbaki pour la généralisation de la plupart des résultats au cas non
commutatif.}

\oldpage[0\textsubscript{I}]{55}

Si $M$, $N$ sont deux $A$-modules, $M'$ (resp. $N'$) un sous-module de $M$
(resp. $N$), nous noterons $\operatorname{Im}(M'\otimes_A N')$ le sous-module
de $M\otimes_A N$, image de l'application canonique
$M'\otimes_A N'\to M\otimes_A N$.
\end{env}

\subsection{Modules plats}
\label{subsection:0.6.1-fr}

\begin{env}[6.1.1]
\label{0.6.1.1-fr}
Soit $M$ un $A$-module. Les propriétés suivantes sont équivalentes~:
\begin{enumerate}
  \item[$a)$] Le foncteur $M\otimes_A N$ en $N$ est exact dans la catégorie
    des $A$-modules~;
  \item[$b)$] $\operatorname{Tor}_i^A(M,N)=0$ pour tout $i>0$ et tout
    $A$-module $N$~;
  \item[$c)$] $\operatorname{Tor}_1^A(M,N)=0$ pour tout $A$-module $N$.
\end{enumerate}

Lorsque $M$ vérifie ces conditions, on dit que $M$ est un
\emph{$A$-module plat}. Il est clair que tout $A$-module libre est plat.

Pour que $M$ soit un $A$-module plat, il suffit que pour tout idéal
$\mathfrak{J}$ de $A$, \emph{de type fini}, l'application canonique
$M\otimes_A\mathfrak{J}\to M\otimes_A A=M$ soit \emph{injective}.
\end{env}

\begin{env}[6.1.2]
\EGAReaderTarget{canon-ref.0.6.1.2-fr}{0.6.1.2-fr}
Toute limite inductive de $A$-modules plats est un $A$-module plat. Pour
qu'une somme directe $\bigoplus_{\lambda\in L}M_\lambda$ de $A$-modules soit
un $A$-module plat, il faut et il suffit que chacun des $A$-modules
$M_\lambda$ soit plat. En particulier, tout $A$-module projectif est plat.

Soit
\[
  0\longrightarrow M'\longrightarrow M\longrightarrow M''\longrightarrow 0
\]
une suite exacte de $A$-modules, telle que $M''$ soit \emph{plat}. Alors, pour
tout $A$-module $N$, la suite
\[
  0\longrightarrow M'\otimes N\longrightarrow M\otimes N
  \longrightarrow M''\otimes N\longrightarrow 0
\]
est exacte. En outre, pour que $M$ soit plat, il faut et il suffit alors que
$M'$ le soit (mais il peut se faire que $M$ et $M'$ soient plats sans que
$M''=M/M'$ le soit).
\end{env}

\begin{env}[6.1.3]
\label{0.6.1.3-fr}
Soient $M$ un $A$-module plat, $N$ un $A$-module quelconque~; pour deux
sous-modules $N'$, $N''$ de $N$, on a alors
\[
  \operatorname{Im}(M\otimes(N'+N''))
  =\operatorname{Im}(M\otimes N')+\operatorname{Im}(M\otimes N''),
\]
\[
  \operatorname{Im}(M\otimes(N'\cap N''))
  =\operatorname{Im}(M\otimes N')\cap\operatorname{Im}(M\otimes N'')
\]
(images prises dans $M\otimes N$).
\end{env}

\begin{env}[6.1.4]
\label{0.6.1.4-fr}
Soient $M$, $N$ deux $A$-modules, $M'$ (resp. $N'$) un sous-module de $M$
(resp. $N$), et supposons que l'un des modules $M/M'$, $N/N'$ soit plat.
Alors on a
\[
  \operatorname{Im}(M'\otimes N')
  =\operatorname{Im}(M'\otimes N)\cap\operatorname{Im}(M\otimes N')
\]
(images dans $M\otimes N$). En particulier, si $\mathfrak{J}$ est un idéal
de $A$ et si $M/M'$ est plat, on a
$\mathfrak{J}M'=M'\cap\mathfrak{J}M$.
\end{env}

\subsection{Changement d'anneaux}
\label{subsection:0.6.2-fr}

Lorsqu'un groupe additif $M$ est muni de plusieurs structures de module par
rapport à des anneaux $A$, $B$, \ldots, au lieu de dire que $M$ est plat en
tant que $A$-module, $B$-module, \ldots, nous dirons parfois aussi que $M$
est \emph{$A$-plat}, \emph{$B$-plat}, \ldots

\begin{env}[6.2.1]
\label{0.6.2.1-fr}
Soient $A$ et $B$ deux anneaux, $M$ un $A$-module, $N$ un
$(A,B)$-bimodule. Si $M$ est plat et si $N$ est $B$-plat, alors
$M\otimes_A N$ est $B$-plat. En particulier, si $M$ et $N$ sont deux
$A$-modules plats, $M\otimes_A N$ est un $A$-module plat. Si $B$ est une
$A$-algèbre et si $M$ est
\oldpage[0\textsubscript{I}]{56}
un $A$-module plat, le $B$-module $M_{(B)}=M\otimes_A B$ est plat. Enfin, si
$B$ est une $A$-algèbre qui est plate en tant que $A$-module, et si $N$ est
un $B$-module plat, alors $N$ est aussi $A$-plat.
\end{env}

\begin{env}[6.2.2]
\label{0.6.2.2-fr}
Soient $A$ un anneau, $B$ une $A$-algèbre plate en tant que $A$-module.
Soient $M$, $N$ deux $A$-modules, tels que $M$ admette une présentation
finie~; alors l'homomorphisme canonique
\[
  \operatorname{Hom}_A(M,N)\otimes_A B
  \longrightarrow
  \operatorname{Hom}_B(M\otimes_A B,N\otimes_A B)
  \tag{6.2.2.1}
\]
(transformant $u\otimes b$ en l'homomorphisme
$m\otimes b'\longmapsto u(m)\otimes b'b$) est un isomorphisme.
\end{env}

\begin{env}[6.2.3]
\label{0.6.2.3-fr}
Soit $(A_\lambda,\varphi_{\mu\lambda})$ un système inductif filtrant
d'anneaux~; soit $A=\varinjlim A_\lambda$. Soit d'autre part, pour chaque
$\lambda$, $M_\lambda$ un $A_\lambda$-module et pour $\lambda\leq\mu$ soit
$\theta_{\mu\lambda}:M_\lambda\to M_\mu$ un
$\varphi_{\mu\lambda}$-homomorphisme, tel que
$(M_\lambda,\theta_{\mu\lambda})$ soit un système inductif~;
$M=\varinjlim M_\lambda$ est alors un $A$-module. Cela étant, si pour tout
$\lambda$, $M_\lambda$ est un $A_\lambda$-module \emph{plat}, alors $M$ est
un $A$-module \emph{plat}. En effet, soit $\mathfrak{J}$ un idéal \emph{de
type fini} de $A$~; par définition de la limite inductive, il existe un
indice $\lambda$ et un idéal $\mathfrak{J}_\lambda$ de $A_\lambda$ tels que
$\mathfrak{J}=\mathfrak{J}_\lambda A$. Si on pose
$\mathfrak{J}'_\mu=\mathfrak{J}_\lambda A_\mu$ pour $\mu\geq\lambda$, on a
aussi $\mathfrak{J}=\varinjlim\mathfrak{J}'_\mu$ (où $\mu$ parcourt les
indices $\geq\lambda$), d'où (le foncteur $\varinjlim$ étant exact et
commutant au produit tensoriel)
\[
  M\otimes_A\mathfrak{J}
  =\varinjlim(M_\mu\otimes_{A_\mu}\mathfrak{J}'_\mu)
  =\varinjlim\mathfrak{J}'_\mu M_\mu
  =\mathfrak{J}M.
\]
\end{env}

\subsection{Localisation de la platitude}
\label{subsection:0.6.3-fr}

\begin{env}[6.3.1]
\label{0.6.3.1-fr}
Si $A$ est un anneau, $S$ une partie multiplicative de $A$, $S^{-1}A$ est
un $A$-module \emph{plat}. En effet, pour tout $A$-module $N$,
$N\otimes_A S^{-1}A$ s'identifie à $S^{-1}N$ (1.2.5) et on sait (1.3.2)
que $S^{-1}N$ est un foncteur exact en $N$.

Si maintenant $M$ est un $A$-module plat,
$S^{-1}M=M\otimes_A S^{-1}A$ est un $S^{-1}A$-module plat (6.2.1), donc
est aussi $A$-plat en vertu de ce qui précède et de (6.2.1). En particulier,
si $P$ est un $S^{-1}A$-module, on peut le considérer comme un $A$-module
isomorphe à $S^{-1}P$~; pour que $P$ soit $A$-plat, il faut et il suffit
qu'il soit $S^{-1}A$-plat.
\end{env}

\begin{env}[6.3.2]
\label{0.6.3.2-fr}
Soient $A$ un anneau, $B$ une $A$-algèbre, $T$ une partie multiplicative de
$B$. Si $P$ est un $B$-module qui est $A$-plat, $T^{-1}P$ est $A$-plat. En
effet, pour tout $A$-module $N$, on a
\[
  (T^{-1}P)\otimes_A N
  =(T^{-1}B\otimes_B P)\otimes_A N
  =T^{-1}B\otimes_B(P\otimes_A N)
  =T^{-1}(P\otimes_A N)~;
\]
or, $T^{-1}(P\otimes_A N)$ est un foncteur exact en $N$, étant composé des
deux foncteurs exacts $P\otimes_A N$ (en $N$) et $T^{-1}Q$ (en $Q$). Si $S$
est une partie multiplicative de $A$ dont l'image dans $B$ est
\emph{contenue dans $T$}, $T^{-1}P$ est égal à $S^{-1}(T^{-1}P)$, donc est
aussi $S^{-1}A$-plat en vertu de (6.3.1).
\end{env}

\begin{env}[6.3.3]
\label{0.6.3.3-fr}
Soient $\varphi:A\to B$ un homomorphisme d'anneaux, $M$ un $B$-module. Les
propriétés suivantes sont équivalentes~:
\begin{enumerate}
  \item[$a)$] $M$ est un $A$-module plat.
  \item[$b)$] Pour tout idéal maximal $\mathfrak{n}$ de $B$,
    $M_\mathfrak{n}$ est un $A$-module plat.
  \item[$c)$] Pour tout idéal maximal $\mathfrak{n}$ de $B$, en posant
    $\mathfrak{m}=\varphi^{-1}(\mathfrak{n})$, $M_\mathfrak{n}$ est un
    $A_\mathfrak{m}$-module plat.
\end{enumerate}

En effet, comme $M_\mathfrak{n}=(M_\mathfrak{n})_\mathfrak{m}$,
l'équivalence de $b)$ et $c)$ résulte de (6.3.1), et le fait que $a)$
entraîne $b)$ est un cas particulier de (6.3.2). Reste à voir que $b)$
entraîne $a)$,
\oldpage[0\textsubscript{I}]{57}
c'est-à-dire que, pour tout homomorphisme injectif $u:N'\to N$ de
$A$-modules, l'homomorphisme
$v=1\otimes u:M\otimes_A N'\to M\otimes_A N$ est injectif. Or, $v$ est aussi
un homomorphisme de $B$-modules, et on sait que pour qu'il soit injectif, il
suffit que pour tout idéal maximal $\mathfrak{n}$ de $B$,
\[
  v_\mathfrak{n}:(M\otimes_A N')_\mathfrak{n}
  \longrightarrow(M\otimes_A N)_\mathfrak{n}
\]
soit injectif. Mais comme
\[
  (M\otimes_A N)_\mathfrak{n}
  =B_\mathfrak{n}\otimes_B(M\otimes_A N)
  =M_\mathfrak{n}\otimes_A N,
\]
$v_\mathfrak{n}$ n'est autre que l'homomorphisme
\[
  1\otimes u:M_\mathfrak{n}\otimes_A N'
  \longrightarrow M_\mathfrak{n}\otimes_A N,
\]
qui est injectif puisque $M_\mathfrak{n}$ est $A$-plat.

En particulier (en faisant $B=A$), pour qu'un $A$-module $M$ soit plat, il
faut et il suffit que $M_\mathfrak{m}$ soit $A_\mathfrak{m}$-plat pour tout
idéal maximal $\mathfrak{m}$ de $A$.
\end{env}

\begin{env}[6.3.4]
\label{0.6.3.4-fr}
Soit $M$ un $A$-module~; si $M$ est plat, et si $f\in A$ n'est pas diviseur
de $0$ dans $A$, $f$ n'annule aucun élément $\neq0$ de $M$, car
l'homomorphisme $m\mapsto f\cdot m$ s'écrit $1\otimes u$, où $u$ est la
multiplication $a\mapsto f\cdot a$ dans $A$ et $M$ est identifié à
$M\otimes_A A$~; si $u$ est injective, il en est donc de même de
$1\otimes u$ puisque $M$ est plat. En particulier, si $A$ est
\emph{intègre}, $M$ est \emph{sans torsion}.

Inversement, supposons $A$ intègre, $M$ sans torsion, et supposons que pour
tout idéal maximal $\mathfrak{m}$ de $A$, $A_\mathfrak{m}$ soit un
\emph{anneau de valuation discrète}~; alors $M$ est $A$-plat. Il suffit en
effet (6.3.3) de prouver que $M_\mathfrak{m}$ est $A_\mathfrak{m}$-plat, et
on peut donc supposer que $A$ est déjà un anneau de valuation discrète. Mais
comme $M$ est limite inductive de ses sous-modules de type fini, et que ces
derniers sont sans torsion, on peut en outre se borner au cas où $M$ est de
type fini (6.1.2). La proposition résulte dans ce cas de ce que $M$ est un
$A$-module libre.

En particulier, si $A$ est un anneau \emph{intègre},
$\varphi:A\to B$ un homomorphisme d'anneaux faisant de $B$ un $A$-module
\emph{plat} et $\neq\{0\}$, $\varphi$ est nécessairement \emph{injectif}.
Inversement, si $B$ est intègre, $A$ un sous-anneau de $B$ et si pour tout
idéal maximal $\mathfrak{m}$ de $A$, $A_\mathfrak{m}$ est un anneau de
valuation discrète, $B$ est $A$-plat.
\end{env}

\subsection{Modules fidèlement plats}
\label{subsection:0.6.4-fr}

\begin{env}[6.4.1]
\label{0.6.4.1-fr}
Pour un $A$-module $M$, les quatre propriétés suivantes sont équivalentes~:
\begin{enumerate}
  \item[$a)$] Pour qu'une suite $N'\to N\to N''$ de $A$-modules soit exacte,
    il faut et il suffit que la suite
    $M\otimes N'\to M\otimes N\to M\otimes N''$ soit exacte~;
  \item[$b)$] $M$ est plat et pour tout $A$-module $N$, la relation
    $M\otimes N=0$ entraîne $N=0$~;
  \item[$c)$] $M$ est plat et pour tout homomorphisme $v:N\to N'$ de
    $A$-modules, la relation $1_M\otimes v=0$ entraîne $v=0$, $1_M$ étant
    l'automorphisme identique de $M$~;
  \item[$d)$] $M$ est plat et pour tout idéal maximal $\mathfrak{m}$ de
    $A$, $\mathfrak{m}M\neq M$.
\end{enumerate}

Lorsque $M$ vérifie ces conditions, on dit que $M$ est un $A$-module
\emph{fidèlement plat}~; $M$ est alors nécessairement un module
\emph{fidèle}. En outre, si $u:N\to N'$ est un homomorphisme de $A$-modules,
pour que $u$ soit injectif (resp. surjectif, bijectif), il faut et il suffit
que $1\otimes u:M\otimes N\to M\otimes N'$ le soit.
\end{env}

\oldpage[0\textsubscript{I}]{58}

\begin{env}[6.4.2]
\label{0.6.4.2-fr}
Un module libre $\neq\{0\}$ est fidèlement plat~; il en est de même de la
somme directe d'un module plat et d'un module fidèlement plat. Si $S$ est
une partie multiplicative de $A$, $S^{-1}A$ n'est un $A$-module fidèlement
plat que si $S$ est formé d'éléments inversibles (donc $S^{-1}A=A$).
\end{env}

\begin{env}[6.4.3]
\label{0.6.4.3-fr}
Soit
\[
  0\longrightarrow M'\longrightarrow M\longrightarrow M''\longrightarrow 0
\]
une suite exacte de $A$-modules~; si $M'$ et $M''$ sont plats, et si l'un
d'eux est fidèlement plat, alors $M$ est fidèlement plat.
\end{env}

\begin{env}[6.4.4]
\label{0.6.4.4-fr}
Soient $A$ et $B$ deux anneaux, $M$ un $A$-module, $N$ un
$(A,B)$-bimodule. Si $M$ est fidèlement plat et si $N$ est un $B$-module
fidèlement plat, alors $M\otimes_A N$ est un $B$-module fidèlement plat. En
particulier, si $M$ et $N$ sont deux $A$-modules fidèlement plats, il en est
de même de $M\otimes_A N$. Si $B$ est une $A$-algèbre et si $M$ est un
$A$-module fidèlement plat, le $B$-module $M_{(B)}$ est fidèlement plat.
\end{env}

\begin{env}[6.4.5]
\label{0.6.4.5-fr}
Si $M$ est un $A$-module fidèlement plat et si $S$ est une partie
multiplicative de $A$, $S^{-1}M$ est un $S^{-1}A$-module fidèlement plat,
puisque $S^{-1}M=M\otimes_A(S^{-1}A)$ (6.4.4). Inversement, si pour tout
idéal maximal $\mathfrak{m}$ de $A$, $M_\mathfrak{m}$ est un
$A_\mathfrak{m}$-module fidèlement plat, $M$ est un $A$-module fidèlement
plat, car $M$ est $A$-plat (6.3.3), et on a
\[
  M_\mathfrak{m}/\mathfrak{m}M_\mathfrak{m}
  =(M\otimes_A A_\mathfrak{m})
  \otimes_{A_\mathfrak{m}}
  (A_\mathfrak{m}/\mathfrak{m}A_\mathfrak{m})
  =M\otimes_A(A/\mathfrak{m})
  =M/\mathfrak{m}M,
\]
donc l'hypothèse entraîne que $M/\mathfrak{m}M\neq0$ pour tout idéal
maximal $\mathfrak{m}$ de $A$, ce qui prouve notre assertion (6.4.1).
\end{env}

\subsection{Restriction des scalaires}
\label{subsection:0.6.5-fr}

\begin{env}[6.5.1]
\label{0.6.5.1-fr}
Soient $A$ un anneau, $\varphi:A\to B$ un homomorphisme d'anneaux faisant
de $B$ une $A$-algèbre. Supposons qu'il existe un $B$-module $N$ qui soit un
$A$-module \emph{fidèlement plat}. Alors, pour tout $A$-module $M$,
l'homomorphisme $x\mapsto 1\otimes x$ de $M$ dans
$B\otimes_A M=M_{(B)}$ est \emph{injectif}. En particulier, $\varphi$ est
injectif~; pour tout idéal $\mathfrak{a}$ de $A$, on a
$\varphi^{-1}(\mathfrak{a}B)=\mathfrak{a}$~; pour tout idéal maximal (resp.
premier) $\mathfrak{m}$ de $A$, il existe un idéal maximal (resp. premier)
$\mathfrak{n}$ de $B$ tel que $\varphi^{-1}(\mathfrak{n})=\mathfrak{m}$.
\end{env}

\begin{env}[6.5.2]
\label{0.6.5.2-fr}
Lorsque la condition de (6.5.1) est remplie, on identifie $A$ à un
sous-anneau de $B$ par $\varphi$ et plus généralement, pour tout $A$-module
$M$, on identifie $M$ à un sous-$A$-module de $M_{(B)}$. On notera que si
$B$ est alors \emph{noethérien}, il en est de même de $A$, car l'application
$\mathfrak{a}\mapsto\mathfrak{a}B$ est une injection croissante de
l'ensemble des idéaux de $A$ dans l'ensemble des idéaux de $B$~;
l'existence d'une suite infinie strictement croissante d'idéaux de $A$
entraînerait donc l'existence d'une suite analogue d'idéaux de $B$.
\end{env}

\subsection{Anneaux fidèlement plats}
\label{subsection:0.6.6-fr}

\begin{env}[6.6.1]
\label{0.6.6.1-fr}
Soit $\varphi:A\to B$ un homomorphisme d'anneaux faisant de $B$ une
$A$-algèbre. Les cinq propriétés suivantes sont équivalentes~:
\begin{enumerate}
  \item[$a)$] $B$ est un $A$-module fidèlement plat (autrement dit,
    $M_{(B)}$ est un foncteur exact et fidèle en $M$).
  \item[$b)$] L'homomorphisme $\varphi$ est injectif et le $A$-module
    $B/\varphi(A)$ est plat.

\oldpage[0\textsubscript{I}]{59}

  \item[$c)$] Le $A$-module $B$ est plat (autrement dit, le foncteur
    $M_{(B)}$ est exact), et pour tout $A$-module $M$, l'homomorphisme
    $x\mapsto1\otimes x$ de $M$ dans $M_{(B)}$ est injectif.
  \item[$d)$] Le $A$-module $B$ est plat et pour tout idéal $\mathfrak a$ de
    $A$, on a $\varphi^{-1}(\mathfrak aB)=\mathfrak a$.
  \item[$e)$] Le $A$-module $B$ est plat et pour tout idéal maximal
    $\mathfrak m$ de $A$, il existe un idéal maximal $\mathfrak n$ de $B$ tel
    que $\varphi^{-1}(\mathfrak n)=\mathfrak m$.
\end{enumerate}

Lorsque ces conditions ont lieu, on identifie $A$ à un sous-anneau de $B$.
\end{env}

\begin{env}[6.6.2]
\label{0.6.6.2-fr}
Soient $A$ un anneau local, $\mathfrak m$ son idéal maximal, $B$ une
$A$-algèbre telle que $\mathfrak mB\neq B$ (ce qui a lieu par exemple lorsque
$B$ est un anneau local et $A\to B$ un homomorphisme local). Si $B$ est un
$A$-module plat, $B$ est un $A$-module fidèlement plat. En effet, comme
$\mathfrak mB\neq B$, il y a un idéal maximal $\mathfrak n$ de $B$ contenant
$\mathfrak mB$~; comme $\varphi^{-1}(\mathfrak n)\cap A$ contient
$\mathfrak m$ et ne contient pas $1$, on a
$\varphi^{-1}(\mathfrak n)=\mathfrak m$ et le critère $e)$ de (6.6.1)
s'applique. Sous les conditions indiquées, on voit donc que si $B$ est
noethérien, il en est de même de $A$ (6.5.2).
\end{env}

\begin{env}[6.6.3]
\label{0.6.6.3-fr}
Soit $B$ une $A$-algèbre qui est un $A$-module fidèlement plat. Pour tout
$A$-module $M$ et tout sous-$A$-module $M'$ de $M$, on a (en identifiant $M$
à un sous-$A$-module de $M_{(B)}$) $M'=M\cap M'_{(B)}$. Pour que $M$ soit un
$A$-module plat (resp. fidèlement plat), il faut et il suffit que $M_{(B)}$
soit un $B$-module plat (resp. fidèlement plat).
\end{env}

\begin{env}[6.6.4]
\label{0.6.6.4-fr}
Soient $B$ une $A$-algèbre, $N$ un $B$-module fidèlement plat. Pour que $B$
soit un $A$-module plat (resp. fidèlement plat), il faut et il suffit que $N$
le soit.

En particulier, soit $C$ une $B$-algèbre~; si l'anneau $C$ est fidèlement
plat sur $B$ et $B$ fidèlement plat sur $A$, alors $C$ est fidèlement plat
sur $A$~; si $C$ est fidèlement plat sur $B$ et sur $A$, alors $B$ est
fidèlement plat sur $A$.
\end{env}

\subsection{Morphismes plats d'espaces annelés}
\label{subsection:0.6.7-fr}

\begin{env}[6.7.1]
\label{0.6.7.1-fr}
Soit $f:X\to Y$ un morphisme d'espaces annelés, et soit $\mathcal F$ un
$\mathcal O_X$-Module. On dit que $\mathcal F$ est $f$-plat (ou $Y$-plat
lorsque aucune confusion n'est à craindre sur $f$) en un point $x\in X$ si
$\mathcal F_x$ est un $\mathcal O_{f(x)}$-module plat~; on dit que
$\mathcal F$ est $f$-plat au-dessus de $y\in Y$ si $\mathcal F$ est $f$-plat
en tous les points $x\in f^{-1}(y)$~; on dit que $\mathcal F$ est $f$-plat
si $\mathcal F$ est $f$-plat en tous les points de $X$. On dit que le
morphisme $f$ est plat en $x\in X$ (resp. plat au-dessus de $y\in Y$, resp.
plat) si $\mathcal O_X$ est $f$-plat en $x$ (resp. $f$-plat au-dessus de $y$,
resp. $f$-plat).
\end{env}

\begin{env}[6.7.2]
\label{0.6.7.2-fr}
Avec les notations de (6.7.1), si $\mathcal F$ est $f$-plat en $x$, pour tout
voisinage ouvert $U$ de $y=f(x)$, le foncteur
$(f^*(\mathcal G)\otimes_{\mathcal O_X}\mathcal F)_x$ en $\mathcal G$ est
exact dans la catégorie des $(\mathcal O_Y|U)$-Modules~; en effet, cette
fibre s'identifie canoniquement à
$\mathcal G_y\otimes_{\mathcal O_y}\mathcal F_x$, et notre assertion résulte
de la définition. En particulier, si $f$ est un morphisme plat, le foncteur
$f^*$ est exact dans la catégorie des $\mathcal O_Y$-Modules.
\end{env}

\begin{env}[6.7.3]
\label{0.6.7.3-fr}
Inversement, supposons le faisceau d'anneaux $\mathcal O_Y$ cohérent, et
supposons que pour tout voisinage ouvert $U$ de $y$, le foncteur
$(f^*(\mathcal G)\otimes_{\mathcal O_X}\mathcal F)_x$ soit exact en
$\mathcal G$ dans la catégorie des $(\mathcal O_Y|U)$-Modules cohérents.
Alors $\mathcal F$ est $f$-plat en $x$. En effet, il suffit de prouver que
pour tout idéal de type fini $\mathfrak J$ de $\mathcal O_y$,
l'homomorphisme canonique
$\mathfrak J\otimes_{\mathcal O_y}\mathcal F_x\to\mathcal F_x$ est injectif
(6.1.1). Or, on sait (5.3.8) qu'il existe alors un voisinage

\oldpage[0\textsubscript{I}]{60}

ouvert $U$ de $y$ et un faisceau cohérent d'idéaux $\mathcal J$ de
$\mathcal O_Y|U$ tel que $\mathcal J_y=\mathfrak J$, d'où la conclusion.
\end{env}

\begin{env}[6.7.4]
\label{0.6.7.4-fr}
Les résultats de (6.1) sur les modules plats se transcrivent aussitôt en
propositions sur les faisceaux $f$-plats en un point~:

Si $0\to\mathcal F'\to\mathcal F\to\mathcal F''\to0$ est une suite exacte
de $\mathcal O_X$-Modules et si $\mathcal F''$ est $f$-plat au point
$x\in X$, alors, pour tout voisinage ouvert $U$ de $y=f(x)$ et tout
$(\mathcal O_Y|U)$-Module $\mathcal G$, la suite
\[
  0\longrightarrow
  (f^*(\mathcal G)\otimes_{\mathcal O_X}\mathcal F')_x
  \longrightarrow
  (f^*(\mathcal G)\otimes_{\mathcal O_X}\mathcal F)_x
  \longrightarrow
  (f^*(\mathcal G)\otimes_{\mathcal O_X}\mathcal F'')_x
  \longrightarrow0
\]
est exacte. Pour que $\mathcal F$ soit $f$-plat en $x$, il faut et il suffit
alors que $\mathcal F'$ le soit. On a des énoncés analogues pour les notions
correspondantes de $\mathcal O_X$-Module $f$-plat au-dessus de $y\in Y$, ou
de $\mathcal O_X$-Module $f$-plat.
\end{env}

\begin{env}[6.7.5]
\label{0.6.7.5-fr}
Soient $f:X\to Y$, $g:Y\to Z$ deux morphismes d'espaces annelés~; soient
$x\in X$, $y=f(x)$, $\mathcal F$ un $\mathcal O_X$-Module. Si $\mathcal F$
est $f$-plat au point $x$ et si le morphisme $g$ est plat au point $y$, alors
$\mathcal F$ est $(g\circ f)$-plat en $x$ (6.2.1). En particulier, si $f$ et
$g$ sont des morphismes plats, $g\circ f$ est plat.
\end{env}

\begin{env}[6.7.6]
\label{0.6.7.6-fr}
Soient $X$, $Y$ deux espaces annelés, $f:X\to Y$ un morphisme plat. Alors
l'homomorphisme canonique de bifoncteurs (4.4.6)
\[
  f^*(\mathcal{H}om_{\mathcal O_Y}(\mathcal F,\mathcal G))
  \longrightarrow
  \mathcal{H}om_{\mathcal O_X}(f^*(\mathcal F),f^*(\mathcal G))
  \tag{6.7.6.1}
  \label{0.6.7.6.1-fr}
\]
est un isomorphisme lorsque $\mathcal F$ admet une présentation finie
(5.2.5).

En effet, la question étant locale, on peut supposer qu'il existe une suite
exacte $\mathcal O_Y^m\to\mathcal O_Y^n\to\mathcal F\to0$. Or, les deux
membres de (6.7.6.1) sont des foncteurs exacts à gauche en $\mathcal F$ en
vertu de l'hypothèse sur $f$~; on est alors ramené à démontrer la proposition
lorsque $\mathcal F=\mathcal O_Y$, cas où elle est triviale.
\end{env}

\begin{env}[6.7.8]
\label{0.6.7.8-fr}
On dit qu'un morphisme $f:X\to Y$ d'espaces annelés est fidèlement plat si
$f$ est surjectif et si, pour tout $x\in X$, $\mathcal O_x$ est un
$\mathcal O_{f(x)}$-module fidèlement plat. Lorsque $X$ et $Y$ sont des
espaces annelés en anneaux locaux (5.5.1), il revient au même de dire que le
morphisme $f$ est surjectif et plat (6.6.2). Lorsque $f$ est fidèlement plat,
$f^*$ est un foncteur exact et fidèle dans la catégorie des
$\mathcal O_Y$-Modules (6.6.1, $a)$) et pour qu'un $\mathcal O_Y$-Module
$\mathcal G$ soit $Y$-plat, il faut et il suffit que $f^*(\mathcal G)$ le
soit (6.6.3).
\end{env}

\section{Anneaux adiques}
\label{section:0.7-fr}

\subsection{Anneaux admissibles}
\label{subsection:0.7.1-fr}

\begin{env}[7.1.1]
\label{0.7.1.1-fr}
Rappelons que dans un anneau topologique $A$ (non nécessairement séparé), on
dit qu'un élément $x$ est topologiquement nilpotent si $0$ est une limite de
la suite $(x^n)_{n\geq0}$. On dit qu'un anneau topologique $A$ est
linéairement topologisé s'il existe un système fondamental de voisinages de
$0$ dans $A$ formé d'idéaux (nécessairement ouverts).
\end{env}

\begin{definition}[7.1.2]
\label{0.7.1.2-fr}
Dans un anneau linéairement topologisé $A$, on dit qu'un idéal $\mathfrak J$
est un idéal de définition si $\mathfrak J$ est ouvert et si, pour tout
voisinage $V$ de $0$, il existe un entier $n>0$ tel

\oldpage[0\textsubscript{I}]{61}

que $\mathfrak J^n\subset V$ (ce qu'on exprime, par abus de langage, en disant
que la suite $(\mathfrak J^n)$ tend vers $0$). On dit qu'un anneau linéairement
topologisé $A$ est préadmissible s'il existe dans $A$ un idéal de définition~;
on dit que $A$ est admissible s'il est préadmissible et si en outre il est
séparé et complet.
\end{definition}

Il est clair que si $\mathfrak J$ est un idéal de définition, $\mathfrak L$ un
idéal ouvert de $A$, $\mathfrak J\cap\mathfrak L$ est encore un idéal de
définition~; les idéaux de définition d'un anneau préadmissible $A$ forment
donc un système fondamental de voisinages de $0$.

\begin{lemma}[7.1.3]
\label{0.7.1.3-fr}
Soit $A$ un anneau linéairement topologisé.
\begin{enumerate}
  \item[{\upshape(i)}] Pour que $x\in A$ soit topologiquement nilpotent, il
    faut et il suffit que pour tout idéal ouvert $\mathfrak J$ de $A$, l'image
    canonique de $x$ dans $A/\mathfrak J$ soit nilpotente. L'ensemble
    $\mathfrak T$ des éléments topologiquement nilpotents de $A$ est un idéal.
  \item[{\upshape(ii)}] Supposons en outre que $A$ soit préadmissible, et soit
    $\mathfrak J$ un idéal de définition de $A$. Pour que $x\in A$ soit
    topologiquement nilpotent, il faut et il suffit que son image canonique dans
    $A/\mathfrak J$ soit nilpotente~; l'idéal $\mathfrak T$ est l'image
    réciproque du nilradical de $A/\mathfrak J$ et est donc ouvert.
\end{enumerate}
\end{lemma}

(i) découle immédiatement des définitions. Pour prouver (ii), il suffit de
remarquer que pour tout voisinage $V$ de $0$ dans $A$, il existe $n>0$ tel que
$\mathfrak J^n\subset V$~; si $x\in A$ est tel que $x^m\in\mathfrak J$, on a
$x^{mq}\in V$ pour $q\geq n$, donc $x$ est topologiquement nilpotent.

\begin{proposition}[7.1.4]
\label{0.7.1.4-fr}
Soient $A$ un anneau préadmissible, $\mathfrak J$ un idéal de définition de
$A$.
\begin{enumerate}
  \item[{\upshape(i)}] Pour qu'un idéal $\mathfrak J'$ de $A$ soit contenu
    dans un idéal de définition, il faut et il suffit qu'il existe un entier
    $n>0$ tel que ${\mathfrak J'}^n\subset\mathfrak J$.
  \item[{\upshape(ii)}] Pour qu'un $x\in A$ soit contenu dans un idéal de
    définition, il faut et il suffit qu'il soit topologiquement nilpotent.
\end{enumerate}
\end{proposition}

(i) Si ${\mathfrak J'}^n\subset\mathfrak J$, pour tout voisinage ouvert $V$ de
$0$ dans $A$, il existe $m$ tel que $\mathfrak J^m\subset V$, donc
${\mathfrak J'}^{mn}\subset V$.

(ii) La condition est évidemment nécessaire~; elle est suffisante, car si elle
est remplie, il existe $n$ tel que $x^n\in\mathfrak J$, donc
$\mathfrak J'=\mathfrak J+Ax$ est un idéal de définition, puisqu'il est ouvert,
et que ${\mathfrak J'}^n\subset\mathfrak J$.

\begin{corollary}[7.1.5]
\label{0.7.1.5-fr}
Dans un anneau préadmissible $A$, un idéal premier ouvert contient tous les
idéaux de définition.
\end{corollary}

\begin{corollary}[7.1.6]
\label{0.7.1.6-fr}
Les notations et hypothèses étant celles de (7.1.4), les propriétés suivantes
d'un idéal $\mathfrak J_0$ de $A$ sont équivalentes~:
\begin{enumerate}
  \item[$a)$] $\mathfrak J_0$ est le plus grand idéal de définition de $A$~;
  \item[$b)$] $\mathfrak J_0$ est un idéal de définition maximal~;
  \item[$c)$] $\mathfrak J_0$ est un idéal de définition tel que l'anneau
    $A/\mathfrak J_0$ soit réduit.
\end{enumerate}
Pour qu'il existe un idéal $\mathfrak J_0$ ayant ces propriétés, il faut et il
suffit que le nilradical de $A/\mathfrak J$ soit nilpotent~; $\mathfrak J_0$
est alors égal à l'idéal $\mathfrak T$ des éléments topologiquement nilpotents
de $A$.
\end{corollary}

Il est clair que $a)$ implique $b)$, et $b)$ implique $c)$ en vertu de
(7.1.4, (ii)) et (7.1.3, (ii))~; pour la même raison, $c)$ entraîne $a)$. La
dernière assertion résulte de (7.1.4, (i)) et (7.1.3, (ii)).

Lorsque $\mathfrak T/\mathfrak J$, nilradical de $A/\mathfrak J$, est
nilpotent, on note $A_{\mathrm{red}}$ l'anneau quotient (réduit)
$A/\mathfrak T$.

\oldpage[0\textsubscript{I}]{62}

\begin{corollary}[7.1.7]
\label{0.7.1.7-fr}
Un anneau préadmissible noethérien admet un plus grand idéal de définition.
\end{corollary}

\begin{corollary}[7.1.8]
\label{0.7.1.8-fr}
Si un anneau préadmissible $A$ est tel que, pour un idéal de définition
$\mathfrak J$, les puissances $\mathfrak J^n$ ($n>0$) forment un système
fondamental de voisinages de $0$, il en est de même des puissances
${\mathfrak J'}^n$ de tout idéal de définition $\mathfrak J'$ de $A$.
\end{corollary}

\begin{definition}[7.1.9]
\label{0.7.1.9-fr}
On dit qu'un anneau préadmissible $A$ est préadique s'il existe un idéal de
définition $\mathfrak J$ de $A$ tel que les $\mathfrak J^n$ forment un système
fondamental de voisinages de $0$ dans $A$ (ou, ce qui revient au même, tel que
les $\mathfrak J^n$ soient ouverts). On appelle anneau adique un anneau
préadique séparé et complet.
\end{definition}

Si $\mathfrak J$ est un idéal de définition d'un anneau préadique (resp.
adique) $A$, on dit encore que $A$ est un anneau
\emph{$\mathfrak J$-préadique} (resp. \emph{$\mathfrak J$-adique}), et que sa
topologie est la topologie \emph{$\mathfrak J$-préadique} (resp.
\emph{$\mathfrak J$-adique}). Plus généralement, si $M$ est un $A$-module, la
topologie sur $M$ ayant pour système fondamental de voisinages de $0$ les
sous-modules $\mathfrak J^nM$ est dite topologie
\emph{$\mathfrak J$-préadique} (resp. \emph{$\mathfrak J$-adique}). En vertu
de (7.1.8), ces topologies sont indépendantes du choix de l'idéal de définition
$\mathfrak J$.

\begin{proposition}[7.1.10]
\label{0.7.1.10-fr}
Soient $A$ un anneau admissible, $\mathfrak J$ un idéal de définition de $A$.
Alors $\mathfrak J$ est contenu dans le radical de $A$.
\end{proposition}

Cet énoncé est équivalent à l'un quelconque des corollaires suivants~:

\begin{corollary}[7.1.11]
\label{0.7.1.11-fr}
Pour tout $x\in\mathfrak J$, $1+x$ est inversible dans $A$.
\end{corollary}

\begin{corollary}[7.1.12]
\label{0.7.1.12-fr}
Pour que $f\in A$ soit inversible dans $A$, il faut et il suffit que son image
canonique dans $A/\mathfrak J$ soit inversible dans $A/\mathfrak J$.
\end{corollary}

\begin{corollary}[7.1.13]
\label{0.7.1.13-fr}
Pour tout $A$-module $M$ de type fini, la relation $M=\mathfrak JM$
(équivalente à $M\otimes_A(A/\mathfrak J)=0$) entraîne $M=0$.
\end{corollary}

\begin{corollary}[7.1.14]
\label{0.7.1.14-fr}
Soit $u:M\to N$ un homomorphisme de $A$-modules, $N$ étant de type fini~;
pour que $u$ soit surjectif, il faut et il suffit que
$u\otimes1:M\otimes_A(A/\mathfrak J)\to N\otimes_A(A/\mathfrak J)$ le soit.
\end{corollary}

En effet, l'équivalence de (7.1.10) et (7.1.11) résulte de Bourbaki,
\emph{Alg.}, chap.~VIII, \textsection6, n\textsuperscript{o}~3, th.~1, et
celle de (7.1.10) et (7.1.13) de \emph{loc. cit.}, th.~2~; le fait que
(7.1.10) entraîne (7.1.14) résulte de \emph{loc. cit.}, cor.~4 de la prop.~6~;
d'autre part, (7.1.14) entraîne (7.1.13) en l'appliquant à l'homomorphisme nul.
Enfin, (7.1.10) entraîne que si $f$ est inversible dans $A/\mathfrak J$, $f$
n'est contenu dans aucun idéal maximal de $A$, donc $f$ est inversible dans
$A$, autrement dit (7.1.10) entraîne (7.1.12)~; et inversement, (7.1.12)
entraîne (7.1.11).

Tout revient donc à démontrer (7.1.11). Or, comme $A$ est séparé et complet et
que la suite $(\mathfrak J^n)$ tend vers $0$, il est immédiat que la série
$\sum_{n=0}^{\infty}(-1)^n x^n$ est convergente dans $A$ et que, si $y$ est
sa somme, on a $y(1+x)=1$.

\subsection{Anneaux adiques et limites projectives}
\label{subsection:0.7.2-fr}

\begin{env}[7.2.1]
\label{0.7.2.1-fr}
Toute limite projective d'anneaux discrets est évidemment un anneau
linéairement topologisé, séparé et complet. Inversement, soit $A$ un anneau
linéairement topologisé, et soit $(\mathfrak J_\lambda)$ un système fondamental
de voisinages ouverts de $0$ dans $A$ formé

\oldpage[0\textsubscript{I}]{63}

d'idéaux. Les applications canoniques
$\varphi_\lambda:A\to A/\mathfrak J_\lambda$ forment un système projectif de
représentations continues et définissent donc une représentation continue
$\varphi:A\to\varprojlim A/\mathfrak J_\lambda$~; si $A$ est séparé,
$\varphi$ est un isomorphisme topologique de $A$ sur un
sous-anneau partout dense de $\varprojlim A/\mathfrak J_\lambda$~; si, en
outre, $A$ est complet, $\varphi$ est un isomorphisme topologique de $A$ sur
$\varprojlim A/\mathfrak J_\lambda$.
\end{env}

\begin{lemma}[7.2.2]
\label{0.7.2.2-fr}
Pour qu'un anneau linéairement topologisé soit admissible, il faut et il suffit
qu'il soit isomorphe à une limite projective
$A=\varprojlim A_\lambda$, où $(A_\lambda,u_{\lambda\mu})$ est un système
projectif d'anneaux discrets ayant pour ensemble d'indices un ensemble ordonné
filtrant $L$ (pour $\leq$) qui admet un plus petit élément noté $0$ et satisfait
aux conditions suivantes~: 1\textsuperscript{o} les
$u_\lambda:A\to A_\lambda$ sont surjectifs~; 2\textsuperscript{o} le noyau
$\mathfrak J_\lambda$ de $u_{0\lambda}:A_\lambda\to A_0$ est nilpotent.
Lorsqu'il en est ainsi, le noyau $\mathfrak J$ de $u_0:A\to A_0$ est égal à
$\varprojlim\mathfrak J_\lambda$.
\end{lemma}

La nécessité de la condition résulte de (7.2.1), en prenant pour
$(\mathfrak J_\lambda)$ un système fondamental de voisinages de $0$ formé
d'idéaux de définition contenus dans l'un d'eux $\mathfrak J_0$ et appliquant
(7.1.4, (i)). La réciproque résulte de la définition d'une limite projective et
de (7.1.2), et la dernière assertion est immédiate.

\begin{env}[7.2.3]
\label{0.7.2.3-fr}
Soient $A$ un anneau topologique admissible, $\mathfrak J$ un idéal de $A$
contenu dans un idéal de définition (autrement dit (7.1.4) tel que
$(\mathfrak J^n)$ tende vers $0$)~; on peut considérer sur $A$ la topologie
d'anneau ayant pour système fondamental de voisinages de $0$ les puissances
$\mathfrak J^n$ ($n>0$)~; nous l'appellerons encore la topologie
\emph{$\mathfrak J$-préadique}. L'hypothèse que $A$ est admissible entraîne que
$\bigcap_n\mathfrak J^n=0$, donc la topologie $\mathfrak J$-préadique sur $A$
est séparée~; soit
$\widehat A=\varprojlim A/\mathfrak J^n$ le complété de $A$ pour cette topologie
(où les $A/\mathfrak J^n$ sont munis de la topologie discrète), et désignons
par $u$ l'homomorphisme d'anneaux $A\to\widehat A$ (non nécessairement
continu), limite projective de la suite d'homomorphismes
$u_n:A\to A/\mathfrak J^n$. D'autre part, la topologie $\mathfrak J$-préadique
sur $A$ est plus fine que la topologie donnée $\mathcal T$ sur $A$~; comme $A$
est séparé et complet pour $\mathcal T$, on peut prolonger par continuité
l'application identique de $A$ (muni de la topologie $\mathfrak J$-préadique)
dans $A$ muni de $\mathcal T$~; cela donne une représentation continue
$v:\widehat A\to A$.
\end{env}

\begin{proposition}[7.2.4]
\label{0.7.2.4-fr}
Si $A$ est un anneau admissible et $\mathfrak J$ est contenu dans un idéal de
définition de $A$, $A$ est séparé et complet pour la topologie
$\mathfrak J$-préadique.
\end{proposition}

En effet, avec les notations de (7.2.3), il est immédiat que $v\circ u$ est
l'application identique de $A$. D'autre part,
$u_n\circ v:\widehat A\to A/\mathfrak J^n$ est le prolongement par continuité
(pour la topologie $\mathfrak J$-préadique sur $A$ et la topologie discrète sur
$A/\mathfrak J^n$) de l'application canonique $u_n$~; autrement dit, c'est
l'application canonique de
$\widehat A=\varprojlim_k A/\mathfrak J^k$ sur $A/\mathfrak J^n$~;
$u\circ v$ est donc limite projective de cette suite d'applications,
c'est-à-dire par définition l'application identique de $\widehat A$~; ceci
démontre la proposition.

\begin{corollary}[7.2.5]
\label{0.7.2.5-fr}
Sous les hypothèses de (7.2.3), les conditions suivantes sont équivalentes~:
\begin{enumerate}
  \item[$a)$] L'homomorphisme $u$ est continu~;

\oldpage[0\textsubscript{I}]{64}

  \item[$b)$] L'homomorphisme $v$ est bicontinu~;
  \item[$c)$] $A$ est un anneau $\mathfrak J$-adique.
\end{enumerate}
\end{corollary}

\begin{corollary}[7.2.6]
\label{0.7.2.6-fr}
Soient $A$ un anneau admissible, $\mathfrak J$ un idéal de définition de $A$.
Pour que $A$ soit noethérien, il faut et il suffit que $A/\mathfrak J$ soit
noethérien et que $\mathfrak J/\mathfrak J^2$ soit un
$(A/\mathfrak J)$-module de type fini.
\end{corollary}

Ces conditions sont évidemment nécessaires. Inversement, supposons-les
vérifiées~; comme en vertu de (7.2.4) $A$ est complet pour la topologie
$\mathfrak J$-préadique, pour qu'il soit noethérien, il faut et il suffit que
l'anneau gradué associé $\operatorname{grad}(A)$ (pour la filtration des
$\mathfrak J^n$) le soit ([I], p.~18-07, th.~4). Or, soient
$a_1,\ldots,a_n$ des éléments de $\mathfrak J$ dont les classes mod.
$\mathfrak J^2$ sont des générateurs de $\mathfrak J/\mathfrak J^2$ en tant
que $A/\mathfrak J$-module. Il est immédiat par récurrence que les classes mod.
$\mathfrak J^{m+1}$ des monômes de degré total $m$ en les $a_i$
$(1\leq i\leq n)$ forment un système de générateurs du
$(A/\mathfrak J)$-module $\mathfrak J^m/\mathfrak J^{m+1}$. On en conclut que
$\operatorname{grad}(A)$ est un anneau isomorphe à un quotient de
$(A/\mathfrak J)[T_1,\ldots,T_n]$ ($T_i$ indéterminées), ce qui achève la
démonstration.

\begin{proposition}[7.2.7]
\label{0.7.2.7-fr}
Soit $(A_i,u_{ij})$ un système projectif $(i\in\mathbf N)$ d'anneaux discrets,
et pour tout entier $i$, soit $\mathfrak J_i$ le noyau dans $A_i$ de
l'homomorphisme $u_{0i}:A_i\to A_0$. On suppose que~:
\begin{enumerate}
  \item[$a)$] Pour $i\leq j$, $u_{ij}$ est surjectif et son noyau est
    $\mathfrak J_j^{i+1}$ (donc $A_i$ est isomorphe à
    $A_j/\mathfrak J_j^{i+1}$).
  \item[$b)$] $\mathfrak J_1/\mathfrak J_1^2$ ($=\mathfrak J_1$) est un
    module de type fini sur $A_0=A_1/\mathfrak J_1$.
\end{enumerate}
Soit $A=\varprojlim_i A_i$, et pour tout entier $n\geq0$, soient $u_n$
l'homomorphisme canonique $A\to A_n$, $\mathfrak J^{(n+1)}\subset A$ son
noyau. Dans ces conditions~:
\begin{enumerate}
  \item[(i)] $A$ est un anneau adique, ayant pour idéal de définition
    $\mathfrak J=\mathfrak J^{(1)}$.
  \item[(ii)] On a $\mathfrak J^{(n)}=\mathfrak J^n$ pour tout $n\geq1$.
  \item[(iii)] $\mathfrak J/\mathfrak J^2$ est isomorphe à
    $\mathfrak J_1=\mathfrak J_1/\mathfrak J_1^2$, et est par suite un module
    de type fini sur $A_0=A/\mathfrak J$.
\end{enumerate}
\end{proposition}

L'hypothèse de surjectivité des $u_{ij}$ entraîne que $u_n$ est surjectif~;
en outre, l'hypothèse $a)$ implique que $\mathfrak J_j^{j+1}=0$, donc $A$ est
un anneau admissible (7.2.2)~; par définition, les $\mathfrak J^{(n)}$ forment
un système fondamental de voisinages de $0$ dans $A$, donc (ii) entraîne (i).
En outre, on a $\mathfrak J=\varprojlim_i\mathfrak J_i$ et les applications
$\mathfrak J\to\mathfrak J_i$ sont surjectives, donc (ii) entraîne (iii), et
on est ramené à prouver (ii). Par définition, $\mathfrak J^{(n)}$ est formé
des éléments $(x_k)_{k\geq0}$ de $A$ tels que $x_k=0$ pour $k<n$, donc
$\mathfrak J^{(n)}\mathfrak J^{(m)}\subset\mathfrak J^{(n+m)}$, autrement dit
les $\mathfrak J^{(n)}$ constituent une \emph{filtration} de $A$. D'autre
part, $\mathfrak J^{(n)}/\mathfrak J^{(n+1)}$ est isomorphe à la projection de
$\mathfrak J^{(n)}$ sur $A_n$~; comme
$\mathfrak J^{(n)}=\varprojlim_{i>n}\mathfrak J_i^n$, cette projection n'est
autre que $\mathfrak J_n^n$, qui est un module sur
$A_0=A_n/\mathfrak J_n$. Soient alors $a_j=(a_{jk})_{k\geq0}$ $r$ éléments de
$\mathfrak J=\mathfrak J^{(1)}$ tels que $a_{11},\ldots,a_{r1}$ forment un
système de générateurs de $\mathfrak J_1$ sur $A_0$~; nous allons voir que
l'ensemble $S_n$ des monômes de degré total $n$ en les $a_j$ engendre l'idéal
$\mathfrak J^{(n)}$ de $A$. Comme $\mathfrak J_i^{i+1}=0$, il est clair tout
d'abord que $S_n\subset\mathfrak J^{(n)}$~; puisque $A$ est complet pour la
filtration $(\mathfrak J^{(m)})$, il suffit de prouver que l'ensemble
$\overline S_n$ des classes mod. $\mathfrak J^{(n+1)}$ des éléments de $S_n$
engendre le module gradué $\operatorname{grad}(\mathfrak J^{(n)})$ sur
l'anneau gradué $\operatorname{grad}(A)$ pour la filtration précédente
([I], p.~18-06, lemme)~; en vertu de la définition de la multiplication dans
$\operatorname{grad}(A)$,

\oldpage[0\textsubscript{I}]{65}

il suffira de prouver que pour tout $m$, $\overline S_m$ est un système de
générateurs du $A_0$-module
$\mathfrak J^{(m)}/\mathfrak J^{(m+1)}$, ou encore que $\mathfrak J_m^m$ est
engendré par les monômes de degré $m$ en les $a_{jm}$ ($1\leq j\leq r$).
Pour cela, il reste à montrer que $\mathfrak J_m$ est (en tant que
$A_m$-module) engendré par les monômes de degré $\leq m$ par rapport aux
$a_{jm}$~; la proposition étant évidente par définition pour $m=1$,
raisonnons par récurrence sur $m$, et soit $\mathfrak J'_m$ le sous-$A_m$-module
de $\mathfrak J_m$ engendré par ces monômes. La relation
$\mathfrak J_{m-1}=\mathfrak J_m/\mathfrak J_m^m$ et l'hypothèse de récurrence
prouvent que $\mathfrak J_m=\mathfrak J'_m+\mathfrak J_m^m$ d'où, puisque
$\mathfrak J_m^{m+1}=0$, l'on tire
$\mathfrak J_m^m={\mathfrak J'_m}^{\,m}$, et finalement
$\mathfrak J_m=\mathfrak J'_m$.

\begin{corollary}[7.2.8]
\label{0.7.2.8-fr}
Sous les conditions de (7.2.7), pour que $A$ soit noethérien, il faut et il
suffit que $A_0$ le soit.
\end{corollary}

Cela résulte aussitôt de (7.2.6).

\begin{proposition}[7.2.9]
\label{0.7.2.9-fr}
Supposons vérifiées les hypothèses de (7.2.7)~: pour tout entier $i$, soit
$M_i$ un $A_i$-module, et, pour $i\leq j$, soit
$v_{ij}:M_j\to M_i$ un $u_{ij}$-homomorphisme, tels que $(M_i,v_{ij})$ soit un
système projectif. Supposons en outre que $M_0$ soit un $A_0$-module de type
fini, que $v_{ij}$ soit surjectif et que son noyau soit
$\mathfrak J_j^{i+1}M_j$. Alors $M=\varprojlim M_i$ est un $A$-module de type
fini, et le noyau du $u_n$-homomorphisme surjectif $v_n:M\to M_n$ est
$\mathfrak J^{n+1}M$ (de sorte que $M_n$ s'identifie à
$M/\mathfrak J^{n+1}M=M\otimes_A(A/\mathfrak J^{n+1})$).
\end{proposition}

Soient $z_h=(z_{hk})_{k\geq0}$ un système de $s$ éléments de $M$ tels que les
$z_{h0}$ ($1\leq h\leq s$) forment un système de générateurs de $M_0$~; on va
montrer que les $z_h$ engendrent le $A$-module $M$. Le $A$-module $M$ est
séparé et complet pour la filtration formée par les $M^{(n)}$, où $M^{(n)}$
est l'ensemble des $y=(y_k)_{k\geq0}$ de $M$ tels que $y_k=0$ pour $k<n$~;
il est clair que l'on a $\mathfrak J^{(n)}M\subset M^{(n)}$ et que
$M^{(n)}/M^{(n+1)}=\mathfrak J_n^nM_n$. On est donc ramené à montrer que les
classes des $z_h$ mod. $M^{(0)}$ engendrent le module gradué
$\operatorname{grad}(M)$ (pour la filtration précédente) sur l'anneau gradué
$\operatorname{grad}(A)$ ([I], p.~18-06, lemme)~; pour cela, on constate
aisément qu'il suffit encore de prouver que les $z_{hn}$ ($1\leq h\leq s$)
engendrent le $A_n$-module $M_n$. On raisonne de nouveau par récurrence sur
$n$, la proposition étant évidente par définition pour $n=0$~; la relation
$M_{n-1}=M_n/\mathfrak J_n^nM_n$ et l'hypothèse de récurrence montrent que si
$M'_n$ est le sous-module de $M_n$ engendré par les $z_{hn}$, on a
$M_n=M'_n+\mathfrak J_n^nM_n$, et comme $\mathfrak J_n$ est nilpotent, cela
entraîne $M_n=M'_n$. Le même raisonnement de passage aux modules gradués
associés montre que l'application canonique de $\mathfrak J^{(n)}M$ dans
$M^{(n)}$ est surjective (donc bijective), autrement dit que
$\mathfrak J^{(n)}M=\mathfrak J^nM$ est le noyau de $M\to M_n$.

\begin{corollary}[7.2.10]
\label{0.7.2.10-fr}
Soit $(N_i,w_{ij})$ un second système projectif de $A_i$-modules vérifiant les
conditions de (7.2.9), et soit $N=\varprojlim N_i$. Il y a correspondance
biunivoque entre les systèmes projectifs $(h_i)$ de $A_i$-homomorphismes
$h_i:M_i\to N_i$ et les homomorphismes de $A$-modules $h:M\to N$ (qui sont
nécessairement continus pour les topologies $\mathfrak J$-adiques).
\end{corollary}

Il est clair que si $h:M\to N$ est un $A$-homomorphisme, on a
$h(\mathfrak J^nM)\subset\mathfrak J^nN$, d'où la continuité de $h$~; par
passage aux quotients, il correspond donc à $h$ un système projectif de
$A_i$-homomorphismes $h_i:M_i\to N_i$, dont $h$ est la limite projective,
d'où le corollaire.

\begin{remark}[7.2.11]
\label{0.7.2.11-fr}
Soit $A$ un anneau adique ayant un idéal de définition $\mathfrak J$ tel que
$\mathfrak J/\mathfrak J^2$ soit un $(A/\mathfrak J)$-module de type fini~;
il est clair que les $A_i=A/\mathfrak J^{i+1}$ vérifient

\oldpage[0\textsubscript{I}]{66}

les conditions de (7.2.7)~; comme $A$ est la limite projective des $A_i$, on
voit que la prop. (7.2.7) donne la description de \emph{tous} les anneaux
adiques du type considéré (et en particulier de tous les anneaux
\emph{adiques noethériens}).
\end{remark}

\begin{example}[7.2.12]
\label{0.7.2.12-fr}
Soient $B$ un anneau, $\mathfrak J$ un idéal de $B$ tel que
$\mathfrak J/\mathfrak J^2$ soit un module de type fini sur
$B/\mathfrak J$ (ou sur $B$, ce qui revient au même)~; posons
$A=\varprojlim_n B/\mathfrak J^{n+1}$~; $A$ est le séparé complété de $B$
muni de la topologie $\mathfrak J$-préadique. Si
$A_n=B/\mathfrak J^{n+1}$, il est immédiat que les $A_n$ vérifient les
conditions de (7.2.7)~; donc $A$ est un anneau adique et si
$\overline{\mathfrak J}$ est l'adhérence dans $A$ de l'image canonique de
$\mathfrak J$, $\overline{\mathfrak J}$ est un idéal de définition de $A$,
$\overline{\mathfrak J^n}$ est l'adhérence de l'image canonique de
$\mathfrak J^n$, $A/\overline{\mathfrak J^n}$ s'identifie à
$B/\mathfrak J^n$ et
$\overline{\mathfrak J}/\overline{\mathfrak J^2}$ est isomorphe à
$\mathfrak J/\mathfrak J^2$ en tant que
$(A/\overline{\mathfrak J})$-module. De même, si $N$ est tel que
$N/\mathfrak JN$ soit un $B$-module de type fini, et si on pose
$M_i=N/\mathfrak J^{i+1}N$, $M=\varprojlim M_i$ est un $A$-module de type
fini, isomorphe au séparé complété de $N$ pour la topologie
$\mathfrak J$-préadique, $\overline{\mathfrak J^n}M$ s'identifie à
l'adhérence de l'image canonique de $\mathfrak J^nN$, et
$M/\overline{\mathfrak J^n}M$ à $N/\mathfrak J^nN$.
\end{example}

\subsection{Anneaux préadiques noethériens}
\label{subsection:0.7.3-fr}

\begin{env}[7.3.1]
\label{0.7.3.1-fr}
Soient $A$ un anneau, $\mathfrak J$ un idéal de $A$, $M$ un $A$-module~;
nous désignerons par
$\widehat A=\varprojlim_n A/\mathfrak J^n$ (resp.
$\widehat M=\varprojlim_u M/\mathfrak J^nM$) le séparé complété de $A$
(resp. $M$) pour la topologie $\mathfrak J$-préadique. Soit
\[
  M'\xrightarrow{u}M\xrightarrow{v}M''\to0
\]
une suite exacte de $A$-modules~; comme
$M/\mathfrak J^nM=M\otimes_A(A/\mathfrak J^n)$, la suite
\[
  M'/\mathfrak J^nM'\xrightarrow{u_n}
  M/\mathfrak J^nM\xrightarrow{v_n}
  M''/\mathfrak J^nM''\to0
\]
est exacte pour tout $n$. En outre, comme
$v(\mathfrak J^nM)=\mathfrak J^nv(M)=\mathfrak J^nM''$,
$\widehat v=\varprojlim v_n$ est surjective (Bourbaki,
\emph{Top. gén.}, chap.~IX, 2\textsuperscript{e} éd., p.~60, cor.~2).
D'autre part, si $z=(z_k)$ est un élément du noyau de $\widehat v$, pour tout
entier $k$, il existe un $z'_k\in M'/\mathfrak J^kM'$ tel que
$u_k(z'_k)=z_k$~; on en conclut qu'il existe
$z'=(z'_n)\in\widehat{M'}$ tel que les $k$ premières composantes de
$\widehat u(z')$ coïncident avec celles de $z$~; autrement dit, l'image par
$\widehat u$ de $\widehat{M'}$ est \emph{dense} dans le noyau de
$\widehat v$.

Si on suppose $A$ \emph{noethérien}, il en est de même de $\widehat A$, en
vertu de (7.2.12), $\mathfrak J/\mathfrak J^2$ étant alors un $A$-module de
type fini. En outre~:
\end{env}

\begin{namedtheorem}[7.3.2]{Théorème de Krull}
\label{0.7.3.2-fr}
Soient $A$ un anneau noethérien, $\mathfrak J$ un idéal de $A$, $M$ un
$A$-module de type fini, $M'$ un sous-module de $M$~; alors la topologie
induite sur $M'$ par la topologie $\mathfrak J$-préadique de $M$ est identique
à la topologie $\mathfrak J$-préadique de $M'$.
\end{namedtheorem}

Cela résulte aussitôt du

\begin{namedlemma}[7.3.2.1]{Lemme d'Artin-Rees}
\label{0.7.3.2.1-fr}
Sous les hypothèses de (7.3.2), il existe un entier $p$ tel que, pour
$n\geq p$, on ait
\[
  M'\cap\mathfrak J^nM=\mathfrak J^{n-p}
  (M'\cap\mathfrak J^pM).
\]
\end{namedlemma}

Pour la démonstration, voir ([I], p.~2-04).

\oldpage[0\textsubscript{I}]{67}

\begin{corollary}[7.3.3]
\label{0.7.3.3-fr}
Sous les hypothèses de (7.3.2), l'application canonique
$M\otimes_A\widehat A\to\widehat M$ est bijective, et le foncteur
$M\otimes_A\widehat A$ est exact en $M$ dans la catégorie des $A$-modules de
type fini~; par suite, le séparé complété $\mathfrak J$-adique $\widehat A$
est un $A$-module plat (6.1.1).
\end{corollary}

Notons d'abord que dans la catégorie des $A$-modules de type fini,
$\widehat M$ est un foncteur \emph{exact} en $M$. Soit en effet
\[
  0\to M'\xrightarrow{u}M\xrightarrow{v}M''\to0
\]
une suite exacte~; on sait déjà que
$\widehat v:\widehat M\to\widehat{M''}$ est surjectif (7.3.1)~; d'autre part,
si $i$ est l'homomorphisme canonique $M\to\widehat M$, il résulte du th. de
Krull que l'adhérence dans $\widehat M$ de $i(u(M'))$ s'identifie au séparé
complété de $M'$ pour la topologie $\mathfrak J$-préadique~; donc
$\widehat u$ est injectif, et en vertu de (7.3.1) l'image de $\widehat u$ est
égale au noyau de $\widehat v$.

Cela étant, l'application canonique $M\otimes_A\widehat A\to\widehat M$
s'obtient en passant à la limite projective sur les applications
\[
  M\otimes_A\widehat A\to
  M\otimes_A(A/\mathfrak J^n)=M/\mathfrak J^nM.
\]
Il est clair que cette application est bijective lorsque $M$ est de la forme
$A^p$. Si $M$ est un $A$-module de type fini, on a une suite exacte
$A^p\to A^q\to M\to0$, d'où, en vertu de l'exactitude \emph{à droite} des
foncteurs $M\otimes_A\widehat A$ et $\widehat M$ (en $M$) dans la catégorie
des $A$-modules de type fini, le diagramme commutatif
\[
\begin{tikzcd}[column sep=large,row sep=large]
  A^p\otimes_A\widehat A \arrow[r] \arrow[d] &
  A^q\otimes_A\widehat A \arrow[r] \arrow[d] &
  M\otimes_A\widehat A \arrow[r] \arrow[d] &
  0 \\
  \widehat A^{\,p} \arrow[r] &
  \widehat A^{\,q} \arrow[r] &
  \widehat M \arrow[r] &
  0
\end{tikzcd}
\]
où les deux lignes sont exactes et les deux premières flèches verticales des
isomorphismes~; on en tire aussitôt la conclusion.

\begin{corollary}[7.3.4]
\label{0.7.3.4-fr}
Soient $A$ un anneau noethérien, $\mathfrak J$ un idéal de $A$, $M$, $N$ deux
$A$-modules de type fini~; on a des isomorphismes canoniques fonctoriels
\[
  (M\otimes_A N)^\wedge
  \xrightarrow{\sim}
  \widehat M\otimes_{\widehat A}\widehat N,
  \qquad
  \bigl(\operatorname{Hom}_A(M,N)\bigr)^\wedge
  \xrightarrow{\sim}
  \operatorname{Hom}_{\widehat A}(\widehat M,\widehat N).
\]
\end{corollary}

Cela résulte de (7.3.3), (6.2.1) et (6.2.2).

\begin{corollary}[7.3.5]
\label{0.7.3.5-fr}
Soient $A$ un anneau noethérien, $\mathfrak J$ un idéal de $A$. Les
conditions suivantes sont équivalentes~:
\begin{enumerate}
  \item[$a)$] $\mathfrak J$ est contenu dans le radical de $A$.
  \item[$b)$] $\widehat A$ est un $A$-module fidèlement plat (6.4.1).
  \item[$c)$] Tout $A$-module de type fini est séparé pour la topologie
  $\mathfrak J$-préadique.
  \item[$d)$] Tout sous-module d'un $A$-module de type fini est fermé pour la
  topologie $\mathfrak J$-préadique.
\end{enumerate}
\end{corollary}

Comme $\widehat A$ est un $A$-module plat, les conditions $b)$ et $c)$ sont
équivalentes, car $b)$ équivaut à dire que si $M$ est un $A$-module de type
fini, l'application canonique
$M\to\widehat M=M\otimes_A\widehat A$ est injective (6.6.1, $c)$). Il est
immédiat que $c)$ entraîne $d)$, car si $N$ est un sous-module d'un $A$-module
$M$ de type fini, $M/N$ est séparé pour la topologie $\mathfrak J$-préadique,
donc $N$ est fermé dans $M$. Montrons que $d)$ implique $a)$~: si
$\mathfrak m$ est un idéal maximal de $A$, $\mathfrak m$ est fermé dans $A$
pour la topologie $\mathfrak J$-préadique, donc
\[
  \mathfrak m=\bigcap_{p\geq0}(\mathfrak m+\mathfrak J^p),
\]
et comme $\mathfrak m+\mathfrak J^p$ est nécessairement égal à $A$ ou à
$\mathfrak m$, on a $\mathfrak m+\mathfrak J^p=\mathfrak m$ pour $p$ assez

\oldpage[0\textsubscript{I}]{68}

grand, d'où $\mathfrak J^p\subset\mathfrak m$ et
$\mathfrak J\subset\mathfrak m$ puisque $\mathfrak m$ est premier. Enfin,
$a)$ entraîne $b)$~: soit en effet $P$ l'adhérence de $\{0\}$ dans un
$A$-module $M$ de type fini, pour la topologie $\mathfrak J$-préadique~; en
vertu du th. de Krull (7.3.2), la topologie induite sur $P$ par la topologie
$\mathfrak J$-préadique de $M$ est la topologie $\mathfrak J$-préadique de
$P$, donc $\mathfrak JP=P$~; comme $P$ est de type fini, il résulte du lemme
de Nakayama que $P=0$ ($\mathfrak J$ étant contenu dans le radical de $A$).

On notera que les conditions de (7.3.5) sont remplies lorsque $A$ est un
\emph{anneau local noethérien} et $\mathfrak J\neq A$ un idéal quelconque de
$A$.

\begin{corollary}[7.3.6]
\label{0.7.3.6-fr}
Si $A$ est un anneau $\mathfrak J$-adique noethérien, tout $A$-module de type
fini est séparé et complet pour la topologie $\mathfrak J$-préadique.
\end{corollary}

Comme on a alors $\widehat A=A$, cela résulte aussitôt de (7.3.3).

On en conclut que la prop. (7.2.9) donne la description de \emph{tous} les
modules de type fini sur un anneau adique noethérien.

\begin{corollary}[7.3.7]
\label{0.7.3.7-fr}
Sous les hypothèses de (7.3.2), le noyau de l'application canonique
$M\to\widehat M=M\otimes_A\widehat A$ est l'ensemble des $x\in M$ annulés par
un élément de $1+\mathfrak J$.
\end{corollary}

En effet, pour que $x\in M$ appartienne à ce noyau, il faut et il suffit que
le séparé complété du sous-module $Ax$ se réduise à 0 (par (7.3.2)), autrement
dit que $x\in\mathfrak Jx$.

\subsection{Modules quasi-finis sur les anneaux locaux}
\label{subsection:0.7.4-fr}

\begin{definition}[7.4.1]
\label{0.7.4.1-fr}
Étant donné un anneau local $A$, d'idéal maximal $\mathfrak m$, on dit qu'un
$A$-module $M$ est quasi-fini (sur $A$) si $M/\mathfrak mM$ est de rang fini
sur le corps résiduel $k=A/\mathfrak m$.
\end{definition}

Lorsque $A$ est noethérien, le séparé complété $\widehat M$ de $M$ pour la
topologie $\mathfrak m$-préadique est alors un
$\widehat A$-\emph{module de type fini}~; en effet, comme
$\mathfrak m/\mathfrak m^2$ est alors un $A$-module de type fini, cela résulte
de (7.2.12) et de l'hypothèse sur $M/\mathfrak mM$.

En particulier, si on suppose de plus que $A$ est \emph{complet} et $M$
\emph{séparé} pour la topologie $\mathfrak m$-préadique (autrement dit,
$\bigcap_n\mathfrak m^nM=0$), $M$ lui-même est un
$A$-\emph{module de type fini}~: en effet, $\widehat M$ est alors un
$A$-module de type fini, et comme $M$ s'identifie à un sous-module de
$\widehat M$, $M$ est lui aussi de type fini (et d'ailleurs identique à son
complété en vertu de (7.3.6)).

\begin{proposition}[7.4.2]
\label{0.7.4.2-fr}
Soient $A$, $B$ deux anneaux locaux, $\mathfrak m$, $\mathfrak n$ leurs idéaux
maximaux, et supposons $B$ noethérien. Soient $\varphi:A\to B$ un
homomorphisme local, $M$ un $B$-module de type fini. Si $M$ est un $A$-module
quasi-fini, les topologies $\mathfrak m$-préadique et $\mathfrak n$-préadique
sur $M$ sont identiques, donc séparées.
\end{proposition}

Remarquons que par hypothèse $M/\mathfrak mM$ est de \emph{longueur finie} en
tant que $A$-module, donc aussi \emph{a fortiori} en tant que $B$-module. On en
conclut que $\mathfrak n$ est le \emph{seul idéal premier} de $B$ contenant
l'annulateur de $M/\mathfrak mM$~: en effet, on se ramène aussitôt, en vertu
de (I.7.4) et (I.7.2), au cas où $M/\mathfrak mM$ est \emph{simple}, donc
nécessairement isomorphe à $B/\mathfrak n$, et notre assertion est évidente
dans ce dernier cas. D'autre part, comme $M$ est un $B$-module de type fini,
les idéaux premiers qui contiennent l'annulateur de $M/\mathfrak mM$ sont ceux
qui contiennent $\mathfrak mB+\mathfrak b$, en désignant par $\mathfrak b$
l'annulateur du $B$-module $M$ (I.7.5). Comme $B$ est noethérien, on en conclut
([III], p.~127, cor.~4) que $\mathfrak mB+\mathfrak b$ est un idéal

\oldpage[0\textsubscript{I}]{69}

de définition de $B$, autrement dit qu'il existe $k>0$ tel que
$\mathfrak n^k\subset\mathfrak mB+\mathfrak b\subset\mathfrak n$~; par suite,
pour tout $h>0$
\[
  \mathfrak n^{hk}M
  \subset(\mathfrak mB+\mathfrak b)^hM
  =\mathfrak m^hM
  \subset\mathfrak n^hM,
\]
ce qui prouve que les topologies $\mathfrak m$-préadique et
$\mathfrak n$-préadique sur $M$ sont les mêmes~; la seconde est d'ailleurs
séparée en vertu de (7.3.5).

\begin{corollary}[7.4.3]
\label{0.7.4.3-fr}
Sous les hypothèses de (7.4.2), si en outre $A$ est noethérien et complet pour
la topologie $\mathfrak m$-préadique, $M$ est un $A$-module de type fini.
\end{corollary}

En effet, $M$ est alors séparé pour la topologie $\mathfrak m$-préadique, et
notre assertion résulte de la remarque faite dans (7.4.1).

\begin{env}[7.4.4]
\label{0.7.4.4-fr}
Le cas d'application de (7.4.2), qui est le plus important, est celui où $B$
lui-même est un $A$-\emph{module quasi-fini}, ce qui revient à dire que
$B/\mathfrak mB$ est une \emph{algèbre de rang fini sur}
$k=A/\mathfrak m$~; cette condition peut d'ailleurs se décomposer en la
conjonction des deux suivantes, en vertu de ce qui précède~:
\begin{enumerate}
  \item[(i)] $\mathfrak mB$ \emph{est un idéal de définition de} $B$~;
  \item[(ii)] $B/\mathfrak n$ \emph{est une extension de rang fini du corps}
  $A/\mathfrak m$.
\end{enumerate}
Lorsqu'il en est ainsi, tout $B$-module de type fini est évidemment un
$A$-module quasi-fini.
\end{env}

\begin{corollary}[7.4.5]
\label{0.7.4.5-fr}
Sous les hypothèses de (7.4.2), si $\mathfrak b$ est l'annulateur du
$B$-module $M$, $B/\mathfrak b$ est un $A$-module quasi-fini.
\end{corollary}

Supposons $M\neq0$ (sinon le corollaire est évident). On peut considérer $M$
comme un module sur l'anneau local noethérien $B/\mathfrak b$~; son annulateur
étant alors réduit à $0$, la démonstration de (7.4.2) montre que
$\mathfrak m(B/\mathfrak b)$ est un idéal de définition de
$B/\mathfrak b$. Par ailleurs, $M/\mathfrak nM$ est un espace vectoriel de
rang fini sur $A/\mathfrak m$, étant un quotient de $M/\mathfrak mM$, qui est
par hypothèse de rang fini sur $A/\mathfrak m$~; comme $M\neq0$, on a
$M\neq\mathfrak nM$ en vertu du lemme de Nakayama~; comme
$M/\mathfrak nM$ est un espace vectoriel $\neq0$ sur $B/\mathfrak n$, le fait
qu'il est de rang fini sur $A/\mathfrak m$ implique que $B/\mathfrak n$ est
aussi de rang fini sur $A/\mathfrak m$~; la conclusion résulte donc de (7.4.4)
appliqué à l'anneau $B/\mathfrak b$.

\subsection{Anneaux de séries formelles restreintes}
\label{subsection:0.7.5-fr}

\begin{env}[7.5.1]
\label{0.7.5.1-fr}
Soient $A$ un anneau topologique, linéairement topologisé, séparé et complet~;
soit $(\mathfrak J_\lambda)$ un système fondamental de voisinages de $0$ dans
$A$ formé d'idéaux (ouverts), de sorte que $A$ s'identifie canoniquement à
$\varprojlim A/\mathfrak J_\lambda$ (7.2.1). Pour tout $\lambda$, soit
$B_\lambda=(A/\mathfrak J_\lambda)[T_1,\ldots,T_r]$, où les $T_i$ sont des
indéterminées~; il est clair que les $B_\lambda$ forment un système projectif
d'anneaux discrets. Nous poserons
$\varprojlim B_\lambda=A\{T_1,\ldots,T_r\}$, et nous allons voir que cet anneau
topologique est indépendant du système fondamental d'idéaux
$(\mathfrak J_\lambda)$ considéré. De façon précise, soit $A'$ le sous-anneau
de l'anneau de séries formelles $A[\![T_1,\ldots,T_r]\!]$ formé des séries
formelles $\sum_\alpha c_\alpha T^\alpha$ (avec
$\alpha=(\alpha_1,\ldots,\alpha_r)\in\mathbf{N}^r$) telles que
$\lim_\alpha c_\alpha=0$ (suivant le filtre des complémentaires des parties
finies de $\mathbf{N}^r$)~; nous dirons que ces séries sont les séries
formelles \emph{restreintes} en les $T_i$, à coefficients dans $A$.

\oldpage[0\textsubscript{I}]{70}

Pour tout voisinage $V$ de $0$ dans $A$, soit $V'$ l'ensemble des
$x=\sum_\alpha c_\alpha T^\alpha\in A'$ telles que $c_\alpha\in V$ pour tout
$\alpha$. On vérifie aussitôt que les $V'$ forment un système fondamental de
voisinages de $0$ définissant sur $A'$ une topologie d'anneau \emph{séparée}~;
nous allons définir canoniquement un \emph{isomorphisme topologique} de
l'anneau $A\{T_1,\ldots,T_r\}$ sur $A'$. Pour tout
$\alpha\in\mathbf{N}^r$ et tout $\lambda$, soit
$\varphi_{\lambda,\alpha}$ l'application de
$(A/\mathfrak J_\lambda)[T_1,\ldots,T_r]$ dans
$A/\mathfrak J_\lambda$ qui, à tout polynôme du premier anneau, fait
correspondre le coefficient de $T^\alpha$ dans ce polynôme. Il est clair que
les $\varphi_{\lambda,\alpha}$ forment un système projectif d'homomorphismes
de $(A/\mathfrak J_\lambda)$-modules, dont la limite projective est un
homomorphisme continu
$\varphi_\alpha:A\{T_1,\ldots,T_r\}\to A$~; nous allons voir que, pour tout
$y\in A\{T_1,\ldots,T_r\}$, la série formelle
$\sum_\alpha\varphi_\alpha(y)T^\alpha$ est \emph{restreinte}. En effet, si
$y_\lambda$ est la composante de $y$ dans $B_\lambda$, et si on désigne par
$H_\lambda$ l'ensemble fini des $\alpha\in\mathbf{N}^r$ pour lesquels les
coefficients du polynôme $y_\lambda$ ne sont pas nuls, on a
$\varphi_{\lambda,\alpha}(y_\mu)\in\mathfrak J_\lambda$ pour
$\mathfrak J_\mu\subset\mathfrak J_\lambda$ et $\alpha\notin H_\lambda$, et
par passage à la limite, $\varphi_\alpha(y)\in\mathfrak J_\lambda$ pour
$\alpha\notin H_\lambda$. On définit donc un homomorphisme d'anneaux
$\varphi:A\{T_1,\ldots,T_r\}\to A'$ en posant
$\varphi(y)=\sum_\alpha\varphi_\alpha(y)T^\alpha$, et il est immédiat que
$\varphi$ est continu. Inversement, si $\theta_\lambda$ est l'homomorphisme
canonique $A\to A/\mathfrak J_\lambda$, pour tout élément
$z=\sum_\alpha c_\alpha T^\alpha\in A'$ et tout $\lambda$, il n'y a qu'un
nombre fini d'indices $\alpha$ tels que $\theta_\lambda(c_\alpha)\neq0$, et
par suite
$\psi_\lambda(z)=\sum_\alpha\theta_\lambda(c_\alpha)T^\alpha$ appartient à
$B_\lambda$~; les $\psi_\lambda$ sont continus et forment un système
projectif d'homomorphismes dont la limite projective est un homomorphisme
continu $\psi:A'\to A\{T_1,\ldots,T_r\}$~; il reste enfin à vérifier que
$\varphi\circ\psi$ et $\psi\circ\varphi$ sont les automorphismes identiques,
ce qui est immédiat.
\end{env}

\begin{env}[7.5.2]
\label{0.7.5.2-fr}
Nous identifierons $A\{T_1,\ldots,T_r\}$ à l'anneau $A'$ des séries formelles
restreintes au moyen des isomorphismes définis dans (7.5.1). Les isomorphismes
canoniques
\[
  ((A/\mathfrak J_\lambda)[T_1,\ldots,T_r])
  [T_{r+1},\ldots,T_s]
  \xrightarrow{\sim}
  (A/\mathfrak J_\lambda)[T_1,\ldots,T_s]
\]
définissent, par passage à la limite projective, un isomorphisme canonique
\[
  \Bigl[(A\{T_1,\ldots,T_r\})\{T_{r+1},\ldots,T_s\}
  \xrightarrow{\sim}
  A\{T_1,\ldots,T_s\}
\]
\end{env}

\begin{env}[7.5.3]
\label{0.7.5.3-fr}
Pour tout homomorphisme continu $u:A\to B$ de $A$ dans un anneau linéairement
topologisé $B$, séparé et complet, et tout système $(b_1,\ldots,b_r)$ de $r$
éléments de $B$, il existe \emph{un homomorphisme continu, et un seul},
$\overline u:A\{T_1,\ldots,T_r\}\to B$, tel que $\overline u(a)=u(a)$ pour
tout $a\in A$ et $\overline u(T_j)=b_j$ pour $1\leq j\leq r$. Il suffit en
effet de prendre
\[
  \overline u\Bigl(\sum_\alpha c_\alpha T^\alpha\Bigr)
  =\sum_\alpha u(c_\alpha)b_1^{\alpha_1}\cdots b_r^{\alpha_r}~;
\]
les vérifications du fait que la famille
$(u(c_\alpha)b_1^{\alpha_1}\cdots b_r^{\alpha_r})$ est sommable dans $B$ et
de la continuité de $\overline u$ sont immédiates et laissées au lecteur. On
notera que cette propriété (pour $B$ et les $b_j$ arbitraires)
\emph{caractérise} l'anneau topologique $A\{T_1,\ldots,T_r\}$ à un
isomorphisme unique près.
\end{env}

\begin{proposition}[7.5.4]
\label{0.7.5.4-fr}
\begin{enumerate}
  \item[(i)] Si $A$ est un anneau admissible, il en est de même de
  $A'=A\{T_1,\ldots,T_r\}$.
  \item[(ii)] Soient $A$ un anneau adique, $\mathfrak J$ un idéal de définition
  de $A$ tel que $\mathfrak J/\mathfrak J^2$ soit de type fini

\oldpage[0\textsubscript{I}]{71}

  sur $A/\mathfrak J$. Si on pose $\mathfrak J'=\mathfrak JA'$, $A'$ est alors
  un anneau $\mathfrak J'$-adique, et
  $\mathfrak J'/\mathfrak J'^{\,2}$ est de type fini sur
  $A'/\mathfrak J'$. Si, en outre, $A$ est noethérien, il en est de même de
  $A'$.
\end{enumerate}
\end{proposition}

\begin{enumerate}
  \item[(i)] Si $\mathfrak J$ est un idéal de $A$, $\mathfrak J'$ l'idéal de
  $A'$ formé des $\sum_\alpha c_\alpha T^\alpha$ tels que
  $c_\alpha\in\mathfrak J$ pour tout $\alpha$, alors
  $(\mathfrak J')^n\subset(\mathfrak J^n)'$~; si $\mathfrak J$ est un idéal de
  définition de $A$, $\mathfrak J'$ est donc un idéal de définition de $A'$.

  \item[(ii)] Posons $A_i=A/\mathfrak J^{i+1}$, et pour $i\leq j$, soit
  $u_{ij}$ l'homomorphisme canonique
  $A/\mathfrak J^{j+1}\to A/\mathfrak J^{i+1}$~; posons
  $A'_i=A_i[T_1,\ldots,T_r]$, et soit $u'_{ij}$ l'homomorphisme
  $A'_j\to A'_i$ ($i\leq j$) obtenu en appliquant $u_{ij}$ aux coefficients
  des polynômes de $A'_j$. Montrons que le système projectif
  $(A'_i,u'_{ij})$ vérifie les conditions de (7.2.7)~; comme $\mathfrak J'$ est
  le noyau de $A'\to A'_0$, cela démontrera la première assertion de (ii). Or,
  il est clair que les $u'_{ij}$ sont surjectifs~; le noyau $\mathfrak J'_i$
  de $u'_{0i}$ est l'ensemble des polynômes de
  $A_i[T_1,\ldots,T_r]$ dont les coefficients sont dans
  $\mathfrak J/\mathfrak J^{i+1}$~; en particulier, $\mathfrak J'_1$ est
  l'ensemble des polynômes de $A_1[T_1,\ldots,T_r]$ dont les coefficients
  sont dans $\mathfrak J/\mathfrak J^2$. Comme $\mathfrak J/\mathfrak J^2$
  est de type fini sur $A_1=A/\mathfrak J^2$, on voit déjà que
  $\mathfrak J'_1/\mathfrak J_1'^{\,2}$ est un module de type fini sur $A'_1$
  (ou, ce qui revient au même, sur
  $A'_0=A'_1/\mathfrak J'_1$). Montrons que le noyau de $u'_{ij}$ est
  $\mathfrak J_j'^{\,i+1}$. Il est évident que
  $\mathfrak J_j'^{\,i+1}$ est contenu dans ce noyau. D'autre part, soient
  $a_1,\ldots,a_m$ des éléments de $\mathfrak J$ dont les classes
  mod.~$\mathfrak J^2$ engendrent $\mathfrak J/\mathfrak J^2$~; on vérifie
  immédiatement que les classes mod.~$\mathfrak J^{j+1}$ des monômes de degré
  $\leq j$ en les $a_k$ ($1\leq k\leq m$) engendrent
  $\mathfrak J/\mathfrak J^{j+1}$, et les classes des monômes de degré $>i$
  et $\leq j$ engendrent donc
  $\mathfrak J^{i+1}/\mathfrak J^{j+1}$~; un monôme en les $T_k$ ayant un tel
  élément comme coefficient est donc produit de $i+1$ éléments de
  $\mathfrak J'_j$, ce qui établit notre assertion. Enfin, si $A$ est
  noethérien, il en est de même de
  $A'/\mathfrak J'=(A/\mathfrak J)[T_1,\ldots,T_r]$, donc $A'$ est noethérien
  (7.2.8).
\end{enumerate}

\begin{proposition}[7.5.5]
\label{0.7.5.5-fr}
Soient $A$ un anneau noethérien $\mathfrak J$-adique, $B$ un anneau topologique
admissible, $\varphi:A\to B$ un homomorphisme continu, faisant de $B$ une
$A$-algèbre. Les conditions suivantes sont équivalentes~:
\begin{enumerate}
  \item[a)] $B$ est noethérien et $\mathfrak JB$-adique, et
  $B/\mathfrak JB$ est une algèbre de type fini sur $A/\mathfrak J$.

  \item[b)] $B$ est topologiquement $A$-isomorphe à $\varprojlim B_n$, où
  $B_n=B_m/\mathfrak J^{n+1}B_m$ pour $m\geq n$, et $B_1$ est une algèbre de
  type fini sur $A_1=A/\mathfrak J^2$.

  \item[c)] $B$ est topologiquement $A$-isomorphe à un quotient d'une algèbre
  de la forme\linebreak $A\{T_1,\ldots,T_r\}$ par un idéal (nécessairement fermé en vertu
  de (7.3.6) et (7.5.4 (ii))).
\end{enumerate}
\end{proposition}

Comme $A$ est noethérien, il en est de même de
$A'=A\{T_1,\ldots,T_r\}$ (7.5.4), donc c) implique que $B$ est noethérien~;
comme $\mathfrak J'=\mathfrak JA'$ est un voisinage ouvert de 0 dans $A'$ tel
que les $\mathfrak J'^{n}$ forment un système fondamental de voisinages de 0,
les images $\mathfrak J^nB$ des $\mathfrak J'^{n}$ forment un système
fondamental de voisinages de 0 dans $B$, et comme on sait que $B$ est séparé et
complet, $B$ est un anneau $\mathfrak JB$-adique. Enfin, $B/\mathfrak JB$ est
une algèbre (sur $A/\mathfrak J$) quotient de
$A'/\mathfrak JA'=(A/\mathfrak J)[T_1,\ldots,T_r]$, donc est de type fini, ce
qui achève de prouver que c) entraîne a).

Si $B$ est $\mathfrak JB$-adique et noethérien, $B$ est isomorphe à
$\varprojlim B_n$, où $B_n=B/\mathfrak J^{n+1}B$ (7.2.11), et
$\mathfrak JB/\mathfrak J^2B$ est un module de type fini sur
$B/\mathfrak JB$. Soit $(a_j)_{1\leq j\leq s}$ un système de générateurs du
$(B/\mathfrak JB)$-module $\mathfrak JB/\mathfrak J^2B$, et
$(c_i)_{1\leq i\leq r}$ un système d'éléments de $B/\mathfrak J^2B$ dont les
classes

\oldpage[0\textsubscript{I}]{72}

mod.~$\mathfrak JB/\mathfrak J^2B$ forment un système de générateurs de la
$(A/\mathfrak J)$-algèbre $B/\mathfrak JB$~; on voit aussitôt que les $c_ia_j$
forment un système de générateurs de la $(A/\mathfrak J^2)$-algèbre
$B/\mathfrak J^2B$, donc a) implique b).

Reste à prouver que b) entraîne c). L'hypothèse entraîne que $B_1$ est un
anneau noethérien, et comme $B_1=B_2/\mathfrak J^2B_2$, on a
$\mathfrak J^2B_1=0$, donc
$\mathfrak JB_1=\mathfrak JB_1/\mathfrak J^2B_1$ est un $B_0$-module de type
fini. Les conditions de (7.2.7) sont donc vérifiées par le système projectif
$(B_n)$ et $B$ est un anneau $\mathfrak JB$-adique. Soit
$(c_i)_{1\leq i\leq r}$ un système fini d'éléments de $B$ dont les classes
mod.~$\mathfrak JB$ engendrent la $(A/\mathfrak J)$-algèbre
$B/\mathfrak JB$, et dont les combinaisons linéaires à coefficients dans
$\mathfrak J$ sont telles que leurs classes mod.~$\mathfrak J^2B$ engendrent
le $B_0$-module $\mathfrak JB/\mathfrak J^2B$. Il existe un
$A$-homomorphisme continu $u$ de $A'=A\{T_1,\ldots,T_r\}$ dans $B$ qui se
réduit à $\varphi$ dans $A$ et est tel que $u(T_i)=c_i$ pour $1\leq i\leq r$
(7.5.3)~; si nous prouvons que $u$ est \emph{surjectif}, c) sera établi, car de
$u(A')=B$ on déduira $u(\mathfrak J^nA')=\mathfrak J^nB$, autrement dit $u$
sera un morphisme strict d'anneaux topologiques et $B$ sera donc isomorphe à un
quotient de $A'$ par un idéal fermé. Or, comme $B$ est complet pour la topologie
$\mathfrak JB$-adique, il suffit ([I], p.~18-07) de montrer que
l'homomorphisme $\operatorname{grad}(A')\to\operatorname{grad}(B)$ déduit
canoniquement de $u$ pour les filtrations $\mathfrak J$-adiques sur $A'$ et
$B$, est surjectif. Mais par définition, les homomorphismes
$A'/\mathfrak JA'\to B/\mathfrak JB$ et
$\mathfrak JA'/\mathfrak J^2A'\to\mathfrak JB/\mathfrak J^2B$ déduits de $u$
sont surjectifs~; par récurrence sur $n$, on en déduit aussitôt qu'il en est de
même de $\mathfrak JA'/\mathfrak J^nA'\to\mathfrak JB/\mathfrak J^nB$, et
\emph{a fortiori} de
$\mathfrak J^nA'/\mathfrak J^{n+1}A'\to
\mathfrak J^nB/\mathfrak J^{n+1}B$, ce qui achève la démonstration.

\subsection{Anneaux complets de fractions}
\label{subsection:0.7.6-fr}

\begin{env}[7.6.1]
\label{0.7.6.1-fr}
Soient $A$ un anneau linéairement topologisé, $(\mathfrak J_\lambda)$ un
système fondamental de voisinages de 0 dans $A$ formé d'idéaux, $S$ une partie
multiplicative de $A$. Soit $u_\lambda$ l'homomorphisme canonique
$A\to A_\lambda=A/\mathfrak J_\lambda$, et pour
$\mathfrak J_\mu\subset\mathfrak J_\lambda$, soit $u_{\lambda\mu}$
l'homomorphisme canonique $A_\mu\to A_\lambda$. Posons
$S_\lambda=u_\lambda(S)$, de sorte que
$u_{\lambda\mu}(S_\mu)=S_\lambda$. Les $u_{\lambda\mu}$ donnent
canoniquement des homomorphismes surjectifs
$S_\mu^{-1}A_\mu\to S_\lambda^{-1}A_\lambda$, pour lesquels ces anneaux
forment un système projectif~; désignons par $A\{S^{-1}\}$ la limite
projective de ce système. Cette définition ne dépend pas du système fondamental
de voisinages $(\mathfrak J_\lambda)$ choisi~; en effet~:
\end{env}

\begin{proposition}[7.6.2]
\label{0.7.6.2-fr}
L'anneau $A\{S^{-1}\}$ est topologiquement isomorphe au séparé complété de
l'anneau $S^{-1}A$ pour la topologie dont un système fondamental de voisinages
de 0 est formé des $S^{-1}\mathfrak J_\lambda$.
\end{proposition}

En effet, si $v_\lambda$ est l'homomorphisme canonique
$S^{-1}A\to S_\lambda^{-1}A_\lambda$ déduit de $u_\lambda$, le noyau de
$v_\lambda$ est $S^{-1}\mathfrak J_\lambda$ et $v_\lambda$ est surjectif,
d'où la proposition (7.2.1).

\begin{corollary}[7.6.3]
\label{0.7.6.3-fr}
Si $S'$ est l'image canonique de $S$ dans le séparé complété $\widehat A$ de
$A$, $A\{S^{-1}\}$ s'identifie canoniquement à
$\widehat A\{{S'}^{-1}\}$.
\end{corollary}

On notera que même si $A$ est séparé et complet, il n'en est pas de même de
$S^{-1}A$ pour la topologie définie par les
$S^{-1}\mathfrak J_\lambda$, comme on le voit par exemple en prenant pour $S$
l'ensemble des $f^n$ ($n\geq0$), où $f$ est topologiquement nilpotent et non
nilpotent~: en effet, $S^{-1}A$ n'est pas réduit à 0 et d'autre part, pour tout
$\lambda$ il existe $n$ tel que $f^n\in\mathfrak J_\lambda$, donc
$1=f^n/f^n\in S^{-1}\mathfrak J_\lambda$ et
$S^{-1}\mathfrak J_\lambda=S^{-1}A$.

\oldpage[0\textsubscript{I}]{73}

\begin{corollary}[7.6.4]
\label{0.7.6.4-fr}
Si, dans $A$, $0$ n'est pas adhérent à $S$, l'anneau $A\{S^{-1}\}$ n'est pas
réduit à $0$.
\end{corollary}

En effet, $0$ n'est pas adhérent à $\{1\}$ dans l'anneau $S^{-1}A$~; sinon,
on aurait $1\in S^{-1}\mathfrak J_\lambda$ pour tout idéal ouvert
$\mathfrak J_\lambda$ de $A$, et il en résulterait que
$\mathfrak J_\lambda\cap S\neq\varnothing$ pour tout $\lambda$, contrairement
à l'hypothèse.

\begin{env}[7.6.5]
\label{0.7.6.5-fr}
Nous dirons que $A\{S^{-1}\}$ est l'\emph{anneau complet des fractions} de $A$
ayant leurs dénominateurs dans $S$. Avec les notations précédentes, il est clair
que l'image réciproque de $S^{-1}\mathfrak J_\lambda$ dans $A$ contient
$\mathfrak J_\lambda$, donc l'application canonique $A\to S^{-1}A$ est
continue, et si on la compose avec l'application canonique
$S^{-1}A\to A\{S^{-1}\}$, on obtient un homomorphisme canonique continu
$A\to A\{S^{-1}\}$, limite projective des homomorphismes
$A\to S_\lambda^{-1}A_\lambda$.
\end{env}

\begin{env}[7.6.6]
\label{0.7.6.6-fr}
Le couple formé de $A\{S^{-1}\}$ et de l'application canonique
$A\to A\{S^{-1}\}$ est caractérisé par la \emph{propriété universelle}
suivante~: tout homomorphisme continu $u$ de $A$ dans un anneau linéairement
topologisé $B$, séparé et complet, tel que $u(S)$ soit formé d'éléments
inversibles dans $B$, se factorise d'une seule manière en
$A\to A\{S^{-1}\}\xrightarrow{u'}B$, où $u'$ est continu. En effet, $u$ se
factorise d'une seule manière en $A\to S^{-1}A\xrightarrow{v'}B$~; comme pour
tout idéal ouvert $\mathfrak K$ de $B$, $u^{-1}(\mathfrak K)$ contient un
$\mathfrak J_\lambda$, ${v'}^{-1}(\mathfrak K)$ contient nécessairement
$S^{-1}\mathfrak J_\lambda$, donc $v'$ est continu~; puisque $B$ est séparé et
complet, $v'$ se factorise d'une seule manière en
$S^{-1}A\to A\{S^{-1}\}\xrightarrow{u'}B$, où $u'$ est continu~; d'où notre
assertion.
\end{env}

\begin{env}[7.6.7]
\label{0.7.6.7-fr}
Soient $B$ un second anneau linéairement topologisé, $T$ une partie
multiplicative de $B$, $\varphi:A\to B$ un homomorphisme continu tel que
$\varphi(S)\subset T$. D'après ce qui précède, l'homomorphisme continu
$A\xrightarrow{\varphi}B\to B\{T^{-1}\}$ se factorise de façon unique en
$A\to A\{S^{-1}\}\xrightarrow{\varphi'}B\{T^{-1}\}$, où $\varphi'$ est
continu. En particulier, si $B=A$ et si $\varphi$ est l'identité, on voit que
pour $S\subset T$ on a un homomorphisme continu
$\rho^{T,S}:A\{S^{-1}\}\to A\{T^{-1}\}$ obtenu par passage au séparé complété
à partir de $S^{-1}A\to T^{-1}A$~; si $U$ est une troisième partie
multiplicative de $A$ telle que $S\subset T\subset U$, on a
$\rho^{U,S}=\rho^{U,T}\circ\rho^{T,S}$.
\end{env}

\begin{env}[7.6.8]
\label{0.7.6.8-fr}
Soient $S_1$, $S_2$ deux parties multiplicatives de $A$, et soit $S_2'$ l'image
canonique de $S_2$ dans $A\{S_1^{-1}\}$~; on a alors un isomorphisme
topologique canonique
$A\{(S_1S_2)^{-1}\}\xrightarrow{\sim}A\{S_1^{-1}\}\{{S_2'}^{-1}\}$,
comme on le voit en partant de l'isomorphisme canonique
$(S_1S_2)^{-1}A\xrightarrow{\sim}{S_2''}^{-1}(S_1^{-1}A)$ (où $S_2''$ est
l'image canonique de $S_2$ dans $S_1^{-1}A$), qui est bicontinu.
\end{env}

\begin{env}[7.6.9]
\label{0.7.6.9-fr}
Soit $\mathfrak a$ un idéal \emph{ouvert} de $A$~; on peut supposer que
$\mathfrak J_\lambda\subset\mathfrak a$ pour tout $\lambda$, et par suite
$S^{-1}\mathfrak J_\lambda\subset S^{-1}\mathfrak a$ dans l'anneau
$S^{-1}A$, autrement dit, $S^{-1}\mathfrak a$ est un idéal \emph{ouvert} de
$S^{-1}A$~; nous désignerons par $\mathfrak a\{S^{-1}\}$ son séparé complété,
égal à $\varprojlim(S^{-1}\mathfrak a/S^{-1}\mathfrak J_\lambda)$, qui est un
idéal \emph{ouvert} de $A\{S^{-1}\}$, isomorphe à l'adhérence de l'image
canonique de $S^{-1}\mathfrak a$. En outre, \emph{l'anneau discret
$A\{S^{-1}\}/\mathfrak a\{S^{-1}\}$ est canoniquement isomorphe à
$S^{-1}A/S^{-1}\mathfrak a=S^{-1}(A/\mathfrak a)$}. Inversement, si
$\mathfrak a'$ est un idéal ouvert de $A\{S^{-1}\}$, $\mathfrak a'$ contient
un idéal de la forme $\mathfrak J_\lambda\{S^{-1}\}$, donc est l'image
réciproque d'un idéal de $S^{-1}A/S^{-1}\mathfrak J_\lambda$, qui est
nécessairement (1.2.6) de la forme $S^{-1}\mathfrak a$, où
$\mathfrak a\supset\mathfrak J_\lambda$. On en conclut que l'on a
$\mathfrak a'=\mathfrak a\{S^{-1}\}$. En particulier (1.2.6)~:
\end{env}

\begin{proposition}[7.6.10]
\label{0.7.6.10-fr}
L'application $\mathfrak p\mapsto\mathfrak p\{S^{-1}\}$ est une bijection
croissante de l'ensemble des idéaux premiers ouverts $\mathfrak p$ de $A$ tels
que $\mathfrak p\cap S=\varnothing$ sur l'ensemble des idéaux premiers ouverts

\oldpage[0\textsubscript{I}]{74}

de $A\{S^{-1}\}$~; en outre, le corps des fractions de
$A\{S^{-1}\}/\mathfrak p\{S^{-1}\}$ est canoniquement isomorphe à celui de
$A/\mathfrak p$.
\end{proposition}

\begin{proposition}[7.6.11]
\label{0.7.6.11-fr}
\begin{enumerate}
  \item[\rm(i)] Si $A$ est un anneau admissible, il en est de même de
    $A'=A\{S^{-1}\}$ et pour tout idéal de définition $\mathfrak J$ de $A$,
    $\mathfrak J'=\mathfrak J\{S^{-1}\}$ est un idéal de définition de $A'$.
  \item[\rm(ii)] Soient $A$ un anneau adique, $\mathfrak J$ un idéal de
    définition de $A$ tel que $\mathfrak J/\mathfrak J^2$ soit de type fini
    sur $A/\mathfrak J$~; alors $A'$ est un anneau $\mathfrak J'$-adique et
    $\mathfrak J'/\mathfrak J'^2$ est de type fini sur
    $A'/\mathfrak J'$. Si en outre $A$ est noethérien, il en est de même de
    $A'$.
\end{enumerate}
\end{proposition}

\begin{enumerate}
  \item[\rm(i)] Si $\mathfrak J$ est un idéal de définition dans $A$, il est
    clair que $S^{-1}\mathfrak J$ est un idéal de définition dans l'anneau
    topologique $S^{-1}A$, car on a
    $(S^{-1}\mathfrak J)^n=S^{-1}\mathfrak J^n$. Soit $A''$ l'anneau séparé
    associé à $S^{-1}A$, $\mathfrak J''$ l'image de $S^{-1}\mathfrak J$ dans
    $A''$~; l'image de $S^{-1}\mathfrak J^n$ est $\mathfrak J''^n$, donc
    $\mathfrak J''^n$ tend vers 0 dans $A''$~; comme $\mathfrak J'$ est
    l'adhérence de $\mathfrak J''$ dans $A'$, $\mathfrak J'^n$ est contenu
    dans l'adhérence de $\mathfrak J''^n$, donc tend vers 0 dans $A'$.
  \item[\rm(ii)] Posons $A_i=A/\mathfrak J^{i+1}$, et pour $i\leq j$, soit
    $u_{ij}$ l'homomorphisme canonique
    $A/\mathfrak J^{j+1}\to A/\mathfrak J^{i+1}$~; soit $S_i$ l'image
    canonique de $S$ dans $A_i$, et posons $A_i'=S_i^{-1}A_i$~; soit enfin
    $u_{ij}':A_j'\to A_i'$ l'homomorphisme déduit canoniquement de $u_{ij}$.
    Montrons que le système projectif $(A_i',u_{ij}')$ vérifie les conditions
    de la prop.~(7.2.7)~: il est clair que les $u_{ij}'$ sont surjectifs~;
    d'autre part, le noyau de $u_{ij}'$ est
    $S_j^{-1}(\mathfrak J^{i+1}/\mathfrak J^{j+1})$ (1.3.2), égal à
    $\mathfrak J_j'^{i+1}$, où
    $\mathfrak J_j'=S_j^{-1}(\mathfrak J/\mathfrak J^{j+1})$~; enfin,
    $\mathfrak J_1'/\mathfrak J_1'^2=S_1^{-1}(\mathfrak J/\mathfrak J^2)$,
    et comme $\mathfrak J/\mathfrak J^2$ est de type fini sur
    $A/\mathfrak J^2$, $\mathfrak J_1'/\mathfrak J_1'^2$ est de type fini
    sur $A_1'$. Enfin, si $A$ est noethérien, il en est de même de
    $A_0'=S_0^{-1}(A/\mathfrak J)$, ce qui achève de démontrer la proposition
    (7.2.8).
\end{enumerate}

\begin{corollary}[7.6.12]
\label{0.7.6.12-fr}
Sous les hypothèses de (7.6.11, (ii)), on a
$(\mathfrak J\{S^{-1}\})^n=\mathfrak J^n\{S^{-1}\}$.
\end{corollary}

Cela résulte en effet de (7.2.7) et de la démonstration de (7.6.11).

\begin{proposition}[7.6.13]
\label{0.7.6.13-fr}
Soient $A$ un anneau adique noethérien, $S$ une partie multiplicative de $A$~;
alors $A\{S^{-1}\}$ est un $A$-module plat.
\end{proposition}

En effet, si $\mathfrak J$ est un idéal de définition de $A$,
$A\{S^{-1}\}$ est le séparé complété de l'anneau noethérien $S^{-1}A$ muni
de la topologie $S^{-1}\mathfrak J$-préadique~; par suite (7.3.3)
$A\{S^{-1}\}$ est un $S^{-1}A$-module plat~; comme $S^{-1}A$ est un
$A$-module plat (6.3.1), la proposition résulte de la transitivité de la
platitude (6.2.1).

\begin{corollary}[7.6.14]
\label{0.7.6.14-fr}
Sous les hypothèses de (7.6.13), soit $S'\subset S$ une seconde partie
multiplicative de $A$~; alors $A\{S^{-1}\}$ est un
$A\{{S'}^{-1}\}$-module plat.
\end{corollary}

En effet (7.6.8), $A\{S^{-1}\}$ s'identifie canoniquement à
$A\{{S'}^{-1}\}\{S_0^{-1}\}$, où $S_0$ est l'image canonique de $S$ dans
$A\{{S'}^{-1}\}$, et $A\{{S'}^{-1}\}$ est noethérien (7.6.11).

\begin{env}[7.6.15]
\label{0.7.6.15-fr}
Pour tout élément $f$ d'un anneau linéairement topologisé $A$, nous
désignerons par $A_{\{f\}}$ l'anneau complet de fractions $A\{S_f^{-1}\}$,
où $S_f$ est l'ensemble multiplicatif des $f^n$ ($n\geq 0$)~; pour tout
idéal ouvert $\mathfrak a$ de $A$, nous écrirons $\mathfrak a_{\{f\}}$ au
lieu de $\mathfrak a\{S_f^{-1}\}$. Si $g$ est un second élément de $A$, on a
un homomorphisme canonique continu $A_{\{f\}}\to A_{\{fg\}}$ (7.6.7).
Lorsque $f$ parcourt une partie multiplicative $S$ de $A$, les $A_{\{f\}}$
forment donc un système inductif filtrant d'anneaux pour les homomorphismes
précédents~; nous poserons
$A_{\{S\}}=\varinjlim_{f\in S}A_{\{f\}}$. Pour tout $f\in S$, on a un
homomorphisme $A_{\{f\}}\to A\{S^{-1}\}$ (7.6.7), et

\oldpage[0\textsubscript{I}]{75}

ces homomorphismes forment un système inductif~; par passage à la limite
inductive, ils définissent donc un homomorphisme canonique
$A_{\{S\}}\to A\{S^{-1}\}$.
\end{env}

\begin{proposition}[7.6.16]
\label{0.7.6.16-fr}
Si $A$ est un anneau noethérien, $A\{S^{-1}\}$ est un module plat sur
$A_{\{S\}}$.
\end{proposition}

En effet (7.6.14), $A\{S^{-1}\}$ est plat sur chacun des anneaux
$A_{\{f\}}$ pour $f\in S$, et la conclusion résulte de (6.2.3).

\begin{proposition}[7.6.17]
\label{0.7.6.17-fr}
Soit $\mathfrak p$ un idéal premier ouvert dans un anneau admissible $A$, et
soit $S=A-\mathfrak p$. Alors les anneaux $A\{S^{-1}\}$ et $A_{\{S\}}$ sont
des anneaux locaux, l'homomorphisme canonique
$A_{\{S\}}\to A\{S^{-1}\}$ est local et les corps résiduels de $A_{\{S\}}$
et $A\{S^{-1}\}$ sont canoniquement isomorphes au corps des fractions de
$A/\mathfrak p$.
\end{proposition}

En effet, soit $\mathfrak J\subset\mathfrak p$ un idéal de définition de
$A$~; on a $S^{-1}\mathfrak J\subset S^{-1}\mathfrak p=\mathfrak p
A_{\mathfrak p}$, donc $A_{\mathfrak p}/S^{-1}\mathfrak J$ est un anneau
local~; on conclut de (7.1.12), (7.6.9) et (7.6.11, (i)) que
$A\{S^{-1}\}$ est un anneau local. Posons
$\mathfrak m=\varinjlim_{f\in S}\mathfrak p_{\{f\}}$, qui est un idéal de
$A_{\{S\}}$~; nous allons voir que tout élément de $A_{\{S\}}$ n'appartenant
pas à $\mathfrak m$ est inversible. En effet, un tel élément est l'image dans
$A_{\{S\}}$ d'un élément $z\in A_{\{f\}}$ n'appartenant pas à
$\mathfrak p_{\{f\}}$, pour un $f\in S$~; son image canonique $z_0$ dans
$A_{\{f\}}/\mathfrak J_{\{f\}}=S_f^{-1}(A/\mathfrak J)$ n'appartient donc
pas à $S_f^{-1}(\mathfrak p/\mathfrak J)$ (7.6.9), ce qui signifie que
$z_0=\overline{x}/\overline{f}^k$, où $x\notin\mathfrak p$ et
$\overline{x},\overline{f}$ sont les classes de $x,f$ mod.~$\mathfrak J$.
Comme $x\in S$, on a $g=xf\in S$, et dans $S_g^{-1}A$,
$y_0=x^{k+1}/g^k$, image canonique de $x/f^k\in S_f^{-1}A$, admet un
inverse $x^{k-1}f^{2k}/g^k$. Cela entraîne \emph{a fortiori} que l'image de
$y_0$ dans $S_g^{-1}A/S_g^{-1}\mathfrak J$ est inversible, donc (7.6.9 et
7.1.12) l'image canonique $y$ de $z$ dans $A_{\{g\}}$ est inversible~;
l'image de $z$ dans $A_{\{S\}}$ (égale à celle de $y$) est par suite
inversible. On voit donc que $A_{\{S\}}$ est un anneau local d'idéal maximal
$\mathfrak m$~; en outre, l'image de $\mathfrak p_{\{f\}}$ dans
$A\{S^{-1}\}$ est contenue dans l'idéal maximal $\mathfrak p\{S^{-1}\}$ de
cet anneau~; \emph{a fortiori}, l'image de $\mathfrak m$ dans
$A\{S^{-1}\}$ est contenue dans $\mathfrak p\{S^{-1}\}$, donc
l'homomorphisme canonique $A_{\{S\}}\to A\{S^{-1}\}$ est local. Enfin,
comme tout élément de $A\{S^{-1}\}/\mathfrak p\{S^{-1}\}$ est image d'un
élément d'un anneau $S_f^{-1}A$ pour un $f\in S$ convenable,
l'homomorphisme $A_{\{S\}}\to A\{S^{-1}\}/\mathfrak p\{S^{-1}\}$ est
surjectif, et donne donc par passage aux quotients un isomorphisme des corps
résiduels.

\begin{corollary}[7.6.18]
\label{0.7.6.18-fr}
Sous les hypothèses de (7.6.17), si on suppose de plus que $A$ est un anneau
adique noethérien, les anneaux locaux $A\{S^{-1}\}$ et $A_{\{S\}}$ sont
noethériens, et $A\{S^{-1}\}$ est un $A_{\{S\}}$-module fidèlement plat.
\end{corollary}

On sait déjà (7.6.11, (ii)) que $A\{S^{-1}\}$ est noethérien et
$A_{\{S\}}$-plat (7.6.16)~; comme l'homomorphisme
$A_{\{S\}}\to A\{S^{-1}\}$ est local, on en conclut que $A\{S^{-1}\}$ est
un $A_{\{S\}}$-module fidèlement plat (6.6.2), et par suite que $A_{\{S\}}$
est noethérien (6.5.2).

\subsection{Produits tensoriels complétés}
\label{subsection:0.7.7-fr}

\begin{env}[7.7.1]
\label{0.7.7.1-fr}
Soient $A$ un anneau linéairement topologisé, $M$, $N$, deux $A$-modules
linéairement topologisés. Soient $\mathfrak J$, $V$, $W$ des voisinages
ouverts de $0$ dans $A$, $M$, $N$ respectivement, qui soient des
$A$-modules, et tels que $\mathfrak JM\subset V$, $\mathfrak JN\subset W$,
de sorte que $M/V$ et $N/W$ peuvent être considérés comme des
$(A/\mathfrak J)$-modules. Lorsque $\mathfrak J$, $V$, $W$ parcourent les
systèmes de voisinages ouverts vérifiant les conditions précédentes, il est
immédiat que les modules
$(M/V)\otimes_{A/\mathfrak J}(N/W)$ forment un système projectif de modules

\oldpage[0\textsubscript{I}]{76}

sur le système projectif d'anneaux $A/\mathfrak J$~; par passage à la limite
projective, on en déduit donc un module sur le séparé complété $\widehat A$
de $A$, que l'on appelle le \emph{produit tensoriel complété} de $M$ et de
$N$ et que l'on note $(M\otimes_A N)^\wedge$. Si l'on remarque que $M/V$
est canoniquement isomorphe à $\widehat M/\overline V$, où $\widehat M$ est
le séparé complété de $M$ et $\overline V$ l'adhérence dans $\widehat M$ de
l'image de $V$, on voit que le produit tensoriel complété
$(M\otimes_A N)^\wedge$ s'identifie canoniquement à
$(\widehat M\otimes_{\widehat A}\widehat N)^\wedge$, que l'on note aussi
$\widehat M\widehat\otimes_{\widehat A}\widehat N$.
\end{env}

\begin{env}[7.7.2]
\label{0.7.7.2-fr}
\begin{sloppypar}
Avec les notations précédentes, les produits tensoriels
$(M/V)\otimes_A(N/W)$ et
$(M/V)\otimes_{A/\mathfrak J}(N/W)$ s'identifient canoniquement~; ils
s'identifient donc aussi à
$(M\otimes_A N)/(\Im(V\otimes_A N)+\Im(M\otimes_A W))$.
On en conclut que $(M\otimes_A N)^\wedge$ est \emph{le séparé complété du
$A$-module $M\otimes_A N$, muni de la topologie pour laquelle les
sous-modules}
\[
  \Im(V\otimes_A N)+\Im(M\otimes_A W)
\]
\emph{forment un système fondamental de voisinages de $0$} ($V$ et $W$
parcourant l'ensemble des sous-modules ouverts de $M$ et $N$
respectivement)~; nous dirons pour abréger que cette topologie est le
\emph{produit tensoriel} des topologies données sur $M$ et $N$.
\end{sloppypar}
\end{env}

\begin{env}[7.7.3]
\label{0.7.7.3-fr}
Soient $M'$, $N'$ deux $A$-modules linéairement topologisés,
$u:M\to M'$, $v:N\to N'$ deux homomorphismes continus~; il est immédiat
que $u\otimes v$ est continu pour les topologies produits tensoriels sur
$M\otimes_A N$ et $M'\otimes_A N'$ respectivement~; par passage au séparé
complété, on en déduit un homomorphisme continu
$(M\otimes_A N)^\wedge\to(M'\otimes_A N')^\wedge$, que nous désignerons
par $u\widehat\otimes v$~; $(M\otimes_A N)^\wedge$ est donc un
\emph{bifoncteur} en $M$ et $N$ dans la catégorie des $A$-modules
linéairement topologisés.
\end{env}

\begin{env}[7.7.4]
\label{0.7.7.4-fr}
On définit de même le produit tensoriel complété d'un nombre fini quelconque
de $A$-modules linéairement topologisés~; il est immédiat que ce produit
possède les propriétés usuelles d'associativité et de commutativité.
\end{env}

\begin{env}[7.7.5]
\label{0.7.7.5-fr}
Si $B$, $C$ sont deux $A$-algèbres topologiques linéairement topologisées,
la topologie produit tensoriel sur $B\otimes_A C$ a pour système fondamental
de voisinages de $0$ les \emph{idéaux}
$\Im(\mathfrak K\otimes_A C)+\Im(B\otimes_A\mathfrak L)$ de l'algèbre
$B\otimes_A C$, $\mathfrak K$ (resp. $\mathfrak L$) parcourant l'ensemble
des idéaux ouverts de $B$ (resp. $C$). Par suite, $(B\otimes_A C)^\wedge$
est muni d'une structure de \emph{$\widehat A$-algèbre topologique}, limite
projective du système projectif de $(A/\mathfrak J)$-algèbres
$(B/\mathfrak K)\otimes_{A/\mathfrak J}(C/\mathfrak L)$ ($\mathfrak J$
idéal ouvert de $A$ tel que $\mathfrak J.B\subset\mathfrak K$,
$\mathfrak J.C\subset\mathfrak L$~; il en existe toujours). On dit que
cette algèbre est le \emph{produit tensoriel complété} des algèbres $B$ et
$C$.
\end{env}

\begin{env}[7.7.6]
\label{0.7.7.6-fr}
Les homomorphismes de $A$-algèbres $b\mapsto b\otimes 1$,
$c\mapsto 1\otimes c$ de $B$ et $C$ dans $B\otimes_A C$ sont continus
quand on munit cette dernière algèbre de la topologie produit tensoriel~;
par composition avec l'homomorphisme canonique de $B\otimes_A C$ dans son
séparé complété, ils donnent donc des \emph{homomorphismes canoniques}
$\rho:B\to(B\otimes_A C)^\wedge$, $\sigma:C\to(B\otimes_A C)^\wedge$.
L'algèbre $(B\otimes_A C)^\wedge$ et les homomorphismes $\rho$ et $\sigma$
possèdent en outre la \emph{propriété universelle suivante}~: pour toute
$A$-algèbre linéairement topologisée, séparée et complète $D$, et tout
couple de $A$-homomorphismes continus $u:B\to D$, $v:C\to D$, il existe un
$A$-homomorphisme continu et un seul
$w:(B\otimes_A C)^\wedge\to D$ et un seul tel que $u=w\circ\rho$ et
$v=w\circ\sigma$. En effet, il existe déjà un $A$-homomorphisme unique
$w_0:B\otimes_A C\to D$ tel que $u(b)=w_0(b\otimes 1)$ et
$v(c)=w_0(1\otimes c)$, et tout revient à prouver que $w_0$ est continu,
car il donnera alors un

\oldpage[0\textsubscript{I}]{77}

homomorphisme continu $w$ par passage au séparé complété. Or, si
$\mathfrak M$ est un idéal ouvert de $D$, il existe par hypothèse des idéaux
ouverts $\mathfrak K\subset B$, $\mathfrak L\subset C$ tels que
$u(\mathfrak K)\subset\mathfrak M$, $v(\mathfrak L)\subset\mathfrak M$~;
l'image par $w_0$ de
$\Im(\mathfrak K\otimes_A C)+\Im(B\otimes_A\mathfrak L)$ est encore
contenue dans $\mathfrak M$, d'où notre assertion.
\end{env}

\begin{proposition}[7.7.7]
\label{0.7.7.7-fr}
Si $B$ et $C$ sont deux $A$-algèbres préadmissibles,
$(B\otimes_A C)^\wedge$ est admissible, et si $\mathfrak K$ (resp.
$\mathfrak L$) est un idéal de définition de $B$ (resp. $C$), l'adhérence
dans $(B\otimes_A C)^\wedge$ de l'image canonique de
$\mathfrak H=\Im(\mathfrak K\otimes_A C)+\Im(B\otimes_A\mathfrak L)$ est
un idéal de définition.
\end{proposition}

Il suffit de montrer que $\mathfrak H^n$ tend vers $0$ pour la topologie
produit tensoriel, ce qui résulte immédiatement de l'inclusion
\[
  \mathfrak H^{2n}\subset
  \Im(\mathfrak K^n\otimes_A C)+\Im(B\otimes_A\mathfrak L^n).
\]

\begin{proposition}[7.7.8]
\label{0.7.7.8-fr}
Soient $A$ un anneau préadique, $\mathfrak J$ un idéal de définition de
$A$, $M$ un $A$-module \emph{de type fini}, muni de la topologie
$\mathfrak J$-préadique. Pour toute $A$-algèbre topologique adique et
noethérienne $B$, $B\otimes_A M$ s'identifie au produit tensoriel complété
$(B\otimes_A M)^\wedge$.
\end{proposition}

Si $\mathfrak K$ est un idéal de définition de $B$, il existe par hypothèse
un entier $m$ tel que $\mathfrak J^mB\subset\mathfrak K$, donc
$\Im(B\otimes_A\mathfrak J^{nm}M)=
\Im(\mathfrak J^{nm}B\otimes_A M)\subset
\Im(\mathfrak K^nB\otimes_A M)=\mathfrak K^n(B\otimes_A M)$~; on en
conclut que sur $B\otimes_A M$, le produit tensoriel des topologies de $B$
et de $M$ est la topologie $\mathfrak K$-préadique. Comme $B\otimes_A M$
est un $B$-module de type fini, la proposition résulte aussitôt de (7.3.6).

\subsection{Topologies sur les modules d'homomorphismes}
\label{subsection:0.7.8-fr}

\begin{env}[7.8.1]
\label{0.7.8.1-fr}
Soient $A$ un anneau $\mathfrak J$-adique noethérien, $M$ et $N$ deux
$A$-modules de type fini, munis de la topologie $\mathfrak J$-préadique~;
on sait (7.3.6) qu'ils sont séparés et complets~; en outre, tout
$A$-homomorphisme $M\to N$ est automatiquement continu, et le $A$-module
$\operatorname{Hom}_A(M,N)$ est de type fini. Pour tout entier $i\geq 0$, posons
$A_i=A/\mathfrak J^{i+1}$, $M_i=M/\mathfrak J^{i+1}M$,
$N_i=N/\mathfrak J^{i+1}N$~; pour $i\leq j$, tout homomorphisme
$u_j:M_j\to N_j$ applique $\mathfrak J^{i+1}M_j$ dans
$\mathfrak J^{i+1}N_j$, donc donne par passage aux quotients un
homomorphisme $u_i:M_i\to N_i$, ce qui définit un homomorphisme canonique
$\operatorname{Hom}_{A_j}(M_j,N_j)\to\operatorname{Hom}_{A_i}(M_i,N_i)$~; en outre, les
$\operatorname{Hom}_{A_i}(M_i,N_i)$ constituent un \emph{système projectif} pour ces
homomorphismes, et il résulte de (7.2.10) qu'il y a un isomorphisme
canonique
\[
  \varphi:\operatorname{Hom}_A(M,N)\longrightarrow
  \varprojlim_i\operatorname{Hom}_{A_i}(M_i,N_i).
\]
En outre~:
\end{env}

\begin{proposition}[7.8.2]
\label{0.7.8.2-fr}
Si $M$ et $N$ sont des modules de type fini sur un anneau $\mathfrak J$-adique
noethérien $A$, les sous-modules $\operatorname{Hom}_A(M,\mathfrak J^{i+1}N)$ forment un
système fondamental de voisinages de $0$ dans $\operatorname{Hom}_A(M,N)$ pour la
topologie $\mathfrak J$-adique, et l'homomorphisme canonique
\[
  \varphi:\operatorname{Hom}_A(M,N)\longrightarrow
  \varprojlim_i\operatorname{Hom}_{A_i}(M_i,N_i)
\]
est un isomorphisme topologique.
\end{proposition}

On peut en effet considérer $M$ comme quotient d'un $A$-module libre $L$ de
type fini, et par suite identifier $\operatorname{Hom}_A(M,N)$ à un sous-module de
$\operatorname{Hom}_A(L,N)$~; dans cette identification,
$\operatorname{Hom}_A(M,\mathfrak J^{i+1}N)$ est l'intersection de $\operatorname{Hom}_A(M,N)$ et de
$\operatorname{Hom}_A(L,\mathfrak J^{i+1}N)$~; comme la topologie induite sur
$\operatorname{Hom}_A(M,N)$ par la topologie $\mathfrak J$-adique de $\operatorname{Hom}_A(L,N)$ est
la topologie $\mathfrak J$-adique (7.3.2), on est ramené à démontrer la
première assertion

\oldpage[0\textsubscript{I}]{78}

pour $M=L=A^m$~; mais alors $\operatorname{Hom}_A(L,N)=N^m$,
$\operatorname{Hom}_A(L,\mathfrak J^{i+1}N)=(\mathfrak J^{i+1}N)^m=
\mathfrak J^{i+1}N^m$ et le résultat est évident. Pour établir la seconde
assertion, notons que l'image de $\operatorname{Hom}_A(M,\mathfrak J^{i+1}N)$ dans
$\operatorname{Hom}_{A_j}(M_j,N_j)$ est nulle pour $j\leq i$, donc $\varphi$ est
continu~; inversement, l'image réciproque dans $\operatorname{Hom}_A(M,N)$ du $0$ de
$\operatorname{Hom}_{A_i}(M_i,N_i)$ est $\operatorname{Hom}_A(M,\mathfrak J^{i+1}N)$, donc $\varphi$
est bicontinu.

Si on suppose seulement que $A$ est un anneau $\mathfrak J$-préadique
noethérien, $M$ et $N$ deux $A$-modules de type fini, séparés pour la
topologie $\mathfrak J$-préadique, la démonstration précédente montre que
la première assertion de (7.8.2) reste valable, et que $\varphi$ est un
isomorphisme topologique de $\operatorname{Hom}_A(M,N)$ sur un sous-module de
$\varprojlim_i\operatorname{Hom}_{A_i}(M_i,N_i)$.

\begin{proposition}[7.8.3]
\label{0.7.8.3-fr}
Sous les hypothèses de (7.8.2), l'ensemble des homomorphismes injectifs
(resp. surjectifs, bijectifs) de $M$ dans $N$ est une partie ouverte de
$\operatorname{Hom}_A(M,N)$.
\end{proposition}

En effet, en vertu de (7.3.5) et (7.1.14), pour que $u$ soit surjectif, il
faut et il suffit que l'homomorphisme correspondant
$u_0:M/\mathfrak JM\to N/\mathfrak JN$ le soit, et l'ensemble des
homomorphismes surjectifs de $M$ dans $N$ est donc image réciproque par
l'application continue $\operatorname{Hom}_A(M,N)\to\operatorname{Hom}_{A_0}(M_0,N_0)$ d'une partie
d'un espace discret. Montrons maintenant que l'ensemble des homomorphismes
injectifs est ouvert~; soit $v$ un tel homomorphisme et posons $M'=v(M)$~;
par le lemme d'Artin--Rees (7.3.2.1), il existe un entier $k\geq 0$ tel que
$M'\cap\mathfrak J^{m+k}N\subset\mathfrak J^mM'$ pour tout $m>0$~; nous
allons voir que pour $w\in\mathfrak J^{k+1}\operatorname{Hom}_A(M,N)$,
$u=v+w$ est injectif, ce qui achèvera la démonstration. En effet, soit
$x\in M$ tel que $u(x)=0$~; prouvons que pour tout $i\geq 0$ la relation
$x\in\mathfrak J^iM$ implique $x\in\mathfrak J^{i+1}M$~; il en résultera
bien
\[
  x\in\bigcap_{i\geq 0}\mathfrak J^iM=(0).
\]
En effet, on a alors $w(x)\in\mathfrak J^{i+k+1}N$, et par ailleurs
$w(x)=-v(x)\in M'$, donc
$v(x)\in M'\cap\mathfrak J^{i+k+1}N\subset\mathfrak J^{i+1}M'$, et comme
$v$ est un isomorphisme de $M$ sur $M'$, $x\in\mathfrak J^{i+1}M$~;
c.q.f.d.

\begin{flushright}
\emph{(A suivre.)}
\end{flushright}

% FIN SOURCE ega1/ega0-1-fr.tex


\clearpage
\setcounter{section}{0}
\setcounter{subsection}{0}
\setcounter{equation}{0}
\setcounter{footnote}{0}

%
% DEBUT SOURCE ega1/chapter1-frontmatter-fr.tex
\clearpage
\oldpage[I]{79}
\phantomsection
\label{chapter:I-fr}

\begin{center}
{\large CHAPITRE PREMIER\par}
\vspace{1.5em}
{\Large\bfseries LE LANGAGE DES SCHÉMAS\par}
\vspace{1.8em}
{\bfseries Sommaire\par}
\end{center}

\medskip
\begin{tabular}{@{}r@{\ }l@{}}
\S~1.  & Schémas affines.\\
\S~2.  & Préschémas et morphismes de préschémas.\\
\S~3.  & Produits de préschémas.\\
\S~4.  & Sous-préschémas et morphismes d'immersion.\\
\S~5.  & Préschémas réduits~; condition de séparation.\\
\S~6.  & Conditions de finitude.\\
\S~7.  & Applications rationnelles.\\
\S~8.  & Les schémas de Chevalley.\\
\S~9.  & Compléments sur les faisceaux quasi-cohérents.\\
\S~10. & Schémas formels.
\end{tabular}

\medskip
Les \S\S~1 à~8 ne font guère que développer un langage, celui qui sera utilisé
dans toute la suite. Notons cependant que, conformément à l'esprit général de
ce Traité, les \S\S~7 et~8 seront moins utilisés que les autres, et de façon
moins essentielle~; on n'a d'ailleurs parlé des schémas de Chevalley que pour
faire le lien avec le langage de Chevalley~\cite{I-1} et Nagata~\cite{I-9}. Le
\S~9 donne des définitions et résultats sur les faisceaux quasi-cohérents, dont
certains ne se bornent plus à une traduction en langage <<~géométrique~>> de
notions connues d'Algèbre commutative, mais sont déjà de nature globale~; ils
seront indispensables, dès les chapitres suivants, dans l'étude globale des
morphismes. Enfin, le \S~10 introduit une généralisation de la notion de
schéma, qui nous servira d'intermédiaire au chapitre~III pour formuler et
démontrer de façon commode les résultats fondamentaux de l'étude
cohomologique des morphismes propres~; par ailleurs, signalons que la notion
de schéma formel semble indispensable pour exprimer certains faits de la
<<~théorie des modules~>> (problèmes de classification des variétés
algébriques). Les résultats du \S~10 ne seront pas utilisés avant le \S~3 du
chapitre~III et il est recommandé d'en omettre la lecture jusque-là.

\clearpage

% FIN SOURCE ega1/chapter1-frontmatter-fr.tex

%
% DEBUT SOURCE ega1/ega1-1-fr.tex
\setcounter{section}{0}
\section{Schémas affines}
\label{section:I.1-fr}

\setcounter{subsection}{0}
\subsection{Le spectre premier d'un anneau}
\label{subsection:I.1.1-fr}

\oldpage[I]{80}
\begin{env}[1.1.1]
\label{I.1.1.1-fr}
\emph{Notations}~: Soient $A$ un anneau (commutatif), $M$ un $A$-module. Dans
ce chapitre et les suivants, nous utiliserons constamment les notations
suivantes~:

\smallskip
\noindent
$\operatorname{Spec}(A)={}$\emph{ensemble des idéaux premiers} de $A$, appelé
aussi \emph{spectre premier} de $A$~; pour un
$x\in X=\operatorname{Spec}(A)$, il sera souvent commode d'écrire
$\mathfrak{j}_x$ au lieu de $x$. Pour que $\operatorname{Spec}(A)$ soit
\emph{vide}, il faut et il suffit que l'anneau $A$ soit réduit à $0$.

\smallskip
\noindent
$A_x=A_{\mathfrak{j}_x}={}$\emph{anneau (local) des fractions $S^{-1}A$}, où
$S=A-\mathfrak{j}_x$.

\smallskip
\noindent
$\mathfrak{m}_x=\mathfrak{j}_xA_{\mathfrak{j}_x}={}$\emph{idéal maximal de
$A_x$}.

\smallskip
\noindent
$k(x)=A_x/\mathfrak{m}_x={}$\emph{corps résiduel de $A_x$}, isomorphe
canoniquement au corps des fractions de l'anneau intègre
$A/\mathfrak{j}_x$, auquel on l'identifie.

\smallskip
\noindent
$f(x)={}$\emph{classe de $f$ mod.~$\mathfrak{j}_x$}, dans
$A/\mathfrak{j}_x\subset k(x)$, pour $f\in A$ et $x\in X$. On dit encore que
$f(x)$ est la \emph{valeur} de $f$ au point $x\in\operatorname{Spec}(A)$~;
les relations $f(x)=0$ et $f\in\mathfrak{j}_x$ sont \emph{équivalentes}.

\smallskip
\noindent
$M_x=M\otimes_AA_x={}$\emph{module des fractions à dénominateurs dans
$A-\mathfrak{j}_x$}.

\smallskip
\noindent
$\mathfrak{r}(E)={}$\emph{racine de l'idéal de $A$ engendré par une partie
$E$ de $A$}.

\smallskip
\noindent
$V(E)={}$\emph{ensemble des $x\in X$ tels que $E\subset\mathfrak{j}_x$}
(ou encore ensemble des $x\in X$ tels que $f(x)=0$ pour tout $f\in E$), pour
$E\subset A$. On a donc
\begin{equation}
  \mathfrak{r}(E)=\bigcap_{x\in V(E)}\mathfrak{j}_x.
  \tag{1.1.1.1}\label{I.1.1.1.1-fr}
\end{equation}

\smallskip
\noindent
$V(f)=V(\{f\})$ pour $f\in A$.

\smallskip
\noindent
$D(f)=X-V(f)={}$\emph{ensemble des $x\in X$ où $f(x)\neq0$}.
\end{env}

\begin{proposition}[1.1.2]
\label{I.1.1.2-fr}
On a les propriétés suivantes~:
\begin{enumerate}
  \item[\rm(i)] $V(0)=X$, $V(1)=\varnothing$.
  \item[\rm(ii)] La relation $E\subset E'$ entraîne $V(E)\supset V(E')$.
  \item[\rm(iii)] Pour toute famille $(E_\lambda)$ de parties de $A$,
  \[
    V\left(\bigcup_\lambda E_\lambda\right)
      =V\left(\sum_\lambda E_\lambda\right)
      =\bigcap_\lambda V(E_\lambda).
  \]
  \item[\rm(iv)] $V(EE')=V(E)\cup V(E')$.
  \item[\rm(v)] $V(E)=V(\mathfrak{r}(E))$.
\end{enumerate}
\end{proposition}

Les propriétés (i), (ii), (iii) sont triviales, et (v) résulte de (ii) et de
la formule (\hyperref[I.1.1.1.1-fr]{1.1.1.1}). Il est évident que
$V(EE')\supset V(E)\cup V(E')$~; inversement, si $x\notin V(E)$ et
$x\notin V(E')$, il existe $f\in E$ et $f'\in E'$ tels que $f(x)\neq0$ et
$f'(x)\neq0$ dans $k(x)$, d'où $f(x)f'(x)\neq0$, autrement dit
$x\notin V(EE')$, ce qui prouve (iv).

La prop. (\hyperref[I.1.1.2-fr]{1.1.2}) montre entre autres que les ensembles
de la forme $V(E)$ (où $E$ parcourt l'ensemble des parties de $A$) sont les
\emph{ensembles fermés} d'une topologie sur $X$, que nous appellerons la
\emph{topologie spectrale}\footnote{L'introduction de cette topologie en
géométrie algébrique est due à Zariski. Aussi est-elle souvent appelée la
<<~topologie de Zariski~>> de $X$.}~; sauf mention expresse du contraire, on
supposera toujours $X=\operatorname{Spec}(A)$ muni de la topologie spectrale.

\oldpage[I]{81}
\begin{env}[1.1.3]
\label{I.1.1.3-fr}
Pour toute partie $Y$ de $X$, nous désignerons par $\mathfrak{j}(Y)$
l'ensemble des $f\in A$ tels que $f(y)=0$ pour tout $y\in Y$~; il revient au
même de dire que $\mathfrak{j}(Y)$ est l'intersection des idéaux premiers
$\mathfrak{j}_y$ pour $y\in Y$. Il est clair que la relation $Y\subset Y'$
entraîne $\mathfrak{j}(Y)\supset\mathfrak{j}(Y')$ et que l'on a
\begin{equation}
  \mathfrak{j}\left(\bigcup_\lambda Y_\lambda\right)
    =\bigcap_\lambda\mathfrak{j}(Y_\lambda)
  \tag{1.1.3.1}\label{I.1.1.3.1-fr}
\end{equation}
pour toute famille $(Y_\lambda)$ de parties de $X$. Enfin on a
\begin{equation}
  \mathfrak{j}(\{x\})=\mathfrak{j}_x.
  \tag{1.1.3.2}\label{I.1.1.3.2-fr}
\end{equation}
\end{env}

\begin{proposition}[1.1.4]
\label{I.1.1.4-fr}
\begin{enumerate}
  \item[\rm(i)] Pour toute partie $E$ de $A$, on a
  $\mathfrak{j}(V(E))=\mathfrak{r}(E)$.
  \item[\rm(ii)] Pour toute partie $Y$ de $X$,
  $V(\mathfrak{j}(Y))=\overline{Y}$, adhérence de $Y$ dans $X$.
\end{enumerate}
\end{proposition}

(i) est conséquence immédiate des définitions et de
(\hyperref[I.1.1.1.1-fr]{1.1.1.1})~; d'autre part,
$V(\mathfrak{j}(Y))$ est fermé et contient $Y$~; inversement, si
$Y\subset V(E)$, on a $f(y)=0$ pour tout $f\in E$ et tout $y\in Y$, donc
$E\subset\mathfrak{j}(Y)$, $V(E)\supset V(\mathfrak{j}(Y))$, ce qui prouve
(ii).

\begin{corollary}[1.1.5]
\label{I.1.1.5-fr}
Les parties fermées de $X=\operatorname{Spec}(A)$ et les idéaux de $A$ égaux
à leurs racines (autrement dit les intersections d'idéaux premiers) se
correspondent biunivoquement par les applications décroissantes
$Y\mapsto\mathfrak{j}(Y)$, $\mathfrak{a}\mapsto V(\mathfrak{a})$~; à la
réunion $Y_1\cup Y_2$ de deux parties fermées correspond
$\mathfrak{j}(Y_1)\cap\mathfrak{j}(Y_2)$, et à l'intersection d'une famille
quelconque $(Y_\lambda)$ de parties fermées correspond la racine de la somme
des $\mathfrak{j}(Y_\lambda)$.
\end{corollary}

\begin{corollary}[1.1.6]
\label{I.1.1.6-fr}
Si $A$ est un anneau noethérien, $X=\operatorname{Spec}(A)$ est un espace
noethérien.
\end{corollary}

On notera que la réciproque de ce corollaire est inexacte, comme le montre
l'exemple d'un anneau intègre non noethérien ayant un seul idéal premier
$\neq\{0\}$, par exemple un anneau de valuation non discrète de rang 1.

Comme exemple d'anneau $A$ dont le spectre n'est pas un espace noethérien, on
peut signaler l'anneau $\mathcal{C}(Y)$ des fonctions continues numériques sur
un espace compact infini $Y$~; on sait qu'en tant qu'ensemble, $Y$ s'identifie
à l'ensemble des idéaux maximaux de $A$, et il est facile de voir que la
topologie induite sur $Y$ par celle de $X=\operatorname{Spec}(A)$ est la
topologie initiale de $Y$. Comme $Y$ n'est pas un espace noethérien, il en est
donc de même de $X$.

\begin{corollary}[1.1.7]
\label{I.1.1.7-fr}
Pour tout $x\in X$, l'adhérence de $\{x\}$ est l'ensemble des $y\in X$ tels
que $\mathfrak{j}_x\subset\mathfrak{j}_y$. Pour que $\{x\}$ soit fermé, il
faut et il suffit que $\mathfrak{j}_x$ soit maximal.
\end{corollary}

\begin{corollary}[1.1.8]
\label{I.1.1.8-fr}
L'espace $X=\operatorname{Spec}(A)$ est un espace de Kolmogoroff.
\end{corollary}

En effet, si $x$, $y$ sont deux points distincts de $X$, on a, soit
$\mathfrak{j}_x\not\subset\mathfrak{j}_y$, soit
$\mathfrak{j}_y\not\subset\mathfrak{j}_x$, donc l'un des points $x$, $y$
n'appartient pas à l'adhérence de l'autre.

\begin{env}[1.1.9]
\label{I.1.1.9-fr}
D'après la prop. (\hyperref[I.1.1.2-fr]{1.1.2}, (iv)), pour deux éléments
$f$, $g$ de $A$, on a
\begin{equation}
  D(fg)=D(f)\cap D(g).
  \tag{1.1.9.1}\label{I.1.1.9.1-fr}
\end{equation}
Notons aussi que la relation $D(f)=D(g)$ signifie, d'après la prop.
(\hyperref[I.1.1.4-fr]{1.1.4}, (i)) et la prop.
(\hyperref[I.1.1.2-fr]{1.1.2}, (v)) que
$\mathfrak{r}(f)=\mathfrak{r}(g)$, ou encore que les idéaux premiers minimaux
contenant $(f)$ et $(g)$ sont les mêmes~; en particulier, il en est ainsi
lorsque $f=ug$, où $u$ est inversible.
\end{env}

\oldpage[I]{82}
\begin{proposition}[1.1.10]
\label{I.1.1.10-fr}
\begin{enumerate}
  \item[\rm(i)] Lorsque $f$ parcourt $A$, les ensembles $D(f)$ forment une
  base de la topologie de $X$.
  \item[\rm(ii)] Pour tout $f\in A$, $D(f)$ est quasi-compact. En particulier
  $X=D(1)$ est quasi-compact.
\end{enumerate}
\end{proposition}

(i) Soit $U$ un ensemble ouvert dans $X$~; par définition, on a
$U=X-V(E)$ où $E$ est une partie de $A$, et
$V(E)=\bigcap_{f\in E}V(f)$, d'où $U=\bigcup_{f\in E}D(f)$.

(ii) D'après (i), il suffit de prouver que si $(f_\lambda)_{\lambda\in L}$
est une famille d'éléments de $A$ telle que
$D(f)\subset\bigcup_{\lambda\in L}D(f_\lambda)$, il existe une partie finie
$J$ de $L$ telle que $D(f)\subset\bigcup_{\lambda\in J}D(f_\lambda)$. Soit
$\mathfrak{a}$ l'idéal de $A$ engendré par les $f_\lambda$~; on a par
hypothèse $V(f)\supset V(\mathfrak{a})$, donc
$\mathfrak{r}(f)\subset\mathfrak{r}(\mathfrak{a})$~; comme
$f\in\mathfrak{r}(f)$, il existe un entier $n\geq0$ tel que
$f^n\in\mathfrak{a}$. Mais alors $f^n$ appartient à l'idéal $\mathfrak{b}$
engendré par une sous-famille finie $(f_\lambda)_{\lambda\in J}$, et on a
$V(f)=V(f^n)\supset V(\mathfrak{b})=\bigcap_{\lambda\in J}V(f_\lambda)$,
c'est-à-dire $D(f)\subset\bigcup_{\lambda\in J}D(f_\lambda)$.

\begin{proposition}[1.1.11]
\label{I.1.1.11-fr}
Pour tout idéal $\mathfrak{a}$ de $A$,
$\operatorname{Spec}(A/\mathfrak{a})$ s'identifie canoniquement au sous-espace
fermé $V(\mathfrak{a})$ de $\operatorname{Spec}(A)$.
\end{proposition}

On sait en effet qu'il y a correspondance biunivoque canonique, respectant la
structure d'ordre de l'inclusion, entre idéaux (resp. idéaux premiers) de
$A/\mathfrak{a}$ et idéaux (resp. idéaux premiers) de $A$ contenant
$\mathfrak{a}$.

Rappelons que l'ensemble $\mathfrak{N}$ des éléments nilpotents de $A$ (ou
\emph{nilradical} de $A$) est un idéal égal à $\mathfrak{r}(0)$, intersection
de tous les idéaux premiers de $A$
(\hyperref[0.1.1.1-fr]{0, 1.1.1}).

\begin{corollary}[1.1.12]
\label{I.1.1.12-fr}
Les espaces topologiques $\operatorname{Spec}(A)$ et
$\operatorname{Spec}(A/\mathfrak{N})$ sont canoniquement isomorphes.
\end{corollary}

\begin{proposition}[1.1.13]
\label{I.1.1.13-fr}
Pour que $X=\operatorname{Spec}(A)$ soit irréductible
(\hyperref[0.2.1.1-fr]{0, 2.1.1}), il faut et il suffit que l'anneau
$A/\mathfrak{N}$ soit intègre (ou, ce qui revient au même, que l'idéal
$\mathfrak{N}$ soit premier).
\end{proposition}

En vertu du cor. (\hyperref[I.1.1.12-fr]{1.1.12}), on peut se borner au cas où
$\mathfrak{N}=0$. Si $X$ est réductible, il existe deux parties fermées
$Y_1$, $Y_2$ distinctes de $X$ et telles que $X=Y_1\cup Y_2$, d'où
$\mathfrak{j}(X)=\mathfrak{j}(Y_1)\cap\mathfrak{j}(Y_2)=0$, les idéaux
$\mathfrak{j}(Y_1)$ et $\mathfrak{j}(Y_2)$ étant distincts de $(0)$
(\hyperref[I.1.1.5-fr]{1.1.5})~; donc $A$ n'est pas intègre. Inversement, si
dans $A$ il y a deux éléments $f\neq0$, $g\neq0$ tels que $fg=0$, on a
$V(f)\neq X$, $V(g)\neq X$ (puisque l'intersection des idéaux premiers de
$A$ est $(0)$), et $X=V(fg)=V(f)\cup V(g)$.

\begin{corollary}[1.1.14]
\label{I.1.1.14-fr}
\begin{enumerate}
  \item[\rm(i)] Dans la correspondance biunivoque entre parties fermées de
  $X=\operatorname{Spec}(A)$ et idéaux de $A$ égaux à leurs racines, les
  parties fermées irréductibles de $X$ correspondent aux idéaux premiers de
  $A$. En particulier, les composantes irréductibles de $X$ correspondent aux
  idéaux premiers minimaux de $A$.
  \item[\rm(ii)] L'application $x\mapsto\overline{\{x\}}$ établit une
  correspondance biunivoque entre $X$ et l'ensemble des parties fermées
  irréductibles de $X$ (autrement dit toute partie fermée irréductible de $X$
  admet un point générique et un seul).
\end{enumerate}
\end{corollary}

(i) résulte aussitôt de (\hyperref[I.1.1.13-fr]{1.1.13}) et
(\hyperref[I.1.1.11-fr]{1.1.11})~; et pour démontrer (ii), on peut, en vertu de
(\hyperref[I.1.1.11-fr]{1.1.11}), se borner au cas où $X$ est irréductible~;
alors, d'après la prop. (\hyperref[I.1.1.13-fr]{1.1.13}), il existe dans $A$
un plus petit idéal premier $\mathfrak{N}$, qui correspond donc à un point
générique
\oldpage[I]{83}
de $X$~; en outre, $X$ n'admet qu'un seul point générique puisque c'est un
espace de Kolmogoroff ((\hyperref[I.1.1.8-fr]{1.1.8}) et
(\hyperref[0.2.1.3-fr]{0, 2.1.3})).

\begin{proposition}[1.1.15]
\label{I.1.1.15-fr}
Si $\mathfrak{J}$ est un idéal de $A$ contenu dans le radical
$\mathfrak{R}(A)$, le seul voisinage de $V(\mathfrak{J})$ dans
$X=\operatorname{Spec}(A)$ est l'espace $X$ tout entier.
\end{proposition}

En effet, tout idéal maximal de $A$ appartient par définition à
$V(\mathfrak{J})$. Comme tout idéal $\mathfrak{a}$ de $A$ est contenu dans un
idéal maximal, on a $V(\mathfrak{a})\cap V(\mathfrak{J})\neq\varnothing$,
d'où la proposition.

\subsection{Propriétés fonctorielles des spectres premiers d'anneaux}
\label{subsection:I.1.2-fr}

\begin{env}[1.2.1]
\label{I.1.2.1-fr}
Soient $A$, $A'$ deux anneaux,
\[
  \varphi:A'\longrightarrow A
\]
un homomorphisme d'anneaux. Pour tout idéal premier
$x=\mathfrak{j}_x\in\operatorname{Spec}(A)=X$, l'anneau
$A'/\varphi^{-1}(\mathfrak{j}_x)$ est canoniquement isomorphe à un
sous-anneau de $A/\mathfrak{j}_x$, donc est intègre, autrement dit
$\varphi^{-1}(\mathfrak{j}_x)$ est un idéal premier de $A'$~; nous le
noterons ${}^a\varphi(x)$, et nous avons ainsi défini une application
\[
  {}^a\varphi:X=\operatorname{Spec}(A)\longrightarrow
  X'=\operatorname{Spec}(A')
\]
(aussi notée $\operatorname{Spec}(\varphi)$) que nous appellerons
l'application associée à l'homomorphisme $\varphi$. Nous désignerons par
$\varphi^x$ l'homomorphisme injectif de
$A'/\varphi^{-1}(\mathfrak{j}_x)$ dans $A/\mathfrak{j}_x$, déduit de
$\varphi$ par passage aux quotients, ainsi que son prolongement canonique en
un monomorphisme de corps
\[
  \varphi^x:K({}^a\varphi(x))\longrightarrow K(x)~;
\]
pour tout $f'\in A'$, on a donc par définition
\begin{equation}
  \varphi^x\bigl(f'({}^a\varphi(x))\bigr)=\bigl(\varphi(f')\bigr)(x)
  \qquad(x\in X).
  \tag{1.2.1.1}\label{I.1.2.1.1-fr}
\end{equation}
\end{env}

\begin{proposition}[1.2.2]
\label{I.1.2.2-fr}
\begin{enumerate}
  \item[\rm(i)] Pour toute partie $E'$ de $A'$, on a
  \begin{equation}
    {}^a\varphi^{-1}(V(E'))=V(\varphi(E'))
    \tag{1.2.2.1}\label{I.1.2.2.1-fr}
  \end{equation}
  et en particulier, pour tout $f'\in A'$,
  \begin{equation}
    {}^a\varphi^{-1}(D(f'))=D(\varphi(f')).
    \tag{1.2.2.2}\label{I.1.2.2.2-fr}
  \end{equation}
  \item[\rm(ii)] Pour tout idéal $\mathfrak{a}$ de $A$, on a
  \begin{equation}
    \overline{{}^a\varphi(V(\mathfrak{a}))}
    =V(\varphi^{-1}(\mathfrak{a})).
    \tag{1.2.2.3}\label{I.1.2.2.3-fr}
  \end{equation}
\end{enumerate}
\end{proposition}

En effet, la relation ${}^a\varphi(x)\in V(E')$ est par définition
équivalente à $E'\subset\varphi^{-1}(\mathfrak{j}_x)$, donc à
$\varphi(E')\subset\mathfrak{j}_x$, et finalement à
$x\in V(\varphi(E'))$, d'où (i). Pour démontrer (ii), on peut supposer
$\mathfrak{a}$ égal à sa racine, puisque
$V(\mathfrak{r}(\mathfrak{a}))=V(\mathfrak{a})$
(\hyperref[I.1.1.2-fr]{1.1.2 (v)}) et
$\varphi^{-1}(\mathfrak{r}(\mathfrak{a}))
=\mathfrak{r}(\varphi^{-1}(\mathfrak{a}))$~; si on pose
$Y=V(\mathfrak{a})$, et
$\mathfrak{a}'=\mathfrak{j}({}^a\varphi(Y))$, on a
$\overline{{}^a\varphi(Y)}=V(\mathfrak{a}')$ (prop.
(\hyperref[I.1.1.4-fr]{1.1.4 (ii)}))~; la relation
$f'\in\mathfrak{a}'$ est par définition équivalente à $f'(x')=0$ pour tout
$x'\in{}^a\varphi(Y)$, donc, en vertu de la formule
(\hyperref[I.1.2.1.1-fr]{1.2.1.1}), elle est aussi équivalente à
$\varphi(f')(x)=0$ pour tout $x\in Y$, ou encore à
$\varphi(f')\in\mathfrak{j}(Y)=\mathfrak{a}$, puisque $\mathfrak{a}$ est
égal à sa racine~; d'où (ii).

\begin{corollary}[1.2.3]
\label{I.1.2.3-fr}
L'application ${}^a\varphi$ est continue.
\end{corollary}

Remarquons que si $A''$ est un troisième anneau, $\varphi'$ un homomorphisme
$A''\to A'$, on a
${}^a(\varphi'\circ\varphi)={}^a\varphi\circ{}^a\varphi'$~; ce résultat et le
cor. (\hyperref[I.1.2.3-fr]{1.2.3}) signifient que
$\operatorname{Spec}(A)$ est un foncteur contravariant en $A$, de la
catégorie des anneaux dans celle des espaces topologiques.

\begin{corollary}[1.2.4]
\label{I.1.2.4-fr}
\oldpage[I]{84}
Supposons que $\varphi$ soit tel que tout $f\in A$ s'écrive
$f=h\varphi(f')$, où $h$ est inversible dans $A$ (ce qui est en particulier
le cas lorsque $\varphi$ est surjectif). Alors ${}^a\varphi$ est un
homéomorphisme de $X$ sur ${}^a\varphi(X)$.
\end{corollary}

Montrons que pour toute partie $E\subset A$, il existe une partie $E'$ de
$A'$ telle que $V(E)=V(\varphi(E'))$~; en vertu de l'axiome $(T_0)$
(\hyperref[I.1.1.8-fr]{1.1.8}) et de la formule
(\hyperref[I.1.2.2.1-fr]{1.2.2.1}), cela entraînera d'abord que
${}^a\varphi$ est injective, puis, toujours en vertu de
(\hyperref[I.1.2.2.1-fr]{1.2.2.1}), que ${}^a\varphi$ est un
homéomorphisme. Or, il suffit pour chaque $f\in E$ de prendre un $f'\in A'$ tel
que $h\varphi(f')=f$ avec $h$ inversible dans $A$~; l'ensemble $E'$ de ces
éléments $f'$ répond à la question.

\begin{env}[1.2.5]
\label{I.1.2.5-fr}
En particulier, lorsque $\varphi$ est l'homomorphisme canonique de $A$ sur un
anneau quotient $A/\mathfrak{a}$, on retrouve
(\hyperref[I.1.1.12-fr]{1.1.12}), et ${}^a\varphi$ est l'injection canonique
de $V(\mathfrak{a})$, identifié à $\operatorname{Spec}(A/\mathfrak{a})$,
dans $X=\operatorname{Spec}(A)$.
\end{env}

Un autre cas particulier de (\hyperref[I.1.2.4-fr]{1.2.4}) donne~:

\begin{corollary}[1.2.6]
\label{I.1.2.6-fr}
Si $S$ est une partie multiplicative de $A$, le spectre
$\operatorname{Spec}(S^{-1}A)$ s'identifie canoniquement (avec sa topologie)
au sous-espace de $X=\operatorname{Spec}(A)$ formé des $x$ tels que
$\mathfrak{j}_x\cap S=\varnothing$.
\end{corollary}

On sait en effet (\hyperref[0.1.2.6-fr]{0, 1.2.6}) que les idéaux premiers de
$S^{-1}A$ sont les idéaux $S^{-1}\mathfrak{j}_x$ tels que
$\mathfrak{j}_x\cap S=\varnothing$, et que l'on a
$\mathfrak{j}_x=({}^S i_A)^{-1}(S^{-1}\mathfrak{j}_x)$. Il suffit donc
d'appliquer à ${}^S i_A$ le cor. (\hyperref[I.1.2.4-fr]{1.2.4}).

\begin{corollary}[1.2.7]
\label{I.1.2.7-fr}
Pour que ${}^a\varphi(X)$ soit partout dense dans $X'$, il faut et il suffit
que tout élément du noyau $\operatorname{Ker}\varphi$ soit nilpotent.
\end{corollary}

En effet, appliquant la formule (\hyperref[I.1.2.2.3-fr]{1.2.2.3}) à l'idéal
$\mathfrak{a}=(0)$, il vient
$\overline{{}^a\varphi(X)}=V(\operatorname{Ker}\varphi)$, et pour que
$V(\operatorname{Ker}\varphi)=X$, il faut et il suffit que
$\operatorname{Ker}\varphi$ soit contenu dans tous les idéaux premiers de
$A$, c'est-à-dire dans le nilradical $\mathfrak{N}$ de $A$.

\subsection{Faisceau associé à un module}
\label{subsection:I.1.3-fr}
\label{I.1.3-fr}

\begin{env}[1.3.1]
\label{I.1.3.1-fr}
Soient $A$ un anneau commutatif, $M$ un $A$-module, $f$ un élément de $A$,
$S_f$ l'ensemble multiplicatif des $f^n$, où $n\geq0$. Rappelons que nous
posons $A_f=S_f^{-1}A$, $M_f=S_f^{-1}M$. Si $S'_f$ est la partie
multiplicative saturée de $A$ formée des $g\in A$ qui divisent un élément de
$S_f$, on sait que $A_f$ et $M_f$ s'identifient canoniquement à
${S'_f}^{-1}A$ et ${S'_f}^{-1}M$
(\hyperref[0.1.4.3-fr]{0, 1.4.3}).
\end{env}

\begin{lemma}[1.3.2]
\label{I.1.3.2-fr}
Les conditions suivantes sont équivalentes~:
\begin{enumerate}
  \item[\rm(a)] $g\in S'_f$~;
  \item[\rm(b)] $S'_g\subset S'_f$~;
  \item[\rm(c)] $f\in\mathfrak{r}(g)$~;
  \item[\rm(d)] $\mathfrak{r}(f)\subset\mathfrak{r}(g)$~;
  \item[\rm(e)] $V(g)\subset V(f)$~;
  \item[\rm(f)] $D(f)\subset D(g)$.
\end{enumerate}
\end{lemma}

Cela résulte immédiatement des définitions et de
(\hyperref[I.1.1.5-fr]{1.1.5}).

\begin{env}[1.3.3]
\label{I.1.3.3-fr}
Si $D(f)=D(g)$, le lemme (\hyperref[I.1.3.2-fr]{1.3.2, b)}) montre que
$M_f=M_g$. Plus généralement, si $D(f)\supset D(g)$, donc
$S'_f\subset S'_g$, on sait (\hyperref[0.1.4.1-fr]{0, 1.4.1}) qu'il existe un
homomorphisme canonique fonctoriel
\[
  \rho_{g,f}:M_f\longrightarrow M_g,
\]
et si $D(f)\supset D(g)\supset D(h)$, on a
(\hyperref[0.1.4.4-fr]{0, 1.4.4})
\begin{equation}
  \rho_{h,g}\circ\rho_{g,f}=\rho_{h,f}.
  \tag{1.3.3.1}\label{I.1.3.3.1-fr}
\end{equation}
\end{env}

\oldpage[I]{85}
Lorsque $f$ parcourt $A-\mathfrak{j}_x$ (pour un $x$ donné dans
$X=\operatorname{Spec}(A)$), les ensembles $S'_f$ constituent un ensemble
filtrant croissant de parties de $A-\mathfrak{j}_x$, car pour deux éléments
$f$, $g$ de $A-\mathfrak{j}_x$, $S'_f$ et $S'_g$ sont contenus dans
$S'_{fg}$~; comme la réunion des $S'_f$ pour $f\in A-\mathfrak{j}_x$ est
$A-\mathfrak{j}_x$, on en conclut (\hyperref[0.1.4.5-fr]{0, 1.4.5}) que le
$A_x$-module $M_x$ s'identifie canoniquement à la \emph{limite inductive}
$\varinjlim M_f$, relativement à la famille d'homomorphismes
$(\rho_{g,f})$. Nous désignerons par
\[
  \rho_x^f:M_f\longrightarrow M_x
\]
l'homomorphisme canonique pour $f\in A-\mathfrak{j}_x$ (ou, ce qui revient au
même, $x\in D(f)$).

\begin{definition}[1.3.4]
\label{I.1.3.4-fr}
On appelle \emph{faisceau structural du spectre premier}
$X=\operatorname{Spec}(A)$ (resp. \emph{faisceau associé au $A$-module}
$M$) et on note $\widetilde{A}$ ou $\mathscr{O}_X$ (resp. $\widetilde{M}$) le
faisceau d'anneaux (resp. le $\widetilde{A}$-Module) associé au préfaisceau
$D(f)\mapsto A_f$ (resp. $D(f)\mapsto M_f$) sur la base $\mathfrak{B}$ de
$X$ formée des $D(f)$ où $f\in A$
((\hyperref[I.1.1.10-fr]{1.1.10}),
(\hyperref[0.3.2.1-fr]{0, 3.2.1} et
\hyperref[0.3.5.6-fr]{3.5.6})).
\end{definition}

On a vu (\hyperref[0.3.2.4-fr]{0, 3.2.4}) que la fibre
$\widetilde{A}_x$ (resp. $\widetilde{M}_x$) \emph{s'identifie à l'anneau}
$A_x$ (resp. \emph{au $A_x$-module} $M_x$)~; nous désignerons par
\[
  \theta_f:A_f\longrightarrow\Gamma(D(f),\widetilde{A})
\]
\[
  \text{(resp. }\theta_f:M_f\longrightarrow
  \Gamma(D(f),\widetilde{M})\text{),}
\]
l'application canonique, de sorte que pour tout $x\in D(f)$ et tout
$\xi\in M_f$, on a
\begin{equation}
  (\theta_f(\xi))_x=\rho_x^f(\xi).
  \tag{1.3.4.1}\label{I.1.3.4.1-fr}
\end{equation}

\begin{proposition}[1.3.5]
\label{I.1.3.5-fr}
$\widetilde{M}$ est un foncteur covariant exact en $M$, de la catégorie des
$A$-modules dans la catégorie des $\widetilde{A}$-Modules.
\end{proposition}

En effet, soient $M$, $N$ deux $A$-modules, $u$ un homomorphisme
$M\to N$~; pour tout $f\in A$, il correspond canoniquement à $u$ un
homomorphisme $u_f$ du $A_f$-module $M_f$ dans le $A_f$-module $N_f$, et le
diagramme (pour $D(g)\subset D(f)$)
\[
  \xymatrix{
    M_f\ar[r]^{u_f}\ar[d]_{\rho_{g,f}} & N_f\ar[d]^{\rho_{g,f}}\\
    M_g\ar[r]_{u_g} & N_g
  }
\]
est commutatif (\hyperref[0.1.4.1-fr]{0, 1.4.1})~; ces homomorphismes
définissent donc un homomorphisme de $\widetilde{A}$-Modules
$\widetilde{u}:\widetilde{M}\to\widetilde{N}$
(\hyperref[0.3.2.3-fr]{0, 3.2.3}). En outre, pour tout $x\in X$,
$\widetilde{u}_x$ est la limite inductive des $u_f$ pour $x\in D(f)$
($f\in A$), et par suite (\hyperref[0.1.4.5-fr]{0, 1.4.5}) si on identifie
canoniquement $\widetilde{M}_x$ et $\widetilde{N}_x$ à $M_x$ et $N_x$
respectivement, $\widetilde{u}_x$ s'identifie à l'homomorphisme $u_x$ déduit
canoniquement de $u$. Si $P$ est un troisième $A$-module, $v$ un
homomorphisme $N\to P$ et $w=v\circ u$, il est immédiat que
$w_x=v_x\circ u_x$, donc $\widetilde{w}=\widetilde{v}\circ\widetilde{u}$. On
a donc bien défini un \emph{foncteur covariant} $\widetilde{M}$ en $M$, de la
catégorie des $A$-modules dans celle des $\widetilde{A}$-Modules. \emph{Ce
foncteur est exact}, car pour tout $x\in X$, $M_x$ est un foncteur exact en
$M$ (\hyperref[0.1.3.2-fr]{0, 1.3.2})~; en outre, on a
$\operatorname{Supp}(M)=\operatorname{Supp}(\widetilde{M})$ en vertu des
définitions des deux membres (\hyperref[0.1.7.1-fr]{0, 1.7.1} et
\hyperref[0.3.1.6-fr]{3.1.6}).

\oldpage[I]{86}
\begin{proposition}[1.3.6]
\label{I.1.3.6-fr}
Pour tout $f\in A$, l'ensemble ouvert $D(f)\subset X$ s'identifie
canoniquement au spectre premier $\operatorname{Spec}(A_f)$, et le faisceau
$\widetilde{M_f}$ associé au $A_f$-module $M_f$ s'identifie canoniquement à
la restriction $\widetilde{M}|D(f)$.
\end{proposition}

La première assertion est un cas particulier de
(\hyperref[I.1.2.6-fr]{1.2.6}). En outre, si $g\in A$ est tel que
$D(g)\subset D(f)$, $M_g$ s'identifie canoniquement au module des fractions
de $M_f$ dont les dénominateurs sont les puissances de l'image canonique de
$g$ dans $A_f$ (\hyperref[0.1.4.6-fr]{0, 1.4.6}). L'identification canonique
de $\widetilde{M_f}$ à $\widetilde{M}|D(f)$ résulte alors des définitions.

\begin{theorem}[1.3.7]
\label{I.1.3.7-fr}
Pour tout $A$-module $M$ et tout $f\in A$, l'homomorphisme
\[
  \theta_f:M_f\longrightarrow\Gamma(D(f),\widetilde{M})
\]
est bijectif (autrement dit, le préfaisceau $D(f)\mapsto M_f$ est un
\emph{faisceau}). En particulier, $M$ s'identifie par $\theta_1$ à
$\Gamma(X,\widetilde{M})$.
\end{theorem}

On notera que, si $M=A$, $\theta_f$ est un homomorphisme de structure
d'anneau~; le th. (\hyperref[I.1.3.7-fr]{1.3.7}) entraînera donc que, si on
identifie les anneaux $A_f$ et $\Gamma(D(f),\widetilde{A})$ au moyen de
$\theta_f$, l'homomorphisme
$\theta_f:M_f\to\Gamma(D(f),\widetilde{M})$ sera un isomorphisme de
\emph{modules}.

Montrons d'abord que $\theta_f$ est \emph{injectif}. En effet, si
$\xi\in M_f$ est tel que $\theta_f(\xi)=0$, cela signifie que pour tout idéal
premier $\mathfrak{p}$ de $A_f$, il existe $h\notin\mathfrak{p}$ tel que
$h\xi=0$~; comme l'annulateur de $\xi$ n'est contenu dans aucun idéal premier
de $A_f$, c'est $A_f$ tout entier, donc $\xi=0$.

Reste à montrer que $\theta_f$ est \emph{surjectif}~; on peut se ramener au
cas où $f=1$, le cas général s'en déduisant en «~localisant~» à l'aide de
(\hyperref[I.1.3.6-fr]{1.3.6}). Soit donc $s$ une section de $\widetilde{M}$
au-dessus de $X$~; en vertu de (\hyperref[I.1.3.4-fr]{1.3.4}) et de
(\hyperref[I.1.1.10-fr]{1.1.10, (ii)}), il existe un recouvrement
\emph{fini} $(D(f_i))_{i\in I}$ de $X$ ($f_i\in A$) tel que, pour tout
$i\in I$, la restriction $s_i=s|D(f_i)$ soit de la forme
$\theta_{f_i}(\xi_i)$, où $\xi_i\in M_{f_i}$. Si $i$, $j$ sont deux indices
de $I$, en écrivant que les restrictions de $s_i$ et de $s_j$ à
$D(f_i)\cap D(f_j)=D(f_i f_j)$ sont égales, il vient par définition de
$\widetilde{M}$
\begin{equation}
  \rho_{f_i f_j,f_i}(\xi_i)=\rho_{f_i f_j,f_j}(\xi_j).
  \tag{1.3.7.1}\label{I.1.3.7.1-fr}
\end{equation}
Par définition, on peut écrire, pour chaque $i\in I$,
$\xi_i=z_i/f_i^{n_i}$, où $z_i\in M$, et comme $I$ est fini, en multipliant
chaque $z_i$ par une puissance de $f_i$, on peut supposer tous les $n_i$
égaux à un même $n$. Alors par définition, (\hyperref[I.1.3.7.1-fr]{1.3.7.1})
signifie qu'il existe un entier $m_{ij}\geq0$ tel que
$(f_i f_j)^{m_{ij}}(f_j^n z_i-f_i^n z_j)=0$, et on peut encore supposer tous
les $m_{ij}$ égaux à un même entier $m$~; remplaçant alors $z_i$ par
$f_i^m z_i$, on est ramené au cas où $m=0$, autrement dit au cas où l'on a
\begin{equation}
  f_j^n z_i=f_i^n z_j
  \tag{1.3.7.2}\label{I.1.3.7.2-fr}
\end{equation}
quels que soient $i$, $j$. Or, on a $D(f_i^n)=D(f_i)$, et comme les $D(f_i)$
forment un recouvrement de $X$, l'idéal engendré par les $f_i^n$ est $A$~;
en d'autres termes, il existe des éléments $g_i\in A$ tels que
$\sum_i g_i f_i^n=1$. Considérons alors l'élément
$z=\sum_i g_i z_i$ de $M$~; d'après (\hyperref[I.1.3.7.2-fr]{1.3.7.2}), on a
$f_i^n z=\sum_j g_j f_i^n z_j=(\sum_j g_j f_j^n)z_i=z_i$, d'où par
définition $\xi_i=z/1$ dans $M_{f_i}$. On en conclut
\oldpage[I]{87}
que $s_i$ est la restriction à $D(f_i)$ de $\theta_1(z)$, ce qui prouve que
$s=\theta_1(z)$ et achève la démonstration.

\begin{corollary}[1.3.8]
\label{I.1.3.8-fr}
Soient $M$, $N$ deux $A$-modules~; l'homomorphisme canonique
$u\mapsto\widetilde{u}$ de $\operatorname{Hom}_A(M,N)$ dans
$\operatorname{Hom}_{\widetilde{A}}(\widetilde{M},\widetilde{N})$ est
bijectif. En particulier, les relations $M=0$ et $\widetilde{M}=0$ sont
équivalentes.
\end{corollary}

Considérons en effet l'homomorphisme canonique $v\mapsto\Gamma(v)$ de
$\operatorname{Hom}_{\widetilde{A}}(\widetilde{M},\widetilde{N})$ dans
$\operatorname{Hom}_{\Gamma(\widetilde{A})}
(\Gamma(\widetilde{M}),\Gamma(\widetilde{N}))$~; ce dernier module s'identifie
canoniquement à $\operatorname{Hom}_A(M,N)$ en vertu du th.
(\hyperref[I.1.3.7-fr]{1.3.7}). Il reste à vérifier que
$u\mapsto\widetilde{u}$ et $v\mapsto\Gamma(v)$ sont réciproques l'un de
l'autre~; or, il est évident que $\Gamma(\widetilde{u})=u$ par définition de
$\widetilde{u}$~; et d'autre part, si on pose $u=\Gamma(v)$ pour
$v\in\operatorname{Hom}_{\widetilde{A}}(\widetilde{M},\widetilde{N})$,
l'application
$w:\Gamma(D(f),\widetilde{M})\to\Gamma(D(f),\widetilde{N})$ déduite
canoniquement de $v$ est telle que le diagramme
\[
  \xymatrix{
    M\ar[r]^u\ar[d]_{\rho_{f,1}} & N\ar[d]^{\rho_{f,1}}\\
    M_f\ar[r]_w & N_f
  }
\]
soit commutatif~; on a donc nécessairement $w=u_f$ pour tout $f\in A$
(\hyperref[0.1.2.4-fr]{0, 1.2.4}), ce qui montre que
$\widetilde{\Gamma(v)}=v$.

\begin{corollary}[1.3.9]
\label{I.1.3.9-fr}
\begin{enumerate}
  \item[\rm(i)] Soit $u$ un homomorphisme d'un $A$-module $M$ dans un
  $A$-module $N$~; alors les faisceaux associés à $\operatorname{Ker}u$,
  $\operatorname{Im}u$, $\operatorname{Coker}u$, sont respectivement
  $\operatorname{Ker}\widetilde{u}$, $\operatorname{Im}\widetilde{u}$,
  $\operatorname{Coker}\widetilde{u}$. En particulier pour que
  $\widetilde{u}$ soit injectif (resp. surjectif, bijectif), il faut et il
  suffit que $u$ le soit.
  \item[\rm(ii)] Si $M$ est une limite inductive (resp. somme directe) d'une
  famille de $A$-modules $(M_\lambda)$, $\widetilde{M}$ est limite inductive
  (resp. somme directe) de la famille $(\widetilde{M_\lambda})$, à un
  isomorphisme canonique près.
\end{enumerate}
\end{corollary}

(i) Il suffit d'appliquer le fait que $\widetilde{M}$ est un foncteur exact en
$M$ (\hyperref[I.1.3.5-fr]{1.3.5}) aux deux suites exactes de $A$-modules
\[
  0\to\operatorname{Ker}u\to M\to\operatorname{Im}u\to0
\]
\[
  0\to\operatorname{Im}u\to N\to\operatorname{Coker}u\to0
\]
La seconde assertion résulte alors du th.
(\hyperref[I.1.3.7-fr]{1.3.7}).

(ii) Soit $(M_\lambda,g_{\mu\lambda})$ un système inductif de $A$-modules, de
limite inductive $M$, et soit $g_\lambda$ l'homomorphisme canonique
$M_\lambda\to M$. Comme on a
$\widetilde{g_{\nu\mu}}\circ\widetilde{g_{\mu\lambda}}
=\widetilde{g_{\nu\lambda}}$ et
$\widetilde{g_\lambda}=\widetilde{g_\mu}\circ
\widetilde{g_{\mu\lambda}}$ pour $\lambda\leq\mu\leq\nu$,
$(\widetilde{M_\lambda},\widetilde{g_{\mu\lambda}})$ est un système
inductif de faisceaux sur $X$, et si on désigne par $h_\lambda$
l'homomorphisme canonique
$\widetilde{M_\lambda}\to\varinjlim\widetilde{M_\lambda}$, il y a un
homomorphisme unique
$v:\varinjlim\widetilde{M_\lambda}\to\widetilde{M}$ tel que
$v\circ h_\lambda=\widetilde{g_\lambda}$. Pour voir que $v$ est bijectif, il
suffit de vérifier que, pour tout $x\in X$, $v_x$ est une bijection de
$(\varinjlim\widetilde{M_\lambda})_x$ sur $\widetilde{M}_x$~; mais
$\widetilde{M}_x=M_x$ et
\[
  (\varinjlim\widetilde{M_\lambda})_x
  =\varinjlim(\widetilde{M_\lambda})_x
  =\varinjlim(M_\lambda)_x=M_x
  \quad(\hyperref[0.1.3.3-fr]{0, 1.3.3}).
\]
D'autre part, il résulte des définitions que $(\widetilde{g_\lambda})_x$ et
$(h_\lambda)_x$ sont tous deux égaux à l'application canonique de
$(M_\lambda)_x$ dans $M_x$~; comme
$(\widetilde{g_\lambda})_x=v_x\circ(h_\lambda)_x$, $v_x$ est l'identité.
\oldpage[I]{88}
Enfin, si $M$ est somme directe de deux $A$-modules $N$, $P$, il est immédiat
que $\widetilde{M}=\widetilde{N}\oplus\widetilde{P}$~; toute somme directe
étant limite inductive de sommes directes finies, les assertions de (ii) sont
démontrées.

On notera que (\hyperref[I.1.3.8-fr]{1.3.8}) prouve que les faisceaux
isomorphes aux faisceaux associés aux $A$-modules forment une
\emph{catégorie abélienne} (T, I, 1.4).

On notera aussi qu'il résulte de (\hyperref[I.1.3.9-fr]{1.3.9}) que si $M$ est
un $A$-module \emph{de type fini}, c'est-à-dire s'il existe un homomorphisme
surjectif $A^n\to M$, alors il existe un homomorphisme surjectif
$\widetilde{A^n}\to\widetilde{M}$, autrement dit le $\widetilde{A}$-Module
$\widetilde{M}$ est \emph{engendré par une famille finie de sections au-dessus
de $X$} (\hyperref[0.5.1.1-fr]{0, 5.1.1}), et réciproquement.

\begin{env}[1.3.10]
\label{I.1.3.10-fr}
Si $N$ est un sous-module d'un $A$-module $M$, l'injection canonique
$j:N\to M$ donne par (\hyperref[I.1.3.9-fr]{1.3.9}) un homomorphisme injectif
$\widetilde{N}\to\widetilde{M}$, qui permet d'identifier canoniquement
$\widetilde{N}$ à un \emph{sous-$\widetilde{A}$-Module} de $\widetilde{M}$~;
nous supposerons toujours faite cette identification. Si $N$ et $P$ sont deux
sous-modules de $M$, on a alors
\begin{equation}
  \widetilde{(N+P)}=\widetilde{N}+\widetilde{P}
  \tag{1.3.10.1}\label{I.1.3.10.1-fr}
\end{equation}
\begin{equation}
  \widetilde{(N\cap P)}=\widetilde{N}\cap\widetilde{P}
  \tag{1.3.10.2}\label{I.1.3.10.2-fr}
\end{equation}
car $N+P$ et $N\cap P$ sont respectivement l'image de l'homomorphisme
canonique $N\oplus P\to M$, et le noyau de l'homomorphisme canonique
$M\to(M/N)\oplus(M/P)$ et il suffit d'appliquer
(\hyperref[I.1.3.9-fr]{1.3.9}).

On conclut de (\hyperref[I.1.3.10.1-fr]{1.3.10.1}) et
(\hyperref[I.1.3.10.2-fr]{1.3.10.2}) que si
$\widetilde{N}=\widetilde{P}$, on a $N=P$.
\end{env}

\begin{corollary}[1.3.11]
\label{I.1.3.11-fr}
Dans la catégorie des faisceaux isomorphes aux faisceaux associés aux
$A$-modules, le foncteur $\Gamma$ est exact.
\end{corollary}

En effet, soit
$\widetilde{M}\xrightarrow{\widetilde{u}}\widetilde{N}
\xrightarrow{\widetilde{v}}\widetilde{P}$ une suite exacte correspondant à
deux homomorphismes $u:M\to N$, $v:N\to P$ de $A$-modules. Si
$Q=\operatorname{Im}u$ et $R=\operatorname{Ker}v$, on a
$\widetilde{Q}=\operatorname{Im}\widetilde{u}
=\operatorname{Ker}\widetilde{v}=\widetilde{R}$
(cor. (\hyperref[I.1.3.9-fr]{1.3.9})), donc $Q=R$.

\begin{corollary}[1.3.12]
\label{I.1.3.12-fr}
Soient $M$, $N$ deux $A$-modules.
\begin{enumerate}
  \item[\rm(i)] Le faisceau associé à $M\otimes_A N$ s'identifie
  canoniquement à $\widetilde{M}\otimes_{\widetilde{A}}\widetilde{N}$.
  \item[\rm(ii)] Si de plus $M$ admet une présentation finie, le faisceau
  associé à $\operatorname{Hom}_A(M,N)$ s'identifie canoniquement à
  $\mathscr{H}\!om_{\widetilde{A}}(\widetilde{M},\widetilde{N})$.
\end{enumerate}
\end{corollary}

(i) Le faisceau
$\mathscr{F}=\widetilde{M}\otimes_{\widetilde{A}}\widetilde{N}$ est associé
au préfaisceau
\[
  U\mapsto\mathscr{F}(U)=\Gamma(U,\widetilde{M})
  \otimes_{\Gamma(U,\widetilde{A})}\Gamma(U,\widetilde{N})
\]
$U$ parcourant la base (\hyperref[I.1.1.10-fr]{1.1.10, (i)}) de $X$ formée
des $D(f)$, où $f\in A$. Or, $\mathscr{F}(D(f))$ s'identifie canoniquement à
$M_f\otimes_{A_f}N_f$ en vertu de (\hyperref[I.1.3.7-fr]{1.3.7}) et
(\hyperref[I.1.3.6-fr]{1.3.6}). On sait par ailleurs que le $A_f$-module
$M_f\otimes_{A_f}N_f$ est canoniquement isomorphe à $(M\otimes_A N)_f$
(\hyperref[0.1.3.4-fr]{0, 1.3.4}), qui lui-même est canoniquement isomorphe à
$\Gamma(D(f),\widetilde{(M\otimes_A N)})$
(\hyperref[I.1.3.7-fr]{1.3.7} et \hyperref[I.1.3.6-fr]{1.3.6}). En outre, on
vérifie aussitôt que les isomorphismes canoniques
\[
  \mathscr{F}(D(f))\simeq\Gamma(D(f),\widetilde{(M\otimes_A N)})
\]
\oldpage[I]{89}
ainsi obtenus vérifient les conditions de compatibilité avec les opérateurs de
restriction (\hyperref[0.1.4.2-fr]{0, 1.4.2}), donc définissent un isomorphisme
canonique fonctoriel
\[
  \widetilde{M}\otimes_{\widetilde{A}}\widetilde{N}
  \simeq\widetilde{(M\otimes_A N)}
\]

(ii) Le faisceau
$\mathscr{G}=\mathscr{H}\!om_{\widetilde{A}}(\widetilde{M},\widetilde{N})$
est associé au préfaisceau
\[
  U\longmapsto\mathscr{G}(U)
  =\operatorname{Hom}_{\widetilde{A}|U}
  (\widetilde{M}|U,\widetilde{N}|U)
\]
$U$ parcourant la base de $X$ formée des $D(f)$. Or,
$\mathscr{G}(D(f))$ s'identifie canoniquement à
$\operatorname{Hom}_{A_f}(M_f,N_f)$
(\hyperref[I.1.3.6-fr]{1.3.6} et \hyperref[I.1.3.8-fr]{1.3.8}), qui
lui-même, en vertu de l'hypothèse sur $M$, s'identifie canoniquement à
$(\operatorname{Hom}_A(M,N))_f$
(\hyperref[0.1.3.5-fr]{0, 1.3.5}). Finalement,
$(\operatorname{Hom}_A(M,N))_f$ s'identifie canoniquement à
$\Gamma(D(f),\widetilde{(\operatorname{Hom}_A(M,N))})$
(\hyperref[I.1.3.6-fr]{1.3.6} et \hyperref[I.1.3.7-fr]{1.3.7}), et les
isomorphismes canoniques
$\mathscr{G}(D(f))\simeq
\Gamma(D(f),\widetilde{(\operatorname{Hom}_A(M,N))})$ ainsi obtenus sont
compatibles avec les opérateurs de restriction
(\hyperref[0.1.4.2-fr]{0, 1.4.2})~; ils définissent donc un isomorphisme
canonique
\[
  \mathscr{H}\!om_{\widetilde{A}}(\widetilde{M},\widetilde{N})
  \simeq\widetilde{(\operatorname{Hom}_A(M,N))}.
\]

\begin{env}[1.3.13]
\label{I.1.3.13-fr}
Soit maintenant $B$ une $A$-algèbre (commutative)~; cela peut s'interpréter
en disant que $B$ est un $A$-module et qu'on s'est donné un élément $e\in B$
et un $A$-homomorphisme
$\varphi:B\otimes_A B\to B$, de sorte que~: $1^\circ$ Les diagrammes
\[
  \xymatrix{
    B\otimes_A B\otimes_A B\ar[r]^{\varphi\otimes 1}
      \ar[d]_{1\otimes\varphi} &
    B\otimes_A B\ar[d]^\varphi & &
    B\otimes_A B\ar[rr]^\sigma\ar[rd]_\varphi & &
    B\otimes_A B\ar[dl]^\varphi\\
    B\otimes_A B\ar[r]_\varphi &
    B & & & B
  }
\]
($\sigma$ symétrie canonique) soient commutatifs~; $2^\circ$
$\varphi(e\otimes x)=\varphi(x\otimes e)=x$. En vertu de
(\hyperref[I.1.3.12-fr]{1.3.12}), l'homomorphisme
$\widetilde{\varphi}:\widetilde{B}\otimes_{\widetilde{A}}\widetilde{B}
\to\widetilde{B}$ de $\widetilde{A}$-Modules vérifie les conditions
analogues, donc définit sur $\widetilde{B}$ une structure de
$\widetilde{A}$-Algèbre. De la même manière, la donnée d'un $B$-module $N$
revient à la donnée d'un $A$-module $N$ et d'un $A$-homomorphisme
$\psi:B\otimes_A N\to N$ tel que le diagramme
\[
  \xymatrix{
    B\otimes_A B\otimes_A N\ar[r]^{\varphi\otimes 1}
      \ar[d]_{1\otimes\psi} &
    B\otimes_A N\ar[d]^\psi\\
    B\otimes_A N\ar[r]_\psi & N
  }
\]
soit commutatif et $\psi(e\otimes n)=n$~; l'homomorphisme
$\widetilde{\psi}:\widetilde{B}\otimes_{\widetilde{A}}\widetilde{N}
\to\widetilde{N}$ vérifiera la condition analogue, et définit donc sur
$\widetilde{N}$ une structure de $\widetilde{B}$-Module.

De la même façon, on voit que si $u:B\to B'$ (resp. $v:N\to N'$) est un
homomorphisme de $A$-algèbres (resp. de $B$-modules), $\widetilde{u}$
(resp. $\widetilde{v}$) est un homomorphisme de $\widetilde{A}$-Algèbres
(resp. de $\widetilde{B}$-Modules), $\operatorname{Ker}\widetilde{u}$ un
$\widetilde{B}$-Idéal (resp. $\operatorname{Ker}\widetilde{v}$,
$\operatorname{Coker}\widetilde{v}$ et $\operatorname{Im}\widetilde{v}$
des $\widetilde{B}$-Modules). Si $N$ est un $B$-module, $\widetilde{N}$ est
un $\widetilde{B}$-Module de type fini si et seulement si $N$ est un
$B$-module de type fini (\hyperref[0.5.2.3-fr]{0, 5.2.3}).
\oldpage[I]{90}
Si $M$, $N$ sont deux $B$-modules, le $\widetilde{B}$-Module
$\widetilde{M}\otimes_{\widetilde{B}}\widetilde{N}$ s'identifie
canoniquement à $\widetilde{(M\otimes_B N)}$~; de même
$\mathscr{H}\!om_{\widetilde{B}}(\widetilde{M},\widetilde{N})$ s'identifie
canoniquement à $\widetilde{(\operatorname{Hom}_B(M,N))}$ lorsque $M$ admet
une présentation finie~; les démonstrations sont les mêmes que dans
(\hyperref[I.1.3.12-fr]{1.3.12}).

Si $\mathfrak{J}$ est un idéal de $B$, $N$ un $B$-module, on a
$\widetilde{(\mathfrak{J}N)}=\widetilde{\mathfrak{J}}\cdot\widetilde{N}$.

Enfin, si $B$ est une $A$-algèbre graduée par des sous-$A$-modules $B_n$
($n\in\mathbb{Z}$), la $\widetilde{A}$-Algèbre $\widetilde{B}$, somme directe
des $\widetilde{A}$-Modules $\widetilde{B_n}$
(\hyperref[I.1.3.9-fr]{1.3.9}) est graduée par ces sous-$\widetilde{A}$-Modules,
l'axiome de la graduation exprimant que l'image de l'homomorphisme
$B_m\otimes_A B_n\to B$ est contenu dans $B_{m+n}$. De même, si $M$ est un
$B$-module gradué par des sous-modules $M_n$, $\widetilde{M}$ est un
$\widetilde{B}$-Module gradué par les $\widetilde{M_n}$.
\end{env}

\begin{env}[1.3.14]
\label{I.1.3.14-fr}
Si $B$ est une $A$-algèbre, $M$ un sous-module de $B$, la
sous-$\widetilde{A}$-Algèbre de $\widetilde{B}$ engendrée par $\widetilde{M}$
(\hyperref[0.4.1.3-fr]{0, 4.1.3}) est la sous-$\widetilde{A}$-Algèbre
$\widetilde{C}$, en désignant par $C$ la sous-algèbre de $B$ engendrée par
$M$. En effet, $C$ est la somme des sous-modules de $B$ images des
homomorphismes $\bigotimes_A^n M\to B$ ($n\geq0$), et il suffit d'appliquer
(\hyperref[I.1.3.9-fr]{1.3.9}) et (\hyperref[I.1.3.12-fr]{1.3.12}).
\end{env}

\subsection{Faisceaux quasi-cohérents sur un spectre premier.}
\label{subsection:I.1.4-fr}

\begin{theorem}[1.4.1]
\label{I.1.4.1-fr}
Soient $X$ le spectre premier d'un anneau $A$, $V$ une partie ouverte
quasi-compacte de $X$, $\mathscr{F}$ un $(\mathscr{O}_X|V)$-Module. Les
quatre conditions suivantes sont équivalentes~:
\begin{enumerate}
  \item[\textit{a})] Il existe un $A$-module $M$ tel que $\mathscr{F}$ soit
  isomorphe à $\widetilde{M}|V$.
  \item[\textit{b})] Il existe un recouvrement ouvert fini $(V_i)$ de $V$ par des
  ensembles de la forme $D(f_i)$ ($f_i\in A$) contenus dans $V$, tels que,
  pour tout $i$, $\mathscr{F}|V_i$ soit isomorphe à un faisceau de la forme
  $\widetilde{M_i}$, où $M_i$ est un $A_{f_i}$-module.
  \item[\textit{c})] Le faisceau $\mathscr{F}$ est quasi-cohérent
  (\hyperref[0.5.1.3-fr]{0, 5.1.3}).
  \item[\textit{d})] Les deux propriétés suivantes ont lieu~:
  \begin{enumerate}
    \item[\textit{d}\,1)] Pour tout $f\in A$ tel que $D(f)\subset V$ et toute
    section $s\in\Gamma(D(f),\mathscr{F})$, il existe un entier $n\geq0$ tel
    que $f^n s$ se prolonge en une section de $\mathscr{F}$ sur $V$.
    \item[\textit{d}\,2)] Pour tout $f\in A$ tel que $D(f)\subset V$ et toute
    section $t\in\Gamma(V,\mathscr{F})$ telle que la restriction de $t$ à
    $D(f)$ soit $0$, il existe un entier $n\geq0$ tel que $f^n t=0$.
  \end{enumerate}
\end{enumerate}
\end{theorem}

(Dans l'écriture des conditions \textit{d}\,1) et \textit{d}\,2), on a tacitement identifié $A$
et $\Gamma(\widetilde{A})$ en vertu de
\hyperref[I.1.3.7-fr]{1.3.7}.)

Le fait que \textit{a}) entraîne \textit{b}) est conséquence immédiate de
(\hyperref[I.1.3.6-fr]{1.3.6}) et du fait que les $D(f_i)$ forment une base
de la topologie de $X$ (\hyperref[I.1.1.10-fr]{1.1.10}). Comme tout
$A$-module est isomorphe au conoyau d'un homomorphisme
$A^{(I)}\to A^{(J)}$, (\hyperref[I.1.3.9-fr]{1.3.9}) prouve que tout faisceau
associé à un $A$-module est quasi-cohérent~; donc \textit{b}) entraîne \textit{c}).
Réciproquement, si $\mathscr{F}$ est quasi-cohérent, tout $x\in V$ possède
un voisinage de la forme $D(f)\subset V$ tel que $\mathscr{F}|D(f)$ soit
isomorphe au conoyau d'un homomorphisme
$\widetilde{A_f}^{(I)}\to\widetilde{A_f}^{(J)}$, donc au faisceau associé au
module $\widetilde{N}$, conoyau de l'homomorphisme
$A_f^{(I)}\to A_f^{(J)}$ correspondant
(\hyperref[I.1.3.8-fr]{1.3.8} et \hyperref[I.1.3.9-fr]{1.3.9})~; comme $V$
est quasi-compact, il est clair que \textit{c}) entraîne \textit{b}).
\oldpage[I]{91}
Pour prouver que \textit{b}) entraîne \textit{d}\,1) et \textit{d}\,2),
supposons d'abord que $V=D(g)$ pour un $g\in A$, et que $\mathscr{F}$ soit
isomorphe à un faisceau $\widetilde{N}$ associé à un $A_g$-module $N$~; en
remplaçant $X$ par $V$ et $A$ par $A_g$
(\hyperref[I.1.3.6-fr]{1.3.6}), on peut se ramener au cas où $g=1$. Alors
$\Gamma(D(f),\widetilde{N})$ et $N_f$ s'identifient canoniquement
(\hyperref[I.1.3.6-fr]{1.3.6} et \hyperref[I.1.3.7-fr]{1.3.7}), donc une
section $s\in\Gamma(D(f),\widetilde{N})$ s'identifie à un élément de la
forme $z/f^n$, où $z\in N$~; la section $f^n s$ s'identifie à l'élément
$z/1$ de $N_f$, et est par suite restriction à $D(f)$ de la section de
$\widetilde{N}$ sur $X$ identifiée à l'élément $z\in N$~; d'où
\textit{d}\,1) dans ce cas. De même, $t\in\Gamma(X,\widetilde{N})$ est
identifié à un élément $z'\in N$, la restriction de $t$ à $D(f)$ est
identifiée à l'image $z'/1$ de $z'$ dans $N_f$, et dire que cette image est
nulle signifie qu'il existe $n\geq0$ tel que $f^n z'=0$ dans $N$, ou, ce
qui revient au même, $f^n t=0$.

Pour achever de prouver que \textit{b}) entraîne \textit{d}\,1) et
\textit{d}\,2), il suffira d'établir le lemme suivant~:

\begin{lemma}[1.4.1.1]
\label{I.1.4.1.1-fr}
Supposons que $V$ soit réunion finie d'ensembles de la forme $D(g_i)$, et
que chacun des faisceaux $\mathscr{F}|D(g_i)$,
$\mathscr{F}|(D(g_i)\cap D(g_j))=\mathscr{F}|D(g_i g_j)$ vérifie
\textit{d}\,1) et \textit{d}\,2)~; alors $\mathscr{F}$ possède les deux
propriétés suivantes~:
\begin{enumerate}
  \item[\textit{d}'\,1)] Pour tout $f\in A$ et toute section
  $s\in\Gamma(D(f)\cap V,\mathscr{F})$, il existe un entier $n\geq0$ tel
  que $f^n s$ se prolonge en une section de $\mathscr{F}$ sur $V$.
  \item[\textit{d}'\,2)] Pour tout $f\in A$ et toute section
  $t\in\Gamma(V,\mathscr{F})$ telle que la restriction de $t$ à
  $D(f)\cap V$ soit $0$, il existe un entier $n\geq0$ tel que $f^n t=0$.
\end{enumerate}
\end{lemma}

Prouvons d'abord \textit{d}'\,2)~: comme
$D(f)\cap D(g_i)=D(fg_i)$, il existe pour chaque $i$ un entier $n_i$ tel
que la restriction de $(fg_i)^{n_i}t$ à $D(g_i)$ soit nulle~: comme
l'image de $g_i$ dans $A_{g_i}$ est inversible, la restriction de
$f^{n_i}t$ à $D(g_i)$ est aussi nulle~; prenant pour $n$ le plus grand des
$n_i$, on a \textit{d}'\,2).

Pour démontrer \textit{d}'\,1), appliquons \textit{d}\,1) au faisceau
$\mathscr{F}|D(g_i)$~: il existe un entier $n_i\geq0$ et une section $s_i'$
de $\mathscr{F}$ sur $D(g_i)$ prolongeant la restriction de
$(fg_i)^{n_i}s$ à $D(fg_i)$~; comme l'image de $g_i$ dans $A_{g_i}$ est
inversible, il y a une section $s_i$ de $\mathscr{F}$ sur $D(g_i)$ telle
que $s_i'=g_i^{n_i}s_i$, et $s_i$ prolonge la restriction de $f^{n_i}s$ à
$D(fg_i)$~; on peut en outre supposer tous les $n_i$ égaux à un même
entier $n$. Par construction, la restriction de $s_i-s_j$ à
$D(f)\cap D(g_i)\cap D(g_j)=D(fg_i g_j)$ est nulle~; d'après
\textit{d}\,2) appliqué au faisceau $\mathscr{F}|D(g_i g_j)$, il existe un
entier $m_{ij}\geq0$ tel que la restriction à $D(g_i g_j)$ de
$(fg_i g_j)^{m_{ij}}(s_i-s_j)$ soit nulle~; comme l'image de $g_i g_j$
dans $A_{g_i g_j}$ est inversible, la restriction de
$f^{m_{ij}}(s_i-s_j)$ à $D(g_i g_j)$ est nulle. On peut alors supposer tous
les $m_{ij}$ égaux à un même entier $m$, et il existe donc une section
$s'\in\Gamma(V,\mathscr{F})$ prolongeant les $f^m s_i$~; cette section
prolonge par suite $f^{n+m}s$, d'où \textit{d}'\,1).

Reste à prouver que \textit{d}\,1) et \textit{d}\,2) entraînent
\textit{a}). Montrons d'abord que \textit{d}\,1) et \textit{d}\,2)
entraînent que ces conditions sont vérifiées pour tout faisceau
$\mathscr{F}|D(g)$, où $g\in A$ est tel que $D(g)\subset V$. C'est évident
pour \textit{d}\,1)~; d'autre part, si $t\in\Gamma(D(g),\mathscr{F})$ est
telle que sa restriction à $D(f)\subset D(g)$ soit nulle, il existe par
\textit{d}\,1) un entier $m\geq0$ tel que $g^m t$ se
\oldpage[I]{92}
prolonge en une section $s$ de $\mathscr{F}$ sur $V$~; appliquant
\textit{d}\,2), on voit qu'il existe un entier $n\geq0$ tel que
$f^n g^m t=0$, et comme l'image de $g$ dans $A_g$ est inversible, $f^n t=0$.

Cela étant, comme $V$ est quasi-compact, le lemme
(\hyperref[I.1.4.1.1-fr]{1.4.1.1}) prouve que les conditions
\textit{d}'\,1) et \textit{d}'\,2) sont vérifiées. Considérons alors le
$A$-module $M=\Gamma(V,\mathscr{F})$, et définissons un homomorphisme de
$\widetilde{A}$-Modules $u:\widetilde{M}\to j_*(\mathscr{F})$, où $j$ est
l'injection canonique $V\to X$. Comme les $D(f)$ forment une base de la
topologie de $X$, il suffit, pour chaque $f\in A$, de définir un
homomorphisme
$u_f:M_f\to\Gamma(D(f),j_*(\mathscr{F}))=\Gamma(D(f)\cap V,\mathscr{F})$,
avec les conditions de compatibilité usuelles
(\hyperref[0.3.2.5-fr]{0, 3.2.5}). Comme l'image canonique de $f$ dans
$A_f$ est inversible, l'homomorphisme de restriction
$M=\Gamma(V,\mathscr{F})\to\Gamma(D(f)\cap V,\mathscr{F})$ se factorise en
$M\longrightarrow M_f\xrightarrow{u_f}\Gamma(D(f)\cap V,\mathscr{F})$
(\hyperref[0.1.2.4-fr]{0, 1.2.4}), et la vérification des conditions de
compatibilité pour $D(g)\subset D(f)$ est immédiate. Cela étant, montrons
que la condition \textit{d}'\,1) (resp. \textit{d}'\,2)) entraîne que chacun
des $u_f$ est surjectif (resp. injectif), ce qui prouvera que $u$ est
\emph{bijectif}, et par suite que $\mathscr{F}$ est la restriction à $V$
d'un $\widetilde{A}$-Module isomorphe à $\widetilde{M}$. Or, si
$s\in\Gamma(D(f)\cap V,\mathscr{F})$, il existe d'après \textit{d}'\,1) un
entier $n\geq0$ tel que $f^n s$ se prolonge en une section $z\in M$~; on a
alors $u_f(z/f^n)=s$, donc $u_f$ est surjectif. De même, si $z\in M$ est tel
que $u_f(z/1)=0$, cela signifie que la restriction à $D(f)\cap V$ de la
section $z$ est nulle~; d'après \textit{d}'\,2), il existe un entier
$n\geq0$ tel que $f^n z=0$, d'où $z/1=0$ dans $M_f$, et $u_f$ est donc
injectif.

\par\hfill C.Q.F.D.\par

\begin{corollary}[1.4.2]
\label{I.1.4.2-fr}
Tout faisceau quasi-cohérent sur un ouvert quasi-compact de $X$ est induit
par un faisceau quasi-cohérent sur $X$.
\end{corollary}

\begin{corollary}[1.4.3]
\label{I.1.4.3-fr}
Toute $\mathscr{O}_X$-Algèbre quasi-cohérente sur
$X=\operatorname{Spec}(A)$ est isomorphe à une $\mathscr{O}_X$-Algèbre de
la forme $\widetilde{B}$, où $B$ est une algèbre sur $A$~; tout
$\widetilde{B}$-Module quasi-cohérent est isomorphe à un
$\widetilde{B}$-Module de la forme $\widetilde{N}$, où $N$ est un
$B$-module.
\end{corollary}

En effet, une $\mathscr{O}_X$-Algèbre quasi-cohérente est un
$\mathscr{O}_X$-Module quasi-cohérent, donc de la forme $\widetilde{B}$, où
$B$ est un $A$-module~; le fait que $B$ soit une $A$-algèbre résulte de la
caractérisation de la structure de $\mathscr{O}_X$-Algèbre à l'aide de
l'homomorphisme
$\widetilde{B}\otimes_{\widetilde{A}}\widetilde{B}\to\widetilde{B}$ de
$\widetilde{A}$-Modules, ainsi que de
(\hyperref[I.1.3.12-fr]{1.3.12}). Si $\mathscr{G}$ est un
$\widetilde{B}$-Module quasi-cohérent, il suffit de montrer que $\mathscr{G}$
est aussi un $\widetilde{A}$-Module quasi-cohérent pour conclure ensuite de
la même façon~; comme la question est locale, on peut, en se restreignant à
un ouvert de $X$ de la forme $D(f)$, supposer que $\mathscr{G}$ est le
conoyau d'un homomorphisme
$\widetilde{B}^{(I)}\to\widetilde{B}^{(J)}$ de $\widetilde{B}$-Modules (et
\emph{a fortiori} de $\widetilde{A}$-Modules)~; la proposition résulte alors
de (\hyperref[I.1.3.8-fr]{1.3.8}) et
(\hyperref[I.1.3.9-fr]{1.3.9}).

\subsection{Faisceaux cohérents sur un spectre premier.}
\label{subsection:I.1.5-fr}

\begin{theorem}[1.5.1]
\label{I.1.5.1-fr}
Soient $A$ un anneau \emph{noethérien}, $X=\operatorname{Spec}(A)$ son spectre
premier, $V$ une partie ouverte de $X$, $\mathscr{F}$ un
$(\mathscr{O}_X|V)$-Module. Les conditions suivantes sont équivalentes~:
\begin{enumerate}
  \item[\textit{a})] $\mathscr{F}$ est cohérent.
  \item[\textit{b})] $\mathscr{F}$ est de type fini et quasi-cohérent.
  \item[\textit{c})] Il existe un $A$-module $M$ de type fini tel que
  $\mathscr{F}$ soit isomorphe au faisceau $\widetilde{M}|V$.
\end{enumerate}
\end{theorem}

\oldpage[I]{93}
\textit{a}) implique trivialement \textit{b}). Pour voir que \textit{b})
implique \textit{c}), remarquons déjà, puisque $V$ est quasi-compact
(\hyperref[0.2.2.3-fr]{0, 2.2.3}), que $\mathscr{F}$ est isomorphe à un
faisceau $\widetilde{N}|V$, où $N$ est un $A$-module
(\hyperref[I.1.4.1-fr]{1.4.1}). Or, on a
$N=\varinjlim M_\lambda$, où $M_\lambda$ parcourt l'ensemble des
sous-$A$-modules de type fini de $N$, d'où
(\hyperref[I.1.3.9-fr]{1.3.9})
$\mathscr{F}=\widetilde{N}|V=\varinjlim\widetilde{M_\lambda}|V$~; mais comme
$\mathscr{F}$ est de type fini, et $V$ quasi-compact, il existe un indice
$\lambda$ tel que $\mathscr{F}=\widetilde{M_\lambda}|V$
(\hyperref[0.5.2.3-fr]{0, 5.2.3}).

Montrons enfin que \textit{c}) entraîne \textit{a}). Il est clair que
$\mathscr{F}$ est alors de type fini
(\hyperref[I.1.3.6-fr]{1.3.6} et \hyperref[I.1.3.9-fr]{1.3.9})~; en outre,
la question étant locale, on peut se borner au cas où $V=D(f)$, $f\in A$.
Comme $A_f$ est noethérien, on voit finalement que tout revient à prouver
que le noyau d'un homomorphisme $\widetilde{A^n}\to\widetilde{M}$, où $M$
est un $A$-module, est de type fini. Or, un tel homomorphisme est de la
forme $\widetilde{u}$, où $u$ est un homomorphisme $A^n\to M$
(\hyperref[I.1.3.8-fr]{1.3.8}), et si $P=\operatorname{Ker}u$, on a
$\widetilde{P}=\operatorname{Ker}\widetilde{u}$
(\hyperref[I.1.3.9-fr]{1.3.9}). Comme $A$ est noethérien, $P$ est de type
fini, ce qui achève la démonstration.

\begin{corollary}[1.5.2]
\label{I.1.5.2-fr}
Sous les hypothèses de (\hyperref[I.1.5.1-fr]{1.5.1}), le faisceau
$\mathscr{O}_X$ est un faisceau cohérent d'anneaux.
\end{corollary}

\begin{corollary}[1.5.3]
\label{I.1.5.3-fr}
Sous les hypothèses de (\hyperref[I.1.5.1-fr]{1.5.1}), tout faisceau
cohérent sur un ouvert de $X$ est induit par un faisceau cohérent sur $X$.
\end{corollary}

\begin{corollary}[1.5.4]
\label{I.1.5.4-fr}
Sous les hypothèses de (\hyperref[I.1.5.1-fr]{1.5.1}), tout
$\mathscr{O}_X$-Module quasi-cohérent $\mathscr{F}$ est limite inductive
des sous-$\mathscr{O}_X$-Modules cohérents de $\mathscr{F}$.
\end{corollary}

En effet, $\mathscr{F}=\widetilde{M}$ où $M$ est un $A$-module, et $M$ est
limite inductive de ses sous-modules de type fini~; on conclut par
(\hyperref[I.1.3.9-fr]{1.3.9}) et
(\hyperref[I.1.5.1-fr]{1.5.1}).

\subsection{Propriétés fonctorielles des faisceaux quasi-cohérents sur un
spectre premier.}
\label{subsection:I.1.6-fr}

\begin{env}[1.6.1]
\label{I.1.6.1-fr}
Soient $A$, $A'$ deux anneaux,
\[
  \varphi:A'\to A
\]
un homomorphisme,
\[
  {}^a\varphi:X=\operatorname{Spec}(A)\to
  X'=\operatorname{Spec}(A')
\]
l'application continue associée à $\varphi$
(\hyperref[I.1.2.1-fr]{1.2.1}). Nous allons définir un
\emph{homomorphisme canonique}
\[
  \widetilde{\varphi}:\mathscr{O}_{X'}\to
  {}^a\varphi_*(\mathscr{O}_X)
\]
de faisceaux d'anneaux. Pour tout $f'\in A'$, posons $f=\varphi(f')$~; on a
${}^a\varphi^{-1}(D(f'))=D(f)$
(\hyperref[I.1.2.2.2-fr]{1.2.2.2}). Les anneaux
$\Gamma(D(f'),\widetilde{A'})$ et $\Gamma(D(f),\widetilde{A})$
s'identifient respectivement à $A'_{f'}$ et $A_f$
(\hyperref[I.1.3.6-fr]{1.3.6} et \hyperref[I.1.3.7-fr]{1.3.7}). Or,
l'homomorphisme $\varphi$ définit canoniquement un homomorphisme
$\varphi_{f'}:A'_{f'}\to A_f$
(\hyperref[0.1.5.1-fr]{0, 1.5.1}), autrement dit on a un homomorphisme
d'anneaux
\[
  \Gamma(D(f'),\widetilde{A'})\to
  \Gamma({}^a\varphi^{-1}(D(f')),\widetilde{A})
  =\Gamma(D(f'),{}^a\varphi_*(\widetilde{A})).
\]

\oldpage[I]{94}
En outre, ces homomorphismes satisfont aux conditions de compatibilité
usuelles~: pour $D(f')\supset D(g')$, le diagramme
\[
  \xymatrix{
    \Gamma(D(f'),\widetilde{A'})\ar[r]\ar[d] &
    \Gamma(D(f'),{}^a\varphi_*(\widetilde{A}))\ar[d]\\
    \Gamma(D(g'),\widetilde{A'})\ar[r] &
    \Gamma(D(g'),{}^a\varphi_*(\widetilde{A}))
  }
\]
est commutatif (\hyperref[0.1.5.1-fr]{0, 1.5.1})~; on a donc bien défini
un homomorphisme de $\mathscr{O}_{X'}$-Algèbres, les $D(f')$ formant une
base de la topologie de $X'$ (\hyperref[0.3.2.3-fr]{0, 3.2.3}). Le couple
$\Phi=({}^a\varphi,\widetilde{\varphi})$ est donc un \emph{morphisme}
d'espaces annelés
\[
  \Phi:(X,\mathscr{O}_X)\to(X',\mathscr{O}_{X'})
\]
(\hyperref[0.4.1.1-fr]{0, 4.1.1}).

Notons en outre que, si on pose $x'={}^a\varphi(x)$, l'homomorphisme
$\widetilde{\varphi}_x^\#$ (\hyperref[0.3.7.1-fr]{0, 3.7.1}) n'est autre
que l'homomorphisme
\[
  \varphi_x:A'_{x'}\to A_x
\]
déduit canoniquement de $\varphi:A'\to A$
(\hyperref[0.1.5.1-fr]{0, 1.5.1}). En effet, tout $z'\in A'_{x'}$
s'écrit $g'/f'$, où $f'$, $g'$ sont dans $A'$ et
$f'\notin\mathfrak{j}_{x'}$~; $D(f')$ est donc un voisinage de $x'$ dans
$X'$, et l'homomorphisme
$\Gamma(D(f'),\widetilde{A'})\to
\Gamma({}^a\varphi^{-1}(D(f')),\widetilde{A})$ déduit de
$\widetilde{\varphi}$ n'est autre que $\varphi_{f'}$~; en considérant la
section $s'\in\Gamma(D(f'),\widetilde{A'})$ correspondant à
$g'/f'\in A'_{f'}$, on obtient bien
$\widetilde{\varphi}_x^\#(z')=\varphi(g')/\varphi(f')$ dans $A_x$.
\end{env}

\begin{example}[1.6.2]
\label{I.1.6.2-fr}
Soient $S$ une partie multiplicative de $A$, $\varphi$ l'homomorphisme
canonique $A\to S^{-1}A$~; on a déjà vu
(\hyperref[I.1.2.6-fr]{1.2.6}) que ${}^a\varphi$ est un
\emph{homéomorphisme} de $Y=\operatorname{Spec}(S^{-1}A)$ sur le
sous-espace de $X=\operatorname{Spec}(A)$ formé des $x$ tels que
$\mathfrak{j}_x\cap S=\varnothing$. En outre, pour tout $x$ appartenant à
ce sous-espace, donc de la forme ${}^a\varphi(y)$ avec $y\in Y$,
l'homomorphisme
$\widetilde{\varphi}_y^\#:\mathscr{O}_x\to\mathscr{O}_y$ est
\emph{bijectif} (\hyperref[0.1.2.6-fr]{0, 1.2.6})~; autrement dit,
$\mathscr{O}_Y$ s'identifie au faisceau induit sur $Y$ par $\mathscr{O}_X$.
\end{example}

\begin{proposition}[1.6.3]
\label{I.1.6.3-fr}
Pour tout $A$-module $M$, il existe un isomorphisme canonique fonctoriel du
$\mathscr{O}_{X'}$-Module $(M_{[\varphi]})^{\sim}$ sur l'image directe
$\Phi_*(\widetilde{M})$.
\end{proposition}

Posons pour abréger $M'=M_{[\varphi]}$, et pour tout $f'\in A'$, posons
encore $f=\varphi(f')$. Les modules de sections
$\Gamma(D(f'),\widetilde{M'})$ et $\Gamma(D(f),\widetilde{M})$
s'identifient respectivement aux modules $M'_{f'}$ et $M_f$ (sur
$A'_{f'}$ et $A_f$ respectivement)~; en outre, le $A'_{f'}$-module
$(M_f)_{[\varphi_{f'}]}$ est canoniquement isomorphe à $M'_{f'}$
(\hyperref[0.1.5.2-fr]{0, 1.5.2}). On a donc un isomorphisme fonctoriel de
$\Gamma(D(f'),\widetilde{A'})$-Modules~:
$\Gamma(D(f'),\widetilde{M'})\xrightarrow{\sim}
\Gamma({}^a\varphi^{-1}(D(f')),\widetilde{M})_{[\varphi_{f'}]}$ et ces
isomorphismes satisfont aux conditions de compatibilité usuelles avec les
restrictions (\hyperref[0.1.5.6-fr]{0, 1.5.6}), donc définissent
l'isomorphisme fonctoriel annoncé. On notera que, de façon précise, si
$u:M_1\to M_2$ est un homomorphisme de $A$-modules, il peut être considéré
comme un homomorphisme $(M_1)_{[\varphi]}\to(M_2)_{[\varphi]}$ de
$A'$-modules~; si on note $u_{[\varphi]}$ cet homomorphisme,
$\Phi_*(\widetilde{u})$ s'identifie à $(u_{[\varphi]})^{\sim}$.

Cette démonstration prouve aussi que pour toute $A$-algèbre $B$,
l'isomorphisme
\oldpage[I]{95}
canonique fonctoriel
$(B_{[\varphi]})^{\sim}\xrightarrow{\sim}\Phi_*(\widetilde{B})$ est un
isomorphisme de $\mathscr{O}_{X'}$-Algèbres~; si $M$ est un $B$-module,
l'isomorphisme canonique fonctoriel
$(M_{[\varphi]})^{\sim}\xrightarrow{\sim}\Phi_*(\widetilde{M})$ est un
isomorphisme de $\Phi_*(\widetilde{B})$-Modules.

\begin{corollary}[1.6.4]
\label{I.1.6.4-fr}
Le foncteur image directe $\Phi_*$ est exact sur la catégorie des
$\mathscr{O}_X$-Modules quasi-cohérents.
\end{corollary}

En effet, il est clair que $M_{[\varphi]}$ est un foncteur exact en $M$ et
$\widetilde{M'}$ un foncteur exact en $M'$ (\hyperref[I.1.3.5-fr]{1.3.5}).

\begin{proposition}[1.6.5]
\label{I.1.6.5-fr}
Soient $N'$ un $A'$-module, $N$ le $A$-module
$N'\otimes_{A'}A_{[\varphi]}$~; il existe un isomorphisme canonique
fonctoriel du $\mathscr{O}_X$-Module $\Phi^*(\widetilde{N'})$ sur
$\widetilde{N}$.
\end{proposition}

Remarquons d'abord que $j:z'\mapsto z'\otimes 1$ est un
$A'$-homomorphisme de $N'$ dans $N_{[\varphi]}$~; en effet, par définition,
pour $f'\in A'$, on a
$(f'z')\otimes 1=z'\otimes\varphi(f')=\varphi(f')(z'\otimes 1)$. On en
déduit (\hyperref[I.1.3.8-fr]{1.3.8}) un homomorphisme
$\widetilde{j}:\widetilde{N'}\to(N_{[\varphi]})^{\sim}$ de
$\mathscr{O}_{X'}$-Modules, et en vertu de
(\hyperref[I.1.6.3-fr]{1.6.3}), on peut considérer que $\widetilde{j}$
applique $\widetilde{N'}$ dans $\Phi_*(\widetilde{N})$. Il correspond
canoniquement à cet homomorphisme $\widetilde{j}$ un homomorphisme
$h=\widetilde{j}^{\#}$ de $\Phi^*(\widetilde{N'})$ dans $\widetilde{N}$
(\hyperref[0.4.4.3-fr]{0, 4.4.3})~; on va voir que pour chaque fibre,
$h_x$ est \emph{bijectif}. Posons $x'={}^a\varphi(x)$ et soit $f'\in A'$
tel que $x'\in D(f')$~; soit $f=\varphi(f')$. L'anneau
$\Gamma(D(f),\widetilde{A})$ s'identifie à $A_f$, les modules
$\Gamma(D(f),\widetilde{N})$ et $\Gamma(D(f'),\widetilde{N'})$ à $N_f$ et
$N'_{f'}$ respectivement~; soient
$s'\in\Gamma(D(f'),\widetilde{N'})$, identifiée à $n'/{f'}^p$ ($n'\in N'$),
$s$ son image par $\widetilde{j}$ dans $\Gamma(D(f),\widetilde{N})$~; $s$
est identifiée à $(n'\otimes 1)/f^p$. Soit d'autre part
$t\in\Gamma(D(f),\widetilde{A})$, identifiée à $g/f^q$ ($g\in A$)~; alors,
par définition, on a
$h_x(s'_{x'}\otimes t_x)=t_x\cdot s_x$
(\hyperref[0.4.4.3-fr]{0, 4.4.3}). Mais on peut identifier canoniquement
$N_f$ à
$N'_{f'}\otimes_{A'_{f'}}(A_f)_{[\varphi_{f'}]}$
(\hyperref[0.1.5.4-fr]{0, 1.5.4})~; $s$ correspond alors à l'élément
$(n'/{f'}^p)\otimes 1$, et la section $y\mapsto t_y\cdot s_y$ à
$(n'/{f'}^p)\otimes(g/f^q)$. Les diagrammes de compatibilité de
(\hyperref[0.1.5.6-fr]{0, 1.5.6}) montrent que $h_x$ n'est autre que
l'isomorphisme canonique
\[
  \label{I.1.6.5.1-fr}
  N'_{x'}\otimes_{A'_{x'}}(A_x)_{[\varphi_x]}
  \xrightarrow{\sim}N_x=(N'\otimes_{A'}A_{[\varphi]})_x.
  \tag{1.6.5.1}
\]

En outre, soit $v:N'_1\to N'_2$ un homomorphisme de $A'$-modules~; comme
$\widetilde{v}_{x'}=v_{x'}$ pour tout $x'\in X'$, il résulte aussitôt de ce
qui précède que $\Phi^*(\widetilde{v})$ s'identifie canoniquement à
$(v\otimes 1)^{\sim}$, ce qui achève la démonstration de
(\hyperref[I.1.6.5-fr]{1.6.5}).

Si $B'$ est une $A'$-algèbre, l'isomorphisme canonique de
$\Phi^*(\widetilde{B'})$ sur
$(B'\otimes_{A'}A_{[\varphi]})^{\sim}$ est un isomorphisme de
$\mathscr{O}_X$-Algèbres~; si en outre $N'$ est un $B'$-module,
l'isomorphisme canonique de $\Phi^*(\widetilde{N'})$ sur
$(N'\otimes_{A'}A_{[\varphi]})^{\sim}$ est un isomorphisme de
$\Phi^*(\widetilde{B'})$-Modules.

\begin{corollary}[1.6.6]
\label{I.1.6.6-fr}
Les sections de $\Phi^*(\widetilde{N'})$ images canoniques des sections $s'$,
où $s'$ parcourt le $A'$-module $\Gamma(\widetilde{N'})$, engendrent le
$A$-module $\Gamma(\Phi^*(\widetilde{N'}))$.
\end{corollary}

En effet, ces images s'identifient aux éléments $z'\otimes 1$ de $N$,
lorsqu'on identifie $N'$ et $N$ à $\Gamma(\widetilde{N'})$ et
$\Gamma(\widetilde{N})$ respectivement (\hyperref[I.1.3.7-fr]{1.3.7}) et
que $z'$ parcourt $N'$.

\begin{env}[1.6.7]
\label{I.1.6.7-fr}
Dans la démonstration de (\hyperref[I.1.6.5-fr]{1.6.5}), on a prouvé en
passant que l'application canonique
(\hyperref[0.4.4.3.2-fr]{0, 4.4.3.2})
$\rho:\widetilde{N'}\to\Phi_*(\Phi^*(\widetilde{N'}))$ n'est autre que
l'homomorphisme $\widetilde{j}$,
\oldpage[I]{96}
où $j:N'\to N'\otimes_{A'}A_{[\varphi]}$ est l'homomorphisme
$z'\mapsto z'\otimes 1$. De même, l'application canonique
(\hyperref[0.4.4.3.3-fr]{0, 4.4.3.3})
$\sigma:\Phi^*(\Phi_*(\widetilde{M}))\to\widetilde{M}$ n'est autre que
$\widetilde{p}$, où
$p:M_{[\varphi]}\otimes_{A'}A_{[\varphi]}\to M$ est l'homomorphisme
canonique qui, à tout produit tensoriel $z\otimes a$ ($z\in M$, $a\in A$)
fait correspondre $a\cdot z$~; cela résulte aussitôt des définitions
(\hyperref[0.3.7.1-fr]{0, 3.7.1},
\hyperref[0.4.4.3-fr]{0, 4.4.3}, et
\hyperref[I.1.3.7-fr]{I, 1.3.7}).

On en conclut (\hyperref[0.4.4.3-fr]{0, 4.4.3} et
\hyperref[0.3.5.4.4-fr]{3.5.4.4}) que si
$v:N'\to M_{[\varphi]}$ est un $A'$-homomorphisme, on a
$\widetilde{v}^{\#}=(v\otimes 1)^{\sim}$.
\end{env}

\begin{env}[1.6.8]
\label{I.1.6.8-fr}
Soient $N'_1$, $N'_2$ deux $A'$-modules, $N'_1$ étant supposé avoir une
\emph{présentation finie}~; il résulte alors de
(\hyperref[I.1.6.7-fr]{1.6.7}) et
(\hyperref[I.1.3.12-fr]{1.3.12, (ii)}) que l'homomorphisme canonique
(\hyperref[0.4.4.6-fr]{0, 4.4.6})
\[
  \Phi^*\bigl(\mathscr{H}\!om_{\widetilde{A'}}
    (\widetilde{N'_1},\widetilde{N'_2})\bigr)
  \longrightarrow
  \mathscr{H}\!om_{\widetilde A}
    \bigl(\Phi^*(\widetilde{N'_1}),\Phi^*(\widetilde{N'_2})\bigr)
\]
n'est autre que $\widetilde{\gamma}$, en désignant par $\gamma$
l'homomorphisme canonique de $A$-modules
\[
  \operatorname{Hom}_{A'}(N'_1,N'_2)\otimes_{A'}A
  \longrightarrow
  \operatorname{Hom}_A(N'_1\otimes_{A'}A,N'_2\otimes_{A'}A).
\]
\end{env}

\begin{env}[1.6.9]
\label{I.1.6.9-fr}
Soient $\mathfrak{J}'$ un idéal de $A'$, $M$ un $A$-module~; comme par
définition $\widetilde{\mathfrak{J}'}\widetilde M$ est l'image de
l'homomorphisme canonique
$\Phi^*(\widetilde{\mathfrak{J}'})\otimes_{\widetilde A}\widetilde M
\to\widetilde M$, il résulte de (\hyperref[I.1.6.5-fr]{1.6.5}) et
(\hyperref[I.1.3.12-fr]{1.3.12, (i)}) que
$\widetilde{\mathfrak{J}'}\widetilde M$ s'identifie canoniquement à
$(\mathfrak{J}'M)^{\sim}$~; en particulier,
$\Phi^*(\widetilde{\mathfrak{J}'})\widetilde A$ s'identifie à
$(\mathfrak{J}'A)^{\sim}$, et en tenant compte de l'exactitude à droite du
foncteur $\Phi^*$, la $\widetilde A$-Algèbre
$\Phi^*((A'/\mathfrak{J}')^{\sim})$ s'identifie à
$(A/\mathfrak{J}'A)^{\sim}$.
\end{env}

\begin{env}[1.6.10]
\label{I.1.6.10-fr}
Soient $A''$ un troisième anneau, $\varphi'$ un homomorphisme
$A''\to A'$, et posons $\varphi''=\varphi\circ\varphi'$. Il résulte
aussitôt des définitions que
${}^a\varphi''=({}^a\varphi')\circ({}^a\varphi)$, et
$\widetilde{\varphi''}=\widetilde\varphi\circ\widetilde{\varphi'}$
(\hyperref[0.1.5.7-fr]{0, 1.5.7}). On en conclut que l'on a
$\Phi''=\Phi'\circ\Phi$~; en d'autres termes,
$(\operatorname{Spec} A,\widetilde A)$ est un \emph{foncteur} en $A$ de la
catégorie des anneaux dans celle des espaces annelés.
\end{env}

\subsection{Caractérisation des morphismes de schémas affines.}
\label{subsection:I.1.7-fr}

\begin{definition}[1.7.1]
\label{I.1.7.1-fr}
On dit qu'un espace annelé $(X,\mathscr{O}_X)$ est un \emph{schéma affine}
s'il est isomorphe à un espace annelé de la forme
$(\operatorname{Spec}(A),\widetilde A)$, où $A$ est un anneau~; on dit alors
que $\Gamma(X,\mathscr{O}_X)$, qui s'identifie canoniquement à l'anneau $A$
(\hyperref[I.1.3.7-fr]{1.3.7}) est l'anneau du schéma affine
$(X,\mathscr{O}_X)$, et on le note $A(X)$ quand aucune confusion n'en résulte.
\end{definition}

Par abus de langage, quand nous parlerons du \emph{schéma affine}
$\operatorname{Spec}(A)$, il s'agira toujours de l'espace annelé
$(\operatorname{Spec}(A),\widetilde A)$.

\begin{env}[1.7.2]
\label{I.1.7.2-fr}
Soient $A$, $B$ deux anneaux, $(X,\mathscr{O}_X)$,
$(Y,\mathscr{O}_Y)$ les schémas affines correspondants sur les spectres
premiers $X=\operatorname{Spec}(A)$, $Y=\operatorname{Spec}(B)$. On a vu
(\hyperref[I.1.6.1-fr]{1.6.1}) qu'à tout homomorphisme d'anneaux
$\varphi:B\to A$ correspond un morphisme
$\Phi=({}^a\varphi,\widetilde\varphi)=\operatorname{Spec}(\varphi):
(X,\mathscr{O}_X)\to(Y,\mathscr{O}_Y)$. On notera que $\varphi$ est
entièrement déterminé par $\Phi$, car on a par définition
\[
  \varphi=\Gamma(\widetilde\varphi):
  \Gamma(\widetilde B)\longrightarrow
  \Gamma({}^a\varphi_*(\widetilde A))=\Gamma(\widetilde A).
\]
\end{env}

\begin{theorem}[1.7.3]
\label{I.1.7.3-fr}
Soient $(X,\mathscr{O}_X)$, $(Y,\mathscr{O}_Y)$ deux schémas affines. Pour
qu'un morphisme d'espaces annelés
$(\psi,\theta):(X,\mathscr{O}_X)\to(Y,\mathscr{O}_Y)$ soit de la forme
$({}^a\varphi,\widetilde\varphi)$, où $\varphi$ est un homomorphisme
d'anneaux $A(Y)\to A(X)$, il faut et il suffit que, pour tout $x\in X$,
$\theta_x^{\#}$ soit un homomorphisme local~:
$\mathscr{O}_{\psi(x)}\to\mathscr{O}_x$.
\end{theorem}

\oldpage[I]{97}
Posons $A=A(X)$, $B=A(Y)$. La condition est nécessaire, car on a vu
(\hyperref[I.1.6.1-fr]{1.6.1}) que $\widetilde\varphi_x^{\#}$ est
l'homomorphisme de $B_{{}^a\varphi(x)}$ dans $A_x$ déduit canoniquement de
$\varphi$, et par définition de
${}^a\varphi(x)=\varphi^{-1}(\mathfrak{j}_x)$, cet homomorphisme est local.

Prouvons que la condition est suffisante. Par définition, $\theta$ est un
homomorphisme $\mathscr{O}_Y\to\psi_*(\mathscr{O}_X)$, et on en déduit
canoniquement un homomorphisme d'anneaux
\[
  \varphi=\Gamma(\theta):
  B=\Gamma(Y,\mathscr{O}_Y)
  \longrightarrow
  \Gamma(Y,\psi_*(\mathscr{O}_X))=\Gamma(X,\mathscr{O}_X)=A.
\]

L'hypothèse sur $\theta_x^{\#}$ permet de déduire de cet homomorphisme, par
passage aux quotients, un monomorphisme $\theta^x$ du corps des restes
$k(\psi(x))$ dans le corps des restes $k(x)$, tel que, pour toute section
$f\in\Gamma(Y,\mathscr{O}_Y)=B$, on ait
$\theta^x(f(\psi(x)))=\varphi(f)(x)$. La relation $f(\psi(x))=0$ est donc
équivalente à $\varphi(f)(x)=0$, ce qui signifie que
$\mathfrak{j}_{\psi(x)}=\mathfrak{j}_{{}^a\varphi(x)}$, et s'écrit encore
$\psi(x)={}^a\varphi(x)$ pour tout $x\in X$, ou $\psi={}^a\varphi$. On sait
aussi que le diagramme
\phantomsection\label{I.1.7.3.diagram-fr}
\[
  \xymatrix{
    B=\Gamma(Y,\mathscr{O}_Y)\ar[r]^\varphi\ar[d] &
    \Gamma(X,\mathscr{O}_X)=A\ar[d]\\
    B_{\psi(x)}\ar[r]_{\theta_x^{\#}} & A_x
  }
\]
est commutatif (\hyperref[0.3.7.2-fr]{0, 3.7.2}), ce qui signifie que
$\theta_x^{\#}$ est égal à l'homomorphisme
$\varphi_x:B_{\psi(x)}\to A_x$ déduit canoniquement de $\varphi$
(\hyperref[0.1.5.1-fr]{0, 1.5.1}). Comme la donnée des $\theta_x^{\#}$
caractérise complètement $\theta^{\#}$, et par suite aussi $\theta$
(\hyperref[0.3.7.1-fr]{0, 3.7.1}), on en conclut que l'on a
$\theta=\widetilde\varphi$, par définition de $\widetilde\varphi$
(\hyperref[I.1.6.1-fr]{1.6.1}).

Nous dirons qu'un morphisme $(\psi,\theta)$ d'espaces annelés satisfaisant à
la condition de (\hyperref[I.1.7.3-fr]{1.7.3}) est un \emph{morphisme de
schémas affines}.

\begin{corollary}[1.7.4]
\label{I.1.7.4-fr}
Si $(X,\mathscr{O}_X)$, $(Y,\mathscr{O}_Y)$ sont deux schémas affines, il
existe un isomorphisme canonique de l'ensemble de morphismes de schémas
affines
$\operatorname{Hom}((X,\mathscr{O}_X),(Y,\mathscr{O}_Y))$ sur l'ensemble
d'homomorphismes d'anneaux de $B$ dans $A$, où
$A=\Gamma(\mathscr{O}_X)$ et $B=\Gamma(\mathscr{O}_Y)$.
\end{corollary}

On peut encore dire que les foncteurs $(\operatorname{Spec}(A),\widetilde A)$
en $A$ et $\Gamma(X,\mathscr{O}_X)$ en $(X,\mathscr{O}_X)$ définissent une
\emph{équivalence} de la catégorie des anneaux commutatifs et de la catégorie
duale de la catégorie des schémas affines (T, I, 1.2).

\begin{corollary}[1.7.5]
\label{I.1.7.5-fr}
Si $\varphi:B\to A$ est surjectif, le morphisme correspondant
$({}^a\varphi,\widetilde\varphi)$ est un monomorphisme d'espaces annelés
\emph{(cf. (\hyperref[I.4.1.7-fr]{4.1.7}))}.
\end{corollary}

On sait en effet que ${}^a\varphi$ est injective
(\hyperref[I.1.2.5-fr]{1.2.5}), et comme $\varphi$ est surjectif, pour tout
$x\in X$, $\varphi_x^{\#}:B_{{}^a\varphi(x)}\to A_x$, qui se déduit de
$\varphi$ par passage aux anneaux de fractions, est aussi surjectif
(\hyperref[0.1.5.1-fr]{0, 1.5.1})~; d'où la conclusion
(\hyperref[0.4.1.1-fr]{0, 4.1.1}).

% FIN SOURCE ega1/ega1-1-fr.tex

%
% DEBUT SOURCE ega1/ega1-2-fr.tex
\section{Préschémas et morphismes de préschémas}
\label{section:I.2-fr}

\subsection{Définition des préschémas.}
\label{subsection:I.2.1-fr}

\begin{env}[2.1.1]
\label{I.2.1.1-fr}
Étant donné un espace annelé $(X,\mathscr{O}_X)$, on dit qu'une partie ouverte
$V$ de $X$ est un \emph{ouvert affine} si l'espace annelé
$(V,\mathscr{O}_X|V)$ est un schéma affine
(\hyperref[I.1.7.1-fr]{1.7.1}).
\end{env}

\begin{definition}[2.1.2]
\label{I.2.1.2-fr}
On appelle \emph{préschéma} un espace annelé $(X,\mathscr{O}_X)$ tel que tout
point de $X$ admette un voisinage ouvert affine.
\end{definition}

\oldpage[I]{98}

\begin{proposition}[2.1.3]
\label{I.2.1.3-fr}
Si $(X,\mathscr{O}_X)$ est un préschéma, les ensembles ouverts affines forment
une base de la topologie de $X$.
\end{proposition}

En effet, si $V$ est un voisinage ouvert quelconque de $x\in X$, il existe par
hypothèse un voisinage ouvert $W$ de $x$ tel que
$(W,\mathscr{O}_X|W)$ soit un schéma affine~; désignons par $A$ son anneau.
Dans l'espace $W$, $V\cap W$ est un voisinage ouvert de $x$~; donc il existe
$f\in A$ tel que $D(f)$ soit un voisinage ouvert de $x$ contenu dans
$V\cap W$ (\hyperref[I.1.1.10-fr]{1.1.10 (i)}). L'espace annelé
$(D(f),\mathscr{O}_X|D(f))$ est alors un schéma affine d'anneau isomorphe à
$A_f$ (\hyperref[I.1.3.6-fr]{1.3.6}), d'où la proposition.

\begin{proposition}[2.1.4]
\label{I.2.1.4-fr}
L'espace sous-jacent d'un préschéma est un espace de Kolmogoroff.
\end{proposition}

En effet, si $x,y$ sont deux points distincts d'un préschéma $X$, il est
évident qu'il existe un voisinage ouvert de l'un de ces points ne contenant
pas l'autre si $x,y$ ne sont pas dans un même ouvert affine~; et s'ils sont
dans un même ouvert affine, cela résulte de
(\hyperref[I.1.1.8-fr]{1.1.8}).

\begin{proposition}[2.1.5]
\label{I.2.1.5-fr}
Si $(X,\mathscr{O}_X)$ est un préschéma, toute partie fermée irréductible de
$X$ admet un point générique et un seul, et l'application
$x\longmapsto\overline{\{x\}}$ est donc une bijection de $X$ sur l'ensemble
de ses parties fermées irréductibles.
\end{proposition}

En effet, si $Y$ est une partie fermée irréductible de $X$ et $y\in Y$, et si
$U$ est un voisinage ouvert affine de $y$ dans $X$, $U\cap Y$ est partout
dense dans $Y$ et est irréductible
(\hyperref[0.2.1.1-fr]{0, 2.1.1} et
\hyperref[0.2.1.4-fr]{2.1.4})~; donc
(\hyperref[I.1.1.14-fr]{1.1.14}), $U\cap Y$ est l'adhérence dans $U$ d'un
point $x$, et par suite $Y\subset\overline{U}$ est l'adhérence de $x$ dans
$X$. L'unicité du point générique de $X$ résulte de
(\hyperref[I.2.1.4-fr]{2.1.4}) et de
(\hyperref[0.2.1.3-fr]{0, 2.1.3}).

\begin{env}[2.1.6]
\label{I.2.1.6-fr}
Si $Y$ est une partie fermée irréductible de $X$, $y$ son point générique,
l'anneau local $\mathscr{O}_y$ se note aussi $\mathscr{O}_{X/Y}$ et s'appelle
l'\emph{anneau local de $X$ le long de $Y$}, ou l'\emph{anneau local de $Y$
dans $X$}.

Si $X$ lui-même est irréductible et si $x$ est son point générique, on dit
encore que $\mathscr{O}_x$ est l'\emph{anneau des fonctions rationnelles sur
$X$} (cf. \hyperref[section:I.7-fr]{§~7}).
\end{env}

\begin{proposition}[2.1.7]
\label{I.2.1.7-fr}
Si $(X,\mathscr{O}_X)$ est un préschéma, pour toute partie ouverte $U$ de $X$,
l'espace annelé $(U,\mathscr{O}_X|U)$ est un préschéma.
\end{proposition}

Cela résulte aussitôt de la déf. (\hyperref[I.2.1.2-fr]{2.1.2}) et de la prop.
(\hyperref[I.2.1.3-fr]{2.1.3}).

On dit que $(U,\mathscr{O}_X|U)$ est le préschéma \emph{induit} sur $U$ par
$(X,\mathscr{O}_X)$, ou la \emph{restriction} de $(X,\mathscr{O}_X)$ à $U$.

\begin{env}[2.1.8]
\label{I.2.1.8-fr}
On dit qu'un préschéma $(X,\mathscr{O}_X)$ est \emph{irréductible} (resp.
\emph{connexe}) si l'espace sous-jacent $X$ est irréductible (resp. connexe).
On dit qu'un préschéma est \emph{intègre} s'il est \emph{irréductible et
réduit} (cf. (\hyperref[I.5.1.4-fr]{5.1.4})). On dit qu'un préschéma
$(X,\mathscr{O}_X)$ est \emph{localement intègre} si tout $x\in X$ admet un
voisinage ouvert $U$ tel que le préschéma induit sur $U$ par
$(X,\mathscr{O}_X)$ soit intègre.
\end{env}

\subsection{Morphismes de préschémas.}
\label{subsection:I.2.2-fr}
\label{I.2.2-fr}

\begin{definition}[2.2.1]
\label{I.2.2.1-fr}
Étant donnés deux préschémas $(X,\mathscr{O}_X)$, $(Y,\mathscr{O}_Y)$, on
appelle \emph{morphisme (de préschémas)} de $(X,\mathscr{O}_X)$ dans
$(Y,\mathscr{O}_Y)$ tout morphisme d'espaces annelés $(\psi,\theta)$ tel que,
pour tout $x\in X$, $\theta_x^\sharp$ soit un homomorphisme local~:
$\mathscr{O}_{\psi(x)}\longrightarrow\mathscr{O}_x$.
\end{definition}

\oldpage[I]{99}

Par passage aux quotients, l'application
$\theta_x^\sharp:\mathscr{O}_{\psi(x)}\to\mathscr{O}_x$ donne donc un
monomorphisme $\theta^x:k(\psi(x))\to k(x)$, qui permet de considérer $k(x)$
comme une \emph{extension} du corps $k(\psi(x))$.

\begin{env}[2.2.2]
\label{I.2.2.2-fr}
Le composé $(\psi'',\theta'')$ de deux morphismes de préschémas
$(\psi,\theta)$ et $(\psi',\theta')$ est encore un morphisme de préschémas,
comme il résulte de la formule
${\theta''}^\sharp=\theta^\sharp\circ\psi^*({\theta'}^\sharp)$
(\hyperref[0.3.5.5-fr]{0, 3.5.5}). On en conclut que les préschémas forment
une \emph{catégorie}~; suivant la notation générale, on écrira
$\operatorname{Hom}(X,Y)$ l'ensemble des morphismes d'un préschéma $X$ dans
un préschéma $Y$.
\end{env}

\begin{example}[2.2.3]
\label{I.2.2.3-fr}
Si $U$ est une partie ouverte de $X$, l'injection canonique
(\hyperref[0.4.1.2-fr]{0, 4.1.2}) du préschéma induit
$(U,\mathscr{O}_X|U)$ dans $(X,\mathscr{O}_X)$ est un morphisme de
préschémas~; c'est d'ailleurs un \emph{monomorphisme} d'espaces annelés (et
\emph{a fortiori} un monomorphisme de préschémas), comme il résulte aussitôt
de (\hyperref[0.4.1.1-fr]{0, 4.1.1}).
\end{example}

\begin{proposition}[2.2.4]
\label{I.2.2.4-fr}
Soient $(X,\mathscr{O}_X)$ un préschéma, $(S,\mathscr{O}_S)$ un schéma affine
d'anneau $A$. Il existe une correspondance biunivoque canonique entre les
morphismes du préschéma $(X,\mathscr{O}_X)$ dans le préschéma
$(S,\mathscr{O}_S)$ et les homomorphismes de $A$ dans l'anneau
$\Gamma(X,\mathscr{O}_X)$.
\end{proposition}

Notons d'abord que, si $(X,\mathscr{O}_X)$ et $(Y,\mathscr{O}_Y)$ sont deux
espaces annelés quelconques, un morphisme $(\psi,\theta)$ de
$(X,\mathscr{O}_X)$ dans $(Y,\mathscr{O}_Y)$ définit canoniquement un
homomorphisme d'anneaux
$\Gamma(\theta):\Gamma(Y,\mathscr{O}_Y)\to
\Gamma(Y,\psi_*(\mathscr{O}_X))=\Gamma(X,\mathscr{O}_X)$. Dans le cas
considéré, tout revient à voir qu'un homomorphisme quelconque
$\varphi:A\to\Gamma(X,\mathscr{O}_X)$ est de la forme $\Gamma(\theta)$ pour
un $\theta$ et un seul. Or, il y a par hypothèse un recouvrement $(V_\alpha)$
de $X$ par des ouverts affines~; par composition de $\varphi$ avec
l'homomorphisme de restriction
$\Gamma(X,\mathscr{O}_X)\to
\Gamma(V_\alpha,\mathscr{O}_X|V_\alpha)$, on obtient un homomorphisme
$\varphi_\alpha:A\to\Gamma(V_\alpha,\mathscr{O}_X|V_\alpha)$, qui correspond
à un morphisme unique $(\psi_\alpha,\theta_\alpha)$ du préschéma
$(V_\alpha,\mathscr{O}_X|V_\alpha)$ dans $(S,\mathscr{O}_S)$, en vertu de
(\hyperref[I.1.7.3-fr]{1.7.3}). En outre, pour tout couple d'indices
$(\alpha,\beta)$, tout point de $V_\alpha\cap V_\beta$ admet un voisinage
ouvert affine $W$ contenu dans $V_\alpha\cap V_\beta$
(\hyperref[I.2.1.3-fr]{2.1.3})~; il est clair qu'en composant $\varphi_\alpha$
et $\varphi_\beta$ avec les homomorphismes de restriction à $W$, on obtient
le même homomorphisme
$\Gamma(S,\mathscr{O}_S)\to\Gamma(W,\mathscr{O}_X|W)$, donc, en vertu des
relations $(\theta_\alpha^\sharp)_x=(\varphi_\alpha)_x$ pour tout
$x\in V_\alpha$ et tout $\alpha$ (\hyperref[I.1.6.1-fr]{1.6.1}), les
restrictions à $W$ des morphismes $(\psi_\alpha,\theta_\alpha)$ et
$(\psi_\beta,\theta_\beta)$ coïncident. On en conclut qu'il y a un morphisme
d'espaces annelés
$(\psi,\theta):(X,\mathscr{O}_X)\to(S,\mathscr{O}_S)$ et un seul dont la
restriction à chaque $V_\alpha$ est $(\psi_\alpha,\theta_\alpha)$, et il est
clair que ce morphisme est un morphisme de préschémas et est tel que
$\Gamma(\theta)=\varphi$.

Soit $u:A\to\Gamma(X,\mathscr{O}_X)$ un homomorphisme d'anneaux, et soit
$v=(\psi,\theta)$ le morphisme $(X,\mathscr{O}_X)\to(S,\mathscr{O}_S)$
correspondant. Pour tout $f\in A$, on a
\begin{equation}
  \psi^{-1}(D(f))=X_{u(f)}
  \tag{2.2.4.1}\label{I.2.2.4.1-fr}
\end{equation}
avec les notations de (\hyperref[0.5.5.2-fr]{0, 5.5.2}) relatives au faisceau
localement libre $\mathscr{O}_X$. Il suffit en effet de vérifier cette formule
lorsque $X$ lui-même est affine, et alors elle n'est autre que
(\hyperref[I.1.2.2.2-fr]{1.2.2.2}).

\begin{proposition}[2.2.5]
\label{I.2.2.5-fr}
Sous les hypothèses de (\hyperref[I.2.2.4-fr]{2.2.4}), soient
$\varphi:A\to\Gamma(X,\mathscr{O}_X)$ un homomorphisme d'anneau,
$f:(X,\mathscr{O}_X)\to(S,\mathscr{O}_S)$ le morphisme de préschémas
correspondant, $\mathscr{G}$ (resp. $\mathscr{F}$) un
$\mathscr{O}_X$-Module (resp. un $\mathscr{O}_S$-Module) quasi-cohérent et
soit $M=\Gamma(S,\mathscr{F})$. Il existe alors une
\oldpage[I]{100}
correspondance biunivoque canonique entre les $f$-morphismes
$\mathscr{F}\to\mathscr{G}$ (\hyperref[0.4.4.1-fr]{0, 4.4.1}) et les
$A$-homomorphismes
$M\to(\Gamma(X,\mathscr{G}))_{[\varphi]}$.
\end{proposition}

En effet, en raisonnant comme dans (\hyperref[I.2.2.4-fr]{2.2.4}), on est
aussitôt ramené au cas où $X$ est affine et la proposition résulte alors de
(\hyperref[I.1.6.3-fr]{1.6.3}) et (\hyperref[I.1.3.8-fr]{1.3.8}).

\begin{env}[2.2.6]
\label{I.2.2.6-fr}
On dit qu'un morphisme de préschémas
$(\psi,\theta):(X,\mathscr{O}_X)\to(Y,\mathscr{O}_Y)$ est \emph{ouvert}
(resp. \emph{fermé}), si pour toute partie ouverte $U$ de $X$ (resp. toute
partie fermée $F$ de $X$), $\psi(U)$ est ouvert dans $Y$ (resp. $\psi(F)$
fermé dans $Y$). On dit que $(\psi,\theta)$ est \emph{dominant} si $\psi(X)$
est dense dans $Y$, \emph{surjectif} si $\psi$ est surjectif. On notera que
ces conditions ne font intervenir que l'application continue $\psi$.
\end{env}

\begin{proposition}[2.2.7]
\label{I.2.2.7-fr}
Soient
\[
  f=(\psi,\theta):(X,\mathscr{O}_X)\to(Y,\mathscr{O}_Y),
  \qquad
  g=(\psi',\theta'):(Y,\mathscr{O}_Y)\to(Z,\mathscr{O}_Z)
\]
deux morphismes de préschémas.
\begin{enumerate}
  \item[\rm(i)] Si $f$ et $g$ sont tous deux ouverts (resp. fermés, dominants,
    surjectifs), il en est de même de $g\circ f$.
  \item[\rm(ii)] Si $f$ est surjectif, et si $g\circ f$ est fermé, $g$ est
    fermé.
  \item[\rm(iii)] Si $g\circ f$ est surjectif, $g$ est surjectif.
\end{enumerate}
\end{proposition}

Les assertions (i) et (iii) sont évidentes. Posons
$g\circ f=(\psi'',\theta'')$. Si $F$ est fermé dans $Y$,
$\psi^{-1}(F)$ est fermé dans $X$, donc
$\psi''(\psi^{-1}(F))$ est fermé dans $Z$~; mais comme $\psi$ est surjectif,
$\psi(\psi^{-1}(F))=F$, donc
$\psi''(\psi^{-1}(F))=\psi'(F)$, ce qui démontre (ii).

\begin{proposition}[2.2.8]
\label{I.2.2.8-fr}
Soient $f=(\psi,\theta)$ un morphisme
$(X,\mathscr{O}_X)\to(Y,\mathscr{O}_Y)$, $(U_\alpha)$ un recouvrement ouvert
de $Y$. Pour que $f$ soit ouvert (resp. fermé, surjectif, dominant), il faut
et il suffit que sa restriction à chacun des préschémas induits
$(\psi^{-1}(U_\alpha),\mathscr{O}_X|\psi^{-1}(U_\alpha))$, considérée comme
un morphisme de ce préschéma induit dans le préschéma induit
$(U_\alpha,\mathscr{O}_Y|U_\alpha)$, soit ouvert (resp. fermé, surjectif,
dominant).
\end{proposition}

La proposition résulte aussitôt des définitions, tenant compte du fait qu'une
partie $F$ de $Y$ est fermée (resp. ouverte, dense) dans $Y$ si et seulement
si chacun des ensembles $F\cap U_\alpha$ est fermé (resp. ouvert, dense) dans
$U_\alpha$.

\begin{env}[2.2.9]
\label{I.2.2.9-fr}
Soient $(X,\mathscr{O}_X)$, $(Y,\mathscr{O}_Y)$ deux préschémas~; on suppose
que $X$ et $Y$ ont un même nombre \emph{fini} de composantes irréductibles
$X_i$ (resp. $Y_i$) $(1\leq i\leq n)$~; soit $\xi_i$ (resp. $\eta_i$) le
point générique de $X_i$ (resp. $Y_i$)
(\hyperref[I.2.1.5-fr]{2.1.5}). On dit qu'un morphisme
\[
  f=(\psi,\theta):(X,\mathscr{O}_X)\to(Y,\mathscr{O}_Y)
\]
est \emph{birationnel} si, pour tout $i$, $\psi^{-1}(\eta_i)=\{\xi_i\}$ et
$\theta_{\xi_i}^{\sharp}:\mathscr{O}_{\eta_i}\to\mathscr{O}_{\xi_i}$ est un
\emph{isomorphisme}. Il est clair qu'un morphisme birationnel est
\emph{dominant} (\hyperref[0.2.1.8-fr]{0, 2.1.8}), donc surjectif s'il est
aussi fermé.
\end{env}

\begin{notation}[2.2.10]
\label{I.2.2.10-fr}
Dans toute la suite de cet ouvrage et lorsque cela ne risquera pas de créer
des confusions, nous \emph{supprimerons} dans la notation d'un préschéma
(resp. d'un morphisme) le faisceau structural (resp. le morphisme de faisceaux
structuraux). Si $U$ est une partie ouverte de l'espace sous-jacent d'un
préschéma $X$, lorsqu'on parlera de $U$ comme d'un préschéma, il s'agira
toujours du préschéma induit sur $U$.
\end{notation}

\subsection{Recollement de préschémas}
\label{subsection:I.2.3-fr}

\begin{env}[2.3.1]
\label{I.2.3.1-fr}
\oldpage[I]{101}
Il résulte de la définition (\hyperref[I.2.1.2-fr]{2.1.2}) que tout espace
annelé obtenu par \emph{recollement} de préschémas
(\hyperref[0.4.1.6-fr]{0, 4.1.6}) est encore un préschéma. En particulier,
puisque tout préschéma admet par définition un recouvrement formé d'ensembles
ouverts affines, on voit que tout préschéma peut s'obtenir par
\emph{recollement de schémas affines}.
\end{env}

\begin{example}[2.3.2]
\label{I.2.3.2-fr}
Soient $K$ un corps, $B=K[s]$, $C=K[t]$ deux anneaux de polynômes à une
indéterminée sur $K$, et soient $X_1=\operatorname{Spec}(B)$,
$X_2=\operatorname{Spec}(C)$, qui sont deux schémas affines isomorphes. Dans
$X_1$ (resp. $X_2$), soit $U_{12}$ (resp. $U_{21}$) l'ouvert affine $D(s)$
(resp. $D(t)$) dont l'anneau $B_s$ (resp. $C_t$) est formé des fractions
rationnelles de la forme $f(s)/s^m$ (resp. $g(t)/t^n$) avec $f\in B$
(resp. $g\in C$). Soit $u_{12}$ l'isomorphisme de préschémas
$U_{21}\to U_{12}$ correspondant (\hyperref[I.2.2.4-fr]{2.2.4}) à
l'isomorphisme de $B_s$ sur $C_t$ qui, à $f(s)/s^m$ fait correspondre la
fraction rationnelle $f(1/t)/(1/t^m)$. On peut recoller $X_1$ et $X_2$ le long
de $U_{12}$ et $U_{21}$ au moyen de $u_{12}$, car il n'y a évidemment pas de
condition de recollement. Nous retrouverons plus tard le préschéma $X$ ainsi
obtenu comme cas particulier d'une méthode de construction générale
(\hyperref[II.2.4.3-fr]{II, 2.4.3}). Montrons seulement ici que $X$ \emph{n'est
pas un schéma affine}~; cela résultera de ce que l'anneau
$\Gamma(X,\mathscr{O}_X)$ est \emph{isomorphe} à $K$, donc a un spectre réduit
à un point. En effet, une section de $\mathscr{O}_X$ au-dessus de $X$ a une
restriction au-dessus de $X_1$ (resp. $X_2$), identifié à un ouvert affine de
$X$, qui est un polynôme $f(s)$ (resp. $g(t)$), et il résulte des définitions
que l'on doit avoir $g(t)=f(1/t)$, ce qui n'est possible que si $f=g\in K$.
\end{example}

\subsection{Schémas locaux}
\label{subsection:I.2.4-fr}

\begin{env}[2.4.1]
\label{I.2.4.1-fr}
On appelle \emph{schéma local} un schéma affine dont l'anneau $A$ est local~;
il existe alors dans $X=\operatorname{Spec}(A)$ un seul \emph{point fermé}
$a$, et pour tout autre point $b\in X$, on a
$a\in\overline{\{b\}}$ (\hyperref[I.1.1.7-fr]{1.1.7}).
\end{env}

Pour tout préschéma $Y$ et tout point $y\in Y$, le schéma local
$\operatorname{Spec}(\mathscr{O}_y)$ est appelé le \emph{schéma local de $Y$
au point $y$}. Soient $V$ un ouvert affine de $Y$ contenant $y$, $B$ l'anneau
du schéma affine $V$~; $\mathscr{O}_y$ s'identifie canoniquement à $B_y$
(\hyperref[I.1.3.4-fr]{1.3.4}), et l'homomorphisme canonique $B\to B_y$
correspond par suite (\hyperref[I.1.6.1-fr]{1.6.1}) à un morphisme de
préschémas $\operatorname{Spec}(\mathscr{O}_y)\to V$. Si on compose ce
morphisme avec l'injection canonique $V\to Y$, on obtient donc un morphisme
$\operatorname{Spec}(\mathscr{O}_y)\to Y$, qui est \emph{indépendant} de
l'ouvert affine $V$ (contenant $y$) choisi~: en effet, si $V'$ est un second
ouvert affine contenant $y$, il existe un troisième ouvert affine $W$
contenant $y$ et tel que $W\subset V\cap V'$
(\hyperref[I.2.1.3-fr]{2.1.3})~; on peut donc se limiter au cas où
$V\subset V'$, et si $B'$ est l'anneau de $V'$, tout revient à remarquer que
le diagramme
\phantomsection\label{I.2.4.1.diagram-fr}
\[
  \xymatrix{
    B'\ar[rr]\ar[dr] & & B\ar[dl]\\
    & \mathscr{O}_y
  }
\]
est commutatif (\hyperref[0.1.5.1-fr]{0, 1.5.1}). Le morphisme
\[
  \operatorname{Spec}(\mathscr{O}_y)\to Y
\]
ainsi défini est dit \emph{canonique}.

\begin{proposition}[2.4.2]
\label{I.2.4.2-fr}
\oldpage[I]{102}
Soit $(Y,\mathscr{O}_Y)$ un préschéma~; pour tout $y\in Y$, soit
$(\psi,\theta)$ le morphisme canonique
$(\operatorname{Spec}(\mathscr{O}_y),\widetilde{\mathscr{O}}_y)
\to(Y,\mathscr{O}_Y)$. Alors $\psi$ est un homéomorphisme de
$\operatorname{Spec}(\mathscr{O}_y)$ sur le sous-espace $S_y$ de $Y$ formé
des $z$ tels que $y\in\overline{\{z\}}$ (autrement dit, des
\emph{générisations} de $y$ (\hyperref[0.2.1.2-fr]{0, 2.1.2}))~; en outre,
si $z=\psi(\mathfrak{p})$,
$\theta_z^{\sharp}:\mathscr{O}_z\to(\mathscr{O}_y)_{\mathfrak{p}}$ est un
isomorphisme~; $(\psi,\theta)$ est donc un monomorphisme d'espaces annelés.
\end{proposition}

Comme l'unique point fermé $a$ de $\operatorname{Spec}(\mathscr{O}_y)$ est
adhérent à tout point de cet espace, et que $\psi(a)=y$, l'image de
$\operatorname{Spec}(\mathscr{O}_y)$ par l'application continue $\psi$ est
contenue dans $S_y$. Comme $S_y$ est contenu dans tout ouvert affine
contenant $y$, on peut se ramener au cas où $Y$ est un schéma affine~; mais
dans ce cas la proposition résulte de
(\hyperref[I.1.6.2-fr]{1.6.2}).

\emph{On voit donc (\hyperref[I.2.1.5-fr]{2.1.5}) qu'il y a correspondance
biunivoque entre $\operatorname{Spec}(\mathscr{O}_y)$ et l'ensemble des
parties fermées irréductibles de $Y$ contenant $y$.}

\begin{corollary}[2.4.3]
\label{I.2.4.3-fr}
Pour que $y\in Y$ soit le point générique d'une composante irréductible de
$Y$, il faut et il suffit que le seul idéal premier de l'anneau local
$\mathscr{O}_y$ soit son idéal maximal (autrement dit, que $\mathscr{O}_y$
soit de dimension zéro).
\end{corollary}

\begin{proposition}[2.4.4]
\label{I.2.4.4-fr}
Soient $(X,\mathscr{O}_X)$ un schéma local d'anneau $A$, $a$ son unique point
fermé, $(Y,\mathscr{O}_Y)$ un préschéma. Tout morphisme
$u=(\psi,\theta):(X,\mathscr{O}_X)\to(Y,\mathscr{O}_Y)$ se factorise de façon
unique en
$X\to\operatorname{Spec}(\mathscr{O}_{\psi(a)})\to Y$, où la seconde flèche
désigne le morphisme canonique, et la première correspond à un homomorphisme
local $\mathscr{O}_{\psi(a)}\to A$. Cela établit une correspondance biunivoque
canonique entre l'ensemble des morphismes
$(X,\mathscr{O}_X)\to(Y,\mathscr{O}_Y)$ et l'ensemble des homomorphismes locaux
$\mathscr{O}_y\to A$ $(y\in Y)$.
\end{proposition}

En effet, pour tout $x\in X$, on a $a\in\overline{\{x\}}$, donc
$\psi(a)\in\overline{\{\psi(x)\}}$, ce qui prouve que $\psi(X)$ est contenu
dans tout ouvert affine contenant $\psi(a)$. On peut donc se ramener au cas où
$(Y,\mathscr{O}_Y)$ est un schéma affine d'anneau $B$, et on a alors
$u=({}^a\varphi,\widetilde{\varphi})$, où
$\varphi\in\operatorname{Hom}(B,A)$ (\hyperref[I.1.7.3-fr]{1.7.3}). En outre,
on a $\varphi^{-1}(\mathfrak{j}_a)=\mathfrak{j}_{\psi(a)}$, et par suite
l'image par $\varphi$ de tout élément de $B-\mathfrak{j}_{\psi(a)}$ est
inversible dans l'anneau local $A$~; la factorisation de l'énoncé résulte donc
de la propriété universelle des anneaux de fractions
(\hyperref[0.1.2.4-fr]{0, 1.2.4}). Inversement, à tout homomorphisme local
$\mathscr{O}_y\to A$ correspond un morphisme unique
$(\psi,\theta):X\to\operatorname{Spec}(\mathscr{O}_y)$ tel que $\psi(a)=y$
(\hyperref[I.1.7.3-fr]{1.7.3}), et en le composant avec le morphisme canonique
$\operatorname{Spec}(\mathscr{O}_y)\to Y$, on obtient un morphisme $X\to Y$,
ce qui achève de démontrer la proposition.

\begin{env}[2.4.5]
\label{I.2.4.5-fr}
Les schémas affines dont l'anneau est un \emph{corps} $K$ ont un espace
sous-jacent réduit à un point. Si $A$ est un anneau local d'idéal maximal
$\mathfrak{m}$, tout homomorphisme local $A\to K$ a un noyau égal à
$\mathfrak{m}$, donc se factorise en
$A\to A/\mathfrak{m}\to K$, où la seconde flèche est un monomorphisme. Les
morphismes $\operatorname{Spec}(K)\to\operatorname{Spec}(A)$ correspondent
donc biunivoquement aux monomorphismes de corps $A/\mathfrak{m}\to K$.
\end{env}

Soit $(Y,\mathscr{O}_Y)$ un préschéma~; pour tout $y\in Y$ et tout idéal
$\mathfrak{a}_y$ de $\mathscr{O}_y$, l'homomorphisme canonique
$\mathscr{O}_y\to\mathscr{O}_y/\mathfrak{a}_y$ définit un morphisme
$\operatorname{Spec}(\mathscr{O}_y/\mathfrak{a}_y)
\to\operatorname{Spec}(\mathscr{O}_y)$~; si on le compose avec le morphisme
canonique $\operatorname{Spec}(\mathscr{O}_y)\to Y$, on obtient un morphisme
$\operatorname{Spec}(\mathscr{O}_y/\mathfrak{a}_y)\to Y$, dit encore
\emph{canonique}. Pour $\mathfrak{a}_y=\mathfrak{m}_y$, c'est-à-dire
$\mathscr{O}_y/\mathfrak{a}_y=k(y)$, la prop.
(\hyperref[I.2.4.4-fr]{2.4.4}) entraîne donc~:

\begin{corollary}[2.4.6]
\label{I.2.4.6-fr}
\oldpage[I]{103}
Soient $(X,\mathscr{O}_X)$ un schéma local dont l'anneau est un corps $K$,
$\xi$ l'unique point de $X$, $(Y,\mathscr{O}_Y)$ un préschéma. Tout morphisme
$u:(X,\mathscr{O}_X)\to(Y,\mathscr{O}_Y)$ se factorise de façon unique en
$X\to\operatorname{Spec}(k(\psi(\xi)))\to Y$, où la seconde flèche désigne le
morphisme canonique, et la première correspond à un monomorphisme
$k(\psi(\xi))\to K$. Cela établit une correspondance biunivoque canonique
entre l'ensemble des morphismes $(X,\mathscr{O}_X)\to(Y,\mathscr{O}_Y)$ et
l'ensemble des monomorphismes $k(y)\to K$ $(y\in Y)$.
\end{corollary}

\begin{corollary}[2.4.7]
\label{I.2.4.7-fr}
Pour tout $y\in Y$, tout morphisme canonique
$\operatorname{Spec}(\mathscr{O}_y/\mathfrak{a}_y)\to Y$ est un monomorphisme
d'espaces annelés.
\end{corollary}

On l'a déjà vu lorsque $\mathfrak{a}_y=0$
(\hyperref[I.2.4.2-fr]{2.4.2}), et il suffit d'appliquer
(\hyperref[I.1.7.5-fr]{1.7.5}).

\begin{remark}[2.4.8]
\label{I.2.4.8-fr}
Soient $X$ un schéma local, $a$ son unique point fermé. Comme tout ouvert
affine contenant $a$ est nécessairement $X$ tout entier, tout
$\mathscr{O}_X$-Module \emph{inversible}
(\hyperref[0.5.4.1-fr]{0, 5.4.1}) est nécessairement \emph{isomorphe à
$\mathscr{O}_X$} (ou, comme on dit encore, est \emph{trivial}). Cette propriété
n'a pas lieu en général pour un schéma affine quelconque
$\operatorname{Spec}(A)$~; on verra au chap.~V que si $A$ est un anneau normal,
elle est vraie lorsque $A$ est \emph{factoriel}.
\end{remark}

\subsection{Préschémas au-dessus d'un préschéma}
\label{subsection:I.2.5-fr}

\begin{definition}[2.5.1]
\label{I.2.5.1-fr}
Étant donné un préschéma $S$, on dit que la donnée d'un préschéma $X$ et d'un
morphisme de préschémas $\varphi:X\to S$ définit un préschéma $X$ \emph{au-dessus
du préschéma $S$}, ou un \emph{$S$-préschéma}~; on dit que $S$ est le
\emph{préschéma de base} du $S$-préschéma $X$. Le morphisme $\varphi$ est appelé
le \emph{morphisme structural} du $S$-préschéma $X$. Lorsque $S$ est un schéma
affine d'anneau $A$, on dit aussi que $X$ muni de $\varphi$ est un préschéma
au-dessus de l'anneau $A$ (ou un $A$-préschéma).
\end{definition}

Il résulte de (\hyperref[I.2.2.4-fr]{2.2.4}) que la donnée d'un préschéma
au-dessus d'un anneau $A$ équivaut à la donnée d'un préschéma
$(X,\mathscr{O}_X)$ dont le faisceau structural $\mathscr{O}_X$ est un faisceau
de $A$-\emph{algèbres}. \emph{Un préschéma quelconque peut donc toujours être
considéré de façon unique comme un $\mathbf{Z}$-préschéma.}

Si $\varphi:X\to S$ est le morphisme structural d'un $S$-préschéma $X$, on dit
qu'un point $x\in X$ est \emph{au-dessus d'un point $s\in S$} si
$\varphi(x)=s$. On dit que $X$ \emph{domine} $S$ si $\varphi$ est un morphisme
dominant (\hyperref[I.2.2.6-fr]{2.2.6}).

\begin{env}[2.5.2]
\label{I.2.5.2-fr}
Soient $X$, $Y$ deux $S$-préschémas~; on dit qu'un morphisme de préschémas
$u:X\to Y$ est un \emph{morphisme de préschémas au-dessus de $S$} (ou un
$S$-\emph{morphisme}) si le diagramme
\phantomsection\label{I.2.5.2.diagram-fr}
\[
  \xymatrix{
    X\ar[rr]^u\ar[dr] & & Y\ar[dl]\\
    & S &
  }
\]
(où les flèches obliques sont les morphismes structuraux) est commutatif~:
cela entraîne que pour tout $s\in S$ et tout $x\in X$ au-dessus de $s$, $u(x)$
doit aussi être au-dessus de $s$.
\end{env}

Cette définition montre aussitôt que le composé de deux $S$-morphismes est un
$S$-morphisme~; les $S$-préschémas forment donc une \emph{catégorie}.

On notera $\operatorname{Hom}_S(X,Y)$ l'ensemble des $S$-morphismes d'un
$S$-préschéma $X$ dans un $S$-préschéma $Y$~; le morphisme identique d'un
$S$-préschéma $X$ sera désigné par $1_X$.

Lorsque $S$ est un schéma affine d'anneau $A$, on dira aussi
$A$-\emph{morphisme} au lieu de $S$-morphisme.

\begin{env}[2.5.3]
\label{I.2.5.3-fr}
\oldpage[I]{104}
Si $X$ est un $S$-préschéma, $v:X'\to X$ un morphisme de préschémas, le
morphisme composé $X'\xrightarrow{v}X\to S$ définit $X'$ comme
$S$-préschéma~; en particulier, tout préschéma induit sur un ouvert $U$ de
$X$ peut être considéré comme un $S$-préschéma au moyen de l'injection
canonique.

Si $u:X\to Y$ est un $S$-morphisme de $S$-préschémas, la restriction de $u$
à tout préschéma induit sur un ouvert $U$ de $X$ est donc un $S$-morphisme
$U\to Y$. Inversement, soit $(U_\alpha)$ un recouvrement ouvert de $X$ et pour
chaque $\alpha$, soit $u_\alpha:U_\alpha\to Y$ un $S$-morphisme~; si, pour tout
couple d'indices $(\alpha,\beta)$, les restrictions de $u_\alpha$ et de
$u_\beta$ à $U_\alpha\cap U_\beta$ coïncident, il existe un $S$-morphisme
$X\to Y$ et un seul dont la restriction à chaque $U_\alpha$ soit $u_\alpha$.

Si $u:X\to Y$ est un $S$-morphisme tel que $u(X)\subset V$, où $V$ est une
partie ouverte de $Y$, alors $u$, considéré comme morphisme de $X$ dans $V$,
est encore un $S$-morphisme.
\end{env}

\begin{env}[2.5.4]
\label{I.2.5.4-fr}
Soit $S'\to S$ un morphisme de préschémas~; pour tout $S'$-préschéma, le
morphisme composé $X\to S'\to S$ définit $X$ comme $S$-préschéma.
Inversement, supposons que $S'$ soit le préschéma induit sur un ouvert de
$S$~; soit $X$ un $S$-préschéma et supposons que le morphisme structural
$f:X\to S$ soit tel que l'on ait $f(X)\subset S'$~; alors on peut considérer
$X$ comme un $S'$-préschéma. Dans ce dernier cas, si $Y$ est un $S$-préschéma
dont le morphisme structural applique aussi l'espace sous-jacent dans $S'$,
tout $S$-morphisme de $X$ dans $Y$ est aussi un $S'$-morphisme.
\end{env}

\begin{env}[2.5.5]
\label{I.2.5.5-fr}
Si $X$ est un $S$-morphisme, $\varphi:X\to S$ le morphisme structural, on
appelle $S$-\emph{section de $X$} un $S$-morphisme de $S$ dans $X$,
c'est-à-dire un morphisme de préschémas $\psi:S\to X$ tel que
$\varphi\circ\psi$ soit l'identité de $S$. On désigne par $\Gamma(X/S)$
l'ensemble des $S$-sections de $X$.
\end{env}

% FIN SOURCE ega1/ega1-2-fr.tex

%
% DEBUT SOURCE ega1/ega1-3-fr.tex
\section{Produit de préschémas}
\label{section:I.3-fr}

\subsection{Somme de préschémas.}
\label{subsection:I.3.1-fr}

\phantomsection\label{I.3.1.1-fr}
Soit $(X_\alpha)$ une famille quelconque de préschémas~; soit $X$ un espace
topologique \emph{somme} des espaces sous-jacents $X_\alpha$~; $X$ est alors
réunion de sous-espaces ouverts deux à deux disjoints $X'_\alpha$, et pour
chaque $\alpha$ il y a un homéomorphisme $\varphi_\alpha$ de $X_\alpha$ sur
$X'_\alpha$. Si on munit chacun des $X'_\alpha$ du faisceau
$(\varphi_\alpha)_*(\mathscr{O}_{X_\alpha})$, il est clair que $X$ devient un
préschéma, qu'on appelle \emph{somme} de la famille de préschémas $(X_\alpha)$
et que l'on note $\bigsqcup_\alpha X_\alpha$. Si $Y$ est un préschéma,
l'application $f\mapsto(f\circ\varphi_\alpha)$ est une \emph{bijection} de
l'ensemble $\operatorname{Hom}(X,Y)$ sur l'ensemble produit
$\prod_\alpha\operatorname{Hom}(X_\alpha,Y)$. En particulier, si les
$X_\alpha$ sont des $S$-préschémas, de morphismes structuraux $\psi_\alpha$,
$X$ est un $S$-préschéma pour l'unique morphisme $\psi:X\to S$ tel que
$\psi\circ\varphi_\alpha=\psi_\alpha$ pour tout $\alpha$. La somme de deux
préschémas $X$, $Y$ se note $X\sqcup Y$. Il est immédiat que si
$X=\operatorname{Spec}(A)$, $Y=\operatorname{Spec}(B)$, $X\sqcup Y$
s'identifie canoniquement à $\operatorname{Spec}(A\times B)$.

\subsection{Produit de préschémas.}
\label{subsection:I.3.2-fr}

\begin{definition}[3.2.1]
\label{I.3.2.1-fr}
Étant donnés deux $S$-préschémas $X$, $Y$, on dit qu'un triplet
$(Z,p_1,p_2)$ formé d'un $S$-préschéma $Z$ et de deux $S$-morphismes
$p_1:Z\to X$, $p_2:Z\to Y$, est un produit des
\oldpage[I]{105}
$S$-préschémas $X$ et $Y$, si, pour tout $S$-préschéma $T$, l'application
$f\mapsto(p_1\circ f,p_2\circ f)$ est une bijection de l'ensemble des
$S$-morphismes de $T$ dans $Z$, sur l'ensemble des couples formés d'un
$S$-morphisme $T\to X$ et d'un $S$-morphisme $T\to Y$ (autrement dit, une
bijection
\[
  \operatorname{Hom}_S(T,Z)\xrightarrow{\sim}
  \operatorname{Hom}_S(T,X)\times\operatorname{Hom}_S(T,Y)).
\]
\end{definition}

Il s'agit donc là de la notion générale de \emph{produit} de deux objets d'une
catégorie, appliquée à la catégorie des $S$-préschémas (T, I, 1.1)~; en
particulier, un produit de deux $S$-préschémas est \emph{unique} à un
$S$-isomorphisme unique près. En raison de cette unicité, on désigne le plus
souvent un produit de deux $S$-préschémas $X$, $Y$ par la notation
$X\times_S Y$ (ou simplement $X\times Y$ si aucune confusion n'est à craindre),
les morphismes $p_1$, $p_2$ (appelés les \emph{projections canoniques} de
$X\times_S Y$ dans $X$ et $Y$ respectivement) étant supprimés de la notation.
Si $g:T\to X$, $h:T\to Y$ sont deux $S$-morphismes, on désignera par
$(g,h)_S$, ou simplement $(g,h)$ le $S$-morphisme
$f:T\to X\times_S Y$ tel que $p_1\circ f=g$, $p_2\circ f=h$. Si $X'$, $Y'$
sont deux $S$-préschémas, $p'_1$, $p'_2$ les projections canoniques de
$X'\times_S Y'$ (supposé exister), $u:X'\to X$, $v:Y'\to Y$ deux
$S$-morphismes, on écrira $u\times_S v$ (ou simplement $u\times v$) le
$S$-morphisme $(u\circ p'_1,v\circ p'_2)_S$ de $X'\times_S Y'$ dans
$X\times_S Y$.

Lorsque $S$ est un schéma affine d'anneau $A$, on remplace souvent $S$ par $A$
dans les notations précédentes.

\begin{proposition}[3.2.2]
\label{I.3.2.2-fr}
Soient $X$, $Y$, $S$ trois schémas affines, $B$, $C$, $A$ leurs anneaux
respectifs. Soient $Z=\operatorname{Spec}(B\otimes_A C)$, $p_1$, $p_2$ les
$S$-morphismes correspondant (\hyperref[I.2.2.4-fr]{2.2.4}) aux
$A$-homomorphismes canoniques $u:b\mapsto b\otimes1$ et
$v:c\mapsto1\otimes c$ de $B$ et $C$ dans $B\otimes_A C$~; alors
$(Z,p_1,p_2)$ est un produit de $X$ et $Y$.
\end{proposition}

En vertu de (\hyperref[I.2.2.4-fr]{2.2.4}), tout revient à vérifier que si, à
tout $A$-homomorphisme $f:B\otimes_A C\to L$ (où $L$ est une $A$-algèbre), on
associe le couple $(f\circ u,f\circ v)$, on définit une bijection
\[
  \operatorname{Hom}_A(B\otimes_A C,L)\xrightarrow{\sim}
  \operatorname{Hom}_A(B,L)\times\operatorname{Hom}_A(C,L)\footnotemark,
\]
\footnotetext{La notation $\operatorname{Hom}_A$ désigne ici l'ensemble des
homomorphismes des $A$-algèbres.} ce qui résulte immédiatement des définitions
et de la relation $b\otimes c=(b\otimes1)(1\otimes c)$.

\begin{corollary}[3.2.3]
\label{I.3.2.3-fr}
Soient $T$ un schéma affine d'anneau $D$,
$\alpha=({}^a\rho,\widetilde{\rho})$
(resp. $\beta=({}^a\sigma,\widetilde{\sigma})$) un $S$-morphisme $T\to X$
(resp. $T\to Y$), où $\rho$ (resp. $\sigma$) est un $A$-homomorphisme de $B$
(resp. $C$) dans $D$~; alors
$(\alpha,\beta)_S=({}^a\tau,\widetilde{\tau})$, où $\tau$ est l'homomorphisme
$B\otimes_A C\to D$ tel que $\tau(b\otimes c)=\rho(b)\sigma(c)$.
\end{corollary}

\begin{proposition}[3.2.4]
\label{I.3.2.4-fr}
Soient $f:S'\to S$ un monomorphisme de préschémas (T, I, 1.1), $X$, $Y$ deux
$S'$-préschémas, qui sont considérés aussi comme $S$-préschémas au moyen de
$f$. Tout produit des $S$-préschémas $X$, $Y$ est alors un produit des
$S'$-préschémas $X$, $Y$ et réciproquement.
\end{proposition}

Soient $\varphi:X\to S'$, $\psi:Y\to S'$ les morphismes structuraux. Si $T$
est un $S$-préschéma, $u:T\to X$, $v:T\to Y$ deux $S$-morphismes, on a par
définition $f\circ\varphi\circ u=f\circ\psi\circ v=\theta$, morphisme
structural de $T$~; l'hypothèse sur $f$ entraîne
$\varphi\circ u=\psi\circ v=\theta'$, et on voit qu'on peut considérer $T$
comme $S'$-préschéma de morphisme structural $\theta'$, $u$ et $v$ comme des
$S'$-morphismes. La conclusion de la proposition en résulte immédiatement,
compte tenu de (\hyperref[I.3.2.1-fr]{3.2.1}).

\begin{corollary}[3.2.5]
\label{I.3.2.5-fr}
Soient $X$, $Y$ deux $S$-préschémas, $\varphi:X\to S$, $\psi:Y\to S$ leurs
morphismes structuraux, $S'$ une partie ouverte de $S$ telle que
$\varphi(X)\subset S'$, $\psi(Y)\subset S'$. Tout produit des $S$-préschémas
$X$, $Y$ est aussi un produit des $S'$-préschémas $X$, $Y$, et réciproquement.
\end{corollary}

\oldpage[I]{106}
Il suffit d'appliquer (\hyperref[I.3.2.4-fr]{3.2.4}) à l'injection canonique
$S'\to S$.

\begin{theorem}[3.2.6]
\label{I.3.2.6-fr}
Étant donnés deux $S$-préschémas $X$, $Y$, il existe un produit
$X\times_S Y$.
\end{theorem}

Nous procéderons en plusieurs étapes.

\begin{lemma}[3.2.6.1]
\label{I.3.2.6.1-fr}
Soient $(Z,p,q)$ un produit de $X$ et $Y$, $U$, $V$ des parties ouvertes de
$X$, $Y$ respectivement. Si on pose
$W=p^{-1}(U)\cap q^{-1}(V)$, le triplet formé de $W$ et des restrictions de
$p$ et $q$ à $W$ \textup{(considérées comme des morphismes $W\to U$,
$W\to V$ respectivement)} est un produit de $U$ et $V$.
\end{lemma}

En effet, si $T$ est un $S$-préschéma, on peut identifier les $S$-morphismes
$T\to W$ et les $S$-morphismes $T\to Z$ appliquant $T$ dans $W$. Si alors
$g:T\to U$, $h:T\to V$ sont deux $S$-morphismes quelconques, on peut les
considérer comme des $S$-morphismes de $T$ dans $X$ et $Y$ respectivement, et
par hypothèse il y a donc un $S$-morphisme et un seul $f:T\to Z$ tel que
$g=p\circ f$, $h=q\circ f$. Comme $p(f(T))\subset U$, $q(f(T))\subset V$,
on a
\[
  f(T)\subset p^{-1}(U)\cap q^{-1}(V)=W,
\]
d'où notre assertion.

\begin{lemma}[3.2.6.2]
\label{I.3.2.6.2-fr}
Soient $Z$ un $S$-préschéma, $p:Z\to X$, $q:Z\to Y$ deux $S$-morphismes,
$(U_\alpha)$ un recouvrement ouvert de $X$, $(V_\lambda)$ un recouvrement
ouvert de $Y$. On suppose que pour tout couple $(\alpha,\lambda)$, le
$S$-préschéma
$W_{\alpha\lambda}=p^{-1}(U_\alpha)\cap q^{-1}(V_\lambda)$ et les
restrictions de $p$ et $q$ à $W_{\alpha\lambda}$ constituent un produit de
$U_\alpha$ et de $V_\lambda$. Alors $(Z,p,q)$ est un produit de $X$ et $Y$.
\end{lemma}

Montrons d'abord que, si $f_1$, $f_2$ sont deux $S$-morphismes $T\to Z$, les
relations $p\circ f_1=p\circ f_2$ et $q\circ f_1=q\circ f_2$ entraînent
$f_1=f_2$. En effet, $Z$ est réunion des $W_{\alpha\lambda}$, donc les
$f_1^{-1}(W_{\alpha\lambda})$ forment un recouvrement ouvert de $T$, et il
en est de même des $f_2^{-1}(W_{\alpha\lambda})$. En outre, on a
\[
\begin{aligned}
  f_1^{-1}(W_{\alpha\lambda})
  &=f_1^{-1}(p^{-1}(U_\alpha))\cap
    f_1^{-1}(q^{-1}(V_\lambda)) \\
  &=f_2^{-1}(p^{-1}(U_\alpha))\cap
    f_2^{-1}(q^{-1}(V_\lambda))
   =f_2^{-1}(W_{\alpha\lambda})
\end{aligned}
\]
par hypothèse, et tout revient à voir que les restrictions de $f_1$ et $f_2$
à $f_1^{-1}(W_{\alpha\lambda})=f_2^{-1}(W_{\alpha\lambda})$ sont
identiques pour tout couple d'indices. Mais comme ces restrictions peuvent
être considérées comme des $S$-morphismes de
$f_1^{-1}(W_{\alpha\lambda})$ dans $W_{\alpha\lambda}$, notre assertion
résulte des hypothèses et de la déf. (\hyperref[I.3.2.1-fr]{3.2.1}).

Supposons maintenant donnés deux $S$-morphismes $g:T\to X$, $h:T\to Y$.
Posons
$T_{\alpha\lambda}=g^{-1}(U_\alpha)\cap h^{-1}(V_\lambda)$~; les
$T_{\alpha\lambda}$ forment un recouvrement ouvert de $T$. Par hypothèse,
il existe un $S$-morphisme $f_{\alpha\lambda}$ tel que
$p\circ f_{\alpha\lambda}$ et $q\circ f_{\alpha\lambda}$ soient les
restrictions respectives de $g$ et $h$ à $T_{\alpha\lambda}$. En outre,
montrons que les restrictions de $f_{\alpha\lambda}$ et
$f_{\beta\mu}$ à
$T_{\alpha\lambda}\cap T_{\beta\mu}$ coïncident, ce qui achèvera de
prouver (\hyperref[I.3.2.6.2-fr]{3.2.6.2}). Or, les images de
$T_{\alpha\lambda}\cap T_{\beta\mu}$ par $f_{\alpha\lambda}$ et par
$f_{\beta\mu}$ sont contenues dans
$W_{\alpha\lambda}\cap W_{\beta\mu}$ par définition. Comme
\[
  W_{\alpha\lambda}\cap W_{\beta\mu}
  =p^{-1}(U_\alpha\cap U_\beta)\cap
   q^{-1}(V_\lambda\cap V_\mu),
\]
il résulte de (\hyperref[I.3.2.6.1-fr]{3.2.6.1}) que
$W_{\alpha\lambda}\cap W_{\beta\mu}$ et les restrictions à ce préschéma
de $p$ et $q$ constituent un produit de
$U_\alpha\cap U_\beta$ et de $V_\lambda\cap V_\mu$. Comme
$p\circ f_{\alpha\lambda}$ et $p\circ f_{\beta\mu}$ coïncident dans
$T_{\alpha\lambda}\cap T_{\beta\mu}$ et qu'il en est de même de
$q\circ f_{\alpha\lambda}$ et $q\circ f_{\beta\mu}$, on voit que
$f_{\alpha\lambda}$ et $f_{\beta\mu}$ coïncident dans
$T_{\alpha\lambda}\cap T_{\beta\mu}$, c.q.f.d.

\oldpage[I]{107}
\begin{lemma}[3.2.6.3]
\label{I.3.2.6.3-fr}
Soient $(U_\alpha)$ un recouvrement ouvert de $X$, $(V_\lambda)$ un
recouvrement ouvert de $Y$, et supposons que pour tout couple
$(\alpha,\lambda)$, il existe un produit de $U_\alpha$ et de $V_\lambda$~;
alors il existe un produit de $X$ et $Y$.
\end{lemma}

Appliquant le lemme (\hyperref[I.3.2.6.1-fr]{3.2.6.1}) aux ouverts
$U_\alpha\cap U_\beta$ et $V_\lambda\cap V_\mu$, on voit qu'il existe un
produit des $S$-préschémas induits respectivement par $X$ et $Y$ sur ces
ouverts~; en outre, l'unicité du produit montre que, si on pose
$i=(\alpha,\lambda)$, $j=(\beta,\mu)$, il y a un isomorphisme canonique
$h_{ij}$ (resp. $h_{ji}$) de ce produit sur un $S$-préschéma $W_{ij}$
(resp. $W_{ji}$) induit par $U_\alpha\times_S V_\lambda$
(resp. $U_\beta\times_S V_\mu$) sur un ouvert~;
$f_{ij}=h_{ij}\circ h_{ji}^{-1}$ est donc un isomorphisme de $W_{ji}$ sur
$W_{ij}$. En outre, pour un troisième couple $k=(\gamma,\nu)$, on a
$f_{ik}=f_{ij}\circ f_{jk}$ dans $W_{ki}\cap W_{kj}$, comme il résulte de
l'application de (\hyperref[I.3.2.6.1-fr]{3.2.6.1}) aux ouverts
$U_\alpha\cap U_\beta\cap U_\gamma$ et
$V_\lambda\cap V_\mu\cap V_\nu$, dans $U_\beta$ et $V_\mu$
respectivement. Il y a par suite un préschéma $Z$, un recouvrement ouvert
$(Z_i)$ de l'espace sous-jacent à $Z$ et pour chaque $i$ un isomorphisme $g_i$
du préschéma induit $Z_i$ sur le préschéma
$U_\alpha\times_S V_\lambda$, de sorte que pour tout couple $(i,j)$, on ait
$f_{ij}=g_i\circ g_j^{-1}$
(\hyperref[I.2.3.1-fr]{2.3.1})~; de plus, on a
$g_i(Z_i\cap Z_j)=W_{ij}$. Si $p_i$, $q_i$, $\theta_i$ sont les projections
et le morphisme structural du $S$-préschéma
$U_\alpha\times_S V_\lambda$, on constate aussitôt que
$p_i\circ g_i=p_j\circ g_j$ dans $Z_i\cap Z_j$, et de même pour les deux
autres morphismes. On peut donc définir des morphismes de préschémas
$p:Z\to X$ (resp. $q:Z\to Y$, $\theta:Z\to S$) par la condition que $p$
(resp. $q$, $\theta$) coïncide avec $p_i\circ g_i$
(resp. $q_i\circ g_i$, $\theta_i\circ g_i$) dans chacun des $Z_i$~; $Z$,
muni de $\theta$, est alors un $S$-préschéma. Montrons maintenant que
$Z'_i=p^{-1}(U_\alpha)\cap q^{-1}(V_\lambda)$ est égal à $Z_i$. Pour tout
indice $j=(\beta,\mu)$, on a
$Z_j\cap Z'_i=g_j^{-1}(p_j^{-1}(U_\alpha)\cap q_j^{-1}(V_\lambda))$. Or,
\[
  p_j^{-1}(U_\alpha)\cap q_j^{-1}(V_\lambda)
  =p_j^{-1}(U_\alpha\cap U_\beta)\cap
   q_j^{-1}(V_\lambda\cap V_\mu)~;
\]
en vertu de (\hyperref[I.3.2.6.1-fr]{3.2.6.1}), les restrictions de $p_j$ et
$q_j$ à
$p_j^{-1}(U_\alpha)\cap q_j^{-1}(V_\lambda)$ définissent sur ce
$S$-préschéma une structure de produit de $U_\alpha\cap U_\beta$ et
$V_\lambda\cap V_\mu$~; mais l'unicité du produit entraîne alors que
$p_j^{-1}(U_\alpha)\cap q_j^{-1}(V_\lambda)=W_{ji}$. On a par suite
$Z_j\cap Z'_i=Z_j\cap Z_i$ pour tout $j$, d'où $Z'_i=Z_i$. On déduit alors
de (\hyperref[I.3.2.6.2-fr]{3.2.6.2}) que $(Z,p,q)$ est un produit de $X$ et
$Y$.

\begin{lemma}[3.2.6.4]
\label{I.3.2.6.4-fr}
Soient $\varphi:X\to S$, $\psi:Y\to S$ les morphismes structuraux de $X$
et $Y$, $(S_i)$ un recouvrement ouvert de $S$, et posons
$X_i=\varphi^{-1}(S_i)$, $Y_i=\psi^{-1}(S_i)$. Si chacun des produits
$X_i\times_S Y_i$ existe, alors $X\times_S Y$ existe.
\end{lemma}

D'après (\hyperref[I.3.2.6.3-fr]{3.2.6.3}), tout revient à prouver que les
produits $X_i\times_S Y_j$ existent quels que soient $i$ et $j$. Posons
$X_{ij}=X_i\cap X_j=\varphi^{-1}(S_i\cap S_j)$,
$Y_{ij}=Y_i\cap Y_j=\psi^{-1}(S_i\cap S_j)$~; en vertu de
(\hyperref[I.3.2.6.1-fr]{3.2.6.1}), le produit
$Z_{ij}=X_{ij}\times_S Y_{ij}$ existe. Notons maintenant que si $T$ est un
$S$-préschéma et si $g:T\to X_i$, $h:T\to Y_j$ sont des $S$-morphismes, on
a nécessairement
$\varphi(g(T))=\psi(h(T))\subset S_i\cap S_j$, d'après la définition d'un
$S$-morphisme, donc $g(T)\subset X_{ij}$ et $h(T)\subset Y_{ij}$~; il est
alors immédiat que $Z_{ij}$ est un produit de $X_i$ et $Y_j$.

\begin{env}[3.2.6.5]
\label{I.3.2.6.5-fr}
Nous pouvons maintenant achever de démontrer le th.
(\hyperref[I.3.2.6-fr]{3.2.6}). Si $S$ est un \emph{schéma affine}, il y a
des recouvrements $(U_\alpha)$, $(V_\lambda)$ de $X$ et $Y$ respectivement,
formés d'ouverts affines~; comme $U_\alpha\times_S V_\lambda$ existe en
vertu de (\hyperref[I.3.2.2-fr]{3.2.2}), il en est de même de
$X\times_S Y$ par (\hyperref[I.3.2.6.3-fr]{3.2.6.3}). Si $S$ est un
préschéma quelconque, il y a un recouvrement $(S_i)$ de $S$ formé d'ouverts
affines. Si $\varphi:X\to S$, $\psi:Y\to S$ sont les morphismes
structuraux, et si on pose $X_i=\varphi^{-1}(S_i)$,
$Y_i=\psi^{-1}(S_i)$, les produits $X_i\times_{S_i}Y_i$ existent d'après ce
qui
\oldpage[I]{108}
précède~; mais alors les produits $X_i\times_S Y_i$ existent aussi
(\hyperref[I.3.2.5-fr]{3.2.5}), donc il en est de même de $X\times_S Y$ par
(\hyperref[I.3.2.6.4-fr]{3.2.6.4}).
\end{env}

\begin{corollary}[3.2.7]
\label{I.3.2.7-fr}
Soient $Z=X\times_S Y$ le produit de deux $S$-préschémas, $p$, $q$ les
projections de $Z$ dans $X$ et $Y$, $\varphi$ (resp. $\psi$) le morphisme
structural de $X$ (resp. $Y$). Soient $S'$ une partie ouverte de $S$, $U$
(resp. $V$) une partie ouverte de $X$ (resp. $Y$) contenue dans
$\varphi^{-1}(S')$ (resp. $\psi^{-1}(S')$). Alors le produit
$U\times_{S'}V$ s'identifie canoniquement au préschéma induit par $Z$ sur
$p^{-1}(U)\cap q^{-1}(V)$ (considéré comme $S'$-préschéma). En outre, si
$f:T\to X$, $g:T\to Y$ sont des $S$-morphismes tels que $f(T)\subset U$,
$g(T)\subset V$, le $S'$-morphisme $(f,g)_{S'}$ s'identifie à la restriction
de $(f,g)_S$ à $p^{-1}(U)\cap q^{-1}(V)$.
\end{corollary}

Cela résulte de (\hyperref[I.3.2.5-fr]{3.2.5}) et
(\hyperref[I.3.2.6.1-fr]{3.2.6.1}).

\begin{env}[3.2.8]
\label{I.3.2.8-fr}
Soient $(X_\alpha)$, $(Y_\lambda)$ deux familles de $S$-préschémas, $X$
(resp. $Y$) la somme de la famille $(X_\alpha)$ (resp. $(Y_\lambda)$)
(\hyperref[subsection:I.3.1-fr]{3.1}). Alors $X\times_S Y$ s'identifie à la
\emph{somme} de la famille $(X_\alpha\times_S Y_\lambda)$~; cela résulte
aussitôt de (\hyperref[I.3.2.6.3-fr]{3.2.6.3}).
\end{env}

\subsection{Propriétés formelles du produit~; changement de préschéma de base.}
\label{subsection:I.3.3-fr}

\begin{env}[3.3.1]
\label{I.3.3.1-fr}
Le lecteur remarquera que toutes les propriétés énoncées dans cette section,
sauf (\hyperref[I.3.3.13-fr]{3.3.13}) et
(\hyperref[I.3.3.15-fr]{3.3.15}), sont valables sans modification dans
\emph{toute catégorie}, chaque fois que les produits qui interviennent dans
les énoncés \emph{existent} (car il est clair que les notions de $S$-objet et
de $S$-morphisme peuvent se définir exactement comme dans
(\hyperref[subsection:I.2.5-fr]{2.5}) pour tout objet $S$ de la catégorie).
\end{env}

\begin{env}[3.3.2]
\label{I.3.3.2-fr}
En premier lieu, $X\times_S Y$ est un \emph{bifoncteur covariant} en $X$ et
$Y$ dans la catégorie des $S$-préschémas~: il suffit en effet de remarquer
que le diagramme
\[
  \xymatrix{
    X\times Y\ar[r]^{f\times 1}\ar[d] &
    X'\times Y\ar[r]^{f'\times 1}\ar[d] &
    X''\times Y\ar[d]\\
    X\ar[r]_f & X'\ar[r]_{f'} & X''
  }
\]
est commutatif.
\end{env}

\begin{proposition}[3.3.3]
\label{I.3.3.3-fr}
Pour tout $S$-préschéma $X$, la première (resp. seconde) projection de
$X\times_S S$ (resp. $S\times_S X$) est un isomorphisme fonctoriel de
$X\times_S S$ (resp. $S\times_S X$) sur $X$, dont l'isomorphisme réciproque
est $(1_X,\varphi)_S$ (resp. $(\varphi,1_X)_S$), en désignant par $\varphi$
le morphisme structural $X\to S$~; on peut donc écrire, à un isomorphisme
canonique près
\[
  X\times_S S=S\times_S X=X.
\]
\end{proposition}

Il suffit de prouver que le triplet $(X,1_X,\varphi)$ est un produit de $X$
et $S$. Or, si $T$ est un $S$-préschéma, le seul $S$-morphisme de $T$ dans
$S$ est nécessairement le morphisme structural $\psi:T\to S$. Si $f$ est
un $S$-morphisme de $T$ dans $X$, on a nécessairement
$\psi=\varphi\circ f$, d'où notre assertion.

\begin{corollary}[3.3.4]
\label{I.3.3.4-fr}
Soient $X$, $Y$ deux $S$-préschémas, $\varphi:X\to S$, $\psi:Y\to S$ leurs
morphismes structuraux. Si on identifie canoniquement $X$ à $X\times_S S$
et $Y$ à $S\times_S Y$, les projections $X\times_S Y\to X$ et
$X\times_S Y\to Y$ s'identifient respectivement à $1_X\times\psi$ et
$\varphi\times 1_Y$.
\end{corollary}

La vérification est immédiate et laissée au lecteur.

\begin{env}[3.3.5]
\label{I.3.3.5-fr}
On peut définir de la même manière que dans
(\hyperref[subsection:I.3.2-fr]{3.2}) le produit d'un
\oldpage[I]{109}
nombre fini quelconque $n$ de $S$-préschémas, l'existence de ces produits
résulte de (\hyperref[I.3.2.6-fr]{3.2.6}) par récurrence sur $n$, en remarquant
que
$(X_1\times_S X_2\times_S\cdots\times_S X_{n-1})\times_S X_n$ satisfait à
la définition du produit. L'unicité du produit entraîne, comme dans toute
catégorie, ses propriétés de \emph{commutativité} et d'\emph{associativité}.
Si, par exemple, $p_1$, $p_2$, $p_3$ désignent les projections de
$X_1\times_S X_2\times_S X_3$, et si on identifie ce préschéma à
$(X_1\times_S X_2)\times_S X_3$, la projection dans $X_1\times_S X_2$ est
identifiée à $(p_1,p_2)_S$.
\end{env}

\begin{env}[3.3.6]
\label{I.3.3.6-fr}
Soient $S$, $S'$ deux préschémas, $\varphi:S'\to S$ un morphisme, qui fait de
$S'$ un $S$-préschéma. Pour tout $S$-préschéma $X$, considérons le produit
$X\times_S S'$, et soient $p$ et $\pi'$ ses projections dans $X$ et $S'$,
respectivement. Muni de $\pi'$, ce produit est un $S'$-préschéma~; quand on le
considère comme tel, on le désigne par $X_{(S')}$ ou $X_{(\varphi)}$, et on dit
que c'est le préschéma obtenu par \emph{extension du préschéma de base} de $S$
à $S'$, au moyen du morphisme $\varphi$, ou l'\emph{image réciproque} de $X$
par $\varphi$. On notera que si $\pi$ est le morphisme structural de $X$,
$\theta$ le morphisme structural de $X\times_S S'$, considéré comme
$S$-préschéma, le diagramme
\phantomsection\label{I.3.3.6.base-change-diagram-fr}
\[
  \xymatrix{
    X\ar[d]_{\pi} &
    X_{(S')}\ar[l]_{p}\ar[ld]_{\theta}\ar[d]^{\pi'}\\
    S & S'\ar[l]^{\varphi}
  }
\]
est commutatif.
\end{env}

\begin{env}[3.3.7]
\label{I.3.3.7-fr}
Avec les notations de (\hyperref[I.3.3.6-fr]{3.3.6}), pour tout
$S$-morphisme $f:X\to Y$, on note encore $f_{(S')}$ le $S'$-morphisme
$f\times_S 1:X_{(S')}\to Y_{(S')}$, et on dit que $f_{(S')}$ est
l'\emph{image réciproque} du morphisme $f$ par $\varphi$. $X_{(S')}$ est donc
un \emph{foncteur covariant} en $X$, de la catégorie des $S$-préschémas dans
celle des $S'$-préschémas.
\end{env}

\begin{env}[3.3.8]
\label{I.3.3.8-fr}
Le préschéma $X_{(S')}$ peut encore être considéré comme solution d'un
\emph{problème d'application universelle}~: tout $S'$-préschéma $T$ est aussi
un $S$-préschéma au moyen de $\varphi$~; tout $S$-morphisme $g:T\to X$
s'écrit alors d'une seule manière $g=p\circ f$, où $f$ est un
$S'$-morphisme $T\to X_{(S')}$, comme il résulte de la définition du produit
appliquée aux $S$-morphismes $f$ et $\psi:T\to S'$ (morphisme structural de
$T$).
\end{env}

\begin{proposition}[3.3.9]
\label{I.3.3.9-fr}
\textup{(«~transitivité de l'extension des préschémas de base~»).} Soient
$S''$ un préschéma, $\varphi':S''\to S'$ un morphisme. Pour tout
$S$-préschéma $X$, il existe un isomorphisme canonique fonctoriel du
$S''$-préschéma $(X_{(\varphi)})_{(\varphi')}$ sur le $S''$-préschéma
$X_{(\varphi\circ\varphi')}$.
\end{proposition}

En effet, soient $T$ un $S''$-préschéma, $\psi$ son morphisme structural,
$g$ un $S$-morphisme de $T$ dans $X$ ($T$ étant considéré comme
$S$-préschéma de morphisme structural $\varphi\circ\varphi'\circ\psi$).
Comme $T$ est aussi un $S'$-préschéma de morphisme structural
$\varphi'\circ\psi$, on peut écrire $g=p\circ g'$, où $g'$ est un
$S'$-morphisme $T\to X_{(\varphi)}$, puis $g'=p'\circ g''$, où $g''$ est un
$S''$-morphisme $T\to(X_{(\varphi)})_{(\varphi')}$~:
\phantomsection\label{I.3.3.9.transitivity-diagram-fr}
\[
  \xymatrix{
    X\ar[d]_{\pi} &
    X_{(\varphi)}\ar[l]_{p}\ar[d]_{\pi'} &
    (X_{(\varphi)})_{(\varphi')}\ar[l]_{p'}\ar[d]_{\pi''}\\
    S & S'\ar[l]^{\varphi} & S''\ar[l]^{\varphi'}
  }
\]
\oldpage[I]{110}
D'où la proposition en raison de l'unicité de la solution d'un problème
d'application universelle.

Ce résultat s'exprime encore en écrivant l'égalité (à un isomorphisme
canonique près) $(X_{(S')})_{(S'')}=X_{(S'')}$, si aucune confusion n'est à
craindre, ou encore
\begin{equation}
  (X\times_S S')\times_{S'}S''=X\times_S S''~;
  \tag{3.3.9.1}\label{I.3.3.9.1-fr}
\end{equation}
le caractère fonctoriel de l'isomorphisme défini dans
(\hyperref[I.3.3.9-fr]{3.3.9}) s'exprime de même par la formule de
\emph{transitivité des images réciproques de morphismes}
\begin{equation}
  (f_{(S')})_{(S'')}=f_{(S'')}
  \tag{3.3.9.2}\label{I.3.3.9.2-fr}
\end{equation}
pour tout $S$-morphisme $f:X\to Y$.

\begin{corollary}[3.3.10]
\label{I.3.3.10-fr}
Si $X$ et $Y$ sont deux $S$-préschémas, il existe un isomorphisme canonique
fonctoriel du $S'$-préschéma $X_{(S')}\times_{S'}Y_{(S')}$ sur le
$S'$-préschéma $(X\times_S Y)_{(S')}$.
\end{corollary}

En effet, on a, à des isomorphismes canoniques près
\[
  (X\times_S S')\times_{S'}(Y\times_S S')
  =X\times_S(Y\times_S S')=(X\times_S Y)\times_S S'
\]
en vertu de (\hyperref[I.3.3.9.1-fr]{3.3.9.1}) et de l'associativité des
produits de $S$-préschémas.

Le caractère fonctoriel de l'isomorphisme défini dans
(\hyperref[I.3.3.10-fr]{3.3.10}) s'exprime par la formule
\begin{equation}
  (u_{(S')},v_{(S')})_{S'}=((u,v)_S)_{(S')}
  \tag{3.3.10.1}\label{I.3.3.10.1-fr}
\end{equation}
pour tout couple de $S$-morphismes $u:T\to X$, $v:T\to Y$.

En d'autres termes, le foncteur image réciproque $X_{(S')}$
\emph{commute à la formation des produits}~; il commute aussi à la formation
des sommes (\hyperref[I.3.2.8-fr]{3.2.8}).

\begin{corollary}[3.3.11]
\label{I.3.3.11-fr}
Soient $Y$ un $S$-préschéma, $f:X\to Y$ un morphisme qui fait de $X$ un
$Y$-préschéma (et par suite aussi un $S$-préschéma). Le préschéma $X_{(S')}$
s'identifie alors au produit $X\times_Y Y_{(S')}$, la projection
$X\times_Y Y_{(S')}\to Y_{(S')}$ s'identifiant à $f_{(S')}$.
\end{corollary}

Soit $\psi:Y\to S$ le morphisme structural de $Y$~; on a le diagramme
commutatif
\phantomsection\label{I.3.3.11.diagram-fr}
\[
  \xymatrix{
    S'\ar[d] &
    Y_{(S')}\ar[l]_{\psi_{(S')}}\ar[d] &
    X_{(S')}\ar[l]_{f_{(S')}}\ar[d]\\
    S & Y\ar[l]^{\psi} & X\ar[l]^{f}
  }
\]
Or $Y_{(S')}$ s'identifie à $S'_{(\psi)}$, et $X_{(S')}$ à
$S'_{(\psi\circ f)}$~; tenant compte de (\hyperref[I.3.3.9-fr]{3.3.9}) et de
(\hyperref[I.3.3.4-fr]{3.3.4}), on en déduit le corollaire.

\begin{env}[3.3.12]
\label{I.3.3.12-fr}
Soient $f:X\to X'$, $g:Y\to Y'$ deux $S$-morphismes qui sont des
\emph{monomorphismes} de préschémas (T, I, 1.1)~; alors $f\times_S g$ est un
\emph{monomorphisme}. En effet, si $p$ et $q$ sont les projections de
$X\times_S Y$, $p'$, $q'$ celles de $X'\times_S Y'$, $u$, $v$ deux
$S$-morphismes $T\to X\times_S Y$, la relation
$(f\times_S g)\circ u=(f\times_S g)\circ v$ entraîne
$p'\circ(f\times_S g)\circ u=p'\circ(f\times_S g)\circ v$, autrement dit
$f\circ p\circ u=f\circ p\circ v$, et comme $f$ est un monomorphisme,
$p\circ u=p\circ v$~; utilisant le fait que $g$ est un monomorphisme, on
obtient de même $q\circ u=q\circ v$, d'où $u=v$.
\oldpage[I]{111}
Il en résulte que pour toute extension $S'\to S$ du préschéma de base,
\phantomsection\label{I.3.3.12.base-change-monomorphism-fr}
\[
  f_{(S')}:X_{(S')}\to Y_{(S')}
\]
est un monomorphisme.
\end{env}

\begin{env}[3.3.13]
\label{I.3.3.13-fr}
Soient $S$, $S'$ deux schémas affines d'anneaux respectifs $A$, $A'$~; un
morphisme $S'\to S$ correspond donc à un homomorphisme d'anneaux $A\to A'$.
Si $X$ est un $S$-préschéma, on notera alors aussi $X_{(A')}$ ou
$X\otimes_A A'$ le $S'$-préschéma $X_{(S')}$~; lorsque $X$ est lui-même
affine d'anneau $B$, $X_{(A')}$ est affine d'anneau
$B_{(A')}=B\otimes_A A'$ obtenu par extension à $A'$ de l'anneau des scalaires
de la $A$-algèbre $B$.
\end{env}

\begin{env}[3.3.14]
\label{I.3.3.14-fr}
Avec les notations de (\hyperref[I.3.3.6-fr]{3.3.6}), pour tout
\emph{$S$-morphisme} $f:S'\to X$, $f'=(f,1_{S'})_S$ est un $S'$-morphisme
$S'\to X'=X_{(S')}$ tel que $p\circ f'=f$, $\pi'\circ f'=1_{S'}$, autrement
dit une \emph{$S'$-section de $X'$}~; et réciproquement si $f'$ est une telle
$S'$-section, $f=p\circ f'$ est un $S$-morphisme $S'\to X$. On définit donc
ainsi une \emph{correspondance biunivoque} canonique
\phantomsection\label{I.3.3.14.correspondence-fr}
\[
  \operatorname{Hom}_S(S',X)\overset{\sim}{\longrightarrow}
  \operatorname{Hom}_{S'}(S',X').
\]
On dit que $f'$ est le \emph{morphisme graphe} de $f$, et on le désigne par
$\Gamma_f$.
\end{env}

\begin{env}[3.3.15]
\label{I.3.3.15-fr}
Étant donné un préschéma $X$, qu'on peut toujours considérer comme un
$\mathbf{Z}$-préschéma, il résulte en particulier de
(\hyperref[I.3.3.14-fr]{3.3.14}) que les $X$-sections de
$X\otimes_{\mathbf{Z}}\mathbf{Z}[T]$ (où $T$ est une indéterminée)
correspondent biunivoquement aux \emph{morphismes}
$\mathbf{Z}[T]\to X$. Montrons que ces $X$-sections correspondent
biunivoquement aussi aux \emph{sections du faisceau structural
$\mathcal{O}_X$ au-dessus de $X$}. En effet, soit $(U_\alpha)$ un recouvrement
de $X$ par des ouverts affines~; soit
$u:X\to X\otimes_{\mathbf{Z}}\mathbf{Z}[T]$ un $X$-morphisme et soit
$u_\alpha$ sa restriction à $U_\alpha$~; si $A_\alpha$ est l'anneau du schéma
affine $U_\alpha$, $U_\alpha\otimes_{\mathbf{Z}}\mathbf{Z}[T]$ est un schéma
affine d'anneau $A_\alpha[T]$ (\hyperref[I.3.2.2-fr]{3.2.2}), et $u_\alpha$
correspond canoniquement à un $A_\alpha$-homomorphisme
$A_\alpha[T]\to A_\alpha$ (1.7.3). Or, un tel homomorphisme est complètement
déterminé par la donnée de l'image de $T$ dans $A_\alpha$, soit
$s_\alpha\in A_\alpha=\Gamma(U_\alpha,\mathcal{O}_X)$, et si on écrit que les
restrictions de $u_\alpha$ et de $u_\beta$ à un ouvert affine
$V\subset U_\alpha\cap U_\beta$ coïncident, on voit aussitôt que $s_\alpha$
et $s_\beta$ coïncident dans $V$~; donc la famille $(s_\alpha)$ est formée des
restrictions aux $U_\alpha$ d'une section $s$ de $\mathcal{O}_X$ au-dessus de
$X$~; et réciproquement, il est clair qu'une telle section définit une famille
$(u_\alpha)$ de morphismes qui sont les restrictions aux $U_\alpha$ d'un
$X$-morphisme $X\to X\otimes_{\mathbf{Z}}\mathbf{Z}[T]$. Ce résultat sera
généralisé dans II, 1.7.12.
\end{env}

\subsection{Points d'un préschéma à valeurs dans un préschéma~; points
géométriques}
\label{subsection:I.3.4-fr}

\begin{env}[3.4.1]
\label{I.3.4.1-fr}
Soit $X$ un préschéma~; pour tout préschéma $T$, on note encore $X(T)$
l'ensemble $\operatorname{Hom}(T,X)$ des morphismes $T\to X$, et les éléments de cet ensemble
sont aussi appelés \emph{points de $X$ à valeurs dans $T$}. Si on associe à
tout morphisme $f:T\to T'$ l'application $u'\mapsto u'\circ f$ de $X(T')$ dans
$X(T)$, on voit que, pour $X$ fixé, $X(T)$ est un \emph{foncteur contravariant
en $T$}, de la catégorie des préschémas dans celle des ensembles. En outre,
tout morphisme de préschémas $g:X\to Y$ définit un homomorphisme fonctoriel
$X(T)\to Y(T)$, faisant correspondre $g\circ v$ à $v\in X(T)$.
\end{env}

\begin{env}[3.4.2]
\label{I.3.4.2-fr}
Étant donnés trois ensembles $P$, $Q$, $R$ et deux applications
$\varphi:P\to R$, $\psi:Q\to R$, on appelle \emph{produit fibré de $P$ et $Q$
au-dessus de $R$} (relatif à $\varphi$ et $\psi$) la partie de
\oldpage[I]{112}
l'ensemble produit $P\times Q$ formé des couples $(p,q)$ tels que
$\varphi(p)=\psi(q)$~; on le note $P\times_R Q$. La définition
(\hyperref[I.3.2.1-fr]{3.2.1}) du produit de $S$-préschémas peut encore
s'interpréter, avec les notations de (\hyperref[I.3.4.1-fr]{3.4.1}), par la
formule
\phantomsection\label{I.3.4.2.1-fr}
\[
  (X\times_S Y)(T)=X(T)\times_{S(T)}Y(T)
  \tag{3.4.2.1}
\]
les applications $X(T)\to S(T)$ et $Y(T)\to S(T)$ correspondant aux morphismes
structuraux $X\to S$ et $Y\to S$.
\end{env}

\begin{env}[3.4.3]
\label{I.3.4.3-fr}
Si l'on se donne un préschéma $S$ et que l'on ne considère que les
$S$-préschémas et les $S$-morphismes, on notera encore $X(T)_S$ l'ensemble
$\operatorname{Hom}_S(T,X)$ des $S$-morphismes $T\to X$, en se permettant de
supprimer l'indice $S$ lorsque aucune confusion n'est possible~; on dit encore
que les éléments de $X(T)_S$ sont les \emph{points} (ou \emph{$S$-points}
lorsqu'il y a à craindre des confusions) \emph{du $S$-préschéma $X$ à valeurs
dans le $S$-préschéma $T$}. En particulier, une \emph{$S$-section} de $X$
n'est autre qu'un \emph{point de $X$ à valeurs dans $S$}. La formule
(\hyperref[I.3.4.2.1-fr]{3.4.2.1}) s'écrit alors aussi
\phantomsection\label{I.3.4.3.1-fr}
\[
  (X\times_S Y)(T)_S=X(T)_S\times Y(T)_S~;
  \tag{3.4.3.1}
\]
plus généralement, si $Z$ est un $S$-préschéma, $X$, $Y$, $T$ des
$Z$-préschémas (donc \emph{ipso facto} des $S$-préschémas), on a
\phantomsection\label{I.3.4.3.2-fr}
\[
  (X\times_Z Y)(T)_S=X(T)_S\times_{Z(T)_S}Y(T)_S.
  \tag{3.4.3.2}
\]

On notera que pour démontrer qu'un triplet $(W,r,s)$ formé d'un $S$-préschéma
$W$ et de deux $S$-morphismes $r:W\to X$, $s:W\to Y$ est un produit de $X$ et
$Y$ (au-dessus de $Z$), il suffit par définition de vérifier que pour
\emph{tout $S$-préschéma $T$}, le diagramme
\phantomsection\label{I.3.4.3.product-diagram-fr}
\[
  \xymatrix{
    W(T)_S\ar[r]^{r'}\ar[d]_{s'} & X(T)_S\ar[d]^{\varphi'}\\
    Y(T)_S\ar[r]_{\psi'} & Z(T)_S
  }
\]
fait de $W(T)_S$ le produit fibré de $X(T)_S$ et $Y(T)_S$ au-dessus de
$Z(T)_S$, $r'$ et $s'$ correspondant à $r$ et $s$, $\varphi'$ et $\psi'$ aux
morphismes structuraux $\varphi:X\to Z$, $\psi:Y\to Z$.
\end{env}

\begin{env}[3.4.4]
\label{I.3.4.4-fr}
Lorsque dans ce qui précède, $T$ (resp. $S$) est un schéma affine d'anneau $B$
(resp. $A$), on remplace $T$ (resp. $S$) par $B$ (resp. $A$) dans les notations
précédentes, et on parle alors de \emph{points de $X$ à valeurs dans l'anneau
$B$}, ou de \emph{points du $A$-préschéma $X$ à valeurs dans la $A$-algèbre
$B$} pour les éléments de $X(B)$ et de $X(B)_A$ respectivement. On notera que
$X(B)$ et $X(B)_A$ sont maintenant des foncteurs \emph{covariants en $B$}. On
écrira de même $X(T)_A$ pour l'ensemble des points du $A$-préschéma $X$ à
valeurs dans le $A$-préschéma $T$.
\end{env}

\begin{env}[3.4.5]
\label{I.3.4.5-fr}
Considérons en particulier le cas où $T$ est de la forme $\operatorname{Spec}(A)$,
où $A$ est un anneau \emph{local}~; les éléments de $X(A)$ correspondent alors
biunivoquement aux homomorphismes \emph{locaux} $\mathcal{O}_x\to A$ pour
$x\in X$ (\hyperref[I.2.4.4-fr]{2.4.4})~; on dit que le point $x$ de l'espace
sous-jacent à $X$ est la \emph{localité} du point de $X$ à valeurs dans $A$
auquel il correspond.

Plus particulièrement, on appellera \emph{points géométriques} d'un préschéma
$X$ les \emph{points de $X$ à valeurs dans un corps $K$}~: la donnée d'un tel
point revient donc à la donnée de sa
\oldpage[I]{113}
localité $x$ dans l'espace sous-jacent à $X$, et d'une \emph{extension} $K$ de
$k(x)$~; $K$ sera appelé le \emph{corps des valeurs} du point géométrique
correspondant, et on dit encore que ce point géométrique est \emph{localisé en
$x$}. On définit ainsi une application $X(K)\to X$, appliquant un point
géométrique à valeurs dans $K$ sur sa localité.

Si $S'=\operatorname{Spec}(K)$ est un $S$-préschéma (autrement dit, si $K$ est
considéré comme extension d'un corps résiduel $k(s)$, où $s\in S$) et si $X$
est un $S$-préschéma, un élément de $X(K)_S$, ou, comme on dit encore, un
\emph{point géométrique de $X$ au-dessus de $s$ à valeurs dans $K$}, consiste
en la donnée d'un $k(s)$-monomorphisme d'un corps résiduel $k(x)$ dans $K$, où
$x$ est un point de $X$ \emph{au-dessus de $s$} (donc $k(x)$ une extension de
$k(s)$).

Plus particulièrement, si $S=\operatorname{Spec}(K)=\{\xi\}$, \emph{les points
géométriques de $X$ à valeurs dans $K$ s'identifient aux points $x\in X$ tels
que $k(x)=K$}~; on dit encore que ces derniers sont les \emph{points du
$K$-préschéma $X$ rationnels sur $K$}~; si $K'$ est une extension de $K$, les
points géométriques de $X$ à valeurs dans $K'$ correspondent donc
biunivoquement aux points de $X'=X_{(K')}$ rationnels sur $K'$
(\hyperref[I.3.3.14-fr]{3.3.14}).
\end{env}

\begin{lemma}[3.4.6]
\label{I.3.4.6-fr}
Soient $X_i$ ($1\leq i\leq n$) des $S$-préschémas, $s$ un point de $S$, $x_i$
($1\leq i\leq n$) un point de $X_i$ au-dessus de $s$. Il existe alors une
extension $K$ de $k(s)$ et un point géométrique du produit
$Y=X_1\times_S X_2\times_S\cdots\times_S X_n$, à valeurs dans $K$, dont les
projections soient localisées aux $x_i$.
\end{lemma}

En effet, il existe des $k(s)$-monomorphismes $k(x_i)\to K$ dans une même
extension $K$ de $k(s)$ (Bourbaki, \emph{Alg.}, chap. V, \S~4, prop. 2). Les
composés $k(s)\to k(x_i)\to K$ sont tous identiques, donc les morphismes
$\operatorname{Spec}(K)\to X_i$ correspondant à $k(x_i)\to K$ sont des
$S$-morphismes, et on en conclut qu'ils définissent un morphisme unique
$\operatorname{Spec}(K)\to Y$. Si $y$ est le point correspondant de $Y$, il
est clair que sa projection dans chacun des $X_i$ est $x_i$.

\begin{proposition}[3.4.7]
\label{I.3.4.7-fr}
Soient $X_i$ ($1\leq i\leq n$) des $S$-préschémas, et pour chaque indice $i$,
soit $x_i$ un point de $X_i$. Pour qu'il existe un point $y$ de
$Y=X_1\times_S X_2\times_S\cdots\times_S X_n$ dont $x_i$ soit la projection
d'indice $i$ pour $1\leq i\leq n$, il faut et il suffit que les $x_i$ soient
au-dessus d'un même point $s$ de $S$.
\end{proposition}

La condition est évidemment nécessaire~; le lemme
(\hyperref[I.3.4.6-fr]{3.4.6}) prouve qu'elle est suffisante.

\phantomsection\label{I.3.4.7.underlying-map-fr}
En d'autres termes, si on désigne par $(X)$ l'ensemble sous-jacent à $X$, on
voit qu'on a une application \emph{surjective} canonique
$(X\times_S Y)\to (X)\times_{(S)}(Y)$~; il faut noter que cette application
\emph{n'est pas injective} en général~; autrement dit, \emph{il peut exister
plusieurs points $z$ distincts dans $X\times_S Y$ ayant mêmes projections
$x\in X$, $y\in Y$}~; c'est ce qu'on voit déjà lorsque $S$, $X$, $Y$ sont des
spectres premiers de corps $k$, $K$, $K'$, car le produit tensoriel
$K\otimes_k K'$ a en général plusieurs idéaux premiers distincts
(cf. \hyperref[I.3.4.9-fr]{3.4.9}).

\begin{corollary}[3.4.8]
\label{I.3.4.8-fr}
Soient $f:X\to Y$ un $S$-morphisme, $f_{(S')}:X_{(S')}\to Y_{(S')}$ le
$S'$-morphisme déduit de $f$ par une extension $S'\to S$ du préschéma de base.
Soit $p$ (resp. $q$) la projection $X_{(S')}\to X$ (resp.
$Y_{(S')}\to Y$)~; pour toute partie $M$ de $X$, on a
\phantomsection\label{I.3.4.8.base-change-image-fr}
\[
  q^{-1}(f(M))=f_{(S')}(p^{-1}(M)).
\]
\end{corollary}

\oldpage[I]{114}
En effet (\hyperref[I.3.3.11-fr]{3.3.11}), $X_{(S')}$ s'identifie au produit
$X\times_Y Y_{(S')}$ grâce au diagramme commutatif
\phantomsection\label{I.3.4.8.proof-diagram-fr}
\[
  \xymatrix{
    X\ar[d]_f & X_{(S')}\ar[l]_p\ar[d]^{f_{(S')}}\\
    Y & Y_{(S')}\ar[l]^q
  }
\]

En vertu de (\hyperref[I.3.4.7-fr]{3.4.7}), la relation $q(y')=f(x)$ pour
$x\in M$, $y'\in Y_{(S')}$ équivaut à l'existence d'un $x'\in X_{(S')}$ tel
que $p(x')=x$ et $f_{(S')}(x')=y'$, d'où le corollaire.

Le lemme (\hyperref[I.3.4.6-fr]{3.4.6}) se précise de la façon suivante~:

\begin{proposition}[3.4.9]
\label{I.3.4.9-fr}
Soient $X$, $Y$ deux $S$-préschémas, $x$ un point de $X$, $y$ un point de $Y$,
au-dessus du même point $s\in S$. L'ensemble des points de $X\times_S Y$ ayant
pour projections $x$ et $y$ est en correspondance biunivoque canonique avec
l'ensemble des types d'extensions composées de $k(x)$ et $k(y)$ considérés
comme extensions de $k(s)$ (Bourbaki, \emph{Alg.}, chap. VIII, \S~8, prop. 2).
\end{proposition}

Soit $p$ (resp. $q$) la projection de $X\times_S Y$ dans $X$ (resp. $Y$) et
soit $E$ le sous-espace $p^{-1}(x)\cap q^{-1}(y)$ de l'espace sous-jacent à
$X\times_S Y$. Notons d'abord que les morphismes
$\operatorname{Spec}(k(x))\to S$ et $\operatorname{Spec}(k(y))\to S$ se
factorisent en
$\operatorname{Spec}(k(x))\to\operatorname{Spec}(k(s))\to S$ et
$\operatorname{Spec}(k(y))\to\operatorname{Spec}(k(s))\to S$~; comme
$\operatorname{Spec}(k(s))\to S$ est un monomorphisme
(\hyperref[I.2.4.7-fr]{2.4.7}), il résulte de
(\hyperref[I.3.2.4-fr]{3.2.4}) que l'on a
\phantomsection\label{I.3.4.9.tensor-product-fibre-fr}
\[
  P=\operatorname{Spec}(k(x))\times_S\operatorname{Spec}(k(y))
  =\operatorname{Spec}(k(x))\times_{\operatorname{Spec}(k(s))}
   \operatorname{Spec}(k(y))
  =\operatorname{Spec}\bigl(k(x)\otimes_{k(s)}k(y)\bigr).
\]

Nous allons définir deux applications $\alpha:P_0\to E$, $\beta:E\to P_0$
réciproques l'une de l'autre ($P_0$ désignant l'ensemble sous-jacent au
préschéma $P$). Si $i:\operatorname{Spec}(k(x))\to X$ et
$j:\operatorname{Spec}(k(y))\to Y$ sont les morphismes canoniques
(\hyperref[I.2.4.5-fr]{2.4.5}), on prendra pour $\alpha$ l'application des
espaces sous-jacents correspondant au morphisme $i\times_S j$. D'autre part,
tout $z\in E$ définit par hypothèse deux $k(s)$-monomorphismes
$k(x)\to k(z)$ et $k(y)\to k(z)$, donc un $k(s)$-monomorphisme produit
tensoriel $k(x)\otimes_{k(s)}k(y)\to k(z)$ et par suite un morphisme
$\operatorname{Spec}(k(z))\to P$~; $\beta(z)$ sera l'image de $z$ dans $P_0$
par ce morphisme. La vérification du fait que $\alpha\circ\beta$ et
$\beta\circ\alpha$ sont les applications identiques découle de
(\hyperref[I.2.4.5-fr]{2.4.5}) et de la définition du produit
(\hyperref[I.3.2.1-fr]{3.2.1}). On sait enfin que $P_0$ est en correspondance
biunivoque avec l'ensemble des types d'extensions composées de $k(x)$ et
$k(y)$ (Bourbaki, \emph{Alg.}, chap. VIII, \S~8, prop. 1).

\subsection{Surjections et injections}
\label{subsection:I.3.5-fr}

\begin{env}[3.5.1]
\label{I.3.5.1-fr}
De façon générale, considérons une propriété $\mathbf{P}$ de morphismes de
préschémas, et les deux propositions suivantes~:
\begin{enumerate}
  \item[(i)] Si $f:X\to X'$, $g:Y\to Y'$ sont deux $S$-morphismes possédant
  la propriété $\mathbf{P}$, $f\times_S g$ possède la propriété $\mathbf{P}$.
  \item[(ii)] Si $f:X\to Y$ est un $S$-morphisme possédant la propriété
  $\mathbf{P}$, tout $S'$-morphisme $f_{(S')}:X_{(S')}\to Y_{(S')}$, déduit
  de $f$ par extension du préschéma de base, possède la propriété
  $\mathbf{P}$.
\end{enumerate}

\oldpage[I]{115}
Comme $f_{(S')}=f\times_S 1_{S'}$, on voit que si pour tout préschéma $X$
l'\emph{identité} $1_X$ possède la propriété $\mathbf{P}$, (i) implique
(ii)~; comme d'autre part, $f\times_S g$ est le morphisme composé
\phantomsection\label{I.3.5.1.product-composition-fr}
\[
  X\times_S Y\xrightarrow{f\times 1_Y}X'\times_S Y
  \xrightarrow{1_{X'}\times g}X'\times_S Y'
\]
on voit que, si le \emph{composé} de deux morphismes possédant la propriété
$\mathbf{P}$, possède aussi cette propriété, alors (ii) implique (i).
\end{env}

Une première application de cette remarque est la

\begin{proposition}[3.5.2]
\label{I.3.5.2-fr}
\begin{enumerate}
  \item[(i)] Si $f:X\to X'$, $g:Y\to Y'$ sont des $S$-morphismes surjectifs,
  $f\times_S g$ est surjectif.
  \item[(ii)] Si $f:X\to Y$ est un $S$-morphisme surjectif, $f_{(S')}$ est
  surjectif pour toute extension $S'$ du préschéma de base.
\end{enumerate}
\end{proposition}

Le composé de deux surjections étant une surjection, il suffit de démontrer
(ii)~; mais cette proposition résulte aussitôt de
(\hyperref[I.3.4.8-fr]{3.4.8}) appliqué à $M=X$.

\begin{proposition}[3.5.3]
\label{I.3.5.3-fr}
Pour qu'un morphisme $f:X\to Y$ soit surjectif, il faut et il suffit que pour
tout corps $K$ et tout morphisme $\operatorname{Spec}(K)\to Y$, il existe une
extension $K'$ de $K$ et un morphisme $\operatorname{Spec}(K')\to X$ rendant
commutatif le diagramme
\phantomsection\label{I.3.5.3.surjectivity-diagram-fr}
\[
  \xymatrix{
    X\ar[d]_f & \operatorname{Spec}(K')\ar[l]\ar[d]\\
    Y & \operatorname{Spec}(K)\ar[l]
  }
\]
\end{proposition}

La condition est suffisante, car pour tout $y\in Y$, il suffit de l'appliquer
à un morphisme $\operatorname{Spec}(K)\to Y$ correspondant à un monomorphisme
$k(y)\to K$, $K$ étant une extension de $k(y)$
(\hyperref[I.2.4.6-fr]{2.4.6}). Inversement, supposons $f$ surjectif, et soit
$y\in Y$ l'image de l'unique point de $\operatorname{Spec}(K)$~; il existe
$x\in X$ tel que $f(x)=y$~; considérons le monomorphisme correspondant
$k(y)\to k(x)$ (\hyperref[I.2.2.1-fr]{2.2.1})~; il suffit alors de prendre
$K'$ extension de $k(y)$ telle qu'il existe des $k(y)$-monomorphismes de
$k(x)$ et de $K$ dans $K'$ (Bourbaki, \emph{Alg.}, chap. V, \S~4, prop. 2)~;
le morphisme $\operatorname{Spec}(K')\to X$ correspondant à $k(x)\to K'$
répond à la question.

Avec le langage introduit dans (\hyperref[I.3.4.5-fr]{3.4.5}), on peut dire
encore que \emph{tout point géométrique de $Y$ à valeurs dans $K$ provient
d'un point géométrique de $X$ à valeurs dans une extension de $K$}.

\begin{definition}[3.5.4]
\label{I.3.5.4-fr}
On dit qu'un morphisme de préschémas $f:X\to Y$ est \emph{universellement
injectif}, ou un \emph{morphisme radiciel}, si pour tout corps $K$,
l'application correspondante $X(K)\to Y(K)$ est injective.
\end{definition}

Il résulte aussitôt des définitions que tout \emph{monomorphisme de
préschémas} (T, 1.1) est radiciel.

\begin{env}[3.5.5]
\label{I.3.5.5-fr}
Pour qu'un morphisme $f:X\to Y$ soit radiciel, il suffit que la condition de
la déf. (\hyperref[I.3.5.4-fr]{3.5.4}) soit vérifiée pour tout corps
\emph{algébriquement clos}. En effet, si $K$ est un corps quelconque, $K'$ une
extension algébriquement close de $K$, le diagramme
\phantomsection\label{I.3.5.5.algebraic-closure-diagram-fr}
\[
  \xymatrix{
    X(K)\ar[r]^\alpha\ar[d]_\varphi & Y(K)\ar[d]^{\varphi'}\\
    X(K')\ar[r]_{\alpha'} & Y(K')
  }
\]

\oldpage[I]{116}
est commutatif, $\varphi$ et $\varphi'$ provenant du morphisme
$\operatorname{Spec}(K')\to\operatorname{Spec}(K)$, $\alpha$ et $\alpha'$
correspondant à $f$. Or, $\varphi$ est injectif et il en est de même de
$\alpha'$ par hypothèse~; donc $\alpha$ est nécessairement injectif.
\end{env}

\begin{proposition}[3.5.6]
\label{I.3.5.6-fr}
Soient $f:X\to Y$, $g:Y\to Z$ deux morphismes de préschémas.
\begin{enumerate}
  \item[(i)] Si $f$ et $g$ sont radiciels, il en est de même de $g\circ f$.
  \item[(ii)] Inversement, si $g\circ f$ est radiciel, il en est de même de
  $f$.
\end{enumerate}
\end{proposition}

Compte tenu de la déf. (\hyperref[I.3.5.4-fr]{3.5.4}), la proposition revient
aux assertions correspondantes pour les applications
$X(K)\to Y(K)\to Z(K)$, qui sont évidentes.

\begin{proposition}[3.5.7]
\label{I.3.5.7-fr}
\begin{enumerate}
  \item[(i)] Si les $S$-morphismes $f:X\to X'$, $g:Y\to Y'$ sont radiciels,
  il en est de même de $f\times_S g$.
  \item[(ii)] Si le $S$-morphisme $f:X\to Y$ est radiciel, il en est de même
  de $f_{(S')}:X_{(S')}\to Y_{(S')}$ pour toute extension $S'\to S$ du
  préschéma de base.
\end{enumerate}
\end{proposition}

Vu (\hyperref[I.3.5.1-fr]{3.5.1}), il suffit de démontrer (i). On a vu
(\hyperref[I.3.4.2.1-fr]{3.4.2.1}) que
\phantomsection\label{I.3.5.7.field-valued-product-fr}
\[
  (X\times_S Y)(K)=X(K)\times_{S(K)}Y(K),\qquad
  (X'\times_S Y')(K)=X'(K)\times_{S(K)}Y'(K)
\]
l'application $(X\times_S Y)(K)\to(X'\times_S Y')(K)$ correspondant à
$f\times_S g$ s'identifie alors à $(u,v)\mapsto(f\circ u,g\circ v)$ et la
proposition en résulte aussitôt.

\begin{proposition}[3.5.8]
\label{I.3.5.8-fr}
Pour qu'un morphisme $f=(\psi,\theta):X\to Y$ soit radiciel, il faut et il
suffit que $\psi$ soit injectif et que, pour tout $x\in X$, le monomorphisme
$\theta^x:k(\psi(x))\to k(x)$ fasse de $k(x)$ une extension radicielle de
$k(\psi(x))$.
\end{proposition}

Supposons $f$ radiciel et montrons d'abord que la relation
$\psi(x_1)=\psi(x_2)=y$ entraîne nécessairement $x_1=x_2$. En effet, il
existe un corps $K$, extension de $k(y)$, et des $k(y)$-monomorphismes
$k(x_1)\to K$, $k(x_2)\to K$ (Bourbaki, \emph{Alg.}, chap. V, \S~4,
prop. 2)~; les morphismes correspondants
$u_1:\operatorname{Spec}(K)\to X$, $u_2:\operatorname{Spec}(K)\to X$ sont
donc tels que $f\circ u_1=f\circ u_2$, donc $u_1=u_2$ par hypothèse, et cela
implique en particulier $x_1=x_2$. Considérons maintenant $k(x)$ comme
extension de $k(\psi(x))$ au moyen de $\theta^x$~: si $k(x)$ n'est pas
extension radicielle de $k(\psi(x))$, il existe deux
$k(\psi(x))$-monomorphismes distincts de $k(x)$ dans une extension
algébriquement close $K$ de $k(\psi(x))$ et les deux morphismes correspondants
$\operatorname{Spec}(K)\to X$ violeraient l'hypothèse. Inversement, compte
tenu de (\hyperref[I.2.4.6-fr]{2.4.6}), il est immédiat que les conditions de
l'énoncé sont suffisantes pour que $f$ soit radiciel.

\begin{corollary}[3.5.9]
\label{I.3.5.9-fr}
Si $A$ est un anneau, $S$ une partie multiplicative de $A$, le morphisme
canonique $\operatorname{Spec}(S^{-1}A)\to\operatorname{Spec}(A)$ est
radiciel.
\end{corollary}

En effet, ce morphisme est un monomorphisme
(\hyperref[I.1.6.2-fr]{1.6.2}).

\begin{corollary}[3.5.10]
\label{I.3.5.10-fr}
Soient $f:X\to Y$ un morphisme radiciel, $g:Y'\to Y$ un morphisme, et soit
$X'=X_{(Y')}=X\times_Y Y'$. Alors le morphisme radiciel $f_{(Y')}$
(\hyperref[I.3.5.7-fr]{3.5.7 (ii)}) est une bijection de l'espace sous-jacent
$X'$ sur $g^{-1}(f(X))$~; en outre, pour tout corps $K$, l'ensemble $X'(K)$
s'identifie au sous-ensemble de $Y'(K)$ image réciproque par l'application
$Y'(K)\to Y(K)$ (correspondant à $g$) du sous-ensemble $X(K)$ de $Y(K)$.
\end{corollary}

La première assertion résulte aussitôt de
(\hyperref[I.3.5.8-fr]{3.5.8}) et de
(\hyperref[I.3.4.8-fr]{3.4.8})~; la seconde, de la commutativité du diagramme
\oldpage[I]{117}
\phantomsection\label{I.3.5.10.base-change-points-diagram-fr}
\[
  \xymatrix{
    X'(K)\ar[r]\ar[d] & Y'(K)\ar[d]\\
    X(K)\ar[r] & Y(K)
  }
\]

\begin{remark}[3.5.11]
\label{I.3.5.11-fr}
Nous dirons qu'un morphisme $f=(\psi,\theta)$ de préschémas est
\emph{injectif} si l'application $\psi$ est injective. Pour qu'un morphisme
$f=(\psi,\theta):X\to Y$ soit radiciel, il faut et il suffit que pour tout
morphisme $Y'\to Y$, le morphisme $f_{(Y')}:X_{(Y')}\to Y'$ soit injectif (ce
qui justifie la terminologie de morphisme \emph{universellement injectif}). En
effet, la condition est nécessaire en vertu de
(\hyperref[I.3.5.7-fr]{3.5.7, (ii)}) et de
(\hyperref[I.3.5.8-fr]{3.5.8}). Inversement, la condition implique d'abord que
$\psi$ est injectif~; si pour un $x\in X$, le monomorphisme
$\theta^x:k(\psi(x))\to k(x)$ n'était pas radiciel, il y aurait une extension
$K$ de $k(\psi(x))$ et deux morphismes distincts
$\operatorname{Spec}(K)\to X$ correspondant au même morphisme
$\operatorname{Spec}(K)\to Y$ (\hyperref[I.3.5.8-fr]{3.5.8}). Mais alors, en
posant $Y'=\operatorname{Spec}(K)$, il y aurait deux $Y'$-sections distinctes
de $X_{(Y')}$ (\hyperref[I.3.3.14-fr]{3.3.14}), ce qui contredit l'hypothèse
que $f_{(Y')}$ est injectif.
\end{remark}

\subsection{Fibres.}
\label{subsection:I.3.6-fr}

\begin{proposition}[3.6.1]
\label{I.3.6.1-fr}
Soient $f:X\to Y$ un morphisme, $y$ un point de $Y$, $\mathfrak{a}_y$ un
idéal de définition de $\mathcal{O}_y$ pour la topologie
$\mathfrak{m}_y$-préadique. La projection
$p:X\times_Y\operatorname{Spec}(\mathcal{O}_y/\mathfrak{a}_y)\to X$ est un
homéomorphisme de l'espace sous-jacent au préschéma
$X\times_Y\operatorname{Spec}(\mathcal{O}_y/\mathfrak{a}_y)$ sur la fibre
$f^{-1}(y)$ munie de la topologie induite par celle de l'espace sous-jacent à
$X$.
\end{proposition}

Comme $\operatorname{Spec}(\mathcal{O}_y/\mathfrak{a}_y)\to Y$ est radiciel
(\hyperref[I.3.5.4-fr]{3.5.4} et \hyperref[I.2.4.7-fr]{2.4.7}) et que
$\operatorname{Spec}(\mathcal{O}_y/\mathfrak{a}_y)$ est réduit à un seul
point, l'idéal $\mathfrak{m}_y/\mathfrak{a}_y$ étant nilpotent par hypothèse
(\hyperref[I.1.1.12-fr]{1.1.12}), on sait déjà
(\hyperref[I.3.5.10-fr]{3.5.10} et \hyperref[I.3.3.4-fr]{3.3.4}) que $p$
identifie en tant \emph{qu'ensemble} l'espace sous-jacent à
$X\times_Y\operatorname{Spec}(\mathcal{O}_y/\mathfrak{a}_y)$ à $f^{-1}(y)$~;
tout revient à démontrer que $p$ est un homéomorphisme. En vertu de
(\hyperref[I.3.2.7-fr]{3.2.7}), la question est locale sur $X$ et $Y$, et on
peut donc supposer que $X=\operatorname{Spec}(B)$,
$Y=\operatorname{Spec}(A)$, $B$ étant une $A$-algèbre. Le morphisme $p$
correspond alors à l'homomorphisme
$1\otimes\varphi:B\to B\otimes_A A'$, où $A'=A_y/\mathfrak{a}_y$ et
$\varphi$ est l'application canonique de $A$ dans $A'$. Or, tout élément de
$B\otimes_A A'$ s'écrit
\phantomsection\label{I.3.6.1.localization-fraction-fr}
\[
  \sum_i b_i\otimes\varphi(a_i)/\varphi(s)
  =\left(\sum_i(a_i b_i\otimes 1)\right)(1\otimes\varphi(s))^{-1},
\]
où $s\notin\mathfrak{j}_y$, et la prop.
(\hyperref[I.1.2.4-fr]{1.2.4}) est applicable.

\begin{env}[3.6.2]
\label{I.3.6.2-fr}
Dans toute la suite de ce Traité, lorsque nous considérerons une fibre
$f^{-1}(y)$ d'un morphisme comme munie d'une structure de $k(y)$-préschéma,
\emph{il s'agira toujours du préschéma obtenu en transportant la structure de
$X\times_Y\operatorname{Spec}(k(y))$ par la projection dans $X$}. Nous
écrirons aussi ce dernier produit $X\otimes_Y k(y)$, ou
$X\otimes_{\mathcal{O}_Y}k(y)$~; plus généralement, si $B$ est une
$\mathcal{O}_y$-algèbre, nous noterons $X\otimes_Y B$ ou
$X\otimes_{\mathcal{O}_Y}B$ le produit
$X\times_Y\operatorname{Spec}(B)$.
\end{env}

Avec la convention précédente, il résulte de
(\hyperref[I.3.5.10-fr]{3.5.10}) que les points de $X$ à valeurs dans une
extension $K$ de $k(y)$ sont identifiés aux \emph{points de $f^{-1}(y)$ à
valeurs dans $K$}.

\begin{env}[3.6.3]
\label{I.3.6.3-fr}
Soient $f:X\to Y$, $g:Y\to Z$ deux morphismes, $h=g\circ f$ leur composé~;
pour tout $z\in Z$, la fibre $h^{-1}(z)$ est un préschéma isomorphe à
\phantomsection\label{I.3.6.3.fibre-composition-fr}
\[
  X\times_Z\operatorname{Spec}(k(z))
  =(X\times_Y Y)\times_Z\operatorname{Spec}(k(z))
  =X\times_Y g^{-1}(z)
\]
\oldpage[I]{118}
En particulier, si $U$ est une partie ouverte de $X$, le préschéma induit sur
$U\cap f^{-1}(y)$ par le préschéma $f^{-1}(y)$ est isomorphe à
$f_U^{-1}(y)$ ($f_U$ étant la restriction de $f$ à $U$).
\end{env}

\begin{proposition}[3.6.4) («~transitivité des fibres~»]
\label{I.3.6.4-fr}
Soient $f:X\to Y$, $g:Y'\to Y$ deux morphismes~; posons
$X'=X\times_Y Y'=X_{(Y')}$, et
$f'=f_{(Y')}:X'\to Y'$. Pour tout $y'\in Y'$, si on pose $y=g(y')$, le
préschéma $f'^{-1}(y')$ est isomorphe à
$f^{-1}(y)\otimes_{k(y)}k(y')$.
\end{proposition}

En effet, cela revient à remarquer que les deux préschémas
$(X\otimes_Y k(y))\otimes_{k(y)}k(y')$ et
$(X\times_Y Y')\otimes_{Y'}k(y')$ sont tous deux canoniquement isomorphes à
$X\times_Y\operatorname{Spec}(k(y'))$ par
(\hyperref[I.3.3.9.1-fr]{3.3.9.1}).

En particulier, si $V$ est un voisinage ouvert de $y$ dans $Y$, et si on
désigne par $f_V$ la restriction de $f$ au préschéma induit sur $f^{-1}(V)$,
les préschémas $f^{-1}(y)$ et $f_V^{-1}(y)$ s'identifient canoniquement.

\begin{proposition}[3.6.5]
\label{I.3.6.5-fr}
Soient $f:X\to Y$ un morphisme, $y$ un point de $Y$, $Z$ le préschéma local
$\operatorname{Spec}(\mathcal{O}_y)$, $p=(\psi,\theta)$ la projection
$X\times_Y Z\to X$~; $p$ est un homéomorphisme de l'espace sous-jacent à
$X\times_Y Z$ sur le sous-espace $f^{-1}(Z)$ de $X$ (lorsque l'espace
sous-jacent à $Z$ est identifié à un sous-espace de $Y$, cf.
(\hyperref[I.2.4.2-fr]{2.4.2})), et pour tout $t\in X\times_Y Z$, si on pose
$z=\psi(t)$, $\theta_t^\sharp$ est un isomorphisme de $\mathcal{O}_z$ sur
$\mathcal{O}_t$.
\end{proposition}

Comme $Z$ (identifié à un sous-espace de $Y$) est contenu dans tout ouvert
affine contenant $y$ (\hyperref[I.2.4.2-fr]{2.4.2}), on peut, comme dans
(\hyperref[I.3.6.1-fr]{3.6.1}) se ramener au cas où
$X=\operatorname{Spec}(A)$ et $Y=\operatorname{Spec}(B)$ sont des schémas
affines, $A$ étant une $B$-algèbre. Alors $X\times_Y Z$ est le spectre premier
de $A\otimes_B B_y$ et cet anneau s'identifie canoniquement à $S^{-1}A$, où
$S$ est l'image de $B-\mathfrak{j}_y$ dans $A$ (0, 1.5.2)~; comme $p$
correspond alors à l'homomorphisme canonique $A\to S^{-1}A$, la proposition
résulte de (\hyperref[I.1.6.2-fr]{1.6.2}).

\subsection{Application : réduction d'un préschéma mod.
\texorpdfstring{$\mathfrak{J}$\protect\footnotemark}{J}}
\label{subsection:I.3.7-fr}
\label{I.3.7-fr}
\footnotetext{Ce numéro, qui fait usage de notions et résultats de la suite du
chap. I\textsuperscript{er} et du chap. II, ne sera pas utilisé par la suite,
et n'est destiné qu'aux lecteurs familiers avec la Géométrie algébrique
classique.}

\begin{env}[3.7.1]
\label{I.3.7.1-fr}
Soient $A$ un anneau, $X$ un $A$-préschéma, $\mathfrak{J}$ un idéal de $A$~;
alors $X_0=X\otimes_A(A/\mathfrak{J})$ est un
$(A/\mathfrak{J})$-préschéma, dont on dit parfois qu'il est déduit de $X$ par
\emph{réduction} mod. $\mathfrak{J}$.
\end{env}

\begin{env}[3.7.2]
\label{I.3.7.2-fr}
Cette terminologie est surtout utilisée lorsque $A$ est un \emph{anneau
local} et $\mathfrak{J}$ son idéal maximal, de sorte que $X_0$ est un
préschéma sur le corps résiduel $k=A/\mathfrak{J}$ de $A$.

Lorsqu'en outre $A$ est intègre, de corps des fractions $K$, on peut aussi
considérer le $K$-préschéma $X'=X\otimes_A K$. Par un abus de langage dont
nous ne ferons pas usage, on disait d'ordinaire jusqu'à présent que $X_0$ est
\emph{déduit de $X'$} par réduction mod. $\mathfrak{J}$. Dans les cas où ce
langage était utilisé, $A$ était un anneau local de dimension 1 (le plus
souvent un anneau de valuation discrète) et il était sous-entendu (de façon
plus ou moins explicite) que le $K$-préschéma $X'$ donné était un
sous-préschéma fermé d'un $K$-préschéma $P'$ (en fait, un espace projectif
type $\mathbf{P}_K^r$, cf. II, 4.1.1), lui-même de la forme
$P'=P\otimes_A K$, où $P$ est un $A$-préschéma donné (en fait, le $A$-schéma
$\mathbf{P}_A^r$, avec les notations de II, 4.1.1). Dans notre langage, la
définition de $X_0$ à partir de $X'$ se formule comme suit~:

On considère le schéma affine $Y=\operatorname{Spec}(A)$, formé de deux
points, l'unique point
\oldpage[I]{119}
fermé $y=\mathfrak{J}$ et le point générique $(0)$, l'ensemble $U$ réduit au
point générique étant donc un ouvert $U=\operatorname{Spec}(K)$ dans $Y$. Si
$X$ est un $A$-préschéma (autrement dit un $Y$-préschéma),
$X\otimes_A K=X'$ n'est autre que le préschéma induit par $X$ sur
$\psi^{-1}(U)$, en désignant par $\psi$ le morphisme structural $X\to Y$. En
particulier, si $\varphi$ est le morphisme structural $P\to Y$, un
sous-préschéma fermé $X'$ de $P'=\varphi^{-1}(U)$ est donc un sous-préschéma
(localement fermé) de $P$. Si $P$ est noethérien (par exemple si $A$ est
noethérien et $P$ de type fini sur $A$), il existe un plus petit
sous-préschéma fermé $X=\overline{X'}$ de $P$ qui majore $X$ (9.5.10), et
$X'$ est le préschéma induit par $X$ sur l'ouvert
$\varphi^{-1}(U)\cap X$, donc est isomorphe à $X\otimes_A K$ (9.5.10).
\emph{L'immersion de $X'$ dans $P'=P\otimes_A K$ permet donc de façon
canonique de considérer $X'$ comme étant de la forme $X'=X\otimes_A K$, où
$X$ est un $A$-préschéma.} On peut alors considérer le préschéma réduit mod.
$\mathfrak{J}$, $X_0=X\otimes_A k$, qui n'est autre d'ailleurs que la fibre
$\psi^{-1}(y)$ du point fermé $y$. Jusqu'à présent, faute d'une terminologie
adéquate, on avait évité d'introduire explicitement le $A$-préschéma $X$. Il
convient cependant de noter que toutes les assertions faites habituellement
sur le préschéma «~réduit mod. $\mathfrak{J}$~» $X_0$ doivent être regardées
comme conséquences d'assertions plus complètes concernant $X$ lui-même, et ne
peuvent être formulées et comprises de façon satisfaisante qu'en les
interprétant ainsi. Il semble d'ailleurs que les hypothèses faites reviennent
toujours à des hypothèses sur $X$ lui-même (indépendamment de la donnée
préalable d'une immersion de $X'$ dans un $\mathbf{P}_K^r$), ce qui permet de
donner des énoncés plus intrinsèques.
\end{env}

\begin{env}[3.7.3]
\label{I.3.7.3-fr}
Signalons enfin un fait très particulier, qui a sans doute contribué à
retarder la clarification conceptuelle de la situation envisagée ici~: si
$A$ est un anneau de valuation discrète et si $X$ est \emph{propre sur $A$}
(ce qui est en fait le cas si $X$ est un sous-préschéma fermé d'un
$\mathbf{P}_A^r$, cf. II, 5.5.4), les points de $X$ à valeurs dans $A$ et les
points de $X'$ à valeurs dans $K$ sont en correspondance biunivoque
(II, 7.3.8). C'est pourquoi on a souvent cru démontrer des résultats
concernant $X'$, alors qu'en réalité on démontrait des énoncés concernant
$X$, et qui restent valables (sous cette forme) lorsqu'on ne suppose plus
l'anneau local de base de dimension 1.
\end{env}

% FIN SOURCE ega1/ega1-3-fr.tex

%
% DEBUT SOURCE ega1/ega1-4-fr.tex
\section{Sous-préschémas et morphismes d'immersion}
\label{section:I.4-fr}

\subsection{Sous-préschémas.}
\label{subsection:I.4.1-fr}

\begin{env}[4.1.1]
\label{I.4.1.1-fr}
Comme la notion de faisceau quasi-cohérent
(\hyperref[0.5.1.3-fr]{0, 5.1.3}) est locale, un
$\mathscr{O}_X$-Module quasi-cohérent $\mathscr{F}$ sur un préschéma $X$ peut
être défini par la condition que, pour tout ouvert affine $V$ de $X$,
$\mathscr{F}|V$ est isomorphe à un faisceau associé à un
$\Gamma(V,\mathscr{O}_X)$-module (\hyperref[I.1.4.1-fr]{1.4.1}). Il est clair
que, sur un préschéma $X$, le faisceau structural $\mathscr{O}_X$ est
quasi-cohérent et que noyaux, conoyaux, images d'homomorphismes de
$\mathscr{O}_X$-Modules quasi-cohérents, limites inductives et sommes directes
de $\mathscr{O}_X$-Modules quasi-cohérents sont quasi-cohérents
(\hyperref[I.1.3.7-fr]{1.3.7} et \hyperref[I.1.3.9-fr]{1.3.9}).
\end{env}

\begin{proposition}[4.1.2]
\label{I.4.1.2-fr}
Soient $X$ un préschéma, $\mathscr{I}$ un faisceau quasi-cohérent d'idéaux
dans $\mathscr{O}_X$. Le support $Y$ du faisceau
$\mathscr{O}_X/\mathscr{I}$ est alors fermé, et si on désigne par
$\mathscr{O}_Y$ la restriction de $\mathscr{O}_X/\mathscr{I}$ à $Y$,
$(Y,\mathscr{O}_Y)$ est un préschéma.
\end{proposition}

\oldpage[I]{120}
Il suffit évidemment (\hyperref[I.2.1.3-fr]{2.1.3}) de considérer le cas où
$X$ est un schéma affine, et de montrer que dans ce cas $Y$ est fermé dans
$X$ et est un \emph{schéma affine}. En effet, si
$X=\operatorname{Spec}(A)$, on a $\mathscr{O}_X=\widetilde{A}$ et
$\mathscr{I}=\widetilde{\mathfrak{I}}$, où $\mathfrak{I}$ est un idéal de
$A$ (\hyperref[I.1.4.1-fr]{1.4.1})~; $Y$ est alors égal à la partie fermée
$V(\mathfrak{I})$ de $X$ et s'identifie au spectre premier de l'anneau
$B=A/\mathfrak{I}$ (\hyperref[I.1.1.11-fr]{1.1.11})~; de plus, si $\varphi$
est l'homomorphisme canonique $A\to B=A/\mathfrak{I}$, l'image directe
${}^a\varphi_*(\widetilde{B})$ s'identifie canoniquement au faisceau
$\widetilde{A}/\widetilde{\mathfrak{I}}=\mathscr{O}_X/\mathscr{I}$
(\hyperref[I.1.6.3-fr]{1.6.3} et \hyperref[I.1.3.9-fr]{1.3.9}), ce qui achève
la démonstration.

Nous dirons que $(Y,\mathscr{O}_Y)$ est le \emph{sous-préschéma} de
$(X,\mathscr{O}_X)$ \emph{défini par le faisceau d'idéaux $\mathscr{I}$}~;
c'est un cas particulier de la notion générale de \emph{sous-préschéma}~:

\begin{definition}[4.1.3]
\label{I.4.1.3-fr}
On dit qu'un espace annelé $(Y,\mathscr{O}_Y)$ est un
\emph{sous-préschéma} d'un préschéma $(X,\mathscr{O}_X)$ si~:
\begin{enumerate}
  \item[$1^\circ$] $Y$ est un sous-espace localement fermé de $X$~;
  \item[$2^\circ$] Si $U$ désigne le plus grand ouvert de $X$ contenant $Y$
  et tel que $Y$ soit fermé dans $U$ (autrement dit, le complémentaire dans
  $X$ de la frontière de $Y$ par rapport à $\overline{Y}$),
  $(Y,\mathscr{O}_Y)$ est un sous-préschéma de
  $(U,\mathscr{O}_X|U)$ défini par un faisceau quasi-cohérent d'idéaux de
  $\mathscr{O}_X|U$.
\end{enumerate}
On dit que le sous-préschéma $(Y,\mathscr{O}_Y)$ de $(X,\mathscr{O}_X)$ est
\emph{fermé} si $Y$ est fermé dans $X$ (auquel cas $U=X$).
\end{definition}

Il résulte aussitôt de cette définition et de
(\hyperref[I.4.1.2-fr]{4.1.2}) que les sous-préschémas fermés de $X$ sont en
\emph{correspondance biunivoque canonique} avec les faisceaux d'idéaux
quasi-cohérents $\mathscr{J}$ de $\mathscr{O}_X$, car le fait que deux tels
faisceaux $\mathscr{J}$, $\mathscr{J}'$ aient même support (fermé) $Y$ et
soient tels que les restrictions de $\mathscr{O}_X/\mathscr{J}$ et de
$\mathscr{O}_X/\mathscr{J}'$ à $Y$ soient identiques, entraîne aussitôt que
$\mathscr{J}'=\mathscr{J}$.

\begin{env}[4.1.4]
\label{I.4.1.4-fr}
Soient $(Y,\mathscr{O}_Y)$ un sous-préschéma de $X$, $U$ le plus grand ouvert
de $X$ contenant $Y$ et dans lequel $Y$ soit fermé, $V$ un ouvert de $X$
contenu dans $U$~; alors $V\cap Y$ est fermé dans $V$. En outre, si $Y$ est
défini par le faisceau quasi-cohérent $\mathscr{J}$ d'idéaux de
$\mathscr{O}_X|U$, $\mathscr{J}|V$ est un faisceau quasi-cohérent d'idéaux de
$\mathscr{O}_X|V$, et il est immédiat que le préschéma induit par $Y$ sur
$Y\cap V$ est le sous-préschéma fermé de $V$ défini par le faisceau d'idéaux
$\mathscr{J}|V$. Inversement~:
\end{env}

\begin{proposition}[4.1.5]
\label{I.4.1.5-fr}
Soit $(Y,\mathscr{O}_Y)$ un espace annelé tel que $Y$ soit un sous-espace de
$X$ et qu'il existe un recouvrement $(V_\alpha)$ de $Y$ par des ouverts de
$X$ tels que pour tout $\alpha$, $Y\cap V_\alpha$ soit fermé dans $V_\alpha$
et que l'espace annelé
$(Y\cap V_\alpha,\mathscr{O}_Y|(Y\cap V_\alpha))$ soit un sous-préschéma
fermé du préschéma induit sur $V_\alpha$ par $X$. Alors
$(Y,\mathscr{O}_Y)$ est un sous-préschéma de $X$.
\end{proposition}

L'hypothèse implique que $Y$ est localement fermé dans $X$ et que le plus
grand ouvert $U$ contenant $Y$ et dans lequel $Y$ est fermé contient tous les
$V_\alpha$~; on peut donc se ramener au cas où $U=X$ et $Y$ est fermé dans
$X$. On définit alors un faisceau quasi-cohérent d'idéaux $\mathscr{J}$ de
$\mathscr{O}_X$ en prenant pour $\mathscr{J}|V_\alpha$ le faisceau d'idéaux de
$\mathscr{O}_X|V_\alpha$ qui définit le sous-préschéma fermé
$(Y\cap V_\alpha,\mathscr{O}_Y|(Y\cap V_\alpha))$ et pour tout ouvert $W$ de
$X$ ne rencontrant pas $Y$, $\mathscr{J}|W=\mathscr{O}_X|W$. On vérifie
immédiatement en vertu de (\hyperref[I.4.1.3-fr]{4.1.3}) et
(\hyperref[I.4.1.4-fr]{4.1.4}) qu'il existe un faisceau d'idéaux $\mathscr{J}$
et un seul satisfaisant à ces conditions et qu'il définit le sous-préschéma
fermé $(Y,\mathscr{O}_Y)$.

En particulier, le préschéma \emph{induit} par $X$ sur un \emph{ouvert} de
$X$ est un \emph{sous-préschéma} de $X$.

\begin{proposition}[4.1.6]
\label{I.4.1.6-fr}
Un sous-préschéma (resp. un sous-préschéma fermé) d'un
sous-\oldpage[I]{121}préschéma (resp. sous-préschéma fermé) de $X$
s'identifie canoniquement à un sous-préschéma (resp. sous-préschéma fermé) de
$X$.
\end{proposition}

Une partie localement fermée d'un sous-espace localement fermé de $X$ étant
un sous-espace localement fermé de $X$, il est clair
(\hyperref[I.4.1.5-fr]{4.1.5}) que la question est locale et qu'on peut donc
supposer $X$ affine~; la proposition résulte alors de l'identification
canonique de $A/\mathfrak{J}'$ et de
$(A/\mathfrak{J})/(\mathfrak{J}'/\mathfrak{J})$ lorsque
$\mathfrak{J}$, $\mathfrak{J}'$ sont deux idéaux d'un anneau $A$ tels que
$\mathfrak{J}\subset\mathfrak{J}'$.

Nous ferons toujours par la suite l'identification précédente.

\begin{env}[4.1.7]
\label{I.4.1.7-fr}
Soit $Y$ un sous-préschéma d'un préschéma $X$, et désignons par $\psi$
l'injection canonique $Y\to X$ des \emph{espaces sous-jacents}~; on sait que
l'image réciproque $\psi^*(\mathscr{O}_X)$ est la restriction
$\mathscr{O}_X|Y$ (\hyperref[0.3.7.1-fr]{0, 3.7.1}). Si, pour tout $y\in Y$,
on désigne par $\omega_y$ l'homomorphisme canonique
$(\mathscr{O}_X)_y\to(\mathscr{O}_Y)_y$, ces homomorphismes sont les
restrictions aux fibres d'un homomorphisme \emph{surjectif} $\omega$ de
faisceaux d'anneaux $\mathscr{O}_X|Y\to\mathscr{O}_Y$~: il suffit en effet de
le vérifier localement sur $Y$, c'est-à-dire que l'on peut supposer $X$
affine et le sous-préschéma $Y$ fermé~; si dans ce cas $\mathscr{I}$ est le
faisceau d'idéaux dans $\mathscr{O}_X$ qui définit $Y$, les $\omega_y$ ne sont
autres que les restrictions aux fibres de l'homomorphisme
$\mathscr{O}_X|Y\to(\mathscr{O}_X/\mathscr{I})|Y$.

Nous avons donc défini un \emph{monomorphisme d'espaces annelés}
(\hyperref[0.4.1.1-fr]{0, 4.1.1}) $j=(\psi,\omega^b)$ qui est évidemment un
morphisme $Y\to X$ de préschémas (\hyperref[I.2.2.1-fr]{2.2.1}), et que nous
appellerons le \emph{morphisme d'injection canonique}.

Si $f:X\to Z$ est un morphisme, nous dirons encore que le morphisme composé
$Y\xrightarrow{j}X\xrightarrow{f}Z$ est la \emph{restriction} de $f$ au
sous-préschéma $Y$.
\end{env}

\begin{env}[4.1.8]
\label{I.4.1.8-fr}
Conformément aux définitions générales (T, I, 1.1), nous dirons qu'un
morphisme de préschémas $f:Z\to X$ est \emph{majoré} par le morphisme
d'injection $j:Y\to X$ d'un sous-préschéma $Y$ de $X$ si $f$ se factorise en
$Z\xrightarrow{g}Y\xrightarrow{j}X$, où $g$ est un morphisme de préschémas~;
$g$ est nécessairement \emph{unique} puisque $j$ est un monomorphisme.
\end{env}

\begin{proposition}[4.1.9]
\label{I.4.1.9-fr}
Pour qu'un morphisme $f:Z\to X$ soit \emph{majoré par un morphisme
d'injection} $j:Y\to X$, il faut et il suffit que $f(Z)\subset Y$ et que,
pour tout $z\in Z$, si on pose $y=f(z)$, l'homomorphisme
$(\mathscr{O}_X)_y\to\mathscr{O}_z$ correspondant à $f$ se factorise en
$(\mathscr{O}_X)_y\to(\mathscr{O}_Y)_y\to\mathscr{O}_z$ (ou, ce qui revient
au même, que le noyau de $(\mathscr{O}_X)_y\to\mathscr{O}_z$ contienne celui
de $(\mathscr{O}_X)_y\to(\mathscr{O}_Y)_y$).
\end{proposition}

Les conditions sont évidemment nécessaires. Pour voir qu'elles sont
suffisantes, on peut se ramener au cas où $Y$ est un sous-préschéma
\emph{fermé} de $X$, en remplaçant au besoin $X$ par un ouvert $U$ tel que
$Y$ soit fermé dans $U$ (\hyperref[I.4.1.3-fr]{4.1.3})~; $Y$ est alors défini
par un faisceau quasi-cohérent $\mathscr{I}$ d'idéaux de $\mathscr{O}_X$.
Posons $f=(\psi,\theta)$, et soit $\mathscr{J}$ le faisceau d'idéaux de
$\psi^*(\mathscr{O}_X)$, noyau de
$\theta^\sharp:\psi^*(\mathscr{O}_X)\to\mathscr{O}_Z$~; compte tenu des
propriétés du foncteur $\psi^*$ (\hyperref[0.3.7.2-fr]{0, 3.7.2}),
l'hypothèse entraîne que pour tout $z\in Z$, on a
$(\psi^*(\mathscr{I}))_z\subset\mathscr{J}_z$, et par suite
$\psi^*(\mathscr{I})\subset\mathscr{J}$. Par suite, $\theta^\sharp$ se
factorise en
\[
  \psi^*(\mathscr{O}_X)\longrightarrow
  \psi^*(\mathscr{O}_X)/\psi^*(\mathscr{I})
  =\psi^*(\mathscr{O}_X/\mathscr{I})
  \xrightarrow{\omega}\mathscr{O}_Z,
\]
la première flèche étant l'homomorphisme canonique. Soit $\psi'$ l'application
continue $Z\to Y$ coïncidant avec $\psi$~; il est clair que l'on a
${\psi'}^*(\mathscr{O}_Y)=\psi^*(\mathscr{O}_X/\mathscr{I})$~; d'autre part,
$\omega$ est évidemment un homomorphisme local, donc
$g=(\psi',\omega^b)$ est un morphisme $Z\to Y$
\oldpage[I]{122}
de préschémas (\hyperref[I.2.2.1-fr]{2.2.1}), qui, en vertu de ce qui précède,
est tel que $f=j\circ g$, d'où la proposition.

\begin{corollary}[4.1.10]
\label{I.4.1.10-fr}
Pour qu'un morphisme d'injection $Z\to X$ soit \emph{majoré par le morphisme
d'injection} $Y\to X$, il faut et il suffit que $Z$ soit un sous-préschéma de
$Y$.
\end{corollary}

On écrit alors $Z\leq Y$, et cette relation est évidemment une
\emph{relation d'ordre} dans l'ensemble des sous-préschémas de $X$.

\subsection{Morphismes d'immersion}
\label{subsection:I.4.2-fr}

\begin{definition}[4.2.1]
\label{I.4.2.1-fr}
On dit qu'un morphisme $f:Y\to X$ est une \emph{immersion} (resp. une
\emph{immersion fermée}, une \emph{immersion ouverte}) s'il se factorise en
$Y\xrightarrow{g}Z\xrightarrow{j}X$, où $g$ est un isomorphisme, $Z$ un
sous-préschéma de $X$ (resp. un sous-préschéma fermé, un sous-préschéma induit
sur un ouvert) et $j$ le morphisme d'injection.
\end{definition}

Le sous-préschéma $Z$ et l'isomorphisme $g$ sont alors déterminés de façon
\emph{unique}, car si $Z'$ est un second sous-préschéma de $X$, $j'$ l'injection
$Z'\to X$ et $g'$ un isomorphisme $Y\to Z'$ tel que
$j\circ g=j'\circ g'$, on en déduit $j'=j\circ g\circ {g'}^{-1}$, d'où
$Z'\leq Z$ (\hyperref[I.4.1.10-fr]{4.1.10}), et on montre de même que
$Z\leq Z'$, donc $Z'=Z$, et comme $j$ est un monomorphisme de préschémas,
$g'=g$.

On dit que $f=j\circ g$ est la \emph{factorisation canonique} de l'immersion
$f$, et le sous-préschéma $Z$ et l'isomorphisme $g$ sont dits
\emph{associés à $f$}.

Il est clair qu'une immersion est un \emph{monomorphisme} de préschémas
(\hyperref[I.4.1.7-fr]{4.1.7}) et \emph{a fortiori} un morphisme
\emph{radiciel} (\hyperref[I.3.5.4-fr]{3.5.4}).

\begin{proposition}[4.2.2]
\label{I.4.2.2-fr}
\begin{enumerate}
  \item[a)] Pour qu'un morphisme $f=(\psi,\theta):Y\to X$ soit une
  \emph{immersion ouverte}, il faut et il suffit que $\psi$ soit un
  homéomorphisme de $Y$ sur une partie ouverte de $X$, et que pour tout
  $y\in Y$, l'homomorphisme
  $\theta_y^\sharp:\mathscr{O}_{\psi(y)}\to\mathscr{O}_y$ soit bijectif.
  \item[b)] Pour qu'un morphisme $f=(\psi,\theta):Y\to X$ soit une
  \emph{immersion} (resp. une \emph{immersion fermée}), il faut et il suffit
  que $\psi$ soit un homéomorphisme de $Y$ sur une partie localement fermée
  (resp. fermée) de $X$, et que pour tout $y\in Y$, l'homomorphisme
  $\theta_y^\sharp:\mathscr{O}_{\psi(y)}\to\mathscr{O}_y$ soit surjectif.
\end{enumerate}
\end{proposition}

\begin{enumerate}
  \item[a)] Les conditions sont évidemment nécessaires. Inversement, si elles
  sont remplies, il est clair que $\theta^\sharp$ est un isomorphisme de
  $\mathscr{O}_Y$ sur $\psi^*(\mathscr{O}_X)$, et
  $\psi^*(\mathscr{O}_X)$ est le faisceau déduit par transport de structure au
  moyen de $\psi^{-1}$ à partir de $\mathscr{O}_X|\psi(Y)$~; d'où la
  conclusion.
  \item[b)] Les conditions étant évidemment nécessaires, prouvons qu'elles
  sont suffisantes. Considérons d'abord le cas particulier où l'on suppose que
  $X$ est un schéma affine et que $Z=\psi(Y)$ est \emph{fermé} dans $X$. On
  sait (\hyperref[0.3.4.6-fr]{0, 3.4.6}) que $\psi_*(\mathscr{O}_Y)$ a alors
  pour support $Z$ et que si l'on désigne par $\mathscr{O}'_Z$ sa restriction
  à $Z$, l'espace annelé $(Z,\mathscr{O}'_Z)$ se déduit de
  $(Y,\mathscr{O}_Y)$ par transport de structure au moyen de
  l'homéomorphisme $\psi$, considéré comme application de $Y$ sur $Z$.
  Montrons que $f_*(\mathscr{O}_Y)=\psi_*(\mathscr{O}_Y)$ est un
  $\mathscr{O}_X$-Module \emph{quasi-cohérent}. En effet, pour tout
  $x\notin Z$, $\psi_*(\mathscr{O}_Y)$, restreint à un voisinage convenable de
  $x$, est nul. Si au contraire $x\in Z$, on a $x=\psi(y)$ pour un $y\in Y$
  bien déterminé~; soit $V$ un voisinage ouvert affine de $y$ dans $Y$~;
  $\psi(V)$ est alors ouvert dans $Z$, donc trace sur $Z$ d'un ouvert $U$ de
  $X$, et la restriction à $U$ de $\psi_*(\mathscr{O}_Y)$ est identique à la
  restriction à $U$ de l'image
\oldpage[I]{123}
  directe $(\psi_V)_*(\mathscr{O}_Y|V)$, où $\psi_V$ est la restriction de
  $\psi$ à $V$. Or, la restriction à $(V,\mathscr{O}_Y|V)$ du morphisme
  $(\psi,\theta)$ est un morphisme de ce préschéma dans
  $(X,\mathscr{O}_X)$, et par suite est de la forme
  $({}^a\varphi,\widetilde{\varphi})$, où $\varphi$ est un homomorphisme de
  l'anneau $A=\Gamma(X,\mathscr{O}_X)$ dans l'anneau
  $\Gamma(V,\mathscr{O}_Y)$ (\hyperref[I.1.7.3-fr]{1.7.3})~; on en conclut
  que $(\psi_V)_*(\mathscr{O}_Y|V)$ est un $\mathscr{O}_X$-Module
  quasi-cohérent (\hyperref[I.1.6.3-fr]{1.6.3}), ce qui prouve notre
  assertion, en raison du caractère local des faisceaux quasi-cohérents. En
  outre, l'hypothèse que $\psi$ est un homéomorphisme entraîne
  (\hyperref[0.3.4.5-fr]{0, 3.4.5}) que pour tout $y\in Y$, $\psi_y$ est un
  isomorphisme
  $(\psi_*(\mathscr{O}_Y))_{\psi(y)}\to\mathscr{O}_y$~; comme le diagramme
  \phantomsection\label{I.4.2.2.theta-diagram-fr}
  \[
    \xymatrix{
      \mathscr{O}_{\psi(y)}\ar[r]^{\theta_{\psi(y)}}
        \ar[d]_{\psi_y\circ\rho_{\psi(y)}} &
      (\psi_*(\mathscr{O}_Y))_{\psi(y)}\ar[d]^{\psi_y}\\
      (\psi^*(\mathscr{O}_X))_y\ar[r]^{\theta_y^\sharp} &
      \mathscr{O}_y
    }
  \]
  est commutatif et que les deux flèches verticales sont des isomorphismes
  (\hyperref[0.3.7.2-fr]{0, 3.7.2}), l'hypothèse que $\theta_y^\sharp$ est
  surjectif entraîne qu'il en est de même de $\theta_{\psi(y)}$. Comme le
  support de $\psi_*(\mathscr{O}_Y)$ est $Z=\psi(Y)$, $\theta$ est un
  homomorphisme \emph{surjectif} de $\mathscr{O}_X=\widetilde{A}$ dans le
  $\mathscr{O}_X$-Module quasi cohérent $f_*(\mathscr{O}_Y)$. Par suite, il
  existe un isomorphisme unique $\omega$ d'un faisceau quotient
  $\widetilde{A}/\widetilde{\mathfrak{J}}$ ($\mathfrak{J}$ idéal de $A$) sur
  $f_*(\mathscr{O}_Y)$ qui, composé avec l'homomorphisme canonique
  $\widetilde{A}\to\widetilde{A}/\widetilde{\mathfrak{J}}$, donne $\theta$
  (\hyperref[I.1.3.8-fr]{1.3.8})~; si $\mathscr{O}_Z$ désigne la restriction
  de $\widetilde{A}/\widetilde{\mathfrak{J}}$ à $Z$,
  $(Z,\mathscr{O}_Z)$ est un sous-préschéma de $(X,\mathscr{O}_X)$, et $f$ se
  factorise en l'injection canonique de ce sous-préschéma dans $X$, et
  l'isomorphisme $(\psi_0,\omega_0)$, où $\psi_0$ est $\psi$ considéré comme
  application de $Y$ sur $Z$, et $\omega_0$ la restriction de $\omega$ à
  $\mathscr{O}_Z$.

  Passons au cas général. Soit $U$ un ensemble ouvert affine dans $X$ tel que
  $U\cap\psi(Y)$ soit fermé dans $U$ et non vide. En restreignant $f$ au
  préschéma induit par $Y$ sur l'ouvert $\psi^{-1}(U)$, et en le considérant
  comme un morphisme de ce préschéma dans le préschéma induit par $X$ sur
  $U$, on est ramené au premier cas~; la restriction de $f$ à
  $\psi^{-1}(U)$ est donc une immersion fermée $\psi^{-1}(U)\to U$, se
  factorisant canoniquement en $j_U\circ g_U$, où $g_U$ est un isomorphisme
  du préschéma $\psi^{-1}(U)$ sur un sous-préschéma $Z_U$ de $U$, et $j_U$
  l'injection canonique $Z_U\to U$. Soit $V$ un second ouvert affine de $X$
  tel que $V\subset U$~; comme la restriction $Z'_V$ de $Z_U$ à $V$ est un
  sous-préschéma du préschéma $V$, la restriction de $f$ à $\psi^{-1}(V)$ se
  factorise en $j'_V\circ g'_V$, où $j'_V$ est l'injection canonique
  $Z'_V\to V$ et $g'_V$ un isomorphisme de $\psi^{-1}(V)$ sur $Z'_V$. Par
  l'unicité de la factorisation canonique d'une immersion
  (\hyperref[I.4.2.1-fr]{4.2.1}), on a nécessairement $Z'_V=Z_V$ et
  $g'_V=g_V$. On en conclut (\hyperref[I.4.1.5-fr]{4.1.5}) qu'il y a un
  sous-préschéma $Z$ de $X$ dont l'espace sous-jacent est $\psi(Y)$ et dont
  la restriction à chaque $U\cap\psi(Y)$ est $Z_U$~; les $g_U$ sont alors les
  restrictions aux $\psi^{-1}(U)$ d'un isomorphisme $g:Y\to Z$ tel que
  $f=j\circ g$, où $j$ est l'injection canonique $Z\to X$.
\end{enumerate}

\begin{corollary}[4.2.3]
\label{I.4.2.3-fr}
Soit $X$ un schéma affine. Pour qu'un morphisme
$f=(\psi,\theta):Y\to X$ soit une immersion fermée, il faut et il suffit que
$Y$ soit un schéma affine et que l'homomorphisme
$\Gamma(\psi):\Gamma(\mathscr{O}_X)\to\Gamma(\mathscr{O}_Y)$ soit surjectif.
\end{corollary}

\begin{corollary}[4.2.4]
\label{I.4.2.4-fr}
\begin{enumerate}
  \item[a)] Soient $f$ un morphisme $Y\to X$, $(V_\lambda)$ un recouvrement
  de $f(Y)$ par des ouverts de $X$. Pour que $f$ soit une immersion (resp. une
  immersion ouverte), il faut et il suffit
\oldpage[I]{124}
  que sa restriction à chacun des préschémas induits $f^{-1}(V_\lambda)$ soit
  une immersion (resp. une immersion ouverte) dans $V_\lambda$.
  \item[b)] Soient $f$ un morphisme $Y\to X$, $(V_\lambda)$ un recouvrement
  ouvert de $X$. Pour que $f$ soit une immersion fermée, il faut et il suffit
  que sa restriction à chacun des préschémas induits $f^{-1}(V_\lambda)$ soit
  une immersion fermée dans $V_\lambda$.
\end{enumerate}
\end{corollary}

Soit $f=(\psi,\theta)$~; dans le cas a), $\theta_y^\sharp$ est surjectif
(resp. bijectif) pour tout $y\in Y$, et dans le cas b) $\theta_y^\sharp$ est
surjectif pour tout $y\in Y$~; il suffit donc de vérifier que $\psi$, dans le
cas a), est un homéomorphisme de $Y$ sur une partie localement fermée (resp.
ouverte) de $X$, et dans le cas b), un homéomorphisme de $Y$ sur une partie
fermée de $X$. Or, $\psi$ est évidemment injectif et transforme tout voisinage
de $y$ dans $Y$ en un voisinage de $\psi(y)$ dans $\psi(Y)$ pour tout $y\in Y$,
en vertu de l'hypothèse~; dans le cas a), $\psi(Y)\cap V_\lambda$ est localement
fermé (resp. ouvert) dans $V_\lambda$, donc $\psi(Y)$ est localement fermé
(resp. ouvert) dans la réunion des $V_\lambda$, et \emph{a fortiori} dans $X$~;
dans le cas b), $\psi(Y)\cap V_\lambda$ est fermé dans $V_\lambda$, donc
$\psi(Y)$ est fermé dans $X$ puisque $X=\bigcup_\lambda V_\lambda$.

\begin{proposition}[4.2.5]
\label{I.4.2.5-fr}
Le composé de deux immersions (resp. de deux immersions ouvertes, de deux
immersions fermées) est une immersion (resp. une immersion ouverte, une
immersion fermée).
\end{proposition}

Cela résulte trivialement de (\hyperref[I.4.1.6-fr]{4.1.6}).

\subsection{Produit d'immersions.}
\label{subsection:I.4.3-fr}

\begin{proposition}[4.3.1]
\label{I.4.3.1-fr}
Soient $\alpha:X'\to X$, $\beta:Y'\to Y$ deux $S$-morphismes~; si $\alpha$ et
$\beta$ sont des immersions (resp. des immersions ouvertes, des immersions
fermées), $\alpha\times_S\beta$ est une immersion (resp. une immersion ouverte,
une immersion fermée). En outre, si $\alpha$ (resp. $\beta$) identifie $X'$
(resp. $Y'$) à un sous-préschéma $X''$ (resp. $Y''$) de $X$ (resp. $Y$),
$\alpha\times_S\beta$ identifie l'espace sous-jacent à $X'\times_S Y'$ au
sous-espace $p^{-1}(X'')\cap q^{-1}(Y'')$ de l'espace sous-jacent à
$X\times_S Y$, en désignant par $p$ et $q$ les projections de $X\times_S Y$
dans $X$ et $Y$ respectivement.
\end{proposition}

Compte tenu de la déf. (\hyperref[I.4.2.1-fr]{4.2.1}), on peut se restreindre
au cas où $X'$ et $Y'$ sont des sous-préschémas, $\alpha$ et $\beta$ les
morphismes d'injection. La proposition a déjà été établie pour les
sous-préschémas induits sur les ouverts (\hyperref[I.3.2.7-fr]{3.2.7})~; comme
tout sous-préschéma est un sous-préschéma fermé d'un préschéma induit sur un
ouvert (\hyperref[I.4.1.3-fr]{4.1.3}), on est ramené au cas où $X'$ et $Y'$ sont
des sous-préschémas \emph{fermés}.

Montrons d'abord qu'on peut supposer que $S$ est \emph{affine}. En effet, soit
$(S_\lambda)$ un recouvrement de $S$ formé d'ouverts affines~; si $\varphi$ et
$\psi$ sont les morphismes structuraux de $X$ et $Y$, soit
$X_\lambda=\varphi^{-1}(S_\lambda)$, $Y_\lambda=\psi^{-1}(S_\lambda)$. La
restriction $X'_\lambda$ (resp. $Y'_\lambda$) de $X'$ (resp. $Y'$) à
$X_\lambda\cap X'$ (resp. $Y_\lambda\cap Y'$) est un sous-préschéma fermé de
$X_\lambda$ (resp. $Y_\lambda$), les préschémas $X_\lambda$, $Y_\lambda$,
$X'_\lambda$, $Y'_\lambda$ peuvent être considérés comme des
$S_\lambda$-préschémas et les produits $X_\lambda\times_S Y_\lambda$ et
$X_\lambda\times_{S_\lambda}Y_\lambda$ (resp.
$X'_\lambda\times_S Y'_\lambda$ et
$X'_\lambda\times_{S_\lambda}Y'_\lambda$) sont identiques
(\hyperref[I.3.2.5-fr]{3.2.5}). Si la proposition est vraie lorsque $S$ est
affine, la restriction de $\alpha\times_S\beta$ à chacun des
$X'_\lambda\times_S Y'_\lambda$ sera donc une immersion
(\hyperref[I.3.2.7-fr]{3.2.7}). Comme le produit
$X'_\lambda\times_S Y'_\mu$ (resp. $X_\lambda\times_S Y_\mu$) s'identifie à
$(X'_\lambda\cap X'_\mu)\times_S(Y'_\lambda\cap Y'_\mu)$ (resp.
$(X_\lambda\cap X_\mu)\times_S(Y_\lambda\cap Y_\mu)$)
(\hyperref[I.3.2.6.4-fr]{3.2.6.4}), la restriction de
$\alpha\times_S\beta$
\oldpage[I]{125}
à chacun des $X'_\lambda\times_S Y'_\mu$ est encore une immersion~; il en est
par suite de même de $\alpha\times_S\beta$ en vertu de
(\hyperref[I.4.2.4-fr]{4.2.4}).

En second lieu, prouvons qu'on peut aussi supposer que $X$ et $Y$ sont
\emph{affines}. En effet, soit $(U_i)$ (resp. $(V_j)$) un recouvrement de $X$
(resp. $Y$) par des ouverts affines, et soit $X'_i$ (resp. $Y'_j$) la
restriction de $X'$ (resp. $Y'$) à $X'\cap U_i$ (resp. $Y'\cap V_j$), qui est
un sous-préschéma fermé de $U_i$ (resp. $V_j$)~; $U_i\times_S V_j$ s'identifie
à la restriction de $X\times_S Y$ à
$p^{-1}(U_i)\cap q^{-1}(V_j)$ (\hyperref[I.3.2.7-fr]{3.2.7})~; et de même, si
$p'$, $q'$ sont les projections de $X'\times_S Y'$,
$X'_i\times_S Y'_j$ s'identifie à la restriction de $X'\times_S Y'$ à
${p'}^{-1}(X'_i)\cap{q'}^{-1}(Y'_j)$. Posons
$\gamma=\alpha\times_S\beta$~; on a par définition
$p\circ\gamma=\alpha\circ p'$, $q\circ\gamma=\beta\circ q'$~; comme
$X'_i=\alpha^{-1}(U_i)$, $Y'_j=\beta^{-1}(V_j)$, on a aussi
${p'}^{-1}(X'_i)=\gamma^{-1}(p^{-1}(U_i))$,
${q'}^{-1}(Y'_j)=\gamma^{-1}(q^{-1}(V_j))$, d'où
\[
  {p'}^{-1}(X'_i)\cap{q'}^{-1}(Y'_j)
  =\gamma^{-1}\bigl(p^{-1}(U_i)\cap q^{-1}(V_j)\bigr)
  =\gamma^{-1}(U_i\times_S V_j)
\]
on conclut comme dans la première partie du raisonnement.

Supposons donc $X$, $Y$, $S$ affines, et soient $B$, $C$, $A$ leurs anneaux
respectifs. Alors $B$ et $C$ sont des $A$-algèbres, $X'$ et $Y'$ des schémas
affines dont les anneaux sont des algèbres quotients $B'$, $C'$ de $B$ et
$C$ respectivement. En outre, on a
$\alpha=({}^a\rho,\widetilde{\rho})$, $\beta=({}^a\sigma,\widetilde{\sigma})$,
où $\rho$ et $\sigma$ sont respectivement les homomorphismes canoniques
$B\to B'$, $C\to C'$ (\hyperref[I.1.7.3-fr]{1.7.3}). Cela étant, on sait que
$X\times_S Y$ (resp. $X'\times_S Y'$) est un schéma affine d'anneau
$B\otimes_A C$ (resp. $B'\otimes_A C'$), et
$\alpha\times_S\beta=({}^a\tau,\widetilde{\tau})$, où $\tau$ est
l'homomorphisme $\rho\otimes\sigma$ de $B\otimes_A C$ dans
$B'\otimes_A C'$ (\hyperref[I.3.2.2-fr]{3.2.2} et
\hyperref[I.3.2.3-fr]{3.2.3})~; comme cet homomorphisme est surjectif,
$\alpha\times_S\beta$ est bien une immersion. En outre, si $\mathfrak{b}$
(resp. $\mathfrak{c}$) est le noyau de $\rho$ (resp. $\sigma$), le noyau de
$\tau$ est $u(\mathfrak{b})+v(\mathfrak{c})$, où $u$ (resp. $v$) est
l'homomorphisme $b\mapsto b\otimes1$ (resp. $c\mapsto1\otimes c$). Comme
$p=({}^a u,\widetilde{u})$ et $q=({}^a v,\widetilde{v})$, ce noyau correspond,
dans le spectre premier de $B\otimes_A C$, à l'ensemble fermé
$p^{-1}(X')\cap q^{-1}(Y')$ (\hyperref[I.1.2.2.1-fr]{1.2.2.1} et
\hyperref[I.1.1.2-fr]{1.1.2 (iii)}), ce qui achève la démonstration.

\begin{corollary}[4.3.2]
\label{I.4.3.2-fr}
Si $f:X\to Y$ est une immersion (resp. une immersion ouverte, une immersion
fermée) et un $S$-morphisme, $f_{(S')}$ est une immersion (resp. une immersion
ouverte, une immersion fermée) pour toute extension $S'\to S$ du préschéma de
base.
\end{corollary}

\subsection{Image réciproque d'un sous-préschéma.}
\label{subsection:I.4.4-fr}

\begin{proposition}[4.4.1]
\label{I.4.4.1-fr}
Soient $f:X\to Y$ un morphisme, $Y'$ un sous-préschéma (resp. un
sous-préschéma fermé, un préschéma induit sur un ouvert) de $Y$, $j:Y'\to Y$
le morphisme d'injection. Alors la projection $p:X\times_Y Y'\to X$ est une
immersion (resp. une immersion fermée, une immersion ouverte)~; le
sous-préschéma de $X$ associé à $p$ a $f^{-1}(Y')$ pour espace sous-jacent~;
en outre, si $j'$ est le morphisme d'injection de ce sous-préschéma dans $X$,
pour qu'un morphisme $h:Z\to X$ soit tel que $f\circ h:Z\to Y$ soit majoré
par $j$, il faut et il suffit que $h$ soit majoré par $j'$.
\end{proposition}

Comme $p=1_X\times_Y j$ (\hyperref[I.3.3.4-fr]{3.3.4}), la première assertion
résulte de (\hyperref[I.4.3.1-fr]{4.3.1})~; la seconde est un cas particulier
de (\hyperref[I.3.5.10-fr]{3.5.10}) (où il faut échanger les rôles de $X$ et
de $Y'$). Enfin, si on a $f\circ h=j\circ h'$, où $h'$ est un morphisme
$Z\to Y'$, il résulte de la définition du produit que l'on a $h=p\circ u$,
où $u$ est un morphisme $Z\to X\times_Y Y'$, d'où la dernière assertion.

Nous dirons que le sous-préschéma de $X$ ainsi défini est l'\emph{image
réciproque} du sous-préschéma $Y'$ de $Y$ par le morphisme $f$, terminologie
qui s'accorde avec celle introduite
\oldpage[I]{126}
de façon plus générale dans (\hyperref[I.3.3.6-fr]{3.3.6}). Quand nous
parlerons de $f^{-1}(Y')$ comme d'un sous-préschéma de $X$, c'est toujours de
ce sous-préschéma qu'il s'agira.

Lorsque les préschémas $f^{-1}(Y')$ et $X$ sont identiques, $j'$ est
l'identité et tout morphisme $h:Z\to X$ est donc majoré par $j'$, donc
\emph{le morphisme $f:X\to Y$ se factorise alors en}
$X\xrightarrow{g}Y'\xrightarrow{j}Y$.

Lorsque $y$ est un point \emph{fermé} de $Y$ et
$Y'=\operatorname{Spec}(k(y))$ le plus petit sous-préschéma fermé de $Y$ ayant
$\{y\}$ comme espace sous-jacent (\hyperref[I.4.1.9-fr]{4.1.9}), le
sous-préschéma fermé $f^{-1}(Y')$ est canoniquement isomorphe à la
\emph{fibre} $f^{-1}(y)$ définie en (\hyperref[I.3.6.2-fr]{3.6.2}), à
laquelle on l'identifiera.

\begin{corollary}[4.4.2]
\label{I.4.4.2-fr}
Soient $f:X\to Y$, $g:Y\to Z$ deux morphismes, $h=g\circ f$ leur composé.
Pour tout sous-préschéma $Z'$ de $Z$, les sous-préschémas
$f^{-1}(g^{-1}(Z'))$ et $h^{-1}(Z')$ de $X$ sont identiques.
\end{corollary}

Cela résulte de l'existence de l'isomorphisme canonique
$X\times_Y(Y\times_Z Z')\xrightarrow{\sim}X\times_Z Z'$
(\hyperref[I.3.3.9.1-fr]{3.3.9.1}).

\begin{corollary}[4.4.3]
\label{I.4.4.3-fr}
Soient $X'$, $X''$ deux sous-préschémas de $X$, $j':X'\to X$,
$j'':X''\to X$ les morphismes d'injection~; alors, ${j'}^{-1}(X'')$ et
${j''}^{-1}(X')$ sont tous deux égaux à la borne inférieure
$\inf(X',X'')$ de $X'$ et $X''$ pour la relation d'ordre entre
sous-préschémas, et canoniquement isomorphée à $X'\times_X X''$.
\end{corollary}

Cela résulte aussitôt de (\hyperref[I.4.4.1-fr]{4.4.1}) et
(\hyperref[I.4.1.10-fr]{4.1.10}).

\begin{corollary}[4.4.4]
\label{I.4.4.4-fr}
Soient $f:X\to Y$ un morphisme, $Y'$, $Y''$ deux sous-préschémas de $Y$~;
on a
\[
  f^{-1}(\inf(Y',Y''))=\inf(f^{-1}(Y'),f^{-1}(Y'')).
\]
\end{corollary}

Cela résulte de l'existence de l'isomorphisme canonique entre
$(X\times_Y Y')\times_X(X\times_Y Y'')$ et
$X\times_Y(Y'\times_Y Y'')$ (\hyperref[I.3.3.9.1-fr]{3.3.9.1}).

\begin{proposition}[4.4.5]
\label{I.4.4.5-fr}
Soient $f:X\to Y$ un morphisme, $Y'$ un sous-préschéma fermé de $Y$ défini
par un faisceau quasi-cohérent $\mathscr{K}$ d'idéaux de $\mathscr{O}_Y$
(\hyperref[I.4.1.3-fr]{4.1.3})~; le sous-préschéma fermé $f^{-1}(Y')$ de $X$
est alors défini par le faisceau quasi cohérent d'idéaux
$f^*(\mathscr{K})\mathscr{O}_X$ de $\mathscr{O}_X$.
\end{proposition}

La question est évidemment locale sur $X$ et $Y$~; il suffit alors de
remarquer que si $B$ est une $A$-algèbre et $\mathfrak{K}$ un idéal de $B$,
on a $A\otimes_B(B/\mathfrak{K})=A/\mathfrak{K}A$, et d'appliquer
(\hyperref[I.1.6.9-fr]{1.6.9}).

\begin{corollary}[4.4.6]
\label{I.4.4.6-fr}
Soient $X'$ un sous-préschéma fermé de $X$ défini par un faisceau quasi
cohérent d'idéaux $\mathscr{J}$ de $\mathscr{O}_X$, $i$ l'injection
$X'\to X$~; pour que la restriction $f\circ i$ de $f$ à $X'$ soit majorée
par l'injection $j:Y'\to Y$ (autrement dit se factorise en $j\circ g$, où
$g$ est un morphisme $X'\to Y'$), il faut et il suffit que
$f^*(\mathscr{K})\mathscr{O}_X\subset\mathscr{J}$.
\end{corollary}

Il suffit d'appliquer à $i$ la prop. (\hyperref[I.4.4.1-fr]{4.4.1}) en
tenant compte de (\hyperref[I.4.4.5-fr]{4.4.5}).

\subsection{Immersions locales et isomorphismes locaux.}
\label{subsection:I.4.5-fr}

\begin{definition}[4.5.1]
\label{I.4.5.1-fr}
Soit $f:X\to Y$ un morphisme de préschémas. On dit que $f$ est une
\emph{immersion locale en un point $x\in X$} s'il existe un voisinage ouvert
$U$ de $x$ dans $X$ et un voisinage ouvert $V$ de $f(x)$ dans $Y$ tels que
la restriction de $f$ au préschéma induit $U$ soit une immersion fermée de
$U$ dans le préschéma induit $V$. On dit que $f$ est une immersion locale si
$f$ est une immersion locale en tout point de $X$.
\end{definition}

\begin{definition}[4.5.2]
\label{I.4.5.2-fr}
On dit qu'un morphisme $f:X\to Y$ est un isomorphisme local en
\oldpage[I]{127}
un point $x\in X$ s'il existe un voisinage ouvert $U$ de $x$ dans $X$ tel
que la restriction de $f$ au préschéma induit $U$ soit une immersion ouverte
de $U$ dans $Y$. On dit que $f$ est un isomorphisme local si $f$ est un
isomorphisme local en tout point de $X$.
\end{definition}

\begin{env}[4.5.3]
\label{I.4.5.3-fr}
Une immersion (resp. une immersion fermée) $f:X\to Y$ peut donc se
caractériser comme une immersion locale telle que $f$ soit un homéomorphisme
de l'espace sous-jacent à $X$ sur une partie (resp. une partie fermée) de
$Y$. Une immersion ouverte $f$ peut se caractériser comme un isomorphisme
local \emph{injectif}.
\end{env}

\begin{proposition}[4.5.4]
\label{I.4.5.4-fr}
Soient $X$ un préschéma irréductible, $f:X\to Y$ un morphisme injectif
dominant. Si $f$ est une immersion locale, $f$ est une immersion et $f(X)$
est ouvert dans $Y$.
\end{proposition}

En effet, soit $x\in X$ et soient $U$ un voisinage ouvert de $x$, $V$ un
voisinage ouvert de $f(x)$ dans $Y$ tels que la restriction de $f$ à $U$
soit une immersion fermée dans $V$~; comme $U$ est dense dans $X$, $f(U)$
est dense dans $Y$ par hypothèse, donc $f(U)=V$ et $f$ est un homéomorphisme
de $U$ sur $V$~; l'hypothèse que $f$ est injectif entraîne que
$f^{-1}(V)=U$, d'où aussitôt la proposition.

\begin{proposition}[4.5.5]
\label{I.4.5.5-fr}
\begin{enumerate}
  \item[\rm(i)] Le composé de deux immersions locales (resp. de deux
    isomorphismes locaux) est une immersion locale (resp. un isomorphisme
    local).
  \item[\rm(ii)] Soient $f:X\to X'$, $g:Y\to Y'$ deux $S$-morphismes. Si
    $f$ et $g$ sont des immersions locales (resp. des isomorphismes locaux),
    il en est de même de $f\times_S g$.
  \item[\rm(iii)] Si un $S$-morphisme $f$ est une immersion locale (resp. un
    isomorphisme local), il en est de même de $f_{(S')}$ pour toute extension
    $S'\to S$ du préschéma de base.
\end{enumerate}
\end{proposition}

D'après (\hyperref[I.3.5.1-fr]{3.5.1}), il suffit de prouver (i) et (ii).

(i) résulte aussitôt de la transitivité des immersions fermées (resp.
ouvertes) (\hyperref[I.4.2.4-fr]{4.2.4}) et du fait que si $f$ est un
homéomorphisme de $X$ sur une partie fermée de $Y$, pour tout ouvert
$U\subset X$, $f(U)$ est ouvert dans $f(X)$, donc il existe une partie
ouverte $V$ de $Y$ telle que $f(U)=V\cap f(X)$, et $f(U)$ est par suite fermé
dans $V$.

Pour démontrer (ii), soient $p$, $q$ les projections de $X\times_S Y$, $p'$,
$q'$ celles de $X'\times_S Y'$. Il existe par hypothèse des voisinages
ouverts $U$, $U'$, $V$, $V'$ de $x=p(z)$, $x'=p'(z')$, $y=q(z)$,
$y'=q'(z')$ respectivement, tels que les restrictions de $f$ et $g$ à $U$
et $V$ respectivement soient des immersions fermées (resp. ouvertes) dans
$U'$ et $V'$ respectivement. Comme l'espace sous-jacent de $U\times_S V$ et
celui de $U'\times_S V'$ s'identifient aux voisinages ouverts
$p^{-1}(U)\cap q^{-1}(V)$ et ${p'}^{-1}(U')\cap{q'}^{-1}(V')$ de $z$ et
$z'$ respectivement (\hyperref[I.3.2.7-fr]{3.2.7}), la proposition résulte
de (\hyperref[I.4.3.1-fr]{4.3.1}).

% FIN SOURCE ega1/ega1-4-fr.tex

%
% DEBUT SOURCE ega1/ega1-5-fr.tex
\section{Préschémas réduits~; condition de séparation}
\label{section:I.5-fr}

\subsection{Préschémas réduits.}
\label{subsection:I.5.1-fr}

\begin{proposition}[5.1.1]
\label{I.5.1.1-fr}
Soient $(X,\mathscr{O}_X)$ un préschéma, $\mathscr{B}$ une
$\mathscr{O}_X$-Algèbre quasi-cohérente. Il existe un
$\mathscr{O}_X$-Module quasi-cohérent et un seul $\mathscr{N}$ dont la fibre
$\mathscr{N}_x$ en tout $x\in X$ soit le nilradical de l'anneau
$\mathscr{B}_x$. Lorsque $X$ est affine, et par suite
$\mathscr{B}=\widetilde{B}$, où $B$ est une algèbre sur $A(X)$, on a
$\mathscr{N}=\widetilde{\mathfrak{N}}$, où $\mathfrak{N}$ est le nilradical
de $B$.
\end{proposition}

\oldpage[I]{128}
La question étant locale, on est ramené à démontrer la dernière assertion. On
sait que $\widetilde{\mathfrak{N}}$ est un $\mathscr{O}_X$-Module
quasi-cohérent (\hyperref[I.1.4.1-fr]{1.4.1}) et que sa fibre au point
$x\in X$ est l'idéal $\mathfrak{N}_x$ de l'anneau de fractions $B_x$~; tout
revient à prouver que le nilradical de $B_x$ est contenu dans
$\mathfrak{N}_x$, l'inclusion opposée étant évidente. Or, soit $z/s$ un
élément du nilradical de $B_x$, avec $z\in B$, $s\notin\mathfrak{j}_x$~; par
hypothèse, il existe un entier $k$ tel que $(z/s)^k=0$, ce qui signifie qu'il
existe $t\notin\mathfrak{j}_x$ tel que $tz^k=0$. On en conclut $(tz)^k=0$, et
par suite $z/s=(tz)/(ts)$ appartient à $\mathfrak{N}_x$.

Nous dirons que le $\mathscr{O}_X$-Module quasi-cohérent $\mathscr{N}$ ainsi
défini est le \emph{Nilradical} de la $\mathscr{O}_X$-Algèbre $\mathscr{B}$~;
nous noterons en particulier $\mathscr{N}_X$ le nilradical de
$\mathscr{O}_X$.

\begin{corollary}[5.1.2]
\label{I.5.1.2-fr}
Soit $X$ un préschéma~; le sous-préschéma fermé de $X$ défini par le faisceau
d'idéaux $\mathscr{N}_X$ est le seul sous-préschéma réduit
(\hyperref[0.4.1.4-fr]{0, 4.1.4}) de $X$ ayant $X$ pour espace sous-jacent~;
c'est aussi le plus petit sous-préschéma de $X$ ayant $X$ pour espace
sous-jacent.
\end{corollary}

Comme le faisceau structural du sous-préschéma fermé défini $Y$ par
$\mathscr{N}_X$ est $\mathscr{O}_X/\mathscr{N}_X$, il est immédiat que $Y$
est réduit et a $X$ pour espace sous-jacent, puisque
$\mathscr{N}_x\ne\mathscr{O}_x$ pour tout $x\in X$. Pour démontrer les autres
assertions, remarquons qu'un sous-préschéma $Z$ de $X$ ayant $X$ pour espace
sous-jacent est défini par un faisceau d'idéaux $\mathscr{J}$
(\hyperref[I.4.1.3-fr]{4.1.3}) tel que
$\mathscr{J}_x\ne\mathscr{O}_x$ pour tout $x\in X$. On peut se borner au cas
où $X$ est affine, soit $X=\operatorname{Spec}(A)$ et
$\mathscr{J}=\widetilde{\mathfrak{J}}$, où $\mathfrak{J}$ est un idéal de
$A$~; alors, pour tout $x\in X$, on a
$\mathfrak{J}\subset\mathfrak{j}_x$, donc $\mathfrak{J}$ est contenu dans
tous les idéaux premiers de $A$, c'est-à-dire dans leur intersection
$\mathfrak{N}$, nilradical de $A$. Cela prouve que $Y$ est le plus petit
sous-préschéma de $X$ ayant $X$ comme espace sous-jacent
(\hyperref[I.4.1.9-fr]{4.1.9})~; en outre, si $Z$ est distinct de $Y$, on a
nécessairement $\mathscr{J}_x\ne\mathscr{N}_x$ pour un $x\in X$ au moins, et
par suite (\hyperref[I.5.1.1-fr]{5.1.1}) $Z$ n'est pas réduit.

\begin{definition}[5.1.3]
\label{I.5.1.3-fr}
On appelle préschéma réduit associé à un préschéma $X$ et on note
$X_{\mathrm{red}}$ l'unique sous-préschéma réduit de $X$ ayant $X$ pour
espace sous-jacent.
\end{definition}

Dire qu'un préschéma $X$ est réduit signifie donc que
$X=X_{\mathrm{red}}$.

\begin{proposition}[5.1.4]
\label{I.5.1.4-fr}
Pour que le spectre premier d'un anneau $A$ soit un préschéma réduit (resp.
intègre) (\hyperref[I.2.1.7-fr]{2.1.7}), il faut et il suffit que $A$ soit un
anneau réduit (resp. intègre).
\end{proposition}

En effet, il résulte aussitôt de (\hyperref[I.5.1.1-fr]{5.1.1}) que la
condition $\mathscr{N}=(0)$ est nécessaire et suffisante pour que
$X=\operatorname{Spec}(A)$ soit réduit~; l'assertion relative aux anneaux
intègres est alors une conséquence de
(\hyperref[I.1.1.13-fr]{1.1.13}).

Comme tout anneau de fractions $\ne\{0\}$ d'un anneau intègre est intègre, il
résulte de (\hyperref[I.5.1.4-fr]{5.1.4}) que pour tout préschéma
\emph{localement intègre} $X$, $\mathscr{O}_x$ est un anneau \emph{intègre}
pour tout $x\in X$. La réciproque est vraie lorsque l'espace sous-jacent à
$X$ est \emph{localement noethérien}~: en effet $X$ est alors réduit, et si
$U$ est un ouvert affine de $X$, qui soit un espace noethérien, $U$ n'a qu'un
nombre fini de composantes irréductibles, donc son anneau $A$ n'a qu'un
nombre fini d'idéaux premiers minimaux
(\hyperref[I.1.1.14-fr]{1.1.14}). Si deux de ces composantes $U_i$ avaient un
point commun $x$, $\mathscr{O}_x$ aurait au moins deux idéaux premiers
minimaux distincts, donc ne serait pas intègre~; les $U_i$ sont par suite des
ouverts deux à deux disjoints, et chacun d'eux est donc intègre.

\begin{env}[5.1.5]
\label{I.5.1.5-fr}
Soit $f=(\psi,\theta):X\to Y$ un morphisme de préschémas~; l'homomor-
\oldpage[I]{129}
phisme $\theta_x^\sharp:\mathscr{O}_{\psi(x)}\to\mathscr{O}_x$ applique tout
élément nilpotent de $\mathscr{O}_{\psi(x)}$ sur un élément nilpotent de
$\mathscr{O}_x$~; par passage aux quotients, on déduit donc de $\theta^\sharp$
un homomorphisme
\[
  \omega:\psi^*(\mathscr{O}_Y/\mathscr{N}_Y)
  \longrightarrow \mathscr{O}_X/\mathscr{N}_X
\]
il est clair que pour tout $x\in X$,
$\omega_x:\mathscr{O}_{\psi(x)}/\mathscr{N}_{\psi(x)}
\to\mathscr{O}_x/\mathscr{N}_x$ est un homomorphisme local, donc
$(\psi,\omega^b)$ est un morphisme de préschémas
$X_{\mathrm{red}}\to Y_{\mathrm{red}}$, que nous noterons $f_{\mathrm{red}}$
et appellerons le morphisme \emph{réduit} associé à $f$. Il est immédiat que
pour deux morphismes $f:X\to Y$, $g:Y\to Z$, on a
$(g\circ f)_{\mathrm{red}}=g_{\mathrm{red}}\circ f_{\mathrm{red}}$, donc on a
défini $X_{\mathrm{red}}$ comme \emph{foncteur covariant} en $X$.

La définition précédente montre que le diagramme
\[
  \xymatrix{
    X_{\mathrm{red}}\ar[r]^{f_{\mathrm{red}}}\ar[d] &
    Y_{\mathrm{red}}\ar[d]\\
    X\ar[r]_f & Y
  }
\]
est commutatif, les flèches verticales étant les morphismes d'injection~; en
d'autres termes, $X_{\mathrm{red}}\to X$ est un morphisme \emph{fonctoriel}.
On notera en particulier que si $X$ est \emph{réduit}, tout morphisme
$f:X\to Y$ se factorise en
$X\xrightarrow{f_{\mathrm{red}}}Y_{\mathrm{red}}\to Y$~; en d'autres termes,
$f$ est \emph{majoré} par le morphisme d'injection
$Y_{\mathrm{red}}\to Y$.
\end{env}

\begin{proposition}[5.1.6]
\label{I.5.1.6-fr}
Soit $f:X\to Y$ un morphisme~; si $f$ est surjectif (resp. radiciel, une
immersion, une immersion fermée, une immersion ouverte, une immersion locale,
un isomorphisme local), il en est de même de $f_{\mathrm{red}}$. Inversement,
si $f_{\mathrm{red}}$ est surjectif (resp. radiciel), il en est de même de
$f$.
\end{proposition}

La proposition est triviale si $f$ est surjectif~; si $f$ est radiciel, elle
résulte de ce que pour tout $x\in X$, le corps $k(x)$ est le même pour les
préschémas $X$ et $X_{\mathrm{red}}$
(\hyperref[I.3.5.8-fr]{3.5.8}). Enfin, si $f=(\psi,\theta)$ est une immersion,
une immersion fermée ou une immersion locale (resp. une immersion ouverte, ou
un isomorphisme local), la proposition résulte de ce que si
$\theta_x^\sharp$ est surjectif (resp. bijectif), il en est de même de
l'homomorphisme obtenu en passant aux quotients par les nilradicaux de
$\mathscr{O}_{\psi(x)}$ et $\mathscr{O}_x$
(\hyperref[I.5.1.2-fr]{5.1.2} et \hyperref[I.4.2.2-fr]{4.2.2})
(cf. (\hyperref[I.5.5.12-fr]{5.5.12})).

\begin{proposition}[5.1.7]
\label{I.5.1.7-fr}
Si $X$, $Y$ sont deux $S$-préschémas, les préschémas
$X_{\mathrm{red}}\times_{S_{\mathrm{red}}}Y_{\mathrm{red}}$ et
$X_{\mathrm{red}}\times_S Y_{\mathrm{red}}$ sont identiques, et
s'identifient canoniquement à un sous-préschéma de $X\times_S Y$ ayant même
espace sous-jacent que ce produit.
\end{proposition}

L'identification canonique de $X_{\mathrm{red}}\times_S Y_{\mathrm{red}}$ à
un sous-préschéma de $X\times_S Y$ ayant même espace sous-jacent résulte de
(\hyperref[I.4.3.1-fr]{4.3.1}). D'autre part, si $\varphi$ et $\psi$ sont les
morphismes structuraux $X_{\mathrm{red}}\to S$, $Y_{\mathrm{red}}\to S$, ils
se factorisent par $S_{\mathrm{red}}$
(\hyperref[I.5.1.5-fr]{5.1.5}), et comme $S_{\mathrm{red}}\to S$ est un
monomorphisme, la première assertion résulte de
(\hyperref[I.3.2.4-fr]{3.2.4}).

\begin{corollary}[5.1.8]
\label{I.5.1.8-fr}
Les préschémas $(X\times_S Y)_{\mathrm{red}}$ et
$(X_{\mathrm{red}}\times_{S_{\mathrm{red}}}Y_{\mathrm{red}})_{\mathrm{red}}$
s'identifient canoniquement.
\end{corollary}

Cela résulte de (\hyperref[I.5.1.2-fr]{5.1.2} et
\hyperref[I.5.1.7-fr]{5.1.7}).

On notera que si $X$ et $Y$ sont des préschémas réduits, il n'en est pas
nécessairement de même de $X\times_S Y$, car le produit tensoriel de deux
algèbres réduites peut avoir des éléments nilpotents.

\begin{proposition}[5.1.9]
\label{I.5.1.9-fr}
\oldpage[I]{130}
Soient $X$ un préschéma, $\mathscr{J}$ un faisceau quasi-cohérent d'idéaux de
$\mathscr{O}_X$ tel que $\mathscr{J}^n=0$ pour un entier $n>0$. Soit $X_0$ le
sous-préschéma fermé $(X,\mathscr{O}_X/\mathscr{J})$ de $X$~; pour que $X$
soit un schéma affine, il faut et il suffit que $X_0$ le soit.
\end{proposition}

La condition étant évidemment nécessaire, prouvons qu'elle est suffisante.
Si on pose $X_k=(X,\mathscr{O}_X/\mathscr{J}^{k+1})$, tout revient à prouver
par récurrence sur $k$ que les $X_k$ sont affines, donc on est ramené au cas
où $\mathscr{J}^2=0$. Posons
\[
  A=\Gamma(X,\mathscr{O}_X),\qquad
  A_0=\Gamma(X_0,\mathscr{O}_{X_0})
      =\Gamma(X,\mathscr{O}_X/\mathscr{J}).
\]
On déduit de l'homomorphisme canonique
$\mathscr{O}_X\to\mathscr{O}_X/\mathscr{J}$ un homomorphisme d'anneaux
$\varphi:A\to A_0$. Nous verrons ci-dessous que $\varphi$ est
\emph{surjectif}, de sorte que la suite
\[
  0\to\Gamma(X,\mathscr{J})\to\Gamma(X,\mathscr{O}_X)
  \to\Gamma(X,\mathscr{O}_X/\mathscr{J})\to 0
  \tag{5.1.9.1}\label{I.5.1.9.1-fr}
\]
est \emph{exacte}. Supposons ce point établi, et montrons que cela entraîne la
proposition. Notons que $\mathfrak{K}=\Gamma(X,\mathscr{J})$ est un idéal de
carré nul dans $A$, et est donc un module sur $A_0=A/\mathfrak{K}$. Par
hypothèse, on a $X_0=\operatorname{Spec}(A_0)$, et comme les espaces
topologiques sous-jacents $X_0$ et $X$ sont identiques,
$\mathfrak{K}=\Gamma(X_0,\mathscr{J})$~; par ailleurs, comme
$\mathscr{J}^2=0$, $\mathscr{J}$ est un
$(\mathscr{O}_X/\mathscr{J})$-Module quasi-cohérent, donc on a
$\mathscr{J}\simeq\widetilde{\mathfrak{K}}$ et
$\mathfrak{K}_x=\mathscr{J}_x$ pour tout $x\in X_0$
(\hyperref[I.1.4.1-fr]{1.4.1}).

Cela étant, soit $X'=\operatorname{Spec}(A)$, et considérons le morphisme
$f=(\psi,\theta):X\to X'$ de préschémas correspondant à l'application
identique $A\to\Gamma(X,\mathscr{O}_X)$
(\hyperref[I.2.2.4-fr]{2.2.4}). Pour tout ouvert affine $V$ dans $X$, le
diagramme
\[
  \xymatrix{
    A\ar[r]\ar[d] & \Gamma(V,\mathscr{O}_X|V)\ar[d]\\
    A_0=A/\mathfrak{K}\ar[r] & \Gamma(V,\mathscr{O}_{X_0}|V)
  }
\]
est commutatif, d'où on conclut que le diagramme
\[
  \xymatrix{
    X' & X\ar[l]_f\\
    X'_0\ar[u]^{j'} & X_0\ar[l]^{f_0}\ar[u]_j
  }
\]
est commutatif, $X'_0$ étant le sous-préschéma fermé de $X'$ défini par le
faisceau quasi-cohérent d'idéaux $\widetilde{\mathfrak{K}}$, $j$, $j'$ les
morphismes d'injection canoniques. Mais puisque $X_0$ est affine, $f_0$ est un
isomorphisme et comme les applications continues sous-jacentes à $j$ et $j'$
sont les applications identiques, on voit tout d'abord que
$\psi:X\to X'$ est un homéomorphisme. En outre, la relation
$\mathfrak{K}_x=\mathscr{J}_x$ montre que la restriction de
$\theta^\sharp:\psi^*(\mathscr{O}_{X'})\to\mathscr{O}_X$ est un
\emph{isomorphisme} de $\psi^*(\widetilde{\mathfrak{K}})$ sur $\mathscr{J}$~;
d'autre part, par passage aux quotients, $\theta^\sharp$ donne un
\emph{isomorphisme}
$\psi^*(\mathscr{O}_{X'}/\widetilde{\mathfrak{K}})
\to\mathscr{O}_X/\mathscr{J}$, puisque $f_0$ est un isomorphisme~; on en
conclut aussitôt par le lemme des 5 (M, I, 1.1) que $\theta^\sharp$ est
lui-même un isomorphisme, donc que $f$ est un \emph{isomorphisme}, et par suite
que $X$ est affine.

Tout revient donc à démontrer l'exactitude de
(\hyperref[I.5.1.9.1-fr]{5.1.9.1}), ce qui résultera de

\oldpage[I]{131}
$H^1(X,\mathscr{J})=0$. Or,
$H^1(X,\mathscr{J})=H^1(X_0,\mathscr{J})$ et nous avons vu que
$\mathscr{J}$ est un $\mathscr{O}_{X_0}$-Module quasi-cohérent. Notre
assertion résultera donc du

\begin{lemma}[5.1.9.2]
\label{I.5.1.9.2-fr}
Si $Y$ est un schéma affine et $\mathscr{F}$ un $\mathscr{O}_Y$-Module
quasi-cohérent, on a $H^1(Y,\mathscr{F})=0$.
\end{lemma}

Ce lemme sera démontré au chap. III, § 1, comme conséquence du théorème plus
général que $H^i(Y,\mathscr{F})=0$ pour tout $i>0$. Pour en donner une
démonstration indépendante, observons que $H^1(Y,\mathscr{F})$ s'identifie au
module $\operatorname{Ext}^1_{\mathscr{O}_Y}
(Y;\mathscr{O}_Y,\mathscr{F})$ des classes d'extensions du
$\mathscr{O}_Y$-Module $\mathscr{O}_Y$ par le $\mathscr{O}_Y$-Module
$\mathscr{F}$ (T, 4.2.3)~; tout revient donc à prouver qu'une telle extension
$\mathscr{G}$ est triviale. Or, pour tout $y\in Y$, il y a un voisinage $V$
de $y$ dans $Y$ tel que $\mathscr{G}|V$ soit isomorphe à
$\mathscr{F}|Y\oplus\mathscr{O}_Y|V$
(\hyperref[0.5.4.9-fr]{0, 5.4.9})~; on en conclut que $\mathscr{G}$ est un
$\mathscr{O}_Y$-Module quasi-cohérent. Si $A$ est l'anneau de $Y$, on a donc
$\mathscr{F}=\widetilde{M}$, $\mathscr{G}=\widetilde{N}$, où $M$ et $N$ sont
des $A$-modules, et par hypothèse $N$ est une extension du $A$-module $A$ par
le $A$-module $M$ (\hyperref[I.1.3.11-fr]{1.3.11}). Comme cette extension est
nécessairement triviale, le lemme est démontré, et par suite aussi
(\hyperref[I.5.1.9-fr]{5.1.9}).

\begin{corollary}[5.1.10]
\label{I.5.1.10-fr}
Soit $X$ un préschéma tel que $\mathscr{N}_X$ soit nilpotent. Pour que $X$
soit un schéma affine, il faut et il suffit que $X_{\mathrm{red}}$ le soit.
\end{corollary}

\subsection{Existence d'un sous-préschéma d'espace sous-jacent donné.}
\label{subsection:I.5.2-fr}

\begin{proposition}[5.2.1]
\label{I.5.2.1-fr}
Pour tout sous-espace localement fermé $Y$ de l'espace sous-jacent à un
préschéma $X$, il existe un sous-préschéma réduit et un seul de $X$ ayant $Y$
pour espace sous-jacent.
\end{proposition}

L'unicité résultant de (\hyperref[I.5.1.2-fr]{5.1.2}), il reste à prouver
l'existence du sous-préschéma en question.

Si $X$ est affine d'anneau $A$, et $Y$ fermé dans $X$, la proposition est
immédiate~: $\mathfrak{j}(Y)$ est le plus grand idéal
$\mathfrak{a}\subset A$ tel que $V(\mathfrak{a})=Y$, et il est égal à sa
racine (\hyperref[I.1.1.4-fr]{1.1.4 (i)}), donc
$A/\mathfrak{j}(Y)$ est un anneau réduit.

Dans le cas général, pour tout ouvert affine $U\subset X$ tel que $U\cap Y$
soit fermé dans $U$, considérons le sous-préschéma fermé $Y_U$ de $U$ défini
par le faisceau d'idéaux associé à l'idéal $\mathfrak{j}(U\cap Y)$ de $A(U)$,
qui est réduit. Montrons que, si $V$ est un ouvert affine de $X$ contenu dans
$U$, $Y_V$ est \emph{induit} par $Y_U$ sur $V\cap Y$~; or, ce préschéma
induit est un sous-préschéma fermé de $V$ qui est réduit et a $V\cap Y$ comme
espace sous-jacent~; l'unicité de $Y_V$ entraîne donc notre assertion.

\begin{proposition}[5.2.2]
\label{I.5.2.2-fr}
Soient $X$ un préschéma réduit, $f:X\to Y$ un morphisme, $Z$ un
sous-préschéma fermé de $Y$ tel que $f(X)\subset Z$~; alors $f$ se factorise
en $X\xrightarrow{g}Z\xrightarrow{j}Y$, où $j$ est le morphisme d'injection.
\end{proposition}

Il résulte de l'hypothèse que le sous-préschéma fermé $f^{-1}(Z)$ de $X$ a
pour espace sous-jacent $X$ tout entier
(\hyperref[I.4.4.1-fr]{4.4.1})~; comme $X$ est réduit, ce sous-préschéma fermé
coïncide avec $X$ (\hyperref[I.5.1.2-fr]{5.1.2}), et la proposition résulte
donc de (\hyperref[I.4.4.1-fr]{4.4.1}).

\begin{corollary}[5.2.3]
\label{I.5.2.3-fr}
Soit $X$ un sous-préschéma réduit d'un préschéma $Y$~; si $Z$ est le
sous-préschéma fermé réduit de $Y$ ayant pour espace sous-jacent
$\overline{X}$, $X$ est un sous-préschéma induit sur un ouvert de $Z$.
\end{corollary}

\oldpage[I]{132}
Il y a en effet un ouvert $U$ de $Y$ tel que
$X=U\cap\overline{X}$~; comme $X$ est un sous-préschéma réduit de $Z$ en
vertu de (\hyperref[I.5.2.2-fr]{5.2.2}), le sous-préschéma $X$ est induit par
$Z$ sur le sous-espace ouvert $X$ en vertu de l'unicité
(\hyperref[I.5.2.1-fr]{5.2.1}).

\begin{corollary}[5.2.4]
\label{I.5.2.4-fr}
Soient $f:X\to Y$ un morphisme, $X'$ (resp. $Y'$) un sous-préschéma fermé de
$X$ (resp. $Y$) défini par un faisceau d'idéaux quasi-cohérent $\mathscr{J}$
(resp. $\mathscr{K}$) de $\mathscr{O}_X$ (resp. $\mathscr{O}_Y$). Supposons
que $X'$ soit réduit et que $f(X')\subset Y'$. Alors on a
$f^*(\mathscr{K})\mathscr{O}_X\subset\mathscr{J}$.
\end{corollary}

Comme la restriction de $f$ à $X'$ se factorise en
$X'\to Y'\to Y$ d'après (\hyperref[I.5.2.2-fr]{5.2.2}), il suffit d'appliquer
(\hyperref[I.4.4.6-fr]{4.4.6}).

\subsection{Diagonale; graphe d'un morphisme.}
\label{subsection:I.5.3-fr}

\begin{env}[5.3.1]
\label{I.5.3.1-fr}
Soit $X$ un $S$-préschéma~; on appelle \emph{morphisme diagonal} de $X$ dans
$X\times_S X$, et on note $\Delta_{X|S}$, ou $\Delta_X$, ou même $\Delta$ si
aucune confusion n'est possible, le $S$-morphisme $(1_X,1_X)_S$, autrement
dit, l'unique $S$-morphisme $\Delta_X$ tel que
\[
  p_1\circ\Delta_X=p_2\circ\Delta_X=1_X
  \tag{5.3.1.1}\label{I.5.3.1.1-fr}
\]
en désignant par $p_1$, $p_2$ les projections de $X\times_S X$
(déf. \hyperref[I.3.2.1-fr]{3.2.1}). Si $f:T\to X$, $g:T\to Y$ sont deux
$S$-morphismes, on vérifie aussitôt que
\[
  (f,g)_S=(f\times_S g)\circ\Delta_{T|S}
  \tag{5.3.1.2}\label{I.5.3.1.2-fr}
\]

Le lecteur observera que la définition précédente et les résultats énoncés
dans les n\textsuperscript{os} (\hyperref[I.5.3.1-fr]{5.3.1}) à (5.3.8) sont
valables dans toute catégorie, \emph{pourvu que les produits qui y figurent
existent dans cette catégorie}.
\end{env}

\begin{proposition}[5.3.2]
\label{I.5.3.2-fr}
Soient $X$, $Y$ deux $S$-préschémas~; si on identifie canoniquement le produit
$(X\times Y)\times(X\times Y)$ à $(X\times X)\times(Y\times Y)$, le morphisme
$\Delta_{X\times Y}$ s'identifie à $\Delta_X\times\Delta_Y$.
\end{proposition}

En effet, si $p_1$, $q_1$ sont les premières projections
$X\times X\to X$, $Y\times Y\to Y$, la première projection
$(X\times Y)\times(X\times Y)\to X\times Y$ s'identifie à $p_1\times q_1$,
et l'on a
\[
  (p_1\times q_1)\circ(\Delta_X\times\Delta_Y)
  =(p_1\circ\Delta_X)\times(q_1\circ\Delta_Y)=1_{X\times Y}
\]
même raisonnement pour les secondes projections.

\begin{corollary}[5.3.4]
\label{I.5.3.4-fr}
Pour toute extension $S'\to S$ du préschéma de base,
$\Delta_{X_{(S')}}$ s'identifie canoniquement à $(\Delta_X)_{(S')}$.
\end{corollary}

Il suffit de remarquer que $(X\times_S X)_{(S')}$ s'identifie canoniquement à
$X_{(S')}\times_{S'}X_{(S')}$
(\hyperref[I.3.3.10-fr]{3.3.10}).

\begin{proposition}[5.3.5]
\label{I.5.3.5-fr}
Soient $X$, $Y$ deux $S$-préschémas, $\varphi:S\to T$ un morphisme, faisant de
tout $S$-préschéma un $T$-préschéma. Soient $f:X\to S$, $Y\to S$ les
morphismes structuraux, $p$, $q$ les projections de $X\times_S Y$,
$\pi=f\circ p=g\circ q$ le morphisme structural $X\times_S Y\to S$. Alors, le
diagramme
\[
  \xymatrix{
    X\times_S Y\ar[r]^{(p,q)_T}\ar[d]_\pi &
    X\times_T Y\ar[d]^{f\times_T g}\\
    S\ar[r]_{\Delta_{S|T}} &
    S\times_T S
  }
  \tag{5.3.5.1}\label{I.5.3.5.1-fr}
\]
\oldpage[I]{133}
est commutatif, et identifie $X\times_S Y$ au produit des
$(S\times_T S)$-préschémas $S$ et $X\times_T Y$, les projections
s'identifiant à $\pi$ et $(p,q)_T$.
\end{proposition}

En vertu de (\hyperref[I.3.4.3-fr]{3.4.3}), on est ramené à démontrer la
proposition correspondante \emph{dans la catégorie des ensembles}, en
remplaçant $X$, $Y$, $S$ par $X(Z)_T$, $Y(Z)_T$, $S(Z)_T$, $Z$ étant un
$T$-préschéma arbitraire. Mais pour la catégorie des ensembles, la vérification
est immédiate et laissée au lecteur.

\begin{corollary}[5.3.6]
\label{I.5.3.6-fr}
Le morphisme $(p,q)_T$ s'identifie (en posant $P=S\times_T S$) à
$1_{X\times_T Y}\times_P\Delta_S$.
\end{corollary}

Cela résulte de (\hyperref[I.5.3.5-fr]{5.3.5}) et
(\hyperref[I.3.3.4-fr]{3.3.4}).

\begin{corollary}[5.3.7]
\label{I.5.3.7-fr}
Si $f:X\to Y$ est un $S$-morphisme, le diagramme
\[
  \xymatrix{
    X\ar[r]^{(1_X,f)_S}\ar[d]_f &
    X\times_S Y\ar[d]^{f\times_S 1_Y}\\
    Y\ar[r]_{\Delta_Y} &
    Y\times_S Y
  }
\]
est commutatif, et identifie $X$ au produit des
$(Y\times_S Y)$-préschémas $Y$ et $X\times_S Y$.
\end{corollary}

Il suffit d'appliquer (\hyperref[I.5.3.5-fr]{5.3.5}) en remplaçant $S$ par
$Y$ et $T$ par $S$, et remarquant que $X\times_Y Y=X$
(\hyperref[I.3.3.3-fr]{3.3.3}).

\begin{proposition}[5.3.8]
\label{I.5.3.8-fr}
Pour que $f:X\to Y$ soit un monomorphisme de préschémas, il faut et il suffit
que $\Delta_{X|Y}$ soit un isomorphisme de $X$ sur $X\times_Y X$.
\end{proposition}

En effet, dire que $f$ est un monomorphisme signifie que pour tout
$Y$-préschéma $Z$, l'application correspondante
$f':X(Z)_Y\to Y(Z)_Y$ est une injection, et comme $Y(Z)_Y$ est réduit à un
élément, cela signifie qu'il en est de même de $X(Z)_Y$. Mais cela s'exprime
aussi en disant que $X(Z)_Y\times X(Z)_Y$ est canoniquement isomorphe à
$X(Z)_Y$, et le premier de ces ensembles étant $(X\times_Y X)(Z)_Y$
(\hyperref[I.3.4.3.1-fr]{3.4.3.1}), cela signifie que $\Delta_{X|Y}$ est un
isomorphisme.

\begin{proposition}[5.3.9]
\label{I.5.3.9-fr}
Le morphisme diagonal $\Delta_X$ est une immersion de $X$ dans
$X\times_S X$.
\end{proposition}

En effet, comme les applications continues $p_1$ et $\Delta_X$ des espaces
sous-jacents sont telles que $p_1\circ\Delta_X$ soit l'identité, $\Delta_X$
est un homéomorphisme de $X$ sur $\Delta_X(X)$. De même, l'homomorphisme
composé
$\mathscr{O}_x\to\mathscr{O}_{\Delta_X(x)}\to\mathscr{O}_x$ des
homomorphismes correspondant à $p_1$ et à $\Delta_X$ étant l'identité,
l'homomorphisme correspondant à $\Delta_X$ est surjectif~; la proposition
résulte donc de (\hyperref[I.4.2.2-fr]{4.2.2}).

On dit que le sous-préschéma de $X\times_S X$ associé à l'immersion
$\Delta_X$ (\hyperref[I.4.2.1-fr]{4.2.1}) est \emph{la diagonale} de
$X\times_S X$.

\begin{corollary}[5.3.10]
\label{I.5.3.10-fr}
Sous les hypothèses de (\hyperref[I.5.3.5-fr]{5.3.5}), $(p,q)_T$ est une
immersion.
\end{corollary}

Cela résulte de (\hyperref[I.5.3.6-fr]{5.3.6}) et de
(\hyperref[I.4.3.1-fr]{4.3.1}).

On dit (sous les hypothèses de (\hyperref[I.5.3.5-fr]{5.3.5})) que $(p,q)_T$
est l'\emph{immersion canonique} de $X\times_S Y$ dans $X\times_T Y$.

\begin{corollary}[5.3.11]
\label{I.5.3.11-fr}
Soient $X$, $Y$ deux $S$-préschémas, $f:X\to Y$ un $S$-morphisme~; alors le
morphisme graphe $\Gamma_f=(1_X,f)_S$ de $f$
(\hyperref[I.3.3.14-fr]{3.3.14}) est une immersion de $X$ dans
$X\times_S Y$.
\end{corollary}

C'est le cas particulier du cor. (\hyperref[I.5.3.10-fr]{5.3.10}) où l'on
remplace $S$ par $Y$ et $T$ par $S$
(cf. (\hyperref[I.5.3.7-fr]{5.3.7})).

\oldpage[I]{134}
Le sous-préschéma de $X\times_S Y$ associé à l'immersion $\Gamma_f$
(\hyperref[I.4.2.1-fr]{4.2.1}) est appelé le \emph{graphe} du morphisme $f$~;
les sous-préschémas de $X\times_S Y$ qui sont des graphes de morphismes
$X\to Y$ se caractérisent par le fait que la restriction à un tel
sous-préschéma $G$ de la projection $p_1:X\times_S Y\to X$ est un
\emph{isomorphisme} $g$ de $G$ sur $X$~: $G$ est alors le graphe du morphisme
$p_2\circ g^{-1}$, où $p_2$ est la projection $X\times_S Y\to Y$.

Lorsqu'on prend en particulier $X=S$, les $S$-morphismes $S\to Y$, qui ne
sont autres que les $S$-sections de $Y$ (\hyperref[I.2.5.5-fr]{2.5.5}) sont
égaux à leurs morphismes graphes~; les sous-préschémas de $Y$ qui sont des
graphes de $S$-sections (autrement dit, ceux qui sont isomorphes à $S$ par la
restriction du morphisme structural $Y\to S$) s'appellent encore les
\emph{images} de ces sections, ou, par abus de langage, les
\emph{$S$-sections} de $Y$.

\begin{corollary}[5.3.12]
\label{I.5.3.12-fr}
Les hypothèses et notations étant celles de
(\hyperref[I.5.3.11-fr]{5.3.11}), pour tout morphisme $g:S'\to S$, soit $f'$
l'image réciproque de $f$ par $g$ (\hyperref[I.3.3.7-fr]{3.3.7})~; alors
$\Gamma_{f'}$ est l'image réciproque de $\Gamma_f$ par $g$.
\end{corollary}

C'est un cas particulier de la formule
(\hyperref[I.3.3.10.1-fr]{3.3.10.1}).

\begin{corollary}[5.3.13]
\label{I.5.3.13-fr}
Soient $f:X\to Y$, $g:Y\to Z$ deux morphismes~; si $g\circ f$ est une
immersion (resp. une immersion locale), il en est de même de $f$.
\end{corollary}

En effet, $f$ se factorise en
\[
  X\xrightarrow{\Gamma_f}X\times_Z Y\xrightarrow{p_2}Y.
\]
D'autre part, $p_2$ s'identifie à $(g\circ f)\times_Z 1_Y$
(\hyperref[I.3.3.4-fr]{3.3.4})~; si $g\circ f$ est une immersion (resp. une
immersion locale), il en est de même de $p_2$
(\hyperref[I.4.3.1-fr]{4.3.1} et \hyperref[I.4.5.5-fr]{4.5.5}), et comme
$\Gamma_f$ est une immersion (\hyperref[I.5.3.11-fr]{5.3.11}), on conclut
par (\hyperref[I.4.2.4-fr]{4.2.4}) (resp.
(\hyperref[I.4.5.5-fr]{4.5.5})).

\begin{corollary}[5.3.14]
\label{I.5.3.14-fr}
Soient $j:X\to Y$, $g:X\to Z$ deux $S$-morphismes. Si $j$ est une immersion
(resp. une immersion locale), il en est de même de $(j,g)_S$.
\end{corollary}

En effet, si $p:Y\times_S Z\to Y$ est la première projection, on a
$j=p\circ(j,g)_S$, et il suffit d'appliquer
(\hyperref[I.5.3.13-fr]{5.3.13}).

\begin{proposition}[5.3.15]
\label{I.5.3.15-fr}
Si $f:X\to Y$ est un $S$-morphisme, le diagramme
\[
  \xymatrix{
    X\ar[r]^{\Delta_X}\ar[d]_f &
    X\times_S X\ar[d]^{f\times_S f}\\
    Y\ar[r]_{\Delta_Y} &
    Y\times_S Y
  }
  \tag{5.3.15.1}\label{I.5.3.15.1-fr}
\]
est commutatif (en d'autres termes $\Delta_X$ est un morphisme fonctoriel
dans la catégorie des préschémas).
\end{proposition}

La vérification est immédiate et laissée au lecteur.

\begin{corollary}[5.3.16]
\label{I.5.3.16-fr}
Si $X$ est un sous-préschéma de $Y$, la diagonale $\Delta_X(X)$ s'identifie à
un sous-préschéma de $\Delta_Y(Y)$, dont l'espace sous-jacent s'identifie à
\[
  \Delta_Y(Y)\cap p_1^{-1}(X)=\Delta_Y(Y)\cap p_2^{-1}(X)
\]
($p_1$, $p_2$ projections de $Y\times_S Y$).
\end{corollary}

Appliquons (\hyperref[I.5.3.15-fr]{5.3.15}) au morphisme d'injection
$f:X\to Y$~; on sait alors que $f\times_S f$ est une immersion, identifiant
l'espace sous-jacent à $X\times_S X$ au sous-espace
$p_1^{-1}(X)\cap p_2^{-1}(X)$ de $Y\times_S Y$
(\hyperref[I.4.3.1-fr]{4.3.1})~; en outre, si
$z\in\Delta_Y(Y)\cap p_1^{-1}(X)$, on a $z=\Delta_Y(y)$

\oldpage[I]{135}
et $y=p_1(z)\in X$, donc $y=f(y)$, et $z=\Delta_Y(f(y))$ appartient à
$\Delta_X(X)$ en vertu de la commutativité du diagramme
(\hyperref[I.5.3.15.1-fr]{5.3.15.1}).

\begin{corollary}[5.3.17]
\label{I.5.3.17-fr}
Soient $f_1:Y\to X$, $f_2:Y\to X$ deux $S$-morphismes, $y$ un point de $Y$
tel que $f_1(y)=f_2(y)=x$ et que les homomorphismes $k(x)\to k(y)$
correspondant à $f_1$ et $f_2$ soient identiques. Alors, si
$f=(f_1,f_2)_S$, le point $f(y)$ appartient à la diagonale
$\Delta_{X|S}(X)$.
\end{corollary}

Les deux homomorphismes $k(x)\to k(y)$ correspondant à $f_i$ ($i=1,2$)
définissent deux $S$-morphismes
$g_i:\operatorname{Spec}(k(y))\to\operatorname{Spec}(k(x))$ tels que les
diagrammes
\[
  \xymatrix{
    \operatorname{Spec}(k(y))\ar[r]^{g_i}\ar[d] &
    \operatorname{Spec}(k(x))\ar[d]\\
    Y\ar[r]_{f_i} & X
  }
\]
soient commutatifs. Le diagramme
\[
  \xymatrix{
    \operatorname{Spec}(k(y))\ar[rr]^{(g_1,g_2)_S}\ar[d] & &
    \operatorname{Spec}(k(x))\times_S\operatorname{Spec}(k(x))\ar[d]\\
    Y\ar[rr]_{(f_1,f_2)_S} & & X\times_S X
  }
\]
est donc aussi commutatif. Or il résulte de l'égalité $g_1=g_2$ que l'image
par $(g_1,g_2)_S$ de l'unique point de $\operatorname{Spec}(k(y))$
appartient à la diagonale de
$\operatorname{Spec}(k(x))\times_S\operatorname{Spec}(k(x))$~; la conclusion
résulte donc de (\hyperref[I.5.3.15-fr]{5.3.15}).

\subsection{Morphismes et préschémas séparés.}
\label{subsection:I.5.4-fr}

\begin{definition}[5.4.1]
\label{I.5.4.1-fr}
On dit qu'un morphisme de préschémas $f:X\to Y$ est \emph{séparé} si le
morphisme diagonal $X\to X\times_Y X$ est une \emph{immersion fermée}~; on
dit alors aussi que $X$ est un \emph{préschéma séparé au-dessus de $Y$}, ou un
\emph{$Y$-schéma}. On dit qu'un préschéma $X$ est séparé s'il est séparé
au-dessus de $\operatorname{Spec}(\mathbf{Z})$~; on dit aussi alors que $X$
est un \emph{schéma} (cf. (\hyperref[I.5.5.7-fr]{5.5.7})).
\end{definition}

En vertu de (\hyperref[I.5.3.9-fr]{5.3.9}), pour que $X$ soit séparé au-dessus
de $Y$, il faut et il suffit que $\Delta_X(X)$ soit un \emph{sous-espace
fermé} de l'espace sous-jacent à $X\times_Y X$.

\begin{proposition}[5.4.2]
\label{I.5.4.2-fr}
Soit $S\to T$ un morphisme séparé. Si $X$ et $Y$ sont deux $S$-préschémas,
l'immersion canonique $X\times_S Y\to X\times_T Y$
(\hyperref[I.5.3.10-fr]{5.3.10}) est fermée.
\end{proposition}

En effet, si on se reporte au diagramme
(\hyperref[I.5.3.5.1-fr]{5.3.5.1}), on voit que $(p,q)_T$ peut être considéré
comme obtenu à partir de $\Delta_{S|T}$ par l'extension
$f\times_T g:X\times_T Y\to S\times_T S$ du préschéma de base $S\times_T S$~;
la proposition résulte alors de (\hyperref[I.4.3.2-fr]{4.3.2}).

\begin{corollary}[5.4.3]
\label{I.5.4.3-fr}
Soient $Y$ un $S$-schéma, $f:X\to Y$ un $S$-morphisme. Alors le morphisme
graphe $\Gamma_f:X\to X\times_S Y$ (\hyperref[I.5.3.11-fr]{5.3.11}) est une
immersion fermée.
\end{corollary}

C'est le cas particulier de (\hyperref[I.5.4.2-fr]{5.4.2}) où l'on remplace
$S$ par $Y$ et $T$ par $S$.

\begin{corollary}[5.4.4]
\label{I.5.4.4-fr}
Soient $f:X\to Y$, $g:Y\to Z$ deux morphismes, $g$ étant séparé. Si $g\circ f$
est une immersion fermée, il en est de même de $f$.
\end{corollary}

La démonstration à partir de (\hyperref[I.5.4.3-fr]{5.4.3}) est la même que
celle de (\hyperref[I.5.3.13-fr]{5.3.13}) à partir de
(\hyperref[I.5.3.11-fr]{5.3.11}).

\oldpage[I]{136}
\begin{corollary}[5.4.5]
\label{I.5.4.5-fr}
Soient $Z$ un $S$-schéma, $j:X\to Y$, $g:X\to Z$ deux $S$-morphismes. Si $j$
est une immersion fermée, il en est de même de
$(j,g)_S:X\to Y\times_S Z$.
\end{corollary}

La démonstration à partir de (\hyperref[I.5.4.4-fr]{5.4.4}) est la même que
celle de (\hyperref[I.5.3.14-fr]{5.3.14}) à partir de
(\hyperref[I.5.3.13-fr]{5.3.13}).

\begin{corollary}[5.4.6]
\label{I.5.4.6-fr}
Si $X$ est un $S$-schéma, toute $S$-section de $X$
(\hyperref[I.2.5.5-fr]{2.5.5}) est une immersion fermée.
\end{corollary}

Si $\varphi:X\to S$ est le morphisme structural, $\psi:S\to X$ une
$S$-section de $X$, il suffit d'appliquer
(\hyperref[I.5.4.5-fr]{5.4.5}) à $\varphi\circ\psi=1_S$.

\begin{corollary}[5.4.7]
\label{I.5.4.7-fr}
Soient $S$ un préschéma intègre, $s$ son point générique, $X$ un $S$-schéma.
Si deux $S$-sections $f$, $g$ de $X$ sont telles que $f(s)=g(s)$, alors
$f=g$.
\end{corollary}

En effet, si $x=f(s)=g(s)$, les homomorphismes $k(x)\to k(s)$ correspondant à
$f$ et $g$ sont nécessairement identiques. Si $h=(f,g)_S$, on en déduit
(\hyperref[I.5.3.17-fr]{5.3.17}) que $h(s)$ appartient à la diagonale
$Z=\Delta_X(X)$~; mais comme $S=\overline{\{s\}}$ et que $Z$ est fermée par
hypothèse, on a $h(S)\subset Z$. Il résulte alors de
(\hyperref[I.5.2.2-fr]{5.2.2}) que $h$ se factorise en
$S\to Z\to X\times_S X$, et on en conclut que $f=g$ par définition de la
diagonale.

\begin{remark}[5.4.8]
\label{I.5.4.8-fr}
Si l'on suppose réciproquement que la conclusion de
(\hyperref[I.5.4.3-fr]{5.4.3}) est vérifiée lorsque $f=1_Y$, on en conclut
que $Y$ est séparé au-dessus de $S$~; de même, si on suppose que la conclusion
de (\hyperref[I.5.4.5-fr]{5.4.5}) s'applique aux deux morphismes
\[
  Y\xrightarrow{\Delta_Y}Y\times_Z Y\xrightarrow{p_1}Y,
\]
on en tire que $\Delta_Y$ est une immersion fermée, donc que $Y$ est séparé
au-dessus de $Z$~; enfin, la validité de la conclusion de
(\hyperref[I.5.4.6-fr]{5.4.6}) pour la $Y$-section $\Delta_Y$ du
$Y$-préschéma $Y\times_S Y\to Y$, implique que $Y$ est séparé au-dessus de
$S$.
\end{remark}

\subsection{Critères de séparation.}
\label{subsection:I.5.5-fr}

\begin{proposition}[5.5.1]
\label{I.5.5.1-fr}
\begin{enumerate}
  \item[\rm(i)] Tout monomorphisme de préschémas (et en particulier toute
    immersion) est un morphisme séparé.
  \item[\rm(ii)] Le composé de deux morphismes séparés est séparé.
  \item[\rm(iii)] Si $f:X\to X'$, $g:Y\to Y'$ sont deux $S$-morphismes
    séparés, $f\times_S g$ est séparé.
  \item[\rm(iv)] Si $f:X\to Y$ est un $S$-morphisme séparé, le
    $S'$-morphisme $f_{(S')}$ est séparé pour toute extension $S'\to S$ du
    préschéma de base.
  \item[\rm(v)] Si le composé $g\circ f$ de deux morphismes est séparé, $f$
    est séparé.
  \item[\rm(vi)] Pour qu'un morphisme $f$ soit séparé, il faut et il suffit
    que $f_{\mathrm{red}}$ (\hyperref[I.5.1.5-fr]{5.1.5}) le soit.
\end{enumerate}
\end{proposition}

(i) résulte aussitôt de (\hyperref[I.5.3.8-fr]{5.3.8}). Si $f:X\to Y$,
$g:Y\to Z$ sont deux morphismes, le diagramme
\[
  \xymatrix{
    X\ar[rr]^{\Delta_{X|Z}}\ar[dr]_{\Delta_{X|Y}} & &
    X\times_Z X\\
    & X\times_Y X\ar[ur]_j
  }
  \tag{5.5.1.1}\label{I.5.5.1.1-fr}
\]
où $j$ désigne l'immersion canonique (\hyperref[I.5.3.10-fr]{5.3.10}) est
commutatif, comme on le vérifie aussitôt. Si $f$ et $g$ sont séparés,
$\Delta_{X|Y}$ est une immersion fermée par définition, et $j$ est une
immersion fermée en vertu de (\hyperref[I.5.4.2-fr]{5.4.2}), donc
$\Delta_{X|Z}$ est une immersion fermée par
(\hyperref[I.4.2.4-fr]{4.2.4}), ce qui

\oldpage[I]{137}
prouve (ii). Vu (i) et (ii), (iii) et (iv) sont équivalentes
(\hyperref[I.3.5.1-fr]{3.5.1}), et il suffit de démontrer (iv). Or,
$X_{(S')}\times_{Y_{(S')}}X_{(S')}$ s'identifie canoniquement à
$(X\times_Y X)\times_Y Y_{(S')}$ en vertu de
(\hyperref[I.3.3.11-fr]{3.3.11}) et de
(\hyperref[I.3.3.9.1-fr]{3.3.9.1}), et on vérifie aussitôt que le morphisme
diagonal $\Delta_{X_{(S')}}$ s'identifie alors à
$\Delta_X\times_Y 1_{Y_{(S')}}$~; la proposition résulte donc de
(\hyperref[I.4.3.1-fr]{4.3.1}).

Pour établir (v), considérons, comme dans
(\hyperref[I.5.3.13-fr]{5.3.13}) la factorisation
$X\xrightarrow{\Gamma_f}X\times_Z Y\xrightarrow{p_2}Y$ de $f$, en remarquant
que $p_2=(g\circ f)\times_Z 1_Y$~; l'hypothèse que $g\circ f$ est séparé
entraîne que $p_2$ est séparé par (iii) et (i), et comme $\Gamma_f$ est une
immersion, $\Gamma_f$ est séparé par (i), donc $f$ est séparé par (ii).
Enfin, pour démontrer (vi), rappelons que les préschémas
$X_{\mathrm{red}}\times_{Y_{\mathrm{red}}}X_{\mathrm{red}}$ et
$X_{\mathrm{red}}\times_Y X_{\mathrm{red}}$ s'identifient canoniquement
(\hyperref[I.5.1.7-fr]{5.1.7})~; si on désigne par $j$ l'injection
$X_{\mathrm{red}}\to X$, le diagramme
\[
  \xymatrix{
    X_{\mathrm{red}}\ar[r]^{\Delta_{X_{\mathrm{red}}}}\ar[d]_j &
    X_{\mathrm{red}}\times_Y X_{\mathrm{red}}\ar[d]^{j\times_Y j}\\
    X\ar[r]_{\Delta_X} & X\times_Y X
  }
\]
est commutatif (\hyperref[I.5.3.15-fr]{5.3.15}), et la proposition résulte de
ce que les flèches verticales sont des homéomorphismes des espaces
sous-jacents (\hyperref[I.4.3.1-fr]{4.3.1}).

\begin{corollary}[5.5.2]
\label{I.5.5.2-fr}
Si $f:X\to Y$ est séparé, la restriction de $f$ à tout sous-préschéma de $X$
est séparée.
\end{corollary}

Cela résulte de (\hyperref[I.5.5.1-fr]{5.5.1, (i) et (ii)}).

\begin{corollary}[5.5.3]
\label{I.5.5.3-fr}
Si $X$, $Y$ sont deux $S$-préschémas tels que $Y$ soit séparé au-dessus de
$S$, $X\times_S Y$ est séparé au-dessus de $X$.
\end{corollary}

C'est un cas particulier de (\hyperref[I.5.5.1-fr]{5.5.1, (iv)}).

\begin{proposition}[5.5.4]
\label{I.5.5.4-fr}
Soit $X$ un préschéma, et supposons que son espace sous-jacent soit réunion
d'une famille finie de parties fermées $X_k$ $(1\leq k\leq n)$~; on considère
pour chaque $k$ le sous-préschéma réduit de $X$ ayant $X_k$ pour espace
sous-jacent (\hyperref[I.5.2.1-fr]{5.2.1}) et on le note encore $X_k$. Soit
$f:X\to Y$ un morphisme, et pour chaque $k$, soit $Y_k$ une partie fermée de
$Y$ telle que $f(X_k)\subset Y_k$~; on note encore $Y_k$ le sous-préschéma
réduit de $Y$ ayant $Y_k$ pour espace sous-jacent, de sorte que la restriction
$X_k\to Y$ de $f$ à $X_k$ se factorise en
$X_k\xrightarrow{f_k}Y_k\to Y$ (\hyperref[I.5.2.2-fr]{5.2.2}). Pour que $f$
soit séparé, il faut et il suffit que les $f_k$ le soient.
\end{proposition}

La nécessité résulte de
(\hyperref[I.5.5.1-fr]{5.5.1, (i), (ii) et (v)}). Inversement, si la condition
de l'énoncé est satisfaite, chacune des restrictions $X_k\to Y$ de $f$ est
séparée (\hyperref[I.5.5.1-fr]{5.5.1, (i) et (ii)})~; si $p_1$, $p_2$ sont les
projections de $X\times_Y X$, le sous-espace $\Delta_{X_k}(X_k)$ s'identifie
au sous-espace $\Delta_X(X)\cap p_1^{-1}(X_k)$ de l'espace sous-jacent à
$X\times_Y X$ (\hyperref[I.5.3.16-fr]{5.3.16})~; ces sous-espaces étant
fermés dans $X\times_Y X$, il en est de même de leur réunion $\Delta_X(X)$.

Supposons en particulier que les $X_k$ soient les
\emph{composantes irréductibles} de $X$~; alors on peut supposer que les $Y_k$
sont des composantes irréductibles de $Y$
(\hyperref[0.2.1.5-fr]{0, 2.1.5})~; la prop.
(\hyperref[I.5.5.4-fr]{5.5.4}) ramène donc dans ce cas la notion de séparation
au cas des préschémas \emph{intègres} (\hyperref[I.2.1.7-fr]{2.1.7}).

\oldpage[I]{138}

\begin{proposition}[5.5.5]
\label{I.5.5.5-fr}
Soit $(Y_\lambda)$ un recouvrement ouvert d'un préschéma $Y$~; pour qu'un
morphisme $f:X\to Y$ soit séparé, il faut et il suffit que chacune de ses
restrictions $f^{-1}(Y_\lambda)\to Y_\lambda$ soit séparée.
\end{proposition}

Si on pose $X_\lambda=f^{-1}(Y_\lambda)$, tout revient, compte tenu de
(\hyperref[I.4.2.4-fr]{4.2.4, b)}) et de l'identité des produits
$X_\lambda\times_Y X_\lambda$ et
$X_\lambda\times_{Y_\lambda}X_\lambda$
(\hyperref[I.3.2.5-fr]{3.2.5}), à prouver que les
$X_\lambda\times_Y X_\lambda$ forment un recouvrement de $X\times_Y X$. Or,
si l'on pose $Y_{\lambda\mu}=Y_\lambda\cap Y_\mu$ et
$X_{\lambda\mu}=X_\lambda\cap X_\mu=f^{-1}(Y_{\lambda\mu})$,
$X_\lambda\times_Y X_\mu$ s'identifie au produit
$X_{\lambda\mu}\times_{Y_{\lambda\mu}}X_{\lambda\mu}$
(\hyperref[I.3.2.6.4-fr]{3.2.6.4}), donc aussi à
$X_{\lambda\mu}\times_Y X_{\lambda\mu}$
(\hyperref[I.3.2.5-fr]{3.2.5}), et finalement à un ouvert de
$X_\lambda\times_Y X_\lambda$, ce qui établit notre assertion
(\hyperref[I.3.2.7-fr]{3.2.7}).

La prop. (\hyperref[I.5.5.4-fr]{5.5.4}) permet, en prenant un recouvrement de
$Y$ par des ouverts affines, de ramener l'étude des morphismes séparés à celle
des morphismes séparés à valeurs dans des schémas affines.

\begin{proposition}[5.5.6]
\label{I.5.5.6-fr}
Soient $Y$ un schéma affine, $X$ un préschéma, $(U_\alpha)$ un recouvrement de
$X$ par des ouverts affines. Pour qu'un morphisme $f:X\to Y$ soit séparé, il
faut et il suffit que, pour tout couple d'indices $(\alpha,\beta)$,
$U_\alpha\cap U_\beta$ soit un ouvert affine, et que l'anneau
$\Gamma(U_\alpha\cap U_\beta,\mathscr{O}_X)$ soit engendré par la réunion des
images canoniques des anneaux $\Gamma(U_\alpha,\mathscr{O}_X)$ et
$\Gamma(U_\beta,\mathscr{O}_X)$.
\end{proposition}

Les $U_\alpha\times_Y U_\beta$ forment un recouvrement ouvert de
$X\times_Y X$ (\hyperref[I.3.2.7-fr]{3.2.7})~; en désignant par $p$ et $q$ les
projections de $X\times_Y X$, on a
\[
  \begin{aligned}
    \Delta_X^{-1}(U_\alpha\times_Y U_\beta)
      &=\Delta_X^{-1}\bigl(p^{-1}(U_\alpha)\cap q^{-1}(U_\beta)\bigr)\\
      &=\Delta_X^{-1}\bigl(p^{-1}(U_\alpha)\bigr)
        \cap\Delta_X^{-1}\bigl(q^{-1}(U_\beta)\bigr)
        =U_\alpha\cap U_\beta~;
  \end{aligned}
\]
tout revient donc à exprimer que la restriction de $\Delta_X$ à
$U_\alpha\cap U_\beta$ est une immersion fermée dans
$U_\alpha\times_Y U_\beta$. Or, cette restriction n'est autre que
$(j_\alpha,j_\beta)_Y$, en désignant par $j_\alpha$ (resp. $j_\beta$) le
morphisme d'injection de $U_\alpha\cap U_\beta$ dans $U_\alpha$ (resp.
$U_\beta$), ainsi qu'il résulte des définitions. Comme
$U_\alpha\times_Y U_\beta$ est un schéma affine dont l'anneau est
canoniquement isomorphe à
$\Gamma(U_\alpha,\mathscr{O}_X)
 \otimes_{\Gamma(Y,\mathscr{O}_Y)}
 \Gamma(U_\beta,\mathscr{O}_X)$
(\hyperref[I.3.2.2-fr]{3.2.2}), on voit que $U_\alpha\cap U_\beta$ doit être
un schéma affine et que l'application
$h_\alpha\otimes h_\beta\to h_\alpha h_\beta$ de l'anneau
$A(U_\alpha\times_Y U_\beta)$ dans
$\Gamma(U_\alpha\cap U_\beta,\mathscr{O}_X)$ doit être surjective
(\hyperref[I.4.2.3-fr]{4.2.3}), ce qui achève la démonstration.

\begin{corollary}[5.5.7]
\label{I.5.5.7-fr}
Un schéma affine est séparé (et est par suite un schéma, ce qui justifie la
terminologie de (\hyperref[I.5.4.1-fr]{5.4.1})).
\end{corollary}

\begin{corollary}[5.5.8]
\label{I.5.5.8-fr}
Soit $Y$ un schéma affine~; pour que $f:X\to Y$ soit un morphisme séparé, il
faut et il suffit que $X$ soit séparé (autrement dit, que $X$ soit un schéma).
\end{corollary}

On observera en effet que le critère de
(\hyperref[I.5.5.6-fr]{5.5.6}) ne dépend pas de $f$.

\begin{corollary}[5.5.9]
\label{I.5.5.9-fr}
Pour qu'un morphisme $f:X\to Y$ soit séparé, il faut que pour tout ouvert $U$
sur lequel $Y$ induit un préschéma séparé, le préschéma induit $f^{-1}(U)$
soit séparé, et il suffit qu'il en soit ainsi pour tout ouvert affine
$U\subset Y$.
\end{corollary}

La nécessité de la condition résulte de
(\hyperref[I.5.5.4-fr]{5.5.4}) et
(\hyperref[I.5.5.1-fr]{5.5.1, (ii)})~; la suffisance résulte de
(\hyperref[I.5.5.4-fr]{5.5.4}) et
(\hyperref[I.5.5.8-fr]{5.5.8}), compte tenu de l'existence de recouvrements
ouverts affines de $Y$.

En particulier, si $X$ et $Y$ sont des schémas affines, \emph{tout} morphisme
$X\to Y$ est séparé.

\begin{proposition}[5.5.10]
\label{I.5.5.10-fr}
Soient $Y$ un schéma, $f:X\to Y$ un morphisme. Pour tout ouvert affine $U$ de
$X$ et tout ouvert affine $V$ de $Y$, $U\cap f^{-1}(V)$ est affine.
\end{proposition}

Soient $p_1$, $p_2$ les projections de $X\times_{\mathbf{Z}}Y$~; le
sous-espace $U\cap f^{-1}(V)$ est l'image par $p_1$ de
$\Gamma_f(X)\cap p_1^{-1}(U)\cap p_2^{-1}(V)$. Or,
$p_1^{-1}(U)\cap p_2^{-1}(V)$ s'identifie à l'espace sous-jacent au

\oldpage[I]{139}
préschéma $U\times_{\mathbf{Z}}V$
(\hyperref[I.3.2.7-fr]{3.2.7}), et est par suite un schéma affine
(\hyperref[I.3.2.2-fr]{3.2.2})~; comme $\Gamma_f(X)$ est fermé dans
$X\times_{\mathbf{Z}}Y$ (\hyperref[I.5.4.3-fr]{5.4.3}),
$\Gamma_f(X)\cap p_1^{-1}(U)\cap p_2^{-1}(V)$ est fermé dans
$U\times_{\mathbf{Z}}V$, et par suite le préschéma induit par le
sous-préschéma de $X\times_{\mathbf{Z}}Y$ associé à $\Gamma_f$
(\hyperref[I.4.2.1-fr]{4.2.1}), sur la partie ouverte
$\Gamma_f(X)\cap p_1^{-1}(U)\cap p_2^{-1}(V)$ de son espace sous-jacent, est
un sous-préschéma fermé d'un schéma affine, donc un schéma affine
(\hyperref[I.4.2.3-fr]{4.2.3}). La proposition résulte alors du fait que
$\Gamma_f$ est une immersion.

\begin{examples}[5.5.11]
\label{I.5.5.11-fr}
Le préschéma de l'exemple (\hyperref[I.2.3.2-fr]{2.3.2})
(<<~droite projective sur un corps $K$~>>) est \emph{séparé}, car pour le
recouvrement $(X_1,X_2)$ de $X$ par des ouverts affines,
$X_1\cap X_2=U_{12}$ est affine et $\Gamma(U_{12},\mathscr{O}_X)$, anneau des
fractions rationnelles de la forme $f(s)/s^m$ avec $f\in K[s]$, est engendré
par $K[s]$ et par $1/s$, donc les conditions de
(\hyperref[I.5.5.6-fr]{5.5.6}) sont vérifiées.

Avec le même choix de $X_1$, $X_2$, $U_{12}$ et $U_{21}$ que dans l'exemple
(\hyperref[I.2.3.2-fr]{2.3.2}), prenons cette fois pour $u_{12}$
l'isomorphisme qui à $f(s)$ fait correspondre $f(t)$~; on obtient cette fois
par recollement un préschéma \emph{intègre non séparé} $X$, car la première
condition de (\hyperref[I.5.5.6-fr]{5.5.6}) est vérifiée, mais non la seconde.
Il est immédiat ici que
$\Gamma(X,\mathscr{O}_X)\to\Gamma(X_1,\mathscr{O}_X)=K[s]$ est un
isomorphisme~; l'isomorphisme réciproque définit un morphisme
$f:X\to\operatorname{Spec}(K[s])$ qui est surjectif, et pour tout
$y\in\operatorname{Spec}(K[s])$ tel que $\mathfrak{j}_y\neq(0)$,
$f^{-1}(y)$ est réduit à un point, mais pour $\mathfrak{j}_y=(0)$,
$f^{-1}(y)$ se compose de \emph{deux} points distincts (on dit que $X$ est la
<<~droite affine sur $K$, où le point $0$ est dédoublé~>>).

On peut aussi donner des exemples où \emph{aucune} des deux conditions de
(\hyperref[I.5.5.6-fr]{5.5.6}) n'est vérifiée. Remarquons d'abord que dans le
spectre premier $Y$ de l'anneau de polynômes $A=K[s,t]$ à deux indéterminées
sur un corps $K$, l'ouvert $U$ réunion de $D(s)$ et de $D(t)$ \emph{n'est pas
un ouvert affine}. En effet, si $z$ est une section de $\mathscr{O}_Y$
au-dessus de $U$, il existe deux entiers $m\geq0$, $n\geq0$ tels que $s^m z$
et $t^n z$ soient les restrictions à $U$ de polynômes en $s$ et $t$
(\hyperref[I.1.4.1-fr]{1.4.1}), ce qui n'est évidemment possible que si la
section $z$ se prolonge en une section au-dessus de $Y$ tout entier,
identifiée à un polynôme en $s$ et $t$. Si $U$ était un ouvert affine, le
morphisme d'injection $U\to Y$ serait donc un isomorphisme
(\hyperref[I.1.7.3-fr]{1.7.3}), ce qui est absurde puisque $U\neq Y$.

Cela étant, prenons deux schémas affines $Y_1$, $Y_2$, spectres premiers des
anneaux $A_1=K[s_1,t_1]$, $A_2=K[s_2,t_2]$~; prenons
$U_{12}=D(s_1)\cup D(t_1)$, $U_{21}=D(s_2)\cup D(t_2)$, et pour $u_{12}$ la
restriction à $U_{21}$ de l'isomorphisme $Y_2\to Y_1$ correspondant à
l'isomorphisme d'anneaux qui à $f(s_1,t_1)$ fait correspondre $f(s_2,t_2)$~;
on a ainsi un exemple où aucune des conditions de
(\hyperref[I.5.5.6-fr]{5.5.6}) n'est satisfaite (le préschéma intègre ainsi
obtenu est dit <<~plan affine sur $K$, où le point $0$ est dédoublé~>>).
\end{examples}

\begin{remark}[5.5.12]
\label{I.5.5.12-fr}
\normalfont Étant donnée une propriété \textbf{P} de morphismes de préschémas,
considérons les propositions suivantes~:
\begin{enumerate}
  \item[\rm(i)] \emph{Toute immersion fermée possède la propriété \textbf{P}.}
  \item[\rm(ii)] \emph{Le composé de deux morphismes possédant la propriété
    \textbf{P} possède la propriété \textbf{P}.}
  \item[\rm(iii)] \emph{Si $f:X\to X'$, $g:Y\to Y'$ sont deux $S$-morphismes
    possédant la propriété \textbf{P}, $f\times_S g$ possède la propriété
    \textbf{P}.}

\oldpage[I]{140}
  \item[\rm(iv)] \emph{Si $f:X\to Y$ est un $S$-morphisme possédant la
    propriété \textbf{P}, tout $S'$-morphisme $f_{(S')}$ obtenu par une
    extension $S'\to S$ du préschéma de base, possède la propriété
    \textbf{P}.}
  \item[\rm(v)] \emph{Si le composé $g\circ f$ de deux morphismes
    $f:X\to Y$, $g:Y\to Z$ possède la propriété \textbf{P}, et si $g$ est
    séparé, $f$ possède la propriété \textbf{P}.}
  \item[\rm(vi)] \emph{Si un morphisme $f:X\to Y$ possède la propriété
    \textbf{P}, il en est de même de $f_{\mathrm{red}}$
    (\hyperref[I.5.1.5-fr]{5.1.5}).}
\end{enumerate}

Dans ces conditions, si on suppose (i) et (ii) vérifiées, (iii) et (iv) sont
\emph{équivalentes}, et (v) et (vi) sont des \emph{conséquences} de (i), (ii)
et (iii).

La première assertion a déjà été démontrée
(\hyperref[I.3.5.1-fr]{3.5.1}). Considérons la factorisation
(\hyperref[I.5.3.13-fr]{5.3.13}) de $f$ en
$X\xrightarrow{\Gamma_f}X\times_Z Y\xrightarrow{p_2}Y$~; la relation
$p_2=(g\circ f)\times_Z 1_Y$ montre que si $g\circ f$ possède la propriété
\textbf{P}, il en est de même de $p_2$ en vertu de (iii)~; si $g$ est séparé,
$\Gamma_f$ est une immersion fermée
(\hyperref[I.5.4.3-fr]{5.4.3}), donc possède aussi la propriété \textbf{P} par
(i)~; enfin, en vertu de (ii), $f$ possède la propriété \textbf{P}.

Enfin, considérons le diagramme commutatif
\[
  \xymatrix{
    X_{\mathrm{red}}\ar[r]^{f_{\mathrm{red}}}\ar[d] &
    Y_{\mathrm{red}}\ar[d]\\
    X\ar[r]_f &
    Y
  }
\]
où les flèches verticales sont des immersions fermées
(\hyperref[I.5.1.5-fr]{5.1.5}), donc possèdent la propriété \textbf{P} par
(i). L'hypothèse que $f$ possède la propriété \textbf{P} entraîne donc par
(ii) que
$X_{\mathrm{red}}\xrightarrow{f_{\mathrm{red}}}Y_{\mathrm{red}}\to Y$
possède la propriété \textbf{P}~; enfin, comme une immersion fermée est
séparée (\hyperref[I.5.5.1-fr]{5.5.1, (i)}), $f_{\mathrm{red}}$ possède la
propriété \textbf{P} en vertu de (v).

On notera que si on considère les propositions~:
\begin{enumerate}
  \item[\rm(i')] \emph{Toute immersion possède la propriété \textbf{P}.}
  \item[\rm(v')] \emph{Si $g\circ f$ possède la propriété \textbf{P}, il en
    est de même de $f$~;}
\end{enumerate}
alors le raisonnement fait ci-dessus montre que (v') est conséquence de (i'),
(ii) et (iii).
\end{remark}

\begin{env}[5.5.13]
\label{I.5.5.13-fr}
On notera que (v) et (vi) sont encore conséquences de (i), (iii) et
\begin{enumerate}
  \item[\rm(ii')] \emph{Si $j:X\to Y$ est une immersion fermée et
    $g:Y\to Z$ un morphisme possédant la propriété \textbf{P}, alors
    $g\circ j$ possède la propriété \textbf{P}.}
\end{enumerate}
De même, (v') est conséquence de (i'), (iii) et
\begin{enumerate}
  \item[\rm(ii'')] \emph{Si $j:X\to Y$ est une immersion et $g:Y\to Z$ un
    morphisme possédant la propriété \textbf{P}, alors $g\circ j$ possède la
    propriété \textbf{P}.}
\end{enumerate}
Cela résulte en effet aussitôt des raisonnements de
(\hyperref[I.5.5.12-fr]{5.5.12}).
\end{env}

% FIN SOURCE ega1/ega1-5-fr.tex

%
% DEBUT SOURCE ega1/ega1-6-fr.tex
\section{Conditions de finitude}
\label{section:I.6-fr}

\subsection{Préschémas noethériens et localement noethériens.}
\label{subsection:I.6.1-fr}

\begin{definition}[6.1.1]
\label{I.6.1.1-fr}
On dit qu'un préschéma $X$ est \emph{noethérien} (resp. \emph{localement
noethérien}) s'il est réunion finie (resp. réunion) d'ouverts affines
$V_\alpha$ tels que l'anneau de chacun des schémas induits sur les $V_\alpha$
soit noethérien.
\end{definition}

Il résulte aussitôt de (\hyperref[I.1.5.2-fr]{1.5.2}) que si $X$ est
localement noethérien, le faisceau structural $\mathscr{O}_X$ est un
\emph{faisceau cohérent d'anneaux}, la question étant locale. \emph{Tout
sous-$\mathscr{O}_X$-}

\oldpage[I]{141}
\emph{Module (resp. tout $\mathscr{O}_X$-Module quotient) quasi-cohérent}
d'un $\mathscr{O}_X$-Module \emph{cohérent} $\mathscr{F}$ est alors
\emph{cohérent}, car la question est de nouveau locale, et il suffit
d'appliquer (\hyperref[I.1.5.1-fr]{1.5.1}),
(\hyperref[I.1.4.1-fr]{1.4.1}) et
(\hyperref[I.1.3.10-fr]{1.3.10}), joints au fait qu'un sous-module (resp.
module quotient) d'un module de type fini sur un anneau noethérien est de
type fini. Plus particulièrement, \emph{tout faisceau d'idéaux
quasi-cohérent} de $\mathscr{O}_X$ est \emph{cohérent}.

Si un préschéma $X$ est réunion finie (resp. réunion) d'ouverts $W_\lambda$
tels que les préschémas induits sur les $W_\lambda$ soient noethériens (resp.
localement noethériens), il est clair que $X$ est noethérien (resp.
localement noethérien).

\begin{proposition}[6.1.2]
\label{I.6.1.2-fr}
Pour qu'un préschéma $X$ soit noethérien, il faut et il suffit qu'il soit
localement noethérien et que son espace sous-jacent soit quasi-compact.
L'espace sous-jacent de $X$ est alors noethérien.
\end{proposition}

\begin{proof}
La première assertion découle aussitôt des définitions et de
(\hyperref[I.1.1.10-fr]{1.1.10 (ii)}). La seconde résulte de
(\hyperref[I.1.1.6-fr]{1.1.6}) et du fait que tout espace réunion finie de
sous-espaces noethériens est noethérien
(\hyperref[0.2.2.3-fr]{0, 2.2.3}).
\end{proof}

\begin{proposition}[6.1.3]
\label{I.6.1.3-fr}
Soit $X$ un schéma affine d'anneau $A$. Les conditions suivantes sont
équivalentes~: a) $X$ est noethérien~; b) $X$ est localement noethérien~;
c) $A$ est noethérien.
\end{proposition}

\begin{proof}
L'équivalence de a) et b) résulte de (\hyperref[I.6.1.2-fr]{6.1.2}) et de ce
que tout schéma affine a un espace sous-jacent quasi-compact
(\hyperref[I.1.1.10-fr]{1.1.10})~; il est clair en outre que c) entraîne a).
Pour voir que a) entraîne c), remarquons qu'il y a un recouvrement fini
$(V_i)$ de $X$ par des ouverts affines tels que l'anneau $A_i$ du préschéma
induit sur $V_i$ soit noethérien. Soit alors $(\mathfrak{a}_n)$ une suite
croissante d'idéaux de $A$~; il lui correspond canoniquement de façon
biunivoque (\hyperref[I.1.3.7-fr]{1.3.7}) une suite croissante
$(\widetilde{\mathfrak{a}}_n)$ de faisceaux d'idéaux dans
$\widetilde{A}\equiv\mathscr{O}_X$~; pour voir que la suite
$(\mathfrak{a}_n)$ est stationnaire, il suffit de prouver que la suite
$(\widetilde{\mathfrak{a}}_n)$ l'est. Or, la restriction
$\widetilde{\mathfrak{a}}_n|V_i$ est un faisceau quasi-cohérent d'idéaux dans
$\mathscr{O}_X|V_i$, étant l'image réciproque de
$\widetilde{\mathfrak{a}}_n$ par l'injection canonique $V_i\to X$
(\hyperref[0.5.1.4-fr]{0, 5.1.4})~; $\widetilde{\mathfrak{a}}_n|V_i$ est
donc de la forme $\widetilde{\mathfrak{a}}_{ni}$, où $\mathfrak{a}_{ni}$ est
un idéal de $A_i$ (\hyperref[I.1.3.7-fr]{1.3.7}). Comme $A_i$ est noethérien,
la suite $(\mathfrak{a}_{ni})$ est stationnaire pour tout $i$, d'où la
proposition.
\end{proof}

On notera que le raisonnement précédent prouve aussi que \emph{si $X$ est un
préschéma noethérien, toute suite croissante de faisceaux cohérents d'idéaux
de $\mathscr{O}_X$ est stationnaire}.

\begin{proposition}[6.1.4]
\label{I.6.1.4-fr}
Tout sous-préschéma d'un préschéma noethérien (resp. localement noethérien)
est noethérien (resp. localement noethérien).
\end{proposition}

\begin{proof}
Il suffit de faire la démonstration pour un préschéma noethérien $X$~; en
outre, on est aussitôt ramené par la définition
(\hyperref[I.6.1.1-fr]{6.1.1}) au cas où $X$ est un schéma affine. Comme tout
sous-préschéma de $X$ est un sous-préschéma fermé d'un préschéma induit sur
un ouvert (\hyperref[I.4.1.3-fr]{4.1.3}), on peut se borner à considérer le
cas d'un sous-préschéma $Y$ fermé ou induit sur un ouvert de $X$. Le cas où
$Y$ est fermé est immédiat, car si $A$ est l'anneau de $X$, on sait que $Y$
est un schéma affine d'anneau $A/\mathfrak{J}$, où $\mathfrak{J}$ est un
idéal de $A$ (\hyperref[I.4.2.3-fr]{4.2.3})~; comme $A$ est noethérien
(\hyperref[I.6.1.3-fr]{6.1.3}), il en est de même de $A/\mathfrak{J}$.

Supposons maintenant $Y$ ouvert dans $X$~; l'espace sous-jacent $Y$ est
noethérien (\hyperref[I.6.1.2-fr]{6.1.2}), donc quasi-compact, et par suite
réunion finie d'ouverts $D(f_i)$ ($f_i\in A$)~;
\oldpage[I]{142}
tout revient à démontrer la proposition lorsque $Y=D(f)$ avec $f\in A$. Mais
alors $Y$ est un schéma affine dont l'anneau est isomorphe à $A_f$
(\hyperref[I.1.3.6-fr]{1.3.6})~; comme $A$ est noethérien
(\hyperref[I.6.1.3-fr]{6.1.3}), il en est de même de $A_f$.
\end{proof}

\begin{env}[6.1.5]
\label{I.6.1.5-fr}
On notera que le \emph{produit} de deux $S$-préschémas noethériens n'est pas
nécessairement noethérien, même si ces préschémas sont affines, car le produit
tensoriel de deux algèbres noethériennes n'est pas nécessairement un anneau
noethérien (cf. (\hyperref[I.6.3.8-fr]{6.3.8})).
\end{env}

\begin{proposition}[6.1.6]
\label{I.6.1.6-fr}
Si $X$ est un préschéma noethérien, le Nilradical $\mathscr{N}_X$ de
$\mathscr{O}_X$ est nilpotent.
\end{proposition}

On peut en effet recouvrir $X$ par un nombre fini d'ouverts affines $U_i$ et il
suffit de prouver qu'il existe des entiers $n_i$ tels que
$(\mathscr{N}_X|U_i)^{n_i}=0$~; si $n$ est le plus grand des $n_i$, on aura
alors $\mathscr{N}_X^n=0$. On est donc ramené au cas où
$X=\operatorname{Spec}(A)$ est affine, $A$ étant un anneau noethérien~; en
vertu de (\hyperref[I.5.1.1-fr]{5.1.1}) et de
(\hyperref[I.1.3.13-fr]{1.3.13}), il suffit d'observer que le nilradical de
$A$ est nilpotent (\cite{I-11}, p.~127, cor.~4).

\begin{corollary}[6.1.7]
\label{I.6.1.7-fr}
Soit $X$ un préschéma noethérien~; pour que $X$ soit un schéma affine, il faut
et il suffit que $X_{\mathrm{red}}$ le soit.
\end{corollary}

Cela résulte de (\hyperref[I.6.1.6-fr]{6.1.6}) et
(\hyperref[I.5.1.10-fr]{5.1.10}).

\begin{lemma}[6.1.8]
\label{I.6.1.8-fr}
Soient $X$ un espace topologique, $x$ un point de $X$, $U$ un voisinage ouvert
de $x$ n'ayant qu'un nombre fini de composantes irréductibles. Alors il existe
un voisinage $V$ de $x$ tel que tout voisinage ouvert de $x$ contenu dans $V$
soit connexe.
\end{lemma}

Soient $U_i$ ($1\leq i\leq m$) les composantes irréductibles de $U$ ne
contenant pas $x$~; le complémentaire dans $U$ de la réunion des $U_i$ est un
voisinage ouvert $V$ de $x$ dans $U$, donc aussi dans $X$~; c'est d'ailleurs
le complémentaire dans $X$ de la réunion des composantes irréductibles de $X$
qui ne contiennent pas $x$ (\hyperref[0.2.1.6-fr]{0, 2.1.6}). Soit alors $W$
un voisinage ouvert de $x$ contenu dans $V$. Les composantes irréductibles de
$W$ sont les traces sur $W$ des composantes irréductibles de $U$ qui
rencontrent $W$ (\hyperref[0.2.1.6-fr]{0, 2.1.6}), donc ces composantes
contiennent $x$~; comme elles sont connexes, il en est de même de $W$.

\begin{corollary}[6.1.9]
\label{I.6.1.9-fr}
Un espace topologique localement noethérien est localement connexe (ce qui
entraîne entre autres que ses composantes connexes sont ouvertes).
\end{corollary}

\begin{proposition}[6.1.10]
\label{I.6.1.10-fr}
Soit $X$ un espace topologique localement noethérien. Les conditions suivantes
sont équivalentes~:
\begin{enumerate}
  \item[a)] Les composantes irréductibles de $X$ sont ouvertes.
  \item[b)] Les composantes irréductibles de $X$ sont identiques à ses
    composantes connexes.
  \item[c)] Les composantes connexes de $X$ sont irréductibles.
  \item[d)] Deux composantes irréductibles distinctes de $X$ ne se rencontrent
    pas.
\end{enumerate}

Enfin, si $X$ est un préschéma, ces conditions équivalent aussi à~:
\begin{enumerate}
  \item[e)] Pour tout $x\in X$, $\operatorname{Spec}(\mathscr{O}_x)$ est
    irréductible (autrement dit le nilradical de $\mathscr{O}_x$ est premier).
\end{enumerate}
\end{proposition}

Il est immédiat que a) entraîne b), car un espace irréductible est connexe, et
a) entraîne que les composantes irréductibles de $X$ sont des ensembles à la
fois ouverts et fermés. Il est trivial que b) entraîne c)~; inversement, un
ensemble fermé $F$ contenant
\oldpage[I]{143}
une composante connexe $C$ de $X$ et distinct de $C$ ne peut être
irréductible, car cet ensemble n'étant pas connexe est réunion de deux
ensembles non vides disjoints à la fois ouverts et fermés dans $F$, donc
fermés dans $X$~; par suite c) entraîne b). On en conclut aussitôt que c)
entraîne d), deux composantes connexes distinctes étant sans point commun.

Nous n'avons pas utilisé jusqu'ici le fait que $X$ est localement noethérien.
Supposons maintenant cette hypothèse réalisée et montrons que d) entraîne
a)~: en vertu de (\hyperref[0.2.1.6-fr]{0, 2.1.6}), on peut se borner au cas
où l'espace $X$ est noethérien, donc n'a qu'un nombre fini de composantes
irréductibles. Comme celles-ci sont fermées et deux à deux disjointes, elles
sont ouvertes.

Enfin l'équivalence de d) et e) est valable sans supposer que l'espace
sous-jacent au préschéma $X$ soit localement noethérien. On peut en effet se
ramener au cas où $X=\operatorname{Spec}(A)$ est affine en vertu de
(\hyperref[0.2.1.6-fr]{0, 2.1.6})~; dire que $x$ n'est contenu que dans une
seule composante irréductible de $X$ signifie alors que $\mathfrak{j}_x$ ne
contient qu'un seul idéal minimal de $A$ (\hyperref[I.1.1.14-fr]{1.1.14}), ce
qui équivaut à dire que $\mathfrak{j}_x\mathscr{O}_x$ ne contient qu'un seul
idéal minimal de $\mathscr{O}_x$, d'où la conclusion.

\begin{corollary}[6.1.11]
\label{I.6.1.11-fr}
Soit $X$ un espace localement noethérien. Pour que $X$ soit irréductible, il
faut et il suffit que $X$ soit connexe et non vide, et que deux composantes
irréductibles distinctes de $X$ ne se rencontrent pas. Si $X$ est un
préschéma, cette dernière condition équivaut à ce que
$\operatorname{Spec}(\mathscr{O}_x)$ est irréductible pour tout $x\in X$.
\end{corollary}

La dernière partie a été vue dans (\hyperref[I.6.1.10-fr]{6.1.10})~; il n'y a
donc à prouver que la suffisance des conditions de la première assertion.
Mais d'après (\hyperref[I.6.1.10-fr]{6.1.10}), ces conditions entraînent que
les composantes irréductibles de $X$ sont ses composantes connexes, et comme
$X$ est connexe et non vide, il est irréductible.

\begin{corollary}[6.1.12]
\label{I.6.1.12-fr}
Soit $X$ un préschéma localement noethérien. Pour que $X$ soit intègre, il
faut et il suffit que $X$ soit connexe et que $\mathscr{O}_x$ soit intègre
pour tout $x\in X$.
\end{corollary}

\begin{proposition}[6.1.13]
\label{I.6.1.13-fr}
Soit $X$ un préschéma localement noethérien, et soit $x\in X$ un point tel que
le nilradical $\mathscr{N}_x$ de $\mathscr{O}_x$ soit premier (resp. que
$\mathscr{O}_x$ soit réduit, resp. intègre)~; alors il existe un voisinage
ouvert $U$ de $x$ qui est irréductible (resp. réduit, resp. intègre).
\end{proposition}

Il suffit de considérer les deux cas où $\mathscr{N}_x$ est premier et où
$\mathscr{N}_x=0$, la troisième hypothèse étant conjonction des deux
premières. Si $\mathscr{N}_x$ est premier, $x$ n'appartient qu'à une seule
composante irréductible $Y$ de $X$ (\hyperref[I.6.1.10-fr]{6.1.10})~; la
réunion des composantes irréductibles de $X$ ne contenant pas $x$ est fermée
(l'ensemble de ces composantes étant localement fini), et le complémentaire
$U$ de cette réunion est donc ouvert et contenu dans $Y$, donc irréductible
(\hyperref[0.2.1.6-fr]{0, 2.1.6}). Si $\mathscr{N}_x=0$, on a aussi
$\mathscr{N}_y=0$ pour tout $y$ dans un voisinage de $x$, car $\mathscr{N}$
est quasi-cohérent (\hyperref[I.5.1.1-fr]{5.1.1}), et par suite cohérent
puisque $X$ est localement noethérien, et la conclusion résulte de
(\hyperref[0.5.2.2-fr]{0, 5.2.2}).

\subsection{Préschémas artiniens.}
\label{subsection:I.6.2-fr}

\begin{definition}[6.2.1]
\label{I.6.2.1-fr}
On dit qu'un préschéma est \emph{artinien} s'il est affine et si son anneau
est artinien.
\end{definition}

\oldpage[I]{144}
\begin{proposition}[6.2.2]
\label{I.6.2.2-fr}
Étant donné un préschéma $X$, les conditions suivantes sont équivalentes~:
\begin{enumerate}
  \item[a)] $X$ est un schéma artinien~;
  \item[b)] $X$ est noethérien et son espace sous-jacent est discret~;
  \item[c)] $X$ est noethérien et les points de son espace sous-jacent sont
    fermés (condition $\mathrm{T}_1$).
\end{enumerate}

Lorsqu'il en est ainsi, l'espace sous-jacent à $X$ est fini, et l'anneau $A$
de $X$ est composé direct des anneaux locaux (artiniens) des points de $X$.
\end{proposition}

On sait que a) entraîne la dernière assertion (\cite{I-13}, p.~205, th.~3),
tout idéal premier de $A$ est alors maximal et est l'image réciproque de
l'idéal maximal de l'un des composants locaux de $A$, donc l'espace $X$ est
fini et discret~; a) entraîne donc b) et b) entraîne évidemment c). Pour voir
que c) entraîne a), montrons d'abord que $X$ est alors fini~; on peut en effet
se ramener au cas où $X$ est affine, et on sait qu'un anneau noethérien dont
tous les idéaux premiers sont maximaux est artinien (\cite{I-13}, p.~203),
d'où notre assertion. L'espace sous-jacent $X$ est alors discret, somme
topologique d'un nombre fini de points $x_i$, et les anneaux locaux
$\mathscr{O}_{x_i}=A_i$ sont artiniens~; il est clair que $X$ est isomorphe au
schéma affine spectre premier de l'anneau $A$ composé direct des $A_i$
(\hyperref[I.1.7.3-fr]{1.7.3}).

\subsection{Morphismes de type fini.}
\label{subsection:I.6.3-fr}

\begin{definition}[6.3.1]
\label{I.6.3.1-fr}
On dit qu'un morphisme $f:X\to Y$ est \emph{de type fini} si $Y$ est réunion
d'une famille $(V_\alpha)$ d'ouverts affines ayant la propriété suivante~:
\begin{enumerate}
  \item[(P)] $f^{-1}(V_\alpha)$ est réunion finie d'ouverts affines
    $U_{\alpha i}$ tels que chacun des anneaux $A(U_{\alpha i})$ soit une
    algèbre de type fini sur $A(V_\alpha)$.
\end{enumerate}
On dit alors aussi que $X$ est un préschéma de type fini sur $Y$, ou un
$Y$-préschéma de type fini.
\end{definition}

\begin{proposition}[6.3.2]
\label{I.6.3.2-fr}
Si $f:X\to Y$ est un morphisme de type fini, tout ouvert affine $W$ de $Y$
possède la propriété (P) de la déf. (\hyperref[I.6.3.1-fr]{6.3.1}).
\end{proposition}

Nous démontrerons d'abord le

\begin{lemma}[6.3.2.1]
\label{I.6.3.2.1-fr}
Si $T\subset Y$ est un ouvert affine, possédant la propriété (P), alors, pour
tout $g\in A(T)$, $D(g)$ possède la propriété (P).
\end{lemma}

En effet, par hypothèse, $f^{-1}(T)$ est réunion finie d'ouverts affines
$Z_j$ tels que $A(Z_j)$ soit une algèbre de type fini sur $A(T)$~; soit
$\varphi_j:A(T)\to A(Z_j)$ l'homomorphisme d'anneaux correspondant à la
restriction de $f$ à $Z_j$ (\hyperref[I.2.2.4-fr]{2.2.4}), et posons
$g_j=\varphi_j(g)$~; on a alors
$f^{-1}(D(g))\cap Z_j=D(g_j)$
(\hyperref[I.1.2.2.2-fr]{1.2.2.2}). Or,
$A(D(g_j))=A(Z_j)_{g_j}=A(Z_j)[1/g_j]$ est de type fini sur $A(Z_j)$ et
\emph{a fortiori} sur $A(T)$ en vertu de l'hypothèse, donc aussi sur
$A(D(g))=A(T)[1/g]$, ce qui démontre le lemme.

Ce lemme étant établi, comme $W$ est quasi-compact
(\hyperref[I.1.1.10-fr]{1.1.10}), il existe un recouvrement fini de $W$ par
des ensembles de la forme $D(g_i)$, où chaque $g_i$ appartient à un anneau
$A(V_{\alpha(i)})$. Chaque $D(g_i)$, étant quasi-compact, est réunion finie
d'ensembles $D(h_{ik})$ où $h_{ik}\in A(W)$~; si
$\varphi_i:A(W)\to A(D(g_i))$ est l'application canonique, on a
$D(h_{ik})=D(\varphi_i(h_{ik}))$ en vertu de
(\hyperref[I.1.2.2.2-fr]{1.2.2.2}). En vertu de
(\hyperref[I.6.3.2.1-fr]{6.3.2.1}), chacun des
$f^{-1}(D(h_{ik}))$ admet un recouvrement fini par des ouverts affines
$U_{ijk}$ tels que $A(U_{ijk})$ soit une algèbre de type fini sur
$A(D(h_{ik}))=A(W)[1/h_{ik}]$, d'où la proposition.

\oldpage[I]{145}
On peut donc dire que la notion de préschéma de type fini sur $Y$ est
\emph{locale sur $Y$}.

\begin{proposition}[6.3.3]
\label{I.6.3.3-fr}
Soient $X$, $Y$ deux schémas affines~; pour que $X$ soit de type fini sur
$Y$, il faut et il suffit que $A(X)$ soit une algèbre de type fini sur
$A(Y)$.
\end{proposition}

La condition étant évidemment suffisante, prouvons qu'elle est nécessaire.
Posons $A=A(Y)$, $B=A(X)$~; en vertu de
(\hyperref[I.6.3.2-fr]{6.3.2}), il existe un recouvrement ouvert affine fini
$(V_i)$ de $X$ tel que chacun des anneaux $A(V_i)$ soit une $A$-algèbre de
type fini. En outre, les $V_i$ étant quasi-compacts, on peut recouvrir chacun
d'eux par un nombre fini d'ouverts de la forme $D(g_{ij})\subset V_i$, où
$g_{ij}\in B$~; si $\varphi_i$ est l'homomorphisme $B\to A(V_i)$ correspondant
à l'injection canonique $V_i\to X$, on a
\[
  B_{g_{ij}}=(A(V_i))_{\varphi_i(g_{ij})}
  =A(V_i)[1/\varphi_i(g_{ij})],
\]
donc $B_{g_{ij}}$ est une $A$-algèbre de type fini. On peut donc se ramener
au cas où $V_i=D(g_i)$ avec $g_i\in B$. Par hypothèse, il existe une partie
finie $F_i$ de $B$ et un entier $n_i\geq 0$ tels que $B_{g_i}$ soit l'algèbre
engendrée sur $A$ par les éléments $b_i/g_i^{n_i}$, où $b_i$ parcourt $F_i$.
Comme les $g_i$ sont en nombre fini, on peut d'ailleurs supposer tous les
$n_i$ égaux à un même entier $n$. En outre, comme les $D(g_i)$ forment un
recouvrement de $X$, l'idéal engendré dans $B$ par les $g_i$ est égal à $B$,
autrement dit, il existe des $h_i\in B$ tels que
$\sum_i h_i g_i=1$. Soit alors $F$ la partie finie de $B$, réunion des $F_i$,
de l'ensemble des $g_i$ et de l'ensemble des $h_i$~; montrons que le
sous-anneau $B'=A[F]$ de $B$ est égal à $B$. Par hypothèse, pour tout $b\in B$
et tout $i$, l'image canonique de $b$ dans $B_{g_i}$ est de la forme
$b'_i/g_i^{m_i}$, où $b'_i\in B'$~; en multipliant les $b'_i$ par des
puissances convenables des $g_i$, on peut encore supposer tous les $m_i$
égaux à un même entier $m$. Par définition des anneaux de fractions, il y a
donc un entier $N$ (dépendant de $b$) tel que $N\geq m$ et $g_i^N b\in B'$ pour
tout $i$~; or, \emph{dans l'anneau $B'$}, les $g_i^N$ engendrent l'idéal $B'$,
puisqu'il en est ainsi des $g_i$ (les $h_i$ appartenant à $B'$)~; il y a donc
des $c_i\in B'$ tels que $\sum_i c_i g_i^N=1$, d'où
$b=\sum_i c_i g_i^N b\in B'$, C.Q.F.D.

\begin{proposition}[6.3.4]
\label{I.6.3.4-fr}
\begin{enumerate}
  \item[(i)] Toute immersion fermée est de type fini.
  \item[(ii)] Le composé de deux morphismes de type fini est de type fini.
  \item[(iii)] Si $f:X\to X'$, $g:Y\to Y'$ sont deux $S$-morphismes de type
    fini, $f\times_S g$ est de type fini.
  \item[(iv)] Si $f:X\to Y$ est un $S$-morphisme de type fini, $f_{(S')}$ est
    de type fini pour toute extension $g:S'\to S$ du préschéma de base.
  \item[(v)] Si le composé $g\circ f$ de deux morphismes est de type fini, et
    si $g$ est séparé, $f$ est de type fini.
  \item[(vi)] Si un morphisme $f$ est de type fini, il en est de même de
    $f_{\mathrm{red}}$.
\end{enumerate}
\end{proposition}

En vertu de (\hyperref[I.5.5.12-fr]{5.5.12}), il suffit de démontrer (i),
(ii) et (iv).

Pour établir (i), on peut se borner au cas d'une injection canonique
$X\to Y$, $X$ étant un sous-préschéma fermé de $Y$~; en outre
(\hyperref[I.6.3.2-fr]{6.3.2}), on peut supposer $Y$ affine, auquel cas $X$
est aussi affine (\hyperref[I.4.2.3-fr]{4.2.3}) et son anneau est isomorphe à
un anneau quotient $A/\mathfrak{J}$, où $A$ est l'anneau de $Y$ et
$\mathfrak{J}$ un idéal de $A$~; comme $A/\mathfrak{J}$ est de type fini sur
$A$, la conclusion en résulte.

Démontrons maintenant (ii). Soient $f:X\to Y$, $g:Y\to Z$ deux morphismes de
type fini, et soit $U$ un ouvert affine de $Z$~; $g^{-1}(U)$ admet un
recouvrement fini par des ouverts affines $V_i$ tels que $A(V_i)$ soit une
algèbre de type fini sur $A(U)$ (\hyperref[I.6.3.2-fr]{6.3.2})~; de même,

\oldpage[I]{146}
chacun des $f^{-1}(V_i)$ admet un recouvrement fini par des ouverts affines
$W_{ij}$ tels que $A(W_{ij})$ soit une algèbre de type fini sur $A(V_i)$, et
par suite aussi une algèbre de type fini sur $A(U)$~; d'où la conclusion.

Enfin, pour démontrer (iv), on peut se borner au cas où $S=Y$~; en effet,
$f_{(S')}$ est aussi égal à $f_{(Y_{(S')})}$, $f$ étant considéré comme un
$Y$-morphisme, et l'extension de la base étant $Y_{(S')}\to Y$
(\hyperref[I.3.3.9-fr]{3.3.9}). Soient alors $p$, $q$ les projections
$X_{(S')}\to X$ et $X_{(S')}\to S'$. Soit $V$ un ouvert affine dans $S$~;
$f^{-1}(V)$ est réunion finie d'ouverts affines $W_i$ dont chacun est tel que
$A(W_i)$ soit une algèbre de type fini sur $A(V)$
(\hyperref[I.6.3.2-fr]{6.3.2}). Soit $V'$ un ouvert affine de $S'$ contenu
dans $g^{-1}(V)$~; comme $f\circ p=g\circ q$, $q^{-1}(V')$ est contenu dans la
réunion des $p^{-1}(W_i)$~; d'autre part, l'intersection
$p^{-1}(W_i)\cap q^{-1}(V')$ s'identifie au produit $W_i\times_V V'$
(\hyperref[I.3.2.7-fr]{3.2.7}), qui est un schéma affine d'anneau isomorphe à
$A(W_i)\otimes_{A(V)}A(V')$ (\hyperref[I.3.2.2-fr]{3.2.2})~; ce dernier étant
par hypothèse une algèbre de type fini sur $A(V')$, la proposition est
démontrée.

\begin{corollary}[6.3.5]
\label{I.6.3.5-fr}
Soit $f:X\to Y$ un morphisme d'immersion. Si l'espace sous-jacent à $Y$
(resp. $X$) est localement noethérien (resp. noethérien), $f$ est de type fini.
\end{corollary}

On peut toujours supposer $Y$ affine (\hyperref[I.6.3.2-fr]{6.3.2})~; si
l'espace sous-jacent à $Y$ est localement noethérien, on peut en outre le
supposer noethérien et alors l'espace sous-jacent à $X$, qui en est un
sous-espace, est noethérien. Autrement dit, on peut supposer $Y$ affine et
l'espace sous-jacent à $X$ noethérien~; il existe alors un recouvrement de $X$
par un nombre fini d'ouverts affines $D(g_i)\subset Y$, où $g_i\in A(Y)$, tels
que $X\cap D(g_i)$ soit fermé dans $D(g_i)$ (donc un schéma affine
(\hyperref[I.4.2.3-fr]{4.2.3})), puisque $X$ est localement fermé dans $Y$
(\hyperref[I.4.1.3-fr]{4.1.3}). Alors $A(X\cap D(g_i))$ est une algèbre de type
fini sur $A(D(g_i))$, d'après (\hyperref[I.6.3.4-fr]{6.3.4, (i)}) et
(\hyperref[I.6.3.3-fr]{6.3.3}), et
$A(D(g_i))=A(Y)_{g_i}=A(Y)[1/g_i]$ est de type fini sur $A(Y)$, ce qui achève
la démonstration.

\begin{corollary}[6.3.6]
\label{I.6.3.6-fr}
Soient $f:X\to Y$, $g:Y\to Z$ deux morphismes. Si $g\circ f$ est de type fini,
et si $X$ est noethérien, ou $X\times_Z Y$ localement noethérien, $f$ est de
type fini.
\end{corollary}

Cela résulte aussitôt de la démonstration de
(\hyperref[I.5.5.12-fr]{5.5.12}) et de (\hyperref[I.6.3.5-fr]{6.3.5}) appliquée
au morphisme d'immersion $\Gamma_f$.

\begin{proposition}[6.3.7]
\label{I.6.3.7-fr}
Soit $f:X\to Y$ un morphisme de type fini~; si $Y$ est noethérien (resp.
localement noethérien), $X$ est noethérien (resp. localement noethérien).
\end{proposition}

On peut se borner à faire la démonstration lorsque $Y$ est noethérien. Alors $Y$
est réunion finie d'ouverts affines $V_i$ tels que les $A(V_i)$ soient des
anneaux noethériens. En vertu de (\hyperref[I.6.3.2-fr]{6.3.2}), chacun des
$f^{-1}(V_i)$ est réunion d'un nombre fini d'ouverts affines $W_{ij}$ tels que
les $A(W_{ij})$ soient des algèbres de type fini sur $A(V_i)$, donc des anneaux
noethériens~; cela prouve que $X$ est noethérien.

\begin{corollary}[6.3.8]
\label{I.6.3.8-fr}
Soit $X$ un préschéma de type fini sur $S$. Pour toute extension de la base
$S'\to S$ telle que $S'$ soit noethérien (resp. localement noethérien),
$X_{(S')}$ est noethérien (resp. localement noethérien).
\end{corollary}

Cela résulte de (\hyperref[I.6.3.7-fr]{6.3.7}), $X_{(S')}$ étant de type fini
sur $S'$ en vertu de (\hyperref[I.6.3.4-fr]{6.3.4, (iv)}).

On peut encore dire que dans un produit $X\times_S Y$ de $S$-préschémas, si
\emph{l'un} des
\oldpage[I]{147}
facteurs $X$, $Y$ est \emph{de type fini} sur $S$ et
\emph{l'autre noethérien} (resp. \emph{localement noethérien}), alors
$X\times_S Y$ est \emph{noethérien} (resp. \emph{localement noethérien}).

\begin{corollary}[6.3.9]
\label{I.6.3.9-fr}
Soit $X$ un préschéma de type fini sur un préschéma localement noethérien $S$.
Alors tout $S$-morphisme $f:X\to Y$ est de type fini.
\end{corollary}

En effet, on peut supposer $S$ noethérien~; si $\varphi:X\to S$,
$\psi:Y\to S$ sont les morphismes structuraux, on a $\varphi=\psi\circ f$, et
$X$ est noethérien en vertu de (\hyperref[I.6.3.7-fr]{6.3.7})~; $f$ est donc
de type fini en vertu de (\hyperref[I.6.3.6-fr]{6.3.6}).

\begin{proposition}[6.3.10]
\label{I.6.3.10-fr}
Soit $f:X\to Y$ un morphisme de type fini. Pour que $f$ soit surjectif, il faut
et il suffit que, pour tout corps algébriquement clos $\Omega$, l'application
$X(\Omega)\to Y(\Omega)$ correspondant à $f$
(\hyperref[I.3.4.1-fr]{3.4.1}) soit surjective.
\end{proposition}

La condition est suffisante, comme on le voit en considérant pour tout
$y\in Y$, une extension algébriquement close $\Omega$ de $k(y)$, et le
diagramme commutatif
\[
  \xymatrix{
    X\ar[dd]_f\\
    & \operatorname{Spec}(\Omega)\ar[ul]\ar[dl]\\
    Y
  }
\]
(\emph{cf.} (\hyperref[I.3.5.3-fr]{3.5.3})). Inversement, supposons $f$
surjective, et soit $g:\{\xi\}=\operatorname{Spec}(\Omega)\to Y$ un morphisme,
$\Omega$ étant un corps algébriquement clos. Si on considère le diagramme
\[
  \xymatrix{
    X\ar[d]_f &
    X_{(\Omega)}\ar[l]\ar[d]^{f_{(\Omega)}}\\
    Y &
    \operatorname{Spec}(\Omega)\ar[l]
  }
\]
il suffit donc de montrer qu'il existe dans $X_{(\Omega)}$ un
\emph{point rationnel sur $\Omega$}
(\hyperref[I.3.3.14-fr]{3.3.14}, \hyperref[I.3.4.3-fr]{3.4.3} et
\hyperref[I.3.4.4-fr]{3.4.4}). Comme $f$ est surjectif, $X_{(\Omega)}$ n'est
pas vide (\hyperref[I.3.5.10-fr]{3.5.10}) et comme $f$ est de type fini, il en
est de même de $f_{(\Omega)}$ (\hyperref[I.6.3.4-fr]{6.3.4, (iv)})~; donc
$X_{(\Omega)}$ contient un ouvert affine non vide $Z$ tel que $A(Z)$ soit une
algèbre de type fini non nulle sur $\Omega$. En vertu du th. des zéros de
Hilbert \cite{I-21}, il existe un $\Omega$-homomorphisme $A(Z)\to\Omega$, donc
une section de $X_{(\Omega)}$ au-dessus de $\operatorname{Spec}(\Omega)$, ce
qui démontre la proposition.

\subsection{Préschémas algébriques.}
\label{subsection:I.6.4-fr}

\begin{definition}[6.4.1]
\label{I.6.4.1-fr}
Étant donné un corps $K$, on appelle \emph{$K$-préschéma algébrique} un
préschéma $X$ de type fini sur $K$~; $K$ est appelé le \emph{corps de base}
de $X$. Si en outre $X$ est un schéma (ou, ce qui revient au même
(\hyperref[I.5.5.8-fr]{5.5.8}), si $X$ est un $K$-schéma), on dit aussi que
$X$ est un \emph{$K$-schéma algébrique}.
\end{definition}

Tout $K$-préschéma algébrique est \emph{noethérien}
(\hyperref[I.6.3.7-fr]{6.3.7}).

\begin{proposition}[6.4.2]
\label{I.6.4.2-fr}
Soit $X$ un $K$-préschéma algébrique. Pour qu'un point $x\in X$ soit fermé,
il faut et il suffit que $k(x)$ soit une extension algébrique de $K$, de degré
fini.
\end{proposition}

On peut supposer $X$ affine, l'anneau $A$ de $X$ étant une $K$-algèbre de type
fini. En effet, les ouverts affines $U$ de $X$ tels que $A(U)$ soit une
$K$-algèbre de type fini forment un recouvrement de $X$
(\hyperref[I.6.3.1-fr]{6.3.1}). Les points fermés de $X$ sont alors les points
tels que $\mathfrak{j}_x$ soit
\oldpage[I]{148}
un idéal maximal de $A$, autrement dit tels que $A/\mathfrak{j}_x$ soit un
corps (nécessairement égal à $k(x)$). Comme $A/\mathfrak{j}_x$ est une
$K$-algèbre de type fini, on voit que si $x$ est fermé, $k(x)$ est un corps
qui est une algèbre de type fini sur $K$, donc nécessairement une $K$-algèbre
de \emph{rang fini} \cite{I-21}. Inversement, si $k(x)$ est de rang fini sur
$K$, il en est de même de $A/\mathfrak{j}_x\subset k(x)$ et comme tout anneau
intègre qui est une $K$-algèbre de rang fini est un corps, on a
$A/\mathfrak{j}_x=k(x)$, donc $x$ est fermé.

\begin{corollary}[6.4.3]
\label{I.6.4.3-fr}
Soient $K$ un corps algébriquement clos, $X$ un $K$-préschéma algébrique~; les
points fermés de $X$ sont alors les points rationnels sur $K$
(\hyperref[I.3.4.4-fr]{3.4.4}) et s'identifient canoniquement aux points de
$X$ à valeurs dans $K$.
\end{corollary}

\begin{proposition}[6.4.4]
\label{I.6.4.4-fr}
Soit $X$ un préschéma algébrique sur un corps $K$. Les propriétés suivantes
sont équivalentes~:
\begin{enumerate}
  \item[\rm(a)] $X$ est artinien.
  \item[\rm(b)] L'espace sous-jacent à $X$ est discret.
  \item[\rm(c)] L'espace sous-jacent à $X$ n'a qu'un nombre fini de points
    fermés.
  \item[\rm(c')] L'espace sous-jacent à $X$ est fini.
  \item[\rm(d)] Les points de $X$ sont fermés.
  \item[\rm(e)] $X$ est isomorphe à $\operatorname{Spec}(A)$, où $A$ est une
    $K$-algèbre de rang fini.
\end{enumerate}
\end{proposition}

Comme $X$ est noethérien, il résulte de (\hyperref[I.6.2.2-fr]{6.2.2}) que
les conditions a), b), d) sont équivalentes et entraînent c) et c')~; par
ailleurs, il est clair que e) entraîne a). Reste à voir que c) entraîne d) et
e)~; on peut se borner au cas où $X$ est affine. Alors $A(X)$ est une
$K$-algèbre de type fini (\hyperref[I.6.3.3-fr]{6.3.3}), donc un anneau de
Jacobson (\cite[p.~3-11 et 3-12]{I-1}), dans lequel il n'y a par hypothèse
qu'un nombre fini d'idéaux maximaux. Comme une intersection finie d'idéaux
premiers ne peut être un idéal premier que si elle est égale à l'un d'eux,
tout idéal premier de $A(X)$ est donc maximal, d'où d). En outre, on sait
alors (\hyperref[I.6.2.2-fr]{6.2.2}) que $A(X)$ est une $K$-algèbre
artinienne de type fini, donc nécessairement de \emph{rang fini} \cite{I-21}.

\begin{env}[6.4.5]
\label{I.6.4.5-fr}
Lorsque les conditions de (\hyperref[I.6.4.4-fr]{6.4.4}) sont satisfaites,
on dit que $X$ est un schéma \emph{fini sur $K$}
(\emph{cf.} (\hyperref[II.6.1.1-fr]{II, 6.1.1})), ou un
\emph{$K$-schéma fini}, de \emph{rang} $[A:K]$ que l'on note aussi
$\operatorname{rg}_K(X)$~; si $X$, $Y$ sont deux schémas finis sur $K$, on a
\[
  \operatorname{rg}_K(X\amalg Y)
  =\operatorname{rg}_K(X)+\operatorname{rg}_K(Y)
  \tag{6.4.5.1}\label{I.6.4.5.1-fr}
\]
\[
  \operatorname{rg}_K(X\times_K Y)
  =\operatorname{rg}_K(X)\operatorname{rg}_K(Y)
  \tag{6.4.5.2}\label{I.6.4.5.2-fr}
\]
comme il résulte de (\hyperref[I.3.2.2-fr]{3.2.2}).
\end{env}

\begin{corollary}[6.4.6]
\label{I.6.4.6-fr}
Soit $X$ un schéma fini sur un corps $K$. Pour toute extension $K'$ de $K$,
$X\otimes_K K'$ est un schéma fini sur $K'$, et son rang sur $K'$ est égal au
rang de $X$ sur $K$.
\end{corollary}

En effet, si $A=A(X)$, on a $[A\otimes_K K':K']=[A:K]$.

\begin{corollary}[6.4.7]
\label{I.6.4.7-fr}
Soit $X$ un schéma fini sur un corps $K$~; on pose
$n=\sum_{x\in X}[k(x):K]_s$
\emph{(on rappelle que si $K'$ est une extension de $K$, $[K':K]_s$ est le
\emph{rang séparable} de $K'$ sur $K$, rang de la plus grande extension
algébrique séparable de $K$ contenue dans $K'$)}~;
\oldpage[I]{149}
alors, pour toute extension algébriquement close $\Omega$ de $K$, l'espace
sous-jacent à $X\otimes_K\Omega$ a exactement $n$ points, qui s'identifient
aux points de $X$ à valeurs dans $\Omega$.
\end{corollary}

On peut évidemment se borner au cas où l'anneau $A=A(X)$ est local
(\hyperref[I.6.2.2-fr]{6.2.2})~; soient $\mathfrak m$ son idéal maximal,
$L=A/\mathfrak m$ son corps résiduel, extension algébrique de $K$. Les points
de $X$ à valeurs dans $\Omega$ correspondent alors biunivoquement aux
$\Omega$-sections de $X\otimes_K\Omega$
(\hyperref[I.3.4.1-fr]{3.4.1} et \hyperref[I.3.3.14-fr]{3.3.14}), et aussi
aux $K$-homomorphismes de $L$ dans $\Omega$
(\hyperref[I.1.7.3-fr]{1.7.3}), d'où la proposition (Bourbaki, \emph{Alg.},
chap.~V, \textsection7, n\textsuperscript{o}~5, prop.~8), compte tenu de
(\hyperref[I.6.4.3-fr]{6.4.3}).

\begin{env}[6.4.8]
\label{I.6.4.8-fr}
Le nombre $n$ défini dans (\hyperref[I.6.4.7-fr]{6.4.7}) s'appelle le
\emph{rang séparable} de $A$ (ou de $X$) sur $K$, ou aussi le
\emph{nombre géométrique de points de $X$}~; il est donc égal au nombre
d'éléments de $X(\Omega)_K$. Il résulte aussitôt de cette définition que,
pour toute extension $K'$ de $K$, $X\otimes_K K'$ a même nombre géométrique
de points que $X$. Si on désigne ce nombre par $n(X)$, il est clair que si
$X$, $Y$ sont deux schémas finis sur $K$, on a
\[
  n(X\amalg Y)=n(X)+n(Y).
  \tag{6.4.8.1}\label{I.6.4.8.1-fr}
\]

Sous les mêmes hypothèses, on a aussi
\[
  n(X\times_K Y)=n(X)n(Y)
  \tag{6.4.8.2}\label{I.6.4.8.2-fr}
\]
comme il résulte aussitôt de l'interprétation de $n(X)$ comme nombre
d'éléments de $X(\Omega)_K$ et de la formule
(\hyperref[I.3.4.3.1-fr]{3.4.3.1}).
\end{env}

\begin{proposition}[6.4.9]
\label{I.6.4.9-fr}
Soient $K$ un corps, $X$, $Y$ deux $K$-préschémas algébriques, $f:X\to Y$ un
$K$-morphisme, $\Omega$ une extension algébriquement close de $K$, de degré
de transcendance infini sur $K$. Pour que $f$ soit surjectif, il faut et il
suffit que l'application $X(\Omega)_K\to Y(\Omega)_K$ correspondant à $f$
(\hyperref[I.3.4.1-fr]{3.4.1}) soit surjective.
\end{proposition}

La nécessité résulte de (\hyperref[I.6.3.10-fr]{6.3.10}), en remarquant que
$f$ est nécessairement de type fini (\hyperref[I.6.3.9-fr]{6.3.9}). Pour voir
que la condition est suffisante, on raisonne comme dans
(\hyperref[I.6.3.10-fr]{6.3.10}), en remarquant que pour tout $y\in Y$,
$k(y)$ est une extension de $K$ de type fini, et par suite est $K$-isomorphe
à un sous-corps de $\Omega$.

\begin{remark}[6.4.10]
\label{I.6.4.10-fr}
Nous verrons au chap.~IV que la conclusion de
(\hyperref[I.6.4.9-fr]{6.4.9}) est encore valable sans hypothèse relative au
degré de transcendance de $\Omega$ sur $K$.
\end{remark}

\begin{proposition}[6.4.11]
\label{I.6.4.11-fr}
Si $f:X\to Y$ est un morphisme de type fini, pour tout $y\in Y$, la fibre
$f^{-1}(y)$ est un préschéma algébrique sur le corps résiduel $k(y)$ et pour
tout $x\in f^{-1}(y)$, $k(x)$ est une extension de type fini de $k(y)$.
\end{proposition}

Comme $f^{-1}(y)=X\otimes_Y k(y)$
(\hyperref[I.3.6.3-fr]{3.6.3}), la proposition résulte de
(\hyperref[I.6.3.4-fr]{6.3.4, (iv)}) et de
(\hyperref[I.6.3.3-fr]{6.3.3}).

\begin{proposition}[6.4.12]
\label{I.6.4.12-fr}
Soient $f:X\to Y$, $g:Y'\to Y$ deux morphismes~; posons
$X'=X\times_Y Y'$ et soit $f'=f_{(Y')}:X'\to Y'$. Soit $y'\in Y'$,
$y=g(y')$~; si la fibre $f^{-1}(y)$ est un schéma algébrique fini sur $k(y)$,
alors la fibre $f'^{-1}(y')$ est un schéma algébrique fini sur $k(y')$, ayant
même rang et même nombre géométrique de points que $f^{-1}(y)$.
\end{proposition}

Compte tenu de la transitivité des fibres
(\hyperref[I.3.6.5-fr]{3.6.5}), cela résulte aussitôt de
(\hyperref[I.6.4.6-fr]{6.4.6}) et
(\hyperref[I.6.4.8-fr]{6.4.8}).

\begin{env}[6.4.13]
\label{I.6.4.13-fr}
La prop. (\hyperref[I.6.4.11-fr]{6.4.11}) montre que les morphismes de type
fini correspondent intuitivement aux «~familles algébriques de variétés
algébriques~», les points de $Y$ jouant
\oldpage[I]{150}
le rôle de «~paramètres~», ce qui donne à ces morphismes une signification
«~géométrique~». Les morphismes qui ne sont pas de type fini interviendront
surtout par la suite dans les questions de «~changement du préschéma de
base~», par exemple par localisation ou complétion.
\end{env}

\subsection{Détermination locale d'un morphisme.}
\label{subsection:I.6.5-fr}

\begin{proposition}[6.5.1]
\label{I.6.5.1-fr}
Soient $X$, $Y$ deux $S$-préschémas, $Y$ étant de type fini sur $S$~; soient
$x\in X$, $y\in Y$ au-dessus d'un même point $s\in S$.
\begin{enumerate}
  \item[(i)] Si deux $S$-morphismes $f=(\psi,\theta)$,
    $f'=(\psi',\theta')$ de $X$ dans $Y$ sont tels que
    $\psi(x)=\psi'(x)=y$, et que les $\mathscr{O}_s$-homomorphismes (locaux)
    $\theta_x^\sharp$ et ${\theta'}_x^\sharp$ de $\mathscr{O}_y$ dans
    $\mathscr{O}_x$ soient identiques, $f$ et $f'$ coïncident dans un
    voisinage ouvert de $x$.
  \item[(ii)] Supposons en outre $S$ localement noethérien. Pour tout
    $\mathscr{O}_s$-homomorphisme local
    $\varphi:\mathscr{O}_y\to\mathscr{O}_x$, il existe un voisinage ouvert
    $U$ de $x$ dans $X$ et un $S$-morphisme $f=(\psi,\theta)$ de $U$ dans $Y$
    tels que $\psi(x)=y$ et $\theta_x^\sharp=\varphi$.
\end{enumerate}
\end{proposition}

\begin{enumerate}
  \item[(i)] La question étant locale sur $S$, $X$ et $Y$, on peut supposer
    $S$, $X$, $Y$ affines d'anneaux respectifs $A$, $B$, $C$, $f$ et $f'$
    étant de la forme $({}^a\varphi,\widetilde{\varphi})$ et
    $({}^a\varphi',\widetilde{\varphi}')$ respectivement, où $\varphi$ et
    $\varphi'$ sont deux $A$-homomorphismes de $C$ dans $B$ tels que
    $\varphi^{-1}(\mathfrak{j}_x)={\varphi'}^{-1}(\mathfrak{j}_x)
    =\mathfrak{j}_y$, et les homomorphismes $\varphi_x$ et $\varphi'_x$ de
    $C_y$ dans $B_x$, déduits de $\varphi$ et $\varphi'$, sont identiques~;
    on peut en outre supposer que $C$ est une $A$-algèbre de type fini. Soient
    $c_i$ ($1\leq i\leq n$) des générateurs de la $A$-algèbre $C$, et posons
    $b_i=\varphi(c_i)$, $b'_i=\varphi'(c_i)$~; par hypothèse, on a
    $b_i/1=b'_i/1$ dans l'anneau de fractions $B_x$ ($1\leq i\leq n$). Cela
    signifie qu'il existe des éléments $s_i\in B-\mathfrak{j}_x$ tels que
    $s_i(b_i-b'_i)=0$ pour $1\leq i\leq n$, et on peut évidemment supposer
    tous les $s_i$ égaux à un même élément $g\in B-\mathfrak{j}_x$. On en
    conclut que l'on a $b_i/1=b'_i/1$ pour $1\leq i\leq n$ dans l'anneau de
    fractions $B_g$~; si $i_g$ est l'homomorphisme canonique $B\to B_g$, on a
    par suite $i_g\circ\varphi=i_g\circ\varphi'$~; donc les restrictions de
    $f$ et $f'$ à $D(g)$ sont identiques.
  \item[(ii)] On peut se ramener à la même situation que dans (i), et supposer
    en outre que l'anneau $A$ est noethérien. Soient $c_i$
    ($1\leq i\leq n$) des générateurs de la $A$-algèbre $C$, et soit
    $\alpha:A[X_1,\ldots,X_n]\to C$ l'homomorphisme de l'algèbre de polynômes
    $A[X_1,\ldots,X_n]$ sur $C$ transformant $X_i$ en $c_i$ pour
    $1\leq i\leq n$. Soit d'autre part $i_y$ l'homomorphisme canonique
    $C\to C_y$, et considérons l'homomorphisme composé
    \[
      \beta:A[X_1,\ldots,X_n]
      \xrightarrow{\alpha} C
      \xrightarrow{i_y} C_y
      \xrightarrow{\varphi} B_x.
    \]
    Désignons par $\mathfrak{a}$ le noyau de $\beta$~; comme $A$ est
    noethérien, il en est de même de $A[X_1,\ldots,X_n]$, et par suite
    $\mathfrak{a}$ admet un système fini de générateurs
    $Q_j(X_1,\ldots,X_n)$ ($1\leq j\leq m$). D'autre part, chacun des
    éléments $\varphi(i_y(c_i))$ peut s'écrire $b_i/s_i$, où $b_i\in B$ et
    $s_i\notin\mathfrak{j}_x$~; on peut en outre supposer tous les $s_i$ égaux
    à un même élément $g\in B-\mathfrak{j}_x$. Cela étant, on a par hypothèse
    $Q_j(b_1/g,\ldots,b_n/g)=0$ dans $B_x$~; posons
    \[
      Q_j(X_1/T,\ldots,X_n/T)
      =P_j(X_1,\ldots,X_n,T)/T^{k_j}
    \]
    où $P_j$ est homogène de degré $k_j$. Soit alors
    $d_j=P_j(b_1,\ldots,b_n,g)\in B$. Par hypothèse, on a $t_jd_j=0$ pour un
    $t_j\in B-\mathfrak{j}_x$ ($1\leq j\leq m$), et on peut évidemment
    supposer tous les $t_j$
\oldpage[I]{151}
    égaux à un même élément $h\in B-\mathfrak{j}_x$~; on en conclut que
    $P_j(hb_1,\ldots,hb_n,hg)=0$ pour $1\leq j\leq m$. Cela étant, considérons
    l'homomorphisme $\rho$ de $A[X_1,\ldots,X_n]$ dans l'anneau de fractions
    $B_{hg}$ qui applique $X_i$ sur $hb_i/hg$ ($1\leq i\leq n$)~; l'image de
    $\mathfrak{a}$ par cet homomorphisme est 0, et il en est \emph{a fortiori}
    de même de l'image par $\rho$ du noyau $\alpha^{-1}(0)$. Donc $\rho$ se
    factorise en
    $A[X_1,\ldots,X_n]\xrightarrow{\alpha}C\xrightarrow{\gamma}B_{hg}$, avec
    $\gamma(c_i)=hb_i/hg$, et il est clair que si $i_x$ est l'homomorphisme
    canonique $B_{hg}\to B_x$, le diagramme
    \[
      \xymatrix{
        C\ar[r]^\gamma\ar[d]_{i_y} & B_{hg}\ar[d]^{i_x}\\
        C_y\ar[r]_\varphi & B_x
      }
      \tag{6.5.1.1}\label{I.6.5.1.1-fr}
    \]
    est commutatif~; on a donc $\varphi=\gamma_x$, et comme $\varphi$ est un
    homomorphisme local, ${}^a\gamma(x)=y$~; $f=({}^a\gamma,\widetilde{\gamma})$
    est donc un $S$-morphisme du voisinage $D(hg)$ de $x$ dans $Y$ qui répond à
    la question.
\end{enumerate}

\begin{corollary}[6.5.2]
\label{I.6.5.2-fr}
Sous les hypothèses de (\hyperref[I.6.5.1-fr]{6.5.1, (ii)}), si en outre $X$
est de type fini sur $S$, on peut supposer le morphisme $f$ de type fini.
\end{corollary}

Cela résulte de (\hyperref[I.6.3.6-fr]{6.3.6}).

\begin{corollary}[6.5.3]
\label{I.6.5.3-fr}
Supposons vérifiées les hypothèses de
(\hyperref[I.6.5.1-fr]{6.5.1, (ii)}), et supposons en outre que $Y$ soit
intègre, et $\varphi$ un homomorphisme injectif. Alors on peut supposer que
$f=({}^a\gamma,\widetilde{\gamma})$, où $\gamma$ est injectif.
\end{corollary}

En effet, on peut supposer $C$ intègre
(\hyperref[I.5.1.4-fr]{5.1.4}), donc $i_y$ injectif~; il résulte alors du
diagramme (\hyperref[I.6.5.1.1-fr]{6.5.1.1}) que $\gamma$ est injectif.

\begin{proposition}[6.5.4]
\label{I.6.5.4-fr}
Soient $f=(\psi,\theta):X\to Y$ un morphisme de type fini, $x$ un point de
$X$, $y=\psi(x)$.
\begin{enumerate}
  \item[(i)] Pour que $f$ soit une immersion locale au point $x$
    (\hyperref[I.4.5.1-fr]{4.5.1}), il faut et il suffit que
    $\theta_x^\sharp:\mathscr{O}_y\to\mathscr{O}_x$ soit surjectif.
  \item[(ii)] On suppose en outre $Y$ localement noethérien. Pour que $f$ soit
    un isomorphisme local au point $x$
    (\hyperref[I.4.5.2-fr]{4.5.2}), il faut et il suffit que
    $\theta_x^\sharp$ soit un isomorphisme.
\end{enumerate}
\end{proposition}

\begin{enumerate}
  \item[(ii)] En vertu de (\hyperref[I.6.5.1-fr]{6.5.1}), il existe alors un
    voisinage ouvert $V$ de $y$ et un morphisme $g:V\to X$ tels que $g\circ f$
    (resp. $f\circ g$) soit défini et coïncide avec l'identité dans un
    voisinage de $x$ (resp. $y$), d'où on tire aisément que $f$ est un
    isomorphisme local.
  \item[(i)] La question étant locale sur $X$ et $Y$, on peut supposer $X$ et
    $Y$ affines, d'anneaux respectifs $A$, $B$~; on a
    $f=({}^a\varphi,\widetilde{\varphi})$, où $\varphi$ est un homomorphisme
    d'anneaux $B\to A$ qui fait de $A$ une $B$-algèbre de type fini~; on a
    $\varphi^{-1}(\mathfrak{j}_x)=\mathfrak{j}_y$ et l'homomorphisme
    $\varphi_x:B_y\to A_x$ déduit de $\varphi$ est surjectif. Soit $(t_i)$
    ($1\leq i\leq n$) un système de générateurs de la $B$-algèbre $A$~;
    l'hypothèse sur $\varphi_x$ implique qu'il existe des $b_i\in B$ et un
    $c\in B-\mathfrak{j}_y$ tels que, dans l'anneau de fractions $A_x$, on ait
    $t_i/1=\varphi(b_i)/\varphi(c)$ pour $1\leq i\leq n$. Par suite
    (\hyperref[I.1.3.3-fr]{1.3.3}), il existe $a\in A-\mathfrak{j}_x$ tel que,
    si l'on pose $g=a\varphi(c)$, on ait aussi
    $t_i/1=a\varphi(b_i)/g$ \emph{dans l'anneau de fractions $A_g$}. Cela étant,
    il existe par hypothèse un polynôme $Q(X_1,\ldots,X_n)$ à coefficients dans
    l'anneau $\varphi(B)$,
\oldpage[I]{152}
    tel que $a=Q(t_1,\ldots,t_n)$~; posons
    \[
      Q(X_1/T,\ldots,X_n/T)
      =P(X_1,\ldots,X_n,T)/T^m,
    \]
    où $P$ est homogène de degré $m$. Dans l'anneau $A_g$, on a
    \[
      a/1
      =a^mP(\varphi(b_1),\ldots,\varphi(b_n),\varphi(c))/g^m
      =a^m\varphi(d)/g^m
    \]
    où $d\in B$. Comme, dans $A_g$,
    $g/1=(a/1)(\varphi(c)/1)$ est inversible par définition, il en est de même
    de $a/1$ et de $\varphi(c)/1$, et on peut donc écrire
    \[
      a/1=(\varphi(d)/1)(\varphi(c)/1)^{-m}.
    \]
    On en conclut que $\varphi(d)/1$ est aussi inversible dans $A_g$. Posons
    alors $h=cd$~; comme $\varphi(h)/1$ est inversible dans $A_g$,
    l'homomorphisme composé
    \[
      B\xrightarrow{\varphi}A\longrightarrow A_g
    \]
    se factorise en
    \[
      B\longrightarrow B_h\xrightarrow{\gamma}A_g
    \]
    (\hyperref[0.1.2.4-fr]{0, 1.2.4}). Montrons que $\gamma$ est
    \emph{surjectif}~; il suffit de vérifier que l'image de $B_h$ dans $A_g$
    contient les $t_i/1$ et $(g/1)^{-1}$. Or, on a
    \[
      (g/1)^{-1}
      =(\varphi(c)/1)^{m-1}(\varphi(d)/1)^{-1}
      =\gamma(c^m/h),
    \]
    et
    \[
      a/1=\gamma(d^{m+1}/h^m),
    \]
    donc
    \[
      (a\varphi(b_i))/1=\gamma(b_i d^{m+1}/h^m),
    \]
    et comme
    $t_i/1=(a\varphi(b_i)/1)(g/1)^{-1}$, notre assertion est démontrée. Le
    choix de $h$ implique que $\psi(D(g))\subset D(h)$, et la restriction de
    $f$ à $D(g)$ est égale à
    $({}^a\gamma,\widetilde{\gamma})$~; comme $\gamma$ est surjectif, cette
    restriction est une immersion fermée de $D(g)$ dans $D(h)$
    (\hyperref[I.4.2.3-fr]{4.2.3}).
\end{enumerate}

\begin{corollary}[6.5.5]
\label{I.6.5.5-fr}
Soit $f=(\psi,\theta):X\to Y$ un morphisme de type fini. On suppose $X$
irréductible, on désigne par $x$ son point générique, et on pose $y=\psi(x)$.
\begin{enumerate}
  \item[(i)] Pour que $f$ soit une immersion locale en un point de $X$, il
    faut et il suffit que
    $\theta_x^\sharp:\mathscr{O}_y\to\mathscr{O}_x$ soit surjectif.
  \item[(ii)] On suppose de plus $Y$ irréductible et localement noethérien.
    Pour que $f$ soit un isomorphisme local en un point de $X$, il faut et il
    suffit que $y$ soit le point générique de $Y$ (ou, ce qui revient au même
    (\hyperref[0.2.1.4-fr]{0, 2.1.4}) que $f$ soit un morphisme
    \emph{dominant}) et que $\theta_x^\sharp$ soit un isomorphisme (autrement
    dit, que $f$ soit \emph{birationnel}
    (\hyperref[I.2.2.9-fr]{2.2.9})).
\end{enumerate}
\end{corollary}

Il est clair que (i) résulte de
(\hyperref[I.6.5.4-fr]{6.5.4, (i)}) compte tenu de ce que tout ouvert non vide
de $X$ contient $x$~; de même (ii) résulte de
(\hyperref[I.6.5.4-fr]{6.5.4, (ii)}).

\subsection{Morphismes quasi-compacts et morphismes localement de type fini.}
\label{subsection:I.6.6-fr}

\begin{definition}[6.6.1]
\label{I.6.6.1-fr}
On dit qu'un morphisme $f:X\to Y$ est \emph{quasi-compact} si l'image
réciproque par $f$ de tout ouvert quasi-compact de $Y$ est quasi-compacte.
\end{definition}

Soit $\mathfrak{B}$ une base de la topologie de $Y$ formée d'ouverts
quasi-compacts (par exemple d'ouverts affines)~; pour que $f$ soit
quasi-compact, il faut et il suffit que l'image réciproque par $f$ de tout
ensemble de $\mathfrak{B}$ soit quasi-compacte (ou, ce qui revient au même,
réunion \emph{finie} d'ouverts affines), car tout ouvert quasi-compact de $Y$
est réunion finie d'ensembles de $\mathfrak{B}$. Par exemple, si $X$ est
\emph{quasi-compact} et $Y$ \emph{affine}, alors \emph{tout} morphisme
$f:X\to Y$ est quasi-compact~: en effet, $X$ est réunion finie d'ensembles
ouverts affines $U_i$, et pour tout ouvert affine $V$ de $Y$,
$U_i\cap f^{-1}(V)$ est affine
(\hyperref[I.5.5.10-fr]{5.5.10}), donc quasi-compact.

Si $f:X\to Y$ est un morphisme quasi-compact, il est clair que pour tout
ouvert $V$ de $Y$, la restriction de $f$ à $f^{-1}(V)$ est un morphisme
quasi-compact $f^{-1}(V)\to V$. Inversement, si $(U_\alpha)$ est un
recouvrement ouvert de $Y$ et $f:X\to Y$ un morphisme tel que les restrictions
$f^{-1}(U_\alpha)\to U_\alpha$ soient quasi-compactes, alors $f$ est
quasi-compact.

\begin{definition}[6.6.2]
\label{I.6.6.2-fr}
On dit qu'un morphisme $f:X\to Y$ est \emph{localement de type fini} si, pour
tout $x\in X$, il existe un voisinage ouvert $U$ de $x$ et un voisinage ouvert
$V\supset f(U)$ de $y$ tels que la
\oldpage[I]{153}
restriction de $f$ à $U$ soit un morphisme de type fini de $U$ dans $V$. On
dit alors aussi que $X$ est un préschéma localement de type fini sur $Y$, ou
un $Y$-préschéma localement de type fini.
\end{definition}

Il résulte aussitôt de (\hyperref[I.6.3.2-fr]{6.3.2}) que si $f$ est
localement de type fini, alors, pour tout ouvert $W$ de $Y$, la restriction
de $f$ à $f^{-1}(W)$ est un morphisme $f^{-1}(W)\to W$ qui est localement de
type fini.

Si $Y$ est localement noethérien et si $X$ est localement de type fini sur
$Y$, $X$ est localement noethérien en vertu de
(\hyperref[I.6.3.7-fr]{6.3.7}).

\begin{proposition}[6.6.3]
\label{I.6.6.3-fr}
Pour qu'un morphisme $f:X\to Y$ soit de type fini, il faut et il suffit qu'il
soit quasi-compact et localement de type fini.
\end{proposition}

La nécessité des conditions est immédiate, vu
(\hyperref[I.6.3.1-fr]{6.3.1}) et la remarque suivant
(\hyperref[I.6.6.1-fr]{6.6.1}). Inversement, supposons ces conditions
satisfaites et soit $U$ un ouvert affine de $Y$, d'anneau $A$~; pour tout
$x\in f^{-1}(U)$, il y a par hypothèse un voisinage
$V(x)\subset f^{-1}(U)$ de $x$ et un voisinage $W(x)\subset U$ de
$y=f(x)$, contenant $f(V(x))$ et tel que la restriction de $f$ à $V(x)$ soit
un morphisme $V(x)\to W(x)$ qui est de type fini. En remplaçant $W(x)$ par
un voisinage $W_1(x)\subset W(x)$ de $x$ de la forme $D(g)$ (avec $g\in A$),
et $V(x)$ par $V(x)\cap f^{-1}(W_1(x))$, on peut supposer que $W(x)$ est de
la forme $D(g)$, donc de type fini sur $U$ (puisque son anneau s'écrit
$A[1/g]$)~; par suite $V(x)$ est de type fini sur $U$. En outre
$f^{-1}(U)$ est quasi-compact par hypothèse, donc réunion d'un nombre fini
d'ouverts $V(x_i)$, ce qui achève la démonstration.

\begin{proposition}[6.6.4]
\label{I.6.6.4-fr}
\begin{enumerate}
  \item[(i)] Une immersion $X\to Y$ est quasi-compacte si elle est fermée,
    ou si l'espace sous-jacent à $Y$ est localement noethérien ou si l'espace
    sous-jacent à $X$ est noethérien.
  \item[(ii)] Le composé de deux morphismes quasi-compacts est quasi-compact.
  \item[(iii)] Si $f:X\to Y$ est un $S$-morphisme quasi-compact, il en est de
    même de $f_{(S')}:X_{(S')}\to Y_{(S')}$ pour toute extension $g:S'\to S$
    du préschéma de base.
  \item[(iv)] Si $f:X\to X'$ et $g:Y\to Y'$ sont deux $S$-morphismes
    quasi-compacts, $f\times_S g$ est quasi-compact.
  \item[(v)] Si le composé de deux morphismes $f:X\to Y$, $g:Y\to Z$ est
    quasi-compact et si $g$ est séparé, ou l'espace sous-jacent à $X$
    localement noethérien, $f$ est quasi-compact.
  \item[(vi)] Pour qu'un morphisme $f$ soit quasi-compact, il faut et il
    suffit que $f_{\mathrm{red}}$ le soit.
\end{enumerate}
\end{proposition}

On notera que (vi) est évident puisque la propriété d'être quasi-compact pour
un morphisme ne dépend que de l'application continue correspondante des
espaces sous-jacents. Démontrons de même la partie de (v) correspondant au cas
où l'espace sous-jacent $X$ est supposé localement noethérien. Posons
$h=g\circ f$, et soit $U$ un ouvert quasi-compact dans $Y$~; $g(U)$ est
quasi-compact (non nécessairement ouvert) dans $Z$, donc contenu dans une
réunion finie d'ouverts quasi-compacts $V_j$
(\hyperref[I.2.1.3-fr]{2.1.3}), et $f^{-1}(U)$ est par suite contenu dans la
réunion des $h^{-1}(V_j)$, qui sont des sous-espaces quasi-compacts de $X$,
donc des sous-espaces noethériens. On en conclut
(\hyperref[0.2.2.3-fr]{0, 2.2.3}) que $f^{-1}(U)$ est un espace noethérien,
et \emph{a fortiori} quasi-compact.

Pour prouver les autres assertions, il suffit de démontrer (i), (ii) et (iii)
(\hyperref[I.5.5.12-fr]{5.5.12}). Or, (ii) est évidente, et (i) résulte de
(\hyperref[I.6.3.5-fr]{6.3.5}) lorsque l'espace $Y$ est localement noethérien
ou l'espace $X$ noethérien et est évidente pour une immersion fermée. Pour
établir (iii),
\oldpage[I]{154}
on peut se borner au cas où $S=Y$ (\hyperref[I.3.3.11-fr]{3.3.11})~; posons
$f'=f_{(S')}$, et soit $U'$ un ouvert quasi-compact dans $S'$. Pour tout
$s'\in U'$, soit $T$ un voisinage ouvert affine de $g(s')$ dans $S$, et soit
$W$ un voisinage ouvert affine de $s'$ contenu dans
$U'\cap g^{-1}(T)$~; il suffira de montrer que $f'^{-1}(W)$ est
quasi-compact~; autrement dit, on peut se ramener à prouver que lorsque $S$
et $S'$ sont affines, l'espace sous-jacent à $X\times_S S'$ est
quasi-compact. Mais comme $X$ est alors par hypothèse réunion finie d'ouverts
affines $V_j$, $X\times_S S'$ est réunion des espaces sous-jacents aux
schémas affines $V_j\times_S S'$
(\hyperref[I.3.2.2-fr]{3.2.2} et \hyperref[I.3.2.7-fr]{3.2.7}), ce qui achève
de prouver la proposition.

On notera aussi que si $X=X'\amalg X''$ est somme de deux préschémas, un
morphisme $f:X\to Y$ est quasi-compact si et seulement si ses restrictions à
$X'$ et $X''$ le sont.

\begin{proposition}[6.6.5]
\label{I.6.6.5-fr}
Soit $f:X\to Y$ un morphisme quasi-compact. Pour que $f$ soit dominant, il
faut et il suffit que pour tout point générique $y$ d'une composante
irréductible de $Y$, $f^{-1}(y)$ contienne le point générique d'une composante
irréductible de $X$.
\end{proposition}

Il est immédiat que la condition est suffisante (sans supposer $f$
quasi-compact). Pour voir qu'elle est nécessaire, considérons un voisinage
ouvert affine $U$ de $y$~; $f^{-1}(U)$ est quasi-compact, donc réunion
\emph{finie} d'ouverts affines $V_i$, et l'hypothèse que $f$ est dominant
implique que $y$ appartient à l'adhérence dans $U$ d'un des $f(V_i)$. On peut
évidemment supposer $X$ et $Y$ réduits~; comme l'adhérence dans $X$ d'une
composante irréductible de $V_i$ est une composante irréductible de $X$
(\hyperref[0.2.1.6-fr]{0, 2.1.6}), on peut remplacer $X$ par $V_i$, $Y$ par
le sous-préschéma fermé réduit de $U$ ayant $\overline{f(V_i)}\cap U$ pour
espace sous-jacent (\hyperref[I.5.2.1-fr]{5.2.1}), et on est ainsi ramené à
démontrer la proposition lorsque $X=\operatorname{Spec}(A)$ et
$Y=\operatorname{Spec}(B)$ sont affines et réduits. Comme $f$ est dominant,
$B$ est alors un sous-anneau de $A$ (\hyperref[I.1.2.7-fr]{1.2.7}), et la
proposition résulte alors du fait que tout idéal premier minimal de $B$ est
l'intersection de $B$ et d'un idéal premier minimal de $A$
(\hyperref[0.1.5.8-fr]{0, 1.5.8}).

\begin{proposition}[6.6.6]
\label{I.6.6.6-fr}
\begin{enumerate}
  \item[(i)] Toute immersion locale est localement de type fini.
  \item[(ii)] Si deux morphismes $f:X\to Y$, $g:Y\to Z$ sont localement de
    type fini, il en est de même de $g\circ f$.
  \item[(iii)] Si $f:X\to Y$ est un $S$-morphisme localement de type fini,
    $f_{(S')}:X_{(S')}\to Y_{(S')}$ est localement de type fini pour toute
    extension $S'\to S$ du préschéma de base.
  \item[(iv)] Si $f:X\to X'$ et $g:Y\to Y'$ sont deux $S$-morphismes
    localement de type fini, $f\times_S g$ est localement de type fini.
  \item[(v)] Si le composé $g\circ f$ de deux morphismes est localement de
    type fini, $f$ est localement de type fini.
  \item[(vi)] Si un morphisme $f$ est localement de type fini, il en est de
    même de $f_{\mathrm{red}}$.
\end{enumerate}
\end{proposition}

En vertu de (\hyperref[I.5.5.12-fr]{5.5.12}), il suffit de démontrer (i), (ii)
et (iii). Si $j:X\to Y$ est une immersion locale, pour tout $x\in X$, il y a
un voisinage ouvert $V$ de $j(x)$ dans $Y$ et un voisinage ouvert $U$ de $x$
dans $X$ tels que la restriction de $j$ à $U$ soit une immersion fermée
$U\to V$ (\hyperref[I.4.5.1-fr]{4.5.1}), donc cette restriction est de type
fini. Pour établir (ii), considérons un point $x\in X$~; il y a par hypothèse
un voisinage ouvert $W$ de $g(f(x))$ et un voisinage ouvert $V$ de $f(x)$
tels que $g(V)\subset W$ et que $V$ soit de type fini sur $W$~; en outre
$f^{-1}(V)$ est localement de type fini sur $V$
(\hyperref[I.6.6.2-fr]{6.6.2}), donc il y a un voisinage ouvert $U$ de $x$ qui
\oldpage[I]{155}
est contenu dans $f^{-1}(V)$ et de type fini sur $V$~; par suite on a
$g(f(U))\subset W$ et $U$ est de type fini sur $W$
(\hyperref[I.6.3.4-fr]{6.3.4, (ii)}). Enfin, pour démontrer (iii), on peut se
borner au cas où $Y=S$ (\hyperref[I.3.3.11-fr]{3.3.11})~; pour tout
$x'\in X'=X_{(S')}$, soient $x$ l'image de $x'$ dans $X$, $s$ l'image de $x$
dans $S$, $T$ un voisinage ouvert de $s$, $T'$ son image réciproque dans $S'$,
$U$ un voisinage ouvert de $x$ dont l'image est contenue dans $T$ et qui est de
type fini sur $T$~; alors $U\times_S T'=U\times_T T'$ est un voisinage ouvert
de $x'$ (\hyperref[I.3.2.7-fr]{3.2.7}) qui est de type fini sur $T'$
(\hyperref[I.6.3.4-fr]{6.3.4, (iv)}).

\begin{corollary}[6.6.7]
\label{I.6.6.7-fr}
Soient $X$, $Y$ deux $S$-préschémas qui sont localement de type fini sur $S$.
Si $S$ est localement noethérien, $X\times_S Y$ est localement noethérien.
\end{corollary}

En effet, $X$ étant localement de type fini sur $S$, est localement
noethérien, et $X\times_S Y$ est localement de type fini sur $X$, donc est
aussi localement noethérien.

\begin{remark}[6.6.8]
\label{I.6.6.8-fr}
La proposition (\hyperref[I.6.3.10-fr]{6.3.10}) et sa démonstration s'étendent
aussitôt au cas où l'on suppose seulement que le morphisme $f$ est localement
de type fini. De même, les propositions
(\hyperref[I.6.4.2-fr]{6.4.2}) et (\hyperref[I.6.4.9-fr]{6.4.9}) restent
valables lorsqu'on suppose que les préschémas $X$, $Y$ qui figurent dans leur
énoncé sont seulement localement de type fini sur le corps $K$.
\end{remark}

% FIN SOURCE ega1/ega1-6-fr.tex

%
% DEBUT SOURCE ega1/ega1-7-fr.tex
\section{Applications rationnelles}
\label{section:I.7-fr}

\subsection{Applications rationnelles et fonctions rationnelles.}
\label{subsection:I.7.1-fr}

\begin{env}[7.1.1]
\label{I.7.1.1-fr}
Soient $X$, $Y$ deux préschémas, $U$, $V$, deux ouverts denses dans $X$, $f$
(resp. $g$) un morphisme de $U$ (resp. $V$) dans $Y$~; nous dirons que $f$ et
$g$ sont
\emph{équivalents} s'ils coïncident dans un ouvert dense dans $U\cap V$.
Comme une intersection finie d'ouverts denses dans $X$ est un ouvert dense
dans $X$, il est clair que cette relation est une
\emph{relation d'équivalence}.
\end{env}

\begin{definition}[7.1.2]
\label{I.7.1.2-fr}
Étant donnés deux préschémas $X$, $Y$, on appelle
\emph{application rationnelle de $X$ dans $Y$} une classe d'équivalence de
morphismes de parties ouvertes denses de $X$ dans $Y$, pour la relation
définie dans (\hyperref[I.7.1.1-fr]{7.1.1}). Si $X$ et $Y$ sont des
$S$-préschémas, on dit qu'une telle classe est une
\emph{$S$-application rationnelle} s'il existe un représentant de cette classe
qui est un $S$-morphisme. On appelle
\emph{$S$-section rationnelle de $X$} toute $S$-application rationnelle de $S$
dans $X$. On appelle
\emph{fonction rationnelle sur un préschéma $X$} toute $X$-section rationnelle
du $X$-préschéma $X\otimes_{\mathbf{Z}}\mathbf{Z}[T]$ (où $T$ est une
indéterminée).

Par abus de langage, lorsqu'il ne sera question que de $S$-préschémas, on dira
«~application rationnelle~» au lieu de «~$S$-application rationnelle~» si
aucune confusion ne peut en résulter.

Soient $f$ une application rationnelle de $X$ dans $Y$, $U$ un ouvert de
$X$~; si $f_1$, $f_2$ sont des morphismes appartenant à la classe $f$, définis
respectivement dans des ouverts denses $V$, $W$ de $X$, les restrictions
$f_1|(U\cap V)$, $f_2|(U\cap W)$ coïncident dans $U\cap V\cap W$ qui est dense
dans $U$~; la classe de morphismes $f$ définit donc une application rationnelle
de $U$ dans $Y$, appelée
\emph{restriction de $f$ à $U$} et notée $f|U$.

Si, à tout $S$-morphisme $f:X\to Y$, on fait correspondre la $S$-application
rationnelle
\oldpage[I]{156}
à laquelle appartient $f$, on définit une application canonique de
$\operatorname{Hom}_S(X,Y)$ dans l'ensemble des $S$-applications rationnelles
de $X$ dans $Y$. On désigne par $\Gamma_{\mathrm{rat}}(X/Y)$ l'ensemble des
$Y$-sections rationnelles de $X$, et on a donc une application canonique
$\Gamma(X/Y)\to\Gamma_{\mathrm{rat}}(X/Y)$. Il est clair en outre que si $X$ et
$Y$ sont deux $S$-préschémas, l'ensemble des $S$-applications rationnelles de
$X$ dans $Y$ s'identifie canoniquement à
$\Gamma_{\mathrm{rat}}((X\times_S Y)/X)$
(\hyperref[I.3.3.14-fr]{3.3.14}).
\end{definition}

\begin{env}[7.1.3]
\label{I.7.1.3-fr}
Il résulte aussitôt de (\hyperref[I.7.1.2-fr]{7.1.2}) et de
(\hyperref[I.3.3.14-fr]{3.3.14}) que les fonctions rationnelles sur $X$
s'identifient canoniquement aux
\emph{classes d'équivalence des sections du faisceau structural
$\mathscr{O}_X$} au-dessus d'ouverts partout denses de $X$, deux telles
sections étant équivalentes si elles coïncident dans un ouvert partout dense
contenu dans l'intersection de leurs ensembles de définition. Il en résulte en
particulier que les fonctions rationnelles sur $X$ forment un
\emph{anneau} $R(X)$.
\end{env}

\begin{env}[7.1.4]
\label{I.7.1.4-fr}
Lorsque $X$ est un préschéma \emph{irréductible}, tout ouvert non vide est dense
dans $X$~; on peut encore dire que les ouverts non vides de $X$ sont les
\emph{voisinages ouverts du point générique $x$} de $X$. Dire que deux
morphismes de parties ouvertes non vides de $X$ dans $Y$ sont équivalents
signifie donc dans ce cas qu'ils ont \emph{même germe} au point $x$. Autrement
dit, les applications rationnelles (resp. $S$-applications rationnelles)
$X\to Y$ s'identifient aux \emph{germes de morphismes} (resp. de
$S$-morphismes) de parties ouvertes non vides de $X$ dans $Y$ au point
générique $x$ de $X$. En particulier~:
\end{env}

\begin{proposition}[7.1.5]
\label{I.7.1.5-fr}
Si $X$ est un préschéma irréductible, l'anneau $R(X)$ des fonctions rationnelles
sur $X$ s'identifie canoniquement à l'anneau local $\mathscr{O}_x$ du point
générique $x$ de $X$. C'est un anneau local de dimension 0, et par suite un
anneau local artinien lorsque $X$ est noethérien~; c'est un corps lorsque $X$
est intègre, et il s'identifie au corps des fractions de $A(X)$ lorsqu'en outre
$X$ est un schéma affine.
\end{proposition}

Vu ce qui précède et l'identification des fonctions rationnelles aux sections
de $\mathscr{O}_X$ au-dessus d'un ouvert partout dense, la première assertion
n'est autre que la définition de la fibre d'un faisceau en un point. Pour les
autres assertions, on peut se borner au cas où $X$ est affine d'anneau $A$~;
alors $\mathfrak{j}_x$ est le nilradical de $A$, et $\mathscr{O}_x$ est donc de
dimension 0~; si $A$ est intègre, $\mathfrak{j}_x=(0)$, et $\mathscr{O}_x$ est
donc le corps des fractions de $A$. Enfin, si $A$ est noethérien, on sait
(\cite{I-11}, p.~127, cor.~4) que $\mathfrak{j}_x$ est nilpotent et
$\mathscr{O}_x=A_{\mathfrak{j}_x}$ artinien.

Si $X$ est \emph{intègre}, l'anneau $\mathscr{O}_z$ est intègre pour tout
$z\in X$~; tout ouvert affine $U$ contenant $z$ contient aussi $x$, et $R(U)$,
égal au corps des fractions de $A(U)$, s'identifie donc à $R(X)$~; on en
conclut que $R(X)$ s'identifie aussi au
\emph{corps des fractions de $\mathscr{O}_z$}~: l'identification canonique de
$\mathscr{O}_z$ à un sous-anneau de $R(X)$ consiste à faire correspondre à tout
germe de section $s\in\mathscr{O}_z$ l'unique fonction rationnelle sur $X$,
classe d'une section de $\mathscr{O}_X$ (nécessairement définie sur un ouvert
partout dense) ayant pour germe $s$ au point $z$.

\begin{env}[7.1.6]
\label{I.7.1.6-fr}
Supposons maintenant que $X$ ait un nombre \emph{fini} de composantes
irréductibles $X_i$ ($1\leq i\leq n$) (ce qui est le cas lorsque l'espace
sous-jacent à $X$ est \emph{noethérien})~; soit $X'_i$ l'ouvert de $X$,
complémentaire par rapport à $X_i$ de la réunion des $X_j\cap X_i$ pour
$j\neq i$~; $X'_i$ est irréductible, son point générique $x_i$ est le point
générique de $X_i$, et les $X'_i$ sont deux à deux disjoints, leur réunion
étant dense dans $X$ (\hyperref[0.2.1.6-fr]{0, 2.1.6}). Pour
\oldpage[I]{157}
tout ouvert partout dense $U$ de $X$, $U_i=U\cap X'_i$ est un ouvert non vide
dense dans $X'_i$, les $U_i$ étant deux à deux sans point commun, donc
$U'=\bigcup_i U_i$ est dense dans $X$. Se donner un morphisme de $U'$ dans $Y$
revient à se donner (arbitrairement) un morphisme de chacun des $U_i$ dans $Y$.
Donc~:
\end{env}

\begin{proposition}[7.1.7]
\label{I.7.1.7-fr}
Soient $X$, $Y$ deux préschémas (resp. $S$-préschémas) tels que $X$ ait un
nombre fini de composantes irréductibles $X_i$, de points génériques $x_i$
($1\leq i\leq n$). Si $R_i$ est l'ensemble des germes de morphismes (resp.
$S$-morphismes) de parties ouvertes de $X$ dans $Y$ au point $x_i$, l'ensemble
des applications rationnelles (resp. $S$-applications rationnelles) de $X$
dans $Y$ s'identifie au produit des $R_i$ ($1\leq i\leq n$).
\end{proposition}

\begin{corollary}[7.1.8]
\label{I.7.1.8-fr}
Soit $X$ un préschéma noethérien. L'anneau des fonctions rationnelles sur $X$
est un anneau artinien, dont les composants locaux sont les anneaux
$\mathscr{O}_{x_i}$ des points génériques $x_i$ des composantes irréductibles
de $X$.
\end{corollary}

\begin{corollary}[7.1.9]
\label{I.7.1.9-fr}
Soient $A$ un anneau noethérien, et soit $X=\operatorname{Spec}(A)$. Si $Q$ est
le complémentaire de la réunion des idéaux premiers minimaux de $A$, l'anneau
des fonctions rationnelles sur $X$ s'identifie canoniquement à l'anneau de
fractions $Q^{-1}A$.
\end{corollary}

Cela résultera du lemme suivant~:

\begin{lemma}[7.1.9.1]
\label{I.7.1.9.1-fr}
Pour qu'un élément $f\in A$ soit tel que $D(f)$ soit partout dense dans $X$, il
faut et il suffit que $f\in Q$~; tout ouvert dense dans $X$ contient un ouvert
de la forme $D(f)$, où $f\in Q$.
\end{lemma}

En effet, supposons ce lemme démontré~; l'anneau des sections
$\Gamma(D(f),\mathscr{O}_X)$ s'identifiant à $A_f$
(\hyperref[I.1.3.6-fr]{1.3.6} et \hyperref[I.1.3.7-fr]{1.3.7}), il résulte du
fait que les $D(f)$ avec $f\in Q$ forment un ensemble cofinal dans l'ensemble
ordonné (pour $\supset$) des ouverts denses dans $X$, et de la déf.
(\hyperref[I.7.1.1-fr]{7.1.1}) que l'anneau des fonctions rationnelles sur $X$
s'identifie à la limite inductive des $A_f$ pour $f\in Q$ (pour la relation de
préordre «~$g$ est multiple de $f$~»), c'est-à-dire à $Q^{-1}A$
(\hyperref[0.1.4.5-fr]{0, 1.4.5}).

Pour démontrer (\hyperref[I.7.1.9.1-fr]{7.1.9.1}), désignons encore par $X_i$
les composantes irréductibles de $X$ ($1\leq i\leq n$)~; si $D(f)$ est dense
dans $X$, $D(f)\cap X_i\neq\varnothing$ pour $1\leq i\leq n$ et
réciproquement~; mais cela signifie que $f\notin\mathfrak{p}_i$ pour
$1\leq i\leq n$, en posant $\mathfrak{p}_i=\mathfrak{j}(X_i)$, et comme les
$\mathfrak{p}_i$ sont les idéaux premiers minimaux de $A$
(\hyperref[I.1.1.14-fr]{1.1.14}), les relations
$f\notin\mathfrak{p}_i$ ($1\leq i\leq n$) équivalent à $f\in Q$, d'où la
première assertion du lemme. D'autre part, si $U$ est un ouvert dense dans $X$,
le complémentaire de $U$ est un ensemble de la forme $V(\mathfrak{a})$, où
$\mathfrak{a}$ est un idéal qui n'est contenu dans aucun des $\mathfrak{p}_i$~;
il n'est donc pas contenu dans leur réunion (\cite{I-10}, p.~13), et il existe
donc un $f\in\mathfrak{a}$ appartenant à $Q$~; d'où $D(f)\subset U$, ce qui
achève la démonstration.

\begin{env}[7.1.10]
\label{I.7.1.10-fr}
Supposons de nouveau $X$ irréductible, de point générique $x$. Comme tout
ouvert non vide $U$ de $X$ contient $x$, et par suite contient aussi tout
$z\in X$ tel que $x\in\overline{\{z\}}$, tout morphisme $U\to Y$ peut se
composer avec le morphisme canonique $\operatorname{Spec}(\mathscr{O}_x)\to X$
(\hyperref[I.2.4.1-fr]{2.4.1})~; et deux morphismes dans $Y$ de deux parties
ouvertes non vides de $X$, qui coïncident dans une partie ouverte non vide de
$X$ donnent par composition le même morphisme
$\operatorname{Spec}(\mathscr{O}_x)\to Y$. Autrement dit, à toute application
rationnelle de $X$ dans $Y$ correspond ainsi un morphisme
$\operatorname{Spec}(\mathscr{O}_x)\to Y$ bien déterminé.
\end{env}

\oldpage[I]{158}
\begin{proposition}[7.1.11]
\label{I.7.1.11-fr}
Soient $X$, $Y$ deux $S$-préschémas~; on suppose $X$ irréductible, de point
générique $x$, et $Y$ de type fini sur $S$. Deux $S$-applications rationnelles
de $X$ dans $Y$, auxquelles correspond le même $S$-morphisme
$\operatorname{Spec}(\mathscr{O}_x)\to Y$ sont identiques. Si on suppose de
plus $S$ localement noethérien, tout $S$-morphisme de
$\operatorname{Spec}(\mathscr{O}_x)$ dans $Y$ correspond à une $S$-application
rationnelle (et une seule) de $X$ dans $Y$.
\end{proposition}

Compte tenu de ce que tout ouvert non vide de $X$ est partout dense, cela
résulte aussitôt de (\hyperref[I.6.5.1-fr]{6.5.1}).

\begin{corollary}[7.1.12]
\label{I.7.1.12-fr}
On suppose $S$ localement noethérien, et les autres hypothèses de
(\hyperref[I.7.1.11-fr]{7.1.11}) satisfaites. Les $S$-applications rationnelles
de $X$ dans $Y$ s'identifient alors aux points du $S$-préschéma $Y$, à valeurs
dans le $S$-préschéma $\operatorname{Spec}(\mathscr{O}_x)$.
\end{corollary}

Cela n'est autre que (\hyperref[I.7.1.11-fr]{7.1.11}), avec la terminologie
introduite dans (\hyperref[I.3.4.1-fr]{3.4.1}).

\begin{corollary}[7.1.13]
\label{I.7.1.13-fr}
On suppose remplies les conditions de (\hyperref[I.7.1.12-fr]{7.1.12}). Soit
$s$ l'image de $x$ dans $S$. La donnée d'une $S$-application rationnelle de
$X$ dans $Y$ équivaut à la donnée d'un point $y$ de $Y$ au-dessus de $s$, et
d'un $\mathscr{O}_s$-homomorphisme local
$\mathscr{O}_y\to\mathscr{O}_x=R(X)$.
\end{corollary}

Cela résulte de (\hyperref[I.7.1.11-fr]{7.1.11}) et de
(\hyperref[I.2.4.4-fr]{2.4.4}).

En particulier~:

\begin{corollary}[7.1.14]
\label{I.7.1.14-fr}
Sous les conditions de (\hyperref[I.7.1.12-fr]{7.1.12}), les $S$-applications
rationnelles de $X$ dans $Y$ ne dépendent (pour $Y$ donné) que du
$S$-préschéma $\operatorname{Spec}(\mathscr{O}_x)$ et en particulier restent
les mêmes lorsqu'on remplace $X$ par
$\operatorname{Spec}(\mathscr{O}_z)$, pour tout $z\in X$.
\end{corollary}

En effet, comme $z\in\overline{\{x\}}$, $x$ est le point générique de
$Z=\operatorname{Spec}(\mathscr{O}_z)$, et
$\mathscr{O}_{X,x}=\mathscr{O}_{Z,x}$.

Lorsque $X$ est intègre, $R(X)=\mathscr{O}_x=k(x)$ est un corps
(\hyperref[I.7.1.5-fr]{7.1.5})~; les corollaires précédents se spécialisent
alors en~:

\begin{corollary}[7.1.15]
\label{I.7.1.15-fr}
On suppose vérifiées les conditions de (\hyperref[I.7.1.12-fr]{7.1.12}) et en
outre que $X$ est intègre. Soit $s$ l'image de $x$ dans $S$. Alors les
$S$-applications rationnelles de $X$ dans $Y$ s'identifient aux points
géométriques de $Y\otimes_S k(s)$ à valeurs dans l'extension $R(X)$ de $k(s)$,
autrement dit chacune d'elles équivaut à la donnée d'un point $y\in Y$
au-dessus de $s$ et d'un $k(s)$-monomorphisme de $k(y)$ dans
$k(x)=R(X)$.
\end{corollary}

Les points de $Y$ au-dessus de $s$ s'identifient en effet à ceux de
$Y\otimes_S k(s)$ (\hyperref[I.3.6.3-fr]{3.6.3}) et les
$\mathscr{O}_s$-homomorphismes locaux $\mathscr{O}_y\to R(X)$ aux
$k(s)$-monomorphismes $k(y)\to R(X)$.

Plus particulièrement~:

\begin{corollary}[7.1.16]
\label{I.7.1.16-fr}
Soient $k$ un corps, $X$, $Y$ deux préschémas algébriques
(\hyperref[I.6.4.1-fr]{6.4.1}) sur $k$~; on suppose en outre $X$ intègre. Alors
les $k$-applications rationnelles de $X$ dans $Y$ s'identifient aux points
géométriques de $Y$ à valeurs dans l'extension $R(X)$ de $k$
(\hyperref[I.3.4.4-fr]{3.4.4}).
\end{corollary}

\subsection{Domaine de définition d'une application rationnelle.}
\label{subsection:I.7.2-fr}

\begin{env}[7.2.1]
\label{I.7.2.1-fr}
Soient $X$, $Y$ deux préschémas, $f$ une application rationnelle de $X$ dans
$Y$. On dit que $f$ est \emph{définie en un point $x\in X$} s'il existe un
ensemble ouvert partout dense $U$ contenant $x$ et un morphisme $U\to Y$
appartenant à la classe d'équivalence $f$. L'ensemble des points $x\in X$ où
$f$ est définie est appelé le \emph{domaine de définition de $f$}~; il est
clair que c'est un ouvert partout dense dans $X$.
\end{env}

\oldpage[I]{159}
\begin{proposition}[7.2.2]
\label{I.7.2.2-fr}
Soient $X$, $Y$ deux $S$-préschémas, tels que $X$ soit réduit et $Y$ séparé
sur $S$. Soit $f$ une $S$-application rationnelle de $X$ dans $Y$, $U_0$ son
domaine de définition. Il existe alors un $S$-morphisme et un seul $U_0\to Y$
appartenant à la classe $f$.
\end{proposition}

Comme pour tout morphisme $U\to Y$ appartenant à la classe $f$, on a
nécessairement $U\subset U_0$, il est clair que la proposition sera
conséquence du

\begin{lemma}[7.2.2.1]
\label{I.7.2.2.1-fr}
Sous les hypothèses de (\hyperref[I.7.2.2-fr]{7.2.2}), soient $U_1$, $U_2$
deux ouverts partout denses de $X$, $f_i:U_i\to Y$ ($i=1,2$) deux
$S$-morphismes tels qu'il existe un ouvert $V\subset U_1\cap U_2$, dense dans
$X$ et dans lequel $f_1$ et $f_2$ coïncident. Alors $f_1$ et $f_2$ coïncident
dans $U_1\cap U_2$.
\end{lemma}

On peut évidemment se borner au cas où $X=U_1=U_2$. Comme $X$ (donc $V$)
est réduit, $X$ est le plus petit sous-préschéma fermé de $X$ majorant $V$
(\hyperref[I.5.2.2-fr]{5.2.2}). Soit
$g=(f_1,f_2)_S:X\to Y\times_S Y$~; comme par hypothèse la diagonale
$T=\Delta_Y(Y)$ est un sous-préschéma fermé de $Y\times_S Y$,
$Z=g^{-1}(T)$ est un sous-préschéma fermé de $X$
(\hyperref[I.4.4.1-fr]{4.4.1}). Si $h:V\to Y$ est la restriction commune de
$f_1$ et $f_2$ à $V$, la restriction de $g$ à $V$ est $g'=(h,h)_S$, qui se
factorise en $g'=\Delta_Y\circ h$~; comme $\Delta_Y^{-1}(T)=Y$, on a
$g'^{-1}(T)=V$, et par suite $Z$ est un sous-préschéma fermé de $X$ induisant
$V$, donc majorant $V$, ce qui entraîne $Z=X$. De la relation
$g^{-1}(T)=X$, on déduit (\hyperref[I.4.4.1-fr]{4.4.1}) que $g$ se factorise
en $\Delta_Y\circ f$, où $f$ est un morphisme $X\to Y$, ce qui entraîne par
définition du morphisme diagonal que $f_1=f_2=f$.

Il est clair que le morphisme $U_0\to Y$ défini dans
(\hyperref[I.7.2.2-fr]{7.2.2}) est l'unique morphisme de la classe $f$ qui
\emph{ne puisse être prolongé} à un morphisme d'une partie ouverte de $X$
contenant strictement $U_0$. \emph{Sous les hypothèses de
(\hyperref[I.7.2.2-fr]{7.2.2})}, on peut donc \emph{identifier} les
applications rationnelles de $X$ dans $Y$ aux morphismes \emph{non
prolongeables} (à des ouverts strictement plus grands) d'ouverts partout
denses de $X$ dans $Y$. Avec cette identification, la prop.
(\hyperref[I.7.2.2-fr]{7.2.2}) entraîne~:

\begin{corollary}[7.2.3]
\label{I.7.2.3-fr}
Les hypothèses sur $X$ et $Y$ étant celles de
(\hyperref[I.7.2.2-fr]{7.2.2}), soit $U$ un ouvert partout dense de $X$. Il
existe une correspondance biunivoque canonique entre les $S$-morphismes de
$U$ dans $Y$ et les $S$-applications rationnelles de $X$ dans $Y$ définies
en tous les points de $U$.
\end{corollary}

En vertu de (\hyperref[I.7.2.2-fr]{7.2.2}), pour tout $S$-morphisme $f$ de
$U$ dans $Y$, il existe en effet une $S$-application rationnelle et une seule
$\overline f$ de $X$ dans $Y$ qui prolonge $f$.

\begin{corollary}[7.2.4]
\label{I.7.2.4-fr}
Soient $S$ un schéma, $X$ un $S$-préschéma réduit, $Y$ un $S$-schéma,
$f:U\to Y$ un $S$-morphisme d'un ouvert dense $U$ de $X$ dans $Y$. Si
$\overline f$ est la $\mathbf{Z}$-application rationnelle de $X$ dans $Y$ qui
prolonge $f$, $\overline f$ est un $S$-morphisme (et est par suite la
$S$-application rationnelle de $X$ dans $Y$ prolongeant $f$).
\end{corollary}

En effet, si $\varphi:X\to S$, $\psi:Y\to S$ sont les morphismes structuraux,
$U_0$ le domaine de définition de $\overline f$, $j$ l'injection $U_0\to X$,
il suffit de prouver que $\psi\circ\overline f=\varphi\circ j$, ce qui résulte
aussitôt de (\hyperref[I.7.2.2.1-fr]{7.2.2.1}), puisque $f$ est un
$S$-morphisme.

\begin{corollary}[7.2.5]
\label{I.7.2.5-fr}
Soient $X$, $Y$ deux $S$-préschémas~; on suppose $X$ réduit, $X$ et $Y$
séparés sur $S$. Soient $p:Y\to X$ un $S$-morphisme (faisant de $Y$ un
$X$-préschéma), $U$ un ouvert partout dense de $X$, $f$ une $U$-section de
$Y$~; alors l'application rationnelle $\overline f$ de $X$ dans $Y$
prolongeant $f$ est une $X$-section rationnelle de $Y$.
\end{corollary}

\oldpage[I]{160}
Il faut prouver que $p\circ\overline f$ est l'identité dans le domaine de
définition de $\overline f$~; puisque $X$ est séparé sur $S$, cela résulte
encore de (\hyperref[I.7.2.2.1-fr]{7.2.2.1}).

\begin{corollary}[7.2.6]
\label{I.7.2.6-fr}
Soient $X$ un préschéma réduit, $U$ un ouvert partout dense de $X$. Il y a
correspondance biunivoque canonique entre les sections de $\mathscr{O}_X$ au-dessus
de $U$ et les fonctions rationnelles $f$ sur $X$ définies en tout point de $U$.
\end{corollary}

Compte tenu de (\hyperref[I.7.2.3-fr]{7.2.3}),
(\hyperref[I.7.1.2-fr]{7.1.2}) et (\hyperref[I.7.1.3-fr]{7.1.3}), il suffit de
remarquer que le $X$-préschéma $X\otimes_{\mathbf Z}\mathbf Z[T]$ est séparé
au-dessus de $X$ (\hyperref[I.5.5.1-fr]{5.5.1, (iv)}).

\begin{corollary}[7.2.7]
\label{I.7.2.7-fr}
Soient $Y$ un préschéma réduit, $f:X\to Y$ un morphisme séparé, $U$ un ouvert
partout dense de $Y$, $g:U\to f^{-1}(U)$ une $U$-section de $f^{-1}(U)$, $Z$
le sous-préschéma réduit de $X$ ayant $\overline{g(U)}$ pour espace sous-jacent
(\hyperref[I.5.2.1-fr]{5.2.1}). Pour que $g$ soit restriction d'une
$Y$-section de $X$ (autrement dit (\hyperref[I.7.2.5-fr]{7.2.5}) pour que
l'application rationnelle de $Y$ dans $X$ prolongeant $g$ soit partout
définie), il faut et il suffit que la restriction de $f$ à $Z$ soit un
isomorphisme de $Z$ sur $Y$.
\end{corollary}

La restriction de $f$ à $f^{-1}(U)$ est un morphisme séparé
(\hyperref[I.5.5.1-fr]{5.5.1, (i)}), donc $g$ est une immersion fermée
(\hyperref[I.5.4.6-fr]{5.4.6}), et par suite
$g(U)=Z\cap f^{-1}(U)$ et le sous-préschéma induit par $Z$ sur l'ouvert
$g(U)$ de $Z$ est identique au sous-préschéma fermé de $f^{-1}(U)$ associé à
$g$ (\hyperref[I.5.2.1-fr]{5.2.1}). Il est clair alors que la condition de
l'énoncé est suffisante, car si elle est remplie et si $f_Z:Z\to Y$ est la
restriction de $f$ à $Z$ et $\overline g:Y\to Z$ l'isomorphisme réciproque,
$\overline g$ prolonge $g$. Inversement, si $g$ est restriction à $U$ d'une
$Y$-section $h$ de $X$, $h$ est une immersion fermée
(\hyperref[I.5.4.6-fr]{5.4.6}), donc $h(Y)$ est fermé, et comme il est contenu
dans $Z$, il est égal à $Z$, et il résulte de
(\hyperref[I.5.2.1-fr]{5.2.1}) que $h$ est nécessairement un isomorphisme de
$Y$ sur le sous-préschéma fermé $Z$ de $X$.

\begin{env}[7.2.8]
\label{I.7.2.8-fr}
Soient $X$, $Y$ deux $S$-préschémas, $X$ étant supposé réduit et $Y$ séparé
sur $S$. Soit $f$ une $S$-application rationnelle de $X$ dans $Y$, et soit
$x$ un point de $X$~; on peut composer $f$ avec le $S$-morphisme canonique
$\operatorname{Spec}(\mathscr{O}_x)\to X$
(\hyperref[I.2.4.1-fr]{2.4.1}) pourvu que la trace sur
$\operatorname{Spec}(\mathscr{O}_x)$ du domaine de définition de $f$ soit
dense dans $\operatorname{Spec}(\mathscr{O}_x)$ (identifié à l'ensemble des
$z\in X$ tels que $x\in\overline{\{z\}}$
(\hyperref[I.2.4.2-fr]{2.4.2})). Ceci aura lieu dans les cas suivants~:
\begin{enumerate}
  \item[1\textsuperscript{o}] $X$ est irréductible (donc intègre), car alors
    le point générique $\xi$ de $X$ est le point générique de
    $\operatorname{Spec}(\mathscr{O}_x)$~; comme le domaine de définition $U$
    de $f$ contient $\xi$, $U\cap\operatorname{Spec}(\mathscr{O}_x)$ contient
    $\xi$, donc est dense dans $\operatorname{Spec}(\mathscr{O}_x)$.
  \item[2\textsuperscript{o}] $X$ est localement noethérien~; notre assertion
    résulte en effet alors du
\end{enumerate}
\end{env}

\begin{lemma}[7.2.8.1]
\label{I.7.2.8.1-fr}
Soient $X$ un préschéma dont l'espace sous-jacent est localement noethérien,
$x$ un point de $X$. Les composantes irréductibles de
$\operatorname{Spec}(\mathscr{O}_x)$ sont les traces sur
$\operatorname{Spec}(\mathscr{O}_x)$ des composantes irréductibles de $X$
contenant $x$. Pour qu'un ouvert $U\subset X$ soit tel que
$U\cap\operatorname{Spec}(\mathscr{O}_x)$ soit dense dans
$\operatorname{Spec}(\mathscr{O}_x)$, il faut et il suffit qu'il rencontre les
composantes irréductibles de $X$ contenant $x$ (ce qui a lieu en particulier
si $U$ est dense dans $X$).
\end{lemma}

La seconde assertion résulte évidemment de la première, et il suffit donc de
démontrer celle-ci. Comme $\operatorname{Spec}(\mathscr{O}_x)$ est contenu
dans tout ouvert affine $U$ contenant $x$, et que les composantes
irréductibles de $U$ contenant $x$ sont les traces sur $U$ des composantes
irréductibles de $X$ contenant $x$
(\hyperref[0.2.1.6-fr]{0, 2.1.6}), on peut supposer $X$ affine d'anneau $A$.
Comme les idéaux premiers de $A_x$ correspondent biunivoquement aux idéaux
premiers
\oldpage[I]{161}
de $A$ contenus dans $\mathfrak{j}_x$
(\hyperref[0.1.2.6-fr]{0, 1.2.6}), les idéaux premiers minimaux de $A_x$
correspondent aux idéaux premiers minimaux de $A$ contenus dans
$\mathfrak{j}_x$, d'où le lemme (\hyperref[I.1.1.14-fr]{1.1.14}).

Cela étant, supposons que l'on soit dans l'un des deux cas précités. Si $U$
est le domaine de définition de la $S$-application rationnelle $f$, désignons
par $f'$ l'application rationnelle de $\operatorname{Spec}(\mathscr{O}_x)$
dans $Y$ qui coïncide (compte tenu de (\hyperref[I.2.4.2-fr]{2.4.2})) avec
$f$ dans $U\cap\operatorname{Spec}(\mathscr{O}_x)$~; nous dirons que cette
application rationnelle est \emph{induite par $f$}.

\begin{proposition}[7.2.9]
\label{I.7.2.9-fr}
Soient $S$ un préschéma localement noethérien, $X$ un $S$-préschéma réduit,
$Y$ un $S$-schéma de type fini. On suppose en outre $X$ irréductible ou
localement noethérien. Soient alors $f$ une $S$-application rationnelle de
$X$ dans $Y$, $x$ un point de $X$. Pour que $f$ soit définie au point $x$,
il faut et il suffit que l'application rationnelle $f'$ de
$\operatorname{Spec}(\mathscr{O}_x)$ dans $Y$, induite par $f$
(\hyperref[I.7.2.8-fr]{7.2.8}), soit un morphisme.
\end{proposition}

La condition étant évidemment nécessaire (puisque
$\operatorname{Spec}(\mathscr{O}_x)$ est contenu dans tout ouvert contenant
$x$), prouvons qu'elle est suffisante. En vertu de
(\hyperref[I.6.5.1-fr]{6.5.1}), il existe un voisinage ouvert $U$ de $x$ dans
$X$ et un $S$-morphisme $g$ de $U$ dans $Y$, induisant $f'$ sur
$\operatorname{Spec}(\mathscr{O}_x)$. Si $X$ est irréductible, $U$ est dense
dans $X$, et en vertu de (\hyperref[I.7.2.3-fr]{7.2.3}) on peut supposer que
$g$ est une $S$-application rationnelle. En outre, le point générique de $X$
appartient à $\operatorname{Spec}(\mathscr{O}_x)$ et au domaine de définition
de $f$, donc $f$ et $g$ coïncident en ce point, et par suite dans un ensemble
ouvert non vide de $X$ (\hyperref[I.6.5.1-fr]{6.5.1}). Mais comme $f$ et $g$
sont des $S$-applications rationnelles, elles sont identiques
(\hyperref[I.7.2.3-fr]{7.2.3}), donc $f$ est définie en $x$.

Si maintenant on suppose $X$ localement noethérien, on peut supposer $U$
noethérien~; il n'y a alors qu'un nombre fini de composantes irréductibles
$X_i$ de $X$ contenant $x$ (\hyperref[I.7.2.8.1-fr]{7.2.8.1}), et on peut
supposer que ce sont les seules rencontrant $U$, en remplaçant au besoin $U$
par un ouvert plus petit (puisqu'il n'y a qu'un nombre fini de composantes
irréductibles de $X$ rencontrant $U$, $U$ étant noethérien). On voit alors
comme ci-dessus que $f$ et $g$ coïncident dans un ouvert non vide de chacun
des $X_i$. Tenant compte du fait que chacun des $X_i$ est contenu dans
$\overline U$, considérons alors le morphisme $f_1$, défini dans un ouvert
dense de $U\cup(X-\overline U)$, égal à $g$ dans $U$ et à $f$ dans
l'intersection de $X-\overline U$ et du domaine de définition de $f$. Comme
$U\cup(X-\overline U)$ est dense dans $X$, $f_1$ et $f$ coïncident dans un
ouvert dense de $X$, et comme $f$ est une application rationnelle, $f$ est
une extension de $f_1$ (\hyperref[I.7.2.3-fr]{7.2.3}), donc est définie au
point $x$.

\subsection{Faisceau des fonctions rationnelles.}
\label{subsection:I.7.3-fr}

\begin{env}[7.3.1]
\label{I.7.3.1-fr}
Soit $X$ un préschéma. Pour tout ouvert $U\subset X$, désignons par $R(U)$
l'anneau des fonctions rationnelles sur $U$
(\hyperref[I.7.1.3-fr]{7.1.3})~; c'est une
$\Gamma(U,\mathscr{O}_X)$-algèbre. En outre, si $V\subset U$ est un second
ouvert de $X$, toute section de $\mathscr{O}_X$ au-dessus d'une partie
ouverte partout dense de $U$ donne par restriction à $V$ une section
au-dessus d'une partie ouverte partout dense de $V$, et si deux sections
coïncident au-dessus d'une partie ouverte partout dense de $U$, leurs
restrictions à $V$ coïncident au-dessus d'une partie ouverte partout dense de
$V$. On définit donc ainsi un di-homomorphisme d'algèbres
$\rho_{V,U}:R(U)\to R(V)$, et il est clair que si $U\supset V\supset W$ sont
trois ouverts de $X$, on a
$\rho_{W,U}=\rho_{W,V}\circ\rho_{V,U}$~; les $R(U)$ définissent donc un
\emph{préfaisceau} d'algèbres sur $X$.
\end{env}

\oldpage[I]{162}
\begin{definition}[7.3.2]
\label{I.7.3.2-fr}
On appelle \emph{faisceau des fonctions rationnelles} sur un préschéma $X$
et on désigne par $\mathscr{R}(X)$ la $\mathscr{O}_X$-Algèbre associée au
préfaisceau formé des $R(U)$.
\end{definition}

Pour tout préschéma $X$ et tout ouvert $U\subset X$, il est clair que le
faisceau induit $\mathscr{R}(X)|U$ n'est autre que $\mathscr{R}(U)$.

\begin{proposition}[7.3.3]
\label{I.7.3.3-fr}
Soit $X$ un préschéma tel que la famille $(X_\lambda)$ de ses composantes
irréductibles soit localement finie (ce qui est en particulier le cas lorsque
l'espace sous-jacent à $X$ est localement noethérien). Alors le
$\mathscr{O}_X$-Module $\mathscr{R}(X)$ est quasi-cohérent, et pour tout
ouvert $U$ de $X$ ne rencontrant qu'un nombre fini de composantes $X_\lambda$,
$R(U)$ est égal à $\Gamma(U,\mathscr{R}(X))$ et s'identifie canoniquement au
composé direct des anneaux locaux des points génériques des $X_\lambda$ telles
que $U\cap X_\lambda\neq\varnothing$.
\end{proposition}

On peut évidemment se limiter au cas où $X$ n'a qu'un nombre fini de
composantes irréductibles $X_i$, de points génériques $x_i$
($1\leq i\leq n$). Le fait que $R(U)$ est canoniquement identifié au composé
direct des $\mathscr{O}_{x_i}=R(X_i)$ tels que
$U\cap X_i\neq\varnothing$ résulte alors de
(\hyperref[I.7.1.7-fr]{7.1.7}). Montrons en outre que le préfaisceau
$U\to R(U)$ vérifie les axiomes des faisceaux, ce qui prouvera que
$R(U)=\Gamma(U,\mathscr{R}(X))$. En effet, il vérifie (F 1) d'après ce qui
précède. Pour voir qu'il satisfait à (F 2), considérons un recouvrement d'un
ouvert $U$ de $X$ par des ouverts $V_\alpha\subset U$~; si les
$s_\alpha\in R(V_\alpha)$ sont telles que les restrictions de $s_\alpha$ et
$s_\beta$ à $V_\alpha\cap V_\beta$ coïncident pour tout couple d'indices, on
en conclut que pour tout indice $i$ tel que $U\cap X_i\neq\varnothing$, les
composantes dans $R(X_i)$ de toutes les $s_\alpha$ telles que
$V_\alpha\cap X_i\neq\varnothing$ sont les mêmes~; désignant par $t_i$ cette
composante, il est clair que l'élément de $R(U)$ ayant les $t_i$ pour
composantes a pour restriction $s_\alpha$ à chaque $V_\alpha$. Enfin, pour
voir que $\mathscr{R}(X)$ est quasi-cohérent, on peut se limiter au cas où
$X=\operatorname{Spec}(A)$ est affine~; en prenant pour $U$ les ouverts
affines de la forme $D(f)$, où $f\in A$, il résulte de ce qui précède et de la
définition (\hyperref[I.1.3.4-fr]{1.3.4}) que l'on a
$\mathscr{R}(X)=\widetilde M$, où $M$ est somme directe des $A$-modules
$A_{x_i}$.

\begin{corollary}[7.3.4]
\label{I.7.3.4-fr}
Soit $X$ un préschéma réduit n'ayant qu'un nombre fini de composantes
irréductibles, et soient $X_i$ ($1\leq i\leq n$) les sous-préschémas fermés
réduits de $X$ ayant pour espaces sous-jacents les composantes irréductibles de
$X$ (\hyperref[I.5.2.1-fr]{5.2.1}). Si $h_i$ est l'injection canonique
$X_i\to X$, $\mathscr{R}(X)$ est alors composée directe des
$\mathscr{O}_X$-Algèbres $(h_i)_*(\mathscr{R}(X_i))$.
\end{corollary}

\begin{corollary}[7.3.5]
\label{I.7.3.5-fr}
Si $X$ est irréductible, tout $\mathscr{R}(X)$-Module quasi-cohérent
$\mathscr{F}$ est un faisceau simple.
\end{corollary}

Il suffit de montrer que tout $x\in X$ admet un voisinage $U$ tel que
$\mathscr{F}|U$ soit un faisceau simple
(\hyperref[0.3.6.2-fr]{0, 3.6.2}), autrement dit on est ramené au cas où $X$
est affine~; on peut en outre supposer que $\mathscr{F}$ est le conoyau d'un
homomorphisme
$(\mathscr{R}(X))^{(I)}\to(\mathscr{R}(X))^{(J)}$
(\hyperref[0.5.1.3-fr]{0, 5.1.3}), et tout revient à voir que
$\mathscr{R}(X)$ est un faisceau simple~; mais cela est évident puisque
$\Gamma(U,\mathscr{R}(X))=R(X)$ pour tout ouvert non vide $U$, $U$ contenant
le point générique de $X$.

\begin{corollary}[7.3.6]
\label{I.7.3.6-fr}
Si $X$ est irréductible, pour tout $\mathscr{O}_X$-Module quasi-cohérent
$\mathscr{F}$, $\mathscr{F}\otimes_{\mathscr{O}_X}\mathscr{R}(X)$ est un
faisceau simple~; si en outre $X$ est réduit (donc intègre),
$\mathscr{F}\otimes_{\mathscr{O}_X}\mathscr{R}(X)$ est isomorphe à un
faisceau de la forme $(\mathscr{R}(X))^{(I)}$.
\end{corollary}

La seconde assertion résulte de ce que $R(X)$ est alors un corps.

\begin{proposition}[7.3.7]
\label{I.7.3.7-fr}
Supposons que le préschéma $X$ soit localement intègre ou localement
\oldpage[I]{163}
noethérien. Alors $\mathscr{R}(X)$ est une $\mathscr{O}_X$-Algèbre
quasi-cohérente~; si en outre $X$ est réduit (ce qui est le cas lorsque $X$
est localement intègre), l'homomorphisme canonique
$\mathscr{O}_X\to\mathscr{R}(X)$ est injectif.
\end{proposition}

La question étant locale, la première assertion résulte de
(\hyperref[I.7.3.3-fr]{7.3.3})~; la seconde résulte aussitôt de
(\hyperref[I.7.2.3-fr]{7.2.3}).

\begin{env}[7.3.8]
\label{I.7.3.8-fr}
Soient $X$, $Y$ deux préschémas ayant chacun un nombre fini de composantes
irréductibles, et soit $f:X\to Y$ un morphisme dont la restriction à
l'ensemble des points génériques des composantes irréductibles de $X$ est une
surjection sur l'ensemble des points génériques des composantes irréductibles
de $Y$. Alors on a
\[
  f^*(\mathscr{R}(Y))=\mathscr{R}(X).
  \tag{7.3.8.1}\label{I.7.3.8.1-fr}
\]
\end{env}

En effet, on est ramené (en vertu de
(\hyperref[I.7.3.3-fr]{7.3.3})) au cas où $X$ et $Y$ sont irréductibles, de
points génériques $x$, $y$, avec $f(x)=y$~; donc
\[
  (f^*(\mathscr{R}(Y)))_x
  =\mathscr{O}_y\otimes_{\mathscr{O}_y}\mathscr{O}_x
  =\mathscr{O}_x
\]
(\hyperref[0.4.3.1-fr]{0, 4.3.1}), ce qui démontre (7.3.8.1) en vertu de
(\hyperref[I.7.3.5-fr]{7.3.5}).

\subsection{Faisceaux de torsion et faisceaux sans torsion.}
\label{subsection:I.7.4-fr}

\begin{env}[7.4.1]
\label{I.7.4.1-fr}
Soit $X$ un préschéma \emph{intègre}. Pour tout $\mathscr{O}_X$-Module
$\mathscr{F}$, l'homomorphisme canonique
$\mathscr{O}_X\to\mathscr{R}(X)$ définit par tensorisation un homomorphisme
(dit encore \emph{canonique})
\[
  \mathscr{F}\to
  \mathscr{F}\otimes_{\mathscr{O}_X}\mathscr{R}(X)
\]
qui, sur chaque fibre, n'est autre que l'homomorphisme $z\to z\otimes 1$ de
$\mathscr{F}_x$ dans
$\mathscr{F}_x\otimes_{\mathscr{O}_x}R(X)$. Le \emph{noyau $\mathscr{T}$}
de cet homomorphisme est un sous-$\mathscr{O}_X$-Module de $\mathscr{F}$,
appelé \emph{faisceau de torsion de $\mathscr{F}$}~; il est quasi-cohérent si
$\mathscr{F}$ est quasi-cohérent
(\hyperref[I.4.1.1-fr]{4.1.1} et \hyperref[I.7.3.6-fr]{7.3.6}). On dit que
$\mathscr{F}$ est \emph{sans torsion} si $\mathscr{T}=0$ et que
$\mathscr{F}$ est un \emph{faisceau de torsion} si
$\mathscr{T}=\mathscr{F}$. Pour tout $\mathscr{O}_X$-Module $\mathscr{F}$,
$\mathscr{F}/\mathscr{T}$ est sans torsion. On déduit de
(\hyperref[I.7.3.5-fr]{7.3.5}) que~:
\end{env}

\begin{proposition}[7.4.2]
\label{I.7.4.2-fr}
Si $X$ est un préschéma intègre, tout $\mathscr{O}_X$-Module quasi-cohérent
sans torsion $\mathscr{F}$ est isomorphe à un sous-faisceau $\mathscr{G}$
d'un faisceau simple de la forme $(\mathscr{R}(X))^{(I)}$, engendré (en tant
que $\mathscr{R}(X)$-Module) par $\mathscr{G}$.
\end{proposition}

Le cardinal de $I$ est appelé le \emph{rang de $\mathscr{F}$}~; pour tout
ouvert affine non vide $U$ de $X$, le rang de $\mathscr{F}$ est égal au rang
de $\Gamma(U,\mathscr{F})$ en tant que $\Gamma(U,\mathscr{O}_X)$-module,
comme on le voit aussitôt en considérant le point générique de $X$, contenu
dans $U$. En particulier~:

\begin{corollary}[7.4.3]
\label{I.7.4.3-fr}
Sur un préschéma intègre $X$, tout $\mathscr{O}_X$-Module quasi-cohérent sans
torsion et de rang $1$ (en particulier tout $\mathscr{O}_X$-Module
inversible) est isomorphe à un sous-$\mathscr{O}_X$-Module de
$\mathscr{R}(X)$, et réciproquement.
\end{corollary}

\begin{corollary}[7.4.4]
\label{I.7.4.4-fr}
Soient $X$ un préschéma intègre, $\mathscr{L}$, $\mathscr{L}'$ deux
$\mathscr{O}_X$-Modules sans torsion, $f$ (resp. $f'$) une section de
$\mathscr{L}$ (resp. $\mathscr{L}'$) au-dessus de $X$. Pour que
$f\otimes f'=0$, il faut et il suffit que l'une des sections $f$, $f'$ soit
nulle.
\end{corollary}

Soit $x$ le point générique de $X$~; on a par hypothèse
$(f\otimes f')_x=f_x\otimes f'_x=0$. Comme $\mathscr{L}_x$ et
$\mathscr{L}'_x$ s'identifient à des sous-$\mathscr{O}_x$-Modules du corps
$\mathscr{O}_x$, la relation précédente entraîne $f_x=0$ ou $f'_x=0$, et par
suite $f=0$ ou $f'=0$ puisque $\mathscr{L}$ et $\mathscr{L}'$ sont sans
torsion (\hyperref[I.7.3.5-fr]{7.3.5}).

\begin{proposition}[7.4.5]
\label{I.7.4.5-fr}
Soient $X$, $Y$ deux préschémas intègres, $f:X\to Y$ un morphisme dominant.
Pour tout $\mathscr{O}_X$-Module quasi-cohérent sans torsion $\mathscr{F}$,
$f_*(\mathscr{F})$ est un $\mathscr{O}_Y$-Module sans torsion.
\end{proposition}

\oldpage[I]{164}
\begin{proof}
Comme $f_*$ est exact à gauche
(\hyperref[0.4.2.1-fr]{0, 4.2.1}), il suffit, en vertu de
(\hyperref[I.7.4.2-fr]{7.4.2}), de prouver la proposition lorsque
$\mathscr{F}=(\mathscr{R}(X))^{(I)}$. Or, tout ouvert non vide $U$ de $Y$
contient le point générique de $Y$, donc $f^{-1}(U)$ contient le point
générique de $X$ (\hyperref[0.2.1.5-fr]{0, 2.1.5}), donc on a alors
\[
  \Gamma(U,f_*(\mathscr{F}))
  =\Gamma(f^{-1}(U),\mathscr{F})=(R(X))^{(I)}~;
\]
autrement dit, $f_*(\mathscr{F})$ est le faisceau simple de fibre
$(R(X))^{(I)}$, considéré comme $\mathscr{R}(Y)$-Module, et il est évidemment
sans torsion.
\end{proof}

\begin{proposition}[7.4.6]
\label{I.7.4.6-fr}
Soient $X$ un préschéma intègre, $x$ son point générique. Pour tout
$\mathscr{O}_X$-Module quasi-cohérent de type fini $\mathscr{F}$, les
conditions suivantes sont équivalentes~: a) $\mathscr{F}$ est un faisceau de
torsion~; b) $\mathscr{F}_x=0$~; c)
$\operatorname{Supp}(\mathscr{F})\neq X$.
\end{proposition}

En vertu de (\hyperref[I.7.3.5-fr]{7.3.5}) et de
(\hyperref[I.7.4.1-fr]{7.4.1}), les relations $\mathscr{F}_x=0$ et
$\mathscr{F}\otimes_{\mathscr{O}_X}\mathscr{R}(X)=0$ sont équivalentes, donc
a) et b) sont équivalentes~; d'autre part,
$\operatorname{Supp}(\mathscr{F})$ est fermé dans $X$
(\hyperref[0.5.2.2-fr]{0, 5.2.2}), et comme tout ouvert non vide de $X$
contient $x$, b) et c) sont équivalentes.

\begin{env}[7.4.7]
\label{I.7.4.7-fr}
On étend (par abus de langage) les définitions de
(\hyperref[I.7.4.1-fr]{7.4.1}) au cas où $X$ est un préschéma
\emph{réduit} n'ayant qu'un nombre \emph{fini} de composantes irréductibles~;
il résulte alors de (\hyperref[I.7.3.4-fr]{7.3.4}) que l'équivalence de a) et
c) dans (\hyperref[I.7.4.6-fr]{7.4.6}) est encore valable pour un tel
préschéma.
\end{env}

% FIN SOURCE ega1/ega1-7-fr.tex

%
% DEBUT SOURCE ega1/ega1-8-fr.tex
\section{Les schémas de Chevalley}
\label{section:I.8-fr}

\subsection{Anneaux locaux apparentés.}
\label{subsection:I.8.1-fr}

Pour tout anneau local $A$, nous noterons $\mathfrak m(A)$ l'idéal maximal de
$A$.

\begin{lemma}[8.1.1]
\label{I.8.1.1-fr}
Soient $A$, $B$ deux anneaux locaux tels que $A\subset B$~; les conditions
suivantes sont équivalentes~:
\begin{enumerate}
  \item[(i)] $\mathfrak m(B)\cap A=\mathfrak m(A)$~;
  \item[(ii)] $\mathfrak m(A)\subset\mathfrak m(B)$~;
  \item[(iii)] $1$ n'appartient pas à l'idéal de $B$ engendré par
    $\mathfrak m(A)$.
\end{enumerate}
\end{lemma}

Il est évident que (i) entraîne (ii) et (ii) entraîne (iii)~; enfin, si (iii)
est vérifiée, $\mathfrak m(B)\cap A$ contient $\mathfrak m(A)$ et ne contient
pas $1$, donc est égal à $\mathfrak m(A)$.

Lorsque les conditions équivalentes de
(\hyperref[I.8.1.1-fr]{8.1.1}) sont remplies, on dit que $B$ \emph{domine}
$A$~; il revient au même de dire que l'injection $A\to B$ est un
homomorphisme \emph{local}. Il est clair que, dans l'ensemble des sous-anneaux
locaux d'un anneau $R$, la relation de domination est une relation d'ordre.

\begin{env}[8.1.2]
\label{I.8.1.2-fr}
Considérons maintenant un \emph{corps} $R$. Pour tout sous-anneau $A$ de $R$,
nous désignerons par $L(A)$ l'ensemble des anneaux locaux $A_\mathfrak p$, où
$\mathfrak p$ parcourt le spectre premier de $A$~; ils sont identifiés à des
sous-anneaux de $R$ contenant $A$. Comme
$\mathfrak p=(\mathfrak p A_\mathfrak p)\cap A$, l'application
$\mathfrak p\mapsto A_\mathfrak p$ de $\operatorname{Spec}(A)$ dans $L(A)$
est bijective.
\end{env}

\begin{lemma}[8.1.3]
\label{I.8.1.3-fr}
Soient $R$ un corps, $A$ un sous-anneau de $R$. Pour qu'un sous-anneau local
$M$ de $R$ domine un anneau $A_\mathfrak p\in L(A)$, il faut et il suffit que
$A\subset M$~; l'anneau local $A_\mathfrak p$ dominé par $M$ est alors unique
et correspond à $\mathfrak p=\mathfrak m(M)\cap A$.
\end{lemma}

En effet, si $M$ domine $A_\mathfrak p$, on a
$\mathfrak m(M)\cap A_\mathfrak p=\mathfrak p A_\mathfrak p$ d'après
(\hyperref[I.8.1.1-fr]{8.1.1}), d'où l'unicité de $\mathfrak p$~; d'autre
part, si $A\subset M$, $\mathfrak m(M)\cap A=\mathfrak p$ est premier dans
$A$, et comme $A-\mathfrak p\subset M$, on a $A_\mathfrak p\subset M$ et
$\mathfrak p A_\mathfrak p\subset\mathfrak m(M)$, donc $M$ domine
$A_\mathfrak p$.

\oldpage[I]{165}
\begin{lemma}[8.1.4]
\label{I.8.1.4-fr}
Soient $R$ un corps, $M$, $N$ deux sous-anneaux locaux de $R$, $P$ le
sous-anneau de $R$ engendré par $M\cup N$. Les conditions suivantes sont
équivalentes~:
\begin{enumerate}
  \item[(i)] Il existe un idéal premier $\mathfrak p$ de $P$ tel que
    $\mathfrak m(M)=\mathfrak p\cap M$,
    $\mathfrak m(N)=\mathfrak p\cap N$.
  \item[(ii)] L'idéal $\mathfrak a$ engendré dans $P$ par
    $\mathfrak m(M)\cup\mathfrak m(N)$ est distinct de $P$.
  \item[(iii)] Il existe un sous-anneau local $Q$ de $R$ dominant à la fois
    $M$ et $N$.
\end{enumerate}
\end{lemma}

Il est clair que (i) implique (ii)~; inversement, si $\mathfrak a\ne P$,
$\mathfrak a$ est contenu dans un idéal maximal $\mathfrak n$ de $P$ et comme
$1\notin\mathfrak n$, $\mathfrak n\cap M$ contient $\mathfrak m(M)$ et est
distinct de $M$, donc $\mathfrak n\cap M=\mathfrak m(M)$ et de même
$\mathfrak n\cap N=\mathfrak m(N)$. Il est clair que si $Q$ domine $M$ et
$N$, on a $P\subset Q$, et
\[
  \mathfrak m(M)=\mathfrak m(Q)\cap M
  =(\mathfrak m(Q)\cap P)\cap M,\qquad
  \mathfrak m(N)=(\mathfrak m(Q)\cap P)\cap N,
\]
donc (iii) entraîne (i)~; la réciproque est évidente en prenant
$Q=P_\mathfrak p$.

Lorsque les conditions de (\hyperref[I.8.1.4-fr]{8.1.4}) sont satisfaites,
on dit, avec C.~Chevalley, que les anneaux locaux $M$ et $N$ sont
\emph{apparentés}.

\begin{proposition}[8.1.5]
\label{I.8.1.5-fr}
Soient $A$, $B$ deux sous-anneaux d'un corps $R$, $C$ le sous-anneau de $R$
engendré par $A\cup B$. Les conditions suivantes sont équivalentes~:
\begin{enumerate}
  \item[(i)] Pour tout anneau local $Q$ contenant $A$ et $B$, on a
    $A_\mathfrak p=B_\mathfrak q$, en posant
    $\mathfrak p=\mathfrak m(Q)\cap A$,
    $\mathfrak q=\mathfrak m(Q)\cap B$.
  \item[(ii)] Pour tout idéal premier $\mathfrak r$ de $C$, on a
    $A_\mathfrak p=B_\mathfrak q$, en posant
    $\mathfrak p=\mathfrak r\cap A$, $\mathfrak q=\mathfrak r\cap B$.
  \item[(iii)] Si $M\in L(A)$ et $N\in L(B)$ sont apparentés, ils sont
    identiques.
  \item[(iv)] On a $L(A)\cap L(B)=L(C)$.
\end{enumerate}
\end{proposition}

Les lemmes (\hyperref[I.8.1.3-fr]{8.1.3}) et
(\hyperref[I.8.1.4-fr]{8.1.4}) prouvent que (i) et (iii) sont équivalentes~;
il est clair que (i) entraîne (ii) en l'appliquant à $Q=C_\mathfrak r$~;
inversement, (ii) entraîne (i), car si $Q$ contient $A\cup B$, il contient
$C$, et si $\mathfrak r=\mathfrak m(Q)\cap C$, on a
$\mathfrak p=\mathfrak r\cap A$ et $\mathfrak q=\mathfrak r\cap B$ d'après
(\hyperref[I.8.1.3-fr]{8.1.3}). Il est immédiat que (iv) implique (i), car si
$Q$ contient $A\cup B$, il domine un anneau local
$C_\mathfrak r\in L(C)$ par (\hyperref[I.8.1.3-fr]{8.1.3})~; on a par
hypothèse $C_\mathfrak r\in L(A)\cap L(B)$, et
(\hyperref[I.8.1.1-fr]{8.1.1}) et (\hyperref[I.8.1.3-fr]{8.1.3}) prouvent que
$C_\mathfrak r=A_\mathfrak p=B_\mathfrak q$. Prouvons enfin que (iii)
entraîne (iv). Soit $Q\in L(C)$~; $Q$ domine un $M\in L(A)$ et un
$N\in L(B)$ (\hyperref[I.8.1.3-fr]{8.1.3}), donc $M$ et $N$, étant
apparentés, sont identiques par hypothèse. Comme on a alors $C\subset M$,
$M$ domine un $Q'\in L(C)$ (\hyperref[I.8.1.3-fr]{8.1.3}), donc $Q$ domine
$Q'$, ce qui (\hyperref[I.8.1.3-fr]{8.1.3}) entraîne nécessairement
$Q=Q'=M$, donc $Q\in L(A)\cap L(B)$. Inversement, si
$Q\in L(A)\cap L(B)$, on a $C\subset Q$, donc
(\hyperref[I.8.1.3-fr]{8.1.3}) $Q$ domine un
$Q''\in L(C)\subset L(A)\cap L(B)$~; $Q$ et $Q''$ étant apparentés sont
identiques, donc $Q''=Q\in L(C)$, ce qui achève la démonstration.

\subsection{Anneaux locaux d'un schéma intègre.}
\label{subsection:I.8.2-fr}

\begin{env}[8.2.1]
\label{I.8.2.1-fr}
Soient $X$ un préschéma \emph{intègre}, $R$ son corps des fonctions
rationnelles, identique à l'anneau local du point générique $a$ de $X$~; pour
tout $x\in X$, on sait que $\mathscr{O}_x$ s'identifie canoniquement à un
sous-anneau de $R$ (\hyperref[I.7.1.5-fr]{7.1.5}), et pour toute fonction
rationnelle $f\in R$, le domaine de définition $\delta(f)$ de $f$ est
l'ensemble ouvert des $x\in X$ tels que $f\in\mathscr{O}_x$. Il en résulte
(\hyperref[I.7.2.6-fr]{7.2.6}) que pour tout ouvert $U\subset X$, on a
\begin{equation}
  \Gamma(U,\mathscr{O}_X)=\bigcap_{x\in U}\mathscr{O}_x.
  \tag{8.2.1.1}\label{I.8.2.1.1-fr}
\end{equation}
\end{env}

\oldpage[I]{166}
\begin{proposition}[8.2.2]
\label{I.8.2.2-fr}
Soient $X$ un préschéma intègre, $R$ son corps des fonctions rationnelles.
Pour que $X$ soit un schéma, il faut et il suffit que la relation
«~$\mathscr{O}_x$ et $\mathscr{O}_y$ sont apparentés~»
(\hyperref[I.8.1.4-fr]{8.1.4}) entre points $x$, $y$ de $X$ implique $x=y$.
\end{proposition}

Supposons cette condition vérifiée, et montrons que $X$ est séparé. Soient
$U$ et $V$ deux ouverts affines distincts de $X$, $A$ et $B$ leurs anneaux,
identifiés à des sous-anneaux de $R$~; $U$ (resp. $V$) s'identifie donc
(\hyperref[I.8.1.2-fr]{8.1.2}) à $L(A)$ (resp. $L(B)$), et l'hypothèse
entraîne (\hyperref[I.8.1.5-fr]{8.1.5}) que si $C$ est le sous-anneau de $R$
engendré par $A\cup B$, $W=U\cap V$ s'identifie à
$L(A)\cap L(B)=L(C)$. En outre, on sait
(\cite{I-1}, p.~5-03, prop.~4 \emph{bis}) que tout sous-anneau $E$ de $R$ est
égal à l'intersection des anneaux locaux appartenant à $L(E)$~; $C$
s'identifie donc à l'intersection des anneaux $\mathscr{O}_z$ pour $z\in W$,
autrement dit (\hyperref[I.8.2.1.1-fr]{8.2.1.1}) à
$\Gamma(W,\mathscr{O}_X)$. Considérons alors le sous-préschéma induit par $X$
sur $W$~; à l'homomorphisme identique
$\varphi:C\to\Gamma(W,\mathscr{O}_X)$ correspond
(\hyperref[I.2.2.4-fr]{2.2.4}) un morphisme
$\Phi=(\psi,\theta):W\to\operatorname{Spec}(C)$~; nous allons voir que $\Phi$
est un \emph{isomorphisme} de préschémas, d'où résultera que $W$ est un ouvert
\emph{affine}. L'identification de $W$ à
$L(C)=\operatorname{Spec}(C)$ montre que $\psi$ est \emph{bijective}. D'autre
part, pour tout $x\in W$, $\theta_x^\sharp$ est l'injection
$C_\mathfrak r\to\mathscr{O}_x$, si
$\mathfrak r=\mathfrak m_x\cap C$, et par définition $C_\mathfrak r$ est
identifié à $\mathscr{O}_x$, donc $\theta_x^\sharp$ est bijective. Il reste
donc à voir que $\psi$ est un \emph{homéomorphisme}, autrement dit que pour
toute partie fermée $F\subset W$, $\psi(F)$ est fermé dans
$\operatorname{Spec}(C)$. Or, $F$ est la trace sur $W$ d'une partie fermée de
$U$, de la forme $V(\mathfrak a)$, où $\mathfrak a$ est un idéal de $A$~;
montrons que $\psi(F)=V(\mathfrak a C)$, ce qui prouvera notre assertion. En
effet, les idéaux premiers de $C$ contenant $\mathfrak a C$ sont les idéaux
premiers de $C$ contenant $\mathfrak a$, donc les idéaux de la forme
$\psi(x)=\mathfrak m_x\cap C$ où
$\mathfrak a\subset\mathfrak m_x$ et $x\in W$~; comme
$\mathfrak a\subset\mathfrak m_x$ équivaut à
$x\in V(\mathfrak a)=W\cap F$ pour $x\in U$, on a bien
$\psi(F)=V(\mathfrak a C)$.

Il s'ensuit que $X$ est séparé, car $U\cap V$ est affine et son anneau $C$ est
engendré par la réunion $A\cup B$ des anneaux de $U$ et $V$
(\hyperref[I.5.5.6-fr]{5.5.6}).

Inversement, supposons $X$ séparé, et soient $x$, $y$ deux points de $X$ tels
que $\mathscr{O}_x$ et $\mathscr{O}_y$ soient apparentés. Soit $U$ (resp.
$V$) un ouvert affine contenant $x$ (resp. $y$), d'anneau $A$ (resp. $B$)~;
on sait alors que $U\cap V$ est affine et que son anneau $C$ est engendré par
$A\cup B$ (\hyperref[I.5.5.6-fr]{5.5.6}). Si
$\mathfrak p=\mathfrak m_x\cap A$, $\mathfrak q=\mathfrak m_y\cap B$, on a
$A_\mathfrak p=\mathscr{O}_x$, $B_\mathfrak q=\mathscr{O}_y$, et comme
$A_\mathfrak p$ et $B_\mathfrak q$ sont apparentés, il existe un idéal premier
$\mathfrak r$ de $C$ tel que $\mathfrak p=\mathfrak r\cap A$,
$\mathfrak q=\mathfrak r\cap B$ (\hyperref[I.8.1.4-fr]{8.1.4}). Mais alors il
existe un point $z\in U\cap V$ tel que
$\mathfrak r=\mathfrak m_z\cap C$ puisque $U\cap V$ est affine, et on a
évidemment $x=z$, et $y=z$, d'où $x=y$.

\begin{corollary}[8.2.3]
\label{I.8.2.3-fr}
Soient $X$ un schéma intègre, $x$, $y$ deux points de $X$. Pour que
$x\in\overline{\{y\}}$, il faut et il suffit que
$\mathscr{O}_x\subset\mathscr{O}_y$, autrement dit que toute fonction
rationnelle définie en $x$ soit définie en $y$.
\end{corollary}

La condition est évidemment nécessaire puisque le domaine de définition
$\delta(f)$ d'une fonction rationnelle $f\in R$ est ouvert~; montrons qu'elle
est suffisante. Si $\mathscr{O}_x\subset\mathscr{O}_y$, il existe un idéal
premier $\mathfrak p$ de $\mathscr{O}_x$ tel que $\mathscr{O}_y$ domine
$(\mathscr{O}_x)_\mathfrak p$ (\hyperref[I.8.1.3-fr]{8.1.3})~; or
(\hyperref[I.2.4.2-fr]{2.4.2}) il existe $z\in X$ tel que
$x\in\overline{\{z\}}$ et que
$\mathscr{O}_z=(\mathscr{O}_x)_\mathfrak p$~; comme $\mathscr{O}_z$ et
$\mathscr{O}_y$ sont apparentés, on a $z=y$ par
(\hyperref[I.8.2.2-fr]{8.2.2}), d'où le corollaire.

\begin{corollary}[8.2.4]
\label{I.8.2.4-fr}
Si $X$ est un schéma intègre, l'application $x\mapsto\mathscr{O}_x$ est
injective~; autrement dit, si $x$, $y$ sont deux points distincts de $X$, il
existe une fonction rationnelle définie en l'un de ces points et non en
l'autre.
\end{corollary}

\oldpage[I]{167}
Cela résulte de (\hyperref[I.8.2.3-fr]{8.2.3}) et de l'axiome
$(T_0)$ (\hyperref[I.2.1.4-fr]{2.1.4}).

\begin{corollary}[8.2.5]
\label{I.8.2.5-fr}
Soit $X$ un schéma intègre dont l'espace sous-jacent est noethérien~;
lorsque $f$ parcourt le corps $R$ des fonctions rationnelles sur $X$, les
ensembles $\delta(f)$ engendrent la topologie de $X$.
\end{corollary}

En effet, toute partie fermée de $X$ est alors réunion finie d'ensembles
fermés irréductibles, c'est-à-dire de la forme $\overline{\{y\}}$
(\hyperref[I.2.1.5-fr]{2.1.5}). Or, si
$x\notin\overline{\{y\}}$, il existe une fonction rationnelle $f$ définie en
$x$ et non en $y$ (\hyperref[I.8.2.3-fr]{8.2.3}), autrement dit, on a
$x\in\delta(f)$ et $\delta(f)$ ne rencontre pas $\overline{\{y\}}$. Le
complémentaire de $\overline{\{y\}}$ est par suite réunion d'ensembles de la
forme $\delta(f)$, et en vertu de la première remarque, tout ouvert de $X$ est
réunion d'intersections finies d'ouverts de la forme $\delta(f)$.

\begin{env}[8.2.6]
\label{I.8.2.6-fr}
Le cor. (\hyperref[I.8.2.5-fr]{8.2.5}) montre que la topologie de $X$ est
entièrement caractérisée par la donnée de la famille d'anneaux locaux
$(\mathscr{O}_x)_{x\in X}$ ayant $R$ pour corps des fractions. Il revient au
même d'ailleurs de dire que les parties fermées de $X$ sont définies de la
façon suivante~: étant donnée une partie finie
$\{x_1,\ldots,x_n\}$ de $X$, on considère l'ensemble des $y\in X$ tels que
$\mathscr{O}_y\subset\mathscr{O}_{x_i}$ pour un indice $i$ au moins, et ces
ensembles (pour tous les choix de $\{x_1,\ldots,x_n\}$) sont les ensembles
fermés de $X$. En outre, une fois connue la topologie de $X$, le faisceau
structural $\mathscr{O}_X$ est aussi bien déterminé par la famille des
$\mathscr{O}_x$, puisque
$\Gamma(U,\mathscr{O}_X)=\bigcap_{x\in U}\mathscr{O}_x$
(\hyperref[I.8.2.1.1-fr]{8.2.1.1}). La famille
$(\mathscr{O}_x)_{x\in X}$ détermine donc complètement le préschéma $X$
lorsque $X$ est un schéma intègre dont l'espace sous-jacent est noethérien.
\end{env}

\begin{proposition}[8.2.7]
\label{I.8.2.7-fr}
Soient $X$, $Y$ deux schémas intègres, $f:X\to Y$ un morphisme dominant
(\hyperref[I.2.2.6-fr]{2.2.6}), $K$ (resp. $L$) le corps des fonctions
rationnelles sur $X$ (resp. $Y$). Alors $L$ s'identifie à un sous-corps de
$K$, et pour tout $x\in X$, $\mathscr{O}_{f(x)}$ est l'unique anneau local de
$Y$ dominé par $\mathscr{O}_x$.
\end{proposition}

En effet, si $f=(\psi,\theta)$ et si $a$ est le point générique de $X$,
$\psi(a)$ est le point générique de $Y$
(\hyperref[0.2.1.5-fr]{0, 2.1.5})~; $\theta_a^\sharp$ est par suite un
monomorphisme du corps $L=\mathscr{O}_{\psi(a)}$ dans le corps
$K=\mathscr{O}_a$. Comme tout ouvert affine non vide $U$ de $Y$ contient
$\psi(a)$, il résulte de (\hyperref[I.2.2.4-fr]{2.2.4}) que l'homomorphisme
$\Gamma(U,\mathscr{O}_Y)\to\Gamma(\psi^{-1}(U),\mathscr{O}_X)$ correspondant
à $f$ est la restriction de $\theta_a^\sharp$ à
$\Gamma(U,\mathscr{O}_Y)$. Donc, pour tout $x\in X$,
$\theta_x^\sharp$ est la restriction à $\mathscr{O}_{\psi(x)}$ de
$\theta_a^\sharp$, et est par suite un monomorphisme. On sait en outre que
$\theta_x^\sharp$ est un homomorphisme local, donc, si on identifie $L$ à un
sous-corps de $K$ par $\theta_a^\sharp$, $\mathscr{O}_{\psi(x)}$ est dominé
par $\mathscr{O}_x$ (\hyperref[I.8.1.1-fr]{8.1.1})~; c'est d'ailleurs le seul
anneau local de $Y$ dominé par $\mathscr{O}_x$ puisque deux anneaux locaux de
$Y$ qui sont apparentés sont identiques
(\hyperref[I.8.2.2-fr]{8.2.2}).

\begin{proposition}[8.2.8]
\label{I.8.2.8-fr}
Soient $X$ un préschéma irréductible, $f:X\to Y$ une immersion locale
(resp. un isomorphisme local)~; on suppose en outre le morphisme $f$ séparé.
Alors $f$ est une immersion (resp. une immersion ouverte).
\end{proposition}

Soit $f=(\psi,\theta)$~; il suffit, dans les deux cas, de prouver que $\psi$
est un \emph{homéomorphisme} de $X$ sur $\psi(X)$
(\hyperref[I.4.5.3-fr]{4.5.3}). Remplaçant $f$ par $f_{\mathrm{red}}$
(\hyperref[I.5.1.6-fr]{5.1.6} et
\hyperref[I.5.5.1-fr]{5.5.1, (vi)}), on peut supposer $X$ et $Y$ réduits. Si
$Y'$ est le sous-préschéma fermé réduit de $Y$ ayant pour espace sous-jacent
$\overline{\psi(X)}$, $f$ se factorise en
$X\xrightarrow{f'}Y'\xrightarrow{j}Y$, où $j$ est l'injection canonique
(\hyperref[I.5.2.2-fr]{5.2.2}). Il résulte de
(\hyperref[I.5.5.1-fr]{5.5.1, (v)}) que $f'$ est encore un morphisme séparé~;
en outre, $f'$ est encore

\oldpage[I]{168}
une immersion locale (resp. un isomorphisme local), car la question étant
locale sur $X$ et $Y$, on peut se borner pour le démontrer au cas où $f$ est
une immersion fermée (resp. une immersion ouverte) et notre assertion découle
alors aussitôt de (\hyperref[I.4.2.2-fr]{4.2.2}).

On peut donc supposer que $f$ est un morphisme \emph{dominant}, ce qui
entraîne que $Y$ est, lui aussi, irréductible
(\hyperref[0.2.1.5-fr]{0, 2.1.5}), donc que $X$ et $Y$ sont tous deux
\emph{intègres}. En outre, la question étant locale sur $Y$, on peut supposer
que $Y$ est un schéma affine~; comme $f$ est séparé, $X$ est un schéma
(\hyperref[I.5.5.1-fr]{5.5.1, (ii)}), et on est finalement dans les
hypothèses de (\hyperref[I.8.2.7-fr]{8.2.7}). Alors, pour tout $x\in X$,
$\theta_x^\sharp$ est injectif~; mais l'hypothèse que $f$ est une immersion
locale implique que $\theta_x^\sharp$ est surjectif
(\hyperref[I.4.2.2-fr]{4.2.2}), donc $\theta_x^\sharp$ est bijectif,
autrement dit (avec l'identification de
(\hyperref[I.8.2.7-fr]{8.2.7})) on a
$\mathscr{O}_{\psi(x)}=\mathscr{O}_x$. Cela implique par
(\hyperref[I.8.2.4-fr]{8.2.4}) que $\psi$ est une application
\emph{injective}, ce qui démontre déjà la proposition lorsque $f$ est un
isomorphisme local (\hyperref[I.4.5.3-fr]{4.5.3}). Lorsqu'on suppose seulement
que $f$ est une immersion locale, pour tout $x\in X$ il existe un voisinage
ouvert $U$ de $x$ dans $X$ et un voisinage ouvert $V$ de $\psi(x)$ dans $Y$
tels que la restriction de $\psi$ à $U$ soit un homéomorphisme de $U$ sur une
partie \emph{fermée} de $V$. Or, $U$ est dense dans $X$, donc $\psi(U)$ est
dense dans $Y$ et \emph{a fortiori} dans $V$, ce qui prouve que
$\psi(U)=V$~; comme $\psi$ est injective, $\psi^{-1}(V)=U$ et ceci achève de
montrer que $\psi$ est un homéomorphisme de $X$ sur $\psi(X)$.

\subsection{Les schémas de Chevalley.}
\label{subsection:I.8.3-fr}

\begin{env}[8.3.1]
\label{I.8.3.1-fr}
Soient $X$ un schéma intègre \emph{noethérien}, $R$ son corps des fonctions
rationnelles~; désignons par $X'$ l'ensemble des sous-anneaux locaux
$\mathscr{O}_x\subset R$, où $x$ parcourt $X$. L'ensemble $X'$ vérifie les
trois conditions suivantes~:
\begin{enumerate}
  \item[(Sch. 1)] \emph{Pour tout $M\in X'$, $R$ est le corps des fractions
    de $M$.}
  \item[(Sch. 2)] \emph{Il existe un ensemble fini de sous-anneaux
    noethériens $A_i$ de $R$ tels que $X'=\bigcup_i L(A_i)$ et que, pour tout
    couple d'indices $i$, $j$, le sous-anneau $A_{ij}$ de $R$ engendré par
    $A_i\cup A_j$ soit une algèbre de type fini sur $A_i$.}
  \item[(Sch. 3)] \emph{Deux éléments $M$, $N$ de $X'$ qui sont apparentés
    sont identiques.}
\end{enumerate}
\end{env}

On a en effet vu dans (\hyperref[I.8.2.1-fr]{8.2.1}) que (Sch. 1) est
satisfaite, et (Sch. 3) résulte de (\hyperref[I.8.2.2-fr]{8.2.2}). Pour
démontrer (Sch. 2), il suffit de recouvrir $X$ par un nombre fini d'ouverts
affines $U_i$, d'anneaux noethériens, et de prendre
$A_i=\Gamma(U_i,\mathscr{O}_X)$~; l'hypothèse que $X$ est un schéma entraîne
que $U_i\cap U_j$ est affine et que
$\Gamma(U_i\cap U_j,\mathscr{O}_X)=A_{ij}$
(\hyperref[I.5.5.6-fr]{5.5.6})~; en outre, comme l'espace $U_i$ est
noethérien, l'immersion $U_i\cap U_j\to U_i$ est de type fini
(\hyperref[I.6.3.5-fr]{6.3.5}), donc $A_{ij}$ est une $A_i$-algèbre de type
fini (\hyperref[I.6.3.3-fr]{6.3.3}).

\begin{env}[8.3.2]
\label{I.8.3.2-fr}
Les structures dont les axiomes sont (Sch. 1), (Sch. 2) et (Sch. 3)
généralisent les «~schémas~» au sens de C.~Chevalley, qui suppose en outre que
$R$ est une extension de type fini d'un corps $K$ et que les $A_i$ sont des
$K$-algèbres de type fini (ce qui rend inutile une partie de (Sch. 2))
\cite{I-1}. Inversement, si on a une telle structure sur un ensemble $X'$,
on peut lui associer un schéma intègre $X$ en utilisant les remarques de
(\hyperref[I.8.2.6-fr]{8.2.6})~: l'espace sous-jacent de $X$ est égal à $X'$
muni de la topologie définie dans (\hyperref[I.8.2.6-fr]{8.2.6}), et du
faisceau $\mathscr{O}_X$

\oldpage[I]{169}
tel que $\Gamma(U,\mathscr{O}_X)=\bigcap_{x\in U}\mathscr{O}_x$ pour tout
ouvert $U\subset X$, avec une définition évidente des homomorphismes de
restriction. Nous laissons au lecteur le soin de vérifier qu'on obtient bien
ainsi un schéma intègre, dont les anneaux locaux sont les éléments de $X'$~;
nous n'utiliserons pas ce résultat par la suite.
\end{env}

% FIN SOURCE ega1/ega1-8-fr.tex

%
% DEBUT SOURCE ega1/ega1-9-fr.tex
\section{Compléments sur les faisceaux quasi-cohérents}
\label{section:I.9-fr}

\subsection{Produit tensoriel de faisceaux quasi-cohérents.}
\label{subsection:I.9.1-fr}

\begin{proposition}[9.1.1]
\label{I.9.1.1-fr}
Soit $X$ un préschéma (resp. un préschéma localement noethérien). Soient
$\mathscr{F}$ et $\mathscr{G}$ deux $\mathscr{O}_X$-Modules quasi-cohérents
(resp. cohérents)~; alors
$\mathscr{F}\otimes_{\mathscr{O}_X}\mathscr{G}$ est quasi-cohérent (resp.
cohérent) et de type fini si $\mathscr{F}$ et $\mathscr{G}$ sont de type
fini. Si $\mathscr{F}$ admet une présentation finie et si $\mathscr{G}$ est
quasi-cohérent (resp. cohérent), alors
$\mathscr{H}\!om_{\mathscr{O}_X}(\mathscr{F},\mathscr{G})$ est quasi-cohérent
(resp. cohérent).
\end{proposition}

La question étant locale, on peut supposer $X$ affine (resp. affine
noethérien)~; en outre, si $\mathscr{F}$ est cohérent, on peut supposer qu'il
est le conoyau d'un homomorphisme
$\mathscr{O}_X^m\to\mathscr{O}_X^n$. Les assertions relatives aux faisceaux
quasi-cohérents résultent alors de
(\hyperref[I.1.3.12-fr]{1.3.12}) et
(\hyperref[I.1.3.9-fr]{1.3.9})~; les assertions relatives aux faisceaux
cohérents résultent de (\hyperref[I.1.5.1-fr]{1.5.1}) et du fait que, si $M$
et $N$ sont des modules de type fini sur un anneau noethérien $A$,
$M\otimes_A N$ et $\operatorname{Hom}_A(M,N)$ sont des $A$-modules de type
fini.

\begin{definition}[9.1.2]
\label{I.9.1.2-fr}
Soient $X$, $Y$ deux $S$-préschémas, $p$, $q$ les projections de
$X\times_S Y$, $\mathscr{F}$ (resp. $\mathscr{G}$) un
$\mathscr{O}_X$-Module (resp. un $\mathscr{O}_Y$-Module) quasi-cohérent. On
appelle produit tensoriel de $\mathscr{F}$ et $\mathscr{G}$ sur
$\mathscr{O}_S$ (ou sur $S$) et on note
$\mathscr{F}\otimes_{\mathscr{O}_S}\mathscr{G}$ (ou
$\mathscr{F}\otimes_S\mathscr{G}$) le produit tensoriel
$p^*(\mathscr{F})\otimes_{\mathscr{O}_{X\times_S Y}}q^*(\mathscr{G})$ sur le
préschéma $X\times_S Y$.
\end{definition}

Si $X_i$ ($1\leq i\leq n$) sont des $S$-préschémas, $\mathscr{F}_i$ un
$\mathscr{O}_{X_i}$-Module quasi-cohérent ($1\leq i\leq n$), on définit de
la même manière le produit tensoriel
$\mathscr{F}_1\otimes_S\mathscr{F}_2\otimes_S\cdots\otimes_S\mathscr{F}_n$
sur le préschéma
$Z=X_1\times_S X_2\times_S\cdots\times_S X_n$~; c'est un
$\mathscr{O}_Z$-module \emph{quasi-cohérent} en vertu de
(\hyperref[I.9.1.1-fr]{9.1.1}) et de
(\hyperref[0.5.1.4-fr]{0, 5.1.4})~; il est \emph{cohérent} si les
$\mathscr{F}_i$ le sont et si $Z$ est \emph{localement noethérien} en vertu
de (\hyperref[I.9.1.1-fr]{9.1.1}),
(\hyperref[0.5.3.11-fr]{0, 5.3.11}) et
(\hyperref[I.6.1.1-fr]{6.1.1}).

On notera que si on prend $X=Y=S$, la définition
(\hyperref[I.9.1.2-fr]{9.1.2}) redonne le produit tensoriel de
$\mathscr{O}_S$-Modules. En outre, comme
$q^*(\mathscr{O}_Y)=\mathscr{O}_{X\times_S Y}$
(\hyperref[0.4.3.4-fr]{0, 4.3.4}), le produit
$\mathscr{F}\otimes_S\mathscr{O}_Y$ s'identifie canoniquement à
$p^*(\mathscr{F})$, et de même
$\mathscr{O}_X\otimes_S\mathscr{G}$ s'identifie canoniquement à
$q^*(\mathscr{G})$. Plus particulièrement, si on prend $Y=S$ et qu'on
désigne par $f$ le morphisme structural $X\to Y$, on a
$\mathscr{O}_X\otimes_Y\mathscr{G}=f^*(\mathscr{G})$~: le produit tensoriel
ordinaire et l'image réciproque apparaissent donc comme des cas particuliers
du produit tensoriel général.

La déf. (\hyperref[I.9.1.2-fr]{9.1.2}) entraîne immédiatement que, pour $X$
et $Y$ fixés, $\mathscr{F}\otimes_S\mathscr{G}$ est un
\emph{bifoncteur covariant additif et exact à droite} en $\mathscr{F}$ et
$\mathscr{G}$.

\begin{proposition}[9.1.3]
\label{I.9.1.3-fr}
Soient $S$, $X$, $Y$ trois schémas affines d'anneaux respectifs, $A$, $B$,
$C$, $B$ et $C$ étant des $A$-algèbres. Soit $M$ (resp. $N$) un $B$-module
(resp. $C$-module), $\mathscr{F}=\widetilde{M}$ (resp.
$\mathscr{G}=\widetilde{N}$) le faisceau quasi-cohérent associé~; alors
$\mathscr{F}\otimes_S\mathscr{G}$ est canoniquement isomorphe au faisceau
associé au $(B\otimes_A C)$-module $M\otimes_A N$.
\end{proposition}

\begin{proof}
\oldpage[I]{170}
En effet, en vertu de (\hyperref[I.1.6.5-fr]{1.6.5}),
$\mathscr{F}\otimes_S\mathscr{G}$ est canoniquement isomorphe au faisceau
associé au $(B\otimes_A C)$-module
\[
  \bigl(M\otimes_B(B\otimes_A C)\bigr)
  \otimes_{B\otimes_A C}
  \bigl((B\otimes_A C)\otimes_C N\bigr)
\]
et en raison des isomorphismes canoniques entre produits tensoriels, ce
dernier est isomorphe à
\[
  M\otimes_B(B\otimes_A C)\otimes_C N
  =(M\otimes_B B)\otimes_A(C\otimes_C N)=M\otimes_A N.
\]
\end{proof}

\begin{proposition}[9.1.4]
\label{I.9.1.4-fr}
Soient $f:T\to X$, $g:T\to Y$ deux $S$-morphismes, $\mathscr{F}$ (resp.
$\mathscr{G}$) un $\mathscr{O}_X$-Module (resp. $\mathscr{O}_Y$-Module)
quasi-cohérent. On a alors
\[
  (f,g)_S^*(\mathscr{F}\otimes_S\mathscr{G})
  =f^*(\mathscr{F})\otimes_{\mathscr{O}_T}g^*(\mathscr{G}).
\]
\end{proposition}

\begin{proof}
Si $p$, $q$ sont les projections de $X\times_S Y$, la formule résulte en
effet des relations $(f,g)_S^*\circ p^*=f^*$ et
$(f,g)_S^*\circ q^*=g^*$ (\hyperref[0.3.5.5-fr]{0, 3.5.5}), et du fait
qu'une image réciproque d'un produit tensoriel de faisceaux algébriques est
le produit tensoriel de leurs images réciproques
(\hyperref[0.4.3.3-fr]{0, 4.3.3}).
\end{proof}

\begin{corollary}[9.1.5]
\label{I.9.1.5-fr}
Soient $f:X\to X'$, $g:Y\to Y'$ deux $S$-morphismes, $\mathscr{F}'$ (resp.
$\mathscr{G}'$) un $\mathscr{O}_{X'}$-Module (resp.
$\mathscr{O}_{Y'}$-Module) quasi-cohérent. On a alors
\[
  (f\times_S g)^*(\mathscr{F}'\otimes_S\mathscr{G}')
  =f^*(\mathscr{F}')\otimes_S g^*(\mathscr{G}').
\]
\end{corollary}

\begin{proof}
Cela résulte de (\hyperref[I.9.1.4-fr]{9.1.4}) et du fait que
$f\times_S g=(f\circ p,g\circ q)_S$, $p$ et $q$ étant les projections de
$X\times_S Y$.
\end{proof}

\begin{corollary}[9.1.6]
\label{I.9.1.6-fr}
Soient $X$, $Y$, $Z$ trois $S$-préschémas, $\mathscr{F}$ (resp.
$\mathscr{G}$, $\mathscr{H}$) un $\mathscr{O}_X$-Module (resp.
$\mathscr{O}_Y$-Module, $\mathscr{O}_Z$-Module) quasi-cohérent~; le faisceau
$\mathscr{F}\otimes_S\mathscr{G}\otimes_S\mathscr{H}$ est l'image
réciproque de
$(\mathscr{F}\otimes_S\mathscr{G})\otimes_S\mathscr{H}$ par l'isomorphisme
canonique de $X\times_S Y\times_S Z$ sur
$(X\times_S Y)\times_S Z$.
\end{corollary}

\begin{proof}
En effet, cet isomorphisme s'écrit $(p_1,p_2)_S\times_S p_3$, en désignant
par $p_1$, $p_2$, $p_3$ les projections de $X\times_S Y\times_S Z$.

De même, l'image réciproque de $\mathscr{G}\otimes_S\mathscr{F}$ par
l'isomorphisme canonique de $X\times_S Y$ sur $Y\times_S X$ est
$\mathscr{F}\otimes_S\mathscr{G}$.
\end{proof}

\begin{corollary}[9.1.7]
\label{I.9.1.7-fr}
Si $X$ est un $S$-préschéma, tout $\mathscr{O}_X$-Module quasi-cohérent
$\mathscr{F}$ est l'image réciproque de
$\mathscr{F}\otimes_S\mathscr{O}_S$ par l'isomorphisme canonique de $X$ sur
$X\times_S S$ (\hyperref[I.3.3.3-fr]{3.3.3}).
\end{corollary}

\begin{proof}
En effet, cet isomorphisme est $(1_X,\varphi)_S$, où $\varphi$ est le
morphisme structural $X\to S$, et le corollaire résulte de
(\hyperref[I.9.1.4-fr]{9.1.4}) et du fait que
$\varphi^*(\mathscr{O}_S)=\mathscr{O}_X$.
\end{proof}

\begin{env}[9.1.8]
\label{I.9.1.8-fr}
Soient $X$ un $S$-préschéma, $\mathscr{F}$ un $\mathscr{O}_X$-Module
quasi-cohérent, $\varphi:S'\to S$ un morphisme~; on désigne par
$\mathscr{F}_{(\varphi)}$ ou $\mathscr{F}_{(S')}$ le faisceau
quasi-cohérent $\mathscr{F}\otimes_S\mathscr{O}_{S'}$, sur
$X\times_S S'=X_{(\varphi)}=X_{(S')}$~; donc
$\mathscr{F}_{(S')}=p^*(\mathscr{F})$, où $p$ est la projection
$X_{(S')}\to X$.
\end{env}

\begin{proposition}[9.1.9]
\label{I.9.1.9-fr}
Soit $\varphi':S''\to S'$ un morphisme. Pour tout
$\mathscr{O}_X$-Module quasi-cohérent $\mathscr{F}$ sur le $S$-préschéma
$X$, $(\mathscr{F}_{(\varphi)})_{(\varphi')}$ est l'image réciproque de
$\mathscr{F}_{(\varphi\circ\varphi')}$ par l'isomorphisme canonique
$(X_{(\varphi)})_{(\varphi')}\xrightarrow{\sim}
X_{(\varphi\circ\varphi')}$ (\hyperref[I.3.3.9-fr]{3.3.9}).
\end{proposition}

\begin{proof}
Cela résulte aussitôt des définitions et de
(\hyperref[I.3.3.9-fr]{3.3.9}), et s'écrit encore
\[
  (\mathscr{F}\otimes_S\mathscr{O}_{S'})
  \otimes_{S'}\mathscr{O}_{S''}
  =\mathscr{F}\otimes_S\mathscr{O}_{S''}.
  \tag{9.1.9.1}\label{I.9.1.9.1-fr}
\]
\end{proof}

\begin{proposition}[9.1.10]
\label{I.9.1.10-fr}
Soient $Y$ un $S$-préschéma, $f:X\to Y$ un $S$-morphisme. Pour tout
$\mathscr{O}_Y$-Module quasi-cohérent $\mathscr{G}$ et tout morphisme
$S'\to S$, on a
\[
  (f_{(S')})^*(\mathscr{G}_{(S')})
  =(f^*(\mathscr{G}))_{(S')}.
\]
\end{proposition}

\begin{proof}
Cela résulte aussitôt de la commutativité du diagramme
\oldpage[I]{171}
\[
  \xymatrix{
    X_{(S')}\ar[r]^{f_{(S')}}\ar[d] &
    Y_{(S')}\ar[d]\\
    X\ar[r]_f &
    Y
  }
\]
\end{proof}

\begin{corollary}[9.1.11]
\label{I.9.1.11-fr}
Soient $X$, $Y$ deux $S$-préschémas, $\mathscr{F}$ (resp. $\mathscr{G}$) un
$\mathscr{O}_X$-Module (resp. $\mathscr{O}_Y$-Module) quasi-cohérent.
L'image réciproque du faisceau
$(\mathscr{F}_{(S')}\otimes_{S'}\mathscr{G}_{(S')})$ par l'isomorphisme
canonique
$(X\times_S Y)_{(S')}\xrightarrow{\sim}
(X_{(S')})\times_{S'}(Y_{(S')})$
(\hyperref[I.3.3.10-fr]{3.3.10}) est égale à
$(\mathscr{F}\otimes_S\mathscr{G})_{(S')}$.
\end{corollary}

\begin{proof}
Si $p$, $q$ sont les projections de $X\times_S Y$, l'isomorphisme en
question n'est autre que $(p_{(S')},q_{(S')})_{S'}$~; le corollaire résulte
des prop. (\hyperref[I.9.1.4-fr]{9.1.4}) et
(\hyperref[I.9.1.10-fr]{9.1.10}).
\end{proof}

\begin{proposition}[9.1.12]
\label{I.9.1.12-fr}
Avec les notations de (\hyperref[I.9.1.2-fr]{9.1.2}) soient $z$ un point de
$X\times_S Y$, $x=p(z)$, $y=q(z)$~; la fibre
$(\mathscr{F}\otimes_S\mathscr{G})_z$ est isomorphe à
$(\mathscr{F}_x\otimes_{\mathscr{O}_x}\mathscr{O}_z)
\otimes_{\mathscr{O}_z}
(\mathscr{G}_y\otimes_{\mathscr{O}_y}\mathscr{O}_z)
=\mathscr{F}_x\otimes_{\mathscr{O}_x}\mathscr{O}_z
\otimes_{\mathscr{O}_y}\mathscr{G}_y$.
\end{proposition}

\begin{proof}
Comme on peut se ramener au cas affine, la proposition résulte de la formule
(\hyperref[I.1.6.5.1-fr]{1.6.5.1}).
\end{proof}

\begin{corollary}[9.1.13]
\label{I.9.1.13-fr}
Si $\mathscr{F}$ et $\mathscr{G}$ sont de type fini, on a
\[
  \operatorname{Supp}(\mathscr{F}\otimes_S\mathscr{G})
  =p^{-1}(\operatorname{Supp}(\mathscr{F}))
  \cap q^{-1}(\operatorname{Supp}(\mathscr{G})).
\]
\end{corollary}

\begin{proof}
Comme $p^*(\mathscr{F})$ et $q^*(\mathscr{G})$ sont de type fini sur
$\mathscr{O}_{X\times_S Y}$, on est ramené, par
(\hyperref[I.9.1.12-fr]{9.1.12}) et
(\hyperref[0.1.7.5-fr]{0, 1.7.5}), au cas particulier où
$\mathscr{G}=\mathscr{O}_Y$, c'est-à-dire à démontrer la formule
\[
  \operatorname{Supp}(p^{-1}(\mathscr{F}))
  =p^{-1}(\operatorname{Supp}(\mathscr{F})).
  \tag{9.1.13.1}\label{I.9.1.13.1-fr}
\]

Le même raisonnement que dans (\hyperref[0.1.7.5-fr]{0, 1.7.5}) ramène à
vérifier que l'on a, pour tout $z\in X\times_S Y$,
$\mathscr{O}_z/\mathfrak{m}_x\mathscr{O}_z\neq0$ (avec $x=p(z)$), ce qui
découle du fait que l'homomorphisme
$\mathscr{O}_x\to\mathscr{O}_z$ est \emph{local} par hypothèse.
\end{proof}

Nous laissons au lecteur le soin d'étendre à un produit d'un nombre
quelconque de facteurs les résultats démontrés dans ce numéro pour deux
facteurs.

\subsection{Image directe d'un faisceau quasi-cohérent.}
\label{subsection:I.9.2-fr}

\begin{proposition}[9.2.1]
\label{I.9.2.1-fr}
Soit $f:X\to Y$ un morphisme de préschémas. On suppose qu'il existe un
recouvrement $(Y_\alpha)$ de $Y$ par des ouverts affines ayant la propriété
suivante~: chacun des $f^{-1}(Y_\alpha)$ admet un recouvrement fini
$(X_{\alpha i})$ par des ouverts affines contenus dans
$f^{-1}(Y_\alpha)$, tel que chacune des intersections
$X_{\alpha i}\cap X_{\alpha j}$ soit elle-même réunion finie d'ouverts
affines. Dans ces conditions, pour tout $\mathscr{O}_X$-Module
quasi-cohérent $\mathscr{F}$, $f_*(\mathscr{F})$ est un
$\mathscr{O}_Y$-Module quasi-cohérent.
\end{proposition}

\begin{proof}
La question étant locale sur $Y$, on peut supposer $Y$ égal à l'un des
$Y_\alpha$, donc supprimer les indices $\alpha$.

\begin{enumerate}
\item[a)] Supposons d'abord que les $X_i\cap X_j$ soient eux-mêmes des
ouverts \emph{affines}. Posons
$\mathscr{F}_i=\mathscr{F}|X_i$,
$\mathscr{F}_{ij}=\mathscr{F}|(X_i\cap X_j)$, et soient
$\mathscr{F}'_i$ et $\mathscr{F}'_{ij}$ les images de $\mathscr{F}_i$ et
$\mathscr{F}_{ij}$ respectivement par les restrictions de $f$ à $X_i$ et
$X_i\cap X_j$~; on sait que les $\mathscr{F}'_i$ et
$\mathscr{F}'_{ij}$ sont quasi-cohérents
(\hyperref[I.1.6.3-fr]{1.6.3}). Posons
$\mathscr{G}=\bigoplus_i\mathscr{F}'_i$,
$\mathscr{H}=\bigoplus_{i,j}\mathscr{F}'_{ij}$~; $\mathscr{G}$ et
$\mathscr{H}$ sont des $\mathscr{O}_Y$-Modules quasi-cohérents~; nous allons
définir un homomorphisme $u:\mathscr{G}\to\mathscr{H}$ tel que
$f_*(\mathscr{F})$ soit le \emph{noyau} de $u$~; il en résultera que
$f_*(\mathscr{F})$ est quasi-cohérent
(\hyperref[I.1.3.9-fr]{1.3.9}). Il suffit de définir $u$
\oldpage[I]{172}
comme homomorphisme de préfaisceaux~; tenant compte des définitions de
$\mathscr{G}$ et $\mathscr{H}$, il suffit donc, pour tout ensemble ouvert
$W\subset Y$, de définir un homomorphisme
\[
  u_W:\bigoplus_i\Gamma(f^{-1}(W)\cap X_i,\mathscr{F})
  \longrightarrow
  \bigoplus_{i,j}\Gamma(f^{-1}(W)\cap X_i\cap X_j,\mathscr{F})
\]
de façon à satisfaire aux conditions de compatibilité usuelles lorsque $W$
varie. Si, pour toute section
$s_i\in\Gamma(f^{-1}(W)\cap X_i,\mathscr{F})$, on désigne par $s_{i|j}$ sa
restriction à $f^{-1}(W)\cap X_i\cap X_j$, on posera
\[
  u_W((s_i))=(s_{i|j}-s_{j|i})
\]
et les conditions de compatibilité sont remplies de façon évidente. Pour
prouver que le noyau $\mathscr{R}$ de $u$ est $f_*(\mathscr{F})$, définissons
un homomorphisme de $f_*(\mathscr{F})$ dans $\mathscr{R}$ en faisant
correspondre à toute section $s\in\Gamma(f^{-1}(W),\mathscr{F})$ la famille
$(s_i)$, où $s_i$ est la restriction de $s$ à $f^{-1}(W)\cap X_i$~; les
axiomes (F 1) et (F 2) des faisceaux (G, II, 1.1) entraînent que cet
homomorphisme est \emph{bijectif}, ce qui termine la démonstration dans ce cas.

\item[b)] Dans le cas général, le même raisonnement s'applique une fois que
l'on a établi que les $\mathscr{F}'_{ij}$ sont quasi-cohérents. Or, par
hypothèse, $X_i\cap X_j$ est réunion finie d'ouverts affines $X_{ijk}$~; et
comme les $X_{ijk}$ sont des ouverts affines \emph{dans un schéma},
l'intersection de deux quelconques d'entre eux est encore un ouvert affine
(\hyperref[I.5.5.6-fr]{5.5.6}). On est donc ramené au premier cas, et
(9.2.1) est donc démontrée.
\end{enumerate}
\end{proof}

\begin{corollary}[9.2.2]
\label{I.9.2.2-fr}
La conclusion de (9.2.1) est valable dans chacun des cas suivants~:
\begin{enumerate}
\item[a)] $f$ est séparé et quasi-compact.
\item[b)] $f$ est séparé et de type fini.
\item[c)] $f$ est quasi-compact et l'espace sous-jacent à $X$ est localement
noethérien.
\end{enumerate}
\end{corollary}

\begin{proof}
Dans le cas a), les $X_{\alpha i}\cap X_{\alpha j}$ sont affines
(\hyperref[I.5.5.6-fr]{5.5.6}). Le cas b) est un cas particulier de a)
(\hyperref[I.6.6.3-fr]{6.6.3}). Enfin, dans le cas c), on peut se ramener au
cas où $Y$ est affine et l'espace sous-jacent à $X$ noethérien~; alors $X$
admet un recouvrement ouvert affine fini $(X_i)$, et les $X_i\cap X_j$, étant
quasi-compacts, sont réunions finies d'ouverts affines
(\hyperref[I.2.1.3-fr]{2.1.3}).
\end{proof}

\subsection{Prolongement des sections de faisceaux quasi-cohérents.}
\label{subsection:I.9.3-fr}

\begin{theorem}[9.3.1]
\label{I.9.3.1-fr}
Soit $X$ un préschéma dont l'espace sous-jacent est noethérien, ou un schéma
dont l'espace sous-jacent est quasi-compact. Soient $\mathscr{L}$ un
$\mathscr{O}_X$-Module inversible
(\hyperref[0.5.4.1-fr]{0, 5.4.1}), $f$ une section de $\mathscr{L}$ au-dessus
de $X$, $X_f$ l'ensemble ouvert des $x\in X$ où $f(x)\neq0$
(\hyperref[0.5.5.1-fr]{0, 5.5.1}), $\mathscr{F}$ un
$\mathscr{O}_X$-Module quasi-cohérent.
\begin{enumerate}
\item[(i)] Si $s\in\Gamma(X,\mathscr{F})$ est telle que $s|X_f=0$, il existe
un entier $n>0$ tel que $s\otimes f^{\otimes n}=0$.
\item[(ii)] Pour toute section $s\in\Gamma(X_f,\mathscr{F})$, il existe un
entier $n>0$ tel que $s\otimes f^{\otimes n}$ se prolonge en une section de
$\mathscr{F}\otimes\mathscr{L}^{\otimes n}$ au-dessus de $X$.
\end{enumerate}
\end{theorem}

\begin{proof}
\begin{enumerate}
\item[(i)] Comme l'espace sous-jacent à $X$ est quasi-compact, donc réunion
finie d'ouverts affines $U_i$ tels que $\mathscr{L}|U_i$ soit isomorphe à
$\mathscr{O}_X|U_i$, on est ramené au cas où $X$ est affine et
$\mathscr{L}=\mathscr{O}_X$. Dans ce cas, $f$ s'identifie à un élément de
$A(X)$ et on a $X_f=D(f)$~; $s$ s'identifie à un élément d'un
$A(X)$-module $M$ et $s|X_f$ à l'élément correspondant de $M_f$, et le
résultat est trivial, compte tenu de la définition d'un module de fractions.
\oldpage[I]{173}

\item[(ii)] De nouveau $X$ est réunion finie d'ouverts affines $U_i$
($1\leq i\leq r$) tels que $\mathscr{L}|U_i\simeq\mathscr{O}_X|U_i$, et pour
chaque $i$,
$(s\otimes f^{\otimes n})|(U_i\cap X_f)$ s'identifie par l'isomorphisme
précédent à
$(f|(U_i\cap X_f))^n(s|(U_i\cap X_f))$. On sait alors
(\hyperref[I.1.4.1-fr]{1.4.1}) qu'il existe un entier $n>0$ tel que pour
chaque $i$, $(s\otimes f^{\otimes n})|(U_i\cap X_f)$ se prolonge en une
section $s_i$ de $\mathscr{F}\otimes\mathscr{L}^{\otimes n}$ au-dessus de
$U_i$. Soit $s_{i|j}$ la restriction de $s_i$ à $U_i\cap U_j$~; on a par
définition $s_{i|j}-s_{j|i}=0$ dans $X_f\cap U_i\cap U_j$. Or, si $X$ est un
espace noethérien, $U_i\cap U_j$ est quasi-compact~; si $X$ est un schéma,
$U_i\cap U_j$ est un ouvert affine
(\hyperref[I.5.5.6-fr]{5.5.6}), donc encore quasi-compact. En vertu de (i),
il existe donc un entier $m$ (indépendant de $i$ et $j$) tel que
$(s_{i|j}-s_{j|i})\otimes f^{\otimes m}=0$. On en conclut aussitôt qu'il
existe une section $s'$ de
$\mathscr{F}\otimes\mathscr{L}^{\otimes(n+m)}$ au-dessus de $X$, induisant
$s_i\otimes f^{\otimes m}$ au-dessus de chaque $U_i$, et induisant par suite
$s\otimes f^{\otimes(n+m)}$ au-dessus de $X_f$.
\end{enumerate}
\end{proof}

Les corollaires qui suivent donnent une interprétation du th. (9.3.1) en un
langage plus algébrique~:

\begin{corollary}[9.3.2]
\label{I.9.3.2-fr}
Les hypothèses étant celles de (9.3.1), considérons l'anneau gradué
$A_* = \Gamma_*(\mathscr{L})$ et le $A_*$-module gradué
$M_* = \Gamma_*(\mathscr{L},\mathscr{F})$
(\hyperref[0.5.4.6-fr]{0, 5.4.6}). Si $f\in A_n$, où $n\in\mathbb{Z}$, alors
on a un isomorphisme canonique
\[
  \Gamma(X_f,\mathscr{F})\xrightarrow{\sim}((M_*)_f)_0
\]
(sous-groupe du module de fractions $(M_*)_f$ formé des éléments de degré
$0$).
\end{corollary}

\begin{corollary}[9.3.3]
\label{I.9.3.3-fr}
On suppose vérifiées les hypothèses de (9.3.1), et on suppose en outre que
$\mathscr{L}=\mathscr{O}_X$. Alors si on pose
$A=\Gamma(X,\mathscr{O}_X)$, $M=\Gamma(X,\mathscr{F})$, le $A_f$-module
$\Gamma(X_f,\mathscr{F})$ est canoniquement isomorphe à $M_f$.
\end{corollary}

\begin{proposition}[9.3.4]
\label{I.9.3.4-fr}
Soient $X$ un préschéma noethérien, $\mathscr{F}$ un
$\mathscr{O}_X$-Module cohérent, $\mathscr{J}$ un faisceau cohérent d'idéaux
dans $\mathscr{O}_X$ tels que le support de $\mathscr{F}$ soit contenu dans
celui de $\mathscr{O}_X/\mathscr{J}$. Alors il existe un entier $n>0$ tel que
$\mathscr{J}^n\mathscr{F}=0$.
\end{proposition}

\begin{proof}
Comme $X$ est réunion finie d'ouverts affines dont les anneaux sont
noethériens, on peut supposer $X$ affine d'anneau $A$ noethérien~; alors
$\mathscr{F}=\widetilde M$, où $M=\Gamma(X,\mathscr{F})$ est un $A$-module
de type fini, et $\mathscr{J}=\widetilde{\mathfrak{J}}$, où
$\mathfrak{J}=\Gamma(X,\mathscr{J})$ est un idéal de $A$
(\hyperref[I.1.4.1-fr]{1.4.1} et \hyperref[I.1.5.1-fr]{1.5.1}). Comme $A$
est noethérien, $\mathfrak{J}$ admet un système fini de générateurs $f_i$
($1\leq i\leq m$). Par hypothèse, toute section de $\mathscr{F}$ au-dessus
de $X$ est nulle dans chacun des $D(f_i)$~; si $s_j$ ($1\leq j\leq q$) sont
des sections de $\mathscr{F}$ engendrant $M$, il existe donc un entier $h$
indépendant de $i$ et $j$, tel que $f_i^h s_j=0$
(\hyperref[I.1.4.1-fr]{1.4.1}), donc $f_i^h s=0$ pour tout $s\in M$. On en
conclut que si $n=mh$, on a $\mathfrak{J}^nM=0$, et par suite le
$\mathscr{O}_X$-Module correspondant
$\mathscr{J}^n\mathscr{F}=\widetilde{\mathfrak{J}^nM}$
(\hyperref[I.1.3.13-fr]{1.3.13}) est nul.
\end{proof}

\begin{corollary}[9.3.5]
\label{I.9.3.5-fr}
Sous les hypothèses de (9.3.4), il existe un sous-préschéma fermé $Y$ de
$X$, dont l'espace sous-jacent est le support de
$\mathscr{O}_X/\mathscr{J}$, et tel que, si $j:Y\to X$ est l'injection
canonique, on ait $\mathscr{F}=j_*(j^*(\mathscr{F}))$.
\end{corollary}

\begin{proof}
Notons d'abord que les supports de $\mathscr{O}_X/\mathscr{J}$ et de
$\mathscr{O}_X/\mathscr{J}^n$ sont les mêmes car si
$\mathscr{J}_x=\mathscr{O}_x$, on a aussi
$\mathscr{J}_x^n=\mathscr{O}_x$, et on a d'autre part
$\mathscr{J}_x^n\subset\mathscr{J}_x$ pour tout $x\in X$. On peut donc en
vertu de (9.3.4) supposer que $\mathscr{J}\mathscr{F}=0$~; on prend alors
pour $Y$ le sous-préschéma fermé de $X$ défini par $\mathscr{J}$, et comme
$\mathscr{F}$ est alors un $(\mathscr{O}_X/\mathscr{J})$-Module, la
conclusion est immédiate.
\end{proof}

\oldpage[I]{174}

\subsection{Prolongement des faisceaux quasi-cohérents.}
\label{subsection:I.9.4-fr}

\begin{env}[9.4.1]
\label{I.9.4.1-fr}
Soient $X$ un espace topologique, $\mathscr{F}$ un faisceau d'ensembles
(resp. de groupes, d'anneaux) sur $X$, $U$ une partie ouverte de $X$,
$\psi:U\to X$ l'injection canonique, $\mathscr{G}$ un sous-faisceau de
$\mathscr{F}|U=\psi^*(\mathscr{F})$. Comme $\psi_*$ est exact à gauche,
$\psi_*(\mathscr{G})$ est un sous-faisceau de
$\psi_*(\psi^*(\mathscr{F}))$~; si on considère l'homomorphisme canonique
$\rho:\mathscr{F}\to\psi_*(\psi^*(\mathscr{F}))$
(\hyperref[0.3.5.3-fr]{0, 3.5.3}), nous désignerons par
$\overline{\mathscr{G}}$ le sous-faisceau
$\rho^{-1}(\psi_*(\mathscr{G}))$ de $\mathscr{F}$. Il résulte immédiatement
des définitions que pour tout ouvert $V$ de $X$,
$\Gamma(V,\overline{\mathscr{G}})$ est formé des sections
$s\in\Gamma(V,\mathscr{F})$ dont la restriction à $V\cap U$ soit une section
de $\mathscr{G}$ au-dessus de $V\cap U$. On a donc
$\overline{\mathscr{G}}|U=\psi^*(\overline{\mathscr{G}})=\mathscr{G}$, et
$\overline{\mathscr{G}}$ est le \emph{plus grand} sous-faisceau de
$\mathscr{F}$ induisant $\mathscr{G}$ sur $U$~; nous dirons que
$\overline{\mathscr{G}}$ est le \emph{prolongement canonique} du
sous-faisceau $\mathscr{G}$ de $\mathscr{F}|U$ en un sous-faisceau de
$\mathscr{F}$.
\end{env}

\begin{proposition}[9.4.2]
\label{I.9.4.2-fr}
Soient $X$ un préschéma, $U$ une partie ouverte de $X$ telle que l'injection
canonique $j:U\to X$ soit un morphisme quasi-compact
\emph{(ce qui sera vérifié pour \emph{tout} $U$ si l'espace sous-jacent $X$
est \emph{localement noethérien}
(\hyperref[I.6.6.4-fr]{6.6.4, (i)}))}. Alors~:
\begin{enumerate}
\item[(i)] Pour tout $(\mathscr{O}_X|U)$-Module quasi-cohérent
  $\mathscr{G}$, $j_*(\mathscr{G})$ est un $\mathscr{O}_X$-Module
  quasi-cohérent et on a
  $j_*(\mathscr{G})|U=j^*(j_*(\mathscr{G}))=\mathscr{G}$.
\item[(ii)] Pour tout $\mathscr{O}_X$-Module quasi-cohérent $\mathscr{F}$ et
  tout sous-$(\mathscr{O}_X|U)$-Module quasi-cohérent $\mathscr{G}$, le
  prolongement canonique $\overline{\mathscr{G}}$ de $\mathscr{G}$ (9.4.1)
  est un sous-$\mathscr{O}_X$-Module quasi-cohérent de $\mathscr{F}$.
\end{enumerate}
\end{proposition}

\begin{proof}
Si $j=(\psi,\theta)$ ($\psi$ étant l'injection $U\to X$ des espaces
sous-jacents), on a par définition $j_*(\mathscr{G})=\psi_*(\mathscr{G})$
pour tout $(\mathscr{O}_X|U)$-Module $\mathscr{G}$, et en outre
$j^*(\mathscr{H})=\psi^*(\mathscr{H})=\mathscr{H}|U$ pour tout
$\mathscr{O}_X$-Module $\mathscr{H}$ en raison de la définition d'un
préschéma induit sur un ouvert. (i) est donc un cas particulier de
(\hyperref[I.9.2.2-fr]{9.2.2, a)})~; pour la même raison,
$j_*(j^*(\mathscr{F}))$ est quasi-cohérent, et comme
$\overline{\mathscr{G}}$ est l'image réciproque de $j_*(\mathscr{G})$ par
l'homomorphisme $\rho:\mathscr{F}\to j_*(j^*(\mathscr{F}))$, (ii) résulte de
(\hyperref[I.4.1.1-fr]{4.1.1}).
\end{proof}

On notera que l'hypothèse que le morphisme $j:U\to X$ est quasi-compact est
aussi vérifiée lorsque l'ouvert $U$ est \emph{quasi-compact} et $X$ un
\emph{schéma}~: en effet, $U$ est alors réunion finie d'ouverts affines
$U_i$, et pour tout ouvert affine $V$ de $X$, $V\cap U_i$ est un ouvert
affine (\hyperref[I.5.5.6-fr]{5.5.6}), donc quasi-compact.

\begin{corollary}[9.4.3]
\label{I.9.4.3-fr}
Soient $X$ un préschéma, $U$ un ouvert quasi-compact de $X$ tel que le
morphisme d'injection $j:U\to X$ soit quasi-compact. Supposons en outre que
tout $\mathscr{O}_X$-Module quasi-cohérent soit limite inductive de ses
sous-$\mathscr{O}_X$-Modules quasi-cohérents de type fini
\emph{(ce qui a lieu lorsque $X$ est un \emph{schéma affine})}. Soient alors
$\mathscr{F}$ un $\mathscr{O}_X$-Module quasi-cohérent, $\mathscr{G}$ un
sous-$(\mathscr{O}_X|U)$-Module quasi-cohérent de type fini de
$\mathscr{F}|U$. Il existe alors un sous-$\mathscr{O}_X$-Module
quasi-cohérent de type fini $\mathscr{G}'$ de $\mathscr{F}$ tel que
$\mathscr{G}'|U=\mathscr{G}$.
\end{corollary}

\begin{proof}
En effet, on a $\mathscr{G}=\overline{\mathscr{G}}|U$, et
$\overline{\mathscr{G}}$ est quasi-cohérent d'après (9.4.2), donc limite
inductive de ses sous-$\mathscr{O}_X$-Modules quasi-cohérents de type fini
$\mathscr{H}_\lambda$. Par suite $\mathscr{G}$ est limite inductive des
$\mathscr{H}_\lambda|U$, donc égal à un des $\mathscr{H}_\lambda|U$
puisqu'il est de type fini (\hyperref[0.5.2.3-fr]{0, 5.2.3}).
\end{proof}

\begin{remark}[9.4.4]
\label{I.9.4.4-fr}
Supposons que pour \emph{tout} ouvert affine $U\subset X$ le morphisme
d'injection $U\to X$ soit quasi-compact. Alors si la conclusion de (9.4.3) a
lieu pour tout ouvert affine $U$ et tout sous-$(\mathscr{O}_X|U)$-Module
quasi-cohérent de type fini $\mathscr{G}$ de $\mathscr{F}|U$,
\oldpage[I]{175}
il en résulte que $\mathscr{F}$ est limite inductive de ses
sous-$\mathscr{O}_X$-Modules quasi-cohérents de type fini. En effet, pour
tout ouvert affine $U\subset X$, on a $\mathscr{F}|U=\widetilde{M}$, où $M$
est un $A(U)$-module, et comme ce dernier est limite inductive de ses
sous-modules de type fini, $\mathscr{F}|U$ est limite inductive de ses
sous-$(\mathscr{O}_X|U)$-Modules quasi-cohérents de type fini
(\hyperref[I.1.3.9-fr]{1.3.9}). Or, par hypothèse, chacun de ces sous-Modules
est induit sur $U$ par un sous-$\mathscr{O}_X$-Module quasi-cohérent de type
fini $\mathscr{G}_{\lambda,U}$ de $\mathscr{F}$. Les sommes finies des
$\mathscr{G}_{\lambda,U}$ sont encore des $\mathscr{O}_X$-Modules
quasi-cohérents de type fini, car la question est locale et le cas où $X$
est affine a été traité dans (\hyperref[I.1.3.10-fr]{1.3.10})~; il est clair
alors que $\mathscr{F}$ est limite inductive de ces sommes finies, d'où notre
assertion.
\end{remark}

\begin{corollary}[9.4.5]
\label{I.9.4.5-fr}
Sous les hypothèses de (9.4.3), pour tout $(\mathscr{O}_X|U)$-Module
quasi-cohérent de type fini $\mathscr{G}$, il existe un
$\mathscr{O}_X$-Module quasi-cohérent de type fini $\mathscr{G}'$ tel que
$\mathscr{G}'|U=\mathscr{G}$.
\end{corollary}

\begin{proof}
Comme $\mathscr{F}=j_*(\mathscr{G})$ est quasi-cohérent (9.4.2) et
$\mathscr{F}|U=\mathscr{G}$, il suffit d'appliquer (9.4.3) à $\mathscr{F}$.
\end{proof}

\begin{lemma}[9.4.6]
\label{I.9.4.6-fr}
Soient $X$ un préschéma, $L$ un ensemble bien ordonné,
$(V_\lambda)_{\lambda\in L}$ un recouvrement de $X$ par des ouverts affines,
$U$ un ouvert de $X$~; pour tout $\lambda\in L$, on pose
$W_\lambda=\bigcup_{\mu<\lambda}V_\mu$. On suppose que~:
\begin{enumerate}
  \item[1\textsuperscript{o}] Pour tout $\lambda\in L$,
    $V_\lambda\cap W_\lambda$ soit quasi-compact~;
  \item[2\textsuperscript{o}] Le morphisme d'immersion $U\to X$ est
    quasi-compact.
\end{enumerate}
Alors, pour tout $\mathscr{O}_X$-Module quasi-cohérent $\mathscr{F}$ et tout
sous-$(\mathscr{O}_X|U)$-Module quasi-cohérent de type fini $\mathscr{G}$ de
$\mathscr{F}|U$, il existe un sous-$\mathscr{O}_X$-Module quasi-cohérent de
type fini $\mathscr{G}'$ de $\mathscr{F}$ tel que
$\mathscr{G}'|U=\mathscr{G}$.
\end{lemma}

\begin{proof}
Posons $U_\lambda=U\cup W_\lambda$~; on va définir par récurrence transfinie
une famille $(\mathscr{G}'_\lambda)$, où $\mathscr{G}'_\lambda$ est un
sous-$(\mathscr{O}_X|U_\lambda)$-Module quasi-cohérent de type fini de
$\mathscr{F}|U_\lambda$, tel que
$\mathscr{G}'_\lambda|U_\mu=\mathscr{G}'_\mu$ pour $\mu<\lambda$ et
$\mathscr{G}'_\lambda|U=\mathscr{G}$. L'unique sous-$\mathscr{O}_X$-Module
$\mathscr{G}'$ de $\mathscr{F}$ tel que
$\mathscr{G}'|U_\lambda=\mathscr{G}'_\lambda$ pour tout $\lambda\in L$
(\hyperref[0.3.3.1-fr]{0, 3.3.1}) répondra à la question. Supposons donc les
$\mathscr{G}'_\mu$ définis et ayant les propriétés précédentes pour
$\mu<\lambda$~; si $\lambda$ n'a pas de prédécesseur on prendra pour
$\mathscr{G}'_\lambda$ l'unique sous-$(\mathscr{O}_X|U_\lambda)$-Module de
$\mathscr{F}|U_\lambda$ tel que
$\mathscr{G}'_\lambda|U_\mu=\mathscr{G}'_\mu$ pour tout $\mu<\lambda$, ce
qui est licite puisque les $U_\mu$ avec $\mu<\lambda$ forment alors un
recouvrement de $U_\lambda$. Si au contraire $\lambda=\mu+1$, on a
$U_\lambda=U_\mu\cup V_\mu$, et il suffira de définir un
sous-$(\mathscr{O}_X|V_\mu)$-Module quasi-cohérent de type fini
$\mathscr{G}''_\mu$ de $\mathscr{F}|V_\mu$ tel que
\[
  \mathscr{G}''_\mu|(U_\mu\cap V_\mu)
  =\mathscr{G}'_\mu|(U_\mu\cap V_\mu)~;
\]
on prendra ensuite pour $\mathscr{G}'_\lambda$ le
sous-$(\mathscr{O}_X|U_\lambda)$-Module de $\mathscr{F}|U_\lambda$ tel que
$\mathscr{G}'_\lambda|U_\mu=\mathscr{G}'_\mu$ et
$\mathscr{G}'_\lambda|V_\mu=\mathscr{G}''_\mu$
(\hyperref[0.3.3.1-fr]{0, 3.3.1}). Or, comme $V_\mu$ est affine, l'existence
de $\mathscr{G}''_\mu$ est assurée par (9.4.3) dès que l'on a démontré que
$U_\mu\cap V_\mu$ est quasi-compact~; mais $U_\mu\cap V_\mu$ est réunion de
$U\cap V_\mu$ et de $W_\mu\cap V_\mu$, qui sont tous deux quasi-compacts en
vertu de l'hypothèse.
\end{proof}

\begin{theorem}[9.4.7]
\label{I.9.4.7-fr}
Soient $X$ un préschéma, $U$ un ouvert de $X$. On suppose vérifiée l'une des
conditions suivantes~:
\begin{enumerate}
  \item[(a)] L'espace sous-jacent à $X$ est localement noethérien.
  \item[(b)] $X$ est un schéma quasi-compact et $U$ un ouvert quasi-compact.
\end{enumerate}
Pour tout $\mathscr{O}_X$-Module quasi-cohérent $\mathscr{F}$ et tout
sous-$(\mathscr{O}_X|U)$-Module quasi-cohérent de type fini $\mathscr{G}$ de
$\mathscr{F}|U$, il existe alors un sous-$\mathscr{O}_X$-Module
quasi-cohérent de type fini $\mathscr{G}'$ de $\mathscr{F}$ tel que
$\mathscr{G}'|U=\mathscr{G}$.
\end{theorem}

\oldpage[I]{176}

\begin{proof}
Soit $(V_\lambda)_{\lambda\in L}$ un recouvrement de $X$ par des ouverts
affines, $L$ étant supposé fini dans le cas b)~; $L$ étant muni d'une
structure d'ensemble bien ordonné, il suffit de vérifier que les conditions
du lemme (9.4.6) sont satisfaites. C'est évident dans l'hypothèse a), les
espaces $V_\lambda$ étant noethériens. Dans l'hypothèse b), les
$V_\lambda\cap V_\mu$ sont affines (\hyperref[I.5.5.6-fr]{5.5.6}), donc
quasi-compacts, et comme $L$ est fini, $V_\lambda\cap W_\lambda$ est
quasi-compact. D'où le théorème.
\end{proof}

\begin{corollary}[9.4.8]
\label{I.9.4.8-fr}
Sous les hypothèses de (9.4.7), pour tout $(\mathscr{O}_X|U)$-Module
quasi-cohérent de type fini $\mathscr{G}$, il existe un
$\mathscr{O}_X$-Module quasi-cohérent de type fini $\mathscr{G}'$ tel que
$\mathscr{G}'|U=\mathscr{G}$.
\end{corollary}

\begin{proof}
Il suffit d'appliquer (9.4.7) à $\mathscr{F}=j_*(\mathscr{G})$, qui est
quasi-cohérent (9.4.2) et tel que $\mathscr{F}|U=\mathscr{G}$.
\end{proof}

\begin{corollary}[9.4.9]
\label{I.9.4.9-fr}
Soit $X$ un préschéma dont l'espace sous-jacent est localement noethérien, ou
un schéma quasi-compact. Alors tout $\mathscr{O}_X$-Module quasi-cohérent est
limite inductive de ses sous-$\mathscr{O}_X$-Modules quasi-cohérents de type
fini.
\end{corollary}

\begin{proof}
Cela résulte de (9.4.7) et de la remarque (9.4.4).
\end{proof}

\begin{corollary}[9.4.10]
\label{I.9.4.10-fr}
Sous les hypothèses de (9.4.9), si un $\mathscr{O}_X$-Module quasi-cohérent
$\mathscr{F}$ est tel que tout sous-$\mathscr{O}_X$-Module quasi-cohérent de
type fini de $\mathscr{F}$ soit engendré par ses sections au-dessus de $X$,
alors $\mathscr{F}$ est engendré par ses sections au-dessus de $X$.
\end{corollary}

\begin{proof}
En effet, soient $U$ un voisinage ouvert affine d'un point $x\in X$, et soit
$s$ une section de $\mathscr{F}$ au-dessus de $U$~; le
sous-$\mathscr{O}_X$-Module $\mathscr{G}$ de $\mathscr{F}|U$ engendré par
$s$ est quasi-cohérent et de type fini, donc il existe un
sous-$\mathscr{O}_X$-Module quasi-cohérent de type fini $\mathscr{G}'$ de
$\mathscr{F}$ tel que $\mathscr{G}'|U=\mathscr{G}$ (9.4.7). Par hypothèse,
il y a donc un nombre fini de sections $t_i$ de $\mathscr{G}'$ au-dessus de
$X$ et des sections $a_i$ de $\mathscr{O}_X$ au-dessus d'un voisinage
$V\subset U$ de $x$ telles que
$s|V=\sum_i a_i\mathbin{\cdot}(t_i|V)$, ce qui démontre le corollaire.
\end{proof}

\subsection{Image fermée d'un préschéma. Adhérence d'un sous-préschéma.}
\label{subsection:I.9.5-fr}

\begin{proposition}[9.5.1]
\label{I.9.5.1-fr}
Soit $f:X\to Y$ un morphisme de préschémas tel que $f_*(\mathscr{O}_X)$ soit
un $\mathscr{O}_Y$-Module quasi-cohérent (ce qui a lieu si $f$ est
quasi-compact et si de plus $f$ est séparé ou $X$ localement noethérien
(\hyperref[I.9.2.2-fr]{9.2.2})). Alors il existe un plus petit sous-préschéma
$Y'$ de $Y$ tel que l'injection canonique $j:Y'\to Y$ majore $f$ (ou, ce qui
revient au même (\hyperref[I.4.4.1-fr]{4.4.1}), tel que le sous-préschéma
$f^{-1}(Y')$ de $X$ soit \emph{identique} à $X$).
\end{proposition}

De façon plus précise~:

\begin{corollary}[9.5.2]
\label{I.9.5.2-fr}
Sous les conditions de (9.5.1), soit $f=(\psi,\theta)$ et soit $\mathscr{J}$
le noyau (quasi-cohérent) de l'homomorphisme
$\theta:\mathscr{O}_Y\to f_*(\mathscr{O}_X)$. Alors le sous-préschéma fermé
$Y'$ de $Y$ défini par $\mathscr{J}$ vérifie les conditions de (9.5.1).
\end{corollary}

\begin{proof}
Comme le foncteur $\psi^*$ est exact, la factorisation canonique
$\theta:\mathscr{O}_Y\to\mathscr{O}_Y/\mathscr{J}
\xrightarrow{\theta'}\psi_*(\mathscr{O}_X)$ donne
(\hyperref[0.3.5.4.3-fr]{0, 3.5.4.3}) une factorisation
\[
  \theta^\sharp:\psi^*(\mathscr{O}_Y)
  \to\psi^*(\mathscr{O}_Y)/\psi^*(\mathscr{J})
  \xrightarrow{{\theta'}^\sharp}\mathscr{O}_X~;
\]
comme pour tout $x\in X$, $\theta_x^\sharp$ est un homomorphisme local, il
en est de même de ${\theta'_x}^\sharp$~; si on désigne par $\psi_0$
l'application continue $\psi$ considérée comme application de $X$ dans
$X'$, par $\theta_0$ la restriction
$\theta'|X':(\mathscr{O}_Y/\mathscr{J})|X'
\to\psi_*(\mathscr{O}_X)|X'=(\psi_0)_*(\mathscr{O}_X)$, on voit donc que
$f_0=(\psi_0,\theta_0)$ est un morphisme de préschémas $X\to X'$
(\hyperref[I.2.2.1-fr]{2.2.1}) tel que $f=j\circ f_0$. Si maintenant $X''$
\oldpage[I]{177}
est un second sous-préschéma fermé de $Y$, défini par un faisceau
quasi-cohérent d'idéaux $\mathscr{J}'$ de $\mathscr{O}_Y$, et tel que
l'injection $j':X''\to Y$ majore $f$, on doit tout d'abord avoir
$X''\supset\psi(X)$, donc $X'\subset X''$ puisque $X''$ est fermé. En outre,
pour tout $y\in X''$, $\theta$ doit se factoriser en
$\mathscr{O}_y\to\mathscr{O}_y/\mathscr{J}'_y
\to(\psi_*(\mathscr{O}_X))_y$, ce qui par définition entraîne
$\mathscr{J}'_y\subset\mathscr{J}_y$, et par suite $X'$ est un
sous-préschéma fermé de $X''$ (\hyperref[I.4.1.10-fr]{4.1.10}).
\end{proof}

\begin{definition}[9.5.3]
\label{I.9.5.3-fr}
Lorsqu'il existe un plus petit sous-préschéma fermé $Y'$ de $Y$ tel que
l'injection canonique $j:Y'\to Y$ majore $f$, on dit que $Y'$ est le
préschéma image fermée de $X$ par le morphisme $f$.
\end{definition}

\begin{proposition}[9.5.4]
\label{I.9.5.4-fr}
Si $f_*(\mathscr{O}_X)$ est un $\mathscr{O}_Y$-Module quasi-cohérent,
l'espace sous-jacent de l'image fermée de $X$ par $f$ est l'adhérence
$\overline{f(X)}$ dans $Y$.
\end{proposition}

\begin{proof}
Comme le support de $f_*(\mathscr{O}_X)$ est contenu dans
$\overline{f(X)}$, on a (avec les notations de
(\hyperref[I.9.5.2-fr]{9.5.2})) $\mathscr{J}_y=\mathscr{O}_y$ pour
$y\notin\overline{f(X)}$, donc le support de
$\mathscr{O}_Y/\mathscr{J}$ est contenu dans $\overline{f(X)}$. En outre, ce
support est fermé et contient $f(X)$~: en effet, si $y\in f(X)$, l'élément
unité de l'anneau $(\psi_*(\mathscr{O}_X))_y$ n'est pas nul, étant le germe
en $y$ de la section
\[
  1\in\Gamma(X,\mathscr{O}_X)=\Gamma(Y,\psi_*(\mathscr{O}_X))~;
\]
comme il est l'image par $\theta$ de l'élément unité de $\mathscr{O}_y$, ce
dernier n'appartient pas à $\mathscr{J}_y$, donc
$\mathscr{O}_y/\mathscr{J}_y\neq0$~; ceci achève la démonstration.
\end{proof}

\begin{proposition}[9.5.5]
\label{I.9.5.5-fr}
\emph{(transitivité des images fermées).} Soient $f:X\to Y$ et $g:Y\to Z$
deux morphismes de préschémas~; on suppose que l'image fermée $Y'$ de $X$
par $f$ existe, et que, si $g'$ est la restriction de $g$ à $Y'$, l'image
fermée $Z'$ de $Y'$ par $g'$ existe. Alors l'image fermée de $X$ par
$g\circ f$ existe et est égale à $Z'$.
\end{proposition}

\begin{proof}
Il suffit (\hyperref[I.9.5.1-fr]{9.5.1}) de montrer que $Z'$ est le plus
petit sous-préschéma fermé $Z_1$ de $Z$ tel que le sous-préschéma fermé
$(g\circ f)^{-1}(Z_1)$ de $X$ (égal à $f^{-1}(g^{-1}(Z_1))$ par
(\hyperref[I.4.4.2-fr]{4.4.2})) soit égal à $X$~; il revient au même de dire
que $Z'$ est le plus petit sous-préschéma fermé de $Z$ tel que $f$ soit
majoré par l'injection $g^{-1}(Z_1)\to Y$
(\hyperref[I.4.4.1-fr]{4.4.1}). Or, en vertu de l'existence de l'image
fermée $Y'$, tout $Z_1$ ayant cette propriété est tel que $g^{-1}(Z_1)$
majore $Y'$, ce qui équivaut à dire que
$j^{-1}(g^{-1}(Z_1))=g'^{-1}(Z_1)=Y'$, en désignant par $j$ l'injection
$Y'\to Y$. Par définition de $Z'$, on en conclut bien que $Z'$ est le plus
petit sous-préschéma fermé de $Z$ vérifiant la condition précédente.
\end{proof}

\begin{corollary}[9.5.6]
\label{I.9.5.6-fr}
Soit $f:X\to Y$ un $S$-morphisme tel que $Y$ soit l'image fermée de $X$ par
$f$. Soit $Z$ un $S$-schéma~; si deux $S$-morphismes $g_1$, $g_2$ de $Y$
dans $Z$ sont tels que $g_1\circ f=g_2\circ f$, alors $g_1=g_2$.
\end{corollary}

\begin{proof}
Soit $h=(g_1,g_2)_S:Y\to Z\times_S Z$~; comme la diagonale
$T=\Delta_Z(Z)$ est un sous-préschéma fermé de $Z\times_S Z$,
$Y'=h^{-1}(T)$ est un sous-préschéma fermé de $Y$
(\hyperref[I.4.4.1-fr]{4.4.1}). Posons $u=g_1\circ f=g_2\circ f$~; on a
alors par définition du produit $h'=h\circ f=(u,u)_S$, donc
$h\circ f=\Delta_Z\circ u$~; comme $\Delta_Z^{-1}(T)=Z$, on a
$h'^{-1}(T)=u^{-1}(Z)=X$, donc $f^{-1}(Y')=X$. On en conclut
(\hyperref[I.4.4.1-fr]{4.4.1}) que l'injection canonique $Y'\to Y$ majore
$f$, donc $Y'=Y$ par hypothèse~; par suite
(\hyperref[I.4.4.1-fr]{4.4.1}), $h$ se factorise en $\Delta_Z\circ v$, où
$v$ est un morphisme $Y\to Z$, ce qui entraîne $g_1=g_2=v$.
\end{proof}

\begin{remark}[9.5.7]
\label{I.9.5.7-fr}
Si $X$ et $Y$ sont des $S$-schémas, la prop.
(\hyperref[I.9.5.6-fr]{9.5.6}) signifie que
\oldpage[I]{178}
lorsque $Y$ est l'image fermée de $X$ par $f$, $f$ est un
\emph{épimorphisme} dans la catégorie des $S$-schémas (T, 1.1). Nous
démontrerons au chap. V que, réciproquement, si l'image fermée $Y'$ de $X$
par $f$ existe et si $f$ est un épimorphisme de $S$-schémas, alors on a
nécessairement $Y'=Y$.
\end{remark}

\begin{proposition}[9.5.8]
\label{I.9.5.8-fr}
Supposons vérifiées les hypothèses de (\hyperref[I.9.5.1-fr]{9.5.1}), et
soit $Y'$ l'image fermée de $X$ par $f$. Pour tout ouvert $V$ de $Y$, soit
$f_V:f^{-1}(V)\to V$ la restriction de $f$~; alors l'image fermée de
$f^{-1}(V)$ par $f_V$ dans $V$ existe et est égale au préschéma induit par
$Y'$ sur l'ouvert $V\cap Y'$ de $Y'$ (autrement dit, au sous-préschéma
$\inf(V,Y')$ de $Y$ (\hyperref[I.4.4.3-fr]{4.4.3})).
\end{proposition}

\begin{proof}
Posons $X'=f^{-1}(V)$~; comme l'image directe de $\mathscr{O}_{X'}$ par
$f_V$ n'est autre que la restriction de $f_*(\mathscr{O}_X)$ à $V$, il est
clair que le noyau $\mathscr{J}'$ de l'homomorphisme
$\mathscr{O}_V\to(f_V)_*(\mathscr{O}_{X'})$ est la restriction de
$\mathscr{J}$ à $V$, d'où aussitôt la proposition.
\end{proof}

On notera que ce résultat s'interprète en disant que la formation de l'image
fermée commute à une extension $Y_1\to Y$ du préschéma de base, qui est une
\emph{immersion ouverte}. Nous verrons au chap. IV qu'il en est de même pour
une extension $Y_1\to Y$ qui est un morphisme \emph{plat}, pourvu que $f$
soit séparé et quasi-compact.

\begin{proposition}[9.5.9]
\label{I.9.5.9-fr}
Soit $f:X\to Y$ un morphisme tel que l'image fermée $Y'$ de $X$ par $f$
existe.
\begin{enumerate}
  \item[(i)] Si $X$ est réduit, il en est de même de $Y'$.
  \item[(ii)] Si on suppose vérifiées les hypothèses de
    (\hyperref[I.9.5.1-fr]{9.5.1}) et si $X$ est irréductible (resp.
    intègre), il en est de même de $Y'$.
\end{enumerate}
\end{proposition}

\begin{proof}
Par hypothèse, le morphisme $f$ se factorise en
$X\xrightarrow{g}Y'\xrightarrow{j}Y$, où $j$ est l'injection canonique.
Comme $X$ est réduit, $g$ se factorise en
$X\xrightarrow{h}Y'_{\mathrm{red}}\xrightarrow{j'}Y'$, où $j'$ est
l'injection canonique (\hyperref[I.5.2.2-fr]{5.2.2}), et il résulte alors de
la définition de $Y'$ que $Y'_{\mathrm{red}}=Y'$. Si de plus les conditions
de (\hyperref[I.9.5.1-fr]{9.5.1}) sont remplies, il résulte de
(\hyperref[I.9.5.4-fr]{9.5.4}) que $f(X)$ est dense dans $Y'$~; si $X$ est
irréductible, il en est donc de même de $Y'$
(\hyperref[0.2.1.5-fr]{0, 2.1.5}). L'assertion relative aux préschémas
intègres résulte de la conjonction des deux autres.
\end{proof}

\begin{proposition}[9.5.10]
\label{I.9.5.10-fr}
Soit $Y$ un sous-préschéma d'un préschéma $X$, tel que l'injection canonique
$i:Y\to X$ soit un morphisme quasi-compact. Il existe alors un plus petit
sous-préschéma fermé $\overline{Y}$ de $X$ majorant $Y$~; son espace
sous-jacent est l'adhérence de celui de $Y$~; ce dernier est ouvert dans son
adhérence, et le préschéma $Y$ est induit sur cet ouvert par $\overline{Y}$.
\end{proposition}

\begin{proof}
Il suffit d'appliquer (\hyperref[I.9.5.1-fr]{9.5.1}) à l'injection $j$, qui
est séparée (\hyperref[I.5.5.1-fr]{5.5.1}) et quasi-compacte par hypothèse~;
(\hyperref[I.9.5.1-fr]{9.5.1}) prouve donc l'existence de $\overline{Y}$ et
(\hyperref[I.9.5.4-fr]{9.5.4}) montre que son espace sous-jacent est
l'adhérence de $Y$ dans $X$~; comme $Y$ est localement fermé dans $X$, il
est ouvert dans $\overline{Y}$, et la dernière assertion provient de
(\hyperref[I.9.5.8-fr]{9.5.8}) appliqué à un ouvert $V$ de $X$ tel que $Y$
soit fermé dans $V$.
\end{proof}

Avec ces notations, si l'injection $V\to X$ est quasi-compacte, et si
$\mathscr{J}$ est le faisceau quasi-cohérent d'idéaux de
$\mathscr{O}_X|V$ définissant le sous-préschéma fermé $Y$ de $V$, il résulte
de (\hyperref[I.9.5.1-fr]{9.5.1}) que le faisceau quasi-cohérent d'idéaux de
$\mathscr{O}_X$ définissant $\overline{Y}$ est le prolongement canonique
(\hyperref[I.9.4.1-fr]{9.4.1}) $\overline{\mathscr{J}}$ de $\mathscr{J}$,
car c'est évidemment le plus grand sous-faisceau d'idéaux quasi-cohérent de
$\mathscr{O}_X$ induisant $\mathscr{J}$ sur $V$.
\oldpage[I]{179}
\begin{corollary}[9.5.11]
\label{I.9.5.11-fr}
Sous les hypothèses de (\hyperref[I.9.5.10-fr]{9.5.10}), toute section de
$\mathscr{O}_{\overline{Y}}$ au-dessus d'un ouvert $V$ de $\overline{Y}$
qui est nulle dans $V\cap Y$ est nulle.
\end{corollary}

\begin{proof}
En vertu de (\hyperref[I.9.5.8-fr]{9.5.8}), on peut se ramener au cas où
$V=\overline{Y}$. Si l'on tient compte de ce que les sections de
$\mathscr{O}_{\overline{Y}}$ au-dessus de $\overline{Y}$ correspondent
canoniquement aux $\overline{Y}$-sections de
$\overline{Y}\otimes_{\mathbb{Z}}\mathbb{Z}[T]$
(\hyperref[I.3.3.15-fr]{3.3.15}) et de ce que ce dernier est séparé sur
$\overline{Y}$, le corollaire apparaît comme un cas particulier de
(\hyperref[I.9.5.6-fr]{9.5.6}).
\end{proof}

Lorsqu'il existe un plus petit sous-préschéma fermé $Y$ de $X$
majorant un sous-préschéma $Y$ de $X$, on dit que $Y'$ est
l'\emph{adhérence} de $Y$ dans $X$, lorsqu'il n'en résulte pas de
confusion.

\subsection{Faisceaux quasi-cohérents d'algèbres~; changement du faisceau
structural.}
\label{subsection:I.9.6-fr}

\begin{proposition}[9.6.1]
\label{I.9.6.1-fr}
Soient $X$ un préschéma, $\mathscr{B}$ une $\mathscr{O}_X$-Algèbre
quasi-cohérente (\hyperref[0.5.1.3-fr]{0, 5.1.3}). Pour qu'un
$\mathscr{B}$-Module $\mathscr{F}$ soit quasi-cohérent (sur l'espace
annelé $(X,\mathscr{B})$), il faut et il suffit que $\mathscr{F}$ soit un
$\mathscr{O}_X$-Module quasi-cohérent.
\end{proposition}

\begin{proof}
Comme la question est locale, on peut supposer $X$ affine d'anneau $A$, et
alors $\mathscr{B}=\widetilde{B}$, où $B$ est une $A$-algèbre
(\hyperref[I.1.4.3-fr]{1.4.3}). Si $\mathscr{F}$ est quasi-cohérent sur
l'espace annelé $(X,\mathscr{B})$, on peut aussi supposer que
$\mathscr{F}$ est le conoyau d'un $\mathscr{B}$-homomorphisme
$\mathscr{B}^{(I)}\to\mathscr{B}^{(J)}$~; comme cet homomorphisme est
aussi un $\mathscr{O}_X$-homomorphisme de $\mathscr{O}_X$-Modules, et que
$\mathscr{B}^{(I)}$ et $\mathscr{B}^{(J)}$ sont des
$\mathscr{O}_X$-Modules quasi-cohérents
(\hyperref[I.1.3.9-fr]{1.3.9, (ii)}), $\mathscr{F}$ est aussi un
$\mathscr{O}_X$-Module quasi-cohérent
(\hyperref[I.1.3.9-fr]{1.3.9, (i)}).

Inversement, si $\mathscr{F}$ est un $\mathscr{O}_X$-Module
quasi-cohérent, on a $\mathscr{F}=\widetilde{M}$, où $M$ est un $B$-module
(\hyperref[I.1.4.3-fr]{1.4.3})~; $M$ est isomorphe à un conoyau de
$B$-homomorphisme $B^{(I)}\to B^{(J)}$, donc $\mathscr{F}$ est un
$\mathscr{B}$-Module isomorphe au conoyau de l'homomorphisme correspondant
$\mathscr{B}^{(I)}\to\mathscr{B}^{(J)}$
(\hyperref[I.1.3.13-fr]{1.3.13}), ce qui achève la démonstration.
\end{proof}

En particulier, si $\mathscr{F}$ et $\mathscr{G}$ sont deux
$\mathscr{B}$-Modules quasi-cohérents,
$\mathscr{F}\otimes_{\mathscr{B}}\mathscr{G}$ est un
$\mathscr{B}$-Module quasi-cohérent~; il en est de même de
$\mathscr{H}\!om_{\mathscr{B}}(\mathscr{F},\mathscr{G})$ lorsqu'on
suppose en outre que $\mathscr{F}$ admet une présentation finie
(\hyperref[I.1.3.13-fr]{1.3.13}).

\begin{env}[9.6.2]
\label{I.9.6.2-fr}
Étant donné un préschéma $X$, nous dirons qu'une
$\mathscr{O}_X$-Algèbre quasi-cohérente $\mathscr{B}$ est de
\emph{type fini} si pour tout $x\in X$, il existe un voisinage ouvert
\emph{affine} $U$ de $x$ tel que $\Gamma(U,\mathscr{B})=B$ soit une
algèbre de type fini sur $\Gamma(U,\mathscr{O}_X)=A$. On a alors
$\mathscr{B}|U=\widetilde{B}$ et pour tout $f\in A$, la
$(\mathscr{O}_X|D(f))$-Algèbre induite $\mathscr{B}|D(f)$ est de type
fini, car elle est isomorphe à $\widetilde{B_f}$, et
$B_f=B\otimes_A A_f$ est évidemment une algèbre de type fini sur $A_f$.
Comme les $D(f)$ forment une base de la topologie de $U$, on en conclut que
si $\mathscr{B}$ est une $\mathscr{O}_X$-Algèbre quasi-cohérente de type
fini, pour tout ouvert $V$ de $X$, $\mathscr{B}|V$ est une
$(\mathscr{O}_X|V)$-Algèbre quasi-cohérente de type fini.
\end{env}

\begin{proposition}[9.6.3]
\label{I.9.6.3-fr}
Soit $X$ un préschéma localement noethérien. Toute
$\mathscr{O}_X$-Algèbre quasi-cohérente de type fini $\mathscr{B}$ est
alors un faisceau cohérent d'anneaux
(\hyperref[0.5.3.7-fr]{0, 5.3.7}).
\end{proposition}

\begin{proof}
On peut de nouveau se limiter au cas où $X$ est un schéma affine d'anneau
noethérien $A$, et où $\mathscr{B}=\widetilde{B}$, $B$ étant une
$A$-algèbre de type fini~; $B$ est donc un anneau noethérien. Cela étant, il
faut prouver que le noyau $\mathscr{N}$ d'un
$\mathscr{B}$-homomorphisme $\mathscr{B}^m\to\mathscr{B}$ est un
$\mathscr{B}$-
\oldpage[I]{180}
Module de type fini~; or, il est isomorphe (en tant que
$\mathscr{B}$-Module) à $\widetilde{N}$, où $N$ est le noyau de
l'homomorphisme correspondant de $B$-modules $B^m\to B$
(\hyperref[I.1.3.13-fr]{1.3.13}). Comme $B$ est noethérien, le sous-module
$N$ de $B^m$ est un $B$-module de type fini, donc il existe un homomorphisme
$B^p\to B^m$ d'image $N$~; la suite $B^p\to B^m\to B$ étant exacte, il
en est de même de la suite
$\mathscr{B}^p\to\mathscr{B}^m\to\mathscr{B}$ correspondante
(\hyperref[I.1.3.5-fr]{1.3.5}) et comme $\mathscr{N}$ est l'image de
$\mathscr{B}^p\to\mathscr{B}^m$
(\hyperref[I.1.3.9-fr]{1.3.9, (i)}), la proposition est démontrée.
\end{proof}

\begin{corollary}[9.6.4]
\label{I.9.6.4-fr}
Sous les hypothèses de (\hyperref[I.9.6.3-fr]{9.6.3}), pour qu'un
$\mathscr{B}$-Module $\mathscr{F}$ soit cohérent, il faut et il suffit
qu'il soit un $\mathscr{O}_X$-Module quasi-cohérent et un
$\mathscr{B}$-Module de type fini. S'il en est ainsi, et si $\mathscr{G}$
est un sous-$\mathscr{B}$-Module ou un $\mathscr{B}$-Module quotient de
$\mathscr{F}$, pour que $\mathscr{G}$ soit un $\mathscr{B}$-Module
cohérent, il faut et il suffit que $\mathscr{G}$ soit un
$\mathscr{O}_X$-Module quasi-cohérent.
\end{corollary}

\begin{proof}
Compte tenu de (\hyperref[I.9.6.1-fr]{9.6.1}), les conditions sur
$\mathscr{F}$ sont évidemment nécessaires~; montrons qu'elles sont
suffisantes. On peut se limiter au cas où $X$ est affine d'anneau
noethérien $A$, $\mathscr{B}=\widetilde{B}$, où $B$ est une $A$-algèbre de
type fini, $\mathscr{F}=\widetilde{M}$, où $M$ est un $B$-module, et où il
existe un $\mathscr{B}$-homomorphisme surjectif
$\mathscr{B}^m\to\mathscr{F}\to 0$. On a alors une suite exacte
correspondante $B^m\to M\to 0$, donc $M$ est un $B$-module de type fini~;
en outre, le noyau $P$ de l'homomorphisme $B^m\to M$ est alors un
$B$-module de type fini, puisque $B$ est noethérien. On en conclut
(\hyperref[I.1.3.13-fr]{1.3.13}) que $\mathscr{F}$ est le conoyau d'un
$\mathscr{B}$-homomorphisme $\mathscr{B}^n\to\mathscr{B}^m$, et est donc
cohérent puisque $\mathscr{B}$ est un faisceau cohérent d'anneaux
(\hyperref[0.5.3.4-fr]{0, 5.3.4}). Le même raisonnement montre qu'un
sous-$\mathscr{B}$-Module (resp. un $\mathscr{B}$-Module quotient)
quasi-cohérent de $\mathscr{F}$ est de type fini, d'où la seconde partie du
corollaire.
\end{proof}

\begin{proposition}[9.6.5]
\label{I.9.6.5-fr}
Soit $X$ un schéma quasi-compact ou un préschéma dont l'espace sous-jacent
est noethérien. Pour toute $\mathscr{O}_X$-Algèbre quasi-cohérente de type
fini $\mathscr{B}$, il existe un sous-$\mathscr{O}_X$-Module quasi-cohérent
de type fini $\mathscr{E}$ de $\mathscr{B}$ tel que $\mathscr{E}$ engendre
(\hyperref[0.4.1.4-fr]{0, 4.1.4}) la $\mathscr{O}_X$-Algèbre
$\mathscr{B}$.
\end{proposition}

\begin{proof}
En effet, il existe par hypothèse un recouvrement fini $(U_i)$ de $X$ formé
d'ouverts affines tels que $\Gamma(U_i,\mathscr{B})=B_i$ soit une algèbre
de type fini sur $\Gamma(U_i,\mathscr{O}_X)=A_i$~; soit $E_i$ un
sous-$A_i$-module de type fini de $B_i$ engendrant la $A_i$-algèbre $B_i$~;
en vertu de (\hyperref[I.9.4.7-fr]{9.4.7}), il existe un
sous-$\mathscr{O}_X$-Module $\mathscr{E}_i$ de $\mathscr{B}$,
quasi-cohérent et de type fini, tel que
$\mathscr{E}_i|U_i=\widetilde{E_i}$. Il est clair que la somme
$\mathscr{E}$ des $\mathscr{E}_i$ répond à la question.
\end{proof}

\begin{proposition}[9.6.6]
\label{I.9.6.6-fr}
Soit $X$ un préschéma dont l'espace sous-jacent est localement noethérien,
ou un schéma quasi-compact. Alors toute $\mathscr{O}_X$-Algèbre
quasi-cohérente $\mathscr{B}$ est limite inductive de ses
sous-$\mathscr{O}_X$-Algèbres quasi-cohérentes de type fini.
\end{proposition}

\begin{proof}
En effet, il résulte de (\hyperref[I.9.4.9-fr]{9.4.9}) que $\mathscr{B}$
est limite inductive (en tant que $\mathscr{O}_X$-Module) de ses
sous-$\mathscr{O}_X$-Modules quasi-cohérents de type fini~; ces derniers
engendrent des sous-$\mathscr{O}_X$-Algèbres quasi-cohérentes de type fini
de $\mathscr{B}$ (\hyperref[I.1.3.14-fr]{1.3.14}), dont $\mathscr{B}$ est
\emph{a fortiori} limite inductive.
\end{proof}

% FIN SOURCE ega1/ega1-9-fr.tex

%
% DEBUT SOURCE ega1/ega1-10-fr.tex
\section{Schémas formels}
\label{section:I.10-fr}

\subsection{Schémas formels affines.}
\label{subsection:I.10.1-fr}

\begin{env}[10.1.1]
\label{I.10.1.1-fr}
Soit $A$ un anneau topologique \emph{admissible}
(\hyperref[0.7.1.2-fr]{0, 7.1.2})~; pour tout idéal de définition
$\mathfrak{J}$ de $A$, $\operatorname{Spec}(A/\mathfrak{J})$ s'identifie
au sous-espace fermé $V(\mathfrak{J})$ de $\operatorname{Spec}(A)$
(\hyperref[I.1.1.11-fr]{1.1.11}), ensemble des idéaux premiers
\emph{ouverts} de $A$~; cet espace topologique ne dépend pas de
\oldpage[I]{181}
l'idéal de définition $\mathfrak{J}$ considéré~; notons-le
$\mathfrak{X}$. Soit $(\mathfrak{J}_\lambda)$ un système fondamental de
voisinages de 0 dans $A$, formé d'idéaux de définition, et pour tout
$\lambda$, soit $\mathscr{O}_\lambda$ le faisceau structural de
$\operatorname{Spec}(A/\mathfrak{J}_\lambda)$~; ce faisceau est induit sur
$\mathfrak{X}$ par
$\widetilde{A}/\widetilde{\mathfrak{J}}_\lambda$ (lequel est nul hors de
$\mathfrak{X}$).

Pour $\mathfrak{J}_\mu\subset\mathfrak{J}_\lambda$, l'homomorphisme
canonique $A/\mathfrak{J}_\mu\to A/\mathfrak{J}_\lambda$ définit donc un
homomorphisme
$u_{\lambda\mu}:\mathscr{O}_\mu\to\mathscr{O}_\lambda$ de faisceaux
d'anneaux (\hyperref[I.1.6.1-fr]{1.6.1}), et
$(\mathscr{O}_\lambda)$ est un \emph{système projectif de faisceaux
d'anneaux} pour ces homomorphismes. Comme la topologie de $\mathfrak{X}$
admet une base formée d'ouverts quasi-compacts, on peut associer à tout
$\mathscr{O}_\lambda$ un \emph{faisceau d'anneaux topologiques
pseudo-discrets} (\hyperref[0.3.8.1-fr]{0, 3.8.1}) qui a
$\mathscr{O}_\lambda$ comme faisceau d'anneaux (sans topologie) sous-jacent,
et que nous noterons encore $\mathscr{O}_\lambda$~; et les
$\mathscr{O}_\lambda$ forment encore un \emph{système projectif de faisceaux
d'anneaux topologiques} (\hyperref[0.3.8.2-fr]{0, 3.8.2}). Nous désignerons
par $\mathscr{O}_{\mathfrak{X}}$ le \emph{faisceau d'anneaux topologiques}
sur $\mathfrak{X}$, limite projective du système
$(\mathscr{O}_\lambda)$~; pour tout ouvert \emph{quasi-compact} $U$ de
$\mathfrak{X}$, $\Gamma(U,\mathscr{O}_{\mathfrak{X}})$ est donc l'anneau
topologique limite projective du système d'anneaux \emph{discrets}
$\Gamma(U,\mathscr{O}_\lambda)$
(\hyperref[0.3.2.6-fr]{0, 3.2.6}).
\end{env}

\begin{definition}[10.1.2]
\label{I.10.1.2-fr}
Étant donné un anneau topologique admissible $A$, on appelle spectre formel
de $A$ et on note $\operatorname{Spf}(A)$ le sous-espace fermé
$\mathfrak{X}$ de $\operatorname{Spec}(A)$ formé des idéaux premiers
ouverts de $A$. On dit qu'un espace topologiquement annelé est un schéma
formel affine s'il est isomorphe à un spectre formel
$\operatorname{Spf}(A)=\mathfrak{X}$ muni du faisceau d'anneaux
topologiques $\mathscr{O}_{\mathfrak{X}}$ limite projective des faisceaux
d'anneaux pseudo-discrets
$(\widetilde{A}/\widetilde{\mathfrak{J}}_\lambda)|\mathfrak{X}$, où
$\mathfrak{J}_\lambda$ parcourt l'ensemble filtrant des idéaux de
définition de $A$.
\end{definition}

Quand nous parlerons d'un \emph{spectre formel}
$\mathfrak{X}=\operatorname{Spf}(A)$ comme d'un schéma formel affine, il
s'agira toujours de l'espace topologiquement annelé
$(\mathfrak{X},\mathscr{O}_{\mathfrak{X}})$ où
$\mathscr{O}_{\mathfrak{X}}$ est défini comme ci-dessus.

On notera que tout \emph{schéma affine} $X=\operatorname{Spec}(A)$ peut être
considéré comme un schéma formel affine d'une seule manière, en considérant
$A$ comme un anneau topologique discret~: les anneaux topologiques
$\Gamma(U,\mathscr{O}_X)$ sont alors discrets lorsque $U$ est
quasi-compact (mais non en général lorsque $U$ est un ouvert quelconque de
$X$).

\begin{proposition}[10.1.3]
\label{I.10.1.3-fr}
Si $\mathfrak{X}=\operatorname{Spf}(A)$, où $A$ est un anneau admissible,
$\Gamma(\mathfrak{X},\mathscr{O}_{\mathfrak{X}})$ est topologiquement
isomorphe à $A$.
\end{proposition}

\begin{proof}
En effet, comme $\mathfrak{X}$ est fermé dans $\operatorname{Spec}(A)$, il
est quasi-compact, et par suite
$\Gamma(\mathfrak{X},\mathscr{O}_{\mathfrak{X}})$ est topologiquement
isomorphe à la limite projective des anneaux discrets
$\Gamma(\mathfrak{X},\mathscr{O}_\lambda)$~; mais
$\Gamma(\mathfrak{X},\mathscr{O}_\lambda)$ est isomorphe à
$A/\mathfrak{J}_\lambda$ (\hyperref[I.1.3.7-fr]{1.3.7})~; comme $A$ est
séparé et complet, il est topologiquement isomorphe à
$\varprojlim A/\mathfrak{J}_\lambda$
(\hyperref[0.7.2.1-fr]{0, 7.2.1}), d'où la proposition.
\end{proof}

\begin{proposition}[10.1.4]
\label{I.10.1.4-fr}
Soient $A$ un anneau admissible, $\mathfrak{X}=\operatorname{Spf}(A)$, et
pour tout $f\in A$, soit
$\mathfrak{D}(f)=D(f)\cap\mathfrak{X}$~; l'espace topologiquement annelé
$(\mathfrak{D}(f),\mathscr{O}_{\mathfrak{X}}|\mathfrak{D}(f))$ est
isomorphe au spectre formel affine
$\operatorname{Spf}(A_{\{f\}})$
(\hyperref[0.7.6.15-fr]{0, 7.6.15}).
\end{proposition}

\begin{proof}
Pour tout idéal de définition $\mathfrak{J}$ de $A$, l'anneau discret
$S_f^{-1}A/S_f^{-1}\mathfrak{J}$ s'identifie canoniquement à
$A_{\{f\}}/\mathfrak{J}_{\{f\}}$
(\hyperref[0.7.6.9-fr]{0, 7.6.9}), donc
(\hyperref[I.1.2.5-fr]{1.2.5} et \hyperref[I.1.2.6-fr]{1.2.6}), l'espace
topologique $\operatorname{Spf}(A_{\{f\}})$ s'identifie canoniquement à
$\mathfrak{D}(f)$. En outre, pour tout ouvert quasi-compact $U$ de
$\mathfrak{X}$ contenu dans $\mathfrak{D}(f)$,
$\Gamma(U,\mathscr{O}_\lambda)$ s'identifie au module des sections du
faisceau structural de
$\operatorname{Spec}(S_f^{-1}A/S_f^{-1}\mathfrak{J}_\lambda)$ au-dessus de
$U$ (\hyperref[I.1.3.6-fr]{1.3.6}), donc, si on pose
$\mathfrak{Y}=\operatorname{Spf}(A_{\{f\}})$,
$\Gamma(U,\mathscr{O}_{\mathfrak{X}})$ s'identifie au module des sections
$\Gamma(U,\mathscr{O}_{\mathfrak{Y}})$, ce qui démontre la proposition.
\end{proof}

\oldpage[I]{182}
\begin{env}[10.1.5]
\label{I.10.1.5-fr}
En tant que faisceau d'anneaux \emph{sans topologie}, le faisceau
structural $\mathscr{O}_{\mathfrak{X}}$ de
$\operatorname{Spf}(A)$ admet pour tout $x\in\mathfrak{X}$ une fibre qui,
en vertu de (\hyperref[I.10.1.4-fr]{10.1.4}), s'identifie à la limite
inductive
$\varinjlim_{f\notin\mathfrak{j}_x}A_{\{f\}}$. Par suite
(\hyperref[0.7.6.17-fr]{0, 7.6.17} et
\hyperref[0.7.6.18-fr]{7.6.18})~:
\end{env}

\begin{proposition}[10.1.6]
\label{I.10.1.6-fr}
Pour tout $x\in\mathfrak{X}=\operatorname{Spf}(A)$, la fibre
$\mathscr{O}_x$ est un anneau local dont le corps résiduel est isomorphe à
$k(x)=A_x/\mathfrak{j}_xA_x$. Si en outre $A$ est adique et noethérien,
$\mathscr{O}_x$ est un anneau noethérien.
\end{proposition}

Comme $k(x)$ n'est pas réduit à 0, on conclut de ce résultat que le
\emph{support} du faisceau d'anneaux $\mathscr{O}_{\mathfrak{X}}$ est
\emph{égal à $\mathfrak{X}$}.

\subsection{Morphismes de schémas formels affines.}
\label{subsection:I.10.2-fr}

\begin{env}[10.2.1]
\label{I.10.2.1-fr}
Soient $A$, $B$ deux anneaux admissibles, et soit
$\varphi:B\to A$ un homomorphisme \emph{continu}. L'application continue
$ {}^a\varphi:\operatorname{Spec}(A)\to\operatorname{Spec}(B)$
(\hyperref[I.1.2.1-fr]{1.2.1}) applique alors
$\mathfrak{X}=\operatorname{Spf}(A)$ dans
$\mathfrak{Y}=\operatorname{Spf}(B)$, car l'image réciproque par
$\varphi$ d'un idéal premier ouvert de $A$ est un idéal premier ouvert de
$B$. D'autre part, pour tout $g\in B$, $\varphi$ définit un homomorphisme
continu
\[
  \Gamma(\mathfrak{D}(g),\mathscr{O}_{\mathfrak{Y}})
  \longrightarrow
  \Gamma(\mathfrak{D}(\varphi(g)),\mathscr{O}_{\mathfrak{X}})
\]
en vertu de (\hyperref[I.10.1.4-fr]{10.1.4}),
(\hyperref[I.10.1.3-fr]{10.1.3}) et
(\hyperref[0.7.6.7-fr]{0, 7.6.7})~; comme ces homomorphismes satisfont aux
conditions de compatibilité pour les restrictions correspondant au
passage de $g$ à un multiple de $g$, et que
$\mathfrak{D}(\varphi(g))={}^a\varphi^{-1}(\mathfrak{D}(g))$, ils
définissent un homomorphisme \emph{continu} de faisceaux d'anneaux
topologiques
$\mathscr{O}_{\mathfrak{Y}}\to
{}^a\varphi_*(\mathscr{O}_{\mathfrak{X}})$
(\hyperref[0.3.2.5-fr]{0, 3.2.5}), que nous noterons encore
$\widetilde{\varphi}$~; on a ainsi obtenu un morphisme
$\Phi=({}^a\varphi,\widetilde{\varphi})$ d'espaces topologiquement annelés
$\mathfrak{X}\to\mathfrak{Y}$. On notera qu'en tant qu'homomorphisme de
faisceaux d'anneaux sans topologie, $\widetilde{\varphi}$ définit un
homomorphisme
$\widetilde{\varphi}_x^\sharp:
\mathscr{O}_{{}^a\varphi(x)}\to\mathscr{O}_x$ sur les fibres, pour tout
$x\in\mathfrak{X}$.
\end{env}

\begin{proposition}[10.2.2]
\label{I.10.2.2-fr}
Soient $A$, $B$ deux anneaux topologiques admissibles, et soient
$\mathfrak{X}=\operatorname{Spf}(A)$,
$\mathfrak{Y}=\operatorname{Spf}(B)$. Pour qu'un morphisme
$u=(\psi,\theta):\mathfrak{X}\to\mathfrak{Y}$ d'espaces topologiquement
annelés soit de la forme $({}^a\varphi,\widetilde{\varphi})$, où
$\varphi$ est un homomorphisme continu d'anneaux $B\to A$, il faut et il
suffit que pour tout $x\in\mathfrak{X}$,
$\theta_x^\sharp$ soit un homomorphisme local
$\mathscr{O}_{\psi(x)}\to\mathscr{O}_x$.
\end{proposition}

\begin{proof}
La condition est nécessaire~: en effet, soit
$\mathfrak{p}=\mathfrak{j}_x\in\operatorname{Spf}(A)$, et soit
$\mathfrak{q}=\varphi^{-1}(\mathfrak{j}_x)$~; si
$g\notin\mathfrak{q}$, on a donc $\varphi(g)\notin\mathfrak{p}$, et il est
immédiat que l'homomorphisme
$B_{\{g\}}\to A_{\{\varphi(g)\}}$ déduit de $\varphi$
(\hyperref[0.7.6.7-fr]{0, 7.6.7}) transforme
$\mathfrak{q}_{\{g\}}$ en une partie de
$\mathfrak{p}_{\{\varphi(g)\}}$~; passant à la limite inductive, on voit
donc (compte tenu de (\hyperref[I.10.1.5-fr]{10.1.5}) et de
(\hyperref[0.7.6.17-fr]{0, 7.6.17})) que
$\widetilde{\varphi}_x^\sharp$ est un homomorphisme local.

Inversement, soit $(\psi,\theta)$ un morphisme vérifiant la condition de
l'énoncé~; en vertu de (\hyperref[I.10.1.3-fr]{10.1.3}), $\theta$ définit
un homomorphisme continu d'anneaux
\[
  \varphi=\Gamma(\theta):
  B=\Gamma(\mathfrak{Y},\mathscr{O}_{\mathfrak{Y}})
  \longrightarrow
  \Gamma(\mathfrak{X},\mathscr{O}_{\mathfrak{X}})=A.
\]
En vertu de l'hypothèse sur $\theta$, pour que la section $\varphi(g)$ de
$\mathscr{O}_{\mathfrak{X}}$ au-dessus de $\mathfrak{X}$ ait un germe
inversible au point $x$, il faut et il suffit que $g$ ait un germe
inversible au point $\psi(x)$. Mais en vertu de
(\hyperref[0.7.6.17-fr]{0, 7.6.17}), les sections de
$\mathscr{O}_{\mathfrak{X}}$ (resp. $\mathscr{O}_{\mathfrak{Y}}$)
au-dessus de $\mathfrak{X}$ (resp. $\mathfrak{Y}$) dont le germe n'est pas
inversible au point $x$ (resp. $\psi(x)$) sont exactement les éléments de
$\mathfrak{j}_x$
\oldpage[I]{183}
(resp. $\mathfrak{j}_{\psi(x)}$)~; la remarque précédente montre donc que
${}^a\varphi=\psi$. Enfin, pour tout $g\in B$, le diagramme
\[
  \xymatrix{
    B=\Gamma(\mathfrak{Y},\mathscr{O}_{\mathfrak{Y}})
      \ar[r]^\varphi\ar[d] &
    \Gamma(\mathfrak{X},\mathscr{O}_{\mathfrak{X}})=A\ar[d]\\
    B_{\{g\}}=\Gamma(\mathfrak{D}(g),\mathscr{O}_{\mathfrak{Y}})
      \ar[r]_{\Gamma(\theta_{\mathfrak{D}(g)})} &
    \Gamma(\mathfrak{D}(\varphi(g)),\mathscr{O}_{\mathfrak{X}})
      =A_{\{\varphi(g)\}}
  }
\]
est commutatif~; d'après la propriété universelle des anneaux complets de
fractions (\hyperref[0.7.6.6-fr]{0, 7.6.6}),
$\theta_{\mathfrak{D}(g)}$ est égale à
$\widetilde{\varphi}_{\mathfrak{D}(g)}$ pour tout $g\in B$, donc
(\hyperref[0.3.2.5-fr]{0, 3.2.5}) on a
$\theta=\widetilde{\varphi}$.
\end{proof}

Nous dirons qu'un morphisme $(\psi,\theta)$ d'espaces topologiquement
annelés vérifiant la condition de (\hyperref[I.10.2.2-fr]{10.2.2}) est un
\emph{morphisme de schémas formels affines}. On peut dire que les foncteurs
$\operatorname{Spf}(A)$ en $A$ et
$\Gamma(\mathfrak{X},\mathscr{O}_{\mathfrak{X}})$ en $\mathfrak{X}$
définissent une \emph{équivalence} de la catégorie des anneaux admissibles
et de la catégorie duale de la catégorie des schémas formels affines
(T, I, 1.2).

\begin{env}[10.2.3]
\label{I.10.2.3-fr}
Comme cas particulier de (\hyperref[I.10.2.2-fr]{10.2.2}), notons que, pour
$f\in A$, l'injection canonique du schéma formel affine induit par
$\mathfrak{X}$ sur $\mathfrak{D}(f)$ correspond à l'homomorphisme continu
canonique $A\to A_{\{f\}}$. Sous les hypothèses de
(\hyperref[I.10.2.2-fr]{10.2.2}), soient $h$ un élément de $B$, $g$ un
élément de $A$, multiple de $\varphi(h)$~; on a alors
$\psi(\mathfrak{D}(g))\subset\mathfrak{D}(h)$~; la restriction de $u$ à
$\mathfrak{D}(g)$, considérée comme morphisme de $\mathfrak{D}(g)$ dans
$\mathfrak{D}(h)$, est l'unique morphisme $v$ rendant commutatif le
diagramme
\[
  \xymatrix{
    \mathfrak{D}(g)\ar[r]^v\ar[d] &
    \mathfrak{D}(h)\ar[d]\\
    \mathfrak{X}\ar[r]_u &
    \mathfrak{Y}
  }
\]
Ce morphisme correspond à l'unique homomorphisme continu
$\varphi':B_{\{h\}}\to A_{\{g\}}$
(\hyperref[0.7.6.7-fr]{0, 7.6.7}) rendant commutatif le diagramme
\[
  \xymatrix{
    A\ar[d] &
    B\ar[l]_\varphi\ar[d]\\
    A_{\{g\}} &
    B_{\{h\}}\ar[l]_{\varphi'}
  }
\]
\end{env}

\subsection{Idéaux de définition d'un schéma formel affine.}
\label{subsection:I.10.3-fr}

\begin{env}[10.3.1]
\label{I.10.3.1-fr}
Soient $A$ un anneau admissible, $\mathfrak{J}$ un idéal ouvert de $A$,
$\mathfrak{X}$ le schéma formel affine $\operatorname{Spf}(A)$. Soit
$(\mathfrak{J}_\lambda)$ l'ensemble des idéaux de définition de $A$
contenus dans $\mathfrak{J}$~; alors
$\widetilde{\mathfrak{J}}/\widetilde{\mathfrak{J}}_\lambda$ est un
faisceau d'idéaux de
$\widetilde{A}/\widetilde{\mathfrak{J}}_\lambda$. Désignons par
$\mathfrak{J}^\Delta$ la limite projective des faisceaux induits sur
$\mathfrak{X}$ par
$\widetilde{\mathfrak{J}}/\widetilde{\mathfrak{J}}_\lambda$, qui
s'identifie à un \emph{faisceau d'idéaux} de
$\mathscr{O}_{\mathfrak{X}}$ (\hyperref[0.3.2.6-fr]{0, 3.2.6}). Pour tout
$f\in A$, $\Gamma(\mathfrak{D}(f),\mathfrak{J}^\Delta)$ est la limite
projective de
$S_f^{-1}\mathfrak{J}/S_f^{-1}\mathfrak{J}_\lambda$, autrement dit
s'identifie à l'idéal ouvert $\mathfrak{J}_{\{f\}}$ de l'anneau
$A_{\{f\}}$ (\hyperref[0.7.6.9-fr]{0, 7.6.9}), et en particulier
$\Gamma(\mathfrak{X},\mathfrak{J}^\Delta)=\mathfrak{J}$~; on en conclut
(les $\mathfrak{D}(f)$ formant une base de la topologie de $\mathfrak{X}$)
que l'on a
\[
  \mathfrak{J}^\Delta|\mathfrak{D}(f)
  =(\mathfrak{J}_{\{f\}})^\Delta.
  \tag{10.3.1.1}\label{I.10.3.1.1-fr}
\]
\end{env}

\oldpage[I]{184}
\begin{env}[10.3.2]
\label{I.10.3.2-fr}
Avec les notations de (\hyperref[I.10.3.1-fr]{10.3.1}), pour tout $f\in A$,
l'application canonique de
$A_{\{f\}}=\Gamma(\mathfrak{D}(f),\mathscr{O}_{\mathfrak{X}})$ dans
$\Gamma\bigl(\mathfrak{D}(f),
(\widetilde{A}/\widetilde{\mathfrak{J}})|\mathfrak{X}\bigr)
=S_f^{-1}A/S_f^{-1}\mathfrak{J}$
est \emph{surjective} et a pour noyau
$\Gamma(\mathfrak{D}(f),\mathfrak{J}^\Delta)=\mathfrak{J}_{\{f\}}$
(\hyperref[0.7.6.9-fr]{0, 7.6.9})~; ces applications définissent donc un
homomorphisme \emph{surjectif} continu, dit \emph{canonique}, du faisceau
d'anneaux topologiques $\mathscr{O}_{\mathfrak{X}}$ sur le faisceau
d'anneaux discrets
$(\widetilde{A}/\widetilde{\mathfrak{J}})|\mathfrak{X}$, dont le noyau est
$\mathfrak{J}^\Delta$~; cet homomorphisme n'est autre d'ailleurs que
$\widetilde{\varphi}$ (\hyperref[I.10.2.1-fr]{10.2.1}), où $\varphi$ est
l'homomorphisme continu $A\to A/\mathfrak{J}$~; le morphisme
$({}^a\varphi,\widetilde{\varphi}):
\operatorname{Spec}(A/\mathfrak{J})\to\mathfrak{X}$ de schémas formels
affines (où ${}^a\varphi$ est d'ailleurs l'homéomorphisme identique de
$\mathfrak{X}$ sur lui-même) est encore dit \emph{canonique}. On a donc,
d'après ce qui précède, un \emph{isomorphisme canonique}
\[
  \mathscr{O}_{\mathfrak{X}}/\mathfrak{J}^\Delta
  \xrightarrow{\sim}
  (\widetilde{A}/\widetilde{\mathfrak{J}})|\mathfrak{X}.
  \tag{10.3.2.1}\label{I.10.3.2.1-fr}
\]

Il est clair (en vertu de
$\Gamma(\mathfrak{X},\mathfrak{J}^\Delta)=\mathfrak{J}$) que
l'application $\mathfrak{J}\to\mathfrak{J}^\Delta$ est
\emph{strictement croissante}~; d'après ce qui précède, pour
$\mathfrak{J}\subset\mathfrak{J}'$, le faisceau
${\mathfrak{J}'}^\Delta/\mathfrak{J}^\Delta$ est canoniquement isomorphe à
$\widetilde{\mathfrak{J}'}/\widetilde{\mathfrak{J}}
=\widetilde{\mathfrak{J}'/\mathfrak{J}}$.
\end{env}

\begin{env}[10.3.3]
\label{I.10.3.3-fr}
Les hypothèses et notations étant toujours celles de
(\hyperref[I.10.3.1-fr]{10.3.1}), nous dirons qu'un faisceau d'idéaux
$\mathscr{J}$ de $\mathscr{O}_{\mathfrak{X}}$ est un
\emph{faisceau d'idéaux de définition de $\mathfrak{X}$} (ou un
\emph{Idéal de définition de $\mathfrak{X}$}) si, pour tout
$x\in\mathfrak{X}$, il existe un voisinage ouvert de $x$ de la forme
$\mathfrak{D}(f)$, où $f\in A$, tel que
$\mathscr{J}|\mathfrak{D}(f)$ soit de la forme $\mathfrak{H}^\Delta$, où
$\mathfrak{H}$ est un idéal de définition de $A_{\{f\}}$.
\end{env}

\begin{proposition}[10.3.4]
\label{I.10.3.4-fr}
Pour tout $f\in A$, tout Idéal de définition de $\mathfrak{X}$ induit un
Idéal de définition de $\mathfrak{D}(f)$.
\end{proposition}

\begin{proof}
Cela résulte de (\hyperref[I.10.3.1.1-fr]{10.3.1.1}).
\end{proof}

\begin{proposition}[10.3.5]
\label{I.10.3.5-fr}
Si $A$ est un anneau admissible, tout Idéal de définition de
$\mathfrak{X}=\operatorname{Spf}(A)$ est de la forme
$\mathfrak{J}^\Delta$, où $\mathfrak{J}$ est un idéal de définition de
$A$, uniquement déterminé.
\end{proposition}

\begin{proof}
En effet, soit $\mathscr{J}$ un Idéal de définition de $\mathfrak{X}$~;
par hypothèse, et puisque $\mathfrak{X}$ est quasi-compact, il y a un nombre
fini d'éléments $f_i\in A$ tels que les $\mathfrak{D}(f_i)$ recouvrent
$\mathfrak{X}$ et que
$\mathscr{J}|\mathfrak{D}(f_i)=\mathfrak{H}_i^\Delta$, où
$\mathfrak{H}_i$ est un idéal de définition de $A_{\{f_i\}}$. Pour tout
$i$, il existe donc un idéal ouvert $\mathfrak{K}_i$ de $A$ tel que
$(\mathfrak{K}_i)_{\{f_i\}}=\mathfrak{H}_i$
(\hyperref[0.7.6.9-fr]{0, 7.6.9})~; soit $\mathfrak{K}$ un idéal de
définition de $A$ contenu dans tous les $\mathfrak{K}_i$. L'image canonique
de $\mathscr{J}/\mathfrak{K}^\Delta$ dans le faisceau structural
$\widetilde{A/\mathfrak{K}}$ de
$\operatorname{Spec}(A/\mathfrak{K})$
(\hyperref[I.10.3.2-fr]{10.3.2}) est donc telle que sa restriction à
$\mathfrak{D}(f_i)$ soit égale à celle de
$\widetilde{\mathfrak{K}_i/\mathfrak{K}}$~; on en conclut que cette image
canonique est un faisceau quasi-cohérent sur
$\operatorname{Spec}(A/\mathfrak{K})$, donc de la forme
$\widetilde{\mathfrak{J}/\mathfrak{K}}$, où $\mathfrak{J}$ est un idéal de
$A$ contenant $\mathfrak{K}$ (\hyperref[I.1.4.1-fr]{1.4.1}) d'où
$\mathscr{J}=\mathfrak{J}^\Delta$
(\hyperref[I.10.3.2-fr]{10.3.2})~; en outre, comme pour tout $i$ il existe
un entier $n_i$ tel que
$\mathfrak{H}_i^{n_i}\subset\mathfrak{K}_{\{f_i\}}$, on aura, en désignant
par $n$ le plus grand des $n_i$,
$(\mathscr{J}/\mathfrak{K}^\Delta)^n=0$, et par suite
(\hyperref[I.10.3.2-fr]{10.3.2})
$(\widetilde{\mathfrak{J}/\mathfrak{K}})^n=0$, d'où finalement
$(\mathfrak{J}/\mathfrak{K})^n=0$
(\hyperref[I.1.3.13-fr]{1.3.13}), ce qui prouve que $\mathfrak{J}$ est un
idéal de définition de $A$ (\hyperref[0.7.1.4-fr]{0, 7.1.4}).
\end{proof}

\begin{proposition}[10.3.6]
\label{I.10.3.6-fr}
Soient $A$ un anneau adique, $\mathfrak{J}$ un idéal de définition de $A$
tel que $\mathfrak{J}/\mathfrak{J}^2$ soit un
$A/\mathfrak{J}$-module de type fini. Pour tout entier $n>0$, on a alors
$(\mathfrak{J}^\Delta)^n=(\mathfrak{J}^n)^\Delta$.
\end{proposition}

\begin{proof}
En effet, pour tout $f\in A$, on a (puisque $\mathfrak{J}^n$ est un idéal
ouvert)
\[
  \bigl(\Gamma(\mathfrak{D}(f),\mathfrak{J}^\Delta)\bigr)^n
  =(\mathfrak{J}_{\{f\}})^n
  =(\mathfrak{J}^n)_{\{f\}}
  =\Gamma\bigl(\mathfrak{D}(f^n),(\mathfrak{J}^n)^\Delta\bigr)
\]
\oldpage[I]{185}
en vertu de (\hyperref[I.10.3.1.1-fr]{10.3.1.1}) et de
(\hyperref[0.7.6.12-fr]{0, 7.6.12}). Comme
$(\mathfrak{J}^\Delta)^n$ est associé au préfaisceau
$U\mapsto(\Gamma(U,\mathfrak{J}^\Delta))^n$
(\hyperref[0.4.1.6-fr]{0, 4.1.6}), le corollaire en résulte, puisque les
$\mathfrak{D}(f)$ forment une base de la topologie de $\mathfrak{X}$.
\end{proof}

\begin{env}[10.3.7]
\label{I.10.3.7-fr}
On dit qu'une famille $(\mathscr{J}_\lambda)$ d'Idéaux de définition de
$\mathfrak{X}$ est un \emph{système fondamental d'Idéaux de définition}
si tout Idéal de définition de $\mathfrak{X}$ contient un des
$\mathscr{J}_\lambda$~; comme
$\mathscr{J}_\lambda=\mathfrak{J}_\lambda^\Delta$, il revient au même de
dire que les $\mathfrak{J}_\lambda$ forment un \emph{système fondamental
de voisinages de 0} dans $A$. Soit $(f_\alpha)$ une famille d'éléments de
$A$ tels que les $\mathfrak{D}(f_\alpha)$ recouvrent $\mathfrak{X}$. Si
$(\mathscr{J}_\lambda)$ est une famille filtrante décroissante d'idéaux de
$\mathscr{O}_{\mathfrak{X}}$ telle que pour tout $\alpha$, la famille
$(\mathscr{J}_\lambda|\mathfrak{D}(f_\alpha))$ soit un système fondamental
d'Idéaux de définition de $\mathfrak{D}(f_\alpha)$, alors
$(\mathscr{J}_\lambda)$ est un système fondamental d'Idéaux de définition
de $\mathfrak{X}$. En effet, pour tout Idéal de définition $\mathscr{J}$
de $\mathfrak{X}$, il y a un recouvrement fini de $\mathfrak{X}$ par des
$\mathfrak{D}(f_i)$ tel que, pour tout $i$,
$\mathscr{J}_{\lambda_i}|\mathfrak{D}(f_i)$ soit un Idéal de définition
de $\mathfrak{D}(f_i)$ contenu dans
$\mathscr{J}|\mathfrak{D}(f_i)$. Si $\mu$ est un indice tel que
$\mathscr{J}_\mu\subset\mathscr{J}_{\lambda_i}$ pour tout $i$, il résulte
de (\hyperref[I.10.3.3-fr]{10.3.3}) que $\mathscr{J}_\mu$ est un Idéal de
définition de $\mathfrak{X}$, évidemment contenu dans $\mathscr{J}$, d'où
notre assertion.
\end{env}

\subsection{Préschémas formels et morphismes de préschémas formels.}
\label{subsection:I.10.4-fr}

\begin{env}[10.4.1]
\label{I.10.4.1-fr}
Étant donné un espace topologiquement annelé $\mathfrak{X}$, on dit qu'un
ouvert $U\subset\mathfrak{X}$ est un \emph{ouvert formel affine} (resp. un
\emph{ouvert formel affine adique}, resp. un \emph{ouvert formel affine
noethérien}) si l'espace topologiquement annelé induit par $\mathfrak{X}$
sur $U$ est un schéma formel affine (resp. un tel schéma dont l'anneau est
adique, resp. adique et noethérien).
\end{env}

\begin{definition}[10.4.2]
\label{I.10.4.2-fr}
On appelle préschéma formel un espace topologiquement annelé $\mathfrak{X}$
dont tout point admet un voisinage ouvert formel affine. On dit que le
préschéma formel $\mathfrak{X}$ est adique (resp. localement noethérien) si
tout point de $\mathfrak{X}$ admet un voisinage ouvert formel affine adique
(resp. noethérien). On dit que $\mathfrak{X}$ est noethérien s'il est
localement noethérien et si son espace sous-jacent est quasi-compact (donc
noethérien).
\end{definition}

\begin{proposition}[10.4.3]
\label{I.10.4.3-fr}
Si $\mathfrak{X}$ est un préschéma formel (resp. localement noethérien),
les ensembles ouverts formels affines (resp. affines noethériens) forment
une base de la topologie de $\mathfrak{X}$.
\end{proposition}

\begin{proof}
Cela résulte de (\hyperref[I.10.4.2-fr]{10.4.2}) et
(\hyperref[I.10.1.4-fr]{10.1.4}) en tenant compte de ce que si $A$ est un
anneau adique noethérien, il en est de même de $A_{\{f\}}$ pour tout
$f\in A$ (\hyperref[0.7.6.11-fr]{0, 7.6.11}).
\end{proof}

\begin{corollary}[10.4.4]
\label{I.10.4.4-fr}
Si $\mathfrak{X}$ est un préschéma formel (resp. un préschéma formel
localement noethérien, resp. noethérien), l'espace topologiquement annelé
induit sur tout ouvert de $\mathfrak{X}$ est encore un préschéma formel
(resp. un préschéma formel localement noethérien, resp. noethérien).
\end{corollary}

\begin{definition}[10.4.5]
\label{I.10.4.5-fr}
Étant donnés deux préschémas formels $\mathfrak{X}$, $\mathfrak{Y}$, on
appelle morphisme (de préschémas formels) de $\mathfrak{X}$ dans
$\mathfrak{Y}$ tout morphisme $(\psi,\theta)$ d'espaces topologiquement
annelés tel que, pour tout $x\in\mathfrak{X}$, $\theta_x^\sharp$ soit un
homomorphisme local
$\mathscr{O}_{\psi(x)}\to\mathscr{O}_x$.
\end{definition}

Il est immédiat que le composé de deux morphismes de préschémas formels est
encore un tel morphisme~; les préschémas formels forment donc une
\emph{catégorie}, et on notera
$\operatorname{Hom}(\mathfrak{X},\mathfrak{Y})$ l'ensemble des morphismes
d'un préschéma formel $\mathfrak{X}$ dans un préschéma formel
$\mathfrak{Y}$.

\oldpage[I]{186}
Si $U$ est une partie ouverte de $\mathfrak{X}$, l'injection canonique
dans $\mathfrak{X}$ du préschéma formel induit par $\mathfrak{X}$ sur $U$
est un morphisme de préschémas formels (et même un
\emph{monomorphisme} d'espaces topologiquement annelés
(\hyperref[0.4.1.1-fr]{0, 4.1.1})).

\begin{proposition}[10.4.6]
\label{I.10.4.6-fr}
Soient $\mathfrak{X}$ un préschéma formel,
$\mathfrak{S}=\operatorname{Spf}(A)$ un schéma formel affine. Il existe
une correspondance biunivoque canonique entre les morphismes du préschéma
formel $\mathfrak{X}$ dans le préschéma formel $\mathfrak{S}$ et les
homomorphismes continus de l'anneau $A$ dans l'anneau topologique
$\Gamma(\mathfrak{X},\mathscr{O}_{\mathfrak{X}})$.
\end{proposition}

La démonstration est la même que celle de
(\hyperref[I.2.2.4-fr]{2.2.4}), en remplaçant «~homomorphisme~» par
«~homomorphisme continu~», «~ouvert affine~» par «~ouvert formel
affine~», et en utilisant (\hyperref[I.10.2.2-fr]{10.2.2}) au lieu de
(\hyperref[I.1.7.3-fr]{1.7.3})~; nous en laissons les détails au lecteur.

\begin{env}[10.4.7]
\label{I.10.4.7-fr}
Étant donné un préschéma formel $\mathfrak{S}$, on dit que la donnée d'un
préschéma formel $\mathfrak{X}$ et d'un morphisme
$\varphi:\mathfrak{X}\to\mathfrak{S}$ définit un préschéma formel
$\mathfrak{X}$ \emph{au-dessus de $\mathfrak{S}$} ou un
\emph{$\mathfrak{S}$-préschéma formel}, $\varphi$ étant appelé le
\emph{morphisme structural} du $\mathfrak{S}$-préschéma
$\mathfrak{X}$. Si $\mathfrak{S}=\operatorname{Spf}(A)$, où $A$ est un
anneau admissible, on dit aussi que le $\mathfrak{S}$-préschéma formel
$\mathfrak{X}$ est un \emph{$A$-préschéma formel} ou un préschéma formel
\emph{au-dessus de $A$}. Un préschéma formel quelconque peut toujours
être considéré comme un préschéma formel au-dessus de $\mathbf{Z}$ (muni
de la topologie discrète).

Si $\mathfrak{X}$, $\mathfrak{Y}$ sont deux $\mathfrak{S}$-préschémas
formels, on dit qu'un morphisme
$u:\mathfrak{X}\to\mathfrak{Y}$ est un
\emph{$\mathfrak{S}$-morphisme} si le diagramme
\[
  \xymatrix{
    \mathfrak{X}\ar[rr]^u\ar[rd] & &
    \mathfrak{Y}\ar[ld]\\
    & \mathfrak{S}
  }
\]
(où les flèches obliques sont les morphismes structuraux) est
commutatif. Avec cette définition, les $\mathfrak{S}$-préschémas formels
(pour $\mathfrak{S}$ fixé) forment une \emph{catégorie}. On désigne par
$\operatorname{Hom}_{\mathfrak{S}}(\mathfrak{X},\mathfrak{Y})$
l'ensemble des $\mathfrak{S}$-morphismes du $\mathfrak{S}$-préschéma
formel $\mathfrak{X}$ dans le $\mathfrak{S}$-préschéma formel
$\mathfrak{Y}$. Lorsque $\mathfrak{S}=\operatorname{Spf}(A)$, on dit
aussi \emph{$A$-morphisme} au lieu de
\emph{$\mathfrak{S}$-morphisme}.
\end{env}

\begin{env}[10.4.8]
\label{I.10.4.8-fr}
Comme tout schéma affine peut être considéré comme un schéma formel affine
(\hyperref[I.10.1.2-fr]{10.1.2}), tout préschéma (usuel) peut être
considéré comme un préschéma formel. Il résulte en outre de
(\hyperref[I.10.4.5-fr]{10.4.5}) que pour les préschémas \emph{usuels},
les morphismes (resp. $S$-morphismes) de préschémas \emph{formels}
coïncident avec les morphismes (resp. $S$-morphismes) définis au
\S~\hyperref[section:I.2-fr]{2}.
\end{env}

\subsection{Idéaux de définition des préschémas formels.}
\label{subsection:I.10.5-fr}

\begin{env}[10.5.1]
\label{I.10.5.1-fr}
Soit $\mathfrak{X}$ un préschéma formel~; on dit qu'un
$\mathscr{O}_{\mathfrak{X}}$-Idéal $\mathscr{J}$ est un
\emph{faisceau d'idéaux de définition} (ou un
\emph{Idéal de définition}) de $\mathfrak{X}$ si tout
$x\in\mathfrak{X}$ possède un voisinage ouvert formel affine $U$ tel que
$\mathscr{J}|U$ soit un Idéal de définition du schéma formel affine
induit par $\mathfrak{X}$ sur $U$
(\hyperref[I.10.3.3-fr]{10.3.3})~; en vertu de
(\hyperref[I.10.3.1.1-fr]{10.3.1.1}) et
(\hyperref[I.10.4.3-fr]{10.4.3}), pour tout ouvert
$V\subset\mathfrak{X}$, $\mathscr{J}|V$ est alors un Idéal de définition
du préschéma formel induit par $\mathfrak{X}$ sur $V$.

On dit qu'une famille $(\mathscr{J}_\lambda)$ d'Idéaux de définition de
$\mathfrak{X}$ est un \emph{système fondamental}
\oldpage[I]{187}
\emph{d'Idéaux de définition} s'il existe un recouvrement
$(U_\alpha)$ de $\mathfrak{X}$ par des ouverts formels affines tel que,
pour tout $\alpha$, la famille des
$\mathscr{J}_\lambda|U_\alpha$ soit un système fondamental d'Idéaux de
définition (\hyperref[I.10.3.6-fr]{10.3.6}) du schéma formel affine
induit par $\mathfrak{X}$ sur $U_\alpha$. Il résulte de la remarque finale
de (\hyperref[I.10.3.7-fr]{10.3.7}) que lorsque $\mathfrak{X}$ est un
schéma formel affine, cette définition coïncide avec la définition donnée
dans (\hyperref[I.10.3.7-fr]{10.3.7}). Pour tout ouvert $V$ de
$\mathfrak{X}$, les restrictions $\mathscr{J}_\lambda|V$ forment alors un
système fondamental d'Idéaux de définition du préschéma formel induit sur
$V$, en vertu de (\hyperref[I.10.3.1.1-fr]{10.3.1.1}). Si
$\mathfrak{X}$ est un préschéma formel \emph{localement noethérien}, et
$\mathscr{J}$ un Idéal de définition de $\mathfrak{X}$, il résulte de
(\hyperref[I.10.3.6-fr]{10.3.6}) que les puissances $\mathscr{J}^n$
forment un système fondamental d'Idéaux de définition de $\mathfrak{X}$.
\end{env}

\begin{env}[10.5.2]
\label{I.10.5.2-fr}
Soient $\mathfrak{X}$ un préschéma formel, $\mathscr{J}$ un Idéal de
définition de $\mathfrak{X}$. Alors l'espace annelé
$(\mathfrak{X},\mathscr{O}_{\mathfrak{X}}/\mathscr{J})$ est un
\emph{préschéma} (usuel), qui est affine (resp. localement noethérien,
resp. noethérien) lorsque $\mathfrak{X}$ est un schéma formel affine
(resp. un préschéma formel localement noethérien, resp. noethérien)~; on
est en effet aussitôt ramené au cas affine, et alors la proposition a déjà
été démontrée dans (\hyperref[I.10.3.2-fr]{10.3.2}). En outre, si
$\theta:\mathscr{O}_{\mathfrak{X}}\to
\mathscr{O}_{\mathfrak{X}}/\mathscr{J}$ est l'homomorphisme canonique,
$u=(1_{\mathfrak{X}},\theta)$ est un \emph{morphisme} (dit
\emph{canonique}) de préschémas formels
$(\mathfrak{X},\mathscr{O}_{\mathfrak{X}}/\mathscr{J})\to
(\mathfrak{X},\mathscr{O}_{\mathfrak{X}})$, car ici encore, cela a été vu
dans le cas affine (\hyperref[I.10.3.2-fr]{10.3.2}), auquel on se ramène
aussitôt.
\end{env}

\begin{proposition}[10.5.3]
\label{I.10.5.3-fr}
Soient $\mathfrak{X}$ un préschéma formel, $(\mathscr{J}_\lambda)$ un
système fondamental d'Idéaux de définition de $\mathfrak{X}$. Alors le
faisceau d'anneaux topologiques $\mathscr{O}_{\mathfrak{X}}$ est limite
projective des faisceaux d'anneaux pseudo-discrets
(\hyperref[0.3.8.1-fr]{0, 3.8.1})
$\mathscr{O}_{\mathfrak{X}}/\mathscr{J}_\lambda$.
\end{proposition}

\begin{proof}
Comme la topologie de $\mathfrak{X}$ admet une base d'ouverts formels
affines quasi-compacts (\hyperref[I.10.4.3-fr]{10.4.3}), on est ramené au
cas affine, où la proposition est conséquence de
(\hyperref[I.10.3.5-fr]{10.3.5}),
(\hyperref[I.10.3.2-fr]{10.3.2}) et de la définition
(\hyperref[I.10.1.1-fr]{10.1.1}).
\end{proof}

Il n'est pas certain que tout préschéma formel admette des Idéaux de
définition. Toutefois~:

\begin{proposition}[10.5.4]
\label{I.10.5.4-fr}
Soit $\mathfrak{X}$ un préschéma formel localement noethérien. Il existe
un plus grand Idéal de définition $\mathscr{T}$ de $\mathfrak{X}$~;
c'est le seul Idéal de définition $\mathscr{J}$ tel que le préschéma
$(\mathfrak{X},\mathscr{O}_{\mathfrak{X}}/\mathscr{J})$ soit réduit. Si
$\mathscr{J}$ est un Idéal de définition de $\mathfrak{X}$,
$\mathscr{T}$ est l'image réciproque par
$\mathscr{O}_{\mathfrak{X}}\to
\mathscr{O}_{\mathfrak{X}}/\mathscr{J}$ du Nilradical de
$\mathscr{O}_{\mathfrak{X}}/\mathscr{J}$.
\end{proposition}

\begin{proof}
Supposons d'abord que $\mathfrak{X}=\operatorname{Spf}(A)$, où $A$ est
un anneau adique noethérien. L'existence et les propriétés de
$\mathscr{T}$ résultent immédiatement de
(\hyperref[I.10.3.4-fr]{10.3.4}) et
(\hyperref[I.5.1.1-fr]{5.1.1}), compte tenu de l'existence et des
propriétés du plus grand idéal de définition de $A$
(\hyperref[0.7.1.6-fr]{0, 7.1.6} et
\hyperref[0.7.1.7-fr]{7.1.7}).

Pour prouver l'existence et les propriétés de $\mathscr{T}$ dans le cas
général, il suffit de montrer que si $U\supset V$ sont deux ouverts
formels affines noethériens de $\mathfrak{X}$, le plus grand Idéal de
définition $\mathscr{T}_U$ de $U$ induit le plus grand Idéal de définition
$\mathscr{T}_V$ de $V$~; mais comme
$(V,(\mathscr{O}_{\mathfrak{X}}|V)/(\mathscr{T}_U|V))$ est réduit, cela
résulte de ce qui précède.
\end{proof}

On désigne par $\mathfrak{X}_{\mathrm{red}}$ le préschéma (usuel) réduit
$(\mathfrak{X},\mathscr{O}_{\mathfrak{X}}/\mathscr{T})$.

\begin{corollary}[10.5.5]
\label{I.10.5.5-fr}
Soient $\mathfrak{X}$ un préschéma formel localement noethérien,
$\mathscr{T}$ le plus grand Idéal de définition de $\mathfrak{X}$~; pour
tout ouvert $V$ de $\mathfrak{X}$, $\mathscr{T}|V$ est le plus grand
Idéal de définition du préschéma formel induit par $\mathfrak{X}$ sur
$V$.
\end{corollary}

\oldpage[I]{188}

\begin{proposition}[10.5.6]
\label{I.10.5.6-fr}
Soient $\mathfrak{X}$, $\mathfrak{Y}$ deux préschémas formels,
$\mathscr{J}$ (resp. $\mathscr{K}$) un Idéal de définition de
$\mathfrak{X}$ (resp. $\mathfrak{Y}$),
$f:\mathfrak{X}\to\mathfrak{Y}$ un morphisme de préschémas formels.
\begin{enumerate}
  \item[\rm{(i)}] Si
    $f^*(\mathscr{K})\mathscr{O}_{\mathfrak{X}}\subset\mathscr{J}$,
    il existe un unique morphisme
    $f':(\mathfrak{X},\mathscr{O}_{\mathfrak{X}}/\mathscr{J})\to
    (\mathfrak{Y},\mathscr{O}_{\mathfrak{Y}}/\mathscr{K})$ de
    préschémas usuels rendant commutatif le diagramme
    \[
      \label{I.10.5.6.1-fr}
      \xymatrix{
        (\mathfrak{X},\mathscr{O}_{\mathfrak{X}})\ar[r]^f &
        (\mathfrak{Y},\mathscr{O}_{\mathfrak{Y}})\\
        (\mathfrak{X},\mathscr{O}_{\mathfrak{X}}/\mathscr{J})
          \ar[r]_{f'}\ar[u] &
        (\mathfrak{Y},\mathscr{O}_{\mathfrak{Y}}/\mathscr{K})\ar[u]
      }
      \tag{10.5.6.1}
    \]
    où les flèches verticales sont les morphismes canoniques.
  \item[\rm{(ii)}] Supposons que
    $\mathfrak{X}=\operatorname{Spf}(A)$,
    $\mathfrak{Y}=\operatorname{Spf}(B)$ soient des schémas formels
    affines, $\mathscr{J}=\mathfrak{J}^{\Delta}$,
    $\mathscr{K}=\mathfrak{K}^{\Delta}$, où $\mathfrak{J}$ (resp.
    $\mathfrak{K}$) est un idéal de définition de $A$ (resp. $B$), et
    $f=({}^a\varphi,\widetilde{\varphi})$, où
    $\varphi:B\to A$ est un homomorphisme continu~; pour que
    $f^*(\mathscr{K})\mathscr{O}_{\mathfrak{X}}\subset\mathscr{J}$,
    il faut et il suffit que
    $\varphi(\mathfrak{K})\subset\mathfrak{J}$, et $f'$ est alors le
    morphisme $({}^a\varphi',\widetilde{\varphi'})$, où
    $\varphi':B/\mathfrak{K}\to A/\mathfrak{J}$ est l'homomorphisme
    déduit de $\varphi$ par passage aux quotients.
\end{enumerate}
\end{proposition}

\begin{proof}
\medskip\noindent
\begin{enumerate}
  \item[\rm{(i)}] Si $f=(\psi,\theta)$, l'hypothèse entraîne que
    l'image par
    $\theta^\sharp:\psi^*(\mathscr{O}_{\mathfrak{Y}})\to
    \mathscr{O}_{\mathfrak{X}}$ du faisceau d'idéaux
    $\psi^*(\mathscr{K})$ de $\psi^*(\mathscr{O}_{\mathfrak{Y}})$ est
    contenue dans $\mathscr{J}$
    (\hyperref[0.4.3.5-fr]{0, 4.3.5}). Par passage aux quotients, on
    déduit donc de $\theta^\sharp$ un homomorphisme de faisceau d'anneaux
    \[
      \omega:\psi^*(\mathscr{O}_{\mathfrak{Y}}/\mathscr{K})
      =\psi^*(\mathscr{O}_{\mathfrak{Y}})/\psi^*(\mathscr{K})
      \to\mathscr{O}_{\mathfrak{X}}/\mathscr{J}~;
    \]
    en outre, comme pour tout $x\in\mathfrak{X}$,
    $\theta_x^\sharp$ est un homomorphisme \emph{local}, il en est de
    même de $\omega_x$. Le morphisme d'espaces annelés
    $(\psi,\omega^b)$ est donc
    (\hyperref[I.2.2.1-fr]{2.2.1}) l'unique morphisme $f'$ d'espaces
    annelés répondant à la question.
  \item[\rm{(ii)}] La correspondance canonique fonctorielle entre
    morphismes de préschémas formels affines et homomorphismes continus
    d'anneaux (\hyperref[I.10.2.2-fr]{10.2.2}) montre que dans le cas
    considéré, la relation
    $f^*(\mathscr{K})\mathscr{O}_{\mathfrak{X}}\subset\mathscr{J}$
    entraîne que l'on a
    $f'=({}^a\varphi',\widetilde{\varphi'})$, où
    $\varphi':B/\mathfrak{K}\to A/\mathfrak{J}$ est l'unique
    homomorphisme rendant commutatif le diagramme
    \[
      \label{I.10.5.6.2-fr}
      \xymatrix{
        B\ar[r]^\varphi\ar[d] & A\ar[d]\\
        B/\mathfrak{K}\ar[r]_{\varphi'} & A/\mathfrak{J}
      }
      \tag{10.5.6.2}
    \]
    L'existence de $\varphi'$ implique donc que
    $\varphi(\mathfrak{K})\subset\mathfrak{J}$. Inversement, si cette
    condition est vérifiée, en désignant par $\varphi'$ l'unique
    homomorphisme rendant commutatif le diagramme
    (\hyperref[I.10.5.6.2-fr]{10.5.6.2}) et posant
    $f'=({}^a\varphi',\widetilde{\varphi'})$, il est clair que le
    diagramme (\hyperref[I.10.5.6.1-fr]{10.5.6.1}) est commutatif~; la
    considération des homomorphismes
    ${}^a\varphi^*(\mathscr{O}_{\mathfrak{Y}})\to
    \mathscr{O}_{\mathfrak{X}}$ et
    ${}^a{\varphi'}^*(\mathscr{O}_{\mathfrak{Y}}/\mathscr{K})\to
    \mathscr{O}_{\mathfrak{X}}/\mathscr{J}$ correspondant à $f$ et
    $f'$ respectivement, montre alors que cela entraîne la relation
    $f^*(\mathscr{K})\mathscr{O}_{\mathfrak{X}}\subset\mathscr{J}$.
\end{enumerate}
\end{proof}

Il est clair que la correspondance $f\mapsto f'$ définie ci-dessus est
\emph{fonctorielle}.

\subsection{Préschémas formels comme limites inductives de préschémas.}
\label{subsection:I.10.6-fr}

\begin{env}[10.6.1]
\label{I.10.6.1-fr}
Soient $\mathfrak{X}$ un préschéma formel,
$(\mathscr{J}_\lambda)$ un système fondamental d'Idéaux de définition de
$\mathfrak{X}$~; pour tout $\lambda$, soit $f_\lambda$ le morphisme
canonique
$(\mathfrak{X},\mathscr{O}_{\mathfrak{X}}/\mathscr{J}_\lambda)
\to\mathfrak{X}$ (\hyperref[I.10.5.2-fr]{10.5.2})~; pour
$\mathscr{J}_\mu\subset\mathscr{J}_\lambda$, l'homomorphisme canonique
$\mathscr{O}_{\mathfrak{X}}/\mathscr{J}_\mu\to
\mathscr{O}_{\mathfrak{X}}/\mathscr{J}_\lambda$ définit un morphisme

\oldpage[I]{189}

canonique
$f_{\mu\lambda}:(\mathfrak{X},
\mathscr{O}_{\mathfrak{X}}/\mathscr{J}_\lambda)\to
(\mathfrak{X},\mathscr{O}_{\mathfrak{X}}/\mathscr{J}_\mu)$ de préschémas
(usuels) tel que l'on ait $f_\lambda=f_\mu\circ f_{\mu\lambda}$. Les
préschémas
$X_\lambda=(\mathfrak{X},
\mathscr{O}_{\mathfrak{X}}/\mathscr{J}_\lambda)$ et les morphismes
$f_{\mu\lambda}$ constituent donc (en vertu de
\hyperref[I.10.4.8-fr]{10.4.8}) un \emph{système inductif} dans la
catégorie des préschémas formels.
\end{env}

\begin{proposition}[10.6.2]
\label{I.10.6.2-fr}
Avec les notations de (\hyperref[I.10.6.1-fr]{10.6.1}), le préschéma
formel $\mathfrak{X}$ et les morphismes $f_\lambda$ constituent une
limite inductive (T, I, 1.8) du système $(X_\lambda,f_{\mu\lambda})$
dans la catégorie des préschémas formels.
\end{proposition}

\begin{proof}
Soit $\mathfrak{Y}$ un préschéma formel, et pour chaque indice $\lambda$,
soit
\[
  g_\lambda=(\psi_\lambda,\theta_\lambda):X_\lambda\to\mathfrak{Y}
\]
un morphisme, tel que l'on ait
$g_\lambda=g_\mu\circ f_{\mu\lambda}$ pour
$\mathscr{J}_\mu\subset\mathscr{J}_\lambda$. Cette dernière condition
et la définition des $X_\lambda$ entraînent d'abord que tous les
$\psi_\lambda$ sont identiques à une même application continue
$\psi:\mathfrak{X}\to\mathfrak{Y}$ des espaces sous-jacents~; en outre,
les homomorphismes
$\theta_\lambda^\sharp:\psi^*(\mathscr{O}_{\mathfrak{Y}})\to
\mathscr{O}_{X_\lambda}=
\mathscr{O}_{\mathfrak{X}}/\mathscr{J}_\lambda$ forment un
\emph{système projectif} d'homomorphismes de faisceaux d'anneaux. Par
passage à la limite projective, on en déduit donc un homomorphisme
\[
  \omega:\psi^*(\mathscr{O}_{\mathfrak{Y}})\to
  \varprojlim\mathscr{O}_{\mathfrak{X}}/\mathscr{J}_\lambda
  =\mathscr{O}_{\mathfrak{X}},
\]
et il est clair que le morphisme $g=(\psi,\omega^b)$ d'\emph{espaces
annelés} est le \emph{seul} rendant commutatif les diagrammes
\[
  \label{I.10.6.2.1-fr}
  \xymatrix{
    X_\lambda\ar[rr]^{g_\lambda}\ar[rd]_{f_\lambda} &&
    \mathfrak{Y}\\
    & \mathfrak{X}\ar[ru]_g
  }
  \tag{10.6.2.1}
\]
Il reste donc à prouver que $g$ est un morphisme de \emph{préschémas
formels}~; la question étant locale sur $\mathfrak{X}$ et
$\mathfrak{Y}$, on peut donc supposer
$\mathfrak{X}=\operatorname{Spf}(A)$,
$\mathfrak{Y}=\operatorname{Spf}(B)$, $A$ et $B$ étant des anneaux
admissibles, avec
$\mathscr{J}_\lambda=\mathfrak{J}_\lambda^{\Delta}$, où
$(\mathfrak{J}_\lambda)$ est un système fondamental d'idéaux de
définition de $A$ (\hyperref[I.10.3.5-fr]{10.3.5})~; comme
$A=\varprojlim A/\mathfrak{J}_\lambda$, l'existence d'un morphisme de
schémas formels affines $g$ rendant commutatifs les diagrammes
(\hyperref[I.10.6.2.1-fr]{10.6.2.1}) résulte alors de la correspondance
biunivoque (\hyperref[I.10.2.2-fr]{10.2.2}) entre morphismes de schémas
formels affines et homomorphismes continus d'anneaux, et de la définition
de la limite projective. Mais l'unicité de $g$ en tant que morphisme
d'espaces annelés montre qu'il coïncide avec le morphisme noté de même au
début de la démonstration.
\end{proof}

La proposition suivante établit, sous certaines conditions
supplémentaires, l'existence de la limite inductive d'un système inductif
donné de préschémas (usuels) dans la catégorie des préschémas formels~:

\begin{proposition}[10.6.3]
\label{I.10.6.3-fr}
Soient $\mathfrak{X}$ un espace topologique,
$(\mathscr{O}_i,u_{ji})$ un système projectif de faisceaux d'anneaux sur
$\mathfrak{X}$, ayant $\mathbf{N}$ pour ensemble d'indices. Soit
$\mathscr{J}_i$ le noyau de
$u_{0i}:\mathscr{O}_i\to\mathscr{O}_0$. On suppose que~:
\begin{enumerate}
  \item[a)] L'espace annelé $(\mathfrak{X},\mathscr{O}_i)$ est un
    préschéma $X_i$.
  \item[b)] Pour tout $x\in X$ et tout $i$, il existe un voisinage
    ouvert $U_i$ de $x$ dans $\mathfrak{X}$ tel que la restriction
    $\mathscr{J}_i|U_i$ soit nilpotente.
  \item[c)] Les homomorphismes $u_{ji}$ sont surjectifs.
\end{enumerate}

\oldpage[I]{190}

Soit $\mathscr{O}_{\mathfrak{X}}$ le faisceau d'anneaux topologiques
limite projective des faisceaux d'anneaux pseudo-discrets
$\mathscr{O}_i$, et soit
$u_i:\mathscr{O}_{\mathfrak{X}}\to\mathscr{O}_i$ l'homomorphisme
canonique. Alors l'espace topologiquement annelé
$(\mathfrak{X},\mathscr{O}_{\mathfrak{X}})$ est un préschéma formel~;
les homomorphismes $u_i$ sont surjectifs~; leurs noyaux
$\mathscr{J}^{(i)}$ forment un système fondamental d'Idéaux de définition
de $\mathfrak{X}$, et $\mathscr{J}^{(0)}$ est la limite projective des
faisceaux d'idéaux $\mathscr{J}_i$.
\end{proposition}

\begin{proof}
Notons d'abord que sur chaque fibre, $u_{ji}$ est un homomorphisme
surjectif et \emph{a fortiori} un homomorphisme local~; donc
$v_{ij}=(1_{\mathfrak{X}},u_{ji})$ est un morphisme de préschémas
$X_j\to X_i$ ($i\geq j$) (\hyperref[I.2.2.1-fr]{2.2.1}). Supposons
d'abord que chaque $X_i$ soit un schéma affine d'anneau $A_i$. Il existe
un homomorphisme \emph{d'anneaux}
$\varphi_{ji}:A_i\to A_j$ tel que
$u_{ji}=\widetilde{\varphi_{ji}}$
(\hyperref[I.1.7.3-fr]{1.7.3})~; par suite
(\hyperref[I.1.6.3-fr]{1.6.3}), le faisceau $\mathscr{O}_j$ est un
$\mathscr{O}_i$-Module quasi-cohérent sur $X_i$ (pour la loi externe
définie par $u_{ji}$), associé à $A_j$ considéré comme $A_i$-module à
l'aide de $\varphi_{ji}$. Pour tout $f\in A_i$, soit
$f'=\varphi_{ji}(f)$~; par hypothèse, les ouverts $D(f)$ et $D(f')$ sont
identiques dans $\mathfrak{X}$, et l'homomorphisme de
$\Gamma(D(f),\mathscr{O}_i)=(A_i)_f$ dans
$\Gamma(D(f),\mathscr{O}_j)=(A_j)_{f'}$ correspondant à $u_{ji}$ n'est
autre que $(\varphi_{ji})_f$ (\hyperref[I.1.6.1-fr]{1.6.1}). Mais
lorsqu'on considère $A_j$ comme $A_i$-module, $(A_j)_{f'}$ est le
$(A_i)_f$-module $(A_j)_f$, donc on a aussi
$u_{ji}=\widetilde{\varphi_{ji}}$ lorsque $\varphi_{ji}$ est cette fois
considéré comme homomorphisme de $A_i$-modules. Alors, comme $u_{ji}$ est
surjectif, on en conclut que $\varphi_{ji}$ l'est aussi
(\hyperref[I.1.3.9-fr]{1.3.9}) et si $\mathfrak{J}_{ji}$ est le noyau de
$\varphi_{ji}$, le noyau de $u_{ji}$ est un $\mathscr{O}_i$-Module
quasi-cohérent égal à $\widetilde{\mathfrak{J}_{ji}}$. En particulier,
on a $\mathscr{J}_i=\widetilde{\mathfrak{J}_i}$, où $\mathfrak{J}_i$ est
le noyau de $\varphi_{0i}:A_i\to A_0$. L'hypothèse b) entraîne que
$\mathscr{J}_i$ est \emph{nilpotent}~: en effet, comme $\mathfrak{X}$
est quasi-compact, on peut recouvrir $\mathfrak{X}$ par un nombre fini
d'ouverts $U_k$ tels que $(\mathscr{J}_i|U_k)^{n_k}=0$ et en prenant pour
$n$ le plus grand des $n_k$, on a $\mathscr{J}_i^n=0$. On en conclut que
$\mathfrak{J}_i$ est nilpotent (\hyperref[I.1.3.13-fr]{1.3.13}). Alors
l'anneau $A=\varprojlim A_i$ est admissible
(\hyperref[0.7.2.2-fr]{0, 7.2.2}), l'homomorphisme canonique
$\varphi_i:A\to A_i$ est surjectif et son noyau
$\mathfrak{J}^{(i)}$ est égal à la limite projective des
$\mathfrak{J}_{ik}$ pour $k\geq i$~; les $\mathfrak{J}^{(i)}$ forment
un système fondamental de voisinages de 0 dans $A$. Les assertions de
(\hyperref[I.10.6.3-fr]{10.6.3}) résultent dans ce cas de
(\hyperref[I.10.1.1-fr]{10.1.1}) et
(\hyperref[I.10.3.2-fr]{10.3.2}),
$(\mathfrak{X},\mathscr{O}_{\mathfrak{X}})$ n'étant autre que
$\operatorname{Spf}(A)$.

Toujours dans ce même cas particulier, notons que si $f=(f_i)$ est un
élément de la limite projective $A=\varprojlim A_i$, tous les ouverts
$D(f_i)$ (ouvert affine dans $X_i$) s'identifient à l'ouvert
$\mathfrak{D}(f)$ de $\mathfrak{X}$, le préschéma induit par $X_i$ sur
$\mathfrak{D}(f)$ s'identifiant donc au schéma affine
$\operatorname{Spec}((A_i)_{f_i})$.

Dans le cas général, remarquons d'abord que pour tout ouvert quasi-compact
$U$ de $\mathfrak{X}$, chacun des $\mathscr{J}_i|U$ est nilpotent, comme
le montre le raisonnement fait ci-dessus. Nous allons voir que pour tout
$x\in\mathfrak{X}$, il y a un voisinage ouvert $U$ de $x$ dans
$\mathfrak{X}$ qui est un \emph{ouvert affine pour tous les $X_i$}. En
effet, prenons $U$ ouvert affine pour $X_0$, et observons que
$\mathscr{O}_{X_0}=\mathscr{O}_{X_i}/\mathscr{J}_i$. Comme
$\mathscr{J}_i|U$ est nilpotent, en vertu de ce qui précède, $U$ est
ouvert affine aussi pour chaque $X_i$ en vertu de
(\hyperref[I.5.1.9-fr]{5.1.9}). Cela étant, pour tout $U$ satisfaisant
aux conditions précédentes, l'étude du cas affine faite plus haut montre
que $(U,\mathscr{O}_{\mathfrak{X}}|U)$ est un préschéma formel dont les
$\mathscr{J}^{(i)}|U$ forment un système fondamental d'idéaux de
définition et $\mathscr{J}^{(0)}|U$ est limite projective des
$\mathscr{J}_i|U$~; d'où la conclusion.
\end{proof}

\begin{corollary}[10.6.4]
\label{I.10.6.4-fr}
Supposons que pour $i\geq j$, le noyau de $u_{ji}$ soit
$\mathscr{J}_i^{j+1}$ et que $\mathscr{J}_1/\mathscr{J}_1^2$

\oldpage[I]{191}

soit de type fini sur
$\mathscr{O}_0=\mathscr{O}_1/\mathscr{J}_1$. Alors $\mathfrak{X}$ est
un préschéma formel adique, et si $\mathscr{J}^{(n)}$ est le noyau de
$\mathscr{O}_{\mathfrak{X}}\to\mathscr{O}_n$, on a
$\mathscr{J}^{(n)}=\mathscr{J}^{n+1}$ et
$\mathscr{J}/\mathscr{J}^2$ est isomorphe à $\mathscr{J}_1$. Si en
outre $X_0$ est localement noethérien (resp. noethérien),
$\mathfrak{X}$ est localement noethérien (resp. noethérien).
\end{corollary}

\begin{proof}
Comme les espaces sous-jacents à $\mathfrak{X}$ et $X_0$ sont les mêmes,
la question est locale et l'on peut supposer tous les $X_i$ affines~;
compte tenu des relations
$\mathscr{J}_{ij}=\widetilde{\mathfrak{J}_{ji}}$ (avec les notations
de (\hyperref[I.10.6.3-fr]{10.6.3})), on est aussitôt ramené aux
assertions correspondantes de
(\hyperref[0.7.2.7-fr]{0, 7.2.7} et
\hyperref[0.7.2.8-fr]{7.2.8}), en notant que
$\mathfrak{J}_1/\mathfrak{J}_1^2$ est alors un $A_0$-module de type
fini (\hyperref[I.1.3.9-fr]{1.3.9}).
\end{proof}

En particulier, \emph{tout préschéma formel localement noethérien
$\mathfrak{X}$} est limite inductive d'une suite $(X_n)$ de préschémas
(usuels) localement noethériens vérifiant les conditions de
(\hyperref[I.10.6.3-fr]{10.6.3}) et
(\hyperref[I.10.6.4-fr]{10.6.4})~: il suffit de considérer un Idéal de
définition $\mathscr{J}$ de $\mathfrak{X}$
(\hyperref[I.10.5.4-fr]{10.5.4}) et de prendre
$X_n=(\mathfrak{X},
\mathscr{O}_{\mathfrak{X}}/\mathscr{J}^{n+1})$
((\hyperref[I.10.5.1-fr]{10.5.1}) et
(\hyperref[I.10.6.2-fr]{10.6.2})).

\begin{corollary}[10.6.5]
\label{I.10.6.5-fr}
Soit $A$ un anneau admissible. Pour que le schéma formel affine
$\mathfrak{X}=\operatorname{Spf}(A)$ soit noethérien, il faut et il
suffit que $A$ soit adique et noethérien.
\end{corollary}

\begin{proof}
La condition est évidemment suffisante. Inversement, supposons que
$\mathfrak{X}$ soit noethérien, et soient $\mathfrak{J}$ un idéal de
définition de $A$, $\mathscr{J}=\mathfrak{J}^\Delta$ l'Idéal de
définition correspondant de $\mathfrak{X}$. Les préschémas (usuels)
$X_n=(\mathfrak{X},
\mathscr{O}_{\mathfrak{X}}/\mathscr{J}^{n+1})$ sont alors affines et
noethériens, donc les anneaux
$A_n=A/\mathfrak{J}^{n+1}$ sont noethériens
(\hyperref[I.6.1.3-fr]{6.1.3}), d'où on conclut que
$\mathfrak{J}/\mathfrak{J}^2$ est un $A/\mathfrak{J}$-module de type
fini. Comme les $\mathscr{J}^n$ forment un système fondamental d'Idéaux
de définition de $\mathfrak{X}$
(\hyperref[I.10.5.1-fr]{10.5.1}), on a
$\mathscr{O}_{\mathfrak{X}}
=\varprojlim(\mathscr{O}_{\mathfrak{X}}/\mathscr{J}^n)$
(\hyperref[I.10.5.3-fr]{10.5.3})~; on en conclut
(\hyperref[I.10.1.3-fr]{10.1.3}) que $A$ est topologiquement isomorphe à
$\varprojlim A/\mathfrak{J}^n$, donc est adique et noethérien
(\hyperref[0.7.2.8-fr]{0, 7.2.8}).
\end{proof}

\begin{remark}[10.6.6]
\label{I.10.6.6-fr}
Avec les notations de (\hyperref[I.10.6.3-fr]{10.6.3}), soit
$\mathscr{F}_i$ un $\mathscr{O}_i$-Module, et supposons donné, pour
$i\geq j$, un $v_{ij}$-morphisme
$\theta_{ji}:\mathscr{F}_i\to\mathscr{F}_j$, de sorte que
$\theta_{kj}\circ\theta_{ji}=\theta_{ki}$ pour $k\leq j\leq i$.
Comme l'application continue sous-jacente à $v_{ij}$ est l'identité,
$\theta_{ji}$ est un homomorphisme de faisceaux de groupes abéliens sur
l'espace $\mathfrak{X}$~; en outre, si $\mathscr{F}$ est la limite
projective du système projectif $(\mathscr{F}_i)$ de faisceaux de groupes
abéliens, le fait que les $\theta_{ji}$ sont des $v_{ij}$-morphismes
permet de définir sur $\mathscr{F}$ une structure de
$\mathscr{O}_{\mathfrak{X}}$-Module par passage à la limite projective~;
muni de cette structure, nous dirons que $\mathscr{F}$ est la
\emph{limite projective} (pour les $\theta_{ji}$) du système de
$\mathscr{O}_i$-Modules $(\mathscr{F}_i)$. Dans le cas particulier où
$v_{ij}^*(\mathscr{F}_i)=\mathscr{F}_j$ et où $\theta_{ji}$ est
l'\emph{identité}, nous dirons pour abréger, que $\mathscr{F}$ est la
limite projective d'un système $(\mathscr{F}_i)$ tel que
$v_{ij}^*(\mathscr{F}_i)=\mathscr{F}_j$ pour $j\leq i$ (sans
mentionner les $\theta_{ji}$).
\end{remark}

\begin{env}[10.6.7]
\label{I.10.6.7-fr}
Soient $\mathfrak{X}$, $\mathfrak{Y}$ deux préschémas formels,
$\mathscr{J}$ (resp. $\mathscr{K}$) un Idéal de définition de
$\mathfrak{X}$ (resp. $\mathfrak{Y}$),
$f:\mathfrak{X}\to\mathfrak{Y}$ un morphisme tel que
$f^*(\mathscr{K})\mathscr{O}_{\mathfrak{X}}\subset\mathscr{J}$. On a
alors pour tout entier $n>0$,
$f^*(\mathscr{K}^n)\mathscr{O}_{\mathfrak{X}}
=(f^*(\mathscr{K})\mathscr{O}_{\mathfrak{X}})^n
\subset\mathscr{J}^n$~; on peut donc
(\hyperref[I.10.5.6-fr]{10.5.6}) déduire de $f$ un morphisme de
préschémas (usuels) $f_n:X_n\to Y_n$, en posant
$X_n=(\mathfrak{X},
\mathscr{O}_{\mathfrak{X}}/\mathscr{J}^{n+1})$,
$Y_n=(\mathfrak{Y},
\mathscr{O}_{\mathfrak{Y}}/\mathscr{K}^{n+1})$, et il résulte
aussitôt des définitions que les diagrammes
\[
  \label{I.10.6.7.1-fr}
  \xymatrix{
    X_m\ar[r]^{f_m}\ar[d] &
    Y_m\ar[d]\\
    X_n\ar[r]_{f_n} &
    Y_n
  }
  \tag{10.6.7.1}
\]

\oldpage[I]{192}

sont commutatifs pour $m\leq n$~; autrement dit, $(f_n)$ est un
\emph{système inductif} de morphismes.
\end{env}

\begin{env}[10.6.8]
\label{I.10.6.8-fr}
Inversement, soit $(X_n)$ (resp. $(Y_n)$) un système inductif de
préschémas (usuels) satisfaisant aux conditions b) et c) de
(\hyperref[I.10.6.3-fr]{10.6.3}), et soit $\mathfrak{X}$ (resp.
$\mathfrak{Y}$) sa limite inductive. Par définition des limites
inductives, toute suite $(f_n)$ de morphismes $X_n\to Y_n$ formant un
système inductif admet une \emph{limite inductive}
$f:\mathfrak{X}\to\mathfrak{Y}$, qui est l'unique morphisme de
préschémas formels rendant commutatifs les diagrammes
\[
  \xymatrix{
    X_n\ar[r]^{f_n}\ar[d] &
    Y_n\ar[d]\\
    \mathfrak{X}\ar[r]_f &
    \mathfrak{Y}
  }
\]
\end{env}

\begin{proposition}[10.6.9]
\label{I.10.6.9-fr}
Soient $\mathfrak{X}$, $\mathfrak{Y}$ deux préschémas formels
localement noethériens, $\mathscr{J}$ (resp. $\mathscr{K}$) un Idéal
de définition de $\mathfrak{X}$ (resp. $\mathfrak{Y}$)~;
l'application $f\mapsto(f_n)$ définie dans
(\hyperref[I.10.6.7-fr]{10.6.7}) est une bijection de l'ensemble des
morphismes $f:\mathfrak{X}\to\mathfrak{Y}$ tels que
$f^*(\mathscr{K})\mathscr{O}_{\mathfrak{X}}\subset\mathscr{J}$, sur
l'ensemble des suites $(f_n)$ de morphismes rendant commutatifs les
diagrammes (\hyperref[I.10.6.7.1-fr]{10.6.7.1}).
\end{proposition}

\begin{proof}
Si $f$ est la limite inductive d'une telle suite, il faut montrer que
$f^*(\mathscr{K})\mathscr{O}_{\mathfrak{X}}\subset\mathscr{J}$. La
question étant locale sur $\mathfrak{X}$ et $\mathfrak{Y}$, on peut se
borner au cas où $\mathfrak{X}=\operatorname{Spf}(A)$,
$\mathfrak{Y}=\operatorname{Spf}(B)$ sont affines, $A$ et $B$ étant
adiques noethériens,
$\mathscr{J}=\mathfrak{J}^\Delta$,
$\mathscr{K}=\mathfrak{K}^\Delta$, où $\mathfrak{J}$ (resp.
$\mathfrak{K}$) est un idéal de définition de $A$ (resp. $B$). On a
alors $X_n=\operatorname{Spec}(A_n)$,
$Y_n=\operatorname{Spec}(B_n)$, avec
$A_n=A/\mathfrak{J}^{n+1}$ et $B_n=B/\mathfrak{K}^{n+1}$, en vertu de
(\hyperref[I.10.3.6-fr]{10.3.6}) et
(\hyperref[I.10.3.2-fr]{10.3.2})~;
$f_n=({}^a\varphi_n,\widetilde{\varphi_n})$, où les homomorphismes
$\varphi_n:B_n\to A_n$ forment un système projectif, donc
$f=({}^a\varphi,\widetilde{\varphi})$, où
$\varphi=\varprojlim\varphi_n$. La commutativité du diagramme
(\hyperref[I.10.6.7.1-fr]{10.6.7.1}) pour $m=0$ donne alors la
condition
$\varphi_n(\mathfrak{K}/\mathfrak{K}^{n+1})
\subset\mathfrak{J}/\mathfrak{J}^{n+1}$ pour tout $n$, donc, en
passant à la limite projective, $\varphi(\mathfrak{K})\subset
\mathfrak{J}$, et cela entraîne
$f^*(\mathscr{K})\mathscr{O}_{\mathfrak{X}}\subset\mathscr{J}$
(\hyperref[I.10.5.6-fr]{10.5.6, (ii)}).
\end{proof}

\begin{corollary}[10.6.10]
\label{I.10.6.10-fr}
Soient $\mathfrak{X}$, $\mathfrak{Y}$ deux préschémas formels
localement noethériens, $\mathscr{T}$ le plus grand Idéal de définition
de $\mathfrak{X}$ (\hyperref[I.10.5.4-fr]{10.5.4}).
\begin{enumerate}
  \item[\rm{(i)}] Pour tout Idéal de définition $\mathscr{K}$ de
  $\mathfrak{Y}$ et tout morphisme
  $f:\mathfrak{X}\to\mathfrak{Y}$, on a
  $f^*(\mathscr{K})\mathscr{O}_{\mathfrak{X}}\subset\mathscr{T}$.
  \item[\rm{(ii)}] Il y a correspondance biunivoque canonique entre
  $\operatorname{Hom}(\mathfrak{X},\mathfrak{Y})$ et l'ensemble des
  suites $(f_n)$ de morphismes rendant commutatifs les diagrammes
  (\hyperref[I.10.6.7.1-fr]{10.6.7.1}), où
  $X_n=(\mathfrak{X},
  \mathscr{O}_{\mathfrak{X}}/\mathscr{T}^{n+1})$,
  $Y_n=(\mathfrak{Y},
  \mathscr{O}_{\mathfrak{Y}}/\mathscr{K}^{n+1})$.
\end{enumerate}
\end{corollary}

\begin{proof}
(ii) résulte aussitôt de (i) et de
(\hyperref[I.10.6.9-fr]{10.6.9}). Pour démontrer (i), on peut se borner
au cas où $\mathfrak{X}=\operatorname{Spf}(A)$,
$\mathfrak{Y}=\operatorname{Spf}(B)$, $A$ et $B$ étant noethériens,
$\mathscr{T}=\mathfrak{T}^\Delta$,
$\mathscr{K}=\mathfrak{K}^\Delta$, où $\mathfrak{T}$ est le plus grand
idéal de définition de $A$ et $\mathfrak{K}$ un idéal de définition de
$B$. Soit $f=({}^a\varphi,\widetilde{\varphi})$, où
$\varphi:B\to A$ est un homomorphisme continu~; comme les éléments de
$\mathfrak{K}$ sont topologiquement nilpotents
(\hyperref[0.7.1.4-fr]{0, 7.1.4, (ii)}), il en est de même de ceux de
$\varphi(\mathfrak{K})$, donc
$\varphi(\mathfrak{K})\subset\mathfrak{T}$ puisque $\mathfrak{T}$ est
l'ensemble des éléments topologiquement nilpotents de $A$
(\hyperref[0.7.1.6-fr]{0, 7.1.6})~; d'où la conclusion en vertu de
(\hyperref[I.10.5.6-fr]{10.5.6, (ii)}).
\end{proof}

\begin{corollary}[10.6.11]
\label{I.10.6.11-fr}
Soient $\mathfrak{S}$, $\mathfrak{X}$, $\mathfrak{Y}$ trois
préschémas formels localement noethériens,
$f:\mathfrak{X}\to\mathfrak{S}$,
$g:\mathfrak{Y}\to\mathfrak{S}$ des morphismes faisant de
$\mathfrak{X}$ et $\mathfrak{Y}$ des $\mathfrak{S}$-préschémas
formels. Soit $\mathscr{J}$ (resp. $\mathscr{K}$, $\mathscr{L}$) un
Idéal de définition de $\mathfrak{S}$ (resp. $\mathfrak{X}$,
$\mathfrak{Y}$), et supposons que
$f^*(\mathscr{J})\mathscr{O}_{\mathfrak{X}}\subset\mathscr{K}$,
$g^*(\mathscr{J})\mathscr{O}_{\mathfrak{Y}}=\mathscr{L}$~; posons
$S_n=(\mathfrak{S},
\mathscr{O}_{\mathfrak{S}}/\mathscr{J}^{n+1})$,
$X_n=(\mathfrak{X},
\mathscr{O}_{\mathfrak{X}}/\mathscr{K}^{n+1})$,
$Y_n=(\mathfrak{Y},
\mathscr{O}_{\mathfrak{Y}}/\mathscr{L}^{n+1})$. Il y a alors
correspondance

\oldpage[I]{193}

biunivoque canonique entre
$\operatorname{Hom}_{\mathfrak{S}}(\mathfrak{X},\mathfrak{Y})$ et
l'ensemble des suites $(u_n)$ de $S_n$-morphismes
$u_n:X_n\to Y_n$ rendant commutatifs les diagrammes
(\hyperref[I.10.6.7.1-fr]{10.6.7.1}).
\end{corollary}

\begin{proof}
Pour tout $\mathfrak{S}$-morphisme
$u:\mathfrak{X}\to\mathfrak{Y}$, on a par définition $f=g\circ u$, donc
\[
  u^*(\mathscr{L})\mathscr{O}_{\mathfrak{X}}
  =u^*(g^*(\mathscr{J})\mathscr{O}_{\mathfrak{Y}})
    \mathscr{O}_{\mathfrak{X}}
  =f^*(\mathscr{J})\mathscr{O}_{\mathfrak{X}}
  \subset\mathscr{K}
\]
le corollaire découle donc de
(\hyperref[I.10.6.9-fr]{10.6.9}).
\end{proof}

On notera que, pour $m\leq n$, la donnée d'un morphisme
$f_n:X_n\to Y_n$ détermine un morphisme et un seul
$f_m:X_m\to Y_m$ rendant commutatif le diagramme
(\hyperref[I.10.6.7.1-fr]{10.6.7.1}), comme on le voit aussitôt en se
ramenant au cas affine~; on a ainsi défini une application
$\varphi_{mn}:\operatorname{Hom}_{S_n}(X_n,Y_n)
\to\operatorname{Hom}_{S_m}(X_m,Y_m)$ et les
$\operatorname{Hom}_{S_n}(X_n,Y_n)$ forment pour les $\varphi_{mn}$ un
\emph{système projectif d'ensembles}~;
(\hyperref[I.10.6.11-fr]{10.6.11}) s'énonce encore en disant qu'il existe
une bijection canonique
\[
  \operatorname{Hom}_{\mathfrak{S}}(\mathfrak{X},\mathfrak{Y})
  \xrightarrow{\sim}
  \varprojlim_n\operatorname{Hom}_{S_n}(X_n,Y_n).
\]

\subsection{Produit de préschémas formels.}
\label{subsection:I.10.7-fr}

\begin{env}[10.7.1]
\label{I.10.7.1-fr}
Soit $\mathfrak{S}$ un préschéma formel~; les $\mathfrak{S}$-préschémas
formels formant une catégorie, on peut définir la notion de
\emph{produit} de $\mathfrak{S}$-préschémas formels.
\end{env}

\begin{proposition}[10.7.2]
\label{I.10.7.2-fr}
Soient $\mathfrak{X}=\operatorname{Spf}(B)$,
$\mathfrak{Y}=\operatorname{Spf}(C)$ deux schémas formels affines
au-dessus d'un schéma formel affine
$\mathfrak{S}=\operatorname{Spf}(A)$. Soient
$\mathfrak{Z}=\operatorname{Spf}(B\widehat{\otimes}_A C)$,
$p_1$, $p_2$ les $\mathfrak{S}$-morphismes correspondant
(\hyperref[I.10.2.2-fr]{10.2.2}) aux $A$-homomorphismes canoniques
(continus) $\rho$ et $\sigma$ de $B$ et $C$ dans
$B\widehat{\otimes}_A C$~; alors $(\mathfrak{Z},p_1,p_2)$ est un produit
des $\mathfrak{S}$-schémas formels affines $\mathfrak{X}$ et
$\mathfrak{Y}$.
\end{proposition}

\begin{proof}
En vertu de (\hyperref[I.10.4.6-fr]{10.4.6}), tout revient à vérifier que
si, à tout $A$-homomorphisme continu
$\varphi:B\widehat{\otimes}_A C\to D$, où $D$ est un anneau admissible qui
est une $A$-algèbre topologique, on associe le couple
$(\varphi\circ\rho,\varphi\circ\sigma)$, on définit une bijection
\[
  \operatorname{Hom}_A(B\widehat{\otimes}_A C,D)
  \xrightarrow{\sim}
  \operatorname{Hom}_A(B,D)\times\operatorname{Hom}_A(C,D)
\]
ce qui n'est autre que la propriété universelle du produit tensoriel
complété (\hyperref[0.7.7.6-fr]{0, 7.7.6}).
\end{proof}

\begin{proposition}[10.7.3]
\label{I.10.7.3-fr}
Étant donnés deux $\mathfrak{S}$-préschémas formels
$\mathfrak{X}$, $\mathfrak{Y}$, le produit
$\mathfrak{X}\times_{\mathfrak{S}}\mathfrak{Y}$ existe.
\end{proposition}

\begin{proof}
La démonstration est identique à celle de
(\hyperref[I.3.2.6-fr]{3.2.6}), en y remplaçant les schémas affines
(resp. les ouverts affines) par les schémas formels affines (resp. les
ouverts formels affines), et la prop.
(\hyperref[I.3.2.2-fr]{3.2.2}) par
(\hyperref[I.10.7.2-fr]{10.7.2}).
\end{proof}

Toutes les propriétés formelles du produit de préschémas
(\hyperref[I.3.2.7-fr]{3.2.7} et
\hyperref[I.3.2.8-fr]{3.2.8},
\hyperref[I.3.3.1-fr]{3.3.1} à
\hyperref[I.3.3.12-fr]{3.3.12}) sont valables sans aucune modification
pour le produit de préschémas formels.

\begin{env}[10.7.4]
\label{I.10.7.4-fr}
Soient $\mathfrak{S}$, $\mathfrak{X}$, $\mathfrak{Y}$ trois préschémas
formels et soient $f:\mathfrak{X}\to\mathfrak{S}$,
$g:\mathfrak{Y}\to\mathfrak{S}$ deux morphismes. Supposons qu'il existe
dans $\mathfrak{S}$, $\mathfrak{X}$, $\mathfrak{Y}$ respectivement,
trois systèmes fondamentaux d'Idéaux de définition
$(\mathscr{J}_\lambda)$, $(\mathscr{K}_\lambda)$,
$(\mathscr{L}_\lambda)$ respectivement, ayant même ensemble d'indices
$I$, tels que
$f^*(\mathscr{J}_\lambda)\mathscr{O}_{\mathfrak{X}}
\subset\mathscr{K}_\lambda$ et
$g^*(\mathscr{J}_\lambda)\mathscr{O}_{\mathfrak{Y}}
\subset\mathscr{L}_\lambda$ pour tout $\lambda$. Posons
$S_\lambda=(\mathfrak{S},
\mathscr{O}_{\mathfrak{S}}/\mathscr{J}_\lambda)$,
$X_\lambda=(\mathfrak{X},
\mathscr{O}_{\mathfrak{X}}/\mathscr{K}_\lambda)$,
$Y_\lambda=(\mathfrak{Y},
\mathscr{O}_{\mathfrak{Y}}/\mathscr{L}_\lambda)$~; pour
$\mathscr{J}_\mu\subset\mathscr{J}_\lambda$,
$\mathscr{K}_\mu\subset\mathscr{K}_\lambda$,
$\mathscr{L}_\mu\subset\mathscr{L}_\lambda$, notons que
$S_\lambda$ (resp. $X_\lambda$, $Y_\lambda$) est un sous-préschéma fermé
de $S_\mu$ (resp. $X_\mu$, $Y_\mu$) ayant même

\oldpage[I]{194}

espace sous-jacent (\hyperref[I.10.6.1-fr]{10.6.1}). Comme
$S_\lambda\to S_\mu$ est un monomorphisme de préschémas, on voit donc
d'abord que les produits
$X_\lambda\times_{S_\lambda}Y_\lambda$ et
$X_\lambda\times_{S_\mu}Y_\lambda$ sont identiques
(\hyperref[I.3.2.4-fr]{3.2.4}), puis que
$X_\lambda\times_{S_\mu}Y_\lambda$ s'identifie à un sous-préschéma fermé
de $X_\mu\times_{S_\mu}Y_\mu$ ayant même espace sous-jacent
(\hyperref[I.4.3.1-fr]{4.3.1}). Cela étant, le produit
$\mathfrak{X}\times_{\mathfrak{S}}\mathfrak{Y}$ est la
\emph{limite inductive} des préschémas usuels
$X_\lambda\times_{S_\lambda}Y_\lambda$~: en effet, on voit comme dans
(\hyperref[I.10.6.2-fr]{10.6.2}) qu'on peut se ramener au cas où
$\mathfrak{S}$, $\mathfrak{X}$ et $\mathfrak{Y}$ sont des schémas formels
affines. Compte tenu de (\hyperref[I.10.5.6-fr]{10.5.6, (ii)}) et de
l'hypothèse sur les systèmes fondamentaux d'Idéaux de définition de
$\mathfrak{S}$, $\mathfrak{X}$ et $\mathfrak{Y}$, on voit aussitôt que
notre assertion résulte de la définition du produit tensoriel complété de
deux algèbres (\hyperref[0.7.7.1-fr]{0, 7.7.1}).

En outre, soit $\mathfrak{Z}$ un $\mathfrak{S}$-préschéma formel,
$(\mathscr{M}_\lambda)$ un système fondamental d'Idéaux de définition de
$\mathfrak{Z}$ ayant $I$ comme ensemble d'indices,
$u:\mathfrak{Z}\to\mathfrak{X}$,
$v:\mathfrak{Z}\to\mathfrak{Y}$ deux $\mathfrak{S}$-morphismes tels que
$u^*(\mathscr{K}_\lambda)\mathscr{O}_{\mathfrak{Z}}
\subset\mathscr{M}_\lambda$ et
$v^*(\mathscr{L}_\lambda)\mathscr{O}_{\mathfrak{Z}}
\subset\mathscr{M}_\lambda$. Si l'on pose
$Z_\lambda=(\mathfrak{Z},
\mathscr{O}_{\mathfrak{Z}}/\mathscr{M}_\lambda)$, et si
$u_\lambda:Z_\lambda\to X_\lambda$ et
$v_\lambda:Z_\lambda\to Y_\lambda$ sont les $S_\lambda$-morphismes
correspondant à $u$ et $v$ (\hyperref[I.10.5.6-fr]{10.5.6}), on vérifie
aussitôt que $(u,v)_{\mathfrak{S}}$ est la limite inductive des
$S_\lambda$-morphismes $(u_\lambda,v_\lambda)_{S_\lambda}$.

Les considérations de ce numéro s'appliquent en particulier lorsque
$\mathfrak{S}$, $\mathfrak{X}$ et $\mathfrak{Y}$ sont localement
noethériens, en prenant comme systèmes fondamentaux d'Idéaux de définition
les systèmes formés des puissances d'un Idéal de définition
(\hyperref[I.10.5.1-fr]{10.5.1}). Mais on notera que
$\mathfrak{X}\times_{\mathfrak{S}}\mathfrak{Y}$ n'est pas nécessairement
localement noethérien (voir toutefois
(\hyperref[I.10.13.5-fr]{10.13.5})).
\end{env}

\subsection{Complété formel d'un préschéma le long d'une partie fermée.}
\label{subsection:I.10.8-fr}

\begin{env}[10.8.1]
\label{I.10.8.1-fr}
Soient $X$ un préschéma (usuel) localement noethérien, $X'$ une partie
fermée de l'espace sous-jacent à $X$~; désignons par $\Phi$ l'ensemble des
faisceaux cohérents d'idéaux $\mathscr{J}$ dans $\mathscr{O}_X$, tels que
le support de $\mathscr{O}_X/\mathscr{J}$ soit $X'$. L'ensemble $\Phi$
n'est pas vide (\hyperref[I.5.2.1-fr]{5.2.1},
\hyperref[I.4.1.4-fr]{4.1.4} et
\hyperref[I.6.1.1-fr]{6.1.1})~; nous l'ordonnerons par la relation
$\supset$.
\end{env}

\begin{lemma}[10.8.2]
\label{I.10.8.2-fr}
L'ensemble ordonné $\Phi$ est filtrant~; si $X$ est noethérien, pour tout
$\mathscr{J}_0\in\Phi$, l'ensemble des puissances
$\mathscr{J}_0^n$ ($n>0$) est cofinal à $\Phi$.
\end{lemma}

\begin{proof}
En effet, si $\mathscr{J}_1$ et $\mathscr{J}_2$ appartiennent à $\Phi$, et
si on pose $\mathscr{J}=\mathscr{J}_1\cap\mathscr{J}_2$,
$\mathscr{J}$ est cohérent puisque $\mathscr{O}_X$ est cohérent
(\hyperref[I.6.1.1-fr]{6.1.1} et
\hyperref[0.5.3.4-fr]{0, 5.3.4}), et l'on a
$\mathscr{J}_x=(\mathscr{J}_1)_x\cap(\mathscr{J}_2)_x$ pour tout
$x\in X$, donc $\mathscr{J}_x=\mathscr{O}_x$ pour $x\notin X'$ et
$\mathscr{J}_x\neq\mathscr{O}_x$ pour $x\in X'$, ce qui prouve que
$\mathscr{J}\in\Phi$. D'autre part, si $X$ est noethérien, et si
$\mathscr{J}_0$ et $\mathscr{J}$ appartiennent à $\Phi$, il existe un
entier $n>0$ tel que
$\mathscr{J}_0^n(\mathscr{O}_X/\mathscr{J})=0$
(\hyperref[I.9.3.4-fr]{9.3.4}), ce qui signifie que
$\mathscr{J}_0^n\subset\mathscr{J}$.
\end{proof}

\begin{env}[10.8.3]
\label{I.10.8.3-fr}
Soit maintenant $\mathscr{F}$ un $\mathscr{O}_X$-Module cohérent~; pour
tout $\mathscr{J}\in\Phi$,
$\mathscr{F}\otimes_{\mathscr{O}_X}(\mathscr{O}_X/\mathscr{J})$ est un
$\mathscr{O}_X$-Module cohérent
(\hyperref[I.9.1.1-fr]{9.1.1}), de support contenu dans $X'$, et que nous
identifierons le plus souvent à sa restriction à $X'$. Lorsque
$\mathscr{J}$ parcourt $\Phi$, ces faisceaux forment un
\emph{système projectif} de faisceaux de groupes abéliens.
\end{env}

\begin{definition}[10.8.4]
\label{I.10.8.4-fr}
Étant donnés une partie fermée $X'$ d'un préschéma localement noethérien
$X$ et un $\mathscr{O}_X$-Module cohérent $\mathscr{F}$, on appelle
\emph{complété de $\mathscr{F}$ le long de $X'$} et on désigne par
$\mathscr{F}_{/X'}$ ou par $\widehat{\mathscr{F}}$ (lorsqu'aucune confusion
n'est possible) la restriction à $X'$ du faisceau

\oldpage[I]{195}

$\varprojlim_{\Phi}
(\mathscr{F}\otimes_{\mathscr{O}_X}
(\mathscr{O}_X/\mathscr{J}))$~; on dit que ses sections au-dessus de $X'$ sont
les \emph{sections formelles de $\mathscr{F}$ le long de $X'$}.
\end{definition}

Il est immédiat que pour tout ouvert $U\subset X$, on a
$(\mathscr{F}|U)_{/(U\cap X')}=(\mathscr{F}_{/X'})|(U\cap X')$.

Par passage à la limite projective, il est clair que
$(\mathscr{O}_X)_{/X'}$ est un faisceau d'anneaux, et que
$\mathscr{F}_{/X'}$ peut être considéré comme un
$(\mathscr{O}_X)_{/X'}$-Module. En outre, comme il existe une base de la
topologie de $X'$ formée d'ouverts quasi-compacts, on peut considérer
$(\mathscr{O}_X)_{/X'}$ (resp. $\mathscr{F}_{/X'}$) comme un
\emph{faisceau d'anneaux topologiques} (resp. de
\emph{groupes topologiques}) limite projective des faisceaux d'anneaux
(resp. groupes) \emph{pseudo-discrets}
$\mathscr{O}_X/\mathscr{J}$ (resp.
$\mathscr{F}\otimes_{\mathscr{O}_X}
(\mathscr{O}_X/\mathscr{J})=\mathscr{F}/\mathscr{J}\mathscr{F}$), et, par
passage à la limite projective, $\mathscr{F}_{/X'}$ devient alors un
\emph{$(\mathscr{O}_X)_{/X'}$-Module topologique}
(\hyperref[0.3.8.1-fr]{0, 3.8.1} et
\hyperref[0.3.8.2-fr]{3.8.2})~; rappelons que pour tout ouvert
\emph{quasi-compact} $U\subset X$,
$\Gamma(U\cap X',(\mathscr{O}_X)_{/X'})$ (resp.
$\Gamma(U\cap X',\mathscr{F}_{/X'})$) est alors limite projective des
anneaux (resp. groupes) discrets $\Gamma(U,\mathscr{O}_X/\mathscr{J})$
(resp. $\Gamma(U,\mathscr{F}/\mathscr{J}\mathscr{F})$).

Si maintenant $u:\mathscr{F}\to\mathscr{G}$ est un homomorphisme de
$\mathscr{O}_X$-Modules, on en déduit canoniquement des homomorphismes
$u_{\mathscr{J}}:
\mathscr{F}\otimes_{\mathscr{O}_X}(\mathscr{O}_X/\mathscr{J})
\to
\mathscr{G}\otimes_{\mathscr{O}_X}(\mathscr{O}_X/\mathscr{J})$
pour tout $\mathscr{J}\in\Phi$, et ces homomorphismes forment un système
projectif. Par passage à la limite projective et restriction à $X'$, ils
donnent donc un $(\mathscr{O}_X)_{/X'}$-homomorphisme continu
$\mathscr{F}_{/X'}\to\mathscr{G}_{/X'}$, noté $u_{/X'}$ ou $\widehat u$,
et appelé le \emph{complété de l'homomorphisme $u$ le long de $X'$}. Il est
clair que si $v:\mathscr{G}\to\mathscr{H}$ est un second homomorphisme de
$\mathscr{O}_X$-Modules, on a
$(v\circ u)_{/X'}=(v_{/X'})\circ(u_{/X'})$ donc
$\mathscr{F}_{/X'}$ est un \emph{foncteur additif covariant} en
$\mathscr{F}$, de la catégorie des $\mathscr{O}_X$-Modules cohérents, à
valeurs dans la catégorie des $(\mathscr{O}_X)_{/X'}$-Modules topologiques.

\begin{proposition}[10.8.5]
\label{I.10.8.5-fr}
Le support de $(\mathscr{O}_X)_{/X'}$ est $X'$~; l'espace topologiquement
annelé $(X',(\mathscr{O}_X)_{/X'})$ est un préschéma formel localement
noethérien, et si $\mathscr{J}\in\Phi$, $\mathscr{J}_{/X'}$ est un Idéal de
définition de ce préschéma formel. Si $X=\operatorname{Spec}(A)$ est un
schéma affine d'anneau noethérien, $\mathscr{J}=\widetilde{\mathfrak{J}}$
où $\mathfrak{J}$ est un idéal de $A$, et $X'=V(\mathfrak{J})$,
$(X',(\mathscr{O}_X)_{/X'})$ s'identifie canoniquement à
$\operatorname{Spf}(\widehat A)$, où $\widehat A$ est le séparé complété de
$A$ pour la topologie $\mathfrak{J}$-préadique.
\end{proposition}

\begin{proof}
On peut évidemment se borner à prouver la dernière assertion. On sait
(\hyperref[0.7.3.3-fr]{0, 7.3.3}) que le séparé complété
$\widehat{\mathfrak{J}}$ de $\mathfrak{J}$ pour la topologie
$\mathfrak{J}$-préadique s'identifie à l'idéal
$\mathfrak{J}\widehat A$ de $\widehat A$, et que $\widehat A$ est un anneau
$\widehat{\mathfrak{J}}$-adique noethérien tel que
$\widehat A/\widehat{\mathfrak{J}}^n=A/\mathfrak{J}^n$
(\hyperref[0.7.2.6-fr]{0, 7.2.6}). Cette dernière relation montre que les
idéaux premiers ouverts de $\widehat A$ sont les idéaux
$\widehat{\mathfrak p}=\mathfrak p\widehat A$, où $\mathfrak p$ est un
idéal premier de $A$ contenant $\mathfrak{J}$, et que l'on a
$\widehat{\mathfrak p}\cap A=\mathfrak p$, d'où
$\operatorname{Spf}(\widehat A)=X'$. Comme
$\mathscr{O}_X/\mathscr{J}^n=(A/\mathfrak{J}^n)^{\sim}$, la proposition
découle aussitôt des définitions.
\end{proof}

On dit que le préschéma formel ainsi défini est le
\emph{complété de $X$ le long de $X'$} et on le note $X_{/X'}$ ou
$\widehat X$ si aucune confusion n'est à craindre. Lorsque l'on prend
$X'=X$, on peut prendre $\mathscr{J}=0$, et on a donc $X_{/X}=X$.

Il est clair que si $U$ est un sous-préschéma induit sur un ouvert de $X$,
$U_{/(U\cap X')}$ s'identifie canoniquement au sous-préschéma formel induit
par $X_{/X'}$ sur l'ouvert $U\cap X'$ de $X'$.

\begin{corollary}[10.8.6]
\label{I.10.8.6-fr}
Le préschéma (usuel) $\widehat X_{\mathrm{red}}$ est l'unique
sous-préschéma réduit de $X$ ayant pour espace sous-jacent $X'$
(\hyperref[I.5.2.1-fr]{5.2.1}). Pour que $\widehat X$ soit noethérien, il
faut et il suffit que $\widehat X_{\mathrm{red}}$ le soit, et il suffit que
$X$ le soit.
\end{corollary}

\oldpage[I]{196}

La détermination de $\widehat X_{\mathrm{red}}$ étant locale
(\hyperref[I.10.5.4-fr]{10.5.4}), on peut encore supposer que $X$ est un
schéma affine d'anneau noethérien~; avec les notations de
(\hyperref[I.10.8.5-fr]{10.8.5}), l'idéal $\mathfrak{T}$ des éléments
topologiquement nilpotents de $\widehat A$ est l'image réciproque par
l'application canonique
$\widehat A\to\widehat A/\widehat{\mathfrak{J}}=A/\mathfrak{J}$ du
nilradical de $A/\mathfrak{J}$
(\hyperref[0.7.1.3-fr]{0, 7.1.3}), donc $\widehat A/\mathfrak{T}$ est
isomorphe au quotient de $A/\mathfrak{J}$ par son nilradical. La première
assertion résulte donc de (\hyperref[I.10.5.4-fr]{10.5.4}) et
(\hyperref[I.5.1.1-fr]{5.1.1}). Si $\widehat X_{\mathrm{red}}$ est
noethérien, son espace sous-jacent $X'$ l'est aussi, donc les
$X'_n=\operatorname{Spec}(\mathscr{O}_X/\mathscr{J}^n)$ sont noethériens
(\hyperref[I.6.1.2-fr]{6.1.2}) et il en est de même de $\widehat X$
(\hyperref[I.10.6.4-fr]{10.6.4})~; la réciproque est immédiate, en vertu
de (\hyperref[I.6.1.2-fr]{6.1.2}).

\begin{env}[10.8.7]
\label{I.10.8.7-fr}
Les homomorphismes canoniques
$\mathscr{O}_X\to\mathscr{O}_X/\mathscr{J}$ (pour
$\mathscr{J}\in\Phi$) forment un système projectif et donnent donc, par
passage à la limite projective, un homomorphisme de faisceau d'anneaux
$\theta:\mathscr{O}_X\to\psi_*((\mathscr{O}_X)_{/X'})
=\varprojlim_\Phi(\mathscr{O}_X/\mathscr{J})$, en désignant par $\psi$
l'injection canonique $X'\to X$ des espaces sous-jacents. Nous désignerons
par $i$ (ou $i_X$) le morphisme (dit canonique)
\[
  (\psi,\theta):X_{/X'}\to X
\]
d'espaces annelés.

Par tensorisation, pour tout $\mathscr{O}_X$-Module cohérent
$\mathscr{F}$, les homomorphismes canoniques
$\mathscr{O}_X\to\mathscr{O}_X/\mathscr{J}$ donnent des homomorphismes
$\mathscr{F}\to
\mathscr{F}\otimes_{\mathscr{O}_X}(\mathscr{O}_X/\mathscr{J})$ de
$\mathscr{O}_X$-Modules qui forment encore un système projectif, et
donnent donc, par passage à la limite projective, un homomorphisme
canonique fonctoriel
$\gamma:\mathscr{F}\to\psi_*(\mathscr{F}_{/X'})$ de
$\mathscr{O}_X$-Modules.
\end{env}

\begin{proposition}[10.8.8]
\label{I.10.8.8-fr}
\medskip\noindent
\begin{enumerate}
  \item[\rm{(i)}] Le foncteur $\mathscr{F}_{/X'}$ (en $\mathscr{F}$) est
    exact.
  \item[\rm{(ii)}] L'homomorphisme fonctoriel
    $\gamma^\sharp:i^*(\mathscr{F})\to\mathscr{F}_{/X'}$ de
    $(\mathscr{O}_X)_{/X'}$-Modules est un isomorphisme.
\end{enumerate}
\end{proposition}

\medskip\noindent
\begin{enumerate}
  \item[\rm{(i)}] Il suffit de prouver que si
    $0\to\mathscr{F}'\to\mathscr{F}\to\mathscr{F}''\to0$ est une suite
    exacte de $\mathscr{O}_X$-Modules cohérents, et $U$ un ouvert affine
    de $X$, d'anneau noethérien $A$, la suite
    \[
      0\to\Gamma(U\cap X',\mathscr{F}'_{/X'})
      \to\Gamma(U\cap X',\mathscr{F}_{/X'})
      \to\Gamma(U\cap X',\mathscr{F}''_{/X'})\to0
    \]
    est exacte. On a alors $\mathscr{F}|U=\widetilde M$,
    $\mathscr{F}'|U=\widetilde{M'}$, $\mathscr{F}''|U=\widetilde{M''}$,
    où $M$, $M'$, $M''$ sont trois $A$-modules de type fini tels que la
    suite $0\to M'\to M\to M''\to0$ soit exacte
    (\hyperref[I.1.5.1-fr]{1.5.1} et
    \hyperref[I.1.3.11-fr]{1.3.11})~; soit $\mathscr{J}\in\Phi$ et soit
    $\mathfrak{J}$ un idéal de $A$ tel que
    $\mathscr{J}|U=\widetilde{\mathfrak{J}}$. On a alors
    \[
      \Gamma(U\cap X',
      \mathscr{F}\otimes_{\mathscr{O}_X}(\mathscr{O}_X/\mathscr{J}^n))
      =M\otimes_A(A/\mathfrak{J}^n)
    \]
    (\hyperref[I.1.3.12-fr]{1.3.12})~; donc, par définition de la limite
    projective, on a
    \[
      \Gamma(U\cap X',\mathscr{F}_{/X'})
      =\varprojlim_n(M\otimes_A(A/\mathfrak{J}^n))=\widehat M
    \]
    séparé complété de $M$ pour la topologie $\mathfrak{J}$-préadique, et
    de même
    \[
      \Gamma(U\cap X',\mathscr{F}'_{/X'})=\widehat{M'},\qquad
      \Gamma(U\cap X',\mathscr{F}''_{/X'})=\widehat{M''}~;
    \]
    notre assertion résulte alors de ce que lorsque $A$ est noethérien,
    le foncteur $\widehat M$ en $M$ est exact sur la catégorie des
    $A$-modules de type fini (\hyperref[0.7.3.3-fr]{0, 7.3.3}).

\oldpage[I]{197}
  \item[\rm{(ii)}] La question étant locale, on peut supposer que l'on a
    une suite exacte
    $\mathscr{O}_X^m\to\mathscr{O}_X^n\to\mathscr{F}\to0$
    (\hyperref[0.5.3.2-fr]{0, 5.3.2})~; comme $\gamma^\sharp$ est
    fonctoriel, et que les foncteurs $i^*(\mathscr{F})$ et
    $\mathscr{F}_{/X'}$ sont exacts à droite (d'après (i) et
    (\hyperref[0.4.3.1-fr]{0, 4.3.1})), on a le diagramme commutatif
    \[
      \xymatrix{
        i^*(\mathscr{O}_X^m)\ar[r]\ar[d]_{\gamma^\sharp} &
        i^*(\mathscr{O}_X^n)\ar[r]\ar[d]_{\gamma^\sharp} &
        i^*(\mathscr{F})\ar[r]\ar[d]_{\gamma^\sharp} &
        0\\
        (\mathscr{O}_X^m)_{/X'}\ar[r] &
        (\mathscr{O}_X^n)_{/X'}\ar[r] &
        \mathscr{F}_{/X'}\ar[r] &
        0
      }
      \tag{10.8.8.1}\label{I.10.8.8.1-fr}
    \]
    dont les lignes sont exactes. En outre, les deux foncteurs
    $i^*(\mathscr{F})$ et $\mathscr{F}_{/X'}$ commutent aux sommes
    directes finies (\hyperref[0.3.2.6-fr]{0, 3.2.6} et
    \hyperref[0.4.3.2-fr]{4.3.2}) et on est donc ramené à démontrer notre
    assertion pour $\mathscr{F}=\mathscr{O}_X$. On a alors
    $i^*(\mathscr{O}_X)=(\mathscr{O}_X)_{/X'}=
    \mathscr{O}_{\widehat X}$
    (\hyperref[0.4.3.4-fr]{0, 4.3.4}), et $\gamma^\sharp$ est un
    homomorphisme de $\mathscr{O}_{\widehat X}$-Modules~; il suffit donc
    de vérifier que $\gamma^\sharp$ transforme la section unité de
    $\mathscr{O}_{\widehat X}$ au-dessus d'un ouvert de $X'$ en elle-même,
    ce qui est immédiat et montre donc que dans ce cas $\gamma^\sharp$ est
    l'identité.
\end{enumerate}

\begin{corollary}[10.8.9]
\label{I.10.8.9-fr}
Le morphisme d'espaces annelés $i:X_{/X'}\to X$ est plat.
\end{corollary}

Cela résulte en effet de (\hyperref[0.6.7.3-fr]{0, 6.7.3}) et de
(\hyperref[I.10.8.8-fr]{10.8.8, (i)}).

\begin{corollary}[10.8.10]
\label{I.10.8.10-fr}
Si $\mathscr{F}$ et $\mathscr{G}$ sont des $\mathscr{O}_X$-Modules
cohérents, il existe des isomorphismes canoniques fonctoriels (en
$\mathscr{F}$ et $\mathscr{G}$)
\[
  (\mathscr{F}_{/X'})\otimes_{(\mathscr{O}_X)_{/X'}}
  (\mathscr{G}_{/X'})
  \xrightarrow{\sim}
  (\mathscr{F}\otimes_{\mathscr{O}_X}\mathscr{G})_{/X'}
  \tag{10.8.10.1}\label{I.10.8.10.1-fr}
\]
\[
  (\mathscr{H}\!om_{\mathscr{O}_X}(\mathscr{F},\mathscr{G}))_{/X'}
  \xrightarrow{\sim}
  \mathscr{H}\!om_{(\mathscr{O}_X)_{/X'}}
  (\mathscr{F}_{/X'},\mathscr{G}_{/X'})
  \tag{10.8.10.2}\label{I.10.8.10.2-fr}
\]
\end{corollary}

Cela résulte de l'identification canonique de $i^*(\mathscr{F})$ et de
$\mathscr{F}_{/X'}$~; l'existence du premier isomorphisme est alors un
résultat valable pour tous les morphismes d'espaces annelés
(\hyperref[0.4.3.3.1-fr]{0, 4.3.3.1}) et celle du second est un résultat
valable pour tous les morphismes plats
(\hyperref[0.6.7.6-fr]{0, 6.7.6}), donc découle de
(\hyperref[I.10.8.9-fr]{10.8.9}).

\begin{proposition}[10.8.11]
\label{I.10.8.11-fr}
Pour tout $\mathscr{O}_X$-Module cohérent $\mathscr{F}$, le noyau de
l'homomorphisme canonique
$\Gamma(X,\mathscr{F})\to\Gamma(X',\mathscr{F}_{/X'})$ déduit de
$\mathscr{F}\to\mathscr{F}_{/X'}$ est formé des sections nulles dans un
voisinage de $X'$.
\end{proposition}

Il résulte de la définition de $\mathscr{F}_{/X'}$ que l'image canonique
d'une telle section est nulle. Réciproquement, si
$s\in\Gamma(X,\mathscr{F})$ a une image nulle dans
$\Gamma(X',\mathscr{F}_{/X'})$, il suffit de voir que tout $x\in X'$ admet
un voisinage dans $X$ dans lequel $s$ est nulle, et on peut donc se ramener
au cas où $X=\operatorname{Spec}(A)$ est affine, $A$ noethérien,
$X'=V(\mathfrak{J})$, où $\mathfrak{J}$ est un idéal de $A$, et
$\mathscr{F}=\widetilde M$, où $M$ est un $A$-module de type fini. Alors
$\Gamma(X',\mathscr{F}_{/X'})$ est le séparé complété $\widehat M$ de $M$
pour la topologie $\mathfrak{J}$-préadique, et l'homomorphisme
$\Gamma(X,\mathscr{F})\to\Gamma(X',\mathscr{F}_{/X'})$ est
l'homomorphisme canonique $M\to\widehat M$. On sait
(\hyperref[0.7.3.7-fr]{0, 7.3.7}) que le noyau de cet homomorphisme est
l'ensemble des $z\in M$ annulés par un élément de $1+\mathfrak{J}$. On a
donc $(1+f)s=0$ pour un $f\in\mathfrak{J}$~; pour tout $x\in X'$ on en
déduit $(1_x+f_x)s_x=0$, et comme $1_x+f_x$ est inversible dans
$\mathscr{O}_x$ ($\mathfrak{J}_x\mathscr{O}_x$ étant contenu dans l'idéal
maximal de $\mathscr{O}_x$), on a $s_x=0$, ce qui démontre la proposition.

\begin{corollary}[10.8.12]
\label{I.10.8.12-fr}
Le support de $\mathscr{F}_{/X'}$ est égal à
$\operatorname{Supp}(\mathscr{F})\cap X'$.
\end{corollary}

Il est clair que $\mathscr{F}_{/X'}$ est un
$(\mathscr{O}_X)_{/X'}$-Module de type fini
(\hyperref[I.10.8.8-fr]{10.8.8, (ii)}) et
(\hyperref[0.5.2.4-fr]{0, 5.2.4}),

\oldpage[I]{198}

donc son support est fermé (\hyperref[0.5.2.2-fr]{0, 5.2.2}) et
évidemment contenu dans $\operatorname{Supp}(\mathscr{F})\cap X'$. Pour
montrer qu'il est égal à ce dernier ensemble, on est aussitôt ramené à
prouver que la relation $\Gamma(X',\mathscr{F}_{/X'})=0$ entraîne
$\operatorname{Supp}(\mathscr{F})\cap X'=\varnothing$~; or cela résulte de
(\hyperref[I.10.8.11-fr]{10.8.11}) et de
(\hyperref[I.1.4.1-fr]{1.4.1}).

\begin{corollary}[10.8.13]
\label{I.10.8.13-fr}
Soit $u:\mathscr{F}\to\mathscr{G}$ un homomorphisme de
$\mathscr{O}_X$-Modules cohérents. Pour que
$u_{/X'}:\mathscr{F}_{/X'}\to\mathscr{G}_{/X'}$ soit nul, il faut et il
suffit que $u$ soit nul dans un voisinage de $X'$.
\end{corollary}

En effet, d'après (\hyperref[I.10.8.8-fr]{10.8.8, (ii)}), $u_{/X'}$
s'identifie à $i^*(u)$, donc si on considère $u$ comme une section au-dessus
de $X$ du faisceau
$\mathscr{H}=\mathscr{H}\!om_{\mathscr{O}_X}(\mathscr{F},\mathscr{G})$,
$u_{/X'}$ est la section de
$i^*(\mathscr{H})=\mathscr{H}_{/X'}$ au-dessus de $X'$ qui lui correspond
canoniquement ((\hyperref[I.10.8.10.2-fr]{10.8.10.2}) et
(\hyperref[0.4.4.6-fr]{0, 4.4.6})). Il suffit donc d'appliquer
(\hyperref[I.10.8.11-fr]{10.8.11}) au $\mathscr{O}_X$-Module cohérent
$\mathscr{H}$.

\begin{corollary}[10.8.14]
\label{I.10.8.14-fr}
Soit $u:\mathscr{F}\to\mathscr{G}$ un homomorphisme de
$\mathscr{O}_X$-Modules cohérents. Pour que $u_{/X'}$ soit un monomorphisme
(resp. un épimorphisme), il faut et il suffit que $u$ soit un monomorphisme
(resp. un épimorphisme) dans un voisinage de $X'$.
\end{corollary}

Soient $\mathscr{P}$ et $\mathscr{N}$ le conoyau et le noyau de $u$, de
sorte qu'on a la suite exacte
\[
  0\to\mathscr{N}\xrightarrow{v}\mathscr{F}\xrightarrow{u}
  \mathscr{G}\xrightarrow{w}\mathscr{P}\to0,
\]
d'où (\hyperref[I.10.8.8-fr]{10.8.8, (i)}) la suite exacte
\[
  0\to\mathscr{N}_{/X'}\xrightarrow{v_{/X'}}
  \mathscr{F}_{/X'}\xrightarrow{u_{/X'}}
  \mathscr{G}_{/X'}\xrightarrow{w_{/X'}}
  \mathscr{P}_{/X'}\to0.
\]
Si $u_{/X'}$ est un monomorphisme (resp. un épimorphisme), on a
$v_{/X'}=0$ (resp. $w_{/X'}=0$), donc il y a un voisinage de $X'$ dans
lequel $v=0$ (resp. $w=0$) en vertu de
(\hyperref[I.10.8.13-fr]{10.8.13}).

\subsection{Prolongement d'un morphisme aux complétés.}
\label{subsection:I.10.9-fr}

\begin{env}[10.9.1]
\label{I.10.9.1-fr}
Soient $X$, $Y$ deux préschémas (usuels) localement noethériens,
$f:X\to Y$ un morphisme, $X'$ (resp. $Y'$) une partie fermée de l'espace
sous-jacent $X$ (resp. $Y$), telles que $f(X')\subset Y'$. Soit
$\mathscr{J}$ (resp. $\mathscr{K}$) un faisceau d'idéaux de
$\mathscr{O}_X$ (resp. $\mathscr{O}_Y$) tel que le support de
$\mathscr{O}_X/\mathscr{J}$ (resp. $\mathscr{O}_Y/\mathscr{K}$) soit
$X'$ (resp. $Y'$) et que
$f^*(\mathscr{K})\mathscr{O}_X\subset\mathscr{J}$~; on notera qu'il
existe toujours de tels faisceaux d'idéaux, car on peut par exemple prendre
pour $\mathscr{J}$ le plus grand faisceau d'idéaux de $\mathscr{O}_X$
définissant un sous-préschéma de $X$ ayant $X'$ pour espace sous-jacent
(\hyperref[I.5.2.1-fr]{5.2.1}), et l'hypothèse $f(X')\subset Y'$ entraîne
alors $f^*(\mathscr{K})\mathscr{O}_X\subset\mathscr{J}$
(\hyperref[I.5.2.4-fr]{5.2.4}). On a donc pour tout entier $n>0$,
$f^*(\mathscr{K}^n)\mathscr{O}_X\subset\mathscr{J}^n$
(\hyperref[0.4.3.5-fr]{0, 4.3.5})~; par suite
(\hyperref[I.4.4.6-fr]{4.4.6}), si on pose
\[
  X'_n=(X',\mathscr{O}_X/\mathscr{J}^{n+1}),\qquad
  Y'_n=(Y',\mathscr{O}_Y/\mathscr{K}^{n+1}),
\]
on déduit de $f$ un morphisme $f_n:X'_n\to Y'_n$, et il est immédiat que
les $f_n$ forment un système inductif. Nous désignerons sa limite inductive
(\hyperref[I.10.6.8-fr]{10.6.8}) par
$\widehat f:X_{/X'}\to Y_{/Y'}$, et nous dirons (par abus de langage) que
$\widehat f$ est le \emph{prolongement de $f$ aux complétés de $X$ et $Y$
le long de $X'$ et $Y'$}. Il est immédiat de vérifier que ce morphisme ne
dépend pas du choix des faisceaux d'idéaux $\mathscr{J}$, $\mathscr{K}$
vérifiant les conditions ci-dessus. Il suffit de le voir en effet lorsque
$X$ et $Y$ sont des schémas affines noethériens d'anneaux $A$, $B$~; alors
$\mathscr{J}=\widetilde{\mathfrak{J}}$,
$\mathscr{K}=\widetilde{\mathfrak{K}}$, où $\mathfrak{J}$ (resp.
$\mathfrak{K}$) est un idéal de $A$ (resp. $B$), $f$ correspond à un
homomorphisme d'anneaux $\varphi:B\to A$ tel que
$\varphi(\mathfrak{K})\subset\mathfrak{J}$
(\hyperref[I.4.4.6-fr]{4.4.6} et
\hyperref[I.1.7.4-fr]{1.7.4})~; \emph{$\widehat f$ est alors le morphisme
qui correspond (\hyperref[I.10.2.2-fr]{10.2.2}) à l'homomorphisme continu
$\widehat\varphi:\widehat B\to\widehat A$}, où $\widehat A$ (resp.
$\widehat B$) est le séparé complété de $A$ (resp. $B$) pour la topologie
$\mathfrak{J}$-préadique (resp. $\mathfrak{K}$-préadique)
(\hyperref[I.10.6.8-fr]{10.6.8})~; et on sait que si on remplace
$\mathscr{J}$ par un autre

\oldpage[I]{199}

faisceau d'idéaux
$\mathscr{J}'=\widetilde{\mathfrak{J}'}$ tel que le support de
$\mathscr{O}_X/\mathscr{J}'$ soit encore $X'$, les topologies
$\mathfrak{J}$-préadique et $\mathfrak{J}'$-préadique sur $A$ sont les
mêmes (\hyperref[I.10.8.2-fr]{10.8.2}).

On notera que, d'après cette définition, l'application continue
$X'\to Y'$ des espaces sous-jacents à $X_{/X'}$ et $Y_{/Y'}$, qui
correspond à $\widehat f$, n'est autre que la restriction à $X'$ de $f$.
\end{env}

\begin{env}[10.9.2]
\label{I.10.9.2-fr}
Il résulte aussitôt de la définition précédente que le diagramme de
morphismes d'espaces annelés
\[
  \label{I.10.9.2.1-fr}
  \xymatrix{
    \widehat X\ar[r]^{\widehat f}\ar[d]_{i_X} &
    \widehat Y\ar[d]^{i_Y}\\
    X\ar[r]_f &
    Y
  }
\]
est commutatif, les flèches verticales étant les morphismes canoniques
(\hyperref[I.10.8.7-fr]{10.8.7}).
\end{env}

\begin{env}[10.9.3]
\label{I.10.9.3-fr}
Soient $Z$ un troisième préschéma, $g:Y\to Z$ un morphisme, $Z'$ une
partie fermée de $Z$ telle que $g(Y')\subset Z'$. Si $\widehat g$ désigne
le complété le long de $Y'$ et $Z'$ du morphisme $g$, il résulte aussitôt
de (\hyperref[I.10.9.1-fr]{10.9.1}) que l'on a
$\widehat{(g\circ f)}=\widehat g\circ\widehat f$.
\end{env}

\begin{proposition}[10.9.4]
\label{I.10.9.4-fr}
Soient $X$, $Y$ deux $S$-préschémas localement noethériens, $Y$ étant de
type fini sur $S$. Soient $f$, $g$, deux $S$-morphismes de $X$ dans $Y$
tels que $f(X')\subset Y'$, $g(X')\subset Y'$. Pour que
$\widehat f=\widehat g$, il faut et il suffit que $f$ et $g$ coïncident
dans un voisinage de $X'$.
\end{proposition}

La condition est évidemment suffisante (sans hypothèse de finitude sur
$Y$). Pour voir qu'elle est nécessaire, remarquons d'abord que l'hypothèse
$\widehat f=\widehat g$ implique $f(x)=g(x)$ pour tout $x\in X'$. D'autre
part, la question étant locale, on peut supposer que $X$ et $Y$ sont des
voisinages ouverts affines respectifs de $x$ et de $y=f(x)=g(x)$,
d'anneaux noethériens, que $S$ est affine et que
$\Gamma(Y,\mathscr{O}_Y)$ est une $\Gamma(S,\mathscr{O}_S)$-algèbre de
type fini (\hyperref[I.6.3.3-fr]{6.3.3}). Alors $f$ et $g$ correspondent à
deux $\Gamma(S,\mathscr{O}_S)$-homomorphismes $\rho$, $\sigma$ de
$\Gamma(Y,\mathscr{O}_Y)$ dans $\Gamma(X,\mathscr{O}_X)$
(\hyperref[I.1.7.3-fr]{1.7.3}), et par hypothèse, les prolongements par
continuité de ces homomorphismes au séparé complété de
$\Gamma(Y,\mathscr{O}_Y)$ sont les mêmes. On conclut de
(\hyperref[I.10.8.11-fr]{10.8.11}) que pour toute section
$s\in\Gamma(Y,\mathscr{O}_Y)$, les sections $\rho(s)$ et $\sigma(s)$
coïncident dans un voisinage de $X'$ (dépendant de $s$)~; comme
$\Gamma(Y,\mathscr{O}_Y)$ est une algèbre de type fini sur
$\Gamma(S,\mathscr{O}_S)$, on en déduit aussitôt qu'il existe un
voisinage $V$ de $X'$ tel que $\rho(s)$ et $\sigma(s)$ coïncident dans
$V$ pour \emph{toute} section $s\in\Gamma(Y,\mathscr{O}_Y)$. Si
$h\in\Gamma(Y,\mathscr{O}_X)$ est tel que $D(h)$ soit un voisinage de $x$
contenu dans $V$, on conclut de ce qui précède et de
(\hyperref[I.1.4.1-fr]{1.4.1, d)}) que $f$ et $g$ coïncident dans $D(h)$.

\begin{proposition}[10.9.5]
\label{I.10.9.5-fr}
Sous les hypothèses de (\hyperref[I.10.9.1-fr]{10.9.1}), pour tout
$\mathscr{O}_Y$-Module cohérent $\mathscr{G}$, il existe un isomorphisme
canonique fonctoriel de $(\mathscr{O}_X)_{/X'}$-Modules
\[
  (f^*(\mathscr{G}))_{/X'}\xrightarrow{\sim}
  \widehat f^*(\mathscr{G}_{/Y'})
\]
\end{proposition}

Si on identifie canoniquement $(f^*(\mathscr{G}))_{/X'}$ à
$i_X^*(f^*(\mathscr{G}))$ et $\widehat f^*(\mathscr{G}_{/Y'})$ à
$\widehat f^*(i_Y^*(\mathscr{G}))$
(\hyperref[I.10.8.8-fr]{10.8.8}), la proposition résulte aussitôt de la
commutativité du diagramme de
(\hyperref[I.10.9.2-fr]{10.9.2}).

\begin{env}[10.9.6]
\label{I.10.9.6-fr}
Soient maintenant $\mathscr{F}$ un $\mathscr{O}_X$-Module cohérent,
$\mathscr{G}$ un $\mathscr{O}_Y$-Module cohérent. Si
$u:\mathscr{G}\to\mathscr{F}$ est un $f$-morphisme de $\mathscr{G}$ dans
$\mathscr{F}$, il lui correspond un $\mathscr{O}_X$-homomorphisme
$u^\sharp:f^*(\mathscr{G})\to\mathscr{F}$, donc par complétion un
$(\mathscr{O}_X)_{/X'}$-homomorphisme continu
$(u^\sharp)_{/X'}:(f^*(\mathscr{G}))_{/X'}\to\mathscr{F}_{/X'}$, et en
vertu de (\hyperref[I.10.9.5-fr]{10.9.5}) il existe un
$\widehat f$-morphisme $v:\mathscr{G}_{/Y'}\to\mathscr{F}_{/X'}$,
\oldpage[I]{200}

et un seul, tel que $v^\sharp=(u^\sharp)_{/X'}$. Si on considère les
triplets $(\mathscr{F},X,X')$ ($\mathscr{F}$ étant un
$\mathscr{O}_X$-Module cohérent et $X'$ une partie fermée de $X$) comme
une \emph{catégorie}, les morphismes
$(\mathscr{F},X,X')\to(\mathscr{G},Y,Y')$ consistant en un morphisme de
préschémas $f:X\to Y$ tel que $f(X')\subset Y'$ et en un
$f$-morphisme $u:\mathscr{G}\to\mathscr{F}$, on peut donc dire que
$(X_{/X'},\mathscr{F}_{/X'})$ est un \emph{foncteur} en
$(\mathscr{F},X,X')$, prenant ses valeurs dans la catégorie des couples
$(\mathfrak{Z},\mathscr{H})$ formés d'un préschéma formel localement
noethérien $\mathfrak{Z}$ et d'un
$\mathscr{O}_{\mathfrak{Z}}$-Module $\mathscr{H}$, les morphismes de
cette dernière catégorie consistant en les couples formés d'un morphisme
$g$ de préschémas formels et d'un $g$-morphisme.
\end{env}

\begin{proposition}[10.9.7]
\label{I.10.9.7-fr}
Soient $S$, $X$, $Y$ trois préschémas localement noethériens,
$g:X\to S$, $h:Y\to S$ deux morphismes, $S'$ une partie fermée de $S$,
$X'$ (resp. $Y'$) une partie fermée de $X$ (resp. $Y$) telle que
$g(X')\subset S'$ (resp. $h(Y')\subset S'$)~; soit
$Z=X\times_S Y$~; supposons $Z$ localement noethérien, et soit
$Z'=p^{-1}(X')\cap q^{-1}(Y')$, où $p$ et $q$ sont les projections de
$X\times_S Y$. Dans ces conditions, le complété $Z_{/Z'}$ s'identifie au
produit des $S_{/S'}$-préschémas formels
$(X_{/X'})\times_{S_{/S'}}(Y_{/Y'})$, les morphismes structuraux
s'identifiant à $\widehat g$ et $\widehat h$, et les projections à
$\widehat p$ et $\widehat q$.
\end{proposition}

Il est immédiat que la question est locale pour $S$, $X$ et $Y$, et on
est donc ramené au cas où $S=\operatorname{Spec}(A)$,
$X=\operatorname{Spec}(B)$, $Y=\operatorname{Spec}(C)$,
$S'=V(\mathfrak{J})$, $X'=V(\mathfrak{K})$, $Y'=V(\mathfrak{L})$, où
$\mathfrak{J}$, $\mathfrak{K}$, $\mathfrak{L}$ sont trois idéaux tels
que $\varphi(\mathfrak{J})\subset\mathfrak{K}$ et
$\psi(\mathfrak{J})\subset\mathfrak{L}$, en désignant par $\varphi$ et
$\psi$ les homomorphismes $A\to B$ et $A\to C$ qui correspondent à $g$
et $h$. Alors on sait que $Z=\operatorname{Spec}(B\otimes_A C)$ et que
$Z'=V(\mathfrak{M})$, où $\mathfrak{M}$ est l'idéal
$\operatorname{Im}(\mathfrak{K}\otimes_A C)
+\operatorname{Im}(B\otimes_A\mathfrak{L})$. La conclusion résulte
(\hyperref[I.10.7.2-fr]{10.7.2}) de ce que le produit tensoriel complété
$(\widehat B\otimes_{\widehat A}\widehat C)^\wedge$ (où $\widehat A$,
$\widehat B$, $\widehat C$ sont respectivement les séparés complétés de
$A$, $B$, $C$ pour les topologies $\mathfrak{J}$-, $\mathfrak{K}$- et
$\mathfrak{L}$-préadiques) est le séparé complété du produit tensoriel
$B\otimes_A C$ pour la topologie $\mathfrak{M}$-préadique
(\hyperref[0.7.7.2-fr]{0, 7.7.2}).

On notera en outre que si $T$ est un $S$-préschéma localement
noethérien, $u:T\to X$, $v:T\to Y$ deux $S$-morphismes, $T'$ une partie
fermée de $T$ telle que $u(T')\subset X'$, $v(T')\subset Y'$, alors le
prolongement aux complétés $((u,v)_S)^\wedge$ s'identifie à
$(\widehat u,\widehat v)_{S_{/S'}}$.

\begin{corollary}[10.9.8]
\label{I.10.9.8-fr}
Soient $X$, $Y$ deux $S$-préschémas localement noethériens tels que
$X\times_S Y$ soit localement noethérien~; soient $S'$ une partie fermée
de $S$, $X'$ (resp. $Y'$) une partie fermée de $X$ (resp. $Y$) dont
l'image dans $S$ est contenue dans $S'$. Pour tout $S$-morphisme
$f:X\to Y$ tel que $f(X')\subset Y'$, le morphisme graphe
$\Gamma_{\widehat f}$ s'identifie au prolongement $(\Gamma_f)^\wedge$ du
morphisme graphe de $f$.
\end{corollary}

\begin{corollary}[10.9.9]
\label{I.10.9.9-fr}
Soient $X$, $Y$ deux préschémas localement noethériens, $f:X\to Y$ un
morphisme, $Y'$ une partie fermée de $Y$, $X'=f^{-1}(Y')$. Alors le
préschéma $X_{/X'}$ s'identifie par le diagramme commutatif
\[
  \label{I.10.9.9.1-fr}
  \xymatrix{
    X\ar[d]_f &
    X_{/X'}\ar[l]\ar[d]^{\widehat f}\\
    Y &
    Y_{/Y'}\ar[l]
  }
\]
au produit $X\times_Y(Y_{/Y'})$ de préschémas formels.
\end{corollary}

Il suffit d'appliquer (\hyperref[I.10.9.7-fr]{10.9.7}) en remplaçant
$S$ et $S'$ par $Y$, $X$ et $X'$ par $X$.
\oldpage[I]{201}

\begin{remark}[10.9.10]
\label{I.10.9.10-fr}
Si $X$ est la somme $X_1\amalg X_2$
(\hyperref[subsection:I.3.1-fr]{3.1}), $X'$ la réunion
$X'_1\cup X'_2$, où $X'_i$ est une partie fermée de $X_i$ ($i=1,2$),
on voit aussitôt que l'on a
$X_{/X'}={X_1}_{/X'_1}\amalg {X_2}_{/X'_2}$.
\end{remark}

\subsection{Application aux faisceaux cohérents sur les schémas formels
affines.}
\label{subsection:I.10.10-fr}

\begin{env}[10.10.1]
\label{I.10.10.1-fr}
Dans tout ce paragraphe, $A$ désignera un \emph{anneau adique
noethérien}, $\mathfrak{J}$ un idéal de définition de $A$. Soit
$X=\operatorname{Spec}(A)$, $\mathfrak{X}=\operatorname{Spf}(A)$, qui
s'identifie à la partie fermée $V(\mathfrak{J})$ de $X$
(\hyperref[I.10.1.2-fr]{10.1.2}). En outre, la définition
(\hyperref[I.10.1.2-fr]{10.1.2}) et la définition
(\hyperref[I.10.8.4-fr]{10.8.4}) montrent que le
\emph{schéma formel affine $\mathfrak{X}$} est identique au
\emph{complété $X_{/\mathfrak{X}}$} du schéma affine $X$ le long de la
partie fermée $\mathfrak{X}$ de son espace sous-jacent. A tout
$\mathscr{O}_X$-Module cohérent $\mathscr{F}$ correspond donc un
$\mathscr{O}_{\mathfrak{X}}$-Module de type fini
$\mathscr{F}_{/\mathfrak{X}}$, qui est d'ailleurs un faisceau de modules
topologiques sur le faisceau d'anneaux topologiques
$\mathscr{O}_{\mathfrak{X}}$. Mais tout $\mathscr{O}_X$-Module cohérent
$\mathscr{F}$ est de la forme $\widetilde M$, où $M$ est un $A$-module
de type fini (\hyperref[I.1.5.1-fr]{1.5.1})~; nous poserons
$(\widetilde M)_{/\mathfrak{X}}=M^\Delta$. En outre, si $u:M\to N$ est
un $A$-homomorphisme de $A$-modules de type fini, il lui correspond un
homomorphisme $\widetilde u:\widetilde M\to\widetilde N$, et par suite
aussi un homomorphisme continu
$\widetilde u_{/\mathfrak{X}}:(\widetilde M)_{/\mathfrak{X}}\to
(\widetilde N)_{/\mathfrak{X}}$, que nous noterons $u^\Delta$. Il est
immédiat que
$(v\circ u)^\Delta=v^\Delta\circ u^\Delta$~; on a ainsi défini un
\emph{foncteur additif covariant $M^\Delta$} de la catégorie des
$A$-modules de type fini dans celle des
$\mathscr{O}_{\mathfrak{X}}$-Modules de type fini. Lorsque $A$ est un
anneau \emph{discret}, on a $M^\Delta=\widetilde M$.
\end{env}

\begin{proposition}[10.10.2]
\label{I.10.10.2-fr}
\medskip\noindent
\begin{enumerate}
  \item[\rm{(i)}] $M^\Delta$ est un foncteur exact en $M$, et il existe
    un isomorphisme canonique fonctoriel de $A$-modules
    $\Gamma(\mathfrak{X},M^\Delta)\xrightarrow{\sim}M$.
  \item[\rm{(ii)}] Si $M$ et $N$ sont deux $A$-modules de type fini, il
    existe des isomorphismes canoniques fonctoriels
    \[
      (M\otimes_A N)^\Delta
      \xrightarrow{\sim}
      M^\Delta\otimes_{\mathscr{O}_{\mathfrak{X}}}N^\Delta
      \tag{10.10.2.1}\label{I.10.10.2.1-fr}
    \]
    \[
      (\operatorname{Hom}_A(M,N))^\Delta
      \xrightarrow{\sim}
      \mathscr{H}\!om_{\mathscr{O}_{\mathfrak{X}}}
      (M^\Delta,N^\Delta)
      \tag{10.10.2.2}\label{I.10.10.2.2-fr}
    \]
  \item[\rm{(iii)}] L'application $u\to u^\Delta$ est un isomorphisme
    fonctoriel
    \[
      \operatorname{Hom}_A(M,N)
      \xrightarrow{\sim}
      \operatorname{Hom}_{\mathscr{O}_{\mathfrak{X}}}
      (M^\Delta,N^\Delta)
      \tag{10.10.2.3}\label{I.10.10.2.3-fr}
    \]
\end{enumerate}
\end{proposition}

L'exactitude de $M^\Delta$ résulte de l'exactitude des foncteurs
$\widetilde M$ (\hyperref[I.1.3.5-fr]{1.3.5}) et
$\mathscr{F}_{/\mathfrak{X}}$ (\hyperref[I.10.8.8-fr]{10.8.8}). Par
définition, $\Gamma(\mathfrak{X},M^\Delta)$ est le séparé complété du
$A$-module $\Gamma(X,\widetilde M)=M$ pour la topologie
$\mathfrak{J}$-préadique~; mais comme $A$ est complet et $M$ de type
fini, on sait (\hyperref[0.7.3.6-fr]{0, 7.3.6}) que $M$ est séparé et
complet, ce qui achève de prouver (i). L'isomorphisme
(\hyperref[I.10.10.2.1-fr]{10.10.2.1}) (resp.
\hyperref[I.10.10.2.2-fr]{10.10.2.2}) provient de la composition des
isomorphismes (\hyperref[I.1.3.12-fr]{1.3.12, (i)}) et
(\hyperref[I.10.8.10.1-fr]{10.8.10.1}) (resp.
(\hyperref[I.1.3.12-fr]{1.3.12, (ii)}) et
(\hyperref[I.10.8.10.2-fr]{10.8.10.2})). Enfin, comme
$\operatorname{Hom}_A(M,N)$ est un $A$-module de type fini, on peut lui
appliquer (i), qui identifie
$\Gamma(\mathfrak{X},(\operatorname{Hom}_A(M,N))^\Delta)$ à
$\operatorname{Hom}_A(M,N)$, et utiliser
(\hyperref[I.10.10.2.2-fr]{10.10.2.2}), qui prouve que
l'homomorphisme (\hyperref[I.10.10.2.3-fr]{10.10.2.3}) est un
isomorphisme.

On déduit de (\hyperref[I.10.10.2-fr]{10.10.2}) toute une série de
conséquences analogues à celles déduites de
(\hyperref[I.1.3.7-fr]{1.3.7}) et
(\hyperref[I.1.3.12-fr]{1.3.12}), que nous laissons au lecteur le soin
de formuler.
\oldpage[I]{202}

Notons que la propriété d'exactitude de $M^\Delta$, appliqué à la suite
exacte
$0\to\mathfrak{J}\to A\to A/\mathfrak{J}\to0$ montre que le faisceau
d'idéaux de $\mathscr{O}_{\mathfrak{X}}$ désigné ici par
$\mathfrak{J}^\Delta$ coïncide avec celui qui avait été noté de la même
façon en (\hyperref[I.10.3.1-fr]{10.3.1}), en vertu de
(\hyperref[I.10.3.2-fr]{10.3.2}).

\begin{proposition}[10.10.3]
\label{I.10.10.3-fr}
Sous les hypothèses de (\hyperref[I.10.10.1-fr]{10.10.1}),
$\mathscr{O}_{\mathfrak{X}}$ est un faisceau cohérent d'anneaux.
\end{proposition}

Si $f\in A$, on sait que $A_{\{f\}}$ est un anneau adique noethérien
(\hyperref[0.7.6.11-fr]{0, 7.6.11}) et comme la question est locale, on
est ramené (\hyperref[I.10.1.4-fr]{10.1.4}) à prouver que le noyau d'un
homomorphisme
$v:\mathscr{O}_{\mathfrak{X}}^n\to\mathscr{O}_{\mathfrak{X}}$ est un
$\mathscr{O}_{\mathfrak{X}}$-Module de type fini. On a alors
$v=u^\Delta$, où $u$ est un $A$-homomorphisme $A^n\to A$
(\hyperref[I.10.10.2-fr]{10.10.2})~; comme $A$ est noethérien, le noyau
de $u$ est de type fini, autrement dit on a un homomorphisme
$A^m\xrightarrow{w}A^n$ tel que la suite
$A^m\xrightarrow{w}A^n\xrightarrow{u}A$ soit exacte. On en conclut
(\hyperref[I.10.10.2-fr]{10.10.2}) que la suite
$\mathscr{O}_{\mathfrak{X}}^m\xrightarrow{w^\Delta}
\mathscr{O}_{\mathfrak{X}}^n\xrightarrow{v}
\mathscr{O}_{\mathfrak{X}}$ est exacte, ce qui prouve que le noyau de
$v$ est de type fini.

\begin{env}[10.10.4]
\label{I.10.10.4-fr}
Avec les notations précédentes, posons
$A_n=A/\mathfrak{J}^{n+1}$, et soit $X_n$ le schéma affine
$\operatorname{Spec}(A_n)=(\mathfrak{X},
\mathscr{O}_{\mathfrak{X}}/\mathscr{J}^{n+1})$,
$\mathscr{J}=\mathfrak{J}^\Delta$ étant le faisceau d'idéaux de
définition de $\mathscr{O}_{\mathfrak{X}}$ correspondant à l'idéal
$\mathfrak{J}$. Soit $u_{mn}$ le morphisme de préschémas
$X_m\to X_n$ correspondant à l'homomorphisme canonique
$A_n\to A_m$ pour $m\leq n$~; le schéma formel $\mathfrak{X}$ est
limite inductive des $X_n$ pour les $u_{mn}$
(\hyperref[I.10.6.3-fr]{10.6.3}).
\end{env}

\begin{proposition}[10.10.5]
\label{I.10.10.5-fr}
Sous les hypothèses de (\hyperref[I.10.10.1-fr]{10.10.1}), soit
$\mathscr{F}$ un $\mathscr{O}_{\mathfrak{X}}$-Module. Les conditions
suivantes sont équivalentes~:
\begin{enumerate}
  \item[a)] $\mathscr{F}$ est un $\mathscr{O}_{\mathfrak{X}}$-Module
    cohérent.
  \item[b)] $\mathscr{F}$ est isomorphe à la limite projective
    (\hyperref[I.10.6.6-fr]{10.6.6}) d'une suite $(\mathscr{F}_n)$ de
    $\mathscr{O}_{X_n}$-Modules cohérents tels que
    $u_{mn}^*(\mathscr{F}_n)=\mathscr{F}_m$.
  \item[c)] Il existe un $A$-module de type fini $M$ (déterminé à un
    isomorphisme canonique près par
    (\hyperref[I.10.10.2-fr]{10.10.2, (i)})) tel que $\mathscr{F}$ soit
    isomorphe à $M^\Delta$.
\end{enumerate}
\end{proposition}

Montrons d'abord que b) implique c). On a
$\mathscr{F}_n=\widetilde{M_n}$, où $M_n$ est un $A_n$-module de type
fini, et l'hypothèse entraîne que
$M_m=M_n\otimes_{A_n}A_m$ pour $m\leq n$
(\hyperref[I.1.6.5-fr]{1.6.5})~; les $M_n$ forment donc un système
projectif pour les di-homomorphismes canoniques $M_n\to M_m$
($m\leq n$), et il résulte aussitôt de la définition des $A_n$ que ce
système projectif vérifie les conditions de
(\hyperref[0.7.2.9-fr]{0, 7.2.9})~; sa limite projective $M$ est par
suite un $A$-module de type fini tel que
$M_n=M\otimes_A A_n$ pour tout $n$. On en déduit que $\mathscr{F}_n$
est induit sur $X_n$ par
$\widetilde M\otimes_{\mathscr{O}_X}
(\mathscr{O}_X/\widetilde{\mathfrak{J}}^{n+1})$, donc
$\mathscr{F}=M^\Delta$ par définition
(\hyperref[I.10.8.4-fr]{10.8.4}).

Inversement, c) entraîne b)~; en effet, si $u_n$ est le morphisme
d'immersion $X_n\to X$,
$u_n^*(\widetilde M)=(M\otimes_A A_n)^{\sim}$ est induit sur $X_n$ par
$\widetilde M\otimes_{\mathscr{O}_X}
(\mathscr{O}_X/\widetilde{\mathfrak{J}}^{n+1})$, et
$M^\Delta=\varprojlim u_n^*(\widetilde M)$ par définition
(\hyperref[I.10.8.4-fr]{10.8.4})~; comme
$u_m=u_n\circ u_{mn}$ pour $m\leq n$, les
$\mathscr{F}_n=u_n^*(\widetilde M)$ vérifient les conditions de b),
d'où notre assertion.

Montrons maintenant que c) implique a)~: en effet, on a par définition
$\mathscr{O}_{\mathfrak{X}}=A^\Delta$~; $M$ étant conoyau d'un
homomorphisme $A^m\to A^n$, il résulte de
(\hyperref[I.10.10.2-fr]{10.10.2}) que $M^\Delta$ est conoyau d'un
homomorphisme
$\mathscr{O}_{\mathfrak{X}}^m\to\mathscr{O}_{\mathfrak{X}}^n$, et
comme le faisceau d'anneaux $\mathscr{O}_{\mathfrak{X}}$ est cohérent
(\hyperref[I.10.10.3-fr]{10.10.3}), il en est de même de $M^\Delta$
(\hyperref[0.5.3.4-fr]{0, 5.3.4}).
\oldpage[I]{203}

Enfin, a) entraîne b). Considéré comme
$\mathscr{O}_{\mathfrak{X}}$-Module, on a
$\mathscr{O}_{X_n}=\mathscr{O}_{\mathfrak{X}}/
\mathscr{J}^{n+1}=A_n^\Delta$~; $\mathscr{F}_n=\mathscr{F}
\otimes_{\mathscr{O}_{\mathfrak{X}}}\mathscr{O}_{X_n}$ est un
$\mathscr{O}_{\mathfrak{X}}$-Module cohérent
(\hyperref[0.5.3.5-fr]{0, 5.3.5}), et comme c'est aussi un
$\mathscr{O}_{X_n}$-Module et que $\mathscr{J}^{n+1}$ est cohérent,
on en conclut que $\mathscr{F}_n$ est un
$\mathscr{O}_{X_n}$-Module cohérent
(\hyperref[0.5.3.10-fr]{0, 5.3.10}), et il est immédiat que
$u_{mn}^*(\mathscr{F}_n)=\mathscr{F}_m$ pour $m\leq n$ (en se
souvenant que l'application continue $X_m\to X_n$ des espaces
sous-jacents est l'identité de $\mathfrak{X}$). Le faisceau
$\mathscr{G}=\varprojlim\mathscr{F}_n$ est donc un
$\mathscr{O}_{\mathfrak{X}}$-Module cohérent, puisqu'on a vu que b)
entraîne a). Les homomorphismes canoniques
$\mathscr{F}\to\mathscr{F}_n$ forment un système projectif, qui par
passage à la limite donne un homomorphisme canonique
$w:\mathscr{F}\to\mathscr{G}$, et tout revient à démontrer que $w$ est
bijectif. La question étant maintenant locale, on peut se borner au cas
où $\mathscr{F}$ est conoyau d'un homomorphisme
$\mathscr{O}_{\mathfrak{X}}^p\to\mathscr{O}_{\mathfrak{X}}^q$~; cet
homomorphisme étant de la forme $v^\Delta$, où $v$ est un homomorphisme
$A^m\to A^n$ (\hyperref[I.10.10.2-fr]{10.10.2}), $\mathscr{F}$ est
isomorphe à $M^\Delta$, où $M=\operatorname{Coker}v$
(\hyperref[I.10.10.2-fr]{10.10.2}). On a alors, en vertu de
(\hyperref[I.10.10.2-fr]{10.10.2}),
$\mathscr{F}_n=M^\Delta\otimes_{\mathscr{O}_{\mathfrak{X}}}A_n^\Delta
=(M\otimes_A A_n)^\Delta$, et comme la topologie $\mathfrak{J}$-adique
sur $M\otimes_A A_n$ est discrète, on a
$(M\otimes_A A_n)^\Delta=(M\otimes_A A_n)^{\sim}$ (en tant que
$\mathscr{O}_{X_n}$-Module)~; on a vu plus haut que
$M^\Delta=\varprojlim\mathscr{F}_n$ et $w$ est donc bien dans ce cas
l'identité.
\par\hfill C.Q.F.D.\par

\begin{corollary}[10.10.6]
\label{I.10.10.6-fr}
Si $\mathscr{F}$ vérifie la condition b) de
(\hyperref[I.10.10.5-fr]{10.10.5}), le système projectif
$(\mathscr{F}_n)$ est isomorphe au système des
$\mathscr{F}\otimes_{\mathscr{O}_{\mathfrak{X}}}\mathscr{O}_{X_n}$.
\end{corollary}

\begin{env}[10.10.7]
\label{I.10.10.7-fr}
Soient maintenant $A$, $B$ deux anneaux adiques noethériens,
$\varphi:B\to A$ un homomorphisme continu~; on désignera par
$\mathfrak{J}$ (resp. $\mathfrak{K}$) un idéal de définition de $A$
(resp. $B$), tels que $\varphi(\mathfrak{K})\subset\mathfrak{J}$, et
on posera $X=\operatorname{Spec}(A)$, $Y=\operatorname{Spec}(B)$,
$\mathfrak{X}=\operatorname{Spf}(A)$,
$\mathfrak{Y}=\operatorname{Spf}(B)$. Soient $f:X\to Y$ le morphisme
de préschémas correspondant à $\varphi$
(\hyperref[I.1.6.1-fr]{1.6.1}), $\widehat f:\mathfrak{X}\to
\mathfrak{Y}$ son prolongement aux complétés
(\hyperref[I.10.9.1-fr]{10.9.1}), qui est aussi le morphisme de
préschémas formels correspondant à $\varphi$
(\hyperref[I.10.2.2-fr]{10.2.2}).
\end{env}

\begin{proposition}[10.10.8]
\label{I.10.10.8-fr}
Pour tout $B$-module $N$ de type fini, il existe un isomorphisme
canonique fonctoriel de $\mathscr{O}_{\mathfrak{X}}$-Modules
\[
  \widehat f^*(N^\Delta)
  \xrightarrow{\sim}
  (N\otimes_B A)^\Delta.
\]
\end{proposition}

En effet, en désignant par $i_X:\mathfrak{X}\to X$ et
$i_Y:\mathfrak{Y}\to Y$ les morphismes canoniques, on a
(\hyperref[I.10.8.8-fr]{10.8.8}), à des isomorphismes canoniques
fonctoriels près, $N^\Delta=i_Y^*(\widetilde N)$ et
\[
  (N\otimes_B A)^\Delta
  =i_X^*((N\otimes_B A)^{\sim})
  =i_X^*(f^*(\widetilde N))
\]
(\hyperref[I.1.6.5-fr]{1.6.5})~; la proposition résulte donc de la
commutativité du diagramme (\hyperref[I.10.9.2-fr]{10.9.2}).

\begin{corollary}[10.10.9]
\label{I.10.10.9-fr}
Pour tout idéal $\mathfrak b$ de $B$, on a
$\widehat f^*(\mathfrak b^\Delta)\mathscr{O}_{\mathfrak{X}}
=(\mathfrak bA)^\Delta$.
\end{corollary}

En effet, soit $j$ l'injection canonique $\mathfrak b\to B$, à laquelle
il correspond l'injection canonique
$j^\Delta:\mathfrak b^\Delta\to\mathscr{O}_{\mathfrak{Y}}$ de
faisceaux de $\mathscr{O}_{\mathfrak{Y}}$-Modules~; par définition,
$\widehat f^*(\mathfrak b^\Delta)\mathscr{O}_{\mathfrak{X}}$ est
l'image de l'homomorphisme
$\widehat f^*(j^\Delta):\widehat f^*(\mathfrak b^\Delta)\to
\mathscr{O}_{\mathfrak{X}}=\widehat f^*(\mathscr{O}_{\mathfrak{Y}})$~;
mais cet homomorphisme s'identifie à
$(j\otimes1)^\Delta:(\mathfrak b\otimes_B A)^\Delta\to
\mathscr{O}_{\mathfrak{X}}=(B\otimes_B A)^\Delta$ d'après
(\hyperref[I.10.10.8-fr]{10.10.8}). Comme l'image de $j\otimes1$ est
l'idéal $\mathfrak bA$ de $A$, l'image de $(j\otimes1)^\Delta$ est
donc $(\mathfrak bA)^\Delta$ en vertu de
(\hyperref[I.10.10.2-fr]{10.10.2}), d'où la conclusion.
\oldpage[I]{204}

\subsection{Faisceaux cohérents sur les préschémas formels.}
\label{subsection:I.10.11-fr}

\begin{proposition}[10.11.1]
\label{I.10.11.1-fr}
Si $\mathfrak{X}$ est un préschéma formel localement noethérien, le
faisceau d'anneaux $\mathscr{O}_{\mathfrak{X}}$ est cohérent et tout
faisceau d'idéaux de définition de $\mathfrak{X}$ est cohérent.
\end{proposition}

La question étant locale, on est ramené au cas d'un schéma formel affine
noethérien, et la proposition résulte donc de
(\hyperref[I.10.10.3-fr]{10.10.3}) et
(\hyperref[I.10.10.5-fr]{10.10.5}).

\begin{env}[10.11.2]
\label{I.10.11.2-fr}
Soient $\mathfrak{X}$ un préschéma formel localement noethérien,
$\mathscr{J}$ un faisceau d'idéaux de définition de $\mathfrak{X}$,
$X_n$ le préschéma (usuel) localement noethérien
$(\mathfrak{X},\mathscr{O}_{\mathfrak{X}}/\mathscr{J}^{n+1})$, de
sorte que $\mathfrak{X}$ est limite inductive de la suite $(X_n)$ pour
les morphismes canoniques $u_{mn}:X_m\to X_n$
(\hyperref[I.10.6.3-fr]{10.6.3}). Avec ces notations~:
\end{env}

\begin{theorem}[10.11.3]
\label{I.10.11.3-fr}
Pour qu'un $\mathscr{O}_{\mathfrak{X}}$-Module $\mathscr{F}$ soit
cohérent, il faut et il suffit qu'il soit isomorphe à une limite
projective d'une suite $(\mathscr{F}_n)$, où $\mathscr{F}_n$ est un
$\mathscr{O}_{X_n}$-Module cohérent, tel que
$u_{mn}^*(\mathscr{F}_n)=\mathscr{F}_m$ pour $m\leq n$
(\hyperref[I.10.6.6-fr]{10.6.6}). Le système projectif
$(\mathscr{F}_n)$ est alors isomorphe au système des
$u_n^*(\mathscr{F})=\mathscr{F}
\otimes_{\mathscr{O}_{\mathfrak{X}}}\mathscr{O}_{X_n}$, $u_n$ étant
le morphisme canonique $X_n\to\mathfrak{X}$.
\end{theorem}

La question étant locale, on est ramené au cas où $\mathfrak{X}$ est
un schéma formel affine noethérien, et le théorème est alors conséquence
de (\hyperref[I.10.10.5-fr]{10.10.5}) et
(\hyperref[I.10.10.6-fr]{10.10.6}).

On peut donc dire que \emph{la donnée d'un
$\mathscr{O}_{\mathfrak{X}}$-Module cohérent équivaut à celle d'un
système projectif $(\mathscr{F}_n)$ de
$\mathscr{O}_{X_n}$-Modules cohérents tels que
$u_{mn}^*(\mathscr{F}_n)=\mathscr{F}_m$ pour $m\leq n$}.

\begin{corollary}[10.11.4]
\label{I.10.11.4-fr}
Si $\mathscr{F}$ et $\mathscr{G}$ sont deux
$\mathscr{O}_{\mathfrak{X}}$-Modules cohérents, on peut (avec les
notations de (\hyperref[I.10.11.3-fr]{10.11.3})) définir un
isomorphisme canonique fonctoriel
\[
  \label{I.10.11.4.1-fr}
  \operatorname{Hom}_{\mathscr{O}_{\mathfrak{X}}}
  (\mathscr{F},\mathscr{G})
  \xrightarrow{\sim}
  \varprojlim_n
  \operatorname{Hom}_{\mathscr{O}_{X_n}}
  (\mathscr{F}_n,\mathscr{G}_n).
  \tag{10.11.4.1}
\]
\end{corollary}

La limite projective du second membre doit s'entendre pour les
applications
$\theta_n\mapsto u_{mn}^*(\theta_n)$ ($m\leq n$) de
$\operatorname{Hom}_{\mathscr{O}_{X_n}}
(\mathscr{F}_n,\mathscr{G}_n)$ dans
$\operatorname{Hom}_{\mathscr{O}_{X_m}}
(\mathscr{F}_m,\mathscr{G}_m)$. L'homomorphisme
(\hyperref[I.10.11.4.1-fr]{10.11.4.1}) fait correspondre à un élément
$\theta\in\operatorname{Hom}_{\mathscr{O}_{\mathfrak{X}}}
(\mathscr{F},\mathscr{G})$ la suite $(u_n^*(\theta))$~; on voit
aussitôt qu'on en définit un homomorphisme réciproque du précédent en
faisant correspondre au système projectif
$(\theta_n)\in\varprojlim_n
\operatorname{Hom}_{\mathscr{O}_{X_n}}
(\mathscr{F}_n,\mathscr{G}_n)$
sa limite projective dans
$\operatorname{Hom}_{\mathscr{O}_{\mathfrak{X}}}
(\mathscr{F},\mathscr{G})$, compte tenu de
(\hyperref[I.10.11.3-fr]{10.11.3}).

\begin{corollary}[10.11.5]
\label{I.10.11.5-fr}
Pour qu'un homomorphisme $\theta:\mathscr{F}\to\mathscr{G}$ soit
surjectif, il faut et il suffit que l'homomorphisme correspondant
$\theta_0=u_0^*(\theta):\mathscr{F}_0\to\mathscr{G}_0$ le soit.
\end{corollary}

La question étant locale, on est ramené au cas où
$\mathfrak{X}=\operatorname{Spf}(A)$, $A$ étant adique noethérien,
$\mathscr{F}=M^\Delta$, $\mathscr{G}=N^\Delta$ et $\theta=u^\Delta$,
où $M$ et $N$ sont des $A$-modules de type fini et $u$ un
homomorphisme $M\to N$~; on a en outre alors
$\theta_0=\widetilde{u_0}$, où $u_0$ est l'homomorphisme
$u\otimes1:M\otimes_A A/\mathfrak{J}\to
N\otimes_A A/\mathfrak{J}$~;
la conclusion résulte de ce que $\theta$ et $u$ (resp. $\theta_0$ et
$u_0$) sont simultanément surjectifs
(\hyperref[I.1.3.9-fr]{1.3.9} et
\hyperref[I.10.10.2-fr]{10.10.2}) et de ce que $u$ et $u_0$ sont
simultanément surjectifs (\hyperref[0.7.1.14-fr]{0, 7.1.14}).

\begin{env}[10.11.6]
\label{I.10.11.6-fr}
Le th. (\hyperref[I.10.11.3-fr]{10.11.3}) montre qu'on peut considérer
tout $\mathscr{O}_{\mathfrak{X}}$-Module cohérent $\mathscr{F}$ comme
un $\mathscr{O}_{\mathfrak{X}}$-Module \emph{topologique}, en le
considérant comme limite projective des faisceaux de groupes
\emph{pseudo-discrets} $\mathscr{F}_n$
(\hyperref[0.3.8.1-fr]{0, 3.8.1}). Il résulte alors de
(\hyperref[I.10.11.4-fr]{10.11.4}) que tout homomorphisme
$u:\mathscr{F}\to\mathscr{G}$ de
$\mathscr{O}_{\mathfrak{X}}$-Modules cohérents est automatiquement
\emph{continu}

\oldpage[I]{205}

(\hyperref[0.3.8.2-fr]{0, 3.8.2}). En outre, si $\mathscr{H}$ est un
sous-$\mathscr{O}_{\mathfrak{X}}$-Module cohérent d'un
$\mathscr{O}_{\mathfrak{X}}$-Module cohérent $\mathscr{F}$, pour tout
ouvert $U\subset\mathfrak{X}$, $\Gamma(U,\mathscr{H})$ est un sous-groupe
\emph{fermé} du groupe topologique $\Gamma(U,\mathscr{F})$ car le foncteur
$\Gamma$ étant exact à gauche, $\Gamma(U,\mathscr{H})$ est le noyau de
l'homomorphisme
$\Gamma(U,\mathscr{F})\to\Gamma(U,\mathscr{F}/\mathscr{H})$, qui est
\emph{continu} d'après ce qui précède, puisque $\mathscr{F}/\mathscr{H}$
est cohérent (\hyperref[0.5.3.4-fr]{0, 5.3.4})~; notre assertion résulte
de ce que $\Gamma(U,\mathscr{F}/\mathscr{H})$ est un groupe topologique
séparé.
\end{env}

\begin{proposition}[10.11.7]
\label{I.10.11.7-fr}
Soient $\mathscr{F}$ et $\mathscr{G}$ deux
$\mathscr{O}_{\mathfrak{X}}$-Modules cohérents. On peut définir (avec les
notations de (\hyperref[I.10.11.3-fr]{10.11.3})) des isomorphismes
canoniques fonctoriels de $\mathscr{O}_{\mathfrak{X}}$-Modules
topologiques (\hyperref[I.10.11.6-fr]{10.11.6})
\[
  \label{I.10.11.7.1-fr}
  \mathscr{F}\otimes_{\mathscr{O}_{\mathfrak{X}}}\mathscr{G}
  \xrightarrow{\sim}
  \varprojlim_n
  (\mathscr{F}_n\otimes_{\mathscr{O}_{X_n}}\mathscr{G}_n)
  \tag{10.11.7.1}
\]
\[
  \label{I.10.11.7.2-fr}
  \mathscr{H}\!om_{\mathscr{O}_{\mathfrak{X}}}
  (\mathscr{F},\mathscr{G})
  \xrightarrow{\sim}
  \varprojlim_n
  \mathscr{H}\!om_{\mathscr{O}_{X_n}}(\mathscr{F}_n,\mathscr{G}_n)
  \tag{10.11.7.2}
\]
\end{proposition}

L'existence de l'isomorphisme
(\hyperref[I.10.11.7.1-fr]{10.11.7.1}) résulte de la formule
\[
  \mathscr{F}_n\otimes_{\mathscr{O}_{X_n}}\mathscr{G}_n
  =
  (\mathscr{F}\otimes_{\mathscr{O}_{\mathfrak{X}}}
  \mathscr{O}_{X_n})
  \otimes_{\mathscr{O}_{X_n}}
  (\mathscr{G}\otimes_{\mathscr{O}_{\mathfrak{X}}}
  \mathscr{O}_{X_n})
  =
  (\mathscr{F}\otimes_{\mathscr{O}_{\mathfrak{X}}}\mathscr{G})
  \otimes_{\mathscr{O}_{\mathfrak{X}}}\mathscr{O}_{X_n}
\]
et de (\hyperref[I.10.11.3-fr]{10.11.3}). L'isomorphisme
(\hyperref[I.10.11.7.2-fr]{10.11.7.2}) où les deux membres sont
considérés comme faisceaux de modules sans topologie, résulte de la
définition des sections de
$\mathscr{H}\!om_{\mathscr{O}_{\mathfrak{X}}}
(\mathscr{F},\mathscr{G})$ et
$\mathscr{H}\!om_{\mathscr{O}_{X_n}}(\mathscr{F}_n,\mathscr{G}_n)$ et de
l'existence de l'isomorphisme
(\hyperref[I.10.11.4.1-fr]{10.11.4.1}), appliqué au préschéma induit sur
un ouvert formel affine noethérien arbitraire de $\mathfrak{X}$. Reste à
prouver que l'isomorphisme
(\hyperref[I.10.11.7.2-fr]{10.11.7.2}) est bicontinu au-dessus d'un
ensemble quasi-compact, et on est donc ramené au cas où
$\mathfrak{X}=\operatorname{Spf}(A)$, $A$ étant adique noethérien, d'où
(\hyperref[I.10.10.5-fr]{10.10.5})
$\mathscr{F}=M^\Delta$, $\mathscr{G}=N^\Delta$, $M$, $N$ étant des
$A$-modules de type fini~; compte tenu de
(\hyperref[I.10.10.2.1-fr]{10.10.2.1}),
(\hyperref[I.10.10.2.3-fr]{10.10.2.3}) et
(\hyperref[I.1.3.12-fr]{1.3.12, (ii)}), on est ramené à montrer que
l'isomorphisme canonique
$\operatorname{Hom}_A(M,N)\xrightarrow{\sim}
\varprojlim_n\operatorname{Hom}_{A_n}(M_n,N_n)$ (avec
$M_n=M\otimes_A A_n$, $N_n=N\otimes_A A_n$) est continu, ce qui a été
prouvé dans (\hyperref[0.7.8.2-fr]{0, 7.8.2}).

\begin{env}[10.11.8]
\label{I.10.11.8-fr}
Comme $\operatorname{Hom}_{\mathscr{O}_{\mathfrak{X}}}
(\mathscr{F},\mathscr{G})$ est le groupe des sections du faisceau de
groupes topologiques
$\mathscr{H}\!om_{\mathscr{O}_{\mathfrak{X}}}
(\mathscr{F},\mathscr{G})$, il est muni d'une topologie de groupe. Si
$\mathfrak{X}$ est \emph{noethérien}, il résulte de
(\hyperref[I.10.11.7.2-fr]{10.11.7.2}) qu'un système fondamental de
voisinages de $0$ dans ce groupe s'obtient en prenant les sous-groupes
$\operatorname{Hom}_{\mathscr{O}_{\mathfrak{X}}}
(\mathscr{F},\mathscr{J}^n\mathscr{G})$ ($n$ arbitraire).
\end{env}

\begin{proposition}[10.11.9]
\label{I.10.11.9-fr}
Soient $\mathfrak{X}$ un préschéma formel noethérien, $\mathscr{F}$ et
$\mathscr{G}$ deux $\mathscr{O}_{\mathfrak{X}}$-Modules cohérents. Dans
le groupe topologique
$\operatorname{Hom}_{\mathscr{O}_{\mathfrak{X}}}
(\mathscr{F},\mathscr{G})$ les homomorphismes surjectifs (resp.
injectifs, bijectifs) forment une partie ouverte.
\end{proposition}

En vertu de (\hyperref[I.10.11.5-fr]{10.11.5}), l'ensemble des
homomorphismes surjectifs dans
$\operatorname{Hom}_{\mathscr{O}_{\mathfrak{X}}}
(\mathscr{F},\mathscr{G})$ est l'image réciproque par l'application
continue
$\operatorname{Hom}_{\mathscr{O}_{\mathfrak{X}}}
(\mathscr{F},\mathscr{G})\to
\operatorname{Hom}_{\mathscr{O}_{X_0}}(\mathscr{F}_0,\mathscr{G}_0)$
d'une partie du groupe discret
$\operatorname{Hom}_{\mathscr{O}_{X_0}}(\mathscr{F}_0,\mathscr{G}_0)$
d'où la première assertion. Pour démontrer la seconde, recouvrons
$\mathfrak{X}$ par un nombre fini d'ouverts formels affines noethériens
$U_i$. Pour que
$\theta\in\operatorname{Hom}_{\mathscr{O}_{\mathfrak{X}}}
(\mathscr{F},\mathscr{G})$ soit injectif, il faut et il suffit que toutes
ses images par les applications (continues) de restriction
$\operatorname{Hom}_{\mathscr{O}_{\mathfrak{X}}}
(\mathscr{F},\mathscr{G})\to
\operatorname{Hom}_{\mathscr{O}_{\mathfrak{X}}|U_i}
(\mathscr{F}|U_i,\mathscr{G}|U_i)$ le soient~; on est donc ramené au cas
affine, et alors cela a déjà été prouvé dans
(\hyperref[0.7.8.3-fr]{0, 7.8.3}).

\oldpage[I]{206}

\subsection{Morphismes adiques de préschémas formels.}
\label{subsection:I.10.12-fr}

\begin{env}[10.12.1]
\label{I.10.12.1-fr}
Soient $\mathfrak{X}$, $\mathfrak{S}$ deux préschémas formels
\emph{localement noethériens}~; nous dirons qu'un morphisme
$f:\mathfrak{X}\to\mathfrak{S}$ est \emph{adique} s'il existe un Idéal
de définition $\mathscr{J}$ de $\mathfrak{S}$ tel que
$\mathscr{K}=f^*(\mathscr{J})\mathscr{O}_{\mathfrak{X}}$ soit un Idéal de
définition de $\mathfrak{X}$~; on dit aussi alors que $\mathfrak{X}$ est
un \emph{$\mathfrak{S}$-préschéma adique} (pour $f$). Lorsqu'il en est
ainsi, pour \emph{tout} Idéal de définition $\mathscr{J}_1$ de
$\mathfrak{S}$,
$\mathscr{K}_1=f^*(\mathscr{J}_1)\mathscr{O}_{\mathfrak{X}}$ est un Idéal
de définition de $\mathfrak{X}$. En effet, la question étant locale, on
peut supposer $\mathfrak{X}$ et $\mathfrak{S}$ affines noethériens~; il
existe donc un entier $n$ tel que $\mathscr{J}^n\subset\mathscr{J}_1$ et
$\mathscr{J}_1^n\subset\mathscr{J}$
(\hyperref[I.10.3.6-fr]{10.3.6} et
\hyperref[0.7.1.4-fr]{0, 7.1.4}), d'où
$\mathscr{K}^n\subset\mathscr{K}_1$ et
$\mathscr{K}_1^n\subset\mathscr{K}$. La première de ces relations prouve
que $\mathscr{K}_1=\mathfrak{K}_1^\Delta$, où $\mathfrak{K}_1$ est un
idéal ouvert de
$A=\Gamma(\mathfrak{X},\mathscr{O}_{\mathfrak{X}})$, et la seconde prouve
que $\mathfrak{K}_1$ est un idéal de définition de $A$
(\hyperref[0.7.1.4-fr]{0, 7.1.4}), d'où notre assertion.
\end{env}

Il résulte aussitôt de ce qui précède que si $\mathfrak{X}$ et
$\mathfrak{Y}$ sont deux $\mathfrak{S}$-préschémas adiques, \emph{tout
$\mathfrak{S}$-morphisme $u:\mathfrak{X}\to\mathfrak{Y}$ est adique}~:
en effet, si $f:\mathfrak{X}\to\mathfrak{S}$,
$g:\mathfrak{Y}\to\mathfrak{S}$ sont les morphismes structuraux, et
$\mathscr{J}$ un Idéal de définition de $\mathfrak{S}$, on a
$f=g\circ u$, donc
$u^*(g^*(\mathscr{J})\mathscr{O}_{\mathfrak{Y}})
\mathscr{O}_{\mathfrak{X}}=f^*(\mathscr{J})
\mathscr{O}_{\mathfrak{X}}$ est un Idéal de définition de
$\mathfrak{X}$, et par hypothèse
$g^*(\mathscr{J})\mathscr{O}_{\mathfrak{Y}}$ est un Idéal de définition
de $\mathfrak{Y}$.

\begin{env}[10.12.2]
\label{I.10.12.2-fr}
Dans ce qui suit, nous supposerons fixé un préschéma formel localement
noethérien $\mathfrak{S}$ et un Idéal de définition $\mathscr{J}$ de
$\mathfrak{S}$~; nous poserons
$S_n=(\mathfrak{S},\mathscr{O}_{\mathfrak{S}}/\mathscr{J}^{n+1})$. Les
$\mathfrak{S}$-préschémas adiques (localement noethériens) forment
évidemment une \emph{catégorie}. Nous dirons qu'un système inductif
$(X_n)$ de $S_n$-préschémas (usuels) localement noethériens est un
\emph{$(S_n)$-système inductif adique} si les morphismes structuraux
$f_n:X_n\to S_n$ sont tels que, pour $m\leq n$, les diagrammes
\[
  \label{I.10.12.2.1-fr}
  \xymatrix{
    X_n \ar[d]_{f_n}
    & X_m \ar[l] \ar[d]^{f_m}\\
    S_n
    & S_m \ar[l]
  }
  \tag{10.12.2.1}
\]
soient commutatifs et \emph{identifient $X_m$ au produit
$X_n\times_{S_n}S_m=(X_n)_{(S_m)}$}. Les systèmes inductifs adiques
forment une \emph{catégorie}~: il suffit en effet de définir un morphisme
$(X_n)\to(Y_n)$ de tels systèmes comme un \emph{système inductif de
$S_n$-morphismes $u_n:X_n\to Y_n$} tel que $u_m$ s'identifie à
$(u_n)_{(S_m)}$ pour $m\leq n$. Cela étant~:
\end{env}

\begin{theorem}[10.12.3]
\label{I.10.12.3-fr}
Il y a équivalence canonique entre la catégorie des
$\mathfrak{S}$-préschémas adiques et la catégorie des
$(S_n)$-systèmes inductifs adiques.
\end{theorem}

L'équivalence en question s'obtient de la façon suivante~: si
$\mathfrak{X}$ est un $\mathfrak{S}$-préschéma adique,
$f:\mathfrak{X}\to\mathfrak{S}$ le morphisme structural,
$\mathscr{K}=f^*(\mathscr{J})\mathscr{O}_{\mathfrak{X}}$ est un Idéal de
définition de $\mathfrak{X}$ et on fait correspondre à $\mathfrak{X}$ le
système inductif des
$X_n=(\mathfrak{X},\mathscr{O}_{\mathfrak{X}}/\mathscr{K}^{n+1})$, le
morphisme structural $f_n:X_n\to S_n$ correspondant à $f$
(\hyperref[I.10.5.6-fr]{10.5.6}). Montrons d'abord que $(X_n)$ est un
\emph{système inductif adique}~: si $f=(\psi,\theta)$, on a
$\psi^*(\mathscr{J})\mathscr{O}_{\mathfrak{X}}=\mathscr{K}$, donc
$\psi^*(\mathscr{J}^n)\mathscr{O}_{\mathfrak{X}}=\mathscr{K}^n$ pour tout
$n$, et (par l'exactitude du foncteur $\psi^*$)
$\mathscr{K}^{m+1}/\mathscr{K}^{n+1}
=\psi^*(\mathscr{J}^{m+1}/\mathscr{J}^{n+1})
(\mathscr{O}_{\mathfrak{X}}/\mathscr{K}^{n+1})$ pour $m\leq n$~; notre
conclusion résulte donc de (\hyperref[I.4.4.5-fr]{4.4.5}). Il est immédiat
en outre de vérifier qu'à un $\mathfrak{S}$-morphisme
$u:\mathfrak{X}\to\mathfrak{Y}$ de $\mathfrak{S}$-préschémas adiques
correspond (avec des notations évidentes)

\oldpage[I]{207}

un système inductif de $S_n$-morphismes $u_n:X_n\to Y_n$ tel que $u_m$
s'identifie à $(u_n)_{(S_m)}$ pour $m\leq n$.

Le fait qu'on a bien défini ainsi une équivalence résultera de la proposition
plus précise suivante~:

\begin{proposition}[10.12.3.1]
\label{I.10.12.3.1-fr}
Soit $(X_n)$ un système inductif de $S_n$-préschémas~; on suppose que les
morphismes structuraux $f_n:X_n\to S_n$ sont tels que les diagrammes
(\hyperref[I.10.12.2.1-fr]{10.12.2.1}) soient commutatifs et identifient
$X_m$ à $X_n\times_{S_n}S_m$ pour $m\leq n$. Alors le système inductif
$(X_n)$ vérifie les conditions b) et c) de
(\hyperref[I.10.6.3-fr]{10.6.3})~; soient $\mathfrak{X}$ sa limite inductive,
$f:\mathfrak{X}\to\mathfrak{S}$ le morphisme limite inductive du système
inductif $(f_n)$. Alors, si $X_0$ est localement noethérien,
$\mathfrak{X}$ est localement noethérien et $f$ est un morphisme adique.
\end{proposition}

Comme le faisceau d'idéaux de $\mathscr{O}_{S_n}$ qui définit le
sous-préschéma $S_m$ de $S_n$ est nilpotent, il en est de même, en vertu de
(\hyperref[I.4.4.5-fr]{4.4.5}) du faisceau d'idéaux de
$\mathscr{O}_{X_n}$ définissant le sous-préschéma $X_m$ de $X_n$, donc les
conditions de (\hyperref[I.10.6.3-fr]{10.6.3}) sont bien vérifiées. La
question est par suite locale sur $\mathfrak{X}$ et $\mathfrak{S}$ et on
peut supposer que $\mathfrak{S}=\operatorname{Spf}(A)$,
$\mathscr{J}=\mathfrak{J}^\Delta$, $A$ étant un anneau $\mathfrak{J}$-adique
noethérien, et $X_n=\operatorname{Spec}(B_n)$~; si
$A_n=A/\mathfrak{J}^{n+1}$, l'hypothèse entraîne que $B_0$ est noethérien et
que si on pose $\mathfrak{J}_n=\mathfrak{J}/\mathfrak{J}^{n+1}$,
$B_m=B_n/\mathfrak{J}_n^{m+1}B_n$. Le noyau de $B_n\to B_0$ est donc
$\mathfrak{K}_n=\mathfrak{J}_nB_n$ et le noyau de $B_n\to B_m$ est
$\mathfrak{K}_n^{m+1}$ pour $m\leq n$~; en outre, comme $A_1$ est
noethérien, $\mathfrak{J}_1$ est de type fini sur $A_1$, donc
$\overline{\mathfrak{K}}_1=\mathfrak{K}_1/\mathfrak{K}_1^2$ est de type fini
sur $B_1$, et \emph{a fortiori} sur
$B_0=B_1/\mathfrak{K}_1$~; le fait que $\mathfrak{X}$ soit noethérien
résulte alors de (\hyperref[I.10.6.4-fr]{10.6.4})~; si
$B=\varprojlim B_n$, on a $\mathfrak{X}=\operatorname{Spf}(B)$, et si
$\mathfrak{K}$ est le noyau de $B\to B_0$,
$B_n=B/\mathfrak{K}^{n+1}$. Si
$\rho_n:A/\mathfrak{J}^{n+1}\to B/\mathfrak{K}^{n+1}$ est l'homomorphisme
correspondant à $f_n$, on a donc
\[
  \mathfrak{K}/\mathfrak{K}^{n+1}
  =(B/\mathfrak{K}^{n+1})\rho_n(\mathfrak{J}/\mathfrak{J}^{n+1})
\]
comme l'homomorphisme $\rho:A\to B$ correspondant à $f$ est égal à
$\varprojlim\rho_n$, l'idéal $\mathfrak{J}B$ de $B$ est dense dans
$\mathfrak{K}$, et comme tout idéal de $B$ est fermé
(\hyperref[0.7.3.5-fr]{0, 7.3.5}), on a $\mathfrak{K}=\mathfrak{J}B$. Si
$\mathscr{K}=\mathfrak{K}^\Delta$, la relation
$f^*(\mathscr{J})\mathscr{O}_{\mathfrak{X}}=\mathscr{K}$ résulte alors de
(\hyperref[I.10.10.9-fr]{10.10.9}) et achève la démonstration.

\begin{env}[10.12.3.2]
\label{I.10.12.3.2-fr}
L'équivalence précédente fournit, pour deux $\mathfrak{S}$-préschémas
adiques $\mathfrak{X}$, $\mathfrak{Y}$, une \emph{bijection canonique}
\[
  \operatorname{Hom}_{\mathfrak{S}}(\mathfrak{X},\mathfrak{Y})
  \xrightarrow{\sim}
  \varprojlim_n\operatorname{Hom}_{S_n}(X_n,Y_n)
\]
la limite projective étant relative aux applications
$u_n\to(u_n)_{(S_m)}$ pour $m\leq n$.
\end{env}

\subsection{Morphismes de type fini.}
\label{subsection:I.10.13-fr}

\begin{proposition}[10.13.1]
\label{I.10.13.1-fr}
Soient $\mathfrak{Y}$ un préschéma formel localement noethérien,
$\mathscr{K}$ un Idéal de définition de $\mathfrak{Y}$,
$f:\mathfrak{X}\to\mathfrak{Y}$ un morphisme de préschémas formels. Les
conditions suivantes sont équivalentes~:
\begin{enumerate}
  \item[a)] $\mathfrak{X}$ est localement noethérien, $f$ est un morphisme
    adique (\hyperref[I.10.12.1-fr]{10.12.1}) et si l'on pose
    $\mathscr{J}=f^*(\mathscr{K})\mathscr{O}_{\mathfrak{X}}$, le morphisme
    $f_0:(\mathfrak{X},\mathscr{O}_{\mathfrak{X}}/\mathscr{J})
    \to(\mathfrak{Y},\mathscr{O}_{\mathfrak{Y}}/\mathscr{K})$ déduit de
    $f$ est de type fini.
  \item[b)] $\mathfrak{X}$ est localement noethérien, et est limite
    inductive d'un $(Y_n)$-système inductif adique $(X_n)$ tel que le
    morphisme $X_0\to Y_0$ soit de type fini.

\oldpage[I]{208}

  \item[c)] Tout point de $\mathfrak{Y}$ possède un voisinage ouvert formel
    affine noethérien $V$ ayant la propriété suivante~:
    \begin{enumerate}
      \item[(\textbf{Q})] $f^{-1}(V)$ est réunion d'une famille finie
        d'ouverts formels affines noethériens $U_i$ tels que l'anneau adique
        noethérien $\Gamma(U_i,\mathscr{O}_{\mathfrak{X}})$ soit
        topologiquement isomorphe au quotient d'une algèbre de séries
        formelles restreintes (\hyperref[0.7.5.1-fr]{0, 7.5.1}) sur
        $\Gamma(V,\mathscr{O}_{\mathfrak{Y}})$, par un idéal (nécessairement
        fermé).
    \end{enumerate}
\end{enumerate}
\end{proposition}

Il est immédiat que a) entraîne b) en vertu de
(\hyperref[I.10.12.3-fr]{10.12.3}). Pour montrer que b) entraîne c), on
peut, puisque la question est locale sur $\mathfrak{Y}$, supposer que
$\mathfrak{Y}=\operatorname{Spf}(B)$, où $B$ est adique noethérien~; soit
$\mathscr{K}=\mathfrak{K}^\Delta$, $\mathfrak{K}$ étant un idéal de
définition de $B$. Comme par hypothèse, $X_0$ est de type fini sur $Y_0$,
$X_0$ est réunion finie d'ouverts affines $U_i$ tels que l'anneau $A_{i0}$
du schéma affine induit par $X_0$ sur $U_i$ soit une algèbre de type fini sur
l'anneau $B/\mathfrak{K}$ de $Y_0$
(\hyperref[I.6.3.2-fr]{6.3.2}). En vertu de
(\hyperref[I.5.1.9-fr]{5.1.9}), $U_i$ est aussi un ouvert affine dans chacun
des préschémas noethériens $X_n$, et si $A_{in}$ est l'anneau du schéma
affine induit par $X_n$ sur $U_i$, l'hypothèse b) entraîne que pour
$m\leq n$, $A_{im}$ est isomorphe à
$A_{in}/\mathfrak{K}^{m+1}A_{in}$. Par suite, le préschéma formel induit sur
$U_i$ par $\mathfrak{X}$ est isomorphe à $\operatorname{Spf}(A_i)$, où
$A_i=\varprojlim_n A_{in}$ (\hyperref[I.10.6.4-fr]{10.6.4})~; $A_i$ est un
anneau $\mathfrak{K}A_i$-adique, et $A_i/\mathfrak{K}A_i$, isomorphe à
$A_{i0}$, est une algèbre de type fini sur $B/\mathfrak{K}$. On en conclut
(\hyperref[0.7.5.5-fr]{0, 7.5.5}) que $A_i$ est topologiquement isomorphe à
un quotient d'une algèbre de séries formelles restreintes sur $B$ (par un
idéal nécessairement fermé, puisqu'une telle algèbre est noethérienne
(\hyperref[0.7.5.4-fr]{0, 7.5.4})).

Pour démontrer que c) entraîne a), on peut se limiter au cas où
$\mathfrak{X}=\operatorname{Spf}(A)$ est aussi affine, $A$ étant un anneau
adique noethérien, isomorphe au quotient d'une algèbre de séries formelles
restreintes sur $B$ par un idéal fermé. Alors
(\hyperref[0.7.5.5-fr]{0, 7.5.5}), $A/\mathfrak{K}A$ est une algèbre de type
fini sur $B/\mathfrak{K}$, et $\mathfrak{K}A=\mathfrak{J}$ est un idéal de
définition de $A$, donc, en vertu de
(\hyperref[I.10.10.9-fr]{10.10.9}), les conditions de a) sont satisfaites.

On notera que si les conditions de la prop.
(\hyperref[I.10.13.1-fr]{10.13.1}) sont remplies, la propriété a) est
valable pour \emph{tout Idéal de définition $\mathscr{K}$} de
$\mathfrak{Y}$ (en vertu de c)), et par suite, dans la propriété b),
\emph{tous les $f_n$ sont des morphismes de type fini}.

\begin{corollary}[10.13.2]
\label{I.10.13.2-fr}
Si les conditions de (\hyperref[I.10.13.1-fr]{10.13.1}) sont vérifiées,
tout ouvert formel affine noethérien $V$ de $\mathfrak{Y}$ possède la
propriété (\textbf{Q}) et si $\mathfrak{Y}$ est noethérien, il en est de
même de $\mathfrak{X}$.
\end{corollary}

Cela résulte aussitôt de (\hyperref[I.10.13.1-fr]{10.13.1}) et de
(\hyperref[I.6.3.2-fr]{6.3.2}).

\begin{definition}[10.13.3]
\label{I.10.13.3-fr}
Lorsque les propriétés équivalentes a), b), c) de
(\hyperref[I.10.13.1-fr]{10.13.1}) sont vérifiées, on dit que le morphisme
$f$ est \emph{de type fini}, ou que $\mathfrak{X}$ est un
\emph{$\mathfrak{Y}$-préschéma formel de type fini}, ou un
\emph{préschéma formel de type fini au-dessus de $\mathfrak{Y}$}.
\end{definition}

\begin{corollary}[10.13.4]
\label{I.10.13.4-fr}
Soient $\mathfrak{X}=\operatorname{Spf}(A)$,
$\mathfrak{Y}=\operatorname{Spf}(B)$ deux schémas formels affines
noethériens~; pour que $\mathfrak{X}$ soit de type fini sur
$\mathfrak{Y}$, il faut et il suffit que l'anneau adique noethérien $A$ soit
isomorphe au quotient d'une algèbre de séries formelles restreintes sur $B$
par un idéal fermé.
\end{corollary}

En effet, avec les notations de (\hyperref[I.10.13.1-fr]{10.13.1}), si
$\mathfrak{X}$ est de type fini sur $\mathfrak{Y}$,
$A/\mathfrak{K}A$ est alors une $(B/\mathfrak{K})$-algèbre de type fini en
vertu de (\hyperref[I.6.3.3-fr]{6.3.3}) et $\mathfrak{K}A$ est un idéal de
définition de $A$ (\hyperref[I.10.10.9-fr]{10.10.9}). On conclut donc par
(\hyperref[0.7.5.5-fr]{0, 7.5.5}).

\begin{proposition}[10.13.5]
\label{I.10.13.5-fr}
\medskip\noindent
\begin{enumerate}
  \item[\rm{(i)}] Le composé de deux morphismes de préschémas formels qui
    sont de type fini est de type fini.

\oldpage[I]{209}

  \item[\rm{(ii)}] Soient $\mathfrak{X}$, $\mathfrak{S}$,
    $\mathfrak{S}'$ trois préschémas formels localement noethériens (resp.
    noethériens), $f:\mathfrak{X}\to\mathfrak{S}$,
    $g:\mathfrak{S}'\to\mathfrak{S}$ deux morphismes. Si $f$ est de type
    fini, $\mathfrak{X}\times_{\mathfrak{S}}\mathfrak{S}'$ est localement
    noethérien (resp. noethérien) et est de type fini sur $\mathfrak{S}'$.
  \item[\rm{(iii)}] Soient $\mathfrak{S}$ un préschéma formel localement
    noethérien, $\mathfrak{X}'$, $\mathfrak{Y}'$ deux
    $\mathfrak{S}$-préschémas formels localement noethériens tels que
    $\mathfrak{X}'\times_{\mathfrak{S}}\mathfrak{Y}'$ soit localement
    noethérien. Si $\mathfrak{X}$, $\mathfrak{Y}$ sont des
    $\mathfrak{S}$-préschémas formels localement noethériens,
    $f:\mathfrak{X}\to\mathfrak{X}'$,
    $g:\mathfrak{Y}\to\mathfrak{Y}'$ deux $\mathfrak{S}$-morphismes de
    type fini, $\mathfrak{X}\times_{\mathfrak{S}}\mathfrak{Y}$ est
    localement noethérien et $f\times_{\mathfrak{S}}g$ est un
    $\mathfrak{S}$-morphisme de type fini.
\end{enumerate}
\end{proposition}

(iii) se déduit de (i) et (ii) par le raisonnement formel de
(\hyperref[I.3.5.1-fr]{3.5.1}) et il suffit donc de prouver (i) et (ii).

Soient $\mathfrak{X}$, $\mathfrak{Y}$, $\mathfrak{Z}$ trois préschémas
formels localement noethériens, $f:\mathfrak{X}\to\mathfrak{Y}$,
$g:\mathfrak{Y}\to\mathfrak{Z}$ deux morphismes de type fini. Si
$\mathscr{L}$ est un Idéal de définition de $\mathfrak{Z}$,
$\mathscr{K}=g^*(\mathscr{L})\mathscr{O}_{\mathfrak{Y}}$ en est un pour
$\mathfrak{Y}$ et
$\mathscr{J}=f^*(g^*(\mathscr{L}))\mathscr{O}_{\mathfrak{X}}$ en est un
pour $\mathfrak{X}$. Posons
$X_0=(\mathfrak{X},\mathscr{O}_{\mathfrak{X}}/\mathscr{J})$,
$Y_0=(\mathfrak{Y},\mathscr{O}_{\mathfrak{Y}}/\mathscr{K})$,
$Z_0=(\mathfrak{Z},\mathscr{O}_{\mathfrak{Z}}/\mathscr{L})$ et soient
$f_0:X_0\to Y_0$, $g_0:Y_0\to Z_0$ les morphismes correspondant à $f$ et
$g$. Comme par hypothèse $f_0$ et $g_0$ sont de type fini, il en est de même
de $g_0\circ f_0$ (\hyperref[I.6.3.4-fr]{6.3.4}) qui correspond à
$g\circ f$~; donc $g\circ f$ est de type fini par
(\hyperref[I.10.13.1-fr]{10.13.1}).

Sous les conditions de (ii), $\mathfrak{S}$ (resp. $\mathfrak{X}$,
$\mathfrak{S}'$) est limite inductive d'une suite $(S_n)$ (resp. $(X_n)$,
$(S'_n)$) de préschémas localement noethériens et on peut supposer
(\hyperref[I.10.13.1-fr]{10.13.1}) que
$X_m=X_n\times_{S_n}S_m$ pour $m\leq n$. Le préschéma formel
$\mathfrak{X}\times_{\mathfrak{S}}\mathfrak{S}'$ est alors limite inductive
des préschémas $X_n\times_{S_n}S'_n$
(\hyperref[I.10.7.4-fr]{10.7.4}), et on a
\[
  X_m\times_{S_m}S'_m
  =(X_n\times_{S_n}S_m)\times_{S_m}S'_m
  =(X_n\times_{S_n}S'_n)\times_{S'_n}S'_m.
\]
En outre, $X_0\times_{S_0}S'_0$ est localement noethérien puisque $X_0$
est de type fini sur $S_0$ (\hyperref[I.6.3.8-fr]{6.3.8}). On en conclut
d'abord (\hyperref[I.10.12.3.1-fr]{10.12.3.1}) que
$\mathfrak{X}\times_{\mathfrak{S}}\mathfrak{S}'$ est localement
noethérien~; en outre, comme $X_0\times_{S_0}S'_0$ est de type fini sur
$S'_0$ (\hyperref[I.6.3.8-fr]{6.3.8}), il résulte de
(\hyperref[I.10.12.3.1-fr]{10.12.3.1}) et de
(\hyperref[I.10.13.1-fr]{10.13.1}) que
$\mathfrak{X}\times_{\mathfrak{S}}\mathfrak{S}'$ est de type fini sur
$\mathfrak{S}'$, ce qui achève de prouver (ii) (l'assertion relative aux
préschémas noethériens étant conséquence immédiate de
(\hyperref[I.6.3.8-fr]{6.3.8})).

\begin{corollary}[10.13.6]
\label{I.10.13.6-fr}
Sous les hypothèses de (\hyperref[I.10.9.9-fr]{10.9.9}), si $f$ est un
morphisme de type fini, il en est de même de son prolongement $\widehat f$
aux complétés.
\end{corollary}

\subsection{Sous-préschémas fermés des préschémas formels.}
\label{subsection:I.10.14-fr}

\begin{proposition}[10.14.1]
\label{I.10.14.1-fr}
Soient $\mathfrak{X}$ un préschéma formel localement noethérien,
$\mathscr{A}$ un faisceau d'idéaux cohérent de
$\mathscr{O}_{\mathfrak{X}}$. Si $\mathfrak{Y}$ est le support (fermé) de
$\mathscr{O}_{\mathfrak{X}}/\mathscr{A}$, l'espace topologiquement annelé
$(\mathfrak{Y},(\mathscr{O}_{\mathfrak{X}}/\mathscr{A})|\mathfrak{Y})$ est
un préschéma formel localement noethérien, qui est noethérien si
$\mathfrak{X}$ l'est.
\end{proposition}

Notons que $\mathscr{O}_{\mathfrak{X}}/\mathscr{A}$ est cohérent en vertu de
(\hyperref[I.10.10.3-fr]{10.10.3}) et
(\hyperref[0.5.3.4-fr]{0, 5.3.4}), donc son support $\mathfrak{Y}$ est fermé
(\hyperref[0.5.2.2-fr]{0, 5.2.2}). Soit $\mathscr{J}$ un Idéal de définition
de $\mathfrak{X}$, et soit
$X_n=(\mathfrak{X},\mathscr{O}_{\mathfrak{X}}/\mathscr{J}^{n+1})$~; le
faisceau d'anneaux $\mathscr{O}_{\mathfrak{X}}/\mathscr{A}$ est limite
projective des faisceaux
$\mathscr{O}_{\mathfrak{X}}/(\mathscr{A}+\mathscr{J}^{n+1})
=(\mathscr{O}_{\mathfrak{X}}/\mathscr{A})
\otimes_{\mathscr{O}_{\mathfrak{X}}}
(\mathscr{O}_{\mathfrak{X}}/\mathscr{J}^{n+1})$
(\hyperref[I.10.11.3-fr]{10.11.3}), qui ont tous pour support
$\mathfrak{Y}$. Le faisceau
$(\mathscr{A}+\mathscr{J}^{n+1})/\mathscr{J}^{n+1}$ est un
$\mathscr{O}_{\mathfrak{X}}$-Module cohérent, puisque
$\mathscr{J}^{n+1}$ est cohérent, donc
$(\mathscr{A}+\mathscr{J}^{n+1})/\mathscr{J}^{n+1}$ est aussi un
$(\mathscr{O}_{\mathfrak{X}}/\mathscr{J}^{n+1})$-Module cohérent
(\hyperref[0.5.3.10-fr]{0, 5.3.10})~; si $Y_n$ est le sous-préschéma fermé
de $X_n$ défini par ce faisceau d'idéaux, il est immédiat que
$(\mathfrak{Y},(\mathscr{O}_{\mathfrak{X}}/\mathscr{A})|\mathfrak{Y})$

\oldpage[I]{210}
est le préschéma formel limite inductive des $Y_n$, et comme les conditions de
(\hyperref[I.10.6.4-fr]{10.6.4}) sont satisfaites, cela prouve que ce préschéma
formel est localement noethérien, et noethérien si $\mathfrak{X}$ l'est (puisque
alors $Y_0$ l'est en vertu de (\hyperref[I.6.1.4-fr]{6.1.4})).

\begin{definition}[10.14.2]
\label{I.10.14.2-fr}
On appelle \emph{sous-préschéma fermé} d'un préschéma formel $\mathfrak{X}$ tout
préschéma formel
$(\mathfrak{Y},(\mathscr{O}_{\mathfrak{X}}/\mathscr{A})|\mathfrak{Y})$ où
$\mathscr{A}$ est un $\mathscr{O}_{\mathfrak{X}}$-Module cohérent~; on dit que ce
préschéma est le sous-préschéma fermé défini par $\mathscr{A}$.
\end{definition}

Il est clair que la correspondance ainsi définie entre
$\mathscr{O}_{\mathfrak{X}}$-Modules cohérents et sous-préschémas fermés de
$\mathfrak{X}$ est biunivoque.

Le morphisme d'espaces topologiquement annelés
$j=(\psi,\theta):\mathfrak{Y}\to\mathfrak{X}$, où $\psi$ est l'injection
$\mathfrak{Y}\to\mathfrak{X}$ et $\theta^\sharp$ l'homomorphisme canonique
$\mathscr{O}_{\mathfrak{X}}\to\mathscr{O}_{\mathfrak{X}}/\mathscr{A}$, est
évidemment (\hyperref[I.10.4.5-fr]{10.4.5}) un morphisme de préschémas formels,
qu'on appelle l'\emph{injection canonique} de $\mathfrak{Y}$ dans $\mathfrak{X}$.
On notera que si $\mathfrak{X}=\operatorname{Spf}(A)$, où $A$ est adique
noethérien, on a $\mathscr{A}=\mathfrak{a}^\Delta$, où $\mathfrak{a}$ est un
idéal de $A$ (\hyperref[I.10.10.5-fr]{10.10.5}), et il résulte aussitôt de ce
qui précède que l'on a alors
$\mathfrak{Y}=\operatorname{Spf}(A/\mathfrak{a})$ à un isomorphisme près, et que
$j$ correspond (\hyperref[I.10.2.2-fr]{10.2.2}) à l'homomorphisme canonique
$A\to A/\mathfrak{a}$.

On dit qu'un morphisme $f:\mathfrak{Z}\to\mathfrak{X}$ de préschémas formels
localement noethériens est une \emph{immersion fermée} si elle se factorise en
$\mathfrak{Z}\xrightarrow{g}\mathfrak{Y}\xrightarrow{j}\mathfrak{X}$, où $g$
est un isomorphisme de $\mathfrak{Z}$ sur un sous-préschéma fermé
$\mathfrak{Y}$ de $\mathfrak{X}$ et $j$ l'injection canonique. Comme $j$ est un
monomorphisme d'espaces annelés, $g$ et $\mathfrak{Y}$ sont nécessairement
\emph{uniques}.

\begin{proposition}[10.14.3]
\label{I.10.14.3-fr}
Une immersion fermée est un morphisme de type fini.
\end{proposition}

On se ramène aussitôt au cas où $\mathfrak{X}$ est un schéma formel affine
$\operatorname{Spf}(A)$ et $\mathfrak{Y}=\operatorname{Spf}(A/\mathfrak{a})$~;
la proposition résulte de (\hyperref[I.10.13.1-fr]{10.13.1, c)}).

\begin{lemma}[10.14.4]
\label{I.10.14.4-fr}
Soit $f:\mathfrak{Y}\to\mathfrak{X}$ un morphisme de préschémas formels
localement noethériens, et soit $(U_\alpha)$ un recouvrement de
$f(\mathfrak{Y})$ par des ouverts formels affines noethériens de
$\mathfrak{X}$, tels que les $f^{-1}(U_\alpha)$ soient des ouverts formels
affines noethériens de $\mathfrak{Y}$. Pour que $f$ soit une immersion fermée,
il faut et il suffit que $f(\mathfrak{Y})$ soit une partie fermée de
$\mathfrak{X}$ et que, pour tout $\alpha$, la restriction de $f$ à
$f^{-1}(U_\alpha)$ corresponde (\hyperref[I.10.4.6-fr]{10.4.6}) à un
homomorphisme surjectif
$\Gamma(U_\alpha,\mathscr{O}_{\mathfrak{X}})\to
\Gamma(f^{-1}(U_\alpha),\mathscr{O}_{\mathfrak{Y}})$.
\end{lemma}

Les conditions sont évidemment nécessaires. Inversement, si elles sont
remplies, et si on désigne par $\mathfrak{a}_\alpha$ le noyau de
$\Gamma(U_\alpha,\mathscr{O}_{\mathfrak{X}})\to
\Gamma(f^{-1}(U_\alpha),\mathscr{O}_{\mathfrak{Y}})$, on définit un faisceau
cohérent d'idéaux $\mathscr{A}$ de $\mathscr{O}_{\mathfrak{X}}$ en prenant
$\mathscr{A}|U_\alpha=\mathfrak{a}_\alpha^\Delta$, et en prenant $\mathscr{A}$
nul dans le complémentaire de la réunion des $U_\alpha$. En effet, puisque
$f(\mathfrak{Y})$ est fermé et que le support de
$\mathfrak{a}_\alpha^\Delta$ est $U_\alpha\cap f(\mathfrak{Y})$, tout revient
à vérifier que $\mathfrak{a}_\alpha^\Delta$ et
$\mathfrak{a}_\beta^\Delta$ induisent le même faisceau sur un ouvert formel
affine noethérien $V\subset U_\alpha\cap U_\beta$. Or, la restriction de $f$
à $f^{-1}(U_\alpha)$ étant une immersion fermée de ce préschéma formel dans
$U_\alpha$, $f^{-1}(V)$ est un ouvert formel affine noethérien dans
$f^{-1}(U_\alpha)$ et la restriction de $f$ à $f^{-1}(V)$ est une immersion
fermée~; si $\mathfrak{b}$ est le noyau de l'homomorphisme surjectif
$\Gamma(V,\mathscr{O}_{\mathfrak{X}})\to
\Gamma(f^{-1}(V),\mathscr{O}_{\mathfrak{Y}})$ correspondant à cette
restriction, il est immédiat (\hyperref[I.10.10.2-fr]{10.10.2}) que
$\mathfrak{a}_\alpha^\Delta$ induit $\mathfrak{b}^\Delta$ sur $V$. Le faisceau
d'idéaux $\mathscr{A}$ étant ainsi défini, il est alors clair que
$f=g\circ j$, où $j:\mathfrak{Z}\to\mathfrak{X}$ est l'injection canonique du
sous-préschéma fermé $\mathfrak{Z}$ de $\mathfrak{X}$ défini par
$\mathscr{A}$, et $g$ un isomorphisme de $\mathfrak{Y}$ sur $\mathfrak{Z}$.

\begin{proposition}[10.14.5]
\label{I.10.14.5-fr}
\begin{enumerate}
  \item[\rm{(i)}] Si $f:\mathfrak{Z}\to\mathfrak{Y}$,
  $g:\mathfrak{Y}\to\mathfrak{X}$ sont des immersions fermées de préschémas
  formels localement noethériens, $g\circ f$ est une immersion fermée.

\oldpage[I]{211}
  \item[\rm{(ii)}] Soient $\mathfrak{X}$, $\mathfrak{Y}$, $\mathfrak{S}$ trois
  préschémas formels localement noethériens,
  $f:\mathfrak{X}\to\mathfrak{S}$ une immersion fermée,
  $g:\mathfrak{Y}\to\mathfrak{S}$ un morphisme. Alors le morphisme
  $\mathfrak{X}\times_{\mathfrak{S}}\mathfrak{Y}\to\mathfrak{Y}$ est une
  immersion fermée.
  \item[\rm{(iii)}] Soient $\mathfrak{S}$ un préschéma formel localement
  noethérien, $\mathfrak{X}'$, $\mathfrak{Y}'$ deux
  $\mathfrak{S}$-préschémas formels localement noethériens tels que
  $\mathfrak{X}'\times_{\mathfrak{S}}\mathfrak{Y}'$ soit localement
  noethérien. Si $\mathfrak{X}$, $\mathfrak{Y}$ sont des
  $\mathfrak{S}$-préschémas formels localement noethériens,
  $f:\mathfrak{X}\to\mathfrak{X}'$, $g:\mathfrak{Y}\to\mathfrak{Y}'$ deux
  $\mathfrak{S}$-morphismes qui sont des immersions fermées, alors
  $f\times_{\mathfrak{S}}g$ est une immersion fermée.
\end{enumerate}
\end{proposition}

En vertu de (\hyperref[I.3.5.1-fr]{3.5.1}), il suffit encore de prouver (i)
et (ii).

Pour démontrer (i), on peut supposer que $\mathfrak{Y}$ (resp.
$\mathfrak{Z}$) est un sous-préschéma fermé de $\mathfrak{X}$ (resp.
$\mathfrak{Y}$) défini par un faisceau cohérent $\mathscr{J}$ (resp.
$\mathscr{K}$) d'idéaux de $\mathscr{O}_{\mathfrak{X}}$ (resp.
$\mathscr{O}_{\mathfrak{Y}}$)~; si $\psi$ est l'injection
$\mathfrak{Y}\to\mathfrak{X}$ des espaces sous-jacents,
$\psi_*(\mathscr{K})$ est un faisceau cohérent d'idéaux de
$\psi_*(\mathscr{O}_{\mathfrak{Y}})
=\mathscr{O}_{\mathfrak{X}}/\mathscr{J}$
(\hyperref[0.5.3.12-fr]{0, 5.3.12}), donc aussi un
$\mathscr{O}_{\mathfrak{X}}$-Module cohérent
(\hyperref[0.5.3.10-fr]{0, 5.3.10})~; le noyau $\mathscr{K}_1$ de
$\mathscr{O}_{\mathfrak{X}}\to
(\mathscr{O}_{\mathfrak{X}}/\mathscr{J})/\psi_*(\mathscr{K})$ est donc un
faisceau cohérent d'idéaux de $\mathscr{O}_{\mathfrak{X}}$
(\hyperref[0.5.3.4-fr]{0, 5.3.4}), et
$\mathscr{O}_{\mathfrak{X}}/\mathscr{K}_1$ est isomorphe à
$\psi_*(\mathscr{O}_{\mathfrak{Y}}/\mathscr{K})$, ce qui prouve que
$\mathfrak{Z}$ est isomorphe à un sous-préschéma fermé de $\mathfrak{X}$.

Pour démontrer (ii), il est immédiat qu'on peut se limiter au cas où
$\mathfrak{S}=\operatorname{Spf}(A)$,
$\mathfrak{X}=\operatorname{Spf}(B)$,
$\mathfrak{Y}=\operatorname{Spf}(C)$, $A$ étant un anneau
$\mathfrak{J}$-adique noethérien, $B=A/\mathfrak{a}$, où $\mathfrak{a}$ est
un idéal de $A$, $C$ une $A$-algèbre topologique adique et noethérienne. Tout
revient à prouver que l'homomorphisme
$C\to C\widehat{\otimes}_A(A/\mathfrak{a})$ est \emph{surjectif}~: or,
$A/\mathfrak{a}$ est un $A$-module de type fini, et sa topologie est la
topologie $\mathfrak{J}$-adique~; il résulte alors de
(\hyperref[0.7.7.8-fr]{0, 7.7.8}) que
$C\widehat{\otimes}_A(A/\mathfrak{a})$ s'identifie à
$C\otimes_A(A/\mathfrak{a})=C/\mathfrak{a}C$, d'où notre assertion.

\begin{corollary}[10.14.6]
\label{I.10.14.6-fr}
Sous les hypothèses de (\hyperref[I.10.14.5-fr]{10.14.5, (ii)}), soient
$p:\mathfrak{X}\times_{\mathfrak{S}}\mathfrak{Y}\to\mathfrak{X}$,
$q:\mathfrak{X}\times_{\mathfrak{S}}\mathfrak{Y}\to\mathfrak{Y}$ les
projections, de sorte que le diagramme
\[
  \xymatrix{
    \mathfrak{X} \ar[d]_{f}
    & \mathfrak{X}\times_{\mathfrak{S}}\mathfrak{Y}
      \ar[l]_{p} \ar[d]^{q}\\
    \mathfrak{S}
    & \mathfrak{Y} \ar[l]^{g}
  }
\]
est commutatif. Pour tout $\mathscr{O}_{\mathfrak{X}}$-Module cohérent
$\mathscr{F}$, on a alors un isomorphisme canonique de
$\mathscr{O}_{\mathfrak{Y}}$-Modules
\[
  \label{I.10.14.6.1-fr}
  u:g^*(f_*(\mathscr{F}))
  \xrightarrow{\sim}
  q_*(p^*(\mathscr{F})).
  \tag{10.14.6.1}
\]
\end{corollary}

Pour définir un homomorphisme
$g^*(f_*(\mathscr{F}))\to q_*(p^*(\mathscr{F}))$, on sait qu'il revient au
même de définir un homomorphisme
$f_*(\mathscr{F})\to g_*(q_*(p^*(\mathscr{F})))
=f_*(p_*(p^*(\mathscr{F})))$
(\hyperref[0.4.4.3-fr]{0, 4.4.3})~: nous prendrons $u=f_*(\rho)$, où $\rho$
est l'homomorphisme canonique
$\mathscr{F}\to p_*(p^*(\mathscr{F}))$
(\hyperref[0.4.4.3-fr]{0, 4.4.3}). Pour voir que $u$ est un isomorphisme, on
se ramène aussitôt au cas où $\mathfrak{S}$, $\mathfrak{X}$,
$\mathfrak{Y}$ sont des spectres formels d'anneaux adiques noethériens $A$,
$B$, $C$, avec les conditions vues ci-dessus dans
(\hyperref[I.10.14.5-fr]{10.14.5, (ii)})~; on a alors
$\mathscr{F}=M^\Delta$, où $M$ est un $(A/\mathfrak{a})$-module de type fini
(\hyperref[I.10.10.5-fr]{10.10.5}), et les deux membres de
(\hyperref[I.10.14.6.1-fr]{10.14.6.1}) s'identifient respectivement, en vertu
de (\hyperref[I.10.10.8-fr]{10.10.8}), à
$(C\otimes_A M)^\Delta$ et à
$((C/\mathfrak{a}C)\otimes_{A/\mathfrak{a}}M)^\Delta$, d'où le corollaire,
puisque
$(C/\mathfrak{a}C)\otimes_{A/\mathfrak{a}}M
=(C\otimes_A(A/\mathfrak{a}))\otimes_{A/\mathfrak{a}}M$ s'identifie
canoniquement à $C\otimes_A M$.

\oldpage[I]{212}
\begin{corollary}[10.14.7]
\label{I.10.14.7-fr}
Soient $X$ un préschéma usuel localement noethérien, $Y$ un sous-préschéma
fermé de $X$, $j$ l'injection canonique $Y\to X$, $X'$ une partie fermée de
$X$ et $Y'=Y\cap X'$~; alors
$\widehat{j}:Y_{/Y'}\to X_{/X'}$ est une immersion fermée, et pour tout
$\mathscr{O}_Y$-Module cohérent $\mathscr{F}$, on a
\[
  \widehat{j}_*(\mathscr{F}_{/Y'})=(j_*(\mathscr{F}))_{/X'}.
\]
\end{corollary}

Comme $Y'=j^{-1}(X')$, il suffit d'utiliser
(\hyperref[I.10.9.9-fr]{10.9.9}) et d'appliquer
(\hyperref[I.10.14.5-fr]{10.14.5}) et
(\hyperref[I.10.14.6-fr]{10.14.6}).

\subsection{Préschémas formels séparés.}
\label{subsection:I.10.15-fr}

\begin{definition}[10.15.1]
\label{I.10.15.1-fr}
Soient $\mathfrak{S}$ un préschéma formel, $\mathfrak{X}$ un
$\mathfrak{S}$-préschéma formel,
$f:\mathfrak{X}\to\mathfrak{S}$ le morphisme structural. On appelle
morphisme diagonal
$\Delta_{\mathfrak{X}|\mathfrak{S}}:\mathfrak{X}\to
\mathfrak{X}\times_{\mathfrak{S}}\mathfrak{X}$ (noté aussi
$\Delta_{\mathfrak{X}}$) le morphisme
$(1_{\mathfrak{X}},1_{\mathfrak{X}})_{\mathfrak{S}}$. On dit que
$\mathfrak{X}$ est séparé sur $\mathfrak{S}$, ou un
$\mathfrak{S}$-schéma formel, ou que $f$ est un morphisme séparé, si l'image
par $\Delta_{\mathfrak{X}}$ de l'espace sous-jacent $\mathfrak{X}$ est une
partie fermée de l'espace sous-jacent à
$\mathfrak{X}\times_{\mathfrak{S}}\mathfrak{X}$. On dit qu'un préschéma
formel $\mathfrak{X}$ est séparé, ou est un schéma formel, s'il est séparé
sur $\mathbf{Z}$.
\end{definition}

\begin{proposition}[10.15.2]
\label{I.10.15.2-fr}
Supposons que les préschémas formels $\mathfrak{S}$, $\mathfrak{X}$ soient
respectivement limites inductives des suites $(S_n)$, $(X_n)$ de
préschémas usuels, et que le morphisme
$f:\mathfrak{X}\to\mathfrak{S}$ soit limite inductive d'une suite de
morphismes $f_n:X_n\to S_n$. Pour que $f$ soit séparé, il faut et il suffit
que le morphisme $f_0:X_0\to S_0$ le soit.
\end{proposition}

En effet, $\Delta_{\mathfrak{X}|\mathfrak{S}}$ est alors limite inductive de
la suite de morphismes $\Delta_{X_n|S_n}$
(\hyperref[I.10.7.4-fr]{10.7.4}), et l'image par
$\Delta_{\mathfrak{X}|\mathfrak{S}}$ de l'espace sous-jacent
$\mathfrak{X}$ (resp. l'espace sous-jacent
$\mathfrak{X}\times_{\mathfrak{S}}\mathfrak{X}$) est identique à l'image
par $\Delta_{X_0|S_0}$ de l'espace sous-jacent $X_0$ (resp. à l'espace
sous-jacent $X_0\times_{S_0}X_0$)~; d'où la conclusion.

\begin{proposition}[10.15.3]
\label{I.10.15.3-fr}
On suppose dans ce qui suit que tous les préschémas formels (resp.
morphismes de préschémas formels) considérés soient limites inductives de
suites de préschémas usuels (resp. de morphismes de préschémas usuels).
\begin{enumerate}
  \item[\rm{(i)}] Le composé de deux morphismes séparés est séparé.
  \item[\rm{(ii)}] Si $f:\mathfrak{X}\to\mathfrak{X}'$,
  $g:\mathfrak{Y}\to\mathfrak{Y}'$ sont deux $\mathfrak{S}$-morphismes
  séparés, $f\times_{\mathfrak{S}}g$ est séparé.
  \item[\rm{(iii)}] Si $f:\mathfrak{X}\to\mathfrak{Y}$ est un
  $\mathfrak{S}$-morphisme séparé, le $\mathfrak{S}'$-morphisme
  $f_{(\mathfrak{S}')}$ est séparé pour toute extension
  $\mathfrak{S}'\to\mathfrak{S}$ du préschéma formel de base.
  \item[\rm{(iv)}] Si le composé $g\circ f$ de deux morphismes est séparé,
  $f$ est séparé.
\end{enumerate}
\emph{(On sous-entend dans cet énoncé que si un même préschéma formel
$\mathfrak{Z}$ intervient plusieurs fois dans une même proposition, on le
considère comme limite inductive de la \emph{même} suite $(Z_n)$ de
préschémas usuels partout où il figure, et les morphismes de $\mathfrak{Z}$
dans un préschéma formel (resp. d'un préschéma formel dans $\mathfrak{Z}$)
comme limites inductives de morphismes des $Z_n$ dans des préschémas usuels
(resp. de préschémas usuels dans les $Z_n$)).}
\end{proposition}

Avec les notations de (\hyperref[I.10.15.2-fr]{10.15.2}), on a en effet
$(g\circ f)_0=g_0\circ f_0$, et
$(f\times_{\mathfrak{S}}g)_0=f_0\times_{S_0}g_0$~; les assertions de
(\hyperref[I.10.15.3-fr]{10.15.3}) sont alors conséquences immédiates de
(\hyperref[I.10.15.2-fr]{10.15.2}) et des assertions correspondantes de
(\hyperref[I.5.5.1-fr]{5.5.1}) pour les préschémas usuels.

Nous laissons au lecteur le soin d'énoncer pour la même espèce de
préschémas formels et de morphismes que dans
(\hyperref[I.10.15.3-fr]{10.15.3}) les propositions correspondant à
(\hyperref[I.5.5.5-fr]{5.5.5}),

\oldpage[I]{213}
(\hyperref[I.5.5.9-fr]{5.5.9}) et
(\hyperref[I.5.5.10-fr]{5.5.10}) (en y remplaçant «~ouvert affine~» par
«~ouvert formel affine vérifiant la condition b) de
(\hyperref[I.10.6.3-fr]{10.6.3})~»).

Un raisonnement analogue montre aussi que tout schéma formel affine
\emph{noethérien} est séparé, ce qui justifie la terminologie.

\begin{proposition}[10.15.4]
\label{I.10.15.4-fr}
Soient $\mathfrak{S}$ un préschéma formel localement noethérien,
$\mathfrak{X}$, $\mathfrak{Y}$ deux $\mathfrak{S}$-préschémas formels
localement noethériens, tels que $\mathfrak{X}$ ou $\mathfrak{Y}$ soit de
type fini sur $\mathfrak{S}$ (de sorte que
$\mathfrak{X}\times_{\mathfrak{S}}\mathfrak{Y}$ est localement noethérien)
et que $\mathfrak{Y}$ soit séparé sur $\mathfrak{S}$. Soit
$f:\mathfrak{X}\to\mathfrak{Y}$ un $\mathfrak{S}$-morphisme~; alors le
morphisme graphe
$\Gamma_f=(1_{\mathfrak{X}},f)_{\mathfrak{S}}:\mathfrak{X}\to
\mathfrak{X}\times_{\mathfrak{S}}\mathfrak{Y}$ est une immersion fermée.
\end{proposition}

On peut supposer que $\mathfrak{S}$ est limite inductive d'une suite
$(S_n)$ de préschémas localement noethériens, $\mathfrak{X}$ (resp.
$\mathfrak{Y}$) limite inductive d'une suite $(X_n)$ (resp. $(Y_n)$) de
$S_n$-préschémas, $f$ limite inductive d'une suite de $S_n$-morphismes
$f_n:X_n\to Y_n$~; alors
$\mathfrak{X}\times_{\mathfrak{S}}\mathfrak{Y}$ est limite inductive de la
suite $(X_n\times_{S_n}Y_n)$ et $\Gamma_f$ de la suite
$(\Gamma_{f_n})$ (\hyperref[I.10.7.4-fr]{10.7.4})~; par hypothèse, $Y_0$ est
séparé sur $S_0$ (\hyperref[I.10.15.2-fr]{10.15.2}), donc l'espace
$\Gamma_{f_0}(X_0)$ est un sous-espace fermé de
$X_0\times_{S_0}Y_0$~; comme les espaces sous-jacents de
$\mathfrak{X}\times_{\mathfrak{S}}\mathfrak{Y}$ (resp.
$\Gamma_f(\mathfrak{X})$) et $X_0\times_{S_0}Y_0$ (resp.
$\Gamma_{f_0}(X_0)$) sont les mêmes, on voit déjà que
$\Gamma_f(\mathfrak{X})$ est un sous-espace \emph{fermé} de
$\mathfrak{X}\times_{\mathfrak{S}}\mathfrak{Y}$. Remarquons maintenant que
lorsque $(U,V)$ parcourt l'ensemble des couples formés d'un ouvert formel
affine noethérien $U$ (resp. $V$) de $\mathfrak{X}$ (resp.
$\mathfrak{Y}$) tels que $f(U)\subset V$, les ouverts
$U\times_{\mathfrak{S}}V$ forment un recouvrement de
$\Gamma_f(\mathfrak{X})$ dans
$\mathfrak{X}\times_{\mathfrak{S}}\mathfrak{Y}$, et si
$f_U:U\to V$ est la restriction de $f$ à $U$,
$\Gamma_{f_U}:U\to U\times_{\mathfrak{S}}V$ est la restriction de
$\Gamma_f$ à $U$. Si nous montrons que $\Gamma_{f_U}$ est une immersion
fermée, il en sera donc de même de $\Gamma_f$
(\hyperref[I.10.14.4-fr]{10.14.4}), autrement dit, on est ramené au cas où
$\mathfrak{S}=\operatorname{Spf}(A)$,
$\mathfrak{X}=\operatorname{Spf}(B)$,
$\mathfrak{Y}=\operatorname{Spf}(C)$ sont affines ($A$, $B$, $C$ adiques
noethériens), $f$ correspondant à un $A$-homomorphisme continu
$\varphi:C\to B$~; $\Gamma_f$ correspond alors à l'unique homomorphisme
continu $\omega:B\widehat{\otimes}_A C\to B$ qui, composé avec les
homomorphismes canoniques $B\to B\widehat{\otimes}_A C$ et
$C\to B\widehat{\otimes}_A C$, donne respectivement l'identité et
$\varphi$. Or, il est clair que $\omega$ est \emph{surjectif}, d'où notre
assertion.

\begin{corollary}[10.15.5]
\label{I.10.15.5-fr}
Soient $\mathfrak{S}$ un préschéma formel localement noethérien,
$\mathfrak{X}$ un $\mathfrak{S}$-préschéma de type fini~; pour que
$\mathfrak{X}$ soit séparé sur $\mathfrak{S}$, il faut et il suffit que le
morphisme diagonal
$\mathfrak{X}\to\mathfrak{X}\times_{\mathfrak{S}}\mathfrak{X}$ soit une
immersion fermée.
\end{corollary}

\begin{proposition}[10.15.6]
\label{I.10.15.6-fr}
Une immersion fermée $j:\mathfrak{Y}\to\mathfrak{X}$ de préschémas formels
localement noethériens est un morphisme séparé.
\end{proposition}

Avec les notations de (\hyperref[I.10.14.2-fr]{10.14.2}),
$j_0:Y_0\to X_0$ est une immersion fermée, donc un morphisme séparé, et il
suffit d'appliquer (\hyperref[I.10.15.2-fr]{10.15.2}).

\begin{proposition}[10.15.7]
\label{I.10.15.7-fr}
Soient $X$ un préschéma (usuel) localement noethérien, $X'$ une partie
fermée de $X$ et $\widehat{X}=X_{/X'}$. Pour que $\widehat{X}$ soit séparé,
il faut et il suffit que $\widehat{X}_{\mathrm{red}}$ le soit, et il suffit
que $X$ le soit.
\end{proposition}

En effet, avec les notations de (\hyperref[I.10.8.5-fr]{10.8.5}), pour que
$\widehat{X}$ soit séparé, il faut et il suffit que $X'_0$ le soit
(\hyperref[I.10.15.2-fr]{10.15.2}), et comme
$\widehat{X}_{\mathrm{red}}=(X'_0)_{\mathrm{red}}$, il est équivalent de
dire que $\widehat{X}_{\mathrm{red}}$ l'est
(\hyperref[I.5.5.1-fr]{5.5.1, (vi)}).

% FIN SOURCE ega1/ega1-10-fr.tex

%
% DEBUT SOURCE ega1/bib-fr.tex
\oldpage[I]{214}
\begin{thebibliography}{22}

\bibitem{I-1}
H.~\textsc{Cartan} et C.~\textsc{Chevalley},
\emph{Séminaire de l'École Normale Supérieure},
8\textsuperscript{e} année (1955-56), \emph{Géométrie algébrique}.

\bibitem{I-2}
H.~\textsc{Cartan} and S.~\textsc{Eilenberg},
\emph{Homological Algebra}, Princeton Math. Series (Princeton University
Press), 1956.

\bibitem{I-3}
W.~L.~\textsc{Chow} and J.~\textsc{Igusa},
\emph{Cohomology theory of varieties over rings},
\emph{Proc. Nat. Acad. Sci. U.S.A.}, t.~XLIV (1958), p.~1244-1248.

\bibitem{I-4}
R.~\textsc{Godement}, \emph{Théorie des faisceaux}, Actual. Scient. et Ind.,
n\textsuperscript{o}~1252, Paris (Hermann), 1958.

\bibitem{I-5}
H.~\textsc{Grauert},
\emph{Ein Theorem der analytischen Garbentheorie und die Modulräume komplexer
Strukturen}, \emph{Publ. Math. Inst. Hautes Études Scient.},
n\textsuperscript{o}~5, 1960.

\bibitem{I-6}
A.~\textsc{Grothendieck}, \emph{Sur quelques points d'algèbre homologique},
\emph{Tôhoku Math. Journ.}, t.~IX (1957), p.~119-221.

\bibitem{I-7}
A.~\textsc{Grothendieck},
\emph{Cohomology theory of abstract algebraic varieties},
\emph{Proc. Intern. Congress of Math.}, p.~103-118, Edinburgh (1958).

\bibitem{I-8}
A.~\textsc{Grothendieck},
\emph{Géométrie formelle et géométrie algébrique},
\emph{Séminaire Bourbaki}, 11\textsuperscript{e} année (1958-59), exposé 182.

\bibitem{I-9}
M.~\textsc{Nagata},
\emph{A general theory of algebraic geometry over Dedekind domains},
\emph{Amer. Math. Journ.}~: I, t.~LXXVIII, p.~78-116 (1956)~; II, t.~LXXX,
p.~382-420 (1958).

\bibitem{I-10}
D.~G.~\textsc{Northcott}, \emph{Ideal theory}, Cambridge Univ. Press, 1953.

\bibitem{I-11}
P.~\textsc{Samuel}, \emph{Commutative algebra} (Notes by D.~Herzig), Cornell
Univ., 1953.

\bibitem{I-12}
P.~\textsc{Samuel}, \emph{Algèbre locale}, Mém. Sci. Math.,
n\textsuperscript{o}~123, Paris, 1953.

\bibitem{I-13}
P.~\textsc{Samuel} and O.~\textsc{Zariski},
\emph{Commutative algebra}, 2 vol., New York (Van Nostrand), 1958-60.

\bibitem{I-14}
J.-P.~\textsc{Serre}, \emph{Faisceaux algébriques cohérents},
\emph{Ann. of Math.}, t.~LXI (1955), p.~197-278.

\bibitem{I-15}
J.-P.~\textsc{Serre},
\emph{Sur la cohomologie des variétés algébriques},
\emph{Journ. de Math.} (9), t.~XXXVI (1957), p.~1-16.

\bibitem{I-16}
J.-P.~\textsc{Serre},
\emph{Géométrie algébrique et géométrie analytique},
\emph{Ann. Inst. Fourier}, t.~VI (1955-56), p.~1-42.

\bibitem{I-17}
J.-P.~\textsc{Serre},
\emph{Sur la dimension homologique des anneaux et des modules noethériens},
\emph{Proc. Intern. Symp. on Alg. Number theory}, p.~176-189, Tokyo-Nikko,
1955.

\bibitem{I-18}
A.~\textsc{Weil}, \emph{Foundations of algebraic geometry}, Amer. Math. Soc.
Coll. Publ., n\textsuperscript{o}~29, 1946.

\bibitem{I-19}
A.~\textsc{Weil}, \emph{Numbers of solutions of equations in finite fields},
\emph{Bull. Amer. Math. Soc.}, t.~LV (1949), p.~497-508.

\bibitem{I-20}
O.~\textsc{Zariski},
\emph{Theory and applications of holomorphic functions on algebraic varieties
over arbitrary ground fields}, Mem. Amer. Math. Soc., n\textsuperscript{o}~5
(1951).

\bibitem{I-21}
O.~\textsc{Zariski}, \emph{A new proof of Hilbert's Nullstellensatz},
\emph{Bull. Amer. Math. Soc.}, t.~LIII (1947), p.~362-368.

\bibitem{I-22}
E.~\textsc{Kähler}, \emph{Geometria Arithmetica},
\emph{Ann. di Mat.} (4), t.~XLV (1958), p.~1-368.

\end{thebibliography}

% FIN SOURCE ega1/bib-fr.tex

%
% DEBUT SOURCE ega2/front-fr.tex
\clearpage
\thispagestyle{plain}
\markboth{CHAPITRE II}{CHAPITRE II}
\oldpage[II]{5}
\phantomsection
\label{chapter:II-fr}

\begin{center}
{\large CHAPITRE II\par}
\vspace{1.5em}
{\Large\bfseries ÉTUDE GLOBALE ÉLÉMENTAIRE\par}
{\Large\bfseries DE QUELQUES CLASSES DE MORPHISMES\par}
\vspace{1.8em}
{\bfseries Sommaire\par}
\end{center}

\medskip
\begin{tabular}{@{}r@{\ }p{0.82\textwidth}@{}}
\S~1. & Morphismes affines.\\
\S~2. & Spectres premiers homogènes.\\
\S~3. & Spectre homogène d'un faisceau d'algèbres graduées.\\
\S~4. & Fibrés projectifs. Faisceaux amples.\\
\S~5. & Morphismes quasi-affines~; morphismes quasi-projectifs~; morphismes
         propres~; morphismes projectifs.\\
\S~6. & Morphismes entiers et morphismes finis.\\
\S~7. & Critères valuatifs.\\
\S~8. & Schémas éclatés~; cônes projetants~; fermeture projective.
\end{tabular}

\medskip
Les diverses classes de morphismes étudiées dans ce chapitre le sont sans
faire grand usage des méthodes cohomologiques~; une étude plus poussée,
utilisant ces dernières méthodes, en sera faite au chapitre~III, où on
utilisera surtout les \S\S~2, 4 et~5 du chapitre~II. Le \S~8 peut être omis en
première lecture~: il donne quelques compléments au formalisme développé dans
les \S\S~1 à~3, se réduisant à des applications faciles de ce formalisme, et
nous en ferons un usage moins constant que des autres résultats de ce
chapitre.

% FIN SOURCE ega2/front-fr.tex

%
% DEBUT SOURCE ega2/ega2-1-fr.tex
\setcounter{section}{0}
\section{Morphismes affines}
\label{section:II.1-fr}

La plupart des résultats de ce paragraphe sont les contreparties
<<~globales~>> de ceux du chapitre~I\textsuperscript{er}, \S~1~; ils ne sont
donc pas essentiellement nouveaux et fournissent simplement un langage
commode pour la suite.

\setcounter{subsection}{0}
\subsection{$S$-préschémas et $\mathcal{O}_S$-Algèbres}
\label{subsection:II.1.1-fr}

\begin{env}[1.1.1]
\label{II.1.1.1-fr}
Soient $S$ un préschéma, $X$ un $S$-préschéma, $f:X\to S$ son morphisme
structural. On sait (\hyperref[0.4.2.4-fr]{0, 4.2.4}) que l'image directe
$f_*(\mathcal{O}_X)$ est une $\mathcal{O}_S$-Algèbre, que nous
\oldpage[II]{6}
noterons $\mathcal{A}(X)$ lorsqu'il n'en résultera pas de confusion~; si $U$
est un ouvert de $S$, on a
\[
  \mathcal{A}(f^{-1}(U))=\mathcal{A}(X)|U.
\]
De même, pour tout $\mathcal{O}_X$-Module $\mathcal{F}$ (resp. toute
$\mathcal{O}_X$-Algèbre $\mathcal{B}$), nous écrirons
$\mathcal{A}(\mathcal{F})$ (resp. $\mathcal{A}(\mathcal{B})$) l'image directe
$f_*(\mathcal{F})$ (resp. $f_*(\mathcal{B})$) qui est un
$\mathcal{A}(X)$-Module (resp. une $\mathcal{A}(X)$-Algèbre), et non seulement
un $\mathcal{O}_S$-Module (resp. $\mathcal{O}_S$-Algèbre).
\end{env}

\begin{env}[1.1.2]
\label{II.1.1.2-fr}
Soient $Y$ un second $S$-préschéma, $g:Y\to S$ son morphisme structural,
$h:X\to Y$ un $S$-morphisme~; on a donc le diagramme commutatif
\[
  \xymatrix{
    X\ar[rr]^h\ar[rd]_f && Y\ar[ld]^g\\
    & S
  }
  \tag{1.1.2.1}\label{II.1.1.2.1-fr}
\]

On a par définition $h=(\psi,\theta)$, où
$\theta:\mathcal{O}_Y\to h_*(\mathcal{O}_X)=\psi_*(\mathcal{O}_X)$ est un
homomorphisme de faisceaux d'anneaux~; on en déduit
(\hyperref[0.4.2.2-fr]{0, 4.2.2}) un homomorphisme de
$\mathcal{O}_S$-Algèbres
\[
  g_*(\theta):g_*(\mathcal{O}_Y)\to
  g_*(h_*(\mathcal{O}_X))=f_*(\mathcal{O}_X),
\]
autrement dit, un homomorphisme de $\mathcal{O}_S$-Algèbres
$\mathcal{A}(Y)\to\mathcal{A}(X)$, que nous noterons aussi
$\mathcal{A}(h)$. Si $h':Y\to Z$ est un second $S$-morphisme, il est immédiat
que $\mathcal{A}(h'\circ h)=\mathcal{A}(h)\circ\mathcal{A}(h')$. Nous avons
donc défini un \emph{foncteur contravariant} $\mathcal{A}(X)$ sur la catégorie
des $S$-préschémas, à valeurs dans la catégorie des
$\mathcal{O}_S$-Algèbres.

Soient maintenant $\mathcal{F}$ un $\mathcal{O}_X$-Module, $\mathcal{G}$ un
$\mathcal{O}_Y$-Module, et $u:\mathcal{G}\to\mathcal{F}$ un $h$-morphisme,
c'est-à-dire (\hyperref[0.4.4.1-fr]{0, 4.4.1}) un homomorphisme de
$\mathcal{O}_Y$-Modules $\mathcal{G}\to h_*(\mathcal{F})$. Alors
\[
  g_*(u):g_*(\mathcal{G})\to g_*(h_*(\mathcal{F}))=f_*(\mathcal{F})
\]
est un homomorphisme $\mathcal{A}(\mathcal{G})\to
\mathcal{A}(\mathcal{F})$ de $\mathcal{O}_S$-Modules, que nous noterons
$\mathcal{A}(u)$~; en outre, le couple
$(\mathcal{A}(h),\mathcal{A}(u))$ constitue un \emph{di-homomorphisme} du
$\mathcal{A}(Y)$-Module $\mathcal{A}(\mathcal{G})$ dans le
$\mathcal{A}(X)$-Module $\mathcal{A}(\mathcal{F})$.
\end{env}

\begin{env}[1.1.3]
\label{II.1.1.3-fr}
Le préschéma $S$ étant fixé, on peut considérer les couples
$(X,\mathcal{F})$, où $X$ est un $S$-préschéma et $\mathcal{F}$ un
$\mathcal{O}_X$-Module, comme formant une \emph{catégorie}, en définissant un
\emph{morphisme} $(X,\mathcal{F})\to(Y,\mathcal{G})$ comme un couple $(h,u)$,
où $h:X\to Y$ est un $S$-morphisme et $u:\mathcal{G}\to\mathcal{F}$ un
$h$-morphisme. On peut alors dire que
$(\mathcal{A}(X),\mathcal{A}(\mathcal{F}))$ est un \emph{foncteur
contravariant} à valeurs dans la catégorie dont les objets sont les couples
formés d'une $\mathcal{O}_S$-Algèbre et d'un Module sur cette Algèbre, et les
morphismes sont les di-homomorphismes.
\end{env}

\subsection{Préschémas affines sur un préschéma}
\label{subsection:II.1.2-fr}

\begin{definition}[1.2.1]
\label{II.1.2.1-fr}
Soient $X$ un $S$-préschéma, $f:X\to S$ son morphisme structural. On dit que
$X$ est \emph{affine au-dessus de $S$} s'il existe un recouvrement
$(S_\alpha)$ de $S$ par des ouverts affines tels que, pour tout $\alpha$, le
préschéma induit par $X$ sur l'ensemble ouvert $f^{-1}(S_\alpha)$ soit affine.
\end{definition}

\begin{example}[1.2.2]
\label{II.1.2.2-fr}
Tout sous-préschéma fermé de $S$ est un $S$-préschéma affine au-dessus de $S$
(\hyperref[I.4.2.3-fr]{I, 4.2.3} et
\hyperref[I.4.2.4-fr]{4.2.4}).
\end{example}

\begin{remark}[1.2.3]
\label{II.1.2.3-fr}
Un préschéma $X$ affine au-dessus de $S$ n'est pas nécessairement un schéma
affine, comme le montre l'exemple $X=S$ (\hyperref[II.1.2.2-fr]{1.2.2}).
D'autre part, si un schéma affine $X$ est un $S$-préschéma, $X$ n'est pas
nécessairement affine au-dessus
\oldpage[II]{7}
de $S$ (voir (\hyperref[II.1.3.3-fr]{1.3.3})). Toutefois, rappelons que si
$S$ est un \emph{schéma}, tout $S$-préschéma qui est un schéma affine est
affine au-dessus de $S$ (\hyperref[I.5.5.10-fr]{I, 5.5.10}).
\end{remark}

\begin{proposition}[1.2.4]
\label{II.1.2.4-fr}
Tout $S$-préschéma qui est affine au-dessus de $S$ est séparé au-dessus de
$S$ (autrement dit, est un $S$-schéma).
\end{proposition}

Cela résulte aussitôt de (\hyperref[I.5.5.5-fr]{I, 5.5.5}) et
(\hyperref[I.5.5.8-fr]{I, 5.5.8}).

\begin{proposition}[1.2.5]
\label{II.1.2.5-fr}
Soient $X$ un $S$-schéma affine au-dessus de $S$, $f:X\to S$ le morphisme
structural. Pour tout ouvert $U\subset S$, $f^{-1}(U)$ est affine au-dessus
de $U$.
\end{proposition}

En vertu de la déf. (\hyperref[II.1.2.1-fr]{1.2.1}), on est ramené au cas où
$S=\operatorname{Spec}(A)$ et $X=\operatorname{Spec}(B)$ sont affines et
alors $f=({}^a\varphi,\widetilde{\varphi})$, où $\varphi:A\to B$ est un
homomorphisme. Comme les $D(g)$, où $g\in A$, forment une base de $S$, on est
ramené au cas où $U=D(g)$~; mais on sait alors que
$f^{-1}(U)=D(\varphi(g))$ (\hyperref[I.1.2.2.2-fr]{I, 1.2.2.2}), d'où la
proposition.

\begin{proposition}[1.2.6]
\label{II.1.2.6-fr}
Soient $X$ un $S$-schéma affine au-dessus de $S$, $f:X\to S$ le morphisme
structural. Pour tout $\mathcal{O}_X$-Module quasi-cohérent $\mathcal{F}$,
$f_*(\mathcal{F})$ est un $\mathcal{O}_S$-Module quasi-cohérent.
\end{proposition}

Compte tenu de (\hyperref[II.1.2.4-fr]{1.2.4}), cela résulte de
(\hyperref[I.9.2.2-fr]{I, 9.2.2, a)}).

En particulier, la $\mathcal{O}_S$-Algèbre
$\mathcal{A}(X)=f_*(\mathcal{O}_X)$ est \emph{quasi-cohérente}.

\begin{proposition}[1.2.7]
\label{II.1.2.7-fr}
Soient $X$ un $S$-schéma affine au-dessus de $S$. Pour tout $S$-préschéma
$Y$, l'application $h\mapsto\mathcal{A}(h)$ de l'ensemble
$\operatorname{Hom}_S(Y,X)$ dans l'ensemble
$\operatorname{Hom}(\mathcal{A}(X),\mathcal{A}(Y))$
(\hyperref[II.1.1.2-fr]{1.1.2}) est bijective.
\end{proposition}

Soient $f:X\to S$, $g:Y\to S$ les morphismes structuraux. Supposons d'abord
$S=\operatorname{Spec}(A)$ et $X=\operatorname{Spec}(B)$ affines~; il faut
prouver que pour tout homomorphisme
$\omega:f_*(\mathcal{O}_X)\to g_*(\mathcal{O}_Y)$ de
$\mathcal{O}_S$-Algèbres, il existe un $S$-morphisme $h:Y\to X$ et un seul
tel que $\mathcal{A}(h)=\omega$. Par définition, pour tout ouvert $U\subset
S$, $\omega$ définit un homomorphisme
\[
  \omega_U=\Gamma(U,\omega):
  \Gamma(f^{-1}(U),\mathcal{O}_X)\to
  \Gamma(g^{-1}(U),\mathcal{O}_Y)
\]
de $\Gamma(U,\mathcal{O}_S)$-algèbres. En particulier, pour $U=S$, cela
donne un homomorphisme
$\varphi:\Gamma(X,\mathcal{O}_X)\to\Gamma(Y,\mathcal{O}_Y)$ de
$\Gamma(S,\mathcal{O}_S)$-algèbres, auquel correspond un $S$-morphisme
$h:Y\to X$ bien déterminé, puisque $X$ est affine
(\hyperref[I.2.2.4-fr]{I, 2.2.4}). Reste à prouver que
$\mathcal{A}(h)=\omega$, autrement dit, que pour tout ouvert $U$ d'une base
de $S$, $\omega_U$ coïncide avec l'homomorphisme d'algèbres $\varphi_U$
correspondant au $S$-morphisme $g^{-1}(U)\to f^{-1}(U)$ restriction de $h$.
On peut se borner au cas où $U=D(\lambda)$, avec $\lambda\in S$~; alors, si
$f=({}^a\rho,\widetilde{\rho})$, où $\rho:A\to B$ est un homomorphisme
d'anneaux, on a $f^{-1}(U)=D(\mu)$, où $\mu=\rho(\lambda)$, et
$\Gamma(f^{-1}(U),\mathcal{O}_X)$ est l'anneau de fractions $B_\mu$~; or,
le diagramme
\[
  \xymatrix{
    B\ar[r]^\varphi\ar[d] & \Gamma(Y,\mathcal{O}_Y)\ar[d]\\
    B_\mu\ar[r]_{\varphi_U} & \Gamma(g^{-1}(U),\mathcal{O}_Y)
  }
\]
est commutatif, et il en est de même du diagramme analogue où $\varphi_U$
est remplacé par $\omega_U$~; l'égalité $\varphi_U=\omega_U$ résulte alors
de la propriété universelle des anneaux de fractions
(\hyperref[0.1.2.4-fr]{0, 1.2.4}).

Passons au cas général, et soit $(S_\alpha)$ un recouvrement de $S$ par des
ouverts affines
\oldpage[II]{8}
tels que les $f^{-1}(S_\alpha)$ soient affines. Alors tout homomorphisme
$\omega:\mathcal{A}(X)\to\mathcal{A}(Y)$ de $\mathcal{O}_S$-Algèbres donne
par restriction une famille d'homomorphismes
\[
  \omega_\alpha:
  \mathcal{A}(f^{-1}(S_\alpha))\to
  \mathcal{A}(g^{-1}(S_\alpha))
\]
de $\mathcal{O}_{S_\alpha}$-Algèbres, donc une famille de
$S_\alpha$-morphismes
$h_\alpha:g^{-1}(S_\alpha)\to f^{-1}(S_\alpha)$ d'après ce qui précède.
Tout revient à voir que pour tout ouvert affine $U$ d'une base de
$S_\alpha\cap S_\beta$, les restrictions de $h_\alpha$ et de $h_\beta$ à
$g^{-1}(U)$ coïncident, ce qui est évident, puisque, en vertu de ce qui
précède, ces restrictions correspondent toutes deux à l'homomorphisme
$\mathcal{A}(X)|U\to\mathcal{A}(Y)|U$ restriction de $\omega$.

\begin{corollary}[1.2.8]
\label{II.1.2.8-fr}
Soient $X$, $Y$ deux $S$-schémas qui sont affines au-dessus de $S$. Pour
qu'un $S$-morphisme $h:Y\to X$ soit un isomorphisme, il faut et il suffit que
$\mathcal{A}(h):\mathcal{A}(X)\to\mathcal{A}(Y)$ soit un isomorphisme.
\end{corollary}

Cela résulte aussitôt de (\hyperref[II.1.2.7-fr]{1.2.7}) et du caractère
fonctoriel de $\mathcal{A}(X)$.

\subsection{Préschéma affine au-dessus de $S$ associé à une
$\mathcal{O}_S$-Algèbre}
\label{subsection:II.1.3-fr}

\begin{proposition}[1.3.1]
\label{II.1.3.1-fr}
Soit $S$ un préschéma. Pour toute $\mathcal{O}_S$-Algèbre quasi-cohérente
$\mathcal{B}$, il existe un préschéma $X$ affine au-dessus de $S$, défini à
un $S$-isomorphisme unique près, tel que $\mathcal{A}(X)=\mathcal{B}$.
\end{proposition}

L'unicité résulte de (\hyperref[II.1.2.8-fr]{1.2.8})~; démontrons l'existence
de $X$. Pour tout ouvert affine $U\subset S$, soit $X_U$ le préschéma
$\operatorname{Spec}(\Gamma(U,\mathcal{B}))$~; comme
$\Gamma(U,\mathcal{B})$ est une $\Gamma(U,\mathcal{O}_S)$-algèbre, $X_U$
est un $S$-préschéma (\hyperref[I.1.6.1-fr]{I, 1.6.1}). En outre, comme
$\mathcal{B}$ est quasi-cohérente, la $\mathcal{O}_S$-Algèbre
$\mathcal{A}(X_U)$ s'identifie canoniquement à $\mathcal{B}|U$
(\hyperref[I.1.3.7-fr]{I, 1.3.7},
\hyperref[I.1.3.13-fr]{1.3.13} et
\hyperref[I.1.6.3-fr]{1.6.3}). Soit $V$ un second ouvert affine de $S$, et
soit $X_{U,V}$ le préschéma induit par $X_U$ sur $f_U^{-1}(U\cap V)$, en
désignant par $f_U$ le morphisme structural $X_U\to S$~; $X_{U,V}$ et
$X_{V,U}$ sont affines sur $U\cap V$
(\hyperref[II.1.2.5-fr]{1.2.5}), et par définition $\mathcal{A}(X_{U,V})$
et $\mathcal{A}(X_{V,U})$ s'identifient canoniquement à
$\mathcal{B}|(U\cap V)$. Il y a donc
(\hyperref[II.1.2.8-fr]{1.2.8}) un $S$-isomorphisme canonique
$\theta_{U,V}:X_{V,U}\to X_{U,V}$~; en outre, si $W$ est un troisième
ouvert affine de $S$, et si $\theta'_{U,V}$, $\theta'_{V,W}$,
$\theta'_{U,W}$ sont les restrictions de $\theta_{U,V}$, $\theta_{V,W}$,
$\theta_{U,W}$ aux images réciproques de $U\cap V\cap W$ dans $X_V$,
$X_W$ et $X_W$ respectivement par les morphismes structuraux, on a
$\theta'_{U,V}\circ\theta'_{V,W}=\theta'_{U,W}$. Il existe par suite un
préschéma $X$, un recouvrement $(T_U)$ de $X$ par des ouverts affines, et
pour chaque $U$ un isomorphisme $\varphi_U:X_U\to T_U$, de sorte que
$\varphi_U$ applique $f_U^{-1}(U\cap V)$ sur $T_U\cap T_V$ et que l'on ait
$\theta_{U,V}=\varphi_U^{-1}\circ\varphi_V$
(\hyperref[I.2.3.1-fr]{I, 2.3.1}). Le morphisme
$g_U=f_U\circ\varphi_U^{-1}$ fait de $T_U$ un $S$-préschéma, et les
morphismes $g_U$ et $g_V$ coïncident dans $T_U\cap T_V$, donc $X$ est un
$S$-préschéma. En outre, il est clair par définition que $X$ est affine
au-dessus de $S$ et que $\mathcal{A}(T_U)=\mathcal{B}|U$, donc
$\mathcal{A}(X)=\mathcal{B}$.

On dit que le $S$-schéma $X$ ainsi défini est \emph{associé à la
$\mathcal{O}_S$-Algèbre $\mathcal{B}$} ou est le \emph{spectre de
$\mathcal{B}$}, et on le note $\operatorname{Spec}(\mathcal{B})$.

\begin{corollary}[1.3.2]
\label{II.1.3.2-fr}
Soient $X$ un préschéma affine au-dessus de $S$, $f:X\to S$ le morphisme
structural. Pour tout ouvert affine $U\subset S$, le préschéma induit sur
$f^{-1}(U)$ est un schéma affine d'anneau
$\Gamma(U,\mathcal{A}(X))$.
\end{corollary}

\oldpage[II]{9}
Comme on peut supposer $X$ associé à une $\mathcal{O}_S$-Algèbre en vertu de
(\hyperref[II.1.2.6-fr]{1.2.6}) et (\hyperref[II.1.3.1-fr]{1.3.1}), le
corollaire résulte de la construction de $X$ décrite dans
(\hyperref[II.1.3.1-fr]{1.3.1}).

\begin{example}[1.3.3]
\label{II.1.3.3-fr}
Soit $S$ le plan affine sur un corps $K$, où le point $0$ a été dédoublé
(\hyperref[I.5.5.11-fr]{I, 5.5.11})~; avec les notations de
(\hyperref[I.5.5.11-fr]{I, 5.5.11}), $S$ est réunion de deux ouverts affines
$Y_1$, $Y_2$~; si $f$ est l'immersion ouverte $Y_1\to S$,
$f^{-1}(Y_2)=Y_1\cap Y_2$ n'est pas un ouvert affine dans $Y_1$
(\emph{loc. cit.}), donc on a un exemple d'un schéma affine qui n'est pas
affine au-dessus de $S$.
\end{example}

\begin{corollary}[1.3.4]
\label{II.1.3.4-fr}
Soit $S$ un schéma affine~; pour qu'un $S$-préschéma $X$ soit affine au-dessus
de $S$, il faut et il suffit que $X$ soit un schéma affine.
\end{corollary}

\begin{corollary}[1.3.5]
\label{II.1.3.5-fr}
Soient $X$ un préschéma affine au-dessus d'un préschéma $S$, $Y$ un
$X$-préschéma. Pour que $Y$ soit affine au-dessus de $S$, il faut et il
suffit que $Y$ soit affine au-dessus de $X$.
\end{corollary}

On est aussitôt ramené au cas où $S$ est un schéma affine, et il en est alors
de même de $X$ (\hyperref[II.1.3.4-fr]{1.3.4})~; les deux conditions de
l'énoncé expriment alors que $Y$ est un schéma affine
(\hyperref[II.1.3.4-fr]{1.3.4}).

\begin{env}[1.3.6]
\label{II.1.3.6-fr}
Soit $X$ un préschéma affine au-dessus de $S$. Pour définir un préschéma $Y$
affine \emph{au-dessus de $X$}, il revient au même, en vertu de
(\hyperref[II.1.3.5-fr]{1.3.5}), de se donner un préschéma $Y$ affine
\emph{au-dessus de $S$}, et un $S$-morphisme $g:Y\to X$~; en d'autres termes
(\hyperref[II.1.3.1-fr]{1.3.1} et
\hyperref[II.1.2.7-fr]{1.2.7}), cela revient à se donner une
$\mathcal{O}_S$-Algèbre quasi-cohérente $\mathcal{B}$ et un homomorphisme
$\mathcal{A}(X)\to\mathcal{B}$ de $\mathcal{O}_S$-Algèbres (que l'on peut
envisager comme définissant sur $\mathcal{B}$ une structure de
$\mathcal{A}(X)$-Algèbre). Si $f:X\to S$ est le morphisme structural, on a
alors $\mathcal{B}=f_*(g_*(\mathcal{O}_Y))$.
\end{env}

\begin{corollary}[1.3.7]
\label{II.1.3.7-fr}
Soit $X$ un préschéma affine au-dessus de $S$~; pour que $X$ soit de type fini
sur $S$, il faut et il suffit que la $\mathcal{O}_S$-Algèbre quasi-cohérente
$\mathcal{A}(X)$ soit de type fini
(\hyperref[I.9.6.2-fr]{I, 9.6.2}).
\end{corollary}

Par définition (\hyperref[I.9.6.2-fr]{I, 9.6.2}), on est ramené au cas où
$S$ est affine~; alors $X$ est un schéma affine
(\hyperref[II.1.3.4-fr]{1.3.4}), et si $S=\operatorname{Spec}(A)$,
$X=\operatorname{Spec}(B)$, $\mathcal{A}(X)$ est la
$\mathcal{O}_S$-Algèbre $\widetilde{B}$~; comme
$\Gamma(U,\widetilde{B})=B$, le corollaire résulte de
(\hyperref[I.9.6.2-fr]{I, 9.6.2}) et
(\hyperref[I.6.3.3-fr]{I, 6.3.3}).

\begin{corollary}[1.3.8]
\label{II.1.3.8-fr}
Soit $X$ un préschéma affine au-dessus de $S$~; pour que $X$ soit réduit il
faut et il suffit que la $\mathcal{O}_S$-Algèbre quasi-cohérente
$\mathcal{A}(X)$ soit réduite
(\hyperref[0.4.1.4-fr]{0, 4.1.4}).
\end{corollary}

En effet, la question est évidemment locale sur $S$~; en vertu de
(\hyperref[II.1.3.2-fr]{1.3.2}), le corollaire résulte de
(\hyperref[I.5.1.1-fr]{I, 5.1.1}) et
(\hyperref[I.5.1.4-fr]{I, 5.1.4}).

\subsection{Faisceaux quasi-cohérents sur un préschéma affine au-dessus de
$S$}
\label{subsection:II.1.4-fr}

\begin{proposition}[1.4.1]
\label{II.1.4.1-fr}
Soient $X$ un préschéma affine au-dessus de $S$, $Y$ un $S$-préschéma,
$\mathcal{F}$ (resp. $\mathcal{G}$) un $\mathcal{O}_X$-Module (resp. un
$\mathcal{O}_Y$-Module) quasi-cohérent. Alors l'application
$(h,u)\mapsto(\mathcal{A}(h),\mathcal{A}(u))$ de l'ensemble des morphismes
$(Y,\mathcal{G})\to(X,\mathcal{F})$ dans l'ensemble des di-homomorphismes
$(\mathcal{A}(X),\mathcal{A}(\mathcal{F}))\to
(\mathcal{A}(Y),\mathcal{A}(\mathcal{G}))$
(\hyperref[II.1.1.2-fr]{1.1.2} et
\hyperref[II.1.1.3-fr]{1.1.3}) est bijective.
\end{proposition}

La démonstration suit exactement la même marche que celle de
(\hyperref[II.1.2.7-fr]{1.2.7}) en utilisant
(\hyperref[I.2.2.5-fr]{I, 2.2.5}) et
(\hyperref[I.2.2.4-fr]{I, 2.2.4}), et les détails en sont laissés au
lecteur.

\begin{corollary}[1.4.2]
\label{II.1.4.2-fr}
Si, en outre des hypothèses de (\hyperref[II.1.4.1-fr]{1.4.1}), on suppose
$Y$ affine au-dessus de $S$, alors, pour que $(h,u)$ soit un isomorphisme,
il faut et il suffit que $(\mathcal{A}(h),\mathcal{A}(u))$ soit un
di-isomorphisme.
\end{corollary}

\oldpage[II]{10}
\begin{proposition}[1.4.3]
\label{II.1.4.3-fr}
Pour tout couple $(\mathcal{B},\mathcal{M})$ formé d'une
$\mathcal{O}_S$-Algèbre quasi-cohérente $\mathcal{B}$ et d'un
$\mathcal{B}$-Module $\mathcal{M}$ quasi-cohérent (en tant que
$\mathcal{O}_S$-Module ou en tant que $\mathcal{B}$-Module, ce qui revient au
même (\hyperref[I.9.6.1-fr]{I, 9.6.1})), il existe un couple
$(X,\mathcal{F})$ formé d'un préschéma $X$ affine au-dessus de $S$ et d'un
$\mathcal{O}_X$-Module quasi-cohérent $\mathcal{F}$, tels que
$\mathcal{A}(X)=\mathcal{B}$ et $\mathcal{A}(\mathcal{F})=\mathcal{M}$~; en
outre ce couple est déterminé à un isomorphisme unique près.
\end{proposition}

L'unicité résulte de (\hyperref[II.1.4.1-fr]{1.4.1}) et
(\hyperref[II.1.4.2-fr]{1.4.2})~; l'existence se démontre comme dans
(\hyperref[II.1.3.1-fr]{1.3.1}), et nous laissons encore les détails au
lecteur. On note $\widetilde{\mathcal{M}}$ le $\mathcal{O}_X$-Module
$\mathcal{F}$, et on dit qu'il est \emph{associé} au $\mathcal{B}$-Module
quasi-cohérent $\mathcal{M}$~; pour tout ouvert affine $U\subset S$,
$\widetilde{\mathcal{M}}|p^{-1}(U)$ (où $p$ est le morphisme structural
$X\to S$) est canoniquement isomorphe à
$(\Gamma(U,\mathcal{M}))^{\sim}$.

\begin{corollary}[1.4.4]
\label{II.1.4.4-fr}
Dans la catégorie des $\mathcal{B}$-Modules quasi-cohérents,
$\widetilde{\mathcal{M}}$ est un foncteur covariant additif exact en
$\mathcal{M}$, qui commute aux limites inductives et aux sommes directes.
\end{corollary}

On est en effet aussitôt ramené au cas où $S$ est affine, et le corollaire
résulte alors de (\hyperref[I.1.3.5-fr]{I, 1.3.5}),
(\hyperref[I.1.3.9-fr]{I, 1.3.9}) et
(\hyperref[I.1.3.11-fr]{I, 1.3.11}).

\begin{corollary}[1.4.5]
\label{II.1.4.5-fr}
Sous les hypothèses de (\hyperref[II.1.4.3-fr]{1.4.3}), pour que
$\widetilde{\mathcal{M}}$ soit un $\mathcal{O}_X$-Module de type fini, il faut
et il suffit que $\mathcal{M}$ soit un $\mathcal{B}$-Module de type fini.
\end{corollary}

On est aussitôt ramené au cas où $S=\operatorname{Spec}(A)$ est un schéma
affine. Alors $\mathcal{B}=\widetilde{B}$, où $B$ est une $A$-algèbre de type
fini (\hyperref[I.9.6.2-fr]{I, 9.6.2}), et
$\mathcal{M}=\widetilde{M}$, où $M$ est un $B$-module
(\hyperref[I.1.3.13-fr]{I, 1.3.13})~; \emph{sur le préschéma $X$},
$\mathcal{O}_X$ est associé à l'anneau $B$ et
$\widetilde{\mathcal{M}}$ au $B$-module $M$~; pour que
$\widetilde{\mathcal{M}}$ soit de type fini, il faut et il suffit donc que
$M$ soit de type fini (\hyperref[I.1.3.13-fr]{I, 1.3.13}), d'où notre
assertion.

\begin{proposition}[1.4.6]
\label{II.1.4.6-fr}
Soient $Y$ un préschéma affine au-dessus de $S$, $X$, $X'$ deux préschémas
affines au-dessus de $Y$ (donc aussi au-dessus de $S$
(\hyperref[II.1.3.5-fr]{1.3.5})). Soient
$\mathcal{B}=\mathcal{A}(Y)$, $\mathcal{A}=\mathcal{A}(X)$,
$\mathcal{A}'=\mathcal{A}(X')$. Alors $X\times_Y X'$ est affine au-dessus de
$Y$ (donc aussi au-dessus de $S$) et $\mathcal{A}(X\times_Y X')$
s'identifie canoniquement à $\mathcal{A}\otimes_{\mathcal{B}}\mathcal{A}'$.
\end{proposition}

En effet (\hyperref[I.9.6.1-fr]{I, 9.6.1})
$\mathcal{A}\otimes_{\mathcal{B}}\mathcal{A}'$ est une
$\mathcal{B}$-Algèbre quasi-cohérente, donc aussi une
$\mathcal{O}_S$-Algèbre quasi-cohérente
(\hyperref[I.9.6.1-fr]{I, 9.6.1})~; soit $Z$ le spectre de
$\mathcal{A}\otimes_{\mathcal{B}}\mathcal{A}'$
(\hyperref[II.1.3.1-fr]{1.3.1}). Les $\mathcal{B}$-homomorphismes canoniques
$\mathcal{A}\to\mathcal{A}\otimes_{\mathcal{B}}\mathcal{A}'$ et
$\mathcal{A}'\to\mathcal{A}\otimes_{\mathcal{B}}\mathcal{A}'$ correspondent
(\hyperref[II.1.2.7-fr]{1.2.7}) à des $Y$-morphismes $p:Z\to X$ et
$p':Z\to X'$. Pour voir que le triplet $(Z,p,p')$ est un produit
$X\times_Y X'$, on peut se borner au cas où $S$ est un schéma affine
d'anneau $C$ (\hyperref[I.3.2.6.4-fr]{I, 3.2.6.4}). Mais alors $Y$, $X$,
$X'$ sont des schémas affines (\hyperref[II.1.3.4-fr]{1.3.4}) dont les
anneaux $B$, $A$, $A'$ sont des $C$-algèbres telles que
$\mathcal{B}=\widetilde{B}$, $\mathcal{A}=\widetilde{A}$,
$\mathcal{A}'=\widetilde{A'}$. On sait alors
(\hyperref[I.1.3.13-fr]{I, 1.3.13}) que
$\mathcal{A}\otimes_{\mathcal{B}}\mathcal{A}'$ s'identifie canoniquement à
la $\mathcal{O}_S$-Algèbre $\widetilde{A\otimes_B A'}$, donc l'anneau
$A(Z)$ s'identifie à $A\otimes_B A'$ et les morphismes $p$, $p'$ correspondent
aux homomorphismes canoniques $A\to A\otimes_B A'$,
$A'\to A\otimes_B A'$. La proposition résulte alors de
(\hyperref[I.3.2.2-fr]{I, 3.2.2}).

\begin{corollary}[1.4.7]
\label{II.1.4.7-fr}
Soit $\mathcal{F}$ (resp. $\mathcal{F}'$) un $\mathcal{O}_X$-Module (resp. un
$\mathcal{O}_{X'}$-Module) quasi-cohérent~; alors
$\mathcal{A}(\mathcal{F}\otimes_Y\mathcal{F}')$ s'identifie canoniquement à
$\mathcal{A}(\mathcal{F})\otimes_{\mathcal{A}(Y)}
\mathcal{A}(\mathcal{F}')$.
\end{corollary}

On sait que $\mathcal{F}\otimes_Y\mathcal{F}'$ est quasi-cohérent sur
$X\times_Y X'$ (\hyperref[I.9.1.2-fr]{I, 9.1.2}). Soient $g:Y\to S$,
$f:X\to Y$, $f':X'\to Y$ les morphismes structuraux, de sorte que le
morphisme structural
\oldpage[II]{11}
$h:Z\to S$ est égal à $g\circ f\circ p$ et à $g\circ f'\circ p'$. On
définit un homomorphisme canonique
\[
  \mathcal{A}(\mathcal{F})\otimes_{\mathcal{A}(Y)}\mathcal{A}(\mathcal{F}')
  \longrightarrow \mathcal{A}(\mathcal{F}\otimes_Y\mathcal{F}')
\]
de la façon suivante : pour tout ouvert $U\subset S$, on a des homomorphismes
canoniques
\[
  \begin{aligned}
  \Gamma(f^{-1}(g^{-1}(U)),\mathcal{F})
  &\longrightarrow \Gamma(h^{-1}(U),p^*(\mathcal{F}))\quad\text{et}\quad
  \Gamma(f^{\prime-1}(g^{-1}(U)),\mathcal{F}')
  \longrightarrow \Gamma(h^{-1}(U),p^{\prime*}(\mathcal{F}')).
  \end{aligned}
\]
(\hyperref[0.4.4.3-fr]{0, 4.4.3}), d'où on déduit un homomorphisme canonique
\[
  \begin{aligned}
  &\Gamma(f^{-1}(g^{-1}(U)),\mathcal{F})
    \otimes_{\Gamma(g^{-1}(U),\mathcal{O}_Y)}
    \Gamma(f^{\prime-1}(g^{-1}(U)),\mathcal{F}')\longrightarrow\\
  &\hspace{17em}
    \Gamma(h^{-1}(U),p^*(\mathcal{F}))
    \otimes_{\Gamma(h^{-1}(U),\mathcal{O}_Z)}
    \Gamma(h^{-1}(U),p^{\prime*}(\mathcal{F}')).
  \end{aligned}
\]

Pour voir qu'on a bien là un isomorphisme de $\mathcal{A}(Z)$-Modules, on
peut se borner au cas où $S$ (et par suite $X$, $X'$, $Y$,
$X\times_Y X'$) sont des schémas affines, et (avec les notations de
(\hyperref[II.1.4.6-fr]{1.4.6})), $\mathcal{F}=\widetilde{M}$,
$\mathcal{F}'=\widetilde{M'}$, $M$ (resp. $M'$) étant un $A$-module (resp. un
$A'$-module). Alors $\mathcal{F}\otimes_Y\mathcal{F}'$ s'identifie au
faisceau sur $X\times_Y X'$ associé au $(A\otimes_B A')$-module
$M\otimes_B M'$ (\hyperref[I.9.1.3-fr]{I, 9.1.3}) et le corollaire résulte de
l'identification canonique des $\mathcal{O}_S$-Modules
$\widetilde{M\otimes_B M'}$ et
$\widetilde{M}\otimes_{\widetilde{B}}\widetilde{M'}$ (où $M$, $M'$ et $B$
sont considérés comme des $C$-modules)
(\hyperref[I.1.3.13-fr]{I, 1.3.13} et
\hyperref[I.1.6.3-fr]{I, 1.6.3}).

Si on applique en particulier (\hyperref[II.1.4.7-fr]{1.4.7}) au cas où
$X=Y$ et $\mathcal{F}'=\mathcal{O}_{X'}$, on voit que le
$\mathcal{A}'$-Module $\mathcal{A}(f^{\prime*}(\mathcal{F}))$ s'identifie à
$\mathcal{A}(\mathcal{F})\otimes_{\mathcal{B}}\mathcal{A}'$.

\begin{env}[1.4.8]
\label{II.1.4.8-fr}
En particulier, lorsque $X=X'=Y$ ($X$ étant affine au-dessus de $S$), on voit
que si $\mathcal{F}$, $\mathcal{G}$ sont deux $\mathcal{O}_X$-Modules
quasi-cohérents, on a
\[
  \mathcal{A}(\mathcal{F}\otimes_{\mathcal{O}_X}\mathcal{G})
  =\mathcal{A}(\mathcal{F})\otimes_{\mathcal{A}(X)}\mathcal{A}(\mathcal{G})
  \tag{1.4.8.1}\label{II.1.4.8.1-fr}
\]
à un isomorphisme canonique fonctoriel près. Si de plus $\mathcal{F}$ admet
une présentation finie, il résulte de
(\hyperref[I.1.6.3-fr]{I, 1.6.3} et
\hyperref[I.1.3.12-fr]{I, 1.3.12}) que
\[
  \mathcal{A}(\mathcal{H}om_X(\mathcal{F},\mathcal{G}))
  =\mathcal{H}om_{\mathcal{A}(X)}
    (\mathcal{A}(\mathcal{F}),\mathcal{A}(\mathcal{G}))
  \tag{1.4.8.2}\label{II.1.4.8.2-fr}
\]
à un isomorphisme canonique près.
\end{env}

\begin{remark}[1.4.9]
\label{II.1.4.9-fr}
Si $X$, $X'$ sont deux préschémas affines au-dessus de $S$, la somme
$X\sqcup X'$ est aussi affine au-dessus de $S$, la somme de deux schémas
affines étant un schéma affine.
\end{remark}

\begin{proposition}[1.4.10]
\label{II.1.4.10-fr}
Soient $S$ un préschéma, $\mathcal{B}$ une $\mathcal{O}_S$-Algèbre
quasi-cohérente, $X=\operatorname{Spec}(\mathcal{B})$. Pour tout faisceau
quasi-cohérent d'idéaux $\mathcal{J}$ de $\mathcal{B}$,
$\widetilde{\mathcal{J}}$ est un faisceau quasi-cohérent d'idéaux de
$\mathcal{O}_X$, et le sous-préschéma fermé $Y$ de $X$ défini par
$\widetilde{\mathcal{J}}$ est canoniquement isomorphe à
$\operatorname{Spec}(\mathcal{B}/\mathcal{J})$.
\end{proposition}

En effet, il résulte aussitôt de (\hyperref[I.4.2.3-fr]{I, 4.2.3}) que $Y$
est affine au-dessus de $S$~; en vertu de
(\hyperref[II.1.3.1-fr]{1.3.1}), on est donc ramené au cas où $S$ est affine,
et la proposition résulte alors aussitôt de
(\hyperref[I.4.1.2-fr]{I, 4.1.2}).

On peut encore exprimer le résultat de
(\hyperref[II.1.4.10-fr]{1.4.10}) en disant que si
$h:\mathcal{B}\to\mathcal{B}'$ est un homomorphisme \emph{surjectif} de
$\mathcal{O}_S$-Algèbres quasi-cohérentes,
${}^a h:\operatorname{Spec}(\mathcal{B}')\to
\operatorname{Spec}(\mathcal{B})$ est une \emph{immersion fermée}.

\oldpage[II]{12}
\begin{proposition}[1.4.11]
\label{II.1.4.11-fr}
Soient $S$ un préschéma, $\mathcal{B}$ une $\mathcal{O}_S$-Algèbre
quasi-cohérente, $X=\operatorname{Spec}(\mathcal{B})$. Pour tout Idéal
quasi-cohérent $\mathcal{K}$ de $\mathcal{O}_S$, on a (en désignant par $f$
le morphisme structural $X\to S$),
$f^*(\mathcal{K})\mathcal{O}_X=\widetilde{\mathcal{K}\mathcal{B}}$ à un
isomorphisme canonique près.
\end{proposition}

La question étant locale sur $S$, on est ramené au cas où
$S=\operatorname{Spec}(A)$ est affine, et dans ce cas la proposition n'est
autre que (\hyperref[I.1.6.9-fr]{I, 1.6.9}).

\subsection{Changement du préschéma de base}
\label{II.1.5-fr}

\begin{proposition}[1.5.1]
\label{II.1.5.1-fr}
Soit $X$ un préschéma affine au-dessus de $S$. Pour toute extension
$g:S'\to S$ du préschéma de base, $X'=X_{(S')}=X\times_S S'$ est affine
au-dessus de $S'$.
\end{proposition}

Si $f'$ est la projection $X'\to S'$, il suffit de prouver que
$f^{\prime-1}(U')$ est un ouvert affine pour tout ouvert affine $U'$ de $S'$
tel que $g(U')$ soit contenu dans un ouvert affine $U$ de $S$
(\hyperref[II.1.2.1-fr]{1.2.1})~; on peut donc se limiter au cas où $S$ et
$S'$ sont affines, et il suffit de prouver que $X'$ est alors un schéma
affine (\hyperref[II.1.3.4-fr]{1.3.4}). Mais alors
(\hyperref[II.1.3.4-fr]{1.3.4}) $X$ est un schéma affine, et si $A$, $A'$ et
$B$ sont les anneaux de $S$, $S'$ et $X$ respectivement, on sait que $X'$ est
un schéma affine d'anneau $A'\otimes_A B$
(\hyperref[I.3.2.2-fr]{I, 3.2.2}).

\begin{corollary}[1.5.2]
\label{II.1.5.2-fr}
Sous les hypothèses de (\hyperref[II.1.5.1-fr]{1.5.1}), soient $f:X\to S$
le morphisme structural, $f':X'\to S'$, $g':X'\to X$ les projections, de
sorte que le diagramme
\[
  \xymatrix{
    X\ar[d] & X'\ar[l]_{g'}\ar[d]^{f'}\\
    S & S'\ar[l]_g
  }
\]
est commutatif. Pour tout $\mathcal{O}_X$-Module quasi-cohérent $\mathcal{F}$,
il existe un isomorphisme canonique de $\mathcal{O}_{S'}$-Modules
\[
  u:g^*(f_*(\mathcal{F}))\xrightarrow{\sim}
  f'_*(g^{\prime*}(\mathcal{F})).
  \tag{1.5.2.1}\label{II.1.5.2.1-fr}
\]
En particulier, il existe un isomorphisme canonique de $\mathcal{A}(X')$ sur
$g^*(\mathcal{A}(X))$.
\end{corollary}

Pour définir $u$, il suffit de définir un homomorphisme
\[
  v:f_*(\mathcal{F})\longrightarrow
  g_*(f'_*(g^{\prime*}(\mathcal{F})))
  =f_*(g'_*(g^{\prime*}(\mathcal{F})))
\]
et de prendre $u=v^\sharp$ (\hyperref[0.4.4.3-fr]{0, 4.4.3}). Nous prendrons
$v=f_*(\rho)$, où $\rho$ est l'homomorphisme canonique
$\mathcal{F}\to g'_*(g^{\prime*}(\mathcal{F}))$
(\hyperref[0.4.4.3-fr]{0, 4.4.3}). Pour prouver que $u$ est un isomorphisme,
on peut se borner au cas où $S$ et $S'$, donc $X$ et $X'$, sont affines~;
avec les notations de (\hyperref[II.1.5.1-fr]{1.5.1}), on a alors
$\mathcal{F}=\widetilde{M}$, où $M$ est un $B$-module. On constate alors
aussitôt que $g^*(f_*(\mathcal{F}))$ et
$f'_*(g^{\prime*}(\mathcal{F}))$ sont tous deux égaux au
$\mathcal{O}_{S'}$-Module associé au $A'$-module $A'\otimes_A M$ (où $M$ est
considéré comme $A$-module), et que $u$ est l'homomorphisme associé à
l'identité (\hyperref[I.1.6.3-fr]{I, 1.6.3},
\hyperref[I.1.6.5-fr]{1.6.5} et \hyperref[I.1.6.7-fr]{1.6.7}).

\begin{remark}[1.5.3]
\label{II.1.5.3-fr}
On se gardera de croire que (\hyperref[II.1.5.2-fr]{1.5.2}) reste valable
lorsque $X$ n'est plus supposé affine au-dessus de $S$, même lorsque
$S'=\operatorname{Spec}(k(s))$ ($s\in S$) et que $S'\to S$ est le morphisme
canonique (\hyperref[I.2.4.5-fr]{I, 2.4.5}) --- auquel cas $X'$ n'est autre
que la \emph{fibre} $f^{-1}(s)$ (\hyperref[I.3.6.2-fr]{I, 3.6.2}). En
d'autres termes, lorsque $X$ n'est pas affine au-dessus de $S$, l'opération
\oldpage[II]{13}
<<~image directe de faisceaux quasi-cohérents~>> ne permute pas à l'opération
de <<~passage aux fibres~>>. Nous verrons cependant au chapitre~III
(\hyperref[III.4.2.4-fr]{III, 4.2.4}) un résultat dans ce sens, de nature
<<~asymptotique~>>, valable pour les faisceaux \emph{cohérents} sur $X$ lorsque
$f$ est propre (\hyperref[II.5.4-fr]{5.4}) et $S$ noethérien.
\end{remark}

\begin{corollary}[1.5.4]
\label{II.1.5.4-fr}
Pour tout préschéma $X$ affine au-dessus de $S$ et tout $s\in S$, la fibre
$f^{-1}(s)$ (où $f$ désigne le morphisme structural $X\to S$) est un schéma
affine.
\end{corollary}

Il suffit d'appliquer (\hyperref[II.1.5.1-fr]{1.5.1}) à
$S'=\operatorname{Spec}(k(s))$ et d'utiliser
(\hyperref[II.1.3.4-fr]{1.3.4}).

\begin{corollary}[1.5.5]
\label{II.1.5.5-fr}
Soient $X$ un $S$-préschéma, $S'$ un préschéma affine au-dessus de $S$~; alors
$X'=X_{(S')}$ est un préschéma affine au-dessus de $X$. En outre, si
$f:X\to S$ est le morphisme structural, il y a un isomorphisme canonique de
$\mathcal{O}_X$-Algèbres
$\mathcal{A}(X')\xrightarrow{\sim}f^*(\mathcal{A}(S'))$, et pour tout
$\mathcal{A}(S')$-Module quasi-cohérent $\mathcal{M}$, un di-isomorphisme
canonique
$f^*(\mathcal{M})\xrightarrow{\sim}
\mathcal{A}(f^{\prime*}(\widetilde{\mathcal{M}}))$, en désignant par
$f'=f_{(S')}$ le morphisme structural $X'\to S'$.
\end{corollary}

Il suffit d'intervertir les rôles de $X$ et de $S'$ dans
(\hyperref[II.1.5.1-fr]{1.5.1}) et
(\hyperref[II.1.5.2-fr]{1.5.2}).

\begin{env}[1.5.6]
\label{II.1.5.6-fr}
Soient maintenant $S$, $S'$ deux préschémas, $q:S'\to S$ un morphisme,
$\mathcal{B}$ (resp. $\mathcal{B}'$) une $\mathcal{O}_S$-Algèbre (resp. une
$\mathcal{O}_{S'}$-Algèbre) quasi-cohérente, $u:\mathcal{B}\to\mathcal{B}'$
un $q$-morphisme (c'est-à-dire un homomorphisme
$\mathcal{B}\to q_*(\mathcal{B}')$ de $\mathcal{O}_S$-Algèbres). Si
$X=\operatorname{Spec}(\mathcal{B})$,
$X'=\operatorname{Spec}(\mathcal{B}')$, on en déduit canoniquement un
morphisme
\[
  v=\operatorname{Spec}(u):X'\to X
\]
tel que le diagramme
\[
  \xymatrix{
    X'\ar[r]^v\ar[d] & X\ar[d]\\
    S'\ar[r]_q & S
  }
  \tag{1.5.6.1}\label{II.1.5.6.1-fr}
\]
soit commutatif (les flèches verticales étant les morphismes structuraux). En
effet, la donnée de $u$ équivaut à celle d'un homomorphisme de
$\mathcal{O}_{S'}$-Algèbres quasi-cohérentes
$u^\sharp:q^*(\mathcal{B})\to\mathcal{B}'$
(\hyperref[0.4.4.3-fr]{0, 4.4.3})~; elle définit donc canoniquement un
$S'$-morphisme
\[
  w:\operatorname{Spec}(\mathcal{B}')\to
  \operatorname{Spec}(q^*(\mathcal{B}))
\]
tel que $\mathcal{A}(w)=u^\sharp$
(\hyperref[II.1.2.7-fr]{1.2.7}). D'autre part, il résulte de
(\hyperref[II.1.5.2-fr]{1.5.2}) que
$\operatorname{Spec}(q^*(\mathcal{B}))$ s'identifie canoniquement à
$X\times_S S'$~; le morphisme $v$ est le composé
$X'\xrightarrow{w}X\times_S S'\xrightarrow{p_1}X$ de $w$ avec la première
projection, et la commutativité de
(\hyperref[II.1.5.6.1-fr]{1.5.6.1}) résulte des définitions. Soient $U$ (resp.
$U'$) un ouvert affine de $S$ (resp. $S'$) tels que $q(U')\subset U$,
$A=\Gamma(U,\mathcal{O}_S)$, $A'=\Gamma(U',\mathcal{O}_{S'})$ leurs anneaux,
$B=\Gamma(U,\mathcal{B})$, $B'=\Gamma(U',\mathcal{B}')$~; la restriction de
$u$ en un $(q|U')$-morphisme
$\mathcal{B}|U\to\mathcal{B}'|U'$ correspond à un di-homomorphisme d'algèbres
$B\to B'$~; si $V$, $V'$ sont les images réciproques de $U$, $U'$ dans $X$,
$X'$ respectivement, par les morphismes structuraux, le morphisme $V'\to V$,
restriction de $v$, correspond (\hyperref[I.1.7.3-fr]{I, 1.7.3}) au
di-homomorphisme précédent.
\end{env}

\begin{env}[1.5.7]
\label{II.1.5.7-fr}
Sous les mêmes hypothèses que dans (\hyperref[II.1.5.6-fr]{1.5.6}), soit
$\mathcal{M}$ un $\mathcal{B}$-Module quasi-cohérent~; il y a alors un
isomorphisme canonique de $\mathcal{O}_{X'}$-Modules
\[
  v^*(\widetilde{\mathcal{M}})\xrightarrow{\sim}
  \bigl(q^*(\mathcal{M})\otimes_{q^*(\mathcal{B})}\mathcal{B}'\bigr)^{\sim}
  \tag{1.5.7.1}\label{II.1.5.7.1-fr}
\]

\oldpage[II]{14}
En effet, l'isomorphisme canonique
(\hyperref[II.1.5.2.1-fr]{1.5.2.1}) fournit un isomorphisme canonique de
$p_1^*(\widetilde{\mathcal{M}})$ sur le faisceau sur
$\operatorname{Spec}(q^*(\mathcal{B}))$ associé au
$q^*(\mathcal{B})$-Module $q^*(\mathcal{M})$, et il suffit ensuite d'appliquer
(\hyperref[II.1.4.7-fr]{1.4.7}).
\end{env}

\subsection{Morphismes affines}
\label{II.1.6-fr}

\begin{env}[1.6.1]
\label{II.1.6.1-fr}
Nous dirons qu'un morphisme $f:X\to Y$ de préschémas est \emph{affine} s'il
définit $X$ comme un préschéma affine au-dessus de $Y$. Les propriétés des
préschémas affines au-dessus d'un autre se traduisent comme suit dans ce
langage~:
\end{env}

\begin{proposition}[1.6.2]
\label{II.1.6.2-fr}
\begin{enumerate}
  \item[\rm(i)] Une immersion fermée est affine.
  \item[\rm(ii)] Le composé de deux morphismes affines est affine.
  \item[\rm(iii)] Si $f:X\to Y$ est un $S$-morphisme affine,
    $f_{(S')}:X_{(S')}\to Y_{(S')}$ est affine pour toute extension $S'\to S$
    de la base.
  \item[\rm(iv)] Si $f:X\to Y$ et $f':X'\to Y'$ sont deux $S$-morphismes
    affines,
    \[
      f\times_S f':X\times_S X'\to Y\times_S Y'
    \]
    est affine.
  \item[\rm(v)] Si $f:X\to Y$ et $g:Y\to Z$ sont deux morphismes tels que
    $g\circ f$ soit affine et $g$ séparé, alors $f$ est affine.
  \item[\rm(vi)] Si $f$ est affine, il en est de même de
    $f_{\mathrm{red}}$.
\end{enumerate}
\end{proposition}

En vertu de (\hyperref[I.5.5.12-fr]{I, 5.5.12}), il suffit de démontrer (i),
(ii) et (iii). Or (i) n'est autre que l'exemple
(\hyperref[II.1.2.2-fr]{1.2.2}), et (ii) n'est autre que
(\hyperref[II.1.3.5-fr]{1.3.5})~; enfin, (iii) découle de
(\hyperref[II.1.5.1-fr]{1.5.1}) puisque $X_{(S')}$ s'identifie au produit
$X\times_Y Y_{(S')}$ (\hyperref[I.3.3.11-fr]{I, 3.3.11}).

\begin{corollary}[1.6.3]
\label{II.1.6.3-fr}
Si $X$ est un schéma affine et $Y$ un schéma, tout morphisme $X\to Y$ est
affine.
\end{corollary}

\begin{proposition}[1.6.4]
\label{II.1.6.4-fr}
Soient $Y$ un préschéma localement noethérien, $f:X\to Y$ un morphisme de type
fini. Pour que $f$ soit affine, il faut et il suffit que $f_{\mathrm{red}}$ le
soit.
\end{proposition}

Vu (\hyperref[II.1.6.2-fr]{1.6.2, (vi)}), nous n'avons à démontrer que la
suffisance de la condition. Il suffit de prouver que si $Y$ est affine et
noethérien, $X$ est affine~; or $Y_{\mathrm{red}}$ est alors affine, donc il
en est de même de $X_{\mathrm{red}}$ par hypothèse. Or, $X$ est noethérien,
donc la conclusion résulte de (\hyperref[I.6.1.7-fr]{I, 6.1.7}).

\subsection{Fibré vectoriel associé à un Faisceau de modules}
\label{II.1.7-fr}

\begin{env}[1.7.1]
\label{II.1.7.1-fr}
Soient $A$ un anneau, $E$ un $A$-module. Rappelons qu'on appelle \emph{algèbre
symétrique} sur $E$ et qu'on note $\mathbf{S}(E)$ (ou $\mathbf{S}_A(E)$)
l'algèbre quotient de l'algèbre tensorielle $\mathbf{T}(E)$ par l'idéal
bilatère engendré par les éléments $x\otimes y-y\otimes x$, où $x,y$
parcourent $E$. L'algèbre $\mathbf{S}(E)$ est caractérisée par la propriété
universelle suivante~: si $\sigma$ est l'application canonique
$E\to\mathbf{S}(E)$ (obtenue par composition de $E\to\mathbf{T}(E)$ et de
l'application canonique $\mathbf{T}(E)\to\mathbf{S}(E)$), toute application
$A$-linéaire $E\to B$, où $B$ est une $A$-algèbre \emph{commutative}, se
factorise de façon unique en
$E\xrightarrow{\sigma}\mathbf{S}(E)\xrightarrow{g}B$, où $g$ est un
$A$-homomorphisme \emph{d'algèbres}. On déduit aussitôt de cette
caractérisation que pour deux $A$-modules $E$, $F$, on a
\[
  \mathbf{S}(E\oplus F)=\mathbf{S}(E)\otimes\mathbf{S}(F)
\]

\oldpage[II]{15}
à un isomorphisme canonique près~; en outre, $\mathbf{S}(E)$ est un foncteur
covariant en $E$, de la catégorie des $A$-modules dans celle des $A$-algèbres
commutatives~; enfin, la caractérisation précédente montre aussi que si
$E=\varinjlim E_\lambda$, on a
$\mathbf{S}(E)=\varinjlim\mathbf{S}(E_\lambda)$ à un isomorphisme canonique
près. Par abus de langage, un produit
$\sigma(x_1)\sigma(x_2)\cdots\sigma(x_n)$, où les $x_i\in E$, se note souvent
$x_1x_2\cdots x_n$ si aucune confusion n'en résulte. L'algèbre
$\mathbf{S}(E)$ est \emph{graduée}, $\mathbf{S}_n(E)$ étant l'ensemble des
combinaisons linéaires des produits de $n$ éléments de $E$ ($n\geq 0$)~;
l'algèbre $\mathbf{S}(A)$ est canoniquement isomorphe à l'algèbre des
polynômes $A[T]$ à une indéterminée et l'algèbre $\mathbf{S}(A^n)$ à
l'algèbre des polynômes à $n$ indéterminées $A[T_1,\ldots,T_n]$.
\end{env}

\begin{env}[1.7.2]
\label{II.1.7.2-fr}
Soit $\varphi$ un homomorphisme d'anneaux $A\to B$. Si $F$ est un $B$-module,
l'application canonique $F\to\mathbf{S}(F)$ donne une application canonique
$F_{[\varphi]}\to\mathbf{S}(F)_{[\varphi]}$, qui se factorise donc en
$F_{[\varphi]}\to\mathbf{S}(F_{[\varphi]})\to
\mathbf{S}(F)_{[\varphi]}$~; l'homomorphisme canonique
$\mathbf{S}(F_{[\varphi]})\to\mathbf{S}(F)_{[\varphi]}$ est surjectif, mais
non nécessairement bijectif. Si $E$ est un $A$-module, tout di-homomorphisme
$E\to F$ (c'est-à-dire tout $A$-homomorphisme $E\to F_{[\varphi]}$) donne donc
canoniquement un $A$-homomorphisme d'algèbres
$\mathbf{S}(E)\to\mathbf{S}(F_{[\varphi]})\to
\mathbf{S}(F)_{[\varphi]}$, c'est-à-dire un di-homomorphisme d'algèbres
$\mathbf{S}(E)\to\mathbf{S}(F)$.

Avec les mêmes notations, pour tout $A$-module $E$,
$\mathbf{S}(E\otimes_A B)$ s'identifie canoniquement à l'algèbre
$\mathbf{S}(E)\otimes_A B$~; cela résulte aussitôt de la propriété universelle
de $\mathbf{S}(E)$ (\hyperref[II.1.7.1-fr]{1.7.1}).
\end{env}

\begin{env}[1.7.3]
\label{II.1.7.3-fr}
Soit $R$ une partie multiplicative de l'anneau $A$~; appliquant
(\hyperref[II.1.7.2-fr]{1.7.2}) à l'anneau $B=R^{-1}A$, et se rappelant que
$R^{-1}E=E\otimes_A R^{-1}A$, on voit que l'on a
$\mathbf{S}(R^{-1}E)=R^{-1}\mathbf{S}(E)$ à un isomorphisme canonique près.
En outre, si $R'\supset R$ est une seconde partie multiplicative de $A$, le
diagramme
\[
  \xymatrix{
    R^{-1}E\ar[r]\ar[d] & {R'}^{-1}E\ar[d]\\
    \mathbf{S}(R^{-1}E)\ar[r] & \mathbf{S}({R'}^{-1}E)
  }
\]
est commutatif.
\end{env}

\begin{env}[1.7.4]
\label{II.1.7.4-fr}
Soit maintenant $(S,\mathcal{A})$ un espace annelé, et soit $\mathcal{E}$ un
$\mathcal{A}$-Module sur $S$. Si, à tout ouvert $U\subset S$ on associe le
$\Gamma(U,\mathcal{A})$-module
$\mathbf{S}(\Gamma(U,\mathcal{E}))$, on définit (vu le caractère fonctoriel
de $\mathbf{S}(E)$ (\hyperref[II.1.7.2-fr]{1.7.2})) un préfaisceau d'algèbres~;
on dit que le faisceau associé, que l'on note $\mathbf{S}(\mathcal{E})$ ou
$\mathbf{S}_{\mathcal{A}}(\mathcal{E})$ est la $\mathcal{A}$-Algèbre
\emph{symétrique} du $\mathcal{A}$-Module $\mathcal{E}$. Il résulte aussitôt
de (\hyperref[II.1.7.1-fr]{1.7.1}) que $\mathbf{S}(\mathcal{E})$ est solution
d'un problème universel~: tout homomorphisme de $\mathcal{A}$-Modules
$\mathcal{E}\to\mathcal{B}$, où $\mathcal{B}$ est une $\mathcal{A}$-Algèbre,
se factorise d'une seule manière en
$\mathcal{E}\to\mathbf{S}(\mathcal{E})\to\mathcal{B}$, la seconde flèche étant
un homomorphisme de $\mathcal{A}$-Algèbres. Il y a donc correspondance
biunivoque entre homomorphismes $\mathcal{E}\to\mathcal{B}$ de
$\mathcal{A}$-Modules et homomorphismes
$\mathbf{S}(\mathcal{E})\to\mathcal{B}$ de $\mathcal{A}$-Algèbres. En
particulier, tout homomorphisme $u:\mathcal{E}\to\mathcal{F}$ de
$\mathcal{A}$-Modules définit un homomorphisme
$\mathbf{S}(u):\mathbf{S}(\mathcal{E})\to\mathbf{S}(\mathcal{F})$ de
$\mathcal{A}$-Algèbres et $\mathbf{S}(\mathcal{E})$ est donc un foncteur
covariant en $\mathcal{E}$.

\oldpage[II]{16}
En vertu de (\hyperref[II.1.7.2-fr]{1.7.2}) et de la permutabilité de
$\mathbf{S}$ avec les limites inductives, on a
$(\mathbf{S}(\mathcal{E}))_s=\mathbf{S}(\mathcal{E}_s)$ pour tout point
$s\in S$. Si $\mathcal{E}$, $\mathcal{F}$ sont deux $\mathcal{A}$-Modules,
$\mathbf{S}(\mathcal{E}\oplus\mathcal{F})$ s'identifie canoniquement à
$\mathbf{S}(\mathcal{E})\otimes_{\mathcal{A}}\mathbf{S}(\mathcal{F})$,
comme on le voit sur les préfaisceaux correspondants.

Notons aussi que $\mathbf{S}(\mathcal{E})$ est une $\mathcal{A}$-Algèbre
graduée, somme directe infinie des $\mathbf{S}_n(\mathcal{E})$, où le
$\mathcal{A}$-Module $\mathbf{S}_n(\mathcal{E})$ est le faisceau associé au
préfaisceau $U\mapsto\mathbf{S}_n(\Gamma(U,\mathcal{E}))$. Si on prend en
particulier $\mathcal{E}=\mathcal{A}$, on voit que
$\mathbf{S}_{\mathcal{A}}(\mathcal{A})$ s'identifie à
$\mathcal{A}[T]=\mathcal{A}\otimes_{\mathbf{Z}}\mathbf{Z}[T]$
($T$ indéterminée, $\mathbf{Z}$ étant considéré comme faisceau simple).
\end{env}

\begin{env}[1.7.5]
\label{II.1.7.5-fr}
Soient $(T,\mathcal{B})$ un second espace annelé, $f$ un morphisme
$(S,\mathcal{A})\to(T,\mathcal{B})$. Si $\mathcal{F}$ est un
$\mathcal{B}$-Module, $\mathbf{S}(f^*(\mathcal{F}))$ s'identifie
canoniquement à $f^*(\mathbf{S}(\mathcal{F}))$~; en effet, si
$f=(\psi,\theta)$, on a par définition
(\hyperref[0.4.3.1-fr]{0, 4.3.1})
\[
  \mathbf{S}(f^*(\mathcal{F}))
  =\mathbf{S}(\psi^*(\mathcal{F})
    \otimes_{\psi^*(\mathcal{B})}\mathcal{A})
  =\mathbf{S}(\psi^*(\mathcal{F}))
    \otimes_{\psi^*(\mathcal{B})}\mathcal{A}
\]
(\hyperref[II.1.7.2-fr]{1.7.2})~; pour tout ouvert $U$ de $S$ et toute
section $h$ de $\mathbf{S}(\psi^*(\mathcal{F}))$ au-dessus de $U$, $h$
coïncide, dans un voisinage $V$ de chaque point $s\in U$, avec un élément de
$\mathbf{S}(\Gamma(V,\psi^*(\mathcal{F})))$~; si on se reporte à la
définition de $\psi^*(\mathcal{F})$
(\hyperref[0.3.7.1-fr]{0, 3.7.1}) et qu'on tient compte de ce que tout élément
de $\mathbf{S}(E)$ pour un module $E$ est combinaison linéaire d'un nombre
fini de produits d'éléments de $E$, on voit qu'il y a un voisinage $W$ de
$\psi(s)$ dans $T$, une section $h'$ de $\mathbf{S}(\mathcal{F})$ au-dessus
de $W$ et un voisinage $V'\subset V\cap\psi^{-1}(W)$ de $s$ tels que $h$
coïncide avec $t\mapsto h'(\psi(t))$ au-dessus de $V'$~; d'où notre
assertion.
\end{env}

\begin{proposition}[1.7.6]
\label{II.1.7.6-fr}
Soient $A$ un anneau, $S=\operatorname{Spec}(A)$ son spectre premier,
$\mathcal{E}=\widetilde{M}$ le $\mathcal{O}_S$-Module associé à un $A$-module
$M$~; alors la $\mathcal{O}_S$-Algèbre $\mathbf{S}(\mathcal{E})$ est associée
à la $A$-algèbre $\mathbf{S}(M)$.
\end{proposition}

\begin{proof}
En effet, pour tout $f\in A$,
$\mathbf{S}(M_f)=(\mathbf{S}(M))_f$
(\hyperref[II.1.7.3-fr]{1.7.3}) et la proposition résulte donc des définitions
(\hyperref[I.1.3.4-fr]{I, 1.3.4}).
\end{proof}

\begin{corollary}[1.7.7]
\label{II.1.7.7-fr}
Si $S$ est un préschéma, $\mathcal{E}$ un $\mathcal{O}_S$-Module
quasi-cohérent, la $\mathcal{O}_S$-Algèbre $\mathbf{S}(\mathcal{E})$ est
quasi-cohérente. Si en outre $\mathcal{E}$ est de type fini, chacun des
$\mathcal{O}_S$-Modules $\mathbf{S}_n(\mathcal{E})$ est de type fini.
\end{corollary}

\begin{proof}
La première assertion est conséquence immédiate de
(\hyperref[II.1.7.6-fr]{1.7.6}) et de
(\hyperref[I.1.4.1-fr]{I, 1.4.1})~; la seconde résulte de ce que, si $E$ est
un $A$-module de type fini, $\mathbf{S}_n(E)$ est un $A$-module de type
fini~; on applique alors (\hyperref[I.1.3.13-fr]{I, 1.3.13}).
\end{proof}

\begin{definition}[1.7.8]
\label{II.1.7.8-fr}
Soit $\mathcal{E}$ un $\mathcal{O}_S$-Module quasi-cohérent. On appelle
\emph{fibré vectoriel sur $S$ défini par $\mathcal{E}$} et on note
$\mathbf{V}(\mathcal{E})$ le spectre
(\hyperref[II.1.3.1-fr]{1.3.1}) de la $\mathcal{O}_S$-Algèbre quasi-cohérente
$\mathbf{S}(\mathcal{E})$.
\end{definition}

En vertu de (\hyperref[II.1.2.7-fr]{1.2.7}), pour tout $S$-préschéma $X$, il
y a correspondance biunivoque canonique entre les $S$-morphismes
$X\to\mathbf{V}(\mathcal{E})$ et les homomorphismes de
$\mathcal{O}_S$-Algèbres $\mathbf{S}(\mathcal{E})\to\mathcal{A}(X)$, donc
aussi entre ces $S$-morphismes et les homomorphismes de
$\mathcal{O}_S$-Modules
$\mathcal{E}\to\mathcal{A}(X)=f_*(\mathcal{O}_X)$ (où $f$ est le morphisme
structural $X\to S$). En particulier~:

\begin{env}[1.7.9]
\label{II.1.7.9-fr}
Prenons pour $X$ un sous-préschéma induit par $S$ sur un ouvert $U\subset S$.
Alors, les $S$-morphismes $U\to\mathbf{V}(\mathcal{E})$ ne sont autres que
les $U$-sections (\hyperref[I.2.5.5-fr]{I, 2.5.5}) du $U$-préschéma induit
par $\mathbf{V}(\mathcal{E})$ sur l'ouvert $p^{-1}(U)$ (où $p$ est le
morphisme structural $\mathbf{V}(\mathcal{E})\to S$). D'après ce qu'on vient
de voir, ces $U$-sections correspondent biunivoquement aux homomorphismes de
$\mathcal{O}_S$-Modules
$\mathcal{E}\to j_*(\mathcal{O}_S|U)$ (où $j$ est l'injection canonique
$U\to S$), ou,

\oldpage[II]{17}
ce qui revient au même (\hyperref[0.4.4.3-fr]{0, 4.4.3}) aux
$(\mathcal{O}_S|U)$-homomorphismes
$j^*(\mathcal{E})=\mathcal{E}|U\to\mathcal{O}_S|U$. En outre, il est
immédiat qu'à la restriction à un ouvert $U'\subset U$ d'un $S$-morphisme
$U\to\mathbf{V}(\mathcal{E})$ correspond la restriction à $U'$ de
l'homomorphisme correspondant
$\mathcal{E}|U\to\mathcal{O}_S|U$. On en conclut que \emph{le faisceau des
germes de $S$-sections} de $\mathbf{V}(\mathcal{E})$ s'identifie
canoniquement au \emph{dual $\check{\mathcal{E}}$} de $\mathcal{E}$.

En particulier, si on prend $X=U=S$, à l'homomorphisme \emph{nul}
$\mathcal{E}\to\mathcal{O}_S$ correspond une $S$-section canonique de
$\mathbf{V}(\mathcal{E})$, dite \emph{$S$-section nulle}
(cf. (\hyperref[II.8.3.3-fr]{8.3.3})).
\end{env}

\begin{env}[1.7.10]
\label{II.1.7.10-fr}
Prenons maintenant pour $X$ le spectre $\{\xi\}$ d'un corps $K$~; le
morphisme structural $f:X\to S$ correspond donc à un monomorphisme
$k(s)\to K$, où $s=f(\xi)$
(\hyperref[I.2.4.6-fr]{I, 2.4.6})~; les $S$-morphismes
$\{\xi\}\to\mathbf{V}(\mathcal{E})$ ne sont donc autres que les
\emph{points géométriques de $\mathbf{V}(\mathcal{E})$ à valeurs dans
l'extension $K$ de $k(s)$} (\hyperref[I.3.4.5-fr]{I, 3.4.5}), points qui
sont localisés aux points de $p^{-1}(s)$. L'ensemble de ces points, qu'on
peut encore appeler \emph{la fibre géométrique rationnelle sur $K$ de
$\mathbf{V}(\mathcal{E})$ au-dessus du point $s$}, s'identifie d'après
(\hyperref[II.1.7.8-fr]{1.7.8}) à l'ensemble des homomorphismes de
$\mathcal{O}_S$-Modules
$\mathcal{E}\to f_*(\mathcal{O}_X)$, ou, ce qui revient au même
(\hyperref[0.4.4.3-fr]{0, 4.4.3}) à l'ensemble des homomorphismes de
$\mathcal{O}_X$-Modules
$f^*(\mathcal{E})\to\mathcal{O}_X=K$. Mais on a par définition
(\hyperref[0.4.3.1-fr]{0, 4.3.1})
$f^*(\mathcal{E})=\mathcal{E}_s\otimes_{\mathcal{O}_s}K
=\mathcal{E}^s\otimes_{k(s)}K$, en posant
$\mathcal{E}^s=\mathcal{E}_s/\mathfrak{m}_s\mathcal{E}_s$~; la fibre
géométrique de $\mathbf{V}(\mathcal{E})$ rationnelle sur $K$ au-dessus de
$s$ s'identifie donc au \emph{dual} du \emph{$K$-espace vectoriel
$\mathcal{E}^s\otimes_{k(s)}K$}~; si $\mathcal{E}^s$ ou $K$ est de
dimension finie sur $k(s)$, ce dual s'identifie aussi à
$(\mathcal{E}^s)^\vee\otimes_{k(s)}K$, en désignant par
$(\mathcal{E}^s)^\vee$ le dual du $k(s)$-espace vectoriel
$\mathcal{E}^s$.
\end{env}

\begin{proposition}[1.7.11]
\label{II.1.7.11-fr}
\begin{enumerate}
  \item[\rm(i)] $\mathbf{V}(\mathcal{E})$ est un foncteur contravariant en
    $\mathcal{E}$ de la catégorie des $\mathcal{O}_S$-Modules
    quasi-cohérents dans la catégorie des $S$-schémas affines.
  \item[\rm(ii)] Si $\mathcal{E}$ est un $\mathcal{O}_S$-Module de type
    fini, $\mathbf{V}(\mathcal{E})$ est de type fini sur $S$.
  \item[\rm(iii)] Si $\mathcal{E}$ et $\mathcal{F}$ sont deux
    $\mathcal{O}_S$-Modules quasi-cohérents,
    $\mathbf{V}(\mathcal{E}\oplus\mathcal{F})$ s'identifie canoniquement à
    $\mathbf{V}(\mathcal{E})\times_S\mathbf{V}(\mathcal{F})$.
  \item[\rm(iv)] Soit $g:S'\to S$ un morphisme~; pour tout
    $\mathcal{O}_S$-Module quasi-cohérent $\mathcal{E}$,
    $\mathbf{V}(g^*(\mathcal{E}))$ s'identifie canoniquement à
    $\mathbf{V}(\mathcal{E})_{(S')}
    =\mathbf{V}(\mathcal{E})\times_S S'$.
  \item[\rm(v)] À un homomorphisme surjectif
    $\mathcal{E}\to\mathcal{F}$ de $\mathcal{O}_S$-Modules quasi-cohérents
    correspond une immersion fermée
    $\mathbf{V}(\mathcal{F})\to\mathbf{V}(\mathcal{E})$.
\end{enumerate}
\end{proposition}

\begin{proof}
(i) est conséquence immédiate de
(\hyperref[II.1.2.7-fr]{1.2.7}), compte tenu de ce que tout homomorphisme de
$\mathcal{O}_S$-Modules $\mathcal{E}\to\mathcal{F}$ définit canoniquement un
homomorphisme de $\mathcal{O}_S$-Algèbres
$\mathbf{S}(\mathcal{E})\to\mathbf{S}(\mathcal{F})$. (ii) découle
immédiatement des définitions (\hyperref[I.6.3.1-fr]{I, 6.3.1}) et du fait
que si $E$ est un $A$-module de type fini, $\mathbf{S}(E)$ est une
$A$-algèbre de type fini. Pour démontrer (iii), il suffit de partir de
l'isomorphisme canonique
$\mathbf{S}(\mathcal{E}\oplus\mathcal{F})
\simeq\mathbf{S}(\mathcal{E})
\otimes_{\mathcal{O}_S}\mathbf{S}(\mathcal{F})$
(\hyperref[II.1.7.4-fr]{1.7.4}) et d'appliquer
(\hyperref[II.1.4.6-fr]{1.4.6}). De même, pour démontrer (iv), il suffit de
partir de l'isomorphisme canonique
$\mathbf{S}(g^*(\mathcal{E}))
\simeq g^*(\mathbf{S}(\mathcal{E}))$
(\hyperref[II.1.7.5-fr]{1.7.5}) et d'appliquer
(\hyperref[II.1.5.2-fr]{1.5.2}). Enfin, pour établir (v), il suffit de
remarquer que si l'homomorphisme $\mathcal{E}\to\mathcal{F}$ est surjectif,
il en est de même de l'homomorphisme correspondant
$\mathbf{S}(\mathcal{E})\to\mathbf{S}(\mathcal{F})$ de
$\mathcal{O}_S$-Algèbres, et la conclusion résulte de
(\hyperref[II.1.4.10-fr]{1.4.10}).
\end{proof}

\begin{env}[1.7.12]
\label{II.1.7.12-fr}
Prenons en particulier $\mathcal{E}=\mathcal{O}_S$~; le préschéma
$\mathbf{V}(\mathcal{O}_S)$ est le $S$-schéma affine, spectre de la
$\mathcal{O}_S$-Algèbre $\mathbf{S}(\mathcal{O}_S)$ qui s'identifie à la
$\mathcal{O}_S$-Algèbre
$\mathcal{O}_S[T]=\mathcal{O}_S\otimes_{\mathbf{Z}}\mathbf{Z}[T]$
\oldpage[II]{18}
($T$ indéterminée)~; c'est évident lorsque
$S=\operatorname{Spec}(\mathbf{Z})$, en vertu de
(\hyperref[II.1.7.6-fr]{1.7.6}), et on passe de là au cas général en
considérant le morphisme structural
$S\to\operatorname{Spec}(\mathbf{Z})$ et utilisant
(\hyperref[II.1.7.11-fr]{1.7.11, (iv)}). En raison de ce résultat, on pose
encore $\mathbf{V}(\mathcal{O}_S)=S[T]$, et on a donc la formule
\[
  S[T]=S\otimes_{\mathbf{Z}}\mathbf{Z}[T].
  \tag{1.7.12.1}\label{II.1.7.12.1-fr}
\]

L'identification du faisceau des germes de $S$-sections de $S[T]$ avec
$\mathcal{O}_S$, déjà vue dans
(\hyperref[I.3.3.15-fr]{I, 3.3.15}), se retrouve ici dans un contexte plus
général, comme cas particulier de (\hyperref[II.1.7.9-fr]{1.7.9}).
\end{env}

\begin{env}[1.7.13]
\label{II.1.7.13-fr}
Pour tout $S$-préschéma $X$, on a vu
(\hyperref[II.1.7.8-fr]{1.7.8}) que
$\operatorname{Hom}_S(X,S[T])$ s'identifie canoniquement à
$\operatorname{Hom}_{\mathcal{O}_S}(\mathcal{O}_S,\mathcal{A}(X))$, lequel
est canoniquement isomorphe à $\Gamma(S,\mathcal{A}(X))$, et est par suite
muni d'une structure d'anneau~; en outre, à tout $S$-morphisme $h:X\to Y$
correspond un morphisme
$\Gamma(\mathcal{A}(h)):\Gamma(S,\mathcal{A}(Y))
\to\Gamma(S,\mathcal{A}(X))$ pour les structures d'anneau
(\hyperref[II.1.1.2-fr]{1.1.2}). Quand on munit
$\operatorname{Hom}_S(X,S[T])$ de la structure d'anneau ainsi définie, on
voit donc qu'on peut considérer $\operatorname{Hom}_S(X,S[T])$ comme un
\emph{foncteur contravariant en $X$}, de la catégorie des $S$-préschémas
dans celle des anneaux. D'autre part,
$\operatorname{Hom}_S(X,\mathbf{V}(\mathcal{E}))$ s'identifie de même à
$\operatorname{Hom}_{\mathcal{O}_S}(\mathcal{E},\mathcal{A}(X))$ (où
$\mathcal{A}(X)$ est considéré comme un
\emph{$\mathcal{O}_S$-Module})~; on peut par suite le munir canoniquement
d'une structure de \emph{module} sur l'anneau
$\operatorname{Hom}_S(X,S[T])$, et on voit comme ci-dessus que le couple
\[
  \bigl(\operatorname{Hom}_S(X,S[T]),
  \operatorname{Hom}_S(X,\mathbf{V}(\mathcal{E}))\bigr)
\]
est un foncteur contravariant en $X$, à valeurs dans la catégorie dont les
éléments sont les couples $(A,M)$ formés d'un anneau $A$ et d'un $A$-module
$M$, les morphismes étant les di-homomorphismes.

Nous interpréterons ces faits en disant que $S[T]$ est un
\emph{$S$-schéma d'anneaux} et que $\mathbf{V}(\mathcal{E})$ est un
\emph{$S$-schéma de modules} sur le $S$-schéma d'anneaux $S[T]$
(cf. chap.~0, \textsection8).
\end{env}

\begin{env}[1.7.14]
\label{II.1.7.14-fr}
Nous allons voir que la structure de $S$-schéma de modules définie sur le
$S$-schéma $\mathbf{V}(\mathcal{E})$ permet de reconstituer le
$\mathcal{O}_S$-Module $\mathcal{E}$ à un isomorphisme unique près~: pour
cela, nous allons montrer que $\mathcal{E}$ est canoniquement isomorphe à un
sous-$\mathcal{O}_S$-Module de
$\mathbf{S}(\mathcal{E})=\mathcal{A}(\mathbf{V}(\mathcal{E}))$, défini au
moyen de cette structure. En effet
(\hyperref[II.1.7.4-fr]{1.7.4}) l'ensemble
$\operatorname{Hom}_{\mathcal{O}_S}(\mathbf{S}(\mathcal{E}),
\mathcal{A}(X))$ des homomorphismes de
\emph{$\mathcal{O}_S$-Algèbres} s'identifie canoniquement à
$\operatorname{Hom}_{\mathcal{O}_S}(\mathcal{E},\mathcal{A}(X))$, ensemble
d'homomorphismes de \emph{$\mathcal{O}_S$-Modules}~: si $h$, $h'$ sont deux
éléments de ce dernier ensemble, $s_i$ ($1\leq i\leq n$) des sections de
$\mathcal{E}$ au-dessus d'un ouvert $U\subset S$, $t$ une section de
$\mathcal{A}(X)$ au-dessus de $U$, on a par définition
\[
  (h+h')(s_1s_2\cdots s_n)
  =\prod_{i=1}^n\bigl(h(s_i)+h'(s_i)\bigr)
\]
et
\[
  (t.h)(s_1s_2\cdots s_n)=t^n\prod_{i=1}^n h(s_i).
\]

Cela étant, si $z$ est une section de $\mathbf{S}(\mathcal{E})$ au-dessus
de $U$, $h\mapsto h(z)$ est une application de
$\operatorname{Hom}_S(X,\mathbf{V}(\mathcal{E}))
=\operatorname{Hom}_{\mathcal{O}_S}(\mathbf{S}(\mathcal{E}),
\mathcal{A}(X))$ dans $\Gamma(U,\mathcal{A}(X))$. Nous allons
\oldpage[II]{19}
montrer que $\mathcal{E}$ s'identifie au sous-Module de
$\mathbf{S}(\mathcal{E})$ tel que, \emph{pour tout ouvert $U\subset S$,
toute section $z$ de ce sous-$\mathcal{O}_S$-Module au-dessus de $U$ et tout
$S$-préschéma $X$, l'application $h\mapsto h(z)$ de
$\operatorname{Hom}_{\mathcal{O}_S|U}(\mathbf{S}(\mathcal{E})|U,
\mathcal{A}(X)|U)$ dans $\Gamma(U,\mathcal{A}(X))$ soit un homomorphisme de
$\Gamma(U,\mathcal{A}(X))$-modules}.

Il est immédiat en effet que $\mathcal{E}$ possède cette propriété~; pour
démontrer la réciproque, on peut se borner à prouver que lorsque
$S=\operatorname{Spec}(A)$, $\mathcal{E}=\widetilde{M}$, une section $z$ de
$\mathbf{S}(\mathcal{E})$ au-dessus de $S$ qui (pour $U=S$) possède la
propriété énoncée ci-dessus, est nécessairement une section de
$\mathcal{E}$~; on a alors $z=\sum_{n=0}^{\infty}z_n$, où
$z_n\in\mathbf{S}_n(M)$, et il s'agit de prouver que $z_n=0$ pour $n\ne1$.
Posons $B=\mathbf{S}(M)$ et prenons pour $X$ le préschéma
$\operatorname{Spec}(B[T])$, où $T$ est une indéterminée. L'ensemble
$\operatorname{Hom}_{\mathcal{O}_S}(\mathbf{S}(\mathcal{E}),
\mathcal{A}(X))$ s'identifie à l'ensemble des homomorphismes d'anneaux
$h:B\to B[T]$ (\hyperref[I.1.3.13-fr]{I, 1.3.13}), et d'après ce qu'on a vu
plus haut, on a
$(T.h)(z)=\sum_{n=0}^{\infty}T^nh(z_n)$~: l'hypothèse sur $z$ implique que
l'on a
$\sum_{n=0}^{\infty}T^nh(z_n)=T.\sum_{n=0}^{\infty}h(z_n)$ pour tout
homomorphisme $h$. Si on prend en particulier pour $h$ l'injection
canonique, il vient
$\sum_{n=0}^{\infty}T^nz_n=T.\sum_{n=0}^{\infty}z_n$, ce qui entraîne bien
la conclusion $z_n=0$ pour $n\ne1$.
\end{env}

\begin{proposition}[1.7.15]
\label{II.1.7.15-fr}
Soit $Y$ un préschéma dont l'espace sous-jacent est noethérien, ou un schéma
quasi-compact. Tout $Y$-schéma affine $X$ de type fini sur $Y$ est
$Y$-isomorphe à un sous-$Y$-schéma fermé d'un $Y$-schéma de la forme
$\mathbf{V}(\mathcal{E})$, où $\mathcal{E}$ est un $\mathcal{O}_Y$-Module
quasi-cohérent de type fini.
\end{proposition}

\begin{proof}
En effet, la $\mathcal{O}_Y$-Algèbre quasi-cohérente $\mathcal{A}(X)$ est de
type fini (\hyperref[II.1.3.7-fr]{1.3.7}). Les hypothèses faites entraînent
que $\mathcal{A}(X)$ est engendrée par un sous-$\mathcal{O}_Y$-Module
quasi-cohérent de type fini $\mathcal{E}$
(\hyperref[I.9.6.5-fr]{I, 9.6.5})~; par définition, cela entraîne que
l'homomorphisme canonique
$\mathbf{S}(\mathcal{E})\to\mathcal{A}(X)$ prolongeant canoniquement
l'injection $\mathcal{E}\to\mathcal{A}(X)$ est \emph{surjectif}~; la
conclusion résulte alors de (\hyperref[II.1.4.10-fr]{1.4.10}).
\end{proof}

\section{Spectres premiers homogènes}
\label{section:II.2-fr}

\subsection{Généralités sur les anneaux et modules gradués.}
\label{subsection:II.2.1-fr}

\begin{notation}[2.1.1]
\label{II.2.1.1-fr}
Étant donné un anneau $S$ \emph{gradué} en degrés positifs, nous désignerons
par $S_n$ la partie de $S$ formée des éléments homogènes de degré $n$
($n\geq0$), par $S_+$ la somme (directe) des $S_n$ tels que $n>0$~; on a
$1\in S_0$, $S_0$ est un sous-anneau de $S$, $S_+$ un idéal gradué de $S$ et
$S$ est somme directe de $S_0$ et de $S_+$. Si $M$ est un \emph{module
gradué} sur $S$ (à degrés positifs ou négatifs), on désignera de même par
$M_n$ le $S_0$-module formé des éléments homogènes de degré $n$ de $M$ (avec
$n\in\mathbf{Z}$).

Pour tout entier $d>0$, on désigne par $S^{(d)}$ la somme directe des
$S_{nd}$~; en considérant les éléments de $S_{nd}$ comme homogènes de degré
$n$, les $S_{nd}$ définissent sur $S^{(d)}$ une structure d'anneau gradué.

Pour tout entier $k$ tel que $0\leq k\leq d-1$, on désignera par
$M^{(d,k)}$ la somme directe
\oldpage[II]{20}
des $M_{nd+k}$ ($n\in\mathbf{Z}$)~; c'est un $S^{(d)}$-module gradué, quand
on considère les éléments de $M_{nd+k}$ comme homogènes de degré $n$. On
écrira $M^{(d)}$ au lieu de $M^{(d,0)}$.

Avec les notations précédentes, pour tout entier $n$ (positif ou négatif),
nous désignerons par $M(n)$ le $S$-module gradué défini par
$(M(n))_k=M_{n+k}$ pour tout $k\in\mathbf{Z}$. En particulier, $S(n)$ sera
un $S$-module gradué tel que $(S(n))_k=S_{n+k}$, en convenant de poser
$S_n=0$ pour $n<0$. On dit qu'un $S$-module gradué $M$ est \emph{libre}
s'il est isomorphe, en tant que \emph{module gradué}, à une somme directe de
modules de la forme $S(n)$~; comme $S(n)$ est un $S$-module monogène,
engendré par l'élément $1$ de $S$ considéré comme élément de degré $-n$, il
revient au même de dire que $M$ admet une \emph{base} sur $S$ formée
d'éléments \emph{homogènes}.

On dit qu'un $S$-module gradué $M$ \emph{admet une présentation finie} s'il
existe une suite exacte $P\to Q\to M\to0$, où $P$ et $Q$ sont des sommes
directes finies de modules de la forme $S(n)$ et les homomorphismes sont de
degré $0$ (cf. (\hyperref[II.2.1.2-fr]{2.1.2})).
\end{notation}

\begin{env}[2.1.2]
\label{II.2.1.2-fr}
Soient $M$, $N$ deux $S$-modules gradués~; on définit sur $M\otimes_S N$ une
structure de $S$-module \emph{gradué} de la façon suivante. Sur le produit
tensoriel $M\otimes_{\mathbf{Z}}N$, on peut définir une structure de
$\mathbf{Z}$-module gradué (où $\mathbf{Z}$ est gradué par
$\mathbf{Z}_0=\mathbf{Z}$, $\mathbf{Z}_n=0$ pour $n\ne0$) en posant
\[
  (M\otimes_{\mathbf{Z}}N)_q
  =\bigoplus_{m+n=q}M_m\otimes_{\mathbf{Z}}N_n
\]
(comme $M$ et $N$ sont respectivement sommes directes des $M_m$ et des
$N_n$, on sait que l'on peut identifier canoniquement
$M\otimes_{\mathbf{Z}}N$ à la somme directe de tous les
$M_m\otimes_{\mathbf{Z}}N_n$). Cela étant, on a
$M\otimes_S N=(M\otimes_{\mathbf{Z}}N)/P$, où $P$ est le
sous-$\mathbf{Z}$-module de $M\otimes_{\mathbf{Z}}N$ engendré par les
éléments $(xs)\otimes y-x\otimes(sy)$ pour $x\in M$, $y\in N$, $s\in S$~;
il est clair que $P$ est un sous-$\mathbf{Z}$-module \emph{gradué} de
$M\otimes_{\mathbf{Z}}N$, et on voit aussitôt qu'on obtient bien en passant
au quotient une structure de $S$-module gradué sur $M\otimes_S N$.

Pour deux $S$-modules gradués $M$, $N$, rappelons qu'un homomorphisme
$u:M\to N$ de $S$-modules est dit \emph{de degré $k$} si
$u(M_j)\subset N_{j+k}$ pour tout $j\in\mathbf{Z}$. Si $H_n$ désigne
l'ensemble de tous les homomorphismes de degré $n$ de $M$ dans $N$, on
désigne par $\operatorname{Hom}_S(M,N)$ la \emph{somme} (directe) des $H_n$
($n\in\mathbf{Z}$) dans le $S$-module $H$ de tous les homomorphismes (de
$S$-module) de $M$ dans $N$~; en général $\operatorname{Hom}_S(M,N)$ n'est
pas égal à ce dernier. Toutefois, on a $H=\operatorname{Hom}_S(M,N)$ lorsque
$M$ est \emph{de type fini}~; on peut en effet supposer alors $M$ engendré
par un nombre fini d'éléments homogènes $x_i$ ($1\leq i\leq n$), et tout
homomorphisme $u\in H$ s'écrit d'une seule manière
$\sum_{k\in\mathbf{Z}}u_k$, où pour chaque $k$, $u_k(x_i)$ est égal au
composant homogène de degré $k+\deg(x_i)$ de $u(x_i)$ ($1\leq i\leq n$), ce
qui entraîne $u_k=0$ sauf pour un nombre fini d'indices~; on a par définition
$u_k\in H_k$, d'où la conclusion.

On dit que les éléments de degré $0$ de $\operatorname{Hom}_S(M,N)$ sont des
\emph{homomorphismes de $S$-modules gradués}. Il est clair que
$S_mH_n\subset H_{m+n}$, donc les $H_n$ définissent sur
$\operatorname{Hom}_S(M,N)$ une structure de $S$-module gradué.

Il résulte aussitôt de ces définitions que l'on a
\[
  \label{II.2.1.2.1-fr}
  M(m)\otimes_S N(n)=(M\otimes_S N)(m+n),
  \tag{2.1.2.1}
\]
\[
  \label{II.2.1.2.2-fr}
  \operatorname{Hom}_S(M(m),N(n))
  =(\operatorname{Hom}_S(M,N))(n-m)
  \tag{2.1.2.2}
\]
pour deux $S$-modules gradués $M$, $N$.
\oldpage[II]{21}
Soient $S$, $S'$ deux anneaux gradués~; un homomorphisme d'\emph{anneaux
gradués} $\varphi:S\to S'$ est un homomorphisme d'anneaux tel que
$\varphi(S_n)\subset S'_n$, pour tout $n\in\mathbf{Z}$ (autrement dit,
$\varphi$ doit être un homomorphisme \emph{de degré $0$} de
$\mathbf{Z}$-modules gradués). La donnée d'un tel homomorphisme définit sur
$S'$ une structure de $S$-module \emph{gradué}~; muni de cette structure et
de sa structure d'anneau gradué, on dit que $S'$ est une \emph{$S$-algèbre
graduée}.

Si alors $M$ est un $S$-module gradué, le produit tensoriel $M\otimes_S S'$
de $S$-modules gradués est muni de façon naturelle d'une structure de
$S'$-module \emph{gradué}, la graduation étant celle définie plus haut.
\end{env}

\begin{lemma}[2.1.3]
\label{II.2.1.3-fr}
Soit $S$ un anneau gradué à degrés positifs. Pour qu'une partie $E$ de $S_+$
formée d'éléments homogènes engendre $S_+$ en tant que $S$-module, il faut et
il suffit que $E$ engendre $S$ en tant que $S_0$-algèbre.
\end{lemma}

\begin{proof}
La condition est évidemment suffisante~; montrons qu'elle est nécessaire.
Soit $E_n$ (resp. $E^n$) l'ensemble des éléments de $E$ de degré égal à $n$
(resp. $\leq n$)~; il suffira de montrer, par récurrence sur $n>0$, que $S_n$
est le $S_0$-module engendré par les éléments de degré $n$ qui sont des
produits d'éléments de $E^n$. Or, cela est évident pour $n=1$ en vertu de
l'hypothèse~; cette dernière montre en outre que
$S_n=\sum_{p=0}^{n-1}S_pE_{n-p}$, et le raisonnement par récurrence est alors
immédiat.
\end{proof}

\begin{corollary}[2.1.4]
\label{II.2.1.4-fr}
Pour que $S_+$ soit un idéal de type fini, il faut et il suffit que $S$ soit
une $S_0$-algèbre de type fini.
\end{corollary}

\begin{proof}
On peut en effet toujours supposer qu'un système fini de générateurs de la
$S_0$-algèbre $S$ (resp. du $S$-idéal $S_+$) est formé d'éléments homogènes,
en remplaçant chacun des générateurs considérés par ses composants homogènes.
\end{proof}

\begin{corollary}[2.1.5]
\label{II.2.1.5-fr}
Pour que $S$ soit noethérien, il faut et il suffit que $S_0$ soit noethérien
et que $S$ soit une $S_0$-algèbre de type fini.
\end{corollary}

\begin{proof}
La condition est évidemment suffisante~; elle est nécessaire, puisque $S_0$
est isomorphe à $S/S_+$ et que $S_+$ doit être un idéal de type fini
(\hyperref[II.2.1.4-fr]{2.1.4}).
\end{proof}

\begin{lemma}[2.1.6]
\label{II.2.1.6-fr}
Soit $S$ un anneau gradué à degrés positifs, qui soit une $S_0$-algèbre de
type fini. Soit $M$ un $S$-module gradué de type fini. Alors~:
\begin{enumerate}
  \item[\rm(i)] Les $M_n$ sont des $S_0$-modules de type fini, et il existe
    un entier $n_0$ tel que $M_n=0$ pour $n\leq n_0$.
  \item[\rm(ii)] Il existe un entier $n_1$ et un entier $h>0$ tel que, pour
    tout entier $n\geq n_1$, on ait $M_{n+h}=S_hM_n$.
  \item[\rm(iii)] Pour tout couple d'entiers $(d,k)$ tels que $d>0$,
    $0\leq k\leq d-1$, $M^{(d,k)}$ est un $S^{(d)}$-module de type fini.
  \item[\rm(iv)] Pour tout entier $d>0$, $S^{(d)}$ est une $S_0$-algèbre de
    type fini.
  \item[\rm(v)] Il existe un entier $h>0$ tel que $S_{mh}=(S_h)^m$ pour tout
    $m>0$.
  \item[\rm(vi)] Pour tout entier $n>0$, il existe un entier $m_0$ tel que
    $S_m\subset S_+^n$ pour tout $m\geq m_0$.
\end{enumerate}
\end{lemma}

\begin{proof}
On peut supposer $S$ engendré (en tant que $S_0$-algèbre) par des éléments
homogènes $f_i$, de degré $h_i$ ($1\leq i\leq r$), et $M$ engendré (en tant
que $S$-module) par des éléments homogènes $x_j$ de degré $k_j$
($1\leq j\leq s$). Il est clair que $M_n$ est formé des combinaisons
\oldpage[II]{22}
linéaires, à coefficients dans $S_0$, des éléments
$f_1^{\alpha_1}\cdots f_r^{\alpha_r}x_j$ tels que les $\alpha_i$ soient des
entiers $\geq 0$ vérifiant $k_j+\sum_i\alpha_i h_i=n$~; pour chaque $j$, il
n'y a qu'un nombre fini de systèmes $(\alpha_i)$ vérifiant cette équation,
puisque les $h_i$ sont $>0$, d'où la première assertion de (i)~; la seconde
est évidente. Soit d'autre part $h$ le p.p.c.m. des $h_i$ et posons
$g_i=f_i^{h/h_i}$ ($1\leq i\leq r$) de sorte que tous les $g_i$ sont de
degré $h$~; soient $z_\mu$ les éléments de $M$ de la forme
$f_1^{\alpha_1}\cdots f_r^{\alpha_r}x_j$ avec
$0\leq\alpha_i<h/h_i$ pour $1\leq i\leq r$~; ces éléments sont en nombre
fini, soit $n_1$ le plus grand de leurs degrés. Il est clair que pour
$n\geq n_1$, tout élément de $M_{n+h}$ est combinaison linéaire des $z_\mu$
dont les coefficients sont des monômes de degré $>0$ par rapport aux $g_i$,
donc on a $M_{n+h}=S_hM_n$, ce qui établit (ii). De la même manière, on voit
(pour tout $d>0$) qu'un élément de $M^{(d,k)}$ est combinaison linéaire, à
coefficients dans $S_0$, d'éléments de la forme
$g^df_1^{\alpha_1}\cdots f_r^{\alpha_r}x_j$ avec
$0\leq\alpha_i<d$, $g$ étant un élément homogène de $S$~; d'où (iii)~; (iv)
résulte alors de (iii) et du lemme (\hyperref[II.2.1.3-fr]{2.1.3}), en
prenant $M=S_+$, puisque $(S_+)^{(d)}=(S^{(d)})_+$. L'assertion de (v) se
déduit de (ii) en prenant $M=S$. Enfin, pour un $n$ donné, il n'y a qu'un
nombre fini de systèmes $(\alpha_i)$ tels que $\alpha_i\geq 0$ et
$\sum_i\alpha_i<n$, donc si $m_0$ est la plus grande valeur de la somme
$\sum_i\alpha_i h_i$ pour ces systèmes, on a $S_m\subset S_+^n$ pour
$m>m_0$, ce qui prouve (vi).
\end{proof}

\begin{corollary}[2.1.7]
\label{II.2.1.7-fr}
Si $S$ est noethérien, il en est de même de $S^{(d)}$ pour tout entier
$d>0$.
\end{corollary}

\begin{proof}
Cela résulte de (\hyperref[II.2.1.5-fr]{2.1.5}) et de
(\hyperref[II.2.1.6-fr]{2.1.6, (iv)}).
\end{proof}

\begin{env}[2.1.8]
\label{II.2.1.8-fr}
Soit $\mathfrak{p}$ un idéal premier gradué de l'anneau gradué $S$~;
$\mathfrak{p}$ est donc somme directe des sous-groupes
$\mathfrak{p}_n=\mathfrak{p}\cap S_n$. Supposons que $\mathfrak{p}$ ne
contienne pas $S_+$. Alors, si $f\in S_+$ n'appartient pas à
$\mathfrak{p}$, la relation $f^nx\in\mathfrak{p}$ est équivalente à
$x\in\mathfrak{p}$~; en particulier, si $f\in S_d$ ($d>0$), pour tout
$x\in S_{m-nd}$, la relation $f^nx\in\mathfrak{p}_m$ équivaut à
$x\in\mathfrak{p}_{m-nd}$.
\end{env}

\begin{proposition}[2.1.9]
\label{II.2.1.9-fr}
Soit $n_0$ un entier $>0$~; pour tout $n\geq n_0$ soit $\mathfrak{p}_n$ un
sous-groupe de $S_n$. Pour qu'il existe un idéal premier gradué
$\mathfrak{p}$ de $S$ ne contenant pas $S_+$ et tel que
$\mathfrak{p}\cap S_n=\mathfrak{p}_n$ pour tout $n\geq n_0$, il faut et il
suffit que les conditions suivantes soient vérifiées~:
\begin{enumerate}
  \item[$1^\circ$] $S_m\mathfrak{p}_n\subset\mathfrak{p}_{m+n}$ pour tout
    $m\geq 0$ et tout $n\geq n_0$.
  \item[$2^\circ$] Pour $m\geq n_0$, $n\geq n_0$, $f\in S_m$, $g\in S_n$,
    la relation $fg\in\mathfrak{p}_{m+n}$ entraîne
    $f\in\mathfrak{p}_m$ ou $g\in\mathfrak{p}_n$.
  \item[$3^\circ$] $\mathfrak{p}_n\neq S_n$ pour un $n\geq n_0$ au moins.
\end{enumerate}
En outre l'idéal premier gradué $\mathfrak{p}$ est alors unique.
\end{proposition}

\begin{proof}
Il est évident que les conditions $1^\circ$ et $2^\circ$ sont nécessaires.
En outre, si $\mathfrak{p}\not\supset S_+$, il existe au moins un $k>0$ tel que
$\mathfrak{p}\cap S_k\neq S_k$~; si $f\in S_k$ n'appartient pas à
$\mathfrak{p}$, la relation $\mathfrak{p}\cap S_n=S_n$ entraîne
$\mathfrak{p}\cap S_{n-mk}=S_{n-mk}$ d'après
(\hyperref[II.2.1.8-fr]{2.1.8})~; donc, si
$\mathfrak{p}\cap S_n=S_n$ à partir d'une certaine valeur de $n$, on aurait
$\mathfrak{p}\supset S_+$ contrairement à l'hypothèse, ce qui prouve que
$3^\circ$ est nécessaire. Inversement, supposons satisfaites les conditions
$1^\circ$, $2^\circ$ et $3^\circ$. Notons que, si pour un entier
$d\geq n_0$, $f\in S_d$ n'appartient pas à $\mathfrak{p}_d$, alors, si
$\mathfrak{p}$ existe, $\mathfrak{p}_m$, pour $m<n_0$, est nécessairement
égal à l'ensemble des $x\in S_m$ tels que
$f^rx\in\mathfrak{p}_{m+rd}$, sauf pour un nombre fini de valeurs de $r$.
Cela prouve déjà que si $\mathfrak{p}$ existe, il est unique. Reste à montrer
que si on définit les $\mathfrak{p}_m$ pour $m<n_0$ par la condition
précédente, $\mathfrak{p}=\sum_{n=0}^{\infty}\mathfrak{p}_n$ est un idéal
premier. Notons d'abord qu'en vertu de $2^\circ$, pour $m\geq n_0$,
$\mathfrak{p}_m$ est aussi défini comme l'ensemble des $x\in S_m$ tels que
$f^rx\in\mathfrak{p}_{m+rd}$, sauf pour un nombre fini de valeurs de $r$.
Cela
\oldpage[II]{23}
étant, si $g\in S_m$, $x\in\mathfrak{p}_n$, on a
$f^rgx\in\mathfrak{p}_{m+n+rd}$ sauf pour un nombre fini de valeurs de $r$,
donc $gx\in\mathfrak{p}_{m+n}$, ce qui prouve que $\mathfrak{p}$ est un
idéal de $S$. Pour établir que cet idéal est premier, autrement dit que
l'anneau $S/\mathfrak{p}$, gradué par les sous-groupes
$S_n/\mathfrak{p}_n$, est un anneau intègre, il suffit (en considérant les
composants de plus haut degré de deux éléments de $S/\mathfrak{p}$) de
prouver que si $x\in S_m$, $y\in S_n$ sont tels que
$x\notin\mathfrak{p}_m$, $y\notin\mathfrak{p}_n$, alors
$xy\notin\mathfrak{p}_{m+n}$. Sinon, pour $r$ assez grand, on aurait
$f^{2r}xy\in\mathfrak{p}_{m+n+2rd}$~; mais on a
$f^ry\notin\mathfrak{p}_{n+rd}$ pour tout $r>0$~; il résulte alors de
$2^\circ$ que, sauf pour un nombre fini de valeurs de $r$, on a
$f^rx\in\mathfrak{p}_{m+rd}$, et on en conclurait que
$x\in\mathfrak{p}_m$ contrairement à l'hypothèse.
\end{proof}

\begin{env}[2.1.10]
\label{II.2.1.10-fr}
Nous dirons qu'une partie $\mathfrak{J}$ de $S_+$ est un
\emph{idéal de $S_+$} si c'est un idéal de $S$, et que $\mathfrak{J}$ est un
\emph{idéal premier gradué de $S_+$} si elle est l'intersection de $S_+$ et
d'un idéal premier gradué de $S$ \emph{ne contenant pas $S_+$} (cet idéal
premier est d'ailleurs unique d'après
(\hyperref[II.2.1.9-fr]{2.1.9})). Si $\mathfrak{J}$ est un idéal de $S_+$,
la \emph{racine de $\mathfrak{J}$ dans $S_+$} est l'ensemble des éléments de
$S_+$ dont une puissance appartient à $\mathfrak{J}$, autrement dit c'est
l'ensemble $r_+(\mathfrak{J})=r(\mathfrak{J})\cap S_+$~; en particulier, la
racine de $0$ dans $S_+$ est encore appelée le \emph{nilradical} de $S_+$ et
notée $\mathfrak{N}_+$~: c'est donc l'ensemble des éléments nilpotents de
$S_+$. Si $\mathfrak{J}$ est un idéal \emph{gradué} de $S_+$, sa racine
$r_+(\mathfrak{J})$ est un idéal gradué~: en passant à l'anneau quotient
$S/\mathfrak{J}$, on peut se ramener au cas $\mathfrak{J}=0$, et tout
revient à voir que si $x=x_h+x_{h+1}+\cdots+x_k$ est nilpotent, il en est de
même des $x_i\in S_i$ ($1\leq h\leq i\leq k$)~; on peut supposer
$x_k\neq 0$ et le composant de plus grand degré de $x^n$ est alors $x_k^n$,
donc $x_k$ est nilpotent, et on raisonne alors par récurrence sur $k$. On dit
que l'anneau gradué $S$ est \emph{essentiellement réduit} si
$\mathfrak{N}_+=0$, autrement dit si $S_+$ ne contient pas d'éléments
nilpotents $\neq 0$.
\end{env}

\begin{env}[2.1.11]
\label{II.2.1.11-fr}
On notera que si, dans l'anneau gradué $S$, un élément $x$ est diviseur de
$0$, il en est de même de son composant de plus haut degré. On dit qu'un
anneau $S$ est \emph{essentiellement intègre} si l'anneau $S_+$
(\emph{sans élément unité}) ne contient pas de diviseur de $0$ et est
$\neq 0$~; il suffit pour cela qu'un élément homogène $\neq 0$ dans $S_+$ ne
soit pas diviseur de $0$ dans cet anneau. Il est clair que si $\mathfrak{p}$
est un idéal premier gradué de $S_+$, $S/\mathfrak{p}$ est essentiellement
intègre.

Soit $S$ un anneau gradué essentiellement intègre, et soit $x_0\in S_0$~;
s'il existe alors \emph{un} élément homogène $f\neq 0$ de $S_+$ tel que
$x_0f=0$, on a $x_0S_+=0$, car on a $(x_0g)f=(x_0f)g=0$ pour tout
$g\in S_+$, et l'hypothèse entraîne donc $x_0g=0$. Pour que $S$ soit
intègre, il faut et il suffit donc que $S_0$ soit intègre et que l'annulateur
de $S_+$ dans $S_0$ soit réduit à $0$.
\end{env}

\subsection{Anneaux de fractions d'un anneau gradué.}
\label{subsection:II.2.2-fr}

\begin{env}[2.2.1]
\label{II.2.2.1-fr}
Soient $S$ un anneau gradué, à degrés positifs, $f$ un élément
\emph{homogène} de $S$, de degré $d>0$~; alors l'anneau de fractions
$S'=S_f$ est gradué, en prenant pour $S'_n$ l'ensemble des $x/f^k$, où
$x\in S_{n+kd}$ avec $k\geq 0$ (on observera ici que $n$ peut prendre des
valeurs négatives arbitraires)~; nous désignerons le sous-anneau
$S'_0=(S_f)_0$ de $S'$ formé des éléments \emph{de degré $0$} par la
notation $S_{(f)}$.

Si $f\in S_d$, les monômes $(f/1)^h$ dans $S_f$ ($h$ entier positif ou
négatif) forment un \emph{système libre} sur l'anneau $S_{(f)}$, et
l'ensemble de leurs combinaisons linéaires n'est autre que
\oldpage[II]{24}
l'anneau $(S^{(d)})_f$, qui est donc \emph{isomorphe à
$S_{(f)}[T,T^{-1}]=S_{(f)}\otimes_{\mathbf{Z}}\mathbf{Z}[T,T^{-1}]$}
(où $T$ est une indéterminée). En effet, si on a une relation
$\sum_{h=-a}^{b}z_h(f/1)^h=0$ avec $z_h=x_h/f^m$, où tous les $x_h$
appartiennent à $S_{md}$, cette relation équivaut par définition à
l'existence d'un $k\geq a$ tel que
$\sum_{h=-a}^{b}f^{h+k}x_h=0$, et comme les degrés des termes de cette somme
sont distincts, on a $f^{h+k}x_h=0$ pour tout $h$, d'où $z_h=0$ pour tout
$h$.

Si $M$ est un $S$-module gradué, $M'=M_f$ est un $S_f$-module gradué,
$M'_n$ étant l'ensemble des $z/f^k$ avec $z\in M_{n+kd}$ ($k\geq 0$)~; on
désigne par $M_{(f)}$ l'ensemble des éléments homogènes de degré $0$ de
$M'$~; il est immédiat que $M_{(f)}$ est un $S_{(f)}$-module et que l'on a
$(M^{(d)})_f=M_{(f)}\otimes_{S_{(f)}}(S^{(d)})_f$.
\end{env}

\begin{lemma}[2.2.2]
\label{II.2.2.2-fr}
Soient $d$, $e$ deux entiers $>0$, $f\in S_d$, $g\in S_e$. Il existe un
isomorphisme canonique d'anneaux
\[
  S_{(fg)}\xrightarrow{\sim}(S_{(f)})_{g^d/f^e}~;
\]
si on identifie canoniquement ces deux anneaux, il existe un isomorphisme
canonique de modules
\[
  M_{(fg)}\xrightarrow{\sim}(M_{(f)})_{g^d/f^e}.
\]
\end{lemma}

\begin{proof}
En effet, $fg$ divise $f^e g^d$, et ce dernier élément divise $(fg)^{de}$,
donc les anneaux gradués $S_{fg}$ et $S_{f^e g^d}$ s'identifient
canoniquement~; d'autre part, $S_{f^e g^d}$ s'identifie aussi à
$(S_{f^e})_{g^d/1}$ (\hyperref[0.1.4.6-fr]{0, I.4.6}), et comme $f^e/1$ est
inversible dans $S_{f^e}$, $S_{f^e g^d}$ s'identifie aussi à
$(S_{f^e})_{g^d/f^e}$. Or l'élément $g^d/f^e$ est de degré $0$ dans
$S_{f^e}$~; on en conclut aussitôt que le sous-anneau de
$(S_{f^e})_{g^d/f^e}$ formé des éléments de degré $0$ est
$(S_{(f^e)})_{g^d/f^e}$, et comme on a évidemment
$S_{(f^e)}=S_{(f)}$, cela démontre la première partie de la proposition~;
la seconde s'établit de même.
\end{proof}

\begin{env}[2.2.3]
\label{II.2.2.3-fr}
Avec les hypothèses de (\hyperref[II.2.2.2-fr]{2.2.2}), il est clair que
l'homomorphisme canonique $S_f\to S_{fg}$
(\hyperref[0.1.4.1-fr]{0, I.4.1}), qui transforme $x/f^k$ en
$g^k x/(fg)^k$ est de degré $0$, donc donne par restriction un
\emph{homomorphisme canonique $S_{(f)}\to S_{(fg)}$}, tel que le diagramme
\[
  \xymatrix{
    & S_{(f)}\ar[dl]\ar[dr]\\
    S_{(fg)}\ar[rr]^-{\sim} && (S_{(f)})_{g^d/f^e}
  }
\]
soit commutatif. On définit de même un homomorphisme canonique
$M_{(f)}\to M_{(fg)}$.
\end{env}

\begin{lemma}[2.2.4]
\label{II.2.2.4-fr}
Si $f$, $g$ sont deux éléments homogènes de $S_+$, l'anneau $S_{(fg)}$ est
engendré par la réunion des images canoniques de $S_{(f)}$ et $S_{(g)}$.
\end{lemma}

\begin{proof}
En vertu de (\hyperref[II.2.2.2-fr]{2.2.2}), il suffit de voir que
$1/(g^d/f^e)=f^{d+e}/(fg)^d$ appartient à l'image canonique de $S_{(g)}$
dans $S_{(fg)}$, ce qui est évident par définition.
\end{proof}

\begin{proposition}[2.2.5]
\label{II.2.2.5-fr}
Soit $d$ un entier $>0$ et soit $f\in S_d$. Alors il existe un isomorphisme
canonique d'anneaux
$S_{(f)}\xrightarrow{\sim}S^{(d)}/(f-1)S^{(d)}$~; si on identifie ces deux
anneaux par cet isomorphisme, il existe un isomorphisme canonique de modules
$M_{(f)}\xrightarrow{\sim}M^{(d)}/(f-1)M^{(d)}$.
\end{proposition}

\begin{proof}
Le premier de ces isomorphismes se définit en faisant correspondre à
$x/f^n$, où $x\in S_{nd}$, l'élément $\overline{x}$, classe de $x$ mod.
$(f-1)S^{(d)}$~; cette application est bien définie, car on a la congruence
$f^h x\equiv x\pmod{(f-1)S^{(d)}}$ pour tout $x\in S^{(d)}$, donc si
$f^h x=0$ pour un $h>0$,
\oldpage[II]{25}
on a $\overline{x}=0$. D'autre part, si $x\in S_{nd}$ est tel que
$x=(f-1)y$ avec
$y=y_{hd}+y_{(h+1)d}+\cdots+y_{kd}$ avec $y_{jd}\in S_{jd}$ et
$y_{hd}\neq 0$, on a nécessairement $h=n$ et $x=-y_{hd}$, ainsi que les
relations $y_{(j+1)d}=fy_{jd}$ pour $h\leq j\leq k-1$, $fy_{kd}=0$, ce qui
donne finalement $f^{k-n+1}x=0$~; à toute classe $\overline{x}$ mod.
$(f-1)S^{(d)}$ d'un élément $x\in S_{nd}$, on peut donc faire correspondre
l'élément $x/f^n$ de $S_{(f)}$, car la remarque qui précède montre que cette
application est bien définie. Il est immédiat que les deux applications ainsi
définies sont des homomorphismes d'anneaux, réciproques l'un de l'autre. On
procède exactement de même pour $M$.
\end{proof}

\begin{corollary}[2.2.6]
\label{II.2.2.6-fr}
Si $S$ est noethérien, il en est de même de $S_{(f)}$ pour tout $f$ homogène
de degré $>0$.
\end{corollary}

\begin{proof}
Cela résulte aussitôt de (\hyperref[II.2.1.7-fr]{2.1.7}) et
(\hyperref[II.2.2.5-fr]{2.2.5}).
\end{proof}

\begin{env}[2.2.7]
\label{II.2.2.7-fr}
Soit $T$ une partie multiplicative de $S_+$ formée d'éléments
\emph{homogènes}~; $T_0=T\cup\{1\}$ est alors une partie multiplicative de
$S$~; comme les éléments de $T_0$ sont homogènes, l'anneau $T_0^{-1}S$ est
encore gradué de façon évidente~; on désignera par $S_{(T)}$ le sous-anneau
de $T_0^{-1}S$ formé des éléments d'ordre $0$, c'est-à-dire des éléments de
la forme $x/h$, où $h\in T$ et $x$ est homogène de degré égal à celui de
$h$. On sait (\hyperref[0.1.4.5-fr]{0, 1.4.5}) que $T_0^{-1}S$ s'identifie
canoniquement à la limite inductive des anneaux $S_f$, où $f$ parcourt $T$
(pour les homomorphismes canoniques $S_f\to S_{fg}$)~; comme cette
identification respecte les degrés, elle identifie $S_{(T)}$ à la
\emph{limite inductive} des $S_{(f)}$ pour $f\in T$. Pour tout $S$-module
gradué $M$, on définit de même le module $M_{(T)}$ (sur l'anneau $S_{(T)}$)
formé des éléments de degré $0$ de $T_0^{-1}M$, et on voit que ce module est
limite inductive des $M_{(f)}$ pour $f\in T$.

Si $\mathfrak{p}$ est un idéal premier gradué de $S_+$, on désignera par
$S_{(\mathfrak{p})}$ et $M_{(\mathfrak{p})}$ l'anneau $S_{(T)}$ et le
module $M_{(T)}$ respectivement, où $T$ est l'ensemble des éléments
\emph{homogènes} de $S_+$ qui n'appartiennent pas à $\mathfrak{p}$.
\end{env}

\subsection{Spectre premier homogène d'un anneau gradué.}
\label{subsection:II.2.3-fr}

\begin{env}[2.3.1]
\label{II.2.3.1-fr}
Étant donné un anneau gradué $S$, à degrés positifs, on appelle
\emph{spectre premier homogène} de $S$ et on désigne par
$\operatorname{Proj}(S)$ l'ensemble des idéaux premiers gradués de $S_+$
(\hyperref[II.2.1.10-fr]{2.1.10}), ou, ce qui revient au même, l'ensemble
des idéaux premiers gradués de $S$ \emph{ne contenant pas} $S_+$~; nous
allons définir une structure de \emph{schéma} ayant
$\operatorname{Proj}(S)$ comme ensemble de base.
\end{env}

\begin{env}[2.3.2]
\label{II.2.3.2-fr}
Pour toute partie $E$ de $S$, soit $V_+(E)$ l'ensemble des idéaux premiers
gradués de $S$ contenant $E$ et ne contenant pas $S_+$~; c'est donc la
partie $V(E)\cap\operatorname{Proj}(S)$ de
$\operatorname{Spec}(S)$. De (\hyperref[I.1.1.2-fr]{I, 1.1.2}) on déduit
donc~:
\[
  \label{II.2.3.2.1-fr}
  V_+(0)=\operatorname{Proj}(S),\quad
  V_+(S)=V_+(S_+)=\varnothing
  \tag{2.3.2.1}
\]
\[
  \label{II.2.3.2.2-fr}
  V_+\left(\bigcup_\lambda E_\lambda\right)
  =\bigcap_\lambda V_+(E_\lambda)
  \tag{2.3.2.2}
\]
\[
  \label{II.2.3.2.3-fr}
  V_+(EE')=V_+(E)\cup V_+(E')
  \tag{2.3.2.3}
\]
On ne change pas $V_+(E)$ en remplaçant $E$ par l'idéal gradué engendré par
$E$~; en outre, si $\mathfrak{J}$ est un idéal gradué de $S$, on a
\[
  \label{II.2.3.2.4-fr}
  V_+(\mathfrak{J})
  =V_+\left(\bigcup_{q\geq n}(\mathfrak{J}\cap S_q)\right)
  \tag{2.3.2.4}
\]
\oldpage[II]{26}
pour tout $n>0$~: en effet, si
$\mathfrak{p}\in\operatorname{Proj}(S)$ contient les éléments homogènes de
$\mathfrak{J}$ de degré $\geq n$, comme par hypothèse il existe un élément
homogène $f\in S_d$ non contenu dans $\mathfrak{p}$, pour tout $m\geq 0$ et
tout $x\in S_m\cap\mathfrak{J}$, on a
$f^r x\in\mathfrak{J}\cap S_{m+rd}$ sauf pour un nombre fini de valeurs de
$r$, donc $f^r x\in\mathfrak{p}\cap S_{m+rd}$, ce qui entraîne
$x\in\mathfrak{p}\cap S_m$
(\hyperref[II.2.1.9-fr]{2.1.9}).

Enfin, on a, pour tout idéal gradué $\mathfrak{J}$ de $S$
\[
  \label{II.2.3.2.5-fr}
  V_+(\mathfrak{J})=V_+(r_+(\mathfrak{J})).
  \tag{2.3.2.5}
\]
\end{env}

\begin{env}[2.3.3]
\label{II.2.3.3-fr}
Par définition, les $V_+(E)$ sont les parties fermées de
$X=\operatorname{Proj}(S)$ pour la topologie induite par la topologie
spectrale de $\operatorname{Spec}(S)$, et qu'on appellera encore
\emph{topologie spectrale} sur $X$. On pose, pour tout $f\in S$
\[
  \label{II.2.3.3.1-fr}
  D_+(f)=D(f)\cap\operatorname{Proj}(S)
  =\operatorname{Proj}(S)\setminus V_+(f)
  \tag{2.3.3.1}
\]
et on a par suite, pour deux éléments $f$, $g$ de $S$
(\hyperref[I.1.1.9.1-fr]{I, 1.1.9.1})
\[
  \label{II.2.3.3.2-fr}
  D_+(fg)=D_+(f)\cap D_+(g).
  \tag{2.3.3.2}
\]
\end{env}

\begin{proposition}[2.3.4]
\label{II.2.3.4-fr}
Lorsque $f$ parcourt l'ensemble des éléments homogènes de $S_+$, les
$D_+(f)$ forment une base de la topologie de
$X=\operatorname{Proj}(S)$.
\end{proposition}

\begin{proof}
Il résulte en effet de (\hyperref[II.2.3.2.2-fr]{2.3.2.2}) et
(\hyperref[II.2.3.2.4-fr]{2.3.2.4}) que toute partie fermée de $X$ est
intersection d'ensembles de la forme $V_+(f)$, où $f$ est homogène de degré
$>0$.
\end{proof}

\begin{env}[2.3.5]
\label{II.2.3.5-fr}
Soit $f$ un élément \emph{homogène} de $S_+$, de degré $d>0$~; pour tout
idéal premier gradué $\mathfrak{p}$ de $S$, ne contenant pas $f$, on sait
que l'ensemble des $x/f^n$, où $x\in\mathfrak{p}$ et $n\geq 0$, est un
idéal premier de l'anneau de fractions $S_f$
(\hyperref[0.1.2.6-fr]{0, 1.2.6})~; sa trace sur $S_{(f)}$ est donc un idéal
premier de cet anneau, que nous désignerons par
$\psi_f(\mathfrak{p})$~: c'est l'ensemble des $x/f^n$ pour $n\geq 0$,
$x\in\mathfrak{p}\cap S_{nd}$. On a ainsi défini une application
\[
  \psi_f:D_+(f)\longrightarrow\operatorname{Spec}(S_{(f)})~;
\]
en outre, si $g\in S_e$ est un second élément homogène de $S_+$, on a un
diagramme commutatif
\[
  \label{II.2.3.5.1-fr}
  \xymatrix{
    D_+(f)\ar[r]^{\psi_f}
    & \operatorname{Spec}(S_{(f)})\\
    D_+(fg)\ar[u]\ar[r]_{\psi_{fg}}
    & \operatorname{Spec}(S_{(fg)})\ar[u]
  }
  \tag{2.3.5.1}
\]
où la flèche verticale de gauche est l'inclusion, et celle de droite est
l'application ${}^a\omega_{fg,f}$ déduite de l'homomorphisme canonique
$\omega=\omega_{fg,f}:S_{(f)}\to S_{(fg)}$
(\hyperref[I.1.2.1-fr]{I, 1.2.1}). En effet, si
$x/f^n\in\omega^{-1}(\psi_{fg}(\mathfrak{p}))$, où
$fg\notin\mathfrak{p}$, on a par définition
$g^n x/(fg)^n\in\psi_{fg}(\mathfrak{p})$, donc
$g^n x\in\mathfrak{p}$ et par suite $x\in\mathfrak{p}$, et la réciproque
est évidente.
\end{env}

\begin{proposition}[2.3.6]
\label{II.2.3.6-fr}
L'application $\psi_f$ est un homéomorphisme de $D_+(f)$ sur
$\operatorname{Spec}(S_{(f)})$.
\end{proposition}

\begin{proof}
En premier lieu, $\psi_f$ est continue~; car si $h\in S_{nd}$ est tel que
$h/f^n\in\psi_f(\mathfrak{p})$, on a par définition
$h\in\mathfrak{p}$ et réciproquement, donc
\[
  \psi_f^{-1}(D(h/f^n))=D_+(hf)
\]
et notre assertion résulte de la formule
(\hyperref[II.2.3.3.2-fr]{2.3.3.2}). En outre, les $D_+(hf)$, où $h$
parcourt les ensembles $S_{nd}$, forment une base de la topologie de
$D_+(f)$, en vertu de (\hyperref[II.2.3.4-fr]{2.3.4}) et de la formule
(\hyperref[II.2.3.3.2-fr]{2.3.3.2})~; ce
\oldpage[II]{27}
qui précède prouve donc, compte tenu de l'axiome $(T_0)$ valable dans
$D_+(f)$ et dans $\operatorname{Spec}(S_{(f)})$, que $\psi_f$ est injective
et que l'application réciproque
$\psi_f(D_+(f))\longrightarrow D_+(f)$ est continue. Enfin, pour voir que
$\psi_f$ est surjective, on remarque que, si $\mathfrak{q}_0$ est un idéal
premier de $S_{(f)}$, et si, pour tout $n>0$, on désigne par
$\mathfrak{p}_n$ l'ensemble des $x\in S_n$ tels que
$x^d/f^n\in\mathfrak{q}_0$, les $\mathfrak{p}_n$ vérifient les conditions
de (\hyperref[II.2.1.9-fr]{2.1.9})~: en effet, si $x\in S_n$, $y\in S_n$
sont tels que $x^d/f^n\in\mathfrak{q}_0$ et
$y^d/f^n\in\mathfrak{q}_0$, on a
$(x+y)^{2d}/f^{2n}\in\mathfrak{q}_0$, d'où
$(x+y)^d/f^n\in\mathfrak{q}_0$ puisque $\mathfrak{q}_0$ est premier~;
cela prouve que les $\mathfrak{p}_n$ sont des sous-groupes des $S_n$, et la
vérification des autres conditions de
(\hyperref[II.2.1.9-fr]{2.1.9}) est immédiate, compte tenu de ce que
$\mathfrak{q}_0$ est premier. Si $\mathfrak{p}$ est l'idéal premier gradué
de $S$ ainsi défini, on a bien
$\psi_f(\mathfrak{p})=\mathfrak{q}_0$, car si $x\in S_{nd}$, les relations
$x/f^n\in\mathfrak{q}_0$ et $x^d/f^{nd}\in\mathfrak{q}_0$ sont
équivalentes, $\mathfrak{q}_0$ étant premier.
\end{proof}

\begin{corollary}[2.3.7]
\label{II.2.3.7-fr}
Pour que $D_+(f)=\varnothing$, il faut et il suffit que $f$ soit nilpotent.
\end{corollary}

\begin{proof}
En effet, pour que $\operatorname{Spec}(S_{(f)})=\varnothing$, il faut et
il suffit que $S_{(f)}=0$, ou encore que $1=0$ dans $S_f$, ce qui signifie
par définition que $f$ est nilpotent.
\end{proof}

\begin{corollary}[2.3.8]
\label{II.2.3.8-fr}
Soit $E$ une partie de $S_+$. Les conditions suivantes sont équivalentes~:
\begin{enumerate}
  \item[\rm a)] $V_+(E)=X=\operatorname{Proj}(S)$.
  \item[\rm b)] Tout élément de $E$ est nilpotent.
  \item[\rm c)] Les composants homogènes de tout élément de $E$ sont
    nilpotents.
\end{enumerate}
\end{corollary}

\begin{proof}
Il est clair que c) entraîne b) et que b) entraîne a). Si $\mathfrak{J}$ est
l'idéal gradué de $S$ engendré par $E$, la condition a) équivaut à
$V_+(\mathfrak{J})=X$~; \emph{a fortiori}, a) entraîne que tout élément
homogène $f\in\mathfrak{J}$ est tel que $V_+(f)=X$, donc $f$ est nilpotent
en vertu de (\hyperref[II.2.3.7-fr]{2.3.7}).
\end{proof}

\begin{corollary}[2.3.9]
\label{II.2.3.9-fr}
Si $\mathfrak{J}$ est un idéal gradué de $S_+$,
$r_+(\mathfrak{J})$ est l'intersection des idéaux premiers gradués de
$S_+$ qui contiennent $\mathfrak{J}$.
\end{corollary}

\begin{proof}
En considérant l'anneau gradué $S/\mathfrak{J}$, on peut se ramener au cas
où $\mathfrak{J}=0$. Il faut prouver que si $f\in S_+$ n'est pas
nilpotent, il y a un idéal premier gradué de $S$ ne contenant pas $f$~; or,
un au moins des composants homogènes de $f$ n'est pas nilpotent, et on peut
donc supposer $f$ homogène~; l'assertion résulte alors de
(\hyperref[II.2.3.7-fr]{2.3.7}).
\end{proof}

\begin{env}[2.3.10]
\label{II.2.3.10-fr}
Pour toute partie $Y$ de $X=\operatorname{Proj}(S)$, soit
$\mathfrak{j}_+(Y)$ l'ensemble des $f\in S_+$ tels que
$Y\subset V_+(f)$~; il revient au même de dire que
$\mathfrak{j}_+(Y)=\mathfrak{j}(Y)\cap S_+$~; $\mathfrak{j}_+(Y)$ est
donc un idéal de $S_+$, égal à sa racine dans $S_+$.
\end{env}

\begin{proposition}[2.3.11]
\label{II.2.3.11-fr}
\begin{enumerate}
  \item[\rm(i)] Pour toute partie $E$ de $S_+$,
    $\mathfrak{j}_+(V_+(E))$ est la racine dans $S_+$ de l'idéal gradué de
    $S_+$ engendré par $E$.
  \item[\rm(ii)] Pour toute partie $Y$ de $X$,
    $V_+(\mathfrak{j}_+(Y))=\overline{Y}$, adhérence de $Y$ dans $X$.
\end{enumerate}
\end{proposition}

\begin{proof}
\begin{enumerate}
  \item[\rm(i)] Si $\mathfrak{J}$ est l'idéal gradué de $S_+$ engendré
    par $E$, on a $V_+(E)=V_+(\mathfrak{J})$ et l'assertion résulte alors de
    (\hyperref[II.2.3.9-fr]{2.3.9}).
  \item[\rm(ii)] Comme
    $V_+(\mathfrak{J})=\bigcap_{f\in\mathfrak{J}}V_+(f)$, la relation
    $Y\subset V_+(\mathfrak{J})$ entraîne $Y\subset V_+(f)$ pour tout
    $f\in\mathfrak{J}$, et par suite
    $\mathfrak{j}_+(Y)\supset\mathfrak{J}$, d'où
    $V_+(\mathfrak{j}_+(Y))\subset V_+(\mathfrak{J})$, ce qui prouve (ii)
    par définition des ensembles fermés.
\end{enumerate}
\end{proof}

\begin{corollary}[2.3.12]
\label{II.2.3.12-fr}
Les parties fermées $Y$ de $X=\operatorname{Proj}(S)$ et les idéaux gradués
de $S_+$ égaux à leur racine dans $S_+$ se correspondent biunivoquement par
les applications décroissantes
$Y\longmapsto\mathfrak{j}_+(Y)$, $\mathfrak{J}\longmapsto
V_+(\mathfrak{J})$~; à la réunion $Y_1\cup Y_2$ de deux parties fermées de
$X$ correspond
\oldpage[II]{28}
$\mathfrak{j}_+(Y_1)\cap\mathfrak{j}_+(Y_2)$, et à l'intersection d'une
famille quelconque $(Y_\lambda)$ de parties fermées correspond la racine
dans $S_+$ de la somme des $\mathfrak{j}_+(Y_\lambda)$.
\end{corollary}

\begin{corollary}[2.3.13]
\label{II.2.3.13-fr}
Soit $\mathfrak{J}$ un idéal gradué dans $S_+$~; pour que
$V_+(\mathfrak{J})=\varnothing$, il faut et il suffit que tout élément de
$S_+$ ait une puissance dans $\mathfrak{J}$.
\end{corollary}

Ce dernier corollaire s'exprime aussi sous l'une des formes équivalentes~:

\begin{corollary}[2.3.14]
\label{II.2.3.14-fr}
Soit $(f_\alpha)$ une famille d'éléments homogènes de $S_+$. Pour que les
$D_+(f_\alpha)$ forment un recouvrement de
$X=\operatorname{Proj}(S)$, il faut et il suffit que tout élément de $S_+$
ait une puissance dans l'idéal engendré par les $f_\alpha$.
\end{corollary}

\begin{corollary}[2.3.15]
\label{II.2.3.15-fr}
Soient $(f_\alpha)$ une famille d'éléments homogènes de $S_+$, $f$ un
élément de $S_+$. Les relations suivantes sont équivalentes~:
\begin{enumerate}
  \item[\rm a)] $D_+(f)\subset\bigcup_\alpha D_+(f_\alpha)$~;
  \item[\rm b)] $V_+(f)\supset\bigcap_\alpha V_+(f_\alpha)$~;
  \item[\rm c)] une puissance de $f$ appartient à l'idéal engendré par les
    $f_\alpha$.
\end{enumerate}
\end{corollary}

\begin{corollary}[2.3.16]
\label{II.2.3.16-fr}
Pour que $X=\operatorname{Proj}(S)$ soit vide, il faut et il suffit que tout
élément de $S_+$ soit nilpotent.
\end{corollary}

\begin{corollary}[2.3.17]
\label{II.2.3.17-fr}
Dans la correspondance biunivoque décrite dans
(\hyperref[II.2.3.12-fr]{2.3.12}), aux parties fermées irréductibles de $X$
correspondent les idéaux gradués premiers dans $S_+$.
\end{corollary}

\begin{proof}
En effet, si $Y=Y_1\cup Y_2$, où $Y_1$ et $Y_2$ sont fermés et distincts de
$Y$, on a
\[
  \mathfrak{j}_+(Y)
  =\mathfrak{j}_+(Y_1)\cap\mathfrak{j}_+(Y_2)
\]
les idéaux $\mathfrak{j}_+(Y_1)$ et $\mathfrak{j}_+(Y_2)$ étant distincts
de $\mathfrak{j}_+(Y)$, donc $\mathfrak{j}_+(Y)$ n'est pas premier.
Inversement, si $\mathfrak{J}$ est un idéal gradué non premier de $S_+$, il
existe deux éléments $f$, $g$ de $S_+$ tels que
$f\notin\mathfrak{J}$, $g\notin\mathfrak{J}$, $fg\in\mathfrak{J}$~; alors
$V_+(f)\not\supset V_+(\mathfrak{J})$,
$V_+(g)\not\supset V_+(\mathfrak{J})$ mais
$V_+(\mathfrak{J})\subset V_+(f)\cup V_+(g)$ par
(\hyperref[II.2.3.2.3-fr]{2.3.2.3})~; on en conclut que
$V_+(\mathfrak{J})$ est réunion des ensembles fermés
$V_+(f)\cap V_+(\mathfrak{J})$ et $V_+(g)\cap V_+(\mathfrak{J})$, qui sont
distincts de $V_+(\mathfrak{J})$.
\end{proof}

\subsection{La structure de schéma sur $\operatorname{Proj}(S)$.}
\label{subsection:II.2.4-fr}

\begin{env}[2.4.1]
\label{II.2.4.1-fr}
Soient $f$, $g$ deux éléments homogènes de $S_+$~; considérons les schémas
affines $Y_f=\operatorname{Spec}(S_{(f)})$,
$Y_g=\operatorname{Spec}(S_{(g)})$ et
$Y_{fg}=\operatorname{Spec}(S_{(fg)})$. En vertu de
(\hyperref[II.2.2.2-fr]{2.2.2}), le morphisme
$w_{fg,f}=({}^a\omega_{fg,f},\widetilde{\omega}_{fg,f})$ de $Y_{fg}$ dans
$Y_f$, correspondant à l'homomorphisme canonique
$\omega_{fg,f}:S_{(f)}\to S_{(fg)}$, est une \emph{immersion ouverte}
(\hyperref[I.1.3.6-fr]{I, 1.3.6}). Au moyen de l'homéomorphisme réciproque
de $\psi_f:D_+(f)\to Y_f$ (\hyperref[II.2.3.6-fr]{2.3.6}), on peut
transporter à $D_+(f)$ la structure de schéma affine de $Y_f$~; en vertu de
la commutativité du diagramme (\hyperref[II.2.3.5.1-fr]{2.3.5.1}), le
schéma affine $D_+(fg)$ s'identifie ainsi au schéma induit sur l'ensemble
ouvert $D_+(fg)$ de l'espace sous-jacent au schéma affine $D_+(f)$. Il est
clair alors (compte tenu de (\hyperref[II.2.3.4-fr]{2.3.4})) que
$X=\operatorname{Proj}(S)$ est muni d'une unique structure de
\emph{préschéma}, dont la restriction à chaque $D_+(f)$ est le schéma
affine qui vient d'être défini. En outre~:
\end{env}

\begin{proposition}[2.4.2]
\label{II.2.4.2-fr}
Le préschéma $\operatorname{Proj}(S)$ est un schéma.
\end{proposition}

\begin{proof}
Il suffit (\hyperref[I.5.5.6-fr]{I, 5.5.6}) de montrer que, quels que soient
$f$, $g$ homogènes dans $S_+$,
$D_+(f)\cap D_+(g)=D_+(fg)$ est affine et que son anneau est engendré par
les images canoniques des anneaux de $D_+(f)$ et $D_+(g)$~; le premier
point est évident par définition, et le second résulte de
(\hyperref[II.2.2.4-fr]{2.2.4}).
\end{proof}
\oldpage[II]{29}
Lorsqu'on parlera du spectre premier homogène $\operatorname{Proj}(S)$ comme
d'un \emph{schéma}, il s'agira toujours de la structure qui vient d'être
définie.

\begin{example}[2.4.3]
\label{II.2.4.3-fr}
Prenons $S=K[T_1,T_2]$, $K$ étant un corps, $T_1$, $T_2$ deux
indéterminées, et $S$ étant gradué par le degré total. Il résulte de
(\hyperref[II.2.3.14-fr]{2.3.14}) que $\operatorname{Proj}(S)$ est réunion
de $D_+(T_1)$ et $D_+(T_2)$~; on voit aussitôt que ces schémas affines sont
canoniquement isomorphes à $K[T]$, et que $\operatorname{Proj}(S)$
s'obtient par le recollement de ces deux schémas affines décrit dans
(\hyperref[I.2.3.2-fr]{I, 2.3.2}) (cf.
(\hyperref[II.7.4.14-fr]{7.4.14})).
\end{example}

\begin{proposition}[2.4.4]
\label{II.2.4.4-fr}
Soit $S$ un anneau gradué à degrés positifs, et soit $X$ le schéma
$\operatorname{Proj}(S)$.
\begin{enumerate}
  \item[\rm(i)] Si $\mathfrak{N}_+$ est le nilradical de $S_+$
    (\hyperref[II.2.1.10-fr]{2.1.10}), le schéma $X_{\mathrm{red}}$ est
    canoniquement isomorphe à $\operatorname{Proj}(S/\mathfrak{N}_+)$~; en
    particulier, si $S$ est essentiellement réduite,
    $\operatorname{Proj}(S)$ est réduit.
  \item[\rm(ii)] Supposons $S$ essentiellement réduit. Alors, pour que $X$
    soit intègre, il faut et il suffit que $S$ soit essentiellement intègre.
\end{enumerate}
\end{proposition}

\begin{proof}
\begin{enumerate}
  \item[\rm(i)] Soit $\overline{S}$ l'anneau gradué
    $S/\mathfrak{N}_+$, et désignons par $x\mapsto\overline{x}$
    l'homomorphisme canonique $S\to\overline{S}$, de degré $0$. Pour tout
    $f\in S_d$ ($d>0$), l'homomorphisme canonique
    $S_f\to\overline{S}_{\overline{f}}$
    (\hyperref[0.1.5.1-fr]{0, 1.5.1}) est surjectif et de degré $0$, donc
    par restriction donne un homomorphisme surjectif
    $S_{(f)}\to\overline{S}_{(\overline{f})}$~; si on suppose
    $f\notin\mathfrak{N}_+$, on vérifie immédiatement que
    $\overline{S}_{(\overline{f})}$ est réduit, et que le noyau de
    l'homomorphisme précédent est le nilradical de $S_{(f)}$, autrement dit
    $\overline{S}_{(\overline{f})}=(S_{(f)})_{\mathrm{red}}$. À cet
    homomorphisme correspond donc une immersion fermée
    $D_+(\overline{f})\to D_+(f)$ qui identifie $D_+(\overline{f})$ à
    $(D_+(f))_{\mathrm{red}}$ (\hyperref[I.5.1.2-fr]{I, 5.1.2}), et est en
    particulier un homéomorphisme des espaces sous-jacents à ces deux
    schémas affines. En outre, si $g\notin\mathfrak{N}_+$ est un second
    élément homogène de $S_+$, le diagramme
    \[
      \xymatrix{
        S_{(f)}\ar[r]\ar[d]
        & \overline{S}_{(\overline{f})}\ar[d]\\
        S_{(fg)}\ar[r]
        & \overline{S}_{(\overline{fg})}
      }
    \]
    est commutatif~; comme en outre les $D_+(f)$, pour $f$ homogène de
    degré $>0$ et $f\notin\mathfrak{N}_+$, forment un recouvrement de
    $X=\operatorname{Proj}(S)$ (\hyperref[II.2.3.7-fr]{2.3.7}), on voit
    que les morphismes $D_+(\overline{f})\to D_+(f)$ sont les restrictions
    d'une immersion fermée
    $\operatorname{Proj}(\overline{S})\to\operatorname{Proj}(S)$ qui est
    un homéomorphisme des espaces sous-jacents~; d'où la conclusion
    (\hyperref[I.5.1.2-fr]{I, 5.1.2}).
  \item[\rm(ii)] Supposons que $S$ soit essentiellement intègre, autrement
    dit que $(0)$ soit un idéal gradué premier de $S_+$ distinct de $S_+$~;
    alors $X$ est réduit en vertu de (i) et irréductible en vertu de
    (\hyperref[II.2.3.17-fr]{2.3.17}). Inversement, supposons que $S$ soit
    essentiellement réduit et $X$ intègre~; alors, pour $f\neq0$ homogène
    dans $S_+$, on a $D_+(f)\neq\varnothing$
    (\hyperref[II.2.3.7-fr]{2.3.7})~; l'hypothèse que $X$ est irréductible
    entraîne que $D_+(f)\cap D_+(g)\neq\varnothing$ pour $f$, $g$
    homogènes et $\neq0$ dans $S_+$~; donc $fg\neq0$ en vertu de
    (\hyperref[II.2.3.3.2-fr]{2.3.3.2}), et on en conclut que $S_+$ n'a pas
    de diviseurs de zéro, d'où la première assertion.
\end{enumerate}
\end{proof}

\begin{env}[2.4.5]
\label{II.2.4.5-fr}
Étant donné un anneau commutatif $A$, rappelons qu'on dit qu'un anneau
gradué $S$ est une \emph{$A$-algèbre graduée} s'il est muni d'une structure
de $A$-algèbre telle que chacun des sous-groupes $S_n$ soit un $A$-module~;
il suffit d'ailleurs pour cela que $S_0$ soit
\oldpage[II]{30}
une $A$-algèbre, autrement dit on définit sur $S$ une structure de
$A$-algèbre graduée en définissant une structure de $A$-algèbre sur $S_0$,
et en posant, pour $\alpha\in A$, $x\in S_n$,
$\alpha\cdot x=(\alpha\cdot1)x$.
\end{env}

\begin{proposition}[2.4.6]
\label{II.2.4.6-fr}
Supposons que $S$ soit une $A$-algèbre graduée. Alors sur
$X=\operatorname{Proj}(S)$, le faisceau structural $\mathcal{O}_X$ est une
$A$-Algèbre ($A$ étant considéré comme faisceau simple sur $X$)~; en
d'autres termes, $X$ est un schéma au-dessus de $\operatorname{Spec}(A)$.
\end{proposition}

\begin{proof}
Il suffit de noter que pour tout $f$ homogène dans $S_+$, $S_{(f)}$ est une
algèbre sur $A$, et que le diagramme
\[
  \xymatrix{
    S_{(f)}\ar[rr] && S_{(fg)}\\
    & A\ar[ul]\ar[ur] &
  }
\]
est commutatif, pour $f$, $g$ homogènes dans $S_+$.
\end{proof}

\begin{proposition}[2.4.7]
\label{II.2.4.7-fr}
Soit $S$ un anneau gradué à degrés positifs.
\begin{enumerate}
  \item[\rm(i)] Pour tout entier $d>0$, il existe un isomorphisme canonique
    du schéma $\operatorname{Proj}(S)$ sur le schéma
    $\operatorname{Proj}(S^{(d)})$.
  \item[\rm(ii)] Soit $S'$ l'anneau gradué tel que $S'_0=\mathbf{Z}$,
    $S'_n=S_n$ (considéré comme $\mathbf{Z}$-module) pour $n>0$. Il existe un
    isomorphisme canonique du schéma $\operatorname{Proj}(S)$ sur le schéma
    $\operatorname{Proj}(S')$.
\end{enumerate}
\end{proposition}

\begin{proof}
\begin{enumerate}
  \item[\rm(i)] Montrons d'abord que l'application
    $\mathfrak{p}\mapsto\mathfrak{p}\cap S^{(d)}$ est une bijection de
    l'ensemble $\operatorname{Proj}(S)$ sur
    $\operatorname{Proj}(S^{(d)})$. En effet, supposons donné un idéal
    premier gradué $\mathfrak{p}'\in\operatorname{Proj}(S^{(d)})$, et posons
    $\mathfrak{p}_{nd}=\mathfrak{p}'\cap S_{nd}$ ($n\geq0$). Pour tout
    $n>0$, non multiple de $d$, définissons $\mathfrak{p}_n$ comme
    l'ensemble des $x\in S_n$ tels que
    $x^d\in\mathfrak{p}_{nd}$~; si $x\in\mathfrak{p}_n$,
    $y\in\mathfrak{p}_n$, on a
    $(x+y)^{2d}\in\mathfrak{p}_{2nd}$, donc
    $(x+y)^d\in\mathfrak{p}_{nd}$ puisque $\mathfrak{p}'$ est premier~; il
    est immédiat que les $\mathfrak{p}_n$ ainsi définis pour tout $n\geq0$
    vérifient les conditions de
    (\hyperref[II.2.1.9-fr]{2.1.9}), donc il existe un idéal premier unique
    $\mathfrak{p}\in\operatorname{Proj}(S)$ tel que
    $\mathfrak{p}\cap S^{(d)}=\mathfrak{p}'$. Comme pour tout $f$ homogène
    dans $S_+$, on a
    $V_+(f)=V_+(f^d)$ (\hyperref[II.2.3.2.3-fr]{2.3.2.3}), on voit que la
    bijection précédente est un homéomorphisme d'espaces topologiques.
    Enfin, avec les mêmes notations, $S_{(f)}$ et
    $S^{(d)}_{(f^d)}$ s'identifient canoniquement
    (\hyperref[II.2.2.2-fr]{2.2.2}), donc $\operatorname{Proj}(S)$ et
    $\operatorname{Proj}(S^{(d)})$ s'identifient canoniquement en tant que
    schémas.
  \item[\rm(ii)] Si, à tout
    $\mathfrak{p}\in\operatorname{Proj}(S)$ on fait correspondre l'unique
    idéal premier $\mathfrak{p}'\in\operatorname{Proj}(S')$ tel que
    $\mathfrak{p}'\cap S_n=\mathfrak{p}\cap S_n$ pour tout $n>0$
    (\hyperref[II.2.1.9-fr]{2.1.9}), il est clair qu'on définit un
    homéomorphisme canonique
    $\operatorname{Proj}(S)\xrightarrow{\sim}\operatorname{Proj}(S')$ des
    espaces sous-jacents, car $V_+(f)$ est le même ensemble pour $S$ et
    $S'$ lorsque $f$ est un élément homogène de $S_+$. Comme en outre
    $S_{(f)}=S'_{(f)}$, $\operatorname{Proj}(S)$ et
    $\operatorname{Proj}(S')$ s'identifient en tant que schémas.
\end{enumerate}
\end{proof}

\begin{corollary}[2.4.8]
\label{II.2.4.8-fr}
Si $S$ est une $A$-algèbre graduée, $S'_A$ la $A$-algèbre graduée telle que
$(S'_A)_0=A$, $(S'_A)_n=S_n$ pour $n>0$, il existe un isomorphisme canonique
de $\operatorname{Proj}(S)$ sur $\operatorname{Proj}(S'_A)$.
\end{corollary}

\begin{proof}
En effet, ces deux schémas sont isomorphes canoniquement à
$\operatorname{Proj}(S')$, avec les notations de
(\hyperref[II.2.4.7-fr]{2.4.7, (ii)}).
\end{proof}

\subsection{Faisceau associé à un module gradué.}
\label{subsection:II.2.5-fr}

\begin{env}[2.5.1]
\label{II.2.5.1-fr}
Soit $M$ un $S$-module gradué. Pour tout $f$ homogène dans $S_+$,
$M_{(f)}$ est un $S_{(f)}$-module, et il lui correspond donc un faisceau
quasi-cohérent associé $\widetilde{M_{(f)}}$ sur le schéma affine
$\operatorname{Spec}(S_{(f)})$, identifié à $D_+(f)$
(\hyperref[I.1.3.4-fr]{I, 1.3.4}).
\end{env}
\oldpage[II]{31}

\begin{proposition}[2.5.2]
\label{II.2.5.2-fr}
Il existe sur $X=\operatorname{Proj}(S)$ un $\mathcal{O}_X$-Module
quasi-cohérent et un seul $\widetilde{M}$ tel que pour tout $f$ homogène
dans $S_+$, on ait
$\Gamma(D_+(f),\widetilde{M})=M_{(f)}$, l'homomorphisme de restriction
$\Gamma(D_+(f),\widetilde{M})\to
\Gamma(D_+(fg),\widetilde{M})$ pour $f$, $g$ homogènes dans $S_+$, étant
l'homomorphisme canonique $M_{(f)}\to M_{(fg)}$
(\hyperref[II.2.2.3-fr]{2.2.3}).
\end{proposition}

\begin{proof}
Supposons $f\in S_d$, $g\in S_e$. Comme $D_+(fg)$ s'identifie au spectre
premier de $(S_{(f)})_{g^d/f^e}$ en vertu de
(\hyperref[II.2.2.2-fr]{2.2.2}), la restriction à $D_+(fg)$ du faisceau
$\widetilde{M_{(f)}}$ sur $D_+(f)$ s'identifie canoniquement au faisceau
associé au module $(M_{(f)})_{g^d/f^e}$
(\hyperref[I.1.3.6-fr]{I, 1.3.6}), donc aussi à
$\widetilde{M_{(fg)}}$ (\hyperref[II.2.2.2-fr]{2.2.2})~; on en conclut
qu'il existe un isomorphisme canonique
\[
  \theta_{g,f}:\widetilde{M_{(f)}}|D_+(fg)
  \xrightarrow{\sim}\widetilde{M_{(g)}}|D_+(fg)
\]
tel que, si $h$ est un troisième élément homogène de $S_+$, on ait
$\theta_{f,h}=\theta_{f,g}\circ\theta_{g,h}$ dans $D_+(fgh)$. Par suite
(\hyperref[0.3.3.1-fr]{0, 3.3.1}) il existe sur $X$ un
$\mathcal{O}_X$-Module quasi-cohérent $\mathcal{F}$, et pour chaque $f$
homogène dans $S_+$, un isomorphisme $\eta_f$ de
$\mathcal{F}|D_+(f)$ sur $\widetilde{M_{(f)}}$ tels que
$\theta_{g,f}=\eta_g\circ\eta_f^{-1}$. Si alors on considère le faisceau
$\mathcal{G}$ associé au préfaisceau (sur la base de la topologie de $X$
formée des $D_+(f)$) défini par $D_+(f)\to M_{(f)}$, avec les
homomorphismes canoniques $M_{(f)}\to M_{(fg)}$ comme homomorphismes de
restriction, ce qui précède prouve que $\mathcal{F}$ et $\mathcal{G}$ sont
isomorphes (compte tenu de (\hyperref[I.1.3.7-fr]{I, 1.3.7}))~; le
faisceau $\mathcal{G}$ sera désigné par $\widetilde{M}$ et vérifie bien les
conditions de l'énoncé. On a en particulier
$\widetilde{S}=\mathcal{O}_X$.
\end{proof}

\begin{definition}[2.5.3]
\label{II.2.5.3-fr}
On dit que le $\mathcal{O}_X$-Module quasi-cohérent $\widetilde{M}$ défini
dans (\hyperref[II.2.5.2-fr]{2.5.2}) est \emph{associé} au $S$-module
gradué $M$.
\end{definition}

Rappelons que les $S$-modules gradués forment une catégorie lorsqu'on
restreint les homomorphismes de modules gradués aux homomorphismes
\emph{de degré $0$}. Avec cette convention~:

\begin{proposition}[2.5.4]
\label{II.2.5.4-fr}
Le foncteur $M\mapsto\widetilde{M}$ est un foncteur covariant additif exact,
de la catégorie des $S$-modules gradués dans celle des
$\mathcal{O}_X$-Modules quasi-cohérents, qui commute aux limites
inductives et aux sommes directes.
\end{proposition}

\begin{proof}
En effet, ces propriétés étant locales, il suffit de les vérifier sur les
faisceaux
$\widetilde{M}|D_+(f)=\widetilde{M_{(f)}}$~; or, le foncteur
$M\mapsto M_f$, le foncteur $N\mapsto N_0$ (dans la catégorie des
$S_f$-modules gradués) et le foncteur $P\mapsto\widetilde{P}$ (dans la
catégorie des $S_{(f)}$-modules) ont tous trois les propriétés d'exactitude
et de permutabilité aux limites inductives et aux sommes directes
(\hyperref[I.1.3.5-fr]{I, 1.3.5} et
\hyperref[I.1.3.9-fr]{I, 1.3.9})~; d'où la proposition.
\end{proof}

On notera $\widetilde{u}$ l'homomorphisme
$\widetilde{M}\to\widetilde{N}$ correspondant à un homomorphisme de degré
$0$, $u:M\to N$. On déduit aussitôt de
(\hyperref[II.2.5.4-fr]{2.5.4}) que les résultats de
(\hyperref[I.1.3.9-fr]{I, 1.3.9} et
\hyperref[I.1.3.10-fr]{I, 1.3.10}) sont encore valables pour les
$S$-modules gradués et les homomorphismes de degré $0$ (avec le sens donné
ici à $\widetilde{M}$), les démonstrations étant purement formelles.

\begin{proposition}[2.5.5]
\label{II.2.5.5-fr}
Pour tout $\mathfrak{p}\in X=\operatorname{Proj}(S)$, on a
$\widetilde{M}_{\mathfrak{p}}=M_{(\mathfrak{p})}$.
\end{proposition}

\begin{proof}
En effet, par définition
$\widetilde{M}_{\mathfrak{p}}=
\varinjlim\Gamma(D_+(f),\widetilde{M})$, où $f$ parcourt l'ensemble des
éléments homogènes de $S_+$ tels que $f\notin\mathfrak{p}$~; comme
$\Gamma(D_+(f),\widetilde{M})=M_{(f)}$, la proposition résulte de la
définition de $M_{(\mathfrak{p})}$
(\hyperref[II.2.2.7-fr]{2.2.7}).
\end{proof}
\oldpage[II]{32}

En particulier, \emph{l'anneau local $(\mathcal{O}_X)_\mathfrak{p}$} n'est
autre que l'anneau $S_{(\mathfrak{p})}$, ensemble des fractions $x/f$, où
$f$ est homogène dans $S_+$ et n'appartient pas à $\mathfrak{p}$, et $x$
homogène est de même degré que $f$.

Plus particulièrement, si $S$ est \emph{essentiellement intègre}, de sorte
que $\operatorname{Proj}(S)=X$ est \emph{intègre}
(\hyperref[II.2.4.4-fr]{2.4.4}), et si $\xi=(0)$ est le point générique de
$X$, le \emph{corps des fonctions rationnelles}
$R(X)=\mathcal{O}_\xi$, est le corps formé des éléments $fg^{-1}$, où $f$
et $g$ sont homogènes de même degré dans $S_+$ et $g\neq0$.

\begin{proposition}[2.5.6]
\label{II.2.5.6-fr}
Si, pour tout $z\in M$ et tout $f$ homogène dans $S_+$, il existe une
puissance de $f$ annulant $z$, on a $\widetilde{M}=0$. Cette condition
suffisante est aussi nécessaire lorsque la $S_0$-algèbre $S$ est engendrée
par l'ensemble $S_1$ des éléments homogènes de degré $1$.
\end{proposition}

\begin{proof}
En effet, la condition $\widetilde{M}=0$ équivaut à $M_{(f)}=0$ pour tout
$f$ homogène dans $S_+$. D'autre part, si $f\in S_d$, dire que $M_{(f)}=0$
signifie que pour tout $z\in M$ homogène et de degré multiple de $d$, il y
a une puissance $f^n$ telle que $f^nz=0$. Si $M_{(f)}=0$ pour tout
$f\in S_1$, la condition de l'énoncé est donc vérifiée pour tout
$f\in S_1$~; elle l'est \emph{a fortiori} pour tout $f$ homogène dans
$S_+$ lorsque $S_1$ engendre $S$, puisque tout élément homogène de $S_+$
est alors combinaison linéaire de produits d'éléments de $S_1$.
\end{proof}

\begin{proposition}[2.5.7]
\label{II.2.5.7-fr}
Soient $d$ un entier $>0$, $f\in S_d$. Alors, pour tout $n\in\mathbf{Z}$,
le $(\mathcal{O}_X|D_+(f))$-Module
$\widetilde{S(nd)}|D_+(f)$ est canoniquement isomorphe à
$\mathcal{O}_X|D_+(f)$.
\end{proposition}

\begin{proof}
En effet, la multiplication par l'élément inversible $(f/1)^n$ de $S_f$
est une bijection de $S_{(f)}=(S_f)_0$ sur
$(S_f)_{nd}=(S_f(nd))_0=(S(nd)_f)_0=S(nd)_{(f)}$~; autrement dit, les
$S_{(f)}$-modules $S_{(f)}$ et $S(nd)_{(f)}$ sont canoniquement
isomorphes, d'où la proposition.
\end{proof}

\begin{corollary}[2.5.8]
\label{II.2.5.8-fr}
Sur l'ensemble ouvert
\[
  U=\bigcup_{f\in S_d}D_+(f),
\]
la restriction du $\mathcal{O}_X$-Module $\widetilde{S(nd)}$ est un
$(\mathcal{O}_X|U)$-Module inversible
(\hyperref[0.5.4.1-fr]{0, 5.4.1}).
\end{corollary}

\begin{corollary}[2.5.9]
\label{II.2.5.9-fr}
Si l'idéal $S_+$ de $S$ est engendré par l'ensemble $S_1$ des éléments
homogènes de degré $1$, les $\mathcal{O}_X$-Modules
$\widetilde{S(n)}$ sont inversibles pour tout $n\in\mathbf{Z}$.
\end{corollary}

\begin{proof}
Il suffit de remarquer que
$X=\bigcup_{f\in S_1}D_+(f)$ en vertu de l'hypothèse
(\hyperref[II.2.3.14-fr]{2.3.14}) et d'appliquer
(\hyperref[II.2.5.8-fr]{2.5.8}) avec $U=X$.
\end{proof}

\begin{env}[2.5.10]
\label{II.2.5.10-fr}
Nous poserons, dans toute la suite de ce paragraphe
\[
\label{II.2.5.10.1-fr}
  \mathcal{O}_X(n)=\widetilde{S(n)}
\tag{2.5.10.1}
\]
pour tout $n\in\mathbf{Z}$, et pour toute partie ouverte $U$ de $X$ et tout
$(\mathcal{O}_X|U)$-Module $\mathcal{F}$
\[
\label{II.2.5.10.2-fr}
  \mathcal{F}(n)=
  \mathcal{F}\otimes_{\mathcal{O}_X|U}(\mathcal{O}_X(n)|U)
\tag{2.5.10.2}
\]
pour tout $n\in\mathbf{Z}$. Si l'idéal $S_+$ est engendré par $S_1$, le
foncteur $\mathcal{F}(n)$ est \emph{exact} en
$\mathcal{F}$ pour tout $n\in\mathbf{Z}$, car $\mathcal{O}_X(n)$ est alors
un $\mathcal{O}_X$-Module inversible.
\end{env}

\begin{env}[2.5.11]
\label{II.2.5.11-fr}
Soient $M$ et $N$ deux $S$-modules gradués. Pour tout $f\in S_d$ ($d>0$)
on définit un homomorphisme canonique fonctoriel de $S_{(f)}$-modules
\[
\label{II.2.5.11.1-fr}
  \lambda_f:M_{(f)}\otimes_{S_{(f)}}N_{(f)}
  \longrightarrow(M\otimes_S N)_{(f)}
\tag{2.5.11.1}
\]
\oldpage[II]{33}

en composant l'homomorphisme
$M_{(f)}\otimes_{S_{(f)}}N_{(f)}\to M_f\otimes_{S_f}N_f$
(provenant des injections $M_{(f)}\to M_f$, $N_{(f)}\to N_f$ et
$S_{(f)}\to S_f$) et l'isomorphisme canonique
$M_f\otimes_{S_f}N_f\xrightarrow{\sim}(M\otimes_S N)_f$
(\hyperref[0.1.3.4-fr]{0, 1.3.4}) et en notant que, par définition du
produit tensoriel de deux modules gradués, ce dernier isomorphisme conserve
les degrés~; pour $x\in M_{md}$, $y\in N_{nd}$ ($m\geq0$, $n\geq0$), on a
donc
\[
  \lambda_f((x/f^m)\otimes(y/f^n))=(x\otimes y)/f^{m+n}.
\]

Il résulte aussitôt de cette définition que si $g\in S_e$ ($e>0$), le
diagramme
\[
  \xymatrix{
    M_{(f)}\otimes_{S_{(f)}}N_{(f)}\ar[r]^{\lambda_f}\ar[d]
    & (M\otimes_S N)_{(f)}\ar[d]\\
    M_{(fg)}\otimes_{S_{(fg)}}N_{(fg)}\ar[r]_{\lambda_{fg}}
    & (M\otimes_S N)_{(fg)}
  }
\]
(où la flèche verticale de droite est l'homomorphisme canonique et celle de
gauche provient des homomorphismes canoniques) est commutatif. On déduit
donc des $\lambda$ un homomorphisme canonique fonctoriel de
$\mathcal{O}_X$-Modules
\[
\label{II.2.5.11.2-fr}
  \lambda:\widetilde{M}\otimes_{\mathcal{O}_X}\widetilde{N}
  \longrightarrow\widetilde{M\otimes_S N}.
\tag{2.5.11.2}
\]

Considérons en particulier deux idéaux gradués $\mathfrak{J}$,
$\mathfrak{K}$ de $S$~; comme $\widetilde{\mathfrak{J}}$ et
$\widetilde{\mathfrak{K}}$ sont des faisceaux d'idéaux de
$\mathcal{O}_X$, on a un homomorphisme canonique
$\widetilde{\mathfrak{J}}\otimes_{\mathcal{O}_X}
\widetilde{\mathfrak{K}}\to\mathcal{O}_X$~; le diagramme
\[
\label{II.2.5.11.3-fr}
  \xymatrix{
    \widetilde{\mathfrak{J}}\otimes_{\mathcal{O}_X}
    \widetilde{\mathfrak{K}}\ar[rr]^\lambda\ar[dr]
    && \widetilde{\mathfrak{J}\otimes_S\mathfrak{K}}\ar[dl]\\
    & \mathcal{O}_X &
  }
\tag{2.5.11.3}
\]
est alors commutatif. On se ramène en effet à le vérifier dans chaque
ouvert $D_+(f)$ ($f$ homogène dans $S_+$) et cela résulte aussitôt de la
définition (\hyperref[II.2.5.11.1-fr]{2.5.11.1}) de $\lambda_f$ et de
(\hyperref[I.1.3.13-fr]{I, 1.3.13}).

Notons enfin que si $M$, $N$, $P$ sont trois $S$-modules gradués, le
diagramme
\[
\label{II.2.5.11.4-fr}
  \xymatrix{
    \widetilde{M}\otimes_{\mathcal{O}_X}\widetilde{N}
    \otimes_{\mathcal{O}_X}\widetilde{P}
    \ar[r]^{\lambda\otimes1}\ar[d]_{1\otimes\lambda}
    & \widetilde{M\otimes_S N}\otimes_{\mathcal{O}_X}\widetilde{P}
    \ar[d]^\lambda\\
    \widetilde{M}\otimes_{\mathcal{O}_X}\widetilde{N\otimes_S P}
    \ar[r]_\lambda
    & \widetilde{M\otimes_S N\otimes_S P}
  }
\tag{2.5.11.4}
\]
\oldpage[II]{34}

est commutatif. Il suffit encore de le vérifier dans chaque $D_+(f)$ et
cela découle aussitôt des définitions et de
(\hyperref[I.1.3.13-fr]{I, 1.3.13}).
\end{env}

\begin{env}[2.5.12]
\label{II.2.5.12-fr}
Sous les hypothèses de (\hyperref[II.2.5.11-fr]{2.5.11}), on définit un
homomorphisme canonique fonctoriel de $S_{(f)}$-modules
\[
\label{II.2.5.12.1-fr}
  \mu_f:\bigl(\operatorname{Hom}_S(M,N)\bigr)_{(f)}
  \longrightarrow
  \operatorname{Hom}_{S_{(f)}}(M_{(f)},N_{(f)})
\tag{2.5.12.1}
\]
en faisant correspondre à $u/f^n$, où $u$ est un homomorphisme de degré
$nd$, l'homomorphisme $M_{(f)}\to N_{(f)}$ qui transforme $x/f^m$
($x\in M_{md}$) en $u(x)/f^{m+n}$. Pour $g\in S_e$ ($e>0$), on a encore
un diagramme commutatif
\[
  \xymatrix{
    \bigl(\operatorname{Hom}_S(M,N)\bigr)_{(f)}
    \ar[r]^{\mu_f}\ar[d]
    & \operatorname{Hom}_{S_{(f)}}(M_{(f)},N_{(f)})\ar[d]\\
    \bigl(\operatorname{Hom}_S(M,N)\bigr)_{(fg)}
    \ar[r]_{\mu_{fg}}
    & \operatorname{Hom}_{S_{(fg)}}(M_{(fg)},N_{(fg)})
  }
\]
(la flèche de gauche étant l'homomorphisme canonique, celle de droite
provenant des homomorphismes canoniques). On en conclut encore (compte
tenu de (\hyperref[I.1.3.8-fr]{I, 1.3.8})) que les $\mu_f$ définissent un
homomorphisme canonique fonctoriel de $\mathcal{O}_X$-Modules
\[
\label{II.2.5.12.2-fr}
  \mu:\widetilde{\operatorname{Hom}_S(M,N)}
  \longrightarrow
  \mathcal{H}om_{\mathcal{O}_X}(\widetilde{M},\widetilde{N}).
\tag{2.5.12.2}
\]
\end{env}

\begin{proposition}[2.5.13]
\label{II.2.5.13-fr}
Supposons l'idéal $S_+$ engendré par $S_1$. Alors l'homomorphisme
$\lambda$ (\hyperref[II.2.5.11.2-fr]{2.5.11.2}) est un isomorphisme~; il
en est de même de l'homomorphisme $\mu$
(\hyperref[II.2.5.12.2-fr]{2.5.12.2}) lorsque le $S$-module gradué $M$
admet une présentation finie (\hyperref[II.2.1.1-fr]{2.1.1}).
\end{proposition}

\begin{proof}
Comme $X$ est réunion des $D_+(f)$ pour $f\in S_1$
(\hyperref[II.2.3.14-fr]{2.3.14}), on est ramené à prouver que $\lambda_f$
et $\mu_f$ sont des isomorphismes, sous les hypothèses envisagées, lorsque
$f$ est homogène et \emph{de degré $1$}. Or, on définit alors une
application $\mathbf{Z}$-bilinéaire
\[
  M_m\times N_n\longrightarrow
  M_{(f)}\otimes_{S_{(f)}}N_{(f)}
\]
en faisant correspondre à $(x,y)$ l'élément
$(x/f^m)\otimes(y/f^n)$ (lorsque $m<0$, nous écrivons $x/f^m$ pour
$f^{-m}x/1$)~; ces applications définissent une application
$\mathbf{Z}$-linéaire
\[
  M\otimes_{\mathbf{Z}}N\longrightarrow
  M_{(f)}\otimes_{S_{(f)}}N_{(f)},
\]
et si $s\in S_q$, cette application transforme $(sx)\otimes y$ en
\[
  (s/f^q)\bigl((x/f^m)\otimes(y/f^n)\bigr)
\]
(pour $x\in M_m$, $y\in N_n$). On en déduit donc un di-homomorphisme de
modules
\[
  \gamma_f:M\otimes_S N\longrightarrow
  M_{(f)}\otimes_{S_{(f)}}N_{(f)},
\]
relatif à l'homomorphisme canonique $S\to S_{(f)}$ (transformant
$s\in S_q$ en $s/f^q$). Supposons en outre que pour un élément
$\sum_i(x_i\otimes y_i)$ de $M\otimes_S N$ ($x_i$, $y_i$ homogènes de
degrés respectifs $m_i$, $n_i$) on ait
\[
  f^r\sum_i(x_i\otimes y_i)=0,
\]
autrement dit $\sum_i(f^rx_i\otimes y_i)=0$. On en déduit en vertu de
(\hyperref[0.1.3.4-fr]{0, 1.3.4}) que l'on a
\[
  \sum_i(f^rx_i/f^{m_i+r})\otimes(y_i/f^{n_i})=0,
\]
c'est-à-dire $\gamma_f(\sum_i(x_i\otimes y_i))=0$. Par suite $\gamma_f$
se factorise en
\[
  M\otimes_S N\longrightarrow(M\otimes_S N)_f
  \xrightarrow{\gamma'_f}
  M_{(f)}\otimes_{S_{(f)}}N_{(f)}~;
\]
si $\lambda'_f$ est la restriction de $\gamma'_f$
\oldpage[II]{35}
à $(M\otimes_S N)_{(f)}$, on vérifie aussitôt que $\lambda_f$ et
$\lambda'_f$ sont des $S_{(f)}$-homomorphismes réciproques, d'où la
première partie de la proposition.

Pour démontrer la seconde partie, supposons que $M$ soit le conoyau d'un
homomorphisme $P\to Q$, de $S$-modules gradués, $P$ et $Q$ étant sommes
directes d'un nombre fini de modules de la forme $S(n)$~; utilisant
l'exactitude à gauche de $\operatorname{Hom}_S(L,N)$ en $L$ et l'exactitude
de $M_{(f)}$ en $M$, on est aussitôt ramené à prouver que $\mu_f$ est un
isomorphisme lorsque $M=S(n)$. Or, pour tout $z$ homogène dans $N$, soit
$u_z$ l'homomorphisme de $S(n)$ dans $N$ tel que $u_z(1)=z$~; on voit
aussitôt que $\eta:z\mapsto u_z$ est un isomorphisme de degré $0$ de
$N(-n)$ sur $\operatorname{Hom}_S(S(n),N)$. Il lui correspond un
isomorphisme
\[
  \eta_f:\bigl(N(-n)\bigr)_{(f)}
  \longrightarrow
  \bigl(\operatorname{Hom}_S(S(n),N)\bigr)_{(f)}.
\]
D'autre part, soit $\eta'_f$ l'isomorphisme
\[
  N_{(f)}\longrightarrow
  \operatorname{Hom}_{S_{(f)}}\bigl(S(n)_{(f)},N_{(f)}\bigr)
\]
qui, à tout $z'\in N_{(f)}$, fait correspondre l'homomorphisme $v_{z'}$ tel
que $v_{z'}(s/f^k)=sz'/f^{n+k}$ (pour
$s\in S_{n+k}=(S(n))_k$). On constate aisément que l'application composée
\[
  \bigl(N(-n)\bigr)_{(f)}
  \xrightarrow{\eta_f}
  \bigl(\operatorname{Hom}_S(S(n),N)\bigr)_{(f)}
  \xrightarrow{\mu_f}
  \operatorname{Hom}_{S_{(f)}}\bigl(S(n)_{(f)},N_{(f)}\bigr)
  \xrightarrow{{\eta'_f}^{-1}}N_{(f)}
\]
est l'isomorphisme $z/f^h\mapsto z/f^{h-n}$ de
$\bigl(N(-n)\bigr)_{(f)}$ sur $N_{(f)}$, donc $\mu_f$ est un isomorphisme.
\end{proof}

Si l'idéal $S_+$ est engendré par $S_1$, on déduit de
(\hyperref[II.2.5.13-fr]{2.5.13}) que pour tout idéal gradué $\mathfrak{J}$
de $S$ et tout $S$-module gradué $M$, on a
\[
\label{II.2.5.13.1-fr}
  \widetilde{\mathfrak{J}}\cdot\widetilde{M}
  =\widetilde{\mathfrak{J}\cdot M}
\tag{2.5.13.1}
\]
à un isomorphisme canonique près~; cela résulte en effet de la commutativité
du diagramme
\[
  \xymatrix{
    \widetilde{\mathfrak{J}}\otimes_{\mathcal{O}_X}\widetilde{M}
    \ar[rr]^\lambda\ar[dr]
    && \widetilde{\mathfrak{J}\otimes_S M}\ar[dl]\\
    & \widetilde{M} &
  }
\]
que l'on vérifie comme celle de
(\hyperref[II.2.5.11.3-fr]{2.5.11.3}).

\begin{corollary}[2.5.14]
\label{II.2.5.14-fr}
Supposons $S$ engendrée par $S_1$. Quels que soient $m$, $n$ dans
$\mathbf{Z}$, on a alors~:
\[
\label{II.2.5.14.1-fr}
  \mathcal{O}_X(m)\otimes_{\mathcal{O}_X}\mathcal{O}_X(n)
  =\mathcal{O}_X(m+n)
\tag{2.5.14.1}
\]
\[
\label{II.2.5.14.2-fr}
  \mathcal{O}_X(n)=\bigl(\mathcal{O}_X(1)\bigr)^{\otimes n}
\tag{2.5.14.2}
\]
à des isomorphismes canoniques près.
\end{corollary}

\begin{proof}
La première formule résulte de (\hyperref[II.2.5.13-fr]{2.5.13}) et de
l'existence de l'isomorphisme canonique de degré $0$,
$S(m)\otimes_S S(n)\xrightarrow{\sim}S(m+n)$, qui, à l'élément
$1\otimes 1$ (où le premier facteur $1$ est dans $(S(m))_{-m}$ et le second
dans $(S(n))_{-n}$) fait correspondre l'élément
$1\in(S(m+n))_{-m-n}$. Il suffit ensuite de démontrer la seconde formule
pour $n=-1$, et en vertu de (\hyperref[II.2.5.13-fr]{2.5.13}), cela revient
à voir que $\operatorname{Hom}_S(S(1),S)$ est canoniquement isomorphe à
$S(-1)$, vérification immédiate en remontant aux définitions
(\hyperref[II.2.1.2-fr]{2.1.2}) et en se souvenant que $S(1)$ est un
$S$-module monogène.
\end{proof}

\oldpage[II]{36}
\begin{corollary}[2.5.15]
\label{II.2.5.15-fr}
Supposons $S$ engendrée par $S_1$. Pour tout $S$-module gradué $M$ et
tout $n\in\mathbf{Z}$, on a
\[
\label{II.2.5.15.1-fr}
  \widetilde{M(n)}=\widetilde{M}(n)
\tag{2.5.15.1}
\]
à un isomorphisme canonique près.
\end{corollary}

\begin{proof}
Cela résulte des définitions (\hyperref[II.2.5.10.2-fr]{2.5.10.2}) et
(\hyperref[II.2.5.10.1-fr]{2.5.10.1}), de la prop.
(\hyperref[II.2.5.13-fr]{2.5.13}) et de l'existence d'un isomorphisme
canonique $M(n)\xrightarrow{\sim}M\otimes_S S(n)$ de degré $0$ qui, à tout
$z\in(M(n))_h=M_{n+h}$, fait correspondre
$z\otimes1\in M_{n+h}\otimes(S(n))_{-n}\subset(M\otimes_S S(n))_h$.
\end{proof}

\begin{env}[2.5.16]
\label{II.2.5.16-fr}
Désignons par $S'$ l'anneau gradué tel que $S'_0=\mathbf{Z}$,
$S'_n=S_n$ pour $n>0$. Alors, si $f\in S_d$ ($d>0$), on a
$(S(n))_{(f)}=(S'(n))_{(f)}$ pour tout $n\in\mathbf{Z}$, car un élément de
$(S'(n))_{(f)}$ est de la forme $x/f^k$ avec $x\in S'_{n+kd}$ ($k>0$), et
on peut toujours prendre $k$ tel que $n+kd\neq0$. Comme
$X=\operatorname{Proj}(S)$ et $X'=\operatorname{Proj}(S')$ s'identifient
canoniquement (\hyperref[II.2.4.7-fr]{2.4.7, (ii)}), on voit que pour tout
$n\in\mathbf{Z}$, $\mathcal{O}_X(n)$ et $\mathcal{O}_{X'}(n)$ sont images
l'un de l'autre par l'identification précédente.

Notons d'autre part que pour tout $d>0$ et tout $n\in\mathbf{Z}$, on a
\[
  (S^{(d)}(n))_h=S_{(n+h)d}=(S(nd))_{hd}
\]
pour $f\in S_d$, on a donc $(S^{(d)}(n))_{(f)}=(S(nd))_{(f)}$. On sait que
les schémas $X=\operatorname{Proj}(S)$ et
$X^{(d)}=\operatorname{Proj}(S^{(d)})$ s'identifient canoniquement
(\hyperref[II.2.4.7-fr]{2.4.7, (i)})~; ce qui précède montre que si la
$S_0$-algèbre $S^{(d)}$ est engendrée par $S_d$, $\mathcal{O}_X(nd)$ et
$\mathcal{O}_{X^{(d)}}(n)$ sont images l'un de l'autre par cette
identification, pour tout $n\in\mathbf{Z}$.
\end{env}

\begin{proposition}[2.5.17]
\label{II.2.5.17-fr}
Soit $d$ un entier $>0$, et soit
$U=\bigcup_{f\in S_d}D_+(f)$. La restriction à $U$ de l'homomorphisme
canonique
$\mathcal{O}_X(nd)\otimes_{\mathcal{O}_X}\mathcal{O}_X(-nd)
\to\mathcal{O}_X$ est un isomorphisme pour tout entier $n$.
\end{proposition}

\begin{proof}
En vertu de (\hyperref[II.2.5.16-fr]{2.5.16}), on peut se borner au cas où
$d=1$, et la conclusion résulte alors de la démonstration de
(\hyperref[II.2.5.13-fr]{2.5.13}).
\end{proof}

\subsection{$S$-module gradué associé à un faisceau sur
$\operatorname{Proj}(S)$.}
\label{subsection:II.2.6-fr}

\emph{On suppose dans tout ce numéro que l'idéal $S_+$ de $S$ est engendré
par l'ensemble $S_1$ des éléments homogènes de degré $1$.}

\begin{env}[2.6.1]
\label{II.2.6.1-fr}
Le $\mathcal{O}_X$-Module $\mathcal{O}_X(1)$ est alors
\emph{inversible} (\hyperref[II.2.5.9-fr]{2.5.9})~; on pose alors, pour tout
$\mathcal{O}_X$-Module $\mathcal{F}$
(\hyperref[0.5.4.6-fr]{0, 5.4.6}),
\[
\label{II.2.6.1.1-fr}
  \Gamma_*(\mathcal{F})
  =\Gamma_*(\mathcal{O}_X(1),\mathcal{F})
  =\bigoplus_{n\in\mathbf{Z}}\Gamma(X,\mathcal{F}(n))
\tag{2.6.1.1}
\]
compte tenu de (\hyperref[II.2.5.14.2-fr]{2.5.14.2}). Rappelons
(\hyperref[0.5.4.6-fr]{0, 5.4.6}) que $\Gamma_*(\mathcal{O}_X)$ est muni
d'une structure d'\emph{anneau gradué}, et $\Gamma_*(\mathcal{F})$ d'une
structure de
\emph{module gradué} sur $\Gamma_*(\mathcal{O}_X)$.

Puisque $\mathcal{O}_X(n)$ est localement libre, $\Gamma_*(\mathcal{F})$
est un foncteur covariant additif
\emph{exact à gauche} en $\mathcal{F}$~; en particulier, si $\mathcal{J}$
est un faisceau d'idéaux de $\mathcal{O}_X$, $\Gamma_*(\mathcal{J})$
s'identifie canoniquement à un
\emph{idéal gradué} de $\Gamma_*(\mathcal{O}_X)$.
\end{env}

\begin{env}[2.6.2]
\label{II.2.6.2-fr}
Soit $M$ un $S$-module gradué~; pour tout $f\in S_d$ ($d>0$),
$x\mapsto x/1$ est un homomorphisme de groupes abéliens
$M_0\to M_{(f)}$, et comme $M_{(f)}$ s'identifie canoni-

\oldpage[II]{37}

quement à $\Gamma(D_+(f),\widetilde{M})$, on obtient ainsi un
homomorphisme de groupes abéliens
\[
  \alpha_0^f:M_0\longrightarrow\Gamma(D_+(f),\widetilde{M}).
\]
Il est clair que, pour tout $g\in S_e$ ($e>0$), le diagramme
\[
  \xymatrix{
    & \Gamma(D_+(f),\widetilde{M})\ar[dd]\\
    M_0\ar[ur]^{\alpha_0^f}\ar[dr]_{\alpha_0^{fg}}\\
    & \Gamma(D_+(fg),\widetilde{M})
  }
\]
est commutatif~; cela signifie que pour tout $x\in M_0$, les sections
$\alpha_0^f(x)$ et $\alpha_0^g(x)$ de $\widetilde{M}$ coïncident dans
$D_+(f)\cap D_+(g)$, et par suite il existe une section unique
$\alpha_0(x)\in\Gamma(X,\widetilde{M})$ dont la restriction à chaque
$D_+(f)$ est $\alpha_0^f(x)$. On a ainsi défini (sans utiliser l'hypothèse
que $S$ est engendrée par $S_1$) un homomorphisme de groupes abéliens
\[
\label{II.2.6.2.1-fr}
  \alpha_0:M_0\longrightarrow\Gamma(X,\widetilde{M}).
\tag{2.6.2.1}
\]

Appliquant ce résultat au $S$-module gradué $M(n)$ (pour tout
$n\in\mathbf{Z}$), on obtient pour chaque $n\in\mathbf{Z}$ un homomorphisme
de groupes abéliens
\[
\label{II.2.6.2.2-fr}
  \alpha_n:M_n=(M(n))_0\longrightarrow\Gamma(X,\widetilde{M}(n))
\tag{2.6.2.2}
\]
(compte tenu de (\hyperref[II.2.5.15-fr]{2.5.15}))~; d'où un homomorphisme
fonctoriel (de degré $0$) de groupes abéliens gradués
\[
\label{II.2.6.2.3-fr}
  \alpha:M\longrightarrow\Gamma_*(\widetilde{M})
\tag{2.6.2.3}
\]
(aussi noté $\alpha_M$) qui, dans chaque $M_n$, coïncide avec $\alpha_n$.

Si on prend en particulier $M=S$, on vérifie aussitôt (compte tenu de la
définition de la multiplication dans $\Gamma_*(\mathcal{O}_X)$
(\hyperref[0.5.4.6-fr]{0, 5.4.6})) que
$\alpha:S\to\Gamma_*(\mathcal{O}_X)$ est un homomorphisme d'anneaux
gradués, et que pour tout $S$-module gradué $M$,
(\hyperref[II.2.6.2.3-fr]{2.6.2.3}) est un di-homomorphisme de modules
gradués.
\end{env}

\begin{proposition}[2.6.3]
\label{II.2.6.3-fr}
Pour tout $f\in S_d$ ($d>0$), $D_+(f)$ est identique à l'ensemble des
$\mathfrak{p}\in X$ où la section $\alpha_d(f)$ de $\mathcal{O}_X(d)$ ne
s'annule pas (\hyperref[0.5.5.2-fr]{0, 5.5.2}).
\end{proposition}

\begin{proof}
Comme $X=\bigcup_{g\in S_1}D_+(g)$ par hypothèse, il suffit de vérifier que
pour tout $g\in S_1$, l'ensemble des $\mathfrak{p}\in D_+(g)$ où
$\alpha_d(f)$ ne s'annule pas est identique à $D_+(fg)$. Or, la restriction
de $\alpha_d(f)$ à $D_+(g)$ est par définition la section correspondant à
l'élément $f/1$ de $(S(d))_{(g)}$~; par l'isomorphisme canonique
$(S(d))_{(g)}\xrightarrow{\sim}S_{(g)}$
(\hyperref[II.2.5.7-fr]{2.5.7}), cette section de $\mathcal{O}_X(d)$
au-dessus de $D_+(g)$ s'identifie à la section de $\mathcal{O}_X$ au-dessus
de $D_+(g)$ qui correspond à l'élément $f/g^d$ de $S_{(g)}$~; dire que cette
section s'annule en $\mathfrak{p}\in D_+(g)$ signifie que
$f/g^d\in\mathfrak{q}$, où $\mathfrak{q}$ est l'idéal premier de $S_{(g)}$
correspondant à $\mathfrak{p}$ (\hyperref[II.2.3.6-fr]{2.3.6})~; par
définition cela veut dire que $f\in\mathfrak{p}$, d'où la proposition.
\end{proof}

\begin{env}[2.6.4]
\label{II.2.6.4-fr}
Soit maintenant $\mathcal{F}$ un $\mathcal{O}_X$ module, et posons
$M=\Gamma_*(\mathcal{F})$~; en vertu de l'existence de l'homomorphisme
d'anneaux gradués $\alpha:S\to\Gamma_*(\mathcal{O}_X)$, $M$ peut être
considéré comme un $S$-module gradué. Pour tout $f\in S_d$ ($d>0$), il
résulte de (\hyperref[II.2.6.3-fr]{2.6.3}) que la restriction à $D_+(f)$ de
la section $\alpha_d(f)$ de $\mathcal{O}_X(d)$ est inversible~; il en est
donc de même de la restriction à $D_+(f)$ de la section $\alpha_d(f^n)$ de
$\mathcal{O}_X(nd)$, pour tout $n>0$. Soit alors
$z\in M_{nd}=\Gamma(X,\mathcal{F}(nd))$ ($n>0$)~; s'il existe un entier
$k\geq0$ tel que la restriction à $D_+(f)$ de $f^kz$, c'est-à-dire la
\oldpage[II]{38}
section
$(z|D_+(f))(\alpha_d(f^k)|D_+(f))$ de $\mathcal{F}((n+k)d)$, soit nulle,
alors, en raison de la remarque précédente, on a aussi $z|D_+(f)=0$. Cela
montre que l'on définit un $S_{(f)}$-homomorphisme
$\beta_f:M_{(f)}\to\Gamma(D_+(f),\mathcal{F})$ en faisant correspondre à
l'élément $z/f^n$ la section
$(z|D_+(f))(\alpha_d(f^n)|D_+(f))^{-1}$ de $\mathcal{F}$ au-dessus de
$D_+(f)$. On vérifie en outre aussitôt que le diagramme
\[
\label{II.2.6.4.1-fr}
  \xymatrix{
    M_{(f)}\ar[r]^{\beta_f}\ar[d]
    & \Gamma(D_+(f),\mathcal{F})\ar[d]\\
    M_{(fg)}\ar[r]_{\beta_{fg}}
    & \Gamma(D_+(fg),\mathcal{F})
  }
\tag{2.6.4.1}
\]
est commutatif pour $g\in S_e$ ($e>0$). Si on se rappelle que $M_{(f)}$
s'identifie canoniquement à $\Gamma(D_+(f),\widetilde{M})$ et que les
$D_+(f)$ forment une base de la topologie de $X$
(\hyperref[II.2.3.4-fr]{2.3.4}), on voit que les $\beta_f$ proviennent d'un
unique homomorphisme canonique de $\mathcal{O}_X$-Modules
\[
\label{II.2.6.4.2-fr}
  \beta:\widetilde{\Gamma_*(\mathcal{F})}
  \longrightarrow\mathcal{F}
\tag{2.6.4.2}
\]
(aussi noté $\beta_{\mathcal{F}}$) qui est évidemment fonctoriel.
\end{env}

\begin{proposition}[2.6.5]
\label{II.2.6.5-fr}
Soient $M$ un $S$-module gradué, $\mathcal{F}$ un $\mathcal{O}_X$-Module~;
alors les homomorphismes composés
\[
\label{II.2.6.5.1-fr}
  \widetilde{M}\xrightarrow{\widetilde{\alpha}}
  \widetilde{\Gamma_*(\widetilde{M})}
  \xrightarrow{\beta}\widetilde{M}
\tag{2.6.5.1}
\]
\[
\label{II.2.6.5.2-fr}
  \Gamma_*(\mathcal{F})\xrightarrow{\alpha}
  \Gamma_*(\widetilde{\Gamma_*(\mathcal{F})})
  \xrightarrow{\Gamma_*(\beta)}\Gamma_*(\mathcal{F})
\tag{2.6.5.2}
\]
sont les isomorphismes identiques.
\end{proposition}

\begin{proof}
La vérification de (\hyperref[II.2.6.5.1-fr]{2.6.5.1}) est locale~: dans un
ouvert $D_+(f)$, elle résulte aussitôt des définitions et de ce que $\beta$,
appliqué à des faisceaux quasi-cohérents, est déterminé par son action sur les
sections au-dessus de $D_+(f)$ (\hyperref[I.1.3.8-fr]{I, 1.3.8}). La
vérification de (\hyperref[II.2.6.5.2-fr]{2.6.5.2}) se fait pour chaque degré
séparément~: si on pose $M=\Gamma_*(\mathcal{F})$, on a
$M_n=\Gamma(X,\mathcal{F}(n))$ et
$(\Gamma_*(\widetilde{M}))_n=\Gamma(X,\widetilde{M}(n))
=\Gamma(X,\widetilde{M(n)})$. Or, si $f\in S_1$ et $z\in M_n$,
$\alpha_n^f(z)$ est l'élément $z/1$ de $(M(n))_{(f)}$, égal à
$(f/1)^n(z/f^n)$~; il lui correspond par $\beta_f$ la section
\[
  ((\alpha_1(f))^n|D_+(f))
  ((z|D_+(f))((\alpha_1(f))^n|D_+(f))^{-1})
\]
au-dessus de $D_+(f)$, c'est-à-dire la restriction de $z$ à $D_+(f)$, ce qui
vérifie (\hyperref[II.2.6.5.2-fr]{2.6.5.2}).
\end{proof}

\subsection{Conditions de finitude.}
\label{subsection:II.2.7-fr}

\begin{proposition}[2.7.1]
\label{II.2.7.1-fr}
\begin{enumerate}
  \item[\rm(i)] Si $S$ est un anneau gradué noethérien,
    $X=\operatorname{Proj}(S)$ est un schéma noethérien.
  \item[\rm(ii)] Si $S$ est une $A$-algèbre graduée de type fini,
    $X=\operatorname{Proj}(S)$ est un schéma de type fini au-dessus de
    $Y=\operatorname{Spec}(A)$.
\end{enumerate}
\end{proposition}

\oldpage[II]{39}
\begin{proof}
\begin{enumerate}
  \item[\rm(i)] Si $S$ est noethérien, l'idéal $S_+$ admet un système fini de
    générateurs homogènes $f_i$ ($1\leq i\leq p$), donc
    (\hyperref[II.2.3.14-fr]{2.3.14}) l'espace sous-jacent $X$ est réunion des
    $D_+(f_i)=\operatorname{Spec}(S_{(f_i)})$, et tout revient à voir que
    chacun des $S_{(f_i)}$ est noethérien, ce qui résulte de
    (\hyperref[II.2.2.6-fr]{2.2.6}).
  \item[\rm(ii)] L'hypothèse entraîne que $S_0$ est une $A$-algèbre de type
    fini et $S$ une $S_0$-algèbre de type fini, donc $S_+$ est un idéal de
    type fini (\hyperref[II.2.1.4-fr]{2.1.4}). On est alors ramené, comme dans
    (i), à prouver que pour $f\in S_d$, $S_{(f)}$ est une $A$-algèbre de type
    fini. En vertu de (\hyperref[II.2.2.5-fr]{2.2.5}), il suffit de montrer
    que $S^{(d)}$ est une $A$-algèbre de type fini, ce qui résulte de
    (\hyperref[II.2.1.6-fr]{2.1.6}).
\end{enumerate}
\end{proof}

\begin{env}[2.7.2]
\label{II.2.7.2-fr}
Nous considérerons dans ce qui suit les conditions de finitude suivantes pour
un $S$-module gradué $M$~:
\begin{enumerate}
  \item[\rm(TF)] Il existe un entier $n$ tel que le sous-module
    $\bigoplus_{k\geq n}M_k$ soit un $S$-module de type fini.
  \item[\rm(TN)] Il existe un entier $n$ tel que $M_k=0$ pour $k\geq n$.
\end{enumerate}

Si $M$ vérifie (TN), on a $M_{(f)}=0$ pour tout $f$ homogène dans $S_+$,
donc $\widetilde{M}=0$.

Soient $M$, $N$ deux $S$-modules gradués~; on dit qu'un homomorphisme
$u:M\to N$ de degré $0$ est (TN)-\emph{injectif} (resp.
(TN)-\emph{surjectif}, (TN)-\emph{bijectif}) s'il existe un entier $n$ tel
que $u_k:M_k\to N_k$ soit injectif (resp. surjectif, bijectif) pour
$k\geq n$. Dire que $u$ est (TN)-injectif (resp. (TN)-surjectif) revient donc
à dire que $\operatorname{Ker}u$ (resp. $\operatorname{Coker}u$) vérifie
(TN). En vertu de (\hyperref[II.2.5.4-fr]{2.5.4}), si $u$ est (TN)-injectif
(resp. (TN)-surjectif, (TN)-bijectif), $\widetilde{u}$ est injectif (resp.
surjectif, bijectif)~; lorsque $u$ est (TN)-bijectif, on dit encore que $u$
est un (TN)-\emph{isomorphisme}.
\end{env}

\begin{proposition}[2.7.3]
\label{II.2.7.3-fr}
Soient $S$ un anneau gradué tel que l'idéal $S_+$ soit de type fini, $M$ un
$S$-module gradué.
\begin{enumerate}
  \item[\rm(i)] Si $M$ vérifie la condition (TF), le $\mathcal{O}_X$-Module
    $\widetilde{M}$ est de type fini.
  \item[\rm(ii)] Supposons que $M$ vérifie (TF)~; pour que $\widetilde{M}=0$,
    il faut et il suffit que $M$ vérifie (TN).
\end{enumerate}
\end{proposition}

\begin{proof}
On vient de voir que la condition (TN) implique $\widetilde{M}=0$. Si $M$
vérifie (TF), le sous-module gradué $M'=\bigoplus_{k\geq n}M_k$ qui est par
hypothèse de type fini, est tel que $M/M'$ satisfasse à (TN)~; on a par suite
$\widetilde{M/M'}=0$, et l'exactitude du foncteur
$M\mapsto\widetilde{M}$ (\hyperref[II.2.5.4-fr]{2.5.4}) entraîne
$\widetilde{M}=\widetilde{M'}$~; pour prouver que $\widetilde{M}$ est de type
fini, on est donc ramené au cas où $M$ est de type fini. La question étant
locale, il suffit de prouver que $M_{(f)}$ est un $S_{(f)}$-module de type
fini (\hyperref[I.1.3.9-fr]{I, 1.3.9})~; mais $M^{(d)}$ est un
$S^{(d)}$-module de type fini (\hyperref[II.2.1.6-fr]{2.1.6, (iii)}) et notre
assertion résulte alors de (\hyperref[II.2.2.5-fr]{2.2.5}).

Supposons maintenant que $M$ vérifie (TF) et que $\widetilde{M}=0$~; alors,
avec les mêmes notations que ci-dessus, on a $\widetilde{M'}=0$, et la
condition (TN) pour $M'$ est équivalente à la condition (TN) pour $M$, donc
pour prouver que $\widetilde{M}=0$ implique que $M$ vérifie (TN), on peut
encore se borner au cas où $M$ est engendré par un nombre fini d'éléments
homogènes $x_i$ ($1\leq i\leq p$)~; soit d'autre part
$(f_j)_{1\leq j\leq q}$ un système de générateurs homogènes de l'idéal $S_+$.
On a par hypothèse $M_{(f_j)}=0$ pour tout $j$, donc il existe un entier $n$
tel que $f_j^nx_i=0$ quels que soient $i$ et $j$. Soit $n_j$ le degré de
$f_j$, et soit $m$ la plus grande valeur de $\sum_jr_jn_j$ pour les systèmes
d'entiers $(r_j)$ en nombre fini tels que $\sum_jr_j\leq nq$~; il est clair
alors

\oldpage[II]{40}
que si $k>m$, on a $S_kx_i=0$ pour tout $i$~; si $h$ est le plus grand des
degrés des $x_i$, on en conclut que $M_k=0$ pour $k>h+m$, ce qui termine la
démonstration.
\end{proof}

\begin{corollary}[2.7.4]
\label{II.2.7.4-fr}
Soit $S$ un anneau gradué tel que l'idéal $S_+$ soit de type fini~; pour que
$X=\operatorname{Proj}(S)=\varnothing$, il faut et il suffit qu'il existe $n$
tel que $S_k=0$ pour $k\geq n$.
\end{corollary}

\begin{proof}
La condition $X=\varnothing$ est en effet équivalente à
$\mathcal{O}_X=\widetilde{S}=0$, et $S$ est un $S$-module monogène.
\end{proof}

\begin{theorem}[2.7.5]
\label{II.2.7.5-fr}
On suppose que l'idéal $S_+$ est engendré par un nombre fini d'éléments
homogènes de degré $1$~; soit $X=\operatorname{Proj}(S)$. Alors, pour tout
$\mathcal{O}_X$-Module quasi-cohérent $\mathcal{F}$, l'homomorphisme
canonique $\beta:\widetilde{\Gamma_*(\mathcal{F})}
\to\mathcal{F}$ (\hyperref[II.2.6.4-fr]{2.6.4}) est un isomorphisme.
\end{theorem}

\begin{proof}
En effet, si $S_+$ est engendré par un nombre fini d'éléments $f_i\in S_1$,
$X$ est réunion des sous-espaces $\operatorname{Spec}(S_{(f_i)})$
(\hyperref[II.2.3.6-fr]{2.3.6}) qui sont quasi-compacts, donc est
quasi-compact~; en outre $X$ est un schéma
(\hyperref[II.2.4.2-fr]{2.4.2})~; d'après
(\hyperref[I.9.3.2-fr]{I, 9.3.2}),
(\hyperref[II.2.5.14.2-fr]{2.5.14.2}) et
(\hyperref[II.2.6.3-fr]{2.6.3}) on a, pour tout $f\in S_d$ ($d>0$), un
isomorphisme canonique
$(\Gamma_*(\mathcal{F}))_{(\alpha_d(f))}\xrightarrow{\sim}
\Gamma(D_+(f),\mathcal{F})$~; d'ailleurs, par définition,
$(\Gamma_*(\mathcal{F}))_{(\alpha_d(f))}$ (où $\Gamma_*(\mathcal{F})$ est
considéré comme $\Gamma_*(\mathcal{O}_X)$-module) n'est autre que
$(\Gamma_*(\mathcal{F}))_{(f)}$ (où $\Gamma_*(\mathcal{F})$ est considéré
comme $S$-module)~; si on se reporte à la définition
(\hyperref[I.9.3.1-fr]{I, 9.3.1}) de l'isomorphisme canonique précédent, on
voit qu'il coïncide avec l'homomorphisme $\beta_f$, d'où le théorème.
\end{proof}

\begin{remark}[2.7.6]
\label{II.2.7.6-fr}
Si on suppose l'anneau gradué $S$ \emph{noethérien}, la condition de
(\hyperref[II.2.7.5-fr]{2.7.5}) est vérifiée \emph{ipso facto} dès que l'on
suppose l'idéal $S_+$ engendré par l'ensemble $S_1$ des éléments homogènes de
degré $1$.
\end{remark}

\begin{corollary}[2.7.7]
\label{II.2.7.7-fr}
Sous les hypothèses de (\hyperref[II.2.7.5-fr]{2.7.5}), tout
$\mathcal{O}_X$-Module quasi-cohérent $\mathcal{F}$ est isomorphe à un
$\mathcal{O}_X$-Module de la forme $\widetilde{M}$, où $M$ est un $S$-module
gradué.
\end{corollary}

\begin{corollary}[2.7.8]
\label{II.2.7.8-fr}
Sous les hypothèses de (\hyperref[II.2.7.5-fr]{2.7.5}), tout
$\mathcal{O}_X$-Module quasi-cohérent de type fini $\mathcal{F}$ est
isomorphe à un $\mathcal{O}_X$-Module de la forme $\widetilde{N}$, où $N$ est
un $S$-module gradué de type fini.
\end{corollary}

\begin{proof}
On peut supposer $\mathcal{F}=\widetilde{M}$, où $M$ est un $S$-module gradué
(\hyperref[II.2.7.7-fr]{2.7.7}). Soit $(f_\lambda)_{\lambda\in L}$ un
système de générateurs homogènes de $M$~; pour toute partie finie $H$ de $L$,
soit $M_H$ le sous-module gradué de $M$ engendré par les $f_\lambda$ tels que
$\lambda\in H$~; il est clair que $M$ est limite inductive de ses sous-modules
$M_H$, donc $\mathcal{F}$ est limite inductive de ses sous-$\mathcal{O}_X$-Modules
$\widetilde{M_H}$ (\hyperref[II.2.5.4-fr]{2.5.4}). Mais comme $\mathcal{F}$
est de type fini et l'espace sous-jacent $X$ quasi-compact, il résulte de
(\hyperref[0.5.2.3-fr]{0, 5.2.3}) que $\mathcal{F}=\widetilde{M_H}$ pour une
partie finie $H$ de $L$.
\end{proof}

\begin{corollary}[2.7.9]
\label{II.2.7.9-fr}
Sous les hypothèses de (\hyperref[II.2.7.5-fr]{2.7.5}), soit $\mathcal{F}$ un
$\mathcal{O}_X$-Module quasi-cohérent de type fini. Il existe alors un entier
$n_0$ tel que, pour tout $n\geq n_0$, $\mathcal{F}(n)$ soit isomorphe à un
quotient d'un $\mathcal{O}_X$-Module de la forme $\mathcal{O}_X^k$ ($k>0$
dépendant de $n$), et par suite est engendré par un nombre fini de ses sections
au-dessus de $X$ (\hyperref[0.5.1.1-fr]{0, 5.1.1}).
\end{corollary}

\begin{proof}
En vertu de (\hyperref[II.2.7.8-fr]{2.7.8}), on peut supposer que
$\mathcal{F}=\widetilde{M}$, où $M$ est quotient d'une somme directe finie de
$S$-modules de la forme $S(m_i)$~; en vertu de
(\hyperref[II.2.5.4-fr]{2.5.4}) on peut donc se borner au cas où $M=S(m)$,
donc $\mathcal{F}(n)=\widetilde{S(m+n)}
=\mathcal{O}_X(m+n)$ (\hyperref[II.2.5.15-fr]{2.5.15}). Il suffira donc de
démontrer le

\oldpage[II]{41}
\begin{lemma}[2.7.9.1]
\label{II.2.7.9.1-fr}
Sous les hypothèses de (\hyperref[II.2.7.5-fr]{2.7.5}), pour tout $n\geq0$,
il existe un entier $k$ (dépendant de $n$) et un homomorphisme surjectif
$\mathcal{O}_X^k\to\mathcal{O}_X(n)$.
\end{lemma}

Il suffit en effet (\hyperref[II.2.7.2-fr]{2.7.2}) de montrer que, pour un $k$
convenable, il y a un homomorphisme (TN)-surjectif $u$ \emph{de degré $0$},
du $S$-module gradué produit $S^k$ dans $S(n)$. Or, on a $(S(n))_0=S_n$, et
par hypothèse $S_h=S_1^h$ pour tout $h>0$, donc
$SS_n=S_n+S_{n+1}+\cdots$. Comme $S_n$ est un $S_0$-module de type fini
(\hyperref[II.2.1.5-fr]{2.1.5} et
\hyperref[II.2.1.6-fr]{2.1.6, (i)}), considérons un système de générateurs
$(a_i)_{1\leq i\leq k}$ de ce module~; l'homomorphisme $u$ fera correspondre
$a_i$ au $i$-ème élément de la base canonique de $S^k$ ($1\leq i\leq k$)~;
comme $\operatorname{Coker}u$ s'identifie alors à
$(S(n))_{-n}+\cdots+(S(n))_{-1}$, $u$ répond bien à la question.
\end{proof}

\begin{corollary}[2.7.10]
\label{II.2.7.10-fr}
Sous les hypothèses de (\hyperref[II.2.7.5-fr]{2.7.5}), soit $\mathcal{F}$ un
$\mathcal{O}_X$-Module quasi-cohérent de type fini. Il existe un entier $n_0$
tel que, pour tout $n\geq n_0$, $\mathcal{F}$ soit isomorphe à un quotient
d'un $\mathcal{O}_X$-Module de la forme $(\mathcal{O}_X(-n))^k$ ($k$
dépendant de $n$).
\end{corollary}

\begin{proposition}[2.7.11]
\label{II.2.7.11-fr}
Supposons vérifiées les hypothèses de
(\hyperref[II.2.7.5-fr]{2.7.5}) et soit $M$ un $S$-module gradué. Alors~:
\begin{enumerate}
  \item[\rm(i)] L'homomorphisme canonique
    $\widetilde{\alpha}:\widetilde{M}\to
    \widetilde{\Gamma_*(\widetilde{M})}$ est un isomorphisme.
  \item[\rm(ii)] Soit $\mathcal{G}$ un sous-$\mathcal{O}_X$-Module
    quasi-cohérent de $\widetilde{M}$, et soit $N$ le sous-$S$-module gradué
    de $M$, image réciproque de $\Gamma_*(\mathcal{G})$ par $\alpha$. Alors
    on a $\widetilde{N}=\mathcal{G}$ ($\widetilde{N}$ étant identifié, en
    vertu de (\hyperref[II.2.5.4-fr]{2.5.4}), à un
    sous-$\mathcal{O}_X$-Module de $\widetilde{M}$).
\end{enumerate}
\end{proposition}

\begin{proof}
Comme $\beta:\widetilde{\Gamma_*(\widetilde{M})}
\to\widetilde{M}$ est un isomorphisme en vertu de
(\hyperref[II.2.7.5-fr]{2.7.5}), $\widetilde{\alpha}$ est l'isomorphisme
réciproque en vertu de (\hyperref[II.2.6.5.1-fr]{2.6.5.1}), d'où (i). Soit
$P$ le sous-module gradué $\alpha(M)$ de $\Gamma_*(\widetilde{M})$~; comme
$M\mapsto\widetilde{M}$ est un foncteur exact
(\hyperref[II.2.5.4-fr]{2.5.4}), l'image de $\widetilde{M}$ par
$\widetilde{\alpha}$ est égale à $\widetilde{P}$, donc, en vertu de (i),
$\widetilde{P}=\widetilde{\Gamma_*(\widetilde{M})}$. Posons
$Q=\Gamma_*(\mathcal{G})\cap P$~; en vertu de ce qui précède et de
(\hyperref[II.2.5.4-fr]{2.5.4}), on a
$\widetilde{Q}=\widetilde{\Gamma_*(\mathcal{G})}$, donc la
restriction de $\beta$ à $\widetilde{Q}$ est un \emph{isomorphisme} de ce
$\mathcal{O}_X$-Module sur $\mathcal{G}$ par
(\hyperref[II.2.7.5-fr]{2.7.5}). Mais par définition de $N$ et
(\hyperref[II.2.5.4-fr]{2.5.4}), la restriction de l'isomorphisme
$\widetilde{\alpha}$ à $\widetilde{N}$ est un isomorphisme de $\widetilde{N}$
sur $\widetilde{Q}$, d'où la conclusion par
(\hyperref[II.2.6.5.1-fr]{2.6.5.1}).
\end{proof}

\subsection{Comportements fonctoriels.}
\label{subsection:II.2.8-fr}

\begin{env}[2.8.1]
\label{II.2.8.1-fr}
Soient $S$, $S'$ deux anneaux gradués à degrés positifs,
$\varphi:S'\to S$ un homomorphisme d'anneaux gradués. Nous désignerons par
$G(\varphi)$ la partie ouverte de $X=\operatorname{Proj}(S)$ complémentaire
de $V_+(\varphi(S'_+))$, ou, ce qui revient au même, la réunion des
$D_+(\varphi(f'))$, où $f'$ parcourt l'ensemble des éléments homogènes de
$S'_+$. La restriction à $G(\varphi)$ de l'application continue
${}^a\varphi$ de $\operatorname{Spec}(S)$ dans $\operatorname{Spec}(S')$
(\hyperref[I.1.2.1-fr]{I, 1.2.1}) est donc une application continue de
$G(\varphi)$ dans $\operatorname{Proj}(S')$, que nous noterons encore
${}^a\varphi$ par abus de langage. Si $f'\in S'_+$ est homogène, on a
\[
\label{II.2.8.1.1-fr}
  {}^a\varphi^{-1}(D_+(f'))=D_+(\varphi(f'))
\tag{2.8.1.1}
\]
compte tenu du fait que ${}^a\varphi$ applique $G(\varphi)$ dans
$\operatorname{Proj}(S')$ et de
(\hyperref[I.1.2.2.2-fr]{I, 1.2.2.2}). D'autre part, l'homomorphisme
$\varphi$ définit canoniquement (avec les mêmes notations) un homomorphisme
d'anneaux gradués $S'_{f'}\to S_f$, d'où, par restriction aux éléments de
degré $0$,

\oldpage[II]{42}
un homomorphisme $S'_{(f')}\to S_{(f)}$ que nous désignerons par
$\varphi_{(f)}$~; il lui correspond
(\hyperref[I.1.6.1-fr]{I, 1.6.1}) un morphisme
$({}^a\varphi_{(f)},\widetilde{\varphi}_{(f)}):
\operatorname{Spec}(S_{(f)})\to\operatorname{Spec}(S'_{(f')})$ de schémas
affines. Si on identifie canoniquement $\operatorname{Spec}(S_{(f)})$ au
schéma induit par $\operatorname{Proj}(S)$ sur $D_+(f)$
(\hyperref[II.2.3.6-fr]{2.3.6}), on a défini un morphisme
$\Phi_f:D_+(f)\to D_+(f')$ et ${}^a\varphi_{(f)}$ s'identifie à la
restriction de ${}^a\varphi$ à $D_+(f)$. D'autre part, il est immédiat que si
$g'$ est un second élément homogène de $S'_+$ et $g=\varphi(g')$, le
diagramme
\[
  \xymatrix{
    D_+(f)\ar[r]^{\Phi_f} & D_+(f')\\
    D_+(fg)\ar[u]\ar[r]_{\Phi_{fg}} & D_+(f'g')\ar[u]
  }
\]
est commutatif, en raison de la commutativité du diagramme
\[
  \xymatrix{
    S'_{(f')}\ar[r]^{\varphi_{(f)}}\ar[d]_{\omega_{f'g',f'}}
    & S_{(f)}\ar[d]^{\omega_{fg,f}}\\
    S'_{(f'g')}\ar[r]_{\varphi_{(fg)}} & S_{(fg)}
  }
\]
Compte tenu de la définition de $G(\varphi)$ et de
(\hyperref[II.2.3.3.2-fr]{2.3.3.2}), on voit donc que~:
\end{env}

\begin{proposition}[2.8.2]
\label{II.2.8.2-fr}
Étant donné un homomorphisme d'anneaux gradués $\varphi:S'\to S$, il existe
un morphisme et un seul $({}^a\varphi,\widetilde{\varphi})$ du préschéma
induit $G(\varphi)$ dans $\operatorname{Proj}(S')$ (dit \emph{associé à
$\varphi$} et noté $\operatorname{Proj}(\varphi)$), tel que pour tout élément
homogène $f'\in S'_+$, la restriction de ce morphisme à
$D_+(\varphi(f'))$ coïncide avec le morphisme associé à l'homomorphisme
$S'_{(f')}\to S_{(\varphi(f'))}$ correspondant à $\varphi$.
\end{proposition}

\begin{proof}
On notera qu'avec les notations précédentes, si $f'\in S'_d$, le diagramme
\[
\label{II.2.8.2.1-fr}
  \xymatrix{
    S'_{(f')}\ar[r]^{\varphi_{(f)}}\ar[d]_{\sim}
    & S_{(f)}\ar[d]^{\sim}\\
    {S'}^{(d)}/(f'-1){S'}^{(d)}\ar[r]
    & S^{(d)}/(f-1)S^{(d)}
  }
\tag{2.8.2.1}
\]
est commutatif (les flèches verticales étant les isomorphismes
(\hyperref[II.2.2.5-fr]{2.2.5})).
\end{proof}

\begin{corollary}[2.8.3]
\label{II.2.8.3-fr}
\begin{enumerate}
  \item[\rm(i)] Le morphisme $\operatorname{Proj}(\varphi)$ est affine.
  \item[\rm(ii)] Si $\operatorname{Ker}(\varphi)$ est nilpotent (et en
    particulier si $\varphi$ est injectif), le morphisme
    $\operatorname{Proj}(\varphi)$ est dominant.
\end{enumerate}
\end{corollary}

\begin{proof}
(i) est conséquence immédiate de (\hyperref[II.2.8.2-fr]{2.8.2}) et de
(\hyperref[II.2.8.1.1-fr]{2.8.1.1}). D'autre part, si
$\operatorname{Ker}(\varphi)$ est nilpotent, pour tout $f'$ homogène dans
$S'_+$, on vérifie aussitôt que $\operatorname{Ker}(\varphi_f)$ (avec
$f=\varphi(f')$) est nilpotent, donc aussi
$\operatorname{Ker}(\varphi_{(f)})$~; la conclusion résulte donc de
(\hyperref[II.2.8.2-fr]{2.8.2}) et de
(\hyperref[I.1.2.7-fr]{I, 1.2.7}).
\end{proof}

\oldpage[II]{43}
On notera qu'il y a en général des morphismes de $\operatorname{Proj}(S)$
dans $\operatorname{Proj}(S')$ qui ne sont pas affines, et ne proviennent
donc pas d'homomorphismes d'anneaux gradués $S'\to S$~; un exemple est le
morphisme structural $\operatorname{Proj}(S)\to\operatorname{Spec}(A)$ quand
$A$ est un corps ($\operatorname{Spec}(A)$ étant identifié à
$\operatorname{Proj}(A[T])$ (cf. (\hyperref[II.3.1.7-fr]{3.1.7})))~; cela
résulte en effet de (\hyperref[I.2.3.2-fr]{I, 2.3.2}).

\begin{env}[2.8.4]
\label{II.2.8.4-fr}
Soient $S''$ un troisième anneau gradué à degrés positifs,
$\varphi':S''\to S'$ un homomorphisme d'anneaux gradués, et posons
$\varphi''=\varphi\circ\varphi'$. En raison de
(\hyperref[II.2.8.1.1-fr]{2.8.1.1}) et de la formule
${}^a\varphi''=({}^a\varphi')\circ({}^a\varphi)$ on vérifie aussitôt que l'on
a $G(\varphi'')\subset G(\varphi)$ et que, si $\Phi$, $\Phi'$ et $\Phi''$
sont les morphismes associés à $\varphi$, $\varphi'$ et $\varphi''$, on a
$\Phi''=\Phi'\circ(\Phi|G(\varphi''))$.
\end{env}

\begin{env}[2.8.5]
\label{II.2.8.5-fr}
Supposons que $S$ (resp. $S'$) soit une $A$-algèbre graduée (resp. une
$A'$-algèbre graduée), et soit $\psi:A'\to A$ un homomorphisme d'anneaux tel
que le diagramme
\[
\label{II.2.8.5.1-fr}
  \xymatrix{
    A'\ar[r]^{\psi}\ar[d] & A\ar[d]\\
    S'\ar[r]_{\varphi} & S
  }
\tag{2.8.5.1}
\]
soit commutatif. On peut alors considérer $G(\varphi)$ et
$\operatorname{Proj}(S')$ comme des schémas au-dessus de
$\operatorname{Spec}(A)$ et $\operatorname{Spec}(A')$ respectivement~; si
$\Phi$ et $\Psi$ sont les morphismes associés respectivement à $\varphi$ et
$\psi$, le diagramme
\[
  \xymatrix{
    G(\varphi)\ar[r]^{\Phi}\ar[d] & \operatorname{Proj}(S')\ar[d]\\
    \operatorname{Spec}(A)\ar[r]_{\Psi} & \operatorname{Spec}(A')
  }
\]
est alors commutatif~: il suffit en effet de le démontrer pour la restriction
de $\Phi$ à $D_+(f)$, où $f=\varphi(f')$, $f'$ homogène dans $S'_+$~; et
alors cela résulte de la commutativité du diagramme
\[
  \xymatrix{
    A'\ar[r]^{\psi}\ar[d] & A\ar[d]\\
    S'_{(f')}\ar[r]_{\varphi_{(f)}} & S_{(f)}
  }
\]
\end{env}

\begin{env}[2.8.6]
\label{II.2.8.6-fr}
Soit maintenant $M$ un $S$-module gradué, et considérons le $S'$-module
$M_{[\varphi]}$, qui est évidemment gradué. Soit $f'$ homogène dans $S'_+$,
et soit $f=\varphi(f')$~; on sait
(\hyperref[0.1.5.2-fr]{0, 1.5.2}) qu'il y a un isomorphisme canonique
$(M_{[\varphi]})_{f'}\xrightarrow{\sim}(M_f)_{[\varphi_f]}$, et il est
immédiat que cet isomorphisme conserve les degrés, donc il y a un isomorphisme
canonique
$(M_{[\varphi]})_{(f')}\xrightarrow{\sim}
(M_{(f)})_{[\varphi_{(f)}]}$. Il correspond canoniquement à cet isomorphisme
un isomorphisme de faisceaux
$\widetilde{M_{[\varphi]}}|D_+(f')
\xrightarrow{\sim}(\Phi_f)_*(\widetilde{M}|D_+(f))$
(\hyperref[II.2.5.2-fr]{2.5.2} et
\hyperref[I.1.6.3-fr]{I, 1.6.3}). En outre,

\oldpage[II]{44}
si $g'$ est un second élément homogène de $S'_+$ et $g=\varphi(g')$, le
diagramme
\[
  \xymatrix{
    (M_{[\varphi]})_{(f')}\ar[r]^{\sim}\ar[d]
    & (M_{(f)})_{[\varphi_{(f)}]}\ar[d]\\
    (M_{[\varphi]})_{(f'g')}\ar[r]^{\sim}
    & (M_{(fg)})_{[\varphi_{(fg)}]}
  }
\]
est commutatif, d'où on conclut aussitôt que l'isomorphisme
\[
  \widetilde{M_{[\varphi]}}|D_+(f'g')
  \xrightarrow{\sim}(\Phi_{fg})_*(\widetilde{M}|D_+(fg))
\]
est la restriction à $D_+(f'g')$ de l'isomorphisme
$\widetilde{M_{[\varphi]}}|D_+(f')
\xrightarrow{\sim}(\Phi_f)_*(\widetilde{M}|D_+(f))$. Comme $\Phi_f$ est la
restriction à $D_+(f)$ du morphisme $\Phi$, on voit donc, compte tenu de
(\hyperref[II.2.8.1.1-fr]{2.8.1.1}), et en posant
$X'=\operatorname{Proj}(S')$~:
\end{env}

\begin{proposition}[2.8.7]
\label{II.2.8.7-fr}
Il existe un isomorphisme canonique fonctoriel du $\mathcal{O}_{X'}$-Module
$\widetilde{M_{[\varphi]}}$ sur le
$\mathcal{O}_{X'}$-Module $\Phi_*(\widetilde{M}|G(\varphi))$.
\end{proposition}

On en déduit aussitôt une application canonique fonctorielle de l'ensemble
des $\varphi$-morphismes $M'\to M$ d'un $S'$-module gradué dans le $S$-module
gradué $M$, dans l'ensemble des $\Phi$-morphismes
$\widetilde{M'}\to\widetilde{M}|G(\varphi)$. Avec les notations de
(\hyperref[II.2.8.4-fr]{2.8.4}), si $M''$ est un $S''$-module gradué, au
composé d'un $\varphi$-morphisme $M'\to M$ et d'un $\varphi'$-morphisme
$M''\to M'$ correspond canoniquement le composé de
$\widetilde{M'}|G(\varphi')\to\widetilde{M}|G(\varphi'')$ et de
$\widetilde{M''}\to\widetilde{M'}|G(\varphi')$.

\begin{proposition}[2.8.8]
\label{II.2.8.8-fr}
Sous les hypothèses de (\hyperref[II.2.8.1-fr]{2.8.1}), soit $M'$ un
$S'$-module gradué. Il existe un homomorphisme canonique fonctoriel $\nu$ du
$(\mathcal{O}_X|G(\varphi))$-Module $\Phi^*(\widetilde{M'})$ dans le
$(\mathcal{O}_X|G(\varphi))$-Module
$\widetilde{M'\otimes_{S'}S}|G(\varphi)$. Si l'idéal $S'_+$
est engendré par $S'_1$, $\nu$ est un isomorphisme.
\end{proposition}

\begin{proof}
En effet, pour $f'\in S'_d$ ($d>0$), on définit un homomorphisme canonique
fonctoriel de $S_{(f)}$-modules (où $f=\varphi(f')$)
\[
\label{II.2.8.8.1-fr}
  \nu_f:M'_{(f')}\otimes_{S'_{(f')}}S_{(f)}
  \longrightarrow(M'\otimes_{S'}S)_{(f)}
\tag{2.8.8.1}
\]
en composant l'homomorphisme
$M'_{(f')}\otimes_{S'_{(f')}}S_{(f)}\to
M'_{f'}\otimes_{S'_{f'}}S_f$ et l'isomorphisme canonique
$M'_{f'}\otimes_{S'_{f'}}S_f\xrightarrow{\sim}(M'\otimes_{S'}S)_f$
(\hyperref[0.1.5.4-fr]{0, 1.5.4}) et notant que ce dernier conserve les
degrés. On vérifie aussitôt la compatibilité des $\nu_f$ avec les opérateurs
de restriction de $D_+(f)$ à $D_+(fg)$ (pour un second $g'\in S'_+$ et
$g=\varphi(g')$), d'où la définition de l'homomorphisme
\[
  \nu:\Phi^*(\widetilde{M'})
  \longrightarrow\widetilde{M'\otimes_{S'}S}|G(\varphi)
\]
compte tenu de (\hyperref[I.1.6.5-fr]{I, 1.6.5}). Pour prouver la seconde
assertion, il suffit de montrer que $\nu_f$ est un isomorphisme pour tout
$f'\in S'_1$, puisque $G(\varphi)$ est alors réunion des
$D_+(\varphi(f'))$. On définit d'abord une application $\mathbf{Z}$-bilinéaire
$M'_m\times S_n\to M'_{(f')}\otimes_{S'_{(f')}}S_{(f)}$ en faisant
correspondre à $(x',s)$ l'élément
$(x'/{f'}^m)\otimes(s/f^n)$ (avec la convention que $x'/{f'}^m$ est
${f'}^{-m}x'/1$

\oldpage[II]{45}
lorsque $m<0$). On constate comme dans la démonstration de
(\hyperref[II.2.5.13-fr]{2.5.13}) que cette application donne naissance à un
di-homomorphisme de modules
\[
  \eta_f:M'\otimes_{S'}S
  \longrightarrow M'_{(f')}\otimes_{S'_{(f')}}S_{(f)}.
\]
En outre, si, pour $r>0$, on a
$f^r\sum_i(x'_i\otimes s_i)=0$, cela s'écrit aussi
$\sum_i({f'}^rx'_i\otimes s_i)=0$, d'où, par
(\hyperref[0.1.5.4-fr]{0, 1.5.4}),
$\sum_i({f'}^rx'_i/{f'}^{m_i+r})\otimes(s_i/f^{n_i})=0$, c'est-à-dire
$\eta_f(\sum_i x'_i\otimes s_i)=0$, ce qui prouve que $\eta_f$ se factorise
en
\[
  M'\otimes_{S'}S\longrightarrow(M'\otimes_{S'}S)_f
  \xrightarrow{\eta'_f}M'_{(f')}\otimes_{S'_{(f')}}S_{(f)}~;
\]
on vérifie enfin que $\eta'_f$ et $\nu_f$ sont des isomorphismes réciproques
l'un de l'autre.

En particulier, il résulte de (\hyperref[II.2.1.2.1-fr]{2.1.2.1}) que l'on a
un homomorphisme canonique
\[
\label{II.2.8.8.2-fr}
  \Phi^*(\mathcal{O}_{X'}(n))\xrightarrow{\sim}
  \mathcal{O}_X(n)|G(\varphi)
\tag{2.8.8.2}
\]
pour tout $n\in\mathbf{Z}$.
\end{proof}

\begin{env}[2.8.9]
\label{II.2.8.9-fr}
Soient $A$, $A'$ deux anneaux, $\psi:A'\to A$ un homomorphisme d'anneaux,
définissant un morphisme
$\Psi:\operatorname{Spec}(A)\to\operatorname{Spec}(A')$. Soit $S'$ une
$A'$-algèbre graduée à degrés positifs, et posons $S=S'\otimes_{A'}A$, qui
est de façon évidente une $A$-algèbre graduée par les $S'_n\otimes_{A'}A$~;
l'application $\varphi:s'\mapsto s'\otimes1$ est alors un homomorphisme
d'anneaux gradués rendant commutatif le diagramme
(\hyperref[II.2.8.5.1-fr]{2.8.5.1}). Comme ici $S_+$ est le $A$-module
engendré par $\varphi(S'_+)$, on a $G(\varphi)=\operatorname{Proj}(S)=X$~;
d'où, en posant $X'=\operatorname{Proj}(S')$, un diagramme commutatif
\[
\label{II.2.8.9.1-fr}
  \xymatrix{
    X\ar[r]^{\Phi}\ar[d]_p & X'\ar[d]\\
    Y\ar[r]_{\Psi} & Y'
  }
\tag{2.8.9.1}
\]

Soit maintenant $M'$ un $S'$-module gradué, et posons
$M=M'\otimes_{A'}A=M'\otimes_{S'}S$. Dans ces conditions~:
\end{env}

\begin{proposition}[2.8.10]
\label{II.2.8.10-fr}
Le diagramme (\hyperref[II.2.8.9.1-fr]{2.8.9.1}) identifie le schéma $X$ au
produit $X'\times_{Y'}Y$~; en outre, l'homomorphisme canonique
$\nu:\Phi^*(\widetilde{M'})\to\widetilde{M}$
(\hyperref[II.2.8.8-fr]{2.8.8}) est un isomorphisme.
\end{proposition}

\begin{proof}
La première assertion sera démontrée si l'on prouve que pour tout $f'$
homogène dans $S'_+$, en posant $f=\varphi(f')$, les restrictions de $\Phi$
et de $p$ à $D_+(f)$ identifient ce schéma au produit
$D_+(f')\times_{Y'}Y$ (\hyperref[I.3.2.6.2-fr]{I, 3.2.6.2})~; autrement
dit, il suffit (\hyperref[I.3.2.2-fr]{I, 3.2.2}) de prouver que $S_{(f)}$
s'identifie canoniquement à $S'_{(f')}\otimes_{A'}A$, ce qui est immédiat
en vertu de l'existence de l'isomorphisme canonique
$S_f\xrightarrow{\sim}S'_{f'}\otimes_{A'}A$ conservant les degrés
(\hyperref[0.1.5.4-fr]{0, 1.5.4}). La seconde assertion résulte de ce que
$M'_{(f')}\otimes_{S'_{(f')}}S_{(f)}$ s'identifie d'après ce qui précède à
$M'_{(f')}\otimes_{A'}A$, et ce dernier à $M_{(f)}$, puisque $M_f$
s'identifie canoniquement à $M'_{f'}\otimes_{A'}A$ par un isomorphisme
conservant les degrés.
\end{proof}

\begin{corollary}[2.8.11]
\label{II.2.8.11-fr}
Pour tout entier $n\in\mathbf{Z}$, $\widetilde{M(n)}$ s'identifie à
$\Phi^*(\widetilde{M'(n)})=\widetilde{M'(n)}\otimes_{Y'}\mathcal{O}_Y$~;
en particulier
$\mathcal{O}_X(n)=\Phi^*(\mathcal{O}_{X'}(n))
=\mathcal{O}_{X'}(n)\otimes_{Y'}\mathcal{O}_Y$.
\end{corollary}

\begin{proof}
Cela résulte de (\hyperref[II.2.8.10-fr]{2.8.10}) et de
(\hyperref[II.2.5.15-fr]{2.5.15}).
\end{proof}

\oldpage[II]{46}

\begin{env}[2.8.12]
\label{II.2.8.12-fr}
Sous les hypothèses de (\hyperref[II.2.8.9-fr]{2.8.9}), pour
$f'\in S'_d$ ($d>0$) et $f=\varphi(f')$, le diagramme
\[
\label{II.2.8.12.1-fr}
  \xymatrix{
    M'_{(f')}\ar[r]^{\sim}\ar[d]
    & M'^{(d)}/(f'-1)M'^{(d)}\ar[d]\\
    M_{(f)}\ar[r]^{\sim}
    & M^{(d)}/(f-1)M^{(d)}
  }
\tag{2.8.12.1}
\]
(cf. (\hyperref[II.2.2.5-fr]{2.2.5})) est commutatif.
\end{env}

\begin{env}[2.8.13]
\label{II.2.8.13-fr}
Gardons les notations et hypothèses de (\hyperref[II.2.8.9-fr]{2.8.9}), et
soit $\mathcal{F}'$ un $\mathcal{O}_{X'}$-Module~; si on pose
$\mathcal{F}=\Phi^*(\mathcal{F}')$, on a, pour tout $n\in\mathbf{Z}$,
$\mathcal{F}(n)=\Phi^*(\mathcal{F}'(n))$ en vertu de
(\hyperref[II.2.8.11-fr]{2.8.11}) et
(\hyperref[0.4.3.3-fr]{0, 4.3.3}). Par suite
(\hyperref[0.3.7.1-fr]{0, 3.7.1}), on a un homomorphisme canonique
\[
  \Gamma(\rho):\Gamma(X',\mathcal{F}'(n))
  \longrightarrow\Gamma(X,\mathcal{F}(n))
\]
ce qui donne un di-homomorphisme canonique de modules gradués
\[
  \Gamma_\bullet(\mathcal{F}')\longrightarrow
  \Gamma_\bullet(\mathcal{F}).
\]

Supposons l'idéal $S_+$ engendré par $S_1$, et
$\mathcal{F}'=\widetilde{M'}$, donc $\mathcal{F}=\widetilde{M}$ avec
$M=M'\otimes_{A'}A$. Si $f'$ est homogène dans $S'_+$ et
$f=\varphi(f')$, on a vu que
$M_{(f)}=M'_{(f')}\otimes_{A'}A$, et le diagramme
\[
  \xymatrix{
    M'_0\ar[r]\ar[d]
    & M'_{(f')}\ar[d]
    & =\Gamma(D_+(f'),\widetilde{M'})\\
    M_0\ar[r]
    & M_{(f)}
    & =\Gamma(D_+(f),\widetilde{M})
  }
\]
est donc commutatif~; on conclut aussitôt de cette remarque et de la
définition de l'homomorphisme $\alpha$
(\hyperref[II.2.6.2-fr]{2.6.2}) que le diagramme
\[
\label{II.2.8.13.1-fr}
  \xymatrix{
    M'\ar[r]^{\alpha_{M'}}\ar[d]
    & \Gamma_\bullet(\widetilde{M'})\ar[d]\\
    M\ar[r]_{\alpha_M}
    & \Gamma_\bullet(\widetilde{M})
  }
\tag{2.8.13.1}
\]
est commutatif. De même, le diagramme
\[
\label{II.2.8.13.2-fr}
  \xymatrix{
    \widetilde{\Gamma_\bullet(\mathcal{F}')}\ar[r]^{\beta_{\mathcal{F}'}}
      \ar[d]
    & \mathcal{F}'\ar[d]\\
    \widetilde{\Gamma_\bullet(\mathcal{F})}\ar[r]_{\beta_{\mathcal{F}}}
    & \mathcal{F}
  }
\tag{2.8.13.2}
\]
est commutatif (la flèche verticale de droite étant le $\Phi$-morphisme
canonique
$\mathcal{F}'\to\Phi^*(\mathcal{F}')=\mathcal{F}$).
\end{env}

\oldpage[II]{47}

\begin{env}[2.8.14]
\label{II.2.8.14-fr}
Conservant toujours les hypothèses et notations de
(\hyperref[II.2.8.9-fr]{2.8.9}), soit $N'$ un second $S'$-module gradué, et
soit $N=N'\otimes_{A'}A$. Il est immédiat que les di-homomorphismes
canoniques $M'\to M$, $N'\to N$ donnent un di-homomorphisme
$M'\otimes_{S'}N'\to M\otimes_S N$ (relatif à l'homomorphisme canonique
d'anneaux $S'\to S$), et par suite aussi un $S$-homomorphisme de degré~$0$
\[
  (M'\otimes_{S'}N')\otimes_{A'}A\longrightarrow M\otimes_S N,
\]
auquel correspond (compte tenu de
(\hyperref[II.2.8.10-fr]{2.8.10})) un homomorphisme de
$\mathcal{O}_X$-Modules
\[
\label{II.2.8.14.1-fr}
  \Phi^*(\widetilde{M'\otimes_{S'}N'})
  \longrightarrow\widetilde{M\otimes_S N}.
\tag{2.8.14.1}
\]

En outre, on vérifie aussitôt que le diagramme
\[
\label{II.2.8.14.2-fr}
  \xymatrix{
    \Phi^*(\widetilde{M'}\otimes_{\mathcal{O}_{X'}}\widetilde{N'})
      \ar[r]^{\sim}\ar[d]_{\Phi^*(\lambda)}
    & \widetilde{M}\otimes_{\mathcal{O}_X}\widetilde{N}
      \ar[d]^{\lambda}
    & =\Phi^*(\widetilde{M'})\otimes_{\mathcal{O}_X}
      \Phi^*(\widetilde{N'})\\
    \Phi^*(\widetilde{M'\otimes_{S'}N'})\ar[r]
    & \widetilde{M\otimes_S N}
  }
\tag{2.8.14.2}
\]
est commutatif, la première ligne étant l'isomorphisme canonique
(\hyperref[0.4.3.3-fr]{0, 4.3.3}). Si l'idéal $S'_+$ est engendré par
$S'_1$, il est clair que $S_+$ est engendré par $S_1$, et les deux flèches
verticales de (\hyperref[II.2.8.14.2-fr]{2.8.14.2}) sont alors des
isomorphismes (\hyperref[II.2.5.13-fr]{2.5.13})~; il en est donc de même de
(\hyperref[II.2.8.14.1-fr]{2.8.14.1}).

On a de même un di-homomorphisme canonique
$\operatorname{Hom}_{S'}(M',N')\to\operatorname{Hom}_S(M,N)$ en faisant
correspondre à un homomorphisme $u'$ de degré $k$ l'homomorphisme
$u'\otimes1$, qui est aussi de degré $k$~; on en déduit encore un
$S$-homomorphisme de degré~$0$
\[
  (\operatorname{Hom}_{S'}(M',N'))\otimes_{A'}A
  \longrightarrow\operatorname{Hom}_S(M,N)
\]
d'où un homomorphisme de $\mathcal{O}_X$-Modules
\[
\label{II.2.8.14.3-fr}
  \Phi^*(\widetilde{\operatorname{Hom}_{S'}(M',N')})
  \longrightarrow\widetilde{\operatorname{Hom}_S(M,N)}.
\tag{2.8.14.3}
\]

En outre, le diagramme
\[
  \xymatrix{
    \Phi^*(\widetilde{\operatorname{Hom}_{S'}(M',N')})\ar[r]
      \ar[d]_{\Phi^*(\mu)}
    & \widetilde{\operatorname{Hom}_S(M,N)}\ar[d]^{\mu}\\
    \Phi^*(\mathcal{H}om_{\mathcal{O}_{X'}}
      (\widetilde{M'},\widetilde{N'}))\ar[r]
    & \mathcal{H}om_{\mathcal{O}_X}(\widetilde{M},\widetilde{N})
  }
\]
est commutatif (la seconde ligne horizontale étant l'homomorphisme canonique
(\hyperref[0.4.4.6-fr]{0, 4.4.6})).
\end{env}

\oldpage[II]{48}

\begin{env}[2.8.15]
\label{II.2.8.15-fr}
Les notations et hypothèses étant celles de
(\hyperref[II.2.8.1-fr]{2.8.1}), il résulte de
(\hyperref[II.2.4.7-fr]{2.4.7}) que l'on ne change pas le morphisme $\Phi$,
à des isomorphismes près, en remplaçant $S_0$ et $S'_0$ par
$\mathbf{Z}$ et $\varphi_0$ par l'application identique, puis en remplaçant
$S$ et $S'$ par $S^{(d)}$ et ${S'}^{(d)}$ respectivement ($d>0$) et
$\varphi$ par sa restriction $\varphi^{(d)}$ à $S^{(d)}$.
\end{env}

\subsection{Sous-schémas fermés d'un schéma $\operatorname{Proj}(S)$.}
\label{subsection:II.2.9-fr}

\begin{env}[2.9.1]
\label{II.2.9.1-fr}
Si $\varphi:S\to S'$ est un homomorphisme d'anneaux gradués, on dit que
$\varphi$ est (TN)-\emph{surjectif} (resp. (TN)-\emph{injectif},
(TN)-\emph{bijectif}) s'il existe un entier $n$ tel que, pour $k\geq n$,
$\varphi_k:S_k\to S'_k$ soit surjectif (resp. injectif, bijectif). Alors il
résulte de (\hyperref[II.2.8.15-fr]{2.8.15}) que l'étude de $\Phi$ peut se
ramener au cas où $\varphi$ est \emph{surjectif} (resp. \emph{injectif},
\emph{bijectif}). Au lieu de dire que $\varphi$ est (TN)-bijectif, on dit aussi
que c'est alors un (TN)-\emph{isomorphisme}.
\end{env}

\begin{proposition}[2.9.2]
\label{II.2.9.2-fr}
Soit $S$ un anneau gradué à degrés positifs et soit
$X=\operatorname{Proj}(S)$.
\begin{enumerate}
  \item[\rm(i)] Si $\varphi:S\to S'$ est un homomorphisme (TN)-surjectif
    d'anneaux gradués, le morphisme correspondant $\Phi$
    (\hyperref[II.2.8.1-fr]{2.8.1}) est défini dans
    $\operatorname{Proj}(S')$ tout entier et est une immersion fermée de
    $\operatorname{Proj}(S')$ dans $X$. Si $\mathfrak{J}$ est le noyau de
    $\varphi$, le sous-schéma fermé de $X$ associé à $\Phi$ est défini par
    le faisceau d'idéaux quasi-cohérent $\widetilde{\mathfrak{J}}$ de
    $\mathcal{O}_X$.
  \item[\rm(ii)] Supposons en outre l'idéal $S_+$ engendré par un nombre
    fini d'éléments homogènes de degré~$1$. Soit $X'$ un sous-schéma fermé
    de $X$, défini par un faisceau quasi-cohérent d'idéaux $\mathcal{J}$ de
    $\mathcal{O}_X$. Soit $\mathfrak{J}$ l'idéal gradué de $S$, image
    réciproque de $\Gamma_\bullet(\mathcal{J})$ par l'homomorphisme canonique
    $\alpha:S\to\Gamma_\bullet(\mathcal{O}_X)$
    (\hyperref[II.2.6.2-fr]{2.6.2}), et posons
    $S'=S/\mathfrak{J}$. Alors $X'$ est le sous-schéma associé à
    l'immersion fermée $\operatorname{Proj}(S')\to X$ correspondant à
    l'homomorphisme canonique d'anneaux gradués $S\to S'$.
\end{enumerate}
\end{proposition}

\begin{proof}
\begin{enumerate}
  \item[\rm(i)] On peut supposer $\varphi$ surjectif
    (\hyperref[II.2.9.1-fr]{2.9.1}). Comme par hypothèse
    $\varphi(S_+)$ engendre $S'_+$, on a
    $G(\varphi)=\operatorname{Proj}(S')$. D'autre part, la seconde assertion
    se vérifie localement sur $X$~; soit donc $f$ un élément homogène de
    $S_+$ et posons $f'=\varphi(f)$. Comme $\varphi$ est un homomorphisme
    surjectif d'anneaux gradués, on vérifie aussitôt que
    $\varphi_{(f')}:S_{(f)}\to S'_{(f')}$ est surjectif et que son noyau est
    $\mathfrak{J}_{(f)}$, ce qui achève de prouver (i)
    (\hyperref[I.4.2.3-fr]{I, 4.2.3}).
  \item[\rm(ii)] En vertu de (i), on est ramené à vérifier que
    l'homomorphisme $\widetilde{j}:\widetilde{\mathfrak{J}}\to\mathcal{O}_X$
    déduit de l'injection canonique $j:\mathfrak{J}\to S$ est un
    isomorphisme de $\widetilde{\mathfrak{J}}$ sur $\mathcal{J}$, ce qui
    résulte de (\hyperref[II.2.7.11-fr]{2.7.11}).
\end{enumerate}
\end{proof}

On notera que $\mathfrak{J}$ est le \emph{plus grand} des idéaux gradués
$\mathfrak{J}'$ de $S$ tels que
$\widetilde{j}(\widetilde{\mathfrak{J}'})=\mathcal{J}$, car on vérifie
immédiatement en remontant aux définitions
(\hyperref[II.2.6.2-fr]{2.6.2}) que cette relation implique
$\alpha(\mathfrak{J}')\subset\Gamma_\bullet(\mathcal{J})$.

\begin{corollary}[2.9.3]
\label{II.2.9.3-fr}
On suppose vérifiées les hypothèses de
(\hyperref[II.2.9.2-fr]{2.9.2}, (i)) et en outre que l'idéal $S_+$ est
engendré par $S_1$~; alors $\Phi^*(\widetilde{S(n)})$ est canoniquement
isomorphe à $\widetilde{S'(n)}$ pour tout $n\in\mathbf{Z}$, et par suite
$\Phi^*(\mathcal{F}(n))$ à $\Phi^*(\mathcal{F})(n)$ pour tout
$\mathcal{O}_X$-Module $\mathcal{F}$.
\end{corollary}

\begin{proof}
C'est un cas particulier de (\hyperref[II.2.8.8-fr]{2.8.8}), compte tenu de
(\hyperref[II.2.5.10.2-fr]{2.5.10.2}).
\end{proof}

\begin{corollary}[2.9.4]
\label{II.2.9.4-fr}
On suppose vérifiées les hypothèses de
(\hyperref[II.2.9.2-fr]{2.9.2}, (ii)). Pour que le sous-préschéma fermé $X'$
de $X$ soit intègre, il faut et il suffit que l'idéal gradué
$\mathfrak{J}$ soit premier dans $S$.
\end{corollary}

\oldpage[II]{49}

\begin{proof}
Comme $X'$ est isomorphe à $\operatorname{Proj}(S/\mathfrak{J})$, la
condition est suffisante en vertu de (\hyperref[II.2.4.4-fr]{2.4.4}). Pour
voir qu'elle est nécessaire, considérons la suite exacte
$0\to\mathcal{J}\to\mathcal{O}_X\to\mathcal{O}_X/\mathcal{J}$, qui donne une
suite exacte
\[
  0\longrightarrow\Gamma_\bullet(\mathcal{J})
  \longrightarrow\Gamma_\bullet(\mathcal{O}_X)
  \longrightarrow\Gamma_\bullet(\mathcal{O}_X/\mathcal{J}).
\]
Il suffit de prouver que si $f\in S_m$, $g\in S_n$ sont tels que l'image dans
$\Gamma_\bullet(\mathcal{O}_X/\mathcal{J})$ de $\alpha_{n+m}(fg)$ est nulle,
alors l'une des images de $\alpha_m(f)$, $\alpha_n(g)$ est nulle. Or, par
définition, ces images sont des sections des
$(\mathcal{O}_X/\mathcal{J})$-Modules inversibles
$\mathcal{L}=(\mathcal{O}_X/\mathcal{J})(m)$ et
$\mathcal{L}'=(\mathcal{O}_X/\mathcal{J})(n)$ au-dessus du schéma intègre
$X'$~; l'hypothèse entraîne que le produit de ces deux sections est nul dans
$\mathcal{L}\otimes\mathcal{L}'$
((\hyperref[II.2.9.3-fr]{2.9.3}) et
(\hyperref[II.2.5.14.1-fr]{2.5.14.1})), donc l'une d'elles est nulle en vertu
de (\hyperref[I.7.4.4-fr]{I, 7.4.4}).
\end{proof}

\begin{corollary}[2.9.5]
\label{II.2.9.5-fr}
Soient $A$ un anneau, $M$ un $A$-module, $S$ une $A$-algèbre graduée
engendrée par l'ensemble $S_1$ des éléments homogènes de degré~$1$,
$u:M\to S_1$ un homomorphisme surjectif de $A$-modules,
$\overline{u}:\mathbf{S}(M)\to S$ l'homomorphisme (de $A$-algèbres) de
l'algèbre symétrique $\mathbf{S}(M)$ de $M$ dans $S$ qui prolonge $u$. Le
morphisme correspondant à $\overline{u}$ est alors une immersion fermée de
$\operatorname{Proj}(S)$ dans $\operatorname{Proj}(\mathbf{S}(M))$.
\end{corollary}

\begin{proof}
En effet, $\overline{u}$ est surjectif par hypothèse et il suffit d'appliquer
(\hyperref[II.2.9.2-fr]{2.9.2}).
\end{proof}

\section{Spectre homogène d'un faisceau d'algèbres graduées}
\label{section:II.3-fr}

\subsection{Spectre homogène d'une $\mathcal{O}_Y$-Algèbre graduée
quasi-cohérente.}
\label{subsection:II.3.1-fr}
\label{II.3.1-fr}

\begin{env}[3.1.1]
\label{II.3.1.1-fr}
Soient $Y$ un préschéma, $\mathcal{S}$ une $\mathcal{O}_Y$-Algèbre graduée,
$\mathcal{M}$ un $\mathcal{S}$-Module gradué. Si $\mathcal{S}$ est
\emph{quasi-cohérente}, chacun de ses composants homogènes $\mathcal{S}_n$ est
un $\mathcal{O}_Y$-Module \emph{quasi-cohérent}, étant l'image de
$\mathcal{S}$ par un homomorphisme de $\mathcal{S}$ dans lui-même
((\hyperref[I.1.3.8-fr]{I, 1.3.8}) et
(\hyperref[I.1.3.9-fr]{I, 1.3.9}))~; de même, si $\mathcal{M}$ est
quasi-cohérent en tant que $\mathcal{O}_Y$-Module, ses composants homogènes
$\mathcal{M}_n$ le sont aussi, et réciproquement. Si $d$ est un entier $>0$,
on désignera par $\mathcal{S}^{(d)}$ la somme directe des
$\mathcal{O}_Y$-Modules $\mathcal{S}_{nd}$ ($n\in\mathbf{Z}$), qui est
quasi-cohérente si $\mathcal{S}$ l'est
(\hyperref[I.1.3.9-fr]{I, 1.3.9})~; pour tout entier $k$ tel que
$0\leq k\leq d-1$, on désignera par $\mathcal{M}^{(d,k)}$ (ou
$\mathcal{M}^{(d)}$ pour $k=0$) la somme directe des
$\mathcal{M}_{nd+k}$ ($n\in\mathbf{Z}$), qui est un
$\mathcal{S}^{(d)}$-Module gradué, quasi-cohérent lorsque $\mathcal{S}$ et
$\mathcal{M}$ sont quasi-cohérents (\hyperref[I.9.6.1-fr]{I, 9.6.1}). On
notera $\mathcal{M}(n)$ le $\mathcal{S}$-Module gradué tel que
$(\mathcal{M}(n))_k=\mathcal{M}_{n+k}$ pour tout $k\in\mathbf{Z}$~; si
$\mathcal{S}$ et $\mathcal{M}$ sont quasi-cohérents, $\mathcal{M}(n)$ est un
$\mathcal{S}$-Module gradué quasi-cohérent
(\hyperref[I.9.6.1-fr]{I, 9.6.1}).

Nous dirons que $\mathcal{M}$ est un $\mathcal{S}$-Module gradué
\emph{de type fini} (resp. qu'il admet une \emph{présentation finie}) si pour
tout $y\in Y$, il existe un voisinage ouvert $U$ de $y$ et des entiers $n_i$
(resp. des entiers $m_i$ et $n_j$) tels qu'il y ait un homomorphisme surjectif
de degré~$0$
\[
  \bigoplus_{i=1}^r(\mathcal{S}(n_i)|U)\longrightarrow\mathcal{M}|U
\]
(resp. que $\mathcal{M}|U$ soit isomorphe au conoyau d'un homomorphisme de
degré~$0$
\[
  \bigoplus_{i=1}^r(\mathcal{S}(m_i)|U)
  \longrightarrow\bigoplus_{j=1}^s(\mathcal{S}(n_j)|U)).
\]

Soient $U$ un ouvert affine de $Y$, $A=\Gamma(U,\mathcal{O}_Y)$ son anneau~;
par hypothèse, la $(\mathcal{O}_Y|U)$-Algèbre graduée $\mathcal{S}|U$ est
isomorphe à $\widetilde{S}$, où $S=\Gamma(U,\mathcal{S})$ est une
$A$-algèbre

\oldpage[II]{50}

graduée (\hyperref[I.1.4.3-fr]{I, 1.4.3})~; posons
$X_U=\operatorname{Proj}(\Gamma(U,\mathcal{S}))$. Soient $U'\subset U$ un
second ouvert affine de $Y$, $A'$ son anneau, $j:U'\to U$ l'injection
canonique, qui correspond à l'homomorphisme de restriction $A\to A'$~; on a
$\mathcal{S}|U'=j^*(\mathcal{S}|U)$, et par suite
$S'=\Gamma(U',\mathcal{S})$ s'identifie canoniquement à $S\otimes_A A'$
(\hyperref[I.1.6.4-fr]{I, 1.6.4}). On en conclut
(\hyperref[II.2.8.10-fr]{2.8.10}) que $X_{U'}$ s'identifie canoniquement à
$X_U\times_U U'$, et par suite aussi à $f_U^{-1}(U')$, en désignant par
$f_U$ le morphisme structural $X_U\to U$
(\hyperref[I.4.4.1-fr]{I, 4.4.1}). Nous désignerons par $\sigma_{U',U}$
l'isomorphisme canonique $f_U^{-1}(U')\xrightarrow{\sim}X_{U'}$ ainsi défini,
par $\rho_{U',U}$ l'immersion ouverte $X_{U'}\to X_U$ obtenue en composant
$\sigma_{U',U}^{-1}$ et l'injection canonique $f_U^{-1}(U')\to X_U$. Il est
immédiat que si $U''\subset U'$ est un troisième ouvert affine de $Y$, on a
$\rho_{U'',U}=\rho_{U',U}\circ\rho_{U'',U'}$.
\end{env}

\begin{proposition}[3.1.2]
\label{II.3.1.2-fr}
Soit $Y$ un préschéma. Pour toute $\mathcal{O}_Y$-Algèbre graduée
quasi-cohérente $\mathcal{S}$ à degrés positifs, il existe un préschéma $X$
au-dessus de $Y$, et un seul à un $Y$-isomorphisme près, ayant la propriété
suivante~: si $f$ est le morphisme structural $X\to Y$, pour tout ouvert
affine $U$ de $Y$, il existe un \emph{isomorphisme} $\eta_U$ du préschéma
induit $f^{-1}(U)$ sur $X_U=\operatorname{Proj}(\Gamma(U,\mathcal{S}))$ tel
que, si $V$ est un second ouvert affine de $Y$ contenu dans $U$, le diagramme
\[
\label{II.3.1.2.1-fr}
  \xymatrix{
    f^{-1}(V)\ar[r]^{\eta_V}\ar[d]
    & X_V\ar[d]^{\rho_{V,U}}\\
    f^{-1}(U)\ar[r]_{\eta_U}
    & X_U
  }
\tag{3.1.2.1}
\]
soit commutatif.
\end{proposition}

\begin{proof}
Pour deux ouverts affines $U$, $V$ de $Y$, soit $X_{U,V}$ le préschéma
induit sur $f_U^{-1}(U\cap V)$ par $X_U$~; nous allons définir un
$Y$-isomorphisme $\theta_{U,V}:X_{V,U}\xrightarrow{\sim}X_{U,V}$. Pour cela,
considérons un ouvert affine $W\subset U\cap V$~: en composant les
isomorphismes
\[
  f_U^{-1}(W)\xrightarrow{\sigma_{W,U}}X_W
  \xrightarrow{\sigma_{W,V}^{-1}}f_V^{-1}(W),
\]
on obtient un isomorphisme $\tau_W$, et on vérifie aussitôt que si
$W'\subset W$ est un ouvert affine, $\tau_{W'}$ est la restriction à
$f_U^{-1}(W')$ de $\tau_W$~; les $\tau_W$ sont donc bien les restrictions
d'un $Y$-isomorphisme $\theta_{V,U}$. En outre, si $U$, $V$, $W$ sont trois
ouverts affines de $Y$, $\theta'_{U,V}$, $\theta'_{V,W}$ et
$\theta'_{U,W}$ les restrictions de $\theta_{U,V}$, $\theta_{V,W}$,
$\theta_{U,W}$ aux images réciproques de $U\cap V\cap W$ dans $X_V$, $X_W$,
$X_W$ respectivement, il résulte des définitions précédentes que l'on a
$\theta'_{U,V}\circ\theta'_{V,W}=\theta'_{U,W}$. L'existence de $X$
vérifiant les propriétés énoncées résulte donc de
(\hyperref[I.2.3.1-fr]{I, 2.3.1})~; son unicité à un $Y$-isomorphisme près
est triviale, compte tenu de (\hyperref[II.3.1.2.1-fr]{3.1.2.1}).
\end{proof}

\begin{env}[3.1.3]
\label{II.3.1.3-fr}
Nous dirons que le préschéma $X$ défini dans
(\hyperref[II.3.1.2-fr]{3.1.2}) est le \emph{spectre homogène} de la
$\mathcal{O}_Y$-Algèbre graduée quasi-cohérente $\mathcal{S}$ et nous le
noterons $\operatorname{Proj}(\mathcal{S})$. Il est immédiat que
$\operatorname{Proj}(\mathcal{S})$ est \emph{séparé au-dessus de $Y$}
((\hyperref[II.2.4.2-fr]{2.4.2}) et
(\hyperref[I.5.5.5-fr]{I, 5.5.5}))~; si $\mathcal{S}$ est une
$\mathcal{O}_Y$-Algèbre \emph{de type fini}
(\hyperref[I.9.6.2-fr]{I, 9.6.2}), $\operatorname{Proj}(\mathcal{S})$ est
\emph{de type fini} sur $Y$ ((\hyperref[II.2.7.1-fr]{2.7.1, (ii)}) et
(\hyperref[I.6.3.1-fr]{I, 6.3.1})).

Si $f$ est le morphisme structural $X\to Y$, il est immédiat que pour tout
préschéma induit par $Y$ sur un ouvert $U$ de $Y$, $f^{-1}(U)$ s'identifie au
spectre homogène $\operatorname{Proj}(\mathcal{S}|U)$.
\end{env}

\begin{proposition}[3.1.4]
\label{II.3.1.4-fr}
Soit $f\in\Gamma(Y,\mathcal{S}_d)$ pour $d>0$. Il existe une partie ouverte
$X_f$ de l'espace sous-jacent à $X=\operatorname{Proj}(\mathcal{S})$ ayant la
propriété suivante~: pour tout ouvert affine $U$ de $Y$,

\oldpage[II]{51}

$X_f\cap\varphi^{-1}(U)=D_+(f|U)$ dans $\varphi^{-1}(U)$ identifié à
$X_U=\operatorname{Proj}(\Gamma(U,\mathcal{S}))$ ($\varphi$ désignant le
morphisme structural $X\to Y$). En outre, le préschéma induit sur $X_f$ est
affine au-dessus de $Y$ et est canoniquement isomorphe à
$\operatorname{Spec}(\mathcal{S}^{(d)}/(f-1)\mathcal{S}^{(d)})$
(\hyperref[II.1.3.1-fr]{1.3.1}).
\end{proposition}

\begin{proof}
On a $f|U\in\Gamma(U,\mathcal{S}_d)=(\Gamma(U,\mathcal{S}))_d$. Si $U$,
$U'$ sont deux ouverts affines de $Y$ tels que $U'\subset U$, $f|U'$ est
l'image de $f|U$ par l'homomorphisme de restriction
\[
  \Gamma(U,\mathcal{S})\longrightarrow\Gamma(U',\mathcal{S}),
\]
donc $D_+(f|U')$ est égal (avec les notations de
(\hyperref[II.3.1.1-fr]{3.1.1})) au préschéma induit sur l'image réciproque
$\rho_{U',U}^{-1}(D_+(f|U))$ dans $X_{U'}$
(\hyperref[II.2.8.1-fr]{2.8.1})~; d'où la première assertion. En outre, le
préschéma induit sur $D_+(f|U)$ par $X_U$ s'identifie canoniquement à
$\operatorname{Spec}((\Gamma(U,\mathcal{S}))_{(f|U)})$, ces identifications
étant compatibles avec les homomorphismes de restriction
(\hyperref[II.2.8.1-fr]{2.8.1})~; la seconde assertion résulte alors de
(\hyperref[II.2.2.5-fr]{2.2.5}) et de la commutativité du diagramme
(\hyperref[II.2.8.2.1-fr]{2.8.2.1}).
\end{proof}

On dira encore que $X_f$ (en tant qu'ouvert dans l'espace sous-jacent $X$)
est l'ensemble des $x\in X$ où $f$ \emph{ne s'annule pas}.

\begin{corollary}[3.1.5]
\label{II.3.1.5-fr}
Si $f\in\Gamma(Y,\mathcal{S}_d)$, $g\in\Gamma(Y,\mathcal{S}_e)$, on a
\[
\label{II.3.1.5.1-fr}
  X_{fg}=X_f\cap X_g.
\tag{3.1.5.1}
\]
\end{corollary}

\begin{proof}
Il suffit en effet de considérer l'intersection des deux membres avec un
ensemble $\varphi^{-1}(U)$, où $U$ est un ouvert affine dans $Y$, et
d'appliquer la formule (\hyperref[II.2.3.3.2-fr]{2.3.3.2}).
\end{proof}

\begin{corollary}[3.1.6]
\label{II.3.1.6-fr}
Soit $(f_\alpha)$ une famille de sections de $\mathcal{S}$ au-dessus de $Y$
telle que $f_\alpha\in\Gamma(Y,\mathcal{S}_{d_\alpha})$~; si le faisceau
d'idéaux de $\mathcal{S}$ engendré par cette famille
(\hyperref[0.5.1.1-fr]{0, 5.1.1}) contient tous les $\mathcal{S}_n$ à partir
d'un certain rang, l'espace sous-jacent $X$ est réunion des $X_{f_\alpha}$.
\end{corollary}

\begin{proof}
En effet, pour tout ouvert affine $U$ de $Y$, $\varphi^{-1}(U)$ est réunion
des $X_{f_\alpha}\cap\varphi^{-1}(U)$
(\hyperref[II.2.3.14-fr]{2.3.14}).
\end{proof}

\begin{corollary}[3.1.7]
\label{II.3.1.7-fr}
Soit $\mathcal{A}$ une $\mathcal{O}_Y$-Algèbre quasi-cohérente~; posons
\[
  \mathcal{S}=\mathcal{A}[T]
  =\mathcal{A}\otimes_{\mathbf{Z}}\mathbf{Z}[T]
\]
où $T$ est une indéterminée ($\mathbf{Z}$ et $\mathbf{Z}[T]$ étant
considérés comme faisceaux simples sur $Y$). Alors
$X=\operatorname{Proj}(\mathcal{S})$ s'identifie canoniquement à
$\operatorname{Spec}(\mathcal{A})$. En particulier,
$\operatorname{Proj}(\mathcal{O}_Y[T])$ s'identifie à $Y$.
\end{corollary}

\begin{proof}
En appliquant (\hyperref[II.3.1.6-fr]{3.1.6}) à l'unique section
$f\in\Gamma(Y,\mathcal{S})$ égale à $T$ en chaque point de $Y$, on voit que
$X_f=X$. En outre, on a ici $d=1$, et
$\mathcal{S}^{(1)}/(f-1)\mathcal{S}^{(1)}
=\mathcal{S}/(f-1)\mathcal{S}$ est canoniquement isomorphe à $\mathcal{A}$,
d'où le corollaire (\hyperref[II.1.2.2-fr]{1.2.2}).
\end{proof}

Soit $g\in\Gamma(Y,\mathcal{O}_Y)$~; si on prend
$\mathcal{S}=\mathcal{O}_Y[T]$, on a donc $g\in\Gamma(Y,\mathcal{S}_0)$~; soit
\[
  h=gT\in\Gamma(Y,\mathcal{S}_1).
\]
Si $X=\operatorname{Proj}(\mathcal{S})$, l'identification canonique définie
dans (\hyperref[II.3.1.7-fr]{3.1.7}) identifie $X_h$ à l'ensemble ouvert
$Y_g$ de $Y$ (au sens de (\hyperref[0.5.5.2-fr]{0, 5.5.2}))~: en effet, on
peut se borner au cas où $Y=\operatorname{Spec}(A)$ est affine, et tout se
ramène alors (compte tenu de (\hyperref[II.2.2.5-fr]{2.2.5})) au fait que
l'anneau de fractions $A_g$ s'identifie canoniquement à
$A[T]/(gT-1)A[T]$ (\hyperref[0.1.2.3-fr]{0, 1.2.3}).

\begin{proposition}[3.1.8]
\label{II.3.1.8-fr}
Soit $\mathcal{S}$ une $\mathcal{O}_Y$-Algèbre graduée quasi-cohérente à
degrés positifs.
\begin{enumerate}
  \item[\rm(i)] Pour tout $d>0$, il existe un $Y$-isomorphisme canonique de
    $\operatorname{Proj}(\mathcal{S})$ sur
    $\operatorname{Proj}(\mathcal{S}^{(d)})$.

\oldpage[II]{52}

  \item[\rm(ii)] Soit $\mathcal{S}'$ la $\mathcal{O}_Y$-Algèbre graduée
    somme directe de $\mathcal{O}_Y$ et des $\mathcal{S}_n$ ($n\geq0$)~;
    alors $\operatorname{Proj}(\mathcal{S}')$ et
    $\operatorname{Proj}(\mathcal{S})$ sont canoniquement $Y$-isomorphes.
  \item[\rm(iii)] Soit $\mathcal{L}$ un $\mathcal{O}_Y$-Module inversible
    (\hyperref[0.5.4.1-fr]{0, 5.4.1}) et soit $\mathcal{S}_{(\mathcal{L})}$ la
    $\mathcal{O}_Y$-Algèbre graduée somme directe des
    $\mathcal{S}_d\otimes\mathcal{L}^{\otimes d}$ ($d\geq0$)~; alors
    $\operatorname{Proj}(\mathcal{S})$ et
    $\operatorname{Proj}(\mathcal{S}_{(\mathcal{L})})$ sont canoniquement
    $Y$-isomorphes.
\end{enumerate}
\end{proposition}

\begin{proof}
Dans chacun des trois cas, il suffit de définir l'isomorphisme localement sur
$Y$, la vérification des compatibilités avec les opérations de restriction
d'un ouvert à un ouvert plus petit étant triviale. On peut donc supposer $Y$
affine, et alors (i) résulte de (\hyperref[II.2.4.7-fr]{2.4.7, (i)}) et (ii)
de (\hyperref[II.2.4.8-fr]{2.4.8}). En ce qui concerne (iii), lorsqu'on
suppose en outre $\mathcal{L}$ isomorphe à $\mathcal{O}_Y$ (ce qui est permis,
la question étant locale sur $Y$), l'isomorphie de
$\operatorname{Proj}(\mathcal{S})$ et
$\operatorname{Proj}(\mathcal{S}_{(\mathcal{L})})$ est évidente~; pour
définir un isomorphisme \emph{canonique}, soient
$Y=\operatorname{Spec}(A)$, $\mathcal{S}=\widetilde{S}$, où $S$ est une
$A$-algèbre graduée, et soit $c$ un générateur du $A$-module libre $L$ tel
que $\mathcal{L}=\widetilde{L}$~; alors pour tout $n>0$,
$x_n\mapsto x_n\otimes c^{\otimes n}$ est un $A$-isomorphisme de $S_n$ sur
$S_n\otimes L^{\otimes n}$, et ces $A$-isomorphismes définissent un
$A$-isomorphisme d'algèbres graduées
\[
  \varphi_c:S\longrightarrow S_{(L)}
  =\bigoplus_{n\geq0}S_n\otimes L^{\otimes n}.
\]
Soit alors $f\in S_+$, homogène de degré $d$~; pour tout $x\in S_{nd}$, on a
\[
  \frac{x\otimes c^{\otimes nd}}{(f\otimes c^{\otimes d})^n}
  =\frac{x\otimes(\varepsilon c)^{\otimes nd}}
  {(f\otimes(\varepsilon c)^{\otimes d})^n}
\]
pour tout élément inversible $\varepsilon\in A$, ce qui montre que
l'isomorphisme $S_{(f)}\to(S_{(L)})_{(f\otimes c^d)}$ déduit de $\varphi_c$
est \emph{indépendant} du générateur $c$ de $L$ considéré et achève la
démonstration.
\end{proof}

\begin{env}[3.1.9]
\label{II.3.1.9-fr}
Rappelons ((\hyperref[0.4.1.3-fr]{0, 4.1.3}) et
(\hyperref[I.1.3.14-fr]{I, 1.3.14})) que pour que la
$\mathcal{O}_Y$-Algèbre graduée quasi-cohérente $\mathcal{S}$ soit
\emph{engendrée par le $\mathcal{O}_Y$-Module $\mathcal{S}_1$}, il faut et il
suffit qu'il existe un recouvrement $(U_\alpha)$ de $Y$ par des ouverts
affines tels que l'algèbre graduée $\Gamma(U_\alpha,\mathcal{S})$ sur
$\Gamma(U_\alpha,\mathcal{S}_0)$ soit engendrée par l'ensemble
$\Gamma(U_\alpha,\mathcal{S}_1)$ de ses éléments homogènes de degré~$1$. Pour
tout ouvert $V$ de $Y$, $\mathcal{S}|V$ est alors engendrée par le
$(\mathcal{O}_Y|V)$-Module $\mathcal{S}_1|V$.
\end{env}

\begin{proposition}[3.1.10]
\label{II.3.1.10-fr}
Supposons qu'il existe un recouvrement ouvert affine fini $(U_i)$ de $Y$ tel
que chacune des algèbres graduées $\Gamma(U_i,\mathcal{S})$ soit de type fini
sur $\Gamma(U_i,\mathcal{O}_Y)$. Alors il existe $d>0$ tel que
$\mathcal{S}^{(d)}$ soit engendrée par $\mathcal{S}_d$, $\mathcal{S}_d$ étant
un $\mathcal{O}_Y$-Module de type fini.
\end{proposition}

\begin{proof}
En effet, il résulte de (\hyperref[II.2.1.6-fr]{2.1.6, (v)}) que pour chaque
$i$, il existe un entier $m_i$ tel que
$\Gamma(U_i,\mathcal{S}_{nm_i})=(\Gamma(U_i,\mathcal{S}_{m_i}))^n$ pour tout
$n>0$~; il suffit de prendre pour $d$ un multiple commun des $m_i$, compte
tenu de (\hyperref[II.2.1.6-fr]{2.1.6, (i)}).
\end{proof}

\begin{corollary}[3.1.11]
\label{II.3.1.11-fr}
Sous les hypothèses de (\hyperref[II.3.1.10-fr]{3.1.10}),
$\operatorname{Proj}(\mathcal{S})$ est $Y$-isomorphe à un spectre homogène
$\operatorname{Proj}(\mathcal{S}')$, où $\mathcal{S}'$ est une
$\mathcal{O}_Y$-Algèbre graduée engendrée par $\mathcal{S}'_1$,
$\mathcal{S}'_1$ étant supposé être un $\mathcal{O}_Y$-Module de type fini.
\end{corollary}

\begin{proof}
Il suffit en effet de prendre $\mathcal{S}'=\mathcal{S}^{(d)}$, $d$ étant
déterminé par la propriété de (\hyperref[II.3.1.10-fr]{3.1.10}), et
d'appliquer (\hyperref[II.3.1.8-fr]{3.1.8, (i)}).
\end{proof}

\begin{env}[3.1.12]
\label{II.3.1.12-fr}
Si $\mathcal{S}$ est une $\mathcal{O}_Y$-Algèbre graduée quasi-cohérente à
degrés positifs, on sait (\hyperref[I.5.1.1-fr]{I, 5.1.1}) que son
\emph{nilradical} $\mathcal{N}$ est un $\mathcal{O}_Y$-Module
quasi-cohérent~; nous dirons que $\mathcal{N}_+=\mathcal{N}\cap\mathcal{S}_+$
est le \emph{nilradical de $\mathcal{S}_+$}~; c'est un
$\mathcal{S}_0$-Module quasi-cohérent gradué, car on se ramène immédiatement
au cas où $Y$ est affine, et la proposition résulte alors de
(\hyperref[II.2.1.10-fr]{2.1.10}). Pour tout $y\in Y$,
$(\mathcal{N}_+)_y$ est alors le nilradical de
$(\mathcal{S}_+)_y=(\mathcal{S}_y)_+$
(\hyperref[I.5.1.1-fr]{I, 5.1.1}). On dit que la
$\mathcal{O}_Y$-Algèbre graduée $\mathcal{S}$ est \emph{essentiellement
réduite} si $\mathcal{N}_+=0$, ce qui équivaut

\oldpage[II]{53}

à dire que $\mathcal{S}_y$
est une $\mathcal{O}_y$-algèbre graduée essentiellement réduite pour tout
$y\in Y$. Pour toute $\mathcal{O}_Y$-Algèbre graduée $\mathcal{S}$,
$\mathcal{S}/\mathcal{N}_+$ est essentiellement réduite.

Nous dirons que $\mathcal{S}$ est \emph{intègre} si $\mathcal{S}_y$ est un
anneau intègre pour tout $y\in Y$ et si en outre
$(\mathcal{S}_y)_+=(\mathcal{S}_+)_y\neq0$ pour tout $y\in Y$.
\end{env}

\begin{proposition}[3.1.13]
\label{II.3.1.13-fr}
Soit $\mathcal{S}$ une $\mathcal{O}_Y$-Algèbre graduée à degrés positifs.
Si $X=\operatorname{Proj}(\mathcal{S})$, le $Y$-schéma $X_{\mathrm{red}}$ est
canoniquement isomorphe à $\operatorname{Proj}(\mathcal{S}/\mathcal{N}_+)$~;
en particulier, si $\mathcal{S}$ est essentiellement réduite, $X$ est réduit.
\end{proposition}

\begin{proof}
Le fait que $X'=\operatorname{Proj}(\mathcal{S}/\mathcal{N}_+)$ est réduit
découle immédiatement de (\hyperref[II.2.4.4-fr]{2.4.4, (i)}), la propriété
étant locale~; en outre, pour tout ouvert affine $U\subset Y$,
${\varphi'}^{-1}(U)$ est égal à $(\varphi^{-1}(U))_{\mathrm{red}}$ (en
désignant par $\varphi$ et $\varphi'$ les morphismes structuraux $X\to Y$,
$X'\to Y$)~; on vérifie aussitôt que les $U$-morphismes canoniques
${\varphi'}^{-1}(U)\to\varphi^{-1}(U)$ sont compatibles avec les opérations
de restriction et définissent par suite une immersion fermée $X'\to X$ qui
est un homéomorphisme des espaces sous-jacents~; d'où la conclusion
(\hyperref[I.5.1.2-fr]{I, 5.1.2}).
\end{proof}

\begin{proposition}[3.1.14]
\label{II.3.1.14-fr}
Soient $Y$ un préschéma intègre, $\mathcal{S}$ une
$\mathcal{O}_Y$-Algèbre graduée quasi-cohérente telle que
$\mathcal{S}_0=\mathcal{O}_Y$.
\begin{enumerate}
  \item[\rm(i)] Si $\mathcal{S}$ est intègre
    (\hyperref[II.3.1.12-fr]{3.1.12}), alors
    $X=\operatorname{Proj}(\mathcal{S})$ est intègre et le morphisme
    structural $\varphi:X\to Y$ est dominant.
  \item[\rm(ii)] Supposons en outre que $\mathcal{S}$ soit essentiellement
    réduite. Alors, inversement, si $X$ est intègre et si $\varphi$ est
    dominant, $\mathcal{S}$ est intègre.
\end{enumerate}
\end{proposition}

\begin{proof}
\begin{enumerate}
  \item[\rm(i)] Si $(U_\alpha)$ est une base de $Y$ formée d'ouverts affines
    non vides, il suffit de démontrer la proposition lorsque $Y$ est remplacé
    par un des $U_\alpha$ et $\mathcal{S}$ par $\mathcal{S}|U_\alpha$~: en
    effet, il en résultera d'une part que les espaces sous-jacents
    $\varphi^{-1}(U_\alpha)$ sont des ouverts irréductibles (donc non vides)
    de $X$ tels que
    $\varphi^{-1}(U_\alpha)\cap\varphi^{-1}(U_\beta)\neq\varnothing$ pour tout
    couple d'indices (puisque $U_\alpha\cap U_\beta$ contient un $U_\gamma$),
    donc que $X$ est irréductible (\hyperref[0.2.1.4-fr]{0, 2.1.4})~; d'autre
    part $X$ sera réduit, puisque c'est là une propriété locale, donc $X$
    sera bien intègre et $\varphi(X)$ dense dans $Y$.

    Supposons donc que $Y=\operatorname{Spec}(A)$ où $A$ est intègre
    (\hyperref[I.5.1.4-fr]{I, 5.1.4}) et $\mathcal{S}=\widetilde{S}$, où $S$
    est une $A$-algèbre graduée~; l'hypothèse est que pour tout $y\in Y$,
    $\widetilde{S}_y=S_y$ est un anneau gradué intègre tel que
    $(S_y)_+\neq0$. Il suffit de prouver que $S$ est un anneau \emph{intègre},
    car alors on aura $S_+\neq0$ et on pourra appliquer
    (\hyperref[II.2.4.4-fr]{2.4.4, (ii)}). Or, soient $f$, $g$ deux éléments
    $\neq0$ de $S$ et supposons que $fg=0$~; pour tout $y\in Y$ on aura donc
    $(f/1)(g/1)=0$ dans $S_y$, donc $f/1=0$ ou $g/1=0$ par hypothèse.
    Supposons par exemple $f/1=0$ dans $S_y$~; cela signifie qu'il existe
    $a\in A$ tel que $a\notin\mathfrak{j}_y$ et $af=0$~; on a alors pour tout
    $z\in Y$, $(a/1)(f/1)=0$ dans l'anneau \emph{intègre} $S_z$, et comme
    $a/1\neq0$ (puisque $A$ est intègre), $f/1=0$, ce qui implique $f=0$.
  \item[\rm(ii)] La question étant locale sur $Y$, on peut encore supposer
    que $Y=\operatorname{Spec}(A)$, $A$ étant intègre, et
    $\mathcal{S}=\widetilde{S}$. Par hypothèse, pour tout $y\in Y$, $(S_y)_+$
    ne contient pas d'élément nilpotent $\neq0$, et il en est de même de
    $(S_0)_y=A_y$ par hypothèse~; donc $S_y$ est un anneau réduit pour tout
    $y\in Y$, et on en conclut d'abord que $S$ lui-même est réduit
    (\hyperref[I.5.1.1-fr]{I, 5.1.1}). L'hypothèse que $X$ est intègre
    entraîne d'autre part que $S$ est essentiellement intègre
    (\hyperref[II.2.4.4-fr]{2.4.4, (ii)}), et tout revient alors à voir que
    l'annulateur $\mathfrak{J}$ de $S_+$ dans $A=S_0$ est réduit

\oldpage[II]{54}

    à $0$
    (\hyperref[II.2.1.11-fr]{2.1.11}). Dans le cas contraire, on aurait
    $(S_h)_+=0$ pour un $h\neq0$ dans $\mathfrak{J}$, donc
    (\hyperref[II.3.1.1-fr]{3.1.1}) $\varphi^{-1}(D(h))=\varnothing$, et
    $\varphi(X)$ ne serait pas dense dans $Y$ contrairement à l'hypothèse
    (puisque $D(h)\neq\varnothing$, $h$ n'étant pas nilpotent).
\end{enumerate}
\end{proof}

\subsection{Faisceau sur $\operatorname{Proj}(\mathcal{S})$ associé à un
$\mathcal{S}$-Module gradué.}
\label{subsection:II.3.2-fr}

\begin{env}[3.2.1]
\label{II.3.2.1-fr}
Soient $Y$ un préschéma, $\mathcal{S}$ une $\mathcal{O}_Y$-Algèbre graduée
quasi-cohérente à degrés positifs, $\mathcal{M}$ un $\mathcal{S}$-module
gradué quasi-cohérent (sur $(Y,\mathcal{O}_Y)$ ou sur l'espace annelé
$(Y,\mathcal{S})$, ce qui revient au même
(\hyperref[I.9.6.1-fr]{I, 9.6.1})). Avec les notations de
(\hyperref[II.3.1.1-fr]{3.1.1}), désignons par $\widetilde{\mathcal{M}}_U$ le
$\mathcal{O}_{X_U}$-Module quasi-cohérent
$\widetilde{\Gamma(U,\mathcal{M})}$~; pour $U'\subset U$,
$\Gamma(U',\mathcal{M})$ s'identifie canoniquement à
$\Gamma(U,\mathcal{M})\otimes_A A'$ (\hyperref[I.1.6.4-fr]{I, 1.6.4})~; donc
on a $\widetilde{\mathcal{M}}_{U'}
=\rho_{U',U}^*(\widetilde{\mathcal{M}}_U)$
(\hyperref[II.2.8.11-fr]{2.8.11}).
\end{env}

\begin{proposition}[3.2.2]
\label{II.3.2.2-fr}
Il existe sur $\operatorname{Proj}(\mathcal{S})=X$ un
$\mathcal{O}_X$-Module quasi-cohérent et un seul $\widetilde{\mathcal{M}}$ tel
que, pour tout ouvert affine $U$ de $Y$, on ait
\[
  \eta_U^*(\widetilde{\Gamma(U,\mathcal{M})})
  =\widetilde{\mathcal{M}}|f^{-1}(U)
\]
(en désignant par $\eta_U$ l'isomorphisme
$f^{-1}(U)\xrightarrow{\sim}\operatorname{Proj}(\Gamma(U,\mathcal{S}))$, où
$f$ est le morphisme structural $X\to Y$).
\end{proposition}

\begin{proof}
Comme $\rho_{U',U}$ s'identifie au morphisme d'injection
$f^{-1}(U')\to f^{-1}(U)$ (\hyperref[II.3.1.2.1-fr]{3.1.2.1}), la proposition
résulte aussitôt de la relation
$\widetilde{\mathcal{M}}_{U'}=\rho_{U',U}^*(\widetilde{\mathcal{M}}_U)$ et du
principe de recollement des faisceaux (\hyperref[0.3.3.1-fr]{0, 3.3.1}).
\end{proof}

On dit que $\widetilde{\mathcal{M}}$ est le $\mathcal{O}_X$-Module associé au
$\mathcal{S}$-Module gradué quasi-cohérent $\mathcal{M}$.

\begin{proposition}[3.2.3]
\label{II.3.2.3-fr}
Soit $\mathcal{M}$ un $\mathcal{S}$-Module gradué quasi-cohérent, et soit
$f\in\Gamma(Y,\mathcal{S}_d)$ ($d>0$). Si $\xi_f$ est l'isomorphisme canonique
de $X_f$ sur le $Y$-préschéma
$Z_f=\operatorname{Spec}(\mathcal{S}^{(d)}/(f-1)\mathcal{S}^{(d)})$
(\hyperref[II.3.1.4-fr]{3.1.4}), $(\xi_f)_*(\widetilde{\mathcal{M}}|X_f)$ est
le $\mathcal{O}_{Z_f}$-Module
$\widetilde{\mathcal{M}^{(d)}/(f-1)\mathcal{M}^{(d)}}$
(\hyperref[II.1.4.3-fr]{1.4.3}).
\end{proposition}

\begin{proof}
La question étant locale sur $Y$, on est aussitôt ramené à
(\hyperref[II.2.2.5-fr]{2.2.5}), compte tenu de la commutativité du diagramme
(\hyperref[II.2.8.12.1-fr]{2.8.12.1}).
\end{proof}

\begin{proposition}[3.2.4]
\label{II.3.2.4-fr}
Le $\mathcal{O}_X$-Module $\widetilde{\mathcal{M}}$ est un foncteur covariant
additif exact en $\mathcal{M}$, de la catégorie des $\mathcal{S}$-Modules
gradués quasi-cohérents dans celle des $\mathcal{O}_X$-Modules
quasi-cohérents, qui commute aux limites inductives et aux sommes directes.
\end{proposition}

\begin{proof}
La question étant locale sur $Y$, se ramène à
(\hyperref[I.1.3.11-fr]{I, 1.3.11} et
\hyperref[I.1.3.9-fr]{I, 1.3.9}) et à
(\hyperref[II.2.5.4-fr]{2.5.4}).
\end{proof}

En particulier, si $\mathcal{N}$ est un sous-$\mathcal{S}$-Module gradué
quasi-cohérent de $\mathcal{M}$, $\widetilde{\mathcal{N}}$ s'identifie
canoniquement à un sous-$\mathcal{O}_X$-Module quasi-cohérent de
$\widetilde{\mathcal{M}}$~; plus particulièrement, pour tout faisceau gradué
quasi-cohérent $\mathcal{J}$ d'idéaux de $\mathcal{S}$,
$\widetilde{\mathcal{J}}$ est un faisceau quasi-cohérent d'idéaux de
$\mathcal{O}_X$.

Si $\mathcal{M}$ est un $\mathcal{S}$-Module gradué quasi-cohérent et
$\mathcal{I}$ un faisceau quasi-cohérent d'idéaux de $\mathcal{O}_Y$,
$\mathcal{I}\mathcal{M}$ est un sous-$\mathcal{S}$-Module gradué
quasi-cohérent de $\mathcal{M}$ et on a
\[
\label{II.3.2.4.1-fr}
  \widetilde{\mathcal{I}\mathcal{M}}
  =\mathcal{I}\cdot\widetilde{\mathcal{M}}
\tag{3.2.4.1}
\]
(le second membre ayant le sens défini en
(\hyperref[0.4.3.5-fr]{0, 4.3.5})). Il suffit en effet de vérifier cette
formule lorsque $Y=\operatorname{Spec}(A)$ est affine,
$\mathcal{S}=\widetilde{S}$, où $S$ est une $A$-algèbre graduée,
$\mathcal{M}=\widetilde{M}$,

\oldpage[II]{55}

où $M$ est un $S$-module gradué, et
$\mathcal{I}=\widetilde{\mathfrak{I}}$, où $\mathfrak{I}$ est un idéal de
$A$. Pour tout $f$ homogène de $S_+$, la restriction à
$D_+(f)=\operatorname{Spec}(S_{(f)})$ du premier membre de
(\hyperref[II.3.2.4.1-fr]{3.2.4.1}) est associée à
$(\mathfrak{I}M)_{(f)}=\mathfrak{I}\cdot M_{(f)}$, et il en est de même de la
restriction du second membre, vu
(\hyperref[I.1.3.13-fr]{I, 1.3.13} et
\hyperref[I.1.6.9-fr]{I, 1.6.9}).

\begin{proposition}[3.2.5]
\label{II.3.2.5-fr}
Soit $f\in\Gamma(Y,\mathcal{S}_d)$ ($d>0$). Sur l'ensemble ouvert $X_f$, le
$(\mathcal{O}_X|X_f)$-Module
$\widetilde{\mathcal{S}(nd)}|X_f$ est canoniquement isomorphe à
$\mathcal{O}_X|X_f$ pour tout $n\in\mathbf{Z}$. En particulier, si la
$\mathcal{O}_Y$-Algèbre $\mathcal{S}$ est engendrée par $\mathcal{S}_1$
(\hyperref[II.3.1.9-fr]{3.1.9}), les $\mathcal{O}_X$-Modules
$\widetilde{\mathcal{S}(n)}$ sont inversibles pour tout $n\in\mathbf{Z}$.
\end{proposition}

\begin{proof}
En effet, pour tout ouvert affine $U$ de $Y$, on a défini dans
(\hyperref[II.2.5.7-fr]{2.5.7}) un isomorphisme canonique de
$\widetilde{\mathcal{S}(nd)}|(X_f\cap\varphi^{-1}(U))$ sur
$\mathcal{O}_X|(X_f\cap\varphi^{-1}(U))$, compte tenu de
(\hyperref[II.3.1.4-fr]{3.1.4}) (où $\varphi$ est le morphisme structural
$X\to Y$)~; il est immédiat de vérifier que ces isomorphismes sont compatibles
avec la restriction de $U$ à un ouvert affine $U'\subset U$, d'où la première
assertion. Pour établir la seconde, il suffit de remarquer que si
$\mathcal{S}$ est engendrée par $\mathcal{S}_1$, il y a un recouvrement
$(U_\alpha)$ de $Y$ par des ouverts affines tels que
$\Gamma(U_\alpha,\mathcal{S})$ soit engendrée par
$\Gamma(U_\alpha,\mathcal{S})_1=\Gamma(U_\alpha,\mathcal{S}_1)$~; on est alors
ramené à appliquer le résultat de (\hyperref[II.2.5.9-fr]{2.5.9}), la
propriété d'être inversible étant locale.
\end{proof}

On posera encore, pour tout $n\in\mathbf{Z}$,
\[
\label{II.3.2.5.1-fr}
  \mathcal{O}_X(n)=\widetilde{\mathcal{S}(n)}
\tag{3.2.5.1}
\]
et pour tout $\mathcal{O}_X$-Module $\mathcal{F}$
\[
\label{II.3.2.5.2-fr}
  \mathcal{F}(n)=\mathcal{F}\otimes_{\mathcal{O}_X}\mathcal{O}_X(n).
\tag{3.2.5.2}
\]
Il résulte aussitôt de ces définitions que pour tout ouvert $U$ de $Y$, on a
\[
  \widetilde{(\mathcal{S}|U)(n)}
  =\mathcal{O}_X(n)|f^{-1}(U)
\]
$f$ étant le morphisme structural $X\to Y$.

\begin{proposition}[3.2.6]
\label{II.3.2.6-fr}
Soient $\mathcal{M}$ et $\mathcal{N}$ deux $\mathcal{S}$-Modules gradués
quasi-cohérents, il existe un homomorphisme canonique fonctoriel (en
$\mathcal{M}$ et $\mathcal{N}$)
\[
\label{II.3.2.6.1-fr}
  \lambda:\widetilde{\mathcal{M}}
  \otimes_{\mathcal{O}_X}\widetilde{\mathcal{N}}
  \longrightarrow\widetilde{\mathcal{M}\otimes_{\mathcal{S}}\mathcal{N}}
\tag{3.2.6.1}
\]
et un homomorphisme canonique fonctoriel (en $\mathcal{M}$ et
$\mathcal{N}$)
\[
\label{II.3.2.6.2-fr}
  \mu:\widetilde{\mathcal{H}om_{\mathcal{S}}(\mathcal{M},\mathcal{N})}
  \longrightarrow
  \mathcal{H}om_{\mathcal{O}_X}
  (\widetilde{\mathcal{M}},\widetilde{\mathcal{N}}).
\tag{3.2.6.2}
\]
En outre, si $\mathcal{S}$ est engendrée par $\mathcal{S}_1$
(\hyperref[II.3.1.9-fr]{3.1.9}), $\lambda$ est un isomorphisme~; si de plus
$\mathcal{M}$ admet une présentation finie
(\hyperref[II.3.1.1-fr]{3.1.1}), $\mu$ est un isomorphisme.
\end{proposition}

\begin{proof}
Les isomorphismes $\lambda$ et $\mu$ ont été définis dans
(\hyperref[II.2.5.11.2-fr]{2.5.11.2}) et
(\hyperref[II.2.5.12.2-fr]{2.5.12.2}) lorsque $Y$ est affine~; ces définitions
étant locales se transportent aussitôt au cas général considéré ici, compte
tenu de (\hyperref[II.2.8.14-fr]{2.8.14}).
\end{proof}

\begin{corollary}[3.2.7]
\label{II.3.2.7-fr}
Si $\mathcal{S}$ est engendrée par $\mathcal{S}_1$, on a, quels que soient
$m$, $n$ dans $\mathbf{Z}$,
\[
\label{II.3.2.7.1-fr}
  \mathcal{O}_X(m)\otimes_{\mathcal{O}_X}\mathcal{O}_X(n)
  =\mathcal{O}_X(m+n)
\tag{3.2.7.1}
\]
\[
\label{II.3.2.7.2-fr}
  \mathcal{O}_X(n)=(\mathcal{O}_X(1))^{\otimes n}
\tag{3.2.7.2}
\]
à des isomorphismes canoniques près.
\end{corollary}

\oldpage[II]{56}

\begin{corollary}[3.2.8]
\label{II.3.2.8-fr}
Supposons $\mathcal{S}$ engendrée par $\mathcal{S}_1$. Pour tout
$\mathcal{S}$-Module gradué $\mathcal{M}$ et tout $n\in\mathbf{Z}$, on a
\[
\label{II.3.2.8.1-fr}
  (\mathcal{M}(n))^{\sim}=\widetilde{\mathcal{M}}(n)
\tag{3.2.8.1}
\]
à un isomorphisme canonique près.
\end{corollary}

\begin{proof}
Cela résulte des propriétés correspondantes pour $Y$ affine
(\hyperref[II.2.5.14-fr]{2.5.14} et
\hyperref[II.2.5.15-fr]{2.5.15}) ainsi que de
(\hyperref[II.2.8.11-fr]{2.8.11}).
\end{proof}

\par\medskip\noindent\textit{Remarques (3.2.9). ---\ }
\label{II.3.2.9-fr}
\begin{enumerate}
  \item[\rm(i)] Si $\mathcal{S}=\mathcal{A}[T]$, $\mathcal{A}$ étant une
    $\mathcal{O}_Y$-Algèbre quasi-cohérente
    (\hyperref[II.3.1.7-fr]{3.1.7}), on vérifie immédiatement que tous les
    $\mathcal{O}_X$-Modules inversibles $\mathcal{O}_X(n)$ sont canoniquement
    isomorphes à $\mathcal{O}_X$.

    En outre, soit $\mathcal{N}$ un $\mathcal{A}$-Module quasi-cohérent, et
    posons $\mathcal{M}=\mathcal{N}\otimes_{\mathcal{A}}\mathcal{A}[T]$. Il
    résulte alors de (\hyperref[II.3.2.3-fr]{3.2.3}) et de
    (\hyperref[II.3.1.7-fr]{3.1.7}) que dans l'identification canonique de
    $X=\operatorname{Proj}(\mathcal{A}[T])$ et de
    $X'=\operatorname{Spec}(\mathcal{A})$, le $\mathcal{O}_X$-Module
    $\widetilde{\mathcal{M}}$ s'identifie au $\mathcal{O}_{X'}$-Module
    $\widetilde{\mathcal{N}}$ associé à $\mathcal{N}$ (au sens de
    (\hyperref[II.1.4.3-fr]{1.4.3})).
  \item[\rm(ii)] Soit $\mathcal{S}$ une $\mathcal{O}_Y$-Algèbre graduée
    quelconque, $\mathcal{S}'$ la $\mathcal{O}_Y$-Algèbre graduée telle que
    $\mathcal{S}'_0=\mathcal{O}_Y$, $\mathcal{S}'_n=\mathcal{S}_n$ pour tout
    $n>0$~; l'isomorphisme canonique de
    $X=\operatorname{Proj}(\mathcal{S})$ sur
    $X'=\operatorname{Proj}(\mathcal{S}')$
    (\hyperref[II.3.1.8-fr]{3.1.8, (ii)}) identifie $\mathcal{O}_X(n)$ et
    $\mathcal{O}_{X'}(n)$ pour tout $n\in\mathbf{Z}$~: cela résulte de la même
    proposition pour le cas affine (\hyperref[II.2.5.16-fr]{2.5.16}) et du
    fait que ces identifications, pour les ouverts affines de $Y$, commutent
    avec les opérations de restriction. De même, soit
    $X^{(d)}=\operatorname{Proj}(\mathcal{S}^{(d)})$~; l'isomorphisme
    canonique de $X$ sur $X^{(d)}$
    (\hyperref[II.3.1.8-fr]{3.1.8, (i)}) identifie $\mathcal{O}_X(nd)$ et
    $\mathcal{O}_{X^{(d)}}(n)$ pour tout $n\in\mathbf{Z}$.
\end{enumerate}

\begin{proposition}[3.2.10]
\label{II.3.2.10-fr}
Soient $\mathcal{L}$ un $\mathcal{O}_Y$-Module inversible, $g$ l'isomorphisme
canonique
$X_{(\mathcal{L})}=\operatorname{Proj}(\mathcal{S}_{(\mathcal{L})})
\to X=\operatorname{Proj}(\mathcal{S})$
(\hyperref[II.3.1.8-fr]{3.1.8, (iii)}). Pour tout entier $n\in\mathbf{Z}$,
$g_*(\mathcal{O}_{X_{(\mathcal{L})}}(n))$ est canoniquement isomorphe à
$\mathcal{O}_X(n)\otimes_Y\mathcal{L}^{\otimes n}$.
\end{proposition}

\begin{proof}
Supposons d'abord $Y$ affine, d'anneau $A$, et
$\mathcal{L}=\widetilde{L}$, où $L$ est un $A$-module libre monogène. Avec les
notations de la démonstration de
(\hyperref[II.3.1.8-fr]{3.1.8, (iii)}) on définit, pour $f\in S_d$, un
isomorphisme de $(S(n))_{(f)}\otimes_A L^{\otimes n}$ sur
$(S_{(L)}(n))_{(f\otimes c^{\otimes d})}$ en faisant correspondre à
$(x/f^k)\otimes c^{\otimes n}$, où $x\in S_{kd+n}$, l'élément
$(x\otimes c^{\otimes(n+kd)})/(f\otimes c^{\otimes d})^k$~; il est immédiat
que cet isomorphisme est indépendant du générateur $c$ choisi pour $L$~; en
outre, les isomorphismes ainsi définis pour chaque $f\in S_+$ sont compatibles
avec les opérations de restriction $D_+(f)\to D_+(fg)$. Enfin, dans le cas
général, on voit sans peine en revenant aux définitions
(\hyperref[II.3.1.1-fr]{3.1.1}) que les isomorphismes ainsi définis pour
chaque ouvert affine $U$ de $Y$ sont compatibles avec le passage de $U$ à un
ouvert affine $U'\subset U$.
\end{proof}

\subsection{$\mathcal{S}$-Module gradué associé à un faisceau sur
$\operatorname{Proj}(\mathcal{S})$.}
\label{subsection:II.3.3-fr}

\emph{On suppose dans tout ce numéro que la $\mathcal{O}_Y$-Algèbre graduée
$\mathcal{S}$ est engendrée par $\mathcal{S}_1$
(\hyperref[II.3.1.9-fr]{3.1.9}).} Rappelons qu'en raison de
(\hyperref[II.3.1.8-fr]{3.1.8, (i)}), cette restriction n'est pas essentielle,
moyennant les conditions de finitude (\hyperref[II.3.1.10-fr]{3.1.10}).

\begin{env}[3.3.1]
\label{II.3.3.1-fr}
Soit $p$ le morphisme structural $X=\operatorname{Proj}(\mathcal{S})\to Y$.
Pour tout $\mathcal{O}_X$-Module $\mathcal{F}$, nous poserons
\[
\label{II.3.3.1.1-fr}
  \Gamma_*(\mathcal{F})
  =\bigoplus_{n\in\mathbf{Z}}p_*(\mathcal{F}(n))
\tag{3.3.1.1}
\]

\oldpage[II]{57}

et en particulier
\[
\label{II.3.3.1.2-fr}
  \Gamma_*(\mathcal{O}_X)
  =\bigoplus_{n\in\mathbf{Z}}p_*(\mathcal{O}_X(n)).
\tag{3.3.1.2}
\]
On sait (\hyperref[0.4.2.2-fr]{0, 4.2.2}) qu'il existe un homomorphisme
canonique
\[
  p_*(\mathcal{F})\otimes_{\mathcal{O}_Y}p_*(\mathcal{G})
  \longrightarrow
  p_*(\mathcal{F}\otimes_{\mathcal{O}_X}\mathcal{G})
\]
pour deux $\mathcal{O}_X$-Modules $\mathcal{F}$ et $\mathcal{G}$~; on déduit
donc de (\hyperref[II.3.2.7.1-fr]{3.2.7.1}) que
$\Gamma_*(\mathcal{O}_X)$ est muni d'une structure de
$\mathcal{O}_Y$-\emph{Algèbre graduée}, et
(\hyperref[II.3.2.5.2-fr]{3.2.5.2}) définit de même sur
$\Gamma_*(\mathcal{F})$ une structure de \emph{Module gradué} sur
$\Gamma_*(\mathcal{O}_X)$.

En vertu de (\hyperref[II.3.2.5-fr]{3.2.5}), et de l'exactitude à gauche du
foncteur $p_*$ (\hyperref[0.4.2.1-fr]{0, 4.2.1}),
$\Gamma_*(\mathcal{F})$ est un foncteur covariant additif et
\emph{exact à gauche} en $\mathcal{F}$ dans la catégorie des
$\mathcal{O}_X$-Modules, à valeurs dans la catégorie des
$\mathcal{O}_Y$-Modules gradués (où les morphismes sont les homomorphismes de
degré $0$). En particulier, si $\mathcal{J}$ est un faisceau d'idéaux dans
$\mathcal{O}_X$, $\Gamma_*(\mathcal{J})$ s'identifie à un
\emph{faisceau gradué d'idéaux} de $\Gamma_*(\mathcal{O}_X)$.
\end{env}

\begin{env}[3.3.2]
\label{II.3.3.2-fr}
Soit $\mathcal{M}$ un $\mathcal{S}$-Module gradué quasi-cohérent. Pour tout
ouvert affine $U$ de $Y$, on a défini en
(\hyperref[II.2.6.2-fr]{2.6.2}) un homomorphisme de groupes abéliens
\[
  \alpha_{0,U}:\Gamma(U,\mathcal{M}_0)
  \longrightarrow\Gamma(p^{-1}(U),\widetilde{\mathcal{M}}).
\]
Il est immédiat que ces homomorphismes commutent aux opérations de restriction
(\hyperref[II.2.8.13.1-fr]{2.8.13.1}) et définissent donc (sans utiliser
l'hypothèse que $\mathcal{S}$ est engendrée par $\mathcal{S}_1$) un
homomorphisme de faisceaux de groupes abéliens
\[
\label{II.3.3.2.1-fr}
  \alpha_0:\mathcal{M}_0\longrightarrow p_*(\widetilde{\mathcal{M}}).
\tag{3.3.2.1}
\]
Appliquant ce résultat à chacun des
$\mathcal{M}_n=(\mathcal{M}(n))_0$, et tenant compte de
(\hyperref[II.3.2.8.1-fr]{3.2.8.1}), on définit un homomorphisme de faisceaux
de groupes abéliens
\[
\label{II.3.3.2.2-fr}
  \alpha_n:\mathcal{M}_n\longrightarrow
  p_*(\widetilde{\mathcal{M}}(n))
  \qquad\text{pour tout }n\in\mathbf{Z},
\tag{3.3.2.2}
\]
d'où un homomorphisme fonctoriel (de degré $0$) de faisceaux gradués de groupes
abéliens
\[
\label{II.3.3.2.3-fr}
  \alpha:\mathcal{M}\longrightarrow\Gamma_*(\widetilde{\mathcal{M}})
\tag{3.3.2.3}
\]
(aussi noté $\alpha_{\mathcal{M}}$).

En prenant en particulier $\mathcal{M}=\mathcal{S}$, on vérifie que
$\alpha:\mathcal{S}\to\Gamma_*(\mathcal{O}_X)$ est un homomorphisme de
$\mathcal{O}_Y$-Algèbres graduées et que
(\hyperref[II.3.3.2.3-fr]{3.3.2.3}) est un di-homomorphisme de Modules
gradués, relatif à cet homomorphisme d'Algèbres graduées.

Remarquons encore qu'à chacun des $\alpha_n$ il correspond
(\hyperref[0.4.4.3-fr]{0, 4.4.3}) un homomorphisme canonique de
$\mathcal{O}_X$-Modules
\[
\label{II.3.3.2.4-fr}
  \alpha_n^\sharp:p^*(\mathcal{M}_n)
  \longrightarrow\widetilde{\mathcal{M}}(n).
\tag{3.3.2.4}
\]
On vérifie sans peine que cet homomorphisme n'est autre que celui qui
correspond fonctoriellement (\hyperref[II.3.2.4-fr]{3.2.4}) à
l'homomorphisme canonique (de degré $0$) de $\mathcal{O}_Y$-Modules gradués
\[
\label{II.3.3.2.5-fr}
  \mathcal{M}_n\otimes_{\mathcal{O}_Y}\mathcal{S}
  \longrightarrow\mathcal{M}(n)
\tag{3.3.2.5}
\]

\oldpage[II]{58}

où la graduation du second membre provient naturellement de celle de
$\mathcal{S}$. On peut en effet se borner au cas où
$Y=\operatorname{Spec}(A)$ est affine, $\mathcal{M}=\widetilde{M}$ et
$\mathcal{S}=\widetilde{S}$, la $A$-algèbre graduée $S$ étant engendrée par
$S_1$, de sorte que lorsque $f$ parcourt $S_1$, les $D_+(f)$ forment un
recouvrement de $X$. Revenant aux définitions
(\hyperref[II.2.6.2-fr]{2.6.2}) on voit alors, compte tenu de
(\hyperref[I.1.6.7-fr]{I, 1.6.7}), que la restriction à $D_+(f)$ de
l'homomorphisme (\hyperref[II.3.3.2.4-fr]{3.3.2.4}) correspond
(\hyperref[I.1.3.8-fr]{I, 1.3.8}) à l'homomorphisme de $S_{(f)}$-modules
$M_n\otimes_A S_{(f)}\to(S(n))_{(f)}$ qui, à $x\otimes1$ ($x\in M_n$), fait
correspondre $x/1$~; cela démontre notre assertion.
\end{env}

\begin{proposition}[3.3.3]
\label{II.3.3.3-fr}
Pour toute section $f\in\Gamma(Y,\mathcal{S}_d)$ ($d>0$), $X_f$ est identique
à l'ensemble des points de $X$ où $\alpha_d(f)$ (considérée comme section de
$\mathcal{O}_X(d)$) ne s'annule pas
(\hyperref[0.5.5.2-fr]{0, 5.5.2}).
\end{proposition}

\begin{proof}
($\alpha_d(f)$ est une section de $p_*(\mathcal{O}_X(d))$ au-dessus de $Y$,
mais par définition une telle section est aussi une section de
$\mathcal{O}_X(d)$ au-dessus de $X$
(\hyperref[0.4.2.1-fr]{0, 4.2.1})). La définition de $X_f$
(\hyperref[II.3.1.4-fr]{3.1.4}) ramène au cas où $Y$ est affine, qui a été
traité dans (\hyperref[II.2.6.3-fr]{2.6.3}).
\end{proof}

\begin{env}[3.3.4]
\label{II.3.3.4-fr}
Nous supposerons désormais, outre l'hypothèse du début de ce numéro, que pour
tout $\mathcal{O}_X$-Module quasi-cohérent $\mathcal{F}$, les
$p_*(\mathcal{F}(n))$ sont \emph{quasi-cohérents} sur $Y$, et par suite
$\Gamma_*(\mathcal{F})=\bigoplus_{n\in\mathbf{Z}}p_*(\mathcal{F}(n))$ est
aussi un $\mathcal{O}_Y$-Module quasi-cohérent
(\hyperref[I.1.4.1-fr]{I, 1.4.1} et
\hyperref[I.1.3.9-fr]{I, 1.3.9})~; cette circonstance se produira toujours en
particulier si $X$ est \emph{de type fini} sur $Y$
(\hyperref[I.9.2.2-fr]{I, 9.2.2}). On en conclut que
$(\Gamma_*(\mathcal{F}))^{\sim}$ est défini et est un
$\mathcal{O}_X$-Module quasi-cohérent. Pour tout ouvert affine $U$ de $Y$, on
a (\hyperref[I.1.3.9-fr]{I, 1.3.9} et
\hyperref[II.2.5.4-fr]{2.5.4})
\begin{align*}
  \left(\Gamma\left(U,\bigoplus_{n\in\mathbf{Z}}
  p_*(\mathcal{F}(n))\right)\right)^{\sim}
  &=\bigoplus_{n\in\mathbf{Z}}
  \bigl(\Gamma(U,p_*(\mathcal{F}(n)))\bigr)^{\sim}\\
  &=\bigoplus_{n\in\mathbf{Z}}
  \bigl(\Gamma(p^{-1}(U),\mathcal{F}(n))\bigr)^{\sim}\\
  &=\left(\bigoplus_{n\in\mathbf{Z}}
  \Gamma(p^{-1}(U),\mathcal{F}(n))\right)^{\sim}\\
  &=\bigl(\Gamma_*(\mathcal{F}|p^{-1}(U))\bigr)^{\sim}
\end{align*}
et par suite (\hyperref[II.2.6.4-fr]{2.6.4}) on a un homomorphisme canonique
\[
  \beta_U:\left(\Gamma\left(U,\bigoplus_{n\in\mathbf{Z}}
  p_*(\mathcal{F}(n))\right)\right)^{\sim}
  \longrightarrow\mathcal{F}|p^{-1}(U).
\]

En outre, la commutativité de
(\hyperref[II.2.8.13.2-fr]{2.8.13.2}) montre que ces homomorphismes commutent
avec les opérations de restriction sur $Y$~; on en déduit donc un
homomorphisme canonique fonctoriel
\[
\label{II.3.3.4.1-fr}
  \beta:(\Gamma_*(\mathcal{F}))^{\sim}\longrightarrow\mathcal{F}
\tag{3.3.4.1}
\]
(aussi noté $\beta_{\mathcal{F}}$) pour les $\mathcal{O}_X$-Modules
quasi-cohérents.
\end{env}

\begin{proposition}[3.3.5]
\label{II.3.3.5-fr}
Soient $\mathcal{M}$ un $\mathcal{S}$-Module gradué quasi-cohérent,
$\mathcal{F}$ un $\mathcal{O}_X$-Module quasi-cohérent~; les homomorphismes
composés
\[
\label{II.3.3.5.1-fr}
  \widetilde{\mathcal{M}}
  \xrightarrow{\widetilde{\alpha}}
  (\Gamma_*(\widetilde{\mathcal{M}}))^{\sim}
  \xrightarrow{\beta}\widetilde{\mathcal{M}}
\tag{3.3.5.1}
\]
\[
\label{II.3.3.5.2-fr}
  \Gamma_*(\mathcal{F})\xrightarrow{\alpha}
  \Gamma_*((\Gamma_*(\mathcal{F}))^{\sim})
  \xrightarrow{\Gamma_*(\beta)}\Gamma_*(\mathcal{F})
\tag{3.3.5.2}
\]
sont les isomorphismes identiques.
\end{proposition}

\begin{proof}
La question est en effet locale sur $Y$, et on est ramené à
(\hyperref[II.2.6.5-fr]{2.6.5}).
\end{proof}

\subsection{Conditions de finitude.}
\label{subsection:II.3.4-fr}

\oldpage[II]{59}

\begin{proposition}[3.4.1]
\label{II.3.4.1-fr}
Soient $Y$ un préschéma, $\mathcal{S}$ une $\mathcal{O}_Y$-Algèbre
quasi-cohérente engendrée par $\mathcal{S}_1$
(\hyperref[II.3.1.9-fr]{3.1.9})~; on suppose en outre $\mathcal{S}_1$ de type
fini. Alors $X=\operatorname{Proj}(\mathcal{S})$ est de type fini au-dessus de
$Y$.
\end{proposition}

\begin{proof}
En effet, la question étant locale sur $Y$, on peut supposer $Y$ affine
d'anneau $A$~; on a alors $\mathcal{S}=\widetilde{S}$, où
$S=\Gamma(Y,\mathcal{S})$, et par hypothèse $S$ est une $A$-algèbre engendrée
par $S_1=\Gamma(Y,\mathcal{S}_1)$, où on peut en outre supposer que $S_1$ est
un $A$-module de type fini
(\hyperref[I.1.3.9-fr]{I, 1.3.9} et
\hyperref[I.1.3.12-fr]{I, 1.3.12}). Par suite $S$ est une $A$-algèbre graduée
de type fini, et on est ramené à
(\hyperref[II.2.7.1-fr]{2.7.1, (ii)}).
\end{proof}

\begin{env}[3.4.2]
\label{II.3.4.2-fr}
Soit $\mathcal{S}$ une $\mathcal{O}_Y$-Algèbre graduée quasi-cohérente~; pour
un $\mathcal{S}$-Module gradué quasi-cohérent $\mathcal{M}$, nous considérerons
les conditions de finitude suivantes~:
\begin{enumerate}
  \item[(\textbf{TF})] \emph{Il existe un entier $n$ tel que le
    $\mathcal{S}$-Module $\bigoplus_{k\geq n}\mathcal{M}_k$ soit de type fini.}
  \item[(\textbf{TN})] \emph{Il existe un entier $n$ tel que
    $\mathcal{M}_k=0$ pour $k\geq n$.}
\end{enumerate}
Si $\mathcal{M}$ vérifie (\textbf{TN}), on a
$\widetilde{\mathcal{M}}=0$, la question étant locale sur $Y$
(\hyperref[II.2.7.2-fr]{2.7.2}).

Soient $\mathcal{M}$, $\mathcal{N}$ deux $\mathcal{S}$-Modules gradués
quasi-cohérents~; on dit qu'un homomorphisme
$u:\mathcal{M}\to\mathcal{N}$ de degré $0$ est
(\textbf{TN})-\emph{injectif} (resp. (\textbf{TN})-\emph{surjectif},
(\textbf{TN})-\emph{bijectif}) s'il existe un entier $n$ tel que
$u_k:\mathcal{M}_k\to\mathcal{N}_k$ soit injectif (resp. surjectif, bijectif)
pour $k\geq n$~; alors
$\widetilde{u}:\widetilde{\mathcal{M}}\to\widetilde{\mathcal{N}}$ est injectif
(resp. surjectif, bijectif) en vertu de
(\hyperref[II.2.7.2-fr]{2.7.2}), la question étant locale sur $Y$, et compte
tenu de (\hyperref[I.1.3.9-fr]{I, 1.3.9})~; lorsque $u$ est
(\textbf{TN})-bijectif, on dit aussi que $u$ est un
(\textbf{TN})-\emph{isomorphisme}.
\end{env}

\begin{proposition}[3.4.3]
\label{II.3.4.3-fr}
Soient $Y$ un préschéma, $\mathcal{S}$ une $\mathcal{O}_Y$-Algèbre graduée
quasi-cohérente engendrée par $\mathcal{S}_1$, $\mathcal{S}_1$ étant supposé
de type fini. Soit $\mathcal{M}$ un $\mathcal{S}$-Module gradué
quasi-cohérent.
\begin{enumerate}
  \item[\rm(i)] Si $\mathcal{M}$ vérifie la condition (\textbf{TF}),
    $\widetilde{\mathcal{M}}$ est de type fini.
  \item[\rm(ii)] Supposons que $\mathcal{M}$ vérifie (\textbf{TF})~; pour que
    $\widetilde{\mathcal{M}}=0$, il faut et il suffit que $\mathcal{M}$
    vérifie (\textbf{TN}).
\end{enumerate}
\end{proposition}

\begin{proof}
Les questions étant locales sur $Y$, on est ramené au cas où $Y$ est affine
d'anneau $A$, $\mathcal{S}=\widetilde{S}$, où $S$ est une $A$-algèbre graduée
telle que l'idéal $S_+$ soit de type fini,
$\mathcal{M}=\widetilde{M}$ où $M$ est un $S$-module gradué~; la proposition
résulte alors de (\hyperref[II.2.7.3-fr]{2.7.3}).
\end{proof}

\begin{theorem}[3.4.4]
\label{II.3.4.4-fr}
Soient $Y$ un préschéma, $\mathcal{S}$ une $\mathcal{O}_Y$-Algèbre graduée
quasi-cohérente engendrée par $\mathcal{S}_1$, $\mathcal{S}_1$ étant supposé
de type fini~; soit $X=\operatorname{Proj}(\mathcal{S})$. Pour tout
$\mathcal{O}_X$-Module quasi-cohérent $\mathcal{F}$, l'homomorphisme canonique
$\beta$ (\hyperref[II.3.3.4-fr]{3.3.4}) est un isomorphisme.
\end{theorem}

\begin{proof}
Notons d'abord que $\beta$ est défini en vertu de
(\hyperref[II.3.4.1-fr]{3.4.1}). Pour voir que $\beta$ est un isomorphisme,
on est ramené au cas où $Y$ est affine d'anneau $A$,
$\mathcal{S}=\widetilde{S}$, où $S$ est une $A$-algèbre graduée engendrée par
$S_1$ et $S_1$ est un $A$-module de type fini. Il suffit alors d'appliquer
(\hyperref[II.2.7.5-fr]{2.7.5}).
\end{proof}

\begin{corollary}[3.4.5]
\label{II.3.4.5-fr}
Sous les hypothèses de (\hyperref[II.3.4.4-fr]{3.4.4}), tout
$\mathcal{O}_X$-Module quasi-cohérent $\mathcal{F}$ est isomorphe à un
$\mathcal{O}_X$-Module de la forme $\widetilde{\mathcal{M}}$, où
$\mathcal{M}$ est un $\mathcal{S}$-Module gradué quasi-cohérent. Si en outre
$\mathcal{F}$ est de type fini, et si on suppose que $Y$ est un schéma
quasi-compact, ou que l'espace sous-jacent à $Y$ est noethérien, on peut
supposer $\mathcal{M}$ de type fini.
\end{corollary}

\oldpage[II]{60}

\begin{proof}
La première assertion résulte aussitôt de
(\hyperref[II.3.4.4-fr]{3.4.4}) en prenant
$\mathcal{M}=\Gamma_*(\mathcal{F})$. Pour établir la seconde, il suffira
d'établir que $\mathcal{M}$ est limite inductive de ses sous-$\mathcal{S}$-Modules
gradués de type fini $\mathcal{N}_\lambda$~: en effet, il en résultera que
$\widetilde{\mathcal{M}}$ est limite inductive des
$\widetilde{\mathcal{N}}_\lambda$
(\hyperref[II.3.2.4-fr]{3.2.4}), donc $\mathcal{F}$ est limite inductive des
$\beta(\widetilde{\mathcal{N}}_\lambda)$~; comme $X$ est quasi-compact
((\hyperref[II.3.4.1-fr]{3.4.1}) et
\hyperref[I.6.3.1-fr]{I, 6.3.1}) et que $\mathcal{F}$ est de type fini,
$\mathcal{F}$ sera nécessairement égal à un des
$\beta(\widetilde{\mathcal{N}}_\lambda)$
(\hyperref[0.5.2.3-fr]{0, 5.2.3}).

Pour définir les $\mathcal{N}_\lambda$ ayant $\mathcal{M}$ pour limite
inductive, il suffit de considérer pour chaque $n\in\mathbf{Z}$ le
$\mathcal{O}_Y$-Module quasi-cohérent $\mathcal{M}_n$, qui est limite
inductive de ses sous-$\mathcal{O}_Y$-Modules
$\mathcal{M}_n^{(\mu_n)}$ de type fini, en raison des hypothèses sur $Y$
(\hyperref[I.9.4.9-fr]{I, 9.4.9})~; il est immédiat que
$\mathcal{P}_{\mu_n}=\mathcal{S}\cdot\mathcal{M}_n^{(\mu_n)}$ est un
$\mathcal{S}$-Module gradué de type fini, et on vérifie aussitôt qu'en prenant
pour $\mathcal{N}_\lambda$ les sommes finies de $\mathcal{S}$-Modules de la
forme $\mathcal{P}_{\mu_n}$, on répond bien à la question.
\end{proof}

\begin{corollary}[3.4.6]
\label{II.3.4.6-fr}
Supposons vérifiées les hypothèses de (\hyperref[II.3.4.4-fr]{3.4.4}), et en
outre que l'espace sous-jacent $Y$ soit quasi-compact~; soit $\mathcal{F}$ un
$\mathcal{O}_X$-Module quasi-cohérent de type fini. Il existe $n_0$ tel que
pour $n\geq n_0$, l'homomorphisme canonique
$\sigma:p^*(p_*(\mathcal{F}(n)))\to\mathcal{F}(n)$
(\hyperref[0.4.4.3-fr]{0, 4.4.3}) soit surjectif.
\end{corollary}

\begin{proof}
En effet, pour tout $y\in Y$, soit $U$ un voisinage ouvert affine de $y$ dans
$Y$. Il existe un entier $n_0(U)$ tel que, pour $n\geq n_0(U)$,
$\mathcal{F}(n)|p^{-1}(U)$ soit engendré par un nombre fini de ses sections
au-dessus de $p^{-1}(U)$ (\hyperref[II.2.7.9-fr]{2.7.9})~; mais ces dernières
sont images canoniques de sections de $p^*(p_*(\mathcal{F}(n)))$ au-dessus de
$p^{-1}(U)$ (\hyperref[0.3.7.1-fr]{0, 3.7.1} et
\hyperref[0.4.4.3-fr]{0, 4.4.3}), donc
$\mathcal{F}(n)|p^{-1}(U)$ est égal à l'image canonique de
$p^*(p_*(\mathcal{F}(n)))|p^{-1}(U)$. Enfin, comme $Y$ est quasi-compact, il
existe un recouvrement fini de $Y$ par des ouverts affines $U_i$, et en
prenant pour $n_0$ le plus grand des $n_0(U_i)$, on achève la démonstration.
\end{proof}

\par\medskip\noindent\textit{Remarques (3.4.7). ---\ }
\label{II.3.4.7-fr}
Si $p=(\psi,\theta):X\to Y$ est un morphisme d'espaces annelés et
$\mathcal{F}$ un $\mathcal{O}_X$-Module, le fait que l'homomorphisme canonique
$\sigma:p^*(p_*(\mathcal{F}))\to\mathcal{F}$ soit surjectif s'explicite de la
façon suivante (\hyperref[0.4.4.1-fr]{0, 4.4.1})~: pour tout $x\in X$ et toute
section $s$ de $\mathcal{F}$ au-dessus d'un voisinage ouvert $V$ de $x$, il
existe un voisinage ouvert $U$ de $p(x)$ dans $Y$, un nombre fini de sections
$t_i$ ($1\leq i\leq m$) de $\mathcal{F}$ au-dessus de $p^{-1}(U)$, un
voisinage $W\subset V\cap p^{-1}(U)$ de $x$ et des sections $a_i$
($1\leq i\leq m$) de $\mathcal{O}_X$ au-dessus de $W$ telles que
\[
  s|W=\sum_i a_i\cdot(t_i|W).
\]
Lorsque $Y$ est un \emph{schéma affine} et $p_*(\mathcal{F})$
\emph{quasi-cohérent}, cette condition équivaut au fait que $\mathcal{F}$ est
\emph{engendré par ses sections au-dessus de $X$}
(\hyperref[0.5.5.1-fr]{0, 5.5.1})~: en effet, si
$Y=\operatorname{Spec}(A)$, on peut supposer que $U=D(f)$ avec $f\in A$~; il
existe un entier $n>0$ et des sections $s_i$ de $\mathcal{F}$ au-dessus de
$X$ telles que $t_i$ soit la restriction à $p^{-1}(U)$ de $s_ig^n$, avec
$g=\theta(f)$ (en appliquant (\hyperref[I.1.4.1-fr]{I, 1.4.1}) à
$p_*(\mathcal{F})$)~; comme $g$ est inversible au-dessus de $p^{-1}(U)$, on a
donc
\[
  s|W=\sum_i b_i\cdot(s_i|W)
\]
avec $b_i=a_i(g|W)^{-n}$, d'où notre assertion. Lorsque $Y$ est affine, le
corollaire (\hyperref[II.3.4.6-fr]{3.4.6}) redonne donc
(\hyperref[II.2.7.9-fr]{2.7.9}).

\oldpage[II]{61}

On en conclut que lorsque $Y$ est un préschéma quelconque, les trois
conditions suivantes, pour un $\mathcal{O}_X$-Module quasi-cohérent
$\mathcal{F}$ tel que $p_*(\mathcal{F})$ soit \emph{quasi-cohérent}, sont
équivalentes~:
\begin{enumerate}
  \item[a)] \emph{L'homomorphisme canonique
    $\sigma:p^*(p_*(\mathcal{F}))\to\mathcal{F}$ est surjectif.}
  \item[b)] \emph{Il existe un $\mathcal{O}_Y$-Module quasi-cohérent
    $\mathcal{G}$ et un homomorphisme surjectif
    $p^*(\mathcal{G})\to\mathcal{F}$.}
  \item[c)] \emph{Pour tout ouvert affine $U$ de $Y$,
    $\mathcal{F}|p^{-1}(U)$ est engendré par ses sections au-dessus de
    $p^{-1}(U)$.}
\end{enumerate}
On vient en effet de prouver l'équivalence de \emph{a)} et \emph{c)}. D'autre
part, il est clair que \emph{a)} entraîne \emph{b)}, $p_*(\mathcal{F})$ étant
par hypothèse quasi-cohérent. Inversement, tout homomorphisme
$u:p^*(\mathcal{G})\to\mathcal{F}$ se factorise en
$p^*(\mathcal{G})\to p^*(p_*(\mathcal{F}))\xrightarrow{\sigma}\mathcal{F}$
(\hyperref[0.3.5.4.4-fr]{0, 3.5.4.4}), donc si $u$ est surjectif il en est de
même de $\sigma$, ce qui prouve que \emph{b)} entraîne \emph{a)}.

\begin{corollary}[3.4.8]
\label{II.3.4.8-fr}
Supposons vérifiées les hypothèses de (\hyperref[II.3.4.4-fr]{3.4.4}) et
supposons en outre que $Y$ soit un schéma quasi-compact ou que l'espace
sous-jacent à $Y$ soit noethérien. Soit $\mathcal{F}$ un
$\mathcal{O}_X$-Module quasi-cohérent de type fini~; il existe alors un entier
$n_0$ tel que pour $n\geq n_0$, $\mathcal{F}$ soit isomorphe à un quotient
d'un $\mathcal{O}_X$-Module de la forme $(p^*(\mathcal{G}))(-n)$ où
$\mathcal{G}$ est un $\mathcal{O}_Y$-Module quasi-cohérent de type fini
(dépendant de $n$).
\end{corollary}

\begin{proof}
Comme le morphisme structural $X\to Y$ est séparé et de type fini,
$p_*(\mathcal{F}(n))$ est quasi-cohérent
(\hyperref[I.9.2.2-fr]{I, 9.2.2, b)}), donc limite inductive de ses
sous-$\mathcal{O}_Y$-Modules quasi-cohérents de type fini, en vertu de
l'hypothèse sur $Y$ (\hyperref[I.9.4.9-fr]{I, 9.4.9}). On en déduit par
(\hyperref[II.3.4.6-fr]{3.4.6}),
(\hyperref[0.4.3.2-fr]{0, 4.3.2}) et
(\hyperref[0.5.2.3-fr]{0, 5.2.3}) que $\mathcal{F}(n)$ est l'image canonique
d'un $\mathcal{O}_X$-Module de la forme $p^*(\mathcal{G})$, où $\mathcal{G}$
est un sous-$\mathcal{O}_Y$-Module quasi-cohérent de type fini de
$p_*(\mathcal{F}(n))$~; le corollaire résulte alors de
(\hyperref[II.3.2.5.2-fr]{3.2.5.2}) et
(\hyperref[II.3.2.7.1-fr]{3.2.7.1}).
\end{proof}

\subsection{Comportements fonctoriels.}
\label{subsection:II.3.5-fr}

\begin{env}[3.5.1]
\label{II.3.5.1-fr}
Soient $Y$ un préschéma, $\mathcal{S}$, $\mathcal{S}'$ deux
$\mathcal{O}_Y$-Algèbres graduées quasi-cohérentes à degrés positifs~; posons
$X=\operatorname{Proj}(\mathcal{S})$, $X'=\operatorname{Proj}(\mathcal{S}')$,
et soient $p$, $p'$ les morphismes structuraux de $X$ et $X'$ dans $Y$. Soit
$\varphi:\mathcal{S}'\to\mathcal{S}$ un $\mathcal{O}_Y$-homomorphisme
d'Algèbres graduées. Pour tout ouvert affine $U$ de $Y$, posons
$S_U=\Gamma(U,\mathcal{S})$, $S'_U=\Gamma(U,\mathcal{S}')$~;
l'homomorphisme $\varphi$ définit un homomorphisme
$\varphi_U:S'_U\to S_U$ de $A_U$-algèbres graduées, en posant
$A_U=\Gamma(U,\mathcal{O}_Y)$. Il lui correspond dans $p^{-1}(U)$ un ensemble
ouvert $G(\varphi_U)$ et un morphisme
$\Phi_U:G(\varphi_U)\to p'^{-1}(U)$
(\hyperref[II.2.8.1-fr]{2.8.1}). En outre, si $V\subset U$ est un ouvert
affine, le diagramme
\[
\label{II.3.5.1.1-fr}
  \xymatrix{
    S'_U\ar[r]^{\varphi_U}\ar[d]
    & S_U\ar[d]\\
    S'_V\ar[r]_{\varphi_V}
    & S_V
  }
\tag{3.5.1.1}
\]
est commutatif, et on vérifie aussitôt, à partir des définitions
(\hyperref[II.2.8.1-fr]{2.8.1}), que l'on a
\[
  G(\varphi_V)=G(\varphi_U)\cap p^{-1}(V)
\]
et que $\Phi_V$ est la restriction de $\Phi_U$ à $G(\varphi_V)$. On a ainsi
défini une partie ouverte $G(\varphi)$ de $X$ telle que
$G(\varphi)\cap p^{-1}(U)=G(\varphi_U)$ pour tout ouvert affine $U\subset Y$,
et un $Y$-morphisme

\oldpage[II]{62}

affine $\Phi:G(\varphi)\to X'$, qui est dit \emph{associé} à $\varphi$ et que
nous noterons $\operatorname{Proj}(\varphi)$. Lorsque, pour tout $y\in Y$, il
y a un voisinage affine $U$ de $y$ tel que le
$\Gamma(U,\mathcal{O}_Y)$-module $\Gamma(U,\mathcal{S}_+)$ soit engendré par
$\varphi(\Gamma(U,\mathcal{S}'_+))$, on a
$G(\varphi_U)=p^{-1}(U)$, et par suite $G(\varphi)=X$.
\end{env}

\begin{proposition}[3.5.2]
\label{II.3.5.2-fr}
\begin{enumerate}
  \item[\rm(i)] Si $\mathcal{M}$ est un $\mathcal{S}$-Module gradué
    quasi-cohérent, il existe un isomorphisme canonique fonctoriel du
    $\mathcal{O}_{X'}$-Module $(\mathcal{M}_{[\varphi]})^{\sim}$ sur le
    $\mathcal{O}_{X'}$-Module
    $\Phi_*(\widetilde{\mathcal{M}}|G(\varphi))$.
  \item[\rm(ii)] Si $\mathcal{M}'$ est un $\mathcal{S}'$-Module gradué
    quasi-cohérent, il existe un homomorphisme canonique fonctoriel $\nu$ du
    $(\mathcal{O}_X|G(\varphi))$-Module
    $\Phi^*(\widetilde{\mathcal{M}'})$ dans le
    $(\mathcal{O}_X|G(\varphi))$-Module
    $(\mathcal{M}'\otimes_{\mathcal{S}'}\mathcal{S})^{\sim}|G(\varphi)$. Si
    $\mathcal{S}'$ est engendrée par $\mathcal{S}'_1$, $\nu$ est un
    isomorphisme.
\end{enumerate}
\end{proposition}

\begin{proof}
Les homomorphismes considérés ont en effet déjà été définis lorsque $Y$ est
affine (\hyperref[II.2.8.7-fr]{2.8.7} et
\hyperref[II.2.8.8-fr]{2.8.8}), et dans le cas général il suffit de vérifier
qu'ils sont compatibles avec les opérations de restriction d'un ouvert
affine de $Y$ à un ouvert plus petit, ce qui résulte aussitôt de la
commutativité de (\hyperref[II.3.5.1.1-fr]{3.5.1.1}).
\end{proof}

En particulier, pour tout $n\in\mathbf{Z}$, on a un homomorphisme canonique
\[
\label{II.3.5.2.1-fr}
  \Phi^*(\mathcal{O}_{X'}(n))\longrightarrow
  \mathcal{O}_X(n)|G(\varphi).
\tag{3.5.2.1}
\]

\begin{proposition}[3.5.3]
\label{II.3.5.3-fr}
Soient $Y$, $Y'$ deux préschémas, $\psi:Y'\to Y$ un morphisme,
$\mathcal{S}$ une $\mathcal{O}_Y$-Algèbre graduée quasi-cohérente, et posons
$\mathcal{S}'=\psi^*(\mathcal{S})$. Alors le $Y'$-schéma
$X'=\operatorname{Proj}(\mathcal{S}')$ s'identifie canoniquement à
$\operatorname{Proj}(\mathcal{S})\times_Y Y'$. En outre, si $\mathcal{M}$ est
un $\mathcal{S}$-Module gradué quasi-cohérent, le $\mathcal{O}_{X'}$-Module
$(\psi^*(\mathcal{M}))^{\sim}$ s'identifie à
$\widetilde{\mathcal{M}}\otimes_Y\mathcal{O}_{Y'}$.
\end{proposition}

\begin{proof}
Notons en premier lieu que $\psi^*(\mathcal{S})$ et
$\psi^*(\mathcal{M})$ sont des $\mathcal{O}_{Y'}$-Modules quasi-cohérents
ainsi que leurs composants homogènes
(\hyperref[0.5.1.4-fr]{0, 5.1.4}). Soient $U$ un ouvert affine de $Y$,
$U'\subset\psi^{-1}(U)$ un ouvert affine de $Y'$, $A$, $A'$ les anneaux de
$U$ et $U'$ respectivement~; on a alors
$\mathcal{S}|U=\widetilde{S}$, où $S$ est une $A$-algèbre graduée, et
$\mathcal{S}'|U'$ s'identifie à $(S\otimes_A A')^{\sim}$
(\hyperref[I.1.6.5-fr]{I, 1.6.5})~; la première assertion résulte alors de
(\hyperref[II.2.8.10-fr]{2.8.10}) et de
(\hyperref[I.3.2.6.2-fr]{I, 3.2.6.2}), car on vérifie aussitôt que la
projection
$\operatorname{Proj}(\mathcal{S}'|U')\to
\operatorname{Proj}(\mathcal{S}|U)$ définie par l'identification précédente
est compatible avec les opérations de restrictions sur $U$ et $U'$ et
définit donc bien un morphisme
$\operatorname{Proj}(\mathcal{S}')\to\operatorname{Proj}(\mathcal{S})$.
Soient maintenant
\[
  q:\operatorname{Proj}(\mathcal{S})\to Y,
  \qquad q':\operatorname{Proj}(\mathcal{S}')\to Y'
\]
les morphismes structuraux~; $q'^{-1}(U')$ s'identifie alors à
$q^{-1}(U)\times_U U'$, et les deux faisceaux
$(\psi^*(\mathcal{M}))^{\sim}|q'^{-1}(U')$ et
$(\widetilde{\mathcal{M}}\otimes_Y\mathcal{O}_{Y'})|q'^{-1}(U')$
s'identifient alors canoniquement tous deux à $(M\otimes_A A')^{\sim}$, où
on a posé $M=\Gamma(U,\mathcal{M})$, en vertu de
(\hyperref[II.2.8.10-fr]{2.8.10}) et de
(\hyperref[I.1.6.5-fr]{I, 1.6.5})~; d'où la seconde assertion, car on vérifie
encore de façon immédiate la compatibilité des identifications précédentes
avec les opérations de restriction.
\end{proof}

\begin{corollary}[3.5.4]
\label{II.3.5.4-fr}
Avec les notations de (\hyperref[II.3.5.3-fr]{3.5.3}),
$\mathcal{O}_{X'}(n)$ s'identifie canoniquement à
$\mathcal{O}_X(n)\otimes_Y\mathcal{O}_{Y'}$ pour tout $n\in\mathbf{Z}$ (avec
$X=\operatorname{Proj}(\mathcal{S})$).
\end{corollary}

\begin{proof}
En effet, avec les notations de (\hyperref[II.3.5.3-fr]{3.5.3}), il est clair
que $\psi^*(\mathcal{S}(n))=\mathcal{S}'(n)$ pour tout $n\in\mathbf{Z}$.
\end{proof}

\begin{env}[3.5.5]
\label{II.3.5.5-fr}
Gardant les notations précédentes, désignons par $\Psi$ la projection
canonique $X'\to X$, et posons $\mathcal{M}'=\psi^*(\mathcal{M})$~; nous
supposerons en outre que $\mathcal{S}$ est engendrée par $\mathcal{S}_1$

\oldpage[II]{63}

et que $X$ est de type fini sur $Y$~; il en résulte alors que $\mathcal{S}'$
est engendrée par $\mathcal{S}'_1$ (comme on le voit en se ramenant au cas où
$Y$ et $Y'$ sont affines) et que $X'$ est de type fini sur $Y'$
(\hyperref[I.6.3.4-fr]{I, 6.3.4}). Soit $\mathcal{F}$ un
$\mathcal{O}_X$-Module et posons $\mathcal{F}'=\Psi^*(\mathcal{F})$~; il
résulte alors de (\hyperref[II.3.5.4-fr]{3.5.4}) et de
(\hyperref[0.4.3.3-fr]{0, 4.3.3}) que l'on a
$\mathcal{F}'(n)=\Psi^*(\mathcal{F}(n))$ pour tout $n\in\mathbf{Z}$. On
définit en outre un $\psi$-homomorphisme canonique
$\theta_n:q_*(\mathcal{F}(n))\to q'_*(\mathcal{F}'(n))$ de la façon
suivante~: vu la commutativité du diagramme
\[
  \xymatrix{
    X\ar[d]_{q}
    & X'\ar[l]_{\Psi}\ar[d]^{q'}\\
    Y
    & Y'\ar[l]^{\psi}
  }
\]
il s'agit de définir un homomorphisme
$q_*(\mathcal{F}(n))\to
\psi_*(q'_*(\Psi^*(\mathcal{F}(n))))
=q_*(\Psi_*(\Psi^*(\mathcal{F}(n))))$, et il suffit de prendre
l'homomorphisme $\theta_n=q_*(\rho_n)$, où $\rho_n$ est l'homomorphisme
canonique
$\mathcal{F}(n)\to\Psi_*(\Psi^*(\mathcal{F}(n)))$
(\hyperref[0.4.4.3-fr]{0, 4.4.3}). Il est immédiat que pour tout ouvert affine
$U$ de $Y$ et tout ouvert affine $U'$ de $Y'$ tel que
$U'\subset\psi^{-1}(U)$, l'homomorphisme $\theta_n$ donne pour les sections
l'homomorphisme canonique (\hyperref[0.3.7.2-fr]{0, 3.7.2})
$\Gamma(q^{-1}(U),\mathcal{F}(n))\to
\Gamma(q'^{-1}(U'),\mathcal{F}'(n))$.

La commutativité de (\hyperref[II.2.8.13.2-fr]{2.8.13.2}) montre alors que si
$\mathcal{F}$ est quasi-cohérent, le diagramme
\[
  \xymatrix{
    \mathcal{F}\ar[r]^{\rho}
    & \mathcal{F}'\\
    (\Gamma_*(\mathcal{F}))^{\sim}
      \ar[u]^{\beta_{\mathcal{F}}}\ar[r]_{\widetilde{\theta}}
    & (\Gamma_*(\mathcal{F}'))^{\sim}
      \ar[u]_{\beta_{\mathcal{F}'}}
  }
\]
est commutatif (la flèche horizontale du haut étant le $\Psi$-morphisme
canonique $\mathcal{F}\to\Psi^*(\mathcal{F})$).

De même, la commutativité de
(\hyperref[II.2.8.13.1-fr]{2.8.13.1}) montre que le diagramme
\[
  \xymatrix{
    \Gamma_*(\widetilde{\mathcal{M}})\ar[r]^{\theta}
    & \Gamma_*(\widetilde{\mathcal{M}'})\\
    \mathcal{M}\ar[r]_{\rho}\ar[u]^{\alpha_{\mathcal{M}}}
    & \mathcal{M}'\ar[u]_{\alpha_{\mathcal{M}'}}
  }
\]
est commutatif (la flèche horizontale du bas étant le $\psi$-morphisme
canonique $\mathcal{M}\to\psi^*(\mathcal{M})$).
\end{env}

\begin{env}[3.5.6]
\label{II.3.5.6-fr}
Considérons maintenant deux préschémas $Y$, $Y'$, un morphisme
$g:Y'\to Y$, une $\mathcal{O}_Y$-Algèbre (resp.
$\mathcal{O}_{Y'}$-Algèbre) graduée quasi-cohérente $\mathcal{S}$ (resp.
$\mathcal{S}'$) et un $g$-morphisme d'Algèbres graduées
$u:\mathcal{S}\to\mathcal{S}'$, c'est-à-dire un
$\mathcal{O}_Y$-homomorphisme d'Algèbres graduées
$\mathcal{S}\to g_*(\mathcal{S}')$~; on sait d'ailleurs qu'il revient au même
de se donner un $\mathcal{O}_{Y'}$-homomorphisme d'Algèbres graduées
$u^\sharp:g^*(\mathcal{S})\to\mathcal{S}'$. On déduit alors canoniquement de
$u^\sharp$ un $Y'$-morphisme
$W=\operatorname{Proj}(u^\sharp):G(u^\sharp)\to
\operatorname{Proj}(g^*(\mathcal{S}))$ où $G(u^\sharp)$ est un ouvert de
$X'=\operatorname{Proj}(\mathcal{S}')$
(\hyperref[II.3.5.1-fr]{3.5.1}). D'autre part,
$X''=\operatorname{Proj}(g^*(\mathcal{S}))$ s'identifie

\oldpage[II]{64}

canoniquement à $X\times_Y Y'$, en posant
$X=\operatorname{Proj}(\mathcal{S})$
(\hyperref[II.3.5.3-fr]{3.5.3})~; composant avec
$\operatorname{Proj}(u^\sharp)$ la première projection
$p:X\times_Y Y'\to X$, on obtient donc un morphisme
$v:G(u^\sharp)\to X$, que nous désignerons par $\operatorname{Proj}(u)$, et
qui est tel que le diagramme
\[
  \xymatrix{
    G(u^\sharp)\ar[r]^{v}\ar[d]
    & X\ar[d]\\
    Y'\ar[r]_{g}
    & Y
  }
\]
soit commutatif.

En outre, pour tout $\mathcal{S}$-Module gradué quasi-cohérent $\mathcal{M}$,
on a un $v$-morphisme canonique
\[
\label{II.3.5.6.1-fr}
  \nu:\widetilde{\mathcal{M}}\longrightarrow
  (g^*(\mathcal{M})\otimes_{g^*(\mathcal{S})}\mathcal{S}')^{\sim}
  |G(u^\sharp).
\tag{3.5.6.1}
\]

En effet, $\nu^\sharp$ s'obtient en composant les homomorphismes
\[
  v^*(\widetilde{\mathcal{M}})
  =W^*(p^*(\widetilde{\mathcal{M}}))
  \longrightarrow W^*((g^*(\mathcal{M}))^{\sim})
  \longrightarrow
  (g^*(\mathcal{M})\otimes_{g^*(\mathcal{S})}\mathcal{S}')^{\sim}
  |G(u^\sharp)
\]
où la première flèche provient de l'isomorphisme
(\hyperref[II.3.5.3-fr]{3.5.3}) et la seconde est l'homomorphisme
(\hyperref[II.3.5.2-fr]{3.5.2, (ii)})~; lorsque $\mathcal{S}$ est engendrée
par $\mathcal{S}_1$, il résulte de (\hyperref[II.3.5.2-fr]{3.5.2}) que
$\nu^\sharp$ est un isomorphisme.

Comme cas particulier de (\hyperref[II.3.5.6.1-fr]{3.5.6.1}), on a, pour
tout $n\in\mathbf{Z}$, un $v$-morphisme canonique
\[
\label{II.3.5.6.2-fr}
  \nu:\mathcal{O}_X(n)\longrightarrow
  \mathcal{O}_{X'}(n)|G(u^\sharp).
\tag{3.5.6.2}
\]
\end{env}

\subsection{Sous-préschémas fermés d'un préschéma
$\operatorname{Proj}(\mathcal{S})$.}
\label{subsection:II.3.6-fr}

\begin{env}[3.6.1]
\label{II.3.6.1-fr}
Soient $Y$ un préschéma,
$\varphi:\mathcal{S}\to\mathcal{S}'$ un homomorphisme de
$\mathcal{O}_Y$-Algèbres graduées quasi-cohérentes, de degré $0$. On dit que
$\varphi$ est (\textbf{TN})-\emph{surjectif} (resp.
(\textbf{TN})-\emph{injectif}, (\textbf{TN})-\emph{bijectif}) s'il existe $n$
tel que, pour $k\geq n$, $\varphi_k:\mathcal{S}_k\to\mathcal{S}'_k$ soit
surjectif (resp. injectif, bijectif). Lorsqu'il en est ainsi on ramène l'étude
du morphisme
$\Phi:\operatorname{Proj}(\mathcal{S}')\to\operatorname{Proj}(\mathcal{S})$
correspondant au cas où $\varphi$ est surjectif (resp. injectif, bijectif).
Cela se démontre comme dans (\hyperref[II.2.9.1-fr]{2.9.1}) (qui en est le cas
particulier correspondant à $Y$ affine) en utilisant
(\hyperref[II.3.1.8-fr]{3.1.8}). Au lieu de dire que $\varphi$ est
(\textbf{TN})-bijectif, on dit aussi que c'est alors un
(\textbf{TN})-isomorphisme.
\end{env}

\begin{proposition}[3.6.2]
\label{II.3.6.2-fr}
Soient $Y$ un préschéma, $\mathcal{S}$ une $\mathcal{O}_Y$-Algèbre graduée
quasi-cohérente, et soit $X=\operatorname{Proj}(\mathcal{S})$.
\begin{enumerate}
  \item[\rm(i)] Si $\varphi:\mathcal{S}\to\mathcal{S}'$ est un
    homomorphisme (\textbf{TN})-surjectif de $\mathcal{O}_Y$-Algèbres
    graduées, le morphisme correspondant
    $\Phi=\operatorname{Proj}(\varphi)$
    (\hyperref[II.3.5.1-fr]{3.5.1}) est défini dans
    $\operatorname{Proj}(\mathcal{S}')$ tout entier et est une immersion
    fermée de $\operatorname{Proj}(\mathcal{S}')$ dans $X$. Si $\mathcal{J}$
    est le noyau de $\varphi$, le sous-préschéma fermé de $X$ associé à
    $\Phi$ est défini par le faisceau d'idéaux quasi-cohérent
    $\widetilde{\mathcal{J}}$ de $\mathcal{O}_X$.
  \item[\rm(ii)] Supposons en outre $\mathcal{S}_0=\mathcal{O}_Y$,
    $\mathcal{S}$ engendrée par $\mathcal{S}_1$ et $\mathcal{S}_1$ de type
    fini. Soit $X'$ un sous-préschéma fermé de
    $X=\operatorname{Proj}(\mathcal{S})$, défini par un faisceau
    quasi-cohérent $\mathcal{I}$ d'idéaux de $\mathcal{O}_X$. Soit
    $\mathcal{J}$

\oldpage[II]{65}

    le faisceau gradué quasi-cohérent d'idéaux de $\mathcal{S}$, image
    réciproque de $\Gamma_*(\mathcal{I})$ par l'homomorphisme canonique
    $\alpha:\mathcal{S}\to\Gamma_*(\mathcal{O}_X)$
    (\hyperref[II.3.3.2-fr]{3.3.2}), et posons
    $\mathcal{S}'=\mathcal{S}/\mathcal{J}$. Alors $X'$ est le sous-préschéma
    associé (\hyperref[I.4.2.1-fr]{I, 4.2.1}) à l'immersion fermée
    $\operatorname{Proj}(\mathcal{S}')\to X$ correspondant à l'homomorphisme
    canonique $\mathcal{S}\to\mathcal{S}'$ de $\mathcal{O}_Y$-Algèbres
    graduées.
\end{enumerate}
\end{proposition}

\begin{proof}
\begin{enumerate}
  \item[\rm(i)] On peut supposer $\varphi$ surjectif
    (\hyperref[II.3.6.1-fr]{3.6.1}). Alors, pour tout ouvert affine $U$ de
    $Y$, $\Gamma(U,\mathcal{S})\to\Gamma(U,\mathcal{S}')$ est surjectif
    (\hyperref[I.1.3.9-fr]{I, 1.3.9}), donc
    (\hyperref[II.3.5.1-fr]{3.5.1}) on a $G(\varphi)=X$. On est
    immédiatement ramené à démontrer la proposition lorsque $Y$ est affine,
    et alors elle découle de (\hyperref[II.2.9.2-fr]{2.9.2, (i)}).
  \item[\rm(ii)] On est ramené à prouver que l'homomorphisme
    $\widetilde{\mathcal{J}}\to\mathcal{O}_X$ déduit de l'injection canonique
    $\mathcal{J}\to\mathcal{S}$ est un isomorphisme de
    $\widetilde{\mathcal{J}}$ sur $\mathcal{I}$~; comme la question est locale
    sur $Y$, on peut supposer $Y$ affine d'anneau $A$, ce qui entraîne
    $\mathcal{S}=\widetilde{S}$, où $S$ est une $A$-algèbre graduée engendrée
    par $S_1$, $S_1$ étant de type fini sur $A$. Il suffit alors d'appliquer
    (\hyperref[II.2.9.2-fr]{2.9.2, (ii)}).
\end{enumerate}
\end{proof}

\begin{corollary}[3.6.3]
\label{II.3.6.3-fr}
Sous les conditions de (\hyperref[II.3.6.2-fr]{3.6.2, (i)}), et en supposant
de plus $\mathcal{S}$ engendrée par $\mathcal{S}_1$,
$\Phi^*(\mathcal{O}_X(n))$ s'identifie canoniquement à
$\mathcal{O}_{X'}(n)$ pour tout $n\in\mathbf{Z}$.
\end{corollary}

\begin{proof}
On a défini un tel isomorphisme canonique lorsque $Y$ est affine
(\hyperref[II.2.9.3-fr]{2.9.3})~; dans le cas général, il suffit de vérifier
que les isomorphismes ainsi définis pour chaque ouvert affine $U$ de $Y$ sont
compatibles avec le passage de $U$ à un ouvert affine $U'\subset U$, ce qui
est immédiat.
\end{proof}

\begin{corollary}[3.6.4]
\label{II.3.6.4-fr}
Soient $Y$ un préschéma, $\mathcal{S}$ une $\mathcal{O}_Y$-Algèbre graduée
quasi-cohérente engendrée par $\mathcal{S}_1$, $\mathcal{M}$ un
$\mathcal{O}_Y$-Module quasi-cohérent, $u$ un $\mathcal{O}_Y$-homomorphisme
surjectif $\mathcal{M}\to\mathcal{S}_1$,
$\overline{u}:\mathbf{S}_{\mathcal{O}_Y}(\mathcal{M})\to\mathcal{S}$
l'homomorphisme de $\mathcal{O}_Y$-Algèbres graduées prolongeant $u$
(\hyperref[II.1.7.4-fr]{1.7.4}). Le morphisme correspondant à
$\overline{u}$ est alors une immersion fermée de
$\operatorname{Proj}(\mathcal{S})$ dans
$\operatorname{Proj}(\mathbf{S}_{\mathcal{O}_Y}(\mathcal{M}))$.
\end{corollary}

\begin{proof}
En effet, $\overline{u}$ est surjectif par hypothèse et on applique
(\hyperref[II.3.6.1-fr]{3.6.1, (i)}).
\end{proof}

\subsection{Morphismes d'un préschéma dans un spectre homogène.}
\label{subsection:II.3.7-fr}

\begin{env}[3.7.1]
\label{II.3.7.1-fr}
Soient $q:X\to Y$ un morphisme de préschémas, $\mathcal{L}$ un
$\mathcal{O}_X$-Module inversible, $\mathcal{S}$ une $\mathcal{O}_Y$-Algèbre
graduée quasi-cohérente à degrés positifs~; $q^*(\mathcal{S})$ est alors une
$\mathcal{O}_X$-Algèbre graduée quasi-cohérente à degrés positifs.
Considérons la $\mathcal{O}_X$-Algèbre graduée quasi-cohérente
$\mathcal{S}'=\bigoplus_{n\geq0}\mathcal{L}^{\otimes n}$ et supposons donné
un $\mathcal{O}_X$-homomorphisme d'Algèbres graduées
\[
  \psi:q^*(\mathcal{S})\longrightarrow
  \mathcal{S}'=\bigoplus_{n\geq0}\mathcal{L}^{\otimes n}
\]
ce qui revient d'ailleurs à se donner un $q$-morphisme d'Algèbres graduées
\[
  \psi^\flat:\mathcal{S}\longrightarrow q_*(\mathcal{S}').
\]
On sait que $\operatorname{Proj}(\mathcal{S}')$ s'identifie canoniquement à
$X$ (\hyperref[II.3.1.7-fr]{3.1.7} et
\hyperref[II.3.1.8-fr]{3.1.8, (iii)})~; on déduit canoniquement de $\psi$ un
ensemble ouvert $G(\psi)$ de $X$ et un $Y$-morphisme
\[
\label{II.3.7.1.1-fr}
  r_{\mathcal{L},\psi}:G(\psi)\longrightarrow
  \operatorname{Proj}(\mathcal{S})=P
\tag{3.7.1.1}
\]

\oldpage[II]{66}

que nous appellerons le morphisme \emph{associé à $\mathcal{L}$ et à $\psi$}~;
rappelons (\hyperref[II.3.5.6-fr]{3.5.6}) que ce morphisme s'obtient en
composant avec la première projection
$\pi:\operatorname{Proj}(q^*(\mathcal{S}))=P\times_YX\to P$ le
$Y$-morphisme
\[
  \tau=\operatorname{Proj}(\psi):G(\psi)\longrightarrow
  \operatorname{Proj}(q^*(\mathcal{S})).
\]
\end{env}

\begin{env}[3.7.2]
\label{II.3.7.2-fr}
Explicitons $r=r_{\mathcal{L},\psi}$ lorsque $Y=\operatorname{Spec}(A)$ est
affine, et par suite $\mathcal{S}=\widetilde{S}$, où $S$ est une $A$-algèbre
graduée à degrés positifs. Supposons d'abord que $X=\operatorname{Spec}(B)$
soit aussi affine et que l'on ait $\mathcal{L}=\widetilde{L}$, où $L$ est un
$B$-module libre de rang $1$. On a alors
$q^*(\mathcal{S})=(S\otimes_AB)^{\sim}$
(\hyperref[I.1.6.5-fr]{I, 1.6.5})~; si $c$ est un générateur de $L$,
$\psi_n:q^*(\mathcal{S}_n)\to\mathcal{L}^{\otimes n}$ correspond à un
homomorphisme
$w_n:s\otimes b\mapsto bv_n(s)c^{\otimes n}$ de $S_n\otimes_AB$ dans
$L^{\otimes n}$, où $v_n:S_n\to B$ est un homomorphisme de $A$-modules, les
$v_n$ constituant un homomorphisme d'algèbres $S\to B$. Soit $f\in S_d$
($d>0$) et posons $g=v_d(f)$~; on a
$\pi^{-1}(D_+(f))=D_+(f\otimes1)$ en vertu de
(\hyperref[II.2.8.10-fr]{2.8.10}) et de l'identification de $D_+(f)$ à
$\operatorname{Spec}(S_{(f)})$ (\hyperref[II.2.3.6-fr]{2.3.6})~; d'autre
part, la formule (\hyperref[II.2.8.1.1-fr]{2.8.1.1}) montre (compte tenu de
l'identification canonique de $X$ et de $\operatorname{Proj}(\mathcal{S}')$)
que
\[
  \tau^{-1}(D_+(f\otimes1))=D(g)
\]
d'où
\[
\label{II.3.7.2.1-fr}
  r^{-1}(D_+(f))=D(g).
\tag{3.7.2.1}
\]

En outre, le morphisme $\tau=\operatorname{Proj}(\psi)$, restreint à $D(g)$,
correspond à l'homomorphisme qui applique
$(s\otimes1)/(f\otimes1)^n$ (pour $s\in S_{nd}$) sur
$v_{nd}(s)/g^n$ (\hyperref[II.2.8.1-fr]{2.8.1}), et la projection $\pi$,
restreinte à $D_+(f\otimes1)$, correspond à l'homomorphisme
$s/f^n\mapsto(s\otimes1)/(f\otimes1)^n$~; on en conclut que $r$, restreint à
$D(g)$, correspond à l'homomorphisme de $A$-algèbres
$\omega:S_{(f)}\to B_g$ tel que $\omega(s/f^n)=v_{nd}(s)/g^n$ pour
$s\in S_{nd}$ ($n>0$). Passant au cas où $X$ est quelconque ($Y$ étant
toujours affine), on obtient donc, compte tenu de
(\hyperref[II.2.8.1-fr]{2.8.1})~:
\end{env}

\begin{proposition}[3.7.3]
\label{II.3.7.3-fr}
Si $Y=\operatorname{Spec}(A)$ est affine et $\mathcal{S}=\widetilde{S}$, où
$S$ est une $A$-algèbre graduée, pour tout
$f\in S_d=\Gamma(Y,\mathcal{S}_d)$, on a
\[
\label{II.3.7.3.1-fr}
  r_{\mathcal{L},\psi}^{-1}(D_+(f))=X_{\psi^\flat(f)}
  \qquad\bigl(\text{où }\psi^\flat(f)\in
  \Gamma(X,\mathcal{L}^{\otimes d})\bigr)
\tag{3.7.3.1}
\]
et la restriction $X_{\psi^\flat(f)}\to D_+(f)=\operatorname{Spec}(S_{(f)})$
de $r_{\mathcal{L},\psi}$ correspond
(\hyperref[I.2.2.4-fr]{I, 2.2.4}) à l'homomorphisme d'algèbres
\[
\label{II.3.7.3.2-fr}
  \psi_{(f)}^\flat:S_{(f)}\longrightarrow
  \Gamma(X_{\psi^\flat(f)},\mathcal{O}_X)
\tag{3.7.3.2}
\]
tel que, pour $s\in S_{nd}=\Gamma(Y,\mathcal{S}_{nd})$,
\[
\label{II.3.7.3.3-fr}
  \psi_{(f)}^\flat(s/f^n)
  =(\psi^\flat(s)|X_{\psi^\flat(f)})
   (\psi^\flat(f)|X_{\psi^\flat(f)})^{-n}.
\tag{3.7.3.3}
\]
\end{proposition}

Nous dirons que $r_{\mathcal{L},\psi}$ est \emph{partout défini} si l'on a
$G(\psi)=X$. Pour qu'il en soit ainsi, il faut et il suffit évidemment que
$G(\psi)\cap q^{-1}(U)=q^{-1}(U)$ pour tout ouvert affine $U\subset Y$~;
autrement dit, la question est \emph{locale sur $Y$}. Si $Y$ est affine,
$G(\psi)$ est réunion des $r^{-1}(D_+(f))$ pour $f$ homogène dans $S_+$
(\hyperref[II.2.8.1-fr]{2.8.1})~; en vertu de
(\hyperref[II.3.7.3.1-fr]{3.7.3.1}), les $X_{\psi^\flat(f)}$ doivent donc
former un \emph{recouvrement} de $X$, autrement dit~:

\begin{corollary}[3.7.4]
\label{II.3.7.4-fr}
Sous les hypothèses de (\hyperref[II.3.7.3-fr]{3.7.3}), pour que
$r_{\mathcal{L},\psi}$ soit partout défini, il faut et il suffit que pour tout
$x\in X$, il existe un entier $n>0$ et une section $s$ de $\mathcal{S}_n$
au-dessus de $Y$ telle que si on pose
$t=\psi^\flat(s)\in\Gamma(X,\mathcal{L}^{\otimes n})$, on ait $t(x)\neq0$.
\end{corollary}

\oldpage[II]{67}

On notera que cette condition est toujours vérifiée si $\psi$ est
(\textbf{TN})-surjectif.

De même, la question de savoir si $r_{\mathcal{L},\psi}$ est \emph{dominant}
est locale sur $Y$, et l'on a~:

\begin{corollary}[3.7.5]
\label{II.3.7.5-fr}
Sous les hypothèses de (\hyperref[II.3.7.3-fr]{3.7.3}), pour que
$r_{\mathcal{L},\psi}$ soit dominant, il faut et il suffit que, pour tout
entier $n>0$, toute section $s\in S_n$ telle que
$\psi^\flat(s)\in\Gamma(X,\mathcal{L}^{\otimes n})$ soit localement
nilpotente, soit elle-même nilpotente.
\end{corollary}

\begin{proof}
On doit en effet exprimer que
$r_{\mathcal{L},\psi}^{-1}(D_+(s))$ n'est vide que si $D_+(s)$ est vide,
et le corollaire résulte donc de
(\hyperref[II.3.7.3.1-fr]{3.7.3.1}) et de
(\hyperref[II.2.3.7-fr]{2.3.7}).
\end{proof}

\begin{proposition}[3.7.6]
\label{II.3.7.6-fr}
Soient $q:X\to Y$ un morphisme, $\mathcal{L}$ un $\mathcal{O}_X$-Module
inversible, $\mathcal{S}$, $\mathcal{S}'$ deux $\mathcal{O}_Y$-Algèbres
graduées quasi-cohérentes,
$u:\mathcal{S}'\to\mathcal{S}$ un homomorphisme d'Algèbres graduées,
$\psi:q^*(\mathcal{S})\to\bigoplus_{n\geq0}\mathcal{L}^{\otimes n}$ un
homomorphisme d'Algèbres graduées,
$\psi'=\psi\circ q^*(u)$ l'homomorphisme composé. Si
$r_{\mathcal{L},\psi'}$ est partout défini, il en est de même de
$r_{\mathcal{L},\psi}$~; si $u$ est (\textbf{TN})-surjectif et si
$r_{\mathcal{L},\psi'}$ est dominant, il en est de même de
$r_{\mathcal{L},\psi}$~; inversement, si $u$ est (\textbf{TN})-injectif et si
$r_{\mathcal{L},\psi}$ est dominant, $r_{\mathcal{L},\psi'}$ est dominant.
\end{proposition}

\begin{proof}
On a en effet $G(\psi')\subset G(\psi)$
(\hyperref[II.2.8.4-fr]{2.8.4}), d'où la première assertion~; si $u$ est
(\textbf{TN})-surjectif,
$\operatorname{Proj}(u):\operatorname{Proj}(\mathcal{S})\to
\operatorname{Proj}(\mathcal{S}')$ est partout défini et est une immersion
fermée~; comme $r_{\mathcal{L},\psi'}$ est composé de
$\operatorname{Proj}(u)$ et de la restriction de $r_{\mathcal{L},\psi}$ à
$G(\psi')$, on en conclut que si $r_{\mathcal{L},\psi'}$ est dominant, il en
est de même de $r_{\mathcal{L},\psi}$. Enfin, si $u$ est
(\textbf{TN})-injectif, on sait que $\operatorname{Proj}(u)$ est un morphisme
dominant de $G(u)$ dans $\operatorname{Proj}(\mathcal{S}')$
(\hyperref[II.2.8.3-fr]{2.8.3})~; comme $G(\psi')$ est l'image réciproque de
$G(u)$ par $r_{\mathcal{L},\psi}$, on voit que si $r_{\mathcal{L},\psi}$ est
dominant, il en est de même de $r_{\mathcal{L},\psi'}$.
\end{proof}

\begin{proposition}[3.7.7]
\label{II.3.7.7-fr}
Soient $Y$ un préschéma quasi-compact, $q:X\to Y$ un morphisme quasi-compact,
$\mathcal{L}$ un $\mathcal{O}_X$-Module inversible, $\mathcal{S}$ une
$\mathcal{O}_Y$-Algèbre graduée quasi-cohérente, limite inductive filtrante
d'un système inductif $(\mathcal{S}^\lambda)$ de $\mathcal{O}_Y$-Algèbres
quasi-cohérentes. Soient
$\varphi_\lambda:\mathcal{S}^\lambda\to\mathcal{S}$ l'homomorphisme
canonique,
$\psi:q^*(\mathcal{S})\to\bigoplus_{n\geq0}\mathcal{L}^{\otimes n}$ un
homomorphisme d'Algèbres graduées, et posons
$\psi_\lambda=\psi\circ q^*(\varphi_\lambda)$. Pour que
$r_{\mathcal{L},\psi}$ soit partout défini, il faut et il suffit qu'il existe
un $\lambda$ tel que $r_{\mathcal{L},\psi_\lambda}$ soit partout défini~;
$r_{\mathcal{L},\psi_\mu}$ est alors partout défini pour $\mu\geq\lambda$.
\end{proposition}

\begin{proof}
La condition est suffisante en vertu de
(\hyperref[II.3.7.6-fr]{3.7.6}). Inversement, supposons
$r_{\mathcal{L},\psi}$ partout défini~; on peut se ramener au cas où $Y$ est
affine, car si pour tout ouvert affine $U\subset Y$, il existe $\lambda(U)$
tel que la restriction de $r_{\mathcal{L},\psi_{\lambda(U)}}$ à $q^{-1}(U)$
soit partout définie, il suffira ($Y$ étant quasi-compact) de recouvrir $Y$
par un nombre fini d'ouverts affines $U_i$, et de prendre
$\lambda\geq\lambda(U_i)$ pour tous les indices $i$, en vertu de
(\hyperref[II.3.7.6-fr]{3.7.6}). Si $Y$ est affine, l'hypothèse entraîne que
pour tout $x\in X$, il y a une section $s^{(x)}$ d'un $S_n$ telle que, si on
pose $t^{(x)}=\psi^\flat(s^{(x)})$, on ait $t^{(x)}(x)\neq0$
($t^{(x)}$ étant considérée comme section de $\mathcal{L}^{\otimes n}$
au-dessus de $X$), ce qui entraîne $t^{(x)}(z)\neq0$ pour tout $z$ dans un
voisinage $V(x)$ de $x$. Couvrons $X$ par un nombre fini de $V(x_i)$ et
soient $s^{(i)}$ les sections correspondantes de $S$~; il existe alors un
$\lambda$ tel que les $s^{(i)}$ soient toutes de la forme
$\varphi_\lambda(s_\lambda^{(i)})$, avec $s_\lambda^{(i)}\in S^\lambda$ pour
tous les $i$~; il résulte donc de (\hyperref[II.3.7.4-fr]{3.7.4}) que
$r_{\mathcal{L},\psi_\lambda}$ est partout défini. La dernière assertion est
conséquence triviale de (\hyperref[II.3.7.6-fr]{3.7.6}).
\end{proof}

\begin{corollary}[3.7.8]
\label{II.3.7.8-fr}
Sous les hypothèses de (\hyperref[II.3.7.7-fr]{3.7.7}), si les
$r_{\mathcal{L},\psi_\lambda}$ sont dominants, il en est de même de
$r_{\mathcal{L},\psi}$~; la réciproque est vraie lorsque les
$\varphi_\lambda$ sont injectifs.
\end{corollary}

\oldpage[II]{68}

\begin{proof}
La seconde assertion est un cas particulier de
(\hyperref[II.3.7.6-fr]{3.7.6})~; d'autre part, pour montrer que
$r_{\mathcal{L},\psi}$ est dominant, on peut se borner au cas où $Y$ est
affine~; si $s\in S$ est telle que $\psi^\flat(s)$ soit localement
nilpotente, comme on peut écrire $s=\varphi_\lambda(s_\lambda)$ pour un
$\lambda$ au moins, on conclut de l'hypothèse et de
(\hyperref[II.3.7.5-fr]{3.7.5}) que $s_\lambda$ est nilpotente, donc il en est
de même de $s$, et le critère de (\hyperref[II.3.7.5-fr]{3.7.5}) s'applique.
\end{proof}

\begin{namedtheorem}[3.7.9]{Remarques}
\label{II.3.7.9-fr}
\begin{enumerate}
  \item[(i)] Avec les notations de (\hyperref[II.3.7.1-fr]{3.7.1}) et en
    tenant compte de (\hyperref[II.3.2.10-fr]{3.2.10}), on a, pour tout
    $n\in\mathbf{Z}$, un homomorphisme canonique
    \[
    \label{II.3.7.9.1-fr}
      \theta:r_{\mathcal{L},\psi}^*(\mathcal{O}_P(n))\longrightarrow
      \mathcal{L}^{\otimes n}
    \tag{3.7.9.1}
    \]
    défini de façon générale dans
    (\hyperref[II.3.5.6.2-fr]{3.5.6.2}). On voit aussitôt que sous les
    conditions de (\hyperref[II.3.7.3-fr]{3.7.3}), la restriction de cet
    homomorphisme à $X_{\psi^\flat(f)}$ s'explicite comme faisant
    correspondre à $s/f^k$ ($s\in S_{n+kd}$) l'élément
    $(\psi^\flat(s)|X_{\psi^\flat(f)})
    (\psi^\flat(f)|X_{\psi^\flat(f)})^{-k}$.
  \item[(ii)] Soit $\mathcal{F}$ un $\mathcal{O}_X$-Module quasi-cohérent, et
    supposons que $q$ soit quasi-compact et séparé, de sorte que pour tout
    $n\geq0$, $q_*(\mathcal{F}\otimes\mathcal{L}^{\otimes n})$ soit un
    $\mathcal{O}_Y$-Module quasi-cohérent
    (\hyperref[I.9.2.2-fr]{I, 9.2.2}). Soit
    $\mathcal{M}'=\bigoplus_{n\geq0}\mathcal{F}\otimes
    \mathcal{L}^{\otimes n}$, qui est un $\mathcal{S}'$-Module gradué
    quasi-cohérent, et considérons son image
    $\mathcal{M}=q_*(\mathcal{M}')=
    \bigoplus_{n\geq0}q_*(\mathcal{F}\otimes\mathcal{L}^{\otimes n})$
    (qui est un $\mathcal{S}$-Module quasi-cohérent, au moyen de
    l'homomorphisme $\psi^\flat$). Nous allons voir qu'il y a un
    homomorphisme canonique de $\mathcal{O}_X$-Modules
    \[
    \label{II.3.7.9.2-fr}
      \xi:r_{\mathcal{L},\psi}^*(\widetilde{\mathcal{M}})
      \longrightarrow\mathcal{F}|G(\psi).
    \tag{3.7.9.2}
    \]

    En effet, on a défini (\hyperref[II.3.5.6.1-fr]{3.5.6.1}) un
    homomorphisme canonique
    \[
    \label{II.3.7.9.3-fr}
      r_{\mathcal{L},\psi}^*(\widetilde{\mathcal{M}})
      \longrightarrow
      (q^*(\mathcal{M})\otimes_{q^*(\mathcal{S})}\mathcal{S}')^{\sim}
      |G(\psi),
    \tag{3.7.9.3}
    \]
    où le second membre est considéré comme un faisceau quasi-cohérent sur
    $\operatorname{Proj}(\mathcal{S}')$. D'autre part, on a un
    homomorphisme canonique
    \[
    \label{II.3.7.9.4-fr}
      q^*(q_*(\mathcal{M}'))\otimes_{q^*(\mathcal{S})}\mathcal{S}'
      \longrightarrow\mathcal{M}'
    \tag{3.7.9.4}
    \]
    pour tout $\mathcal{S}'$-Module gradué quasi-cohérent $\mathcal{M}'$~:
    pour tout ouvert $U$ de $X$, toute section $t'$ de
    $q^*(q_*(\mathcal{M}'_h))$ au-dessus de $U$ et toute section $b'$ de
    $\mathcal{S}'_k$ au-dessus de $U$, on fait correspondre à $t'\otimes b'$ la
    section $b'\sigma(t')$ de $\mathcal{M}'_{h+k}$, où $\sigma(t')$ est la
    section de $\mathcal{M}'_h$ au-dessus de $U$ qui correspond canoniquement
    (\hyperref[0.4.4.3-fr]{0, 4.4.3}) à $t'$. On en conclut un homomorphisme
    canonique
    \[
    \label{II.3.7.9.5-fr}
      (q^*(q_*(\mathcal{M}'))\otimes_{q^*(\mathcal{S})}\mathcal{S}')^{\sim}
      |G(\psi)\longrightarrow\widetilde{\mathcal{M}'}|G(\psi)
    \tag{3.7.9.5}
    \]
    et comme finalement $\widetilde{\mathcal{M}'}$ s'identifie
    canoniquement à $\mathcal{F}$
    (\hyperref[II.3.2.9-fr]{3.2.9, (i)}), on obtient bien l'homomorphisme
    canonique cherché.

    Sous les conditions de (\hyperref[II.3.7.3-fr]{3.7.3}), la restriction
    de cet homomorphisme à $X_{\psi^\flat(f)}$ s'explicite de la façon
    suivante~: étant donnée une section $t_{nd}$ de
    $\mathcal{F}\otimes\mathcal{L}^{\otimes nd}$ au-dessus de $X$, si
    $t'_{nd}$ est la section $t_{nd}$ considérée comme section de
    $q_*(\mathcal{F}\otimes\mathcal{L}^{\otimes nd})$ au-dessus de $Y$, à
    l'élément $t'_{nd}/f^n$ il correspond la section
    $(t_{nd}|X_{\psi^\flat(f)})
    (\psi^\flat(f)|X_{\psi^\flat(f)})^{-n}$ de $\mathcal{F}$ au-dessus de
    $X_{\psi^\flat(f)}$.
\end{enumerate}
\end{namedtheorem}

\oldpage[II]{69}

\subsection{Critères d'immersion dans un spectre homogène.}
\label{subsection:II.3.8-fr}

\begin{env}[3.8.1]
\label{II.3.8.1-fr}
Avec les notations de (\hyperref[II.3.7.1-fr]{3.7.1}), la question de savoir
si $r_{\mathcal{L},\psi}$ est une immersion (resp. une immersion ouverte, une
immersion fermée) est évidemment \emph{locale sur $Y$}.
\end{env}

\begin{proposition}[3.8.2]
\label{II.3.8.2-fr}
Sous les hypothèses de (\hyperref[II.3.7.3-fr]{3.7.3}), pour que
$r_{\mathcal{L},\psi}$ soit partout définie et soit une immersion, il faut et
il suffit qu'il existe une famille de sections $s_\alpha\in S_{n_\alpha}$
($n_\alpha>0$) telle que, si on pose
$f_\alpha=\psi^\flat(s_\alpha)$, les conditions suivantes soient vérifiées~:
\begin{enumerate}
  \item[(i)] Les $X_{f_\alpha}$ forment un recouvrement de $X$.
  \item[(ii)] Les $X_{f_\alpha}$ sont des ouverts affines.
  \item[(iii)] Pour tout $\alpha$ et tout
    $t\in\Gamma(X_{f_\alpha},\mathcal{O}_X)$, il existe un entier $m>0$ et
    un $s\in S_{mn_\alpha}$ tels que
    $t=(\psi^\flat(s)|X_{f_\alpha})(f_\alpha|X_{f_\alpha})^{-m}$.
\end{enumerate}

Pour que $r_{\mathcal{L},\psi}$ soit partout définie et soit une immersion
ouverte, il faut et il suffit qu'il existe une famille $(s_\alpha)$ vérifiant
les conditions (i), (ii), (iii) et~:
\begin{enumerate}
  \item[(iv)] Pour tout $n>0$ et tout $s\in S_{nn_\alpha}$ tel que
    $\psi^\flat(s)|X_{f_\alpha}=0$, il existe un entier $k>0$ tel que
    $s_\alpha^ks=0$.
\end{enumerate}

Pour que $r_{\mathcal{L},\psi}$ soit partout définie et soit une immersion
fermée, il faut et il suffit qu'il existe une famille $(s_\alpha)$ vérifiant
(i), (ii), (iii) et~:
\begin{enumerate}
  \item[(v)] Les $D_+(s_\alpha)$ forment un recouvrement de
    $P=\operatorname{Proj}(S)$.
\end{enumerate}
\end{proposition}

\begin{proof}
En effet, pour que $r$ soit une immersion (resp. une immersion fermée), il
faut et il suffit qu'il existe un recouvrement de $r(G(\psi))$ (resp. de
$P$) par des $D_+(s_\alpha)$ tel que si on pose
$V_\alpha=r^{-1}(D_+(s_\alpha))$, la restriction de $r$ à $U_\alpha$ soit
une immersion fermée de $V_\alpha$ dans $D_+(s_\alpha)$
(\hyperref[I.4.2.4-fr]{I, 4.2.4}). La condition (i) exprime à la fois que
$r$ est partout définie et que les $D_+(s_\alpha)$ recouvrent $r(X)$,
d'après (\hyperref[II.3.7.3.1-fr]{3.7.3.1})~; comme $D_+(s_\alpha)$ est
affine, les conditions (ii) et (iii) expriment que la restriction de $r$ à
$X_{f_\alpha}$ est une immersion fermée dans $D_+(s_\alpha)$, d'après
(\hyperref[I.4.2.3-fr]{I, 4.2.3})~; enfin, comme (iii) et (iv) expriment que
$\psi_{(s_\alpha)}^\flat$ est bijective (notation de
(\hyperref[II.3.7.3.2-fr]{3.7.3.2})), (ii), (iii) et (iv) expriment que la
restriction de $r$ à $X_{f_\alpha}$ est un isomorphisme sur
$D_+(s_\alpha)$ pour tout $\alpha$, donc (i), (ii), (iii) et (iv) expriment
que $r$ est une immersion ouverte.
\end{proof}

\begin{corollary}[3.8.3]
\label{II.3.8.3-fr}
Sous les hypothèses de (\hyperref[II.3.7.6-fr]{3.7.6}), si
$r_{\mathcal{L},\psi'}$ est partout définie et est une immersion, il en est
de même de $r_{\mathcal{L},\psi}$. Si on suppose en outre que $u$ est
(\textbf{TN})-surjectif et si $r_{\mathcal{L},\psi'}$ est une immersion
ouverte (resp. fermée), il en est de même de $r_{\mathcal{L},\psi}$.
\end{corollary}

\begin{proof}
En effet, d'après (\hyperref[II.3.8.2-fr]{3.8.2}), il y a une famille de
$s'_\alpha\in S'_{n_\alpha}$ telle que si on pose
$f_\alpha=\psi'^\flat(s'_\alpha)$, les conditions (i), (ii), (iii) soient
vérifiées. Or, si on pose $s_\alpha=u(s'_\alpha)$, on a aussi
$f_\alpha=\psi^\flat(s_\alpha)$, et si on a
$t=(\psi'^\flat(s')|X_{f_\alpha})(f_\alpha|X_{f_\alpha})^{-m}$, on a aussi
$t=(\psi^\flat(s)|X_{f_\alpha})(f_\alpha|X_{f_\alpha})^{-m}$ en posant
$s=u(s')$, d'où la première assertion. La seconde résulte aussitôt de ce que
$\operatorname{Proj}(u)$ est une immersion fermée.
\end{proof}

\begin{proposition}[3.8.4]
\label{II.3.8.4-fr}
Supposons vérifiées les hypothèses de
(\hyperref[II.3.7.7-fr]{3.7.7}) et en outre que $q:X\to Y$ soit un morphisme
de type fini. Alors, pour que $r_{\mathcal{L},\psi}$ soit partout définie et
soit une immersion, il faut et il suffit qu'il existe $\lambda$ tel que
$r_{\mathcal{L},\psi_\lambda}$ soit partout définie et soit une immersion et
dans ce cas $r_{\mathcal{L},\psi_\mu}$ est partout définie et est une
immersion pour $\mu\geq\lambda$.
\end{proposition}

\oldpage[II]{70}

\begin{proof}
Compte tenu de (\hyperref[II.3.8.3-fr]{3.8.3}), il suffit de démontrer que si
$r_{\mathcal{L},\psi}$ est partout définie et est une immersion, il en est de
même de $r_{\mathcal{L},\psi_\lambda}$ pour un $\lambda$ au moins. Par le
même raisonnement que dans (\hyperref[II.3.7.7-fr]{3.7.7}), utilisant la
quasi-compacité de $Y$, on se ramène d'abord au cas où $Y$ est affine. Comme
$X$ est quasi-compact, (\hyperref[II.3.8.2-fr]{3.8.2}) montre l'existence
d'une famille \emph{finie} $(s_i)$ d'éléments de $S$ ($s_i\in S_{n_i}$)
ayant les propriétés (i), (ii) et (iii). Le morphisme $X_{f_i}\to Y$ (où
$f_i=\psi^\flat(s_i)$) est de type fini~: en effet, c'est un morphisme de
schémas affines, donc il est quasi-compact
(\hyperref[I.6.6.1-fr]{I, 6.6.1}), et localement de type fini puisque $q$ est
de type fini (\hyperref[I.6.3.2-fr]{I, 6.3.2}), et la conclusion résulte de
(\hyperref[I.6.6.3-fr]{I, 6.6.3}). L'anneau $B_i$ de $X_{f_i}$ est par suite
une $A$-algèbre de type fini (\hyperref[I.6.3.3-fr]{I, 6.3.3})~; soit
$(t_{ij})$ une famille de générateurs de cette algèbre. Il y a par hypothèse
des éléments $s'_{ij}\in S_{m_{ij}n_i}$ tels que
\[
  t_{ij}=(\psi^\flat(s'_{ij})|X_{f_i})
  (\psi^\flat(s_i)|X_{f_i})^{-m_{ij}}.
\]
Il y a par hypothèse un $\lambda$ et des éléments
$s_{i\lambda}\in S_{n_i}^\lambda$,
$s'_{ij\lambda}\in S_{m_{ij}n_i}^\lambda$ dont les images par
$\varphi_\lambda$ sont respectivement $s_i$ et $s'_{ij}$~; il est clair que
$r_{\mathcal{L},\psi_\lambda}$ vérifie les conditions (i), (ii) et (iii) de
(\hyperref[II.3.8.2-fr]{3.8.2}).
\end{proof}

\begin{proposition}[3.8.5]
\label{II.3.8.5-fr}
Soient $Y$ un schéma quasi-compact, ou un préschéma dont l'espace sous-jacent
est noethérien, $q:X\to Y$ un morphisme de type fini, $\mathcal{L}$ un
$\mathcal{O}_X$-Module inversible, $\mathcal{S}$ une $\mathcal{O}_Y$-Algèbre
graduée quasi-cohérente,
$\psi:q^*(\mathcal{S})\to\bigoplus_{n\geq0}\mathcal{L}^{\otimes n}$ un
homomorphisme d'Algèbres graduées. Pour que $r_{\mathcal{L},\psi}$ soit
partout défini et soit une immersion, il faut et il suffit qu'il existe un
entier $n>0$ et un sous-$\mathcal{O}_Y$-Module quasi-cohérent de type fini
$\mathcal{E}$ de $\mathcal{S}_n$, tels que~:
\begin{enumerate}
  \item[a)] l'homomorphisme
    $\psi_n\circ q^*(j_n):q^*(\mathcal{E})\to\mathcal{L}^{\otimes n}$
    (où $j_n:\mathcal{E}\to\mathcal{S}$ est l'injection canonique) soit
    surjectif~;
  \item[b)] si on désigne par $\mathcal{S}'$ la sous-$\mathcal{O}_Y$-Algèbre
    (graduée) de $\mathcal{S}$ engendrée par $\mathcal{E}$, par $\psi'$
    l'homomorphisme $\psi\circ q^*(j')$ ($j'$ injection de $\mathcal{S}'$
    dans $\mathcal{S}$), $r_{\mathcal{L},\psi'}$ soit partout défini et soit
    une immersion.
\end{enumerate}

Lorsqu'il en est ainsi, tout sous-$\mathcal{O}_Y$-Module quasi-cohérent
$\mathcal{E}'$ de $\mathcal{S}_n$, contenant $\mathcal{E}$, possède la même
propriété, et il en est de même de l'image de $\bigotimes^k\mathcal{E}$ dans
$\mathcal{S}_{kn}$ pour tout $k>0$.
\end{proposition}

\begin{proof}
La suffisance de la condition et les deux dernières assertions sont des cas
particuliers de (\hyperref[II.3.8.3-fr]{3.8.3}), compte tenu de l'isomorphie
canonique entre $\operatorname{Proj}(\mathcal{S})$ et
$\operatorname{Proj}(\mathcal{S}^{(k)})$
(\hyperref[II.3.1.8-fr]{3.1.8}).

Soit $(U_i)$ un recouvrement ouvert fini de $Y$ par des ouverts affines et
posons $A_i=A(U_i)$. Comme $q^{-1}(U_i)$ est compact, l'hypothèse que
$r_{\mathcal{L},\psi}$ est une immersion partout définie et
(\hyperref[II.3.8.2-fr]{3.8.2}) entraînent l'existence d'une famille finie
$(s_{ij})$ d'éléments de
$S^{(i)}=\Gamma(U_i,\mathcal{S})$ ($s_{ij}\in S_{n_{ij}}^{(i)}$) ayant les
propriétés (i), (ii), (iii). On voit comme dans la démonstration de
(\hyperref[II.3.8.4-fr]{3.8.4}) que le morphisme $X_{f_{ij}}\to U_i$ (où
$f_{ij}=\psi^\flat(s_{ij})$) est de type fini, et l'anneau $B_{ij}$ de
$X_{f_{ij}}$ est par suite une $A_i$-algèbre de type fini, ayant un système de
générateurs de la forme
$(\psi^\flat(t_{ijk})|X_{f_{ij}})(f_{ij}|X_{f_{ij}})^{-m_{ijk}}$, où
$t_{ijk}\in S_{m_{ijk}n_{ij}}^{(i)}$. Soit $n$ un multiple commun des
$m_{ijk}n_{ij}$~; remplaçant (pour chaque $(i,j,k)$) $s_{ij}$ par une
puissance $s_{ij}^\rho$ telle que $\rho m_{ijk}n_{ij}=n$, et multipliant
$t_{ijk}$ par $s_{ij}^{\rho-m_{ijk}}$, on peut supposer que pour chaque $i$
tous les $s_{ij}$ et $t_{ijk}$ appartiennent à $S_n^{(i)}$ et que
$m_{ijk}=1$. Soit $E_i$ le sous-$A_i$-Module de $S_n^{(i)}$ engendré par ces
éléments (pour $i$ fixé). Il existe un sous-$\mathcal{O}_Y$-Module cohérent de
type

\oldpage[II]{71}

fini $\mathcal{E}_i$ de $\mathcal{S}_n$ tel que
$\mathcal{E}_i|U_i=(E_i)^{\sim}$
(\hyperref[I.9.4.7-fr]{I, 9.4.7}). Il est clair que le
sous-$\mathcal{O}_Y$-Module $\mathcal{E}$ de $\mathcal{S}_n$, somme des
$\mathcal{E}_i$, répond à la question (chaque section $f_{ij}$ étant telle
que pour tout $x\in X_{f_{ij}}$, il y a un voisinage affine
$V\subset X_{f_{ij}}$ de $x$ tel que $f|V$ soit une base de
$\Gamma(V,\mathcal{L}^{\otimes n})$).
\end{proof}

\section{Fibrés projectifs. Faisceaux amples}
\label{section:II.4-fr}

\subsection{Définition des fibrés projectifs.}
\label{subsection:II.4.1-fr}

\begin{definition}[4.1.1]
\label{II.4.1.1-fr}
Soient $Y$ un préschéma, $\mathcal{E}$ un $\mathcal{O}_Y$-Module
quasi-cohérent, $\mathbf{S}_{\mathcal{O}_Y}(\mathcal{E})$ la
$\mathcal{O}_Y$-Algèbre symétrique de $\mathcal{E}$
(\hyperref[II.1.7.4-fr]{1.7.4}) qui est quasi-cohérente
(\hyperref[II.1.7.7-fr]{1.7.7}). On appelle \emph{fibré projectif sur $Y$
défini par $\mathcal{E}$} et on note $\mathbf{P}(\mathcal{E})$ le $Y$-schéma
$P=\operatorname{Proj}(\mathbf{S}_{\mathcal{O}_Y}(\mathcal{E}))$. Le
$\mathcal{O}_P$-Module $\mathcal{O}_P(1)$ est appelé le \emph{faisceau
fondamental sur $P$}.
\end{definition}

Lorsque $Y$ est affine d'anneau $A$, on a $\mathcal{E}=\widetilde{E}$, où
$E$ est un $A$-module, et on écrit alors $\mathbf{P}(E)$ au lieu de
$\mathbf{P}(\widetilde{E})$.

Lorsque l'on prend $\mathcal{E}=\mathcal{O}_Y^n$, on écrit
$\mathbf{P}_Y^{n-1}$ au lieu de $\mathbf{P}(\mathcal{E})$~; si de plus $Y$
est affine d'anneau $A$, on écrit aussi $\mathbf{P}_A^{n-1}$ au lieu de
$\mathbf{P}_Y^{n-1}$. Comme
$\mathbf{S}_{\mathcal{O}_Y}(\mathcal{O}_Y)$ s'identifie canoniquement à
$\mathcal{O}_Y[T]$ (\hyperref[II.1.7.4-fr]{1.7.4}),
$\mathbf{P}_Y^0$ s'identifie canoniquement à $Y$
(\hyperref[II.3.1.7-fr]{3.1.7})~; l'exemple
(\hyperref[II.2.4.3-fr]{2.4.3}) n'est autre que $\mathbf{P}_K^1$.

\begin{env}[4.1.2]
\label{II.4.1.2-fr}
Soient $\mathcal{E}$, $\mathcal{F}$ deux $\mathcal{O}_Y$-Modules
quasi-cohérents~; $u:\mathcal{E}\to\mathcal{F}$ un
$\mathcal{O}_Y$-homomorphisme~; il lui correspond canoniquement un
homomorphisme
$\mathbf{S}(u):\mathbf{S}_{\mathcal{O}_Y}(\mathcal{E})\to
\mathbf{S}_{\mathcal{O}_Y}(\mathcal{F})$ de $\mathcal{O}_Y$-Algèbres
graduées (\hyperref[II.1.7.4-fr]{1.7.4}). Si $u$ est \emph{surjectif}, il en
est de même de $\mathbf{S}(u)$, et par suite
(\hyperref[II.3.6.2-fr]{3.6.2, (i)}),
$\operatorname{Proj}(\mathbf{S}(u))$ est une immersion fermée
$\mathbf{P}(\mathcal{F})\to\mathbf{P}(\mathcal{E})$, que nous noterons
$\mathbf{P}(u)$. On peut donc dire que $\mathbf{P}(\mathcal{E})$ est un
\emph{foncteur contravariant en $\mathcal{E}$}, à condition de ne prendre
pour morphismes des $\mathcal{O}_Y$-Modules quasi-cohérents que les
homomorphismes surjectifs.

En supposant toujours $u$ surjectif et posant
$P=\mathbf{P}(\mathcal{E})$, $Q=\mathbf{P}(\mathcal{F})$ et
$j=\mathbf{P}(u)$, on a, à un isomorphisme près,
\[
\label{II.4.1.2.1-fr}
  j^*(\mathcal{O}_P(n))=\mathcal{O}_Q(n)
  \qquad\text{pour tout }n\in\mathbf{Z}
\tag{4.1.2.1}
\]
comme il résulte de (\hyperref[II.3.6.3-fr]{3.6.3}).
\end{env}

\begin{env}[4.1.3]
\label{II.4.1.3-fr}
Soit maintenant $\psi:Y'\to Y$ un morphisme, et soit
$\mathcal{E}'=\psi^*(\mathcal{E})$~; on a alors
$\mathbf{S}_{\mathcal{O}_{Y'}}(\mathcal{E}')=
\psi^*(\mathbf{S}_{\mathcal{O}_Y}(\mathcal{E}))$
(\hyperref[II.1.7.5-fr]{1.7.5})~; par suite
(\hyperref[II.3.5.3-fr]{3.5.3}), on a
\[
\label{II.4.1.3.1-fr}
  \mathbf{P}(\psi^*(\mathcal{E}))=
  \mathbf{P}(\mathcal{E})\times_YY'
\tag{4.1.3.1}
\]
à un isomorphisme canonique près~; en outre, on a évidemment
\[
  \psi^*((\mathbf{S}_{\mathcal{O}_Y}(\mathcal{E}))(n))
  =(\mathbf{S}_{\mathcal{O}_{Y'}}(\mathcal{E}'))(n)
\]
pour tout $n\in\mathbf{Z}$, donc, en posant $P=\mathbf{P}(\mathcal{E})$,
$P'=\mathbf{P}(\mathcal{E}')$, on a
(\hyperref[II.3.5.4-fr]{3.5.4}), à un isomorphisme près,
\[
\label{II.4.1.3.2-fr}
  \mathcal{O}_{P'}(n)=\mathcal{O}_P(n)\otimes_Y\mathcal{O}_{Y'}
  \qquad\text{pour tout }n\in\mathbf{Z}.
\tag{4.1.3.2}
\]
\end{env}

\oldpage[II]{72}

\begin{proposition}[4.1.4]
\label{II.4.1.4-fr}
Soit $\mathcal{L}$ un $\mathcal{O}_Y$-Module inversible. Pour tout
$\mathcal{O}_Y$-Module quasi-cohérent $\mathcal{E}$, il existe un
$Y$-isomorphisme canonique
$i_{\mathcal{L}}:\mathbf{P}(\mathcal{E})\xrightarrow{\sim}
\mathbf{P}(\mathcal{E}\otimes\mathcal{L})$~; en outre, si l'on pose
$P=\mathbf{P}(\mathcal{E})$, $Q=\mathbf{P}(\mathcal{E}\otimes\mathcal{L})$,
$i_{\mathcal{L}}^*(\mathcal{O}_Q(n))$ est canoniquement isomorphe à
$\mathcal{O}_P(n)\otimes_Y\mathcal{L}^{\otimes n}$ pour tout
$n\in\mathbf{Z}$.
\end{proposition}

\begin{proof}
Remarquons d'abord que si $A$ est un anneau, $E$ un $A$-module, $L$ un
$A$-module \emph{monogène libre}, on définit canoniquement un homomorphisme
de $A$-modules
\[
  \mathbf{S}_n(E\otimes L)\longrightarrow
  \mathbf{S}_n(E)\otimes L^{\otimes n}
\]
en faisant correspondre à
$(x_1\otimes y_1)\cdots(x_n\otimes y_n)$ l'élément
\[
  (x_1x_2\cdots x_n)\otimes(y_1\otimes y_2\otimes\cdots\otimes y_n)
  \qquad(x_i\in E,\ y_i\in L\text{ pour }1\leq i\leq n)~;
\]
on vérifie immédiatement (en se ramenant au cas où $L=A$) que cet
homomorphisme est en fait un isomorphisme. On en conclut un isomorphisme
canonique de $A$-algèbres graduées
$\mathbf{S}_A(E\otimes L)\xrightarrow{\sim}
\bigoplus_{n\geq0}\mathbf{S}_n(E)\otimes L^{\otimes n}$. En revenant aux
conditions de (4.1.4), les remarques précédentes permettent de définir un
isomorphisme canonique de $\mathcal{O}_Y$-Algèbres graduées
\[
\label{II.4.1.4.1-fr}
  \mathbf{S}_{\mathcal{O}_Y}(\mathcal{E}\otimes_{\mathcal{O}_Y}\mathcal{L})
  \xrightarrow{\sim}
  \bigoplus_{n\geq0}\mathbf{S}_n(\mathcal{E})
  \otimes_{\mathcal{O}_Y}\mathcal{L}^{\otimes n}
\tag{4.1.4.1}
\]
en définissant cet isomorphisme comme un isomorphisme de préfaisceaux et
tenant compte de (\hyperref[II.1.7.4-fr]{1.7.4}),
(\hyperref[I.1.3.9-fr]{I, 1.3.9}) et
(\hyperref[I.1.3.12-fr]{I, 1.3.12}). La proposition résulte alors de
(\hyperref[II.3.1.8-fr]{3.1.8, (iii)}) et
(\hyperref[II.3.2.10-fr]{3.2.10}).
\end{proof}

\begin{env}[4.1.5]
\label{II.4.1.5-fr}
Avec les hypothèses de (\hyperref[II.4.1.1-fr]{4.1.1}), posons
$P=\mathbf{P}(\mathcal{E})$, et désignons par $p$ le morphisme structural
$P\to Y$. Comme on a par définition
$\mathcal{E}=(\mathbf{S}_{\mathcal{O}_Y}(\mathcal{E}))_1$, on a un
homomorphisme canonique
$\alpha_1:\mathcal{E}\to p_*(\mathcal{O}_P(1))$
(\hyperref[II.3.3.2.2-fr]{3.3.2.2}) et par suite aussi
(\hyperref[0.4.4.3-fr]{0, 4.4.3}) un homomorphisme canonique
\[
\label{II.4.1.5.1-fr}
  \alpha_1^\sharp:p^*(\mathcal{E})\longrightarrow\mathcal{O}_P(1).
\tag{4.1.5.1}
\]
\end{env}

\begin{proposition}[4.1.6]
\label{II.4.1.6-fr}
L'homomorphisme canonique (\hyperref[II.4.1.5.1-fr]{4.1.5.1}) est surjectif.
\end{proposition}

\begin{proof}
En effet, on a vu dans (\hyperref[II.3.3.2-fr]{3.3.2}) que
$\alpha_1^\sharp$ correspond fonctoriellement à l'homomorphisme canonique
\[
  \mathcal{E}\otimes_{\mathcal{O}_Y}
  \mathbf{S}_{\mathcal{O}_Y}(\mathcal{E})
  \longrightarrow
  (\mathbf{S}_{\mathcal{O}_Y}(\mathcal{E}))(1)~;
\]
comme par définition $\mathcal{E}$ engendre
$\mathbf{S}_{\mathcal{O}_Y}(\mathcal{E})$, cet homomorphisme est surjectif,
d'où la conclusion en vertu de (\hyperref[II.3.2.4-fr]{3.2.4}).
\end{proof}

\subsection{Morphismes d'un préschéma dans un fibré projectif.}
\label{subsection:II.4.2-fr}

\begin{env}[4.2.1]
\label{II.4.2.1-fr}
Gardant les notations de (\hyperref[II.4.1.5-fr]{4.1.5}), soient $X$ un
$Y$-préschéma, $q:X\to Y$ le morphisme structural, et soit $r:X\to P$ un
$Y$-\emph{morphisme}, de sorte que l'on a le diagramme commutatif
\[
  \xymatrix{
    P \ar[d]_p & X \ar[l]_r \ar[dl]^q
    \\Y
  }
\]

\oldpage[II]{73}

Comme le foncteur $r^*$ est exact à droite
(\hyperref[0.4.3.1-fr]{0, 4.3.1}), on déduit de l'homomorphisme
(\hyperref[II.4.1.5.1-fr]{4.1.5.1}), qui est surjectif
(\hyperref[II.4.1.6-fr]{4.1.6}), un homomorphisme surjectif
\[
  r^*(\alpha_1^\sharp):r^*(p^*(\mathcal{E}))\longrightarrow
  r^*(\mathcal{O}_P(1)).
\]
Mais $r^*(p^*(\mathcal{E}))=q^*(\mathcal{E})$, et
$r^*(\mathcal{O}_P(1))$ est localement isomorphe à
$r^*(\mathcal{O}_P)=\mathcal{O}_X$, autrement dit c'est un faisceau
inversible $\mathcal{L}_r$ sur $\mathcal{O}_X$ et on a ainsi défini à partir
de $r$ un $\mathcal{O}_X$-homomorphisme canonique surjectif
\[
\label{II.4.2.1.1-fr}
  \varphi_r:q^*(\mathcal{E})\longrightarrow\mathcal{L}_r.
\tag{4.2.1.1}
\]

Lorsque $Y=\operatorname{Spec}(A)$ est affine, et
$\mathcal{E}=\widetilde{E}$ on explicite encore cet homomorphisme de la façon
suivante~: étant donné $f\in E$, il résulte d'abord de
(\hyperref[II.2.6.3-fr]{2.6.3}) que l'on a
\[
\label{II.4.2.1.2-fr}
  r^{-1}(D_+(f))=X_{\varphi_r^\flat(f)}.
\tag{4.2.1.2}
\]
D'autre part, soit $V$ un ouvert affine de $X$ contenu dans
$r^{-1}(D_+(f))$, et soit $B$ son anneau, qui est une $A$-algèbre~; posons
$S=\mathbf{S}_A(E)$~; la restriction de $r$ à $V$ correspond à un
$A$-homomorphisme $\omega:S_{(f)}\to B$, on a
$q^*(\mathcal{E})|V=\widetilde{E\otimes_A B}$ et
$\mathcal{L}_r|V=\widetilde{L_r}$, où
$L_r=(S(1))_{(f)}\otimes_{S_{(f)}}B_{[\omega]}$
(\hyperref[I.1.6.5-fr]{I, 1.6.5}). La restriction de $\varphi_r$ à
$q^*(\mathcal{E})|V$ correspond au $B$-homomorphisme
$u:E\otimes_A B\to L_r$ qui transforme $x\otimes1$ en
$(x/1)\otimes f=(f/1)\otimes\omega(x/f)$. L'extension canonique de
$\varphi_r$ en un homomorphisme de $\mathcal{O}_X$-Algèbres
\[
  \psi_r:q^*(\mathbf{S}(\mathcal{E}))=
  \mathbf{S}(q^*(\mathcal{E}))\longrightarrow
  \mathbf{S}(\mathcal{L}_r)=\bigoplus_{n\geq0}\mathcal{L}_r^{\otimes n}
\]
est donc telle que la restriction de $\psi_r$ à
$q^*(\mathbf{S}_n(\mathcal{E}))|V$ correspond à l'homomorphisme
$\mathbf{S}_n(E\otimes_A B)=\mathbf{S}_n(E)\otimes_A B\to L_r^{\otimes n}$
qui transforme $s\otimes1$ en
$(f/1)^{\otimes n}\otimes\omega(s/f^n)$.
\end{env}

\begin{env}[4.2.2]
\label{II.4.2.2-fr}
Inversement, supposons donnés un morphisme $q:X\to Y$, un
$\mathcal{O}_X$-Module inversible $\mathcal{L}$ et un
$\mathcal{O}_Y$-Module quasi-cohérent $\mathcal{E}$~; à un homomorphisme
$\varphi:q^*(\mathcal{E})\to\mathcal{L}$ correspond canoniquement un
homomorphisme de $\mathcal{O}_X$-Algèbres quasi-cohérentes
\[
  \psi:\mathbf{S}(q^*(\mathcal{E}))=q^*(\mathbf{S}(\mathcal{E}))
  \longrightarrow\bigoplus_{n\geq0}\mathcal{L}^{\otimes n},
\]
et par suite (\hyperref[II.3.7.1-fr]{3.7.1}) un $Y$-morphisme
$r_{\mathcal{L},\psi}:G(\psi)\to
\operatorname{Proj}(\mathbf{S}(\mathcal{E}))=\mathbf{P}(\mathcal{E})$, que
nous noterons $r_{\mathcal{L},\varphi}$. Si $\varphi$ est surjectif, il en est
de même de $\psi$, donc (\hyperref[II.3.7.4-fr]{3.7.4})
$r_{\mathcal{L},\varphi}$ est partout défini. En outre, avec les notations de
(\hyperref[II.4.2.1-fr]{4.2.1}) et de (4.2.2)~:
\end{env}

\begin{proposition}[4.2.3]
\label{II.4.2.3-fr}
Étant donnés un morphisme $q:X\to Y$ et un $\mathcal{O}_Y$-Module
quasi-cohérent $\mathcal{E}$, les applications
$r\mapsto(\mathcal{L}_r,\varphi_r)$ et
$(\mathcal{L},\varphi)\mapsto r_{\mathcal{L},\varphi}$ mettent en
correspondance biunivoque l'ensemble des $Y$-morphismes
$r:X\to\mathbf{P}(\mathcal{E})$ et l'ensemble des classes d'équivalence des
couples $(\mathcal{L},\varphi)$ formés d'un $\mathcal{O}_X$-Module inversible
$\mathcal{L}$ et d'un homomorphisme surjectif
$\varphi:q^*(\mathcal{E})\to\mathcal{L}$, deux couples
$(\mathcal{L},\varphi)$, $(\mathcal{L}',\varphi')$ étant équivalents s'il
existe un $\mathcal{O}_X$-isomorphisme
$\tau:\mathcal{L}\xrightarrow{\sim}\mathcal{L}'$ tel que
$\varphi'=\tau\circ\varphi$.
\end{proposition}

\begin{proof}
Partons d'abord d'un $Y$-morphisme $r:X\to\mathbf{P}(\mathcal{E})$, formons
$\mathcal{L}_r$ et $\varphi_r$ (\hyperref[II.4.2.1-fr]{4.2.1}), puis
$r'=r_{\mathcal{L}_r,\varphi_r}$~; il résulte aussitôt de
(\hyperref[II.4.2.1-fr]{4.2.1}) et de
(\hyperref[II.3.7.2-fr]{3.7.2}) que les morphismes $r$ et $r'$ sont
identiques (en prenant pour générateur de $\mathcal{L}_r$ l'élément
$(f/1)\otimes1$ pour définir les homomorphismes $v_n$ de
(\hyperref[II.3.7.2-fr]{3.7.2})). Inversement, partons d'un couple
$(\mathcal{L},\varphi)$ et formons

\oldpage[II]{74}

$r''=r_{\mathcal{L},\varphi}$, puis $\mathcal{L}_{r''}$ et
$\varphi_{r''}$~; montrons qu'il y a un isomorphisme canonique
$\tau:\mathcal{L}_{r''}\xrightarrow{\sim}\mathcal{L}$ tel que
$\varphi=\tau\circ\varphi_{r''}$~; pour le définir on peut se placer dans le
cas où $Y=\operatorname{Spec}(A)$ et $X=\operatorname{Spec}(B)$ sont affines,
et (avec les notations de (\hyperref[II.4.2.1-fr]{4.2.1}) et de
(\hyperref[II.3.7.2-fr]{3.7.2})) faire correspondre à tout élément
$(x/1)\otimes1$ de $L_{r''}$ (avec $x\in E$) l'élément $v_1(x)c$ de $L$. On
vérifie aussitôt que $\tau$ ne dépend pas du générateur choisi $c$ de $L$~;
comme $v_1$ est par hypothèse surjectif, pour prouver que $\tau$ est un
isomorphisme, il suffit de voir que si $x/1=0$ dans $(S(1))_{(f)}$, on a
$v_1(x)/1=0$ dans $B_g$~; mais la première relation signifie que $f^nx=0$
dans $\mathbf{S}_{n+1}(E)$ pour un certain $n$, et on en déduit que
$v_{n+1}(f^nx)=g^nv_1(x)=0$ dans $B$, d'où la conclusion. Enfin, il est
immédiat que pour deux couples équivalents $(\mathcal{L},\varphi)$,
$(\mathcal{L}',\varphi')$, on a
$r_{\mathcal{L},\varphi}=r_{\mathcal{L}',\varphi'}$.
\end{proof}

En particulier, pour $X=Y$~:

\begin{theorem}[4.2.4]
\label{II.4.2.4-fr}
L'ensemble des $Y$-sections de $\mathbf{P}(\mathcal{E})$ est en
correspondance biunivoque canonique avec l'ensemble des
sous-$\mathcal{O}_Y$-Modules quasi-cohérents $\mathcal{F}$ de $\mathcal{E}$
tels que $\mathcal{E}/\mathcal{F}$ soit inversible.
\end{theorem}

On notera que cette propriété correspond à la définition classique de
l'\og espace projectif\fg{} comme ensemble des hyperplans d'un espace
vectoriel (le cas classique correspondant à $Y=\operatorname{Spec}(K)$, où
$K$ est un corps, et $\mathcal{E}=\widetilde{E}$, $E$ étant un $K$-espace
vectoriel de dimension finie~; les $\mathcal{F}$ ayant la propriété énoncée
dans (\hyperref[II.4.2.4-fr]{4.2.4}) correspondent alors aux hyperplans de
$E$, et on sait d'autre part que les $Y$-sections de
$\mathbf{P}(\mathcal{E})$ sont alors les points de
$\mathbf{P}(\mathcal{E})$ rationnels sur $K$
(\hyperref[I.3.4.5-fr]{I, 3.4.5})).

\begin{remark}[4.2.5]
\label{II.4.2.5-fr}
Comme il y a correspondance biunivoque canonique entre $Y$-morphismes de $X$
dans $P$ et leurs morphismes graphes, $X$-sections de $P\times_YX$
(\hyperref[I.3.3.14-fr]{I, 3.3.14}), on voit qu'inversement
(\hyperref[II.4.2.3-fr]{4.2.3}) peut se déduire de
(\hyperref[II.4.2.4-fr]{4.2.4}). Désignons par
$\operatorname{Hyp}_Y(X,\mathcal{E})$ l'ensemble des
sous-$\mathcal{O}_X$-Modules quasi-cohérents $\mathcal{F}$ de
$\mathcal{E}\otimes_Y\mathcal{O}_X=q^*(\mathcal{E})$ tels que
$q^*(\mathcal{E})/\mathcal{F}$ soit un $\mathcal{O}_X$-Module inversible. Si
$g:X'\to X$ est un $Y$-morphisme, il résulte du fait que $g^*$ est exact à
droite que l'on a
\[
  g^*(q^*(\mathcal{E})/\mathcal{F})=
  g^*(q^*(\mathcal{E}))/g^*(\mathcal{F}),
\]
donc que ce dernier faisceau est inversible, et par suite
$\operatorname{Hyp}_Y(X,\mathcal{E})$ est un foncteur contravariant dans la
catégorie des $Y$-préschémas. On peut alors interpréter le th.
(\hyperref[II.4.2.4-fr]{4.2.4}) comme définissant un isomorphisme canonique
des foncteurs $\operatorname{Hom}_Y(X,\mathbf{P}(\mathcal{E}))$ et
$\operatorname{Hyp}_Y(X,\mathcal{E})$ contravariants en $X$ variable dans la
catégorie des $Y$-préschémas. Ceci fournit aussi une caractérisation du fibré
projectif $P=\mathbf{P}(\mathcal{E})$ par la propriété universelle suivante,
plus proche de l'intuition géométrique que les constructions des \S\S~2 et
3~: pour tout morphisme $q:X\to Y$ et tout $\mathcal{O}_X$-Module inversible
$\mathcal{L}$, quotient de $\mathcal{E}\otimes_{\mathcal{O}_Y}\mathcal{O}_X$,
il existe un $Y$-morphisme unique $r:X\to P$ tel que
$\mathcal{L}=r^*(\mathcal{O}_P(1))$.

On verra plus tard comment on peut de la même manière définir, entre autres,
les schémas \og grassmanniennes\fg{}.
\end{remark}

\begin{corollary}[4.2.6]
\label{II.4.2.6-fr}
Supposons que tout $\mathcal{O}_Y$-Module inversible soit trivial
(\hyperref[I.2.4.8-fr]{I, 2.4.8}). Soit $V$ le groupe
$\operatorname{Hom}_{\mathcal{O}_Y}(\mathcal{E},\mathcal{O}_Y)$, considéré
comme module sur l'anneau $A=\Gamma(Y,\mathcal{O}_Y)$, et soit $V^*$ le
sous-ensemble de $V$ formé des homomorphismes surjectifs. Alors l'ensemble
des $Y$-sections de $\mathbf{P}(\mathcal{E})$ s'identifie canoniquement à
$V^*/A^*$, où $A^*$ est le groupe des unités de $A$.
\end{corollary}

\oldpage[II]{75}

En particulier~:
\begin{enumerate}
  \item[1\textsuperscript{o}] Le corollaire
  (\hyperref[II.4.2.6-fr]{4.2.6}) s'applique lorsque $Y$ est un
  \emph{schéma local} (\hyperref[I.2.4.8-fr]{I, 2.4.8}). Soient $Y$ un
  préschéma quelconque, $y$ un point de $Y$,
  $Y'=\operatorname{Spec}(k(y))$~; la fibre $p^{-1}(y)$ de
  $\mathbf{P}(\mathcal{E})$ est, en vertu de
  (\hyperref[II.4.1.3.1-fr]{4.1.3.1}), identifiée à
  $\mathbf{P}(\mathcal{E}^y)$, où
  $\mathcal{E}^y=\mathcal{E}_y\otimes_{\mathcal{O}_y}k(y)
  =\mathcal{E}_y/\mathfrak{m}_y\mathcal{E}_y$ est considéré comme espace
  vectoriel sur $k(y)$. Plus généralement, si $K$ est une extension de
  $k(y)$, $p^{-1}(y)\otimes_{k(y)}K$ s'identifie à
  $\mathbf{P}(\mathcal{E}^y\otimes_{k(y)}K)$. Le cor.
  (\hyperref[II.4.2.6-fr]{4.2.6}) montre donc que l'ensemble des
  \emph{points géométriques de $\mathbf{P}(\mathcal{E})$ à valeurs dans
  l'extension $K$ de $k(y)$} (\hyperref[I.3.4.5-fr]{I, 3.4.5}), que l'on
  peut encore appeler \emph{fibre géométrique rationnelle sur $K$ de
  $\mathbf{P}(\mathcal{E})$ au-dessus du point $y$}, s'identifie à
  l'\emph{espace projectif} associé au \emph{dual} du $K$-espace vectoriel
  $\mathcal{E}^y\otimes_{k(y)}K$.
  \item[2\textsuperscript{o}] Supposons que $Y$ soit affine d'anneau $A$, et
  en outre que tout $\mathcal{O}_Y$-Module inversible soit trivial~; prenons
  en outre $\mathcal{E}=\mathcal{O}_Y^n$~; alors, dans
  (\hyperref[II.4.2.6-fr]{4.2.6}), $V$ s'identifie à $A^n$
  (\hyperref[I.1.3.8-fr]{I, 1.3.8}), $V^*$ à l'ensemble des systèmes
  $(f_i)_{1\leq i\leq n}$ d'éléments de $A$ engendrant l'idéal $A$~; deux
  tels systèmes définissent la même $Y$-section de
  $\mathbf{P}_Y^{n-1}=\mathbf{P}_A^{n-1}$, autrement dit le même point de
  $\mathbf{P}_A^{n-1}$ à valeurs dans $A$ si et seulement si l'un se déduit
  de l'autre par multiplication par un élément inversible de $A$.
\end{enumerate}

Ces propriétés justifient la terminologie de \og fibré projectif\fg{} pour
$\mathbf{P}(\mathcal{E})$. On notera que la définition que nous obtenons
ainsi de l'\og espace projectif\fg{} est en fait duale de la définition
classique~; cela nous est imposé par la nécessité de pouvoir définir
$\mathbf{P}(\mathcal{E})$ pour un $\mathcal{O}_Y$-Module quasi-cohérent
quelconque $\mathcal{E}$, non nécessairement localement libre.

\begin{remark}[4.2.7]
\label{II.4.2.7-fr}
Nous verrons au chapitre~V que si $Y$ est localement noethérien et connexe,
et si $\mathcal{E}$ est localement libre, alors, en posant
$P=\mathbf{P}(\mathcal{E})$, tout $\mathcal{O}_P$-Module inversible
$\mathcal{L}$ est isomorphe à un $\mathcal{O}_P$-Module de la forme
$\mathcal{L}'\otimes_{\mathcal{O}_Y}\mathcal{O}_P(m)$, où $\mathcal{L}'$ est
un $\mathcal{O}_Y$-Module inversible, bien déterminé à un isomorphisme près,
et $m$ un entier bien déterminé. En d'autres termes,
$\mathrm{H}^1(P,\mathcal{O}_P^*)$ est isomorphe à
$\mathbf{Z}\times\mathrm{H}^1(Y,\mathcal{O}_Y^*)$
(\hyperref[0.5.4.7-fr]{0, 5.4.7}). Nous verrons aussi
(\hyperref[III.2.1.14-fr]{III, 2.1.14}, compte tenu de
(\hyperref[0.5.4.10-fr]{0, 5.4.10})) que
$p_*(\mathcal{L}^{\otimes m})=0$ si $m<0$ et que
$p_*(\mathcal{L}^{\otimes m})$ est isomorphe à
$\mathcal{L}'\otimes_{\mathcal{O}_Y}
(\mathbf{S}_{\mathcal{O}_Y}(\mathcal{E}))_m$ si $m\geq0$. Si $\mathcal{F}$
est un $\mathcal{O}_Y$-Module quasi-cohérent, tout $Y$-morphisme
$\mathbf{P}(\mathcal{E})\to\mathbf{P}(\mathcal{F})$ est donc déterminé par
la donnée d'un $\mathcal{O}_Y$-Module inversible $\mathcal{L}'$, d'un entier
$m\geq0$ et d'un $\mathcal{O}_Y$-homomorphisme
$\psi:\mathcal{F}\to\mathcal{L}'\otimes_{\mathcal{O}_Y}
(\mathbf{S}_{\mathcal{O}_Y}(\mathcal{E}))_m$ tel que l'homomorphisme
correspondant $\psi^\sharp$ de $\mathcal{O}_{\mathbf{P}(\mathcal{F})}$-Modules
soit surjectif. Nous verrons aussi que si le $Y$-morphisme considéré est un
isomorphisme, alors $m=1$ et $\mathcal{F}$ est isomorphe à
$\mathcal{E}\otimes_{\mathcal{O}_Y}\mathcal{L}'$ (réciproque de
(\hyperref[II.4.1.4-fr]{4.1.4})). Ceci permettra de déterminer le faisceau
des germes d'automorphismes de $\mathbf{P}(\mathcal{E})$ comme quotient du
faisceau de groupes $\mathcal{A}ut(\mathcal{E})$ (localement isomorphe à
$\operatorname{GL}(n,\mathcal{O}_Y)$ si $\mathcal{E}$ est de rang $n$) par
$\mathcal{O}_Y^*$.
\end{remark}

\begin{env}[4.2.8]
\label{II.4.2.8-fr}
Conservant les notations de (\hyperref[II.4.2.1-fr]{4.2.1}), soit
$u:X'\to X$ un morphisme~; si le $Y$-morphisme $r:X\to P$ correspond à
l'homomorphisme $\varphi:q^*(\mathcal{E})\to\mathcal{L}$, le $Y$-morphisme
$r\circ u$ correspond à
$u^*(\varphi):u^*(q^*(\mathcal{E}))\to u^*(\mathcal{L})$, comme il résulte
aussitôt des définitions.
\end{env}

\begin{env}[4.2.9]
\label{II.4.2.9-fr}
Soient $\mathcal{E}$, $\mathcal{F}$ deux $\mathcal{O}_Y$-Modules
quasi-cohérents, $v:\mathcal{E}\to\mathcal{F}$ un homomorphisme surjectif,
$j=\mathbf{P}(v)$ l'immersion fermée correspondante
$\mathbf{P}(\mathcal{F})\to\mathbf{P}(\mathcal{E})$
(\hyperref[II.4.1.2-fr]{4.1.2}). Si le $Y$-morphisme
$r:X\to\mathbf{P}(\mathcal{F})$ correspond à l'homomorphisme
$\varphi:q^*(\mathcal{F})\to\mathcal{L}$, le

\oldpage[II]{76}

$Y$-morphisme $j\circ r$ correspond à
\[
  q^*(\mathcal{E})\xrightarrow{q^*(v)}q^*(\mathcal{F})
  \xrightarrow{\varphi}\mathcal{L}~;
\]
cela résulte encore de la définition donnée en
(\hyperref[II.4.2.1-fr]{4.2.1}).
\end{env}

\begin{env}[4.2.10]
\label{II.4.2.10-fr}
Soit $\psi:Y'\to Y$ un morphisme, et posons
$\mathcal{E}'=\psi^*(\mathcal{E})$. Si le $Y$-morphisme $r:X\to P$ correspond
à l'homomorphisme $\varphi:q^*(\mathcal{E})\to\mathcal{L}$, le
$Y'$-morphisme
\[
  r_{(Y')}:X_{(Y')}\longrightarrow P'=\mathbf{P}(\mathcal{E}')
\]
correspond à
$\varphi_{(Y')}:q_{(Y')}^*(\mathcal{E}')=
  q^*(\mathcal{E})\otimes_Y\mathcal{O}_{Y'}\to
\mathcal{L}\otimes_Y\mathcal{O}_{Y'}$. En effet, en vertu de
(\hyperref[II.4.1.3.1-fr]{4.1.3.1}), on a le diagramme commutatif
\[
  \xymatrix{
    Y' \ar[d]
    & P'=P_{(Y')} \ar[l]_{p_{(Y')}} \ar[d]^{u}
    & X_{(Y')} \ar[l]_{r_{(Y')}} \ar[d]^{v}
    \\Y
    & P \ar[l]_{p}
    & X \ar[l]_{r}
  }
\]
D'après (\hyperref[II.4.1.3.2-fr]{4.1.3.2}), on a
\[
  (r_{(Y')})^*(\mathcal{O}_{P'}(1))
  =(r_{(Y')})^*(u^*(\mathcal{O}_P(1)))
  =v^*(r^*(\mathcal{O}_P(1)))
  =v^*(\mathcal{L})=\mathcal{L}\otimes_Y\mathcal{O}_{Y'}~;
\]
d'autre part, $u^*(\alpha_1^\sharp)$ n'est autre que l'homomorphisme
canonique
$\alpha_1^\sharp:(p_{(Y')})^*(\mathcal{E}')\to\mathcal{O}_{P'}(1)$~; cela
se voit en explicitant les homomorphismes canoniques $\alpha_1^\sharp$ dans
$P$ et $P'$ comme dans (\hyperref[II.4.1.6-fr]{4.1.6}). D'où notre assertion.
\end{env}

\subsection{Le morphisme de Segre.}
\label{subsection:II.4.3-fr}

\begin{env}[4.3.1]
\label{II.4.3.1-fr}
Soient $Y$ un préschéma, $\mathcal{E}$, $\mathcal{F}$ deux
$\mathcal{O}_Y$-Modules quasi-cohérents~; posons
$P_1=\mathbf{P}(\mathcal{E})$, $P_2=\mathbf{P}(\mathcal{F})$, et désignons
par $p_1:P_1\to Y$, $p_2:P_2\to Y$ les morphismes structuraux. Soit
$Q=P_1\times_YP_2$, et soient $q_1:Q\to P_1$, $q_2:Q\to P_2$ les
projections canoniques~; le $\mathcal{O}_Q$-Module
\[
  \mathcal{L}=\mathcal{O}_{P_1}(1)\otimes_Y\mathcal{O}_{P_2}(1)
  :=q_1^*(\mathcal{O}_{P_1}(1))\otimes_{\mathcal{O}_Q}
  q_2^*(\mathcal{O}_{P_2}(1))
\]
est inversible comme produit tensoriel de deux $\mathcal{O}_Q$-Modules
inversibles (\hyperref[0.5.4.4-fr]{0, 5.4.4}). D'autre part, si
$r=p_1\circ q_1=p_2\circ q_2$ est le morphisme structural $Q\to Y$, on a
\[
  r^*(\mathcal{E}\otimes_{\mathcal{O}_Y}\mathcal{F})
  =q_1^*(p_1^*(\mathcal{E}))\otimes_{\mathcal{O}_Q}
  q_2^*(p_2^*(\mathcal{F}))
\]
(\hyperref[0.4.3.3-fr]{0, 4.3.3})~; les homomorphismes canoniques surjectifs
(\hyperref[II.4.1.5.1-fr]{4.1.5.1})
$p_1^*(\mathcal{E})\to\mathcal{O}_{P_1}(1)$,
$p_2^*(\mathcal{F})\to\mathcal{O}_{P_2}(1)$ donnent donc, par produit
tensoriel, un homomorphisme canonique
\[
\label{II.4.3.1.1-fr}
  s:r^*(\mathcal{E}\otimes_{\mathcal{O}_Y}\mathcal{F})\longrightarrow
  \mathcal{L}
\tag{4.3.1.1}
\]
évidemment surjectif~; on en déduit donc
(\hyperref[II.4.2.2-fr]{4.2.2}) un morphisme canonique, dit
\emph{morphisme de Segre}~:
\[
\label{II.4.3.1.2-fr}
  \varsigma:\mathbf{P}(\mathcal{E})\times_Y\mathbf{P}(\mathcal{F})
  \longrightarrow
  \mathbf{P}(\mathcal{E}\otimes_{\mathcal{O}_Y}\mathcal{F}).
\tag{4.3.1.2}
\]
Explicitons le morphisme $\varsigma$, lorsque $Y=\operatorname{Spec}(A)$ est
affine, $\mathcal{E}=\widetilde{E}$, $\mathcal{F}=\widetilde{F}$, $E$ et $F$
étant deux $A$-modules, d'où
$\mathcal{E}\otimes_{\mathcal{O}_Y}\mathcal{F}=\widetilde{E\otimes_A F}$
(\hyperref[I.1.3.12-fr]{I, 1.3.12})~; posons
$R=\mathbf{S}_A(E)$, $S=\mathbf{S}_A(F)$,
$T=\mathbf{S}_A(E\otimes_A F)$~; soient $f\in E$, $g\in F$, et considérons
l'ouvert affine
\[
  D_+(f)\times_YD_+(g)=\operatorname{Spec}(B)
\]

\oldpage[II]{77}

de $Q$, où $B=R_{(f)}\otimes_AS_{(g)}$~; la restriction de $\mathcal{L}$ à
cet ouvert affine est $\widetilde{L}$, où
\[
  L=(R(1))_{(f)}\otimes_A(S(1))_{(g)}
\]
et l'élément $c=(f/1)\otimes(g/1)$ est un générateur de $L$ considéré comme
$B$-module libre (\hyperref[II.2.5.7-fr]{2.5.7}). L'homomorphisme
(\hyperref[II.4.3.1.1-fr]{4.3.1.1}) correspond à l'homomorphisme
\[
  (x\otimes y)\otimes b\longmapsto b((x/1)\otimes(y/1))
\]
de $(E\otimes_A F)\otimes_AB$ dans $L$. Avec les notations de
(\hyperref[II.3.7.2-fr]{3.7.2}), on a donc
$v_1(x\otimes y)=(x/f)\otimes(y/g)$~; la restriction de $\varsigma$ à
$D_+(f)\times_YD_+(g)$ est un morphisme de ce schéma affine dans
$D_+(f\otimes g)$, correspondant à l'homomorphisme d'anneaux
$\omega:T_{(f\otimes g)}\to R_{(f)}\otimes_AS_{(g)}$ tel que
\[
\label{II.4.3.1.3-fr}
  \omega((x\otimes y)/(f\otimes g))=(x/f)\otimes(y/g)
\tag{4.3.1.3}
\]
pour $x\in E$ et $y\in F$.
\end{env}

\begin{env}[4.3.2]
\label{II.4.3.2-fr}
Il résulte de (\hyperref[II.4.2.3-fr]{4.2.3}) que l'on a un isomorphisme
canonique
\[
\label{II.4.3.2.1-fr}
  \tau:\varsigma^*(\mathcal{O}_P(1))\xrightarrow{\sim}
  \mathcal{O}_{P_1}(1)\otimes_Y\mathcal{O}_{P_2}(1)
\tag{4.3.2.1}
\]
où on a posé
$P=\mathbf{P}(\mathcal{E}\otimes_{\mathcal{O}_Y}\mathcal{F})$. Montrons
que, pour $x\in\Gamma(Y,\mathcal{E})$ et
$y\in\Gamma(Y,\mathcal{F})$, on a
\[
\label{II.4.3.2.2-fr}
  \tau(\alpha_1(x\otimes y))=\alpha_1(x)\otimes\alpha_1(y).
\tag{4.3.2.2}
\]
En effet, on est ramené au cas où $Y$ est affine, et on a alors, avec les
notations de (\hyperref[II.4.3.1-fr]{4.3.1}) et de
(\hyperref[II.2.6.2-fr]{2.6.2}),
$\alpha_1^{f\otimes g}(x\otimes y)=(x\otimes y)/1$ dans
$(T(1))_{(f\otimes g)}$, $\alpha_1^f(x)=x/1$ dans $(R(1))_{(f)}$ et
$\alpha_1^g(y)=y/1$ dans $(S(1))_{(g)}$. La définition de $\tau$ donnée dans
(\hyperref[II.4.2.3-fr]{4.2.3}) et le calcul de $v_1$ fait dans
(\hyperref[II.4.3.1-fr]{4.3.1}) prouvent alors aussitôt
(\hyperref[II.4.3.2.2-fr]{4.3.2.2}). On en déduit la formule
\[
\label{II.4.3.2.3-fr}
  \varsigma^{-1}(P_{x\otimes y})=(P_1)_x\times_Y(P_2)_y
\tag{4.3.2.3}
\]
avec les notations de (\hyperref[II.3.1.4-fr]{3.1.4}). En effet, compte tenu
de (\hyperref[II.3.3.3-fr]{3.3.3}), la formule
(\hyperref[II.4.3.2.2-fr]{4.3.2.2}) ramène (en revenant au cas affine, à
l'aide de (\hyperref[I.3.2.7-fr]{I, 3.2.7} et
(\hyperref[I.3.2.3-fr]{3.2.3})) à prouver le lemme suivant~:

\begin{lemma}[4.3.2.4]
\label{II.4.3.2.4-fr}
Soient $B$, $B'$ deux $A$-algèbres, et soient
$Y=\operatorname{Spec}(A)$, $Z=\operatorname{Spec}(B)$,
$Z'=\operatorname{Spec}(B')$~; quels que soient $t\in B$, $t'\in B'$, on a
$D(t\otimes t')=D(t)\times_YD(t')$.
\end{lemma}

\begin{proof}
En effet, si $p:Z\times_YZ'\to Z$ et $p':Z\times_YZ'\to Z'$ sont les
projections canoniques, il résulte de
(\hyperref[I.1.2.2.2-fr]{I, 1.2.2.2}) que l'on a
$p^{-1}(D(t))=D(t\otimes1)$ et
$p'^{-1}(D(t'))=D(1\otimes t')$~; la conclusion résulte de
(\hyperref[I.3.2.7-fr]{I, 3.2.7}) et
(\hyperref[I.1.1.9.1-fr]{I, 1.1.9.1}), puisque
$(t\otimes1)(1\otimes t')=t\otimes t'$.
\end{proof}
\end{env}

\begin{proposition}[4.3.3]
\label{II.4.3.3-fr}
Le morphisme de Segre est une immersion fermée.
\end{proposition}

\begin{proof}
En effet, la question étant locale sur $Y$, on est ramené au cas où $Y$ est
affine. Avec les notations de (\hyperref[II.4.3.1-fr]{4.3.1}) et
(\hyperref[II.4.3.2-fr]{4.3.2}), les $D_+(f\otimes g)$ forment alors une base
de la topologie de $P$, puisque les $f\otimes g$ engendrent $T$ lorsque $f$
parcourt $E$ et $g$ parcourt $F$. D'autre part, on a en vertu de
(\hyperref[II.4.3.2.3-fr]{4.3.2.3})
$\varsigma^{-1}(D_+(f\otimes g))=D_+(f)\times_YD_+(g)$. Il suffit donc
(\hyperref[I.4.2.4-fr]{I, 4.2.4}) de prouver que la restriction de
$\varsigma$ à $D_+(f)\times_YD_+(g)$ est une immersion fermée dans
$D_+(f\otimes g)$. Or, avec les mêmes notations, la formule
(\hyperref[II.4.3.1.3-fr]{4.3.1.3}) montre que $\omega$ est surjectif, ce qui
termine la démonstration.
\end{proof}

\begin{env}[4.3.4]
\label{II.4.3.4-fr}
Le morphisme de Segre est fonctoriel en $\mathcal{E}$ et $\mathcal{F}$,
lorsqu'on restreint les

\oldpage[II]{78}

homomorphismes de $\mathcal{O}_Y$-Modules quasi-cohérents aux homomorphismes
surjectifs. Il faut en effet montrer que si
$\mathcal{E}\to\mathcal{E}'$ est un $\mathcal{O}_Y$-homomorphisme surjectif,
le diagramme
\[
  \xymatrix{
    \mathbf{P}(\mathcal{E}')\times_Y\mathbf{P}(\mathcal{F})
      \ar[r]^{j\times1} \ar[d]_{\varsigma}
    & \mathbf{P}(\mathcal{E})\times_Y\mathbf{P}(\mathcal{F})
      \ar[d]^{\varsigma}
    \\\mathbf{P}(\mathcal{E}'\otimes\mathcal{F}) \ar[r]
    & \mathbf{P}(\mathcal{E}\otimes\mathcal{F})
  }
\]
est commutatif, $j$ désignant l'immersion fermée canonique
$\mathbf{P}(\mathcal{E}')\to\mathbf{P}(\mathcal{E})$. Posons
$P_1'=\mathbf{P}(\mathcal{E}')$ et gardons par ailleurs les notations de
(\hyperref[II.4.3.1-fr]{4.3.1})~; $j\times1$ est une immersion fermée
(\hyperref[I.4.3.1-fr]{I, 4.3.1}) et on a, à des isomorphismes près,
\[
  (j\times1)^*(\mathcal{O}_{P_1}(1)\otimes\mathcal{O}_{P_2}(1))
  =j^*(\mathcal{O}_{P_1}(1))\otimes\mathcal{O}_{P_2}(1)
  =\mathcal{O}_{P_1'}(1)\otimes\mathcal{O}_{P_2}(1)
\]
en vertu de (\hyperref[II.4.1.2.1-fr]{4.1.2.1}) et de
(\hyperref[I.9.1.5-fr]{I, 9.1.5})~; notre assertion résulte donc de
(\hyperref[II.4.2.8-fr]{4.2.8}) et
(\hyperref[II.4.2.9-fr]{4.2.9}).
\end{env}

\begin{env}[4.3.5]
\label{II.4.3.5-fr}
Avec les notations de (\hyperref[II.4.3.1-fr]{4.3.1}), soit
$\psi:Y'\to Y$ un morphisme, et posons
$\mathcal{E}'=\psi^*(\mathcal{E})$, $\mathcal{F}'=\psi^*(\mathcal{F})$~;
alors le morphisme de Segre
$\mathbf{P}(\mathcal{E}')\times\mathbf{P}(\mathcal{F}')
\to\mathbf{P}(\mathcal{E}'\otimes\mathcal{F}')$ s'identifie à
$\varsigma_{(Y')}$. En effet, gardant les notations de
(\hyperref[II.4.3.1-fr]{4.3.1}), posons par ailleurs
$P_1'=\mathbf{P}(\mathcal{E}')$, $P_2'=\mathbf{P}(\mathcal{F}')$~; on sait
(\hyperref[II.4.1.3.1-fr]{4.1.3.1}) que $P_i'$ s'identifie à $(P_i)_{(Y')}$
($i=1,2$), donc le morphisme structural $P_1'\times P_2'\to Y'$ s'identifie
à $r_{(Y')}$. D'autre part,
$\mathcal{E}'\otimes\mathcal{F}'$ s'identifie à
$\psi^*(\mathcal{E}\otimes\mathcal{F})$, donc
$\mathbf{P}(\mathcal{E}'\otimes\mathcal{F}')$ s'identifie à
$(\mathbf{P}(\mathcal{E}\otimes\mathcal{F}))_{(Y')}$
(\hyperref[II.4.1.3.1-fr]{4.1.3.1}). Enfin,
$\mathcal{O}_{P_1'}(1)\otimes_{Y'}\mathcal{O}_{P_2'}(1)=\mathcal{L}'$
s'identifie à $\mathcal{L}\otimes_Y\mathcal{O}_{Y'}$, en vertu de
(\hyperref[II.4.1.3.2-fr]{4.1.3.2}) et de
(\hyperref[I.9.1.11-fr]{I, 9.1.11}). L'homomorphisme canonique
$(r_{(Y')})^*(\mathcal{E}'\otimes\mathcal{F}')\to\mathcal{L}'$ s'identifie
alors à $s_{(Y')}$, et notre assertion résulte de
(\hyperref[II.4.2.10-fr]{4.2.10}).
\end{env}

\begin{remark}[4.3.6]
\label{II.4.3.6-fr}
Le préschéma somme de $\mathbf{P}(\mathcal{E})$ et
$\mathbf{P}(\mathcal{F})$ est de même canoniquement isomorphe à un
sous-préschéma fermé de $\mathbf{P}(\mathcal{E}\oplus\mathcal{F})$. En
effet, les homomorphismes surjectifs
$\mathcal{E}\oplus\mathcal{F}\to\mathcal{E}$,
$\mathcal{E}\oplus\mathcal{F}\to\mathcal{F}$ correspondent à des immersions
fermées
$\mathbf{P}(\mathcal{E})\to\mathbf{P}(\mathcal{E}\oplus\mathcal{F})$,
$\mathbf{P}(\mathcal{F})\to\mathbf{P}(\mathcal{E}\oplus\mathcal{F})$~;
tout revient à voir que les espaces sous-jacents des sous-préschémas fermés
de $\mathbf{P}(\mathcal{E}\oplus\mathcal{F})$ ainsi obtenus sont sans point
commun. La question étant locale sur $Y$, on peut adopter les notations de
(\hyperref[II.4.3.1-fr]{4.3.1})~; or $\mathbf{S}_n(E)$ et
$\mathbf{S}_n(F)$ s'identifient à des sous-modules de
$\mathbf{S}_n(E\oplus F)$ d'intersection réduite à $0$~; et si
$\mathfrak{p}$ est un idéal premier gradué de $\mathbf{S}(E)$ tel que
$\mathfrak{p}\cap\mathbf{S}_n(E)\neq\mathbf{S}_n(E)$ pour tout $n\geq0$,
il lui correspond dans $\mathbf{S}(E\oplus F)$ un idéal premier gradué dont
la trace sur $\mathbf{S}_n(E)$ est
$\mathfrak{p}\cap\mathbf{S}_n(E)$, mais qui contient $\mathbf{S}_+(F)$,
comme on le voit aussitôt~; deux points de
$\operatorname{Proj}(\mathbf{S}(E))$ et
$\operatorname{Proj}(\mathbf{S}(F))$ respectivement ne peuvent donc avoir
même image dans $\operatorname{Proj}(\mathbf{S}(E\oplus F))$.
\end{remark}

\subsection{Immersions dans les fibrés projectifs. Faisceaux très amples}
\label{subsection:II.4.4-fr}

\begin{proposition}[4.4.1]
\label{II.4.4.1-fr}
Soient $Y$ un schéma quasi-compact, ou un préschéma dont l'espace sous-jacent
est noethérien, $q:X\to Y$ un morphisme de type fini, $\mathcal{L}$ un
$\mathcal{O}_X$-Module inversible.
\begin{enumerate}
  \item[\rm{(i)}] Soient $\mathcal{S}$ une $\mathcal{O}_Y$-Algèbre graduée
  quasi-cohérente à degrés positifs,
  $\psi:q^*(\mathcal{S})\to\bigoplus_{n\geq0}\mathcal{L}^{\otimes n}$ un
  homomorphisme d'Algèbres graduées. Pour que $r_{\mathcal{L},\psi}$ soit
  partout défini et soit une immersion, il faut et

\oldpage[II]{79}

  il suffit qu'il existe un entier $n\geq0$, un sous-$\mathcal{O}_Y$-Module
  quasi-cohérent de type fini $\mathcal{E}$ de $\mathcal{S}_n$, tels que
  l'homomorphisme
  \[
    \psi'=\psi_n\circ q^*(j):q^*(\mathcal{E})
      \longrightarrow\mathcal{L}^{\otimes n}=\mathcal{L}'
  \]
  ($j$ étant l'injection $\mathcal{E}\to\mathcal{S}_n$) soit surjectif et que
  le morphisme
  $r_{\mathcal{L}',\psi'}:X\to\mathbf{P}(\mathcal{E})$ soit une immersion.
  \item[\rm{(ii)}] Soient $\mathcal{F}$ un $\mathcal{O}_Y$-Module
  quasi-cohérent, $\varphi:q^*(\mathcal{F})\to\mathcal{L}$ un homomorphisme
  surjectif. Pour que le morphisme $r_{\mathcal{L},\varphi}$ soit une
  immersion $X\to\mathbf{P}(\mathcal{F})$, il faut et il suffit qu'il existe
  un sous-$\mathcal{O}_Y$-Module quasi-cohérent de type fini $\mathcal{E}$ de
  $\mathcal{F}$ tel que l'homomorphisme
  $\varphi'=\varphi\circ q^*(j):q^*(\mathcal{E})\to\mathcal{L}$ (où
  $j:\mathcal{E}\to\mathcal{F}$ est l'injection canonique) soit surjectif et
  que $r_{\mathcal{L},\varphi'}:X\to\mathbf{P}(\mathcal{E})$ soit une
  immersion.
\end{enumerate}
\end{proposition}

\begin{proof}
\begin{enumerate}
  \item[\rm{(i)}] Le fait que $r_{\mathcal{L},\psi}$ soit partout défini et
  soit une immersion équivaut d'après
  (\hyperref[II.3.8.5-fr]{3.8.5}) à l'existence de $n\geq0$ et de
  $\mathcal{E}$ tels que, si $\mathcal{S}'$ est la sous-algèbre de
  $\mathcal{S}$ engendrée par $\mathcal{E}$, l'homomorphisme
  $q^*(\mathcal{E})\to\mathcal{L}^{\otimes n}$ soit surjectif et que
  $r_{\mathcal{L}',\psi'}:X\to\operatorname{Proj}(\mathcal{S}')$ soit partout
  défini et soit une immersion. On a d'autre part un homomorphisme canonique
  surjectif $\mathbf{S}(\mathcal{E})\to\mathcal{S}'$ auquel correspond une
  immersion fermée
  $\operatorname{Proj}(\mathcal{S}')\to\mathbf{P}(\mathcal{E})$
  (\hyperref[II.3.6.2-fr]{3.6.2})~; d'où la conclusion.
  \item[\rm{(ii)}] Comme $\mathcal{F}$ est limite inductive de ses
  sous-Modules quasi-cohérents de type fini $\mathcal{E}_\lambda$
  (\hyperref[I.9.4.9-fr]{I, 9.4.9}), $\mathbf{S}(\mathcal{F})$ est limite
  inductive des $\mathbf{S}(\mathcal{E}_\lambda)$~; la conclusion résulte
  alors de (\hyperref[II.3.8.4-fr]{3.8.4}), en observant que, dans la
  démonstration de (\hyperref[II.3.8.4-fr]{3.8.4}), on peut prendre tous les
  $n_i$ égaux à $1$~: en effet ($Y$ étant supposé affine), si
  $r=r_{\mathcal{L},\varphi}$ est une immersion, $r(X)$ est un sous-espace
  quasi-compact de $\mathbf{P}(\mathcal{F})$ que l'on peut recouvrir par un
  nombre fini d'ouverts de $\mathbf{P}(\mathcal{F})$ de la forme $D_+(f)$ avec
  $f\in F$, tels que $D_+(f)\cap r(X)$ soit fermé.
\end{enumerate}
\end{proof}

\begin{definition}[4.4.2]
\label{II.4.4.2-fr}
Soient $Y$ un préschéma, $q:X\to Y$ un morphisme. On dit qu'un
$\mathcal{O}_X$-Module inversible $\mathcal{L}$ est \emph{très ample pour
$q$} (ou \emph{très ample pour $Y$}, ou simplement \emph{très ample} si
aucune confusion n'en résulte) s'il existe un $\mathcal{O}_Y$-Module
quasi-cohérent $\mathcal{E}$ et une $Y$-immersion $i$ de $X$ dans
$P=\mathbf{P}(\mathcal{E})$ tels que $\mathcal{L}$ soit isomorphe à
$i^*(\mathcal{O}_P(1))$.
\end{definition}

Il revient au même (\hyperref[II.4.2.3-fr]{4.2.3}) de dire qu'il existe un
$\mathcal{O}_Y$-Module quasi-cohérent $\mathcal{E}$ et un homomorphisme
surjectif $\varphi:q^*(\mathcal{E})\to\mathcal{L}$ tels que
$r_{\mathcal{L},\varphi}:X\to\mathbf{P}(\mathcal{E})$ soit une immersion.

On notera que l'existence d'un $\mathcal{O}_X$-Module très ample relativement
à $Y$ implique que $q$ est séparé
((\hyperref[II.3.1.3-fr]{3.1.3}) et
(\hyperref[I.5.5.1-fr]{I, 5.5.1, (i) et (ii)})).

\begin{corollary}[4.4.3]
\label{II.4.4.3-fr}
Supposons qu'il existe une $\mathcal{O}_Y$-Algèbre graduée quasi-cohérente
$\mathcal{S}$ engendrée par $\mathcal{S}_1$, et une $Y$-immersion
$i:X\to P=\operatorname{Proj}(\mathcal{S})$ telle que $\mathcal{L}$ soit
isomorphe à $i^*(\mathcal{O}_P(1))$~; alors $\mathcal{L}$ est très ample
relativement à $q$.
\end{corollary}

\begin{proof}
En effet, si $\mathcal{F}=\mathcal{S}_1$, l'homomorphisme canonique
$\mathbf{S}(\mathcal{F})\to\mathcal{S}$ est surjectif, donc, en composant
l'immersion fermée correspondante
$\operatorname{Proj}(\mathcal{S})\to\mathbf{P}(\mathcal{F})$
(\hyperref[II.3.6.2-fr]{3.6.2}) et l'immersion $i$, on obtient une immersion
$j:X\to\mathbf{P}(\mathcal{F})=P'$ telle que $\mathcal{L}$ soit isomorphe à
$j^*(\mathcal{O}_{P'}(1))$ (\hyperref[II.3.6.3-fr]{3.6.3}).
\end{proof}

\begin{proposition}[4.4.4]
\label{II.4.4.4-fr}
Soient $q:X\to Y$ un morphisme quasi-compact, $\mathcal{L}$ un
$\mathcal{O}_X$-Module inversible. Les propriétés suivantes sont
équivalentes~:
\begin{enumerate}
  \item[\rm{a)}] $\mathcal{L}$ est très ample relativement à $q$.
  \item[\rm{b)}] $q_*(\mathcal{L})$ est quasi-cohérent, l'homomorphisme
  canonique $\sigma:q^*(q_*(\mathcal{L}))\to\mathcal{L}$ est surjectif, et le
  morphisme
  $r_{\mathcal{L},\sigma}:X\to\mathbf{P}(q_*(\mathcal{L}))$ est une immersion.
\end{enumerate}
\end{proposition}

\begin{proof}
Comme $q$ est quasi-compact, on sait que $q_*(\mathcal{L})$ est quasi-cohérent
si $q$ est séparé (\hyperref[I.9.2.2-fr]{I, 9.2.2}).

\oldpage[II]{80}

On sait (\hyperref[II.3.4.7-fr]{3.4.7}) que l'existence d'un homomorphisme
surjectif $\varphi:q^*(\mathcal{E})\to\mathcal{L}$ ($\mathcal{E}$ étant un
$\mathcal{O}_Y$-Module quasi-cohérent) entraîne que $\sigma$ est surjectif~;
en outre, à la factorisation
$q^*(\mathcal{E})\to q^*(q_*(\mathcal{L}))\xrightarrow{\sigma}\mathcal{L}$
de $\varphi$ correspond canoniquement une factorisation
\[
  q^*(\mathbf{S}(\mathcal{E}))
    \longrightarrow q^*(\mathbf{S}(q_*(\mathcal{L})))
    \longrightarrow\bigoplus_{n\geq0}\mathcal{L}^{\otimes n}
\]
donc (\hyperref[II.3.8.3-fr]{3.8.3}) l'hypothèse que
$r_{\mathcal{L},\varphi}$ est une immersion entraîne qu'il en est de même de
$j=r_{\mathcal{L},\sigma}$~; en outre
(\hyperref[II.4.2.4-fr]{4.2.4}), $\mathcal{L}$ est isomorphe à
$j^*(\mathcal{O}_{P'}(1))$ en posant
$P'=\mathbf{P}(q_*(\mathcal{L}))$. On voit donc que a) et b) sont
équivalents.
\end{proof}

\begin{corollary}[4.4.5]
\label{II.4.4.5-fr}
Supposons $q$ quasi-compact. Pour que $\mathcal{L}$ soit très ample
relativement à $Y$, il faut et il suffit qu'il existe un recouvrement ouvert
$(U_\alpha)$ de $Y$ tel que
$\mathcal{L}|q^{-1}(U_\alpha)$ soit très ample relativement à $U_\alpha$ pour
tout $\alpha$.
\end{corollary}

\begin{proof}
En effet, la condition b) de (\hyperref[II.4.4.4-fr]{4.4.4}) est locale sur
$Y$.
\end{proof}

\begin{proposition}[4.4.6]
\label{II.4.4.6-fr}
Soient $Y$ un schéma quasi-compact, ou un préschéma dont l'espace sous-jacent
est noethérien, $q:X\to Y$ un morphisme de type fini, $\mathcal{L}$ un
$\mathcal{O}_X$-Module inversible. Les conditions a) et b) de
(\hyperref[II.4.4.4-fr]{4.4.4}) sont alors équivalentes aux suivantes~:
\begin{enumerate}
  \item[\rm{a')}] Il existe un $\mathcal{O}_Y$-Module quasi-cohérent de type
  fini $\mathcal{E}$ et un homomorphisme surjectif
  $\varphi:q^*(\mathcal{E})\to\mathcal{L}$ tel que
  $r_{\mathcal{L},\varphi}$ soit une immersion.
  \item[\rm{b')}] Il existe un sous-$\mathcal{O}_Y$-Module cohérent de type
  fini $\mathcal{E}$ de $q_*(\mathcal{L})$ ayant les propriétés énoncées dans
  a').
\end{enumerate}
\end{proposition}

\begin{proof}
Il est clair que a') ou b') implique a)~; d'autre part a) implique a') en
vertu de (\hyperref[II.4.4.1-fr]{4.4.1}) et de même b) implique b').
\end{proof}

\begin{corollary}[4.4.7]
\label{II.4.4.7-fr}
Supposons que $Y$ soit un schéma quasi-compact ou un préschéma noethérien. Si
$\mathcal{L}$ est très ample pour $q$, il existe une $\mathcal{O}_Y$-Algèbre
graduée quasi-cohérente $\mathcal{S}$, telle que $\mathcal{S}_1$ soit de type
fini et engendre $\mathcal{S}$, et une $Y$-immersion ouverte dominante
$i:X\to P=\operatorname{Proj}(\mathcal{S})$ telle que $\mathcal{L}$ soit
isomorphe à $i^*(\mathcal{O}_P(1))$.
\end{corollary}

\begin{proof}
En effet, la condition b') de (\hyperref[II.4.4.6-fr]{4.4.6}) est vérifiée~;
le morphisme structural $p:\mathbf{P}(\mathcal{E})=P'\to Y$ est alors séparé
et de type fini (\hyperref[II.3.1.3-fr]{3.1.3}), donc $P'$ est un schéma
quasi-compact (resp. un préschéma noethérien) si $Y$ est un schéma
quasi-compact (resp. un préschéma noethérien). Soit $Z$ l'adhérence
(\hyperref[I.9.5.11-fr]{I, 9.5.11}) du sous-préschéma $X'$ de $P'$ associé à
l'immersion $j=r_{\mathcal{L},\varphi}$ de $X$ dans $P'$~; il est clair que
$j$ se factorise en l'injection canonique $Z\to P'$ et en une immersion
ouverte dominante $i:X\to Z$. Mais $Z$ s'identifie à un préschéma
$\operatorname{Proj}(\mathcal{S})$, où $\mathcal{S}$ est une
$\mathcal{O}_Y$-Algèbre graduée quotient de $\mathbf{S}(\mathcal{E})$ par un
faisceau quasi-cohérent gradué d'idéaux
(\hyperref[II.3.6.2-fr]{3.6.2}), et il est clair que $\mathcal{S}_1$ est de
type fini et engendre $\mathcal{S}$~; en outre $\mathcal{O}_Z(1)$ est l'image
réciproque de $\mathcal{O}_{P'}(1)$ par l'injection canonique
(\hyperref[II.3.6.3-fr]{3.6.3}), donc
$\mathcal{L}=i^*(\mathcal{O}_Z(1))$.
\end{proof}

\begin{proposition}[4.4.8]
\label{II.4.4.8-fr}
Soient $q:X\to Y$ un morphisme, $\mathcal{L}$ un $\mathcal{O}_X$-Module très
ample relativement à $q$, $\mathcal{L}'$ un $\mathcal{O}_X$-Module inversible
tel qu'il existe un $\mathcal{O}_Y$-Module quasi-cohérent $\mathcal{E}'$ et un
homomorphisme surjectif $q^*(\mathcal{E}')\to\mathcal{L}'$. Alors
$\mathcal{L}\otimes_{\mathcal{O}_X}\mathcal{L}'$ est très ample relativement
à $q$.
\end{proposition}

\begin{proof}
En effet, l'hypothèse entraîne l'existence d'un $Y$-morphisme
$r':X\to P'=\mathbf{P}(\mathcal{E}')$ tel que
$\mathcal{L}'=r'^*(\mathcal{O}_{P'}(1))$
(\hyperref[II.4.2.1-fr]{4.2.1}). Il y a par hypothèse un
$\mathcal{O}_Y$-Module quasi-cohérent $\mathcal{E}$ et une

\oldpage[II]{81}

$Y$-immersion $r:X\to P=\mathbf{P}(\mathcal{E})$ telle que
$\mathcal{L}=r^*(\mathcal{O}_P(1))$. Posons
$Q=\mathbf{P}(\mathcal{E}\otimes\mathcal{E}')$ et considérons le morphisme de
Segre $\varsigma:P\times_YP'\to Q$ (\hyperref[II.4.3.1-fr]{4.3.1}). Comme
$r$ est une immersion, il en est de même de
$(r,r')_Y:X\to P\times_YP'$ (\hyperref[I.5.3.14-fr]{I, 5.3.14})~; mais
$\varsigma$ étant une immersion (\hyperref[II.4.3.3-fr]{4.3.3}), il en est de
même de
$r'':X\xrightarrow{(r,r')_Y}P\times_YP'\xrightarrow{\varsigma}Q$. D'autre
part (\hyperref[II.4.3.2.1-fr]{4.3.2.1}),
$\varsigma^*(\mathcal{O}_Q(1))$ est isomorphe à
$\mathcal{O}_P(1)\otimes_Y\mathcal{O}_{P'}(1)$, donc
(\hyperref[I.9.1.4-fr]{I, 9.1.4}), $r''^*(\mathcal{O}_Q(1))$ est isomorphe à
$\mathcal{L}\otimes\mathcal{L}'$, ce qui démontre la proposition.
\end{proof}

\begin{corollary}[4.4.9]
\label{II.4.4.9-fr}
Soit $q:X\to Y$ un morphisme.
\begin{enumerate}
  \item[\rm{(i)}] Soient $\mathcal{L}$ un $\mathcal{O}_X$-Module inversible,
  $\mathcal{K}$ un $\mathcal{O}_Y$-Module inversible. Pour que $\mathcal{L}$
  soit très ample relativement à $q$, il faut et il suffit que
  $\mathcal{L}\otimes q^*(\mathcal{K})$ le soit.
  \item[\rm{(ii)}] Si $\mathcal{L}$ et $\mathcal{L}'$ sont deux
  $\mathcal{O}_X$-Modules très amples relativement à $q$, il en est de même
  de $\mathcal{L}\otimes\mathcal{L}'$~; en particulier
  $\mathcal{L}^{\otimes n}$ est très ample relativement à $q$ pour tout
  $n>0$.
\end{enumerate}
\end{corollary}

\begin{proof}
(ii) est conséquence immédiate de (\hyperref[II.4.4.8-fr]{4.4.8}), ainsi que
la nécessité de la condition (i)~; d'autre part, si
$\mathcal{L}\otimes q^*(\mathcal{K})$ est très ample, il en est de même de
$(\mathcal{L}\otimes q^*(\mathcal{K}))\otimes q^*(\mathcal{K}^{-1})$
d'après ce qui précède, et ce dernier $\mathcal{O}_X$-Module est isomorphe à
$\mathcal{L}$ (\hyperref[0.5.4.3-fr]{0, 5.4.3} et
\hyperref[0.5.4.5-fr]{5.4.5}).
\end{proof}

\begin{proposition}[4.4.10]
\label{II.4.4.10-fr}
\begin{enumerate}
  \item[\rm{(i)}] Pour tout préschéma $Y$, tout $\mathcal{O}_Y$-Module
  inversible $\mathcal{L}$ est très ample relativement au morphisme identique
  $1_Y$.
  \item[\rm{(i bis)}] Soient $f:X\to Y$ un morphisme, $j:X'\to X$ une
  immersion. Si $\mathcal{L}$ est un $\mathcal{O}_X$-Module très ample
  relativement à $f$, $j^*(\mathcal{L})$ est très ample relativement à
  $f\circ j$.
  \item[\rm{(ii)}] Soient $Z$ un préschéma quasi-compact, $f:X\to Y$ un
  morphisme de type fini, $g:Y\to Z$ un morphisme quasi-compact,
  $\mathcal{L}$ un $\mathcal{O}_X$-Module très ample relativement à $f$,
  $\mathcal{K}$ un $\mathcal{O}_Y$-Module très ample relativement à $g$.
  Alors il existe un entier $n_0>0$ tel que
  $\mathcal{L}\otimes f^*(\mathcal{K}^{\otimes n})$ soit très ample
  relativement à $g\circ f$ pour tout $n\geq n_0$.
  \item[\rm{(iii)}] Soient $f:X\to Y$, $g:Y'\to Y$ deux morphismes, et posons
  $X'=X_{(Y')}$. Si $\mathcal{L}$ est un $\mathcal{O}_X$-Module très ample
  relativement à $f$, $\mathcal{L}'=\mathcal{L}\otimes_Y\mathcal{O}_{Y'}$
  est un $\mathcal{O}_{X'}$-Module très ample relativement à $f_{(Y')}$.
  \item[\rm{(iv)}] Soient $f_i:X_i\to Y_i$ ($i=1,2$) deux $S$-morphismes.
  Si $\mathcal{L}_i$ est un $\mathcal{O}_{X_i}$-Module très ample relativement
  à $f_i$ ($i=1,2$), $\mathcal{L}_1\otimes_S\mathcal{L}_2$ est très ample
  relativement à $f_1\times_Sf_2$.
  \item[\rm{(v)}] Soient $f:X\to Y$, $g:Y\to Z$ deux morphismes. Si un
  $\mathcal{O}_X$-Module $\mathcal{L}$ est très ample relativement à
  $g\circ f$, alors $\mathcal{L}$ est très ample relativement à $f$.
  \item[\rm{(vi)}] Soient $f:X\to Y$ un morphisme, $j$ l'injection canonique
  $X_{\mathrm{red}}\to X$. Si $\mathcal{L}$ est un $\mathcal{O}_X$-Module
  très ample relativement à $f$, alors $j^*(\mathcal{L})$ est très ample
  relativement à $f_{\mathrm{red}}$.
\end{enumerate}
\end{proposition}

\begin{proof}
La propriété (i bis) découle aussitôt de la définition
(\hyperref[II.4.4.2-fr]{4.4.2}) et il est immédiat que (vi) se déduit
formellement de (i bis) et de (v) par un raisonnement calqué sur celui de
(\hyperref[I.5.5.12-fr]{I, 5.5.12}), que nous laissons au lecteur. Pour
démontrer (v), considérons, comme dans
(\hyperref[I.5.5.12-fr]{I, 5.5.12}), la factorisation
$X\xrightarrow{\Gamma_f}X\times_ZY\xrightarrow{p_2}Y$, en notant que
$p_2=(g\circ f)\times1_Y$. Il résulte de l'hypothèse et de (i) et (iv) que
$\mathcal{L}\otimes_{\mathcal{O}_Z}\mathcal{O}_Y$ est très ample pour $p_2$~;
d'autre part, on a
$\mathcal{L}=\Gamma_f^*(\mathcal{L}\otimes_{\mathcal{O}_Z}\mathcal{O}_Y)$
(\hyperref[I.9.1.4-fr]{I, 9.1.4}), et $\Gamma_f$ est une immersion
(\hyperref[I.5.3.11-fr]{I, 5.3.11})~; on peut donc appliquer (i bis).

\oldpage[II]{82}

Pour prouver (i), on applique la définition
(\hyperref[II.4.4.2-fr]{4.4.2}) avec $\mathcal{E}=\mathcal{L}$, en notant
qu'alors $\mathbf{P}(\mathcal{E})$ s'identifie à $Y$
(\hyperref[II.4.1.4-fr]{4.1.4}).

Démontrons (iii). Il existe un $\mathcal{O}_Y$-Module quasi-cohérent
$\mathcal{E}$ et une $Y$-immersion $i:X\to\mathbf{P}(\mathcal{E})=P$ telle
que $\mathcal{L}=i^*(\mathcal{O}_P(1))$~; alors, si on pose
$\mathcal{E}'=g^*(\mathcal{E})$, $\mathcal{E}'$ est un
$\mathcal{O}_{Y'}$-Module quasi-cohérent et on a
$P'=\mathbf{P}(\mathcal{E}')=P_{(Y')}$
(\hyperref[II.4.1.3.1-fr]{4.1.3.1}), $i_{(Y')}$ est une immersion de
$X_{(Y')}$ dans $P'$ (\hyperref[I.4.3.2-fr]{I, 4.3.2}) et $\mathcal{L}'$ est
isomorphe à $(i_{(Y')})^*(\mathcal{O}_{P'}(1))$
(\hyperref[II.4.2.10-fr]{4.2.10}).

Pour prouver (iv), remarquons qu'il y a par hypothèse une $Y_i$-immersion
$r_i:X_i\to P_i=\mathbf{P}(\mathcal{E}_i)$, où $\mathcal{E}_i$ est un
$\mathcal{O}_{Y_i}$-Module quasi-cohérent, et
$\mathcal{L}_i=r_i^*(\mathcal{O}_{P_i}(1))$ ($i=1,2$)~;
$r_1\times_Sr_2$ est une $S$-immersion de $X_1\times_SX_2$ dans
$P_1\times_SP_2$ (\hyperref[I.4.3.1-fr]{I, 4.3.1}) et l'image réciproque de
$\mathcal{O}_{P_1}(1)\otimes_S\mathcal{O}_{P_2}(1)$ par cette immersion est
$\mathcal{L}_1\otimes_S\mathcal{L}_2$. D'autre part, posons
$T=Y_1\times_SY_2$, et soient $p_1,p_2$ les projections de $T$ dans $Y_1,Y_2$
respectivement. Si on pose
$P_i'=\mathbf{P}(p_i^*(\mathcal{E}_i))$ ($i=1,2$), on a par
(\hyperref[II.4.1.3.1-fr]{4.1.3.1}), $P_i'=P_i\times_{Y_i}T$ d'où
\[
  P_1'\times_TP_2'
  =(P_1\times_{Y_1}T)\times_T(P_2\times_{Y_2}T)
  =P_1\times_{Y_1}(T\times_{Y_2}P_2)
  =P_1\times_{Y_1}(Y_1\times_SP_2)
  =P_1\times_SP_2
\]
à un isomorphisme canonique près. De même, on a
$\mathcal{O}_{P_i'}(1)=\mathcal{O}_{P_i}(1)\otimes_{Y_i}\mathcal{O}_T$
(\hyperref[II.4.1.3.2-fr]{4.1.3.2}), et un calcul analogue (basé notamment
sur (\hyperref[I.9.1.9.1-fr]{I, 9.1.9.1} et
\hyperref[I.9.1.2-fr]{9.1.2})) montre que dans l'identification précédente,
$\mathcal{O}_{P_1'}(1)\otimes_T\mathcal{O}_{P_2'}(1)$ s'identifie à
$\mathcal{O}_{P_1}(1)\otimes_S\mathcal{O}_{P_2}(1)$. On peut donc considérer
$r_1\times_Sr_2$ comme une $T$-immersion de $X_1\times_SX_2$ dans
$P_1'\times_TP_2'$, l'image réciproque de
$\mathcal{O}_{P_1'}(1)\otimes_T\mathcal{O}_{P_2'}(1)$ par cette immersion
étant $\mathcal{L}_1\otimes_S\mathcal{L}_2$. On termine alors le raisonnement
comme dans (\hyperref[II.4.4.8-fr]{4.4.8}) en utilisant le morphisme de Segre.

Il reste à démontrer (ii). On peut d'abord se restreindre au cas où $Z$ est un
schéma affine, car en général il y a un recouvrement fini $(U_i)$ de $Z$ par
des ouverts affines~; si la proposition est démontrée pour
$\mathcal{K}|g^{-1}(U_i)$, $\mathcal{L}|f^{-1}(g^{-1}(U_i))$ et un entier
$n_i$, il suffira de prendre pour $n_0$ le plus grand des $n_i$ pour démontrer
la proposition pour $\mathcal{K}$ et $\mathcal{L}$
(\hyperref[II.4.4.5-fr]{4.4.5}). L'hypothèse implique que $f$ et $g$ sont des
morphismes séparés, donc $X$ et $Y$ sont des schémas quasi-compacts.

Il y a une immersion $r:X\to P=\mathbf{P}(\mathcal{E})$, où $\mathcal{E}$ est
un $\mathcal{O}_Y$-Module quasi-cohérent de type fini et
$\mathcal{L}=r^*(\mathcal{O}_P(1))$, en vertu de
(\hyperref[II.4.4.6-fr]{4.4.6}). Nous allons voir qu'il existe un
$\mathcal{O}_P$-Module $\mathcal{M}$ très ample relativement au morphisme
composé $P\xrightarrow{h}Y\xrightarrow{g}Z$, tel que $\mathcal{O}_P(1)$ soit
isomorphe à $\mathcal{M}\otimes_Y\mathcal{K}^{\otimes(-m)}$ pour un entier
$m$. Pour $n\geq m+1$,
$\mathcal{O}_P(1)\otimes_Y\mathcal{K}^{\otimes n}$ sera donc très ample pour
$Z$ en vertu de l'hypothèse et de (iv) appliqué aux morphismes $h:P\to Y$ et
$1_Y$~; comme $r$ est une immersion et que
$\mathcal{L}\otimes f^*(\mathcal{K}^{\otimes n})
=r^*(\mathcal{O}_P(1)\otimes_Y\mathcal{K}^{\otimes n})$, la conclusion
résultera de (i bis). Pour établir notre assertion sur $\mathcal{O}_P(1)$,
nous utiliserons le lemme suivant~:

\begin{lemma}[4.4.10.1]
\label{II.4.4.10.1-fr}
Soient $Z$ un schéma quasi-compact ou un préschéma dont l'espace sous-jacent
est noethérien, $g:Y\to Z$ un morphisme quasi-compact, $\mathcal{K}$ un
$\mathcal{O}_Y$-Module inversible très ample pour $g$, $\mathcal{E}$ un
$\mathcal{O}_Y$-Module quasi-cohérent de type fini. Il existe alors un entier
$m_0$ tel que, pour tout $m\geq m_0$, $\mathcal{E}$ soit isomorphe à un
quotient d'un $\mathcal{O}_Y$-Module de la forme
$g^*(\mathcal{F})\otimes\mathcal{K}^{\otimes(-m)}$, où $\mathcal{F}$ est un
$\mathcal{O}_Z$-Module quasi-cohérent de type fini (dépendant de $m$).
\end{lemma}

Ce lemme sera démontré en (\hyperref[II.4.5.10.1-fr]{4.5.10.1})~; le lecteur
pourra vérifier que (\hyperref[II.4.4.10-fr]{4.4.10}) n'est utilisé nulle part
dans le n\textsuperscript{o}~4.5.

\oldpage[II]{83}

Ce lemme étant admis, il existe une immersion fermée $j_1$ de $P$ dans
\[
  P_1=\mathbf{P}(g^*(\mathcal{F})\otimes\mathcal{K}^{\otimes(-m)})
\]
telle que $\mathcal{O}_P(1)$ soit isomorphe à
$j_1^*(\mathcal{O}_{P_1}(1))$ (\hyperref[II.4.1.2-fr]{4.1.2}). D'autre part,
il existe un isomorphisme de $P_1$ sur
$P_2=\mathbf{P}(g^*(\mathcal{F}))$, identifiant
$\mathcal{O}_{P_1}(1)$ à
$\mathcal{O}_{P_2}(1)\otimes_Y\mathcal{K}^{\otimes(-m)}$
(\hyperref[II.4.1.4-fr]{4.1.4})~; on a donc une immersion fermée
$j_2:P\to P_2$ telle que $\mathcal{O}_P(1)$ soit isomorphe à
$j_2^*(\mathcal{O}_{P_2}(1))\otimes_Y\mathcal{K}^{\otimes(-m)}$. Enfin,
$P_2$ s'identifie à $P_3\times_ZY$, où $P_3=\mathbf{P}(\mathcal{F})$, et
$\mathcal{O}_{P_2}(1)$ à
$\mathcal{O}_{P_3}(1)\otimes_Z\mathcal{O}_Y$
(\hyperref[II.4.1.3-fr]{4.1.3}). Par définition,
$\mathcal{O}_{P_3}(1)$ est très ample pour $Z$~; comme il en est de même de
$\mathcal{K}$, on conclut de (iv) que
$\mathcal{O}_{P_2}(1)\otimes_Y\mathcal{K}$ est très ample pour $Z$~; il en
est de même de
$\mathcal{M}=j_2^*(\mathcal{O}_{P_2}(1)\otimes_Y\mathcal{K})$ en vertu de
(i bis), et $\mathcal{O}_P(1)$ est isomorphe à
$\mathcal{M}\otimes_Y\mathcal{K}^{\otimes(-m-1)}$, ce qui achève la
démonstration.
\end{proof}

\begin{proposition}[4.4.11]
\label{II.4.4.11-fr}
Soient $f:X\to Y$, $f':X'\to Y$ deux morphismes, $X''$ le préschéma somme
$X\sqcup X'$, $f''$ le morphisme $X''\to Y$ qui coïncide avec $f$ (resp.
$f'$) dans $X$ (resp. $X'$). Soit $\mathcal{L}$ (resp. $\mathcal{L}'$) un
$\mathcal{O}_X$-Module (resp. un $\mathcal{O}_{X'}$-Module) inversible, et
soit $\mathcal{L}''$ le $\mathcal{O}_{X''}$-Module inversible qui coïncide
avec $\mathcal{L}$ dans $X$ et avec $\mathcal{L}'$ dans $X'$. Pour que
$\mathcal{L}''$ soit très ample relativement à $f''$, il faut et il suffit
que $\mathcal{L}$ soit très ample relativement à $f$ et $\mathcal{L}'$ très
ample relativement à $f'$.
\end{proposition}

\begin{proof}
On est aussitôt ramené au cas où $Y$ est affine. Si $\mathcal{L}''$ est très
ample il en est de même de $\mathcal{L}$ et $\mathcal{L}'$ en vertu de
(\hyperref[II.4.4.10-fr]{4.4.10, (i bis)}). Inversement, si $\mathcal{L}$ et
$\mathcal{L}'$ sont très amples, il résulte aussitôt de la définition
(\hyperref[II.4.4.2-fr]{4.4.2}) et de
(\hyperref[II.4.3.6-fr]{4.3.6}) que $\mathcal{L}''$ est très ample.
\end{proof}

\subsection{Faisceaux amples}
\label{subsection:II.4.5-fr}

\begin{env}[4.5.1]
\label{II.4.5.1-fr}
Étant donnés un préschéma $X$ et un $\mathcal{O}_X$-Module inversible
$\mathcal{L}$, nous poserons, pour tout $\mathcal{O}_X$-Module $\mathcal{F}$
(lorsque aucune confusion ne sera possible sur $\mathcal{L}$),
$\mathcal{F}(n)=\mathcal{F}\otimes_{\mathcal{O}_X}\mathcal{L}^{\otimes n}$
($n\in\mathbf{Z}$)~; nous poserons d'autre part
\[
  S=\bigoplus_{n\geq0}\Gamma(X,\mathcal{L}^{\otimes n})
\]
(sous-anneau gradué de l'anneau $\Gamma_\bullet(\mathcal{L})$ défini dans
(\hyperref[0.5.4.6-fr]{0, 5.4.6})). Si on considère $X$ comme un
$\mathbf{Z}$-préschéma et qu'on désigne par $p$ le morphisme structural
$X\to\operatorname{Spec}(\mathbf{Z})$, il y a correspondance biunivoque
entre les homomorphismes
\[
  p^*(\widetilde{S})\longrightarrow
  \bigoplus_{n\geq0}\mathcal{L}^{\otimes n}
\]
de $\mathcal{O}_X$-Algèbres graduées et les endomorphismes de l'anneau gradué
$S$ (\hyperref[I.2.2.5-fr]{I, 2.2.5})~; l'homomorphisme
\[
  \varepsilon:p^*(\widetilde{S})\longrightarrow
  \bigoplus_{n\geq0}\mathcal{L}^{\otimes n}
\]
qui correspond ainsi à l'automorphisme identique de $S$ sera dit canonique.
Il lui correspond (\hyperref[II.3.7.1-fr]{3.7.1}) un morphisme
$G(\varepsilon)\to\operatorname{Proj}(S)$ qui sera aussi dit canonique.
\end{env}

\begin{theorem}[4.5.2]
\label{II.4.5.2-fr}
Soient $X$ un schéma quasi-compact ou un préschéma dont l'espace sous-jacent
est noethérien, $\mathcal{L}$ un $\mathcal{O}_X$-Module inversible, $S$
l'anneau gradué
$\bigoplus_{n\geq0}\Gamma(X,\mathcal{L}^{\otimes n})$. Les conditions
suivantes sont équivalentes~:
\begin{enumerate}
  \item[\rm{a)}] Lorsque $f$ parcourt l'ensemble des éléments homogènes de
  $S_+$, les $X_f$ forment une base de la topologie de $X$.
  \item[\rm{a')}] Lorsque $f$ parcourt l'ensemble des éléments homogènes de
  $S_+$, ceux des $X_f$ qui sont affines forment un recouvrement de $X$.
  \item[\rm{b)}] Le morphisme canonique
  $G(\varepsilon)\to\operatorname{Proj}(S)$
  (\hyperref[II.4.5.1-fr]{4.5.1}) est partout défini et est une immersion
  ouverte dominante.

\oldpage[II]{84}

  \item[\rm{b')}] Le morphisme canonique
  $G(\varepsilon)\to\operatorname{Proj}(S)$ est partout défini et est un
  homéomorphisme de l'espace sous-jacent à $X$ sur un sous-espace de
  $\operatorname{Proj}(S)$.
  \item[\rm{c)}] Pour tout $\mathcal{O}_X$-Module quasi-cohérent
  $\mathcal{F}$, si on désigne par $\mathcal{F}_n$ le
  sous-$\mathcal{O}_X$-Module de $\mathcal{F}(n)$ engendré par les sections de
  $\mathcal{F}(n)$ au-dessus de $X$, alors $\mathcal{F}$ est la somme des
  sous-$\mathcal{O}_X$-Modules $\mathcal{F}_n(-n)$, pour les entiers $n>0$.
  \item[\rm{c')}] La propriété c) a lieu pour tout faisceau quasi-cohérent
  d'idéaux de $\mathcal{O}_X$.
\end{enumerate}
En outre, si $(f_\alpha)$ est une famille d'éléments homogènes de $S_+$,
telle que les $X_{f_\alpha}$ soient affines, alors la restriction à
$\bigcup_\alpha X_{f_\alpha}$ du morphisme canonique
$X\to\operatorname{Proj}(S)$ est un isomorphisme de
$\bigcup_\alpha X_{f_\alpha}$ sur
$\bigcup_\alpha(\operatorname{Proj}(S))_{f_\alpha}$.
\end{theorem}

\begin{proof}
Il est clair que b) implique b'), et b') implique a) en vertu de la formule
(\hyperref[II.3.7.3.1-fr]{3.7.3.1}) (tenant compte de ce que
$\varepsilon^\flat$ est l'identité). La condition a) implique a'), car tout
$x\in X$ a un voisinage affine $U$ tel que
$\mathcal{L}|U$ soit isomorphe à $\mathcal{O}_X|U$~; si
$f\in\Gamma(X,\mathcal{L}^{\otimes n})$ est telle que $x\in X_f\subset U$,
$X_f$ est aussi l'ensemble des $x'\in U$ tels que $(f|U)(x')\neq0$, et c'est
par suite un ouvert affine (\hyperref[I.1.3.6-fr]{I, 1.3.6}). Pour prouver
que a') entraîne b), il suffit de prouver la dernière assertion de l'énoncé,
et de vérifier de plus que si $X=\bigcup_\alpha X_{f_\alpha}$, la condition
(iv) de (\hyperref[II.3.8.2-fr]{3.8.2}) est satisfaite. Ce dernier point
résulte aussitôt de (\hyperref[I.9.3.1-fr]{I, 9.3.1, (i)}). Quant à la
dernière assertion de (\hyperref[II.4.5.2-fr]{4.5.2}), comme
$X_{f_\alpha}$ est l'image réciproque de
$(\operatorname{Proj}(S))_{f_\alpha}$ par
$G(\varepsilon)\to\operatorname{Proj}(S)$, il suffit d'appliquer
(\hyperref[I.9.3.2-fr]{I, 9.3.2}). Donc a), a'), b), b') sont équivalentes.

Pour montrer que a') entraîne c), notons que si $X_f$ est affine (avec
$f\in S_k$), $\mathcal{F}|X_f$ est engendré par ses sections au-dessus de
$X_f$ (\hyperref[I.1.3.9-fr]{I, 1.3.9})~; d'autre part
(\hyperref[I.9.3.1-fr]{I, 9.3.1, (ii)}) une telle section $s$ est de la forme
$(t|X_f)\otimes(f|X_f)^{-m}$ où
$t\in\Gamma(X,\mathcal{F}(km))$~; par définition, $t$ est aussi une section
de $\mathcal{F}_{km}$, donc $s$ est bien une section de
$\mathcal{F}_{km}(-km)$ au-dessus de $X_f$, ce qui prouve c). Il est clair
que c) implique c'), et il reste à montrer que c') entraîne a). Or, soit $U$
un voisinage ouvert de $x\in X$, et soit $\mathcal{J}$ un faisceau d'idéaux
quasi-cohérent de $\mathcal{O}_X$ définissant un sous-préschéma fermé de $X$
ayant $X-U$ comme espace sous-jacent
(\hyperref[I.5.2.1-fr]{I, 5.2.1}). L'hypothèse c') entraîne qu'il existe un
entier $n>0$ et une section $f$ de $\mathcal{J}(n)$ au-dessus de $X$ telle
que $f(x)\neq0$. Mais on a évidemment $f\in S_n$, et $x\in X_f\subset U$, ce
qui prouve a).
\end{proof}

Lorsque $X$ est un préschéma dont l'espace sous-jacent est noethérien, les
conditions équivalentes de (\hyperref[II.4.5.2-fr]{4.5.2}) impliquent que
$X$ est un schéma, puisqu'il est isomorphe à un sous-préschéma du schéma
$S=\operatorname{Proj}(A)$ en vertu de
(\hyperref[II.4.5.2-fr]{4.5.2, b)}).

\begin{definition}[4.5.3]
\label{II.4.5.3-fr}
On dit qu'un $\mathcal{O}_X$-Module inversible $\mathcal{L}$ est ample si
$X$ est un schéma quasi-compact et si les conditions équivalentes de
(\hyperref[II.4.5.2-fr]{4.5.2}) sont vérifiées.
\end{definition}

Il résulte évidemment du critère (\hyperref[II.4.5.2-fr]{4.5.2, a)}) que si
$\mathcal{L}$ est un $\mathcal{O}_X$-Module ample, alors, pour tout ouvert
$U$ de $X$, $\mathcal{L}|U$ est un $(\mathcal{O}_X|U)$-Module ample.

Il résulte de la démonstration de (\hyperref[II.4.5.2-fr]{4.5.2}) que les
$X_f$ affines forment déjà une base de la topologie de $X$. De plus~:

\begin{corollary}[4.5.4]
\label{II.4.5.4-fr}
Soit $\mathcal{L}$ un $\mathcal{O}_X$-Module ample. Pour tout sous-espace fini
$Z$ de $X$ et tout voisinage $U$ de $Z$, il existe un entier $n$ et un
$f\in\Gamma(X,\mathcal{L}^{\otimes n})$ tels que $X_f$ soit un voisinage
affine de $Z$ contenu dans $U$.
\end{corollary}

\begin{proof}
\oldpage[II]{85}

En vertu de (\hyperref[II.4.5.2-fr]{4.5.2, b)}), on peut se borner à prouver
que pour toute partie finie $Z'$ de $\operatorname{Proj}(S)$ et tout
voisinage ouvert $U$ de $Z'$, il existe un élément homogène $f\in S_+$ tel
que $Z'\subset(\operatorname{Proj}(S))_f\subset U$
(\hyperref[II.2.4.1-fr]{2.4.1}). Or, par définition, l'ensemble fermé $Y$,
complémentaire de $U$ dans $\operatorname{Proj}(S)$, est de la forme
$V_+(\mathfrak{I})$, où $\mathfrak{I}$ est un idéal gradué de $S$, ne
contenant pas $S_+$ (\hyperref[II.2.3.2-fr]{2.3.2})~; d'autre part, les
points de $Z'$ sont par définition des idéaux premiers gradués
$\mathfrak{p}_i$ de $S_+$ ne contenant pas $\mathfrak{I}$
(\hyperref[II.2.3.1-fr]{2.3.1}). Il existe donc un élément
$f\in\mathfrak{I}$ n'appartenant à aucun des $\mathfrak{p}_i$ (Bourbaki,
\emph{Alg. comm.}, chap.~II, \S~1, n\textsuperscript{o}~1, prop.~2), et comme
les $\mathfrak{p}_i$ sont gradués, le raisonnement fait \emph{loc. cit.}
montre qu'on peut même supposer $f$ homogène~; cet élément répond alors à la
question.
\end{proof}

\begin{proposition}[4.5.5]
\label{II.4.5.5-fr}
Supposons que $X$ soit un schéma quasi-compact, ou un préschéma dont l'espace
sous-jacent est noethérien. Les conditions a) à c') de
(\hyperref[II.4.5.2-fr]{4.5.2}) sont aussi équivalentes aux suivantes~:
\begin{enumerate}
  \item[\rm{d)}] Pour tout $\mathcal{O}_X$-Module quasi-cohérent
  $\mathcal{F}$ de type fini, il existe un entier $n_0$ tel que, pour tout
  $n\geq n_0$, $\mathcal{F}(n)$ soit engendré par ses sections au-dessus de
  $X$.
  \item[\rm{d')}] Pour tout $\mathcal{O}_X$-Module quasi-cohérent
  $\mathcal{F}$ de type fini, il existe deux entiers $n>0$, $k>0$ tels que
  $\mathcal{F}$ soit isomorphe à un quotient du $\mathcal{O}_X$-Module
  $\mathcal{L}^{\otimes(-n)}\otimes\mathcal{O}_X^k$.
  \item[\rm{d'')}] La propriété d') a lieu pour tout faisceau d'idéaux
  quasi-cohérent de type fini de $\mathcal{O}_X$.
\end{enumerate}
\end{proposition}

\begin{proof}
Comme $X$ est quasi-compact, si un $\mathcal{O}_X$-Module quasi-cohérent de
type fini $\mathcal{F}$ est tel que $\mathcal{F}(n)$ (qui est de type fini)
est engendré par ses sections au-dessus de $X$, $\mathcal{F}(n)$ est alors
engendré par un nombre fini de ces sections
(\hyperref[0.5.2.3-fr]{0, 5.2.3}), donc d) implique d') et il est clair que
d') implique d''). Comme tout $\mathcal{O}_X$-Module quasi-cohérent
$\mathcal{G}$ est limite inductive de ses sous-$\mathcal{O}_X$-Modules de
type fini (\hyperref[I.9.4.9-fr]{I, 9.4.9}), pour vérifier la condition c')
de (\hyperref[II.4.5.2-fr]{4.5.2}), il suffit de le faire pour un faisceau
quasi-cohérent d'idéaux de $\mathcal{O}_X$ qui est de type fini, et d'')
entraîne donc c'). Reste à voir que si $\mathcal{L}$ est ample la propriété
d) est vérifiée. Considérons un recouvrement fini de $X$ par des
$X_{f_i}$ ($f_i\in S_{n_i}$) qu'on peut supposer affines~; en remplaçant les
$f_i$ par des puissances convenables (ce qui ne change pas les $X_{f_i}$),
on peut supposer tous les $n_i$ égaux à un même entier $m$. Le faisceau
$\mathcal{F}|X_{f_i}$, étant de type fini par hypothèse, est engendré par un
nombre fini de ses sections $h_{ij}$ au-dessus de $X_{f_i}$
(\hyperref[I.1.3.13-fr]{I, 1.3.13})~; il existe donc un entier $k_0$ tel que
la section $h_{ij}\otimes f_i^{\otimes k_0}$ se prolonge en une section de
$\mathcal{F}(k_0m)$ au-dessus de $X$ pour tout couple $(i,j)$
(\hyperref[I.9.3.1-fr]{I, 9.3.1}). \emph{A fortiori} les
$h_{ij}\otimes f_i^{\otimes k}$ se prolongent en sections de
$\mathcal{F}(km)$ au-dessus de $X$ pour tout $k\geq k_0$, et pour ces valeurs
de $k$, $\mathcal{F}(km)$ est donc engendré par ses sections au-dessus de
$X$. Pour tout $p$ tel que $0<p<m$, $\mathcal{F}(p)$ est aussi de type fini,
donc il y a un entier $k_p$ tel que
$\mathcal{F}(p)(km)=\mathcal{F}(p+km)$ soit engendré par ses sections
au-dessus de $X$ pour $k\geq k_p$. Prenant $n_0$ plus grand que tous les
$k_pm$, on en conclut que $\mathcal{F}(n)$ est engendré par ses sections
au-dessus de $X$ pour tout $n\geq n_0$, car un tel $n$ s'écrit $n=km+p$ avec
$k\geq k_p$ et $0\leq p<m$.
\end{proof}

\begin{proposition}[4.5.6]
\label{II.4.5.6-fr}
Soient $X$ un schéma quasi-compact, $\mathcal{L}$ un $\mathcal{O}_X$-Module
inversible.
\begin{enumerate}
  \item[\rm{(i)}] Soit $n$ un entier $>0$. Pour que $\mathcal{L}$ soit
  ample, il faut et il suffit que $\mathcal{L}^{\otimes n}$ soit ample.
  \item[\rm{(ii)}] Soit $\mathcal{L}'$ un $\mathcal{O}_X$-Module inversible
  tel que, pour tout $x\in X$, il existe un entier $n>0$

\oldpage[II]{86}

  et une section $s'$ de $\mathcal{L}'^{\otimes n}$ au-dessus de $X$ tels
  que $s'(x)\neq0$. Alors, si $\mathcal{L}$ est ample, il en est de même de
  $\mathcal{L}\otimes\mathcal{L}'$.
\end{enumerate}
\end{proposition}

\begin{proof}
La propriété (i) est conséquence évidente du critère a) de
(\hyperref[II.4.5.2-fr]{4.5.2}) puisque $X_{f^{\otimes n}}=X_f$. D'autre
part, si $\mathcal{L}$ est ample, pour tout $x\in X$ et tout voisinage $U$ de
$x$, il existe $m>0$ et $f\in\Gamma(X,\mathcal{L}^{\otimes m})$ tel que
$x\in X_f\subset U$ (\hyperref[II.4.5.2-fr]{4.5.2, a)})~; si
$f'\in\Gamma(X,\mathcal{L}'^{\otimes n})$ est tel que $f'(x)\neq0$, on aura
$s(x)\neq0$ pour
$s=f^{\otimes n}\otimes f'^{\otimes m}\in
\Gamma(X,(\mathcal{L}\otimes\mathcal{L}')^{\otimes mn})$, donc
$x\in X_s\subset X_f\subset U$, ce qui prouve que
$\mathcal{L}\otimes\mathcal{L}'$ est ample
(\hyperref[II.4.5.2-fr]{4.5.2, a)}).
\end{proof}

\begin{corollary}[4.5.7]
\label{II.4.5.7-fr}
Le produit tensoriel de deux $\mathcal{O}_X$-Modules amples est ample.
\end{corollary}

\begin{corollary}[4.5.8]
\label{II.4.5.8-fr}
Soient $\mathcal{L}$ un $\mathcal{O}_X$-Module ample, $\mathcal{L}'$ un
$\mathcal{O}_X$-Module inversible~; il existe alors un entier $n_0>0$ tel
que $\mathcal{L}^{\otimes n}\otimes\mathcal{L}'$ soit ample et engendré par
ses sections au-dessus de $X$ pour $n\geq n_0$.
\end{corollary}

\begin{proof}
En effet, il résulte de (\hyperref[II.4.5.5-fr]{4.5.5}) qu'il existe un entier
$m_0$ tel que $\mathcal{L}^{\otimes m}\otimes\mathcal{L}'$ soit engendré par
ses sections au-dessus de $X$ pour $m\geq m_0$~; d'après
(\hyperref[II.4.5.6-fr]{4.5.6}) on peut donc prendre $n_0=m_0+1$.
\end{proof}

\begin{remark}[4.5.9]
\label{II.4.5.9-fr}
Soit $P=\mathrm{H}^1(X,\mathcal{O}_X^*)$ le groupe des classes de
$\mathcal{O}_X$-Modules inversibles (\hyperref[0.5.4.7-fr]{0, 5.4.7}), et
soit $P^+$ la partie de $P$ formée des classes de faisceaux amples. Supposons
que $P^+$ soit non vide. Alors il résulte de
(\hyperref[II.4.5.7-fr]{4.5.7}) et
(\hyperref[II.4.5.8-fr]{4.5.8}) que l'on a
\[
  P^++P^+\subset P^+
  \qquad\text{et}\qquad
  P^+-P^+=P
\]
autrement dit $P^+\cup\{0\}$ est l'ensemble des éléments positifs dans $P$
pour une structure de préordre sur $P$ compatible avec sa structure de
groupe, et qui est même archimédienne en vertu de
(\hyperref[II.4.5.8-fr]{4.5.8}). C'est pourquoi on dit parfois «~faisceau
positif~» au lieu de faisceau ample et «~faisceau négatif~» pour l'inverse
d'un tel faisceau (terminologie que nous ne suivrons pas).
\end{remark}

\begin{proposition}[4.5.10]
\label{II.4.5.10-fr}
Soient $Y$ un schéma affine, $q:X\to Y$ un morphisme séparé quasi-compact,
$\mathcal{L}$ un $\mathcal{O}_X$-Module inversible.
\begin{enumerate}
  \item[\rm{(i)}] Si $\mathcal{L}$ est très ample relativement à $q$, alors
  $\mathcal{L}$ est ample.
  \item[\rm{(ii)}] Supposons en outre que le morphisme $q$ soit de type fini.
  Alors, pour que $\mathcal{L}$ soit ample, il faut et il suffit qu'il possède
  l'une des propriétés équivalentes suivantes~:
  \begin{enumerate}
    \item[\rm{e)}] Il existe $n_0>0$ tel que pour tout entier $n\geq n_0$,
    $\mathcal{L}^{\otimes n}$ soit très ample relativement à $q$.
    \item[\rm{e')}] Il existe $n>0$ tel que $\mathcal{L}^{\otimes n}$ soit
    très ample relativement à $q$.
  \end{enumerate}
\end{enumerate}
\end{proposition}

\begin{proof}
La première assertion résulte de la définition
(\hyperref[II.4.4.2-fr]{4.4.2}) d'un $\mathcal{O}_X$-Module très ample~: si
$A$ est l'anneau de $Y$, il y a un $A$-module $E$ et un homomorphisme
surjectif
\[
  \psi:q^*((\mathbf{S}(E))^{\sim})\longrightarrow
  \bigoplus_{n\geq0}\mathcal{L}^{\otimes n}
\]
tel que $i=r_{\mathcal{L},\psi}$ soit une immersion partout définie
$X\to P=\mathbf{P}(\widetilde{E})$ et que l'on ait
$\mathcal{L}=i^*(\mathcal{O}_P(1))$~; comme les $D_+(f)$ pour $f$ homogène
dans $(\mathbf{S}(E))_+$ forment une base de la topologie de $P$, et que
$i^{-1}(D_+(f))=X_{\psi^\flat(f)}$ par
(\hyperref[II.3.7.3.1-fr]{3.7.3.1}), on voit que la condition a) de
(\hyperref[II.4.5.2-fr]{4.5.2}) est vérifiée, donc $\mathcal{L}$ est ample.

Prouvons maintenant que si $q$ est de type fini et si $\mathcal{L}$ est
ample, il vérifie la condition e). En premier lieu, il résulte du critère b)
de (\hyperref[II.4.5.2-fr]{4.5.2}) et de
(\hyperref[II.4.4.1-fr]{4.4.1, (i)}) qu'il existe

\oldpage[II]{87}

un entier $k_0>0$ tel que $\mathcal{L}^{\otimes k_0}$ soit très ample
relativement à $q$. D'autre part, en vertu de
(\hyperref[II.4.5.5-fr]{4.5.5}), il existe un entier $m_0$ tel que, pour
$m\geq m_0$, $\mathcal{L}^{\otimes m}$ soit engendré par ses sections
au-dessus de $X$. Posons $n_0=k_0+m_0$~; si $n\geq n_0$, on peut écrire
$n=k_0+m$ avec $m\geq m_0$, d'où
$\mathcal{L}^{\otimes n}
=\mathcal{L}^{\otimes k_0}\otimes\mathcal{L}^{\otimes m}$. Comme
$\mathcal{L}^{\otimes m}$ est engendré par ses sections au-dessus de $X$, il
résulte du critère (\hyperref[II.4.4.8-fr]{4.4.8}) et de
(\hyperref[II.3.4.7-fr]{3.4.7}) que $\mathcal{L}^{\otimes n}$ est très ample
relativement à $q$. Enfin, il est clair que e) entraîne e'), et e') implique
que $\mathcal{L}$ est ample en vertu de (i) et de
(\hyperref[II.4.5.6-fr]{4.5.6, (i)}).

\begin{env}[4.5.10.1]
\label{II.4.5.10.1-fr}
\emph{Démonstration du lemme (\hyperref[II.4.4.10.1-fr]{4.4.10.1}).} Posons
$\mathcal{E}(n)=\mathcal{E}\otimes\mathcal{K}^{\otimes n}$~; comme $g$ est
séparé (\hyperref[II.4.4.2-fr]{4.4.2}), le raisonnement de
(\hyperref[II.3.4.8-fr]{3.4.8}) s'applique et ramène à voir que
l'homomorphisme canonique
$g^*(g_*(\mathcal{E}(n)))\to\mathcal{E}(n)$ est surjectif pour $n$ assez
grand. En outre, comme $Z$ est quasi-compact, le raisonnement de
(\hyperref[II.3.4.6-fr]{3.4.6}) ramène à prouver la proposition dans le cas
où $Z$ est affine. Or, $\mathcal{K}$ est alors ample en vertu de
(\hyperref[II.4.5.10-fr]{4.5.10, (i)}), et la conclusion résulte de
(\hyperref[II.4.5.5-fr]{4.5.5, d')}).
\end{env}
\end{proof}

\begin{corollary}[4.5.11]
\label{II.4.5.11-fr}
Soient $Y$ un schéma affine, $q:X\to Y$ un morphisme séparé de type fini,
$\mathcal{L}$ un $\mathcal{O}_X$-Module ample, $\mathcal{L}'$ un
$\mathcal{O}_X$-Module inversible. Il existe un entier $n_0$ tel que, pour
$n\geq n_0$, $\mathcal{L}^{\otimes n}\otimes\mathcal{L}'$ soit très ample
relativement à $q$.
\end{corollary}

\begin{proof}
En effet, il existe $m_0$ tel que pour $m\geq m_0$,
$\mathcal{L}^{\otimes m}\otimes\mathcal{L}'$ soit engendré par ses sections
au-dessus de $X$ (\hyperref[II.4.5.8-fr]{4.5.8})~; d'autre part, il existe
$k_0$ tel que $\mathcal{L}^{\otimes k}$ soit très ample relativement à $q$
pour $k\geq k_0$. Par suite
$\mathcal{L}^{\otimes(k+m_0)}\otimes\mathcal{L}'$ est très ample pour
$k\geq k_0$ ((\hyperref[II.4.4.8-fr]{4.4.8}) et
(\hyperref[II.3.4.7-fr]{3.4.7})).
\end{proof}

\begin{remark}[4.5.12]
\label{II.4.5.12-fr}
On ignore si l'hypothèse qu'un $\mathcal{O}_X$-Module $\mathcal{L}$ est tel
que $\mathcal{L}^{\otimes n}$ soit très ample (relativement à $q$) entraîne
la même conclusion pour $\mathcal{L}^{\otimes(n+1)}$.
\end{remark}

\begin{proposition}[4.5.13]
\label{II.4.5.13-fr}
Soient $X$ un préschéma quasi-compact, $Z$ un sous-préschéma fermé de $X$
défini par un faisceau quasi-cohérent nilpotent $\mathcal{J}$ d'idéaux de
$\mathcal{O}_X$, $j$ l'injection canonique $Z\to X$. Pour qu'un
$\mathcal{O}_X$-Module inversible $\mathcal{L}$ soit ample, il faut et il
suffit que $\mathcal{L}'=j^*(\mathcal{L})$ soit un
$\mathcal{O}_Z$-Module ample.
\end{proposition}

\begin{proof}
La condition est nécessaire. En effet, pour toute section $f$ de
$\mathcal{L}^{\otimes n}$ au-dessus de $X$, soit $f'$ son image canonique
$f\otimes1$, section de
$\mathcal{L}'^{\otimes n}=\mathcal{L}^{\otimes n}
\otimes_{\mathcal{O}_X}(\mathcal{O}_X/\mathcal{J})$ au-dessus de l'espace
$Z$ (qui est identique à $X$)~; il est clair que $X_f=Z_{f'}$, donc le
critère a) de (\hyperref[II.4.5.2-fr]{4.5.2}) montre que $\mathcal{L}'$ est
ample.

Pour voir que la condition est suffisante, notons d'abord qu'on peut se
ramener au cas où $\mathcal{J}^2=0$, en considérant la suite (finie) des
préschémas
$X_k=(X,\mathcal{O}_X/\mathcal{J}^{k+1})$ dont chacun est sous-préschéma
fermé du suivant, défini par un faisceau d'idéaux de carré nul. D'ailleurs
$X$ est un schéma puisque $X_{\mathrm{red}}$ l'est par hypothèse
((\hyperref[II.4.5.3-fr]{4.5.3}) et
(\hyperref[I.5.5.1-fr]{I, 5.5.1})). Le critère a) de
(\hyperref[II.4.5.2-fr]{4.5.2}) montre qu'il suffira de prouver le

\begin{lemma}[4.5.13.1]
\label{II.4.5.13.1-fr}
Sous les hypothèses de (\hyperref[II.4.5.13-fr]{4.5.13}), supposons de plus
$\mathcal{J}$ de carré nul~; $\mathcal{L}$ étant un
$\mathcal{O}_X$-Module inversible, soit $g$ une section de
$\mathcal{L}'^{\otimes n}$ au-dessus de $Z$ telle que $Z_g$ soit affine.
Alors il existe un entier $m>0$ tel que $g^{\otimes m}$ soit l'image
canonique d'une section $f$ de $\mathcal{L}^{\otimes nm}$ au-dessus de $X$.
\end{lemma}

On a la suite exacte de $\mathcal{O}_X$-Modules
\[
  0\longrightarrow\mathcal{J}(n)\longrightarrow
  \mathcal{O}_X(n)=\mathcal{L}^{\otimes n}\longrightarrow
  \mathcal{O}_Z(n)=\mathcal{L}'^{\otimes n}\longrightarrow0
\]

\oldpage[II]{88}

puisque $\mathcal{F}(n)$ est un foncteur exact en $\mathcal{F}$~; d'où la
suite exacte de cohomologie
\[
  0\longrightarrow\Gamma(X,\mathcal{J}(n))\longrightarrow
  \Gamma(X,\mathcal{L}^{\otimes n})\longrightarrow
  \Gamma(X,\mathcal{L}'^{\otimes n})\xrightarrow{\partial}
  \mathrm{H}^1(X,\mathcal{J}(n))
\]
qui associe en particulier à $g$ un élément
$\partial g\in\mathrm{H}^1(X,\mathcal{J}(n))$.

Notons que puisque $\mathcal{J}^2=0$, $\mathcal{J}$ peut être considéré
comme un $\mathcal{O}_Z$-Module quasi-cohérent et l'on a, pour tout $k$,
$\mathcal{L}'^{\otimes k}\otimes_{\mathcal{O}_Z}\mathcal{J}(n)
=\mathcal{J}(n+k)$~; pour toute section
$s\in\Gamma(X,\mathcal{L}'^{\otimes k})$ la multiplication tensorielle par
$s$ est donc un homomorphisme $\mathcal{J}(n)\to\mathcal{J}(n+k)$ de
$\mathcal{O}_Z$-Modules, qui donne par suite un homomorphisme
$\mathrm{H}^i(X,\mathcal{J}(n))\to
\mathrm{H}^i(X,\mathcal{J}(n+k))$ de groupes de cohomologie.

Cela étant, nous allons voir que l'on a
\begin{equation}
\label{II.4.5.13.2-fr}
  g^{\otimes m}\otimes\partial g=0
\tag{4.5.13.2}
\end{equation}
pour un $m>0$ assez grand. En effet, $Z_g$ est un ouvert affine de $Z$ et
on a donc $\mathrm{H}^1(Z_g,\mathcal{J}(n))=0$ quand $\mathcal{J}(n)$ est
considéré comme un $\mathcal{O}_Z$-Module
(\hyperref[I.5.1.9.2-fr]{I, 5.1.9.2}). En particulier, si on pose
$g'=g|Z_g$, et si on considère son image par l'application
$\partial:\Gamma(Z_g,\mathcal{L}'^{\otimes n})\to
\mathrm{H}^1(Z_g,\mathcal{J}(n))$, on a $\partial g'=0$. Pour expliciter
cette relation, observons qu'en dimension $1$ la cohomologie d'un faisceau
de groupes abéliens est la même que sa cohomologie de Čech (G, II, 5.9)~;
pour former $\partial g$, il faut donc considérer un recouvrement ouvert
$(U_\alpha)$ assez fin de $X$, que l'on peut supposer fini et formé
d'ouverts affines, prendre pour chaque $\alpha$ une section
$g_\alpha\in\Gamma(U_\alpha,\mathcal{L}^{\otimes n})$ dont l'image
canonique dans $\Gamma(U_\alpha,\mathcal{L}'^{\otimes n})$ est
$g|U_\alpha$, et considérer la classe du cocycle
$(g_{\alpha|\beta}-g_{\beta|\alpha})$, $g_{\alpha|\beta}$ étant la
restriction de $g_\alpha$ à $U_\alpha\cap U_\beta$ (ce cocycle étant à
valeurs dans $\mathcal{J}(n)$). On peut en outre supposer que $\partial g'$
est calculé de la même manière à l'aide du recouvrement formé des
$U_\alpha\cap Z_g$ et des restrictions $g_\alpha|(U_\alpha\cap Z_g)$ (en
remplaçant au besoin $(U_\alpha)$ par un recouvrement plus fin)~; la
relation $\partial g'=0$ signifie alors qu'il existe pour chaque $\alpha$
une section $h_\alpha\in\Gamma(U_\alpha\cap Z_g,\mathcal{J}(n))$ telle que
\[
  (g_{\alpha|\beta}-g_{\beta|\alpha})
  |(U_\alpha\cap U_\beta\cap Z_g)
  =h_{\alpha|\beta}-h_{\beta|\alpha}
\]
en désignant par $h_{\alpha|\beta}$ la restriction de $h_\alpha$ à
$U_\alpha\cap U_\beta\cap Z_g$ (G, II, 5.11). Cela étant, il y a un entier
$m>0$ tel que $g^{\otimes m}\otimes h_\alpha$ soit la restriction à
$U_\alpha\cap Z_g$ d'une section
$t_\alpha\in\Gamma(U_\alpha,\mathcal{J}(n+nm))$ pour tout $\alpha$
(\hyperref[I.9.3.1-fr]{I, 9.3.1})~; on a donc
\[
  g^{\otimes m}\otimes(g_{\alpha|\beta}-g_{\beta|\alpha})
  =t_{\alpha|\beta}-t_{\beta|\alpha}
\]
pour tout couple d'indices, ce qui prouve
(\hyperref[II.4.5.13.2-fr]{4.5.13.2}).

Remarquons d'autre part que si
$s\in\Gamma(X,\mathcal{O}_Z(p))$,
$t\in\Gamma(X,\mathcal{O}_Z(q))$, on a, dans le groupe
$\mathrm{H}^1(X,\mathcal{J}(p+q))$
\begin{equation}
\label{II.4.5.13.3-fr}
  \partial(s\otimes t)=(\partial s)\otimes t+s\otimes(\partial t).
\tag{4.5.13.3}
\end{equation}

En effet, on peut encore, pour calculer les deux membres, considérer un
recouvrement ouvert $(U_\alpha)$ de $X$, pour chaque $\alpha$ une section
$s_\alpha\in\Gamma(U_\alpha,\mathcal{O}_X(p))$ (resp.
$t_\alpha\in\Gamma(U_\alpha,\mathcal{O}_X(q))$) dont l'image canonique dans
$\Gamma(U_\alpha,\mathcal{O}_Z(p))$ (resp.
$\Gamma(U_\alpha,\mathcal{O}_Z(q))$ soit $s|U_\alpha$ (resp.
$t|U_\alpha$)~; la relation
(\hyperref[II.4.5.13.3-fr]{4.5.13.3}) résulte alors des relations
\[
  (s_{\alpha|\beta}\otimes t_{\alpha|\beta})
  -(s_{\beta|\alpha}\otimes t_{\beta|\alpha})
  =(s_{\alpha|\beta}-s_{\beta|\alpha})\otimes t_{\alpha|\beta}
  +s_{\beta|\alpha}\otimes(t_{\alpha|\beta}-t_{\beta|\alpha})
\]
avec les mêmes notations que ci-dessus. Par récurrence sur $k$, on a donc
\begin{equation}
\label{II.4.5.13.4-fr}
  \partial(g^{\otimes k})=(kg^{\otimes(k-1)})\otimes(\partial g)
\tag{4.5.13.4}
\end{equation}

\oldpage[II]{89}

et on conclut de (\hyperref[II.4.5.13.2-fr]{4.5.13.2}) et
(\hyperref[II.4.5.13.4-fr]{4.5.13.4}) que l'on a
$\partial(g^{\otimes(m+1)})=0$~; $g^{\otimes(m+1)}$ est donc l'image
canonique d'une section $f$ de $\mathcal{L}^{\otimes n(m+1)}$ au-dessus de
$X$, ce qui achève de démontrer (\hyperref[II.4.5.13-fr]{4.5.13}).
\end{proof}

\begin{corollary}[4.5.14]
\label{II.4.5.14-fr}
Soient $X$ un schéma noethérien, $j$ l'injection canonique
$X_{\mathrm{red}}\to X$. Pour qu'un $\mathcal{O}_X$-Module inversible
$\mathcal{L}$ soit ample, il faut et il suffit que $j^*(\mathcal{L})$ soit
un $\mathcal{O}_{X_{\mathrm{red}}}$-Module ample.
\end{corollary}

Cela résulte de (\hyperref[I.6.1.6-fr]{I, 6.1.6}).

\subsection{Faisceaux relativement amples}
\label{II.4.6-fr}

\begin{definition}[4.6.1]
\label{II.4.6.1-fr}
Soient $f:X\to Y$ un morphisme quasi-compact, $\mathcal{L}$ un
$\mathcal{O}_X$-Module inversible. On dit que $\mathcal{L}$ est ample
relativement à $f$, ou relativement à $Y$, ou $f$-ample, ou $Y$-ample (ou
même ample si aucune confusion n'en résulte avec la notion définie dans
(\hyperref[II.4.5.3-fr]{4.5.3})) s'il existe un recouvrement $(U_\alpha)$
de $Y$ formé d'ouverts affines tel que si on pose
$X_\alpha=f^{-1}(U_\alpha)$, $\mathcal{L}|X_\alpha$ soit un
$\mathcal{O}_{X_\alpha}$-Module ample pour tout $\alpha$.
\end{definition}

L'existence d'un $\mathcal{O}_X$-Module $f$-ample entraîne que $f$ est
nécessairement séparé ((\hyperref[II.4.5.3-fr]{4.5.3}) et
(\hyperref[I.5.5.5-fr]{I, 5.5.5})).

\begin{proposition}[4.6.2]
\label{II.4.6.2-fr}
Soit $f:X\to Y$ un morphisme quasi-compact, $\mathcal{L}$ un
$\mathcal{O}_X$-Module inversible. Si $\mathcal{L}$ est très ample
relativement à $f$, il est ample relativement à $f$.
\end{proposition}

Cela résulte du caractère local (sur $Y$) de la notion de faisceau très
ample (\hyperref[II.4.4.5-fr]{4.4.5}), de la définition
(\hyperref[II.4.6.1-fr]{4.6.1}) et du critère
(\hyperref[II.4.5.10-fr]{4.5.10, (i)}).

\begin{proposition}[4.6.3]
\label{II.4.6.3-fr}
Soient $f:X\to Y$ un morphisme quasi-compact, $\mathcal{L}$ un
$\mathcal{O}_X$-Module inversible, et soit $\mathcal{S}$ la
$\mathcal{O}_Y$-Algèbre graduée
$\bigoplus_{n\geq0}f_*(\mathcal{L}^{\otimes n})$. Les conditions suivantes
sont équivalentes~:
\begin{enumerate}
  \item[\rm{a)}] $\mathcal{L}$ est $f$-ample.
  \item[\rm{b)}] $\mathcal{S}$ est quasi-cohérente et l'homomorphisme
  canonique $\sigma:f^*(\mathcal{S})\to
  \bigoplus_{n\geq0}\mathcal{L}^{\otimes n}$
  (\hyperref[0.4.4.3-fr]{0, 4.4.3}) est tel que le $Y$-morphisme
  $r_{\mathcal{L},\sigma}:G(\sigma)\to\operatorname{Proj}(\mathcal{S})=P$
  soit partout défini et soit une immersion ouverte dominante.
  \item[\rm{b')}] Le morphisme $f$ est séparé, le $Y$-morphisme
  $r_{\mathcal{L},\sigma}$ est partout défini et est un homéomorphisme de
  l'espace sous-jacent à $X$ sur un sous-espace de
  $\operatorname{Proj}(\mathcal{S})$.
\end{enumerate}
En outre, lorsqu'il en est ainsi, pour tout $n\in\mathbf{Z}$,
l'homomorphisme canonique
\begin{equation}
\label{II.4.6.3.1-fr}
  r_{\mathcal{L},\sigma}^*(\mathcal{O}_P(n))
  \longrightarrow\mathcal{L}^{\otimes n}
\tag{4.6.3.1}
\end{equation}
défini dans (\hyperref[II.3.7.9.1-fr]{3.7.9.1}), est un isomorphisme.

Enfin, pour tout $\mathcal{O}_X$-Module quasi-cohérent $\mathcal{F}$, si
l'on pose
$\mathcal{M}=\bigoplus_{n\geq0}f_*(\mathcal{F}\otimes
\mathcal{L}^{\otimes n})$, l'homomorphisme canonique
\begin{equation}
\label{II.4.6.3.2-fr}
  r_{\mathcal{L},\sigma}^*(\widetilde{\mathcal{M}})
  \longrightarrow\mathcal{F}
\tag{4.6.3.2}
\end{equation}
défini dans (\hyperref[II.3.7.9.2-fr]{3.7.9.2}), est un isomorphisme.
\end{proposition}

On a remarqué que a) implique que $f$ est séparé, donc $\mathcal{S}$ est
quasi-cohérente (\hyperref[I.9.2.2-fr]{I, 9.2.2, a)}). Comme le fait que
$r_{\mathcal{L},\sigma}$ soit une immersion partout définie est de
caractère local sur $Y$, pour prouver que a) entraîne b), on peut supposer
$Y$ affine et $\mathcal{L}$ ample~;

\oldpage[II]{90}

l'assertion résulte alors de (\hyperref[II.4.5.2-fr]{4.5.2, b)}). Il est
clair que b) entraîne b')~; enfin, pour prouver que b') entraîne a), il
suffit de considérer un recouvrement de $Y$ par des ouverts affines
$U_\alpha$ et d'appliquer le critère
(\hyperref[II.4.5.2-fr]{4.5.2, b')}) à chaque faisceau
$\mathcal{L}|f^{-1}(U_\alpha)$.

Pour les deux dernières assertions, on utilise le fait que
$\sigma^\flat$ est ici l'identité, et l'explicitation des homomorphismes
(\hyperref[II.3.7.9.1-fr]{3.7.9.1}) et
(\hyperref[II.3.7.9.2-fr]{3.7.9.2})~; il en résulte aussitôt que
(\hyperref[II.4.6.3.1-fr]{4.6.3.1}) est un isomorphisme. Quant à
(\hyperref[II.4.6.3.2-fr]{4.6.3.2}), on peut se ramener au cas où $Y$ est
affine, donc $\mathcal{L}$ ample~; il est clair que l'homomorphisme
(\hyperref[II.4.6.3.2-fr]{4.6.3.2}) est injectif, et le critère
(\hyperref[II.4.5.2-fr]{4.5.2, c)}) montre qu'il est surjectif, d'où la
conclusion.

\begin{corollary}[4.6.4]
\label{II.4.6.4-fr}
Soit $(U_\alpha)$ un recouvrement ouvert de $Y$. Pour que $\mathcal{L}$
soit ample relativement à $Y$, il faut et il suffit que
$\mathcal{L}|f^{-1}(U_\alpha)$ soit ample relativement à $U_\alpha$ pour
tout $\alpha$.
\end{corollary}

La condition b) est en effet locale sur $Y$.

\begin{corollary}[4.6.5]
\label{II.4.6.5-fr}
Soit $\mathcal{K}$ un $\mathcal{O}_Y$-Module inversible. Pour que
$\mathcal{L}$ soit $Y$-ample, il faut et il suffit que
$\mathcal{L}\otimes f^*(\mathcal{K})$ le soit.
\end{corollary}

C'est une conséquence évidente de (\hyperref[II.4.6.4-fr]{4.6.4}) en
prenant les $U_\alpha$ tels que $\mathcal{K}|U_\alpha$ soit isomorphe à
$\mathcal{O}_Y|U_\alpha$ pour tout $\alpha$.

\begin{corollary}[4.6.6]
\label{II.4.6.6-fr}
Supposons $Y$ affine~; pour que $\mathcal{L}$ soit $Y$-ample, il faut et il
suffit que $\mathcal{L}$ soit ample.
\end{corollary}

C'est une conséquence immédiate de la définition
(\hyperref[II.4.6.1-fr]{4.6.1}), et des critères
(\hyperref[II.4.6.3-fr]{4.6.3, b)}) et
(\hyperref[II.4.5.2-fr]{4.5.2, b)}), car ici
$\operatorname{Proj}(\mathcal{S})
=\operatorname{Proj}(\Gamma(Y,\mathcal{S}))$ par définition.

\begin{corollary}[4.6.7]
\label{II.4.6.7-fr}
Soit $f:X\to Y$ un morphisme quasi-compact. Supposons qu'il existe un
$\mathcal{O}_Y$-Module quasi-cohérent $\mathcal{E}$ et un $Y$-morphisme
$g:X\to P=\mathbf{P}(\mathcal{E})$ qui soit un homéomorphisme de l'espace
sous-jacent à $X$ sur un sous-espace de $P$~; alors
$\mathcal{L}=g^*(\mathcal{O}_P(1))$ est $Y$-ample.
\end{corollary}

On peut en effet supposer $Y$ affine~; le corollaire résulte alors du
critère (\hyperref[II.4.5.2-fr]{4.5.2, a)}), de la formule
(\hyperref[II.3.7.3.1-fr]{3.7.3.1}) et de
(\hyperref[II.4.2.3-fr]{4.2.3}).

\begin{proposition}[4.6.8]
\label{II.4.6.8-fr}
Soient $X$ un schéma quasi-compact ou un préschéma dont l'espace
sous-jacent est noethérien, $f:X\to Y$ un morphisme séparé quasi-compact.
Pour qu'un $\mathcal{O}_X$-Module inversible $\mathcal{L}$ soit $f$-ample,
il faut et il suffit que l'une des conditions équivalentes suivantes soit
vérifiée~:
\begin{enumerate}
  \item[\rm{c)}] Pour tout $\mathcal{O}_X$-Module quasi-cohérent
  $\mathcal{F}$ de type fini, il existe un entier $n_0>0$ tel que, pour
  tout $n\geq n_0$, l'homomorphisme canonique
  $\sigma:f^*(f_*(\mathcal{F}\otimes\mathcal{L}^{\otimes n}))\to
  \mathcal{F}\otimes\mathcal{L}^{\otimes n}$ soit surjectif.
  \item[\rm{c')}] Pour tout faisceau d'idéaux quasi-cohérent de type fini
  $\mathcal{J}$ de $\mathcal{O}_X$, il existe un entier $n>0$ tel que
  l'homomorphisme canonique
  $\sigma:f^*(f_*(\mathcal{J}\otimes\mathcal{L}^{\otimes n}))\to
  \mathcal{J}\otimes\mathcal{L}^{\otimes n}$ soit surjectif.
\end{enumerate}
\end{proposition}

Comme $X$ est quasi-compact, il en est de même de $f(X)$ et il existe donc
un recouvrement fini $(U_i)$ de $f(X)$ par des ouverts affines de $Y$. Pour
démontrer la condition c) lorsque $\mathcal{L}$ est $f$-ample, on peut
remplacer $Y$ par les $U_i$ et $X$ par les $f^{-1}(U_i)$, car si on obtient
pour chaque $i$ un entier $n_i$ tel que c) ait lieu (pour $U_i$,
$f^{-1}(U_i)$ et $\mathcal{L}|f^{-1}(U_i)$) dès que $n\geq n_i$, il
suffira de prendre pour $n_0$ le plus grand des $n_i$ pour obtenir c) pour
$Y$, $X$ et $\mathcal{L}$. Or lorsque $Y$ est affine, la condition c)
découle de (\hyperref[II.4.5.5-fr]{4.5.5, d)}), compte tenu de
(\hyperref[II.4.6.6-fr]{4.6.6}). Il est trivial que c) entraîne c'). Enfin,
pour prouver que c') implique que $\mathcal{L}$ est $f$-ample, on peut
encore se borner au cas où $Y$ est affine~: en effet, tout faisceau
d'idéaux quasi-cohérent de type fini $\mathcal{J}_i$ de
$\mathcal{O}_X|f^{-1}(U_i)$ est la restriction d'un faisceau cohérent
d'idéaux de type fini de $\mathcal{O}_X$
(\hyperref[I.9.4.7-fr]{I, 9.4.7}) et l'hypothèse c') entraîne que

\oldpage[II]{91}

$\mathcal{J}_i\otimes(\mathcal{L}^{\otimes n}|f^{-1}(U_i))$ est engendré
par ses sections (compte tenu de (\hyperref[I.9.2.2-fr]{I, 9.2.2}) et de
(\hyperref[II.3.4.7-fr]{3.4.7}))~; il suffit donc d'appliquer le critère
(\hyperref[II.4.5.5-fr]{4.5.5, d'')}).

\begin{proposition}[4.6.9]
\label{II.4.6.9-fr}
Soient $f:X\to Y$ un morphisme quasi-compact, $\mathcal{L}$ un
$\mathcal{O}_X$-Module inversible.
\begin{enumerate}
  \item[\rm{(i)}] Soit $n$ un entier $>0$. Pour que $\mathcal{L}$ soit
  $f$-ample, il faut et il suffit que $\mathcal{L}^{\otimes n}$ soit
  $f$-ample.
  \item[\rm{(ii)}] Soit $\mathcal{L}'$ un $\mathcal{O}_X$-Module
  inversible, et supposons qu'il existe un entier $n>0$ tel que
  l'homomorphisme canonique
  $\sigma:f^*(f_*(\mathcal{L}'^{\otimes n}))\to
  \mathcal{L}'^{\otimes n}$ soit surjectif. Alors, si $\mathcal{L}$ est
  $f$-ample, il en est de même de $\mathcal{L}\otimes\mathcal{L}'$.
\end{enumerate}
\end{proposition}

On se ramène en effet aussitôt au cas où $Y$ est affine, et la proposition
est alors conséquence immédiate de (\hyperref[II.4.5.6-fr]{4.5.6}).

\begin{corollary}[4.6.10]
\label{II.4.6.10-fr}
Le produit tensoriel de deux $\mathcal{O}_X$-Modules $f$-amples est
$f$-ample.
\end{corollary}

\begin{proposition}[4.6.11]
\label{II.4.6.11-fr}
Soient $Y$ un préschéma quasi-compact, $f:X\to Y$ un morphisme de type
fini, $\mathcal{L}$ un $\mathcal{O}_X$-Module inversible. Pour que
$\mathcal{L}$ soit $f$-ample, il faut et il suffit qu'il possède une des
propriétés équivalentes suivantes~:
\begin{enumerate}
  \item[\rm{d)}] Il existe $n_0>0$ tel que, pour tout entier $n\geq n_0$,
  $\mathcal{L}^{\otimes n}$ soit très ample relativement à $f$.
  \item[\rm{d')}] Il existe $n>0$ tel que $\mathcal{L}^{\otimes n}$ soit
  très ample relativement à $f$.
\end{enumerate}
\end{proposition}

Si $\mathcal{L}$ est ample relativement à $f$, il y a un recouvrement fini
$(U_i)$ de $Y$ par des ouverts affines tel que les
$\mathcal{L}|f^{-1}(U_i)$ soient amples. On en conclut
(\hyperref[II.4.5.10-fr]{4.5.10}) qu'il existe un entier $n_0$ tel que
$\mathcal{L}^{\otimes n}|f^{-1}(U_i)$ soit très ample relativement à
$f^{-1}(U_i)\to U_i$ pour tout $n\geq n_0$ et tout $i$, donc
$\mathcal{L}^{\otimes n}$ est très ample relativement à $f$
(\hyperref[II.4.4.5-fr]{4.4.5}). Inversement, d') entraîne déjà que
$\mathcal{L}^{\otimes n}$ est $f$-ample
(\hyperref[II.4.6.2-fr]{4.6.2}), donc il en est de même de $\mathcal{L}$
(\hyperref[II.4.6.9-fr]{4.6.9, (i)}).

\begin{corollary}[4.6.12]
\label{II.4.6.12-fr}
Soient $Y$ un préschéma quasi-compact, $f:X\to Y$ un morphisme de type
fini, $\mathcal{L}$, $\mathcal{L}'$ deux $\mathcal{O}_X$-Modules
inversibles. Si $\mathcal{L}$ est $f$-ample, il existe $n_0$ tel que
$\mathcal{L}^{\otimes n}\otimes\mathcal{L}'$ soit très ample relativement
à $f$ pour tout $n\geq n_0$.
\end{corollary}

On raisonne comme dans (\hyperref[II.4.6.11-fr]{4.6.11}) en utilisant un
recouvrement ouvert affine fini de $Y$ et
(\hyperref[II.4.5.11-fr]{4.5.11}).

\begin{proposition}[4.6.13]
\label{II.4.6.13-fr}
\begin{enumerate}
  \item[\rm{(i)}] Pour tout préschéma $Y$, tout
  $\mathcal{O}_Y$-Module inversible $\mathcal{L}$ est ample relativement
  au morphisme identique $1_Y$.
  \item[\rm{(i bis)}] Soient $f:X\to Y$ un morphisme quasi-compact,
  $j:X'\to X$ un morphisme quasi-compact qui est un homéomorphisme de
  l'espace sous-jacent à $X'$ sur un sous-espace de $X$. Si $\mathcal{L}$
  est un $\mathcal{O}_X$-Module ample relativement à $f$,
  $j^*(\mathcal{L})$ est ample relativement à $f\circ j$.
  \item[\rm{(ii)}] Soient $Z$ un préschéma quasi-compact, $f:X\to Y$,
  $g:Y\to Z$ deux morphismes quasi-compacts, $\mathcal{L}$ un
  $\mathcal{O}_X$-Module ample relativement à $f$, $\mathcal{K}$ un
  $\mathcal{O}_Y$-Module ample relativement à $g$. Alors il existe un
  entier $n_0>0$ tel que
  $\mathcal{L}\otimes f^*(\mathcal{K}^{\otimes n})$ soit ample relativement
  à $g\circ f$ pour tout $n\geq n_0$.
  \item[\rm{(iii)}] Soient $f:X\to Y$ un morphisme quasi-compact,
  $g:Y'\to Y$ un morphisme, et posons $X'=X_{(Y')}$. Si $\mathcal{L}$ est
  un $\mathcal{O}_X$-Module ample relativement à $f$,
  $\mathcal{L}'=\mathcal{L}\otimes_Y\mathcal{O}_{Y'}$ est un
  $\mathcal{O}_{X'}$-Module ample relativement à $f_{(Y')}$.
  \item[\rm{(iv)}] Soient $f_i:X_i\to Y_i$ ($i=1,2$) deux
  $S$-morphismes quasi-compacts. Si $\mathcal{L}_i$ est un
  $\mathcal{O}_{X_i}$-Module ample relativement à $f_i$ ($i=1,2$),
  $\mathcal{L}_1\otimes_S\mathcal{L}_2$ est ample relativement à
  $f_1\times_S f_2$.
  \item[\rm{(v)}] Soient $f:X\to Y$, $g:Y\to Z$ deux morphismes tels que
  $g\circ f$ soit quasi-compact. Si un $\mathcal{O}_X$-Module
  $\mathcal{L}$ est ample relativement à $g\circ f$, et si $g$ est séparé
  ou l'espace sous-jacent à $X$ localement noethérien, alors $\mathcal{L}$
  est ample relativement à $f$.

  \oldpage[II]{92}

  \item[\rm{(vi)}] Soient $f:X\to Y$ un morphisme quasi-compact, $j$
  l'injection canonique $X_{\mathrm{red}}\to X$. Si $\mathcal{L}$ est un
  $\mathcal{O}_X$-Module ample relativement à $f$, alors
  $j^*(\mathcal{L})$ est ample relativement à $f_{\mathrm{red}}$.
\end{enumerate}
\end{proposition}

Notons d'abord que (v) et (vi) se dérivent de (i), (i bis) et (iv) par le
même raisonnement que dans (\hyperref[II.4.4.10-fr]{4.4.10}), en utilisant
(\hyperref[II.4.6.4-fr]{4.6.4}) au lieu de
(\hyperref[II.4.4.5-fr]{4.4.5})~; nous laissons le détail du raisonnement au
lecteur. L'assertion (i) est trivialement conséquence de
(\hyperref[II.4.4.10-fr]{4.4.10, (i)}) et de
(\hyperref[II.4.6.2-fr]{4.6.2}). Pour démontrer (i bis), (iii) et (iv), nous
utiliserons le lemme suivant~:

\begin{lemma}[4.6.13.1]
\label{II.4.6.13.1-fr}
\begin{enumerate}
  \item[\rm{(i)}] Soient $u:Z\to S$ un morphisme, $\mathcal{L}$ un
  $\mathcal{O}_S$-Module inversible, $s$ une section de $\mathcal{L}$
  au-dessus de $S$, $s'$ la section de $u^*(\mathcal{L})=\mathcal{L}'$
  au-dessus de $Z$ qui lui correspond canoniquement. On a alors
  $Z_{s'}=u^{-1}(S_s)$.
  \item[\rm{(ii)}] Soient $Z$, $Z'$ deux $S$-préschémas, $p$, $p'$ les
  projections de $T=Z\times_S Z'$, $\mathcal{L}$ (resp.
  $\mathcal{L}'$) un $\mathcal{O}_Z$-Module (resp. un
  $\mathcal{O}_{Z'}$-Module) inversible, $t$ (resp. $t'$) une section de
  $\mathcal{L}$ (resp. $\mathcal{L}'$) au-dessus de $Z$ (resp. $Z'$), $s$
  (resp. $s'$) la section de $p^*(\mathcal{L})$ (resp.
  $p'^*(\mathcal{L}')$) au-dessus de $Z\times_S Z'$ qui lui correspond.
  Alors on a $T_{s\otimes s'}=Z_t\times_S Z'_{t'}$.
\end{enumerate}
\end{lemma}

Il résulte des définitions que l'on peut se ramener au cas où tous les
préschémas considérés sont affines. En outre, dans (i), on peut supposer
$\mathcal{L}=\mathcal{O}_S$~; l'assertion (i) résulte alors aussitôt de
(\hyperref[I.1.2.2.2-fr]{I, 1.2.2.2}). De même, dans (ii), on peut se
limiter au cas où $\mathcal{L}=\mathcal{O}_Z$,
$\mathcal{L}'=\mathcal{O}_{Z'}$, et alors l'assertion se réduit au lemme
(\hyperref[II.4.3.2.4-fr]{4.3.2.4}).

Démontrons alors (i bis). On peut supposer $Y$ affine
(\hyperref[II.4.6.4-fr]{4.6.4}), donc $\mathcal{L}$ ample
(\hyperref[II.4.6.6-fr]{4.6.6})~; lorsque $s$ parcourt l'ensemble réunion
des $\Gamma(X,\mathcal{L}^{\otimes n})$ ($n>0$), les $X_s$ forment une
base de la topologie de $X$ (\hyperref[II.4.5.2-fr]{4.5.2, a)}), donc par
hypothèse les $j^{-1}(X_s)$ forment une base de la topologie de $X'$~; on
conclut donc par le lemme
(\hyperref[II.4.6.13.1-fr]{4.6.13.1, (i)}) et par
(\hyperref[II.4.5.2-fr]{4.5.2, a)}) que $j^*(\mathcal{L})$ est ample.

Démontrons ensuite (iii). On peut de nouveau supposer $Y$ et $Y'$ affines
(\hyperref[II.4.6.4-fr]{4.6.4}), d'où il résulte que la projection
$h:X'\to X$ est affine (\hyperref[II.1.5.5-fr]{1.5.5}). Comme
$\mathcal{L}$ est ample (\hyperref[II.4.6.6-fr]{4.6.6}), lorsque $s$
parcourt l'ensemble des sections des $\mathcal{L}^{\otimes n}$ ($n>0$)
au-dessus de $X$ telles que $X_s$ soit affine, les $X_s$ recouvrent $X$
(\hyperref[II.4.5.2-fr]{4.5.2, a')})~; donc les $h^{-1}(X_s)$ sont affines
(\hyperref[II.1.2.5-fr]{1.2.5}) et recouvrent $X'$~; il résulte donc encore
du lemme (\hyperref[II.4.6.13.1-fr]{4.6.13.1, (i)}) et de
(\hyperref[II.4.5.2-fr]{4.5.2, a')}) que $\mathcal{L}'$ est ample, le
morphisme $f_{(Y')}$ étant quasi-compact
(\hyperref[I.6.6.4-fr]{I, 6.6.4, (iii)}).

Pour prouver (iv), notons d'abord que $f_1\times_S f_2$ est quasi-compact
(\hyperref[I.6.6.4-fr]{I, 6.6.4, (iv)}). On peut en outre supposer $S$,
$Y_1$ et $Y_2$ affines ((\hyperref[II.4.6.4-fr]{4.6.4}) et
(\hyperref[I.3.2.7-fr]{I, 3.2.7})), donc $\mathcal{L}_i$ ample ($i=1,2$)
(\hyperref[II.4.6.6-fr]{4.6.6}). Les ouverts
$(X_1)_{s_1}\times_S(X_2)_{s_2}$ forment un recouvrement de
$X_1\times_S X_2$ lorsque $s_i$ parcourt les sections des
$\mathcal{L}_i^{\otimes n_i}$ telles que $(X_i)_{s_i}$ soit affine
(\hyperref[II.4.5.2-fr]{4.5.2, a')}). D'ailleurs, en remplaçant $s_1$ et
$s_2$ par des puissances convenables, ce qui ne change pas les
$(X_i)_{s_i}$, on peut toujours supposer $n_1=n_2$. On déduit alors de
(\hyperref[II.4.6.13.1-fr]{4.6.13.1, (ii)}) et de
(\hyperref[II.4.5.2-fr]{4.5.2, a')}) que
$\mathcal{L}_1\otimes_S\mathcal{L}_2$ est ample, d'où l'assertion, puisque
$Y_1\times_S Y_2$ est affine (\hyperref[II.4.6.6-fr]{4.6.6}).

Reste à prouver (ii). Par le même raisonnement que dans
(\hyperref[II.4.4.10-fr]{4.4.10}), mais utilisant ici
(\hyperref[II.4.6.4-fr]{4.6.4}), on peut se borner au cas où $Z$ est
affine. Comme $\mathcal{K}$ est alors ample, et $Y$ quasi-compact, il existe
un nombre fini de sections $s_i\in\Gamma(Y,\mathcal{K}^{\otimes k_i})$
telles que les $Y_{s_i}$

\oldpage[II]{93}

soient affines et recouvrent $Y$
(\hyperref[II.4.5.2-fr]{4.5.2, a')})~;
remplaçant les $s_i$ par des puissances convenables, on peut en outre
supposer tous les $k_i$ égaux à un même entier $k$. Soient $s'_i$ les
sections de $f^*(\mathcal{K}^{\otimes k})$ au-dessus de $X$ correspondant
canoniquement aux $s_i$, de sorte que les
$X_{s'_i}=f^{-1}(Y_{s_i})$ (4.6.1.13, (i))
recouvrent $X$. Comme $\mathcal{L}|X_{s'_i}$ est ample
((\hyperref[II.4.6.4-fr]{4.6.4}) et
(\hyperref[II.4.6.6-fr]{4.6.6})), il existe pour chaque $i$ des sections en
nombre fini $t_{ij}\in\Gamma(X,\mathcal{L}^{\otimes n_{ij}})$ tels que les
$X_{t_{ij}}$ soient affines, contenus dans $X_{s'_i}$ et recouvrent
$X_{s'_i}$ (\hyperref[II.4.5.2-fr]{4.5.2, a')})~; on peut d'ailleurs
supposer tous les $n_{ij}$ égaux à un même nombre $n$. Cela étant, $X$ est
séparé et quasi-compact, donc il existe un entier $m>0$, et pour tout
$(i,j)$ une section
\[
  u_{ij}\in\Gamma(X,\mathcal{L}^{\otimes n}\otimes
  f^*(\mathcal{K}^{\otimes mk}))
\]
tels que $t_{ij}\otimes s_i'^{\otimes m}$ soit la restriction à $X_{s'_i}$
de $u_{ij}$ (\hyperref[I.9.3.1-fr]{I, 9.3.1})~; en outre on a
$X_{u_{ij}}=X_{t_{ij}}$, donc les $X_{u_{ij}}$ sont affines et recouvrent
$X$. On peut d'ailleurs supposer que $m$ est de la forme $nr$~; si on pose
$n_0=rk$, on voit donc (\hyperref[II.4.5.2-fr]{4.5.2, a')}) que
$\mathcal{L}\otimes_{\mathcal{O}_X}f^*(\mathcal{K}^{\otimes n_0})$ est
ample. En outre,
il existe $h_0>0$ tel que $\mathcal{K}^{\otimes h}$ soit engendré par ses
sections au-dessus de $Y$ dès que $h\geq h_0$
(\hyperref[II.4.5.5-fr]{4.5.5})~; a fortiori,
$f^*(\mathcal{K}^{\otimes h})$ est engendré par ses sections au-dessus de
$X$ pour $h\geq h_0$, par définition des images réciproques
((\hyperref[0.3.7.1-fr]{0, 3.7.1}) et
(\hyperref[II.4.4.1-fr]{4.4.1})). On en conclut que
$\mathcal{L}\otimes f^*(\mathcal{K}^{\otimes n_0+h})$ est ample dès que
$h\geq h_0$ (\hyperref[II.4.5.6-fr]{4.5.6}), ce qui achève la
démonstration.

\begin{remark}[4.6.14]
\label{II.4.6.14-fr}
Sous les conditions de (ii), on se gardera de croire que
$\mathcal{L}\otimes f^*(\mathcal{K})$ soit ample pour $g\circ f$~; en effet,
comme $\mathcal{L}\otimes f^*(\mathcal{K}^{-1})$ est aussi ample pour $f$
(\hyperref[II.4.6.5-fr]{4.6.5}), on en conclurait que $\mathcal{L}$ serait
ample pour $g\circ f$~; prenant en particulier pour $g$ le morphisme
identique, tout $\mathcal{O}_X$-Module inversible serait ample pour $f$, ce
qui n'est pas le cas en général (voir (\hyperref[II.5.1.6-fr]{5.1.6}),
(\hyperref[II.5.3.4-fr]{5.3.4, (i)}) et
(\hyperref[II.5.3.1-fr]{5.3.1})).
\end{remark}

\begin{proposition}[4.6.15]
\label{II.4.6.15-fr}
Soient $f:X\to Y$ un morphisme quasi-compact, $\mathcal{J}$ un faisceau
quasi-cohérent d'idéaux localement nilpotent de $\mathcal{O}_X$, $Z$ le
sous-préschéma fermé de $X$ défini par $\mathcal{J}$, $j:Z\to X$
l'injection canonique. Pour qu'un $\mathcal{O}_X$-Module inversible
$\mathcal{L}$ soit ample pour $f$, il faut et il suffit que
$j^*(\mathcal{L})$ soit ample pour $f\circ j$.
\end{proposition}

En effet, la question étant locale sur $Y$
(\hyperref[II.4.6.4-fr]{4.6.4}), on peut supposer $Y$ affine~; $X$ étant
alors quasi-compact, on peut supposer $\mathcal{J}$ nilpotent. Compte tenu
de (\hyperref[II.4.6.6-fr]{4.6.6}), la proposition n'est autre alors que
(\hyperref[II.4.5.13-fr]{4.5.13}).

\begin{corollary}[4.6.16]
\label{II.4.6.16-fr}
Soient $X$ un préschéma localement noethérien, $f:X\to Y$ un morphisme
quasi-compact. Pour qu'un $\mathcal{O}_X$-Module inversible $\mathcal{L}$
soit ample pour $f$, il faut et il suffit que son image réciproque
$\mathcal{L}'$ par l'injection canonique $X_{\mathrm{red}}\to X$ soit ample
pour $f_{\mathrm{red}}$.
\end{corollary}

On a déjà vu que la condition est nécessaire
(\hyperref[II.4.6.13-fr]{4.6.13, (vi)})~; inversement, si elle est remplie,
on peut se borner, pour prouver que $\mathcal{L}$ est ample pour $f$, au
cas où $Y$ est affine (\hyperref[II.4.6.4-fr]{4.6.4})~; alors
$Y_{\mathrm{red}}$ est aussi affine, donc $\mathcal{L}'$ est ample
(\hyperref[II.4.6.6-fr]{4.6.6}), et il en est de même de $\mathcal{L}$ en
vertu de (\hyperref[II.4.5.13-fr]{4.5.13}), puisque alors $X$ est
noethérien et $X_{\mathrm{red}}$ un sous-préschéma fermé de $X$ défini par
un faisceau quasi-cohérent nilpotent d'idéaux
(\hyperref[I.6.1.6-fr]{I, 6.1.6}).

\begin{proposition}[4.6.17]
\label{II.4.6.17-fr}
Les notations et hypothèses étant celles de
(\hyperref[II.4.4.11-fr]{4.4.11}), pour que $\mathcal{L}''$ soit ample
relativement à $f''$, il faut et il suffit que $\mathcal{L}$ soit ample
relativement à $f$ et $\mathcal{L}'$ ample relativement à $f'$.
\end{proposition}

\oldpage[II]{94}

La nécessité de la condition résulte de
(\hyperref[II.4.6.13-fr]{4.6.13, (i bis)}). Pour voir que la condition est
suffisante, on peut se borner au cas où $Y$ est affine, et alors le fait
que $\mathcal{L}''$ est ample résulte du critère
(\hyperref[II.4.5.2-fr]{4.5.2, a)}) appliqué à $\mathcal{L}$,
$\mathcal{L}'$ et $\mathcal{L}''$, en observant qu'une section de
$\mathcal{L}$ au-dessus de $X$ se prolonge (par $0$) en une section de
$\mathcal{L}''$ au-dessus de $X''$.

\begin{proposition}[4.6.18]
\label{II.4.6.18-fr}
Soient $Y$ un préschéma quasi-compact, $\mathcal{S}$ une
$\mathcal{O}_Y$-Algèbre graduée quasi-cohérente de type fini,
$X=\operatorname{Proj}(\mathcal{S})$, $f:X\to Y$ le morphisme structural.
Alors $f$ est de type fini, et il existe un entier $d>0$ tel que
$\mathcal{O}_X(d)$ soit inversible et $f$-ample.
\end{proposition}

En vertu de (\hyperref[II.3.1.10-fr]{3.1.10}), il existe un entier $d>0$
tel que $\mathcal{S}^{(d)}$ soit engendrée par $\mathcal{S}_d$. On sait que,
dans l'isomorphisme canonique entre $X$ et
$X^{(d)}=\operatorname{Proj}(\mathcal{S}^{(d)})$, $\mathcal{O}_X(d)$
s'identifie à $\mathcal{O}_{X^{(d)}}(1)$
(\hyperref[II.3.2.9-fr]{3.2.9, (ii)}). On voit donc qu'on est ramené au cas
où $\mathcal{S}$ est engendrée par $\mathcal{S}_1$~; la proposition résulte
alors de (\hyperref[II.4.4.3-fr]{4.4.3}) et de
(\hyperref[II.4.6.2-fr]{4.6.2}) (compte tenu de ce que $f$ est un morphisme
de type fini (\hyperref[II.3.4.1-fr]{3.4.1})).

\section{Morphismes quasi-affines~; morphismes quasi-projectifs. Morphismes
propres~; morphismes projectifs}
\label{section:II.5-fr}

\subsection{Morphismes quasi-affines.}
\label{subsection:II.5.1-fr}

\begin{definition}[5.1.1]
\label{II.5.1.1-fr}
On appelle \emph{schéma quasi-affine} un schéma isomorphe à un sous-schéma
induit sur un ouvert quasi-compact d'un schéma affine. On dit qu'un morphisme
$f:X\to Y$ est \emph{quasi-affine}, ou encore que $X$ (considéré comme
$Y$-préschéma au moyen de $f$) est un \emph{$Y$-schéma quasi-affine}, s'il
existe un recouvrement $(U_\alpha)$ de $Y$ par des ouverts affines tels que
les $f^{-1}(U_\alpha)$ soient des schémas quasi-affines.
\end{definition}

Il est clair qu'un morphisme quasi-affine est \emph{séparé}
(\hyperref[I.5.5.5-fr]{I, 5.5.5} et \hyperref[I.5.5.8-fr]{5.5.8}) et
\emph{quasi-compact} (\hyperref[I.6.6.1-fr]{I, 6.6.1})~; tout morphisme
affine est évidemment quasi-affine.

Rappelons que pour tout préschéma $X$, si on pose
$A=\Gamma(X,\mathcal{O}_X)$, l'homomorphisme identique
$A\to A=\Gamma(X,\mathcal{O}_X)$ définit un morphisme
$X\to\operatorname{Spec}(A)$ dit \emph{canonique}
(\hyperref[I.2.2.4-fr]{I, 2.2.4})~; ce n'est autre d'ailleurs que le morphisme
canonique défini dans (\hyperref[II.4.5.1-fr]{4.5.1}) pour le cas particulier
où $\mathcal{L}=\mathcal{O}_X$, si on se souvient que
$\operatorname{Proj}(A[T])$ s'identifie canoniquement à
$\operatorname{Spec}(A)$ (\hyperref[II.3.1.7-fr]{3.1.7}).

\begin{proposition}[5.1.2]
\label{II.5.1.2-fr}
Soient $X$ un schéma quasi-compact ou un préschéma dont l'espace sous-jacent
est noethérien, $A$ l'anneau $\Gamma(X,\mathcal{O}_X)$. Les conditions
suivantes sont équivalentes~:
\begin{enumerate}
  \item[\rm{a)}] $X$ est un schéma quasi-affine.
  \item[\rm{b)}] Le morphisme canonique
  $u:X\to\operatorname{Spec}(A)$ est une immersion ouverte.
  \item[\rm{b')}] Le morphisme canonique
  $u:X\to\operatorname{Spec}(A)$ est un homéomorphisme de $X$ sur un
  sous-espace de l'espace sous-jacent à $\operatorname{Spec}(A)$.
  \item[\rm{c)}] Le $\mathcal{O}_X$-Module $\mathcal{O}_X$ est très ample
  relativement à $u$ (\hyperref[II.4.4.2-fr]{4.4.2}).
  \item[\rm{c')}] Le $\mathcal{O}_X$-Module $\mathcal{O}_X$ est ample
  (\hyperref[II.4.5.1-fr]{4.5.1}).
  \item[\rm{d)}] Lorsque $f$ parcourt $A$, les $X_f$ forment une base de la
  topologie de $X$.
  \item[\rm{d')}] Lorsque $f$ parcourt $A$, ceux des $X_f$ qui sont affines
  forment un recouvrement de $X$.

\oldpage[II]{95}

  \item[\rm{e)}] Tout $\mathcal{O}_X$-Module quasi-cohérent est engendré par
  ses sections au-dessus de $X$.
  \item[\rm{e')}] Tout faisceau quasi-cohérent d'idéaux de type fini de
  $\mathcal{O}_X$ est engendré par ses sections au-dessus de $X$.
\end{enumerate}
\end{proposition}

Il est clair que b) entraîne a), et a) entraîne c) en vertu du critère b) de
(\hyperref[II.4.4.4-fr]{4.4.4}) appliqué au morphisme identique (compte tenu
de la remarque précédant l'énoncé)~; d'autre part c) entraîne c')
(\hyperref[II.4.5.10-fr]{4.5.10, (i)}), et c'), b) et b') sont équivalents par
(\hyperref[II.4.5.2-fr]{4.5.2, critères b) et b')}). Enfin c') équivaut à
chacun des critères d), d'), e), e') en vertu de
(\hyperref[II.4.5.2-fr]{4.5.2, critères a), a'), c)}) et
(\hyperref[II.4.5.5-fr]{4.5.5, critère d'')}).

On observera en outre qu'avec les notations précédentes, ceux des $X_f$ qui
sont affines forment une \emph{base} de la topologie de $X$, et que le
morphisme canonique $u$ est \emph{dominant}
(\hyperref[II.4.5.2-fr]{4.5.2}).

\begin{corollary}[5.1.3]
\label{II.5.1.3-fr}
Soit $X$ un préschéma quasi-compact. S'il existe un morphisme $v:X\to Y$ de
$X$ dans un schéma affine $Y$, qui soit un homéomorphisme de $X$ sur un
sous-espace ouvert de $Y$, alors $X$ est quasi-affine.
\end{corollary}

En effet, il existe une famille $(g_\alpha)$ de sections de $\mathcal{O}_Y$
au-dessus de $Y$ telles que les $D(g_\alpha)$ forment une base de la
topologie de $v(X)$~; si $v=(\psi,\theta)$ et si l'on pose
$f_\alpha=\theta(g_\alpha)$, on a
$X_{f_\alpha}=\psi^{-1}(D(g_\alpha))$
(\hyperref[I.2.2.4.1-fr]{I, 2.2.4.1}), donc les $X_{f_\alpha}$ forment une
base de la topologie de $X$, et le critère d) de
(\hyperref[II.5.1.2-fr]{5.1.2}) est vérifié.

\begin{corollary}[5.1.4]
\label{II.5.1.4-fr}
Si $X$ est un schéma quasi-affine, tout $\mathcal{O}_X$-Module inversible est
très ample (relativement au morphisme canonique) et a fortiori ample.
\end{corollary}

En effet, un tel Module $\mathcal{L}$ est engendré par ses sections au-dessus
de $X$ (\hyperref[II.5.1.2-fr]{5.1.2, e)}), donc
$\mathcal{L}\otimes\mathcal{O}_X=\mathcal{L}$ est très ample
(\hyperref[II.4.4.8-fr]{4.4.8}).

\begin{corollary}[5.1.5]
\label{II.5.1.5-fr}
Soit $X$ un préschéma quasi-compact. S'il existe un $\mathcal{O}_X$-Module
inversible $\mathcal{L}$ tel que $\mathcal{L}$ et $\mathcal{L}^{-1}$ soient
amples, $X$ est un schéma quasi-affine.
\end{corollary}

En effet, $\mathcal{O}_X=\mathcal{L}\otimes\mathcal{L}^{-1}$ est alors ample
(\hyperref[II.4.5.7-fr]{4.5.7}).

\begin{proposition}[5.1.6]
\label{II.5.1.6-fr}
Soit $f:X\to Y$ un morphisme quasi-compact. Les conditions suivantes sont
équivalentes~:
\begin{enumerate}
  \item[\rm{a)}] $f$ est quasi-affine.
  \item[\rm{b)}] La $\mathcal{O}_Y$-Algèbre
  $f_*(\mathcal{O}_X)=\mathcal{A}(X)$ est quasi-cohérente et le morphisme
  canonique $X\to\operatorname{Spec}(\mathcal{A}(X))$ correspondant à
  l'homomorphisme identique $\mathcal{A}(X)\to\mathcal{A}(X)$
  (\hyperref[II.1.2.7-fr]{1.2.7}) est une immersion ouverte.
  \item[\rm{b')}] La $\mathcal{O}_Y$-Algèbre $\mathcal{A}(X)$ est
  quasi-cohérente et le morphisme canonique
  $X\to\operatorname{Spec}(\mathcal{A}(X))$ est un homéomorphisme de $X$
  sur un sous-espace de $\operatorname{Spec}(\mathcal{A}(X))$.
  \item[\rm{c)}] Le $\mathcal{O}_X$-Module $\mathcal{O}_X$ est très ample
  pour $f$.
  \item[\rm{c')}] Le $\mathcal{O}_X$-Module $\mathcal{O}_X$ est ample pour
  $f$.
  \item[\rm{d)}] Le morphisme $f$ est séparé et, pour tout
  $\mathcal{O}_X$-Module quasi-cohérent $\mathcal{F}$, l'homomorphisme
  canonique $\sigma:f^*(f_*(\mathcal{F}))\to\mathcal{F}$
  (\hyperref[0.4.4.3-fr]{0, 4.4.3}) est surjectif.
\end{enumerate}
En outre, lorsque $f$ est quasi-affine, tout $\mathcal{O}_X$-Module
inversible $\mathcal{L}$ est très ample relativement à $f$.
\end{proposition}

L'équivalence de a) et c') résulte du caractère local sur $Y$ de la
$f$-amplitude (\hyperref[II.4.6.4-fr]{4.6.4}), de la définition
(\hyperref[II.5.1.1-fr]{5.1.1}) et du critère
(\hyperref[II.5.1.2-fr]{5.1.2, c')}). Les autres propriétés sont locales sur
$Y$

\oldpage[II]{96}

et résultent donc aussitôt de (\hyperref[II.5.1.2-fr]{5.1.2}) et
(\hyperref[II.5.1.4-fr]{5.1.4}), compte tenu de ce que
$f_*(\mathcal{F})$ est quasi-cohérent lorsque $f$ est séparé
(\hyperref[I.9.2.2-fr]{I, 9.2.2, a)}).

\begin{corollary}[5.1.7]
\label{II.5.1.7-fr}
Soit $f:X\to Y$ un morphisme quasi-affine. Pour tout ouvert $U$ de $Y$, la
restriction $f^{-1}(U)\to U$ de $f$ est quasi-affine.
\end{corollary}

\begin{corollary}[5.1.8]
\label{II.5.1.8-fr}
Soient $Y$ un schéma affine, $f:X\to Y$ un morphisme quasi-compact. Pour que
$f$ soit quasi-affine, il faut et il suffit que $X$ soit un schéma
quasi-affine.
\end{corollary}

C'est une conséquence immédiate de (\hyperref[II.5.1.6-fr]{5.1.6}) et
(\hyperref[II.4.6.6-fr]{4.6.6}).

\begin{corollary}[5.1.9]
\label{II.5.1.9-fr}
Soient $Y$ un schéma quasi-compact ou un préschéma dont l'espace sous-jacent
est noethérien, $f:X\to Y$ un morphisme de type fini. Si $f$ est
quasi-affine, il existe une sous-$\mathcal{O}_Y$-Algèbre quasi-cohérente
$\mathcal{B}$ de $\mathcal{A}(X)=f_*(\mathcal{O}_X)$, de type fini
(\hyperref[I.9.6.2-fr]{I, 9.6.2}), telle que le morphisme
$X\to\operatorname{Spec}(\mathcal{B})$ correspondant à l'injection canonique
$\mathcal{B}\to\mathcal{A}(X)$ (\hyperref[II.1.2.7-fr]{1.2.7}) soit une
immersion. En outre, toute sous-$\mathcal{O}_Y$-Algèbre quasi-cohérente de
type fini $\mathcal{B}'$ de $\mathcal{A}(X)$, contenant $\mathcal{B}$,
possède la même propriété.
\end{corollary}

En effet, $\mathcal{A}(X)$ est limite inductive de ses
sous-$\mathcal{O}_Y$-Algèbres quasi-cohérentes de type fini
(\hyperref[I.9.6.5-fr]{I, 9.6.5})~; le résultat est alors un cas particulier
de (\hyperref[II.3.8.4-fr]{3.8.4}), compte tenu de l'identification de
$\operatorname{Spec}(\mathcal{A}(X))$ et de
$\operatorname{Proj}(\mathcal{A}(X)[T])$
(\hyperref[II.3.1.7-fr]{3.1.7}).

\begin{proposition}[5.1.10]
\label{II.5.1.10-fr}
\begin{enumerate}
  \item[\rm{(i)}] Un morphisme quasi-compact $X\to Y$ qui est un
  homéomorphisme de l'espace sous-jacent à $X$ sur un sous-espace de
  l'espace sous-jacent à $Y$ (en particulier une immersion fermée) est
  quasi-affine.
  \item[\rm{(ii)}] Le composé de deux morphismes quasi-affines est
  quasi-affine.
  \item[\rm{(iii)}] Si $f:X\to Y$ est un $S$-morphisme quasi-affine,
  $f_{(S')}:X_{(S')}\to Y_{(S')}$ est un morphisme quasi-affine pour toute
  extension $S'\to S$ du préschéma de base.
  \item[\rm{(iv)}] Si $f:X\to Y$ et $g:X'\to Y'$ sont deux $S$-morphismes
  quasi-affines, $f\times_S g$ est quasi-affine.
  \item[\rm{(v)}] Si $f:X\to Y$, $g:Y\to Z$ sont deux morphismes tels que
  $g\circ f$ soit quasi-affine, et si $g$ est séparé ou l'espace sous-jacent
  à $X$ localement noethérien, $f$ est quasi-affine.
  \item[\rm{(vi)}] Si $f$ est un morphisme quasi-affine, il en est de même
  de $f_{\mathrm{red}}$.
\end{enumerate}
\end{proposition}

Compte tenu du critère (\hyperref[II.5.1.6-fr]{5.1.6, c')}), (i), (iii),
(iv), (v) et (vi) découlent aussitôt de
(\hyperref[II.4.6.13-fr]{4.6.13, (i bis), (iii), (iv), (v) et (vi)})
respectivement. Pour démontrer (ii), on peut se borner au cas où $Z$ est
affine, et alors l'assertion résulte directement de
(\hyperref[II.4.6.13-fr]{4.6.13, (ii)}), appliqué à
$\mathcal{L}=\mathcal{O}_X$ et $\mathcal{K}=\mathcal{O}_Y$.

\begin{remark}[5.1.11]
\label{II.5.1.11-fr}
Soient $f:X\to Y$, $g:Y\to Z$ deux morphismes tels que $X\times_ZY$ soit
localement noethérien. Alors l'immersion graphe
$\Gamma_f:X\to X\times_ZY$ est quasi-affine, étant quasi-compacte
(\hyperref[I.6.3.5-fr]{I, 6.3.5}), et
(\hyperref[I.5.5.12-fr]{I, 5.5.12}) montre que, dans (v), la conclusion reste
valable si on supprime l'hypothèse que $g$ est séparé.
\end{remark}

\begin{proposition}[5.1.12]
\label{II.5.1.12-fr}
Soient $f:X\to Y$ un morphisme quasi-compact, $g:X'\to X$ un morphisme
quasi-affine. Si $\mathcal{L}$ est un $\mathcal{O}_X$-Module ample pour
$f$, $g^*(\mathcal{L})$ est un $\mathcal{O}_{X'}$-Module ample pour
$f\circ g$.
\end{proposition}

En effet, comme $\mathcal{O}_{X'}$ est très ample pour $g$, et que la
question est locale sur $Y$ (\hyperref[II.4.6.4-fr]{4.6.4}), il résulte de
(\hyperref[II.4.6.13-fr]{4.6.13, (ii)}) qu'il existe (pour $Y$ affine) un
entier $n$ tel que $g^*(\mathcal{L})^{\otimes n}$ soit ample pour
$f\circ g$, donc $g^*(\mathcal{L})$ est ample pour $f\circ g$
(\hyperref[II.4.6.9-fr]{4.6.9}).

\oldpage[II]{97}

\subsection{Le critère de Serre.}
\label{subsection:II.5.2-fr}

\begin{theorem}[5.2.1]
\label{II.5.2.1-fr}
\emph{(Critère de Serre.)} Soit $X$ un schéma quasi-compact ou un préschéma
dont l'espace sous-jacent est noethérien. Les conditions suivantes sont
équivalentes~:
\begin{enumerate}
  \item[\rm{a)}] $X$ est un schéma affine.
  \item[\rm{b)}] Il existe une famille d'éléments
  $f_\alpha\in A=\Gamma(X,\mathcal{O}_X)$ tels que les $X_{f_\alpha}$ soient
  affines et que l'idéal engendré par les $f_\alpha$ dans $A$ soit égal à
  $A$.
  \item[\rm{c)}] Le foncteur $\Gamma(X,\mathcal{F})$ est exact en
  $\mathcal{F}$ dans la catégorie des $\mathcal{O}_X$-Modules
  quasi-cohérents~; autrement dit, si
  \[
    0\to\mathcal{F}'\to\mathcal{F}\to\mathcal{F}''\to0 \tag{*}
  \]
  est une suite exacte de $\mathcal{O}_X$-Modules quasi-cohérents, la suite
  \[
    0\to\Gamma(X,\mathcal{F}')\to\Gamma(X,\mathcal{F})
      \to\Gamma(X,\mathcal{F}'')\to0
  \]
  est exacte.
  \item[\rm{c')}] La propriété c) a lieu pour toute suite exacte (*) de
  $\mathcal{O}_X$-Modules quasi-cohérents telle que $\mathcal{F}$ soit
  isomorphe à un sous-$\mathcal{O}_X$-Module d'un produit fini
  $\mathcal{O}_X^n$.
  \item[\rm{d)}] $\mathrm{H}^1(X,\mathcal{F})=0$ pour tout
  $\mathcal{O}_X$-Module quasi-cohérent $\mathcal{F}$.
  \item[\rm{d')}] $\mathrm{H}^1(X,\mathcal{J})=0$ pour tout faisceau
  quasi-cohérent d'idéaux $\mathcal{J}$ de $\mathcal{O}_X$.
\end{enumerate}
\end{theorem}

Il est évident que a) implique b)~; b) implique d'autre part que les
$X_{f_\alpha}$ recouvrent $X$, puisque par hypothèse la section $1$ est
combinaison linéaire des $f_\alpha$ et que les $D(f_\alpha)$ recouvrent
alors $\operatorname{Spec}(A)$. La dernière assertion de
(\hyperref[II.4.5.2-fr]{4.5.2}) implique donc que
$X\to\operatorname{Spec}(A)$ est un isomorphisme.

On sait que a) implique c) (\hyperref[I.1.3.11-fr]{I, 1.3.11}) et c)
entraîne trivialement c'). Prouvons que c') implique b). Tout d'abord, c')
entraîne que, pour tout point fermé $x\in X$ et tout voisinage ouvert $U$ de
$x$, il existe $f\in A$ tel que $x\in X_f\subset X-U$. Soit en effet
$\mathcal{J}$ (resp. $\mathcal{J}'$) le faisceau quasi-cohérent d'idéaux de
$\mathcal{O}_X$ définissant le sous-préschéma fermé réduit de $X$ ayant pour
espace sous-jacent $X-U$ (resp. $(X-U)\cup\{x\}$)
(\hyperref[I.5.2.1-fr]{I, 5.2.1})~; il est clair que l'on a
$\mathcal{J}'\subset\mathcal{J}$ et que
$\mathcal{J}''=\mathcal{J}/\mathcal{J}'$ est un $\mathcal{O}_X$-Module
quasi-cohérent ayant pour support $\{x\}$ et tel que
$\mathcal{J}''_x=k(x)$. L'hypothèse c') appliquée à la suite exacte
$0\to\mathcal{J}'\to\mathcal{J}\to\mathcal{J}''\to0$ montre que
$\Gamma(X,\mathcal{J})\to\Gamma(X,\mathcal{J}'')$ est surjectif. La section
de $\mathcal{J}''$ dont le germe en $x$ est $1_x$ est donc l'image d'une
section $f\in\Gamma(X,\mathcal{J})\subset\Gamma(X,\mathcal{O}_X)$, et on a
par définition $f(x)=1_x$ et $f(y)=0$ dans $X-U$, ce qui établit notre
assertion. D'ailleurs, si $U$ est affine, il en est de même de $X_f$
(\hyperref[I.1.3.6-fr]{I, 1.3.6}), donc la réunion des $X_f$ qui sont affines
($f\in A$) est un ensemble ouvert $Z$ contenant tous les points fermés de
$X$~; comme $X$ est un espace de Kolmogoroff quasi-compact, on a
nécessairement $Z=X$ (\hyperref[0.2.1.3-fr]{0, 2.1.3}). Puisque $X$ est
quasi-compact, il y a un nombre fini d'éléments $f_i\in A$
($1\leq i\leq n$) tels que les $X_{f_i}$ soient affines et recouvrent $X$.
Considérons alors l'homomorphisme $\mathcal{O}_X^n\to\mathcal{O}_X$ défini
par les sections $f_i$ (\hyperref[0.5.1.1-fr]{0, 5.1.1})~; comme, pour tout
$x\in X$, un au moins des $(f_i)_x$ est inversible, cet homomorphisme est
surjectif, et on a donc une suite exacte
$0\to\mathcal{R}\to\mathcal{O}_X^n\to\mathcal{O}_X\to0$, où $\mathcal{R}$
est un sous-$\mathcal{O}_X$-Module quasi-cohérent de $\mathcal{O}_X^n$. Il
résulte alors

\oldpage[II]{98}

de c') que l'homomorphisme correspondant
$\Gamma(X,\mathcal{O}_X^n)\to\Gamma(X,\mathcal{O}_X)$ est surjectif, ce qui
démontre b).

Enfin, a) implique d) (\hyperref[I.5.1.9.2-fr]{I, 5.1.9.2}) et d) implique
trivialement d'). Reste à montrer que d') entraîne c'). Or, si
$\mathcal{F}'$ est un sous-$\mathcal{O}_X$-Module quasi-cohérent de
$\mathcal{O}_X^n$, la filtration
$0\subset\mathcal{O}_X\subset\mathcal{O}_X^2\subset\cdots\subset
\mathcal{O}_X^n$ définit sur $\mathcal{F}'$ une filtration formée des
$\mathcal{F}'_k=\mathcal{O}_X^k\cap\mathcal{F}'$ ($0\leq k\leq n$), qui
sont des $\mathcal{O}_X$-Modules quasi-cohérents
(\hyperref[I.4.1.1-fr]{I, 4.1.1}), et
$\mathcal{F}'_{k+1}/\mathcal{F}'_k$ est isomorphe à un
sous-$\mathcal{O}_X$-Module quasi-cohérent de
$\mathcal{O}_X^{k+1}/\mathcal{O}_X^k=\mathcal{O}_X$, c'est-à-dire à un
faisceau quasi-cohérent d'idéaux de $\mathcal{O}_X$. L'hypothèse d')
entraîne donc
$\mathrm{H}^1(X,\mathcal{F}'_{k+1}/\mathcal{F}'_k)=0$~; la suite exacte de
cohomologie
$\mathrm{H}^1(X,\mathcal{F}'_k)\to\mathrm{H}^1(X,\mathcal{F}'_{k+1})
\to\mathrm{H}^1(X,\mathcal{F}'_{k+1}/\mathcal{F}'_k)=0$ permet alors de
prouver par récurrence sur $k$ que
$\mathrm{H}^1(X,\mathcal{F}'_k)=0$ pour tout $k$. C.Q.F.D.

\begin{remark}[5.2.1.1]
\label{II.5.2.1.1-fr}
Lorsque $X$ est un préschéma noethérien, on peut dans l'énoncé de c') et de
d'), remplacer «~quasi-cohérent~» par «~cohérent~». En effet, dans la
démonstration du fait que c') entraîne b), $\mathcal{J}$ et $\mathcal{J}'$
sont alors des faisceaux cohérents d'idéaux, et, d'autre part, tout
sous-Module quasi-cohérent d'un Module cohérent est cohérent
(\hyperref[I.6.1.1-fr]{I, 6.1.1})~; d'où la conclusion.
\end{remark}

\begin{corollary}[5.2.2]
\label{II.5.2.2-fr}
Soit $f:X\to Y$ un morphisme séparé quasi-compact. Les conditions suivantes
sont équivalentes~:
\begin{enumerate}
  \item[\rm{a)}] $f$ est un morphisme affine.
  \item[\rm{b)}] Le foncteur $f_*$ est exact dans la catégorie des
  $\mathcal{O}_X$-Modules quasi-cohérents.
  \item[\rm{c)}] Pour tout $\mathcal{O}_X$-Module quasi-cohérent
  $\mathcal{F}$, on a $\mathrm{R}^1f_*(\mathcal{F})=0$.
  \item[\rm{c')}] Pour tout faisceau d'idéaux quasi-cohérent $\mathcal{J}$
  de $\mathcal{O}_X$, on a $\mathrm{R}^1f_*(\mathcal{J})=0$.
\end{enumerate}
\end{corollary}

Toutes ces conditions étant locales sur $Y$, par définition du foncteur
$\mathrm{R}^1f_*$ (T, 3.7.3), on peut supposer que $Y$ est affine. Si $f$
est affine, $X$ est alors affine et la propriété b) n'est autre que
(\hyperref[I.1.6.4-fr]{I, 1.6.4}). Inversement, montrons que b) implique
a)~: pour tout $\mathcal{O}_X$-Module quasi-cohérent $\mathcal{F}$,
$f_*(\mathcal{F})$ est un $\mathcal{O}_Y$-Module quasi-cohérent
(\hyperref[I.9.2.2-fr]{I, 9.2.2, a)}). Par hypothèse le foncteur
$f_*(\mathcal{F})$ est exact en $\mathcal{F}$, et le foncteur
$\Gamma(Y,\mathcal{G})$ est exact en $\mathcal{G}$ (dans la catégorie des
$\mathcal{O}_Y$-Modules quasi-cohérents) puisque $Y$ est affine
(\hyperref[I.1.3.11-fr]{I, 1.3.11})~; donc
$\Gamma(Y,f_*(\mathcal{F}))=\Gamma(X,\mathcal{F})$ est exact en
$\mathcal{F}$, ce qui prouve notre assertion en vertu de
(\hyperref[II.5.2.1-fr]{5.2.1, c)}).

Si $f$ est affine, $f^{-1}(U)$ est affine pour tout ouvert affine $U$ de
$Y$ (\hyperref[II.1.3.2-fr]{1.3.2}), donc
$\mathrm{H}^1(f^{-1}(U),\mathcal{F})=0$
(\hyperref[II.5.2.1-fr]{5.2.1, d)}), ce qui par définition entraîne
$\mathrm{R}^1f_*(\mathcal{F})=0$. Supposons enfin vérifiée la condition
c')~; la suite exacte des termes de bas degré dans la suite spectrale de
Leray (G, II, 4.17.1 et I, 4.5.1) donne en particulier la suite exacte
\[
  0\to\mathrm{H}^1(Y,f_*(\mathcal{J}))\to
  \mathrm{H}^1(X,\mathcal{J})\to
  \mathrm{H}^0(Y,\mathrm{R}^1f_*(\mathcal{J})).
\]
Comme $Y$ est affine et $f_*(\mathcal{J})$ quasi-cohérent
(\hyperref[I.9.2.2-fr]{I, 9.2.2, a)}), on a
$\mathrm{H}^1(Y,f_*(\mathcal{J}))=0$
(\hyperref[II.5.2.1-fr]{5.2.1})~; l'hypothèse c') entraîne donc
$\mathrm{H}^1(X,\mathcal{J})=0$, et on conclut par
(\hyperref[II.5.2.1-fr]{5.2.1}) que $X$ est un schéma affine.

\begin{corollary}[5.2.3]
\label{II.5.2.3-fr}
Si $f:X\to Y$ est un morphisme affine, alors, pour tout
$\mathcal{O}_X$-Module quasi-cohérent $\mathcal{F}$, l'homomorphisme
canonique
$\mathrm{H}^1(Y,f_*(\mathcal{F}))\to\mathrm{H}^1(X,\mathcal{F})$ est
bijectif.
\end{corollary}

\oldpage[II]{99}

En effet, on a la suite exacte
\[
  0\to\mathrm{H}^1(Y,f_*(\mathcal{F}))\to
  \mathrm{H}^1(X,\mathcal{F})\to
  \mathrm{H}^0(Y,\mathrm{R}^1f_*(\mathcal{F}))
\]
des termes de bas degré de la suite spectrale de Leray, et la conclusion
résulte de (\hyperref[II.5.2.2-fr]{5.2.2}).

\begin{remark}[5.2.4]
\label{II.5.2.4-fr}
Au chapitre III, § 1, nous prouverons que si $X$ est affine, on a
$\mathrm{H}^i(X,\mathcal{F})=0$ pour tout $i>0$ et tout
$\mathcal{O}_X$-Module quasi-cohérent $\mathcal{F}$.
\end{remark}

\subsection{Morphismes quasi-projectifs.}
\label{subsection:II.5.3-fr}

\begin{definition}[5.3.1]
\label{II.5.3.1-fr}
On dit qu'un morphisme $f:X\to Y$ est \emph{quasi-projectif}, ou que $X$
(considéré comme $Y$-préschéma au moyen de $f$) est \emph{quasi-projectif
sur $Y$}, ou que $X$ est un \emph{$Y$-schéma quasi-projectif}, si $f$ est de
type fini et s'il existe un $\mathcal{O}_X$-Module inversible qui est
$f$-ample.
\end{definition}

On notera que cette notion n'est pas locale sur $Y$~: les contre-exemples de
Nagata [26] et de Hironaka montrent que, même si $X$ et $Y$ sont des schémas
algébriques non singuliers sur un corps algébriquement clos, tout point de
$Y$ peut avoir un voisinage affine $U$ tel que $f^{-1}(U)$ soit
quasi-projectif sur $U$, sans que $f$ soit quasi-projectif.

On notera qu'un morphisme quasi-projectif est nécessairement séparé
(\hyperref[II.4.6.1-fr]{4.6.1}). Lorsque $Y$ est quasi-compact, il revient au
même de dire que $f$ est quasi-projectif ou que $f$ est de type fini et qu'il
existe un $\mathcal{O}_X$-Module très ample relativement à $f$
(\hyperref[II.4.6.2-fr]{4.6.2} et \hyperref[II.4.6.11-fr]{4.6.11}). De plus~:

\begin{proposition}[5.3.2]
\label{II.5.3.2-fr}
Soit $Y$ un schéma quasi-compact ou un préschéma dont l'espace sous-jacent
est noethérien, et soit $X$ un $Y$-préschéma. Les conditions suivantes sont
équivalentes~:
\begin{enumerate}
  \item[\rm{a)}] $X$ est un $Y$-schéma quasi-projectif.
  \item[\rm{b)}] $X$ est de type fini sur $Y$, et il existe un
  $\mathcal{O}_Y$-Module quasi-cohérent de type fini $\mathcal{E}$ tel que
  $X$ soit $Y$-isomorphe à un sous-préschéma de $\mathbf{P}(\mathcal{E})$.
  \item[\rm{c)}] $X$ est de type fini sur $Y$, et il existe une
  $\mathcal{O}_Y$-Algèbre graduée quasi-cohérente $\mathcal{S}$ telle que
  $\mathcal{S}_1$ soit de type fini et engendre $\mathcal{S}$ et que $X$
  soit $Y$-isomorphe à un sous-préschéma induit sur un ouvert partout dense
  de $\operatorname{Proj}(\mathcal{S})$.
\end{enumerate}
\end{proposition}

Cela résulte aussitôt de la remarque précédente et de
(\hyperref[II.4.4.3-fr]{4.4.3}), (\hyperref[II.4.4.6-fr]{4.4.6}) et
(\hyperref[II.4.4.7-fr]{4.4.7}).

On notera que lorsque $Y$ est un préschéma noethérien, on peut, dans les
conditions b) et c) de (\hyperref[II.5.3.2-fr]{5.3.2}), supprimer
l'hypothèse que $X$ est de type fini sur $Y$, qui est alors automatiquement
vérifiée (\hyperref[I.6.3.5-fr]{I, 6.3.5}).

\begin{corollary}[5.3.3]
\label{II.5.3.3-fr}
Soit $Y$ un schéma quasi-compact tel qu'il existe un $\mathcal{O}_Y$-Module
ample $\mathcal{L}$ (\hyperref[II.4.5.3-fr]{4.5.3}). Pour qu'un $Y$-schéma
$X$ soit quasi-projectif, il faut et il suffit qu'il soit de type fini sur
$Y$ et isomorphe à un sous-$Y$-schéma d'un fibré projectif de la forme
$\mathbf{P}_Y^r$.
\end{corollary}

En effet, si $\mathcal{E}$ est un $\mathcal{O}_Y$-Module quasi-cohérent de
type fini, $\mathcal{E}$ est isomorphe à un quotient d'un
$\mathcal{O}_Y$-Module
$\mathcal{L}^{\otimes(-n)}\otimes_{\mathcal{O}_Y}\mathcal{O}_Y^k$
(\hyperref[II.4.5.5-fr]{4.5.5}), donc $\mathbf{P}(\mathcal{E})$ est
isomorphe à un sous-schéma fermé de $\mathbf{P}_Y^{k-1}$
(\hyperref[II.4.1.2-fr]{4.1.2} et \hyperref[II.4.1.4-fr]{4.1.4}).

\begin{proposition}[5.3.4]
\label{II.5.3.4-fr}
\begin{enumerate}
  \item[\rm{(i)}] Un morphisme quasi-affine de type fini (et en particulier
  une immersion quasi-compacte, ou un morphisme affine de type fini) est
  quasi-projectif.
  \item[\rm{(ii)}] Si $f:X\to Y$ et $g:Y\to Z$ sont quasi-projectifs et si
  $Z$ est quasi-compact, $g\circ f$ est quasi-projectif.

\oldpage[II]{100}

  \item[\rm{(iii)}] Si $f:X\to Y$ est un $S$-morphisme quasi-projectif,
  $f_{(S')}:X_{(S')}\to Y_{(S')}$ est quasi-projectif pour toute extension
  $S'\to S$ du préschéma de base.
  \item[\rm{(iv)}] Si $f:X\to Y$ et $g:X'\to Y'$ sont deux $S$-morphismes
  quasi-projectifs, $f\times_Sg$ est quasi-projectif.
  \item[\rm{(v)}] Si $f:X\to Y$, $g:Y\to Z$ sont deux morphismes tels que
  $g\circ f$ soit quasi-projectif, et si $g$ est séparé ou $X$ localement
  noethérien, alors $f$ est quasi-projectif.
  \item[\rm{(vi)}] Si $f$ est un morphisme quasi-projectif, il en est de
  même de $f_{\mathrm{red}}$.
\end{enumerate}
\end{proposition}

(i) résulte de (\hyperref[II.5.1.6-fr]{5.1.6}) et
(\hyperref[II.5.1.10-fr]{5.1.10, (i)}). Les autres assertions sont des
conséquences immédiates de la définition (\hyperref[II.5.3.1-fr]{5.3.1}),
des propriétés des morphismes de type fini (\hyperref[I.6.3.4-fr]{I, 6.3.4})
et de (\hyperref[II.4.6.13-fr]{4.6.13}).

\begin{remark}[5.3.5]
\label{II.5.3.5-fr}
On notera qu'il peut se faire que $f_{\mathrm{red}}$ soit quasi-projectif
sans que $f$ le soit, même en supposant que $Y$ est le spectre d'une algèbre
de rang fini sur $\mathbf{C}$ et que $f$ est propre.
\end{remark}

\begin{corollary}[5.3.6]
\label{II.5.3.6-fr}
Si $X$ et $X'$ sont deux $Y$-schémas quasi-projectifs, $X\amalg X'$ est un
$Y$-schéma quasi-projectif.
\end{corollary}

Cela résulte de (\hyperref[II.4.6.18-fr]{4.6.18}).

\subsection{Morphismes propres et morphismes universellement fermés.}
\label{subsection:II.5.4-fr}
\label{II.5.4-fr}

\begin{definition}[5.4.1]
\label{II.5.4.1-fr}
On dit qu'un morphisme de préschémas $f:X\to Y$ est \emph{propre} s'il
vérifie les deux conditions suivantes~:
\begin{enumerate}
  \item[\rm{a)}] $f$ est séparé et de type fini.
  \item[\rm{b)}] Pour tout préschéma $Y'$ et tout morphisme $Y'\to Y$, la
  projection $f_{(Y')}:X\times_YY'\to Y'$ est un morphisme fermé
  (\hyperref[I.2.2.6-fr]{I, 2.2.6}).
\end{enumerate}
Lorsqu'il en est ainsi, on dit aussi que $X$ (considéré comme $Y$-préschéma
de morphisme structural $f$) est propre au-dessus de $Y$.
\end{definition}

Il est immédiat que les conditions a) et b) sont locales sur $Y$. Pour
vérifier que l'image d'une partie fermée $Z$ de $X\times_YY'$ par la
projection $q:X\times_YY'\to Y'$ est fermée dans $Y'$, il suffit de voir que
$q(Z)\cap U'$ est fermé dans $U'$ pour tout ouvert affine $U'$ de $Y'$~;
comme $q(Z)\cap U'=q(Z\cap q^{-1}(U'))$ et que $q^{-1}(U')$ s'identifie à
$X\times_YU'$ (\hyperref[I.4.4.1-fr]{I, 4.4.1}), on voit que pour vérifier
la condition b) de la déf. (\hyperref[II.5.4.1-fr]{5.4.1}), on peut se
limiter au cas où $Y'$ est un schéma affine. On verra plus loin
(\hyperref[II.5.6.3-fr]{5.6.3}) que si $Y$ est localement noethérien, on
peut même se borner à vérifier b) lorsque $Y'$ est de type fini sur $Y$.

Il est clair que tout morphisme propre est \emph{fermé}.

\begin{proposition}[5.4.2]
\label{II.5.4.2-fr}
\begin{enumerate}
  \item[\rm{(i)}] Une immersion fermée est un morphisme propre.
  \item[\rm{(ii)}] Le composé de deux morphismes propres est propre.
  \item[\rm{(iii)}] Si $X$, $Y$ sont des $S$-préschémas,
  $f:X\to Y$ un $S$-morphisme propre, alors
  $f_{(S')}:X_{(S')}\to Y_{(S')}$ est propre pour toute extension
  $S'\to S$ du préschéma de base.
  \item[\rm{(iv)}] Si $f:X\to Y$ et $g:X'\to Y'$ sont deux $S$-morphismes
  propres, le $S$-morphisme
  $f\times_Sg:X\times_SX'\to Y\times_SY'$ est propre.
\end{enumerate}
\end{proposition}

\oldpage[II]{101}

Il suffit de démontrer (i), (ii) et (iii) (\hyperref[I.3.5.1-fr]{I, 3.5.1}).
Dans chacun de ces trois cas, la vérification de la condition a) de
(\hyperref[II.5.4.1-fr]{5.4.1}) découle de résultats antérieurs
(\hyperref[I.5.5.1-fr]{I, 5.5.1} et \hyperref[I.6.3.4-fr]{6.3.4})~; reste à
vérifier la condition b). C'est immédiat dans le cas (i), car si $X\to Y$
est une immersion fermée, il en est de même de
$X\times_YY'\to Y\times_YY'=Y'$ (\hyperref[I.4.3.2-fr]{I, 4.3.2} et
\hyperref[II.3.3.3-fr]{3.3.3}). Pour démontrer (ii), considérons deux
morphismes propres $X\to Y$, $Y\to Z$, et un morphisme $Z'\to Z$. On peut
écrire $X\times_ZZ'=X\times_Y(Y\times_ZZ')$
(\hyperref[I.3.3.9.1-fr]{I, 3.3.9.1}), et par suite la projection
$X\times_ZZ'\to Z'$ se factorise en
$X\times_Y(Y\times_ZZ')\to Y\times_ZZ'\to Z'$. Compte tenu de la remarque
du début, (ii) résulte donc de ce que le composé de deux morphismes fermés
est fermé. Enfin, pour tout morphisme $S'\to S$, $X_{(S')}$ s'identifie à
$X\times_YY_{(S')}$ (\hyperref[I.3.3.11-fr]{I, 3.3.11})~; pour tout
morphisme $Z\to Y_{(S')}$, on peut écrire
\[
  X_{(S')}\times_{Y_{(S')}}Z
  =(X\times_YY_{(S')})\times_{Y_{(S')}}Z=X\times_YZ~;
\]
$X\times_YZ\to Z$ est fermé par hypothèse, donc cela prouve (iii).

\begin{corollary}[5.4.3]
\label{II.5.4.3-fr}
Soient $f:X\to Y$, $g:Y\to Z$ deux morphismes tels que $g\circ f$ soit
propre.
\begin{enumerate}
  \item[\rm{(i)}] Si $g$ est séparé, $f$ est propre.
  \item[\rm{(ii)}] Si $g$ est séparé et de type fini, et si $f$ est
  surjectif, $g$ est propre.
\end{enumerate}
\end{corollary}

(i) découle de (\hyperref[II.5.4.2-fr]{5.4.2}) par le procédé général
(\hyperref[I.5.5.12-fr]{I, 5.5.12}). Pour démontrer (ii), il n'y a à
vérifier que la condition b) de la déf. (\hyperref[II.5.4.1-fr]{5.4.1}).
Pour tout morphisme $Z'\to Z$, le diagramme
\[
  \xymatrix{
    X\times_ZZ'\ar[r]^{f\times1_{Z'}}\ar[dr]_p &
    Y\times_ZZ'\ar[d]^{p'}\\
    & Z'
  }
\]
(où $p$ et $p'$ sont les projections) est commutatif
(\hyperref[I.3.2.1-fr]{I, 3.2.1})~; en outre, $f\times1_{Z'}$ est surjectif
puisqu'il en est ainsi de $f$ (\hyperref[I.3.5.2-fr]{I, 3.5.2}), et $p$ est
un morphisme fermé par hypothèse. Toute partie fermée $F$ de
$Y\times_ZZ'$ est alors l'image par $f\times1_{Z'}$ d'une partie fermée $E$
de $X\times_ZZ'$, donc $p'(F)=p(E)$ est fermé dans $Z'$ par hypothèse, d'où
le corollaire.

\begin{corollary}[5.4.4]
\label{II.5.4.4-fr}
Si $X$ est un préschéma propre au-dessus de $Y$, et $\mathcal{S}$ une
$\mathcal{O}_Y$-Algèbre quasi-cohérente, tout $Y$-morphisme
$f:X\to\operatorname{Proj}(\mathcal{S})$ est propre (et a fortiori fermé).
\end{corollary}

En effet, le morphisme structural
$p:\operatorname{Proj}(\mathcal{S})\to Y$ est séparé, et $p\circ f$ est
propre par hypothèse.

\begin{corollary}[5.4.5]
\label{II.5.4.5-fr}
Soit $f:X\to Y$ un morphisme séparé de type fini. Soient
$(X_i)_{1\leq i\leq n}$ (resp. $(Y_i)_{1\leq i\leq n}$) une famille finie
de sous-préschémas fermés de $X$ (resp. $Y$), $j_i$ (resp. $h_i$)
l'injection canonique $X_i\to X$ (resp. $Y_i\to Y$). Supposons que l'espace
sous-jacent $X$ soit réunion des $X_i$, et que pour tout $i$, il y ait un
morphisme $f_i:X_i\to Y_i$ tel que le diagramme
\[
  \xymatrix{
    X_i\ar[r]^{f_i}\ar[d]_{j_i} & Y_i\ar[d]^{h_i}\\
    X\ar[r]_f & Y
  }
\]
soit commutatif. Alors, pour que $f$ soit propre, il faut et il suffit que
chacun des $f_i$ le soit.
\end{corollary}

\oldpage[II]{102}

Si $f$ est propre, il en est de même de $f\circ j_i$, puisque $j_i$ est une
immersion fermée (\hyperref[II.5.4.2-fr]{5.4.2})~; comme $h_i$ est une
immersion fermée, donc un morphisme séparé, $f_i$ est propre par
(\hyperref[II.5.4.3-fr]{5.4.3}). Supposons inversement que chacun des $f_i$
soit propre, et considérons le préschéma $Z$ somme des $X_i$, soit $u$ le
morphisme $Z\to X$ qui se réduit à $j_i$ dans chacun des $X_i$. La
restriction de $f\circ u$ à chaque $X_i$ est égale à
$f\circ j_i=h_i\circ f_i$, donc est propre puisqu'il en est ainsi de $h_i$
et de $f_i$ (\hyperref[II.5.4.2-fr]{5.4.2})~; il résulte alors aussitôt de
la déf. (\hyperref[II.5.4.1-fr]{5.4.1}) que $f\circ u$ est propre. Mais
comme $u$ est surjectif par hypothèse, on conclut que $f$ est propre par
(\hyperref[II.5.4.3-fr]{5.4.3}).

\begin{corollary}[5.4.6]
\label{II.5.4.6-fr}
Soit $f:X\to Y$ un morphisme séparé et de type fini~; pour que $f$ soit
propre, il faut et il suffit que
$f_{\mathrm{red}}:X_{\mathrm{red}}\to Y_{\mathrm{red}}$ le soit.
\end{corollary}

C'est un cas particulier de (\hyperref[II.5.4.5-fr]{5.4.5}) avec $n=1$,
$X_1=X_{\mathrm{red}}$ et $Y_1=Y_{\mathrm{red}}$
(\hyperref[I.5.1.5-fr]{I, 5.1.5}).

\begin{env}[5.4.7]
\label{II.5.4.7-fr}
Si $X$ et $Y$ sont des préschémas noethériens et $f:X\to Y$ un morphisme
séparé de type fini, on peut, pour vérifier que $f$ est propre, se ramener à
le faire pour des morphismes dominants de préschémas intègres. En effet,
soient $X_i$ ($1\leq i\leq n$) les composantes irréductibles en nombre fini
de $X$ et considérons pour chaque $i$ l'unique sous-préschéma fermé réduit de
$X$ ayant $X_i$ pour espace sous-jacent, que nous désignerons encore par
$X_i$ (\hyperref[I.5.2.1-fr]{I, 5.2.1}). Soit d'autre part $Y_i$ l'unique
sous-préschéma fermé réduit de $Y$ ayant $\overline{f(X_i)}$ pour espace
sous-jacent. Si $g_i$ (resp. $h_i$) est le morphisme d'injection $X_i\to X$
(resp. $Y_i\to Y$), on en conclut que $f\circ g_i=h_i\circ f_i$, où $f_i$
est un morphisme dominant $X_i\to Y_i$
(\hyperref[I.5.2.2-fr]{I, 5.2.2})~; on est alors dans les conditions
d'application de (\hyperref[II.5.4.5-fr]{5.4.5}), et pour que $f$ soit
propre, il faut et il suffit que les $f_i$ le soient.
\end{env}

\begin{corollary}[5.4.8]
\label{II.5.4.8-fr}
Soient $X$, $Y$ des $S$-préschémas séparés de type fini sur $S$, et soit
$f:X\to Y$ un $S$-morphisme. Pour que $f$ soit propre, il faut et il suffit
que, pour tout $S$-préschéma $S'$, le morphisme
$f\times_S1_{S'}:X\times_SS'\to Y\times_SS'$ soit fermé.
\end{corollary}

Notons d'abord que si $g:X\to S$ et $h:Y\to S$ sont les morphismes
structuraux, on a par définition $g=h\circ f$, donc $f$ est séparé et de
type fini (\hyperref[I.5.5.1-fr]{I, 5.5.1} et
\hyperref[I.6.3.4-fr]{6.3.4}). Si $f$ est propre, il en est de même de
$f\times_S1_{S'}$ (\hyperref[II.5.4.2-fr]{5.4.2})~; a fortiori,
$f\times_S1_{S'}$ est fermé. Inversement, supposons vérifiée la condition de
l'énoncé et soit $Y'$ un $Y$-préschéma~; $Y'$ peut aussi être considéré
comme un $S$-préschéma, et comme $Y\to S$ est séparé, $X\times_YY'$
s'identifie à un sous-préschéma fermé de $X\times_SY'$
(\hyperref[I.5.4.2-fr]{I, 5.4.2}). Dans le diagramme commutatif
\[
  \xymatrix{
    X\times_YY'\ar[r]^{f\times1_{Y'}}\ar[d] &
    Y\times_YY'=Y'\ar[d]\\
    X\times_SY'\ar[r]_{f\times_S1_{Y'}} & Y\times_SY'
  }
\]
les flèches verticales sont des immersions fermées~; il en résulte aussitôt
que si $f\times_S1_{Y'}$ est un morphisme fermé, il en est de même de
$f\times_Y1_{Y'}$.

\begin{remark}[5.4.9]
\label{II.5.4.9-fr}
Nous dirons qu'un morphisme $f:X\to Y$ est \emph{universellement fermé} s'il
satisfait à la condition b) de la définition
(\hyperref[II.5.4.1-fr]{5.4.1}). Le lecteur observera que

\oldpage[II]{103}

dans les propositions (\hyperref[II.5.4.2-fr]{5.4.2}) à
(\hyperref[II.5.4.8-fr]{5.4.8}), on peut partout remplacer «~propre~» par
«~universellement fermé~» sans modifier la validité des résultats (et dans
les hypothèses de (\hyperref[II.5.4.3-fr]{5.4.3}),
(\hyperref[II.5.4.5-fr]{5.4.5}), (\hyperref[II.5.4.6-fr]{5.4.6}) et
(\hyperref[II.5.4.8-fr]{5.4.8}), on peut omettre les conditions de
finitude).
\end{remark}

\begin{env}[5.4.10]
\label{II.5.4.10-fr}
Soit $f:X\to Y$ un morphisme de type fini. Nous dirons qu'une partie fermée
$Z$ de $X$ est \emph{propre sur $Y$} (ou \emph{$Y$-propre}, ou \emph{propre
pour $f$}) si la restriction de $f$ à un sous-préschéma fermé de $X$,
d'espace sous-jacent $Z$ (\hyperref[I.5.2.1-fr]{I, 5.2.1}), est propre.
Comme cette restriction est alors séparée, il résulte de
(\hyperref[II.5.4.6-fr]{5.4.6}) et (\hyperref[I.5.5.1-fr]{I, 5.5.1, (vi)})
que la propriété précédente ne dépend pas du sous-préschéma fermé de $X$
ayant $Z$ pour espace sous-jacent. Si $g:X'\to X$ est un morphisme propre,
$g^{-1}(Z)$ est une partie propre de $X'$~: si $T$ est un sous-préschéma de
$X$ ayant pour espace sous-jacent $Z$, il suffit en effet de remarquer que
la restriction de $g$ au sous-préschéma fermé $g^{-1}(T)$ de $X'$ est un
morphisme propre $g^{-1}(T)\to T$ par
(\hyperref[II.5.4.2-fr]{5.4.2, (iii)}), et d'appliquer ensuite
(\hyperref[II.5.4.2-fr]{5.4.2, (ii)}). D'autre part, si $X''$ est un
$Y$-schéma de type fini et $u:X\to X''$ un $Y$-morphisme, $u(Z)$ est une
partie propre de $X''$~; en effet, prenons pour $T$ le sous-préschéma fermé
réduit de $X$ ayant $Z$ pour espace sous-jacent~; la restriction de $f$ à
$T$ étant propre, il en est de même de la restriction de $u$ à $T$
(\hyperref[II.5.4.3-fr]{5.4.3, (i)}), donc $u(Z)$ est fermé dans $X''$~;
soit $T''$ un sous-préschéma fermé de $X''$ ayant $u(Z)$ pour espace
sous-jacent (\hyperref[I.5.2.1-fr]{I, 5.2.1}), de sorte que $u|T$ se
factorise en
\[
  T\xrightarrow{v}T''\xrightarrow{j}X'',
\]
où $j$ est l'injection canonique (\hyperref[I.5.2.2-fr]{I, 5.2.2}) et $v$
est donc propre et surjectif (\hyperref[II.5.4.5-fr]{5.4.5})~; si $g$ est la
restriction à $T''$ du morphisme structural $X''\to Y$, $g$ est séparé et
de type fini, et on a $f|T=g\circ v$~; il résulte donc de
(\hyperref[II.5.4.3-fr]{5.4.3, (ii)}) que $g$ est propre, d'où notre
assertion.
\end{env}

Il résulte en particulier de ces remarques que si $Z$ est une partie
$Y$-propre de $X$, alors~:
\begin{enumerate}
  \item[\rm{1\textsuperscript{o}}] Pour tout sous-préschéma fermé $X'$ de
  $X$, $Z\cap X'$ est une partie $Y$-propre de $X'$.
  \item[\rm{2\textsuperscript{o}}] Si $X$ est un sous-préschéma d'un
  $Y$-schéma de type fini $X''$, $Z$ est aussi une partie $Y$-propre de
  $X''$ (et en particulier est fermée dans $X''$).
\end{enumerate}

\subsection{Morphismes projectifs.}
\label{subsection:II.5.5-fr}

\begin{proposition}[5.5.1]
\label{II.5.5.1-fr}
Soit $X$ un $Y$-préschéma. Les conditions suivantes sont équivalentes~:
\begin{enumerate}
  \item[\rm{a)}] $X$ est $Y$-isomorphe à un sous-préschéma fermé d'un fibré
  projectif $\mathbf{P}(\mathcal{E})$, où $\mathcal{E}$ est un
  $\mathcal{O}_Y$-Module quasi-cohérent de type fini.
  \item[\rm{b)}] Il existe une $\mathcal{O}_Y$-Algèbre graduée
  quasi-cohérente $\mathcal{S}$, telle que $\mathcal{S}_1$ soit de type fini
  et engendre $\mathcal{S}$, et que $X$ soit $Y$-isomorphe à
  $\operatorname{Proj}(\mathcal{S})$.
\end{enumerate}
\end{proposition}

La condition a) entraîne b) en vertu de
(\hyperref[II.3.6.2-fr]{3.6.2, (ii)})~: si $\mathcal{J}$ est un faisceau
gradué quasi-cohérent d'idéaux de $\mathbf{S}(\mathcal{E})$, la
$\mathcal{O}_Y$-Algèbre graduée quasi-cohérente
$\mathcal{S}=\mathbf{S}(\mathcal{E})/\mathcal{J}$ est engendrée par
$\mathcal{S}_1$ et ce dernier, image canonique de $\mathcal{E}$, est un
$\mathcal{O}_Y$-Module de type fini. La condition b) implique a) en vertu de
(\hyperref[II.3.6.2-fr]{3.6.2}) appliqué au cas où
$\mathcal{M}\to\mathcal{S}_1$ est l'application identique.

\oldpage[II]{104}

\begin{definition}[5.5.2]
\label{II.5.5.2-fr}
On dit qu'un $Y$-préschéma $X$ est \emph{projectif sur $Y$} ou est un
\emph{$Y$-schéma projectif} s'il vérifie les conditions équivalentes a) et b)
de (\hyperref[II.5.5.1-fr]{5.5.1}). On dit qu'un morphisme $f:X\to Y$ est
\emph{projectif} s'il fait de $X$ un $Y$-schéma projectif.
\end{definition}

Il est clair que si $f:X\to Y$ est projectif, il existe un
$\mathcal{O}_X$-Module \emph{très ample relativement à $f$}
(\hyperref[II.4.4.2-fr]{4.4.2}).

\begin{theorem}[5.5.3]
\label{II.5.5.3-fr}
\begin{enumerate}
  \item[\rm{(i)}] Tout morphisme projectif est quasi-projectif et propre.
  \item[\rm{(ii)}] Inversement, soit $Y$ un schéma quasi-compact ou un
  préschéma dont l'espace sous-jacent est noethérien~; alors tout morphisme
  $f:X\to Y$ qui est quasi-projectif et propre est projectif.
\end{enumerate}
\end{theorem}

(i) Il est clair que si $f:X\to Y$ est projectif, il est de type fini et
quasi-projectif (donc en particulier séparé)~; d'autre part, il résulte
aussitôt de (\hyperref[II.5.5.1-fr]{5.5.1, b)}) et de
(\hyperref[II.3.5.3-fr]{3.5.3}) que si $f$ est projectif, il en est de même de
$f\times_Y1_{Y'}:X\times_YY'\to Y'$ pour tout morphisme $Y'\to Y$. Pour
montrer que $f$ est universellement fermé, tout revient donc à prouver qu'un
morphisme projectif $f$ est \emph{fermé}. La question étant locale sur $Y$, on
peut supposer que $Y=\operatorname{Spec}(A)$, donc
(\hyperref[II.5.5.1-fr]{5.5.1}) $X=\operatorname{Proj}(S)$, où $S$ est une
$A$-algèbre graduée engendrée par un nombre fini d'éléments de $S_1$. Pour
tout $y\in Y$, la fibre $f^{-1}(y)$ s'identifie à
$\operatorname{Proj}(S)\times_Y\operatorname{Spec}(k(y))$
(\hyperref[I.3.6.1-fr]{I, 3.6.1}), donc à
$\operatorname{Proj}(S\otimes_Ak(y))$
(\hyperref[II.2.8.10-fr]{2.8.10})~; par suite, $f^{-1}(y)$ est vide si et
seulement si $S\otimes_Ak(y)$ vérifie la condition (TN)
(\hyperref[II.2.7.4-fr]{2.7.4}), autrement dit $S_n\otimes_Ak(y)=0$ pour $n$
assez grand. Mais comme $(S_n)_y$ est un $\mathcal{O}_y$-module de type fini,
la condition précédente signifie que $(S_n)_y=0$ pour $n$ assez grand, en
vertu du lemme de Nakayama. Si $\mathfrak a_n$ est l'annulateur dans $A$ du
$A$-module $S_n$, la condition précédente signifie encore que
$\mathfrak a_n\not\subset\mathfrak j_y$ pour $n$ assez grand
(\hyperref[0.1.7.4-fr]{0, 1.7.4}). Or, comme $S_nS_1=S_{n+1}$ par hypothèse,
on a $\mathfrak a_n\subset\mathfrak a_{n+1}$, et si $\mathfrak a$ est la
réunion des $\mathfrak a_n$, on voit donc que $f(X)=V(\mathfrak a)$, ce qui
prouve que $f(X)$ est fermé dans $Y$. Si maintenant $X'$ est une partie
fermée quelconque de $X$, il existe un sous-préschéma fermé de $X$ ayant $X'$
pour espace sous-jacent (\hyperref[I.5.2.1-fr]{I, 5.2.1}) et il est clair
(\hyperref[II.5.5.1-fr]{5.5.1, a)}) que le morphisme
$X'\to X\xrightarrow{f}Y$ est projectif, donc $f(X')$ est fermé dans $Y$.

(ii) L'hypothèse sur $Y$ et le fait que $f$ est quasi-projectif entraînent
l'existence d'un $\mathcal{O}_Y$-Module quasi-cohérent de type fini
$\mathcal{E}$ et d'une $Y$-immersion $j:X\to\mathbf{P}(\mathcal{E})$
(\hyperref[II.5.3.2-fr]{5.3.2}). Mais comme $f$ est propre, $j$ est fermée par
(\hyperref[II.5.4.4-fr]{5.4.4}), donc $f$ est projectif.

\par\medskip\noindent\textit{Remarques (5.5.4). ---\ }
\label{II.5.5.4-fr}
\begin{enumerate}
  \item[\rm{(i)}] Soit $f:X\to Y$ un morphisme tel que~:
  1\textsuperscript{o} $f$ soit propre~; 2\textsuperscript{o} il existe un
  $\mathcal{O}_X$-Module $\mathcal{L}$ très ample relativement à $f$~;
  3\textsuperscript{o} le $\mathcal{O}_Y$-Module quasi-cohérent
  $\mathcal{E}=f_*(\mathcal{L})$ soit de type fini. Alors $f$ est un morphisme
  projectif~: en effet (\hyperref[II.4.4.4-fr]{4.4.4}), il y a alors une
  $Y$-immersion $r:X\to\mathbf{P}(\mathcal{E})$, et comme $f$ est propre, $r$
  est une immersion fermée (\hyperref[II.5.4.4-fr]{5.4.4}). On verra au
  chapitre III, \S~3, que lorsque $Y$ est localement noethérien, la condition
  3\textsuperscript{o} ci-dessus est une conséquence des deux autres, donc les
  conditions 1\textsuperscript{o} et 2\textsuperscript{o} caractérisent dans
  ce cas les morphismes projectifs et si $Y$ est quasi-compact, on peut
  remplacer la condition 2\textsuperscript{o} par l'hypothèse de l'existence
  d'un $\mathcal{O}_X$-Module ample pour $f$
  (\hyperref[II.4.6.11-fr]{4.6.11}).
  \item[\rm{(ii)}] Soit $Y$ un schéma quasi-compact tel qu'il existe un
  $\mathcal{O}_Y$-Module ample. Pour qu'un $Y$-schéma $X$ soit projectif, il
  faut et il suffit qu'il soit $Y$-isomorphe à un sous-$Y$-schéma fermé d'un
  fibré projectif de la forme $\mathbf{P}_Y^r$. La condition est évidemment
  suffisante.

\oldpage[II]{105}

  Inversement, si $X$ est projectif sur $Y$, il est quasi-projectif, donc il
  existe une $Y$-immersion $j$ de $X$ dans un $\mathbf{P}_Y^r$
  (\hyperref[II.5.3.3-fr]{5.3.3}) qui est fermée par
  (\hyperref[II.5.4.4-fr]{5.4.4}) et
  (\hyperref[II.5.5.3-fr]{5.5.3}).
  \item[\rm{(iii)}] Le raisonnement de
  (\hyperref[II.5.5.3-fr]{5.5.3}) montre que, pour tout préschéma $Y$, et tout
  entier $r\geq0$, le morphisme structural $\mathbf{P}_Y^r\to Y$ est
  surjectif, car si on pose
  $\mathcal{S}=\mathbf{S}_{\mathcal{O}_Y}(\mathcal{O}_Y^{r+1})$, on a
  évidemment $\mathcal{S}_y=\mathbf{S}_{k(y)}((k(y))^{r+1})$
  (\hyperref[II.1.7.3-fr]{1.7.3}), donc $(\mathcal{S}_n)_y\neq0$ pour tout
  $y\in Y$ et tout $n\geq0$.
  \item[\rm{(iv)}] Il résulte des exemples de Nagata \cite{II-26} qu'il y a
  des morphismes propres qui ne sont pas quasi-projectifs.
\end{enumerate}

\begin{proposition}[5.5.5]
\label{II.5.5.5-fr}
\begin{enumerate}
  \item[\rm{(i)}] Une immersion fermée est un morphisme projectif.
  \item[\rm{(ii)}] Si $f:X\to Y$ et $g:Y\to Z$ sont des morphismes
  projectifs, et si $Z$ est un schéma quasi-compact ou un préschéma dont
  l'espace sous-jacent est noethérien, $g\circ f$ est projectif.
  \item[\rm{(iii)}] Si $f:X\to Y$ est un $S$-morphisme projectif,
  $f_{(S')}:X_{(S')}\to Y_{(S')}$ est projectif pour toute extension
  $S'\to S$ du préschéma de base.
  \item[\rm{(iv)}] Si $f:X\to Y$ et $g:X'\to Y'$ sont des $S$-morphismes
  projectifs, il en est de même de $f\times_Sg$.
  \item[\rm{(v)}] Si $g\circ f$ est un morphisme projectif et si $g$ est
  séparé, $f$ est projectif.
  \item[\rm{(vi)}] Si $f$ est projectif, il en est de même de
  $f_{\mathrm{red}}$.
\end{enumerate}
\end{proposition}

(i) résulte aussitôt de (\hyperref[II.3.1.7-fr]{3.1.7}). Il faut ici
démontrer (iii) et (iv) séparément, à cause de la restriction introduite sur
$Z$ dans (ii) (cf. (\hyperref[I.3.5.1-fr]{I, 3.5.1})). Pour prouver (iii), on
se ramène au cas où $S=Y$ (\hyperref[I.3.3.11-fr]{I, 3.3.11}) et l'assertion
résulte alors aussitôt de (\hyperref[II.5.5.1-fr]{5.5.1, b)}) et de
(\hyperref[II.3.5.3-fr]{3.5.3}). Pour démontrer (iv), on est aussitôt ramené
au cas où $X=\mathbf{P}(\mathcal{E})$, $X'=\mathbf{P}(\mathcal{E}')$,
$\mathcal{E}$ (resp. $\mathcal{E}'$) étant un $\mathcal{O}_Y$-Module (resp.
un $\mathcal{O}_{Y'}$-Module) quasi-cohérent de type fini. Soient $p$, $p'$ les
projections canoniques de $T=Y\times_SY'$ dans $Y$ et $Y'$ respectivement~;
d'après (\hyperref[II.4.1.3.1-fr]{4.1.3.1}), on a
$\mathbf{P}(p^*(\mathcal{E}))=\mathbf{P}(\mathcal{E})\times_YT$,
$\mathbf{P}(p'^*(\mathcal{E}'))=\mathbf{P}(\mathcal{E}')\times_{Y'}T$, d'où
\begin{align*}
\mathbf{P}(p^*(\mathcal{E}))\times_T\mathbf{P}(p'^*(\mathcal{E}'))
&=(\mathbf{P}(\mathcal{E})\times_YT)\times_T
(T\times_{Y'}\mathbf{P}(\mathcal{E}'))\\
&=\mathbf{P}(\mathcal{E})\times_Y
(T\times_{Y'}\mathbf{P}(\mathcal{E}'))
=\mathbf{P}(\mathcal{E})\times_S\mathbf{P}(\mathcal{E}')
\end{align*}
en remplaçant $T$ par $Y\times_SY'$ et utilisant
(\hyperref[I.3.3.9.1-fr]{I, 3.3.9.1}). Or, $p^*(\mathcal{E})$ et
$p'^*(\mathcal{E}')$ sont de type fini sur $T$
(\hyperref[0.5.2.4-fr]{0, 5.2.4}), donc il en est de même de
$p^*(\mathcal{E})\otimes_{\mathcal{O}_T}p'^*(\mathcal{E}')$~; comme
$\mathbf{P}(p^*(\mathcal{E}))\times_T\mathbf{P}(p'^*(\mathcal{E}'))$
s'identifie à un sous-préschéma fermé de
$\mathbf{P}(p^*(\mathcal{E})\otimes_{\mathcal{O}_T}p'^*(\mathcal{E}'))$
(\hyperref[II.4.3.3-fr]{4.3.3}), cela achève de prouver (iv). Pour obtenir (v)
et (vi), on peut appliquer (\hyperref[I.5.5.13-fr]{I, 5.5.13}), car tout
sous-préschéma fermé d'un $Y$-schéma projectif est un $Y$-schéma projectif par
(\hyperref[II.5.5.1-fr]{5.5.1, a)}).

Il reste à prouver (ii)~; en vertu de l'hypothèse sur $Z$, cela résulte de
(\hyperref[II.5.5.3-fr]{5.5.3}), de
(\hyperref[II.5.3.4-fr]{5.3.4, (ii)}) et
(\hyperref[II.5.4.2-fr]{5.4.2, (ii)}).

\begin{proposition}[5.5.6]
\label{II.5.5.6-fr}
Si $X$ et $X'$ sont deux $Y$-schémas projectifs, $X\amalg X'$ est un
$Y$-schéma projectif.
\end{proposition}

C'est une conséquence évidente de (\hyperref[II.5.5.2-fr]{5.5.2}) et
(\hyperref[II.4.3.6-fr]{4.3.6}).

\begin{proposition}[5.5.7]
\label{II.5.5.7-fr}
Soit $X$ un $Y$-schéma projectif, $\mathcal{L}$ un $\mathcal{O}_X$-Module
$Y$-ample~; pour toute section $f$ de $\mathcal{L}$ au-dessus de $X$, $X_f$
est affine sur $Y$.
\end{proposition}

\oldpage[II]{106}

La question étant locale sur $Y$, on peut supposer
$Y=\operatorname{Spec}(A)$~; en outre
$X_{f^{\otimes n}}=X_f$, donc en remplaçant $\mathcal{L}$ par un
$\mathcal{L}^{\otimes n}$ convenable, on peut supposer $\mathcal{L}$ très
ample pour le morphisme structural $q:X\to Y$
(\hyperref[II.4.6.11-fr]{4.6.11}). L'homomorphisme canonique
$\sigma:q^*(q_*(\mathcal{L}))\to\mathcal{L}$ est donc surjectif et le
morphisme correspondant
\[
r=r_{\mathcal{L},\sigma}:X\to P=\mathbf{P}(q_*(\mathcal{L}))
\]
une immersion telle que $\mathcal{L}=r^*(\mathcal{O}_P(1))$
(\hyperref[II.4.4.4-fr]{4.4.4})~; en outre, comme $X$ est propre sur $Y$,
l'immersion $r$ est fermée (\hyperref[II.5.4.4-fr]{5.4.4}). Par ailleurs, on
a par définition $f\in\Gamma(Y,q_*(\mathcal{L}))$ et $\sigma^\flat$ est
l'identité de $q_*(\mathcal{L})$~; il résulte alors de la formule
(\hyperref[II.3.7.3.1-fr]{3.7.3.1}) que l'on a $X_f=r^{-1}(D_+(f))$~; $X_f$
est par suite un sous-préschéma fermé du schéma affine $D_+(f)$ et est donc
bien un schéma affine.

Lorsqu'on prend en particulier $Y=X$, on obtient (compte tenu de
(\hyperref[II.4.6.13-fr]{4.6.13, (i)})) le corollaire suivant, dont la
démonstration directe est d'ailleurs immédiate~:

\begin{corollary}[5.5.8]
\label{II.5.5.8-fr}
Soit $X$ un préschéma, $\mathcal{L}$ un $\mathcal{O}_X$-Module inversible.
Pour toute section $f$ de $\mathcal{L}$ au-dessus de $X$, $X_f$ est affine
sur $X$ (et par suite un schéma affine lorsque $X$ est un schéma affine).
\end{corollary}

\subsection{Le lemme de Chow.}
\label{subsection:II.5.6-fr}

\begin{theorem}[5.6.1]
\label{II.5.6.1-fr}
(Lemme de Chow). Soit $S$ un préschéma, $X$ un $S$-schéma de type fini. On
suppose vérifiée l'une des hypothèses suivantes~:
\begin{enumerate}
  \item[\rm{a)}] $S$ est noethérien.
  \item[\rm{b)}] $S$ est un schéma quasi-compact, et $X$ a un nombre fini de
  composantes irréductibles.
\end{enumerate}
Dans ces conditions~:
\begin{enumerate}
  \item[\rm{(i)}] Il existe un $S$-schéma $X'$ quasi-projectif sur $S$ et un
  $S$-morphisme $f:X'\to X$, projectif et surjectif.
  \item[\rm{(ii)}] On peut prendre $X'$ et $f$ tels qu'il existe un ouvert
  $U\subset X$ pour lequel $U'=f^{-1}(U)$ soit dense dans $X'$ et la
  restriction de $f$ à $U'$ un isomorphisme $U'\simeq U$.
  \item[\rm{(iii)}] Si $X$ est réduit (resp. irréductible, intègre), on peut
  supposer $X'$ réduit (resp. irréductible, intègre).
\end{enumerate}
\end{theorem}

La démonstration se fait en plusieurs étapes.

A) On peut d'abord se ramener au cas où $X$ est irréductible. En effet, dans
l'hypothèse a), $X$ est noethérien, donc, dans les deux hypothèses, les
composantes irréductibles $X_i$ de $X$ sont en nombre fini. Si le théorème est
démontré pour chacun des sous-préschémas fermés réduits de $X$ ayant les $X_i$
pour espaces sous-jacents, et si $X'_i$ et $f_i:X'_i\to X_i$ sont le préschéma
et le morphisme correspondant à $X_i$, le préschéma $X'$ somme des $X'_i$, et
le morphisme $f:X'\to X$ dont la restriction à chaque $X'_i$ est $j_i\circ f_i$
($j_i$ injection canonique $X_i\to X$), répondent à la question. Il est
immédiat en effet que $X'$ est réduit si les $X'_i$ le sont~; d'autre part on
satisfait à (ii) en prenant pour $U$ la réunion des ensembles
$U_i\cap(\bigcup_{j\neq i}X_j)^{\complement}$. Enfin, comme les $X'_i$ sont
quasi-projectifs sur $S$, il en est de même de $X'$

\oldpage[II]{107}

(\hyperref[II.5.3.6-fr]{5.3.6})~; de même, les morphismes $X'_i\to X$ sont
projectifs par (\hyperref[II.5.5.5-fr]{5.5.5, (i) et (ii)}), donc $f$ est
projectif (\hyperref[II.5.5.6-fr]{5.5.6}), et est évidemment surjectif par
définition.

B) Supposons maintenant $X$ irréductible. Comme le morphisme structural
$r:X\to S$ est de type fini, il y a un recouvrement fini $(S_i)$ de $S$ par
des ouverts affines, et pour chaque $i$ un recouvrement fini $(T_{ij})$ de
$r^{-1}(S_i)$ par des ouverts affines, le morphisme $T_{ij}\to S_i$ étant
affine et de type fini, donc quasi-projectif
(\hyperref[II.5.3.4-fr]{5.3.4, (i)})~; comme dans chacune des hypothèses a),
b), l'immersion $S_i\to S$ est quasi-compacte, c'est un morphisme
quasi-projectif (\hyperref[II.5.3.4-fr]{5.3.4, (ii)}), donc la restriction de
$r$ à $T_{ij}$ est un morphisme quasi-projectif
(\hyperref[II.5.3.4-fr]{5.3.4, (iii)}). Désignons les $T_{ij}$ par $U_k$
$(1\leq k\leq n)$. Il existe, pour chaque indice $k$ une immersion ouverte
$\varphi_k:U_k\to P_k$, où $P_k$ est projectif sur $S$
(\hyperref[II.5.3.2-fr]{5.3.2} et \hyperref[II.5.5.2-fr]{5.5.2}). Soit
$U=\bigcap_kU_k$~; comme $X$ est irréductible et les $U_k$ non vides, $U$ est
non vide, et par suite partout dense dans $X$~; les restrictions à $U$ des
$\varphi_k$ définissent un morphisme
\[
\varphi:U\to P=P_1\times_SP_2\times_S\cdots\times_SP_n
\]
de sorte que les diagrammes
\[
\tag{5.6.1.1}
\label{II.5.6.1.1-fr}
\xymatrix{
U\ar[r]^\varphi\ar[d]_{j_k} & P\ar[d]^{p_k}\\
U_k\ar[r]_{\varphi_k} & P_k
}
\]
soient commutatifs, $j_k$ étant l'injection canonique $U\to U_k$ et $p_k$ la
projection canonique $P\to P_k$. Si $j$ est l'injection canonique $U\to X$,
le morphisme $\psi=(j,\varphi)_S:U\to X\times_SP$ est une immersion
(\hyperref[I.5.3.14-fr]{I, 5.3.14}). Dans l'hypothèse a), $X\times_SP$ est
localement noethérien (\hyperref[II.3.4.1-fr]{3.4.1} et
\hyperref[I.6.3.7-fr]{I, 6.3.7 et 6.3.8})~; dans l'hypothèse b),
$X\times_SP$ est un schéma quasi-compact
(\hyperref[I.5.5.1-fr]{I, 5.5.1 et 6.6.4})~; dans les deux cas l'adhérence
$X'$ dans $X\times_SP$ du sous-préschéma $Z$ associé à $\psi$ (et dont
$\psi(U)$ est l'espace sous-jacent) existe, et $\psi$ se factorise en
\[
\tag{5.6.1.2}
\label{II.5.6.1.2-fr}
\psi:U\xrightarrow{\psi'}X'\xrightarrow{h}X\times_SP
\]
où $\psi'$ est une immersion ouverte et $h$ une immersion fermée
(\hyperref[I.9.5.10-fr]{I, 9.5.10}). Soient
$q_1:X\times_SP\to X$ et $q_2:X\times_SP\to P$ les projections canoniques~;
nous poserons
\[
\tag{5.6.1.3}
\label{II.5.6.1.3-fr}
f:X'\xrightarrow{h}X\times_SP\xrightarrow{q_1}X
\]
et
\[
\tag{5.6.1.4}
\label{II.5.6.1.4-fr}
g:X'\xrightarrow{h}X\times_SP\xrightarrow{q_2}P.
\]
Nous allons voir que $X'$ et $f$ répondent à la question.

C) Montrons d'abord que $f$ est projectif et surjectif, et que la restriction
de $f$ à $U'=f^{-1}(U)$ est un isomorphisme de $U'$ sur $U$. Comme les $P_k$
sont projectifs sur $S$, il en est de même de $P$
(\hyperref[II.5.5.5-fr]{5.5.5, (iv)}), donc $X\times_SP$ est projectif sur
$X$ (\hyperref[II.5.5.5-fr]{5.5.5, (iii)}), et il en est de même de $X'$ qui
est un sous-préschéma fermé de $X\times_SP$. D'autre part, on a
$f\circ\psi'=q_1\circ(h\circ\psi')=q_1\circ\psi=j$, donc $f(X')$ contient
l'ensemble ouvert partout dense $U$ de $X$~; mais $f$ est un morphisme fermé
(\hyperref[II.5.5.3-fr]{5.5.3}), donc $f(X')=X$. Notons maintenant que
$q_1^{-1}(U)=U\times_SP$ est induit sur un ouvert de $X\times_SP$, et par
définition le préschéma $U'=h^{-1}(U\times_SP)$ est induit par $X'$ sur
l'ouvert $U'$~; c'est par suite l'adhérence par

\oldpage[II]{108}

rapport à $U\times_SP$ du préschéma $Z$
(\hyperref[I.9.5.8-fr]{I, 9.5.8}). Or l'immersion $\psi$ se factorise en
\[
U\xrightarrow{\Gamma_\varphi}U\times_SP\xrightarrow{j\times1}X\times_SP,
\]
et comme $P$ est séparé sur $S$, le morphisme graphe $\Gamma_\varphi$ est une
immersion fermée (\hyperref[I.5.4.3-fr]{I, 5.4.3}), donc $Z$ est un
sous-préschéma fermé de $U\times_SP$, d'où on conclut que $U'=Z$. Comme
d'ailleurs $\psi$ est une immersion, la restriction de $f$ à $U'$ est un
isomorphisme sur $U$, réciproque de $\psi'$~; enfin, par définition de $X'$,
$U'$ est dense dans $X'$.

D) Prouvons maintenant que $g$ est une immersion, ce qui établira que $X'$ est
quasi-projectif sur $S$, puisque $P$ est projectif sur $S$. Posons
\begin{align*}
V_k&=\varphi_k(U_k) &&\text{(ouvert de $P_k$)},\\
W_k&=p_k^{-1}(V_k) &&\text{(ouvert de $P$)},\\
U'_k&=f^{-1}(U_k),\qquad U''_k=g^{-1}(W_k)
&&\text{(ouverts de $X'$)}.
\end{align*}
Il est clair que les $U'_k$ forment un recouvrement ouvert de $X'$~; nous
allons d'abord voir qu'il en est de même des $U''_k$, en prouvant que
$U'_k\subset U''_k$. Il suffira pour cela de montrer que le diagramme
\[
\tag{5.6.1.5}
\label{II.5.6.1.5-fr}
\xymatrix{
U'_k\ar[r]^{g|U'_k}\ar[d]_{f|U'_k} & P\ar[d]^{p_k}\\
U_k\ar[r]_{\varphi_k} & P_k
}
\]
est commutatif. Or le préschéma $U'_k=h^{-1}(U_k\times_SP)$ est induit par
$X'$ sur l'ouvert $U'_k$, et est donc l'adhérence de $Z=U'\subset U'_k$ par
rapport à $U'_k$ (\hyperref[I.9.5.8-fr]{I, 9.5.8}). Pour montrer la
commutativité du diagramme (\hyperref[II.5.6.1.5-fr]{5.6.1.5}), il suffit
donc, puisque $P_k$ est un $S$-schéma, de montrer qu'en composant ce diagramme
avec l'injection canonique $U'\to U'_k$ (ou, ce qui revient au même en vertu de
l'isomorphisme de $U'$ et $U$, avec $\psi$), on obtient un diagramme commutatif
(\hyperref[I.9.5.6-fr]{I, 9.5.6}). Mais par définition le diagramme que l'on
obtient alors n'est autre que (\hyperref[II.5.6.1.1-fr]{5.6.1.1}), d'où notre
assertion.

Les $W_k$ forment donc un recouvrement ouvert de $g(X')$~; pour prouver que
$g$ est une immersion, il suffit de voir que chacune des restrictions
$g|U''_k$ est une immersion dans $W_k$
(\hyperref[I.4.2.4-fr]{I, 4.2.4}). Pour cela, considérons le morphisme
$u_k:W_k\xrightarrow{p_k}V_k\xrightarrow{\varphi_k^{-1}}U_k\to X$~; comme
$X$ est séparé sur $S$, le morphisme graphe
$\Gamma_{u_k}:W_k\to X\times_SW_k$ est une immersion fermée
(\hyperref[I.5.4.3-fr]{I, 5.4.3}), donc le graphe
$T_k=\Gamma_{u_k}(W_k)$ est un sous-préschéma fermé de $X\times_SW_k$~; si
nous montrons que ce sous-préschéma majore $U'$, il majorera aussi le
sous-préschéma induit par $X'$ sur l'ouvert $X''_k$ de $X'$, par
(\hyperref[I.9.5.8-fr]{I, 9.5.8}). Comme la restriction de $q_2$ à $T_k$ est
un isomorphisme sur $W_k$, la restriction de $g$ à $X''_k$ sera une immersion
dans $W_k$, et notre assertion sera démontrée. Soit $v_k$ l'injection
canonique $U'\to X\times_SW_k$~; nous devons montrer qu'il existe un morphisme
$w_k:U'\to W_k$ tel que $v_k=\Gamma_{u_k}\circ w_k$. Par définition du
produit, il suffit de prouver que
$q_1\circ v_k=u_k\circ q_2\circ v_k$
(\hyperref[I.3.2.1-fr]{I, 3.2.1}), ou, en composant à droite

\oldpage[II]{109}

avec l'isomorphisme $\psi':U\to U'$, que
$q_1\circ\psi=u_k\circ q_2\circ\psi$. Mais comme $q_1\circ\psi=j$,
$q_2\circ\psi=\varphi$, notre assertion résulte de la commutativité du
diagramme (\hyperref[II.5.6.1.1-fr]{5.6.1.1}), compte tenu de la définition de
$u_k$.

E) Il est clair que comme $U$, et par suite $U'$, est irréductible, il en est
de même de $X'$ dans la construction précédente, et le morphisme $f$ est donc
birationnel (\hyperref[I.2.2.9-fr]{I, 2.2.9}). Si en outre $X$ est réduit, il
en est de même de $U'$, donc $X'$ est réduit
(\hyperref[I.9.5.9-fr]{I, 9.5.9}). Cela achève la démonstration.

\begin{corollary}[5.6.2]
\label{II.5.6.2-fr}
On suppose vérifiée l'une des hypothèses a), b) de
(\hyperref[II.5.6.1-fr]{5.6.1}). Pour que $X$ soit propre sur $S$, il faut et
il suffit qu'il existe un schéma $X'$ projectif sur $S$ et un $S$-morphisme
surjectif $f:X'\to X$ (qui est alors projectif par
(\hyperref[II.5.5.5-fr]{5.5.5, (v)})). Lorsqu'il en est ainsi, on peut en
outre choisir $f$ de sorte qu'il existe un ouvert dense $U$ de $X$ tel que la
restriction de $f$ à $f^{-1}(U)$ soit un isomorphisme
$f^{-1}(U)\simeq U$ et que $f^{-1}(U)$ soit dense dans $X'$. Si en outre $X$
est irréductible (resp. réduit) on peut supposer qu'il en est de même de
$X'$~; lorsque $X$ et $X'$ sont irréductibles, $f$ est un morphisme
birationnel.
\end{corollary}

La condition est suffisante en vertu de (\hyperref[II.5.5.3-fr]{5.5.3}) et
(\hyperref[II.5.4.3-fr]{5.4.3, (ii)}). Elle est nécessaire, car avec les
notations de (\hyperref[II.5.6.1-fr]{5.6.1}), si $X$ est propre sur $S$, $X'$
est propre sur $S$, car il est projectif sur $X$, donc propre sur $X$
(\hyperref[II.5.5.3-fr]{5.5.3}), et notre assertion résulte de
(\hyperref[II.5.4.2-fr]{5.4.2, (ii)})~; en outre, comme $X'$ est
quasi-projectif sur $S$, il est projectif sur $S$ en vertu de
(\hyperref[II.5.5.3-fr]{5.5.3}).

\begin{corollary}[5.6.3]
\label{II.5.6.3-fr}
Soit $S$ un préschéma localement noethérien, $X$ un $S$-schéma de type fini
sur $S$, $f_0:X\to S$ son morphisme structural. Pour que $X$ soit propre sur
$S$, il faut et il suffit que pour tout morphisme de type fini $S'\to S$,
$(f_0)_{(S')}:X_{(S')}\to S'$ soit un morphisme fermé. Il suffit même que
cette condition soit vérifiée pour tout $S$-préschéma de la forme
$S'=S\otimes_{\mathbf Z}\mathbf Z[T_1,\ldots,T_n]$ ($T_i$ indéterminées).
\end{corollary}

La condition étant évidemment nécessaire, montrons qu'elle est suffisante.
La question étant locale sur $S$ et $S'$ (\hyperref[II.5.4.1-fr]{5.4.1}), on
peut supposer $S$ et $S'$ affines et noethériens. En vertu du lemme de Chow,
il existe un $S$-schéma projectif $P$, une immersion $j:X'\to P$ et un
morphisme projectif et surjectif $f:X'\to X$, tels que le diagramme
\[
\xymatrix{
X\ar[d]_{f_0} & X'\ar[l]_f\ar[d]^j\\
S & P\ar[l]_r
}
\]
soit commutatif. Comme $P$ est de type fini sur $S$, la première hypothèse
entraîne que la projection $q_2:X\times_SP\to P$ est un morphisme fermé. Mais
l'immersion $j$ est composée de $q_2$ et du morphisme $f\times1$ de
$X'\times_SP$ dans $X\times_SP$~; or $f$, étant projectif, est propre
(\hyperref[II.5.5.3-fr]{5.5.3}), donc $f\times1$ est fermé. On en conclut que
$j$ est une immersion fermée, et par suite est propre
(\hyperref[II.5.4.2-fr]{5.4.2, (i)}). En outre le morphisme structural
$r:P\to S$ est projectif, donc propre (\hyperref[II.5.5.3-fr]{5.5.3}), donc
$f_0\circ f=r\circ j$ est propre
(\hyperref[II.5.4.2-fr]{5.4.2, (ii)})~; finalement, comme $f$ est surjectif,
$f_0$ est propre en vertu de (\hyperref[II.5.4.3-fr]{5.4.3}).

Pour prouver la proposition en utilisant seulement la seconde hypothèse, il
suffit de montrer qu'elle entraîne la première. Or, si $S'$ est affine et de
type fini sur $S=\operatorname{Spec}(A)$,

\oldpage[II]{110}

on a $S'=\operatorname{Spec}(A[c_1,\ldots,c_n])$
(\hyperref[I.6.3.3-fr]{I, 6.3.3}), et $S'$ est donc isomorphe à un
sous-préschéma fermé de $S''=\operatorname{Spec}(A[T_1,\ldots,T_n])$ ($T_i$
indéterminées). Dans le diagramme commutatif
\[
\xymatrix{
X\times_SS'\ar[r]^{1_X\times j}\ar[d]_{(f_0)_{(S')}} &
X\times_SS''\ar[d]^{(f_0)_{(S'')}}\\
S'\ar[r]_j & S''
}
\]
$j$ et $1_X\times j$ sont des immersions fermées
(\hyperref[I.4.3.1-fr]{I, 4.3.1}) et $(f_0)_{(S'')}$ est fermé par hypothèse~;
il en est donc de même de $(f_0)_{(S')}$. 

\section{Morphismes entiers et morphismes finis}
\label{section:II.6-fr}

\subsection{Préschémas entiers sur un autre.}
\label{subsection:II.6.1-fr}

\begin{definition}[6.1.1]
\label{II.6.1.1-fr}
Soit $X$ un $S$-préschéma, $f:X\to S$ son morphisme structural. On dit que
$X$ est \emph{entier au-dessus de $S$}, ou que $f$ est un \emph{morphisme
entier}, s'il existe un recouvrement $(S_\alpha)$ de $S$ par des ouverts
affines tels que, pour tout $\alpha$, le préschéma induit $f^{-1}(S_\alpha)$
soit un schéma affine, dont l'anneau $B_\alpha$ est une algèbre entière sur
l'anneau $A_\alpha$ de $S_\alpha$. On dit que $X$ est \emph{fini au-dessus de
$S$}, ou que $f$ est un \emph{morphisme fini} si $X$ est entier sur $S$ et de
type fini sur $S$.
\end{definition}

Lorsque $S$ est affine d'anneau $A$, on dira aussi «~entier (resp. fini) sur
$A$~» au lieu de «~entier (resp. fini) sur $S$~».

\begin{env}[6.1.2]
\label{II.6.1.2-fr}
Il est clair que si $X$ est entier au-dessus de $S$, il est affine au-dessus
de $S$. Pour qu'un préschéma $X$ affine au-dessus de $S$ soit entier (resp.
fini) au-dessus de $S$, il faut et il suffit que la $\mathcal{O}_S$-Algèbre
quasi-cohérente associée $\mathcal{A}(X)$ soit telle qu'il existe un
recouvrement $(S_\alpha)$ de $S$ par des ouverts affines ayant la propriété
que pour tout $\alpha$, $\Gamma(S_\alpha,\mathcal{A}(X))$ est une algèbre
entière (resp. entière et de type fini) sur
$\Gamma(S_\alpha,\mathcal{O}_S)$. Une $\mathcal{O}_S$-Algèbre quasi-cohérente
ayant cette propriété est dite \emph{entière} (resp. \emph{finie}) sur
$\mathcal{O}_S$. Se donner un préschéma entier (resp. fini) au-dessus de $S$
revient donc (\hyperref[II.1.3.1-fr]{1.3.1}) à se donner une
$\mathcal{O}_S$-Algèbre quasi-cohérente entière (resp. finie) sur
$\mathcal{O}_S$. On notera qu'une $\mathcal{O}_S$-Algèbre quasi-cohérente
$\mathcal{B}$ est finie si et seulement si c'est un $\mathcal{O}_S$-Module de
type fini (\hyperref[I.1.3.9-fr]{I, 1.3.9})~; il revient au même de dire que
$\mathcal{B}$ est une $\mathcal{O}_S$-Algèbre entière et de type fini, car une
algèbre entière et de type fini sur un anneau $A$ est un $A$-module de type
fini.
\end{env}

\begin{proposition}[6.1.3]
\label{II.6.1.3-fr}
Soit $S$ un préschéma localement noethérien. Pour qu'un préschéma $X$ affine
au-dessus de $S$ soit fini au-dessus de $S$, il faut et il suffit que la
$\mathcal{O}_S$-Algèbre $\mathcal{A}(X)$ soit cohérente.
\end{proposition}

Compte tenu de la remarque précédente, cela revient à remarquer que si $S$ est
localement noethérien, les $\mathcal{O}_S$-Modules quasi-cohérents de type fini
ne sont autres que les $\mathcal{O}_S$-Modules cohérents
(\hyperref[I.1.5.1-fr]{I, 1.5.1}).

\begin{proposition}[6.1.4]
\label{II.6.1.4-fr}
Soit $X$ un préschéma entier (resp. fini) au-dessus de $S$, $f:X\to S$ le
morphisme structural. Alors, pour tout ouvert affine $U\subset S$, d'anneau
$A$, $f^{-1}(U)$ est un schéma affine dont l'anneau $B$ est une algèbre
entière (resp. finie) sur $A$.
\end{proposition}

\oldpage[II]{111}

Nous démontrerons d'abord le lemme suivant~:

\begin{lemma}[6.1.4.1]
\label{II.6.1.4.1-fr}
Soient $A$ un anneau, $M$ un $A$-module, $(g_i)_{1\leq i\leq m}$ un système
fini d'éléments de $A$ tels que les $D(g_i)$ ($1\leq i\leq m$) recouvrent
$\operatorname{Spec}(A)$. Si, pour tout $i$, $M_{g_i}$ est un
$A_{g_i}$-module de type fini, alors $M$ est un $A$-module de type fini.
\end{lemma}

On peut en effet supposer que $M_{g_i}$ admet un système de générateurs fini
$(m_{ij}/g_i^n)$ avec $m_{ij}\in M$, $n$ étant le même pour tous les indices
$i$. Montrons que les $m_{ij}$ forment un système de générateurs de $M$. Soit
$M'$ le sous-$A$-module de $M$ engendré par ces éléments, et soit $m$ un
élément de $M$. Par hypothèse, pour chaque $i$, il existe des $a_{ij}\in A$ et
un entier $p$ (indépendant de $i$) tels que, dans $M_{g_i}$,
$m/1=(\sum_j a_{ij}m_{ij})/g_i^p$~; cela entraîne qu'il existe un entier
$r\geq p$ tel que, pour tout $i$, on ait $g_i^r m\in M'$. Or, comme les
$D(g_i^r)=D(g_i)$ recouvrent $\operatorname{Spec}(A)$, l'idéal de $A$
engendré par les $g_i^r$ est égal à $A$, autrement dit il existe des éléments
$a_i\in A$ tels que $\sum_i a_i g_i^r=1$~; par suite
$m=(\sum_i a_i g_i^r)m\in M'$, d'où le lemme.

Cela étant, on sait déjà (\hyperref[II.1.3.2-fr]{1.3.2}) que $f^{-1}(U)$ est
affine. Si $\varphi$ est l'homomorphisme $A\to B$ correspondant à $f$, il
existe un recouvrement fini de $U$ par des ouverts $D(g_i)$ ($g_i\in A$) tels
que, si $h_i=\varphi(g_i)$, $B_{h_i}$ soit une algèbre entière (resp. entière
finie) sur $A_{g_i}$. En effet, il y a un recouvrement de $U$ par des ouverts
affines $V_\alpha\subset U$ tels que si
$A_\alpha=A(V_\alpha)$, $B_\alpha=A(f^{-1}(V_\alpha))$ soit une algèbre
entière (resp. finie) sur $A_\alpha$. Tout $x\in U$ appartient à un
$V_\alpha$, donc il existe $g\in A$ tel que $x\in D(g)\subset V_\alpha$~; si
$g_\alpha$ est l'image de $g$ dans $A_\alpha$, on a
$A(D(g))=A_g=(A_\alpha)_{g_\alpha}$~; soit $h=\varphi(g)$, et soit $h_\alpha$
l'image de $g_\alpha$ dans $B_\alpha$~; on a
\[
A(D(h))=B_h=(B_\alpha)_{h_\alpha}
\]
et comme $B_\alpha$ est entière (resp. finie) sur $A_\alpha$,
$(B_\alpha)_{h_\alpha}$ est entière (resp. finie) sur
$(A_\alpha)_{g_\alpha}$. Il suffit maintenant d'utiliser le fait que $U$ est
quasi-compact pour obtenir le recouvrement cherché.

Si on suppose d'abord que les $B_{h_i}$ sont entières finies sur les
$A_{g_i}$, comme, en tant que $A_{g_i}$-module, $B_{h_i}$ s'écrit aussi
$B_{g_i}$, le lemme (\hyperref[II.6.1.4.1-fr]{6.1.4.1}) montre que dans ce cas
$B$ est un $A$-module de type fini.

Supposons seulement maintenant chaque $B_{h_i}$ entière sur $A_{g_i}$~; soit
$b\in B$, et soit $C$ la sous-$A$-algèbre de $B$ engendrée par $b$. Pour tout
$i$, $C_{h_i}$ est l'algèbre sur $A_{g_i}$ engendrée par $b/1$ dans
$B_{h_i}$~; il résulte de l'hypothèse que chaque $C_{h_i}$ est un
$A_{g_i}$-module de type fini, donc
(\hyperref[II.6.1.4.1-fr]{6.1.4.1}) $C$ est un $A$-module de type fini, ce qui
prouve que $B$ est entière sur $A$.

\begin{proposition}[6.1.5]
\label{II.6.1.5-fr}
\begin{enumerate}
\item[(i)] Une immersion fermée est finie (et a fortiori entière).
\item[(ii)] Le composé de deux morphismes finis (resp. entiers) est fini
(resp. entier).
\item[(iii)] Si $f:X\to Y$ est un $S$-morphisme fini (resp. entier),
$f_{(S')}:X_{(S')}\to Y_{(S')}$ est fini (resp. entier) pour toute extension
$S'\to S$ de la base.
\item[(iv)] Si $f:X\to Y$ et $g:X'\to Y'$ sont deux $S$-morphismes finis
(resp. entiers), il en est de même de
$f\times_S g:X\times_S X'\to Y\times_S Y'$.
\item[(v)] Si $f:X\to Y$ et $g:Y\to Z$ sont deux morphismes tels que $g\circ
f$ soit fini (resp. entier) et $g$ séparé, alors $f$ est fini (resp. entier).
\item[(vi)] Si $f:X\to Y$ est un morphisme fini (resp. entier), il en est de
même de $f_{\mathrm{red}}$.
\end{enumerate}
\end{proposition}

\oldpage[II]{112}

En vertu de (\hyperref[I.5.5.12-fr]{I, 5.5.12}), il suffit de démontrer (i),
(ii) et (iii). Pour prouver qu'une immersion fermée $X\to S$ est finie, on
peut se borner au cas où $S=\operatorname{Spec}(A)$, et tout revient alors à
remarquer qu'un anneau quotient $A/\mathfrak{J}$ est un $A$-module monogène.
Pour prouver que le composé de deux morphismes finis (resp. entiers) $X\to Y$,
$Y\to Z$ est fini (resp. entier), on peut encore supposer $Z$ (et par suite
$X$, $Y$ (\hyperref[II.1.3.4-fr]{1.3.4})) affines, et alors l'assertion
équivaut à dire que si $B$ est une $A$-algèbre finie (resp. entière) et $C$
une $B$-algèbre finie (resp. entière), alors $C$ est une $A$-algèbre finie
(resp. entière), ce qui est immédiat. Enfin, pour démontrer (iii), on peut se
borner au cas où $S=Y$, puisque $X_{(S')}$ s'identifie à
$X\times_Y Y_{(S')}$ (\hyperref[I.3.3.11-fr]{I, 3.3.11})~; on peut en outre
supposer que $S=\operatorname{Spec}(A)$, $S'=\operatorname{Spec}(A')$~; alors
$X$ est affine d'anneau $B$ (\hyperref[II.1.3.4-fr]{1.3.4}), $X_{(S')}$
affine d'anneau $A'\otimes_A B$, et il suffit de remarquer que si $B$ est une
$A$-algèbre finie (resp. entière), $A'\otimes_A B$ est une $A'$-algèbre finie
(resp. entière).

Notons aussi que si $X$ et $Y$ sont deux $S$-préschémas qui sont finis (resp.
entiers) au-dessus de $S$, leur \emph{somme} $X\amalg Y$ est un préschéma fini
(resp. entier) au-dessus de $S$, car cela revient à voir que si $B$ et $C$
sont deux $A$-algèbres finies (resp. entières) sur $A$, il en est de même de
$B\times C$.

\begin{corollary}[6.1.6]
\label{II.6.1.6-fr}
Si $X$ est un préschéma entier (resp. fini) au-dessus de $S$, alors, pour tout
ouvert $U\subset S$, $f^{-1}(U)$ est entier (resp. fini) au-dessus de $U$.
\end{corollary}

C'est un cas particulier de (\hyperref[II.6.1.5-fr]{6.1.5, (iii)}).

\begin{corollary}[6.1.7]
\label{II.6.1.7-fr}
Soit $f:X\to Y$ un morphisme fini. Alors, pour tout $y\in Y$, la fibre
$f^{-1}(y)$ est un schéma algébrique fini sur $k(y)$, et a fortiori son
espace sous-jacent est discret et fini.
\end{corollary}

En effet, en tant que $k(y)$-préschéma, $f^{-1}(y)$ est identifié à
$X\times_Y\operatorname{Spec}(k(y))$ (\hyperref[I.3.6.1-fr]{I, 3.6.1}), qui
est fini au-dessus de $\operatorname{Spec}(k(y))$
(\hyperref[II.6.1.5-fr]{6.1.5, (iii)})~; c'est donc un schéma affine dont
l'anneau est une algèbre de rang fini sur $k(y)$
(\hyperref[II.6.1.4-fr]{6.1.4}). La proposition résulte alors de
(\hyperref[I.6.4.4-fr]{I, 6.4.4}).

\begin{corollary}[6.1.8]
\label{II.6.1.8-fr}
Soient $X$, $S$ deux préschémas intègres, $f:X\to S$ un morphisme dominant. Si
$f$ est entier (resp. fini), le corps $R(X)$ des fonctions rationnelles sur
$X$ est algébrique (resp. algébrique de degré fini) sur le corps $R(S)$ des
fonctions rationnelles sur $S$.
\end{corollary}

En effet, soit $s$ le point générique de $S$~; le $k(s)$-préschéma
$f^{-1}(s)$ est entier (resp. fini) au-dessus de
$\operatorname{Spec}(k(s))$ (\hyperref[II.6.1.5-fr]{6.1.5, (iii)}), et
contient par hypothèse le point générique $x$ de $X$~; l'anneau local de $x$
dans $f^{-1}(s)$ étant égal à $k(x)$ (\hyperref[I.3.6.5-fr]{I, 3.6.5}) est un
anneau local d'une algèbre entière (resp. finie) sur $k(s)$
(\hyperref[II.6.1.4-fr]{6.1.4}), d'où le corollaire.

\begin{remark}[6.1.9]
\label{II.6.1.9-fr}
L'hypothèse que $g$ est \emph{séparé} est essentielle pour la validité de
(\hyperref[II.6.1.5-fr]{6.1.5, (v)})~: en effet, si $Y$ n'est pas séparé sur
$Z$, l'identité $1_Y$ est le morphisme composé
$Y\xrightarrow{\Delta_Y}Y\times_ZY\xrightarrow{p_1}Y$, mais $\Delta_Y$ n'est
pas un morphisme entier, comme il résulte de
(\hyperref[II.6.1.10-fr]{6.1.10})~:
\end{remark}

\begin{proposition}[6.1.10]
\label{II.6.1.10-fr}
Tout morphisme entier est universellement fermé.
\end{proposition}

Soit $f:X\to Y$ un morphisme entier~; en vertu de
(\hyperref[II.6.1.5-fr]{6.1.5, (iii)}), il suffit de montrer que $f$ est
fermé. Soit $Z$ une partie fermée de $X$~; il existe un sous-préschéma de $X$
ayant

\oldpage[II]{113}

pour espace sous-jacent $Z$ (\hyperref[I.5.2.1-fr]{I, 5.2.1}), et il résulte
donc de (\hyperref[II.6.1.5-fr]{6.1.5, (i) et (ii)}) qu'on peut se borner à
prouver que $f(X)$ est fermé dans $Y$. En vertu de
(\hyperref[II.6.1.5-fr]{6.1.5, (vii)}), on peut supposer $X$ et $Y$ réduits~;
en outre, si $T$ est le sous-préschéma fermé réduit de $Y$ ayant pour espace
sous-jacent $\overline{f(X)}$ (\hyperref[I.5.2.1-fr]{I, 5.2.1}), on sait que
$f$ se factorise en $X\xrightarrow{g}T\xrightarrow{j}Y$, où $j$ est le
morphisme d'injection (\hyperref[I.5.2.2-fr]{I, 5.2.2}), et comme $j$ est
séparé (\hyperref[I.5.5.1-fr]{I, 5.5.1, (i)}), il résulte de
(\hyperref[II.6.1.5-fr]{6.1.5, (v)}) que $g$ est un morphisme entier. On peut
donc supposer que $f(X)$ est dense dans $Y$. Enfin, la question étant locale
sur $Y$, on peut se borner au cas où $Y=\operatorname{Spec}(A)$. Alors
$X=\operatorname{Spec}(B)$, où $B$ est une $A$-algèbre entière sur $A$
(\hyperref[II.6.1.4-fr]{6.1.4})~; en outre $A$ est réduit
(\hyperref[I.5.1.4-fr]{I, 5.1.4}), et l'hypothèse que $f(X)$ est dense dans
$Y$ entraîne que l'homomorphisme $\varphi:A\to B$ correspondant à $f$ est
injectif (\hyperref[I.1.2.7-fr]{I, 1.2.7}). Dire que dans ces conditions
$f(X)=Y$ signifie que tout idéal premier de $A$ est la trace sur $A$ d'un
idéal premier de $B$, ce qui n'est autre que le premier théorème de
Cohen-Seidenberg (\cite[t.~I, p.~257, th.~3]{I-13}).

\begin{corollary}[6.1.11]
\label{II.6.1.11-fr}
Tout morphisme fini $f:X\to Y$ est projectif.
\end{corollary}

En effet, comme $f$ est affine, $\mathcal{O}_X$ est un
$\mathcal{O}_X$-Module très ample relativement à $f$
(\hyperref[II.5.1.2-fr]{5.1.2})~; en outre $f_*(\mathcal{O}_X)$ est un
$\mathcal{O}_Y$-Module quasi-cohérent de type fini
(\hyperref[II.6.1.2-fr]{6.1.2})~; enfin $f$ est séparé, de type fini et
universellement fermé (\hyperref[II.6.1.10-fr]{6.1.10}), et on est donc dans
les conditions d'applications du critère
(\hyperref[II.5.5.4-fr]{5.5.4, (ii)}).

\begin{proposition}[6.1.12]
\label{II.6.1.12-fr}
Soient $f:X'\to X$ un morphisme fini, et soit
$\mathcal{B}=f_*(\mathcal{O}_{X'})$ (qui est une $\mathcal{O}_X$-Algèbre
quasi-cohérente et un $\mathcal{O}_X$-Module de type fini). Soit
$\mathcal{F}'$ un $\mathcal{O}_{X'}$-Module quasi-cohérent~; pour que
$\mathcal{F}'$ soit localement libre, de rang $r$, il faut et il suffit que
$f_*(\mathcal{F}')$ soit un $\mathcal{B}$-Module localement libre de rang
$r$.
\end{proposition}

Il est clair que si $f_*(\mathcal{F}')|U$ est isomorphe à
$\mathcal{B}^r|U$ ($U$ ouvert de $X$), $\mathcal{F}'|f^{-1}(U)$ est isomorphe
à $\mathcal{O}_{X'}^r|f^{-1}(U)$
(\hyperref[II.1.4.2-fr]{1.4.2}). Inversement, supposons que $\mathcal{F}'$
soit localement libre de rang $r$ et montrons que $f_*(\mathcal{F}')$ est
localement isomorphe à $\mathcal{B}^r$ en tant que $\mathcal{B}$-Module. Soit
$x$ un point de $X$~; lorsque $U$ parcourt un système fondamental de
voisinages affines de $x$, $f^{-1}(U)$ parcourt un système fondamental de
voisinages affines (\hyperref[II.1.2.5-fr]{1.2.5}) de l'ensemble fini
$f^{-1}(x)$, puisque $f$ est fermé
(\hyperref[II.6.1.10-fr]{6.1.10}). La proposition résulte alors du lemme
suivant~:

\begin{lemma}[6.1.12.1]
\label{II.6.1.12.1-fr}
Soient $Y$ un préschéma, $\mathcal{E}$ un $\mathcal{O}_Y$-Module localement
libre de rang $r$, $Z$ une partie finie de $Y$ contenue dans un ouvert affine
$V$. Il existe alors un voisinage $U\subset V$ de $Z$ tel que
$\mathcal{E}|U$ soit isomorphe à $\mathcal{O}_Y^r|U$.
\end{lemma}

On peut évidemment supposer $Y$ affine~; pour tout $z_i\in Z$, il existe dans
l'adhérence $\overline{\{z_i\}}$ au moins un point fermé $z'_i$
(\hyperref[0.2.1.3-fr]{0, 2.1.3})~; si $Z'$ est l'ensemble des $z'_i$, tout
voisinage de $Z'$ est un voisinage de $Z$, et on peut donc supposer $Z$ fermé
dans $Y$ et discret. Considérons le sous-préschéma fermé réduit de $Y$ ayant
$Z$ pour espace sous-jacent (\hyperref[I.5.2.1-fr]{I, 5.2.1}) et soit
$j:Z\to Y$ l'injection canonique~; $j^*(\mathcal{E})=\mathcal{E}\otimes_Y
\mathcal{O}_Z$ est localement libre de rang $r$ sur le schéma discret $Z$,
et par suite est isomorphe à $\mathcal{O}_Z^r$~; autrement dit il existe $r$
sections $s_i$ ($1\leq i\leq r$) de $\mathcal{E}\otimes_Y\mathcal{O}_Z$
au-dessus de $Z$ telles que l'homomorphisme
$\mathcal{O}_Z^r\to\mathcal{E}\otimes_Y\mathcal{O}_Z$ défini par ces sections
soit bijectif. Mais $Y=\operatorname{Spec}(A)$ est affine, $Z$ est défini par
un idéal $\mathfrak{J}$ de $A$ et on a $\mathcal{E}=\widetilde{M}$, où $M$
est un $A$-module~; les $s_i$ sont des éléments

\oldpage[II]{114}

de $M\otimes_A(A/\mathfrak{J})$ et sont donc images de $r$ éléments
$t_i\in M=\Gamma(Y,\mathcal{E})$. Pour tout $z_j\in Z$ il y a alors un
voisinage $V_j$ de $z_j$ tel que les restrictions des $t_i$ à $V_j$
définissent un isomorphisme
$\mathcal{O}_Y^r|V_j\to\mathcal{E}|V_j$
(\hyperref[0.5.5.4-fr]{0, 5.5.4})~; le voisinage $U$ réunion des $V_j$ répond
donc à la question.

\begin{proposition}[6.1.13]
\label{II.6.1.13-fr}
Soient $g:X'\to X$ un morphisme entier de préschémas, $Y$ un préschéma normal
localement intègre, $f$ une application rationnelle de $Y$ dans $X'$ telle
que $g\circ f$ soit une application rationnelle partout définie
(\hyperref[I.7.2.1-fr]{I, 7.2.1})~; alors $f$ est partout définie.
\end{proposition}

Si $f_1$ et $f_2$ sont deux morphismes (d'ouverts denses de $Y$ dans $X'$) de
la classe $f$, il est clair que $g\circ f_1$ et $g\circ f_2$ sont des
morphismes équivalents, ce qui justifie la notation $g\circ f$ pour leur
classe d'équivalence. On rappelle aussi que lorsqu'on suppose en outre que
$Y$ est localement noethérien, l'hypothèse que $Y$ est normal entraîne déjà
que $Y$ est localement intègre (\hyperref[I.6.1.13-fr]{I, 6.1.13}).

Pour démontrer (\hyperref[II.6.1.13-fr]{6.1.13}), notons d'abord que la
question est locale sur $Y$ et on peut par suite supposer qu'il existe dans
la classe de $g\circ f$ un morphisme $h:Y\to X$. Considérons l'image
réciproque $Y'=X'_{(h)}=X'_{(Y)}$, et notons que le morphisme
$g'=g_{(Y)}:Y'\to Y$ est entier
(\hyperref[II.6.1.5-fr]{6.1.5, (iii)}). Vu la correspondance entre
applications rationnelles de $Y$ dans $X'$ et $Y$-sections rationnelles de
$Y'$ (\hyperref[I.7.1.2-fr]{I, 7.1.2}), on voit qu'on est ramené à prouver le
cas particulier de (\hyperref[II.6.1.13-fr]{6.1.13}) où $X=Y$, autrement dit
le

\begin{corollary}[6.1.14]
\label{II.6.1.14-fr}
Soient $X$ un préschéma normal localement intègre, $g:X'\to X$ un morphisme
entier, $f$ une $X$-section rationnelle de $X'$. Alors $f$ est partout
définie.
\end{corollary}

La question étant locale sur $X$, on peut supposer $X$ intègre, et $f$
s'identifie alors à un morphisme d'un ouvert $U$ de $X$ dans $X'$
(\hyperref[I.7.2.2-fr]{I, 7.2.2}) qui est une $U$-section de $g^{-1}(U)$.
Comme $g$ est séparé, $f$ est une immersion fermée de $U$ dans $g^{-1}(U)$
(\hyperref[I.5.4.6-fr]{I, 5.4.6})~; soit $Z$ le sous-préschéma fermé de
$g^{-1}(U)$ associé à $f$ (\hyperref[I.4.2.1-fr]{I, 4.2.1}), qui est
isomorphe à $U$, donc intègre~; soit $X_1$ le sous-préschéma réduit de $X'$
ayant pour espace sous-jacent l'adhérence $\overline{Z}$ de $Z$ dans $X'$
(\hyperref[I.5.2.1-fr]{I, 5.2.1})~; $Z$ est alors un sous-préschéma induit
sur un ouvert de $X_1$ (\hyperref[I.5.2.3-fr]{I, 5.2.3}) et comme il est
irréductible, il en est de même de $X_1$, qui est par suite intègre. Le
morphisme $f$ peut alors être considéré comme une $X$-section rationnelle de
$X_1$~; comme la restriction de $g$ à $X_1$ est un morphisme entier
(\hyperref[II.6.1.5-fr]{6.1.5, (i) et (ii)}), on est finalement ramené à
prouver (\hyperref[II.6.1.14-fr]{6.1.14}) dans le cas particulier où
$X'=X_1$, autrement dit, le

\begin{corollary}[6.1.15]
\label{II.6.1.15-fr}
Soient $X$ un préschéma intègre normal, $X'$ un préschéma intègre,
$g:X'\to X$ un morphisme entier. S'il existe une $X$-section rationnelle $f$
de $X'$, $g$ est un isomorphisme.
\end{corollary}

La question étant locale sur $X$, on peut supposer $X$ affine d'anneau
intègre $A$, et alors $X'$ est affine d'anneau $A'$ entier sur $A$
(\hyperref[II.6.1.4-fr]{6.1.4}) et intègre~; en outre, le raisonnement de
(\hyperref[II.6.1.14-fr]{6.1.14}) montre qu'il y a un ouvert dense dans $X$
isomorphe à un ouvert dense de $X'$, donc $A$ et $A'$ ont même corps de
fractions. D'autre part, en vertu de
(\hyperref[I.8.2.1.1-fr]{I, 8.2.1.1}) et de l'hypothèse que les
$\mathcal{O}_x$ sont intégralement clos, l'anneau $A$ est intégralement clos,
donc $A'=A$, ce qui achève la démonstration de
(\hyperref[II.6.1.13-fr]{6.1.13}).

\subsection{Morphismes quasi-finis.}
\label{subsection:II.6.2-fr}

\begin{proposition}[6.2.1]
\label{II.6.2.1-fr}
Soient $f:X\to Y$ un morphisme localement de type fini, $x$ un point de $X$.
Les conditions suivantes sont équivalentes~:

\oldpage[II]{115}

\begin{enumerate}
\item[a)] Le point $x$ est isolé dans sa fibre $f^{-1}(f(x))$.
\item[b)] L'anneau $\mathcal{O}_x$ est un $\mathcal{O}_{f(x)}$-module
quasi-fini (\hyperref[0.7.4.1-fr]{0, 7.4.1}).
\end{enumerate}
\end{proposition}

La question étant évidemment locale sur $X$ et sur $Y$, on peut supposer
$X=\operatorname{Spec}(A)$ et $Y=\operatorname{Spec}(B)$ affines, $A$ étant
une $B$-algèbre de type fini (\hyperref[I.6.3.3-fr]{I, 6.3.3}). En outre, on
peut remplacer $X$ par $X\times_Y\operatorname{Spec}(\mathcal{O}_{f(x)})$
sans changer la fibre $f^{-1}(f(x))$ ni l'anneau local $\mathcal{O}_x$
(\hyperref[I.3.6.5-fr]{I, 3.6.5})~; on peut donc supposer que $B$ est un
anneau local (égal à $\mathcal{O}_{f(x)}$)~; si $\mathfrak{n}$ est l'idéal
maximal de $B$, $f^{-1}(f(x))$ est alors un schéma affine d'anneau
$A/\mathfrak{n}A$, de type fini sur $k(f(x))=B/\mathfrak{n}$
(\hyperref[I.6.4.11-fr]{I, 6.4.11}). Cela étant, si a) est vérifiée, on peut
en outre supposer que $f^{-1}(f(x))$ est réduite au point $x$~; donc
$A/\mathfrak{n}A$ est de rang fini sur $B/\mathfrak{n}$
(\hyperref[I.6.4.4-fr]{I, 6.4.4}), autrement dit $A$ est un $B$-module
quasi-fini. Inversement, s'il en est ainsi, $f^{-1}(f(x))$ est un schéma
affine artinien, donc discret (\hyperref[I.6.4.4-fr]{I, 6.4.4})~; par suite
$x$ est isolé dans sa fibre, et cela montre que b) entraîne a).

\begin{corollary}[6.2.2]
\label{II.6.2.2-fr}
Soit $f:X\to Y$ un morphisme de type fini. Les conditions suivantes sont
équivalentes~:
\begin{enumerate}
\item[a)] Tout point $x\in X$ est isolé dans sa fibre $f^{-1}(f(x))$
(autrement dit le sous-espace $f^{-1}(f(x))$ est discret).
\item[b)] Pour tout $x\in X$, le préschéma $f^{-1}(f(x))$ est un
$k(f(x))$-préschéma fini.
\item[c)] Pour tout $x\in X$, l'anneau $\mathcal{O}_x$ est un
$\mathcal{O}_{f(x)}$-module quasi-fini.
\end{enumerate}
\end{corollary}

L'équivalence de a) et c) résulte de
(\hyperref[II.6.2.1-fr]{6.2.1}). D'autre part, comme $f^{-1}(f(x))$ est un
$k(f(x))$-préschéma algébrique (\hyperref[I.6.4.11-fr]{I, 6.4.11}),
l'équivalence de a) et b) résulte de (\hyperref[I.6.4.4-fr]{I, 6.4.4}).

\begin{definition}[6.2.3]
\label{II.6.2.3-fr}
Lorsque $f:X\to Y$ est un morphisme de type fini vérifiant les conditions
équivalentes de (\hyperref[II.6.2.2-fr]{6.2.2}), on dit que $f$ est
\emph{quasi-fini}, ou que $X$ est quasi-fini sur $Y$.
\end{definition}

Il est clair que tout morphisme fini est quasi-fini
(\hyperref[II.6.1.8-fr]{6.1.8}).

\begin{proposition}[6.2.4]
\label{II.6.2.4-fr}
\begin{enumerate}
\item[(i)] Une immersion $X\to Y$, qui est fermée, ou telle que $X$ soit
noethérien, est un morphisme quasi-fini.
\item[(ii)] Si $f:X\to Y$ et $g:Y\to Z$ sont des morphismes quasi-finis,
$g\circ f$ est quasi-fini.
\item[(iii)] Si $X$ et $Y$ sont des $S$-préschémas, $f:X\to Y$ un
$S$-morphisme quasi-fini, alors $f_{(S')}:X_{(S')}\to Y_{(S')}$ est quasi-fini
pour toute extension $g:S'\to S$ de la base.
\item[(iv)] Si $f:X\to Y$ et $g:X'\to Y'$ sont deux $S$-morphismes
quasi-finis,
\[
f\times_S g:X\times_S Y\to X'\times_S Y'
\]
est quasi-fini.
\item[(v)] Soient $f:X\to Y$ et $g:Y\to Z$ deux morphismes tels que $g\circ
f$ soit quasi-fini~; alors, si en outre $g$ est séparé, ou $X$ noethérien, ou
$X\times_ZY$ localement noethérien, $f$ est quasi-fini.
\item[(vi)] Si $f$ est quasi-fini, il en est de même de $f_{\mathrm{red}}$.
\end{enumerate}
\end{proposition}

Si $f:X\to Y$ est une immersion, toute fibre est réduite à un point, et
l'assertion (i) résulte donc de (\hyperref[I.6.3.4-fr]{I, 6.3.4 (i) et
6.3.5}). Pour prouver (ii), on remarque d'abord que $h=g\circ f$ est de type
fini (\hyperref[I.6.3.4-fr]{I, 6.3.4 (ii)})~; en outre, si $z=h(x)$ et
$y=f(x)$, $y$ est isolé dans $g^{-1}(z)$, donc il existe un voisinage ouvert
$V$ de $y$ dans $Y$ ne rencontrant aucun point de $g^{-1}(z)$ distinct de
$y$~; $f^{-1}(V)$ est donc un voisinage ouvert de $x$ ne rencontrant aucun

\oldpage[II]{116}

des $f^{-1}(y')$, où $y'\neq y$ est dans $g^{-1}(z)$~; comme $x$ est isolé
dans $f^{-1}(y)$, il est isolé dans
$h^{-1}(z)=f^{-1}(g^{-1}(z))$. Pour prouver (iii), on peut se limiter au cas
où $Y=S$ (\hyperref[I.3.3.11-fr]{I, 3.3.11})~; on remarque encore que tout
d'abord $f'=f_{(S')}$ est de type fini
(\hyperref[I.6.3.4-fr]{I, 6.3.4, (iii)})~; d'autre part, si
$x'\in X'=X_{(S')}$ et si on pose $y'=f'(x')$ et $y=g(y')$,
$f'^{-1}(y')$ s'identifie à
$f^{-1}(y)\otimes_{k(y)}k(y')$ (\hyperref[I.3.6.5-fr]{I, 3.6.5})~; comme
$f^{-1}(y)$ est de rang fini sur $k(y)$ par hypothèse, $f'^{-1}(y')$ est de
rang fini sur $k(y')$, donc discret. Les assertions (iv), (v), (vi) se
déduisent de (i), (ii), (iii) par la méthode générale
(\hyperref[I.5.5.12-fr]{I, 5.5.12}), sauf lorsque les hypothèses dans (v) sont
autres que l'hypothèse «~$g$ séparé~»~; pour traiter ce cas, on remarque
d'abord que si $x$ est isolé dans
$f^{-1}(g^{-1}(g(f(x))))$, il l'est a fortiori dans $f^{-1}(f(x))$~; le fait
que $f$ est de type fini résulte alors de
(\hyperref[I.6.3.6-fr]{I, 6.3.6}).

\begin{proposition}[6.2.5]
\label{II.6.2.5-fr}
Soient $A$ un anneau local noethérien complet, $X$ un $A$-schéma localement
de type fini, $x$ un point de $X$ au-dessus du point fermé $y$ de
$Y=\operatorname{Spec}(A)$, et supposons $x$ isolé dans sa fibre
$f^{-1}(y)$ ($f$ étant le morphisme structural $X\to Y$). Alors
$\mathcal{O}_x$ est un $A$-module de type fini et $X$ est $Y$-isomorphe à la
somme (I, 3.1) de
$X'=\operatorname{Spec}(\mathcal{O}_x)$ (qui est un $Y$-schéma fini) et d'un
$A$-schéma $X''$.
\end{proposition}

Il résulte de (\hyperref[II.6.2.1-fr]{6.2.1}) que $\mathcal{O}_x$ est un
$A$-module quasi-fini. Comme $\mathcal{O}_x$ est noethérien
(\hyperref[I.6.3.7-fr]{I, 6.3.7}) et l'homomorphisme
$A\to\mathcal{O}_x$ local, l'hypothèse que $A$ est complet entraîne que
$\mathcal{O}_x$ est un $A$-module de type fini
(\hyperref[0.7.4.3-fr]{0, 7.4.3}). Soit
$X'=\operatorname{Spec}(\mathcal{O}_x)$ le schéma local de $X$ au point $x$
(\hyperref[I.2.4.1-fr]{I, 2.4.1}), et soit $g:X'\to X$ le morphisme
canonique. Comme le composé $X'\xrightarrow{g}X\xrightarrow{f}Y$ est fini
(\hyperref[II.6.1.1-fr]{6.1.1}) et que $f$ est séparé, $g$ est fini
(\hyperref[II.6.1.5-fr]{6.1.5, (v)}), donc $g(X')$ est fermé dans $X$
(\hyperref[II.6.1.10-fr]{6.1.10})~; d'autre part, comme $g$ est de type fini,
$g$ est un isomorphisme local au point fermé $x'$ de $X'$ en vertu de la
définition de $g$ et de (\hyperref[I.6.5.4-fr]{I, 6.5.4})~; mais comme $X'$
est le seul voisinage ouvert de $x'$, cela signifie que $g$ est une immersion
ouverte, donc $g(X')$ est aussi ouvert dans $X$, ce qui achève la
démonstration.

\begin{corollary}[6.2.6]
\label{II.6.2.6-fr}
Soient $A$ un anneau local noethérien complet, $Y=\operatorname{Spec}(A)$,
$f:X\to Y$ un morphisme séparé et quasi-fini. Alors $X$ est $Y$-isomorphe à
une somme $X'\amalg X''$, où $X'$ est un $Y$-schéma fini et $X''$ un
$Y$-schéma quasi-fini tel que, si $y$ est le point fermé de $Y$,
$X''\cap f^{-1}(y)=\varnothing$.
\end{corollary}

En effet, la fibre $f^{-1}(y)$ est finie et discrète par hypothèse, et le
corollaire résulte donc, par récurrence sur le nombre de points de cette
fibre, de (\hyperref[II.6.2.5-fr]{6.2.5}).

\begin{remark}[6.2.7]
\label{II.6.2.7-fr}
Au chapitre V, nous verrons que si $Y$ est localement noethérien, un morphisme
quasi-fini séparé $X\to Y$ est nécessairement quasi-affine.
\end{remark}

\subsection{Fermeture intégrale d'un préschéma.}
\label{subsection:II.6.3-fr}

\begin{proposition}[6.3.1]
\label{II.6.3.1-fr}
Soient $(X,\mathcal{A})$ un espace annelé, $\mathcal{B}$ une
$\mathcal{A}$-Algèbre (commutative), $f$ une section de $\mathcal{B}$
au-dessus de $X$. Les propriétés suivantes sont équivalentes~:
\begin{enumerate}
\item[a)] Le sous-$\mathcal{A}$-Module de $\mathcal{B}$ engendré par les
$f^n$ pour $n\geq 0$ (\hyperref[0.5.1.1-fr]{0, 5.1.1}) est de type fini.
\item[b)] Il existe une sous-$\mathcal{A}$-Algèbre $\mathcal{C}$ de
$\mathcal{B}$, qui est un $\mathcal{A}$-Module de type fini, telle que
$f\in\Gamma(X,\mathcal{C})$.
\item[c)] Pour tout $x\in X$, $f_x$ est entier sur la fibre
$\mathcal{A}_x$.
\end{enumerate}
\end{proposition}

\oldpage[II]{117}

Comme le sous-$\mathcal{A}$-Module de $\mathcal{B}$ engendré par les $f^n$ est
une $\mathcal{A}$-Algèbre, il est clair que a) entraîne b). D'autre part, b)
implique que pour tout $x\in X$, le $\mathcal{A}_x$-module $\mathcal{C}_x$ est
de type fini, ce qui entraîne que tout élément de l'algèbre $\mathcal{C}_x$ et
en particulier $f_x$, est entier sur $\mathcal{A}_x$. Enfin, si pour un
$x\in X$, on a une relation de la forme
\[
f_x^n+(a_1)_xf_x^{n-1}+\ldots+(a_n)_x=0
\]
les $a_i$ ($1\leq i\leq n$) étant des sections de $\mathcal{A}$ au-dessus
d'un voisinage ouvert $U$ de $x$, la section
$f^n|U+a_1\cdot f^{n-1}|U+\ldots+a_n$ est nulle au-dessus d'un voisinage
$V\subset U$ de $x$, d'où résulte aussitôt que tous les $f^k|V$ ($k\geq0$)
sont combinaisons linéaires à coefficients dans $\Gamma(V,\mathcal{A})$ des
$f^j|V$ pour $0\leq j\leq n-1$~; on en conclut que c) entraîne a).

Lorsque les conditions équivalentes de (\hyperref[II.6.3.1-fr]{6.3.1}) sont
remplies on dit que la section $f$ est \emph{entière sur $\mathcal{A}$}.

\begin{corollary}[6.3.2]
\label{II.6.3.2-fr}
Sous les hypothèses de (\hyperref[II.6.3.1-fr]{6.3.1}), il existe un
sous-$\mathcal{A}$-Module (unique) $\mathcal{A}'$ de $\mathcal{B}$ tel que pour
tout $x\in X$, $\mathcal{A}'_x$ soit l'ensemble des germes
$f_x\in\mathcal{B}_x$ qui sont entiers sur $\mathcal{A}_x$. Pour tout ouvert
$U\subset X$, les sections de $\mathcal{A}'$ au-dessus de $U$ sont les
sections $f\in\Gamma(U,\mathcal{B})$ qui sont entières sur
$\mathcal{A}|U$.
\end{corollary}

L'existence de $\mathcal{A}'$ est immédiate, en prenant pour
$\Gamma(U,\mathcal{A}')$ l'ensemble des $f\in\Gamma(U,\mathcal{B})$ tels que
$f_x$ soit entier sur $\mathcal{A}_x$ pour tout $x\in U$. La seconde
assertion résulte aussitôt de (\hyperref[II.6.3.1-fr]{6.3.1}).

Il est clair que $\mathcal{A}'$ est une sous-$\mathcal{A}$-Algèbre de
$\mathcal{B}$~; on dit que c'est la \emph{fermeture intégrale de
$\mathcal{A}$ dans $\mathcal{B}$}.

\begin{env}[6.3.3]
\label{II.6.3.3-fr}
Soient $(X,\mathcal{A})$, $(Y,\mathcal{B})$ deux espaces annelés,
\[
g=(\psi,\theta):X\to Y
\]
un morphisme. Soit $\mathcal{C}$ (resp. $\mathcal{D}$) une
$\mathcal{A}$-Algèbre (resp. une $\mathcal{B}$-Algèbre) et soit
\[
u:\mathcal{D}\to\mathcal{C}
\]
un $g$-morphisme (\hyperref[0.4.4.1-fr]{0, 4.4.1}). Alors, si
$\mathcal{A}'$ (resp. $\mathcal{B}'$) est la fermeture intégrale de
$\mathcal{A}$ (resp. $\mathcal{B}$) dans $\mathcal{C}$ (resp.
$\mathcal{D}$), la \emph{restriction} de $u$ à $\mathcal{B}'$ est un
$g$-morphisme
\[
u':\mathcal{B}'\to\mathcal{A}'.
\]

En effet, si $j$ est l'injection canonique $\mathcal{B}'\to\mathcal{D}$, il
suffit de montrer que
\[
v=u^\sharp\circ g^*(j):g^*(\mathcal{B}')\to\mathcal{C}
\]
applique $g^*(\mathcal{B}')$ dans $\mathcal{A}'$. Or, un élément de
$(g^*(\mathcal{B}'))_x=\mathcal{B}'_{\psi(x)}\otimes_{\mathcal{B}_{\psi(x)}}
\mathcal{A}_x$ est entier sur $\mathcal{A}_x$ par définition de
$\mathcal{B}'$, donc il en est de même de son image par $v_x$, ce qui établit
notre assertion.
\end{env}

\begin{proposition}[6.3.4]
\label{II.6.3.4-fr}
Soient $X$ un préschéma, $\mathcal{A}$ une $\mathcal{O}_X$-Algèbre
quasi-cohérente. La fermeture intégrale $\mathcal{O}'_X$ de
$\mathcal{O}_X$ dans $\mathcal{A}$ est alors une $\mathcal{O}_X$-Algèbre
quasi-cohérente, et pour tout ouvert affine $U$ de $X$,
$\Gamma(U,\mathcal{O}'_X)$ est la fermeture intégrale de
$\Gamma(U,\mathcal{O}_X)$ dans $\Gamma(U,\mathcal{A})$.
\end{proposition}

On peut se borner au cas où $X=\operatorname{Spec}(B)$ est affine et
$\mathcal{A}=\widetilde{A}$, $A$ étant une

\oldpage[II]{118}

$B$-algèbre~; soit $B'$ la fermeture intégrale de $B$ dans $A$. Tout revient
à voir que pour tout $x\in X$, un élément de $A_x$, entier sur $B_x$,
appartient nécessairement à $B'_x$, ce qui résulte du fait que, pour un anneau
commutatif $C$, les opérations de fermeture intégrale dans une $C$-algèbre et
de passage à un anneau de fractions (pour une partie multiplicative de $C$)
commutent (\cite[t.~I, p.~261 et 257]{I-13}).

Le $X$-schéma $X'=\operatorname{Spec}(\mathcal{O}'_X)$ est alors appelé la
\emph{fermeture intégrale de $X$ relativement à $\mathcal{A}$} (ou
\emph{relativement à $\operatorname{Spec}(\mathcal{A})$})~; il est clair que
$X'$ est \emph{entier} sur $X$ (\hyperref[II.6.1.2-fr]{6.1.2}).

On déduit aussitôt de (\hyperref[II.6.3.4-fr]{6.3.4}) que si $f:X'\to X$ est
le morphisme structural, alors, pour tout ouvert $U$ de $X$, $f^{-1}(U)$ est
la \emph{fermeture intégrale du préschéma induit par $X$ sur $U$, relativement
à $\mathcal{A}|U$}.

\begin{env}[6.3.5]
\label{II.6.3.5-fr}
Soient $X$, $Y$ deux préschémas, $f:X\to Y$ un morphisme, $\mathcal{A}$
(resp. $\mathcal{B}$) une $\mathcal{O}_X$-Algèbre (resp. une
$\mathcal{O}_Y$-Algèbre) quasi-cohérente, et soit
$u:\mathcal{B}\to\mathcal{A}$ un $f$-morphisme. On a vu
(\hyperref[II.6.3.3-fr]{6.3.3}) qu'on en déduit un $f$-morphisme
$u':\mathcal{O}'_Y\to\mathcal{O}'_X$, où $\mathcal{O}'_X$ (resp.
$\mathcal{O}'_Y$) est la fermeture intégrale de $\mathcal{O}_X$ (resp.
$\mathcal{O}_Y$) dans $\mathcal{A}$ (resp. $\mathcal{B}$). Par suite, si
$X'$ (resp. $Y'$) est la fermeture intégrale de $X$ (resp. $Y$) relativement
à $\mathcal{A}$ (resp. $\mathcal{B}$), on déduit canoniquement de $u$ un
morphisme $f'=\operatorname{Spec}(u'):X'\to Y'$
(\hyperref[II.1.5.6-fr]{1.5.6}) rendant commutatif le diagramme
\[
\label{II.6.3.5.1-fr}
\xymatrix{
X'\ar[r]^{f'}\ar[d]&Y'\ar[d]\\
X\ar[r]_{f}&Y
}
\tag{6.3.5.1}
\]
\end{env}

\begin{env}[6.3.6]
\label{II.6.3.6-fr}
Supposons que $X$ n'ait qu'un nombre \emph{fini} de composantes irréductibles
$X_i$ ($1\leq i\leq r$), de points génériques $\xi_i$, et considérons en
particulier la fermeture intégrale de $X$ relativement à une
$\mathcal{R}(X)$-Algèbre $\mathcal{A}$ quasi-cohérente (en tant que
$\mathcal{O}_X$-Algèbre ou en tant que $\mathcal{R}(X)$-Algèbre, ce qui
revient au même). On sait (\hyperref[I.7.3.5-fr]{I, 7.3.5}) que
$\mathcal{A}$ est composée directe de $r$ $\mathcal{O}_X$-Algèbres
quasi-cohérentes $\mathcal{A}_i$, le support de $\mathcal{A}_i$ étant contenu
dans $X_i$, et le faisceau induit par $\mathcal{A}_i$ sur $X_i$ étant un
faisceau simple dont la fibre $A_i$ est une algèbre sur
$\mathcal{O}_{\xi_i}$. Il est clair alors (\hyperref[II.6.3.4-fr]{6.3.4}) que
la fermeture intégrale $\mathcal{O}'_X$ de $\mathcal{O}_X$ dans
$\mathcal{A}$ est \emph{composée directe} des fermetures intégrales
$\mathcal{O}^{(i)}_X$ de $\mathcal{O}_X$ dans chacune des $\mathcal{A}_i$, et
par suite la fermeture intégrale $X'=\operatorname{Spec}(\mathcal{O}'_X)$ de
$X$ relativement à $\mathcal{A}$ est un $X$-schéma \emph{somme} des
$\operatorname{Spec}(\mathcal{O}^{(i)}_X)=X'_i$ ($1\leq i\leq r$).

Supposons en outre que la $\mathcal{O}_X$-Algèbre $\mathcal{A}$ soit
\emph{réduite}, ou, ce qui revient au même, que chacune des algèbres $A_i$
soit réduite, et puisse par suite être considérée comme une algèbre sur le
\emph{corps} $k(\xi_i)$ (égal au corps des fonctions rationnelles du
sous-préschéma réduit de $X$ ayant $X_i$ pour espace sous-jacent)~; alors
(\hyperref[II.1.3.8-fr]{1.3.8}) chacun des $X'_i$ est un $X$-schéma
\emph{réduit} et $X'$ est aussi la fermeture intégrale de $X_{\mathrm{red}}$.
Supposons de plus que chacune des algèbres $A_i$ soit \emph{composée directe
d'un nombre fini de corps} $K_{ij}$ ($1\leq j\leq s_i$)~; si
$\mathcal{K}_{ij}$ est la sous-Algèbre de $\mathcal{A}_i$ correspondant à
$K_{ij}$, il est clair que $\mathcal{O}^{(i)}_X$ est \emph{composée directe}
des fermetures intégrales $\mathcal{O}^{(ij)}_X$ de $\mathcal{O}_X$ dans
chacune des $\mathcal{K}_{ij}$. Par suite, $X'_i$ est alors \emph{somme} des
$X$-schémas $X'_{ij}=\operatorname{Spec}(\mathcal{O}^{(ij)}_X)$
($1\leq j\leq s_i$). En outre, sous ces hypothèses et avec ces notations~:
\end{env}

\oldpage[II]{119}

\begin{proposition}[6.3.7]
\label{II.6.3.7-fr}
Chacun des $X'_{ij}$ est un $X$-préschéma intègre et normal, et son corps des
fonctions rationnelles $R(X'_{ij})$ s'identifie canoniquement à la fermeture
algébrique $K'_{ij}$ de $k(\xi_i)$ dans $K_{ij}$.
\end{proposition}

En vertu de ce qui précède, on peut supposer $X$ intègre, donc $r=1$, et
$s_1=1$, de sorte que l'unique algèbre $A_1$ est un corps $K$~; soit $\xi$ le
point générique de $X$, et soit $f:X'\to X$ le morphisme structural. Pour tout
ouvert affine non vide $U$ de $X$, $f^{-1}(U)$ s'identifie au spectre de la
fermeture intégrale $B'_U$ dans le corps $K$ de l'anneau intègre
$B_U=\Gamma(U,\mathcal{O}_X)$ (\hyperref[II.6.3.4-fr]{6.3.4})~; comme l'anneau
$B'_U$ est intègre et intégralement clos, il en est de même des anneaux locaux
des points de son spectre, donc $f^{-1}(U)$ est par définition un schéma
\emph{intègre et normal} ((\hyperref[0.4.1.4-fr]{0, 4.1.4}) et
(\hyperref[I.5.1.4-fr]{I, 5.1.4})). En outre, comme $(0)$ est le seul idéal
premier de $B'_U$ au-dessus de l'idéal premier $(0)$ de $B_U$
(\cite[t.~I, p.~259]{I-13}), $f^{-1}(\xi)$ se réduit à un seul point $\xi'$,
et $k(\xi')$ est le corps des fractions $K'$ de $B'_U$, qui n'est autre que la
fermeture algébrique de $k(\xi)$ dans $K$. Enfin, $X'$ est irréductible, car
lorsque $U$ parcourt l'ensemble des ouverts affines non vides de $X$, les
$f^{-1}(U)$ constituent un recouvrement ouvert de $X'$ formé d'ouverts
irréductibles~; en outre l'intersection $f^{-1}(U\cap V)$ de deux de ces
ouverts contient $\xi'$, donc n'est pas vide, et on conclut par
(\hyperref[0.2.1.4-fr]{0, 2.1.4}).

\begin{corollary}[6.3.8]
\label{II.6.3.8-fr}
Soit $X$ un préschéma réduit n'ayant qu'un nombre fini de composantes
irréductibles $X_i$ ($1\leq i\leq r$), et soit $\xi_i$ le point générique de
$X_i$. La fermeture intégrale $X'$ de $X$ relativement à $\mathcal{R}(X)$ est
somme de $r$ $X$-schémas $X'_i$ qui sont intègres et normaux. Si $f:X'\to X$
est le morphisme structural, $f^{-1}(\xi_i)$ se réduit au point générique
$\xi'_i$ de $X'_i$ et on a $k(\xi'_i)=k(\xi_i)$, autrement dit $f$ est
birationnel.
\end{corollary}

On dit dans ce cas que $X'$ est le \emph{normalisé} du préschéma réduit $X$~;
on notera que $f$, étant birationnel et entier, est \emph{surjectif}
(\hyperref[II.6.1.10-fr]{6.1.10}). Pour que l'on ait alors $X'=X$, il faut et
il suffit que $X$ soit normal. Lorsque $X$ est un préschéma intègre, il résulte
de (\hyperref[II.6.3.8-fr]{6.3.8}) que son normalisé $X'$ est intègre.

\begin{env}[6.3.9]
\label{II.6.3.9-fr}
Soient $X$, $Y$ deux préschémas intègres, $f:X\to Y$ un morphisme dominant,
$L=R(X)$, $K=R(Y)$ les corps des fonctions rationnelles de $X$ et $Y$~; il
correspond canoniquement à $f$ une injection $K\to L$, et lorsqu'on identifie
$K$ (resp. $L$) au faisceau simple $\mathcal{R}(Y)$ (resp.
$\mathcal{R}(X)$), cette injection est un $f$-morphisme. Soit $K_1$ (resp.
$L_1$) une extension de $K$ (resp. $L$) et supposons donné un monomorphisme
$K_1\to L_1$ tel que le diagramme
\[
\xymatrix{
K_1\ar[r]&L_1\\
K\ar[u]\ar[r]&L\ar[u]
}
\]
soit commutatif~; si $K_1$ (resp. $L_1$) est considéré comme faisceau simple
sur $Y$ (resp. $X$), donc comme une $\mathcal{R}(Y)$-Algèbre (resp. une
$\mathcal{R}(X)$-Algèbre), cela signifie que $K_1\to L_1$ est un
$f$-morphisme. Cela étant, si $X'$ (resp. $Y'$) est la fermeture intégrale de
$X$ (resp. $Y$) relativement à $L_1$ (resp. $K_1$), $X'$ (resp. $Y'$) est un
préschéma intègre et normal (\hyperref[II.6.3.6-fr]{6.3.6}) dont le corps des
fonctions rationnelles s'identifie canoniquement à la fermeture algébrique

\oldpage[II]{120}

$L'$ (resp. $K'$) de $L$ (resp. $K$) dans $L_1$ (resp. $K_1$), et il existe
un morphisme canonique (nécessairement dominant) $f':X'\to Y'$ rendant
commutatif le diagramme (\hyperref[II.6.3.5.1-fr]{6.3.5.1}).

Le cas le plus important est celui où on prend $L_1=L$, $K_1$ étant donc une
extension de $K$ contenue dans $L$, et où on suppose $X$ intègre et
\emph{normal}, donc $X'=X$. Ce qui précède montre alors que lorsque $X$ est
normal, et que $Y'$ est la fermeture intégrale de $Y$ relative à un corps
$K_1\subset L=R(X)$, tout morphisme dominant $f:X\to Y$ se factorise en
\[
f:X\xrightarrow{f'}Y'\to Y
\]
où $f'$ est dominant~; d'ailleurs, lorsque le monomorphisme $K_1\to L$ est
donné, $f'$ est nécessairement unique, comme on le voit en se ramenant au cas
où $X$ et $Y$ sont affines. On voit donc que pour la donnée de $Y$, $L$ et
d'un $K$-monomorphisme $K_1\to L$, la fermeture intégrale $Y'$ de $Y$
relativement à $K_1$ est solution d'un \emph{problème universel}.
\end{env}

\begin{remark}[6.3.10]
\label{II.6.3.10-fr}
Reprenons les hypothèses de (\hyperref[II.6.3.6-fr]{6.3.6}), en supposant de
plus que chacune des algèbres $A_i$ soit de \emph{rang fini} sur $k(\xi_i)$
(ce qui entraîne que $A_i$ est composée directe d'un nombre fini de corps)~;
on peut affirmer dans certains cas que le morphisme structural $X'\to X$ est
non seulement entier, mais même \emph{fini}. Bornons-nous au cas où $X$ est
\emph{réduit}~; la question étant locale sur $X$, on peut en outre supposer
que $X$ est affine d'anneau $C$, et que $C$ n'a qu'un nombre fini d'idéaux
minimaux $\mathfrak{p}_i$ ($1\leq i\leq r$) tels que les
$C_i=C/\mathfrak{p}_i$ soient intègres~; $X'$ sera alors fini sur $X$ si la
fermeture intégrale de chaque $C_i$ dans toute extension de degré fini de son
corps des fractions est un $C$-module de type fini
(\hyperref[II.6.3.4-fr]{6.3.4}). On sait que cette condition est toujours
vérifiée lorsque $C$ est une \emph{algèbre de type fini sur un corps}
(\cite[t.~I, p.~267, th.~9]{I-13}), ou \emph{sur $\mathbf{Z}$}
(\cite[t.~I, p.~93, th.~3]{I-9}), ou \emph{sur un anneau local noethérien
complet} ([25], p.~298). On en conclut que $X'\to X$ sera un morphisme
fini lorsque $X$ est un schéma de \emph{type fini sur un corps}, ou sur
$\mathbf{Z}$, ou sur un anneau local noethérien complet.
\end{remark}

\subsection{Déterminant d'un endomorphisme de $\mathcal{O}_X$-Module.}
\label{subsection:II.6.4-fr}

\begin{env}[6.4.1]
\label{II.6.4.1-fr}
Soient $A$ un anneau (commutatif), $E$ un $A$-module libre de rang $n$, $u$
un endomorphisme de $E$~; rappelons que pour définir le \emph{polynôme
caractéristique} de $u$, on considère l'endomorphisme $u\otimes1$ du
$A[T]$-module libre de rang $n$, $E\otimes_AA[T]$ ($T$ indéterminée), et on
pose
\[
\label{II.6.4.1.1-fr}
P(u,T)=\det(T\cdot I-(u\otimes1))
\tag{6.4.1.1}
\]
($I$ automorphisme identique de $E\otimes_AA[T]$). On a
\[
\label{II.6.4.1.2-fr}
P(u,T)=T^n-\sigma_1(u)T^{n-1}+\ldots+(-1)^n\sigma_n(u)
\tag{6.4.1.2}
\]
où $\sigma_i(u)$ est un élément de $A$, égal à un polynôme homogène de degré
$i$ (à coefficients entiers) par rapport aux éléments de la matrice de $u$
relative à une base quelconque de $E$~; on dit que les $\sigma_i(u)$ sont les
\emph{fonctions symétriques élémentaires de $u$}, et on a en particulier
$\sigma_1(u)=\operatorname{Tr}u$ et $\sigma_n(u)=\det u$. Rappelons qu'en
vertu du th. de Hamilton-Cayley, on a
\[
\label{II.6.4.1.3-fr}
P(u,u)=u^n-\sigma_1(u)u^{n-1}+\ldots+(-1)^n\sigma_n(u)=0
\tag{6.4.1.3}
\]

\oldpage[II]{121}

ce qui s'écrit aussi
\[
\label{II.6.4.1.4-fr}
(\det u)\cdot I_E=uQ(u)
\tag{6.4.1.4}
\]
($I_E$ automorphisme identique de $E$), avec
\[
\label{II.6.4.1.5-fr}
Q(u)=(-1)^{n+1}(u^{n-1}-\sigma_1(u)u^{n-2}+\ldots+(-1)^{n-1}\sigma_{n-1}(u)).
\tag{6.4.1.5}
\]

Soit $\varphi:A\to B$ un homomorphisme d'anneaux, faisant de $B$ une
$A$-algèbre~; considérons le $B$-module
$E_{(B)}=E\otimes_AB$ qui est libre de rang $n$, et l'extension $u\otimes1$
de $u$ à un endomorphisme de $E_{(B)}$~; il est immédiat que l'on a
$\sigma_i(u\otimes1)=\varphi(\sigma_i(u))$ pour tout indice $i$.
\end{env}

\begin{env}[6.4.2]
\label{II.6.4.2-fr}
Supposons maintenant que $A$ soit un anneau \emph{intègre}, de corps des
fractions $K$, et $E$ un $A$-module de \emph{type fini} (mais non
nécessairement libre). Soit $n$ le \emph{rang} de $E$, c'est-à-dire le rang du
$K$-module libre $E\otimes_AK$~; à tout endomorphisme $u$ de $E$ correspond
canoniquement l'endomorphisme $u\otimes1$ de $E\otimes_AK$~; par abus de
langage, on appelle encore \emph{polynôme caractéristique} de $u$ et on note
$P(u,T)$ le polynôme $P(u\otimes1,T)$, dont les coefficients
$\sigma_i(u\otimes1)$ (qui appartiennent à $K$) seront encore notés
$\sigma_i(u)$ et appelés les \emph{fonctions symétriques élémentaires de
$u$}~; en particulier $\det u=\det(u\otimes1)$ par définition. Avec ces
notations, les formules (\hyperref[II.6.4.1.3-fr]{6.4.1.3}) à
(\hyperref[II.6.4.1.5-fr]{6.4.1.5}) ont un sens et sont encore valables,
lorsque l'on interprète $u^j$ comme l'homomorphisme $E\to E\otimes_AK$
composé de l'endomorphisme $u^j\otimes1=(u\otimes1)^j$ de $E\otimes_AK$ et
de l'homomorphisme canonique $x\mapsto x\otimes1$.

Si $F$ est le sous-module de torsion de $E$, et si $E_0=E/F$, on a
$u(F)\subset F$, donc, par passage aux quotients, $u$ donne un endomorphisme
$u_0$ de $E_0$~; en outre $E\otimes_AK$ s'identifie à $E_0\otimes_AK$ et
$u\otimes1$ à $u_0\otimes1$, donc $\sigma_i(u)=\sigma_i(u_0)$ pour
$1\leq i\leq n$.

Lorsque $E$ est \emph{sans torsion}, $E$ s'identifie canoniquement à un
sous-$A$-module de $E\otimes_AK$, et la relation $u\otimes1=0$ est
équivalente à $u=0$. Lorsque $E$ est un $A$-module libre, les deux définitions
des $\sigma_i(u)$ données dans (\hyperref[II.6.4.1-fr]{6.4.1}) et
(\hyperref[II.6.4.2-fr]{6.4.2}) coïncident d'après ce qui précède, ce qui
justifie les notations adoptées.

On notera aussi que lorsque $E$ est un \emph{module de torsion}, autrement dit
$E_0=\{0\}$, l'algèbre extérieure de $E_0$ est réduite à $K$ et le déterminant
de l'unique endomorphisme $u_0$ de $E_0$ est \emph{égal à $1$}.
\end{env}

\begin{proposition}[6.4.3]
\label{II.6.4.3-fr}
Soient $A$ un anneau intègre, $E$ un $A$-module de type fini, $u$ un
endomorphisme de $E$~; alors les fonctions symétriques élémentaires
$\sigma_i(u)$ de $u$ (et en particulier $\det u$) sont des éléments de $K$
entiers sur $A$.
\end{proposition}

Soit $A'$ la clôture intégrale de $A$~; comme $A'[T]$ est un anneau
intégralement clos ([24], p.~99), c'est la clôture intégrale de
$A[T]$ dans son corps des fractions $K(T)$. Remplaçant $u$ par
$T\cdot I-u\otimes1$ et $A$ par $A[T]$, on voit qu'on est ramené à prouver que
$\det u$ est entier sur $A$. Si $n$ est le rang de $E$, on a
$\det u=\det(\bigwedge^n u)$ et
$(\bigwedge^n u)\otimes1=\bigwedge^n(u\otimes1)$, donc on peut supposer
$n=1$. Mais alors l'application $u\mapsto\det u$ est un homomorphisme du
$A$-module $\operatorname{Hom}_A(E,E)$ dans $K$~; comme $E$ est de type fini,
$\operatorname{Hom}_A(E,E)$ est isomorphe à un sous-module du $A$-module de
type fini $E^n$ (si $E$ admet un système de $n$ générateurs)~; donc les
éléments $\det u$ appartiennent à un sous-$A$-module de $K$ de type fini, et
sont par suite entiers sur $A$.

\oldpage[II]{122}

\begin{corollary}[6.4.4]
\label{II.6.4.4-fr}
Sous les hypothèses de (\hyperref[II.6.4.3-fr]{6.4.3}), si on suppose en outre
$A$ normal, les $\sigma_i(u)$ (en particulier $\operatorname{Tr}u$ et
$\det u$) appartiennent à $A$.
\end{corollary}

\begin{proposition}[6.4.5]
\label{II.6.4.5-fr}
Soient $A$ un anneau intègre, $E$ un $A$-module de type fini, de rang $n$, $u$
un endomorphisme de $E$ tel que les $\sigma_i(u)$ appartiennent à $A$. Pour
que $u$ soit un automorphisme de $E$, il est nécessaire que $\det u$ soit
inversible dans $A$~; cette condition est suffisante lorsque $E$ est sans
torsion.
\end{proposition}

La condition est \emph{suffisante}, car l'hypothèse et les formules
(\hyperref[II.6.4.1.4-fr]{6.4.1.4}) et
(\hyperref[II.6.4.1.5-fr]{6.4.1.5}) (valables dans $E$, et non seulement dans
$E\otimes_AK$, puisque $E$ est sans torsion) prouvent que
$(\det u)^{-1}Q(u)$ est l'inverse de $u$.

La condition est \emph{nécessaire}, car si $u$ est inversible, il résulte de
(\hyperref[II.6.4.3-fr]{6.4.3}) que $\det(u^{-1})$ appartient à la clôture
intégrale $A'$ de $A$ dans son corps des fractions $K$ et est évidemment
inverse de $\det u$ dans $A'$. Notre assertion résulte alors du~:

\begin{lemma}[6.4.5.1]
\label{II.6.4.5.1-fr}
Soit $A$ un sous-anneau d'un anneau $A'$ tel que $A'$ soit entier sur $A$. Si
un élément $x\in A$ est inversible dans $A'$, il est aussi inversible dans
$A$.
\end{lemma}

Dans le cas contraire, $x$ appartiendrait à un idéal maximal $\mathfrak{m}$
de $A$ et il résulte du premier th. de Cohen-Seidenberg
(\cite[t.~I, p.~257, th.~3]{I-13}) qu'il y aurait un idéal maximal
$\mathfrak{m}'$ de $A'$ tel que $\mathfrak{m}=\mathfrak{m}'\cap A$~; on aurait
donc $x\in\mathfrak{m}'$, ce qui est absurde.

\begin{corollary}[6.4.6]
\label{II.6.4.6-fr}
Soient $A$ un anneau intègre et intégralement clos, $E$ un $A$-module de type
fini sans torsion, $u$ un endomorphisme de $E$. Pour que $u$ soit un
automorphisme de $E$, il faut et il suffit que $\det u$ soit inversible dans
$A$.
\end{corollary}

Cela résulte de (\hyperref[II.6.4.4-fr]{6.4.4}) et
(\hyperref[II.6.4.5-fr]{6.4.5}).

\begin{remark}[6.4.7]
\label{II.6.4.7-fr}
Nous aurons besoin plus loin d'une généralisation des résultats précédents.
Considérons un anneau noethérien \emph{réduit} $A$~; soient
$\mathfrak{p}_\alpha$ ($1\leq\alpha\leq r$) ses idéaux minimaux, $K_\alpha$
le corps des fractions de l'anneau intègre $A/\mathfrak{p}_\alpha$, $K$
l'anneau total des fractions de $A$, \emph{composé direct} des corps
$K_\alpha$. Soit $E$ un $A$-module de type fini, et \emph{supposons que}
$E\otimes_AK$ soit un $K$-module libre de rang $n$ (ce qui ici n'est plus
conséquence des autres hypothèses)~; il revient au même de dire que tous les
$K_\alpha$-espaces vectoriels $E\otimes_AK_\alpha=E_\alpha$ ont \emph{même
dimension $n$}~; si alors $u$ est un endomorphisme de $E$, on pose encore
$P(u,T)=P(u\otimes1,T)$ et $\sigma_j(u)=\sigma_j(u\otimes1)$, et en
particulier $\det u=\det(u\otimes1)$~; les $\sigma_j(u)$ sont donc des
éléments de $K$. Il est immédiat que $E\otimes_AK$ est somme directe des
$E_\alpha$ et que ces derniers sont stables par $u\otimes1$~; la restriction
de $u\otimes1$ à $E_\alpha$ n'est autre d'ailleurs que l'extension $u_\alpha$
de $u$ à ce $K_\alpha$-espace vectoriel~; on en conclut que $\sigma_j(u)$ est
l'élément de $K$ dont les composantes dans les $K_\alpha$ sont les
$\sigma_j(u_\alpha)$. Comme la fermeture intégrale de $A$ dans $K$ est
\emph{composée directe} des fermetures intégrales de $A$ dans les
$K_\alpha$, les $\sigma_j(u)$ sont \emph{entiers} sur $A$.

\begin{lemma}[6.4.7.1]
\label{II.6.4.7.1-fr}
La sous-$A$-algèbre de $K$ engendrée par tous les éléments $\sigma_j(u)$
($1\leq j\leq n$) lorsque $u$ parcourt $\operatorname{Hom}_A(E,E)$, est un
$A$-module de type fini.
\end{lemma}

Il suffit de démontrer que la sous-$A[T]$-algèbre de $K[T]$ engendrée par les
$P(u,T)$ est un $A[T]$-module de type fini, car si les $F_i(T)$
($1\leq i\leq m$) forment un système de générateurs de ce $A[T]$-module, les
coefficients des $F_i(T)$ sont entiers sur $A$, donc engendrent une
$A$-algèbre qui est un $A$-module de type fini
(\cite[t.~I, p.~255, th.~1]{I-13}). On peut donc

\oldpage[II]{123}

remplacer $A$ par $A[T]$ (qui est noethérien) et $E$ par
$E\otimes_AA[T]=E'$ qui est tel que
$E'\otimes_{A[T]}K[T]=E\otimes_AK[T]$ soit un $K[T]$-module libre de rang $n$.
Revenant aux notations initiales, on voit donc qu'il suffit de prouver que le
$A$-module engendré par les éléments $\det u$, où $u$ parcourt
$\operatorname{Hom}_A(E,E)$, est de type fini~; \emph{a fortiori} (puisque
tout sous-module d'un $A$-module de type fini est de type fini) il suffit de
prouver que lorsque $v$ parcourt l'ensemble des endomorphismes de
$\bigwedge^nE$, le $A$-module engendré par les $\det v$ est de type fini~;
autrement dit, on est encore ramené au cas où $n=1$. Mais alors la proposition
résulte de ce que $\operatorname{Hom}_A(E,E)$ est un $A$-module de type fini
et de ce que $v\mapsto\det v$ est un homomorphisme de ce $A$-module dans $K$.

Soit $F$ le noyau de l'homomorphisme canonique $E\to E\otimes_AK$ et soit
$E_0=E/F$~; on voit comme ci-dessus que $E\otimes_AK$ s'identifie à
$E_0\otimes_AK$, que $u(F)\subset F$, et que si $u_0$ est l'endomorphisme de
$E_0$ déduit de $u$ par passage aux quotients, $u\otimes1$ s'identifie à
$u_0\otimes1$, et $\sigma_j(u)=\sigma_j(u_0)$ pour tout $j$. \emph{Si l'on a
$F=0$}, les formules (\hyperref[II.6.4.1.3-fr]{6.4.1.3}) et
(\hyperref[II.6.4.1.5-fr]{6.4.1.5}) ont un sens et sont valables lorsque $E$
est identifié à un sous-module de $E\otimes_AK$ et les $u^j$ à des
homomorphismes $E\to E\otimes_AK$~; par suite, la prop.
(\hyperref[II.6.4.5-fr]{6.4.5}) s'étend aussi à ce cas, avec sa démonstration.
\end{remark}

\begin{env}[6.4.8]
\label{II.6.4.8-fr}
Soient $(X,\mathcal{A})$ un espace annelé, $\mathcal{E}$ un
$\mathcal{A}$-Module \emph{localement libre} (de rang fini), $u$ un
endomorphisme de $\mathcal{E}$. Il y a par hypothèse une base $\mathfrak{B}$
de la topologie de $X$ telle que pour tout $V\in\mathfrak{B}$,
$\mathcal{E}|V$ soit isomorphe à $\mathcal{A}^n|V$ (pour un $n$ pouvant varier
avec $V$). Soit $u$ un endomorphisme de $\mathcal{E}$~; pour tout
$V\in\mathfrak{B}$, $u_V$ est donc un endomorphisme du
$\Gamma(V,\mathcal{A})$-module $\Gamma(V,\mathcal{E})$, qui est libre par
hypothèse~; le déterminant $\det u_V$ est donc défini et appartient à
$\Gamma(V,\mathcal{A})$. En outre, si $e_1,\ldots,e_n$ forment une base de
$\Gamma(V,\mathcal{E})$, leurs restrictions à tout ouvert $W\subset V$
forment une base de $\Gamma(W,\mathcal{E})$ sur $\Gamma(W,\mathcal{A})$,
donc $\det u_W$ est la restriction de $\det u_V$ à $W$. Il existe donc une
section et une seule de $\mathcal{A}$ au-dessus de $X$, que nous noterons
$\det u$ et appellerons le \emph{déterminant de $u$}, telle que la restriction
de $\det u$ à tout $V\in\mathfrak{B}$ soit $\det u_V$. Il est clair que pour
tout $x\in X$, on a $(\det u)_x=\det u_x$~; pour deux endomorphismes $u$, $v$
de $\mathcal{E}$, on a
\[
\label{II.6.4.8.1-fr}
\det(u\circ v)=(\det u)(\det v)
\tag{6.4.8.1}
\]
ainsi que
\[
\label{II.6.4.8.2-fr}
\det(1_{\mathcal{E}})=1_{\mathcal{A}}
\tag{6.4.8.2}
\]
et, si le rang de $\mathcal{E}$ est constant (ce qui sera le cas
(\hyperref[0.5.4.1-fr]{0, 5.4.1}) pour $X$ connexe) et égal à $n$
\[
\label{II.6.4.8.3-fr}
\det(s\cdot u)=s^n\det u
\tag{6.4.8.3}
\]
pour tout $s\in\Gamma(X,\mathcal{A})$ (on notera que
$\det(0)=0_{\mathcal{A}}$ si $n\geq1$, mais $\det(0)=1_{\mathcal{A}}$ pour
$n=0$). En outre, pour que $u$ soit un \emph{automorphisme} de
$\mathcal{E}$, il faut et il suffit que $\det u$ soit \emph{inversible} dans
$\Gamma(X,\mathcal{A})$.

Si le rang de $\mathcal{E}$ est constant, on peut définir de la même manière
les fonctions symétriques élémentaires $\sigma_i(u)$, qui sont des éléments
de $\Gamma(X,\mathcal{A})$~; on a encore les relations
(\hyperref[II.6.4.1.3-fr]{6.4.1.3}) à
(\hyperref[II.6.4.1.5-fr]{6.4.1.5}).

\oldpage[II]{124}

On a ainsi défini un homomorphisme
$\det:\operatorname{Hom}_{\mathcal{A}}(\mathcal{E},\mathcal{E})
\to\Gamma(X,\mathcal{A})$ de monoïdes multiplicatifs~; si on note que
$\operatorname{Hom}_{\mathcal{A}}(\mathcal{E},\mathcal{E})=
\Gamma(X,\mathcal{H}om_{\mathcal{A}}(\mathcal{E},\mathcal{E}))$ par
définition, on voit qu'on peut dans cette définition remplacer $X$ par un
ouvert quelconque $U$ de $X$, ce qui définit immédiatement un homomorphisme
$\det:\mathcal{H}om_{\mathcal{A}}(\mathcal{E},\mathcal{E})\to\mathcal{A}$ de
faisceaux de monoïdes multiplicatifs. Lorsque $\mathcal{E}$ a un rang
constant, on définit de même des homomorphismes
$\sigma_i:\mathcal{H}om_{\mathcal{A}}(\mathcal{E},\mathcal{E})
\to\mathcal{A}$ de faisceaux d'ensembles~; pour $i=1$, l'homomorphisme
$\sigma_1=\operatorname{Tr}$ est un homomorphisme de $\mathcal{A}$-Modules.

Soit $(Y,\mathcal{B})$ un second espace annelé, et soit
$f:(X,\mathcal{A})\to(Y,\mathcal{B})$ un morphisme d'espaces annelés~; si
$\mathcal{F}$ est un $\mathcal{B}$-Module localement libre,
$f^*(\mathcal{F})$ est un $\mathcal{A}$-Module localement libre (qui a même
rang que $\mathcal{F}$ si ce dernier est de rang constant)
(\hyperref[0.5.4.5-fr]{0, 5.4.5}). Pour tout endomorphisme $v$ de
$\mathcal{F}$, $f^*(v)$ est un endomorphisme de $f^*(\mathcal{F})$, et il
résulte aussitôt des définitions que $\det f^*(v)$ est la section de
$\mathcal{A}=f^*(\mathcal{B})$ au-dessus de $X$ qui correspond canoniquement
à $\det v\in\Gamma(Y,\mathcal{B})$. On peut encore dire que l'homomorphisme
$f^*(\det):f^*(\mathcal{H}om_{\mathcal{B}}(\mathcal{F},\mathcal{F}))
\to f^*(\mathcal{B})=\mathcal{A}$ est le composé
\[
f^*(\mathcal{H}om_{\mathcal{B}}(\mathcal{F},\mathcal{F}))
\xrightarrow{\gamma^\sharp}
\mathcal{H}om_{\mathcal{A}}(f^*(\mathcal{F}),f^*(\mathcal{F}))
\xrightarrow{\det}\mathcal{A}
\]
(\hyperref[0.4.4.6-fr]{0, 4.4.6}). On a des résultats analogues pour les
$\sigma_i$.
\end{env}

\begin{env}[6.4.9]
\label{II.6.4.9-fr}
Supposons maintenant que $X$ soit un \emph{préschéma localement intègre}, de
sorte que le faisceau $\mathcal{R}(X)$ des fonctions rationnelles sur $X$ est
un faisceau de corps localement simple (\hyperref[I.7.3.4-fr]{I, 7.3.4}),
quasi-cohérent en tant que $\mathcal{O}_X$-Module. Si $\mathcal{E}$ est un
$\mathcal{O}_X$-Module quasi-cohérent \emph{de type fini},
$\mathcal{E}'=\mathcal{E}\otimes_{\mathcal{O}_X}\mathcal{R}(X)$ est alors un
$\mathcal{R}(X)$-Module localement libre
(\hyperref[I.7.3.6-fr]{I, 7.3.6})~; pour tout endomorphisme $u$ de
$\mathcal{E}$, $u\otimes1_{\mathcal{R}(X)}$ est alors un endomorphisme de
$\mathcal{E}'$, et $\det(u\otimes1)$ une section de $\mathcal{R}(X)$ au-dessus
de $X$, que l'on appelle encore \emph{déterminant de $u$} et que l'on note
encore $\det u$. Il résulte de (\hyperref[II.6.4.3-fr]{6.4.3}) que $\det u$ est
une section de la \emph{fermeture intégrale} de $\mathcal{O}_X$ dans
$\mathcal{R}(X)$ (\hyperref[II.6.3.2-fr]{6.3.2})~; si en outre $X$ est
\emph{normal}, $\det u$ est donc une \emph{section de $\mathcal{O}_X$}
au-dessus de $X$, et si on suppose de plus que $\mathcal{E}$ est \emph{sans
torsion}, pour que $u$ soit un automorphisme de $\mathcal{E}$, il faut et il
suffit que $\det u$ soit \emph{inversible} en vertu de
(\hyperref[II.6.4.6-fr]{6.4.6}). Les formules
(\hyperref[II.6.4.8.1-fr]{6.4.8.1}) à
(\hyperref[II.6.4.8.3-fr]{6.4.8.3}) sont encore valables~; de l'homomorphisme
$u\mapsto\det u$, appliqué aux modules de sections de
$\mathcal{H}om_{\mathcal{O}_X}(\mathcal{E},\mathcal{E})$, on déduit un
homomorphisme de faisceaux
$\det:\mathcal{H}om_{\mathcal{O}_X}(\mathcal{E},\mathcal{E})
\to\mathcal{R}(X)$, qui prend ses valeurs dans $\mathcal{O}_X$ quand $X$ est
normal. On a des définitions et résultats analogues pour les autres fonctions
symétriques élémentaires $\sigma_j(u)$, lorsque $\mathcal{E}'$ est de
\emph{rang constant}~; si de plus $X$ est \emph{normal}, les $\sigma_j(u)$
sont des \emph{sections de $\mathcal{O}_X$} au-dessus de $X$.

Enfin, soient $X$, $Y$ deux préschémas intègres, et soit $f:X\to Y$ un
morphisme \emph{dominant}. On sait qu'il existe alors un homomorphisme
canonique $f^*(\mathcal{R}(Y))\to\mathcal{R}(X)$
(\hyperref[I.7.3.8-fr]{I, 7.3.8}\footnote{Voir \emph{Errata et Addenda},
liste 1.}), d'où résulte, pour tout
$\mathcal{O}_Y$-Module quasi-cohérent de type fini $\mathcal{F}$, un
homomorphisme canonique
$\theta:f^*(\mathcal{F}\otimes_{\mathcal{O}_Y}\mathcal{R}(Y))
\to f^*(\mathcal{F})\otimes_{\mathcal{O}_X}\mathcal{R}(X)$. Si $v$ est un
endo-

\oldpage[II]{125}

morphisme de $\mathcal{F}$,
$f^*(v\otimes1_{\mathcal{R}(Y)})$ est un endomorphisme de
$f^*(\mathcal{F}\otimes_{\mathcal{O}_Y}\mathcal{R}(Y))$, et on a un diagramme
commutatif
\[
\xymatrix{
f^*(\mathcal{F}\otimes_{\mathcal{O}_Y}\mathcal{R}(Y))
  \ar[r]^{f^*(v\otimes1)}\ar[d]_{\theta}
& f^*(\mathcal{F}\otimes_{\mathcal{O}_Y}\mathcal{R}(Y))\ar[d]^{\theta}\\
f^*(\mathcal{F})\otimes_{\mathcal{O}_X}\mathcal{R}(X)
  \ar[r]_{f^*(v)\otimes1}
& f^*(\mathcal{F})\otimes_{\mathcal{O}_X}\mathcal{R}(X).
}
\]
On en conclut aisément que $\det f^*(v)$ est l'image canonique, par
l'homomorphisme $f^*(\mathcal{R}(Y))\to\mathcal{R}(X)$, de la section
$\det v$ de $\mathcal{R}(Y)$~; on est en effet aussitôt ramené au cas où
$X=\operatorname{Spec}(A)$ et $Y=\operatorname{Spec}(B)$ sont affines, $A$ et
$B$ étant des anneaux intègres de corps de fractions respectifs $K$ et $L$,
l'homomorphisme $B\to A$ étant injectif et s'étendant donc en un monomorphisme
$L\to K$~; si $\mathcal{F}=\widetilde{M}$, où $M$ est un $B$-module de type
fini, le rang sur $L$ de $M\otimes_B L$ est \emph{égal} à celui de
$(M\otimes_B A)\otimes_A K$ sur $K$, et $\det((u\otimes1)\otimes1)$ est
l'image dans $K$ de $\det(u\otimes1)$ pour tout endomorphisme $u$ de $M$,
d'où la conclusion.
\end{env}

\begin{env}[6.4.10]
\label{II.6.4.10-fr}
Supposons enfin que $X$ soit un \emph{préschéma localement noethérien réduit},
de sorte que le faisceau $\mathcal{R}(X)$ des fonctions rationnelles sur $X$
est encore un $\mathcal{O}_X$-Module quasi-cohérent
(\hyperref[I.7.3.4-fr]{I, 7.3.4})~; soit d'autre part $\mathcal{E}$ un
$\mathcal{O}_X$-Module \emph{cohérent} tel que
$\mathcal{E}'=\mathcal{E}\otimes_{\mathcal{O}_X}\mathcal{R}(X)$ soit
\emph{localement libre} (de rang fini). En vertu de
(\hyperref[II.6.4.7-fr]{6.4.7}), si $\mathcal{E}'$ est de rang constant, on
peut, pour tout endomorphisme $u$ de $\mathcal{E}$, définir les $\sigma_j(u)$,
qui sont des sections de $\mathcal{R}(X)$ au-dessus de $X$. Lorsqu'on ne
suppose plus $\mathcal{E}'$ de rang constant, on peut encore définir
l'homomorphisme
$\det:\mathcal{H}om_{\mathcal{O}_X}(\mathcal{E},\mathcal{E})
\to\mathcal{R}(X)$.
\end{env}

\subsection{Norme d'un faisceau inversible.}
\label{subsection:II.6.5-fr}

\begin{env}[6.5.1]
\label{II.6.5.1-fr}
Soit $(X,\mathcal{A})$ un espace annelé et soit $\mathcal{B}$ une
$\mathcal{A}$-Algèbre (commutative). Le $\mathcal{A}$-Module $\mathcal{B}$
s'identifie canoniquement à un sous-$\mathcal{A}$-Module de
$\mathcal{H}om_{\mathcal{A}}(\mathcal{B},\mathcal{B})$, une section $f$ de
$\mathcal{B}$ au-dessus d'un ouvert $U$ de $X$ s'identifiant à la
multiplication par cette section. Si $(X,\mathcal{A})$ et $\mathcal{B}$
vérifient l'une des conditions énumérées dans
(\hyperref[II.6.4.8-fr]{6.4.8}), (\hyperref[II.6.4.9-fr]{6.4.9}) ou
(\hyperref[II.6.4.10-fr]{6.4.10}), on peut donc définir $\det(f)$ (et dans
certains cas les $\sigma_j(f)$) qui sont des sections de $\mathcal{A}$ ou de
$\mathcal{R}(X)$ au-dessus de $U$, que nous appellerons la \emph{norme de
$f$} (resp. les \emph{fonctions symétriques élémentaires de $f$}) et noterons
$N_{\mathcal{B}/\mathcal{A}}(f)$. Nous allons supposer vérifiée \emph{l'une}
des conditions suivantes~:
\begin{enumerate}
\item[(I)] \emph{$\mathcal{B}$ est un $\mathcal{A}$-Module localement libre
(de rang fini).}
\item[(II)] \emph{$(X,\mathcal{A})$ est un préschéma localement noethérien
réduit, $\mathcal{B}$ est un $\mathcal{A}$-Module cohérent tel que
$\mathcal{B}\otimes_{\mathcal{A}}\mathcal{R}(X)$ soit un
$\mathcal{R}(X)$-Module localement libre, et pour toute section
$f\in\Gamma(U,\mathcal{B})$ au-dessus d'un ouvert $U\subset X$,
$N_{\mathcal{B}/\mathcal{A}}(f)$ est une section de $\mathcal{A}$ au-dessus
de $U$.}
\end{enumerate}

\oldpage[II]{126}

On notera que cette dernière condition est automatiquement vérifiée lorsque
le préschéma localement noethérien $X$ est \emph{normal}
(\hyperref[II.6.4.9-fr]{6.4.9}).

D'autre part, l'hypothèse que
$\mathcal{B}\otimes_{\mathcal{A}}\mathcal{R}(X)$ soit localement libre peut
encore s'exprimer de la façon suivante~: désignons par $X_\alpha$ les
sous-préschémas fermés réduits de $X$ ayant pour espaces sous-jacents les
composantes irréductibles de $X$ (\hyperref[I.5.2.1-fr]{I, 5.2.1}), qui sont
donc des préschémas intègres localement noethériens. Tout $x\in X$ appartient
à un nombre fini des sous-espaces $X_\alpha$~; d'autre part,
$\mathcal{B}\otimes_{\mathcal{A}}\mathcal{R}(X_\alpha)$ est un
$\mathcal{R}(X_\alpha)$-Module localement libre de rang constant $k_\alpha$
(\hyperref[I.7.3.6-fr]{I, 7.3.6})~; dire que
$\mathcal{B}\otimes_{\mathcal{A}}\mathcal{R}(X)$ est un
$\mathcal{R}(X)$-Module localement libre signifie que, pour tout $x\in X$,
\emph{les rangs $k_\alpha$ correspondant aux indices tels que $x\in X_\alpha$
sont égaux}. La question est en effet locale, et on est ramené au cas
$X=\operatorname{Spec}(C)$, où $C$ est un anneau noethérien réduit, et
$\mathcal{B}=\widetilde{D}$, où $D$ est une $C$-algèbre qui est un $C$-module
de type fini~; si $\mathfrak{p}_i$ ($1\leq i\leq m$) sont les idéaux premiers
minimaux de $C$, l'anneau total des fractions $L$ de $C$ est composé direct
des corps des fractions $K_i$ des anneaux intègres $C/\mathfrak{p}_i$, et
$D\otimes_C L$ est somme directe des $D\otimes_C K_i$, d'où la conclusion.

Il est clair que sous les hypothèses (I) ou (II), on a défini ainsi un
\emph{homomorphisme de faisceaux de monoïdes multiplicatifs}
$N_{\mathcal{B}/\mathcal{A}}:\mathcal{B}\to\mathcal{A}$, aussi noté $N$ si
aucune confusion n'en résulte, et appelé l'homomorphisme \emph{norme}. Pour
deux sections $f$, $g$ de $\mathcal{B}$ au-dessus d'un même ouvert $U$, on a
donc
\[
\label{II.6.5.1.1-fr}
N_{\mathcal{B}/\mathcal{A}}(fg)
=N_{\mathcal{B}/\mathcal{A}}(f)N_{\mathcal{B}/\mathcal{A}}(g)
\tag{6.5.1.1}
\]
pour les sections correspondantes de $\mathcal{A}$ au-dessus de $U$~;
\[
\label{II.6.5.1.2-fr}
N_{\mathcal{B}/\mathcal{A}}(1_{\mathcal{B}})=1_{\mathcal{A}}~;
\tag{6.5.1.2}
\]
enfin, pour toute section $s$ de $\mathcal{A}$ au-dessus de $U$
\[
\label{II.6.5.1.3-fr}
N_{\mathcal{B}/\mathcal{A}}(s\cdot1_{\mathcal{B}})=s^n
\tag{6.5.1.3}
\]
si le rang de $\mathcal{B}$ est constant et égal à $n$ (pour
$s=0_{\mathcal{A}}$, cette formule donne
$N(0_{\mathcal{B}})=0_{\mathcal{A}}$ si $n\geq1$,
$N(0_{\mathcal{B}})=N(1_{\mathcal{B}})=1_{\mathcal{A}}$ si $n=0$).

Dans l'hypothèse (I), pour que $f\in\Gamma(U,\mathcal{B})$ soit inversible,
il faut et il suffit que $N(f)\in\Gamma(U,\mathcal{A})$ le soit. Dans
l'hypothèse (II), cette condition est nécessaire~; elle est suffisante (par
(\hyperref[II.6.4.7-fr]{6.4.7})) lorsqu'on suppose que
$\mathcal{B}\to\mathcal{B}\otimes_{\mathcal{A}}\mathcal{R}(X)$ est
\emph{injectif} et que l'hypothèse plus restrictive suivante est vérifiée~:
\begin{enumerate}
\item[(II \emph{bis})] \emph{$(X,\mathcal{A})$ est un préschéma localement
noethérien réduit, $\mathcal{B}$ est un $\mathcal{A}$-Module cohérent tel que
$\mathcal{B}\otimes_{\mathcal{A}}\mathcal{R}(X)$ soit un
$\mathcal{R}(X)$-Module localement libre, et pour toute section
$f\in\Gamma(U,\mathcal{B})$ au-dessus d'un ouvert tel que
$\mathcal{B}\otimes_{\mathcal{A}}\mathcal{R}(X)|U$ soit de rang constant $n$
sur $\mathcal{R}(X)|U$, les $\sigma_j(f)$ ($1\leq j\leq n$) sont des sections
de $\mathcal{A}$ au-dessus de $U$.}
\end{enumerate}
(On notera encore que cette condition est vérifiée lorsque $X$ est
\emph{normal}.)
\end{env}

\begin{env}[6.5.2]
\label{II.6.5.2-fr}
Supposons vérifiée une des hypothèses (I), (II) de
(\hyperref[II.6.5.1-fr]{6.5.1}), et soit $\mathcal{L}'$ un
$\mathcal{B}$-Module \emph{inversible}. Nous allons lui associer canoniquement
(à un isomorphisme unique près) un $\mathcal{A}$-Module \emph{inversible}
$\mathcal{L}$ de la façon suivante. Désignons par $\mathcal{A}^*$ (resp.
$\mathcal{B}^*$) le sous-faisceau de $\mathcal{A}$ (resp. $\mathcal{B}$) tel
que $\Gamma(U,\mathcal{A}^*)$ (resp. $\Gamma(U,\mathcal{B}^*)$) soit
l'ensemble des

\oldpage[II]{127}

éléments inversibles de $\Gamma(U,\mathcal{A})$ (resp.
$\Gamma(U,\mathcal{B})$) pour tout ouvert $U\subset X$~; ce sont des faisceaux
de groupes multiplicatifs, et $N_{\mathcal{B}/\mathcal{A}}$, restreint à
$\mathcal{B}^*$, est un \emph{homomorphisme}
$\mathcal{B}^*\to\mathcal{A}^*$ de faisceaux de groupes
(\hyperref[II.6.5.1-fr]{6.5.1}). Soit $\mathfrak{L}$ l'ensemble des couples
$(U_\lambda,\eta_\lambda)$, ayant la propriété suivante~: $U_\lambda$ est un
ouvert de $X$, et $\eta_\lambda$ est un isomorphisme
$\mathcal{L}'|U_\lambda\simeq\mathcal{B}|U_\lambda$ de
$(\mathcal{B}|U_\lambda)$-Modules. Par hypothèse, les $U_\lambda$ forment un
recouvrement de $X$~; pour deux indices quelconques $\lambda$, $\mu$, on pose
$\omega_{\lambda\mu}=(\eta_\lambda|U_\lambda\cap U_\mu)
\circ(\eta_\mu|U_\lambda\cap U_\mu)^{-1}$, automorphisme de
$\mathcal{B}|U_\lambda\cap U_\mu$, qui s'identifie canoniquement à une section
de $\mathcal{B}^*$ au-dessus de $U_\lambda\cap U_\mu$, et
$(\omega_{\lambda\mu})$ est un $1$-cocycle du recouvrement
$\mathfrak{U}=(U_\lambda)$ à valeurs dans $\mathcal{B}^*$
(\hyperref[0.5.4.7-fr]{0, 5.4.7}). Le fait que
$N_{\mathcal{B}/\mathcal{A}}:\mathcal{B}^*\to\mathcal{A}^*$ est un
homomorphisme entraîne alors que
$(N_{\mathcal{B}/\mathcal{A}}\omega_{\lambda\mu})$ est un $1$-cocycle de
$\mathfrak{U}$ à valeurs dans $\mathcal{A}^*$, et il lui correspond donc (à un
isomorphisme unique près) un $\mathcal{A}$-Module inversible $\mathcal{L}$
(\hyperref[0.5.4.7-fr]{0, 5.4.7}) que nous désignerons par
$N_{\mathcal{B}/\mathcal{A}}(\mathcal{L}')$ et appellerons \emph{la norme du
$\mathcal{B}$-Module inversible $\mathcal{L}'$}.

Soit $\mathfrak{M}$ une partie de $\mathfrak{L}$ telle que les $U_\lambda$
correspondants forment encore un recouvrement de $X$, et soit $\mathfrak{B}$
ce recouvrement~; la restriction du cocycle $(\omega_{\lambda\mu})$ à
$\mathfrak{B}$ définit un $1$-cocycle
$(N_{\mathcal{B}/\mathcal{A}}\omega_{\lambda\mu})$ de $\mathfrak{B}$ à
valeurs dans $\mathcal{A}^*$, restriction du $1$-cocycle
$(N_{\mathcal{B}/\mathcal{A}}\omega_{\lambda\mu})$ de $\mathfrak{U}$~; il est
clair qu'il y a un isomorphisme canonique du $\mathcal{A}$-Module inversible
défini par ce $1$-cocycle de $\mathfrak{B}$ sur
$N_{\mathcal{B}/\mathcal{A}}(\mathcal{L}')$, ce qui permet de définir ce
$\mathcal{A}$-Module inversible par un sous-recouvrement quelconque de
$\mathfrak{U}$. Cette possibilité montre aussitôt que, si
$\mathcal{L}'_1$, $\mathcal{L}'_2$ sont deux $\mathcal{B}$-Modules inversibles,
on a, en vertu de (\hyperref[II.6.5.1.1-fr]{6.5.1.1}) et
(\hyperref[II.6.5.1.2-fr]{6.5.1.2}),
\[
\label{II.6.5.2.1-fr}
N(\mathcal{L}'_1\otimes_{\mathcal{B}}\mathcal{L}'_2)
=N(\mathcal{L}'_1)\otimes_{\mathcal{A}}N(\mathcal{L}'_2)
\tag{6.5.2.1}
\]
et
\[
\label{II.6.5.2.2-fr}
N_{\mathcal{B}/\mathcal{A}}(\mathcal{B})=\mathcal{A}
\tag{6.5.2.2}
\]
ainsi que
\[
\label{II.6.5.2.3-fr}
N((\mathcal{L}')^{-1})=(N(\mathcal{L}'))^{-1}
\tag{6.5.2.3}
\]
à des isomorphismes canoniques près. En outre il résulte de
(\hyperref[II.6.5.1.3-fr]{6.5.1.3}) que si $\mathcal{L}$ est un
$\mathcal{A}$-Module inversible et si $\mathcal{B}$ est de rang constant $n$
sur $\mathcal{A}$ dans le cas (I) (resp.
$\mathcal{B}\otimes_{\mathcal{A}}\mathcal{R}(X)$ de rang constant $n$ sur
$\mathcal{R}(X)$ dans le cas (II)), on a, à un isomorphisme canonique près
\[
\label{II.6.5.2.4-fr}
N_{\mathcal{B}/\mathcal{A}}(\mathcal{L}\otimes_{\mathcal{A}}\mathcal{B})
=\mathcal{L}^{\otimes n}.
\tag{6.5.2.4}
\]
\end{env}

\begin{env}[6.5.3]
\label{II.6.5.3-fr}
Montrons maintenant que $N_{\mathcal{B}/\mathcal{A}}(\mathcal{L}')$ est un
\emph{foncteur covariant} dans la catégorie des $\mathcal{B}$-Modules
inversibles. Soit en effet $h':\mathcal{L}'_1\to\mathcal{L}'_2$ un
homomorphisme de $\mathcal{B}$-Modules inversibles, et soit
$\mathfrak{B}=(U_\lambda)$ un recouvrement ouvert de $X$ tel que pour tout
$\lambda$, on ait des $(\mathcal{B}|U_\lambda)$-isomorphismes
$\eta^{(1)}_\lambda:\mathcal{L}'_1|U_\lambda\simeq\mathcal{B}|U_\lambda$ et
$\eta^{(2)}_\lambda:\mathcal{L}'_2|U_\lambda\simeq\mathcal{B}|U_\lambda$~;
il y a donc pour chaque $\lambda$ un endomorphisme $h'_\lambda$ de
$\mathcal{B}|U_\lambda$ tel que
$h'_\lambda\circ\eta^{(1)}_\lambda
=\eta^{(2)}_\lambda\circ(h'|U_\lambda)$, et on peut évidemment identifier
$h'_\lambda$ à une section de $\mathcal{B}$ au-dessus de $U_\lambda$
(\hyperref[0.5.1.1-fr]{0, 5.1.1}). Donc, pour tout couple $(\lambda,\mu)$
d'indices, les restrictions à $U_\lambda\cap U_\mu$ de
$(\eta^{(2)}_\lambda)^{-1}\circ h'_\lambda\circ\eta^{(1)}_\lambda$ et de
$(\eta^{(2)}_\mu)^{-1}\circ h'_\mu\circ\eta^{(1)}_\mu$ coïncident. On en
déduit pour les $1$-cocycles $(\omega^{(1)}_{\lambda\mu})$ et
$(\omega^{(2)}_{\lambda\mu})$ à valeurs dans $\mathcal{B}^*$ correspondant à
$\mathcal{L}'_1$ et $\mathcal{L}'_2$, la relation
\[
\omega^{(2)}_{\lambda\mu}h'_\mu
=h'_\lambda\omega^{(1)}_{\lambda\mu}.
\]

\oldpage[II]{128}

Si on pose $h_\lambda=N(h'_\lambda)$, on aura donc les relations analogues
\[
N(\omega^{(2)}_{\lambda\mu})h_\mu
=h_\lambda N(\omega^{(1)}_{\lambda\mu})
\]
et par suite les $h_\lambda$ définissent un homomorphisme
$N(\mathcal{L}'_1)\to N(\mathcal{L}'_2)$ que nous désignerons par
$N_{\mathcal{B}/\mathcal{A}}(h')$ ou $N(h')$. Dans l'hypothèse (I), pour que
$h'$ soit un \emph{isomorphisme}, il faut et il suffit (puisqu'il s'agit d'une
question locale) que $N_{\mathcal{B}/\mathcal{A}}(h')$ le soit. Dans
l'hypothèse (II), cette condition est encore nécessaire~; elle est suffisante
si l'hypothèse (II \emph{bis}) est vérifiée et si
$\mathcal{B}\to\mathcal{B}\otimes_{\mathcal{A}}\mathcal{R}(X)$ est injectif.

Prenons en particulier $\mathcal{L}'_1=\mathcal{B}$~; les homomorphismes
$\mathcal{B}\to\mathcal{L}'$ s'identifient alors
(\hyperref[0.5.1.1-fr]{0, 5.1.1}) aux sections de $\mathcal{L}'$ au-dessus de
$X$, d'où une application canonique
\[
N_{\mathcal{B}/\mathcal{A}}:\Gamma(X,\mathcal{L}')
\to\Gamma(X,N_{\mathcal{B}/\mathcal{A}}(\mathcal{L}')).
\]
Il résulte encore de (\hyperref[II.6.5.1.1-fr]{6.5.1.1}) que si
$f'_1\in\Gamma(X,\mathcal{L}'_1)$, $f'_2\in\Gamma(X,\mathcal{L}'_2)$, on a
\[
\label{II.6.5.3.1-fr}
N(f'_1\otimes f'_2)=N(f'_1)\otimes N(f'_2).
\tag{6.5.3.1}
\]
Pour tout $\mathcal{A}$-Module inversible $\mathcal{L}$ et toute section
$f\in\Gamma(X,\mathcal{L})$, on a, compte tenu de
(\hyperref[II.6.5.2.4-fr]{6.5.2.4})
\[
\label{II.6.5.3.2-fr}
N_{\mathcal{B}/\mathcal{A}}(f\otimes1_{\mathcal{B}})=f^{\otimes n}
\tag{6.5.3.2}
\]
lorsque $\mathcal{B}$ est de rang constant $n$ dans l'hypothèse (I) (resp.
lorsque $\mathcal{B}\otimes_{\mathcal{A}}\mathcal{R}(X)$ est de rang constant
$n$ dans l'hypothèse (II)). Enfin, pour que l'homomorphisme
$\mathcal{B}\to\mathcal{L}'$ correspondant à une section $f'$ de
$\mathcal{L}'$ au-dessus de $X$ soit un isomorphisme, il faut et il suffit que
$f'_x$ soit une base de $\mathcal{L}'_x$ pour tout $x\in X$~; dans l'hypothèse
(I), cette condition équivaut donc à dire que $(N(f'))_x$ est une base de
$(N(\mathcal{L}'))_x$ pour tout $x$~; dans l'hypothèse (II), cette condition
est encore nécessaire, et elle est suffisante lorsque $\mathcal{B}$ vérifie
(II \emph{bis}) et que
$\mathcal{B}\to\mathcal{B}\otimes_{\mathcal{A}}\mathcal{R}(X)$ est injectif.
\end{env}

\begin{env}[6.5.4]
\label{II.6.5.4-fr}
Soient $(X,\mathcal{A})$, $(X',\mathcal{A}')$ deux espaces annelés,
$f:X'\to X$ un morphisme, $\mathcal{B}$ une $\mathcal{A}$-Algèbre,
$\mathcal{B}'=f^*(\mathcal{B})$ la $\mathcal{A}'$-Algèbre image réciproque. On
suppose vérifiée l'une des hypothèses suivantes~:
\begin{enumerate}
\item[$1^\circ$] $\mathcal{B}$ vérifie l'hypothèse (I) de
(\hyperref[II.6.5.1-fr]{6.5.1}).
\item[$2^\circ$] $(X,\mathcal{A})$ et $\mathcal{B}$ vérifient l'hypothèse
(II) de (\hyperref[II.6.5.1-fr]{6.5.1}), $(X',\mathcal{A}')$ est un préschéma
localement noethérien réduit, et si l'on désigne par $X_\alpha$ et
$X'_\beta$ les sous-préschémas fermés réduits respectifs de $X$ et $X'$ ayant
pour espaces sous-jacents les composantes irréductibles de ces espaces, la
restriction de $f$ à chaque $X'_\beta$ est un morphisme \emph{dominant} de
$X'_\beta$ dans un $X_\alpha$.
\end{enumerate}
Sous ces conditions, $\mathcal{B}'$ vérifie respectivement l'hypothèse (I) ou
l'hypothèse (II) de (\hyperref[II.6.5.1-fr]{6.5.1})~; la première assertion
est immédiate. Pour établir la seconde, il suffit de voir que pour tout
$x'\in X'$, les rangs des
$\mathcal{B}'\otimes_{\mathcal{O}_{X'}}\mathcal{R}(X'_\beta)$ pour les $\beta$
tels que $x'\in X'_\beta$ sont \emph{les mêmes}. Or, si la restriction de $f$
à $X'_\beta$ est un morphisme dominant dans $X_\alpha$, le rang de
$\mathcal{B}'\otimes_{\mathcal{O}_{X'}}\mathcal{R}(X'_\beta)$ est égal à celui
de $\mathcal{B}\otimes_{\mathcal{O}_X}\mathcal{R}(X_\alpha)$ (on le voit
aussitôt en se ramenant au cas affine comme dans
(\hyperref[II.6.4.9-fr]{6.4.9})), d'où notre assertion en vertu de l'hypothèse
(II) et de (\hyperref[II.6.5.1-fr]{6.5.1}).

\oldpage[II]{129}

Cela étant, il résulte de (\hyperref[II.6.4.8-fr]{6.4.8}),
(\hyperref[II.6.4.9-fr]{6.4.9}) et (\hyperref[II.6.4.10-fr]{6.4.10}) que si
$s$ est une section de $\mathcal{B}$ au-dessus d'un ouvert $U\subset X$,
$s'$ la section correspondante de $\mathcal{B}'$ au-dessus de $f^{-1}(U)$,
$N_{\mathcal{B}'/\mathcal{A}'}(s')$ est la section de $\mathcal{A}'$ au-dessus
de $f^{-1}(U)$ qui correspond à la section
$N_{\mathcal{B}/\mathcal{A}}(s)$ de $\mathcal{A}$ au-dessus de $U$.

Si $\mathcal{M}$ est un $\mathcal{B}$-Module inversible, on déduit de ce qui
précède que si $\mathcal{M}'=f^*(\mathcal{M})$ (qui est un
$\mathcal{B}'$-Module inversible), on a
$N_{\mathcal{B}'/\mathcal{A}'}(\mathcal{M}')
=f^*(N_{\mathcal{B}/\mathcal{A}}(\mathcal{M}))$ à un isomorphisme canonique
près.
\end{env}

\begin{env}[6.5.5]
\label{II.6.5.5-fr}
Supposons désormais que $(X,\mathcal{A})$ soit un \emph{préschéma}. La donnée
d'une $\mathcal{A}$-Algèbre quasi-cohérente $\mathcal{B}$, qui est un
$\mathcal{A}$-Module \emph{de type fini} équivaut alors, comme on sait, à
celle d'un morphisme \emph{fini} $g:X'\to X$ tel que
$g_*(\mathcal{O}_{X'})=\mathcal{B}$, défini à un $X$-isomorphisme près
(\hyperref[II.6.1.2-fr]{6.1.2} et \hyperref[II.1.3.1-fr]{1.3.1}). En outre,
se donner un $\mathcal{O}_{X'}$-Module quasi-cohérent $\mathcal{F}'$ équivaut
à se donner un $\mathcal{B}$-Module quasi-cohérent $\mathcal{F}$ tel que
$g_*(\mathcal{F}')=\mathcal{F}$ (\hyperref[II.1.4.3-fr]{1.4.3}), et pour que
$\mathcal{F}'$ soit inversible, il faut et il suffit que $\mathcal{F}$ le soit
(\hyperref[II.6.1.12-fr]{6.1.12}). Pour traduire les résultats précédents en
termes du morphisme fini $g$, il faudra donc supposer, soit que
$g_*(\mathcal{O}_{X'})$ est un $\mathcal{O}_X$-Module \emph{localement libre}
(de type fini), soit que $(X,\mathcal{O}_X)$ et $g_*(\mathcal{O}_{X'})$
vérifient l'hypothèse (II). Pour tout $\mathcal{O}_{X'}$-Module inversible
$\mathcal{L}'$, on posera alors
\[
\label{II.6.5.5.1-fr}
N_{X'/X}(\mathcal{L}')
=N_{g_*(\mathcal{O}_{X'})/\mathcal{O}_X}(g_*(\mathcal{L}'))
\tag{6.5.5.1}
\]
et on dira que c'est la \emph{norme} (relative à $g$) de $\mathcal{L}'$. De
même, si $h':\mathcal{L}'_1\to\mathcal{L}'_2$ est un homomorphisme de
$\mathcal{O}_{X'}$-Modules inversibles, on posera
\[
\label{II.6.5.5.2-fr}
N_{X'/X}(h')
=N_{g_*(\mathcal{O}_{X'})/\mathcal{O}_X}(g_*(h')):
N_{X'/X}(\mathcal{L}'_1)\to N_{X'/X}(\mathcal{L}'_2).
\tag{6.5.5.2}
\]
En particulier, pour $\mathcal{L}'_1=\mathcal{O}_{X'}$, on obtient ainsi une
application canonique
\[
\label{II.6.5.5.3-fr}
N_{X'/X}:\Gamma(X',\mathcal{L}')
\to\Gamma(X,N_{X'/X}(\mathcal{L}')).
\tag{6.5.5.3}
\]
Nous laisserons au lecteur la plupart de ces traductions, et nous bornerons à
expliciter les suivantes~:
\end{env}

\begin{proposition}[6.5.6]
\label{II.6.5.6-fr}
Soit $g:X'\to X$ un morphisme fini, et supposons, soit que
$g_*(\mathcal{O}_{X'})$ est un $\mathcal{O}_X$-Module localement libre, soit
que $(X,\mathcal{O}_X)$ et $g_*(\mathcal{O}_{X'})$ vérifient (II \emph{bis})
(ce qui sera le cas en particulier lorsque $X$ est localement noethérien et
normal). Pour qu'un homomorphisme $h':\mathcal{L}'_1\to\mathcal{L}'_2$ de
$\mathcal{O}_{X'}$-Modules inversibles soit un isomorphisme, il faut et il
suffit, dans la première hypothèse, que $N_{X'/X}(h')$ soit un isomorphisme~;
dans la seconde hypothèse, cette condition est encore nécessaire, et elle est
suffisante lorsque l'homomorphisme
$g_*(\mathcal{O}_{X'})\to
g_*(\mathcal{O}_{X'})\otimes_{\mathcal{O}_X}\mathcal{R}(X)$ est injectif.
\end{proposition}

On notera qu'on utilise ici le fait que, pour que $g_*(h')$ soit un
isomorphisme, il faut et il suffit que $h'$ en soit un
(\hyperref[II.1.4.2-fr]{1.4.2}).

\begin{corollary}[6.5.7]
\label{II.6.5.7-fr}
Soit $g:X'\to X$ un morphisme fini, et supposons, soit que
$g_*(\mathcal{O}_{X'})$ est un $\mathcal{O}_X$-Module localement libre, soit
que $(X,\mathcal{O}_X)$ et $g_*(\mathcal{O}_{X'})$ vérifient (II \emph{bis})
et que
$g_*(\mathcal{O}_{X'})\to
g_*(\mathcal{O}_{X'})\otimes_{\mathcal{O}_X}\mathcal{R}(X)$ est injectif.
Soient $\mathcal{L}'$ un $\mathcal{O}_{X'}$-Module inversible, $f'$ une
section de $\mathcal{L}'$ au-dessus de $X'$, $f=N_{X'/X}(f')$ la section de
$\mathcal{L}=N_{X'/X}(\mathcal{L}')$ au-dessus de $X$ qui lui correspond
(\hyperref[II.6.5.5.3-fr]{6.5.5.3}).

\oldpage[II]{130}

Alors on a $g(X'-X'_{f'})=X-X_f$ et $X_f$ est le plus grand ouvert $U$ de
$X$ tel que $g^{-1}(U)\subset X'_{f'}$.
\end{corollary}

En effet, $g(X'-X'_{f'})$ est fermé dans $X$
(\hyperref[II.6.1.10-fr]{6.1.10})~; il suffit donc de prouver la dernière
assertion. Or la relation $U\subset X_f$ équivaut au fait que l'homomorphisme
$\mathcal{O}_X|U\to\mathcal{L}|U$ défini par $f|U$ est un isomorphisme. En
vertu de (\hyperref[II.6.5.6-fr]{6.5.6}), cela équivaut à dire que
l'homomorphisme
$\mathcal{O}_{X'}|g^{-1}(U)\to\mathcal{L}'|g^{-1}(U)$ défini par
$f'|g^{-1}(U)$ est un isomorphisme, c'est-à-dire à la relation
$g^{-1}(U)\subset X'_{f'}$.

\begin{proposition}[6.5.8]
\label{II.6.5.8-fr}
Soient $g:X'\to X$ un morphisme fini, $f:Y\to X$ un morphisme~; soient
$Y'=X'_{(Y)}$, $g'=g_{(Y)}$, $f'=f_{(X')}$ de sorte que l'on a le diagramme
commutatif
\[
\xymatrix{
X'\ar[d]_g & Y'\ar[l]_{f'}\ar[d]^{g'}\\
X & Y\ar[l]^f
}
\]
On suppose, soit que $g_*(\mathcal{O}_{X'})$ est localement libre, soit que
$(X,\mathcal{O}_X)$ et $g_*(\mathcal{O}_{X'})$ vérifient (II), que $Y$ est un
préschéma localement noethérien réduit, et que la restriction de $f$ à toute
composante irréductible de $Y$ est un morphisme dominant dans une composante
irréductible de $X$. Alors, pour tout $\mathcal{O}_{X'}$-Module inversible
$\mathcal{L}'$, on a
\[
N_{Y'/Y}({f'}^*(\mathcal{L}'))=f^*(N_{X'/X}(\mathcal{L}'))
\]
à un isomorphisme canonique près.
\end{proposition}

Notons que l'on a $f^*(g_*(\mathcal{L}'))=g'_*({f'}^*(\mathcal{L}'))$ en
vertu de (\hyperref[II.1.5.2-fr]{1.5.2}), et en particulier
$g'_*(\mathcal{O}_{Y'})=f^*(g_*(\mathcal{O}_{X'}))$~; si
$g_*(\mathcal{O}_{X'})$ est localement libre, il en est donc de même de
$g'_*(\mathcal{O}_{Y'})$. La conclusion résulte alors des définitions et de
(\hyperref[II.6.5.4-fr]{6.5.4}).

\begin{remark}[6.5.9]
\label{II.6.5.9-fr}
Nous généraliserons ultérieurement la notion de norme développée ci-dessus,
et la mettrons en relation avec la notion d'image directe d'un diviseur.
\end{remark}

\subsection{Application : critères d'amplitude.}
\label{subsection:II.6.6-fr}

\begin{proposition}[6.6.1]
\label{II.6.6.1-fr}
Soient $Y$ un préschéma, $f:X\to Y$ un morphisme quasi-compact,
$g:X'\to X$ un morphisme fini et surjectif. On suppose, soit que
$g_*(\mathcal{O}_{X'})$ est un $\mathcal{O}_X$-Module localement libre, soit
que $(X,\mathcal{O}_X)$ et $g_*(\mathcal{O}_{X'})$ vérifient (II
\emph{bis}). Alors, pour tout $\mathcal{O}_{X'}$-Module inversible
$\mathcal{L}'$, ample pour $f\circ g$, $N_{X'/X}(\mathcal{L}')=\mathcal{L}$
est ample pour $f$.
\end{proposition}

On peut supposer $Y$ affine (\hyperref[II.4.6.4-fr]{4.6.4}), et alors, en
vertu de (\hyperref[II.4.6.6-fr]{4.6.6}), l'énoncé est équivalent au

\begin{corollary}[6.6.2]
\label{II.6.6.2-fr}
Soient $X$ un préschéma quasi-compact, $g:X'\to X$ un morphisme fini
surjectif, tel que $g_*(\mathcal{O}_{X'})$ soit un $\mathcal{O}_X$-Module
localement libre, ou que $(X,\mathcal{O}_X)$ et $g_*(\mathcal{O}_{X'})$
vérifient (II \emph{bis}). Alors, pour tout $\mathcal{O}_{X'}$-Module ample
$\mathcal{L}'$, $\mathcal{L}=N_{X'/X}(\mathcal{L}')$ est ample.
\end{corollary}

Dans la seconde hypothèse, on peut supposer en outre que l'homomorphisme
canonique
$g_*(\mathcal{O}_{X'})\to
g_*(\mathcal{O}_{X'})\otimes_{\mathcal{O}_X}\mathcal{R}(X)$ est
\emph{injectif}. En effet, dans le cas contraire, soit $\mathcal{T}$

\oldpage[II]{131}

le noyau de cet homomorphisme, qui est un Idéal cohérent de
$g_*(\mathcal{O}_{X'})=\mathcal{B}$
(\hyperref[I.6.1.1-fr]{I, 6.1.1}), et posons
$X''=\operatorname{Spec}(\mathcal{B}/\mathcal{T})$~; on a donc un diagramme
commutatif
\[
\xymatrix{
X''\ar[rr]^h\ar[dr]_{g'} && X'\ar[dl]^g\\
& X
}
\]
où $h$ est une immersion fermée (\hyperref[II.1.4.10-fr]{1.4.10}). En outre,
on sait que le support de $\mathcal{T}$ est un ensemble fermé
(\hyperref[0.5.2.2-fr]{0, 5.2.2}) rare dans $X$
(\hyperref[I.7.4.6-fr]{I, 7.4.6}), d'où on conclut que pour le point générique
$x$ d'une composante irréductible de $X$, il y a un voisinage ouvert affine
$U$ de $x$ tel que
$\mathcal{B}|U=(\mathcal{B}/\mathcal{T})|U$. Comme $g$ est par hypothèse
surjectif, on en conclut que $x\in g'(X'')$~; $g'$ est donc dominant, et étant
un morphisme fini, il est \emph{surjectif}
(\hyperref[II.6.1.10-fr]{6.1.10})~; on a par définition
\[
g'_*(\mathcal{O}_{X''})\otimes_{\mathcal{O}_X}\mathcal{R}(X)
=(\mathcal{B}/\mathcal{T})\otimes_{\mathcal{O}_X}\mathcal{R}(X)
=g_*(\mathcal{O}_{X'})\otimes_{\mathcal{O}_X}\mathcal{R}(X),
\]
donc $(X,\mathcal{O}_X)$ et $g'_*(\mathcal{O}_{X''})$ vérifient
(II \emph{bis}), et en outre
\[
g'_*(\mathcal{O}_{X''})\longrightarrow
g'_*(\mathcal{O}_{X''})\otimes_{\mathcal{O}_X}\mathcal{R}(X)
\]
est injectif. Enfin, $h^*(\mathcal{L}')=\mathcal{L}''$ est un
$\mathcal{O}_{X''}$-Module ample
(\hyperref[II.4.6.13-fr]{4.6.13, (i \emph{bis})}), et on a
$N_{X''/X}(\mathcal{L}'')=N_{X'/X}(\mathcal{L}')$. En effet, pour définir
ces deux $\mathcal{O}_X$-Modules inversibles on peut utiliser un même
recouvrement ouvert affine $(U_\lambda)$ de $X$ tel que les restrictions de
$g_*(\mathcal{L}')$ et $g'_*(\mathcal{L}'')$ à $U_\lambda$ soient
respectivement isomorphes à $\mathcal{B}|U_\lambda$ et
$(\mathcal{B}/\mathcal{T})|U_\lambda$. On voit aussitôt qu'à tout
isomorphisme
\[
\eta_\lambda:g_*(\mathcal{L}')|U_\lambda\longrightarrow
\mathcal{B}|U_\lambda
\]
correspond canoniquement un isomorphisme
\[
\eta'_\lambda:g'_*(\mathcal{L}'')|U_\lambda\longrightarrow
(\mathcal{B}/\mathcal{T})|U_\lambda,
\]
de sorte que, si $(\omega_{\lambda\mu})$ et
$(\omega'_{\lambda\mu})$ sont les $1$-cocycles correspondant aux systèmes
d'isomorphismes $(\eta_\lambda)$ et $(\eta'_\lambda)$
(\hyperref[II.6.5.2-fr]{6.5.2}), $\omega'_{\lambda\mu}$ soit l'image
canonique dans
$\Gamma(U_\lambda\cap U_\mu,\mathcal{B}/\mathcal{T})$ de
$\omega_{\lambda\mu}\in\Gamma(U_\lambda\cap U_\mu,\mathcal{B})$. En vertu
de la définition de $\mathcal{T}$, on en conclut que
\[
N_{\mathcal{B}/\mathcal{A}}\omega_{\lambda\mu}
=N_{(\mathcal{B}/\mathcal{T})/\mathcal{A}}\omega'_{\lambda\mu}
\]
(avec $\mathcal{A}=\mathcal{O}_X$), d'où l'égalité énoncée.

Supposons donc l'homomorphisme
$g_*(\mathcal{O}_{X'})\to
g_*(\mathcal{O}_{X'})\otimes_{\mathcal{O}_X}\mathcal{R}(X)$ injectif
lorsqu'on est dans l'hypothèse (II \emph{bis}). Il suffit de prouver que
lorsque $f$ parcourt les sections de $\mathcal{L}^{\otimes n}$ ($n>0$)
au-dessus de $X$, les $X_f$ forment une base de la topologie de $X$
(\hyperref[II.4.5.2-fr]{4.5.2}). Or, soit $x\in X$, et soit $U$ un voisinage
quelconque de $x$~; comme $g^{-1}(x)$ est fini
(\hyperref[II.6.1.7-fr]{6.1.7}) et que $\mathcal{L}'$ est ample, il existe un
entier $n>0$ et une section $f'$ de ${\mathcal{L}'}^{\otimes n}$ au-dessus de
$X'$, tels que $X'_{f'}$ soit un voisinage de $g^{-1}(x)$ contenu dans
$g^{-1}(U)$ (\hyperref[II.4.5.4-fr]{4.5.4}). Comme
\[
\mathcal{L}^{\otimes n}=N_{X'/X}({\mathcal{L}'}^{\otimes n}),
\]
il suffit de prendre $f=N_{X'/X}(f')$~; en effet, alors
$X-X_f=g(X'-X'_{f'})$ (\hyperref[II.6.5.7-fr]{6.5.7}), donc
$x\in X_f\subset U$.

\begin{corollary}[6.6.3]
\label{II.6.6.3-fr}
Sous les hypothèses de (\hyperref[II.6.6.1-fr]{6.6.1}), pour qu'un
$\mathcal{O}_X$-Module inversible $\mathcal{L}$ soit ample pour $f$, il faut
et il suffit que $\mathcal{L}'=g^*(\mathcal{L})$ soit ample pour $f\circ g$.
\end{corollary}

La condition est nécessaire, puisque $g$ est affine
(\hyperref[II.5.1.12-fr]{5.1.12}). Pour démontrer que la condition est
suffisante, on peut supposer $Y$ affine (\hyperref[II.4.6.4-fr]{4.6.4}), donc
$X$ et $X'$ quasi-compacts et $\mathcal{L}'$ ample
(\hyperref[II.4.6.6-fr]{4.6.6}), et il faut montrer que $\mathcal{L}$ est
ample. Or, l'ensemble des points

\oldpage[II]{132}

$x\in X$ au voisinage desquels $g_*(\mathcal{O}_{X'})$ (resp.
$g_*(\mathcal{O}_{X'})\otimes_{\mathcal{O}_X}\mathcal{R}(X)$) a un rang
donné $n$ dans la première (resp. la seconde) hypothèse, est à la fois ouvert
et fermé dans $X$, donc $X$ est le préschéma somme d'un nombre fini de ces
ouverts, et on peut par suite supposer qu'il est égal à l'un d'eux
(\hyperref[II.4.6.17-fr]{4.6.17}). Mais on a alors
$N_{X'/X}(\mathcal{L}')=\mathcal{L}^{\otimes n}$, donc
$\mathcal{L}^{\otimes n}$ est ample en vertu de
(\hyperref[II.6.6.2-fr]{6.6.2}), et il en est donc de même de $\mathcal{L}$
(\hyperref[II.4.5.6-fr]{4.5.6}).

\begin{corollary}[6.6.4]
\label{II.6.6.4-fr}
Supposons vérifiées les hypothèses de
(\hyperref[II.6.6.1-fr]{6.6.1}) et supposons en outre que
$f:X\to Y$ soit de type fini. Alors, pour que $f$ soit quasi-projectif, il
faut et il suffit que $f\circ g$ le soit. Si on suppose de plus que $Y$ est
un schéma quasi-compact ou un préschéma dont l'espace sous-jacent est
noethérien, alors, pour que $f$ soit projectif, il faut et il suffit que
$f\circ g$ le soit.
\end{corollary}

L'hypothèse entraîne que $f\circ g$ est de type fini. Compte tenu de la
définition des morphismes quasi-projectifs
(\hyperref[II.5.3.1-fr]{5.3.1}), la première assertion résulte de
(\hyperref[II.6.6.1-fr]{6.6.1}) et
(\hyperref[II.6.6.3-fr]{6.6.3}). Compte tenu de ce résultat et de
(\hyperref[II.5.5.3-fr]{5.5.3, (ii)}), il reste à prouver que lorsque $f$ est
supposé quasi-projectif, pour que $f$ soit propre, il faut et il suffit que
$f\circ g$ le soit. Or $f$ est alors séparé
(\hyperref[II.5.3.1-fr]{5.3.1}) et de type fini~; comme $g$ est surjectif,
notre assertion résulte de (\hyperref[II.5.4.2-fr]{5.4.2, (ii)}) et
(\hyperref[II.5.4.3-fr]{5.4.3, (ii)}).

En particulier~:

\begin{corollary}[6.6.5]
\label{II.6.6.5-fr}
Soient $X$ un préschéma de type fini sur un corps $K$, $K'$ une extension de
degré fini de $K$. Pour que $X$ soit projectif (resp. quasi-projectif) sur
$K$, il faut et il suffit que $X'=X\otimes_K K'$ soit projectif (resp.
quasi-projectif) sur $K'$.
\end{corollary}

La condition est en effet nécessaire
(\hyperref[II.5.3.4-fr]{5.3.4, (iii)} et
\hyperref[II.5.5.5-fr]{5.5.5, (iii)}). Inversement, supposons-la vérifiée,
et soit $g:X'\to X$ la projection canonique. Il est clair que $g$ est un
morphisme fini (\hyperref[II.6.1.5-fr]{6.1.5, (iii)}) et surjectif
(\hyperref[I.3.5.2-fr]{I, 3.5.2, (ii)}). En outre,
$g_*(\mathcal{O}_{X'})$ est un $\mathcal{O}_X$-Module localement libre, étant
isomorphe à $\mathcal{O}_X\otimes_K K'$
(\hyperref[II.1.5.2-fr]{1.5.2}). Il résulte alors de l'hypothèse et de
(\hyperref[II.6.1.11-fr]{6.1.11}) et
(\hyperref[II.5.5.5-fr]{5.5.5, (ii)}) que $X'$ est projectif (resp.
quasi-projectif) \emph{sur $K$}~; on déduit alors de
(\hyperref[II.6.6.4-fr]{6.6.4}) que $X$ est projectif (resp.
quasi-projectif) sur $K$.

Au chapitre V, nous montrerons que l'énoncé
(\hyperref[II.6.6.5-fr]{6.6.5}) reste valable lorsque $K'$ est une extension
\emph{quelconque} de $K$.

La fin de ce numéro est consacrée à la démonstration du critère
(\hyperref[II.6.6.11-fr]{6.6.11}), qui est un raffinement assez technique de
(\hyperref[II.6.6.1-fr]{6.6.1})~; elle peut être omise en première lecture.

\begin{lemma}[6.6.6]
\label{II.6.6.6-fr}
Soient $X$ un préschéma noethérien réduit, $\mathcal{E}$ un
$\mathcal{O}_X$-Module cohérent tel que
$\mathcal{E}\otimes_{\mathcal{O}_X}\mathcal{R}(X)$ soit un
$\mathcal{R}(X)$-Module localement libre de rang constant $n$. Il existe
alors un préschéma noethérien réduit $Z$ et un morphisme fini birationnel
$h:Z\to X$ ayant la propriété suivante~: les morphismes de faisceaux
d'ensembles
$\sigma_i:\mathcal{H}om_{\mathcal{O}_X}(\mathcal{E},\mathcal{E})
\to\mathcal{R}(X)$ ($1\leq i\leq n$)
(cf. \hyperref[II.6.4.10-fr]{6.4.10}) appliquent
$\mathcal{H}om_{\mathcal{O}_X}(\mathcal{E},\mathcal{E})$ dans la
$\mathcal{O}_X$-Algèbre cohérente $h_*(\mathcal{O}_Z)$.
\end{lemma}

Considérons en effet un ouvert affine $U$ de $X$, d'anneau $A(U)=A$~; soit
$E=\Gamma(U,\mathcal{E})$, et soit $C_U$ la sous-algèbre de $R(U)$ engendrée
par les $\sigma_i(u)$ lorsque $u$ parcourt $\operatorname{Hom}_A(E,E)$~; on
a vu (\hyperref[II.6.4.7.1-fr]{6.4.7.1}) que cette $A$-algèbre est de
\emph{rang fini}. En outre, il est clair que la formation des algèbres $C_U$
commute avec les opérations de restriction d'un ouvert affine $U$ à un ouvert
affine $U'\subset U$. On a donc défini ainsi une sous-$\mathcal{O}_X$-Algèbre
\emph{finie} $\mathcal{C}$ de $\mathcal{R}(X)$ telle que
$\Gamma(U,\mathcal{C})=C_U$ pour tout ouvert affine $U$ de $X$. On

\oldpage[II]{133}

prendra $Z=\operatorname{Spec}(\mathcal{C})$, et pour $h$ le morphisme
structural, qui est donc fini (\hyperref[II.6.1.2-fr]{6.1.2})~; comme
$\mathcal{C}$ est réduite, $Z$ est un préschéma noethérien réduit
(\hyperref[II.1.3.8-fr]{1.3.8}). Enfin, l'anneau total des fractions de $C_U$
est $R(U)$ par définition, et comme $C_U$ est contenu dans la fermeture
intégrale de $A(U)$ dans $R(U)$, il y a correspondance biunivoque entre idéaux
premiers minimaux de $A(U)$ et idéaux premiers minimaux de $C_U$
(\cite[t.~I, p.~259]{I-13}), ce qui prouve que $h$ est birationnel et achève
la démonstration.

\begin{corollary}[6.6.7]
\label{II.6.6.7-fr}
Sous les hypothèses de (\hyperref[II.6.6.6-fr]{6.6.6}), soit $W$ un ouvert
de $X$ tel que, pour tout $x\in W$, ou bien $X$ est normal au point $x$, ou
bien $\mathcal{E}_x$ est un $\mathcal{O}_x$-Module libre. Alors on peut
supposer $h$ défini de sorte que la restriction de $h$ à $h^{-1}(W)$ soit un
isomorphisme de $h^{-1}(W)$ sur $W$.
\end{corollary}

En effet, l'une ou l'autre hypothèse entraîne que si $U\subset W$ est un
ouvert affine, on a, avec les notations de
(\hyperref[II.6.6.6-fr]{6.6.6}), $(\sigma_i(u))_x\in A_x$ pour tout
$x\in U$ (\hyperref[II.6.4.3-fr]{6.4.3}), donc $\sigma_i(u)\in A$, et la
conclusion résulte de la définition de $h$ donnée dans
(\hyperref[II.6.6.6-fr]{6.6.6}).

\begin{env}[6.6.8]
\label{II.6.6.8-fr}
Soient $X$ un préschéma noethérien réduit, $g:X'\to X$ un morphisme fini
surjectif, de sorte que $\mathcal{B}=g_*(\mathcal{O}_{X'})$ est une
$\mathcal{O}_X$-Algèbre \emph{cohérente}~; supposons en outre que
$\mathcal{B}\otimes_{\mathcal{O}_X}\mathcal{R}(X)$ soit un
$\mathcal{R}(X)$-Module \emph{localement libre de rang constant $n$}. On peut
alors appliquer le lemme (\hyperref[II.6.6.6-fr]{6.6.6}) en prenant
$\mathcal{E}=\mathcal{B}$, d'où, avec les notations de
(\hyperref[II.6.6.6-fr]{6.6.6}), un homomorphisme de faisceaux de monoïdes
multiplicatifs
$\sigma_n:\mathcal{H}om_{\mathcal{O}_X}(\mathcal{B},\mathcal{B})
\to h_*(\mathcal{O}_Z)$, et en composant cet homomorphisme avec
l'homomorphisme canonique
$\mathcal{B}\to
\mathcal{H}om_{\mathcal{O}_X}(\mathcal{B},\mathcal{B})$
(\hyperref[II.6.5.1-fr]{6.5.1}), on obtient donc un homomorphisme de
faisceaux de monoïdes multiplicatifs~:
\[
\label{II.6.6.8.1-fr}
N':\mathcal{B}=g_*(\mathcal{O}_{X'})\longrightarrow
h_*(\mathcal{O}_Z)=\mathcal{C}.
\tag{6.6.8.1}
\]

Cela étant, pour tout $\mathcal{O}_{X'}$-Module inversible $\mathcal{L}'$,
$g_*(\mathcal{L}')$ est un $\mathcal{B}$-Module inversible
(\hyperref[II.6.1.12-fr]{6.1.12}), et la méthode de
(\hyperref[II.6.5.2-fr]{6.5.2}) permet d'associer fonctoriellement à
$\mathcal{L}'$ un $\mathcal{C}$-Module inversible que nous noterons
$N'(g_*(\mathcal{L}'))$.
\end{env}

\begin{lemma}[6.6.9]
\label{II.6.6.9-fr}
Soient $X$ un préschéma noethérien réduit, $g:X'\to X$ un morphisme fini
surjectif tel que
$g_*(\mathcal{O}_{X'})\otimes_{\mathcal{O}_X}\mathcal{R}(X)$ soit un
$\mathcal{R}(X)$-Module localement libre de rang constant. Il existe alors
un préschéma noethérien réduit $Z$ et un morphisme fini birationnel
$h:Z\to X$ ayant la propriété suivante~: pour tout
$\mathcal{O}_{X'}$-Module ample $\mathcal{L}'$, le $\mathcal{O}_Z$-Module
inversible $\mathcal{M}$ tel que
$h_*(\mathcal{M})=N'(g_*(\mathcal{L}'))$ (notations de
\hyperref[II.6.6.8-fr]{6.6.8}) est ample.
\end{lemma}

Supposons d'abord que l'homomorphisme
$\mathcal{B}\to
\mathcal{B}\otimes_{\mathcal{O}_X}\mathcal{R}(X)$ soit \emph{injectif}.
Définissons $Z$ et $h$ comme dans (\hyperref[II.6.6.6-fr]{6.6.6}) (avec
$\mathcal{E}=g_*(\mathcal{O}_{X'})$). Soit $z\in Z$~; il faut montrer qu'il
existe un entier $m>0$ et une section $t$ de $\mathcal{M}^{\otimes m}$
au-dessus de $Z$ telle que $Z_t$ soit un ouvert affine contenant $z$
(\hyperref[II.4.5.2-fr]{4.5.2}). Soient $x=h(z)$, $U$ un ouvert affine de
$X$ contenant $x$~; $h^{-1}(U)$ est alors un voisinage ouvert affine de $z$,
et il suffira de trouver $t$ tel que $z\in Z_t\subset h^{-1}(U)$, car $Z_t$
sera alors nécessairement affine (\hyperref[II.5.5.8-fr]{5.5.8}). Il existe
par hypothèse un entier $n>0$ et une section $s'$ de
${\mathcal{L}'}^{\otimes n}$ au-dessus de $X'$ telle que l'on ait
\[
\label{II.6.6.9.1-fr}
g^{-1}(x)\subset X'_{s'}\subset g^{-1}(U)
\tag{6.6.9.1}
\]
en vertu de (\hyperref[II.4.5.4-fr]{4.5.4}). Par définition, $s'$ est aussi
une section de $g_*({\mathcal{L}'}^{\otimes n})$ au-dessus de $X$, et il lui
correspond comme dans (\hyperref[II.6.5.2-fr]{6.5.2}) une section
$s=N'(s')$ de $N'(g_*({\mathcal{L}'}^{\otimes n}))$ au-dessus

\oldpage[II]{134}

de $X$. Nous allons voir que si $t$ est la section $s$ considérée comme
section de $\mathcal{M}^{\otimes n}$ au-dessus de $Z$, $t$ répond à la
question. Posons en effet
\[
\label{II.6.6.9.2-fr}
V=X-g(X'-X'_{s'})
\tag{6.6.9.2}
\]
qui est un ouvert de $X$ contenant $x$ et contenu dans $U$, en vertu de
(\hyperref[II.6.6.9.1-fr]{6.6.9.1}) et
(\hyperref[II.6.1.10-fr]{6.1.10}). Nous allons montrer que l'on a
\[
\label{II.6.6.9.3-fr}
h^{-1}(V)\subset Z_t\subset h^{-1}(U),
\tag{6.6.9.3}
\]
ce qui achèvera la démonstration. Il revient au même de dire que l'ensemble
$T$ des $y\in X$ tels que $s_y$ soit inversible, contient $V$ et est contenu
dans $U$. Pour cela, considérons d'abord un ouvert affine $W$ contenu dans
$V$~; $g^{-1}(W)$ est ouvert affine dans $X'$ et en vertu de
(\hyperref[II.6.6.9.2-fr]{6.6.9.2}), $s'_{y'}$ est inversible pour tout
$y'\in g^{-1}(W)$~; en vertu des hypothèses sur $X$ et $\mathcal{B}$, on
peut appliquer les résultats de (\hyperref[II.6.4.7-fr]{6.4.7}) et on voit
que si $y=g(y')$, $s_y$ est inversible, autrement dit $V\subset T$. D'autre
part, il résulte aussi de (\hyperref[II.6.4.7-fr]{6.4.7}) que, inversement,
si $s_y$ est inversible, il en est de même de $s'_{y'}$, ce qui implique
$y'\in g^{-1}(U)$ par (\hyperref[II.6.6.9.1-fr]{6.6.9.1}) et par suite
$y\in U$, d'où $T\subset U$ dans ce cas.

On passe de là au cas général par le même raisonnement que dans
(\hyperref[II.6.6.2-fr]{6.6.2}), remplaçant $X'$ par $X''$ tel que
\[
g'_*(\mathcal{O}_{X''})\longrightarrow
g'_*(\mathcal{O}_{X''})\otimes_{\mathcal{O}_X}\mathcal{R}(X)
\]
soit injectif, et $\mathcal{L}'$ par un $\mathcal{O}_{X''}$-Module ample
$\mathcal{L}''$ tel que
$N'(g_*(\mathcal{L}'))=N'(g'_*(\mathcal{L}''))$. Le lemme
(\hyperref[II.6.6.9-fr]{6.6.9}) est donc démontré dans tous les cas (avec un
choix convenable de $h$).

\begin{corollary}[6.6.10]
\label{II.6.6.10-fr}
Supposons vérifiées les hypothèses de
(\hyperref[II.6.6.9-fr]{6.6.9})~; pour tout $\mathcal{O}_X$-Module inversible
$\mathcal{L}$ tel que $g^*(\mathcal{L})$ soit ample, $h^*(\mathcal{L})$ est
ample.
\end{corollary}

En effet, si on pose $\mathcal{L}'=g^*(\mathcal{L})$, on a alors
$g_*(\mathcal{L}')=\mathcal{L}\otimes_{\mathcal{O}_X}\mathcal{B}$
(\hyperref[0.5.4.10-fr]{0, 5.4.10}), donc
\[
N'(g_*(\mathcal{L}'))=
(\mathcal{L}\otimes_{\mathcal{O}_X}\mathcal{C})^{\otimes n}
\]
(par le même raisonnement que pour
\hyperref[II.6.5.2.4-fr]{6.5.2.4}). On en conclut que l'on a
$\mathcal{M}=(h^*(\mathcal{L}))^{\otimes n}$, et comme $\mathcal{M}$ est
ample, il en est de même de $h^*(\mathcal{L})$
(\hyperref[II.4.5.6-fr]{4.5.6}).

\begin{proposition}[6.6.11]
\label{II.6.6.11-fr}
Soient $Y$ un schéma affine, $X$ un préschéma noethérien réduit,
$f:X\to Y$ un morphisme quasi-compact, $g:X'\to X$ un morphisme fini
surjectif. Soit $W$ une partie ouverte de $X$ telle que, pour tout $x\in W$,
ou bien $X$ est normal au point $x$, ou bien il existe un voisinage ouvert
$T\subset W$ de $x$ tel que $(g_*(\mathcal{O}_{X'}))|T$ soit un
$(\mathcal{O}_X|T)$-Module localement libre. Il existe alors un
$Y$-préschéma réduit $Z$ et un $Y$-morphisme fini et birationnel
$h:Z\to X$ tel que la restriction de $h$ à $h^{-1}(W)$ soit un isomorphisme
$h^{-1}(W)\xrightarrow{\sim}W$ et qui possède la propriété suivante~: pour
tout $\mathcal{O}_X$-Module inversible $\mathcal{L}$ tel que
$g^*(\mathcal{L})$ soit ample relativement à $f\circ g$, $h^*(\mathcal{L})$
est ample relativement à $f\circ h$.
\end{proposition}

Comme $Y$ est affine, $g^*(\mathcal{L})$ est ample, et il s'agit alors de
prouver que pour un choix convenable de $h$, $h^*(\mathcal{L})$ est ample
(\hyperref[II.4.6.6-fr]{4.6.6}). Nous allons voir qu'on peut remplacer $g$
par un morphisme fini surjectif $g':X''\to X$ tel que
${g'}^*(\mathcal{L})$ soit ample et que
$g'_*(\mathcal{O}_{X''})\otimes_{\mathcal{O}_X}\mathcal{R}(X)$ soit un
$\mathcal{R}(X)$-Module \emph{localement libre de rang constant}~; on sera
alors ramené aux conditions de (\hyperref[II.6.6.10-fr]{6.6.10}) et la
proposition sera démontrée.

Pour cela, soit $\mathcal{B}=g_*(\mathcal{O}_{X'})$~; désignons par $X_i$
($1\leq i\leq n$) les sous-préschémas fermés réduits de $X$ qui ont pour
espaces sous-jacents les composantes irréductibles de $X$
(\hyperref[I.5.2.1-fr]{I, 5.2.1})~; ils sont intègres par hypothèse. Soient
$X'_i$ le sous-préschéma fermé

\oldpage[II]{135}

$g^{-1}(X_i)$ de $X'$, $g_i:X'_i\to X_i$ le morphisme $g$ restreint à
$X'_i$, qui est fini (\hyperref[II.6.1.5-fr]{6.1.5, (iii)}) et surjectif~;
soit $k_i$ le \emph{rang} de la $\mathcal{O}_{X_i}$-Algèbre
$\mathcal{B}_i=(g_i)_*(\mathcal{O}_{X'_i})$. Comme
$\mathcal{B}_i\otimes_{\mathcal{O}_{X_i}}\mathcal{R}(X_i)$ est un
préfaisceau constant (\hyperref[I.7.3.5-fr]{I, 7.3.5}), le rang $k_i$ est
aussi le rang de la $(\mathcal{O}_X|U)$-Algèbre $\mathcal{B}|U$ pour tout
ouvert $U$ de $X$ ne rencontrant que la \emph{seule} composante irréductible
$X_i$. Si $T$ est un ensemble ouvert dans $X$ tel que $\mathcal{B}|T$ soit
isomorphe à $\mathcal{O}_X^m|T$, il résulte de la remarque précédente que les
nombres $k_i$ sont égaux à $m$ pour tous les indices $i$ tels que
$T\cap X_i\neq\varnothing$. Soit alors $U$ l'ouvert de $X$ formé des points
au voisinage desquels $\mathcal{B}$ est un $\mathcal{O}_X$-Module localement
libre, et soient $U_j$ ($1\leq j\leq s$) ses composantes connexes, qui sont
ouvertes dans $X$ et en nombre fini (puisque $U$ est noethérien)~; désignons
par $V_j$ le sous-préschéma fermé de $X'$, \emph{adhérence} du sous-préschéma
induit sur l'ouvert $g^{-1}(U_j)$
(\hyperref[I.9.5.11-fr]{I, 9.5.11}). D'après ce qui précède, pour tous les
indices $i$ tels que $X_i\cap U_j\neq\varnothing$, les rangs $k_i$ sont tous
égaux à un même entier $m_j$~; on notera d'ailleurs qu'un même $X_i$ ne peut
rencontrer deux $U_j$ d'indices distincts. Soient $i_\lambda$ les indices
$i$ tels que $X_i\cap U=\varnothing$. Considérons le produit $k$ de tous les
$k_i$, posons $n_i=k/k_i$, et soit $X''$ le préschéma défini comme suit.
Pour chaque $j$ ($1\leq j\leq s$), on considère $k/m_j$ préschémas
isomorphes à $V_j$, et pour chaque $\lambda$, $k/k_{i_\lambda}$ préschémas
isomorphes à $X'_{i_\lambda}$~; $X''$ est la \emph{somme} de tous ces
préschémas. On définit un morphisme $g'':X''\to X'$, se réduisant à
l'injection canonique dans chacun des summandes de $X''$~; il est clair que
$g''$ est un morphisme fini dominant, donc surjectif (puisqu'un morphisme
fini est fermé (\hyperref[II.6.1.10-fr]{6.1.10}))~; posons
$g'=g\circ g''$, qui est un morphisme fini surjectif $X''\to X$~; on a
${g'}^*(\mathcal{L})={g''}^*(g^*(\mathcal{L}))$, donc
${g'}^*(\mathcal{L})$ est un $\mathcal{O}_{X''}$-Module ample
(\hyperref[II.5.1.12-fr]{5.1.12}). Il est clair alors que pour ce nouveau
préschéma $X''$, les rangs définis comme les $k_i$ pour $X'$ sont tous
\emph{égaux à $k$}~; compte tenu de
(\hyperref[I.7.3.3-fr]{I, 7.3.3}), on en conclut aussitôt que pour tout ouvert
affine $T$ de $X$,
\[
\bigl(g'_*(\mathcal{O}_{X''})
\otimes_{\mathcal{O}_X}\mathcal{R}(X)\bigr)|T
\]
est un $(\mathcal{R}(X)|T)$-Module isomorphe à
$(\mathcal{R}(X)|T)^k$.

\begin{corollary}[6.6.12]
\label{II.6.6.12-fr}
Si, dans l'énoncé de (\hyperref[II.6.6.11-fr]{6.6.11}), on a $W=X$, alors
pour qu'un $\mathcal{O}_X$-Module inversible $\mathcal{L}$ soit ample
relativement à $f$, il faut et il suffit que $g^*(\mathcal{L})$ soit ample
relativement à $f\circ g$.
\end{corollary}

\begin{remark}[6.6.13]
\label{II.6.6.13-fr}
Au chapitre III, nous verrons que si $Y$ est noethérien, $f$ de type fini, et
si la restriction de $f$ au sous-préschéma fermé réduit de $X$ ayant $X-W$
pour espace sous-jacent est \emph{propre}, alors la conclusion de
(\hyperref[II.6.6.12-fr]{6.6.12}) est encore valable. Mais nous donnerons au
chapitre V des exemples de schémas algébriques $X$ sur un corps $K$ (le
morphisme structural $X\to\operatorname{Spec}(K)$ n'étant pas propre) dont
le normalisé $X'$ est quasi-affine, mais qui n'est pas quasi-affine (de sorte
que $\mathcal{O}_X$ n'est pas ample, bien que $\mathcal{O}_{X'}$ le soit
(\hyperref[II.5.1.12-fr]{5.1.12}) et que le morphisme $X'\to X$ soit fini
surjectif (\hyperref[II.6.3.10-fr]{6.3.10})). Nous allons voir au numéro
suivant que cette circonstance ne peut se produire lorsqu'on remplace
«~quasi-affine~» par «~affine~».
\end{remark}

\subsection{Le théorème de Chevalley.}
\label{subsection:II.6.7-fr}

Nous allons établir (à l'aide du critère de Serre
(\hyperref[II.5.2.1-fr]{5.2.1})) le théorème suivant, qui a été démontré par
C.~Chevalley par d'autres méthodes, dans le cas des schémas
\emph{algébriques}.

\oldpage[II]{136}

\begin{theorem}[6.7.1]
\label{II.6.7.1-fr}
Soient $X$ un schéma affine, $Y$ un préschéma noethérien, $f:X\to Y$ un
morphisme fini surjectif. Alors $Y$ est un schéma affine.
\end{theorem}

Il est clair que
$f_{\mathrm{red}}:X_{\mathrm{red}}\to Y_{\mathrm{red}}$ est fini
(\hyperref[II.6.1.5-fr]{6.1.5, (vi)})~; comme $X_{\mathrm{red}}$ est un
schéma affine, et qu'il revient au même de dire que $Y$ est affine ou que
$Y_{\mathrm{red}}$ est affine, puisque $Y$ est noethérien
(\hyperref[I.6.1.7-fr]{I, 6.1.7}), on voit qu'on peut supposer $Y$
\emph{réduit}. Pour toute partie fermée $Y'$ de $Y$, il y a alors un seul
sous-préschéma réduit de $Y$ ayant $Y'$ comme espace sous-jacent
(\hyperref[I.5.1.2-fr]{I, 5.1.2})~; son image réciproque $f^{-1}(Y')$,
canoniquement isomorphe à $X\times_Y Y'$
(\hyperref[I.4.4.1-fr]{I, 4.4.1}), est affine comme sous-préschéma fermé de
$X$, et la restriction de $f$ à $f^{-1}(Y')$, qui s'identifie à
$f\times_Y 1_{Y'}$, est un morphisme fini surjectif
(\hyperref[II.6.1.5-fr]{6.1.5, (iii)}). En vertu du principe de récurrence
noethérienne (\hyperref[0.2.2.2-fr]{0, 2.2.2}), on peut donc (compte tenu de
\hyperref[I.6.1.7-fr]{I, 6.1.7}) se ramener à démontrer le théorème sous
l'hypothèse que pour \emph{toute partie fermée $Y'\neq Y$}, tout
sous-préschéma fermé de $Y$ ayant $Y'$ pour espace sous-jacent est affine.
On en conclut que, \emph{pour tout $\mathcal{O}_Y$-Module cohérent
$\mathcal{F}$ dont le support (fermé) $Z$ est distinct de $Y$, on a
$\mathrm{H}^1(Y,\mathcal{F})=0$}. En effet, il existe un sous-préschéma fermé
$Y'$ de $Y$ ayant $Z$ pour espace sous-jacent et tel que, si $j:Y'\to Y$ est
l'injection canonique, on ait
$\mathcal{F}=j_*(j^*(\mathcal{F}))$
(\hyperref[I.9.3.5-fr]{I, 9.3.5})~; par suite
(\hyperref[II.5.2.3-fr]{5.2.3}),
$\mathrm{H}^1(Y,\mathcal{F})
=\mathrm{H}^1(Y',j^*(\mathcal{F}))=0$ d'après
(\hyperref[I.5.1.9.2-fr]{I, 5.1.9.2}).

Supposons d'abord que $Y$ ne soit pas irréductible, et soit $Y'$ une
composante irréductible de $Y$, et $Y''=Y-Y'$~; on désignera encore par
$Y'$ le sous-préschéma fermé réduit de $Y$ ayant $Y'$ pour espace
sous-jacent, et par $j$ l'injection canonique $Y'\to Y$. Soit $\mathcal{F}$
un $\mathcal{O}_Y$-Module cohérent, et considérons l'homomorphisme canonique
\[
\rho:\mathcal{F}\longrightarrow
\mathcal{F}'=j_*(j^*(\mathcal{F}))
\]
(\hyperref[0.4.4.3-fr]{0, 4.4.3})~; $\mathcal{F}'$ est un
$\mathcal{O}_Y$-Module cohérent en vertu de
(\hyperref[0.5.3.10-fr]{0, 5.3.10}) et
(\hyperref[0.5.3.12-fr]{0, 5.3.12}), car on a
$j_*(\mathcal{O}_{Y'})=\mathcal{O}_Y/\mathcal{J}$, en désignant par
$\mathcal{J}$ le faisceau d'idéaux de $\mathcal{O}_Y$ définissant le
sous-préschéma $Y'$. Par suite $\mathcal{G}=\operatorname{Ker}\rho$ et
$\mathcal{K}=\operatorname{Im}\rho$ sont aussi des $\mathcal{O}_Y$-Modules
cohérents (\hyperref[0.5.3.4-fr]{0, 5.3.4})~; or par définition la fibre
$\mathcal{F}'_y$ de $\mathcal{F}'$ au point générique $y$ de $Y'$ est égale
à $\mathcal{F}_y$, car $y$ est intérieur à $Y'$ et on a donc
$\mathcal{J}_y=0$, puisque $Y$ est réduit. On en conclut que $y$ n'est pas
contenu dans le support (fermé) de $\mathcal{G}$~; par ailleurs, le support de
$\mathcal{F}'$ (et \emph{a fortiori} celui de $\mathcal{K}$) est contenu
dans $Y'$~; autrement dit les supports de $\mathcal{G}$ et de $\mathcal{K}$
sont \emph{distincts de $Y$}. On en déduit que
$\mathrm{H}^1(Y,\mathcal{G})=\mathrm{H}^1(Y,\mathcal{K})=0$, et la suite
exacte de cohomologie appliquée à la suite exacte
$0\to\mathcal{G}\to\mathcal{F}\to\mathcal{K}\to0$ montre que
$\mathrm{H}^1(Y,\mathcal{F})=0$. On conclut alors par le critère de Serre
(\hyperref[II.5.2.1-fr]{5.2.1}).

Supposons donc $Y$ irréductible, et par suite \emph{intègre}. On peut aussi
supposer que $X$ est \emph{intègre}~: en effet, si on désigne par $X_i$ les
sous-préschémas fermés réduits de $X$ ayant pour espaces sous-jacents les
composantes irréductibles de $X$ (\hyperref[I.5.2.1-fr]{I, 5.2.1}) et par
$f_i$ la restriction de $f$ à $X_i$, un des $f_i$ au moins est dominant, et
comme c'est un morphisme fini (\hyperref[II.6.1.5-fr]{6.1.5}), il est
surjectif (\hyperref[II.6.1.10-fr]{6.1.10})~; comme d'autre part $X_i$ est un
schéma affine, on voit qu'on peut remplacer $X$ par $X_i$ dans l'énoncé.

\begin{lemma}[6.7.1.1]
\label{II.6.7.1.1-fr}
Soient $X$, $Y$ deux préschémas intègres noethériens, $x$ (resp. $y$) le
point générique de $X$ (resp. $Y$), $f:X\to Y$ un morphisme fini surjectif.
Soit $\mathcal{L}$ un $\mathcal{O}_X$-Module inversible tel qu'il existe un
voisinage ouvert affine $U$ de $y$ et une section
$g\in\Gamma(X,\mathcal{L})$ pour lesquels

\oldpage[II]{137}

$x\in X_g\subset f^{-1}(U)$. Alors il existe deux entiers $m>0$, $n>0$, un
homomorphisme
$u:\mathcal{O}_Y^m\to f_*(\mathcal{L}^{\otimes n})$ et un voisinage ouvert
$V$ de $y$ tels que la restriction $u|V$ soit un isomorphisme de
$\mathcal{O}_Y^m|V$ sur $f_*(\mathcal{L}^{\otimes n})|V$.
\end{lemma}

Soient $C$ l'anneau (intègre) de $U$, $k=\mathcal{O}_y$ son corps des
fractions~: comme $f$ est fini, $U'=f^{-1}(U)$ est affine
(\hyperref[II.1.3.2-fr]{1.3.2})~; soit $D$ son anneau (intègre) de corps des
fractions $K=\mathcal{O}_x$~; par hypothèse $D$ est un $C$-module de type
fini (\hyperref[II.6.1.4-fr]{6.1.4}), donc $K$ est une extension de rang fini
de $k$. La fibre
\[
f^{-1}(y)=X\times_Y\operatorname{Spec}(k(y))
=X\times_Y\operatorname{Spec}(\mathcal{O}_y)
\]
s'identifie à $\operatorname{Spec}(K)$
(\hyperref[I.3.6.5-fr]{I, 3.6.6})~; soient $s_i$ ($1\leq i\leq m$) des
éléments de $D$ formant une base de $K$ sur $k$. Il existe $n>0$ tel que les
sections $(s_i|X_g)g^{\otimes n}$ de $\mathcal{L}^{\otimes n}$ au-dessus de
$X_g$ se prolongent en sections $b_i$ ($1\leq i\leq m$) de
$\mathcal{L}^{\otimes n}$ au-dessus de $X$
(\hyperref[I.9.3.1-fr]{I, 9.3.1}). Les $b_i$ sont aussi par définition des
sections de $f_*(\mathcal{L}^{\otimes n})$ au-dessus de $Y$ et définissent
donc un homomorphisme
$u:\mathcal{O}_Y^m\to f_*(\mathcal{L}^{\otimes n})$
(\hyperref[0.5.1.1-fr]{0, 5.1.1})~; nous allons voir que $u$ répond à la
question. On a $\mathcal{L}^{\otimes n}|U'=\widetilde{M}$, où $M$ est un
$D$-module de type fini, donc si $\varphi$ est l'injection $C\to D$
correspondant au morphisme $f^{-1}(U)\to U$ restriction de $f$,
$M_{[\varphi]}$ est un $C$-module de type fini~; par suite
\[
f_*(\mathcal{L}^{\otimes n})|U=(M_{[\varphi]})^\sim
\]
(\hyperref[I.1.6.3-fr]{I, 1.6.3}) est cohérent et comme $U$ est un ouvert
affine quelconque de $Y$, $f_*(\mathcal{L}^{\otimes n})$ est cohérent~; en
outre $u|U=\widetilde{\theta}$, où $\theta$ est un $C$-homomorphisme
$C^m\to M_{[\varphi]}$, et $u_y=\theta_y$ est l'homomorphisme
\[
\theta\otimes1:K^m=C^m\otimes K\longrightarrow
M_{[\varphi]}\otimes K,
\]
or ce dernier est par définition un \emph{isomorphisme}, car les $(b_i)_x$
forment une base de $M_{[\varphi]}\otimes K$ sur $k$, $g_x$ étant par
hypothèse un générateur de $(\mathcal{L}^{\otimes n})_x$. On en conclut que
les supports des $\mathcal{O}_Y$-Modules $\operatorname{Ker}u$ et
$\operatorname{Coker}u$ ne contiennent pas $y$~; comme ces
$\mathcal{O}_Y$-Modules sont cohérents
(\hyperref[0.5.3.3-fr]{0, 5.3.3}), leurs supports sont fermés
(\hyperref[0.5.2.2-fr]{0, 5.2.2}), d'où le lemme.

Cela étant, les hypothèses du lemme
(\hyperref[II.6.7.1.1-fr]{6.7.1.1}) sont remplies dans le cas que nous
considérons en prenant $\mathcal{L}=\mathcal{O}_X$, puisque $X$ est affine
(\hyperref[I.1.1.10-fr]{I, 1.1.10})~; nous poserons
$\mathcal{A}=\mathcal{O}_Y$, $\mathcal{B}=f_*(\mathcal{O}_X)$. En vertu du
critère de Serre (\hyperref[II.5.2.1.1-fr]{5.2.1.1}), il suffit de prouver
que pour tout $\mathcal{O}_Y$-Module cohérent $\mathcal{F}$, on a
$\mathrm{H}^1(Y,\mathcal{F})=0$~; il suffit même de le prouver lorsque
$\mathcal{F}\subset\mathcal{O}_Y$, ce qui entraîne que $\mathcal{F}$ est
sans torsion puisque $Y$ est intègre~; en fait nous allons montrer que
$\mathrm{H}^1(Y,\mathcal{F})=0$ pour \emph{tout} $\mathcal{O}_Y$-Module
cohérent \emph{sans torsion} $\mathcal{F}$. Or, l'homomorphisme
$u:\mathcal{A}^m\to\mathcal{B}$ définit un homomorphisme
\[
v:\mathcal{G}
=\mathcal{H}om_{\mathcal{A}}(\mathcal{B},\mathcal{F})
\longrightarrow
\mathcal{H}om_{\mathcal{A}}(\mathcal{A}^m,\mathcal{F})
=\mathcal{F}^m.
\]
Montrons d'abord que $v$ est \emph{injectif}~: par hypothèse,
$\mathcal{T}=\operatorname{Coker}u$ a un support ne rencontrant pas $V$, donc
est un $\mathcal{O}_Y$-Module de torsion
(\hyperref[I.7.4.6-fr]{I, 7.4.6})~; la suite exacte
\[
\mathcal{A}^m\longrightarrow\mathcal{B}\longrightarrow\mathcal{T}
\longrightarrow0
\]
donne, par exactitude à gauche du foncteur
$\mathcal{H}om_{\mathcal{A}}$, la suite exacte
\[
0\longrightarrow
\mathcal{H}om_{\mathcal{A}}(\mathcal{T},\mathcal{F})
\longrightarrow\mathcal{G}\xrightarrow{v}\mathcal{F}^m.
\]
Mais comme $\mathcal{F}$ est sans torsion, on a
$\mathcal{H}om_{\mathcal{A}}(\mathcal{T},\mathcal{F})=0$
(\hyperref[0.5.2.6-fr]{0, 5.2.6}), d'où notre assertion. On a donc la suite
exacte
\[
0\longrightarrow\mathcal{G}\longrightarrow\mathcal{F}^m
\longrightarrow\operatorname{Coker}v\longrightarrow0
\]

\oldpage[II]{138}

où $\mathcal{G}$ et $\operatorname{Coker}v$ sont des
$\mathcal{O}_Y$-Modules cohérents
(\hyperref[0.5.3.4-fr]{0, 5.3.4} et
\hyperref[0.5.3.5-fr]{5.3.5})~; en vertu de la suite exacte de cohomologie,
il suffira de montrer que
$\mathrm{H}^1(Y,\mathcal{G})
=\mathrm{H}^1(Y,\operatorname{Coker}v)=0$ pour en déduire
$\mathrm{H}^1(Y,\mathcal{F}^m)
=(\mathrm{H}^1(Y,\mathcal{F}))^m=0$, et par suite
$\mathrm{H}^1(Y,\mathcal{F})=0$. Or la restriction $v|V$ est un isomorphisme,
donc le support de $\operatorname{Coker}v$ est distinct de $Y$, d'où
$\mathrm{H}^1(Y,\operatorname{Coker}v)=0$ par l'hypothèse du début. D'autre
part, $\mathcal{G}$ est un $\mathcal{B}$-Module cohérent
(\hyperref[I.9.6.4-fr]{I, 9.6.4})~; comme $X$ est affine au-dessus de $Y$, il
existe un $\mathcal{O}_X$-Module quasi-cohérent $\mathcal{K}$ tel que
$\mathcal{G}$ soit isomorphe à $f_*(\mathcal{K})$
(\hyperref[II.1.4.3-fr]{1.4.3})~; comme $X$ est affine, on a
$\mathrm{H}^1(X,\mathcal{K})=0$
(\hyperref[I.5.1.9.2-fr]{I, 5.1.9.2}), donc aussi
$\mathrm{H}^1(Y,\mathcal{G})=0$ par
(\hyperref[II.5.2.3-fr]{5.2.3}), ce qui achève la démonstration du théorème
(\hyperref[II.6.7.1-fr]{6.7.1}).

\begin{corollary}[6.7.2]
\label{II.6.7.2-fr}
Soient $X$ un préschéma noethérien, $(X_i)_{1\leq i\leq n}$ un recouvrement
fini de l'espace $X$ formé d'ensembles fermés. Pour que $X$ soit affine, il
faut et il suffit que pour chaque $i$, il existe un sous-préschéma fermé de
$X$, ayant $X_i$ pour espace sous-jacent et qui soit affine.
\end{corollary}

En effet, s'il en est ainsi, soit $X'$ le schéma \emph{somme} des $X_i$~; il
est clair que $X'$ est affine, et on définit un morphisme surjectif
$f:X'\to X$ en prenant pour restriction de $f$ à $X_i$ l'injection
canonique. Tout revient à vérifier que $f$ est \emph{fini}, en vertu de
(\hyperref[II.6.7.1-fr]{6.7.1}), et cela a été vu en
(\hyperref[II.6.1.5-fr]{6.1.5}).

\begin{corollary}[6.7.3]
\label{II.6.7.3-fr}
Pour qu'un préschéma noethérien $X$ soit affine, il faut et il suffit que les
sous-préschémas fermés réduits ayant pour espaces sous-jacents les composantes
irréductibles de $X$ soient affines.
\end{corollary}

\section{Critères valuatifs}
\label{section:II.7-fr}

Nous donnons dans ce paragraphe des critères valuatifs de séparation et de
propreté pour un morphisme, c'est-à-dire des critères qui font intervenir un
schéma auxiliaire variable de la forme $\operatorname{Spec}(A)$, où $A$ est
un anneau de valuation. Moyennant des hypothèses «~noethériennes~»
convenables, on peut dans ces critères se restreindre au cas où $A$ est un
anneau de valuation \emph{discrète}. Ce sera sans doute le seul cas que nous
aurons à appliquer par la suite, et nous n'avons introduit les anneaux de
valuation quelconques, dans le cas général, que pour faire le lien avec des
développements classiques.

\subsection{Rappels sur les anneaux de valuation.}
\label{subsection:II.7.1-fr}

\begin{env}[7.1.1]
\label{II.7.1.1-fr}
Parmi les diverses propriétés équivalentes qui caractérisent les anneaux de
valuation, nous utiliserons la suivante~: un anneau $A$ est dit
\emph{anneau de valuation} si c'est un anneau intègre qui n'est pas un
corps, et si, dans l'ensemble des anneaux locaux contenus dans le corps des
fractions $K$ de $A$ et distincts de $K$, $A$ est \emph{maximal} pour la
relation de domination (\hyperref[I.8.1.1-fr]{I, 8.1.1}). Rappelons qu'un
anneau de valuation est \emph{intégralement clos}. Si $A$ est un anneau de
valuation, $A_{\mathfrak{p}}$ est aussi un anneau de valuation pour tout
idéal premier $\mathfrak{p}\neq0$ de $A$.
\end{env}

\begin{env}[7.1.2]
\label{II.7.1.2-fr}
Soient $K$ un corps, $A$ un sous-anneau local de $K$ qui n'est pas un corps~;

\oldpage[II]{139}

il existe alors un anneau de valuation ayant $K$ pour corps des fractions et
qui domine $A$ (\cite[p.~1-07, lemme~2]{I-1}).

D'autre part, soient $B$ un anneau de valuation, $k$ son corps des restes,
$K$ son corps des fractions, $L$ une extension de $k$. Alors il existe un
anneau de valuation \emph{complet} $C$ qui domine $B$ et dont le corps des
restes est $L$. En effet, $L$ est extension algébrique d'une extension
transcendante pure $L'=k(T_\mu)_{\mu\in M}$~; on sait qu'on peut prolonger la
valuation de $K$ correspondant à $B$ à une valuation de
$K'=K(T_\mu)_{\mu\in M}$ de sorte que $L'$ soit le corps des restes de cette
valuation (\cite[p.~98]{II-24})~; remplaçant $B$ par le complété de l'anneau
de cette valuation prolongée, on voit qu'on peut se limiter au cas où $B$ est
complet et $L$ est une clôture algébrique de $k$. Si $\overline{K}$ est une
clôture algébrique de $K$, on peut alors prolonger à $\overline{K}$ la
valuation qui définit $B$, et le corps des restes correspondant est une
clôture algébrique de $k$, comme on le voit en relevant dans $\overline{K}$
les coefficients d'un polynôme unitaire de $k[T]$. On est donc finalement
ramené au cas où $L=k$ et il suffit alors de prendre pour $C$ le complété de
$B$ pour répondre à la question.
\end{env}

\begin{env}[7.1.3]
\label{II.7.1.3-fr}
Soient $K$ un corps, $A$ un sous-anneau de $K$~; la fermeture intégrale $A'$
de $A$ dans $K$ est l'intersection des anneaux de valuation contenant $A$ et
ayant $K$ pour corps des fractions
(\cite[p.~51, th.~2]{I-11}). La proposition
(\hyperref[II.7.1.2-fr]{7.1.2}) s'exprime sous la forme géométrique
équivalente~:
\end{env}

\begin{proposition}[7.1.4]
\label{II.7.1.4-fr}
Soient $Y$ un préschéma, $p:X\to Y$ un morphisme, $x$ un point de $X$,
$y=p(x)$, $y'\neq y$ une spécialisation
(\hyperref[0.2.1.2-fr]{0, 2.1.2}) de $y$. Il existe alors un schéma local
$Y'$, spectre d'un anneau de valuation, et un morphisme séparé $f:Y'\to Y$
tel que, si $a$ désigne l'unique point fermé de $Y'$ et $b$ le point
générique de $Y'$, on ait $f(a)=y'$ et $f(b)=y$. On peut en outre supposer
vérifiée l'une des deux propriétés additionnelles suivantes~:
\begin{enumerate}
\item[\rm{(i)}] $Y'$ est le spectre d'un anneau de valuation \emph{complet}
dont le corps des restes est algébriquement clos, et il existe un
$k(y)$-homomorphisme $k(x)\to k(b)$.
\item[\rm{(ii)}] Il existe un $k(y)$-isomorphisme
$k(x)\xrightarrow{\sim}k(b)$.
\end{enumerate}
\end{proposition}

Soit $Y_1$ le sous-préschéma fermé réduit de $Y$ ayant
$\overline{\{y\}}$ comme espace sous-jacent
(\hyperref[I.5.2.1-fr]{I, 5.2.1}), et soit $X_1$ le sous-préschéma fermé
image réciproque $p^{-1}(Y_1)$~; comme $y'\in\overline{\{y\}}$ par hypothèse
et que $k(x)$ est le même dans $X$ et dans $X_1$, on peut supposer $Y$
\emph{intègre}, de point générique $y$~; $\mathcal{O}_{y'}$ est alors un
anneau local intègre qui n'est pas un corps, et dont le corps des fractions
est $\mathcal{O}_y=k(y)$, et $k(x)$ est une extension de $k(y)$. Pour
réaliser les conditions $f(a)=y'$ et $f(b)=y$ ainsi que la condition
additionnelle (i) (resp. (ii)), on prendra
$Y'=\operatorname{Spec}(A')$, où $A'$ est un anneau de valuation dominant
$\mathcal{O}_{y'}$ et qui est complet et a un corps des restes
algébriquement clos extension de $k(x)$ (resp. un anneau de valuation $A'$
dominant $\mathcal{O}_{y'}$ et dont $k(x)$ est le corps des fractions)~;
l'existence de $A'$ est garantie par
(\hyperref[II.7.1.2-fr]{7.1.2}).

\begin{env}[7.1.5]
\label{II.7.1.5-fr}
Rappelons qu'un anneau local $A$ est dit \emph{de dimension $1$} s'il existe
un idéal premier distinct de l'idéal maximal $\mathfrak{m}$, et si tout idéal
premier de $A$ distinct de $\mathfrak{m}$ est un idéal premier
\emph{minimal}~; lorsque $A$ est \emph{intègre}, il revient au même de dire
que $\mathfrak{m}$ et $(0)$ sont les seuls idéaux premiers et
$\mathfrak{m}\neq(0)$~; autrement dit, $Y=\operatorname{Spec}(A)$ est
réduit à deux

\oldpage[II]{140}

points $a$, $b$~: $a$ est l'unique point \emph{fermé}, on a
$\mathfrak{j}_a=\mathfrak{m}$ et $k(a)=k$ est le \emph{corps résiduel}
$k=A/\mathfrak{m}$~; $b$ est le \emph{point générique} de $Y$,
$\mathfrak{j}_b=(0)$, l'ensemble $\{b\}$ étant l'unique ensemble ouvert de
$Y$ distinct de $\varnothing$ et de $Y$ (ensemble ouvert qui est donc
\emph{partout dense}), et $k(b)=K$ est le \emph{corps des fractions} de $A$.
\end{env}

\begin{env}[7.1.6]
\label{II.7.1.6-fr}
Pour un anneau local $A$, noethérien et de dimension $1$, on sait que les
conditions suivantes sont équivalentes
(\cite[p.~2-08 et 17-01]{I-1})~:
\begin{enumerate}
\item[\rm{a)}] $A$ est normal~;
\item[\rm{b)}] $A$ est régulier~;
\item[\rm{c)}] $A$ est un anneau de valuation~;
\end{enumerate}
en outre, $A$ est alors un \emph{anneau de valuation discrète}. Les prop.
(\hyperref[II.7.1.2-fr]{7.1.2}) et
(\hyperref[II.7.1.3-fr]{7.1.3}) ont alors les analogues suivants pour les
anneaux de valuation discrète~:
\end{env}

\begin{proposition}[7.1.7]
\label{II.7.1.7-fr}
Soient $A$ un anneau local intègre noethérien qui n'est pas un corps, $K$ son
corps des fractions, $L$ une extension de type fini de $K$~; il existe alors
un anneau de valuation discrète ayant $L$ pour corps des fractions et
dominant $A$.
\end{proposition}

Supposons d'abord $L=K$. Soient $\mathfrak{m}$ l'idéal maximal de $A$,
$(x_1,\ldots,x_n)$ un système de générateurs non nuls de $\mathfrak{m}$,
$B$ le sous-anneau
$A[x_2/x_1,\ldots,x_n/x_1]$ de $K$, qui est noethérien. Il est immédiat que
l'idéal $\mathfrak{m}B$ de $B$ est identique à l'idéal principal $x_1B$~;
si $\mathfrak{p}$ est un idéal premier minimal de $x_1B$, on sait que
$\mathfrak{p}$ est de rang $1$ (\cite[t.~I, p.~277]{I-13})~; autrement dit
$B_{\mathfrak{p}}$ est un anneau local noethérien \emph{de dimension $1$}~;
il est clair que
$\mathfrak{p}B_{\mathfrak{p}}\cap A$ est un idéal de $A$ contenant
$\mathfrak{m}$ et ne contenant pas $1$, donc égal à $\mathfrak{m}$, et par
suite $B_{\mathfrak{p}}$ \emph{domine} $A$
(\hyperref[I.8.1.1-fr]{I, 8.1.1}). Il résulte du th. de Krull-Akizuki
(\cite[p.~293]{II-25}) que la clôture intégrale $C$ de $B_{\mathfrak{p}}$
est un anneau noethérien (bien que $C$ ne soit pas nécessairement un
$B_{\mathfrak{p}}$-module de type fini)~; si $\mathfrak{n}$ est un idéal
maximal de $C$, $C_{\mathfrak{n}}$ est un anneau local noethérien normal de
dimension $1$ (\cite[p.~295]{II-25}), donc un anneau de valuation discrète,
qui domine $B_{\mathfrak{p}}$ et \emph{a fortiori} $A$.

Si maintenant $L$ est une extension de type fini de $K$, on peut, d'après ce
qui précède, se limiter au cas où $A$ est déjà un anneau de valuation
discrète. Soit $w$ une valuation de $K$ associée à $A$~; il existe une
valuation discrète $w'$ de $L$ qui \emph{prolonge} $w$~: on se ramène en
effet par récurrence sur le nombre de générateurs de $L$ au cas où
$L=K(\alpha)$, et alors la proposition est classique
(\cite[p.~106]{II-24}).

\begin{corollary}[7.1.8]
\label{II.7.1.8-fr}
Soient $A$ un anneau intègre noethérien, $K$ son corps des fractions, $L$ une
extension de type fini de $K$. Alors la fermeture intégrale de $A$ dans $L$
est l'intersection des anneaux de valuation discrète ayant $L$ pour corps des
fractions et contenant $A$.
\end{corollary}

En effet, un tel anneau de valuation discrète, étant normal, contient
\emph{a fortiori} tout élément de $L$ entier sur $A$. Il suffit donc de
prouver que si $x\in L$ n'est pas entier sur $A$, il existe un anneau de
valuation discrète $C$ ayant $L$ pour corps des fractions, contenant $A$ et
ne contenant pas $x$. L'hypothèse sur $x$ signifie en effet que l'on a
$x\notin B=A[1/x]$, autrement dit $1/x$ n'est pas inversible dans l'anneau
noethérien $B$. Il y a donc un idéal premier $\mathfrak{p}$ de $B$ contenant
$1/x$. L'anneau local intègre $B_{\mathfrak{p}}$ est noethérien et contenu
dans $L$, qui est une extension de type fini du corps des fractions de
$B_{\mathfrak{p}}$ (ce dernier contenant $K$). En vertu de
(\hyperref[II.7.1.7-fr]{7.1.7}), il y a donc un anneau de valuation discrète
$C$, dominant $B_{\mathfrak{p}}$ et ayant $L$ pour corps des fractions~;
comme $1/x\in\mathfrak{p}B_{\mathfrak{p}}$ appartient à l'idéal maximal de
$C$, on a $x\notin C$, ce qui conclut la démonstration.

La forme géométrique de (\hyperref[II.7.1.7-fr]{7.1.7}) est la suivante~:

\oldpage[II]{141}
\begin{proposition}[7.1.9]
\label{II.7.1.9-fr}
Soient $Y$ un préschéma localement noethérien, $p:X\to Y$ un morphisme
localement de type fini, $x$ un point de $X$, $y=p(x)$, $y'\neq y$ une
spécialisation de $y$. Il existe alors un schéma local $Y'$, spectre d'un
anneau de valuation discrète, un morphisme séparé $f:Y'\to Y$ et une
$Y$-application rationnelle $g$ de $Y'$ dans $X$, tels que, si $a$ désigne le
point fermé de $Y'$ et $b$ son point générique, on ait $f(a)=y'$,
$f(b)=y$, $g(b)=x$, et que dans le diagramme commutatif
\[
\tag{7.1.9.1}
\label{II.7.1.9.1-fr}
\begin{tikzcd}[column sep=large,row sep=large]
& k(x) \arrow[dl,"\gamma"'] \\
k(b) & k(y) \arrow[u,"\pi"] \arrow[l,"\varphi"']
\end{tikzcd}
\]
(où $\pi$, $\varphi$, $\gamma$ sont les homomorphismes correspondant
respectivement à $p$, $f$ et $g$), $\gamma$ soit une bijection.
\end{proposition}

Comme dans (\hyperref[II.7.1.4-fr]{7.1.4}), on peut se ramener au cas où $Y$
est intègre de point générique $y$ (compte tenu de
(\hyperref[I.6.3.4-fr]{I, 6.3.4, (iv)})) et comme la question est locale sur
$X$ et $Y$, on peut supposer $p$ de type fini~; on est alors dans la situation
de (\hyperref[II.7.1.4-fr]{7.1.4}), avec la propriété supplémentaire que
$k(x)$ est une extension \emph{de type fini} de $k(y)$
(\hyperref[I.6.4.11-fr]{I, 6.4.11}) et que $\mathcal{O}_{y'}$ est noethérien~;
cela permet d'appliquer (\hyperref[II.7.1.7-fr]{7.1.7}) et de prendre
$Y'=\operatorname{Spec}(A')$, où $A'$ est un anneau de valuation discrète
dominant $\mathcal{O}_{y'}$ et dont $k(x)$ est le corps des fractions. On a
donc ainsi défini un diagramme commutatif
(\hyperref[II.7.1.9.1-fr]{7.1.9.1}), où $\gamma$ est une bijection, et $\pi$
et $\varphi$ correspondent aux morphismes $p$ et $f$. En outre, comme $X$ et
$Y$ sont localement noethériens (\hyperref[I.6.6.2-fr]{I, 6.6.2}) et que $Y'$
est intègre, il y a une $Y$-application rationnelle $g$ et une seule de $Y'$
dans $X$ à laquelle correspond l'isomorphisme $\gamma$
(\hyperref[I.7.1.15-fr]{I, 7.1.15}), ce qui achève la démonstration.

\subsection{Critère valuatif de séparation}
\label{subsection:II.7.2-fr}

\begin{proposition}[7.2.1]
\label{II.7.2.1-fr}
Soient $X$, $Y$ deux préschémas, $f:X\to Y$ un morphisme quasi-compact. Les
deux conditions suivantes sont équivalentes~:
\begin{enumerate}
\item[\rm a)] Le morphisme $f$ est fermé.
\item[\rm b)] Pour tout $x\in X$ et toute spécialisation $y'$ de $y=f(x)$,
distincte de $y$, il existe une spécialisation $x'$ de $x$ telle que
$f(x')=x$.
\end{enumerate}
\end{proposition}

La condition b) exprime que $f(\overline{\{x\}})=\overline{\{y\}}$ et est donc
conséquence de a). Pour montrer que b) entraîne a), considérons une partie
fermée $X'$ de l'espace sous-jacent $X$~; soit $Y'=\overline{f(X')}$ et
montrons que $Y'=f(X')$. Considérons les sous-préschémas fermés réduits de
$X$ et $Y$ respectivement ayant pour espaces sous-jacents $X'$ et $Y'$
(\hyperref[I.5.2.1-fr]{I, 5.2.1})~; il y a alors un morphisme $f':X'\to Y'$ tel
que le diagramme
\[
\begin{tikzcd}[column sep=large,row sep=large]
X' \arrow[r,"f'"] \arrow[d] & Y' \arrow[d] \\
X \arrow[r,"f"'] & Y
\end{tikzcd}
\]
soit commutatif (\hyperref[I.5.2.2-fr]{I, 5.2.2}), et comme $f$ est
quasi-compact, il en est de même de $f'$. On est donc ramené à prouver que si
$f$ est un morphisme quasi-compact et \emph{dominant}, alors la
\oldpage[II]{142}
condition b) implique que $f(X)=Y$. Or, soit $y'$ un point de $Y$ et soit $y$
le point générique d'une composante irréductible de $Y$ contenant $y'$~; en
vertu de b), il suffit de montrer que $f^{-1}(y)$ n'est pas vide. Mais on sait
que cette propriété est conséquence du fait que $f$ est dominant et
quasi-compact (\hyperref[I.6.6.5-fr]{I, 6.6.5}).

\begin{corollary}[7.2.2]
\label{II.7.2.2-fr}
Soit $f:X\to Y$ un morphisme d'immersion quasi-compact. Pour que l'espace
sous-jacent $X$ soit fermé dans $Y$, il faut et il suffit qu'il contienne toute
spécialisation (dans $Y$) de chacun de ses points.
\end{corollary}

\begin{proposition}[7.2.3]
\label{II.7.2.3-fr}
Soient $Y$ un préschéma (resp. un préschéma localement noethérien), $f:X\to Y$
un morphisme (resp. un morphisme localement de type fini). Les conditions
suivantes sont équivalentes~:
\begin{enumerate}
\item[\rm a)] $f$ est séparé.
\item[\rm b)] Le morphisme diagonal $X\to X\times_YX$ est quasi-compact, et
pour tout $Y$-préschéma de la forme $Y'=\operatorname{Spec}(A)$, où $A$ est un
anneau de valuation (resp. un anneau de valuation discrète), deux
$Y$-morphismes de $Y'$ dans $X$ qui coïncident au point générique de $Y'$ sont
égaux.
\item[\rm c)] Le morphisme diagonal $X\to X\times_YX$ est quasi-compact, et
pour tout $Y$-préschéma de la forme $Y'=\operatorname{Spec}(A)$, où $A$ est un
anneau de valuation (resp. un anneau de valuation discrète), deux
$Y'$-sections de $X'=X_{(Y')}$ qui coïncident au point générique de $Y'$ sont
égales.
\end{enumerate}
\end{proposition}

L'équivalence de b) et de c) résulte de la correspondance biunivoque entre
$Y$-morphismes de $Y'$ dans $X$ et $Y'$-sections de $X'$
(\hyperref[I.3.3.14-fr]{I, 3.3.14}). Si $X$ est séparé sur $Y$, la condition b)
est vérifiée en vertu de (\hyperref[I.7.2.2.1-fr]{I, 7.2.2.1}), puisque $Y'$
est intègre. Il reste à prouver que b) entraîne que le morphisme diagonal
$\Delta:X\to X\times_YX$ est fermé, et il revient au même de montrer qu'il
vérifie le critère de (\hyperref[II.7.2.2-fr]{7.2.2}). Or, soit $z$ un point de
la diagonale $\Delta(X)$, $z'\neq z$ une spécialisation de $z$ dans
$X\times_YX$. Il existe alors (\hyperref[II.7.1.4-fr]{7.1.4}) un anneau de
valuation $A$ et un morphisme $f$ de $Y'=\operatorname{Spec}(A)$ dans
$X\times_YX$, qui applique le point fermé $a$ de $Y'$ sur $z'$ et le point
générique $b$ de $Y'$ sur $z$~; ce morphisme fait de $Y'$ un
$(X\times_YX)$-préschéma, et \emph{a fortiori} un $Y$-préschéma. Si on compose
avec $f$ les deux projections de $X\times_YX$, on obtient deux $Y$-morphismes
$g_1$, $g_2$ de $Y'$ dans $X$, qui par hypothèse coïncident au point $b$~; ils
sont donc égaux à un même morphisme $g$, ce qui signifie
(\hyperref[I.5.3.1-fr]{I, 5.3.1}) que $f$ se factorise en $f=\Delta\circ g$, et
par suite $z'\in\Delta(X)$. Lorsqu'on suppose $Y$ localement noethérien et $f$
localement de type fini, $X\times_YX$ est localement noethérien
(\hyperref[I.6.6.7-fr]{I, 6.6.7})~; on peut alors reprendre le raisonnement
précédent en y supposant que $A$ est un anneau de valuation discrète, en vertu
de (\hyperref[II.7.1.9-fr]{7.1.9}).

\begin{remark}[7.2.4]
\label{II.7.2.4-fr}
\begin{enumerate}
\item[\rm (i)] L'hypothèse que le morphisme $\Delta$ est quasi-compact est
toujours vérifiée lorsque $Y$ est localement noethérien et $f$ localement de
type fini, car $X\times_YX$ est alors localement noethérien
(\hyperref[I.6.6.4-fr]{I, 6.6.4, (i)}). Dans le cas général, elle signifie aussi
que pour tout recouvrement $(U_\alpha)$ de $X$ par des ouverts affines, les
ensembles $U_\alpha\cap U_\beta$ sont \emph{quasi-compacts}.
\item[\rm (ii)] Pour que $f$ soit séparé, il \emph{suffit} que la condition b)
ou la condition c) soit vérifiée pour un anneau de valuation $A$ qui est
\emph{complet} et dont le corps des restes est \emph{algébriquement clos}~;
cela résulte en effet de la démonstration de
(\hyperref[II.7.2.3-fr]{7.2.3}) et de
(\hyperref[II.7.1.4-fr]{7.1.4}).
\end{enumerate}
\end{remark}

\oldpage[II]{143}
\subsection{Critère valuatif de propreté}
\label{subsection:II.7.3-fr}

\begin{proposition}[7.3.1]
\label{II.7.3.1-fr}
Soient $A$ un anneau de valuation, $Y=\operatorname{Spec}(A)$, $b$ le point
générique de $Y$, $X$ un schéma intègre, $f:X\to Y$ un morphisme fermé tel
que $f^{-1}(b)$ se réduise à un point $x$ et que l'homomorphisme correspondant
$k(b)\to k(x)$ soit bijectif. Alors $f$ est un isomorphisme.
\end{proposition}

Comme $f$ est fermé et dominant, on a $f(X)=Y$~; il suffit
(\hyperref[I.4.2.2-fr]{I, 4.2.2}) de prouver que pour tout $y'\neq b$ dans $Y$,
il existe un \emph{seul} point $x'$ tel que $f(x')=y'$ et que l'homomorphisme
correspondant $\mathcal{O}_{y'}\to\mathcal{O}_{x'}$ soit bijectif, car $f$ sera
alors un homéomorphisme. Or, si $f(x')=y'$, $\mathcal{O}_{x'}$ est un anneau
local contenu dans $K=k(x)=k(b)$ et dominant $\mathcal{O}_{y'}$~; ce dernier
est l'anneau local $A_{y'}$, donc est un anneau de valuation
(\hyperref[II.7.1.1-fr]{7.1.1}) ayant $K$ pour corps des fractions. Par
ailleurs on a $\mathcal{O}_{x'}\neq K$ puisque $x'$ n'est pas le point
générique de $X$ (\hyperref[0.2.1.3-fr]{0, 2.1.3})~; on en conclut que
$\mathcal{O}_{x'}=\mathcal{O}_{y'}$. Comme $X$ est un schéma intègre, la
relation $\mathcal{O}_{x'}=\mathcal{O}_{x''}$ entraîne $x'=x''$
(\hyperref[I.8.2.2-fr]{I, 8.2.2}), ce qui achève la démonstration.

\begin{env}[7.3.2]
\label{II.7.3.2-fr}
Soient $A$ un anneau de valuation, $Y=\operatorname{Spec}(A)$, $b$ le point
générique de $Y$, de sorte que $\mathcal{O}_b=k(b)$ est égal à $K$, corps des
fractions de $A$~; soit $f:X\to Y$ un morphisme. On sait
(\hyperref[I.7.1.4-fr]{I, 7.1.4}) que les $Y$-sections \emph{rationnelles} de
$X$ sont en correspondance biunivoque avec les \emph{germes} de $Y$-sections
(définies dans des voisinages de $b$) au point $b$, d'où une application
canonique
\[
\tag{7.3.2.1}
\label{II.7.3.2.1-fr}
\Gamma_{\mathrm{rat}}(X/Y)\longrightarrow
\Gamma(f^{-1}(b)/\operatorname{Spec}(K))
\]
les éléments de $\Gamma(f^{-1}(b)/\operatorname{Spec}(K))$ s'identifiant
d'ailleurs, par définition (\hyperref[I.3.4.5-fr]{I, 3.4.5}) aux \emph{points
de $f^{-1}(b)=X\otimes_AK$ rationnels sur $K$}. Lorsque $f$ est \emph{séparé},
il résulte de (\hyperref[I.5.4.7-fr]{I, 5.4.7}) que l'application
(\hyperref[II.7.3.2.1-fr]{7.3.2.1}) est \emph{injective}, puisque $Y$ est un
schéma intègre.

Composant (\hyperref[II.7.3.2.1-fr]{7.3.2.1}) avec l'application canonique
$\Gamma(X/Y)\to\Gamma_{\mathrm{rat}}(X/Y)$
(\hyperref[I.7.1.2-fr]{I, 7.1.2}), on obtient une application canonique
\[
\tag{7.3.2.2}
\label{II.7.3.2.2-fr}
\Gamma(X/Y)\longrightarrow \Gamma(f^{-1}(b)/\operatorname{Spec}(K)).
\]
Lorsque $f$ est \emph{séparé}, cette application est encore \emph{injective}
(\hyperref[I.5.4.7-fr]{I, 5.4.7}).
\end{env}

\begin{proposition}[7.3.3]
\label{II.7.3.3-fr}
Soient $A$ un anneau de valuation de corps des fractions $K$,
$Y=\operatorname{Spec}(A)$, $b$ le point générique de $Y$, $f:X\to Y$ un
morphisme \emph{séparé} et \emph{fermé}. Alors l'application canonique
(\hyperref[II.7.3.2.2-fr]{7.3.2.2}) est bijective (ce qui revient à dire
qu'elle est \emph{surjective} et entraîne que les $Y$-sections rationnelles de
$X$ sont \emph{partout définies}).
\end{proposition}

Soit donc $x$ un point de $f^{-1}(b)$, \emph{rationnel sur $K$}. Comme $f$ est
séparé, il en est de même du morphisme
$f^{-1}(b)\to\operatorname{Spec}(K)$ correspondant à $f$
(\hyperref[I.5.5.1-fr]{I, 5.5.1, (iv)}), et toute section de $f^{-1}(b)$ étant
une immersion fermée (\hyperref[I.5.4.6-fr]{I, 5.4.6}), $\{x\}$ est fermé
\emph{dans $f^{-1}(b)$}. Considérons le sous-préschéma fermé réduit $X'$ de
$X$ ayant pour espace sous-jacent l'adhérence $\overline{\{x\}}$ de $\{x\}$
dans $X$. Il est clair que la restriction de $f$ à $X'$ vérifie les hypothèses
de (\hyperref[II.7.3.1-fr]{7.3.1}), donc est un \emph{isomorphisme} de $X'$ sur
$Y$, dont l'isomorphisme réciproque est la $Y$-section de $X$ cherchée.

\begin{env}[7.3.4]
\label{II.7.3.4-fr}
Pour énoncer les deux résultats suivants, nous nous servirons d'une
terminologie qui sera justifiée et discutée dans le chapitre IV~: si $F$ est
une partie
\oldpage[II]{144}
d'un préschéma $Y$, on appelle \emph{codimension} de $F$ dans $Y$ et on note
$\operatorname{codim}_YF$ la borne inférieure des entiers
$\dim(\mathcal{O}_z)$, où $z$ parcourt $F$.
\end{env}

\begin{corollary}[7.3.5]
\label{II.7.3.5-fr}
Soient $Y$ un préschéma localement noethérien réduit, $N$ l'ensemble des
points $y\in Y$ où $Y$ n'est pas régulier
(\hyperref[0.4.1.4-fr]{0, 4.1.4})~; on suppose que
$\operatorname{codim}_YN\geq2$. Soit $f:X\to Y$ un morphisme de type fini,
séparé et fermé et soit $g$ une $Y$-section rationnelle de $X$~; si $Y'$ est
l'ensemble des points de $Y$ où $g$ n'est pas définie (ensemble qui est
\emph{fermé} (\hyperref[I.7.2.1-fr]{I, 7.2.1})) on a
$\operatorname{codim}_YY'\geq2$.
\end{corollary}

Il suffit de prouver que $g$ est définie en tout point $z\in Y$ tel que
$\dim\mathcal{O}_z\leq1$. Si $\dim\mathcal{O}_z=0$, $z$ est point générique
d'une composante irréductible de $Y$ (\hyperref[I.1.1.14-fr]{I, 1.1.14}), donc
appartient à tout ouvert partout dense de $Y$, et en particulier au domaine de
définition de $g$. Supposons donc que $\dim\mathcal{O}_z=1$~;
$\mathcal{O}_z$ est alors par hypothèse un anneau local noethérien régulier,
donc (\hyperref[II.7.1.6-fr]{7.1.6}) un \emph{anneau de valuation discrète}.
Soit $Z=\operatorname{Spec}(\mathcal{O}_z)$~; comme $U=Y-Y'$ est partout
dense, $U\cap Z$ n'est pas vide (\hyperref[I.2.4.2-fr]{I, 2.4.2})~; soit $g'$
l'application rationnelle de $Z$ dans $X$ induite par $g$
(\hyperref[I.7.2.8-fr]{I, 7.2.8})~; il suffit de montrer que $g'$ est un
\emph{morphisme} (\hyperref[I.7.2.9-fr]{I, 7.2.9}). Or, $g'$ peut être
considérée comme une $Z$-section rationnelle du $Z$-préschéma
$f^{-1}(Z)=X\times_YZ$~; il est clair que le morphisme $f^{-1}(Z)\to Z$
correspondant à $f$ est fermé, et il résulte de
(\hyperref[I.5.5.1-fr]{I, 5.5.1, (i)}) qu'il est séparé~; on conclut alors de
(\hyperref[II.7.3.3-fr]{7.3.3}) que $g'$ est partout définie~; comme $Z$ est
réduit et $X$ séparé sur $Y$, $g'$ est un morphisme
(\hyperref[I.7.2.2-fr]{I, 7.2.2}).

\begin{corollary}[7.3.6]
\label{II.7.3.6-fr}
Soient $S$ un préschéma localement noethérien, $X$ et $Y$ deux
$S$-préschémas~; on suppose $Y$ réduit, et en outre que l'ensemble $N$ des
points $y\in Y$ où $Y$ n'est pas régulier soit tel que
$\operatorname{codim}_YN\geq2$~; on suppose enfin que le morphisme structural
$X\to S$ soit propre. Soit $f$ une $S$-application rationnelle de $Y$ dans
$X$ et soit $Y'$ l'ensemble des points de $Y$ où $f$ n'est pas définie~;
alors on a $\operatorname{codim}_YY'\geq2$.
\end{corollary}

On sait (\hyperref[I.7.1.2-fr]{I, 7.1.2}) qu'on peut identifier les
$S$-applications rationnelles de $Y$ dans $X$ aux $Y$-sections rationnelles
de $X\times_SY$~; comme le morphisme structural $X\times_SY\to Y$ est fermé
(\hyperref[II.5.4.1-fr]{5.4.1}), on peut appliquer
(\hyperref[II.7.3.5-fr]{7.3.5}), d'où le corollaire.

\begin{remark}[7.3.7]
\label{II.7.3.7-fr}
Les hypothèses faites sur $Y$ dans (\hyperref[II.7.3.5-fr]{7.3.5}) et
(\hyperref[II.7.3.6-fr]{7.3.6}) seront en particulier réalisées lorsque $Y$
est \emph{normal} (\hyperref[0.4.1.4-fr]{0, 4.1.4}), en vertu de
(\hyperref[II.7.1.6-fr]{7.1.6}).
\end{remark}

Nous pouvons caractériser les morphismes universellement fermés (resp.
propres) par une réciproque de (\hyperref[II.7.3.3-fr]{7.3.3})~:

\begin{theorem}[7.3.8]
\label{II.7.3.8-fr}
Soient $Y$ un préschéma (resp. un préschéma localement noethérien), $f:X\to Y$
un morphisme quasi-compact séparé (resp. de type fini). Les conditions
suivantes sont équivalentes~:
\begin{enumerate}
\item[\rm a)] $f$ est universellement fermé (resp. propre).
\item[\rm b)] Pour tout $Y$-schéma de la forme
$Y'=\operatorname{Spec}(A)$, où $A$ est un anneau de valuation (resp. un anneau
de valuation discrète) de corps des fractions $K$, l'application canonique
\[
\operatorname{Hom}_Y(Y',X)\longrightarrow
\operatorname{Hom}_Y(\operatorname{Spec}(K),X)
\]
correspondant à l'injection canonique $A\to K$, est surjective (resp.
bijective).
\oldpage[II]{145}
\item[\rm c)] Pour tout $Y$-schéma de la forme $Y'=\operatorname{Spec}(A)$,
où $A$ est un anneau de valuation (resp. un anneau de valuation discrète)
l'application canonique (\hyperref[II.7.3.2.2-fr]{7.3.2.2}) relative au
$Y'$-préschéma $X_{(Y')}$ est surjective (resp. bijective).
\end{enumerate}
\end{theorem}

L'équivalence de b) et c) résulte aussitôt de
(\hyperref[I.3.3.14-fr]{I, 3.3.14})~; a) implique c), car a) entraîne, dans
l'une ou l'autre hypothèse, que $f_{(Y')}$ est séparé
(\hyperref[I.5.5.1-fr]{I, 5.5.1, (iv)}) et fermé, et il suffit d'appliquer
(\hyperref[II.7.3.3-fr]{7.3.3}). Reste à prouver que b) entraîne a).
Plaçons-nous d'abord dans le cas où $Y$ est quelconque, $f$ séparé et
quasi-compact. Si la condition b) est vérifiée par $f$, elle l'est aussi par
$f_{(Y'')}:X_{(Y'')}\to Y''$, où $Y''$ est un $Y$-préschéma quelconque, en
vertu de l'équivalence de b) et c), et du fait que
$X_{(Y'')}\times_{Y''}Y'=X\times_YY'$ pour tout morphisme $Y'\to Y''$
(\hyperref[I.3.3.9.1-fr]{I, 3.3.9.1})~; comme en outre $f_{(Y'')}$ est séparé
et quasi-compact quand $f$ l'est
(\hyperref[I.5.5.1-fr]{I, 5.5.1, (iv)} et
\hyperref[I.6.6.4-fr]{6.6.4, (iii)}), on est ramené à prouver que b) entraîne
que $f$ est \emph{fermé}. Pour cela, il suffit de vérifier la condition b) de
(\hyperref[II.7.2.1-fr]{7.2.1}). Soient donc $x\in X$, $y'$ une spécialisation
de $y=f(x)$, distincte de $y$~; en vertu de
(\hyperref[II.7.1.4-fr]{7.1.4}), il y a un schéma $Y'$, spectre d'un anneau de
valuation, et un morphisme séparé $g:Y'\to Y$ tels que, si $a$ désigne le point
fermé et $b$ le point générique de $Y'$, on ait $g(a)=y'$, $g(b)=y$, et qu'il
existe un $k(y)$-homomorphisme $k(x)\to k(b)$. Ce dernier correspond
canoniquement à un $Y$-morphisme $\operatorname{Spec}(k(b))\to X$
(\hyperref[I.2.4.6-fr]{I, 2.4.6}), et il résulte donc de b) qu'il existe un
$Y$-morphisme $h:Y'\to X$ auquel correspond le morphisme précédent. On a
alors $h(b)=x$~; si l'on pose $h(a)=x'$, $x'$ est spécialisation de $x$, et
l'on a $f(x')=f(h(a))=g(a)=y'$.

Si maintenant $Y$ est localement noethérien et $f$ de type fini, l'hypothèse
b) entraîne d'abord que $f$ est \emph{séparé}, en vertu de
(\hyperref[II.7.2.3-fr]{7.2.3}), le morphisme diagonal $X\to X\times_YX$
étant quasi-compact (\hyperref[II.7.2.4-fr]{7.2.4}). En outre, pour vérifier
que $f$ est propre, il suffit de montrer que $f_{(Y'')}:X_{(Y'')}\to Y''$ est
\emph{fermé} pour tout $Y$-préschéma $Y''$ \emph{de type fini}, compte tenu de
(\hyperref[II.5.6.3-fr]{5.6.3}). Comme alors $Y''$ est localement noethérien,
on peut reprendre le raisonnement fait dans le premier cas en prenant pour
$Y'$ un spectre d'anneau de valuation discrète, et en appliquant
(\hyperref[II.7.1.9-fr]{7.1.9}) au lieu de
(\hyperref[II.7.1.4-fr]{7.1.4}).

\begin{remark}[7.3.9]
\label{II.7.3.9-fr}
\begin{enumerate}
\item[\rm (i)] Lorsque $Y$ est un préschéma quelconque et $f$ un morphisme
séparé, pour que $f$ soit universellement fermé, il \emph{suffit} que la
condition b) ou la condition c) soit vérifiée pour les anneaux de valuation
$A$ \emph{complets} et dont le corps des restes est \emph{algébriquement
clos}~; cela résulte en effet de la démonstration et de
(\hyperref[II.7.1.4-fr]{7.1.4}).
\item[\rm (ii)] On déduit du critère c) de
(\hyperref[II.7.3.8-fr]{7.3.8}) une nouvelle démonstration du fait qu'un
morphisme projectif $X\to Y$ est fermé (\hyperref[II.5.5.3-fr]{5.5.3}), plus
proche des méthodes classiques. On peut en effet supposer $Y$ affine, et par
suite $X$ identifié à un sous-préschéma fermé d'un fibré projectif
$\mathbf{P}_Y^n$ (\hyperref[II.5.3.3-fr]{5.3.3})~; pour prouver que $X\to Y$
est fermé, il suffit de vérifier qu'il en est ainsi du morphisme structural
$\mathbf{P}_Y^n\to Y$, et le critère c) de
(\hyperref[II.7.3.8-fr]{7.3.8}), joint à
(\hyperref[II.4.1.3.1-fr]{4.1.3.1}), prouve qu'on est ramené à démontrer le
fait suivant~: \emph{si $Y$ est le spectre d'un anneau de valuation $A$, de
corps des fractions $K$, tout point de $\mathbf{P}_Y^n$ à valeurs dans $K$
provient (par restriction au point générique de $Y$) d'un point de
$\mathbf{P}_Y^n$ à valeurs dans $A$}. Or, tout $\mathcal{O}_Y$-Module
inversible est trivial (\hyperref[I.2.4.8-fr]{I, 2.4.8})~; donc, il résulte de
(\hyperref[II.4.2.6-fr]{4.2.6}) qu'un point de $\mathbf{P}_Y^n$ à valeurs dans
$K$ s'identifie à une classe d'éléments $(\zeta c_0,\zeta c_1,\ldots,\zeta
c_n)$ de $K$, où $\zeta\neq0$ et les $c_i$ sont des éléments non tous nuls de
$K$. Or, en multipliant les $c_i$ par un élément de $A$ de
\oldpage[II]{146}
valuation convenable, on peut supposer que les $c_i$ appartiennent tous à
$A$, et que l'un au moins est inversible. Mais alors
(\hyperref[II.4.2.6-fr]{4.2.6}), le système $(c_0,\ldots,c_n)$ définit aussi
un point de $\mathbf{P}_Y^n$ à valeurs dans $A$, ce qui démontre notre
assertion.
\item[\rm (iii)] Les critères (\hyperref[II.7.2.3-fr]{7.2.3}) et
(\hyperref[II.7.3.8-fr]{7.3.8}) sont surtout commodes quand on considère la
donnée d'un $Y$-préschéma $X$ comme équivalente à la donnée du foncteur
\[
X(Y')=\operatorname{Hom}_Y(Y',X)
\]
en un $Y$-préschéma $Y'$~; ces critères nous permettront par exemple de
prouver que sous certaines conditions les «~schémas de Picard~» sont propres.
\end{enumerate}
\end{remark}

\begin{corollary}[7.3.10]
\label{II.7.3.10-fr}
Soient $Y$ un schéma intègre (resp. un schéma intègre localement noethérien),
$X$ un schéma intègre, $f:X\to Y$ un morphisme dominant.
\begin{enumerate}
\item[\rm (i)] Si $f$ est quasi-compact et universellement fermé, tout anneau
de valuation dont le corps des fractions est le corps $R(X)$ des fonctions
rationnelles sur $X$, et qui domine un anneau local de $Y$, domine aussi un
anneau local de $X$.
\item[\rm (ii)] Inversement, supposons que $f$ soit de type fini, et que la
propriété énoncée dans (i) soit vérifiée par tout anneau de valuation (resp.
tout anneau de valuation discrète) ayant $R(X)$ pour corps des fractions.
Alors $f$ est propre.
\end{enumerate}
\end{corollary}

Notons d'abord que les hypothèses impliquent dans tous les cas que $f$ est
séparé (\hyperref[I.5.5.9-fr]{I, 5.5.9}).
\begin{enumerate}
\item[\rm (i)] Soient $K=R(Y)$, $L=R(X)$, $y$ un point de $Y$, $A$ un anneau
de valuation ayant $L$ pour corps des fractions et dominant
$\mathcal{O}_y$~; l'injection $\mathcal{O}_y\to A$ définit donc un morphisme
$h$ de $Y'=\operatorname{Spec}(A)$ dans $Y$
(\hyperref[I.2.4.4-fr]{I, 2.4.4}) tel que $h(a)=y$, en désignant par $a$ le
point fermé de $Y'$~; en outre, si $\eta$ est le point générique de $Y$, qui
est aussi celui de $\operatorname{Spec}(\mathcal{O}_y)$, on a $h(b)=\eta$, en
désignant par $b$ le point générique de $Y'$ (puisque $K\subset L$ par
hypothèse). Si $\xi$ est le point générique de $X$, on a $k(\xi)=k(b)=L$ par
hypothèse, d'où un $Y$-morphisme $g:\operatorname{Spec}(L)\to X$ tel que
$g(b)=\xi$~; en vertu de (\hyperref[II.7.3.8-fr]{7.3.8}), $g$ provient d'un
$Y$-morphisme $g':Y'\to X$. Si $x=g'(a)$, il est clair que $A$ domine
$\mathcal{O}_x$.
\item[\rm (ii)] La question étant locale sur $Y$, on peut toujours supposer
que $Y$ est affine (resp. affine et noethérien). Comme $f$ est de type fini,
on peut appliquer dans les deux cas le lemme de Chow
(\hyperref[II.5.6.1-fr]{5.6.1}). Il y a donc un morphisme projectif
$p:P\to Y$, un morphisme d'immersion $j:X'\to P$ et un morphisme projectif,
surjectif et birationnel $g:X'\to X$ ($X'$ étant intègre) tels que le diagramme
\[
\begin{tikzcd}[column sep=large,row sep=large]
P \arrow[d,"p"'] & X' \arrow[l,"j"'] \arrow[d,"g"] \\
Y & X \arrow[l,"f"]
\end{tikzcd}
\]
soit commutatif. Il suffit de prouver que $j$ est une immersion \emph{fermée},
car alors $f\circ g=p\circ j$ sera un morphisme projectif, donc propre, et
comme $g$ est surjectif, $f$ sera aussi propre
(\hyperref[II.5.4.3-fr]{5.4.3}). Soit $Z$ le sous-préschéma fermé réduit de
$P$ ayant $\overline{j(X')}$ comme espace sous-jacent
(\hyperref[I.5.2.1-fr]{I, 5.2.1})~; comme $X'$ est intègre, $j$ se factorise
en $i\circ h$, où $i:Z\to P$ est l'injection canonique, $h:X'\to Z$ une
immersion ouverte dominante (\hyperref[I.5.2.3-fr]{I, 5.2.3}), et $Z$ est
intègre~;

\oldpage[II]{147}
en outre $Z$ est projectif sur $Y$, et on voit qu'on peut se borner au cas où
$P$ est \emph{intègre} et $j$ \emph{dominant et birationnel}, et tout revient à
voir que $j$ est \emph{surjectif}. Or, soit $z\in P$~; $\mathcal{O}_z$ est un
anneau local intègre (resp. intègre et noethérien) dont le corps des fractions
est
\[
L=R(P)=R(X')=R(X).
\]
On peut se borner au cas où $z$ n'est pas le point générique de $P$. Il y a par
suite (\hyperref[II.7.1.2-fr]{7.1.2} et
\hyperref[II.7.1.7-fr]{7.1.7}) un anneau de valuation (resp. un anneau de
valuation discrète) $A$, ayant $L$ pour corps des fractions et dominant
$\mathcal{O}_z$. \emph{A fortiori}, $A$ domine $\mathcal{O}_y$, en posant
$y=p(z)$, et par hypothèse il y a donc un $x\in X$ tel que $A$ domine
$\mathcal{O}_x$. Comme $g$ est propre, la première partie de la démonstration
prouve que $A$ domine aussi un $\mathcal{O}_{x'}$, pour un $x'\in X'$~; il
s'ensuit que $\mathcal{O}_z$ et $\mathcal{O}_{j(x')}=\mathcal{O}_{x'}$ sont
apparentés (\hyperref[I.8.1.4-fr]{I, 8.1.4}), et, comme $P$ est un schéma, cela
entraîne $z=j(x')$ (\hyperref[I.8.2.2-fr]{I, 8.2.2}) et achève la démonstration.
\end{enumerate}

\begin{corollary}[7.3.11]
\label{II.7.3.11-fr}
Soient $X$, $Y$ deux schémas intègres, $f:X\to Y$ un morphisme dominant,
quasi-compact et universellement fermé. Supposons en outre $Y$ affine d'anneau
(intègre) $B$. Alors $\Gamma(X,\mathcal{O}_X)$ est canoniquement isomorphe à
un sous-anneau de la fermeture intégrale de $B$ dans $R(X)$.
\end{corollary}

En effet (\hyperref[I.8.2.1.1-fr]{I, 8.2.1.1}),
$\Gamma(X,\mathcal{O}_X)$ s'identifie à l'intersection des $\mathcal{O}_x$ pour
$x\in X$~; en vertu de (\hyperref[II.7.3.10-fr]{7.3.10}), de
(\hyperref[II.7.1.2-fr]{7.1.2}) et (\hyperref[II.7.1.3-fr]{7.1.3}),
$\Gamma(X,\mathcal{O}_X)$ est par suite contenu dans l'intersection des
anneaux de valuation contenant $B$ et ayant $R(X)$ pour corps des fractions~;
la conclusion résulte alors de (\hyperref[II.7.1.3-fr]{7.1.3}).

\begin{remark}[7.3.12]
\label{II.7.3.12-fr}
Sous les hypothèses de (\hyperref[II.7.3.11-fr]{7.3.11}), et lorsqu'on suppose
que $R(X)$ est une extension de type fini de $R(Y)$, on pourra dans de nombreux
cas conclure que $\Gamma(X,\mathcal{O}_X)$ est un \emph{module de type fini}
sur l'anneau $B=\Gamma(Y,\mathcal{O}_Y)$. Ce sera le cas par exemple lorsque
$B$ est une \emph{algèbre de type fini sur un corps}, car on sait alors que la
fermeture intégrale de $B$ dans une extension de type fini de son corps des
fractions est un $B$-module de type fini ([13], t.~I, p.~267, th.~9)~; la
conclusion résulte alors de (\hyperref[II.7.3.11-fr]{7.3.11}) et du fait que
$B$ est noethérien.

En particulier, \emph{un schéma $X$ propre et affine sur un corps $K$ est
fini}. En effet, en vertu de (\hyperref[II.1.6.4-fr]{1.6.4}),
(\hyperref[II.5.4.6-fr]{5.4.6}) et
(\hyperref[I.6.4.4-fr]{I, 6.4.4, c)}), on peut se borner au cas où $X$ est
\emph{réduit}. En outre, il suffira de prouver que chacun des sous-préschémas
fermés de $X$ ayant pour espace sous-jacent une composante irréductible de $X$
(en nombre fini) est fini sur $K$, si bien que (tenant compte de
(\hyperref[II.5.4.5-fr]{5.4.5})) on est finalement ramené au cas où $X$ est
\emph{intègre}. Mais alors le résultat découle des remarques faites ci-dessus.

Au chapitre III, nous retrouverons cette dernière proposition par d'autres
méthodes et comme conséquence de résultats plus généraux, montrant que si
$f:X\to Y$ est propre et $Y$ localement noethérien, $f_*(\mathcal{F})$ est
cohérent pour tout $\mathcal{O}_X$-Module cohérent $\mathcal{F}$
(\hyperref[III.4.4.2-fr]{III, 4.4.2}).

Notons enfin que le critère (\hyperref[II.7.3.10-fr]{7.3.10}) est pris comme
\emph{définition} des morphismes propres en Géométrie algébrique classique.
Nous ne l'avons mentionné que pour cette raison, le critère
(\hyperref[II.7.3.8-fr]{7.3.8}) semblant plus maniable dans toutes les
applications à notre connaissance.
\end{remark}

\oldpage[II]{148}
\subsection{Courbes algébriques et corps de fonctions de dimension 1}
\label{subsection:II.7.4-fr}

Le but de ce numéro est de montrer comment se formule en langage des schémas
la notion classique de courbe algébrique (telle qu'elle est exposée par
exemple dans le livre de C.~Chevalley [23]). Dans tout le numéro, \emph{on
désigne par $k$ un corps, tous les schémas considérés sont des $k$-schémas de
type fini, et tous les morphismes des $k$-morphismes}.

\begin{proposition}[7.4.1]
\label{II.7.4.1-fr}
Soit $X$ un préschéma de type fini sur $k$ (donc noethérien)~; soient $x_i$
$(1\leq i\leq n)$ les points génériques des composantes irréductibles $X_i$ de
$X$, et soit $K_i=k(x_i)$ $(1\leq i\leq n)$. Les conditions suivantes sont
équivalentes~:
\begin{enumerate}
\item[\rm a)] Chacun des $K_i$ est une extension de $k$ dont le degré de
transcendance est égal à 1.
\item[\rm b)] Pour tout point fermé $x$ de $X$, l'anneau local
$\mathcal{O}_x$ est de dimension 1 (\hyperref[II.7.1.5-fr]{7.1.5}).
\item[\rm c)] Les parties fermées irréductibles de $X$ distinctes des $X_i$
sont les points fermés de $X$.
\end{enumerate}
\end{proposition}

Comme $X$ est quasi-compact, toute partie fermée irréductible $F$ de $X$
contient un point fermé (\hyperref[0.2.1.3-fr]{0, 2.1.3}). En vertu de
(\hyperref[I.2.4.2-fr]{I, 2.4.2}), il y a correspondance biunivoque entre les
idéaux premiers de $\mathcal{O}_x$ et les parties fermées irréductibles de $X$
contenant $x$ (\hyperref[I.1.1.14-fr]{I, 1.1.14})~; l'équivalence de b) et c)
en découle aussitôt. D'autre part, si $\mathfrak{p}_\alpha$
$(1\leq\alpha\leq r)$ sont les idéaux premiers minimaux de l'anneau local
noethérien $\mathcal{O}_x$, les anneaux locaux
$\mathcal{O}_x/\mathfrak{p}_\alpha$ sont intègres, et ont pour corps des
fractions ceux des $K_i$ tels que $x\in X_i$. En outre, on sait ([1],
p.~4-06, th.~2) que la dimension d'une $k$-algèbre locale intègre de type fini
est égale au degré de transcendance sur $k$ de son corps des fractions. Enfin,
la dimension de $\mathcal{O}_x$ est borne supérieure des dimensions des
$\mathcal{O}_x/\mathfrak{p}_\alpha$~; or, la condition a) implique que ces
dimensions sont égales à 1, donc a) implique b)~; inversement, si
$\mathcal{O}_x$ est de dimension 1, aucun des $\mathfrak{p}_\alpha$ ne peut
être égal à l'idéal maximal de $\mathcal{O}_x$, sans quoi $\mathcal{O}_x$
serait de dimension 0~; donc chacun des
$\mathcal{O}_x/\mathfrak{p}_\alpha$ est de dimension 1, ce qui montre que b)
entraîne a).

On notera que sous les conditions de (\hyperref[II.7.4.1-fr]{7.4.1}),
l'ensemble $X$ est \emph{vide ou infini}, comme il résulte aussitôt de
(\hyperref[I.6.4.4-fr]{I, 6.4.4}).

\begin{definition}[7.4.2]
\label{II.7.4.2-fr}
On appelle courbe algébrique sur $k$ un schéma algébrique non vide sur $k$
vérifiant les conditions de (\hyperref[II.7.4.1-fr]{7.4.1}).
\end{definition}

Dans le langage de la dimension, qui sera introduit au chapitre IV, cela
s'exprimera en disant qu'une courbe algébrique sur $k$ est un $k$-schéma
algébrique non vide dont \emph{toutes les composantes irréductibles sont de
dimension 1}.

On notera que si $X$ est une courbe algébrique sur $k$, les sous-préschémas
fermés réduits $X_i$ $(1\leq i\leq n)$ de $X$ ayant pour espaces sous-jacents
les composantes irréductibles de $X$ sont des courbes algébriques sur $k$.

\begin{corollary}[7.4.3]
\label{II.7.4.3-fr}
Soit $X$ une courbe algébrique irréductible. Le seul point non fermé de $X$ est
son point générique. Les parties fermées de $X$ distinctes de $X$ sont les
ensembles finis de points fermés~; ce sont aussi les seules parties de $X$ qui
ne soient pas partout denses.
\end{corollary}

Si un point $x\in X$ n'est pas fermé, son adhérence dans $X$ est une partie
fermée irréductible de $X$, donc nécessairement $X$ tout entier en vertu de
(\hyperref[II.7.4.1-fr]{7.4.1}), et par suite $x$ est le point générique de
$X$. Une partie fermée $F$ de $X$, distincte de $X$, ne peut contenir

\oldpage[II]{149}
le point générique de $X$, donc tous ses points sont fermés (dans $X$ et
\emph{a fortiori} dans $F$)~; en considérant le sous-préschéma fermé réduit de
$X$ ayant $F$ pour espace sous-jacent (\hyperref[I.5.2.1-fr]{I, 5.2.1}), il
résulte donc de (\hyperref[I.6.2.2-fr]{I, 6.2.2}) que $F$ est fini et discret.
L'adhérence dans $X$ d'une partie infinie de $X$ est donc nécessairement égale
à $X$.

Lorsque $X$ est une courbe algébrique quelconque, en appliquant
(\hyperref[II.7.4.3-fr]{7.4.3}) aux composantes irréductibles de $X$, on voit
que les seuls points non fermés de $X$ sont les points génériques de ces
composantes.

\begin{corollary}[7.4.4]
\label{II.7.4.4-fr}
Soient $X$ et $Y$ deux courbes algébriques irréductibles sur $k$, $f:X\to Y$
un $k$-morphisme. Pour que $f$ soit dominant, il faut et il suffit que
$f^{-1}(y)$ soit fini pour tout $y\in Y$.
\end{corollary}

En effet, si $f$ n'est pas dominant, $f(X)$ est nécessairement une partie
\emph{finie} de $Y$ en vertu de (\hyperref[II.7.4.3-fr]{7.4.3}), donc il n'est
pas possible que $f^{-1}(y)$ soit fini pour tout point de $Y$, sinon $X$ serait
fini, ce qui est absurde (\hyperref[II.7.4.1-fr]{7.4.1}). Inversement, si $f$
est dominant, pour tout $y\in Y$ distinct du point générique $\eta$ de $Y$,
$f^{-1}(y)$ est fermé dans $X$ puisque $\{y\}$ est fermé dans $Y$
(\hyperref[II.7.4.3-fr]{7.4.3})~; d'autre part, par hypothèse, $f^{-1}(y)$ ne
contient pas le point générique $\xi$ de $X$, donc est fini en vertu de
(\hyperref[II.7.4.3-fr]{7.4.3}). Enfin, pour voir que lorsque $f$ est dominant,
$f^{-1}(\eta)$ est fini, on note que la fibre $f^{-1}(\eta)$ est un schéma
irréductible de type fini sur $k(\eta)$, et de point générique $\xi$
(\hyperref[I.6.3.9-fr]{I, 6.3.9} et
\hyperref[I.6.4.11-fr]{6.4.11}). Comme $k(\xi)$ et $k(\eta)$ sont des
extensions de type fini de $k$, de même degré de transcendance 1, $k(\xi)$ est
nécessairement une extension de degré fini de $k(\eta)$, donc $\xi$ est fermé
dans $f^{-1}(\eta)$ (\hyperref[I.6.4.2-fr]{I, 6.4.2}), et $f^{-1}(\eta)$ est
par suite réduit au point $\xi$.

Nous verrons au chapitre III qu'un morphisme \emph{propre} $f:X\to Y$ de
préschémas noethériens, tel que $f^{-1}(y)$ soit fini pour tout $y\in Y$, est
nécessairement \emph{fini}~; il s'ensuivra donc de
(\hyperref[II.7.4.4-fr]{7.4.4}) qu'un morphisme dominant propre d'une courbe
algébrique irréductible dans une courbe algébrique est \emph{fini}.

\begin{corollary}[7.4.5]
\label{II.7.4.5-fr}
Soit $X$ une courbe algébrique sur $k$. Pour que $X$ soit régulière, il faut et
il suffit que $X$ soit normale, ou encore que les anneaux locaux de ses points
fermés soient des anneaux de valuation discrète.
\end{corollary}

Cela résulte aussitôt de la condition b) de
(\hyperref[II.7.4.1-fr]{7.4.1}) et de
(\hyperref[II.7.1.6-fr]{7.1.6}).

\begin{corollary}[7.4.6]
\label{II.7.4.6-fr}
Soient $X$ une courbe algébrique réduite, $\mathcal{A}$ une
$\mathcal{R}(X)$-Algèbre cohérente réduite~; alors la fermeture intégrale $X'$
de $X$ relativement à $\mathcal{A}$ (\hyperref[II.6.3.4-fr]{6.3.4}) est une
courbe algébrique normale, et le morphisme canonique $X'\to X$ est fini.
\end{corollary}

Le fait que $X'\to X$ soit fini résulte de
(\hyperref[II.6.3.10-fr]{6.3.10})~; $X'$ est donc un $k$-schéma algébrique~;
en outre, si $x_i$ $(1\leq i\leq n)$ sont les points génériques des composantes
irréductibles de $X$, $x'_j$ $(1\leq j\leq m)$ ceux des composantes
irréductibles de $X'$, chacun des $k(x'_j)$ est extension algébrique finie de
l'un des $k(x_i)$ (\hyperref[II.6.3.6-fr]{6.3.6}), donc a un degré de
transcendance 1 sur $k$. $X'$ est donc bien une courbe algébrique sur $k$, et
en outre on sait que $X'$ est somme d'un nombre fini de schémas intègres et
normaux (\hyperref[II.6.3.6-fr]{6.3.6} et
\hyperref[II.6.3.7-fr]{6.3.7}).

\begin{env}[7.4.7]
\label{II.7.4.7-fr}
On dit qu'une courbe algébrique $X$ sur $k$ est \emph{complète} si elle est
\emph{propre} sur $k$.
\end{env}

\begin{corollary}[7.4.8]
\label{II.7.4.8-fr}
Pour qu'une courbe algébrique réduite $X$ sur $k$ soit complète, il faut et il
suffit que sa normalisée $X'$ le soit.
\end{corollary}

\oldpage[II]{150}
En effet, le morphisme canonique $f:X'\to X$ est alors fini
(\hyperref[II.7.4.6-fr]{7.4.6}), donc propre
(\hyperref[II.6.1.11-fr]{6.1.11}) et surjectif
(\hyperref[II.6.3.8-fr]{6.3.8})~; si $g:X\to\operatorname{Spec}(k)$ est le
morphisme structural, $g$ et $g\circ f$ sont donc propres simultanément,
comme il résulte de (\hyperref[II.5.4.2-fr]{5.4.2, (iii)}) et
(\hyperref[II.5.4.3-fr]{5.4.3, (ii)}), $g$ étant séparé par hypothèse.

\begin{proposition}[7.4.9]
\label{II.7.4.9-fr}
Soient $X$ une courbe algébrique normale sur $k$, $Y$ un $k$-schéma algébrique
propre sur $k$. Toute $k$-application rationnelle de $X$ dans $Y$ est alors
partout définie, autrement dit est un morphisme.
\end{proposition}

En effet, il résulte de (\hyperref[II.7.3.7-fr]{7.3.7}) qu'aux points $x\in X$
où une telle application n'est pas définie, la dimension de $\mathcal{O}_x$
devrait être $\geq2$, donc l'ensemble de ces points est vide~; la dernière
assertion provient de (\hyperref[I.7.2.3-fr]{I, 7.2.3}).

\begin{corollary}[7.4.10]
\label{II.7.4.10-fr}
Une courbe algébrique normale sur $k$ est quasi-projective sur $k$.
\end{corollary}

Comme $X$ est somme d'un nombre fini de courbes algébriques intègres et
normales (\hyperref[II.6.3.8-fr]{6.3.8}), on peut se borner au cas où $X$ est
intègre (\hyperref[II.5.3.6-fr]{5.3.6}). Comme $X$ est quasi-compact, il est
recouvert par un nombre fini d'ouverts affines $U_i$ $(1\leq i\leq n)$, et
comme chacun de ces derniers est de type fini sur $k$, il existe un entier
$n_i$ et une $k$-immersion $f_i:U_i\to\mathbf{P}_k^{n_i}$
(\hyperref[II.5.3.3-fr]{5.3.3} et
\hyperref[II.5.3.4-fr]{5.3.4, (i)}). Comme $U_i$ est dense dans $X$, il résulte
de (\hyperref[II.7.4.9-fr]{7.4.9}) que $f_i$ se prolonge en un $k$-morphisme
$g_i:X\to\mathbf{P}_k^{n_i}$, d'où un $k$-morphisme
$g=(g_1,\ldots,g_n)_k$ de $X$ dans le produit $P$ des
$\mathbf{P}_k^{n_i}$ sur $k$. En outre, pour chaque indice $i$, comme la
restriction de $g_i$ à $U_i$ est une immersion, il en est de même de la
restriction de $g$ à $U_i$ (\hyperref[I.5.3.14-fr]{I, 5.3.14}). Comme les
$U_i$ recouvrent $X$ et que $g$ est séparé
(\hyperref[I.5.5.1-fr]{I, 5.5.1, (v)}), $g$ est une immersion de $X$ dans $P$
(\hyperref[I.8.2.8-fr]{I, 8.2.8}). Comme le morphisme de Segre
(\hyperref[II.4.3.3-fr]{4.3.3}) fournit une immersion de $P$ dans un
$\mathbf{P}_k^N$, cela achève de prouver que $X$ est quasi-projectif.

\begin{corollary}[7.4.11]
\label{II.7.4.11-fr}
Une courbe algébrique normale $X$ est isomorphe au schéma induit sur un ouvert
partout dense d'une courbe algébrique normale et complète $\widehat{X}$,
déterminée à un isomorphisme unique près.
\end{corollary}

Si $X_1$, $X_2$ sont deux courbes normales et complètes, il résulte aussitôt
de (\hyperref[II.7.4.9-fr]{7.4.9}) que tout isomorphisme d'un ouvert $U_1$
dense dans $X_1$, sur un ouvert $U_2$ dense dans $X_2$, se prolonge de façon
unique en un isomorphisme de $X_1$ sur $X_2$~; d'où l'assertion d'unicité. Pour
prouver l'existence de $\widehat{X}$, il suffit de remarquer que l'on peut
considérer $X$ comme un sous-schéma d'un fibré projectif $\mathbf{P}_k^n$
(\hyperref[II.7.4.10-fr]{7.4.10}). Soit $\overline{X}$ l'adhérence de $X$ dans
$\mathbf{P}_k^n$ (\hyperref[I.9.5.11-fr]{I, 9.5.11})~; comme $X$ est induit
par $\overline{X}$ sur un ouvert dense dans $\overline{X}$
(\hyperref[I.9.5.10-fr]{I, 9.5.10}), les points génériques $x_i$ des
composantes irréductibles de $X$ sont aussi ceux des composantes irréductibles
de $\overline{X}$, et les $k(x_i)$ sont les mêmes pour ces deux schémas, donc
(\hyperref[II.7.4.1-fr]{7.4.1}) $\overline{X}$ est une courbe algébrique sur
$k$, qui est réduite (\hyperref[I.9.5.9-fr]{I, 9.5.9}) et projective sur $k$
(\hyperref[II.5.5.1-fr]{5.5.1}), donc complète
(\hyperref[II.5.5.3-fr]{5.5.3}). Prenons alors pour $\widehat{X}$ la
\emph{normalisée} de $\overline{X}$, qui est encore complète
(\hyperref[II.7.4.8-fr]{7.4.8})~; en outre, si
$h:\widehat{X}\to\overline{X}$ est le morphisme canonique, la restriction de
$h$ à $h^{-1}(X)$ est un isomorphisme sur $X$ puisque $X$ est normale
(\hyperref[II.6.3.4-fr]{6.3.4}), et comme $h^{-1}(X)$ contient les points
génériques des composantes irréductibles de $\widehat{X}$
(\hyperref[II.6.3.8-fr]{6.3.8}), il est dense dans $\widehat{X}$, ce qui achève
la démonstration.

\oldpage[II]{151}

\begin{remark}[7.4.12]
\label{II.7.4.12-fr}
Nous montrerons au chapitre V que la conclusion de
(\hyperref[II.7.4.10-fr]{7.4.10}) est encore valable sans supposer la courbe
normale (ni même réduite)~; nous montrerons aussi que pour qu'une courbe
algébrique (réduite ou non) soit affine, il faut et il suffit que ses
composantes irréductibles (réduites) ne soient pas complètes.
\end{remark}

\begin{corollary}[7.4.13]
\label{II.7.4.13-fr}
Soient $X$ une courbe normale irréductible de corps $R(X)=K$, $Y$ une courbe
intègre complète de corps $L=R(Y)$. Il y a une correspondance biunivoque
canonique entre les $k$-morphismes dominants $X\to Y$ et les
$k$-monomorphismes $L\to K$.
\end{corollary}

En vertu de (\hyperref[II.7.4.9-fr]{7.4.9}), les $k$-applications rationnelles
de $X$ dans $Y$ s'identifient aux $k$-morphismes $u:X\to Y$. Les morphismes
$u$ dominants étant caractérisés par le fait que $u(x)=y$ (en désignant par
$x$ et $y$ les points génériques respectifs de $X$ et de $Y$), le corollaire
résulte de ces remarques et de (\hyperref[I.7.1.13-fr]{I, 7.1.13}).

\begin{env}[7.4.14]
\label{II.7.4.14-fr}
On peut préciser le résultat de (\hyperref[II.7.4.13-fr]{7.4.13}) lorsqu'on
prend pour $Y$ la \emph{droite projective}
$\mathbf{P}_k^1=\operatorname{Proj}(k[T_0,T])$, $T_0$ et $T$ étant deux
indéterminées. C'est bien là un schéma intègre
(\hyperref[II.2.4.4-fr]{2.4.4}), et le schéma induit sur l'ouvert $D_+(T_0)$
de $Y$ est isomorphe à $\operatorname{Spec}(k[T])$
(\hyperref[II.2.3.6-fr]{2.3.6}), donc le point générique de $Y$ est l'idéal
$(0)$ de $k[T]$ et son corps de fonctions rationnelles $k(T)$, ce qui montre
que $Y$ est une courbe algébrique complète sur $k$. En outre, le seul idéal
premier gradué de $S=k[T_0,T]$ contenant $T_0$ et distinct de $S_+$ est
l'idéal principal $(T_0)$, donc le complémentaire de $D_+(T_0)$ dans
$Y=\mathbf{P}_k^1$ est réduit à \emph{un point fermé}, dit «~point à
l'infini~» que nous noterons $\infty$ (pour une étude générale des rapports
entre fibrés vectoriels et fibrés projectifs, voir
\hyperref[subsection:II.8.4-fr]{8.4}). Avec ces notations~:
\end{env}

\begin{corollary}[7.4.15]
\label{II.7.4.15-fr}
Soit $X$ une courbe normale irréductible de corps $R(X)=K$. Il existe une
application bijective canonique de $K$ sur l'ensemble des morphismes $u$ de
$X$ dans $\mathbf{P}_k^1$, distincts du morphisme constant de valeur
$\infty$. Pour que $u$ soit dominant, il faut et il suffit que l'élément
correspondant de $K$ soit transcendant sur $k$.
\end{corollary}

Cet énoncé résulte immédiatement de (\hyperref[II.7.4.9-fr]{7.4.9}) et du

\begin{lemma}[7.4.15.1]
\label{II.7.4.15.1-fr}
Soit $X$ un préschéma intègre sur $k$, et soit $K=R(X)$ son corps des fonctions
rationnelles. Il existe une application bijective canonique de l'ensemble
$K$ sur l'ensemble des applications rationnelles $u$ de $X$ dans
$\mathbf{P}_k^1$, distinctes du morphisme constant de valeur $\infty$. Pour
qu'une telle application rationnelle soit dominante, il faut et il suffit que
l'élément correspondant de $K$ soit transcendant sur $k$.
\end{lemma}

Tout d'abord, les applications rationnelles de $X$ dans $\mathbf{P}_k^1$
correspondent biunivoquement aux points de $\mathbf{P}_k^1$ à valeurs dans
l'extension $K$ de $k$ (\hyperref[I.7.1.12-fr]{I, 7.1.12}). Si un tel point
est localisé (\hyperref[I.3.4.5-fr]{I, 3.4.5}) au point générique de
$\mathbf{P}_k^1$, l'application rationnelle correspondante est évidemment
dominante. Dans le cas contraire, comme tout point de $\mathbf{P}_k^1$
distinct du point générique est fermé (\hyperref[II.7.4.3-fr]{7.4.3}),
l'image du domaine de définition $U$ de $u$ par l'unique morphisme
$U\to\mathbf{P}_k^1$ de la classe $u$ (\hyperref[I.7.2.2-fr]{I, 7.2.2}),
est réduite à un point fermé $y$ de $\mathbf{P}_k^1$, et ce morphisme (qui
n'est pas nécessairement partout défini dans $X$) n'est donc pas dominant~;
par abus de langage, on dit alors que l'application rationnelle $u$ est
«~constante, de valeur $y$~». Il reste à mettre en correspondance biunivoque
les points de $\mathbf{P}_k^1$ à valeurs dans $K$, de localité
(\hyperref[I.3.4.5-fr]{I, 3.4.5}) distincte de $\infty$, et les éléments de
$K$, et à vérifier que la localité d'un tel point est le point générique de
$\mathbf{P}_k^1$ si et seulement s'il correspond à un

\oldpage[II]{152}
élément transcendant sur $k$. Or, cette vérification est immédiate sur
(\hyperref[II.4.2.6-fr]{4.2.6, exemple 1\textsuperscript{o}}).

\begin{corollary}[7.4.16]
\label{II.7.4.16-fr}
Soient $X$, $Y$ deux courbes algébriques sur $k$, normales, complètes et
irréductibles~; soient $K=R(X)$, $L=R(Y)$ leurs corps. Il existe une
correspondance biunivoque canonique entre l'ensemble des $k$-isomorphismes
$X\xrightarrow{\sim}Y$ et l'ensemble des $k$-isomorphismes
$L\xrightarrow{\sim}K$.
\end{corollary}

C'est une conséquence évidente de (\hyperref[II.7.4.13-fr]{7.4.13}).

\begin{env}[7.4.17]
\label{II.7.4.17-fr}
Le corollaire (\hyperref[II.7.4.16-fr]{7.4.16}) montre qu'une courbe algébrique
sur $k$, normale, complète et irréductible, est \emph{déterminée par son corps
de fonctions rationnelles $K$ à un isomorphisme unique près}~; par définition,
$K$ est une extension de type fini de $k$, de degré de transcendance $1$ (ce
qu'on appelle classiquement un \emph{corps de fonctions algébriques d'une
variable}). De plus~:
\end{env}

\begin{proposition}[7.4.18]
\label{II.7.4.18-fr}
Pour toute extension $K$ de $k$, de type fini et de degré de transcendance
$1$, il existe une courbe algébrique normale, complète et irréductible $X$
telle que $R(X)=K$ (déterminée à isomorphisme unique près). L'ensemble des
anneaux locaux de $X$ s'identifie (\hyperref[I.8.2.1-fr]{I, 8.2.1}) à
l'ensemble formé de $K$ et des anneaux de valuation contenant $k$ et ayant
$K$ comme corps des fractions.
\end{proposition}

En effet, $K$ est une extension de degré fini d'une extension transcendante
pure $k(T)$ de $k$, qui s'identifie, comme on l'a vu, au corps des fonctions
rationnelles de la droite projective $Y=\mathbf{P}_k^1$. Soit $X$ la
\emph{fermeture intégrale} de $Y$ relativement à $K$
(\hyperref[II.6.3.4-fr]{6.3.4})~; $X$ est une courbe algébrique normale de
corps $K$ (\hyperref[II.6.3.7-fr]{6.3.7}), et elle est complète puisque le
morphisme $X\to Y$ est fini (\hyperref[II.7.4.6-fr]{7.4.6}). Les anneaux
locaux $\mathcal{O}_x$ de $X$ sont~: le corps $K$ lorsque $x$ est le point
générique~; si $x$ est distinct du point générique, $\mathcal{O}_x$ est un
anneau de valuation discrète contenant $k$ et ayant $K$ comme corps des
fractions (\hyperref[II.7.4.5-fr]{7.4.5}). Inversement, soit $A$ un tel
anneau~; comme le morphisme $X\to\operatorname{Spec}(k)$ est propre et que
$A$ domine $k$, il domine un anneau local $\mathcal{O}_x$ de $X$
(\hyperref[II.7.3.10-fr]{7.3.10})~; ce dernier étant un anneau de valuation
ayant $K$ comme corps des fractions, est donc nécessairement égal à $A$.

\begin{namedtheorem}[7.4.19]{Remarques}
\label{II.7.4.19-fr}
Il résulte de (\hyperref[II.7.4.16-fr]{7.4.16}) et
(\hyperref[II.7.4.18-fr]{7.4.18}) que \emph{la donnée d'une courbe algébrique
sur $k$, normale, complète et irréductible, est essentiellement équivalente à
la donnée d'une extension $K$ de $k$, de type fini et de degré de
transcendance $1$}. On notera que si $k'$ est une extension du corps de base
$k$, $X\otimes_k k'$ sera encore une courbe algébrique complète sur $k'$
(\hyperref[II.5.4.2-fr]{5.4.2, (iii)}), mais en général elle ne sera ni réduite
ni irréductible. Elle le sera cependant si $K$ est une extension séparable de
$k$ et si $k$ est algébriquement fermé dans $K$ (ce qui s'exprime, dans une
terminologie classique que nous ne suivrons pas, en disant que $K$ est une
«~extension régulière~» de $k$). Mais même dans ce cas, il peut se faire que
$X\otimes_k k'$ ne soit pas normale. Le lecteur trouvera des détails sur ces
questions dans le chapitre IV.
\end{namedtheorem}

\section{Schémas éclatés~; cônes projetants~; fermeture projective}
\label{section:II.8-fr}

\subsection{Préschémas éclatés}
\label{subsection:II.8.1-fr}

\begin{env}[8.1.1]
\label{II.8.1.1-fr}
Soient $Y$ un préschéma, et pour tout entier $n\geq0$, soit $\mathcal{I}_n$ un
Idéal quasi-cohérent de $\mathcal{O}_Y$~; on suppose vérifiées les conditions
suivantes~:
\begin{equation}
\label{II.8.1.1.1-fr}
  \mathcal{I}_0=\mathcal{O}_Y,\qquad
  \mathcal{I}_n\subset\mathcal{I}_m\quad\text{pour }m\leq n
\tag{8.1.1.1}
\end{equation}
\begin{equation}
\label{II.8.1.1.2-fr}
  \mathcal{I}_m\mathcal{I}_n\subset\mathcal{I}_{m+n}
  \quad\text{quels que soient }m,n.
\tag{8.1.1.2}
\end{equation}

\oldpage[II]{153}
On notera que ces hypothèses entraînent
\begin{equation}
\label{II.8.1.1.3-fr}
  \mathcal{I}_1^n\subset\mathcal{I}_n.
\tag{8.1.1.3}
\end{equation}

Posons
\begin{equation}
\label{II.8.1.1.4-fr}
  \mathcal{S}=\bigoplus_{n\geq0}\mathcal{I}_n.
\tag{8.1.1.4}
\end{equation}

Il résulte de (\hyperref[II.8.1.1.1-fr]{8.1.1.1}) et
(\hyperref[II.8.1.1.2-fr]{8.1.1.2}) que $\mathcal{S}$ est une
$\mathcal{O}_Y$-Algèbre graduée quasi-cohérente, et définit donc un $Y$-schéma
$X=\operatorname{Proj}(\mathcal{S})$. Si $\mathcal{J}$ est un Idéal
\emph{inversible} de $\mathcal{O}_Y$,
$\mathcal{I}_n\otimes_{\mathcal{O}_Y}\mathcal{J}^{\otimes n}$ s'identifie
canoniquement à $\mathcal{I}_n\mathcal{J}^n$. Si on remplace donc les
$\mathcal{I}_n$ par les $\mathcal{I}_n\mathcal{J}^n$, ce qui remplace
$\mathcal{S}$ par une $\mathcal{O}_Y$-Algèbre quasi-cohérente
$\mathcal{S}_{(\mathcal{J})}$,
$X_{(\mathcal{J})}=\operatorname{Proj}(\mathcal{S}_{(\mathcal{J})})$ est
canoniquement isomorphe à $X$ (\hyperref[II.3.1.8-fr]{3.1.8}).
\end{env}

\begin{env}[8.1.2]
\label{II.8.1.2-fr}
Supposons $Y$ \emph{localement intègre}, de sorte que le faisceau
$\mathcal{R}(Y)$ des fonctions rationnelles est une $\mathcal{O}_Y$-Algèbre
quasi-cohérente (\hyperref[I.7.3.7-fr]{I, 7.3.7}). Nous dirons qu'un
sous-$\mathcal{O}_Y$-Module $\mathcal{I}$ de $\mathcal{R}(Y)$ est un
\emph{Idéal fractionnaire} de $\mathcal{R}(Y)$ s'il est de \emph{type fini}
(\hyperref[0.5.2.1-fr]{0, 5.2.1}). Supposons donné, pour tout $n\geq0$, un
Idéal fractionnaire quasi-cohérent $\mathcal{I}_n$ de $\mathcal{R}(Y)$, tel
que $\mathcal{I}_0=\mathcal{O}_Y$, et que la condition
(\hyperref[II.8.1.1.2-fr]{8.1.1.2}) soit vérifiée (mais non nécessairement la
seconde condition (\hyperref[II.8.1.1.1-fr]{8.1.1.1}))~; on peut encore alors
définir par la formule (\hyperref[II.8.1.1.4-fr]{8.1.1.4}) une
$\mathcal{O}_Y$-Algèbre graduée quasi-cohérente, et le $Y$-schéma correspondant
$X=\operatorname{Proj}(\mathcal{S})$~; on aura encore un isomorphisme
canonique de $X$ sur $X_{(\mathcal{J})}$ pour tout Idéal fractionnaire
inversible $\mathcal{J}$ de $\mathcal{R}(Y)$.
\end{env}

\begin{definition}[8.1.3]
\label{II.8.1.3-fr}
Soit $Y$ un préschéma (resp. un préschéma localement intègre), et soit
$\mathcal{I}$ un Idéal quasi-cohérent de $\mathcal{O}_Y$ (resp. un Idéal
fractionnaire quasi-cohérent de $\mathcal{R}(Y)$). On dit que le $Y$-schéma
$X=\operatorname{Proj}(\bigoplus_{n\geq0}\mathcal{I}^n)$ est obtenu en
\emph{faisant éclater l'Idéal $\mathcal{I}$}, ou est le \emph{préschéma éclaté
de $Y$ relativement à $\mathcal{I}$}. Lorsque $\mathcal{I}$ est un Idéal
quasi-cohérent de $\mathcal{O}_Y$ et $Y'$ le sous-préschéma fermé de $Y$ défini
par $\mathcal{I}$, on dit aussi que $X$ est le $Y$-schéma obtenu en faisant
éclater $Y'$.
\end{definition}

Par définition, $\mathcal{S}=\bigoplus_{n\geq0}\mathcal{I}^n$ est alors
engendrée par $\mathcal{S}_1=\mathcal{I}$~; si $\mathcal{I}$ est un
$\mathcal{O}_Y$-Module de \emph{type fini}, $X$ est donc \emph{projectif} sur
$Y$ (\hyperref[II.5.5.2-fr]{5.5.2}). Sans hypothèse sur $\mathcal{I}$, le
$\mathcal{O}_X$-Module $\mathcal{O}_X(1)$ est \emph{inversible}
(\hyperref[II.3.2.5-fr]{3.2.5}) et \emph{très ample} en vertu de
(\hyperref[II.4.4.3-fr]{4.4.3}) pour le morphisme structural $X\to Y$.

On notera que si $f:X\to Y$ est le morphisme structural, la restriction de
$f$ à $f^{-1}(Y-Y')$ est un \emph{isomorphisme} sur $Y-Y'$ lorsque
$\mathcal{I}$ est un Idéal de $\mathcal{O}_Y$ et $Y'$ le sous-préschéma fermé
qu'il définit~: en effet, la question étant locale sur $Y$, il suffit de
supposer $\mathcal{I}=\mathcal{O}_Y$, et notre assertion résulte de
(\hyperref[II.3.1.7-fr]{3.1.7}).

Lorsqu'on remplace $\mathcal{I}$ par $\mathcal{I}^d$ ($d>0$), le $Y$-schéma
éclaté $X$ est remplacé par un $Y$-schéma canoniquement isomorphe $X'$
(\hyperref[II.8.1.1-fr]{8.1.1})~; de même, pour tout Idéal (resp. Idéal
fractionnaire) inversible $\mathcal{J}$, le préschéma éclaté
$X_{(\mathcal{J})}$ relativement à l'Idéal $\mathcal{I}\mathcal{J}$ est
canoniquement isomorphe à $X$ (\hyperref[II.8.1.1-fr]{8.1.1}).

En particulier, lorsque $\mathcal{I}$ est un Idéal (resp. un Idéal
fractionnaire) \emph{inversible}, le $Y$-schéma obtenu en faisant éclater
$\mathcal{I}$ est \emph{isomorphe} à $Y$ (\hyperref[II.3.1.7-fr]{3.1.7}).

\begin{proposition}[8.1.4]
\label{II.8.1.4-fr}
Soit $Y$ un préschéma intègre.
\begin{enumerate}
\item[\rm(i)] Pour toute suite $(\mathcal{I}_n)$ d'Idéaux fractionnaires
quasi-cohérents de $\mathcal{R}(Y)$ vérifiant
(\hyperref[II.8.1.1.2-fr]{8.1.1.2})

\oldpage[II]{154}
et tels que $\mathcal{I}_0=\mathcal{O}_Y$, le $Y$-schéma
$X=\operatorname{Proj}(\bigoplus_{n\geq0}\mathcal{I}_n)$ est intègre et le
morphisme structural $f:X\to Y$ dominant.
\item[\rm(ii)] Soit $\mathcal{I}$ un Idéal fractionnaire quasi-cohérent de
$\mathcal{R}(Y)$, et soit $X$ le $Y$-schéma éclaté de $Y$ relativement à
$\mathcal{I}$. Si $\mathcal{I}\neq0$, le morphisme structural $f:X\to Y$ est
alors birationnel et surjectif.
\end{enumerate}
\end{proposition}

\begin{enumerate}
\item[\rm(i)] résulte de ce que $\mathcal{S}=\bigoplus_{n\geq0}\mathcal{I}_n$
est une $\mathcal{O}_Y$-Algèbre \emph{intègre}
(\hyperref[II.3.1.12-fr]{3.1.12} et \hyperref[II.3.1.14-fr]{3.1.14}) puisque
pour tout $y\in Y$, $\mathcal{O}_y$ est un anneau intègre
(\hyperref[I.5.1.4-fr]{I, 5.1.4}).
\item[\rm(ii)] D'après (i), $X$ est intègre~; si en outre, $x$ et $y$ sont les
points génériques de $X$ et $Y$, on a $f(x)=y$, et il faut prouver que $k(x)$
est de rang $1$ sur $k(y)$. Or $x$ est aussi le point générique de la fibre
$f^{-1}(y)$~; si $\psi$ est le morphisme canonique $Z\to Y$, où
$Z=\operatorname{Spec}(k(y))$, le préschéma $f^{-1}(y)$ s'identifie à
$\operatorname{Proj}(\mathcal{S}')$, où
$\mathcal{S}'=\psi^*(\mathcal{S})$ (\hyperref[II.3.5.3-fr]{3.5.3}). Mais il est
clair que $\mathcal{S}'=\bigoplus_{n\geq0}(\mathcal{I}_y)^n$, et comme
$\mathcal{I}$ est un Idéal fractionnaire quasi-cohérent non nul de
$\mathcal{R}(Y)$, $\mathcal{I}_y\neq0$
(\hyperref[I.7.3.6-fr]{I, 7.3.6}), d'où $\mathcal{I}_y=k(y)$~;
$\operatorname{Proj}(\mathcal{S}')$ s'identifie donc à
$\operatorname{Spec}(k(y))$ (\hyperref[II.3.1.7-fr]{3.1.7}), d'où la
conclusion.
\end{enumerate}

Nous démontrerons une \emph{réciproque} de
(\hyperref[II.8.1.4-fr]{8.1.4}) dans
(\hyperref[III.2.3.8-fr]{III, 2.3.8}).

\begin{env}[8.1.5]
\label{II.8.1.5-fr}
Revenons à la situation et aux notations de
(\hyperref[II.8.1.1-fr]{8.1.1}). Par définition, les homomorphismes
d'injection $\mathcal{I}_{n+1}\to\mathcal{I}_n$
(\hyperref[II.8.1.1.1-fr]{8.1.1.1}) définissent pour chaque $k\in\mathbf{Z}$ un
homomorphisme injectif de degré $0$ de $\mathcal{S}$-Modules gradués
\begin{equation}
\label{II.8.1.5.1-fr}
  u_k:\mathcal{S}_+(k+1)\to\mathcal{S}(k)~;
\tag{8.1.5.1}
\end{equation}
comme $\mathcal{S}_+(k+1)$ et $\mathcal{S}(k+1)$ sont canoniquement
(TN)-isomorphes, il correspond canoniquement à $u_k$ un homomorphisme injectif
de $\mathcal{O}_X$-Modules (\hyperref[II.3.4.2-fr]{3.4.2})~:
\begin{equation}
\label{II.8.1.5.2-fr}
  \widetilde{u}_k:\mathcal{O}_X(k+1)\to\mathcal{O}_X(k).
\tag{8.1.5.2}
\end{equation}

Rappelons d'autre part (\hyperref[II.3.2.6-fr]{3.2.6}) que l'on a défini des
homomorphismes canoniques
\begin{equation}
\label{II.8.1.5.3-fr}
  \lambda:\mathcal{O}_X(h)\otimes_{\mathcal{O}_X}\mathcal{O}_X(k)
  \to\mathcal{O}_X(h+k)
\tag{8.1.5.3}
\end{equation}
et comme le diagramme
\[
\xymatrix{
  \mathcal{S}(h)\otimes_{\mathcal{S}}\mathcal{S}(k)
    \otimes_{\mathcal{S}}\mathcal{S}(l)\ar[r]\ar[d]
  & \mathcal{S}(h+k)\otimes_{\mathcal{S}}\mathcal{S}(l)\ar[d]\\
  \mathcal{S}(h)\otimes_{\mathcal{S}}\mathcal{S}(k+l)\ar[r]
  & \mathcal{S}(h+k+l)
}
\]
est commutatif, il résulte de la fonctorialité des $\lambda$
(\hyperref[II.3.2.6-fr]{3.2.6}) que les homomorphismes
(\hyperref[II.8.1.5.3-fr]{8.1.5.3}) définissent sur
\begin{equation}
\label{II.8.1.5.4-fr}
  \mathcal{S}_X=\bigoplus_{n\in\mathbf{Z}}\mathcal{O}_X(n)
\tag{8.1.5.4}
\end{equation}

\oldpage[II]{155}
une structure de $\mathcal{O}_X$-Algèbre graduée quasi-cohérente. En outre, le
diagramme
\[
\xymatrix{
  \mathcal{S}(h)\otimes_{\mathcal{S}}\mathcal{S}_+(k+1)\ar[r]\ar[d]_{1\otimes u_k}
  & \mathcal{S}_+(h+k+1)\ar[d]^{u_{k+h}}\\
  \mathcal{S}(h)\otimes_{\mathcal{S}}\mathcal{S}(k)\ar[r]
  & \mathcal{S}(h+k)
}
\]
est commutatif~; la fonctorialité des $\lambda$ montre alors que l'on a un
diagramme commutatif
\begin{equation}
\label{II.8.1.5.5-fr}
\xymatrix{
  \mathcal{O}_X(h)\otimes_{\mathcal{O}_X}\mathcal{O}_X(k+1)
    \ar[r]^{\lambda}\ar[d]_{1\otimes\widetilde{u}_k}
  & \mathcal{O}_X(h+k+1)\ar[d]^{\widetilde{u}_{k+h}}\\
  \mathcal{O}_X(h)\otimes_{\mathcal{O}_X}\mathcal{O}_X(k)
    \ar[r]^{\lambda}
  & \mathcal{O}_X(h+k)
}
\tag{8.1.5.5}
\end{equation}
les flèches horizontales étant les homomorphismes canoniques. On peut donc
dire que les $\widetilde{u}_k$ définissent un \emph{homomorphisme injectif} (de
degré $0$) \emph{de $\mathcal{S}_X$-Modules gradués}
\begin{equation}
\label{II.8.1.5.6-fr}
  \widetilde{u}:\mathcal{S}_X(1)\to\mathcal{S}_X.
\tag{8.1.5.6}
\end{equation}
\end{env}

\begin{env}[8.1.6]
\label{II.8.1.6-fr}
Gardant les notations de (\hyperref[II.8.1.5-fr]{8.1.5}), remarquons maintenant
que, pour $n\geq0$, l'homomorphisme composé
$\widetilde{v}_n=\widetilde{u}_{n-1}\widetilde{u}_{n-2}\cdots\widetilde{u}_0$
est un homomorphisme \emph{injectif} $\mathcal{O}_X(n)\to\mathcal{O}_X$~;
nous désignerons par $\mathcal{I}_{n,X}$ son image, qui est donc un Idéal
quasi-cohérent de $\mathcal{O}_X$, \emph{isomorphe} à $\mathcal{O}_X(n)$. En
outre, le diagramme
\[
\xymatrix{
  \mathcal{O}_X(m)\otimes_{\mathcal{O}_X}\mathcal{O}_X(n)
    \ar[r]^{\lambda}\ar[d]_{\widetilde{v}_m\otimes\widetilde{v}_n}
  & \mathcal{O}_X(m+n)\ar[d]^{\widetilde{v}_{m+n}}\\
  \mathcal{O}_X\ar[r]_{\mathrm{id}}
  & \mathcal{O}_X
}
\]
est commutatif pour $m\geq0$, $n\geq0$. On en conclut les inclusions suivantes
\begin{equation}
\label{II.8.1.6.1-fr}
  \mathcal{I}_{0,X}=\mathcal{O}_X,\qquad
  \mathcal{I}_{n,X}\subset\mathcal{I}_{m,X}
  \quad\text{pour }0\leq m\leq n
\tag{8.1.6.1}
\end{equation}
\begin{equation}
\label{II.8.1.6.2-fr}
  \mathcal{I}_{m,X}\mathcal{I}_{n,X}\subset\mathcal{I}_{m+n,X}
  \quad\text{pour }m\geq0,\ n\geq0.
\tag{8.1.6.2}
\end{equation}
\end{env}

\oldpage[II]{156}

\begin{proposition}[8.1.7]
\label{II.8.1.7-fr}
Soient $Y$ un préschéma, $\mathcal{I}$ un Idéal quasi-cohérent de
$\mathcal{O}_Y$, $X=\operatorname{Proj}(\bigoplus_{n\geq0}\mathcal{I}^n)$ le
$Y$-schéma obtenu en faisant éclater $\mathcal{I}$. On a alors, pour tout
$n>0$, un isomorphisme canonique
\begin{equation}
\label{II.8.1.7.1-fr}
  \mathcal{O}_X(n)\xrightarrow{\sim}\mathcal{I}^n\mathcal{O}_X
  =\mathcal{I}_{n,X}
\tag{8.1.7.1}
\end{equation}
(cf. (\hyperref[0.4.3.5-fr]{0, 4.3.5})), et par suite
$\mathcal{I}^n\mathcal{O}_X$ est un $\mathcal{O}_X$-Module inversible très
ample si $n>0$.
\end{proposition}

La dernière assertion est immédiate, puisque $\mathcal{O}_X(1)$ est inversible
(\hyperref[II.3.2.5-fr]{3.2.5}) et très ample pour $Y$ par définition
(\hyperref[II.4.4.3-fr]{4.4.3} et \hyperref[II.4.4.9-fr]{4.4.9}). D'autre
part, par définition, l'image de $v_n$ n'est autre que
$\mathcal{I}^n\mathcal{S}$, et (\hyperref[II.8.1.7.1-fr]{8.1.7.1}) résulte donc
de l'exactitude du foncteur $\widetilde{\mathcal{M}}$
(\hyperref[II.3.2.4-fr]{3.2.4}) et de la formule
(\hyperref[II.3.2.4.1-fr]{3.2.4.1}).

\begin{corollary}[8.1.8]
\label{II.8.1.8-fr}
Sous les hypothèses de (\hyperref[II.8.1.7-fr]{8.1.7}), si $f:X\to Y$ est le
morphisme structural et $Y'$ le sous-préschéma fermé de $Y$ défini par
$\mathcal{I}$, le sous-préschéma fermé $X'=f^{-1}(Y')$ de $X$ est défini par
$\mathcal{I}\mathcal{O}_X$ (isomorphe canoniquement à $\mathcal{O}_X(1)$),
d'où une suite exacte canonique
\begin{equation}
\label{II.8.1.8.1-fr}
  0\to\mathcal{O}_X(1)\to\mathcal{O}_X\to\mathcal{O}_{X'}\to0.
\tag{8.1.8.1}
\end{equation}
\end{corollary}

Cela résulte de (\hyperref[II.8.1.7.1-fr]{8.1.7.1}) et de
(\hyperref[I.4.4.5-fr]{I, 4.4.5}).

\begin{env}[8.1.9]
\label{II.8.1.9-fr}
Sous les hypothèses de (\hyperref[II.8.1.7-fr]{8.1.7}), on peut préciser la
structure des $\mathcal{I}_{n,X}$. Remarquons en effet que l'homomorphisme
\[
  \widetilde{u}_{-1}:\mathcal{O}_X\to\mathcal{O}_X(-1)
\]
correspond canoniquement à une section $s$ de $\mathcal{O}_X(-1)$ au-dessus de
$X$, que nous appellerons la \emph{section canonique} (relative à
$\mathcal{I}$) (\hyperref[0.5.1.1-fr]{0, 5.1.1}). D'autre part, dans le
diagramme (\hyperref[II.8.1.5.5-fr]{8.1.5.5}), les flèches horizontales sont
des isomorphismes (\hyperref[II.3.2.7-fr]{3.2.7})~; en remplaçant dans ce
diagramme $h$ par $k$ et $k$ par $-1$, on obtient
$\widetilde{u}_k=1_k\otimes\widetilde{u}_{-1}$ ($1_h$ désignant l'identité de
$\mathcal{O}_X(h)$), autrement dit l'homomorphisme $\widetilde{u}_k$ n'est
autre que la \emph{multiplication tensorielle par la section canonique $s$}
(pour tout $k\in\mathbf{Z}$). L'homomorphisme $\widetilde{u}$
(\hyperref[II.8.1.5.6-fr]{8.1.5.6}) s'interprète donc aussi de la même manière.

Par suite, pour tout $n\geq0$, l'homomorphisme
$\widetilde{v}_n:\mathcal{O}_X(n)\to\mathcal{O}_X$ n'est autre que la
multiplication tensorielle par $s^{\otimes n}$~; on en déduit~:
\end{env}

\begin{corollary}[8.1.10]
\label{II.8.1.10-fr}
Avec les notations de (\hyperref[II.8.1.8-fr]{8.1.8}), l'espace sous-jacent à
$X'$ est l'ensemble des $x\in X$ tels que $s(x)=0$, $s$ désignant la section
canonique de $\mathcal{O}_X(-1)$.
\end{corollary}

En effet, si $c_x$ est un générateur de la fibre $(\mathcal{O}_X(1))_x$ en un
point $x$, $s_x\otimes c_x$ s'identifie canoniquement à un générateur de la
fibre de $\mathcal{I}_{1,X}$ au point $x$, et est donc inversible si et
seulement si l'on a
$s_x\notin\mathfrak{m}_x(\mathcal{O}_X(-1))_x$, autrement dit si
$s(x)\neq0$.

\begin{proposition}[8.1.11]
\label{II.8.1.11-fr}
Soient $Y$ un préschéma intègre, $\mathcal{I}$ un Idéal fractionnaire
quasi-cohérent de $\mathcal{R}(Y)$, $X$ le $Y$-schéma obtenu en faisant éclater
$\mathcal{I}$. Alors $\mathcal{I}\mathcal{O}_X$ est un
$\mathcal{O}_X$-Module inversible très ample pour $Y$.
\end{proposition}

La question étant locale sur $Y$ (\hyperref[II.4.4.5-fr]{4.4.5}), on peut se
borner au cas où $Y=\operatorname{Spec}(A)$, où $A$ est un anneau intègre de
corps des fractions $K$ et $\mathcal{I}=\widetilde{\mathfrak{I}}$,
$\mathfrak{I}$ étant un idéal fractionnaire de $K$~; il existe par suite un
élément $a\neq0$ de $A$ tel que $a\mathfrak{I}\subset A$. Posons
$S=\bigoplus_{n\geq0}\mathfrak{I}^n$~; l'application $x\mapsto ax$ est un
$A$-isomorphisme de $\mathfrak{I}^{n+1}=(S(1))_n$ sur
$a\mathfrak{I}^{n+1}=a\mathfrak{I}S_n\subset\mathfrak{I}^n=S_n$,

\oldpage[II]{157}
donc définit un (TN)-isomorphisme de degré $0$ de $S$-modules gradués
$S_+(1)\to a\mathfrak{I}S$. Par ailleurs $x\mapsto a^{-1}x$ est un
isomorphisme de degré $0$ de $S$-modules gradués
$a\mathfrak{I}S\xrightarrow{\sim}\mathfrak{I}S$. On obtient donc ainsi par
composition (\hyperref[II.3.2.4-fr]{3.2.4}) un isomorphisme de
$\mathcal{O}_X$-Modules
$\mathcal{O}_X(1)\xrightarrow{\sim}\mathcal{I}\mathcal{O}_X$, et comme $S$
est engendrée par $S_1=\mathfrak{I}$, $\mathcal{O}_X(1)$ est inversible
(\hyperref[II.3.2.5-fr]{3.2.5}) et très ample
(\hyperref[II.4.4.3-fr]{4.4.3} et \hyperref[II.4.4.9-fr]{4.4.9}), d'où notre
assertion.

\subsection{Résultats préliminaires sur la localisation dans les anneaux gradués}
\label{subsection:II.8.2-fr}

\begin{env}[8.2.1]
\label{II.8.2.1-fr}
Soit $S$ un anneau gradué, où nous ne supposons pas pour le moment les degrés
positifs. On posera
\begin{equation}
\label{II.8.2.1.1-fr}
  S^{\geq}=\bigoplus_{n\geq0}S_n,
  \qquad
  S^{\leq}=\bigoplus_{n\leq0}S_n,
\tag{8.2.1.1}
\end{equation}
qui sont des sous-anneaux gradués de $S$, à degrés respectivement tous
positifs et tous négatifs. Si $f$ est un élément homogène de degré $d$
(positif ou négatif) de $S$, l'anneau des fractions $S_f=S'$ est encore muni
d'une structure d'anneau gradué, en prenant pour $S'_n$ ($n\in\mathbf{Z}$)
l'ensemble des $x/f^k$, où $x\in S_{n+kd}$ ($k\geq0$)~; on posera encore
$S_{(f)}=S'_0$, et on écrira $S_f^{\geq}$ et $S_f^{\leq}$ pour $S'^{\geq}$ et
$S'^{\leq}$ respectivement. Si $d>0$, on a
\begin{equation}
\label{II.8.2.1.2-fr}
  (S^{\geq})_f=S_f
\tag{8.2.1.2}
\end{equation}
puisque, si $x\in S_{n+kd}$, avec $n+kd<0$, on peut écrire
$x/f^k=xf^h/f^{h+k}$ et que $n+(h+k)d>0$ pour $h$ assez grand et $>0$. On en
conclut par définition que
\begin{equation}
\label{II.8.2.1.3-fr}
  (S^{\geq})_{(f)}=(S_f^{\geq})_0=S_{(f)}.
\tag{8.2.1.3}
\end{equation}

Si $M$ est un $S$-module gradué, on posera de même
\begin{equation}
\label{II.8.2.1.4-fr}
  M^{\geq}=\bigoplus_{n\geq0}M_n,
  \qquad
  M^{\leq}=\bigoplus_{n\leq0}M_n
\tag{8.2.1.4}
\end{equation}
qui sont respectivement un $S^{\geq}$-module gradué et un $S^{\leq}$-module
gradué et ont pour intersection le $S_0$-module $M_0$. Si $f\in S_d$, on
définit encore $M_f$ comme $S_f$-module gradué en prenant comme éléments de
degré $n$ les $z/f^k$, où $z\in M_{n+kd}$ ($k\geq0$)~; on désignera par
$M_{(f)}$ l'ensemble des éléments de degré $0$ de $M_f$, qui est un
$S_{(f)}$-module, et on écrira $M_f^{\geq}$ et $M_f^{\leq}$ au lieu de
$(M_f)^{\geq}$ et $(M_f)^{\leq}$ respectivement. Si $d>0$, on voit comme
ci-dessus que
\begin{equation}
\label{II.8.2.1.5-fr}
  (M^{\geq})_f=M_f
\tag{8.2.1.5}
\end{equation}
et
\begin{equation}
\label{II.8.2.1.6-fr}
  (M^{\geq})_{(f)}=(M_f^{\geq})_0=M_{(f)}.
\tag{8.2.1.6}
\end{equation}
\end{env}

\begin{env}[8.2.2]
\label{II.8.2.2-fr}
Soit $z$ une indéterminée, que nous nommerons \emph{variable
d'homogénéisation}. Si $S$ est un anneau gradué (à degrés positifs ou
négatifs), l'algèbre de polynômes\footnote{Il ne saurait ici y avoir de
confusion avec l'usage de la notation $\widehat{S}$ pour désigner le séparé
complété d'un anneau.}
\begin{equation}
\label{II.8.2.2.1-fr}
  \widehat{S}=S[z]
\tag{8.2.2.1}
\end{equation}

\oldpage[II]{158}
est une $S$-algèbre graduée, quand on prend pour degré de $fz^n$ ($n\geq0$) où
$f$ est homogène,
\begin{equation}
\label{II.8.2.2.2-fr}
  \deg(fz^n)=n+\deg f.
\tag{8.2.2.2}
\end{equation}
\end{env}

\begin{lemma}[8.2.3]
\label{II.8.2.3-fr}
\begin{enumerate}
\item[\rm(i)] On a des isomorphismes canoniques d'anneaux (non gradués)
\begin{equation}
\label{II.8.2.3.1-fr}
  \widehat{S}_{(z)}\xrightarrow{\sim}
  \widehat{S}/(z-1)\widehat{S}\xrightarrow{\sim}S.
\tag{8.2.3.1}
\end{equation}
\item[\rm(ii)] On a un isomorphisme canonique d'anneaux (non gradués)
\begin{equation}
\label{II.8.2.3.2-fr}
  \widehat{S}_{(f)}\xrightarrow{\sim}S_f^{\leq}
\tag{8.2.3.2}
\end{equation}
pour tout $f\in S_d$, avec $d>0$.
\end{enumerate}
\end{lemma}

Le premier des isomorphismes de (\hyperref[II.8.2.3.1-fr]{8.2.3.1}) a été
défini dans (\hyperref[II.2.2.5-fr]{2.2.5}) et le second est trivial~;
l'isomorphisme $\widehat{S}_{(z)}\xrightarrow{\sim}S$ ainsi défini fait donc
correspondre à $xz^n/z^{n+k}$, où $\deg(x)=k$ ($k\geq-n$) l'élément $x$.
L'homomorphisme (\hyperref[II.8.2.3.2-fr]{8.2.3.2}) fait correspondre à
$xz^n/f^k$, où $\deg(x)=kd-n$, l'élément $x/f^k$, de degré $-n$ dans
$S_f^{\leq}$, et il est encore évident qu'on a bien ici un isomorphisme.

\begin{env}[8.2.4]
\label{II.8.2.4-fr}
Soit $M$ un $S$-module gradué. Il est clair que le $S$-module
\begin{equation}
\label{II.8.2.4.1-fr}
  \widehat{M}=M\otimes_S\widehat{S}=M\otimes_S S[z]
\tag{8.2.4.1}
\end{equation}
est somme directe des $S$-modules $M\otimes_S Sz^n$, donc des groupes abéliens
$M_k\otimes_S Sz^n$ ($k\in\mathbf{Z}$, $n\geq0$)~; on définit sur
$\widehat{M}$ une structure de $\widehat{S}$-module gradué en prenant
\begin{equation}
\label{II.8.2.4.2-fr}
  \deg(x\otimes z^n)=n+\deg x
\tag{8.2.4.2}
\end{equation}
pour tout $x$ homogène dans $M$. Nous laissons au lecteur le soin de démontrer
l'analogue de (\hyperref[II.8.2.3-fr]{8.2.3})~:
\end{env}

\begin{lemma}[8.2.5]
\label{II.8.2.5-fr}
\begin{enumerate}
\item[\rm(i)] On a un di-isomorphisme canonique de modules (non gradués)
\begin{equation}
\label{II.8.2.5.1-fr}
  \widehat{M}_{(z)}\xrightarrow{\sim}M.
\tag{8.2.5.1}
\end{equation}
\item[\rm(ii)] Pour tout $f\in S_d$ ($d>0$), on a un di-isomorphisme de
modules (non gradués)
\begin{equation}
\label{II.8.2.5.2-fr}
  \widehat{M}_{(f)}\xrightarrow{\sim}M_f^{\leq}.
\tag{8.2.5.2}
\end{equation}
\end{enumerate}
\end{lemma}

\begin{env}[8.2.6]
\label{II.8.2.6-fr}
Soit $S$ un anneau gradué à degrés \emph{positifs}, et considérons dans $S$
la suite décroissante d'idéaux gradués
\begin{equation}
\label{II.8.2.6.1-fr}
  S_{[n]}=\bigoplus_{m\geq n}S_m\qquad(n\geq0)
\tag{8.2.6.1}
\end{equation}
(en particulier on a $S_{[0]}=S$, $S_{[1]}=S_+$). Comme il est clair que
$S_{[m]}S_{[n]}\subset S_{[m+n]}$, on peut définir un \emph{anneau gradué}
$S^\natural$ en prenant
\begin{equation}
\label{II.8.2.6.2-fr}
  S^\natural=\bigoplus_{n\geq0}S_n^\natural
  \quad\text{avec}\quad S_n^\natural=S_{[n]}.
\tag{8.2.6.2}
\end{equation}
$S_0^\natural$ est donc l'anneau $S$ considéré comme anneau \emph{non gradué},
et $S^\natural$ est par suite une $S_0^\natural$-algèbre. Pour tout élément
homogène $f\in S_d$ ($d>0$), nous désignerons par $f^\natural$ l'élément
$f$ considéré comme appartenant à $S_{[d]}=S_d^\natural$. Avec ces
notations~:
\end{env}

\oldpage[II]{159}

\begin{lemma}[8.2.7]
\label{II.8.2.7-fr}
Soient $S$ un anneau gradué à degrés positifs, $f$ un élément homogène de
$S_d$ ($d>0$). On a des isomorphismes canoniques d'anneaux
\begin{equation}
\label{II.8.2.7.1-fr}
  S_f\xrightarrow{\sim}\bigoplus_{n\in\mathbf{Z}}S(n)_{(f)}
\tag{8.2.7.1}
\end{equation}
\begin{equation}
\label{II.8.2.7.2-fr}
  (S_f^{\geq})_{f/1}\xrightarrow{\sim}S_f
\tag{8.2.7.2}
\end{equation}
\begin{equation}
\label{II.8.2.7.3-fr}
  S^\natural_{(f^\natural)}\xrightarrow{\sim}S_f^{\geq}
\tag{8.2.7.3}
\end{equation}
dont les deux premiers sont des isomorphismes d'anneaux gradués.
\end{lemma}

Il est immédiat par définition que l'on a
$(S_f)_n=(S(n)_f)_0$, d'où l'isomorphisme
(\hyperref[II.8.2.7.1-fr]{8.2.7.1}), qui n'est autre que l'identité. D'autre
part, comme $f/1$ est inversible dans $S_f$, il y a un isomorphisme canonique
$S_f\xrightarrow{\sim}(S_f^{\geq})_{f/1}=(S_f)_{f/1}$ en vertu de
(\hyperref[II.8.2.1.2-fr]{8.2.1.2}) appliqué à $S_f$~; l'isomorphisme
réciproque est par définition l'isomorphisme
(\hyperref[II.8.2.7.2-fr]{8.2.7.2}). Enfin, si
$x=\sum_{m\geq n}y_m$ est un élément de $S_{[n]}$, avec $n=kd$, on fait
correspondre à l'élément $x/(f^\natural)^k$ l'élément
$\sum_m y_m/f^k$ de $S_f^{\geq}$, et on vérifie aussitôt que cela définit un
isomorphisme (\hyperref[II.8.2.7.3-fr]{8.2.7.3}).

\begin{env}[8.2.8]
\label{II.8.2.8-fr}
Si $M$ est un $S$-module gradué, on pose de même pour tout $n\in\mathbf{Z}$
\begin{equation}
\label{II.8.2.8.1-fr}
  M_{[n]}=\bigoplus_{m\geq n}M_m
\tag{8.2.8.1}
\end{equation}
et comme $S_{[m]}M_{[n]}\subset M_{[m+n]}$ ($m\geq0$), on peut définir un
$S^\natural$-module gradué $M^\natural$ en prenant
\begin{equation}
\label{II.8.2.8.2-fr}
  M^\natural=\bigoplus_{n\in\mathbf{Z}}M_n^\natural,
  \quad\text{avec}\quad M_n^\natural=M_{[n]}.
\tag{8.2.8.2}
\end{equation}
\end{env}

Nous laissons encore au lecteur la démonstration du~:

\begin{lemma}[8.2.9]
\label{II.8.2.9-fr}
Avec les notations de (\hyperref[II.8.2.7-fr]{8.2.7}) et
(\hyperref[II.8.2.8-fr]{8.2.8}), on a les di-isomorphismes canoniques de
modules
\begin{equation}
\label{II.8.2.9.1-fr}
  M_f\xrightarrow{\sim}\bigoplus_{n\in\mathbf{Z}}M(n)_{(f)}
\tag{8.2.9.1}
\end{equation}
\begin{equation}
\label{II.8.2.9.2-fr}
  (M_f^{\geq})_{f/1}\xrightarrow{\sim}M_f
\tag{8.2.9.2}
\end{equation}
\begin{equation}
\label{II.8.2.9.3-fr}
  M^\natural_{(f^\natural)}\xrightarrow{\sim}M_f^{\geq}
\tag{8.2.9.3}
\end{equation}
dont les deux premiers sont des di-isomorphismes de modules gradués.
\end{lemma}

\begin{lemma}[8.2.10]
\label{II.8.2.10-fr}
Soit $S$ un anneau gradué à degrés positifs.
\begin{enumerate}
\item[\rm(i)] Pour que $S^\natural$ soit une $S_0^\natural$-algèbre de type
fini (resp. noethérienne) il faut et il suffit que $S$ soit une $S_0$-algèbre
de type fini (resp. noethérienne).
\item[\rm(ii)] Pour que $S_{n+1}^\natural=S_1^\natural S_n^\natural$ pour
$n\geq n_0$, il faut et il suffit que $S_{n+1}=S_1S_n$ pour $n\geq n_0$.
\item[\rm(iii)] Pour que $S_n^\natural=S_1^{\natural n}$ pour $n\geq n_0$, il
faut et il suffit que $S_n=S_1^n$ pour $n\geq n_0$.
\item[\rm(iv)] Si $(f_\alpha)$ est un ensemble d'éléments homogènes de $S_+$
telle que $S_+$ soit la racine dans $S_+$ de l'idéal de $S_+$ engendré par
les $f_\alpha$, alors $S_+^\natural$ est la racine dans $S_+^\natural$ de
l'idéal de $S_+^\natural$ engendré par les $f_\alpha^\natural$.
\end{enumerate}
\end{lemma}

\begin{enumerate}
\item[\rm(i)] Si $S^\natural$ est une $S_0^\natural$-algèbre de type fini,
$S_+=S_1^\natural$ est un module de type fini sur $S=S_0^\natural$ par
(\hyperref[II.2.1.6-fr]{2.1.6}, (i)), donc $S$ est une $S_0$-algèbre de type
fini (\hyperref[II.2.1.4-fr]{2.1.4})~; si $S^\natural$ est un anneau
noethérien, il en est de même de $S_0^\natural=S$
(\hyperref[II.2.1.5-fr]{2.1.5}). Inversement, si $S$ est une $S_0$-algèbre

\oldpage[II]{160}
de type fini, on sait (\hyperref[II.2.1.6-fr]{2.1.6}, (ii)) qu'il existe
$h>0$ et $m_0>0$ tels que $S_{n+h}=S_hS_n$ pour $n\geq m_0$~; on peut
évidemment supposer $m_0\geq h$. En outre, les $S_m$ sont des $S_0$-modules
de type fini (\hyperref[II.2.1.6-fr]{2.1.6}, (i)). Cela étant, si
$n\geq m_0+h$, on a $S_n^\natural=S_hS_{n-h}^\natural
=S_h^\natural S_{n-h}^\natural$~; et si $m<m_0+h$, on a, en posant
$E=S_{m_0}+\cdots+S_{m_0+h-1}$,
\[
  S_m^\natural=S_m+\cdots+S_{m_0+h-1}+S_hE+S_h^2E+\cdots.
\]
Pour $1\leq m\leq m_0$, soit $G_m$ la réunion de systèmes finis de
générateurs des $S_0$-modules $S_i$, pour $m\leq i\leq m_0+h-1$,
considérée comme partie de $S_{[m]}$. Pour $m_0+1\leq m\leq m_0+h-1$, soit
de même $G_m$ la réunion de systèmes finis de générateurs des $S_0$-modules
$S_i$, pour $m\leq i\leq m_0+h-1$, et de $S_hE$, considérée comme partie de
$S_{[m]}$. Il est clair que l'on a $S_m^\natural=S_0^\natural G_m$ pour
$1\leq m\leq m_0+h-1$, et par suite la réunion $G$ des $G_m$ pour
$1\leq m\leq m_0+h-1$ est un système de générateurs de la
$S_0^\natural$-algèbre $S^\natural$. On en conclut que si $S=S_0^\natural$
est un anneau noethérien, il en est de même de $S^\natural$.

\item[\rm(ii)] Il est clair que si $S_{n+1}=S_1S_n$ pour $n\geq n_0$, on a
$S_{n+1}^\natural=S_1S_n^\natural$, et \emph{a fortiori}
$S_{n+1}^\natural=S_1^\natural S_n^\natural$ pour $n\geq n_0$.
Inversement, cette dernière relation s'écrit
\[
  S_{n+1}+S_{n+2}+\cdots=(S_1+S_2+\cdots)(S_n+S_{n+1}+\cdots)
\]
et la comparaison des termes de degré $n+1$ (dans $S$) aux deux membres donne
$S_{n+1}=S_1S_n$.

\item[\rm(iii)] Si $S_n=S_1^n$ pour $n\geq n_0$, on a
$S_n^\natural=S_1^n+S_1^{n+1}+\cdots$~; comme $S_1^\natural$ contient
$S_1+S_1^2+\cdots$, on a $S_n^\natural\subset S_1^{\natural n}$, et par suite
$S_n^\natural=S_1^{\natural n}$ pour $n\geq n_0$. Inversement, les seuls
termes de $S_1^{\natural n}=(S_1+S_2+\cdots)^n$ qui soient de degré $n$ dans
$S$ sont ceux de $S_1^n$~; la relation $S_n^\natural=S_1^{\natural n}$
implique donc $S_n=S_1^n$.

\item[\rm(iv)] Il suffit de prouver que si un élément $g\in S_{k+h}$ est
considéré comme un élément de $S_k^\natural$ ($k>0$, $h\geq0$), alors il
existe un entier $n>0$ tel que dans $S_{kn}^\natural$, $g^n$ soit combinaison
linéaire des $f_\alpha^\natural$ à coefficients dans $S^\natural$. Par
hypothèse, il y a un entier $m_0$ tel que pour $m\geq m_0$, on ait, \emph{dans
$S$}, $g^m=\sum_\alpha c_{\alpha m}f_\alpha$, les indices $\alpha$ figurant
dans cette formule étant \emph{indépendants de $m$}~; en outre, on peut
évidemment supposer les $c_{\alpha m}$ homogènes, avec
\[
  \deg(c_{\alpha m})=m(k+h)-\deg f_\alpha
\]
dans $S$. Prenons alors $m_0$ assez grand pour que l'on ait
$km_0>\deg f_\alpha$ pour tous les $f_\alpha$ qui figurent dans $g^{m_0}$~;
pour tout $\alpha$, soit $c'_{\alpha m}$ l'élément $c_{\alpha m}$ considéré
comme de degré $km-\deg(f_\alpha)$ dans $S^\natural$~; on a alors, dans
$S^\natural$, $g^m=\sum_\alpha c'_{\alpha m}f_\alpha^\natural$, ce qui
termine la démonstration.
\end{enumerate}

\begin{env}[8.2.11]
\label{II.8.2.11-fr}
Considérons la $S_0$-algèbre graduée
\begin{equation}
\label{II.8.2.11.1-fr}
  S^\natural\otimes_S S_0=S^\natural/S_+S^\natural
  =\bigoplus_{n\geq0}S_{[n]}/S_+S_{[n]}.
\tag{8.2.11.1}
\end{equation}
Comme $S_n$ est un $S_0$-module quotient de $S_{[n]}/S_+S_{[n]}$, on a un
homomorphisme canonique de $S_0$-algèbres graduées
\begin{equation}
\label{II.8.2.11.2-fr}
  S^\natural\otimes_S S_0\to S
\tag{8.2.11.2}
\end{equation}
qui est évidemment \emph{surjectif}, et correspond par suite
(\hyperref[II.2.9.2-fr]{2.9.2}) à une \emph{immersion fermée} canonique
\begin{equation}
\label{II.8.2.11.3-fr}
  \operatorname{Proj}(S)\to
  \operatorname{Proj}(S^\natural\otimes_S S_0).
\tag{8.2.11.3}
\end{equation}
\end{env}

\oldpage[II]{161}

\begin{proposition}[8.2.12]
\label{II.8.2.12-fr}
Le morphisme canonique (\hyperref[II.8.2.11.3-fr]{8.2.11.3}) est bijectif.
Pour que l'homomorphisme (\hyperref[II.8.2.11.2-fr]{8.2.11.2}) soit
(TN)-bijectif, il faut et il suffit qu'il existe $n_0$ tel que
$S_{n+1}=S_1S_n$ pour $n\geq n_0$. Si cette dernière condition est vérifiée,
(\hyperref[II.8.2.11.3-fr]{8.2.11.3}) est un isomorphisme~; la réciproque est
vraie lorsque $S$ est noethérien.
\end{proposition}

Pour démontrer la première assertion, il suffit
(\hyperref[II.2.8.3-fr]{2.8.3}) de prouver que le noyau $\mathfrak{I}$ de
l'homomorphisme (\hyperref[II.8.2.11.2-fr]{8.2.11.2}) est formé d'éléments
\emph{nilpotents}. Or si $f\in S_{[n]}$ est un élément dont la classe mod.
$S_+S_{[n]}$ appartient à ce noyau, cela signifie que $f\in S_{[n+1]}$~;
l'élément $f^{n+1}$, considéré comme appartenant à $S_{[n(n+1)]}$, appartient
alors à $S_+S_{[n(n+1)]}$, puisqu'il s'écrit $f\cdot f^n$~; donc la classe de
$f^{n+1}$ mod. $S_+S_{[n(n+1)]}$ est nulle, ce qui prouve notre assertion.
Comme l'hypothèse $S_{n+1}=S_1S_n$ pour $n\geq n_0$ équivaut à
$S_{n+1}^{\natural}=S_1^{\natural}S_n^{\natural}$ pour $n\geq n_0$
(\hyperref[II.8.2.10-fr]{8.2.10}, (ii)), cette hypothèse équivaut par définition
au fait que (\hyperref[II.8.2.11.2-fr]{8.2.11.2}) est (TN)-injectif, donc
(TN)-bijectif, et alors (\hyperref[II.8.2.11.3-fr]{8.2.11.3}) est un
isomorphisme en vertu de (\hyperref[II.2.9.1-fr]{2.9.1}). Inversement, si
(\hyperref[II.8.2.11.3-fr]{8.2.11.3}) est un isomorphisme, le faisceau
$\widetilde{\mathfrak{I}}$ sur
$\operatorname{Proj}(S^{\natural}\otimes_S S_0)$ est nul
(\hyperref[II.2.9.2-fr]{2.9.2}, (i))~; comme
$S^{\natural}\otimes_S S_0$ est noethérien comme quotient de $S^{\natural}$
(\hyperref[II.8.2.10-fr]{8.2.10}, (i)), on conclut de
(\hyperref[II.2.7.3-fr]{2.7.3}) que $\mathfrak{I}$ vérifie la condition (TN),
donc $S_{n+1}^{\natural}=S_1^{\natural}S_n^{\natural}$ pour $n\geq n_0$, et
cela achève la démonstration en vertu de
(\hyperref[II.8.2.10-fr]{8.2.10}, (ii)).

\begin{env}[8.2.13]
\label{II.8.2.13-fr}
Considérons maintenant les injections canoniques $(S_+)^n\to S_{[n]}$, qui
définissent un homomorphisme injectif de degré $0$ d'anneaux gradués
\begin{equation}
\label{II.8.2.13.1-fr}
  \bigoplus_{n\geq0}(S_+)^n\to S^{\natural}.
\tag{8.2.13.1}
\end{equation}
\end{env}

\begin{proposition}[8.2.14]
\label{II.8.2.14-fr}
Pour que l'homomorphisme (\hyperref[II.8.2.13.1-fr]{8.2.13.1}) soit un
(TN)-isomorphisme, il faut et il suffit qu'il existe $n_0$ tel que
$S_n=S_1^n$ pour tout $n\geq n_0$. Lorsqu'il en est ainsi, le morphisme
correspondant à (\hyperref[II.8.2.13.1-fr]{8.2.13.1}) est partout défini et
est un isomorphisme
\[
  \operatorname{Proj}(S^{\natural})\xrightarrow{\sim}
  \operatorname{Proj}\left(\bigoplus_{n\geq0}(S_+)^n\right)~;
\]
la réciproque est vraie lorsque $S$ est noethérien.
\end{proposition}

Les deux premières assertions sont évidentes, vu
(\hyperref[II.8.2.10-fr]{8.2.10}, (iii)) et
(\hyperref[II.2.9.1-fr]{2.9.1}). La troisième résultera de
(\hyperref[II.8.2.10-fr]{8.2.10}, (i) et (iii)) et du lemme suivant~:

\begin{lemma}[8.2.14.1]
\label{II.8.2.14.1-fr}
Soit $T$ un anneau gradué en degrés positifs qui est une $T_0$-algèbre de type
fini. Si le morphisme correspondant à l'homomorphisme injectif
$\bigoplus_{n\geq0}T_1^n\to T$ est partout défini et est un isomorphisme
\[
  \operatorname{Proj}(T)\to
  \operatorname{Proj}\left(\bigoplus_{n\geq0}T_1^n\right),
\]
il existe $n_0$ tel que $T_n=T_1^n$ pour $n\geq n_0$.
\end{lemma}

En effet, soient $g_i$ ($1\leq i\leq r$) des générateurs du $T_0$-module
$T_1$. L'hypothèse entraîne d'abord que les $D_+(g_i)$ recouvrent
$\operatorname{Proj}(T)$ (\hyperref[II.2.8.1-fr]{2.8.1}). Soit
$(h_j)_{1\leq j\leq s}$ un système d'éléments homogènes de $T_+$, avec
$\deg(h_j)=n_j$, qui forment avec les $g_i$ un système de générateurs de
l'idéal $T_+$, ou (ce qui revient au même
(\hyperref[II.2.1.3-fr]{2.1.3})) un système de générateurs de $T$ en tant que
$T_0$-algèbre~; si on pose $T'=\bigoplus_{n\geq0}T_1^n$, l'élément
$h_j/g_i^{n_j}$ de l'anneau $T_{(g_i)}$ doit par hypothèse appartenir au
sous-anneau $T'_{(g_i)}$, donc il existe un entier $k$ tel que
$T_1^k h_j\subset T_1^{k+n_j}$ pour tout $j$. On en conclut par récurrence
sur $r$ que $T_1^k h_j^r\subset T'$ pour tout $r\geq1$, et par définition des
$h_j$, on a donc $T_1^kT\subset T'$. Il y a d'autre part pour tout $j$ un
entier $m_j$ tel que $h_j^{m_j}$ appartienne à l'idéal de $T$ engendré par
les $g_i$ (\hyperref[II.2.3.14-fr]{2.3.14}), donc $h_j^{m_j}\in T_1T$, et

\oldpage[II]{162}
$h_j^{m_jk}\in T_1^kT\subset T'$. Il y a par suite un entier $m_0\geq k$ tel
que $h_j^m\in T_1^{mn_j}$ pour $m\geq m_0$. Cela étant, si $q$ est le plus
grand des entiers $n_j$, le nombre $n_0=qsm_0+k$ répond à la question. En
effet, un élément de $S_n$, pour $n\geq n_0$, est somme de monômes appartenant
à $T_1^\alpha u$, où $u$ est un produit de puissances des $h_j$~; si
$\alpha\geq k$, il résulte de ce qui précède que $T_1^\alpha u\subset T_1^n$~;
dans le cas contraire, un des exposants des $h_j$ est $\geq m_0$, donc
$u\in T_1^\beta v$, où $\beta\geq k$ et $v$ est encore un produit de
puissances des $h_j$~; on est alors ramené au cas précédent, et on voit
finalement que $T_1^\alpha u\subset T_1^n$ dans tous les cas.

\begin{remark}[8.2.15]
\label{II.8.2.15-fr}
La condition $S_n=S_1^n$ pour $n\geq n_0$ entraîne évidemment
$S_{n+1}=S_1S_n$ pour $n\geq n_0$, mais la réciproque est inexacte, même
quand on suppose $S$ noethérien. Par exemple, soient $K$ un corps,
$A=K[x]$, $B=K[y]/y^2K[y]$, où $x$ et $y$ sont deux indéterminées, $x$ étant
prise de degré $1$ et $y$ de degré $2$, et soit $S=A\otimes_KB$, de sorte que
$S$ est une algèbre graduée sur $K$ ayant une base formée des éléments $1$,
$x^n$ ($n\geq1$) et $x^ny$ ($n\geq0$). Il est immédiat que
$S_{n+1}=S_1S_n$ pour $n\geq2$, mais $S_1^n=Kx^n$ tandis que
$S_n=Kx^n+Kx^{n-2}y$ pour $n\geq2$.
\end{remark}

\subsection{Cônes projetants}
\label{subsection:II.8.3-fr}

\begin{env}[8.3.1]
\label{II.8.3.1-fr}
Soit $Y$ un préschéma~; dans tout ce numéro, il ne sera question que de
\emph{$Y$-préschémas} et de \emph{$Y$-morphismes}. Soit $\mathcal{S}$ une
$\mathcal{O}_Y$-Algèbre quasi-cohérente graduée à \emph{degrés positifs}~;
\emph{nous supposerons en outre que l'on a $\mathcal{S}_0=\mathcal{O}_Y$}.
Conformément aux notations introduites en
(\hyperref[II.8.2.2-fr]{8.2.2}), nous poserons
\begin{equation}
\label{II.8.3.1.1-fr}
  \widehat{\mathcal{S}}=\mathcal{S}[z]
  =\mathcal{S}\otimes_{\mathcal{O}_Y}\mathcal{O}_Y[z]
\tag{8.3.1.1}
\end{equation}
que l'on considère comme $\mathcal{O}_Y$-Algèbre graduée à degrés positifs en
définissant les degrés par la formule
(\hyperref[II.8.2.2.2-fr]{8.2.2.2}), de sorte que pour tout ouvert affine $U$
de $Y$, on a
\[
  \Gamma(U,\widehat{\mathcal{S}})=(\Gamma(U,\mathcal{S}))[z].
\]
Dans ce qui suit, nous poserons
\begin{equation}
\label{II.8.3.1.2-fr}
  X=\operatorname{Proj}(\mathcal{S}),\qquad
  C=\operatorname{Spec}(\mathcal{S}),\qquad
  \widehat{C}=\operatorname{Proj}(\widehat{\mathcal{S}})
\tag{8.3.1.2}
\end{equation}
(où dans la définition de $C$, $\mathcal{S}$ est considérée comme
$\mathcal{O}_Y$-Algèbre non graduée), et nous dirons que $C$ (resp.
$\widehat{C}$) est le \emph{cône affine} (resp. \emph{cône projectif}) défini
par $\mathcal{S}$~; on dira aussi «~cône~» au lieu de «~cône affine~». Par
abus de langage, nous dirons encore que $C$ (resp. $\widehat{C}$) est le
\emph{cône projetant affine de $X$} (resp. le \emph{cône projetant projectif
de $X$}), étant sous-entendu que le préschéma $X$ est donné sous la forme
$\operatorname{Proj}(\mathcal{S})$~; enfin, nous dirons que $\widehat{C}$ est
la \emph{fermeture projective} de $C$ (étant entendu que la donnée de
$\mathcal{S}$ figure dans la structure de $C$).
\end{env}

\begin{proposition}[8.3.2]
\label{II.8.3.2-fr}
Il existe des $Y$-morphismes canoniques
\begin{equation}
\label{II.8.3.2.1-fr}
  Y\xrightarrow{\varepsilon}C\xrightarrow{i}\widehat{C}
\tag{8.3.2.1}
\end{equation}
\begin{equation}
\label{II.8.3.2.2-fr}
  X\xrightarrow{j}\widehat{C}
\tag{8.3.2.2}
\end{equation}

\oldpage[II]{163}
tels que $\varepsilon$ et $j$ soient des immersions fermées, et $i$ un
morphisme affine, qui est une immersion ouverte dominante, pour laquelle
\begin{equation}
\label{II.8.3.2.3-fr}
  i(C)=\widehat{C}-j(X)~;
\tag{8.3.2.3}
\end{equation}
en outre $\widehat{C}$ est le plus petit sous-préschéma fermé de $\widehat{C}$
majorant $i(C)$.
\end{proposition}

Pour définir $i$, on considère l'ouvert de $\widehat{C}$,
\begin{equation}
\label{II.8.3.2.4-fr}
  \widehat{C}_z=
  \operatorname{Spec}(\widehat{\mathcal{S}}/(z-1)\widehat{\mathcal{S}})
\tag{8.3.2.4}
\end{equation}
(\hyperref[II.3.1.4-fr]{3.1.4}), $z$ s'identifiant canoniquement à une
section de $\widehat{\mathcal{S}}$ au-dessus de $Y$. L'isomorphisme
$i:C\xrightarrow{\sim}\widehat{C}_z$ correspond alors à l'isomorphisme
canonique (\hyperref[II.8.2.3.1-fr]{8.2.3.1})
\[
  \widehat{\mathcal{S}}/(z-1)\widehat{\mathcal{S}}
  \xrightarrow{\sim}\mathcal{S}.
\]
Le morphisme $\varepsilon$ correspond à l'homomorphisme d'augmentation
$\mathcal{S}\to\mathcal{S}_0=\mathcal{O}_Y$ ayant pour noyau
$\mathcal{S}_+$ (\hyperref[II.1.2.7-fr]{1.2.7}) et comme ce dernier est
surjectif, $\varepsilon$ est une immersion fermée
(\hyperref[II.1.4.10-fr]{1.4.10}). Enfin, $j$ correspond de même
(\hyperref[II.3.5.1-fr]{3.5.1}) à l'homomorphisme surjectif de degré $0$,
$\widehat{\mathcal{S}}\to\mathcal{S}$, qui se réduit à l'identité dans
$\mathcal{S}$ et est nul dans $z\widehat{\mathcal{S}}$, qui est son noyau~;
$j$ est partout défini et est une immersion fermée en vertu de
(\hyperref[II.3.6.2-fr]{3.6.2}).

Pour démontrer les autres assertions de
(\hyperref[II.8.3.2-fr]{8.3.2}), on peut évidemment se borner au cas où
$Y=\operatorname{Spec}(A)$ est affine, $\mathcal{S}=\widetilde{S}$, où $S$
est une $A$-algèbre graduée, d'où
$\widehat{\mathcal{S}}=(\widehat{S})^{\sim}$~; les éléments homogènes $f$ de
$S_+$ s'identifient alors à des sections de $\widehat{\mathcal{S}}$ au-dessus
de $Y$, et l'ouvert de $\widehat{C}$ noté $D_+(f)$ dans
(\hyperref[II.2.3.3-fr]{2.3.3}) s'écrit aussi $\widehat{C}_f$
(\hyperref[II.3.1.4-fr]{3.1.4})~; de même l'ouvert de $C$ noté $D(f)$ dans
(\hyperref[I.1.1.1-fr]{I, 1.1.1}) s'écrit aussi $C_f$
(\hyperref[0.5.5.2-fr]{0, 5.5.2}). Cela étant, il résulte de
(\hyperref[II.2.3.14-fr]{2.3.14}) et de la définition de $\widehat{S}$ que
dans le cas présent, les ouverts $\widehat{C}_z=i(C)$ et $\widehat{C}_f$
($f$ homogène dans $S_+$) constituent un \emph{recouvrement} de
$\widehat{C}$. En outre, on a, avec ces notations,
\begin{equation}
\label{II.8.3.2.5-fr}
  i^{-1}(\widehat{C}_f)=C_f~;
\tag{8.3.2.5}
\end{equation}
en effet,
$\widehat{C}_f\cap i(C)=\widehat{C}_f\cap\widehat{C}_z
=\widehat{C}_{fz}=\operatorname{Spec}(\widehat{S}_{(fz)})$. Or, si
$d=\deg(f)$, $\widehat{S}_{(fz)}$ est canoniquement isomorphe à
$(\widehat{S}_{(z)})_{f/z^d}$ (\hyperref[II.2.2.2-fr]{2.2.2}), et il résulte
de la définition de l'isomorphisme
(\hyperref[II.8.2.3.1-fr]{8.2.3.1}) que l'image de
$(\widehat{S}_{(z)})_{f/z^d}$ par l'isomorphisme correspondant des anneaux de
fractions est exactement $S_f$. Comme $C_f=\operatorname{Spec}(S_f)$, cela
prouve (\hyperref[II.8.3.2.5-fr]{8.3.2.5}) et montre en même temps que le
morphisme $i$ est affine~; en outre, la restriction de $i$ à $C_f$,
considérée comme morphisme dans $\widehat{C}_f$, correspond
(\hyperref[I.1.7.3-fr]{I, 1.7.3}) à l'homomorphisme canonique
$\widehat{S}_{(f)}\to\widehat{S}_{(fz)}$ et en vertu de ce qui précède et de
(\hyperref[II.8.2.3.2-fr]{8.2.3.2}), on peut énoncer le résultat suivant~:

\begin{env}[8.3.2.6]
\label{II.8.3.2.6-fr}
Si $Y=\operatorname{Spec}(A)$ est affine, et $\mathcal{S}=\widetilde{S}$,
alors, pour tout $f$ homogène dans $S_+$, $\widehat{C}_f$ s'identifie
canoniquement à $\operatorname{Spec}(S_f^{\leq})$, et le morphisme
$C_f\to\widehat{C}_f$ restriction de $i$ correspond alors à l'injection
canonique $S_f^{\leq}\to S_f$.
\end{env}

Notons maintenant que (pour $Y$ affine) le complémentaire de
$\widehat{C}_z$ dans $\widehat{C}=\operatorname{Proj}(\widehat{S})$

\oldpage[II]{164}
est, par définition, l'ensemble des idéaux premiers gradués de $\widehat{S}$
contenant $z$, c'est-à-dire $j(X)$ par définition de $j$, ce qui démontre
(\hyperref[II.8.3.2.3-fr]{8.3.2.3}).

Pour prouver enfin la dernière assertion de
(\hyperref[II.8.3.2-fr]{8.3.2}), on peut toujours supposer $Y$ affine. Avec
les notations précédentes, remarquons que dans l'anneau $\widehat{S}$, $z$
n'est pas diviseur de $0$~; comme $\overline{i(C)}=\widehat{C}$, il suffira
de prouver le

\begin{lemma}[8.3.2.7]
\label{II.8.3.2.7-fr}
Soient $T$ un anneau gradué à degrés positifs,
$Z=\operatorname{Proj}(T)$, $g$ un élément homogène de $T$ de degré $d>0$.
Si $g$ n'est pas diviseur de $0$ dans $T$, $Z$ est le plus petit
sous-préschéma fermé de $Z$ qui majore $Z_g=D_+(g)$.
\end{lemma}

En vertu de (\hyperref[I.4.1.9-fr]{I, 4.1.9}) la question est locale sur $Z$~;
pour tout élément homogène $h\in T_e$ ($e>0$), il suffit donc de prouver que
$Z_h$ est le plus petit sous-préschéma fermé de $Z_h$ qui majore $Z_{gh}$~;
il résulte des définitions et de (\hyperref[I.4.3.2-fr]{I, 4.3.2}) que cette
condition équivaut au fait que l'homomorphisme canonique
$T_{(h)}\to T_{(gh)}$ est \emph{injectif}. Or cet homomorphisme s'identifie à
l'homomorphisme canonique
$T_{(h)}\to(T_{(h)})_{g^e/h^d}$ (\hyperref[II.2.2.3-fr]{2.2.3}). Mais comme
$g^e$ n'est pas diviseur de $0$ dans $T$, $g^e/h^d$ ne l'est pas dans $T_h$
(ni \emph{a fortiori} dans $T_{(h)}$), car la relation
$(g^e/h^d)(t/h^m)=0$ avec $t\in T$ et $m>0$ entraîne l'existence d'un $n>0$
tel que $h^ng^et=0$, d'où $h^nt=0$, et par suite $t/h^m=0$ dans $T_h$. Cela
achève donc la démonstration (\hyperref[0.1.2.2-fr]{0, 1.2.2}).

\begin{env}[8.3.3]
\label{II.8.3.3-fr}
On identifiera souvent le cône affine $C$ au sous-préschéma induit par le cône
projectif $\widehat{C}$ sur l'ouvert $i(C)$, au moyen de l'immersion ouverte
$i$. Le sous-préschéma fermé de $C$, associé à l'immersion fermée
$\varepsilon$, est appelé le \emph{préschéma des sommets de $C$}~; on dit
aussi que $\varepsilon$, qui est une $Y$-section de $C$, est la \emph{section
sommet} ou la \emph{section nulle de $C$}~; on peut identifier $Y$ au
préschéma des sommets de $C$ au moyen de $\varepsilon$. D'ailleurs
$i\circ\varepsilon$ est une $Y$-section de $\widehat{C}$, donc aussi une
immersion fermée (\hyperref[I.5.4.6-fr]{I, 5.4.6}), correspondant à
l'homomorphisme canonique surjectif de degré $0$,
$\widehat{\mathcal{S}}=\mathcal{S}[z]\to\mathcal{O}_Y[z]$
(cf. (\hyperref[II.3.1.7-fr]{3.1.7})), de noyau
$\mathcal{S}_+[z]=\mathcal{S}_+\widehat{\mathcal{S}}$~; le sous-préschéma de
$\widehat{C}$ associé à cette immersion fermée est encore appelé le
\emph{préschéma des sommets de $\widehat{C}$}, et $i\circ\varepsilon$ la
\emph{section sommet de $\widehat{C}$}~; il peut s'identifier par
$i\circ\varepsilon$ à $Y$. Enfin, le sous-préschéma fermé de $\widehat{C}$
associé à $j$ est appelé le \emph{lieu à l'infini de $\widehat{C}$}, et peut
être identifié à $X$ au moyen de $j$.
\end{env}

\begin{env}[8.3.4]
\label{II.8.3.4-fr}
Les sous-préschémas de $C$ (resp. $\widehat{C}$) induits respectivement sur
les \emph{ouverts}
\begin{equation}
\label{II.8.3.4.1-fr}
  E=C-\varepsilon(Y),\qquad
  \widehat{E}=\widehat{C}-i(\varepsilon(Y))
\tag{8.3.4.1}
\end{equation}
sont appelés respectivement (par abus de langage) le \emph{cône affine
épointé} et le \emph{cône projectif épointé} définis par $\mathcal{S}$~; on
notera qu'en dépit de cette terminologie, $E$ \emph{n'est pas nécessairement
affine} sur $Y$, ni $\widehat{E}$ projectif sur $Y$
(cf. (\hyperref[II.8.4.3-fr]{8.4.3})). Lorsqu'on identifie $C$ à $i(C)$, on a
donc pour les espaces sous-jacents
\begin{equation}
\label{II.8.3.4.2-fr}
  C\cup\widehat{E}=\widehat{C},\qquad C\cap\widehat{E}=E
\tag{8.3.4.2}
\end{equation}
de sorte que $\widehat{C}$ peut être considéré comme obtenu par
\emph{recollement} des sous-préschémas ouverts $C$ et $\widehat{E}$~; en
outre, en vertu de (\hyperref[II.8.3.2.3-fr]{8.3.2.3}),
\begin{equation}
\label{II.8.3.4.3-fr}
  E=\widehat{E}-j(X).
\tag{8.3.4.3}
\end{equation}

\oldpage[II]{165}
Lorsque $Y=\operatorname{Spec}(A)$ est affine, on a, avec les notations de
(\hyperref[II.8.3.2-fr]{8.3.2})
\begin{equation}
\label{II.8.3.4.4-fr}
  E=\bigcup C_f,\qquad
  \widehat{E}=\bigcup\widehat{C}_f,\qquad
  C_f=C\cap\widehat{C}_f
\tag{8.3.4.4}
\end{equation}
$f$ parcourant l'ensemble des éléments homogènes de $S_+$ (ou seulement une
partie $M$ de cet ensemble engendrant un idéal de $S_+$ dont la racine dans
$S_+$ est $S_+$ lui-même, autrement dit, telle que les $X_f$ pour $f\in M$
recouvrent $X$ (\hyperref[II.2.3.14-fr]{2.3.14})). Le recollement de $C$ et de
$\widehat{C}_f$ le long de $C_f$ est alors déterminé par des morphismes
d'injection $C_f\to C$, $C_f\to\widehat{C}_f$, qui, ainsi qu'on l'a vu
(\hyperref[II.8.3.2.6-fr]{8.3.2.6}) correspondent respectivement aux
homomorphismes canoniques $S\to S_f$, $S_f^{\leq}\to S_f$.
\end{env}

\begin{proposition}[8.3.5]
\label{II.8.3.5-fr}
Avec les notations de (\hyperref[II.8.3.1-fr]{8.3.1}) et
(\hyperref[II.8.3.4-fr]{8.3.4}), le morphisme associé
(\hyperref[II.3.5.1-fr]{3.5.1}) à l'injection canonique
$\varphi:\mathcal{S}\to\widehat{\mathcal{S}}=\mathcal{S}[z]$ est un morphisme
affine surjectif (dit \emph{rétraction canonique})
\begin{equation}
\label{II.8.3.5.1-fr}
  p:\widehat{E}\to X
\tag{8.3.5.1}
\end{equation}
tel que l'on ait
\begin{equation}
\label{II.8.3.5.2-fr}
  p\circ j=1_X.
\tag{8.3.5.2}
\end{equation}
\end{proposition}

Pour démontrer la proposition, on peut se borner au cas où $Y$ est affine.
Compte tenu de l'expression (\hyperref[II.8.3.4.4-fr]{8.3.4.4}) de
$\widehat{E}$, le fait que le domaine de définition $G(\varphi)$ de $p$ est
égal à $\widehat{E}$ résultera de la première des assertions suivantes~:

\begin{env}[8.3.5.3]
\label{II.8.3.5.3-fr}
Si $Y=\operatorname{Spec}(A)$ est affine, et $\mathcal{S}=\widetilde{S}$, on
a, pour tout $f\in S_+$ homogène
\begin{equation}
\label{II.8.3.5.4-fr}
  p^{-1}(X_f)=\widehat{C}_f
\tag{8.3.5.4}
\end{equation}
et la restriction de $p$ à
$\widehat{C}_f=\operatorname{Spec}(S_f^{\leq})$, considérée comme morphisme de
$\widehat{C}_f$ dans $X_f$, correspond à l'injection canonique
$S_{(f)}\to S_f^{\leq}$. Si de plus $f\in S_1$, $\widehat{C}_f$ est isomorphe
à $X_f\otimes_{\mathbf{Z}}\mathbf{Z}[T]$ ($T$ indéterminée).
\end{env}

En effet, la formule (\hyperref[II.8.3.5.4-fr]{8.3.5.4}) n'est qu'un cas
particulier de (\hyperref[II.2.8.1.1-fr]{2.8.1.1}) et la seconde assertion
n'est autre que la définition de $\operatorname{Proj}(\varphi)$ lorsque $Y$
est affine (\hyperref[II.2.8.1-fr]{2.8.1}). Alors la formule
(\hyperref[II.8.3.5.2-fr]{8.3.5.2}) et le fait que $p$ est surjectif résultent
de ce que le composé $\mathcal{S}\to\widehat{\mathcal{S}}\to\mathcal{S}$ des
homomorphismes canoniques est l'identité dans $\mathcal{S}$. Enfin, la
dernière assertion de (\hyperref[II.8.3.5.3-fr]{8.3.5.3}) provient de ce que
$S_f^{\leq}$ est isomorphe à $S_{(f)}[T]$ lorsque $f\in S_1$
(\hyperref[II.2.2.1-fr]{2.2.1}).

\begin{corollary}[8.3.6]
\label{II.8.3.6-fr}
La restriction
\begin{equation}
\label{II.8.3.6.1-fr}
  \pi:E\to X
\tag{8.3.6.1}
\end{equation}
de $p$ à $E$ est un morphisme affine surjectif. Si $Y$ est affine et $f$
homogène dans $S_+$, on a
\begin{equation}
\label{II.8.3.6.2-fr}
  \pi^{-1}(X_f)=C_f
\tag{8.3.6.2}
\end{equation}
et la restriction de $\pi$ à $C_f$ correspond à l'injection canonique
$S_{(f)}\to S_f$. Si en outre $f\in S_1$, $C_f$ est isomorphe à
$X_f\otimes_{\mathbf{Z}}\mathbf{Z}[T,T^{-1}]$ ($T$ indéterminée).
\end{corollary}

La formule (\hyperref[II.8.3.6.2-fr]{8.3.6.2}) résulte aussitôt de
(\hyperref[II.8.3.5.3-fr]{8.3.5.3}) et
(\hyperref[II.8.3.2.5-fr]{8.3.2.5}), et démontre la surjectivité de $\pi$~; on
a vu par ailleurs que l'immersion $i$, restreinte à $C_f$, correspondait

\oldpage[II]{166}
à l'injection $S_f^{\leq}\to S_f$ (\hyperref[II.8.3.2-fr]{8.3.2}). Enfin, la
dernière assertion est conséquence de ce que pour $f\in S_1$, $S_f$ est
isomorphe à $S_{(f)}[T,T^{-1}]$ (\hyperref[II.2.2.1-fr]{2.2.1}).

\begin{remark}[8.3.7]
\label{II.8.3.7-fr}
Lorsque $Y$ est affine, les éléments de l'espace sous-jacent à $E$ sont les
idéaux premiers $\mathfrak{p}$ (non nécessairement gradués) de $S$ ne
contenant pas $S_+$, en vertu de la définition de l'immersion $\varepsilon$
(\hyperref[II.8.3.2-fr]{8.3.2}). Pour un tel idéal $\mathfrak{p}$, les
$\mathfrak{p}\cap S_n$ vérifient de façon évidente les conditions de
(\hyperref[II.2.1.9-fr]{2.1.9}), donc il existe un seul idéal premier
\emph{gradué} $\mathfrak{q}$ de $S$ tel que
$\mathfrak{q}\cap S_n=\mathfrak{p}\cap S_n$ pour tout $n$~; l'application
$\pi:E\to X$ des espaces sous-jacents s'interprète alors grâce à la relation
\begin{equation}
\label{II.8.3.7.1-fr}
  \pi(\mathfrak{p})=\mathfrak{q}.
\tag{8.3.7.1}
\end{equation}
En effet, pour vérifier cette relation, il suffit de considérer un $f$
homogène dans $S_+$ tel que $\mathfrak{p}\in D(f)$, et d'observer que
$\mathfrak{q}_{(f)}$ est l'image réciproque de $\mathfrak{p}_f$ par
l'injection $S_{(f)}\to S_f$.
\end{remark}

\begin{corollary}[8.3.8]
\label{II.8.3.8-fr}
Si $\mathcal{S}$ est engendrée par $\mathcal{S}_1$, les morphismes $p$ et
$\pi$ sont de type fini~; pour tout $x\in X$, la fibre $p^{-1}(x)$ est
isomorphe à $\operatorname{Spec}(k(x)[T])$ et la fibre $\pi^{-1}(x)$ est
isomorphe à $\operatorname{Spec}(k(x)[T,T^{-1}])$.
\end{corollary}

Cela résulte aussitôt de (\hyperref[II.8.3.5-fr]{8.3.5}) et
(\hyperref[II.8.3.6-fr]{8.3.6}) en observant que, lorsque $Y$ est affine et
$S$ engendrée par $S_1$, les $X_f$, pour $f\in S_1$, forment un recouvrement
de $X$ (\hyperref[II.2.3.14-fr]{2.3.14}).

\begin{remark}[8.3.9]
\label{II.8.3.9-fr}
Le cône épointé affine correspondant à la $\mathcal{O}_Y$-Algèbre graduée
$\mathcal{O}_Y[T]$ ($T$ indéterminée) s'identifie à
$G_m=\operatorname{Spec}(\mathcal{O}_Y[T,T^{-1}])$, puisqu'il n'est autre que
$C_T$, comme on l'a vu dans (\hyperref[II.8.3.2-fr]{8.3.2}) (voir
(\hyperref[II.8.4.4-fr]{8.4.4}) pour un résultat plus général). Ce préschéma
est canoniquement muni d'une structure de «~\emph{$Y$-schéma en groupes
commutatif}~». Il s'agit là d'une notion qui sera développée en détail plus
tard, et que l'on peut ici résumer rapidement comme suit. Un $Y$-schéma en
groupes est un $Y$-schéma $G$, muni de deux $Y$-morphismes
$p:G\times_YG\to G$ et $s:G\to G$ qui satisfont aux conditions formellement
analogues aux axiomes de la loi de composition et de l'application de
symétrie dans un groupe~: le diagramme
\[
  \xymatrix{
    G\times G\times G \ar[r]^{p\times1}\ar[d]_{1\times p}
      & G\times G \ar[d]^{p} \\
    G\times G \ar[r]_{p} & G
  }
\]
doit être commutatif («~associativité~») et l'on doit avoir une condition qui
correspond dans les groupes au fait que les applications
\[
  (x,y)\mapsto(x,x^{-1},y)\mapsto(x,x^{-1}y)\mapsto x(x^{-1}y)
\]
et
\[
  (x,y)\mapsto(x,x^{-1},y)\mapsto(x,yx^{-1})\mapsto(yx^{-1})x
\]
doivent toutes deux se réduire à $(x,y)\mapsto y$~; la suite de morphismes
correspondant par exemple à la première application composée est
\[
  G\times G\xrightarrow{(1,s)\times1}G\times G\times G
  \xrightarrow{1\times p}G\times G\xrightarrow{p}G
\]
et le lecteur écrira de même la seconde suite.

\oldpage[II]{167}
Il est immédiat (\hyperref[I.3.4.3-fr]{I, 3.4.3}) que la donnée d'une structure
de $Y$-schéma en groupes sur un $Y$-schéma $G$ équivaut à la donnée, pour
\emph{tout} $Y$-préschéma $Z$, d'une structure de \emph{groupe} sur
l'ensemble $\operatorname{Hom}_Y(Z,G)$, ces structures devant être telles que
pour tout $Y$-morphisme $Z\to Z'$, l'application correspondante
$\operatorname{Hom}_Y(Z',G)\to\operatorname{Hom}_Y(Z,G)$ soit un
homomorphisme de groupes. Dans le cas particulier de $G_m$ que nous
considérons ici, $\operatorname{Hom}_Y(Z,G)$ s'identifie à l'ensemble des
$Z$-sections de $Z\times_YG_m$ (\hyperref[I.3.3.14-fr]{I, 3.3.14}), donc à
l'ensemble des $Z$-sections de
$\operatorname{Spec}(\mathcal{O}_Z[T,T^{-1}])$~; finalement, le même
raisonnement que dans (\hyperref[I.3.3.15-fr]{I, 3.3.15}) montre que cet
ensemble s'identifie canoniquement à l'ensemble des éléments \emph{inversibles}
de l'anneau $\Gamma(Z,\mathcal{O}_Z)$, et la structure de groupe sur cet
ensemble est la structure provenant de la multiplication dans l'anneau
$\Gamma(Z,\mathcal{O}_Z)$. Le lecteur pourra vérifier que les morphismes $p$
et $s$ considérés plus haut sont obtenus de la façon suivante~: ils
correspondent d'après (\hyperref[II.1.2.7-fr]{1.2.7}) et
(\hyperref[II.1.4.6-fr]{1.4.6}) à des homomorphismes de $\mathcal{O}_Y$-Algèbres
\begin{align*}
  \pi &: \mathcal{O}_Y[T,T^{-1}]
  \to\mathcal{O}_Y[T,T^{-1},T',T'^{-1}],\\
  \sigma &: \mathcal{O}_Y[T,T^{-1}]\to\mathcal{O}_Y[T,T^{-1}]
\end{align*}
et sont entièrement définis par les données de $\pi(T)=TT'$ et
$\sigma(T)=T^{-1}$.

Cela étant, $G_m$ peut être considéré comme un «~\emph{domaine d'opérateurs
universel}~» pour tout \emph{cône affine}
$C=\operatorname{Spec}(\mathcal{S})$, où $\mathcal{S}$ est une
$\mathcal{O}_Y$-Algèbre graduée quasi-cohérente à degrés positifs. Cela
signifie qu'on peut définir canoniquement un $Y$-morphisme
$G_m\times_YC\to C$ qui a les propriétés formelles d'une loi externe d'un
ensemble muni d'un groupe d'opérateurs~; ou encore, comme ci-dessus pour les
schémas en groupes, on se donne pour tout $Y$-préschéma $Z$ une loi externe
sur $\operatorname{Hom}_Y(Z,C)$, ayant le groupe
$\operatorname{Hom}_Y(Z,G_m)$ comme ensemble d'opérateurs, avec les axiomes
usuels des ensembles munis d'un groupe d'opérateurs, et une condition de
compatibilité avec les $Y$-morphismes $Z\to Z'$. Dans le cas présent, le
morphisme $G_m\times_YC\to C$ est défini par la donnée de l'homomorphisme de
$\mathcal{O}_Y$-Algèbres
\[
  \mathcal{S}\to
  \mathcal{S}\otimes_{\mathcal{O}_Y}\mathcal{O}_Y[T,T^{-1}]
  =\mathcal{S}[T,T^{-1}]
\]
qui, à toute section $s_n\in\Gamma(U,\mathcal{S}_n)$ ($U$ ouvert de $Y$) fait
correspondre la section
$s_nT^n\in\Gamma(U,\mathcal{S}\otimes_{\mathcal{O}_Y}
\mathcal{O}_Y[T,T^{-1}])$.

Inversement, supposons donnée une $\mathcal{O}_Y$-Algèbre quasi-cohérente
$\mathcal{S}$ \emph{non graduée a priori}, et, sur
$C=\operatorname{Spec}(\mathcal{S})$, une structure de «~$Y$-schéma en
ensembles \emph{munis de groupes d'opérateurs}~» ayant pour domaine
d'opérateurs le $Y$-schéma en groupes $G_m$~; alors on en déduit
canoniquement une \emph{graduation} de $\mathcal{O}_Y$-Algèbre sur
$\mathcal{S}$. En effet, la donnée d'un $Y$-morphisme $G_m\times_YC\to C$
équivaut à celle d'un homomorphisme de $\mathcal{O}_Y$-Algèbres
$\psi:\mathcal{S}\to\mathcal{S}[T,T^{-1}]$, qui s'écrira donc
$\psi=\sum_{n\in\mathbf{Z}}\psi_nT^n$, où les
$\psi_n:\mathcal{S}\to\mathcal{S}$ sont des homomorphismes de
$\mathcal{O}_Y$-Modules (avec $\psi_n(s)=0$ sauf pour un nombre fini d'indices
pour toute section $s\in\Gamma(U,\mathcal{S})$, $U$ ouvert dans $Y$). On peut
alors vérifier que les axiomes des ensembles munis d'un groupe d'opérateurs
entraînent que les $\psi_n(\mathcal{S})=\mathcal{S}_n$ définissent une
graduation (à degrés positifs ou négatifs) de $\mathcal{O}_Y$-Algèbre sur
$\mathcal{S}$, les $\psi_n$ étant les projecteurs correspondants. On a ainsi
une notion de structure de «~\emph{cône affine}~» sur tout $Y$-schéma affine,
définie de façon «~géométrique~» sans appel à une graduation préalable.

\oldpage[II]{168}
Nous ne développerons pas davantage ce point de vue ici et laissons au
lecteur le soin de formuler de façon précise les définitions et résultats qui
correspondent aux indications précédentes.
\end{remark}

\subsection{Fermeture projective d'un fibré vectoriel}
\label{subsection:II.8.4-fr}

\begin{env}[8.4.1]
\label{II.8.4.1-fr}
Soient $Y$ un préschéma, $\mathcal{E}$ un $\mathcal{O}_Y$-Module
quasi-cohérent. Si on prend pour $\mathcal{S}$ la $\mathcal{O}_Y$-Algèbre
graduée $\mathbf{S}_{\mathcal{O}_Y}(\mathcal{E})$, la définition
(\hyperref[II.8.3.1.1-fr]{8.3.1.1}) montre que
$\widehat{\mathcal{S}}$ s'identifie à
$\mathbf{S}_{\mathcal{O}_Y}(\mathcal{E}\oplus\mathcal{O}_Y)$. Le cône affine
$\operatorname{Spec}(\mathcal{S})$ défini par $\mathcal{S}$ étant alors par
définition $\mathbf{V}(\mathcal{E})$, et $\operatorname{Proj}(\mathcal{S})$
étant par définition $\mathbf{P}(\mathcal{E})$, on voit donc que~:
\end{env}

\begin{proposition}[8.4.2]
\label{II.8.4.2-fr}
La fermeture projective d'un fibré vectoriel $\mathbf{V}(\mathcal{E})$ sur
$Y$ est canoniquement isomorphe à
$\mathbf{P}(\mathcal{E}\oplus\mathcal{O}_Y)$, et le lieu à l'infini de cette
dernière est canoniquement isomorphe à $\mathbf{P}(\mathcal{E})$.
\end{proposition}

\begin{remark}[8.4.3]
\label{II.8.4.3-fr}
Prenons par exemple $\mathcal{E}=\mathcal{O}_Y^r$ avec $r\geq2$~; alors les
cônes épointés $E$, $\widehat{E}$ définis par $\mathcal{S}$ ne sont ni affines
ni projectifs sur $Y$ si $Y\neq\varnothing$. La seconde assertion est
immédiate, car $\widehat{C}=\mathbf{P}(\mathcal{O}_Y^{r+1})$ est projectif sur
$Y$ et les espaces sous-jacents à $E$ et $\widehat{E}$ sont des ouverts non
fermés dans $\widehat{C}$, donc les immersions canoniques
$E\to\widehat{E}$ et $\widehat{E}\to\widehat{C}$ ne sont pas projectives
(\hyperref[II.5.5.3-fr]{5.5.3}), et on conclut par
(\hyperref[II.5.5.5-fr]{5.5.5, (v)}). D'autre part, supposant
$Y=\operatorname{Spec}(A)$ affine et par exemple $r=2$, on a
$C=\operatorname{Spec}(A[T_1,T_2])$ et $E$ est alors le préschéma induit par
$C$ sur l'ouvert $D(T_1)\cup D(T_2)$~; or nous avons vu que ce dernier n'est
pas affine (\hyperref[I.5.5.11-fr]{I, 5.5.11})~; 
\emph{a fortiori} $\widehat{E}$ ne peut être affine, puisque $E$ est l'ouvert
où ne s'annule pas la section $z$ au-dessus de $\widehat{E}$
(\hyperref[II.8.3.2-fr]{8.3.2}).

Toutefois~:
\end{remark}

\begin{proposition}[8.4.4]
\label{II.8.4.4-fr}
Si $\mathcal{L}$ est un $\mathcal{O}_Y$-Module inversible, on a des
isomorphismes canoniques pour les cônes épointés $E$ et $\widehat{E}$
correspondants à $C=\mathbf{V}(\mathcal{L})$,
\begin{equation}
\label{II.8.4.4.1-fr}
  \operatorname{Spec}\left(
    \bigoplus_{n\in\mathbf{Z}}\mathcal{L}^{\otimes n}
  \right)\xrightarrow{\sim}E
\tag{8.4.4.1}
\end{equation}
\begin{equation}
\label{II.8.4.4.2-fr}
  \mathbf{V}(\mathcal{L}^{-1})\xrightarrow{\sim}\widehat{E}.
\tag{8.4.4.2}
\end{equation}

En outre, il existe un isomorphisme canonique de la fermeture projective de
$\mathbf{V}(\mathcal{L})$ sur celle de $\mathbf{V}(\mathcal{L}^{-1})$,
transformant la section nulle (resp. le lieu à l'infini) de la première en le
lieu à l'infini (resp. la section nulle) de la seconde.
\end{proposition}

\begin{proof}
On a ici
$\mathcal{S}=\bigoplus_{n\geq0}\mathcal{L}^{\otimes n}$~; l'injection
canonique
\[
  \mathcal{S}\to\bigoplus_{n\in\mathbf{Z}}\mathcal{L}^{\otimes n}
\]
définit un morphisme dominant canonique
\begin{equation}
\label{II.8.4.4.3-fr}
  \operatorname{Spec}\left(
    \bigoplus_{n\in\mathbf{Z}}\mathcal{L}^{\otimes n}
  \right)
  \to\mathbf{V}(\mathcal{L})
  =\operatorname{Spec}\left(
    \bigoplus_{n\geq0}\mathcal{L}^{\otimes n}
  \right)
\tag{8.4.4.3}
\end{equation}
et il suffit de prouver que ce morphisme est un isomorphisme du schéma
$\operatorname{Spec}(\bigoplus_{n\in\mathbf{Z}}\mathcal{L}^{\otimes n})$
sur $E$. La question étant locale sur $Y$, on peut supposer que
$Y=\operatorname{Spec}(A)$ est affine

\oldpage[II]{169}
et $\mathcal{L}=\mathcal{O}_Y$, donc $\mathcal{S}=(A[T])^{\sim}$ et
$\bigoplus_{n\in\mathbf{Z}}\mathcal{L}^{\otimes n}
=(A[T,T^{-1}])^{\sim}$. Or $A[T,T^{-1}]$ est l'anneau de fractions
$A[T]_T$ de $A[T]$, donc (\hyperref[II.8.4.4.3-fr]{8.4.4.3}) identifie le
premier membre au préschéma induit par $C=\mathbf{V}(\mathcal{L})$ sur
l'ouvert $D(T)$~; le complémentaire $V(T)$ de cet ouvert dans $C$ est
l'espace sous-jacent au sous-préschéma fermé de $C$ défini par l'idéal
$TA[T]$, c'est-à-dire la section nulle de $C$, donc $E=D(T)$.

L'isomorphisme (\hyperref[II.8.4.4.2-fr]{8.4.4.2}) sera d'autre part
conséquence de la dernière assertion, $\mathbf{V}(\mathcal{L}^{-1})$ étant le
complémentaire du lieu à l'infini de sa fermeture projective et $\widehat{E}$
le complémentaire de la section nulle de la fermeture projective de
$C=\mathbf{V}(\mathcal{L})$. Or, ces fermetures projectives sont
respectivement $\mathbf{P}(\mathcal{L}^{-1}\oplus\mathcal{O}_Y)$ et
$\mathbf{P}(\mathcal{L}\oplus\mathcal{O}_Y)$~; mais on peut écrire
$\mathcal{L}\oplus\mathcal{O}_Y
=\mathcal{L}\otimes(\mathcal{L}^{-1}\oplus\mathcal{O}_Y)$. L'existence de
l'isomorphisme canonique cherché résulte alors de
(\hyperref[II.4.1.4-fr]{4.1.4}), et tout revient à voir que cet isomorphisme
échange sections nulles et lieux à l'infini. On peut pour cela se borner au
cas où $Y=\operatorname{Spec}(A)$ est affine,
$L=Ac$, $L^{-1}=Ac'$, l'isomorphisme canonique
$L\otimes L^{-1}\to A$ faisant correspondre à $c\otimes c'$ l'élément $1$
de $A$. Alors $\mathbf{S}(L\oplus A)$ est produit tensoriel de $A[z]$ et de
$\bigoplus_{n\geq0}Ac^{\otimes n}$, $\mathbf{S}(L^{-1}\oplus A)$ produit
tensoriel de $A[z]$ et de $\bigoplus_{n\geq0}Ac'^{\otimes n}$, et
l'isomorphisme défini dans (\hyperref[II.4.1.4-fr]{4.1.4}) fait correspondre à
$z^h\otimes c'^{\otimes(n-h)}$ l'élément
$z^{n-h}\otimes c^{\otimes h}$. Or, dans
$\mathbf{P}(\mathcal{L}^{-1}\oplus\mathcal{O}_Y)$, le lieu à l'infini est
l'ensemble des points où s'annule la section $z$ et la section nulle est
l'ensemble des points où s'annule la section $c'$~; comme on a des définitions
analogues pour $\mathbf{P}(\mathcal{L}\oplus\mathcal{O}_Y)$, notre conclusion
résulte aussitôt de l'explicitation qui précède.
\end{proof}

\subsection{Comportements fonctoriels}
\label{subsection:II.8.5-fr}

\begin{env}[8.5.1]
\label{II.8.5.1-fr}
Soient $Y$, $Y'$ deux préschémas, $q:Y'\to Y$ un morphisme,
$\mathcal{S}$ (resp. $\mathcal{S}'$) une $\mathcal{O}_Y$-Algèbre
(resp. une $\mathcal{O}_{Y'}$-Algèbre) graduée quasi-cohérente à degrés
\emph{positifs}. Considérons un $q$-morphisme d'Algèbres graduées
\begin{equation}
\label{II.8.5.1.1-fr}
  \varphi:\mathcal{S}\to\mathcal{S}'.
\tag{8.5.1.1}
\end{equation}
On sait (\hyperref[II.1.5.6-fr]{1.5.6}) qu'il lui correspond canoniquement un
morphisme
\[
  \Phi=\operatorname{Spec}(\varphi):
  \operatorname{Spec}(\mathcal{S}')\to\operatorname{Spec}(\mathcal{S})
\]
tel que le diagramme
\begin{equation}
\label{II.8.5.1.2-fr}
  \xymatrix{
    C'\ar[r]^{\Phi}\ar[d] & C\ar[d]\\
    Y'\ar[r]_{q} & Y
  }
\tag{8.5.1.2}
\end{equation}
où on a posé $C=\operatorname{Spec}(\mathcal{S})$,
$C'=\operatorname{Spec}(\mathcal{S}')$, soit commutatif. 
\emph{Supposons en outre que
$\mathcal{S}_0=\mathcal{O}_Y$, $\mathcal{S}'_0=\mathcal{O}_{Y'}$}~; soient
$\varepsilon:Y\to C$ et $\varepsilon':Y'\to C'$ les immersions canoniques
(\hyperref[II.8.3.2-fr]{8.3.2})~; on a alors un diagramme commutatif
\begin{equation}
\label{II.8.5.1.3-fr}
  \xymatrix{
    Y'\ar[r]^{q}\ar[d]_{\varepsilon'} & Y\ar[d]^{\varepsilon}\\
    C'\ar[r]_{\Phi} & C
  }
\tag{8.5.1.3}
\end{equation}

\oldpage[II]{170}
qui correspond au diagramme
\[
  \xymatrix{
    \mathcal{S}\ar[r]^{\varphi}\ar[d] & \mathcal{S}'\ar[d]\\
    \mathcal{O}_Y\ar[r] & \mathcal{O}_{Y'}
  }
\]
où les flèches verticales sont les homomorphismes d'augmentation, et dont la
commutativité résulte de l'hypothèse que $\varphi$ est supposé être un
homomorphisme d'Algèbres \emph{graduées}.
\end{env}

\begin{proposition}[8.5.2]
\label{II.8.5.2-fr}
Si $E$ (resp. $E'$) est le cône épointé affine défini par $\mathcal{S}$
(resp. $\mathcal{S}'$), on a $\Phi^{-1}(E)\subset E'$~; si en outre
$\operatorname{Proj}(\varphi):G(\varphi)\to\operatorname{Proj}(\mathcal{S})$
est partout défini (autrement dit si
$G(\varphi)=\operatorname{Proj}(\mathcal{S}')$), alors on a
$\Phi^{-1}(E)=E'$, et réciproquement.
\end{proposition}

\begin{proof}
La première assertion résulte de la commutativité de
(\hyperref[II.8.5.1.3-fr]{8.5.1.3}). Pour démontrer la seconde, on peut se
borner au cas où $Y=\operatorname{Spec}(A)$ et
$Y'=\operatorname{Spec}(A')$ sont affines,
$\mathcal{S}=\widetilde{S}$, $\mathcal{S}'=\widetilde{S'}$. Pour tout $f$
homogène dans $S_+$, si on pose $f'=\varphi(f)$, on a
$\Phi^{-1}(C_f)=C'_{f'}$ (\hyperref[I.2.2.4.1-fr]{I, 2.2.4.1})~; dire que
$G(\varphi)=\operatorname{Proj}(S')$ signifie que dans $S'_+$ la racine de
l'idéal engendré par les $f'=\varphi(f)$ est $S'_+$ lui-même
(\hyperref[II.2.8.1-fr]{2.8.1}) et (\hyperref[II.2.3.14-fr]{2.3.14}), et cela
équivaut encore à dire que les $C'_{f'}$ recouvrent $E'$
(\hyperref[II.8.3.4.4-fr]{8.3.4.4}).
\end{proof}

\begin{env}[8.5.3]
\label{II.8.5.3-fr}
Le $q$-morphisme $\varphi$ se prolonge canoniquement en un $q$-morphisme
d'Algèbres graduées
\begin{equation}
\label{II.8.5.3.1-fr}
  \widehat{\varphi}:\widehat{\mathcal{S}}
  \to\widehat{\mathcal{S}'}
\tag{8.5.3.1}
\end{equation}
en posant $\widehat{\varphi}(z)=z$. On en déduit par suite un morphisme
\[
  \widehat{\Phi}=\operatorname{Proj}(\widehat{\varphi}):
  G(\widehat{\varphi})\to\widehat{C}
  =\operatorname{Proj}(\widehat{\mathcal{S}})
\]
tel que le diagramme
\[
  \xymatrix{
    G(\widehat{\varphi})\ar[r]^{\widehat{\Phi}}\ar[d] &
    \widehat{C}\ar[d]\\
    Y'\ar[r]_{q} & Y
  }
\]
soit commutatif (\hyperref[II.3.5.6-fr]{3.5.6}). Il résulte aussitôt des
définitions que si on désigne par $i:C\to\widehat{C}$ et
$i':C'\to\widehat{C'}$ les immersions ouvertes canoniques
(\hyperref[II.8.3.2-fr]{8.3.2}), on a
$i'(C')\subset G(\widehat{\varphi})$ et le diagramme
\begin{equation}
\label{II.8.5.3.2-fr}
  \xymatrix{
    C'\ar[r]^{\Phi}\ar[d]_{i'} & C\ar[d]^{i}\\
    G(\widehat{\varphi})\ar[r]_{\widehat{\Phi}} & \widehat{C}
  }
\tag{8.5.3.2}
\end{equation}
est commutatif. Enfin, si on pose
$X=\operatorname{Proj}(\mathcal{S})$, $X'=\operatorname{Proj}(\mathcal{S}')$,
et si $j:X\to\widehat{C}$, $j':X'\to\widehat{C'}$ sont les immersions
fermées canoniques (\hyperref[II.8.3.2-fr]{8.3.2}), il résulte encore des
définitions de ces immersions que l'on a
$j'(G(\varphi))\subset G(\widehat{\varphi})$ et que le diagramme

\oldpage[II]{171}
\begin{equation}
\label{II.8.5.3.3-fr}
  \xymatrix{
    G(\varphi)\ar[r]^{\operatorname{Proj}(\varphi)}\ar[d]_{j'} &
    X\ar[d]^{j}\\
    G(\widehat{\varphi})\ar[r]_{\widehat{\Phi}} & \widehat{C}
  }
\tag{8.5.3.3}
\end{equation}
est commutatif.
\end{env}

\begin{proposition}[8.5.4]
\label{II.8.5.4-fr}
Si $\widehat{E}$ (resp. $\widehat{E'}$) est le cône épointé projectif défini
par $\mathcal{S}$ (resp. $\mathcal{S}'$), on a
$\widehat{\Phi}^{-1}(\widehat{E})\subset\widehat{E'}$~; en outre, si
$p:\widehat{E}\to X$ et $p':\widehat{E'}\to X'$ sont les rétractions
canoniques, on a
$p'(\widehat{\Phi}^{-1}(\widehat{E}))\subset G(\varphi)$, et le diagramme
\begin{equation}
\label{II.8.5.4.1-fr}
  \xymatrix{
    \widehat{\Phi}^{-1}(\widehat{E})
      \ar[r]^{\widehat{\Phi}}\ar[d]_{p'} &
    \widehat{E}\ar[d]^{p}\\
    G(\varphi)\ar[r]_{\operatorname{Proj}(\varphi)} & X
  }
\tag{8.5.4.1}
\end{equation}
est commutatif. Si $\operatorname{Proj}(\varphi)$ est partout défini, il en
est de même de $\widehat{\Phi}$, et on a
$\widehat{\Phi}^{-1}(\widehat{E})=\widehat{E'}$.
\end{proposition}

\begin{proof}
La première assertion résulte de la commutativité des diagrammes
(\hyperref[II.8.5.1.3-fr]{8.5.1.3}) et
(\hyperref[II.8.5.3.2-fr]{8.5.3.2}), et les deux suivantes de la définition
des rétractions canoniques (\hyperref[II.8.3.5-fr]{8.3.5}) et de la définition
de $\widehat{\varphi}$. D'autre part, pour voir que $\widehat{\Phi}$ est
partout défini lorsqu'il en est ainsi de $\operatorname{Proj}(\varphi)$, on
peut se borner à considérer le cas où $Y=\operatorname{Spec}(A)$ et
$Y'=\operatorname{Spec}(A')$ sont affines,
$\mathcal{S}=\widetilde{S}$, $\mathcal{S}'=\widetilde{S'}$~; l'hypothèse est
que lorsque $f$ parcourt l'ensemble des éléments homogènes de $S_+$, l'idéal
engendré dans $S'_+$ par les $\varphi(f)$ a pour racine dans $S'_+$ l'idéal
$S'_+$ lui-même~; on en conclut aussitôt que la racine dans $(S'[z])_+$ de
l'idéal engendré par $z$ et les $\varphi(f)$ est $(S'[z])_+$ lui-même, d'où
notre assertion~; cela prouve en même temps que $\widehat{E'}$ est réunion des
$\widehat{C'}_{\varphi(f)}$, donc égal à
$\widehat{\Phi}^{-1}(\widehat{E})$.
\end{proof}

\begin{corollary}[8.5.5]
\label{II.8.5.5-fr}
Lorsque $\operatorname{Proj}(\varphi)$ est partout défini, l'image réciproque
par $\widehat{\Phi}$ de l'espace sous-jacent au lieu à l'infini (resp. au
préschéma des sommets) de $\widehat{C}$ est l'espace sous-jacent au lieu à
l'infini (resp. au préschéma des sommets) de $\widehat{C'}$.
\end{corollary}

\begin{proof}
Cela résulte aussitôt de (\hyperref[II.8.5.4-fr]{8.5.4}) et
(\hyperref[II.8.5.2-fr]{8.5.2}), compte tenu des relations
(\hyperref[II.8.3.4.1-fr]{8.3.4.1}) et
(\hyperref[II.8.3.4.2-fr]{8.3.4.2}).
\end{proof}

\subsection{Un isomorphisme canonique pour les cônes épointés}
\label{subsection:II.8.6-fr}

\begin{env}[8.6.1]
\label{II.8.6.1-fr}
Soient $Y$ un préschéma, $\mathcal{S}$ une $\mathcal{O}_Y$-Algèbre graduée
quasi-cohérente à degrés positifs \emph{telle que
$\mathcal{S}_0=\mathcal{O}_Y$}, $X$ le $Y$-schéma
$\operatorname{Proj}(\mathcal{S})$. Nous allons appliquer les résultats de
(\hyperref[subsection:II.8.5-fr]{8.5}) au cas où on prend $Y'=X$,
$q:X\to Y$ étant le morphisme structural~; soit
\begin{equation}
\label{II.8.6.1.1-fr}
  \mathcal{S}_X=\bigoplus_{n\in\mathbf{Z}}\mathcal{O}_X(n)
\tag{8.6.1.1}
\end{equation}

\oldpage[II]{172}
qui est une $\mathcal{O}_X$-Algèbre graduée quasi-cohérente, la multiplication
y étant définie au moyen des homomorphismes canoniques
(\hyperref[II.3.2.6.1-fr]{3.2.6.1})
\[
  \mathcal{O}_X(m)\otimes_{\mathcal{O}_X}\mathcal{O}_X(n)
  \to\mathcal{O}_X(m+n)
\]
dont l'associativité est garantie par le diagramme commutatif
(\hyperref[II.2.5.11.4-fr]{2.5.11.4}). Prenons pour $\mathcal{S}'$ la
sous-$\mathcal{O}_X$-Algèbre graduée quasi-cohérente
$\mathcal{S}_X^{\geq}=\bigoplus_{n\geq0}\mathcal{O}_X(n)$ de
$\mathcal{S}_X$, à degrés positifs.

Enfin, nous considérerons le $q$-morphisme canonique
\begin{equation}
\label{II.8.6.1.2-fr}
  \alpha:\mathcal{S}\to\mathcal{S}_X^{\geq}
\tag{8.6.1.2}
\end{equation}
défini dans (\hyperref[II.3.3.2.3-fr]{3.3.2.3}) comme un homomorphisme
$\mathcal{S}\to q_*(\mathcal{S}_X)$, mais qui applique évidemment
$\mathcal{S}$ dans $q_*(\mathcal{S}_X^{\geq})$. Nous poserons
\begin{equation}
\label{II.8.6.1.3-fr}
  C_X=\operatorname{Spec}(\mathcal{S}_X^{\geq}),\qquad
  \widehat{C}_X=\operatorname{Proj}(\mathcal{S}_X^{\geq}[z]),\qquad
  X'=\operatorname{Proj}(\mathcal{S}_X^{\geq})
\tag{8.6.1.3}
\end{equation}
et nous désignerons par $E_X$ et $\widehat{E}_X$ le cône épointé affine et le
cône épointé projectif correspondants~; on notera
$\varepsilon_X:X\to C_X$, $i_X:C_X\to\widehat{C}_X$,
$j_X:X'\to\widehat{C}_X$, $p_X:\widehat{E}_X\to X'$,
$\pi_X:E_X\to X'$ les morphismes canoniques définis dans
(\hyperref[subsection:II.8.3-fr]{8.3}).
\end{env}

\begin{proposition}[8.6.2]
\label{II.8.6.2-fr}
Le morphisme structural $u:X'\to X$ est un isomorphisme, et le morphisme
$\operatorname{Proj}(\alpha)$ est partout défini et identique à $u$. Le
morphisme
$\operatorname{Proj}(\widehat{\alpha}):\widehat{C}_X\to\widehat{C}$ est
partout défini, et ses restrictions à $\widehat{E}_X$ et $E_X$ sont des
isomorphismes sur $\widehat{E}$ et $E$ respectivement. Enfin, lorsqu'on
identifie $X'$ et $X$ au moyen de $u$, les morphismes $p_X$ et $\pi_X$
s'identifient aux morphismes structuraux des $X$-préschémas $\widehat{E}_X$ et
$E_X$.
\end{proposition}

\begin{proof}
On peut évidemment se borner au cas où $Y=\operatorname{Spec}(A)$ est affine
et $\mathcal{S}=\widetilde{S}$~; alors $X$ est réunion des ouverts affines
$X_f$, où $f$ parcourt l'ensemble des éléments homogènes de $S_+$, l'anneau
de $X_f$ étant $S_{(f)}$. Il résulte en outre de
(\hyperref[II.8.2.7.1-fr]{8.2.7.1}) que l'on a
\begin{equation}
\label{II.8.6.2.1-fr}
  \Gamma(X_f,\mathcal{S}_X^{\geq})=S_f^{\geq}.
\tag{8.6.2.1}
\end{equation}

On a donc $u^{-1}(X_f)=\operatorname{Proj}(S_f^{\geq})$. Mais si
$f\in S_d$ ($d>0$), $\operatorname{Proj}(S_f^{\geq})$ est canoniquement
isomorphe à $\operatorname{Proj}((S_f^{\geq})^{(d)})$
(\hyperref[II.2.4.7-fr]{2.4.7}), et d'autre part
$(S_f^{\geq})^{(d)}=(S^{(d)})_f^{\geq}$ s'identifie à $S_{(f)}[T]$
(\hyperref[II.2.2.1-fr]{2.2.1}) par l'application $T\to f/1$~; on en conclut
(\hyperref[II.3.1.7-fr]{3.1.7}) que le morphisme structural
$u^{-1}(X_f)\to X_f$ est un isomorphisme, d'où la première assertion. Pour
démontrer la seconde, remarquons que la restriction
$u^{-1}(X_f)\cap G(\alpha)\to X=\operatorname{Proj}(S)$ de
$\operatorname{Proj}(\alpha)$ correspond à l'application canonique
$x\to x/1$ de $S$ dans $S_f^{\geq}$
(\hyperref[II.2.6.2-fr]{2.6.2})~; on en déduit tout d'abord que
$G(\alpha)=X'$, puis, en tenant compte de ce que
$u^{-1}(X_f)=(u^{-1}(X_f))_{f/1}$, il résulte de
(\hyperref[II.2.8.1.1-fr]{2.8.1.1}) que l'image de $u^{-1}(X_f)$ par
$\operatorname{Proj}(\alpha)$ est contenue dans $X_f$, et la restriction de
$\operatorname{Proj}(\alpha)$ à $u^{-1}(X_f)$, considérée comme morphisme
dans $X_f=\operatorname{Spec}(S_{(f)})$, est bien identique à celle de $u$.
Enfin, la formule (\hyperref[II.8.3.5.4-fr]{8.3.5.4}), appliquée à $p_X$ au
lieu de $p$, montre que
$p_X^{-1}(u^{-1}(X_f))=
\operatorname{Spec}((S_f^{\geq})_{f/1}^{\leq})$, et cet ouvert est, d'après
(\hyperref[II.8.5.4.1-fr]{8.5.4.1}), l'image réciproque par
$\operatorname{Proj}(\widehat{\alpha})$ de
$p^{-1}(X_f)=\operatorname{Spec}(S_f^{\leq})$
(\hyperref[II.8.3.5.3-fr]{8.3.5.3}). Compte tenu de
(\hyperref[II.8.2.3.2-fr]{8.2.3.2}), la restriction de
$\operatorname{Proj}(\widehat{\alpha})$ à
$p_X^{-1}(u^{-1}(X_f))$ correspond à l'isomorphisme réciproque de
(\hyperref[II.8.2.7.2-fr]{8.2.7.2}), restreint à $S_f^{\leq}$, d'où le
troisième point~; la dernière assertion est évidente par définition.

\oldpage[II]{173}
On notera aussi qu'il résulte du diagramme commutatif
(\hyperref[II.8.5.3.2-fr]{8.5.3.2}) que \emph{la restriction à $C_X$ de
$\operatorname{Proj}(\widehat{\alpha})$ n'est autre que le morphisme
$\operatorname{Spec}(\alpha)$}.
\end{proof}

\begin{corollary}[8.6.3]
\label{II.8.6.3-fr}
Considérés comme $X$-schémas, $\widehat{E}_X$ est canoniquement isomorphe à
$\operatorname{Spec}(\mathcal{S}_X^{\leq})$, et $E_X$ à
$\operatorname{Spec}(\mathcal{S}_X)$.
\end{corollary}

\begin{proof}
Comme on sait que les morphismes $p_X$ et $\pi_X$ sont affines
(\hyperref[II.8.3.5-fr]{8.3.5} et
\hyperref[II.8.3.6-fr]{8.3.6}), il suffit (vu
(\hyperref[II.1.3.1-fr]{1.3.1})) de vérifier le corollaire lorsque
$Y=\operatorname{Spec}(A)$ est affine et $\mathcal{S}=\widetilde{S}$. La
première assertion résulte de l'existence de l'isomorphisme canonique
(\hyperref[II.8.2.7.2-fr]{8.2.7.2})
$(S_f^{\geq})_{f/1}^{\leq}\xrightarrow{\sim}S_f^{\leq}$ et du fait que ces
isomorphismes sont compatibles avec le passage de $f$ à $fg$ ($f$ et $g$
homogènes dans $S_+$). De même, la formule
(\hyperref[II.8.3.6.2-fr]{8.3.6.2}), appliquée à $\pi_X$ au lieu de $\pi$,
montre que
$\pi_X^{-1}(u^{-1}(X_f))=
\operatorname{Spec}((S_f^{\geq})_{f/1})$ pour $f$ homogène dans $S_+$, et la
seconde assertion découle de l'existence de l'isomorphisme canonique
(\hyperref[II.8.2.7.2-fr]{8.2.7.2})
$(S_f^{\geq})_{f/1}\xrightarrow{\sim}S_f$.

On peut ainsi dire que $\widehat{C}_X$, considéré comme $X$-schéma, s'obtient
par \emph{recollement} des deux $X$-schémas affines sur $X$,
$C_X=\operatorname{Spec}(\mathcal{S}_X^{\geq})$ et
$\widehat{E}_X=\operatorname{Spec}(\mathcal{S}_X^{\leq})$, dont
l'intersection est l'ouvert $E_X=\operatorname{Spec}(\mathcal{S}_X)$.
\end{proof}

\begin{corollary}[8.6.4]
\label{II.8.6.4-fr}
Supposons que $\mathcal{O}_X(1)$ soit un $\mathcal{O}_X$-Module inversible et
que $\mathcal{S}_X$ soit isomorphe à
$\bigoplus_{n\in\mathbf{Z}}(\mathcal{O}_X(1))^{\otimes n}$ (ce qui sera en
particulier le cas si $\mathcal{S}$ est engendrée par $\mathcal{S}_1$
(\hyperref[II.3.2.5-fr]{3.2.5} et
\hyperref[II.3.2.7-fr]{3.2.7})). Alors le cône projectif épointé
$\widehat{E}$ s'identifie au fibré vectoriel de rang un
$\mathbf{V}(\mathcal{O}_X(-1))$ sur $X$, et le cône affine épointé $E$ au
sous-préschéma de ce fibré vectoriel induit sur le complémentaire de la
section nulle. Avec cette identification, la rétraction canonique
$\widehat{E}\to X$ s'identifie au morphisme structural du $X$-schéma
$\mathbf{V}(\mathcal{O}_X(-1))$. Enfin, il existe un $Y$-morphisme canonique
$\mathbf{V}(\mathcal{O}_X(1))\to C$, dont la restriction au complémentaire de
la section nulle de $\mathbf{V}(\mathcal{O}_X(1))$ est un isomorphisme de ce
complémentaire sur le cône affine épointé $E$.
\end{corollary}

\begin{proof}
En effet, si on pose $\mathcal{L}=\mathcal{O}_X(1)$,
$\mathcal{S}_X^{\geq}$ est alors identique à
$\mathbf{S}_{\mathcal{O}_X}(\mathcal{L})$, donc $\widehat{E}_X$ s'identifie
canoniquement à $\mathbf{V}(\mathcal{L}^{-1})$ en vertu de
(\hyperref[II.8.6.3-fr]{8.6.3}) et $C_X$ à $\mathbf{V}(\mathcal{L})$. Le
morphisme $\mathbf{V}(\mathcal{L})\to C$ est la restriction de
$\operatorname{Proj}(\widehat{\alpha})$, et les assertions du corollaire sont
donc des cas particuliers de (\hyperref[II.8.6.2-fr]{8.6.2}).
\end{proof}

On notera que l'image réciproque par le morphisme
$\mathbf{V}(\mathcal{O}_X(1))\to C$ de l'espace sous-jacent au préschéma des
sommets de $C$ est l'espace sous-jacent à la section nulle de
$\mathbf{V}(\mathcal{O}_X(1))$ (\hyperref[II.8.5.5-fr]{8.5.5})~; mais en
général les sous-préschémas correspondants de $C$ et de
$\mathbf{V}(\mathcal{O}_X(1))$ ne sont pas isomorphes. Cette question va être
étudiée ci-dessous.

\subsection{Éclatement de cônes projetants}
\label{subsection:II.8.7-fr}

\begin{env}[8.7.1]
\label{II.8.7.1-fr}
Sous les conditions de (\hyperref[II.8.6.1-fr]{8.6.1}), on a, en posant
$r=\operatorname{Proj}(\widehat{\alpha})$, un diagramme commutatif
\begin{equation}
\label{II.8.7.1.1-fr}
  \xymatrix{
    X\ar[r]^{i_X\circ\varepsilon_X}\ar[d]_{q} &
    \widehat{C}_X\ar[d]^{r}\\
    Y\ar[r]_{i\circ\varepsilon} & \widehat{C}
  }
\tag{8.7.1.1}
\end{equation}

\oldpage[II]{174}
en vertu de (\hyperref[II.8.5.1.3-fr]{8.5.1.3}) et
(\hyperref[II.8.5.3.2-fr]{8.5.3.2})~; en outre, la restriction de $r$ au
complémentaire $\widehat{C}_X-i_X(\varepsilon_X(X))$ de la section nulle est
un \emph{isomorphisme} sur le complémentaire
$\widehat{C}-i(\varepsilon(Y))$ de la section nulle en vertu de
(\hyperref[II.8.6.2-fr]{8.6.2}). Si on suppose pour simplifier $Y$ affine,
$\mathcal{S}$ de type fini et engendré par $\mathcal{S}_1$, $X$ est projectif
sur $Y$ et $\widehat{C}_X$ projectif sur $X$
(\hyperref[II.5.5.1-fr]{5.5.1}), donc $\widehat{C}_X$ est projectif sur $Y$
(\hyperref[II.5.5.5-fr]{5.5.5, (ii)}), et \emph{a fortiori} sur
$\widehat{C}$ (\hyperref[II.5.5.5-fr]{5.5.5, (v)}). On a ainsi un
$Y$-morphisme projectif $r:\widehat{C}_X\to\widehat{C}$ (dont la restriction
à $C_X$ est un $Y$-morphisme projectif $C_X\to C$) qui \emph{contracte $X$
dans $Y$} et induit un \emph{isomorphisme} quand on le restreint aux
\emph{complémentaires de $X$ et de $Y$}. On a donc une relation entre $C_X$
et $C$, analogue à celle qui a lieu entre un préschéma éclaté et son préschéma
de départ (\hyperref[II.8.1.3-fr]{8.1.3}). Nous allons effectivement montrer
qu'on peut identifier $C_X$ au spectre homogène d'une
$\mathcal{O}_C$-Algèbre graduée.
\end{env}

\begin{env}[8.7.2]
\label{II.8.7.2-fr}
Gardant les notations de (\hyperref[II.8.6.1-fr]{8.6.1}), considérons pour
chaque $n\geqslant 0$, l'Idéal quasi-cohérent
\begin{equation}
\label{II.8.7.2.1-fr}
  \mathcal{S}_{[n]}=\bigoplus_{m\geqslant n}\mathcal{S}_m
\tag{8.7.2.1}
\end{equation}
de la $\mathcal{O}_Y$-Algèbre graduée $\mathcal{S}$. Il est clair que l'on a
\begin{equation}
\label{II.8.7.2.2-fr}
  \mathcal{S}_{[0]}=\mathcal{S},\qquad
  \mathcal{S}_{[n]}\subset\mathcal{S}_{[m]}
  \qquad\text{pour }m\leqslant n
\tag{8.7.2.2}
\end{equation}
\begin{equation}
\label{II.8.7.2.3-fr}
  \mathcal{S}_{[n]}\mathcal{S}_{[m]}\subset\mathcal{S}_{[m+n]}.
\tag{8.7.2.3}
\end{equation}

Considérons le $\mathcal{O}_C$-Module associé à $\mathcal{S}_{[n]}$, qui est un
Idéal quasi-cohérent de $\mathcal{O}_C=\widetilde{\mathcal{S}}$
(\hyperref[II.1.4.4-fr]{1.4.4})
\begin{equation}
\label{II.8.7.2.4-fr}
  \mathcal{I}_n=(\mathcal{S}_{[n]})^{\sim}.
\tag{8.7.2.4}
\end{equation}

On déduit alors de (\hyperref[II.8.7.2.2-fr]{8.7.2.2}) et
(\hyperref[II.8.7.2.3-fr]{8.7.2.3}), utilisant
(\hyperref[II.1.4.4-fr]{1.4.4}) et (\hyperref[II.1.4.8.1-fr]{1.4.8.1}),
les formules analogues
\begin{equation}
\label{II.8.7.2.5-fr}
  \mathcal{I}_0=\mathcal{O}_C,\qquad
  \mathcal{I}_n\subset\mathcal{I}_m
  \qquad\text{pour }m\leqslant n
\tag{8.7.2.5}
\end{equation}
\begin{equation}
\label{II.8.7.2.6-fr}
  \mathcal{I}_n\mathcal{I}_m\subset\mathcal{I}_{n+m}.
\tag{8.7.2.6}
\end{equation}

On est donc dans les conditions de (\hyperref[II.8.1.1-fr]{8.1.1}), ce qui
conduit à introduire la $\mathcal{O}_C$-Algèbre graduée quasi-cohérente
\begin{equation}
\label{II.8.7.2.7-fr}
  \mathcal{S}^{\natural}
  =\bigoplus_{n\geqslant 0}\mathcal{I}_n
  =\left(\bigoplus_{n\geqslant 0}\mathcal{S}_{[n]}\right)^{\sim}.
\tag{8.7.2.7}
\end{equation}
\end{env}

\begin{proposition}[8.7.3]
\label{II.8.7.3-fr}
Il existe un $C$-isomorphisme canonique
\begin{equation}
\label{II.8.7.3.1-fr}
  h:C_X\xrightarrow{\sim}\operatorname{Proj}(\mathcal{S}^{\natural}).
\tag{8.7.3.1}
\end{equation}
\end{proposition}

\begin{proof}
Supposons d'abord $Y=\operatorname{Spec}(A)$ affine, d'où
$\mathcal{S}=\widetilde{S}$, où $S$ est une $A$-algèbre graduée à degrés
positifs et $C=\operatorname{Spec}(S)$. La définition
(\hyperref[II.8.2.6.2-fr]{8.2.7.4}) montre que l'on a alors, avec les
notations de (\hyperref[II.8.2.6-fr]{8.2.6}),
$\mathcal{S}^{\natural}=(S^{\natural})^{\sim}$. Pour définir
(\hyperref[II.8.7.3.1-fr]{8.7.3.1}), considérons un élément homogène
$f\in S_d$ ($d>0$) et l'élément correspondant $f^{\natural}\in S^{\natural}$
(\hyperref[II.8.2.6-fr]{8.2.6})~; le $S$-isomorphisme
(\hyperref[II.8.2.7.3-fr]{8.2.7.3}) définit donc un $C$-isomorphisme
\begin{equation}
\label{II.8.7.3.2-fr}
  \operatorname{Spec}(S_f^{\geqslant})
  \xrightarrow{\sim}
  \operatorname{Spec}(S_{(f^{\natural})}^{\natural}).
\tag{8.7.3.2}
\end{equation}

\oldpage[II]{175}
Mais avec les notations de (\hyperref[II.8.6.2-fr]{8.6.2}), si
$v:C_X\to X$ est le morphisme structural, il résulte de
(\hyperref[II.8.6.2.1-fr]{8.6.2.1}) que l'on a
$v^{-1}(X_f)=\operatorname{Spec}(S_f^{\geqslant})$. On a d'autre part
$\operatorname{Spec}(S_{(f^{\natural})}^{\natural})=D_+(f^{\natural})$, si
bien que (\hyperref[II.8.7.3.2-fr]{8.7.3.2}) définit un isomorphisme
$v^{-1}(X_f)\to D_+(f^{\natural})$. En outre, si $g\in S_e$ ($e>0$), le
diagramme
\[
  \xymatrix{
    v^{-1}(X_{fg})\ar[r]^{\sim}\ar[d] &
    D_+(f^{\natural}g^{\natural})\ar[d]\\
    v^{-1}(X_f)\ar[r]^{\sim} & D_+(f^{\natural})
  }
\]
est commutatif, comme il résulte aussitôt de la définition de l'isomorphisme
(\hyperref[II.8.2.7.3-fr]{8.2.7.3}). Enfin par définition $S_+$ est engendré
par les $f$ homogènes, donc il résulte de
(\hyperref[II.8.2.10-fr]{8.2.10, (iv)}) et de
(\hyperref[II.2.3.14-fr]{2.3.14}) que les $D_+(f^{\natural})$ forment un
recouvrement de $\operatorname{Proj}(S^{\natural})$ et les $v^{-1}(X_f)$ un
recouvrement de $C_X$, puisque les $X_f$ forment un recouvrement de $X$~;
cela achève donc de définir dans ce cas l'isomorphisme
(\hyperref[II.8.7.3.1-fr]{8.7.3.1}).

Pour démontrer (\hyperref[II.8.7.3-fr]{8.7.3}) dans le cas général, il suffit
de voir que si $U$, $U'$ sont deux ouverts affines de $Y$ tels que
$U'\subset U$, d'anneaux $A$ et $A'$, et si on pose
$\mathcal{S}|U=\widetilde{S}$, $\mathcal{S}|U'=\widetilde{S'}$, le diagramme
\begin{equation}
\label{II.8.7.3.3-fr}
  \xymatrix{
    C_{U'}\ar[r]\ar[d] & \operatorname{Proj}(S'^{\natural})\ar[d]\\
    C_U\ar[r] & \operatorname{Proj}(S^{\natural})
  }
\tag{8.7.3.3}
\end{equation}
est commutatif. Mais $S'$ s'identifie canoniquement à $S\otimes_A A'$, donc
$S'^{\natural}$ à
\[
  S^{\natural}\otimes_S S'=S^{\natural}\otimes_A A'~;
\]
on a donc
$\operatorname{Proj}(S'^{\natural})=\operatorname{Proj}(S^{\natural})
\times_U U'$ (\hyperref[II.2.8.10-fr]{2.8.10})~; de même, si
$X=\operatorname{Proj}(S)$, $X'=\operatorname{Proj}(S')$, on a
$X'=X\times_U U'$ et
$\mathcal{S}_{X'}=\mathcal{S}_X\otimes_{\mathcal{O}_U}\mathcal{O}_{U'}$
(\hyperref[II.3.5.4-fr]{3.5.4}), autrement dit
$\mathcal{S}_{X'}=j^*(\mathcal{S}_X)$, où $j$ est la projection $X'\to X$.
On a par suite (\hyperref[II.1.5.2-fr]{1.5.2})
$C_{U'}=C_U\times_X X'=C_U\times_U U'$, et la commutativité de
(\hyperref[II.8.7.3.3-fr]{8.7.3.3}) est alors immédiate.
\end{proof}

\begin{remark}[8.7.4]
\label{II.8.7.4-fr}
\begin{enumerate}
  \item[\rm{(i)}] La fin du raisonnement de
  (\hyperref[II.8.7.3-fr]{8.7.3}) se généralise aussitôt de la façon
  suivante. Soient $g:Y'\to Y$ un morphisme,
  $\mathcal{S}'=g^*(\mathcal{S})$, $X'=\operatorname{Proj}(\mathcal{S}')$~;
  on a alors un diagramme commutatif
  \begin{equation}
  \label{II.8.7.4.1-fr}
    \xymatrix{
      C_{X'}\ar[r]\ar[d] &
      \operatorname{Proj}(\mathcal{S}'^{\natural})\ar[d]\\
      C_X\ar[r] & \operatorname{Proj}(\mathcal{S}^{\natural})
    }
  \tag{8.7.4.1}
  \end{equation}

  D'autre part, soit $\varphi:\mathcal{S}''\to\mathcal{S}$ un homomorphisme
  de $\mathcal{O}_Y$-Algèbres graduées, tel que si on pose
  $X''=\operatorname{Proj}(\mathcal{S}'')$,
  $u=\operatorname{Proj}(\varphi):X\to X''$ soit partout défini~; on a
  d'autre part

  \oldpage[II]{176}
  un $Y$-morphisme $v:C\to C''$ (avec
  $C''=\operatorname{Spec}(\mathcal{S}'')$) tel que
  $\mathcal{A}(v)=\varphi$, et comme $\varphi$ est un homomorphisme
  d'Algèbres graduées, on déduit de $\varphi$ un $v$-morphisme d'Algèbres
  graduées
  $\psi:\mathcal{S}''^{\natural}\to\mathcal{S}^{\natural}$
  (\hyperref[II.1.4.1-fr]{1.4.1}). En outre, il résulte aisément de
  (\hyperref[II.8.2.10-fr]{8.2.10, (iv)}) et de l'hypothèse sur $\varphi$,
  que $\operatorname{Proj}(\psi)$ est partout défini. Enfin, compte tenu de
  (\hyperref[II.3.5.6.1-fr]{3.5.6.1}), on a un $u$-morphisme canonique
  $\mathcal{S}_{X''}\to\mathcal{S}_X$, d'où
  (\hyperref[II.1.5.6-fr]{1.5.6}) un morphisme $w:C_{X''}\to C_X$. Cela
  étant, le diagramme
  \begin{equation}
  \label{II.8.7.4.2-fr}
    \xymatrix{
      C_{X''}\ar[r]^{\sim}\ar[d]_{w} &
      \operatorname{Proj}(\mathcal{S}''^{\natural})
        \ar[d]^{\operatorname{Proj}(\psi)}\\
      C_X\ar[r]^{\sim} & \operatorname{Proj}(\mathcal{S}^{\natural})
    }
  \tag{8.7.4.2}
  \end{equation}
  est commutatif, comme on le vérifie aussitôt en se ramenant au cas où $Y$
  est affine.

  \item[\rm{(ii)}] Notons qu'en vertu de
  (\hyperref[II.8.7.2.5-fr]{8.7.2.5}) et
  (\hyperref[II.8.7.2.6-fr]{8.7.2.6}), on a
  $\mathcal{I}_1^m\subset\mathcal{I}_m\subset\mathcal{I}_1$ pour tout $m>0$.
  Or, par définition, $\mathcal{I}_1=(\mathcal{S}_+)^{\sim}$, donc
  $\mathcal{I}_1$ définit dans $C$ le sous-préschéma fermé
  $\varepsilon(Y)$ ((\hyperref[II.1.4.10-fr]{1.4.10}) et
  (\hyperref[II.8.3.2-fr]{8.3.2}))~; on en conclut que pour tout $m>0$,
  \emph{le support de $\mathcal{O}_C/\mathcal{I}_m$ est contenu dans l'espace
  sous-jacent au préschéma des sommets $\varepsilon(Y)$}~; dans l'image
  réciproque du cône épointé affine $E$, le morphisme structural
  $\operatorname{Proj}(\mathcal{S}^{\natural})\to C$ se réduit donc à un
  \emph{isomorphisme} (comme il résulte d'ailleurs de
  (\hyperref[II.8.7.3-fr]{8.7.3}) et
  (\hyperref[II.8.7.1-fr]{8.7.1})). En outre, identifiant canoniquement $C$ à
  un ouvert de $\widehat{C}$ (\hyperref[II.8.3.3-fr]{8.3.3}), on peut
  évidemment prolonger les Idéaux $\mathcal{I}_m$ de $\mathcal{O}_C$ en des
  Idéaux $\mathcal{J}_m$ de $\mathcal{O}_{\widehat{C}}$, par la condition de
  coïncider avec $\mathcal{O}_{\widehat{C}}$ dans l'ouvert $\widehat{E}$ de
  $\widehat{C}$. Si on pose
  $\mathcal{T}=\bigoplus_{n\geqslant 0}\mathcal{J}_m$, qui est une
  $\mathcal{O}_{\widehat{C}}$-Algèbre graduée quasi-cohérente, on peut
  prolonger l'isomorphisme (\hyperref[II.8.7.3.1-fr]{8.7.3.1}) en un
  $\widehat{C}$-isomorphisme
  \begin{equation}
  \label{II.8.7.4.3-fr}
    \widehat{C}_X\xrightarrow{\sim}\operatorname{Proj}(\mathcal{T}).
  \tag{8.7.4.3}
  \end{equation}

  En effet, au-dessus de $\widehat{E}$, il résulte de ce qui précède que
  $\operatorname{Proj}(\mathcal{T})$ s'identifie canoniquement à
  $\widehat{E}$, et on définit donc l'isomorphisme
  (\hyperref[II.8.7.4.3-fr]{8.7.4.3}) au-dessus de $\widehat{E}$ en lui
  imposant de coïncider avec l'isomorphisme canonique
  $\widehat{E}_X\to\widehat{E}$ (\hyperref[II.8.6.2-fr]{8.6.2})~; il est clair
  qu'alors cet isomorphisme et (\hyperref[II.8.7.3.1-fr]{8.7.3.1}) coïncident
  au-dessus de $E$.
\end{enumerate}
\end{remark}

\begin{corollary}[8.7.5]
\label{II.8.7.5-fr}
Supposons qu'il existe $n_0>0$ tel que
\begin{equation}
\label{II.8.7.5.1-fr}
  \mathcal{S}_{n+1}=\mathcal{S}_1\mathcal{S}_n
  \qquad\text{pour }n\geqslant n_0.
\tag{8.7.5.1}
\end{equation}
Alors le sous-préschéma des sommets de $C_X$ (isomorphe à $X$) est l'image
réciproque par le morphisme canonique $r:C_X\to C$ du sous-préschéma des
sommets de $C$ (isomorphe à $Y$). Inversement, si cette propriété a lieu et
si on suppose de plus $Y$ noethérien et $\mathcal{S}$ de type fini, il existe
$n_0>0$ tel que l'on ait (\hyperref[II.8.7.5.1-fr]{8.7.5.1}).
\end{corollary}

\begin{proof}
La première assertion étant locale sur $Y$, on peut pour l'établir supposer
$Y=\operatorname{Spec}(A)$ affine, et $\mathcal{S}=\widetilde{S}$, où $S$ est
une $A$-algèbre graduée à degrés positifs. Elle résulte alors de
(\hyperref[II.8.2.12-fr]{8.2.12}), car on a
$\operatorname{Proj}(S^{\natural}\otimes_S S_0)=
C_X\times_C\varepsilon(Y)$ (en vertu de l'identification
(\hyperref[II.8.7.3.1-fr]{8.7.3.1})), c'est-à-dire que ce préschéma est
l'image réciproque de $\varepsilon(Y)$ dans $C_X$
(\hyperref[I.4.4.1-fr]{I, 4.4.1}). La réciproque résulte aussi de
(\hyperref[II.8.2.12-fr]{8.2.12}) lorsque $Y$ est affine noethérien et $S$ de
type fini.

\oldpage[II]{177}
Si $Y$ est noethérien (non nécessairement affine) et $\mathcal{S}$ de type
fini, il y a un recouvrement fini de $Y$ par des ouverts affines noethériens
$U_i$, et on déduit alors de ce qui précède que pour tout $i$, il y a un
entier $n_i$ tel que
$\mathcal{S}_{n+1}|U_i=(\mathcal{S}_1|U_i)(\mathcal{S}_n|U_i)$ dès que
$n\geqslant n_i$~; le plus grand des $n_i$ vérifie donc
(\hyperref[II.8.7.5.1-fr]{8.7.5.1}).
\end{proof}

\begin{env}[8.7.6]
\label{II.8.7.6-fr}
Considérons maintenant le $C$-préschéma $Z$ obtenu en faisant \emph{éclater}
dans le cône affine $C$ le \emph{sous-préschéma des sommets}
$\varepsilon(Y)$~; par définition (\hyperref[II.8.1.3-fr]{8.1.3}) c'est donc
le préschéma
$\operatorname{Proj}(\bigoplus_{n\geqslant 0}\mathcal{S}_+^n)$~; l'injection
canonique
\begin{equation}
\label{II.8.7.6.1-fr}
  \iota:\bigoplus_{n\geqslant 0}\mathcal{S}_+^n
  \longrightarrow\mathcal{S}^{\natural}
\tag{8.7.6.1}
\end{equation}
définit (grâce à l'identification (\hyperref[II.8.7.3-fr]{8.7.3})) un
$C$-morphisme dominant canonique
\begin{equation}
\label{II.8.7.6.2-fr}
  G(\iota)\longrightarrow Z
\tag{8.7.6.2}
\end{equation}
où $G(\iota)$ est un ouvert de $C_X$ (\hyperref[II.3.5.1-fr]{3.5.1})~; on
notera qu'on peut avoir $G(\iota)\ne C_X$, comme le montre l'exemple où
$Y=\operatorname{Spec}(K)$, $K$ étant un corps, $\mathcal{S}=\widetilde{S}$
avec $S=K[\mathbf{y}]$, où $\mathbf{y}$ est une indéterminée \emph{de degré
2}~; si $R_n$ désigne l'ensemble $(S_+)^n$, considéré comme partie de
$S_{[n]}=S_n^{\natural}$, $S_+^{\natural}$ n'est pas la racine dans
$S^{\natural}$ de l'idéal engendré par la réunion des $R_n$ (cf.
(\hyperref[II.2.3.14-fr]{2.3.14})).
\end{env}

\begin{corollary}[8.7.7]
\label{II.8.7.7-fr}
Supposons qu'il existe $n_0>0$ tel que
\begin{equation}
\label{II.8.7.7.1-fr}
  \mathcal{S}_n=\mathcal{S}_1^n
  \qquad\text{pour }n\geqslant n_0.
\tag{8.7.7.1}
\end{equation}
Alors le morphisme canonique (\hyperref[II.8.7.6.2-fr]{8.7.6.2}) est partout
défini et est un isomorphisme $C_X\xrightarrow{\sim}Z$. Inversement, si cette
propriété a lieu et si on suppose de plus $Y$ noethérien et $\mathcal{S}$ de
type fini, il existe $n_0$ tel que l'on ait
(\hyperref[II.8.7.7.1-fr]{8.7.7.1}).
\end{corollary}

\begin{proof}
En effet, la première assertion est locale sur $Y$, et découle donc de
(\hyperref[II.8.2.14-fr]{8.2.14})~; il en est de même de la réciproque, en
raisonnant comme dans (\hyperref[II.8.7.5-fr]{8.7.5}).
\end{proof}

\begin{remark}[8.7.8]
\label{II.8.7.8-fr}
Comme la condition (\hyperref[II.8.7.7.1-fr]{8.7.7.1}) implique
(\hyperref[II.8.7.5.1-fr]{8.7.5.1}), on voit que lorsqu'elle est vérifiée,
non seulement $C_X$ s'identifie au préschéma obtenu en faisant éclater le
sommet (identifié à $Y$) du cône affine $C$, mais encore le sommet (identifié
à $X$) de $C_X$ s'identifie au sous-préschéma fermé image réciproque du
sommet $Y$ de $C$. En outre, l'hypothèse
(\hyperref[II.8.7.7.1-fr]{8.7.7.1}) implique que sur
$X=\operatorname{Proj}(\mathcal{S})$, les $\mathcal{O}_X$-Modules
$\mathcal{O}_X(n)$ sont inversibles (\hyperref[II.3.2.5-fr]{3.2.5} et
\hyperref[II.3.2.9-fr]{3.2.9}) et que l'on a
$\mathcal{O}_X(n)=\mathcal{L}^{\otimes n}$ avec
$\mathcal{L}=\mathcal{O}_X(1)$ (\hyperref[II.3.2.7-fr]{3.2.7} et
\hyperref[II.3.2.9-fr]{3.2.9})~; par définition
(\hyperref[II.8.6.1.1-fr]{8.6.1.1}), $C_X$ est donc le \emph{fibré vectoriel}
$\mathbf{V}(\mathcal{L})$ sur $X$, et son sommet la \emph{section nulle} de ce
fibré vectoriel.
\end{remark}

\subsection{Faisceaux amples et contractions}
\label{subsection:II.8.8-fr}

\begin{env}[8.8.1]
\label{II.8.8.1-fr}
Soient $Y$ un préschéma, $f:X\to Y$ un morphisme \emph{séparé et
quasi-compact}, $\mathcal{L}$ un $\mathcal{O}_X$-Module inversible \emph{ample
relativement à $f$}. Considérons la $\mathcal{O}_Y$-Algèbre graduée à degrés
positifs
\begin{equation}
\label{II.8.8.1.1-fr}
  \mathcal{S}=\mathcal{O}_Y\oplus
  \bigoplus_{n\geqslant 1}f_*(\mathcal{L}^{\otimes n})
\tag{8.8.1.1}
\end{equation}

\oldpage[II]{178}
qui est quasi-cohérente (\hyperref[I.9.2.2-fr]{I, 9.2.2, a)}). On a un
homomorphisme canonique de $\mathcal{O}_X$-Algèbres graduées
\begin{equation}
\label{II.8.8.1.2-fr}
  \tau:f^*(\mathcal{S})\longrightarrow
  \bigoplus_{n\geqslant 0}\mathcal{L}^{\otimes n}
\tag{8.8.1.2}
\end{equation}
qui, pour les degrés $\geqslant 1$, coïncide avec l'homomorphisme canonique
$\sigma:f^*(f_*(\mathcal{L}^{\otimes n}))\to\mathcal{L}^{\otimes n}$
(\hyperref[0.4.4.3-fr]{0, 4.4.3}), et pour le degré $0$ est l'identité.
L'hypothèse que $\mathcal{L}$ est $f$-ample entraîne alors
(\hyperref[II.4.6.3-fr]{4.6.3} et \hyperref[II.3.6.1-fr]{3.6.1}) que le
$Y$-morphisme correspondant
\begin{equation}
\label{II.8.8.1.3-fr}
  r=r_{\mathcal{L},\tau}:X\longrightarrow
  P=\operatorname{Proj}(\mathcal{S})
\tag{8.8.1.3}
\end{equation}
est partout défini et est une \emph{immersion ouverte dominante}, et que l'on
a
\begin{equation}
\label{II.8.8.1.4-fr}
  r^*(\mathcal{O}_P(n))=\mathcal{L}^{\otimes n}
  \qquad\text{pour tout }n\in\mathbf{Z}.
\tag{8.8.1.4}
\end{equation}
\end{env}

\begin{proposition}[8.8.2]
\label{II.8.8.2-fr}
Soit $C=\operatorname{Spec}(\mathcal{S})$ \emph{le cône affine défini par
$\mathcal{S}$}~; si $\mathcal{L}$ est $f$-ample, il existe un $Y$-morphisme
canonique
\begin{equation}
\label{II.8.8.2.1-fr}
  g:V=\mathbf{V}(\mathcal{L})\longrightarrow C
\tag{8.8.2.1}
\end{equation}
tel que le diagramme
\begin{equation}
\label{II.8.8.2.2-fr}
  \xymatrix{
    X\ar[r]^{j}\ar[d]_{f} &
    \mathbf{V}(\mathcal{L})\ar[r]^{\pi}\ar[d]^{g} & X\ar[d]^{f}\\
    Y\ar[r]_{\varepsilon} & C\ar[r]_{\psi} & Y
  }
\tag{8.8.2.2}
\end{equation}
soit commutatif, $\psi$ et $\pi$ étant les morphismes structuraux, $j$ et
$\varepsilon$ les immersions canoniques appliquant respectivement $X$ et $Y$
sur la section nulle de $\mathbf{V}(\mathcal{L})$ et le préschéma des sommets
de $C$. En outre, la restriction de $g$ à
$\mathbf{V}(\mathcal{L})-j(X)$ est une immersion ouverte
\begin{equation}
\label{II.8.8.2.3-fr}
  \mathbf{V}(\mathcal{L})-j(X)\longrightarrow
  E=C-\varepsilon(Y)
\tag{8.8.2.3}
\end{equation}
dans le cône épointé affine $E$ correspondant à $\mathcal{S}$.
\end{proposition}

\begin{proof}
Avec les notations de (\hyperref[II.8.8.1-fr]{8.8.1}), soient
$\mathcal{S}_P^{\geqslant}=\bigoplus_{n\geqslant 0}\mathcal{O}_P(n)$, et
$C_P=\operatorname{Spec}(\mathcal{S}_P^{\geqslant})$. On sait
(\hyperref[II.8.6.2-fr]{8.6.2}) qu'on a un morphisme canonique
$h=\operatorname{Spec}(\alpha):C_P\to C$ tel que le diagramme
\begin{equation}
\label{II.8.8.2.4-fr}
  \xymatrix{
    C_P\ar[r]\ar[d]_{h} & P\ar[d]^{p}\\
    C\ar[r]_{\psi} & Y
  }
\tag{8.8.2.4}
\end{equation}
soit commutatif~; en outre, si $\varepsilon_P:P\to C_P$ est l'immersion
canonique, le diagramme
\begin{equation}
\label{II.8.8.2.5-fr}
  \xymatrix{
    P\ar[r]^{\varepsilon_P}\ar[d]_{p} & C_P\ar[d]^{h}\\
    Y\ar[r]_{\varepsilon} & C
  }
\tag{8.8.2.5}
\end{equation}
est commutatif (\hyperref[II.8.7.1.1-fr]{8.7.1.1}), et enfin, la restriction
de $h$ au cône épointé affine $E_P$ est un \emph{isomorphisme}
$E_P\xrightarrow{\sim}E$ (\hyperref[II.8.6.2-fr]{8.6.2}). D'autre part, il
résulte de (\hyperref[II.8.8.1.4-fr]{8.8.1.4}) que l'on a
\[
  r^*(\mathcal{S}_P^{\geqslant})=
  \mathbf{S}_{\mathcal{O}_X}(\mathcal{L})
\]

\oldpage[II]{179}
et par suite on a un $P$-morphisme canonique
$q:\mathbf{V}(\mathcal{L})\to C_P$, le diagramme commutatif
\begin{equation}
\label{II.8.8.2.6-fr}
  \xymatrix{
    \mathbf{V}(\mathcal{L})\ar[r]^{\pi}\ar[d]_{q} & X\ar[d]^{r}\\
    C_P\ar[r] & P
  }
\tag{8.8.2.6}
\end{equation}
identifiant $\mathbf{V}(\mathcal{L})$ au produit $C_P\times_P X$
(\hyperref[II.1.5.2-fr]{1.5.2})~; comme $r$ est une immersion ouverte, il en
est donc de même de $q$ (\hyperref[I.4.3.2-fr]{I, 4.3.2}). En outre, la
restriction de $q$ à $\mathbf{V}(\mathcal{L})-j(X)$ applique ce préschéma dans
$E_P$ par (\hyperref[II.8.5.2-fr]{8.5.2}), et le diagramme
\begin{equation}
\label{II.8.8.2.7-fr}
  \xymatrix{
    X\ar[r]^{j}\ar[d]_{r} & \mathbf{V}(\mathcal{L})\ar[d]^{q}\\
    P\ar[r]_{\varepsilon_P} & C_P
  }
\tag{8.8.2.7}
\end{equation}
est commutatif (cas particulier de
(\hyperref[II.8.5.1.3-fr]{8.5.1.3})). Les assertions de
(\hyperref[II.8.8.2-fr]{8.8.2}) résultent aussitôt de ces faits, en prenant
pour $g$ le morphisme composé $h\circ q$.
\end{proof}

\begin{remark}[8.8.3]
\label{II.8.8.3-fr}
Supposons en outre que $Y$ soit un préschéma \emph{noethérien} et $f$ un
morphisme \emph{propre}. Comme $r$ est alors \emph{propre}
(\hyperref[II.5.4.4-fr]{5.4.4}), donc fermé, et que c'est une immersion ouverte
dominante, $r$ est nécessairement un \emph{isomorphisme}
$X\xrightarrow{\sim}P$. On verra en outre au chapitre III
(\hyperref[III.2.3.5.1-fr]{III, 2.3.5.1}), que $\mathcal{S}$ est alors
nécessairement une $\mathcal{O}_Y$-Algèbre \emph{de type fini}. Alors il en
résulte que $\mathcal{S}^{\natural}$ est une
$\mathcal{S}_0^{\natural}$-Algèbre \emph{de type fini}
(\hyperref[II.8.2.10-fr]{8.2.10, (i)} et
\hyperref[II.8.7.2.7-fr]{8.7.2.7})~; comme $C_P$ est $C$-isomorphe à
$\operatorname{Proj}(\mathcal{S}^{\natural})$
(\hyperref[II.8.7.3-fr]{8.7.3}), on voit que le morphisme $h:C_P\to C$ est
\emph{projectif}~; comme le morphisme $r$ est un isomorphisme, il en est de
même de $q:\mathbf{V}(\mathcal{L})\to C_P$, et on en conclut que le morphisme
$g:\mathbf{V}(\mathcal{L})\to C$ est \emph{projectif}. En outre, comme la
restriction de $h$ à $E_P$ est un isomorphisme sur $E$ et que $q$ est un
isomorphisme, la restriction (\hyperref[II.8.8.2.3-fr]{8.8.2.3}) de $g$ est un
\emph{isomorphisme}
$\mathbf{V}(\mathcal{L})-j(X)\xrightarrow{\sim}E$.

Si on suppose de plus que $\mathcal{L}$ est \emph{très ample} pour $f$, on
verra aussi au chapitre III (\hyperref[III.2.3.5.1-fr]{III, 2.3.5.1}), qu'il
existe un entier $n_0>0$ tel que
$\mathcal{S}_n=\mathcal{S}_1^n$ pour $n\geqslant n_0$. On en conclura par
(\hyperref[II.8.7.7-fr]{8.7.7}) que $\mathbf{V}(\mathcal{L})$ s'identifie au
préschéma $Z$ obtenu en \emph{faisant éclater dans le cône affine $C$ le
sous-préschéma des sommets} (identifié à $Y$), et que la \emph{section nulle}
de $\mathbf{V}(\mathcal{L})$ (identifiée à $Y$) est \emph{l'image réciproque}
du sous-préschéma des sommets $Y$ de $C$. Une partie des résultats précédents
peut en fait s'établir sans hypothèse noethérienne~:
\end{remark}

\begin{corollary}[8.8.4]
\label{II.8.8.4-fr}
Soient $Y$ un préschéma (resp. un schéma quasi-compact), $f:X\to Y$ un
morphisme propre, $\mathcal{L}$ un $\mathcal{O}_X$-Module inversible ample pour
$f$. Alors le morphisme (\hyperref[II.8.8.2.1-fr]{8.8.2.1}) est propre (resp.
projectif) et sa restriction (\hyperref[II.8.8.2.3-fr]{8.8.2.3}) est un
isomorphisme.
\end{corollary}

\begin{proof}
Pour prouver que $g$ est propre, on peut se ramener au cas où $Y$ est affine,
et il suffit donc de considérer le cas où $Y$ est un schéma quasi-compact. Le
même raisonnement que dans (\hyperref[II.8.8.3-fr]{8.8.3}) montre d'abord que
$r$ est un \emph{isomorphisme} $X\xrightarrow{\sim}P$~; par suite $q$ est aussi
un isomorphisme, et comme la restriction de $h$ à $E_P$ est un isomorphisme
$E_P\xrightarrow{\sim}E$, on voit déjà que
(\hyperref[II.8.8.2.3-fr]{8.8.2.3}) est un isomorphisme. Reste à prouver que
$g$ est \emph{projectif}.

Comme $f$ est de type fini par hypothèse, on peut appliquer
(\hyperref[II.3.8.5-fr]{3.8.5}) à l'homomorphisme

\oldpage[II]{180}
$\tau$ de (\hyperref[II.8.8.1.2-fr]{8.8.1.2})~: il y a par suite un entier
$d>0$, un sous-$\mathcal{O}_Y$-Module quasi-cohérent de type fini
$\mathcal{E}$ de $\mathcal{S}_d$ tel que si $\mathcal{S}'$ est la
sous-$\mathcal{O}_Y$-Algèbre de $\mathcal{S}$ engendrée par $\mathcal{E}$ et
$\tau'=\tau\circ f^*(\varphi)$ ($\varphi$ injection canonique
$\mathcal{S}'\to\mathcal{S}$), $r'=r_{\mathcal{L},\tau'}$ soit une immersion
\[
  X\longrightarrow P'=\operatorname{Proj}(\mathcal{S}').
\]
En outre, comme $\varphi$ est injectif, $r'$ est aussi une \emph{immersion
dominante} (\hyperref[II.3.7.6-fr]{3.7.6})~; le même raisonnement que pour $r$
montre donc que $r'$ est une \emph{immersion fermée surjective}~; comme $r'$ se
factorise en
$X\xrightarrow{r}\operatorname{Proj}(\mathcal{S})
\xrightarrow{\Phi}\operatorname{Proj}(\mathcal{S}')$, où
$\Phi=\operatorname{Proj}(\varphi)$, on en conclut que $\Phi$ est aussi une
\emph{immersion fermée surjective}. Mais cela entraîne que $\Phi$ est un
\emph{isomorphisme}~; en effet, on peut se ramener au cas où
$Y=\operatorname{Spec}(A)$ est affine, $\mathcal{S}=\widetilde{S}$,
$\mathcal{S}'=\widetilde{S'}$, où $S$ est une $A$-algèbre graduée, $S'$ une
sous-algèbre graduée de $S$. Pour tout élément homogène $t\in S'$,
$S'_{(t)}$ est alors un sous-anneau de $S_{(t)}$~; si on revient à la
définition de $\operatorname{Proj}(\varphi)$
(\hyperref[II.2.8.1-fr]{2.8.1}), on voit qu'on est ramené à prouver que si
$B'$ est un sous-anneau d'un anneau $B$, et si le morphisme
$\operatorname{Spec}(B)\to\operatorname{Spec}(B')$ correspondant à l'injection
canonique $B'\to B$ est une immersion fermée, ce morphisme est nécessairement
un \emph{isomorphisme}~; mais cela résulte de
(\hyperref[I.4.2.3-fr]{I, 4.2.3}).

On a en outre
$\Phi^*(\mathcal{O}_{P'}(n))=\mathcal{O}_P(n)$
(\hyperref[II.3.5.2-fr]{3.5.2, (ii)} et
\hyperref[II.3.5.4-fr]{3.5.4}), donc
$r'^*(\mathcal{O}_{P'}(n))$ est isomorphe à $\mathcal{L}^{\otimes n}$
(\hyperref[II.4.6.3-fr]{4.6.3}). Posons
$\mathcal{S}''=\mathcal{S}'^{(d)}$, de sorte
(\hyperref[II.3.1.8-fr]{3.1.8, (i)}) que $X$ s'identifie canoniquement à
$P''=\operatorname{Proj}(\mathcal{S}'')$ et
$\mathcal{L}''=\mathcal{L}^{\otimes d}$ à $\mathcal{O}_{P''}(1)$
(\hyperref[II.3.2.9-fr]{3.2.9, (ii)}).

Cela étant, si $C''=\operatorname{Spec}(\mathcal{S}'')$,
$\mathcal{S}_{P''}^{\geqslant}=
\bigoplus_{n\geqslant 0}\mathcal{O}_{P''}(n)$ s'identifie à
$\bigoplus_{n\geqslant 0}\mathcal{L}''^{\otimes n}$, donc
$C_{P''}=\operatorname{Spec}(\mathcal{S}_{P''}^{\geqslant})$ à
$\mathbf{V}(\mathcal{L}'')$~; on sait d'autre part
(\hyperref[II.8.7.3-fr]{8.7.3}) que $C_{P''}$ est $C''$-isomorphe à
$\operatorname{Proj}(\mathcal{S}''^{\natural})$~; par définition de
$\mathcal{S}''$, $\mathcal{S}''^{\natural}$ est engendré par
$\mathcal{S}_1''^{\natural}$ et $\mathcal{S}_1''^{\natural}$ est de type fini
sur $\mathcal{S}_0''^{\natural}=\mathcal{S}''$
(\hyperref[II.8.2.10-fr]{8.2.10, (i) et (iii)}), donc
$\operatorname{Proj}(\mathcal{S}''^{\natural})$ est \emph{projectif} sur
$C''$ (\hyperref[II.5.5.1-fr]{5.5.1}). Considérons alors le diagramme
\begin{equation}
\label{II.8.8.4.1-fr}
  \xymatrix{
    \mathbf{V}(\mathcal{L})\ar[r]^{g}\ar[d]_{u} &
    \operatorname{Spec}(\mathcal{S})=C\ar[d]^{v}\\
    \mathbf{V}(\mathcal{L}'')\ar[r]_{g''} &
    \operatorname{Spec}(\mathcal{S}'')=C''
  }
\tag{8.8.4.1}
\end{equation}
où $g$ et $g''$ correspondent par (\hyperref[II.1.5.6-fr]{1.5.6}) aux
$f$-morphismes canoniques
\[
  \mathcal{S}\longrightarrow
  \bigoplus_{n\geqslant 0}\mathcal{L}^{\otimes n}
  \quad\text{et}\quad
  \mathcal{S}''\longrightarrow
  \bigoplus_{n\geqslant 0}\mathcal{L}''^{\otimes n}
\]
(\hyperref[II.3.3.2.3-fr]{3.3.2.3}) (voir
(\hyperref[II.8.8.5-fr]{8.8.5}) ci-dessous), $v$ au morphisme d'inclusion
$\mathcal{S}''\to\mathcal{S}$, et $u$ au morphisme d'inclusion
$\bigoplus_{n\geqslant 0}\mathcal{L}^{\otimes nd}\to
\bigoplus_{n\geqslant 0}\mathcal{L}^{\otimes n}$~; il est immédiat
(\hyperref[II.3.3.2-fr]{3.3.2}) que ce diagramme est commutatif. Nous venons
de voir que $g''$ est un morphisme projectif~; d'autre part, $u$ est un
morphisme \emph{fini}. En effet, la question étant locale sur $X$, on peut
supposer que $X$ est affine d'anneau $A$ et que $\mathcal{L}=\mathcal{O}_X$~;
tout revient alors à remarquer que l'anneau $A[T]$ est un module de type fini
sur son sous-anneau $A[T^d]$ ($T$ indéterminée). Comme $Y$ est un schéma
quasi-compact, et que $C''$ est affine sur $Y$, $C''$ est aussi un schéma
quasi-compact,

\oldpage[II]{181}
et par suite $g''\circ u$ est un morphisme projectif
(\hyperref[II.5.5.5-fr]{5.5.5, (ii)})~; par la commutativité de
(\hyperref[II.8.8.4.1-fr]{8.8.4.1}), il en est de même de $v\circ g$, et comme
$v$ est affine, donc séparé, on en conclut finalement que $g$ est projectif
(\hyperref[II.5.5.5-fr]{5.5.5, (v)}).
\end{proof}

\begin{env}[8.8.5]
\label{II.8.8.5-fr}
Reprenons la situation de (\hyperref[II.8.8.1-fr]{8.8.1}). Nous allons voir
qu'on peut donner du morphisme $g:\mathbf{V}(\mathcal{L})\to C$ une autre
définition valable pour tout $\mathcal{O}_X$-Module inversible $\mathcal{L}$
(non nécessairement ample). Pour cela, considérons le $f$-morphisme
\begin{equation}
\label{II.8.8.5.1-fr}
  \tau^\flat:\mathcal{S}\longrightarrow
  \bigoplus_{n\geqslant 0}\mathcal{L}^{\otimes n}
\tag{8.8.5.1}
\end{equation}
correspondant au morphisme $\tau$ de
(\hyperref[II.8.8.1.2-fr]{8.8.1.2}). On en déduit
(\hyperref[II.1.5.6-fr]{1.5.6}) un morphisme $g':V\to C$ tel que, si
$\pi:V\to X$ et $\psi:C\to Y$ sont les morphismes structuraux, les diagrammes
\begin{equation}
\label{II.8.8.5.2-fr}
  \xymatrix{
    X\ar[d]_{f} & V\ar[l]^{\pi}\ar[d]^{g'}\\
    Y & C\ar[l]_{\psi}
  }
  \qquad
  \xymatrix{
    X\ar[r]^{j}\ar[d]_{f} & V\ar[d]^{g'}\\
    Y\ar[r]_{\varepsilon} & C
  }
\tag{8.8.5.2}
\end{equation}
soient commutatifs
(\hyperref[II.8.5.1.2-fr]{8.5.1.2} et
\hyperref[II.8.5.1.3-fr]{8.5.1.3}). Nous allons montrer que (si l'on suppose
$\mathcal{L}$ ample pour $f$) les morphismes $g$ et $g'$ sont identiques.

En effet, la question est locale sur $Y$, et on peut donc supposer que
$Y=\operatorname{Spec}(A)$ est affine, et (en raison de
(\hyperref[II.8.8.1.3-fr]{8.8.1.3})) identifier $X$ à un ouvert de
$P=\operatorname{Proj}(S)$, où
\[
  S=A\oplus\bigoplus_{n\geqslant 1}
  \Gamma(X,\mathcal{L}^{\otimes n})~;
\]
on déduit alors de (\hyperref[II.8.8.1.4-fr]{8.8.1.4}) que
$\Gamma(X,\mathcal{O}_P(n))=\Gamma(X,\mathcal{L}^{\otimes n})$ pour tout
$n\in\mathbf{Z}$. Tenant compte de la définition de
$h=\operatorname{Spec}(\alpha)$, où $\alpha$ est le $p$-morphisme canonique
$\widetilde{S}\to\mathcal{S}_P^{\geqslant}$
(\hyperref[II.8.6.1.2-fr]{8.6.1.2}), il faut vérifier que la restriction à
$X$ de
$\alpha^\sharp:p^*(\widetilde{S})\to\mathcal{S}_P^{\geqslant}$ est
identique à $\tau$. Compte tenu de
(\hyperref[0.4.4.3-fr]{0, 4.4.3}), on est ramené à voir que si on compose
l'homomorphisme canonique
$\alpha_n:S_n\to\Gamma(P,\mathcal{O}_P(n))$ avec l'homomorphisme de
restriction
\[
  \Gamma(P,\mathcal{O}_P(n))\longrightarrow
  \Gamma(X,\mathcal{O}_P(n))=
  \Gamma(X,\mathcal{L}^{\otimes n}),
\]
on obtient l'identité pour tout $n>0$~; or, cela résulte aussitôt de la
définition de l'algèbre $S$ et de celle de $\alpha_n$
(\hyperref[II.2.6.2-fr]{2.6.2}).
\end{env}

\begin{proposition}[8.8.6]
\label{II.8.8.6-fr}
Supposons (avec les notations de (\hyperref[II.8.8.5-fr]{8.8.5})) que si
l'on pose $f=(f_0,\lambda)$, l'homomorphisme
$\lambda:\mathcal{O}_Y\to f_*(\mathcal{O}_X)$ soit bijectif~; alors~:
\begin{enumerate}
  \item[\rm{(i)}] Si on pose $g=(g_0,\mu)$,
  $\mu:\mathcal{O}_C\to g_*(\mathcal{O}_V)$ est un isomorphisme.
  \item[\rm{(ii)}] Si $X$ est intègre (resp. localement intègre et normal),
  $C$ est intègre (resp. normal).
\end{enumerate}
\end{proposition}

\begin{proof}
En effet, le $f$-morphisme $\tau^\flat$ est alors un \emph{isomorphisme}
\[
  \tau^\flat:\mathcal{S}=\psi_*(\mathcal{O}_C)
  \longrightarrow f_*(\pi_*(\mathcal{O}_V))
  =\psi_*(g_*(\mathcal{O}_V))
\]
et le $Y$-morphisme $g$ peut être considéré comme celui pour lequel
l'homomorphisme $\mathcal{A}(g)$
(\hyperref[II.1.1.2-fr]{1.1.2}) est égal à $\tau^\flat$. Pour voir que
$\mu$ est un isomorphisme de $\mathcal{O}_C$-Modules, il suffit
(\hyperref[II.1.4.2-fr]{1.4.2}) de voir que
\[
  \mathcal{A}(\mu):\psi_*(\mathcal{O}_C)
  \longrightarrow\psi_*(g_*(\mathcal{O}_V))
\]
est un isomorphisme. Mais par définition
(\hyperref[II.1.1.2-fr]{1.1.2}), on a
$\mathcal{A}(\mu)=\mathcal{A}(g)$, d'où la conclusion de (i).

Pour prouver (ii), on peut se ramener au cas où $Y$ est affine, donc
$\mathcal{S}=\widetilde{S}$, avec

\oldpage[II]{182}
$S=\bigoplus_{n\geqslant 0}\Gamma(X,\mathcal{L}^{\otimes n})$~;
l'hypothèse que $X$ est intègre entraîne que l'anneau $S$ est intègre
(\hyperref[I.7.4.4-fr]{I, 7.4.4}), donc il en est de même de $C$
(\hyperref[I.5.1.4-fr]{I, 5.1.4}). Pour démontrer que $C$ est normal, nous
utiliserons le lemme suivant~:

\begin{lemma}[8.8.6.1]
\label{II.8.8.6.1-fr}
Soit $Z$ un préschéma intègre normal. Alors l'anneau
$\Gamma(Z,\mathcal{O}_Z)$ est intègre et intégralement clos.
\end{lemma}

\begin{proof}
En effet, il résulte de
(\hyperref[I.8.2.1.1-fr]{I, 8.2.1.1}) que
$\Gamma(Z,\mathcal{O}_Z)$ est l'intersection, dans le corps des fonctions
rationnelles $R(Z)$, des anneaux intégralement clos $\mathcal{O}_z$ pour
$z\in Z$.
\end{proof}

Cela étant, montrons d'abord que $V$ est \emph{localement intègre et normal}~;
on peut pour cela se borner au cas où $X=\operatorname{Spec}(A)$ est affine,
d'anneau $A$ intègre et intégralement clos
(\hyperref[II.6.3.8-fr]{6.3.8}), et où
$\mathcal{L}=\mathcal{O}_X$. Comme alors
$V=\operatorname{Spec}(A[T])$ et que $A[T]$ est intègre et intégralement clos
([24], p.~99), cela établit notre assertion. D'autre part, pour tout ouvert
affine $U$ de $C$, $g^{-1}(U)$ est quasi-compact puisque le morphisme $g$ est
quasi-compact~; puisque $V$ est localement intègre, les composantes connexes
de $g^{-1}(U)$ sont des préschémas intègres ouverts dans $g^{-1}(U)$, donc en
nombre fini, et comme $V$ est normal, ces préschémas sont normaux
(\hyperref[II.6.3.8-fr]{6.3.8}). Alors $\Gamma(U,\mathcal{O}_C)$, qui est égal
à $\Gamma(g^{-1}(U),\mathcal{O}_V)$ en vertu de (i), est composé direct d'un
nombre fini d'anneaux intègres et intégralement clos
(\hyperref[II.8.8.6.1-fr]{8.8.6.1}), ce qui montre que $C$ est normal
(\hyperref[II.6.3.4-fr]{6.3.4}).
\end{proof}

\subsection{Le critère d'amplitude de Grauert : énoncé}
\label{subsection:II.8.9-fr}

Nous nous proposons de montrer que les propriétés démontrées dans
(\hyperref[II.8.8.2-fr]{8.8.2}) \emph{caractérisent} les
$\mathcal{O}_X$-Modules amples pour $f$, et de façon précise, de démontrer le
critère suivant~:

\begin{theorem}[8.9.1]
\label{II.8.9.1-fr}
\emph{(critère de Grauert).} Soient $Y$ un préschéma,
$p:X\to Y$ un morphisme séparé et quasi-compact, $\mathcal{L}$ un
$\mathcal{O}_X$-Module inversible. Pour que $\mathcal{L}$ soit ample
relativement à $p$, il faut et il suffit qu'il existe un $Y$-préschéma $C$,
une $Y$-section $\varepsilon:Y\to C$ de $C$, et un $Y$-morphisme
$q:\mathbf{V}(\mathcal{L})\to C$, ayant les propriétés suivantes~:
\begin{enumerate}
  \item[\rm{(i)}] Le diagramme
  \begin{equation}
  \label{II.8.9.1.1-fr}
    \xymatrix{
      X\ar[r]^{j}\ar[d]_{p} & \mathbf{V}(\mathcal{L})\ar[d]^{q}\\
      Y\ar[r]_{\varepsilon} & C
    }
  \tag{8.9.1.1}
  \end{equation}
  où $j$ est la section nulle du fibré vectoriel $\mathbf{V}(\mathcal{L})$,
  est commutatif.
  \item[\rm{(ii)}] La restriction de $q$ à
  $\mathbf{V}(\mathcal{L})-j(X)$ est une immersion ouverte quasi-compacte
  \[
    \mathbf{V}(\mathcal{L})-j(X)\longrightarrow C
  \]
  dont l'image ne rencontre pas $\varepsilon(Y)$.
\end{enumerate}
\end{theorem}

On notera que si $C$ est \emph{séparé} sur $Y$, on peut, dans la condition
(ii), supprimer l'hypothèse que l'immersion ouverte est quasi-compacte~; pour
voir en effet que cette dernière propriété est conséquence des autres
conditions, on peut se borner au cas où $Y$ est affine, et l'assertion découle
alors de (\hyperref[I.5.5.1-fr]{I, 5.5.1, (i)} et
\hyperref[I.5.5.10-fr]{5.5.10}). On peut aussi supprimer

\oldpage[II]{183}
la même hypothèse lorsque l'on suppose $X$ noethérien, car alors $V$ est aussi
noethérien et l'assertion découle cette fois de
(\hyperref[I.6.3.5-fr]{I, 6.3.5}).

\begin{corollary}[8.9.2]
\label{II.8.9.2-fr}
Si le morphisme $p:X\to Y$ est propre, on peut dans l'énoncé de
(\hyperref[II.8.9.1-fr]{8.9.1}), supposer $q$ propre, et remplacer
«~immersion ouverte~» par «~isomorphisme~».
\end{corollary}

De façon imagée, on peut dire (lorsque $p:X\to Y$ est propre) que
$\mathcal{L}$ est ample relativement à $p$ si et seulement si on peut
«~contracter~» la section nulle du fibré vectoriel $\mathbf{V}(\mathcal{L})$
en le préschéma de base $Y$. Un cas particulier important est celui où $Y$
est le spectre d'un corps, et où l'opération de «~contraction~» consiste donc
à contracter la section nulle de $\mathbf{V}(\mathcal{L})$ \emph{en un seul
point}.

\begin{env}[8.9.3]
\label{II.8.9.3-fr}
La nécessité des conditions de (\hyperref[II.8.9.1-fr]{8.9.1}) et le
corollaire (\hyperref[II.8.9.2-fr]{8.9.2}) découlent aussitôt de
(\hyperref[II.8.8.2-fr]{8.8.2}) et
(\hyperref[II.8.8.4-fr]{8.8.4}).

Pour démontrer la suffisance de (\hyperref[II.8.9.1-fr]{8.9.1}), nous
considérerons une situation un peu plus générale. Pour cela, posons (avec les
notations de (\hyperref[II.8.8.2-fr]{8.8.2}))
\[
  \mathcal{S}'=\bigoplus_{n\geqslant 0}\mathcal{L}^{\otimes n}
\]
et
\[
  V=\mathbf{V}(\mathcal{L})=\operatorname{Spec}(\mathcal{S}').
\]
Le sous-préschéma fermé $j(X)$, section nulle de $\mathbf{V}(\mathcal{L})$,
est défini par le faisceau quasi-cohérent d'idéaux
$\mathcal{J}=(\mathcal{S}'_+)^{\sim}$ de $\mathcal{O}_V$
(\hyperref[II.1.4.10-fr]{1.4.10}). Ce $\mathcal{O}_V$-Module est
\emph{inversible}, car la question est locale sur $X$ et cela revient à
remarquer que l'idéal $TA[T]$ dans un anneau de polynômes $A[T]$ est un
$A[T]$-module libre monogène. En outre, il est immédiat (toujours parce que la
question est locale sur $X$) que
\[
  \mathcal{L}=j^*(\mathcal{J})
\]
et
\[
  j_*(\mathcal{L})=\mathcal{J}/\mathcal{J}^2.
\]
D'autre part, si
\[
  \pi:\mathbf{V}(\mathcal{L})\longrightarrow X
\]
est le morphisme structural, on a
$\pi_*(\mathcal{J})=\mathcal{S}'_+$ et
$\pi_*(\mathcal{J}/\mathcal{J}^2)=\mathcal{L}$~; on a donc deux
homomorphismes canoniques
$\mathcal{L}\to\pi_*(\mathcal{J})\to\mathcal{L}$, le premier étant
l'injection canonique $\mathcal{L}\to\mathcal{S}'_+$, le second la projection
canonique de $\mathcal{S}'_+$ sur $\mathcal{S}'_1=\mathcal{L}$, et leur
composé est l'identité. On peut d'ailleurs plonger canoniquement
\[
  \pi_*(\mathcal{J})=\mathcal{S}'_+
  =\bigoplus_{n\geqslant 1}\mathcal{L}^{\otimes n}
\]
dans le \emph{produit}
\[
  \prod_{n\geqslant 1}\mathcal{L}^{\otimes n}
  =\varprojlim_n\pi_*(\mathcal{J}/\mathcal{J}^{n+1})
\]
(puisque
$\pi_*(\mathcal{J}/\mathcal{J}^{n+1})=
\mathcal{L}\oplus\mathcal{L}^{\otimes2}\oplus\cdots\oplus
\mathcal{L}^{\otimes n}$), et on a donc encore deux homomorphismes canoniques
\begin{equation}
\label{II.8.9.3.1-fr}
  \mathcal{L}\longrightarrow
  \varprojlim_n\pi_*(\mathcal{J}/\mathcal{J}^{n+1})
  \longrightarrow\mathcal{L}
\tag{8.9.3.1}
\end{equation}
dont le composé est l'identité.

Cela étant, la généralisation de (\hyperref[II.8.9.1-fr]{8.9.1}) que nous
allons prouver est la suivante~:
\end{env}

\begin{proposition}[8.9.4]
\label{II.8.9.4-fr}
Soient $Y$ un préschéma, $V$ un $Y$-préschéma, $X$ un sous-préschéma fermé de
$V$, défini par un Idéal $\mathcal{J}$ de $\mathcal{O}_V$, qui est un
$\mathcal{O}_V$-Module inversible~; si $j:X\to V$ est

\oldpage[II]{184}
l'injection canonique, on pose
$\mathcal{L}=j^*(\mathcal{J})=
\mathcal{J}\otimes_{\mathcal{O}_V}\mathcal{O}_X$, de sorte que
$j_*(\mathcal{L})=\mathcal{J}/\mathcal{J}^2$. On suppose que le morphisme
structural $p:X\to Y$ est séparé et quasi-compact, et que les conditions
suivantes sont remplies~:
\begin{enumerate}
  \item[\rm{(i)}] Il existe un $Y$-morphisme de type fini $\pi:V\to X$ tel
  que $\pi\circ j=1_X$, et par suite
  $\pi_*(\mathcal{J}/\mathcal{J}^2)=\mathcal{L}$.
  \item[\rm{(ii)}] Il existe un homomorphisme de $\mathcal{O}_X$-Modules
  \[
    \varphi:\mathcal{L}\longrightarrow
    \varprojlim_n\pi_*(\mathcal{J}/\mathcal{J}^{n+1})
  \]
  tel que le composé
  \[
    \mathcal{L}\xrightarrow{\varphi}
    \varprojlim_n\pi_*(\mathcal{J}/\mathcal{J}^{n+1})
    \xrightarrow{\alpha}\pi_*(\mathcal{J}/\mathcal{J}^2)=\mathcal{L}
  \]
  (où $\alpha$ est l'homomorphisme canonique) soit l'identité.
  \item[\rm{(iii)}] Il existe un $Y$-préschéma $C$, une $Y$-section
  $\varepsilon$ de $C$ et un $Y$-morphisme $q:V\to C$ tels que le diagramme
  \begin{equation}
  \label{II.8.9.4.1-fr}
    \xymatrix{
      X\ar[r]^{j}\ar[d]_{p} & V\ar[d]^{q}\\
      Y\ar[r]_{\varepsilon} & C
    }
  \tag{8.9.4.1}
  \end{equation}
  soit commutatif.
  \item[\rm{(iv)}] La restriction de $q$ à $W=V-j(X)$ est une immersion
  ouverte quasi-compacte dans $C$, dont l'image ne rencontre pas
  $\varepsilon(Y)$.
\end{enumerate}
Dans ces conditions, $\mathcal{L}$ est ample relativement à $p$.
\end{proposition}

\subsection{Le critère d'amplitude de Grauert : démonstration}
\label{subsection:II.8.10-fr}

\begin{lemma}[8.10.1]
\label{II.8.10.1-fr}
Soient $\pi:V\to X$ un morphisme, $j:X\to V$ une $X$-section de $V$ qui soit
une immersion fermée, $\mathcal{J}$ l'Idéal quasi-cohérent de
$\mathcal{O}_V$ définissant le sous-préschéma fermé de $V$ associé à $j$.
\begin{enumerate}
  \item[\rm{(i)}] Pour tout $n\geqslant 0$,
  $\pi_*(\mathcal{O}_V/\mathcal{J}^{n+1})$ et
  $\pi_*(\mathcal{J}/\mathcal{J}^{n+1})$ sont des
  $\mathcal{O}_X$-Modules quasi-cohérents, et on a
  $\pi_*(\mathcal{O}_V/\mathcal{J})=\mathcal{O}_X$,
  $\pi_*(\mathcal{J}/\mathcal{J}^2)=j^*(\mathcal{J})$.
  \item[\rm{(ii)}] Si $X=\{\xi\}=\operatorname{Spec}(k)$, où $k$ est un
  corps, $\varprojlim_n\pi_*(\mathcal{O}_V/\mathcal{J}^{n+1})$ est isomorphe
  au séparé complété de l'anneau local $\mathcal{O}_{j(\xi)}$ pour la
  topologie $\mathfrak{m}_{j(\xi)}$-préadique.
  \item[\rm{(iii)}] Supposons que $\mathcal{J}$ soit un
  $\mathcal{O}_V$-Module inversible (ce qui entraîne que
  \[
    \mathcal{L}=j^*(\mathcal{J})=
    \pi_*(\mathcal{J}/\mathcal{J}^2)
  \]
  est un $\mathcal{O}_X$-Module inversible) et qu'il existe un homomorphisme
  \[
    \varphi:\mathcal{L}\longrightarrow
    \varprojlim_n\pi_*(\mathcal{J}/\mathcal{J}^{n+1})
  \]
  tel que le composé
  \[
    \mathcal{L}\xrightarrow{\varphi}
    \varprojlim_n\pi_*(\mathcal{J}/\mathcal{J}^{n+1})
    \xrightarrow{\alpha}\pi_*(\mathcal{J}/\mathcal{J}^2)
  \]
  ($\alpha$ homomorphisme canonique) soit l'identité. Si on pose
  $\mathcal{S}=\bigoplus_{n\geqslant0}\mathcal{L}^{\otimes n}$, on déduit
  alors canoniquement de $\varphi$ un isomorphisme de
  $\mathcal{O}_X$-Algèbres du complété $\widehat{\mathcal{S}}$ de
  $\mathcal{S}$ relatif à sa filtration canonique (complété qui est isomorphe
  au produit $\prod_{n\geqslant0}\mathcal{L}^{\otimes n}$) sur
  $\varprojlim_n\pi_*(\mathcal{O}_V/\mathcal{J}^{n+1})$.
\end{enumerate}
\end{lemma}

\begin{proof}
Notons en premier lieu que le support du $\mathcal{O}_V$-Module
$\mathcal{O}_V/\mathcal{J}^{n+1}$ est $j(X)$ et celui de
$\mathcal{J}/\mathcal{J}^{n+1}$ est contenu dans $j(X)$. Dans le cas (ii),
$j(X)$ est un point fermé $j(\xi)$ de $V$,

\oldpage[II]{185}
et par définition $\pi_*(\mathcal{O}_V/\mathcal{J}^{n+1})$ est la fibre de
$\mathcal{O}_V/\mathcal{J}^{n+1}$ au point $j(\xi)$, c'est-à-dire, en posant
$C=\mathcal{O}_{j(\xi)}$ et en désignant par $\mathfrak{m}$ l'idéal maximal
de $C$, le $C$-module $C/\mathfrak{m}^{n+1}$~; l'assertion (ii) est alors
évidente.

Pour démontrer (i), notons que la question est locale sur $X$~; on peut donc
se restreindre au cas où $X$ est affine. Soit $U$ un ouvert affine dans $V$~;
$j(X)\cap U$ est un ouvert affine dans $j(X)$, donc
$U_0=\pi(j(X)\cap U)$, qui lui est isomorphe, un ouvert affine dans $X$~;
pour tout ouvert affine $W_0\subset U_0$ dans $X$,
$W=\pi^{-1}(W_0)\cap U$ est un ouvert affine dans $V$, puisque $X$ est un
schéma (\hyperref[I.5.5.10-fr]{I, 5.5.10})~; en particulier
$U'=U\cap\pi^{-1}(U_0)$ est un ouvert affine dans $V$ et on a évidemment
$\pi(U')=U_0$ et $j(U_0)=j(X)\cap U$. Cela étant, par définition
\[
  \Gamma(W_0,\pi_*(\mathcal{O}_V/\mathcal{J}^{n+1}))
  =\Gamma(\pi^{-1}(W_0),\mathcal{O}_V/\mathcal{J}^{n+1})~;
\]
mais comme tout point de $\pi^{-1}(W_0)$ n'appartenant pas à $j(W_0)$ a un
voisinage ouvert dans $\pi^{-1}(W_0)$ ne rencontrant pas $j(X)$ et dans
lequel $\mathcal{O}_V/\mathcal{J}^{n+1}$ est donc nul, il est clair que les
sections de $\mathcal{O}_V/\mathcal{J}^{n+1}$ au-dessus de $\pi^{-1}(W_0)$ et
au-dessus de $W$ se correspondent biunivoquement. Autrement dit, si $\pi'$ est
la restriction de $\pi$ à $U'$, les $(\mathcal{O}_X|U_0)$-Modules
\[
  \pi_*(\mathcal{O}_V/\mathcal{J}^{n+1})|U_0
  \quad\hbox{et}\quad
  \pi'_*\bigl((\mathcal{O}_V/\mathcal{J}^{n+1})|U'\bigr)
\]
sont identiques. Comme $U'$ et $U_0$ sont affines et que les $U_0$ recouvrent
$X$, on en conclut (\hyperref[I.1.6.3-fr]{I, 1.6.3}) que
$\pi_*(\mathcal{O}_V/\mathcal{J}^{n+1})$ est quasi-cohérent, et la
démonstration est identique pour
$\pi_*(\mathcal{J}/\mathcal{J}^{n+1})$.

Pour démontrer enfin (iii), notons que $\mathcal{S}$ n'est autre que
$\mathbf{S}_{\mathcal{O}_X}(\mathcal{L})$~; on déduit donc canoniquement de
$\varphi$ un homomorphisme de $\mathcal{O}_X$-Algèbres
\[
  \psi:\mathcal{S}\longrightarrow
  \varprojlim_n\pi_*(\mathcal{O}_V/\mathcal{J}^{n+1})
\]
(\hyperref[II.1.7.4-fr]{1.7.4})~; en outre, cet homomorphisme applique
$\mathcal{L}^{\otimes n}$ dans
$\varprojlim_m\pi_*(\mathcal{J}^n/\mathcal{J}^{m+1})$, donc est continu pour
les topologies considérées, et se prolonge bien par suite en un homomorphisme
\[
  \widehat{\psi}:\widehat{\mathcal{S}}\longrightarrow
  \varprojlim_n\pi_*(\mathcal{O}_V/\mathcal{J}^{n+1}).
\]
Pour voir qu'il s'agit d'un isomorphisme, on peut, comme dans la démonstration
de (i), se réduire au cas où $X=\operatorname{Spec}(A)$ et
$V=\operatorname{Spec}(B)$ sont affines, avec
$\mathcal{J}=\widetilde{\mathfrak{J}}$, où $\mathfrak{J}$ est un idéal de
$B$~; il correspond à $\pi$ une injection $A\to B$, identifiant $A$ à un
sous-anneau de $B$ supplémentaire de $\mathfrak{J}$, et $\mathcal{L}$ (resp.
$\pi_*(\mathcal{O}_V/\mathcal{J}^{n+1})$) est le $\mathcal{O}_X$-Module
quasi-cohérent associé au $A$-module $L=\mathfrak{J}/\mathfrak{J}^2$ (resp.
$B/\mathfrak{J}^{n+1}$). Comme $\mathcal{J}$ est un
$\mathcal{O}_V$-Module \emph{inversible}, on peut en outre supposer que
$\mathfrak{J}=Bt$, où $t$ est non diviseur de 0 dans $B$. De la relation
$B=A\oplus Bt$ on déduit alors, pour tout $n>0$,
\[
  B=A\oplus At\oplus At^2\oplus\cdots\oplus At^n\oplus Bt^{n+1}
\]
donc on a un $A$-isomorphisme canonique de l'anneau de séries formelles
$A[[T]]$ sur $C=\varprojlim B/\mathfrak{J}^{n+1}$ faisant correspondre $t$ à
$T$. D'autre part, on a $L=A\bar t$, où $\bar t$ est la classe de $t$
mod. $Bt^2$, et l'homomorphisme $\varphi$ applique par hypothèse $\bar t$ sur
un élément $t'\in C$ congru à $t$ mod. $Ct^2$. On en déduit par récurrence
sur $n$ que
\[
  A\oplus At'\oplus\cdots\oplus At'{}^n\oplus Ct^{n+1}
  =A\oplus At\oplus\cdots\oplus At^n\oplus Ct^{n+1}
\]
ce qui prouve que l'homomorphisme $\widehat{\psi}$ correspond bien à un
isomorphisme de $\prod_{n\geqslant0}L^{\otimes n}$ sur $C$.
\end{proof}

\begin{lemma}[8.10.2]
\label{II.8.10.2-fr}
Sous les hypothèses de (\hyperref[II.8.10.1-fr]{8.10.1}), soit
$g:X'\to X$ un morphisme,

\oldpage[II]{186}
soit $V'=V\times_X X'$, et soient $\pi':V'\to X'$,
$g':V'\to V$ les projections canoniques, de sorte qu'on a le diagramme
commutatif
\[
  \xymatrix{
    V\ar[d]_{\pi} & V'\ar[l]^{g'}\ar[d]^{\pi'}\\
    X & X'\ar[l]_{g}
  }
\]
Alors $j'=j\times1_{X'}$ est une $X'$-section de $V'$ qui est une immersion
fermée, et $\mathcal{J}'=(g')^*(\mathcal{J})\mathcal{O}_{V'}$ est l'Idéal
quasi-cohérent de $\mathcal{O}_{V'}$ définissant le sous-préschéma fermé de
$V'$ associé à $j'$. En outre, on a
\[
  \pi'_*(\mathcal{O}_{V'}/\mathcal{J}'^{n+1})
  =g^*(\pi_*(\mathcal{O}_V/\mathcal{J}^{n+1})).
\]
Enfin, $\mathcal{J}'$ est un $\mathcal{O}_{V'}$-Module canoniquement isomorphe
à $(g')^*(\mathcal{J})$ et est en particulier inversible si $\mathcal{J}$ est
un $\mathcal{O}_V$-Module inversible.
\end{lemma}

\begin{proof}
Le fait que $j'$ est une immersion fermée découle de
(\hyperref[I.4.3.1-fr]{I, 4.3.1}), et c'est une $X'$-section de $V'$ par
fonctorialité de l'extension du préschéma de base. En outre, si $Z$ (resp.
$Z'$) est le sous-préschéma fermé de $V$ (resp. $V'$) associé à $j$ (resp.
$j'$), on a $Z'=g'{}^{-1}(Z)$
(\hyperref[I.4.3.1-fr]{I, 4.3.1}), et la seconde assertion résulte alors de
(\hyperref[I.4.4.5-fr]{I, 4.4.5}). Pour démontrer les autres assertions, on
voit comme dans (\hyperref[II.8.10.1-fr]{8.10.1}) qu'on peut se ramener au cas
où $X$, $V$ et $X'$ (donc aussi $V'$) sont affines~; gardons les notations de
la démonstration de (\hyperref[II.8.10.1-fr]{8.10.1}) et soit
$X'=\operatorname{Spec}(A')$. On a alors $V'=\operatorname{Spec}(B')$ où
$B'=B\otimes_A A'$, et $\mathcal{J}'=\widetilde{\mathfrak{J}'}$, avec
$\mathfrak{J}'=\operatorname{Im}(\mathfrak{J}\otimes_A A')$. On a donc
\[
  B'/\mathfrak{J}'^{n+1}=(B/\mathfrak{J}^{n+1})\otimes_A A'~;
\]
en outre, comme $\mathfrak{J}$ est facteur direct (en tant que $A$-module) de
$B$, $\mathfrak{J}\otimes_A A'$ est facteur direct (en tant que $A'$-module)
de $B'$, donc s'identifie canoniquement à $\mathfrak{J}'$.
\end{proof}

\begin{corollary}[8.10.3]
\label{II.8.10.3-fr}
Supposons vérifiées les hypothèses de (\hyperref[II.8.10.1-fr]{8.10.1}) et
supposons en outre que $\pi$ soit de type fini et que $\mathcal{J}$ soit un
$\mathcal{O}_V$-Module inversible. Alors, pour tout $x\in X$, l'anneau local
au point $j(x)$ de la fibre $\pi^{-1}(x)$ est un anneau régulier (donc
intègre) de dimension 1, dont le complété est isomorphe à l'anneau de séries
formelles $k(x)[[T]]$ ($T$ indéterminée)~; en outre il n'existe qu'une seule
composante irréductible de $\pi^{-1}(x)$ contenant $j(x)$.
\end{corollary}

\begin{proof}
Comme $\pi^{-1}(x)=V\times_X\operatorname{Spec}(k(x))$, on est ramené par
(\hyperref[II.8.10.2-fr]{8.10.2}) au cas où $X$ est le spectre d'un corps
$K$. Comme $\pi$ est de type fini
(\hyperref[I.6.3.4-fr]{I, 6.3.4, (iv)}), $\mathcal{O}_{j(x)}$ est alors un
anneau local noethérien, donc séparé pour la topologie
$\mathfrak{m}_{j(x)}$-préadique
(\hyperref[0.7.3.5-fr]{0, 7.3.5})~; il résulte de
(\hyperref[II.8.10.1-fr]{8.10.1, (ii) et (iii)}) que le complété de cet anneau
est isomorphe à $K[[T]]$, et par suite $\mathcal{O}_{j(x)}$ est régulier et de
dimension 1 ([1], p.~17-01, th.~1)~; enfin, puisque $\mathcal{O}_{j(x)}$ est
intègre, $j(x)$ n'appartient qu'à une seule des composantes irréductibles (en
nombre fini) de $V$ (\hyperref[I.5.1.4-fr]{I, 5.1.4}).
\end{proof}

\begin{corollary}[8.10.4]
\label{II.8.10.4-fr}
Supposons vérifiées les hypothèses de (\hyperref[II.8.10.1-fr]{8.10.1}) et en
outre supposons que $\mathcal{J}$ soit un $\mathcal{O}_V$-Module inversible.
Soit $W=V-j(X)$~; pour tout Idéal quasi-cohérent $\mathcal{K}$ de
$\mathcal{O}_X$, posons
$\mathcal{K}_V=\pi^*(\mathcal{K})\mathcal{O}_V$ et
$\mathcal{K}_W=\mathcal{K}_V|W$. Alors $\mathcal{K}_V$ est le plus grand
Idéal quasi-cohérent de $\mathcal{O}_V$ dont la restriction à $W$ soit
$\mathcal{K}_W$.
\end{corollary}

\begin{proof}
En effet, on voit comme dans (\hyperref[II.8.10.1-fr]{8.10.1}) que la question
est locale sur $X$ et sur $V$~; on peut donc garder les notations de la
démonstration de (\hyperref[II.8.10.1-fr]{8.10.1}), avec
$\mathfrak{J}=Bt$, où $t$ n'est pas diviseur de 0 dans $B$. En outre, on a
$W=\operatorname{Spec}(B_t)$ et $\mathcal{K}=\widetilde{\mathfrak{K}}$, où
$\mathfrak{K}$ est un idéal de $A$~; d'où
$\pi^*(\mathcal{K})\mathcal{O}_V=(\mathfrak{K}.B)^{\sim}$
(\hyperref[I.1.6.9-fr]{I, 1.6.9}),
$\mathcal{K}_W=(\mathfrak{K}.B_t)^{\sim}$, et le plus grand idéal

\oldpage[II]{187}
de $B$ dont l'image canonique dans $B_t$ soit $\mathfrak{K}.B_t$ est l'image
réciproque de $\mathfrak{K}.B_t$, c'est-à-dire l'ensemble des $s\in B$ tels
que pour un entier $n>0$, on ait $t^ns\in\mathfrak{K}.B$. Il faut montrer que
cette dernière relation entraîne $s\in\mathfrak{K}.B$, ou encore que l'image
canonique de $t$ n'est pas diviseur de 0 dans
$B/\mathfrak{K}B=(A/\mathfrak{K})\otimes_A B$, ce qui résulte de
(\hyperref[II.8.10.2-fr]{8.10.2}) appliqué à
$X'=\operatorname{Spec}(A/\mathfrak{K})$.
\end{proof}

\begin{corollary}[8.10.5]
\label{II.8.10.5-fr}
Supposons vérifiées les hypothèses de (\hyperref[II.8.10.3-fr]{8.10.3})~;
soient $W=V-j(X)$, $x$ un point de $X$, $\mathcal{K}$ un Idéal
quasi-cohérent de $\mathcal{O}_X$, $z$ le point générique de la composante
irréductible de $\pi^{-1}(x)$ contenant $j(x)$
(\hyperref[II.8.10.3-fr]{8.10.3}).
\begin{enumerate}
  \item[\rm{(i)}] Soit $g$ une section de $\mathcal{O}_V$ au-dessus de $V$
  telle que $g|W$ soit une section de $\mathcal{K}_W$ au-dessus de $W$
  (notations de (\hyperref[II.8.10.4-fr]{8.10.4})). Alors $g$ est une section
  de $\mathcal{K}_V$~; si de plus $g(z)\neq0$, et si, pour tout entier $m>0$,
  on désigne par $g_m^x$ le germe au point $x$ de l'image canonique $g_m$ de
  $g$ dans
  $\Gamma(X,\pi_*(\mathcal{O}_V/\mathcal{J}^{m+1}))$, alors il existe un
  entier $m>0$ tel que l'image de $g_m^x$ dans
  \[
    (\pi_*(\mathcal{O}_V/\mathcal{J}^{m+1}))_x
    \otimes_{\mathcal{O}_x}k(x)
  \]
  soit $\neq0$.
  \item[\rm{(ii)}] Supposons en outre que les conditions de
  (\hyperref[II.8.10.1-fr]{8.10.1, (iii)}) soient remplies. Alors, s'il existe
  une section $g$ de $\mathcal{K}_V$ au-dessus de $V$ telle que $g(z)\neq0$,
  il existe un entier $n\geqslant0$ et une section $f$ de
  $\mathcal{K}.\mathcal{L}^{\otimes n}=
  \mathcal{K}\otimes\mathcal{L}^{\otimes n}\subset\mathcal{L}^{\otimes n}$
  telle que $f(x)\neq0$. Si $g$ est une section de $\mathcal{J}$, on peut
  prendre $n>0$.
\end{enumerate}
\end{corollary}

\begin{proof}
\begin{enumerate}
  \item[\rm{(i)}] Comme l'Idéal de $\mathcal{O}_W$ engendré par $g|W$ est
  contenu dans $\mathcal{K}_W$ par hypothèse, l'Idéal de $\mathcal{O}_V$
  engendré par $g$ est contenu dans $\mathcal{K}_V$ par
  (\hyperref[II.8.10.4-fr]{8.10.4}), autrement dit $g$ est une section de
  $\mathcal{K}_V$. Pour démontrer la seconde assertion de (i), on peut encore
  supposer $X$ et $V$ affines et conserver les notations de
  (\hyperref[II.8.10.1-fr]{8.10.1})~; la fibre $\pi^{-1}(x)$ est alors affine
  d'anneau $B'=B\otimes_A k(x)$, et il existe dans $B'$ un élément $t'$ non
  diviseur de 0 tel que $B'=k(x)\oplus B't'$. Comme $j(x)$ est une
  spécialisation de $z$ et que $g(z)\neq0$, on a nécessairement
  $g_{j(x)}\neq0$. Mais $\mathcal{O}_{j(x)}$ est un anneau local séparé
  (\hyperref[II.8.10.3-fr]{8.10.3}), qui se plonge donc dans son complété, et
  l'image de $g$ dans ce complété n'est par suite pas nulle. Or, ce complété
  est isomorphe à $\varprojlim_n(B'/B't'^{n+1})$
  (\hyperref[II.8.10.3-fr]{8.10.3})~; si $g'=g\otimes1\in B'$, il existe donc
  un entier $m$ tel que $g'\notin B't'^{m+1}$, ou encore l'image $g'_m$ de
  $g'$ dans $B'/B't'^{m+1}$ n'est pas nulle. Mais comme $g'_m$ n'est autre que
  l'image de $g_m^x$, notre assertion est démontrée.
  \item[\rm{(ii)}] En vertu de
  (\hyperref[II.8.10.1-fr]{8.10.1, (iii)})
  $\pi_*(\mathcal{O}_V/\mathcal{J}^{m+1})$ est isomorphe à la somme directe
  des $\mathcal{L}^{\otimes k}$ pour $0\leqslant k\leqslant m$~; désignons
  par $f_k$ la section de $\mathcal{L}^{\otimes k}$ au-dessus de $X$,
  composante de l'élément de
  $\bigoplus_{k=0}^m\Gamma(X,\mathcal{L}^{\otimes k})$ qui correspond à $g_m$
  par cet isomorphisme. Choisissant $m$ comme dans (i), il y a donc un indice
  $k$ tel que $f_k(x)\neq0$ d'après (i). Pour voir que $f_k$ est une section
  de $\mathcal{K}\mathcal{L}^{\otimes k}$, il suffit de considérer comme
  ci-dessus le cas où $X$ et $V$ sont affines, et cela résulte aussitôt alors
  de ce que $g\in\mathfrak{K}.B$ (notations de
  (\hyperref[II.8.10.4-fr]{8.10.4})). La dernière assertion résulte de ce que
  l'hypothèse $g\in\Gamma(V,\mathcal{J})$ entraîne $f_0=0$.
\end{enumerate}
\end{proof}

\begin{env}[8.10.6]
\label{II.8.10.6-fr}
\emph{Démonstration de (\hyperref[II.8.9.4-fr]{8.9.4}).} La question est
locale sur $Y$ (\hyperref[II.4.6.4-fr]{4.6.4})~; comme $\varepsilon$ est une
$Y$-section, on peut donc remplacer $C$ par un voisinage ouvert affine $U$
d'un point de $\varepsilon(Y)$ tel que $\varepsilon(Y)\cap U$ soit fermé dans
$U$. Autrement dit, on peut supposer $C$ affine, et $Y$ sous-préschéma fermé
de $C$ (donc affine) défini par un faisceau quasi-

\oldpage[II]{188}
cohérent $\mathcal{I}$ d'idéaux de $\mathcal{O}_C$. Comme $p$ est séparé et
quasi-compact, $X$ est alors un schéma quasi-compact, et on est ramené à
prouver que $\mathcal{L}$ est ample
(\hyperref[II.4.6.4-fr]{4.6.4}). En vertu du critère
(\hyperref[II.4.5.2-fr]{4.5.2, a)}), on doit donc prouver qui ce suit~: pour
tout Idéal quasi-cohérent $\mathcal{K}$ de $\mathcal{O}_X$ et tout point
$x\in X$ n'appartenant pas au support de $\mathcal{O}_X/\mathcal{K}$, il
existe un entier $n>0$ et une section $f$ de
$\mathcal{K}\otimes\mathcal{L}^{\otimes n}$ au-dessus de $X$ telle que
$f(x)\neq0$.

Pour cela, posons
\[
  \mathcal{K}_V=\pi^*(\mathcal{K})\mathcal{O}_V
\]
et
\[
  \mathcal{K}_W=\mathcal{K}_V|W,\qquad\hbox{où}\qquad W=V-j(X)~;
\]
comme la restriction de $q$ à $W$ est une immersion quasi-compacte dans $C$,
il résulte de (\hyperref[I.9.4.2-fr]{I, 9.4.2}) que $\mathcal{K}_W$ est la
restriction à $W$ d'un Idéal quasi-cohérent $\mathcal{K}'_V$ de
$\mathcal{O}_V$, de la forme
\[
  \mathcal{K}'_V=q^*(\mathcal{K}_C)\mathcal{O}_V
\]
où $\mathcal{K}_C$ est un Idéal quasi-cohérent de $\mathcal{O}_C$. En outre,
comme par hypothèse $q^{-1}(Y)\subset j(X)$ et que $Y$ est défini par l'Idéal
$\mathcal{I}$, la restriction à $W$ de $q^*(\mathcal{I})\mathcal{O}_V$ est
identique à celle de $\mathcal{O}_V$, donc $\mathcal{K}_W$ est aussi la
restriction à $W$ de $q^*(\mathcal{I}\mathcal{K}_C)\mathcal{O}_V$, et l'on
peut donc supposer que l'on a $\mathcal{K}_C\subset\mathcal{I}$, d'où
\begin{equation}
\label{II.8.10.6.1-fr}
  \mathcal{K}'_V\subset q^*(\mathcal{I})\mathcal{O}_V\subset\mathcal{J}
\tag{8.10.6.1}
\end{equation}
compte tenu de (\hyperref[I.4.4.6-fr]{I, 4.4.6}) et de la commutativité de
(\hyperref[II.8.9.4.1-fr]{8.9.4.1}). En outre, on déduit de
(\hyperref[II.8.10.4-fr]{8.10.4}) que l'on a
\begin{equation}
\label{II.8.10.6.2-fr}
  \mathcal{K}'_V\subset\mathcal{K}_V.
\tag{8.10.6.2}
\end{equation}

Cela étant, il résulte de (\hyperref[II.8.10.3-fr]{8.10.3}) que $j(x)$
appartient à une seule composante irréductible de $\pi^{-1}(x)$~; soit $z$ le
point générique de cette composante, et posons $z'=q(z)$. En vertu de
(\hyperref[II.8.10.5-fr]{8.10.5}), la démonstration sera achevée (compte tenu
de (\hyperref[II.8.10.6.1-fr]{8.10.6.1}) et
(\hyperref[II.8.10.6.2-fr]{8.10.6.2})) si nous prouvons l'existence d'une
section $g$ de $\mathcal{K}'_V$ au-dessus de $V$ telle que $g(z)\neq0$. Or,
par hypothèse, $\mathcal{K}$ a une restriction égale à celle de
$\mathcal{O}_X$ dans un voisinage ouvert de $x$~; d'autre part, il résulte de
(\hyperref[II.8.10.3-fr]{8.10.3}) que l'on a $z\neq j(x)$, donc $z\in W$, et
par suite $(\mathcal{K}_W)_z=\mathcal{O}_{V,z}$, d'où par définition
$(\mathcal{K}_C)_{z'}=\mathcal{O}_{C,z'}$. Puisque $C$ est affine, il y a donc
une section $g'$ de $\mathcal{K}_C$ au-dessus de $C$ telle que
$g'(z')\neq0$, et en prenant pour $g$ la section de $\mathcal{K}'_V$
correspondant canoniquement à $g'$, on a bien $g(z)\neq0$, ce qui achève la
démonstration.
\end{env}

\begin{remark}[8.10.7]
\label{II.8.10.7-fr}
Nous ignorons si, dans l'énoncé de (\hyperref[II.8.9.4-fr]{8.9.4}), la
condition (ii) est ou non superflue. En tout cas, la conclusion ne subsiste
plus si l'on ne suppose pas l'existence d'un $Y$-morphisme $\pi:V\to X$ tel
que $\pi\circ j=1_X$~; indiquons brièvement comment on peut en effet former un
contre-exemple, dont les détails ne pourront être développés que plus tard.
On prend $Y=\operatorname{Spec}(k)$ où $k$ est un corps,
$C=\operatorname{Spec}(A)$, où $A=k[T_1,T_2]$, la $Y$-section $\varepsilon$
correspondant à l'homomorphisme d'augmentation $A\to k$. On désigne par $C'$
le schéma déduit de $C$ par éclatement du point fermé $a=\varepsilon(Y)$ de
$C$~; si $D$ est l'image réciproque de $a$ dans $C'$, on considère dans $D$
un point fermé $b$,

\oldpage[II]{189}
et on désigne par $V$ le schéma déduit de $C'$ par éclatement de $b$~; $X$ est
le sous-préschéma fermé de $V$, image réciproque de $a$ par le morphisme
structural $q:V\to C$. On montre que $X$ est réunion de deux composantes
irréductibles $X_1$, $X_2$, où $X_1$ est l'image réciproque de $b$ dans $V$.
Il est immédiat que l'Idéal $\mathcal{J}$ de $\mathcal{O}_V$ qui définit $X$
est encore inversible, mais on prouve que
$j^*(\mathcal{J})=\mathcal{L}$ ($j$ injection canonique $X\to V$) n'est pas
ample, par la considération du «~degré~» de l'image réciproque de
$\mathcal{L}$ sur $X_1$, qui devrait être $>0$ si $\mathcal{L}$ était ample,
et qu'on montre (par un calcul élémentaire d'intersections) être égal à 0.
\end{remark}

\subsection{Unicité des contractions}
\label{subsection:II.8.11-fr}

\begin{lemma}[8.11.1]
\label{II.8.11.1-fr}
Soient $U$, $V$ deux préschémas,
$h=(h_0,\lambda):U\to V$ un morphisme surjectif. On suppose que~:
\begin{enumerate}
  \item[$1^\circ$]
  $\lambda:\mathcal{O}_V\to h_*(\mathcal{O}_U)
  =(h_0)_*(\mathcal{O}_U)$ est un isomorphisme.
  \item[$2^\circ$] L'espace sous-jacent à $V$ s'identifie à l'espace quotient
  de l'espace sous-jacent à $U$ par la relation
  $R:h_0(x)=h_0(y)$ (condition toujours vérifiée lorsque le morphisme $h$ est
  ouvert, ou fermé, ou a fortiori lorsque $h$ est propre).
\end{enumerate}
Alors, pour tout préschéma $W$, l'application
\begin{equation}
\label{II.8.11.1.1-fr}
  \operatorname{Hom}(V,W)\longrightarrow\operatorname{Hom}(U,W)
\tag{8.11.1.1}
\end{equation}
qui, à tout morphisme $v=(v_0,\nu)$ de $V$ dans $W$, fait correspondre le
morphisme $u=v\circ h=(u_0,\mu)$, est une bijection de
$\operatorname{Hom}(V,W)$ sur l'ensemble des $u$ tels que $u_0$ soit
constante sur toute fibre $h_0^{-1}(x)$.
\end{lemma}

\begin{proof}
Il est clair que si $u=v\circ h$, donc $u_0=v_0\circ h_0$, $u_0$ est
constante sur tout ensemble $h_0^{-1}(x)$. Inversement, si $u$ a cette
propriété, montrons qu'il existe un $v\in\operatorname{Hom}(V,W)$ et un seul
tel que $u=v\circ h$. L'existence et l'unicité de l'application continue
$v_0:V\to W$ telle que $u_0=v_0\circ h_0$ résultent de l'hypothèse, puisque
$h_0$ s'identifie à l'application canonique de $U$ sur $U/R$. On peut d'autre
part, en remplaçant au besoin $V$ par un préschéma isomorphe, supposer que
$\lambda$ est l'identité~; par hypothèse, $\mu$ est alors un homomorphisme
\[
  \mu:\mathcal{O}_W\longrightarrow(u_0)_*(\mathcal{O}_U)
  =(v_0)_*((h_0)_*(\mathcal{O}_U))
\]
tel que l'homomorphisme correspondant
$\mu^\sharp:u_0^*(\mathcal{O}_W)\to\mathcal{O}_U$ soit local sur chaque
fibre. Comme
$(v_0)_*((h_0)_*(\mathcal{O}_U))=(v_0)_*(\mathcal{O}_V)$, on doit
nécessairement avoir $\nu=\mu$, et tout revient à voir que l'homomorphisme
correspondant $\nu^\sharp:v_0^*(\mathcal{O}_W)\to\mathcal{O}_V$ est local sur
chaque fibre. Or, tout $y\in V$ est de la forme $h_0(x)$ pour un $x\in U$~;
posons $z=v_0(y)=u_0(x)$. Alors
(\hyperref[0.3.5.5-fr]{0, 3.5.5}) l'homomorphisme $\mu_x^\sharp$ se factorise
en
\[
  \mu_x^\sharp:\mathcal{O}_z
  \xrightarrow{\nu_y^\sharp}\mathcal{O}_y
  \xrightarrow{\lambda_x^\sharp}\mathcal{O}_x.
\]
Par hypothèse $\lambda_x^\sharp$ et $\mu_x^\sharp$ sont des homomorphismes
locaux~; donc $\lambda_x^\sharp$ transforme tout élément inversible de
$\mathcal{O}_y$ en un élément inversible de $\mathcal{O}_x$~; si
$\nu_y^\sharp$ transformait un élément non inversible de $\mathcal{O}_z$ en
un élément inversible de $\mathcal{O}_y$, $\mu_x^\sharp$ transformerait cet
élément de $\mathcal{O}_z$ en un élément inversible de $\mathcal{O}_x$,
contrairement à l'hypothèse, d'où le lemme.
\end{proof}

\begin{corollary}[8.11.2]
\label{II.8.11.2-fr}
Soient $U$ un préschéma intègre, $V$ un préschéma normal~; alors tout morphisme
$h:U\to V$ qui est universellement fermé, birationnel et radiciel est un
isomorphisme.
\end{corollary}

\begin{proof}
Si $h=(h_0,\lambda)$, il résulte des hypothèses que $h_0$ est injectif et
fermé, et que $h_0(U)$

\oldpage[II]{190}
est dense dans $V$, donc $h_0$ est un \emph{homéomorphisme} de $U$ sur $V$.
Pour démontrer le corollaire, il suffira de voir que
$\lambda:\mathcal{O}_V\to(h_0)_*(\mathcal{O}_U)$ est un isomorphisme~: on
pourra appliquer alors (\hyperref[II.8.11.1-fr]{8.11.1}) qui prouvera que
l'application (\hyperref[II.8.11.1.1-fr]{8.11.1.1}) est bijective (les fibres
$h_0^{-1}(x)$ étant chacune réduite à un seul point)~; donc $h$ sera un
isomorphisme. La question étant évidemment locale sur $V$, on peut supposer
que $V=\operatorname{Spec}(A)$ est affine d'anneau intègre et intégralement
clos (\hyperref[II.8.8.6.1-fr]{8.8.6.1})~; $h$ correspond alors
(\hyperref[I.2.2.4-fr]{I, 2.2.4}) à un homomorphisme
$\varphi:A\to\Gamma(U,\mathcal{O}_U)$ et tout revient à voir que $\varphi$
est un isomorphisme. Or, si $K$ est le corps des fractions de $A$,
$\Gamma(U,\mathcal{O}_U)$ a par hypothèse $K$ pour corps des fractions et $A$
est un sous-anneau de $\Gamma(U,\mathcal{O}_U)$, $\varphi$ étant l'injection
canonique (\hyperref[I.8.2.7-fr]{I, 8.2.7}). Comme le morphisme $h$ vérifie
les hypothèses de (\hyperref[II.7.3.11-fr]{7.3.11}),
$\Gamma(U,\mathcal{O}_U)$ est un sous-anneau de la fermeture intégrale de $A$
dans $K$, donc est identique à $A$ par hypothèse.
\end{proof}

\begin{remark}[8.11.3]
\label{II.8.11.3-fr}
On verra au chapitre III (\hyperref[III.4.4.11-fr]{III, 4.4.11}) que lorsque
$V$ est un préschéma \emph{localement noethérien}, tout morphisme $h:U\to V$
qui est propre et quasi-fini (en particulier tout morphisme vérifiant les
hypothèses de (\hyperref[II.8.11.2-fr]{8.11.2})) est nécessairement
\emph{fini}. La conclusion de (\hyperref[II.8.11.2-fr]{8.11.2}) résulte donc
dans ce cas de (\hyperref[II.6.1.15-fr]{6.1.15}).
\end{remark}

\begin{env}[8.11.4]
\label{II.8.11.4-fr}
Nous allons maintenant voir que, dans le critère de Grauert
(\hyperref[II.8.9.1-fr]{8.9.1}), on peut souvent affirmer que le préschéma
$C$ et la «~contraction~» $q$ sont déterminés de façon essentiellement
unique.
\end{env}

\begin{lemma}[8.11.5]
\label{II.8.11.5-fr}
Soient $Y$ un préschéma, $p:X\to Y$ un morphisme propre, $\mathcal{L}$ un
$\mathcal{O}_X$-Module inversible $p$-ample, $C$ un $Y$-préschéma,
$\varepsilon:Y\to C$ une $Y$-section,
$q:V=\mathbf{V}(\mathcal{L})\to C$ un $Y$-morphisme, tels que le diagramme
(\hyperref[II.8.9.1.1-fr]{8.9.1.1}) soit commutatif. On suppose en outre que,
si $p=(p_0,\theta)$,
$\theta:\mathcal{O}_Y\to p_*(\mathcal{O}_X)$ est un isomorphisme. Soient alors
$\mathcal{S}'=\bigoplus_{n\geqslant0}p_*(\mathcal{L}^{\otimes n})$,
$C'=\operatorname{Spec}(\mathcal{S}')$, et
$q':\mathbf{V}(\mathcal{L})\to C'$ le $Y$-morphisme canonique
(\hyperref[II.8.8.5-fr]{8.8.5}). Il existe un $Y$-morphisme $u:C'\to C$ et
un seul tel que $q=u\circ q'$.
\end{lemma}

\begin{proof}
L'hypothèse sur $\theta$ entraîne en particulier que $p$ est surjectif~;
comme en vertu de (\hyperref[II.8.8.4-fr]{8.8.4}) la restriction de $q'$ à
$\mathbf{V}(\mathcal{L})-j(X)$ est un \emph{isomorphisme} sur
$C'-\varepsilon'(Y)$ ($\varepsilon'$ étant la section sommet de $C'$), il
résulte de (\hyperref[II.8.8.4-fr]{8.8.4}) que $q'$ est \emph{propre et
surjectif}~; en outre, en vertu de (\hyperref[II.8.8.6-fr]{8.8.6}), si on pose
$q'=(q'_0,\tau)$,
$\tau:\mathcal{O}_{C'}\to q'_*(\mathcal{O}_V)$ est un isomorphisme. On est
donc dans les conditions d'application du lemme
(\hyperref[II.8.11.1-fr]{8.11.1}), et le lemme sera démontré si l'on prouve
que $q$ est constant sur toute fibre $q'{}^{-1}(z')$, où $z'\in C'$. Or, la
condition est trivialement vérifiée pour $z'\notin\varepsilon'(Y)$. D'autre
part, si $z'\in\varepsilon'(Y)$, il existe un $y\in Y$ et un seul tel que
$z'=\varepsilon'(y)$~; et en vertu de la commutativité de
(\hyperref[II.8.8.5.2-fr]{8.8.5.2}) et le fait que $q'$ applique
$\mathbf{V}(\mathcal{L})-j(X)$ dans $C'-\varepsilon'(Y)$,
$q'{}^{-1}(z')=j(p^{-1}(y))$~; la commutativité du diagramme
(\hyperref[II.8.9.1.1-fr]{8.9.1.1}) établit donc notre assertion.
\end{proof}

\begin{corollary}[8.11.6]
\label{II.8.11.6-fr}
Sous les hypothèses de (\hyperref[II.8.11.5-fr]{8.11.5}), supposons outre que
$q$ soit propre et que la restriction de $q$ à
$\mathbf{V}(\mathcal{L})-j(X)$ soit un isomorphisme sur
$C-\varepsilon(Y)$. Alors le morphisme $u$ est universellement fermé,
surjectif et radiciel, et sa restriction à $C'-\varepsilon'(Y)$ est un
isomorphisme sur $C-\varepsilon(Y)$.
\end{corollary}

\begin{proof}
Comme $q'$ est un isomorphisme de $\mathbf{V}(\mathcal{L})-j(X)$ sur
$C'-\varepsilon'(Y)$ (\hyperref[II.8.8.4-fr]{8.8.4}), la dernière assertion
résulte aussitôt de $q=u\circ q'$. En outre, la commutativité des diagrammes

\oldpage[II]{191}
(\hyperref[II.8.8.5.2-fr]{8.8.5.2}) et
(\hyperref[II.8.9.1.1-fr]{8.9.1.1}) montre que la restriction de $u$ au
sous-préschéma fermé $\varepsilon'(Y)$ de $C'$ est un isomorphisme sur le
sous-préschéma fermé $\varepsilon(Y)$ de $C$, d'où résulte aussitôt que pour
tout $z'\in\varepsilon'(Y)$, si $z=u(z')$, $u$ définit un isomorphisme de
$k(z)$ sur $k(z')$. Ces remarques prouvent que $u$ est bijectif et radiciel~;
en outre, si $\psi$ et $\psi'$ sont les morphismes structuraux $C\to Y$,
$C'\to Y$, on a $\psi'=\psi\circ u$ et comme $\psi'$ est séparé
(\hyperref[II.1.2.4-fr]{1.2.4}), il en est de même de $u$
(\hyperref[I.5.5.1-fr]{I, 5.5.1, (v)}). On a vu d'autre part dans la
démonstration du lemme (\hyperref[II.8.11.5-fr]{8.11.5}) que $q'$ est
surjectif~; comme $q=u\circ q'$ est propre, on conclut finalement de
(\hyperref[II.5.4.3-fr]{5.4.3}) et (\hyperref[II.5.4.9-fr]{5.4.9}) que $u$
est universellement fermé.
\end{proof}

\begin{proposition}[8.11.7]
\label{II.8.11.7-fr}
Soient $Y$ un préschéma, $X$ un préschéma intègre, $p:X\to Y$ un morphisme
propre, $\mathcal{L}$ un $\mathcal{O}_X$-Module inversible $p$-ample, $C$ un
$Y$-préschéma normal, $\varepsilon:Y\to C$ une $Y$-section,
$q:V=\mathbf{V}(\mathcal{L})\to C$ un $Y$-morphisme, tels que le diagramme
(\hyperref[II.8.9.1.1-fr]{8.9.1.1}) soit commutatif. On suppose en outre que,
si $p=(p_0,\theta)$, $\theta:\mathcal{O}_Y\to p_*(\mathcal{O}_X)$ soit un
isomorphisme. Soient alors
$\mathcal{S}'=\bigoplus_{n\geqslant0}p_*(\mathcal{L}^{\otimes n})$,
$C'=\operatorname{Spec}(\mathcal{S}')$, et
$q':\mathbf{V}(\mathcal{L})\to C'$ le $Y$-morphisme canonique
(\hyperref[II.8.8.5-fr]{8.8.5}). Alors l'unique $Y$-morphisme $u:C'\to C$ tel
que $q=u\circ q'$ est un isomorphisme.
\end{proposition}

\begin{proof}
Il résulte de (\hyperref[II.8.8.6-fr]{8.8.6}) que $C'$ est intègre~; comme
$u$ est un homéomorphisme des espaces sous-jacents $C'\to C$ ($u$ étant
bijectif et fermé d'après (\hyperref[II.8.11.6-fr]{8.11.6})), $C$ est
irréductible, donc intègre, et puisque la restriction de $u$ à un ouvert non
vide de $C'$ est un isomorphisme sur un ouvert de $C$, $u$ est birationnel.
Puisque $C$ est supposé normal, il suffit d'appliquer
(\hyperref[II.8.11.2-fr]{8.11.2}) pour obtenir la conclusion.
\end{proof}

\begin{remark}[8.11.8]
\label{II.8.11.8-fr}
\begin{enumerate}
  \item[\rm{(i)}] Notons que l'hypothèse que $C$ est normal implique qu'il en
  est de même de $X$. En effet, $C'=\operatorname{Spec}(\mathcal{S}')$ est
  alors normal, étant isomorphe à $C$, et intègre en vertu de
  (\hyperref[II.8.8.6-fr]{8.8.6})~; on en conclut que
  $\operatorname{Proj}(\mathcal{S}')$ est \emph{normal}. En effet, la question
  est locale sur $Y$~; si $Y$ est affine, $\mathcal{S}'=\widetilde{S'}$,
  l'anneau $S'=\Gamma(C',\mathcal{O}_{C'})$ est intègre et intégralement clos
  (\hyperref[II.8.8.6.1-fr]{8.8.6.1}), donc, pour tout élément homogène
  $f\in S'_+$, l'anneau gradué $S'_f$ est intègre et intégralement clos
  ([13], t.~I, p.~257 et 261), et par suite aussi l'anneau $S'_{(f)}$ de ses
  termes de degré 0, car l'intersection de $S'_f$ et du corps des fractions de
  $S'_{(f)}$ est égale à $S'_{(f)}$~; ce qui prouve notre assertion
  (\hyperref[II.6.3.4-fr]{6.3.4}). Enfin, comme $X$ est isomorphe à un
  sous-préschéma ouvert de $\operatorname{Proj}(\mathcal{S}')$
  (\hyperref[II.8.8.1-fr]{8.8.1}), $X$ est bien normal. On peut donc exprimer
  la prop. (\hyperref[II.8.11.7-fr]{8.11.7}) sous la forme suivante~:
  \emph{Si $X$ est intègre et normal,
  $p=(p_0,\theta):X\to Y$ un morphisme propre tel que
  $\theta:\mathcal{O}_Y\to p_*(\mathcal{O}_X)$ soit un isomorphisme, alors,
  pour tout $\mathcal{O}_X$-Module $p$-ample $\mathcal{L}$, il existe une
  façon et une seule de contracter la section nulle de
  $V=\mathbf{V}(\mathcal{L})$ de façon à obtenir un $Y$-schéma normal $C$ et
  un $Y$-morphisme propre $q:V\to C$.}
  \item[\rm{(ii)}] Lorsque $p$ est propre, l'hypothèse
  $p_*(\mathcal{O}_X)=\mathcal{O}_Y$ peut être considérée comme une hypothèse
  auxiliaire, ne constituant pas une vraie restriction de la généralité du
  résultat. En effet, si elle n'est pas vérifiée, il suffit de remplacer $Y$
  par le $Y$-schéma
  $Y'=\operatorname{Spec}(p_*(\mathcal{O}_X))$, et de considérer $X$ comme un
  $Y'$-schéma. Nous reviendrons sur cette méthode générale au chapitre III,
  \S~4.
\end{enumerate}
\end{remark}
\subsection{Faisceaux quasi-cohérents sur les cônes projetants}
\label{subsection:II.8.12-fr}

\begin{env}[8.12.1]
\label{II.8.12.1-fr}
Reprenons les hypothèses et notations de
(\hyperref[II.8.3.1-fr]{8.3.1}). Soit $\mathcal{M}$ un
$\mathcal{S}$-Module gradué quasi-cohérent~; pour éviter toute confusion,
nous désignerons par $\widetilde{\mathcal{M}}$ le $\mathcal{O}_C$-Module

\oldpage[II]{192}
quasi-cohérent associé à $\mathcal{M}$
(\hyperref[II.1.4.3-fr]{1.4.3}) lorsque $\mathcal{M}$ est considéré comme
$\mathcal{S}$-Module \emph{non gradué}, et par
$\operatorname{Proj}_0(\mathcal{M})$ le $\mathcal{O}_X$-Module
quasi-cohérent associé à $\mathcal{M}$, $\mathcal{M}$ étant cette fois
considéré comme $\mathcal{S}$-Module gradué (autrement dit le
$\mathcal{O}_X$-Module noté $\widetilde{\mathcal{M}}$ dans
(\hyperref[II.3.2.2-fr]{3.2.2})). Nous poserons de plus
\begin{equation}
\label{II.8.12.1.1-fr}
  \mathcal{M}_X=\operatorname{Proj}(\mathcal{M})
  =\bigoplus_{n\in\mathbf{Z}}\operatorname{Proj}_0(\mathcal{M}(n))~;
\tag{8.12.1.1}
\end{equation}
la $\mathcal{O}_X$-Algèbre graduée quasi-cohérente $\mathcal{S}_X$ étant
définie par (\hyperref[II.8.6.1.1-fr]{8.6.1.1}),
$\operatorname{Proj}(\mathcal{M})$ est muni d'une structure de
$\mathcal{S}_X$-Module gradué (quasi-cohérent), au moyen des homomorphismes
canoniques (\hyperref[II.3.2.6.1-fr]{3.2.6.1})
\begin{equation}
\label{II.8.12.1.2-fr}
  \mathcal{O}_X(m)\otimes_{\mathcal{O}_X}
  \operatorname{Proj}_0(\mathcal{M}(n))
  \to \operatorname{Proj}_0(\mathcal{S}(m)\otimes_{\mathcal{S}}\mathcal{M}(n))
  \to \operatorname{Proj}_0(\mathcal{M}(m+n))
\tag{8.12.1.2}
\end{equation}
la vérification des axiomes des Modules se faisant à l'aide du diagramme
commutatif (\hyperref[II.2.5.11.4-fr]{2.5.11.4}).

Si $Y=\operatorname{Spec}(A)$ est affine, $\mathcal{S}=\widetilde{S}$ et
$\mathcal{M}=\widetilde{M}$, où $S$ est une $A$-algèbre graduée et $M$ un
$S$-module gradué, alors, pour tout élément homogène $f\in S_+$, on a
\begin{equation}
\label{II.8.12.1.3-fr}
  \Gamma(X_f,\operatorname{Proj}(\widetilde{M}))=M_f
\tag{8.12.1.3}
\end{equation}
en vertu des définitions et de (\hyperref[II.8.2.9.1-fr]{8.2.9.1}).

Considérons maintenant le $\widehat{\mathcal{S}}$-Module gradué
quasi-cohérent
\begin{equation}
\label{II.8.12.1.4-fr}
  \widehat{\mathcal{M}}
  =\mathcal{M}\otimes_{\mathcal{S}}\widehat{\mathcal{S}}
\tag{8.12.1.4}
\end{equation}
($\widehat{\mathcal{S}}$ étant défini par
(\hyperref[II.8.3.1.1-fr]{8.3.1.1}))~; on en déduit un
$\mathcal{O}_{\widehat{C}}$-Module gradué quasi-cohérent
$\operatorname{Proj}_0(\widehat{\mathcal{M}})$, que nous noterons aussi
\begin{equation}
\label{II.8.12.1.5-fr}
  \mathcal{M}^{\square}=\operatorname{Proj}_0(\widehat{\mathcal{M}}).
\tag{8.12.1.5}
\end{equation}
Il est clair (\hyperref[II.3.2.4-fr]{3.2.4}) que
$\mathcal{M}^{\square}$ est un foncteur additif exact en $\mathcal{M}$,
commutant aux sommes directes et aux limites inductives.
\end{env}

\begin{proposition}[8.12.2]
\label{II.8.12.2-fr}
Avec les notations de (\hyperref[II.8.3.2-fr]{8.3.2}), on a des
isomorphismes canoniques fonctoriels
\begin{equation}
\label{II.8.12.2.1-fr}
  i^*(\mathcal{M}^{\square})\xrightarrow{\sim}\widetilde{\mathcal{M}},
  \qquad
  j^*(\mathcal{M}^{\square})\xrightarrow{\sim}
  \operatorname{Proj}_0(\mathcal{M}).
\tag{8.12.2.1}
\end{equation}
\end{proposition}

\begin{proof}
En effet, $i^*(\mathcal{M}^{\square})$ s'identifie canoniquement à
$\left(\widehat{\mathcal{M}}/(z-1)\widehat{\mathcal{M}}\right)^{\sim}$
sur
$\operatorname{Spec}(\widehat{\mathcal{S}}/(z-1)\widehat{\mathcal{S}})$
en vertu de (\hyperref[II.3.2.3-fr]{3.2.3})~; le premier des isomorphismes
canoniques (\hyperref[II.8.12.2.1-fr]{8.12.2.1}) se déduit alors aussitôt
(\hyperref[II.1.4.1-fr]{1.4.1}) de l'isomorphisme canonique
$\widehat{\mathcal{M}}/(z-1)\widehat{\mathcal{M}}
\xrightarrow{\sim}\mathcal{M}$. D'autre part, l'immersion canonique
$j:X\to\widehat{C}$ correspond à l'homomorphisme canonique
$\widehat{\mathcal{S}}\to\mathcal{S}$ de noyau
$z\widehat{\mathcal{S}}$ (\hyperref[II.8.3.2-fr]{8.3.2})~; le second
homomorphisme (\hyperref[II.8.12.2.1-fr]{8.12.2.1}) est le cas particulier
de l'homomorphisme canonique (\hyperref[II.3.5.2-fr]{3.5.2, (ii)}), du fait
que l'on a ici
$\widehat{\mathcal{M}}\otimes_{\widehat{\mathcal{S}}}\mathcal{S}
=\mathcal{M}$~; pour vérifier que c'est un isomorphisme, on peut se borner au
cas où $Y=\operatorname{Spec}(A)$ est affine,
$\mathcal{S}=\widetilde{S}$, $\mathcal{M}=\widetilde{M}$~; en se reportant à
(\hyperref[II.2.8.8-fr]{2.8.8}), la vérification que pour tout $f$ homogène
dans $S_+$, l'homomorphisme précédent, restreint à $X_f$, se réduit à un
isomorphisme, est alors immédiate.
\end{proof}

\oldpage[II]{193}
Par abus de langage, on dira encore, en raison de l'existence du premier
isomorphisme (\hyperref[II.8.12.2.1-fr]{8.12.2.1}), que
$\mathcal{M}^{\square}$ est la \emph{fermeture projective} du
$\mathcal{O}_C$-Module $\widetilde{\mathcal{M}}$ (étant sous-entendu que dans
la donnée du $\mathcal{O}_C$-Module $\widetilde{\mathcal{M}}$ figure la
graduation du $\mathcal{S}$-Module $\mathcal{M}$).

\begin{env}[8.12.3]
\label{II.8.12.3-fr}
Avec les notations de (\hyperref[II.8.3.5-fr]{8.3.5}), on a un
homomorphisme canonique fonctoriel
\begin{equation}
\label{II.8.12.3.1-fr}
  p^*(\operatorname{Proj}_0(\mathcal{M}))
  \to\mathcal{M}^{\square}|\widehat{E}.
\tag{8.12.3.1}
\end{equation}
En effet, c'est un cas particulier de l'homomorphisme $u^{\sharp}$ défini de
façon générale dans (\hyperref[II.3.5.6-fr]{3.5.6}). Si
$Y=\operatorname{Spec}(A)$ est affine, $\mathcal{S}=\widetilde{S}$,
$\mathcal{M}=\widetilde{M}$, on voit, en se reportant à
(\hyperref[II.2.8.8-fr]{2.8.8}), que la restriction de
(\hyperref[II.8.12.3.1-fr]{8.12.3.1}) à
$p^{-1}(X_f)=\widehat{C}_f$ (pour un $f$ homogène dans $S_+$) correspond à
l'homomorphisme canonique
\begin{equation}
\label{II.8.12.3.2-fr}
  M_{(f)}\otimes_{S_{(f)}}S_f^{\leq}\to M_f^{\leq}
\tag{8.12.3.2}
\end{equation}
compte tenu de (\hyperref[II.8.2.3.2-fr]{8.2.3.2}) et
(\hyperref[II.8.2.5.2-fr]{8.2.5.2}).
\end{env}

\begin{env}[8.12.4]
\label{II.8.12.4-fr}
Plaçons-nous dans les hypothèses de (\hyperref[II.8.5.1-fr]{8.5.1}), dont
nous conservons les notations. Il résulte de
(\hyperref[II.1.5.6-fr]{1.5.6}) que pour tout $\mathcal{S}$-Module gradué
quasi-cohérent $\mathcal{M}$, on a d'une part un isomorphisme canonique
\begin{equation}
\label{II.8.12.4.1-fr}
  \Phi^*(\widetilde{\mathcal{M}})
  \xrightarrow{\sim}
  \left(q^*(\mathcal{M})\otimes_{q^*(\mathcal{S})}\mathcal{S}'\right)^{\sim}
\tag{8.12.4.1}
\end{equation}
de $\mathcal{O}_{C'}$-Modules~; et d'autre part
(\hyperref[II.3.5.6-fr]{3.5.6}) entraîne l'existence d'un
$\operatorname{Proj}(\varphi)$-morphisme canonique
\begin{equation}
\label{II.8.12.4.2-fr}
  \operatorname{Proj}_0(\mathcal{M})\to
  \left(\operatorname{Proj}_0
  (q^*(\mathcal{M})\otimes_{q^*(\mathcal{S})}\mathcal{S}')\right)
  |G(\varphi)
\tag{8.12.4.2}
\end{equation}
et aussi d'un $\widehat{\Phi}$-morphisme canonique
\begin{equation}
\label{II.8.12.4.3-fr}
  \operatorname{Proj}_0(\widehat{\mathcal{M}})\to
  \left(\operatorname{Proj}_0
  (q^*(\widehat{\mathcal{M}})
  \otimes_{q^*(\widehat{\mathcal{S}})}\widehat{\mathcal{S}'})\right)
  |G(\widehat{\varphi}).
\tag{8.12.4.3}
\end{equation}
\end{env}

\begin{env}[8.12.5]
\label{II.8.12.5-fr}
Considérons maintenant la situation de (\hyperref[II.8.6.1-fr]{8.6.1}),
avec les mêmes notations~; on prend donc dans ce qui précède $Y'=X$, le
morphisme $q:X\to Y$ étant le morphisme structural, et $\varphi$ est le
$q$-morphisme canonique (\hyperref[II.8.6.1.2-fr]{8.6.1.2}). On a alors un
isomorphisme canonique
\begin{equation}
\label{II.8.12.5.1-fr}
  q^*(\mathcal{M})\otimes_{q^*(\mathcal{S})}\mathcal{S}_X^{\geq}
  \xrightarrow{\sim}\mathcal{M}_X^{\geq}
\tag{8.12.5.1}
\end{equation}
en posant
$\mathcal{M}_X^{\geq}=\bigoplus_{n\geq0}\operatorname{Proj}_0(\mathcal{M}(n))$.
On peut en effet se borner au cas où $Y=\operatorname{Spec}(A)$ est affine,
$\mathcal{S}=\widetilde{S}$ et $\mathcal{M}=\widetilde{M}$, et définir
l'isomorphisme (\hyperref[II.8.12.5.1-fr]{8.12.5.1}) dans chacun des ouverts
affines $X_f$ ($f$ homogène dans $S_+$), en vérifiant la compatibilité avec le
passage à un multiple homogène de $f$. Or, la restriction à $X_f$ du premier
membre de (\hyperref[II.8.12.5.1-fr]{8.12.5.1}) est
$\widetilde{M}'=
\left((M\otimes_A S_{(f)})
\otimes_{S\otimes_A S_{(f)}}S_f^{\geq}\right)^{\sim}$ par
(\hyperref[II.8.6.2.1-fr]{8.6.2.1})~; comme on a un isomorphisme canonique de
$M\otimes_A S_{(f)}$ sur $M\otimes_S(S\otimes_A S_{(f)})$, on en déduit un
isomorphisme de $\widetilde{M}'$ sur
$\left(M\otimes_S S_f^{\geq}\right)^{\sim}$, et ce dernier est
canoniquement isomorphe par (\hyperref[II.8.2.9.1-fr]{8.2.9.1}) à la
restriction à $X_f$ du second membre de
(\hyperref[II.8.12.5.1-fr]{8.12.5.1}), et satisfait aux conditions de
compatibilité requises.

\oldpage[II]{194}
Remplaçant $\mathcal{M}$ par $\widehat{\mathcal{M}}$, $\mathcal{S}$ par
$\widehat{\mathcal{S}}$ et $\mathcal{S}_X$ par
$(\mathcal{S}_X^{\geq})^{\wedge}$ dans le raisonnement précédent, on a de
même un isomorphisme canonique
\begin{equation}
\label{II.8.12.5.2-fr}
  q^*(\widehat{\mathcal{M}})
  \otimes_{q^*(\widehat{\mathcal{S}})}(\mathcal{S}_X^{\geq})^{\wedge}
  \xrightarrow{\sim}(\mathcal{M}_X^{\geq})^{\wedge}.
\tag{8.12.5.2}
\end{equation}
Si l'on se souvient (\hyperref[II.8.6.2-fr]{8.6.2}) que le morphisme
structural $u:\operatorname{Proj}(\mathcal{S}_X^{\geq})\to X$ est un
isomorphisme, on déduit d'abord de ce qui précède que l'on a un
$u$-isomorphisme canonique
\begin{equation}
\label{II.8.12.5.3-fr}
  \operatorname{Proj}_0\mathcal{M}
  \xrightarrow{\sim}\operatorname{Proj}_0(\mathcal{M}_X^{\geq})
\tag{8.12.5.3}
\end{equation}
comme cas particulier de (\hyperref[II.8.12.4.2-fr]{8.12.4.2}). On constate
en effet, avec les notations de la démonstration de
(\hyperref[II.8.6.2-fr]{8.6.2}), que cela revient à voir que
l'homomorphisme canonique
$M_{(f)}\otimes_{S_{(f)}}(S_f^{\geq})^{(d)}\to(M_f^{\geq})^{(d)}$ est un
isomorphisme lorsque $f\in S_d$, ce qui est immédiat.

En second lieu, l'isomorphisme (\hyperref[II.8.12.5.2-fr]{8.12.5.2}) permet,
en appliquant cette fois (\hyperref[II.8.12.4.3-fr]{8.12.4.3}) au morphisme
canonique $r=\operatorname{Proj}(\widehat{\alpha}):\widehat{C}_X\to
\widehat{C}$, d'obtenir un $r$-morphisme canonique
\begin{equation}
\label{II.8.12.5.4-fr}
  \mathcal{M}^{\square}\to(\mathcal{M}_X^{\geq})^{\square}.
\tag{8.12.5.4}
\end{equation}
Rappelons maintenant (\hyperref[II.8.6.2-fr]{8.6.2}) que les restrictions de
$r$ aux cônes épointés $\widehat{E}_X$ et $E_X$ sont des \emph{isomorphismes}
sur $\widehat{E}$ et $E$ respectivement. En outre~:
\end{env}

\begin{proposition}[8.12.6]
\label{II.8.12.6-fr}
Les restrictions à $\widehat{E}_X$ et à $E_X$ du $r$-morphisme canonique
(\hyperref[II.8.12.5.4-fr]{8.12.5.4}) sont des isomorphismes
\begin{equation}
\label{II.8.12.6.1-fr}
  \mathcal{M}^{\square}|\widehat{E}
  \xrightarrow{\sim}(\mathcal{M}_X^{\geq})^{\square}|\widehat{E}_X
\tag{8.12.6.1}
\end{equation}
\begin{equation}
\label{II.8.12.6.2-fr}
  \widetilde{\mathcal{M}}|E
  \xrightarrow{\sim}(\mathcal{M}_X^{\geq})^{\sim}|E_X.
\tag{8.12.6.2}
\end{equation}
\end{proposition}

\begin{proof}
On se ramène au cas où $Y$ est affine comme dans la démonstration de
(\hyperref[II.8.6.2-fr]{8.6.2})~; avec les notations de cette dernière, et
en se ramenant aux définitions (\hyperref[II.2.8.8-fr]{2.8.8}), on doit
montrer que l'homomorphisme canonique
\[
  \widehat{M}_{(f)}\otimes_{\widehat{S}_{(f)}}
  (S_f^{\geq})_{(f/1)}^{\wedge}
  \to(M\otimes_S S_f^{\geq})_{(f/1)}^{\wedge}
\]
est un isomorphisme~; mais en vertu de
(\hyperref[II.8.2.3.2-fr]{8.2.3.2}) et
(\hyperref[II.8.2.5.2-fr]{8.2.5.2}), le premier membre s'identifie
canoniquement à
$M_f^{\leq}\otimes_{S_f^{\leq}}(S_f^{\geq})_{f/1}^{\leq}$, donc à
$M_f^{\leq}$ en vertu de (\hyperref[II.8.2.7.2-fr]{8.2.7.2}), et le second à
$(M_f^{\geq})_{f/1}^{\leq}$, donc aussi à $M_f^{\leq}$ en vertu de
(\hyperref[II.8.2.9.2-fr]{8.2.9.2}), d'où la conclusion en ce qui concerne
(\hyperref[II.8.12.6.1-fr]{8.12.6.1})~; la formule
(\hyperref[II.8.12.6.2-fr]{8.12.6.2}) résulte alors de
(\hyperref[II.8.12.6.1-fr]{8.12.6.1}) et de
(\hyperref[II.8.12.2.1-fr]{8.12.2.1}).
\end{proof}

\begin{corollary}[8.12.7]
\label{II.8.12.7-fr}
Avec les identifications de (\hyperref[II.8.6.3-fr]{8.6.3}), la restriction
de $(\mathcal{M}_X^{\geq})^{\square}$ à $\widehat{E}_X$ s'identifie à
$(\mathcal{M}_X^{\leq})^{\sim}$ et la restriction de
$(\mathcal{M}_X^{\geq})^{\square}$ à $E_X$ s'identifie à
$\widetilde{\mathcal{M}}_X$.
\end{corollary}

\begin{proof}
On est en effet ramené au cas affine, et cela résulte de l'identification de
$(M_f^{\geq})_{f/1}^{\leq}$ à $M_f^{\leq}$ et de
$(M_f^{\geq})_{f/1}$ à $M_f$ (\hyperref[II.8.2.9.2-fr]{8.2.9.2}).
\end{proof}

\begin{proposition}[8.12.8]
\label{II.8.12.8-fr}
Sous les hypothèses de (\hyperref[II.8.6.4-fr]{8.6.4}), l'homomorphisme
canonique (\hyperref[II.8.12.3.1-fr]{8.12.3.1}) est un isomorphisme.
\end{proposition}

\begin{proof}
Compte tenu du fait que $\operatorname{Proj}(\mathcal{S}_X^{\geq})\to X$
est un isomorphisme (\hyperref[II.8.6.2-fr]{8.6.2}) et des

\oldpage[II]{195}
isomorphismes (\hyperref[II.8.12.5.4-fr]{8.12.5.4}) et
(\hyperref[II.8.12.6.1-fr]{8.12.6.1}), on est ramené à démontrer la
proposition correspondante pour l'homomorphisme canonique
$p_X^*(\operatorname{Proj}_0(\mathcal{M}_X^{\geq}))
\to(\mathcal{M}_X^{\geq})^{\square}|\widehat{E}_X$, autrement dit, on est
ramené au cas où $\mathcal{S}_1$ est un $\mathcal{O}_Y$-Module inversible et
où $\mathcal{S}$ est engendrée par $\mathcal{S}_1$. Avec les notations de
(\hyperref[II.8.12.3-fr]{8.12.3}), on a alors, pour un $f\in S_1$,
$S_f^{\leq}=S_{(f)}[1/f]$ et l'homomorphisme canonique
$M_{(f)}\otimes_{S_{(f)}}S_f^{\leq}\to M_f^{\leq}$ est un isomorphisme par
définition de $M_f^{\leq}$.
\end{proof}

\begin{env}[8.12.9]
\label{II.8.12.9-fr}
Considérons maintenant les $\mathcal{S}$-Modules quasi-cohérents
\[
  \mathcal{M}_{[n]}=\bigoplus_{m\geqslant n}\mathcal{M}_m
\]
et (avec les notations de (\hyperref[II.8.7.2-fr]{8.7.2})) le
$\mathcal{S}^{\natural}$-Module gradué quasi-cohérent
\begin{equation}
\label{II.8.12.9.1-fr}
  \mathcal{M}^{\natural}
  =\left(\bigoplus_{n\geqslant0}\mathcal{M}_{[n]}\right)^{\sim}.
\tag{8.12.9.1}
\end{equation}
On a vu (\hyperref[II.8.7.3-fr]{8.7.3}) qu'il existe un $C$-isomorphisme
canonique $h:C_X\xrightarrow{\sim}\operatorname{Proj}(\mathcal{S}^{\natural})$.
En outre~:
\end{env}

\begin{proposition}[8.12.10]
\label{II.8.12.10-fr}
Il existe un $h$-isomorphisme canonique
\begin{equation}
\label{II.8.12.10.1-fr}
  \operatorname{Proj}_0(\mathcal{M}^{\natural})
  \xrightarrow{\sim}\widetilde{\mathcal{M}}_X.
\tag{8.12.10.1}
\end{equation}
\end{proposition}

\begin{proof}
On raisonne comme dans (\hyperref[II.8.7.3-fr]{8.7.3}), en utilisant cette
fois l'existence du di-isomorphisme (\hyperref[II.8.2.9.3-fr]{8.2.9.3}) au
lieu de (\hyperref[II.8.2.7.3-fr]{8.2.7.3}). Nous laissons les détails au
lecteur.
\end{proof}

\subsection{Fermetures projectives de sous-faisceaux et de sous-schémas
fermés}
\label{subsection:II.8.13-fr}

\begin{env}[8.13.1]
\label{II.8.13.1-fr}
Les hypothèses et notations étant celles de
(\hyperref[II.8.12.1-fr]{8.12.1}), considérons un sous-$\mathcal{S}$-Module
quasi-cohérent $\mathcal{N}$ de $\mathcal{M}$, \emph{non nécessairement
gradué}. On peut donc considérer le $\mathcal{O}_C$-Module quasi-cohérent
$\widetilde{\mathcal{N}}$ associé à $\mathcal{N}$, qui est un
sous-$\mathcal{O}_C$-Module de $\widetilde{\mathcal{M}}$. On a vu d'autre
part (\hyperref[II.8.12.2.1-fr]{8.12.2.1}) que
$\widetilde{\mathcal{M}}$ s'identifie à la restriction de
$\mathcal{M}^{\square}$ à $C$. Comme l'injection canonique
$i:C\to\widehat{C}$ est un morphisme affine
(\hyperref[II.8.3.2-fr]{8.3.2}), et \emph{a fortiori} quasi-compact, le
\emph{prolongement canonique} $(\widetilde{\mathcal{N}})^{-}$, plus grand
sous-$\mathcal{O}_{\widehat{C}}$-Module contenu dans
$\mathcal{M}^{\square}$ et induisant $\widetilde{\mathcal{N}}$ sur $C$, est
un $\mathcal{O}_{\widehat{C}}$-Module quasi-cohérent
(\hyperref[I.9.4.2-fr]{I, 9.4.2}). Nous allons en donner une description
explicite à l'aide d'un $\widehat{\mathcal{S}}$-Module gradué.
\end{env}

\begin{env}[8.13.2]
\label{II.8.13.2-fr}
Pour cela, considérons, pour tout entier $n\geqslant0$, l'homomorphisme
$\bigoplus_{i\leqslant n}\mathcal{M}_i\to\mathcal{M}$ qui, pour tout ouvert
$U$ de $Y$, fait correspondre à la famille
\[
  (s_i)\in\bigoplus_{i\leqslant n}\Gamma(U,\mathcal{M}_i)
\]
la section $\sum_i s_i\in\Gamma(U,\mathcal{M})$. Désignons par
$\mathcal{N}'_n$ l'image réciproque de $\mathcal{N}$ par cet homomorphisme,
qui est un sous-$\mathcal{O}_Y$-Module quasi-cohérent de
$\bigoplus_{i\leqslant n}\mathcal{M}_i$. Considérons maintenant
l'homomorphisme
$\bigoplus_{i\leqslant n}\mathcal{M}_i\to
\widehat{\mathcal{M}}=\mathcal{M}[z]$ qui fait correspondre à $(s_i)$ la
section
$\sum_{i\leqslant n}s_i z^{n-i}\in\Gamma(U,\widehat{\mathcal{M}}_n)$, et
soit $\mathcal{N}_n$ l'image de $\mathcal{N}'_n$ par cet homomorphisme~; on
vérifie aussitôt que
$\overline{\mathcal{N}}=\bigoplus_{n\geqslant0}\mathcal{N}_n$ est un
sous-$\widehat{\mathcal{S}}$-Module gradué (quasi-cohérent) de
$\widehat{\mathcal{M}}$~; on dit que $\overline{\mathcal{N}}$ est déduit de
$\mathcal{N}$ par \emph{homogénéisation}, à l'aide de la «~variable
homogénéisante~» $z$. On notera

\oldpage[II]{196}
que si $\mathcal{N}$ est déjà un sous-$\mathcal{S}$-Module \emph{gradué} de
$\mathcal{M}$, alors $\overline{\mathcal{N}}$ s'identifie à la somme directe
des composantes $\widehat{\mathcal{N}}_n$ de degré $n\geqslant0$ dans
$\widehat{\mathcal{N}}=\mathcal{N}[z]$.
\end{env}

\begin{proposition}[8.13.3]
\label{II.8.13.3-fr}
Le $\mathcal{O}_{\widehat{C}}$-Module
$\operatorname{Proj}_0(\overline{\mathcal{N}})$ est le prolongement canonique
$(\widetilde{\mathcal{N}})^{-}$ de $\widetilde{\mathcal{N}}$ à
$\widehat{C}$.
\end{proposition}

\begin{proof}
La question est locale sur $Y$ et $\widehat{C}$ en vertu de la définition du
prolongement canonique (\hyperref[I.9.4.1-fr]{I, 9.4.1}). On peut donc déjà
supposer que $Y=\operatorname{Spec}(A)$ est affine, avec
$\mathcal{S}=\widetilde{S}$, $\mathcal{M}=\widetilde{M}$ et
$\mathcal{N}=\widetilde{N}$, où $N$ est un sous-$S$-module non
nécessairement gradué de $M$. En outre
(\hyperref[II.8.3.2.6-fr]{8.3.2.6}) $\widehat{C}$ est réunion des ouverts
affines $\widehat{C}_z=C$ et
$\widehat{C}_f=\operatorname{Spec}(S_f^{\leq})$ ($f$ homogène dans $S_+$).
Il suffira donc de montrer que~: $1^\circ$ la restriction de
$\operatorname{Proj}_0(\overline{\mathcal{N}})$ à $C$ est
$\widetilde{\mathcal{N}}$~; $2^\circ$ la restriction de
$\operatorname{Proj}_0(\overline{\mathcal{N}})$ à chaque $\widehat{C}_f$ est
le prolongement canonique de la restriction de $\mathcal{N}$ à
$C\cap\widehat{C}_f=\operatorname{Spec}(S_f)$
(\hyperref[II.8.3.2.6-fr]{8.3.2.6}). Pour le premier point, notons que
$\operatorname{Proj}_0(\overline{\mathcal{N}})|C$ s'identifie à
$(\overline{N}_{(z)})^{\sim}$
(\hyperref[II.8.3.2.4-fr]{8.3.2.4})~; or $\overline{N}_{(z)}$ s'identifie
canoniquement (\hyperref[II.2.2.5-fr]{2.2.5}) à l'image de
$\overline{N}$ dans $\widehat{M}/(z-1)\widehat{M}$, et par l'isomorphisme
canonique de ce dernier sur $M$ (\hyperref[II.8.2.5-fr]{8.2.5}), cette image
s'identifie à $N$, en vertu de la définition de $\overline{N}$ donnée dans
(\hyperref[II.8.13.2-fr]{8.13.2}).

Pour prouver le second point, notons que l'injection
$i:C\cap\widehat{C}_f\to\widehat{C}_f$ correspond à l'injection canonique
$S_f^{\leq}\to S_f$ (\hyperref[II.8.3.2.6-fr]{8.3.2.6})~; d'autre part, on
a $\Gamma(\widehat{C}_f,\mathcal{M}^{\square})=M_f^{\leq}$,
$\Gamma(\widehat{C}_f,i_*(\widetilde{\mathcal{N}}))=N_f$ et (par
(\hyperref[II.8.12.2.1-fr]{8.12.2.1}))
$\Gamma(\widehat{C}_f,i_*(i^*(\mathcal{M}^{\square})))=M_f$. Compte tenu de
(\hyperref[I.9.4.2-fr]{I, 9.4.2}), on est donc ramené à montrer que
$\overline{N}_{(f)}\subset\widehat{M}_{(f)}=M_f^{\leq}$ s'identifie
canoniquement à l'image réciproque de $N_f$, par l'injection canonique
$M_f^{\leq}\to M_f$. En effet, soit $d=\deg(f)>0$, et supposons qu'un
élément $(\sum_{k\leqslant md}x_k)/f^m$ de $M_f$ (avec $x_k\in M_k$) soit de
la forme $y/f^m$ avec $y\in N$. En multipliant $y$ et les $x_k$ par un même
$f^h$ convenable, on peut déjà supposer que $\sum_{k\leqslant md}x_k=y$.
Mais dans l'identification (\hyperref[II.8.2.5.2-fr]{8.2.5.2}),
$(\sum_{k\leqslant md}x_k)/f^m$ correspond à
$\sum_{k\leqslant md}x_kz^{md-k}/f^m$, et cet élément appartient bien à
$\overline{N}_{(f)}$, puisque $\sum_{k\leqslant md}x_k\in N$~; la réciproque
est évidente.
\end{proof}

\begin{remark}[8.13.4]
\label{II.8.13.4-fr}
\begin{enumerate}
  \item[\rm{(i)}] Le cas d'application le plus important de
  (\hyperref[II.8.13.3-fr]{8.13.3}) est celui où
  $\mathcal{M}=\mathcal{S}$, $\widetilde{\mathcal{N}}$ étant donc un faisceau
  quasi-cohérent \emph{arbitraire} d'idéaux $\mathcal{J}$ de
  $\mathcal{O}_C$ (\hyperref[II.1.4.3-fr]{1.4.3}), correspondant
  biunivoquement à un \emph{sous-préschéma fermé} $Z$ de $C$. Alors le
  prolongement canonique $\overline{\mathcal{J}}$ de $\mathcal{J}$ est le
  faisceau quasi-cohérent d'idéaux de $\mathcal{O}_{\widehat{C}}$ qui définit
  l'\emph{adhérence} $\overline{Z}$ de $Z$ dans $\widehat{C}$
  (\hyperref[I.9.5.10-fr]{I, 9.5.10})~; la prop.
  (\hyperref[II.8.13.3-fr]{8.13.3}) donne un moyen canonique de définir
  $\overline{Z}$ à l'aide d'un Idéal gradué dans
  $\widehat{\mathcal{S}}=\mathcal{S}[z]$.
  \item[\rm{(ii)}] Supposons pour simplifier que $Y$ soit affine, et gardons
  les notations de la démonstration de
  (\hyperref[II.8.13.3-fr]{8.13.3}). Pour tout $x\in N$ non nul, soit $d(x)$
  le plus grand des degrés des composantes homogènes $x_i$ de $x$ dans $M$~;
  $\overline{N}$ est par définition le sous-module de $\widehat{M}$ formé de
  $0$ et des éléments de la forme
  $h(x,k)=z^k\sum_{i\leqslant d(x)}x_i z^{d(x)-i}$ ($k$ entier
  $\geqslant0$)~; il est donc engendré, en tant que module sur
  $\widehat{S}=S[z]$, par les éléments de la forme
  \[
    h(x,0)=\sum_{i\leqslant d(x)}x_i z^{d(x)-i}.
  \]

  \oldpage[II]{197}
  On dit que $h(x,0)$ est déduit de $x$ par \emph{homogénéisation} à l'aide
  de la «~variable homogénéisante~» $z$. Mais comme $h(x,0)$ ne dépend pas
  additivement de $x$ (ni \emph{a fortiori} $S$-linéairement), \emph{on se
  gardera de croire} (même lorsque $M=S$) que les $h(x,0)$ parcourent un
  \emph{système de générateurs} du $\widehat{S}$-module gradué
  $\overline{N}$ lorsqu'on fait parcourir à $x$ un \emph{système de
  générateurs} du $S$-module $N$. Il en est cependant bien ainsi dans le cas
  (seul considéré dans la géométrie algébrique élémentaire) où $N$ est un
  $S$-module \emph{libre monogène}, car si $t$ est une base de $N$, $h(t,0)$
  engendre le $\widehat{S}$-module $\overline{N}$.
\end{enumerate}
\end{remark}

\subsection{Compléments sur les faisceaux associés aux
$\mathcal{S}$-Modules gradués}
\label{subsection:II.8.14-fr}

\begin{env}[8.14.1]
\label{II.8.14.1-fr}
Soient $Y$ un préschéma, $\mathcal{S}$ une $\mathcal{O}_Y$-Algèbre graduée
quasi-cohérente à \emph{degrés positifs},
$X=\operatorname{Proj}(\mathcal{S})$, $q:X\to Y$ le morphisme structural
(séparé en vertu de (\hyperref[II.3.1.3-fr]{3.1.3})). Conservons les
notations de (\hyperref[II.8.12.1-fr]{8.12.1}), de sorte que nous avons
défini un foncteur
$\mathcal{M}_X=\operatorname{Proj}(\mathcal{M})$ en $\mathcal{M}$, de la
catégorie des $\mathcal{S}$-Modules gradués quasi-cohérents dans celle des
$\mathcal{S}_X$-Modules gradués quasi-cohérents~; il est clair en outre
(\hyperref[II.3.2.4-fr]{3.2.4}) que c'est un foncteur \emph{additif exact},
commutant aux limites inductives.

Notons en outre qu'il résulte aussitôt de la définition
(\hyperref[II.8.12.1.1-fr]{8.12.1.1}) que l'on a
\begin{equation}
\label{II.8.14.1.1-fr}
  \operatorname{Proj}(\mathcal{M}(n))
  =(\operatorname{Proj}(\mathcal{M}))(n)
  \qquad\text{pour tout }n\in\mathbf{Z}.
\tag{8.14.1.1}
\end{equation}
\end{env}

\begin{env}[8.14.2]
\label{II.8.14.2-fr}
Nous allons d'abord étendre aux $\mathcal{S}_X$-Modules de la forme
$\operatorname{Proj}(\mathcal{M})$ les homomorphismes canoniques $\lambda$
et $\mu$ définis dans (\hyperref[II.3.2.6-fr]{3.2.6}). Pour cela, notons
que, quels que soient $m\in\mathbf{Z}$ et $n\in\mathbf{Z}$, on a, en vertu
de (\hyperref[II.2.1.2.1-fr]{2.1.2.1}), un homomorphisme canonique de
$\mathcal{O}_X$-Modules
\begin{equation}
\label{II.8.14.2.1-fr}
  \lambda_{mn}:\operatorname{Proj}_0(\mathcal{M}(m))
  \otimes_{\mathcal{O}_X}\operatorname{Proj}_0(\mathcal{N}(n))
  \to\operatorname{Proj}_0
  ((\mathcal{M}\otimes_{\mathcal{S}}\mathcal{N})(m+n))
\tag{8.14.2.1}
\end{equation}
et de même, en vertu de (\hyperref[II.2.1.2.2-fr]{2.1.2.2}), un
homomorphisme canonique de $\mathcal{O}_X$-Modules
\begin{equation}
\label{II.8.14.2.2-fr}
  \mu_{mn}:\operatorname{Proj}_0
  ((\mathcal{H}om_{\mathcal{S}}(\mathcal{M},\mathcal{N}))(n-m))
  \to\mathcal{H}om_{\mathcal{O}_X}
  (\operatorname{Proj}_0(\mathcal{M}(m)),
   \operatorname{Proj}_0(\mathcal{N}(n)))
\tag{8.14.2.2}
\end{equation}
quels que soient les $\mathcal{S}$-Modules gradués quasi-cohérents
$\mathcal{M}$, $\mathcal{N}$. On en déduit un homomorphisme
\[
  \mu_k:\operatorname{Proj}_0
  ((\mathcal{H}om_{\mathcal{S}}(\mathcal{M},\mathcal{N}))(k))
  \to
  (\mathcal{H}om_{\mathcal{S}_X}
  (\operatorname{Proj}(\mathcal{M}),\operatorname{Proj}(\mathcal{N})))_k
\]
obtenu en faisant correspondre à tout
$u\in\Gamma(U,\operatorname{Proj}_0
((\mathcal{H}om_{\mathcal{S}}(\mathcal{M},\mathcal{N}))(k)))$
l'homomorphisme $\mu_k(u)$, de degré $k$, de $\mathbf{Z}$-modules gradués
$\Gamma(U,\operatorname{Proj}(\mathcal{M}))\to
\Gamma(U,\operatorname{Proj}(\mathcal{N}))$ ($U$ ouvert de $X$) qui dans
chaque $\Gamma(U,\operatorname{Proj}_0(\mathcal{M}(m)))$ coïncide avec
$\mu_{m,m+k}(u)$~; en outre, en revenant à la définition des $\mu_{mn}$
(\hyperref[II.2.5.12.1-fr]{2.5.12.1}), on voit aussitôt que $\mu_k(u)$ est
en fait un homomorphisme de degré $k$ de
$\Gamma(U,\mathcal{S}_X)$-modules gradués, et en outre que les $\mu_k$
définissent un homomorphisme de $\mathcal{S}_X$-Modules gradués
\begin{equation}
\label{II.8.14.2.3-fr}
  \operatorname{Proj}
  (\mathcal{H}om_{\mathcal{S}}(\mathcal{M},\mathcal{N}))
  \to\mathcal{H}om_{\mathcal{S}_X}
  (\operatorname{Proj}(\mathcal{M}),\operatorname{Proj}(\mathcal{N})).
\tag{8.14.2.3}
\end{equation}

De même, compte tenu du diagramme d'associativité
(\hyperref[II.2.5.11.4-fr]{2.5.11.4}), les homomorphismes
(\hyperref[II.8.14.2.1-fr]{8.14.2.1}) donnent un homomorphisme de
$\mathcal{S}_X$-Modules gradués
\begin{equation}
\label{II.8.14.2.4-fr}
  \lambda:\operatorname{Proj}(\mathcal{M})
  \otimes_{\mathcal{S}_X}\operatorname{Proj}(\mathcal{N})
  \to\operatorname{Proj}(\mathcal{M}\otimes_{\mathcal{S}}\mathcal{N}).
\tag{8.14.2.4}
\end{equation}
\end{env}

\oldpage[II]{198}
\begin{proposition}[8.14.3]
\label{II.8.14.3-fr}
L'homomorphisme (\hyperref[II.8.14.2.4-fr]{8.14.2.4}) est bijectif~; il en
est de même de (\hyperref[II.8.14.2.3-fr]{8.14.2.3}) lorsque le
$\mathcal{S}$-Module gradué $\mathcal{M}$ admet une présentation finie
(\hyperref[II.3.1.1-fr]{3.1.1}).
\end{proposition}

\begin{proof}
La question est évidemment locale sur $X$ et sur $Y$~; on peut donc supposer
$Y=\operatorname{Spec}(A)$ affine, $\mathcal{S}=\widetilde{S}$,
$\mathcal{M}=\widetilde{M}$, $\mathcal{N}=\widetilde{N}$, où $S$ est une
$A$-algèbre graduée à degrés positifs, $M$ et $N$ deux $S$-modules gradués.
Si $f$ est un élément homogène de $S_+$, les homomorphismes
(\hyperref[II.8.14.2.1-fr]{8.14.2.1}) et
(\hyperref[II.8.14.2.2-fr]{8.14.2.2}), restreints à l'ouvert affine
$D_+(f)$, correspondent aux homomorphismes canoniques
(\hyperref[II.2.5.11.1-fr]{2.5.11.1}) et
(\hyperref[II.2.5.12.1-fr]{2.5.12.1})~:
\begin{align*}
  M(m)_{(f)}\otimes_{S_{(f)}}N(n)_{(f)}
  &\to(M\otimes_S N)(m+n)_{(f)},\\
  (\operatorname{Hom}_S(M,N))(n-m)_{(f)}
  &\to\operatorname{Hom}_{S_{(f)}}(M(m)_{(f)},N(n)_{(f)}).
\end{align*}
Si on se reporte aux définitions de ces homomorphismes, on voit donc (compte
tenu de (\hyperref[II.8.2.9.1-fr]{8.2.9.1})) que la restriction de
(\hyperref[II.8.14.2.4-fr]{8.14.2.4}) à $D_+(f)$ correspond à
l'homomorphisme canonique
\[
  M_f\otimes_{S_f}N_f\to(M\otimes_S N)_f
\]
défini dans (\hyperref[0.1.3.4-fr]{0, 1.3.4}), et on sait que ce dernier est
un isomorphisme. De même, la restriction de
(\hyperref[II.8.14.2.3-fr]{8.14.2.3}) à $D_+(f)$ correspond à
l'homomorphisme canonique (\hyperref[0.1.3.5-fr]{0, 1.3.5})
\[
  (\operatorname{Hom}_S(M,N))_f\to
  \operatorname{Hom}_{S_f}(M_f,N_f)
\]
compte tenu de ce que, $M$ étant de type fini, le module
$\operatorname{Hom}_S(M,N)$, somme directe des sous-groupes formés des
homomorphismes \emph{homogènes} de $S$-modules
(\hyperref[II.2.1.2-fr]{2.1.2}), coïncide avec l'ensemble de \emph{tous} les
homomorphismes $M\to N$ de $S$-modules. L'hypothèse que $M$ admet une
présentation finie entraîne alors (\hyperref[0.1.3.5-fr]{0, 1.3.5}) que
l'homomorphisme canonique considéré est bien un isomorphisme.
\end{proof}

\begin{proposition}[8.14.4]
\label{II.8.14.4-fr}
Si $U$ est un ouvert quasi-compact de $X$, il existe un entier $d$ tel que
pour tout entier $n$ multiple de $d$, $\mathcal{O}_X(n)|U$ soit inversible,
son inverse étant $\mathcal{O}_X(-n)|U$.
\end{proposition}

\begin{proof}
Comme $q(U)$ est quasi-compact, il est recouvert par un nombre fini d'ouverts
affines $V_i$, donc tout $x\in U$ est contenu dans un ouvert affine de la
forme $D_+(f)$, où $f$ est un élément homogène de degré $>0$ de l'un des
anneaux $\Gamma(V_i,\mathcal{S})$. Comme $U$ est quasi-compact, on peut le
recouvrir par un nombre fini de tels ouverts $D_+(f_j)$~; soit $d$ un
multiple commun des degrés des $f_j$. Le nombre $d$ répond à la question en
vertu de (\hyperref[II.2.5.17-fr]{2.5.17}).
\end{proof}

\begin{env}[8.14.5]
\label{II.8.14.5-fr}
Avec les hypothèses et notations de (\hyperref[II.8.14.1-fr]{8.14.1}), on a
défini dans (\hyperref[II.3.3.2-fr]{3.3.2}) des homomorphismes canoniques de
$\mathcal{O}_Y$-Modules
\begin{equation}
\label{II.8.14.5.1-fr}
  \alpha_n:\mathcal{M}_n\to
  q_*(\operatorname{Proj}_0(\mathcal{M}(n)))
  \qquad(n\in\mathbf{Z})
\tag{8.14.5.1}
\end{equation}
Généralisant les notations de (\hyperref[II.3.3.1-fr]{3.3.1}), nous
poserons, pour tout $\mathcal{S}_X$-Module gradué $\mathcal{F}$,
\begin{equation}
\label{II.8.14.5.2-fr}
  \Gamma_*(\mathcal{F})=
  \bigoplus_{n\in\mathbf{Z}}q_*(\mathcal{F}_n).
\tag{8.14.5.2}
\end{equation}
En particulier,
$\Gamma_*(\mathcal{S}_X)=\bigoplus_{n\in\mathbf{Z}}q_*(\mathcal{O}_X(n))$
est la $\mathcal{O}_Y$-Algèbre graduée notée
$\Gamma_*(\mathcal{O}_X)$ dans
(\hyperref[II.3.3.1.2-fr]{3.3.1.2})~; il est clair que
$\Gamma_*(\mathcal{F})$ est un $\Gamma_*(\mathcal{S}_X)$-Module gradué
(\hyperref[0.4.2.2-fr]{0, 4.2.2}). Lorsque

\oldpage[II]{199}
dans les homomorphismes (\hyperref[II.8.14.5.1-fr]{8.14.5.1}), on prend
$\mathcal{M}=\mathcal{S}$, on obtient l'homomorphisme de
$\mathcal{O}_Y$-Algèbres graduées
\begin{equation}
\label{II.8.14.5.3-fr}
  \alpha:\mathcal{S}\to\Gamma_*(\mathcal{S}_X)
\tag{8.14.5.3}
\end{equation}
déjà défini dans (\hyperref[II.3.3.2-fr]{3.3.2}), et qui fait donc de
$\Gamma_*(\mathcal{F})$ un $\mathcal{S}$-Module gradué~; les homomorphismes
(\hyperref[II.8.14.5.1-fr]{8.14.5.1}) définissent alors un homomorphisme (de
degré $0$) de $\mathcal{S}$-Modules gradués
\begin{equation}
\label{II.8.14.5.4-fr}
  \alpha:\mathcal{M}\to\Gamma_*(\operatorname{Proj}(\mathcal{M})).
\tag{8.14.5.4}
\end{equation}
\end{env}

\begin{env}[8.14.6]
\label{II.8.14.6-fr}
En général, pour un $\mathcal{S}_X$-Module gradué quasi-cohérent
$\mathcal{F}$, il n'est pas certain que le $\mathcal{S}$-Module gradué
$\Gamma_*(\mathcal{F})$ soit nécessairement quasi-cohérent. Considérons un
ouvert $X'$ de $X$ tel que la restriction $q':X'\to Y$ de $q$ à $X'$ soit
un morphisme \emph{quasi-compact}. Comme $q'$ est en outre séparé,
$q'_*(\mathcal{F}')$ est alors un $\mathcal{O}_Y$-Module quasi-cohérent pour
tout $\mathcal{O}_{X'}$-Module quasi-cohérent $\mathcal{F}'$
(\hyperref[I.9.2.2-fr]{I, 9.2.2, b)}). Nous poserons
\begin{equation}
\label{II.8.14.6.1-fr}
  \mathcal{S}_{X'}=\mathcal{S}_X|X'
  =\bigoplus_{n\in\mathbf{Z}}\mathcal{O}_X(n)|X'
\tag{8.14.6.1}
\end{equation}
et, pour tout $\mathcal{S}_{X'}$-Module gradué $\mathcal{F}'$,
\begin{equation}
\label{II.8.14.6.2-fr}
  \Gamma'_*(\mathcal{F}')=
  \bigoplus_{n\in\mathbf{Z}}q'_*(\mathcal{F}'_n).
\tag{8.14.6.2}
\end{equation}
La remarque précédente montre alors que si $\mathcal{F}'$ est un
$\mathcal{S}_{X'}$-Module quasi-cohérent, $\Gamma'_*(\mathcal{F}')$ est un
$\mathcal{S}$-Module gradué quasi-cohérent
(\hyperref[I.9.6.1-fr]{I, 9.6.1}).

On notera d'ailleurs que l'injection canonique $j:X'\to X$ est
\emph{quasi-compacte} puisque $q'=q\circ j$ l'est et que $q$ est séparé
(\hyperref[I.6.6.4-fr]{I, 6.6.4, (v)}). Par suite
$\mathcal{F}=j_*(\mathcal{F}')$ est un $\mathcal{S}_X$-Module gradué
quasi-cohérent pour tout $\mathcal{S}_{X'}$-Module gradué quasi-cohérent
$\mathcal{F}'$, et il résulte des définitions précédentes que l'on a
\begin{equation}
\label{II.8.14.6.3-fr}
  \Gamma'_*(\mathcal{F}')=\Gamma_*(\mathcal{F}).
\tag{8.14.6.3}
\end{equation}
Avec les mêmes hypothèses sur $X'$, pour tout $\mathcal{S}$-Module gradué
quasi-cohérent $\mathcal{M}$, nous poserons
\begin{equation}
\label{II.8.14.6.4-fr}
  \operatorname{Proj}'(\mathcal{M})
  =\operatorname{Proj}(\mathcal{M})|X'
\tag{8.14.6.4}
\end{equation}
qui est un $\mathcal{S}_{X'}$-Module gradué quasi-cohérent. L'homomorphisme
canonique
\[
  \operatorname{Proj}(\mathcal{M})\to
  j_*(\operatorname{Proj}'(\mathcal{M}))
\]
(\hyperref[0.4.4.3-fr]{0, 4.4.3}) donne donc un homomorphisme canonique
$\Gamma_*(\operatorname{Proj}(\mathcal{M}))\to
\Gamma'_*(\operatorname{Proj}'(\mathcal{M}))$ de $\mathcal{S}$-Modules
gradués, et par suite, en le composant avec
(\hyperref[II.8.14.5.4-fr]{8.14.5.4}), on obtient un homomorphisme canonique
fonctoriel (de degré $0$) de $\mathcal{S}$-Modules gradués quasi-cohérents
\begin{equation}
\label{II.8.14.6.5-fr}
  \alpha':\mathcal{M}\to
  \Gamma'_*(\operatorname{Proj}'(\mathcal{M})).
\tag{8.14.6.5}
\end{equation}
\end{env}

\begin{env}[8.14.7]
\label{II.8.14.7-fr}
Conservons sur $X'$ les hypothèses de
(\hyperref[II.8.14.6-fr]{8.14.6}), et soit $\mathcal{F}'$ un
$\mathcal{S}_{X'}$-Module gradué quasi-cohérent, de sorte que
$\operatorname{Proj}'(\Gamma'_*(\mathcal{F}'))$ est aussi un
$\mathcal{S}_{X'}$-Module gradué quasi-cohérent. Nous allons définir un
homomorphisme canonique fonctoriel (de degré $0$) de
$\mathcal{S}_{X'}$-Modules gradués
\begin{equation}
\label{II.8.14.7.1-fr}
  \beta':\operatorname{Proj}'(\Gamma'_*(\mathcal{F}'))\to\mathcal{F}'.
\tag{8.14.7.1}
\end{equation}

\oldpage[II]{200}
Supposons d'abord $Y=\operatorname{Spec}(A)$ affine,
$\mathcal{S}=\widetilde{S}$, où $S$ est une $A$-algèbre graduée à degrés
positifs~; alors $\Gamma'_*(\mathcal{F}')=\widetilde{M}$, où
$M=\bigoplus_{n\in\mathbf{Z}}\Gamma(X',\mathcal{F}'_n)$ est un $S$-module
gradué. Soit $f\in S_d$ tel que $D_+(f)\subset X'$~; par définition
(\hyperref[II.2.6.2-fr]{2.6.2}), $\alpha_d(f)$ restreinte à $D_+(f)$ est la
section de $\mathcal{O}_X(d)$ au-dessus de $D_+(f)$ correspondant à
l'élément $f/1$ de $(S(d))_{(f)}$ et est par suite inversible~; il en est donc
de même de $\alpha_d(f^n)$ pour tout $n>0$. On en conclut aussitôt que
l'on définit un $S_f$-homomorphisme (de degré $0$) de modules gradués
$\beta_f:M_f\to\Gamma(D_+(f),\mathcal{F}')$ en faisant correspondre à tout
élément $z/f^n\in M_f$ (où $z\in M$), la section
$(z|D_+(f))(\alpha_d(f^n)|D_+(f))^{-1}$ de $\mathcal{F}'$ au-dessus de
$D_+(f)$. En outre, on a un diagramme commutatif correspondant à
(\hyperref[II.2.6.4.1-fr]{2.6.4.1}), d'où la définition de $\beta'$ dans ce
cas. Pour passer au cas général, il faut considérer une $A$-algèbre $A'$,
la $A'$-algèbre graduée $S'=S\otimes_A A'$, et utiliser un diagramme
commutatif analogue à (\hyperref[II.2.8.13.2-fr]{2.8.13.2})~; nous laissons
les détails aux lecteurs.
\end{env}

\begin{proposition}[8.14.8]
\label{II.8.14.8-fr}
Si $X'$ est un ouvert de $X=\operatorname{Proj}(\mathcal{S})$ tel que
$q':X'\to Y$ soit quasi-compact, l'homomorphisme $\beta'$ défini dans
(\hyperref[II.8.14.7-fr]{8.14.7}) est bijectif.
\end{proposition}

\begin{proof}
On peut évidemment se restreindre au cas où $Y$ est affine et tout revient à
prouver (avec les notations de (\hyperref[II.8.14.7-fr]{8.14.7})) que
l'homomorphisme $\beta_f:M_f\to\Gamma(D_+(f),\mathcal{F}')$ est un
isomorphisme. Or, lorsqu'on remplace $f$ par une de ses puissances, on ne
change pas $D_+(f)$ ni $\beta_f$~; comme $X'$ est \emph{quasi-compact} par
hypothèse, on peut donc toujours supposer, en vertu de
(\hyperref[II.8.14.4-fr]{8.14.4}), que le faisceau $\mathcal{O}_X(d)$ est
\emph{inversible}. Comme $X'$ est un schéma (puisque $q'$ est séparé), la
proposition n'est autre alors que (\hyperref[I.9.3.1-fr]{I, 9.3.1}).
\end{proof}

\begin{corollary}[8.14.9]
\label{II.8.14.9-fr}
Sous les hypothèses de (\hyperref[II.8.14.8-fr]{8.14.8}), tout
$\mathcal{S}_{X'}$-Module gradué quasi-cohérent $\mathcal{F}'$ est isomorphe
à un $\mathcal{S}_{X'}$-Module gradué de la forme
$\operatorname{Proj}'(\mathcal{M})$, où $\mathcal{M}$ est un
$\mathcal{S}$-Module gradué quasi-cohérent. Si en outre $\mathcal{F}'$ est de
type fini, et si on suppose que $Y$ est un schéma quasi-compact ou un
préschéma dont l'espace sous-jacent est noethérien, on peut supposer
$\mathcal{M}$ de type fini.
\end{corollary}

\begin{proof}
La démonstration à partir de (\hyperref[II.8.14.8-fr]{8.14.8}) suit
exactement la marche de celle de (\hyperref[II.3.4.5-fr]{3.4.5}) à partir de
(\hyperref[II.3.4.4-fr]{3.4.4}), et nous en laissons les détails au lecteur.
\end{proof}

\begin{proposition}[8.14.10]
\label{II.8.14.10-fr}
Sous les hypothèses de (\hyperref[II.8.14.7-fr]{8.14.7}), soient
$\mathcal{M}$ un $\mathcal{S}$-Module gradué quasi-cohérent,
$\mathcal{F}'$ un $\mathcal{S}_{X'}$-Module gradué quasi-cohérent~; les
homomorphismes composés
\begin{equation}
\label{II.8.14.10.1-fr}
  \operatorname{Proj}'(\mathcal{M})
  \xrightarrow{\operatorname{Proj}'(\alpha')}
  \operatorname{Proj}'(\Gamma'_*(\operatorname{Proj}'(\mathcal{M})))
  \xrightarrow{\beta'}\operatorname{Proj}'(\mathcal{M})
\tag{8.14.10.1}
\end{equation}
\begin{equation}
\label{II.8.14.10.2-fr}
  \Gamma'_*(\mathcal{F}')
  \xrightarrow{\alpha'}
  \Gamma'_*(\operatorname{Proj}'(\Gamma'_*(\mathcal{F}')))
  \xrightarrow{\Gamma'_*(\beta')}\Gamma'_*(\mathcal{F}')
\tag{8.14.10.2}
\end{equation}
sont les isomorphismes identiques.
\end{proposition}

\begin{proof}
La question est locale sur $Y$ et la vérification se fait comme dans
(\hyperref[II.2.6.5-fr]{2.6.5})~; nous en laissons les détails au lecteur.
\end{proof}

\begin{remark}[8.14.11]
\label{II.8.14.11-fr}
Au chapitre III (\hyperref[III.2.3.1-fr]{III, 2.3.1}), nous verrons que
lorsque $Y$ est \emph{localement noethérien} et $\mathcal{S}$ une
$\mathcal{O}_Y$-Algèbre graduée quasi-cohérente \emph{de type fini} (auquel
cas

\oldpage[II]{201}
on peut prendre $X'=X$), alors l'homomorphisme $\alpha$
(\hyperref[II.8.14.5.4-fr]{8.14.5.4}) est \textbf{(TN)}-\emph{bijectif} pour
tout $\mathcal{S}$-Module gradué quasi-cohérent $\mathcal{M}$ vérifiant la
condition \textbf{(TF)}.
\end{remark}

\begin{remark}[8.14.12]
\label{II.8.14.12-fr}
La situation décrite dans (\hyperref[II.8.14.4-fr]{8.14.4}) est un cas
particulier de la suivante. Soient $X$ un espace annelé, $\mathcal{S}$ une
$\mathcal{O}_X$-Algèbre graduée (à degrés positifs et négatifs)~; supposons
qu'il existe un entier $d>0$ tel que $\mathcal{S}_d$ et
$\mathcal{S}_{-d}$ soient \emph{inversibles}, l'homomorphisme canonique
\begin{equation}
\label{II.8.14.12.1-fr}
  \mathcal{S}_d\otimes_{\mathcal{O}_X}\mathcal{S}_{-d}
  \to\mathcal{O}_X
\tag{8.14.12.1}
\end{equation}
étant un \emph{isomorphisme} (de sorte que $\mathcal{S}_{-d}$ s'identifie à
$\mathcal{S}_d^{-1}$). On dit alors que la $\mathcal{O}_X$-Algèbre graduée
$\mathcal{S}$ est \emph{périodique de période $d$}. Cette dénomination
provient de la propriété suivante~: \emph{sous les hypothèses précédentes,
pour tout $\mathcal{S}$-Module gradué $\mathcal{F}$, l'homomorphisme
canonique}
\begin{equation}
\label{II.8.14.12.2-fr}
  \mathcal{S}_d\otimes\mathcal{F}_n\to\mathcal{F}_{n+d}
\tag{8.14.12.2}
\end{equation}
\emph{est un isomorphisme pour tout $n\in\mathbf{Z}$.} En effet, la question
est locale sur $X$ et on peut donc supposer que $\mathcal{S}_d$ possède une
section \emph{inversible} $s$ au-dessus de $X$, son inverse $s'$ étant une
section de $\mathcal{S}_{-d}$. L'homomorphisme
$\mathcal{F}_{n+d}\to\mathcal{S}_d\otimes\mathcal{F}_n$ qui, à toute section
$z\in\Gamma(U,\mathcal{F}_{n+d})$ fait correspondre la section
$(s|U)\otimes(s'|U)z$ de $\mathcal{S}_d\otimes\mathcal{F}_n$ sur $U$, est
alors réciproque de (\hyperref[II.8.14.12.2-fr]{8.14.12.2}), d'où notre
assertion. On en déduit pour tout $k\in\mathbf{Z}$ un isomorphisme canonique
\[
  (\mathcal{S}_d)^{\otimes k}\otimes\mathcal{F}_n
  \xrightarrow{\sim}\mathcal{F}_{n+kd}.
\]
Par suite, \emph{un $\mathcal{S}$-Module gradué $\mathcal{F}$ est connu
lorsqu'on connaît les $\mathcal{S}_0$-Modules $\mathcal{F}_i$
$(0\leqslant i\leqslant d-1)$ et les homomorphismes canoniques}
\[
  \mathcal{S}_i\otimes\mathcal{F}_j\to\mathcal{F}_{i+j}
  \qquad\text{pour }0\leqslant i,j\leqslant d-1
\]
(en posant
$\mathcal{F}_{i+j}=\mathcal{S}_d\otimes_{\mathcal{S}_0}
\mathcal{F}_{i+j-d}$ lorsque $i+j\geqslant d$). Bien entendu, pour que ces
homomorphismes définissent bien une structure de $\mathcal{S}$-Module sur la
somme directe des $(\mathcal{S}_d)^{\otimes k}\otimes\mathcal{F}_i$
$(k\in\mathbf{Z},\ 0\leqslant i\leqslant d-1)$, ils doivent satisfaire à des
conditions d'associativité que nous n'expliciterons pas.

Dans le cas où $d=1$ (qui est celui considéré dans
(\hyperref[subsection:II.3.3-fr]{3.3})) on peut donc dire que la catégorie des
$\mathcal{S}$-Modules gradués (resp. quasi-cohérents si $X$ est un préschéma
et $\mathcal{S}$ est quasi-cohérente) est \emph{équivalente} à la catégorie
des $\mathcal{S}_0$-Modules quelconques (resp. quasi-cohérents)~; c'est dans
ce sens qu'on peut considérer les développements de ce numéro comme
généralisant ceux du \S~\hyperref[section:II.3-fr]{3}. En outre, on voit que,
sous des conditions de finitude convenables, les résultats de ce numéro
(joints à (\hyperref[II.8.14.11-fr]{8.14.11})) ramènent dans une certaine
mesure l'étude des Algèbres graduées quasi-cohérentes sur un préschéma, et des
Modules gradués «~mod. \textbf{(TN)}~» sur de telles Algèbres, à l'étude du
cas particulier où les Algèbres considérées sont \emph{périodiques} (et où
la condition \textbf{(TN)} pour $\mathcal{M}$
(\hyperref[II.3.4.2-fr]{3.4.2}) implique donc $\mathcal{M}=0$).
\end{remark}

\begin{remark}[8.14.13]
\label{II.8.14.13-fr}
Sous les hypothèses de (\hyperref[II.8.14.1-fr]{8.14.1}), soit $d$ un entier
$>0$~; on a défini un $Y$-isomorphisme canonique $h$ de $X$ sur
$X^{(d)}=\operatorname{Proj}(\mathcal{S}^{(d)})$
(\hyperref[II.3.1.8-fr]{3.1.8}). Pour tout

\oldpage[II]{202}
$\mathcal{S}$-Module gradué quasi-cohérent $\mathcal{M}$ et tout entier $k$
tel que $0\leqslant k\leqslant d-1$, on a aussi un $h$-isomorphisme canonique
(avec les notations de (\hyperref[II.3.1.1-fr]{3.1.1}))
\begin{equation}
\label{II.8.14.13.1-fr}
  (\operatorname{Proj}(\mathcal{M}))^{(d,k)}
  \xleftarrow{\sim}\operatorname{Proj}(\mathcal{M}^{(d,k)}).
\tag{8.14.13.1}
\end{equation}

Supposons d'abord $Y=\operatorname{Spec}(A)$ affine,
$\mathcal{S}=\widetilde{S}$, $\mathcal{M}=\widetilde{M}$, où $S$ est une
$A$-algèbre graduée à degrés positifs et $M$ un $S$-module gradué. On sait
que pour tout $f\in S_e$ ($e>0$) $h$ applique $D_+(f)$ sur $D_+(f^d)$ et
correspond à l'isomorphisme canonique
$S_{(f^d)}\xrightarrow{\sim}S_{(f)}$
(\hyperref[II.2.2.2-fr]{2.2.2}). La restriction de
(\hyperref[II.8.14.13.1-fr]{8.14.13.1}) à $D_+(f^d)$ correspond alors au
di-isomorphisme canonique $M_{f^d}\xrightarrow{\sim}M_f$, restreint aux
éléments de $M_{f^d}$ de degré congru à $k$ (mod. $d$). Nous laissons au
lecteur le soin de vérifier que ces isomorphismes sont compatibles avec le
passage de $f$ à un de ses multiples homogènes $fg$, puis qu'on a une
compatibilité analogue avec le passage de $S$ à une $A'$-algèbre graduée
$S'=S\otimes_A A'$, lorsque $A'$ est une $A$-algèbre.

En particulier, on a ainsi un $h$-isomorphisme
\begin{equation}
\label{II.8.14.13.2-fr}
  (\mathcal{S}^{(d)})_{X^{(d)}}
  \xrightarrow{\sim}(\mathcal{S}_X)^{(d)}
\tag{8.14.13.2}
\end{equation}
qui respecte les structures multiplicatives des deux membres, et grâce auquel
(\hyperref[II.8.14.13.1-fr]{8.14.13.1}) devient un $h$-di-isomorphisme d'un
$(\mathcal{S}^{(d)})_{X^{(d)}}$-Module gradué sur un
$(\mathcal{S}_X)^{(d)}$-Module gradué. De même, on a un $h$-isomorphisme
\begin{equation}
\label{II.8.14.13.3-fr}
  \operatorname{Proj}_0(\mathcal{S}^{(d,k)}(n))
  \xrightarrow{\sim}\mathcal{O}_X(nd+k)
\tag{8.14.13.3}
\end{equation}
qui complète le résultat de (\hyperref[II.3.2.9-fr]{3.2.9, (ii)}).

On déduit immédiatement de l'isomorphisme
(\hyperref[II.8.14.13.1-fr]{8.14.13.1}) un isomorphisme de
$\mathcal{S}^{(d)}$-Modules gradués
\begin{equation}
\label{II.8.14.13.4-fr}
  \Gamma_*^{(d)}(\operatorname{Proj}(\mathcal{M}^{(d,k)}))
  \xrightarrow{\sim}
  \Gamma_*((\operatorname{Proj}(\mathcal{M}))^{(d,k)})
\tag{8.14.13.4}
\end{equation}
où $\Gamma_*^{(d)}$ correspond au morphisme structural
$q^{(d)}:X^{(d)}\to Y$~; il est immédiat de vérifier que l'homomorphisme
canonique $\alpha$ (\hyperref[II.8.14.5.4-fr]{8.14.5.4}) et l'homomorphisme
analogue $\alpha^{(d)}$ pour $X^{(d)}$ rendent commutatif le diagramme
\begin{equation}
\label{II.8.14.13.5-fr}
  \xymatrix{
    & \mathcal{M}^{(d,k)}\ar[dl]_{\alpha^{(d)}}\ar[dr]^{\alpha} & \\
    \Gamma_*^{(d)}(\operatorname{Proj}(\mathcal{M}^{(d,k)}))
      \ar[rr]^{\sim}
      && \Gamma_*((\operatorname{Proj}(\mathcal{M}))^{(d,k)})
  }
\tag{8.14.13.5}
\end{equation}
la vérification se faisant en supposant $Y$ affine, et calculant les
restrictions des images par $\alpha^{(d)}$ et $\alpha$ d'un même élément de
$M^{(d,k)}$ aux ouverts $D_+(f^d)$ et $D_+(f)$ (avec les mêmes notations que
ci-dessus).
\end{remark}

\begin{proposition}[8.14.14]
\label{II.8.14.14-fr}
Soient $Y$ un préschéma quasi-compact, $\mathcal{S}$ une
$\mathcal{O}_Y$-Algèbre graduée quasi-cohérente de type fini, $\mathcal{M}$
un $\mathcal{S}$-Module gradué quasi-cohérent vérifiant la condition
\textbf{(TF)}~; soit $X=\operatorname{Proj}(\mathcal{S})$. Alors
$\mathcal{S}_X$ est une $\mathcal{O}_X$-Algèbre graduée périodique
(\hyperref[II.8.14.12-fr]{8.14.12}), et il existe une

\oldpage[II]{203}
période $d$ de $\mathcal{S}_X$ telle que les
$(\operatorname{Proj}(\mathcal{M}))^{(d,k)}$
$(0\leqslant k\leqslant d-1)$ soient des
$(\mathcal{S}_X)^{(d)}$-Modules de type fini.
\end{proposition}

\begin{proof}
En effet, (\hyperref[II.3.1.10-fr]{3.1.10}) prouve qu'il existe $d$ tel que
$\mathcal{S}^{(d)}$ soit engendrée par
$\mathcal{S}_d=(\mathcal{S}^{(d)})_1$, ce dernier étant un
$\mathcal{S}_0$-Module de type fini. Pour démontrer la première assertion, on
peut donc, en vertu de (\hyperref[II.8.14.13.2-fr]{8.14.13.2}), se borner au
cas où $d=1$, et alors la proposition résulte de
(\hyperref[II.3.2.7-fr]{3.2.7}). En outre, compte tenu de
(\hyperref[II.8.14.13.1-fr]{8.14.13.1}), la seconde assertion est conséquence
de (\hyperref[II.2.1.6-fr]{2.1.6, (iii)}) et de
(\hyperref[II.3.4.3-fr]{3.4.3}).
\end{proof}

% FIN SOURCE ega2/ega2-1-fr.tex

%
% DEBUT SOURCE ega2/ega2-bibliography-fr.tex
\clearpage
\oldpage[II]{205}
\phantomsection
\addcontentsline{toc}{section}{Bibliographie (suite)}
\renewcommand{\bibname}{Bibliographie \textit{(suite)}}

\begin{thebibliography}{26}

\bibitem[23]{II-23}
C.~Chevalley,
\emph{Introduction to the theory of algebraic functions of one variable}.
Math. Surveys, n\textsuperscript{o}~6, Amer. Math. Soc., New York, 1951.

\bibitem[24]{II-24}
P.~Jaffard,
\emph{Les systèmes d'idéaux}, Paris (Dunod), 1960.

\bibitem[25]{II-25}
M.~Nagata,
\emph{On the derived normal rings of Noetherian integral domains},
Mem. Coll. Sci. Kyoto, sér.~A, t.~XXIX (1955), p.~293--303.

\bibitem[26]{II-26}
M.~Nagata,
\emph{Existence theorems for non projective complete algebraic varieties},
Ill. J. Math., t.~II (1958), p.~490--498.

\end{thebibliography}

% FIN SOURCE ega2/ega2-bibliography-fr.tex

%
% DEBUT SOURCE ega2/ega2-indexes-and-contents-fr.tex
\clearpage
\oldpage[II]{207}
\markboth{A. GROTHENDIECK --- Chap. II}{A. GROTHENDIECK --- Chap. II}
\phantomsection
\addcontentsline{toc}{section}{Index des notations}
\section*{Index des notations}
\thispagestyle{plain}

\begingroup
\small
\setlength{\parindent}{0pt}
\setlength{\parskip}{2pt}
\newcommand{\egaShProj}{\mathcal{P}\!\operatorname{roj}}

$\mathcal{A}(X),\mathcal{A}(\mathcal{F}),\mathcal{A}(\mathcal{B})$
($X$ $S$-préschéma, $\mathcal{F}$ $\mathcal{O}_X$-Module,
$\mathcal{B}$ $\mathcal{O}_X$-Algèbre) : \hyperref[II.1.1.1-fr]{1.1.1}.

$\mathcal{A}(h)$ ($h$ $S$-morphisme) : \hyperref[II.1.1.2-fr]{1.1.2}.

$\Spec(\mathcal{B})$ ($\mathcal{B}$ $\mathcal{O}_S$-Algèbre) : \hyperref[II.1.3.1-fr]{1.3.1}.

$\widetilde{\mathcal{M}}$ ($\mathcal{M}$ $\mathcal{B}$-Module,
$\mathcal{B}$ $\mathcal{O}_S$-Algèbre) : \hyperref[II.1.4.3-fr]{1.4.3}.

$\mathbf{T}(E),\mathbf{S}(E),\mathbf{S}_A(E)$ ($E$ $A$-module) : \hyperref[II.1.7.1-fr]{1.7.1}.

$\mathbf{S}(\mathcal{E}),\mathbf{S}_{\mathcal{A}}(\mathcal{E})$
($\mathcal{E}$ $\mathcal{A}$-Module) : \hyperref[II.1.7.4-fr]{1.7.4}.

$\mathbf{V}(\mathcal{E})$ ($\mathcal{E}$ $\mathcal{O}_S$-Module quasi-cohérent) : \hyperref[II.1.7.8-fr]{1.7.8}.

$S_n,S_+,M_n,S^{(d)},M^{(d,k)},M^{(d)}$
($S$ anneau gradué, $M$ $S$-module gradué) : \hyperref[II.2.1.1-fr]{2.1.1}.

$M(n),S(n)$ ($S$ anneau gradué, $M$ $S$-module gradué) : \hyperref[II.2.1.1-fr]{2.1.1}.

$\Hom_S(M,N)$ ($S$ anneau gradué, $M,N$ $S$-modules gradués) : \hyperref[II.2.1.2-fr]{2.1.2}.

$\mathfrak N_+,\mathfrak r_+(\mathfrak J)$
($\mathfrak J$ idéal de $S_+$) : \hyperref[II.2.1.10-fr]{2.1.10}.

$S_{(f)},M_{(f)}$ ($f$ élément homogène de $S_+$) : \hyperref[II.2.2.1-fr]{2.2.1}.

$S_{(T)},M_{(T)},S_{(\mathfrak p)},M_{(\mathfrak p)}$
($T$ partie multiplicative de $S_+$ formée d'éléments homogènes,
$\mathfrak p$ idéal premier gradué de $S_+$) : \hyperref[II.2.2.7-fr]{2.2.7}.

$\Proj(S)$ ($S$ anneau gradué) : \hyperref[II.2.3.1-fr]{2.3.1}.

$V_+(E)$ ($E$ partie de $S_+$) : \hyperref[II.2.3.2-fr]{2.3.2}.

$D_+(f)$ ($f$ élément de $S$) : \hyperref[II.2.3.3-fr]{2.3.3}.

$\mathfrak j_+(Y)$ ($Y$ partie de $\Proj(S)$) : \hyperref[II.2.3.10-fr]{2.3.10}.

$\widetilde M$ ($M$ $S$-module gradué) : \hyperref[II.2.5.2-fr]{2.5.2}.

$\widetilde u$ ($u$ homomorphisme de degré zéro) : \hyperref[II.2.5.4-fr]{2.5.4}.

$\mathcal{O}_X(n),\mathcal{F}(n)$
($X=\Proj(S)$, $\mathcal{F}$ $\mathcal{O}_X$-Module) : \hyperref[II.2.5.10-fr]{2.5.10}.

$\lambda$ (homomorphisme canonique) : \hyperref[II.2.5.11-fr]{2.5.11}.

$\mu$ (homomorphisme canonique) : \hyperref[II.2.5.12-fr]{2.5.12}.

$X^{(d)}$ ($X=\Proj(S)$, $d$ entier $>0$) : \hyperref[II.2.5.16-fr]{2.5.16}.

$\Gamma_*(\mathcal{F})$ ($\mathcal{F}$ $\mathcal{O}_X$-Module,
$X=\Proj(S)$) : \hyperref[II.2.6.1-fr]{2.6.1}.

$\alpha_n,\alpha,\alpha_{\mathcal{M}}$ (homomorphismes canoniques) : \hyperref[II.2.6.2-fr]{2.6.2}.

$\beta,\beta_{\mathcal{F}}$ (homomorphismes canoniques) : \hyperref[II.2.6.4-fr]{2.6.4}.

$G(\vphi),{}^a\vphi$ (par abus de langage)
($\vphi$ homomorphisme d'anneaux gradués) : \hyperref[II.2.8.1-fr]{2.8.1}.

$\Proj(\vphi)$ ($\vphi$ homomorphisme d'anneaux gradués) : \hyperref[II.2.8.2-fr]{2.8.2}.

$\mathcal{S}^{(d)},\mathcal{M}^{(d,k)},\mathcal{M}^{(d)},\mathcal{M}(n)$
($\mathcal{S}$ $\mathcal{O}_Y$-Algèbre graduée,
$\mathcal{M}$ $\mathcal{S}$-Module gradué) : \hyperref[II.3.1.1-fr]{3.1.1}.

$\Proj(\mathcal{S})$ ($\mathcal{S}$ $\mathcal{O}_Y$-Algèbre graduée quasi-cohérente) : \hyperref[II.3.1.3-fr]{3.1.3}.

$X_f$ ($X=\Proj(\mathcal{S})$, $f$ section de $\mathcal{S}_d$ au-dessus de $Y$) : \hyperref[II.3.1.4-fr]{3.1.4}.

$\widetilde{\mathcal{M}}$ ($\mathcal{M}$ $\mathcal{S}$-Module gradué) : \hyperref[II.3.2.2-fr]{3.2.2}.

$\mathcal{O}_X(n),\mathcal{F}(n)$
($X=\Proj(\mathcal{S})$, $\mathcal{F}$ $\mathcal{O}_X$-Module) : \hyperref[II.3.2.5-fr]{3.2.5}.

$\lambda,\mu$ (homomorphismes canoniques) : \hyperref[II.3.2.6-fr]{3.2.6}.

$\Gamma_*(\mathcal{F})$ ($\mathcal{F}$ $\mathcal{O}_X$-Module,
$X=\Proj(\mathcal{S})$) : \hyperref[II.3.3.1-fr]{3.3.1}.

$\alpha_n,\alpha,\alpha_{\mathcal{M}}$ (homomorphismes canoniques) : \hyperref[II.3.3.2-fr]{3.3.2}.

\oldpage[II]{208}
\markboth{A. GROTHENDIECK --- Chap. II}{A. GROTHENDIECK --- Chap. II}
$\beta,\beta_{\mathcal{F}}$ (homomorphismes canoniques) : \hyperref[II.3.3.4-fr]{3.3.4}.

$G(\vphi),\Proj(\vphi)$
($\vphi$ homomorphisme de $\mathcal{O}_Y$-Algèbres graduées) : \hyperref[II.3.5.1-fr]{3.5.1}.

$\Proj(u)$ ($u$ $g$-morphisme d'une $\mathcal{O}_Y$-Algèbre graduée
dans une $\mathcal{O}_{Y'}$-Algèbre graduée) : \hyperref[II.3.5.6-fr]{3.5.6}.

$r_{\mathcal{L},\psi}$ ($\mathcal{L}$ $\mathcal{O}_X$-Module inversible) : \hyperref[II.3.7.1-fr]{3.7.1}.

$\mathbf{P}(\mathcal{E}),\mathbf{P}(E),\mathbf{P}_Y^n,\mathbf{P}_A^n$ : \hyperref[II.4.1.1-fr]{4.1.1}.

$\mathbf{P}(u)$ ($u$ homomorphisme surjectif de $\mathcal{O}_Y$-Modules) : \hyperref[II.4.1.2-fr]{4.1.2}.

$\Hyp_Y(X,\mathcal{E})$ ($\mathcal{E}$ $\mathcal{O}_Y$-Module quasi-cohérent) : \hyperref[II.4.2.5-fr]{4.2.5}.

$\varsigma$ (morphisme de Segre) : \hyperref[II.4.3.1-fr]{4.3.1}.

$\mathcal{F}(n)$ ($\mathcal{F}$ $\mathcal{O}_X$-Module, $X$ muni d'un
$\mathcal{O}_X$-Module inversible) : \hyperref[II.4.5.1-fr]{4.5.1}.

$\varepsilon$ (automorphisme identique) : \hyperref[II.4.5.1-fr]{4.5.1}.

$P(u,T),\sigma_i(u)$ ($u$ endomorphisme de module libre de base finie) : \hyperref[II.6.4.1-fr]{6.4.1}.

$P(u,T),\sigma_i(u),\det u$
($u$ endomorphisme de module de type fini sur un anneau intègre ou un
anneau noethérien réduit) : \hyperref[II.6.4.2-fr]{6.4.2} et \hyperref[II.6.4.7-fr]{6.4.7}.

$\sigma_i(u),\det u$ ($u$ endomorphisme de $\mathcal{A}$-Module localement libre) : \hyperref[II.6.4.8-fr]{6.4.8}.

$\sigma_i,\det$ (homomorphismes de faisceaux) : \hyperref[II.6.4.8-fr]{6.4.8}, \hyperref[II.6.4.9-fr]{6.4.9} et \hyperref[II.6.4.10-fr]{6.4.10}.

$\sigma_i(u),\det u$
($u$ endomorphisme de $\mathcal{O}_X$-Module de type fini sur un préschéma
localement intègre ou un préschéma localement noethérien réduit) :
\hyperref[II.6.4.9-fr]{6.4.9} et \hyperref[II.6.4.10-fr]{6.4.10}.

$N_{\mathcal{B}/\mathcal{A}}(f)$
($f$ section de la $\mathcal{B}$-Algèbre $\mathcal{A}$) : \hyperref[II.6.5.1-fr]{6.5.1}.

$N_{\mathcal{B}/\mathcal{A}}(\mathcal{L}')$
($\mathcal{L}'$ $\mathcal{B}$-Module inversible) : \hyperref[II.6.5.2-fr]{6.5.2}.

$N_{\mathcal{B}/\mathcal{A}}(h')$
($h'$ homomorphisme de $\mathcal{B}$-modules inversibles) : \hyperref[II.6.5.3-fr]{6.5.3}.

$N_{X'/X}(\mathcal{L}')$
($X'$ fini au-dessus de $X$, $\mathcal{L}'$ $\mathcal{O}_{X'}$-Module inversible) : \hyperref[II.6.5.5-fr]{6.5.5}.

$N_{X'/X}(h')$
($h'$ homomorphisme de $\mathcal{O}_{X'}$-Modules inversibles) : \hyperref[II.6.5.5-fr]{6.5.5}.

$\mathcal{S}_X$ ($X=\Proj(\mathcal{S})$, $\mathcal{S}$ somme d'Idéaux de
$\mathcal{O}_Y$) : \hyperref[II.8.1.5-fr]{8.1.5}.

$S^\geq,S^\leq,S_f^\geq,S_f^\leq,
M^\geq,M^\leq,M_f^\geq,M_f^\leq$
($S$ anneau gradué, $M$ $S$-module gradué,
$f$ élément homogène de $S_+$) : \hyperref[II.8.2.1-fr]{8.2.1}.

$\widehat S$ ($S$ anneau gradué) : \hyperref[II.8.2.2-fr]{8.2.2}.

$\widehat M$ ($M$ $S$-module gradué) : \hyperref[II.8.2.4-fr]{8.2.4}.

$S_{[n]},S^\natural,f^\natural$
($S$ anneau gradué, $f$ homogène dans $S$) : \hyperref[II.8.2.6-fr]{8.2.6}.

$M_{[n]},M^\natural$ ($M$ $S$-module gradué) : \hyperref[II.8.2.8-fr]{8.2.8}.

$\widehat{\mathcal{S}}$ ($\mathcal{S}$ $\mathcal{O}_Y$-Algèbre graduée) : \hyperref[II.8.3.1-fr]{3.8.1}.

$G_m$ (schéma en groupes) : \hyperref[II.8.3.9-fr]{8.3.9}.

$\mathcal{S}_X$ ($X=\Proj(\mathcal{S})$, $\mathcal{S}$ $\mathcal{O}_Y$-Algèbre graduée) : \hyperref[II.8.6.1-fr]{8.6.1}.

$\mathcal{S}_X,C_X,\widehat C_X,E_X,\widehat E_X$
($X=\Proj(\mathcal{S})$, $\mathcal{S}$ $\mathcal{O}_Y$-Algèbre graduée
avec $\mathcal{S}_0=\mathcal{O}_Y$) : \hyperref[II.8.6.1-fr]{8.6.1}.

$\mathcal{S}_{[n]},\mathcal{S}^{\natural}$
($\mathcal{S}$ $\mathcal{O}_Y$-Algèbre graduée) : \hyperref[II.8.7.2-fr]{8.7.2}.

$\egaShProj_0(\mathcal{M}),\egaShProj(\mathcal{M}),\mathcal{M}_X,
\widehat{\mathcal{M}},\mathcal{M}^{\square}$
($\mathcal{M}$ $\mathcal{S}$-Module gradué) : \hyperref[II.8.12.1-fr]{8.12.1}.

$\mathcal{M}_X^\geq$
($X=\Proj(\mathcal{S})$, $\mathcal{M}$ $\mathcal{S}$-Module gradué) : \hyperref[II.8.12.5-fr]{8.12.5}.

$\mathcal{M}_{[n]},\mathcal{M}^{\natural}$
($\mathcal{M}$ $\mathcal{S}$-Module gradué) : \hyperref[II.8.12.9-fr]{8.12.9}.

$(\widetilde{\mathcal{N}})^-$
($\mathcal{N}$ sous-$\mathcal{S}$-Module d'un $\mathcal{S}$-Module gradué) : \hyperref[II.8.13.1-fr]{8.13.1}.

$\overline{\mathcal{N}}$
($\mathcal{N}$ sous-$\mathcal{S}$-Module d'un $\mathcal{S}$-Module gradué) : \hyperref[II.8.13.2-fr]{8.13.2}.

$\Gamma_*(\mathcal{F})$
($X=\Proj(\mathcal{S})$, $\mathcal{F}$ $\mathcal{S}_X$-Module gradué) : \hyperref[II.8.14.5-fr]{8.14.5}.

$\mathcal{S}_{X'},\Gamma'_*(\mathcal{F}'),\egaShProj'(\mathcal{M}),
\alpha',\beta'$
($X'$ ouvert de $X=\Proj(\mathcal{S})$, $\mathcal{F}'$ $\mathcal{S}_{X'}$-Module gradué,
$\mathcal{M}$ $\mathcal{S}$-Module gradué) : \hyperref[II.8.14.6-fr]{8.14.6}.

\endgroup

\clearpage
\oldpage[II]{209}
\markboth{A. GROTHENDIECK --- Chap. II}{A. GROTHENDIECK --- Chap. II}
\thispagestyle{plain}
\phantomsection
\addcontentsline{toc}{section}{Index terminologique}
\section*{Index terminologique}

\begingroup
\small
\setlength{\parindent}{0pt}
\setlength{\parskip}{2pt}
\newcommand{\egaTerme}[2]{\textbf{#1} : #2.\par}

\egaTerme{Affine (morphisme)}{\hyperref[II.1.6.1-fr]{1.6.1}}
\egaTerme{Affine (préschéma) au-dessus d'un préschéma}{\hyperref[II.1.2.1-fr]{1.2.1}}
\egaTerme{Algèbre graduée sur un anneau gradué}{\hyperref[II.2.1.2-fr]{2.1.2}}
\egaTerme{Algèbre symétrique d'un $A$-module}{\hyperref[II.1.7.1-fr]{1.7.1}}
\egaTerme{$\mathcal{A}$-Algèbre symétrique d'un $\mathcal{A}$-Module}{\hyperref[II.1.7.4-fr]{1.7.4}}
\egaTerme{$\mathcal{O}_Y$-Algèbre graduée essentiellement réduite}{\hyperref[II.3.1.12-fr]{3.1.12}}
\egaTerme{$\mathcal{O}_Y$-Algèbre graduée intègre}{\hyperref[II.3.1.12-fr]{3.1.12}}
\egaTerme{Ample ($\mathcal{O}_X$-Module inversible)}{\hyperref[II.4.5.3-fr]{4.5.3}}
\egaTerme{Ample relativement à $f$, $f$-ample, $Y$-ample ($\mathcal{O}_X$-Module inversible)}{\hyperref[II.4.6.1-fr]{4.6.1}}
\egaTerme{Anneau gradué essentiellement intègre}{\hyperref[II.2.1.11-fr]{2.1.11}}
\egaTerme{Anneau gradué essentiellement réduit}{\hyperref[II.2.1.10-fr]{2.1.10}}
\egaTerme{Associé à un $\mathcal{B}$-Module (faisceau)}{\hyperref[II.1.4.3-fr]{1.4.3}}
\egaTerme{Associé à un $S$-module gradué (faisceau)}{\hyperref[II.2.5.3-fr]{2.5.3}}
\egaTerme{Associé à un $\mathcal{S}$-Module gradué (faisceau)}{\hyperref[II.3.2.2-fr]{3.2.2}}
\egaTerme{Associé à une $\mathcal{O}_S$-Algèbre ($S$-schéma)}{\hyperref[II.1.3.1-fr]{1.3.1}}

\egaTerme{Condition (TF), condition (TN)}{\hyperref[II.2.7.2-fr]{2.7.2}}
\egaTerme{Condition (TF), condition (TN)}{\hyperref[II.3.4.2-fr]{3.4.2}}
\egaTerme{Cône affine, cône projectif défini par une $\mathcal{O}_Y$-Algèbre graduée}{\hyperref[II.8.3.1-fr]{8.3.1}}
\egaTerme{Cône affine épointé, cône projectif épointé défini par une $\mathcal{O}_Y$-Algèbre graduée}{\hyperref[II.8.3.4-fr]{8.3.4}}
\egaTerme{Cône projetant affine, cône projetant projectif d'un préschéma $\Proj(\mathcal{S})$}{\hyperref[II.8.3.1-fr]{8.3.1}}
\egaTerme{Courbe algébrique sur un corps}{\hyperref[II.7.4.2-fr]{7.4.2}}
\egaTerme{Courbe algébrique complète}{\hyperref[II.7.4.7-fr]{7.4.7}}
\egaTerme{Critère de Serre}{\hyperref[II.5.2.1-fr]{5.2.1}}

\egaTerme{Déterminant d'un endomorphisme d'un module de type fini sur un anneau intègre}{\hyperref[II.6.4.2-fr]{6.4.2}}
\egaTerme{Déterminant d'un endomorphisme d'un module de type fini sur un anneau noethérien réduit}{\hyperref[II.6.4.7-fr]{6.4.7}}
\egaTerme{Déterminant d'un endomorphisme d'un $\mathcal{A}$-Module localement libre}{\hyperref[II.6.4.8-fr]{6.4.8}}
\egaTerme{Déterminant d'un endomorphisme d'un $\mathcal{O}_X$-Module de type fini sur un préschéma $X$ localement intègre}{\hyperref[II.6.4.9-fr]{6.4.9}}
\egaTerme{Déterminant d'un endomorphisme d'un $\mathcal{O}_X$-Module de type fini sur un préschéma $X$ localement noethérien et réduit}{\hyperref[II.6.4.10-fr]{6.4.10}}

\egaTerme{Éclaté (préschéma)}{\hyperref[II.8.1.3-fr]{8.1.3}}
\egaTerme{Entier (morphisme)}{\hyperref[II.6.1.1-fr]{6.1.1}}
\egaTerme{Entier (préschéma) au-dessus d'un préschéma}{\hyperref[II.6.1.1-fr]{6.1.1}}
\egaTerme{Entière ($\mathcal{O}_S$-Algèbre quasi-cohérente)}{\hyperref[II.6.1.2-fr]{6.1.2}}
\egaTerme{Entière (section d'une $\mathcal{A}$-Algèbre) sur $\mathcal{A}$}{\hyperref[II.6.3.1-fr]{6.3.1}}

\egaTerme{Faisceau fondamental d'un fibré projectif}{\hyperref[II.4.1.1-fr]{4.1.1}}
\egaTerme{Fermeture intégrale d'un faisceau d'anneaux $\mathcal{A}$ dans une $\mathcal{A}$-Algèbre}{\hyperref[II.6.3.2-fr]{6.3.2}}
\egaTerme{Fermeture intégrale d'un préschéma $X$ relativement à une $\mathcal{O}_X$-Algèbre}{\hyperref[II.6.3.4-fr]{6.3.4}}
\egaTerme{Fermeture projective d'un cône affine}{\hyperref[II.8.3.1-fr]{8.3.1}}

\oldpage[II]{210}
\markboth{A. GROTHENDIECK --- Chap. II}{A. GROTHENDIECK --- Chap. II}
\egaTerme{Fibre géométrique rationnelle d'un fibré projectif sur une extension $K$ de $k(y)$}{\hyperref[II.4.2.6-fr]{4.2.6}}
\egaTerme{Fibré projectif défini par un $\mathcal{O}_Y$-Module}{\hyperref[II.4.1.1-fr]{4.1.1}}
\egaTerme{Fibré vectoriel défini par un $\mathcal{O}_S$-Module}{\hyperref[II.1.7.8-fr]{1.7.8}}
\egaTerme{Fini (morphisme)}{\hyperref[II.6.1.1-fr]{6.1.1}}
\egaTerme{Fini (préschéma) au-dessus d'un préschéma}{\hyperref[II.6.1.1-fr]{6.1.1}}
\egaTerme{Finie ($\mathcal{O}_S$-Algèbre quasi-cohérente)}{\hyperref[II.6.1.2-fr]{6.1.2}}
\egaTerme{Fonctions symétriques élémentaires d'un endomorphisme d'un module localement libre}{\hyperref[II.6.4.1-fr]{6.4.1}}
\egaTerme{Fonctions symétriques élémentaires d'un endomorphisme d'un module de type fini sur un anneau intègre}{\hyperref[II.6.4.2-fr]{6.4.2}}
\egaTerme{Fonctions symétriques élémentaires d'un endomorphisme d'un module de type fini sur un anneau noethérien réduit}{\hyperref[II.6.4.7-fr]{6.4.7}}
\egaTerme{Fonctions symétriques élémentaires d'un endomorphisme d'un $\mathcal{A}$-Module localement libre de rang constant}{\hyperref[II.6.4.8-fr]{6.4.8}}
\egaTerme{Fonctions symétriques élémentaires d'un endomorphisme d'un $\mathcal{O}_X$-Module de type fini sur un préschéma $X$ localement intègre}{\hyperref[II.6.4.9-fr]{6.4.9}}
\egaTerme{Fonctions symétriques élémentaires d'un endomorphisme d'un $\mathcal{O}_X$-Module de type fini sur un préschéma $X$ localement noethérien réduit}{\hyperref[II.6.4.10-fr]{6.4.10}}
\egaTerme{Fonctions symétriques élémentaires d'une section d'une $\mathcal{A}$-Algèbre}{\hyperref[II.6.5.1-fr]{6.5.1}}

\egaTerme{Homogénisé d'un sous-$\mathcal{S}$-Module d'un $\mathcal{S}$-Module gradué}{\hyperref[II.8.13.2-fr]{8.13.2}}
\egaTerme{Homomorphisme d'anneaux gradués}{\hyperref[II.2.1.2-fr]{2.1.2}}
\egaTerme{Homomorphisme de degré $k$ de modules gradués}{\hyperref[II.2.1.2-fr]{2.1.2}}

\egaTerme{Idéal de $S_+$, idéal premier gradué de $S_+$ ($S$ anneau gradué)}{\hyperref[II.2.1.10-fr]{2.1.10}}

\egaTerme{Lemme de Chow}{\hyperref[II.5.6.1-fr]{5.6.1}}
\egaTerme{Lieu à l'infini d'un cône projectif}{\hyperref[II.8.3.3-fr]{8.3.3}}

\egaTerme{Morphisme canonique $G(\varepsilon)\to\Proj(S)$, $S=\bigoplus_{n\geq0}\Gamma(X,\mathcal{L}^{\otimes n})$}{\hyperref[II.4.5.1-fr]{4.5.1}}
\egaTerme{Morphisme canonique $X\to\Spec(\Gamma(X,\mathcal{O}_X))$}{\hyperref[II.5.1.1-fr]{5.1.1}}
\egaTerme{Morphisme de Segre}{\hyperref[II.4.3.1-fr]{4.3.1}}

\egaTerme{Nilradical de $S_+$ ($S$ anneau gradué)}{\hyperref[II.2.1.10-fr]{2.1.10}}
\egaTerme{Nilradical de $\mathcal{S}_+$ ($\mathcal{S}$ $\mathcal{O}_Y$-Algèbre graduée quasi-cohérente)}{\hyperref[II.3.1.12-fr]{3.1.12}}
\egaTerme{Norme d'une section d'une $\mathcal{A}$-Algèbre}{\hyperref[II.6.5.1-fr]{6.5.1}}
\egaTerme{Norme d'un $\mathcal{B}$-Module inversible}{\hyperref[II.6.5.2-fr]{6.5.2}}
\egaTerme{Norme d'une section d'un $\mathcal{B}$-Module inversible}{\hyperref[II.6.5.3-fr]{6.5.3}}
\egaTerme{Norme d'un $\mathcal{O}_{X'}$-Module inversible}{\hyperref[II.6.5.5-fr]{6.5.5}}

\egaTerme{Périodique (Algèbre graduée)}{\hyperref[II.8.14.12-fr]{8.14.12}}
\egaTerme{Polynôme caractéristique d'un endomorphisme d'un module de type fini sur un anneau intègre}{\hyperref[II.6.4.2-fr]{6.4.2}}
\egaTerme{Polynôme caractéristique d'un endomorphisme d'un module de type fini sur un anneau noethérien réduit}{\hyperref[II.6.4.7-fr]{6.4.7}}
\egaTerme{Préschéma des sommets d'un cône affine (projectif)}{\hyperref[II.8.3.3-fr]{8.3.3}}
\egaTerme{Présentation finie (module gradué ayant une)}{\hyperref[II.2.1.1-fr]{2.1.1}}
\egaTerme{Présentation finie ($\mathcal{S}$-Module gradué ayant une)}{\hyperref[II.3.1.1-fr]{3.1.1}}
\egaTerme{Produit tensoriel gradué de deux modules gradués}{\hyperref[II.2.1.2-fr]{2.1.2}}
\egaTerme{Projectif (morphisme)}{\hyperref[II.5.5.2-fr]{5.5.2}}
\egaTerme{Projectif (préschéma) au-dessus d'un préschéma}{\hyperref[II.5.5.2-fr]{5.5.2}}
\egaTerme{Propre (morphisme)}{\hyperref[II.5.4.1-fr]{5.4.1}}
\egaTerme{Propre (partie)}{\hyperref[II.5.4.10-fr]{5.4.10}}
\egaTerme{Propre (préschéma) au-dessus d'un préschéma}{\hyperref[II.5.4.1-fr]{5.4.1}}

\egaTerme{Quasi-affine (morphisme)}{\hyperref[II.5.1.1-fr]{5.1.1}}
\egaTerme{Quasi-affine (préschéma) au-dessus d'un préschéma}{\hyperref[II.5.1.1-fr]{5.1.1}}
\egaTerme{Quasi-affine (schéma)}{\hyperref[II.5.1.1-fr]{5.1.1}}

\oldpage[II]{211}
\markboth{CHAPITRE II}{CHAPITRE II}
\egaTerme{Quasi-fini (morphisme)}{\hyperref[II.6.2.3-fr]{6.2.3}}
\egaTerme{Quasi-fini (préschéma) au-dessus d'un préschéma}{\hyperref[II.6.2.3-fr]{6.2.3}}
\egaTerme{Quasi-projectif (morphisme)}{\hyperref[II.5.3.1-fr]{5.3.1}}
\egaTerme{Quasi-projectif (préschéma) au-dessus d'un préschéma}{\hyperref[II.5.3.1-fr]{5.3.1}}

\egaTerme{Racine d'un idéal de $S_+$ ($S$ anneau gradué)}{\hyperref[II.2.1.10-fr]{2.1.10}}
\egaTerme{Rétraction canonique d'un cône projectif épointé sur son lieu à l'infini}{\hyperref[II.8.3.5-fr]{8.3.5}}

\egaTerme{Section canonique de $\mathcal{O}_X(-1)$ ($X$ préschéma éclaté)}{\hyperref[II.8.1.9-fr]{8.1.9}}
\egaTerme{Section nulle d'un cône affine}{\hyperref[II.8.3.3-fr]{8.3.3}}
\egaTerme{Section nulle d'un fibré vectoriel}{\hyperref[II.1.7.9-fr]{1.7.9}}
\egaTerme{Section sommet d'un cône affine (projectif)}{\hyperref[II.8.3.3-fr]{8.3.3}}
\egaTerme{Spectre d'une $\mathcal{O}_S$-Algèbre}{\hyperref[II.1.3.1-fr]{1.3.1}}
\egaTerme{Spectre premier homogène d'un anneau gradué}{\hyperref[II.2.3.1-fr]{2.3.1}}
\egaTerme{Spectre homogène d'une $\mathcal{O}_S$-Algèbre graduée quasi-cohérente}{\hyperref[II.3.1.3-fr]{3.1.3}}

\egaTerme{(TN)-injectif, (TN)-surjectif, (TN)-bijectif (homomorphisme de modules gradués)}{\hyperref[II.2.7.2-fr]{2.7.2}}
\egaTerme{(TN)-isomorphisme de modules gradués}{\hyperref[II.2.7.2-fr]{2.7.2}}
\egaTerme{(TN)-injectif, (TN)-surjectif, (TN)-bijectif (homomorphisme d'anneaux gradués)}{\hyperref[II.2.9.1-fr]{2.9.1}}
\egaTerme{(TN)-isomorphisme d'anneaux gradués}{\hyperref[II.2.9.1-fr]{2.9.1}}
\egaTerme{(TN)-injectif, (TN)-surjectif, (TN)-bijectif (homomorphisme de $\mathcal{S}$-Modules gradués)}{\hyperref[II.3.4.2-fr]{3.4.2}}
\egaTerme{(TN)-isomorphisme de $\mathcal{S}$-Modules gradués}{\hyperref[II.3.4.2-fr]{3.4.2}}
\egaTerme{(TN)-injectif, (TN)-surjectif, (TN)-bijectif (homomorphisme de $\mathcal{O}_Y$-Algèbres graduées)}{\hyperref[II.3.6.1-fr]{3.6.1}}
\egaTerme{(TN)-isomorphisme de $\mathcal{O}_Y$-Algèbres graduées}{\hyperref[II.3.6.1-fr]{3.6.1}}
\egaTerme{Topologie spectrale sur $\Proj(S)$}{\hyperref[II.2.3.3-fr]{2.3.3}}
\egaTerme{Très ample pour $q$, très ample pour $Y$ ($\mathcal{O}_X$-Module inversible)}{\hyperref[II.4.4.2-fr]{4.4.2}}
\egaTerme{Type fini ($\mathcal{S}$-Module gradué de)}{\hyperref[II.3.1.1-fr]{3.1.1}}

\egaTerme{Universellement fermé (morphisme)}{\hyperref[II.5.4.9-fr]{5.4.9}}

\endgroup

\clearpage
\oldpage[II]{212}
\thispagestyle{empty}
\mbox{}

\clearpage
\oldpage[II]{213}
\markboth{TABLE DES MATIÈRES}{TABLE DES MATIÈRES}
\phantomsection
\addcontentsline{toc}{section}{Table des matières originale}
\section*{Table des matières}
\thispagestyle{plain}

\begingroup
\small
\renewcommand{\arraystretch}{1.04}
\begin{longtable}{@{}p{0.09\textwidth}p{0.69\textwidth}r@{}}
\textbf{Unité} & \textbf{Titre} & \textbf{Page originale}\\
\hline
\textbf{II} & \textbf{Étude globale élémentaire de quelques classes de morphismes} & \textbf{5}\\
\hyperref[section:II.1-fr]{\textsection1} & Morphismes affines & 5\\
\hyperref[subsection:II.1.1-fr]{1.1} & $S$-préschémas et $\mathcal{O}_S$-Algèbres & 5\\
\hyperref[subsection:II.1.2-fr]{1.2} & Préschémas affines sur un préschéma & 6\\
\hyperref[subsection:II.1.3-fr]{1.3} & Préschéma affine au-dessus de $S$ associé à une $\mathcal{O}_S$-Algèbre & 8\\
\hyperref[subsection:II.1.4-fr]{1.4} & Faisceaux quasi-cohérents sur un préschéma affine au-dessus de $S$ & 9\\
\hyperref[II.1.5-fr]{1.5} & Changement du préschéma de base & 12\\
\hyperref[II.1.6-fr]{1.6} & Morphismes affines & 14\\
\hyperref[II.1.7-fr]{1.7} & Fibré vectoriel associé à un faisceau de modules & 14\\
\hyperref[section:II.2-fr]{\textsection2} & Spectres premiers homogènes & 19\\
\hyperref[subsection:II.2.1-fr]{2.1} & Généralités sur les anneaux et modules gradués & 19\\
\hyperref[subsection:II.2.2-fr]{2.2} & Anneaux de fractions d'un anneau gradué & 23\\
\hyperref[subsection:II.2.3-fr]{2.3} & Spectre premier homogène d'un anneau gradué & 25\\
\hyperref[subsection:II.2.4-fr]{2.4} & La structure de schéma sur $\Proj(S)$ & 28\\
\hyperref[subsection:II.2.5-fr]{2.5} & Faisceau associé à un module gradué & 30\\
\hyperref[subsection:II.2.6-fr]{2.6} & $S$-module gradué associé à un faisceau sur $\Proj(S)$ & 36\\
\hyperref[subsection:II.2.7-fr]{2.7} & Conditions de finitude & 38\\
\hyperref[subsection:II.2.8-fr]{2.8} & Comportements fonctoriels & 41\\
\hyperref[subsection:II.2.9-fr]{2.9} & Sous-préschémas fermés d'un schéma $\Proj(S)$ & 48\\
\hyperref[section:II.3-fr]{\textsection3} & Spectre homogène d'un faisceau d'algèbres graduées & 49\\
\hyperref[subsection:II.3.1-fr]{3.1} & Spectre homogène d'une $\mathcal{O}_Y$-Algèbre graduée quasi-cohérente & 49\\
\hyperref[subsection:II.3.2-fr]{3.2} & Faisceau sur $\Proj(\mathcal{S})$ associé à un $\mathcal{S}$-Module gradué & 54\\
\hyperref[subsection:II.3.3-fr]{3.3} & $\mathcal{S}$-Module gradué associé à un faisceau sur $\Proj(\mathcal{S})$ & 56\\
\hyperref[subsection:II.3.4-fr]{3.4} & Conditions de finitude & 59\\
\hyperref[subsection:II.3.5-fr]{3.5} & Comportements fonctoriels & 61\\
\hyperref[subsection:II.3.6-fr]{3.6} & Sous-préschémas fermés d'un préschéma $\Proj(\mathcal{S})$ & 64\\
\hyperref[subsection:II.3.7-fr]{3.7} & Morphismes d'un préschéma dans un spectre homogène & 65\\
\hyperref[subsection:II.3.8-fr]{3.8} & Critères d'immersion dans un spectre homogène & 69\\
\end{longtable}
\oldpage[II]{214}
\markboth{A. GROTHENDIECK --- Chap. II}{A. GROTHENDIECK --- Chap. II}
\begin{longtable}{@{}p{0.09\textwidth}p{0.69\textwidth}r@{}}
\hyperref[section:II.4-fr]{\textsection4} & Fibrés projectifs. Faisceaux amples & 71\\
\hyperref[subsection:II.4.1-fr]{4.1} & Définition des fibrés projectifs & 71\\
\hyperref[subsection:II.4.2-fr]{4.2} & Morphismes d'un préschéma dans un fibré projectif & 72\\
\hyperref[subsection:II.4.3-fr]{4.3} & Le morphisme de Segre & 76\\
\hyperref[subsection:II.4.4-fr]{4.4} & Immersions dans les fibrés projectifs. Faisceaux très amples & 78\\
\hyperref[subsection:II.4.5-fr]{4.5} & Faisceaux amples & 83\\
\hyperref[II.4.6-fr]{4.6} & Faisceaux relativement amples & 89\\
\hyperref[section:II.5-fr]{\textsection5} & Morphismes quasi-affines; morphismes quasi-projectifs; morphismes propres; morphismes projectifs & 94\\
\hyperref[subsection:II.5.1-fr]{5.1} & Morphismes quasi-affines & 94\\
\hyperref[subsection:II.5.2-fr]{5.2} & Le critère de Serre & 97\\
\hyperref[subsection:II.5.3-fr]{5.3} & Morphismes quasi-projectifs & 99\\
\hyperref[subsection:II.5.4-fr]{5.4} & Morphismes propres et morphismes universellement fermés & 100\\
\hyperref[subsection:II.5.5-fr]{5.5} & Morphismes projectifs & 103\\
\hyperref[subsection:II.5.6-fr]{5.6} & Le lemme de Chow & 106\\
\hyperref[section:II.6-fr]{\textsection6} & Morphismes entiers et morphismes finis & 110\\
\hyperref[subsection:II.6.1-fr]{6.1} & Préschémas entiers sur un autre & 110\\
\hyperref[subsection:II.6.2-fr]{6.2} & Morphismes quasi-finis & 114\\
\hyperref[subsection:II.6.3-fr]{6.3} & Fermeture intégrale d'un préschéma & 116\\
\hyperref[subsection:II.6.4-fr]{6.4} & Déterminant d'un endomorphisme de $\mathcal{O}_X$-Module & 120\\
\hyperref[subsection:II.6.5-fr]{6.5} & Norme d'un faisceau inversible & 125\\
\hyperref[subsection:II.6.6-fr]{6.6} & Application : critères d'amplitude & 130\\
\hyperref[subsection:II.6.7-fr]{6.7} & Le théorème de Chevalley & 135\\
\hyperref[section:II.7-fr]{\textsection7} & Critères valuatifs & 138\\
\hyperref[subsection:II.7.1-fr]{7.1} & Rappels sur les anneaux de valuation & 138\\
\hyperref[subsection:II.7.2-fr]{7.2} & Critère valuatif de séparation & 141\\
\hyperref[subsection:II.7.3-fr]{7.3} & Critère valuatif de propreté & 143\\
\hyperref[subsection:II.7.4-fr]{7.4} & Courbes algébriques et corps de fonctions de dimension 1 & 148\\
\hyperref[section:II.8-fr]{\textsection8} & Schémas éclatés; cônes projetants; fermeture projective & 152\\
\hyperref[subsection:II.8.1-fr]{8.1} & Préschémas éclatés & 152\\
\hyperref[subsection:II.8.2-fr]{8.2} & Résultats préliminaires sur la localisation dans les anneaux gradués & 157\\
\hyperref[subsection:II.8.3-fr]{8.3} & Cônes projetants & 162\\
\hyperref[subsection:II.8.4-fr]{8.4} & Fermeture projective d'un fibré vectoriel & 168\\
\hyperref[subsection:II.8.5-fr]{8.5} & Comportements fonctoriels & 169\\
\hyperref[subsection:II.8.6-fr]{8.6} & Un isomorphisme canonique pour les cônes épointés & 171\\
\end{longtable}
\oldpage[II]{215}
\markboth{TABLE DES MATIÈRES}{TABLE DES MATIÈRES}
\begin{longtable}{@{}p{0.09\textwidth}p{0.69\textwidth}r@{}}
\hyperref[subsection:II.8.7-fr]{8.7} & Éclatement des cônes projetants & 173\\
\hyperref[subsection:II.8.8-fr]{8.8} & Faisceaux amples et contractions & 177\\
\hyperref[subsection:II.8.9-fr]{8.9} & Le critère d'amplitude de Grauert : énoncé & 182\\
\hyperref[subsection:II.8.10-fr]{8.10} & Le critère d'amplitude de Grauert : démonstration & 184\\
\hyperref[subsection:II.8.11-fr]{8.11} & Unicité des contractions & 189\\
\hyperref[subsection:II.8.12-fr]{8.12} & Faisceaux quasi-cohérents sur les cônes projetants & 191\\
\hyperref[subsection:II.8.13-fr]{8.13} & Fermeture projective de sous-faisceaux et de sous-schémas fermés & 195\\
\hyperref[subsection:II.8.14-fr]{8.14} & Compléments sur les faisceaux associés aux $\mathcal{S}$-Modules gradués & 197\\
\multicolumn{2}{@{}l}{Bibliographie (suite)} & 205\\
\multicolumn{2}{@{}l}{Index des notations} & 207\\
\multicolumn{2}{@{}l}{Index terminologique} & 209\\
\multicolumn{2}{@{}l}{Errata et Addenda (Liste 1)} & 217\\
\end{longtable}

\begin{flushright}
\emph{Reçu le 22 juillet 1960.}
\end{flushright}
\endgroup

% FIN SOURCE ega2/ega2-indexes-and-contents-fr.tex

%
% DEBUT SOURCE ega2/ega2-errata-addenda-fr.tex
\clearpage
\oldpage[II]{217}
\markboth{ERRATA ET ADDENDA, LISTE 1}{ERRATA ET ADDENDA, LISTE 1}
\phantomsection
\addcontentsline{toc}{section}{Errata et Addenda (Liste 1)}
\section*{Errata et Addenda}
\thispagestyle{plain}
\begin{center}
\large (Liste 1)
\end{center}

\subsection*{Chapitre 0}

\paragraph{(0, 2.1.3).}
Ligne 15 du bas de la p.~21, remplacer «~inter-section~» par
«~intersection~».

\paragraph{(0, 2.1.6).}
Ligne 13 de la p.~22, remplacer
\[
  Z_i\cap C\left(\bigcup_{j\ne i}Z_j\right)
  \quad\hbox{par}\quad
  \bigcup_{j\ne i}Z_j.
\]

\paragraph{(0, 3.1.6).}
Ligne 17 du bas de la p.~25, remplacer
$\Gamma(U,\sh{F})$ par $\Gamma(V,\sh{F})$.

\paragraph{(0, 4.1.1).}
Ligne 12 de la p.~36, remplacer $\sh{G}\to\sh{F}$ par
$\sh{B}\to\sh{A}$.

\paragraph{(0, 4.2.1).}
Ligne 11 du bas de la p.~39, remplacer «~faisceaux~» par «~faisceau~».

\paragraph{(0, 5.1.2).}
Après la ligne 14 de la p.~45, ajouter :

Si $\sh{G}$ est un $\sh{O}_Y$-Module engendré par ses sections au-dessus
de $Y$, alors $f^*(\sh{G})$ est engendré par ses sections au-dessus de
$X$, car $f^*$ est un foncteur exact à droite.

\paragraph{(0, 5.2.4).}
Ligne 18 du bas de la p.~46, remplacer $f_U$ par $f_U^*$.

\paragraph{(0, 5.3.4).}
Avant la ligne 16 du bas de la p.~47, ajouter :

Si
\[
  \sh{A}\to\sh{B}\to\sh{C}\to\sh{D}\to\sh{E}
\]
est une suite exacte de $\sh{O}_X$-Modules, et si
$\sh{A},\sh{B},\sh{D},\sh{E}$ sont cohérents, alors $\sh{C}$ est
cohérent.

\paragraph{(0, 5.4.1).}
Après la ligne 2 de la p.~49, ajouter :

Supposons $\sh{O}_X$ cohérent, et soit $\sh{F}$ un $\sh{O}_X$-Module
cohérent. Alors, si en un point $x\in X$, $\sh{F}_x$ est un
$\sh{O}_x$-module libre de rang $n$, il existe un voisinage $U$ de $x$
tel que $\sh{F}|U$ soit localement libre de rang $n$; en effet
$\sh{F}_x$ est alors isomorphe à $\sh{O}_x^n$, et la proposition résulte
de (0, 5.2.7).

\paragraph{(0, 6.6.2).}
Remplacer les lignes 10 et 11 de la p.~59 par :

En effet, cela résulte de (0, 6.4.1, d)).

\paragraph{(0, 6.7.1).}
Avant la ligne 10 du bas de la p.~59, ajouter :

Si $f$ est un morphisme plat, on dit encore que $X$ est \emph{plat sur
$Y$}, ou \emph{$Y$-plat}.

\paragraph{(0, 7.2.9).}
Ligne 11 du bas de la p.~65, remplacer $M_n$ par $M_{n-1}$.

\subsection*{Chapitre I}

\paragraph{(I, 1.2.1).}
Ligne 18 de la p.~83, remplacer (deux fois) $K$ par $k$.

\paragraph{(I, 1.2.7).}
Ligne 16 du bas de la p.~84, remplacer $A$ par $A'$ (deux fois) et
$\mathfrak N$ par $\mathfrak N'$; ligne 17 du bas, remplacer $X$ par
$X'$.

\paragraph{(I, 1.7.3) et (I, 2.2.4).}
Les énoncés donnés dans ces numéros peuvent (selon une remarque due à
J.~Tate) être généralisés comme suit :

\subsection*{1.8. Morphismes d'espaces annelés en anneaux locaux dans les schémas affines}
\label{ega2:addendum:subsection:I.1.8-fr}

\begin{proposition}[1.8.1]
\label{ega2:addendum:I.1.8.1-fr}
Soient $(S,\sh{O}_S)$ un schéma affine, $(X,\sh{O}_X)$ un espace annelé
en anneaux locaux. Il y a alors une bijection canonique de l'ensemble
des homomorphismes d'anneaux
\oldpage[II]{218}
\markboth{A. GROTHENDIECK --- Chap. II}{A. GROTHENDIECK --- Chap. II}
\[
  \Gamma(S,\sh{O}_S)\to\Gamma(X,\sh{O}_X)
\]
sur l'ensemble des morphismes d'espaces annelés
$(\psi,\theta):(X,\sh{O}_X)\to(S,\sh{O}_S)$ tels que, pour tout
$x\in X$, $\theta_x^\sharp:\sh{O}_{\psi(x)}\to\sh{O}_x$ soit un
homomorphisme local.
\end{proposition}

Notons d'abord que si $(X,\sh{O}_X)$, $(S,\sh{O}_S)$ sont deux espaces
annelés quelconques, un morphisme $(\psi,\theta)$ de $(X,\sh{O}_X)$
dans $(S,\sh{O}_S)$ définit canoniquement un homomorphisme d'anneaux
$\Gamma(\theta):\Gamma(S,\sh{O}_S)\to\Gamma(X,\sh{O}_X)$; d'où une
première application
\begin{equation}
\label{ega2:addendum:I.1.8.1.1-fr}
  \rho:\Hom((X,\sh{O}_X),(S,\sh{O}_S))
  \to\Hom(\Gamma(S,\sh{O}_S),\Gamma(X,\sh{O}_X)).
\tag{1.8.1.1}
\end{equation}
Inversement, sous les hypothèses de l'énoncé, posons
$A=\Gamma(S,\sh{O}_S)$, et considérons un homomorphisme d'anneaux
$\vphi:A\to\Gamma(X,\sh{O}_X)$. Pour tout $x\in X$, il est clair que
l'ensemble des $f\in A$ tels que $\vphi(f)(x)=0$ est un idéal premier
de $A$, puisque $\sh{O}_x/\mathfrak m_x=\kres(x)$ est un corps; c'est
donc un élément de $S=\Spec(A)$, que nous noterons encore
${}^a\vphi(x)$. En outre, pour tout $f\in A$, on a par définition
(0, 5.5.2)
\[
  {}^a\vphi^{-1}(D(f))=X_f,
\]
ce qui prouve que ${}^a\vphi$ est une application continue $X\to S$.
Définissons ensuite un homomorphisme
\[
  \widetilde{\vphi}:\sh{O}_S\to{}^a\vphi_*(\sh{O}_X)
\]
de $\sh{O}_S$-Modules; pour tout $f\in A$, on a
$\Gamma(D(f),\sh{O}_S)=A_f$ (I, 1.3.6); pour tout $s\in A$, on fera
correspondre à $s/f\in A_f$ l'élément
\[
  (\vphi(s)|X_f)(\vphi(f)|X_f)^{-1}
  \in\Gamma(X_f,\sh{O}_X)
  =\Gamma(D(f),{}^a\vphi_*(\sh{O}_X)),
\]
et on vérifie aussitôt (par passage de $D(f)$ à $D(fg)$) que cela
définit bien un homomorphisme de $\sh{O}_S$-Modules, d'où un morphisme
$({}^a\vphi,\widetilde{\vphi})$ d'espaces annelés. En outre, avec les
mêmes notations, et en posant pour simplifier $y={}^a\vphi(x)$, on voit
aussitôt (0, 3.7.1) que l'on a
\[
  \widetilde{\vphi}_x^\sharp(s_y/f_y)
  =(\vphi(s)_x)(\vphi(f)_x)^{-1};
\]
comme la relation $s_y\in\mathfrak m_y$ est par définition équivalente
à $\vphi(s)_x\in\mathfrak m_x$, on voit que
$\widetilde{\vphi}_x^\sharp$ est un homomorphisme local
$\sh{O}_y\to\sh{O}_x$, et on a ainsi défini une seconde application
\[
  \sigma:\Hom(\Gamma(S,\sh{O}_S),\Gamma(X,\sh{O}_X))\to\mathfrak L,
\]
où $\mathfrak L$ est l'ensemble des morphismes
$(\psi,\theta):(X,\sh{O}_X)\to(S,\sh{O}_S)$ tels que
$\theta_x^\sharp$ soit local pour tout $x\in X$.

Il reste à prouver que $\sigma$ et $\rho$ (restreint à $\mathfrak L$)
sont réciproques l'une de l'autre; or, la définition de
$\widetilde{\vphi}$ montre aussitôt que
$\Gamma(\widetilde{\vphi})=\vphi$, et par suite $\rho\circ\sigma$ est
l'identité. Pour voir que $\sigma\circ\rho$ est l'identité, partons
d'un morphisme $(\psi,\theta)\in\mathfrak L$ et soit
$\vphi=\Gamma(\theta)$; l'hypothèse sur $\theta_x^\sharp$ permet de
déduire de cet homomorphisme, par passage aux quotients, un
monomorphisme
\[
  \theta^x:\kres(\psi(x))\to\kres(x)
\]
tel que pour toute section $f\in A=\Gamma(S,\sh{O}_S)$, on ait
$\theta^x(f(\psi(x)))=\vphi(f)(x)$; la relation $f(\psi(x))=0$ est donc
équivalente à $\vphi(f)(x)=0$, ce qui prouve déjà que
${}^a\vphi=\psi$. D'autre part, les définitions entraînent que le
diagramme
\[
  \xymatrix{
    A\ar[r]^\vphi\ar[d] &
    \Gamma(X,\sh{O}_X)\ar[d]\\
    A_{\psi(x)}\ar[r]^{\theta_x^\sharp} &
    \sh{O}_x
  }
\]
est commutatif, et il en est de même du diagramme analogue où
$\theta_x^\sharp$ est remplacé par
$\widetilde{\vphi}_x^\sharp$, d'où
$\widetilde{\vphi}_x^\sharp=\theta_x^\sharp$ (0, 1.2.4), et par suite
$\widetilde{\vphi}=\theta$.

\begin{env}[1.8.2]
\label{ega2:addendum:I.1.8.2-fr}
Lorsque $(X,\sh{O}_X)$ et $(Y,\sh{O}_Y)$ sont des espaces annelés en
anneaux locaux, nous aurons à considérer les morphismes
$(\psi,\theta):(X,\sh{O}_X)\to(Y,\sh{O}_Y)$ tels que, pour tout
$x\in X$, $\theta_x^\sharp$ soit un homomorphisme local
$\sh{O}_{\psi(x)}\to\sh{O}_x$.

Lorsque nous parlerons désormais
\oldpage[II]{219}
\markboth{ERRATA ET ADDENDA, LISTE 1}{ERRATA ET ADDENDA, LISTE 1}
de \emph{morphisme d'espaces annelés
en anneaux locaux}, c'est toujours d'un tel morphisme qu'il s'agira;
avec cette définition des morphismes, il est clair que les espaces
annelés en anneaux locaux forment une catégorie; pour deux objets
$X,Y$ de cette catégorie, $\Hom(X,Y)$ désignera donc l'ensemble des
morphismes d'espaces annelés en anneaux locaux de $X$ dans $Y$
(ensemble noté $\mathfrak L$ dans (1.8.1)); lorsque nous aurons à
considérer l'ensemble des morphismes d'espaces annelés de $X$ dans
$Y$, nous le noterons $\Hom_{\mathrm{an}}(X,Y)$ pour éviter toute
confusion. L'application (1.8.1.1) s'écrit donc
\begin{equation}
\label{ega2:addendum:I.1.8.2.1-fr}
  \rho:\Hom_{\mathrm{an}}(X,Y)
  \to\Hom(\Gamma(Y,\sh{O}_Y),\Gamma(X,\sh{O}_X)),
\tag{1.8.2.1}
\end{equation}
et sa restriction
\begin{equation}
\label{ega2:addendum:I.1.8.2.2-fr}
  \rho':\Hom(X,Y)
  \to\Hom(\Gamma(Y,\sh{O}_Y),\Gamma(X,\sh{O}_X))
\tag{1.8.2.2}
\end{equation}
est une application fonctorielle en $X$ et $Y$ dans la catégorie des
espaces annelés en anneaux locaux.
\end{env}

\begin{corollary}[1.8.3]
\label{ega2:addendum:I.1.8.3-fr}
Soit $(Y,\sh{O}_Y)$ un espace annelé en anneaux locaux. Pour que $Y$
soit un schéma affine, il faut et il suffit que pour tout espace annelé
en anneaux locaux $(X,\sh{O}_X)$, l'application (1.8.2.2) soit
bijective.
\end{corollary}

La proposition (1.8.1) montre que la condition est nécessaire.
Inversement, si on la suppose vérifiée et si l'on pose
$A=\Gamma(Y,\sh{O}_Y)$, il résulte de l'hypothèse et de (1.8.1) que les
foncteurs $X\mapsto\Hom(X,Y)$ et
$X\mapsto\Hom(X,\Spec(A))$, de la catégorie des espaces annelés en
anneaux locaux dans celle des ensembles, sont isomorphes. On sait que
cela entraîne l'existence d'un isomorphisme canonique
$X\to\Spec(A)$ (cf.~0, 8).

\begin{env}[1.8.4]
\label{ega2:addendum:I.1.8.4-fr}
Soit $S=\Spec(A)$ un schéma affine; désignons par $(S',A')$ l'espace
annelé dont l'espace sous-jacent est réduit à un point et le faisceau
structural $A'$ est le faisceau (nécessairement simple) sur $S'$ défini
par l'anneau $A$. Soit $\pi:S\to S'$ l'unique application de $S$ dans
$S'$; notons d'autre part que pour tout ouvert $U$ de $S$, on a une
application canonique
$\Gamma(S',A')=\Gamma(S,\sh{O}_S)\to\Gamma(U,\sh{O}_S)$ qui définit
donc un $\pi$-morphisme $\iota:A'\to\sh{O}_S$ de faisceaux d'anneaux.
On a ainsi défini canoniquement un morphisme d'espaces annelés
$i=(\pi,\iota):(S,\sh{O}_S)\to(S',A')$. Pour tout $A$-module $M$, nous
désignerons par $M'$ le faisceau simple sur $S'$ défini par $M$, qui
est évidemment un $A'$-Module. Il est clair que l'on a
$i_*(\widetilde M)=M'$ (I, 1.3.7).
\end{env}

\begin{lemma}[1.8.5]
\label{ega2:addendum:I.1.8.5-fr}
Avec les notations de (1.8.4), pour tout $A$-module $M$, le
$\sh{O}_S$-homomorphisme canonique fonctoriel (0, 4.4.3.3)
\begin{equation}
\label{ega2:addendum:I.1.8.5.1-fr}
  i^*(i_*(\widetilde M))\to\widetilde M
\tag{1.8.5.1}
\end{equation}
est un isomorphisme.
\end{lemma}

En effet, les deux membres de (1.8.5.1) sont exacts à droite (le
foncteur $M\mapsto i_*(\widetilde M)$ étant évidemment exact) et
commutent aux sommes directes; en considérant $M$ comme conoyau d'un
homomorphisme $A^{(I)}\to A^{(J)}$, on se ramène à prouver le lemme
dans le cas $M=A$, et il est évident dans ce cas.

\begin{corollary}[1.8.6]
\label{ega2:addendum:I.1.8.6-fr}
Soient $(X,\sh{O}_X)$ un espace annelé, $u:X\to S$ un morphisme
d'espaces
\oldpage[II]{220}
\markboth{A. GROTHENDIECK --- Chap. II}{A. GROTHENDIECK --- Chap. II}
annelés. Pour tout $A$-module $M$, on a (avec les notations de (1.8.4)) un
isomorphisme canonique fonctoriel de $\sh{O}_X$-Modules
\begin{equation}
\label{ega2:addendum:I.1.8.6.1-fr}
  u^*(\widetilde M)\isoto u^*(i^*(M')).
\tag{1.8.6.1}
\end{equation}
\end{corollary}

\begin{corollary}[1.8.7]
\label{ega2:addendum:I.1.8.7-fr}
Sous les hypothèses de (1.8.6), on a, pour tout $A$-module $M$ et tout
$\sh{O}_X$-Module $\sh{F}$, un isomorphisme canonique fonctoriel en
$M$ et $\sh{F}$
\begin{equation}
\label{ega2:addendum:I.1.8.7.1-fr}
  \Hom_{\sh{O}_S}(\widetilde M,u_*(\sh{F}))
  \isoto\Hom_A(M,\Gamma(X,\sh{F})).
\tag{1.8.7.1}
\end{equation}
\end{corollary}

On a en effet, en vertu de (0, 4.4.3) et du lemme (1.8.5), un
isomorphisme canonique de bifoncteurs
\[
  \Hom_{\sh{O}_S}(\widetilde M,u_*(\sh{F}))
  \isoto\Hom_{A'}(M',i_*(u_*(\sh{F}))),
\]
et il est clair que le second membre n'est autre que
$\Hom_A(M,\Gamma(X,\sh{F}))$. On notera que l'homomorphisme canonique
(1.8.7.1) fait correspondre à tout $\sh{O}_S$-homomorphisme
$h:\widetilde M\to u_*(\sh{F})$ (autrement dit, à tout $u$-morphisme
$\widetilde M\to\sh{F}$) le $A$-homomorphisme
\[
  \Gamma(h):M\to\Gamma(S,u_*(\sh{F}))=\Gamma(X,\sh{F}).
\]

\begin{env}[1.8.8]
\label{ega2:addendum:I.1.8.8-fr}
Avec les notations de (1.8.4), il est clair (0, 4.1.1) que tout
morphisme d'espaces annelés $(\psi,\theta):X\to S'$ équivaut à la
donnée d'un homomorphisme d'anneaux
$A\to\Gamma(X,\sh{O}_X)$. On peut donc encore interpréter (1.8.1)
comme définissant une bijection canonique
\[
  \Hom(X,S)\isoto\Hom(X,S')
\]
(où bien entendu il s'agit au second membre de morphismes d'espaces
annelés, puisqu'en général $A$ n'est pas un anneau local). Plus
généralement, si $X,Y$ sont deux espaces annelés en anneaux locaux et
si $(Y',A')$ est l'espace annelé dont l'espace sous-jacent est réduit à
un point et dont le faisceau d'anneaux $A'$ est le faisceau simple
défini par l'anneau $\Gamma(Y,\sh{O}_Y)$, on peut interpréter
(1.8.2.1) comme une application
\begin{equation}
\label{ega2:addendum:I.1.8.8.1-fr}
  \rho:\Hom_{\mathrm{an}}(X,Y)\to\Hom(X,Y').
\tag{1.8.8.1}
\end{equation}
Le résultat de (1.8.3) s'interprète donc en disant que les schémas
affines sont caractérisés parmi les espaces annelés en anneaux locaux
comme ceux pour lesquels la restriction de $\rho$ à $\Hom(X,Y)$ :
\begin{equation}
\label{ega2:addendum:I.1.8.8.2-fr}
  \rho':\Hom(X,Y)\to\Hom(X,Y')
\tag{1.8.8.2}
\end{equation}
est bijective pour tout espace annelé en anneaux locaux $X$. Dans un
chapitre ultérieur, nous généraliserons cette définition, ce qui
permettra d'associer à un espace annelé quelconque $Z$ (et non plus
seulement à un espace annelé dont l'espace sous-jacent est réduit à un
point) un espace annelé en anneaux locaux que nous appellerons encore
$\Spec(Z)$; cela sera le point de départ d'une théorie «~relative~» des
préschémas au-dessus d'espaces annelés quelconques, étendant les
résultats du chap.~I.
\end{env}

\begin{env}[1.8.9]
\label{ega2:addendum:I.1.8.9-fr}
On peut considérer les couples $(X,\sh{F})$ formés d'un espace annelé
en anneaux locaux $X$ et d'un $\sh{O}_X$-Module $\sh{F}$ comme formant
une catégorie, un morphisme $(X,\sh{F})\to(Y,\sh{G})$ de cette catégorie
étant un couple $(u,h)$ formé d'un morphisme d'espaces
\oldpage[II]{221}
\markboth{ERRATA ET ADDENDA, LISTE 1}{ERRATA ET ADDENDA, LISTE 1}
annelés en anneaux locaux $u:X\to Y$ et d'un $u$-morphisme
$h:\sh{G}\to\sh{F}$ de Modules; ces
morphismes (pour $(X,\sh{F})$ et $(Y,\sh{G})$ fixés) forment un ensemble
que nous noterons $\Hom((X,\sh{F}),(Y,\sh{G}))$; l'application
$(u,h)\mapsto(\rho'(u),\Gamma(h))$ est une application canonique
\begin{equation}
\label{ega2:addendum:I.1.8.9.1-fr}
\begin{split}
  \Hom((X,\sh{F}),(Y,\sh{G}))\to
  \Hom\bigl(
    &(\Gamma(Y,\sh{O}_Y),\Gamma(Y,\sh{G})),\\
    &(\Gamma(X,\sh{O}_X),\Gamma(X,\sh{F}))
  \bigr),
\end{split}
\tag{1.8.9.1}
\end{equation}
fonctorielle en $(X,\sh{F})$ et $(Y,\sh{G})$, le second membre étant
l'ensemble des di-homomorphismes correspondant aux anneaux et aux
modules considérés (0, 1.0.2).
\end{env}

\begin{corollary}[1.8.10]
\label{ega2:addendum:I.1.8.10-fr}
Soient $Y$ un espace annelé en anneaux locaux, $\sh{G}$ un
$\sh{O}_Y$-Module. Pour que $Y$ soit un schéma affine et $\sh{G}$ un
$\sh{O}_Y$-Module quasi-cohérent, il faut et il suffit que pour tout
couple $(X,\sh{F})$ formé d'un espace annelé en anneaux locaux $X$ et
d'un $\sh{O}_X$-Module $\sh{F}$, l'application canonique (1.8.8.1) soit
bijective.
\end{corollary}

Nous laissons au lecteur le soin de développer le raisonnement, calqué
sur celui de (1.8.3), et où on utilise (1.8.1) et (1.8.7).

\begin{remark}[1.8.11]
\label{ega2:addendum:I.1.8.11-fr}
Les énoncés (I, 1.7.3), (I, 1.7.4) et (I, 2.2.4) sont des cas
particuliers de (1.8.1), ainsi que la définition (I, 1.6.1); de même
(I, 2.2.5) résulte de (1.8.7). La proposition (1.8.7) entraîne aussi
(I, 1.6.3) (et par suite (I, 1.6.4)) comme cas particulier, car si $X$
est affine et $\Gamma(X,\sh{F})=N$, les foncteurs
\[
  M\longmapsto
  \Hom_{\sh{O}_S}(\widetilde M,u_*(\widetilde N))
  \quad\hbox{et}\quad
  M\longmapsto
  \Hom_{\sh{O}_S}(\widetilde M,(N_{[\vphi]})\supertilde)
\]
(où $\vphi:A\to\Gamma(X,\sh{O}_X)$ correspond à $u$) sont isomorphes
par (1.8.7) et (I, 1.3.8). Enfin, (I, 1.6.5) (et par suite (I, 1.6.6))
résulte de (1.8.6) et du fait que pour tout $f\in A$, les $A_f$-modules
\[
  N'\otimes_{A'}A_f
  \quad\hbox{et}\quad
  (N'\otimes_{A'}A)_f
\]
(notations de (I, 1.6.5)) sont canoniquement isomorphes.
\end{remark}

\paragraph{Insérer après (I, 3.2.8) :}

\begin{env}[3.2.9]
\label{ega2:addendum:I.3.2.9-fr}
Il résulte de (1.8.1) que l'on peut préciser (I, 3.2.2) de la façon
suivante : $Z=\Spec(B\otimes_A C)$ est non seulement un produit de
$X=\Spec(B)$ et de $Y=\Spec(C)$ dans la catégorie des $S$-préschémas,
mais aussi dans la catégorie des espaces annelés en anneaux locaux,
au-dessus de $S$ (avec une définition des $S$-morphismes calquée sur
celle de (I, 2.5.2)). La démonstration de (I, 3.2.6) prouve alors en
fait que pour deux $S$-préschémas quelconques $X,Y$, le préschéma
$X\times_S Y$ est non seulement le produit de $X$ et $Y$ dans la
catégorie des $S$-préschémas, mais encore dans la catégorie des espaces
annelés en anneaux locaux au-dessus du préschéma $S$.
\end{env}

\paragraph{(I, 4.4.5).}
Ligne 14 du bas de la p.~126, remplacer $B$ par $A$; ligne 13 du bas,
remplacer «~$A$-algèbre~» par «~$B$-algèbre~».

\paragraph{(I, 5.3.5).}
Ligne 2 du bas de la p.~132, remplacer $Y\to S$ par $g:Y\to S$.

\paragraph{(I, 6.3.2.1).}
Ligne 6 du bas de la p.~144, lire :
\[
  D(g_i)\subset W.
\]

\paragraph{(I, 7.3.8).}
Remplacer le texte actuel, qui est erroné, par le suivant :

Soient $X,Y$ deux préschémas intègres, ce qui entraîne que
$\sh{R}(X)$ (resp. $\sh{R}(Y)$) est un $\sh{O}_X$-Module (resp. un
$\sh{O}_Y$-Module) quasi-cohérent (I, 7.3.3). Soit $f:X\to Y$ un
morphisme dominant; alors il existe un homomorphisme canonique de
$\sh{O}_X$-Modules
\oldpage[II]{222}
\markboth{A. GROTHENDIECK --- Chap. II}{A. GROTHENDIECK --- Chap. II}
\begin{equation}
\label{ega2:addendum:I.7.3.8.1-fr}
  \tau:f^*(\sh{R}(Y))\to\sh{R}(X).
\tag{7.3.8.1}
\end{equation}

Supposons d'abord $X=\Spec(A)$ et $Y=\Spec(B)$ affines d'anneaux
intègres $A$ et $B$, $f$ correspondant donc à un homomorphisme injectif
$B\to A$, qui se prolonge en un monomorphisme $L\to K$ du corps des
fractions $L$ de $B$ dans le corps des fractions $K$ de $A$.
L'homomorphisme (7.3.8.1) correspond alors à l'homomorphisme canonique
\[
  L\otimes_B A\to K
\]
(I, 1.6.5). Dans le cas général, pour tout couple d'ouverts affines non
vides $U\subset X$, $V\subset Y$ tels que $f(U)\subset V$, on définit
de la façon précédente un homomorphisme $\tau_{U,V}$ et on constate
aussitôt que si $U'\subset U$, $V'\subset V$,
$f(U')\subset V'$, $\tau_{U,V}$ prolonge $\tau_{U',V'}$, d'où notre
assertion. Si $x,y$ sont les points génériques de $X$ et $Y$
respectivement, on a $f(x)=y$,
\[
  \bigl(f^*(\sh{R}(Y))\bigr)_x
  =\sh{O}_y\otimes_{\sh{O}_y}\sh{O}_x
  =\sh{O}_x
\]
(0, 4.3.1), et $\tau_x$ est donc un isomorphisme.

\paragraph{(I, 9.5.2).}
Dans les 3 dernières lignes de la p.~176 et les 5 premières de la
p.~177, remplacer partout $X'$ par $Y'$ et $X''$ par $Y''$.

\paragraph{(I, 10.5.4).}
Ligne 11 du bas de la p.~187, remplacer (10.3.4) par (10.3.5).

\paragraph{(I, 10.6.3).}
Ligne 3 du bas de la p.~189, remplacer $X$ par $\mathfrak X$.

\paragraph{(I, 10.14.2).}
Ligne 5 de la p.~210, remplacer «~$\sh{O}_{\mathfrak X}$-Module cohérent~» par
«~Idéal cohérent de $\sh{O}_{\mathfrak X}$~»; ligne 7, remplacer
«~$\sh{O}_{\mathfrak X}$-Modules cohérents~» par «~Idéaux cohérents de
$\sh{O}_{\mathfrak X}$~».

% FIN SOURCE ega2/ega2-errata-addenda-fr.tex

%
% DEBUT SOURCE ega3/ega3-1-fr.tex
\clearpage
\thispagestyle{plain}
\markboth{PRÉLIMINAIRES}{PRÉLIMINAIRES}
\oldpage[0\textsubscript{III}]{5}
\phantomsection
\label{chapter:0III-fr}

\begin{center}
{\large CHAPITRE 0 (suite)\footnotemark\par}
\vspace{1.5em}
{\Large\bfseries PRÉLIMINAIRES\par}
\end{center}

\footnotetext{Pour faciliter la recherche des références, on renverra
désormais aux paragraphes du chapitre 0 publiés avec le chapitre I par le
signe $0_{\mathrm I}$.}

\setcounter{section}{7}
\section{Foncteurs représentables}
\label{section:0.8-fr}

\setcounter{subsection}{0}
\subsection{Foncteurs représentables}
\label{subsection:0.8.1-fr}

\begin{env}[8.1.1]
\label{0.8.1.1-fr}
Nous désignerons par $\mathrm{Ens}$ la catégorie des ensembles. Soit $C$ une
catégorie~; pour deux objets $X$, $Y$ de $C$, nous poserons
$h_X(Y)=\operatorname{Hom}(Y,X)$~; pour tout morphisme $u:Y\to Y'$ dans $C$,
nous désignerons par $h_X(u)$ l'application $v\mapsto vu$ de
$\operatorname{Hom}(Y',X)$ dans $\operatorname{Hom}(Y,X)$. Il est immédiat
qu'avec ces définitions, $h_X:C\to\mathrm{Ens}$ est un \emph{foncteur
contravariant}, c'est-à-dire un objet de la catégorie, notée
$\operatorname{Hom}(C^0,\mathrm{Ens})$, des foncteurs covariants de la
catégorie $C^0$, duale de la catégorie $C$, dans la catégorie
$\mathrm{Ens}$ (T, 1.7, d) et [29]).
\end{env}

\begin{env}[8.1.2]
\label{0.8.1.2-fr}
Soit maintenant $w:X\to X'$ un morphisme dans $C$~; pour tout $Y\in C$ et
tout $v\in\operatorname{Hom}(Y,X)=h_X(Y)$, on a
$wv\in\operatorname{Hom}(Y,X')=h_{X'}(Y)$~; désignons par $h_w(Y)$
l'application $v\mapsto wv$ de $h_X(Y)$ dans $h_{X'}(Y)$. Il est immédiat
que pour tout morphisme $u:Y\to Y'$ dans $C$, le diagramme
\[
  \xymatrix{
    h_X(Y')\ar[r]^{h_X(u)}\ar[d]_{h_w(Y')} &
    h_X(Y)\ar[d]^{h_w(Y)}\\
    h_{X'}(Y')\ar[r]_{h_{X'}(u)} &
    h_{X'}(Y)
  }
\]
est commutatif~; autrement dit, $h_w$ est un \emph{morphisme fonctoriel}
$h_X\to h_{X'}$ (T, 1.2), ou encore un morphisme dans la catégorie
$\operatorname{Hom}(C^0,\mathrm{Ens})$ (T, 1.7, d)). Les définitions de
$h_X$ et de $h_w$ constituent donc la définition d'un \emph{foncteur
covariant canonique}
\[
  h:C\to\operatorname{Hom}(C^0,\mathrm{Ens}),\qquad X\mapsto h_X.
  \tag{8.1.2.1}
\]
\end{env}

\begin{env}[8.1.3]
\label{0.8.1.3-fr}
Soient $X$ un objet de $C$, $F$ un foncteur contravariant de $C$ dans
$\mathrm{Ens}$ (objet de $\operatorname{Hom}(C^0,\mathrm{Ens})$). Soit
$g:h_X\to F$ un \emph{morphisme fonctoriel}~: pour tout $Y\in C$,
\oldpage[0\textsubscript{III}]{6}
$g(Y)$ est donc une application $h_X(Y)\to F(Y)$ telle que pour tout morphisme
$u:Y\to Y'$ dans $C$, le diagramme
\[
  \xymatrix{
    h_X(Y')\ar[r]^{h_X(u)}\ar[d]_{g(Y')} &
    h_X(Y)\ar[d]^{g(Y)}\\
    F(Y')\ar[r]_{F(u)} & F(Y)
  }
  \tag{8.1.3.1}\label{0.8.1.3.1-fr}
\]
soit commutatif. En particulier, on a une application
$g(X):h_X(X)=\operatorname{Hom}(X,X)\to F(X)$, d'où un élément
\[
  \alpha(g)=(g(X))(1_X)\in F(X)
  \tag{8.1.3.2}
\]
et par suite une application canonique
\[
  \alpha:\operatorname{Hom}(h_X,F)\to F(X).
  \tag{8.1.3.3}
\]

Inversement, considérons un élément $\xi\in F(X)$~; pour tout morphisme
$v:Y\to X$ dans $C$, $F(v)$ est une application $F(X)\to F(Y)$~;
considérons l'application
\[
  v\mapsto(F(v))(\xi)
  \tag{8.1.3.4}
\]
de $h_X(Y)$ dans $F(Y)$~; si on désigne par $(\beta(\xi))(Y)$ cette
application,
\[
  \beta(\xi):h_X\to F
  \tag{8.1.3.5}
\]
est un \emph{morphisme fonctoriel}, car on a pour tout morphisme
$u:Y\to Y'$ dans $C$,
$(F(vu))(\xi)=(F(v)\circ F(u))(\xi)$, ce qui vérifie la commutativité de
(\hyperref[0.8.1.3.1-fr]{8.1.3.1}) pour $g=\beta(\xi)$. On a ainsi défini
une application canonique
\[
  \beta:F(X)\to\operatorname{Hom}(h_X,F).
  \tag{8.1.3.6}
\]
\end{env}

\begin{proposition}[8.1.4]
\label{0.8.1.4-fr}
Les applications $\alpha$ et $\beta$ sont des bijections réciproques l'une de
l'autre.
\end{proposition}

Calculons $\alpha(\beta(\xi))$ pour $\xi\in F(X)$~; pour tout $Y\in C$,
$(\beta(\xi))(Y)$ est l'application
$g_1(Y):v\mapsto(F(v))(\xi)$ de $h_X(Y)$ dans $F(Y)$. On a donc
\[
  \alpha(\beta(\xi))=(g_1(X))(1_X)=(F(1_X))(\xi)
  =1_{F(X)}(\xi)=\xi.
\]
Calculons maintenant $\beta(\alpha(g))$ pour
$g\in\operatorname{Hom}(h_X,F)$~; pour tout $Y\in C$,
$(\beta(\alpha(g)))(Y)$ est l'application
$v\mapsto(F(v))((g(X))(1_X))$~; en vertu de la commutativité de
(\hyperref[0.8.1.3.1-fr]{8.1.3.1}), cette application n'est autre que
$v\mapsto(g(Y))((h_X(v))(1_X))=(g(Y))(v)$ par définition de $h_X(v)$,
autrement dit, elle est égale à $g(Y)$, ce qui démontre la proposition.

\begin{env}[8.1.5]
\label{0.8.1.5-fr}
Rappelons qu'une \emph{sous-catégorie} $C'$ d'une catégorie $C$ est définie
par la condition que ses objets soient des objets de $C$, et que si $X'$,
$Y'$ sont deux objets de $C'$, l'ensemble
$\operatorname{Hom}_{C'}(X',Y')$ des morphismes $X'\to Y'$ \emph{dans
$C'$} est une partie de l'ensemble $\operatorname{Hom}_C(X',Y')$ des
morphismes $X'\to Y'$ \emph{dans $C$}, l'application canonique de
<<~composition des morphismes~>>
\[
  \operatorname{Hom}_{C'}(X',Y')\times\operatorname{Hom}_{C'}(Y',Z')
  \to\operatorname{Hom}_{C'}(X',Z')
\]
\oldpage[0\textsubscript{III}]{7}
étant la restriction de l'application canonique
\[
  \operatorname{Hom}_C(X',Y')\times\operatorname{Hom}_C(Y',Z')
  \to\operatorname{Hom}_C(X',Z').
\]

On dit que $C'$ est une \emph{sous-catégorie pleine} de $C$ si
$\operatorname{Hom}_{C'}(X',Y')=\operatorname{Hom}_C(X',Y')$ pour tout
couple d'objets de $C'$. La sous-catégorie $C''$ de $C$ formée des objets de
$C$ isomorphes aux objets de $C'$ est encore alors une sous-catégorie pleine
de $C$, \emph{équivalente} (T, 1.2) à $C'$ comme on le vérifie sans peine.

Un foncteur covariant $F:C_1\to C_2$ est dit \emph{pleinement fidèle} si,
pour tout couple d'objets $X_1$, $Y_1$ de $C_1$, l'application
$u\mapsto F(u)$ de $\operatorname{Hom}(X_1,Y_1)$ dans
$\operatorname{Hom}(F(X_1),F(Y_1))$ est \emph{bijective}~; cela entraîne que
la sous-catégorie $F(C_1)$ de $C_2$ est \emph{pleine}. En outre, si deux
objets $X_1$, $X'_1$ ont même image $X_2$, il existe un isomorphisme unique
$u:X_1\to X'_1$ tel que $F(u)=1_{X_2}$. Pour tout objet $X_2$ de $F(C_1)$,
soit alors $G(X_2)$ un des objets $X_1$ de $C_1$ tel que $F(X_1)=X_2$
($G$ étant défini au moyen de l'axiome de choix)~; pour tout morphisme
$v:X_2\to Y_2$ dans $F(C_1)$, $G(v)$ sera l'unique morphisme
$u:G(X_2)\to G(Y_2)$ tel que $F(u)=v$~; $G$ est alors un \emph{foncteur} de
$F(C_1)$ dans $C_1$~; $FG$ est le foncteur identique dans $F(C_1)$, et ce
qui précède montre qu'il existe un isomorphisme de foncteurs
$\varphi:1_{C_1}\to GF$ tel que $F$, $G$, $\varphi$ et l'identité
$1_{F(C_1)}\to FG$ définissent une \emph{équivalence} de la catégorie $C_1$
et de la sous-catégorie pleine $F(C_1)$ de $C_2$ (T, 1.2).
\end{env}

\begin{env}[8.1.6]
\label{0.8.1.6-fr}
Appliquons la prop. (\hyperref[0.8.1.4-fr]{8.1.4}) au cas où le foncteur $F$
est $h_{X'}$, $X'$ étant un objet quelconque de $C$~; l'application
$\beta:\operatorname{Hom}(X,X')\to\operatorname{Hom}(h_X,h_{X'})$ n'est
autre ici que l'application $w\mapsto h_w$ définie dans
(\hyperref[0.8.1.2-fr]{8.1.2})~; cette application étant \emph{bijective},
on voit, avec la terminologie de (\hyperref[0.8.1.5-fr]{8.1.5}), que~:
\end{env}

\begin{proposition}[8.1.7]
\label{0.8.1.7-fr}
Le foncteur canonique
$h:C\to\operatorname{Hom}(C^0,\mathrm{Ens})$ est pleinement fidèle.
\end{proposition}

\begin{env}[8.1.8]
\label{0.8.1.8-fr}
Soit $F$ un foncteur contravariant de $C$ dans $\mathrm{Ens}$~; on dit que
$F$ est \emph{représentable} s'il existe un objet $X\in C$ tel que $F$ soit
\emph{isomorphe} à $h_X$~; il résulte de
(\hyperref[0.8.1.7-fr]{8.1.7}) que la donnée d'un $X\in C$ et d'un
isomorphisme de foncteurs $g:h_X\to F$ détermine $X$ à un isomorphisme
unique près. La prop. (\hyperref[0.8.1.7-fr]{8.1.7}) signifie encore que $h$
définit une \emph{équivalence} de $C$ et de la sous-catégorie pleine de
$\operatorname{Hom}(C^0,\mathrm{Ens})$ formée des \emph{foncteurs
contravariants représentables}. Il résulte d'ailleurs de
(\hyperref[0.8.1.4-fr]{8.1.4}) que la donnée d'un morphisme fonctoriel
$g:h_X\to F$ équivaut à celle d'un élément $\xi\in F(X)$~; dire que $g$ est
un \emph{isomorphisme} équivaut pour $\xi$ à la condition suivante~:
\emph{pour tout objet $Y$ de $C$ l'application
$v\mapsto(F(v))(\xi)$ de $\operatorname{Hom}(Y,X)$ dans $F(Y)$ est
bijective}. Lorsque $\xi$ vérifie cette condition, on dira que le couple
$(X,\xi)$ \emph{représente} le foncteur représentable $F$. Par abus de
langage, on dira aussi que l'objet $X\in C$ représente $F$ s'il existe
$\xi\in F(X)$ tel que $(X,\xi)$ représente $F$, autrement dit si $h_X$ est
isomorphe à $F$.

Soient $F$, $F'$ deux foncteurs contravariants représentables de $C$ dans
$\mathrm{Ens}$, $h_X\to F$ et $h_{X'}\to F'$ deux isomorphismes de
foncteurs. Alors il résulte de (\hyperref[0.8.1.6-fr]{8.1.6}) qu'il y a une
correspondance biunivoque canonique entre $\operatorname{Hom}(X,X')$ et
l'ensemble $\operatorname{Hom}(F,F')$ des morphismes fonctoriels $F\to F'$.
\end{env}

\begin{env}[8.1.9]
\label{0.8.1.9-fr}
\textit{Exemples. I : limites projectives.} La notion de foncteur
contravariant représentable couvre en particulier la notion <<~duale~>> de
la notion usuelle de <<~solution d'un
\oldpage[0\textsubscript{III}]{8}
problème universel~>>. Plus généralement, nous allons voir que la notion de
\emph{limite projective} est un cas particulier de celle de foncteur
représentable. Rappelons que dans une catégorie $C$, on définit un
\emph{système projectif} par la donnée d'un ensemble préordonné $I$, d'une
famille $(A_\alpha)_{\alpha\in I}$ d'objets de $C$, et, pour tout couple
d'indices $(\alpha,\beta)$ tel que $\alpha\leq\beta$, d'un morphisme
$u_{\alpha\beta}:A_\beta\to A_\alpha$. Une \emph{limite projective} de ce
système dans $C$ est constituée par un objet $B$ de $C$ (noté
$\varprojlim A_\alpha$), et, pour chaque $\alpha\in I$, un morphisme
$u_\alpha:B\to A_\alpha$, tels que~: 1\textsuperscript{o}
$u_\alpha=u_{\alpha\beta}u_\beta$ pour $\alpha\leq\beta$~;
2\textsuperscript{o} Pour tout objet $X$ de $C$ et toute famille
$(v_\alpha)_{\alpha\in I}$ de morphismes $v_\alpha:X\to A_\alpha$, telle
que $v_\alpha=u_{\alpha\beta}v_\beta$ pour $\alpha\leq\beta$, il existe un
morphisme unique $v:X\to B$ (noté $\varprojlim v_\alpha$) tel que
$v_\alpha=u_\alpha v$ pour tout $\alpha\in I$ (T, 1.8). Ceci s'interprète de
la façon suivante~: les $u_{\alpha\beta}$ définissent canoniquement des
applications
\[
  \overline{u}_{\alpha\beta}:\operatorname{Hom}(X,A_\beta)
  \to\operatorname{Hom}(X,A_\alpha)
\]
qui définissent un \emph{système projectif} d'ensembles
$(\operatorname{Hom}(X,A_\alpha),\overline{u}_{\alpha\beta})$, et
$(v_\alpha)$ est par définition un élément de l'ensemble
$\varprojlim_\alpha\operatorname{Hom}(X,A_\alpha)$~; il est clair que
$X\mapsto\varprojlim_\alpha\operatorname{Hom}(X,A_\alpha)$ est un
\emph{foncteur contravariant} de $C$ dans $\mathrm{Ens}$, et l'existence de
la limite projective $B$ équivaut à dire que
$(v_\alpha)\mapsto\varprojlim v_\alpha$ est un \emph{isomorphisme} de
foncteurs en $X$
\[
  \varprojlim_\alpha\operatorname{Hom}(X,A_\alpha)
  \xrightarrow{\sim}\operatorname{Hom}(X,B)
  \tag{8.1.9.1}
\]
autrement dit que le foncteur
$X\mapsto\varprojlim_\alpha\operatorname{Hom}(X,A_\alpha)$ est
\emph{représentable}.
\end{env}

\begin{env}[8.1.10]
\label{0.8.1.10-fr}
\textit{Exemples. II : Objet final.} Soient $C$ une catégorie, $\{a\}$ un
ensemble réduit à un seul élément. Considérons le foncteur contravariant
$F:C\to\mathrm{Ens}$ qui, à tout objet $X$ de $C$ fait correspondre
l'ensemble $\{a\}$ et à tout morphisme $X\to X'$ dans $C$ l'unique
application $\{a\}\to\{a\}$. Dire que ce foncteur est \emph{représentable}
signifie qu'il existe un objet $e\in C$ tel que pour tout $Y\in C$,
$\operatorname{Hom}(Y,e)=h_e(Y)$ soit \emph{réduit à un élément}~; on dit
que $e$ est un \emph{objet final} de $C$, et il est clair que deux objets
finaux de $C$ sont isomorphes (ce qui permet de définir, en général à l'aide
de l'axiome de choix, un objet final de $C$ qu'on note alors $e_C$). Par
exemple, dans la catégorie $\mathrm{Ens}$, les objets finaux sont les
ensembles réduits à un élément~; dans la catégorie des \emph{algèbres
augmentées} sur un corps $K$ (où les morphismes sont les homomorphismes
d'algèbres compatibles avec les augmentations), $K$ est un objet final~;
dans la catégorie des \emph{$S$-préschémas} (I, 2.5.1), $S$ est un objet
final.
\end{env}

\begin{env}[8.1.11]
\label{0.8.1.11-fr}
Pour deux objets $X$, $Y$ d'une catégorie $C$, posons
$h'_X(Y)=\operatorname{Hom}(X,Y)$ et pour tout morphisme $u:Y\to Y'$, soit
$h'_X(u)$ l'application $v\mapsto uv$ de $\operatorname{Hom}(X,Y)$ dans
$\operatorname{Hom}(X,Y')$~; $h'_X$ est alors un \emph{foncteur covariant}
$C\to\mathrm{Ens}$, d'où l'on déduit comme dans
(\hyperref[0.8.1.2-fr]{8.1.2}) la définition d'un foncteur covariant
canonique $h':C^0\to\operatorname{Hom}(C,\mathrm{Ens})$~; un \emph{foncteur
covariant} $F$ de $C$ dans $\mathrm{Ens}$, autrement dit un objet de
$\operatorname{Hom}(C,\mathrm{Ens})$ est alors dit \emph{représentable}
s'il existe un objet $X\in C$ (nécessairement unique à isomorphisme unique
près) tel que $F$ soit \emph{isomorphe à $h'_X$}~; nous laissons au lecteur
le soin de développer les considérations <<~duales~>> des précédentes pour
cette notion, qui couvre cette fois celle de \emph{limite inductive}, et en
particulier la notion usuelle de <<~solution de problème universel~>>.
\end{env}

\subsection{Structures algébriques dans les catégories}
\label{subsection:0.8.2-fr}

\begin{env}[8.2.1]
\label{0.8.2.1-fr}
\oldpage[0\textsubscript{III}]{9}
Étant donnés deux foncteurs contravariants $F$, $F'$ de $C$ dans
$\mathrm{Ens}$, rappelons que pour tout objet $Y\in C$, on pose
$(F\times F')(Y)=F(Y)\times F'(Y)$, et pour tout morphisme $u:Y\to Y'$ dans
$C$, on pose $(F\times F')(u)=F(u)\times F'(u)$ qui est l'application
$(t,t')\mapsto((F(u))(t),(F'(u))(t'))$ de $F(Y')\times F'(Y')$ dans
$F(Y)\times F'(Y)$~; $F\times F':C\to\mathrm{Ens}$ est donc un
\emph{foncteur contravariant} (qui n'est autre d'ailleurs que le
\emph{produit} des objets $F$, $F'$ dans la catégorie
$\operatorname{Hom}(C^0,\mathrm{Ens})$). Étant donné un objet $X\in C$, nous
appellerons \emph{loi de composition interne sur $X$} un morphisme fonctoriel
\[
  \gamma_X:h_X\times h_X\to h_X.
  \tag{8.2.1.1}
\]

Autrement dit (T, 1.2), pour tout objet $Y\in C$, $\gamma_X(Y)$ est une
application $h_X(Y)\times h_X(Y)\to h_X(Y)$ (donc par définition une
\emph{loi de composition interne} sur l'ensemble $h_X(Y)$) soumise à la
condition que, pour tout morphisme $u:Y\to Y'$ dans $C$, le diagramme
\[
  \xymatrix{
    h_X(Y')\times h_X(Y')
      \ar[rr]^{h_X(u)\times h_X(u)}\ar[dd]_{\gamma_X(Y')} &&
    h_X(Y)\times h_X(Y)\ar[dd]^{\gamma_X(Y)}\\\\
    h_X(Y')\ar[rr]_{h_X(u)} && h_X(Y)
  }
\]
soit commutatif~; cela signifie que pour les lois de composition
$\gamma_X(Y)$ et $\gamma_X(Y')$, $h_X(u)$ est un \emph{homomorphisme} de
$h_X(Y')$ dans $h_X(Y)$.

De la même façon, étant donnés deux objets $Z$, $X$ de $C$, on appelle
\emph{loi de composition externe sur $X$, ayant $Z$ comme domaine
d'opérateurs}, un morphisme fonctoriel
\[
  \omega_{X,Z}:h_Z\times h_X\to h_X.
  \tag{8.2.1.2}
\]

On voit comme ci-dessus que pour tout $Y\in C$, $\omega_{X,Z}(Y)$ est une loi
de composition externe sur $h_X(Y)$, ayant $h_Z(Y)$ comme domaine
d'opérateurs, et telle que pour tout morphisme $u:Y\to Y'$, $h_X(u)$ et
$h_Z(u)$ forment un \emph{dihomomorphisme} de
$(h_Z(Y'),h_X(Y'))$ dans $(h_Z(Y),h_X(Y))$.
\end{env}

\begin{env}[8.2.2]
\label{0.8.2.2-fr}
Soit $X'$ un second objet de $C$, et supposons donnée sur $X'$ une loi de
composition interne $\gamma_{X'}$~; nous dirons qu'un morphisme
$w:X\to X'$ dans $C$ est un \emph{homomorphisme} pour ces lois de
composition, si pour tout $Y\in C$, $h_w(Y):h_X(Y)\to h_{X'}(Y)$ est un
homomorphisme pour les lois de composition $\gamma_X(Y)$ et
$\gamma_{X'}(Y)$. Si $X''$ est un troisième
\oldpage[0\textsubscript{III}]{10}
objet de $C$ muni d'une loi de composition interne $\gamma_{X''}$, et
$w':X'\to X''$ un morphisme dans $C$ qui est un homomorphisme pour
$\gamma_{X'}$ et $\gamma_{X''}$, il est clair que le morphisme
$w'w:X\to X''$ est un homomorphisme pour les lois de composition $\gamma_X$
et $\gamma_{X''}$. Un isomorphisme $w:X\xrightarrow{\sim}X'$ dans $C$ est
appelé \emph{isomorphisme pour les lois de composition $\gamma_X$ et
$\gamma_{X'}$} si $w$ est un homomorphisme pour ces lois de composition, et
si son morphisme réciproque $w^{-1}$ est un homomorphisme pour les lois de
composition $\gamma_{X'}$ et $\gamma_X$.

On définit de la même manière les \emph{dihomomorphismes} pour les couples
d'objets de $C$ munis de lois de composition externes.
\end{env}

\begin{env}[8.2.3]
\label{0.8.2.3-fr}
Lorsqu'une loi de composition interne $\gamma_X$ sur un objet $X\in C$ est
telle que $\gamma_X(Y)$ soit une \emph{loi de groupe} sur $h_X(Y)$ pour
\emph{tout} $Y\in C$, on dit que $X$, muni de cette loi, est un
\emph{$C$-groupe} ou un \emph{$C$-objet en groupes}. On définit de même les
\emph{$C$-anneaux}, \emph{$C$-modules}, etc.
\end{env}

\begin{env}[8.2.4]
\label{0.8.2.4-fr}
Supposons que le \emph{produit} $X\times X$ d'un objet $X\in C$ par lui-même
existe dans $C$~; par définition, on a alors
$h_{X\times X}=h_X\times h_X$ à un isomorphisme canonique près, puisqu'il
s'agit d'un cas particulier de limite projective
(\hyperref[0.8.1.9-fr]{8.1.9})~; une loi de composition interne sur $X$ peut
donc être considérée comme un morphisme fonctoriel
$\gamma_X:h_{X\times X}\to h_X$, et détermine donc canoniquement
(\hyperref[0.8.1.6-fr]{8.1.6}) un élément
$c_X\in\operatorname{Hom}(X\times X,X)$ tel que $h_{c_X}=\gamma_X$~; dans
ce cas, la donnée d'une loi de composition interne sur $X$ est donc
équivalente à celle d'un morphisme $X\times X\to X$~; lorsque $C$ est la
catégorie $\mathrm{Ens}$, on retrouve la notion classique de loi de
composition interne sur un ensemble. On a un résultat analogue pour une loi
de composition externe lorsque le produit $Z\times X$ existe dans $C$.
\end{env}

\begin{env}[8.2.5]
\label{0.8.2.5-fr}
Avec les notations précédentes, supposons en outre que $X\times X\times X$
existe dans $C$~; la caractérisation du produit comme objet représentant un
foncteur (\hyperref[0.8.1.9-fr]{8.1.9}) entraîne l'existence d'isomorphismes
canoniques
\[
  (X\times X)\times X\xrightarrow{\sim}X\times X\times X
  \xrightarrow{\sim}X\times(X\times X)~;
\]
si on identifie canoniquement $X\times X\times X$ à $(X\times X)\times X$,
l'application $\gamma_X(Y)\times 1_{h_X(Y)}$ s'identifie à
$h_{c_X\times 1_X}(Y)$ pour tout $Y\in C$. Il est par suite équivalent de
dire que pour tout $Y\in C$, la loi interne $\gamma_X(Y)$ est
\emph{associative}, ou que le diagramme d'applications
\[
  \xymatrix{
    h_X(Y)\times h_X(Y)\times h_X(Y)
      \ar[rr]^{\gamma_X(Y)\times 1}\ar[dd]_{1\times\gamma_X(Y)} &&
    h_X(Y)\times h_X(Y)\ar[dd]^{\gamma_X(Y)}\\\\
    h_X(Y)\times h_X(Y)\ar[rr]_{\gamma_X(Y)} && h_X(Y)
  }
\]
\oldpage[0\textsubscript{III}]{11}
est commutatif, ou que le diagramme de morphismes
\[
  \xymatrix{
    X\times X\times X
      \ar[rr]^{c_X\times 1_X}\ar[dd]_{1_X\times c_X} &&
    X\times X\ar[dd]^{c_X}\\\\
    X\times X\ar[rr]_{c_X} && X
  }
\]
est commutatif.
\end{env}

\begin{env}[8.2.6]
\label{0.8.2.6-fr}
Sous les hypothèses de (\hyperref[0.8.2.5-fr]{8.2.5}), si on veut exprimer
que pour tout $Y\in C$, la loi interne $\gamma_X(Y)$ est une \emph{loi de
groupe}, il faut d'une part exprimer qu'elle est associative, et de l'autre
qu'il existe une application $\alpha_X(Y):h_X(Y)\to h_X(Y)$ ayant les
propriétés de l'\emph{inverse} dans un groupe~; comme pour tout morphisme
$u:Y\to Y'$ dans $C$, on a vu que $h_X(u)$ doit être un homomorphisme de
groupes $h_X(Y')\to h_X(Y)$, on voit d'abord que $\alpha_X:h_X\to h_X$ doit
être un \emph{morphisme fonctoriel}. On peut d'autre part exprimer les
propriétés caractéristiques de l'inverse $s\mapsto s^{-1}$ dans un groupe
$G$ sans faire intervenir l'élément neutre~: il suffit d'écrire que les deux
applications composées
\[
  (s,t)\mapsto(s,s^{-1},t)\mapsto(s,s^{-1}t)\mapsto s(s^{-1}t),
\]
\[
  (s,t)\mapsto(s,s^{-1},t)\mapsto(s,ts^{-1})\mapsto(ts^{-1})s
\]
sont égales à la seconde projection $(s,t)\mapsto t$ de $G\times G$ dans
$G$. En vertu de (\hyperref[0.8.1.3-fr]{8.1.3}), on a
$\alpha_X=h_{a_X}$, où $a_X\in\operatorname{Hom}(X,X)$~; la première
condition précédente exprime alors que le morphisme composé
\[
  X\times X\xrightarrow{(1_X,a_X)\times 1_X}X\times X\times X
  \xrightarrow{1_X\times c_X}X\times X\xrightarrow{c_X}X
\]
est la seconde projection $X\times X\to X$ dans $C$, et la seconde condition
se traduit de même.
\end{env}

\begin{env}[8.2.7]
\label{0.8.2.7-fr}
Supposons maintenant qu'il existe dans $C$ un \emph{objet final} $e$
(\hyperref[0.8.1.10-fr]{8.1.10}). Supposons toujours que $\gamma_X(Y)$ soit
une loi de groupe sur $h_X(Y)$ pour tout $Y\in C$, et désignons par
$\eta_X(Y)$ l'élément neutre de $\gamma_X(Y)$. Comme, pour tout morphisme
$u:Y\to Y'$ dans $C$, $h_X(u)$ est un homomorphisme de groupes, on a
$\eta_X(Y)=(h_X(u))(\eta_X(Y'))$~; prenant en particulier $Y'=e$, auquel cas
$u$ est l'unique élément $\varepsilon$ de $\operatorname{Hom}(Y,e)$, on voit
que l'élément $\eta_X(e)$ détermine complètement $\eta_X(Y)$ pour tout
$Y\in C$. Posons $e_X=\eta_X(X)$, élément neutre du groupe
$h_X(X)=\operatorname{Hom}(X,X)$~; la commutativité du diagramme
\[
  \xymatrix{
    h_X(e)\ar[r]^{h_X(\varepsilon)}\ar[d]_{h_{e_X}(e)} &
    h_X(Y)\ar[d]^{h_{e_X}(Y)}\\
    h_X(e)\ar[r]_{h_X(\varepsilon)} & h_X(Y)
  }
\]
(cf. \hyperref[0.8.1.2-fr]{8.1.2}) montre que, dans l'ensemble $h_X(Y)$,
l'application $h_{e_X}(Y)$ n'est autre que
\oldpage[0\textsubscript{III}]{12}
$s\mapsto\eta_X(Y)$ transformant tout élément en l'élément neutre. On
vérifie alors que le fait que $\eta_X(Y)$ soit élément neutre de
$\gamma_X(Y)$ pour tout $Y\in C$ équivaut à dire que le morphisme composé
\[
  X\xrightarrow{(1_X,1_X)}X\times X
  \xrightarrow{1_X\times e_X}X\times X\xrightarrow{c_X}X,
\]
et l'analogue où on permute $1_X$ et $e_X$, sont tous deux \emph{égaux} à
$1_X$.
\end{env}

\begin{env}[8.2.8]
\label{0.8.2.8-fr}
On pourrait bien entendu multiplier sans peine les exemples de structures
algébriques dans les catégories. L'exemple des groupes a été traité avec
assez de détails, mais par la suite nous laisserons généralement au lecteur
le soin de développer des considérations analogues dans les exemples de
structures algébriques que nous rencontrerons.
\end{env}

\section{Ensembles constructibles}
\label{section:0.9-fr}

\subsection{Ensembles constructibles}
\label{subsection:0.9.1-fr}

\begin{definition}[9.1.1]
\label{0.9.1.1-fr}
On dit qu'une application continue $f:X\to Y$ est \emph{quasi-compacte} si
pour tout ouvert quasi-compact $U$ de $Y$, $f^{-1}(U)$ est quasi-compact. On
dit qu'une partie $Z$ d'un espace topologique $X$ est \emph{rétrocompacte
dans $X$} si l'injection canonique $Z\to X$ est quasi-compacte, autrement dit
si pour tout ouvert quasi-compact $U$ de $X$, $U\cap Z$ est quasi-compact.
\end{definition}

Une partie \emph{fermée} de $X$ est rétrocompacte dans $X$, mais une partie
quasi-compacte de $X$ n'est pas nécessairement rétrocompacte dans $X$. Si
$X$ est quasi-compact, toute partie ouverte rétrocompacte dans $X$ est
quasi-compacte. Il est clair que toute réunion \emph{finie} d'ensembles
rétrocompacts dans $X$ est rétrocompacte dans $X$, toute réunion finie
d'ensembles quasi-compacts étant quasi-compacte. Toute intersection finie
d'ouverts rétrocompacts dans $X$ est un ouvert rétrocompact dans $X$. Dans un
espace \emph{localement noethérien} $X$, tout ensemble quasi-compact est un
sous-espace noethérien, et par suite \emph{toute partie} de $X$ est
rétrocompacte dans $X$.

\begin{definition}[9.1.2]
\label{0.9.1.2-fr}
Étant donné un espace topologique $X$, on dit qu'une partie de $X$ est
\emph{constructible} si elle appartient au plus petit ensemble de parties
$\mathfrak{F}$ de $X$ contenant toutes les parties ouvertes rétrocompactes de
$\mathfrak{F}$ et stable par intersection finie et passage au complémentaire
(ce qui implique que $\mathfrak{F}$ est aussi stable par réunion finie).
\end{definition}

\begin{proposition}[9.1.3]
\label{0.9.1.3-fr}
Pour qu'une partie de $X$ soit constructible, il faut et il suffit qu'elle
soit réunion finie d'ensembles de la forme $U\cap\complement{V}$, où $U$ et
$V$ sont des ouverts rétrocompacts dans $X$.
\end{proposition}

Il est clair que la condition est suffisante. Pour voir qu'elle est
nécessaire, considérons l'ensemble $\mathfrak{G}$ des réunions finies
d'ensembles de la forme $U\cap\complement{V}$ où $U$ et $V$ sont ouverts
rétrocompacts dans $X$~; il suffit de voir que tout complémentaire d'un
ensemble de $\mathfrak{G}$ appartient à $\mathfrak{G}$. Soit donc
$Z=\bigcup_{i\in I}(U_i\cap\complement{V_i})$, où $I$ est fini, $U_i$ et
$V_i$ ouverts rétrocompacts dans $X$~; on a
$\complement{Z}=\bigcap_{i\in I}(V_i\cup\complement{U_i})$, donc
$\complement{Z}$ est réunion finie d'ensembles qui sont intersections d'un
certain nombre des $V_i$ et d'un certain nombre de $\complement{U_i}$, donc
de la forme
\oldpage[0\textsubscript{III}]{13}
$V\cap\complement{U}$, où $U$ est réunion d'un certain nombre des $U_i$ et
$V$ intersection d'un certain nombre des $V_i$~; mais on a remarqué plus
haut que les réunions et intersections finies d'ouverts rétrocompacts dans
$X$ sont des ouverts rétrocompacts dans $X$, d'où la conclusion.

\begin{corollary}[9.1.4]
\label{0.9.1.4-fr}
Toute partie constructible de $X$ est rétrocompacte dans $X$.
\end{corollary}

Il suffit de montrer que si $U$ et $V$ sont ouverts rétrocompacts dans $X$,
$U\cap\complement{V}$ est rétrocompact dans $X$~; or, si $W$ est ouvert
quasi-compact dans $X$, $W\cap U\cap\complement{V}$ est fermé dans l'espace
quasi-compact $W\cap U$, donc est quasi-compact.

En particulier~:

\begin{corollary}[9.1.5]
\label{0.9.1.5-fr}
Pour qu'une partie ouverte $U$ de $X$ soit constructible, il faut et il
suffit qu'elle soit rétrocompacte dans $X$. Pour qu'une partie fermée $F$ de
$X$ soit constructible, il faut et il suffit que l'ouvert $\complement{F}$
soit rétrocompact.
\end{corollary}

\begin{env}[9.1.6]
\label{0.9.1.6-fr}
Un cas important est celui où toute partie ouverte quasi-compacte de $X$ est
rétrocompacte, autrement dit, où l'intersection de deux parties ouvertes
quasi-compactes de $X$ est quasi-compacte
(cf. \hyperref[I.5.5.6-fr]{I, 5.5.6}). Lorsque $X$ lui-même est
quasi-compact, cela signifie que les parties ouvertes rétrocompactes dans
$X$ sont identiques aux parties ouvertes quasi-compactes de $X$, et les
parties constructibles de $X$ aux réunions finies d'ensembles de la forme
$U\cap\complement{V}$, où $U$ et $V$ sont ouverts quasi-compacts.
\end{env}

\begin{corollary}[9.1.7]
\label{0.9.1.7-fr}
Pour qu'une partie d'un espace noethérien $X$ soit constructible, il faut et
il suffit qu'elle soit réunion finie de parties localement fermées de $X$.
\end{corollary}

\begin{proposition}[9.1.8]
\label{0.9.1.8-fr}
Soient $X$ un espace topologique, $U$ une partie ouverte de $X$.
\begin{enumerate}
  \item[(i)] Si $T$ est une partie constructible de $X$, $T\cap U$ est une
    partie constructible de $U$.
  \item[(ii)] Supposons en outre $U$ rétrocompact dans $X$. Pour qu'une
    partie $Z$ de $U$ soit constructible dans $X$, il faut et il suffit
    qu'elle soit constructible dans $U$.
\end{enumerate}
\end{proposition}

\begin{enumerate}
  \item[(i)] Utilisant (\hyperref[0.9.1.3-fr]{9.1.3}), on est ramené à
    montrer que si $T$ est ouvert rétrocompact dans $X$, $T\cap U$ est ouvert
    rétrocompact dans $U$, autrement dit, pour tout ouvert quasi-compact
    $W\subset U$, $T\cap U\cap W=T\cap W$ est quasi-compact, ce qui résulte
    aussitôt de l'hypothèse.
  \item[(ii)] La condition étant nécessaire en vertu de (i), il reste à
    démontrer qu'elle est suffisante. Compte tenu de
    (\hyperref[0.9.1.3-fr]{9.1.3}), il suffit de considérer le cas où $Z$ est
    ouvert rétrocompact \emph{dans $U$}, car il s'ensuivra alors que
    $U\setminus Z$ est constructible dans $X$, et si $Z$, $Z'$ sont deux
    ouverts rétrocompacts \emph{dans $U$},
    $Z\cap(U\setminus Z')$ sera bien constructible dans $X$. Or, si $W$ est
    ouvert quasi-compact dans $X$ et $Z$ ouvert rétrocompact dans $U$, on a
    $Z\cap W=Z\cap(W\cap U)$ et par hypothèse $W\cap U$ est ouvert
    quasi-compact dans $U$~; donc $W\cap Z$ est bien quasi-compact, et par
    suite $Z$ est ouvert rétrocompact dans $X$, et \emph{a fortiori}
    constructible dans $X$.
\end{enumerate}

\begin{corollary}[9.1.9]
\label{0.9.1.9-fr}
Soient $X$ un espace topologique, $(U_i)_{i\in I}$ un recouvrement fini de
$X$ formé d'ensembles ouverts rétrocompacts dans $X$. Pour qu'une partie $Z$
de $X$ soit constructible dans $X$, il faut et il suffit que pour tout
$i\in I$, $Z\cap U_i$ soit constructible dans $U_i$.
\end{corollary}

\begin{env}[9.1.10]
\label{0.9.1.10-fr}
Supposons en particulier que $X$ soit quasi-compact et que tout point
\oldpage[0\textsubscript{III}]{14}
de $X$ admette un système fondamental de voisinages ouverts rétrocompacts
dans $X$ (et \emph{a fortiori} quasi-compacts)~; alors la condition pour une
partie $Z$ de $X$ d'être constructible dans $X$ est de nature \emph{locale},
autrement dit, il faut et il suffit que pour tout $x\in X$, il existe un
voisinage ouvert $V$ de $x$ tel que $V\cap Z$ soit constructible dans $V$. En
effet, si cette condition est vérifiée, il existe pour tout $x\in X$ un
voisinage ouvert $V$ de $x$ \emph{rétrocompact dans $X$} et tel que
$V\cap Z$ soit constructible dans $V$, en vertu de l'hypothèse sur $X$ et de
(\hyperref[0.9.1.8-fr]{9.1.8, (i)})~; il suffit alors de recouvrir $X$ par un
nombre fini de ces voisinages, et d'appliquer
(\hyperref[0.9.1.9-fr]{9.1.9}).
\end{env}

\begin{definition}[9.1.11]
\label{0.9.1.11-fr}
Soit $X$ un espace topologique. On dit qu'une partie $T$ de $X$ est
\emph{localement constructible dans $X$} si, pour tout $x\in X$, il existe
un voisinage ouvert $V$ de $x$ tel que $T\cap V$ soit constructible dans $V$.
\end{definition}

Il résulte aussitôt de (\hyperref[0.9.1.8-fr]{9.1.8, (i)}) que si $V$ est
tel que $V\cap T$ soit constructible dans $V$, alors pour tout ouvert
$W\subset V$, $W\cap T$ est constructible dans $W$. Si $T$ est localement
constructible dans $X$, alors, pour tout ouvert $U$ de $X$, $T\cap U$ est
localement constructible dans $U$, comme il résulte de la remarque
précédente. Cette même remarque montre que l'ensemble des parties localement
constructibles dans $X$ est stable par réunion finie et intersection finie~;
il est clair, d'autre part, qu'il est aussi stable par passage aux
complémentaires.

\begin{proposition}[9.1.12]
\label{0.9.1.12-fr}
Soit $X$ un espace topologique. Tout ensemble constructible dans $X$ est
localement constructible dans $X$. La réciproque est vraie si $X$ est
quasi-compact et si sa topologie admet une base formée d'ensembles
rétrocompacts dans $X$.
\end{proposition}

La première assertion résulte de la définition
(\hyperref[0.9.1.11-fr]{9.1.11}) et la seconde de
(\hyperref[0.9.1.10-fr]{9.1.10}).

\begin{corollary}[9.1.13]
\label{0.9.1.13-fr}
Soit $X$ un espace topologique dont la topologie admet une base formée
d'ensembles rétrocompacts dans $X$. Alors toute partie $T$ localement
constructible dans $X$ est rétrocompacte dans $X$.
\end{corollary}

En effet, soit $U$ un ensemble ouvert quasi-compact dans $X$~; $T\cap U$ est
localement constructible dans $U$, donc constructible dans $U$ en vertu de
(\hyperref[0.9.1.12-fr]{9.1.12}), et par suite quasi-compact en vertu de
(\hyperref[0.9.1.4-fr]{9.1.4}).

\subsection{Ensembles constructibles dans les espaces noethériens}
\label{subsection:0.9.2-fr}

\begin{env}[9.2.1]
\label{0.9.2.1-fr}
On a vu (\hyperref[0.9.1.7-fr]{9.1.7}) que dans un espace noethérien $X$, les
parties constructibles dans $X$ sont les \emph{réunions finies de parties
localement fermées de $X$}.

L'image réciproque d'un ensemble constructible dans $X$ par une application
continue d'un espace noethérien $X'$ dans $X$ est constructible dans $X'$.
Si $Y$ est une partie constructible d'un espace noethérien $X$, les parties
de $Y$ qui sont constructibles en tant que sous-espaces de $Y$ sont
identiques à celles qui sont constructibles en tant que sous-espaces de $X$.
\end{env}

\begin{proposition}[9.2.2]
\label{0.9.2.2-fr}
Soient $X$ un espace irréductible noethérien, $E$ une partie constructible de
$X$. Pour que $E$ soit partout dense dans $X$, il faut et il suffit que $E$
contienne une partie ouverte non vide de $X$.
\end{proposition}

\oldpage[0\textsubscript{III}]{15}
La condition est évidemment suffisante, tout ensemble ouvert non vide étant
dense dans $X$. Inversement, soit
$E=\bigcup_{i=1}^n(U_i\cap F_i)$ une partie constructible de $X$, les $U_i$
étant ouverts non vides et les $F_i$ fermés dans $X$~; on a donc
$\overline{E}\subset\bigcup_i F_i$. Par suite, si $\overline{E}=X$, $X$ est
égal à l'un des $F_i$, donc $E\supset U_i$, ce qui achève la démonstration.

Lorsque $X$ admet un point générique $x$
(\hyperref[0.2.1.2-fr]{$0_{\mathrm I}$, 2.1.2}), la condition de
(\hyperref[0.9.2.2-fr]{9.2.2}) équivaut à la relation $x\in E$.

\begin{proposition}[9.2.3]
\label{0.9.2.3-fr}
Soit $X$ un espace noethérien. Pour qu'une partie $E$ de $X$ soit
constructible, il faut et il suffit que, pour toute partie fermée
irréductible $Y$ de $X$, $E\cap Y$ soit rare dans $Y$ ou contienne une partie
ouverte non vide de $Y$.
\end{proposition}

La nécessité de la condition provient de ce que $E\cap Y$ doit être une
partie constructible de $Y$ et de (\hyperref[0.9.2.2-fr]{9.2.2}), car une
partie non dense de $Y$ est nécessairement rare dans l'espace irréductible
$Y$ (\hyperref[0.2.1.1-fr]{$0_{\mathrm I}$, 2.1.1}). Pour prouver que la
condition est suffisante, appliquons le principe de récurrence noethérienne
(\hyperref[0.2.2.2-fr]{$0_{\mathrm I}$, 2.2.2}) à l'ensemble
$\mathfrak{F}$ des parties fermées $Y$ de $X$ telles que $Y\cap E$ soit
constructible (par rapport à $Y$ ou par rapport à $X$, ce qui revient au
même)~: on peut donc supposer que pour toute partie fermée $Y\ne X$ de $X$,
$E\cap Y$ est constructible. Supposons d'abord que $X$ ne soit pas
irréductible, et soient $X_i$ ($1\leq i\leq m$) ses composantes
irréductibles, nécessairement en nombre fini
(\hyperref[0.2.2.5-fr]{$0_{\mathrm I}$, 2.2.5})~; par hypothèse, les
$E\cap X_i$ sont constructibles, donc aussi leur réunion $E$. Supposons
ensuite que $X$ soit irréductible~; alors, par hypothèse, ou bien $E$ est
rare, donc $\overline{E}\ne X$ et $E=E\cap\overline{E}$ est constructible~;
ou bien $E$ contient un ouvert non vide $U$ de $E$, donc est réunion de $U$
et de $E\cap(X\setminus U)$~; mais $X\setminus U$ est un ensemble fermé
distinct de $X$, donc $E\cap(X\setminus U)$ est constructible~; $E$
lui-même est par suite constructible, ce qui achève la démonstration.

\begin{corollary}[9.2.4]
\label{0.9.2.4-fr}
Soient $X$ un espace noethérien, $(E_\alpha)$ une famille filtrante
croissante de parties constructibles de $X$, telle que~:
\begin{enumerate}
  \item[$1^\circ$] $X$ est réunion de la famille $(E_\alpha)$.
  \item[$2^\circ$] Toute partie fermée irréductible de $X$ est contenue dans
    l'adhérence d'un des $E_\alpha$.
\end{enumerate}
Alors il existe un indice $\alpha$ tel que $X=E_\alpha$.

Lorsque toute partie fermée irréductible de $X$ admet un point générique,
l'hypothèse $2^\circ$ peut être supprimée.
\end{corollary}

Appliquons le principe de récurrence noethérienne
(\hyperref[0.2.2.2-fr]{$0_{\mathrm I}$, 2.2.2}) à l'ensemble
$\mathfrak{M}$ des parties fermées de $X$ contenues dans l'un des $E_\alpha$
au moins~; on peut donc supposer que toute partie fermée $Y\ne X$ de $X$ est
contenue dans un des $E_\alpha$. La proposition est évidente si $X$ n'est pas
irréductible, car chacune des composantes irréductibles $X_i$ de $X$
($1\leq i\leq m$) est contenue dans un $E_{\alpha_i}$, et il existe un
$E_\alpha$ contenant tous les $E_{\alpha_i}$. Supposons donc $X$
irréductible. Par hypothèse, il existe $\beta$ tel que
$X=\overline{E_\beta}$, donc (\hyperref[0.9.2.2-fr]{9.2.2}) $E_\beta$
contient un ouvert non vide $U$ de $X$. Mais alors l'ensemble fermé
$X\setminus U$ est contenu dans un $E_\gamma$, et il suffit de prendre
$E_\alpha$ contenant $E_\beta$ et $E_\gamma$. Lorsque toute partie fermée
irréductible $Y$ de $X$
\oldpage[0\textsubscript{III}]{16}
admet un point générique $y$, il existe $\alpha$ tel que $y\in E_\alpha$,
donc $Y=\overline{\{y\}}\subset\overline{E_\alpha}$, et la condition
$2^\circ$ est conséquence de $1^\circ$.

\begin{proposition}[9.2.5]
\label{0.9.2.5-fr}
Soient $X$ un espace noethérien, $x$ un point de $X$, $E$ une partie
constructible de $X$. Pour que $E$ soit un voisinage de $x$, il faut et il
suffit que pour toute partie fermée irréductible $Y$ de $X$ contenant $x$,
$E\cap Y$ soit dense dans $Y$ (s'il existe un point générique $y$ de $Y$,
cela signifie aussi (\hyperref[0.9.2.2-fr]{9.2.2}) que $y\in E$).
\end{proposition}

La condition est évidemment nécessaire~; prouvons qu'elle est suffisante.
Appliquant le principe de récurrence noethérienne à l'ensemble $\mathfrak{M}$
des parties fermées $Y$ de $X$ contenant $x$ et telles que $E\cap Y$ soit un
voisinage de $x$ dans $Y$, on peut supposer que toute partie fermée $Y\ne X$
de $X$ contenant $x$ appartient à $\mathfrak{M}$. Si $X$ n'est pas
irréductible, chacune des composantes irréductibles $X_i$ de $X$ contenant
$x$ est distincte de $X$, donc $E\cap X_i$ est un voisinage de $x$ par
rapport à $X_i$~; par suite, $E$ est un voisinage de $x$ dans la réunion des
composantes irréductibles de $X$ contenant $x$, et comme cette réunion est un
voisinage de $x$ dans $X$, il en est de même de $E$. Si $X$ est
irréductible, $E$ est dense dans $X$ par hypothèse, donc contient une partie
ouverte non vide $U$ de $X$ (\hyperref[0.9.2.2-fr]{9.2.2})~; la proposition
est alors évidente si $x\in U$~; sinon, $x$ est par hypothèse intérieur à
$E\cap(X\setminus U)$ par rapport à $X\setminus U$, donc l'adhérence dans
$X$ de $X\setminus E$ ne contient pas $x$, et le complémentaire de cette
adhérence est un voisinage de $x$ contenu dans $E$, ce qui achève la
démonstration.

\begin{corollary}[9.2.6]
\label{0.9.2.6-fr}
Soient $X$ un espace noethérien, $E$ une partie de $X$. Pour que $E$ soit un
ensemble ouvert dans $X$, il faut et il suffit que pour toute partie fermée
irréductible $Y$ de $X$ rencontrant $E$, $E\cap Y$ contienne une partie
ouverte non vide de $Y$.
\end{corollary}

La condition est évidemment nécessaire~; inversement, si elle est vérifiée,
elle implique que $E$ est constructible en vertu de
(\hyperref[0.9.2.3-fr]{9.2.3}). En outre,
(\hyperref[0.9.2.5-fr]{9.2.5}) montre que $E$ est alors voisinage de chacun
de ses points, d'où la conclusion.

\subsection{Fonctions constructibles}
\label{subsection:0.9.3-fr}

\begin{definition}[9.3.1]
\label{0.9.3.1-fr}
Soit $h$ une application d'un espace topologique $X$ dans un ensemble $T$.
On dit que $h$ est \emph{constructible} si $h^{-1}(t)$ est constructible pour
tout $t\in T$, et vide sauf pour un nombre fini de valeurs de $t$~; pour
toute partie $S$ de $T$, $h^{-1}(S)$ est alors constructible. On dit que $h$
est \emph{localement constructible} si tout $x\in X$ possède un voisinage
ouvert $V$ tel que $h|V$ soit constructible.
\end{definition}

Toute fonction constructible est localement constructible~; la réciproque
est vraie quand $X$ est quasi-compact et admet une base formée d'ensembles
ouverts rétrocompacts dans $X$ (en particulier quand $X$ est noethérien).

\begin{proposition}[9.3.2]
\label{0.9.3.2-fr}
Soit $h$ une application d'un espace noethérien $X$ dans un ensemble $T$.
Pour que $h$ soit constructible, il faut et il suffit que pour toute partie
fermée irréductible $Y$ de $X$, il existe une partie non vide $U$ de $Y$,
ouverte par rapport à $Y$, et dans laquelle $h$ soit constante.
\end{proposition}

La condition est nécessaire~: en effet, par hypothèse, $h$ ne prend dans $Y$
qu'un nombre fini de valeurs $t_i$, et chacun des ensembles
$h^{-1}(t_i)\cap Y$ est constructible dans $Y$
(\hyperref[0.9.2.1-fr]{9.2.1})~; comme ils ne peuvent être tous des parties
rares de l'espace $Y$, un d'eux au moins contient un ensemble ouvert non vide
(\hyperref[0.9.2.3-fr]{9.2.3}).

\oldpage[0\textsubscript{III}]{17}
Pour voir que la condition est suffisante, appliquons le principe de
récurrence noethérienne à l'ensemble $\mathfrak{M}$ des parties fermées $Y$
de $X$ telles que la restriction $h|Y$ soit constructible~; on peut donc
supposer que pour toute partie fermée $Y\ne X$ de $X$, $h|Y$ est
constructible. Si $X$ n'est pas irréductible, la restriction de $h$ à
chacune des composantes irréductibles $X_i$ de $X$ (en nombre fini) est donc
constructible, et il résulte alors aussitôt de la définition
(\hyperref[0.9.3.1-fr]{9.3.1}) que $h$ est constructible. Si $X$ est
irréductible, il existe par hypothèse une partie ouverte non vide $U$ de $X$
dans laquelle $h$ est constante~; d'autre part, la restriction de $h$ à
$X\setminus U$ est constructible par hypothèse, et il en résulte aussitôt
que $h$ est constructible.

\begin{corollary}[9.3.3]
\label{0.9.3.3-fr}
Soit $X$ un espace noethérien dans lequel toute partie fermée irréductible
admet un point générique. Si $h$ est une application de $X$ dans un ensemble
$T$ telle que, pour tout $t\in T$, $h^{-1}(t)$ soit constructible, alors $h$
est constructible.
\end{corollary}

En effet, si $Y$ est une partie fermée irréductible de $X$ et $y$ son point
générique, $Y\cap h^{-1}(h(y))$ est constructible et contient $y$, donc
(\hyperref[0.9.2.2-fr]{9.2.2}) cet ensemble contient une partie ouverte non
vide de $Y$, et il suffit d'appliquer (\hyperref[0.9.3.2-fr]{9.3.2}).

\begin{proposition}[9.3.4]
\label{0.9.3.4-fr}
Soient $X$ un espace noethérien dans lequel toute partie fermée irréductible
admet un point générique, $h$ une application constructible de $X$ dans un
ensemble ordonné. Pour que $h$ soit semi-continue supérieurement dans $X$,
il faut et il suffit que pour tout $x\in X$ et toute générisation
(\hyperref[0.2.1.2-fr]{$0_{\mathrm I}$, 2.1.2}) $x'$ de $x$, on ait
$h(x')\leq h(x)$.
\end{proposition}

La fonction $h$ ne prend qu'un nombre fini de valeurs~; dire qu'elle est
semi-continue supérieurement signifie donc que pour tout $x\in X$,
l'ensemble $E$ des $y\in X$ tels que $h(y)\leq h(x)$ est un voisinage de $x$.
Par hypothèse, $E$ est une partie constructible de $X$~; d'autre part, dire
qu'une partie fermée irréductible $Y$ de $X$ contient $x$ signifie que son
point générique $y$ est une générisation de $x$~; la conclusion résulte
alors de (\hyperref[0.9.2.5-fr]{9.2.5}).

\section{Compléments sur les modules plats}
\label{section:0.10-fr}

Pour les démonstrations des propriétés énoncées sans démonstration dans les
n\textsuperscript{os} (\hyperref[subsection:0.10.1-fr]{10.1}) et
(\hyperref[subsection:0.10.2-fr]{10.2}), nous renvoyons le lecteur à Bourbaki,
\emph{Alg. comm.}, chap. II et III.

\subsection{Relations entre modules plats et modules libres}
\label{subsection:0.10.1-fr}

\begin{env}[10.1.1]
\label{0.10.1.1-fr}
Soient $A$ un anneau, $\mathfrak{J}$ un idéal de $A$, $M$ un $A$-module~; pour
tout entier $p\geq0$, on a un homomorphisme canonique de
$(A/\mathfrak{J})$-modules
\[
\label{0.10.1.1.1-fr}
  \varphi_p:(M/\mathfrak{J}M)\otimes_{A/\mathfrak{J}}
  (\mathfrak{J}^p/\mathfrak{J}^{p+1})
  \longrightarrow \mathfrak{J}^pM/\mathfrak{J}^{p+1}M
  \tag{10.1.1.1}
\]
qui est évidemment \emph{surjectif}. Nous désignerons par
$\operatorname{gr}(A)=\bigoplus_{p\geq0}\mathfrak{J}^p/\mathfrak{J}^{p+1}$ l'anneau gradué
associé à $A$ filtré par les $\mathfrak{J}^p$, par
$\operatorname{gr}(M)=\bigoplus_{p\geq0}\mathfrak{J}^pM/\mathfrak{J}^{p+1}M$ le
$\operatorname{gr}(A)$-module gradué associé à $M$ filtré par les $\mathfrak{J}^pM$~; on a
donc $\operatorname{gr}_p(A)=\mathfrak{J}^p/\mathfrak{J}^{p+1}$,
$\operatorname{gr}_p(M)=\mathfrak{J}^pM/\mathfrak{J}^{p+1}M$~; les $\varphi_p$ définissent
un homomorphisme \emph{surjectif} de $\operatorname{gr}(A)$-modules gradués
\[
\label{0.10.1.1.2-fr}
  \varphi:\operatorname{gr}_0(M)\otimes_{\operatorname{gr}_0(A)}
  \operatorname{gr}(A)\longrightarrow\operatorname{gr}(M).
  \tag{10.1.1.2}
\]
\end{env}

\oldpage[0\textsubscript{III}]{18}
\begin{env}[10.1.2]
\label{0.10.1.2-fr}
Supposons vérifiée l'une des hypothèses suivantes~:
\begin{enumerate}
  \item[(i)] $\mathfrak{J}$ est nilpotent~;
  \item[(ii)] $A$ est noethérien, $\mathfrak{J}$ est contenu dans le radical de
    $A$, et $M$ est de type fini.
\end{enumerate}
Alors les propriétés suivantes sont équivalentes~:
\begin{enumerate}
  \item[$a)$] $M$ est un $A$-module libre.
  \item[$b)$] $M/\mathfrak{J}M=M\otimes_A(A/\mathfrak{J})$ est un
    $(A/\mathfrak{J})$-module libre, et
    $\operatorname{Tor}_1^A(M,A/\mathfrak{J})=0$.
  \item[$c)$] $M/\mathfrak{J}M$ est un $(A/\mathfrak{J})$-module libre et
    l'homomorphisme canonique (\hyperref[0.10.1.1.2-fr]{10.1.1.2}) est
    injectif (donc bijectif).
\end{enumerate}
\end{env}

\begin{env}[10.1.3]
\label{0.10.1.3-fr}
Supposons que $A/\mathfrak{J}$ soit un \emph{corps} (autrement dit que
$\mathfrak{J}$ soit maximal), et que l'une des hypothèses (i), (ii) de
(\hyperref[0.10.1.2-fr]{10.1.2}) soit vérifiée. Alors les propriétés suivantes
sont équivalentes~:
\begin{enumerate}
  \item[$a)$] $M$ est un $A$-module libre.
  \item[$b)$] $M$ est un $A$-module projectif.
  \item[$c)$] $M$ est un $A$-module plat.
  \item[$d)$] $\operatorname{Tor}_1^A(M,A/\mathfrak{J})=0$.
  \item[$e)$] L'homomorphisme canonique
    (\hyperref[0.10.1.1.2-fr]{10.1.1.2}) est bijectif.
\end{enumerate}
Ce résultat s'appliquera en particulier dans les deux cas suivants~:
\begin{enumerate}
  \item[(i)] $M$ est un module \emph{quelconque} sur un anneau local $A$ dont
    l'idéal maximal $\mathfrak{J}$ est \emph{nilpotent} (par exemple un anneau
    local artinien).
  \item[(ii)] $M$ est un module \emph{de type fini} sur un anneau local
    \emph{noethérien}.
\end{enumerate}
\end{env}

\subsection{Critères locaux de platitude}
\label{subsection:0.10.2-fr}

\begin{env}[10.2.1]
\label{0.10.2.1-fr}
Les hypothèses et notations étant celles de
(\hyperref[0.10.1.1-fr]{10.1.1}), considérons les conditions suivantes~:
\begin{enumerate}
  \item[$a)$] $M$ est un $A$-module plat.
  \item[$b)$] $M/\mathfrak{J}M$ est un $(A/\mathfrak{J})$-module plat et
    $\operatorname{Tor}_1^A(M,A/\mathfrak{J})=0$.
  \item[$c)$] $M/\mathfrak{J}M$ est un $(A/\mathfrak{J})$-module plat et
    l'homomorphisme canonique (\hyperref[0.10.1.1.2-fr]{10.1.1.2}) est
    bijectif.
  \item[$d)$] Pour tout $n\geq1$, $M/\mathfrak{J}^nM$ est un
    $(A/\mathfrak{J}^n)$-module plat.
\end{enumerate}
On a alors les implications
\[
  a\Rightarrow b\Rightarrow c\Rightarrow d
\]
et si $\mathfrak{J}$ est \emph{nilpotent}, les quatre conditions $a)$, $b)$,
$c)$, $d)$ sont \emph{équivalentes}. Il en est de même si $A$ est noethérien
et en outre si $M$ est \emph{idéalement séparé}, c'est-à-dire que pour tout
idéal $\mathfrak{a}$ de $A$, le $A$-module $\mathfrak{a}\otimes_A M$ est
\emph{séparé} pour la topologie $\mathfrak{J}$-pré-adique.
\end{env}

\begin{env}[10.2.2]
\label{0.10.2.2-fr}
Soient $A$ un anneau noethérien, $B$ une $A$-algèbre commutative noethérienne,
$\mathfrak{J}$ un idéal de $A$ tel que $\mathfrak{J}B$ soit contenu dans le
radical de $B$, $M$ un $B$-module de type fini. Alors, lorsque $M$ est
considéré comme $A$-module, les conditions $a)$, $b)$, $c)$, $d)$ de
(\hyperref[0.10.2.1-fr]{10.2.1}) sont équivalentes.

Ce résultat s'applique surtout lorsque $A$ et $B$ sont des anneaux
\emph{locaux} noethériens, l'homomorphisme $A\to B$ un homomorphisme
\emph{local}. Plus particulièrement, si $\mathfrak{J}$ est alors l'idéal
\emph{maximal} de $A$, on peut, dans les conditions $b)$ et $c)$, supprimer
l'hypothèse que $M/\mathfrak{J}M$ est plat, qui est automatiquement vérifiée,
et la condition $d)$ signifie que les modules $M/\mathfrak{J}^nM$ sont
\emph{libres} sur les $A/\mathfrak{J}^n$.
\end{env}

\oldpage[0\textsubscript{III}]{19}
\begin{env}[10.2.3]
\label{0.10.2.3-fr}
Les hypothèses sur $A$, $B$, $\mathfrak{J}$, $M$ étant celles formulées au
début de (\hyperref[0.10.2.2-fr]{10.2.2}), soient $\widehat A$ le séparé
complété de $A$ pour la topologie $\mathfrak{J}$-pré-adique, $\widehat M$ le
séparé complété de $M$ pour la topologie $\mathfrak{J}B$-pré-adique. Alors,
pour que $M$ soit un $A$-module plat, il faut et il suffit que $\widehat M$
soit un $\widehat A$-module plat.
\end{env}

\begin{env}[10.2.4]
\label{0.10.2.4-fr}
Soient $\rho:A\to B$ un homomorphisme local d'anneaux locaux noethériens, $k$
le corps résiduel de $A$, $M$, $N$ deux $B$-modules de type fini, $N$ étant
supposé être $A$-\emph{plat}. Soit $u:M\to N$ un $B$-homomorphisme. Alors les
conditions suivantes sont équivalentes~:
\begin{enumerate}
  \item[$a)$] $u$ est injectif et $\operatorname{Coker}(u)$ est un $A$-module plat.
  \item[$b)$] $u\otimes1:M\otimes_A k\to N\otimes_A k$ est injectif.
\end{enumerate}
\end{env}

\begin{env}[10.2.5]
\label{0.10.2.5-fr}
Soient $\rho:A\to B$, $\sigma:B\to C$ des homomorphismes locaux d'anneaux
locaux noethériens, $k$ le corps résiduel de $A$, $M$ un $C$-module de type
fini. On suppose que $B$ est un $A$-module \emph{plat} Alors les conditions
suivantes sont équivalentes~:
\begin{enumerate}
  \item[$a)$] $M$ est un $B$-module plat.
  \item[$b)$] $M$ est un $A$-module plat, et $M\otimes_A k$ est un
    $(B\otimes_A k)$-module plat.
\end{enumerate}
\end{env}

\begin{proposition}[10.2.6]
\label{0.10.2.6-fr}
Soient $A$, $B$ deux anneaux locaux noethériens, $\rho:A\to B$ un
homomorphisme local, $\mathfrak{J}$ un idéal de $B$ contenu dans l'idéal
maximal, $M$ un $B$-module de type fini. Supposons que pour tout $n\geq0$,
$M_n=M/\mathfrak{J}^{n+1}M$ soit un $A$-module plat. Alors $M$ est un
$A$-module plat.
\end{proposition}

Il faut prouver que pour tout homomorphisme injectif $u:N'\to N$ de
$A$-modules de type fini, $v=1\otimes u:M\otimes_A N'\to M\otimes_A N$ est
injectif. Or, $M\otimes_A N'$ et $M\otimes_A N$ sont des $B$-modules de type
fini, donc séparés pour la topologie $\mathfrak{J}$-pré-adique
(\hyperref[0.7.3.5-fr]{$0_{\mathrm I}$, 7.3.5})~; il suffit donc de prouver que
l'homomorphisme
$\widehat v:\widehat{(M\otimes_A N')}\to\widehat{(M\otimes_A N)}$ pour les
séparés complétés est injectif. Or, on a $\widehat v=\varprojlim v_n$, où
$v_n$ est l'homomorphisme $1\otimes u:M_n\otimes_A N'\to M_n\otimes_A N$~;
comme par hypothèse $M_n$ est $A$-plat, $v_n$ est injectif pour tout $n$, donc
il en est de même de $v$, le foncteur $\varprojlim$ étant exact à gauche.

\begin{corollary}[10.2.7]
\label{0.10.2.7-fr}
Soient $A$ un anneau noethérien, $B$ un anneau local noethérien,
$\rho:A\to B$ un homomorphisme, $f$ un élément de l'idéal maximal de $B$, $M$
un $B$-module de type fini. Supposons que l'homothétie
$f_M:x\mapsto fx$ de $M$ soit injective et que $M/fM$ soit un $A$-module
plat. Alors $M$ est un $A$-module plat.
\end{corollary}

Posons $M_i=f^iM$ pour $i\geq0$~; comme $f_M$ est injective,
$M_i/M_{i+1}$ est isomorphe à $M/fM$, donc $A$-plat pour tout $i\geq0$~; de
la suite exacte
\[
  0\longrightarrow M_i/M_{i+1}\longrightarrow M/M_{i+1}
  \longrightarrow M/M_i\longrightarrow0
\]
on tire par récurrence sur $i$ que $M/M_i$ est $A$-\emph{plat} pour tout
$i\geq0$ (\hyperref[0.6.1.2-fr]{$0_{\mathrm I}$, 6.1.2})~; on peut donc
appliquer (\hyperref[0.10.2.6-fr]{10.2.6}). On peut aussi raisonner directement
comme suit~: pour tout $A$-module $N$ de type fini, $M\otimes_A N$ est un
$B$-module de type fini~; comme $f$ appartient au radical $\mathfrak n$ de
$B$, la topologie $(f)$-adique sur $M\otimes_A N$ est plus fine que la
\oldpage[0\textsubscript{III}]{20}
topologie $\mathfrak n$-adique, et on sait que cette dernière est \emph{séparée}
(\hyperref[0.7.3.5-fr]{$0_{\mathrm I}$, 7.3.5}). D'ailleurs, comme $M/M_i$ est
$A$-plat, on a
\[
  f^i(M\otimes_A N)=\operatorname{Im}(M_i\otimes_A N\to M\otimes_A N)
  =\operatorname{Ker}(M\otimes_A N\to(M/M_i)\otimes_A N)
\]
(\hyperref[0.6.1.2-fr]{$0_{\mathrm I}$, 6.1.2}). Soient alors $N$ un
$A$-module de type fini, $N'$ un sous-module de $N$, $j:N'\to N$ l'injection
canonique~; dans le diagramme commutatif
\[
\xymatrix{
  M\otimes_A N'\ar[r]\ar[d]_{1_M\otimes j} &
  (M/M_i)\otimes_A N'\ar[d]^{1_{M/M_i}\otimes j}\\
  M\otimes_A N\ar[r] & (M/M_i)\otimes_A N
}
\]
$1_{M/M_i}\otimes j$ est injectif puisque $M/M_i$ est $A$-plat~; on en
conclut que
\[
  \operatorname{Ker}(M\otimes_A N'\to M\otimes_A N)
  \subset\operatorname{Ker}(M\otimes_A N'\to(M/M_i)\otimes_A N')
\]
quel que soit $i$~; puisque l'intersection des seconds membres est réduite à
$0$ comme on l'a vu plus haut, il en est de même du premier membre, et par
suite $M$ est $A$-plat.

\begin{proposition}[10.2.8]
\label{0.10.2.8-fr}
Soient $A$ un anneau noethérien réduit, $M$ un $A$-module de type fini. On
suppose que pour toute $A$-algèbre $B$, qui est un anneau de valuation
discrète, $M\otimes_A B$ soit un $B$-module plat (donc libre
(\hyperref[0.10.1.3-fr]{10.1.3})). Alors $M$ est un $A$-module plat.
\end{proposition}

On sait que pour que $M$ soit plat, il faut et il suffit que pour tout idéal
maximal $\mathfrak m$ de $A$, $M_{\mathfrak m}$ soit un
$A_{\mathfrak m}$-module plat
(\hyperref[0.6.3.3-fr]{$0_{\mathrm I}$, 6.3.3})~; on peut donc se borner au cas
où $A$ est \emph{local} (\hyperref[0.1.2.8-fr]{$0_{\mathrm I}$, 1.2.8}).
Soient alors $\mathfrak m$ l'idéal maximal de $A$, $\mathfrak p_i$
$(1\leq i\leq r)$ ses idéaux premiers minimaux, $k$ le corps résiduel
$A/\mathfrak m$. On sait (\hyperref[II.7.1.7-fr]{II, 7.1.7}) qu'il existe
pour chaque $i$ un anneau de valuation discrète $B_i$, ayant même corps des
fractions $K_i$ que l'anneau intègre $A/\mathfrak p_i$, et dominant ce
dernier. Posons $M_i=M\otimes_A B_i$. Par hypothèse, $M_i$ est libre sur
$B_i$, donc on a, en désignant par $k_i$ le corps résiduel de $B_i$
\[
\label{0.10.2.8.1-fr}
  \operatorname{rg}_{k_i}(M_i\otimes_{B_i}k_i)
  =\operatorname{rg}_{K_i}(M_i\otimes_{B_i}K_i).
  \tag{10.2.8.1}
\]
Mais il est clair que l'homomorphisme composé
$A\to A/\mathfrak p_i\to B_i$ est local, donc $k$ est une extension de $k_i$,
et l'on a
$M_i\otimes_{B_i}k_i=M\otimes_A k_i=(M\otimes_A k)\otimes_k k_i$, et par
ailleurs $M_i\otimes_{B_i}K_i=M\otimes_A K_i$. L'égalité
(\hyperref[0.10.2.8.1-fr]{10.2.8.1}) entraîne donc
\[
  \operatorname{rg}_k(M\otimes_A k)=\operatorname{rg}_{K_i}(M\otimes_A K_i)
  \qquad\text{pour }1\leq i\leq r
\]
et comme $A$ est réduit, on sait que cette condition entraîne que $M$ est un
$A$-module \emph{libre} (Bourbaki, \emph{Alg. comm.}, chap. II, \S~3,
n\textsuperscript{o}~2, prop. 7).

\subsection{Existence d'extensions plates d'anneaux locaux}
\label{subsection:0.10.3-fr}

\begin{proposition}[10.3.1]
\label{0.10.3.1-fr}
Soient $A$ un anneau local noethérien, $\mathfrak J$ son idéal maximal,
$k=A/\mathfrak J$ son corps résiduel. Soit $K$ une extension du corps $k$. Il
existe un homomorphisme local de $A$ dans un anneau local noethérien $B$, tel
que $B/\mathfrak J B$ soit $k$-isomorphe à $K$, et que $B$ soit un
$A$-module plat.
\end{proposition}

Nous démontrerons cette proposition en plusieurs étapes.

\begin{env}[10.3.1.1]
\label{0.10.3.1.1-fr}
Supposons d'abord que $K=k(T)$, où $T$ est une indéterminée. Dans l'anneau de
polynômes $A'=A[T]$, considérons l'idéal premier
$\mathfrak J'=\mathfrak J A'$ formé des
\oldpage[0\textsubscript{III}]{21}
polynômes ayant leurs coefficients dans l'idéal $\mathfrak J$~; il est clair
que $A'/\mathfrak J'$ est canoniquement isomorphe à $k[T]$. Montrons que
l'anneau de fractions $B=A'_{\mathfrak J'}$ répond à la question~; c'est
évidemment un anneau local noethérien dont l'idéal maximal
$\mathfrak L=\mathfrak J B$. En outre,
$B/\mathfrak L=(A'/\mathfrak J')_{\mathfrak J'}=(k[T])_{\mathfrak J'}$
n'est autre que le corps des fractions $K$ de $k[T]$. Enfin, $B$ est un
$A'$-module plat et $A'$ un $A$-module libre, donc $B$ est un $A$-module plat
(\hyperref[0.6.2.1-fr]{$0_{\mathrm I}$, 6.2.1}).
\end{env}

\begin{env}[10.3.1.2]
\label{0.10.3.1.2-fr}
Supposons ensuite que $K=k(t)=k[t]$, où $t$ est algébrique sur $k$~; soit
$f\in k[T]$ le polynôme minimal de $t$~; il existe un polynôme unitaire
$F\in A[T]$ dont l'image canonique dans $k[T]$ soit $f$. Posons encore
$A'=A[T]$, et soit $\mathfrak J'$ l'idéal $\mathfrak J A'+(F)$ dans $A'$.
Nous allons voir que l'anneau quotient $B=A'/(F)$ répond cette fois à la
question. Tout d'abord, il est clair que $B$ est un $A$-module \emph{libre},
donc plat. L'anneau $A'/\mathfrak J'$ est isomorphe à
\[
  (A'/\mathfrak J A')/((\mathfrak J A'+(F))/\mathfrak J A')
  =k[T]/(f)=K~;
\]
l'image $\mathfrak L$ de $\mathfrak J'$ dans $B$ est donc maximal et on a
évidemment $\mathfrak L=\mathfrak J B$. Enfin, $B$ est un anneau semi-local,
puisqu'il est un $A$-module de type fini (Bourbaki, \emph{Alg. comm.}, chap.
IV, \S~2, n\textsuperscript{o}~5, cor. 3 de la prop. 9), et ses idéaux
maximaux sont en correspondance biunivoque avec ceux de $B/\mathfrak J B$
([13], vol. I, p. 259)~; ce qui précède prouve donc que $B$ est un anneau
local.
\end{env}

\begin{lemma}[10.3.1.3]
\label{0.10.3.1.3-fr}
Soit $(A_\lambda,f_{\mu\lambda})$ un système inductif filtrant d'anneaux
locaux, tel que les $f_{\mu\lambda}$ soient des homomorphismes locaux~; soit
$\mathfrak m_\lambda$ l'idéal maximal de $A_\lambda$, et soit
$K_\lambda=A_\lambda/\mathfrak m_\lambda$. Alors
$A'=\varinjlim A_\lambda$ est un anneau local dont
$\mathfrak m'=\varinjlim\mathfrak m_\lambda$ est l'idéal maximal, et
$K=\varinjlim K_\lambda$ le corps résiduel. En outre, si
$\mathfrak m_\mu=\mathfrak m_\lambda A_\mu$ pour $\lambda<\mu$, on a
$\mathfrak m'=\mathfrak m_\lambda A'$ pour tout $\lambda$. Si, de plus, pour
$\lambda<\mu$, $A_\mu$ est un $A_\lambda$-module plat, et si tous les
$A_\lambda$ sont noethériens, alors $A'$ est noethérien et est un
$A_\lambda$-module plat pour tout $\lambda$.
\end{lemma}

Comme $f_{\mu\lambda}(\mathfrak m_\lambda)\subset\mathfrak m_\mu$ pour
$\lambda<\mu$ par hypothèse, les $\mathfrak m_\lambda$ forment un système
inductif, et sa limite $\mathfrak m'$ est évidemment un idéal de $A'$. En
outre, si $x'\notin\mathfrak m'$, il existe $\lambda$ tel que
$x'=f_\lambda(x_\lambda)$ pour un $x_\lambda\in A_\lambda$
($f_\lambda:A_\lambda\to A'$ désignant l'homomorphisme canonique)~; puisque
$x'\notin\mathfrak m'$, on a nécessairement
$x_\lambda\notin\mathfrak m_\lambda$, donc $x_\lambda$ admet un inverse
$y_\lambda$ dans $A_\lambda$, et $y'=f_\lambda(y_\lambda)$ est l'inverse de
$x'$ dans $A'$, ce qui prouve que $A'$ est un anneau local d'idéal maximal
$\mathfrak m'$~; l'assertion relative à $K$ résulte aussitôt du fait que
$\varinjlim$ est un foncteur exact. L'hypothèse
$\mathfrak m_\mu=\mathfrak m_\lambda A_\mu$ signifie que l'application
canonique
$\mathfrak m_\lambda\otimes_{A_\lambda}A_\mu\to\mathfrak m_\mu$ est
surjective~; la relation $\mathfrak m'=\mathfrak m_\lambda A'$ résulte donc
encore de l'exactitude du foncteur $\varinjlim$ et du fait qu'il commute avec
le produit tensoriel.

Supposons maintenant que pour $\lambda<\mu$, on ait
$\mathfrak m_\mu=\mathfrak m_\lambda A_\mu$ et que $A_\mu$ soit un
$A_\lambda$-module plat. Alors $A'$ est un $A_\lambda$-module plat pour tout
$\lambda$, en vertu de (\hyperref[0.6.2.3-fr]{$0_{\mathrm I}$, 6.2.3})~;
comme $A'$ et $A_\lambda$ sont des anneaux locaux et que
$\mathfrak m'=\mathfrak m_\lambda A'$, $A'$ est même un
$A_\lambda$-module \emph{fidèlement plat}
(\hyperref[0.6.6.2-fr]{$0_{\mathrm I}$, 6.6.2}). Supposons enfin, en outre,
les $A_\lambda$ \emph{noethériens}~; les topologies
$\mathfrak m_\lambda$-préadiques sont alors séparées
(\hyperref[0.7.3.5-fr]{$0_{\mathrm I}$, 7.3.5})~; montrons qu'il en résulte
d'abord que sur $A'$ la topologie $\mathfrak m'$-adique est \emph{séparée}.
En effet, si $x'\in A'$ appartient à tous les $\mathfrak m'^{n}$ $(n>0)$,
il est l'image d'un $x_\mu\in A_\mu$ pour un certain indice $\mu$, et comme
l'image réciproque dans $A_\mu$ de
$\mathfrak m'^{n}=\mathfrak m_\mu^n A'$ est $\mathfrak m_\mu^n$
(\hyperref[0.6.6.1-fr]{$0_{\mathrm I}$, 6.6.1}), $x_\mu$ appartient à tous
les $\mathfrak m_\mu^n$, donc $x_\mu=0$ par hypothèse, et par conséquent
$x'=0$. Notons $\widehat A'$ le complété de $A'$ pour la topologie
$\mathfrak m'$-adique~; ce qui précède montre que l'on a
$A'\subset\widehat A'$. Nous allons montrer que $\widehat A'$ est
\emph{noethérien} et $A_\lambda$-\emph{plat} pour tout $\lambda$~; il en
\oldpage[0\textsubscript{III}]{22}
résultera que $\widehat A'$ est $A'$-plat
(\hyperref[0.6.2.3-fr]{$0_{\mathrm I}$, 6.2.3}), et comme
$\mathfrak m'\widehat A'\ne\widehat A'$, $\widehat A'$ est un $A'$-module
fidèlement plat (\hyperref[0.6.6.2-fr]{$0_{\mathrm I}$, 6.6.2}), d'où on
conclura finalement que $A'$ est \emph{noethérien}
(\hyperref[0.6.5.2-fr]{$0_{\mathrm I}$, 6.5.2}), ce qui achèvera la
démonstration du lemme.

On a $\widehat A'=\varprojlim_n A'/\mathfrak m'^{n}$~; en raison de ce que
$A'$ est $A_\lambda$-plat, on a
\[
  \mathfrak m'^{n}/\mathfrak m'^{n+1}
  =(\mathfrak m_\lambda^n/\mathfrak m_\lambda^{n+1})
    \otimes_{A_\lambda}A'
  =(\mathfrak m_\lambda^n/\mathfrak m_\lambda^{n+1})
    \otimes_{K_\lambda}(K_\lambda\otimes_{A_\lambda}A')
  =(\mathfrak m_\lambda^n/\mathfrak m_\lambda^{n+1})
    \otimes_{K_\lambda}K~;
\]
comme $\mathfrak m_\lambda^n/\mathfrak m_\lambda^{n+1}$ est un
$K_\lambda$-espace vectoriel de dimension finie,
$\mathfrak m'^{n}/\mathfrak m'^{n+1}$ est un $K$-espace vectoriel de
dimension finie pour tout $n\geq0$. Il résulte donc de
(\hyperref[0.7.2.12-fr]{$0_{\mathrm I}$, 7.2.12}) et
(\hyperref[0.7.2.8-fr]{$0_{\mathrm I}$, 7.2.8}) que $\widehat A'$ est
noethérien. On sait en outre que l'idéal maximal de $\widehat A'$ est
$\mathfrak m'\widehat A'$ et que
$\widehat A'/\mathfrak m'^{n}\widehat A'$ est isomorphe à
$A'/\mathfrak m'^{n}$~; comme
$A'/\mathfrak m'^{n}=(A_\lambda/\mathfrak m_\lambda^n)
\otimes_{A_\lambda}A'$, $A'/\mathfrak m'^{n}$ est un
$(A_\lambda/\mathfrak m_\lambda^n)$-module plat
(\hyperref[0.6.2.1-fr]{$0_{\mathrm I}$, 6.2.1})~; le critère
(\hyperref[0.10.2.2-fr]{10.2.2}) est donc applicable à la
$A_\lambda$-algèbre noethérienne $\widehat A'$, et montre que $\widehat A'$
est $A_\lambda$-plat. C.Q.F.D.

\begin{env}[10.3.1.4]
\label{0.10.3.1.4-fr}
Abordons maintenant le cas général. Il existe un ordinal $\gamma$ et pour
tout ordinal $\lambda\leq\gamma$ un sous-corps $k_\lambda$ de $K$ contenant
$k$, tels que~: $1^\circ$ Pour tout $\lambda<\gamma$, $k_{\lambda+1}$ soit
une extension de $k_\lambda$ engendrée par un seul élément~; $2^\circ$ Pour
tout ordinal $\mu$ sans prédécesseur,
$k_\mu=\bigcup_{\lambda<\mu}k_\lambda$~; $3^\circ$ $K=k_\gamma$. Il suffit en
effet de considérer une bijection $\xi\mapsto t_\xi$ de l'ensemble des
ordinaux $\xi\leq\beta$ (pour un $\beta$ convenable) sur $K$, de définir
$k_\lambda$ par induction transfinie (pour $\lambda\leq\beta$) comme la
réunion des $k_\mu$ pour $\mu<\lambda$ si $\lambda$ n'a pas de prédécesseur,
et, si $\lambda=\nu+1$, comme $k_\nu(t_\xi)$, où $\xi$ est le plus petit
ordinal tel que $t_\xi\notin k_\nu$~; $\gamma$ est alors par définition le plus
petit ordinal $\leq\beta$ tel que $k_\gamma=K$.

Cela étant, nous allons définir, par récurrence transfinie, une famille
d'anneaux locaux noethériens $A_\lambda$ pour $\lambda\leq\gamma$, et des
homomorphismes locaux $f_{\mu\lambda}:A_\lambda\to A_\mu$ pour
$\lambda\leq\mu$, vérifiant les conditions suivantes~:
\begin{enumerate}
  \item[(i)] $(A_\lambda,f_{\mu\lambda})$ est un système inductif et $A_0=A$.
  \item[(ii)] Pour tout $\lambda$, on a un $k$-isomorphisme
    $A_\lambda/\mathfrak J A_\lambda\xrightarrow{\sim}k_\lambda$.
  \item[(iii)] Pour $\lambda\leq\mu$, $A_\mu$ est un $A_\lambda$-module plat.
\end{enumerate}

Supposons donc les $A_\lambda$ et les $f_{\mu\lambda}$ définis pour
$\lambda<\mu<\xi$, et supposons en premier lieu que $\xi=\zeta+1$, de sorte
que $k_\xi=k_\zeta(t)$. Si $t$ est transcendant sur $k_\zeta$, on définit
$A_\xi$ suivant le procédé de (\hyperref[0.10.3.1.1-fr]{10.3.1.1}) comme
égal à $(A_\zeta[t])_{\mathfrak J A_\zeta[t]}$~;
$f_{\xi\zeta}$ est l'application canonique, et pour $\lambda<\zeta$, on prend
$f_{\xi\lambda}=f_{\xi\zeta}\circ f_{\zeta\lambda}$~; la vérification des
conditions (i) à (iii) est alors immédiate, vu ce qui a été démontré en
(\hyperref[0.10.3.1.1-fr]{10.3.1.1}). Supposons ensuite $t$ algébrique, et
soient $h$ son polynôme minimal dans $k_\zeta[T]$, $H$ un polynôme unitaire de
$A_\zeta[T]$ dont l'image dans $k_\zeta[T]$ est $h$~; on prend alors $A_\xi$
égal à $A_\zeta[T]/(H)$, les $f_{\xi\lambda}$ se définissant comme
précédemment~; la vérification des conditions (i) à (iii) résulte alors de
ce qui a été vu en (\hyperref[0.10.3.1.2-fr]{10.3.1.2}).

Supposons maintenant que $\xi$ n'ait pas de prédécesseur~; on prend alors
pour $A_\xi$ la limite inductive du système inductif d'anneaux locaux
$(A_\lambda,f_{\mu\lambda})$ pour $\lambda<\xi$~; $f_{\xi\lambda}$ est définie
comme l'application canonique pour $\lambda<\xi$. Le fait que $A_\xi$ soit
local noethérien, que les $f_{\xi\lambda}$ soient des homomorphismes locaux,
et les conditions (i) à (iii) pour $\lambda\leq\xi$
\oldpage[0\textsubscript{III}]{23}
résultent alors de l'hypothèse de récurrence et du lemme
(\hyperref[0.10.3.1.3-fr]{10.3.1.3}). Cette construction faite, il est clair
que l'anneau $B=A_\gamma$ vérifie l'énoncé de
(\hyperref[0.10.3.1-fr]{10.3.1}).
\end{env}

On notera qu'en vertu de (\hyperref[0.10.2.1-fr]{10.2.1, $c)$}), on a un
isomorphisme canonique
\[
\label{0.10.3.1.5-fr}
  \operatorname{gr}(A)\otimes_k K\xrightarrow{\sim}\operatorname{gr}(B).
  \tag{10.3.1.5}
\]

D'autre part, on peut remplacer $B$ par son complété $\mathfrak J B$-adique
$\widehat B$ sans changer les conclusions de
(\hyperref[0.10.3.1-fr]{10.3.1}), puisque $\widehat B$ est un $B$-module plat
(\hyperref[0.7.3.3-fr]{$0_{\mathrm I}$, 7.3.3}), donc un $A$-module plat
(\hyperref[0.6.2.1-fr]{$0_{\mathrm I}$, 6.2.1}).

On a en outre démontré le

\begin{corollary}[10.3.2]
\label{0.10.3.2-fr}
Si $K$ est une extension de degré fini, on peut supposer que $B$ est une
$A$-algèbre finie.
\end{corollary}

\section{Compléments d'algèbre homologique}
\label{section:0.11-fr}

\subsection{Rappels sur les suites spectrales}
\label{subsection:0.11.1-fr}

\begin{env}[11.1.1]
\label{0.11.1.1-fr}
Nous utiliserons dans la suite une notion de suite spectrale plus générale que
celle définie dans (T, 2.4)~; gardant les notations de (T, 2.4), nous
appellerons \emph{suite spectrale} dans une catégorie abélienne $\mathcal C$ un
système E formé des éléments suivants~:
\begin{enumerate}
  \item[(a)] Une famille $(E_r^{pq})$ d'objets de $\mathcal C$ définis pour
    $p\in\mathbb Z$, $q\in\mathbb Z$ et $r\geq2$.
  \item[(b)] Une famille de morphismes
    $d_r^{pq}:E_r^{pq}\to E_r^{p+r,q-r+1}$ tels que
    $d_r^{p+r,q-r+1}d_r^{pq}=0$. On pose
    $Z_{r+1}(E_r^{pq})=\operatorname{Ker}(d_r^{pq})$,
    $B_{r+1}(E_r^{pq})=\operatorname{Im}(d_r^{p-r,q+r-1})$, de sorte que
    \[
      B_{r+1}(E_r^{pq})\subset Z_{r+1}(E_r^{pq})\subset E_r^{pq}.
    \]
  \item[(c)] Une famille d'isomorphismes
    $\alpha_r^{pq}:Z_{r+1}(E_r^{pq})/B_{r+1}(E_r^{pq})
    \xrightarrow{\sim}E_{r+1}^{pq}$.
\end{enumerate}

On définit alors pour $k\geq r+1$, par récurrence sur $k$, les sous-objets
$B_k(E_r^{pq})$ et $Z_k(E_r^{pq})$ comme images réciproques, par le morphisme
canonique $E_r^{pq}\to E_r^{pq}/B_{r+1}(E_r^{pq})$, des sous-objets de ce
quotient identifié par $\alpha_r^{pq}$ aux sous-objets $B_k(E_{r+1}^{pq})$ et
$Z_k(E_{r+1}^{pq})$ respectivement. Il est clair que l'on a alors, à un
isomorphisme près
\[
\label{0.11.1.1.1-fr}
  Z_k(E_r^{pq})/B_k(E_r^{pq})=E_k^{pq}
  \qquad\text{pour }k\geq r+1,
  \tag{11.1.1.1}
\]
et, si on pose encore $B_r(E_r^{pq})=0$ et $Z_r(E_r^{pq})=E_r^{pq}$, on a
les relations d'inclusion
\[
\label{0.11.1.1.2-fr}
  0=B_r(E_r^{pq})\subset B_{r+1}(E_r^{pq})\subset B_{r+2}(E_r^{pq})
  \subset\cdots\subset Z_{r+2}(E_r^{pq})\subset Z_{r+1}(E_r^{pq})
  \subset Z_r(E_r^{pq})=E_r^{pq}.
  \tag{11.1.1.2}
\]

Les autres éléments de la donnée de E sont alors~:
\begin{enumerate}
  \item[(d)] Deux sous-objets $B_\infty(E_2^{pq})$ et
    $Z_\infty(E_2^{pq})$ de $E_2^{pq}$ tels que l'on ait
    $B_\infty(E_2^{pq})\subset Z_\infty(E_2^{pq})$ et, pour tout $k\geq2$,
    \[
      B_k(E_2^{pq})\subset B_\infty(E_2^{pq})
      \qquad\text{et}\qquad
      Z_\infty(E_2^{pq})\subset Z_k(E_2^{pq}).
    \]
    On pose
    \[
    \label{0.11.1.1.3-fr}
      E_\infty^{pq}=Z_\infty(E_2^{pq})/B_\infty(E_2^{pq}).
      \tag{11.1.1.3}
    \]
  \item[(e)]
    \oldpage[0\textsubscript{III}]{24}
    Une famille $(E^n)$ d'objets de $\mathcal C$, dont chacun est muni d'une
    \emph{filtration décroissante} $(F^p(E^n))_{p\in\mathbb Z}$. On désigne
    comme d'ordinaire par $\operatorname{gr}(E^n)$ l'objet gradué associé à
    l'objet filtré $E^n$, somme directe des
    $\operatorname{gr}_p(E^n)=F^p(E^n)/F^{p+1}(E^n)$.
  \item[(f)] Pour tout couple $(p,q)\in\mathbb Z\times\mathbb Z$, un
    isomorphisme
    $\beta^{pq}:E_\infty^{pq}\xrightarrow{\sim}
    \operatorname{gr}_p(E^{p+q})$.
\end{enumerate}

La famille $(E^n)$, sans les filtrations, est appelée \emph{l'aboutissement}
de la suite spectrale E.

Supposons que la catégorie $\mathcal C$ admette des sommes directes infinies,
ou que pour tout $r\geq2$ et tout $n\in\mathbb Z$, les couples $(p,q)$ tels
que $p+q=n$ et $E_r^{pq}\ne0$ soient en nombre fini (il suffit qu'il en soit
ainsi pour $r=2$). Alors les
\[
  E_r^{(n)}=\sum_{p+q=n}E_r^{pq}
\]
sont définis, et en désignant par $d_r^{(n)}$ le morphisme
$E_r^{(n)}\to E_r^{(n+1)}$ dont la restriction à $E_r^{pq}$ est $d_r^{pq}$
pour chaque couple $(p,q)$ tel que $p+q=n$, on a
$d_r^{(n+1)}\circ d_r^{(n)}=0$, autrement dit $(E_r^{(n)})_{n\in\mathbb Z}$
est un \emph{complexe $E_r^{(\bullet)}$} dans $\mathcal C$, à opérateur de
dérivation de degré $+1$, et il résulte de c) que
$\mathrm H^n(E_r^{(\bullet)})$ est isomorphe à $E_{r+1}^{(n)}$ pour tout
$r\geq2$.
\end{env}

\begin{env}[11.1.2]
\label{0.11.1.2-fr}
Un \emph{morphisme $u:E\to E'$} d'une suite spectrale E dans une suite
spectrale $E'=(E_r^{\prime pq},E^{\prime n})$ consiste en des systèmes de
morphismes $u_r^{pq}:E_r^{pq}\to E_r^{\prime pq}$,
$u^n:E^n\to E^{\prime n}$, les $u^n$ étant compatibles avec les filtrations
de $E^n$ et $E^{\prime n}$, les diagrammes
\[
  \xymatrix{
    E_r^{pq}\ar[r]^-{d_r^{pq}}\ar[d]_{u_r^{pq}} &
    E_r^{p+r,q-r+1}\ar[d]^{u_r^{p+r,q-r+1}}\\
    E_r^{\prime pq}\ar[r]_{d_r^{\prime pq}} &
    E_r^{\prime p+r,q-r+1}
  }
\]
étant commutatifs~; en outre, par passage aux quotients, $u_r^{pq}$ donne un
morphisme
\[
  \overline u_r^{pq}:
  Z_{r+1}(E_r^{pq})/B_{r+1}(E_r^{pq})\longrightarrow
  Z_{r+1}(E_r^{\prime pq})/B_{r+1}(E_r^{\prime pq})
\]
et on doit avoir
$\alpha_r^{\prime pq}\circ\overline u_r^{pq}
=u_{r+1}^{pq}\circ\alpha_r^{pq}$~; enfin, on doit avoir
$u_2^{pq}(B_\infty(E_2^{pq}))\subset B_\infty(E_2^{\prime pq})$,
$u_2^{pq}(Z_\infty(E_2^{pq}))\subset Z_\infty(E_2^{\prime pq})$~; par
passage aux quotients, $u_2^{pq}$ donne alors un morphisme
$u_\infty^{pq}:E_\infty^{pq}\to E_\infty^{\prime pq}$, et le diagramme
\[
  \xymatrix{
    E_\infty^{pq}\ar[r]^-{u_\infty^{pq}}\ar[d]_{\beta^{pq}} &
    E_\infty^{\prime pq}\ar[d]^{\beta^{\prime pq}}\\
    \operatorname{gr}_p(E^{p+q})
      \ar[r]_{\operatorname{gr}_p(u^{p+q})} &
    \operatorname{gr}_p(E^{\prime p+q})
  }
\]
doit être commutatif.

Les définitions précédentes montrent, par récurrence sur $r$, que si les
$u_2^{pq}$ sont des \emph{isomorphismes}, il en est de même des $u_r^{pq}$
pour $r\geq2$~; \emph{si l'on sait en outre que}
$u_2^{pq}(B_\infty(E_2^{pq}))=B_\infty(E_2^{\prime pq})$ et
$u_2^{pq}(Z_\infty(E_2^{pq}))=Z_\infty(E_2^{\prime pq})$ \emph{et que les
$u^n$ sont des isomorphismes}, alors on peut conclure que $u$ est un
\emph{isomorphisme}.
\end{env}

\begin{env}[11.1.3]
\label{0.11.1.3-fr}
\oldpage[0\textsubscript{III}]{25}
Rappelons que si $(F^p(X))_{p\in\mathbb Z}$ est une filtration (décroissante)
d'un objet $X\in\mathcal C$, on dit que cette filtration est \emph{séparée}
si $\inf(F^p(X))=0$, \emph{discrète} s'il existe un $p$ tel que $F^p(X)=0$,
\emph{exhaustive} (ou \emph{co-séparée}) si $\sup(F^p(X))=X$,
\emph{co-discrète} s'il existe $p$ tel que $F^p(X)=X$.

Nous dirons qu'une suite spectrale $E=(E_r^{pq},E^n)$ est \emph{faiblement
convergente} si l'on a
\[
  B_\infty(E_2^{pq})=\sup_k(B_k(E_2^{pq})),\qquad
  Z_\infty(E_2^{pq})=\inf_k(Z_k(E_2^{pq}))
\]
(autrement dit les objets $B_\infty(E_2^{pq})$ et $Z_\infty(E_2^{pq})$
sont déterminés par les données a) à c) de la suite spectrale E). Nous dirons
que la suite spectrale E est \emph{régulière} si elle est faiblement
convergente et si en outre~:
\begin{enumerate}
  \item[$1^\circ$] Pour tout couple $(p,q)$, la suite décroissante
    $(Z_k(E_2^{pq}))_{k\geq2}$ est \emph{stationnaire}~; l'hypothèse que E
    est faiblement convergente entraîne alors
    $Z_\infty(E_2^{pq})=Z_k(E_2^{pq})$ pour $k$ assez grand (dépendant de
    $p$ et $q$).
  \item[$2^\circ$] Pour tout $n$, la filtration
    $(F^p(E^n))_{p\in\mathbb Z}$ de $E^n$ est \emph{discrète} et
    \emph{exhaustive}.
\end{enumerate}

On dit que la suite spectrale E est \emph{co-régulière} si elle est
faiblement convergente et si en outre~:
\begin{enumerate}
  \item[$3^\circ$] Pour tout couple $(p,q)$, la suite croissante
    $(B_k(E_2^{pq}))_{k\geq2}$ est \emph{stationnaire}, ce qui entraîne
    $B_\infty(E_2^{pq})=B_k(E_2^{pq})$, et par suite
    $E_\infty^{pq}=\inf E_k^{pq}$.
  \item[$4^\circ$] Pour tout $n$, la filtration de $E^n$ est
    \emph{co-discrète}.
\end{enumerate}

Enfin, on dit que E est \emph{birégulière} si elle est à la fois régulière
et co-régulière, autrement dit si on a les conditions suivantes~:
\begin{enumerate}
  \item[(a)] Pour tout couple $(p,q)$, les suites
    $(B_k(E_2^{pq}))_{k\geq2}$ et $(Z_k(E_2^{pq}))_{k\geq2}$ sont
    \emph{stationnaires} et l'on a
    $B_\infty(E_2^{pq})=B_k(E_2^{pq})$ et
    $Z_\infty(E_2^{pq})=Z_k(E_2^{pq})$ pour $k$ assez grand (ce qui entraîne
    $E_\infty^{pq}=E_k^{pq}$).
  \item[(b)] Pour tout $n$, la filtration $(F^p(E^n))_{p\in\mathbb Z}$ est
    \emph{discrète} et \emph{co-discrète} (ce qu'on exprime aussi en disant
    qu'elle est \emph{finie}).
\end{enumerate}

Les suites spectrales définies dans (T, 2.4) sont donc les suites spectrales
birégulières.
\end{env}

\begin{env}[11.1.4]
\label{0.11.1.4-fr}
Supposons que dans la catégorie $\mathcal C$, les limites inductives filtrantes
existent et que le foncteur $\varinjlim$ soit \emph{exact} (ce qui équivaut à
dire que l'axiome AB 5) de (T, 1.5) est vérifié (cf. T, 1.8)). La condition
que la filtration $(F^p(X))_{p\in\mathbb Z}$ d'un objet $X\in\mathcal C$ est
exhaustive s'écrit alors aussi
$\varinjlim_{p\to-\infty}F^p(X)=X$. Si une suite spectrale E est faiblement
convergente, on a
$B_\infty(E_2^{pq})=\varinjlim_{k\to\infty}B_k(E_2^{pq})$~; si en outre
$u:E\to E'$ est un morphisme de E dans une suite spectrale $E'$ de
$\mathcal C$ faiblement convergente, on a
$u_2^{pq}(B_\infty(E_2^{pq}))=B_\infty(E_2^{\prime pq})$ en raison de
l'exactitude de $\varinjlim$. De plus~:
\end{env}

\begin{proposition}[11.1.5]
\label{0.11.1.5-fr}
Soient $\mathcal C$ une catégorie abélienne dans laquelle les limites
inductives filtrantes sont exactes, E, $E'$ deux suites spectrales régulières
de $\mathcal C$, $u:E\to E'$ un morphisme de suites spectrales. Si les
$u_2^{pq}$ sont des isomorphismes, il en est de même de $u$.
\end{proposition}

\begin{proof}
Nous savons déjà (\hyperref[0.11.1.2-fr]{11.1.2}) que les $u_r^{pq}$ sont des
isomorphismes et que
\[
  u_2^{pq}(B_\infty(E_2^{pq}))=B_\infty(E_2^{\prime pq})~;
\]
\oldpage[0\textsubscript{III}]{26}
l'hypothèse que E et $E'$ sont régulières entraîne aussi
$u_2^{pq}(Z_\infty(E_2^{pq}))=Z_\infty(E_2^{\prime pq})$, et comme
$u_2^{pq}$ est un isomorphisme, il en est de même de $u_\infty^{pq}$~; on en
conclut donc que $\operatorname{gr}_p(u^{p+q})$ est aussi un isomorphisme.
Mais comme les filtrations des $E^n$ et des $E^{\prime n}$ sont discrètes et
exhaustives, cela entraîne que les $u^n$ sont aussi des isomorphismes
(Bourbaki, \emph{Alg. comm.}, chap. III, \S~2, n\textsuperscript{o} 8,
th. 1).
\end{proof}

\begin{env}[11.1.6]
\label{0.11.1.6-fr}
Il résulte de (\hyperref[0.11.1.1.2-fr]{11.1.1.2}) et de la définition
(\hyperref[0.11.1.1.3-fr]{11.1.1.3}) que si, dans une suite spectrale E, on a
$E_r^{pq}=0$, on a $E_k^{pq}=0$ pour $k\geq r$ et $E_\infty^{pq}=0$. On dit
qu'une suite spectrale est \emph{dégénérée} s'il existe un entier $r\geq2$ et,
pour tout entier $n\in\mathbb Z$, un entier $q(n)$ tel que
$E_r^{n-q,q}=0$ pour tout $q\ne q(n)$. On déduit d'abord de la remarque
précédente que l'on a aussi $E_k^{n-q,q}=0$ pour $k\geq r$ (y compris
$k=\infty$) et $q\ne q(n)$. En outre, la définition de $E_{r+1}^{pq}$ montre
que l'on a
$E_{r+1}^{n-q(n),q(n)}=E_r^{n-q(n),q(n)}$~; si E est faiblement convergente,
on a donc aussi
$E_\infty^{n-q(n),q(n)}=E_r^{n-q(n),q(n)}$~; autrement dit, pour tout
$n\in\mathbb Z$, $\operatorname{gr}_p(E^n)=0$ pour $p\ne q(n)$ et
$\operatorname{gr}_{q(n)}(E^n)=E_r^{n-q(n),q(n)}$. Si en outre la filtration
de $E^n$ est discrète et exhaustive, la suite E est régulière et on a
$E^n=E_r^{n-q(n),q(n)}$ à un isomorphisme près.
\end{env}

\begin{env}[11.1.7]
\label{0.11.1.7-fr}
Supposons que dans la catégorie $\mathcal C$ les limites inductives filtrantes
existent et soient exactes, et soit $(E_\lambda,u_{\mu\lambda})$ un système
inductif (suivant un ensemble d'indices filtrant) de suites spectrales de
$\mathcal C$. Alors la \emph{limite inductive} de ce système inductif existe
dans la catégorie additive des suites spectrales d'objets de $\mathcal C$~:
il suffit pour le voir de définir $E_r^{pq}$, $d_r^{pq}$, $\alpha_r^{pq}$,
$B_\infty(E_2^{pq})$, $Z_\infty(E_2^{pq})$, $E^n$, $F^p(E^n)$ et
$\beta^{pq}$ comme limites inductives respectives de $E_{r,\lambda}^{pq}$,
$d_{r,\lambda}^{pq}$, $\alpha_{r,\lambda}^{pq}$,
$B_\infty(E_{2,\lambda}^{pq})$, $Z_\infty(E_{2,\lambda}^{pq})$,
$E_\lambda^n$, $F^p(E_\lambda^n)$ et $\beta_\lambda^{pq}$~; la vérification
des conditions de (\hyperref[0.11.1.1-fr]{11.1.1}) résulte de l'exactitude du
foncteur $\varinjlim$ dans $\mathcal C$.
\end{env}

\begin{remark}[11.1.8]
\label{0.11.1.8-fr}
Supposons que la catégorie $\mathcal C$ soit la catégorie des $A$-modules sur
un anneau \emph{noethérien} $A$ (resp. sur un anneau $A$). Alors, les
définitions de (\hyperref[0.11.1.1-fr]{11.1.1}) montrent que si pour un $r$
donné, les $E_r^{pq}$ sont des $A$-modules \emph{de type fini} (resp. de
\emph{longueur finie}), il en est de même de tous les modules $E_s^{pq}$ pour
$s\geq r$, ainsi que des $E_\infty^{pq}$. Si en outre la filtration de
l'aboutissement $(E^n)$ est discrète et co-discrète pour tout $n$, on en
conclut que chacun des $E^n$ est aussi un $A$-module \emph{de type fini}
(resp. de \emph{longueur finie}).
\end{remark}

\begin{env}[11.1.9]
\label{0.11.1.9-fr}
Nous aurons à considérer des conditions assurant qu'une suite spectrale E est
birégulière de façon «~uniforme~» en $p+q=n$. On utilisera alors le lemme
suivant~:
\end{env}

\begin{lemma}[11.1.10]
\label{0.11.1.10-fr}
Soit $(E_r^{pq})$ une famille d'objets de $\mathcal C$ liés par les données
a), b), c) de (\hyperref[0.11.1.1-fr]{11.1.1}). Pour un entier fixé $n$, les
propriétés suivantes sont équivalentes~:
\begin{enumerate}
  \item[(a)] Il existe un entier $r(n)$ tel que pour $r\geq r(n)$,
    $p+q=n$ ou $p+q=n-1$, les morphismes $d_r^{pq}$ soient tous nuls.
  \item[(b)] Il existe un entier $r(n)$ tel que pour $p+q=n$ ou $p+q=n+1$,
    on ait $B_r(E_2^{pq})=B_s(E_2^{pq})$ pour $s\geq r\geq r(n)$.
  \item[(c)] Il existe un entier $r(n)$ tel que pour $p+q=n$ ou $p+q=n-1$,
    on ait $Z_r(E_2^{pq})=Z_s(E_2^{pq})$ pour $s\geq r\geq r(n)$.
  \item[(d)] Il existe un entier $r(n)$ tel que pour $p+q=n$, on ait
    $B_r(E_2^{pq})=B_s(E_2^{pq})$ et $Z_r(E_2^{pq})=Z_s(E_2^{pq})$ pour
    $s\geq r\geq r(n)$.
\end{enumerate}
\end{lemma}

\begin{proof}
\oldpage[0\textsubscript{III}]{27}
En effet, d'après les conditions a), b), c) de
(\hyperref[0.11.1.1-fr]{11.1.1}), dire que
$Z_{r+1}(E_2^{pq})=Z_r(E_2^{pq})$ équivaut à dire que $d_r^{pq}=0$, et dire
que
$B_r(E_2^{p+r,q-r+1})=B_{r+1}(E_2^{p+r,q-r+1})$ équivaut aussi à dire que
$d_r^{pq}=0$~; le lemme résulte aussitôt de cette remarque.
\end{proof}
\subsection{La suite spectrale d'un complexe filtré}
\label{subsection:0.11.2-fr}

\begin{env}[11.2.1]
\label{0.11.2.1-fr}
Étant donnée une catégorie abélienne $\mathcal C$, nous conviendrons de
désigner par des notations telles que $K^\bullet$ les \emph{complexes}
$(K^i)_{i\in\mathbb Z}$ d'objets de $\mathcal C$ dans lesquels l'opérateur de
dérivation est de degré $+1$, par des notations telles que $K_\bullet$ les
complexes $(K_i)_{i\in\mathbb Z}$ d'objets de $\mathcal C$ dans lesquels
l'opérateur de dérivation est de degré $-1$. À tout complexe
$K^\bullet=(K^i)$ dont l'opérateur de dérivation $d$ est de degré $+1$, on
peut associer un complexe $K'_\bullet=(K'_i)$ en posant $K'_i=K^{-i}$,
l'opérateur de dérivation $K'_i\to K'_{i-1}$ étant l'opérateur
$d:K^{-i}\to K^{-i+1}$~; et \emph{vice versa}, ce qui permettra suivant les
circonstances de considérer l'un ou l'autre type de complexes et de traduire
tout résultat pour l'un des types en résultats pour l'autre. Nous désignerons
de même par des notations telles que $K^{\bullet\bullet}=(K^{ij})$ (resp.
$K_{\bullet\bullet}=(K_{ij})$) des \emph{bicomplexes} d'objets de
$\mathcal C$ dans lesquels les \emph{deux} opérateurs de dérivation sont de
degré $+1$ (resp. $-1$)~; on passe encore de l'un à l'autre type en changeant
les signes des indices, et on a des notations et remarques analogues pour des
multicomplexes quelconques. La notation $K^\bullet$ ou $K_\bullet$ sera aussi
utilisée pour des \emph{objets gradués} de $\mathcal C$, de type $\mathbb Z$,
qui ne sont pas nécessairement des complexes (ou que l'on peut considérer
comme tels pour des opérateurs de dérivation nuls)~; par exemple, nous
écrirons
$\mathrm H^\bullet(K^\bullet)=(\mathrm H^i(K^\bullet))_{i\in\mathbb Z}$ la
\emph{cohomologie} d'un complexe $K^\bullet$ dont l'opérateur de dérivation
est de degré $+1$, par
$\mathrm H_\bullet(K_\bullet)=(\mathrm H_i(K_\bullet))_{i\in\mathbb Z}$
l'\emph{homologie} d'un complexe $K_\bullet$ dont l'opérateur de dérivation
est de degré $-1$~; quand on passe de $K^\bullet$ à $K'_\bullet$ par
l'opération décrite ci-dessus, on a
$\mathrm H_i(K'_\bullet)=\mathrm H^{-i}(K^\bullet)$.

Rappelons à ce propos que pour un complexe $K^\bullet$ (resp. $K_\bullet$),
nous écrirons en général
$Z^i(K^\bullet)=\operatorname{Ker}(K^i\to K^{i+1})$ («~objet des cocycles~»)
et $B^i(K^\bullet)=\operatorname{Im}(K^{i-1}\to K^i)$ («~objet des
cobords~») (resp.
$Z_i(K_\bullet)=\operatorname{Ker}(K_i\to K_{i-1})$ («~objet des cycles~»)
et $B_i(K_\bullet)=\operatorname{Im}(K_{i+1}\to K_i)$ («~objet des
bords~»)) de sorte que
$\mathrm H^i(K^\bullet)=Z^i(K^\bullet)/B^i(K^\bullet)$ (resp.
$\mathrm H_i(K_\bullet)=Z_i(K_\bullet)/B_i(K_\bullet)$).

Si $K^\bullet=(K^i)$ (resp. $K_\bullet=(K_i)$) est un complexe dans
$\mathcal C$, et $T:\mathcal C\to\mathcal C'$ un foncteur de $\mathcal C$
dans une catégorie abélienne $\mathcal C'$, nous désignerons par
$T(K^\bullet)$ (resp. $T(K_\bullet)$) le complexe $(T(K^i))$ (resp.
$(T(K_i))$) dans $\mathcal C'$.

Nous ne revenons pas sur la définition des $\partial$-foncteurs (T, 2.1),
sauf pour signaler que nous dirons aussi $\partial$-foncteur au lieu de
$\partial^*$-foncteur lorsque le morphisme $\partial$ diminue le degré d'une
unité, le contexte devant chaque fois préciser ce point s'il peut y avoir
doute.

Enfin, nous dirons qu'un \emph{objet gradué} $(A_i)_{i\in\mathbb Z}$ de
$\mathcal C$ est \emph{limité inférieurement} (resp. \emph{supérieurement})
s'il existe $i_0$ tel que $A_i=0$ pour $i<i_0$ (resp. $i>i_0$).
\end{env}

\begin{env}[11.2.2]
\label{0.11.2.2-fr}
Soit $K^\bullet$ un complexe de $\mathcal C$ dont l'opérateur de dérivation
$d$ est de degré $+1$, et supposons-le muni d'une \emph{filtration}
$F(K^\bullet)=(F^p(K^\bullet))_{p\in\mathbb Z}$ formé de sous-
\oldpage[0\textsubscript{III}]{28}
objets gradués de $K^\bullet$, autrement dit
$F^p(K^\bullet)=(K^i\cap F^p(K^\bullet))_{i\in\mathbb Z}$~; en outre, on
suppose que $d(F^p(K^\bullet))\subset F^p(K^\bullet)$ pour tout
$p\in\mathbb Z$. Rappelons alors rapidement comment on définit
\emph{fonctoriellement} une suite spectrale $E(K^\bullet)$ à partir de
$K^\bullet$ (M, XV, 4 et G, I, 4.3). Pour $r\geq2$, le morphisme canonique
$F^p(K^\bullet)/F^{p+r}(K^\bullet)\to
F^p(K^\bullet)/F^{p+1}(K^\bullet)$ définit un morphisme pour la cohomologie
\[
  \mathrm H^{p+q}(F^p(K^\bullet)/F^{p+r}(K^\bullet))
  \longrightarrow
  \mathrm H^{p+q}(F^p(K^\bullet)/F^{p+1}(K^\bullet)).
\]
On désigne par $Z_r^{pq}(K^\bullet)$ l'image de ce morphisme. De même, de la
suite exacte
\[
  0\to F^p(K^\bullet)/F^{p+1}(K^\bullet)
  \to F^{p-r+1}(K^\bullet)/F^{p+1}(K^\bullet)
  \to F^{p-r+1}(K^\bullet)/F^p(K^\bullet)\to0
\]
on déduit par la suite exacte de cohomologie un morphisme
\[
  \mathrm H^{p+q-1}(F^{p-r+1}(K^\bullet)/F^p(K^\bullet))
  \longrightarrow
  \mathrm H^{p+q}(F^p(K^\bullet)/F^{p+1}(K^\bullet))
\]
et on désigne par $B_r^{pq}(K^\bullet)$ l'image de ce morphisme~; on montre
que $B_r^{pq}(K^\bullet)\subset Z_r^{pq}(K^\bullet)$ et on prend
$E_r^{pq}(K^\bullet)=Z_r^{pq}(K^\bullet)/B_r^{pq}(K^\bullet)$~; nous ne
préciserons pas la définition des $d_r^{pq}$, ni des $\alpha_r^{pq}$.

On notera ici que tous les $Z_r^{pq}(K^\bullet)$ et
$B_r^{pq}(K^\bullet)$, pour $p$, $q$ fixés, sont des sous-objets du même
objet
$\mathrm H^{p+q}(F^p(K^\bullet)/F^{p+1}(K^\bullet))$, que l'on note
$Z_1^{pq}(K^\bullet)$~; on pose $B_1^{pq}(K^\bullet)=0$, de sorte que les
définitions précédentes pour $Z_r^{pq}(K^\bullet)$ et
$B_r^{pq}(K^\bullet)$ s'appliquent aussi pour $r=1$~; on pose encore
$E_1^{pq}(K^\bullet)=Z_1^{pq}(K^\bullet)$. On définit encore les $d_1^{pq}$
et $\alpha_1^{pq}$ de sorte que les conditions de
(\hyperref[0.11.1.1-fr]{11.1.1}) soient vérifiées pour $r=1$. On définit
d'autre part les sous-objets $Z_\infty^{pq}(K^\bullet)$, image du morphisme
\[
  \mathrm H^{p+q}(F^p(K^\bullet))\longrightarrow
  \mathrm H^{p+q}(F^p(K^\bullet)/F^{p+1}(K^\bullet))
  =E_1^{pq}(K^\bullet)
\]
et $B_\infty^{pq}(K^\bullet)$, image du morphisme
\[
  \mathrm H^{p+q-1}(K^\bullet/F^p(K^\bullet))\longrightarrow
  \mathrm H^{p+q}(F^p(K^\bullet)/F^{p+1}(K^\bullet))
  =E_1^{pq}(K^\bullet)
\]
déduit comme ci-dessus d'une suite exacte de cohomologie. On prend pour
$Z_\infty(E_2^{pq}(K^\bullet))$ et $B_\infty(E_2^{pq}(K^\bullet))$ les
images canoniques dans $E_2^{pq}(K^\bullet)$ de
$Z_\infty^{pq}(K^\bullet)$ et $B_\infty^{pq}(K^\bullet)$.

Enfin, on désigne par $F^p(\mathrm H^n(K^\bullet))$ l'image dans
$\mathrm H^n(K^\bullet)$ du morphisme
$\mathrm H^n(F^p(K^\bullet))\to\mathrm H^n(K^\bullet)$ provenant de
l'injection canonique $F^p(K^\bullet)\to K^\bullet$~; par la suite exacte de
cohomologie, c'est aussi le noyau du morphisme
$\mathrm H^n(K^\bullet)\to\mathrm H^n(K^\bullet/F^p(K^\bullet))$. On
définit ainsi une filtration sur
$E^n(K^\bullet)=\mathrm H^n(K^\bullet)$~; nous ne donnerons pas non plus ici
la définition des isomorphismes $\beta^{pq}$.
\end{env}

\begin{env}[11.2.3]
\label{0.11.2.3-fr}
Le caractère \emph{fonctoriel} de $E(K^\bullet)$ doit s'entendre de la façon
suivante~: étant donné deux complexes \emph{filtrés} $K^\bullet$,
$K^{\prime\bullet}$ de $\mathcal C$, et un morphisme de complexes
$u:K^\bullet\to K^{\prime\bullet}$ \emph{compatible avec les filtrations}, on
en déduit de façon évidente les morphismes $u_r^{pq}$ (pour $r\geq1$) et
$u^n$, et on montre que ces morphismes sont compatibles avec les $d_r^{pq}$,
$\alpha_r^{pq}$ et $\beta^{pq}$ au sens de
(\hyperref[0.11.1.2-fr]{11.1.2}), donc définissent bien un morphisme
$E(u):E(K^\bullet)\to E(K^{\prime\bullet})$ de suites spectrales. En outre,
on montre que si $u$ et $v$ sont des morphismes
$K^\bullet\to K^{\prime\bullet}$ du type précédent, \emph{homotopes d'ordre
$\leq k$}, alors $u_r^{pq}=v_r^{pq}$ pour $r>k$ et $u^n=v^n$ pour tout $n$
(M, XV, 3.1).
\end{env}

\begin{env}[11.2.4]
\label{0.11.2.4-fr}
\oldpage[0\textsubscript{III}]{29}
Supposons que dans $\mathcal C$ les limites inductives filtrantes soient
exactes. Alors, si la filtration $(F^p(K^\bullet))$ de $K^\bullet$ est
\emph{exhaustive}, il en est de même de la filtration
$(F^p(\mathrm H^n(K^\bullet)))$ pour tout $n$, car on a par hypothèse
$K^\bullet=\varinjlim_{p\to-\infty}F^p(K^\bullet)$ et l'hypothèse sur
$\mathcal C$ entraîne que la cohomologie commute aux limites inductives. En
outre, on a, pour la même raison,
$B_\infty(E_2^{pq}(K^\bullet))=\sup_k B_k(E_2^{pq}(K^\bullet))$. On dit que
la filtration $(F^p(K^\bullet))$ de $K^\bullet$ est \emph{régulière} si pour
tout $n$ il existe un entier $u(n)$ tel que
$\mathrm H^n(F^p(K^\bullet))=0$ pour $p>u(n)$. Il en est ainsi en particulier
lorsque la filtration de $K^\bullet$ est \emph{discrète}. Lorsque la
filtration de $K^\bullet$ est régulière et exhaustive, et que les limites
inductives filtrantes dans $\mathcal C$ sont exactes, on montre (M, XV, 4)
que la suite spectrale $E(K^\bullet)$ est \emph{régulière}.
\end{env}

\subsection{Les suites spectrales d'un bicomplexe}
\label{subsection:0.11.3-fr}

\begin{env}[11.3.1]
\label{0.11.3.1-fr}
En ce qui concerne les conventions relatives aux bicomplexes, nous suivons
celles de (T, 2.4) plutôt que celles de (M), les deux dérivations $d'$, $d''$
(de degré $+1$) d'un tel bicomplexe
$K^{\bullet\bullet}=(K^{ij})$ étant donc supposées \emph{permutables}.
Supposons vérifiées l'une des deux conditions suivantes~:
\begin{enumerate}
  \item[$1^\circ$] Les sommes directes infinies existent dans $\mathcal C$.
  \item[$2^\circ$] Pour tout $n\in\mathbb Z$, il n'y a qu'un nombre
    \emph{fini} de couples $(p,q)$ tels que $p+q=n$ et $K^{pq}\ne0$.
\end{enumerate}
Alors, le bicomplexe $K^{\bullet\bullet}$ définit un complexe (simple)
$(K^{\prime n})_{n\in\mathbb Z}$, avec
$K^{\prime n}=\sum_{i+j=n}K^{ij}$, l'opérateur de dérivation $d$ (de degré
$+1$) de ce complexe étant donné par
$dx=d'x+(-1)^i d''x$ pour $x\in K^{ij}$. Lorsque nous parlerons par la suite
du complexe (simple) \emph{défini par un bicomplexe $K^{\bullet\bullet}$}, il
sera toujours sous-entendu que l'une des conditions précédentes est satisfaite.
On adoptera des conventions analogues pour les multicomplexes.

On désigne par $K^{i,\bullet}$ (resp. $K^{\bullet,j}$) le complexe simple
$(K^{ij})_{j\in\mathbb Z}$ (resp. $(K^{ij})_{i\in\mathbb Z}$), par
$Z_{\mathrm{II}}^p(K^{i,\bullet})$, $B_{\mathrm{II}}^p(K^{i,\bullet})$,
$\mathrm H_{\mathrm{II}}^p(K^{i,\bullet})$ (resp.
$Z_{\mathrm I}^p(K^{\bullet,j})$, $B_{\mathrm I}^p(K^{\bullet,j})$,
$\mathrm H_{\mathrm I}^p(K^{\bullet,j})$) ses $p$-èmes objets des cocycles,
des cobords et de cohomologie respectivement~; la dérivation
$d':K^{i,\bullet}\to K^{i+1,\bullet}$ est un morphisme de complexes, qui
donne donc un opérateur dans les cocycles, les cobords et la cohomologie,
\[
\begin{aligned}
  d'&:Z_{\mathrm{II}}^p(K^{i,\bullet})
      \longrightarrow Z_{\mathrm{II}}^p(K^{i+1,\bullet}),\\
  d'&:B_{\mathrm{II}}^p(K^{i,\bullet})
      \longrightarrow B_{\mathrm{II}}^p(K^{i+1,\bullet}),\\
  d'&:\mathrm H_{\mathrm{II}}^p(K^{i,\bullet})
      \longrightarrow \mathrm H_{\mathrm{II}}^p(K^{i+1,\bullet}),
\end{aligned}
\]
et il est clair que pour ces opérateurs,
$(Z_{\mathrm{II}}^p(K^{i,\bullet}))_{i\in\mathbb Z}$,
$(B_{\mathrm{II}}^p(K^{i,\bullet}))_{i\in\mathbb Z}$ et
$(\mathrm H_{\mathrm{II}}^p(K^{i,\bullet}))_{i\in\mathbb Z}$ sont des
complexes~; nous désignerons le complexe
$(\mathrm H_{\mathrm{II}}^p(K^{i,\bullet}))_{i\in\mathbb Z}$ par
$\mathrm H_{\mathrm{II}}^p(K^{\bullet\bullet})$, ses $q$-èmes objets de
cocycles, de cobords et de cohomologie par
$Z_{\mathrm I}^q(\mathrm H_{\mathrm{II}}^p(K^{\bullet\bullet}))$,
$B_{\mathrm I}^q(\mathrm H_{\mathrm{II}}^p(K^{\bullet\bullet}))$ et
$\mathrm H_{\mathrm I}^q(\mathrm H_{\mathrm{II}}^p(K^{\bullet\bullet}))$.
On définit de même les complexes
$\mathrm H_{\mathrm I}^p(K^{\bullet\bullet})$ et leurs objets de cohomologie
$\mathrm H_{\mathrm{II}}^q(\mathrm H_{\mathrm I}^p(K^{\bullet\bullet}))$.
Rappelons d'autre part que $\mathrm H^n(K^{\bullet\bullet})$ désigne le
$n$-ème objet de cohomologie du complexe (simple) défini par
$K^{\bullet\bullet}$.
\end{env}

\begin{env}[11.3.2]
\label{0.11.3.2-fr}
Sur le complexe défini par un bicomplexe $K^{\bullet\bullet}$, on peut
considérer deux filtrations canoniques
$(F_{\mathrm I}^p(K^{\bullet\bullet}))$ et
$(F_{\mathrm{II}}^p(K^{\bullet\bullet}))$ données par
\[
\label{0.11.3.2.1-fr}
  F_{\mathrm I}^p(K^{\bullet\bullet})
  =\left(\sum_{\substack{i+j=n\\i\geq p}}K^{ij}\right)_{n\in\mathbb Z},
  \qquad
  F_{\mathrm{II}}^p(K^{\bullet\bullet})
  =\left(\sum_{\substack{i+j=n\\j\geq p}}K^{ij}\right)_{n\in\mathbb Z}.
  \tag{11.3.2.1}
\]
\oldpage[0\textsubscript{III}]{30}
qui, par définition, sont bien des sous-objets gradués du complexe (simple)
défini par $K^{\bullet\bullet}$, et font donc de ce complexe un complexe
filtré~; par ailleurs, il est clair que ces filtrations sont exhaustives et
séparées.

Il correspond à chacune de ces filtrations une suite spectrale
(\hyperref[0.11.2.2-fr]{11.2.2})~; nous désignerons par
$'E(K^{\bullet\bullet})$ et $''E(K^{\bullet\bullet})$ les suites spectrales
correspondant à $(F_{\mathrm I}^p(K^{\bullet\bullet}))$ et
$(F_{\mathrm{II}}^p(K^{\bullet\bullet}))$ respectivement, dites
\emph{suites spectrales du bicomplexe $K^{\bullet\bullet}$}, et ayant toutes
deux pour aboutissement la cohomologie
$(\mathrm H^n(K^{\bullet\bullet}))$. On montre en outre (M, XV, 6) que l'on a
\[
\label{0.11.3.2.2-fr}
  'E_2^{pq}(K^{\bullet\bullet})
  =\mathrm H_{\mathrm I}^p(\mathrm H_{\mathrm{II}}^q(K^{\bullet\bullet})),
  \qquad
  ''E_2^{pq}(K^{\bullet\bullet})
  =\mathrm H_{\mathrm{II}}^q(\mathrm H_{\mathrm I}^p(K^{\bullet\bullet})).
  \tag{11.3.2.2}
\]

Tout morphisme $u:K^{\bullet\bullet}\to K^{\prime\bullet\bullet}$ de
bicomplexes est \emph{ipso facto} compatible avec les filtrations de même
type de $K^{\bullet\bullet}$ et $K^{\prime\bullet\bullet}$, donc définit un
morphisme pour chacune des deux suites spectrales~; en outre, deux morphismes
\emph{homotopes} définissent une homotopie \emph{d'ordre $\leq1$} des
complexes (simples) filtrés correspondant, donc le \emph{même} morphisme pour
chacune des deux suites spectrales (M, XV, 6.1).
\end{env}

\begin{proposition}[11.3.3]
\label{0.11.3.3-fr}
Soit $K^{\bullet\bullet}=(K^{ij})$ un bicomplexe dans une catégorie abélienne
$\mathcal C$.
\begin{enumerate}
  \item[(i)] S'il existe $i_0$ et $j_0$ tels que $K^{ij}=0$ pour $i<i_0$ ou
    $j<j_0$ (resp. $i>i_0$ ou $j>j_0$), les deux suites spectrales
    $'E(K^{\bullet\bullet})$ et $''E(K^{\bullet\bullet})$ sont
    birégulières.
  \item[(ii)] S'il existe $i_0$ et $i_1$ tels que $K^{ij}=0$ pour $i<i_0$ ou
    $i>i_1$ (resp. s'il existe $j_0$ et $j_1$ tels que $K^{ij}=0$ pour
    $j<j_0$ ou $j>j_1$) les deux suites spectrales
    $'E(K^{\bullet\bullet})$ et $''E(K^{\bullet\bullet})$ sont
    birégulières.
\end{enumerate}
Supposons en outre que dans $\mathcal C$ les limites inductives filtrantes
existent et soient exactes. Alors~:
\begin{enumerate}
  \item[(iii)] S'il existe $i_0$ tel que $K^{ij}=0$ pour $i>i_0$ (resp. s'il
    existe $j_0$ tel que $K^{ij}=0$ pour $j<j_0$), la suite
    $'E(K^{\bullet\bullet})$ est régulière.
  \item[(iv)] S'il existe $i_0$ tel que $K^{ij}=0$ pour $i<i_0$ (resp. s'il
    existe $j_0$ tel que $K^{ij}=0$ pour $j>j_0$), la suite
    $''E(K^{\bullet\bullet})$ est régulière.
\end{enumerate}
\end{proposition}

\begin{proof}
La proposition résulte aussitôt des définitions
(\hyperref[0.11.1.3-fr]{11.1.3}) et de
(\hyperref[0.11.2.4-fr]{11.2.4}), ainsi que des observations suivantes
relatives à la filtration $F_{\mathrm I}$ (et des observations analogues qu'on
en déduit pour $F_{\mathrm{II}}$ en échangeant les rôles des deux indices
dans $K^{\bullet\bullet}$)~:
\begin{enumerate}
  \item[$1^\circ$] S'il existe $i_0$ tel que $K^{ij}=0$ pour $i>i_0$, la
    filtration $F_{\mathrm I}(K^{\bullet\bullet})$ est discrète.
  \item[$2^\circ$] S'il existe $i_0$ tel que $K^{ij}=0$ pour $i<i_0$, la
    filtration $F_{\mathrm I}(K^{\bullet\bullet})$ est co-discrète. On en
    déduit aussitôt qu'il en est de même de la filtration correspondante
    $F_{\mathrm I}(\mathrm H^n(K^{\bullet\bullet}))$ pour tout $n$~; en
    outre, la définition de $B_r^{pq}$ correspondant à la filtration
    $F_{\mathrm I}(K^{\bullet\bullet})$
    (\hyperref[0.11.2.2-fr]{11.2.2}) montre que pour tout couple $(p,q)$, la
    suite $(B_r^{pq})_{r\geq2}$ est stationnaire.
  \item[$3^\circ$] S'il existe $j_0$ tel que $K^{ij}=0$ pour $j<j_0$, on a
    \[
      F_{\mathrm I}^{p+r}(K^{\bullet\bullet})
      \cap\left(\sum_{i+j=n}K^{ij}\right)=0
    \]
    dès que $p+r+j_0>n$, donc
    $Z_r^{pq}=Z_\infty(E_2^{pq})$ pour $r>q-j_0+1$~; d'autre part,
    $\mathrm H^n(F_{\mathrm I}^p(K^{\bullet\bullet}))=0$ pour
    $p>n-j_0+1$.
  \item[$4^\circ$] S'il existe $j_0$ tel que $K^{ij}=0$ pour $j>j_0$, on a
    \[
      F_{\mathrm I}^{p-r+1}(K^{\bullet\bullet})
      \cap\left(\sum_{i+j=n}K^{ij}\right)=\sum_{i+j=n}K^{ij}.
    \]
    \oldpage[0\textsubscript{III}]{31}
    dès que $p-r+1+j_0<n$, donc
    $B_r^{pq}=B_\infty(E_2^{pq})$ pour $r<j_0-q+1$~; d'autre part,
    $\mathrm H^n(F_{\mathrm I}^p(K^{\bullet\bullet}))
    =\mathrm H^n(K^{\bullet\bullet})$ pour $p+j_0<n-1$.
\end{enumerate}
\end{proof}

\begin{env}[11.3.4]
\label{0.11.3.4-fr}
Supposons que le bicomplexe $K^{\bullet\bullet}=(K^{ij})$ soit tel que
$K^{ij}=0$ pour $i<0$ ou $j<0$. On sait qu'on peut alors définir pour tout
$p\in\mathbb Z$ un «~edge-homomorphisme~» canonique
\[
\label{0.11.3.4.1-fr}
  'E_2^{p0}(K^{\bullet\bullet})\longrightarrow
  \mathrm H^p(K^{\bullet\bullet})
  \tag{11.3.4.1}
\]
(M, XV, 6). Rappelons rapidement que cela est dû, d'une part à ce que l'on a
dans la suite spectrale $'E(K^{\bullet\bullet})$,
$Z_r^{p0}=Z_{\mathrm I}^p(Z_{\mathrm{II}}^0(K^{\bullet\bullet}))$ pour
$2\leq r\leq+\infty$, et de l'autre au fait que
$\mathrm H^p(F_{\mathrm I}^{p+1}(K^{\bullet\bullet}))=0$, si bien que
l'isomorphisme
$\beta^{p0}:'E_\infty^{p0}\xrightarrow{\sim}
\mathrm H^p(F_{\mathrm I}^p)/\mathrm H^p(F_{\mathrm I}^{p+1})$ donne un
homomorphisme
$'E_\infty^{p0}\to\mathrm H^p(F_{\mathrm I}^p(K^{\bullet\bullet}))
\to\mathrm H^p(K^{\bullet\bullet})$~; l'égalité de tous les $Z_r^{p0}$ permet
alors de définir des homomorphismes canoniques
$'E_r^{p0}\to'E_s^{p0}$ pour $r\leq s$, et en particulier un homomorphisme
$'E_2^{p0}\to'E_\infty^{p0}$, d'où par composition l'edge-homomorphisme
$'E_2^{p0}\to\mathrm H^p(K^{\bullet\bullet})$~; en outre, on vérifie aussitôt
que, à la classe mod. $B_2^{p0}$ d'un élément
$z\in Z_{\mathrm{II}}^0(K^{\bullet\bullet})\subset K^{p0}$ tel que $d'z=0$,
l'edge-homomorphisme ainsi défini fait correspondre dans $'E_\infty^{p0}$ la
classe de $z$ mod. $B_\infty^{p0}$, puis à cette dernière la classe de
cohomologie de $z$ dans $\mathrm H^p(K^{\bullet\bullet})$. On voit donc
finalement que \emph{l'edge-homomorphisme
(\hyperref[0.11.3.4.1-fr]{11.3.4.1}) provient, par passage à la cohomologie,
de l'injection canonique}
$Z_{\mathrm{II}}^0(K^{\bullet\bullet})\to K^{\bullet\bullet}$ (où
$K^{\bullet\bullet}$ est considéré comme complexe simple). On interprète
naturellement de même l'edge-homomorphisme
\[
\label{0.11.3.4.2-fr}
  ''E_2^{p0}(K^{\bullet\bullet})\longrightarrow
  \mathrm H^p(K^{\bullet\bullet})
  \tag{11.3.4.2}
\]
comme provenant de l'injection canonique
$Z_{\mathrm I}^0(K^{\bullet\bullet})\to K^{\bullet\bullet}$.
\end{env}

\begin{env}[11.3.5]
\label{0.11.3.5-fr}
Soit maintenant $K_{\bullet\bullet}=(K_{ij})$ un bicomplexe de
$\mathcal C$ dont les deux opérateurs de dérivation sont de degré $-1$. Nous
écrirons alors $K_{i,\bullet}$ (resp. $K_{\bullet,j}$) le complexe simple
$(K_{ij})_{j\in\mathbb Z}$ (resp. $(K_{ij})_{i\in\mathbb Z}$),
$\mathrm H_p^{\mathrm{II}}(K_{i,\bullet})$ (resp.
$\mathrm H_p^{\mathrm I}(K_{\bullet,j})$) son $p$-ème objet d'homologie,
$\mathrm H_p^{\mathrm{II}}(K_{\bullet\bullet})$ (resp.
$\mathrm H_p^{\mathrm I}(K_{\bullet\bullet})$) le complexe
$(\mathrm H_p^{\mathrm{II}}(K_{i,\bullet}))_{i\in\mathbb Z}$ (resp.
$(\mathrm H_p^{\mathrm I}(K_{\bullet,j}))_{j\in\mathbb Z}$),
$\mathrm H_q^{\mathrm I}(\mathrm H_p^{\mathrm{II}}(K_{\bullet\bullet}))$
(resp.
$\mathrm H_q^{\mathrm{II}}(\mathrm H_p^{\mathrm I}(K_{\bullet\bullet}))$)
son $q$-ème objet d'homologie~; notations analogues pour les objets des cycles
et objets des bords~; enfin, $\mathrm H_n(K_{\bullet\bullet})$ désignera
(lorsqu'il existe) le $n$-ème objet d'homologie du complexe simple (à
opérateur de dérivation de degré $-1$) défini par $K_{\bullet\bullet}$.

Soit $K^{\prime\bullet\bullet}=(K^{\prime ij})$ avec
$K^{\prime ij}=K_{-i,-j}$ le bicomplexe à opérateurs de dérivation de degrés
$+1$ associé à $K_{\bullet\bullet}$~; par définition, les
\emph{suites spectrales} de $K_{\bullet\bullet}$ sont celles de
$K^{\prime\bullet\bullet}$, que l'on écrit $'E(K_{\bullet\bullet})$ et
$''E(K_{\bullet\bullet})$, où l'on change toutefois les notations, en posant
\[
  'E_{pq}^r(K_{\bullet\bullet})
  ='E_r^{-p,-q}(K^{\prime\bullet\bullet}),\qquad
  ''E_{pq}^r(K_{\bullet\bullet})
  =''E_r^{-p,-q}(K^{\prime\bullet\bullet})
\]
pour $2\leq r\leq\infty$. Avec ces notations, on a
\[
  'E_{pq}^2(K_{\bullet\bullet})
  =\mathrm H_p^{\mathrm I}(\mathrm H_q^{\mathrm{II}}(K_{\bullet\bullet})),
  \qquad
  ''E_{pq}^2(K_{\bullet\bullet})
  =\mathrm H_q^{\mathrm{II}}(\mathrm H_p^{\mathrm I}(K_{\bullet\bullet})).
\]
Pour éviter des erreurs de signe, il sera en général préférable, pour les
relations entre ces suites spectrales et leur aboutissement, de revenir au
complexe $K^{\prime\bullet\bullet}$. Notons toutefois les critères
correspondant à (\hyperref[0.11.3.3-fr]{11.3.3})~:
\end{env}

\begin{env}[11.3.6]
\label{0.11.3.6-fr}
Les suites spectrales $'E(K_{\bullet\bullet})$ et
$''E(K_{\bullet\bullet})$ sont \emph{birégulières} dans les cas suivants~:
\begin{enumerate}
  \item[(a)] Il existe $i_0$ et $j_0$ tels que $K_{ij}=0$ pour $i>i_0$ ou
    pour $j>j_0$ (resp. pour $i<i_0$
    \oldpage[0\textsubscript{III}]{32}
    ou pour $j<j_0$).
  \item[(b)] Il existe $i_0$ et $i_1$ tels que $K_{ij}=0$ pour $i<i_0$ et
    $i>i_1$.
  \item[(c)] Il existe $j_0$ et $j_1$ tels que $K_{ij}=0$ pour $j<j_0$ et
    $j>j_1$.
\end{enumerate}

La suite $'E(K_{\bullet\bullet})$ est régulière s'il existe $i_0$ tel que
$K_{ij}=0$ pour $i<i_0$, ou s'il existe $j_0$ tel que $K_{ij}=0$ pour
$j>j_0$.

La suite $''E(K_{\bullet\bullet})$ est régulière s'il existe $i_0$ tel que
$K_{ij}=0$ pour $i>i_0$, ou s'il existe $j_0$ tel que $K_{ij}=0$ pour
$j<j_0$.
\end{env}

\subsection{Hypercohomologie d'un foncteur par rapport à un complexe
$K^\bullet$}
\label{subsection:0.11.4-fr}

\begin{env}[11.4.1]
\label{0.11.4.1-fr}
Soit $\mathcal C$ une catégorie abélienne~; rappelons que l'on appelle
\emph{résolution droite} (ou \emph{cohomologique}) d'un objet $A$ de
$\mathcal C$ un complexe d'objets de $\mathcal C$, dont l'opérateur de
dérivation est de degré $+1$,
\[
  0\longrightarrow L^0\longrightarrow L^1\longrightarrow L^2
  \longrightarrow\cdots
\]
muni d'un morphisme $\varepsilon:A\to L^0$ dit \emph{augmentation} de la
résolution (et que l'on peut considérer comme un morphisme de complexes
\[
\xymatrix{
  0\ar[r] & A\ar[r]\ar[d]_{\varepsilon} & 0\ar[r]\ar[d]
    & 0\ar[r]\ar[d] & \cdots\\
  0\ar[r] & L^0\ar[r] & L^1\ar[r] & L^2\ar[r] & \cdots
})
\]
tel que la suite
\[
  0\longrightarrow A\xrightarrow{\varepsilon}L^0\longrightarrow L^1
  \longrightarrow\cdots
\]
soit exacte~; de même une \emph{résolution gauche} (ou \emph{homologique})
de $A$ est un complexe
$0\leftarrow L_0\leftarrow L_1\leftarrow\cdots$ d'objets de $\mathcal C$
dont l'opérateur de dérivation est de degré $-1$, muni d'une augmentation
$\varepsilon:L_0\to A$, de sorte que la suite
\[
  0\longleftarrow A\xleftarrow{\varepsilon}L_0\longleftarrow L_1
  \longleftarrow\cdots
\]
soit exacte.

Lorsqu'une résolution droite $(L_i)_{i\geq0}$ d'un objet $A$ est telle que
$L_i=0$ pour $i\geq n+1$, on dit que cette résolution est de
\emph{longueur $\leq n$}. On définit de même une résolution gauche de
longueur $\leq n$. Une résolution qui est de longueur $\leq n$ pour un entier
$n$ est dite \emph{finie}.

Une résolution de $A$ est dite \emph{projective} (resp. \emph{injective}) si
les objets de $\mathcal C$ autres que $A$ qui la composent sont
\emph{projectifs} (resp. \emph{injectifs}). Lorsque $\mathcal C$ est la
catégorie des modules (à gauche par exemple) sur un anneau, on dira de même
qu'une résolution de $A$ est \emph{plate} (resp. \emph{libre}) lorsque les
modules autres que $A$ qui la composent sont \emph{plats} (resp.
\emph{libres}).
\end{env}

\begin{env}[11.4.2]
\label{0.11.4.2-fr}
Soit $K^\bullet=(K^i)_{i\in\mathbb Z}$ un complexe d'objets de $\mathcal C$,
dont l'opérateur de dérivation est de degré $+1$.

On appelle \emph{résolution de Cartan-Eilenberg droite} de $K^\bullet$ le
couple formé d'un bicomplexe $L^{\bullet\bullet}=(L^{ij})$ à opérateurs de
dérivation de degré $+1$, avec $L^{ij}=0$ pour $j<0$, et d'un morphisme de
complexes simples $\varepsilon:K^\bullet\to L^{\bullet,0}$, de façon que les
conditions suivantes soient remplies~:
\oldpage[0\textsubscript{III}]{33}
\begin{enumerate}
  \item[(i)] Pour chaque indice $i$, les suites
    \[
    \begin{aligned}
      0\longrightarrow{}&K^i\xrightarrow{\varepsilon}L^{i0}
        \longrightarrow L^{i1}\longrightarrow\cdots\\
      0\longrightarrow{}&\mathrm B^i(K^\bullet)\xrightarrow{\varepsilon}
        \mathrm B_{\mathrm I}^i(L^{\bullet,0})
        \longrightarrow\mathrm B_{\mathrm I}^i(L^{\bullet,1})
        \longrightarrow\cdots\\
      0\longrightarrow{}&\mathrm Z^i(K^\bullet)\xrightarrow{\varepsilon}
        \mathrm Z_{\mathrm I}^i(L^{\bullet,0})
        \longrightarrow\mathrm Z_{\mathrm I}^i(L^{\bullet,1})
        \longrightarrow\cdots\\
      0\longrightarrow{}&\mathrm H^i(K^\bullet)\xrightarrow{\varepsilon}
        \mathrm H_{\mathrm I}^i(L^{\bullet,0})
        \longrightarrow\mathrm H_{\mathrm I}^i(L^{\bullet,1})
        \longrightarrow\cdots
    \end{aligned}
    \]
    sont exactes, autrement dit $(L^{i,\bullet})$,
    $(\mathrm B_{\mathrm I}^i(L^{\bullet,\bullet}))$,
    $(\mathrm Z_{\mathrm I}^i(L^{\bullet,\bullet}))$ et
    $(\mathrm H_{\mathrm I}^i(L^{\bullet,\bullet}))$ sont respectivement
    des \emph{résolutions} de $K^i$, $\mathrm B^i(K^\bullet)$,
    $\mathrm Z^i(K^\bullet)$ et $\mathrm H^i(K^\bullet)$.
  \item[(ii)] Pour chaque $j$, le complexe simple $L^{\bullet,j}$ est
    \emph{scindé}, autrement dit, les suites exactes
    \[
      \label{0.11.4.2.1-fr}
      0\longrightarrow\mathrm B_{\mathrm I}^i(L^{\bullet,j})
      \longrightarrow\mathrm Z_{\mathrm I}^i(L^{\bullet,j})
      \longrightarrow\mathrm H_{\mathrm I}^i(L^{\bullet,j})
      \longrightarrow0
      \tag{11.4.2.1}
    \]
    \[
      \label{0.11.4.2.2-fr}
      0\longrightarrow\mathrm Z_{\mathrm I}^i(L^{\bullet,j})
      \longrightarrow L^{ij}
      \longrightarrow\mathrm B_{\mathrm I}^{i+1}(L^{\bullet,j})
      \longrightarrow0
      \tag{11.4.2.2}
    \]
    sont scindées.
\end{enumerate}

On prouve (M, XVII, 1.2) que si tout objet de $\mathcal C$ est sous-objet
d'un objet injectif, tout complexe $K^\bullet$ de $\mathcal C$ admet une
résolution de Cartan-Eilenberg \emph{injective}, c'est-à-dire formée d'objets
injectifs $L^{ij}$ (la condition (ii) ci-dessus entraîne alors que les
$\mathrm B_{\mathrm I}^i(L^{\bullet,j})$,
$\mathrm Z_{\mathrm I}^i(L^{\bullet,j})$ et
$\mathrm H_{\mathrm I}^i(L^{\bullet,j})$ sont aussi des objets injectifs).
En outre, pour tout morphisme $f:K^\bullet\to K'^\bullet$ de complexes de
$\mathcal C$, toute résolution de Cartan-Eilenberg $L^{\bullet\bullet}$ de
$K^\bullet$ et toute résolution injective de Cartan-Eilenberg
$L'^{\bullet\bullet}$ de $K'^\bullet$, il existe un morphisme de bicomplexes
$F:L^{\bullet\bullet}\to L'^{\bullet\bullet}$ compatible avec $f$ et les
augmentations, et si $f$ et $g$ sont deux morphismes homotopes de $K^\bullet$
dans $K'^\bullet$, les morphismes correspondants de $L^{\bullet\bullet}$
dans $L'^{\bullet\bullet}$ sont homotopes (\emph{loc. cit.}).

Lorsque $K^\bullet$ est \emph{limité inférieurement} (resp.
\emph{supérieurement}), on peut prendre $L^{\bullet\bullet}$ tel que
$L^{ij}=0$ pour $i<i_0$ (resp. $i>i_0$) si $K^i=0$ pour $i<i_0$ (resp.
$i>i_0$) (M, XVII, 1.3).

Supposons d'autre part qu'il existe un entier $n$ tel que tout objet de
$\mathcal C$ admette une \emph{résolution injective de longueur $\leq n$}~;
alors on peut supposer que l'on a $L^{ij}=0$ pour $j>n$ (M, XVII, 1.4).
\end{env}

\begin{env}[11.4.3]
\label{0.11.4.3-fr}
Soit maintenant $T$ un \emph{foncteur covariant additif} de $\mathcal C$
dans une catégorie abélienne $\mathcal C'$. Étant donnés un complexe
$K^\bullet$ de $\mathcal C$ et une résolution de Cartan-Eilenberg injective
$L^{\bullet\bullet}$ de $K^\bullet$, supposons que le complexe (simple) défini
par le bicomplexe $T(L^{\bullet\bullet})$ existe
(cf. \hyperref[0.11.3.1-fr]{11.3.1})~; alors les deux suites spectrales
$'E(T(L^{\bullet\bullet}))$ et $''E(T(L^{\bullet\bullet}))$ de ce bicomplexe
sont dites \emph{suites spectrales d'hypercohomologie de $T$ par rapport au
complexe $K^\bullet$}~; en vertu de
(\hyperref[0.11.4.2-fr]{11.4.2}) et
(\hyperref[0.11.3.2-fr]{11.3.2}), elles ne dépendent effectivement que de
$K^\bullet$ et non de la résolution de Cartan-Eilenberg injective
$L^{\bullet\bullet}$ choisie~; en outre, elles dépendent
\emph{fonctoriellement} de $K^\bullet$. Elles ont un même aboutissement
$\mathrm H^\bullet(T(L^{\bullet\bullet}))$ appelé encore
l'\emph{hypercohomologie de $T$ par rapport à $K^\bullet$}, et noté
$\mathrm R^\bullet T(K^\bullet)$. On montre que les termes $E_2$ des deux
suites spectrales précédentes sont donnés par
\[
  \label{0.11.4.3.1-fr}
  'E_2^{pq}=\mathrm H^p(\mathrm R^qT(K^\bullet))
  \tag{11.4.3.1}
\]
\[
  \label{0.11.4.3.2-fr}
  ''E_2^{pq}=\mathrm R^pT(\mathrm H^q(K^\bullet))
  \tag{11.4.3.2}
\]
\oldpage[0\textsubscript{III}]{34}
$\mathrm R^pT$ désignant comme d'ordinaire le $p$-ième \emph{foncteur
dérivé} de $T$ pour $p\in\mathbb Z$~; $\mathrm R^qT(K^\bullet)$ désigne le
complexe $(\mathrm R^qT(K^i))_{i\in\mathbb Z}$.

\phantomsection\label{0.11.4.4-fr}
Sauf mention expresse du contraire, nous supposerons désormais que tout objet
de $\mathcal C$ est sous-objet d'un objet injectif de $\mathcal C$, de sorte
que les résolutions de Cartan-Eilenberg injectives existent pour tout complexe
de $\mathcal C$. Comme $L^{ij}=0$ pour $j<0$, les critères de
(\hyperref[0.11.3.3-fr]{11.3.3}) montrent que les \emph{deux} suites
spectrales d'hypercohomologie de $T$ par rapport à $K^\bullet$ existent et
sont \emph{birégulières} dans chacun des deux cas suivants~:
\begin{enumerate}
  \item[$1^\circ$] $K^\bullet$ est limité inférieurement~;
  \item[$2^\circ$] Tout objet de $\mathcal C$ admet une résolution injective
    de longueur au plus égale à un entier $n$ (indépendant de l'objet
    considéré).
\end{enumerate}
En effet, dans le premier cas, on peut supposer
(\hyperref[0.11.4.2-fr]{11.4.2}) qu'il existe $i_0$ tel que $L^{ij}=0$ pour
$i<i_0$ et dans le second qu'il existe $j_1$ tel que $L^{ij}=0$ pour $j>j_1$~;
dans chacun des deux cas, il est clair en outre que pour $n$ donné, il n'y a
qu'un nombre fini de couples $(i,j)$ tels que $L^{ij}\neq0$ et $i+j=n$, ce
qui établit nos assertions.

Lorsqu'on suppose que dans $\mathcal C'$ les limites inductives filtrantes
existent et sont exactes (ce qui implique en particulier l'existence dans
$\mathcal C'$ des sommes directes infinies), alors le complexe défini par le
bicomplexe $T(L^{\bullet\bullet})$ existe, et le critère
(\hyperref[0.11.3.3-fr]{11.3.3}) montre que la suite
$'E(T(L^{\bullet\bullet}))$ est toujours \emph{régulière}.

Lorsque $K^\bullet$ est un complexe dont tous les termes $K^i$ sont nuls sauf
un seul $K^{i_0}$, $\mathrm R^nT(K^\bullet)$ est isomorphe à
$\mathrm R^{n-i_0}T(K^{i_0})$, ainsi qu'il résulte aussitôt des définitions en
prenant une résolution de Cartan-Eilenberg $L^{\bullet\bullet}$ telle que
$L^{ij}=0$ pour $i\neq i_0$.

Si $K^\bullet$ et $K'^\bullet$ sont deux complexes de $\mathcal C$, $f$, $g$
deux morphismes \emph{homotopes} de $K^\bullet$ dans $K'^\bullet$, alors les
morphismes $\mathrm R^\bullet T(K^\bullet)\to
\mathrm R^\bullet T(K'^\bullet)$ déduits de $f$ et $g$ sont identiques, et il
en est de même pour les morphismes des suites spectrales de cohomologie.
\end{env}

\begin{proposition}[11.4.5]
\label{0.11.4.5-fr}
Supposons que dans $\mathcal C'$ les limites inductives filtrantes existent et
soient exactes. Si $\mathrm R^nT(K^i)=0$ pour tout $n>0$ et tout
$i\in\mathbb Z$, on a des isomorphismes fonctoriels
\[
  \label{0.11.4.5.1-fr}
  \mathrm R^iT(K^\bullet)\xrightarrow{\sim}
  \mathrm H^i(T(K^\bullet))\qquad(i\in\mathbb Z).
  \tag{11.4.5.1}
\]
\end{proposition}

\begin{proof}
En effet, les seuls termes $E_2$ non nuls de la première suite spectrale
(\hyperref[0.11.4.3.1-fr]{11.4.3.1}) sont alors
$'E_2^{p0}=\mathrm H^p(T(K^\bullet))$~; autrement dit, cette suite est
\emph{dégénérée}~; comme elle est régulière
(\hyperref[0.11.4.4-fr]{11.4.4}), la conclusion résulte de
(\hyperref[0.11.1.6-fr]{11.1.6}).
\end{proof}

\begin{env}[11.4.6]
\label{0.11.4.6-fr}
Considérons maintenant, par exemple, un \emph{bifoncteur covariant}
$(M,N)\mapsto T(M,N)$ de $\mathcal C\times\mathcal C'$ dans $\mathcal C''$,
$\mathcal C$, $\mathcal C'$, $\mathcal C''$ étant trois catégories
abéliennes~; on suppose, pour simplifier, $T$ additif en chacun de ses
arguments, et en outre que tout objet de $\mathcal C$ et tout objet de
$\mathcal C'$ sont sous-objets d'un objet injectif, et que les limites
inductives filtrantes existent dans $\mathcal C''$ et sont exactes. On définit
alors l'\emph{hypercohomologie de $T$ par rapport à deux complexes}
$K^\bullet$, $K'^\bullet$ de $\mathcal C$ et $\mathcal C'$ respectivement, à
opérateurs de dérivation de degré $+1$, en prenant pour $K^\bullet$ (resp.
$K'^\bullet$) une résolution de Cartan-Eilenberg injective
$L^{\bullet\bullet}$ (resp. $L'^{\bullet\bullet}$)~; alors
$T(L^{\bullet\bullet},L'^{\bullet\bullet})$ est un quadricomplexe de
$\mathcal C''$, que l'on considère comme \emph{bicomplexe} de $\mathcal C''$
en prenant pour degrés de $T(L^{ij},L'^{hk})$ les entiers $i+h$ et $j+k$.
L'\emph{hypercohomologie de $T$ par rapport à $K^\bullet$ et $K'^\bullet$}
est par définition la cohomologie
$\mathrm H^\bullet(T(L^{\bullet\bullet},L'^{\bullet\bullet}))$ de ce
bicomplexe (autrement dit, celle du complexe simple associé) notée
$\mathrm R^\bullet T(K^\bullet,K'^\bullet)$~; elle est l'aboutissement de
deux suites spectrales dont les termes $E_2$ sont donnés par
\oldpage[0\textsubscript{III}]{35}
\[
  \label{0.11.4.6.1-fr}
  'E_2^{pq}=\mathrm H^p(\mathrm R^qT(K^\bullet,K'^\bullet))
  \tag{11.4.6.1}
\]
\[
  \label{0.11.4.6.2-fr}
  ''E_2^{pq}=\sum_{q'+q''=q}
  \mathrm R^pT(\mathrm H^{q'}(K^\bullet),\mathrm H^{q''}(K'^\bullet))
  \tag{11.4.6.2}
\]
(cf. M, XVII, 2).

Ici $\mathrm R^qT(K^\bullet,K'^\bullet)$ est le bicomplexe
$(\mathrm R^qT(K^i,K'^j))_{(i,j)\in\mathbb Z\times\mathbb Z}$ et le second
membre de (\hyperref[0.11.4.6.1-fr]{11.4.6.1}) est sa cohomologie quand on le
considère comme complexe simple.

En outre, la première suite spectrale est toujours régulière, et les deux
suites spectrales sont birégulières lorsqu'il existe $n$ tel que tout objet de
$\mathcal C$ et tout objet de $\mathcal C'$ admettent une résolution injective
de longueur $\leq n$, ou lorsque $K^\bullet$ et $K'^\bullet$ sont limités
inférieurement~; dans ce dernier cas, on peut en outre omettre l'hypothèse que
les limites inductives existent dans $\mathcal C$ et $\mathcal C'$.

Si $K_1^\bullet$, $K_1'^\bullet$ sont deux autres complexes de $\mathcal C$ et
$\mathcal C'$ respectivement, $f$, $g$ deux morphismes homotopes de
$K^\bullet$ dans $K_1^\bullet$, $f'$, $g'$ deux morphismes homotopes de
$K'^\bullet$ dans $K_1'^\bullet$, alors les morphismes
$\mathrm R^\bullet T(K^\bullet,K'^\bullet)\to
\mathrm R^\bullet T(K_1^\bullet,K_1'^\bullet)$ déduits de $f$ et $f'$ d'une
part, de $g$ et $g'$ d'autre part, sont identiques, et il en est de même pour
les morphismes des suites spectrales d'hypercohomologie.

On généralise aisément à un multifoncteur covariant additif quelconque.
\end{env}

\begin{proposition}[11.4.7]
\label{0.11.4.7-fr}
Supposons que pour tout objet injectif $I$ de $\mathcal C$ (resp. $I'$ de
$\mathcal C'$), $A'\mapsto T(I,A')$ (resp. $A\mapsto T(A,I')$) soit un
foncteur exact. Alors, avec les notations de
(\hyperref[0.11.4.6-fr]{11.4.6}), on a des isomorphismes canoniques
\[
  \label{0.11.4.7.1-fr}
  \mathrm R^\bullet T(K^\bullet,K'^\bullet)
  \xrightarrow{\sim}\mathrm H^\bullet(T(L^{\bullet\bullet},K'^\bullet))
  \xrightarrow{\sim}\mathrm H^\bullet(T(K^\bullet,L'^{\bullet\bullet}))
  \tag{11.4.7.1}
\]
où les deux derniers termes sont la cohomologie des complexes simples définis
par les tricomplexes $T(L^{\bullet\bullet},K'^\bullet)$ et
$T(K^\bullet,L'^{\bullet\bullet})$ respectivement.
\end{proposition}

\begin{proof}
Définissons par exemple le premier de ces isomorphismes. Le quadricomplexe
$T(L^{\bullet\bullet},L'^{\bullet\bullet})$ peut être considéré comme un
\emph{bicomplexe}, en prenant comme degrés de $T(L^{ij},L'^{hk})$ les nombres
$i+j+h$ et $k$. Comme pour chaque $h$, $L'^{h,\bullet}$ est une
\emph{résolution} de $K'^h$, on a, pour ce bicomplexe, en vertu de l'hypothèse
sur $T$,
$\mathrm H_{\mathrm{II}}^q(T(L^{\bullet\bullet},L'^{\bullet\bullet}))=0$
pour $q\neq0$ et
$\mathrm H_{\mathrm{II}}^0(T(L^{\bullet\bullet},L'^{\bullet\bullet}))
=T(L^{\bullet\bullet},K'^\bullet)$~; la première suite spectrale de ce
bicomplexe est donc \emph{dégénérée}~; comme $L'^{hk}=0$ pour $k<0$, cette
suite est en outre régulière (\hyperref[0.11.3.3-fr]{11.3.3}), et la
conclusion résulte donc de (\hyperref[0.11.1.6-fr]{11.1.6}).
\end{proof}

On a des résultats analogues pour un multifoncteur covariant en un nombre
quelconque $n$ d'arguments~: dans le calcul de l'hypercohomologie, il n'est pas
nécessaire de remplacer \emph{tous} les complexes par une résolution de
Cartan-Eilenberg, mais seulement $n-1$ d'entre eux, pourvu que, lorsqu'on fixe
$n-1$ arguments quelconques en leur donnant pour valeurs des objets
\emph{injectifs}, le foncteur covariant en l'argument restant soit
\emph{exact}.

\subsection{Passage à la limite inductive dans l'hypercohomologie}
\label{subsection:0.11.5-fr}

\begin{lemma}[11.5.1]
\label{0.11.5.1-fr}
Soit $K^\bullet=(K^i)_{i\in\mathbb Z}$ un complexe de $\mathcal C$, et pour
tout entier $r\in\mathbb Z$, soit $K_{(r)}^\bullet$ le complexe tel que
$K_{(r)}^i=0$ pour $i<r$, $K_{(r)}^i=K^i$ pour $i\geq r$. Soit $T$ un
foncteur covariant additif
\oldpage[0\textsubscript{III}]{36}
de $\mathcal C$ dans $\mathcal C'$, permutant aux
limites inductives (on suppose que les limites inductives filtrantes existent
et sont exactes dans $\mathcal C$ et $\mathcal C'$). Alors
$\mathrm R^\bullet T(K^\bullet)$ est isomorphe à la limite inductive
\[
  \varinjlim_{r\to-\infty}\mathrm R^\bullet T(K_{(r)}^\bullet)
\]
lorsque $r$ tend vers $-\infty$.
\end{lemma}

La construction d'une résolution injective de Cartan-Eilenberg de $K^\bullet$
se fait en choisissant arbitrairement, pour chaque $i$, une résolution
injective $(X_{\mathrm B}^{ij})_{j\geq0}$ de $\mathrm B^i(K^\bullet)$ et une
résolution injective $(X_{\mathrm H}^{ij})_{j\geq0}$ de
$\mathrm H^i(K^\bullet)$~; cela fait, la méthode de construction montre que
la résolution injective $(L^{ij})_{j\geq0}$ de $K^i$ et les opérateurs de
dérivation $L^{i,j}\to L^{i+1,j}$ \emph{ne dépendent que des résolutions}
$(X_{\mathrm B}^{i,\bullet})$, $(X_{\mathrm H}^{i,\bullet})$ et
$(X_{\mathrm B}^{i+1,\bullet})$ (M, XVII, 1.2). Or, il est clair que l'on a
$\mathrm B^i(K_{(r)}^\bullet)=\mathrm B^i(K^\bullet)$ et
$\mathrm H^i(K_{(r)}^\bullet)=\mathrm H^i(K^\bullet)$ pour $i\geq r+1$. On
a, d'autre part, pour chaque $r$, une injection canonique
$\varphi_{r-1,r}^\bullet:K_{(r)}^\bullet\to K_{(r-1)}^\bullet$,
$\varphi_{r-1,r}^i$ étant l'identité pour $i\neq r-1$. La remarque précédente
montre que, si $L^{\bullet\bullet}=(L^{ij})$ est une résolution injective de
Cartan-Eilenberg de $K^\bullet$, on peut pour chaque $r$ définir une résolution
injective de Cartan-Eilenberg
$L_{(r)}^{\bullet\bullet}=(L_{(r)}^{ij})$ de $K_{(r)}^\bullet$ telle que
$L_{(r)}^{ij}=0$ pour $i<r$ et $L_{(r)}^{ij}=L^{ij}$ pour $i\geq r+1$. On peut
d'autre part définir un morphisme de bicomplexes
$\Phi_{r-1,r}^{\bullet\bullet}:L_{(r)}^{\bullet\bullet}\to
L_{(r-1)}^{\bullet\bullet}$ correspondant à $\varphi_{r-1,r}^\bullet$, et la
méthode de définition de ce morphisme (\emph{loc. cit.}) montre encore qu'on
peut le construire de sorte que $\Phi_{r-1,r}^{ij}$ soit l'identité pour
$i\neq r$ et $i\neq r-1$. On a ainsi défini un système inductif
$(L_{(r)}^{\bullet\bullet})$ de bicomplexes de $\mathcal C$, dont
$L^{\bullet\bullet}$ est évidemment la limite inductive lorsque $r$ tend vers
$-\infty$~; en raison de la permutabilité des sommes directes et des limites
inductives, le complexe simple associé à $L^{\bullet\bullet}$ est aussi limite
inductive du complexe simple associé à $L_{(r)}^{\bullet\bullet}$. Comme $T$
permute par hypothèse aux limites inductives et qu'il en est de même de la
cohomologie (en raison de l'exactitude du foncteur $\varinjlim$), on a bien
\[
  \mathrm H^\bullet(T(L^{\bullet\bullet}))=
  \varinjlim\mathrm H^\bullet(T(L_{(r)}^{\bullet\bullet}))
\]
à un isomorphisme près.

Le lemme (\hyperref[0.11.5.1-fr]{11.5.1}) permet d'étendre à des complexes
$K^\bullet$ quelconques, par passage à la limite inductive, des propriétés
valables pour les complexes \emph{limités inférieurement}. Comme premier
exemple, nous prouverons~:

\begin{proposition}[11.5.2]
\label{0.11.5.2-fr}
Sous les hypothèses de (\hyperref[0.11.5.1-fr]{11.5.1}) concernant
$\mathcal C$, $\mathcal C'$ et $T$, $\mathrm R^\bullet T(K^\bullet)$ est un
foncteur cohomologique dans la catégorie abélienne des complexes de
$\mathcal C$.
\end{proposition}

\begin{proof}
Montrons qu'on peut se ramener au cas des complexes limités inférieurement~:
si on a une suite exacte
\[
  0\longrightarrow K'^\bullet\longrightarrow K^\bullet
  \longrightarrow K''^\bullet\longrightarrow0
\]
de complexes, on en déduit évidemment pour chaque $r$ une suite exacte
\[
  0\longrightarrow K_{(r)}'^\bullet\longrightarrow K_{(r)}^\bullet
  \longrightarrow K_{(r)}''^\bullet\longrightarrow0,
\]
d'où par hypothèse, une suite exacte
\[
  \cdots\longrightarrow\mathrm R^nT(K_{(r)}'^\bullet)
  \longrightarrow\mathrm R^nT(K_{(r)}^\bullet)
  \longrightarrow\mathrm R^nT(K_{(r)}''^\bullet)
  \xrightarrow{\partial}\mathrm R^{n+1}T(K_{(r)}'^\bullet)
  \longrightarrow\cdots
\]
ces suites exactes formant un système inductif~; le lemme
(\hyperref[0.11.5.1-fr]{11.5.1}) et l'exactitude du foncteur $\varinjlim$
démontrent que l'on a une suite exacte
\[
  \cdots\longrightarrow\mathrm R^nT(K'^\bullet)
  \longrightarrow\mathrm R^nT(K^\bullet)
  \longrightarrow\mathrm R^nT(K''^\bullet)
  \xrightarrow{\partial}\mathrm R^{n+1}T(K'^\bullet)
  \longrightarrow\cdots
\]
\end{proof}

Pour traiter le cas des complexes limités inférieurement, on peut se borner à
ceux tels que $K^i=0$ pour $i<0$~; ils forment évidemment une catégorie
abélienne $\mathcal K$.

\begin{lemma}[11.5.2.1]
\label{0.11.5.2.1-fr}
Dans $\mathcal K$, soit $\mathfrak J$ l'ensemble des complexes
$Q^\bullet=(Q^i)_{i\geq0}$ ayant les
\oldpage[0\textsubscript{III}]{37}
propriétés suivantes~:
\begin{enumerate}
  \item[$1^\circ$] Tout $Q^i$ est un objet injectif de $\mathcal C$~;
  \item[$2^\circ$] Pour tout $i\geq0$, on a
    $\mathrm Z^i(Q^\bullet)=\mathrm B^i(Q^\bullet)$, et
    $\mathrm Z^i(Q^\bullet)$ est facteur direct de $Q^i$.
\end{enumerate}
Alors~:
\begin{enumerate}
  \item[(i)] Tout $Q^\bullet\in\mathfrak J$ est un objet injectif de
    $\mathcal K$.
  \item[(ii)] Tout objet de $\mathcal K$ est isomorphe à un sous-complexe
    d'un complexe appartenant à $\mathfrak J$.
\end{enumerate}
\end{lemma}

\begin{proof}
\textup{(i)} Soient $A^\bullet=(A^i)$ un objet de $\mathcal K$,
$A'^\bullet=(A'^i)$ un sous-objet de $A^\bullet$,
$Q^\bullet=(Q^i)$ un objet de $\mathfrak J$, et supposons donné un morphisme
$f=(f^i):A'^\bullet\to Q^\bullet$, qu'il s'agit de prolonger en un morphisme
$g=(g^i):A^\bullet\to Q^\bullet$. Nous utiliserons le langage de la catégorie
des modules pour simplifier (cf. [27]).

Identifions $Q^i$ à
$\mathrm B^i(Q^\bullet)\oplus\mathrm B^{i+1}(Q^\bullet)$~; procédons par
récurrence sur $i$, et supposons donc les $g^j$ définis pour $j<i$,
compatibles avec les opérateurs de dérivation $d^j:A^j\to A^{j+1}$ et
$d'^j:Q^j\to Q^{j+1}$ pour $j<i-1$ et tels en outre que~:
\begin{enumerate}
  \item[$1^\circ$] $g^{i-1}(\mathrm Z^{i-1}(A^\bullet))\subset
    \mathrm Z^{i-1}(Q^\bullet)$~;
  \item[$2^\circ$] Si on pose $C^j=(d^j)^{-1}(A'^{j+1})$ pour tout $j$, alors
    $d'^{i-1}\circ g^{i-1}$ coïncide avec $f^i\circ d^{i-1}$ sur $C^{i-1}$.
\end{enumerate}
Le morphisme $f^i:A'^i\to Q^i$ donne par composition avec les projections
deux morphismes
$f'^i:A'^i\to\mathrm B^i(Q^\bullet)$ et
$f''^i:A'^i\to\mathrm B^{i+1}(Q^\bullet)$. Comme
$d'^{i-1}\circ g^{i-1}$ applique $A^{i-1}$ dans
$\mathrm B^i(Q^\bullet)$ et s'annule dans $\mathrm Z^{i-1}(A^\bullet)$, il
définit un morphisme
$h^i:\mathrm B^i(A^\bullet)\to\mathrm B^i(Q^\bullet)$, et puisque
$d'^{i-1}\circ g^{i-1}$ coïncide avec $f^i\circ d^{i-1}$ sur $C^{i-1}$,
$h^i$ coïncide avec $f'^i$ dans
$\mathrm B^i(A^\bullet)\cap A'^i$. Comme $\mathrm B^i(Q^\bullet)$, facteur
direct de $Q^i$, est injectif, il y a un morphisme
$g'^i:A^i\to\mathrm B^i(Q^\bullet)$ qui coïncide avec $h^i$ dans
$\mathrm B^i(A^\bullet)$ et avec $f'^i$ dans $A'^i$. Considérons d'autre part
le morphisme
$f''^{i+1}\circ d^i:C^i\to\mathrm B^{i+1}(Q^\bullet)$, qui s'annule dans
$\mathrm Z^i(A^\bullet)$~; comme $\mathrm B^{i+1}(Q^\bullet)$ est injectif,
il y a un morphisme $g''^i:A^i\to\mathrm B^{i+1}(Q^\bullet)$, qui coïncide
avec $f''^{i+1}\circ d^i$ dans $C^i$ et avec $0$ dans
$\mathrm Z^i(A^\bullet)$. Il suffit alors de prendre $g^i=g'^i+g''^i$ pour
pouvoir poursuivre la récurrence.

\textup{(ii)} Pour plonger $A^\bullet=(A^i)$ dans un complexe appartenant à
$\mathfrak J$, on prend pour chaque $i\geq1$ un objet injectif $Q'^i$ de
$\mathcal C$ tel qu'il existe une injection $f'^i:A^i\to Q'^i$. Posons alors
\[
  Q^i=0\quad(i<0),\qquad Q^0=Q'^1,\qquad
  Q^i=Q'^i\oplus Q'^{i+1}\quad(i\geq1),
\]
avec l'opérateur de dérivation évident. Posons
$f^i=f'^i+(f'^{i+1}\circ d^i)$ pour tout $i$ (avec $f'^i=0$ pour $i\leq0$)~;
il est immédiat que $f^i$ est injectif pour tout $i$, et que les $f^i$ sont
compatibles avec les opérateurs de dérivation.
\end{proof}

\begin{corollary}[11.5.2.2]
\label{0.11.5.2.2-fr}
Tout objet $K^\bullet$ de $\mathcal K$ admet une résolution droite formée
d'objets de $\mathfrak J$. Si $L^{\bullet\bullet}$ est une telle résolution,
pour toute résolution $L'^{\bullet\bullet}$ de $K^\bullet$, formée d'objets
de $\mathcal K$, il y a un morphisme de bicomplexes
$F:L'^{\bullet\bullet}\to L^{\bullet\bullet}$ compatible avec les
augmentations, et deux tels morphismes $F$, $F'$ sont homotopes.
\end{corollary}

\begin{proof}
Ce n'est autre que (M, V, 1.1 a)) appliqué à la catégorie abélienne
$\mathcal K$.
\end{proof}

\begin{env}[11.5.2.3]
\label{0.11.5.2.3-fr}
Ces préliminaires étant posés, considérons une résolution de
Cartan-Eilenberg injective $L'^{\bullet\bullet}$ de $K^\bullet$ et une
résolution $L^{\bullet\bullet}$ de $K^\bullet$ formée d'objets de
$\mathfrak J$, et montrons que l'on a un isomorphisme
\[
  \mathrm H^\bullet(T(L'^{\bullet\bullet}))\xrightarrow{\sim}
  \mathrm H^\bullet(T(L^{\bullet\bullet})).
\]
On déduit en effet de (\hyperref[0.11.5.2.2-fr]{11.5.2.2}) un morphisme de
bicomplexes $T(L'^{\bullet\bullet})\to T(L^{\bullet\bullet})$, et par suite
un morphisme
$'E(T(L'^{\bullet\bullet}))\to{}'E(T(L^{\bullet\bullet}))$ des premières
suites spectrales de ces bicomplexes. Comme en vertu de
(\hyperref[0.11.3.3-fr]{11.3.3}), ces suites spectrales sont régulières, il
suffit (\hyperref[0.11.1.5-fr]{11.1.5}) de voir que le morphisme précédent est
un isomorphisme pour les termes $E_2$, ou encore que
$\mathrm H_{\mathrm{II}}^q(T(L^{i,\bullet}))$ est égal à
$\mathrm R^qT(K^i)$~; comme $L^{i,\bullet}$ est une résolution droite de
$K^i$, on est ramené à prouver le
\end{env}

\oldpage[0\textsubscript{III}]{38}
\begin{lemma}[11.5.2.4]
\label{0.11.5.2.4-fr}
Si $L^\bullet=(L^i)_{i\in\mathbb Z}$ est une résolution droite d'un objet $A$
de $\mathcal C$ telle que $\mathrm R^nT(L^i)=0$ pour tout $i\in\mathbb Z$ et
tout $n>0$, alors on a
$\mathrm H^\bullet T(L^\bullet)=\mathrm R^\bullet T(A)$.
\end{lemma}

\begin{proof}
C'est un cas particulier de (T, 2.5.2).
\end{proof}

\begin{env}[11.5.2.5]
\label{0.11.5.2.5-fr}
La démonstration de (\hyperref[0.11.5.2-fr]{11.5.2}) se termine maintenant de
façon immédiate, car $L'^{\bullet\bullet}$ est une résolution injective de
$K^\bullet$ dans la catégorie abélienne $\mathcal K$, autrement dit,
$K^\bullet\mapsto\mathrm H^\bullet(T(L'^{\bullet\bullet}))$ n'est autre que
le foncteur cohomologique dérivé droit de $T$ dans la catégorie $\mathcal K$
(T, 2.3).
\end{env}

\begin{proposition}[11.5.3]
\label{0.11.5.3-fr}
Sous les hypothèses de (\hyperref[0.11.5.1-fr]{11.5.1}) concernant
$\mathcal C$, $\mathcal C'$ et $T$, soit
$L^{\bullet\bullet}=(L^{ij})$ un bicomplexe tel que $L^{ij}=0$ pour $j<0$ et
que, pour tout $i$, $L^{i,\bullet}$ soit une résolution de $K^i$~; supposons
enfin que $\mathrm R^nT(L^{ij})=0$ pour tout couple $(i,j)$ et tout $n>0$.
Alors il existe un isomorphisme fonctoriel
\[
  \label{0.11.5.3.1-fr}
  \mathrm R^\bullet T(K^\bullet)\xrightarrow{\sim}
  \mathrm H^\bullet(T(L^{\bullet\bullet})).
  \tag{11.5.3.1}
\]
\end{proposition}

\begin{proof}
Soit $L_{(r)}^{\bullet\bullet}=(L_{(r)}^{ij})$ le bicomplexe tel que
$L_{(r)}^{ij}=0$ pour $i<r$, $L_{(r)}^{ij}=L^{ij}$ pour $i\geq r$~; il est
immédiat que $L^{\bullet\bullet}$ est limite inductive de
$L_{(r)}^{\bullet\bullet}$ pour $r$ tendant vers $-\infty$~; en vertu de
l'hypothèse sur $T$ et de (\hyperref[0.11.5.1-fr]{11.5.1}), il suffit donc de
prouver la proposition lorsque $K^\bullet$ est limité inférieurement, par
exemple $K^i=0$ pour $i<0$, et $L^{ij}=0$ pour $i<0$. Soit alors
$L'^{\bullet\bullet}=(L'^{ij})$ une résolution droite de $K^\bullet$ formée
d'objets de $\mathfrak J$ (\hyperref[0.11.5.2.2-fr]{11.5.2.2})~; il y a un
morphisme de bicomplexes
$L^{\bullet\bullet}\to L'^{\bullet\bullet}$ compatible avec les
augmentations, d'où un morphisme
$'E(T(L^{\bullet\bullet}))\to{}'E(T(L'^{\bullet\bullet}))$ pour les premières
suites spectrales~; le lemme (\hyperref[0.11.5.2.4-fr]{11.5.2.4}) montre,
comme dans (\hyperref[0.11.5.2.3-fr]{11.5.2.3}), que ce morphisme est un
\emph{isomorphisme}, d'où la conclusion.
\end{proof}

\begin{remark}[11.5.4]
\label{0.11.5.4-fr}
Les raisonnements précédents prouvent que les conclusions de
(\hyperref[0.11.5.2-fr]{11.5.2}) et
(\hyperref[0.11.5.3-fr]{11.5.3}) sont valables dans la catégorie des
complexes $K^\bullet$ limités inférieurement, \emph{sans supposer que $T$
permute aux limites inductives filtrantes}. En outre, quand on considère
seulement la catégorie $\mathcal K$ des complexes $K^\bullet$ tels que
$K^i=0$ pour $i<0$, la caractérisation de $\mathrm R^\bullet T(K^\bullet)$
comme le système des \emph{foncteurs dérivés droits} de $T$ dans $\mathcal K$
montre que ce foncteur cohomologique est \emph{universel} (T, 2.3).

Un autre cas où (\hyperref[0.11.5.2-fr]{11.5.2}) est valable sans faire
d'hypothèse supplémentaire sur $T$ est le cas où il existe un entier $m>0$
tel que tout objet de $\mathcal C$ admette une résolution injective de longueur
$\leq m$. En effet, dans la démonstration de
(\hyperref[0.11.5.1-fr]{11.5.1}), toutes les résolutions injectives d'objets
de $\mathcal C$ qui interviennent peuvent être prises de longueur $\leq m$,
d'où il résulte aussitôt que les termes de degré total $n$ du bicomplexe
$T(L_{(r)}^{\bullet\bullet})$ sont égaux à ceux de
$T(L^{\bullet\bullet})$ et en nombre fini, dès que $r$ est assez grand, ce qui
entraîne que pour tout $n$,
$\mathrm H^n(T(L^{\bullet\bullet}))=
\mathrm H^n(T(L_{(r)}^{\bullet\bullet}))$ dès que $r$ est assez grand. Avec
les notations de (\hyperref[0.11.5.2-fr]{11.5.2}), on a donc aussi
$\mathrm R^nT(K_{(r)}^\bullet)=\mathrm R^nT(K^\bullet)$ pour $r$ assez grand
(dépendant de $n$) et de même pour $K'^\bullet$ et $K''^\bullet$, d'où la
conclusion. De la même manière, (\hyperref[0.11.5.3-fr]{11.5.3}) est valable
sans condition supplémentaire sur $T$ lorsque $\mathcal C$ vérifie
l'hypothèse précédente et que l'on suppose que les résolutions
$L^{i,\bullet}$ sont de longueur $\leq m$.
\end{remark}

\begin{env}[11.5.5]
\label{0.11.5.5-fr}
Le résultat de (\hyperref[0.11.5.2-fr]{11.5.2}) se généralise aux
multifoncteurs covariants. Considérons par exemple la situation de
(\hyperref[0.11.4.6-fr]{11.4.6}), où l'on suppose que dans $\mathcal C$,
$\mathcal C'$ et $\mathcal C''$
\oldpage[0\textsubscript{III}]{39}
les limites inductives filtrantes existent et sont exactes, et que $T$ commute
aux limites inductives. Alors $\mathrm R^\bullet T(K^\bullet,K'^\bullet)$ est
un bifoncteur cohomologique en chacun des complexes $K^\bullet$,
$K'^\bullet$~; pour le voir, on se ramène comme dans
(\hyperref[0.11.5.2-fr]{11.5.2}) au cas où $K^\bullet$ et $K'^\bullet$ sont
limités inférieurement~; prenant alors pour $K^\bullet$ et $K'^\bullet$ des
résolutions injectives du type décrit dans
(\hyperref[0.11.5.2.2-fr]{11.5.2.2}), on est ramené à la propriété générale
démontrée dans (M, V, 4.1).
\end{env}

\begin{env}[11.5.6]
\label{0.11.5.6-fr}
De même, les résultats de (\hyperref[0.11.4.7-fr]{11.4.7}) et
(\hyperref[0.11.5.3-fr]{11.5.3}) se généralisent comme suit. Supposons (sous
les hypothèses de (\hyperref[0.11.5.5-fr]{11.5.5})) que l'on ait deux
bicomplexes $L^{\bullet\bullet}=(L^{ij})$,
$L'^{\bullet\bullet}=(L'^{ij})$ tels que $L^{ij}=0$ et $L'^{ij}=0$ pour
$j<0$, que pour tout $i$, $L^{i,\bullet}$ soit une résolution de $K^i$ et
$L'^{i,\bullet}$ une résolution de $K'^i$, et enfin que
$\mathrm R^nT(L^{ij},L'^{hk})=0$ pour $n>0$ et pour tout système d'indices
$(i,j,h,k)$. Alors on a un isomorphisme fonctoriel en $K^\bullet$ et
$K'^\bullet$
\[
  \label{0.11.5.6.1-fr}
  \mathrm R^\bullet T(K^\bullet,K'^\bullet)\xrightarrow{\sim}
  \mathrm H^\bullet(T(L^{\bullet\bullet},L'^{\bullet\bullet})).
  \tag{11.5.6.1}
\]
Cela s'établit comme dans (\hyperref[0.11.5.3-fr]{11.5.3}) en se ramenant au
cas où $K^\bullet$ et $K'^\bullet$ sont limités inférieurement.

Supposons en outre que pour tout couple $(i,j)$ et pour tout couple $(h,k)$,
les foncteurs $A\mapsto T(A,L'^{hk})$ et $A'\mapsto T(L^{ij},A')$ soient
exacts dans $\mathcal C$ et $\mathcal C'$ respectivement. Alors on a aussi des
isomorphismes fonctoriels
\[
  \label{0.11.5.6.2-fr}
  \mathrm R^\bullet T(K^\bullet,K'^\bullet)
  \xrightarrow{\sim}\mathrm H^\bullet(T(L^{\bullet\bullet},K'^\bullet))
  \xrightarrow{\sim}\mathrm H^\bullet(T(K^\bullet,L'^{\bullet\bullet})).
  \tag{11.5.6.2}
\]
La démonstration est semblable à celle de
(\hyperref[0.11.4.7-fr]{11.4.7}).
\end{env}

\begin{env}[11.5.7]
\label{0.11.5.7-fr}
On notera encore que les résultats de (\hyperref[0.11.5.5-fr]{11.5.5}) et
(\hyperref[0.11.5.6-fr]{11.5.6}) sont valables sans supposer que $T$ permute
aux limites inductives, pourvu que l'on se restreigne aux complexes
$K^\bullet$, $K'^\bullet$ limités inférieurement, ou que l'on suppose que tout
objet de $\mathcal C$ (resp. $\mathcal C'$) admet une résolution injective de
longueur bornée, et que dans (\hyperref[0.11.5.6-fr]{11.5.6}) les bicomplexes
$L^{\bullet\bullet}$ et $L'^{\bullet\bullet}$ ont leur second degré limité
supérieurement.
\end{env}

\subsection{Hyperhomologie d'un foncteur par rapport à un complexe
$K_\bullet$}
\label{subsection:0.11.6-fr}

\begin{env}[11.6.1]
\label{0.11.6.1-fr}
Soient $\mathcal C$ une catégorie abélienne,
$K_\bullet=(K_i)_{i\in\mathbb Z}$ un complexe d'objets de $\mathcal C$, dont
l'opérateur de dérivation est de degré $-1$. Une \emph{résolution de
Cartan-Eilenberg gauche} de $K_\bullet$ est formée d'un bicomplexe
$L_{\bullet\bullet}=(L_{ij})$ à opérateurs de dérivation de degré $-1$ avec
$L_{ij}=0$ pour $j<0$, et d'un morphisme de complexes simples
$\varepsilon:L_{\bullet,0}\to K_\bullet$, de façon que les conditions obtenues
à partir de celles de (\hyperref[0.11.4.2-fr]{11.4.2}) par «~renversement des
flèches~» soient remplies. Si tout objet de $\mathcal C$ est quotient d'un
objet projectif, tout complexe $K_\bullet$ de $\mathcal C$ admet une
résolution de Cartan-Eilenberg \emph{projective}, c'est-à-dire formée d'objets
projectifs $L_{ij}$, avec des propriétés fonctorielles semblables à celles de
(\hyperref[0.11.4.2-fr]{11.4.2}). En outre, si $K_\bullet$ est limité
inférieurement (resp. supérieurement), on peut supposer que $L_{ij}=0$ pour
$i<i_0$ (resp. $i>i_0$) si $K_i=0$ pour $i<i_0$ (resp. $i>i_0$). Si tout
objet de $\mathcal C$ admet une résolution projective de longueur $\leq n$,
on peut supposer que $L_{ij}=0$ pour $j>n$.
\end{env}

\begin{env}[11.6.2]
\label{0.11.6.2-fr}
\oldpage[0\textsubscript{III}]{40}
Supposons que $T$ soit un foncteur covariant additif de $\mathcal C$ dans une
catégorie abélienne $\mathcal C'$. La définition de
\emph{l'hyperhomologie} $L_\bullet T(K_\bullet)$ et des \emph{suites
spectrales d'hyperhomologie} de $T$ par rapport à un complexe $K_\bullet$ de
$\mathcal C$ (lorsqu'elles existent) se fait encore à partir de
(\hyperref[0.11.4.3-fr]{11.4.3}) par «~renversement des flèches~», les termes
$E_2$ des deux suites spectrales ainsi obtenues étant
\[
\label{0.11.6.2.1-fr}
\tag{11.6.2.1}
{}'E^2_{pq}=H_p(L_qT(K_\bullet))
\]
\[
\label{0.11.6.2.2-fr}
\tag{11.6.2.2}
{}''E^2_{pq}=L_pT(H_q(K_\bullet)),
\]
où $L_pT$ désigne le $p$-ième dérivé gauche de $T$ pour $p\geq 0$, et $0$
pour $p<0$, $L_qT(K_\bullet)$ le complexe
$(L_qT(K_i))_{i\in\mathbb Z}$.

Les propriétés de l'hyperhomologie ne se déduisent pas toutes par simple
«~renversement des flèches~» de celles de l'hypercohomologie (à moins qu'on
ne fasse des hypothèses supplémentaires du type (AB~5*) de (T, 1.5) sur la
catégorie $\mathcal C'$) en raison des conditions de régularité des deux
suites spectrales précédentes, auxquelles il faut cette fois appliquer les
critères de (\hyperref[0.11.3.4-fr]{11.3.4}). Ces derniers montrent que
lorsqu'on suppose que dans $\mathcal C'$ les limites inductives filtrantes
existent et sont exactes, alors le complexe défini par le bicomplexe
$T(L_{\bullet\bullet})$ existe, et la seconde suite spectrale
$\,{}''E(T(L_{\bullet\bullet}))$ est cette fois régulière. Si l'on suppose,
soit que $K_\bullet$ soit limité inférieurement, soit qu'il existe un entier
$n$ tel que tout objet de $\mathcal C$ admette une résolution projective de
longueur $\leq n$, alors les deux suites spectrales d'hyperhomologie existent
(sans hypothèse sur $\mathcal C'$) et sont birégulières.
\end{env}

\begin{proposition}[11.6.3]
\label{0.11.6.3-fr}
Soient $\mathcal C$, $\mathcal C'$ deux catégories abéliennes, $T$ un foncteur
covariant additif de $\mathcal C$ dans $\mathcal C'$. Alors :
\begin{enumerate}
\item[{\rm(i)}] \emph{L'hyperhomologie $L_\bullet T(K_\bullet)$ est un
foncteur homologique dans la catégorie abélienne des complexes de
$\mathcal C$ limités inférieurement.}
\item[{\rm(ii)}] \emph{Soit $K_\bullet$ un complexe de $\mathcal C$ limité
inférieurement. Si $L_nT(K_i)=0$ pour tout $n>0$ et tout $i\in\mathbb Z$, on
a des isomorphismes fonctoriels}
\[
\label{0.11.6.3.1-fr}
\tag{11.6.3.1}
L_iT(K_\bullet)\simeq H_i(T(K_\bullet))\qquad(i\in\mathbb Z).
\]
\item[{\rm(iii)}] \emph{Soit $K_\bullet$ un complexe de $\mathcal C$ limité
inférieurement. Soit $L_{\bullet\bullet}=(L_{ij})$ un bicomplexe tel que
$L_{ij}=0$ pour $j<0$ et que, pour tout $i$, $L_{i,\bullet}$ soit une
résolution de $K_i$; supposons enfin que $L_nT(L_{ij})=0$ pour tout couple
$(i,j)$ et tout $n>0$. Alors il existe un isomorphisme fonctoriel}
\[
\label{0.11.6.3.2-fr}
\tag{11.6.3.2}
L_\bullet T(K_\bullet)\simeq H_\bullet(T(L_{\bullet\bullet})).
\]
\end{enumerate}
Les démonstrations se font comme celles de
(\hyperref[0.11.5.2-fr]{11.5.2}), (\hyperref[0.11.4.5-fr]{11.4.5}) et
(\hyperref[0.11.5.3-fr]{11.5.3}) dans le cas des complexes limités
inférieurement. Nous laissons les détails de ces raisonnements au lecteur.
\end{proposition}

\begin{env}[11.6.4]
\label{0.11.6.4-fr}
On a des résultats tout à fait analogues pour les multifoncteurs covariants
additifs en chacun des arguments. Par exemple pour un bifoncteur $T$, on a
les deux suites spectrales d'hyperhomologie de termes $E_2$ donnés par
\oldpage[0\textsubscript{III}]{41}
\[
\label{0.11.6.4.1-fr}
\tag{11.6.4.1}
{}'E^2_{pq}=H_p(L_qT(K_\bullet,K'_\bullet))
\]
\[
\label{0.11.6.4.2-fr}
\tag{11.6.4.2}
{}''E^2_{pq}=
\sum_{q'+q''=q}L_pT(H_{q'}(K_\bullet),H_{q''}(K'_\bullet)).
\]
Ici encore, c'est la seconde suite spectrale qui est régulière, les deux
suites étant birégulières lorsqu'il s'agit de complexes $K_\bullet$,
$K'_\bullet$ limités inférieurement, ou quand les objets des catégories
abéliennes que l'on considère ont des résolutions projectives de longueur
fixée.

En outre :
\end{env}

\begin{proposition}[11.6.5]
\label{0.11.6.5-fr}
Soient $\mathcal C$, $\mathcal C'$, $\mathcal C''$ trois catégories
abéliennes, $T$ un bifoncteur covariant biadditif de
$\mathcal C\times\mathcal C'$ dans $\mathcal C''$.
\begin{enumerate}
\item[{\rm(i)}] \emph{$L_\bullet T(K_\bullet,K'_\bullet)$ est un bifoncteur
homologique en chacun des complexes limités inférieurement $K_\bullet$,
$K'_\bullet$ (formés respectivement d'objets de $\mathcal C$ et d'objets de
$\mathcal C'$).}
\item[{\rm(ii)}] \emph{Supposons $K_\bullet$ et $K'_\bullet$ limités
inférieurement. Soient $L_{\bullet\bullet}=(L_{ij})$,
$L'_{\bullet\bullet}=(L'_{ij})$ deux bicomplexes tels que $L_{ij}=0$ et
$L'_{ij}=0$ pour $j<0$, que pour tout $i$, $L_{i,\bullet}$ soit une
résolution de $K_i$ et $L'_{i,\bullet}$ une résolution de $K'_i$ et enfin que
$L_nT(L_{ij},L'_{hk})=0$ pour $n>0$ et tout système $(i,j,h,k)$. Alors on a
un isomorphisme fonctoriel}
\[
\label{0.11.6.5.1-fr}
\tag{11.6.5.1}
L_\bullet T(K_\bullet,K'_\bullet)
\simeq H_\bullet(T(L_{\bullet\bullet},L'_{\bullet\bullet})).
\]
\item[{\rm(iii)}] \emph{Supposons en outre que pour tout couple $(i,j)$ et
tout couple $(h,k)$, les foncteurs
\[
A\longmapsto T(A,L'_{hk}),\qquad
A'\longmapsto T(L_{ij},A')
\]
soient exacts dans $\mathcal C$ et $\mathcal C'$ respectivement. Alors on a
des isomorphismes fonctoriels}
\[
\label{0.11.6.5.2-fr}
\tag{11.6.5.2}
L_\bullet T(K_\bullet,K'_\bullet)
\simeq H_\bullet(T(L_{\bullet\bullet},K'_\bullet))
\simeq H_\bullet(T(K_\bullet,L'_{\bullet\bullet})).
\]
\end{enumerate}
Les démonstrations sont analogues à celles de
(\hyperref[0.11.5.5-fr]{11.5.5}) et (\hyperref[0.11.5.6-fr]{11.5.6}).
\end{proposition}

\subsection{Hyperhomologie d'un foncteur par rapport à un bicomplexe
$K_{\bullet\bullet}$}
\label{subsection:0.11.7-fr}

\begin{env}[11.7.1]
\label{0.11.7.1-fr}
Soit $\mathcal C$ une catégorie abélienne dans laquelle tout objet est
quotient d'un objet projectif. Considérons un bicomplexe
$K_{\bullet\bullet}=(K_{ij})$ formé d'objets de $\mathcal C$, et dont les deux
degrés sont limités inférieurement; on peut toujours se borner au cas où
$K_{ij}=0$ pour $i<0$ ou $j<0$, et c'est ce que nous ferons désormais. On
peut considérer $K_{\bullet\bullet}$ comme un complexe (simple) formé
d'objets $K_{i,\bullet}=(K_{ij})_{j\geq 0}$ de la catégorie abélienne
$\mathcal K$ des complexes à degrés positifs d'objets de $\mathcal C$. Il
résulte du lemme (\hyperref[0.11.5.2.1-fr]{11.5.2.1}) (ou plutôt du lemme
«~dual~» obtenu par «~renversement des flèches~») et de (M, V, 2.2) que
$K_{\bullet\bullet}$ admet une \emph{résolution projective de
Cartan-Eilenberg} dans la catégorie $\mathcal K$; une telle résolution est un
tricomplexe $M_{\bullet\bullet\bullet}=(M_{ijk})$ de $\mathcal C$, à degrés
tous $\geq 0$, formé d'objets projectifs, tel que pour tout $i$,
$M_{i,\bullet\bullet}$, $\mathrm B_i^{\mathrm I}(M_{\bullet\bullet\bullet})$,
$\mathrm Z_i^{\mathrm I}(M_{\bullet\bullet\bullet})$,
$H_i^{\mathrm I}(M_{\bullet\bullet\bullet})$ constituent des résolutions
projectives de $K_{i,\bullet}$,
$\mathrm B_i^{\mathrm I}(K_{\bullet\bullet})$,
$\mathrm Z_i^{\mathrm I}(K_{\bullet\bullet})$,
$H_i^{\mathrm I}(K_{\bullet\bullet})$ respectivement dans la catégorie
$\mathcal K$; en particulier, pour tout couple $(i,j)$,
$M_{ij,\bullet}$ est une résolution projective de $K_{ij}$ dans
$\mathcal C$.
\end{env}

\begin{proposition}[11.7.2]
\label{0.11.7.2-fr}
Soit $T$ un foncteur covariant additif de $\mathcal C$ dans une catégorie
abélienne $\mathcal C'$. Avec les notations de
(\hyperref[0.11.7.1-fr]{11.7.1}), \emph{l'homologie
$H_\bullet(T(M_{\bullet\bullet\bullet}))$ du complexe simple associé au
tricomplexe $T(M_{\bullet\bullet\bullet})$ est canoniquement isomorphe à
l'hyperhomologie $L_\bullet T(K_{\bullet\bullet})$ du complexe
\oldpage[0\textsubscript{III}]{42}
simple associé à $K_{\bullet\bullet}$
(\hyperref[0.11.6.2-fr]{11.6.2}), et est
l'aboutissement commun de six suites spectrales birégulières notées
$\mathop{}^{(t)}E$ (avec $t=a,b,a',b',c$ ou $d$), dont les termes $E_2$ sont
donnés par les formules}
\[
\begin{aligned}
{}^{(a)}E^2_{pq}&=L_pT(H_q(K_{\bullet\bullet})),&
{}^{(b)}E^2_{pq}&=H_p(L_q^{\mathrm{II}}T(K_{\bullet\bullet})),\\
{}^{(a')}E^2_{pq}&=L_pT(H_q^{\mathrm I}(K_{\bullet\bullet})),&
{}^{(b')}E^2_{pq}&=H_p(L_qT(K_{\bullet\bullet})),\\
{}^{(c)}E^2_{pq}&=L_pT(H_q^{\mathrm{II}}(K_{\bullet\bullet})),&
{}^{(d)}E^2_{pq}&=H_p(L_q^{\mathrm I}T(K_{\bullet\bullet})).
\end{aligned}
\]
\end{proposition}

On rappelle que nous utilisons la notation $F(A_\bullet)$ pour désigner le
complexe des objets $F(A_i)$ pour tout complexe $A_\bullet=(A_i)$; par
exemple $L_q^{\mathrm{II}}T(K_{\bullet\bullet})$ désigne le complexe
$(L_q^{\mathrm{II}}T(K_{i,\bullet}))_{i\geq0}$, où
$L_q^{\mathrm{II}}T(K_{i,\bullet})$ est l'hyperhomologie d'indice $q$ du
foncteur $T$ par rapport au complexe simple $K_{i,\bullet}$.

\begin{proof}
Désignons par $L_\bullet$ le complexe simple associé à
$K_{\bullet\bullet}$, de sorte que
\[
L_i=\bigoplus_{r+s=i}K_{rs},
\]
et posons
\[
N_{ij}=\bigoplus_{r+s=i}M_{rsj};
\]
il est clair que pour chaque $i$, $N_{i,\bullet}$ est une résolution
projective de $L_i$ dans $\mathcal C$; il résulte donc de
(\hyperref[0.11.6.3-fr]{11.6.3}) et
(\hyperref[0.11.6.4-fr]{11.6.4}) que l'on a un isomorphisme fonctoriel
\[
L_\bullet T(L_\bullet)\simeq H_\bullet(T(N_{\bullet\bullet}));
\]
comme le complexe simple associé au bicomplexe $T(N_{\bullet\bullet})$ est
aussi associé au tricomplexe $T(M_{\bullet\bullet\bullet})$, cela prouve la
première assertion de l'énoncé.

En outre, $L_\bullet T(L_\bullet)$ est l'aboutissement des deux suites
spectrales d'hyperhomologie (\hyperref[0.11.6.2-fr]{11.6.2}) de $T$
relativement au complexe simple $L_\bullet$, qui ne sont autres que les suites
$\mathop{}^{(b')}E$ et $\mathop{}^{(a)}E$ respectivement.

Considérons maintenant $M_{\bullet\bullet\bullet}$ comme un bicomplexe
$U_{\bullet\bullet}$ avec
\[
U_{ij}=\bigoplus_{r+s=j}M_{irs};
\]
$H_\bullet(T(M_{\bullet\bullet\bullet}))$ est encore l'aboutissement des deux
suites spectrales du bicomplexe $T(U_{\bullet\bullet})$. Or, pour tout $i$,
$M_{i,\bullet\bullet}$ est un bicomplexe vérifiant les conditions de
(\hyperref[0.11.6.3-fr]{11.6.3}) relativement au complexe simple
$K_{i,\bullet}$; donc
$H_q^{\mathrm{II}}(T(U_{i,\bullet}))=
L_q^{\mathrm{II}}T(K_{i,\bullet})$, et la première suite spectrale de
$T(U_{\bullet\bullet})$ n'est autre que $\mathop{}^{(b)}E$. D'autre part, pour
tout $r$, $M_{\bullet,r,\bullet}$ est une résolution de Cartan-Eilenberg du
complexe simple $K_{\bullet,r}$, le calcul fait dans (M, XV, 2) montre que
\[
H_q^{\mathrm I}(T(M_{\bullet,rs}))
=T(H_q^{\mathrm I}(M_{\bullet,rs})),
\]
d'où
\[
H_q^{\mathrm I}(T(U_{\bullet,j}))
=\bigoplus_{r+s=j}T(H_q^{\mathrm I}(M_{\bullet,rs}));
\]
autrement dit, le complexe simple
$H_q^{\mathrm I}(T(U_{\bullet\bullet}))$ n'est autre que le complexe simple
associé au bicomplexe
$T(H_q^{\mathrm I}(M_{\bullet\bullet\bullet}))$. Or, pour tout $q$,
$H_q^{\mathrm I}(M_{\bullet\bullet\bullet})$ est une résolution projective du
complexe simple $H_q^{\mathrm I}(K_{\bullet\bullet})$ dans la catégorie
$\mathcal K$; appliquant (\hyperref[0.11.6.3-fr]{11.6.3}), on voit que l'on a
\[
H_p^{\mathrm{II}}
\bigl(T(H_q^{\mathrm I}(M_{\bullet\bullet\bullet}))\bigr)
=L_pT(H_q^{\mathrm I}(K_{\bullet\bullet})),
\]
et on obtient ainsi la suite $\mathop{}^{(a')}E$. Enfin, les suites
$\mathop{}^{(c)}E$ et $\mathop{}^{(d)}E$ s'obtiennent en intervertissant les
rôles des indices $i$ et $j$ dans la définition du tricomplexe
$M_{\bullet\bullet\bullet}$ et en appliquant à ce nouveau tricomplexe les
raisonnements qui précèdent.
\end{proof}

On dit que $L_\bullet T(K_{\bullet\bullet})$ est \emph{l'hyperhomologie} de
$T$ relative au bicomplexe $K_{\bullet\bullet}$.

\begin{env}[11.7.3]
\label{0.11.7.3-fr}
\emph{Remarques}.
\begin{enumerate}
\item[{\rm(i)}] Il résulte de (\hyperref[0.11.6.3-fr]{11.6.3}) que
$L_\bullet T(K_{\bullet\bullet})$ est un foncteur homologique dans la
catégorie des bicomplexes $K_{\bullet\bullet}$ de $\mathcal C$ limités
inférieurement en chacun de leurs degrés.
\oldpage[0\textsubscript{III}]{43}
\item[{\rm(ii)}] Soit $M_{\bullet\bullet\bullet}$ un tricomplexe de
$\mathcal C$ tel que pour chaque couple $(r,s)$,
$M_{rs,\bullet}$ soit une résolution de $K_{rs}$ et que
$L_nT(M_{ijk})=0$ pour tous les triplets $(i,j,k)$ et tout $n>0$. Alors on a
un isomorphisme
\[
L_\bullet T(K_{\bullet\bullet})
\simeq H_\bullet(T(M_{\bullet\bullet\bullet}));
\]
en effet, avec les notations de la démonstration de
(\hyperref[0.11.7.2-fr]{11.7.2}), $N_{i,\bullet}$ est une résolution de
$L_i$ telle que $L_nT(N_{ij})=0$ pour $n>0$ et tout couple $(i,j)$, et il
suffit d'appliquer (\hyperref[0.11.6.3-fr]{11.6.3}, (iii)).
\item[{\rm(iii)}] On généralise aussitôt les résultats de
(\hyperref[0.11.7.2-fr]{11.7.2}) aux multifoncteurs covariants; par exemple,
soient $\mathcal C'$ une seconde catégorie abélienne dans laquelle tout objet
est quotient d'un objet projectif, $K'_{\bullet\bullet}$ un bicomplexe de
$\mathcal C'$ dont les deux degrés sont limités inférieurement, et $T$ un
bifoncteur covariant additif de $\mathcal C\times\mathcal C'$ dans une
catégorie abélienne $\mathcal C''$. Si $L_\bullet$ et $L'_\bullet$ sont les
complexes simples associés à $K_{\bullet\bullet}$ et
$K'_{\bullet\bullet}$ respectivement, on définit l'hyperhomologie de $T$ par
rapport aux deux bicomplexes $K_{\bullet\bullet}$,
$K'_{\bullet\bullet}$ comme l'hyperhomologie
$L_\bullet T(L_\bullet,L'_\bullet)$; appliquant
(\hyperref[0.11.6.4-fr]{11.6.4}) et
(\hyperref[0.11.6.5-fr]{11.6.5}), on a de même que dans
(\hyperref[0.11.7.2-fr]{11.7.2}) six suites spectrales aboutissant à cette
hyperhomologie, et que nous laissons le soin d'écrire au lecteur.
\end{enumerate}
\end{env}

\subsection{Compléments sur la cohomologie des complexes simpliciaux}
\label{subsection:0.11.8-fr}

\begin{env}[11.8.1]
\label{0.11.8.1-fr}
Soient $A$ un ensemble fini, $\Sigma(A)$ l'ensemble des suites finies
$\sigma=(\alpha_0,\ldots,\alpha_h)$ d'éléments de $A$ («~simplexes~» de
$A$); on pose $|\sigma|=\{\alpha_0,\ldots,\alpha_h\}$; on rappelle que le
complexe de chaînes $\mathrm C_\bullet(A)$ est le groupe abélien libre
gradué engendré par les éléments de $\Sigma(A)$,
$(\alpha_0,\ldots,\alpha_h)$ étant de degré $h$, avec une différentielle
définie par
\[
d(\alpha_0,\ldots,\alpha_h)=
\sum_{j=0}^{h}(-1)^j
(\alpha_0,\ldots,\widehat{\alpha_j},\ldots,\alpha_h).
\]
Le sous-groupe $\mathrm D_\bullet(A)$ de $\mathrm C_\bullet(A)$ engendré par
les chaînes $\sigma=(\alpha_0,\ldots,\alpha_h)$ pour lesquelles deux des
$\alpha_i$ sont égaux, et par les chaînes
$\pi(\sigma)-\varepsilon_\pi\sigma$, où pour toute permutation $\pi$,
\[
\pi(\sigma)=(\alpha_{\pi(0)},\ldots,\alpha_{\pi(h)})
\]
et $\varepsilon_\pi$ est la signature de $\pi$, est un sous-complexe de
$\mathrm C_\bullet(A)$ dont les éléments sont dits \emph{chaînes
dégénérées}; on pose
$\mathrm L_\bullet(A)=\mathrm C_\bullet(A)/\mathrm D_\bullet(A)$, et on a un
homomorphisme naturel de complexes
\[
p:\mathrm C_\bullet(A)\longrightarrow\mathrm L_\bullet(A).
\]
On définit d'autre part un homomorphisme de complexes
\[
j:\mathrm L_\bullet(A)\longrightarrow\mathrm C_\bullet(A)
\]
de la façon suivante : on ordonne totalement $A$; à la classe
mod.~$\mathrm D_\bullet(A)$ d'un simplexe
$\sigma=(\alpha_0,\ldots,\alpha_h)$, on fait correspondre $0$ si deux des
$\alpha_i$ sont égaux, et la suite $(\beta_0,\ldots,\beta_h)$ des $\alpha_i$
rangés dans l'ordre croissant dans le cas contraire. Il est clair que $p\circ
j$ est l'identité de $\mathrm L_\bullet(A)$.
\end{env}

\begin{env}[11.8.2]
\label{0.11.8.2-fr}
Soit $B$ un second ensemble fini; si $d'$, $d''$ sont les différentielles de
$\mathrm C_\bullet(A)$ et $\mathrm C_\bullet(B)$, le complexe produit
tensoriel $\mathrm C_\bullet(A)\otimes\mathrm C_\bullet(B)$ peut être
considéré comme le groupe abélien libre engendré par les éléments de
$\Sigma(A)\times\Sigma(B)$, avec la différentielle
\[
d(\sigma,\tau)=(d'\sigma,\tau)+(-1)^{h+1}(\sigma,d''\tau)
\quad\text{si }\operatorname{Card}|\sigma|=h+1.
\]
Les homomorphismes naturels
$\mathrm C_\bullet(A)\to\mathrm L_\bullet(A)$,
$\mathrm C_\bullet(B)\to\mathrm L_\bullet(B)$ définissent un homomorphisme
\[
p:\mathrm C_\bullet(A)\otimes\mathrm C_\bullet(B)
\longrightarrow\mathrm L_\bullet(A)\otimes\mathrm L_\bullet(B),
\]
ce dernier produit tensoriel étant isomorphe à
\[
\frac{\mathrm C_\bullet(A)\otimes\mathrm C_\bullet(B)}
{\mathrm C_\bullet(A)\otimes\mathrm D_\bullet(B)
+\mathrm D_\bullet(A)\otimes\mathrm C_\bullet(B)}.
\]
De même, à l'aide des homomorphismes
$\mathrm L_\bullet(A)\to\mathrm C_\bullet(A)$ et
$\mathrm L_\bullet(B)\to\mathrm C_\bullet(B)$ définis dans
(\hyperref[0.11.8.1-fr]{11.8.1}) (à l'aide d'ordres totaux sur $A$ et $B$),
on définit un homomorphisme
\[
j:\mathrm L_\bullet(A)\otimes\mathrm L_\bullet(B)
\longrightarrow\mathrm C_\bullet(A)\otimes\mathrm C_\bullet(B)
\]
tel que $p\circ j$ soit l'identité.
\end{env}

\oldpage[0\textsubscript{III}]{44}
Avec ces notations :
\begin{proposition}[11.8.3]
\label{0.11.8.3-fr}
\emph{Il existe une homotopie
\[
h:\mathrm C_\bullet(A)\otimes\mathrm C_\bullet(B)
\longrightarrow\mathrm C_\bullet(A)\otimes\mathrm C_\bullet(B),
\]
telle que $h(\sigma,\tau)$ soit combinaison linéaire de couples de simplexes
$(\sigma_i,\tau_i)$ avec $|\sigma_i|\subset|\sigma|$,
$|\tau_i|\subset|\tau|$, et que l'on ait pour $f=j\circ p$,}
\[
\label{0.11.8.3.1-fr}
\tag{11.8.3.1}
f-1=h\circ d+d\circ h.
\]
\end{proposition}

\begin{proof}
Il suffit de définir $h$ sur chaque couple $(\sigma,\tau)$ de simplexes, en
raisonnant par récurrence sur la somme des degrés de $\sigma$ et $\tau$, car
on peut prendre $h=0$ lorsque cette somme est $0$. Soit
\[
\omega=f(\sigma,\tau)-(\sigma,\tau)-h(d(\sigma,\tau));
\]
en vertu de l'hypothèse de récurrence et de la définition de $d$, on a
\[
\omega\in\mathrm C_\bullet(|\sigma|)\otimes\mathrm C_\bullet(|\tau|).
\]
On a
\[
d\omega=f(d(\sigma,\tau))-d(\sigma,\tau)-d(h(d(\sigma,\tau)))
=h(d(d(\sigma,\tau)))=0
\]
en vertu de (\hyperref[0.11.8.3.1-fr]{11.8.3.1}) et de l'hypothèse de
récurrence. Or, on a
$H_q(\mathrm C_\bullet(A))=0$ pour $q>0$ (G, I, 3.7.4), donc aussi
\[
H_q(\mathrm C_\bullet(A)\otimes\mathrm C_\bullet(B))=0\quad(q>0),
\]
en vertu de la formule de Künneth (G, I, 5.5.2). Appliquant cette remarque en
remplaçant $A$ par $|\sigma|$ et $B$ par $|\tau|$, on voit qu'il existe un
élément
\[
\omega'\in\mathrm C_\bullet(|\sigma|)\otimes\mathrm C_\bullet(|\tau|)
\]
tel que $\omega=d\omega'$; prenant $h(\sigma,\tau)=\omega'$, on vérifie
(\hyperref[0.11.8.3.1-fr]{11.8.3.1}) pour le couple $(\sigma,\tau)$ et la
récurrence peut se poursuivre.
\end{proof}

\begin{env}[11.8.4]
\label{0.11.8.4-fr}
Les notations étant celles de (\hyperref[0.11.8.2-fr]{11.8.2}), nous poserons
$(\sigma,\tau)\leq(\sigma',\tau')$ si $|\sigma|\subset|\sigma'|$ et
$|\tau|\subset|\tau'|$. Un \emph{système de coefficients} $\Gamma$ sur
$\Sigma(A)\times\Sigma(B)$ est constitué par une famille
$(\Gamma_{\sigma,\tau})$ de groupes abéliens, où
$\Gamma_{\sigma,\tau}$ ne dépend que des ensembles $|\sigma|$ et $|\tau|$,
et une famille d'homomorphismes
\[
\Gamma_{\sigma,\tau}\longrightarrow\Gamma_{\sigma',\tau'}
\quad\text{pour }(\sigma,\tau)\leq(\sigma',\tau'),
\]
formant un système inductif pour cette relation de préordre. On définit alors
un complexe de cochaînes $\mathrm C^\bullet(A,B;\Gamma)$ comme l'ensemble
des familles $\lambda=(\lambda(\sigma,\tau))$, où $(\sigma,\tau)$ parcourt
$\Sigma(A)\times\Sigma(B)$, avec
$\lambda(\sigma,\tau)\in\Gamma_{\sigma,\tau}$ pour tout couple
$(\sigma,\tau)$. La différentielle est donnée de la façon suivante : si
\[
d(\sigma,\tau)=\sum_i\pm(\sigma_i,\tau_i),
\]
on a $|\sigma_i|\subset|\sigma|$, $|\tau_i|\subset|\tau|$ pour tout $i$, et
l'on prend
\[
d\lambda(\sigma,\tau)=\sum_i\pm\lambda_i(\sigma_i,\tau_i),
\]
où $\lambda_i(\sigma_i,\tau_i)$ désigne l'image canonique de
$\lambda(\sigma_i,\tau_i)$ dans $\Gamma_{\sigma,\tau}$.

Nous dirons qu'une cochaîne $\lambda\in\mathrm C^\bullet(A,B;\Gamma)$ est
\emph{bi-alternée} si $\lambda(\sigma,\tau)=0$ lorsque l'un des deux
simplexes $\sigma$, $\tau$ a deux termes égaux, et si l'on a
\[
\lambda(\pi(\sigma),\tau)=\varepsilon_\pi\lambda(\sigma,\tau),
\qquad
\lambda(\sigma,\pi'(\tau))=\varepsilon_{\pi'}\lambda(\sigma,\tau)
\]
pour des permutations quelconques $\pi$, $\pi'$ des indices. Il est clair
que ces cochaînes engendrent un sous-complexe
$\mathrm L^\bullet(A,B;\Gamma)$ de $\mathrm C^\bullet(A,B;\Gamma)$.
\end{env}

\begin{proposition}[11.8.5]
\label{0.11.8.5-fr}
\emph{L'injection canonique
\[
\mathrm L^\bullet(A,B;\Gamma)\longrightarrow
\mathrm C^\bullet(A,B;\Gamma)
\]
définit un isomorphisme pour la cohomologie de ces deux complexes.}
\end{proposition}

\begin{proof}
Notons que si $p$ et $j$ ont le sens défini dans
(\hyperref[0.11.8.2-fr]{11.8.2}), les applications
\[
{}'p:\lambda\longmapsto\lambda\circ p,\qquad
{}'j:\lambda\longmapsto\lambda\circ j
\]
sont définies dans $\mathrm L^\bullet(A,B;\Gamma)$ et
$\mathrm C^\bullet(A,B;\Gamma)$ respectivement, la première n'étant autre
que l'injection canonique. Comme $\mathop{}'j\circ{}'p$ est l'identité, il
suffit de démontrer que $\mathop{}'p\circ{}'j$ est homotope à l'identité; or,
d'après (\hyperref[0.11.8.3-fr]{11.8.3}),
$\mathop{}'h:\lambda\mapsto\lambda\circ h$ est définie dans
$\mathrm C^\bullet(A,B;\Gamma)$ et on peut donc transposer l'identité
(\hyperref[0.11.8.3.1-fr]{11.8.3.1}), qui fournit le résultat cherché.
\end{proof}

\begin{env}[11.8.6]
\label{0.11.8.6-fr}
\oldpage[0\textsubscript{III}]{45}
La proposition (\hyperref[0.11.8.5-fr]{11.8.5}) ramène le calcul de la
cohomologie de $\mathrm L^\bullet(A,B;\Gamma)$ à celui de la cohomologie de
$\mathrm C^\bullet(A,B;\Gamma)$. Rappelons d'autre part que cette dernière
est, en vertu du th. d'Eilenberg-Zilber (G, I, 3.10.2), canoniquement
isomorphe à la cohomologie du complexe de cochaînes défini comme suit : on
forme le complexe de chaînes $\mathrm P_\bullet(A,B)$, constitué par les
combinaisons linéaires des $(\sigma,\tau)\in\Sigma(A)\times\Sigma(B)$ tels
que $\sigma$ et $\tau$ aient le même degré; la différentielle de ce complexe
est donnée par
\[
d(\sigma,\tau)=\sum_{j,k}(-1)^{j+k}(\sigma_j,\tau_k)
\]
si
\[
d\sigma=\sum_j(-1)^j\sigma_j,\qquad
d\tau=\sum_k(-1)^k\tau_k;
\]
on a alors deux homomorphismes canoniques de complexes
\[
f:\mathrm P_\bullet(A,B)\longrightarrow
\mathrm C_\bullet(A)\otimes\mathrm C_\bullet(B),\qquad
g:\mathrm C_\bullet(A)\otimes\mathrm C_\bullet(B)
\longrightarrow\mathrm P_\bullet(A,B),
\]
et on démontre (\emph{loc. cit.}) qu'il y a des homotopies $h$, $h'$ telles
que
\[
f\circ g-1=d\circ h+h\circ d,\qquad
g\circ f-1=d\circ h'+h'\circ d.
\]
En outre, on a
\[
f(\sigma,\tau)\in\mathrm C_\bullet(|\sigma|)
\otimes\mathrm C_\bullet(|\tau|),\qquad
g(\sigma,\tau)\in\mathrm P_\bullet(|\sigma|,|\tau|)
\]
et les homotopies $h$, $h'$ peuvent être prises telles que
\[
h(\sigma,\tau)\in\mathrm C_\bullet(|\sigma|)
\otimes\mathrm C_\bullet(|\tau|),\qquad
h'(\sigma,\tau)\in\mathrm P_\bullet(|\sigma|,|\tau|).
\]
Ce point provient de ce que la définition de $h(\sigma,\tau)$ et
$h'(\sigma,\tau)$ peut se faire par récurrence sur la somme des degrés de
$\sigma$ et $\tau$, et de ce que les
\[
H^q(\mathrm C_\bullet(|\sigma|)\otimes\mathrm C_\bullet(|\tau|))
\quad\text{et}\quad
H^q(\mathrm P_\bullet(|\sigma|,|\tau|))
\]
sont nuls pour $q>0$ (\emph{loc. cit.}); on raisonne alors comme dans
(\hyperref[0.11.8.3-fr]{11.8.3}) et la conclusion en résulte.

On définit alors $\mathrm P^\bullet(A,B;\Gamma)$ comme l'ensemble des
familles $\lambda=(\lambda(\sigma,\tau))$ où $(\sigma,\tau)$ parcourt les
couples dont les termes ont même degré, avec
$\lambda(\sigma,\tau)\in\Gamma_{\sigma,\tau}$, et comme on a
\[
d\sigma=\sum_j(-1)^j\sigma_j,\qquad
d\tau=\sum_k(-1)^k\tau_k,
\]
\[
d\lambda(\sigma,\tau)=
\sum_{j,k}(-1)^{j+k}\lambda(\sigma_j,\tau_k)
\]
est défini et donne la différentielle du complexe
$\mathrm P^\bullet(A,B;\Gamma)$. Cela étant, les applications
\[
{}'f:\lambda\mapsto\lambda\circ f,\quad
{}'g:\lambda\mapsto\lambda\circ g,\quad
{}'h:\lambda\mapsto\lambda\circ h,\quad
{}'h':\lambda\mapsto\lambda\circ h'
\]
sont toutes définies en vertu des remarques qui précèdent;
$\mathop{}'f\circ{}'g$ et $\mathop{}'g\circ{}'f$ sont donc homotopes à
l'identité, d'où l'isomorphisme cherché entre la cohomologie de
$\mathrm C^\bullet(A,B;\Gamma)$ et celle de
$\mathrm P^\bullet(A,B;\Gamma)$.
\end{env}

\begin{env}[11.8.7]
\label{0.11.8.7-fr}
\emph{Remarque}. Le même raisonnement que dans
(\hyperref[0.11.8.3-fr]{11.8.3}), mais appliqué à
$\mathrm C_\bullet(A)$ et $\mathrm L_\bullet(A)$, montre que si $j$ et $p$
sont définis comme dans (\hyperref[0.11.8.1-fr]{11.8.1}), $f=j\circ p$
vérifie encore une relation (\hyperref[0.11.8.3.1-fr]{11.8.3.1}), avec
$|h(\sigma)|\subset|\sigma|$, d'où on déduit comme dans
(\hyperref[0.11.8.5-fr]{11.8.5}) un isomorphisme de la cohomologie de
$\mathrm L^\bullet(A;\Gamma)$ sur celle de
$\mathrm C^\bullet(A;\Gamma)$, ces deux complexes étant définis de façon
évidente. C'est le résultat dont la démonstration est esquissée dans
(G, I, 3.8.1).
\end{env}

\begin{env}[11.8.8]
\label{0.11.8.8-fr}
Reprenons maintenant les notations et hypothèses de
(\hyperref[0.11.8.2-fr]{11.8.2}), et considérons un complexe
$\Gamma^\bullet=(\Gamma^k)$ de systèmes de coefficients sur
$\Sigma(A)\times\Sigma(B)$ : pour chaque $(\sigma,\tau)$, les
$\Gamma^k_{\sigma,\tau}$ forment donc un complexe de groupes abéliens
$(k\in\mathbb Z)$, et on a les diagrammes commutatifs
\[
\xymatrix{
\Gamma^k_{\sigma,\tau}\ar[r]\ar[d] &
\Gamma^{k+1}_{\sigma,\tau}\ar[d]\\
\Gamma^k_{\sigma',\tau'}\ar[r] &
\Gamma^{k+1}_{\sigma',\tau'}
}
\]
\oldpage[0\textsubscript{III}]{46}
pour $(\sigma,\tau)\leq(\sigma',\tau')$. Alors on vérifie aussitôt que
\[
\mathrm C^\bullet(A,B;\Gamma^\bullet)
=\bigl(\mathrm C^h(A,B;\Gamma^k)\bigr)_{(h,k)\in\mathbb Z\times\mathbb Z}
\]
est un bicomplexe de groupes abéliens, et
\[
\mathrm L^\bullet(A,B;\Gamma^\bullet)
=\bigl(\mathrm L^h(A,B;\Gamma^k)\bigr)
\]
en est un sous-bicomplexe.
\end{env}

\begin{proposition}[11.8.9]
\label{0.11.8.9-fr}
\emph{L'injection canonique
\[
\mathrm L^\bullet(A,B;\Gamma^\bullet)\longrightarrow
\mathrm C^\bullet(A,B;\Gamma^\bullet)
\]
définit un isomorphisme pour la cohomologie de ces deux bicomplexes.}
\end{proposition}

\begin{proof}
Posons $\mathrm C^{\bullet\bullet}=
\mathrm C^\bullet(A,B;\Gamma^\bullet)$ et
$\mathrm L^{\bullet\bullet}=
\mathrm L^\bullet(A,B;\Gamma^\bullet)$ pour simplifier, et notons que
puisque $\mathrm C^{hk}=\mathrm L^{hk}=0$ pour $h<0$, les secondes suites
spectrales de ces bicomplexes sont régulières
(\hyperref[0.11.3.3-fr]{11.3.3}); l'homomorphisme
$\mathrm L^{\bullet\bullet}\to\mathrm C^{\bullet\bullet}$ fournit donc un
morphisme de suites spectrales
\[
{}''E(\mathrm L^{\bullet\bullet})\longrightarrow
{}''E(\mathrm C^{\bullet\bullet})
\]
qui, pour les termes $E_2$, se réduit à
\[
\label{0.11.8.9.1-fr}
\tag{11.8.9.1}
H^p_{\mathrm{II}}(H^q_{\mathrm I}(\mathrm L^{\bullet\bullet}))
\longrightarrow
H^p_{\mathrm{II}}(H^q_{\mathrm I}(\mathrm C^{\bullet\bullet})).
\]
Mais pour tout $k\in\mathbb Z$, il résulte de
(\hyperref[0.11.8.3-fr]{11.8.3}) que l'homomorphisme
$H^q_{\mathrm I}(\mathrm L^{\bullet,k})\to
H^q_{\mathrm I}(\mathrm C^{\bullet,k})$ est bijectif; la conclusion résulte
donc de (\hyperref[0.11.1.5-fr]{11.1.5}).
\end{proof}

\begin{env}[11.8.10]
\label{0.11.8.10-fr}
De même, avec les notations de (\hyperref[0.11.8.6-fr]{11.8.6}), on a des
homomorphismes canoniques de bicomplexes
\[
\mathrm C^\bullet(A,B;\Gamma^\bullet)\longrightarrow
\mathrm P^\bullet(A,B;\Gamma^\bullet)
\]
(avec des notations évidentes), et le même raisonnement que dans
(\hyperref[0.11.8.9-fr]{11.8.9}), basé cette fois sur
(\hyperref[0.11.8.6-fr]{11.8.6}), montre que cet homomorphisme donne encore
un isomorphisme en cohomologie.
\end{env}

\subsection{Un lemme sur les complexes de type fini}
\label{subsection:0.11.9-fr}

\begin{proposition}[11.9.1]
\label{0.11.9.1-fr}
Soient $\mathcal C$ une catégorie abélienne, $K'$ et $K''$ des parties de
l'ensemble des objets de $\mathcal C$, telles que $K''\subset K'$, et
vérifiant les conditions suivantes :
\begin{enumerate}
\item[{\rm(i)}] \emph{Pour tout objet $A'\in K'$ et tout épimorphisme
$u:A\to A'$ dans $\mathcal C$, il existe un objet $B\in K''$ et un morphisme
$v:B\to A$ tels que $u\circ v$ soit un épimorphisme.}
\item[{\rm(ii)}] \emph{Pour tout couple d'objets $A\in K'$, $B\in K'$ et
tout épimorphisme $u:A\to B$, $\operatorname{Ker}(u)$ appartient à $K'$.}
\item[{\rm(iii)}] \emph{Le produit de deux objets de $K''$ appartient à
$K''$.}
\end{enumerate}
\emph{Soit $P_\bullet=(P_i)_{i\in\mathbb Z}$ un complexe dans $\mathcal C$,
tel que $H_i(P_\bullet)\in K'$ pour tout $i$, et qu'il existe $d$ tel que
$H_i(P_\bullet)=0$ pour $i<d$. Alors il existe un complexe
$Q_\bullet=(Q_i)_{i\in\mathbb Z}$ dans $\mathcal C$, tel que $Q_i\in K''$
pour tout $i$ et $Q_i=0$ pour $i<d$, et un morphisme
$u:Q_\bullet\to P_\bullet$ de complexes tel que le morphisme correspondant
$H_\bullet(Q_\bullet)\to H_\bullet(P_\bullet)$ soit un isomorphisme.}
\end{proposition}

\begin{proof}
Démontrons d'abord la conséquence suivante de la propriété (i) :

\emph{(i bis) Soient $u:C\to B$ un épimorphisme dans $\mathcal C$, $A$ un
objet de $K'$, $v:A\to B$ un morphisme dans $\mathcal C$; il existe alors un
objet $D\in K''$, un épimorphisme $u':D\to A$ et un morphisme $v':D\to C$
tels que le diagramme}
\[
\xymatrix{
D\ar[r]^{u'}\ar[d]_{v'} & A\ar[d]^{v}\\
C\ar[r]_{u} & B
}
\tag{11.9.1.1}\label{0.11.9.1.1-fr}
\]
\emph{soit commutatif.}

Considérons en effet le produit fibré $C\times_B A$ dans $\mathcal C$ et les
projections canoniques $p:C\times_B A\to C$, $q:C\times_B A\to A$, rendant
commutatif le diagramme
\oldpage[0\textsubscript{III}]{47}
\[
\xymatrix{
C\times_B A\ar[r]^{q}\ar[d]_{p} & A\ar[d]^{v}\\
C\ar[r]_{u} & B
}
\tag{11.9.1.2}\label{0.11.9.1.2-fr}
\]
On sait ([27], p.~1-12) que le conoyau de $q$ est le quotient de $A$ par
$v^{-1}(u(C))$; comme $u$ est un épimorphisme, $u(C)=B$ et
$v^{-1}(u(C))=A$, donc $q$ est un épimorphisme; il suffit alors d'appliquer
(i) à l'épimorphisme $q:C\times_B A\to A$ : il y a un objet $D\in K''$ et un
morphisme $w:D\to C\times_B A$ tels que $q\circ w$ soit un épimorphisme; on
prendra $u'=q\circ w$, $v'=p\circ w$.

Cela étant, pour démontrer la proposition, procédons par récurrence.
Supposons, pour un $i\geq d-1$, construits, pour $j\leq i$, les objets $Q_j$,
les morphismes $d_j:Q_j\to Q_{j-1}$ et les morphismes $u_j:Q_j\to P_j$, de
sorte que $Q_j=0$ pour $j<d$, que $d_{j-1}\circ d_j=0$ et
$d_j\circ u_j=u_{j-1}\circ d_j$ pour $j\leq i$; en outre, nous supposons
vérifiées les conditions suivantes :
\begin{enumerate}
\item[$(\mathrm I_i)$] On a $Q_j\in K''$ pour $j\leq i$ et
$B_j(Q_\bullet)\in K'$ pour $j<i$.
\item[$(\mathrm{II}_i)$] Pour $j<i$, l'homomorphisme
$H_j(Q_\bullet)\to H_j(P_\bullet)$ déduit de la famille $(u_k)_{k\leq i}$
est un isomorphisme.
\item[$(\mathrm{III}_i)$] Le morphisme composé
\[
v_i:Z_i(Q_\bullet)\longrightarrow Z_i(P_\bullet)
\longrightarrow H_i(P_\bullet)
\]
(où la flèche de gauche est la restriction de $u_i$ et la flèche de droite
le morphisme canonique) est un épimorphisme.
\end{enumerate}
Notons que, d'après (ii), $Z_i(Q_\bullet)$, noyau de l'épimorphisme
$Q_i\to B_{i-1}(Q_\bullet)$, appartient à $K'$ en vertu de l'hypothèse
$(\mathrm I_i)$. On tire encore de (ii) que
$N_i=\operatorname{Ker}(v_i)$ appartient aussi à $K'$, compte tenu de
l'hypothèse $(\mathrm{III}_i)$. En vertu de (i bis), il existe un
$Q'_{i+1}\in K''$, un épimorphisme $d'_{i+1}:Q'_{i+1}\to N_i$ et un
morphisme $u'_{i+1}:Q'_{i+1}\to P_{i+1}$, tels que le diagramme
\[
\xymatrix{
Q'_{i+1}\ar[r]^{d'_{i+1}}\ar[d]_{u'_{i+1}} &
N_i\ar[d]^{u_i}\\
P_{i+1}\ar[r]_{d_{i+1}} & B_i(P_\bullet)
}
\tag{11.9.1.3}\label{0.11.9.1.3-fr}
\]
soit commutatif.

Comme le morphisme canonique $Z_{i+1}(P_\bullet)\to H_{i+1}(P_\bullet)$ est
un épimorphisme et que $H_{i+1}(P_\bullet)\in K'$ par hypothèse, il résulte
de (i) qu'il existe un objet $Q''_{i+1}\in K''$ et un morphisme
$u''_{i+1}:Q''_{i+1}\to Z_{i+1}(P_\bullet)$ tels que le composé
$Q''_{i+1}\to Z_{i+1}(P_\bullet)\to H_{i+1}(P_\bullet)$ soit un
épimorphisme. Si on prend $d''_{i+1}:Q''_{i+1}\to N_i$ égal à $0$, le
diagramme
\[
\xymatrix{
Q''_{i+1}\ar[r]^{d''_{i+1}}\ar[d]_{u''_{i+1}} &
N_i\ar[d]^{u_i}\\
Z_{i+1}(P_\bullet)\ar[r]_{d_{i+1}} & P_i
}
\tag{11.9.1.4}\label{0.11.9.1.4-fr}
\]
est commutatif, la flèche horizontale inférieure étant $0$.
\oldpage[0\textsubscript{III}]{48}
Prenons alors
\[
Q_{i+1}=Q'_{i+1}\times Q''_{i+1},\qquad
d_{i+1}=d'_{i+1}+d''_{i+1},\qquad
u_{i+1}=u'_{i+1}+u''_{i+1}.
\]
Comme $d_{i+1}(Q_{i+1})=d'_{i+1}(Q'_{i+1})=N_i\subset Z_i(Q_\bullet)$, on a
$d_i\circ d_{i+1}=0$ et, avec les notations usuelles,
$B_i(Q_\bullet)=N_i$, ce qui vérifie $(\mathrm I_{i+1})$. La commutativité
des diagrammes (\hyperref[0.11.9.1.3-fr]{11.9.1.3}) et
(\hyperref[0.11.9.1.4-fr]{11.9.1.4}) montre que
$d_{i+1}\circ u_{i+1}=u_i\circ d_{i+1}$. Par définition de $N_i$, le
morphisme
\[
H_i(Q_\bullet)=Z_i(Q_\bullet)/N_i\longrightarrow H_i(P_\bullet)
\]
déduit du système des $u_k$ ($k\leq i+1$) est le morphisme déduit de $v_i$
par passage aux quotients, donc c'est un isomorphisme puisque $v_i$ est un
épimorphisme, d'où $(\mathrm{II}_{i+1})$. Enfin, on a
$Q''_{i+1}\subset Z_{i+1}(Q_\bullet)$ par définition; le choix de
$u''_{i+1}$ montre que le morphisme
\[
v_{i+1}:Z_{i+1}(Q_\bullet)\longrightarrow Z_{i+1}(P_\bullet)
\longrightarrow H_{i+1}(P_\bullet)
\]
est un épimorphisme, sa restriction à $Q''_{i+1}$ l'étant déjà, d'où
$(\mathrm{III}_{i+1})$. La récurrence peut donc se poursuivre, et la
proposition est démontrée.
\end{proof}

\begin{corollary}[11.9.2]
\label{0.11.9.2-fr}
Soient $A$ un anneau noethérien (non nécessairement commutatif),
$P_\bullet=(P_i)_{i\in\mathbb Z}$ un complexe de $A$-modules à droite. On
suppose que les $H_i(P_\bullet)$ sont des $A$-modules de type fini et qu'il
existe $d$ tel que $H_i(P_\bullet)=0$ pour $i<d$. Alors il existe un complexe
$Q_\bullet=(Q_i)_{i\in\mathbb Z}$ formé de $A$-modules à droite libres de
rang fini, tel que $Q_i=0$ pour $i<d$, et un homomorphisme
$u:Q_\bullet\to P_\bullet$ de complexes, tel que l'homomorphisme
$H_\bullet(Q_\bullet)\to H_\bullet(P_\bullet)$ correspondant à $u$ soit
bijectif.
\end{corollary}

\begin{proof}
On applique (\hyperref[0.11.9.1-fr]{11.9.1}) en prenant pour $\mathcal C$ la
catégorie des $A$-modules à droite, pour $K'$ (resp. $K''$) l'ensemble des
$A$-modules de type fini (resp. l'ensemble des $A$-modules libres de rang
fini); la vérification des conditions (i), (ii) et (iii) de
(\hyperref[0.11.9.1-fr]{11.9.1}) est immédiate, compte tenu de l'hypothèse
que $A$ est noethérien.
\end{proof}

\begin{env}[11.9.3]
\label{0.11.9.3-fr}
\emph{Remarques}.
\begin{enumerate}
\item[{\rm(i)}] Sous les conditions de
(\hyperref[0.11.9.2-fr]{11.9.2}), supposons en outre que les $P_i$ soient des
$A$-modules à droite plats. Alors, pour tout $A$-module à gauche $M$,
l'homomorphisme de complexes
\[
u\otimes 1:Q_\bullet\otimes_A M\longrightarrow P_\bullet\otimes_A M
\]
définit encore un isomorphisme
$H_\bullet(Q_\bullet\otimes_A M)\simeq
H_\bullet(P_\bullet\otimes_A M)$ de l'homologie comme nous le verrons au
chap.~III.
\item[{\rm(ii)}] La conclusion de (\hyperref[0.11.9.2-fr]{11.9.2}) n'est
plus nécessairement exacte lorsqu'on ne suppose pas $A$ noethérien; en
effet, en l'appliquant à un complexe réduit à $0$ sauf pour un seul terme, on
en conclurait que tout $A$-module à gauche de type fini admet une résolution
par des modules libres de type fini, ce qui n'est pas vrai en général
(cf. Bourbaki, \emph{Alg. comm.}, chap.~I, \S~2, exerc.~6).

Toutefois, au lieu de supposer $A$ noethérien, on peut supposer seulement que
les $H_i(P_\bullet)$ ont une $\infty$-présentation finie (cf. chap.~IV).
\end{enumerate}
\end{env}

\subsection{Caractéristique d'Euler-Poincaré d'un complexe de modules de
longueur finie}
\label{subsection:0.11.10-fr}

\begin{env}[11.10.1]
\label{0.11.10.1-fr}
Soient $A$ un anneau (non nécessairement commutatif),
\[
\label{0.11.10.1.1-fr}
\tag{11.10.1.1}
M^\bullet:0\longrightarrow M^0\longrightarrow M^1
\longrightarrow\cdots\longrightarrow M^n\longrightarrow0
\]
un complexe de $A$-modules à gauche de longueur finie. On appelle
\emph{caractéristique d'Euler-Poincaré} de ce complexe le nombre
\oldpage[0\textsubscript{III}]{49}
\[
\label{0.11.10.1.2-fr}
\tag{11.10.1.2}
\chi(M^\bullet)=\sum_{i=0}^{n}(-1)^i\operatorname{long}(M^i).
\]
\end{env}

\begin{proposition}[11.10.2]
\label{0.11.10.2-fr}
\emph{Pour tout complexe fini $M^\bullet$ de $A$-modules à gauche de longueur
finie, on a
\[
\chi(M^\bullet)=\chi(H^\bullet(M^\bullet))
\]
($H^\bullet(M^\bullet)$ étant considéré comme un complexe pour la dérivation
triviale). En particulier, si la suite
(\hyperref[0.11.10.1.1-fr]{11.10.1.1}) est exacte, on a
$\chi(M^\bullet)=0$.}
\end{proposition}

\begin{proof}
Posons pour abréger $B^i=B^i(M^\bullet)$, $Z^i=Z^i(M^\bullet)$,
$H^i=H^i(M^\bullet)=Z^i/B^i$; les $B^i$, $Z^i$, $H^i$ sont de longueur
finie. Des suites exactes
\[
0\longrightarrow B^i\longrightarrow Z^i\longrightarrow H^i
\longrightarrow0,
\qquad
0\longrightarrow Z^i\longrightarrow M^i\longrightarrow B^{i+1}
\longrightarrow0
\]
on tire les relations
\[
\operatorname{long}(Z^i)=\operatorname{long}(H^i)+
\operatorname{long}(B^i),
\qquad
\operatorname{long}(M^i)=\operatorname{long}(Z^i)+
\operatorname{long}(B^{i+1}),
\]
d'où
\[
\operatorname{long}(M^i)-\operatorname{long}(H^i)=
\operatorname{long}(B^{i+1})+\operatorname{long}(B^i).
\]
Multiplions cette relation par $(-1)^i$ et faisons la somme des relations
obtenues pour $0\leq i\leq n$, en notant que
$B^0=B^{n+1}=0$, il vient l'égalité cherchée.
\end{proof}

\begin{corollary}[11.10.3]
\label{0.11.10.3-fr}
Soit $E=(E_r^{pq})$ une suite spectrale dans la catégorie des modules sur un
anneau $A$. On suppose que les $E_2^{pq}$ sont des $A$-modules de longueur
finie et qu'il n'y a qu'un nombre fini de couples $(p,q)$ tels que
$E_2^{pq}\neq0$. Alors les caractéristiques d'Euler-Poincaré
$\chi(E_r^\bullet)$ de tous les complexes
$E_r^\bullet=(E_r^{(n)})_{n\in\mathbb Z}$
(\hyperref[0.11.1.1-fr]{11.1.1}) sont toutes égales. Si, en outre, la suite
$E$ est faiblement convergente et si on pose
\[
E_\infty^{(n)}=\bigoplus_{p+q=n}E_\infty^{pq}\quad(n\in\mathbb Z),
\]
on a aussi $\chi(E_\infty^\bullet)=\chi(E_2^\bullet)$,
$E_\infty^\bullet=(E_\infty^{(n)})_{n\in\mathbb Z}$ étant considéré comme
un complexe à dérivation triviale.
\end{corollary}

\begin{proof}
Notons d'abord que si $E_2^{pq}=0$, on a $E_r^{pq}=0$ pour
$2\leq r<+\infty$, donc tous les complexes $E_r^\bullet$ sont finis et
formés de $A$-modules de longueur finie; la relation
$\chi(E_r^\bullet)=\chi(E_{r+1}^\bullet)$ pour tout $r$ fini résulte donc de
(\hyperref[0.11.10.2-fr]{11.10.2}) et de l'isomorphie entre
$H^\bullet(E_r^\bullet)$ et $E_{r+1}^\bullet$ (en tant que complexes à
dérivation triviale). L'hypothèse que les $E_r^{pq}$ sont de longueur finie
entraîne que pour tout couple $(p,q)$, les suites
$(B_k(E_2^{pq}))_{k\geq2}$ et $(Z_k(E_2^{pq}))_{k\geq2}$ sont stationnaires;
l'hypothèse que $E$ est faiblement convergente et que $E_2^{pq}=0$, sauf
pour un nombre fini de couples $(p,q)$, entraîne donc qu'il existe un entier
$r\geq2$ tel que $E_\infty^{pq}=E_r^{pq}$ pour tout couple $(p,q)$; d'où
l'assertion relative à $\chi(E_\infty^\bullet)$.
\end{proof}

\section{Compléments sur la cohomologie des faisceaux}
\label{section:0.12-fr}

\subsection{Cohomologie des faisceaux de modules sur les espaces annelés}
\label{subsection:0.12.1-fr}

\begin{env}[12.1.1]
\label{0.12.1.1-fr}
Soit $(X,\mathcal O_X)$ un espace annelé. Rappelons que pour tout
$\mathcal O_X$-Module $\mathcal F$, on définit la cohomologie
$H^\bullet(X,\mathcal F)$, qui est un foncteur cohomologique universel
(T, 2.2) de la catégorie $C(X)$ des $\mathcal O_X$-Modules dans la catégorie
des groupes abéliens; c'est le foncteur dérivé du foncteur exact à gauche
$\mathcal F\to\Gamma(X,\mathcal F)$. Le foncteur $H^\bullet(X,\mathcal F)$
est isomorphe
\oldpage[0\textsubscript{III}]{50}
à la restriction à la catégorie $C(X)$ du foncteur cohomologique défini de
même sur la catégorie des \emph{faisceaux de groupes abéliens} sur $X$
(G, II, 7.2.1).
\end{env}

\begin{env}[12.1.2]
\label{0.12.1.2-fr}
Posons $A=\Gamma(X,\mathcal O_X)$. Comme tout élément de $A$ définit un
endomorphisme du groupe abélien $\Gamma(X,\mathcal F)$, il définit par
fonctorialité un endomorphisme de $\delta$-foncteur de
$H^\bullet(X,\mathcal F)$; ces endomorphismes définissent sur chacun des
$H^p(X,\mathcal F)$ une structure de $A$-module, et l'opérateur $\partial$
est $A$-linéaire. En outre, pour deux entiers positifs quelconques $p$, $q$,
et deux $\mathcal O_X$-Modules $\mathcal F$, $\mathcal G$, on a un
homomorphisme de $A$-modules, dit \emph{cup-produit}
\[
\label{0.12.1.2.1-fr}
\tag{12.1.2.1}
H^p(X,\mathcal F)\otimes_A H^q(X,\mathcal G)
\longrightarrow H^{p+q}(X,\mathcal F\otimes_{\mathcal O_X}\mathcal G)
\]
(G, II, 6.6). Ces homomorphismes font de la somme directe $S$ des
$H^p(X,\mathcal O_X)$ (pour $p\geq0$) une $A$-algèbre graduée
anticommutative et de la somme directe des $H^p(X,\mathcal F)$ un $S$-module
gradué.

Pour tout \emph{recouvrement ouvert} $\mathfrak U$ de $X$, nous désignerons
toujours par $\check C^\bullet(\mathfrak U,\mathcal F)$ (contrairement à
(G, II, 5.1)) le complexe des cochaînes \emph{alternées} du nerf de
$\mathfrak U$ à valeurs dans le système de coefficients
$\Gamma(U_\alpha,\mathcal F)$. Il est clair que
$\check C^\bullet(\mathfrak U,\mathcal F)$ est un $A$-module gradué, donc les
groupes de cohomologie $\check H^p(\mathfrak U,\mathcal F)$ de ce complexe
sont munis d'une structure de $A$-module; en outre, les applications
canoniques
\[
\check H^p(\mathfrak U,\mathcal F)\longrightarrow H^p(X,\mathcal F)
\]
(G, II, 5.4) sont $A$-linéaires.
\end{env}

\begin{env}[12.1.3]
\label{0.12.1.3-fr}
Soit $(X',\mathcal O_{X'})$ un second espace annelé, et soit
$f=(\psi,\theta)$ un morphisme de $X'$ dans $X$.

Posons $A'=\Gamma(X',\mathcal O_{X'})$; le $\psi$-morphisme $\theta$ définit
canoniquement un homomorphisme d'anneaux $A\to A'$. Soient $\mathcal F$ un
$\mathcal O_X$-Module, $\mathcal F'$ un $\mathcal O_{X'}$-Module; pour tout
$f$-morphisme $u:\mathcal F\to\mathcal F'$
(\hyperref[0.4.4.1-fr]{0\textsubscript{I}, 4.4.1}), nous allons voir qu'on
peut définir pour tout $p\geq0$, un \emph{di-homomorphisme}
\[
\label{0.12.1.3.1-fr}
\tag{12.1.3.1}
u_p:H^p(X,\mathcal F)\longrightarrow H^p(X',\mathcal F').
\]
En effet, comme $\psi^*$ est exact dans la catégorie des faisceaux de groupes
abéliens sur $X$, $\mathcal F\to H^\bullet(X',\psi^*(\mathcal F))$ est un
$\delta$-foncteur dans cette catégorie, et on sait qu'on a un homomorphisme
canonique de $\delta$-foncteurs
\[
\label{0.12.1.3.2-fr}
\tag{12.1.3.2}
H^\bullet(X,\mathcal F)\longrightarrow H^\bullet(X',\psi^*(\mathcal F))
\]
uniquement déterminé par la condition de se réduire à l'homomorphisme
canonique $\Gamma(X,\mathcal F)\to\Gamma(X',\psi^*(\mathcal F))$ en degré
0 (T, 3.2.2). En outre, tout élément de $A$ détermine un endomorphisme $\mu$
de $\Gamma(X,\mathcal F)$ et un endomorphisme $\mu'$ de
$\Gamma(X',\psi^*(\mathcal F))$ tels que le diagramme
\[
\label{0.12.1.3.3-fr}
\tag{12.1.3.3}
\xymatrix{
\Gamma(X,\mathcal F)\ar[r]\ar[d]_\mu
&\Gamma(X',\psi^*(\mathcal F))\ar[d]^{\mu'}\\
\Gamma(X,\mathcal F)\ar[r]
&\Gamma(X',\psi^*(\mathcal F))
}
\]
soit commutatif; par la propriété d'unicité de prolongement des morphismes
pour les foncteurs cohomologiques universels (T, 2.2), on en déduit des
prolongements uniques de $\mu$ et $\mu'$ à la cohomologie rendant commutatifs
les diagrammes analogues à (\hyperref[0.12.1.3.3-fr]{12.1.3.3}), ce qui
signifie que (\hyperref[0.12.1.3.2-fr]{12.1.3.2}) est un homomorphisme de
$A$-modules. Notons maintenant que l'on a
\[
f^*(\mathcal F)=\psi^*(\mathcal F)
\otimes_{\psi^*(\mathcal O_X)}\mathcal O_{X'},
\]
et que l'on a donc un di-homomorphisme canonique
$\psi^*(\mathcal F)\to f^*(\mathcal F)$ du
$\psi^*(\mathcal O_X)$-Module $\psi^*(\mathcal F)$ dans le
$\mathcal O_{X'}$-Module $f^*(\mathcal F)$. Par fonctorialité, on en déduit
donc un di-homomorphisme fonctoriel
\[
\label{0.12.1.3.4-fr}
\tag{12.1.3.4}
H^p(X',\psi^*(\mathcal F))\longrightarrow H^p(X',f^*(\mathcal F))
\]
les anneaux correspondants étant $A$ et $A'$; en composant ce
di-homomorphisme avec (\hyperref[0.12.1.3.2-fr]{12.1.3.2}), on obtient un
di-homomorphisme canonique fonctoriel en $\mathcal F$
\[
\label{0.12.1.3.5-fr}
\tag{12.1.3.5}
\theta_p:H^p(X,\mathcal F)\longrightarrow H^p(X',f^*(\mathcal F)).
\]
Enfin, par fonctorialité, on déduit de l'homomorphisme
$u^\sharp:f^*(\mathcal F)\to\mathcal F'$ un homomorphisme de $A'$-modules
$H^p(X',f^*(\mathcal F))\to H^p(X',\mathcal F')$, qui, composé avec
(\hyperref[0.12.1.3.5-fr]{12.1.3.5}), donne
(\hyperref[0.12.1.3.1-fr]{12.1.3.1}).

Soit $f'=(\psi',\theta'):X''\to X'$ un second morphisme d'espaces annelés,
$f''=ff'$ le morphisme composé. Compte tenu de la permutabilité du foncteur
$\psi'^*$ et du produit tensoriel
(\hyperref[0.4.3.3-fr]{0\textsubscript{I}, 4.3.3}), on vérifie aussitôt que le
composé du di-homomorphisme
$H^p(X',f^*(\mathcal F))\to H^p(X'',f'^*(f^*(\mathcal F)))$ et de
(\hyperref[0.12.1.3.5-fr]{12.1.3.5}) est le di-homomorphisme correspondant
$H^p(X,\mathcal F)\to H^p(X'',f''^*(\mathcal F))$.
\end{env}

\begin{env}[12.1.4]
\label{0.12.1.4-fr}
Une définition directe de l'homomorphisme
(\hyperref[0.12.1.3.2-fr]{12.1.3.2}) peut s'obtenir de la façon suivante : on
considère une résolution injective $\mathcal L^\bullet=(\mathcal L^i)$ de
$\mathcal F$ formée de faisceaux de groupes abéliens sur $X$; comme le
foncteur $\psi^*$ est exact, $\psi^*(\mathcal L^\bullet)$ est une résolution
de $\psi^*(\mathcal F)$ formée de faisceaux sur $X'$. Si
$\mathcal L'^\bullet=(\mathcal L'^i)$ est une résolution injective de
$\psi^*(\mathcal F)$ dans la catégorie des faisceaux de groupes abéliens sur
$X'$, il y a donc un morphisme
$\psi^*(\mathcal L^\bullet)\to\mathcal L'^\bullet$ de complexes de faisceaux
de groupes abéliens, compatible avec les augmentations (M, V, 1.1.a)), bien
déterminé à une homotopie près. On en déduit des homomorphismes
\[
\Gamma(X,\mathcal L^\bullet)\longrightarrow
\Gamma(X',\psi^*(\mathcal L^\bullet))\longrightarrow
\Gamma(X',\mathcal L'^\bullet)
\]
de complexes de groupes abéliens, dont le composé, par passage à la
cohomologie donne un morphisme de $\delta$-foncteurs
$H^\bullet(X,\mathcal F)\to H^\bullet(X',\psi^*(\mathcal F))$; comme il
coïncide avec (\hyperref[0.12.1.3.2-fr]{12.1.3.2}) en degré 0, il lui est
identique (T, 2.2).

Considérons maintenant un \emph{recouvrement ouvert}
$\mathfrak U=(U_\alpha)$ de $X$, et soit
$\mathfrak U'=(f^{-1}(U_\alpha))$ le recouvrement ouvert de $X'$, image
réciproque de $\mathfrak U$. Les homomorphismes canoniques
$\Gamma(V,\mathcal F)\to\Gamma(f^{-1}(V),f^*(\mathcal F))$ pour tout ouvert
$V$ de $X$ définissent aussitôt (cf. G, II, 5.1) un homomorphisme de complexes
$\check C^\bullet(\mathfrak U,\mathcal F)\to
\check C^\bullet(\mathfrak U',f^*(\mathcal F))$, d'où des homomorphismes
canoniques
\[
\label{0.12.1.4.1-fr}
\tag{12.1.4.1}
\theta_p:\check H^p(\mathfrak U,\mathcal F)
\longrightarrow\check H^p(\mathfrak U',f^*(\mathcal F)).
\]

\oldpage[0\textsubscript{III}]{52}
En outre, on a des diagrammes commutatifs
\[
\label{0.12.1.4.2-fr}
\tag{12.1.4.2}
\xymatrix{
\check H^p(\mathfrak U,\mathcal F)\ar[r]^{\theta_p}\ar[d]
&\check H^p(\mathfrak U',f^*(\mathcal F))\ar[d]\\
H^p(X,\mathcal F)\ar[r]_{\theta_p}
&H^p(X',f^*(\mathcal F))
}
\]
où les flèches verticales sont les homomorphismes canoniques de (G, II, 5.2).
Pour établir la commutativité de
(\hyperref[0.12.1.4.2-fr]{12.1.4.2}), considérons le complexe de faisceaux de
cochaînes (alternées) de $\mathcal F$ relatives à $\mathfrak U$,
$\mathcal C^\bullet(\mathfrak U,\mathcal F)$, tel que
$\Gamma(X,\mathcal C^\bullet(\mathfrak U,\mathcal F))=
\check C^\bullet(\mathfrak U,\mathcal F)$ (G, II, 5.2). Les homomorphismes
canoniques $\Gamma(V,\mathcal F)\to
\Gamma(f^{-1}(V),\psi^*(\mathcal F))$ définissent alors un $\psi$-morphisme
$\mathcal C^\bullet(\mathfrak U,\mathcal F)\to
\mathcal C^\bullet(\mathfrak U',\psi^*(\mathcal F))$, et on a, avec les
notations ci-dessus, un diagramme commutatif
\[
\xymatrix{
\Gamma(X,\mathcal C^\bullet(\mathfrak U,\mathcal F))\ar[r]\ar[d]
&\Gamma(X',\mathcal C^\bullet(\mathfrak U',\psi^*(\mathcal F)))\ar[d]\\
\Gamma(X,\mathcal L^\bullet)\ar[r]
&\Gamma(X',\psi^*(\mathcal L^\bullet))\ar[r]
&\Gamma(X',\mathcal L'^\bullet)
}
\]
qui, en passant à la cohomologie, donne des diagrammes commutatifs
\[
\xymatrix{
\check H^p(\mathfrak U,\mathcal F)\ar[r]\ar[d]
&\check H^p(\mathfrak U',\psi^*(\mathcal F))\ar[d]\\
H^p(X,\mathcal F)\ar[r]
&H^p(X',\psi^*(\mathcal F))
}
\]
où les flèches verticales sont les homomorphismes canoniques de (G, II, 5.2).
Il suffit alors de combiner ces diagrammes avec les diagrammes commutatifs
\[
\xymatrix{
\check H^p(\mathfrak U',\psi^*(\mathcal F))\ar[r]\ar[d]
&\check H^p(\mathfrak U',f^*(\mathcal F))\ar[d]\\
H^p(X',\psi^*(\mathcal F))\ar[r]
&H^p(X',f^*(\mathcal F))
}
\]
\oldpage[0\textsubscript{III}]{53}
qui proviennent de l'homomorphisme $\psi^*(\mathcal F)\to f^*(\mathcal F)$
et du caractère fonctoriel des homomorphismes canoniques de (G, II, 5.2),
pour obtenir la commutativité de
(\hyperref[0.12.1.4.2-fr]{12.1.4.2}).

On notera que si $A=\Gamma(X,\mathcal O_X)$,
$A'=\Gamma(X',\mathcal O_{X'})$, l'homomorphisme (12.1.4.3) est un
di-homomorphisme de modules
correspondant aux anneaux $A$ et $A'$. On a une propriété de
\emph{transitivité} de (\hyperref[0.12.1.4.1-fr]{12.1.4.1}) pour le composé de
deux morphismes, analogue à la transitivité de
(\hyperref[0.12.1.3.5-fr]{12.1.3.5}). Enfin, notons que dans les définitions
précédentes, au lieu d'une résolution injective $\mathcal L^\bullet$ de
$\mathcal F$, on aurait aussi bien pu partir d'une résolution telle que
$H^p(X,\mathcal L^i)=0$ pour tout $i$ et tout $p>0$ (G, II, 4.7.1).
\end{env}

\begin{env}[12.1.5]
\label{0.12.1.5-fr}
Soient $\mathcal F$, $\mathcal G$, $\mathcal H$ trois
$\mathcal O_X$-Modules, et considérons un $\mathcal O_X$-homomorphisme
$u:\mathcal F\otimes_{\mathcal O_X}\mathcal G\to\mathcal H$, qui donne pour
la cohomologie des homomorphismes
\[
\label{0.12.1.5.1-fr}
\tag{12.1.5.1}
H^p(X,\mathcal F)\otimes_A H^q(X,\mathcal G)
\longrightarrow H^{p+q}(X,\mathcal H)
\]
déduits du cup-produit (\hyperref[0.12.1.2.1-fr]{12.1.2.1}). Montrons qu'avec
les hypothèses et notations de (\hyperref[0.12.1.3-fr]{12.1.3}), on a des
diagrammes commutatifs
\[
\label{0.12.1.5.2-fr}
\tag{12.1.5.2}
\xymatrix@C=10pt{
H^p(X,\mathcal F)\otimes_A H^q(X,\mathcal G)\ar[r]\ar[d]
&H^{p+q}(X,\mathcal H)\ar[d]\\
H^p(X',f^*(\mathcal F))\otimes_{A'}H^q(X',f^*(\mathcal G))\ar[r]
&H^{p+q}(X',f^*(\mathcal H))
}
\]
où les flèches verticales proviennent des homomorphismes canoniques
(\hyperref[0.12.1.3.5-fr]{12.1.3.5}). Pour cela, rappelons que
(\hyperref[0.12.1.5.1-fr]{12.1.5.1}) peut s'obtenir en partant des résolutions
\emph{canoniques} (G, II, 4.3) $\mathcal L^\bullet$, $\mathcal M^\bullet$,
$\mathcal N^\bullet$ de $\mathcal F$, $\mathcal G$, $\mathcal H$
respectivement (qui sont formées de $\mathcal O_X$-Modules), de l'application
linéaire $\mathcal L^\bullet\otimes_{\mathcal O_X}\mathcal M^\bullet
\to\mathcal N^\bullet$ de complexes de $\mathcal O_X$-Modules correspondant
à $u$, qui fournit un homomorphisme de complexes de $A$-modules
$\Gamma(X,\mathcal L^\bullet)\otimes_A\Gamma(X,\mathcal M^\bullet)
\to\Gamma(X,\mathcal N^\bullet)$, et en passant à la cohomologie des
homomorphismes
\[
H^p(\Gamma(X,\mathcal L^\bullet))\otimes_A
H^q(\Gamma(X,\mathcal M^\bullet))
\longrightarrow H^{p+q}(\Gamma(X,\mathcal N^\bullet))
\]
(G, II, 6.6). Or, on a évidemment un diagramme commutatif
\[
\label{0.12.1.5.3-fr}
\tag{12.1.5.3}
\xymatrix@C=8pt{
\Gamma(X,\mathcal L^\bullet)\otimes_A\Gamma(X,\mathcal M^\bullet)
\ar[r]\ar[d]
&\Gamma(X,\mathcal N^\bullet)\ar[d]\\
\Gamma(X',\psi^*(\mathcal L^\bullet))
\otimes_{\Gamma(X',\psi^*(\mathcal O_X))}
\Gamma(X',\psi^*(\mathcal M^\bullet))\ar[r]
&\Gamma(X',\psi^*(\mathcal N^\bullet))
}
\]
\oldpage[0\textsubscript{III}]{54}
qui donne en passant à la cohomologie des diagrammes commutatifs
\[
\label{0.12.1.5.4-fr}
\tag{12.1.5.4}
\xymatrix@C=7pt{
H^p(X,\mathcal F)\otimes_A H^q(X,\mathcal G)\ar[r]\ar[d]
&H^{p+q}(X,\mathcal H)\ar[d]\\
H^p(\Gamma(X',\psi^*(\mathcal L^\bullet)))
\otimes_{\Gamma(X',\psi^*(\mathcal O_X))}
H^q(\Gamma(X',\psi^*(\mathcal M^\bullet)))\ar[r]
&H^{p+q}(\Gamma(X',\psi^*(\mathcal N^\bullet)))
}
\]
Mais comme $\psi^*(\mathcal L^\bullet)$, $\psi^*(\mathcal M^\bullet)$ et
$\psi^*(\mathcal N^\bullet)$ sont des \emph{résolutions} de
$\psi^*(\mathcal F)$, $\psi^*(\mathcal G)$, $\psi^*(\mathcal H)$
respectivement, on a un diagramme commutatif (G, II, 6.6.1)
\[
\label{0.12.1.5.5-fr}
\tag{12.1.5.5}
\xymatrix@C=7pt{
H^p(\Gamma(X',\psi^*(\mathcal L^\bullet)))
\otimes_{\Gamma(X',\psi^*(\mathcal O_X))}
H^q(\Gamma(X',\psi^*(\mathcal M^\bullet)))\ar[r]\ar[d]
&H^{p+q}(\Gamma(X',\psi^*(\mathcal N^\bullet)))\ar[d]\\
H^p(X',\psi^*(\mathcal F))
\otimes_{\Gamma(X',\psi^*(\mathcal O_X))}
H^q(X',\psi^*(\mathcal G))\ar[r]
&H^{p+q}(X',\psi^*(\mathcal H))
}
\]
Enfin, par fonctorialité, on a un diagramme commutatif
\[
\label{0.12.1.5.6-fr}
\tag{12.1.5.6}
\xymatrix@C=7pt{
H^p(X',\psi^*(\mathcal F))
\otimes_{\Gamma(X',\psi^*(\mathcal O_X))}
H^q(X',\psi^*(\mathcal G))\ar[r]\ar[d]
&H^{p+q}(X',\psi^*(\mathcal H))\ar[d]\\
H^p(X',f^*(\mathcal F))\otimes_{A'}H^q(X',f^*(\mathcal G))\ar[r]
&H^{p+q}(X',f^*(\mathcal H))
}
\]
et par combinaison des trois diagrammes
(\hyperref[0.12.1.5.4-fr]{12.1.5.4}),
(\hyperref[0.12.1.5.5-fr]{12.1.5.5}) et
(\hyperref[0.12.1.5.6-fr]{12.1.5.6}), on obtient le diagramme commutatif
(\hyperref[0.12.1.5.2-fr]{12.1.5.2}) cherché.
\end{env}

\begin{remark}[12.1.6]
\label{0.12.1.6-fr}
\oldpage[0\textsubscript{III}]{55}
Avec les notations de (\hyperref[0.12.1.3-fr]{12.1.3}), supposons que l'on
ait un diagramme commutatif
\[
\label{0.12.1.6.1-fr}
\tag{12.1.6.1}
\xymatrix{
0\ar[r]&\mathcal F\ar[r]^r\ar[d]_u
&\mathcal G\ar[r]^s\ar[d]_v
&\mathcal H\ar[r]\ar[d]_w&0\\
0\ar[r]&\mathcal F'\ar[r]_{r'}
&\mathcal G'\ar[r]_{s'}
&\mathcal H'\ar[r]&0
}
\]
où $r$, $s$ sont des homomorphismes de $\mathcal O_X$-Modules, $r'$, $s'$ des
homomorphismes de $\mathcal O_{X'}$-Modules, $u$, $v$, $w$ des $f$-morphismes
et les lignes sont \emph{exactes}. On en déduit alors un diagramme commutatif
\[
\label{0.12.1.6.2-fr}
\tag{12.1.6.2}
\xymatrix@C=6pt{
\cdots\ar[r]&H^p(X,\mathcal F)\ar[r]\ar[d]_{u_p}
&H^p(X,\mathcal G)\ar[r]\ar[d]_{v_p}
&H^p(X,\mathcal H)\ar[r]^\partial\ar[d]_{w_p}
&H^{p+1}(X,\mathcal F)\ar[r]\ar[d]_{u_{p+1}}&\cdots\\
\cdots\ar[r]&H^p(X',\mathcal F')\ar[r]
&H^p(X',\mathcal G')\ar[r]
&H^p(X',\mathcal H')\ar[r]_\partial
&H^{p+1}(X',\mathcal F')\ar[r]&\cdots
}
\]
En effet, (\hyperref[0.12.1.6.1-fr]{12.1.6.1}) se factorise en
\[
\xymatrix@R=14pt{
0\ar[r]&\mathcal F\ar[r]\ar[d]
&\mathcal G\ar[r]\ar[d]
&\mathcal H\ar[r]\ar[d]&0\\
0\ar[r]&\psi^*(\mathcal F)\ar[r]\ar[d]
&\psi^*(\mathcal G)\ar[r]\ar[d]
&\psi^*(\mathcal H)\ar[r]\ar[d]&0\\
0\ar[r]&\mathcal F'\ar[r]
&\mathcal G'\ar[r]
&\mathcal H'\ar[r]&0
}
\]
où la ligne du milieu est \emph{exacte}
(\hyperref[0.3.7.2-fr]{0\textsubscript{I}, 3.7.2}) et il suffit d'utiliser le
fait que (\hyperref[0.12.1.3.2-fr]{12.1.3.2}) est un homomorphisme de
$\delta$-foncteurs et que les $H^p(X',\mathcal F')$ forment un
$\delta$-foncteur en $\mathcal F'$.
\end{remark}

\begin{env}[12.1.7]
\label{0.12.1.7-fr}
Les hypothèses et notations étant celles de
(\hyperref[0.12.1.3-fr]{12.1.3}), considérons maintenant le cas où
$\mathcal F=f_*(\mathcal F')=\psi_*(\mathcal F')$; nous allons voir que le
di-homomorphisme défini dans (\hyperref[0.12.1.3-fr]{12.1.3})
\[
\label{0.12.1.7.1-fr}
\tag{12.1.7.1}
H^p(X,f_*(\mathcal F'))\longrightarrow H^p(X',\mathcal F')
\]
peut s'obtenir (à un automorphisme près de $H^p(X',\mathcal F')$) comme
\emph{edge-homomorphisme d'une suite spectrale} du foncteur composé
$\mathcal F'\leadsto\Gamma(X,\psi_*(\mathcal F'))$ (T, 2.4). La description
de l'homomorphisme (\hyperref[0.12.1.7.1-fr]{12.1.7.1}) donnée dans
(\hyperref[0.12.1.4-fr]{12.1.4}) montre ici qu'on peut obtenir cet
homomorphisme de la façon suivante : on considère des résolutions injectives
$\mathcal L^\bullet$ et $\mathcal L'^\bullet$ de $\psi_*(\mathcal F')$ et de
$\mathcal F'$ respectivement, puis on prend un homomorphisme de complexes
$v:\psi^*(\mathcal L^\bullet)\to\mathcal L'^\bullet$ «~au-dessus~» de
l'homomorphisme canonique $\psi^*(\psi_*(\mathcal F'))\to\mathcal F'$;
\oldpage[0\textsubscript{III}]{56}
on note ensuite que l'on a
$\Gamma(X',\mathcal L'^\bullet)=\Gamma(X,\psi_*(\mathcal L'^\bullet))$ et
que l'homomorphisme composé
\[
\Gamma(X,\mathcal L^\bullet)\longrightarrow
\Gamma(X',\psi^*(\mathcal L^\bullet))
\xrightarrow{\Gamma(v)}\Gamma(X',\mathcal L'^\bullet)
\]
n'est autre que
\[
\label{0.12.1.7.2-fr}
\tag{12.1.7.2}
\Gamma(v^\flat):\Gamma(X,\mathcal L^\bullet)
\longrightarrow\Gamma(X,\psi_*(\mathcal L'^\bullet))
\]
(\hyperref[0.3.7.1-fr]{0\textsubscript{I}, 3.7.1}), et
(\hyperref[0.12.1.7.1-fr]{12.1.7.1}) s'obtient par passage à la cohomologie
dans (\hyperref[0.12.1.7.2-fr]{12.1.7.2}). D'autre part, les suites spectrales
du foncteur composé $\mathcal F'\leadsto\Gamma(X,\psi_*(\mathcal F'))$
s'obtiennent en considérant une résolution injective de Cartan-Eilenberg
$\mathcal M^{\bullet\bullet}=(\mathcal M^{ij})$ du complexe
$\psi_*(\mathcal L'^\bullet)$ dans la catégorie des faisceaux de groupes
abéliens sur $X$; les suites spectrales en question sont celles du bicomplexe
$\Gamma(X,\mathcal M^{\bullet\bullet})$ (qui sont birégulières puisque
$\mathcal M^{ij}=0$ pour $i<0$ ou $j<0$). Or, la première suite spectrale de
ce bicomplexe \emph{dégénère}, car les faisceaux
$\psi_*(\mathcal L'^i)$ sont flasques (G, II, 3.1.1), donc
$H_{\mathrm{II}}^q(\Gamma(X,\mathcal M^{i,\bullet}))
=H^q(\psi_*(\mathcal L'^i))=0$ pour $q>0$ (G, II, 4.4.3); on a donc des
edge-homomorphismes \emph{bijectifs} (\hyperref[0.11.1.6-fr]{11.1.6})
\[
\label{0.12.1.7.3-fr}
\tag{12.1.7.3}
'E_2^{i0}=H^i(H_{\mathrm{II}}^0(\Gamma(X,\mathcal M^{\bullet\bullet})))
\longrightarrow H^i(\Gamma(X,\mathcal M^{\bullet\bullet}))
\]
et on sait (\hyperref[0.11.3.4-fr]{11.3.4}) que cet homomorphisme provient,
par passage à la cohomologie, de l'augmentation
\[
\label{0.12.1.7.4-fr}
\tag{12.1.7.4}
\Gamma(X,\psi_*(\mathcal L'^\bullet))
\longrightarrow\Gamma(X,\mathcal M^{\bullet\bullet})
\]
laquelle provient elle-même de l'augmentation
$\eta:\psi_*(\mathcal L'^\bullet)\to\mathcal M^{\bullet\bullet}$. D'autre
part, pour la seconde suite spectrale, on a des edge-homomorphismes
\[
\label{0.12.1.7.5-fr}
\tag{12.1.7.5}
''E_2^{i0}=H^i(H_{\mathrm I}^0(\Gamma(X,\mathcal M^{\bullet\bullet})))
\longrightarrow H^i(\Gamma(X,\mathcal M^{\bullet\bullet}))
\]
provenant (\hyperref[0.11.3.4-fr]{11.3.4}) par passage à la cohomologie, de
l'homomorphisme de complexes
$Z_{\mathrm I}^0(\Gamma(X,\mathcal M^{\bullet\bullet}))
\to\Gamma(X,\mathcal M^{\bullet\bullet})$. Or, comme $\psi_*$ est exact à
gauche, la suite
\[
0\longrightarrow\psi_*(\mathcal F')\longrightarrow
\psi_*(\mathcal L'^0)\longrightarrow\psi_*(\mathcal L'^1)
\]
est exacte; par définition d'une résolution de Cartan-Eilenberg
(\hyperref[0.11.4.2-fr]{11.4.2}), on peut donc prendre
$B_{\mathrm I}^0(\mathcal M^{\bullet\bullet})=0$,
$Z_{\mathrm I}^0(\mathcal M^{\bullet\bullet})=\mathcal L^\bullet$; comme le
diagramme
\[
\xymatrix@C=50pt@R=24pt{
\mathcal L^0\ar[r]^{i^0}&\mathcal M^{00}\\
\psi_*(\mathcal F')\ar[u]^\varepsilon
\ar[r]_{\varepsilon'}\ar[ur]^{\varepsilon''}
&\psi_*(\mathcal L'^0)\ar[u]_{\eta^0}
}
\]
est commutatif, l'injection de complexes
$i:\mathcal L^\bullet\to\mathcal M^{\bullet\bullet}$ est compatible avec les
augmentations $\varepsilon$ et $\varepsilon''$. On a ainsi deux
homomorphismes de complexes de $\mathcal L^\bullet$ dans
$\mathcal M^{\bullet\bullet}$
\[
\xymatrix@C=52pt@R=22pt{
\mathcal L^\bullet\ar[r]^i\ar[dr]_{v^\flat}
&\mathcal M^{\bullet\bullet}\\
&\psi_*(\mathcal L'^\bullet)\ar[u]_\eta
}
\]
\oldpage[0\textsubscript{III}]{57}
compatibles avec les augmentations $\varepsilon$ et $\varepsilon''$; comme
$\mathcal L^\bullet$ est une \emph{résolution} injective et que
$\mathcal M^{\bullet\bullet}$ est formée de faisceaux injectifs, il résulte de
(M, V, 1.1a)) que ces deux homomorphismes sont \emph{homotopes}; il en est
donc de même des deux homomorphismes correspondants
$\Gamma(X,\mathcal L^\bullet)\to\Gamma(X,\mathcal M^{\bullet\bullet})$, et en
passant à la cohomologie, on obtient donc le \emph{même} homomorphisme; en
d'autres termes, on a bien montré que le edge-homomorphisme
(\hyperref[0.12.1.7.5-fr]{12.1.7.5}), qui s'écrit
$H^p(X,\psi_*(\mathcal F'))\to H^p(\Gamma(X,\mathcal M^{\bullet\bullet}))$
est composé de (\hyperref[0.12.1.7.1-fr]{12.1.7.1}) et de
(\hyperref[0.12.1.7.3-fr]{12.1.7.3}), qui s'écrit
$H^p(X',\mathcal F')\to H^p(\Gamma(X,\mathcal M^{\bullet\bullet}))$ et qu'on
a vu être un \emph{isomorphisme}; d'où notre assertion.
\end{env}

\subsection{Images directes supérieures}
\label{subsection:0.12.2-fr}

\begin{env}[12.2.1]
\label{0.12.2.1-fr}
Soient $(X,\mathcal O_X)$, $(Y,\mathcal O_Y)$ deux espaces annelés,
$f=(\psi,\theta)$ un morphisme de $X$ dans $Y$, qui définit le foncteur
image directe
\[
f_*:C(X)\longrightarrow C(Y),
\]
identique d'ailleurs à la restriction à $C(X)$ du foncteur $\psi_*$ défini
dans la catégorie des faisceaux de groupes abéliens sur $X$. Ce dernier
foncteur est additif et exact à gauche, et comme tout faisceau de groupes
abéliens sur $X$ est isomorphe à un sous-faisceau d'un faisceau injectif de
groupes abéliens, on définit les foncteurs dérivés droits
$\mathcal F\leadsto R^p\psi_*(\mathcal F)$ du foncteur $\psi_*$; les
$R^p\psi_*(\mathcal F)$ sont des faisceaux de groupes abéliens sur $Y$, et
les $R^p\psi_*$ forment un foncteur cohomologique universel (T, 2.3).

En outre, le faisceau $R^p\psi_*(\mathcal F)$ est le faisceau associé au
préfaisceau $V\leadsto H^p(f^{-1}(V),\mathcal F)$ (T, 3.7.2). Si maintenant
on suppose que $\mathcal F$ est un $\mathcal O_X$-Module,
$H^p(f^{-1}(V),\mathcal F)$ est naturellement muni d'une structure de
$\Gamma(f^{-1}(V),\mathcal O_X)$-module, donc de
$\Gamma(V,\psi_*(\mathcal O_X))$-module, et la donnée de l'homomorphisme
$\theta:\mathcal O_Y\to\psi_*(\mathcal O_X)$ permet d'en déduire une
structure de $\Gamma(V,\mathcal O_Y)$-module. Pour les structures ainsi
définies, il est clair que la restriction d'un ouvert $V$ à un ouvert
$V'\subset V$ définit un di-homomorphisme, et cela permet donc de définir sur
chacun des $R^p\psi_*(\mathcal F)$ une structure de $\mathcal O_Y$-Module;
c'est ce $\mathcal O_Y$-Module que nous noterons $R^p f_*(\mathcal F)$,
$R^p f_*$ étant ainsi défini comme un foncteur additif de $C(X)$ dans
$C(Y)$. En outre, les $R^p f_*$ forment un $\partial$-foncteur, car si
\[
0\longrightarrow\mathcal F'\longrightarrow\mathcal F
\longrightarrow\mathcal F''\longrightarrow0
\]
est une suite exacte de $\mathcal O_X$-Modules, la description des
$R^p\psi_*$ et de la structure de $\mathcal O_Y$-Module sur
$R^p\psi_*(\mathcal F)$ donnée ci-dessus montre aussitôt que
l'homomorphisme
\[
\partial:R^p\psi_*(\mathcal F'')\longrightarrow
R^{p+1}\psi_*(\mathcal F')
\]
est dans ce cas un homomorphisme de $\mathcal O_Y$-Modules. Enfin, les
$R^p f_*$ s'identifient aux foncteurs dérivés droits de $f_*$ : en effet,
tout $\mathcal O_X$-Module admet une résolution injective formée de
$\mathcal O_X$-Modules, et comme une telle résolution est formée de faisceaux
flasques de groupes abéliens (G, II, 7.1), elle peut servir à calculer les
$R^p\psi_*(\mathcal F)$, puisque $R^n\psi_*(\mathcal G)=0$ pour $n\geq1$ et
pour tout faisceau flasque $\mathcal G$ (T, 2.4.1, Remarque 3, et cor. de la
prop. 3.3.2). On conclut donc que les $R^p f_*$ forment un foncteur
cohomologique universel de $C(X)$ dans $C(Y)$ (T, 2.3).
\end{env}

\begin{env}[12.2.2]
\label{0.12.2.2-fr}
Soient $\mathcal F$ et $\mathcal G$ deux $\mathcal O_X$-Modules. Avec les
notations de (\hyperref[0.12.2.1-fr]{12.2.1}), pour tout ouvert $V$ de $Y$,
on a l'homomorphisme de cup-produit
(\hyperref[0.12.1.2.1-fr]{12.1.2.1})
\[
H^p(f^{-1}(V),\mathcal F)
\otimes_{\Gamma(f^{-1}(V),\mathcal O_X)}
H^q(f^{-1}(V),\mathcal G)
\longrightarrow
H^{p+q}(f^{-1}(V),\mathcal F\otimes_{\mathcal O_X}\mathcal G)
\]
\oldpage[0\textsubscript{III}]{58}
et il résulte aussitôt de la définition du cup-produit (G, II, 6.6) que ces
homomorphismes commutent au passage de $V$ à un sous-espace ouvert $V'$ de
$V$. D'autre part, on a un homomorphisme d'anneaux
\[
\Gamma(V,\mathcal O_Y)\longrightarrow\Gamma(V,\psi_*(\mathcal O_X))
=\Gamma(f^{-1}(V),\mathcal O_X)
\]
provenant de $\theta$, d'où un homomorphisme canonique de produits tensoriels
\[
\begin{aligned}
&H^p(f^{-1}(V),\mathcal F)\otimes_{\Gamma(V,\mathcal O_Y)}
H^q(f^{-1}(V),\mathcal G)\\
&\qquad\longrightarrow
H^p(f^{-1}(V),\mathcal F)
\otimes_{\Gamma(f^{-1}(V),\mathcal O_X)}
H^q(f^{-1}(V),\mathcal G)
\end{aligned}
\]
qui est lui aussi compatible avec la restriction de $V'$ à $V$. Par
composition, on obtient donc un homomorphisme de
$\Gamma(V,\mathcal O_Y)$-modules, qui définit un homomorphisme canonique
fonctoriel en $\mathcal F$ et $\mathcal G$ pour les faisceaux associés aux
préfaisceaux considérés :
\[
\label{0.12.2.2.1-fr}
\tag{12.2.2.1}
R^p f_*(\mathcal F)\otimes_{\mathcal O_Y}R^q f_*(\mathcal G)
\longrightarrow
R^{p+q}f_*(\mathcal F\otimes_{\mathcal O_X}\mathcal G).
\]
On notera que pour $p=q=0$, cet homomorphisme se réduit à
(\hyperref[0.4.2.2.1-fr]{0\textsubscript{I}, 4.2.2.1}).
\end{env}

\begin{proposition}[12.2.3]
\label{0.12.2.3-fr}
\emph{Pour tout $\mathcal O_X$-Module $\mathcal F$ et tout
$\mathcal O_Y$-Module localement libre de rang fini $\mathcal L$, on a des
isomorphismes canoniques fonctoriels
\[
\label{0.12.2.3.1-fr}
\tag{12.2.3.1}
R^p f_*(\mathcal F)\otimes_{\mathcal O_Y}\mathcal L
\xrightarrow{\sim}
R^p f_*(\mathcal F\otimes_{\mathcal O_X}f^*(\mathcal L)).
\]}
\end{proposition}

\begin{proof}
L'homomorphisme (\hyperref[0.12.2.3.1-fr]{12.2.3.1}) s'obtient en composant
l'homomorphisme, cas particulier de
(\hyperref[0.12.2.2.1-fr]{12.2.2.1}) :
\[
\label{0.12.2.3.2-fr}
\tag{12.2.3.2}
R^p f_*(\mathcal F)\otimes_{\mathcal O_Y}f_*(f^*(\mathcal L))
\longrightarrow
R^p f_*(\mathcal F\otimes_{\mathcal O_X}f^*(\mathcal L))
\]
avec l'homomorphisme du premier membre de
(\hyperref[0.12.2.3.1-fr]{12.2.3.1}) dans celui de
(\hyperref[0.12.2.3.2-fr]{12.2.3.2}), provenant de l'homomorphisme canonique
(\hyperref[0.4.4.3.2-fr]{0\textsubscript{I}, 4.4.3.2}). Pour vérifier que
(\hyperref[0.12.2.3.1-fr]{12.2.3.1}) est un isomorphisme lorsque $\mathcal L$
est localement libre, on peut aussitôt se ramener au cas
$\mathcal L=\mathcal O_Y$, la question étant locale sur $Y$, et les foncteurs
envisagés étant additifs en $\mathcal L$. Mais alors, la proposition se
ramène, vu la définition de (\hyperref[0.12.2.2.1-fr]{12.2.2.1}), à la
vérification du fait que l'homomorphisme correspondant de préfaisceaux est
bijectif, ce qui est immédiat en vertu de la relation
$f^*(\mathcal O_Y)=\mathcal O_X$.
\end{proof}

\begin{env}[12.2.4]
\label{0.12.2.4-fr}
Soient $(Z,\mathcal O_Z)$ un troisième espace annelé, $g:Y\to Z$ un
morphisme d'espaces annelés. On sait (G, II, 7.1 et 3.1.1) que pour tout
$\mathcal O_X$-Module injectif $\mathcal G$, $f_*(\mathcal G)$ est un
faisceau flasque de groupes abéliens, et par suite
(\hyperref[0.12.2.1-fr]{12.2.1}) on a
$R^p g_*(f_*(\mathcal G))=0$ pour tout $p>0$. Il s'ensuit (T, 2.4.1) que la
suite spectrale de Leray des foncteurs composés est applicable au foncteur
composé $g_*f_*$ : il y a une suite spectrale birégulière dont
l'aboutissement est le foncteur $R^\bullet h_*$, où $h=g\circ f$, et dont le
terme $E_2$ est donné par
\[
\label{0.12.2.4.1-fr}
\tag{12.2.4.1}
E_2^{pq}=R^p g_*(R^q f_*(\mathcal F)).
\]
\end{env}

\begin{env}[12.2.5]
\label{0.12.2.5-fr}
Sous les conditions de (\hyperref[0.12.2.4-fr]{12.2.4}), nous allons définir
directement des homomorphismes canoniques de $\mathcal O_Z$-Modules
\[
\label{0.12.2.5.1-fr}
\tag{12.2.5.1}
R^n g_*(f_*(\mathcal F))\longrightarrow R^n h_*(\mathcal F)
\]
\[
\label{0.12.2.5.2-fr}
\tag{12.2.5.2}
R^n h_*(\mathcal F)\longrightarrow g_*(R^n f_*(\mathcal F))
\]
\oldpage[0\textsubscript{III}]{59}
qu'on pourrait identifier aux «~edge-homomorphismes~» de la suite spectrale
de Leray (cf. (\hyperref[0.12.1.7-fr]{12.1.7})). Il suffit d'opérer sur les
préfaisceaux auxquels sont associés les faisceaux images supérieures
(\hyperref[0.12.2.1-fr]{12.2.1}). Pour cela, considérons un ouvert quelconque
$W$ de $Z$ et son image réciproque $g^{-1}(W)$ dans $Y$; on a un
di-homomorphisme canonique
\[
\label{0.12.2.5.3-fr}
\tag{12.2.5.3}
H^n(g^{-1}(W),f_*(\mathcal F))\longrightarrow
H^n(f^{-1}(g^{-1}(W)),f^*(f_*(\mathcal F)))
\]
les anneaux correspondants étant $\Gamma(g^{-1}(W),\mathcal O_Y)$ et
$\Gamma(h^{-1}(W),\mathcal O_X)$; d'autre part, l'homomorphisme canonique
(\hyperref[0.4.4.3.3-fr]{0\textsubscript{I}, 4.4.3.3}) fournit par
fonctorialité des homomorphismes canoniques
\[
\label{0.12.2.5.4-fr}
\tag{12.2.5.4}
H^n(h^{-1}(W),f^*(f_*(\mathcal F)))\longrightarrow
H^n(h^{-1}(W),\mathcal F)
\]
qui sont des homomorphismes de $\Gamma(h^{-1}(W),\mathcal O_X)$-modules.
Compte tenu de l'homomorphisme d'anneaux
$\Gamma(W,\mathcal O_Z)\to\Gamma(g^{-1}(W),\mathcal O_Y)$, on voit qu'en
composant (\hyperref[0.12.2.5.4-fr]{12.2.5.4}) et
(\hyperref[0.12.2.5.3-fr]{12.2.5.3}), on obtient un homomorphisme de
préfaisceaux, qui fournit l'homomorphisme de faisceaux
(\hyperref[0.12.2.5.1-fr]{12.2.5.1}).

La définition de (\hyperref[0.12.2.5.2-fr]{12.2.5.2}) est encore plus
simple; par définition, $R^n h_*(\mathcal F)$ est associé au préfaisceau
$W\leadsto H^n(f^{-1}(g^{-1}(W)),\mathcal F)$ et $R^n f_*(\mathcal F)$ au
préfaisceau $V\leadsto H^n(f^{-1}(V),\mathcal F)$; on a donc un homomorphisme
canonique
\[
H^n(f^{-1}(g^{-1}(W)),\mathcal F)\longrightarrow
\Gamma(g^{-1}(W),R^n f_*(\mathcal F)),
\]
et il est immédiat que ces homomorphismes définissent un homomorphisme de
préfaisceaux, qui à son tour définit
(\hyperref[0.12.2.5.2-fr]{12.2.5.2}).
\end{env}

\begin{env}[12.2.6]
\label{0.12.2.6-fr}
Sous les hypothèses de (\hyperref[0.12.2.4-fr]{12.2.4}), soient
$\mathcal F$, $\mathcal G$, $\mathcal K$ trois $\mathcal O_X$-Modules et
$u:\mathcal F\otimes\mathcal G\to\mathcal K$ un
$\mathcal O_X$-homomorphisme. On a alors des diagrammes commutatifs
\[
\xymatrix@C=10pt@R=24pt{
R^p g_*(f_*(\mathcal F))\otimes_{\mathcal O_Z}R^q g_*(f_*(\mathcal G))
\ar[r]\ar[d]&R^{p+q}g_*(f_*(\mathcal K))\ar[d]\\
R^p h_*(\mathcal F)\otimes_{\mathcal O_Z}R^q h_*(\mathcal G)
\ar[r]&R^{p+q}h_*(\mathcal K)
}
\tag{12.2.6.1}\label{0.12.2.6.1-fr}
\]
et
\[
\xymatrix@C=10pt@R=24pt{
R^p h_*(\mathcal F)\otimes_{\mathcal O_Z}R^q h_*(\mathcal G)
\ar[r]\ar[d]&R^{p+q}h_*(\mathcal K)\ar[d]\\
g_*(R^p f_*(\mathcal F))\otimes_{\mathcal O_Z}g_*(R^q f_*(\mathcal G))
\ar[r]&g_*(R^{p+q}f_*(\mathcal K))
}
\tag{12.2.6.2}\label{0.12.2.6.2-fr}
\]
\oldpage[0\textsubscript{III}]{60}
où les flèches horizontales proviennent de
(\hyperref[0.12.2.2.1-fr]{12.2.2.1}) (la dernière combinée avec
(\hyperref[0.4.2.2.1-fr]{0\textsubscript{I}, 4.2.2.1})) et les flèches
verticales des homomorphismes (\hyperref[0.12.2.5.1-fr]{12.2.5.1}) et
(\hyperref[0.12.2.5.2-fr]{12.2.5.2}) respectivement.

Il suffit en effet de le vérifier pour les homomorphismes correspondants de
préfaisceaux; si on revient aux définitions données dans
(\hyperref[0.12.2.2-fr]{12.2.2}) et
(\hyperref[0.12.2.5-fr]{12.2.5}) pour ces homomorphismes, on est aussitôt
ramené, pour (\hyperref[0.12.2.6.1-fr]{12.2.6.1}), aux diagrammes commutatifs
(\hyperref[0.12.1.5.2-fr]{12.1.5.2}); la vérification est encore plus simple
pour (\hyperref[0.12.2.6.2-fr]{12.2.6.2}).
\end{env}

\subsection{Compléments sur les foncteurs Ext de faisceaux}
\label{subsection:0.12.3-fr}

\begin{env}[12.3.1]
\label{0.12.3.1-fr}
Considérons un espace annelé $(X,\mathcal O_X)$; nous ne reviendrons pas sur
la définition et les principales propriétés des bifoncteurs
$\operatorname{Ext}_{\mathcal O_X}^p(X;\mathcal F,\mathcal G)$ de la
catégorie des $\mathcal O_X$-Modules dans celle des
$\Gamma(X,\mathcal O_X)$-modules, et
$\mathcal{E}xt_{\mathcal O_X}^p(\mathcal F,\mathcal G)$ de la catégorie des
$\mathcal O_X$-Modules dans elle-même, ni sur la suite spectrale birégulière
$E(\mathcal F,\mathcal G)$ qui les relie (T, 4.2 et G, II, 7.3).
\end{env}

\begin{env}[12.3.2]
\label{0.12.3.2-fr}
On définit de la même manière que dans (M, XIV, 1) la notion d'extension d'un
$\mathcal O_X$-Module $\mathcal F$ par un $\mathcal O_X$-Module $\mathcal G$
et la loi de composition entre classes d'extensions équivalentes : les
raisonnements faits pour les modules s'adaptent en effet de façon évidente à
une catégorie abélienne quelconque. La seconde démonstration de (M, XIV, 1.1),
qui n'utilise que l'existence de plongements dans les objets injectifs, est
alors encore valable pour la catégorie des $\mathcal O_X$-Modules, et montre
donc que $\operatorname{Ext}_{\mathcal O_X}^1(X;\mathcal F,\mathcal G)$
s'identifie canoniquement au groupe abélien des classes d'extensions de
$\mathcal F$ par $\mathcal G$.
\end{env}

\begin{proposition}[12.3.3]
\label{0.12.3.3-fr}
\emph{Soit $(X,\mathcal O_X)$ un espace annelé tel que le faisceau d'anneaux
$\mathcal O_X$ soit cohérent. Alors, pour tout couple de
$\mathcal O_X$-Modules cohérents $\mathcal F$, $\mathcal G$ et tout
$p\geq0$, $\mathcal{E}xt_{\mathcal O_X}^p(\mathcal F,\mathcal G)$ est un
$\mathcal O_X$-Module cohérent.}
\end{proposition}

\begin{proof}
Notons que les $\mathcal{E}xt_{\mathcal O_X}^p(\mathcal F,\mathcal G)$
forment un foncteur cohomologique contravariant en $\mathcal F$. Puisque
$\mathcal F$ est cohérent, il existe pour tout $p$ et tout point $x\in X$ un
voisinage ouvert $U$ de $X$ et une suite exacte de
$(\mathcal O_X|U)$-Modules
\[
0\longrightarrow\mathcal R\longrightarrow\mathcal L_{p-1}
\longrightarrow\cdots\longrightarrow\mathcal L_0
\longrightarrow\mathcal F|U\longrightarrow0
\]
où chacun des $\mathcal L_i$ ($0\leq i\leq p-1$) est isomorphe à un
$\mathcal O_X^{n_i}|U$ et $\mathcal R$ est cohérent : cela résulte par
récurrence sur $p$ de (\hyperref[0.5.3.2-fr]{0\textsubscript{I}, 5.3.2}) et
(\hyperref[0.5.3.4-fr]{0\textsubscript{I}, 5.3.4}), vu l'hypothèse que
$\mathcal O_X$ est cohérent.

Notons maintenant que, pour $p\geq1$, on a
\[
\operatorname{Ext}_{\mathcal O_X|U}^p(\mathcal L|U,\mathcal G|U)=0
\]
pour tout $\mathcal O_X$-Module $\mathcal L$ tel que $\mathcal L|U$ soit
isomorphe à un $\mathcal O_X^n|U$ (T, 4.2.3); le raisonnement de (M, V, 7.2)
s'applique donc au foncteur cohomologique contravariant
\[
\mathcal F\leadsto
\mathcal{H}om_{\mathcal O_X|U}(\mathcal F|U,\mathcal G|U),
\]
et donne une suite exacte
\[
\mathcal{H}om_{\mathcal O_X|U}(\mathcal L_{p-1},\mathcal G|U)
\longrightarrow
\mathcal{H}om_{\mathcal O_X|U}(\mathcal R,\mathcal G|U)
\longrightarrow
\mathcal{E}xt_{\mathcal O_X|U}^p(\mathcal F|U,\mathcal G|U)
\longrightarrow0
\]
et comme les deux premiers termes de cette suite sont des
$(\mathcal O_X|U)$-Modules cohérents
(\hyperref[0.5.3.5-fr]{0\textsubscript{I}, 5.3.5}), il en est de même du
troisième (\hyperref[0.5.3.4-fr]{0\textsubscript{I}, 5.3.4}).
\end{proof}

\oldpage[0\textsubscript{III}]{61}
\begin{proposition}[12.3.4]
\label{0.12.3.4-fr}
\emph{Soit $f:X\to Y$ un morphisme plat d'espaces annelés, et soient
$\mathcal F$, $\mathcal G$ deux $\mathcal O_Y$-Modules.}
\begin{enumerate}
\item[{\rm(i)}] \emph{Il existe un homomorphisme de bifoncteurs
cohomologiques
\[
\label{0.12.3.4.1-fr}
\tag{12.3.4.1}
f^*(\mathcal{E}xt_{\mathcal O_Y}^{\bullet}(\mathcal F,\mathcal G))
\xrightarrow{\sim}
\mathcal{E}xt_{\mathcal O_X}^{\bullet}(f^*(\mathcal F),f^*(\mathcal G))
\]
se réduisant en degré $0$ à l'homomorphisme canonique
(\hyperref[0.4.4.6-fr]{0\textsubscript{I}, 4.4.6}).}

\item[{\rm(ii)}] \emph{Il existe un morphisme canonique de suites spectrales
\[
\label{0.12.3.4.2-fr}
\tag{12.3.4.2}
E(\mathcal F,\mathcal G)\longrightarrow
E(f^*(\mathcal F),f^*(\mathcal G))
\]
qui, pour les termes $E_2$, se réduit aux homomorphismes
\[
\label{0.12.3.4.3-fr}
\tag{12.3.4.3}
H^p(Y,\mathcal{E}xt_{\mathcal O_Y}^q(\mathcal F,\mathcal G))
\longrightarrow
H^p(X,\mathcal{E}xt_{\mathcal O_X}^q(f^*(\mathcal F),f^*(\mathcal G)))
\]
déduits de (\hyperref[0.12.3.4.1-fr]{12.3.4.1}) et de
(\hyperref[0.12.1.3.1-fr]{12.1.3.1}).}
\end{enumerate}
\end{proposition}

\begin{proof}
\begin{enumerate}
\item[{\rm(i)}]
Comme $f^*$ est un foncteur exact dans la catégorie des
$\mathcal O_Y$-Modules (\hyperref[0.6.7.2-fr]{0\textsubscript{I}, 6.7.2}),
les foncteurs
\[
\mathcal G\leadsto
f^*(\mathcal{H}om_{\mathcal O_Y}(\mathcal F,\mathcal G))
\quad\hbox{et}\quad
\mathcal G\leadsto
\mathcal{H}om_{\mathcal O_X}(f^*(\mathcal F),f^*(\mathcal G))
\]
sont exacts à gauche; on déduit canoniquement de
(\hyperref[0.4.4.6-fr]{0\textsubscript{I}, 4.4.6}) un homomorphisme de leurs
foncteurs dérivés.

Pour calculer ces derniers, on prend une résolution injective
$\mathcal L^\bullet=(\mathcal L^i)$ de $\mathcal G$, et on a donc des
morphismes
\[
\mathcal H^p(f^*(\mathcal{H}om_{\mathcal O_Y}
(\mathcal F,\mathcal L^\bullet)))
\longrightarrow
\mathcal H^p(\mathcal{H}om_{\mathcal O_X}
(f^*(\mathcal F),f^*(\mathcal L^\bullet)))
\]
de cohomologies de complexes de faisceaux. D'ailleurs, en vertu de
l'exactitude de $f^*$, on a
\[
\mathcal H^p(f^*(\mathcal{H}om_{\mathcal O_Y}
(\mathcal F,\mathcal L^\bullet)))
=f^*(\mathcal H^p(\mathcal{H}om_{\mathcal O_Y}
(\mathcal F,\mathcal L^\bullet)))
=f^*(\mathcal{E}xt_{\mathcal O_Y}^p(\mathcal F,\mathcal G))
\]
par définition. D'autre part, l'exactitude de $f^*$ entraîne que
$f^*(\mathcal L^\bullet)$ est une résolution de $f^*(\mathcal G)$; si
$\mathcal L'^\bullet=(\mathcal L'^{i})$ est une résolution injective de
$f^*(\mathcal G)$ dans la catégorie des $\mathcal O_X$-Modules, il y a donc
un homomorphisme de complexes
$f^*(\mathcal L^\bullet)\to\mathcal L'^\bullet$, déterminé à une homotopie
près, et qui définit par suite un homomorphisme bien déterminé en cohomologie;
composant cet homomorphisme avec l'homomorphisme défini plus haut, on obtient
(\hyperref[0.12.3.4.1-fr]{12.3.4.1}).

\item[{\rm(ii)}]
Avec les notations précédentes, on a un homomorphisme de complexes de
faisceaux de $\mathcal O_X$-Modules
\[
f^*(\mathcal{H}om_{\mathcal O_Y}(\mathcal F,\mathcal L^\bullet))
\longrightarrow
\mathcal{H}om_{\mathcal O_X}(f^*(\mathcal F),\mathcal L'^\bullet).
\]
Soit $\mathcal M^{\bullet\bullet}$ une résolution injective de
Cartan-Eilenberg du complexe
$\mathcal{H}om_{\mathcal O_Y}(\mathcal F,\mathcal L^\bullet)$ dans la
catégorie des $\mathcal O_Y$-Modules; alors, en vertu de l'exactitude du
foncteur $f^*$, $f^*(\mathcal M^{\bullet\bullet})$ est une résolution de
Cartan-Eilenberg du complexe
$f^*(\mathcal{H}om_{\mathcal O_Y}(\mathcal F,\mathcal L^\bullet))$; si
$\mathcal M'^{\bullet\bullet}$ est une résolution injective de
Cartan-Eilenberg du complexe
$\mathcal{H}om_{\mathcal O_X}(f^*(\mathcal F),\mathcal L'^\bullet)$, il y a
donc (\hyperref[0.11.4.2-fr]{11.4.2}), un homomorphisme (déterminé à
homotopie près)
\[
f^*(\mathcal M^{\bullet\bullet})\longrightarrow
\mathcal M'^{\bullet\bullet}
\]
compatible avec l'homomorphisme considéré ci-dessus, autrement dit un
$f$-morphisme $\mathcal M^{\bullet\bullet}\to
\mathcal M'^{\bullet\bullet}$ de bicomplexes de faisceaux. On en déduit un
di-homomorphisme
$\Gamma(Y,\mathcal M^{\bullet\bullet})\to
\Gamma(X,\mathcal M'^{\bullet\bullet})$ de bicomplexes de modules, déterminé
à homotopie près, et un morphisme bien déterminé de suites spectrales
(\hyperref[0.11.3.2-fr]{11.3.2}), qui n'est autre que le morphisme
(\hyperref[0.12.3.4.2-fr]{12.3.4.2}) cherché, la caractérisation de
(\hyperref[0.12.3.4.3-fr]{12.3.4.3}) se déduisant aussitôt des définitions.
\end{enumerate}
\end{proof}

\begin{proposition}[12.3.5]
\label{0.12.3.5-fr}
\emph{Sous les hypothèses de (\hyperref[0.12.3.4-fr]{12.3.4}), supposons en
outre le faisceau d'anneaux $\mathcal O_Y$ cohérent; alors, pour tout
$\mathcal O_Y$-Module cohérent $\mathcal F$, les homomorphismes canoniques
(\hyperref[0.12.3.4.1-fr]{12.3.4.1}) sont bijectifs.}
\end{proposition}

\begin{proof}
\oldpage[0\textsubscript{III}]{62}
La question étant locale sur $Y$, on peut supposer qu'il existe une suite
exacte
\[
0\longrightarrow\mathcal R\longrightarrow\mathcal O_Y^n
\longrightarrow\mathcal F\longrightarrow0,
\]
et $\mathcal R$ est alors aussi un $\mathcal O_Y$-Module cohérent
(\hyperref[0.5.3.4-fr]{0\textsubscript{I}, 5.3.4}). Pour prouver que les
homomorphismes
\[
f^*(\mathcal{E}xt_{\mathcal O_Y}^p(\mathcal F,\mathcal G))
\longrightarrow
\mathcal{E}xt_{\mathcal O_X}^p(f^*(\mathcal F),f^*(\mathcal G))
\]
sont bijectifs, raisonnons par récurrence sur $p$, la proposition résultant de
(\hyperref[0.6.7.6.1-fr]{0\textsubscript{I}, 6.7.6.1}) lorsque $p=0$. Or, on
a le diagramme commutatif
\[
\hbox{\tiny$
\xymatrix@C=2pt@R=24pt{
f^*(\mathcal{E}xt_{\mathcal O_Y}^{p-1}(\mathcal O_Y^n,\mathcal G))
\ar[r]\ar[d]&
f^*(\mathcal{E}xt_{\mathcal O_Y}^{p-1}(\mathcal R,\mathcal G))
\ar[r]^\partial\ar[d]&
f^*(\mathcal{E}xt_{\mathcal O_Y}^{p}(\mathcal F,\mathcal G))
\ar[r]\ar[d]&
f^*(\mathcal{E}xt_{\mathcal O_Y}^{p}(\mathcal O_Y^n,\mathcal G))
\ar[d]\\
\mathcal{E}xt_{\mathcal O_X}^{p-1}(\mathcal O_X^n,f^*(\mathcal G))
\ar[r]&
\mathcal{E}xt_{\mathcal O_X}^{p-1}(f^*(\mathcal R),f^*(\mathcal G))
\ar[r]^\partial&
\mathcal{E}xt_{\mathcal O_X}^{p}(f^*(\mathcal F),f^*(\mathcal G))
\ar[r]&
\mathcal{E}xt_{\mathcal O_X}^{p}(\mathcal O_X^n,f^*(\mathcal G)).
}
$}
\]
puisque $f^*(\mathcal O_Y)=\mathcal O_X$; comme $f^*$ est exact, les deux
lignes sont exactes. En outre, on a
\[
\mathcal{E}xt_{\mathcal O_Y}^p(\mathcal O_Y^n,\mathcal G)=0
\quad\hbox{pour tout }p>0
\quad\hbox{et de même}\quad
\mathcal{E}xt_{\mathcal O_X}^p(\mathcal O_X^n,f^*(\mathcal G))=0
\quad\hbox{pour tout }p>0
\]
(T, 4.2.3). Vu l'hypothèse de récurrence, les deux premières flèches
verticales du diagramme précédent sont des isomorphismes, et les termes de
droite sont 0, donc
\[
f^*(\mathcal{E}xt_{\mathcal O_Y}^p(\mathcal F,\mathcal G))
\longrightarrow
\mathcal{E}xt_{\mathcal O_X}^p(f^*(\mathcal F),f^*(\mathcal G))
\]
est un isomorphisme.
\end{proof}

\subsection{Hypercohomologie du foncteur image directe}
\label{subsection:0.12.4-fr}

\begin{env}[12.4.1]
\label{0.12.4.1-fr}
Soient $(X,\mathcal O_X)$, $(Y,\mathcal O_Y)$ deux espaces annelés,
$f:X\to Y$ un morphisme d'espaces annelés. On peut prendre
l'\emph{hypercohomologie} de $f_*$ par rapport à un complexe quelconque de
$\mathcal O_X$-Modules
$\mathcal K^\bullet=(\mathcal K^i)_{i\in\mathbb Z}$
(\hyperref[0.11.4.4-fr]{11.4.4}), car dans la catégorie abélienne des
$\mathcal O_Y$-Modules, les limites inductives filtrantes existent et sont
exactes (T, 3.1.1). Les $\mathcal O_Y$-Modules d'hypercohomologie
$R^p f_*(\mathcal K^\bullet)$ se noteront aussi
$\mathcal H^p(f,\mathcal K^\bullet)$ ou
$\mathcal H_f^p(\mathcal K^\bullet)$. Rappelons que
$\mathcal H^\bullet(f,\mathcal K^\bullet)$ est la cohomologie du bicomplexe
de $\mathcal O_Y$-Modules $f_*(\mathcal L^{\bullet\bullet})$, où
$\mathcal L^{\bullet\bullet}$ est une résolution injective de
Cartan-Eilenberg de $\mathcal K^\bullet$ dans la catégorie des
$\mathcal O_X$-Modules~; $\mathcal H^\bullet(f,\mathcal K^\bullet)$ est
l'aboutissement de deux suites spectrales
${}'E(f,\mathcal K^\bullet)$ et ${}''E(f,\mathcal K^\bullet)$ dont les
termes $E_2$ sont donnés par
\begin{equation}
\label{0.12.4.1.1-fr}
\tag{12.4.1.1}
{}'E_2^{pq}=\mathcal H^p(\mathcal H^q(f,\mathcal K^\bullet))
\end{equation}
\begin{equation}
\label{0.12.4.1.2-fr}
\tag{12.4.1.2}
{}''E_2^{pq}=\mathcal H^p(f,\mathcal H^q(\mathcal K^\bullet))
\quad(=R^p f_*(\mathcal H^q(\mathcal K^\bullet))).
\end{equation}
On a, dans ces formules, adopté la notation générale $T(A^\bullet)$ pour
le transformé d'un complexe par un foncteur
(\hyperref[0.11.2.1-fr]{11.2.1}), et on écrit
$\mathcal H^p(f,\mathcal F)$ au lieu de $R^p f_*(\mathcal F)$ pour un
$\mathcal O_X$-Module $\mathcal F$. Rappelons encore que la suite
${}'E(f,\mathcal K^\bullet)$ est toujours \emph{régulière}~; les deux
suites spectrales ${}'E(f,\mathcal K^\bullet)$ et
${}''E(f,\mathcal K^\bullet)$ sont \emph{birégulières} lorsque
$\mathcal K^\bullet$ est limité inférieurement,
\oldpage[0\textsubscript{III}]{63}
ou lorsqu'il existe un entier $m$ tel que tout $\mathcal O_X$-Module
admette une résolution flasque de longueur $\leq m$
(\hyperref[0.11.4.4-fr]{11.4.4}).
\end{env}

\begin{env}[12.4.2]
\label{0.12.4.2-fr}
Nous désignerons de même par $\mathbf H^\bullet(X,\mathcal K^\bullet)$
l'hypercohomologie du foncteur $\Gamma$ par rapport à un complexe
$\mathcal K^\bullet$ de $\mathcal O_X$-Modules~; les
$\mathbf H^p(X,\mathcal K^\bullet)$ sont donc des
$\Gamma(X,\mathcal O_X)$-modules. On peut d'ailleurs considérer
$\mathbf H^\bullet(X,\mathcal K^\bullet)$ comme un cas particulier de
$\mathcal H^\bullet(f,\mathcal K^\bullet)$, où $f$ est un morphisme de
$(X,\mathcal O_X)$ sur un espace annelé réduit à un point muni de l'anneau
$\Gamma(X,\mathcal O_X)$.

Pour tout ouvert $V$ de $X$, nous écrirons
$\mathbf H^\bullet(V,\mathcal K^\bullet)$ au lieu de
$\mathbf H^\bullet(V,\mathcal K^\bullet|V)$.
\end{env}

\begin{proposition}[12.4.3]
\label{0.12.4.3-fr}
\emph{Pour tout entier $p\in\mathbb Z$, le $\mathcal O_Y$-Module
$\mathcal H^p(f,\mathcal K^\bullet)$ est canoniquement isomorphe au
faisceau associé au préfaisceau
$U\mapsto\mathbf H^p(f^{-1}(U),\mathcal K^\bullet)$ sur $Y$.}
\end{proposition}

\begin{proof}
En effet, avec les notations de
(\hyperref[0.12.4.1-fr]{12.4.1}), le faisceau de cohomologie
$\mathcal H^p(f_*(\mathcal L^{\bullet\bullet}))$ est associé au préfaisceau
\[
U\longmapsto H^p(\Gamma(U,f_*(\mathcal L^{\bullet\bullet})))
=H^p(\Gamma(f^{-1}(U),\mathcal L^{\bullet\bullet})).
\]
Mais il est clair que
$\mathcal L^{\bullet\bullet}|f^{-1}(U)$ est une résolution injective de
Cartan-Eilenberg de $\mathcal K^\bullet|f^{-1}(U)$ (T, 3.1.3), donc
\[
H^p(\Gamma(f^{-1}(U),\mathcal L^{\bullet\bullet}))
=\mathbf H^p(f^{-1}(U),\mathcal K^\bullet)
\]
par définition.
\end{proof}

\begin{proposition}[12.4.4]
\label{0.12.4.4-fr}
\emph{L'hypercohomologie $\mathcal H^\bullet(f,\mathcal K^\bullet)$ est
un foncteur cohomologique en $\mathcal K^\bullet$ dans chacun des cas
suivants~:}
\begin{enumerate}
\item[\emph{a)}]
\emph{$\mathcal K^\bullet$ varie dans la catégorie des complexes limités
inférieurement.}
\item[\emph{b)}]
\emph{Il existe un entier $m$ tel que tout $\mathcal O_X$-Module admette
une résolution flasque de longueur $\leq m$.}
\item[\emph{c)}]
\emph{$X$ est un espace noethérien.}
\end{enumerate}
\end{proposition}

\begin{proof}
Les cas a) et b) sont des cas particuliers de
(\hyperref[0.11.5.4-fr]{11.5.4}). D'autre part, le cas c) se déduit de
(\hyperref[0.11.5.2-fr]{11.5.2}), car on sait que dans ce cas, le foncteur
$f_*$ permute aux limites inductives (G, II, 3.10.1).
\end{proof}

\begin{env}[12.4.5]
\label{0.12.4.5-fr}
Considérons maintenant un recouvrement ouvert
$\mathfrak U=(U_\alpha)$ de $X$, et pour tout complexe de
\emph{préfaisceaux} $\mathcal K^\bullet=(\mathcal K^j)$ sur $X$, le
bicomplexe $\check C^\bullet(\mathfrak U,\mathcal K^\bullet)$, dont le
composant d'indices $(i,j)$ est
$\check C^i(\mathfrak U,\mathcal K^j)$, groupe des $i$-cochaînes alternées
du nerf de $\mathfrak U$ à valeurs dans $\mathcal K^j$ (G, II, 5.1). Nous
dirons que la cohomologie de ce bicomplexe est
l'\emph{hypercohomologie du recouvrement $\mathfrak U$ à coefficients dans
$\mathcal K^\bullet$}, et nous la noterons
\[
\mathbf H^\bullet(\mathfrak U,\mathcal K^\bullet)
=H^\bullet(\check C^\bullet(\mathfrak U,\mathcal K^\bullet)).
\]
La suite spectrale de Leray d'un recouvrement (T, 3.8.1 et G, II, 5.9.1)
se généralise comme suit à l'hypercohomologie~:
\end{env}

\begin{proposition}[12.4.6]
\label{0.12.4.6-fr}
\emph{Soit $\mathcal K^\bullet=(\mathcal K^i)$ un complexe de
$\mathcal O_X$-Modules. Il existe un foncteur spectral régulier en
$\mathcal K^\bullet$ ayant pour aboutissement l'hypercohomologie
$\mathbf H^\bullet(X,\mathcal K^\bullet)$, et dont le terme $E_2$ est donné
par}
\begin{equation}
\label{0.12.4.6.1-fr}
\tag{12.4.6.1}
E_2^{pq}=\check H^p(\mathfrak U,h^q(\mathcal K^\bullet)),
\end{equation}
\emph{où $h^q(\mathcal K^\bullet)$ désigne le complexe de préfaisceaux
$V\mapsto H^q(V,\mathcal K^\bullet)$ sur $X$. La suite spectrale précédente
est birégulière si $\mathcal K^\bullet$ est limité inférieurement.}
\end{proposition}

\begin{proof}
Considérons une résolution injective de Cartan-Eilenberg
$\mathcal L^{\bullet\bullet}$ de $\mathcal K^\bullet$, et le tricomplexe
$\check C^\bullet(\mathfrak U,\mathcal L^{\bullet\bullet})
=(\check C^i(\mathfrak U,\mathcal L^{jk}))$~; considérons d'abord ce
tricomplexe comme un bicomplexe pour les degrés $i$ et $j+k$. Comme $i$ ne
prend que des valeurs $\geq0$, la seconde suite spectrale de ce bicomplexe
est régulière (\hyperref[0.11.3.3-fr]{11.3.3}) et dégénérée, car on a
$\check H^q(\mathfrak U,\mathcal L^{jk})=0$ pour tout $q>0$, les
$\mathcal O_X$-Modules $\mathcal L^{jk}$ étant des faisceaux
\oldpage[0\textsubscript{III}]{64}
flasques (G, II, 5.2.3). On a par suite
(\hyperref[0.11.1.6-fr]{11.1.6}) un isomorphisme canonique
\[
H^n(\check C^\bullet(\mathfrak U,\mathcal L^{\bullet\bullet}))
\xrightarrow{\sim}
H^n(\Gamma(X,\mathcal L^{\bullet\bullet}))
\]
(en vertu de (G, II, 5.2.2)), donc par définition
(\hyperref[0.12.4.2-fr]{12.4.2}) un isomorphisme
\[
H^n(\check C^\bullet(\mathfrak U,\mathcal L^{\bullet\bullet}))
\xrightarrow{\sim}\mathbf H^n(X,\mathcal K^\bullet).
\]
Considérons d'autre part le tricomplexe
$\check C^\bullet(\mathfrak U,\mathcal L^{\bullet\bullet})$ comme un
bicomplexe pour les degrés $i+j$ et $k$. Comme $k$ ne prend que des valeurs
$\geq0$, la première suite spectrale de ce bicomplexe est toujours
régulière~; elle est birégulière si $\mathcal L^{jk}=0$ pour $j<j_0$,
c'est-à-dire lorsque $\mathcal K^\bullet$ est limité inférieurement
(\hyperref[0.11.3.3-fr]{11.3.3}). Cette suite spectrale est la suite
cherchée~; en effet, pour tout $j$, $\mathcal L^{j,\bullet}$ est une
résolution injective de $\mathcal K^j$~; par suite,
$H^q(\check C^i(\mathfrak U,\mathcal L^{j,\bullet}))$ n'est autre que le
complexe de cochaînes
$\check C^i(\mathfrak U,h^q(\mathcal K^j))$, ce qui termine la
démonstration.
\end{proof}

\begin{corollary}[12.4.7]
\label{0.12.4.7-fr}
\emph{Si, pour tout simplexe $\sigma$ du nerf de $\mathfrak U$, et pour
tout entier $i$, on a
$H^q(U_\sigma,\mathcal K^i)=0$ pour $q>0$, alors on a un isomorphisme
canonique}
\begin{equation}
\label{0.12.4.7.1-fr}
\tag{12.4.7.1}
\mathbf H^\bullet(\mathfrak U,\mathcal K^\bullet)
\xrightarrow{\sim}\mathbf H^\bullet(X,\mathcal K^\bullet).
\end{equation}
\end{corollary}

\begin{proof}
En effet, l'hypothèse entraîne que
$\check C^\bullet(\mathfrak U,h^q(\mathcal K^i))=0$ pour $q>0$, donc
$E_2^{pq}=0$ pour $q>0$~; la suite
(\hyperref[0.12.4.6.1-fr]{12.4.6.1}) étant dégénérée et régulière, la
conclusion résulte de la définition
(\hyperref[0.12.4.5-fr]{12.4.5}) de
$\mathbf H^\bullet(\mathfrak U,\mathcal K^\bullet)$ et de
(\hyperref[0.11.1.6-fr]{11.1.6}).
\end{proof}

\begin{env}[12.4.8]
\label{0.12.4.8-fr}
Soit $(X',\mathcal O_{X'})$ un second espace annelé, et soit
$f=(\psi,\theta)$ un morphisme de $X'$ dans $X$. Par la même méthode que
dans (\hyperref[0.12.1.3-fr]{12.1.3}) et
(\hyperref[0.12.1.4-fr]{12.1.4}), on définit un di-homomorphisme pour
l'hypercohomologie d'un complexe $\mathcal K^\bullet$ de
$\mathcal O_X$-Modules
\begin{equation}
\label{0.12.4.8.1-fr}
\tag{12.4.8.1}
\mathbf H^p(X,\mathcal K^\bullet)
\longrightarrow\mathbf H^p(X',f^*(\mathcal K^\bullet)).
\end{equation}
On part d'une résolution injective de Cartan-Eilenberg
$\mathcal L^{\bullet\bullet}$ de $\mathcal K^\bullet$, et comme $\psi^*$
est exact, $\psi^*(\mathcal L^{\bullet\bullet})$ est une résolution de
Cartan-Eilenberg de $\psi^*(\mathcal K^\bullet)$ dans la catégorie des
$\psi^*(\mathcal O_X)$-Modules~; il y a alors un morphisme
$\psi^*(\mathcal L^{\bullet\bullet})\to
\mathcal L'^{\bullet\bullet}$, où $\mathcal L'^{\bullet\bullet}$ est une
résolution injective de Cartan-Eilenberg de
$\psi^*(\mathcal K^\bullet)$, et on en déduit un morphisme pour la
cohomologie~:
\[
\mathbf H^\bullet(X,\mathcal K^\bullet)
\longrightarrow\mathbf H^\bullet(X',\psi^*(\mathcal K^\bullet))~;
\]
par composition avec le morphisme déduit par fonctorialité de
$\psi^*(\mathcal K^\bullet)\to f^*(\mathcal K^\bullet)$, on obtient le
morphisme (\hyperref[0.12.4.8.1-fr]{12.4.8.1}) cherché.

Partant de (\hyperref[0.12.4.8.1-fr]{12.4.8.1}) et de
(\hyperref[0.12.4.3-fr]{12.4.3}), on peut alors, en raisonnant comme dans
(\hyperref[0.12.2.5-fr]{12.2.5}), définir, pour deux morphismes
$f:X\to Y$, $g:Y\to Z$ d'espaces annelés, des homomorphismes pour
l'hypercohomologie d'un complexe $\mathcal K^\bullet$ de
$\mathcal O_X$-Modules
\begin{equation}
\label{0.12.4.8.2-fr}
\tag{12.4.8.2}
\mathcal H^n(g,f_*(\mathcal K^\bullet))
\longrightarrow\mathcal H^n(h,\mathcal K^\bullet)
\end{equation}
\begin{equation}
\label{0.12.4.8.3-fr}
\tag{12.4.8.3}
\mathcal H^n(h,\mathcal K^\bullet)
\longrightarrow g_*(\mathcal H^n(f,\mathcal K^\bullet)).
\end{equation}
Nous laissons au lecteur le détail des définitions.
\end{env}

\section{Limites projectives en algèbre homologique}
\label{section:0.13-fr}

\subsection{La condition de Mittag-Leffler}
\label{subsection:0.13.1-fr}

\begin{env}[13.1.1]
\label{0.13.1.1-fr}
Soit $\mathcal C$ une catégorie abélienne dans laquelle les produits infinis
existent (axiome AB~3* de T, 1.5))~; alors la borne inférieure d'une famille
de sous-objets d'un objet de $\mathcal C$ existe, et tout système projectif
d'objets de $\mathcal C$ admet une limite projective, qui est un foncteur
exact à gauche du système projectif considéré (T, 1.8). Soit
$(A_\alpha,f_{\alpha\beta})$
\oldpage[0\textsubscript{III}]{65}
un système projectif d'objets de $\mathcal C$ dont l'ensemble d'indices $I$
est filtrant à droite~; soient $A=\varprojlim A_\alpha$ et, pour tout
$\alpha\in I$, soit $f_\alpha:A\to A_\alpha$ le morphisme canonique. Pour
tout $\alpha\in I$, les $f_{\alpha\beta}(A_\beta)$ pour
$\alpha\leq\beta$ forment une famille filtrante décroissante de sous-objets
de $A_\alpha$~; le sous-objet
\[
A'_\alpha=\inf_{\beta\geq\alpha}f_{\alpha\beta}(A_\beta)
\]
est dit sous-objet des «~images universelles~» dans $A_\alpha$~; il est
clair que $f_\alpha(A)\subset A'_\alpha$ et
$f_{\alpha\beta}(A'_\beta)\subset A'_\alpha$ pour
$\alpha\leq\beta$~; donc
$(A'_\alpha,f_{\alpha\beta}|A'_\beta)$ est un système projectif et
$A=\varprojlim A'_\alpha$.
\end{env}

\begin{env}[13.1.2]
\label{0.13.1.2-fr}
Étant donné un système projectif $(A_\alpha,f_{\alpha\beta})$ dans
$\mathcal C$, on appelle \emph{condition de Mittag-Leffler} la condition
suivante~:
\begin{itemize}
\item[\textbf{(ML)}]
\emph{Pour tout indice $\alpha$, il existe $\beta\geq\alpha$ tel que, pour
tout $\gamma\geq\beta$, on ait
$f_{\alpha\gamma}(A_\gamma)=f_{\alpha\beta}(A_\beta)$.}
\end{itemize}
Il est clair que si les $f_{\alpha\beta}$ sont des \emph{épimorphismes}, la
condition (ML) est vérifiée. Inversement, si (ML) est vérifiée, et si pour
tout $\alpha\in I$, $A'_\alpha$ est le sous-objet des «~images
universelles~» dans $A_\alpha$, la restriction de $f_{\alpha\beta}$ à
$A'_\beta$ est un \emph{épimorphisme}
$A'_\beta\to A'_\alpha$ pour $\alpha\leq\beta$~: en effet, si
$\gamma\geq\beta$ est tel que
$f_{\beta\delta}(A_\delta)=f_{\beta\gamma}(A_\gamma)$ pour
$\delta\geq\gamma$, on a
$A'_\beta=f_{\beta\gamma}(A_\gamma)$ et cela entraîne d'autre part
$f_{\alpha\delta}(A_\delta)=f_{\alpha\gamma}(A_\gamma)$ pour
$\delta\geq\gamma$, donc
$A'_\alpha=f_{\alpha\gamma}(A_\gamma)=f_{\alpha\beta}(A'_\beta)$.

Notons aussi que la condition (ML) est vérifiée lorsque les objets
$A_\alpha$ sont \emph{artiniens} dans $\mathcal C$, c'est-à-dire que toute
famille de sous-objets de $A_\alpha$ admet un élément minimal~: un élément
minimal de la famille filtrante décroissante
$(f_{\alpha\beta}(A_\beta))$ de sous-objets de $A_\alpha$ est en effet
alors nécessairement le plus petit de ces sous-objets.
\end{env}

\begin{remark}[13.1.3]
\label{0.13.1.3-fr}
La condition (ML) peut se formuler également lorsque $\mathcal C$ est par
exemple la catégorie des ensembles~; on peut alors encore définir le
sous-ensemble des «~images universelles~» de $A_\alpha$ et les remarques
faites à ce sujet dans (\hyperref[0.13.1.1-fr]{13.1.1}) et
(\hyperref[0.13.1.2-fr]{13.1.2}) restent valables.
\end{remark}

\subsection{La condition de Mittag-Leffler pour les groupes abéliens}
\label{subsection:0.13.2-fr}

\begin{proposition}[13.2.1]
\label{0.13.2.1-fr}
\emph{Soit}
\[
0\longrightarrow A_\alpha\xrightarrow{u_\alpha}B_\alpha
\xrightarrow{v_\alpha}C_\alpha\longrightarrow0
\]
\emph{une suite exacte de systèmes projectifs de groupes abéliens (relatifs
à un même ensemble filtrant d'indices $I$).}
\begin{enumerate}
\item[\emph{(i)}]
\emph{Si $(B_\alpha)$ vérifie (ML), il en est de même de $(C_\alpha)$.}
\item[\emph{(ii)}]
\emph{Si $(A_\alpha)$ et $(C_\alpha)$ vérifient (ML), il en est de même de
$(B_\alpha)$.}
\end{enumerate}
\end{proposition}

\begin{proof}
Soient $(f_{\alpha\beta})$, $(g_{\alpha\beta})$, $(h_{\alpha\beta})$ les
systèmes d'homomorphismes définissant les systèmes projectifs
$(A_\alpha)$, $(B_\alpha)$, $(C_\alpha)$ respectivement.

(i) Supposons que
$g_{\alpha\beta}(B_\beta)=g_{\alpha\lambda}(B_\lambda)$ pour
$\lambda\geq\beta$~; comme $v_\beta$ et $v_\lambda$ sont surjectives, on a
\[
h_{\alpha\beta}(C_\beta)
=v_\alpha(g_{\alpha\beta}(B_\beta))
=v_\alpha(g_{\alpha\lambda}(B_\lambda))
=h_{\alpha\lambda}(C_\lambda)
\]
pour $\lambda\geq\beta$.

(ii) Soit $\alpha\in I$, et soit $\beta\geq\alpha$ un indice tel que pour
$\lambda\geq\beta$, on ait
$f_{\alpha\beta}(A_\beta)=f_{\alpha\lambda}(A_\lambda)$~; soit d'autre
part $\gamma\geq\beta$ un indice tel que, pour $\lambda\geq\gamma$, on ait
$h_{\beta\gamma}(C_\gamma)=h_{\beta\lambda}(C_\lambda)$. Soit alors
$y_\alpha$ un élément de $g_{\alpha\gamma}(B_\gamma)$~; on a donc
$y_\alpha=g_{\alpha\gamma}(y_\gamma)$ avec $y_\gamma\in B_\gamma$~;
posons $y_\beta=g_{\beta\gamma}(y_\gamma)$, de sorte que
$v_\beta(y_\beta)=h_{\beta\gamma}(v_\gamma(y_\gamma))$. Pour tout
$\lambda\geq\gamma$, il existe par hypothèse $y_\lambda\in B_\lambda$ tel
que
$h_{\beta\gamma}(v_\gamma(y_\gamma))
=h_{\beta\lambda}(v_\lambda(y_\lambda))
=v_\beta(g_{\beta\lambda}(y_\lambda))$, d'où
$v_\beta(y_\beta-g_{\beta\lambda}(y_\lambda))=0$, et par suite il existe
$x_\beta\in A_\beta$
\oldpage[0\textsubscript{III}]{66}
tel que
$y_\beta=g_{\beta\lambda}(y_\lambda)+u_\beta(x_\beta)$. On en déduit
\[
y_\alpha=g_{\alpha\lambda}(y_\lambda)
+u_\alpha(f_{\alpha\beta}(x_\beta))~;
\]
mais comme $\lambda\geq\beta$, il existe $x_\lambda\in A_\lambda$ tel que
$f_{\alpha\beta}(x_\beta)=f_{\alpha\lambda}(x_\lambda)$, et finalement
\[
y_\alpha
=g_{\alpha\lambda}(y_\lambda)+u_\alpha(f_{\alpha\lambda}(x_\lambda))
=g_{\alpha\lambda}(y_\lambda+u_\lambda(x_\lambda))
\in g_{\alpha\lambda}(B_\lambda),
\]
ce qui achève la démonstration.
\end{proof}

\begin{proposition}[13.2.2]
\label{0.13.2.2-fr}
\emph{Soit $I$ un ensemble ordonné filtrant ayant une partie cofinale
dénombrable. Soit}
\[
0\longrightarrow A_\alpha\xrightarrow{u_\alpha}B_\alpha
\xrightarrow{v_\alpha}C_\alpha\longrightarrow0
\]
\emph{une suite exacte de systèmes projectifs de groupes abéliens ayant
$I$ pour ensemble d'indices. Si $(A_\alpha)$ satisfait à la condition
(ML), la suite}
\[
0\longrightarrow\varprojlim A_\alpha
\longrightarrow\varprojlim B_\alpha
\longrightarrow\varprojlim C_\alpha\longrightarrow0
\]
\emph{est exacte.}
\end{proposition}

\begin{proof}
Tout revient à prouver que l'homomorphisme
\[
v=\varprojlim v_\alpha:\varprojlim B_\alpha\longrightarrow
\varprojlim C_\alpha
\]
est surjectif. Soit $z=(z_\alpha)$ un élément de
$\varprojlim C_\alpha$, et posons
$E_\alpha=v_\alpha^{-1}(z_\alpha)$~; il est clair que les $E_\alpha$
forment un système projectif d'ensembles non vides pour les restrictions
des homomorphismes $g_{\alpha\beta}:B_\beta\to B_\alpha$. Montrons que ce
système projectif vérifie la condition (ML)~; identifiant $A_\alpha$ à une
partie de $B_\alpha$ par $u_\alpha$, pour tout $\alpha\in I$, il existe
$\beta\geq\alpha$ tel que
$g_{\alpha\beta}(A_\beta)=g_{\alpha\lambda}(A_\lambda)$ pour
$\lambda\geq\beta$~; montrons que l'on a aussi
$g_{\alpha\beta}(E_\beta)=g_{\alpha\lambda}(E_\lambda)$ pour
$\lambda\geq\beta$. En effet, prenons un $y_\lambda\in E_\lambda$ et
posons $y_\beta=g_{\beta\lambda}(y_\lambda)$,
$y_\alpha=g_{\alpha\lambda}(y_\lambda)$~; soit
$y'_\alpha\in g_{\alpha\beta}(E_\beta)$, de sorte que
$y'_\alpha=g_{\alpha\beta}(y'_\beta)$ pour un $y'_\beta\in E_\beta$~;
on a $y'_\beta-y_\beta=x_\beta\in A_\beta$, et par hypothèse il existe
$x_\lambda\in A_\lambda$ tel que
$g_{\alpha\beta}(x_\beta)=g_{\alpha\lambda}(x_\lambda)$~; donc
\[
y'_\alpha=g_{\alpha\beta}(y_\beta)+g_{\alpha\beta}(x_\beta)
=g_{\alpha\lambda}(y_\lambda)+g_{\alpha\lambda}(x_\lambda)
=g_{\alpha\lambda}(y_\lambda+x_\lambda)
\in g_{\alpha\lambda}(E_\lambda),
\]
ce qui démontre notre assertion. Cela étant, on sait (Bourbaki,
\emph{Top. gén.}, chap.~II, 3\textsuperscript{e} éd., §~3, th.~1) que sous
les hypothèses faites sur $I$, un système projectif d'ensembles non vides
vérifiant (ML) a une limite projective non vide~; par suite, il existe un
point $y=(y_\alpha)\in\varprojlim E_\alpha$, et comme
$v_\alpha(y_\alpha)=z_\alpha$ par définition pour tout $\alpha$, on a
$z=v(y)$, C.Q.F.D.
\end{proof}

\begin{proposition}[13.2.3]
\label{0.13.2.3-fr}
\emph{Les hypothèses sur $I$ étant celles de
(\hyperref[0.13.2.2-fr]{13.2.2}), soit
$(K_\alpha^\bullet)_{\alpha\in I}$ un système projectif de complexes de
groupes abéliens
$K_\alpha^\bullet=(K_\alpha^n)_{n\in\mathbb Z}$ dont l'opérateur de
dérivation est de degré $+1$. Pour chaque $n$, il existe un homomorphisme
canonique fonctoriel}
\begin{equation}
\label{0.13.2.3.1-fr}
\tag{13.2.3.1}
h_n:H^n(\varprojlim K_\alpha^\bullet)
\longrightarrow\varprojlim H^n(K_\alpha^\bullet).
\end{equation}
\emph{Si, pour tout degré $n$, le système projectif de groupes abéliens
$(K_\alpha^n)_{\alpha\in I}$ vérifie (ML), alors tous les homomorphismes
$h_n$ sont surjectifs. Si en outre, pour un degré $n$, le système
projectif $(H^{n-1}(K_\alpha^\bullet))_{\alpha\in I}$ vérifie (ML),
l'homomorphisme $h_n$ est bijectif.}
\end{proposition}

\begin{proof}
Posons, pour tout $n$, $K^n=\varprojlim K_\alpha^n$~; la définition des
homomorphismes $h_n$ provient de la commutativité des diagrammes
\[
\xymatrix{
\cdots\ar[r]&K^{n-1}\ar[r]\ar[d]&K^n\ar[r]\ar[d]
&K^{n+1}\ar[r]\ar[d]&\cdots\\
\cdots\ar[r]&K_\alpha^{n-1}\ar[r]&K_\alpha^n\ar[r]
&K_\alpha^{n+1}\ar[r]&\cdots
}
\]
\oldpage[0\textsubscript{III}]{67}
les opérateurs de dérivation dans $K^\bullet$ étant limites projectives
des opérateurs correspondants dans les $K_\alpha^\bullet$.

Considérons les suites exactes
\[
(*_n)\qquad
0\longrightarrow \mathrm B^n(K_\alpha^\bullet)
\longrightarrow \mathrm Z^n(K_\alpha^\bullet)
\longrightarrow H^n(K_\alpha^\bullet)\longrightarrow0
\]
\[
(**_n)\qquad
0\longrightarrow \mathrm Z^{n-1}(K_\alpha^\bullet)
\longrightarrow K_\alpha^{n-1}
\longrightarrow \mathrm B^n(K_\alpha^\bullet)\longrightarrow0.
\]
L'hypothèse et la prop.
(\hyperref[0.13.2.1-fr]{13.2.1}, (i)) montrent que le système projectif
$(\mathrm B^n(K_\alpha^\bullet))_{\alpha\in I}$ vérifie (ML) pour tout
$n$~; il résulte donc de (\hyperref[0.13.2.2-fr]{13.2.2}) que la suite
\[
(***_n)\qquad
0\longrightarrow\varprojlim_\alpha\mathrm B^n(K_\alpha^\bullet)
\longrightarrow\varprojlim_\alpha\mathrm Z^n(K_\alpha^\bullet)
\longrightarrow\varprojlim_\alpha H^n(K_\alpha^\bullet)
\longrightarrow0
\]
est exacte. Or il est clair que
$\varprojlim_\alpha\mathrm B^n(K_\alpha^\bullet)$ s'identifie à un
sous-groupe de $K^{n+1}$ contenant $\mathrm B^n(K^\bullet)$ et que
$\varprojlim_\alpha\mathrm Z^n(K_\alpha^\bullet)$ s'identifie à un
sous-groupe de $\mathrm Z^n(K^\bullet)$~; par suite, $h_n$ est surjectif.
Si maintenant, on suppose en outre que le système projectif
$(H^{n-1}(K_\alpha^\bullet))_{\alpha\in I}$ vérifie (ML), les suites
exactes $(*_{n-1})$ et la prop.
(\hyperref[0.13.2.1-fr]{13.2.1}, (ii)) montrent que le système projectif
$(\mathrm Z^{n-1}(K_\alpha^\bullet))_{\alpha\in I}$ vérifie (ML)~; mais
alors, (\hyperref[0.13.2.2-fr]{13.2.2}) appliquée aux suites exactes
$(**_n)$ montre que la suite
\[
0\longrightarrow
\varprojlim_\alpha\mathrm Z^{n-1}(K_\alpha^\bullet)
\longrightarrow K^{n-1}\xrightarrow{u}
\varprojlim_\alpha\mathrm B^n(K_\alpha^\bullet)
\longrightarrow0
\]
est exacte~; comme
$\varprojlim_\alpha\mathrm B^n(K_\alpha^\bullet)
\supset\mathrm B^n(K^\bullet)$, et que la composée de l'injection
$\varprojlim_\alpha\mathrm B^n(K_\alpha^\bullet)\to K^n$ et de $u$ est
l'opérateur de dérivation $K^{n-1}\to K^n$, le fait que $u$ soit surjectif
entraîne
$\varprojlim_\alpha\mathrm B^n(K_\alpha^\bullet)
=\mathrm B^n(K^\bullet)$, donc $h_n$ est injectif. C.Q.F.D.
\end{proof}

\begin{env}[13.2.4]
\label{0.13.2.4-fr}
\emph{Remarques} (13.2.4).
\begin{enumerate}
\item[\emph{(i)}]
Le raisonnement de (\hyperref[0.13.2.2-fr]{13.2.2}) (cf. Bourbaki,
\emph{loc. cit.}) montre que la conclusion de cette proposition reste
valable lorsqu'on suppose seulement que les $A_\alpha$ peuvent être munis
de structures d'espaces \emph{métrisables complets}, dans lesquels les
translations sont des homéomorphismes, que les applications
$f_{\alpha\beta}:A_\beta\to A_\alpha$ définissant le système projectif
$(A_\alpha)$ sont \emph{uniformément continues} pour les distances
considérées, et enfin que le système $(A_\alpha)$ vérifie la condition
\begin{itemize}
\item[\textbf{(ML')}]
\emph{Pour tout indice $\alpha$, il existe $\beta\geq\alpha$ tel que pour
tout $\gamma\geq\beta$, $f_{\alpha\gamma}(A_\gamma)$ est dense dans
$f_{\alpha\beta}(A_\beta)$.}
\end{itemize}

Cela permet d'apporter un complément analogue à
(\hyperref[0.13.2.3-fr]{13.2.3})~: supposons que
$K_\alpha^n=0$ pour $n<0$ et pour tout $\alpha$~; supposons de plus que
$(K_\alpha^n)_{\alpha\in I}$ vérifie (ML) pour $n\geq0$ et que les
$A_\alpha=H^0(K_\alpha^\bullet)$ puissent être munis de structures
d'espaces métriques vérifiant les propriétés ci-dessus. Alors les
conclusions de (\hyperref[0.13.2.3-fr]{13.2.3}) sont inchangées pour
$n\geq2$, et en outre $h_1$ est \emph{bijectif}, car le raisonnement de
(\hyperref[0.13.2.2-fr]{13.2.2}) montre encore que
$(\mathrm B^1(K_\alpha^\bullet))_{\alpha\in I}$ vérifie (ML), que la suite
$(***_1)$ est exacte, et enfin, en vertu de ce qui précède, que
$\varprojlim\mathrm B^1(K_\alpha^\bullet)=\mathrm B^1(K^\bullet)$. On a
ainsi établi entre autres les assertions de (T, 3.10.2).

\item[\emph{(ii)}]
Il est possible d'introduire les foncteurs dérivés à droite
$\varinjlim^{(n)}$ du foncteur $\varinjlim$, et d'obtenir des énoncés
plus complets que les précédents [28].
\end{enumerate}
\end{env}

\subsection{Application : cohomologie d'une limite projective de faisceaux}
\label{subsection:0.13.3-fr}

\begin{proposition}[13.3.1]
\label{0.13.3.1-fr}
\oldpage[0\textsubscript{III}]{68}
\emph{Soient $X$ un espace topologique,
$(\mathcal F_k)_{k\in\mathbb N}$ un système projectif de faisceaux de
groupes abéliens sur $X$, et soit
$\mathcal F=\varprojlim_k\mathcal F_k$. On suppose vérifiées les conditions
suivantes :}
\begin{enumerate}
\item[\emph{(i)}]
\emph{Il existe une base $\mathfrak B$ de la topologie de $X$ telle que,
pour tout $U\in\mathfrak B$ et tout $i\geq0$, le système projectif
$(H^i(U,\mathcal F_k))_{k\in\mathbb N}$ vérifie (ML).}
\item[\emph{(ii)}]
\emph{Pour tout $x\in X$ et tout $i>0$, on a}
\[
\varinjlim_U\bigl(\varprojlim_k H^i(U,\mathcal F_k)\bigr)=0
\]
\emph{lorsque $U$ parcourt l'ensemble des voisinages de $x$ appartenant à
$\mathfrak B$.}
\item[\emph{(iii)}]
\emph{Les homomorphismes
$u_{hk}:\mathcal F_k\to\mathcal F_h$ $(h\leq k)$ définissant le système
projectif $(\mathcal F_k)$ sont surjectifs.}
\end{enumerate}
\emph{Dans ces conditions, pour tout $i>0$, l'homomorphisme canonique}
\[
h_i:H^i(X,\mathcal F)\longrightarrow
\varprojlim_k H^i(X,\mathcal F_k)
\]
\emph{est surjectif~; si en outre, pour une valeur de $i$, le système
projectif $(H^{i-1}(X,\mathcal F_k))_{k\in\mathbb N}$ vérifie (ML), $h_i$
est bijectif.}
\end{proposition}

\begin{proof}
\emph{a)} Nous allons d'abord supposer que les $\mathcal F_k$ sont flasques
ainsi que les noyaux $\mathcal N_{hk}$ des $u_{hk}$~; nous allons montrer
alors que la condition (iii) de l'énoncé entraîne que
$H^i(X,\mathcal F)=0$ pour $i>0$. Il suffira de prouver que pour tout ouvert
$U$ de $X$ et tout recouvrement $\mathfrak U$ de $U$ par des ouverts de
$U$, on a $\check H^i(\mathfrak U,\mathcal F)=0$ pour $i>0$. Il en résultera
en effet tout d'abord que pour la cohomologie de Čech, on a
$\check H^i(U,\mathcal F)=0$ pour tout $i>0$, puis (en vertu de
(G, II, 5.9.2) appliqué à l'ensemble de tous les ouverts de $X$) que
$H^i(X,\mathcal F)=0$ pour tout $i>0$. Comme les $\mathcal F_k$ sont
flasques, on a $\check H^i(\mathfrak U,\mathcal F_k)=0$ pour $i>0$
(G, II, 5.2.3)~; considérons pour chaque $k$ le complexe
$\check C^\bullet(\mathfrak U,\mathcal F_k)$ des cochaînes alternées du nerf
du recouvrement $\mathfrak U$ (G, II, 5.1), qui forment évidemment un
système projectif de complexes de groupes abéliens. Montrons que toutes les
applications
$\check C^\bullet(\mathfrak U,\mathcal F_k)\to
\check C^\bullet(\mathfrak U,\mathcal F_h)$ $(h\leq k)$ sont surjectives.
Il suffit évidemment, par définition, de montrer que pour tout ouvert $V$
de $X$, l'application
$\Gamma(V,\mathcal F_k)\to\Gamma(V,\mathcal F_h)$ est surjective~; mais la
suite
\[
0\longrightarrow\mathcal N_{hk}\longrightarrow\mathcal F_k
\longrightarrow\mathcal F_h\longrightarrow0
\]
étant exacte par hypothèse, donne la suite exacte de cohomologie
\[
\Gamma(V,\mathcal F_k)\longrightarrow\Gamma(V,\mathcal F_h)
\longrightarrow H^1(V,\mathcal N_{hk})=0
\]
puisque $\mathcal N_{hk}$ est flasque. Le système projectif
$(\check C^\bullet(\mathfrak U,\mathcal F_k))_{k\in\mathbb N}$ vérifie donc
(ML)~; il en est de même de
$(\check H^i(\mathfrak U,\mathcal F_k))_{k\in\mathbb N}$ pour tout
$i\geq0$, car cela est trivial pour $i>0$, et comme
$\check H^0(\mathfrak U,\mathcal F_k)=\Gamma(U,\mathcal F_k)$
(G, II, 5.2.2), la condition (ML) est aussi remplie pour $i=0$ d'après ce
qui précède. On peut donc appliquer
(\hyperref[0.13.2.3-fr]{13.2.3}), qui montre que
\[
\check H^i(\mathfrak U,\mathcal F)
=\varprojlim_k\check H^i(\mathfrak U,\mathcal F_k)=0
\quad\hbox{pour tout }i>0.
\]

\emph{b)} Passons au cas général, et considérons pour chaque $k\in\mathbb N$,
la résolution canonique
$\mathcal C^\bullet(X,\mathcal F_k)
=(\mathcal C^i(X,\mathcal F_k))_{i\geq0}$ de $\mathcal F_k$ par des
faisceaux flasques (G, II, 4.3). Pour chaque $i\geq0$, il est clair que
$(\mathcal C^i(X,\mathcal F_k))_{k\in\mathbb N}$ est un système projectif
de faisceaux flasques~; montrons qu'il satisfait aux conditions de
\emph{a)}.
\oldpage[0\textsubscript{III}]{69}
En effet, si $\mathcal N_{hk}$ est le noyau de $u_{hk}$ pour $h\leq k$, la
suite
\[
0\longrightarrow\mathcal N_{hk}\longrightarrow\mathcal F_k
\longrightarrow\mathcal F_h\longrightarrow0
\]
est exacte d'après (iii), et notre assertion résulte de ce que le foncteur
$\mathcal A\mapsto\mathcal C^i(X,\mathcal A)$ est exact (G, II, 4.3).
Soit
$\mathcal G^i=\varprojlim_k\mathcal C^i(X,\mathcal F_k)$~; on a donc
$H^j(X,\mathcal G^i)=0$ pour $j>0$ et $i\geq0$ en vertu de \emph{a)}. Nous
allons montrer que $\mathcal G^\bullet=(\mathcal G^i)_{i\geq0}$ est une
résolution du faisceau $\mathcal F$~; comme
$H^j(X,\mathcal G^i)=0$ pour $j>0$, la cohomologie $H^\bullet(X,\mathcal F)$
sera égale à $H^\bullet(\Gamma(X,\mathcal G^\bullet))$ (G, II, 4.7.1).

Il est clair que, par passage à la limite projective, on déduit des suites
exactes
\[
0\longrightarrow\mathcal F_k
\longrightarrow\mathcal C^0(X,\mathcal F_k)
\longrightarrow\mathcal C^1(X,\mathcal F_k)\longrightarrow\cdots
\]
un complexe de faisceaux de groupes abéliens
\[
0\longrightarrow\mathcal F\longrightarrow\mathcal G^0
\longrightarrow\mathcal G^1\longrightarrow\cdots.
\]
Pour prouver notre assertion, il faut établir que
$\mathcal H^i(\mathcal G^\bullet)=0$ pour $i>0$. Ce faisceau est engendré
par le préfaisceau
$U\mapsto H^i(\Gamma(U,\mathcal G^\bullet))$ (G, II, 4.1)~; or, le complexe
$\Gamma(U,\mathcal G^\bullet)$ est la limite projective du système
projectif de complexes de groupes abéliens
\[
(\Gamma(U,\mathcal C^\bullet(X,\mathcal F_k)))_{k\in\mathbb N}
\]
(0\textsubscript{I}, 3.2.6). On a vu dans \emph{a)} que pour chaque
$i\geq0$, les applications
$\Gamma(U,\mathcal C^i(X,\mathcal F_k))\to
\Gamma(U,\mathcal C^i(X,\mathcal F_h))$ $(h\leq k)$ sont surjectives~;
d'autre part, on a
$H^i(U,\mathcal F_k)=
H^i(\Gamma(U,\mathcal C^\bullet(X,\mathcal F_k)))$, la résolution canonique
$\mathcal C^\bullet(U,\mathcal F_k|U)$ étant induite sur $U$ par
$\mathcal C^\bullet(X,\mathcal F_k)$~; en vertu de l'hypothèse (i), pour
tout $U\in\mathfrak B$, on peut appliquer
(\hyperref[0.13.2.3-fr]{13.2.3}) au système projectif de complexes
$(\Gamma(U,\mathcal C^\bullet(X,\mathcal F_k)))_{k\in\mathbb N}$, et on a
donc
\[
H^i(\Gamma(U,\mathcal G^\bullet))
=\varprojlim_k H^i(U,\mathcal F_k)
\quad\hbox{pour tout }i\geq0.
\]
L'hypothèse (ii) prouve bien alors, par définition, que les faisceaux
$\mathcal H^i(\mathcal G^\bullet)$ sont nuls pour $i>0$.

On a alors pour tout $i\geq0$,
\[
H^i(X,\mathcal F)=H^i(\Gamma(X,\mathcal G^\bullet))
\quad\hbox{et}\quad
\Gamma(X,\mathcal G^\bullet)
=\varprojlim_k\Gamma(X,\mathcal C^\bullet(X,\mathcal F_k)).
\]
On vient de remarquer que les applications
$\Gamma(X,\mathcal C^i(X,\mathcal F_k))\to
\Gamma(X,\mathcal C^i(X,\mathcal F_h))$ $(h\leq k)$ sont toutes
surjectives~; la conclusion résulte donc encore de
(\hyperref[0.13.2.3-fr]{13.2.3}).
\end{proof}

\begin{env}[13.3.2]
\label{0.13.3.2-fr}
\emph{Remarques} (13.3.2).
\begin{enumerate}
\item[\emph{(i)}]
L'énoncé (\hyperref[0.13.3.1-fr]{13.3.1}) n'a d'intérêt que pour $i>0$,
puisque pour $i=0$, $h_i$ est toujours un isomorphisme sans hypothèse
(0\textsubscript{I}, 3.2.6).
\item[\emph{(ii)}]
Les conditions (i) et (ii) de (\hyperref[0.13.3.1-fr]{13.3.1}) seront en
particulier vérifiées si $H^i(U,\mathcal F_k)=0$ pour tout $k$, tout $i>0$
et tout $U\in\mathfrak B$, et si pour $U\in\mathfrak B$, les applications
$\Gamma(U,\mathcal F_k)\to\Gamma(U,\mathcal F_h)$ sont surjectives. Ce sera
le cas le plus fréquent d'application de
(\hyperref[0.13.3.1-fr]{13.3.1}).
\end{enumerate}
\end{env}

\subsection{Condition de Mittag-Leffler et objets gradués associés aux systèmes projectifs}
\label{subsection:0.13.4-fr}
\label{0.13.4-fr}

\begin{env}[13.4.1]
\label{0.13.4.1-fr}
Soit $\mathbf A=(A_k,u_{kh})_{k\in\mathbb Z}$ un système projectif dans une
catégorie abélienne $C$~; nous dirons qu'il est \emph{limité inférieurement}
s'il existe $k_0$ tel que $A_k=0$ pour $k<k_0$. Nous définirons sur chaque
$A_k$ une \emph{filtration} $(\mathrm F^p(A_k))_{p\in\mathbb Z}$ par les
formules
\begin{equation}
\label{0.13.4.1.1-fr}
\tag{13.4.1.1}
\mathrm F^p(A_k)=
\begin{cases}
\operatorname{Ker}(A_k\longrightarrow A_{p-1})&\text{pour }p\leq k+1,\\
0&\text{pour }p\geq k+1.
\end{cases}
\end{equation}
\oldpage[0\textsubscript{III}]{70}
On a donc par hypothèse
$\mathrm F^{k_0}(A_k)=A_k$ et $\mathrm F^{k+1}(A_k)=0$, autrement dit la
filtration considérée est \emph{finie}
(\hyperref[0.11.1.3-fr]{11.1.3}). Les objets gradués associés à cette
filtration sont donc
\[
\operatorname{gr}^p(A_k)=
\operatorname{Ker}(A_k\longrightarrow A_{p-1})
/\operatorname{Ker}(A_k\longrightarrow A_p)
\]
et par suite, $\operatorname{gr}^p(A_k)$ est isomorphe à l'image par
$A_k\to A_p$ de $\operatorname{Ker}(A_k\to A_{p-1})$~; en vertu de la
transitivité des morphismes définissant un système projectif, on a donc
\begin{equation}
\label{0.13.4.1.2-fr}
\tag{13.4.1.2}
\operatorname{gr}^p(A_k)=
\operatorname{Ker}(A_p\longrightarrow A_{p-1})
\cap\operatorname{Im}(A_k\longrightarrow A_p)
\end{equation}
mais comme, en vertu de
(\hyperref[0.13.4.1.1-fr]{13.4.1.1}), on a
$\operatorname{Ker}(A_p\to A_{p-1})=\operatorname{gr}^p(A_p)$, on a aussi
\begin{equation}
\label{0.13.4.1.3-fr}
\tag{13.4.1.3}
\operatorname{gr}^p(A_k)=
\operatorname{gr}^p(A_p)\cap\operatorname{Im}(A_k\longrightarrow A_p).
\end{equation}
Les définitions précédentes montrent en outre que l'on a pour $k\leq h$
\[
u_{kh}(\mathrm F^p(A_h))\subset\mathrm F^p(A_k)
\]
et par suite que les $\operatorname{gr}^p(u_{kh})$ définissent un
\emph{système projectif}
$(\operatorname{gr}^p(A_k))_{k\in\mathbb Z}$ pour tout $p\in\mathbb Z$.
\end{env}

\begin{env}[13.4.2]
\label{0.13.4.2-fr}
Nous dirons que le système projectif $\mathbf A$ est
\emph{essentiellement constant} si les morphismes
$A_{k+1}\to A_k$ sont des \emph{isomorphismes} pour $k$ assez grand. Nous
dirons que le système projectif $\mathbf A$ est \emph{strict} si les
morphismes $A_i\to A_j$ $(j\leq i)$ sont des \emph{épimorphismes}.
Lorsque $\mathbf A$ est strict, il résulte de
(\hyperref[0.13.4.1.3-fr]{13.4.1.3}) que pour $p\leq k\leq h$, le morphisme
canonique $\operatorname{gr}^p(A_h)\to\operatorname{gr}^p(A_k)$ est un
\emph{isomorphisme}, autrement dit, le système projectif
$(\operatorname{gr}^p(A_k))_{k\in\mathbb Z}$ est essentiellement constant.
La suite des objets $\operatorname{gr}^p(A_p)$ (identifiés à
$\varprojlim_k\operatorname{gr}^p(A_k)$ pour tout $p$) se note alors
$\operatorname{gr}^\bullet(\mathbf A)$ et s'appelle \emph{l'objet gradué
associé au système projectif strict} $\mathbf A=(A_k)$.

Si l'on suppose maintenant que le système projectif $\mathbf A$ (limité
inférieurement) vérifie (ML), on sait
(\hyperref[0.13.1.2-fr]{13.1.2}) que le système projectif
$\mathbf A'=(A'_k)$ des objets d'«~images universelles~» est \emph{strict},
et il est par ailleurs limité inférieurement~; l'objet gradué
$\operatorname{gr}^\bullet(\mathbf A')$ associé à $\mathbf A'$ est alors
encore appelé l'objet gradué \emph{associé à} $\mathbf A$ et noté
$\operatorname{gr}^\bullet(\mathbf A)$.
\end{env}

\begin{proposition}[13.4.3]
\label{0.13.4.3-fr}
\emph{Soit $\mathbf A=(A_k)_{k\in\mathbb Z}$ un système projectif limité
inférieurement dans une catégorie abélienne $C$. Les deux conditions
suivantes sont équivalentes :}
\begin{enumerate}
\item[\emph{a)}] \emph{$\mathbf A$ vérifie la condition (ML).}
\item[\emph{b)}] \emph{Pour tout $p\in\mathbb Z$, le système projectif
$(\operatorname{gr}^p(A_k))_{k\in\mathbb Z}$ est essentiellement constant.}
\end{enumerate}
\emph{En outre, lorsque ces conditions sont satisfaites, on a pour tout
$p\in\mathbb Z$ un isomorphisme canonique}
\begin{equation}
\label{0.13.4.3.1-fr}
\tag{13.4.3.1}
\operatorname{gr}^p(\mathbf A)\xrightarrow{\sim}
\varprojlim_k\operatorname{gr}^p(A_k).
\end{equation}
\end{proposition}

\begin{proof}
Il résulte aussitôt de
(\hyperref[0.13.4.1.2-fr]{13.4.1.2}) que \emph{a)} implique \emph{b)}~; la
même formule appliquée au système projectif $\mathbf A'$ (notations de
(\hyperref[0.13.4.2-fr]{13.4.2})) donne l'isomorphisme
(\hyperref[0.13.4.3.1-fr]{13.4.3.1}) par définition. Pour $k\leq h$,
posons $A_{kh}=\operatorname{Im}(A_h\to A_k)$~; si $k\leq h\leq j$, on a
$A_{kj}\subset A_{kh}\subset A_k$. Munissons $A_{kh}$ de la filtration
induite par $(\mathrm F^p(A_k))$~; on vérifie aussitôt, en vertu de la
transitivité des morphismes définissant $\mathbf A$, que cette filtration
est aussi la filtration quotient de $(\mathrm F^p(A_h))$~; par suite, on a
\begin{equation}
\label{0.13.4.3.2-fr}
\tag{13.4.3.2}
\operatorname{gr}^p(A_{kh})=
\operatorname{Im}\bigl(\operatorname{gr}^p(A_h)\longrightarrow
\operatorname{gr}^p(A_k)\bigr).
\end{equation}
\oldpage[0\textsubscript{III}]{71}
Cela étant, supposons \emph{b)} vérifiée~; pour tout $p\in\mathbb Z$, et
tout $k\geq p$, il existe un entier $\mathrm L(p,k)$ tel que le second
membre de (\hyperref[0.13.4.3.2-fr]{13.4.3.2}) soit constant pour
$h\geq\mathrm L(p,k)$~; comme $\operatorname{gr}^p(A_k)=0$ pour $p<k_0$,
il n'y a (pour $k$ donné) qu'un nombre fini de $\mathrm L(p,k)$ non nuls
lorsque $p$ parcourt l'ensemble des entiers $\leq k$. Soit
$\mathrm L(k)=m$ le plus grand de ces entiers~; pour tout $h\geq m$, on a
$A_{kh}\subset A_{km}$, et par définition de $m$, l'injection canonique
$A_{kh}\to A_{km}$ définit un \emph{isomorphisme}
$\operatorname{gr}^\bullet(A_{kh})\xrightarrow{\sim}
\operatorname{gr}^\bullet(A_{km})$~; comme les filtrations sont
\emph{finies}, on en conclut que l'injection précédente est elle-même
bijective (Bourbaki, \emph{Alg. comm.}, chap.~III, §~2,
n\textsuperscript{o}~8, th.~1), ce qui prouve que $\mathbf A$ vérifie
(ML).
\end{proof}

\begin{env}[13.4.4]
\label{0.13.4.4-fr}
Supposons que dans $C$ la limite projective
$A=\varprojlim_k A_k$ existe. Dans les définitions de
(\hyperref[0.13.4.1-fr]{13.4.1}), on peut alors remplacer $A_k$ par $A$,
et la filtration ainsi définie sur $A$ est encore telle que
\begin{equation}
\label{0.13.4.4.1-fr}
\tag{13.4.4.1}
\operatorname{gr}^p(A)=\operatorname{gr}^p(A_p)
\cap\operatorname{Im}(A\longrightarrow A_p).
\end{equation}
\end{env}

\begin{corollary}[13.4.5]
\label{0.13.4.5-fr}
\emph{Supposons que $C$ soit la catégorie des groupes abéliens. Si le
système projectif $\mathbf A$ vérifie (ML) et si
$A=\varprojlim_k A_k$, on a pour tout $p\in\mathbb Z$ un isomorphisme
canonique}
\begin{equation}
\label{0.13.4.5.1-fr}
\tag{13.4.5.1}
\operatorname{gr}^p(A)\xrightarrow{\sim}
\varprojlim_k\operatorname{gr}^p(A_k).
\end{equation}
\end{corollary}
\begin{proof}
En effet, on a
$\operatorname{Im}(A_k\to A_p)=\operatorname{Im}(A\to A_p)$ dès que $k$
est assez grand (Bourbaki, \emph{Top. gén.}, chap.~II,
3\textsuperscript{e} éd., §~3, n\textsuperscript{o}~5, th.~1), et la
conclusion résulte de
(\hyperref[0.13.4.1.3-fr]{13.4.1.3}) et
(\hyperref[0.13.4.4.1-fr]{13.4.4.1}).
\end{proof}

\subsection{Limites projectives de suites spectrales de complexes filtrés}
\label{subsection:0.13.5-fr}
\begin{env}[13.5.1]
\label{0.13.5.1-fr}
Soient $C$ une catégorie abélienne, $X^\bullet$ un complexe d'objets de
$C$ muni d'une \emph{filtration}
$(\mathrm F^p(X^\bullet))_{p\in\mathbb Z}$ telle que
$\mathrm F^{p_0}(X^\bullet)=X^\bullet$ pour un indice $p_0$.
Considérons pour chaque $k\in\mathbb Z$ le complexe
$X_k^\bullet=X^\bullet/\mathrm F^{k+1}(X^\bullet)$~; il est canoniquement
muni de la filtration formée des
$\mathrm F^p(X_k^\bullet)=
\mathrm F^p(X^\bullet)/\mathrm F^{k+1}(X^\bullet)$ pour $p\leq k$ et des
$\mathrm F^p(X_k^\bullet)=0$ pour $p\geq k+1$. En outre, on a des
morphismes canoniques $X_{k+1}^\bullet\to X_k^\bullet$, qui font de
$\mathbf X^\bullet=(X_k^\bullet)_{k\in\mathbb Z}$ un \emph{système
projectif de complexes filtrés d'objets de $C$}. On notera que ce système
projectif est \emph{strict} et tel que $X_k^\bullet=0$ pour $k<p_0$.
\end{env}
\begin{env}[13.5.2]
\label{0.13.5.2-fr}
Considérons plus généralement un système projectif \emph{strict}
$\mathbf X^\bullet=(X_k^\bullet)_{k\in\mathbb Z}$ de complexes d'objets de
$C$, \emph{limité inférieurement}~; considérons sur chaque
$X_k^\bullet$ la filtration définie dans
(\hyperref[0.13.4.1-fr]{13.4.1}) (en se plaçant dans la catégorie abélienne
des complexes de $C$ limités inférieurement). Les
$X_k^\bullet\to X_p^\bullet$ $(p\leq k)$ deviennent des morphismes de
complexes \emph{filtrés}, à filtrations finies. Le caractère fonctoriel des
suites spectrales des complexes filtrés
(\hyperref[0.11.2.3-fr]{11.2.3}) montre que les morphismes de définition du
système projectif $\mathbf X^\bullet$ fournissent des morphismes faisant de
$\mathrm E(\mathbf X^\bullet)=(\mathrm E(X_k^\bullet))$ un \emph{système
projectif de suites spectrales}.
\end{env}
\begin{lemma}[13.5.3]
\label{0.13.5.3-fr}
\emph{Supposons que le système projectif
$\mathbf X^\bullet=(X_k^\bullet)_{k\in\mathbb Z}$ de complexes filtrés soit
obtenu comme dans (\hyperref[0.13.5.2-fr]{13.5.2}). Alors :}
\begin{enumerate}
\item[\emph{a)}]\emph{Pour $r\geq p-p_0$, on a
$\mathrm B_r^{pq}(X_k^\bullet)=\mathrm B_\infty^{pq}(X_k^\bullet)$ pour
tout $k\in\mathbb Z$.}
\item[\emph{b)}]\emph{Pour $k+1\leq p+r$, on a
$\mathrm Z_r^{pq}(X_k^\bullet)=\mathrm Z_\infty^{pq}(X_k^\bullet)$.}
\item[\emph{c)}]\emph{Pour $k+1\geq p+r$, les morphismes
$\mathrm Z_r^{pq}(X_h^\bullet)\to\mathrm Z_r^{pq}(X_k^\bullet)$ et
$\mathrm B_r^{pq}(X_h^\bullet)\to\mathrm B_r^{pq}(X_k^\bullet)$ sont des
isomorphismes pour tout $h\geq k$.}
\end{enumerate}
\end{lemma}
\oldpage[0\textsubscript{III}]{72}
Ces trois propriétés résultent aussitôt des définitions de
(\hyperref[0.11.2.2-fr]{11.2.2}), en tenant compte de ce que
$\mathrm F^{p-r+1}(X_k^\bullet)=X_k^\bullet$ pour $p-r<p_0$.

\begin{env}[13.5.4]
\label{0.13.5.4-fr}
Supposons vérifiées les hypothèses de
(\hyperref[0.13.5.3-fr]{13.5.3}). Alors, pour $p,q,r$ fixés ($r$ fini), les
systèmes projectifs
\[
(\mathrm Z_r^{pq}(X_k^\bullet))_{k\in\mathbb Z},\quad
(\mathrm B_r^{pq}(X_k^\bullet))_{k\in\mathbb Z},\quad
(\mathrm E_r^{pq}(X_k^\bullet))_{k\in\mathbb Z}
\]
sont \emph{essentiellement constants}~; on désignera par
$\mathrm Z_r^{pq}(\mathbf X^\bullet)$,
$\mathrm B_r^{pq}(\mathbf X^\bullet)$ et
$\mathrm E_r^{pq}(\mathbf X^\bullet)=
\mathrm Z_r^{pq}(\mathbf X^\bullet)/
\mathrm B_r^{pq}(\mathbf X^\bullet)$ leurs limites projectives
respectives. Les $\mathrm Z_r^{pq}(\mathbf X^\bullet)$ et
$\mathrm B_r^{pq}(\mathbf X^\bullet)$ s'identifient canoniquement à des
sous-objets de $\mathrm E_1^{pq}(\mathbf X^\bullet)$. La définition des
$d_r^{pq}$ (M, XV, 1) montre que ces morphismes (relatifs aux
$X_k^\bullet$) sont aussi essentiellement constants, et par suite
définissent des morphismes
\begin{equation}
\label{0.13.5.4.1-fr}
\tag{13.5.4.1}
d_r^{pq}:\mathrm E_r^{pq}(\mathbf X^\bullet)
\longrightarrow\mathrm E_r^{p+r,q-r+1}(\mathbf X^\bullet)
\end{equation}
tels que $d_r^{p+r,q-r+1}\circ d_r^{pq}=0$~; en outre, on a des
isomorphismes canoniques de $\operatorname{Ker}(d_r^{pq})$ sur
$\mathrm Z_{r+1}^{pq}(\mathbf X^\bullet)/
\mathrm B_r^{pq}(\mathbf X^\bullet)$ et de
$\operatorname{Im}(d_r^{pq})$ sur
$\mathrm B_{r+1}^{p+r,q-r+1}(\mathbf X^\bullet)/
\mathrm B_r^{p+r,q-r+1}(\mathbf X^\bullet)$.
\end{env}

\begin{lemma}[13.5.5]
\label{0.13.5.5-fr}
\emph{Sous les hypothèses de
(\hyperref[0.13.5.3-fr]{13.5.3}), on a, pour $s\geq r>p-p_0$, un
monomorphisme canonique}
\begin{equation}
\label{0.13.5.5.1-fr}
\tag{13.5.5.1}
i:\mathrm E_s^{pq}(\mathbf X^\bullet)
\longrightarrow\mathrm E_r^{pq}(\mathbf X^\bullet)
\end{equation}
\emph{et un isomorphisme canonique}
\begin{equation}
\label{0.13.5.5.2-fr}
\tag{13.5.5.2}
j_r:\mathrm E_r^{pq}(\mathbf X^\bullet)
\xrightarrow{\sim}\mathrm E_\infty^{pq}(X_{p+r-1}^\bullet)
\end{equation}
\emph{tels que le diagramme}
\begin{equation}
\label{0.13.5.5.3-fr}
\tag{13.5.5.3}
\xymatrix{
\mathrm E_s^{pq}(\mathbf X^\bullet)\ar[r]^{j_s}\ar[d]_i
&\mathrm E_\infty^{pq}(X_{p+s-1}^\bullet)\ar[d]\\
\mathrm E_r^{pq}(\mathbf X^\bullet)\ar[r]^{j_r}
&\mathrm E_\infty^{pq}(X_{p+r-1}^\bullet)
}
\end{equation}
\emph{soit commutatif} (la flèche verticale de droite provenant du
morphisme $X_{p+s-1}^\bullet\to X_{p+r-1}^\bullet$).
\end{lemma}
\begin{proof}
L'existence de $i$ provient de ce que
$\mathrm B_r^{pq}(X_k^\bullet)=\mathrm B_\infty^{pq}(X_k^\bullet)$ pour
$r>p-p_0$ (\hyperref[0.13.5.3-fr]{13.5.3}, a)~; on a
$\mathrm Z_r^{pq}(X_k^\bullet)=\mathrm Z_\infty^{pq}(X_k^\bullet)$ pour
$k+1\leq p+r$ (\hyperref[0.13.5.3-fr]{13.5.3}, b), d'où en particulier
$\mathrm Z_r^{pq}(X_{p+r-1}^\bullet)=
\mathrm Z_\infty^{pq}(X_{p+r-1}^\bullet)$ et d'autre part
$\mathrm Z_r^{pq}(X_{p+r-1}^\bullet)$ et
$\mathrm B_r^{pq}(X_{p+r-1}^\bullet)$ s'identifient canoniquement à
$\mathrm Z_r^{pq}(\mathbf X^\bullet)$ et
$\mathrm B_r^{pq}(\mathbf X^\bullet)$ en vertu de
(\hyperref[0.13.5.3-fr]{13.5.3}, c)), d'où l'existence de $j_r$ et la
commutativité de
(\hyperref[0.13.5.5.1-fr]{13.5.5.1}).
\end{proof}

\begin{corollary}[13.5.6]
\label{0.13.5.6-fr}
\emph{Sous les hypothèses de
(\hyperref[0.13.5.3-fr]{13.5.3}), si l'une des limites projectives
$\varprojlim_r\mathrm E_r^{pq}(\mathbf X^\bullet)$,
$\varprojlim_k\mathrm E_\infty^{pq}(X_k^\bullet)$ existe, il en est de
même de l'autre, et on a un isomorphisme canonique}
\begin{equation}
\label{0.13.5.6.1-fr}
\tag{13.5.6.1}
j_\infty:\varprojlim_r\mathrm E_r^{pq}(\mathbf X^\bullet)
\xrightarrow{\sim}
\varprojlim_k\mathrm E_\infty^{pq}(X_k^\bullet).
\end{equation}
\emph{En outre, pour que le système projectif
$(\mathrm E_r^{pq}(\mathbf X^\bullet))_{r\in\mathbb Z}$ soit
essentiellement constant
(\hyperref[0.13.4.2-fr]{13.4.2}), il faut et il suffit que le système
projectif $(\mathrm E_\infty^{pq}(X_k^\bullet))_{k\in\mathbb Z}$ le soit.}
\end{corollary}

\begin{env}[13.5.7]
\label{0.13.5.7-fr}
On note $\mathrm B_\infty^{pq}(\mathbf X^\bullet)$ et
$\mathrm Z_\infty^{pq}(\mathbf X^\bullet)$ les sous-objets de
$\mathrm E_1^{pq}(\mathbf X^\bullet)$ égaux respectivement à
$\mathrm B_r^{pq}(\mathbf X^\bullet)$ pour $r>p-p_0$ et à
$\inf_r\mathrm Z_r^{pq}(\mathbf X^\bullet)$ (quand ce dernier existe), de
sorte que $\varprojlim_r\mathrm E_r^{pq}(\mathbf X^\bullet)$
\oldpage[0\textsubscript{III}]{73}
s'identifie canoniquement à
$\mathrm E_\infty^{pq}(\mathbf X^\bullet)=
\mathrm Z_\infty^{pq}(\mathbf X^\bullet)/
\mathrm B_\infty^{pq}(\mathbf X^\bullet)$. On remarquera que les objets
$\mathrm Z_r^{pq}(\mathbf X^\bullet)$,
$\mathrm B_r^{pq}(\mathbf X^\bullet)$,
$\mathrm E_r^{pq}(\mathbf X^\bullet)$ $(1\leq r\leq+\infty)$ et
$d_r^{pq}$ dépendent \emph{fonctoriellement} du système projectif
$\mathbf X^\bullet$ soumis aux restrictions de
(\hyperref[0.13.5.5-fr]{13.5.5}), et que les morphismes définis dans
(\hyperref[0.13.5.5-fr]{13.5.5}) et
(\hyperref[0.13.5.6-fr]{13.5.6}) sont fonctoriels.
\end{env}

\subsection{Suite spectrale d'un foncteur relative à un objet muni d'une filtration finie}
\label{subsection:0.13.6-fr}

\begin{env}[13.6.1]
\label{0.13.6.1-fr}
Soient $C$, $C'$ deux catégories abéliennes, $\mathrm T:C\to C'$ un
foncteur additif covariant. Supposons que tout objet de $C$ soit isomorphe
à un sous-objet d'un objet injectif, de sorte que les foncteurs dérivés
droits $\mathrm R^p\mathrm T$ $(p\geq0)$ existent.
\end{env}

\begin{lemma}[13.6.2]
\label{0.13.6.2-fr}
\emph{Soit $A$ un objet de $C$, muni d'une filtration finie
$(\mathrm F^i(A))_{i\in\mathbb Z}$. Il existe une résolution injective
$X^\bullet=(X^j)_{j\geq0}$ de $A$ munie d'une filtration finie
$(\mathrm F^i(X^\bullet))_{i\in\mathbb Z}$ telle que la relation
$\mathrm F^i(A)=A$ (resp. $\mathrm F^i(A)=0$) entraîne
$\mathrm F^i(X^\bullet)=X^\bullet$ (resp. $\mathrm F^i(X^\bullet)=0$) et
que, pour tout $i\in\mathbb Z$, $\mathrm F^i(X^\bullet)$ soit une
résolution injective de $\mathrm F^i(A)$.}
\end{lemma}
\begin{proof}
Soit $p$ (resp. $q>p$) le plus grand indice tel que
$\mathrm F^p(A)=A$ (resp. le plus petit indice pour lequel
$\mathrm F^q(A)=0$). On raisonne par récurrence sur $q-p$, le lemme étant
évident pour $q-p=1$. Ayant formé une résolution injective
${X''}^\bullet$ de $A/\mathrm F^{q-1}(A)$ ayant les propriétés voulues,
on considère la suite exacte
\[
0\to\mathrm F^{q-1}(A)\to A\to A/\mathrm F^{q-1}(A)\to0,
\]
on prend une résolution injective ${X'}^\bullet$ de
$\mathrm F^{q-1}(A)$, puis on détermine une résolution injective
$X^\bullet$ de $A$ de façon à avoir une suite exacte
\[
0\to {X'}^\bullet\to X^\bullet\to {X''}^\bullet\to0
\]
compatible avec la précédente (M, V, 2.2)~; il est clair que
$X^\bullet$ répond à la question.
\end{proof}

\begin{corollary}[13.6.3]
\label{0.13.6.3-fr}
\emph{Soit $B$ un second objet de $C$, muni d'une filtration finie
$(\mathrm F^i(B))_{i\in\mathbb Z}$, $s$ un entier, et soit $u:A\to B$
un morphisme tel que $u(\mathrm F^i(A))\subset\mathrm F^{i+s}(B)$ pour
tout $i\in\mathbb Z$. Si $Y^\bullet=(Y^j)_{j\geq0}$ est une résolution
injective de $B$ munie d'une filtration
$(\mathrm F^i(Y^\bullet))_{i\in\mathbb Z}$ ayant les propriétés énoncées
en (\hyperref[0.13.6.2-fr]{13.6.2}), il existe un morphisme
$v:X^\bullet\to Y^\bullet$ compatible avec $u$ et tel que
$v(\mathrm F^i(X^\bullet))\subset\mathrm F^{i+s}(Y^\bullet)$ pour tout
$i\in\mathbb Z$. En outre, deux tels morphismes $v$, $v'$ sont
homotopes.}
\end{corollary}

Cela résulte aussitôt par récurrence sur $q-p$ de la construction
précédente et de (M, V, 2.3).

\begin{env}[13.6.4]
\label{0.13.6.4-fr}
Sous les hypothèses de (\hyperref[0.13.6.2-fr]{13.6.2}), considérons
maintenant le complexe $\mathrm T(X^\bullet)$ dans $C'$, qui est
évidemment filtré par les complexes
$\mathrm T(\mathrm F^i(X^\bullet))$, puisque
$\mathrm F^i(X^\bullet)$ est facteur direct de $X^\bullet$. Il résulte de
(\hyperref[0.13.6.3-fr]{13.6.3}) que la \emph{suite spectrale} de ce
complexe filtré ne dépend que de l'objet filtré $A$, à un isomorphisme
près. Son aboutissement est la cohomologie $\mathrm R^\bullet\mathrm T(A)$,
avec la filtration
\begin{equation}
\label{0.13.6.4.1-fr}
\tag{13.6.4.1}
\begin{split}
\mathrm F^p(\mathrm R^n\mathrm T(A))
&=\operatorname{Im}\bigl(\mathrm R^n\mathrm T(\mathrm F^p(A))
\longrightarrow\mathrm R^n\mathrm T(A)\bigr)\\
&=\operatorname{Ker}\bigl(\mathrm R^n\mathrm T(A)
\longrightarrow\mathrm R^n\mathrm T(A/\mathrm F^p(A))\bigr)
\end{split}
\end{equation}
(\hyperref[0.11.2.2-fr]{11.2.2}), et son terme $\mathrm E_1$ est donné
par
\begin{equation}
\label{0.13.6.4.2-fr}
\tag{13.6.4.2}
\mathrm E_1^{pq}=\mathrm R^{p+q}\mathrm T(\operatorname{gr}^p(A))
\end{equation}
$\operatorname{gr}^p(A)$ désignant comme d'ordinaire
$\mathrm F^p(A)/\mathrm F^{p+1}(A)$. Il est clair, d'après
(\hyperref[0.11.2.2-fr]{11.2.2}), que la filtration de l'aboutissement
est finie, et que pour $p$, $q$ donnés, les suites des
\oldpage[0\textsubscript{III}]{74}
$\mathrm B_r^{pq}(A)=\mathrm B_r^{pq}(\mathrm T(X^\bullet))$ et
$\mathrm Z_r^{pq}(A)=\mathrm Z_r^{pq}(\mathrm T(X^\bullet))$ sont
stationnaires, donc la suite spectrale précédente est \emph{birégulière}
(\hyperref[0.11.1.3-fr]{11.1.3}). Nous noterons cette suite
$\mathrm E(A)=(\mathrm E_r^{pq}(A))$ et nous dirons que c'est la
\emph{suite spectrale du foncteur $\mathrm T$ relative à l'objet filtré
$A$}.
\end{env}

\begin{env}[13.6.5]
\label{0.13.6.5-fr}
Supposons maintenant vérifiées les hypothèses de
(\hyperref[0.13.6.3-fr]{13.6.3}), dont nous conservons les notations.
Comme $\mathrm F^i(X^\bullet)$ (resp. $\mathrm F^i(Y^\bullet)$) est
facteur direct de $X^\bullet$ (resp. $Y^\bullet$), on a
\[
(\mathrm T v)(\mathrm T(\mathrm F^i(X^\bullet)))
\subset\mathrm T(\mathrm F^{i+s}(Y^\bullet))
\]
pour tout $i\in\mathbb Z$~; les définitions de
(\hyperref[0.11.2.2-fr]{11.2.2}) montrent alors que pour
$1\leq r\leq+\infty$, $\mathrm T v$ définit un morphisme
$\mathrm B_r^{pq}(\mathrm T(X^\bullet))\to
\mathrm B_r^{p+s,q-s}(\mathrm T(Y^\bullet))$ et un morphisme
$\mathrm Z_r^{pq}(\mathrm T(X^\bullet))\to
\mathrm Z_r^{p+s,q-s}(\mathrm T(Y^\bullet))$, d'où un morphisme
\[
w_r:\mathrm E_r^{pq}(A)\to\mathrm E_r^{p+s,q-s}(B)~;
\]
de même, on a pour l'aboutissement des morphismes
$u_n:\mathrm R^n\mathrm T(A)\to\mathrm R^n\mathrm T(B)$ tels que
$u_n(\mathrm F^p(\mathrm R^n\mathrm T(A)))\subset
\mathrm F^{p+s}(\mathrm R^n\mathrm T(B))$.

La définition des $d_r^{pq}$ (M, XV, 1) montre en outre que les
diagrammes
\[
\xymatrix{
\mathrm E_r^{pq}(A)\ar[r]^{d_r^{pq}}\ar[d]_{w_r}
&\mathrm E_r^{p+r,q-r+1}(A)\ar[d]^{w_r}\\
\mathrm E_r^{p+s,q-s}(B)\ar[r]_{d_r^{p+s,q-s}}
&\mathrm E_r^{p+r+s,q-r-s+1}(B)
}
\]
sont commutatifs~; on en déduit un diagramme commutatif analogue pour
les isomorphismes $\alpha_r^{pq}$, que nous laisserons au lecteur le
soin d'expliciter. Enfin (\emph{loc. cit.}), on a aussi des diagrammes
commutatifs pour les aboutissements
\[
\xymatrix{
\mathrm E_\infty^{pq}(A)\ar[r]^{\beta^{pq}}\ar[d]_{w_\infty}
&\operatorname{gr}^p(\mathrm R^{p+q}\mathrm T(A))
  \ar[d]^{u_{p+q}}\\
\mathrm E_\infty^{p+s,q-s}(B)\ar[r]_{\beta^{p+s,q-s}}
&\operatorname{gr}^{p+s}(\mathrm R^{p+q}\mathrm T(B)).
}
\]
\end{env}

\begin{env}[13.6.6]
\label{0.13.6.6-fr}
Supposons en particulier qu'il existe un anneau $S$, muni d'une
filtration $(\mathrm F^i(S))_{i\in\mathbb Z}$, et un homomorphisme
d'anneaux
\begin{equation}
\label{0.13.6.6.1-fr}
\tag{13.6.6.1}
h:S\to\operatorname{Hom}_C(A,A)
\end{equation}
\oldpage[0\textsubscript{III}]{75}
tel que pour tout $t\in\mathrm F^i(S)$, on ait
$h_t(\mathrm F^j(A))\subset\mathrm F^{i+j}(A)$ pour tout couple $i$,
$j$. Nous dirons pour abréger que $A$ est alors muni d'une structure de
\emph{$S$-$C$-module filtré} sur l'anneau filtré $S$. Par passage aux
objets gradués associés, tout $h_t$, pour $t\in\mathrm F^j(S)$, définit
un endomorphisme gradué $\overline h_t$ de
$\operatorname{gr}^\bullet(A)$, homogène de degré $j$~; en outre, ce
morphisme ne dépend que de la classe de $t$ dans
$\operatorname{gr}^j(S)$, et on définit ainsi un homomorphisme d'anneaux
gradués
\[
\overline h:\operatorname{gr}^\bullet(S)\to
\operatorname{Hom}_C^g(\operatorname{gr}^\bullet(A),
\operatorname{gr}^\bullet(A))
\]
où le second membre est l'anneau des endomorphismes gradués de
$\operatorname{gr}^\bullet(A)$. Nous dirons que
$\operatorname{gr}^\bullet(A)$ est muni d'une structure de
\emph{$\operatorname{gr}^\bullet(S)$-$C$-module gradué}. Il résulte alors
de (\hyperref[0.13.6.5-fr]{13.6.5}) que pour
$1\leq r\leq+\infty$, tout
$\overline t\in\operatorname{gr}^j(S)$ définit canoniquement dans les
objets bigradués
$(\mathrm B_r^{pq}(A))_{(p,q)\in\mathbb Z\times\mathbb Z}$,
$(\mathrm Z_r^{pq}(A))_{(p,q)\in\mathbb Z\times\mathbb Z}$ et
$\mathrm E_r(A)=(\mathrm E_r^{pq}(A))_{(p,q)\in\mathbb Z\times\mathbb Z}$
des endomorphismes bigradués de degrés $(s,-s)$~; dans
$\mathrm E_r(A)$ (pour $r$ fini), cet endomorphisme commute à
l'endomorphisme bigradué défini par les $d_r^{pq}$. Comme ces
endomorphismes vérifient les conditions habituelles d'associativité et
de distributivité par rapport à l'addition dans
$\operatorname{gr}^\bullet(S)$ et dans les objets bigradués considérés,
nous dirons pour abréger que ces derniers sont des
\emph{$\operatorname{gr}^\bullet(S)$-$C'$-modules bigradués}~; il est
immédiat que les $\alpha_r^{pq}$ définissent un isomorphisme pour ce type
de structures. Pour tout entier $n$, on notera
$\mathrm B_r^{(n)}(A)$ (resp. $\mathrm Z_r^{(n)}(A)$,
$\mathrm E_r^{(n)}(A)$) le sous-objet gradué de
$\mathrm B_r^\bullet(A)$ (resp. $\mathrm Z_r^\bullet(A)$,
$\mathrm E_r^\bullet(A)$) formé des $\mathrm B_r^{pq}(A)$ (resp.
$\mathrm Z_r^{pq}(A)$, $\mathrm E_r^{pq}(A)$) tels que $p+q=n$ (pour
$1\leq r\leq+\infty$)~; il est immédiat que ce sont des
\emph{$\operatorname{gr}^\bullet(S)$-$C'$-modules gradués}. Enfin, tout
$\overline t\in\operatorname{gr}^j(S)$ définit pour tout $n$ un
endomorphisme gradué de degré $j$ dans l'objet gradué
$\operatorname{gr}^\bullet(\mathrm R^n\mathrm T(A))$, qui est ainsi muni
d'une structure de
$\operatorname{gr}^\bullet(S)$-$C'$-module gradué~; les
$\beta^{pq}$ (pour $p+q=n$) définissent un isomorphisme de
$\mathrm E^{(n)}(A)$ sur
$\operatorname{gr}^\bullet(\mathrm R^n\mathrm T(A))$ pour cette espèce
de structure.

On notera que lorsque $C'$ est la catégorie des groupes abéliens, les
structures de $S$-$C'$-module (resp. de
$\operatorname{gr}^\bullet(S)$-$C'$-module gradué ou bigradué) ne sont
autres que les structures usuelles de $S$-module (resp.
$\operatorname{gr}^\bullet(S)$-module gradué, bigradué).
\end{env}

\subsection{Foncteurs dérivés d'une limite projective d'arguments}
\label{subsection:0.13.7-fr}

\begin{env}[13.7.1]
\label{0.13.7.1-fr}
Soient $C$, $C'$ deux catégories abéliennes, $C$ étant supposée telle
que tout objet de $C$ soit sous-objet d'un objet injectif~; soit
$\mathrm T:C\to C'$ un foncteur additif covariant. Considérons un
système projectif strict $\mathbf A=(A_k)_{k\in\mathbb Z}$ dans $C$,
\emph{limité inférieurement}~; de façon précise, nous supposerons que
$A_k=0$ pour $k<k_0$. Nous associons canoniquement à ce système une
filtration $(\mathrm F^p(A_k))_{p\in\mathbb Z}$ sur chaque $A_k$, par les
formules (\hyperref[0.13.4.1.1-fr]{13.4.1.1}), et comme il s'agit d'un
système projectif strict, les morphismes canoniques
\begin{equation}
\label{0.13.7.1.1-fr}
\tag{13.7.1.1}
\mathrm F^i(A_h)/\mathrm F^j(A_h)\to
\mathrm F^i(A_k)/\mathrm F^j(A_k)\qquad(h\geq k)
\end{equation}
pour $i\leq j\leq k+1$ sont des isomorphismes. Rappelons en outre que
l'on a $\mathrm F^{k_0}(A_k)=A_k$ et $\mathrm F^{k+1}(A_k)=0$ pour tout
$k$.
\end{env}

\begin{env}[13.7.2]
\label{0.13.7.2-fr}
Construisons maintenant pour chaque $k$ une résolution injective
$X_k^\bullet=(X_k^j)_{j\geq0}$ de $A_k$ ayant les propriétés de
(\hyperref[0.13.6.2-fr]{13.6.2}). Les morphismes canoniques
$A_{k+1}\to A_k$ permettent
(\hyperref[0.13.6.3-fr]{13.6.3}) de définir pour chaque $k$ un morphisme
de complexes $X_{k+1}^\bullet\to X_k^\bullet$ compatibles avec les
filtrations, et faisant de
$\mathbf X^\bullet=(X_k^\bullet)_{k\in\mathbb Z}$ un système projectif
de complexes.
\oldpage[0\textsubscript{III}]{76}
On peut en outre supposer que ce système projectif est strict. Pour
cela, on observe qu'en vertu de l'isomorphisme
(\hyperref[0.13.7.1.1-fr]{13.7.1.1}), $A_k$ est isomorphe à
$A_{k+1}/\mathrm F^{k+1}(A_{k+1})$~; on peut donc prendre, dans la
construction de $X_{k+1}^\bullet$, la résolution injective de
$A_{k+1}/\mathrm F^{k+1}(A_{k+1})$ égale à $X_k^\bullet$, et il résulte
de (M, V, 2.3) que la construction du morphisme de complexes
$X_{k+1}^\bullet\to X_k^\bullet$ peut se faire de sorte que ce morphisme
fournisse par passage aux quotients un isomorphisme
\[
X_{k+1}^\bullet/\mathrm F^{k+1}(X_{k+1}^\bullet)
\xrightarrow{\sim}X_k^\bullet
\]
respectant les filtrations, ce qui est la condition de
(\hyperref[0.13.5.1-fr]{13.5.1}).
\end{env}

\begin{env}[13.7.3]
\label{0.13.7.3-fr}
Par construction, le système projectif $\mathrm T(\mathbf X^\bullet)$
des complexes $\mathrm T(X_k^\bullet)$ satisfait aux hypothèses de
(\hyperref[0.13.5.3-fr]{13.5.3}). Les résultats de
(\hyperref[0.13.5.4-fr]{13.5.4}),
(\hyperref[0.13.5.5-fr]{13.5.5}) et
(\hyperref[0.13.5.6-fr]{13.5.6}) sont donc applicables aux suites
spectrales $\mathrm E(\mathrm T(X_k^\bullet))=\mathrm E(A_k)$~; nous
écrirons $\mathrm E_r^{pq}(\mathbf A)$ au lieu de
$\mathrm E_r^{pq}(\mathrm T(\mathbf X^\bullet))$ pour
$1\leq r\leq+\infty$ (cf. (\hyperref[0.13.5.7-fr]{13.5.7}) pour
$r=+\infty$) et de même pour les notations analogues. On notera en
particulier que l'on a
\begin{equation}
\label{0.13.7.3.1-fr}
\tag{13.7.3.1}
\mathrm E_1^{pq}(\mathbf A)=
\mathrm R^{p+q}\mathrm T(\operatorname{gr}^p(\mathbf A))
\end{equation}
en vertu de (\hyperref[0.13.6.4.2-fr]{13.6.4.2}) et du fait que le
système $(\operatorname{gr}^p(A_k))$ est essentiellement constant.

Ces résultats et (\hyperref[0.13.4.3-fr]{13.4.3}) donnent la proposition
suivante, démontrée d'abord par Shih Weishu par une méthode différente
(inédite) :
\end{env}

\begin{proposition}[13.7.4]
\label{0.13.7.4-fr}
\emph{(Shih). --- Soit $n$ un entier. Les deux conditions suivantes sont
équivalentes :}
\begin{enumerate}
\item[\emph{a)}] \emph{Pour tout couple $(p,q)$ tel que $p+q=n$, le
système projectif $(\mathrm E_r^{pq}(\mathbf A))_{r\geq2}$ est
essentiellement constant.}
\item[\emph{b)}] \emph{Le système projectif
$\mathrm R^n\mathrm T(\mathbf A)=
(\mathrm R^n\mathrm T(A_k))_{k\in\mathbb Z}$ vérifie (ML).}
\end{enumerate}
\emph{En outre, lorsque ces conditions sont remplies, on a un
isomorphisme canonique}
\begin{equation}
\label{0.13.7.4.1-fr}
\tag{13.7.4.1}
\operatorname{gr}^p(\mathrm R^n\mathrm T(\mathbf A))
\xrightarrow{\sim}\mathrm E_\infty^{p,n-p}(\mathbf A)
\qquad\emph{pour tout $p\in\mathbb Z$.}
\end{equation}
\end{proposition}
\begin{proof}
En effet, en vertu de (\hyperref[0.13.5.6-fr]{13.5.6}), la condition
\emph{a)} équivaut à dire que le système projectif
$(\mathrm E_\infty^{pq}(A_k))_{k\in\mathbb Z}$ est essentiellement
constant pour $p+q=n$, et d'autre part
$\operatorname{gr}^p(\mathrm R^n\mathrm T(A_k))$ est canoniquement
isomorphe à $\mathrm E_\infty^{p,n-p}(A_k)$, donc il résulte de
(\hyperref[0.13.4.3-fr]{13.4.3}) que \emph{a)} et \emph{b)} sont
équivalentes~; l'isomorphisme
(\hyperref[0.13.7.4.1-fr]{13.7.4.1}) n'est autre que
(\hyperref[0.13.5.6.1-fr]{13.5.6.1}) appliqué au cas envisagé ici.
\end{proof}

\begin{corollary}[13.7.5]
\label{0.13.7.5-fr}
\emph{Soit $\mathcal F=(\mathcal F_k)_{k\in\mathbb N}$ un système
projectif de faisceaux de groupes abéliens satisfaisant aux conditions
\emph{(i)}, \emph{(ii)} et \emph{(iii)} de
(\hyperref[0.13.3.1-fr]{13.3.1}) et soit
$\mathcal F=\varprojlim_k\mathcal F_k$. Supposons que, pour le foncteur
$\mathcal G\mapsto\Gamma(X,\mathcal G)$, le système projectif
$(\mathrm E_r^{pq}(\mathcal F))_{r\in\mathbb Z}$ soit essentiellement
constant pour tout couple $(p,q)$ tel que $p+q=n$ ou $p+q=n+1$.
Considérons sur $H^{n+1}(X,\mathcal F)$ la filtration définie par
$\mathrm F^p(H^{n+1}(X,\mathcal F))=
\operatorname{Ker}(H^{n+1}(X,\mathcal F)\to
H^{n+1}(X,\mathcal F_{p-1}))$. On a alors un isomorphisme canonique}
\begin{equation}
\label{0.13.7.5.1-fr}
\tag{13.7.5.1}
\operatorname{gr}^p(H^{n+1}(X,\mathcal F))
\xrightarrow{\sim}\mathrm E_\infty^{p,n-p+1}(\mathcal F)
\qquad\emph{pour tout $p\in\mathbb Z$.}
\end{equation}
\end{corollary}
\begin{proof}
\oldpage[0\textsubscript{III}]{77}
Il résulte de (\hyperref[0.13.7.4-fr]{13.7.4}) appliqué au cas où $C$
est la catégorie des faisceaux de groupes abéliens sur $X$, $C'$ la
catégorie des groupes abéliens, et $\mathrm T=\Gamma$, que l'on a un
isomorphisme canonique
$\operatorname{gr}^p(\mathrm R^{n+1}\Gamma(\mathcal F))
\xrightarrow{\sim}\mathrm E_\infty^{p,n-p+1}(\mathcal F)$ pour tout
$p\in\mathbb Z$. D'autre part, comme en vertu de
(\hyperref[0.13.7.4-fr]{13.7.4}), le système projectif
$(H^n(X,\mathcal F_k))_{k\in\mathbb Z}$ vérifie (ML), on déduit de
(\hyperref[0.13.3.1-fr]{13.3.1}) un isomorphisme canonique
\begin{equation}
\label{0.13.7.5.1b-fr}
\tag{13.7.5.1}
H^{n+1}(X,\mathcal F)\xrightarrow{\sim}
\varprojlim_k H^{n+1}(X,\mathcal F_k).
\end{equation}

Comme le système projectif $\mathrm R^{n+1}\Gamma(\mathcal F)$ vérifie
(ML) en vertu de (\hyperref[0.13.7.4-fr]{13.7.4}), on a un isomorphisme
canonique
\[
\operatorname{gr}^p(\mathrm R^{n+1}\Gamma(\mathcal F))
\xrightarrow{\sim}
\varprojlim_k\operatorname{gr}^p(H^{n+1}(X,\mathcal F_k))
\]
(\hyperref[0.13.4.3-fr]{13.4.3}), et un isomorphisme canonique
\[
\varprojlim_k\operatorname{gr}^p(H^{n+1}(X,\mathcal F_k))
\xrightarrow{\sim}
\operatorname{gr}^p(\varprojlim_k H^{n+1}(X,\mathcal F_k))
\]
(\hyperref[0.13.4.5-fr]{13.4.5}). Tout revient donc à voir que
l'isomorphisme (\hyperref[0.13.7.5.1-fr]{13.7.5.1}) est compatible avec
les filtrations des deux membres~; mais cela résulte aussitôt des
définitions et de la commutativité du diagramme
\[
\xymatrix{
H^{n+1}(X,\mathcal F)\ar[r]^{\sim}\ar[dr]
&\varprojlim_k H^{n+1}(X,\mathcal F_k)\ar[d]\\
&H^{n+1}(X,\mathcal F_{p-1})
}
\]
pour tout $p$.
\end{proof}

\begin{env}[13.7.6]
\label{0.13.7.6-fr}
Soit $S$ un anneau muni d'une filtration
$(\mathrm F^i(S))_{i\in\mathbb Z}$ telle que $\mathrm F^0(S)=S$ (donc
$\operatorname{gr}^i(S)=0$ pour $i<0$). Supposons que chacun des $A_k$,
muni de la filtration définie dans
(\hyperref[0.13.7.1-fr]{13.7.1}), soit un $S$-$C$-module filtré
(\hyperref[0.13.6.6-fr]{13.6.6}), les morphismes $A_h\to A_k$ pour
$k\leq h$ étant des morphismes pour la structure de $S$-$C$-module
filtré~; nous dirons pour abréger que $\mathbf A$ est un \emph{système
projectif de $S$-$C$-modules filtrés}. Alors il est immédiat que les
morphismes $\mathrm B_r^{pq}(A_h)\to\mathrm B_r^{pq}(A_k)$ et
$\mathrm Z_r^{pq}(A_h)\to\mathrm Z_r^{pq}(A_k)$ pour $k\leq h$,
$1\leq r\leq+\infty$, sont des morphismes pour les structures de
$\operatorname{gr}^\bullet(S)$-$C'$-module bigradué
(\hyperref[0.13.6.5-fr]{13.6.5}), et que les familles
$(\mathrm Z_r^{pq}(\mathbf A))$, $(\mathrm B_r^{pq}(\mathbf A))$ et
$(\mathrm E_r^{pq}(\mathbf A))$ sont des
$\operatorname{gr}^\bullet(S)$-$C'$-modules bigradués pour $r$ fini,
les deux premiers étant des sous-modules de
$(\mathrm E_1^{pq}(\mathbf A))$. On notera encore
$\mathrm Z_r^{(n)}(\mathbf A)$, $\mathrm B_r^{(n)}(\mathbf A)$,
$\mathrm E_r^{(n)}(\mathbf A)$ les sous-objets respectifs des précédents
obtenus en ne prenant que les termes tels que $p+q=n$~; ce sont des
$\operatorname{gr}^\bullet(S)$-$C'$-modules gradués.

Lorsque le système $(\mathrm E_r^{pq}(\mathbf A))_{r\in\mathbb Z}$ est
essentiellement constant, $(\mathrm E_\infty^{pq}(\mathbf A))$ est donc
aussi un $\operatorname{gr}^\bullet(S)$-$C'$-module bigradué, et chaque
$\mathrm E_\infty^{(n)}(\mathbf A)$ un
$\operatorname{gr}^\bullet(S)$-$C'$-module gradué. En outre, les
$\beta^{p,n-p}:\mathrm E_\infty^{p,n-p}(A_k)\xrightarrow{\sim}
\operatorname{gr}^p(\mathrm R^n\mathrm T(A_k))$ constituent pour chaque
$k$ un isomorphisme pour la structure de
$\operatorname{gr}^\bullet(S)$-$C'$-module gradué de
$\mathrm E_\infty^{(n)}(A_k)$ sur
$\operatorname{gr}^\bullet(\mathrm R^n\mathrm T(A_k))$~; si on est dans
les conditions précédentes,
$\beta^{p,n-p}:\mathrm E_\infty^{(n)}(\mathbf A)\xrightarrow{\sim}
\varprojlim_k\operatorname{gr}^\bullet(\mathrm R^n\mathrm T(A_k))$
sera donc aussi un isomorphisme pour ces structures et il en est
évidemment de même de l'isomorphisme canonique
$\operatorname{gr}^\bullet(\mathrm R^n\mathrm T(\mathbf A))
\xrightarrow{\sim}
\varprojlim_k\operatorname{gr}^\bullet(\mathrm R^n\mathrm T(A_k))$,
donc les isomorphismes
(\hyperref[0.13.7.4.1-fr]{13.7.4.1}) constituent un isomorphisme pour les
structures de $\operatorname{gr}^\bullet(S)$-$C'$-module gradué.
\end{env}

\begin{proposition}[13.7.7]
\label{0.13.7.7-fr}
Soit $S$ un anneau noethérien $\mathfrak J$-adique. Supposons que $C$ soit
une catégorie abélienne dont tout objet est isomorphe à un sous-objet d'un
objet injectif, et soit $\mathrm T$ un foncteur covariant additif de $C$ dans
la catégorie des groupes abéliens. Soit
$\mathbf A=(A_k)_{k\in\mathbb Z}$ un
\oldpage[0\textsubscript{III}]{78}
système projectif strict de $S$-$C$-modules filtrés (pour la filtration
$\mathfrak J$-adique sur $S$) limité inférieurement. On suppose que pour un
entier $n$ donné, la condition suivante soit vérifiée~:
\begin{itemize}
\item[$(\mathrm F_n)$] Le $\operatorname{gr}^\bullet(S)$-module gradué
\[
\mathrm E_1^{(m)}(\mathbf A)=
\bigl(\mathrm R^m\mathrm T(\operatorname{gr}^p(\mathbf A))\bigr)_{p\in\mathbb Z}
\]
de (\hyperref[0.13.7.3.1-fr]{13.7.3.1}) est de type fini pour $m=n$ et
$m=n+1$.
\end{itemize}
Dans ces conditions~:
\begin{enumerate}
\item[\rm(i)] Les systèmes projectifs
$(\mathrm R^n\mathrm T(A_k))_{k\in\mathbb Z}$ et
$(\mathrm R^{n+1}\mathrm T(A_k))_{k\in\mathbb Z}$ vérifient (ML).
\item[\rm(ii)] Si on pose
\[
\mathrm R'^{\,n}\mathrm T(\mathbf A)=
\varprojlim_k\mathrm R^n\mathrm T(A_k),
\]
$\mathrm R'^{\,n}\mathrm T(\mathbf A)$ est un $S$-module de type fini.
\item[\rm(iii)] La filtration définie par
\[
\mathrm F^p(\mathrm R'^{\,n}\mathrm T(\mathbf A))=
\operatorname{Ker}\bigl(\mathrm R'^{\,n}\mathrm T(\mathbf A)
\longrightarrow\mathrm R^n\mathrm T(A_{p-1})\bigr)
\quad(p\in\mathbb Z)
\]
sur $\mathrm R'^{\,n}\mathrm T(\mathbf A)$ est $\mathfrak J$-bonne
(c'est-à-dire que
$\mathfrak J\mathrm F^p(\mathrm R'^{\,n}\mathrm T(\mathbf A))
\subset\mathrm F^{p+1}(\mathrm R'^{\,n}\mathrm T(\mathbf A))$ pour tout
$p$, l'égalité des deux membres ayant lieu dès que $p$ est assez grand). En
particulier, la topologie sur $\mathrm R'^{\,n}\mathrm T(\mathbf A)$ définie
par cette filtration est identique à la topologie $\mathfrak J$-adique.
\item[\rm(iv)] Le système projectif
$(\mathrm E_r^{pq}(\mathbf A))_{r\in\mathbb Z}$ est essentiellement constant
pour $p+q=n$ et $p+q=n+1$, $\mathrm E_\infty^{pq}(\mathbf A)$ est donc défini
(\hyperref[0.13.5.7-fr]{13.5.7}) et on a un isomorphisme canonique de
$\operatorname{gr}^\bullet(S)$-modules gradués
\begin{equation}
\label{0.13.7.7.1-fr}
\tag{13.7.7.1}
\operatorname{gr}^p(\mathrm R'^{\,n}\mathrm T(\mathbf A))
\xrightarrow{\sim}\mathrm E_\infty^{p,n-p}(\mathbf A)
\quad(p\in\mathbb Z).
\end{equation}
\end{enumerate}
\end{proposition}

\begin{proof}
On notera que l'isomorphisme (\hyperref[0.13.7.7.1-fr]{13.7.7.1})
permettra de noter $\mathrm R^n\mathrm T(\mathbf A)$, par abus de langage, la
limite projective $\mathrm R'^{\,n}\mathrm T(\mathbf A)$ du système projectif
$\mathrm R^n\mathrm T(A_k)$, compte tenu des isomorphismes
(\hyperref[0.13.7.4.1-fr]{13.7.4.1}).

Comme l'anneau gradué $\operatorname{gr}^\bullet(S)$ est noethérien
(Bourbaki, \emph{Alg. comm.}, chap.~III, \S~2, n\textsuperscript{o}~9,
cor.~5 du th.~2), la suite croissante des sous-$\operatorname{gr}^\bullet(S)$-
modules gradués $\mathrm B_r^{(m)}(\mathbf A)$ de
$\mathrm E_1^{(m)}(\mathbf A)$ (\hyperref[0.13.6.6-fr]{13.6.6}) est
stationnaire pour $m=n$ et $m=n+1$, et par suite la condition b) de
(\hyperref[0.11.1.10-fr]{11.1.10}) est vérifiée. Il s'ensuit que la condition
a) de (\hyperref[0.13.7.4-fr]{13.7.4}) est remplie pour $n$ et pour $n+1$,
et cela prouve déjà (i). En outre, les isomorphismes
(\hyperref[0.13.7.4.1-fr]{13.7.4.1}) (compte tenu des remarques de
(\hyperref[0.13.7.6-fr]{13.7.6})) montrent que
$\operatorname{gr}^\bullet(\mathrm R^n\mathrm T(\mathbf A))$ est un
$\operatorname{gr}^\bullet(S)$-module gradué isomorphe à
\[
\mathrm E_\infty^{(n)}(\mathbf A)=
\mathrm Z_\infty^{(n)}(\mathbf A)/\mathrm B_\infty^{(n)}(\mathbf A)~;
\]
comme $\mathrm Z_\infty^{(n)}(\mathbf A)$ est un sous-module de
$\mathrm E_1^{(n)}(\mathbf A)$, il est de type fini, et il en est donc de
même de $\mathrm E_\infty^{(n)}(\mathbf A)$. En outre, pour la filtration
$(\mathrm F^p(\mathrm R'^{\,n}\mathrm T(\mathbf A)))$, il résulte de
(\hyperref[0.13.4.5-fr]{13.4.5}) que
$\operatorname{gr}^\bullet(\mathrm R^n\mathrm T(\mathbf A))$ et
$\operatorname{gr}^\bullet(\mathrm R'^{\,n}\mathrm T(\mathbf A))$ sont des
$\operatorname{gr}^\bullet(S)$-modules isomorphes, ce qui démontre (iv). Les
assertions (ii) et (iii) seront enfin conséquences des résultats précédents
et du lemme suivant~:
\end{proof}

\begin{lemma}[13.7.7.2]
\label{0.13.7.7.2-fr}
Soient $S$ un anneau noethérien $\mathfrak J$-adique, $M$ un $S$-module muni
d'une filtration co-discrète $(\mathrm F^p(M))_{p\in\mathbb Z}$ telle que
$\mathfrak J\mathrm F^p(M)\subset\mathrm F^{p+1}(M)$ (ce qui exprime que
$M$ est un module filtré sur l'anneau $S$ filtré par la filtration
$\mathfrak J$-adique). Supposons en outre $M$ séparé pour la topologie définie
par la filtration $(\mathrm F^p(M))$. Alors les conditions suivantes sont
équivalentes~:
\begin{enumerate}
\item[\rm a)] $M$ est un $S$-module de type fini et $(\mathrm F^p(M))$ une
filtration $\mathfrak J$-bonne.
\item[\rm b)] $\operatorname{gr}^\bullet(M)$ est un
$\operatorname{gr}^\bullet(S)$-module de type fini.
\item[\rm c)] Les $\operatorname{gr}^p(M)$ sont des $S$-modules de type fini
et pour $p$ assez grand les homomorphismes canoniques
\begin{equation}
\label{0.13.7.7.3-fr}
\tag{13.7.7.3}
\mathfrak J\otimes_S\operatorname{gr}^p(M)\longrightarrow
\operatorname{gr}^{p+1}(M)
\end{equation}
\oldpage[0\textsubscript{III}]{79}
(déduits de
$\mathfrak J\otimes_S\mathrm F^p(M)\to\mathrm F^{p+1}(M)$, compte tenu de ce
que l'image de l'homomorphisme composé
\[
\mathfrak J\otimes_S\mathrm F^{p+1}(M)\longrightarrow
\mathfrak J\otimes_S\mathrm F^p(M)\longrightarrow\mathrm F^{p+1}(M)
\]
est $\mathfrak J\mathrm F^{p+1}(M)\subset\mathrm F^{p+2}(M)$) sont
surjectifs.
\end{enumerate}
Pour la démonstration, voir Bourbaki, \emph{Alg. comm.}, chap.~III, \S~3,
n\textsuperscript{o}~1, prop.~3.
\end{lemma}

\begin{env}[13.7.7.4]
\label{0.13.7.7.4-fr}
Pour appliquer le lemme (\hyperref[0.13.7.7.2-fr]{13.7.7.2}), il reste à
observer que la topologie définie sur
$\mathrm R'^{\,n}\mathrm T(\mathbf A)$ par la filtration considérée fait de
$\mathrm R'^{\,n}\mathrm T(\mathbf A)$ un $S$-module séparé et complet,
cette topologie étant celle de la limite projective des groupes discrets
$\mathrm R^n\mathrm T(A_k)$~; d'autre part, si $A_k=0$ pour $k<k_0$, on a
aussi $\mathrm R^n\mathrm T(A_k)=0$ pour $k<k_0$, donc
$\mathrm F^{k_0}(\mathrm R'^{\,n}\mathrm T(\mathbf A))=
\mathrm R'^{\,n}\mathrm T(\mathbf A)$, et on est bien dans les conditions
d'application du lemme.
\end{env}

\begin{corollary}[13.7.8]
\label{0.13.7.8-fr}
Si l'hypothèse $(\mathrm F_n)$ est vérifiée, on a, pour tout élément $f\in S$,
un isomorphisme canonique
\begin{equation}
\label{0.13.7.8.1-fr}
\tag{13.7.8.1}
\varprojlim_k\bigl((\mathrm R^n\mathrm T(A_k))_f\bigr)
\xrightarrow{\sim}\mathrm R^n\mathrm T(\mathbf A)\otimes_S S_{\{f\}}.
\end{equation}
\end{corollary}

\begin{proof}
En effet, $\mathrm R^n\mathrm T(\mathbf A)$ est un $S$-module de type fini,
$S_{\{f\}}$ une $S$-algèbre adique noethérienne
(\hyperref[0.7.6.11-fr]{0\textsubscript{I}, 7.6.11}), séparée complétée de
$S_f$ pour la topologie $\mathfrak J$-préadique
(\hyperref[0.7.6.2-fr]{0\textsubscript{I}, 7.6.2}). On conclut de
(\hyperref[0.7.7.8-fr]{0\textsubscript{I}, 7.7.8}) et
(\hyperref[0.7.7.1-fr]{0\textsubscript{I}, 7.7.1}) que
$\mathrm R^n\mathrm T(\mathbf A)\otimes_S S_{\{f\}}$ est isomorphe au séparé
complété de $\mathrm R^n\mathrm T(\mathbf A)\otimes_S S_f$ pour la topologie
$\mathfrak J$-préadique~; un système fondamental de voisinages de $0$ pour
cette topologie est
$(\mathfrak J^p\mathrm R^n\mathrm T(\mathbf A))\otimes_S S_f$, donc
$\mathrm F^p(\mathrm R^n\mathrm T(\mathbf A))\otimes_S S_f$ est aussi un tel
système~; ce dernier est le noyau de l'application canonique
$(\mathrm R^n\mathrm T(\mathbf A))_f\to
(\mathrm R^n\mathrm T(A_{p-1}))_f$ et par suite le groupe séparé associé à
$\mathrm R^n\mathrm T(\mathbf A)\otimes_S S_f$ s'identifie à un sous-groupe
$G$ de $\varprojlim_k((\mathrm R^n\mathrm T(A_k))_f)$. Mais le système
projectif $((\mathrm R^n\mathrm T(A_k))_f)$ vérifie évidemment la condition
(ML), et l'image de $(\mathrm R^n\mathrm T(\mathbf A))_f$ dans chacun des
$(\mathrm R^n\mathrm T(A_k))_f$ est égale à l'image commune des
$(\mathrm R^n\mathrm T(A_h))_f$ pour $h\geq k$ assez grand. On en conclut
aussitôt que $G$ est \emph{partout dense} dans
$\varprojlim_k((\mathrm R^n\mathrm T(A_k))_f)$, et comme ce dernier groupe est
séparé et complet, le corollaire est démontré.
\end{proof}

\begin{flushright}
\textit{(A suivre.)}
\end{flushright}

\clearpage
\thispagestyle{empty}
\null
\clearpage
\thispagestyle{plain}
\markboth{ÉTUDE COHOMOLOGIQUE DES FAISCEAUX COHÉRENTS}{CHAPITRE III}
\oldpage[III]{81}
\phantomsection
\label{chapter:III-fr}

\begin{center}
{\large CHAPITRE III\par}
\vspace{1.5em}
{\Large\bfseries ÉTUDE COHOMOLOGIQUE\par}
{\Large\bfseries DES FAISCEAUX COHÉRENTS\par}
\vspace{1.8em}
{\bfseries Sommaire\par}
\end{center}

\medskip
\begin{tabular}{@{}r@{\ }p{0.82\textwidth}@{}}
\S~1. & Cohomologie des schémas affines.\\
\S~2. & Étude cohomologique des morphismes projectifs.\\
\S~3. & Le théorème de finitude pour les morphismes propres.\\
\S~4. & Le théorème fondamental des morphismes propres~; applications.\\
\S~5. & Un théorème d'existence de faisceaux algébriques cohérents.\\
\S~6. & Foncteurs Tor et Ext locaux et globaux~; formule de Künneth.\\
\S~7. & Étude du changement de base dans les foncteurs cohomologiques
covariants de Modules.\\
\S~8. & Le théorème de dualité sur les fibrés projectifs.\\
\S~9. & Cohomologie relative et cohomologie locale~; dualité locale.\\
\S~10. & Relations entre cohomologie projective et cohomologie locale.
Technique de complétion formelle le long d'un diviseur.\\
\S~11. & Groupes de Picard globaux et locaux\footnotemark.
\end{tabular}

\medskip
Ce chapitre donne les théorèmes fondamentaux sur la cohomologie des faisceaux
algébriques cohérents, à l'exception de ceux découlant de la théorie des
résidus (théorèmes de dualité), qui feront l'objet d'un chapitre ultérieur.
Parmi les premiers, il y a essentiellement six théorèmes fondamentaux,
faisant l'objet des six premiers paragraphes du présent chapitre. Ces
résultats seront dans la suite des outils essentiels, même dans des questions
qui ne sont pas de nature proprement cohomologique~; le lecteur en verra les
premiers exemples dès le \S~4. Le \S~7 donne des résultats de nature plus
technique, mais d'un usage constant dans les applications. Enfin, dans les
\S\S~8 à 11, nous développerons certains résultats, liés à la dualité des
faisceaux cohérents, particulièrement importants pour les applications, et
qui peuvent s'exposer antérieurement à la théorie générale des résidus.

\footnotetext{Le chapitre IV ne dépend pas des \S\S~8 à 11, et sera sans doute
publié avant ces derniers.}

\oldpage[III]{82}
Le contenu des \S\S~1 et 2 est dû à J.-P. Serre, et le lecteur constatera que
nous n'avons eu qu'à suivre l'exposé de (FAC). Les \S\S~8 et 9 sont également
inspirés par (FAC) (les transpositions nécessitées par les contextes
différents étant toutefois moins évidentes). Enfin, comme nous l'avons dit
dans l'Introduction, le \S~4 doit être considéré comme la mise en forme, en
langage moderne, du «~théorème d'invariance~» fondamental de la «~théorie des
fonctions holomorphes~» de Zariski.

Signalons enfin que les résultats du n\textsuperscript{o}~3.4 (et les
propositions préliminaires de (0, 13.4 à 13.7)) ne seront pas utilisés dans la
suite du chap.~III et peuvent donc être omis en première lecture.

\setcounter{section}{0}
\section{Cohomologie des schémas affines}
\label{section:III.1-fr}

\subsection{Rappels sur le complexe de l'algèbre extérieure}
\label{subsection:III.1.1-fr}

\begin{env}[1.1.1]
\label{III.1.1.1-fr}
Soient $A$ un anneau, $\mathbf f=(f_i)_{1\leq i\leq r}$ un système de $r$
éléments de $A$. Le \emph{complexe de l'algèbre extérieure
$K_\bullet(\mathbf f)$} correspondant à $\mathbf f$ est un complexe de chaînes
(G, I, 2.2) se définissant de la façon suivante~: le $A$-module gradué
$K_\bullet(\mathbf f)$ est égal à l'\emph{algèbre extérieure
$\bigwedge(A^r)$}, graduée de la façon usuelle, et l'opérateur bord est la
\emph{multiplication intérieure $i_{\mathbf f}$} par $\mathbf f$ considéré
comme élément du dual $(A^r)^\vee$~; on rappelle que $i_{\mathbf f}$ est une
\emph{antidérivation} de degré $-1$ de $\bigwedge(A^r)$, et que si
$(\mathbf e_i)_{1\leq i\leq r}$ est la base canonique de $A^r$, on a
$i_{\mathbf f}(\mathbf e_i)=f_i$~; la vérification de la condition
$i_{\mathbf f}\circ i_{\mathbf f}=0$ est immédiate.

Une définition équivalente est la suivante~: pour chaque $i$, on considère
un complexe de chaînes $K_\bullet(f_i)$ défini comme suit~:
$K_0(f_i)=K_1(f_i)=A$, $K_n(f_i)=0$ pour $n\neq0,1$~; l'opérateur bord est
défini par la condition que $d_1:A\to A$ est la \emph{multiplication par
$f_i$}. On prend alors pour $K_\bullet(\mathbf f)$ le \emph{produit tensoriel
$K_\bullet(f_1)\otimes K_\bullet(f_2)\otimes\cdots\otimes K_\bullet(f_r)$}
(G, I, 2.7) muni de son degré total~; la vérification de l'isomorphisme de ce
complexe et du complexe défini plus haut est immédiate.
\end{env}

\begin{env}[1.1.2]
\label{III.1.1.2-fr}
Pour tout $A$-module $M$, on définit le \emph{complexe de chaînes}
\begin{equation}
\label{III.1.1.2.1-fr}
\tag{1.1.2.1}
K_\bullet(\mathbf f,M)=K_\bullet(\mathbf f)\otimes_A M
\end{equation}
et le \emph{complexe de cochaînes} (G, I, 2.2)
\begin{equation}
\label{III.1.1.2.2-fr}
\tag{1.1.2.2}
K^\bullet(\mathbf f,M)=\operatorname{Hom}_A(K_\bullet(\mathbf f),M).
\end{equation}
Si $g$ est une $k$-cochaîne de ce dernier complexe, et si on pose
\[
g(i_1,\ldots,i_k)=g(\mathbf e_{i_1}\wedge\cdots\wedge\mathbf e_{i_k}),
\]
$g$ s'identifie à une application \emph{alternée} de $[1,r]^k$ dans $M$, et
il résulte des définitions précédentes que l'on a
\begin{equation}
\label{III.1.1.2.3-fr}
\tag{1.1.2.3}
d^kg(i_1,i_2,\ldots,i_{k+1})=
\sum_{h=1}^{k+1}(-1)^{h-1}f_{i_h}
g(i_1,\ldots,\widehat{i_h},\ldots,i_{k+1}).
\end{equation}
\end{env}

\oldpage[III]{83}
\begin{env}[1.1.3]
\label{III.1.1.3-fr}
Des complexes précédents, on déduit comme d'ordinaire les \emph{$A$-modules
d'homologie et de cohomologie} (G, I, 2.2)
\begin{align}
\label{III.1.1.3.1-fr}
\tag{1.1.3.1}
H_\bullet(\mathbf f,M)&=H_\bullet(K_\bullet(\mathbf f,M)),\\
\label{III.1.1.3.2-fr}
\tag{1.1.3.2}
H^\bullet(\mathbf f,M)&=H^\bullet(K^\bullet(\mathbf f,M)).
\end{align}
On définit d'ailleurs un $A$-isomorphisme
$K_i(\mathbf f,M)\xrightarrow{\sim}K^{r-i}(\mathbf f,M)$ en faisant
correspondre à toute chaîne
$z=\sum(\mathbf e_{i_1}\wedge\cdots\wedge\mathbf e_{i_k})\otimes
z_{i_1\ldots i_k}$ la cochaîne $g_z$ telle que
$g_z(j_1,\ldots,j_{r-k})=\varepsilon z_{i_1\ldots i_k}$, où
$(j_h)_{1\leq h\leq r-k}$ est la suite strictement croissante complémentaire
de la suite strictement croissante $(i_h)_{1\leq h\leq k}$ dans $[1,r]$ et
$\varepsilon=(-1)^\nu$, $\nu$ étant le nombre d'inversions de la permutation
$i_1,\ldots,i_k,j_1,\ldots,j_{r-k}$ de $[1,r]$. On vérifie que
$g_{dz}=d(g_z)$, ce qui donne un isomorphisme
\begin{equation}
\label{III.1.1.3.3-fr}
\tag{1.1.3.3}
H^i(\mathbf f,M)\xrightarrow{\sim}H_{r-i}(\mathbf f,M)
\qquad\text{pour }0\leq i\leq r.
\end{equation}
Dans ce chapitre, nous aurons surtout à considérer les modules de cohomologie
$H^\bullet(\mathbf f,M)$.

Pour un $\mathbf f$ donné, il est immédiat (G, I, 2.1) que
$M\mapsto H^\bullet(\mathbf f,M)$ est un \emph{foncteur cohomologique}
(T, II, 2.1), de la catégorie des $A$-modules dans celle des $A$-modules
gradués, nul pour les degrés $<0$ et $>r$. En outre, on a
\begin{equation}
\label{III.1.1.3.4-fr}
\tag{1.1.3.4}
H^0(\mathbf f,M)=\operatorname{Hom}_A(A/(\mathbf f),M)
\end{equation}
en désignant par $(\mathbf f)$ l'idéal de $A$ engendré par
$f_1,\ldots,f_r$~; cela résulte aussitôt de
(\hyperref[III.1.1.2.3-fr]{1.1.2.3}), et il est clair que
$H^0(\mathbf f,M)$ s'identifie au sous-module de $M$ \emph{annulé par
$(\mathbf f)$}. De même, on a d'après
(\hyperref[III.1.1.2.3-fr]{1.1.2.3})
\begin{equation}
\label{III.1.1.3.5-fr}
\tag{1.1.3.5}
H^r(\mathbf f,M)=M/\Bigl(\sum_{i=1}^r f_iM\Bigr)
=(A/(\mathbf f))\otimes_A M.
\end{equation}
Nous utiliserons le résultat connu suivant, dont nous rappellerons une
démonstration pour être complet~:
\end{env}

\begin{proposition}[1.1.4]
\label{III.1.1.4-fr}
Soient $A$ un anneau, $\mathbf f=(f_i)_{1\leq i\leq r}$ une famille finie
d'éléments de $A$, $M$ un $A$-module. Si, pour $1\leq i\leq r$, l'homothétie
$z\mapsto f_i z$ dans $M_{i-1}=M/(f_1M+\cdots+f_{i-1}M)$ est injective, on a
$H^i(\mathbf f,M)=0$ pour $i\neq r$.
\end{proposition}

\begin{proof}
Il suffit en effet de prouver que $H_i(\mathbf f,M)=0$ pour tout $i>0$ en
vertu de (\hyperref[III.1.1.3.3-fr]{1.1.3.3}). Raisonnons par récurrence sur
$r$, le cas $r=0$ étant trivial. Posons
$\mathbf f'=(f_i)_{1\leq i\leq r-1}$~; cette famille vérifie les conditions de
l'énoncé, donc, si on pose $L_\bullet=K_\bullet(\mathbf f',M)$, on a
$H_i(L_\bullet)=0$ pour $i>0$ par hypothèse, et
$H_0(L_\bullet)=M_{r-1}$ en vertu de
(\hyperref[III.1.1.3.3-fr]{1.1.3.3}) et
(\hyperref[III.1.1.3.5-fr]{1.1.3.5}). Posons pour abréger
$K_\bullet=K_\bullet(f_r)=K_0\oplus K_1$, avec $K_0=K_1=A$,
$d_1:K_1\to K_0$ étant la multiplication par $f_r$~; on a par définition
(\hyperref[III.1.1.1-fr]{1.1.1})
$K_\bullet(\mathbf f,M)=K_\bullet\otimes_A L_\bullet$. Or, on a le lemme
suivant~:

\begin{lemma}[1.1.4.1]
\label{III.1.1.4.1-fr}
Soit $K_\bullet$ un complexe de chaînes formé de $A$-modules libres, nuls sauf
en dimensions $0$ et $1$. Pour tout complexe de chaînes $L_\bullet$ formé de
$A$-modules, on a une suite exacte
\[
0\longrightarrow H_0(K_\bullet\otimes H_p(L_\bullet))
\longrightarrow H_p(K_\bullet\otimes L_\bullet)
\longrightarrow H_1(K_\bullet\otimes H_{p-1}(L_\bullet))
\longrightarrow0
\]
pour tout indice $p$.
\end{lemma}

\oldpage[III]{84}
C'est un cas particulier d'une suite exacte de termes de bas degré de la
suite spectrale de Künneth (M, XVII, 5.2 a) et G, I, 5.5.2)~; on peut le
démontrer directement de la façon suivante. Considérons $K_0$ et $K_1$ comme
des complexes de chaînes (nuls en dimensions $\neq0$ et $\neq1$
respectivement)~; on a alors une suite exacte de complexes
\[
0\longrightarrow K_0\otimes L_\bullet
\longrightarrow K_\bullet\otimes L_\bullet
\longrightarrow K_1\otimes L_\bullet\longrightarrow0
\]
à laquelle nous pouvons appliquer la suite exacte d'homologie
\[
\cdots\longrightarrow H_{p+1}(K_1\otimes L_\bullet)
\xrightarrow{\partial}H_p(K_0\otimes L_\bullet)
\longrightarrow H_p(K_\bullet\otimes L_\bullet)
\longrightarrow H_p(K_1\otimes L_\bullet)
\xrightarrow{\partial}H_{p-1}(K_0\otimes L_\bullet)
\longrightarrow\cdots.
\]
Mais il est évident que
$H_p(K_0\otimes L_\bullet)=K_0\otimes H_p(L_\bullet)$ et
$H_p(K_1\otimes L_\bullet)=K_1\otimes H_{p-1}(L_\bullet)$ pour tout $p$~; en
outre, on vérifie aussitôt que l'opérateur
$\partial:K_1\otimes H_p(L_\bullet)\to K_0\otimes H_p(L_\bullet)$ n'est autre
que $d_1\otimes1$~; le lemme résulte donc de la suite exacte précédente et de
la définition de $H_0(K_\bullet\otimes H_p(L_\bullet))$ et de
$H_1(K_\bullet\otimes H_{p-1}(L_\bullet))$.

Ce lemme étant établi, la fin de la démonstration de
(\hyperref[III.1.1.4-fr]{1.1.4}) est immédiate~: l'hypothèse de récurrence et
ce lemme (\hyperref[III.1.1.4.1-fr]{1.1.4.1}) donnent
$H_p(K_\bullet\otimes L_\bullet)=0$ pour $p\geq2$~; en outre, si on prouve que
$H_1(K_\bullet,H_0(L_\bullet))=0$, on déduira aussi
$H_1(K_\bullet\otimes L_\bullet)=0$ du lemme
(\hyperref[III.1.1.4.1-fr]{1.1.4.1})~; mais par définition,
$H_1(K_\bullet,H_0(L_\bullet))$ n'est autre que le noyau de l'homothétie
$z\mapsto f_rz$ dans $M_{r-1}$, et comme par hypothèse ce noyau est nul, cela
achève la démonstration.
\end{proof}

\begin{env}[1.1.5]
\label{III.1.1.5-fr}
Soit $\mathbf g=(g_i)_{1\leq i\leq r}$ une seconde suite de $r$ éléments de
$A$, et posons $\mathbf f\mathbf g=(f_ig_i)_{1\leq i\leq r}$. On peut définir
un homomorphisme canonique de complexes
\begin{equation}
\label{III.1.1.5.1-fr}
\tag{1.1.5.1}
\varphi_{\mathbf g}:K_\bullet(\mathbf f\mathbf g)\longrightarrow
K_\bullet(\mathbf f)
\end{equation}
comme l'extension canonique à l'algèbre extérieure $\bigwedge(A^r)$ de
l'application $A$-linéaire
$(x_1,\ldots,x_r)\mapsto(g_1x_1,\ldots,g_rx_r)$ de $A^r$ dans lui-même. Pour
voir qu'on a bien un homomorphisme de complexes, il suffit de remarquer, de
façon générale, que si $u:E\to F$ est une application $A$-linéaire,
$\mathbf x\in F^\vee$ et $\mathbf y={}^tu(\mathbf x)\in E^\vee$, alors on a
la formule
\begin{equation}
\label{III.1.1.5.2-fr}
\tag{1.1.5.2}
(\bigwedge u)\circ i_{\mathbf y}=i_{\mathbf x}\circ(\bigwedge u)~;
\end{equation}
en effet, les deux membres sont des antidérivations de $\bigwedge F$, et il
suffit de vérifier qu'ils coïncident dans $F$, ce qui résulte aussitôt des
définitions.

Lorsqu'on identifie $K_\bullet(\mathbf f)$ au produit tensoriel des
$K_\bullet(f_i)$ (\hyperref[III.1.1.1-fr]{1.1.1}), $\varphi_{\mathbf g}$ est
le produit tensoriel des $\varphi_{g_i}$, où $\varphi_{g_i}$ se réduit à
l'identité en degré $0$ et à la multiplication par $g_i$ en degré $1$.
\end{env}

\begin{env}[1.1.6]
\label{III.1.1.6-fr}
En particulier, pour tout couple d'entiers $m,n$ tels que $0\leq n\leq m$,
on a des homomorphismes de complexes
\begin{equation}
\label{III.1.1.6.1-fr}
\tag{1.1.6.1}
\varphi_{\mathbf f^{m-n}}:K_\bullet(\mathbf f^m)\longrightarrow
K_\bullet(\mathbf f^n)
\end{equation}
et par suite des homomorphismes
\begin{align}
\label{III.1.1.6.2-fr}
\tag{1.1.6.2}
\varphi_{\mathbf f^{m-n}} &:K^\bullet(\mathbf f^n,M)\longrightarrow
K^\bullet(\mathbf f^m,M),\\
\label{III.1.1.6.3-fr}
\tag{1.1.6.3}
\varphi_{\mathbf f^{m-n}} &:H^\bullet(\mathbf f^n,M)\longrightarrow
H^\bullet(\mathbf f^m,M).
\end{align}

\oldpage[III]{85}
Ces derniers vérifient évidemment la condition de transitivité
$\varphi_{\mathbf f^{m-p}}=
\varphi_{\mathbf f^{m-n}}\circ\varphi_{\mathbf f^{n-p}}$ pour $p\leq n\leq m$~;
ils définissent donc deux \emph{systèmes inductifs} de $A$-modules~; on posera
\begin{align}
\label{III.1.1.6.4-fr}
\tag{1.1.6.4}
C^\bullet((\mathbf f),M)&=\varinjlim_n K^\bullet(\mathbf f^n,M),\\
\label{III.1.1.6.5-fr}
\tag{1.1.6.5}
H^\bullet((\mathbf f),M)&=H^\bullet(C^\bullet((\mathbf f),M))
=\varinjlim_n H^\bullet(\mathbf f^n,M),
\end{align}
la dernière égalité provenant du fait que le passage à la limite inductive
permute au foncteur $H^\bullet$ (G, I, 2.1). On verra ultérieurement
(\hyperref[III.1.4.3-fr]{1.4.3}) que $H^\bullet((\mathbf f),M)$ ne dépend en
fait que de l'\emph{idéal $(\mathbf f)$} dans $A$ (et même de la topologie
$(\mathbf f)$-préadique sur $A$), ce qui justifie les notations.

Il est clair que $M\mapsto C^\bullet((\mathbf f),M)$ est un foncteur
$A$-linéaire exact, et $M\mapsto H^\bullet((\mathbf f),M)$ un foncteur
cohomologique.
\end{env}

\begin{env}[1.1.7]
\label{III.1.1.7-fr}
Soient $\mathbf f=(f_i)\in A^r$, $\mathbf g=(g_i)\in A^r$~; désignons par
$e_{\mathbf g}$ la multiplication à gauche par le vecteur
$\mathbf g\in A^r$ dans l'algèbre extérieure $\bigwedge(A^r)$~; on sait que
l'on a la \emph{formule d'homotopie}
\begin{equation}
\label{III.1.1.7.1-fr}
\tag{1.1.7.1}
i_{\mathbf f}e_{\mathbf g}+e_{\mathbf g}i_{\mathbf f}
=\langle\mathbf g,\mathbf f\rangle 1
\end{equation}
dans le $A$-module $A^r$ ($1$ désignant l'automorphisme identique de $A^r$)~;
cette relation signifie aussi que, \emph{dans le complexe
$K_\bullet(\mathbf f)$}, on a
\begin{equation}
\label{III.1.1.7.2-fr}
\tag{1.1.7.2}
de_{\mathbf g}+e_{\mathbf g}d=\langle\mathbf g,\mathbf f\rangle 1.
\end{equation}
Si l'idéal $(\mathbf f)$ est égal à $A$, il existe $\mathbf g\in A^r$ tel que
$\langle\mathbf g,\mathbf f\rangle=\sum_{i=1}^r g_if_i=1$. Par suite
(G, I, 2.4)~:
\end{env}

\begin{proposition}[1.1.8]
\label{III.1.1.8-fr}
Supposons que l'idéal $(\mathbf f)$ engendré par les $f_i$ soit égal à $A$.
Alors le complexe $K_\bullet(\mathbf f)$ est homotopiquement trivial, et il en
est donc de même des complexes $K_\bullet(\mathbf f,M)$ et
$K^\bullet(\mathbf f,M)$ pour tout $A$-module $M$.
\end{proposition}

\begin{corollary}[1.1.9]
\label{III.1.1.9-fr}
Si $(\mathbf f)=A$, on a $H^\bullet(\mathbf f,M)=0$ et
$H^\bullet((\mathbf f),M)=0$ pour tout $A$-module $M$.
\end{corollary}

\begin{proof}
En effet, on a alors $(\mathbf f^n)=A$ pour tout $n$.
\end{proof}

\begin{remark}[1.1.10]
\label{III.1.1.10-fr}
Avec les mêmes notations que ci-dessus, soient $X=\operatorname{Spec}(A)$,
$Y$ le sous-schéma fermé de $X$ défini par l'idéal $(\mathbf f)$. Nous
prouverons au \S~9 que $H^\bullet((\mathbf f),M)$ est isomorphe à la
cohomologie $H_Y^\bullet(X,\widetilde M)$ correspondant à l'antifiltre $\Phi$
des parties fermées de $Y$ (T, 3.2). Nous montrerons aussi que la prop.
(\hyperref[III.1.2.3-fr]{1.2.3}) appliquée à $X$ et à
$\mathcal F=\widetilde M$, est un cas particulier d'une suite exacte de
cohomologie
\[
\cdots\longrightarrow H_Y^p(X,\mathcal F)\longrightarrow H^p(X,\mathcal F)
\longrightarrow H^p(X-Y,\mathcal F)\longrightarrow H_Y^{p+1}(X,\mathcal F)
\longrightarrow\cdots.
\]
\end{remark}

\subsection{Cohomologie de Čech d'un recouvrement ouvert}
\label{subsection:III.1.2-fr}

\begin{notation}[1.2.1]
\label{III.1.2.1-fr}
Dans ce numéro, on notera~:
\begin{itemize}
\item $X$ un préschéma~;
\item $\mathcal F$ un $\mathcal O_X$-Module quasi-cohérent~;
\oldpage[III]{86}
\item $A=\Gamma(X,\mathcal O_X)$, $M=\Gamma(X,\mathcal F)$~;
\item $\mathbf f=(f_i)_{1\leq i\leq r}$ un système fini d'éléments de $A$~;
\item $U_i=X_{f_i}$ l'ensemble ouvert
(\hyperref[0.5.5.2-fr]{0\textsubscript{I}, 5.5.2}) des $x\in X$ tels que
$f_i(x)\neq0$~;
\item $U=\bigcup_{i=1}^r U_i$~;
\item $\mathfrak U$ le recouvrement $(U_i)_{1\leq i\leq r}$ de $U$.
\end{itemize}
\end{notation}

\begin{env}[1.2.2]
\label{III.1.2.2-fr}
Supposons que $X$ soit, ou bien un préschéma dont l'espace sous-jacent est
\emph{noethérien}, ou bien un \emph{schéma} dont l'espace sous-jacent soit
\emph{quasi-compact}. On sait alors (\hyperref[I.9.3.3-fr]{I, 9.3.3}) que
l'on a $\Gamma(U_i,\mathcal F)=M_{f_i}$. Nous poserons
\[
U_{i_0i_1\ldots i_p}=\bigcap_{k=0}^p U_{i_k}
=X_{f_{i_0}f_{i_1}\cdots f_{i_p}}
\]
(\hyperref[0.5.5.3-fr]{0\textsubscript{I}, 5.5.3})~; on a donc aussi
\begin{equation}
\label{III.1.2.2.1-fr}
\tag{1.2.2.1}
\Gamma(U_{i_0i_1\ldots i_p},\mathcal F)=M_{f_{i_0}f_{i_1}\cdots f_{i_p}}.
\end{equation}
Or (\hyperref[0.1.6.1-fr]{0\textsubscript{I}, 1.6.1})
$M_{f_{i_0}f_{i_1}\cdots f_{i_p}}$ s'identifie à la limite inductive
$\varinjlim_n M_{i_0i_1\ldots i_p}^{(n)}$, où le système inductif est formé
des $M_{i_0i_1\ldots i_p}^{(n)}=M$, l'homomorphisme
$\varphi_{nm}:M_{i_0i_1\ldots i_p}^{(m)}\to
M_{i_0i_1\ldots i_p}^{(n)}$ étant la multiplication par
$(f_{i_0}f_{i_1}\cdots f_{i_p})^{n-m}$ pour $m\leq n$. Désignons par
$C_n^p(M)$ l'ensemble des applications \emph{alternées} de $[1,r]^{p+1}$
dans $M$ (pour tout $n$)~; ces $A$-modules forment encore un système inductif
pour les $\varphi_{nm}$. Si $C^p(\mathfrak U,\mathcal F)$ est le groupe des
$p$-cochaînes \emph{alternées} de Čech relatif au recouvrement $\mathfrak U$,
à coefficients dans $\mathcal F$ (G, II, 5.1), il résulte de ce qui précède
que l'on peut écrire
\begin{equation}
\label{III.1.2.2.2-fr}
\tag{1.2.2.2}
C^p(\mathfrak U,\mathcal F)=\varinjlim_n C_n^p(M).
\end{equation}
Or, avec les notations de (\hyperref[III.1.1.2-fr]{1.1.2}), $C_n^p(M)$
s'identifie à $K^{p+1}(\mathbf f^n,M)$, et l'application $\varphi_{nm}$
s'identifie à l'application $\varphi_{\mathbf f^{n-m}}$ définie dans
(\hyperref[III.1.1.6-fr]{1.1.6}). On a donc, pour tout $p\geq0$, un
isomorphisme canonique fonctoriel en $\mathcal F$
\begin{equation}
\label{III.1.2.2.3-fr}
\tag{1.2.2.3}
C^p(\mathfrak U,\mathcal F)\xrightarrow{\sim}C^{p+1}((\mathbf f),M).
\end{equation}
En outre, la formule (\hyperref[III.1.1.2.3-fr]{1.1.2.3}) et la définition de
la cohomologie d'un recouvrement (G, II, 5.1) montrent que les isomorphismes
(\hyperref[III.1.2.2.3-fr]{1.2.2.3}) sont compatibles avec les opérateurs
cobords.
\end{env}

\begin{proposition}[1.2.3]
\label{III.1.2.3-fr}
Si $X$ est un préschéma dont l'espace sous-jacent est noethérien, ou un
schéma dont l'espace sous-jacent est quasi-compact, il existe un isomorphisme
canonique fonctoriel en $\mathcal F$
\begin{equation}
\label{III.1.2.3.1-fr}
\tag{1.2.3.1}
H^p(\mathfrak U,\mathcal F)\xrightarrow{\sim}H^{p+1}((\mathbf f),M)
\qquad\text{pour tout }p\geq1.
\end{equation}
On a en outre une suite exacte fonctorielle en $\mathcal F$
\begin{equation}
\label{III.1.2.3.2-fr}
\tag{1.2.3.2}
0\longrightarrow H^0((\mathbf f),M)\longrightarrow M
\longrightarrow H^0(\mathfrak U,\mathcal F)
\longrightarrow H^1((\mathbf f),M)\longrightarrow0.
\end{equation}
\end{proposition}

\oldpage[III]{87}
\begin{proof}
Les relations (\hyperref[III.1.2.3.1-fr]{1.2.3.1}) sont en effet conséquences
immédiates de ce qu'on a vu dans (\hyperref[III.1.2.2-fr]{1.2.2}). D'autre
part, on a
$C^0(\mathfrak U,\mathcal F)=C^1((\mathbf f),M)$~;
$H^0(\mathfrak U,\mathcal F)$ s'identifie par suite au sous-groupe des
$1$-cocycles de $C^1((\mathbf f),M)$~; comme
$M=C^0((\mathbf f),M)$, la suite exacte
(\hyperref[III.1.2.3.2-fr]{1.2.3.2}) n'est autre que celle qui résulte de la
définition des groupes de cohomologie $H^0((\mathbf f),M)$ et
$H^1((\mathbf f),M)$.
\end{proof}

\begin{corollary}[1.2.4]
\label{III.1.2.4-fr}
Supposons que les $X_{f_i}$ soient quasi-compacts et qu'il existe des
$g_i\in\Gamma(U,\mathcal F)$ tels que
$\sum_i g_i(f_i\vert U)=1\vert U$. Alors, pour tout
$(\mathcal O_X\vert U)$-Module quasi-cohérent $\mathcal G$, on a
$H^p(\mathfrak U,\mathcal G)=0$ pour $p>0$~; si en outre $U=X$,
l'homomorphisme canonique (\hyperref[III.1.2.3.2-fr]{1.2.3.2})
$M\to H^0(\mathfrak U,\mathcal F)$ est bijectif.
\end{corollary}

\begin{proof}
Comme par hypothèse les $U_i=X_{f_i}$ sont quasi-compacts, il en est de même
de $U$, et on peut donc se borner au cas où $U=X$~; l'hypothèse entraîne alors
$H^p((\mathbf f),M)=0$ pour tout $p\geq0$
(\hyperref[III.1.1.9-fr]{1.1.9}). Le corollaire résulte alors aussitôt de
(\hyperref[III.1.2.3.1-fr]{1.2.3.1}) et
(\hyperref[III.1.2.3.2-fr]{1.2.3.2}).
\end{proof}

On observera que puisque
$H^0(\mathfrak U,\mathcal F)=H^0(U,\mathcal F)$ (G, II, 5.2.2), on a ainsi
démontré à nouveau (I, 1.3.7) comme cas particulier.

\begin{remark}[1.2.5]
\label{III.1.2.5-fr}
Supposons que $X$ soit un \emph{schéma affine}~; alors les
$U_i=X_{f_i}=D(f_i)$ sont des ouverts affines, ainsi que les
$U_{i_0i_1\ldots i_p}$ (mais $U$ n'est pas nécessairement affine). Dans ce
cas, les foncteurs $\Gamma(X,\mathcal F)$ et
$\Gamma(U_{i_0i_1\ldots i_p},\mathcal F)$ sont exacts en $\mathcal F$
(I, 1.3.11). Si on a une suite exacte
$0\to\mathcal F'\to\mathcal F\to\mathcal F''\to0$ de
$\mathcal O_X$-Modules quasi-cohérents, la suite des complexes
\[
0\longrightarrow C^\bullet(\mathfrak U,\mathcal F')
\longrightarrow C^\bullet(\mathfrak U,\mathcal F)
\longrightarrow C^\bullet(\mathfrak U,\mathcal F'')
\longrightarrow0
\]
est exacte, et donne donc une suite exacte de cohomologie
\[
\cdots\longrightarrow H^p(\mathfrak U,\mathcal F')
\longrightarrow H^p(\mathfrak U,\mathcal F)
\longrightarrow H^p(\mathfrak U,\mathcal F'')
\xrightarrow{\partial}H^{p+1}(\mathfrak U,\mathcal F')
\longrightarrow\cdots.
\]
D'autre part, si on pose $M'=\Gamma(X,\mathcal F')$,
$M''=\Gamma(X,\mathcal F'')$, la suite
$0\to M'\to M\to M''\to0$ est exacte~; comme
$C^\bullet((\mathbf f),M)$ est un foncteur exact en $M$, on a aussi la suite
exacte de cohomologie
\[
\cdots\longrightarrow H^p((\mathbf f),M')
\longrightarrow H^p((\mathbf f),M)
\longrightarrow H^p((\mathbf f),M'')
\xrightarrow{\partial}H^{p+1}((\mathbf f),M')
\longrightarrow\cdots.
\]
Cela étant, comme le diagramme
\[
\xymatrix{
0\ar[r]&C^\bullet(\mathfrak U,\mathcal F')\ar[r]\ar[d]
&C^\bullet(\mathfrak U,\mathcal F)\ar[r]\ar[d]
&C^\bullet(\mathfrak U,\mathcal F'')\ar[r]\ar[d]&0\\
0\ar[r]&C^\bullet((\mathbf f),M')\ar[r]
&C^\bullet((\mathbf f),M)\ar[r]
&C^\bullet((\mathbf f),M'')\ar[r]&0
}
\]
\oldpage[III]{88}
est commutatif, on en conclut que les diagrammes
\begin{equation}
\label{III.1.2.5.1-fr}
\tag{1.2.5.1}
\xymatrix{
H^p(\mathfrak U,\mathcal F'')\ar[r]^{\partial}\ar[d]
&H^{p+1}(\mathfrak U,\mathcal F')\ar[d]\\
H^{p+1}((\mathbf f),M'')\ar[r]^{\partial}
&H^{p+2}((\mathbf f),M')
}
\end{equation}
sont commutatifs pour tout $p$ (G, I, 2.1.1).
\end{remark}

\subsection{Cohomologie d'un schéma affine}
\label{subsection:III.1.3-fr}

\begin{theorem}[1.3.1]
\label{III.1.3.1-fr}
Soit $X$ un schéma affine. Pour tout $\mathcal O_X$-Module quasi-cohérent
$\mathcal F$, on a $H^p(X,\mathcal F)=0$ pour tout $p>0$.
\end{theorem}

\begin{proof}
Soit $\mathfrak U$ un recouvrement fini de $X$ par des ouverts affines
$X_{f_i}=D(f_i)$ ($1\leq i\leq r$)~; on sait qu'alors l'idéal engendré dans
$A=\Gamma(X,\mathcal O_X)$ par les $f_i$ est égal à $A$. On conclut donc de
(\hyperref[III.1.2.4-fr]{1.2.4}) que l'on a
$H^p(\mathfrak U,\mathcal F)=0$ pour $p>0$. Comme il y a des recouvrements
finis de $X$ par des ouverts affines qui sont arbitrairement fins (I, 1.1.10),
la définition de la cohomologie de Čech (G, II, 5.8) montre que l'on a aussi
$\check H^p(X,\mathcal F)=0$ pour $p>0$. Mais ceci s'applique aussi à tout
préschéma $X_f$ pour $f\in A$ (I, 1.3.6), donc
$\check H^p(X_f,\mathcal F)=0$ pour $p>0$. Comme l'on a
$X_f\cap X_g=X_{fg}$, on en déduit que l'on a aussi
$H^p(X,\mathcal F)=0$ pour tout $p>0$, en vertu de (G, II, 5.9.2).
\end{proof}

\begin{corollary}[1.3.2]
\label{III.1.3.2-fr}
Soient $Y$ un préschéma, $f:X\to Y$ un morphisme affine (II, 1.6.1). Pour
tout $\mathcal O_X$-Module quasi-cohérent $\mathcal F$, on a
$R^qf_*(\mathcal F)=0$ pour $q>0$.
\end{corollary}

\begin{proof}
En effet, par définition $R^qf_*(\mathcal F)$ est le $\mathcal O_Y$-Module
associé au préfaisceau $U\mapsto H^q(f^{-1}(U),\mathcal F)$, $U$ parcourant
les ouverts de $Y$. Mais les ouverts affines forment une base de $Y$, et pour
un tel ouvert $U$, $f^{-1}(U)$ est affine (II, 1.3.2), donc
$H^q(f^{-1}(U),\mathcal F)=0$ par
(\hyperref[III.1.3.1-fr]{1.3.1}), ce qui démontre le corollaire.
\end{proof}

\begin{corollary}[1.3.3]
\label{III.1.3.3-fr}
Soient $Y$ un préschéma, $f:X\to Y$ un morphisme affine. Pour tout
$\mathcal O_X$-Module quasi-cohérent $\mathcal F$, l'homomorphisme canonique
$H^p(Y,f_*(\mathcal F))\to H^p(X,\mathcal F)$ (0, 12.1.3.1) est bijectif
pour tout $p$.
\end{corollary}

\begin{proof}
En effet, il suffit (en vertu de (0, 12.1.7)) de montrer que les
edge-homomorphismes
$''E_2^{p0}=H^p(Y,f_*(\mathcal F))\to H^p(X,\mathcal F)$ de la seconde suite
spectrale du foncteur composé $\Gamma f_*$ sont bijectifs. Mais le terme
$E_2$ de cette suite est donné par
$''E_2^{pq}=H^p(Y,R^qf_*(\mathcal F))$ (G, II, 4.17.1), donc il résulte de
(\hyperref[III.1.3.2-fr]{1.3.2}) que $''E_2^{pq}=0$ pour $q>0$, et la suite
spectrale dégénère~; d'où notre assertion (0, 11.1.6).
\end{proof}

\begin{corollary}[1.3.4]
\label{III.1.3.4-fr}
Soient $f:X\to Y$ un morphisme affine, $g:Y\to Z$ un morphisme.
\oldpage[III]{89}
Pour tout $\mathcal O_X$-Module quasi-cohérent $\mathcal F$,
l'homomorphisme canonique
$R^pg_*(f_*(\mathcal F))\to R^p(g\circ f)_*(\mathcal F)$ (0, 12.2.5.1)
est bijectif pour tout $p$.
\end{corollary}

\begin{proof}
Il suffit de remarquer que, d'après
(\hyperref[III.1.3.3-fr]{1.3.3}), pour tout ouvert affine $W$ de $Z$,
l'homomorphisme canonique
$H^p(g^{-1}(W),f_*(\mathcal F))\to
H^p(f^{-1}(g^{-1}(W)),\mathcal F)$ est bijectif~; cela prouve que
l'homomorphisme de préfaisceaux définissant l'homomorphisme canonique
$R^pg_*(f_*(\mathcal F))\to R^p(g\circ f)_*(\mathcal F)$ est bijectif
(0, 12.2.5).
\end{proof}

\subsection{Application à la cohomologie des préschémas quelconques}
\label{subsection:III.1.4-fr}

\begin{proposition}[1.4.1]
\label{III.1.4.1-fr}
Soient $X$ un schéma, $\mathfrak U=(U_\alpha)$ un recouvrement de $X$ par
des ouverts affines. Pour tout $\mathcal O_X$-Module quasi-cohérent
$\mathcal F$, les modules de cohomologie $H^\bullet(X,\mathcal F)$ et
$H^\bullet(\mathfrak U,\mathcal F)$ (sur $\Gamma(X,\mathcal O_X)$) sont
canoniquement isomorphes.
\end{proposition}

\begin{proof}
En effet, comme $X$ est un schéma, toute intersection finie $V$ d'ouverts du
recouvrement $\mathfrak U$ est affine (I, 5.5.6), donc
$H^q(V,\mathcal F)=0$ pour $q\geq1$ en vertu de
(\hyperref[III.1.3.1-fr]{1.3.1}). La proposition résulte alors du th. de
Leray (G, II, 5.4.1).
\end{proof}

\begin{remark}[1.4.2]
\label{III.1.4.2-fr}
On notera que la conclusion de (\hyperref[III.1.4.1-fr]{1.4.1}) est encore
valable lorsque les intersections finies des ensembles $U_\alpha$ sont
affines, même lorsque l'on ne suppose pas nécessairement que $X$ soit un
schéma.
\end{remark}

\begin{corollary}[1.4.3]
\label{III.1.4.3-fr}
Soient $X$ un schéma dont l'espace sous-jacent est quasi-compact,
$A=\Gamma(X,\mathcal O_X)$, $\mathbf f=(f_i)_{1\leq i\leq r}$ une suite
finie d'éléments de $A$ tels que les $X_{f_i}$ (notations de
(\hyperref[III.1.2.1-fr]{1.2.1})) soient affines. Alors (avec les notations
de (\hyperref[III.1.2.1-fr]{1.2.1})), pour tout $\mathcal O_X$-Module
quasi-cohérent $\mathcal F$, on a un isomorphisme canonique fonctoriel en
$\mathcal F$
\begin{equation}
\label{III.1.4.3.1-fr}
\tag{1.4.3.1}
H^q(U,\mathcal F)\xrightarrow{\sim}H^{q+1}((\mathbf f),M)
\qquad\text{pour }q\geq1
\end{equation}
et une suite exacte fonctorielle en $\mathcal F$
\begin{equation}
\label{III.1.4.3.2-fr}
\tag{1.4.3.2}
0\longrightarrow H^0((\mathbf f),M)\longrightarrow M
\longrightarrow H^0(U,\mathcal F)
\longrightarrow H^1((\mathbf f),M)\longrightarrow0.
\end{equation}
\end{corollary}

\begin{proof}
Cela résulte en effet aussitôt de
(\hyperref[III.1.4.1-fr]{1.4.1}) et
(\hyperref[III.1.2.3-fr]{1.2.3}).
\end{proof}

\begin{env}[1.4.4]
\label{III.1.4.4-fr}
Si $X$ est un \emph{schéma affine}, il résulte de
(\hyperref[III.1.2.5-fr]{1.2.5}) et
(\hyperref[III.1.4.1-fr]{1.4.1}) que pour tout $q\geq0$, les diagrammes
\begin{equation}
\label{III.1.4.4.1-fr}
\tag{1.4.4.1}
\xymatrix{
H^q(U,\mathcal F'')\ar[r]^{\partial}\ar[d]
&H^{q+1}(U,\mathcal F')\ar[d]\\
H^{q+1}((\mathbf f),M'')\ar[r]^{\partial}
&H^{q+2}((\mathbf f),M')
}
\end{equation}
correspondant à une suite exacte
$0\to\mathcal F'\to\mathcal F\to\mathcal F''\to0$ de
$\mathcal O_X$-Modules quasi-cohérents (avec les notations de
(\hyperref[III.1.2.5-fr]{1.2.5})), sont commutatifs.
\end{env}

\oldpage[III]{90}
\begin{proposition}[1.4.5]
\label{III.1.4.5-fr}
Soient $X$ un schéma quasi-compact, $\mathcal L$ un $\mathcal O_X$-Module
inversible, et considérons l'anneau gradué
$A_\bullet=\Gamma_\bullet(\mathcal L)$ (0\textsubscript{I}, 5.4.6)~; alors
\[
H^\bullet(\mathcal F,\mathcal L)
=\bigoplus_{n\in\mathbb Z}H^\bullet(X,\mathcal F\otimes\mathcal L^{\otimes n})
\]
est un $A_\bullet$-module gradué, et pour tout $f\in A_n$, on a un
isomorphisme canonique
\begin{equation}
\label{III.1.4.5.1-fr}
\tag{1.4.5.1}
H^\bullet(X_f,\mathcal F)
\xrightarrow{\sim}(H^\bullet(\mathcal F,\mathcal L))_{(f)}
\end{equation}
de $(A_\bullet)_{(f)}$-modules.
\end{proposition}

\begin{proof}
Comme $X$ est un schéma quasi-compact, on peut calculer la cohomologie de
tous les $\mathcal O_X$-Modules
$\mathcal F\otimes\mathcal L^{\otimes n}$ à l'aide d'un même recouvrement
fini $\mathfrak U=(U_i)$ par des ouverts affines tels que la restriction
$\mathcal L\vert U_i$ soit isomorphe à $\mathcal O_X\vert U_i$ pour chaque
$i$ (\hyperref[III.1.4.1-fr]{1.4.1}). Il est alors immédiat que les
$U_i\cap X_f$ sont des ouverts affines (I, 1.3.6), et l'on peut donc aussi
calculer la cohomologie
$H^\bullet(X_f,\mathcal F\otimes\mathcal L^{\otimes n})$ à l'aide du
recouvrement $\mathfrak U\vert X_f=(U_i\cap X_f)$
(\hyperref[III.1.4.1-fr]{1.4.1}). Il est immédiat que pour tout $f\in A_n$,
la multiplication par $f$ définit un homomorphisme
\[
C^\bullet(\mathfrak U,\mathcal F\otimes\mathcal L^{\otimes m})
\longrightarrow
C^\bullet(\mathfrak U,\mathcal F\otimes\mathcal L^{\otimes(m+n)}),
\]
d'où un homomorphisme
\[
H^\bullet(\mathfrak U,\mathcal F\otimes\mathcal L^{\otimes m})
\longrightarrow
H^\bullet(\mathfrak U,\mathcal F\otimes\mathcal L^{\otimes(m+n)}),
\]
ce qui établit la première assertion. D'autre part, pour $f\in A_n$ donné,
il résulte de (I, 9.3.2) que l'on a un isomorphisme de complexes de
$(A_\bullet)_{(f)}$-modules
\[
C^\bullet(\mathfrak U\vert X_f,\mathcal F)
\xrightarrow{\sim}
\left(C^\bullet\left(\mathfrak U,
\bigoplus_{n\in\mathbb Z}\mathcal F\otimes\mathcal L^{\otimes n}
\right)\right)_{(f)}
\]
compte tenu de (I, 1.3.9, (ii)). Passant à la cohomologie de ces deux
complexes, on en déduit l'isomorphisme
(\hyperref[III.1.4.5.1-fr]{1.4.5.1}), en se souvenant que le foncteur
$M\mapsto M_{(f)}$ est exact dans la catégorie des $A_\bullet$-modules
gradués.
\end{proof}

\begin{corollary}[1.4.6]
\label{III.1.4.6-fr}
On suppose vérifiées les hypothèses de
(\hyperref[III.1.4.5-fr]{1.4.5}) et on suppose en outre que
$\mathcal L=\mathcal O_X$. Si on pose $A=\Gamma(X,\mathcal O_X)$, alors,
pour tout $f\in A$, on a un isomorphisme canonique
\[
H^\bullet(X_f,\mathcal F)
\xrightarrow{\sim}(H^\bullet(X,\mathcal F))_f
\]
de $A_f$-modules.
\end{corollary}

\begin{corollary}[1.4.7]
\label{III.1.4.7-fr}
Soient $X$ un schéma quasi-compact, $f$ un élément de
$\Gamma(X,\mathcal O_X)$.
\begin{enumerate}
\item[(i)] Supposons l'ouvert $X_f$ affine. Alors, pour tout
$\mathcal O_X$-Module quasi-cohérent $\mathcal F$, tout $i>0$ et tout
$\xi\in H^i(X,\mathcal F)$, il existe un entier $n>0$ tel que $f^n\xi=0$.
\item[(ii)] Inversement, supposons que $X_f$ soit quasi-compact et que pour
tout faisceau quasi-cohérent $\mathcal J$ d'idéaux de $\mathcal O_X$ et tout
$\zeta\in H^1(X,\mathcal J)$, il existe $n>0$ tel que $f^n\zeta=0$. Alors
$X_f$ est affine.
\end{enumerate}
\end{corollary}

\begin{proof}
\begin{enumerate}
\item[(i)] Si $X_f$ est affine, on a $H^i(X_f,\mathcal F)=0$ pour tout
$i>0$ (\hyperref[III.1.3.1-fr]{1.3.1}), donc l'assertion résulte directement
de (\hyperref[III.1.4.6-fr]{1.4.6}).
\item[(ii)] En vertu du critère de Serre (II, 5.2.1), il suffit de prouver
que pour tout Idéal quasi-cohérent $\mathcal K$ de
$\mathcal O_X\vert X_f$, on a $H^1(X_f,\mathcal K)=0$. Comme $X_f$ est un
ouvert quasi-compact dans un schéma quasi-compact $X$, il existe un Idéal
quasi-cohérent $\mathcal J$ de $\mathcal O_X$ tel que
$\mathcal K=\mathcal J\vert X_f$ (I, 9.4.2). En vertu de
(\hyperref[III.1.4.6-fr]{1.4.6}), on a
$H^1(X_f,\mathcal K)=(H^1(X,\mathcal J))_f$, et l'hypothèse entraîne que le
second membre est nul, d'où la conclusion.
\end{enumerate}
\end{proof}

\begin{remark}[1.4.8]
\label{III.1.4.8-fr}
On notera que (\hyperref[III.1.4.7-fr]{1.4.7}, (i)) redonne une
démonstration plus simple de la relation (II, 4.5.13.2).
\end{remark}

\oldpage[III]{91}
\begin{lemma}[1.4.9]
\label{III.1.4.9-fr}
Soient $X$ un schéma quasi-compact,
$\mathfrak U=(U_i)_{1\leq i\leq n}$ un recouvrement fini de $X$ par des
ouverts affines, $\mathcal F$ un $\mathcal O_X$-Module quasi-cohérent. Le
complexe de faisceaux $\mathcal C^\bullet(\mathfrak U,\mathcal F)$ défini
par le recouvrement $\mathfrak U$ (G, II, 5.2) est alors un
$\mathcal O_X$-Module quasi-cohérent.
\end{lemma}

\begin{proof}
Il résulte des définitions (G, II, 5.2) que
$\mathcal C^p(\mathfrak U,\mathcal F)$ est somme directe des faisceaux
images directes des $\mathcal F\vert U_{i_0\ldots i_p}$ par l'injection
canonique $U_{i_0\ldots i_p}\to X$. Or, l'hypothèse que $X$ est un schéma
entraîne que ces injections sont des morphismes affines (I, 5.5.6), donc les
$\mathcal C^p(\mathfrak U,\mathcal F)$ sont quasi-cohérents (II, 1.2.6).
\end{proof}

\begin{proposition}[1.4.10]
\label{III.1.4.10-fr}
Soit $u:X\to Y$ un morphisme séparé et quasi-compact. Pour tout
$\mathcal O_X$-Module quasi-cohérent $\mathcal F$, les
$R^qu_*(\mathcal F)$ sont des $\mathcal O_Y$-Modules quasi-cohérents.
\end{proposition}

\begin{proof}
La question étant locale sur $Y$, on peut supposer $Y$ affine. Alors $X$ est
réunion finie d'ouverts affines $U_i$ ($1\leq i\leq n$)~; soit
$\mathfrak U$ le recouvrement $(U_i)$. En outre, comme $Y$ est un schéma,
il résulte de (I, 5.5.10) que pour tout ouvert affine $V\subset Y$,
l'injection canonique $u^{-1}(V)\to X$ est un morphisme affine~; on en
conclut ((\hyperref[III.1.4.1-fr]{1.4.1}) et (G, II, 5.2)) que l'on a un
isomorphisme canonique
\begin{equation}
\label{III.1.4.10.1-fr}
\tag{1.4.10.1}
H^\bullet(u^{-1}(V),\mathcal F)
\xrightarrow{\sim}H^\bullet(\Gamma(V,\mathcal K^\bullet))
\end{equation}
où l'on a posé
$\mathcal K^\bullet=u_*(\mathcal C^\bullet(\mathfrak U,\mathcal F))$. En
vertu de (\hyperref[III.1.4.9-fr]{1.4.9}) et (I, 9.2.2),
$\mathcal K^\bullet$ est un $\mathcal O_Y$-Module quasi-cohérent~; par
ailleurs, il constitue un \emph{complexe de faisceaux} puisqu'il en est
ainsi de $\mathcal C^\bullet(\mathfrak U,\mathcal F)$. Il résulte alors de la
définition de la cohomologie
$\mathcal H^\bullet(\mathcal K^\bullet)$ (G, II, 4.1) que cette dernière est
formée de $\mathcal O_Y$-Modules quasi-cohérents (I, 4.1.1). Comme (pour $V$
affine dans $Y$) le foncteur $\Gamma(V,\mathcal G)$ est exact en
$\mathcal G$ dans la catégorie des $\mathcal O_Y$-Modules quasi-cohérents,
on a (G, II, 4.1)
\begin{equation}
\label{III.1.4.10.2-fr}
\tag{1.4.10.2}
H^\bullet(\Gamma(V,\mathcal K^\bullet))
=\Gamma(V,\mathcal H^\bullet(\mathcal K^\bullet)).
\end{equation}
Notons enfin qu'il résulte de la définition de l'homomorphisme canonique
\[
H^\bullet(\mathfrak U,\mathcal F)\longrightarrow H^\bullet(X,\mathcal F)
\]
donnée dans (G, II, 5.2), que si $V'\subset V$ est un second ouvert affine
de $Y$, le diagramme
\[
\xymatrix{
H^\bullet(u^{-1}(V),\mathcal F)\ar[r]^{\sim}\ar[d]
&H^\bullet(\Gamma(V,\mathcal K^\bullet))\ar[d]\\
H^\bullet(u^{-1}(V'),\mathcal F)\ar[r]^{\sim}
&H^\bullet(\Gamma(V',\mathcal K^\bullet))
}
\]
est commutatif. On conclut donc de ce qui précède que les isomorphismes
(\hyperref[III.1.4.10.1-fr]{1.4.10.1}) définissent un isomorphisme de
$\mathcal O_Y$-Modules
\begin{equation}
\label{III.1.4.10.3-fr}
\tag{1.4.10.3}
R^\bullet u_*(\mathcal F)
\xrightarrow{\sim}\mathcal H^\bullet(\mathcal K^\bullet)
\end{equation}
\oldpage[III]{92}
et par suite $R^\bullet u_*(\mathcal F)$ est quasi-cohérent.
\end{proof}

Il résulte en outre de (\hyperref[III.1.4.10.3-fr]{1.4.10.3}),
(\hyperref[III.1.4.10.2-fr]{1.4.10.2}) et
(\hyperref[III.1.4.10.1-fr]{1.4.10.1}) que~:

\begin{corollary}[1.4.11]
\label{III.1.4.11-fr}
Sous les hypothèses de (\hyperref[III.1.4.10-fr]{1.4.10}), pour tout ouvert
affine $V$ de $Y$, l'homomorphisme canonique
\begin{equation}
\label{III.1.4.11.1-fr}
\tag{1.4.11.1}
H^q(u^{-1}(V),\mathcal F)\longrightarrow
\Gamma(V,R^qu_*(\mathcal F))
\end{equation}
est un isomorphisme pour tout $q\geq0$.
\end{corollary}

\begin{corollary}[1.4.12]
\label{III.1.4.12-fr}
Supposons vérifiées les hypothèses de
(\hyperref[III.1.4.10-fr]{1.4.10}) et supposons en outre que $Y$ soit
quasi-compact. Alors il existe un entier $r>0$ tel que pour tout
$\mathcal O_X$-Module quasi-cohérent $\mathcal F$ et tout entier $q>r$, on
ait $R^qu_*(\mathcal F)=0$. Si $Y$ est affine, on peut prendre pour $r$ un
entier tel qu'il existe un recouvrement de $X$ formé de $r$ ouverts affines.
\end{corollary}

\begin{proof}
Comme on peut recouvrir $Y$ par un nombre fini d'ouverts affines, on est
ramené à démontrer la seconde assertion, en vertu de
(\hyperref[III.1.4.11-fr]{1.4.11}). Or, si $\mathfrak U$ est un
recouvrement de $X$ par $r$ ouverts affines, on a
$H^q(\mathfrak U,\mathcal F)=0$ pour $q>r$, puisque les cochaînes de
$C^q(\mathfrak U,\mathcal F)$ sont alternées~; la conclusion résulte donc de
(\hyperref[III.1.4.1-fr]{1.4.1}).
\end{proof}

\begin{corollary}[1.4.13]
\label{III.1.4.13-fr}
Supposons vérifiées les hypothèses de
(\hyperref[III.1.4.10-fr]{1.4.10}) et en outre supposons que
$Y=\operatorname{Spec}(A)$ soit affine. Alors, pour tout
$\mathcal O_X$-Module quasi-cohérent $\mathcal F$ et tout $f\in A$, on a
\[
\Gamma(Y_f,R^qu_*(\mathcal F))
=(\Gamma(Y,R^qu_*(\mathcal F)))_f
\]
à un isomorphisme canonique près.
\end{corollary}

\begin{proof}
En effet, cela résulte de ce que $R^qu_*(\mathcal F)$ est un
$\mathcal O_Y$-Module quasi-cohérent (I, 1.3.7).
\end{proof}

\begin{proposition}[1.4.14]
\label{III.1.4.14-fr}
Soient $f:X\to Y$ un morphisme séparé quasi-compact, $g:Y\to Z$ un
morphisme affine. Pour tout $\mathcal O_X$-Module quasi-cohérent
$\mathcal F$, l'homomorphisme canonique
$R^p(g\circ f)_*(\mathcal F)\to g_*(R^pf_*(\mathcal F))$ (0, 12.2.5.2) est
bijectif pour tout $p$.
\end{proposition}

\begin{proof}
En effet, pour tout ouvert affine $W$ de $Z$, $g^{-1}(W)$ est un ouvert
affine de $Y$. L'homomorphisme de préfaisceaux définissant l'homomorphisme
canonique
\[
R^p(g\circ f)_*(\mathcal F)\longrightarrow g_*(R^pf_*(\mathcal F))
\]
(0, 12.2.5) est donc bijectif en vertu de
(\hyperref[III.1.4.11-fr]{1.4.11}).
\end{proof}

\begin{proposition}[1.4.15]
\label{III.1.4.15-fr}
Soient $u:X\to Y$ un morphisme séparé de type fini, $v:Y'\to Y$ un
morphisme plat de préschémas (0\textsubscript{I}, 6.7.1)~; soit
$u'=u_{(Y')}$, de sorte qu'on a le diagramme commutatif
\begin{equation}
\label{III.1.4.15.1-fr}
\tag{1.4.15.1}
\xymatrix{
X\ar[d]_u & X'=X_{(Y')}\ar[l]_{v'}\ar[d]^{u'}\\
Y & Y'\ar[l]_v
}
\end{equation}
Alors, pour tout $\mathcal O_X$-Module quasi-cohérent $\mathcal F$,
$R^qu_*'(\mathcal F')$, où
$\mathcal F'=v'^*(\mathcal F)=\mathcal F\otimes_{\mathcal O_Y}\mathcal O_{Y'}$,
est canoniquement isomorphe à
$R^qu_*(\mathcal F)\otimes_{\mathcal O_Y}\mathcal O_{Y'}
=v^*(R^qu_*(\mathcal F))$ pour tout $q\geq0$.
\end{proposition}

\oldpage[III]{93}
\begin{proof}
L'homomorphisme canonique
$\rho:\mathcal F\to v_*'(v'^*(\mathcal F))$ (0\textsubscript{I}, 4.4.3.2)
définit par fonctorialité un homomorphisme
\begin{equation}
\label{III.1.4.15.2-fr}
\tag{1.4.15.2}
R^qu_*(\mathcal F)\longrightarrow R^qu_*(v_*'(\mathcal F')).
\end{equation}
D'autre part, on a, en posant $w=u\circ v'=v\circ u'$, les homomorphismes
canoniques (0, 12.2.5.1 et 12.2.5.2)
\begin{equation}
\label{III.1.4.15.3-fr}
\tag{1.4.15.3}
R^qu_*(v_*'(\mathcal F'))\longrightarrow R^qw_*(\mathcal F')
\longrightarrow v_*(R^qu_*'(\mathcal F')).
\end{equation}
Composant (\hyperref[III.1.4.15.3-fr]{1.4.15.3}) et
(\hyperref[III.1.4.15.2-fr]{1.4.15.2}), on a un homomorphisme
\[
\psi:R^qu_*(\mathcal F)\longrightarrow v_*(R^qu_*'(\mathcal F'))
\]
et finalement on en déduit l'homomorphisme canonique (dont la définition ne
fait intervenir \emph{aucune hypothèse sur $v$})
\begin{equation}
\label{III.1.4.15.4-fr}
\tag{1.4.15.4}
\psi^\sharp:v^*(R^qu_*(\mathcal F))\longrightarrow R^qu_*'(\mathcal F')
\end{equation}
dont il s'agit de prouver que c'est un isomorphisme lorsque $v$ est
\emph{plat}. Il est clair que la question est locale sur $Y$ et $Y'$ et l'on
peut donc supposer que $Y=\operatorname{Spec}(A)$,
$Y'=\operatorname{Spec}(B)$~; nous utiliserons en outre le

\begin{lemma}[1.4.15.5]
\label{III.1.4.15.5-fr}
Soient $\varphi:A\to B$ un homomorphisme d'anneaux,
$Y=\operatorname{Spec}(A)$, $X=\operatorname{Spec}(B)$,
$f:X\to Y$ le morphisme correspondant à $\varphi$, $M$ un $B$-module. Pour
que le $\mathcal O_X$-Module $\widetilde M$ soit $f$-plat
(0\textsubscript{I}, 6.7.1), il faut et il suffit que $M$ soit un $A$-module
plat. En particulier, pour que le morphisme $f$ soit plat, il faut et il
suffit que $B$ soit un $A$-module plat.
\end{lemma}

En effet, cela résulte de la définition (0\textsubscript{I}, 6.7.1) et de
(0\textsubscript{I}, 6.3.3), compte tenu de (I, 1.3.4).

Cela étant, il résulte de
(\hyperref[III.1.4.11.1-fr]{1.4.11.1}) et des définitions des
homomorphismes (\hyperref[III.1.4.15.3-fr]{1.4.15.3})
(cf. 0, 12.2.5) que $\psi$ correspond alors au morphisme composé
\[
H^q(X,\mathcal F)\xrightarrow{\rho_q}
H^q(X,v_*'(v'^*(\mathcal F)))\xrightarrow{\theta_q}
H^q(X',v'^*(v_*'(v'^*(\mathcal F))))\xrightarrow{\sigma_q}
H^q(X',v'^*(\mathcal F))
\]
où $\rho_q$ et $\sigma_q$ sont les homomorphismes correspondant dans la
cohomologie aux morphismes canoniques $\rho$ et
$\sigma:v'^*(v_*'(\mathcal G'))\to\mathcal G'$, et $\theta_q$ est le
$\varphi$-morphisme (0, 12.1.3.1) relatif au $\mathcal O_X$-Module
$v_*'(v'^*(\mathcal F))$. Mais par fonctorialité de $\theta_q$, on a le
diagramme commutatif
\[
\xymatrix{
H^q(X,\mathcal F)\ar[rr]^{\rho_q}\ar[dd]_{\theta_q}
&&H^q(X,v_*'(v'^*(\mathcal F)))\ar[dd]^{\theta_q}\\\\
H^q(X',v'^*(\mathcal F))\ar[rr]^{v'^*(\rho_q)}
&&H^q(X',v'^*(v_*'(v'^*(\mathcal F))))
}
\]
\oldpage[III]{94}
et comme par définition (0\textsubscript{I}, 4.4.3), $v'^*(\rho)$ est
l'inverse de $\sigma$, on voit que le morphisme composé considéré plus haut
n'est autre finalement que $\theta_q$~; $\psi^\sharp$ est par suite le
$B$-homomorphisme associé
\[
H^q(X,\mathcal F)\otimes_A B\longrightarrow H^q(X',\mathcal F').
\]
Comme $u$ est de type fini, $X$ est réunion finie d'ouverts affines $U_i$
($1\leq i\leq r$)~; soit $\mathfrak U$ le recouvrement $(U_i)$. D'ailleurs
$v$ est un morphisme affine, donc il en est de même de $v'$ (II, 1.6.2,
(iii)), et les $U_i'=v'^{-1}(U_i)$ forment par suite un recouvrement ouvert
affine $\mathfrak U'$ de $X'$. On sait alors (0, 12.1.4.2) que le diagramme
\[
\xymatrix{
H^q(\mathfrak U,\mathcal F)\ar[r]^{\theta_q}\ar[d]
&H^q(\mathfrak U',\mathcal F')\ar[d]\\
H^q(X,\mathcal F)\ar[r]^{\theta_q}&H^q(X',\mathcal F')
}
\]
est commutatif, et les flèches verticales sont des isomorphismes puisque
$X$ et $X'$ sont des schémas (\hyperref[III.1.4.1-fr]{1.4.1}). Il suffit par
suite de prouver que le $\varphi$-morphisme canonique
$\theta_q:H^q(\mathfrak U,\mathcal F)\to H^q(\mathfrak U',\mathcal F')$
est tel que le $B$-homomorphisme associé
\[
H^q(\mathfrak U,\mathcal F)\otimes_A B
\longrightarrow H^q(\mathfrak U',\mathcal F')
\]
soit un isomorphisme. Pour toute suite
$\mathbf s=(i_k)_{0\leq k\leq p}$ de $p+1$ indices de $[1,r]$, posons
\[
U_{\mathbf s}=\bigcap_{k=0}^pU_{i_k},\qquad
U_{\mathbf s}'=\bigcap_{k=0}^pU_{i_k}'=v'^{-1}(U_{\mathbf s}),\qquad
M_{\mathbf s}=\Gamma(U_{\mathbf s},\mathcal F),\qquad
M_{\mathbf s}'=\Gamma(U_{\mathbf s}',\mathcal F').
\]
L'application canonique
$M_{\mathbf s}\otimes_A B\to M_{\mathbf s}'$ est un isomorphisme
(I, 1.6.5), donc l'application canonique
$C^p(\mathfrak U,\mathcal F)\otimes_A B\to
C^p(\mathfrak U',\mathcal F')$ est un isomorphisme, par lequel
$d\otimes1$ s'identifie à l'opérateur cobord
$C^p(\mathfrak U',\mathcal F')\to C^{p+1}(\mathfrak U',\mathcal F')$.
Comme $B$ est un $A$-module \emph{plat}, il résulte aussitôt de la définition
des modules de cohomologie que l'application canonique
$H^q(\mathfrak U,\mathcal F)\otimes_A B\to
H^q(\mathfrak U',\mathcal F')$ est un isomorphisme
(0\textsubscript{I}, 6.1.1). Ce résultat sera généralisé plus tard (\S~6).
\end{proof}

\begin{corollary}[1.4.16]
\label{III.1.4.16-fr}
Soient $A$ un anneau, $X$ un $A$-schéma de type fini, $B$ une $A$-algèbre
fidèlement plate sur $A$. Pour que $X$ soit affine, il faut et il suffit que
$X\otimes_A B$ le soit.
\end{corollary}

\begin{proof}
La condition est évidemment nécessaire (I, 3.2.2)~; montrons qu'elle est
suffisante. Comme $X$ est séparé sur $A$ et que le morphisme
$\operatorname{Spec}(B)\to\operatorname{Spec}(A)$ est plat, il résulte de
(1.4.1) que l'on a
\begin{equation}
\label{III.1.4.16.1-fr}
\tag{1.4.16.1}
H^i(X\otimes_A B,\mathcal F\otimes_A B)
=H^i(X,\mathcal F)\otimes_A B
\end{equation}
pour tout $i\geq0$ et tout $\mathcal O_X$-Module quasi-cohérent $\mathcal F$.
Si $X\otimes_A B$ est affine, le premier membre de
(\hyperref[III.1.4.16.1-fr]{1.4.16.1}) est nul pour $i=1$, donc il en est de
même de $H^1(X,\mathcal F)$ puisque $B$ est un $A$-module fidèlement plat.
Comme $X$ est un schéma quasi-compact, on conclut par le critère de Serre
(II, 5.2.1).
\end{proof}

\begin{proposition}[1.4.17]
\label{III.1.4.17-fr}
Soient $X$ un préschéma,
$0\to\mathcal F\xrightarrow{u}\mathcal G\xrightarrow{v}\mathcal H\to0$
une suite exacte de $\mathcal O_X$-Modules. Si $\mathcal F$ et $\mathcal H$
sont quasi-cohérents, il en est de même de $\mathcal G$.
\end{proposition}

\oldpage[III]{95}
\begin{proof}
La question étant locale sur $X$, on peut supposer
$X=\operatorname{Spec}(A)$ affine et il suffit alors de prouver que
$\mathcal G$ vérifie les conditions d~1) et d~2) de (I, 1.4.1) (avec
$V=X$). La vérification de d~2) est immédiate, car si
$t\in\Gamma(X,\mathcal G)$ a une restriction nulle à $D(f)$, il en est de
même de son image $v(t)\in\Gamma(X,\mathcal H)$~; il y a donc $m>0$ tel que
$f^mv(t)=v(f^mt)=0$ (I, 1.4.1), et comme $\Gamma$ est exact à gauche,
$f^mt=u(s)$, où $s\in\Gamma(X,\mathcal F)$~; comme $u$ est injectif, la
restriction de $s$ à $D(f)$ est nulle, d'où (I, 1.4.1) l'existence d'un
entier $n>0$ tel que $f^ns=0$~; on en déduit finalement
$f^{m+n}t=u(f^ns)=0$.

Vérifions maintenant d~2), et soit $t'\in\Gamma(D(f),\mathcal G)$~; comme
$\mathcal H$ est quasi-cohérent, il existe un entier $m$ tel que
$f^mv(t')=v(f^mt')$ se prolonge en une section
$z\in\Gamma(X,\mathcal H)$ (I, 1.4.1). Mais en vertu de
(\hyperref[III.1.3.1-fr]{1.3.1}) (ou de (I, 5.1.9.2)) appliqué au
$\mathcal O_X$-Module quasi-cohérent $\mathcal F$, la suite
$\Gamma(X,\mathcal G)\to\Gamma(X,\mathcal H)\to0$ est exacte, donc il existe
$t\in\Gamma(X,\mathcal G)$ tel que $z=v(t)$~; on voit donc que
$v(f^mt'-t'')=0$, en désignant par $t''$ la restriction de $t$ à $D(f)$~;
on a donc $f^mt'-t''=u(s')$, où $s'\in\Gamma(D(f),\mathcal F)$. Mais comme
$\mathcal F$ est quasi-cohérent, il existe un entier $n>0$ tel que $f^ns'$
se prolonge en une section $s\in\Gamma(X,\mathcal F)$~; comme
$f^{m+n}t'-f^nt''=u(f^ns')$, on voit que $f^{m+n}t'$ est la restriction à
$D(f)$ de la section $f^nt+u(f^ns)\in\Gamma(X,\mathcal G)$, ce qui achève
la démonstration.
\end{proof}

\section{Étude cohomologique des morphismes projectifs}
\label{section:III.2-fr}

\subsection{Calculs explicites de certains groupes de cohomologie}
\label{subsection:III.2.1-fr}

\begin{env}[2.1.1]
\label{III.2.1.1-fr}
Soient $X$ un préschéma, $\mathcal L$ un $\mathcal O_X$-Module inversible~;
considérons l'anneau gradué (0\textsubscript{I}, 5.4.6)
\begin{equation}
\label{III.2.1.1.1-fr}
\tag{2.1.1.1}
S=\Gamma_\bullet(X,\mathcal L)
=\bigoplus_{n\in\mathbb Z}\Gamma(X,\mathcal L^{\otimes n}).
\end{equation}
Soit $(f_i)_{1\leq i\leq r}$ une famille finie d'éléments \emph{homogènes}
de $S$, soit $f_i\in S_{d_i}$~; posons
$U_i=X_{f_i}$, $U=\bigcup_iU_i$, et désignons par $\mathfrak U$ le
recouvrement $(U_i)$ de $U$. Pour tout $\mathcal O_X$-Module quasi-cohérent
$\mathcal F$, on posera
\begin{align}
\label{III.2.1.1.2-fr}
\tag{2.1.1.2}
H^\bullet(U,\mathcal F(*))
&=\bigoplus_{n\in\mathbb Z}
H^\bullet(U,\mathcal F\otimes\mathcal L^{\otimes n}),\\
\label{III.2.1.1.3-fr}
\tag{2.1.1.3}
H^\bullet(\mathfrak U,\mathcal F(*))
&=\bigoplus_{n\in\mathbb Z}
H^\bullet(\mathfrak U,\mathcal F\otimes\mathcal L^{\otimes n}).
\end{align}
Les groupes abéliens (\hyperref[III.2.1.1.2-fr]{2.1.1.2}) et
(\hyperref[III.2.1.1.3-fr]{2.1.1.3}) sont \emph{bigradués}, en prenant
\[
(H^\bullet(U,\mathcal F(*)))_{mn}
=H^m(U,\mathcal F\otimes\mathcal L^{\otimes n})
\]
et une définition analogue pour
(\hyperref[III.2.1.1.3-fr]{2.1.1.3}). Pour le deuxième degré, il est clair
que ces groupes sont des $S$-modules gradués, comme il résulte par exemple du
fait que $\mathcal F\mapsto H^m(U,\mathcal F)$ et
$\mathcal F\mapsto H^m(\mathfrak U,\mathcal F)$ sont des foncteurs.
\end{env}

\begin{env}[2.1.2]
\label{III.2.1.2-fr}
Considérons maintenant le $S$-module gradué (0\textsubscript{I}, 5.4.6)
\begin{equation}
\label{III.2.1.2.1-fr}
\tag{2.1.2.1}
M=\Gamma_\bullet(\mathcal L,\mathcal F)
=H^0(X,\mathcal F(*))
=\bigoplus_{n\in\mathbb Z}
\Gamma(X,\mathcal F\otimes\mathcal L^{\otimes n}).
\end{equation}
\oldpage[III]{96}
Si $X$ est un préschéma dont l'espace sous-jacent est noethérien, ou un
schéma quasi-compact, il résulte de (I, 9.3.1) qu'en posant comme d'ordinaire
$U_{i_0i_1\ldots i_p}=\bigcap_{k=0}^pU_{i_k}$, on a, à un isomorphisme
canonique près,
\[
\Gamma(U_{i_0i_1\ldots i_p},\mathcal F(*))
=H^0(U_{i_0i_1\ldots i_p},\mathcal F(*))
=M_{f_{i_0}f_{i_1}\cdots f_{i_p}}.
\]
On peut encore, avec les notations de
(\hyperref[III.1.2.2-fr]{1.2.2}), identifier
$M_{f_{i_0}f_{i_1}\cdots f_{i_p}}$ à
$\varinjlim_n M_{i_0i_1\ldots i_p}^{(n)}$. Cette identification est un
isomorphisme de $S$-modules gradués, si l'on définit le degré d'un élément
homogène $z\in\varinjlim_n M_{i_0i_1\ldots i_p}^{(n)}$ de la façon
suivante~: $z$ est l'image canonique d'un élément homogène
$x\in M_{i_0i_1\ldots i_p}^{(n)}=M$, de degré $m$~; on prend alors pour
degré de $z$ le nombre
$m-n(d_{i_0}+d_{i_1}+\cdots+d_{i_p})$. Tenant compte de la définition des
homomorphismes
$\varphi_{kn}:M_{i_0i_1\ldots i_p}^{(n)}\to
M_{i_0i_1\ldots i_p}^{(k)}$ (\hyperref[III.1.2.2-fr]{1.2.2}), on voit
aussitôt que cette définition ne dépend pas du «~représentant~» $x$ de $z$
que l'on a considéré. Désignant comme dans
(\hyperref[III.1.2.2-fr]{1.2.2}) par $C_n^p(M)$ l'ensemble des applications
alternées de $[1,r]^{p+1}$ dans $M$ (pour tout $n$), on définit de la même
manière que ci-dessus une structure de $S$-module gradué sur
$\varinjlim_n C_n^p(M)$~; on a encore comme dans
(\hyperref[III.1.2.2-fr]{1.2.2})
\begin{equation}
\label{III.2.1.2.2-fr}
\tag{2.1.2.2}
C^p(\mathfrak U,\mathcal F(*))=\varinjlim_n C_n^p(M)
\end{equation}
l'isomorphisme des deux membres \emph{respectant les degrés}. On a alors,
comme dans (\hyperref[III.1.2.2-fr]{1.2.2})
\begin{equation}
\label{III.2.1.2.3-fr}
\tag{2.1.2.3}
C^p(\mathfrak U,\mathcal F(*))
=C^{p+1}((\mathbf f),M)
=\varinjlim_n K^{p+1}(\mathbf f^n,M)
\end{equation}
l'isomorphisme conservant les degrés~: le degré d'un élément de
$\varinjlim_n K^{p+1}(\mathbf f^n,M)$, image canonique d'une cochaîne
$g\in K^{p+1}(\mathbf f^n,M)$ dont les valeurs
$g(i_0,\ldots,i_p)$ sont dans une même composante homogène $M_m$ de $M$,
est $m-n(d_{i_0}+\cdots+d_{i_p})$, et il est indépendant du choix de cette
cochaîne comme représentant de l'élément considéré.
\end{env}

Comme les isomorphismes précédents sont compatibles avec les opérateurs
cobords, on en conclut, comme dans
(\hyperref[III.1.2.2-fr]{1.2.2}) que l'on a~:

\begin{proposition}[2.1.3]
\label{III.2.1.3-fr}
Soit $X$ un préschéma dont l'espace sous-jacent est noethérien, ou un schéma
quasi-compact. Il existe un isomorphisme canonique fonctoriel en $\mathcal F$
\begin{equation}
\label{III.2.1.3.1-fr}
\tag{2.1.3.1}
H^p(\mathfrak U,\mathcal F(*))
\xrightarrow{\sim}H^{p+1}((\mathbf f),M)
\qquad\text{pour tout }p\geq1.
\end{equation}
On a en outre une suite exacte fonctorielle en $\mathcal F$
\begin{equation}
\label{III.2.1.3.2-fr}
\tag{2.1.3.2}
0\longrightarrow H^0((\mathbf f),M)\longrightarrow M
\longrightarrow H^0(\mathfrak U,\mathcal F(*))
\longrightarrow H^1((\mathbf f),M)\longrightarrow0.
\end{equation}
De plus, tous les homomorphismes introduits sont de degré $0$ pour les
structures de $S$-modules gradués ($S$ étant l'anneau
(\hyperref[III.2.1.1.1-fr]{2.1.1.1})).
\end{proposition}
\providecommand{\bbGamma}{\boldsymbol{\Gamma}}
\providecommand{\shProj}{\mathcal{P}\mathrm{roj}}

\oldpage[III]{97}
\begin{corollary}[2.1.4]
\label{III.2.1.4-fr}
Si $X$ est un schéma quasi-compact et si les $U_i=X_{f_i}$ sont affines,
il existe un isomorphisme canonique fonctoriel en $\mathcal F$, de degré $0$,
\begin{equation}
\label{III.2.1.4.1-fr}
\tag{2.1.4.1}
H^p(U,\mathcal F(*))\xrightarrow{\sim}H^{p+1}((\mathbf f),M)
\qquad\text{pour }p\geq1,
\end{equation}
et une suite exacte fonctorielle en $\mathcal F$
\begin{equation}
\label{III.2.1.4.2-fr}
\tag{2.1.4.2}
0\longrightarrow H^0((\mathbf f),M)\longrightarrow M
\longrightarrow H^0(U,\mathcal F(*))
\longrightarrow H^1((\mathbf f),M)\longrightarrow0,
\end{equation}
où tous les homomorphismes sont de degré $0$.
\end{corollary}

\begin{proof}
Il suffit en effet d'appliquer (1.4.1) au résultat de
(\hyperref[III.2.1.3-fr]{2.1.3}).
\end{proof}

La proposition «~locale~» analogue à
(\hyperref[III.2.1.3-fr]{2.1.3}) est la suivante~:

\begin{proposition}[2.1.5]
\label{III.2.1.5-fr}
Soient $S$ un anneau gradué à degrés positifs,
$f_i$ $(1\leq i\leq r)$ un élément homogène de $S_+$ de degré $d_i$,
$M$ un $S$-module gradué. Soit $X=\operatorname{Proj}(S)$ le spectre
premier homogène de $S$ et posons $U_i=D_+(f_i)$,
$U=\bigcup_iU_i$,
\[
H^\bullet(U,\widetilde M(*))
=\bigoplus_{n\in\mathbb Z}H^\bullet(U,(M(n))^\sim).
\]
Il existe alors des isomorphismes canoniques fonctoriels en $M$, de degré
$0$ pour les structures de $S$-module gradué
\begin{equation}
\label{III.2.1.5.1-fr}
\tag{2.1.5.1}
H^p(U,\widetilde M(*))\xrightarrow{\sim}H^{p+1}((\mathbf f),M)
\qquad\text{pour }p\geq1,
\end{equation}
et une suite exacte fonctorielle en $M$
\begin{equation}
\label{III.2.1.5.2-fr}
\tag{2.1.5.2}
0\longrightarrow H^0((\mathbf f),M)\longrightarrow M
\longrightarrow H^0(U,\widetilde M(*))
\longrightarrow H^1((\mathbf f),M)\longrightarrow0,
\end{equation}
où tous les homomorphismes sont de degré $0$.
\end{proposition}

\begin{proof}
En effet, on a
\[
\Gamma(U_{i_0i_1\ldots i_p},(M(n))^\sim)
=(M_{f_{i_0}\cdots f_{i_p}})_n
\]
par définition (II, 2.5.2), donc
$\Gamma(U_{i_0i_1\ldots i_p},\widetilde M(*))
=M_{f_{i_0}\cdots f_{i_p}}$. Le reste du raisonnement est alors le même
que pour démontrer (\hyperref[III.2.1.4-fr]{2.1.4}), tenant compte de ce
que $X$ est un schéma.
\end{proof}

\begin{env}[2.1.6]
\label{III.2.1.6-fr}
\emph{Remarques.} ---
(i) Sous les conditions de (\hyperref[III.2.1.5-fr]{2.1.5}), les foncteurs
$\Gamma(U_{i_0i_1\ldots i_p},\widetilde M(*))$ sont exacts en $M$, en
vertu de (0\textsubscript{I}, 1.3.2)~; le même raisonnement que dans
(1.2.5) montre alors que si
$0\to M'\to M\to M''\to0$ est une suite exacte de $S$-modules gradués
(où les homomorphismes sont de degré $0$), on a des diagrammes
commutatifs pour tout $p\geq0$
\begin{equation}
\label{III.2.1.6.1-fr}
\tag{2.1.6.1}
\xymatrix{
H^p(U,\widetilde {M''}(*))\ar[r]^-\partial\ar[d] &
H^{p+1}(U,\widetilde {M'}(*))\ar[d]\\
H^{p+1}((\mathbf f),M'')\ar[r]^-\partial &
H^{p+2}((\mathbf f),M').
}
\end{equation}

(ii) La prop. (\hyperref[III.2.1.5-fr]{2.1.5}) sera surtout intéressante
lorsque $S$ sera une $A$-algèbre engendrée par un nombre fini d'éléments de
degré $1$, $A$ étant supposé noethérien~; en
\oldpage[III]{98}
effet, lorsqu'il en est ainsi, tout $\mathcal O_X$-Module quasi-cohérent est
de la forme $\widetilde M$ (II, 2.7.7).
\end{env}

\begin{env}[2.1.7]
\label{III.2.1.7-fr}
Nous allons appliquer (\hyperref[III.2.1.5-fr]{2.1.5}) dans le cas
$S=A[T_0,\ldots,T_r]$, où $A$ est un anneau quelconque, les $T_i$ sont des
indéterminées, avec $M=S$ et $f_i=T_i$. On est donc essentiellement ramené
à calculer $H^\bullet((\mathbf T),S)$, où
$\mathbf T=(T_i)_{0\leq i\leq r}$.
\end{env}

\begin{lemma}[2.1.8]
\label{III.2.1.8-fr}
Si $S=A[T_0,\ldots,T_r]$, on a, avec
$\mathbf T=(T_i)_{0\leq i\leq r}$,
\begin{align}
\label{III.2.1.8.1-fr}
\tag{2.1.8.1}
H^i(\mathbf T^n,S)&=0 &&\text{si }i\neq r+1,\\
\label{III.2.1.8.2-fr}
\tag{2.1.8.2}
H^{r+1}(\mathbf T^n,S)&=S/(\mathbf T^n).
\end{align}
Le $A$-module $H^{r+1}(\mathbf T^n,S)$ a donc une base sur $A$ formée des
classes mod. $(\mathbf T^n)$ des monômes
\[
T_{\mathbf p}=T_0^{p_0}\cdots T_r^{p_r},
\qquad \mathbf p=(p_0,\ldots,p_r),\quad 0\leq p_i<n
\quad\text{pour tout }i.
\]
\end{lemma}

\begin{proof}
C'est une conséquence immédiate de (1.1.3.5) et de la prop. (1.1.4), dont
les hypothèses sont trivialement vérifiées.
\end{proof}

\begin{env}[2.1.9]
\label{III.2.1.9-fr}
Passons à la limite inductive sur $n$~; les relations
(\hyperref[III.2.1.8.1-fr]{2.1.8.1}) donnent
$H^i((\mathbf T),S)=0$ pour $i\neq r+1$. Pour $i=r+1$, le système inductif
est formé des $S/(\mathbf T^n)$, l'homomorphisme
\[
\varphi_{nm}:S/(\mathbf T^n)\longrightarrow S/(\mathbf T^m)
\qquad(0\leq n\leq m)
\]
étant la multiplication par $(T_0\cdots T_r)^{m-n}$. Pour
$n\geq\sup(p_i)_{0\leq i\leq r}$, désignons par
$\xi_{\mathbf p}^{(n)}=\xi_{p_0,\ldots,p_r}^{(n)}$ la classe de
$T_0^{n-p_0}\cdots T_r^{n-p_r}\pmod{(\mathbf T^n)}$~; on a alors
$\varphi_{nm}(\xi_{\mathbf p}^{(n)})=\xi_{\mathbf p}^{(m)}$, et ces éléments
ont donc même image canonique
$\xi_{\mathbf p}=\xi_{p_0,\ldots,p_r}$ dans la limite inductive
$H^{r+1}((\mathbf T),S)$~; en vertu de la définition du degré donnée dans
(\hyperref[III.2.1.2-fr]{2.1.2}), le degré de $\xi_{\mathbf p}$ est donc
égal à $-|\mathbf p|=-(p_0+p_1+\cdots+p_r)$. Il est clair que les
$\xi_{\mathbf p}^{(n)}$ pour $0<p_i\leq n$ et $0\leq i\leq r$ forment une
base de $S/(\mathbf T^n)$. On déduit donc aussitôt de
(\hyperref[III.2.1.8-fr]{2.1.8})~:
\end{env}

\begin{corollary}[2.1.10]
\label{III.2.1.10-fr}
Avec les notations de (\hyperref[III.2.1.8-fr]{2.1.8}), on a
\begin{equation}
\label{III.2.1.10.1-fr}
\tag{2.1.10.1}
H^i((\mathbf T),S)=0\qquad\text{pour }i\neq r+1,
\end{equation}
et $H^{r+1}((\mathbf T),S)$ est un $A$-module libre dont une base est formée
des éléments $\xi_{p_0,\ldots,p_r}$ tels que $p_i>0$ pour $0\leq i\leq r$.
\end{corollary}

\begin{env}[2.1.11]
\label{III.2.1.11-fr}
\emph{Remarque.} --- Soit $N$ un $A$-module quelconque et soit
$M=S\otimes_A N$~; le raisonnement de
(\hyperref[III.2.1.8-fr]{2.1.8}) montre que l'on a plus généralement
\begin{align}
\label{III.2.1.11.1-fr}
\tag{2.1.11.1}
H^i(\mathbf T^n,M)&=0 &&\text{si }i\neq r+1,\\
\label{III.2.1.11.2-fr}
\tag{2.1.11.2}
H^{r+1}(\mathbf T^n,M)&=(S/(\mathbf T^n))\otimes_A N,
\end{align}
car la dernière formule se déduit directement de (1.1.3.5), et d'autre
part il est clair que
\[
M/(T_0^nM+\cdots+T_{i-1}^nM)
=\bigl(S/(T_0^nS+\cdots+T_{i-1}^nS)\bigr)\otimes_A N,
\]
l'idéal $T_0^nS+\cdots+T_{i-1}^nS$ étant facteur direct dans le
$A$-module $S$~; cela permet d'appliquer (1.1.4) à $M$, et on obtient ainsi
(\hyperref[III.2.1.11.1-fr]{2.1.11.1}).
\end{env}

Combinant (\hyperref[III.2.1.10-fr]{2.1.10}) et
(\hyperref[III.2.1.5-fr]{2.1.5}), on obtient~:

\begin{proposition}[2.1.12]
\label{III.2.1.12-fr}
Soient $A$ un anneau, $r$ un entier $>0$, et
$X=\mathbf P_A^r$ (II, 4.1.1). Alors~:
\begin{enumerate}
\item[{\rm(i)}] On a $H^i(X,\mathcal O_X(*))=0$ pour $i\neq0,r$.
\item[{\rm(ii)}] L'homomorphisme canonique
$\alpha:S\to H^0(X,\mathcal O_X(*))$ (II, 2.6.2) est bijectif.
\oldpage[III]{99}
\item[{\rm(iii)}] $H^r(X,\mathcal O_X(*))$ est un $A$-module libre ayant
une base formée d'éléments $\xi_{p_0,\ldots,p_r}$, où $p_i>0$ pour
$0\leq i\leq r$, $\xi_{p_0,\ldots,p_r}$ étant de degré
$-|\mathbf p|=-(p_0+\cdots+p_r)$, et le produit
$T_i\xi_{p_0,\ldots,p_r}$ étant
$\xi_{p_0,\ldots,p_i-1,\ldots,p_r}$.
\end{enumerate}
\end{proposition}

\begin{proof}
Remarquons en effet que, dans la suite exacte
(\hyperref[III.2.1.5.2-fr]{2.1.5.2}) appliquée à
$M=S=A[T_0,\ldots,T_r]$, on a
$H^0((\mathbf T),S)=0$ et $H^1((\mathbf T),S)=0$ d'après
(\hyperref[III.2.1.10.1-fr]{2.1.10.1}), et que la prop.
(\hyperref[III.2.1.5-fr]{2.1.5}) s'applique à $U=X$, puisque $X$ est
réunion des $D_+(T_i)$ (II, 2.3.14). Il reste à identifier l'application
$S\to H^0(X,\mathcal O_X(*))$ de la suite exacte
(\hyperref[III.2.1.5.2-fr]{2.1.5.2}) et l'application canonique $\alpha$~;
mais cela résulte de l'identification canonique de
$H^0(U,\mathcal O_X(*))$ et de $H^0(\mathfrak U,\mathcal O_X(*))$.
\end{proof}

\begin{corollary}[2.1.13]
\label{III.2.1.13-fr}
Les seules valeurs de $(i,n)$ pour lesquelles on puisse avoir
$H^i(X,\mathcal O_X(n))\neq0$ sont les suivantes~:
$i=0$ et $n\geq0$, $i=r$ et $n\leq-(r+1)$.
\end{corollary}

On notera que si $A\neq0$, on a effectivement
$H^i(X,\mathcal O_X(n))\neq0$ pour les couples énumérés dans
(\hyperref[III.2.1.13-fr]{2.1.13})~; cela résulte de
(\hyperref[III.2.1.12-fr]{2.1.12}), puisque $S_n$ est alors $\neq0$ pour
tous les degrés $n\geq0$.

Dans les applications qui seront faites dans ce chapitre, nous utiliserons
surtout le résultat moins précis~:

\begin{corollary}[2.1.14]
\label{III.2.1.14-fr}
Les $A$-modules $H^i(X,\mathcal O_X(n))$ sont libres de type fini~; si
$i>0$, ils sont nuls pour $n>0$.
\end{corollary}

\begin{proposition}[2.1.15]
\label{III.2.1.15-fr}
Soient $Y$ un préschéma, $\mathcal E$ un $\mathcal O_Y$-Module localement
libre de rang $r+1$, $X=\mathbf P(\mathcal E)$ le fibré projectif défini
par $\mathcal E$, $f:X\to Y$ le morphisme structural. Les seules valeurs de
$i$ et $n$ pour lesquelles $R^if_*(\mathcal O_X(n))\neq0$ sont
$i=0$ et $n\geq0$, $i=r$ et $n\leq-(r+1)$~; en outre, l'homomorphisme
canonique (II, 3.3.2)
\[
\alpha:S_{\mathcal O_Y}(\mathcal E)
\longrightarrow\Gamma_*(\mathcal O_X)
=R^0f_*(\mathcal O_X(*))
=\bigoplus_{n\in\mathbb Z}f_*(\mathcal O_X(n))
\]
est un isomorphisme.
\end{proposition}

\begin{proof}
La question étant locale sur $Y$, on peut supposer $Y$ affine d'anneau
$A$ et $\mathcal E=\widetilde E$, où $E=A^{r+1}$~; on est alors
immédiatement ramené à (\hyperref[III.2.1.12-fr]{2.1.12}), compte tenu de
(1.4.11).
\end{proof}

\begin{env}[2.1.16]
\label{III.2.1.16-fr}
\emph{Remarque.} --- Nous compléterons plus tard les résultats de
(\hyperref[III.2.1.15-fr]{2.1.15}) en démontrant les propositions
suivantes~: posons
\[
\boldsymbol\omega=f^*\left(\bigwedge^{r+1}\mathcal E\right)(-r-1),
\]
qui est un $\mathcal O_X$-Module inversible. Alors~:
\begin{enumerate}
\item[{\rm(i)}] On a un isomorphisme canonique
\begin{equation}
\label{III.2.1.16.1-fr}
\tag{2.1.16.1}
\rho:R^rf_*(\boldsymbol\omega)\xleftarrow{\sim}\mathcal O_Y.
\end{equation}
\item[{\rm(ii)}] L'accouplement par cup-produit (0, 12.2.2)
\begin{equation}
\label{III.2.1.16.2-fr}
\tag{2.1.16.2}
R^rf_*(\mathcal O_X(n))\times R^0f_*(\boldsymbol\omega(-n))
\longrightarrow R^rf_*(\boldsymbol\omega)
\end{equation}
composé avec l'isomorphisme $\rho^{-1}$, définit un \emph{isomorphisme}
de $R^rf_*(\mathcal O_X(n))$ sur le \emph{dual} du $\mathcal O_Y$-Module
localement libre
\[
R^0f_*(\boldsymbol\omega(-n))
=\left(\bigwedge^{r+1}\mathcal E\right)
\otimes_{\mathcal O_Y}\bigl(S_{\mathcal O_Y}(\mathcal E)\bigr)_{-n}.
\]
\end{enumerate}
\end{env}
\subsection{Le théorème fondamental des morphismes projectifs}
\label{subsection:III.2.2-fr}
\oldpage[III]{100}

\begin{theorem}[2.2.1]
\label{III.2.2.1-fr}
(Serre). Soient $Y$ un préschéma noethérien, $f:X\to Y$ un morphisme
propre, $\mathcal L$ un $\mathcal O_X$-Module inversible ample pour $f$.
Pour tout $\mathcal O_X$-Module $\mathcal F$, posons
\[
\mathcal F(n)=\mathcal F\otimes_{\mathcal O_X}\mathcal L^{\otimes n}
\qquad\text{pour tout }n\in\mathbb Z.
\]
Alors, pour tout $\mathcal O_X$-Module cohérent $\mathcal F$~:
\begin{enumerate}
\item[{\rm(i)}] Les $R^qf_*(\mathcal F)$ sont des
$\mathcal O_Y$-Modules cohérents.
\item[{\rm(ii)}] Il existe un entier $N$ tel que pour $n\geq N$, on ait
\[
R^qf_*(\mathcal F(n))=0\qquad\text{pour tout }q>0.
\]
\item[{\rm(iii)}] Il existe un entier $N$ tel que, pour $n\geq N$,
l'homomorphisme canonique
\[
f^*(f_*(\mathcal F(n)))\longrightarrow\mathcal F(n)
\]
soit surjectif.
\end{enumerate}
\end{theorem}

\begin{proof}
Notons d'abord que si le théorème est vrai quand on y remplace
$\mathcal L$ par $\mathcal L^{\otimes d}$ $(d>0)$, il est vrai sous sa
forme initiale. En effet, on peut alors écrire
\[
\mathcal F(n)
=(\mathcal F\otimes\mathcal L^{\otimes r})
\otimes\mathcal L^{\otimes hd}
\]
avec un $h>0$ et $0\leq r<d$, et par hypothèse pour chaque $r$ il y a
un entier $N_r$ tel que pour $h\geq N_r$ les propriétés (ii) et (iii)
aient lieu pour le $\mathcal O_X$-Module
$\mathcal F\otimes\mathcal L^{\otimes r}$~; prenant pour $N$ le plus grand
des $dN_r$, (ii) et (iii) auront lieu pour $n\geq N$. On peut donc
supposer $\mathcal L$ très ample relativement à $f$ (II, 4.6.11)~; il
existe par suite une $Y$-immersion ouverte dominante
$i:X\to P$, où $P=\operatorname{Proj}(\mathcal S)$, où $\mathcal S$ est
une $\mathcal O_Y$-Algèbre quasi-cohérente graduée à degrés positifs, dans
laquelle $\mathcal S_1$ est de type fini et engendre $\mathcal S$~; en
outre, $\mathcal L$ est isomorphe à $i^*(\mathcal O_P(1))$ (II, 4.4.7).
Mais comme $f$ est propre, il en est de même de $i$ (II, 5.4.4), donc
$i$ est un isomorphisme $X\xrightarrow{\sim}P$. On peut donc se borner au
cas où $X=\operatorname{Proj}(\mathcal S)$ et
$\mathcal L=\mathcal O_X(1)$. Le th. (2.2.1) est alors conséquence de la
\end{proof}

\begin{proposition}[2.2.2]
\label{III.2.2.2-fr}
Soient $A$ un anneau noethérien, $S$ une $A$-algèbre graduée à degrés
positifs, dans laquelle $S_1$ est un $A$-module ayant un système de $r+1$
générateurs, et qui engendre l'algèbre $S$. Soit
$X=\operatorname{Proj}(S)$. Pour tout $\mathcal O_X$-Module cohérent
$\mathcal F$~:
\begin{enumerate}
\item[{\rm(i)}] Les $A$-modules $H^q(X,\mathcal F)$ sont de type fini.
\item[{\rm(ii)}] On a $H^q(X,\mathcal F)=0$ pour $q>r$.
\item[{\rm(iii)}] Il existe un entier $N$ tel que pour $n\geq N$, on ait
$H^q(X,\mathcal F(n))=0$ pour tout $q>0$.
\item[{\rm(iv)}] Il existe un entier $N$ tel que pour $n\geq N$,
$\mathcal F(n)$ soit engendré par ses sections au-dessus de $X$.
\end{enumerate}
\end{proposition}

\begin{proof}
Montrons d'abord comment (2.2.2) entraîne (2.2.1)~: dans (2.2.1)
(ramené au cas particulier $X=\operatorname{Proj}(\mathcal S)$ considéré
ci-dessus), $Y$ est quasi-compact, donc peut être recouvert par un nombre
fini d'ouverts affines, d'anneaux noethériens, tels que la restriction de
$\mathcal S_1$ à chacun de ces ouverts $U_\alpha$ soit engendrée par un
nombre fini de sections de $\mathcal S_1$ au-dessus de $U_\alpha$. Si on
suppose (2.2.2) démontré, il suffira alors de prendre pour $N$ dans les
parties (ii) et (iii) de (2.2.1) le plus grand des entiers analogues
correspondant aux $U_\alpha$ (tenant compte de (1.4.11) et de
(II, 3.4.7)).

Pour prouver (2.2.2), remarquons que $X$ s'identifie à un sous-schéma
fermé de $P=\mathbf P_A^r$ (II, 3.6.2)~; en outre, si $j:X\to P$ est
l'injection canonique, $j_*(\mathcal F)$ est un
$\mathcal O_P$-Module cohérent et on a
$j_*(\mathcal F(n))=(j_*(\mathcal F))(n)$ (II, 3.4.5 et 3.5.2).
Compte tenu de (G, II, cor. du th. 4.9.1), on est donc ramené à prouver
(2.2.2) dans le cas particulier où
\[
X=\mathbf P_A^r\quad\text{et}\quad S=A[T_0,\ldots,T_r].
\]

\oldpage[III]{101}
Comme $X$ est recouvert par les ouverts affines $D_+(T_i)$ en nombre
$r+1$, (ii) résulte de (1.4.12). Notons d'autre part que (iv) a déjà été
démontré (II, 2.7.9).

Nous allons prouver simultanément (i) et (iii). Notons que ces assertions
sont vraies pour $\mathcal F=\mathcal O_X(m)$
(\hyperref[III.2.1.13-fr]{2.1.13})~; elles le sont donc aussi lorsque
$\mathcal F$ est somme directe d'un nombre fini de
$\mathcal O_X$-Modules de la forme $\mathcal O_X(m_j)$. D'autre part,
(i) et (iii) sont vraies trivialement pour $q>r$ en vertu de (ii). Nous
allons procéder par \emph{récurrence descendante} sur $q$. On sait que
$\mathcal F$ est isomorphe à un quotient d'une somme directe $\mathcal E$
d'un nombre fini de faisceaux $\mathcal O_X(m_j)$ (II, 2.7.10)~;
autrement dit, on a une suite exacte
\[
0\longrightarrow\mathcal R\longrightarrow\mathcal E
\longrightarrow\mathcal F\longrightarrow0,
\]
où $\mathcal R$ est cohérent (0\textsubscript{I}, 5.3.3) et où
$\mathcal E$ vérifie (i) et (iii). Comme $\mathcal F(n)$ est un foncteur
exact en $\mathcal F$, on a aussi la suite exacte
\[
0\longrightarrow\mathcal R(n)\longrightarrow\mathcal E(n)
\longrightarrow\mathcal F(n)\longrightarrow0
\]
pour tout $n\in\mathbb Z$. On en déduit la suite exacte de cohomologie
\[
H^{q-1}(X,\mathcal E(n))\longrightarrow
H^{q-1}(X,\mathcal F(n))\longrightarrow
H^q(X,\mathcal R(n)).
\]
Comme $\mathcal E(n)$ est somme directe des $\mathcal O_X(n+m_j)$
(II, 2.5.14), $H^{q-1}(X,\mathcal E(n))$ est de type fini, et il en est
de même de $H^q(X,\mathcal R(n))$ par l'hypothèse de récurrence~; comme
$A$ est noethérien, on en conclut que
$H^{q-1}(X,\mathcal F(n))$ est de type fini pour tout $n\in\mathbb Z$, et
en particulier pour $n=0$. D'autre part, par l'hypothèse de récurrence,
il existe un entier $N$ tel que pour $n\geq N$ on ait
$H^q(X,\mathcal R(n))=0$~; par ailleurs, on peut supposer aussi $N$ choisi
tel que $H^{q-1}(X,\mathcal E(n))=0$ pour $n\geq N$, puisque $\mathcal E$
vérifie (iii)~; on en conclut que
$H^{q-1}(X,\mathcal F(n))=0$ pour $n\geq N$, ce qui achève la
démonstration.
\end{proof}

\begin{corollary}[2.2.3]
\label{III.2.2.3-fr}
Sous les hypothèses de (\hyperref[III.2.2.1-fr]{2.2.1}), soit
$\mathcal F\to\mathcal G\to\mathcal H$ une suite exacte de
$\mathcal O_X$-Modules cohérents. Il existe alors un entier $N$ tel que
pour $n\geq N$, la suite
\[
f_*(\mathcal F(n))\longrightarrow f_*(\mathcal G(n))
\longrightarrow f_*(\mathcal H(n))
\]
soit exacte.
\end{corollary}

\begin{proof}
Soient $\mathcal F'$, $\mathcal G'$, $\mathcal G''$ le noyau, l'image et
le conoyau de $\mathcal F\to\mathcal G$~; $\mathcal G'$ est le noyau et
$\mathcal G''$ l'image de $\mathcal G\to\mathcal H$, soit $\mathcal H''$
le conoyau de cet homomorphisme~; tous ces $\mathcal O_X$-Modules sont
cohérents (0\textsubscript{I}, 5.3.4). Comme $\mathcal F(n)$ est un
foncteur exact en $\mathcal F$, il suffit de prouver que pour $n$ assez
grand, chacune des suites
\begin{align*}
0&\longrightarrow f_*(\mathcal F'(n))
\longrightarrow f_*(\mathcal F(n))
\longrightarrow f_*(\mathcal G'(n))\longrightarrow0,\\
0&\longrightarrow f_*(\mathcal G'(n))
\longrightarrow f_*(\mathcal G(n))
\longrightarrow f_*(\mathcal G''(n))\longrightarrow0,\\
0&\longrightarrow f_*(\mathcal G''(n))
\longrightarrow f_*(\mathcal H(n))
\longrightarrow f_*(\mathcal H''(n))\longrightarrow0
\end{align*}
est exacte~; par suite, on peut supposer que
$0\to\mathcal F\to\mathcal G\to\mathcal H\to0$ est exacte. On a alors
la suite exacte de cohomologie
\[
0\longrightarrow f_*(\mathcal F(n))
\longrightarrow f_*(\mathcal G(n))
\longrightarrow f_*(\mathcal H(n))
\longrightarrow R^1f_*(\mathcal F(n))\longrightarrow\cdots
\]
et la conclusion résulte de (\hyperref[III.2.2.1-fr]{2.2.1}, (ii)).
\end{proof}

\begin{corollary}[2.2.4]
\label{III.2.2.4-fr}
Soient $Y$ un préschéma noethérien, $f:X\to Y$ un morphisme de type fini,
$\mathcal L$ un $\mathcal O_X$-Module inversible ample pour $f$~; pour
tout $\mathcal O_X$-Module $\mathcal F$, on pose
$\mathcal F(n)=\mathcal F\otimes_{\mathcal O_X}\mathcal L^{\otimes n}$
(pour $n\in\mathbb Z$). Soit
$\mathcal F\to\mathcal G\to\mathcal H$ une suite exacte de
$\mathcal O_X$-Modules cohérents, telle que les supports de $\mathcal F$
et de $\mathcal H$ soient propres sur $Y$ (II, 5.4.10). Il existe alors
un entier $N$ tel que, pour $n\geq N$, la suite
\[
f_*(\mathcal F(n))\longrightarrow f_*(\mathcal G(n))
\longrightarrow f_*(\mathcal H(n))
\]
soit exacte.
\end{corollary}

\begin{proof}
Le même raisonnement qu'au début de
(\hyperref[III.2.2.1-fr]{2.2.1}) montre que si le corollaire est vrai
pour $\mathcal L^{\otimes d}$ $(d>0)$, il l'est aussi pour $\mathcal L$~;
on peut donc se borner au cas où $\mathcal L$ est très ample pour $f$
(II, 4.6.11), et par suite on peut identifier $X$ à un ouvert dans un
$Y$-schéma $Z=\operatorname{Proj}(\mathcal S)$, où $\mathcal S$ est une
$\mathcal O_Y$-Algèbre quasi-cohérente graduée à degrés positifs, dans
laquelle $\mathcal S_1$ est de type fini et engendre $\mathcal S$, de
sorte que $\mathcal L=i^*(\mathcal O_Z(1))$, où $i$ est l'immersion
canonique $X\to Z$ (II, 4.4.7). Cela étant, comme
$\operatorname{Supp}(\mathcal G)$ est fermé dans $X$ et contenu dans
$\operatorname{Supp}(\mathcal F)\cap\operatorname{Supp}(\mathcal H)$, il
est propre sur $Y$~; les supports de $\mathcal F$, $\mathcal G$,
$\mathcal H$ sont donc \emph{fermés dans $Z$} (II, 5.4.10). Les faisceaux
$\mathcal F'=i_*(\mathcal F)$, $\mathcal G'=i_*(\mathcal G)$,
$\mathcal H'=i_*(\mathcal H)$ sont donc des $\mathcal O_Z$-Modules
cohérents, et la suite $\mathcal F'\to\mathcal G'\to\mathcal H'$ est
exacte~; en outre, si $g:Z\to Y$ est le morphisme structural, on a
$f=g\circ i$, et il est clair que
$\mathcal F'(n)=i_*(\mathcal F(n))$, et de même pour $\mathcal G'$ et
$\mathcal H'$~; la conclusion résulte donc de
(\hyperref[III.2.2.3-fr]{2.2.3}) appliqué à
$\mathcal F'$, $\mathcal G'$, $\mathcal H'$.
\end{proof}

\begin{env}[2.2.5]
\label{III.2.2.5-fr}
\emph{Remarques.} ---
(i) L'assertion (i) de (\hyperref[III.2.2.1-fr]{2.2.1}) est encore vraie
lorsqu'on suppose seulement que $Y$ est \emph{localement noethérien}~; en
effet, la propriété est évidemment locale sur $Y$~; d'autre part, les
hypothèses de (\hyperref[III.2.2.1-fr]{2.2.1}) impliquent que pour tout
ouvert $U\subset Y$, la restriction de $f$ à $f^{-1}(U)$ est un morphisme
projectif $f^{-1}(U)\to U$ (II, 5.5.5, (iii)) et
$\mathcal L|f^{-1}(U)$ est ample pour ce morphisme (II, 4.6.4).

(ii) L'assertion (iii) de (\hyperref[III.2.2.1-fr]{2.2.1}) est encore
valable, comme on l'a vu, lorsque l'on suppose seulement que $X$ est un
schéma quasi-compact ou un préschéma dont l'espace sous-jacent est
noethérien, et $f:X\to Y$ un morphisme quasi-compact (II, 4.6.8). Mais il
faut noter que même lorsqu'on suppose que $Y$ est le \emph{spectre d'un
corps} $K$ et que $f$ est \emph{quasi-projectif}, l'assertion (ii) de
(\hyperref[III.2.2.1-fr]{2.2.1}) n'est plus nécessairement vérifiée. Par
exemple, soit $X'=\operatorname{Spec}(K[T_0,\ldots,T_r])$ et soit $X$ la
réunion des ouverts affines $D(T_i)$ de $X'$ $(0\leq i\leq r)$~; comme
l'immersion $X\to X'$ est quasi-compacte, le morphisme structural
$f:X\to Y$ est quasi-affine (II, 5.1.10), donc $\mathcal O_X$ est très
ample pour $f$ (II, 5.1.6). Mais l'anneau
$\Gamma(X,\mathcal O_X)$ s'identifie à l'intersection des anneaux de
fractions $(K[T_0,\ldots,T_r])_{T_i}$ pour $0\leq i\leq r$
(I, 8.2.1.1), c'est-à-dire à $K[T_0,\ldots,T_r]$. Par suite, il résulte
des formules (1.4.3.1) et (1.1.3.5) que l'on a
$H^r(X,\mathcal O_X^{\otimes n})=H^r(X,\mathcal O_X)=A\neq0$ pour tout
$n$.
\end{env}

\subsection{Application aux faisceaux gradués d'algèbres et de modules}
\label{subsection:III.2.3-fr}

\begin{theorem}[2.3.1]
\label{III.2.3.1-fr}
Soient $Y$ un préschéma noethérien, $\mathcal S$ une
$\mathcal O_Y$-Algèbre graduée à degrés positifs, quasi-cohérente et de
type fini, $X=\operatorname{Proj}(\mathcal S)$,
$q:X\to Y$ le morphisme structural, $\mathcal M$ un $\mathcal S$-Module
gradué quasi-cohérent vérifiant la condition \textbf{(TF)}. Alors il existe
un entier $N$ tel que, pour $n\geq N$, l'homomorphisme canonique
(II, 8.14.5.1)
\[
\alpha_n:\mathcal M_n\longrightarrow
q_*(\shProj_0(\mathcal M(n)))
=q_*((\shProj(\mathcal M))_n)
\]
\oldpage[III]{103}
soit bijectif. En d'autres termes, l'homomorphisme canonique
\[
\alpha:\mathcal M\longrightarrow\bbGamma_*(\shProj(\mathcal M))
\]
est un \textbf{(TN)}-isomorphisme.
\end{theorem}

\begin{proof}
On peut se borner au cas où $\mathcal M$ est un $\mathcal S$-module de
type fini (II, 3.4.2). Comme $Y$ est quasi-compact, il existe un entier
$d>0$ tel que $\mathcal S^{(d)}$ soit engendrée par le
$\mathcal O_Y$-Module quasi-cohérent $\mathcal S_d$, ce dernier étant de
type fini (II, 3.1.10), donc cohérent puisque $Y$ est noethérien.
Remarquons maintenant que $\mathcal M$ est somme directe des
$\mathcal M^{(d,k)}$ pour $0\leq k<d$ et que chacun des
$\mathcal M^{(d,k)}$ est un $\mathcal S^{(d)}$-Module quasi-cohérent de
type fini, ainsi qu'il résulte de (II, 2.1.6, (iii)), la question étant
locale sur $Y$. Or, il suffit évidemment de prouver que chacun des
homomorphismes canoniques
\[
\alpha:\mathcal M^{(d,k)}
\longrightarrow\bbGamma_*((\shProj(\mathcal M))^{(d,k)})
\]
est un \textbf{(TN)}-isomorphisme. Compte tenu de (II, 8.14.13), et
notamment du diagramme (8.14.13.4), on voit qu'on est ramené à prouver le
théorème lorsque $\mathcal S$ est engendrée par $\mathcal S_1$ et que
$\mathcal S_1$ est un $\mathcal O_Y$-Module cohérent.

Comme $Y$ est noethérien, le même raisonnement qu'au début de
(\hyperref[III.2.2.2-fr]{2.2.2}) montre qu'on peut se borner au cas où
$Y=\operatorname{Spec}(A)$, $\mathcal S=\widetilde S$,
$\mathcal M=\widetilde M$, $A$ étant un anneau noethérien, $S_1$ un
$A$-module de type fini et $M$ un $S$-module gradué de type fini.
Montrons qu'il suffit alors de prouver le théorème lorsque $M=S$. En effet,
dans le cas général, on a une suite exacte
\[
L'\longrightarrow L\longrightarrow M\longrightarrow0,
\]
où $L$ et $L'$ sont des sommes directes de modules gradués de la forme
$S(m)$. Si le résultat est vrai pour $M=S$, il l'est aussi pour $M=S(m)$,
donc pour $L$ et $L'$. Considérons alors le diagramme commutatif
\[
\xymatrix{
\widetilde {L'_n}\ar[r]\ar[d]_{\alpha_n} &
\widetilde {L_n}\ar[r]\ar[d]^{\alpha_n} &
\widetilde {M_n}\ar[r]\ar[d]^{\alpha_n} & 0\\
q_*(\widetilde {L'}(n))\ar[r] &
q_*(\widetilde L(n))\ar[r] &
q_*(\widetilde M(n))\ar[r] & 0 .
}
\]
La deuxième ligne est exacte en vertu de
(\hyperref[III.2.2.3-fr]{2.2.3}), dès que $n$ est assez grand~; comme il
en est de même de la première et que les deux flèches verticales de gauche
sont des isomorphismes, il en est de même de la troisième.

Cela étant, pour prouver le théorème lorsque $M=S$, supposons d'abord que
$S=A[T_0,\ldots,T_r]$ ($T_i$ indéterminées)~; dans ce cas, notre assertion
n'est autre que (\hyperref[III.2.1.11-fr]{2.1.11}, (ii)). Dans le cas
général, $S$ s'identifie à un quotient d'un anneau
$S'=A[T_0,\ldots,T_r]$ par un idéal gradué, donc $X$ à un sous-schéma
fermé de $X'=\mathbf P_A^r$ (II, 2.9.2). Si $j$ est l'injection canonique
$X\to X'$, $j_*(\widetilde S(n))$ n'est autre que le
$\mathcal O_{X'}$-Module $(\shProj(\widetilde S))(n)$ où $S$ est considéré
comme un $S'$-module gradué~; cela résulte en effet aussitôt de
(II, 2.8.7). Comme $j_*(\widetilde S(n))$ est un
$\mathcal O_{X'}$-Module vérifiant \textbf{(TF)}, l'homomorphisme canonique
$\alpha_n:S_n\to\Gamma(X',j_*(\widetilde S(n)))$ est bijectif pour $n$
assez grand, en vertu de ce qui précède~; cela achève la démonstration,
puisque
$\Gamma(X',j_*(\widetilde S(n)))=\Gamma(X,\widetilde S(n))$.
\end{proof}

\oldpage[III]{104}
\begin{corollary}[2.3.2]
\label{III.2.3.2-fr}
Sous les hypothèses de (\hyperref[III.2.3.1-fr]{2.3.1}), soit
\[
\mathcal S_X=\bigoplus_{n\in\mathbb Z}\mathcal O_X(n),
\]
et soit $\mathcal F$ un $\mathcal S_X$-Module gradué quasi-cohérent de
type fini. Alors $\bbGamma_*(\mathcal F)$ vérifie la condition
\textbf{(TF)}.
\end{corollary}

\begin{proof}
On a vu dans la démonstration de
(\hyperref[III.2.3.1-fr]{2.3.1}) que $X$, qui est isomorphe à
$\operatorname{Proj}(\mathcal S^{(d)})$ (II, 3.1.8), est de type fini sur
$Y$ (II, 3.4.1). Il résulte alors de (II, 8.14.9) que $\mathcal F$ est
isomorphe à un $\mathcal S_X$-Module gradué de la forme
$\shProj(\mathcal M)$, où $\mathcal M$ est un $\mathcal S$-Module gradué
quasi-cohérent de type fini. En vertu de
(\hyperref[III.2.3.1-fr]{2.3.1}), $\bbGamma_*(\mathcal F)$ est
\textbf{(TN)}-isomorphe à $\mathcal M$ et par suite vérifie
\textbf{(TF)}.
\end{proof}

\begin{env}[2.3.3]
\label{III.2.3.3-fr}
\emph{Scholie.} --- Soient $Y$ un préschéma noethérien, $\mathcal S$ une
$\mathcal O_Y$-Algèbre graduée vérifiant les conditions de
(\hyperref[III.2.3.1-fr]{2.3.1}) et $X=\operatorname{Proj}(\mathcal S)$.
Soient $K_{\mathcal S}$ la catégorie abélienne des $\mathcal S$-Modules
gradués quasi-cohérents vérifiant \textbf{(TF)}, $K'_{\mathcal S}$ la
sous-catégorie de $K_{\mathcal S}$ formée des $\mathcal S$-Modules
vérifiant \textbf{(TN)}~; enfin, soit $K_X$ la catégorie des
$\mathcal S_X$-Modules gradués quasi-cohérents de type fini $\mathcal F$
(ce qui revient à dire, puisque $\mathcal S_X$ est périodique
(II, 8.14.4 et 8.14.12), que les $\mathcal F_i$ sont des
$\mathcal O_X$-Modules cohérents). Alors les foncteurs
\[
\mathcal M\longmapsto\shProj(\mathcal M)
\quad\text{dans }K_{\mathcal S},
\qquad
\mathcal F\longmapsto\bbGamma_*(\mathcal F)
\quad\text{dans }K_X
\]
définissent, en vertu de (II, 8.14.8 et 8.14.10) et de
(\hyperref[III.2.3.2-fr]{2.3.2}), une équivalence (T, I, 1.2) de la
catégorie quotient $K_{\mathcal S}/K'_{\mathcal S}$ (T, I, 1.11) avec la
catégorie $K_X$. Lorsque $\mathcal S$ est engendrée par $\mathcal S_1$,
on peut remplacer $K_X$ par la catégorie des $\mathcal O_X$-Modules
cohérents (II, 8.14.12).
\end{env}

\begin{proposition}[2.3.4]
\label{III.2.3.4-fr}
Soit $Y$ un préschéma noethérien.
\begin{enumerate}
\item[{\rm(i)}] Soit $\mathcal S$ une $\mathcal O_Y$-Algèbre graduée à
degrés positifs, quasi-cohérente et de type fini. Soient
$X=\operatorname{Proj}(\mathcal S)$, et
\[
\mathcal S_X=\shProj(\mathcal S)
=\bigoplus_{n\in\mathbb Z}\mathcal O_X(n).
\]
Alors $\mathcal S_X$ est une $\mathcal O_X$-Algèbre graduée périodique
(II, 8.14.12) dont les composants homogènes
$(\mathcal S_X)_n=\mathcal O_X(n)$ sont des $\mathcal O_X$-Modules
cohérents, et si $d>0$ est une période de $\mathcal S_X$,
$(\mathcal S_X)_d=\mathcal O_X(d)$ est un $\mathcal O_X$-Module
inversible $Y$-ample. En outre, l'homomorphisme canonique
\[
\alpha:\mathcal S\longrightarrow\bbGamma_*(\mathcal S_X)
\]
est un \textbf{(TN)}-isomorphisme.
\item[{\rm(ii)}] Inversement, soit $q:X\to Y$ un morphisme projectif, et
soit $\mathcal S'$ une $\mathcal O_X$-Algèbre graduée, dont les composants
homogènes $\mathcal S'_n$ $(n\in\mathbb Z)$ sont des
$\mathcal O_X$-Modules cohérents, et qui admet une période $d>0$ telle que
$\mathcal S'_d$ soit un $\mathcal O_X$-Module inversible ample pour $q$.
Alors
\[
\mathcal S=\bigoplus_{n\geq0}q_*(\mathcal S'_n)
\]
est une $\mathcal O_Y$-Algèbre graduée à degrés positifs quasi-cohérente
et de type fini, et il existe un $Y$-isomorphisme
\[
r:X\xrightarrow{\sim}\operatorname{Proj}(\mathcal S)
\]
tel que $r^*(\shProj(\mathcal S))$ soit isomorphe (en tant que
$\mathcal O_X$-Algèbre graduée) à $\mathcal S'$.
\end{enumerate}
\end{proposition}

\begin{proof}
(i) Toutes les assertions ont pratiquement déjà été démontrées, la
dernière n'étant autre qu'un cas particulier de
(\hyperref[III.2.3.2-fr]{2.3.2}). Le fait que $\mathcal S_X$ est
périodique a été vu en (II, 8.14.14) et le fait qu'il y a une période
$d>0$ telle que $\mathcal O_X(d)$ soit inversible et $Y$-ample n'est
autre que (II, 4.6.18). Enfin, pour $0\leq k<d$,
$(\mathcal S_X)^{(d,k)}$ est un $(\mathcal S_X)^{(d)}$-Module de type fini
(II, 8.14.14), donc chacun des $(\mathcal S_X)_n$ est un
$\mathcal O_X$-Module quasi-cohérent de type fini, en vertu de
(II, 2.1.6, (ii))~; la question étant locale, comme $\mathcal O_X$ est
cohérent, il en est de même des $(\mathcal S_X)_n$.

(ii) Quitte à remplacer la période $d$ par un de ses multiples, on peut
supposer que $\mathcal L=\mathcal S'_d$ est un $\mathcal O_X$-Module
\emph{très ample} relativement à $q$ (II, 4.6.11). On a en outre
\[
\mathcal S'^{(d)}=\bigoplus_{n\in\mathbb Z}\mathcal L^{\otimes n}
\quad\text{par hypothèse, donc}\quad
\mathcal S^{(d)}=\bigoplus_{n\geq0}q_*(\mathcal L^{\otimes n})~;
\]
on sait (II, 3.1.8 et 3.2.9) qu'il y a un $Y$-isomorphisme $s$ de
$X'=\operatorname{Proj}(\mathcal S)$ sur
$X''=\operatorname{Proj}(\mathcal S^{(d)})$ tel que
\oldpage[III]{105}
\[
s^*(\mathcal O_{X''}(n))=\mathcal O_{X'}(nd).
\]
On établira donc l'existence d'un $Y$-isomorphisme $X\xrightarrow{\sim}X'$
si l'on prouve la

\begin{env}[2.3.4.1]
\label{III.2.3.4.1-fr}
Soient $Y$ un préschéma noethérien, $q:X\to Y$ un morphisme projectif,
$\mathcal L$ un $\mathcal O_X$-Module inversible très ample pour $q$.
Alors $\mathcal S=\bigoplus_{n\geq0}q_*(\mathcal L^{\otimes n})$ est une
$\mathcal O_Y$-Algèbre graduée quasi-cohérente de type fini, telle que
$\mathcal S_n=\mathcal S_1^n$ pour $n$ assez grand, et il existe un
$Y$-isomorphisme $r:X\xrightarrow{\sim}P=\operatorname{Proj}(\mathcal S)$
tel que $\mathcal L=r^*(\mathcal O_P(1))$.
\end{env}

Comme $q$ est un morphisme projectif, il résulte de (II, 5.4.4 et 4.4.7)
qu'il existe un $Y$-isomorphisme
$r':X\xrightarrow{\sim}P'=\operatorname{Proj}(\mathcal T)$, où
$\mathcal T$ est une $\mathcal O_Y$-Algèbre quasi-cohérente,
$\mathcal T_1$ est de type fini et engendre $\mathcal T$, et
$\mathcal L=r'^*(\mathcal O_{P'}(1))$. Si $q':P'\to Y$ est le morphisme
structural, on a
\[
\mathcal S=\bigoplus_{n\geq0}q'_*(\mathcal O_{P'}(n)).
\]
En vertu de (\hyperref[III.2.3.1-fr]{2.3.1}), l'homomorphisme canonique
$\alpha_n:\mathcal T_n\to\mathcal S_n=q'_*(\mathcal O_{P'}(n))$ est
bijectif pour $n$ assez grand~; comme $\mathcal T_n=\mathcal T_1^n$, on a
\emph{a fortiori} $\mathcal S_n=\mathcal S_1^n$ dès que $n$ est assez
grand. En outre, l'homomorphisme canonique
$\alpha:\mathcal T\to\mathcal S$ de $\mathcal O_Y$-Algèbres graduées est
un \textbf{(TN)}-isomorphisme~; par suite
\phantomsection\label{III.2.3.5.1-fr}
\[
\Phi=\operatorname{Proj}(\alpha):
\operatorname{Proj}(\mathcal S)\longrightarrow
\operatorname{Proj}(\mathcal T)
\]
est un isomorphisme (II, 3.6.1). En outre, il résulte de (II, 3.5.2) et
du fait que les $\mathcal T$-Modules gradués obtenus à partir des
$\mathcal S(n)$ par restriction des scalaires sont \textbf{(TN)}-isomorphes
aux $\mathcal T(n)$ que
\[
\Phi_*(\mathcal O_P(n))=\mathcal O_{P'}(n)
\]
pour tout $n$ (II, 3.4.2). Pour achever de prouver
(\hyperref[III.2.3.4.1-fr]{2.3.4.1}), il reste à montrer que $\mathcal S$
est une $\mathcal O_Y$-Algèbre de type fini~; or les
$\mathcal S_n=q'_*(\mathcal O_{P'}(n))$ sont des
$\mathcal O_Y$-Modules cohérents en vertu de
(\hyperref[III.2.2.1-fr]{2.2.1}), et si
$\mathcal S_n=\mathcal S_1^n$ pour $n\geq n_0$, $\mathcal S$ est
engendrée par $\bigoplus_{i\leq n_0}\mathcal S_i$, qui est cohérent~;
d'où notre assertion (I, 9.6.2).

Revenons à la démonstration de
(\hyperref[III.2.3.4-fr]{2.3.4}), dont nous reprenons les notations. Nous
avons démontré l'existence d'un $Y$-isomorphisme $r'':X\xrightarrow{\sim}X''$
tel que $r''_*(\mathcal L^{\otimes n})=\mathcal O_{X''}(n)$ pour tout
$n\in\mathbb Z$~; nous désignerons par $q''$ le morphisme structural
$X''\to Y$. Notons maintenant que $\mathcal S'$ est somme directe des
$\mathcal S'^{(d)}$-Modules gradués $\mathcal S'^{(d,k)}$~; chacun de ces
derniers est un $\mathcal S'^{(d)}$-Module quasi-cohérent de type fini, en
vertu de la périodicité de $\mathcal S'$ et de l'hypothèse que les
$\mathcal S'_n$ sont des $\mathcal O_X$-Modules de type fini
(II, 8.14.12). Posons
$\mathcal F^{(k)}=r''_*(\mathcal S'^{(d,k)})$, de sorte que les
$\mathcal F^{(k)}$ sont des $\mathcal S_{X''}$-Modules gradués
quasi-cohérents de type fini~; par suite (II, 8.14.8), l'homomorphisme
canonique
\[
\beta:\shProj(\bbGamma_*(\mathcal F^{(k)}))
\longrightarrow\mathcal F^{(k)}
\]
est un isomorphisme de $\mathcal S_{X''}$-Modules. Mais on a
$q''_*((\mathcal F^{(k)})_n)=q_*((\mathcal S'^{(d,k)})_n)$, et pour
$n\geq0$ ce dernier $\mathcal O_Y$-Module est par définition égal à
$(\mathcal S^{(d,k)})_n$. Autrement dit, l'injection canonique
$\mathcal S^{(d,k)}\to\bbGamma_*(\mathcal F^{(k)})$ est un
\textbf{(TN)}-isomorphisme~; donc
\[
\shProj(\mathcal S^{(d,k)})
=\shProj(\bbGamma_*(\mathcal F^{(k)})),
\]
et par suite
$r''{}^*(\shProj(\mathcal S^{(d,k)}))\simeq\mathcal S'^{(d,k)}$.
Il reste à remarquer que, à un isomorphisme canonique près,
\[
\shProj(\mathcal S^{(d,k)})
=s_*((\shProj(\mathcal S))^{(d,k)})
\]
(II, 8.14.13.1) pour avoir démontré l'isomorphisme de
$r^*(\shProj(\mathcal S))$ et de $\mathcal S'$.

Enfin, en vertu de (\hyperref[III.2.3.2-fr]{2.3.2}), chacun des
$\bbGamma_*(\mathcal F^{(k)})$ vérifie la condition \textbf{(TF)}, donc il
en est de même de chacun des $\mathcal S^{(d,k)}$~; en outre, les
$\mathcal S_n=q_*(\mathcal S'_n)$ sont cohérents en vertu de
(\hyperref[III.2.2.1-fr]{2.2.1}), puisque les $\mathcal S'_n$ sont
cohérents. On en conclut aussitôt que les $\mathcal S^{(d,k)}$ sont des
$\mathcal S^{(d)}$-Modules de type fini~; comme
(\hyperref[III.2.3.4.1-fr]{2.3.4.1}) montre que $\mathcal S^{(d)}$ est une
$\mathcal O_Y$-Algèbre de type fini, il en est de même de $\mathcal S$.
\end{proof}

\oldpage[III]{106}
\begin{proposition}[2.3.5]
\label{III.2.3.5-fr}
Soient $Y$ un préschéma intègre noethérien, $X$ un préschéma intègre,
$f:X\to Y$ un morphisme projectif birationnel. Il existe alors un Idéal
fractionnaire cohérent $\mathcal J\subset\mathcal R(Y)$ (II, 8.1.2) tel
que $X$ soit $Y$-isomorphe au préschéma obtenu en faisant éclater
$\mathcal J$ (II, 8.1.3). En outre, il existe un ouvert $U$ de $Y$ tel que
la restriction de $f$ à $f^{-1}(U)$ soit un isomorphisme de $f^{-1}(U)$
sur $U$ (cf. I, 6.5.5), et que $\mathcal J|U$ soit inversible.
\end{proposition}

\begin{proof}
Comme il existe un $\mathcal O_X$-Module inversible $\mathcal L$ très
ample pour $f$ (II, 4.4.2 et 5.3.2), on peut appliquer
(\hyperref[III.2.3.4.1-fr]{2.3.4.1}), et on voit que $X$ s'identifie à
$\operatorname{Proj}(\mathcal S)$, où
$\mathcal S=\bigoplus_{n\geq0}f_*(\mathcal L^{\otimes n})$. On sait en
outre que les $f_*(\mathcal L^{\otimes n})$ sont des
$\mathcal O_Y$-Modules sans torsion (I, 7.4.5), donc il en est de même du
$\mathcal O_Y$-Module $\mathcal S$, et par suite $\mathcal S$ s'identifie
canoniquement à un sous-$\mathcal O_Y$-Module de
$\mathcal S\otimes_{\mathcal O_Y}\mathcal R(Y)$ (I, 7.4.1)~; ce dernier
est un faisceau \emph{simple} (I, 7.3.6) qui est connu lorsqu'on connaît
sa restriction à un ouvert non vide, par exemple à un ouvert non vide
$U'\subset U$ tel que $\mathcal L|f^{-1}(U')$ soit isomorphe à
$\mathcal O_X|f^{-1}(U')$. Comme par hypothèse les
$f_*(\mathcal L^{\otimes n})|U'$ sont alors isomorphes à
$\mathcal O_Y|U'$, on voit que
$\mathcal S\otimes_{\mathcal O_Y}\mathcal R(Y)$ est un
$\mathcal R(Y)$-Module isomorphe à $\mathcal R(Y)[T]$, où $T$ est une
indéterminée, et $\mathcal S$ est \textbf{(TN)}-isomorphe à la
sous-$\mathcal O_Y$-Algèbre engendrée par l'image canonique de
$f_*(\mathcal L)$ dans $\mathcal S\otimes_{\mathcal O_Y}\mathcal R(Y)$
(\hyperref[III.2.3.4.1-fr]{2.3.4.1})~; mais si on identifie ce dernier
faisceau à $\mathcal R(Y)[T]$, l'image de $f_*(\mathcal L)$ s'identifie à
$\mathcal J T$, où $\mathcal J$ est un sous-$\mathcal O_Y$-Module
cohérent (\hyperref[III.2.2.1-fr]{2.2.1}) de $\mathcal R(Y)$, dont la
restriction à $U'$ est isomorphe à $\mathcal O_Y|U'$, et qui par suite
est tel que $\mathcal J|U$ soit inversible. On voit alors que $\mathcal S$
est \textbf{(TN)}-isomorphe à $\bigoplus_{n\geq0}\mathcal J^n$, ce qui
achève la démonstration.
\end{proof}

\begin{corollary}[2.3.6]
\label{III.2.3.6-fr}
Sous les hypothèses de (\hyperref[III.2.3.5-fr]{2.3.5}), supposons en
outre que, pour tout sous-$\mathcal O_Y$-Module cohérent
$\mathcal J\neq0$ de $\mathcal R(Y)$, il existe un
$\mathcal O_Y$-Module inversible $\mathcal L$ tel que
\[
\Gamma(Y,\mathcal L\otimes_{\mathcal O_Y}
\mathcal{H}om(\mathcal J,\mathcal O_Y))\neq0~;
\]
alors, dans l'énoncé de (\hyperref[III.2.3.5-fr]{2.3.5}), on peut supposer
que $\mathcal J$ est un Idéal de $\mathcal O_Y$. Cette condition
supplémentaire est toujours vérifiée s'il existe un
$\mathcal O_Y$-Module ample.
\end{corollary}

\begin{proof}
En effet, on a (0\textsubscript{I}, 5.4.2)
\[
\mathcal L\otimes\mathcal{H}om(\mathcal J,\mathcal O_Y)
=\mathcal{H}om(\mathcal L^{-1},\mathcal{H}om(\mathcal J,\mathcal O_Y))
=\mathcal{H}om(\mathcal J\otimes\mathcal L^{-1},\mathcal O_Y)~;
\]
l'hypothèse signifie donc qu'il y a un homomorphisme non nul $u$ de
$\mathcal J\otimes\mathcal L^{-1}$ dans $\mathcal O_Y$. Comme, pour tout
$y\in Y$, $(\mathcal J\otimes\mathcal L^{-1})_y$ s'identifie à un
sous-$\mathcal O_y$-Module du corps des fractions $\mathcal R(Y)_y$ de
$\mathcal O_y$ (I, 7.1.5), $u_y$ est nécessairement injectif, donc $u$
est un isomorphisme de $\mathcal J\otimes\mathcal L^{-1}$ sur un Idéal
$\mathcal J'$ de $\mathcal O_Y$. Mais comme
\[
\operatorname{Proj}\left(\bigoplus_{n\geq0}\mathcal J^n\right)
\quad\text{et}\quad
\operatorname{Proj}\left(
\bigoplus_{n\geq0}(\mathcal J\otimes\mathcal L^{-1})^n\right)
\]
sont $Y$-isomorphes (II, 3.1.8), cela prouve la première assertion du
corollaire. Pour démontrer la seconde, notons que
$\mathcal F=\mathcal{H}om_{\mathcal O_Y}(\mathcal J,\mathcal O_Y)$ est
cohérent et $\neq0$, puisqu'il existe un ouvert $U$ de $Y$ tel que
$\mathcal J|U$ soit inversible. Si $\mathcal L$ est un
$\mathcal O_Y$-Module ample, il existe un entier $n$ tel que
$\mathcal F(n)=\mathcal F\otimes\mathcal L^{\otimes n}$ soit engendré par
ses sections au-dessus de $Y$ (II, 4.5.5)~; \emph{a fortiori}, on a
$\Gamma(Y,\mathcal F(n))\neq0$, d'où la conclusion.
\end{proof}

\begin{corollary}[2.3.7]
\label{III.2.3.7-fr}
Soient $X$ et $Y$ deux schémas intègres, projectifs sur un corps $k$, et
soit $f:X\to Y$ un $k$-morphisme birationnel. Alors $X$ est $k$-isomorphe
à un $Y$-schéma obtenu en faisant éclater un sous-schéma fermé $Y'$
(non nécessairement réduit) de $Y$.
\end{corollary}

\oldpage[III]{107}
\begin{proof}
En effet, $f$ est projectif (II, 5.5.5, (v)) et comme $Y$ est projectif
sur $k$, la condition supplémentaire de
(\hyperref[III.2.3.6-fr]{2.3.6}) est vérifiée~; il suffit alors de
considérer le sous-schéma fermé $Y'$ de $Y$ défini par l'Idéal cohérent
$\mathcal J$ du cor. (\hyperref[III.2.3.6-fr]{2.3.6}).
\end{proof}

\begin{env}[2.3.8]
\label{III.2.3.8-fr}
\emph{Remarque.} Au chap. IV, en étudiant la notion de diviseur, nous
verrons que si, dans l'énoncé de
(\hyperref[III.2.3.5-fr]{2.3.5}), on suppose que les anneaux
$\mathcal O_y$ $(y\in Y)$ sont factoriels (ce qui est le cas par exemple
si $Y$ est non singulier), alors $X$ peut se déduire de $Y$ en faisant
éclater un sous-préschéma fermé $Y'$ de $Y$ dont l'espace sous-jacent est
contenu dans $Y-U$.
\end{env}

\subsection{Une généralisation du théorème fondamental}
\label{subsection:III.2.4-fr}

\begin{theorem}[2.4.1]
\label{III.2.4.1-fr}
Soient $Y$ un préschéma noethérien, $\mathcal S$ une
$\mathcal O_Y$-Algèbre quasi-cohérente de type fini. Soient
$f:X\to Y$ un morphisme projectif, $\mathcal S'=f^*(\mathcal S)$,
$\mathcal M$ un $\mathcal S'$-Module quasi-cohérent de type fini.
Alors~:
\begin{enumerate}
\item[(i)] Pour tout $p\in\mathbb Z$, $R^pf_*(\mathcal M)$ est un
$\mathcal S$-Module de type fini.
\item[(ii)] Soit de plus $\mathcal L$ un $\mathcal O_X$-Module inversible
ample pour $f$, et posons
$\mathcal M(n)=\mathcal M\otimes\mathcal L^{\otimes n}$ pour tout
$n\in\mathbb Z$. Il existe un entier $N$ tel que, pour $n\geq N$, on ait
\[
R^pf_*(\mathcal M(n))=0
\tag{2.4.1.1}\label{III.2.4.1.1-fr}
\]
pour tout $p>0$, et que l'homomorphisme canonique
\[
f^*(f_*(\mathcal M(n)))\longrightarrow\mathcal M(n)
\]
(0\textsubscript{I}, 4.4.3) soit surjectif.
\end{enumerate}
\end{theorem}

\begin{proof}
Posons $Y'=\operatorname{Spec}(\mathcal S)$,
$X'=\operatorname{Spec}(\mathcal S')$ de sorte que
$X'=X\times_Y Y'$ (II, 1.5.5)~; soient $g:Y'\to Y$,
$g':X'\to X$ les morphismes structuraux, qui sont affines par
définition, et $f'=f_{(Y')}:X'\to Y'$~; on a donc un diagramme
commutatif
\[
\xymatrix{
X\ar[d]_f & X'\ar[l]_{g'}\ar[d]^{f'}\ar[dl]^{h}\\
Y & Y'\ar[l]^{g}
}
\]
et le morphisme $f'$ est projectif (II, 5.5.5, (iii))~; posons
$h=f\circ g'=g\circ f'$.

(i) Soit $\widetilde{\mathcal M}$ le $\mathcal O_{X'}$-Module associé au
$\mathcal S'$-Module quasi-cohérent $\mathcal M$, quand $X'$ est
considéré comme un $X$-schéma affine (II, 1.4.3), de sorte que l'on a
$\mathcal M=g'_*(\widetilde{\mathcal M})$~; comme $\mathcal M$ est un
$\mathcal S'$-Module de type fini, $\widetilde{\mathcal M}$ est un
$\mathcal O_{X'}$-Module de type fini (II, 1.4.5)~; comme $h$ est de type
fini, puisque $g$ et $f'$ le sont (II, 1.3.7 et I, 6.3.4, (ii)), $X'$ est
noethérien (I, 6.3.7) et $\widetilde{\mathcal M}$ est par suite cohérent.
Cela étant, comme $g'$ est affine, l'homomorphisme canonique
$R^pf_*(\mathcal M)\to R^ph_*(\widetilde{\mathcal M})$ est bijectif
(\hyperref[III.1.3.4-fr]{1.3.4}). En outre, cet homomorphisme est un
homomorphisme de $\mathcal S$-Modules~; en effet, de l'homomorphisme
canonique
\[
g'_*(\mathcal O_{X'})\otimes_{\mathcal O_X}
g'_*(\widetilde{\mathcal M})\longrightarrow
g'_*(\widetilde{\mathcal M})
\tag{2.4.1.2}\label{III.2.4.1.2-fr}
\]
qui définit la structure de $\mathcal S'$-Module de $\mathcal M$ (en se
rappelant que $\mathcal S'=g'_*(\mathcal O_{X'})$), on déduit
canoniquement un homomorphisme
\[
f_*(g'_*(\mathcal O_{X'}))\otimes
R^pf_*(g'_*(\widetilde{\mathcal M}))\longrightarrow
R^pf_*(g'_*(\widetilde{\mathcal M})).
\]

\oldpage[III]{108}
(0, 12.2.2), et comme (\hyperref[III.2.4.1.2-fr]{2.4.1.2}) provient
lui-même (par application de (0\textsubscript{I}, 4.2.2.1)) de
l'homomorphisme
\[
\mathcal O_{X'}\otimes\widetilde{\mathcal M}
\longrightarrow\widetilde{\mathcal M}
\]
définissant la structure de $\mathcal O_{X'}$-Module de
$\widetilde{\mathcal M}$, le diagramme
\[
\xymatrix{
f_*(g'_*(\mathcal O_{X'}))\otimes
R^pf_*(g'_*(\widetilde{\mathcal M}))
\ar[r]\ar[d] &
R^pf_*(g'_*(\widetilde{\mathcal M}))\ar[d]\\
h_*(\mathcal O_{X'})\otimes R^ph_*(\widetilde{\mathcal M})
\ar[r] & R^ph_*(\widetilde{\mathcal M})
}
\]
est commutatif (0, 12.2.6)~; composant les flèches horizontales avec
l'homomorphisme provenant de l'homomorphisme canonique
\[
\mathcal S\longrightarrow f_*(f^*(\mathcal S))=f_*(\mathcal S')
=f_*(g'_*(\mathcal O_{X'}))=h_*(\mathcal O_{X'}),
\]
on obtient notre assertion. D'autre part, puisque $g$ est affine et que
$f'$ est séparé et quasi-compact, l'homomorphisme canonique
$R^ph_*(\widetilde{\mathcal M})\to
g_*(R^pf'_*(\widetilde{\mathcal M}))$ est bijectif
(\hyperref[III.1.4.14-fr]{1.4.14}), et on démontre comme ci-dessus que
c'est un isomorphisme de $\mathcal S$-Modules (en utilisant cette fois la
commutativité de (0, 12.2.6.2)). Or, $f'$ étant projectif et
$\widetilde{\mathcal M}$ cohérent,
$R^pf'_*(\widetilde{\mathcal M})$ est un
$\mathcal O_{Y'}$-Module cohérent en vertu de
(\hyperref[III.2.2.1-fr]{2.2.1})~; on en conclut que
$g_*(R^pf'_*(\widetilde{\mathcal M}))$ est un $\mathcal S$-Module de type
fini (II, 1.4.5).

(ii) Soit $\mathcal L'=g'^*(\mathcal L)$, qui est un
$\mathcal O_{X'}$-Module inversible~; pour tout $n\in\mathbb Z$, on a
\[
g'_*(\widetilde{\mathcal M}\otimes\mathcal L'^{\otimes n})
=g'_*(\widetilde{\mathcal M})\otimes\mathcal L^{\otimes n}
=\mathcal M(n)
\]
(0\textsubscript{I}, 5.4.10) à un isomorphisme près~; on peut appliquer
à $\widetilde{\mathcal M}\otimes\mathcal L'^{\otimes n}$ le
raisonnement fait dans (i) pour $\widetilde{\mathcal M}$, qui prouve que
\[
R^pf_*(g'_*(\widetilde{\mathcal M}\otimes\mathcal L'^{\otimes n}))
\quad\hbox{est isomorphe à}\quad
g_*(R^pf'_*(\widetilde{\mathcal M}\otimes\mathcal L'^{\otimes n})).
\]
Or $\mathcal L'$ est ample pour $f'$ (II, 4.6.13, (iii)), donc il résulte
de (\hyperref[III.2.2.1-fr]{2.2.1}) qu'il existe un entier $N$ tel que
\[
R^pf'_*(\widetilde{\mathcal M}\otimes\mathcal L'^{\otimes n})=0
\]
pour tout $p$ et tout $n\geq N$, ce qui prouve
(\hyperref[III.2.4.1.1-fr]{2.4.1.1}). Enfin, il résulte encore de
(\hyperref[III.2.2.1-fr]{2.2.1}) qu'on peut supposer $N$ tel que pour
$n\geq N$, l'homomorphisme canonique
\[
f'^*(f'_*(\widetilde{\mathcal M}\otimes\mathcal L'^{\otimes n}))
\longrightarrow\widetilde{\mathcal M}\otimes\mathcal L'^{\otimes n}
\]
soit surjectif~; comme $g'_*$ est un foncteur exact (II, 1.4.4),
l'homomorphisme correspondant
\[
g'_*(f'^*(f'_*(\widetilde{\mathcal M}\otimes
\mathcal L'^{\otimes n})))
\longrightarrow
g'_*(\widetilde{\mathcal M}\otimes\mathcal L'^{\otimes n})
=\mathcal M(n)
\]
est surjectif. Or, on a
\[
g'_*(f'^*(f'_*(\widetilde{\mathcal M}\otimes
\mathcal L'^{\otimes n})))
=f^*(g_*(f'_*(\widetilde{\mathcal M}\otimes
\mathcal L'^{\otimes n})))
\]
(II, 1.5.2) et comme $g_*\circ f'_* = f_*\circ g'_*$, on voit
finalement que l'on a
\[
g'_*(f'^*(f'_*(\widetilde{\mathcal M}\otimes\mathcal L'^{\otimes n})))
=f^*(f_*(g'_*(\widetilde{\mathcal M}\otimes
\mathcal L'^{\otimes n})))=f^*(f_*(\mathcal M(n))),
\]
ce qui achève la démonstration.
\end{proof}

\begin{env}[2.4.2]
\label{III.2.4.2-fr}
Nous aurons en particulier à appliquer
(\hyperref[III.2.4.1-fr]{2.4.1}) lorsque $\mathcal S$ est une
$\mathcal O_Y$-Algèbre graduée à degrés positifs,
$\mathcal M=\bigoplus_{k\in\mathbb Z}\mathcal M_k$ un
$\mathcal S'$-Module gradué. Alors (avec les

\oldpage[III]{109}
mêmes hypothèses de finitude sur $\mathcal S$ et $\mathcal M$) on conclut
de (\hyperref[III.2.4.1-fr]{2.4.1}) que
\[
\bigoplus_{k\in\mathbb Z}R^pf_*(\mathcal M_k)
\]
est un $\mathcal S$-Module de type fini pour tout $p$, et (sous les
hypothèses supplémentaires de (\hyperref[III.2.4.1-fr]{2.4.1}), (ii))
qu'il existe $N$ tel que pour $n\geq N$, on ait
$R^pf_*(\mathcal M_k(n))=0$ pour tout $p>0$ et tout $k\in\mathbb Z$, et
que l'homomorphisme canonique
$f^*(f_*(\mathcal M_k(n)))\to\mathcal M_k(n)$ soit surjectif pour tout
$k\in\mathbb Z$.
\end{env}

\subsection{Caractéristique d'Euler-Poincaré et polynôme de Hilbert}
\label{subsection:III.2.5-fr}

\begin{env}[2.5.1]
\label{III.2.5.1-fr}
Soient $A$ un anneau artinien, $X$ un $A$-schéma projectif sur
$Y=\operatorname{Spec}(A)$. Pour tout $\mathcal O_X$-Module cohérent
$\mathcal F$, les $H^i(X,\mathcal F)$ $(i\geq0)$ sont des $A$-modules de
type fini (\hyperref[III.2.2.1-fr]{2.2.1}), donc ici de longueur finie
puisque $A$ est artinien. On sait en outre
(\hyperref[III.2.2.1-fr]{2.2.1}) que $H^i(X,\mathcal F)=0$ sauf pour un
nombre fini de valeurs de $i\geq0$~; le nombre entier
\[
\chi_A(\mathcal F)=\sum_{i=0}^{\infty}(-1)^i
\operatorname{long}(H^i(X,\mathcal F))
\tag{2.5.1.1}\label{III.2.5.1.1-fr}
\]
est donc défini pour tout $\mathcal O_X$-Module cohérent $\mathcal F$.
Lorsque $A$ est un anneau local artinien, on dit que $\chi_A(\mathcal F)$
est la caractéristique d'Euler-Poincaré de $\mathcal F$ (par rapport à
l'anneau $A$). Pour $\mathcal F=\mathcal O_X$, on dit que
$\chi_A(\mathcal O_X)$ est le genre arithmétique de $X$ (par rapport à
$A$).
\end{env}

\begin{proposition}[2.5.2]
\label{III.2.5.2-fr}
Soit
\[
0\longrightarrow\mathcal F'\longrightarrow\mathcal F
\longrightarrow\mathcal F''\longrightarrow0
\]
une suite exacte de $\mathcal O_X$-Modules cohérents~; on a alors
\[
\chi_A(\mathcal F)=\chi_A(\mathcal F')+\chi_A(\mathcal F'').
\tag{2.5.2.1}\label{III.2.5.2.1-fr}
\]
\end{proposition}

\begin{proof}
Comme les modules de cohomologie de $\mathcal F$, $\mathcal F'$,
$\mathcal F''$ sont nuls sauf un nombre fini d'entre eux, il y a un
entier $r>0$ tel que la suite exacte de cohomologie s'écrive
\[
0\to H^0(X,\mathcal F')\to H^0(X,\mathcal F)
\to H^0(X,\mathcal F'')\to H^1(X,\mathcal F')\to\cdots
\]
\[
\cdots\to H^r(X,\mathcal F')\to H^r(X,\mathcal F)
\to H^r(X,\mathcal F'')\to0.
\]
Or, on sait que dans une suite exacte de $A$-modules de longueur finie,
ayant des 0 aux deux extrémités, la somme alternée des longueurs est
nulle (0, 11.10.1)~; appliquant ce résultat, on trouve immédiatement la
formule (\hyperref[III.2.5.2.1-fr]{2.5.2.1}).
\end{proof}

On notera que le résultat de (\hyperref[III.2.5.2-fr]{2.5.2}) s'applique
toutes les fois que l'on sait que $X$ est un $A$-schéma quasi-compact et
que les $A$-modules $H^i(X,\mathcal F)$ sont de type fini pour tout
$\mathcal O_X$-Module cohérent $\mathcal F$
(\hyperref[III.1.4.12-fr]{1.4.12}).

\begin{theorem}[2.5.3]
\label{III.2.5.3-fr}
Soient $A$ un anneau local artinien, $X$ un schéma projectif sur
$Y=\operatorname{Spec}(A)$, $\mathcal L$ un $\mathcal O_X$-Module
inversible très ample relativement à $Y$, $\mathcal F$ un
$\mathcal O_X$-Module cohérent~; on pose
$\mathcal F(n)=\mathcal F\otimes_{\mathcal O_X}\mathcal L^{\otimes n}$
pour tout $n\in\mathbb Z$.
\begin{enumerate}
\item[(i)] Il existe un polynôme unique $P\in\mathbb Q[T]$ tel que
$\chi_A(\mathcal F(n))=P(n)$ pour tout $n\in\mathbb Z$ (on dit que $P$
est le polynôme de Hilbert de $\mathcal F$ par rapport à $A$).
\item[(ii)] Pour $n$ assez grand, on a
$\chi_A(\mathcal F(n))=\operatorname{long}_A
\Gamma(X,\mathcal F(n))$.
\item[(iii)] Le terme de plus haut degré de $\chi_A(\mathcal F(n))$ a un
coefficient $\geq0$.
\end{enumerate}
\end{theorem}

Ajoutons qu'au chap. IV, dans le paragraphe consacré à la notion de
dimension, nous démontrerons en outre que le degré de
$\chi_A(\mathcal F(n))$ est égal à la dimension du support de
$\mathcal F$.

\begin{proof}
Comme on a $H^i(X,\mathcal F(n))=0$ pour tout $i>0$ dès que $n$ est
assez grand (\hyperref[III.2.2.1-fr]{2.2.1}),

\oldpage[III]{110}
on a
\[
\chi_A(\mathcal F(n))=\operatorname{long}H^0(X,\mathcal F(n))
=\operatorname{long}\Gamma(X,\mathcal F(n))
\]
pour $n$ assez grand, d'où (ii)~; cela entraîne
$\chi_A(\mathcal F(n))\geq0$ pour $n$ assez grand, et (iii) résulte donc
de (i)~; comme d'ailleurs l'assertion d'unicité de (i) est immédiate, il
reste à prouver l'existence du polynôme $P$.

Montrons d'abord qu'on peut supposer que $\mathfrak m\mathcal F=0$, où
$\mathfrak m$ est l'idéal maximal de $A$. En effet, il existe un entier
$s>0$ tel que $\mathfrak m^s=0$, et $\mathcal F(n)$ admet donc une
filtration finie
\[
\mathcal F(n)\supset\mathfrak m\mathcal F(n)\supset\cdots
\supset\mathfrak m^{s-1}\mathcal F(n)\supset0.
\]
Par récurrence, on déduit de
(\hyperref[III.2.5.2.1-fr]{2.5.2.1}) que
\[
\chi_A(\mathcal F(n))=\sum_{k=1}^{s}
\chi_A(\mathfrak m^{k-1}\mathcal F(n)/\mathfrak m^k\mathcal F(n))~;
\]
comme
$\mathfrak m^{k-1}\mathcal F(n)/\mathfrak m^k\mathcal F(n)
=\mathcal F'_k(n)$, où
$\mathcal F'_k=\mathfrak m^{k-1}\mathcal F/\mathfrak m^k\mathcal F$,
cela prouve notre assertion.

Supposons donc $\mathfrak m\mathcal F=0$~; si $X'$ est le sous-schéma
fermé de $X$, image réciproque par le morphisme structural
$X\to\operatorname{Spec}(A)$ de l'unique point fermé de
$\operatorname{Spec}(A)$, et $j:X'\to X$ l'injection canonique, on a
$\mathcal F=j_*(\mathcal F')$, où $\mathcal F'$ est un
$\mathcal O_{X'}$-Module cohérent~; $X'$ est un schéma projectif sur
$\operatorname{Spec}(K)$, où $K=A/\mathfrak m$. Si
$\mathcal L'=j^*(\mathcal L)$, $\mathcal L'$ est très ample relativement
à $\operatorname{Spec}(K)$ (II, 4.4.10), et on a
$\mathcal F(n)=j_*(\mathcal F'(n))$, où
$\mathcal F'(n)=\mathcal F'\otimes_{\mathcal O_{X'}}
\mathcal L'^{\otimes n}$ (0\textsubscript{I}, 5.4.10). On en conclut que
$\chi_A(\mathcal F(n))=\chi_K(\mathcal F'(n))$ (G, II, 4.9.1), et on est
donc ramené au cas où $A$ est un corps.

Notons maintenant que $X$ peut être considéré comme un sous-schéma
fermé de $P=\mathbf P_A^r$ pour un $r$ convenable (II, 5.5.4, (ii))~;
si $i:X\to P$ est l'injection canonique, on voit comme ci-dessus que l'on
a $\chi_A(\mathcal F(n))=\chi_A(i_*(\mathcal F)(n))$, de sorte que l'on
peut se borner au cas où $X=\mathbf P_A^r=\operatorname{Proj}(S)$ avec
$S=A[T_0,\ldots,T_r]$, $A$ étant un corps.

Cela étant, on a $\mathcal F=\widetilde M$, où $M$ est un $S$-module
gradué de type fini (II, 2.7.8)~; il existe par suite une résolution finie
de $M$ par des $S$-modules gradués libres de type fini
\[
0\to L_q\to L_{q-1}\to\cdots\to L_1\to M\to0
\]
en vertu du th. des syzygies de Hilbert (M, VIII, 6.5)~; comme
$M\mapsto\widetilde M$ est un foncteur exact en $M$ (II, 2.5.4), on a
aussi une suite exacte
\[
0\to\widetilde L_q\to\widetilde L_{q-1}\to\cdots
\to\widetilde L_1\to\widetilde M\to0
\]
et par suite, pour tout $n\in\mathbb Z$, la suite
\[
0\to\widetilde L_q(n)\to\widetilde L_{q-1}(n)\to\cdots
\to\widetilde L_1(n)\to\widetilde M(n)\to0
\]
est exacte~; appliquant par récurrence sur $q$ la prop. (2.5.1), il vient
\[
\chi_A(\widetilde M(n))=\sum_{j=1}^{q}(-1)^{j+1}
\chi_A(\widetilde L_j(n))
\]
et pour prouver (i), on est donc ramené au cas où $M$ est libre et gradué
de type fini, donc au cas où $M=S(h)$ pour un $h\in\mathbb Z$. Comme
alors $\widetilde M(n)=(M(n))^\sim=(S(n+h))^\sim$ (II, 2.5.15), on voit
finalement que le théorème résultera du
\end{proof}

\oldpage[III]{111}
\begin{lemma}[2.5.3.1]
\label{III.2.5.3.1-fr}
Soient $A$ un corps, $r$ un entier $>0$, et $X=\mathbf P_A^r$~; on a
alors
\[
\chi_A(\mathcal O_X(n))=\binom{n+r}{r}
\]
pour tout $n\in\mathbb Z$.
\end{lemma}

\begin{proof}
En effet, pour $n\geq0$, on a
$\chi_A(\mathcal O_X(n))=\operatorname{long}H^0(X,\mathcal O_X(n))$, qui
est le nombre de monômes en les $T_i$ de degré total $n$, c'est-à-dire
$\binom{n+r}{r}$ (\hyperref[III.2.1.12-fr]{2.1.12}). Pour
$n\leq-r-1$, on a de même
\[
\chi_A(\mathcal O_X(n))=(-1)^r
\operatorname{long}H^r(X,\mathcal O_X(n))~;
\]
si $n=-r-h$, la dimension de $H^r(X,\mathcal O_X(n))$ sur $A$ est le
nombre des suites $(p_i)_{0\leq i\leq r}$ d'entiers $p_i>0$ tels que
$\sum_{i=0}^{r}p_i=r+h$ (\hyperref[III.2.1.12-fr]{2.1.12}), ou encore le
nombre des suites d'entiers $q_i\geq0$ $(0\leq i\leq r)$ tels que
$\sum_{i=0}^{r}q_i=h-1$~; c'est donc le nombre
\[
\binom{h+r-1}{r}=(-1)^r\binom{n+r}{r}.
\]
Enfin, pour $-r\leq n<0$, on a $\binom{n+r}{r}=0$ et d'autre part
$H^i(X,\mathcal O_X(n))=0$ pour tout $i\geq0$
(\hyperref[III.2.1.12-fr]{2.1.12}), ce qui prouve le lemme.
\end{proof}

\begin{corollary}[2.5.4]
\label{III.2.5.4-fr}
Soient $A$ un anneau local artinien, $S$ une $A$-algèbre graduée de type
fini engendrée par $S_1$, $M$ un $S$-module gradué de type fini,
$X=\operatorname{Proj}(S)$. On a alors
\[
\chi_A(\widetilde M(n))=\operatorname{long}M_n
\]
pour $n$ assez grand.
\end{corollary}

\begin{proof}
Cela résulte de ce que $M_n$ et $\Gamma(X,\widetilde M(n))$ sont
isomorphes pour $n$ assez grand
(\hyperref[III.2.3.1-fr]{2.3.1}).
\end{proof}

\subsection{Application : critères d'amplitude}
\label{subsection:III.2.6-fr}

\begin{proposition}[2.6.1]
\label{III.2.6.1-fr}
Soient $Y$ un préschéma noethérien, $f:X\to Y$ un morphisme propre,
$\mathcal L$ un $\mathcal O_X$-Module inversible. Les conditions
suivantes sont équivalentes~:
\begin{enumerate}
\item[a)] $\mathcal L$ est ample pour $f$.
\item[b)] Pour tout $\mathcal O_X$-Module cohérent $\mathcal F$, il
existe un entier $N$ tel que pour $n\geq N$, on ait
\[
R^qf_*(\mathcal F\otimes\mathcal L^{\otimes n})=0
\]
pour tout $q>0$.
\item[c)] Pour tout Idéal cohérent $\mathcal J$ de $\mathcal O_X$, il
existe un entier $N$ tel que pour $n\geq N$ on ait
\[
R^1f_*(\mathcal J\otimes\mathcal L^{\otimes n})=0.
\]
\end{enumerate}
\end{proposition}

\begin{proof}
On a vu que a) entraîne b) (\hyperref[III.2.2.1-fr]{2.2.1}, (ii)). Il
est trivial que b) entraîne c), et il reste à prouver que c) implique
a). On peut se borner au cas où $Y$ est affine (II, 4.6.4), et prouver
dans ce cas que $\mathcal L$ est ample~; il suffira de montrer que lorsque
$h$ parcourt l'ensemble des sections des $\mathcal L^{\otimes n}$
$(n>0)$ au-dessus de $X$, ceux des $X_h$ qui sont affines forment un
recouvrement de $X$ (II, 4.5.2). Pour cela, montrons que pour tout point
fermé $x$ de $X$ et tout voisinage ouvert affine $U$ de $x$, il existe un
$n$ et un $h\in\Gamma(X,\mathcal L^{\otimes n})$ tels que
$x\in X_h\subset U$~; $X_h$ sera nécessairement affine (I, 1.3.6) et la
réunion de ces $X_h$ sera un ouvert de $X$ contenant tous les points
fermés de $X$, et par suite $X$ lui-même puisque $X$ est noethérien
(I, 6.3.7 et 0\textsubscript{I}, 2.1.3). Soit $\mathcal J$ (resp.
$\mathcal J'$) le faisceau quasi-cohérent d'idéaux de $\mathcal O_X$
définissant le sous-préschéma fermé réduit de $X$ ayant pour espace
sous-jacent $X-U$ (resp. $(X-U)\cup\{x\}$) (I, 5.2.1)~; il est clair que
$\mathcal J$ et $\mathcal J'$ sont cohérents (I, 6.1.1), que
$\mathcal J'\subset\mathcal J$ et que
$\mathcal J''=\mathcal J/\mathcal J'$ est un $\mathcal O_X$-Module
cohérent (0\textsubscript{I}, 5.3.3) ayant pour support $\{x\}$ et tel
que $\mathcal J''_x=k(x)$. Comme $\mathcal L$ est localement libre, la
suite
\[
0\to\mathcal J'\otimes\mathcal L^{\otimes n}
\to\mathcal J\otimes\mathcal L^{\otimes n}
\to\mathcal J''\otimes\mathcal L^{\otimes n}\to0
\]
est exacte pour tout $n$, et par hypothèse il existe $n$ assez grand tel
que $H^1(X,\mathcal J'\otimes\mathcal L^{\otimes n})=0$~; la suite
exacte de cohomologie

\oldpage[III]{112}
prouve donc que l'homomorphisme
\[
\Gamma(X,\mathcal J\otimes\mathcal L^{\otimes n})
\longrightarrow
\Gamma(X,\mathcal J''\otimes\mathcal L^{\otimes n})
\]
est surjectif. Une section $g$ de
$\mathcal J''\otimes\mathcal L^{\otimes n}$ au-dessus de $X$ telle que
$g(x)\neq0$ est donc l'image d'une section
$h\in\Gamma(X,\mathcal J\otimes\mathcal L^{\otimes n})
\subset\Gamma(X,\mathcal L^{\otimes n})$ (car en vertu de
(0\textsubscript{I}, 5.4.1),
$\mathcal J\otimes\mathcal L^{\otimes n}$ est un sous-$\mathcal O_X$-Module
de $\mathcal L^{\otimes n}$)~; on a par définition $h(x)\neq0$ et
$h(z)=0$ pour $z\notin U$, ce qui achève la démonstration.
\end{proof}

\begin{proposition}[2.6.2]
\label{III.2.6.2-fr}
Soient $Y$ un préschéma noethérien, $f:X\to Y$ un morphisme de type
fini, $g:X'\to X$ un morphisme fini surjectif, $\mathcal L$ un
$\mathcal O_X$-Module inversible et $\mathcal L'=g^*(\mathcal L)$. On
suppose vérifiée la condition suivante~: Il existe une partie $Z$ de
$X$, propre sur $Y$ (II, 5.4.10) telle que pour tout $x\in X-Z$, ou bien
$X$ est normal au point $x$, ou bien $(g_*(\mathcal O_{X'}))_x$ est un
$\mathcal O_x$-module libre. Dans ces conditions, pour que $\mathcal L$
soit ample pour $f$, il faut et il suffit que $\mathcal L'$ soit ample
pour $f\circ g$.
\end{proposition}

\begin{proof}
\begin{env}[2.6.2.1]
\label{III.2.6.2.1-fr}
Puisque $g$ est affine, la condition est nécessaire (II, 5.1.12). Pour
voir qu'elle est suffisante, on peut supposer $Y$ affine (II, 4.6.4).
Montrons en outre qu'on peut se borner au cas où $X$ est réduit. En
effet, soit $j:X_{\mathrm{red}}\to X$ l'injection canonique, et posons
$X_1=X_{\mathrm{red}}$, $X'_1=X'\times_X X_1$, de sorte que l'on a le
diagramme commutatif
\[
\xymatrix{
X'\ar[d]_g & X'_1\ar[l]_{j'}\ar[d]^{g_1\, .}\\
X & X_1\ar[l]^{j}
}
\tag{2.6.2.2}\label{III.2.6.2.2-fr}
\]
Le morphisme $f\circ j$ est alors de type fini (I, 6.3.4) et $g_1$ est
un morphisme fini (II, 6.1.5, (iii))~; si $\mathcal L'$ est ample pour
$f\circ g$, $j'^*(\mathcal L')$ est ample pour $f\circ g\circ j'$ puisque
$j'$ est une immersion fermée (II, 5.1.12 et I, 4.3.2). Si on pose
$Z_1=j^{-1}(Z)$, $Z_1$ est propre sur $Y$ (II, 5.4.10)~; d'autre part,
si $X$ est normal en un point $x$, il en est évidemment de même de
$X_{\mathrm{red}}$~; enfin, si $(g_*(\mathcal O_{X'}))_x$ est un
$\mathcal O_x$-module libre, il résulte aussitôt de (II, 1.5.2) que
$((g_1)_*(\mathcal O_{X'_1}))_x$ est un $\mathcal O_{X_1,x}$-module
libre. Enfin, comme $X$ est noethérien (I, 6.3.7), si
$j^*(\mathcal L)$ est ample, $\mathcal L$ est ample (II, 4.5.14), et
comme $j'^*(\mathcal L')=g_1^*(j^*(\mathcal L))$, cela achève la
réduction annoncée. Nous supposons donc désormais $Y$ affine et $X$
réduit.
\end{env}

Les hypothèses de (II, 6.6.11) étant alors vérifiées, il existe un
$Y$-préschéma réduit $X_2$ et un $Y$-morphisme $h:X_2\to X$ fini et
birationnel tel que la restriction de $h$ à $h^{-1}(X-Z)$ soit un
isomorphisme sur $X-Z$ et que $h^*(\mathcal L)$ soit ample. Remplaçant
$X'$ par $X_2$, on voit qu'on est ramené à démontrer la proposition en
supposant en outre que $g$ possède les propriétés qui viennent d'être
énumérées pour $h$. Nous désignerons encore par $Z$ un sous-préschéma de
$X$ ayant $Z$ pour espace sous-jacent, qui est propre sur $Y$
(II, 5.4.10).

\begin{env}[2.6.2.3]
\label{III.2.6.2.3-fr}
Soient maintenant $X_1$ un sous-préschéma fermé de $X$,
$j:X_1\to X$ l'injection canonique,
$X'_1=g^{-1}(X_1)=X'\times_X X_1$ son image réciproque,
$j':X'_1\to X'$ l'injection canonique, de sorte que l'on a le diagramme
commutatif (\hyperref[III.2.6.2.2-fr]{2.6.2.2})~; posons
$\mathcal L_1=j^*(\mathcal L)$,
$\mathcal L'_1=j'^*(\mathcal L')=g_1^*(\mathcal L_1)$, de sorte que
$\mathcal L'_1$ est ample pour $f\circ g\circ j'$ (II, 5.1.12). Si on
pose $Z_1=j^{-1}(Z)$, le sous-préschéma fermé $Z_1$ de $X_1$ est propre
sur $Y$ (II, 5.4.2, (ii)). En d'autres termes, les hypothèses de
(\hyperref[III.2.6.2-fr]{2.6.2}) sont vérifiées pour $X_1$,
$\mathcal L_1$, $g_1$ et $Z_1$.
\end{env}

\oldpage[III]{113}
Ceci va nous permettre de démontrer
(\hyperref[III.2.6.2-fr]{2.6.2}) par récurrence noethérienne
(0\textsubscript{I}, 2.2.2) dans le cas où la restriction de $g$ à
$g^{-1}(X-Z)$ est un isomorphisme sur $X-Z$ (ce qui est suffisant pour
notre propos, comme on l'a vu en
(\hyperref[III.2.6.2.1-fr]{2.6.2.1}))~: il suffira d'établir que si, pour
tout sous-préschéma fermé $X_1$ de $X$, dont l'espace sous-jacent est
$\neq X$, la conclusion de (\hyperref[III.2.6.2-fr]{2.6.2}) est vraie
pour le faisceau $\mathcal L_1$, alors elle est vraie aussi pour le
faisceau $\mathcal L$.

\begin{env}[2.6.2.4]
\label{III.2.6.2.4-fr}
Soient alors $\mathcal A=\mathcal O_X$, $\mathcal B=g_*(\mathcal O_{X'})$,
de sorte que $\mathcal B$ est une sous-$\mathcal A$-Algèbre de
$\mathcal R(X)$, qui est un $\mathcal A$-Module cohérent~; en outre, la
restriction $\mathcal B|(X-Z)$ est égale à $\mathcal A|(X-Z)$. Soit
$\mathcal K$ le conducteur de $\mathcal B$ sur $\mathcal A$, c'est-à-dire
le plus grand sous-$\mathcal A$-Module quasi-cohérent de $\mathcal A$ tel
que $\mathcal B\mathcal K\subset\mathcal A$ (ou encore l'annulateur du
$\mathcal A$-Module $\mathcal B/\mathcal A$
(0\textsubscript{I}, 5.3.7)), ce qui entraîne
$\mathcal B\mathcal K=\mathcal K$. Il est clair que
$\mathcal K_x=\mathcal A_x$ en tous les points admettant un voisinage
$W_x$ tel que $g$ soit un isomorphisme de $g^{-1}(W_x)$ sur $W_x$, et en
particulier en tous les points de $X-Z$ et dans un voisinage de tout
point générique d'une composante irréductible de $X$. Considérons alors
le sous-préschéma fermé $Z_1=\operatorname{Spec}(\mathcal A/\mathcal K)$
de $X$ défini par $\mathcal K$~; il est encore propre sur $Y$ car le
sous-espace $Z_1$ est fermé dans $Z$ (II, 5.4.10). De plus, la définition
de $\mathcal K$ montre que
$\mathcal B|(X-Z_1)=\mathcal A|(X-Z_1)$~; on voit donc qu'on peut toujours
se ramener au cas où
$Z=\operatorname{Spec}(\mathcal A/\mathcal K)$, et comme on a vu que
$X-Z_1$ est un ouvert non vide de $X$, on peut toujours supposer que
l'espace $Z$ est distinct de $X$.
\end{env}

\begin{env}[2.6.2.5]
\label{III.2.6.2.5-fr}
Considérons $X'$ comme égal à $\operatorname{Spec}(\mathcal B)$ (puisque
$g$ est affine) et soit $\mathcal K'=\widetilde{\mathcal K}$, Idéal
cohérent de $\mathcal O_{X'}$ tel que $g_*(\mathcal K')=\mathcal K$
(II, 1.4.1)~; le sous-préschéma fermé
$Z'=g^{-1}(Z)=Z\times_X X'$ de $X'$ est défini par $\mathcal K'$ et égal
à $\operatorname{Spec}(\mathcal B/\mathcal K)$ (II, 1.4.10)~; comme
$h:Z'\to Z$ est un morphisme fini (II, 6.1.5, (iii)), $Z'$ est propre
sur $Y$ (II, 6.1.11 et 5.4.2, (ii)).

Cela posé, il nous faut prouver que pour tout $x\in X$ et tout voisinage
ouvert $U$ de $x$, il existe une section $s$ d'un
$\mathcal L^{\otimes n}$ $(n>0)$ au-dessus de $X$ telle que
$x\in X_s\subset U$ (II, 4.5.2)~; nous distinguerons deux cas~:

$1^\circ$ On a $x\in X-Z$~; on peut évidemment supposer alors que l'on a
aussi $U\subset X-Z$, donc l'ouvert $U'=g^{-1}(U)$ ne rencontre pas
$Z'$. Comme $\mathcal L'$ est ample par hypothèse, il existe un $n>0$ et
une section $s'$ de $\mathcal L'^{\otimes n}$ au-dessus de $X'$ telle
que $x'=g^{-1}(x)\in X'_{s'}\subset g^{-1}(U)$ (II, 4.5.2). En outre, on
peut supposer que $\mathcal K'\otimes\mathcal L'^{\otimes n}$ soit
engendré par ses sections au-dessus de $X'$ (II, 4.5.5), donc, comme
$\mathcal K'_{x'}=\mathcal O_{x'}$, il y a une de ces sections $s''$
telle que $s''(x')\neq0$~; en la multipliant par $s'$ (ce qui revient à
remplacer $n$ par $2n$), on voit qu'on peut supposer aussi que
$x'\in X'_{s''}\subset g^{-1}(U)$. Cela étant, il résulte de
(0\textsubscript{I}, 5.4.10) que l'on a un isomorphisme canonique
\[
\Gamma(X,\mathcal K\otimes\mathcal L^{\otimes n})
\xrightarrow{\sim}
\Gamma(X',\mathcal K'\otimes\mathcal L'^{\otimes n}).
\]
La section $s$ de $\mathcal K\otimes\mathcal L^{\otimes n}$ qui
correspond à $s''$ par cet isomorphisme a évidemment les propriétés
voulues.

$2^\circ$ On a $x\in Z$. Soit $\mathcal J$ l'Idéal cohérent de
$\mathcal O_X$ définissant le sous-préschéma fermé réduit de $X$ ayant
pour espace sous-jacent $X-U$, et considérons dans $\mathcal B$ les deux

\oldpage[III]{114}
Idéaux cohérents $\mathcal J\mathcal B$ et
$\mathcal J\mathcal K=\mathcal J(\mathcal K\mathcal B)
=\mathcal K(\mathcal J\mathcal B)$, de sorte que l'on a le diagramme
d'inclusions
\[
\xymatrix@R=1.5em{
\mathcal J\mathcal B\ar@{^{(}->}[r] & \mathcal B\\
\mathcal J\ar@{^{(}->}[r]\ar@{^{(}->}[u] &
\mathcal A\ar@{^{(}->}[u]\\
\mathcal J\mathcal K\mathcal B=\mathcal J\mathcal K
\ar@{^{(}->}[r]\ar@{^{(}->}[u] &
\mathcal K\ar@{^{(}->}[u]
}
\tag{2.6.2.6}\label{III.2.6.2.6-fr}
\]
Soit $\mathcal J'$ l'Idéal cohérent $(\mathcal J\mathcal B)^\sim$ de
$\mathcal O_{X'}$, de sorte que $\mathcal J\mathcal B=g_*(\mathcal J')$,
$\mathcal J'\mathcal K'=(\mathcal J\mathcal K\mathcal B)^\sim$, et par
suite
\[
\mathcal J'/\mathcal J'\mathcal K'
=(\mathcal J\mathcal B/\mathcal J\mathcal K\mathcal B)^\sim
\]
(II, 1.4.4). Comme $\mathcal J|V=\mathcal J\mathcal K|V$ pour tout
ouvert $V$ ne rencontrant pas $Z$, on voit que le support de
$\mathcal J'/\mathcal J'\mathcal K'$ est contenu dans $Z'$. Comme $Z'$ est
propre sur $Y$, on peut appliquer (\hyperref[III.2.2.4-fr]{2.2.4}) et on
voit que pour $n$ assez grand, l'application canonique
\[
\Gamma(X',\mathcal J'\otimes\mathcal L'^{\otimes n})
\longrightarrow
\Gamma(X',(\mathcal J'/\mathcal J'\mathcal K')
\otimes\mathcal L'^{\otimes n})
\]
est surjective. Mais en vertu de (0\textsubscript{I}, 5.4.10), on en
conclut que l'application canonique
\[
\Gamma(X,\mathcal J\mathcal B\otimes\mathcal L^{\otimes n})
\longrightarrow
\Gamma(X,(\mathcal J\mathcal B/\mathcal J\mathcal K\mathcal B)
\otimes\mathcal L^{\otimes n})
\]
est surjective.

Cela étant, soient $i:Z\to X$ l'injection canonique,
$i':Z'\to X'$ l'injection canonique, de sorte qu'on a le diagramme
commutatif
\[
\xymatrix{
X'\ar[d]_g & Z'\ar[l]_{i'}\ar[d]^{h}\\
X & Z\ar[l]^{i}
}
\]
Soient $\mathcal M=i^*(\mathcal L)$, $\mathcal M'=i'^*(\mathcal L')$~;
comme $\mathcal L'$ est ample, $\mathcal M'$ est ample (II, 5.1.12), et
d'autre part $\mathcal M'=h^*(\mathcal M)$~; on conclut donc de
l'hypothèse de récurrence noethérienne (puisque $Z\neq X$) que
$\mathcal M$ est ample. Par suite
$i^*(\mathcal J/\mathcal J\mathcal K)\otimes\mathcal M^{\otimes n}$ est
engendré par ses sections au-dessus de $Z$ pour $n$ assez grand
(II, 4.5.5). Comme
$\mathcal J/\mathcal J\mathcal K=i_*(i^*(\mathcal J/\mathcal J\mathcal K))$,
on déduit encore de (0\textsubscript{I}, 5.4.10) qu'il existe une section
$s$ de $(\mathcal J/\mathcal J\mathcal K)\otimes
\mathcal L^{\otimes n}$ au-dessus de $X$ (pour un $n$ assez grand) telle
que $s(x)\neq0$, puisque l'on a $\mathcal J_x=\mathcal A_x$ par
définition de $\mathcal J$ et $\mathcal K_x\neq\mathcal A_x$ par
hypothèse. Le diagramme (\hyperref[III.2.6.2.6-fr]{2.6.2.6}) montre que
$s$ est aussi une section de
$(\mathcal J\mathcal B/\mathcal J\mathcal K\mathcal B)
\otimes\mathcal L^{\otimes n}$ au-dessus de $X$, donc $s$ est l'image
canonique d'une section $t$ de
$(\mathcal J\mathcal B)\otimes\mathcal L^{\otimes n}$ au-dessus de $X$.
Mais par définition, l'image canonique $s$ de $t$ mod.
$(\mathcal J\mathcal K\mathcal B)\otimes\mathcal L^{\otimes n}$ est
dans
\[
(\mathcal J\otimes\mathcal L^{\otimes n})/
(\mathcal J\mathcal K\mathcal B\otimes\mathcal L^{\otimes n})
\]
donc, en vertu de (\hyperref[III.2.6.2.6-fr]{2.6.2.6}), cela implique que
$t$ est une section de $\mathcal J\otimes\mathcal L^{\otimes n}$
au-dessus de $X$, et \emph{a fortiori} une section de
$\mathcal L^{\otimes n}$. On a vu ci-dessus que $t(x)\neq0$, donc
$x\in X_t$, et par définition de $\mathcal J$, $t(y)=0$ dans $X-U$ qui
est le support de $\mathcal O_X/\mathcal J$~; donc $X_t\subset U$, ce qui
achève la démonstration.
\end{env}
\end{proof}

\oldpage[III]{115}
\begin{env}[2.6.3]
\label{III.2.6.3-fr}
\emph{Remarque.} Lorsque $X$ est propre sur $Y$, on peut démontrer
(\hyperref[III.2.6.2-fr]{2.6.2}) plus simplement, en raisonnant comme dans
le th. de Chevalley (II, 6.7.1), à l'aide de
(\hyperref[III.2.6.1-fr]{2.6.1}) et du lemme (II, 6.7.1.1).
\end{env}

\section{Le théorème de finitude pour les morphismes propres}
\label{section:III.3-fr}

\subsection{Le lemme de dévissage}
\label{subsection:III.3.1-fr}

\begin{definition}[3.1.1]
\label{III.3.1.1-fr}
Soit $K$ une catégorie abélienne. On dit qu'un sous-ensemble $K'$ de
l'ensemble des objets de $K$ est \emph{exact} si $0\in K'$ et si, pour toute
suite exacte
\[
0\longrightarrow A'\longrightarrow A\longrightarrow A''\longrightarrow0
\]
dans $K$ telle que deux des objets $A$, $A'$, $A''$ soit dans $K'$, alors le
troisième est aussi dans $K'$.
\end{definition}

\begin{theorem}[3.1.2]
\label{III.3.1.2-fr}
Soit $X$ un préschéma noethérien~; on désigne par $K$ la catégorie abélienne
des $\mathcal O_X$-Modules cohérents. Soient $K'$ un sous-ensemble exact de
$K$, $X'$ une partie fermée de l'espace sous-jacent à $X$. On suppose que
pour toute partie fermée irréductible $Y$ de $X'$, de point générique $y$, il
existe un $\mathcal O_X$-Module $\mathcal G\in K'$ tel que $\mathcal G_y$ soit
un $k(y)$-espace vectoriel de dimension $1$. Alors tout $\mathcal O_X$-Module
cohérent de support contenu dans $X'$ appartient à $K'$ (et en particulier,
si $X'=X$, on a $K'=K$).
\end{theorem}

\begin{proof}
Considérons la propriété suivante $P(Y)$ d'une partie fermée $Y$ de $X'$~:
tout $\mathcal O_X$-Module cohérent de support contenu dans $Y$ appartient à
$K'$. En vertu du principe de récurrence noethérienne (0\textsubscript{I},
2.2.2), on voit qu'on est ramené à démontrer que \emph{si $Y$ est une partie
fermée de $X'$ telle que la propriété $P(Y')$ soit vraie pour toute partie
fermée $Y'$ de $Y$, distincte de $Y$, alors $P(Y)$ est vraie}.

Soit donc $\mathcal F\in K$ à support contenu dans $Y$ et prouvons que
$\mathcal F\in K'$. Désignons encore par $Y$ le sous-préschéma fermé réduit de
$X$ ayant $Y$ pour espace sous-jacent (I, 5.2.1)~; il est défini par un Idéal
cohérent $\mathcal J$ de $\mathcal O_X$. On sait (I, 9.3.4) qu'il existe un
entier $n>0$ tel que $\mathcal J^n\mathcal F=0$~; pour $1\leq k\leq n$, on a
donc une suite exacte
\[
0\longrightarrow
\mathcal J^{k-1}\mathcal F/\mathcal J^k\mathcal F
\longrightarrow \mathcal F/\mathcal J^k\mathcal F
\longrightarrow \mathcal F/\mathcal J^{k-1}\mathcal F
\longrightarrow0
\]
de $\mathcal O_X$-Modules cohérents (0\textsubscript{I}, 5.3.6 et 5.3.3)~;
comme $K'$ est exact, on voit, par récurrence sur $k$, qu'il suffit de montrer
que chacun des
$\mathcal F_k=\mathcal J^{k-1}\mathcal F/\mathcal J^k\mathcal F$ est dans
$K'$. On est donc ramené à prouver que $\mathcal F\in K'$ sous l'hypothèse
supplémentaire que $\mathcal J\mathcal F=0$~; il revient au même de dire que
$\mathcal F=j_*(j^*(\mathcal F))$, où $j$ est l'injection canonique
$Y\to X$. Distinguons maintenant deux cas~:
\begin{enumerate}
\item[a)] $Y$ est \emph{réductible}. Soient $Y=Y'\cup Y''$, $Y'$ et $Y''$
étant des parties fermées de $Y$, distinctes de $Y$~; soient encore $Y'$,
$Y''$ les sous-préschémas fermés réduits de $X$ ayant pour espaces
sous-jacents $Y'$, $Y''$ respectivement, qui sont définis respectivement par
les Idéaux cohérents $\mathcal J'$, $\mathcal J''$ de $\mathcal O_X$. Posons
\[
\mathcal F'=\mathcal F\otimes_{\mathcal O_X}
(\mathcal O_X/\mathcal J'),\qquad
\mathcal F''=\mathcal F\otimes_{\mathcal O_X}
(\mathcal O_X/\mathcal J'').
\]
Les homomorphismes canoniques $\mathcal F\to\mathcal F'$,
$\mathcal F\to\mathcal F''$ définissent donc un homomorphisme
$u:\mathcal F\to\mathcal F'\oplus\mathcal F''$. Montrons que pour tout
$z\notin Y'\cap Y''$, l'homomorphisme
$u_z:\mathcal F_z\to\mathcal F'_z\oplus\mathcal F''_z$ est \emph{bijectif}.
En effet, on a $\mathcal J'\cap\mathcal J''=\mathcal J$, car la question est
locale et

\oldpage[III]{116}
l'égalité précédente résulte de (I, 5.2.1 et 1.1.5)~; si $z\notin Y''$, on a
donc $\mathcal J'_z=\mathcal J_z$, d'où $\mathcal F'_z=\mathcal F_z$ et
$\mathcal F''_z=0$, ce qui établit notre assertion dans ce cas~; on raisonne
de même si $z\notin Y'$. Par suite, le noyau et le conoyau de $u$, qui sont
dans $K$ (0\textsubscript{I}, 5.3.4) ont leurs supports dans $Y'\cap Y''$,
et sont donc dans $K'$ par hypothèse~; pour la même raison, $\mathcal F'$ et
$\mathcal F''$ sont dans $K'$, donc aussi $\mathcal F'\oplus\mathcal F''$,
puisque $K'$ est exact. La conclusion résulte alors de la considération des
deux suites exactes
\[
0\longrightarrow\operatorname{Im}u\longrightarrow
\mathcal F'\oplus\mathcal F''\longrightarrow
\operatorname{Coker}u\longrightarrow0,
\]
\[
0\longrightarrow\operatorname{Ker}u\longrightarrow\mathcal F
\longrightarrow\operatorname{Im}u\longrightarrow0
\]
et de l'hypothèse que $K'$ est exact.

\item[b)] $Y$ est irréductible, et par suite le sous-préschéma $Y$ de $X$ est
\emph{intègre}. Si $y$ est son point générique, on a
$(\mathcal O_Y)_y=k(y)$, et comme $j^*(\mathcal F)$ est un
$\mathcal O_Y$-Module cohérent, $\mathcal F_y=(j^*(\mathcal F))_y$ est un
$k(y)$-espace vectoriel de dimension finie $m$. Par hypothèse, il y a un
$\mathcal O_X$-Module cohérent $\mathcal G$ (nécessairement de support $Y$)
tel que $\mathcal G_y$ soit un $k(y)$-espace vectoriel de dimension $1$. Par
suite, il y a un $k(y)$-isomorphisme
$(\mathcal G_y)^m\xrightarrow{\sim}\mathcal F_y$, qui est aussi un
$\mathcal O_Y$-isomorphisme, et comme $\mathcal G^m$ et $\mathcal F$ sont
cohérents, il existe un voisinage ouvert $W$ de $y$ dans $X$ et un
isomorphisme $\mathcal G^m|W\xrightarrow{\sim}\mathcal F|W$
(0\textsubscript{I}, 5.2.7). Soit $\mathcal H$ le graphe de cet isomorphisme,
qui est un sous-$(\mathcal O_X|W)$-Module cohérent de
$(\mathcal G^m\oplus\mathcal F)|W$, canoniquement isomorphe à
$\mathcal G^m|W$ et à $\mathcal F|W$~; il existe donc un sous-$\mathcal O_X$-
Module cohérent $\mathcal H_0$ de $\mathcal G^m\oplus\mathcal F$, induisant
$\mathcal H$ sur $W$ et $0$ sur $X-Y$ puisque $\mathcal G^m$ et
$\mathcal F$ ont pour support $Y$ (I, 9.4.7). Les restrictions
$v:\mathcal H_0\to\mathcal G^m$ et $w:\mathcal H_0\to\mathcal F$ des
projections canoniques de $\mathcal G^m\oplus\mathcal F$ sont alors des
homomorphismes de $\mathcal O_X$-Modules cohérents, qui, dans $W$ et dans
$X-Y$, se réduisent à des isomorphismes~; autrement dit, les noyaux et
conoyaux de $v$ et $w$ ont leur support dans l'ensemble fermé
$Y-(Y\cap W)$, distinct de $Y$. Ils appartiennent donc à $K'$~; d'autre
part, on a $\mathcal G^m\in K'$ puisque $\mathcal G\in K'$ et que $K'$ est
exact. On en conclut successivement, par l'exactitude de $K'$, que
$\mathcal H_0\in K'$ puis que $\mathcal F\in K'$. C.Q.F.D.
\end{enumerate}
\end{proof}

\begin{corollary}[3.1.3]
\label{III.3.1.3-fr}
Supposons que le sous-ensemble exact $K'$ de $K$ ait en outre la propriété
que tout facteur direct cohérent d'un $\mathcal O_X$-Module cohérent
$\mathcal M\in K'$ appartienne encore à $K'$. Dans ces conditions, la
conclusion de (\hyperref[III.3.1.2-fr]{3.1.2}) subsiste encore lorsque la
condition «~$\mathcal G_y$ est un $k(y)$-espace vectoriel de dimension $1$~»
est remplacée par $\mathcal G_y\neq0$ (ce qui équivaut à
$\operatorname{Supp}(\mathcal G)=Y$).
\end{corollary}

\begin{proof}
En effet, le raisonnement de (\hyperref[III.3.1.2-fr]{3.1.2}) ne doit être
modifié que dans le cas b)~; cette fois $\mathcal G_y$ est un $k(y)$-espace
vectoriel de dimension $q>0$, et on a par suite un $\mathcal O_Y$-
isomorphisme $(\mathcal G_y)^m\xrightarrow{\sim}(\mathcal F_y)^q$~; la fin
du raisonnement de (\hyperref[III.3.1.2-fr]{3.1.2}) prouve alors que
$\mathcal F^q\in K'$, et l'hypothèse additionnelle sur $K'$ implique que
$\mathcal F\in K'$.
\end{proof}

\subsection{Le théorème de finitude : cas des schémas usuels}
\label{subsection:III.3.2-fr}

\begin{theorem}[3.2.1]
\label{III.3.2.1-fr}
Soient $Y$ un préschéma localement noethérien, $f:X\to Y$ un morphisme
propre. Pour tout $\mathcal O_X$-Module cohérent $\mathcal F$, les
$\mathcal O_Y$-Modules $R^qf_*(\mathcal F)$ sont cohérents pour $q\geq0$.
\end{theorem}

\begin{proof}
La question étant locale sur $Y$, on peut supposer $Y$ noethérien, donc $X$
noethérien (I, 6.3.7). Les $\mathcal O_X$-Modules cohérents $\mathcal F$ pour
lesquels la conclusion du th. (\hyperref[III.3.2.1-fr]{3.2.1}) est vraie
forment un sous-ensemble exact $K'$ de la catégorie $K$ des
$\mathcal O_X$-Modules cohérents.

\oldpage[III]{117}
En effet, soit
$0\to\mathcal F'\to\mathcal F\to\mathcal F''\to0$ une suite exacte de
$\mathcal O_X$-Modules cohérents~; supposons par exemple que $\mathcal F'$
et $\mathcal F''$ appartiennent à $K'$~; on a la suite exacte de cohomologie
\[
R^{q-1}f_*(\mathcal F'')\xrightarrow{\partial}R^qf_*(\mathcal F')
\longrightarrow R^qf_*(\mathcal F)\longrightarrow R^qf_*(\mathcal F'')
\xrightarrow{\partial}R^{q+1}f_*(\mathcal F')
\]
dans laquelle par hypothèse les quatre termes extrêmes sont cohérents~; il
en est donc de même du terme médian $R^qf_*(\mathcal F)$ par
(0\textsubscript{I}, 5.3.4 et 5.3.3). On montre de la même manière que
lorsque $\mathcal F$ et $\mathcal F'$ (resp. $\mathcal F$ et
$\mathcal F''$) sont dans $K'$, il en est de même de $\mathcal F''$ (resp.
$\mathcal F'$). De plus, tout \emph{facteur direct} cohérent $\mathcal F'$
d'un $\mathcal O_X$-Module $\mathcal F\in K'$ appartient aussi à $K'$~: en
effet, $R^qf_*(\mathcal F')$ est alors facteur direct de
$R^qf_*(\mathcal F)$ (G, II, 4.4.4), donc est de type fini, et comme il est
quasi-cohérent (\hyperref[III.1.4.10-fr]{1.4.10}), il est cohérent, $Y$
étant noethérien. En vertu de (\hyperref[III.3.1.3-fr]{3.1.3}), on est
ramené à prouver que lorsque $X$ est \emph{irréductible} de point générique
$x$, il existe un $\mathcal O_X$-Module cohérent $\mathcal F$ appartenant à
$K'$, tel que $\mathcal F_x\neq0$~: en effet, si ce point est établi, on
pourra l'appliquer à tout sous-préschéma fermé irréductible $Y$ de $X$, car
si $j:Y\to X$ est l'injection canonique, $f\circ j$ est propre (II, 5.4.2),
et si $\mathcal G$ est un $\mathcal O_Y$-Module cohérent de support $Y$,
$j_*(\mathcal G)$ est un $\mathcal O_X$-Module cohérent tel que
\[
R^q(f\circ j)_*(\mathcal G)=R^qf_*(j_*(\mathcal G))
\]
(G, II, 4.9.1), donc on est bien dans les conditions d'application de
(\hyperref[III.3.1.3-fr]{3.1.3}).

Or, en vertu du lemme de Chow (II, 5.6.2), il existe un préschéma
irréductible $X'$ et un morphisme \emph{projectif} et surjectif
$g:X'\to X$ tel que $f\circ g:X'\to Y$ soit \emph{projectif}. Il existe un
$\mathcal O_{X'}$-Module $\mathcal L$ ample pour $g$ (II, 5.3.1)~;
appliquons le th. fondamental des morphismes projectifs
(\hyperref[III.2.2.1-fr]{2.2.1}) à $g:X'\to X$ et à $\mathcal L$~: il
existe donc un entier $n$ tel que
$\mathcal F=g_*(\mathcal O_{X'}(n))$ soit un $\mathcal O_X$-Module cohérent
et $R^qg_*(\mathcal O_{X'}(n))=0$ pour tout $q>0$~; en outre, comme
\[
g^*(g_*(\mathcal O_{X'}(n)))\longrightarrow\mathcal O_{X'}(n)
\]
est surjectif pour $n$ assez grand (\hyperref[III.2.2.1-fr]{2.2.1}), on
voit qu'on peut supposer que, au point générique $x$ de $X$, on a
$\mathcal F_x\neq0$ (II, 3.4.7). D'autre part, comme $f\circ g$ est
projectif et $Y$ noethérien, les
$R^p(f\circ g)_*(\mathcal O_{X'}(n))$ sont \emph{cohérents}
(\hyperref[III.2.2.1-fr]{2.2.1}). Cela étant,
$R^*(f\circ g)_*(\mathcal O_{X'}(n))$ est l'aboutissement d'une suite
spectrale de Leray, dont le terme $E_2$ est donné par
\[
E_2^{pq}=R^pf_*(R^qg_*(\mathcal O_{X'}(n)))~;
\]
ce qui précède montre que cette suite spectrale est dégénérée, et on sait
alors (0, 11.1.6) que $E_2^{p0}=R^pf_*(\mathcal F)$ est isomorphe à
$R^p(f\circ g)_*(\mathcal O_{X'}(n))$, ce qui achève la démonstration.
\end{proof}

\begin{corollary}[3.2.2]
\label{III.3.2.2-fr}
Soit $Y$ un préschéma localement noethérien. Pour tout morphisme propre
$f:X\to Y$, l'image directe par $f$ de tout $\mathcal O_X$-Module cohérent
est un $\mathcal O_Y$-Module cohérent.
\end{corollary}

\begin{corollary}[3.2.3]
\label{III.3.2.3-fr}
Soient $A$ un anneau noethérien, $X$ un schéma propre sur $A$~; pour tout
$\mathcal O_X$-Module cohérent $\mathcal F$, les $H^p(X,\mathcal F)$ sont
des $A$-modules de type fini, et il existe un entier $r>0$ tel que pour tout
$\mathcal O_X$-Module cohérent $\mathcal F$ et tout $p>r$,
$H^p(X,\mathcal F)=0$.
\end{corollary}

\begin{proof}
La seconde assertion a déjà été démontrée
(\hyperref[III.1.4.12-fr]{1.4.12})~; la première résulte du th. de finitude
(\hyperref[III.3.2.1-fr]{3.2.1}), compte tenu de
(\hyperref[III.1.4.11-fr]{1.4.11}).
\end{proof}

En particulier, si $X$ est un \emph{schéma algébrique propre} sur un corps
$k$, alors, pour tout $\mathcal O_X$-Module cohérent $\mathcal F$, les
$H^p(X,\mathcal F)$ sont des $k$-espaces vectoriels de \emph{dimension
finie}.

\begin{corollary}[3.2.4]
\label{III.3.2.4-fr}
Soient $Y$ un préschéma localement noethérien, $f:X\to Y$ un morphisme de
type fini. Pour tout $\mathcal O_X$-Module cohérent $\mathcal F$ dont le
support est propre sur $Y$ (II, 5.4.10), les $\mathcal O_Y$-Modules
$R^qf_*(\mathcal F)$ sont cohérents.
\end{corollary}

\oldpage[III]{118}
\begin{proof}
La question étant locale sur $Y$, on peut supposer $Y$ noethérien, et il en
est donc de même de $X$ (I, 6.3.7). Par hypothèse, tout sous-préschéma fermé
$Z$ de $X$ dont l'espace sous-jacent est $\operatorname{Supp}(\mathcal F)$
est propre sur $Y$, autrement dit, si $j:Z\to X$ est l'injection canonique,
$f\circ j:Z\to Y$ est propre. Or, on peut supposer $Z$ tel que
$\mathcal F=j_*(\mathcal G)$, où $\mathcal G=j^*(\mathcal F)$ est un
$\mathcal O_Z$-Module cohérent (I, 9.3.5)~; comme on a
$R^qf_*(\mathcal F)=R^q(f\circ j)_*(\mathcal G)$ par
(\hyperref[III.1.3.4-fr]{1.3.4}), la conclusion résulte aussitôt de
(\hyperref[III.3.2.1-fr]{3.2.1}).
\end{proof}

\subsection{Généralisation du théorème de finitude (schémas usuels)}
\label{subsection:III.3.3-fr}

\begin{proposition}[3.3.1]
\label{III.3.3.1-fr}
Soient $Y$ un préschéma noethérien, $\mathcal S$ une
$\mathcal O_Y$-Algèbre graduée à degrés positifs, quasi-cohérente et de type
fini, $Y'=\operatorname{Proj}(\mathcal S)$ et $g:Y'\to Y$ le morphisme
structural. Soient $f:X\to Y$ un morphisme propre,
$\mathcal S'=f^*(\mathcal S)$,
$\mathcal M=\bigoplus_{k\in\mathbb Z}\mathcal M_k$ un $\mathcal S'$-Module
gradué quasi-cohérent et de type fini. Alors les
\[
R^pf_*(\mathcal M)=\bigoplus_{k\in\mathbb Z}R^pf_*(\mathcal M_k)
\]
sont des $\mathcal S$-Modules gradués de type fini pour tout $p$. Supposons
en outre que $\mathcal S$ soit engendrée par $\mathcal S_1$~; alors, pour
chaque $p\in\mathbb Z$, il existe un entier $k_p$ tel que pour tout
$k\geq k_p$ et tout $r\geq0$, on ait
\[
R^pf_*(\mathcal M_{k+r})=\mathcal S_rR^pf_*(\mathcal M_k).
\tag{3.3.1.1}\label{III.3.3.1.1-fr}
\]
\end{proposition}

\begin{proof}
La première assertion est identique à l'énoncé de
(\hyperref[III.2.4.1-fr]{2.4.1}, (i)), où on a simplement remplacé
«~morphisme projectif~» par «~morphisme propre~». Or, dans la démonstration
de (\hyperref[III.2.4.1-fr]{2.4.1}, (i)), l'hypothèse sur $f$ a été utilisée
uniquement pour montrer (avec les notations de cette démonstration) que
$R^pf'_*(\widetilde{\mathcal M})$ est un $\mathcal O_{Y'}$-Module cohérent.
Avec les hypothèses de (\hyperref[III.3.3.1-fr]{3.3.1}), $f'$ est propre
(II, 5.4.2, (iii)), donc on peut reprendre sans changement toute la
démonstration de (\hyperref[III.2.4.1-fr]{2.4.1}, (i)), grâce au théorème de
finitude (\hyperref[III.3.2.1-fr]{3.2.1}).

Quant à la seconde assertion, il suffit de remarquer qu'il y a un
recouvrement ouvert affine fini $(U_i)$ de $Y$ tel que les restrictions à
$U_i$ des deux membres de
(\hyperref[III.3.3.1.1-fr]{3.3.1.1}) soient égales pour tout
$k\geq k_{p,i}$ (II, 2.1.6, (ii))~; il suffit de prendre pour $k_p$ le plus
grand des $k_{p,i}$.
\end{proof}

\begin{corollary}[3.3.2]
\label{III.3.3.2-fr}
Soient $A$ un anneau noethérien, $\mathfrak m$ un idéal de $A$, $X$ un
$A$-schéma propre, $\mathcal F$ un $\mathcal O_X$-Module cohérent. Alors,
pour tout $p\geq0$, la somme directe
$\bigoplus_{k\geq0}H^p(X,\mathfrak m^k\mathcal F)$ est un module de type fini
sur l'anneau $S=\bigoplus_{k\geq0}\mathfrak m^k$~; en particulier, il existe
un entier $k_p\geq0$ tel que pour tout $k\geq k_p$, et tout $r\geq0$, on ait
\[
H^p(X,\mathfrak m^{k+r}\mathcal F)
=\mathfrak m^rH^p(X,\mathfrak m^k\mathcal F).
\tag{3.3.3.1}\label{III.3.3.2.1-fr}
\]
\end{corollary}

\begin{proof}
Il suffit d'appliquer (\hyperref[III.3.3.2-fr]{3.3.2}) avec
$Y=\operatorname{Spec}(A)$, $\mathcal S=\widetilde S$,
$\mathcal M_k=\mathfrak m^k\mathcal F$, compte tenu de
(\hyperref[III.1.4.11-fr]{1.4.11}).
\end{proof}

Il convient de rappeler que la structure de $S$-module de
$\bigoplus_{k\geq0}H^p(X,\mathfrak m^k\mathcal F)$ s'obtient en considérant,
pour tout $a\in\mathfrak m^r$, l'application
\[
H^p(X,\mathfrak m^k\mathcal F)
\longrightarrow H^p(X,\mathfrak m^{k+r}\mathcal F)
\]
qui provient par passage à la cohomologie de la multiplication
$\mathfrak m^k\mathcal F\to\mathfrak m^{k+r}\mathcal F$ définie par $a$
(\hyperref[III.2.4.1-fr]{2.4.1}).

\oldpage[III]{119}
\subsection{Le théorème de finitude : cas des schémas formels}
\label{subsection:III.3.4-fr}

Les résultats de cette section (sauf la définition
(\hyperref[III.3.4.1-fr]{3.4.1})) ne seront pas utilisés dans la suite de ce
chapitre.

\begin{env}[3.4.1]
\label{III.3.4.1-fr}
Soient $\mathfrak X$ et $\mathfrak S$ deux préschémas formels localement
noethériens (I, 10.4.2), $f:\mathfrak X\to\mathfrak S$ un morphisme de
préschémas formels. Nous dirons que $f$ est un morphisme \emph{propre} s'il
vérifie les conditions suivantes~:
\begin{enumerate}
\item[$1^\circ$] $f$ est un morphisme de type fini (I, 10.13.3).
\item[$2^\circ$] Si $\mathcal K$ est un Idéal de définition de $\mathfrak S$
et si l'on pose
$\mathcal J=f^*(\mathcal K)\mathcal O_{\mathfrak X}$,
$X_0=(\mathfrak X,\mathcal O_{\mathfrak X}/\mathcal J)$,
$S_0=(\mathfrak S,\mathcal O_{\mathfrak S}/\mathcal K)$, le morphisme
$f_0:X_0\to Y_0$ déduit de $f$ (I, 10.5.6) est propre.
\end{enumerate}

Il est immédiat que cette définition ne dépend pas de l'Idéal de définition
$\mathcal K$ de $\mathfrak S$ considéré~; en effet, si $\mathcal K'$ est un
second Idéal de définition tel que $\mathcal K'\subset\mathcal K$, et si on
pose $\mathcal J'=f^*(\mathcal K')\mathcal O_{\mathfrak X}$,
$X'_0=(\mathfrak X,\mathcal O_{\mathfrak X}/\mathcal J')$,
$S'_0=(\mathfrak S,\mathcal O_{\mathfrak S}/\mathcal K')$, le morphisme
$f'_0:X'_0\to S'_0$ déduit de $f$ est tel que le diagramme
\[
\xymatrix{
X_0\ar[r]^{f_0}\ar[d]_i & S_0\ar[d]^j\\
X'_0\ar[r]_{f'_0} & S'_0
}
\]
est commutatif, $i$ et $j$ étant des immersions surjectives~; il revient donc
au même de dire que $f_0$ ou $f'_0$ est propre, en vertu de (II, 5.4.5).

On notera que, pour tout $n\geq0$, si on pose
$X_n=(\mathfrak X,\mathcal O_{\mathfrak X}/\mathcal J^{n+1})$,
$S_n=(\mathfrak S,\mathcal O_{\mathfrak S}/\mathcal K^{n+1})$, le morphisme
$f_n:X_n\to Y_n$ déduit de $f$ (I, 10.5.6) est propre pour tout $n$ dès qu'il
l'est pour $n=0$ (II, 5.4.6).

Si $g:Y\to Z$ est un morphisme propre de préschémas usuels, localement
noethériens, $Z'$ une partie fermée de $Z$, $Y'$ une partie fermée de $Y$
telle que $g(Y')\subset Z'$, le prolongement
$\widehat g:Y_{/Y'}\to Z_{/Z'}$ de $g$ aux complétés (I, 10.9.1) est un
morphisme propre de préschémas formels, comme il résulte de la définition et
de (II, 5.4.5).

Soient $\mathfrak X$ et $\mathfrak S$ deux préschémas formels localement
noethériens, $f:\mathfrak X\to\mathfrak S$ un morphisme \emph{de type fini}
(I, 10.13.3)~; les notations étant les mêmes que ci-dessus, on dit qu'une
partie $Z$ de l'espace sous-jacent à $\mathfrak X$ est \emph{propre sur}
$\mathfrak S$ (ou propre pour $f$) si, considérée comme partie de $X_0$, $Z$
est \emph{propre sur $S_0$} (II, 5.4.10). Toutes les propriétés des parties
propres de préschémas usuels énoncées dans (II, 5.4.10) sont encore valables
pour les parties propres de préschémas formels comme il résulte aussitôt des
définitions.
\end{env}

\begin{theorem}[3.4.2]
\label{III.3.4.2-fr}
Soient $\mathfrak X$, $\mathfrak Y$ deux préschémas formels localement
noethériens, $f:\mathfrak X\to\mathfrak Y$ un morphisme propre. Pour tout
$\mathcal O_{\mathfrak X}$-Module cohérent $\mathcal F$, les
$\mathcal O_{\mathfrak Y}$-Modules $R^qf_*(\mathcal F)$ sont cohérents pour
tout $q\geq0$.
\end{theorem}

Soient $\mathcal J$ un Idéal de définition de $\mathfrak Y$,
$\mathcal K=f^*(\mathcal J)\mathcal O_{\mathfrak X}$, et considérons les
$\mathcal O_{\mathfrak X}$-Modules
\[
\mathcal F_k=\mathcal F\otimes_{\mathcal O_{\mathfrak Y}}
(\mathcal O_{\mathfrak Y}/\mathcal J^{k+1})
=\mathcal F/\mathcal K^{k+1}\mathcal F\qquad(k\geq0)
\tag{3.4.2.1}\label{III.3.4.2.1-fr}
\]
qui forment évidemment un \emph{système projectif} de
$\mathcal O_{\mathfrak X}$-Modules topologiques, tel que

\oldpage[III]{120}
\[
\mathcal F=\varprojlim_k\mathcal F_k
\]
par (I, 10.11.3). D'autre part, il résultera de
(\hyperref[III.3.4.2-fr]{3.4.2}) que chacun des $R^qf_*(\mathcal F)$, étant
cohérent, est naturellement muni d'une structure de
$\mathcal O_{\mathfrak Y}$-Module topologique (I, 10.11.6), et il en est de
même des $R^qf_*(\mathcal F_k)$. Aux homomorphismes canoniques
$\mathcal F\to\mathcal F_k=\mathcal F/\mathcal K^{k+1}\mathcal F$
correspondent canoniquement des homomorphismes
\[
R^qf_*(\mathcal F)\longrightarrow R^qf_*(\mathcal F_k)
\]
qui sont nécessairement continus pour les structures de
$\mathcal O_{\mathfrak Y}$-Module topologique précédentes (I, 10.11.6), et
forment un système projectif, donnant à la limite un homomorphisme canonique
fonctoriel
\[
R^qf_*(\mathcal F)\longrightarrow\varprojlim_kR^qf_*(\mathcal F_k)
\tag{3.4.2.2}\label{III.3.4.2.2-fr}
\]
qui sera un homomorphisme continu de $\mathcal O_{\mathfrak Y}$-Modules
topologiques. Nous allons démontrer en même temps que
(\hyperref[III.3.4.2-fr]{3.4.2}) le

\begin{corollary}[3.4.3]
\label{III.3.4.3-fr}
Chacun des homomorphismes (\hyperref[III.3.4.2.2-fr]{3.4.2.2}) est un
isomorphisme topologique. En outre, si $\mathfrak Y$ est noethérien, le
système projectif
$(R^qf_*(\mathcal F/\mathcal K^{k+1}\mathcal F))_{k\geq0}$ satisfait à la
condition (ML) (0, 13.1.1).
\end{corollary}

Nous commencerons par établir (\hyperref[III.3.4.2-fr]{3.4.2}) et
(\hyperref[III.3.4.3-fr]{3.4.3}) lorsque $Y$ est un schéma affine
formel noethérien (I, 10.4.1)~:

\begin{corollary}[3.4.4]
\label{III.3.4.4-fr}
Sous les hypothèses de (\hyperref[III.3.4.2-fr]{3.4.2}), supposons en outre
que $\mathfrak Y=\operatorname{Spf}(A)$, où $A$ est un anneau adique
noethérien. Soit $\mathfrak J$ un idéal de définition de $A$, et posons
$\mathcal F_k=\mathcal F/\mathfrak J^{k+1}\mathcal F$ pour $k\geq0$. Alors
les $H^n(\mathfrak X,\mathcal F)$ sont des $A$-modules de type fini~; le
système projectif $(H^n(\mathfrak X,\mathcal F_k))_{k\geq0}$ satisfait à la
condition (ML) pour tout $n$~; si on pose
\[
N_{n,k}=\operatorname{Ker}\bigl(H^n(\mathfrak X,\mathcal F)
\longrightarrow H^n(\mathfrak X,\mathcal F_k)\bigr)
\tag{3.4.4.1}\label{III.3.4.4.1-fr}
\]
(égal aussi à
$\operatorname{Im}(H^n(\mathfrak X,\mathfrak J^{k+1}\mathcal F)
\to H^n(\mathfrak X,\mathcal F))$ par la suite exacte de cohomologie), les
$N_{k,n}$ définissent sur $H^n(\mathfrak X,\mathcal F)$ une filtration
$\mathfrak J$-bonne (0, 13.7.7)~; enfin, l'homomorphisme canonique
\[
H^n(\mathfrak X,\mathcal F)
\longrightarrow\varprojlim_kH^n(\mathfrak X,\mathcal F_k)
\tag{3.4.4.2}\label{III.3.4.4.2-fr}
\]
est un isomorphisme topologique pour tout $n$ (le premier membre étant muni
de la topologie $\mathfrak J$-adique, les $H^n(\mathfrak X,\mathcal F_k)$ de
la topologie discrète).
\end{corollary}

Posons
\[
S=\operatorname{gr}(A)=\bigoplus_{k\geq0}
\mathfrak J^k/\mathfrak J^{k+1},\qquad
\mathcal M=\operatorname{gr}(\mathcal F)=\bigoplus_{k\geq0}
\mathfrak J^k\mathcal F/\mathfrak J^{k+1}\mathcal F.
\tag{3.4.4.3}\label{III.3.4.4.3-fr}
\]
On sait que $\mathfrak J^\Delta$ est un Idéal de définition de $\mathfrak Y$
(I, 10.3.1)~; soient
$\mathcal K=f^*(\mathfrak J^\Delta)\mathcal O_{\mathfrak X}$,
$X_0=(\mathfrak X,\mathcal O_{\mathfrak X}/\mathcal K)$,
$Y_0=(\mathfrak Y,\mathcal O_{\mathfrak Y}/\mathfrak J^\Delta)
=\operatorname{Spec}(A_0)$ avec $A_0=A/\mathfrak J$. Il est clair que les
$\mathcal M_k=\mathfrak J^k\mathcal F/\mathfrak J^{k+1}\mathcal F$ sont des
$\mathcal O_{X_0}$-Modules cohérents (I, 10.11.3). Considérons d'autre part la
$\mathcal O_{X_0}$-Algèbre graduée quasi-cohérente
\[
\mathcal S=\mathcal O_{X_0}\otimes_{A_0}S
=\operatorname{gr}(\mathcal O_{\mathfrak X})
=\bigoplus_{k\geq0}\mathcal K^k/\mathcal K^{k+1}.
\tag{3.4.4.4}\label{III.3.4.4.4-fr}
\]

L'hypothèse que $\mathcal F$ est un $\mathcal O_{\mathfrak X}$-Module de type
fini entraîne d'abord que $\mathcal M$ est

\oldpage[III]{121}
un $\mathcal S$-Module gradué \emph{de type fini}. En effet, la question est
locale sur $\mathfrak X$, et on peut donc supposer pour la traiter que
$\mathfrak X=\operatorname{Spf}(B)$, où $B$ est un anneau adique noethérien,
et $\mathcal F=N^\Delta$, où $N$ est un $B$-module de type fini
(I, 10.10.5)~; on a en outre $X_0=\operatorname{Spec}(B_0)$ où
$B_0=B/\mathfrak J B$, et les $\mathcal O_{X_0}$-Modules quasi-cohérents
$\mathcal S$ et $\mathcal M$ sont respectivement égaux à $\widetilde{S'}$ et
$\widetilde{M'}$, où
\[
S'=\bigoplus_{k\geq0}
((\mathfrak J^k/\mathfrak J^{k+1})\otimes_{A_0}B_0),\qquad
M'=\bigoplus_{k\geq0}
((\mathfrak J^k/\mathfrak J^{k+1})\otimes_{A_0}N_0),
\]
avec $N_0=N/\mathfrak JN$~; on a donc évidemment
$M'=S'\otimes_{B_0}N_0$, et comme $N_0$ est un $B_0$-module de type fini,
$M'$ est un $S'$-module de type fini, d'où notre assertion (I, 1.3.13).

Comme le morphisme $f_0:X_0\to Y_0$ est \emph{propre} par hypothèse, on peut
appliquer (\hyperref[III.3.3.2-fr]{3.3.2}) à $\mathcal S$, $\mathcal M$ et au
morphisme $f_0$~; compte tenu de (\hyperref[III.1.4.11-fr]{1.4.11}), on en
conclut que pour \emph{tout $n\geq0$},
$\bigoplus_{k\geq0}H^n(X_0,\mathcal M_k)$ est un $S$-module gradué \emph{de
type fini}. Cela prouve que la condition $(\mathrm F_n)$ de (0, 13.7.7) est
vérifiée pour \emph{tout $n\geq0$}, lorsqu'on considère le système projectif
strict $(\mathcal F/\mathfrak J^k\mathcal F)_{k\geq0}$ de faisceaux de
groupes abéliens sur $X_0$, munis chacun de sa structure naturelle de
«~$A$-module filtré~». On peut donc appliquer (0, 13.7.7), qui prouve que~:
\begin{enumerate}
\item[$1^\circ$] Le système projectif
$(H^n(\mathfrak X,\mathcal F_k))_{k\geq0}$ vérifie la condition (ML).
\item[$2^\circ$] Si
$H'^{,n}=\varprojlim_kH^n(\mathfrak X,\mathcal F_k)$, $H'^{,n}$ est un
$A$-module de type fini.
\item[$3^\circ$] La filtration définie sur $H'^{,n}$ par les noyaux des
homomorphismes canoniques $H'^{,n}\to H^n(\mathfrak X,\mathcal F_k)$ est
$\mathfrak J$-bonne.
\end{enumerate}

Notons d'autre part que si on pose
$X_k=(\mathfrak X,\mathcal O_{\mathfrak X}/\mathcal K^{k+1})$,
$\mathcal F_k$ est un $\mathcal O_{X_k}$-Module cohérent (I, 10.11.3) et si
$U$ est un ouvert affine dans $X_0$, $U$ est aussi un ouvert affine dans
chacun des $X_k$ (I, 5.1.9), donc $H^n(U,\mathcal F_k)=0$ pour tout $n>0$ et
tout $k$ (\hyperref[III.1.3.1-fr]{1.3.1}) et
$H^0(U,\mathcal F_k)\to H^0(U,\mathcal F_h)$ est surjectif pour $h\leq k$
(I, 1.3.9). On est donc dans les conditions de (0, 13.3.2) et l'application
de (0, 13.3.1) prouve que $H'^{,n}$ s'identifie canoniquement à
$H^n(\mathfrak X,\varprojlim_k\mathcal F_k)=H^n(\mathfrak X,\mathcal F)$~;
cela achève de prouver (\hyperref[III.3.4.4-fr]{3.4.4}).

\begin{env}[3.4.5]
\label{III.3.4.5-fr}
Revenons maintenant à la démonstration de (\hyperref[III.3.4.2-fr]{3.4.2})
et (\hyperref[III.3.4.3-fr]{3.4.3}). Prouvons d'abord ces propositions dans
le cas $\mathfrak Y=\operatorname{Spf}(A)$ envisagé dans
(\hyperref[III.3.4.4-fr]{3.4.4})~; pour cela, pour tout $g\in A$, appliquons
(\hyperref[III.3.4.4-fr]{3.4.4}) au schéma formel affine noethérien induit
sur l'ouvert $\mathfrak Y_g=\mathfrak D(g)$ de $\mathfrak Y$, qui est égal à
$\operatorname{Spf}(A_{\{g\}})$, et au préschéma formel induit par
$\mathfrak X$ sur $f^{-1}(\mathfrak Y_g)$~; notons que $\mathfrak Y_g$ est
aussi un ouvert affine dans le préschéma
$Y_k=(\mathfrak Y,\mathcal O_{\mathfrak Y}/(\mathfrak J^\Delta)^{k+1})$,
et comme $\mathcal F_k$ est un $\mathcal O_{X_k}$-Module cohérent, on a
\[
H^n(f^{-1}(\mathfrak Y_g),\mathcal F_k)
=\Gamma(\mathfrak Y_g,R^nf_*(\mathcal F_k))
\]
pour tout $k\geq0$ en vertu de
(\hyperref[III.1.4.11-fr]{1.4.11}). L'homomorphisme canonique
\[
H^n(f^{-1}(\mathfrak Y_g),\mathcal F)
\longrightarrow\varprojlim_k\Gamma(\mathfrak Y_g,R^nf_*(\mathcal F_k))
\]
est un isomorphisme~; mais on a (0\textsubscript{I}, 3.2.6)
\[
\varprojlim_k\Gamma(\mathfrak Y_g,R^nf_*(\mathcal F_k))
=\Gamma(\mathfrak Y_g,\varprojlim_kR^nf_*(\mathcal F_k))
\]

\oldpage[III]{122}
et comme le faisceau $R^nf_*(\mathcal F)$ est associé au préfaisceau
$\mathfrak Y_g\mapsto H^n(f^{-1}(\mathfrak Y_g),\mathcal F)$ sur les
$\mathfrak Y_g$ (0\textsubscript{I}, 3.2.1), on a bien montré que
l'homomorphisme (\hyperref[III.3.4.2.2-fr]{3.4.2.2}) est bijectif. Prouvons
ensuite que $R^nf_*(\mathcal F)$ est un $\mathcal O_{\mathfrak Y}$-Module
cohérent, et de façon plus précise que l'on a
\[
R^nf_*(\mathcal F)=(H^n(\mathfrak X,\mathcal F))^\Delta.
\tag{3.4.5.1}\label{III.3.4.5.1-fr}
\]
Avec les notations précédentes, on a, puisque $\mathcal F_k$ est un
$\mathcal O_{X_k}$-Module cohérent (\hyperref[III.1.4.13-fr]{1.4.13})
\[
\Gamma(\mathfrak Y_g,R^nf_*(\mathcal F_k))
=(\Gamma(\mathfrak Y,R^nf_*(\mathcal F_k)))_g
=(H^n(\mathfrak X,\mathcal F_k))_g.
\]
Or, les $H^n(\mathfrak X,\mathcal F_k)$ forment un système projectif
vérifiant (ML) et leur limite projective $H^n(\mathfrak X,\mathcal F)$ est
un $A$-module de type fini. On en conclut (0, 13.7.8) que l'on a
\[
\varprojlim_k(H^n(\mathfrak X,\mathcal F_k))_g
=H^n(\mathfrak X,\mathcal F)\otimes_AA_{\{g\}}
=\Gamma(\mathfrak Y_g,(H^n(\mathfrak X,\mathcal F))^\Delta)
\]
compte tenu de (I, 10.10.8) appliqué à $A$ et $A_{\{g\}}$~; ceci démontre
(\hyperref[III.3.4.5.1-fr]{3.4.5.1}) puisque
\[
\Gamma(\mathfrak Y_g,R^nf_*(\mathcal F))
=\varprojlim_k\Gamma(\mathfrak Y_g,R^nf_*(\mathcal F_k)).
\]
Comme (\hyperref[III.3.4.2.2-fr]{3.4.2.2}) est alors un isomorphisme de
$\mathcal O_{\mathfrak Y}$-Modules cohérents, c'est nécessairement un
isomorphisme \emph{topologique} (I, 10.11.6). Enfin, il résulte des relations
$R^nf_*(\mathcal F_k)=(H^n(\mathfrak X,\mathcal F_k))^\Delta$ que le système
projectif $(R^nf_*(\mathcal F_k))_{k\geq0}$ vérifie (ML) (I, 10.10.2).

Une fois (\hyperref[III.3.4.2-fr]{3.4.2}) et
(\hyperref[III.3.4.3-fr]{3.4.3}) démontrés dans le cas où le préschéma formel
$\mathfrak Y$ est affine noethérien, il est immédiat de passer de là au cas
général pour (\hyperref[III.3.4.2-fr]{3.4.2}) et la première assertion de
(\hyperref[III.3.4.3-fr]{3.4.3}), qui sont locales sur $\mathfrak Y$. Quant à
la seconde assertion de (\hyperref[III.3.4.3-fr]{3.4.3}), il suffit,
$\mathfrak Y$ étant noethérien, de le recouvrir par un nombre fini d'ouverts
affines noethériens $U_i$ et de remarquer que les restrictions du système
projectif $(R^qf_*(\mathcal F_k))$ à chacun des $U_i$ vérifient (ML).
\end{env}

Nous avons en outre démontré en cours de route~:

\begin{corollary}[3.4.6]
\label{III.3.4.6-fr}
Sous les hypothèses de (\hyperref[III.3.4.4-fr]{3.4.4}), l'homomorphisme
canonique
\[
H^q(\mathfrak X,\mathcal F)
\longrightarrow\Gamma(\mathfrak Y,R^qf_*(\mathcal F))
\tag{3.4.6.1}\label{III.3.4.6.1-fr}
\]
est bijectif.
\end{corollary}

\section{Le théorème fondamental des morphismes propres. Applications}
\label{section:III.4-fr}

\subsection{Le théorème fondamental}
\label{subsection:III.4.1-fr}

\begin{env}[4.1.1]
\label{III.4.1.1-fr}
Soient $X$, $Y$ deux préschémas usuels noethériens, $f:X\to Y$ un
morphisme propre, $Y'$ une partie fermée de $Y$, $X'$ son image réciproque
$f^{-1}(Y')$. Nous désignerons par $\widehat X$ et $\widehat Y$ les
préschémas formels $X_{/X'}$ et $Y_{/Y'}$, complétés de $X$ et $Y$ le long
de $X'$ et $Y'$ respectivement (I, 10.8.5), par $\widehat f$ le prolongement
de $f$ à ces complétés (I, 10.9.1) qui est un morphisme
$\widehat X\to\widehat Y$ de préschémas formels. Pour tout
$\mathcal O_X$-Module cohérent $\mathcal F$, nous

\oldpage[III]{123}

désignerons par $\widehat{\mathcal F}$ son complété $\mathcal F_{/X'}$ le
long de $X'$ (I, 10.8.4) qui est un $\mathcal O_{\widehat X}$-Module
cohérent (I, 10.8.8).
\end{env}

\begin{env}[4.1.2]
\label{III.4.1.2-fr}
Soit $\mathcal J$ un Idéal cohérent de $\mathcal O_Y$ tel que
$\operatorname{Supp}(\mathcal O_Y/\mathcal J)=Y'$ (I, 5.2.1)~; on sait
(I, 4.4.5) que
\[
\mathcal K=f^*(\mathcal J)\mathcal O_X
\]
est un Idéal cohérent de $\mathcal O_X$ tel que
\[
\operatorname{Supp}(\mathcal O_X/\mathcal K)=X'.
\]
Nous considérerons pour tout $k\geq0$ les $\mathcal O_X$-Modules cohérents
\[
\mathcal F_k
=\mathcal F\otimes_{\mathcal O_Y}(\mathcal O_Y/\mathcal J^{k+1})
=\mathcal F/\mathcal K^{k+1}\mathcal F.
\]
Les $\mathcal O_Y$-Modules $R^nf_*(\mathcal F)$ et
$R^nf_*(\mathcal F_k)$ sont cohérents pour tout $n$
(\hyperref[III.3.2.1-fr]{3.2.1}). Pour tout $k\geq0$ et tout $n$,
l'homomorphisme canonique $\mathcal F\to\mathcal F_k$ définit par
fonctorialité un homomorphisme
\[
R^nf_*(\mathcal F)\longrightarrow R^nf_*(\mathcal F_k).
\tag{4.1.2.1}\label{III.4.1.2.1-fr}
\]
En outre, comme $\mathcal F_k$ est un
$\mathcal O_X/\mathcal K^{k+1}$-Module, $R^nf_*(\mathcal F_k)$ est un
$\mathcal O_Y/\mathcal J^{k+1}$-Module (0, 12.2.1) et on déduit donc de
(\hyperref[III.4.1.2.1-fr]{4.1.2.1}) un homomorphisme
\[
R^nf_*(\mathcal F)\otimes_{\mathcal O_Y}
(\mathcal O_Y/\mathcal J^{k+1})
\longrightarrow R^nf_*(\mathcal F_k).
\tag{4.1.2.2}\label{III.4.1.2.2-fr}
\]
Les deux membres de (\hyperref[III.4.1.2.2-fr]{4.1.2.2}) forment deux
systèmes projectifs, et la limite projective du premier membre n'est autre
que le complété $(R^nf_*(\mathcal F))_{/Y'}$, que nous noterons
$\widehat{R^nf_*(\mathcal F)}$. En outre, il est immédiat que les
homomorphismes (\hyperref[III.4.1.2.2-fr]{4.1.2.2}) forment un système
projectif, d'où par passage à la limite un homomorphisme canonique
\[
\varphi_n:\widehat{R^nf_*(\mathcal F)}
\longrightarrow\varprojlim_kR^nf_*(\mathcal F_k).
\tag{4.1.2.3}\label{III.4.1.2.3-fr}
\]
D'ailleurs (\hyperref[III.4.1.2.2-fr]{4.1.2.2}) est un homomorphisme de
$(\mathcal O_Y/\mathcal J^{k+1})$-Modules, et par suite (I, 10.8.3) peut
être considéré comme un homomorphisme continu de
$\mathcal O_{\widehat Y}$-Modules topologiques pseudo-discrets
(0\textsubscript{I}, 3.8.1). L'homomorphisme $\varphi_n$ est par suite un
homomorphisme continu de $\mathcal O_{\widehat Y}$-Modules topologiques.
\end{env}

\begin{env}[4.1.3]
\label{III.4.1.3-fr}
Soit $i:\widehat X\to X$ le morphisme canonique d'espaces annelés défini
dans (I, 10.8.7), de sorte que l'on a le diagramme commutatif
\[
\xymatrix{
X_k\ar[r]^{h_k}\ar[dr]_{i_k}&\widehat X\ar[d]^i\\
&X
}
\tag{4.1.3.1}\label{III.4.1.3.1-fr}
\]
où $X_k$ est le sous-préschéma fermé de $X$ défini par l'Idéal
$\mathcal K^{k+1}$, $i_k$ l'injection canonique, $h_k$ le morphisme
d'espaces annelés correspondant à l'identité dans les espaces sous-jacents
et à l'homomorphisme canonique
$\mathcal O_{\widehat X}\to\mathcal O_X/\mathcal K^{k+1}$ (I, 10.5.2).
En outre, on a $\widehat{\mathcal F}=i^*(\mathcal F)$ (I, 10.8.8) à un
isomorphisme canonique près. On sait que l'on a
\[
H^n(X_k,i_k^*(\mathcal F_k))=H^n(X,\mathcal F_k)
\tag{4.1.3.2}\label{III.4.1.3.2-fr}
\]

\oldpage[III]{124}

à un isomorphisme canonique près, puisque
$\mathcal F_k=(i_k)_*(i_k^*(\mathcal F_k))$ (G, II, 4.9.1)~;
l'homomorphisme canonique
$H^n(\widehat X,\widehat{\mathcal F})\to
H^n(X_k,h_k^*(\widehat{\mathcal F}))$ (0, 12.1.3.5) s'écrit donc aussi
\[
H^n(\widehat X,\widehat{\mathcal F})\longrightarrow H^n(X,\mathcal F_k)
\tag{4.1.3.3}\label{III.4.1.3.3-fr}
\]
et ces homomorphismes forment évidemment un système projectif, d'où par
passage à la limite, un homomorphisme canonique
\[
\psi_{n,X}:H^n(\widehat X,\widehat{\mathcal F})
\longrightarrow\varprojlim_kH^n(X,\mathcal F_k).
\tag{4.1.3.4}\label{III.4.1.3.4-fr}
\]
Remplaçant $X$ par un ouvert de la forme $f^{-1}(V)$, où $V$ est un
ouvert affine de $Y$, on a, compte tenu de
(\hyperref[III.1.4.11-fr]{1.4.11}), des homomorphismes
\[
\psi_{n,V}:H^n(\widehat X\cap f^{-1}(V),\widehat{\mathcal F})
\longrightarrow\varprojlim_k\Gamma(V,R^nf_*(\mathcal F_k))~;
\tag{4.1.3.5}\label{III.4.1.3.5-fr}
\]
ces homomorphismes commutent de façon évidente à la restriction de $V$ à
un ouvert affine plus petit, donc définissent finalement un homomorphisme
canonique de faisceaux
\[
\psi_n:R^n\widehat f_*(\widehat{\mathcal F})
\longrightarrow\varprojlim_kR^nf_*(\mathcal F_k).
\tag{4.1.3.6}\label{III.4.1.3.6-fr}
\]
\end{env}

\begin{env}[4.1.4]
\label{III.4.1.4-fr}
Soit enfin $j:\widehat Y\to Y$ le morphisme canonique d'espaces annelés
(I, 10.8.7)~; comme $R^nf_*(\mathcal F)$ est un $\mathcal O_Y$-Module
cohérent (\hyperref[III.3.2.1-fr]{3.2.1}), on a
$j^*(R^nf_*(\mathcal F))=\widehat{R^nf_*(\mathcal F)}$ à un isomorphisme
canonique près (I, 10.8.8) et on a donc un homomorphisme canonique
\[
\rho_n:\widehat{R^nf_*(\mathcal F)}=j^*(R^nf_*(\mathcal F))
\longrightarrow R^n\widehat f_*(i^*(\mathcal F))
=R^n\widehat f_*(\widehat{\mathcal F})
\tag{4.1.4.1}\label{III.4.1.4.1-fr}
\]
défini de façon générale pour les espaces annelés (voir la démonstration
de (\hyperref[III.1.4.15-fr]{1.4.15})). Montrons que le diagramme
\[
\xymatrix{
\widehat{R^nf_*(\mathcal F)}\ar[rr]^{\rho_n}\ar[dr]_{\varphi_n}
&&R^n\widehat f_*(\widehat{\mathcal F})\ar[dl]^{\psi_n}\\
&\displaystyle\varprojlim_kR^nf_*(\mathcal F_k)&
}
\tag{4.1.4.2}\label{III.4.1.4.2-fr}
\]
est commutatif. Il suffit évidemment de prouver la commutativité du
diagramme correspondant d'homomorphismes de préfaisceaux, donc on peut se
borner au cas où $Y$ est affine, et tout revient à prouver que le diagramme
\[
\xymatrix{
\widehat{H^n(X,\mathcal F)}\ar[rr]^{\rho_n}\ar[dr]_{\varphi_n}
&&H^n(\widehat X,\widehat{\mathcal F})\ar[dl]^{\psi_{n,X}}\\
&\displaystyle\varprojlim_kH^n(X,\mathcal F_k)&
}
\tag{4.1.4.3}\label{III.4.1.4.3-fr}
\]
est commutatif. Mais la commutativité de
(\hyperref[III.4.1.3.1-fr]{4.1.3.1}) et les relations vues dans
(\hyperref[III.4.1.3-fr]{4.1.3}) entre les groupes de cohomologie donnent
aussitôt le diagramme commutatif

\oldpage[III]{125}

\[
\xymatrix{
H^n(X,\mathcal F)\ar[r]\ar[dr]&
H^n(\widehat X,\widehat{\mathcal F})
=H^n(\widehat X,i^*(\mathcal F))\ar[d]\\
&H^n(X,\mathcal F_k)=H^n(X_k,i_k^*(\mathcal F))
}
\]
d'où on déduit aussitôt la commutativité de
(\hyperref[III.4.1.4.3-fr]{4.1.4.3}).
\end{env}

\begin{theorem}[4.1.5]
\label{III.4.1.5-fr}
Soient $f:X\to Y$ un morphisme propre de préschémas noethériens, $Y'$ une
partie fermée de $Y$, $X'=f^{-1}(Y')$. Alors, pour tout
$\mathcal O_X$-Module cohérent $\mathcal F$,
$R^n\widehat f_*(\widehat{\mathcal F})$ est un
$\mathcal O_{\widehat X}$-Module cohérent et les homomorphismes
$\varphi_n$, $\psi_n$ et $\rho_n$ du diagramme
(\hyperref[III.4.1.4.2-fr]{4.1.4.2}) sont des isomorphismes topologiques.
\end{theorem}

\begin{proof}
Il suffira évidemment de prouver que $\varphi_n$ et $\psi_n$ sont des
isomorphismes~; comme $R^nf_*(\mathcal F)$ est cohérent
(\hyperref[III.3.2.1-fr]{3.2.1}), il en résultera que
$\widehat{R^nf_*(\mathcal F)}$ est cohérent (I, 10.8.8) et la bicontinuité
de $\varphi_n$, $\psi_n$ et $\rho_n$ est alors automatique (I, 10.11.6).
\end{proof}

\begin{env}[4.1.6]
\label{III.4.1.6-fr}
\emph{Remarques.}
\begin{enumerate}
\item[(i)] Si on pose
$\widehat{\mathcal F}_k=
\widehat{\mathcal F}/\widehat{\mathcal K}^{k+1}
\widehat{\mathcal F}$, il est immédiat que
$\mathcal F_k=i_*(\widehat{\mathcal F}_k)$, et l'homomorphisme canonique
(\hyperref[III.4.1.3.6-fr]{4.1.3.6}) n'est autre que l'homomorphisme déjà
défini en (\hyperref[III.3.4.2.2-fr]{3.4.2.2})
\[
R^n\widehat f_*(\widehat{\mathcal F})
\longrightarrow\varprojlim_kR^n\widehat f_*(\widehat{\mathcal F}_k)~;
\tag{4.1.6.1}\label{III.4.1.6.1-fr}
\]
par suite le fait que $\psi_n$ soit un isomorphisme est un cas particulier
de (\hyperref[III.3.4.3-fr]{3.4.3}). Mais nous en donnerons ci-dessous une
démonstration directe, évitant les considérations délicates sur les limites
projectives de suites spectrales (0, 13.7), sur lesquelles repose le
théorème général (\hyperref[III.3.4.3-fr]{3.4.3}).

\item[(ii)] Compte tenu du fait que les $\psi_n$ sont des isomorphismes, il
est équivalent de dire que les $\varphi_n$ ou les
$\rho_n=\psi_n^{-1}\circ\varphi_n$ sont des isomorphismes. Le th.
(\hyperref[III.4.1.5-fr]{4.1.5}) exprime entre autres que la formation des
$R^nf_*$ commute à la complétion et pourra s'appeler le premier théorème de
comparaison de la théorie «~algébrique~» à la théorie «~formelle~».
\end{enumerate}
Nous commencerons par établir la forme affine de
(\hyperref[III.4.1.5-fr]{4.1.5})~:
\end{env}

\begin{corollary}[4.1.7]
\label{III.4.1.7-fr}
Les hypothèses étant celles de (\hyperref[III.4.1.5-fr]{4.1.5}), supposons
en outre $Y=\operatorname{Spec}(A)$, où $A$ est noethérien, et
$\mathcal J=\widetilde{\mathfrak J}$, où $\mathfrak J$ est un idéal de
$A$, de sorte que
$\mathcal F_k=\mathcal F/\mathfrak J^{k+1}\mathcal F$.
L'homomorphisme canonique
\[
\varphi_n:\widehat{H^n(X,\mathcal F)}
\longrightarrow\varprojlim_kH^n(X,\mathcal F_k)
\tag{4.1.7.1}\label{III.4.1.7.1-fr}
\]
(où le premier membre est le séparé complété de $H^n(X,\mathcal F)$ pour
la topologie $\mathfrak J$-préadique) est un isomorphisme. Le système
projectif $(H^n(X,\mathcal F_k))_{k\geq0}$ vérifie pour tout $n$ la
condition (ML), et l'homomorphisme canonique
\[
\psi_n:H^n(\widehat X,\widehat{\mathcal F})
\longrightarrow\varprojlim_kH^n(X,\mathcal F_k)
\tag{4.1.7.2}\label{III.4.1.7.2-fr}
\]

\oldpage[III]{126}

est un isomorphisme. Enfin, la filtration sur $H^n(X,\mathcal F)$, définie
par les noyaux des homomorphismes canoniques
$H^n(X,\mathcal F)\to H^n(X,\mathcal F_k)$, est $\mathfrak J$-bonne
(0, 13.7.7) et $\varphi_n$ est un isomorphisme topologique.
\footnote{La démonstration qui suit, plus simple que la démonstration
initiale, et le complément relatif à la filtration de
$H^n(X,\mathcal F)$, nous ont été communiqués par J.-P. Serre.}
\end{corollary}

\begin{proof}
L'entier $n\geq0$ étant fixé dans cette démonstration, nous poserons pour
simplifier
\[
\mathrm H=H^n(X,\mathcal F),\qquad
\mathrm H_k=H^n(X,\mathcal F_k)
\tag{4.1.7.3}\label{III.4.1.7.3-fr}
\]
et
\[
\mathrm R_k=\operatorname{Ker}(\mathrm H\longrightarrow\mathrm H_k),
\qquad\text{sous-$A$-module de $\mathrm H$.}
\tag{4.1.7.4}\label{III.4.1.7.4-fr}
\]
La suite exacte de cohomologie
\[
H^n(X,\mathfrak J^{k+1}\mathcal F)\longrightarrow H^n(X,\mathcal F)
\longrightarrow H^n(X,\mathcal F_k)\xrightarrow{\partial}
H^{n+1}(X,\mathfrak J^{k+1}\mathcal F)
\longrightarrow H^{n+1}(X,\mathcal F)
\]
montre que l'on a aussi
\[
\mathrm R_k=\operatorname{Im}\bigl(
H^n(X,\mathfrak J^{k+1}\mathcal F)\longrightarrow H^n(X,\mathcal F)
\bigr)~;
\]
nous poserons
\[
\begin{aligned}
\mathrm Q_k
&=\operatorname{Ker}\bigl(H^{n+1}(X,\mathfrak J^{k+1}\mathcal F)
\longrightarrow H^{n+1}(X,\mathcal F)\bigr)\\
&=\operatorname{Im}\bigl(H^n(X,\mathcal F_k)
\longrightarrow H^{n+1}(X,\mathfrak J^{k+1}\mathcal F)\bigr).
\end{aligned}
\tag{4.1.7.5}\label{III.4.1.7.5-fr}
\]
On a donc la suite exacte
\[
0\longrightarrow\mathrm R_k\longrightarrow\mathrm H
\longrightarrow\mathrm H_k\longrightarrow\mathrm Q_k\longrightarrow0.
\tag{4.1.7.6}\label{III.4.1.7.6-fr}
\]

\begin{env}[4.1.7.7]
\label{III.4.1.7.7-fr}
Soit $x$ un élément de $\mathfrak J^m$ ($m\geq0$)~; la multiplication par
$x$ dans $\mathfrak J^k\mathcal F$ est un homomorphisme
$\mathfrak J^k\mathcal F\to\mathfrak J^{k+m}\mathcal F$, et donne lieu
par suite à un homomorphisme
\[
\mu_{x,m}:H^n(X,\mathfrak J^k\mathcal F)
\longrightarrow H^n(X,\mathfrak J^{k+m}\mathcal F).
\tag{4.1.7.8}\label{III.4.1.7.8-fr}
\]
Si on désigne par $S$ la $A$-algèbre graduée
$\bigoplus_{k\geq0}\mathfrak J^k$, on sait que les multiplications
$\mu_{x,m}$ définissent sur
$\mathrm E=\bigoplus_{k\geq0}H^n(X,\mathfrak J^k\mathcal F)$ une
structure de module gradué de type fini sur l'anneau gradué $S$
(\hyperref[III.3.3.2-fr]{3.3.2}), qui est noethérien (II, 2.1.5).
\end{env}

\begin{lemma}[4.1.7.9]
\label{III.4.1.7.9-fr}
Les sous-modules $\mathrm R_k$ de $\mathrm H$ définissent sur $\mathrm H$
une filtration $\mathfrak J$-bonne.
\end{lemma}

\begin{proof}
En premier lieu, montrons que l'on a
\[
\mathfrak J^m\mathrm R_k\subset\mathrm R_{k+m},
\tag{4.1.7.10}\label{III.4.1.7.10-fr}
\]
la multiplication dans $\mathrm H=H^n(X,\mathcal F)$ par un élément
$x\in\mathfrak J^m$ étant donc l'application $\mu_{x,0}$. Pour tout
$x\in\mathfrak J^m$, le diagramme
\[
\xymatrix{
\mathfrak J^{k+1}\mathcal F\ar[r]^x\ar[d]&
\mathfrak J^{k+m+1}\mathcal F\ar[d]\\
\mathcal F\ar[r]_x&\mathcal F
}
\tag{4.1.7.11}\label{III.4.1.7.11-fr}
\]

\oldpage[III]{127}

(où les flèches horizontales sont la multiplication par $x$, et les
flèches verticales les injections canoniques) est commutatif~; donc le
diagramme correspondant
\[
\xymatrix{
H^n(X,\mathfrak J^{k+1}\mathcal F)\ar[r]^{\mu_{x,m}}\ar[d]&
H^n(X,\mathfrak J^{k+m+1}\mathcal F)\ar[d]\\
H^n(X,\mathcal F)\ar[r]_{\mu_{x,0}}&H^n(X,\mathcal F)
}
\]
est commutatif, ce qui, compte tenu de l'interprétation de $\mathrm R_k$
comme image de
$H^n(X,\mathfrak J^{k+1}\mathcal F)\to H^n(X,\mathcal F)$, prouve
(\hyperref[III.4.1.7.10-fr]{4.1.7.10}) et montre en outre que le $S$-module
gradué $\mathrm R=\bigoplus_{k\geq0}\mathrm R_k$ est un quotient du
sous-$S$-module
$\mathrm M=\bigoplus_{k\geq0}H^n(X,\mathfrak J^{k+1}\mathcal F)$ de
$\mathrm E$~; la remarque faite plus haut montre donc que $\mathrm R$ est
un $S$-module de type fini, ce qui équivaut à l'assertion de
(\hyperref[III.4.1.7.9-fr]{4.1.7.9}) (Bourbaki, \emph{Alg. comm.},
chap. III, §~3, n\textsuperscript{o}~1, th.~1).
\end{proof}

\begin{env}[4.1.7.12]
\label{III.4.1.7.12-fr}
Considérons maintenant le $S$-module gradué
\[
\mathrm N=\bigoplus_{k\geq0}H^{n+1}(X,\mathfrak J^{k+1}\mathcal F)
\]
défini comme dans (4.7.1.8)~; c'est encore un $S$-module de type fini en
vertu de (\hyperref[III.3.3.2-fr]{3.3.2})~; on a
$\mathrm Q_k\subset\mathrm N_k$ pour tout $k$ par
(\hyperref[III.4.1.7.5-fr]{4.1.7.5}) et le diagramme
(\hyperref[III.4.1.7.11-fr]{4.1.7.11}) où on remplace $n$ par $n+1$,
montre que
$S_m\mathrm Q_k=\mathfrak J^m\mathrm Q_k\subset\mathrm Q_{k+m}$.
Autrement dit, $\mathrm Q$ est un sous-$S$-module gradué de $\mathrm N$,
et par suite est de type fini.
\end{env}

\begin{env}[4.1.7.13]
\label{III.4.1.7.13-fr}
Désignons par $\alpha_m$ l'injection canonique
$\mathfrak J^m\to A$, qu'on peut écrire $S_m\to S_0$. Comme
$\mathfrak J^{k+1}\mathcal F_k=0$, le $A$-module
$H^n(X,\mathcal F_k)$ est annulé par $\mathfrak J^{k+1}$~; comme
$\mathrm Q_k$ est l'image du $A$-homomorphisme
$H^n(X,\mathcal F_k)\to H^{n+1}(X,\mathfrak J^{k+1}\mathcal F)$,
$\mathrm Q_k$, en tant que $A$-module, est aussi annulé par
$\mathfrak J^{k+1}$~; cela signifie encore que, dans le $S$-module
$\mathrm Q$, on a
\[
\alpha_{k+1}(S_{k+1})\mathrm Q_k=0.
\tag{4.1.7.14}\label{III.4.1.7.14-fr}
\]
Comme $\mathrm Q$ est un $S$-module de type fini, il existe un entier
$k_0$ et un entier $h$ tels que
$\mathrm Q_{k+h}=S_h\mathrm Q_k$ pour $k\geq k_0$ (II, 2.1.6, (ii))~;
on déduit de cette relation et de
(\hyperref[III.4.1.7.14-fr]{4.1.7.14}) qu'il existe un entier $r>0$ tel
que
\[
\alpha_r(S_r)\mathrm Q=0.
\tag{4.1.7.15}\label{III.4.1.7.15-fr}
\]
\end{env}

\begin{env}[4.1.7.16]
\label{III.4.1.7.16-fr}
Notons maintenant que l'injection canonique
$\mathfrak J^{k+m}\mathcal F\to\mathfrak J^k\mathcal F$ donne en passant
à la cohomologie un $A$-homomorphisme
\[
\nu_m:H^{n+1}(X,\mathfrak J^{k+m}\mathcal F)
\longrightarrow H^{n+1}(X,\mathfrak J^k\mathcal F)
\tag{4.1.7.17}\label{III.4.1.7.17-fr}
\]
et, pour tout $x\in\mathfrak J^m$, on a évidemment la factorisation
\[
\xymatrix@C=4.5em{
H^{n+1}(X,\mathfrak J^k\mathcal F)\ar[r]^{\mu_{x,m}}&
H^{n+1}(X,\mathfrak J^{k+m}\mathcal F)\ar[r]^{\nu_m}&
H^{n+1}(X,\mathfrak J^k\mathcal F)
}
\tag{4.1.7.18}\label{III.4.1.7.18-fr}
\]
dont le composé est $\mu_{x,0}$.
\end{env}

\oldpage[III]{128}

d'où on conclut que, pour tout sous-$A$-module $P$ de
$H^{n+1}(X,\mathfrak J^k\mathcal F)$, on a, dans le $S$-module
$\mathrm N$,
\[
\nu_m(S_mP)=\alpha_m(S_m)P.
\tag{4.1.7.19}\label{III.4.1.7.19-fr}
\]

\begin{lemma}[4.1.7.20]
\label{III.4.1.7.20-fr}
Il existe un entier $m>0$ tel que $\nu_m(\mathrm Q_{k+m})=0$ pour tout
$k\geq k_0$.
\end{lemma}

\begin{proof}
Prenons en effet pour $m$ un multiple de $h$ qui soit $\geq r$~; comme
$\mathrm Q_{k+m}=S_m\mathrm Q_k$ pour $k\geq k_0$, on a en vertu de
(\hyperref[III.4.1.7.19-fr]{4.1.7.19}) et
(\hyperref[III.4.1.7.15-fr]{4.1.7.15})
\[
\nu_m(\mathrm Q_{k+m})=\alpha_m(S_m)\mathrm Q_k
\subset\alpha_r(S_r)\mathrm Q_k=0.
\]
\end{proof}

\begin{env}[4.1.7.21]
\label{III.4.1.7.21-fr}
Remarquons que du diagramme commutatif
\[
\xymatrix@C=2.5em{
H^n(X,\mathcal F)\ar[r]&H^n(X,\mathcal F_k)\ar[r]^-\partial&
H^{n+1}(X,\mathfrak J^{k+1}\mathcal F)\ar[r]&H^{n+1}(X,\mathcal F)\\
H^n(X,\mathcal F)\ar[r]\ar[u]&H^n(X,\mathcal F_{k+m})\ar[r]^-\partial\ar[u]&
H^{n+1}(X,\mathfrak J^{k+m+1}\mathcal F)\ar[r]\ar[u]&
H^{n+1}(X,\mathcal F)\ar[u]
}
\]
provenant lui-même du diagramme commutatif
\[
\xymatrix{
0\ar[r]&\mathfrak J^{k+1}\mathcal F\ar[r]&\mathcal F\ar[r]&
\mathcal F_k\ar[r]&0\\
0\ar[r]&\mathfrak J^{k+m+1}\mathcal F\ar[r]\ar[u]&
\mathcal F\ar[r]\ar[u]&\mathcal F_{k+m}\ar[r]\ar[u]&0
}
\]
où les flèches verticales sont les applications canoniques, on déduit un
diagramme commutatif
\[
\xymatrix@C=2.2em{
0\ar[r]&\mathrm R_k\ar[r]&\mathrm H\ar[r]&\mathrm H_k\ar[r]&
\mathrm Q_k\ar[r]&0\\
0\ar[r]&\mathrm R_{k+m}\ar[r]\ar[u]&\mathrm H\ar[r]\ar@{=}[u]&
\mathrm H_{k+m}\ar[r]\ar[u]&\mathrm Q_{k+m}\ar[r]\ar[u]^{\nu_m}&0
}
\]
où les lignes sont exactes. Comme la dernière flèche verticale est nulle
pour $k\geq k_0$ (\hyperref[III.4.1.7.20-fr]{4.1.7.20}), l'image de
$\mathrm H_{k+m}$ dans $\mathrm H_k$ est contenue dans
$\operatorname{Ker}(\mathrm H_k\to\mathrm Q_k)
=\operatorname{Im}(\mathrm H\to\mathrm H_k)$, mais par ailleurs elle
contient $\operatorname{Im}(\mathrm H\to\mathrm H_k)$ par la commutativité
du diagramme, donc elle lui est égale~; il en est donc de même des images
dans $\mathrm H_k$ des $\mathrm H_{k'}$ pour $k'\geq k+m$, ce qui démontre
la condition (ML) pour le système projectif $(\mathrm H_k)_{k\geq0}$. En
outre, pour tout ouvert affine $U$ de $X$, on a
$H^i(U,\mathcal F_k)=0$ pour $i>0$
(\hyperref[III.1.3.1-fr]{1.3.1}), et pour $m>0$, l'application
$H^0(U,\mathcal F_{k+m})\to H^0(U,\mathcal F_k)$ est surjective
(I, 1.3.9). On peut donc appli-

\oldpage[III]{129}

quer (0, 13.3.1), et l'homomorphisme canonique
$H^n(\widehat X,\widehat{\mathcal F})\to
\varprojlim_kH^n(X,\mathcal F_k)$ est bijectif pour tout $n\geq0$.

Comme le système projectif $(\mathrm H/\mathrm R_k)_{k\geq0}$ est strict,
on peut passer à la limite projective dans les suites exactes
\[
0\longrightarrow\mathrm H/\mathrm R_k\longrightarrow\mathrm H_k
\longrightarrow\mathrm Q_k\longrightarrow0
\tag{4.1.7.22}\label{III.4.1.7.22-fr}
\]
(0, 13.2.2)~; comme $\nu_m(\mathrm Q_{k+m})=0$, on a
$\varprojlim_k\mathrm Q_k=0$, d'où un isomorphisme topologique
\[
\varprojlim_k(\mathrm H/\mathrm R_k)\xrightarrow{\sim}
\varprojlim_k\mathrm H_k.
\]
Mais comme la filtration $(\mathrm R_k)$ de $\mathrm H$ est
$\mathfrak J$-bonne, elle définit sur $\mathrm H$ la topologie
$\mathfrak J$-préadique~; donc $\varprojlim_k(\mathrm H/\mathrm R_k)$ est
le séparé complété de $\mathrm H$ pour la topologie
$\mathfrak J$-préadique, ce qui achève de démontrer
(\hyperref[III.4.1.7-fr]{4.1.7}).
\end{env}
\end{proof}

\begin{env}[4.1.8]
\label{III.4.1.8-fr}
Passons enfin à la démonstration de (\hyperref[III.4.1.5-fr]{4.1.5})~:
pour tout ouvert affine $V$ de $Y$,
$\Gamma(V,\widehat{R^nf_*(\mathcal F)})$ est le séparé complété de
$\Gamma(V,R^nf_*(\mathcal F))$ pour la topologie $\mathfrak J$-préadique
(si $\mathcal J|V=\widetilde{\mathfrak J}$) puisque
$R^nf_*(\mathcal F)$ est un $\mathcal O_Y$-Module cohérent (I, 10.8.4),
et
\[
\Gamma\bigl(V,\varprojlim_kR^nf_*(\mathcal F_k)\bigr)
=\varprojlim_k\Gamma(V,R^nf_*(\mathcal F_k))
\]
(0\textsubscript{I}, 3.2.6)~; le fait que $\varphi_n$ soit un isomorphisme
topologique résulte alors de (\hyperref[III.4.1.7-fr]{4.1.7}) et de
(\hyperref[III.1.4.11-fr]{1.4.11}). D'autre part (toujours en vertu de
(\hyperref[III.1.4.11-fr]{1.4.11})), il résulte de
(\hyperref[III.4.1.7-fr]{4.1.7}) que l'homomorphisme $\psi_{n,V}$ de
(\hyperref[III.4.1.3.3-fr]{4.1.3.3}) est un isomorphisme, donc $\psi_n$
est un isomorphisme par définition de
$R^n\widehat f_*(\widehat{\mathcal F})$.
\end{env}

\begin{corollary}[4.1.9]
\label{III.4.1.9-fr}
Sous les hypothèses de (\hyperref[III.4.1.4-fr]{4.1.4}), pour tout ouvert
affine $V$ de $Y$, l'homomorphisme canonique
\[
H^n(\widehat X\cap f^{-1}(V),\widehat{\mathcal F})
\longrightarrow\Gamma(\widehat Y\cap V,
R^n\widehat f_*(\widehat{\mathcal F}))
\]
est bijectif.
\end{corollary}

\begin{env}[4.1.10]
\label{III.4.1.10-fr}
\emph{Remarque.}
Soit $f:X\to Y$ un morphisme de type fini de préschémas (usuels)
noethériens, et soit $\mathcal F$ un $\mathcal O_X$-Module cohérent dont le
support soit propre sur $Y$ (II, 5.4.10). On sait alors
(\hyperref[III.3.2.4-fr]{3.2.4}) que $R^nf_*(\mathcal F)$ est un
$\mathcal O_Y$-Module cohérent pour tout $n\geq0$. En outre, on peut
toujours supposer que $\mathcal F=u_*(\mathcal G)$, où
$\mathcal G=u^*(\mathcal F)$ est un $\mathcal O_Z$-Module cohérent, $Z$
désignant un sous-préschéma fermé convenable de $X$ dont l'espace
sous-jacent est $\operatorname{Supp}(\mathcal F)$, et $u:Z\to X$
l'injection canonique (I, 9.3.5). Si on pose
\[
\mathcal G_k=\mathcal G\otimes_{\mathcal O_Y}
(\mathcal O_Y/\mathcal J^{k+1}),
\]
on a $\mathcal G_k=u^*(\mathcal F_k)$,
$R^nf_*(\mathcal F_k)=R^n(f\circ u)_*(\mathcal G_k)$ et
$R^nf_*(\mathcal F)=R^n(f\circ u)_*(\mathcal G)$
(\hyperref[III.1.3.4-fr]{1.3.4}), et enfin, compte tenu de (I, 10.9.5),
\[
R^n\widehat f_*(\widehat{\mathcal F})
=R^n\widehat{(f\circ u)}_*(\widehat{\mathcal G}).
\]
On peut alors appliquer (\hyperref[III.4.1.5-fr]{4.1.5}) à $\mathcal G$ et
au morphisme propre $f\circ u$, et on en conclut que sous ces hypothèses,
les résultats de (\hyperref[III.4.1.5-fr]{4.1.5}) sont valables pour
$\mathcal F$ et $f$.
\end{env}

\subsection{Cas particuliers et variantes}
\label{subsection:III.4.2-fr}

La forme la plus utile du th. de comparaison
(\hyperref[III.4.1.5-fr]{4.1.5}) est la suivante~:

\begin{proposition}[4.2.1]
\label{III.4.2.1-fr}
Soient $Y$ un préschéma localement noethérien, $f:X\to Y$ un morphisme
propre, $\mathcal F$ un $\mathcal O_X$-Module cohérent. Alors, pour tout
$y\in Y$ et tout $p$, $(R^pf_*(\mathcal F))_y$ est

\oldpage[III]{130}

un $\mathcal O_y$-module de type fini, donc séparé pour la topologie
$\mathfrak m_y$-préadique, et on a un isomorphisme topologique canonique
\[
\widehat{(R^pf_*(\mathcal F))_y}
\xrightarrow{\sim}
\varprojlim_kH^p\bigl(f^{-1}(y),
\mathcal F\otimes_{\mathcal O_Y}(\mathcal O_y/\mathfrak m_y^k)\bigr)
\tag{4.2.1.1}\label{III.4.2.1.1-fr}
\]
où le premier membre est le complété de $(R^pf_*(\mathcal F))_y$ pour la
topologie $\mathfrak m_y$-préadique, et au second membre $f^{-1}(y)$ est
considéré, pour tout $k\geq0$, comme espace sous-jacent au préschéma
$X\times_Y\operatorname{Spec}(\mathcal O_y/\mathfrak m_y^k)$
(I, 3.6.1).
\end{proposition}

\begin{proof}
Comme $\mathcal O_y$ est un anneau local noethérien et
$(R^pf_*(\mathcal F))_y$ un $\mathcal O_y$-module de type fini
(\hyperref[III.3.2.1-fr]{3.2.1}), la topologie $\mathfrak m_y$-préadique
sur $(R^pf_*(\mathcal F))_y$ est séparée (0\textsubscript{I}, 7.3.5). Les
autres assertions sont conséquences de (\hyperref[III.4.1.7-fr]{4.1.7})
lorsque $Y$ est noethérien et le point $y$ fermé, en remplaçant $Y$ par un
voisinage affine de $y$ et prenant $Y'=\{y\}$, compte tenu de
(G, II, 4.9.1). Dans le cas général, posons
\[
Y_1=\operatorname{Spec}(\mathcal O_y),\qquad
X_1=X\times_YY_1,\qquad
\mathcal F_1=\mathcal F\otimes_{\mathcal O_Y}\mathcal O_{Y_1}
\]
et soit $f_1=f\times1_{Y_1}:X_1\to Y_1$~; $Y_1$ est noethérien et $f_1$
propre (II, 5.4.2, (iii)) et $\mathcal F_1$ est cohérent
(0\textsubscript{I}, 5.3.11). Soit $y_1$ l'unique point fermé de $Y_1$~;
la proposition est valable pour $f_1$, $\mathcal F_1$ et $y_1$~; on a
$\mathcal O_{y_1}=\mathcal O_y$, $f_1^{-1}(y_1)=f^{-1}(y)$ (I, 3.6.5),
les préschémas
$X\times_Y\operatorname{Spec}(\mathcal O_y/\mathfrak m_y^k)$ et
$X_1\times_{Y_1}\operatorname{Spec}(\mathcal O_y/\mathfrak m_y^k)$ étant
canoniquement identifiés (I, 3.3.9)~; en outre,
$\mathcal F_1\otimes_{\mathcal O_{Y_1}}
(\mathcal O_y/\mathfrak m_y^k)$ s'identifie à
$\mathcal F\otimes_{\mathcal O_Y}(\mathcal O_y/\mathfrak m_y^k)$
(I, 9.1.6). Il reste à voir que $R^pf_{1*}(\mathcal F_1)$ est
canoniquement isomorphe à
$R^pf_*(\mathcal F)\otimes_{\mathcal O_Y}\mathcal O_{Y_1}$, ce qui
résulte de (\hyperref[III.1.4.15-fr]{1.4.15}), le morphisme local
$\operatorname{Spec}(\mathcal O_y)\to Y$ étant plat
(0\textsubscript{I}, 6.7.1 et I, 2.4.2).
\end{proof}

Le corollaire suivant utilise la terminologie de la théorie de la dimension
(chap. IV) et ne sera pas appliqué avant le chap. IV.

\begin{corollary}[4.2.2]
\label{III.4.2.2-fr}
Soient $Y$ un préschéma localement noethérien, $f:X\to Y$ un morphisme
propre, $y$ un point de $Y$, $r$ la dimension de $f^{-1}(y)$. Alors, pour
tout $\mathcal O_X$-Module cohérent $\mathcal F$, les faisceaux
$R^pf_*(\mathcal F)$ sont nuls au voisinage de $y$ pour tout $p>r$.
\end{corollary}

\begin{proof}
En effet, on a alors
\[
H^p\bigl(f^{-1}(y),
\mathcal F\otimes_{\mathcal O_Y}(\mathcal O_y/\mathfrak m_y^k)\bigr)=0
\]
(G, II, 4.15.2) pour tout $k$, donc
(\hyperref[III.4.2.1-fr]{4.2.1}) le séparé complété de
$(R^pf_*(\mathcal F))_y$ pour la topologie $\mathfrak m_y$-préadique est
nul, et comme cette topologie est séparée, on a aussi
$(R^pf_*(\mathcal F))_y=0$~; d'où la conclusion, puisque
$R^pf_*(\mathcal F)$ est cohérent (0\textsubscript{I}, 5.2.2).
\end{proof}

\begin{env}[4.2.3]
\label{III.4.2.3-fr}
Le résultat (\hyperref[III.4.2.1-fr]{4.2.1}) est surtout employé pour
$p=0$~; on obtient donc le corollaire suivant~:
\end{env}

\begin{corollary}[4.2.4]
\label{III.4.2.4-fr}
Sous les hypothèses de (\hyperref[III.4.2.1-fr]{4.2.1}), on a un
isomorphisme canonique topologique
\[
\widehat{(f_*(\mathcal F))_y}
\xrightarrow{\sim}
\varprojlim_k\Gamma\bigl(f^{-1}(y),
\mathcal F_y/\mathfrak m_y^k\mathcal F_y\bigr).
\]
\end{corollary}

\subsection{Le théorème de connexion de Zariski}
\label{subsection:III.4.3-fr}
\label{III.4.3-fr}

Les résultats de ce numéro et du suivant généralisent des théorèmes bien
connus de Zariski, et peuvent tous se déduire de
(\hyperref[III.4.2.4-fr]{4.2.4}). Ils sont conséquences du th. suivant~:

\begin{theorem}[4.3.1]
\label{III.4.3.1-fr}
\textnormal{(théorème de connexion).}
Soient $Y$ un préschéma localement noethérien, $f:X\to Y$ un morphisme
propre. Alors $\mathcal A(X)=f_*(\mathcal O_X)$ est une
$\mathcal O_Y$-Algèbre cohérente.

\oldpage[III]{131}

Soit $Y'$ le $Y$-schéma fini au-dessus de $Y$ tel que
$\mathcal A(Y')=\mathcal A(X)$, qui est déterminé à un $Y$-isomorphisme
près (II, 1.3.1 et 6.1.3)~; si $f'=\mathcal A(e)$ est le $Y$-morphisme
$X\to Y'$ déduit de l'isomorphisme identique
$e:\mathcal A(Y')\to\mathcal A(X)$ (II, 1.2.7), $f'$ est propre,
$f'_*(\mathcal O_X)$ est isomorphe à $\mathcal O_{Y'}$, et les fibres
$f'^{-1}(y')$ du morphisme $f'$ sont connexes et non vides pour tout
$y'\in Y'$.
\end{theorem}

\begin{proof}
Soit $g:Y'\to Y$ le morphisme structural. Pour prouver que l'homomorphisme
$\theta:\mathcal O_{Y'}\to f'_*(\mathcal O_X)$ entrant dans la définition
du morphisme $f'$ est bijectif, il suffit, puisque $Y'$ est affine sur $Y$,
de prouver que
\[
g_*(\theta):g_*(\mathcal O_{Y'})\longrightarrow
g_*(f'_*(\mathcal O_X))=f_*(\mathcal O_X)
\]
est l'identité (II, 1.4.2)~; mais cela résulte des définitions puisque
$g_*(\mathcal O_{Y'})=\mathcal A(Y')$ et
$f_*(\mathcal O_X)=\mathcal A(X)$. Le fait que $\mathcal A(X)$ est
cohérent est un cas particulier du th. de finitude
(\hyperref[III.3.2.1-fr]{3.2.1}). Comme $f$ est propre et $g$ séparé,
$f'$ est propre (II, 5.4.3, (i))~; pour terminer la démonstration de
(\hyperref[III.4.3.1-fr]{4.3.1}), il suffit donc d'en prouver le
\end{proof}

\begin{corollary}[4.3.2]
\label{III.4.3.2-fr}
Sous les hypothèses de (\hyperref[III.4.3.1-fr]{4.3.1}), supposons de plus
que $f_*(\mathcal O_X)$ soit isomorphe à $\mathcal O_Y$. Alors les fibres
$f^{-1}(y)$ de $f$ sont connexes et non vides pour tout $y\in Y$.
\end{corollary}

\begin{proof}
L'hypothèse que $f_*(\mathcal O_X)$ est isomorphe à $\mathcal O_Y$
entraîne déjà que $f$ est dominant, donc surjectif puisque $f$ est une
application fermée. On peut se ramener, comme dans
(\hyperref[III.4.2.1-fr]{4.2.1}), au cas où $y$ est fermé dans $Y$~;
$f^{-1}(y)$ étant un espace noethérien, a un nombre fini de composantes
connexes, et c'est l'espace sous-jacent du complété $\widehat X$ le long de
$f^{-1}(y)$. Si $Z_i$ ($1\leq i\leq n$) sont ces composantes connexes, il
est clair que $\Gamma(\widehat X,\mathcal O_{\widehat X})$ est composé
direct des anneaux $\Gamma(Z_i,\mathcal O_{\widehat X})$, et chacun de ces
derniers n'est pas réduit à $0$, puisque la section unité est distincte de
$0$ en chaque point de $\widehat X$. Or, si on applique
(\hyperref[III.4.1.5-fr]{4.1.5}) à $\mathcal F=\mathcal O_X$, dont le
complété le long de $f^{-1}(y)$ est $\mathcal O_{\widehat X}$, on voit que
$\Gamma(\widehat X,\mathcal O_{\widehat X})$ est isomorphe au séparé
complété $\mathfrak m_y$-adique $\widehat{\mathcal O}_y$ de l'anneau local
$\mathcal O_y$~; c'est donc un anneau local qui ne peut être composé
direct de plusieurs anneaux non réduits à $0$ (sans quoi il aurait
plusieurs idéaux maximaux distincts). On a donc $n=1$, ce qui démontre le
corollaire.
\end{proof}

\begin{corollary}[4.3.3]
\label{III.4.3.3-fr}
Sous les hypothèses de (\hyperref[III.4.3.1-fr]{4.3.1}), pour tout $y\in Y$,
l'ensemble des composantes connexes de la fibre $f^{-1}(y)$ est en
correspondance biunivoque avec l'ensemble fini des points de la fibre
$g^{-1}(y)$, où $g:Y'\to Y$ est le morphisme structural (autrement dit,
l'ensemble des idéaux maximaux de $(f_*(\mathcal O_X))_y$).
\end{corollary}

\begin{proof}
Puisque $Y'$ est fini sur $Y$, on sait en effet que $g^{-1}(y)$ est un
espace fini discret (II, 6.1.7). Comme
$f^{-1}(y)=f'^{-1}(g^{-1}(y))$, le corollaire résulte de cette remarque et
de (\hyperref[III.4.3.1-fr]{4.3.1}).
\end{proof}

On a ainsi une interprétation remarquable du $Y$-préschéma $Y'$ défini
dans (\hyperref[III.4.3.1-fr]{4.3.1}). La factorisation $f=g\circ f'$ du
morphisme propre $f$ est analogue à la factorisation obtenue par K.~Stein
pour les applications holomorphes d'espaces analytiques, et nous
l'appellerons par la suite la \emph{factorisation de Stein} de $f$.

\begin{remark}[4.3.4]
\label{III.4.3.4-fr}
Soit $k$ une extension du corps $k(y)$~: si le préschéma
\[
f^{-1}(y)\otimes_{k(y)}k=X\times_Y\operatorname{Spec}(k)
\]
est connexe, il en est de même de $f^{-1}(y)$, qui en est l'image par un
morphisme de projection (I, 3.4.7). Nous dirons que, pour un morphisme
$f:X\to Y$ de préschémas et un point $y\in Y$, la fibre $f^{-1}(y)$ est
\emph{géométriquement connexe}, si pour toute extension $k$ de $k(y)$, le
préschéma
$f^{-1}(y)\otimes_{k(y)}k=X\times_Y\operatorname{Spec}(k)$ est connexe.

\oldpage[III]{132}

Sous les hypothèses de (\hyperref[III.4.3.2-fr]{4.3.2}), on peut alors en
renforcer la conclusion~: les fibres $f^{-1}(y)$ sont en effet
\emph{géométriquement connexes}. Pour le voir, observons que pour toute
extension $k$ de $k(y)$, il existe un anneau local noethérien $A$ et un
homomorphisme local $\varphi:\mathcal O_y\to A$ qui fait de $A$ un
$\mathcal O_y$-module \emph{plat} et tel que le corps résiduel de $A$ soit
$k(y)$-isomorphe à $k$ (0, 10.3.1). Soit alors
$Y_1=\operatorname{Spec}(A)$ et soit $h:Y_1\to Y$ le morphisme local
correspondant à $\varphi$, transformant l'unique point fermée $y_1$ de $Y_1$
en $y$ (I, 2.4.1)~; posons
$X_1=X\times_YY_1$, et $f_1=f\times1_{Y_1}$~; $f_1$ est propre
(II, 5.4.2, (iii)) et $f_1^{-1}(y_1)$ est un $k(y_1)$-préschéma isomorphe à
$X\times_Y\operatorname{Spec}(k)$. Tout revient donc à montrer que
$f_{1*}(\mathcal O_{X_1})=\mathcal O_{Y_1}$ pour pouvoir appliquer
(\hyperref[III.4.3.2-fr]{4.3.2}) à $f_1$. Or $g$ est un morphisme plat,
comme il résulte de (I, 2.4.2) et de
(\hyperref[III.1.4.15-fr]{1.4.15.5})~; on a donc
\[
f_{1*}(\mathcal O_{X_1})
=h^*(f_*(\mathcal O_X))
=h^*(\mathcal O_Y)
=\mathcal O_{Y_1}
\]
en vertu de (\hyperref[III.1.4.15-fr]{1.4.15}) appliqué pour $q=0$.

Dans le cas général (\hyperref[III.4.3.1-fr]{4.3.1}), le même raisonnement
montre que l'on a (avec les notations de
(\hyperref[III.4.3.1-fr]{4.3.1}))
\[
f_{1*}(\mathcal O_{X_1})=h^*(g_*(\mathcal O_{Y'})),
\]
et la factorisation de Stein $f_1=g_1\circ f'_1$ de $f_1$ est telle que
$g_1=g\times1_{Y_1}$ (II, 1.5.2), le $Y_1$-schéma fini correspondant étant
$Y'_1=Y'\times_YY_1$. Tenant compte de la transitivité des fibres
(I, 3.6.4), on voit donc que le nombre des composantes connexes de
$f_1^{-1}(y_1)$ est, en vertu de
(\hyperref[III.4.3.3-fr]{4.3.3}), égal au nombre d'éléments de
\[
g_1^{-1}(y_1)=g^{-1}(y)\otimes_{k(y)}k.
\]
Si on prend pour $k$ une extension algébriquement close de $k(y)$, ce nombre
est indépendant de l'extension algébriquement close considérée et égal au
\emph{nombre géométrique de points} de $g^{-1}(y)$ (I, 6.4.7), ou encore à
la \emph{somme des rangs séparables}
\[
\sum_i[k(y'_i):k(y)]_s
\]
où $y'_i$ parcourt l'ensemble fini $g^{-1}(y)$. On dit encore que ce nombre
est le \emph{nombre géométrique de composantes connexes} de $f^{-1}(y)$. On
notera que les $k(y'_i)$ ne sont autres que les corps résiduels de l'anneau
semi-local $(f_*(\mathcal O_X))_y$.
\end{remark}

\begin{proposition}[4.3.5]
\label{III.4.3.5-fr}
Soient $X$ et $Y$ deux préschémas localement noethériens intègres et
$f:X\to Y$ un morphisme propre dominant. Pour tout $y\in Y$, le nombre de
composantes connexes de $f^{-1}(y)$ est au plus égal au nombre des idéaux
maximaux de la fermeture intégrale $\mathcal O'_y$ de $\mathcal O_y$ dans le
corps des fonctions rationnelles $\mathrm R(X)$.
\end{proposition}

\begin{proof}
En effet, pour tout ouvert $U$ de $Y$,
\[
\Gamma(U,f_*(\mathcal O_X))=\Gamma(f^{-1}(U),\mathcal O_X)
\]
est l'intersection des anneaux locaux $\mathcal O_x$ tels que
$x\in f^{-1}(U)$ (I, 8.2.1.1). On en conclut aussitôt que la fibre
$(f_*(\mathcal O_X))_y$ est un sous-anneau de $\mathrm R(X)$ contenant
$\mathcal O_y$. En outre, comme $f_*(\mathcal O_X)$ est un
$\mathcal O_Y$-Module cohérent, $(f_*(\mathcal O_X))_y$ est un
$\mathcal O_y$-module de type fini, donc contenu dans $\mathcal O'_y$~; on
sait ([13], vol. I, p. 257 et 259) que tout idéal maximal d'un tel anneau
$A$ est l'intersection de $A$ et d'un idéal maximal de $\mathcal O'_y$,
d'où la proposition.
\end{proof}

\begin{definition}[4.3.6]
\label{III.4.3.6-fr}
On dit qu'un anneau local intègre est \emph{unibranche} si sa clôture
intégrale est un anneau local. On dit qu'un point $y$ d'un préschéma intègre
$Y$ est unibranche si l'anneau local $\mathcal O_y$ est unibranche (ce qui
est en particulier le cas lorsque $Y$ est normal au point $y$).
\end{definition}

Soit $A$ un anneau local intègre, et soit $K$ son corps des fractions~; pour
que $A$ soit unibranche, il faut et il suffit que tout sous-anneau $A_1$ de
$K$, contenant $A$ et qui est une $A$-algèbre finie, soit un anneau local.
En effet, soit $A'$ la clôture intégrale de $A$~; il résulte du premier
théorème de Cohen-Seidenberg (Bourbaki, \emph{Alg. comm.}, chap. V, \S~2,
n\textsuperscript{o}~1, th.~1) que tout idéal maximal de $A_1$ est la trace
d'un idéal maximal de $A'$, donc si $A'$ est local, il en est de même de
$A_1$. Réciproquement, $A'$ est limite inductive

\oldpage[III]{133}

de la famille filtrante croissante des sous-$A$-algèbres finies $A_\alpha$
de $A'$, et si chacune des $A_\alpha$ est un anneau local, l'idéal maximal
de $A_\alpha$ est la trace sur $A_\alpha$ de celui de $A_\beta$ pour
$A_\alpha\subset A_\beta$ par le même raisonnement que ci-dessus, donc $A'$
est un anneau local (0, 10.3.1.3).

On notera que si le complété d'un anneau local noethérien $A$ est intègre
(ce qu'on exprime en disant que $A$ est \emph{analytiquement intègre}), $A$
est unibranche. Soient en effet $\mathfrak m$ l'idéal maximal de $A$, $K$
son corps des fractions, $K'$ le corps des fractions de $\widehat A$~; on
a donc
\[
K'=K\otimes_A\widehat A.
\]
Soit $A'_F$ une sous-$A$-algèbre finie de $K$. Le sous-anneau $B_F$ de $K'$
engendré par $\widehat A$ et $A'_F$ est isomorphe à
$A'_F\otimes_A\widehat A$~; c'est un $\widehat A$-module de type fini,
complété de $A'_F$ pour la topologie $\mathfrak m$-adique
(0\textsubscript{I}, 7.3.3 et 7.3.6). Comme $A'_F$ est un anneau semi-local
(Bourbaki, \emph{Alg. comm.}, chap. IV, \S~2, n\textsuperscript{o}~5,
cor. 3 de la prop. 9) et que son complété est intègre, $A'_F$ ne peut avoir
qu'un seul idéal maximal $\mathfrak m'_F$ et on a
$\mathfrak m'_F\cap A=\mathfrak m$~; d'où notre assertion.

\begin{corollary}[4.3.7]
\label{III.4.3.7-fr}
Sous les hypothèses de (\hyperref[III.4.3.5-fr]{4.3.5}), supposons que la
fermeture algébrique de $\mathrm R(Y)$ dans $\mathrm R(X)$ soit de degré
séparable $n$ et que $y\in Y$ soit unibranche. Alors la fibre $f^{-1}(y)$ a
au plus $n$ composantes connexes. En particulier si la fermeture algébrique
de $\mathrm R(Y)$ dans $\mathrm R(X)$ est radicielle sur $\mathrm R(X)$,
$f^{-1}(y)$ est connexe.
\end{corollary}

\begin{proof}
En effet, soit $\mathcal O'_y$ la clôture intégrale de $\mathcal O_y$~; la
fermeture intégrale $\mathcal O''_y$ de $\mathcal O_y$ dans $\mathrm R(X)$
est aussi celle de $\mathcal O'_y$~; mais on sait que si $\mathcal O'_y$ est
un anneau local, $\mathcal O''_y$ est un anneau semi-local dont le nombre
d'idéaux maximaux est au plus égal à $n$ ([13], vol. I, p. 289, th. 22).

Ce corollaire est essentiellement la forme sous laquelle Zariski énonce son
«~théorème de connexion~» pour les schémas algébriques.
\end{proof}

\begin{remark}[4.3.8]
\label{III.4.3.8-fr}
Si on ajoute aux hypothèses de (\hyperref[III.4.3.7-fr]{4.3.7}) l'hypothèse
que $Y$ est normal au point $y$, la fibre $f^{-1}(y)$ est
\emph{géométriquement connexe}, puisque (avec les notations de
(\hyperref[III.4.3.4-fr]{4.3.4})) $g^{-1}(y)$ est réduit à un point $y'$ et
que $k(y')$ est radiciel sur $k(y)$.
\end{remark}

\begin{definition}[4.3.9]
\label{III.4.3.9-fr}
Étant donné un préschéma localement noethérien $Y$, on dit qu'un morphisme de
type fini $f:X\to Y$ est \emph{universellement ouvert} si, pour tout
préschéma irréductible localement noethérien $Y'$, et tout morphisme dominant
$g:Y'\to Y$, toute composante irréductible de $X'=X\times_YY'$ domine $Y'$.
\end{definition}

Si $Y$ est irréductible, cela revient à dire que si $\eta$, $\eta'$ sont les
points génériques de $Y$ et $Y'$ respectivement (de sorte que
$g(\eta')=\eta$), et si on pose $f'=f_{(Y')}$, toute composante irréductible
de $X'$ rencontre $f'^{-1}(\eta')$ (0\textsubscript{I}, 2.1.8)~; cela
implique donc que pour tout ouvert $U$ de $Y$, le morphisme
$f^{-1}(U)\to U$, restriction de $f$, est universellement ouvert.

\begin{corollary}[4.3.10]
\label{III.4.3.10-fr}
Soient $X$, $Y$ deux préschémas localement noethériens intègres.
$f:X\to Y$ un morphisme propre dominant et universellement ouvert. Si la
fermeture algébrique de $\mathrm R(Y)$ dans $\mathrm R(X)$ est radicielle
sur $\mathrm R(Y)$, toute fibre $f^{-1}(y)$ ($y\in Y$) est géométriquement
connexe.
\end{corollary}

\begin{proof}
On peut se borner au cas où $Y=\operatorname{Spec}(B)$, $B$ étant un anneau
intègre noethérien. Il résulte alors de (II, 7.1.7) qu'il existe un anneau
local noethérien intégralement clos $A$ qui domine $\mathcal O_y$ et a
$\mathrm R(Y)$ pour corps des fractions. Soit
$Y'=\operatorname{Spec}(A)$, et soit $h:Y'\to Y$ le morphisme correspondant
à l'injection canonique $B\to A$, qui est birationnel (donc dominant)~; en
outre, si $y'$ est l'unique point fermé de $Y'$, on a $h(y')=y$. Soient
$X'=X\times_YY'$, $f'=f\times1_{Y'}$~; désignons par $\eta$, $\eta'$,
$\xi$ les points génériques de $Y$, $Y'$ et $X$ respectivement, de sorte
que $f(\xi)=\eta$ et $h^{-1}(\eta)=\{\eta'\}$~; en outre,
$k(\eta)=k(\eta')=\mathrm R(Y)$, donc $f'^{-1}(\eta')$ est isomorphe à
$f^{-1}(\eta)$ (I, 3.6.4) et en particulier, puisque $\xi$ est le point
générique de $f^{-1}(\eta)$ (0\textsubscript{I}, 2.1.8),
$f'^{-1}(\eta')$ a un seul point générique. Mais par hypothèse, toute
composante irréductible de $X'$ a son point générique dans
$f'^{-1}(\eta')$, donc $X'$ est nécessairement irréductible, son point
générique $\xi'$ est le point générique de $f'^{-1}(\eta')$ et on a
$k(\xi')=k(\xi)$.

\oldpage[III]{134}

Posons $X''=X'_{\mathrm{red}}$, $f''=f'_{\mathrm{red}}$~; $X''$ est donc
intègre et noethérien, $f''$ est propre (II, 5.4.6) et les espaces
sous-jacents aux fibres $f'^{-1}(y')$ et $f''^{-1}(y')$ sont les mêmes~; en
outre, $\mathrm R(X'')=k(\xi')=\mathrm R(X)$, donc $f''$ vérifie les
hypothèses de (\hyperref[III.4.3.8-fr]{4.3.8}), et $f''^{-1}(y')$ est
géométriquement connexe. Soit maintenant $k$ une extension quelconque de
$k(y)$~; il existe une extension $k_1$ de $k(y)$ dont $k(y')$ et $k$ peuvent
être considérés comme des sous-extensions (Bourbaki, \emph{Alg.}, chap. V,
\S~4, prop. 2). Par hypothèse,
$f''^{-1}(y')\times_{Y'}\operatorname{Spec}(k_1)$ est connexe, et il a même
préschéma réduit que
$f'^{-1}(y')\times_{Y'}\operatorname{Spec}(k_1)$ (I, 5.1.8), donc ce dernier
est connexe, et comme il est isomorphe à
$f^{-1}(y)\times_Y\operatorname{Spec}(k_1)$ (I, 3.6.4), on en conclut que ce
dernier est connexe~; \emph{a fortiori}, il en est de même de
$f^{-1}(y)\times_Y\operatorname{Spec}(k)$ d'après la remarque du début de
(\hyperref[III.4.3.4-fr]{4.3.4}), ce qui achève la démonstration.
\end{proof}

\begin{env}[4.3.11]
\label{III.4.3.11-fr}
\emph{Remarques} (4.3.11). ---
\begin{enumerate}
\item[(i)] Le raisonnement précédent est dû en substance à Zariski [20], à
cela près qu'il peut prendre pour $A$ la clôture intégrale de
$\mathcal O_y$, celle-ci étant un anneau noethérien pour les anneaux locaux
de la géométrie algébrique classique. D'autre part, Zariski prouve que si
$Y$ est la variété de Chow d'un espace projectif $\mathbf P_k^r$ sur un
corps $k$, et si $X$ est la partie fermée de $\mathbf P_k^r\times_kY$ qui
définit la correspondance de Chow entre $\mathbf P_k^r$ et $Y$, alors la
projection $X\to Y$ est un morphisme universellement ouvert
(\emph{loc. cit.}, lemme de la p. 82). Il semble bien que ce soit la seule
propriété formelle des «~coordonnées de Chow~» dont on se soit servi dans
certaines applications~; il y a par suite intérêt dans une telle situation,
à substituer le langage~: fibres d'un morphisme propre (éventuellement
supposé universellement ouvert ou soumis à d'autres restrictions analogues
de régularité locale) au langage~: spécialisation de cycles dans l'espace
projectif.

\item[(ii)] Au chap. IV, nous verrons qu'un morphisme universellement ouvert
$f:X\to Y$ peut encore être défini de la façon suivante (qui justifie la
terminologie)~: pour tout morphisme $Z\to Y$, le morphisme
$f_{(Z)}:X_{(Z)}\to Z$ est ouvert. On peut montrer en outre que si $f$
vérifie les hypothèses de (\hyperref[III.4.3.10-fr]{4.3.10}), alors, si
$y,y'$ sont deux points de $Y$ tels que $y$ soit une spécialisation de $y'$,
le nombre géométrique de composantes connexes de $f^{-1}(y)$ est au plus
égal à celui des composantes connexes de $f^{-1}(y')$.
\end{enumerate}
\end{env}

\begin{corollary}[4.3.12]
\label{III.4.3.12-fr}
Sous les hypothèses de (\hyperref[III.4.3.5-fr]{4.3.5}), supposons en outre
$\mathrm R(Y)$ algébriquement fermé dans $\mathrm R(X)$, et soit $y$ un
point normal de $Y$. Alors $f^{-1}(y)$ est géométriquement connexe, et il
existe un voisinage ouvert $U$ de $y$ dans $Y$ tel que
$f_*(\mathcal O_X|f^{-1}(U))$ soit isomorphe à $\mathcal O_Y|U$. Si plus
particulièrement, on suppose $Y$ normal (et $\mathrm R(Y)$ algébriquement
fermé dans $\mathrm R(X)$), $f_*(\mathcal O_X)$ est isomorphe à
$\mathcal O_Y$.
\end{corollary}

\begin{proof}
La première assertion relative à $f^{-1}(y)$ est un cas particulier de
(\hyperref[III.4.3.8-fr]{4.3.8}). On en

\oldpage[III]{135}

déduit que si $X\xrightarrow{f'}Y'\xrightarrow{g}Y$ est la factorisation
de Stein de $f$ (\hyperref[III.4.3.3-fr]{4.3.3}), $g^{-1}(y)$ se réduit à
un seul point $y'$~; en outre, on a
\[
\mathcal O_y\subset\mathcal O_{y'}=(f_*(\mathcal O_X))_y
\subset\mathrm R(X),
\]
et comme $\mathcal O_{y'}$ est fini sur $\mathcal O_y$ (et \emph{a
fortiori} sur $\mathrm R(Y)$), il est contenu dans $\mathrm R(Y)$ en vertu
de l'hypothèse~; comme $y$ est normal, on a nécessairement
$\mathcal O_{y'}=\mathcal O_y$~; on en conclut que $g$ est un isomorphisme
local au point $y'$ (I, 6.5.4), ce qui achève de prouver la première partie
du corollaire. La seconde résulte de la première, car l'hypothèse
supplémentaire entraîne que $g$ est bijective et un isomorphisme local au
voisinage de tout point de $Y'$, donc un isomorphisme.
\end{proof}

Le fait que (\hyperref[III.4.3.7-fr]{4.3.7}) soit établi dans le cadre des
schémas permet des applications telles que la suivante~:

\begin{proposition}[4.3.13]
\label{III.4.3.13-fr}
Soient $A$ un anneau local noethérien unibranche, $\mathfrak a$ un idéal de
définition de $A$, $A_0=A/\mathfrak a$,
$S=\operatorname{gr}_{\mathfrak a}(A)$ l'anneau gradué associé à $A$ pour
la filtration $\mathfrak a$-préadique~; $S$ est une $A_0$-algèbre graduée
engendrée par $S_1$, $S_1$ étant un $A_0$-module de type fini. Alors
$\operatorname{Proj}(S)$ est un $A_0$-schéma connexe.
\end{proposition}

\begin{proof}
Soit $\mathfrak m$ l'idéal maximal de $A$~; $Y=\operatorname{Spec}(A)$ est
un schéma intègre dont le point $y$ correspondant à $\mathfrak m$ est
l'unique point fermé. Par hypothèse, on a
$\mathfrak m^p\subset\mathfrak a\subset\mathfrak m$ pour un entier $p$,
donc $V(\mathfrak a)=\{\mathfrak m\}$. Soit
$S'=\bigoplus_{n\geq0}\mathfrak a^n$, et soit
$X=\operatorname{Proj}(S')$, qui est le $Y$-schéma obtenu en faisant éclater
l'idéal $\mathfrak a$~; $X$ est intègre et le morphisme structural
$f:X\to Y$ est birationnel (II, 8.1.4) et évidemment projectif. Par suite,
(\hyperref[III.4.3.7-fr]{4.3.7}) est applicable et montre que $f^{-1}(y)$
est connexe~; mais l'espace $f^{-1}(y)$ est sous-jacent à
$\operatorname{Proj}(S'\otimes_AA_0)$ (I, 3.6.1 et II, 2.8.10)~; comme
$S'\otimes_AA_0=S$ par définition, la proposition est démontrée.
\end{proof}


\subsection{Le «~main theorem~» de Zariski}
\label{subsection:III.4.4-fr}

\begin{proposition}[4.4.1]
\label{III.4.4.1-fr}
Soient $Y$ un préschéma localement noethérien, $f:X\to Y$ un morphisme
propre. Soit $X'$ l'ensemble des points $x\in X$ qui sont isolés dans leur
fibre $f^{-1}(f(x))$. Alors l'ensemble $X'$ est ouvert dans $X$, et si
$f=gf'$ est la factorisation de Stein de $f$
(\hyperref[III.4.3.3-fr]{4.3.3}), la restriction de $f'$ à $X'$ est un
isomorphisme de $X'$ sur un sous-préschéma induit sur un ouvert $U$ de $Y'$,
et on a $X'=f'^{-1}(U)$.
\end{proposition}

\begin{proof}
Comme $g^{-1}(f(x))$ est fini et discret
(\hyperref[III.4.3.3-fr]{4.3.3} et II, 6.1.7), pour que $x$ soit isolé dans
$f^{-1}(f(x))$, il faut et il suffit qu'il le soit dans
$f'^{-1}(f'(x))$~; on peut donc se borner au cas où $f'=f$, donc
$f_*(\mathcal O_X)=\mathcal O_Y$. Alors, si $x\in X'$,
$f^{-1}(f(x))$, qui est connexe (\hyperref[III.4.3.2-fr]{4.3.2}), est
nécessairement réduit au point $x$. Comme $f$ est fermé, pour tout voisinage
ouvert $V$ de $x$ dans $X$, $f(X-V)$ est fermé dans $Y$ et ne contient pas
$y=f(x)$, puisque $f^{-1}(y)=\{x\}$~; si $U$ est le complémentaire de
$f(X-V)$ dans $Y$, on a $f^{-1}(U)\subset V$, et on en conclut que les
images réciproques par $f$ d'un système fondamental de voisinages ouverts de
$y$ forment un système fondamental de voisinages ouverts de $x$. L'hypothèse
$f_*(\mathcal O_X)=\mathcal O_Y$ et la définition de l'image directe d'un
faisceau (0\textsubscript{I}, 3.4.1 et 4.2.1) entraînent alors que, si
$f=(\psi,\theta)$, l'homomorphisme
$\theta_y^\#:\mathcal O_y\to\mathcal O_x$ est un isomorphisme. On en conclut
qu'il existe un voisinage ouvert $V$ de $x$ et un voisinage ouvert $U$ de
$y$ tels que la restriction de $f$ à $V$ soit un isomorphisme de $V$ sur
$U$ (I, 6.5.4)~; en outre, d'après ce qu'on vient de voir,

\oldpage[III]{136}

on peut supposer que $f^{-1}(U)=V$, d'où on conclut aussitôt, par définition,
que $V\subset X'$, ce qui achève la démonstration.
\end{proof}

La proposition suivante a été démontrée par Chevalley dans le cas des schémas
algébriques~:

\begin{proposition}[4.4.2]
\label{III.4.4.2-fr}
Soient $Y$ un préschéma localement noethérien, $f:X\to Y$ un morphisme. Les
conditions suivantes sont équivalentes~:
\begin{enumerate}
\item[a)] $f$ est fini.
\item[b)] $f$ est affine et propre.
\item[c)] $f$ est propre et, pour tout $y\in Y$, $f^{-1}(y)$ est un ensemble
fini.
\end{enumerate}
\end{proposition}

\begin{proof}
On sait que a) entraîne b) (II, 6.1.2 et 6.1.11). Si $f$ est propre et affine,
il en est de même du morphisme
$f^{-1}(y)\to\operatorname{Spec}(k(y))$ (II, 1.6.2, (iii) et 5.4.2, (iii)),
et le th. de finitude (\hyperref[III.3.2.1-fr]{3.2.1}) appliqué au faisceau
structural de $f^{-1}(y)$, montre que $f^{-1}(y)=\operatorname{Spec}(A)$, où
$A$ est une $k(y)$-algèbre finie~; donc $f^{-1}(y)$ est un ensemble fini
(II, 6.1.7), et on voit que b) entraîne c). Enfin, comme $f^{-1}(y)$ est un
préschéma algébrique sur $k(y)$, l'hypothèse que l'ensemble $f^{-1}(y)$ est
fini entraîne que l'espace $f^{-1}(y)$ est discret (I, 6.4.4). Avec les
notations de (\hyperref[III.4.4.1-fr]{4.4.1}), on a donc $X'=X$, et
$f':X\to Y'$ est un isomorphisme~; comme $g$ est un morphisme fini, on voit
que c) entraîne a).
\end{proof}

\begin{theorem}[4.4.3]
\label{III.4.4.3-fr}
\textnormal{(«~Main theorem~» de Zariski).}
Soient $Y$ un préschéma noethérien, $f:X\to Y$ un morphisme quasi-projectif,
$X'$ l'ensemble des points $x\in X$ qui sont isolés dans leur fibre
$f^{-1}(f(x))$. Alors $X'$ est une partie ouverte de $X$, et le
sous-préschéma induit $X'$ est isomorphe à un préschéma induit sur une partie
ouverte d'un $Y$-préschéma $Y'$ fini sur $Y$.
\end{theorem}

\begin{proof}
L'hypothèse entraîne qu'il existe un $Y$-préschéma projectif $Z$ tel que $X$
soit $Y$-isomorphe à un sous-préschéma induit sur un ouvert de $Z$
(II, 5.3.2 et 5.5.1). On est donc ramené à démontrer le théorème lorsque $f$
est un morphisme projectif, donc propre (II, 5.5.3), et il résulte alors
aussitôt de (\hyperref[III.4.4.1-fr]{4.4.1}).
\end{proof}

\begin{remark}[4.4.4]
\label{III.4.4.4-fr}
Si $X$ est réduit (resp. irréductible et $X'$ non vide), on peut supposer,
dans l'énoncé de (\hyperref[III.4.4.3-fr]{4.4.3}), que $Y'$ est réduit (resp.
irréductible). En effet, on peut toujours remplacer $Y'$ par le
sous-préschéma adhérence $\overline{X'}$ de $X'$ dans $Y'$ (I, 9.5.11 et
II, 6.1.5, (i) et (ii)), et on sait que si $X'$ est réduit, il en est de même
de $\overline{X'}$ (I, 9.5.9, (i))~; par ailleurs, si $X'$ n'est pas vide, il
est irréductible si $X$ l'est, et $\overline{X'}$ est alors aussi
irréductible.
\end{remark}

\begin{corollary}[4.4.5]
\label{III.4.4.5-fr}
Soient $Y$ un schéma localement noethérien, $f:X\to Y$ un morphisme de type
fini, $x$ un point de $X$ isolé dans sa fibre $f^{-1}(f(x))$. Alors il existe
un voisinage ouvert de $x$ dans $X$ qui est isomorphe à une partie ouverte
d'un $Y$-préschéma fini sur $Y$.
\end{corollary}

\begin{proof}
Soient en effet $y=f(x)$, $U$ un voisinage ouvert affine de $y$ dans $Y$,
$V$ un voisinage ouvert affine de $x$ dans $X$, contenu dans $f^{-1}(U)$.
Comme $Y$ est séparé, l'injection $U\to Y$ est affine (II, 1.6.3), et comme
$V$ est affine sur $U$ (\emph{ibid.}), la restriction de $f$ à $V$ est un
morphisme affine $V\to Y$ (II, 1.6.2, (ii))~; \emph{a fortiori}, cette
restriction est un morphisme quasi-projectif puisqu'il est de type fini
(I, 6.3.5 et II, 5.3.4, (i)). Il suffit alors d'appliquer à cette restriction
le th. (\hyperref[III.4.4.3-fr]{4.4.3}).
\end{proof}

Le cor. (\hyperref[III.4.4.5-fr]{4.4.5}) s'énonce dans le langage de
l'algèbre commutative~:

\oldpage[III]{137}

\begin{corollary}[4.4.6]
\label{III.4.4.6-fr}
Soient $A$ un anneau noethérien, $B$ une $A$-algèbre de type fini,
$\mathfrak q$ un idéal premier de $B$, $\mathfrak p$ son image réciproque
dans $A$. On suppose que $\mathfrak q$ soit à la fois maximal et minimal dans
l'ensemble des idéaux premiers de $B$ dont l'image réciproque est
$\mathfrak p$. Alors il existe $g\in B-\mathfrak q$, une $A$-algèbre finie
$A'$ et un élément $f'\in A'$ tels que les $A$-algèbres $B_g$ et $A'_{f'}$
soient isomorphes.
\end{corollary}

\begin{proof}
Il suffit en effet d'appliquer (\hyperref[III.4.4.5-fr]{4.4.5}) à
$Y=\operatorname{Spec}(A)$ et $X=\operatorname{Spec}(B)$, l'hypothèse sur
$\mathfrak q$ signifiant exactement qu'il est isolé dans sa fibre
(I, 1.1.7).
\end{proof}

On en déduit le résultat suivant, moins général en apparence~:

\begin{corollary}[4.4.7]
\label{III.4.4.7-fr}
Soient $A$ un anneau local noethérien, $B$ une $A$-algèbre de type fini,
$\mathfrak n$ un idéal premier de $B$ dont l'image réciproque dans $A$ est
l'idéal maximal $\mathfrak m$. On suppose que $\mathfrak n$ est maximal dans
$B$ et est minimal dans l'ensemble des idéaux premiers de $B$ dont l'image
réciproque est $\mathfrak m$ (ce qui signifie aussi que
$B\mathfrak m$ est primaire pour $\mathfrak n$). Alors il existe une
$A$-algèbre finie $A'$ et un idéal maximal $\mathfrak m'$ de $A'$ (dont
$\mathfrak m$ est l'image réciproque dans $A$) tels que $B_{\mathfrak n}$
soit isomorphe à la $A$-algèbre $A'_{\mathfrak m'}$.
\end{corollary}

Le cas particulier suivant de (\hyperref[III.4.4.7-fr]{4.4.7}) est aussi
parfois appelé «~Main Theorem~»~:

\begin{corollary}[4.4.8]
\label{III.4.4.8-fr}
Sous les conditions de (\hyperref[III.4.4.7-fr]{4.4.7}), supposons en outre
$A$ et $B$ intègres et ayant même corps des fractions $K$. Alors, si $A$ est
intégralement clos, on a $B=A$.
\end{corollary}

\begin{proof}
En effet, la Remarque (\hyperref[III.4.4.4-fr]{4.4.4}) montre que l'on peut
supposer, dans l'application de (\hyperref[III.4.4.7-fr]{4.4.7}) que $A'$ est
intègre et a $K$ pour corps des fractions~; l'hypothèse sur $A$ entraîne alors
$A'=A$, donc $B_{\mathfrak n}=A$~; comme on a
$A\subset B\subset B_{\mathfrak n}$, on en conclut bien $A=B$.
\end{proof}

L'énoncé (\hyperref[III.4.4.8-fr]{4.4.8}) est la forme donnée par Zariski à
son «~Main theorem~» (étendu aux anneaux locaux noethériens intègres
quelconques).

Les corollaires précédents étaient des variantes de nature locale de
(\hyperref[III.4.4.3-fr]{4.4.3}), qui est un résultat global. Voici une autre
conséquence de nature globale~:

\begin{corollary}[4.4.9]
\label{III.4.4.9-fr}
Soient $Y$ un préschéma intègre localement noethérien, $f:X\to Y$ un
morphisme séparé, de type fini et birationnel. Supposons en outre $Y$ normal
et toutes les fibres $f^{-1}(y)$ finies pour $y\in Y$. Alors $f$ est une
immersion ouverte~; si en outre $f$ est fermé (et en particulier si $f$ est
propre), $f$ est un isomorphisme.
\end{corollary}

\begin{proof}
Soit en effet $x\in X$, et posons $y=f(x)$. Comme $f^{-1}(y)$ est un schéma
algébrique sur $k(y)$, l'hypothèse qu'il est fini entraîne qu'il est discret
(I, 6.4.4)~; en outre $\mathcal O_y$ est intégralement clos et
$\mathcal O_x$ et $\mathcal O_y$ ont même corps des fractions (I, 7.1.5).
On peut donc appliquer (\hyperref[III.4.4.8-fr]{4.4.8}), et si
$f=(\psi,\theta)$, l'homomorphisme
$\theta_y^\#:\mathcal O_y\to\mathcal O_x$ est bijectif~; on en conclut
(I, 6.5.4) que $f$ est un isomorphisme local. Mais comme $f$ est séparé et
$X$ intègre, $f$ est une immersion ouverte (I, 8.2.8). La dernière assertion
résulte de ce que $f$ est dominant.
\end{proof}

\begin{proposition}[4.4.10]
\label{III.4.4.10-fr}
Soient $Y$ un préschéma localement noethérien, $f:X\to Y$ un morphisme
localement de type fini. L'ensemble $X'$ des $x\in X$ isolés dans leur fibre
$f^{-1}(f(x))$ est ouvert dans $X$.
\end{proposition}

\begin{proof}
La question étant locale sur $X$ et $Y$, on peut supposer $X$ et $Y$ affines
noethériens et $f$ de type fini~; $f$ est alors un morphisme affine de type
fini, donc quasi-projectif (II, 5.3.4, (i)), et il suffit d'appliquer
(\hyperref[III.4.4.3-fr]{4.4.3}).
\end{proof}

\oldpage[III]{138}

\begin{corollary}[4.4.11]
\label{III.4.4.11-fr}
Soient $Y$ un préschéma localement noethérien, $f:X\to Y$ un morphisme
propre. L'ensemble $U$ des points $y\in Y$ tels que $f^{-1}(y)$ soit discret
est ouvert dans $Y$, et le morphisme $f^{-1}(U)\to U$ restriction de $f$ est
fini. En particulier, un morphisme propre et quasi-fini $X\to Y$ est fini.
\end{corollary}

\begin{proof}
En effet, le complémentaire de $U$ dans $Y$ est l'image par $f$ de $X-X'$ qui
est fermé dans $X$ en vertu de (\hyperref[III.4.4.10-fr]{4.4.10})~; comme $f$
est une application fermée, $U$ est ouvert. En outre, il résulte de
(II, 6.2.2) que $f^{-1}(y)$ est fini pour tout $y\in U$~; comme le morphisme
$f^{-1}(U)\to U$, restriction de $f$ est propre (II, 5.4.1), il est fini en
vertu de (\hyperref[III.4.4.2-fr]{4.4.2}).
\end{proof}

\begin{env}[4.4.12]
\label{III.4.4.12-fr}
\emph{Remarques} (4.4.12). ---
\begin{enumerate}
\item[(i)] Comme on l'a annoncé dans (II, 6.2.7), nous montrerons au chap. V
que si $Y$ est localement noethérien, tout morphisme quasi-fini et séparé
$f:X\to Y$ est quasi-affine, donc quasi-projectif. Il s'ensuivra que, dans le
Main Theorem (\hyperref[III.4.4.3-fr]{4.4.3}) la conclusion reste valable
lorsqu'on suppose seulement $f$ séparé et de type fini. En effet, il résulte
de (\hyperref[III.4.4.10-fr]{4.4.10}) que $X'$ est ouvert dans $X$, et comme
$X$ est localement noethérien, la restriction de $f$ à $X'$ est encore de
type fini (I, 6.3.5), donc quasi-fini par définition de $X'$, et évidemment
séparé~; on peut donc appliquer (\hyperref[III.4.4.3-fr]{4.4.3}) à cette
restriction, d'où la conclusion.

\item[(ii)] Nous donnerons au chap. IV une démonstration plus élémentaire de
(\hyperref[III.4.4.10-fr]{4.4.10}), utilisant la théorie de la dimension.
\end{enumerate}
\end{env}

\subsection{Complétés de modules d'homomorphismes}
\label{subsection:III.4.5-fr}

\begin{proposition}[4.5.1]
\label{III.4.5.1-fr}
Soient $A$ un anneau noethérien, $\mathfrak J$ un idéal de $A$, $X$ un
$A$-préschéma de type fini, $\mathcal F$, $\mathcal G$ deux
$\mathcal O_X$-Modules cohérents dont les supports ont une intersection
propre sur $Y=\operatorname{Spec}(A)$ (II, 5.4.10). Alors, pour tout entier
$n\geq0$, $\operatorname{Ext}^n_{\mathcal O_X}(X;\mathcal F,\mathcal G)$ est un $A$-module
de type fini, et son séparé complété pour la topologie $\mathfrak J$-préadique
s'identifie canoniquement (avec les notations de
(\hyperref[III.4.1.7-fr]{4.1.7})) à
$\operatorname{Ext}^n_{\mathcal O_{\widehat X}}(\widehat X;\widehat{\mathcal F},
\widehat{\mathcal G})$.
\end{proposition}

\begin{proof}
On sait (T, 4.2) qu'il existe une suite spectrale birégulière
$\mathrm E(\mathcal F,\mathcal G)$ dont l'aboutissement est
$\operatorname{Ext}_{\mathcal O_X}(X;\mathcal F,\mathcal G)$ et dont les termes
$\mathrm E_2$ sont donnés par
\[
\mathrm E_2^{pq}=\mathrm H^p\bigl(X,
\mathcal{E}\!xt^q_{\mathcal O_X}(\mathcal F,\mathcal G)\bigr).
\]
On sait que $\mathcal{E}\!xt^q_{\mathcal O_X}(\mathcal F,\mathcal G)$ est un
$\mathcal O_X$-Module cohérent (0, 12.3.3) dont le support est contenu dans
l'intersection de ceux de $\mathcal F$ et $\mathcal G$ (T, 4.2.2), et est
par suite propre sur $Y$ (II, 5.4.10). On conclut de
(\hyperref[III.3.2.4-fr]{3.2.4}) que les $\mathrm E_2^{pq}$ sont des
$A$-modules de type fini, et par suite (0, 11.1.8) il en est de même de tous
les termes $\mathrm E_r^{pq}$ de la suite spectrale et de son aboutissement.
D'autre part, si $i:\widehat X\to X$ est le morphisme canonique,
$\widehat{\mathcal F}$ et $\widehat{\mathcal G}$ s'identifient canoniquement
à $i^*(\mathcal F)$ et $i^*(\mathcal G)$, et $i$ est plat (I, 10.8.8 et
10.8.9). On sait alors (0, 12.3.4) qu'il existe, pour tout $q\geq0$, un
$i$-morphisme canonique
\[
u_q:\mathcal{E}\!xt^q_{\mathcal O_X}(\mathcal F,\mathcal G)
\longrightarrow
\mathcal{E}\!xt^q_{\mathcal O_{\widehat X}}
(\widehat{\mathcal F},\widehat{\mathcal G}),
\]
et que le $\mathcal O_{\widehat X}$-homomorphisme correspondant
\[
u_q^\#:i^*\bigl(\mathcal{E}\!xt^q_{\mathcal O_X}(\mathcal F,\mathcal G)\bigr)
\longrightarrow
\mathcal{E}\!xt^q_{\mathcal O_{\widehat X}}
(\widehat{\mathcal F},\widehat{\mathcal G})
\]
est un isomorphisme (0, 12.3.5)~; autrement dit (I, 10.8.8),
$\mathcal{E}\!xt^q_{\mathcal O_{\widehat X}}
(\widehat{\mathcal F},\widehat{\mathcal G})$ s'identifie canoniquement au
complété $(\mathcal{E}\!xt^q_{\mathcal O_X}(\mathcal F,\mathcal G))^\wedge$ (avec les
notations de (\hyperref[III.4.1.7-fr]{4.1.7})). On conclut alors du th. de
comparaison (\hyperref[III.4.1.10-fr]{4.1.10}) que pour tout $p\geq0$,
\[
\mathrm H^p\bigl(\widehat X,
\mathcal{E}\!xt^q_{\mathcal O_{\widehat X}}
(\widehat{\mathcal F},\widehat{\mathcal G})\bigr)
\]
s'identifie canoniquement au séparé complété de
$\mathrm H^p\bigl(X,\mathcal{E}\!xt^q_{\mathcal O_X}
(\mathcal F,\mathcal G)\bigr)$ pour la topologie $\mathfrak J$-préadique.

\oldpage[III]{139}

Si on désigne par $\mathrm E(\widehat{\mathcal F},
\widehat{\mathcal G})$ la suite spectrale birégulière définie dans (T, 4.2)
relative à $\widehat{\mathcal F}$ et $\widehat{\mathcal G}$, on voit donc que
si $\widehat A$ désigne le séparé complété de $A$ pour la topologie
$\mathfrak J$-préadique, on a, à un isomorphisme canonique près,
\[
\mathrm E_2^{pq}(\widehat{\mathcal F},\widehat{\mathcal G})
=\mathrm E_2^{pq}(\mathcal F,\mathcal G)\otimes_A\widehat A
\qquad (0\textsubscript{I}, 7.3.3).
\]

Cela étant, on sait que la donnée du morphisme plat $i$ définit un
homomorphisme canonique de suites spectrales
\[
\varphi:\mathrm E(\mathcal F,\mathcal G)\longrightarrow
\mathrm E(\widehat{\mathcal F},\widehat{\mathcal G})
=\mathrm E(i^*(\mathcal F),i^*(\mathcal G))
\]
qui, pour les termes $\mathrm E_2$ (resp. l'aboutissement) se réduit à
l'homomorphisme
\[
\varphi_2^{pq}:\mathrm H^p\bigl(X,
\mathcal{E}\!xt^q_{\mathcal O_X}(\mathcal F,\mathcal G)\bigr)
\longrightarrow
\mathrm H^p\bigl(\widehat X,
\mathcal{E}\!xt^q_{\mathcal O_{\widehat X}}
(\widehat{\mathcal F},\widehat{\mathcal G})\bigr)
\]
\[
\text{(resp. }\varphi^n:\operatorname{Ext}^n_{\mathcal O_X}
(X;\mathcal F,\mathcal G)\longrightarrow
\operatorname{Ext}^n_{\mathcal O_{\widehat X}}
(\widehat X;\widehat{\mathcal F},\widehat{\mathcal G})\text{)}
\]
déduit de $u_q$ (resp. $u_0$) par fonctorialité (0, 12.3.4). Par
tensorisation avec $\widehat A$, les $\varphi_r^{pq}$ et $\varphi^n$ donnent
des homomorphismes de $\widehat A$-modules
\[
\psi_r^{pq}:\mathrm E_r^{pq}(\mathcal F,\mathcal G)
\otimes_A\widehat A\longrightarrow
\mathrm E_r^{pq}(\widehat{\mathcal F},\widehat{\mathcal G})
\]
\[
\psi^n:\operatorname{Ext}^n_{\mathcal O_X}(X;\mathcal F,\mathcal G)
\otimes_A\widehat A\longrightarrow
\operatorname{Ext}^n_{\mathcal O_{\widehat X}}
(\widehat X;\widehat{\mathcal F},\widehat{\mathcal G}).
\]
Comme $\widehat A$ est un $A$-module plat (0\textsubscript{I}, 7.3.3), les
$\widehat A$-modules
$\mathrm E_r^{pq}(\mathcal F,\mathcal G)\otimes_A\widehat A$ forment une
suite spectrale birégulière d'aboutissement les
$\operatorname{Ext}^n_{\mathcal O_X}(X;\mathcal F,\mathcal G)\otimes_A\widehat A$, et les
$\psi_r^{pq}$ et $\psi^n$ un morphisme de suites spectrales. Comme les
$\psi_2^{pq}$ sont des isomorphismes, il en est de même des $\psi^n$
(0, 11.1.5).
\end{proof}

\begin{corollary}[4.5.2]
\label{III.4.5.2-fr}
Sous les hypothèses de (\hyperref[III.4.5.1-fr]{4.5.1}), supposons en outre
que $A$ soit un anneau $\mathfrak J$-adique noethérien. Alors, pour tout
entier $n\geq0$, $\operatorname{Ext}^n_{\mathcal O_X}(X;\mathcal F,\mathcal G)$
s'identifie canoniquement à
$\operatorname{Ext}^n_{\mathcal O_{\widehat X}}
(\widehat X;\widehat{\mathcal F},\widehat{\mathcal G})$.
\end{corollary}

\begin{proof}
Il suffit de remarquer que $\operatorname{Ext}^n_{\mathcal O_X}
(X;\mathcal F,\mathcal G)$, étant un $A$-module de type fini, est séparé et
complet pour la topologie $\mathfrak J$-préadique (0\textsubscript{I}, 7.3.6).
\end{proof}

Le cas particulier $n=0$ de (\hyperref[III.4.5.1-fr]{4.5.1}) s'énonce de la
façon suivante~:

\begin{corollary}[4.5.3]
\label{III.4.5.3-fr}
Sous les hypothèses de (\hyperref[III.4.5.1-fr]{4.5.1}), pour tout
homomorphisme $u:\mathcal F\to\mathcal G$, désignons par
$\widehat u:\widehat{\mathcal F}\to\widehat{\mathcal G}$ l'homomorphisme
complété (I, 10.8.4). Alors on a un isomorphisme canonique
\[
\label{III.4.5.3.1-fr}
\left(\operatorname{Hom}_{\mathcal O_X}(\mathcal F,\mathcal G)\right)^\wedge
\xrightarrow{\sim}
\operatorname{Hom}_{\mathcal O_{\widehat X}}
(\widehat{\mathcal F},\widehat{\mathcal G}),
\tag{4.5.3.1}
\]
où le premier membre est le séparé complété pour la topologie
$\mathfrak J$-préadique du $A$-module
$\operatorname{Hom}_{\mathcal O_X}(\mathcal F,\mathcal G)$, cet isomorphisme étant obtenu
par passage aux séparés complétés à partir de l'homomorphisme
$u\mapsto\widehat u$.
\end{corollary}

\subsection{Relations entre morphismes formels et morphismes usuels}

\begin{proposition}[4.6.1]
\label{III.4.6.1-fr}
Soient $Y$ un préschéma localement noethérien, $f:X\to Y$ un morphisme
propre, $\mathcal F$ un $\mathcal O_X$-Module cohérent et $f$-plat, $y$ un
point de $Y$. Supposons que pour un entier $n$, on ait
\[
H^n\bigl(f^{-1}(y),\mathcal F\otimes_{\mathcal O_Y}k(y)\bigr)=0.
\]
Alors il existe un voisinage $U$ de $y$ dans $Y$ tel que
$R^nf_*(\mathcal F)|U=0$, et pour tout entier $p\geq0$, l'homomorphisme
canonique
\[
\bigl(R^{n-1}f_*(\mathcal F)\bigr)_y
\longrightarrow
H^{n-1}\bigl(f^{-1}(y),
\mathcal F\otimes_{\mathcal O_Y}
(\mathcal O_y/\mathfrak m_y^{p+1})\bigr)
\]
de (\hyperref[III.4.2.1.1-fr]{4.2.1.1}) est surjectif.
\end{proposition}

\begin{proof}
Comme $R^nf_*(\mathcal F)$ est un $\mathcal O_Y$-Module cohérent
(\hyperref[III.3.2.1-fr]{3.2.1}), la première assertion de la proposition
sera établie si l'on prouve que l'on a
$(R^nf_*(\mathcal F))_y=0$ (0\textsubscript{I}, 5.2.2)~; en vertu de
(\hyperref[III.4.2.1-fr]{4.2.1}), il suffira de prouver que
\[
H^n\bigl(f^{-1}(y),\mathcal F\otimes_{\mathcal O_Y}
(\mathcal O_y/\mathfrak m_y^{p+1})\bigr)=0
\]
pour tout $p$. C'est vrai par hypothèse pour $p=0$~; nous allons le démontrer
par récurrence sur $p$. Posons
\[
X_p=X\times_Y\operatorname{Spec}
(\mathcal O_y/\mathfrak m_y^{p+1}),
\]
de sorte que $X_{p-1}$ est un sous-préschéma fermé de $X_p$, ayant même
espace sous-jacent (I, 3.6.1)~; l'hypothèse de récurrence
\[
H^n\bigl(X_{p-1},\mathcal F\otimes_{\mathcal O_Y}
(\mathcal O_y/\mathfrak m_y^p)\bigr)=0
\]
entraîne donc
\[
H^n\bigl(X_p,\mathcal F\otimes_{\mathcal O_Y}
(\mathcal O_y/\mathfrak m_y^p)\bigr)=0~;
\]
d'autre part, la suite exacte de cohomologie donne, à partir de la suite
exacte
\[
0\longrightarrow
\mathfrak m_y^p\mathcal F/\mathfrak m_y^{p+1}\mathcal F
\longrightarrow
\mathcal F/\mathfrak m_y^{p+1}\mathcal F
\longrightarrow
\mathcal F/\mathfrak m_y^p\mathcal F
\longrightarrow0
\]
de $\mathcal O_{X_p}$-Modules, la suite exacte
\[
H^n\bigl(X_p,
\mathfrak m_y^p\mathcal F/\mathfrak m_y^{p+1}\mathcal F\bigr)
\longrightarrow
H^n\bigl(X_p,\mathcal F/\mathfrak m_y^{p+1}\mathcal F\bigr)
\longrightarrow
H^n\bigl(X_p,\mathcal F/\mathfrak m_y^p\mathcal F\bigr)
\]
et il suffira de montrer que l'on a
\[
\label{III.4.6.1.1-fr}
H^n\bigl(X_p,
\mathfrak m_y^p\mathcal F/\mathfrak m_y^{p+1}\mathcal F\bigr)=0.
\tag{4.6.1.1}
\]
car alors $H^n(X_p,\mathcal F/\mathfrak m_y^{p+1}\mathcal F)$ sera un
sous-module de $H^n(X_p,\mathcal F/\mathfrak m_y^p\mathcal F)$, donc $0$ en
vertu de l'hypothèse de récurrence.

\oldpage[III]{140}

Notons maintenant que la fibre
$Z=f^{-1}(y)=X\times_Y\operatorname{Spec}(k(y))$ est un sous-préschéma fermé
de $X_p$, et que
$\mathfrak m_y^p\mathcal F/\mathfrak m_y^{p+1}\mathcal F$ est annulé par
$\mathfrak m_y$, donc peut être considéré comme un $\mathcal O_Z$-Module, de
sorte que
\[
H^n\bigl(Z,\mathfrak m_y^p\mathcal F/
\mathfrak m_y^{p+1}\mathcal F\bigr)
=H^n\bigl(X_p,\mathfrak m_y^p\mathcal F/
\mathfrak m_y^{p+1}\mathcal F\bigr).
\]
Cela étant, nous allons montrer que le $\mathcal O_Z$-homomorphisme canonique
\[
\label{III.4.6.1.2-fr}
(\mathcal F/\mathfrak m_y\mathcal F)
\otimes_{k(y)}(\mathfrak m_y^p/\mathfrak m_y^{p+1})
\longrightarrow
\mathfrak m_y^p\mathcal F/\mathfrak m_y^{p+1}\mathcal F
\tag{4.6.1.2}
\]
est bijectif~; cela établi, il en résultera, puisque
$\mathfrak m_y^p/\mathfrak m_y^{p+1}$ est un $k(y)$-module libre, que l'on a
\[
\begin{split}
H^n\bigl(Z,\mathfrak m_y^p\mathcal F/
\mathfrak m_y^{p+1}\mathcal F\bigr)
&=H^n\bigl(Z,\mathcal F/\mathfrak m_y\mathcal F\bigr)
\otimes_{k(y)}(\mathfrak m_y^p/\mathfrak m_y^{p+1})\\
&=0
\end{split}
\]
(0, 12.2.3), puisque
$H^n(Z,\mathcal F/\mathfrak m_y\mathcal F)=0$ par hypothèse, d'où
(\hyperref[III.4.6.1.1-fr]{4.6.1.1}). Pour établir la première assertion, il
reste donc à prouver que (\hyperref[III.4.6.1.2-fr]{4.6.1.2}) est bijectif~;
comme la question est ponctuelle sur $X$ et que $\mathcal F_x$ est un
$\mathcal O_y$-module plat par hypothèse pour tout $x\in f^{-1}(y)$, il
suffit d'appliquer (0, 10.2.1, c)), car
$\mathcal F_x/\mathfrak m_y\mathcal F_x$ est un module plat sur le corps
$k(y)=\mathcal O_y/\mathfrak m_y$.

Pour démontrer la seconde assertion de (\hyperref[III.4.6.1-fr]{4.6.1}), on
se ramène aussitôt, comme dans (\hyperref[III.4.2.1-fr]{4.2.1}), au cas où
$Y$ est affine et $y$ fermé. Notons que
(\hyperref[III.4.6.1.1-fr]{4.6.1.1}) donne, par un raisonnement analogue,
pour tout $k>0$, la relation
\[
\label{III.4.6.1.3-fr}
H^n\bigl(X_{k+1},
\mathfrak m_y^k\mathcal F/
\mathfrak m_y^{k+p+1}\mathcal F\bigr)=0.
\tag{4.6.1.3}
\]

\oldpage[III]{141}

d'où on déduit, par (\hyperref[III.4.2.1-fr]{4.2.1}), que l'on a aussi
\[
\label{III.4.6.1.4-fr}
\bigl(R^nf_*(\mathfrak m_y^k\mathcal F)\bigr)_y=0.
\tag{4.6.1.4}
\]
Cela étant, on tire de la suite exacte de cohomologie l'exactitude de la
suite
\[
\bigl(R^{n-1}f_*(\mathcal F)\bigr)_y
\longrightarrow
\bigl(R^{n-1}f_*(\mathcal F/\mathfrak m_y^p\mathcal F)\bigr)_y
\longrightarrow
\bigl(R^nf_*(\mathfrak m_y^p\mathcal F)\bigr)_y=0
\]
et comme $y$ est fermé et $Y$ affine, on a (\hyperref[III.1.4.11-fr]{1.4.11})
\[
R^{n-1}f_*(\mathcal F/\mathfrak m_y^p\mathcal F)
=\bigl(H^{n-1}(X,\mathcal F/\mathfrak m_y^p\mathcal F)\bigr)^\sim
=\bigl(H^{n-1}(f^{-1}(y),
\mathcal F/\mathfrak m_y^p\mathcal F)\bigr)^\sim
\]
(G, II, 4.9.1)~; or
$H^{n-1}(f^{-1}(y),\mathcal F/\mathfrak m_y^p\mathcal F)$ est un
$(\mathcal O_y/\mathfrak m_y^p)$-module, d'où
\[
\bigl(R^{n-1}f_*(\mathcal F/\mathfrak m_y^p\mathcal F)\bigr)_y
=H^{n-1}(f^{-1}(y),\mathcal F/\mathfrak m_y^p\mathcal F)
\]
et cela achève de prouver (\hyperref[III.4.6.1-fr]{4.6.1}).
\end{proof}

\begin{corollary}[4.6.2]
\label{III.4.6.2-fr}
Soient $Y$ un préschéma localement noethérien, $f:X\to Y$ un morphisme propre
et plat, $\mathcal F$, $\mathcal G$ deux $\mathcal O_X$-Modules localement
libres, $y$ un point de $Y$. Posons
\[
X_y=f^{-1}(y)=X\otimes_Yk(y),\qquad
\mathcal F_y=\mathcal F\otimes_{\mathcal O_Y}k(y),\qquad
\mathcal G_y=\mathcal G\otimes_{\mathcal O_Y}k(y),
\]
et supposons que l'on ait
\[
\label{III.4.6.2.1-fr}
H^1\bigl(X_y,
\mathcal{H}\!om_{\mathcal O_{X_y}}(\mathcal F_y,\mathcal G_y)\bigr)=0.
\tag{4.6.2.1}
\]
Alors, pour tout homomorphisme $u_0:\mathcal F_y\to\mathcal G_y$, il existe
un voisinage ouvert $U$ de $y$ et un homomorphisme
\[
u:\mathcal F|f^{-1}(U)\longrightarrow\mathcal G|f^{-1}(U)
\]
tels que $u_0$ soit égal à l'homomorphisme $u\otimes1$.
\end{corollary}

\begin{proof}
En effet, l'hypothèse permet d'appliquer (\hyperref[III.4.6.1-fr]{4.6.1}) au
$\mathcal O_X$-Module cohérent
$\mathcal H=\mathcal{H}\!om_{\mathcal O_X}(\mathcal F,\mathcal G)$ pour
$n=1$ et $p=0$, car $\mathcal H$ est localement libre et \emph{a fortiori}
$f$-plat, et le $\mathcal O_{X_y}$-Module
$\mathcal H\otimes_{\mathcal O_Y}k(y)$ s'identifie alors à
$\mathcal{H}\!om_{\mathcal O_{X_y}}(\mathcal F_y,\mathcal G_y)$
(0\textsubscript{I}, 6.2.2). On peut supposer $Y=\operatorname{Spec}(A)$
affine, et alors (\hyperref[III.1.4.11-fr]{1.4.11})
\[
R^0f_*(\mathcal H)=
\bigl(\operatorname{Hom}_{\mathcal O_X}(\mathcal F,\mathcal G)\bigr)^\sim,
\]
donc
\[
\bigl(R^0f_*(\mathcal H)\bigr)_y
=\operatorname{Hom}_{\mathcal O_X}(\mathcal F,\mathcal G)
\otimes_A\mathcal O_y~;
\]
l'homomorphisme canonique
\[
\operatorname{Hom}_{\mathcal O_X}(\mathcal F,\mathcal G)
\otimes_A\mathcal O_y
\longrightarrow
\operatorname{Hom}_{\mathcal O_{X_y}}(\mathcal F_y,\mathcal G_y)
\]
étant surjectif par (\hyperref[III.4.6.1-fr]{4.6.1}), cela établit le
corollaire, tout élément de
$\operatorname{Hom}_{\mathcal O_X}(\mathcal F,\mathcal G)
\otimes_A\mathcal O_y$ pouvant toujours se mettre sous la forme
$u\otimes(1/s)$, où $s\notin\mathfrak m_y$ est un élément de $A$.
\end{proof}

Ce corollaire peut se compléter par le suivant~:

\begin{corollary}[4.6.3]
\label{III.4.6.3-fr}
Sous les hypothèses de (\hyperref[III.4.6.2-fr]{4.6.2}), si $u_0$ est
injectif (resp. surjectif, bijectif), on peut supposer qu'il en est de même de
$u$.
\end{corollary}

\begin{proof}
On peut se borner au cas où $U=Y$. Il suffit de prouver que si $u_0$ est
injectif (resp. surjectif), $\operatorname{Ker}u_x=0$ (resp.
$\operatorname{Coker}u_x=0$) pour tout $x\in f^{-1}(y)$~: en effet,
$\operatorname{Ker}u$ et $\operatorname{Coker}u$ sont des
$\mathcal O_X$-Modules cohérents (0\textsubscript{I}, 5.3.4), donc il
existera un voisinage $V$ de $f^{-1}(y)$ dans $X$ tel que la restriction de
$\operatorname{Ker}u$ (resp. $\operatorname{Coker}u$) à $V$ soit $0$
(0\textsubscript{I}, 5.2.2)~; comme $f$ est fermé, il existera un voisinage
$U'\subset U$ de $y$ tel que $f^{-1}(U')\subset V$, et
(\hyperref[III.4.6.3-fr]{4.6.3}) sera démontré. Par hypothèse,
\[
u_x\otimes1:
\mathcal F_x\otimes_{\mathcal O_y}k(y)
\longrightarrow
\mathcal G_x\otimes_{\mathcal O_y}k(y)
\]
est injectif (resp. surjectif), $\mathcal F_x$ et $\mathcal G_x$ sont des
$\mathcal O_x$-modules libres de type fini et $\mathcal O_x$ est un
$\mathcal O_y$-module plat. Lorsque l'on suppose $u_x\otimes1$ injectif, le
fait que $u_x$ soit injectif résulte de (0, 10.2.4). Lorsque l'on suppose
$u_x\otimes1$ surjectif, \emph{a fortiori} l'homomorphisme
\[
\mathcal F_x/\mathfrak m_x\mathcal F_x
\longrightarrow
\mathcal G_x/\mathfrak m_x\mathcal G_x,
\]
qui s'en déduit par passage aux quotients, est surjectif~; comme
$\mathcal G_x$ est un $\mathcal O_x$-module de type fini et que
$\mathcal O_x$ est un anneau local d'idéal maximal $\mathfrak m_x$, la
conclusion résulte du lemme de Nakayama (Bourbaki, \emph{Alg.}, chap. VIII,
\S\,6, no~3, cor. 4 de la prop. 6).
\end{proof}

On déduit en particulier de (\hyperref[III.4.6.3-fr]{4.6.3})~:

\begin{corollary}[4.6.4]
\label{III.4.6.4-fr}
Soient $Y$ un préschéma localement noethérien, $f:X\to Y$ un morphisme propre
et plat, $y$ un point de $Y$, $X_y=X\otimes_Yk(y)$. Soit $\mathcal E_0$ un
$\mathcal O_{X_y}$-Module localement libre tel que
\[
\label{III.4.6.4.1-fr}
H^1\bigl(X_y,
\mathcal{H}\!om_{\mathcal O_{X_y}}(\mathcal E_0,\mathcal E_0)\bigr)=0.
\tag{4.6.4.1}
\]
Soient $\mathcal F$, $\mathcal G$ deux $\mathcal O_X$-Modules localement
libres tels que $\mathcal F_y$ et $\mathcal G_y$ (avec les notations de
(\hyperref[III.4.6.2-fr]{4.6.2})) soient isomorphes à $\mathcal E_0$. Alors
il existe un voisinage ouvert $U$ de $y$ tel que $\mathcal F|f^{-1}(U)$ et
$\mathcal G|f^{-1}(U)$ soient isomorphes.
\end{corollary}

Plus particulièrement~:

\begin{corollary}[4.6.5]
\label{III.4.6.5-fr}
Sous les hypothèses de (\hyperref[III.4.6.4-fr]{4.6.4}) sur $f$, $X$, $Y$,
supposons que $H^1(X_y,\mathcal O_{X_y})=0$. Si $\mathcal F$ et $\mathcal G$
sont deux $\mathcal O_X$-Modules inversibles tels que $\mathcal F_y$ et
$\mathcal G_y$ soient isomorphes, il existe un voisinage ouvert $U$ de $y$
tel que $\mathcal F|f^{-1}(U)$ et $\mathcal G|f^{-1}(U)$ soient isomorphes.
\end{corollary}

\begin{proof}
Il suffit d'appliquer (\hyperref[III.4.6.4-fr]{4.6.4}) aux modules
$\mathcal F^{-1}\otimes\mathcal G$ et $\mathcal O_X$.
\end{proof}

\begin{remarks}[4.6.6]
\label{III.4.6.6-fr}
\begin{enumerate}
\item[(i)] En utilisant (\hyperref[III.4.6.5-fr]{4.6.5}), nous établirons au
chap. V la classification des faisceaux inversibles sur un fibré projectif,
annoncée dans (II, 4.2.7).
\item[(ii)] Le résultat de (\hyperref[III.4.6.1-fr]{4.6.1}) apparaîtra au
\S\,7 comme conséquence de propositions plus générales.
\end{enumerate}
\end{remarks}

\begin{proposition}[4.6.7]
\label{III.4.6.7-fr}
Soient $Z$ un préschéma localement noethérien, $X$, $Y$ deux $Z$-préschémas
tels que les morphismes structuraux $g:X\to Z$, $h:Y\to Z$ soient propres.
Soient $f:X\to Y$ un $Z$-morphisme, $z$ un point de $Z$, et soit
\[
f_z=f\times_Z1:
X\otimes_Zk(z)\longrightarrow Y\otimes_Zk(z).
\]
\begin{enumerate}
\item[(i)] Si $f_z$ est un morphisme fini (resp. une immersion fermée), il
existe un voisinage ouvert $U$ de $z$ tel que le morphisme
$g^{-1}(U)\to h^{-1}(U)$, restriction de $f$, soit un morphisme fini (resp.
une immersion fermée).
\item[(ii)] On suppose de plus que $g$ soit un morphisme plat. Alors, si
$f_z$ est un isomorphisme, il existe un voisinage ouvert $U$ de $z$ tel que
le morphisme $g^{-1}(U)\to h^{-1}(U)$, restriction de $f$, soit un
isomorphisme.
\end{enumerate}
\end{proposition}

\begin{proof}
Dans les deux cas, il suffira de prouver que pour tout $y\in h^{-1}(z)$ il
existe un voisinage $V_y$ de $y$ tel que la restriction
$f^{-1}(V_y)\to V_y$ de $f$ soit un morphisme fini (resp. une immersion
fermée, un isomorphisme)~; il en résultera alors en effet que si $V$ est la
réunion des $V_y$, la restriction $f^{-1}(V)\to V$ de $f$ est un morphisme
fini (resp. une immersion fermée, un isomorphisme) (II, 6.1.1 et I, 4.2.4).
Comme $h$ est un morphisme fermé, il existera un voisinage ouvert $U$ de $z$
tel que $h^{-1}(U)\subset V$, et la proposition sera démontrée.

\noindent\textit{(i)} Notons tout d'abord que $f$ est un morphisme propre
(II, 5.4.3)~; si on suppose $f_z$ fini, l'existence pour tout
$y\in h^{-1}(z)$ d'un voisinage $V_y$ tel que $f^{-1}(V_y)\to V_y$ soit fini
résulte de (\hyperref[III.4.4.11-fr]{4.4.11}).

\oldpage[III]{143}

Pour traiter le cas où $f_z$ est une immersion fermée, on peut donc déjà
supposer que le morphisme $f$ est fini, donc que
$X=\operatorname{Spec}(\mathcal B)$, où $\mathcal B$ est une
$\mathcal O_Y$-Algèbre cohérente, le morphisme $f$ correspondant (II, 1.2.7)
à l'homomorphisme canonique $u:\mathcal O_Y\to\mathcal B$. Si l'on prouve
que pour tout $y\in h^{-1}(z)$, l'homomorphisme
$u_y:\mathcal O_y\to\mathcal B_y$ est surjectif, il en résultera que pour un
voisinage $V_y$ de $y$, $u|V_y$ sera surjectif, le faisceau
$\operatorname{Coker}u$ étant cohérent (0\textsubscript{I}, 5.3.4 et 5.2.2).
Cela étant, le morphisme fini $f_z$ correspond à l'homomorphisme
\[
v=u\otimes1:
\mathcal O_Y\otimes_{\mathcal O_Z}k(z)
\longrightarrow
\mathcal B\otimes_{\mathcal O_Z}k(z),
\]
et l'hypothèse que $f_z$ est une immersion fermée entraîne que
l'homomorphisme
\[
u_y\otimes1:
\mathcal O_y\otimes_{\mathcal O_z}k(z)
\longrightarrow
\mathcal B_y\otimes_{\mathcal O_z}k(z)
\]
est surjectif. Comme $\mathcal B_y$ est un $\mathcal O_y$-module de type
fini et $\mathcal O_y$ un anneau local noethérien, la conclusion résulte
comme dans (\hyperref[III.4.6.3-fr]{4.6.3}) du lemme de Nakayama.

\noindent\textit{(ii)} Le même raisonnement que ci-dessus montre qu'il
suffit cette fois de prouver que $u_y$ est bijectif, sachant que
$u_y\otimes1$ est bijectif.

Cela résultera du lemme suivant~:
\end{proof}

\begin{lemma}[4.6.7.1]
\label{III.4.6.7.1-fr}
Soient $A$, $B$ deux anneaux locaux noethériens,
$\rho:A\to B$ un homomorphisme local, $u:N\to M$ un homomorphisme de
$B$-modules. On suppose que $M$ est un $A$-module plat, $N$ un $B$-module
de type fini et que
\[
u\otimes1:N\otimes_Ak\longrightarrow M\otimes_Ak
\]
(où $k$ est le corps résiduel de $A$) est injectif. Alors $N$ est un
$A$-module plat et $u$ est injectif.
\end{lemma}

\begin{proof}
Pour établir la première assertion, il faut montrer que pour tout couple de
$A$-modules de type fini $P$, $Q$ et tout $A$-homomorphisme injectif
$v:P\to Q$,
$1_N\otimes v:N\otimes_AP\to N\otimes_AQ$ est injectif. Or, on a le
diagramme commutatif
\[
\xymatrix{
N\otimes_AP\ar[r]^{1_N\otimes v}\ar[d]_{u\otimes1_P}
  &N\otimes_AQ\ar[d]^{u\otimes1_Q}\\
M\otimes_AP\ar[r]_{1_M\otimes v}
  &M\otimes_AQ
}
\]
et comme $1_M\otimes v$ est injectif par hypothèse, il suffira de prouver
qu'il en est de même de $u\otimes1_P$. Soit $\mathfrak m$ l'idéal maximal
de $A$~; la filtration $\mathfrak m$-adique sur le $A$-module
$N\otimes_AP$ est aussi sa filtration $\mathfrak mB$-adique en tant que
$B$-module~; la topologie définie par cette filtration est donc séparée,
puisque $B$ est noethérien, que $\mathfrak mB$ est contenu dans le radical de
$B$, et que $N\otimes_AP$ est un $B$-module de type fini, $N$ étant un
$B$-module de type fini et $P$ un $A$-module de type fini
(0\textsubscript{I}, 7.3.5). Il suffit donc de prouver que l'homomorphisme
\[
\operatorname{gr}(u\otimes1_P):
\operatorname{gr}_\bullet(N\otimes_AP)
\longrightarrow
\operatorname{gr}_\bullet(M\otimes_AP)
\]
(où les modules gradués sont relatifs aux filtrations $\mathfrak m$-adiques)
est injectif (Bourbaki, \emph{Alg. comm.}, chap. III, \S\,2, no~8, cor. 1 du
th. 1). Notons maintenant que puisque $M$ est un $A$-module plat, les
homomorphismes
\[
M\otimes_A(\mathfrak m^nP)
\longrightarrow
\mathfrak m^n(M\otimes_AP)
\]
sont bijectifs~; il en est donc de même de l'homomorphisme canonique
\[
\varphi_M:
\operatorname{gr}_0(M)\otimes_A\operatorname{gr}_\bullet(P)
\longrightarrow
\operatorname{gr}_\bullet(M\otimes_AP).
\]

\oldpage[III]{144}

Or, on a un diagramme commutatif
\[
\xymatrix{
\operatorname{gr}_0(N)\otimes_A\operatorname{gr}_\bullet(P)
  \ar[r]^{\operatorname{gr}_0(u)\otimes1}\ar[d]_{\varphi_N}
  &\operatorname{gr}_0(M)\otimes_A\operatorname{gr}_\bullet(P)
  \ar[d]^{\varphi_M}\\
\operatorname{gr}_\bullet(N\otimes_AP)
  \ar[r]_{\operatorname{gr}(u\otimes1)}
  &\operatorname{gr}_\bullet(M\otimes_AP)
}
\]
dans lequel $\varphi_M$ est bijectif, $\varphi_N$ surjectif~; en outre,
$\operatorname{gr}_0(u)$ est injectif par hypothèse, et comme
\[
\begin{split}
\operatorname{gr}_0(N)\otimes_A\operatorname{gr}_\bullet(P)
&=\operatorname{gr}_0(N)\otimes_k\operatorname{gr}_\bullet(P),\\
\operatorname{gr}_0(M)\otimes_A\operatorname{gr}_\bullet(P)
&=\operatorname{gr}_0(M)\otimes_k\operatorname{gr}_\bullet(P),
\end{split}
\]
$\operatorname{gr}_0(u)\otimes1$ est aussi injectif. On en conclut que
$\operatorname{gr}(u\otimes1)$ est injectif, ce qui achève de démontrer la
première assertion. La seconde se déduit du raisonnement précédent en faisant
$P=A$.
\end{proof}

\begin{proposition}[4.6.8]
\label{III.4.6.8-fr}
Soient $Z$ un préschéma localement noethérien, $X$, $Y$, deux $Z$-préschémas
tels que les morphismes structuraux $g:X\to Z$, $h:Y\to Z$ soient propres,
$Z'$ une partie fermée de $Z$, $X'=g^{-1}(Z')$, $Y'=h^{-1}(Z')$ ses images
réciproques,
\[
\widehat X=X_{/X'},\qquad
\widehat Y=Y_{/Y'},\qquad
\widehat Z=Z_{/Z'}
\]
les complétés formels de $X$, $Y$, $Z$ le long de ces parties fermées,
$f:X\to Y$ un $Z$-morphisme, $\widehat f:\widehat X\to\widehat Y$ son
prolongement aux complétés. Pour que $\widehat f$ soit un isomorphisme (resp.
une immersion fermée), il faut et il suffit qu'il existe un voisinage ouvert
$U$ de $Z'$ tel que le morphisme $g^{-1}(U)\to h^{-1}(U)$, restriction de
$f$, soit un isomorphisme (resp. une immersion fermée).
\end{proposition}

\begin{proof}
La suffisance de la condition est immédiate (I, 10.14.7). Pour en montrer la
nécessité, il suffit encore de prouver que pour tout $y\in Y'$, il existe un
voisinage ouvert $V_y$ de $y$ tel que la restriction
$f^{-1}(V_y)\to V_y$ de $f$ soit un isomorphisme (resp. une immersion
fermée), par le même raisonnement que dans
(\hyperref[III.4.6.7-fr]{4.6.7}). On est ainsi ramené au cas où $Y=Z$,
$Y=\operatorname{Spec}(A)$ étant affine noethérien. Par hypothèse
(I, 10.9.1 et 10.14.2) la fibre $f^{-1}(y)$ est réduite à un point pour
$y\in Y'$, donc comme $f$ est propre (II, 5.4.3), il existe un voisinage
ouvert $U$ de $y$ tel que la restriction $f^{-1}(U)\to U$ de $f$ soit un
morphisme fini (\hyperref[III.4.4.11-fr]{4.4.11}). On peut donc déjà supposer
que $f$ soit un morphisme fini, donc $X=\operatorname{Spec}(B)$, où $B$ est
une $A$-algèbre finie sur $A$. Si $Y'=V(\mathfrak J)$, on a alors
\[
\widehat Y=\operatorname{Spf}(\widehat A),\qquad
\widehat X=\operatorname{Spf}(\widehat B),
\]
$\widehat A$ étant le séparé complété de $A$ pour la topologie
$\mathfrak J$-préadique, $\widehat B$ le séparé complété de $B$ pour la
topologie $\mathfrak JB$-préadique, ou (ce qui revient au même), le séparé
complété du $A$-module $B$ pour la topologie $\mathfrak J$-préadique~; en
outre, $\widehat f$ est le morphisme de schémas formels affines correspondant
au prolongement continu $\widehat\varphi:\widehat A\to\widehat B$ de
l'homomorphisme canonique d'anneaux $\varphi:A\to B$, et l'hypothèse est que
$\widehat\varphi$ est surjectif (resp. bijectif) (I, 10.14.2). Or,
$\widehat\varphi$ est aussi le prolongement continu de $\varphi$ considéré
comme homomorphisme de $A$-modules~; on sait alors (I, 10.8.14) qu'il existe
un voisinage ouvert $U$ de $Y'$ tel que la restriction à $U$ de
l'homomorphisme $\widetilde\varphi:\widetilde A\to\widetilde B$ de
$\mathcal O_Y$-Modules soit surjectif (resp. bijectif), ce qui achève la
démonstration.
\end{proof}

\oldpage[III]{145}

\subsection{Un critère d'amplitude}

\begin{theorem}[4.7.1]
\label{III.4.7.1-fr}
Soient $Y$ un préschéma localement noethérien, $f:X\to Y$ un morphisme
propre, $\mathcal L$ un $\mathcal O_X$-Module inversible, $y$ un point de
$Y$, $X_y=X\otimes_Yk(y)=f^{-1}(y)$, $g$ la projection de $X_y$ dans $X$.
Si $\mathcal L_y=g^*(\mathcal L)=\mathcal L\otimes_{\mathcal O_Y}k(y)$ est
ample sur $X_y$, il existe un voisinage ouvert $U$ de $y$ dans $Y$ tel que
$\mathcal L|f^{-1}(U)$ soit ample pour la restriction de $f$ à $f^{-1}(U)$.
\end{theorem}

\begin{proof}
\noindent\textit{I)} Posons
\[
Y'=\operatorname{Spec}(\mathcal O_y),\qquad
X'=X\times_YY',\qquad
\mathcal L'=\mathcal L\otimes_{\mathcal O_Y}\mathcal O_y~;
\]
nous allons d'abord prouver que $\mathcal L'$ est ample pour
$f'=f_{(Y')}$. On a le diagramme commutatif
\[
\xymatrix{
X\ar[d]_f & X'\ar[l]\ar[d]_{f'} & X_y\ar[l]\ar[d]_{f_y}\\
Y & Y'\ar[l] & \operatorname{Spec}(k(y))\ar[l]
}
\]
Comme $f'$ est propre (II, 5.4.2, (iii)) et $\mathcal O_y$ noethérien, on
voit qu'on peut se borner au cas où
$Y=Y'=\operatorname{Spec}(\mathcal O_y)$, donc $X=X'$, supposer que
$\mathcal L_y$ est ample pour $f_y$ et prouver que $\mathcal L$ est ample
pour $f$ (II, 4.6.6). Nous allons appliquer le critère
(\hyperref[III.2.6.1-fr]{2.6.1, c)}) et montrer en fait que pour tout
$\mathcal O_X$-Module cohérent $\mathcal F$, il existe un entier $N$ tel que
\[
H^1(X,\mathcal F(n))=0\quad\text{pour tout }n\geq N,
\qquad
\mathcal F(n)=\mathcal F\otimes\mathcal L^{\otimes n}.
\]
Notons que $y$ est un point fermé de $Y$ correspondant à l'idéal maximal
$\mathfrak m$ de $\mathcal O_y$~; $X_y$ est donc un sous-préschéma fermé de
$X$ défini par l'Idéal cohérent
$\mathcal J=f^*(\widetilde{\mathfrak m})\mathcal O_X=
\mathfrak m\mathcal O_X$ de $\mathcal O_X$ (I, 4.4.5), et $g$ l'injection
canonique. Considérons alors la $k(y)$-algèbre graduée
\[
S=\operatorname{gr}(\mathcal O_y)
=\bigoplus_{j\geq0}\mathfrak m^j/\mathfrak m^{j+1},
\]
qui est de type fini puisque $\mathcal O_y$ est noethérien~; la
$\mathcal O_X$-Algèbre $\mathcal S=f^*(\widetilde S)$ est donc
quasi-cohérente et de type fini, et elle est évidemment annulée par
$\mathcal J$, donc si on pose $\mathcal S_y=g^*(\mathcal S)$,
$\mathcal S_y$ est une $\mathcal O_{X_y}$-Algèbre quasi-cohérente de type
fini, et $\mathcal S=g_*(\mathcal S_y)$. Posons d'autre part
\[
\mathcal M_j=\mathfrak m^j\mathcal F/\mathfrak m^{j+1}\mathcal F,
\qquad
\mathcal M=\bigoplus_{j\geq0}\mathcal M_j
=\operatorname{gr}(\mathcal F)~;
\]
comme $\mathcal F$ est cohérent, $\mathcal M$ est un $\mathcal S$-Module
quasi-cohérent de type fini (0, 10.1.1) qui est aussi annulé par
$\mathcal J$, de sorte que si l'on pose
\[
\mathcal M'_j=g^*(\mathcal M_j),\qquad
\mathcal M'=g^*(\mathcal M)=\bigoplus_{j\geq0}\mathcal M'_j,
\]
$\mathcal M'$ est un $\mathcal S_y$-Module gradué quasi-cohérent de type
fini tel que $\mathcal M=g_*(\mathcal M')$. En outre, si on pose
$\mathcal M'_j(n)=\mathcal M'_j\otimes\mathcal L_y^{\otimes n}$, on a
$\mathcal M'_j(n)=g^*(\mathcal M_j(n))$. Cela étant, $f_y$ est propre
(II, 5.4.2, (iii)) et $\mathcal L_y$ est ample, donc $f_y$ est projectif
(II, 5.5.4 et 4.6.11), et on peut appliquer à
$\operatorname{Spec}(k(y))$, $f_y$, $\mathcal S_y$, $\mathcal L_y$ et
$\mathcal M'$ le théorème (\hyperref[III.2.4.1-fr]{2.4.1, (ii)})~: il
existe un entier $N$ tel que pour $n\geq N$, on ait
$H^q(X_y,\mathcal M'_j(n))=0$ pour tout $q>0$ et tout $j$~; par suite, on a
aussi $H^q(X,\mathcal M_j(n))=0$ pour tout $q>0$ et tout $j$
(G, II, 4.9.1). Posons alors
\[
\mathcal F(n)_j=\mathcal F(n)/\mathfrak m^{j+1}\mathcal F(n),
\]
de sorte que $\mathcal F(n)_{j-1}=\mathcal F(n)_j/\mathcal M_j(n)$ pour
$j\geq1$ et $\mathcal F(n)_0=\mathcal M_0(n)$. On a
$H^1(X,\mathcal F(n)_0)=0$, et, par la suite exacte de cohomologie,
$H^1(X,\mathcal F(n)_j)=H^1(X,\mathcal F(n)_{j-1})$ pour tout $j\geq1$,
donc $H^1(X,\mathcal F(n)_j)=0$ pour tout $j\geq0$. On conclut donc de
(\hyperref[III.4.2.1-fr]{4.2.1}) que $H^1(X,\mathcal F(n))=0$, ce qui
achève de prouver notre assertion.

\oldpage[III]{146}

\noindent\textit{II)} Revenons aux notations du début de la démonstration,
et remarquons qu'on peut toujours supposer $Y=\operatorname{Spec}(A)$
affine~; comme $f'$ est de type fini et $\mathcal L'$ ample pour $f'$, il
existe un entier $m>0$ tel que $\mathcal L'^{\otimes m}$ soit très ample pour
$f'$ (II, 4.6.11)~; remplaçant au besoin $\mathcal L$ par
$\mathcal L^{\otimes m}$, on peut se borner à considérer le cas où
$\mathcal L'$ est très ample pour $f'$, et à prouver que
$\mathcal L|f^{-1}(U)$ est alors très ample pour $f$. Comme $f'$ est propre,
il existe alors une $Y'$-immersion fermée
$j:X'\to P=\mathbf P_{Y'}^r$ pour un entier $r>0$ convenable, telle que
$\mathcal L'$ soit isomorphe à $j^*(\mathcal O_P(1))$ (II, 5.5.4, (ii))~;
cette immersion correspond canoniquement à un $\mathcal O_{X'}$-homomorphisme
surjectif
\[
u:\mathcal O_{X'}^{r+1}\longrightarrow\mathcal L'
\]
(II, 4.2.3). Ce dernier correspond (0\textsubscript{I}, 5.1.1) à la donnée
de $r+1$ sections $s'_i$ ($0\leq i\leq r$) de $\mathcal L'$ au-dessus de
$X'$ qui engendrent ce $\mathcal O_{X'}$-Module. Ces sections sont aussi par
définition des sections de $f'_*(\mathcal L')$ au-dessus de $Y'$~; on a
\[
f'_*(\mathcal L')=f_*(\mathcal L)\otimes_{\mathcal O_Y}\mathcal O_{Y'}
\]
(0\textsubscript{I}, 5.4.10), $Y$ est affine et $\mathcal O_y$ est l'anneau
local en l'idéal premier $\mathfrak j_y$ de $A$, donc on a
$s'_i=s''_i/t_i$, où les $s''_i$ sont des sections de $f_*(\mathcal L)$
au-dessus de $Y$ et les $t_i$ des éléments de $A$ n'appartenant pas à
$\mathfrak j_y$~; on en conclut qu'il existe un voisinage ouvert affine $V$
de $y$ dans $Y$ et des sections $s_i$ de $f_*(\mathcal L)|V$ telles que
$s'_i=s_i/1$ (on rappelle que l'espace $Y'$ est contenu dans $V$, cf. I,
2.4.2). Les $s_i$ sont alors des sections de $\mathcal L$ au-dessus de
$f^{-1}(V)$, définissant donc un homomorphisme
\[
v:\bigl(\mathcal O_X|f^{-1}(V)\bigr)^{r+1}
\longrightarrow\mathcal L|f^{-1}(V)
\]
qui, par hypothèse, est surjectif en tous les points de $f^{-1}(y)$~; comme
$\operatorname{Coker}(v)$ est cohérent (0\textsubscript{I}, 5.3.4), son
support est fermé (0\textsubscript{I}, 5.2.2) et par suite il existe un
voisinage ouvert $W\subset f^{-1}(V)$ de $f^{-1}(y)$ tel que la restriction
de $v$ à $W$ soit un homomorphisme surjectif. Puisque le morphisme $f$ est
fermé, on peut supposer que $W$ est de la forme $f^{-1}(U)$, où $U$ est un
voisinage ouvert de $y$, et la conclusion résulte alors de (II, 4.2.3).
\end{proof}

\subsection{Morphismes finis de préschémas formels}

\begin{proposition}[4.8.1]
\label{III.4.8.1-fr}
Soient $\mathfrak Y$ un préschéma formel localement noethérien, $\mathcal K$ un
Idéal de définition de $\mathfrak Y$, $f:\mathfrak X\to\mathfrak Y$ un
morphisme de préschémas formels. Les conditions suivantes sont équivalentes~:
\begin{enumerate}
\item[a)] $\mathfrak X$ est localement noethérien, $f$ est un morphisme adique
  (I, 10.12.1) et si l'on pose
  $\mathcal J=f^*(\mathcal K)\mathcal O_{\mathfrak X}$, le morphisme
  \[
  f_0:(\mathfrak X,\mathcal O_{\mathfrak X}/\mathcal J)
  \longrightarrow
  (\mathfrak Y,\mathcal O_{\mathfrak Y}/\mathcal K)
  \]
  déduit de $f$ est fini.
\item[b)] $\mathfrak X$ est localement noethérien et est limite inductive d'un
  $(Y_n)$-système inductif adique $(X_n)$ tel que le morphisme
  $X_0\to Y_0$ soit fini.
\item[c)] Tout point de $\mathfrak Y$ possède un voisinage ouvert formel affine
  noethérien $V$ tel que $f^{-1}(V)$ soit un ouvert formel affine et que
  $\Gamma(f^{-1}(V),\mathcal O_{\mathfrak X})$ soit un
  $\Gamma(V,\mathcal O_{\mathfrak Y})$-module de type fini.
\end{enumerate}
\end{proposition}

\begin{proof}
Il est immédiat que a) entraîne b) en vertu de (I, 10.12.3). Pour voir que b)
entraîne c), on peut supposer que
$\mathfrak Y=\operatorname{Spf}(B)$, où $B$ est adique noethérien et
$\mathcal K=\mathfrak R^\Delta$, où $\mathfrak R$ est un idéal de définition
de $B$. Par hypothèse, $X_0$ est un schéma affine dont l'anneau $A_0$ est un
$B/\mathfrak R$-module de type fini (II, 6.1.3). En vertu de (I, 5.1.9), chacun
des $X_n$ est un schéma affine, et si $A_n$ est son anneau, l'hypothèse b)
entraîne que pour $m\leq n$, $A_m$ est isomorphe à
$A_n/\mathfrak R^{m+1}A_n$. On en déduit que $\mathfrak X$ est isomorphe à
$\operatorname{Spf}(A)$, où $A=\varprojlim A_n$~; on conclut en vertu de
(0\textsubscript{I}, 7.2.9). Enfin, pour prouver que c)

\oldpage[III]{147}

entraîne a), on peut se borner encore au cas où
$\mathfrak Y=\operatorname{Spf}(B)$, $\mathfrak X=\operatorname{Spf}(A)$,
$A$ étant une $B$-algèbre finie~; comme $A/\mathfrak RA$ est alors une
$B/\mathfrak R$-algèbre finie, il résulte de (I, 10.10.9) que les conditions
de a) sont satisfaites.
\end{proof}

\begin{definition}[4.8.2]
\label{III.4.8.2-fr}
Lorsque les propriétés équivalentes a), b), c) de
(\hyperref[III.4.8.1-fr]{4.8.1}) sont vérifiées, on dit que le morphisme $f$
est \emph{fini}, ou que $\mathfrak X$ est un \emph{$\mathfrak Y$-préschéma
formel fini}, ou un préschéma formel \emph{fini au-dessus de $\mathfrak Y$}.
\end{definition}

\begin{proposition}[4.8.3]
\label{III.4.8.3-fr}
\begin{enumerate}
\item[(i)] Une immersion fermée de préschémas formels localement noethériens
  est un morphisme fini.
\item[(ii)] Le composé de deux morphismes finis de préschémas formels
  localement noethériens est un morphisme fini.
\item[(iii)] Soient $\mathfrak X$, $\mathfrak Y$, $\mathfrak S$ trois
  préschémas formels localement noethériens, $f:\mathfrak X\to\mathfrak S$ un
  morphisme fini, $g:\mathfrak Y\to\mathfrak S$ un morphisme~; alors le
  morphisme
  $\mathfrak X\times_{\mathfrak S}\mathfrak Y\to\mathfrak Y$ est fini.
\item[(iv)] Soient $\mathfrak S$ un préschéma formel localement noethérien,
  $\mathfrak X'$, $\mathfrak Y'$ deux préschémas formels localement
  noethériens tels que $\mathfrak X'\times_{\mathfrak S}\mathfrak Y'$ soit
  localement noethérien. Si $\mathfrak X$, $\mathfrak Y$ sont des
  $\mathfrak S$-préschémas formels localement noethériens,
  $f:\mathfrak X\to\mathfrak X'$, $g:\mathfrak Y\to\mathfrak Y'$ deux
  $\mathfrak S$-morphismes finis, alors $f\times_{\mathfrak S}g$ est un
  morphisme fini.
\item[(v)] Soient $f:\mathfrak X\to\mathfrak Y$,
  $g:\mathfrak Y\to\mathfrak Z$ deux morphismes de préschémas formels
  localement noethériens tels que $g$ soit de type fini et séparé~; alors, si
  $g\circ f$ est un morphisme fini, $f$ est un morphisme fini.
\end{enumerate}
\end{proposition}

\begin{proof}
(i) est trivial, et les autres assertions se ramènent aussitôt aux
propositions correspondantes pour les morphismes de préschémas usuels
(II, 6.1.5) à l'aide du critère a) de
(\hyperref[III.4.8.1-fr]{4.8.1})~; nous laissons les détails au lecteur, sur le
modèle de (I, 10.13.5).
\end{proof}

\begin{corollary}[4.8.4]
\label{III.4.8.4-fr}
Sous les hypothèses de (I, 10.9.9), si $f$ est un morphisme fini, il en est de
même de son prolongement $\widehat f$ aux complétés.
\end{corollary}

\begin{corollary}[4.8.5]
\label{III.4.8.5-fr}
Si $\mathfrak X$ est un préschéma formel fini au-dessus de $\mathfrak Y$,
$f:\mathfrak X\to\mathfrak Y$ le morphisme structural, alors, pour tout ouvert
$U\subset\mathfrak Y$, $f^{-1}(U)$ est fini au-dessus de $U$.
\end{corollary}

\begin{proposition}[4.8.6]
\label{III.4.8.6-fr}
Si $f:\mathfrak X\to\mathfrak Y$ est un morphisme fini de préschémas formels
localement noethériens, $f_*(\mathcal O_{\mathfrak X})$ est une
$\mathcal O_{\mathfrak Y}$-Algèbre cohérente.
\end{proposition}

\begin{proof}
On peut considérer $f$ comme limite inductive d'un système inductif $(f_n)$ de
morphismes $f_n:X_n\to Y_n$~; montrons que les $f_n$ sont des morphismes finis
et que $f_*(\mathcal O_{\mathfrak X})$ est isomorphe à la limite projective des
$(f_n)_*(\mathcal O_{X_n})$, ce qui établira notre assertion (I, 10.10.5). Il
suffit de se borner au cas où $\mathfrak Y=\operatorname{Spf}(B)$,
$\mathfrak X=\operatorname{Spf}(A)$, et de remarquer que si $\mathfrak R$ est
un idéal de définition de $B$ et $A$ un $B$-module de type fini,
$A/\mathfrak R^{n+1}A$ est un module de type fini sur
$B/\mathfrak R^{n+1}B$, et que $A$ est limite projective des
$A/\mathfrak R^{n+1}A$.
\end{proof}

Réciproquement~:

\begin{proposition}[4.8.7]
\label{III.4.8.7-fr}
Soient $\mathfrak Y$ un préschéma formel localement noethérien, $\mathcal A$
une $\mathcal O_{\mathfrak Y}$-Algèbre cohérente. Il existe un préschéma
formel $\mathfrak X$ fini au-dessus de $\mathfrak Y$, défini à un
$\mathfrak Y$-isomorphisme unique près, et tel que
$f_*(\mathcal O_{\mathfrak X})=\mathcal A$, $f:\mathfrak X\to\mathfrak Y$
étant le morphisme structural.
\end{proposition}

\begin{proof}
Soit $\mathcal K$ un Idéal de définition de $\mathfrak Y$, et posons
\[
Y_n=(\mathfrak Y,\mathcal O_{\mathfrak Y}/\mathcal K^{n+1}),\qquad
\mathcal A_n=\mathcal A/\mathcal K^{n+1}\mathcal A~;
\]
il est clair que $\mathcal A_n$ est une $\mathcal O_{Y_n}$-Algèbre finie et
définit donc un

\oldpage[III]{148}

$Y_n$-préschéma fini $X_n=\operatorname{Spec}(\mathcal A_n)$ (II, 6.1.3)~;
pour $m\leq n$, l'homomorphisme canonique surjectif
$h_{mn}:\mathcal A_n\to\mathcal A_m$ définit un morphisme
$u_{mn}:X_m\to X_n$ tel que le diagramme
\[
\xymatrix{
X_n\ar[d]_{f_n} & X_m\ar[l]_{u_{mn}}\ar[d]^{f_m}\\
Y_n & Y_m\ar[l]
}
\]
($f_n$ étant le morphisme structural) soit commutatif et identifie $X_m$ au
produit $X_n\times_{Y_n}Y_m$, comme on le voit aussitôt (II, 1.4.6). Le
préschéma formel $\mathfrak X$, limite inductive du système inductif $(X_n)$
est alors localement noethérien et tel que le morphisme structural
$f:\mathfrak X\to\mathfrak Y$, limite inductive du système $(f_n)$, soit fini
(\hyperref[III.4.8.1-fr]{4.8.1} et II, 10.12.3.1)~; on a vu en outre dans la
démonstration de (\hyperref[III.4.8.6-fr]{4.8.6}) que
$f_*(\mathcal O_{\mathfrak X})$ est limite projective des $\mathcal A_n$, donc
égale à $\mathcal A$ (I, 10.10.6). Quant à l'assertion d'unicité, elle est
conséquence du résultat plus général suivant~:
\end{proof}

\begin{proposition}[4.8.8]
\label{III.4.8.8-fr}
Soient $\mathfrak Y$ un préschéma formel localement noethérien,
$\mathfrak X$, $\mathfrak X'$ deux $\mathfrak Y$-préschémas formels finis
au-dessus de $\mathfrak Y$, $f:\mathfrak X\to\mathfrak Y$,
$f':\mathfrak X'\to\mathfrak Y$ les morphismes structuraux. Il existe une
bijection canonique de
\[
\operatorname{Hom}_{\mathfrak Y}(\mathfrak X,\mathfrak X')
\quad\hbox{sur}\quad
\operatorname{Hom}_{\mathcal O_{\mathfrak Y}}
\bigl(f'_*(\mathcal O_{\mathfrak X'}),f_*(\mathcal O_{\mathfrak X})\bigr)
\footnotemark.
\]
\footnotetext{La dernière expression désigne l'ensemble des homomorphismes de
$\mathcal O_{\mathfrak Y}$-Algèbres
$f'_*(\mathcal O_{\mathfrak X'})\to f_*(\mathcal O_{\mathfrak X})$.}
\end{proposition}

\begin{proof}
La définition de cette application $h\to\mathcal A(h)$ est la même que dans
(II, 1.1.2), et pour voir qu'elle est bijective, on est aussitôt ramené au cas
où $\mathfrak Y=\operatorname{Spf}(B)$ est un schéma formel affine noethérien.
Mais alors $\mathfrak X=\operatorname{Spf}(A)$,
$\mathfrak X'=\operatorname{Spf}(A')$, où $A$ et $A'$ sont deux $B$-algèbres
finies et $f_*(\mathcal O_{\mathfrak X})=A^\Delta$,
$f'_*(\mathcal O_{\mathfrak X'})=A'^\Delta$. La conclusion résulte alors de
la correspondance biunivoque, d'une part entre les $\mathfrak Y$-morphismes
$\mathfrak X\to\mathfrak X'$ et les $B$-homomorphismes (nécessairement
continus) $A'\to A$ qui sont des homomorphismes d'algèbres (I, 10.2.2), et
d'autre part entre les homomorphismes de $B$-modules $A'\to A$ et les
homomorphismes de $\mathcal O_{\mathfrak Y}$-Modules
$A'^\Delta\to A^\Delta$ (I, 10.10.2.3).
\end{proof}

\begin{corollary}[4.8.9]
\label{III.4.8.9-fr}
Dans la correspondance biunivoque canonique définie dans
(\hyperref[III.4.8.8-fr]{4.8.8}), les immersions fermées
$\mathfrak X\to\mathfrak X'$ correspondent aux homomorphismes surjectifs de
$\mathcal O_{\mathfrak Y}$-Algèbres
$f'_*(\mathcal O_{\mathfrak X'})\to f_*(\mathcal O_{\mathfrak X})$.
\end{corollary}

\begin{proof}
La question étant encore locale sur $\mathfrak Y$, on est ramené à la
définition des immersions fermées de préschémas formels localement noethériens
(I, 10.14.2).
\end{proof}

\begin{corollary}[4.8.10]
\label{III.4.8.10-fr}
Les notations et hypothèses étant celles de
(\hyperref[III.4.8.1-fr]{4.8.1}), pour qu'un morphisme adique $f$ soit une
immersion fermée, il faut et il suffit que $f_0$ soit une immersion fermée (de
préschémas usuels).
\end{corollary}

\begin{proof}
Cela résulte aussitôt de (\hyperref[III.4.8.9-fr]{4.8.9}) et de la condition de
surjectivité pour un homomorphisme de $\mathcal O_{\mathfrak Y}$-Modules
cohérents (I, 10.11.5).
\end{proof}

\begin{proposition}[4.8.11]
\label{III.4.8.11-fr}
Pour qu'un morphisme $f:\mathfrak X\to\mathfrak Y$ de préschémas formels localement

\oldpage[III]{149}

noethériens soit fini, il faut et il suffit qu'il soit propre et ait ses fibres
$f^{-1}(y)$ finies (pour tout $y\in\mathfrak Y$).
\end{proposition}

\begin{proof}
Grâce aux définitions (3.4.1 et 4.8.2), on est aussitôt ramené à la même proposition
pour $f_0$ (notations de (\hyperref[III.4.8.1-fr]{4.8.1})), ce qui n'est autre que
(4.4.2).
\end{proof}

\section{Un théorème d'existence de faisceaux algébriques cohérents}
\label{section:III.5-fr}

\subsection{Énoncé du théorème}
\label{subsection:III.5.1-fr}

\begin{env}[5.1.1]
\label{III.5.1.1-fr}
Soient $A$ un anneau adique noethérien, $\mathfrak J$ un idéal de définition de
$A$, de sorte que $A$ est séparé et complet pour la topologie $\mathfrak J$-adique.
Si $Y=\operatorname{Spec}(A)$, le schéma formel affine $\operatorname{Spf}(A)$
s'identifie au complété $\widehat Y$ de $Y$ le long de la partie fermée
$Y'=\mathrm V(\mathfrak J)$ (I, 10.10.1). Soient $X$ un $Y$-préschéma (usuel) de
type fini, $f:X\to Y$ le morphisme structural~; nous désignerons par $\widehat X$
le complété de $X$ le long de la partie fermée $X'=f^{-1}(Y')$, ou encore le
$\widehat Y$-préschéma formel $X\times_Y\widehat Y$~; par
$\widehat f:\widehat X\to\widehat Y$ le prolongement de $f$ aux complétés~; enfin,
pour tout $\mathcal O_X$-Module cohérent $\mathcal F$, nous noterons
$\widehat{\mathcal F}$ son complété, $\mathcal F_{/X'}$, qui est un
$\mathcal O_{\widehat X}$-Module cohérent.
\end{env}

\begin{proposition}[5.1.2]
\label{III.5.1.2-fr}
Les hypothèses et notations étant celles de
(\hyperref[III.5.1.1-fr]{5.1.1}), soit $\mathcal F$ un $\mathcal O_X$-Module
cohérent dont le support est propre sur $Y$ (II, 5.4.10). Les homomorphismes
canoniques (4.1.4)
\[
\rho_i:\mathrm H^i(X,\mathcal F)\longrightarrow
\mathrm H^i(\widehat X,\widehat{\mathcal F})
\]
sont alors des isomorphismes.
\end{proposition}

\begin{proof}
Comme $\mathrm H^i(X,\mathcal F)$ est un $A$-module de type fini (3.2.4), donc
identique à son séparé complété pour la topologie $\mathfrak J$-préadique
($0_{\mathrm I}$, 7.3.6), la proposition n'est qu'un cas particulier de (4.1.10).
\end{proof}

Rappelons que les isomorphismes canoniques $\rho_i$ commutent aux cobords pour
toute suite exacte
$0\to\mathcal F'\to\mathcal F\to\mathcal F''\to0$ de $\mathcal O_X$-Modules
cohérents (($0$, 12.1.6) et (I, 10.8.9)).

\begin{corollary}[5.1.3]
\label{III.5.1.3-fr}
Soient $\mathcal F$, $\mathcal G$ deux $\mathcal O_X$-Modules cohérents tels que
l'intersection de leurs supports soit propre sur $Y$. Alors l'homomorphisme canonique
\begin{equation}
\operatorname{Hom}_{\mathcal O_X}(\mathcal F,\mathcal G)
\longrightarrow
\operatorname{Hom}_{\mathcal O_{\widehat X}}
(\widehat{\mathcal F},\widehat{\mathcal G})
\tag{5.1.3.1}
\label{III.5.1.3.1.eq-fr}
\end{equation}
qui, à tout homomorphisme $u:\mathcal F\to\mathcal G$, fait correspondre son
complété $\widehat u:\widehat{\mathcal F}\to\widehat{\mathcal G}$, est un
isomorphisme. De plus, lorsque le morphisme $f$ est fermé, pour que $\widehat u$ soit
injectif (resp. surjectif), il faut et il suffit que $u$ le soit.
\end{corollary}

\begin{proof}
La première assertion est un cas particulier de (4.5.3), dû encore au fait que le
premier membre de (\ref{III.5.1.3.1.eq-fr}) est un $A$-module de type fini, donc
identique à son séparé complété. Pour démontrer la seconde, notons en vertu de
(I, 10.8.14) que $\widehat u$ est injectif (resp. surjectif) si et seulement s'il existe
un voisinage de $X'$ dans lequel $u$ soit injectif (resp. surjectif). La conclusion
résulte donc du lemme suivant~:
\end{proof}

\oldpage[III]{150}

\begin{lemma}[5.1.3.1]
\label{III.5.1.3.1-fr}
Sous les hypothèses de (\hyperref[III.5.1.1-fr]{5.1.1}), si on suppose en outre le
morphisme $f$ fermé, tout voisinage de $X'$ dans $X$ est identique à $X$.
\end{lemma}

\begin{proof}
Tout d'abord, on peut se ramener au cas où $f(X)=Y$. En effet, par hypothèse,
$f(X)$ est une partie fermée $Z$ de $Y$~; on peut en outre remplacer $f$ par
$f_{\mathrm{red}}$ (I, 6.3.4), et supposer par suite $X$ et $Y$ réduits~; on peut alors
remplacer $Y$ par le sous-préschéma fermé réduit de $Y$ ayant $Z$ pour espace
sous-jacent (I, 5.2.2), car tout idéal de $A$ est fermé, et tout anneau quotient de
$A$ est donc adique et noethérien. On a alors $f(X')=Y'$~; si $V$ est un voisinage
ouvert de $X'$ dans $X$, $f(X-V)$ est fermé dans $Y$ par hypothèse, et ne rencontre
pas $Y'$~; mais cela est impossible à moins que $X-V$ ne soit vide, puisque
$\mathfrak J$ est contenu dans le radical de $A$ (I, 1.1.15 et $0_{\mathrm I}$,
7.1.10), d'où la conclusion.
\end{proof}

Lorsqu'on se borne aux $\mathcal O_X$-Modules cohérents dont le support est propre
sur $Y$, (\hyperref[III.5.1.3-fr]{5.1.3}) peut s'énoncer, dans le langage des
catégories, en disant que le foncteur
$\mathcal F\mapsto\widehat{\mathcal F}$ est pleinement fidèle de la catégorie des
$\mathcal O_X$-Modules du type précédent, dans la catégorie des
$\mathcal O_{\widehat X}$-Modules cohérents, et établit par suite une équivalence de
la première de ces catégories avec une sous-catégorie pleine de la seconde
($0$, 8.1.6). Le théorème d'existence va prouver que lorsque $X$ est propre sur $Y$,
cette sous-catégorie est en fait la catégorie de tous les
$\mathcal O_{\widehat X}$-Modules cohérents. De façon précise~:

\begin{theorem}[5.1.4]
\label{III.5.1.4-fr}
Soient $A$ un anneau adique noethérien, $Y=\operatorname{Spec}(A)$,
$\mathfrak J$ un idéal de définition de $A$, $Y'=\mathrm V(\mathfrak J)$,
$f:X\to Y$ un morphisme séparé de type fini, $X'=f^{-1}(Y')$. Soient
$\widehat Y=Y_{/Y'}=\operatorname{Spf}(A)$, $\widehat X=X_{/X'}$ les complétés de
$Y$ et $X$ le long de $Y'$ et $X'$, $\widehat f:\widehat X\to\widehat Y$ le
prolongement de $f$ aux complétés~; alors, le foncteur
$\mathcal F\mapsto\mathcal F_{/X'}=\widehat{\mathcal F}$ est une équivalence de la
catégorie des $\mathcal O_X$-Modules cohérents de support propre sur
$\operatorname{Spec}(A)$, avec la catégorie des $\mathcal O_{\widehat X}$-Modules
cohérents de support propre sur $\operatorname{Spf}(A)$.
\end{theorem}

En d'autres termes, compte tenu de (\hyperref[III.5.1.3-fr]{5.1.3})~:

\begin{corollary}[5.1.5]
\label{III.5.1.5-fr}
Pour qu'un $\mathcal O_{\widehat X}$-Module soit isomorphe au complété d'un
$\mathcal O_X$-Module cohérent et de support propre sur $\operatorname{Spec}(A)$,
il faut et il suffit qu'il soit cohérent et de support propre sur
$\operatorname{Spf}(A)$.
\end{corollary}

Le cas le plus important est le

\begin{corollary}[5.1.6]
\label{III.5.1.6-fr}
Supposons $X$ propre sur $Y=\operatorname{Spec}(A)$. Alors le foncteur
$\mathcal F\mapsto\widehat{\mathcal F}$ est une équivalence de la catégorie des
$\mathcal O_X$-Modules cohérents et de la catégorie des
$\mathcal O_{\widehat X}$-Modules cohérents.
\end{corollary}

\begin{env}[5.1.7]
\label{III.5.1.7-fr}
\emph{Scholie.}
Si on tient compte de la caractérisation des faisceaux cohérents sur les préschémas
formels (I, 10.11.3), on voit que sous les conditions de
(\hyperref[III.5.1.1-fr]{5.1.1}), la donnée d'un $\mathcal O_X$-Module cohérent de
support propre sur $\operatorname{Spec}(A)$ équivaut (en posant
$Y_n=\operatorname{Spec}(A/\mathfrak J^{n+1})$ et $X_n=X\times_Y Y_n$) à la donnée
d'un système projectif de $\mathcal O_{X_n}$-Modules cohérents $(\mathcal F_n)$ tel
que pour $m\leq n$ on ait
\[
\mathcal F_m=\mathcal F_n\otimes_{\mathcal O_{Y_n}}\mathcal O_{Y_m}
\quad\text{(ou encore }\mathcal F_m=\mathcal F_n/\mathfrak J^{m+1}\mathcal F_n\text{)}
\]
et que le support de $\mathcal F_0$ soit une partie de $X_0$ propre sur $Y_0$. Au
moyen de (I, 10.11.4), on interprète de même les homomorphismes de
$\mathcal O_X$-Modules cohérents comme des homomorphismes de systèmes projectifs de
$\mathcal O_{X_n}$-Modules cohérents.

\oldpage[III]{151}

Dans tous les cas d'application connus, $A$ est en fait un anneau adique local
noethérien, donc les $Y_n$ sont des spectres d'anneaux artiniens locaux, et les
résultats de ce paragraphe et des précédents réduisent dans une large mesure la
géométrie algébrique sur un anneau adique local noethérien, à la géométrie algébrique
sur des anneaux locaux artiniens.
\end{env}

\begin{corollary}[5.1.8]
\label{III.5.1.8-fr}
Sous les conditions de (\hyperref[III.5.1.4-fr]{5.1.4}), l'application
$Z\mapsto\widehat Z=Z_{/(Z\cap X')}$ est une bijection de l'ensemble des
sous-préschémas fermés $Z$ de $X$, propres sur $Y$, sur l'ensemble des
sous-préschémas formels fermés de $\widehat X$, propres sur $\widehat Y$.
\end{corollary}

\begin{proof}
En effet, un sous-préschéma formel fermé de $\widehat X$ est de la forme
$(T,(\mathcal O_{\widehat X}/\mathcal A)|T)$, où $\mathcal A$ est un Idéal
cohérent de $\mathcal O_{\widehat X}$ (I, 10.14.2)~; si $T$ est propre sur
$\widehat Y$, il résulte de (\hyperref[III.5.1.4-fr]{5.1.4}) que
$\mathcal O_{\widehat X}/\mathcal A$ est isomorphe à un
$\mathcal O_{\widehat X}$-Module de la forme $\widehat{\mathcal F}$, où
$\mathcal F$ est un $\mathcal O_X$-Module cohérent de support propre sur $Y$~; en
outre, il résulte de (\hyperref[III.5.1.3-fr]{5.1.3}) que l'homomorphisme canonique
$\mathcal O_{\widehat X}\to\mathcal O_{\widehat X}/\mathcal A$ est de la forme
$\widehat u$, où $u:\mathcal O_X\to\mathcal F$ est un homomorphisme surjectif de
$\mathcal O_X$-Modules. Donc $\mathcal F$ est de la forme
$\mathcal O_X/\mathcal N$, où $\mathcal N$ est un Idéal cohérent de
$\mathcal O_X$, et $\mathcal A=\widehat{\mathcal N}$ (I, 10.8.8), d'où la
conclusion (I, 10.14.7).
\end{proof}

\subsection{Démonstration du théorème d'existence : cas projectif et quasi-projectif}
\label{subsection:III.5.2-fr}

\begin{env}[5.2.1]
\label{III.5.2.1-fr}
Sous les conditions de (\hyperref[III.5.1.4-fr]{5.1.4}), nous dirons provisoirement
qu'un $\mathcal O_{\widehat X}$-Module cohérent est algébrisable s'il est isomorphe
à un complété $\widehat{\mathcal F}$ d'un $\mathcal O_X$-Module cohérent
$\mathcal F$ de support propre sur $Y$.
\end{env}

\begin{lemma}[5.2.2]
\label{III.5.2.2-fr}
Soient $\mathcal F'$, $\mathcal G'$ deux $\mathcal O_{\widehat X}$-Modules
algébrisables. Pour tout homomorphisme $u:\mathcal F'\to\mathcal G'$,
$\operatorname{Ker}(u)$, $\operatorname{Im}(u)$ et $\operatorname{Coker}(u)$ sont
algébrisables.
\end{lemma}

\begin{proof}
En effet, si $\mathcal F'=\widehat{\mathcal F}$,
$\mathcal G'=\widehat{\mathcal G}$, où $\mathcal F$ et $\mathcal G$ sont des
$\mathcal O_X$-Modules cohérents de supports propres sur $Y$, on a
$u=\widehat v$, où $v:\mathcal F\to\mathcal G$ est un homomorphisme
(\hyperref[III.5.1.3-fr]{5.1.3}). En vertu de l'exactitude du foncteur
$\mathcal F\mapsto\widehat{\mathcal F}$,
$\operatorname{Ker}(\widehat v)$ est isomorphe à
$\widehat{\operatorname{Ker}(v)}$ et comme le support de
$\operatorname{Ker}(v)$ est contenu dans celui de $\mathcal F$, on voit que
$\operatorname{Ker}(u)$ est algébrisable~; démonstration analogue pour
$\operatorname{Im}(u)$ et $\operatorname{Coker}(u)$.
\end{proof}

\begin{proposition}[5.2.3]
\label{III.5.2.3-fr}
Soient $A$ un anneau adique noethérien, $\mathfrak J$ un idéal de définition de
$A$, $\mathfrak Y=\operatorname{Spf}(A)$,
$f:\mathfrak X\to\mathfrak Y$ un morphisme propre de préschémas formels. On pose
$Y_k=\operatorname{Spec}(A/\mathfrak J^{k+1})$,
$X_k=\mathfrak X\times_{\mathfrak Y}Y_k$, et pour tout
$\mathcal O_{\mathfrak X}$-Module $\mathcal F$,
\[
\mathcal F_k=\mathcal F\otimes_{\mathcal O_{\mathfrak X}}\mathcal O_{X_k}
=\mathcal F/\mathfrak J^{k+1}\mathcal F.
\]
Soit $\mathcal L$ un $\mathcal O_{\mathfrak X}$-Module inversible, et supposons que
$\mathcal L_0=\mathcal L/\mathfrak J\mathcal L$ soit un
$\mathcal O_{X_0}$-Module ample~; pour tout $\mathcal O_{\mathfrak X}$-Module
$\mathcal F$ et tout entier $n$, posons
$\mathcal F(n)=\mathcal F\otimes\mathcal L^{\otimes n}$. Alors, pour tout
$\mathcal O_{\mathfrak X}$-Module cohérent $\mathcal F$, il existe un entier $n_0$
tel que, pour tout $n\geq n_0$, les propriétés suivantes aient lieu~:
\begin{enumerate}
\item[(i)] L'homomorphisme canonique
$\mathrm H^0(\mathfrak X,\mathcal F(n))\to
\mathrm H^0(\mathfrak X,\mathcal F_k(n))$ est surjectif pour tout $k\geq0$.
\item[(ii)] On a $\mathrm H^q(\mathfrak X,\mathcal F(n))=0$ pour tout $q>0$.
\end{enumerate}
\end{proposition}

\begin{proof}
On sait que les espaces sous-jacents à $\mathfrak X$ et à $X_0$ sont les mêmes~;
les faisceaux
\[
\mathcal M_k=\mathfrak J^k\mathcal F/\mathfrak J^{k+1}\mathcal F
\]
étant annulés par $\mathfrak J$, peuvent être considérés comme des
$\mathcal O_{X_0}$-Modules cohérents ($0_{\mathrm I}$, 5.3.10)~; en outre, si on
pose $\mathcal M_k(n)=\mathcal M_k\otimes_{\mathcal O_{X_0}}
\mathcal L_0^{\otimes n}$, on voit aussitôt que
\[
\mathcal M_k(n)=\mathfrak J^k\mathcal F(n)/
\mathfrak J^{k+1}\mathcal F(n).
\]
Notons que, puisque $\mathcal L_0$ est ample pour $f_0:X_0\to Y_0$, et

\oldpage[III]{152}

que $f_0$ est propre, on en conclut que $f_0$ est projectif (II, 5.5.4). Soit
\[
S=\bigoplus_{k\geq0}\mathfrak J^k/\mathfrak J^{k+1}
\]
la $A$-algèbre graduée associée à la filtration $\mathfrak J$-adique de $A$, qui est
de type fini puisque $A$ est noethérien~; si on pose
$\mathcal S'=f_0^*(\widetilde S)$,
$\mathcal M=\bigoplus_{k\geq0}\mathcal M_k$, $\mathcal S'$ est une
$\mathcal O_{X_0}$-Algèbre quasi-cohérente de type fini, et $\mathcal M$ un
$\mathcal S'$-Module gradué quasi-cohérent et de type fini (puisque $\mathcal F_0$
est cohérent et engendre le $\mathcal S'$-Module $\mathcal M$). On est donc dans les
conditions d'application du théorème (2.4.1, (ii)), et on en conclut qu'il existe
$n_0$ tel que, pour $n\geq n_0$ et pour tout $k$, on ait
\begin{equation}
\mathrm H^q(X_0,\mathcal M_k(n))=0\qquad\text{pour tout }q>0.
\tag{5.2.3.1}
\label{III.5.2.3.1-fr}
\end{equation}
On a donc aussi $\mathrm H^q(\mathfrak X,\mathcal M_k(n))=0$ pour $q>0$ et
$n\geq n_0$, $\mathcal M_k(n)$ étant cette fois considéré comme
$\mathcal O_{\mathfrak X}$-Module. Appliquant la suite exacte de cohomologie à
\[
0\to
\frac{\mathfrak J^k\mathcal F(n)}{\mathfrak J^{k+1}\mathcal F(n)}
\to
\frac{\mathfrak J^h\mathcal F(n)}{\mathfrak J^{k+1}\mathcal F(n)}
\to
\frac{\mathfrak J^h\mathcal F(n)}{\mathfrak J^k\mathcal F(n)}
\to0,
\]
on en déduit tout d'abord que pour $0\leq h<k$, $n\geq n_0$ et $q>0$, on a par
récurrence sur $k-h$
\begin{equation}
\mathrm H^q\left(\mathfrak X,
\frac{\mathfrak J^h\mathcal F(n)}{\mathfrak J^k\mathcal F(n)}\right)=0
\tag{5.2.3.2}
\label{III.5.2.3.2-fr}
\end{equation}
et en particulier pour $h=0$
\begin{equation}
\mathrm H^q(\mathfrak X,\mathcal F_k(n))=0
\qquad\text{pour }n\geq n_0,\ k\geq0\text{ et }q>0.
\tag{5.2.3.3}
\label{III.5.2.3.3-fr}
\end{equation}
Une autre portion de la suite exacte de cohomologie, pour $h=0$, donne la suite
exacte
\begin{equation}
\mathrm H^0(\mathfrak X,\mathcal F_{k+1}(n))\to
\mathrm H^0(\mathfrak X,\mathcal F_k(n))\to
\mathrm H^1\left(\mathfrak X,
\frac{\mathfrak J^k\mathcal F(n)}{\mathfrak J^{k+1}\mathcal F(n)}\right)=0
\tag{5.2.3.4}
\label{III.5.2.3.4-fr}
\end{equation}
d'où on déduit que pour $h\leq k$, l'application canonique
\begin{equation}
\mathrm H^0(\mathfrak X,\mathcal F_k(n))\longrightarrow
\mathrm H^0(\mathfrak X,\mathcal F_h(n))
\tag{5.2.3.5}
\label{III.5.2.3.5-fr}
\end{equation}
est surjective. Pour tout $q$, le système projectif
$(\mathrm H^q(\mathfrak X,\mathcal F_k(n)))_{k\geq0}$ vérifie donc la condition
(ML) pour $n\geq n_0$. Par ailleurs, tout ouvert formel affine $U$ de $\mathfrak X$
est aussi un ouvert affine dans chacun des $X_k$ (I, 10.5.2), donc on a
$\mathrm H^q(U,\mathcal F_k(n))=0$ pour tout $q>0$ (I, 1.3.1), et
$\mathrm H^0(U,\mathcal F_k(n))\to\mathrm H^0(U,\mathcal F_h(n))$ est surjective
pour $h\leq k$ (I, 1.3.9). Les conditions d'application de ($0$, 13.3.1) sont par
suite remplies, et on en conclut que, pour $n\geq n_0$~:
\begin{enumerate}
\item[$1^\circ$] Pour tout $q>0$,
\[
\mathrm H^q(\mathfrak X,\mathcal F(n))\longrightarrow
\varprojlim_k\mathrm H^q(\mathfrak X,\mathcal F_k(n))
\]
est bijectif, donc, en vertu de (\ref{III.5.2.3.3-fr}),
$\mathrm H^q(\mathfrak X,\mathcal F(n))=0$.
\item[$2^\circ$] L'homomorphisme
\[
\mathrm H^0(\mathfrak X,\mathcal F(n))\longrightarrow
\varprojlim_k\mathrm H^0(\mathfrak X,\mathcal F_k(n))
\]
est bijectif~; par ailleurs, comme les homomorphismes
(\ref{III.5.2.3.5-fr}) sont surjectifs, il en est de même de chacun des
homomorphismes
\[
\varprojlim_k\mathrm H^0(\mathfrak X,\mathcal F_k(n))
\longrightarrow\mathrm H^0(\mathfrak X,\mathcal F_h(n)),
\]
\end{enumerate}
ce qui achève la démonstration.
\end{proof}

\begin{corollary}[5.2.4]
\label{III.5.2.4-fr}
Les hypothèses étant celles de (\hyperref[III.5.2.3-fr]{5.2.3}), pour tout
$\mathcal O_{\mathfrak X}$-Module cohérent $\mathcal F$, il existe un entier $N$
tel que pour $n\geq N$, $\mathcal F(n)$ soit engendré par ses sections au-dessus

\oldpage[III]{153}

de $\mathfrak X$~; en d'autres termes, $\mathcal F$ est isomorphe au quotient d'un
$\mathcal O_{\mathfrak X}$-Module de la forme
$(\mathcal O_{\mathfrak X}(-n))^k$.
\end{corollary}

\begin{proof}
Comme $X_0$ est noethérien, il résulte de l'hypothèse sur $\mathcal L_0$ et de
(II, 4.5.5) qu'il existe $n_0$ tel que, pour $n\geq n_0$, $\mathcal F_0(n)$ soit
engendré par ses sections au-dessus de $\mathfrak X$~; par ailleurs, on peut supposer
$n_0$ pris assez grand pour que l'homomorphisme
\[
\Gamma(\mathfrak X,\mathcal F(n))\longrightarrow
\Gamma(\mathfrak X,\mathcal F_0(n))
\]
soit surjectif pour $n\geq n_0$ (\hyperref[III.5.2.3-fr]{5.2.3}). Il existe donc
un nombre fini de sections $s_i\in\Gamma(\mathfrak X,\mathcal F(n))$ dont les
images dans $\Gamma(\mathfrak X,\mathcal F_0(n))$ engendrent $\mathcal F_0(n)$
($0_{\mathrm I}$, 5.2.3). Comme $\mathfrak J$ est contenu dans l'idéal maximal de
l'anneau local en tout point de $\mathfrak X$, il résulte du lemme de Nakayama,
appliqué à ces anneaux locaux, que les $s_i$ engendrent $\mathcal F(n)$
($0_{\mathrm I}$, 5.1.1).
\end{proof}

\begin{env}[5.2.5]
\label{III.5.2.5-fr}
\emph{Démonstration du théorème d'existence : cas projectif.}
Les notations étant celles de (\hyperref[III.5.1.4-fr]{5.1.4}), supposons $f$
projectif, de sorte qu'il existe un $\mathcal O_X$-Module ample $\mathcal L$
(I, 5.5.4). Par définition, $X_n=\widehat X\times_{\widehat Y}Y_n$ est égal au
sous-préschéma fermé $X\times_Y Y_n=f^{-1}(Y_n)$ de $X$~; si
$\mathcal L'=\mathcal L\otimes_{\mathcal O_X}\mathcal O_{\widehat X}$ est le
complété de $\mathcal L$, on a donc
$\mathcal L'_0=\mathcal L/\mathfrak J\mathcal L$, considéré comme
$\mathcal O_{X_0}$-Module~; on sait que $\mathcal L'_0$ est ample
(II, 4.6.13, (i bis)). On peut donc appliquer à $\mathcal L'$ et à tout
$\mathcal O_{\widehat X}$-Module cohérent $\mathcal F$ le cor.
(\hyperref[III.5.2.4-fr]{5.2.4})~; on voit donc que $\mathcal F$ est isomorphe à
un quotient de
\[
\mathcal G=(\mathcal L'^{\otimes(-n)})^k
\]
pour des entiers $n>0$ et $k>0$ convenables. Or, il est clair que $\mathcal G$ est
le complété de $(\mathcal L^{\otimes(-n)})^k$ (I, 10.8.10), donc est
algébrisable. Considérons ensuite l'homomorphisme canonique surjectif
$u:\mathcal G\to\mathcal F$, et soit $\mathcal H=\operatorname{Ker}(u)$, qui est
un $\mathcal O_{\widehat X}$-Module cohérent ($0_{\mathrm I}$, 5.3.4). On voit de
même qu'il existe un $\mathcal O_{\widehat X}$-Module algébrisable $\mathcal K$ et
un homomorphisme $v:\mathcal K\to\mathcal G$ tel que
$\mathcal H=\operatorname{Im}(v)$. On a alors
$\mathcal F=\operatorname{Coker}(v)$, et $\mathcal F$ est algébrisable en vertu de
(\hyperref[III.5.2.2-fr]{5.2.2}).
\end{env}

\begin{env}[5.2.6]
\label{III.5.2.6-fr}
\emph{Démonstration du théorème d'existence; cas quasi-projectif.}
Les notations étant toujours celles de (\hyperref[III.5.1.4-fr]{5.1.4}), supposons
maintenant que $f$ soit quasi-projectif. Il existe alors un morphisme projectif
$g:Z\to Y$ tel que $X$ s'identifie au $Y$-préschéma induit sur un ouvert de $Z$
(II, 5.3.2)~; si on pose $Z'=g^{-1}(Y')$, on a $X'=X\cap Z'$. Par suite, le
complété $\widehat X=X_{/X'}$ s'identifie au préschéma formel induit par le complété
$\widehat Z=Z_{/Z'}$ sur l'ouvert $X\cap Z'$ de $\widehat Z$ (I, 10.8.5). Soit
$\mathcal F'$ un $\mathcal O_{\widehat X}$-Module cohérent, dont le support $T'$ est
propre sur $\widehat Y$~; cela signifie par définition qu'il existe un sous-préschéma
fermé de $X'$, ayant $T'\subset X'$ comme espace sous-jacent, tel que la restriction
$T'\to Y'$ de $f$ soit propre~; on en conclut que $T'$ est propre sur $Y$, donc
fermé dans $Z'$ (II, 5.4.10). Il en résulte que $\mathcal F'$ est le faisceau induit
sur $\widehat X$ par le $\mathcal O_{\widehat Z}$-Module $\mathcal G'$ obtenu par
recollement de $\mathcal F'$ (défini sur l'ouvert $\widehat X$ de $\widehat Z$) et
du faisceau $0$ sur l'ouvert $\widehat Z-T'$ de $\widehat Z$, ces deux faisceaux
coïncidant dans l'ouvert intersection $\widehat X-T'$. Il est clair que
$\mathcal G'$ est cohérent~; en vertu de
(\hyperref[III.5.2.5-fr]{5.2.5}), il existe un $\mathcal O_Z$-Module cohérent
$\mathcal G$ tel que $\mathcal G'=\widehat{\mathcal G}$~; soit $T$ le support de
$\mathcal G$, de sorte que $T'=T\cap Z'$ (I, 10.8.12). Si $h$ est la restriction
de $g$ au sous-préschéma fermé réduit de $Z$ ayant $T$ comme espace sous-jacent, on
a donc $T'=h^{-1}(Y')=T\cap g^{-1}(Y')$, et par suite $X\cap T$ est un ouvert de
$T$ contenant $T'$.

\oldpage[III]{154}

Comme $h$ est propre (II, 5.4.2), donc fermé, il résulte de
(\hyperref[III.5.1.3.1-fr]{5.1.3.1}) que $X\cap T=T$, autrement dit $T\subset X$,
et comme $T$ est fermé dans $Z$, $T$ est propre sur $Y$. Si $\mathcal F$ est le
faisceau induit sur $X$ par $\mathcal G$, son complété $\widehat{\mathcal F}$ est
induit sur $\widehat X$ par $\widehat{\mathcal G}$ (I, 10.8.4), donc est égal à
$\mathcal F'$, ce qui achève la démonstration.
\end{env}

\subsection{Démonstration du théorème d'existence : cas général}
\label{subsection:III.5.3-fr}

\begin{lemma}[5.3.1]
\label{III.5.3.1-fr}
Sous les conditions de (\hyperref[III.5.1.4-fr]{5.1.4}), si
\[
0\to\mathcal H\to\mathcal F\to\mathcal G\to0
\]
est une suite exacte de $\mathcal O_{\widehat X}$-Modules cohérents telle que
$\mathcal G$ et $\mathcal H$ soient algébrisables, alors $\mathcal F$ est
algébrisable.
\end{lemma}

\begin{proof}
En effet, supposons que $\mathcal G=\widehat{\mathcal B}$,
$\mathcal H=\widehat{\mathcal C}$, $\mathcal B$ et $\mathcal C$ étant des
$\mathcal O_X$-Modules cohérents à supports propres sur $Y$~; $\mathcal F$ définit
canoniquement un élément du $A$-module
\[
\operatorname{Ext}^1_{\mathcal O_{\widehat X}}
(\widehat X;\widehat{\mathcal B},\widehat{\mathcal C})
\]
($0$, 12.3.2), et les hypothèses entraînent que ce $A$-module est canoniquement
isomorphe à
\[
\operatorname{Ext}^1_{\mathcal O_X}(X;\mathcal B,\mathcal C)
\]
(4.5.2)~; il existe donc une suite exacte
$0\to\mathcal C\to\mathcal A\to\mathcal B\to0$ de $\mathcal O_X$-Modules
cohérents telle que l'image canonique de l'élément de
$\operatorname{Ext}^1_{\mathcal O_X}(X;\mathcal B,\mathcal C)$ correspondant à
$\mathcal A$ soit l'élément de
$\operatorname{Ext}^1_{\mathcal O_{\widehat X}}
(\widehat X;\widehat{\mathcal B},\widehat{\mathcal C})$ correspondant à
$\mathcal F$. Mais par définition (compte tenu de (I, 10.8.8, (iii))), cela signifie
que $\mathcal F$ est isomorphe à $\widehat{\mathcal A}$, d'où le lemme, car
$\operatorname{Supp}(\mathcal A)$ est contenu dans la réunion de
$\operatorname{Supp}(\mathcal B)$ et $\operatorname{Supp}(\mathcal C)$, donc est
propre sur $Y$.
\end{proof}

\begin{corollary}[5.3.2]
\label{III.5.3.2-fr}
Sous les conditions de (\hyperref[III.5.1.1-fr]{5.1.1}), soit
$u:\mathcal F\to\mathcal G$ un homomorphisme de
$\mathcal O_{\widehat X}$-Modules cohérents~; si $\mathcal G$,
$\operatorname{Ker}(u)$ et $\operatorname{Coker}(u)$ sont algébrisables, il en est
de même de $\mathcal F$.
\end{corollary}

\begin{proof}
Le lemme (\hyperref[III.5.2.2-fr]{5.2.2}) appliqué à l'homomorphisme
$\mathcal G\to\operatorname{Coker}(u)$ montre en effet que
$\operatorname{Im}(u)$ est algébrisable, et il suffit ensuite d'appliquer le lemme
(\hyperref[III.5.3.1-fr]{5.3.1}) à la suite exacte
\[
0\to\operatorname{Ker}(u)\to\mathcal F\to\operatorname{Im}(u)\to0.
\]
\end{proof}

\begin{lemma}[5.3.3]
\label{III.5.3.3-fr}
Sous les conditions de (\hyperref[III.5.1.1-fr]{5.1.1}), soient $h:Z\to Y$ un
morphisme de type fini, $\widehat Z$ le complété de $Z$ le long de
$Z'=h^{-1}(Y')$, $g:Z\to X$ un $Y$-morphisme propre,
$\widehat g:\widehat Z\to\widehat X$ son prolongement aux complétés. Pour tout
$\mathcal O_{\widehat Z}$-Module algébrisable $\mathcal F'$,
$\widehat g_*(\mathcal F')$ est un $\mathcal O_{\widehat X}$-Module algébrisable.
\end{lemma}

\begin{proof}
En effet, si $\mathcal F$ est un $\mathcal O_Z$-Module cohérent tel que
$\mathcal F'=\widehat{\mathcal F}$, il résulte du premier th. de comparaison
(4.1.5) que $\widehat g_*(\widehat{\mathcal F})$ est isomorphe au complété de
$g_*(\mathcal F)$.
\end{proof}

\begin{lemma}[5.3.4]
\label{III.5.3.4-fr}
Soient $X$ un schéma (usuel) noethérien, $X'$ une partie fermée de $X$,
$f:Z\to X$ un morphisme propre, $Z'=f^{-1}(X')$,
$\widehat X=X_{/X'}$, $\widehat Z=Z_{/Z'}$,
$\widehat f:\widehat Z\to\widehat X$ le prolongement de $f$ aux complétés. Soit
$\mathcal M$ un Idéal cohérent de $\mathcal O_X$ tel que, si
$U=X-\operatorname{Supp}(\mathcal O_X/\mathcal M)$, la restriction
$f^{-1}(U)\to U$ de $f$ soit un isomorphisme. Alors, pour tout
$\mathcal O_{\widehat X}$-Module cohérent $\mathcal F$, il existe un entier $n>0$
tel que le noyau et le conoyau de l'homomorphisme canonique
\[
\rho_{\mathcal F}:\mathcal F\longrightarrow
\widehat f_*\bigl(\widehat f^{,*}(\mathcal F)\bigr)
\]
($0_{\mathrm I}$, 4.4.3) soient annulés par $\widehat{\mathcal M}^{,n}$.
\end{lemma}

\begin{proof}
On peut se borner au cas où $X=\operatorname{Spec}(B)$ où $B$ est un anneau
noethérien, donc $X'=\mathrm V(\mathfrak K)$, où $\mathfrak K$ est un idéal de
$B$. Nous allons voir qu'on peut se ramener au cas où $B$ est un anneau adique
noethérien et $\mathfrak K$ un idéal de définition de $B$. Soit en effet $B_1$ le
séparé complété de $B$ pour la topologie $\mathfrak K$-préadique~; si
$\mathfrak K_1=\mathfrak K B_1$, $B_1$ est donc un

\oldpage[III]{155}

anneau adique noethérien dont $\mathfrak K_1$ est un idéal de définition. Posons
$X_1=\operatorname{Spec}(B_1)$ et soit $h:X_1\to X$ le morphisme correspondant à
l'homomorphisme canonique $B\to B_1$~; si $X'_1=h^{-1}(X')$, on a donc
$X'_1=\mathrm V(\mathfrak K_1)$. Posons enfin
$Z_1=Z\times_X X_1=Z_{(X_1)}$, $f_1=f_{(X_1)}:Z_1\to X_1$, qui est un morphisme
propre (II, 5.4.2), et désignons par $\widehat X_1$ le complété de $X_1$ le long
de $X'_1$, par
$\widehat Z_1=Z_1\times_{X_1}\widehat X_1$ le complété de $Z_1$ le long de
$Z'_1=f_1^{-1}(X'_1)$, par $\widehat f_1$ le prolongement de $f_1$ aux complétés.
Il est immédiat que le prolongement $\widehat h:\widehat X_1\to\widehat X$ de $h$
aux complétés est un isomorphisme, correspondant à l'application identique de $B_1$
(I, 10.9.1)~; on en conclut que l'homomorphisme correspondant
$\widehat Z_1\to\widehat Z$ est aussi un isomorphisme, ces isomorphismes identifiant
$\widehat f_1$ et $\widehat f$. Enfin,
$\mathcal M_1=h^*(\mathcal M)$ est un Idéal cohérent de $\mathcal O_{X_1}$ et
\[
\operatorname{Supp}(\mathcal O_{X_1}/\mathcal M_1)
=h^{-1}(\operatorname{Supp}(\mathcal O_X/\mathcal M))
\]
(I, 9.1.13), donc, si
$U_1=X_1-\operatorname{Supp}(\mathcal O_{X_1}/\mathcal M_1)$, on a
$U_1=h^{-1}(U)$, d'où résulte aussitôt que la restriction
$f_1^{-1}(U_1)\to U_1$ de $f_1$ est un isomorphisme (I, 3.2.7)~; en outre, les
complétés $\widehat{\mathcal M}$ et $\widehat{\mathcal M}_1$ s'identifient par
$\widehat h$ (I, 10.9.5). Toutes les hypothèses de
(\hyperref[III.5.3.4-fr]{5.3.4}) sont donc remplies par $X_1$, $X'_1$, $f_1$ et
$\mathcal M_1$, et on peut donc désormais supposer $B$ adique noethérien et
$\mathfrak K$ un idéal de définition de $B$. On a alors
$\widehat X=\operatorname{Spf}(B)$, et $\mathcal F=N^\Delta$, où $N$ est un
$B$-module de type fini, d'où $\mathcal F=\widehat{\mathcal G}$, où $\mathcal G$
est le $\mathcal O_X$-Module cohérent $\widetilde N$ (I, 10.10.5), et par suite
\[
\widehat f^{,*}(\mathcal F)=\widehat{f^*(\mathcal G)}
\]
(I, 10.9.5). En outre, en vertu du premier théorème de comparaison (4.1.5),
\[
\widehat f_*\bigl(\widehat{f^*(\mathcal G)}\bigr)
\simeq\widehat{f_*\bigl(f^*(\mathcal G)\bigr)},
\]
et l'homomorphisme canonique $\rho_{\mathcal F}$ n'est autre que
$\widehat{\rho}_{\mathcal G}$ en vertu de
(\hyperref[III.5.1.3-fr]{5.1.3}). Or, le noyau $\mathcal P$ et le conoyau
$\mathcal R$ de
\[
\rho_{\mathcal G}:\mathcal G\longrightarrow f_*\bigl(f^*(\mathcal G)\bigr)
\]
sont des $\mathcal O_X$-Modules cohérents, et par hypothèse leurs restrictions à
$U$ sont évidemment nulles. Il existe par suite un entier $n>0$ tel que
$\mathcal M^n\mathcal P=\mathcal M^n\mathcal R=0$ (I, 9.3.4)~; on en conclut que
$\widehat{\mathcal M}^{,n}\widehat{\mathcal P}=
\widehat{\mathcal M}^{,n}\widehat{\mathcal R}=0$
(I, 10.8.8 et 10.8.10).
\end{proof}

\begin{env}[5.3.5]
\label{III.5.3.5-fr}
\emph{Fin de la démonstration du théorème d'existence.}
Les hypothèses étant celles de (\hyperref[III.5.1.4-fr]{5.1.4}), nous allons
utiliser le principe de récurrence noethérienne ($0_{\mathrm I}$, 2.2.2) en supposant
donc le théorème vrai pour tout sous-préschéma fermé $T$ de $X$ dont l'espace
sous-jacent est distinct de $X$ (le complété $\widehat T$ étant bien entendu le
complété de $T$ le long de $T\cap X'$). On peut supposer $X$ non vide. Comme $f$
est séparé et de type fini, on peut appliquer le lemme de Chow (II, 5.6.1)~: il
existe donc un $Y$-schéma $Z$ et un $Y$-morphisme $g:Z\to X$ tels que le morphisme
structural $h:Z\to Y$ soit quasi-projectif, le morphisme $g$ projectif et surjectif,
et en outre un ouvert non vide $U$ de $X$ tel que la restriction
$g^{-1}(U)\to U$ soit un isomorphisme. Soit $\mathcal M$ un Idéal cohérent de
$\mathcal O_X$ définissant un sous-préschéma fermé d'espace sous-jacent $X-U$
(I, 5.2.2), et soit $\mathcal F$ un $\mathcal O_{\widehat X}$-Module cohérent dont
le support $E$ est propre sur $Y$~; désignons par $\widehat Z$ le complété de $Z$
le long de $h^{-1}(Y')$, par $\widehat g:\widehat Z\to\widehat X$ le prolongement
de $g$ aux complétés. Alors $\widehat g^{,*}(\mathcal F)$ est un
$\mathcal O_{\widehat Z}$-Module cohérent dont le support est contenu dans
$g^{-1}(E)$ et est par suite propre sur $Y$, puisque $g$ est projectif, donc propre
(II, 5.4.6). Comme $h$ est quasi-projectif, $\widehat g^{,*}(\mathcal F)$ est
algébrisable en vertu de (\hyperref[III.5.2.6-fr]{5.2.6}). On en conclut

\oldpage[III]{156}

que $\widehat g_*\bigl(\widehat g^{,*}(\mathcal F)\bigr)$ est un
$\mathcal O_{\widehat X}$-Module algébrisable
(\hyperref[III.5.3.3-fr]{5.3.3}) puisque $g$ est propre. On peut maintenant
appliquer à $\mathcal F$ et à $g$ le résultat de
(\hyperref[III.5.3.4-fr]{5.3.4})~: le noyau $\mathcal P$ et le conoyau
$\mathcal R$ de l'homomorphisme
$\rho_{\mathcal F}:\mathcal F\to
\widehat g_*\bigl(\widehat g^{,*}(\mathcal F)\bigr)$ sont annulés par une
puissance $\widehat{\mathcal M}^{,n}$~; soient $T$ le sous-préschéma fermé de $X$
défini par $\mathcal M^n$, ayant $X-U$ pour espace sous-jacent,
$j:T\to X$ l'injection canonique, de sorte que le prolongement aux complétés
$\widehat j:\widehat T\to\widehat X$ est l'injection canonique (I, 10.14.7). On
peut donc écrire
$\mathcal P=\widehat j_*\bigl(\widehat j^{,*}(\mathcal P)\bigr)$ et
$\mathcal R=\widehat j_*\bigl(\widehat j^{,*}(\mathcal R)\bigr)$ et comme $U$
n'est pas vide, il résulte de l'hypothèse de récurrence que
$\widehat j^{,*}(\mathcal P)$ et $\widehat j^{,*}(\mathcal R)$ sont des
$\mathcal O_{\widehat T}$-Modules algébrisables~; en vertu de
(\hyperref[III.5.3.3-fr]{5.3.3}), $\mathcal P$ et $\mathcal R$ sont
algébrisables, et on peut alors appliquer (\hyperref[III.5.3.2-fr]{5.3.2}), qui
prouve finalement que $\mathcal F$ est algébrisable.

\hfill C.Q.F.D.
\end{env}

\subsection{Application : comparaison de morphismes de schémas usuels et de morphismes de schémas formels. Schémas formels algébrisables}
\label{subsection:III.5.4-fr}

\begin{theorem}[5.4.1]
\label{III.5.4.1-fr}
Soient $A$ un anneau adique noethérien, $\mathfrak J$ un idéal de définition de
$A$, $S=\operatorname{Spec}(A)$, $S'=\mathrm V(\mathfrak J)$. Soient
$u:X\to S$ un morphisme propre, $v:Y\to S$ un morphisme séparé de type fini, et
soient $\widehat S$, $\widehat X$, $\widehat Y$ les complétés de $S$, $X$, $Y$ le
long de $S'$, $u^{-1}(S')$, $v^{-1}(S')$ respectivement. Si, pour tout
$S$-morphisme $f:X\to Y$, $\widehat f:\widehat X\to\widehat Y$ est le
prolongement de $f$ aux complétés, l'application $f\mapsto\widehat f$ est une
bijection
\[
\operatorname{Hom}_S(X,Y)\xrightarrow{\ \sim\ }
\operatorname{Hom}_{\widehat S}(\widehat X,\widehat Y).
\]
\end{theorem}

\begin{proof}
Montrons d'abord que $f\mapsto\widehat f$ est injective. Supposons en effet que deux
$S$-morphismes $f$, $g$ de $X$ dans $Y$ soient tels que $\widehat f=\widehat g$.
On sait alors (I, 10.9.4) qu'il existe un voisinage ouvert $V$ de
$X'=u^{-1}(S')$ dans lequel $f$ et $g$ coïncident. Or, comme $u$ est une application
fermée, on a $V=X$ (\hyperref[III.5.1.3.1-fr]{5.1.3.1}), d'où $f=g$.

Prouvons maintenant que $f\mapsto\widehat f$ est surjective, et soit donc $h$ un
$\widehat S$-morphisme $\widehat X\to\widehat Y$. Soit $Z=X\times_S Y$, et
désignons par $p:Z\to X$ et $q:Z\to Y$ les projections canoniques~; $Z$ est de
type fini sur $S$ (I, 6.3.4), donc noethérien~; désignons par $\widehat Z$ son
complété le long de $Z'=p^{-1}(u^{-1}(S'))$~; on sait que $\widehat Z$ s'identifie
canoniquement à $\widehat X\times_{\widehat S}\widehat Y$, les projections
$\widehat Z\to\widehat X$ et $\widehat Z\to\widehat Y$ s'identifiant aux
prolongements $\widehat p$ et $\widehat q$ (I, 10.9.7). Comme $Y$ est séparé sur
$S$, $\widehat Y$ est séparé sur $\widehat S$ (I, 10.15.7), donc le morphisme
graphe
\[
\Gamma_h=(1_{\widehat X},h):\widehat X\to\widehat Z
\]
est une immersion fermée (I, 10.15.4). Soient $\mathfrak T$ le sous-préschéma
formel fermé de $\widehat Z$ associé à cette immersion, et
$j:\mathfrak T\to\widehat Z$ l'injection canonique, de sorte que
$\Gamma_h=j\circ w$, où $w:\widehat X\to\mathfrak T$ est un isomorphisme
(I, 10.14.3) dont l'isomorphisme réciproque est $\widehat p\circ j$~; en outre,
$\mathfrak T$ est évidemment propre sur $\widehat S$, puisque $\widehat X$ l'est~;
on en conclut (\hyperref[III.5.1.8-fr]{5.1.8}) qu'il existe un sous-préschéma
fermé $T$ de $Z$ tel que
$\mathfrak T=\widehat T=T_{/(T\cap Z')}$, et que $j=\widehat i$, où $i$ est
l'injection canonique $T\to Z$ (I, 10.14.7). Alors $p\circ i:T\to X$ est un
isomorphisme, car il en est ainsi de
$\widehat{p\circ i}=\widehat p\circ\widehat i$ par hypothèse, et il suffit
d'appliquer

\oldpage[III]{157}

(4.6.8) en notant comme ci-dessus que $S$ est le seul voisinage de $S'$ dans $S$.
Soit $g:X\to T$ l'isomorphisme réciproque de $p\circ i$, et posons
$f=q\circ i\circ g$, qui est un morphisme $X\to Y$ dont par définition le graphe
$\Gamma_f=i\circ g$. Comme $\widehat g$ est l'isomorphisme réciproque de
$\widehat{p\circ i}=w$, on a
\[
\widehat{\Gamma_f}=\widehat i\circ\widehat g=j\circ w=\Gamma_h.
\]
Mais on sait que $\widehat{\Gamma_f}=\Gamma_{\widehat f}$ (I, 10.9.8), d'où
finalement $h=\widehat f$, ce qui achève la démonstration.
\end{proof}

On peut donc dire, dans le langage des catégories, que le foncteur
$X\mapsto\widehat X$ est pleinement fidèle ($0$, 8.1.6) de la catégorie des
schémas propres sur $\operatorname{Spec}(A)$ dans la catégorie des schémas formels
propres sur $\operatorname{Spf}(A)$, pour tout anneau adique noethérien $A$~; il
établit par suite une équivalence entre la première de ces catégories et une
sous-catégorie de la seconde~; les objets de cette dernière seront appelés
\emph{schémas formels algébrisables}. Pour un tel schéma $\mathfrak X$, il existe
un schéma usuel $X$, propre sur $\operatorname{Spec}(A)$, déterminé à isomorphisme
unique près, tel que $\mathfrak X$ soit isomorphe à $\widehat X$.

\begin{env}[5.4.2]
\label{III.5.4.2-fr}
\emph{Scholie.}
Avec les notations de (\hyperref[III.5.4.1-fr]{5.4.1}), posons
$S_n=\operatorname{Spec}(A/\mathfrak J^{n+1})$,
$X_n=X\times_S S_n$, $Y_n=Y\times_S S_n$. Il résulte de
(\hyperref[III.5.4.1-fr]{5.4.1}) et de (I, 10.12.3) que se donner un
$S$-morphisme $f:X\to Y$ équivaut à se donner un $(S_n)$-système inductif adique
(I, 10.12.2) de $S_n$-morphismes $f_n:X_n\to Y_n$.
\end{env}

\begin{remark}[5.4.3]
\label{III.5.4.3-fr}
Contrairement à ce que pourrait suggérer le th. d'existence
(\hyperref[III.5.1.6-fr]{5.1.6}), il y a des schémas formels propres sur
$\operatorname{Spf}(A)$ et qui ne sont pas algébrisables (tout comme il y a des
espaces analytiques compacts qui ne proviennent pas de variétés algébriques
complexes). Nous rencontrerons plus tard de tels schémas dans la «~théorie des
modules~», qui traite précisément (lorsque le corps de base est $\mathbb C$) des
variations infinitésimales de la structure complexe d'une variété algébrique
complète, et on sait que de telles variations peuvent donner naissance à des
variétés analytiques qui ne sont pas algébriques.
\end{remark}

\begin{proposition}[5.4.4]
\label{III.5.4.4-fr}
Soient $A$ un anneau adique noethérien, $\mathfrak S=\operatorname{Spf}(A)$,
$g:\mathfrak X\to\mathfrak S$, $h:\mathfrak Y\to\mathfrak S$ deux morphismes
propres de schémas formels, $f:\mathfrak X\to\mathfrak Y$ un
$\mathfrak S$-morphisme. Si $f$ est fini et si $\mathfrak Y$ est algébrisable,
alors $\mathfrak X$ est algébrisable.
\end{proposition}

\begin{proof}
On notera que les hypothèses sur $g$ et $h$ entraînent déjà que $f$ est propre
(3.4.1), et pour que $f$ soit fini, il suffit que pour tout $y\in\mathfrak Y$, la
fibre $f^{-1}(y)$ soit finie (\hyperref[III.4.8.11-fr]{4.8.11}). L'hypothèse
entraîne que $\mathcal B=f_*(\mathcal O_{\mathfrak X})$ est une
$\mathcal O_{\mathfrak Y}$-Algèbre cohérente
(\hyperref[III.4.8.6-fr]{4.8.6}), donc il résulte du th. d'existence que, si
$\mathfrak Y=\widehat Y$ et $h=\widehat w$, où
$w:Y\to\operatorname{Spec}(A)$ est un morphisme propre de schémas usuels, il
existe une $\mathcal O_Y$-Algèbre cohérente $\mathcal C$ telle que
$\mathcal B=\widehat{\mathcal C}$. Soient $X=\operatorname{Spec}(\mathcal C)$, et
$u:X\to Y$ le morphisme structural~; alors, il résulte aussitôt de la définition de
$\mathfrak X$ à partir de $\mathcal B$
(\hyperref[III.4.8.7-fr]{4.8.7}) que $\mathfrak X$ est canoniquement isomorphe à
$\widehat X$ et que $f$ s'identifie à $\widehat u$ (il suffit pour le voir de
considérer le cas où $Y$ est affine).

On notera que (\hyperref[III.5.1.8-fr]{5.1.8}) est un cas particulier de
(\hyperref[III.5.4.4-fr]{5.4.4}).
\end{proof}

\begin{theorem}[5.4.5]
\label{III.5.4.5-fr}
Soient $A$ un anneau adique noethérien, $\mathfrak J$ un idéal de définition de
$A$, $S=\operatorname{Spec}(A)$,
$\mathfrak S=\widehat S=\operatorname{Spf}(A)$,
$f:\mathfrak X\to\mathfrak S$ un morphisme propre de schémas formels. On pose
$S_k=\operatorname{Spec}(A/\mathfrak J^{k+1})$,
$X_k=\mathfrak X\times_{\mathfrak S}S_k$, et pour tout
$\mathcal O_{\mathfrak X}$-Module $\mathcal F$,
$\mathcal F_k=\mathcal F\otimes_{\mathcal O_{\mathfrak X}}\mathcal O_{X_k}
=\mathcal F/\mathfrak J^{k+1}\mathcal F$.

\oldpage[III]{158}

Soit $\mathcal L$ un $\mathcal O_{\mathfrak X}$-Module inversible, et supposons
que $\mathcal L_0=\mathcal L/\mathfrak J\mathcal L$ soit un
$\mathcal O_{X_0}$-Module ample. Alors $\mathfrak X$ est algébrisable, et si
$X$ est un $S$-schéma propre tel que $\mathfrak X$ soit isomorphe à
$\widehat X$, il existe un $\mathcal O_X$-Module ample $\mathcal M$ tel que
$\mathcal L$ soit isomorphe à $\widehat{\mathcal M}$ (ce qui entraîne que
$X$ est projectif sur $S$).
\end{theorem}

\begin{proof}
Appliquons (\hyperref[III.5.2.3-fr]{5.2.3}) à
$\mathcal F=\mathcal O_{\mathfrak X}$~: il existe donc un entier $n_0$ tel que
pour $n\geq n_0$, l'homomorphisme canonique
$\Gamma(\mathfrak X,\mathcal L^{\otimes n})\to
\Gamma(X_0,\mathcal L_0^{\otimes n})$ soit surjectif. On peut supposer
$n\geq n_0$ choisi assez grand pour que $\mathcal L_0^{\otimes n}$ soit très
ample pour $S_0$ (II, 4.5.10). Comme le morphisme $f_0:X_0\to S_0$ est propre,
$\Gamma(X_0,\mathcal L_0^{\otimes n})$ est un $A$-module de type fini (3.2.1),
donc il existe un sous-$A$-module de type fini $E$ de
$\Gamma(\mathfrak X,\mathcal L^{\otimes n})$ dont l'image dans
$\Gamma(X_0,\mathcal L_0^{\otimes n})$ soit ce dernier module tout entier.
Cela étant, pour tout $k\geq0$, considérons l'homomorphisme
\[
u_k:E/\mathfrak J^{k+1}E\longrightarrow
\Gamma(X_k,\mathcal L_k^{\otimes n})
\]
déduit de l'injection canonique
$E\to\Gamma(\mathfrak X,\mathcal L^{\otimes n})$. Notons que
$(f_k)_*(\mathcal L_k^{\otimes n})$ est quasi-cohérent, et comme
$\Gamma(S_k,(f_k)_*(\mathcal L_k^{\otimes n}))
=\Gamma(X_k,\mathcal L_k^{\otimes n})$, $u_k$ définit un homomorphisme
\[
\widetilde u_k:(E/\mathfrak J^{k+1}E)^\sim
\longrightarrow(f_k)_*(\mathcal L_k^{\otimes n}),
\]
et par suite aussi un homomorphisme
\[
\widetilde u_k^\sharp:
f_k^*((E/\mathfrak J^{k+1}E)^\sim)\longrightarrow
\mathcal L_k^{\otimes n}.
\]
D'ailleurs, si on pose
$\mathcal G_k=f_k^*((E/\mathfrak J^{k+1}E)^\sim)$, on a
$\mathcal G_0=\mathcal G_k/\mathfrak J\mathcal G_k$ (I, 9.1.5), donc
$\widetilde u_0^\sharp:\mathcal G_k/\mathfrak J\mathcal G_k
\to\mathcal L_k^{\otimes n}/\mathfrak J\mathcal L_k^{\otimes n}$ se déduit
de $\widetilde u_k^\sharp$ par passage aux quotients. Or, par définition de
$E$, $\widetilde u_0^\sharp$ n'est autre que l'homomorphisme canonique
\[
\sigma:f_0^*((f_0)_*(\mathcal L_0^{\otimes n}))\longrightarrow
\mathcal L_0^{\otimes n},
\]
et l'hypothèse que $\mathcal L_0^{\otimes n}$ est très ample entraîne que
$\widetilde u_0^\sharp$ est surjectif (II, 4.4.3)~; on déduit alors du lemme
de Nakayama que chacun des $\widetilde u_k^\sharp$ est aussi surjectif.
Chacun des $\widetilde u_k^\sharp$ définit donc (II, 4.2.2) un
$S_k$-morphisme
\[
g_k:X_k\longrightarrow P_k=\mathbf P(E/\mathfrak J^{k+1}E),
\]
et comme $P_h=P_k\times_{S_k}S_h$ pour $h\leq k$ en vertu de (II, 4.1.3),
$(g_k)$ est un $(S_k)$-système inductif adique (I, 10.12.2) en vertu des
relations entre les $\widetilde u_k^\sharp$ et de (II, 4.2.10). Les $g_k$
définissent donc un $\mathfrak S$-morphisme de schémas formels
$g:\mathfrak X\to\mathfrak P$, où $\mathfrak P$ est la limite inductive du
système $(P_k)$, ou encore le complété $\widehat P$, où $P=\mathbf P(E)$. De
plus, l'hypothèse que $\mathcal L_0^{\otimes n}$ est très ample entraîne que
$g_0$ est une immersion fermée (II, 4.4.3)~; on en conclut que $g$ est une
immersion fermée de schémas formels (4.8.10), donc $\mathfrak X$ est
algébrisable (5.1.8). Le fait que $\mathcal L$ soit isomorphe au complété
$\widehat{\mathcal M}$ d'un $\mathcal O_X$-Module inversible résulte alors du
th. d'existence (5.1.6). En outre, $\mathcal L^{\otimes n}$ est alors le
complété de $\mathcal M^{\otimes n}$ (I, 10.8.10), et les homomorphismes
$\widetilde u_k^\sharp$ définissent un homomorphisme bien déterminé
\[
v:f^*(\widetilde E)\longrightarrow\mathcal M^{\otimes n}
\]
(5.1.7)~; d'ailleurs, comme $\widetilde u_k^\sharp$ est surjectif, il en est de
même de $\widehat v$ (I, 10.11.5), donc de $v$ (5.1.3)~; en outre, le morphisme
$r:X\to P$ défini par $v$ (II, 4.2.2) a pour prolongement aux complétés $g$,
et comme $g$ est une immersion fermée, il en est de même de $r$, par (5.1.8)
et (5.4.1)~; on en conclut que $\mathcal M^{\otimes n}$ est très ample
(II, 4.4.6) et $\mathcal M$ est ample (II, 4.5.10).
\end{proof}

\begin{remark}[5.4.6]
\label{III.5.4.6-fr}
Soient $A$ un anneau adique noethérien,
$\mathfrak S=\operatorname{Spf}(A)$,
$f:\mathfrak X\to\mathfrak S$ un morphisme propre de schémas formels. Soit
$\mathcal N$ l'Idéal cohérent de $\mathcal O_{\mathfrak X}$ tel que pour tout
ouvert formel affine $U$ de $\mathfrak X$, $\Gamma(U,\mathcal N)$ soit le
nilradical de $\Gamma(U,\mathcal O_{\mathfrak X})$~; l'existence de cet Idéal
résulte facilement de (I, 10.10.2) et du fait que tout homomorphisme d'anneaux
$B\to C$ applique le nilradical de $B$ dans celui de $C$. Soit
$\mathfrak X'$ le sous-schéma formel fermé de $\mathfrak X$ défini par
$\mathcal N$ (I, 10.14.2)~; il serait intéressant de savoir si, lorsque
$\mathfrak X'$ est algébrisable, $\mathfrak X$ lui-même est algébrisable. On
parviendrait sans doute à une solution de ce problème si on savait classifier
(par exemple au moyen d'invariants de nature

\oldpage[III]{159}

cohomologique), les \emph{extensions} d'un faisceau structural
$\mathcal O_{\mathfrak X}$ (pour un préschéma usuel ou un préschéma formel) par
un Idéal de carré nul, autrement dit les $\mathcal O_{\mathfrak X}$-Algèbres
$\mathcal A$ telles que $\mathcal O_{\mathfrak X}$ soit isomorphe à
$\mathcal A/\mathcal J$, où $\mathcal J$ est un Idéal de carré nul de
$\mathcal A$.
\end{remark}

\subsection{Une décomposition de certains schémas}
\label{subsection:III.5.5-fr}

\begin{proposition}[5.5.1]
\label{III.5.5.1-fr}
Soient $A$ un anneau adique noethérien, $\mathfrak J$ un idéal de définition
de $A$, $Y=\operatorname{Spec}(A)$. Soit $f:X\to Y$ un morphisme séparé de
type fini~; on pose $Y_0=\operatorname{Spec}(A/\mathfrak J)$,
$X_0=X\times_Y Y_0=f^{-1}(Y_0)$. Soit $Z_0$ une partie ouverte de $X_0$,
propre sur $Y_0$~; il existe alors dans $X$ une partie ouverte et fermée $Z$,
propre sur $Y$ et telle que $Z\cap X_0=Z_0$.
\end{proposition}

\begin{proof}
Par hypothèse, il y a un ouvert $T$ de $X$ tel que $T\cap X_0=Z_0$~; soit
$\widehat T$ le complété le long de $Z_0$ du schéma induit par $X$ sur
l'ouvert $T$~; le support de $\mathcal O_{\widehat T}$ étant $Z_0$, qui est
propre sur $Y_0$, $\widehat T$ est propre sur
$\widehat Y=\operatorname{Spf}(A)$ (3.4.1). Il résulte de (5.1.8) qu'il existe
un sous-schéma fermé $Z$ de $T$ propre sur $Y$, tel que, si $i:Z\to T$ est
l'injection canonique, $\widehat i:\widehat Z\to\widehat T$ soit un
isomorphisme ($\widehat Z$ étant le complété de $Z$ le long de $Z_0$). On en
conclut (4.6.8) qu'il existe dans $T$ un voisinage ouvert $V$ de $Z_0$ tel que
la restriction $i^{-1}(V)\to V$ de $i$ soit un isomorphisme. Mais
$i^{-1}(V)$ est un voisinage de $Z_0$ dans $Z$, donc est nécessairement
identique à $Z$ (5.1.3.1). On en conclut que $Z$ est ouvert dans $T$, donc
dans $X$, ce qui achève la démonstration.
\end{proof}

\begin{corollary}[5.5.2]
\label{III.5.5.2-fr}
Si $X_0$ est propre sur $Y_0$, $X$ est réunion de deux parties ouvertes
disjointes $Z$ et $Z'$ telles que $Z$ soit propre sur $Y$ et contienne $X_0$~;
en outre, toute partie fermée $P$ de $X$, propre sur $Y$, est contenue dans
$Z$.
\end{corollary}

\begin{proof}
La dernière assertion résulte de ce que $P\cap Z'$ étant fermé dans $P$ est
propre sur $Y$~; si $P\cap Z'$ n'était pas vide, $f(P\cap Z')$ serait fermé
non vide dans $Y$, donc rencontrerait $Y_0$ (5.1.3.1), ce qui contredit la
définition de $Z$.
\end{proof}

\begin{flushright}
\emph{(A suivre.)}
\end{flushright}

% La page imprimée 504 est blanche.
\clearpage
\thispagestyle{empty}
\null
\clearpage
\thispagestyle{plain}

\oldpage[III]{161}
\section*{BIBLIOGRAPHIE (suite)}

\begin{flushleft}
[27] P. Gabriel, \emph{Des catégories abéliennes}, thèse, Paris (1961).

\medskip
[28] J. E. Roos, \emph{Sur les foncteurs dérivés de $\varprojlim$.
Applications}, C. R. Acad. Sci. Paris, t. CCLII, 1961, p. 3702--3704.

\medskip
[29] H. Cartan, \emph{Séminaire de l'École Normale Supérieure},
13\textsuperscript{e} année (1960--1961), exposé n\textsuperscript{o} 11.
\end{flushleft}

\clearpage
\oldpage[III]{162}
\section*{INDEX DES NOTATIONS}

\begin{flushleft}
$\mathbf{Ens}$ : 0, 8.1.1.

$h_X(Y)$, $h_X(u)$ ($X$, $Y$ objets d'une catégorie $C$, $u$ morphisme de
$C$) : 0, 8.1.1.

$h_w(u)$ ($u$, $w$ morphismes de $C$) : 0, 8.1.2.

$Z_k(E_r^{pq})$, $B_k(E_r^{pq})$ ($k\geq r$) : 0, 11.1.1.

$Z_\infty(E_2^{pq})$, $B_\infty(E_2^{pq})$ : 0, 11.1.1.

$K^\bullet$, $K_\bullet$ (complexes) : 0, 11.2.1.

$K^{\bullet\bullet}$, $K_{\bullet\bullet}$ (bicomplexes) : 0, 11.2.1.

$B^i(K^\bullet)$, $Z^i(K^\bullet)$, $H^i(K^\bullet)$,
$H^\bullet(K^\bullet)$ : 0, 11.2.1.

$B_i(K_\bullet)$, $Z_i(K_\bullet)$, $H_i(K_\bullet)$,
$H_\bullet(K_\bullet)$ : 0, 11.2.1.

$T(K^\bullet)$, $T(K_\bullet)$ ($T$ foncteur) : 0, 11.2.1.

$E(K^\bullet)$ ($K^\bullet$ complexe filtré) : 0, 11.2.2.

$K^{i,\bullet}$, $K^{\bullet,j}$, $Z_{\mathrm{II}}^p(K^{i,\bullet})$,
$B_{\mathrm{II}}^p(K^{i,\bullet})$, $H_{\mathrm{II}}^p(K^{i,\bullet})$,
$Z_{\mathrm I}^p(K^{\bullet,j})$, $B_{\mathrm I}^p(K^{\bullet,j})$,
$H_{\mathrm I}^p(K^{\bullet,j})$ : 0, 11.3.1.

$H_{\mathrm{II}}^p(K^{\bullet\bullet})$,
$Z_{\mathrm I}^q(H_{\mathrm{II}}^p(K^{\bullet\bullet}))$,
$B_{\mathrm I}^q(H_{\mathrm{II}}^p(K^{\bullet\bullet}))$,
$H_{\mathrm I}^q(H_{\mathrm{II}}^p(K^{\bullet\bullet}))$,
$H^n(K^{\bullet\bullet})$ : 0, 11.3.1.

$F_{\mathrm I}^p(K^{\bullet\bullet})$,
$F_{\mathrm{II}}^p(K^{\bullet\bullet})$, ${}'E(K^{\bullet\bullet})$,
${}''E(K^{\bullet\bullet})$ : 0, 11.3.2.

$K_{i,\bullet}$, $K_{\bullet,j}$, $H_p^{\mathrm{II}}(K_{i,\bullet})$,
$H_p^{\mathrm I}(K_{\bullet,j})$, $H_p^{\mathrm{II}}(K_{\bullet\bullet})$,
$H_p^{\mathrm I}(K_{\bullet\bullet})$,
$H_q^{\mathrm I}(H_p^{\mathrm{II}}(K_{\bullet\bullet}))$,
$H_q^{\mathrm{II}}(H_p^{\mathrm I}(K_{\bullet\bullet}))$,
$H_n(K_{\bullet\bullet})$ : 0, 11.3.5.

$\mathrm R^\bullet T(K^\bullet)$ ($T$ foncteur covariant additif) : 0, 11.4.3.

$\mathrm R^\bullet T(K^\bullet,K'^{\bullet})$ ($T$ bifoncteur covariant
additif) : 0, 11.4.6.

$\mathrm L_\bullet T(K_\bullet)$ ($T$ foncteur covariant additif) : 0, 11.6.2.

$\mathrm L_\bullet T(K_\bullet,K'_\bullet)$ ($T$ bifoncteur covariant
additif) : 0, 11.6.6.

$\mathrm L_\bullet T(K_{\bullet\bullet})$ ($T$ foncteur covariant additif) :
0, 11.7.2.

$|\sigma|$, $\Sigma(A)$, $\mathrm C_\bullet(A)$, $\mathrm L_\bullet(A)$ ($A$ ensemble fini,
$\sigma$ simplexe de $A$) : 0, 11.8.1.

$\mathrm C^\bullet(A,B;\mathscr S)$, $\mathrm L^\bullet(A,B;\mathscr S)$ ($A$, $B$ ensembles
finis, $\mathscr S$ système de coefficients) : 0, 11.8.4.

$\mathrm P^\bullet(A,B,\mathscr S)$ ($A$, $B$ ensembles finis, $\mathscr S$ système
de coefficients) : 0, 11.8.6.

$\mathrm C^\bullet(A,B;\mathscr S^\bullet)$,
$\mathrm L^\bullet(A,B;\mathscr S^\bullet)$ ($A$, $B$ ensembles finis,
$\mathscr S^\bullet$ complexe de systèmes de coefficients) : 0, 11.8.8.

$\mathrm P^\bullet(A,B;\mathscr S^\bullet)$ ($A$, $B$ ensembles finis,
$\mathscr S^\bullet$ complexe de systèmes de coefficients) : 0, 11.8.10.

$\chi(M^\bullet)$ ($M^\bullet$ complexe fini de modules de longueur finie) :
0, 11.10.1.

$\check C^\bullet(\mathfrak U,\mathcal F)$ ($\mathcal F$ faisceau de groupes abéliens
sur $X$, $\mathfrak U$ recouvrement ouvert de $X$) : 0, 12.1.2.

$\mathrm R^pf_*(\mathcal F)$ ($f:X\to Y$ morphisme d'espaces annelés,
$\mathcal F$ faisceau de $\mathcal O_X$-Modules) : 0, 12.2.1.

$\mathcal H^p(f,\mathcal K^\bullet)$,
$\mathcal H_f^p(\mathcal K^\bullet)$, ${}'E(f,\mathcal K^\bullet)$,
${}''E(f,\mathcal K^\bullet)$ ($f:X\to Y$ morphisme d'espaces
annelés, $\mathcal K^\bullet$ complexe de $\mathcal O_X$-Modules) : 0, 12.4.1.

$\mathbf H^\bullet(X,\mathcal K^\bullet)$,
$\mathbf H^\bullet(V,\mathcal K^\bullet)$ ($X$ espace annelé,
$\mathcal K^\bullet$ complexe de $\mathcal O_X$-Modules, $V$ ouvert de $X$) :
0, 12.4.2.

$\operatorname{gr}^\bullet(A)$ ($A$ système projectif vérifiant (ML)) :
0, 13.4.2.

$E(A)$ ($A$ objet filtré) : 0, 13.6.4.

$K_\bullet(\mathbf f)$, $i_{\mathbf f}$ ($\mathbf f$ famille finie d'éléments
d'un anneau) : III, 1.1.1.

$K^\bullet(\mathbf f,M)$, $K_\bullet(\mathbf f,M)$ ($\mathbf f$ famille
finie d'éléments d'un anneau $A$, $M$ $A$-module) : III, 1.1.2.

$H^\bullet(\mathbf f,M)$, $H_\bullet(\mathbf f,M)$ ($\mathbf f$ famille
finie d'éléments d'un anneau $A$, $M$ $A$-module) : III, 1.1.3.

$C^\bullet((\mathbf f),M)$, $H^\bullet((\mathbf f),M)$ ($\mathbf f$ famille
finie d'éléments d'un anneau $A$, $M$ $A$-module) : III, 1.1.6.

$H^\bullet(U,\mathcal F^{(*)})$,
$H^\bullet(\mathfrak U,\mathcal F^{(*)})$ ($\mathcal F$
$\mathcal O_X$-Module quasi-cohérent, $U$ ouvert de $X$, $\mathfrak U$
recouvrement ouvert de $X$) : III, 2.1.1.
\end{flushleft}

\clearpage
\oldpage[III]{163}
\section*{INDEX TERMINOLOGIQUE}

\begin{flushleft}
Aboutissement d'une suite spectrale : 0, 11.1.1.

Algébrisable ($\mathcal O_X$-Module) : III, 5.2.1.

Algébrisable (schéma formel) : III, 5.4.2.

Analytiquement intègre (anneau) : III, 4.3.6.

Application quasi-compacte : 0, 9.1.1.

Augmentation d'une résolution : 0, 11.4.1.

Caractéristique d'Euler-Poincaré d'un complexe : 0, 11.10.1.

Caractéristique d'Euler-Poincaré d'un $\mathcal O_X$-Module cohérent ($X$
projectif sur $\operatorname{Spec}(A)$, $A$ anneau artinien) : III, 2.5.1.

Cochaîne bi-alternée : 0, 11.8.4.

Complexe défini par un bicomplexe : 0, 11.3.1.

Complexe de l'algèbre extérieure : III, 1.1.1.

Condition (ML) : 0, 13.1.2.

Constructible (partie, ensemble) : 0, 9.1.2.

Constructible (fonction) : 0, 9.3.1.

Cup-produit : 0, 12.1.2.

Dihomomorphisme pour les structures algébriques sur les catégories : 0, 8.2.1.

Exact (sous-ensemble) dans une catégorie abélienne : III, 3.1.1.

Essentiellement constant (système projectif) : 0, 13.4.2.

Factorisation de Stein : III, 4.3.3.

Filtration co-discrète, co-séparée, discrète, exhaustive, finie, séparée :
0, 11.1.3.

Final (objet) d'une catégorie : 0, 8.1.10.

Fini (morphisme) de préschémas formels : III, 4.8.2.

Fini (préschéma formel) au-dessus de $\mathfrak Y$,
$\mathfrak Y$-fini (préschéma formel) : III, 4.8.2.

Foncteur covariant canonique $C\to\operatorname{Hom}(C^o,\mathbf{Ens})$ :
0, 8.1.2.

Genre arithmétique : III, 2.5.1.

Géométriquement connexe (fibre) : III, 4.3.4.

Homomorphisme de structures algébriques sur les catégories : 0, 8.2.1.

Hypercohomologie d'un foncteur par rapport à un complexe $K^\bullet$ :
0, 11.4.3.

Hypercohomologie d'un bifoncteur par rapport à deux complexes
$K^\bullet$, $K'^{\bullet}$ : 0, 11.4.6.

Hyperhomologie d'un foncteur par rapport à un complexe $K_\bullet$ :
0, 11.6.2.

Hyperhomologie d'un bifoncteur par rapport à deux complexes $K_\bullet$,
$K'_\bullet$ : 0, 11.6.6.

Hyperhomologie d'un foncteur par rapport à un bicomplexe : 0, 11.7.2.

Localement constructible (ensemble) : 0, 9.1.11.

Loi de composition externe (interne) sur les objets d'une catégorie :
0, 8.2.1.

\oldpage[III]{164}

$S$-$C$-module filtré, $\operatorname{gr}^\bullet(S)$-$C$-module gradué,
$\operatorname{gr}^{\bullet\bullet}(S)$-$C$-module bigradué : 0, 13.6.6.

Morphisme de suites spectrales : 0, 11.1.2.

Nombre géométrique de composantes connexes d'une fibre : III, 4.3.4.

$C$-objet en groupes, $C$-groupe, $C$-anneau, $C$-module : 0, 8.2.3.

Objet des bords, objet des cobords, objet des cocycles, objet des cycles :
0, 11.2.1.

Objet des images universelles : 0, 13.1.1.

Objet gradué associé à un système projectif satisfaisant à (ML) : 0, 13.4.2.

Objet gradué limité inférieurement (supérieurement) : 0, 11.2.1.

Partie propre (sur $\mathfrak Y$) d'un préschéma formel : III, 3.4.1.

Pleine (sous-catégorie) : 0, 8.1.5.

Pleinement fidèle (foncteur) : 0, 8.1.5.

Polynôme de Hilbert : III, 2.5.3.

Propre (morphisme) de préschémas formels : III, 3.4.1.

Représentable (foncteur) : 0, 8.1.8.

Résolution cohomologique, droite, gauche, homologique, injective, libre,
plate, projective : 0, 11.4.1.

Résolution de Cartan-Eilenberg droite, injective : 0, 11.4.2.

Résolution de Cartan-Eilenberg gauche, projective : 0, 11.6.1.

Rétrocompact (ensemble) : 0, 9.1.1.

Strict (système projectif) : 0, 13.4.2.

Suite spectrale dans une catégorie abélienne : 0, 11.1.1.

Suite spectrale faiblement convergente, régulière, corégulière, birégulière :
0, 11.1.3.

Suite spectrale dégénérée : 0, 11.1.6.

Suite spectrale d'un foncteur relativement à un objet filtré : 0, 13.6.4.

Suites spectrales d'un bicomplexe : 0, 11.3.2 et 11.3.5.

Suites spectrales d'hypercohomologie d'un foncteur par rapport à un complexe :
0, 11.4.3.

Suites spectrales d'hyperhomologie d'un foncteur par rapport à un complexe :
0, 11.6.2.

Suites spectrales d'hyperhomologie d'un foncteur par rapport à un bicomplexe :
0, 11.7.2.

Système de coefficients : 0, 11.8.4.

Unibranche (anneau, point) : III, 4.3.6.

Universellement ouvert (morphisme) : III, 4.3.9.
\end{flushleft}

\clearpage
\oldpage[III]{165}
\section*{TABLE DES MATIÈRES}

\begin{flushleft}
\textsc{Chapitre 0.} --- \textbf{Préliminaires (suite)} \dotfill 5

\medskip
\S~8. Foncteurs représentables \dotfill 5

8.1. Foncteurs représentables \dotfill 5

8.2. Structures algébriques dans les catégories \dotfill 9

\medskip
\S~9. Ensembles constructibles \dotfill 12

9.1. Ensembles constructibles \dotfill 12

9.2. Ensembles constructibles dans les espaces noethériens \dotfill 14

9.3. Fonctions constructibles \dotfill 16

\medskip
\S~10. Compléments sur les modules plats \dotfill 17

10.1. Relations entre modules plats et modules libres \dotfill 17

10.2. Critères locaux de platitude \dotfill 18

10.3. Existence d'extensions plates d'anneaux locaux \dotfill 20

\medskip
\S~11. Compléments d'algèbre homologique \dotfill 23

11.1. Rappels sur les suites spectrales \dotfill 23

11.2. La suite spectrale d'un complexe filtré \dotfill 27

11.3. Les suites spectrales d'un bicomplexe \dotfill 29

11.4. Hypercohomologie d'un foncteur par rapport à un complexe $K^\bullet$
\dotfill 32

11.5. Passage à la limite inductive dans l'hypercohomologie \dotfill 35

11.6. Hyperhomologie d'un foncteur par rapport à un complexe $K_\bullet$
\dotfill 39

11.7. Hyperhomologie d'un foncteur par rapport à un bicomplexe
$K_{\bullet\bullet}$ \dotfill 41

11.8. Compléments sur la cohomologie des complexes simpliciaux \dotfill 43

11.9. Un lemme sur les complexes de type fini \dotfill 46

11.10. Caractéristique d'Euler-Poincaré d'un complexe de modules de longueur
finie \dotfill 48

\medskip
\S~12. Compléments sur la cohomologie des faisceaux \dotfill 49

12.1. Cohomologie des faisceaux de modules sur les espaces annelés \dotfill 49

12.2. Images directes supérieures \dotfill 57

12.3. Compléments sur les foncteurs Ext de faisceaux \dotfill 60

12.4. Hypercohomologie du foncteur image directe \dotfill 62

\oldpage[III]{166}

\S~13. Limites projectives en algèbre homologique \dotfill 64

13.1. La condition de Mittag-Leffler \dotfill 64

13.2. La condition de Mittag-Leffler pour les groupes abéliens \dotfill 65

13.3. Application~: cohomologie d'une limite projective de faisceaux \dotfill 68

13.4. Condition de Mittag-Leffler et objets gradués associés aux systèmes
projectifs \dotfill 69

13.5. Limites projectives de suites spectrales de complexes filtrés \dotfill 71

13.6. Suite spectrale d'un foncteur relative à un objet muni d'une filtration
finie \dotfill 73

13.7. Foncteurs dérivés d'une limite projective d'arguments \dotfill 75

\medskip
\textsc{Chapitre III.} --- \textbf{Étude cohomologique des faisceaux
cohérents} \dotfill 81

\medskip
\S~1. Cohomologie des schémas affines \dotfill 82

1.1. Rappels sur le complexe de l'algèbre extérieure \dotfill 82

1.2. Cohomologie de Čech d'un recouvrement ouvert \dotfill 85

1.3. Cohomologie d'un schéma affine \dotfill 88

1.4. Application à la cohomologie des préschémas quelconques \dotfill 89

\medskip
\S~2. Étude cohomologique des morphismes projectifs \dotfill 95

2.1. Calculs explicites de certains groupes de cohomologie \dotfill 95

2.2. Le théorème fondamental des morphismes projectifs \dotfill 100

2.3. Application aux faisceaux gradués d'algèbres et de modules \dotfill 102

2.4. Une généralisation du théorème fondamental \dotfill 107

2.5. Caractéristique d'Euler-Poincaré et polynôme de Hilbert \dotfill 109

2.6. Application~: critères d'amplitude \dotfill 111

\medskip
\S~3. Le théorème de finitude pour les morphismes propres \dotfill 115

3.1. Le lemme de dévissage \dotfill 115

3.2. Le théorème de finitude~: cas des schémas usuels \dotfill 116

3.3. Généralisation du théorème de finitude (schémas usuels) \dotfill 118

3.4. Le théorème de finitude~: cas des schémas formels \dotfill 119

\medskip
\S~4. Le théorème fondamental des morphismes propres. Applications \dotfill 122

4.1. Le théorème fondamental \dotfill 122

4.2. Cas particuliers et variantes \dotfill 129

4.3. Le théorème de connexion de Zariski \dotfill 130

4.4. Le «~main theorem~» de Zariski \dotfill 135

4.5. Complétés de modules d'homomorphismes \dotfill 138

4.6. Relations entre morphismes formels et morphismes usuels \dotfill 139

4.7. Un critère d'amplitude \dotfill 145

4.8. Morphismes finis de préschémas formels \dotfill 146

\oldpage[III]{167}

\S~5. Un théorème d'existence de faisceaux algébriques cohérents \dotfill 149

5.1. Énoncé du théorème \dotfill 149

5.2. Démonstration du théorème d'existence~: cas projectif et quasi-projectif
\dotfill 151

5.3. Démonstration du théorème d'existence~: cas général \dotfill 154

5.4. Application~: comparaison de morphismes de schémas usuels et de
morphismes de schémas formels. Schémas formels algébrisables \dotfill 156

5.5. Une décomposition de certains schémas \dotfill 159

\medskip
\textsc{Bibliographie (suite)} \dotfill 161

\textsc{Index des notations} \dotfill 162

\textsc{Index terminologique} \dotfill 163

\begin{flushright}
\emph{Reçu le 28 juin 1961.}
\end{flushright}
\end{flushleft}

% FIN SOURCE ega3/ega3-1-fr.tex


% EGA III, seconde partie (sections 6--7), admitted from the sealed
% EGA III, seconde partie.
\clearpage
\setcounter{section}{5}
\setcounter{subsection}{0}
\setcounter{equation}{0}
\setcounter{footnote}{0}
%
% DEBUT SOURCE ega3/ega3-6-fr.tex
\chapter*{CHAPITRE III (suite)\\[1.2ex]
\Large ÉTUDE COHOMOLOGIQUE DES FAISCEAUX COHÉRENTS}

\section{Foncteurs « Tor » locaux et globaux; formule de Künneth}
\label{section:III.6.fr}

\subsection{Introduction}
\label{subsection:III.6.1.fr}

\oldpage[III]{137}

\begin{env}[6.1.1]
Soient $f:X\to Y$ un morphisme de préschémas, $\mathscr F$ un
$\mathscr O_X$-Module quasi-cohérent. Dans l'étude des « images directes
supérieures » $\mathrm R^n f_*(\mathscr F)$, on est amené à considérer le
problème général suivant : étant donné un morphisme « changement de base »
$g:Y'\to Y$, on pose
\[
  X'=X_{(Y')}=X\times_Y Y',\qquad
  \mathscr F'=\mathscr F\otimes_Y\mathscr O_{Y'},\qquad
  f'=f_{(Y')}:X'\to Y',
\]
et l'on se propose d'avoir des renseignements sur les images directes
supérieures $\mathrm R^n f'_*(\mathscr F')$ (en supposant connus les
$\mathrm R^n f_*(\mathscr F)$). On voit aisément (cf. par exemple (7.7.2))
que l'on est ramené à étudier les variations de
$\mathrm R^n f_*(\mathscr F\otimes_{\mathscr O_Y}\mathscr G)$ pour un
$\mathscr O_Y$-Module quasi-cohérent variable $\mathscr G$, autrement dit le
foncteur
\[
  \mathscr G\longmapsto
  \mathrm R^n f_*(\mathscr F\otimes_{\mathscr O_Y}\mathscr G).
\]
Si $\mathscr F$ est plat sur $Y$, le foncteur
$\mathscr G\mapsto\mathscr F\otimes_{\mathscr O_Y}\mathscr G$ est exact
(0\textsubscript{I}, 6.7.4) et par suite le foncteur composé
$\mathscr G\mapsto\mathrm R^n f_*(\mathscr F\otimes_{\mathscr O_Y}\mathscr G)$
est encore un foncteur cohomologique. Mais il n'en est plus de même dans le
cas général; pour pouvoir appliquer les méthodes cohomologiques, on est
conduit à substituer à
$\mathscr G\mapsto\mathrm R^n f_*(\mathscr F\otimes_{\mathscr O_Y}\mathscr G)$
d'autres foncteurs, qui cette fois sont toujours des foncteurs
cohomologiques. Ces foncteurs, qui généralisent les foncteurs « Tor » de la
théorie des modules, sont définis aux n\textsuperscript{os} 6.3 à 6.7; il y a
d'ailleurs deux telles généralisations, l'une « locale » et l'autre
« globale », reliées par des suites spectrales qui seront discutées au
n\textsuperscript{o} 6.7; comme application de ces suites spectrales, on
obtient en particulier, sous certaines conditions, une « formule de Künneth »
exprimant
\[
  \mathrm R^n(f_1\times f_2)_*(\mathscr F_1\otimes_Y\mathscr F_2)
\]
à l'aide des images directes supérieures
$\mathrm R^p f_{1*}(\mathscr F_1)$ et $\mathrm R^q f_{2*}(\mathscr F_2)$.
D'autres suites spectrales (6.8) généralisent les suites spectrales
d'associativité du foncteur « Tor » de modules; enfin, le problème du
changement de base conduit lui aussi à des suites spectrales (6.9).
\end{env}

\begin{env}[6.1.2]
En outre, on constate (6.10) que les foncteurs cohomologiques
$\mathscr G\mapsto\mathscr T_\bullet(\mathscr G)$ ainsi définis sont
localement (sur $Y$) du type
\[
  \mathscr G\longmapsto
  \mathscr H_\bullet(\mathscr L_\bullet\otimes_Y\mathscr G),
\]
où $\mathscr L_\bullet$ est un complexe de $\mathscr O_Y$-modules localement
libres (défini à une homotopie près) et $\mathscr H_\bullet$ l'homologie. Il y
a alors intérêt à oublier la situation particulière qui a donné naissance à
$\mathscr T_\bullet$, et à étudier de

\oldpage[III]{138}

façon générale les foncteurs de la forme précédente
$\mathscr G\mapsto\mathscr H_\bullet(\mathscr L_\bullet\otimes_Y\mathscr G)$
(où l'on fait en outre le cas échéant des hypothèses de finitude appropriées
sur les $\mathscr L_i$ ou les $\mathscr H_i(\mathscr L_\bullet)$) : c'est ce
qui est l'objet du § 7, dont la lecture, pour l'essentiel, est indépendante du
§ 6. Les propriétés les plus importantes de ces foncteurs concernent les
propriétés d'exactitude d'une composante $\mathscr T_i$ de
$\mathscr T_\bullet$; on donnera divers critères permettant d'établir de
telles propriétés; comme application, on obtiendra des conditions permettant
d'affirmer (avec les notations de (6.1.1)) que le foncteur
$\mathscr G\mapsto\mathrm R^n f_*(\mathscr F\otimes_{\mathscr O_Y}\mathscr G)$
est exact (ce qu'on exprimera en disant que $\mathscr F$ est
cohomologiquement plat sur $Y$ en dimension $n$). Une autre propriété
importante pour les composantes $\mathscr T_i$ de
$\mathscr T_\bullet(\mathscr G)=
\mathscr H_\bullet(\mathscr L_\bullet\otimes_Y\mathscr G)$ est une propriété
de semi-continuité de la fonction
$y\mapsto\dim_{k(y)}(T_i(k(y)))$; lorsque $\mathscr T_i$ est exact, cette
propriété est remplacée par une propriété de continuité, la réciproque étant
d'ailleurs vraie d'après Grauert lorsque $Y$ est réduit (7.8.4).
\end{env}

\begin{env}[6.1.3]
Dans les §§ 6 et 7, nous avons systématiquement fait usage de
l'hypercohomologie, en prenant partout comme arguments des complexes de
faisceaux au lieu de faisceaux, bien que la nécessité de ce point de vue
n'apparaisse que dans des chapitres ultérieurs. Le formalisme cohomologique
développé à cette occasion deviendra d'ailleurs plus transparent dans le
chapitre de ce Traité qui sera consacré à la mise au point d'une algèbre des
foncteurs cohomologiques de faisceaux cohérents, incluant le formalisme de la
dualité. Mais cela demandera des développements qui sortent du cadre du
présent chapitre.
\end{env}

\begin{env}[6.1.4]
Pour abréger, étant donnés deux complexes $K^\bullet$, $K'^\bullet$ dans une
catégorie abélienne $C$, nous dirons qu'un morphisme de complexes
$f:K^\bullet\to K'^\bullet$ est un homotopisme s'il existe un morphisme
$g:K'^\bullet\to K^\bullet$ tel que les morphismes composés $f\circ g$ et
$g\circ f$ soient tous deux homotopes à l'identité (par abus de langage,
lorsqu'il existe un tel homotopisme, on dira aussi que $K^\bullet$ et
$K'^\bullet$ sont homotopes). Lorsqu'on peut définir l'hypercohomologie d'un
foncteur covariant additif $T$ de $C$ dans une catégorie abélienne $C'$, par
rapport à un complexe de $C$ (0, 11.4.3), il est immédiat qu'un homotopisme
$K^\bullet\to K'^\bullet$ de complexes de $C$ définit canoniquement un
isomorphisme
\[
  \mathrm R^\bullet T(K^\bullet)\xrightarrow{\sim}
  \mathrm R^\bullet T(K'^\bullet)
\]
pour l'hypercohomologie (loc. cit.).
\end{env}

\subsection{Hypercohomologie des complexes de modules sur un préschéma}
\label{subsection:III.6.2.fr}

\begin{env}[6.2.1]
Soient $X$ un préschéma,
$\mathscr K^\bullet=(\mathscr K^i)_{i\in\mathbf Z}$ un complexe de
$\mathscr O_X$-Modules dont l'opérateur de dérivation est de degré $+1$.
Rappelons que pour tout morphisme $f:X\to Y$ de préschémas, on a défini
(0, 12.4.1) les $\mathscr O_Y$-Modules d'hypercohomologie
$\mathscr H^n(f,\mathscr K^\bullet)$ (aussi notés
$\mathscr H_f^n(\mathscr K^\bullet)$ ou
$\mathrm R^n f_*(\mathscr K^\bullet)$) pour tout $n\in\mathbf Z$;
l'hypercohomologie $\mathscr H^\bullet(f,\mathscr K^\bullet)$ est
l'aboutissement des deux foncteurs spectraux
$'\mathscr E(f,\mathscr K^\bullet)$ et
$''\mathscr E(f,\mathscr K^\bullet)$, dont les termes $\mathscr E_2$ sont
donnés par
\begin{align}
  '\mathscr E_2^{pq}
    &=\mathscr H^p\bigl(\mathscr H^q(f,\mathscr K^\bullet)\bigr),
    \tag{6.2.1.1}\label{III.6.2.1.1.fr}\\
  ''\mathscr E_2^{pq}
    &=\mathscr H^p\bigl(f,\mathscr H^q(\mathscr K^\bullet)\bigr)
      =\mathrm R^p f_*\bigl(\mathscr H^q(\mathscr K^\bullet)\bigr).
    \tag{6.2.1.2}\label{III.6.2.1.2.fr}
\end{align}

\oldpage[III]{139}

où $\mathscr H^q(f,\mathscr K^\bullet)$ est le complexe dont le composant de
degré $i$ est
$\mathscr H^q(f,\mathscr K^i)=\mathrm R^q f_*(\mathscr K^i)$ (loc. cit.).
Rappelons aussi que lorsque $Y$ est réduit à un point, on note
l'hypercohomologie correspondante $\mathbf H^\bullet(X,\mathscr K^\bullet)$
(qui est formée de modules sur $\Gamma(X,\mathscr O_X)$ indépendants du
préschéma ponctuel $Y$ considéré); lorsque $Y=X$ et $f=1_X$, on a
$\mathscr H^n(f,\mathscr K^\bullet)=\mathscr H^n(\mathscr K^\bullet)$
(cohomologie du complexe $\mathscr K^\bullet$); lorsque
$\mathscr K^i=0$, sauf pour $i=i_0$, on a
\[
  \mathscr H^n(f,\mathscr K^\bullet)
    =\mathrm R^{n-i_0}f_*(\mathscr K^{i_0}).
\]
La suite spectrale $'\mathscr E(f,\mathscr K^\bullet)$ est toujours
régulière; les deux suites spectrales sont birégulières lorsque
$\mathscr K^\bullet$ est limité inférieurement (0, 12.4.1).

Tout homotopisme $h:\mathscr K^\bullet\to\mathscr K'^\bullet$ de complexes de
$\mathscr O_X$-Modules (6.1.4) donne un isomorphisme
$\mathscr H^\bullet(f,\mathscr K^\bullet)\xrightarrow{\sim}
\mathscr H^\bullet(f,\mathscr K'^\bullet)$ pour l'hypercohomologie. Il en est
de même lorsqu'on suppose seulement que
$\mathscr H^\bullet(h):\mathscr H^\bullet(\mathscr K^\bullet)\to
\mathscr H^\bullet(\mathscr K'^\bullet)$ est un isomorphisme et que
$\mathscr K^\bullet$ et $\mathscr K'^\bullet$ sont limités inférieurement,
comme il résulte aussitôt de (0, 11.1.5) appliqué à la suite spectrale
(6.2.1.2) et à l'analogue pour $\mathscr K'^\bullet$. Enfin, pour tout
recouvrement ouvert $\mathfrak U=(U_\alpha)$ de $X$, on a aussi défini
(0, 12.4.5) l'hypercohomologie
$\mathbf H^\bullet(\mathfrak U,\mathscr K^\bullet)$ comme la cohomologie du
bicomplexe $C^\bullet(\mathfrak U,\mathscr K^\bullet)$ (dont le composant
d'indices $(i,j)$ est par définition
$C^i(\mathfrak U,\mathscr K^j)$); les
$\mathbf H^n(\mathfrak U,\mathscr K^\bullet)$ sont encore des modules sur
$\Gamma(X,\mathscr O_X)$.
\end{env}

\begin{proposition}[6.2.2]
Soient $X$ un schéma, $\mathfrak U=(U_\alpha)$ un recouvrement de $X$ par des
ouverts affines. Pour tout complexe $\mathscr K^\bullet$ de
$\mathscr O_X$-Modules quasi-cohérents, les modules d'hypercohomologie
$\mathbf H^\bullet(X,\mathscr K^\bullet)$ et
$\mathbf H^\bullet(\mathfrak U,\mathscr K^\bullet)$ sont canoniquement
isomorphes.
\end{proposition}

\begin{proof}
En effet, toute intersection finie $V$ d'ouverts du recouvrement
$\mathfrak U$ est affine (I, 5.5.6), donc
$\mathrm H^q(V,\mathscr K^i)=0$ pour tout $i$ et tout $q>0$ (1.3.1); la
proposition est donc un cas particulier de (0, 12.4.7).
\end{proof}

\begin{proposition}[6.2.3]
Soit $f:X\to Y$ un morphisme quasi-compact et séparé de préschémas. Pour tout
complexe $\mathscr K^\bullet$ de $\mathscr O_X$-Modules quasi-cohérents, les
$\mathscr O_Y$-Modules $\mathscr H^n(f,\mathscr K^\bullet)$ sont
quasi-cohérents.
\end{proposition}

\begin{proof}
Comme les $\mathscr H^q(f,\mathscr K^i)=\mathrm R^q f_*(\mathscr K_i)$ sont
des $\mathscr O_Y$-Modules quasi-cohérents (1.4.10), il en est de même de
$'\mathscr E_2^{pq}$, qui, d'après (6.2.1.1), est quotient d'un noyau
d'homomorphisme de Modules quasi-cohérents par une image d'un tel
homomorphisme (I, 4.1.1). Pour la même raison, tous les
$\mathscr O_Y$-Modules $'\mathscr E_r^{pq}$,
$B_k('\mathscr E_2^{pq})$, $Z_k('\mathscr E_2^{pq})$ de la première suite
spectrale sont quasi-cohérents. La régularité de la suite spectrale
$'\mathscr E(f,\mathscr K^\bullet)$ entraîne que
$Z_\infty('\mathscr E_2^{pq})$ est égal à un des
$Z_k('\mathscr E_2^{pq})$, donc est quasi-cohérent, et il en est de même de
$B_\infty('\mathscr E_2^{pq})=\varinjlim_k B_k('\mathscr E_2^{pq})$
(0, 11.2.4 et I, 4.1.1); les $'\mathscr E_\infty^{pq}$ sont donc aussi
quasi-cohérents. La suite spectrale précédente étant régulière, la filtration
des $F^p(\mathscr H^n(f,\mathscr K^\bullet))$ est discrète et exhaustive;
autrement dit, le $\mathscr O_Y$-Module
$\mathscr H^n(f,\mathscr K^\bullet)$ est réunion d'une suite croissante
$(\mathscr G_k)_{k\geq0}$ de $\mathscr O_Y$-Modules telle que
$\mathscr G_0=0$ et que chaque $\mathscr G_k/\mathscr G_{k-1}$ soit égal à un
des $\mathscr O_Y$-Modules $'\mathscr E_\infty^{pq}$, donc soit
quasi-cohérent. Par récurrence sur $k$, on en déduit que les $\mathscr G_k$
sont quasi-cohérents (1.4.17), et comme
$\mathscr H^n(f,\mathscr K^\bullet)=\varinjlim_k\mathscr G_k$, la proposition
est démontrée (I, 4.1.1).
\end{proof}

\oldpage[III]{140}

\begin{corollary}[6.2.4]
Sous les hypothèses de (6.2.3), pour tout ouvert affine $V$ de $Y$,
l'homomorphisme canonique
\[
  \mathbf H^n(f^{-1}(V),\mathscr K^\bullet)
  \longrightarrow
  \Gamma(V,\mathscr H^n(f,\mathscr K^\bullet))
  \tag{6.2.4.1}\label{III.6.2.4.1.fr}
\]
est bijectif pour tout $n\in\mathbf Z$.
\end{corollary}

\begin{proof}
La démonstration est la même que celle de (1.4.11), en utilisant (6.2.2),
remplaçant $\mathscr F$ par $\mathscr K^\bullet$,
$\mathscr K^\bullet$ par
$f_*(C^\bullet(\mathfrak U,\mathscr K^\bullet))$,
$\mathscr H^\bullet(\mathscr K^\bullet)$ par
$\mathscr H^\bullet(f,\mathscr K^\bullet)$, et notant que ce dernier est un
$\mathscr O_Y$-Module quasi-cohérent par (6.2.3).
\end{proof}

\begin{proposition}[6.2.5]
Soient $Y$ un préschéma localement noethérien, $f:X\to Y$ un morphisme
propre, $\mathscr K^\bullet$ un complexe de $\mathscr O_X$-Modules tel que les
$\mathscr O_X$-Modules $\mathscr H^q(\mathscr K^\bullet)$ soient cohérents.
Alors les $\mathscr O_Y$-Modules $\mathscr H^n(f,\mathscr K^\bullet)$ sont
cohérents.
\end{proposition}

\begin{proof}
La question étant locale sur $Y$, on peut se borner au cas où $Y$ est
noethérien et affine, et il s'agit donc, en vertu de (6.2.4), de prouver que
les $\mathbf H^n(X,\mathscr K^\bullet)$ sont des
$\Gamma(Y,\mathscr O_Y)$-modules de type fini. On a alors
$\mathbf H^\bullet(X,\mathscr K^\bullet)=
\mathbf H^\bullet(\mathfrak U,\mathscr K^\bullet)$ (6.2.2), où l'on peut
supposer que $\mathfrak U$ est fini, puisque $X$ est quasi-compact. Les
cochaînes de chaque complexe $C^\bullet(\mathfrak U,\mathscr K^j)$ étant
alternées par définition, il y a un entier $r>0$ tel que
$C^i(\mathfrak U,\mathscr K^j)=0$ pour $i<0$ et $i>r$; on en conclut
(0, 11.3.3) que les deux suites spectrales du bicomplexe
$C^\bullet(\mathfrak U,\mathscr K^\bullet)$ sont birégulières. Comme les
intersections des ensembles de $\mathfrak U$ sont des ouverts affines
(I, 5.5.6), chaque foncteur
$\mathscr F\mapsto C^i(\mathfrak U,\mathscr F)$ est exact dans la catégorie
des $\mathscr O_X$-Modules quasi-cohérents; donc
$\mathrm H_I^q(C^i(\mathfrak U,\mathscr K^\bullet))
=C^i(\mathfrak U,\mathscr H^q(\mathscr K^\bullet))$, et les termes
$E_2$ de la seconde suite spectrale de
$C^\bullet(\mathfrak U,\mathscr K^\bullet)$ sont donnés (0, 11.3.2) par
\[
  ''E_2^{pq}
  =\mathrm H^p(C^\bullet(\mathfrak U,\mathscr H^q(\mathscr K^\bullet)))
  =\mathbf H^p(\mathfrak U,\mathscr H^q(\mathscr K^\bullet))
  =\mathrm H^p(X,\mathscr H^q(\mathscr K^\bullet))
\]
en vertu de (1.4.1); puisque $f$ est propre, ce sont des
$\Gamma(Y,\mathscr O_Y)$-modules de type fini (3.2.1). La suite spectrale
$''E(C^\bullet(\mathfrak U,\mathscr K^\bullet))$ étant birégulière, on en
déduit bien que les $\mathbf H^n(X,\mathscr K^\bullet)$ sont des
$\Gamma(Y,\mathscr O_Y)$-modules de type fini (0, 11.1.8).

A fortiori, si $\mathscr K^\bullet$ est un complexe de
$\mathscr O_X$-Modules cohérents, les $\mathscr O_Y$-Modules
$\mathscr H^n(f,\mathscr K^\bullet)$ sont cohérents sous les hypothèses de
(6.2.5) relatives à $Y$ et $f$ (0\textsubscript{I}, 5.3.4).
\end{proof}

\begin{env}[6.2.6]
L'hypercohomologie $\mathscr H^\bullet(f,\mathscr K^\bullet)$ est un foncteur
cohomologique dans la catégorie des complexes de $\mathscr O_X$-Modules
limités inférieurement (0, 12.4.4). C'est un foncteur cohomologique dans la
catégorie de tous les complexes de $\mathscr O_X$-Modules lorsque le
morphisme $f$ est quasi-compact et l'espace sous-jacent à $X$ localement
noethérien : en effet, il résulte alors de (G, II, 3.10.1) que $f_*$ permute
aux limites inductives (la question étant locale sur $Y$), et l'on peut
appliquer (0, 11.5.2).

Enfin, si $f$ est séparé, $\mathscr H^\bullet(f,\mathscr K^\bullet)$ est un
foncteur cohomologique dans la catégorie des complexes de
$\mathscr O_X$-Modules quasi-cohérents. C'est immédiat lorsque $Y$ est
affine, car alors $X$ est un schéma, donc, en vertu de l'isomorphisme
canonique (6.2.2), on est ramené à voir que
$\mathscr K^\bullet\mapsto\mathbf H^\bullet(\mathfrak U,\mathscr K^\bullet)$
est un foncteur cohomologique dans la catégorie des complexes de
$\mathscr O_X$-Modules quasi-cohérents, ce qui est immédiat puisque le
foncteur $\mathscr K^\bullet\mapsto
C^\bullet(\mathfrak U,\mathscr K^\bullet)$ est exact dans cette catégorie
(I, 1.3.7). Dans le cas général, pour tout ouvert affine $V$

\oldpage[III]{141}

de $Y$, $f^{-1}(Y)$ est un schéma, et pour appliquer ce qui précède, il
suffit de vérifier que pour une suite exacte
\[
  0\longrightarrow\mathscr K'^\bullet
  \longrightarrow\mathscr K^\bullet
  \longrightarrow\mathscr K''^\bullet
  \longrightarrow0
\]
de complexes de $\mathscr O_X$-Modules quasi-cohérents, l'homomorphisme
\[
  \partial:
  \mathscr H^n(f,\mathscr K''^\bullet)|V
  \longrightarrow
  \mathscr H^{n+1}(f,\mathscr K'^\bullet)|V
\]
ne dépend pas du recouvrement ouvert affine $\mathfrak U$ de $f^{-1}(V)$
utilisé pour le définir. Mais cela résulte de ce que, si $\mathfrak U'$ est
un recouvrement ouvert affine plus fin que $\mathfrak U$, le diagramme
\[
  \xymatrix@C=34pt@R=11pt{
    \mathbf H^\bullet(\mathfrak U,\mathscr K^\bullet)
      \ar[rrd]^-{\sim}\ar[dd]^{\iota}
    && \\
    &&
    \mathbf H^\bullet(f^{-1}(V),\mathscr K^\bullet)
    \\
    \mathbf H^\bullet(\mathfrak U',\mathscr K^\bullet)
      \ar[rru]_-{\sim}
    &&
  }
\]
d'isomorphismes canoniques est commutatif, ainsi que le diagramme
\[
  \xymatrix@C=46pt@R=25pt{
    \mathbf H^n(\mathfrak U,\mathscr K''^\bullet)
      \ar[r]^-{\partial}\ar[d]^{\iota}
    &
    \mathbf H^{n+1}(\mathfrak U,\mathscr K'^\bullet)
      \ar[d]^{\iota}
    \\
    \mathbf H^n(\mathfrak U',\mathscr K''^\bullet)
      \ar[r]_-{\partial}
    &
    \mathbf H^{n+1}(\mathfrak U',\mathscr K'^\bullet).
  }
\]

Lorsque l'une des conditions précédentes est remplie et que
$\mathscr K^\bullet\to\mathscr K'^\bullet$ est un homotopisme (6.1.4),
l'isomorphisme correspondant
$\mathscr H^\bullet(f,\mathscr K^\bullet)\xrightarrow{\sim}
\mathscr H^\bullet(f,\mathscr K'^\bullet)$ est alors un isomorphisme de
$\partial$-foncteurs (0, 11.4.4).
\end{env}

\begin{env}[6.2.7]
Tout ce qui précède s'applique naturellement sans changement (sinon de
notations) à un complexe $\mathscr K_\bullet$ de $\mathscr O_X$-Modules
quasi-cohérents dont l'opérateur de dérivation est de degré $-1$; il suffit
de considérer le complexe $\mathscr K^\bullet=(\mathscr K^i)$ où
$\mathscr K^i=\mathscr K_{-i}$ pour tout $i\in\mathbf Z$.
\end{env}

\subsection{Hypertor de deux complexes de modules}
\label{subsection:III.6.3.fr}

\begin{env}[6.3.1]
Soient $A$ un anneau commutatif, $P_\bullet$, $Q_\bullet$ deux complexes de
$A$-modules dont les opérateurs de dérivation sont de degré $-1$; soit
$L_{\bullet\bullet}$ (resp. $M_{\bullet\bullet}$) une résolution projective
de Cartan-Eilenberg de $P_\bullet$ (resp. $Q_\bullet$) (0, 11.6.1);
$L_{\bullet\bullet}\otimes_A M_{\bullet\bullet}$ est alors (pour la somme
des premiers degrés et la somme des seconds degrés) un bicomplexe (à
opérateurs de dérivation de degré $-1$), dont l'homologie
$H_\bullet(L_{\bullet\bullet}\otimes_A M_{\bullet\bullet})$ ne dépend pas des
résolutions de Cartan-Eilenberg $L_{\bullet\bullet}$,
$M_{\bullet\bullet}$ choisies, et est par définition l'hyperhomologie du
bifoncteur $P_\bullet\otimes_A Q_\bullet$ en $P_\bullet$ et $Q_\bullet$
(0, 11.6.5). Nous poserons par définition
\[
  \mathbf{Tor}_n^A(P_\bullet,Q_\bullet)
  =H_n(L_{\bullet\bullet}\otimes_A M_{\bullet\bullet})
  \tag{6.3.1.1}\label{III.6.3.1.1.fr}
\]

\oldpage[III]{142}

et nous dirons que cet $A$-module est l'hypertor d'indice $n$ des deux
complexes $P_\bullet$, $Q_\bullet$. On sait que dans la catégorie des
complexes de $A$-modules limités inférieurement, les
$\mathbf{Tor}_n^A(P_\bullet,Q_\bullet)$ forment un bifoncteur homologique en
$P_\bullet$, $Q_\bullet$ (0, 11.6.5). En outre :
\end{env}

\begin{proposition}[6.3.2]
Le bifoncteur $\mathbf{Tor}_\bullet^A(P_\bullet,Q_\bullet)$ est
l'aboutissement commun de deux bifoncteurs spectraux
$'E(P_\bullet,Q_\bullet)$, $''E(P_\bullet,Q_\bullet)$, dont les termes $E_2$
sont
\begin{align}
  'E^2_{pq}
    &=H_p\bigl(\operatorname{Tor}_q^A(P_\bullet,Q_\bullet)\bigr),
    \tag{6.3.2.1}\label{III.6.3.2.1.fr}\\
  ''E^2_{pq}
    &=\bigoplus_{q'+q''=q}
      \operatorname{Tor}_p^A\bigl(H_{q'}(P_\bullet),H_{q''}(Q_\bullet)\bigr).
    \tag{6.3.2.2}\label{III.6.3.2.2.fr}
\end{align}
où, dans (6.3.2.1), $\operatorname{Tor}_q^A(P_\bullet,Q_\bullet)$ désigne le
bicomplexe formé des $A$-modules $\operatorname{Tor}_q^A(P_i,Q_j)$. La suite
spectrale (6.3.2.2) est toujours régulière; si $P_\bullet$ et $Q_\bullet$
sont limités inférieurement, ou si $A$ est de dimension cohomologique finie,
les deux suites spectrales (6.3.2.1) et (6.3.2.2) sont birégulières.
\end{proposition}

\begin{proof}
Cela résulte de (0, 11.6.5), car lorsque $A$ est de dimension cohomologique
finie $n$, tout $A$-module admet une résolution projective de longueur $n$
(M, VI, 2.1).
\end{proof}

\begin{corollary}[6.3.3]
Soient $P'_\bullet$, $Q'_\bullet$ deux complexes de $A$-modules,
$u:P_\bullet\to P'_\bullet$, $v:Q_\bullet\to Q'_\bullet$ deux homomorphismes
de complexes. Si les homomorphismes
$H_\bullet(u):H_\bullet(P_\bullet)\to H_\bullet(P'_\bullet)$,
$H_\bullet(v):H_\bullet(Q_\bullet)\to H_\bullet(Q'_\bullet)$ déduits
respectivement de $u$ et $v$ sont bijectifs, alors l'homomorphisme
$\mathbf{Tor}_\bullet^A(P_\bullet,Q_\bullet)\to
\mathbf{Tor}_\bullet^A(P'_\bullet,Q'_\bullet)$ déduit de $u$ et $v$ est
bijectif.
\end{corollary}

\begin{proof}
En effet, l'homomorphisme de suites spectrales
$''E(P_\bullet,Q_\bullet)\to''E(P'_\bullet,Q'_\bullet)$ déduit de $u$ et $v$
est alors un isomorphisme pour les termes $E_2$ et la conclusion résulte de
ce que ces suites sont régulières en vertu de (6.3.2) (0, 11.1.5).
\end{proof}

\begin{proposition}[6.3.4]
Soient $P_\bullet$, $Q_\bullet$ deux complexes de $A$-modules, limités
inférieurement. Soit $L_{\bullet\bullet}$ (resp. $M_{\bullet\bullet}$) un
bicomplexe formé de $A$-modules plats, tel que pour tout $i$,
$L_{i,\bullet}$ (resp. $M_{i,\bullet}$) soit une résolution de $P_i$
(resp. $Q_i$). On a alors des isomorphismes canoniques
\[
  \mathbf{Tor}_\bullet^A(P_\bullet,Q_\bullet)
  \xrightarrow{\sim}H_\bullet(L_{\bullet\bullet}\otimes_AQ_\bullet)
  \xrightarrow{\sim}H_\bullet(P_\bullet\otimes_AM_{\bullet\bullet})
  \xrightarrow{\sim}H_\bullet(L_{\bullet\bullet}\otimes_AM_{\bullet\bullet})
  \tag{6.3.4.1}\label{III.6.3.4.1.fr}
\]
\end{proposition}

\begin{proof}
Cela résulte de (0, 11.6.5, (ii) et (iii)) et de la définition des
$A$-modules plats.
\end{proof}

\begin{namedtheorem}[6.3.5]{Remarques}
(i) Avec les notations de (6.3.1), les bicomplexes
$L_{\bullet\bullet}\otimes_A M_{\bullet\bullet}$ et
$M_{\bullet\bullet}\otimes_A L_{\bullet\bullet}$ sont canoniquement
isomorphes, d'où un isomorphisme canonique
$\mathbf{Tor}_\bullet^A(P_\bullet,Q_\bullet)\xrightarrow{\sim}
\mathbf{Tor}_\bullet^A(Q_\bullet,P_\bullet)$.

(ii) Si $F$ et $G$ sont deux $A$-modules, $P_\bullet$ et $Q_\bullet$ les
complexes de $A$-modules réduits à $F$ et $G$ respectivement en degré 0 et
nuls dans les autres degrés, alors deux résolutions projectives $L_\bullet$,
$M_\bullet$ de $F$ et $G$ respectivement peuvent être considérées comme des
résolutions de Cartan-Eilenberg de $P_\bullet$ et $Q_\bullet$ en les
complétant par des zéros. On a par suite dans ce cas
$\mathbf{Tor}_\bullet^A(P_\bullet,Q_\bullet)
=\operatorname{Tor}_\bullet^A(F,G)$.
\end{namedtheorem}

\begin{proposition}[6.3.6]
Soient $(P_\bullet^\lambda)$, $(Q_\bullet^\mu)$ deux systèmes inductifs
filtrants de complexes de $A$-modules; on a un isomorphisme canonique
\[
  \varinjlim_{\lambda,\mu}
  \mathbf{Tor}_\bullet^A(P_\bullet^\lambda,Q_\bullet^\mu)
  \xrightarrow{\sim}
  \mathbf{Tor}_\bullet^A
  \left(\varinjlim_\lambda P_\bullet^\lambda,
        \varinjlim_\mu Q_\bullet^\mu\right).
  \tag{6.3.6.1}\label{III.6.3.6.1.fr}
\]
\end{proposition}

\begin{proof}
Posons $P_\bullet=\varinjlim_\lambda P_\bullet^\lambda$,
$Q_\bullet=\varinjlim_\mu Q_\bullet^\mu$; par fonctorialité, il est clair que
les $\mathbf{Tor}_\bullet^A(P_\bullet^\lambda,Q_\bullet^\mu)$ forment un
système inductif et que les applications
$\mathbf{Tor}_\bullet^A(P_\bullet^\lambda,Q_\bullet^\mu)\to
\mathbf{Tor}_\bullet^A(P_\bullet,Q_\bullet)$ déduites

\oldpage[III]{143}

des applications canoniques $P_\bullet^\lambda\to P_\bullet$,
$Q_\bullet^\mu\to Q_\bullet$, forment un système inductif d'homomorphismes,
d'où un homomorphisme canonique (6.3.6.1), et plus généralement un
homomorphisme canonique
\[
  \varinjlim_{\lambda,\mu}''E(P_\bullet^\lambda,Q_\bullet^\mu)
  \longrightarrow ''E(P_\bullet,Q_\bullet)
\]
dont (6.3.6.1) est l'homomorphisme des aboutissements. En outre, la suite
spectrale $''E(P_\bullet,Q_\bullet)$ est régulière (6.3.2), et il en est de
même de la suite spectrale
$\varinjlim_{\lambda,\mu}''E(P_\bullet^\lambda,Q_\bullet^\mu)$, comme il
résulte des définitions (0, 11.1.7) et de la démonstration de (0, 11.3.3);
pour démontrer que (6.3.6.1) est bijectif, il suffit donc (0, 11.1.5) de
prouver que l'homomorphisme
\[
  \varinjlim_{\lambda,\mu}''E(P_\bullet^\lambda,Q_\bullet^\mu)
  \longrightarrow ''E(P_\bullet,Q_\bullet)
  \tag{6.3.6.2}\label{III.6.3.6.2.fr}
\]
est bijectif pour les termes $E^2$. Comme le foncteur $H_\bullet$ commute à
la limite inductive des complexes de modules, on est finalement ramené à
prouver que pour deux systèmes inductifs filtrants $(F^\lambda)$,
$(G^\mu)$ de $A$-modules, l'homomorphisme canonique
\[
  \varinjlim_{\lambda,\mu}\operatorname{Tor}_\bullet^A(F^\lambda,G^\mu)
  \longrightarrow
  \operatorname{Tor}_\bullet^A
  \left(\varinjlim_\lambda F^\lambda,
        \varinjlim_\mu G^\mu\right)
\]
est bijectif. Pour cela, considérons pour chaque $F^\lambda$ la résolution
libre canonique
\[
  L_\bullet^\lambda:\quad
  \cdots\longrightarrow L_{i+1}^\lambda
  \longrightarrow L_i^\lambda\longrightarrow\cdots
  \longrightarrow L_1^\lambda\longrightarrow L_0^\lambda\longrightarrow0
\]
où $L_0^\lambda$ est le $A$-module des combinaisons linéaires formelles
d'éléments de $F^\lambda$ et $L_{i+1}^\lambda$ le $A$-module des combinaisons
linéaires formelles d'éléments de
$\operatorname{Ker}(L_i^\lambda\to L_{i-1}^\lambda)$; on vérifie
immédiatement que les $L_\bullet^\lambda$ forment un système inductif de
complexes, et si l'on pose
$F=\varinjlim_\lambda F^\lambda$, $L_i=\varinjlim_\lambda L_i^\lambda$, les
$L_i$ forment une résolution $L_\bullet$ de $F$, le foncteur
$\varinjlim$ étant exact; en outre, les $L_i$, limites inductives de
$A$-modules libres, sont plats (0\textsubscript{I}, 6.1.2). On considère de
même pour chaque $\mu$ la résolution libre canonique $M_\bullet^\mu$ de
$G^\mu$, et $M_\bullet=\varinjlim_\mu M_\bullet^\mu$ est une résolution plate
de $G=\varinjlim_\mu G^\mu$. On a alors
$\operatorname{Tor}_\bullet^A(\varinjlim_\lambda F^\lambda,
\varinjlim_\mu G^\mu)=H_\bullet(L_\bullet\otimes_A M_\bullet)$ en vertu de
(6.3.5) et (6.3.4); mais
\[
  H_\bullet(L_\bullet\otimes_A M_\bullet)
  =\varinjlim_{\lambda,\mu}
   H_\bullet(L_\bullet^\lambda\otimes_A M_\bullet^\mu)
\]
puisque $H_\bullet$ commute aux limites inductives de complexes de modules;
comme $H_\bullet(L_\bullet^\lambda\otimes_A M_\bullet^\mu)
=\operatorname{Tor}_\bullet^A(F^\lambda,G^\mu)$, cela termine la
démonstration.

Lorsqu'on suppose qu'il existe $i_0$ tel que
$P_i^\lambda=Q_i^\mu=0$ pour $i<i_0$ quels que soient $\lambda$ et $\mu$, on
démontre de la même manière que l'homomorphisme canonique
\[
  \varinjlim_{\lambda,\mu}'E(P_\bullet^\lambda,Q_\bullet^\mu)
  \longrightarrow'E(P_\bullet,Q_\bullet)
  \tag{6.3.6.3}\label{III.6.3.6.3.fr}
\]
est bijectif.
\end{proof}

\begin{proposition}[6.3.7]
Supposons $P_\bullet$ et $Q_\bullet$ limités inférieurement. Si le complexe
$P_\bullet$ est formé de $A$-modules plats, on a un $A$-isomorphisme
canonique de $\partial$-foncteurs en $Q_\bullet$
\[
  \mathbf{Tor}_\bullet^A(P_\bullet,Q_\bullet)
  \xrightarrow{\sim}H_\bullet(P_\bullet\otimes_A Q_\bullet).
  \tag{6.3.7.1}\label{III.6.3.7.1.fr}
\]
\end{proposition}

\begin{proof}
En effet, la suite spectrale (6.3.2.1) est birégulière et dégénérée, et
l'existence de l'isomorphisme (6.3.7.1) résulte de (0, 11.1.6). En outre, en
calculant l'hypertor à partir d'une résolution projective de Cartan-Eilenberg
de $P_\bullet$ (6.3.4), on voit aussitôt que l'isomorphisme ainsi défini est
un isomorphisme de $\partial$-foncteurs en $Q_\bullet$.
\end{proof}
\begin{env}[6.3.8]
\oldpage[III]{144}
Soit $\rho:A\to A'$ un homomorphisme d'anneaux. Nous nous proposons de
définir un $A$-homomorphisme factoriel de degré 0 canoniquement associé à
$\rho$ :
\[
  \rho_{P_\bullet,Q_\bullet}:
  \mathbf{Tor}_\bullet^A(P_\bullet,Q_\bullet)
  \longrightarrow
  \mathbf{Tor}_\bullet^{A'}
  (P_\bullet\otimes_A A',Q_\bullet\otimes_A A').
  \tag{6.3.8.1}\label{III.6.3.8.1.fr}
\]

Pour cela, considérons une résolution de Cartan-Eilenberg projective
$L_{\bullet\bullet}$ de $P_\bullet$; considérons d'autre part une résolution
de Cartan-Eilenberg projective $L'_{\bullet\bullet}$ de
$P_\bullet\otimes_A A'$. Nous allons voir qu'on peut définir un
$A'$-homomorphisme de complexes
$L_{\bullet\bullet}\otimes_A A'\to L'_{\bullet\bullet}$, déterminé à
homotopie près. En effet, la construction de $L_{\bullet\bullet}$ est
entièrement déterminée lorsqu'on se donne (arbitrairement) pour chaque $i$,
une résolution projective $(X_{ij}^{\mathrm B})_{j\geqslant0}$ de
$\mathrm B_i(P_\bullet)$ et une résolution projective
$(X_{ij}^{\mathrm H})_{j\geqslant0}$ de $\mathrm H_i(P_\bullet)$, qui sont
respectivement égales à $\mathrm B_i^{\mathrm I}(L_{\bullet\bullet})$ et
$\mathrm H_i^{\mathrm I}(L_{\bullet\bullet})$; on en déduit successivement
\[
  \mathrm Z_i^{\mathrm I}(L_{\bullet\bullet})
  =\mathrm H_i^{\mathrm I}(L_{\bullet\bullet})
   \oplus\mathrm B_i^{\mathrm I}(L_{\bullet\bullet}),
  \qquad
  L_{i,\bullet}=\mathrm Z_i^{\mathrm I}(L_{\bullet\bullet})
   \oplus\mathrm B_{i-1}^{\mathrm I}(L_{\bullet\bullet}).
\]
Cela étant, $X_{i,\bullet}^{\mathrm B}\otimes_A A'$ n'est plus en général
une résolution de $\mathrm B_i(P_\bullet)\otimes_A A'$, mais est encore un
complexe formé de $A'$-modules projectifs, et il y a donc un
$A'$-homomorphisme
\[
  X_{i,\bullet}^{\mathrm B}\otimes_A A'
  \longrightarrow \mathrm B_i^{\mathrm I}(L'_{\bullet\bullet})
\]
compatible avec les augmentations, et déterminé à homotopie près
(M, V, 1.1). On a de même un $A'$-homomorphisme
\[
  X_{i,\bullet}^{\mathrm H}\otimes_A A'
  \longrightarrow \mathrm H_i^{\mathrm I}(L'_{\bullet\bullet})
\]
déterminé à homotopie près, d'où l'on déduit, par la construction rappelée
plus haut, un $A'$-homomorphisme
$L_{i,\bullet}\otimes_A A'\to L'_{i,\bullet}$ pour tout $i$; ces
homomorphismes (pour $i\in\mathbf Z$) sont compatibles avec les opérateurs de
dérivation $L_{i,\bullet}\to L_{i-1,\bullet}$ et les analogues pour
$L'_{\bullet\bullet}$, en vertu de la même construction, et ils constituent
donc le $A'$-homomorphisme
$L_{\bullet\bullet}\otimes_A A'\to L'_{\bullet\bullet}$ cherché.

Pour définir (6.3.8.1), il suffit alors de considérer de même une résolution
de Cartan-Eilenberg projective $M_{\bullet\bullet}$ (resp.
$M'_{\bullet\bullet}$) de $Q_\bullet$ (resp.
$Q_\bullet\otimes_A A'$), et un $A'$-homomorphisme
$M_{\bullet\bullet}\otimes_A A'\to M'_{\bullet\bullet}$. On déduit de ces
homomorphismes un $A'$-homomorphisme
\[
  (L_{\bullet\bullet}\otimes_A A')\otimes_{A'}
  (M_{\bullet\bullet}\otimes_A A')
  \longrightarrow
  L'_{\bullet\bullet}\otimes_{A'}M'_{\bullet\bullet},
\]
puis par composition un $A$-homomorphisme de bicomplexes
\[
  L_{\bullet\bullet}\otimes_A M_{\bullet\bullet}
  \longrightarrow
  L'_{\bullet\bullet}\otimes_{A'}M'_{\bullet\bullet},
\]
et en passant à l'homologie on obtient (6.3.8.1), qui est bien défini
puisqu'il provient d'un morphisme de complexes défini à homotopie près.

Si $\rho':A'\to A''$ est un second homomorphisme d'anneaux, et
$\rho'':A\to A''$ l'homomorphisme composé $\rho'\circ\rho$, il est clair que
\[
  \rho''_{P_\bullet,Q_\bullet}
  =\rho'_{P'_\bullet,Q'_\bullet}\circ\rho_{P_\bullet,Q_\bullet},
  \qquad
  P'_\bullet=P_\bullet\otimes_A A',
  \quad Q'_\bullet=Q_\bullet\otimes_A A'.
\]
Notons encore que le morphisme de bicomplexes
$L_{\bullet\bullet}\otimes_A M_{\bullet\bullet}
\to L'_{\bullet\bullet}\otimes_{A'}M'_{\bullet\bullet}$ considéré
ci-dessus définit des morphismes factoriels (en $P_\bullet$ et $Q_\bullet$)
de suites spectrales
\[
  'E_{pq}^r(P_\bullet,Q_\bullet)
  \longrightarrow
  'E_{pq}^r(P_\bullet\otimes_A A',Q_\bullet\otimes_A A')
\]
et
\[
  ''E_{pq}^r(P_\bullet,Q_\bullet)
  \longrightarrow
  ''E_{pq}^r(P_\bullet\otimes_A A',Q_\bullet\otimes_A A'),
\]
indépendants des résolutions de Cartan-Eilenberg considérées, et ayant aussi
la propriété de transitivité précédente.
\end{env}

\begin{proposition}[6.3.9]
Soit $\rho:A\to A'$ un homomorphisme d'anneaux tel que $A'$ soit un
$A$-module plat. On a alors des isomorphismes canoniques factoriels
\[
  \mathbf{Tor}_\bullet^{A'}
  (P_\bullet\otimes_A A',Q_\bullet\otimes_A A')
  \xrightarrow{\sim}
  \mathbf{Tor}_\bullet^A(P_\bullet,Q_\bullet)\otimes_A A'.
  \tag{6.3.9.1}\label{III.6.3.9.1.fr}
\]
\[
  \left\{
  \begin{aligned}
  'E(P_\bullet\otimes_A A',Q_\bullet\otimes_A A')
    &\xrightarrow{\sim}'E(P_\bullet,Q_\bullet)\otimes_A A',\\
  ''E(P_\bullet\otimes_A A',Q_\bullet\otimes_A A')
    &\xrightarrow{\sim}''E(P_\bullet,Q_\bullet)\otimes_A A'.
  \end{aligned}
  \right.
  \tag{6.3.9.2}\label{III.6.3.9.2.fr}
\]
\end{proposition}

\begin{proof}
\oldpage[III]{145}
En effet, vu l'exactitude du foncteur $M\otimes_A A'$ en $M$,
$L_{\bullet\bullet}\otimes_A A'$ et $M_{\bullet\bullet}\otimes_A A'$ sont
alors des résolutions projectives de Cartan-Eilenberg de
$P_\bullet\otimes_A A'$ et $Q_\bullet\otimes_A A'$ respectivement, d'où la
conclusion.
\end{proof}

\begin{env}[6.3.10]
Soit $\rho:A\to A'$ un homomorphisme d'anneaux; pour tout complexe
$P'_\bullet$ de $A'$-modules, $P'_{\bullet[\rho]}$ est un complexe de
$A$-modules; en outre, l'application identique
$P'_{\bullet[\rho]}\to P'_\bullet$ peut être considérée comme composée des
applications canoniques
\[
  P'_{\bullet[\rho]}\longrightarrow
  P'_{\bullet[\rho]}\otimes_A A'
  \xrightarrow{\mu}P'_\bullet,
\]
où $\mu$ est le $A'$-homomorphisme $\mu(x\otimes a')=a'x$. Si $Q'_\bullet$
est un second complexe de $A'$-modules, on a donc des homomorphismes
canoniques factoriels de degré 0
\[
  \mathbf{Tor}_\bullet^A(P'_{\bullet[\rho]},Q'_{\bullet[\rho]})
  \longrightarrow
  \mathbf{Tor}_\bullet^{A'}
  (P'_{\bullet[\rho]}\otimes_A A',Q'_{\bullet[\rho]}\otimes_A A')
  \longrightarrow
  \mathbf{Tor}_\bullet^{A'}(P'_\bullet,Q'_\bullet),
  \tag{6.3.10.1}\label{III.6.3.10.1.fr}
\]
où la première flèche est le $A$-homomorphisme défini dans (6.3.8) et la
seconde se déduit des $A'$-homomorphismes
$P'_{\bullet[\rho]}\otimes_A A'\to P'_\bullet$ et
$Q'_{\bullet[\rho]}\otimes_A A'\to Q'_\bullet$ par fonctorialité. On a des
homomorphismes analogues pour les suites spectrales de (6.3.2), et des
propriétés évidentes de transitivité, que nous laissons au lecteur le soin
d'énoncer.
\end{env}

\begin{proposition}[6.3.11]
Soit $\rho:A\to A'$ un homomorphisme d'anneaux faisant de $A'$ un
$A$-module plat. Pour tout complexe $P'_\bullet$ de $A'$-modules et tout
complexe $Q_\bullet$ de $A$-modules limités inférieurement, on a un
isomorphisme canonique factoriel
\[
  \mathbf{Tor}_\bullet^A(P'_{\bullet[\rho]},Q_\bullet)
  \xrightarrow{\sim}
  \mathbf{Tor}_\bullet^{A'}(P'_\bullet,Q_\bullet\otimes_A A').
  \tag{6.3.11.1}\label{III.6.3.11.1.fr}
\]
\end{proposition}

\begin{proof}
En effet, si $M_{\bullet\bullet}$ est une résolution projective de
Cartan-Eilenberg de $Q_\bullet$, $M_{\bullet\bullet}\otimes_A A'$ est une
résolution projective de Cartan-Eilenberg de $Q_\bullet\otimes_A A'$, et
l'on a, à un isomorphisme canonique près,
\[
  P'_{\bullet[\rho]}\otimes_A M_{\bullet\bullet}
  =P'_\bullet\otimes_{A'}(M_{\bullet\bullet}\otimes_A A');
\]
la conclusion résulte de (6.3.4).
\end{proof}

\begin{namedtheorem}[6.3.12]{Remarque}
Soit $(A^\lambda)$ un système inductif filtrant d'anneaux, et soient
$(P_\bullet^\lambda)$, $(Q_\bullet^\lambda)$ deux systèmes inductifs de
complexes de $(A^\lambda)$-modules; on a alors un isomorphisme canonique
généralisant (6.3.6.1)
\[
  \varinjlim_\lambda
  \mathbf{Tor}_\bullet^{A^\lambda}
  (P_\bullet^\lambda,Q_\bullet^\lambda)
  \xrightarrow{\sim}
  \mathbf{Tor}_\bullet^A(P_\bullet,Q_\bullet),
  \tag{6.3.12.1}\label{III.6.3.12.1.fr}
\]
où $A=\varinjlim A^\lambda$, $P_\bullet=\varinjlim P_\bullet^\lambda$,
$Q_\bullet=\varinjlim Q_\bullet^\lambda$. Une fois définis les
homomorphismes
\[
  \mathbf{Tor}_\bullet^{A^\lambda}
  (P_\bullet^\lambda,Q_\bullet^\lambda)
  \longrightarrow
  \mathbf{Tor}_\bullet^{A^\mu}
  (P_\bullet^\mu,Q_\bullet^\mu)
\]
pour $\lambda\leqslant\mu$, à l'aide de (6.3.10), la démonstration est celle
de (6.3.6).
\end{namedtheorem}

\begin{proposition}[6.3.13]
Soient $S$ une partie multiplicative de $A$, $P_\bullet$ et $Q_\bullet$ deux
complexes de $A$-modules, dans lesquels les homothéties définies par les
éléments de $S$ soient bijectives, de sorte que, si $A'=S^{-1}A$,
$P_\bullet$ et $Q_\bullet$ sont formés de $A'$-modules. Alors on a un
isomorphisme canonique
\[
  \mathbf{Tor}_\bullet^A(P_\bullet,Q_\bullet)
  \xrightarrow{\sim}
  \mathbf{Tor}_\bullet^{A'}(P_\bullet,Q_\bullet).
\]
\end{proposition}

\begin{proof}
En effet, l'hypothèse entraîne que les homomorphismes canoniques
$P_\bullet\to P_\bullet\otimes_A A'$,
$Q_\bullet\to Q_\bullet\otimes_A A'$ sont bijectifs. D'autre part, la
fonctorialité de l'hypertor montre que tout $s\in S$ définit une homothétie
bijective dans $\mathbf{Tor}_\bullet^A(P_\bullet,Q_\bullet)$, et par suite
\[
  \mathbf{Tor}_\bullet^A(P_\bullet,Q_\bullet)
  \longrightarrow
  \mathbf{Tor}_\bullet^A(P_\bullet,Q_\bullet)\otimes_A A'
\]

\oldpage[III]{146}

est aussi un homomorphisme bijectif. Comme $A'$ est un $A$-module plat, la
conclusion résulte de (6.3.9), et on a de même des isomorphismes canoniques
pour les suites spectrales.
\end{proof}

\subsection{Foncteurs hypertor locaux de complexes de Modules quasi-cohérents : cas des schémas affines}

\begin{env}[6.4.1]
Soient $S$ un schéma affine d'anneau $A$, $X$, $Y$ deux $S$-schémas affines
d'anneaux $B$, $C$ respectivement, de sorte que $B$ et $C$ sont des algèbres
sur $A$. Tout complexe $\mathscr P_\bullet$ (resp. $\mathscr Q_\bullet$) de
$\mathscr O_X$-Modules (resp. $\mathscr O_Y$-Modules) quasi-cohérents est de
la forme $\widetilde{P_\bullet}$ (resp. $\widetilde{Q_\bullet}$), où
$P_\bullet$ (resp. $Q_\bullet$) est un complexe de $B$-modules (resp.
$C$-modules) (I, 1.3.7 et 1.3.8). On peut évidemment considérer $P_\bullet$
et $Q_\bullet$ comme des complexes de $A$-modules et former les
$\mathbf{Tor}_n^A(P_\bullet,Q_\bullet)$; en outre, en vertu du caractère
bifonctoriel de $\mathbf{Tor}_n^A(P_\bullet,Q_\bullet)$, les $A$-algèbres
$B$ et $C$ opèrent dans ce $A$-module, et ces opérations en font un
$(B,C)$-bimodule, ou, ce qui revient au même, un module sur
$B\otimes_A C=A(X\times_S Y)$. On a donc défini de la sorte un
$\mathscr O_{X\times_S Y}$-Module quasi-cohérent
\[
  \mathscr{T}\!or_n^{\mathscr O_S}(\mathscr P_\bullet,\mathscr Q_\bullet)
  =\bigl(\mathbf{Tor}_n^A(P_\bullet,Q_\bullet)\bigr)^{\sim}
  \tag{6.4.1.1}\label{III.6.4.1.1.fr}
\]
que l'on appelle l'\emph{hypertor local d'indice $n$} des complexes
$\mathscr P_\bullet$ et $\mathscr Q_\bullet$ et que l'on note aussi
$\mathscr{T}\!or_n^S(\mathscr P_\bullet,\mathscr Q_\bullet)$.
\end{env}

\begin{namedtheorem}[6.4.2]{Lemme}
Avec les notations de (6.4.1), supposons que l'anneau $A$ soit de la forme
$R^{-1}A'$, où $A'$ est un anneau et $R$ une partie multiplicative de $A'$.
Soit $S'=\operatorname{Spec}(A')$, de sorte que $X$ et $Y$ peuvent être
considérés comme des $S'$-préschémas et que l'on a
$X\times_{S'}Y=X\times_S Y$ (I, 1.6.2 et 3.2.4). On a alors
\[
  \mathscr{T}\!or_\bullet^{S'}(\mathscr P_\bullet,\mathscr Q_\bullet)
  =\mathscr{T}\!or_\bullet^S(\mathscr P_\bullet,\mathscr Q_\bullet).
\]
\end{namedtheorem}

\begin{proof}
Cela résulte de la formule (6.4.1.1) et de (6.3.13).
\end{proof}

\begin{env}[6.4.3]
Avec les notations et hypothèses de (6.4.1), soient
$\mathscr F=\widetilde F$ un $\mathscr O_X$-Module quasi-cohérent,
$\mathscr G=\widetilde G$ un $\mathscr O_Y$-Module quasi-cohérent;
considérant $\mathscr F$ et $\mathscr G$ comme des complexes de Modules, on
notera $\mathscr{T}\!or_n^{\mathscr O_S}(\mathscr F,\mathscr G)$ ou
$\mathscr{T}\!or_n^S(\mathscr F,\mathscr G)$ leur hypertor d'indice $n$; il
résulte de (6.3.5 (ii)) que l'on a
\[
  \mathscr{T}\!or_n^{\mathscr O_S}(\mathscr F,\mathscr G)
  =\bigl(\operatorname{Tor}_n^A(F,G)\bigr)^{\sim}.
  \tag{6.4.3.1}\label{III.6.4.3.1.fr}
\]

Revenons alors au cas général de deux complexes de Modules quasi-cohérents
$\mathscr P_\bullet$, $\mathscr Q_\bullet$. Les formules (6.4.1.1) et
(6.4.3.1) montrent, compte tenu de la prop. (6.3.2), que
$\mathscr{T}\!or_\bullet^S(\mathscr P_\bullet,\mathscr Q_\bullet)$ est
l'aboutissement de deux suites spectrales
$'\mathscr E(\mathscr P_\bullet,\mathscr Q_\bullet)$,
$''\mathscr E(\mathscr P_\bullet,\mathscr Q_\bullet)$, dont les termes
$E_2$ sont donnés par
\[
  '\mathscr E_{pq}^2
  =\mathscr H_p\bigl(\mathscr{T}\!or_q^S
    (\mathscr P_\bullet,\mathscr Q_\bullet)\bigr)
  \tag{6.4.3.2}\label{III.6.4.3.2.fr}
\]
\[
  ''\mathscr E_{pq}^2
  =\bigoplus_{q'+q''=q}
  \mathscr{T}\!or_p^S
  \bigl(\mathscr H_{q'}(\mathscr P_\bullet),
        \mathscr H_{q''}(\mathscr Q_\bullet)\bigr)
  \tag{6.4.3.3}\label{III.6.4.3.3.fr}
\]
où $\mathscr{T}\!or_q^S(\mathscr P_\bullet,\mathscr Q_\bullet)$ est le
bicomplexe de $\mathscr O_{X\times_S Y}$-Modules quasi-cohérents
$\mathscr{T}\!or_q^S(\mathscr P_i,\mathscr Q_j)$.
\end{env}

\begin{env}[6.4.4]
Considérons maintenant deux autres schémas affines
$X^{(1)}=\operatorname{Spec}(B^{(1)})$,
$Y^{(1)}=\operatorname{Spec}(C^{(1)})$, où $B^{(1)}$ et $C^{(1)}$ sont des
$A$-algèbres, et supposons donnés deux

\oldpage[III]{147}
S-morphismes $u:X^{(1)}\to X$, $v:Y^{(1)}\to Y$, correspondant à des
$A$-homomorphismes $\varphi:B\to B^{(1)}$, $\psi:C\to C^{(1)}$.
Considérons les complexes
$u^*(\mathscr P_\bullet)=(P_\bullet\otimes_B B^{(1)})^{\sim}$ de
$\mathscr O_{X^{(1)}}$-Modules,
$v^*(\mathscr Q_\bullet)=(Q_\bullet\otimes_C C^{(1)})^{\sim}$ de
$\mathscr O_{Y^{(1)}}$-Modules. Les $A$-homomorphismes canoniques
\[
  P_\bullet\longrightarrow P_\bullet\otimes_B B^{(1)},
  \qquad
  Q_\bullet\longrightarrow Q_\bullet\otimes_C C^{(1)}
\]
donnent par fonctorialité un $A$-homomorphisme
\[
  \mathbf{Tor}_\bullet^A(P_\bullet,Q_\bullet)
  \longrightarrow
  \mathbf{Tor}_\bullet^A
  (P_\bullet\otimes_B B^{(1)},Q_\bullet\otimes_C C^{(1)});
\]
en outre, toujours par fonctorialité, cet homomorphisme est en fait un
homomorphisme de $(B\otimes_A C)$-modules. D'où l'on conclut que l'on a défini
ainsi un $(u\times_S v)$-morphisme
\[
  \theta:\mathscr{T}\!or_\bullet^S(\mathscr P_\bullet,\mathscr Q_\bullet)
  \longrightarrow
  \mathscr{T}\!or_\bullet^S
  (u^*(\mathscr P_\bullet),v^*(\mathscr Q_\bullet))
  \tag{6.4.4.1}\label{III.6.4.4.1.fr}
\]
et par suite, un homomorphisme de
$\mathscr O_{X^{(1)}\times_S Y^{(1)}}$-Modules
\[
  \theta^\sharp:(u\times_S v)^*
  \bigl(\mathscr{T}\!or_\bullet^S(\mathscr P_\bullet,\mathscr Q_\bullet)\bigr)
  \longrightarrow
  \mathscr{T}\!or_\bullet^S
  (u^*(\mathscr P_\bullet),v^*(\mathscr Q_\bullet))
  \tag{6.4.4.2}\label{III.6.4.4.2.fr}
\]
qui est évidemment un morphisme de bi-$\partial$-foncteurs dans les catégories
de Modules quasi-cohérents limités inférieurement.

L'homomorphisme (6.4.4.2) n'est pas nécessairement bijectif; toutefois :
\end{env}

\begin{namedtheorem}[6.4.5]{Lemme}
Avec les notations de (6.4.4), supposons que $u$ et $v$ soient des immersions
ouvertes; alors l'homomorphisme (6.4.4.2) est bijectif.
\end{namedtheorem}

\begin{proof}
Identifions $X^{(1)}$ (resp. $Y^{(1)}$) à un ouvert de $X$ (resp. $Y$);
$X^{(1)}$ (resp. $Y^{(1)}$) est alors réunion d'ouverts de la forme
$\mathrm D(f)$ (resp. $\mathrm D(g)$), où $f\in B$ (resp. $g\in C$), et les
préschémas induits $\mathrm D(\varphi(f))$ et $\mathrm D(f)$ (resp.
$\mathrm D(\psi(g))$ et $\mathrm D(g)$) sont isomorphes. Il suffira de prouver
le lemme lorsque $X^{(1)}$ (resp. $Y^{(1)}$) est de la forme $\mathrm D(f)$
(resp. $\mathrm D(g)$); en effet, si ce point est établi, et si l'on revient au
cas général, il suffira de prouver que la restriction de $\theta^\sharp$ à
chaque ouvert $\mathrm D(f)\times_S\mathrm D(g)$ est un isomorphisme; or, si
$u_1:\mathrm D(f)\to X^{(1)}$, $v_1:\mathrm D(g)\to Y^{(1)}$ sont les
injections canoniques, la restriction précédente n'est autre que
$(u_1\times_S v_1)^*(\theta^\sharp)$; mais il est immédiat, en vertu des
définitions (6.4.4) et de ($0_{\mathrm I}$, 4.4.8), qu'en la composant avec
l'homomorphisme canonique
\[
  (u_1\times_S v_1)^*
  \bigl(\mathscr{T}\!or_\bullet^S
    (u^*(\mathscr P_\bullet),v^*(\mathscr Q_\bullet))\bigr)
  \longrightarrow
  \mathscr{T}\!or_\bullet^S
  ((u')^*(\mathscr P_\bullet),(v')^*(\mathscr Q_\bullet))
  \tag{6.4.5.1}\label{III.6.4.5.1.fr}
\]
où $u'=u\circ u_1$ et $v'=v\circ v_1$, on obtient l'homomorphisme canonique
\[
  (u'\times_S v')^*
  \bigl(\mathscr{T}\!or_\bullet^S(\mathscr P_\bullet,\mathscr Q_\bullet)\bigr)
  \longrightarrow
  \mathscr{T}\!or_\bullet^S
  ((u')^*(\mathscr P_\bullet),(v')^*(\mathscr Q_\bullet))
  \tag{6.4.5.2}\label{III.6.4.5.2.fr}
\]
et si l'on sait que (6.4.5.1) et (6.4.5.2) sont des isomorphismes, il en
résultera qu'il en est de même de $(u_1\times_S v_1)^*(\theta^\sharp)$.

Supposons donc que $X^{(1)}=\mathrm D(f)$ et $Y^{(1)}=\mathrm D(g)$, de sorte
que $B^{(1)}=B_f$ et $C^{(1)}=C_g$; $u^*(\mathscr P_\bullet)$ (resp.
$v^*(\mathscr Q_\bullet)$) s'identifie alors à
$((P_\bullet)_f)^{\sim}$ (resp. $((Q_\bullet)_g)^{\sim}$); d'autre part,
$X^{(1)}\times_S Y^{(1)}$ s'identifie au sous-schéma ouvert
$\mathrm D(f\otimes g)$ de
$X\times_S Y=\operatorname{Spec}(B\otimes_A C)$ (II, 4.3.2.4); il s'agit de
prouver que l'homomorphisme
\[
  \bigl(\mathbf{Tor}_\bullet^A(P_\bullet,Q_\bullet)\bigr)_{f\otimes g}
  \longrightarrow
  \mathbf{Tor}_\bullet^A((P_\bullet)_f,(Q_\bullet)_g)
  \tag{6.4.5.3}\label{III.6.4.5.3.fr}
\]

\oldpage[III]{148}
déduit par fonctorialité des homomorphismes canoniques
$P_\bullet\to(P_\bullet)_f$, $Q_\bullet\to(Q_\bullet)_g$, est bijectif. Or
($0_{\mathrm I}$, 1.6.1), on peut écrire
$(P_\bullet)_f=\varinjlim P_\bullet^{(n)}$, où les $P_\bullet^{(n)}$ sont tous
des complexes de $B$-modules identiques à $P_\bullet$, l'application
$P_\bullet^{(m)}\to P_\bullet^{(n)}$ pour $m\leqslant n$ étant la
multiplication par $f^{n-m}$; on a un résultat analogue pour $Q_\bullet$ en
remplaçant $f$ par $g$; d'autre part, il est clair que l'homomorphisme
\[
  \mathbf{Tor}_\bullet^A(P_\bullet^{(m)},Q_\bullet^{(m)})
  \longrightarrow
  \mathbf{Tor}_\bullet^A(P_\bullet^{(n)},Q_\bullet^{(n)})
\]
correspondant aux homomorphismes
$P_\bullet^{(m)}\to P_\bullet^{(n)}$ et
$Q_\bullet^{(m)}\to Q_\bullet^{(n)}$ est par définition la multiplication
par $(f\otimes g)^{n-m}$. La conclusion résulte alors de ($0_{\mathrm I}$,
1.6.1) appliqué au premier membre de (6.4.5.3) et de (6.3.6).
\end{proof}

\begin{env}[6.4.6]
Avec les notations de (6.4.4), on définit de même des homomorphismes canoniques
de foncteurs spectraux
\[
  \left\{
  \begin{aligned}
  (u\times_S v)^*\bigl('\mathscr E(\mathscr P_\bullet,\mathscr Q_\bullet)\bigr)
    &\longrightarrow
      '\mathscr E(u^*(\mathscr P_\bullet),v^*(\mathscr Q_\bullet)),\\
  (u\times_S v)^*\bigl(''\mathscr E(\mathscr P_\bullet,\mathscr Q_\bullet)\bigr)
    &\longrightarrow
      ''\mathscr E(u^*(\mathscr P_\bullet),v^*(\mathscr Q_\bullet)).
  \end{aligned}
  \right.
  \tag{6.4.6.1}\label{III.6.4.6.1.fr}
\]
et le raisonnement de (6.4.5) montre que lorsque $u$ et $v$ sont des
immersions ouvertes, les homomorphismes (6.4.6.1) sont bijectifs : en effet,
compte tenu de (6.3.6.2) et (6.3.6.3), il prouve que c'est un isomorphisme pour
les termes $E^2$, et (6.4.5) montre que c'est un isomorphisme pour les
aboutissements; on conclut donc à l'aide de (0, 11.1.2) et (0, 11.2.4).
\end{env}

\subsection{Foncteurs hypertor locaux de complexes de Modules quasi-cohérents : cas général}

\begin{env}[6.5.1]
Considérons maintenant un préschéma $S$ quelconque et deux $S$-préschémas
quelconques $X$, $Y$; soit $\mathscr P_\bullet$ (resp. $\mathscr Q_\bullet$)
un complexe de $\mathscr O_X$-Modules (resp. $\mathscr O_Y$-Modules)
quasi-cohérents. Posons $Z=X\times_S Y$; nous allons définir des
$\mathscr O_Z$-Modules quasi-cohérents
$\mathscr{T}\!or_n^S(\mathscr P_\bullet,\mathscr Q_\bullet)$ dits
\emph{hypertor locaux} de $\mathscr P_\bullet$ et $\mathscr Q_\bullet$ qui se
réduiront à ceux déjà définis dans (6.4) lorsque $S$, $X$ et $Y$ sont affines.

Lorsque $\mathscr P_\bullet$ et $\mathscr Q_\bullet$ se réduisent
respectivement à leurs termes de degré 0, $\mathscr F$ et $\mathscr G$ (les
autres étant nuls), on écrira
$\mathscr{T}\!or_n^S(\mathscr F,\mathscr G)$ au lieu de
$\mathscr{T}\!or_n^S(\mathscr P_\bullet,\mathscr Q_\bullet)$.
\end{env}

\begin{env}[6.5.2]
Supposons d'abord $S$ affine, et soient $(X_\lambda)$, $(Y_\mu)$ des
recouvrements de $X$ et $Y$ respectivement par des ouverts affines; alors les
$Z_{\lambda\mu}=X_\lambda\times_S Y_\mu$ forment un recouvrement ouvert affine
de $Z$. Posons
$\mathscr P_{\lambda\bullet}=\mathscr P_\bullet|X_\lambda$,
$\mathscr Q_{\mu\bullet}=\mathscr Q_\bullet|Y_\mu$; nous avons donc pour tout
couple $(\lambda,\mu)$ un $\mathscr O_{Z_{\lambda\mu}}$-Module quasi-cohérent
\[
  \mathscr F_{\lambda\mu}
  =\mathscr{T}\!or_n^S
   (\mathscr P_{\lambda\bullet},\mathscr Q_{\mu\bullet}),
\]
et il faut montrer que les $\mathscr F_{\lambda\mu}$ vérifient la condition de
recollement ($0_{\mathrm I}$, 3.3.1). Pour cela, il suffit de vérifier que pour
tout ouvert affine $U\subset X_\lambda\cap X_{\lambda'}$ (resp.
$V\subset Y_\mu\cap Y_{\mu'}$), les restrictions de
$\mathscr F_{\lambda\mu}$ et de $\mathscr F_{\lambda'\mu'}$ à $U\times_S V$
sont canoniquement isomorphes; mais cela découle aussitôt de l'existence
d'isomorphismes canoniques de ces restrictions sur
$\mathscr{T}\!or_n^S(\mathscr P_\bullet|U,\mathscr Q_\bullet|V)$ (6.4.5). En
outre, il résulte aussitôt de cette définition et de (6.4.5) que le
$\mathscr O_Z$-Module ainsi défini ne dépend pas (à un isomorphisme près) des
recouvrements ouverts consi\oldpage[III]{149}dérés;
nous le noterons donc
$\mathscr{T}\!or_n^S(\mathscr P_\bullet,\mathscr Q_\bullet)$; il résulte enfin
de (6.4.5) que pour toute partie ouverte $U$ (resp. $V$) de $X$ (resp. $Y$), la
restriction de $\mathscr{T}\!or_n^S(\mathscr P_\bullet,\mathscr Q_\bullet)$ à
$U\times_S V$ est canoniquement isomorphe à
$\mathscr{T}\!or_n^S(\mathscr P_\bullet|U,\mathscr Q_\bullet|V)$.
\end{env}

\begin{env}[6.5.3]
Passons maintenant au cas général où $S$ est quelconque, et soit $(S_\alpha)$
un recouvrement de $S$ formé d'ouverts affines; désignons par $X_\alpha$
(resp. $Y_\alpha$) l'image réciproque de $S_\alpha$ dans $X$ (resp. $Y$); il
faut encore prouver que les faisceaux
\[
  \mathscr{T}\!or_n^{S_\alpha}
  (\mathscr P_\bullet|X_\alpha,\mathscr Q_\bullet|Y_\alpha)
  =\mathscr G_\alpha
\]
vérifient la condition de recollement. Il suffit de définir, pour tout ouvert
affine $T$ contenu dans $S_\alpha\cap S_\beta$, des isomorphismes canoniques des
restrictions de $\mathscr G_\alpha$ et $\mathscr G_\beta$ à $U\times_S V$ (en
désignant par $U$ et $V$ les images réciproques de $T$ dans $X$ et $Y$
respectivement) sur
$\mathscr{T}\!or_n^T(\mathscr P_\bullet|U,\mathscr Q_\bullet|V)$; on peut, en
outre, se borner au cas où $T$ s'écrit à la fois $\mathrm D(f_\alpha)$ et
$\mathrm D(f_\beta)$, $f_\alpha$ (resp. $f_\beta$) étant une section de
$\mathscr O_S$ au-dessus de $S_\alpha$ (resp. $S_\beta$); mais alors
$\mathscr{T}\!or_n^T(\mathscr P_\bullet|U,\mathscr Q_\bullet|V)$ est
canoniquement isomorphe à
$\mathscr{T}\!or_n^{S_\alpha}(\mathscr P_\bullet|U,\mathscr Q_\bullet|V)$
d'une part, à
$\mathscr{T}\!or_n^{S_\beta}(\mathscr P_\bullet|U,\mathscr Q_\bullet|V)$
d'autre part, en vertu de (6.4.2); comme on vient de définir des isomorphismes
canoniques de $\mathscr G_\alpha$ sur
$\mathscr{T}\!or_n^{S_\alpha}(\mathscr P_\bullet|U,\mathscr Q_\bullet|V)$ et
de $\mathscr G_\beta$ sur
$\mathscr{T}\!or_n^{S_\beta}(\mathscr P_\bullet|U,\mathscr Q_\bullet|V)$
(6.5.2), cela achève de définir le $\mathscr O_Z$-Module
$\mathscr{T}\!or_n^S(\mathscr P_\bullet,\mathscr Q_\bullet)$. En outre, pour
toute partie ouverte $U$ (resp. $V$) de $X$ (resp. $Y$),
$\mathscr{T}\!or_n^S(\mathscr P_\bullet|U,\mathscr Q_\bullet|V)$ est
canoniquement isomorphe à la restriction de
$\mathscr{T}\!or_n^S(\mathscr P_\bullet,\mathscr Q_\bullet)$ à $U\times_S V$.

Il est immédiat que l'on a ainsi défini (dans les catégories de complexes de
Modules quasi-cohérents limités inférieurement) un bi-$\partial$-foncteur
$\mathscr{T}\!or_\bullet^S(\mathscr P_\bullet,\mathscr Q_\bullet)$ à valeurs
dans la catégorie des $\mathscr O_Z$-Modules, car il est clair que la question
est locale sur $X$, $Y$ et $S$, en vertu de (6.4.5) et de la remarque que
(6.4.4.2) est un morphisme de bi-$\partial$-foncteurs. On notera que si
$\mathscr P_\bullet$ et $\mathscr Q_\bullet$ sont réduits respectivement à
leurs termes de degré 0, $\mathscr F$ et $\mathscr G$,
$\mathscr{T}\!or_0^S(\mathscr F,\mathscr G)$ n'est autre, en vertu de
(6.4.1.1), que le produit tensoriel externe $\mathscr F\otimes_S\mathscr G$
défini dans (I, 9.1.2); cela résulte en effet de (I, 9.1.3).
\end{env}

\begin{env}[6.5.4]
Il résulte de la construction précédente et des remarques faites dans (6.4.6)
que $\mathscr{T}\!or_\bullet^S(\mathscr P_\bullet,\mathscr Q_\bullet)$ est
l'aboutissement de deux foncteurs spectraux,
$'\mathscr E(\mathscr P_\bullet,\mathscr Q_\bullet)$,
$''\mathscr E(\mathscr P_\bullet,\mathscr Q_\bullet)$, de termes $E_2$ égaux à
\[
  '\mathscr E_{pq}^2
  =\mathscr H_p\bigl(\mathscr{T}\!or_q^S
    (\mathscr P_\bullet,\mathscr Q_\bullet)\bigr)
  \tag{6.5.4.1}\label{III.6.5.4.1.fr}
\]
\[
  ''\mathscr E_{pq}^2
  =\bigoplus_{q'+q''=q}
   \mathscr{T}\!or_p^S
   \bigl(\mathscr H_{q'}(\mathscr P_\bullet),
         \mathscr H_{q''}(\mathscr Q_\bullet)\bigr).
  \tag{6.5.4.2}\label{III.6.5.4.2.fr}
\]
La suite spectrale (6.5.4.2) est toujours régulière; les deux suites spectrales
sont birégulières si $\mathscr P_\bullet$ et $\mathscr Q_\bullet$ sont limités
inférieurement. Un autre cas où les deux suites précédentes sont birégulières
est le suivant :
\end{env}

\begin{env}[6.5.5]
Nous dirons que sur un espace topologique $T$ un faisceau d'anneaux
$\mathscr A$ est \emph{de dimension cohomologique $\leqslant n$} si, pour tout
$t\in T$ l'anneau $\mathscr A_t$ est de dimension cohomologique $\leqslant n$;
on dira alors aussi que l'\emph{espace annelé} $(T,\mathscr A)$ est
\emph{de dimension cohomologique $\leqslant n$}. On dira qu'un faisceau
d'anneaux (resp. un espace annelé) est de dimension cohomologique \emph{finie}
s'il existe un entier $n$ tel qu'il soit de dimension cohomologique
$\leqslant n$. On notera que si les $\mathscr A_t$ sont des anneaux locaux
(commutatifs) noethériens,

\oldpage[III]{150}
dire qu'ils sont de dimension cohomologique $\leqslant n$ signifie qu'ils sont
réguliers et de dimension (de Krull) $\leqslant n$ ($0_{\mathrm{IV}}$, 17.3.1).
Avec la terminologie de la théorie de la dimension que nous introduirons au
chap. IV, il revient au même de dire qu'un préschéma localement noethérien $T$
est de dimension cohomologique $\leqslant n$, ou de dire qu'il est régulier
($0_{\mathrm I}$, 4.1.4) et de dimension $\leqslant n$; cela signifie que pour
tout ouvert affine $U$ de $T$, l'anneau $\Gamma(U,\mathscr O_T)$ est de
dimension cohomologique $\leqslant n$ ($0_{\mathrm{IV}}$, 17.2.6). Cela étant,
cette dernière remarque, jointe à (6.3.2), prouve que si $S$ est localement
noethérien et de dimension cohomologique finie, les suites spectrales
$'\mathscr E(\mathscr P_\bullet,\mathscr Q_\bullet)$ et
$''\mathscr E(\mathscr P_\bullet,\mathscr Q_\bullet)$ sont birégulières.

Il est clair que
$\mathscr{T}\!or_\bullet^S(\mathscr P_\bullet,\mathscr Q_\bullet)$ se
transforme en
$\mathscr{T}\!or_\bullet^S(\mathscr Q_\bullet,\mathscr P_\bullet)$ (à un
isomorphisme près) par l'isomorphisme canonique de $X\times_S Y$ sur
$Y\times_S X$.
\end{env}

\begin{proposition}[6.5.6]
Soit $(\mathscr P_{\alpha\bullet})$ un système inductif filtrant de complexes
de $\mathscr O_X$-Modules quasi-cohérents; il existe alors un isomorphisme
canonique
\[
  \varinjlim_\alpha
  \mathscr{T}\!or_\bullet^S
  (\mathscr P_{\alpha\bullet},\mathscr Q_\bullet)
  \xrightarrow{\sim}
  \mathscr{T}\!or_\bullet^S
  \left(\varinjlim_\alpha\mathscr P_{\alpha\bullet},\mathscr Q_\bullet\right).
  \tag{6.5.6.1}\label{III.6.5.6.1.fr}
\]
\end{proposition}

\begin{proof}
La question étant locale sur $S$, $X$ et $Y$, on peut supposer $S$, $X$, $Y$
affines et la proposition se réduit alors à (6.3.6).
\end{proof}

\begin{namedtheorem}[6.5.7]{Remarques}
\textnormal{(i)} Considérons en particulier le cas où $S=X=Y$,
$\mathscr P_\bullet$ et $\mathscr Q_\bullet$ étant donc deux complexes de
$\mathscr O_S$-Modules quasi-cohérents; alors les
$\mathscr{T}\!or_n^S(\mathscr P_\bullet,\mathscr Q_\bullet)$ sont des
$\mathscr O_S$-Modules quasi-cohérents; en outre, pour tout point $z\in S$, il
résulte de (6.5.6) que l'on a un isomorphisme canonique
\[
  \bigl(\mathscr{T}\!or_n^S
    (\mathscr P_\bullet,\mathscr Q_\bullet)\bigr)_z
  \xrightarrow{\sim}
  \mathbf{Tor}_n^{\mathscr O_z}
  \bigl((\mathscr P_\bullet)_z,(\mathscr Q_\bullet)_z\bigr)
  \tag{6.5.7.1}\label{III.6.5.7.1.fr}
\]
car la question est locale et on est ramené au cas des modules, en vertu de
(6.4.1.1).

\textnormal{(ii)} On peut généraliser la définition des hypertor au cas de
deux complexes de $\mathscr O_X$-Modules $\mathscr P_\bullet$,
$\mathscr Q_\bullet$ sur un même espace annelé $(X,\mathscr O_X)$; pour tout
ouvert $U$ de $X$, posons en effet
$A(U)=\Gamma(U,\mathscr O_X)$,
$P_\bullet(U)=\Gamma(U,\mathscr P_\bullet)$,
$Q_\bullet(U)=\Gamma(U,\mathscr Q_\bullet)$; les $A(U)$-modules
\[
  \mathbf{Tor}_n^{A(U)}(P_\bullet(U),Q_\bullet(U))
\]
forment alors un préfaisceau sur $X$, et l'on désigne par
$\mathscr{T}\!or_n^X(\mathscr P_\bullet,\mathscr Q_\bullet)$ le
$\mathscr O_X$-Module associé à ce préfaisceau. Lorsque $X$ est un préschéma,
il résulte de (6.3.12) que ce $\mathscr O_X$-Module est canoniquement isomorphe
à l'hypertor défini ci-dessus. Nous ne développerons pas davantage cette
généralisation.
\end{namedtheorem}

\begin{proposition}[6.5.8]
Soient $X$, $Y$ deux $S$-préschémas, $\mathscr F$ (resp. $\mathscr G$) un
$\mathscr O_X$-Module (resp. un $\mathscr O_Y$-Module) quasi-cohérent. Si
$\mathscr F$ ou $\mathscr G$ est $S$-plat, on a
$\mathscr{T}\!or_n^S(\mathscr F,\mathscr G)=0$ pour $n\neq0$.
\end{proposition}

\begin{proof}
La question étant locale sur $X$ et $Y$, on peut supposer $X$, $Y$ et $S$
affines, d'anneaux respectifs $B$, $C$, $A$, et
$\mathscr F=\widetilde M$, $\mathscr G=\widetilde N$, $M$ (resp. $N$) étant
un $B$-module (resp. un $C$-module). Supposons par exemple que $\mathscr F$
soit $S$-plat, ce qui signifie que pour tout $s\in S$, $M_s$ est un
$A_s$-module plat ($0_{\mathrm I}$, 6.7.1); par suite $M$ est un $A$-module plat
($0_{\mathrm I}$, 6.3.3), et l'on sait que
$\operatorname{Tor}_n^A(M,N)=0$ pour $n>0$ et pour tout $C$-module $N$
($0_{\mathrm I}$, 6.1.1), d'où la conclusion par (6.4.1.1).
\end{proof}

\begin{namedtheorem}[6.5.9]{Corollaire}
Soient $X$, $Y$ deux $S$-préschémas, $\mathscr P_\bullet$ (resp.
$\mathscr Q_\bullet$) un complexe de

\oldpage[III]{151}
$\mathscr O_X$-Modules (resp. de $\mathscr O_Y$-Modules) quasi-cohérents limité
inférieurement. Supposons que tous les $\mathscr P_i$ soient $S$-plats. Alors
il existe un isomorphisme canonique de $\partial$-foncteurs en
$\mathscr Q_\bullet$
\[
  \mathscr{T}\!or_\bullet^S(\mathscr P_\bullet,\mathscr Q_\bullet)
  \xrightarrow{\sim}
  \mathscr H_\bullet(\mathscr P_\bullet\otimes_S\mathscr Q_\bullet).
  \tag{6.5.9.1}\label{III.6.5.9.1.fr}
\]
\end{namedtheorem}

\begin{proof}
Ce n'est autre que (6.3.7) lorsque $S$, $X$, $Y$ sont affines; on passe de là
au cas général par les raisonnements de (6.5.2) et (6.5.3).
\end{proof}

\begin{namedtheorem}[6.5.10]{Corollaire}
Supposons que $X$ soit plat sur $S$ ($0_{\mathrm I}$, 6.7.1), que
$\mathscr P_\bullet$ et $\mathscr Q_\bullet$ soient limités inférieurement et
que tous les $\mathscr P_i$ soient des $\mathscr O_X$-Modules localement
libres (non nécessairement de type fini). Alors l'homomorphisme (6.5.9.1) est
bijectif.
\end{namedtheorem}

\begin{proof}
En effet, l'hypothèse de (6.5.9) est remplie, la platitude étant une propriété
ponctuelle sur $X$ par définition et toute somme directe de modules plats étant
un module plat ($0_{\mathrm I}$, 6.1.2).
\end{proof}

\begin{proposition}[6.5.11]
Soient $X'$, $Y'$ deux $S$-préschémas, $f:X\to X'$, $g:Y\to Y'$ deux
$S$-morphismes affines. Soit $\mathscr P_\bullet$ (resp.
$\mathscr Q_\bullet$) un complexe de $\mathscr O_X$-Modules (resp. de
$\mathscr O_Y$-Modules) quasi-cohérents; on a alors un isomorphisme canonique
fonctoriel
\[
  (f\times_S g)_*
  \bigl(\mathscr{T}\!or_\bullet^S
    (\mathscr P_\bullet,\mathscr Q_\bullet)\bigr)
  \xrightarrow{\sim}
  \mathscr{T}\!or_\bullet^S
  (f_*(\mathscr P_\bullet),g_*(\mathscr Q_\bullet)).
  \tag{6.5.11.1}\label{III.6.5.11.1.fr}
\]
\end{proposition}

\begin{proof}
Comme $f$ et $g$ sont affines, $f_*(\mathscr P_\bullet)$ et
$g_*(\mathscr Q_\bullet)$ sont des complexes de Modules quasi-cohérents
(II, 1.2.6), et si l'on pose $Z'=X'\times_S Y'$, les deux membres de (6.5.11.1)
sont des $\mathscr O_{Z'}$-Modules quasi-cohérents (6.5.1); on se ramène
aisément au cas où $S$, $X'$ et $Y'$ sont affines; mais alors il en est de même
par hypothèse de $X$ et $Y$ et la vérification résulte aussitôt de (6.4.1.1) et
(I, 1.6.3).
\end{proof}

\begin{namedtheorem}[6.5.12]{Remarque}
Soient $X'$, $Y'$ deux $S$-préschémas et supposons que $X'$ soit $S$-plat;
soient $X$ un sous-préschéma fermé de $X'$, $i:X\to X'$ l'injection canonique,
$\mathscr P_\bullet$ (resp. $\mathscr Q_\bullet$) un complexe de
$\mathscr O_X$-Modules (resp. de $\mathscr O_{Y'}$-Modules) quasi-cohérents,
limité inférieurement. Soit enfin $\mathscr L'_{\bullet\bullet}$ une résolution
de $i_*(\mathscr P_\bullet)$ formé de $\mathscr O_{X'}$-Modules localement
libres, telle que tout point de $X'$ ait un voisinage ouvert affine $U$ pour
lequel $\mathscr L'_{j,\bullet}|U$ soit une résolution libre de
$i_*(\mathscr P_j)|U$ pour tout $j$. On a alors un isomorphisme canonique
\[
  (i\times_S 1)_*
  \bigl(\mathscr{T}\!or_\bullet^S
    (\mathscr P_\bullet,\mathscr Q_\bullet)\bigr)
  \xrightarrow{\sim}
  \mathscr H_\bullet
  (\mathscr L'_{\bullet\bullet}\otimes_S\mathscr Q_\bullet).
  \tag{6.5.12.1}\label{III.6.5.12.1.fr}
\]

Si $S$, $X'$, $Y'$ sont affines et si $\mathscr L'_{j,\bullet}$ est une
résolution libre de $i_*(\mathscr P_j)$ pour tout $j$, on est ramené, en vertu
de (6.5.11), au cas où $X'=X$ est $S$-plat, et il suffit d'appliquer (6.3.4).
Dans le cas général, on définit localement l'isomorphisme (6.5.12.1), et il
s'agit de vérifier que cette définition donne bien un isomorphisme global.
Pour cela, il faut se reporter à la définition du premier isomorphisme
(6.3.4.1) qui provient d'un isomorphisme de suites spectrales (0, 11.6.5 et
11.5.3), obtenu lui-même à partir d'un morphisme de bicomplexes
$L_{\bullet\bullet}\to L''_{\bullet\bullet}$, où $L_{\bullet\bullet}$ est la
résolution donnée de $P_\bullet$, $L''_{\bullet\bullet}$ une résolution
projective de $P_\bullet$ dans la catégorie des complexes de $A$-modules
limités

\oldpage[III]{152}
inférieurement (cf. (0, 11.5.2.2)); notre assertion résulte de ce que
l'isomorphisme (6.3.4.1) ne dépend pas de la résolution projective
$L''_{\bullet\bullet}$ choisie, en vertu de l'existence d'un homotopisme entre
deux telles résolutions (M, V, 1.2).
\end{namedtheorem}

\begin{proposition}[6.5.13]
Soient $X$, $Y$ deux $S$-préschémas, et supposons vérifiée l'une des conditions
suivantes :
\begin{enumerate}
  \item[(i)] $X$ et $Z=X\times_S Y$ sont localement noethériens et $X$ est plat
    sur $S$.
  \item[(ii)] $S$ et $X$ sont localement noethériens et $Y$ est de type fini
    sur $S$.
\end{enumerate}
Soit $\mathscr P_\bullet$ (resp. $\mathscr Q_\bullet$) un complexe de
$\mathscr O_X$-Modules (resp. de $\mathscr O_Y$-Modules) quasi-cohérents,
limité inférieurement. On suppose en outre que, pour tout $n$,
$\mathscr H_n(\mathscr P_\bullet)$ (resp.
$\mathscr H_n(\mathscr Q_\bullet)$) est un $\mathscr O_X$-Module (resp. un
$\mathscr O_Y$-Module) de type fini. Alors les
$\mathscr{T}\!or_n^S(\mathscr P_\bullet,\mathscr Q_\bullet)$ sont des
$\mathscr O_Z$-Modules cohérents.
\end{proposition}

\begin{proof}
Comme $\mathscr P_\bullet$ et $\mathscr Q_\bullet$ sont limités
inférieurement, la suite spectrale
$''\mathscr E(\mathscr P_\bullet,\mathscr Q_\bullet)$ est birégulière (6.5.4),
et en vertu de (0, 11.1.8), il suffit (puisque dans les deux cas (i), (ii), $Z$
est localement noethérien) de prouver que les termes
$''\mathscr E_{pq}^2$ sont cohérents. L'hypothèse sur les
$\mathscr H_n(\mathscr P_\bullet)$ et
$\mathscr H_n(\mathscr Q_\bullet)$ et l'expression (6.5.4.2) des
$''\mathscr E_{pq}^2$ montrent donc que la proposition est équivalente à son
cas particulier correspondant à $\mathscr P_\bullet$ et
$\mathscr Q_\bullet$ réduits à leurs termes de degré 0, autrement dit à son
\end{proof}

\begin{corollary}[6.5.14]
Supposons vérifiée l'une des conditions (i), (ii) de (6.5.13), et soit
$\mathscr F$ (resp. $\mathscr G$) un $\mathscr O_X$-Module (resp. un
$\mathscr O_Y$-Module) quasi-cohérent de type fini; alors les
$\mathscr{T}\!or_n^S(\mathscr F,\mathscr G)$ sont des $\mathscr O_Z$-Modules
cohérents.
\end{corollary}

\begin{proof}
La question étant locale sur $X$ et $Y$, on peut supposer $S$, $X$, $Y$
affines.

\textnormal{(i)} Sous les hypothèses de (i), $S$, $X$ et $Z$ sont noethériens.
Il existe donc une résolution localement libre $\mathscr L_\bullet$ de
$\mathscr F$ formée de $\mathscr O_X$-Modules de type fini (I, 1.3.7); comme
$X$ est plat sur $S$, il résulte de (6.3.4) que l'on a
\[
  \mathscr{T}\!or_n^S(\mathscr F,\mathscr G)
  =\mathscr H_n(\mathscr L_\bullet\otimes_S\mathscr G);
\]
or, les $\mathscr L_i\otimes_S\mathscr G$ sont des $\mathscr O_Z$-Modules
quasi-cohérents de type fini (I, 9.1.1), donc cohérents. On en conclut que
$\mathscr H_n(\mathscr L_\bullet\otimes_S\mathscr G)$ est cohérent
($0_{\mathrm I}$, 5.3.4).

\textnormal{(ii)} Supposons maintenant vérifiées les conditions (ii). Comme
l'anneau $A(Y)$ est quotient d'une $A(S)$-algèbre de polynômes à un nombre fini
d'indéterminées (I, 6.3.3), $Y$ est un sous-$S$-préschéma fermé d'un
$S$-préschéma affine $Y'$, plat et de type fini sur $S$; $Y'$ étant noethérien
(I, 6.3.7), il existe une résolution localement libre $\mathscr M_\bullet$ de
$j_*(\mathscr G)$ par des $\mathscr O_{Y'}$-Modules de type fini
($j:Y\to Y'$ étant l'injection canonique); en vertu de (6.5.12),
$\mathscr{T}\!or_n^S(\mathscr F,\mathscr G)$ est l'image réciproque par
$1\times j$ du $\mathscr O_{Z'}$-Module
$\mathscr H_n(\mathscr F\otimes_S\mathscr M_\bullet)$ (où
$Z'=X\times_S Y'$); on voit comme dans (i) que les
$\mathscr F\otimes_S\mathscr M_i$ sont des $\mathscr O_{Z'}$-Modules
cohérents, et l'on en tire encore la conclusion par ($0_{\mathrm I}$, 5.3.4).
\end{proof}

\begin{env}[6.5.15]
La théorie développée ci-dessus pour deux complexes $\mathscr P_\bullet$,
$\mathscr Q_\bullet$ de faisceaux quasi-cohérents sur deux $S$-préschémas
$X$, $Y$ se généralise sans peine au cas où l'on considère un nombre fini
quelconque de $S$-préschémas $X^{(i)}$ ($1\leqslant i\leqslant m$) et sur
chaque $X^{(i)}$ un complexe $\mathscr P_\bullet^{(i)}$ de
$\mathscr O_{X^{(i)}}$-Modules quasi-cohérents; si
$Z=X^{(1)}\times_S X^{(2)}\times_S\cdots\times_S X^{(m)}$, on définit ainsi un
$\mathscr O_Z$-Module quasi-cohérent
$\mathscr{T}\!or_q^S(\mathscr P_\bullet^{(1)},\mathscr P_\bullet^{(2)},
\ldots,\mathscr P_\bullet^{(m)})$. Nous laissons au lecteur le soin de
développer la théorie dans ce cas général et nous nous bornerons à écrire,
pour référence ultérieure, le terme $E_2$ de la seconde suite spectrale
(régulière) dont l'aboutissement est
$\mathscr{T}\!or_\bullet^S(\mathscr P_\bullet^{(1)},
\mathscr P_\bullet^{(2)},\ldots,\mathscr P_\bullet^{(m)})$ :

\oldpage[III]{153}

\[
  ''\mathscr E_{pq}^2
  =\bigoplus_{q_1+q_2+\cdots+q_m=q}
   \mathscr{T}\!or_p^S
   \bigl(\mathscr H_{q_1}(\mathscr P_\bullet^{(1)}),\ldots,
         \mathscr H_{q_m}(\mathscr P_\bullet^{(m)})\bigr).
  \tag{6.5.15.1}\label{III.6.5.15.1.fr}
\]
Nous étudierons dans (6.8) les suites spectrales d'associativité auxquelles
donnent lieu ces foncteurs hypertor d'un nombre quelconque de complexes.
\end{env}

\subsection{Foncteurs hypertor globaux de complexes de Modules quasi-cohérents
et suites spectrales de Künneth : cas de la base affine}
\label{subsection:III.6.6.fr}

\begin{env}[6.6.1]
Considérons un schéma affine $S=\operatorname{Spec}(A)$ et deux $S$-schémas
quasi-compacts $X^{(i)}$ ($i=1,2$); soit $\mathscr P_\bullet^{(i)}$ un complexe
de $\mathscr O_{X^{(i)}}$-Modules quasi-cohérents, limité inférieurement, dont
l'opérateur de dérivation est de degré $-1$ ($i=1,2$). Considérons d'autre
part un recouvrement fini
$\mathfrak U^{(i)}=(U_\alpha^{(i)})$ de $X^{(i)}$ par des ouverts affines;
soit
\[
  X=X^{(1)}\times_S X^{(2)},
\]
qui est un $S$-schéma quasi-compact (I, 5.5.1 et 6.6.4), et soit
$\mathfrak U=\mathfrak U^{(1)}\times_S\mathfrak U^{(2)}$ le recouvrement de
$X$ formé des ouverts affines
$U_\alpha^{(1)}\times_S U_\beta^{(2)}$. Pour tout couple d'entiers
$p\leqslant0$, $q\in\mathbf Z$, le groupe des $(-p)$-cochaînes alternées
\[
  C^{-p}(\mathfrak U^{(i)},\mathscr P_q^{(i)})
\]
du recouvrement $\mathfrak U^{(i)}$ à coefficients dans le faisceau
$\mathscr P_q^{(i)}$ (G, II, 5.1) est un $A$-module; pour $p>0$, on posera
$C^p(\mathfrak U^{(i)},\mathscr P_q^{(i)})=0$; on a ainsi défini un
bicomplexe
\[
  C^\bullet(\mathfrak U^{(i)},\mathscr P_\bullet^{(i)})
\]
de $A$-modules, dont les deux opérateurs de dérivation sont de degré $-1$. Il
résulte des définitions ($0_{\mathrm I}$, 12.1.2) que le $A$-module
d'homologie
\[
  \mathrm H_n\bigl(C^\bullet
  (\mathfrak U^{(i)},\mathscr P_\bullet^{(i)})\bigr)
\]
de ce bicomplexe (considéré à l'ordinaire comme un complexe simple pour le
degré total) n'est autre que le $A$-module d'hypercohomologie
\[
  \mathbf H^{-n}
  (\mathfrak U^{(i)},\mathscr P^{(i)\bullet}),
\]
où $\mathscr P^{(i)\bullet}$ est le complexe à opérateur de dérivation de
degré $+1$ obtenu en prenant $\mathscr P_{-q}^{(i)}$ pour composante de degré
$q$; par abus de notation, nous l'écrirons
$\mathbf H^{-n}(\mathfrak U^{(i)},\mathscr P_\bullet^{(i)})$. Il résulte alors
de (6.2.2) que ce $A$-module est canoniquement isomorphe au $A$-module
d'hypercohomologie
$\mathbf H^{-n}(X^{(i)},\mathscr P^{(i)\bullet})$, que nous écrirons de même
$\mathbf H^{-n}(X^{(i)},\mathscr P_\bullet^{(i)})$; il ne dépend donc pas du
recouvrement fini $\mathfrak U^{(i)}$ choisi.
\end{env}

\begin{env}[6.6.2]
Nous allons appliquer aux deux bicomplexes de $A$-modules
\[
  L_{\bullet\bullet}^{(i)}
  =C^\bullet(\mathfrak U^{(i)},\mathscr P_\bullet^{(i)})
  \qquad(i=1,2)
\]
et au bifoncteur covariant
$L_{\bullet\bullet}^{(1)}\otimes_A L_{\bullet\bullet}^{(2)}$ en ces deux
bicomplexes, la théorie générale de l'hyperhomologie des foncteurs par rapport
aux bicomplexes ($0_{\mathrm I}$, 11.7.4). Comme les cochaînes considérées sont
alternées et les recouvrements $\mathfrak U^{(i)}$ finis, on notera que les
modules $C^{-p}(\mathfrak U^{(i)},\mathscr P_q^{(i)})$ ne sont $\ne0$ que pour
un nombre fini (indépendant de $q$) de valeurs de $p$, et en particulier les
deux degrés de chacun des $L_{\bullet\bullet}^{(i)}$ sont limités
inférieurement. Nous désignerons par
\[
  \mathbf{Tor}_n^A(L_{\bullet\bullet}^{(1)},L_{\bullet\bullet}^{(2)})
\]
ou
\[
  \mathbf{Tor}_n^S
  (\mathfrak U^{(1)},\mathfrak U^{(2)};
   \mathscr P_\bullet^{(1)},\mathscr P_\bullet^{(2)})
\]
le $n$-ième module d'hyperhomologie de
$L_{\bullet\bullet}^{(1)}\otimes_A L_{\bullet\bullet}^{(2)}$, que nous
appellerons l'\emph{hypertor} d'indice $n$ de
$\mathscr P_\bullet^{(1)}$ et $\mathscr P_\bullet^{(2)}$, relatif aux
recouvrements $\mathfrak U^{(1)}$ et $\mathfrak U^{(2)}$. Lorsque
$\mathscr P_\bullet^{(1)}$ et $\mathscr P_\bullet^{(2)}$ sont réduits à leurs
termes de degré 0, $\mathscr F^{(1)}$ et $\mathscr F^{(2)}$, on écrit
\[
  \mathbf{Tor}_n^S
  (\mathfrak U^{(1)},\mathfrak U^{(2)};
   \mathscr F^{(1)},\mathscr F^{(2)})
\]
leur hypertor. On désignera par
\[
  \mathbf{Tor}_n^S
  (\mathfrak U^{(1)},\mathfrak U^{(2)};
   \mathscr P_\bullet^{(1)},\mathscr P_\bullet^{(2)}),
\]
suivant les conventions générales, le bicomplexe dont le composant d'indices
$(j,k)$ est
\[
  \mathbf{Tor}_n^S
  (\mathfrak U^{(1)},\mathfrak U^{(2)};
   \mathscr P_j^{(1)},\mathscr P_k^{(2)}).
\]

Comme $L_{\bullet\bullet}^{(i)}$ est un foncteur exact en
$\mathscr P_\bullet^{(i)}$, puisque les intersections des ensembles
\oldpage[III]{154}
de $\mathfrak U^{(i)}$ sont affines (I, 5.5.6 et 1.3.11),
$\mathbf{Tor}_\bullet^S
(\mathfrak U^{(1)},\mathfrak U^{(2)};
 \mathscr P_\bullet^{(1)},\mathscr P_\bullet^{(2)})$ est un
bi-$\partial$-foncteur covariant en $\mathscr P_\bullet^{(1)}$,
$\mathscr P_\bullet^{(2)}$, à valeurs dans la catégorie des $A$-modules
($0_{\mathrm I}$, 11.7.3). En outre, on sait ($0_{\mathrm I}$, 11.7.2) que ce
bifoncteur est l'aboutissement commun de six foncteurs spectraux biréguliers,
que nous désignerons par la notation
\[
  {}^{(t)}\mathscr E^S
  (\mathfrak U^{(1)},\mathfrak U^{(2)};
   \mathscr P_\bullet^{(1)},\mathscr P_\bullet^{(2)})
\]
ou
\[
  {}^{(t)}\mathscr E
  (\mathfrak U^{(1)},\mathfrak U^{(2)};
   \mathscr P_\bullet^{(1)},\mathscr P_\bullet^{(2)}),
\]
où $t$ doit être remplacé par une des lettres $a$, $b$, $a'$, $b'$, $c$,
$d$, et dont les termes $E_2$ sont les suivants :
\[
\begin{aligned}
  {}^{(a)}\mathscr E_{pq}^2
  &=\bigoplus_{q_1+q_2=q}
    \mathbf{Tor}_p^A
    \bigl(\mathrm H_{q_1}(L_{\bullet\bullet}^{(1)}),
          \mathrm H_{q_2}(L_{\bullet\bullet}^{(2)})\bigr),\\
  {}^{(b)}\mathscr E_{pq}^2
  &=\mathrm H_p\bigl(
    \mathbf{Tor}_q^{A,\mathrm{II}}
    (L_{\bullet\bullet}^{(1)},L_{\bullet\bullet}^{(2)})\bigr),\\
  {}^{(a')}\mathscr E_{pq}^2
  &=\bigoplus_{q_1+q_2=q}
    \mathbf{Tor}_p^A
    \bigl(\mathrm H_{q_1}^{\mathrm I}(L_{\bullet\bullet}^{(1)}),
          \mathrm H_{q_2}^{\mathrm I}(L_{\bullet\bullet}^{(2)})\bigr),\\
  {}^{(b')}\mathscr E_{pq}^2
  &=\mathrm H_p\bigl(
    \mathbf{Tor}_q^A
    (L_{\bullet\bullet}^{(1)},L_{\bullet\bullet}^{(2)})\bigr),\\
  {}^{(c)}\mathscr E_{pq}^2
  &=\bigoplus_{q_1+q_2=q}
    \mathbf{Tor}_p^A
    \bigl(\mathrm H_{q_1}^{\mathrm{II}}(L_{\bullet\bullet}^{(1)}),
          \mathrm H_{q_2}^{\mathrm{II}}(L_{\bullet\bullet}^{(2)})\bigr),\\
  {}^{(d)}\mathscr E_{pq}^2
  &=\mathrm H_p\bigl(
    \mathbf{Tor}_q^{A,\mathrm I}
    (L_{\bullet\bullet}^{(1)},L_{\bullet\bullet}^{(2)})\bigr),
\end{aligned}
\]
où les notations sont conformes à celles de la théorie générale de
l'hyperhomologie. Nous allons dans ce qui suit expliciter davantage ces
termes initiaux.
\end{env}

\begin{env}[6.6.3]
\emph{Suites spectrales $(a)$ et $(a')$.} Nous avons vu en (6.6.1) que le
module d'homologie $\mathrm H_n(L_{\bullet\bullet}^{(i)})$ du bicomplexe
$L_{\bullet\bullet}^{(i)}$ était égal à
$\mathbf H^{-n}(X^{(i)},\mathscr P_\bullet^{(i)})$; donc
\[
  {}^{(a)}\mathscr E_{pq}^2
  =\bigoplus_{q_1+q_2=q}
   \mathbf{Tor}_p^A
   \bigl(\mathbf H^{-q_1}(X^{(1)},\mathscr P_\bullet^{(1)}),
         \mathbf H^{-q_2}(X^{(2)},\mathscr P_\bullet^{(2)})\bigr).
\]
Par définition, le complexe
$\mathrm H_n^{\mathrm I}(L_{\bullet\bullet}^{(i)})$ a pour terme de degré $k$
le module d'homologie
$\mathrm H_n(C^\bullet(\mathfrak U^{(i)},\mathscr P_k^{(i)}))$, c'est-à-dire,
par définition, le module de cohomologie
$\mathbf H^{-n}(\mathfrak U^{(i)},\mathscr P_k^{(i)})$; on sait (1.4.1) que
ce module est canoniquement isomorphe à
$\mathbf H^{-n}(X^{(i)},\mathscr P_k^{(i)})$; donc
\[
  {}^{(a')}\mathscr E_{pq}^2
  =\bigoplus_{q_1+q_2=q}
   \mathbf{Tor}_p^A
   \bigl(\mathbf H^{-q_1}(X^{(1)},\mathscr P_\bullet^{(1)}),
         \mathbf H^{-q_2}(X^{(2)},\mathscr P_\bullet^{(2)})\bigr).
\]
\end{env}

\begin{env}[6.6.4]
\emph{Suites spectrales $(b)$ et $(b')$.} Par définition,
$\mathbf{Tor}_q^{A,\mathrm{II}}
(L_{\bullet\bullet}^{(1)},L_{\bullet\bullet}^{(2)})$ est un bicomplexe dont
le terme de degré $(h,k)$ est le $A$-module
\[
  \mathbf{Tor}_q^A
  \bigl(C^{-h}(\mathfrak U^{(1)},\mathscr P_\bullet^{(1)}),
        C^{-k}(\mathfrak U^{(2)},\mathscr P_\bullet^{(2)})\bigr).
\]
Soit $\Phi^{(i)}$ l'ensemble d'indices de $\mathfrak U^{(i)}$; par
définition, le complexe de modules
$C^r(\mathfrak U^{(i)},\mathscr P_\bullet^{(i)})$ ($r\geqslant0$) est somme
directe des complexes
$\Gamma(U_\rho^{(i)},\mathscr P_\bullet^{(i)})$, où $U_\rho^{(i)}$ est
l'intersection des $U_\xi^{(i)}$ pour $\xi\in\rho$, et $\rho$ parcourt
$\mathfrak P(\Phi^{(i)})$; donc le $A$-module
\[
  \mathbf{Tor}_q^A
  \bigl(C^{-h}(\mathfrak U^{(1)},\mathscr P_\bullet^{(1)}),
        C^{-k}(\mathfrak U^{(2)},\mathscr P_\bullet^{(2)})\bigr)
\]
est somme directe des $A$-modules
\[
  \mathbf{Tor}_q^A
  \bigl(\Gamma(U_\sigma^{(1)},\mathscr P_\bullet^{(1)}),
        \Gamma(U_\tau^{(2)},\mathscr P_\bullet^{(2)})\bigr),
\]
où $\sigma$ (resp. $\tau$) parcourt les éléments de
$\mathfrak P(\Phi^{(1)})$ (resp. $\mathfrak P(\Phi^{(2)})$) tels que
$\operatorname{Card}(\sigma)=-(h+1)$ (resp.
$\operatorname{Card}(\tau)=-(k+1)$). Comme $X^{(1)}$ et $X^{(2)}$ sont des
schémas, les $U_\rho^{(i)}$ sont affines, donc on a (6.4.1.1)
\[
  \mathbf{Tor}_q^A
  \bigl(\Gamma(U_\sigma^{(1)},\mathscr P_\bullet^{(1)}),
        \Gamma(U_\tau^{(2)},\mathscr P_\bullet^{(2)})\bigr)
  =\Gamma\bigl(U_\sigma^{(1)}\times_S U_\tau^{(2)},
    \mathscr{T}\!or_q^S
    (\mathscr P_\bullet^{(1)},\mathscr P_\bullet^{(2)})\bigr).
\]
\oldpage[III]{155}
On voit donc que ${}^{(b)}\mathscr E_{pq}^2$ est le $(-p)$-ème module de
cohomologie du complexe
\[
  L^\bullet(\Phi^{(1)},\Phi^{(2)};\mathscr S)
\]
des cochaînes bi-alternées sur $\Phi^{(1)}$ et $\Phi^{(2)}$ à valeurs dans le
système de coefficients
\[
  \mathscr S:(\sigma,\tau)\leadsto
  \Gamma\bigl(U_\sigma^{(1)}\times_S U_\tau^{(2)},
    \mathscr{T}\!or_q^S
    (\mathscr P_\bullet^{(1)},\mathscr P_\bullet^{(2)})\bigr)
\]
($0_{\mathrm I}$, 11.8.4). On sait alors ($0_{\mathrm I}$, 11.8.5 et 11.8.6)
que la cohomologie de ce complexe est la même que celle du complexe
$C^\bullet(\Phi^{(1)},\Phi^{(2)};\mathscr S)$ de toutes les cochaînes sur
$\Phi^{(1)}$ et $\Phi^{(2)}$ à valeurs dans $\mathscr S$, et aussi que celle du
complexe $P^\bullet(\Phi^{(1)},\Phi^{(2)};\mathscr S)$, dont les éléments sont
les combinaisons linéaires des
\[
  \lambda(\sigma,\tau)\in
  \Gamma\bigl(U_\sigma^{(1)}\times_S U_\tau^{(2)},
    \mathscr{T}\!or_q^S
    (\mathscr P_\bullet^{(1)},\mathscr P_\bullet^{(2)})\bigr),
\]
où $\sigma=(\alpha_0,\ldots,\alpha_h)$ et
$\tau=(\beta_0,\ldots,\beta_h)$ sont des suites ayant même nombre d'éléments.
Mais on a alors
\[
  U_\sigma^{(1)}\times_S U_\tau^{(2)}
  =(U_{\alpha_0}^{(1)}\times_S U_{\beta_0}^{(2)})\cap\cdots\cap
    (U_{\alpha_h}^{(1)}\times_S U_{\beta_h}^{(2)})
  \qquad\text{(I, 3.2.7)}.
\]
Si l'on désigne par $\mathfrak U$ le recouvrement de
$Z=X^{(1)}\times_S X^{(2)}$ par les ouverts affines
$U_\alpha^{(1)}\times_S U_\beta^{(2)}$, on voit finalement, compte tenu de ce
que $X^{(1)}\times_S X^{(2)}$ est un schéma, que l'on a, en vertu de (1.3.1),
\[
  {}^{(b)}\mathscr E_{pq}^2
  =\mathbf H^{-p}\bigl(X^{(1)}\times_S X^{(2)},
    \mathscr{T}\!or_q^S
    (\mathscr P_\bullet^{(1)},\mathscr P_\bullet^{(2)})\bigr).
\]

En second lieu,
$\mathbf{Tor}_q^A(L_{\bullet\bullet}^{(1)},L_{\bullet\bullet}^{(2)})$ est un
bicomplexe dont le terme de degré $(h,k)$ est la somme directe des $A$-modules
\[
  \mathbf{Tor}_q^A
  \bigl(C^{-h_1}(\mathfrak U^{(1)},\mathscr P_{k_1}^{(1)}),
        C^{-h_2}(\mathfrak U^{(2)},\mathscr P_{k_2}^{(2)})\bigr)
\]
tels que $h_1+h_2=h$ et $k_1+k_2=k$; explicitant les modules
$C^r(\mathfrak U^{(i)},\mathscr P_s^{(i)})$ comme ci-dessus, on voit encore
que ce terme est somme directe des $A$-modules
\[
  \Gamma\bigl(U_\sigma^{(1)}\times_S U_\tau^{(2)},
    \mathscr{T}\!or_q^S
    (\mathscr P_{k_1}^{(1)},\mathscr P_{k_2}^{(2)})\bigr),
\]
où $k_1+k_2=k$, et $\sigma$ (resp. $\tau$) parcourt les éléments de
$\mathfrak P(\Phi^{(1)})$ (resp. $\mathfrak P(\Phi^{(2)})$) tels que
$\operatorname{Card}(\sigma)+\operatorname{Card}(\tau)=-h-2$. Le terme
${}^{(b')}\mathscr E_{pq}^2$ que nous calculons est le $(-p)$-ème module de
cohomologie d'un bicomplexe $N^{\bullet\bullet}=(N^{hk})$, où le complexe
simple $N^{\bullet k}$ est le complexe de cochaînes bi-alternées sur
$\Phi^{(1)}$ et $\Phi^{(2)}$, à valeurs dans le système de coefficients
\[
  \mathscr S_k:(\sigma,\tau)\leadsto
  \Gamma\left(U_\sigma^{(1)}\times_S U_\tau^{(2)},
    \bigoplus_{k_1+k_2=k}
    \mathscr{T}\!or_q^S
    (\mathscr P_{-k_1}^{(1)},\mathscr P_{-k_2}^{(2)})\right),
\]
ces systèmes de coefficients formant un complexe $\mathscr S^\bullet$ où la
différentielle provient de celle du complexe simple associé au bicomplexe
$\mathscr{T}\!or_q^S
(\mathscr P_\bullet^{(1)},\mathscr P_\bullet^{(2)})$. On sait que la
cohomologie de $N^{\bullet\bullet}$ est la même que celle du bicomplexe
$C^\bullet(\Phi^{(1)},\Phi^{(2)};\mathscr S^\bullet)$
($0_{\mathrm I}$, 11.8.9), et aussi la même que celle du bicomplexe
$P^\bullet(\Phi^{(1)},\Phi^{(2)};\mathscr S^\bullet)$, où les éléments de
degrés $(h,k)$ sont les combinaisons linéaires des
\[
  \lambda(\sigma,\tau)\in
  \Gamma\left(U_\sigma^{(1)}\times_S U_\tau^{(2)},
    \bigoplus_{k_1+k_2=k}
    \mathscr{T}\!or_q^S
    (\mathscr P_{-k_1}^{(1)},\mathscr P_{-k_2}^{(2)})\right),
\]
$\sigma=(\alpha_0,\ldots,\alpha_h)$,
$\tau=(\beta_0,\ldots,\beta_h)$ étant des suites ayant même nombre
d'éléments ($0_{\mathrm I}$, 11.8.10). On voit alors comme ci-dessus que
${}^{(b')}\mathscr E_{pq}^2$ est le $(-p)$-ème module de cohomologie du
bicomplexe $C^\bullet(\mathfrak U,\mathscr Q^\bullet)$, où
$\mathscr Q^\bullet$ est le complexe simple associé au bicomplexe
$\mathscr{T}\!or_q^S
(\mathscr P_\bullet^{(1)},\mathscr P_\bullet^{(2)})$ de
$\mathscr O_Z$-Modules. Avec les conventions faites dans (6.6.1), on a donc
\[
  {}^{(b')}\mathscr E_{pq}^2
  =\mathbf H^{-p}\bigl(X^{(1)}\times_S X^{(2)},
    \mathscr{T}\!or_q^S
    (\mathscr P_\bullet^{(1)},\mathscr P_\bullet^{(2)})\bigr).
\]
\end{env}

\oldpage[III]{156}
\begin{env}[6.6.5]
\emph{Suites spectrales $(c)$ et $(d)$.} Par définition, le complexe
$\mathrm H_n^{\mathrm{II}}(L_{\bullet\bullet}^{(i)})$ a pour terme de degré
$h$ le $A$-module
$C^{-h}(\mathfrak U^{(i)},\mathscr H_n(\mathscr P_\bullet^{(i)}))$, en vertu
de l'exactitude du foncteur $C^{-h}$. On a donc, par définition de l'hypertor
de deux Modules relatif à deux recouvrements (6.6.2)
\[
  {}^{(c)}\mathscr E_{pq}^2
  =\bigoplus_{q_1+q_2=q}
   \mathbf{Tor}_p^S
   \bigl(\mathfrak U^{(1)},\mathfrak U^{(2)};
         \mathscr H_{q_1}(\mathscr P_\bullet^{(1)}),
         \mathscr H_{q_2}(\mathscr P_\bullet^{(2)})\bigr).
\]
Enfin, par définition,
$\mathbf{Tor}_q^{A,\mathrm I}
(L_{\bullet\bullet}^{(1)},L_{\bullet\bullet}^{(2)})$ est un bicomplexe dont
le terme de degré $(h,k)$ est le $A$-module
$\mathbf{Tor}_q^S
(\mathfrak U^{(1)},\mathfrak U^{(2)};
 \mathscr P_h^{(1)},\mathscr P_k^{(2)})$. On a donc
\[
  {}^{(d)}\mathscr E_{pq}^2
  =\mathrm H_p\bigl(
    \mathbf{Tor}_q^S
    (\mathfrak U^{(1)},\mathfrak U^{(2)};
     \mathscr P_\bullet^{(1)},\mathscr P_\bullet^{(2)})\bigr).
\]
\end{env}

\begin{env}[6.6.6]
La théorie de l'hyperhomologie des foncteurs de bicomplexes
($0_{\mathrm I}$, 11.7.3) montre, comme dans (6.3.4), que, pour toute
résolution \emph{plate} de Cartan-Eilenberg $M_{\bullet\bullet}^{(i)}$ de
$L_{\bullet\bullet}^{(i)}$ (dans la catégorie des complexes de modules
limités inférieurement) ($i=1,2$), on a des isomorphismes canoniques de
bi-$\partial$-foncteurs
\[
\begin{aligned}
  \mathbf{Tor}_\bullet^S
  (\mathfrak U^{(1)},\mathfrak U^{(2)};
   \mathscr P_\bullet^{(1)},\mathscr P_\bullet^{(2)})
  &\xrightarrow{\sim}
   \mathrm H_\bullet
   (M_{\bullet\bullet}^{(1)}\otimes_A M_{\bullet\bullet}^{(2)})\\
  &\xrightarrow{\sim}
   \mathrm H_\bullet
   (M_{\bullet\bullet}^{(1)}\otimes_A L_{\bullet\bullet}^{(2)})
   \xrightarrow{\sim}
   \mathrm H_\bullet
   (L_{\bullet\bullet}^{(1)}\otimes_A M_{\bullet\bullet}^{(2)}).
\end{aligned}
\tag{6.6.6.1}\label{III.6.6.6.1.fr}
\]
\end{env}

\begin{env}[6.6.7]
Nous allons maintenant montrer que l'hypertor global défini dans (6.6.2), et
les six suites spectrales correspondantes, ne dépendent pas des recouvrements
ouverts affines finis $\mathfrak U^{(i)}$ qui ont servi à les définir (à des
isomorphismes canoniques près). Il suffira pour cela de montrer que si
$\mathfrak V^{(i)}$ sont deux autres recouvrements de même nature, tels que
$\mathfrak V^{(i)}$ soit \emph{plus fin} que $\mathfrak U^{(i)}$ pour
$i=1,2$, alors on a des \emph{isomorphismes canoniques} de foncteurs spectraux
\[
  {}^{(t)}\mathscr E
  (\mathfrak U^{(1)},\mathfrak U^{(2)};
   \mathscr P_\bullet^{(1)},\mathscr P_\bullet^{(2)})
  \xrightarrow{\sim}
  {}^{(t)}\mathscr E
  (\mathfrak V^{(1)},\mathfrak V^{(2)};
   \mathscr P_\bullet^{(1)},\mathscr P_\bullet^{(2)})
  \tag{6.6.7, $t$}\label{III.6.6.7.t.fr}
\]
où $t$ est remplacé par $a$, $b$, $a'$, $b'$, $c$ ou $d$.

Or, on a pour $i=1,2$ des homomorphismes de bicomplexes
\[
  C^\bullet(\mathfrak U^{(i)},\mathscr P_\bullet^{(i)})
  \longrightarrow
  C^\bullet(\mathfrak V^{(i)},\mathscr P_\bullet^{(i)})
\]
bien définis à homotopies près (G, II, 5.7.1); il en résulte déjà des
homomorphismes (6.6.7, $t$) canoniquement définis et compatibles avec les
opérateurs bords dans les aboutissements ($0_{\mathrm I}$, 11.3.2). En outre,
le calcul des termes $E_2$ des suites spectrales $(a)$, $(b)$, $(a')$, $(b')$,
montre que pour ces suites spectrales l'homomorphisme (6.6.7, $t$) est un
isomorphisme sur les termes $E_2$; comme ces suites spectrales sont
birégulières, on voit que (6.6.7, $t$) est un isomorphisme pour ces trois
foncteurs spectraux, donc un \emph{isomorphisme de bi-$\partial$-foncteurs}
pour leur aboutissement commun ($0_{\mathrm I}$, 11.1.5).

En particulier, pour des $\mathscr O_{X^{(i)}}$-Modules quasi-cohérents
$\mathscr F^{(i)}$ ($i=1,2$), l'homomorphisme canonique
\[
  \mathbf{Tor}_\bullet^S
  (\mathfrak U^{(1)},\mathfrak U^{(2)};
   \mathscr F^{(1)},\mathscr F^{(2)})
  \longrightarrow
  \mathbf{Tor}_\bullet^S
  (\mathfrak V^{(1)},\mathfrak V^{(2)};
   \mathscr F^{(1)},\mathscr F^{(2)})
\]
est bijectif; vu le calcul de (6.6.5), on voit que (6.6.7, $t$) est aussi un
isomorphisme des termes $E_2$ pour $t=c$ et $t=d$. On conclut comme ci-dessus
que (6.6.7, $t$) est aussi un isomorphisme de suites spectrales pour $t=c$ et
$t=d$.
\oldpage[III]{157}
On peut considérer que les isomorphismes (6.6.7, $t$) définissent des systèmes
inductifs de foncteurs spectraux sur l'ensemble filtrant des couples
$(\mathfrak U^{(1)},\mathfrak U^{(2)})$ de recouvrements ouverts affines finis
de $X^{(1)}$ et $X^{(2)}$. Nous désignerons par
\[
  {}^{(t)}\mathscr E
  (X^{(1)},X^{(2)};
   \mathscr P_\bullet^{(1)},\mathscr P_\bullet^{(2)})
  \quad\text{ou}\quad
  {}^{(t)}\mathscr E^S
  (X^{(1)},X^{(2)};
   \mathscr P_\bullet^{(1)},\mathscr P_\bullet^{(2)})
\]
la limite inductive de ce système, par
$\mathbf{Tor}_\bullet^S
(X^{(1)},X^{(2)};\mathscr P_\bullet^{(1)},\mathscr P_\bullet^{(2)})$
l'aboutissement de ce foncteur spectral, que nous appellerons l'\emph{hypertor
global} des deux complexes $\mathscr P_\bullet^{(1)}$ et
$\mathscr P_\bullet^{(2)}$; si $\mathscr P_\bullet^{(1)}$ et
$\mathscr P_\bullet^{(2)}$ sont réduits à leurs termes de degré 0,
$\mathscr F^{(1)}$ et $\mathscr F^{(2)}$, nous écrirons
\[
  \mathbf{Tor}_\bullet^S
  (X^{(1)},X^{(2)};\mathscr F^{(1)},\mathscr F^{(2)}),
\]
et conformément aux conventions générales,
\[
  \mathbf{Tor}_q^S
  (X^{(1)},X^{(2)};\mathscr P_\bullet^{(1)},\mathscr P_\bullet^{(2)})
\]
sera donc le bicomplexe des
\[
  \mathbf{Tor}_q^S
  (X^{(1)},X^{(2)};\mathscr P_h^{(1)},\mathscr P_k^{(2)}).
\]
\end{env}

\begin{env}[6.6.8]
Les hypothèses étant celles de (6.6.1), considérons maintenant deux
$S$-morphismes $f_i:X^{(i)}\to Y^{(i)}$, où
$Y^{(i)}=\operatorname{Spec}(B_i)$ est un $S$-schéma \emph{affine}, $B_i$
étant donc une $A$-algèbre ($i=1,2$); cela définit donc un $A$-homomorphisme
$B_i\to\Gamma(X^{(i)},\mathscr O_{X^{(i)}})$ (I, 2.2.4), et par suite chacun
des $L_{\bullet\bullet}^{(i)}$ définis dans (6.6.2) est un bicomplexe de
$B_i$-modules; on en conclut que
$L_{\bullet\bullet}^{(1)}\otimes_A L_{\bullet\bullet}^{(2)}$ est un
quadricomplexe de $(B_1\otimes_A B_2)$-modules, et ses six foncteurs spectraux
d'hyperhomologie peuvent donc être considérés comme prenant leurs valeurs
dans la catégorie des suites spectrales de $(B_1\otimes_A B_2)$-modules. Si
l'on pose
$Y=Y^{(1)}\times_S Y^{(2)}=\operatorname{Spec}(B_1\otimes_A B_2)$, on peut
considérer les $\mathscr O_Y$-Modules quasi-cohérents associés à ces modules
(I, 1.3.4); nous noterons
${}^{(t)}\mathscr E
(f_1,f_2;\mathscr P_\bullet^{(1)},\mathscr P_\bullet^{(2)})$
(pour $t=a$, $b$, $a'$, $b'$, $c$ ou $d$) les six suites spectrales
\[
  \bigl({}^{(t)}\mathscr E
  (X^{(1)},X^{(2)};
   \mathscr P_\bullet^{(1)},\mathscr P_\bullet^{(2)})\bigr)^{\sim}
\]
de $\mathscr O_Y$-Modules, et
$\mathfrak{Tor}_\bullet^S
(f_1,f_2;\mathscr P_\bullet^{(1)},\mathscr P_\bullet^{(2)})$ leur
aboutissement commun
\[
  \bigl(\mathbf{Tor}_\bullet^S
  (X^{(1)},X^{(2)};
   \mathscr P_\bullet^{(1)},\mathscr P_\bullet^{(2)})\bigr)^{\sim}.
\]
On le notera
$\mathfrak{Tor}_\bullet^S(f_1,f_2;\mathscr F^{(1)},\mathscr F^{(2)})$
lorsque $\mathscr P_\bullet^{(i)}$ est réduit à son terme de degré 0,
$\mathscr F^{(i)}$ ($i=1,2$).
\end{env}

\subsection{Foncteurs hypertor globaux de complexes de Modules quasi-cohérents
et suites spectrales de Künneth : cas général.}

\begin{env}[6.7.1]
Nous allons maintenant généraliser les définitions de (6.6.8) au cas où $S$
est un préschéma quelconque, $X^{(i)}$, $Y^{(i)}$ des $S$-préschémas et
$f_i:X^{(i)}\to Y^{(i)}$ des morphismes \emph{séparés et quasi-compacts}. Il
s'agit alors, pour tout couple de complexes $\mathscr P_\bullet^{(i)}$ de
$\mathscr O_{X^{(i)}}$-Modules quasi-cohérents, limités inférieurement
($i=1,2$), de définir pour tout $n$ un $\mathscr O_Y$-Module quasi-cohérent
$\mathfrak{Tor}_n^S
(f_1,f_2;\mathscr P_\bullet^{(1)},\mathscr P_\bullet^{(2)})$ ainsi que 6
foncteurs spectraux, se réduisant aux définitions de (6.6.8) lorsque $S$,
$Y^{(1)}$ et $Y^{(2)}$ sont affines (on pose
$Y=Y^{(1)}\times_S Y^{(2)}$).

Supposons d'abord $S=\operatorname{Spec}(A)$ \emph{affine}, mais $Y^{(1)}$
et $Y^{(2)}$ quelconques; soit $W^{(i)}$ un ouvert affine de $Y^{(i)}$;
$f_i^{-1}(W^{(i)})$ est alors un $S$-schéma quasi-compact,
$W=W^{(1)}\times_S W^{(2)}$ un ouvert affine de $Y$; soit
$f'_i:f_i^{-1}(W^{(i)})\to W^{(i)}$ la restriction de $f_i$, et
$\mathscr P'_\bullet{}^{(i)}$ la restriction
$\mathscr P_\bullet^{(i)}|f_i^{-1}(W^{(i)})$ ($i=1,2$). On a alors d'après
(6.6.8) les suites spectrales
${}^{(t)}\mathscr E
(f'_1,f'_2;\mathscr P'_\bullet{}^{(1)},\mathscr P'_\bullet{}^{(2)})$
de $(\mathscr O_Y|W)$-Modules quasi-cohérents, et il s'agit de vérifier
qu'elles satisfont aux conditions de recollement ($0_{\mathrm I}$, 3.3.1).
On est aussitôt ramené au cas où $Y^{(i)}=\operatorname{Spec}(B_i)$ est affine
et où $W^{(i)}=D(g_i)$, où $g_i\in B_i$, de sorte que
$W=D(g_1\otimes g_2)$
\oldpage[III]{158}
dans $Y=\operatorname{Spec}(B_1\otimes_A B_2)$ (II, 4.3.2.1); si
$X'^{(i)}=f_i^{-1}(W^{(i)})$, il s'agit d'établir un isomorphisme canonique
de foncteurs spectraux
\[
  {}^{(t)}\mathscr E
  (X'^{(1)},X'^{(2)};
   \mathscr P'_\bullet{}^{(1)},\mathscr P'_\bullet{}^{(2)})
  \xrightarrow{\sim}
  {}^{(t)}\mathscr E
  (X^{(1)},X^{(2)};
   \mathscr P_\bullet^{(1)},\mathscr P_\bullet^{(2)})
  \otimes_B B_g
  \tag{6.7.1.1}\label{III.6.7.1.1.fr}
\]
où l'on a posé $B=B_1\otimes_A B_2$ et $g=g_1\otimes g_2$. Pour cela,
partons de recouvrements ouverts affines finis $\mathfrak U^{(i)}$ de
$X^{(i)}$ ($i=1,2$), et soit $\mathfrak U'^{(i)}$ la trace de
$\mathfrak U^{(i)}$ sur $X'^{(i)}$, qui est encore formé d'ouverts affines
(I, 5.5.10); de façon précise, on a
\[
  C^\bullet(\mathfrak U'^{(i)},\mathscr P'_\bullet{}^{(i)})
  =C^\bullet(\mathfrak U^{(i)},\mathscr P_\bullet^{(i)})
   \otimes_{B_i}(B_i)_{g_i}.
\]
Si on pose
$L'_{\bullet\bullet}{}^{(i)}
=C^\bullet(\mathfrak U'^{(i)},\mathscr P'_\bullet{}^{(i)})$, on a donc
\[
  L'_{\bullet\bullet}{}^{(1)}\otimes_A L'_{\bullet\bullet}{}^{(2)}
  =\bigl(L_{\bullet\bullet}^{(1)}\otimes_{B_1}(B_1)_{g_1}\bigr)
   \otimes_A
   \bigl(L_{\bullet\bullet}^{(2)}\otimes_{B_2}(B_2)_{g_2}\bigr);
\]
comme on a $B_g=(B_1)_{g_1}\otimes_A(B_2)_{g_2}$ à un isomorphisme
canonique près, on a, à un isomorphisme canonique près,
\[
  L'_{\bullet\bullet}{}^{(1)}\otimes_A L'_{\bullet\bullet}{}^{(2)}
  =(L_{\bullet\bullet}^{(1)}\otimes_A L_{\bullet\bullet}^{(2)})
    \otimes_B B_g.
\]
Si $M_{\bullet\bullet\bullet}^{(i)}$ est une résolution projective de
Cartan-Eilenberg de $L_{\bullet\bullet}^{(i)}$, qu'on peut supposer formée de
$B_i$-modules, il résulte du fait que $(B_i)_{g_i}$ est \emph{plat} sur $B_i$
que
$M'_{\bullet\bullet\bullet}{}^{(i)}
=M_{\bullet\bullet\bullet}^{(i)}\otimes_{B_i}(B_i)_{g_i}$ est une résolution
projective de Cartan-Eilenberg du bicomplexe
$L'_{\bullet\bullet}{}^{(i)}$; en outre, on a
\[
  M'_{\bullet\bullet\bullet}{}^{(1)}\otimes_A M'_{\bullet\bullet\bullet}{}^{(2)}
  =(M_{\bullet\bullet\bullet}^{(1)}\otimes_A M_{\bullet\bullet\bullet}^{(2)})
    \otimes_B B_g.
\]
L'isomorphisme cherché (6.7.1.1) résulte alors aussitôt des définitions de
l'hyperhomologie d'un bicomplexe, et de l'exactitude du foncteur
$G\otimes_B B_g$ en le $B$-module $G$.
\end{env}

\begin{env}[6.7.2]
Supposons maintenant $S$ quelconque, et soient
$u_i:Y^{(i)}\to S$ les morphismes structuraux ($i=1,2$). Soit $(S_\alpha)$ un
recouvrement ouvert affine de $S$, posons
$Y_\alpha^{(i)}=u_i^{-1}(S_\alpha)$,
$X_\alpha^{(i)}=f_i^{-1}(Y_\alpha^{(i)})$, et soit
$f_{i\alpha}:X_\alpha^{(i)}\to Y_\alpha^{(i)}$ la restriction de $f_i$, qui
est un morphisme \emph{séparé et quasi-compact}. Les
$Y_\alpha=Y_\alpha^{(1)}\times_{S_\alpha}Y_\alpha^{(2)}$ forment un
recouvrement ouvert de $Y$, et sur chaque $Y_\alpha$ sont définis par (6.7.1)
des foncteurs spectraux
\[
  {}^{(t)}\mathscr E^{S_\alpha}
  \bigl(f_{1\alpha},f_{2\alpha};
    \mathscr P_\bullet^{(1)}|X_\alpha^{(1)},
    \mathscr P_\bullet^{(2)}|X_\alpha^{(2)}\bigr);
\]
il s'agit encore de montrer que ces foncteurs vérifient les conditions de
recollement.

On se ramène aussitôt à la situation suivante :
$S=\operatorname{Spec}(A)$ est affine, $S'=D(h)$, où $h\in A$, et
$u_i(Y^{(i)})\subset S'$; on peut en outre supposer
$Y^{(i)}=\operatorname{Spec}(B_i)$ affine; il s'agit de définir des
isomorphismes canoniques
\[
  {}^{(t)}\mathscr E^S
  (X^{(1)},X^{(2)};
   \mathscr P_\bullet^{(1)},\mathscr P_\bullet^{(2)})
  \xrightarrow{\sim}
  {}^{(t)}\mathscr E^{S'}
  (X^{(1)},X^{(2)};
   \mathscr P_\bullet^{(1)},\mathscr P_\bullet^{(2)}).
  \tag{6.7.2.1}\label{III.6.7.2.1.fr}
\]
Or, avec les notations de (6.6.2), les $L_{\bullet\bullet}^{(i)}$ sont formés
de $A_h$-modules, et l'on a donc
$L_{\bullet\bullet}^{(1)}\otimes_{A_h}L_{\bullet\bullet}^{(2)}
=L_{\bullet\bullet}^{(1)}\otimes_A L_{\bullet\bullet}^{(2)}$ à un
isomorphisme canonique près; comme on peut prendre une résolution projective
de Cartan-Eilenberg $M_{\bullet\bullet\bullet}^{(i)}$ de
$L_{\bullet\bullet}^{(i)}$ formée de $A_h$-modules, cela donne aussitôt
l'isomorphisme canonique cherché.
\end{env}

Nous avons en résumé démontré le

\begin{theorem}[6.7.3]
Soient $S$ un préschéma, $f_i:X_i\to Y_i$ un $S$-morphisme séparé et
quasi-compact de $S$-préschémas, $\mathscr P_\bullet^{(i)}$ un complexe de
$\mathscr O_{X^{(i)}}$-Modules quasi-cohérents limité inférieurement
($i=1,2$); on pose $Y=Y^{(1)}\times_S Y^{(2)}$. Il existe un
bi-$\partial$-foncteur
  $\mathfrak{Tor}_\bullet^S
(f_1,f_2;\mathscr P_\bullet^{(1)},\mathscr P_\bullet^{(2)})$
à valeurs dans la catégorie des $\mathscr O_Y$-Modules quasi-cohérents, tel
que si $V^{(i)}$ est un ouvert affine de $Y^{(i)}$ ($i=1,2$) et
$V=V^{(1)}\times_S V^{(2)}$, on ait
\[
  \mathfrak{Tor}_\bullet^S
  (f_1,f_2;\mathscr P_\bullet^{(1)},\mathscr P_\bullet^{(2)})|V
  =\Bigl(\mathbf{Tor}_\bullet^S
    \bigl(f_1^{-1}(V^{(1)}),f_2^{-1}(V^{(2)});
      \mathscr P_\bullet^{(1)}|f_1^{-1}(V^{(1)}),
      \mathscr P_\bullet^{(2)}|f_2^{-1}(V^{(2)})\bigr)\Bigr)^{\sim}.
\]
\oldpage[III]{159}
Ce bifoncteur est l'aboutissement de six foncteurs spectraux biréguliers
\[
  {}^{(t)}\mathscr E
  (f_1,f_2;\mathscr P_\bullet^{(1)},\mathscr P_\bullet^{(2)})
  \qquad(t=a,b,a',b',c,d)
\]
dont les termes $E_2$ sont donnés par
\[
\begin{aligned}
  {}^{(a)}\mathscr E_{pq}^2
  &=\bigoplus_{q_1+q_2=q}
    \mathscr{T}\!or_p^S
    \bigl(\mathscr H^{-q_1}(f_1,\mathscr P_\bullet^{(1)}),
          \mathscr H^{-q_2}(f_2,\mathscr P_\bullet^{(2)})\bigr),\\
  {}^{(b)}\mathscr E_{pq}^2
  &=\mathscr H^{-p}\bigl(f_1\times_S f_2,
    \mathfrak{Tor}_q^S
    (\mathscr P_\bullet^{(1)},\mathscr P_\bullet^{(2)})\bigr),\\
  {}^{(a')}\mathscr E_{pq}^2
  &=\bigoplus_{q_1+q_2=q}
    \mathfrak{Tor}_p^S
    \bigl(\mathscr H^{-q_1}(f_1,\mathscr P_\bullet^{(1)}),
          \mathscr H^{-q_2}(f_2,\mathscr P_\bullet^{(2)})\bigr),\\
  {}^{(b')}\mathscr E_{pq}^2
  &=\mathscr H^{-p}\bigl(f_1\times_S f_2,
    \mathscr{T}\!or_q^S
    (\mathscr P_\bullet^{(1)},\mathscr P_\bullet^{(2)})\bigr),\\
  {}^{(c)}\mathscr E_{pq}^2
  &=\bigoplus_{q_1+q_2=q}
    \mathfrak{Tor}_p^S
    \bigl(f_1,f_2;
      \mathscr H_{q_1}(\mathscr P_\bullet^{(1)}),
      \mathscr H_{q_2}(\mathscr P_\bullet^{(2)})\bigr),\\
  {}^{(d)}\mathscr E_{pq}^2
  &=\mathscr H_p\bigl(\mathfrak{Tor}_q^S
    (f_1,f_2;\mathscr P_\bullet^{(1)},\mathscr P_\bullet^{(2)})\bigr).
\end{aligned}
\]
On dit que les suites spectrales $(a)$ et $(b)$ sont les \emph{suites
spectrales de Künneth}.

On notera que les suites spectrales $(a)$ et $(a')$ (resp. $(b)$ et $(b')$)
sont identiques lorsque $\mathscr P_\bullet^{(1)}$ et
$\mathscr P_\bullet^{(2)}$ se réduisent à leurs termes de degré 0; dans ce cas,
les suites $(c)$ et $(d)$ sont dégénérées et sont donc sans intérêt.
\end{theorem}

\begin{remark}[6.7.4]
Les hypertor globaux que nous avons définis ci-dessus comprennent comme cas
particuliers, à la fois les \emph{Modules d'hypercohomologie} définis dans
(6.2.1) et les \emph{hypertor locaux} définis dans (6.5.3). Montrons que l'on
a, pour tout morphisme $f:X\to Y$ \emph{quasi-compact et séparé} et tout
complexe $\mathscr P_\bullet$ de $\mathscr O_X$-Modules quasi-cohérents,
limité inférieurement, un isomorphisme canonique de $\partial$-foncteurs en
$\mathscr P_\bullet$
\[
  \mathfrak{Tor}_n^Y(f,1_Y;\mathscr P_\bullet,\mathscr O_Y)
  \xrightarrow{\sim}
  \mathscr H^{-n}(f,\mathscr P_\bullet)
  \qquad\text{(pour tout $n\in\mathbf Z$)}.
  \tag{6.7.4.1}\label{III.6.7.4.1.fr}
\]
En effet, les méthodes de recollement de (6.7.2) ramènent aussitôt au cas où
$Y$ est affine; on peut alors, en vertu de (6.2.2), calculer les deux membres
de (6.7.4.1) à l'aide d'un même recouvrement fini $\mathfrak U$ de $X$ par des
ouverts affines, et (pour le premier membre) du recouvrement de $Y$ formé de
$Y$ lui-même; avec les notations de (6.6.2), le bicomplexe
$L_{\bullet\bullet}^{(2)}$ est alors réduit à son terme de degrés $(0,0)$,
égal à $A$, et la conclusion résulte de ($0_{\mathrm I}$, 11.7.5). Pour une
généralisation de ce résultat, voir (6.7.7); mais on notera que lorsque dans
le premier membre de (6.7.4.1), on remplace $\mathscr O_Y$ par un
$\mathscr O_Y$-Module quasi-cohérent \emph{quelconque} $\mathscr F$, on n'a
plus en général un isomorphisme avec
$\mathscr H^{-n}(f,\mathscr P_\bullet\otimes_{\mathscr O_Y}\mathscr F)$, bien
que, dans le calcul précédent, le bicomplexe
$L_{\bullet\bullet}^{(1)}\otimes_A L_{\bullet\bullet}^{(2)}$ s'identifie
encore au bicomplexe
$C^\bullet(\mathfrak U,\mathscr P_\bullet\otimes_Y\mathscr F)$.

D'autre part, on a un isomorphisme canonique de bi-$\partial$-foncteurs
\[
  \mathfrak{Tor}_\bullet^S
  (1_{X^{(1)}},1_{X^{(2)}};
   \mathscr P_\bullet^{(1)},\mathscr P_\bullet^{(2)})
  \xrightarrow{\sim}
  \mathscr{T}\!or_\bullet^S
  (\mathscr P_\bullet^{(1)},\mathscr P_\bullet^{(2)}).
  \tag{6.7.4.2}\label{III.6.7.4.2.fr}
\]
En effet, on se ramène encore, par (6.7.1) et (6.7.2), au cas où $S$ et les
$X^{(i)}$ sont affines; dans le calcul du premier membre de (6.7.4.2), on peut
alors prendre pour
\oldpage[III]{160}
recouvrement $\mathfrak U^{(i)}$ la famille réduite au seul élément $X^{(i)}$,
de sorte qu'avec les notations de (6.6.2), $L_{\bullet\bullet}^{(i)}$ se réduit
à $\Gamma(X^{(i)},\mathscr P_\bullet^{(i)})$ (considéré comme bicomplexe dont
les termes de premier degré $\ne0$ sont nuls), et l'égalité des deux membres
de (6.7.4.2) résulte de (6.4.1.1) et (6.3.1).
\end{remark}

\begin{proposition}[6.7.5]
Soit $u:\mathscr P_\bullet^{(1)}\to\mathscr Q_\bullet^{(1)}$ un
homomorphisme de complexes de $\mathscr O_{X^{(1)}}$-Modules quasi-cohérents,
limités inférieurement, tel que l'homomorphisme
\[
  \mathscr H_\bullet(u):
  \mathscr H_\bullet(\mathscr P_\bullet^{(1)})
  \longrightarrow
  \mathscr H_\bullet(\mathscr Q_\bullet^{(1)})
\]
déduit de $u$ soit un isomorphisme. Alors les homomorphismes
\[
  {}^{(t)}\mathscr E
  (f_1,f_2;\mathscr P_\bullet^{(1)},\mathscr P_\bullet^{(2)})
  \longrightarrow
  {}^{(t)}\mathscr E
  (f_1,f_2;\mathscr Q_\bullet^{(1)},\mathscr P_\bullet^{(2)})
\]
déduits de $u$ sont des isomorphismes pour $t=a$, $t=b$ et $t=c$.
\end{proposition}

\begin{proof}
L'assertion relative à la suite spectrale $(c)$ résulte de ce que cette suite
est birégulière et de ce que l'homomorphisme considéré est un isomorphisme pour
les termes $E_2$ par hypothèse ($0_{\mathrm I}$, 11.1.5). Ceci montre déjà que
\[
  \mathfrak{Tor}_\bullet^S
  (f_1,f_2;\mathscr P_\bullet^{(1)},\mathscr P_\bullet^{(2)})
  \longrightarrow
  \mathfrak{Tor}_\bullet^S
  (f_1,f_2;\mathscr Q_\bullet^{(1)},\mathscr P_\bullet^{(2)})
\]
est un isomorphisme. Appliquant les relations (6.7.4.1) et (6.7.4.2) on voit
d'abord que les homomorphismes
\[
  \mathscr H^{-n}(f_1,\mathscr P_\bullet^{(1)})
  \longrightarrow
  \mathscr H^{-n}(f_1,\mathscr Q_\bullet^{(1)})
  \quad\text{et}\quad
  \mathscr{T}\!or_n^S
  (\mathscr P_\bullet^{(1)},\mathscr P_\bullet^{(2)})
  \longrightarrow
  \mathscr{T}\!or_n^S
  (\mathscr Q_\bullet^{(1)},\mathscr P_\bullet^{(2)})
\]
déduits de $u$ sont des isomorphismes. L'assertion relative aux suites $(a)$
et $(b)$ résulte alors de ce que ces suites sont birégulières (6.7.3) et que
les homomorphismes considérés sont bijectifs pour les termes $E_2$
($0_{\mathrm I}$, 11.1.5).

Notons d'autre part que, si
$u:\mathscr P_\bullet^{(1)}\to\mathscr Q_\bullet^{(1)}$ est un
\emph{homotopisme}, on en déduit des isomorphismes canoniques
\[
  {}^{(t)}\mathscr E
  (f_1,f_2;\mathscr P_\bullet^{(1)},\mathscr P_\bullet^{(2)})
  \longrightarrow
  {}^{(t)}\mathscr E
  (f_1,f_2;\mathscr Q_\bullet^{(1)},\mathscr P_\bullet^{(2)})
\]
pour les six suites spectrales. En effet, si $S$ et les $Y^{(i)}$ sont affines,
on déduit de $u$ un homotopisme de bicomplexes
\[
  C^\bullet(\mathfrak U^{(1)},\mathscr P_\bullet^{(1)})
  \longrightarrow
  C^\bullet(\mathfrak U^{(1)},\mathscr Q_\bullet^{(1)}),
\]
et la proposition résulte de la théorie générale de l'hyperhomologie
($0_{\mathrm I}$, 11.3.2); le passage au cas général se fait par recollement,
en utilisant le fait que, d'un homotopisme de complexes, on déduit un
homotopisme de résolutions projectives de Cartan-Eilenberg de ces complexes
(M, XVII, 1.2).
\end{proof}

\begin{proposition}[6.7.6]
Supposons que le complexe $\mathscr P_\bullet^{(1)}$ ou le complexe
$\mathscr P_\bullet^{(2)}$ soit formé de Modules $S$-plats (les deux complexes
étant limités inférieurement). On a alors un isomorphisme canonique de
bi-$\partial$-foncteurs
\[
  \mathfrak{Tor}_n^S
  (f_1,f_2;\mathscr P_\bullet^{(1)},\mathscr P_\bullet^{(2)})
  \xrightarrow{\sim}
  \mathscr H^{-n}\bigl(f_1\times_S f_2,
    \mathscr P_\bullet^{(1)}\otimes_S\mathscr P_\bullet^{(2)}\bigr).
  \tag{6.7.6.1}\label{III.6.7.6.1.fr}
\]
\end{proposition}

\begin{proof}
Supposons d'abord $S$, $Y^{(1)}$ et $Y^{(2)}$ affines, de sorte qu'on est
dans la situation de (6.6.2), dont nous conservons les notations. Supposons par
exemple que $\mathscr P_\bullet^{(1)}$ soit formé de Modules $S$-plats, et
calculons l'hypertor en utilisant la remarque (6.6.6) : c'est donc l'homologie
de $L_{\bullet\bullet}^{(1)}\otimes_A M_{\bullet\bullet\bullet}^{(2)}$, où
$M_{\bullet\bullet\bullet}^{(2)}$ est une résolution projective de
Cartan-Eilenberg de $L_{\bullet\bullet}^{(2)}$, au sens de
($0_{\mathrm I}$, 11.7.1). D'autre part, les modules
$L_{ij}^{(1)}$ sont \emph{plats} sur $A$ en vertu de l'hypothèse (I, 4.15.1);
on déduit alors de ($0_{\mathrm I}$, 11.7.5) un isomorphisme canonique
\[
  \mathbf{Tor}_\bullet^S
  (\mathfrak U^{(1)},\mathfrak U^{(2)};
   \mathscr P_\bullet^{(1)},\mathscr P_\bullet^{(2)})
  \xrightarrow{\sim}
  \mathscr H_\bullet
  (L_{\bullet\bullet}^{(1)}\otimes_A L_{\bullet\bullet}^{(2)}).
  \tag{6.7.6.2}\label{III.6.7.6.2.fr}
\]
On a d'autre part un homomorphisme naturel de bicomplexes de
$L_{\bullet\bullet}^{(1)}\otimes_A L_{\bullet\bullet}^{(2)}$ dans
$C^\bullet(\mathfrak U,\mathscr Q_\bullet)$, où $\mathfrak U$ est le
recouvrement de $Z=X^{(1)}\times_S X^{(2)}$ par les ouverts affines
\oldpage[III]{161}
$U_\alpha^{(1)}\times_S U_\beta^{(2)}$ et
$\mathscr Q_\bullet
=\mathscr P_\bullet^{(1)}\otimes_S\mathscr P_\bullet^{(2)}$ (considéré comme
complexe simple pour le degré total); en effet, la définition de cet
homomorphisme a en substance été donnée au cours du calcul de la suite $(b')$
dans (6.6.4), pour $q=0$; il suffit simplement (en gardant les notations de
(6.6.4)) de tenir compte de ce qu'il y a d'une part un homomorphisme naturel
du complexe $N^{\bullet k}$ dans le complexe
$C^\bullet(\Phi^{(1)},\Phi^{(2)};\mathscr S_k)$
($0_{\mathrm I}$, 11.8.5), d'autre part un homomorphisme naturel de ce dernier
complexe dans le complexe
$P^\bullet(\Phi^{(1)},\Phi^{(2)};\mathscr S_k)$
($0_{\mathrm I}$, 11.8.6), et enfin un homomorphisme naturel de ce dernier
complexe de cochaînes dans le sous-complexe des cochaînes alternées
($0_{\mathrm I}$, 11.8.7). Par ailleurs, l'homomorphisme de bicomplexes ainsi
défini donne un isomorphisme en homologie, comme on l'a vu en (6.6.4); on a
donc, en composant avec (6.7.6.2), obtenu un isomorphisme
\[
  \mathbf{Tor}_n^S
  (\mathfrak U^{(1)},\mathfrak U^{(2)};
   \mathscr P_\bullet^{(1)},\mathscr P_\bullet^{(2)})
  \xrightarrow{\sim}
  \mathbf H^{-n}
  (\mathfrak U,\mathscr P_\bullet^{(1)}\otimes_S\mathscr P_\bullet^{(2)}).
  \tag{6.7.6.3}\label{III.6.7.6.3.fr}
\]
Il faut ensuite prouver que l'isomorphisme ainsi défini ne dépend pas des
recouvrements ouverts choisis (le second membre de (6.7.6.3) étant
canoniquement isomorphe à
$\mathbf H^{-n}
(X^{(1)}\times_S X^{(2)},
 \mathscr P_\bullet^{(1)}\otimes_S\mathscr P_\bullet^{(2)})$ par (6.2.2));
cela se fait à l'aide de (6.6.7) en remarquant (avec les notations de (6.6.7))
que l'on a un diagramme commutatif à homotopismes près
\[
\xymatrix@C=18pt@R=28pt{
  C^\bullet(\mathfrak U^{(1)},\mathscr P_\bullet^{(1)})
  \otimes_A
  C^\bullet(\mathfrak U^{(2)},\mathscr P_\bullet^{(2)})
  \ar[r]\ar[d]
  & C^\bullet(\mathfrak U,\mathscr Q_\bullet)\ar[d]
  \\
  C^\bullet(\mathfrak V^{(1)},\mathscr P_\bullet^{(1)})
  \otimes_A
  C^\bullet(\mathfrak V^{(2)},\mathscr P_\bullet^{(2)})
  \ar[r]
  & C^\bullet(\mathfrak V,\mathscr Q_\bullet)
}
\]
où les flèches horizontales sont les homomorphismes définis ci-dessus. Enfin,
il faut passer au cas général par recollement, ce qui se fait sans difficulté
comme dans (6.7.1) et (6.7.2); nous laissons les détails au lecteur.
\end{proof}

\begin{proposition}[6.7.7]
Supposons que $\mathscr P_\bullet^{(1)}$ et
$\mathscr P_\bullet^{(2)}$ soient limités inférieurement, et que tous les
Modules $\mathscr H^{-n}(f_1,\mathscr P_\bullet^{(1)})$ ou tous les Modules
$\mathscr H^{-n}(f_2,\mathscr P_\bullet^{(2)})$ soient $S$-plats. On a alors un
isomorphisme canonique de bi-$\partial$-foncteurs ($n$ parcourant $\mathbf Z$)
\[
  \mathfrak{Tor}_n^S
  (f_1,f_2;\mathscr P_\bullet^{(1)},\mathscr P_\bullet^{(2)})
  \xrightarrow{\sim}
  \bigoplus_{q_1+q_2=n}
  \mathscr H^{-q_1}(f_1,\mathscr P_\bullet^{(1)})
  \otimes_S
  \mathscr H^{-q_2}(f_2,\mathscr P_\bullet^{(2)}).
  \tag{6.7.7.1}\label{III.6.7.7.1.fr}
\]
\end{proposition}

\begin{proof}
En effet, vu (6.5.8), la suite spectrale $(a)$ de (6.7.3) est dégénérée, et
la proposition résulte aussitôt de ($0_{\mathrm I}$, 11.1.6), cette suite
étant birégulière (6.7.3).
\end{proof}

\begin{theorem}[6.7.8]
Supposons que : $1^\circ$ les complexes $\mathscr P_\bullet^{(1)}$ et
$\mathscr P_\bullet^{(2)}$ soient limités inférieurement; $2^\circ$ le
complexe $\mathscr P_\bullet^{(1)}$ ou le complexe
$\mathscr P_\bullet^{(2)}$ soit formé de Modules $S$-plats; $3^\circ$ tous les
\oldpage[III]{162}
Modules $\mathscr H^{-n}(f_1,\mathscr P_\bullet^{(1)})$ ou tous les Modules
$\mathscr H^{-n}(f_2,\mathscr P_\bullet^{(2)})$ soient $S$-plats. On a alors un
isomorphisme canonique de bi-$\partial$-foncteurs ($n$ parcourant $\mathbf Z$)
\[
  \mathscr H^n\bigl(f_1\times_S f_2,
    \mathscr P_\bullet^{(1)}\otimes_S\mathscr P_\bullet^{(2)}\bigr)
  \xrightarrow{\sim}
  \bigoplus_{n_1+n_2=n}
  \mathscr H^{n_1}(f_1,\mathscr P_\bullet^{(1)})
  \otimes_S
  \mathscr H^{n_2}(f_2,\mathscr P_\bullet^{(2)}).
  \tag{6.7.8.1}\label{III.6.7.8.1.fr}
\]
(« formule de Künneth »).
\end{theorem}

\begin{proof}
Cela résulte de (6.7.6) et (6.7.7).

Lorsque $S$, $Y^{(1)}$ et $Y^{(2)}$ sont affines, l'isomorphisme réciproque de
(6.7.8.1) se déduit (avec les notations de (6.7.6)) de l'homomorphisme de
bicomplexes
\[
  C^\bullet(\mathfrak U^{(1)},\mathscr P_\bullet^{(1)})
  \otimes_A
  C^\bullet(\mathfrak U^{(2)},\mathscr P_\bullet^{(2)})
  \longrightarrow
  C^\bullet(\mathfrak U,\mathscr Q_\bullet)
\]
par le procédé défini dans (G, I, 2.7), comme il résulte de (G, I, 5.5).
\end{proof}

\begin{proposition}[6.7.9]
Supposons vérifiées les trois conditions suivantes : $1^\circ$ $S$,
$Y^{(1)}$ et $Y^{(2)}$ sont localement noethériens, $f_1$ et $f_2$ sont
propres, $Y^{(1)}$ ou $Y^{(2)}$ de type fini sur $S$; $2^\circ$
$\mathscr P_\bullet^{(1)}$ et $\mathscr P_\bullet^{(2)}$ sont limités
inférieurement; $3^\circ$ pour tout $n\in\mathbf Z$,
$\mathscr H_n(\mathscr P_\bullet^{(i)})$ est un Module cohérent ($i=1,2$).
Dans ces conditions,
$\mathfrak{Tor}_n^S
(f_1,f_2;\mathscr P_\bullet^{(1)},\mathscr P_\bullet^{(2)})$ est un
$\mathscr O_Y$-Module cohérent (avec
$Y=Y^{(1)}\times_S Y^{(2)}$).
\end{proposition}

\begin{proof}
Il résulte de (6.5.13) que les hypertor locaux
$\mathscr{T}\!or_n^S(\mathscr P_\bullet^{(1)},\mathscr P_\bullet^{(2)})$ sont
des $\mathscr O_X$-Modules cohérents
($X=X^{(1)}\times_S X^{(2)}$ étant localement noethérien, car un des
$X^{(i)}$ est par hypothèse de type fini sur $S$ (I, 6.3.4 et 6.3.8)). Comme
$Y$ est localement noethérien et $f_1\times_S f_2$ propre (II, 5.4.2), il
résulte de (6.2.5) que les termes ${}^{(b)}\mathscr E_{pq}^2$ de (6.7.3) sont
des $\mathscr O_Y$-Modules cohérents. Comme toutes les suites spectrales de
(6.7.3) sont birégulières en vertu de l'hypothèse $2^\circ$, on conclut par
($0_{\mathrm I}$, 11.1.8).
\end{proof}

\begin{env}[6.7.10]
Soient maintenant $Y'^{(i)}$ deux $S$-préschémas,
$v_i:Y'^{(i)}\to Y^{(i)}$ deux $S$-morphismes ($i=1,2$),
$v=v_1\times_S v_2$ leur produit, qui est un $S$-morphisme $Y'\to Y$, où
l'on pose $Y'=Y'^{(1)}\times_S Y'^{(2)}$. Considérons d'autre part, pour
$i=1,2$, un $S$-préschéma $X'^{(i)}$, et deux $S$-morphismes
$u_i:X'^{(i)}\to X^{(i)}$, $f'_i:X'^{(i)}\to Y'^{(i)}$ de sorte que les
diagrammes
\[
\xymatrix@C=36pt@R=28pt{
  X'^{(i)}\ar[r]^{u_i}\ar[d]_{f'_i}
  & X^{(i)}\ar[d]^{f_i}
  \\
  Y'^{(i)}\ar[r]_{v_i}
  & Y^{(i)}
}
\tag{6.7.10.1}\label{III.6.7.10.1.fr}
\]
soient commutatifs, les morphismes $f'_i$ étant \emph{séparés et
quasi-compacts}. On a alors des $\mathscr O_{Y'}$-homomorphismes canoniques de
foncteurs spectraux
\[
  v^*\bigl({}^{(t)}\mathscr E
  (f_1,f_2;\mathscr P_\bullet^{(1)},\mathscr P_\bullet^{(2)})\bigr)
  \longrightarrow
  {}^{(t)}\mathscr E
  (f'_1,f'_2;u_1^*(\mathscr P_\bullet^{(1)}),
                 u_2^*(\mathscr P_\bullet^{(2)}))
  \tag{6.7.10.2}\label{III.6.7.10.2.fr}
\]
pour $t=a$, $a'$, $b$, $b'$, $c$, $d$. Pour les définir, supposons d'abord
$S=\operatorname{Spec}(A)$,
$Y^{(i)}=\operatorname{Spec}(B_i)$,
$Y'^{(i)}=\operatorname{Spec}(B'_i)$ affines; les $X^{(i)}$ et $X'^{(i)}$
sont alors des schémas quasi-compacts. Pour calculer les suites spectrales
${}^{(t)}\mathscr E
(f_1,f_2;\mathscr P_\bullet^{(1)},\mathscr P_\bullet^{(2)})$, nous
considérons comme dans (6.6.1) des recouvrements finis
$\mathfrak U^{(i)}$ par des ouverts affines de $X^{(i)}$ ($i=1,2$); pour
calculer
${}^{(t)}\mathscr E
(f'_1,f'_2;u_1^*(\mathscr P_\bullet^{(1)}),
u_2^*(\mathscr P_\bullet^{(2)}))$, nous considérons des recouvrements finis
$\mathfrak U'^{(i)}$

\oldpage[III]{163}
de $X'^{(i)}$ par des ouverts affines, \emph{plus fins} respectivement que les
recouvrements $u_i^{-1}(\mathfrak U^{(i)})$ ($i=1,2$). Il est clair que le
bicomplexe
$C^\bullet(\mathfrak U^{(i)},\mathscr P_\bullet^{(i)})
=L_{\bullet\bullet}^{(i)}$ peut être considéré canoniquement comme un
sous-bicomplexe de
$C^\bullet(u_i^{-1}(\mathfrak U^{(i)}),u_i^*(\mathscr P_\bullet^{(i)}))$
($0_{\mathrm I}$, 4.4.3.2); en outre, en choisissant une application
simpliciale (G, II, 5.7) de $\mathfrak U'^{(i)}$ dans
$u_i^{-1}(\mathfrak U^{(i)})$, on définit un homomorphisme de bicomplexes
\[
  C^\bullet(u_i^{-1}(\mathfrak U^{(i)}),u_i^*(\mathscr P_\bullet^{(i)}))
  \longrightarrow
  C^\bullet(\mathfrak U'^{(i)},u_i^*(\mathscr P_\bullet^{(i)}))
\]
d'où, par composition, un homomorphisme de bicomplexes
\[
  L_{\bullet\bullet}^{(i)}
  \longrightarrow
  L_{\bullet\bullet}^{\prime(i)}
  =C^\bullet(\mathfrak U'^{(i)},u_i^*(\mathscr P_\bullet^{(i)})).
\]
En outre, cet homomorphisme est remplacé par un homomorphisme homotope quand
on change d'application simpliciale (G, II, 5.7.1); on a ainsi un
homomorphisme bien défini de foncteurs spectraux :
\[
  {}^{(t)}\mathscr E
  (\mathfrak U^{(1)},\mathfrak U^{(2)};
   \mathscr P_\bullet^{(1)},\mathscr P_\bullet^{(2)})
  \longrightarrow
  {}^{(t)}\mathscr E
  (\mathfrak U'^{(1)},\mathfrak U'^{(2)};
   u_1^*(\mathscr P_\bullet^{(1)}),u_2^*(\mathscr P_\bullet^{(2)})).
  \tag{6.7.10.3}\label{III.6.7.10.3.fr}
\]
On vérifie aussitôt que si $\mathfrak V^{(i)}$ est un recouvrement affine
fini de $X^{(i)}$ plus fin que $\mathfrak U^{(i)}$,
$\mathfrak V'^{(i)}$ un recouvrement affine fini de $X'^{(i)}$, plus fin que
$u_i^{-1}(\mathfrak V^{(i)})$ et que $\mathfrak U'^{(i)}$, le diagramme
\[
\xymatrix@C=34pt@R=30pt{
  C^\bullet(\mathfrak U^{(i)},\mathscr P_\bullet^{(i)})
  \ar[r]\ar[d]
  & C^\bullet(\mathfrak V^{(i)},\mathscr P_\bullet^{(i)})\ar[d]
  \\
  C^\bullet(\mathfrak U'^{(i)},u_i^*(\mathscr P_\bullet^{(i)}))
  \ar[r]
  & C^\bullet(\mathfrak V'^{(i)},u_i^*(\mathscr P_\bullet^{(i)}))
}
\]
est commutatif, ce qui implique que l'homomorphisme (6.7.10.3) ne dépend pas
essentiellement des recouvrements $\mathfrak U^{(i)}$ et
$\mathfrak U'^{(i)}$ considérés. On a donc en fait défini un homomorphisme de
$A$-modules
\[
  {}^{(t)}E
  (X^{(1)},X^{(2)};\mathscr P_\bullet^{(1)},\mathscr P_\bullet^{(2)})
  \longrightarrow
  {}^{(t)}E
  (X'^{(1)},X'^{(2)};
   u_1^*(\mathscr P_\bullet^{(1)}),u_2^*(\mathscr P_\bullet^{(2)})).
  \tag{6.7.10.4}\label{III.6.7.10.4.fr}
\]
mais il est clair par définition des $u_i^*(\mathscr P_\bullet^{(i)})$ et en
vertu de la commutativité de (6.7.10.1) que cet homomorphisme est aussi un
homomorphisme de $(B_1\otimes_A B_2)$-modules; comme le second membre de
(6.7.10.4) est formé de $(B'_1\otimes_A B'_2)$-modules, on déduit
canoniquement de (6.7.10.4) un homomorphisme de
$(B'_1\otimes_A B'_2)$-modules
\[
  {}^{(t)}E
  (X^{(1)},X^{(2)};\mathscr P_\bullet^{(1)},\mathscr P_\bullet^{(2)})
  \otimes_{B_1\otimes_A B_2}(B'_1\otimes_A B'_2)
  \longrightarrow
  {}^{(t)}E
  (X'^{(1)},X'^{(2)};
   u_1^*(\mathscr P_\bullet^{(1)}),u_2^*(\mathscr P_\bullet^{(2)}))
  \tag{6.7.10.5}\label{III.6.7.10.5.fr}
\]
ce qui, compte tenu de (I, 1.6.5) n'est autre que l'homomorphisme cherché
(6.7.10.2) dans le cas particulier considéré.

Il reste à passer au cas général en suivant les recollements de (6.7.1) et
(6.7.2); le second passage est immédiat; en ce qui concerne le premier, on
considère comme dans (6.7.1) des éléments $g_i\in B_i$, et leurs images
$g'_i\in B'_i$, le produit tensoriel $g=g_1\otimes g_2$ dans
$B=B_1\otimes_A B_2$ et son image $g'=g'_1\otimes g'_2$ dans
$B'=B'_1\otimes_A B'_2$, et tout revient à utiliser

\oldpage[III]{164}
l'isomorphisme canonique
\[
  (M\otimes_B B')_{g'}
  \xrightarrow{\sim}
  M_g\otimes_{B_g}B'_{g'}
\]
($0_{\mathrm I}$, 1.5.4); nous laissons les détails au lecteur.
\end{env}

\begin{env}[6.7.11]
La théorie des hypertor globaux, développée ci-dessus pour deux
$S$-morphismes $X^{(i)}\to Y^{(i)}$ et deux complexes
$\mathscr P_\bullet^{(i)}$ de Modules quasi-cohérents limités inférieurement,
s'étend aussitôt au cas général suivant : on a un préschéma $S$, une famille
finie de $S$-préschémas $Y^{(i)}$ ($i\in I$), une famille finie de
$S$-morphismes séparés et quasi-compacts
$f_i:X^{(i)}\to Y^{(i)}$, et pour chaque $i$ un complexe de
$\mathscr O_{X^{(i)}}$-Modules quasi-cohérents
$\mathscr P_\bullet^{(i)}$ limités inférieurement. Si $Y$ est le produit des
$S$-préschémas $Y^{(i)}$, on définit alors pour chaque entier
$n\in\mathbf Z$, un $\mathscr O_Y$-Module quasi-cohérent
$\mathfrak{Tor}_n^S((f_i)_{i\in I};
(\mathscr P_\bullet^{(i)})_{i\in I})$,
ces Modules formant un $\partial$-foncteur covariant en chacun des complexes
$\mathscr P_\bullet^{(i)}$; en outre, ce foncteur est l'aboutissement commun
de six foncteurs spectraux
${}^{(t)}\mathscr E((f_i)_{i\in I};(\mathscr P_\bullet^{(i)})_{i\in I})$.
Nous laissons au lecteur le soin de répéter pour ce cas général les
définitions et les raisonnements faits ci-dessus pour $I=\{1,2\}$. Notons
simplement que lorsque $I$ se réduit à un seul élément, on retrouve
l'hypercohomologie $\mathscr H^\bullet(f,\mathscr P_\bullet)$ définie dans
(6.2.7) (comme on l'a déjà observé dans (6.7.4)). Lorsque $I$ est l'intervalle
$1\leq i\leq m$ de $\mathbf N$, nous écrirons
\[
  \mathfrak{Tor}_n^S
  (f_1,\ldots,f_m;
   \mathscr P_\bullet^{(1)},\ldots,\mathscr P_\bullet^{(m)})
  \quad\text{pour}\quad
  \mathfrak{Tor}_n^S
  ((f_i)_{i\in I};(\mathscr P_\bullet^{(i)})_{i\in I}).
\]
\end{env}

\begin{proposition}[6.7.12]
Les notations étant celles de (6.7.11), soit $J$ une partie de $I$ telle que,
pour $i\in I-J$, on ait $X^{(i)}=Y^{(i)}=S$, $f_i$ étant réduit à l'identité,
et $\mathscr P_\bullet^{(i)}$ égal au complexe réduit au terme de degré $0$
égal à $\mathscr O_S$. Il y a alors un isomorphisme canonique de
$\partial$-foncteurs
\[
  \mathfrak{Tor}_\bullet^S
  ((f_i)_{i\in I};(\mathscr P_\bullet^{(i)})_{i\in I})
  \xrightarrow{\sim}
  \mathfrak{Tor}_\bullet^S
  ((f_i)_{i\in J};(\mathscr P_\bullet^{(i)})_{i\in J}).
  \tag{6.7.12.1}\label{III.6.7.12.1.fr}
\]
\end{proposition}

\begin{proof}
On peut se borner à définir cet isomorphisme lorsque $S$ et les $Y^{(i)}$
sont affines, le recollement se faisant comme d'ordinaire. Pour $i\in I-J$,
on peut prendre le recouvrement $\mathfrak U^{(i)}$ formé du seul ensemble
$S$, et alors
$L_{\bullet\bullet}^{(i)}=C^\bullet(\mathfrak U^{(i)},
\mathscr P_\bullet^{(i)})$ est réduit à son seul terme de degrés $(0,0)$,
égal à $\Gamma(S,\mathscr O_S)=A(S)$; l'isomorphisme (6.7.12.1) est alors
évident.
\end{proof}

\begin{remark}[6.7.13]
Les notations étant celles de (6.7.3), considérons le $S$-isomorphisme
canonique
$Y^{(1)}\times_S Y^{(2)}\to Y^{(2)}\times_S Y^{(1)}$ (I, 3.3.5); alors
l'image par cet isomorphisme de
$\mathfrak{Tor}_\bullet^S(f_1,f_2;
\mathscr P_\bullet^{(1)},\mathscr P_\bullet^{(2)})$ est
$\mathfrak{Tor}_\bullet^S(f_2,f_1;
\mathscr P_\bullet^{(2)},\mathscr P_\bullet^{(1)})$; la question étant
locale, on est ramené au cas envisagé dans (6.6.2), et si on désigne par
$M_{\bullet\bullet\bullet}^{(i)}$ une résolution projective de
Cartan-Eilenberg de $L_{\bullet\bullet}^{(i)}$ ($i=1,2$), l'isomorphisme
considéré transforme
$M_{\bullet\bullet\bullet}^{(1)}\otimes_A
M_{\bullet\bullet\bullet}^{(2)}$ en
$M_{\bullet\bullet\bullet}^{(2)}\otimes_A
M_{\bullet\bullet\bullet}^{(1)}$, d'où notre assertion en considérant
l'homologie des complexes simples associés à ces tricomplexes.
\end{remark}

\subsection{Les suites spectrales d'associativité des hypertor globaux.}

\begin{env}[6.8.1]
Les hypothèses et notations étant celles de (6.7.11) (et en particulier les
$\mathscr P_\bullet^{(i)}$ étant supposés limités inférieurement, supposons
donnée une partition $(I_j)_{j\in J}$ de l'ensemble d'indices $I$; nous nous
proposons de donner une relation d'« associativité » entre les hypertor
$\mathfrak{Tor}_n^S((f_i)_{i\in I};(\mathscr P_\bullet^{(i)})_{i\in I})$ et
chacun des hypertor « partiels »
\[
  \mathcal T_{\bullet j}
  =\mathfrak{Tor}_\bullet^S
  ((f_i)_{i\in I_j};(\mathscr P_\bullet^{(i)})_{i\in I_j}).
\]

\oldpage[III]{165}
Pour simplifier l'écriture, nous nous bornerons au cas où $I$ est
l'intervalle $1\leq i\leq m$, et où la partition $(I_j)$ se compose des deux
intervalles $\{1,2,\ldots,r\}$ et $\{r+1,\ldots,m\}$.
\end{env}

\begin{proposition}[6.8.2]
Il existe un foncteur spectral canonique birégulier (dit « foncteur spectral
d'associativité ») noté
\[
  {}^{(e)}\mathscr E^S
  (f_1,\ldots,f_m;
   \mathscr P_\bullet^{(1)},\ldots,\mathscr P_\bullet^{(m)})
\]
(ou simplement ${}^{(e)}\mathscr E(f_1,\ldots,f_m;
\mathscr P_\bullet^{(1)},\ldots,\mathscr P_\bullet^{(m)})$) dont
l'aboutissement est
$\mathfrak{Tor}_\bullet^S(f_1,\ldots,f_m;
\mathscr P_\bullet^{(1)},\ldots,\mathscr P_\bullet^{(m)})$, et dont le terme
$E_2$ est donné par
\[
\begin{split}
  {}^{(e)}\mathscr E_{pq}^2
  =\bigoplus_{q_1+q_2=q}
  \mathscr{T}\!or_p^S\Bigl(
  &\mathfrak{Tor}_{q_1}^S
  (f_1,\ldots,f_r;
   \mathscr P_\bullet^{(1)},\ldots,\mathscr P_\bullet^{(r)}),\\
  &\mathfrak{Tor}_{q_2}^S
  (f_{r+1},\ldots,f_m;
   \mathscr P_\bullet^{(r+1)},\ldots,\mathscr P_\bullet^{(m)})
  \Bigr).
\end{split}
\]
\end{proposition}

\begin{proof}
Dans cet énoncé, on a identifié canoniquement $Y$ au produit
$Z^{(1)}\times_S Z^{(2)}$, où
$Z^{(1)}=Y^{(1)}\times_S Y^{(2)}\times_S\cdots\times_S Y^{(r)}$, et
$Z^{(2)}=Y^{(r+1)}\times_S\cdots\times_S Y^{(m)}$. Nous nous bornerons au
cas où $S$ et les $Y^{(i)}$ sont affines; on passe de ce cas particulier au
cas général par les méthodes développées dans (6.7.1) et (6.7.2), et nous
laissons les détails du raisonnement (sans difficulté) au lecteur. Nous
démontrerons donc le corollaire suivant.
\end{proof}

\begin{corollary}[6.8.3]
Soient $A$ un anneau, $S=\operatorname{Spec}(A)$, $X^{(i)}$
($1\leq i\leq m$) des $S$-schémas quasi-compacts et, pour chaque $i$, soit
$\mathscr P_\bullet^{(i)}$ un complexe de
$\mathscr O_{X^{(i)}}$-Modules quasi-cohérents limité inférieurement. Il
existe un foncteur spectral canonique birégulier ayant pour aboutissement
\[
  \mathbf{Tor}_\bullet^S
  (X^{(1)},\ldots,X^{(m)};
   \mathscr P_\bullet^{(1)},\ldots,\mathscr P_\bullet^{(m)})
\]
et dont le terme $E_2$ est donné par
\[
\begin{split}
  {}^{(e)}E_{pq}^2
  =\bigoplus_{q_1+q_2=q}
  \operatorname{Tor}_p^A\Bigl(
  &\mathbf{Tor}_{q_1}^S
  (X^{(1)},\ldots,X^{(r)};
   \mathscr P_\bullet^{(1)},\ldots,\mathscr P_\bullet^{(r)}),\\
  &\mathbf{Tor}_{q_2}^S
  (X^{(r+1)},\ldots,X^{(m)};
   \mathscr P_\bullet^{(r+1)},\ldots,\mathscr P_\bullet^{(m)})
  \Bigr).
\end{split}
\]
\end{corollary}

\begin{proof}
Suivant la définition donnée en (6.6.2), le calcul de l'hypertor considéré se
fait en prenant pour chaque $i$ un recouvrement ouvert affine fini
$\mathfrak U^{(i)}$ de $X^{(i)}$, en considérant les bicomplexes
$L_{\bullet\bullet}^{(i)}=C^\bullet(\mathfrak U^{(i)},
\mathscr P_\bullet^{(i)})$, une résolution projective
$M_{\bullet\bullet\bullet}^{(i)}$ de Cartan-Eilenberg de chacun de ces
bicomplexes (au sens de (0, 11.7.1)), le produit tensoriel
\[
  M_{\bullet\bullet\bullet}
  =\bigotimes_{i=1}^m M_{\bullet\bullet\bullet}^{(i)}
\]
de ces tricomplexes, et en prenant l'homologie de
$M_{\bullet\bullet\bullet}$. Considérons $M_{\bullet\bullet\bullet}$ comme
un complexe simple $N_\bullet$, produit tensoriel des deux complexes simples
\[
  N'_\bullet=\bigotimes_{i=1}^r M_{\bullet\bullet\bullet}^{(i)},
  \qquad
  N''_\bullet=\bigotimes_{i=r+1}^m M_{\bullet\bullet\bullet}^{(i)},
\]
où $N'_\bullet$ et $N''_\bullet$ sont gradués par la somme des degrés totaux
des $M_{\bullet\bullet\bullet}^{(i)}$. En outre, les $A$-modules des
complexes $N'_\bullet$ et $N''_\bullet$ sont projectifs, donc il résulte de
(6.5.9) que l'on a
\[
  H_\bullet(M_{\bullet\bullet\bullet})
  =\mathbf{Tor}_\bullet^A(N'_\bullet,N''_\bullet);
\]
la suite spectrale cherchée n'est autre alors que la suite (6.3.2.2)
appliquée aux complexes $N'_\bullet$ et $N''_\bullet$, compte tenu de
l'interprétation des modules d'homologie de ces complexes qui résulte de ce
qui précède (quand on applique les remarques du début à chacun des produits
partiels
$X^{(1)}\times_S\cdots\times_S X^{(r)}$ et
$X^{(r+1)}\times_S\cdots\times_S X^{(m)}$). Enfin, les propriétés de
régularité résultent de (6.3.2) et du fait que, les
$\mathscr P_\bullet^{(i)}$ étant limités inférieurement, il en est de même
des $M_{\bullet\bullet\bullet}^{(i)}$; par suite, $N'_\bullet$ et
$N''_\bullet$ sont limités inférieurement.
\end{proof}

\oldpage[III]{166}
\subsection{Les suites spectrales de changement de base dans les hypertor
globaux.}

\begin{env}[6.9.1]
Les hypothèses et notations étant toujours celles de (6.7.11) (et en
particulier les $\mathscr P_\bullet^{(i)}$ étant supposés limités
inférieurement), considérons un morphisme $g:S'\to S$ de préschémas, et
posons
$Y'^{(i)}=Y_{(S')}^{(i)}$, $X'^{(i)}=X_{(S')}^{(i)}$ et
$\mathscr P_\bullet'^{(i)}=
\mathscr P_\bullet^{(i)}\otimes_{\mathscr O_S}\mathscr O_{S'}$,
$\mathscr P_\bullet'^{(i)}$ étant donc un complexe de
$\mathscr O_{X'^{(i)}}$-Modules quasi-cohérents; soit
$f'_i=(f_i)_{(S')}:X'^{(i)}\to Y'^{(i)}$, qui est un morphisme séparé et
quasi-compact (I, 5.5.1 et 6.6.4). Nous nous proposons d'étudier les relations
entre les $\mathscr O_{Y'}$-Modules quasi-cohérents
$\mathfrak{Tor}_n^{S'}((f'_i)_{i\in I};
(\mathscr P_\bullet'^{(i)})_{i\in I})$ et
$\mathfrak{Tor}_n^S((f_i)_{i\in I};
(\mathscr P_\bullet^{(i)})_{i\in I})
\otimes_{\mathscr O_S}\mathscr O_{S'}$, où
$Y'=Y\times_S S'=Y_{(S')}$. Un cas particulièrement simple est le suivant,
qui se réduit à (1.4.15) lorsque $I$ se réduit à un seul élément et
$\mathscr P_\bullet$ à un seul module :
\end{env}

\begin{proposition}[6.9.2]
Si le morphisme $g:S'\to S$ est plat, on a un isomorphisme canonique de
$\partial$-foncteurs (en les $\mathscr P_\bullet^{(i)}$) :
\[
  \mathfrak{Tor}_\bullet^S
  ((f_i)_{i\in I};(\mathscr P_\bullet^{(i)})_{i\in I})
  \otimes_{\mathscr O_S}\mathscr O_{S'}
  \xrightarrow{\sim}
  \mathfrak{Tor}_\bullet^{S'}
  ((f'_i)_{i\in I};(\mathscr P_\bullet'^{(i)})_{i\in I}).
  \tag{6.9.2.1}\label{III.6.9.2.1.fr}
\]
\end{proposition}

\begin{proof}
On peut encore se borner au cas où $S$, $S'$ et les $Y^{(i)}$ sont affines,
le recollement se faisant suivant les méthodes de (6.7.1) et (6.7.2). Soient
$S=\operatorname{Spec}(A)$, $S'=\operatorname{Spec}(A')$, et prenons pour
chaque $i$ un recouvrement ouvert affine $\mathfrak U^{(i)}$ de $X^{(i)}$;
si $u_i:X'^{(i)}\to X^{(i)}$ est la projection canonique,
$u_i^{-1}(\mathfrak U^{(i)})$ est un recouvrement ouvert affine de
$X'^{(i)}$ (II, 1.5.5), que nous noterons $\mathfrak U'^{(i)}$; il est clair
alors que
\[
  C^\bullet(\mathfrak U'^{(i)},\mathscr P_\bullet'^{(i)})
  =C^\bullet(\mathfrak U^{(i)},\mathscr P_\bullet^{(i)})\otimes_A A',
\]
et l'existence de l'isomorphisme (6.9.2.1) est immédiate, car si
$M_{\bullet\bullet\bullet}^{(i)}$ est une résolution projective de
Cartan-Eilenberg de
$L_{\bullet\bullet}^{(i)}=C^\bullet(\mathfrak U^{(i)},
\mathscr P_\bullet^{(i)})$ au sens de (0, 11.7.1), formée de $A$-modules,
$M_{\bullet\bullet\bullet}^{(i)}\otimes_A A'$ est une résolution projective
de Cartan-Eilenberg (au même sens) de
$L_{\bullet\bullet}^{(i)}\otimes_A A'$ formée de $A'$-modules, en vertu de
l'hypothèse que $A'$ est un $A$-module plat; cette même hypothèse montre en
outre que
\[
  H_\bullet\left(
    \bigotimes_{i=1}^m
    (M_{\bullet\bullet\bullet}^{(i)}\otimes_A A')
  \right)
  =H_\bullet\left(
    \bigotimes_{i=1}^m M_{\bullet\bullet\bullet}^{(i)}
  \right)\otimes_A A'.
\]
\end{proof}

\begin{env}[6.9.2.2]
On notera que lorsque $I$ est réduit au seul élément $1$, la formule
(6.9.2.1) se déduit directement de (6.7.7), appliqué en prenant
$Y_2=X_2=S'$, $f_2=1_{S'}$, et le complexe $\mathscr P_\bullet^{(2)}$
réduit à son terme de degré $0$, égal à $\mathscr O_{S'}$; on sait alors que
l'hypercohomologie $\mathscr H^n(1_{S'},\mathscr O_{S'})$ est nulle pour tout
$n\ne0$ et se réduit à $\mathscr O_{S'}$ pour $n=0$ (6.2.1).
\end{env}

Dans le cas général, nous allons introduire à la place de $\mathscr O_{S'}$
un complexe $\mathscr Q'_\bullet$ de $\mathscr O_{S'}$-Modules
quasi-cohérents limités inférieurement, de sorte que si, pour simplifier, on
prend $I=\{1,2,\ldots,m\}$, on peut considérer le $\partial$-foncteur
\[
  \mathfrak{Tor}_\bullet^S
  (f_1,\ldots,f_m,1_{S'};
   \mathscr P_\bullet^{(1)},\ldots,\mathscr P_\bullet^{(m)},
   \mathscr Q'_\bullet).
\]

\begin{proposition}[6.9.3]
Il existe trois foncteurs spectraux canoniques biréguliers notés
${}^{(t)}\mathscr E(f_1,\ldots,f_m;
\mathscr P_\bullet^{(1)},\ldots,\mathscr P_\bullet^{(m)},
\mathscr Q'_\bullet)$ (avec $t=e$, $f$ ou $f'$) ayant pour aboutissement
commun
$\mathfrak{Tor}_\bullet^S(f_1,\ldots,f_m,1_{S'};
\mathscr P_\bullet^{(1)},\ldots,\mathscr P_\bullet^{(m)},
\mathscr Q'_\bullet)$ et dont les termes $E_2$ sont respectivement

\oldpage[III]{167}
\[
\begin{split}
  {}^{(e)}\mathscr E_{pq}^2
  =\bigoplus_{q'+q''=q}
  \mathscr{T}\!or_p^S\Bigl(
  &\mathfrak{Tor}_{q'}^S
  (f_1,\ldots,f_m;
   \mathscr P_\bullet^{(1)},\ldots,\mathscr P_\bullet^{(m)}),\\
  &\mathscr H_{q''}(\mathscr Q'_\bullet)
  \Bigr),
\end{split}
\]
\[
\begin{split}
  {}^{(f)}\mathscr E_{pq}^2
  =\bigoplus_{q_1+q_2+\cdots+q_{m+1}=q}
  \mathscr{T}\!or_p^{S'}\Bigl(
  &\mathfrak{Tor}_{q_1}^S
  (f_1,1_{S'};\mathscr P_\bullet^{(1)},\mathscr O_{S'}),\ldots,\\
  &\mathfrak{Tor}_{q_m}^S
  (f_m,1_{S'};\mathscr P_\bullet^{(m)},\mathscr O_{S'}),
  \mathscr H_{q_{m+1}}(\mathscr Q'_\bullet)
  \Bigr),
\end{split}
\]
\[
\begin{split}
  {}^{(f')}\mathscr E_{pq}^2
  =\bigoplus_{q_1+\cdots+q_m=q}
  \mathfrak{Tor}_p^{S'}\Bigl(
  &f'_1,\ldots,f'_m,1_{S'};\\
  &\mathscr{T}\!or_{q_1}^S
  (\mathscr P_\bullet^{(1)},\mathscr O_{S'}),\ldots,
  \mathscr{T}\!or_{q_m}^S
  (\mathscr P_\bullet^{(m)},\mathscr O_{S'}),
  \mathscr Q'_\bullet
  \Bigr).
\end{split}
\]
\end{proposition}

\begin{proof}
La suite $(e)$ n'est autre que la suite d'associativité de (6.8.2) pour
$r=m$. Pour définir les deux autres suites spectrales, on va encore se limiter
au cas où $S$, $S'$ et les $Y^{(i)}$ sont affines, le passage au cas général
se faisant par les méthodes de (6.7.1) et (6.7.2) et étant laissé au lecteur.
Nous démontrerons donc le
\end{proof}

\begin{corollary}[6.9.4]
Soient $A$ un anneau, $A'$ une $A$-algèbre,
$S=\operatorname{Spec}(A)$, $S'=\operatorname{Spec}(A')$,
$X^{(i)}$ ($1\leq i\leq m$) des $S$-schémas quasi-compacts et, pour chaque
$i$, soit $X'^{(i)}=X_{(S')}^{(i)}$, qui est un $S'$-schéma quasi-compact.
Pour chaque $i$, soit $\mathscr P_\bullet^{(i)}$ un complexe de
$\mathscr O_{X^{(i)}}$-Modules quasi-cohérents; soit enfin $Q'_\bullet$ un
complexe de $A'$-modules, ces complexes étant limités inférieurement. Il
existe trois foncteurs spectraux biréguliers en les
$\mathscr P_\bullet^{(i)}$ et en $Q'_\bullet$, ayant pour aboutissement commun
\[
  \mathbf{Tor}_\bullet^S
  (X^{(1)},\ldots,X^{(m)},S';
   \mathscr P_\bullet^{(1)},\ldots,\mathscr P_\bullet^{(m)},
   \widetilde{Q'_\bullet})
\]
et dont les termes $E_2$ sont respectivement
\[
\begin{split}
  {}^{(e)}E_{pq}^2
  =\bigoplus_{q'+q''=q}
  \operatorname{Tor}_p^A\Bigl(
  &\mathbf{Tor}_{q'}^S
  (X^{(1)},\ldots,X^{(m)};
   \mathscr P_\bullet^{(1)},\ldots,\mathscr P_\bullet^{(m)}),\\
  &H_{q''}(Q'_\bullet)
  \Bigr),
\end{split}
\]
\[
\begin{split}
  {}^{(f)}E_{pq}^2
  =\bigoplus_{q_1+q_2+\cdots+q_{m+1}=q}
  \operatorname{Tor}_p^{A'}\Bigl(
  &\mathbf{Tor}_{q_1}^S
  (X^{(1)},S';\mathscr P_\bullet^{(1)},\mathscr O_{S'}),\ldots,\\
  &\mathbf{Tor}_{q_m}^S
  (X^{(m)},S';\mathscr P_\bullet^{(m)},\mathscr O_{S'}),
  H_{q_{m+1}}(Q'_\bullet)
  \Bigr),
\end{split}
\]
\[
\begin{split}
  {}^{(f')}E_{pq}^2
  =\bigoplus_{q_1+\cdots+q_m=q}
  \mathbf{Tor}_p^{S'}\Bigl(
  &X'^{(1)},\ldots,X'^{(m)},S';\\
  &\mathscr{T}\!or_{q_1}^S
  (\mathscr P_\bullet^{(1)},\mathscr O_{S'}),\ldots,
  \mathscr{T}\!or_{q_m}^S
  (\mathscr P_\bullet^{(m)},\mathscr O_{S'}),
  \widetilde{Q'_\bullet}
  \Bigr).
\end{split}
\]
\end{corollary}

\begin{proof}
Nous ne reviendrons pas sur le premier de ces foncteurs spectraux, qui a été
traité dans (6.8.3) et n'est inclus ici que pour mémoire. Pour définir les
autres, considérons pour chaque $i$ un recouvrement ouvert affine fini
$\mathfrak U^{(i)}$ de $X^{(i)}$ et, si
$u_i:X'^{(i)}\to X^{(i)}$ est la projection canonique, le recouvrement ouvert
affine fini correspondant
$\mathfrak U'^{(i)}=u_i^{-1}(\mathfrak U^{(i)})$.

En vertu de (6.6.6),
$\mathbf{Tor}_\bullet^S(X^{(1)},\ldots,X^{(m)},S';
\mathscr P_\bullet^{(1)},\ldots,\mathscr P_\bullet^{(m)},
\widetilde{Q'_\bullet})$ s'obtient en prenant pour $1\leq i\leq m$ une
résolution projective de Cartan-Eilenberg
$M_{\bullet\bullet\bullet}^{(i)}$ de
$L_{\bullet\bullet}^{(i)}=C^\bullet(\mathfrak U^{(i)},
\mathscr P_\bullet^{(i)})$ (au sens de (0, 11.7.1)), considérant le
tricomplexe
\[
  M_{\bullet\bullet\bullet}
  =M_{\bullet\bullet\bullet}^{(1)}\otimes_A
   M_{\bullet\bullet\bullet}^{(2)}\otimes_A\cdots\otimes_A
   M_{\bullet\bullet\bullet}^{(m)}\otimes_A Q'_\bullet
\]
(où $Q'_\bullet$ est considéré comme un tricomplexe dont les deux derniers
degrés se réduisent à $0$), et en en prenant l'homologie. Si l'on pose
$M_{\bullet\bullet\bullet}'^{(i)}=
M_{\bullet\bullet\bullet}^{(i)}\otimes_A A'$, on a (en se rappelant que
$Q'_\bullet$ est un complexe de $A'$-modules)
\[
  M_{\bullet\bullet\bullet}
  =M_{\bullet\bullet\bullet}'^{(1)}\otimes_{A'}
   M_{\bullet\bullet\bullet}'^{(2)}\otimes_{A'}\cdots\otimes_{A'}
   M_{\bullet\bullet\bullet}'^{(m)}\otimes_{A'} Q'_\bullet.
\]
Or, considérons chacun des complexes
$M_{\bullet\bullet\bullet}'^{(i)}$ comme un complexe simple (pour son degré
total) et notons que ce complexe est formé de $A'$-modules projectifs; il
résulte de (6.3.7) (étendu à un nombre quelconque de complexes) que
$H_\bullet(M_{\bullet\bullet\bullet})$ est aussi égal à
\[
  \operatorname{Tor}_\bullet^{A'}
  (M_{\bullet\bullet\bullet}'^{(1)},\ldots,
   M_{\bullet\bullet\bullet}'^{(m)},Q'_\bullet);
\]
c'est donc (6.5.15) l'aboutissement d'une suite spectrale ayant les propriétés
de régularité voulues (les trois degrés de
$M_{\bullet\bullet\bullet}^{(i)}$ étant limités inférieurement lorsque
$\mathscr P_\bullet^{(i)}$ est limité inférieurement) et dont le terme $E_2$
est donné par
\[
  E_{pq}^2
  =\bigoplus_{q_1+\cdots+q_{m+1}=q}
  \operatorname{Tor}_p^{A'}
  (H_{q_1}(M_{\bullet\bullet\bullet}'^{(1)}),\ldots,
   H_{q_m}(M_{\bullet\bullet\bullet}'^{(m)}),H_{q_{m+1}}(Q'_\bullet)).
\]

\oldpage[III]{168}
On a d'ailleurs
\[
  H_{q_i}(M_{\bullet\bullet\bullet}'^{(i)})
  =H_{q_i}(M_{\bullet\bullet\bullet}^{(i)}\otimes_A A')
  =\mathbf{Tor}_{q_i}^S
  (X^{(i)},S';\mathscr P_\bullet^{(i)},\mathscr O_{S'})
\]
en vertu de la définition des hypertor globaux, ce qui donne la suite $(f)$
cherchée. On peut d'autre part considérer
$M_{\bullet\bullet\bullet}'^{(i)}$ comme un bicomplexe dans lequel le premier
degré est la somme du premier et du second degré du tricomplexe
$M_{\bullet\bullet\bullet}^{(i)}$, le second degré étant le troisième degré
de ce tricomplexe; comme les modules formant les
$M_{\bullet\bullet\bullet}'^{(i)}$ sont des $A'$-modules projectifs, la théorie
générale de l'hyperhomologie montre que l'homologie du bicomplexe
\[
  M_{\bullet\bullet\bullet}'^{(1)}\otimes_{A'}
  M_{\bullet\bullet\bullet}'^{(2)}\otimes_{A'}\cdots\otimes_{A'}
  M_{\bullet\bullet\bullet}'^{(m)}\otimes_{A'} Q'_\bullet
\]
est canoniquement isomorphe à son hyperhomologie (0, 11.6.5); c'est donc
l'aboutissement d'une suite spectrale de terme $E_2$ égal à
\[
  E_{pq}^2
  =\bigoplus_{q_1+\cdots+q_{m+1}=q}
  \operatorname{Tor}_p^{A'}
  (H_{q_1}^{\mathrm{II}}(M_{\bullet\bullet\bullet}'^{(1)}),\ldots,
   H_{q_m}^{\mathrm{II}}(M_{\bullet\bullet\bullet}'^{(m)}),
   H_{q_{m+1}}^{\mathrm{II}}(Q'_\bullet)).
\]
Or, comme le second degré de $Q'_\bullet$ se réduit à $0$, on a
$H_n^{\mathrm{II}}(Q'_\bullet)=0$ pour $n\ne0$ et
$H_0^{\mathrm{II}}(Q'_\bullet)=Q'_\bullet$; la formule précédente s'écrit
aussi
\[
  E_{pq}^2
  =\bigoplus_{q_1+\cdots+q_m=q}
  \operatorname{Tor}_p^{A'}
  (H_{q_1}^{\mathrm{II}}(M_{\bullet\bullet\bullet}'^{(1)}),\ldots,
   H_{q_m}^{\mathrm{II}}(M_{\bullet\bullet\bullet}'^{(m)}),Q'_\bullet).
\]
En outre, on a
\[
  H_{q_i}^{\mathrm{II}}(M_{\bullet\bullet\bullet}'^{(i)})
  =H_{q_i}^{\mathrm{II}}(M_{\bullet\bullet\bullet}^{(i)}\otimes_A A')
  =\operatorname{Tor}_{q_i}^A(L_{\bullet\bullet}^{(i)},A')
\]
en vertu de (6.3.4); mais
$L_{-j,k}^{(i)}=C^j(\mathfrak U^{(i)},\mathscr P_k^{(i)})$, somme directe des
$\Gamma(V,\mathscr P_k^{(i)})$, où $V$ parcourt les intersections (affines) de
$j+1$ ensembles du recouvrement $\mathfrak U^{(i)}$; si
$V'=u_i^{-1}(V)$, $V'$ est affine dans $X'^{(i)}$ et il résulte de (6.4.1.1)
que l'on a
\[
  \Gamma\bigl(V',\mathscr{T}\!or_{q_i}^S
  (\mathscr P_k^{(i)},\mathscr O_{S'})\bigr)
  =\operatorname{Tor}_{q_i}^A(\Gamma(V,\mathscr P_k^{(i)}),A')
\]
d'où pour le bicomplexe
$H_{q_i}^{\mathrm{II}}(M_{\bullet\bullet\bullet}'^{(i)})$ l'expression
\[
  C^\bullet\bigl(\mathfrak U'^{(i)},
  \mathscr{T}\!or_{q_i}^S(\mathscr P_\bullet^{(i)},\mathscr O_{S'})\bigr),
\]
ce qui donne finalement l'expression cherchée pour le terme $E_2$ de la
suite $(f')$. Le fait que cette suite soit birégulière sous les conditions
indiquées se vérifie comme d'ordinaire, tenant compte de ce que, si
$\mathscr P_\bullet^{(i)}$ est limité inférieurement, tous les degrés de
$M_{\bullet\bullet\bullet}'^{(i)}$ sont limités inférieurement.
\end{proof}

\begin{remark}[6.9.5]
On voit comme dans (6.7.6) que le remplacement des
$\mathscr P_\bullet^{(i)}$ et de $\mathscr Q'_\bullet$ par des complexes qui
leur sont respectivement homotopes ne change pas les suites $(e)$, $(f)$ et
$(f')$ à un isomorphisme canonique près. En outre, pour la suite $(f)$, des
homomorphismes
$\mathscr P_\bullet^{(i)}\to\mathscr R_\bullet^{(i)}$,
$\mathscr Q'_\bullet\to\mathscr T'_\bullet$ de complexes qui donnent des
isomorphismes en homologie
$\mathscr H_\bullet(\mathscr P_\bullet^{(i)})\xrightarrow{\sim}
\mathscr H_\bullet(\mathscr R_\bullet^{(i)})$,
$\mathscr H_\bullet(\mathscr Q'_\bullet)\xrightarrow{\sim}
\mathscr H_\bullet(\mathscr T'_\bullet)$ fournissent un isomorphisme de suites
spectrales
\[
\begin{split}
  {}^{(f)}\mathscr E
  (f_1,\ldots,f_m;
   \mathscr P_\bullet^{(1)},\ldots,\mathscr P_\bullet^{(m)},
   \mathscr Q'_\bullet)
  \xrightarrow{\sim}\;&
  {}^{(f)}\mathscr E
  (f_1,\ldots,f_m;\\
  &\mathscr R_\bullet^{(1)},\ldots,\mathscr R_\bullet^{(m)},
   \mathscr T'_\bullet);
\end{split}
\]
la démonstration est la même que pour (6.7.6) en tenant compte du résultat
de (6.7.6) et de la régularité de la suite $(f)$.
\end{remark}

\begin{corollary}[6.9.6]
Sous les conditions de (6.9.1), supposons que :
\begin{enumerate}
  \item[$1^\circ$] Les complexes $\mathscr P_\bullet^{(i)}$ sont formés de
  Modules plats sur $S$, et les $\mathscr O_Y$-Modules
  \[
    \mathfrak{Tor}_n^S
    (f_1,\ldots,f_m;
     \mathscr P_\bullet^{(1)},\ldots,\mathscr P_\bullet^{(m)})
  \]
  sont plats sur $S$.
  \item[$2^\circ$] Les $\mathscr P_\bullet^{(i)}$ et
  $\mathscr Q'_\bullet$ sont limités inférieurement.
\end{enumerate}

\oldpage[III]{169}
On a alors, en posant
$\mathscr P_\bullet'^{(i)}=
\mathscr P_\bullet^{(i)}\otimes_{\mathscr O_S}\mathscr O_{S'}$, des
isomorphismes canoniques fonctoriels
\[
\begin{split}
  &\mathfrak{Tor}_n^{S'}
  (f'_1,\ldots,f'_m,1_{S'};
   \mathscr P_\bullet'^{(1)},\ldots,
   \mathscr P_\bullet'^{(m)},\mathscr Q'_\bullet)
  \xrightarrow{\sim}\\
  &\qquad
  \bigoplus_{n'+n''=n}
  \mathfrak{Tor}_{n'}^S
  (f_1,\ldots,f_m;
   \mathscr P_\bullet^{(1)},\ldots,
   \mathscr P_\bullet^{(m)})
  \otimes_{\mathscr O_S}\mathscr H_{n''}(\mathscr Q'_\bullet).
\end{split}
\tag{6.9.6.1}\label{III.6.9.6.1.fr}
\]
En particulier, pour $\mathscr Q'_\bullet$ réduit à un seul terme
$\mathscr F'$ de degré $0$, on a des isomorphismes canoniques fonctoriels
\[
\begin{split}
  &\mathfrak{Tor}_n^{S'}
  (f'_1,\ldots,f'_m,1_{S'};
   \mathscr P_\bullet'^{(1)},\ldots,
   \mathscr P_\bullet'^{(m)},\mathscr F')
  \xrightarrow{\sim}\\
  &\qquad
  \mathfrak{Tor}_n^S
  (f_1,\ldots,f_m;
   \mathscr P_\bullet^{(1)},\ldots,
   \mathscr P_\bullet^{(m)})
  \otimes_{\mathscr O_S}\mathscr F'
\end{split}
\tag{6.9.6.2}\label{III.6.9.6.2.fr}
\]
et plus particulièrement, pour $\mathscr F'=\mathscr O_{S'}$,
\[
\begin{split}
  &\mathfrak{Tor}_n^{S'}
  (f'_1,\ldots,f'_m;
   \mathscr P_\bullet'^{(1)},\ldots,
   \mathscr P_\bullet'^{(m)})
  \xrightarrow{\sim}\\
  &\qquad
  \mathfrak{Tor}_n^S
  (f_1,\ldots,f_m;
   \mathscr P_\bullet^{(1)},\ldots,
   \mathscr P_\bullet^{(m)})
  \otimes_{\mathscr O_S}\mathscr O_{S'}.
\end{split}
\tag{6.9.6.3}\label{III.6.9.6.3.fr}
\]
\end{corollary}

\begin{proof}
L'hypothèse de platitude sur les Modules composant les
$\mathscr P_\bullet^{(i)}$ entraîne que les complexes
$\mathscr{T}\!or_{q_i}^S(\mathscr P_\bullet^{(i)},\mathscr O_{S'})$ sont nuls
pour $q_i\ne0$ (6.5.8). La suite $(f')$ est donc dégénérée; l'hypothèse
$2^\circ$ entraîne d'ailleurs qu'elle est birégulière (6.9.3), donc le
edge-homomorphisme
\[
\begin{split}
  &\mathfrak{Tor}_n^S
  (f_1,\ldots,f_m,1_{S'};
   \mathscr P_\bullet^{(1)},\ldots,
   \mathscr P_\bullet^{(m)},\mathscr Q'_\bullet)
  \longrightarrow {}^{(f')}\mathscr E_{n0}^2\\
  &\qquad=
  \mathfrak{Tor}_n^{S'}
  (f'_1,\ldots,f'_m,1_{S'};
   \mathscr P_\bullet'^{(1)},\ldots,
   \mathscr P_\bullet'^{(m)},\mathscr Q'_\bullet)
\end{split}
\tag{6.9.6.4}\label{III.6.9.6.4.fr}
\]
est bijectif (0, 11.1.6). L'hypothèse de platitude sur les Modules
\[
  \mathfrak{Tor}_n^S
  (f_1,\ldots,f_m;
   \mathscr P_\bullet^{(1)},\ldots,\mathscr P_\bullet^{(m)})
\]
entraîne que ${}^{(e)}\mathscr E_{pq}^2=0$ pour $p\ne0$ (6.5.8). La suite
$(e)$ est donc aussi dégénérée, et comme elle est birégulière, le
edge-homomorphisme
\[
\begin{split}
  &\mathfrak{Tor}_n^S
  (f_1,\ldots,f_m,1_{S'};
   \mathscr P_\bullet^{(1)},\ldots,
   \mathscr P_\bullet^{(m)},\mathscr Q'_\bullet)
  \longrightarrow {}^{(e)}\mathscr E_{0n}^2\\
  &\qquad=
  \bigoplus_{n'+n''=n}
  \mathfrak{Tor}_{n'}^S
  (f_1,\ldots,f_m;
   \mathscr P_\bullet^{(1)},\ldots,
   \mathscr P_\bullet^{(m)})
  \otimes_{\mathscr O_S}\mathscr H_{n''}(\mathscr Q'_\bullet)
\end{split}
\tag{6.9.6.5}\label{III.6.9.6.5.fr}
\]
est bijectif (0, 11.1.6); d'où, en combinant les deux isomorphismes
précédents, l'isomorphisme (6.9.6.1). L'isomorphisme (6.9.6.2) s'en déduit
trivialement, puisque l'on a alors
$\mathscr H_q(\mathscr Q'_\bullet)=0$ si $q\ne0$ et
$\mathscr H_0(\mathscr Q'_\bullet)=\mathscr F'$. Enfin, le cas
$\mathscr F'=\mathscr O_{S'}$ dans (6.9.6.2) donne l'isomorphisme (6.9.6.3),
compte tenu de (6.7.12).
\end{proof}

\begin{corollary}[6.9.7]
Sous les conditions de (6.9.1), supposons $S$ et $S'$ affines, et supposons
donné pour chaque $i$ un entier $d_i$ ($1\leq i\leq m$). Il existe alors un
entier $N$ ne dépendant que de $S$, des $X^{(i)}$ et des $d_i$, ayant la
propriété suivante : pour tout entier $n_0$, on a des isomorphismes canoniques
(6.9.6.3) pour $n\leq n_0$ et pour tout système de complexes
$\mathscr P_\bullet^{(i)}$ vérifiant les conditions suivantes :
\begin{enumerate}
  \item[$1^\circ$] $\mathscr P_k^{(i)}=0$ pour $k<d_i$;
  \item[$2^\circ$] $\mathscr P_k^{(i)}$ est plat sur $S$ pour $k<n_0+N$;
  \item[$3^\circ$]
  $\mathfrak{Tor}_q^S(f_1,\ldots,f_m;
  \mathscr P_\bullet^{(1)},\ldots,\mathscr P_\bullet^{(m)})$ est plat sur
  $S$ pour $q<n_0+N$.
\end{enumerate}
\end{corollary}

\begin{proof}
Supposons $\mathscr P_k^{(i)}$ plat sur $S$ pour $k<r$; alors
$\mathscr{T}\!or_{q_i}^S(\mathscr P_k^{(i)},\mathscr O_{S'})=0$ pour $k<r$
et $q_i\ne0$; calculons
\[
\begin{split}
  \mathfrak{Tor}_p^{S'}\Bigl(
  &f'_1,\ldots,f'_m;
  \mathscr{T}\!or_{q_1}^S(\mathscr P_\bullet^{(1)},\mathscr O_{S'}),
  \ldots,\\
  &\mathscr{T}\!or_{q_m}^S(\mathscr P_\bullet^{(m)},\mathscr O_{S'})
  \Bigr).
\end{split}
\tag{6.9.7.1}\label{III.6.9.7.1.fr}
\]

\oldpage[III]{170}
par la méthode de (6.6.2), à l'aide de l'image réciproque d'un recouvrement
affine fixe $\mathfrak U^{(i)}$ de $X^{(i)}$ ($1\leq i\leq m$)
(indépendant de $S'$ et des $\mathscr P_\bullet^{(i)}$); les termes
$L_{jk}^{(i)}$ sont nuls pour $j<-N_i$ (ne dépendant que de
$\mathfrak U^{(i)}$); si l'un des $q_i$ n'est pas nul, le complexe simple dont
l'homologie de degré $p$ est (6.9.7.1) a ses termes nuls pour tous les degrés
\[
  <r+\sum_{j\ne i}d_j-\sum_{i=1}^m N_i,
\]
donc (6.9.7.1) est nul pour $p<r-N$, en désignant par $N$ le plus grand des
nombres
\[
  \sum_{i=1}^m N_i-\sum_{j\ne i}d_j.
\]
On en conclut que l'on a ${}^{(f')}\mathscr E_{pq}^2=0$ pour $q\ne0$ et
$p<r-N$; comme d'autre part ${}^{(f')}\mathscr E_{pq}^2=0$ pour $q<0$, on
voit que le edge-homomorphisme (6.9.6.4) est bijectif pour $n<r-N$
(M, XV, 5.6) (pour $\mathscr Q'_\bullet=\mathscr O_{S'}$). En second lieu, si
\[
  \mathfrak{Tor}_q^S
  (f_1,\ldots,f_m;
   \mathscr P_\bullet^{(1)},\ldots,\mathscr P_\bullet^{(m)})
\]
est plat sur $S$ pour $q<r$, on a ${}^{(e)}\mathscr E_{pq}^2=0$ pour $p\ne0$
et $q<r$; par ailleurs ${}^{(e)}\mathscr E_{pq}^2=0$ pour $p<0$, donc le
edge-homomorphisme (6.9.6.5) est bijectif pour $n<r$, ce qui achève la
démonstration.
\end{proof}

Le cas le plus important de (6.9.3) dans les applications est celui où $m=1$,
$\mathscr Q'_\bullet$ étant réduit à un seul terme $\mathscr F'$ de degré
$0$; nous l'énoncerons à nouveau dans ce cas en vue de références
ultérieures\footnote{Le cas traité dans (6.9.8), et en particulier les suites
spectrales (6.9.8.3) et (6.9.8.4), nous avaient été signalés en 1957 par
J.-P. Serre.} :

\begin{proposition}[6.9.8]
Soient $S$ un préschéma, $g:S'\to S$ un morphisme,
$f:X\to Y$ un $S$-morphisme séparé et quasi-compact de $S$-préschémas,
$\mathscr P_\bullet$ un complexe de $\mathscr O_X$-Modules quasi-cohérents
limité inférieurement, $\mathscr F'$ un $\mathscr O_{S'}$-Module
quasi-cohérent. Il existe deux foncteurs spectraux biréguliers en
$\mathscr P_\bullet$ et $\mathscr F'$, à valeurs dans la catégorie des
$\mathscr O_{Y_{(S')}}$-Modules quasi-cohérents, ayant même aboutissement
$\mathfrak{Tor}_\bullet^S(f,1_{S'};\mathscr P_\bullet,\mathscr F')$, et dont
les termes $E_2$ sont
\[
  {}'\mathscr E_{pq}^2
  =\mathscr{T}\!or_p^S
  (\mathscr H^{-q}(f,\mathscr P_\bullet),\mathscr F')
  \tag{6.9.8.1}\label{III.6.9.8.1.fr}
\]
et
\[
  {}''\mathscr E_{pq}^2
  =\mathscr H^{-p}
  (f',\mathscr{T}\!or_q^S(\mathscr P_\bullet,\mathscr F')),
  \tag{6.9.8.2}\label{III.6.9.8.2.fr}
\]
où $f'=f_{(S')}:X_{(S')}\to Y_{(S')}$. 
\end{proposition}

\begin{proof}
Les suites en question peuvent aussi s'obtenir, non en partant de (6.9.3),
mais des suites $(a)$ et $(b')$ de (6.7.3) pour
$X^{(1)}=X$, $Y^{(1)}=Y$, $X^{(2)}=Y^{(2)}=S'$, $f_1=f$,
$f_2=1_{S'}$. Lorsque $S=S'=Y$, $Y$ étant affine, on obtient deux suites
spectrales de termes $E_2$ égaux à
\[
  {}'\mathscr E_{pq}^2
  =\mathscr{T}\!or_p^Y
  (\mathscr H^{-q}(f,\mathscr P_\bullet),\mathscr F)
  \tag{6.9.8.3}\label{III.6.9.8.3.fr}
\]
\[
  {}''\mathscr E_{pq}^2
  =\mathscr H^{-p}
  (f,\mathscr{T}\!or_q^Y(\mathscr P_\bullet,\mathscr F))
  \tag{6.9.8.4}\label{III.6.9.8.4.fr}
\]
aboutissant (en vertu de (6.7.6)) à l'hypercohomologie
$\mathscr H^\bullet(f,\mathscr P_\bullet\otimes_Y\mathscr F)$ du foncteur
$f_*$ par rapport au complexe
$\mathscr P_\bullet\otimes_Y\mathscr F$ de $\mathscr O_X$-Modules, pour tout
$\mathscr O_Y$-Module quasi-cohérent et $Y$-plat $\mathscr F$ (ou pour tout
$\mathscr O_Y$-Module quasi-cohérent $\mathscr F$ lorsque
$\mathscr P_\bullet$ est formé de $\mathscr O_X$-Modules $Y$-plats), qui sont
distinctes de celles de (6.2.1).
\end{proof}

\begin{corollary}[6.9.9]
Sous les conditions de (6.9.8), supposons que le complexe
$\mathscr P_\bullet$ soit limité inférieurement, formé de Modules plats sur
$S$, et que les $\mathscr O_Y$-Modules
$\mathscr H^n(f,\mathscr P_\bullet)$ soient plats sur $S$.

\oldpage[III]{171}
On a alors des isomorphismes canoniques fonctoriels
\[
  \mathfrak{Tor}_n^{S'}
  (f',1_{S'};\mathscr P'_\bullet,\mathscr F')
  \xrightarrow{\sim}
  \mathscr H^{-n}(f,\mathscr P_\bullet)
  \otimes_{\mathscr O_S}\mathscr F'
  \tag{6.9.9.1}\label{III.6.9.9.1.fr}
\]
où
$\mathscr P'_\bullet=\mathscr P_\bullet\otimes_{\mathscr O_S}
\mathscr O_{S'}$; en particulier, pour $\mathscr F'=\mathscr O_{S'}$, on a
des isomorphismes canoniques fonctoriels
\[
  \mathscr H^n(f',\mathscr P'_\bullet)
  \xrightarrow{\sim}
  \mathscr H^n(f,\mathscr P_\bullet)
  \otimes_{\mathscr O_S}\mathscr O_{S'}.
  \tag{6.9.9.2}\label{III.6.9.9.2.fr}
\]
C'est le cas particulier $m=1$ de (6.9.6). Plus particulièrement :
\end{corollary}

\begin{corollary}[6.9.10]
Soient $S$ un préschéma, $f:X\to Y$ un $S$-morphisme séparé et
quasi-compact de $S$-préschémas, $\mathscr P_\bullet$ un complexe limité
inférieurement, formé de $\mathscr O_X$-Modules quasi-cohérents, plats sur
$S$. On suppose en outre que les $\mathscr O_Y$-Modules
$\mathscr H^n(f,\mathscr P_\bullet)$ soient plats sur $S$. Pour tout $s\in S$,
notons $X_s$ et $Y_s$ les fibres
$X\otimes_S k(s)$, $Y\otimes_S k(s)$,
$f_s:X_s\to Y_s$ le morphisme $f\times_S1$,
$\mathscr P_\bullet^s$ le complexe
$\mathscr P_\bullet\otimes_{\mathscr O_S}k(s)$ de
$\mathscr O_{X_s}$-Modules. On a alors des isomorphismes canoniques
fonctoriels
\[
  \mathscr H^n(f_s,\mathscr P_\bullet^s)
  \xrightarrow{\sim}
  \mathscr H^n(f,\mathscr P_\bullet)\otimes_{\mathscr O_S}k(s).
  \tag{6.9.10.1}\label{III.6.9.10.1.fr}
\]
\end{corollary}

On a donc, moyennant des hypothèses de platitude convenables, un cas où la
formation des foncteurs dérivés $R^nf_*(\mathscr F)$ « commute au passage aux
fibres », que nous retrouverons par une autre méthode au §~7.

\subsection{Structure locale de certains foncteurs cohomologiques.}

\begin{proposition}[6.10.1]
Soient $S=\operatorname{Spec}(A)$ un schéma affine,
$Y^{(i)}$ ($1\leq i\leq n$) une famille finie de $S$-schémas affines, plats
sur $S$; pour chaque $i$, soit $f_i:X^{(i)}\to Y^{(i)}$ un $S$-morphisme
séparé et quasi-compact, et soit $\mathscr P_\bullet^{(i)}$ un complexe de
$\mathscr O_{X^{(i)}}$-Modules quasi-cohérents limité inférieurement. Soit
$Y$ le produit des $S$-schémas $Y^{(i)}$. Il existe un complexe
$\mathscr R_\bullet$ de $\mathscr O_Y$-Modules quasi-cohérents et plats sur
$S$, ayant la propriété suivante : pour tout $S$-schéma affine $S'$ et tout
complexe de $\mathscr O_{S'}$-Modules quasi-cohérents
$\mathscr Q'_\bullet$ limité inférieurement, il y a un isomorphisme
\[
  \mathbf{Tor}_\bullet^S
  (f_1,\ldots,f_n,1_{S'};
   \mathscr P_\bullet^{(1)},\ldots,\mathscr P_\bullet^{(n)},
   \mathscr Q'_\bullet)
  \xrightarrow{\sim}
  \mathscr H_\bullet
  (\mathscr R_\bullet\otimes_{\mathscr O_S}\mathscr Q'_\bullet)
  \tag{6.10.1.1}\label{III.6.10.1.1.fr}
\]
qui est un isomorphisme de $\partial$-foncteurs en $\mathscr Q'_\bullet$. En
outre, pour tout $S$-morphisme $u:S''\to S'$ de $S$-schémas affines, le
diagramme
\[
\xymatrix@C=12pt@R=30pt{
  {\begin{gathered}
    \mathbf{Tor}_\bullet^S
    (f_1,\ldots,f_n,1_{S'};\\
    \mathscr P_\bullet^{(1)},\ldots,\mathscr P_\bullet^{(n)},
    \mathscr Q'_\bullet)
  \end{gathered}}
  \ar[r]^-{\sim}\ar[d]
  &
  {\mathscr H_\bullet
    (\mathscr R_\bullet\otimes_{\mathscr O_S}\mathscr Q'_\bullet)}
  \ar[d]
  \\
  {\begin{gathered}
    \mathbf{Tor}_\bullet^S
    (f_1,\ldots,f_n,1_{S''};\\
    \mathscr P_\bullet^{(1)},\ldots,\mathscr P_\bullet^{(n)},
    u^*(\mathscr Q'_\bullet))
  \end{gathered}}
  \ar[r]^-{\sim}
  &
  {\mathscr H_\bullet
    (\mathscr R_\bullet\otimes_{\mathscr O_S}u^*(\mathscr Q'_\bullet))}
}
\tag{6.10.1.2}\label{III.6.10.1.2.fr}
\]
(où les flèches verticales sont les $(1_Y\times_Su)$-morphismes canoniques
définis dans (6.7.10)) est commutatif.
\end{proposition}

\begin{proof}
\oldpage[III]{172}
Calculons les hypertor par la méthode de (6.6.2), en tenant compte de la remarque
(6.6.6); avec les notations de (6.6.2), chaque
$L_{\bullet\bullet}^{(i)}$ est un bicomplexe de $A_i$-modules, en désignant par
$A_i$ l'anneau de $Y^{(i)}$; il admet donc une résolution de Cartan--Eilenberg
projective $M_{\bullet\bullet\bullet}^{(i)}$ (au sens de (0, 11.7.1)) formée de
$A_i$-modules, et en vertu de (6.6.6), le premier membre de (6.10.1.1) est
canoniquement isomorphe à
\[
  \mathscr H_\bullet
  (M_{\bullet\bullet\bullet}\otimes_A Q'_\bullet),
\]
où
\[
  M_{\bullet\bullet\bullet}
  =
  M_{\bullet\bullet\bullet}^{(1)}
  \otimes_A M_{\bullet\bullet\bullet}^{(2)}
  \otimes_A\cdots\otimes_A M_{\bullet\bullet\bullet}^{(n)}
\]
et $\mathscr Q'_\bullet=\widetilde{Q'_\bullet}$. Comme, par hypothèse, les
anneaux $A_i$ sont des $A$-modules plats, les
$M_{\bullet\bullet\bullet}^{(i)}$ sont des tricomplexes de $A$-modules plats
($0_{\mathrm I}$, 6.2.1), et il en est de même de
$M_{\bullet\bullet\bullet}$; en outre, si $B$ est l'anneau de $Y$, produit
tensoriel des $A_i$, $M_{\bullet\bullet\bullet}$ est un tricomplexe de
$B$-modules; le complexe
\[
  \mathscr R_\bullet=(M_{\bullet\bullet\bullet})^\sim
\]
de $\mathscr O_Y$-Modules (où $M_{\bullet\bullet\bullet}$ est considéré comme
complexe simple) répond donc à la question, comme il résulte aisément de
(6.7.10).
\end{proof}

\begin{corollary}[6.10.2]
Dans l'énoncé de (6.10.1), on peut supposer $\mathscr R_\bullet$ limité
inférieurement. Lorsque les $\mathscr P_\bullet^{(i)}$ sont limités
supérieurement et les $Y^{(i)}$ de dimension cohomologique finie, on peut
supposer $\mathscr R_\bullet$ limité supérieurement.
\end{corollary}

\begin{proof}
La première assertion résulte de ce que les trois degrés de chacun des
$M_{\bullet\bullet\bullet}^{(i)}$ sont limités inférieurement; d'autre part, si
les anneaux $A_i$ sont de dimension cohomologique finie, le troisième degré de
chacun des $M_{\bullet\bullet\bullet}^{(i)}$ ne prend qu'un nombre fini de
valeurs, et il en est de même par construction de son premier degré (6.6.2);
comme son second degré est limité supérieurement s'il en est ainsi du degré de
$\mathscr P_\bullet^{(i)}$ (6.6.2), la seconde assertion en résulte aussitôt.
\end{proof}

\begin{remarks}[6.10.3]
\textnormal{(i)} Avec les notations de (6.10.1),
$\mathscr H_\bullet(\mathscr R_\bullet\otimes_S\mathscr Q'_\bullet)$ est
isomorphe à
$\mathbf{Tor}_\bullet(\mathscr R_\bullet,\mathscr Q'_\bullet)$ puisque
$\mathscr R_\bullet$ est formé de $\mathscr O_Y$-Modules $S$-plats (6.5.9); il
est donc (6.5.4) l'aboutissement d'une suite spectrale régulière de terme
$E_2$ donné par
\[
  \mathscr E_{pq}^2
  =
  \bigoplus_{q'+q''=q}
  \mathbf{Tor}_p^S
  \bigl(
    \mathscr H_{q'}(\mathscr R_\bullet),
    \mathscr H_{q''}(\mathscr Q'_\bullet)
  \bigr)
  \tag{6.10.3.1}\label{III.6.10.3.1.fr}
\]
qui n'est autre que la suite spectrale $(e)$ du changement de base (6.9.3).

\medskip
\textnormal{(ii)} Soit $\mathscr R'_\bullet$ un second complexe de
$\mathscr O_Y$-Modules quasi-cohérents, plat sur $S$, et soit
$g:\mathscr R'_\bullet\to\mathscr R_\bullet$ un homomorphisme de complexes tel
que
\[
  \mathscr H_\bullet(g):
  \mathscr H_\bullet(\mathscr R'_\bullet)
  \longrightarrow
  \mathscr H_\bullet(\mathscr R_\bullet)
\]
soit un isomorphisme. Alors, en vertu de (6.3.3) et (6.5.9), on déduit de $g$
un isomorphisme de $\partial$-foncteurs en $\mathscr Q'_\bullet$
\[
  \mathscr H_\bullet
  (\mathscr R'_\bullet\otimes_S\mathscr Q'_\bullet)
  \xrightarrow{\sim}
  \mathscr H_\bullet
  (\mathscr R_\bullet\otimes_S\mathscr Q'_\bullet)
\]
tel que le diagramme
\[
\xymatrix@C=28pt@R=30pt{
  \mathscr H_\bullet
  (\mathscr R'_\bullet\otimes_S\mathscr Q'_\bullet)
  \ar[r]^-{\sim}\ar[d]
  &
  \mathscr H_\bullet
  (\mathscr R_\bullet\otimes_S\mathscr Q'_\bullet)
  \ar[d]
  \\
  \mathscr H_\bullet
  (\mathscr R'_\bullet\otimes_Su^*(\mathscr Q'_\bullet))
  \ar[r]^-{\sim}
  &
  \mathscr H_\bullet
  (\mathscr R_\bullet\otimes_Su^*(\mathscr Q'_\bullet))
}
\]
\par\medskip\noindent
soit commutatif. Cela prouve donc que le complexe $\mathscr R_\bullet$ n'est
pas entièrement déterminé par les propriétés de (6.10.1).

\medskip
\oldpage[III]{173}
\textnormal{(iii)} Dans la démonstration de (6.10.1), on peut supposer les
$M_{\bullet\bullet\bullet}^{(i)}$ formés de $A_i$-modules libres (comme il
résulte aisément de la démonstration de (0, 11.5.2.1) « dualisée »); les
\[
  M_{\bullet\bullet\bullet}^{(i)}\otimes_{A_i}B
\]
sont alors formés de $B$-modules libres, et comme
$M_{\bullet\bullet\bullet}$ est égal à leur produit tensoriel sur $B$, on voit
qu'on peut supposer en outre dans (6.10.1) que $\mathscr R_\bullet$ est associé
à un complexe de $B$-modules libres. En outre, en vertu de (M, XVII, 1.2), le
tricomplexe $M_{\bullet\bullet\bullet}$ dépend fonctoriellement de chacun des
bicomplexes $L_{\bullet\bullet}^{(i)}$ (donc de chacun des
$\mathscr P_\bullet^{(i)}$, lorsqu'on fixe un recouvrement fini de chacun des
$X^{(i)}$), les « morphismes » de tricomplexes devant être ici entendus comme
les classes d'homomorphismes pour la relation d'homotopie; d'ailleurs, le
remplacement d'un recouvrement de $X^{(i)}$ par un recouvrement plus fin
donnant lieu pour les $L_{\bullet\bullet}^{(i)}$ à des homomorphismes définis
précisément à homotopie près (6.6.8), on voit finalement qu'avec la convention
précédente pour les morphismes, le tricomplexe
$M_{\bullet\bullet\bullet}$ est un foncteur en chacun des
$\mathscr P_\bullet^{(i)}$. Nous préciserons cette dépendance fonctorielle, et
notamment le comportement de $\mathscr R_\bullet$ relativement à des suites
exactes
\[
  0\longrightarrow\mathscr P_\bullet^{\prime(i)}
  \longrightarrow\mathscr P_\bullet^{(i)}
  \longrightarrow\mathscr P_\bullet^{\prime\prime(i)}
  \longrightarrow0
\]
de complexes, dans le chapitre consacré à une algèbre générale de foncteurs
cohomologiques, mentionné dans (6.1.3).
\end{remarks}

\begin{scholium}[6.10.4]
Le fait que $\mathscr R_\bullet$ est formé de $\mathscr O_Y$-Modules
$S$-plats entraîne aisément que
\[
  \mathscr H_\bullet
  (\mathscr R_\bullet\otimes_S\mathscr Q'_\bullet)
\]
est un foncteur homologique en $\mathscr Q'_\bullet$ (voir le raisonnement de
(7.7.1)). C'est cette propriété qui, ainsi qu'on l'a mentionné en (6.1.1), est
la motivation de l'introduction de l'hypertor. Posons en effet
\[
  X=X^{(1)}\times_SX^{(2)}\times_S\cdots\times_SX^{(n)},
  \qquad
  f=f_1\times_Sf_2\times_S\cdots\times_Sf_n,
\]
\[
  \mathscr P_\bullet
  =
  \mathscr P_\bullet^{(1)}
  \otimes_S\mathscr P_\bullet^{(2)}
  \otimes_S\cdots\otimes_S\mathscr P_\bullet^{(n)},
\]
\[
  X'=X\times_SS',
  \qquad
  Y'=Y\times_SS',
  \qquad
  f'=f\times_S1_{S'}.
\]
Les problèmes de changement de base amènent à étudier l'hypercohomologie
\[
  \mathscr H^\bullet_{f'}
  (\mathscr P_\bullet\otimes_S\mathscr N')
\]
en tant que foncteur par rapport au $\mathscr O_{S'}$-Module quasi-cohérent
$\mathscr N'$, ou encore l'hypercohomologie
\[
  \mathscr H^\bullet_f
  (\mathscr P_\bullet\otimes_S\mathscr N)
\]
comme foncteur en le $\mathscr O_S$-Module quasi-cohérent $\mathscr N$. Lorsque
les $\mathscr P_\bullet^{(i)}$ (donc aussi $\mathscr P_\bullet$) sont
$S$-plats, il résulte de ce qui précède et de (6.7.6) que ce foncteur est bien
un foncteur cohomologique en $\mathscr N$; mais il n'en est plus de même
lorsqu'on ne fait plus l'hypothèse de platitude sur les
$\mathscr P_\bullet^{(i)}$, et l'on ne peut plus alors aborder l'étude de
\[
  \mathscr H^\bullet_f
  (\mathscr P_\bullet\otimes_S\mathscr N)
\]
par les méthodes usuelles de l'Algèbre homologique.

Nous aurons toutefois surtout à utiliser le cas où $n=1$, $Y=S$ et où
$\mathscr P_\bullet$ est formé de $\mathscr O_X$-Modules $Y$-plats. On a dans
ce cas le
\end{scholium}

\begin{theorem}[6.10.5]
Soient $Y=\operatorname{Spec}(A)$ un schéma affine noethérien,
$f:X\to Y$ un morphisme propre, $\mathscr P_\bullet$ un complexe de
$\mathscr O_X$-Modules cohérents, plats sur $Y$, limité inférieurement. Il
existe alors un complexe $\mathscr L_\bullet$ de $\mathscr O_Y$-Modules,
limité inférieurement, dont les termes $\mathscr L_i$ sont des
$\mathscr O_Y$-Modules de la forme $\mathscr O_Y^{m_i}$, et un isomorphisme
\[
  \mathscr H^\bullet
  (f,\mathscr P_\bullet\otimes_Y\mathscr Q_\bullet)
  \xrightarrow{\sim}
  \mathscr H_\bullet
  (\mathscr L_\bullet\otimes_Y\mathscr Q_\bullet)
  \tag{6.10.5.1}\label{III.6.10.5.1.fr}
\]
de $\partial$-foncteurs en le complexe $\mathscr Q_\bullet$ de
$\mathscr O_Y$-Modules quasi-cohérents, limité inférieurement. En outre, pour
tout morphisme $u:Y'\to Y$, on a, en posant
\[
  X'=X_{(Y')},
  \qquad
  f'=f_{(Y')},
  \qquad
  \mathscr P'_\bullet
  =
  \mathscr P_\bullet\otimes_Y\mathscr O_{Y'},
  \qquad
  \mathscr L'_\bullet=u^*(\mathscr L_\bullet),
\]
\oldpage[III]{174}
(qui est un complexe de $\mathscr O_{Y'}$-Modules localement libres de type
fini), un isomorphisme
\[
  \mathscr H^\bullet
  (f',\mathscr P'_\bullet\otimes_{Y'}\mathscr Q'_\bullet)
  \xrightarrow{\sim}
  \mathscr H_\bullet
  (\mathscr L'_\bullet\otimes_{Y'}\mathscr Q'_\bullet)
  \tag{6.10.5.2}\label{III.6.10.5.2.fr}
\]
de $\partial$-foncteurs en le complexe $\mathscr Q'_\bullet$ de
$\mathscr O_{Y'}$-Modules quasi-cohérents, limité inférieurement, de sorte que
le diagramme
\[
\xymatrix@C=25pt@R=30pt{
  \mathscr H^\bullet
  (f,\mathscr P_\bullet\otimes_Y\mathscr Q_\bullet)
  \ar[r]^-{\sim}\ar[d]
  &
  \mathscr H_\bullet
  (\mathscr L_\bullet\otimes_Y\mathscr Q_\bullet)
  \ar[d]
  \\
  \mathscr H^\bullet
  (f',\mathscr P'_\bullet\otimes_{Y'}u^*(\mathscr Q_\bullet))
  \ar[r]^-{\sim}
  &
  \mathscr H_\bullet
  (\mathscr L'_\bullet\otimes_{Y'}u^*(\mathscr Q_\bullet))
}
\tag{6.10.5.3}\label{III.6.10.5.3.fr}
\]
soit commutatif.
\end{theorem}

\begin{proof}
L'application de (6.10.1) donne d'abord un complexe limité inférieurement
(6.10.2) $\mathscr R_\bullet$ de $\mathscr O_Y$-Modules quasi-cohérents et
$Y$-libres (6.10.3, (iii)) et un isomorphisme
\[
  \mathscr H^\bullet
  (f,\mathscr P_\bullet\otimes_Y\mathscr Q_\bullet)
  \xrightarrow{\sim}
  \mathscr H_\bullet
  (\mathscr R_\bullet\otimes_Y\mathscr Q_\bullet)
  \tag{6.10.5.4}\label{III.6.10.5.4.fr}
\]
de $\partial$-foncteurs en $\mathscr Q_\bullet$, mais \emph{a priori} les
termes de $\mathscr R_\bullet$ ne sont pas nécessairement des
$\mathscr O_Y$-Modules de type fini. Mais si l'on applique (6.10.5.4) au cas
où $\mathscr Q_\bullet$ est un complexe réduit à un seul terme
$\mathscr O_Y$, on voit que $\mathscr H_\bullet(\mathscr R_\bullet)$ est
isomorphe à $\mathscr H^\bullet(f,\mathscr P_\bullet)$, et est par suite formé
de $\mathscr O_Y$-Modules cohérents (6.2.5). On sait alors (0, 11.9.2) qu'il
existe un complexe $\mathscr L_\bullet$ limité inférieurement, formé de
$\mathscr O_Y$-Modules associés à des $A$-modules libres de type fini, et un
homomorphisme
\[
  \mathscr L_\bullet\longrightarrow\mathscr R_\bullet,
\]
tels que l'homomorphisme correspondant pour l'homologie
\[
  \mathscr H_\bullet(\mathscr L_\bullet)
  \longrightarrow
  \mathscr H_\bullet(\mathscr R_\bullet)
\]
soit bijectif; d'où l'isomorphisme (6.10.5.1), en vertu de (6.10.3, (ii)).
Les autres assertions de (6.10.5) résultent de (6.10.1) et (6.10.3, (ii))
lorsque $Y'$ est affine; dans le cas général, il suffit de vérifier que
lorsqu'on considère un recouvrement $(V_\alpha)$ de $Y'$ par des ouverts
affines, et l'isomorphisme correspondant (6.10.5.2) relatif à chacun des
$V_\alpha$, les restrictions à un ouvert affine
$W\subset V_\alpha\cap V_\beta$ des isomorphismes correspondant à
$V_\alpha$ et à $V_\beta$ coïncident avec l'isomorphisme correspondant à
$W$, ce qui résulte de la commutativité du diagramme (6.10.1.2) appliqué aux
injections canoniques $W\to V_\alpha$ et $W\to V_\beta$.
\end{proof}

\begin{remark}[6.10.6]
Dans les chapitres suivants, nous appliquerons surtout (6.10.5) au cas où
$\mathscr P_\bullet$ est réduit à un seul $\mathscr O_X$-Module cohérent
$\mathscr F$, plat sur $Y$. Comme on a alors
\[
  \mathscr H_n(\mathscr L_\bullet)
  =
  \mathscr H^{-n}(f,\mathscr F)
  =
  R^{-n}f_*(\mathscr F)
  \qquad(6.2.1),
\]
on voit que les $\mathscr H_n(\mathscr L_\bullet)$ sont nuls pour $n>0$;
nous verrons plus loin (7.7.12, (i)) qu'on peut alors supposer que
$\mathscr L_\bullet$ n'a que des termes de degrés $\leqslant0$ (donc en
nombre fini), à condition de remplacer l'hypothèse que les $\mathscr L_i$
sont associés à des $A$-modules libres de type fini par celle que les
$\mathscr L_i$ sont localement libres de type fini.

Le complexe $\mathscr L_\bullet$ correspondant à un tel
$\mathscr O_X$-Module $\mathscr F$ ne paraît posséder aucune
\oldpage[III]{175}
propriété particulière, en dehors de la restriction précédente sur les
degrés. On peut alors se demander si inversement, étant donné un complexe
$\mathscr L_\bullet$ formé de $\mathscr O_Y$-Modules associés à des
$A$-modules projectifs de type fini, limité inférieurement et dont les termes
de degré $>0$ sont nuls, il existe un $Y$-schéma $X$, projectif et plat sur
$Y$ et un $\mathscr O_X$-Module localement libre $\mathscr F$, tels qu'il y
ait un isomorphisme
\[
  \mathscr H^\bullet
  (f,\mathscr F\otimes_Y\mathscr Q_\bullet)
  \xrightarrow{\sim}
  \mathscr H_\bullet
  (\mathscr L_\bullet\otimes_Y\mathscr Q_\bullet)
\]
fonctoriel en $\mathscr Q_\bullet$. L'intérêt d'un tel résultat serait de
réduire complètement la théorie cohomologique des Modules cohérents et
$Y$-plats sur des $Y$-schémas propres, à la théorie « à homotopie près » des
complexes de $A$-modules projectifs de type fini sur un anneau noethérien
$A$.
\end{remark}

% FIN SOURCE ega3/ega3-6-fr.tex

%
% DEBUT SOURCE ega3/ega3-7-fr.tex
\section{Étude du changement de base dans les foncteurs homologiques covariants de Modules}
\label{section:III.7.fr}

\subsection{Foncteurs de $A$-modules}
\label{subsection:III.7.1.fr}

\begin{env}[7.1.1]
Étant donné un anneau $A$ (non nécessairement commutatif), nous noterons
$\mathbf{Ab}_A$ la catégorie des $A$-modules à gauche, et nous noterons
simplement $\mathbf{Ab}$ la catégorie des $\mathbf Z$-modules, identiques aux
groupes commutatifs. Soit
\[
  T:\mathbf{Ab}_A\longrightarrow\mathbf{Ab}
\]
un foncteur covariant additif, et soit $M$ un $(A,A)$-bimodule; $T(M)$ est
alors muni de façon naturelle d'une structure de $A$-module à droite. En
effet, pour tout $a\in A$, notons $h_{a,M}$ (ou simplement $h_a$)
l'endomorphisme $x\mapsto xa$ du $A$-module à gauche $M$. Par hypothèse,
$T(h_a)$ est un endomorphisme du $\mathbf Z$-module $T(M)$; en outre, comme
$T$ est un foncteur covariant additif, on a, pour $a\in A$, $b\in A$,
\[
  T(h_{ab})
  =
  T(h_b\circ h_a)
  =
  T(h_b)\circ T(h_a),
  \qquad\text{et}\qquad
  T(h_{a+b})
  =
  T(h_a+h_b)
  =
  T(h_a)+T(h_b);
\]
cela prouve que l'application
\[
  (a,y)\longmapsto T(h_a)(y)
\]
est une loi externe de $A$-module à droite sur $T(M)$. En particulier,
$T(A_s)$ est un $A$-module à droite.
\end{env}

\begin{env}[7.1.2]
Lorsque $A$ est un anneau commutatif, il résulte de (7.1.1) que pour tout
$A$-module $M$, $T(M)$ est naturellement muni d'une structure de $A$-module;
en outre, si $u:M\to N$ est un homomorphisme de $A$-modules, on a, pour tout
$a\in A$,
\[
  u\circ h_{a,M}=h_{a,N}\circ u,
\]
d'où
\[
  T(u)\circ T(h_{a,M})
  =
  T(h_{a,N})\circ T(u),
\]
ce qui prouve que $T(u):T(M)\to T(N)$ est un homomorphisme de $A$-modules; on
voit donc que $T$ peut être considéré comme un foncteur covariant additif de
la catégorie $\mathbf{Ab}_A$ dans elle-même. De façon précise, on a ainsi
défini une équivalence canonique entre la catégorie des foncteurs covariants
additifs $\mathbf{Ab}_A\to\mathbf{Ab}$ et la catégorie des foncteurs
covariants $A$-linéaires
\[
  T:\mathbf{Ab}_A\longrightarrow\mathbf{Ab}_A,
\]
c'est-à-dire tels que
\[
  T(h_{a,M})=h_{a,T(M)}
  \qquad\text{pour tout }a\in A.
\]
Comme le foncteur d'inclusion
\[
  I:\mathbf{Ab}_A\longrightarrow\mathbf{Ab},
\]
faisant correspondre à tout $A$-module le $\mathbf Z$-module sous-jacent, est
exact et fidèle, les propriétés d'exactitude des deux foncteurs associés par
l'équivalence précédente sont les mêmes.
\end{env}

\begin{env}[7.1.3]
L'anneau $A$ étant toujours supposé commutatif, soit $B$ une $A$-algèbre (non
nécessairement commutative), et soit $\rho:A\to B$ l'homomorphisme d'anneaux
correspondant à cette structure d'algèbre; cet homomorphisme définit un
foncteur covariant
\oldpage[III]{176}
additif
\[
  \rho_*:M\longmapsto M_{[\rho]}
\]
de la catégorie $\mathbf{Ab}_B$ des $B$-modules à gauche dans la catégorie
$\mathbf{Ab}_A$ des $A$-modules. Par composition, on en déduit un foncteur
\[
  T_{(B)}:
  \mathbf{Ab}_B
  \xrightarrow{\rho_*}
  \mathbf{Ab}_A
  \xrightarrow{T}
  \mathbf{Ab},
\]
évidemment covariant et additif, que nous noterons aussi $T^{(B)}$ (pour des
raisons typographiques) ou $T\otimes_A B$, et que nous dirons obtenu à partir
de $T$ par extension des scalaires de $A$ à $B$. Bien entendu, si $B$ est
commutative, on peut considérer $T_{(B)}$ comme un foncteur de
$\mathbf{Ab}_B$ dans elle-même (7.1.2). Lorsque $B$ est commutative, et $C$
une $B$-algèbre, on voit aussitôt que
\[
  T_{(C)}=(T_{(B)})_{(C)};
\]
il est immédiat que l'extension des scalaires est fonctorielle et additive en
$T$; en outre, lorsque $T$ commute aux limites inductives ou aux sommes
directes (resp. est exact à gauche, exact à droite, exact), il en est de même
de $T_{(B)}$: en effet, $\rho_*$ est exact et commute aux limites inductives
et aux sommes directes.
\end{env}

\begin{env}[7.1.4]
Supposons toujours $A$ commutatif, et soit
$T:\mathbf{Ab}_A\to\mathbf{Ab}_A$ un foncteur covariant additif
$A$-linéaire, commutant aux limites inductives. Alors, pour toute partie
multiplicative $S$ de $A$ et tout $A$-module $M$, on a un isomorphisme
canonique fonctoriel de $A$-modules
\[
  T(S^{-1}M)\xrightarrow{\sim}S^{-1}T(M).
  \tag{7.1.4.1}\label{III.7.1.4.1.fr}
\]
Supposons en effet d'abord que $S$ soit l'ensemble des puissances $f^n$
$(n\geqslant0)$ d'un élément $f\in A$. On sait alors que
\[
  M_f=\varinjlim M_n,
\]
où $(M_n,\varphi_{nm})$ est le système inductif des $A$-modules $M_n=M$,
avec
\[
  \varphi_{nm}:z\longmapsto f^{\,n-m}z
  \qquad(0_{\mathrm I}, 1.6.1);
\]
d'où dans ce cas l'isomorphisme (7.1.4.1) en vertu de l'hypothèse sur $T$; si
ensuite $S$ est quelconque, $S^{-1}M$ est limite inductive des $M_f$, pour
$f\in S$ ($0_{\mathrm I}$, 1.4.5) et l'on conclut de la même façon. En outre,
la fonctorialité de l'isomorphisme (7.1.4.1) montre que c'est un isomorphisme
de $S^{-1}A$-modules, et qu'on peut donc écrire, à un isomorphisme canonique
près
\[
  T_{(S^{-1}A)}(S^{-1}M)
  =
  S^{-1}T(M)
  =
  T(S^{-1}M).
  \tag{7.1.4.2}\label{III.7.1.4.2.fr}
\]
Lorsque $S=A-\mathfrak p$ est le complémentaire d'un idéal premier
$\mathfrak p$ de $A$, on écrit $T_{\mathfrak p}$ au lieu de
$T_{(A_{\mathfrak p})}$.
\end{env}

\begin{proposition}[7.1.5]
Sous les hypothèses de (7.1.4), si $T_{\mathfrak m}$ est exact à gauche
(resp. exact à droite, exact) pour tout idéal maximal $\mathfrak m$ de $A$,
$T$ est exact à gauche (resp. exact à droite, exact).
\end{proposition}

\begin{proof}
On sait en effet que lorsque deux sous-modules $N$, $P$ d'un $A$-module $M$
sont tels que $N_{\mathfrak m}=P_{\mathfrak m}$ pour tout idéal maximal
$\mathfrak m$ de $A$, on a $N=P$ (Bourbaki, \emph{Alg. comm.}, chap. II,
\S~3, no~3, th.~1).
\end{proof}

\subsection{Caractérisations du foncteur produit tensoriel}
\label{subsection:III.7.2.fr}

\begin{env}[7.2.1]
Soient $A$ un anneau (non nécessairement commutatif), $M$ (resp. $N$) un
$A$-module à gauche (resp. à droite), $P$ un $\mathbf Z$-module. Rappelons
que la donnée d'un $\mathbf Z$-homomorphisme
\[
  v:N\otimes_A M\longrightarrow P
\]
équivaut à la donnée d'une application $\mathbf Z$-bilinéaire
\[
  u:N\times M\longrightarrow P
\]
telle que
\[
  u(ta,x)=u(t,ax)
  \qquad\text{pour }a\in A,\ t\in N,\ x\in M,
\]
les deux applications étant reliées par $v(t\otimes x)=u(t,x)$. D'autre part,
la donnée de $u$ équivaut à celle d'un
\oldpage[III]{177}
$\mathbf Z$-homomorphisme $x\mapsto f_x$ de $M$ dans
$\operatorname{Hom}_{\mathbf Z}(N,P)$ tel que
\[
  f_{ax}(t)=f_x(ta)
  \qquad\text{pour }a\in A,\ t\in N,\ x\in M,
\]
les deux applications étant reliées par $u(t,x)=f_x(t)$.
\end{env}

\begin{env}[7.2.2]
Soit $T:\mathbf{Ab}_A\to\mathbf{Ab}$ un foncteur covariant additif. Nous
allons définir pour tout $A$-module à gauche $M$, un homomorphisme canonique
fonctoriel en $M$, de $\mathbf Z$-modules
\[
  t_M:T(A_s)\otimes_A M\longrightarrow T(M).
  \tag{7.2.2.1}\label{III.7.2.2.1.fr}
\]
Il suffira pour cela, en vertu de (7.2.1), de définir un
$\mathbf Z$-homomorphisme $x\mapsto t'_M(x)$ de $M$ dans
$\operatorname{Hom}_{\mathbf Z}(T(A_s),T(M))$, tel que l'on ait
\[
  t'_M(ax)(y)=t'_M(x)(ya)
  \qquad\text{pour }a\in A,\ x\in M,\ y\in T(A_s).
\]
Notons pour cela que
$\operatorname{Hom}_{\mathbf Z}(T(A_s),T(M))$ est canoniquement muni d'une
structure de $A$-module à gauche provenant de la structure de $A$-module à
droite de $T(A_s)$, la loi externe étant telle que si
$a\in A$, $v\in\operatorname{Hom}_{\mathbf Z}(T(A_s),T(M))$,
\[
  (a\cdot v)(y)=v(ya)
  \qquad\text{pour }y\in T(A_s).
\]
Cela étant, nous définirons $t'_M$ comme composé des deux homomorphismes
canoniques
\[
  M\xrightarrow{\sim}\operatorname{Hom}_A(A_s,M)
  \xrightarrow{T}
  \operatorname{Hom}_{\mathbf Z}(T(A_s),T(M)),
\]
la seconde flèche étant l'application $u\mapsto T(u)$, la première
l'isomorphisme canonique de $A$-modules $x\mapsto\theta_x$ tel que
$\theta_x(\xi)=\xi x$ pour $\xi\in A$, $x\in M$. On a
$\theta_{ax}=\theta_x\circ h_a$, donc
\[
  T(\theta_{ax})
  =
  T(\theta_x\circ h_a)
  =
  T(\theta_x)\circ T(h_a)
\]
et par suite, pour $y\in T(A_s)$,
\[
  T(\theta_{ax})(y)
  =
  T(\theta_x)(T(h_a)(y))
  =
  T(\theta_x)(ya)
\]
par définition de la loi externe sur $T(A_s)$, ce qui prouve l'existence de
$t_M$; il est immédiat de vérifier que cet homomorphisme est fonctoriel en
$M$, c'est-à-dire que pour tout homomorphisme $w:M\to M'$ de $A$-modules à
gauche, le diagramme
\[
\xymatrix@C=30pt@R=30pt{
  T(A_s)\otimes_A M
  \ar[r]^-{t_M}\ar[d]_{1\otimes w}
  &
  T(M)\ar[d]^{T(w)}
  \\
  T(A_s)\otimes_A M'
  \ar[r]_-{t_{M'}}
  &
  T(M')
}
\tag{7.2.2.2}\label{III.7.2.2.2.fr}
\]
est commutatif.

La fonctorialité de l'homomorphisme (7.2.2.1) montre que lorsque $A$ est
commutatif, c'est un homomorphisme de $A$-modules (cf. (7.1.2)).
\end{env}

\begin{env}[7.2.3]
Lorsque $A$ est commutatif, on peut plus généralement définir un
homomorphisme canonique de $A$-modules
\[
  T(N)\otimes_A M\longrightarrow T(N\otimes_A M)
  \tag{7.2.3.1}\label{III.7.2.3.1.fr}
\]
pour tout $A$-module $N$; il suffit dans la construction de (7.2.2) de
remplacer l'homomorphisme $\theta_x$ par l'homomorphisme de $A$-modules
\[
  N\longrightarrow N\otimes_A M
\]
qui à tout $y\in N$ fait correspondre $y\otimes x$. Il est immédiat que cet
homomorphisme est fonctoriel en $M$ et $N$.

\oldpage[III]{178}
En particulier, si $B$ est une $A$-algèbre (non nécessairement commutative),
on a un homomorphisme fonctoriel en $M$
\[
  (T(M))_{(B)}
  =
  T(M)\otimes_A B
  \longrightarrow
  T(M\otimes_A B)
  =
  T_{(B)}(M_{(B)})
  \tag{7.2.3.2}\label{III.7.2.3.2.fr}
\]
qui, en vertu de la fonctorialité de (7.2.3.1) en $M$, est un homomorphisme
de $B$-modules.

On a en outre le diagramme commutatif
\[
\xymatrix@C=30pt@R=30pt{
  T(A)\otimes_A M
  \ar[r]^-{t_M}\ar[d]
  &
  T(M)\ar[d]
  \\
  T_{(B)}(B_s)\otimes_B M_{(B)}
  \ar[r]_-{t_{M_{(B)}}}
  &
  T_{(B)}(M_{(B)})
}
\tag{7.2.3.3}\label{III.7.2.3.3.fr}
\]
où la flèche verticale de droite est l'homomorphisme composé
\[
  T(M)
  \longrightarrow
  T(M)\otimes_A B
  \longrightarrow
  T(M\otimes_A B)
  =
  T_{(B)}(M_{(B)})
\]
de (7.2.3.2) et de l'homomorphisme canonique; quant à la flèche verticale de
gauche de (7.2.3.3), c'est l'homomorphisme
\[
  T(A)\otimes_A M
  \longrightarrow
  T_{(B)}(B_s)\otimes_B(B\otimes_A M)
  =
  T_{(B)}(B_s)\otimes_A M
\]
où
\[
  T(A)\longrightarrow T_{(B)}(B_s)=T(B)
\]
est $T(\rho)$, $\rho$ étant considéré comme homomorphisme de $A$-modules
$A\to B_{[\rho]}$.
\end{env}

\begin{lemma}[7.2.4]
Si $T$ est un foncteur covariant additif de $\mathbf{Ab}_A$ dans
$\mathbf{Ab}$, commutant avec les sommes directes, l'homomorphisme canonique
$t_L$ (7.2.2.1) est un isomorphisme pour tout $A$-module libre $L$.
\end{lemma}

\begin{proof}
En effet, on a
\[
  L=\bigoplus_{\alpha\in I}L_\alpha
\]
où $L_\alpha$ est isomorphe à $A_s$ pour tout $\alpha\in I$; la définition de
$t_M$ donnée dans (7.2.2) montre que
\[
  t_L=\bigoplus_{\alpha\in I}t_{L_\alpha},
\]
puisque
\[
  T:
  \operatorname{Hom}_A(A_s,L)
  \longrightarrow
  \operatorname{Hom}_{\mathbf Z}(T(A_s),T(L))
\]
est somme directe des applications $\mathbf Z$-linéaires
\[
  T_\alpha:
  \operatorname{Hom}_A(A_s,L_\alpha)
  \longrightarrow
  \operatorname{Hom}_{\mathbf Z}(T(A_s),T(L_\alpha))
\]
en vertu de l'hypothèse sur $T$. On est donc ramené à prouver le lemme pour
$L=A_s$; mais $t_L$ n'est autre alors que l'isomorphisme canonique
\[
  T(A_s)\otimes_A A_s\xrightarrow{\sim}T(A_s)
\]
valable pour tout $A$-module à droite.
\end{proof}

\begin{proposition}[7.2.5]
Soit $T$ un foncteur covariant additif de $\mathbf{Ab}_A$ dans
$\mathbf{Ab}$, commutant aux sommes directes. Les conditions suivantes sont
équivalentes :
\begin{enumerate}
  \item[a)] $T$ est exact à droite.
  \item[b)] L'homomorphisme canonique $t_M$ (7.2.2.1) est un isomorphisme pour
    tout $A$-module à gauche $M$.
  \item[$b'$)] $T$ est semi-exact et l'homomorphisme $t_M$ est surjectif pour
    tout $A$-module à gauche $M$.
  \item[c)] $T$ est isomorphe à un foncteur en $M$ de la forme
    $N\otimes_A M$, où $N$ est un $A$-module à droite.
\end{enumerate}
\end{proposition}

\begin{proof}
Il est clair que $b)$ implique $c)$ et que $c)$ implique $a)$; montrons que
$a)$ implique $b)$.
\oldpage[III]{179}
Posons
\[
  T'(M)=T(A_s)\otimes_A M
\]
pour tout $A$-module à gauche $M$. Il existe une suite exacte
\[
  L'\longrightarrow L\longrightarrow M\longrightarrow0,
\]
où $L$ et $L'$ sont deux $A$-modules à gauche libres; comme $T$ et $T'$ sont
exacts à droite, on a donc le diagramme commutatif
\[
\xymatrix@C=22pt@R=25pt{
  T'(L')\ar[r]\ar[d]^{t_{L'}}
  &
  T'(L)\ar[r]\ar[d]^{t_L}
  &
  T'(M)\ar[r]\ar[d]^{t_M}
  &
  0
  \\
  T(L')\ar[r]
  &
  T(L)\ar[r]
  &
  T(M)\ar[r]
  &
  0
}
\]
où les deux lignes sont exactes; comme $t_L$ et $t_{L'}$ sont des
isomorphismes en vertu de (7.2.4), il en est de même de $t_M$ par le lemme des
cinq. Enfin, il est clair que $b)$ entraîne $b')$. Pour montrer que $b')$
entraîne $a)$, il suffit de prouver le lemme suivant.
\end{proof}

\begin{lemma}[7.2.5.1]
Soient $\mathcal K$, $\mathcal K'$ deux catégories abéliennes, $F$, $G$ deux
foncteurs covariants additifs de $\mathcal K$ dans $\mathcal K'$,
$f:F\to G$ un morphisme fonctoriel (T, I, 1.2) tel que, pour tout objet $E$
de la catégorie $\mathcal K$,
\[
  f_E:F(E)\longrightarrow G(E)
\]
soit un épimorphisme. Alors, si $F$ est exact à droite et $G$ semi-exact,
$G$ est exact à droite.
\end{lemma}

\begin{proof}
Tout revient en effet à démontrer que pour tout épimorphisme
$v:E'\to E$ dans $\mathcal K$,
\[
  G(v):G(E')\longrightarrow G(E)
\]
est un épimorphisme; or, on a le diagramme commutatif
\[
\xymatrix@C=35pt@R=30pt{
  F(E')\ar[r]^{F(v)}\ar[d]_{f_{E'}}
  &
  F(E)\ar[d]^{f_E}
  \\
  G(E')\ar[r]_{G(v)}
  &
  G(E)
}
\]
dans lequel $F(v)$, $f_{E'}$ et $f_E$ sont des épimorphismes; il en est donc
de même de $G(v)$.
\end{proof}

\begin{remark}[7.2.6]
Pour tout $A$-module à droite $N$, posons
\[
  T_N(M)=N\otimes_A M
\]
pour tout $A$-module à gauche $M$, de sorte que $T_N$ est un foncteur
covariant additif de $\mathbf{Ab}_A$ dans $\mathbf{Ab}$, exact à droite et
commutant aux sommes directes. Si on identifie canoniquement $T_N(A_s)$ à
$N$, on vérifie aussitôt que l'homomorphisme correspondant (7.2.2.1) devient
l'identité. On en conclut que le $A$-module à droite $N$ de l'énoncé de
(7.2.5, $c)$) est déterminé à un isomorphisme unique près et est
canoniquement isomorphe à $T(A_s)$. On peut encore dire que les morphismes
fonctoriels
\[
  T\longmapsto T(A_s),
  \qquad
  N\longmapsto T_N
\]
constituent une équivalence (T, I, 1.2) de la catégorie des $A$-modules à
droite et de la catégorie des foncteurs covariants additifs
$\mathbf{Ab}_A\to\mathbf{Ab}$ qui sont exacts à droite et commutent aux
sommes directes.
\end{remark}

\begin{proposition}[7.2.7]
Soient $A$ un anneau artinien à gauche, dont le quotient par son radical
$\mathfrak m$ est un corps $k$. Soit $T$ un foncteur covariant additif de
$\mathbf{Ab}_A$ dans $\mathbf{Ab}$, commutant aux sommes directes. Les
conditions de (7.2.5) sont alors aussi équivalentes à
\begin{enumerate}
  \item[d)] $T$ est semi-exact et l'homomorphisme
    \[
      T(\varepsilon):T(A_s)\longrightarrow T(k)
    \]
    déduit de l'homomorphisme canonique $\varepsilon:A_s\to k$ est surjectif.
\end{enumerate}
\end{proposition}

\begin{proof}
\oldpage[III]{180}
Il est clair que la condition $b')$ de (7.2.5) entraîne $d)$; prouvons que
$d)$ entraîne $b')$.

Il existe un entier $n$ tel que $\mathfrak m^n=0$; posons, pour tout
$A$-module $M$, $M_h=\mathfrak m^hM$; nous prouverons par récurrence
descendante sur $h$ que $t_{M_h}$ est surjectif. La proposition est évidente
pour $h=n$; pour $h<n$, on a une suite exacte
\[
  0\longrightarrow M_{h+1}
  \longrightarrow M_h
  \longrightarrow M_h/M_{h+1}
  \longrightarrow 0
\]
et l'hypothèse de récurrence entraîne que $t_{M_{h+1}}$ est surjectif.
D'autre part, $M_h/M_{h+1}$ est annulé par $\mathfrak m$ et est donc un
$(A/\mathfrak m)$-module, autrement dit est somme directe de $A$-modules
isomorphes à $k$. Pour prouver que $t_{M_h/M_{h+1}}$ est surjectif, il suffit
donc de prouver que $t_k$ l'est, puisque $T$ commute aux sommes directes. Or,
en vertu de la commutativité du diagramme
\[
\xymatrix@C=35pt@R=30pt{
  T(A_s)\otimes_A A_s
  \ar[r]^-{t_{A_s}}\ar[d]_{1\otimes\varepsilon}
  &
  T(A_s)\ar[d]^{T(\varepsilon)}
  \\
  T(A_s)\otimes_A k
  \ar[r]_-{t_k}
  &
  T(k)
}
\]
et de (7.2.4), l'hypothèse $d)$ entraîne que $t_k$ est bien surjectif. Pour
terminer la démonstration, il suffira de prouver que si l'on a une suite
exacte
\[
  0\longrightarrow M'
  \longrightarrow M
  \xrightarrow{v}M''
  \longrightarrow 0
\]
de $A$-modules, telle que $t_{M'}$ et $t_{M''}$ sont surjectifs, alors $t_M$
est surjectif. Or, on a un diagramme commutatif
\[
\xymatrix@C=27pt@R=30pt{
  T'(M')\ar[r]\ar[d]_{t_{M'}}
  &
  T'(M)\ar[r]\ar[d]_{t_M}
  &
  T'(M'')\ar[r]\ar[d]_{t_{M''}}
  &
  0\ar[d]
  \\
  T(M')\ar[r]
  &
  T(M)\ar[r]_-{T(v)}
  &
  T(M'')\ar[r]
  &
  \operatorname{Coker}(T(v))
}
\]
dans lequel les deux lignes sont exactes, en vertu de l'hypothèse que $T$ est
semi-exact. Comme par l'hypothèse de récurrence $t_{M'}$ et $t_{M''}$ sont
des épimorphismes et que la dernière flèche verticale est un monomorphisme,
le lemme des cinq (M, I, 1.1) montre que $t_M$ est un épimorphisme.
\end{proof}

\subsection{Critères d'exactitude des foncteurs homologiques de modules}
\label{subsection:III.7.3.fr}

\begin{proposition}[7.3.1]
Soient $A$ un anneau (non nécessairement commutatif), $T_\bullet$ un foncteur
homologique covariant (T, II, 2.1) de la catégorie $\mathbf{Ab}_A$ dans la
catégorie $\mathbf{Ab}$, commutant aux sommes directes. Soit $p$ un entier tel
que $T_p$ et $T_{p-1}$ soient définis. Les conditions suivantes sont
équivalentes :
\begin{enumerate}
  \item[a)] $T_p$ est exact à droite.
  \item[b)] $T_{p-1}$ est exact à gauche.
  \oldpage[III]{181}
  \item[c)] Pour tout $A$-module à gauche $M$, l'homomorphisme fonctoriel
    canonique (7.2.2.1)
    \[
      T_p(A_s)\otimes_A M\longrightarrow T_p(M)
      \tag{7.3.1.1}\label{III.7.3.1.1.fr}
    \]
    est un isomorphisme.
  \item[d)] Pour tout $A$-module à gauche $M$, l'homomorphisme (7.3.1.1) est
    un épimorphisme.
  \item[e)] $T_p$ est isomorphe à un foncteur
    $M\mapsto N\otimes_A M$, où $N$ est un $A$-module à droite.
\end{enumerate}
Si en outre les conditions de (7.2.7) sur $A$ et $\mathfrak m$ sont vérifiées,
les conditions précédentes sont aussi équivalentes à
\begin{enumerate}
  \item[f)] L'homomorphisme canonique
    \[
      T_p(\varepsilon):T_p(A_s)\longrightarrow T_p(k)
    \]
    est un épimorphisme.
\end{enumerate}
\end{proposition}

\begin{proof}
Comme par définition d'un foncteur homologique, $T_i$ est semi-exact pour tous
les $i$ tels que $T_i$ soit défini, et que, pour toute suite exacte
\[
  0\longrightarrow M'\xrightarrow{u}M\xrightarrow{v}M''\longrightarrow0,
\]
on a
\[
  \operatorname{Ker}(T_{i-1}(u))
  =
  \operatorname{Coker}(T_i(v)),
\]
il est clair que $a)$ et $b)$ sont équivalentes et les autres assertions
résultent trivialement de (7.2.5) et (7.2.7).
\end{proof}

\begin{corollary}[7.3.2]
Soit $A$ un anneau commutatif. Avec les notations de (7.3.1), supposons $T_p$
exact à droite. Si $f\in A$ n'appartient à l'annulateur d'aucun élément
$\ne0$ d'un $A$-module $M$, alors $f$ n'appartient à l'annulateur d'aucun
élément $\ne0$ de $T_{p-1}(M)$. En particulier, si $A$ est intègre, le
$A$-module $T_{p-1}(A)$ est sans torsion.
\end{corollary}

\begin{proof}
En effet, si $h_f$ désigne l'homothétie $x\mapsto fx$ de $M$, l'hypothèse
signifie que $h_f$ est injectif; il en est donc de même de $T_{p-1}(h_f)$ par
la condition $b)$ de (7.3.1).
\end{proof}

\begin{proposition}[7.3.3]
Soient $A$ un anneau, $T_\bullet$ un foncteur homologique covariant de
$\mathbf{Ab}_A$ dans $\mathbf{Ab}$, commutant aux sommes directes. Soit $p$ un
entier tel que $T_{p-1}$, $T_p$ et $T_{p+1}$ soient définis. Les conditions
suivantes sont équivalentes :
\begin{enumerate}
  \item[a)] $T_p$ est exact.
  \item[b)] $T_{p+1}$ et $T_p$ sont exacts à droite.
  \item[c)] $T_p$ et $T_{p-1}$ sont exacts à gauche.
  \item[d)] $T_{p+1}$ est exact à droite et $T_{p-1}$ est exact à gauche.
  \item[e)] Pour tout $A$-module $M$, les homomorphismes canoniques
    \[
      T_i(A_s)\otimes_A M\longrightarrow T_i(M)
      \tag{7.3.3.1}\label{III.7.3.3.1.fr}
    \]
    sont des isomorphismes pour $i=p$ et $i=p+1$.
  \item[$e'$)] Pour tout $A$-module $M$, les homomorphismes canoniques
    (7.3.3.1) sont des épimorphismes pour $i=p$ et $i=p+1$.
  \item[f)] Pour tout $A$-module $M$, l'homomorphisme (7.3.3.1) est un
    isomorphisme pour $i=p$ et $T_p(A_s)$ est un $A$-module à droite plat.
  \item[$f'$)] Pour tout $A$-module $M$, l'homomorphisme (7.3.3.1) est un
    épimorphisme pour $i=p$ et $T_p(A_s)$ est un $A$-module à droite plat.
\end{enumerate}
\end{proposition}

\begin{proof}
L'équivalence des conditions $a)$, $b)$, $c)$, $d)$ résulte de l'équivalence
des conditions $a)$ et $b)$ de (7.3.1). L'équivalence de $b)$, $e)$ et $e')$
résulte de l'équivalence de $a)$, $c)$ et $d)$ dans (7.3.1). Enfin, dire que
$T_p(A_s)$ est plat signifie que le foncteur
\[
  M\mapsto T_p(A_s)\otimes_A M
\]
est exact à gauche; l'équivalence de $a)$, $f)$ et $f')$ résulte encore de
l'équivalence de $a)$, $c)$, $d)$ dans (7.3.1).
\end{proof}

\oldpage[III]{182}

\begin{corollary}[7.3.4]
Supposons $A$ commutatif, $T_p$ exact et supposons en outre que $T_p(A)$ soit
un $A$-module de présentation finie. Alors la fonction
\[
  x\mapsto\operatorname{rang}_{\kappa(x)}(T_p(\kappa(x)))
\]
est localement constante dans $X=\operatorname{Spec}(A)$, donc constante si
$\operatorname{Spec}(A)$ est connexe.
\end{corollary}

\begin{proof}
En effet, comme $T_p(A)$ est un $A$-module plat en vertu de (7.3.3, $f)$), il
est projectif de type fini, et $\bigl(T_p(A)\bigr)^{\sim}$ est donc un
$\mathscr O_X$-Module localement libre (Bourbaki, \emph{Alg. comm.}, chap. II,
\S 5, n\textsuperscript{o} 2, th. 1); on a en outre
\[
  T_p(\kappa(x))=T_p(A)\otimes_A\kappa(x)
\]
(7.3.3, $e)$) et l'on sait que le rang au point $x$ du $A$-module $T_p(A)$ est
localement constant (\emph{loc. cit.}), d'où le corollaire.
\end{proof}

\begin{proposition}[7.3.5]
Supposons que $A$ soit un anneau artinien à gauche dont le quotient par son
radical $\mathfrak m$ est un corps $k$. Alors les conditions de (7.3.3) sont
encore équivalentes à chacune des suivantes :
\begin{enumerate}
  \item[g)] L'homomorphisme canonique
    \[
      T_i(\varepsilon):T_i(A_s)\longrightarrow T_i(k)
    \]
    est un épimorphisme pour $i=p$ et $i=p+1$.
  \item[h)] $T_p(\varepsilon)$ est un épimorphisme et $T_p(A_s)$ est un
    $A$-module à droite plat (ou, ce qui revient au même (Bourbaki,
    \emph{Alg. comm.}, chap. II, \S 3, n\textsuperscript{o} 2, cor. 2 de la
    prop. 5) un $A$-module libre).
\end{enumerate}
Supposons en outre que $A$ soit commutatif et le $A$-module $T_p(k)$ de
longueur finie $d$. Alors les conditions précédentes équivalent aussi à
chacune des suivantes :
\begin{enumerate}
  \item[i)] Pour tout $A$-module $M$ de longueur finie, on a
    \[
      \operatorname{long}(T_p(M))=d\mathbin{.}\operatorname{long}(M).
      \tag{7.3.5.1}\label{III.7.3.5.1.fr}
    \]
  \item[j)] On a
    \[
      \operatorname{long}(T_p(A))=d\mathbin{.}\operatorname{long}(A).
      \tag{7.3.5.2}\label{III.7.3.5.2.fr}
    \]
\end{enumerate}
\end{proposition}

L'équivalence de $g)$ et $h)$ avec les conditions de (7.3.3) découle aussitôt
de (7.2.7). Pour démontrer les autres assertions, nous utiliserons le lemme
suivant :

\begin{lemma}[7.3.5.3]
Soient $\mathcal K$, $\mathcal K'$ deux catégories abéliennes,
$F:\mathcal K\to\mathcal K'$ un foncteur
covariant additif; on suppose que $F$ soit semi-exact, et que, pour tout objet
simple $S$ de $\mathcal K$, $F(S)$ soit un objet de longueur finie dans
$\mathcal K'$. Alors, pour tout objet $E$ de longueur finie dans $\mathcal K$,
$F(E)$ est de longueur finie dans $\mathcal K'$. Pour toute suite exacte
\[
  0\longrightarrow E'\xrightarrow{u}E\xrightarrow{v}E''\longrightarrow0
\]
d'objets de longueur finie dans $\mathcal K$, on a
\[
  \operatorname{long}F(E)
  \leq
  \operatorname{long}F(E')+\operatorname{long}F(E'')
  \tag{7.3.5.4}\label{III.7.3.5.4.fr}
\]
et pour que les deux membres de (7.3.5.4) soient égaux, il faut et il suffit
que la suite
\[
  0\longrightarrow F(E')\longrightarrow F(E)\longrightarrow F(E'')
  \longrightarrow0
\]
soit exacte.
\end{lemma}

\begin{proof}
En effet, la suite
\[
  F(E')\xrightarrow{F(u)}F(E)\xrightarrow{F(v)}F(E'')
\]
est exacte par hypothèse; si l'on suppose $F(E')$ et $F(E'')$ de longueur
finie, il en est de même de $\operatorname{Im}(F(u))$ et de
$\operatorname{Im}(F(v))$ et comme
$\operatorname{Ker}(F(v))=\operatorname{Im}(F(u))$, $F(E)$ est de longueur
finie et l'on a
\[
  \operatorname{long}F(E)
  =
  \operatorname{long}\operatorname{Im}(F(u))
  +
  \operatorname{long}\operatorname{Im}(F(v))
  \leq
  \operatorname{long}F(E')+\operatorname{long}F(E'').
  \tag{7.3.5.5}\label{III.7.3.5.5.fr}
\]
\oldpage[III]{183}
Par récurrence sur la longueur de $E$, cela prouve déjà la première assertion;
en outre, les deux membres de (7.3.5.5) ne peuvent être égaux que si
\[
  \operatorname{long}\operatorname{Im}(F(u))=\operatorname{long}F(E')
\]
(ce qui équivaut à
$\operatorname{long}\operatorname{Ker}(F(u))=0$, ou
$\operatorname{Ker}(F(u))=0$) et
\[
  \operatorname{long}\operatorname{Im}(F(v))=\operatorname{long}F(E'')
\]
(ce qui équivaut à
$\operatorname{long}\operatorname{Coker}(F(v))=0$, ou
$\operatorname{Coker}(F(v))=0$).
\end{proof}

Notons maintenant que si $M$ est un $A$-module de longueur finie ($A$ étant
commutatif), les quotients d'une suite de Jordan-Hölder de $M$ sont
nécessairement isomorphes au $A$-module $k$, donc on déduit de (7.3.5.4), par
récurrence sur la longueur de $M$
\[
  \operatorname{long}T_p(M)\leq d\mathbin{.}\operatorname{long}(M).
  \tag{7.3.5.6}\label{III.7.3.5.6.fr}
\]
En outre, il résulte de (7.3.5.3) que si $T_p$ est exact, on a l'égalité
(7.3.5.1); donc la condition $a)$ de (7.3.3) entraîne $i)$; il est clair que
$i)$ entraîne $j)$, et il reste à prouver le

\begin{lemma}[7.3.5.7]
La relation
\[
  \operatorname{long}T_p(A)=d\mathbin{.}\operatorname{long}A
\]
entraîne que $T_p(\varepsilon)$ est un épimorphisme et que $T_p(A)$ est un
$A$-module plat.
\end{lemma}

\begin{proof}
En effet, partant de la suite exacte
\[
  0\longrightarrow\mathfrak m\longrightarrow A\longrightarrow k
  \longrightarrow0,
\]
il résulte de (7.3.5.4) et (7.3.5.6) que l'on a
\[
  \operatorname{long}T_p(A)
  \leq
  \operatorname{long}T_p(\mathfrak m)+\operatorname{long}T_p(k)
  \leq
  d\bigl(\operatorname{long}\mathfrak m+\operatorname{long}k\bigr)
  =d\mathbin{.}\operatorname{long}A
\]
et que l'égalité ne peut avoir lieu (7.3.5.3) que si la suite
\[
  0\longrightarrow T_p(\mathfrak m)\longrightarrow T_p(A)
  \longrightarrow T_p(k)\longrightarrow0
  \tag{7.3.5.8}\label{III.7.3.5.8.fr}
\]
est exacte. En vertu de (7.2.7) et (7.2.5), $T_p$ est isomorphe à un foncteur
$M\mapsto N\otimes_A M$, et l'exactitude de la suite (7.3.5.8) montre, en vertu
de la suite exacte des $\operatorname{Tor}$, que l'on a
$\operatorname{Tor}_1^A(N,k)=0$. On en conclut que $N=T_p(A)$ est un
$A$-module plat (0, 10.1.3).
\end{proof}

\begin{lemma}[7.3.6]
Soient $A$ un anneau, $T_\bullet$ un foncteur homologique covariant de
$\mathbf{Ab}_A$ dans $\mathbf{Ab}$, commutant aux sommes directes. Supposons
$T_p$ et $T_{p+1}$ définis, et $T_p$ exact à gauche. Pour que $T_{p+1}$ soit
exact, il faut et il suffit que $T_{p+1}(A_s)$ soit un $A$-module à droite
plat.
\end{lemma}

\begin{proof}
En effet, on sait d'après (7.3.1) que l'homomorphisme canonique
\[
  T_{p+1}(A_s)\otimes_A M\longrightarrow T_{p+1}(M)
\]
est un isomorphisme de foncteurs; il suffit d'appliquer la définition d'un
$A$-module plat.
\end{proof}

\begin{proposition}[7.3.7]
Soient $A$ un anneau, $T_\bullet$ un foncteur homologique covariant de
$\mathbf{Ab}_A$ dans $\mathbf{Ab}$, commutant aux sommes directes. On suppose
qu'il existe $i_0$ tel que $T_i$ soit exact pour $i\leq i_0$. Alors, pour tout
entier $p>i_0$, les conditions suivantes sont équivalentes :
\begin{enumerate}
  \item[a)] $T_q$ est exact pour $q\leq p$.
  \item[b)] $T_q(A_s)$ est un $A$-module à droite plat pour $q\leq p$.
  \item[c)] Pour tout $A$-module $M$, l'homomorphisme canonique
    \[
      T_q(A_s)\otimes_A M\longrightarrow T_q(M)
    \]
    est surjectif pour $q\leq p+1$.
\end{enumerate}
\end{proposition}

\begin{proof}
L'équivalence de $a)$ et $b)$ résulte de (7.3.6) par récurrence sur $q$,
puisque $T_{i_0}$ est exact par hypothèse; l'équivalence de $a)$ et $c)$
résulte de l'équivalence des conditions $a)$ et $e')$ dans (7.3.3).
\end{proof}

\oldpage[III]{184}

\begin{env}[7.3.8]
Si $A$ est un anneau commutatif, $B$ une $A$-algèbre (non nécessairement
commutative), $T_\bullet$ un foncteur homologique covariant de
$\mathbf{Ab}_A$ dans $\mathbf{Ab}$, il résulte des définitions (7.1.3) que le
foncteur de $\mathbf{Ab}_B$ dans $\mathbf{Ab}$ obtenu par extension des
scalaires de $A$ à $B$, et que nous noterons
\[
  T_\bullet^{(B)}=(T_i^{(B)}),
\]
est encore un foncteur homologique.
\end{env}

\begin{corollary}[7.3.9]
Supposons que $T_\bullet$ vérifie les conditions générales de (7.3.7) et
commute aux limites inductives, et en outre que $A$ soit un anneau intègre et
tous les $T_n(A)$ des $A$-modules de présentation finie. Alors, pour tout
entier $N$, il existe un $f\in A-\{0\}$ tel que le foncteur
\[
  T_p^{(A_f)}:\mathbf{Ab}_{A_f}\longrightarrow\mathbf{Ab}
\]
soit exact pour $p\leq N$.
\end{corollary}

\begin{proof}
Par hypothèse, $T_i$ est exact pour $i\leq i_0$, donc $T_i(A)$ est plat pour
ces valeurs de $i$. En vertu de (7.3.7, $b)$), il suffit de prendre $f$ tel que
\[
  T_p^{(A_f)}(A_f)=T_p(A_f)
\]
soit un $A_f$-module libre pour $i_0<p\leq N$. Or, on a
$T_p(A_f)=(T_p(A))_f$ puisque $T_p$ commute aux limites inductives (7.1.4). Si
$x$ est le point générique de $\operatorname{Spec}(A)$, $(T_p(A))_x$ est un
espace vectoriel de dimension finie sur le corps des fractions de $A$. Comme
chaque $T_p(A)$ est de présentation finie, il existe bien un $f$ ayant la
propriété voulue (Bourbaki, \emph{Alg. comm.}, chap. II, \S 5,
n\textsuperscript{o} 1, cor. de la prop. 2).

On notera que s'il n'y a qu'un nombre fini d'indices $i$ tels que $T_i\ne0$,
il existe $f\in A-\{0\}$ tel que tous les $T_p^{(A_f)}$ soient exacts.
\end{proof}

\begin{corollary}[7.3.10]
Supposons que $T_\bullet$ vérifie les conditions générales de (7.3.7) et
commute aux limites inductives, que $A$ soit commutatif et noethérien et les
$T_n(A)$ des $A$-modules de type fini. Alors, pour tout entier $N$, il existe
un ouvert dense $U$ de $\operatorname{Spec}(A)$ tel que, pour tout $p\leq N$,
la fonction
\[
  x\mapsto\operatorname{rang}_{\kappa(x)}(T_p(\kappa(x)))
\]
soit constante dans $U$.
\end{corollary}

\begin{proof}
Soit $\mathfrak p$ un idéal premier minimal de $A$; par hypothèse, l'anneau
$B=A/\mathfrak p$ est intègre et $\operatorname{Spec}(B)$ s'identifie à une
composante irréductible de l'espace topologique $\operatorname{Spec}(A)$. Nous
allons montrer, par récurrence sur $p\leq N$ qu'il existe
$f_p\in B-\{0\}$ tel que si on pose $B'=B_{f_p}$, $T_i^{(B')}$ soit exact et
les $T_i(B')$ des $B'$-modules de type fini pour $i\leq p$. La proposition est
vraie pour $p\leq i_0$ en vertu de l'hypothèse, en prenant $f_p=1$ (donc
$B'=B=A/\mathfrak p$) car $T_p$ étant alors exact, $T_p(B)$ est isomorphe à
$T_p(A)/T_p(\mathfrak p)$, donc est un $A$-module (et \emph{a fortiori} un
$B$-module) de type fini. Raisonnons par récurrence sur $p$; $f_p$ est l'image
canonique dans $B$ d'un élément $g_p\in A$, et si l'on pose $A'=A_{g_p}$, on a
$B'=A'/\mathfrak p'$ où $\mathfrak p'$ est un idéal premier minimal de $A'$,
égal à $\mathfrak p_{g_p}$. Comme on a
\[
  T_i(A_{g_p})=(T_i(A))_{g_p},
\]
les $T_i(A')$ sont des $A'$-modules de type fini, donc le foncteur
$T_\bullet^{(A')}$ vérifie les mêmes hypothèses que $T_\bullet$, mais en
remplaçant $i_0$ par $p$. On peut donc se borner au cas où $A'=A$, et où $T_p$
est exact; la suite exacte
\[
  0\longrightarrow\mathfrak p\longrightarrow A\longrightarrow A/\mathfrak p
  \longrightarrow0
\]
donne alors la suite exacte
\[
  T_{p+1}(A)\longrightarrow T_{p+1}(A/\mathfrak p)
  \xrightarrow{\partial}T_p(\mathfrak p)\longrightarrow T_p(A),
\]
et comme $T_p$ est exact, la dernière flèche de droite est injective, donc
$T_{p+1}(A/\mathfrak p)$ est un quotient de $T_{p+1}(A)$ et est par suite de
type fini. Notons maintenant que le raisonnement de (7.3.9) n'a utilisé le fait
que les $T_p(A)$ sont de type fini que pour $p\leq N$; on peut donc l'appliquer
à l'anneau intègre $B$, au foncteur $T_\bullet^{(B)}$ et à $N=p+1$, ce qui
achève le raisonnement par récurrence. Cela étant, il y a un
$f_N\in B-\{0\}$ tel que $\operatorname{Spec}(B_{f_N})$ soit un ouvert $V$
partout dense dans $\operatorname{Spec}(B)$ ne rencontrant aucune autre
composante irréductible de $\operatorname{Spec}(A)$. Si la proposition est
\oldpage[III]{185}
démontrée pour $B_{f_N}$, on aura un ouvert $W$ partout dense dans $V$ dans
lequel les fonctions de l'énoncé seront constantes, puisque
$A_x=(B_{f_N})_x$ pour tout $x\in W$. En faisant le même raisonnement pour
toute composante irréductible de $\operatorname{Spec}(A)$, le corollaire sera
démontré. On peut donc se borner au cas où $A$ est intègre; le raisonnement de
(7.3.9) prouve alors l'existence d'un $f\in A-\{0\}$ tel que les $T_p(A_f)$
soient des $A_f$-modules libres de type fini pour $p\leq N$, ce qui entraîne
la conclusion de (7.3.10) en vertu de (7.3.4).
\end{proof}

\begin{proposition}[7.3.11]
Soient $A$ un anneau commutatif local, $k$ son corps résiduel, $T_\bullet$ un
foncteur homologique covariant de $\mathbf{Ab}_A$ dans $\mathbf{Ab}$,
commutant aux sommes directes. On suppose qu'il existe $i_0$ tel que $T_i$ soit
exact pour $i\leq i_0$, et que tous les $T_n(A)$ soient des $A$-modules de
présentation finie. Alors les conditions équivalentes $a)$, $b)$, $c)$ de
(7.3.7) impliquent les deux suivantes, et leur sont équivalentes lorsque
l'anneau est en outre réduit :
\begin{enumerate}
  \item[d)] Pour tout $x\in\operatorname{Spec}(A)$, on a
    \[
      \operatorname{rang}_{\kappa(x)}T_q(\kappa(x))
      =
      \operatorname{rang}_kT_q(k)
      \quad\text{pour }q\leq p.
    \]
  \item[$d'$)] Pour tout point générique $x_j$ d'une composante irréductible
    de $\operatorname{Spec}(A)$, on a
    \[
      \operatorname{rang}_{\kappa(x_j)}T_q(\kappa(x_j))
      =
      \operatorname{rang}_kT_q(k)
      \quad\text{pour }q\leq p.
    \]
\end{enumerate}
\end{proposition}

\begin{proof}
Comme $T_q(A)$ est un $A$-module de présentation finie, la condition $b)$ de
(7.3.7) équivaut à dire que $T_q(A)$ est un $A$-module libre pour $q\leq p$
(Bourbaki, \emph{Alg. comm.}, chap. II, \S 3, n\textsuperscript{o} 2, cor. 2
de la prop. 5); la condition $c)$ implique que
\[
  T_q(\kappa(x))=T_q(A)\otimes_A\kappa(x)
  \quad\text{pour }q\leq p,
\]
donc les conditions équivalentes de (7.3.7) impliquent $d)$, et il est trivial
que $d)$ entraîne $d')$. Reste à prouver que $d')$ implique $a)$ lorsque $A$
est réduit. Raisonnons par récurrence sur $q\leq p$, puisque $T_q$ est exact
pour $q\leq i_0$. Supposons donc $T_k$ exact pour $k\leq q<p$ et montrons que
$T_{q+1}(A)$ est un $A$-module libre. En vertu de l'hypothèse de récurrence,
$T_{q+1}(A)\otimes_A M$ est isomorphe à $T_{q+1}(M)$ pour tout $A$-module $M$,
par la condition $c)$ de (7.3.7) et (7.3.3); appliquant cette propriété à
$M=\kappa(x_j)$ et $M=k$, on trouve, en vertu de l'hypothèse $d')$, que
\[
  \operatorname{rang}_{\kappa(x_j)}
  \bigl(T_{q+1}(A)\otimes_A\kappa(x_j)\bigr)
  =
  \operatorname{rang}_kT_{q+1}(k)
\]
pour tout $i$; mais cela implique que $T_{q+1}(A)$ est libre (Bourbaki,
\emph{Alg. comm.}, chap. II, \S 3, n\textsuperscript{o} 2, prop. 7), ce qui
achève la démonstration.
\end{proof}

Les résultats précédents seront considérablement améliorés pour les foncteurs
homologiques de type particulier que nous allons étudier dans (7.4); on
obtiendra en effet des critères d'exactitude ne faisant intervenir qu'un seul
des $T_p$.

\subsection{Critères d'exactitude pour les foncteurs
$H_\bullet(P_\bullet\otimes_A M)$}
\label{subsection:III.7.4.fr}

\begin{env}[7.4.1]
Soient $A$ un anneau (non nécessairement commutatif), $P_\bullet$ un complexe
de $A$-modules à droite plats. Comme le foncteur
$M\mapsto P_k\otimes_A M$ est alors exact dans $\mathbf{Ab}_A$ pour tout $k$,
le $\partial$-foncteur
\[
  T_\bullet(M)=H_\bullet(P_\bullet\otimes_A M)
  \tag{7.4.1.1}\label{III.7.4.1.1.fr}
\]
\oldpage[III]{186}
est un foncteur homologique de $\mathbf{Ab}_A$ dans $\mathbf{Ab}$,
évidemment $A$-linéaire lorsque $A$ est commutatif (7.1.2), et commutant aux
limites inductives.

Si $A$ est commutatif, alors, pour toute $A$-algèbre $B$, le foncteur
homologique $T_\bullet^{(B)}$ (7.3.8) est donné par définition par
\[
  T_\bullet^{(B)}(N)
  =
  H_\bullet(P_\bullet\otimes_A N_{[\rho]})
  \tag{7.4.1.2}\label{III.7.4.1.2.fr}
\]
où $\rho:A\to B$ est l'homomorphisme définissant la structure d'algèbre de
$B$; comme on peut aussi écrire
\[
  P_\bullet\otimes_A N_{[\rho]}
  =P_\bullet\otimes_A(B\otimes_B N)_{[\rho]}
  =(P_\bullet\otimes_A B)\otimes_B N,
\]
on voit que l'on a
\[
  T_\bullet^{(B)}(N)=H_\bullet(P'_\bullet\otimes_B N)
  \tag{7.4.1.3}\label{III.7.4.1.3.fr}
\]
pour tout $B$-module $N$, $P'_\bullet$ étant le complexe
$P_\bullet\otimes_A B$ de $B$-modules plats (0\textsubscript{I}, 6.2.1).
\end{env}

\begin{proposition}[7.4.2]
Sous les conditions générales de (7.4.1), et pour un entier $p\in\mathbf Z$
donné, les propriétés suivantes sont équivalentes :
\begin{enumerate}
  \item[a)] $T_p$ est exact à gauche (ou, ce qui revient au même, $T_{p+1}$
    est exact à droite).
  \item[b)]
    \[
      Z'_p(P_\bullet)
      =
      \operatorname{Coker}(P_{p+1}\longrightarrow P_p)
    \]
    est un $A$-module à droite plat.
  \item[c)] Il existe un complexe $P'_\bullet$ de $A$-modules à droite plats
    tel que la différentielle
    \[
      d'_{p+1}:P'_{p+1}\longrightarrow P'_p
    \]
    soit nulle, et un isomorphisme de foncteurs homologiques de
    $H_\bullet(P_\bullet\otimes_A M)$ sur
    $H_\bullet(P'_\bullet\otimes_A M)$.
\end{enumerate}
\end{proposition}

\begin{proof}
Par définition, on a une suite exacte fonctorielle en $M$
\[
  0\longrightarrow T_p(M)
  \longrightarrow Z'_p(P_\bullet\otimes M)
  \longrightarrow P_{p-1}\otimes M
\]
où
\[
  Z'_p(P_\bullet\otimes M)
  =
  \operatorname{Coker}(P_{p+1}\otimes M\longrightarrow P_p\otimes M)
  =
  Z'_p(P_\bullet)\otimes M
\]
en vertu de l'exactitude à droite du produit tensoriel. Pour tout homomorphisme
$f:M\to N$, on a donc un diagramme commutatif
\[
\xymatrix@C=30pt@R=38pt{
  0\ar[r]
  &T_p(M)\ar[r]\ar[d]_u
  &Z'_p(P_\bullet)\otimes M\ar[r]\ar[d]_v
  &P_{p-1}\otimes M\ar[d]_w
  \\
  0\ar[r]
  &T_p(N)\ar[r]
  &Z'_p(P_\bullet)\otimes N\ar[r]
  &P_{p-1}\otimes N
}
\tag{7.4.2.1}\label{III.7.4.2.1.fr}
\]
dont les lignes sont exactes. Si $f$ est un monomorphisme, il en est de même
de $w$ puisque $P_{p-1}$ est plat; si $T_p$ est exact à gauche, $u$ est lui
aussi un monomorphisme; on en conclut que $v$ est un monomorphisme, ce qui
entraîne que $Z'_p(P_\bullet)$ est plat. Inversement, s'il en est ainsi, $v$
est un monomorphisme pour tout monomorphisme $f:M\to N$, donc le diagramme
(7.4.2.1) montre que $u$ est un monomorphisme, et par suite $T_p$ (qui est déjà
semi-exact) est exact à gauche. Ainsi $a)$ et $b)$ sont équivalentes. Il est
immédiat que $c)$ entraîne $a)$, car si
$d'_{p+1}:P'_{p+1}\to P'_p$ est nulle, et
$0\to M'\to M\to M''\to0$ est une suite exacte de $A$-modules, l'opérateur
bord dans la suite exacte
\[
  H_{p+1}(P'_\bullet\otimes M'')
  \xrightarrow{\partial}H_p(P'_\bullet\otimes M')
  \longrightarrow H_p(P'_\bullet\otimes M)
\]
\oldpage[III]{187}
est nul par définition (M, IV, 1), donc $T_p$ est exact à gauche. Montrons
inversement que $b)$ implique $c)$. Si
\[
  Z_{p+1}(P_\bullet)=\operatorname{Ker}(P_{p+1}\longrightarrow P_p),
\]
on a une suite exacte
\[
  0\longrightarrow Z_{p+1}(P_\bullet)\longrightarrow P_{p+1}
  \longrightarrow Z'_p(P_\bullet)\longrightarrow0
\]
dans laquelle $P_{p+1}$ et $Z'_p(P_\bullet)$ sont plats, donc
$Z_{p+1}(P_\bullet)$ est plat ($0_{\mathrm I}$, 6.1.2). On va prendre
\[
  P'_i=P_i\quad\text{pour }i\ne p\text{ et }i\ne p+1,
  \qquad
  P'_p=Z'_p(P_\bullet),
  \qquad
  P'_{p+1}=Z_{p+1}(P_\bullet);
\]
pour différentielle $d'_i:P'_i\to P'_{i-1}$, on prendra celle du complexe
$P_\bullet$ pour $i\ne p$ et $i\ne p+1$, $0$ pour $i=p+1$ et pour $i=p$
l'homomorphisme $Z'_p(P_\bullet)\to P_{p-1}$ déduit de $d_p$ par passage au
quotient. Comme les $P_i$ sont plats, on a
\[
  Z'_i(P_\bullet\otimes M)=Z'_i(P_\bullet)\otimes M,
  \qquad
  Z_i(P_\bullet\otimes M)=Z_i(P_\bullet)\otimes M
\]
et
\[
  B_i(P_\bullet\otimes M)=B_i(P_\bullet)\otimes M
\]
(en posant $B_i(P_\bullet)=\operatorname{Im}(P_{i+1}\to P_i)$); on en conclut
aussitôt pour tout $M$ des isomorphismes fonctoriels
\[
  H_i(P_\bullet\otimes M)\xrightarrow{\sim}H_i(P'_\bullet\otimes M)
\]
pour tout $i$, et la vérification du fait qu'il s'agit d'un isomorphisme de
$\partial$-foncteurs découle sans peine de la définition de $\partial$
(M, IV, 1).
\end{proof}

On remarquera que les conditions de (7.4.2) impliquent aussi que
$B_p(P_\bullet)$ est plat, car on a une suite exacte
\[
  0\longrightarrow B_p(P_\bullet)\longrightarrow P_p
  \longrightarrow Z'_p(P_\bullet)\longrightarrow0,
\]
dans laquelle $P_p$ et $Z'_p(P_\bullet)$ sont plats ($0_{\mathrm I}$, 6.1.2).

\begin{corollary}[7.4.3]
Supposons que $A$ soit un anneau noethérien régulier de dimension $1$
(autrement dit un produit d'anneaux de Dedekind
($0_{\mathrm{IV}}$, 17.1.3 et 17.3.7), par exemple un anneau principal). Alors,
pour que $T_p$ soit exact à gauche, il faut et il suffit que $T_p(A)$ soit un
$A$-module plat. Pour que $T_p$ soit exact, il faut et il suffit que $T_p(A)$
et $T_{p-1}(A)$ soient des $A$-modules plats.
\end{corollary}

\begin{proof}
Rappelons que pour un module $M$ sur un anneau de Dedekind, il revient au même
de dire que $M$ est plat ou qu'il est sans torsion ($0_{\mathrm I}$, 6.3.3 et
6.3.4); sous les hypothèses de (7.4.3), tout sous-module d'un $A$-module plat
est donc plat.

La seconde assertion de (7.4.3) résulte de la première, puisque dire que $T_p$
est exact signifie que $T_p$ et $T_{p-1}$ sont exacts à gauche. Pour démontrer
la première assertion, notons que l'on a une suite exacte
\[
  0\longrightarrow H_p(P_\bullet)\longrightarrow Z'_p(P_\bullet)
  \longrightarrow B_{p-1}(P_\bullet)\longrightarrow0
\]
dans laquelle $B_{p-1}(P_\bullet)$ est un $A$-module plat, en tant que
sous-module du $A$-module plat $P_{p-1}$. Il est donc équivalent de dire que
$H_p(P_\bullet)$ est plat ou que $Z'_p(P_\bullet)$ est plat
($0_{\mathrm I}$, 6.1.2).
\end{proof}

Les applications les plus importantes de (7.4.2) sont les suivantes :

\begin{proposition}[7.4.4]
Soient $A$ un anneau noethérien, $P_\bullet$ un complexe de $A$-modules plats;
on suppose, soit que les $P_i$ sont de type fini, soit que les
$H_i(P_\bullet)$ sont des $A$-modules de type fini et qu'il existe $i_0$ tel
que $H_i(P_\bullet)=0$ pour $i<i_0$. Soit $T$ le foncteur homologique défini
par (7.4.1.1). Alors l'ensemble $U$ des
$y\in\operatorname{Spec}(A)$ tels que $(T_p)_y$ (7.1.4) soit exact à droite
(resp. exact à gauche, exact) est ouvert dans $\operatorname{Spec}(A)$.
\end{proposition}

\begin{proof}
Dans la seconde hypothèse sur $P_\bullet$, on peut remplacer $P_\bullet$ par
un complexe $P'_\bullet$ de
\oldpage[III]{188}
$A$-modules libres de type fini pour lequel le foncteur
$H_\bullet(P'_\bullet\otimes_A M)$ est isomorphe (en tant que
$\partial$-foncteur) à $T_\bullet(M)$ (0, 11.9.3). On peut donc toujours se
ramener à la première hypothèse et dans ce cas les $Z'_i(P_\bullet)$ sont de
type fini; en outre, on peut se borner à démontrer les assertions relatives à
l'exactitude à gauche (cf. (7.4.2, $a)$). Cela étant, soit $x\in U$; comme le
foncteur $M\mapsto M_x$ est exact, on a
\[
  (Z'_p(P_\bullet))_x=Z'_p((P_\bullet)_x),
\]
et (en tenant compte de (7.4.1.3)) l'hypothèse entraîne, en vertu de
(7.4.2, $b)$) que $(Z'_p(P_\bullet))_x$ est un $A_x$-module plat, donc libre
puisqu'il est de type fini et que $A_x$ est un anneau local noethérien
(Bourbaki, \emph{Alg. comm.}, chap. II, \S 3, n\textsuperscript{o} 2, cor. 2
de la prop. 5). On en conclut qu'il y a un $f\in A$ tel que
$(Z'_p(P_\bullet))_f$ soit libre sur $A_f$ (Bourbaki, \emph{Alg. comm.},
chap. II, \S 5, n\textsuperscript{o} 1, cor. de la prop. 2), et \emph{a
fortiori} $(Z'_p(P_\bullet))_y$ est libre sur $A_y$ pour tout $y\in D(f)$, ce
qui achève la démonstration en vertu de (7.4.2, $b)$).
\end{proof}

\begin{corollary}[7.4.5]
Sous les hypothèses de (7.4.4), supposons en outre que $A$ soit intègre. Alors
l'ensemble $U$ des $x\in\operatorname{Spec}(A)$ tels que $(T_p)_x$ soit exact
est ouvert non vide.
\end{corollary}

\begin{proof}
Il suffit en effet de prouver que $(T_p)_x$ est exact pour le point générique
$x$ de $\operatorname{Spec}(A)$, ce qui est immédiat puisque $A_x$ est un
corps, donc tout foncteur additif sur $\mathbf{Ab}_{A_x}$ est exact.
\end{proof}

\begin{proposition}[7.4.6]
Sous les hypothèses générales de (7.4.4), les conditions $a)$, $b)$ et $c)$ de
(7.4.2) sont aussi équivalentes à :
\begin{enumerate}
  \item[d)] Il existe un $A$-module $Q$ et un isomorphisme fonctoriel
    \[
      T_p(M)\xrightarrow{\sim}\operatorname{Hom}_A(Q,M).
      \tag{7.4.6.1}\label{III.7.4.6.1.fr}
    \]
    En outre, le $A$-module $Q$ est déterminé à isomorphisme unique près par
    cette propriété, et il est de type fini.
\end{enumerate}
\end{proposition}

\begin{proof}
L'unicité de $Q$ est un cas particulier de l'unicité d'un objet représentatif
d'un foncteur représentable (0, 8.1.5). Il est clair que le second membre de
(7.4.6.1) est exact à gauche. Inversement, pour démontrer l'existence de $Q$,
lorsque $T_p$ est exact à gauche, on peut d'abord, comme dans (7.4.4), se
ramener au cas où les $P_i$ sont plats et de type fini, donc (puisque $A$ est
noethérien) projectifs de type fini (Bourbaki, \emph{Alg. comm.}, chap. II,
\S 5, n\textsuperscript{o} 2, cor. du th. 1). Le dual $\check P_i$ de $P_i$
est alors aussi un $A$-module projectif de type fini, $P_i$ est canoniquement
isomorphe au dual de $\check P_i$, et l'homomorphisme canonique
\[
  P_i\otimes_A M\longrightarrow\operatorname{Hom}_A(\check P_i,M)
\]
est bijectif (Bourbaki, \emph{Alg.}, chap. II, 3\textsuperscript{e} éd.,
\S 4, n\textsuperscript{o} 2, prop. 2). On sait d'autre part (7.4.2, $c)$)
qu'on peut supposer que $d_{p+1}:P_{p+1}\to P_p$ est nulle, donc que l'on a une
suite exacte
\[
  0\longrightarrow T_p(M)\xrightarrow{u}P_p\otimes M
  \xrightarrow{v}P_{p-1}\otimes M
\]
où $v=d_p\otimes1$. Posons alors $Q'=\operatorname{Ker}(d_p)$, de sorte qu'on
a la suite exacte
\[
  0\longrightarrow Q'\xrightarrow{w}P_p\xrightarrow{d_p}P_{p-1},
\]
d'où, par transposition, la suite exacte
\[
  \check P_{p-1}\xrightarrow{{}^t d_p}\check P_p
  \xrightarrow{{}^t w}\check Q'\longrightarrow0.
\]
Nous allons voir que
\[
  Q=\check Q'=\operatorname{Coker}({}^t d_p)
\]
répond à la question. En effet, on a la suite exacte
\[
  0\longrightarrow\operatorname{Hom}_A(Q,M)
  \longrightarrow\operatorname{Hom}_A(\check P_p,M)
  \xrightarrow{v'}\operatorname{Hom}_A(\check P_{p-1},M)
\]
\oldpage[III]{189}
où $v'=\operatorname{Hom}({}^t d_p,1)$; lorsqu'on identifie canoniquement
$P_i\otimes M$ à $\operatorname{Hom}(\check P_i,M)$, $v'$ s'identifie donc à
$v=d_p\otimes1$, et l'on a par suite l'isomorphisme fonctoriel
\[
  T_p(M)\xrightarrow{\sim}\operatorname{Hom}_A(Q,M)
\]
cherché. En outre, $Q$, étant un quotient de $\check P_p$, est de type fini.
\end{proof}

\begin{proposition}[7.4.7]
Supposons vérifiées les conditions générales de (7.4.4). Alors, pour tout
$A$-module $M$ de type fini :
\begin{enumerate}
  \item[(i)] Les $T_i(M)$ sont des $A$-modules de type fini.
  \item[(ii)] Pour tout idéal $\mathfrak m$ de $A$, l'homomorphisme canonique
    \[
      \widehat{T_i(M)}
      \longrightarrow
      \varprojlim_n T_i\bigl(M\otimes_A(A/\mathfrak m^{n+1})\bigr)
      \tag{7.4.7.1}\label{III.7.4.7.1.fr}
    \]
    (où le premier membre est le séparé complété de $T_i(M)$ pour la topologie
    $\mathfrak m$-préadique) est bijectif.
\end{enumerate}
\end{proposition}

\begin{proof}
Comme dans (7.4.4), on se ramène d'abord au cas où les $P_i$ sont de type fini;
$A$ étant noethérien, les sous-modules des $P_i\otimes_A M$ sont de type fini,
d'où trivialement l'assertion (i). Quant à l'assertion (ii), elle résulte plus
généralement du lemme suivant :
\end{proof}

\begin{lemma}[7.4.7.2]
Soient $A$ un anneau noethérien, $u:E\to F$ un homomorphisme de $A$-modules de
type fini. Pour tout $A$-module de type fini, posons
\[
  K(M)=\operatorname{Ker}(u\otimes1_M),
  \qquad
  C(M)=\operatorname{Coker}(u\otimes1_M);
\]
alors les homomorphismes canoniques
\[
  \widehat{K(M)}\longrightarrow\varprojlim_n K(M_n),
  \qquad
  \widehat{C(M)}\longrightarrow\varprojlim_n C(M_n)
  \tag{7.4.7.3}\label{III.7.4.7.3.fr}
\]
(où l'on a posé
$M_n=M\otimes_A(A/\mathfrak m^{n+1})=M/\mathfrak m^{n+1}M$) sont bijectifs
pour tout idéal $\mathfrak m$ de $A$.
\end{lemma}

\begin{proof}
Comme $E\otimes M$ et $F\otimes M$ sont de type fini, et que le foncteur
$M\mapsto\widehat M$ est exact dans la catégorie des $A$-modules de type fini
($0_{\mathrm I}$, 7.3.3), $\widehat{K(M)}$ et $\widehat{C(M)}$ sont
respectivement le noyau et le conoyau de
\[
  \widehat{(u\otimes1)}:
  \widehat{(E\otimes M)}\longrightarrow\widehat{(F\otimes M)}.
\]
L'exactitude à gauche du foncteur $\varprojlim$ montre donc que
$\widehat{K(M)}=\varprojlim_n K(M_n)$; d'autre part, l'exactitude à droite du
produit tensoriel prouve que
$C(M_n)=C(M)\otimes_A(A/\mathfrak m^{n+1})$, donc
$\widehat{C(M)}=\varprojlim_n C(M_n)$ par définition.
\end{proof}

\begin{remark}[7.4.8]
Compte tenu de (6.10.5) et (6.10.6), on voit que, moyennant une hypothèse de
platitude supplémentaire, (7.4.7) redonne le fait que (4.1.7.1) est un
isomorphisme, c'est-à-dire l'essentiel du « premier théorème de comparaison »
pour les morphismes propres; en outre, l'énoncé s'applique non plus seulement
à un $\mathscr O_X$-Module cohérent, mais à un complexe de tels Modules. Il
serait intéressant d'obtenir un énoncé comprenant à la fois (7.4.7) et
(4.1.7.1) comme cas particuliers. On notera que lorsque les $P_i$ sont de type
fini, la démonstration de (7.4.7) n'utilise pas le fait que ce sont des modules
plats; il y aurait lieu d'examiner si la conclusion de (7.4.7) est encore
valable lorsque les $P_i$ ne sont pas supposés plats ni de type fini, mais que
les $H_i(P_\bullet)$ sont supposés de type fini pour tout $i$ et nuls pour
$i<i_0$. Est-il alors possible de remplacer $P_\bullet$ par un complexe
$P'_\bullet$ de $A$-modules de type fini tel que les foncteurs
$H_\bullet(P_\bullet\otimes M)$ et $H_\bullet(P'_\bullet\otimes M)$ (qui ne
sont plus homologiques) soient encore isomorphes ?
\end{remark}

\oldpage[III]{190}
\subsection{Cas des anneaux locaux noethériens}
\label{subsection:III.7.5.fr}

\begin{env}[7.5.1]
Soient $A$ un anneau local noethérien, $\mathfrak m$ son idéal maximal, et pour
tout $A$-module $M$, désignons par $\widehat M$ son séparé complété pour la
topologie $\mathfrak m$-préadique, isomorphe à
\[
  \varprojlim_n\bigl(M\otimes_A(A/\mathfrak m^{n+1})\bigr)
  =
  \varprojlim_n\bigl(M/\mathfrak m^{n+1}M\bigr).
\]
Soit $T$ un foncteur covariant additif de $\mathbf{Ab}_A$ dans $\mathbf{Ab}$;
les homomorphismes canoniques (7.2.3.1)
\[
  T(M)\otimes_A(A/\mathfrak m^{n+1})
  \longrightarrow
  T\bigl(M\otimes_A(A/\mathfrak m^{n+1})\bigr)
\]
forment évidemment un système projectif de $A$-homomorphismes, qui donnent donc
à la limite un $\widehat A$-homomorphisme fonctoriel en $M$
\[
  \widehat{T(M)}
  \longrightarrow
  \varprojlim_n T(M_n)
  \tag{7.5.1.1}\label{III.7.5.1.1.fr}
\]
où l'on a posé
\[
  M_n=M\otimes_A(A/\mathfrak m^{n+1}),
  \qquad
  A_n=A/\mathfrak m^{n+1}.
\]
\end{env}

\begin{proposition}[7.5.2]
Soient $A$ un anneau local noethérien d'idéal maximal $\mathfrak m$,
$k=A/\mathfrak m$ son corps résiduel, $T$ un foncteur covariant additif de
$\mathbf{Ab}_A$ dans $\mathbf{Ab}$, semi-exact et commutant aux limites
inductives. On suppose en outre que pour tout $A$-module de type fini $M$,
$T(M)$ est un $A$-module de type fini et que l'homomorphisme canonique
(7.5.1.1) est un isomorphisme. Sous ces conditions, les propriétés suivantes
sont équivalentes :
\begin{enumerate}
  \item[a)] $T$ est exact à droite.
  \item[b)] Pour tout $n$, le foncteur $N\mapsto T(N)$ est exact à droite dans
    la catégorie des $A_n$-modules de type fini (ce qui revient à dire que $T$
    est exact à droite dans la catégorie des $A$-modules de longueur finie).
  \item[c)] L'homomorphisme canonique
    \[
      T(\varepsilon):T(A)\longrightarrow T(k)
    \]
    est surjectif.
  \item[d)] Pour tout $n$ assez grand, l'homomorphisme canonique
    \[
      T(A_n)\longrightarrow T(k)
    \]
    est surjectif.
\end{enumerate}
\end{proposition}

\begin{proof}
Il est clair que a) entraîne b). Montrons que b) entraîne a), c'est-à-dire que
si $u:M\to N$ est un épimorphisme de $A$-modules, $T(u)$ est un épimorphisme.
Comme $T$ commute aux limites inductives et que le foncteur $\varinjlim$ est
exact dans la catégorie des modules (pour des ensembles d'indices filtrants),
on peut se borner au cas où $M$ et $N$ sont de type fini. Comme $T(M)$ et
$T(N)$ sont alors de type fini, et que $A$ est un anneau local noethérien, il
suffit de montrer que
\[
  \widehat{T(u)}:
  \widehat{T(M)}
  \longrightarrow
  \widehat{T(N)}
\]
est surjectif ($0_{\mathrm I}$, 7.3.5 et $0_{\mathrm I}$, 6.4.1). Par hypothèse,
$\widehat{T(M)}$ et $\widehat{T(N)}$ sont respectivement $\varprojlim_n T(M_n)$
et $\varprojlim_n T(N_n)$, donc $\widehat{T(u)}$ est limite du système projectif
d'homomorphismes
\[
  T(u\otimes1_{A_n}):T(M_n)\longrightarrow T(N_n).
\]
Or, b) signifie que ces homomorphismes sont surjectifs; en outre, $T(M_n)$ est
un $A_n$-module de type fini, et $A_n$ est un anneau artinien par hypothèse; on
en conclut que $\widehat{T(u)}$ est surjectif ($0_{\mathrm I}$, 13.1.2 et
13.2.2). Il est clair que a) entraîne c), et comme $T(\varepsilon)$ se factorise
en
\[
  T(A)\longrightarrow T(A_n)\longrightarrow T(k),
\]
c) implique d); enfin, il résulte de (7.2.7) que b) et d) sont équivalentes
puisque $T$ est semi-exact dans $\mathbf{Ab}_{A_n}$, ce qui achève la
démonstration.
\end{proof}

\begin{corollary}[7.5.3]
Sous les hypothèses générales de (7.5.2), si $T(k)=0$, alors $T(M)=0$ pour tout
$A$-module $M$.
\end{corollary}

\begin{proof}
\oldpage[III]{191}
Comme $k$ est le seul $A$-module simple, on déduit de (7.3.5.4) que $T(E)=0$
pour tout $A$-module de longueur finie $E$. Si maintenant $M$ est de type fini,
$\widehat{T(M)}$ est isomorphe à $\varprojlim_n T(M_n)$, et comme les $M_n$ sont
de longueur finie, on a $\widehat{T(M)}=0$; comme $T(M)$ est de type fini par
hypothèse, il est isomorphe à un sous-module de $\widehat{T(M)}$
($0_{\mathrm I}$, 7.3.5), donc on a $T(M)=0$. Enfin, pour un $A$-module $M$
quelconque, $T(M)$ est limite inductive des $T(N_\alpha)$ pour les sous-modules
de type fini $N_\alpha$ de $M$, ce qui achève la démonstration.
\end{proof}

\begin{proposition}[7.5.4]
Soient $A$ un anneau local noethérien d'idéal maximal $\mathfrak m$,
$k=A/\mathfrak m$ son corps résiduel, $T_\bullet$ un foncteur homologique de
$\mathbf{Ab}_A$ dans $\mathbf{Ab}$, commutant aux limites inductives. On suppose
en outre que pour tout $i$ et tout $A$-module $M$ de type fini, $T_i(M)$ est de
type fini et l'homomorphisme canonique
\[
  \widehat{T_i(M)}
  \longrightarrow
  \varprojlim_n T_i(M_n)
\]
est bijectif. Pour un entier $p$ donné, les conditions suivantes sont alors
équivalentes :
\begin{enumerate}
  \item[a)] $T_p$ est exact.
  \item[b)] $T_p$ est exact à droite, et $T_p(A)$ est un $A$-module libre.
  \item[c)] Les homomorphismes canoniques
    \[
      T_{p+1}(A)\longrightarrow T_{p+1}(k)
      \quad\text{et}\quad
      T_p(A)\longrightarrow T_p(k)
    \]
    sont surjectifs.
  \item[d)] Pour tout $n$, les homomorphismes canoniques
    \[
      T_{p+1}(A_n)\longrightarrow T_{p+1}(k)
      \quad\text{et}\quad
      T_p(A_n)\longrightarrow T_p(k)
    \]
    sont surjectifs.
  \item[e)] Pour tout $n$, le foncteur $N\mapsto T_p(N)$ est exact dans la
    catégorie des $A_n$-modules de type fini.
\end{enumerate}
\end{proposition}

\begin{proof}
On sait (7.3.3) que a) équivaut à dire que $T_{p+1}$ et $T_p$ sont exacts à
droite; comme $T_\bullet$ est un foncteur homologique dans la catégorie
$\mathbf{Ab}_{A_n}$, le même raisonnement que dans (7.3.1) montre que e)
équivaut à dire que $T_p$ et $T_{p+1}$ sont exacts à droite dans la catégorie
des $A_n$-modules de type fini. On déduit donc de (7.5.2) que a) et e) sont
équivalentes; l'équivalence de a), c) et d) résulte aussi de (7.5.2). Enfin, on
sait que tout $A$-module plat de type fini est libre (Bourbaki,
\emph{Alg. comm.}, chap. II, \S 3, n\textsuperscript{o} 2, cor. 2 de la prop.
5); l'équivalence de a) et b) résulte alors de (7.3.1) et (7.3.3).
\end{proof}

\begin{corollary}[7.5.5]
Supposons vérifiées les conditions générales de (7.5.4).
\begin{enumerate}
  \item[(i)] Si $T_p(k)=0$, on a $T_p=0$, $T_{p+1}$ est exact à droite et
    $T_{p-1}$ est exact à gauche.
  \item[(ii)] Si $T_{p-1}(k)=T_{p+1}(k)=0$, $T_p$ est exact,
    l'homomorphisme canonique
    \[
      T_p(A)\otimes_A M\longrightarrow T_p(M)
    \]
    est bijectif et $T_p(A)$ est un $A$-module libre.
\end{enumerate}
\end{corollary}

\begin{proof}
(i) découle aussitôt de (7.5.3) puisque $T_p$ est semi-exact, la dernière
assertion résultant de la définition d'un foncteur homologique. On conclut
aussitôt de (i) les deux premières assertions de (ii), compte tenu de (7.3.3);
le fait que $T_p(A)$ soit libre résulte de (7.5.4).
\end{proof}

\begin{corollary}[7.5.6]
Supposons vérifiées les hypothèses générales de (7.5.4), et supposons de plus
que $A$ soit un anneau de valuation discrète.
\begin{enumerate}
  \item[(i)] Pour que $T_p$ soit exact à droite, il faut et il suffit que
    $T_{p-1}(A)$ soit un $A$-module libre.
  \item[(ii)] Pour que $T_p$ soit exact, il faut et il suffit que $T_p(A)$ et
    $T_{p-1}(A)$ soient des $A$-modules libres.
\end{enumerate}
\end{corollary}

\begin{proof}
\oldpage[III]{192}
Il est clair que (i) implique (ii) (cf. (7.3.3)). Pour prouver (i), remarquons
que si $f$ est un générateur de l'idéal maximal de $A$ (« uniformisante » de
$A$), pour qu'un $A$-module de type fini $M$ soit libre (ou plat, ce qui revient
au même), il faut et il suffit que l'homothétie $h_f:x\mapsto fx$ de $M$ soit
injective, car cela équivaut ici à dire que $M$ est sans torsion
($0_{\mathrm I}$, 6.3.4). Considérons alors la suite exacte
\[
  0\longrightarrow A\xrightarrow{h_f}A\longrightarrow k\longrightarrow0,
\]
qui fournit la suite exacte d'homologie
\[
  T_p(A)
  \longrightarrow T_p(k)
  \longrightarrow T_{p-1}(A)
  \xrightarrow{h_f}T_{p-1}(A).
\]
On voit que $T_{p-1}(A)$ est libre si et seulement si
$T_p(A)\to T_p(k)$ est surjectif; la conclusion résulte alors de (7.5.2).
\end{proof}

\begin{remark}[7.5.7]
On notera que, en vertu de (7.4.7), les hypothèses générales de (7.4.4)
entraînent que le foncteur homologique $T_\bullet$ défini par (7.4.1.1)
satisfait aux hypothèses générales de (7.5.4). Dans ce cas, (7.5.6) est donc
contenu dans (7.4.3).
\end{remark}

\subsection{Descente des propriétés d'exactitude. Théorème de semi-continuité
et critère d'exactitude de Grauert}
\label{subsection:III.7.6.fr}

\begin{proposition}[7.6.1]
Sous les conditions de (7.4.1), soit $B$ une $A$-algèbre commutative. Si $T_p$
est exact à droite (resp. exact à gauche, exact) il en est de même de
$T_p^{(B)}$; la réciproque est vraie lorsque $B$ est un $A$-module fidèlement
plat.
\end{proposition}

\begin{proof}
La première assertion est un cas particulier d'une assertion triviale de
(7.1.3). Inversement, supposons d'abord que $B$ soit un $A$-module plat. On a
alors, pour tout $A$-module $M$,
\[
  H_\bullet\bigl(P_\bullet\otimes_A(M\otimes_A B)\bigr)
  =
  \bigl(H_\bullet(P_\bullet\otimes_A M)\bigr)\otimes_A B,
\]
ce qui s'écrit aussi, pour tout $p$,
\[
  T_p(M)\otimes_A B
  =
  T_p^{(B)}(M_{(B)})
  \tag{7.6.1.1}\label{III.7.6.1.1.fr}
\]
à un isomorphisme canonique près. Supposons $T_p^{(B)}$ exact à droite (resp.
exact à gauche, exact); comme $M\mapsto M_{(B)}$ est un foncteur exact, le
premier membre de (7.6.1.1) est un foncteur exact à droite (resp. exact à
gauche, exact) en $M$; si maintenant $B$ est fidèlement plat sur $A$, on en
déduit que $T_p$ a la même propriété d'exactitude ($0_{\mathrm I}$, 6.4.1).
\end{proof}

\begin{proposition}[7.6.2]
Sous les conditions de (7.4.1), on suppose en outre que $A$ soit un anneau
noethérien réduit et que les $P_i$ soient des $A$-modules de type fini. Pour que
$T_p$ soit exact à droite (resp. exact à gauche, exact), il faut et il suffit
que, pour toute $A$-algèbre $B$ qui est un anneau de valuation discrète,
$T_p^{(B)}$ le soit.
\end{proposition}

\begin{proof}
En vertu de (7.3.1) et (7.3.3), on peut se borner à considérer l'exactitude à
droite, et il n'y a naturellement qu'à prouver la suffisance de la condition
(7.6.1). En vertu de (7.4.2), il suffit de montrer que
$Z'_{p-1}(P_\bullet)$ est un $A$-module plat; comme $P_{p-1}$ est de type fini,
$Z'_{p-1}(P_\bullet)$ est aussi de type fini; le critère ($0$, 10.2.8) montre
qu'il suffit alors que $Z'_{p-1}(P_\bullet)\otimes_A B$ soit un $B$-module plat
pour toute $A$-algèbre $B$ qui est un anneau de valuation discrète. Or, comme
$P_\bullet$ est un complexe de $A$-modules plats, on a
\[
  Z'_{p-1}(P_\bullet)\otimes_A B
  =
  Z'_{p-1}(P_\bullet\otimes_A B);
\]
\oldpage[III]{193}
$P_\bullet\otimes_A B$ est un complexe de $B$-modules plats
($0_{\mathrm I}$, 6.2.1), et pour tout $B$-module $N$, on a
\[
  H_\bullet(P_\bullet\otimes_A N)
  =
  H_\bullet\bigl((P_\bullet\otimes_A B)\otimes_B N\bigr),
\]
donc
\[
  T_p^{(B)}(N)
  =
  H_p\bigl((P_\bullet\otimes_A B)\otimes_B N\bigr);
\]
appliquant (7.4.2) à $T_p^{(B)}$, on voit que l'hypothèse que $T_p^{(B)}$ est
exact à droite est équivalente au fait que $Z'_{p-1}(P_\bullet\otimes_A B)$ est
un $B$-module plat.

Le critère précédent amène à étudier de plus près le cas des anneaux de
valuation discrète :
\end{proof}

\begin{proposition}[7.6.3]
Sous les conditions de (7.4.1), supposons que $A$ soit un anneau noethérien
régulier de dimension 1 (autrement dit, que $A$ est noethérien et que, pour tout
$x\in\operatorname{Spec}(A)$, $A_x$ est un corps ou un anneau de valuation
discrète). Alors, pour tout entier $p$ et tout $A$-module $M$, on a une suite
exacte canonique fonctorielle en $M$
\[
  0
  \longrightarrow T_p(A)\otimes_A M
  \xrightarrow{t_M}T_p(M)
  \longrightarrow \operatorname{Tor}_1^A\bigl(T_{p-1}(A),M\bigr)
  \longrightarrow0.
  \tag{7.6.3.1}\label{III.7.6.3.1.fr}
\]
\end{proposition}

\begin{proof}
Dans ce qui suit, nous supprimerons pour simplifier la mention du complexe
$P_\bullet$ dans les notations homologiques usuelles $H_p(P_\bullet)$,
$B_p(P_\bullet)$, $Z_p(P_\bullet)$ et $Z'_p(P_\bullet)$. On a les trois suites
exactes
\begin{gather*}
  0\longrightarrow H_p\longrightarrow Z'_p\longrightarrow B_{p-1}
  \longrightarrow0,\\
  0\longrightarrow B_{p-1}\longrightarrow Z_{p-1}\longrightarrow H_{p-1}
  \longrightarrow0,\\
  0\longrightarrow Z_{p-1}\longrightarrow P_{p-1}\longrightarrow B_{p-2}
  \longrightarrow0.
\end{gather*}
Comme $P_{p-1}$ et $P_{p-2}$ sont plats, il en est de même de leurs
sous-modules respectifs $B_{p-1}$, $Z_{p-1}$ et $B_{p-2}$ puisqu'il y a
identité entre $A_x$-modules plats et $A_x$-modules sans torsion (pour tout
$x\in\operatorname{Spec}(A)$); par tensorisation avec $M$, on a donc les suites
exactes
\[
  0=\operatorname{Tor}_1^A(B_{p-1},M)
  \longrightarrow H_p\otimes M
  \longrightarrow Z'_p\otimes M
  \xrightarrow{u}B_{p-1}\otimes M
  \longrightarrow0
  \tag{7.6.3.2}\label{III.7.6.3.2.fr}
\]
\[
  0=\operatorname{Tor}_1^A(Z_{p-1},M)
  \longrightarrow\operatorname{Tor}_1^A(H_{p-1},M)
  \longrightarrow B_{p-1}\otimes M
  \xrightarrow{v}Z_{p-1}\otimes M
  \tag{7.6.3.3}\label{III.7.6.3.3.fr}
\]
\[
  0=\operatorname{Tor}_1^A(B_{p-2},M)
  \longrightarrow Z_{p-1}\otimes M
  \xrightarrow{w}P_{p-1}\otimes M.
  \tag{7.6.3.4}\label{III.7.6.3.4.fr}
\]
Par définition,
$T_p(M)=\operatorname{Ker}(d_p\otimes1)/\operatorname{Im}(d_{p+1}\otimes1)$;
c'est donc le noyau de l'homomorphisme
\[
  (P_p\otimes M)/\operatorname{Im}(d_{p+1}\otimes1)
  \longrightarrow P_{p-1}\otimes M
\]
obtenu à partir de $d_p\otimes1$ par passage au quotient, homomorphisme qui
s'écrit aussi $Z'_p\otimes M\to P_{p-1}\otimes M$ par définition de
$Z'_p=P_p/B_p$; or, cet homomorphisme peut être considéré comme le composé
\[
  Z'_p\otimes M
  \xrightarrow{u}B_{p-1}\otimes M
  \xrightarrow{v}Z_{p-1}\otimes M
  \xrightarrow{w}P_{p-1}\otimes M.
\]
Comme $w$ est injectif d'après (7.6.3.4), on a une suite exacte
\[
  0\longrightarrow\operatorname{Ker}u
  \longrightarrow T_p(M)
  \longrightarrow\operatorname{Ker}v
  \longrightarrow0,
\]
qui n'est autre que (7.6.3.1), compte tenu de (7.6.3.2) et (7.6.3.3) et de ce
que $H_p=T_p(A)$ par définition.
\end{proof}

\begin{remarks}[7.6.4]
\begin{enumerate}
  \item[(i)] $H_\bullet(P_\bullet\otimes_A M)$ est l'homologie du bicomplexe
    $P_\bullet\otimes_A M$, où $M$ est considéré comme un complexe réduit à son
    terme de degré 0; elle est par suite (6.3.6 et 6.3.2) l'aboutissement de la
    suite spectrale régulière dont le terme $E_2$ est
    \[
      E^2_{pq}
      =
      \operatorname{Tor}_p^A\bigl(H_q(P_\bullet),M\bigr)
      =
      \operatorname{Tor}_p^A\bigl(T_q(A),M\bigr).
    \]
    \oldpage[III]{194}
    Or, l'hypothèse sur l'anneau $A$ entraîne que
    $\operatorname{Tor}_p^A(E,F)=0$ pour $p\geq2$ et pour des $A$-modules
    quelconques ($0_{\mathrm{IV}}$, 17.2.2); on sait (M, XV) que cela entraîne
    l'exactitude de la suite
    \[
      0\longrightarrow E^2_{0,q}
      \longrightarrow H_q(P_\bullet\otimes_A M)
      \longrightarrow E^2_{1,q-1}
      \longrightarrow0
    \]
    qui n'est autre que (7.6.3.1).
  \item[(ii)] Compte tenu de (7.3.1), la suite exacte (7.6.3.1) redonne comme
    cas particulier le résultat de (7.4.3).
\end{enumerate}
\end{remarks}

\begin{corollary}[7.6.5]
Sous les conditions de (7.4.1), supposons que $A$ soit un anneau de valuation
discrète, de corps des fractions $K$, de corps résiduel $k$, et que les $T_i(A)$
soient des $A$-modules de type fini. On a alors
\[
  \operatorname{rang}_k T_p(k)
  \geq
  \operatorname{rang}_k\bigl(T_p(A)\otimes_A k\bigr)
  \geq
  \operatorname{rang}_A T_p(A)
  =
  \operatorname{rang}_K T_p(K).
  \tag{7.6.5.1}\label{III.7.6.5.1.fr}
\]
En outre, pour que les termes extrêmes de cette inégalité soient égaux, il faut
et il suffit que $T_p$ soit exact, ou encore que $T_p(A)$ et $T_{p-1}(A)$ soient
des $A$-modules libres.
\end{corollary}

\begin{proof}
Faisons en effet $M=k$ dans la suite exacte (7.6.3.1); il vient, puisqu'il
s'agit d'espaces vectoriels sur $k$
\[
  \operatorname{rang}_k T_p(k)
  =
  \operatorname{rang}_k\bigl(T_p(A)\otimes_A k\bigr)
  +
  \operatorname{rang}_k\bigl(\operatorname{Tor}_1^A(T_{p-1}(A),k)\bigr).
\]
D'autre part, comme $T_p(A)$ est un module de type fini sur l'anneau local
intègre $A$, on a (Bourbaki, \emph{Alg. comm.}, chap. II, \S 3,
n\textsuperscript{o} 2, cor. 1 de la prop. 4)
\[
  \operatorname{rang}_k\bigl(T_p(A)\otimes_A k\bigr)
  \geq
  \operatorname{rang}_A T_p(A)
  =
  \operatorname{rang}_K\bigl(T_p(A)\otimes_A K\bigr)
  \tag{7.6.5.2}\label{III.7.6.5.2.fr}
\]
et en outre les deux membres de (7.6.5.2) sont égaux si et seulement si $T_p(A)$
est un $A$-module libre (\emph{loc. cit.}, prop. 7). On notera d'ailleurs que
puisque $K$ est un $A$-module plat, on a par définition
\[
  T_p(A)\otimes_A K
  =
  H_p(P_\bullet)\otimes_A K
  =
  H_p(P_\bullet\otimes_A K)
  =
  T_p(K).
\]
On a donc bien l'inégalité (7.6.5.1) et on voit en outre que l'égalité n'est
possible que si : 1\textsuperscript{o} $T_p(A)$ est libre;
2\textsuperscript{o} $\operatorname{Tor}_1^A(T_{p-1}(A),k)=0$, condition qui
équivaut, comme on sait ($0$, 10.1.3), au fait que $T_{p-1}(A)$ est un
$A$-module libre. Enfin, comme les $T_i(A)$ sont des $A$-modules de type fini,
il revient au même de dire qu'ils sont plats ou libres (Bourbaki,
\emph{Alg. comm.}, chap. II, \S 3, n\textsuperscript{o} 2, cor. 2 de la prop.
5), et on conclut par (7.4.3).
\end{proof}

\begin{env}[7.6.6]
Les hypothèses étant toujours celles de (7.4.1), nous poserons, pour tout
$x\in\operatorname{Spec}(A)$,
\[
  d_p(x)=d_p^T(x)=\operatorname{rang}_{\kappa(x)}T_p(\kappa(x)).
  \tag{7.6.6.1}\label{III.7.6.6.1.fr}
\]
\end{env}

\begin{lemma}[7.6.7]
Soit $\varphi:A\to A'$ un homomorphisme d'anneaux, et soit
\[
  f={}^a\!\varphi:
  \operatorname{Spec}(A')\longrightarrow\operatorname{Spec}(A)
\]
l'application correspondante (I, 1.2.1). Si l'on pose
$T'_\bullet=T_\bullet^{(A')}$ (7.1.3), on a
\[
  d_p^{T'}=d_p^T\circ f.
  \tag{7.6.7.1}\label{III.7.6.7.1.fr}
\]
\end{lemma}

\begin{proof}
\oldpage[III]{195}
En effet, pour tout $x'\in\operatorname{Spec}(A')$, on a, en posant
$x=f(x')$,
\[
  H_\bullet(P_\bullet\otimes_A\kappa(x'))
  =
  H_\bullet\bigl((P_\bullet\otimes_A\kappa(x))
    \otimes_{\kappa(x)}\kappa(x')\bigr)
  =
  H_\bullet(P_\bullet\otimes_A\kappa(x))
    \otimes_{\kappa(x)}\kappa(x'),
\]
puisque $\kappa(x')$ est plat sur $\kappa(x)$, d'où la relation (7.6.7.1).
\end{proof}

\begin{lemma}[7.6.8]
Si l'anneau $A$ est noethérien et le complexe $P_\bullet$ formé de $A$-modules
de type fini, la fonction $x\mapsto d_p^T(x)$ sur $\operatorname{Spec}(A)$ est
constructible.
\end{lemma}

\begin{proof}
Il faut prouver que pour toute partie fermée irréductible $Y$ de
$X=\operatorname{Spec}(A)$, il existe un ouvert non vide $U$ de $Y$ dans lequel
$d_p$ est constante ($0$, 9.2.2); comme
$Y=\operatorname{Spec}(A/\mathfrak a)$, où $\mathfrak a$ est un idéal de $A$ tel
que $A/\mathfrak a$ soit réduit, on peut, en vertu de (7.6.7), se borner au cas
où $Y=X$ et où $A$ est un anneau noethérien intègre; mais alors l'assertion
résulte de (7.4.5).
\end{proof}

\begin{theorem}[7.6.9]
Soient $A$ un anneau noethérien, $P_\bullet$ un complexe de $A$-modules plats de
type fini, $T_\bullet(M)=H_\bullet(P_\bullet\otimes_A M)$ le foncteur homologique
défini par $P_\bullet$; pour tout $x\in\operatorname{Spec}(A)$, soit
\[
  d_p(x)=\operatorname{rang}_{\kappa(x)}T_p(\kappa(x)).
\]
Alors :
\begin{enumerate}
  \item[(i)] La fonction $d_p$ est constructible et semi-continue
    supérieurement dans $\operatorname{Spec}(A)$.
  \item[(ii)] Si $T_p$ est exact, $d_p$ est continue (donc localement constante)
    dans $\operatorname{Spec}(A)$; la réciproque est vraie lorsque l'anneau $A$
    est réduit.
\end{enumerate}
\end{theorem}

\begin{proof}
\begin{enumerate}
  \item[(i)] La première assertion a été démontrée dans (7.6.8). Pour prouver
    la seconde, il suffit donc ($0$, 9.2.4) de montrer que si $x'\ne x$ est une
    générisation de $x$ dans $\operatorname{Spec}(A)$, on a
    $d_p(x')\leq d_p(x)$. Or, il existe alors un anneau de valuation discrète
    $B$ et un morphisme
    \[
      f:\operatorname{Spec}(B)\longrightarrow\operatorname{Spec}(A)
    \]
    tels que, si $a$ désigne le point fermé de $\operatorname{Spec}(B)$ et $b$
    son point générique, on ait $f(a)=x$ et $f(b)=y$ (II, 7.1.9). En vertu de
    la formule (7.6.7.1), on voit qu'on est ramené à démontrer l'inégalité
    $d_p(a)\geq d_p(b)$ dans $\operatorname{Spec}(B)$; mais cela n'est autre que
    l'inégalité (7.6.5.1)\footnote{Le principe de démonstration de (i) par
    réduction au cas d'un anneau de valuation discrète nous a été communiqué
    oralement par Hironaka.}.
  \item[(ii)] La première assertion a déjà été démontrée (7.3.4). Pour démontrer
    la réciproque, utilisons le critère valuatif (7.6.2); compte tenu de la
    formule (7.6.7.1), on est donc ramené au cas où $A$ est un anneau de
    valuation discrète; mais comme $\operatorname{Spec}(A)$ ne comporte alors
    que deux points, l'hypothèse que $d_p$ est constante implique bien que $T_p$
    est exact, en vertu de (7.6.5).
\end{enumerate}
\end{proof}

\begin{corollary}[7.6.10]
Soient $A$ un anneau noethérien, $\mathfrak p_i$ $(1\leq i\leq r)$ ses idéaux
premiers minimaux, $k_i$ le corps résiduel de $A_{\mathfrak p_i}$
$(1\leq i\leq r)$.
\begin{enumerate}
  \item[(i)] Pour tout $x\in\operatorname{Spec}(A)$, il existe un indice $i$
    tel que
    \[
      d_p(x)\geq\operatorname{rang}_{k_i}T_p(k_i).
      \tag{7.6.10.1}\label{III.7.6.10.1.fr}
    \]
\end{enumerate}
En particulier, si $A$ est intègre et si $K$ est son corps des fractions, on a
\[
  d_p(x)\geq\operatorname{rang}_K T_p(K)
  \tag{7.6.10.2}\label{III.7.6.10.2.fr}
\]
pour tout $x\in\operatorname{Spec}(A)$.
\oldpage[III]{196}
\begin{enumerate}
  \item[(ii)] Supposons en outre que $A$ soit local et réduit, et soit $k$ son
    corps résiduel. Alors, pour que $T_p$ soit exact, il faut et il suffit que
    l'on ait
    \[
      \operatorname{rang}_k T_p(k)
      =
      \operatorname{rang}_{k_i}T_p(k_i)
      \qquad\text{pour }1\leq i\leq r.
      \tag{7.6.10.3}\label{III.7.6.10.3.fr}
    \]
\end{enumerate}
\end{corollary}

\begin{proof}
(i) est immédiat puisque tout voisinage de $x$ contient un des $\mathfrak p_i$,
et il suffit d'appliquer la définition de la semi-continuité. D'autre part, si
$A$ est local, le seul voisinage dans $\operatorname{Spec}(A)$ de l'idéal
maximal $\mathfrak m$ est $\operatorname{Spec}(A)$ tout entier, donc on a
\[
  d_p(x)\leq\operatorname{rang}_k T_p(k)
\]
pour tout $x\in\operatorname{Spec}(A)$; cette relation, jointe à (i), montre que
la condition (7.6.10.3) entraîne que $d_p(x)$ est constante dans
$\operatorname{Spec}(A)$, et par suite que $T_p$ est exact en vertu de
(7.6.9, (ii)); la réciproque est évidente en vertu de (7.6.9, (ii)).
\end{proof}

\begin{remark}[7.6.11]
On peut se demander si l'assertion de (7.6.9, (i)) ne peut être renforcée par
l'inégalité
\[
  \operatorname{rang}_{\kappa(x)}T_p(\kappa(x))
  \geq
  \operatorname{rang}_{\kappa(x)}
    \bigl(T_p(A)\otimes_A\kappa(x)\bigr)
  \tag{7.6.11.1}\label{III.7.6.11.1.fr}
\]
pour tout $x\in\operatorname{Spec}(A)$, qui a lieu effectivement lorsque $A$
est un anneau de valuation discrète et $x$ son idéal maximal (7.6.5).
Bornons-nous au cas où $A$ est un anneau local noethérien, d'idéal maximal
$\mathfrak m$ et de corps résiduel $k$. Alors, les conditions suivantes sont
équivalentes :
\begin{enumerate}
  \item[a)] Pour tout complexe $P_\bullet$ de $A$-modules plats de type fini,
    on a
    \[
      \operatorname{rang}_k T_p(k)
      \geq
      \operatorname{rang}_k\bigl(T_p(A)\otimes_A k\bigr)
      \quad\text{pour tout entier }p.
      \tag{7.6.11.2}\label{III.7.6.11.2.fr}
    \]
  \item[b)] Pour tout $A$-module $M$ de type fini, on a
    \[
      \operatorname{rang}_k(M\otimes_A k)
      \geq
      \operatorname{rang}_k(\check M\otimes_A k).
      \tag{7.6.11.3}\label{III.7.6.11.3.fr}
    \]
  \item[c)] Pour tout $A$-module $N$ de type fini, on a
    \[
      \operatorname{rang}_k\bigl(\operatorname{Tor}_1^A(N,k)\bigr)
      \geq
      \operatorname{rang}_k\bigl(\operatorname{Tor}_2^A(N,k)\bigr).
      \tag{7.6.11.4}\label{III.7.6.11.4.fr}
    \]
\end{enumerate}
On notera qu'il revient au même, par décalage (M, V, 7.2), de dire que l'on a,
pour tout $i\geq1$,
\[
  \operatorname{rang}_k\bigl(\operatorname{Tor}_i^A(N,k)\bigr)
  \geq
  \operatorname{rang}_k\bigl(\operatorname{Tor}_{i+1}^A(N,k)\bigr).
  \tag{7.6.11.5}\label{III.7.6.11.5.fr}
\]
Donnons rapidement quelques indications sur la démonstration. Pour voir que a)
entraîne b), on considère une suite exacte
\[
  L_1\xrightarrow{d}L_0\longrightarrow M\longrightarrow0
\]
où $L_0$ et $L_1$ sont libres de type fini, et on applique a) au complexe
\[
  P_1\xrightarrow{{}^t d}P_0,
  \qquad
  P_0=\check L_1,
  \quad
  P_1=\check L_0,
\]
les autres termes étant nuls; on a alors $T_1(A)=\check M$ et
\[
  T_1(k)
  =
  \operatorname{Hom}_A(M,k)
  =
  \operatorname{Hom}_k(M/\mathfrak mM,k),
\]
autrement dit $T_1(k)$ est le dual de l'espace vectoriel $M\otimes_A k$, et a
donc même rang que ce dernier. Pour prouver que b) entraîne c), nous établissons
d'abord le lemme suivant :

\begin{lemma}[7.6.11.6]
Étant donné un complexe
\[
  \cdots\longrightarrow0\longrightarrow
  P_1\xrightarrow{d}P_0
  \longrightarrow0\longrightarrow\cdots
\]
de $A$-modules plats, on a une suite exacte
\[
  0\longrightarrow\operatorname{Tor}_2^A(Z'_0,k)
  \longrightarrow T_1(A)\otimes_A k
  \longrightarrow T_1(k)
  \longrightarrow\operatorname{Tor}_1^A(Z'_0,k)
  \longrightarrow0.
  \tag{7.6.11.7}\label{III.7.6.11.7.fr}
\]
\end{lemma}

\begin{proof}
\oldpage[III]{197}
En effet, partant de la suite exacte
$0\to Z_1\to P_1\to B_0\to0$, on en tire la suite exacte
\[
  0\longrightarrow\operatorname{Tor}_1^A(B_0,k)
  \longrightarrow Z_1\otimes k
  \longrightarrow P_1\otimes k
  \xrightarrow{u}B_0\otimes k
  \longrightarrow0.
\]
De la suite exacte $0\to B_0\to Z_0\to Z'_0\to0$, on tire, puisque $Z_0=P_0$
est plat,
$\operatorname{Tor}_1^A(B_0,k)=\operatorname{Tor}_2^A(Z'_0,k)$; par définition,
on a $Z_1=T_1(A)$; enfin, on a $T_1(k)=\operatorname{Ker}(d\otimes1)$, et
$d\otimes1$ se factorise en
\[
  P_1\otimes k
  \xrightarrow{u}B_0\otimes k
  \xrightarrow{v}Z_0\otimes k
  =P_0\otimes k;
\]
on a $T_1(k)=u^{-1}(R)$, où $R=\operatorname{Ker}v$, et comme $u$ est
surjectif, $R=u(T_1(k))$; enfin,
$R=\operatorname{Tor}_1^A(Z'_0,k)$ par définition de $v$, ce qui achève
d'établir la suite exacte (7.6.11.7).
\end{proof}

Pour déduire alors c) de b), on considère une suite exacte
\[
  L_1\xrightarrow{d}L_0\longrightarrow N\longrightarrow0,
\]
où $L_0$ et $L_1$ sont des modules libres de type fini; considérons le foncteur
$T$ associé au complexe formé de $L_1$ et $L_0$; comme $L_0$ et $L_1$ sont
libres, ils s'identifient à leurs biduals; donc si
\[
  M=\operatorname{Coker}({}^t d),
  \qquad
  T_1(A)=\operatorname{Ker}(d)=\check M;
\]
par ailleurs, $M\otimes_A k=\operatorname{Coker}({}^t d\otimes1_k)$ a même rang
sur $k$ que $\operatorname{Ker}(d\otimes1_k)$. L'hypothèse b) entraîne par suite
que
\[
  \operatorname{rang}_k\bigl(T_1(A)\otimes_A k\bigr)
  \leq
  \operatorname{rang}_k\bigl(T_1(k)\bigr);
\]
comme $Z'_0=N$, l'inégalité (7.6.11.4) résulte donc de la suite exacte
(7.6.11.7).

Enfin, pour prouver que c) entraîne a), appliquons (7.6.11.6) en remplaçant
$P_0$ et $P_1$ par $P_p$ et $P_{p-1}$; l'hypothèse c) appliquée au module
$Z'_p$ donne $\operatorname{rang}_kR\geq\operatorname{rang}_kS$, où
\[
  R=\operatorname{Ker}(P_p\otimes k\longrightarrow P_{p-1}\otimes k)
  \quad\text{et}\quad
  S=Z_p\otimes k.
\]
Or, si l'on factorise $d_{p+1}:P_{p+1}\to P_p$ en
\[
  P_{p+1}\xrightarrow{v}Z_p\xrightarrow{j}P_p,
\]
on a
\[
  \operatorname{Im}(d_{p+1}\otimes1)
  =
  (j\otimes1)\bigl(\operatorname{Im}(v\otimes1)\bigr).
\]
Comme
\[
  T_p(A)\otimes_A k
  =(Z_p/B_p)\otimes_A k
  =(Z_p\otimes k)/\operatorname{Im}(v\otimes1),
\]
et
\[
  T_p(k)=R/\operatorname{Im}(d_{p+1}\otimes1),
\]
on en conclut bien l'inégalité (7.6.11.2).

Cela étant, supposons que l'anneau local $A$ soit régulier de dimension $n$;
on sait alors [17] que le $A$-module
$\operatorname{Tor}_i^A(k,k)$ est isomorphe à la puissance extérieure
\[
  \bigwedge\nolimits^i(\mathfrak m/\mathfrak m^2);
\]
on voit donc que la condition (7.6.11.4) n'est pas vérifiée pour $N=k$, dès que
$n\geq4$.

Par contre, si l'anneau local intègre $A$ est tel que tout $A$-module réflexif
de type fini soit libre (ce qui est le cas lorsque $A$ est un anneau régulier
de dimension 2), la condition (7.6.11.3) est vérifiée : en effet, on sait que le
dual $\check M$ d'un $A$-module $M$ de type fini est réflexif, donc libre, et
par suite
\[
  \operatorname{rang}_k(\check M\otimes_A k)
  =
  \operatorname{rang}_K(\check M)
  =
  \operatorname{rang}_K(M),
\]
($K$ corps des fractions de $A$); par ailleurs, on sait que toute base sur $k$
de $M\otimes_A k$ est formée d'images d'un système de générateurs de $M$
(Bourbaki, \emph{Alg. comm.}, chap. II, \S 3, n\textsuperscript{o} 2, cor. 2 de
la prop. 4), donc
$\operatorname{rang}_K(M)\leq\operatorname{rang}_k(M\otimes_A k)$, ce qui
prouve notre assertion.
\end{remark}

\subsection{Application aux morphismes propres : I. La propriété d'échange}
\label{subsection:III.7.7.fr}

Les trois numéros qui suivent sont, pour l'essentiel, des traductions, dans le
langage des morphismes de préschémas, des résultats des numéros précédents.

\begin{env}[7.7.1]
Soit $f:X\to Y$ un morphisme quasi-compact et séparé de préschémas, et soit
$\mathscr P_\bullet$ un complexe de $\mathscr O_X$-Modules quasi-cohérents, dont
l'opérateur de dérivation est de degré $-1$; supposons en outre que les
$\mathscr O_X$-Modules $\mathscr P_i$ sont $Y$-plats ($0_{\mathrm I}$, 6.7.1).
\oldpage[III]{198}
Nous allons considérer le $\partial$-foncteur
$\mathscr M\mapsto\mathscr T_\bullet(\mathscr M)$ (aussi noté
$\mathscr T_\bullet(\mathscr P_\bullet,\mathscr M)$) dans la catégorie des
$\mathscr O_Y$-Modules quasi-cohérents, à valeurs dans la catégorie des
$\mathscr O_Y$-Modules quasi-cohérents (en vertu de (6.2.3)), défini par
\[
  \mathscr T_n(\mathscr P_\bullet,\mathscr M)
  =\mathscr T_n(\mathscr M)
  =\mathscr H^{-n}
    \bigl(f,\mathscr P^\bullet\otimes_{\mathscr O_Y}\mathscr M\bigr)
  \qquad\text{pour }n\in\mathbf Z,
  \tag{7.7.1.1}\label{III.7.7.1.1.fr}
\]
où $\mathscr P^\bullet$ est le complexe dont le terme de degré $j$ est
$\mathscr P_{-j}$, l'opérateur de dérivation étant donc alors de degré $+1$. Le foncteur
$\mathscr T_\bullet$ ainsi défini est un \emph{foncteur homologique} en
$\mathscr M$ (6.2.6).
\end{env}

\begin{env}[7.7.2]
Soit $g:Y'\to Y$ un morphisme, et posons
\[
  X'=X_{(Y')}=X\times_Y Y',
  \qquad
  f'=f_{(Y')}:X'\longrightarrow Y',
\]
qui est un morphisme quasi-compact et séparé; soit d'autre part
\[
  \mathscr P'_\bullet
  =\mathscr P_\bullet\otimes_{\mathscr O_Y}\mathscr O_{Y'};
\]
c'est un complexe de $\mathscr O_{X'}$-Modules quasi-cohérents qui sont
$Y'$-plats en vertu de (I, 9.1.12) et ($0_{\mathrm I}$, 6.2.1). Nous poserons
(avec les mêmes conventions sur les degrés)
\[
  \mathscr T^{Y'}_\bullet(\mathscr M')
  =\mathscr H^\bullet
    \bigl(f',\mathscr P'^\bullet\otimes_{\mathscr O_{Y'}}\mathscr M'\bigr)
  =\mathscr H^\bullet
    \bigl(f',\mathscr P^\bullet\otimes_{\mathscr O_Y}\mathscr M'\bigr)
  \tag{7.7.2.1}\label{III.7.7.2.1.fr}
\]
qui est un foncteur homologique en le $\mathscr O_{Y'}$-Module quasi-cohérent
$\mathscr M'$. Lorsque $Y'$ est un schéma affine d'anneau $A'$, on écrira
$\mathscr T^{A'}_\bullet$ au lieu de $\mathscr T^{Y'}_\bullet$; pour tout
$A'$-module $M'$, on a alors
\[
  \mathscr T^{A'}_\bullet(\widetilde{M'})
  =\bigl(\Gamma(Y',\mathscr T^{Y'}_\bullet(\widetilde{M'}))\bigr)^\sim;
\]
on posera
\[
  T^{A'}_\bullet(M')
  =\Gamma(Y',\mathscr T^{Y'}_\bullet(\widetilde{M'})),
\]
qui est un foncteur homologique de $A'$-modules, à valeurs dans la catégorie des
$A'$-modules.

On observera que si $Y=\operatorname{Spec}(A)$ est aussi affine, le foncteur de
$A'$-modules $T^{A'}_\bullet$ coïncide avec le foncteur obtenu par extension des
scalaires de $A$ à $A'$ à partir du foncteur homologique de $A$-modules
$T^A_\bullet$ (7.1.3) : en effet, soit $g:Y'\to Y$ le morphisme correspondant à
l'homomorphisme d'anneaux $A\to A'$, et soit $g':X'\to X$ le morphisme
correspondant qui est affine (II, 1.6.2); si $\mathfrak U$ est un recouvrement
ouvert affine de $X$, $\mathfrak U'=g'^{-1}(\mathfrak U)$ est un recouvrement
ouvert affine de $X'$; en vertu de (6.2.2), tout revient à voir que
\[
  C^\bullet
    \bigl(\mathfrak U,
      \mathscr P_\bullet\otimes_{\mathscr O_Y}g_*(\mathscr M')\bigr)
  =C^\bullet
    \bigl(\mathfrak U',
      \mathscr P_\bullet\otimes_{\mathscr O_Y}\mathscr M'\bigr),
\]
et finalement, que pour tout ouvert affine $U$ de $X$, en posant
$U'=g'^{-1}(U)$, on a
\[
  \Gamma
    \bigl(U,\mathscr P_\bullet\otimes_{\mathscr O_Y}g_*(\mathscr M')\bigr)
  =\Gamma
    \bigl(U',\mathscr P_\bullet\otimes_{\mathscr O_Y}\mathscr M'\bigr),
\]
ce qui est trivial (I, 1.3 et 3.2).

En particulier, si $U$ est un ouvert de $Y$, on a, pour tout
$\mathscr O_Y$-Module quasi-cohérent $\mathscr M$
\[
  \mathscr T^U_\bullet(\mathscr M|U)
  =\bigl(\mathscr T_\bullet(\mathscr M)\bigr)|U.
  \tag{7.7.2.2}\label{III.7.7.2.2.fr}
\]
\end{env}

\begin{env}[7.7.3]
Pour tout $\mathscr O_Y$-Module quasi-cohérent $\mathscr M$, on a un
homomorphisme canonique, fonctoriel en $\mathscr M$ :
\[
  \mathscr T_p(\mathscr O_Y)\otimes_{\mathscr O_Y}\mathscr M
  \longrightarrow
  \mathscr T_p(\mathscr M).
  \tag{7.7.3.1}\label{III.7.7.3.1.fr}
\]
En effet, si $Y$ est affine, cet homomorphisme a été défini en (7.2.2); cette
définition s'étend sans peine au cas général, en remarquant que si $U$, $V$ sont
deux ouverts affines de $Y$ tels que $V\subset U$, le diagramme
\oldpage[III]{199}
\[
\xymatrix@C=18pt@R=32pt{
  {\begin{gathered}
    \bigl(\mathscr T_p(\mathscr O_Y)
      \otimes_{\mathscr O_Y}\mathscr M\bigr)|U\\
    =\mathscr T^U_p(\mathscr O_Y|U)
      \otimes_{\mathscr O_Y|U}(\mathscr M|U)
  \end{gathered}}
  \ar[r]\ar[d]
  &
  {\begin{gathered}
    \mathscr T^U_p(\mathscr M|U)\\
    =\bigl(\mathscr T_p(\mathscr M)\bigr)|U
  \end{gathered}}
  \ar[d]
  \\
  {\begin{gathered}
    \bigl(\mathscr T_p(\mathscr O_Y)
      \otimes_{\mathscr O_Y}\mathscr M\bigr)|V\\
    =\mathscr T^V_p(\mathscr O_Y|V)
      \otimes_{\mathscr O_Y|V}(\mathscr M|V)
  \end{gathered}}
  \ar[r]
  &
  {\begin{gathered}
    \mathscr T^V_p(\mathscr M|V)\\
    =\bigl(\mathscr T_p(\mathscr M)\bigr)|V
  \end{gathered}}
}
\]
est commutatif par (7.2.3.3).

Pour tout morphisme $g:Y'\to Y$ on a un homomorphisme canonique
\[
  \mathscr T_p(\mathscr O_Y)\otimes_{\mathscr O_Y}\mathscr O_{Y'}
  \longrightarrow
  \mathscr T^{Y'}_p(\mathscr O_{Y'}),
  \tag{7.7.3.2}\label{III.7.7.3.2.fr}
\]
qui n'est autre que le cas particulier de (6.7.11.2) (pour les aboutissements)
dans le cas où $S=Y$, $v_1=f$, $v_2=1_Y$, $\mathscr P^{(2)}_\bullet$ réduit au
seul terme $\mathscr M$ de degré 0.

Lorsque $Y=\operatorname{Spec}(A)$, $Y'=\operatorname{Spec}(A')$ sont affines,
(7.7.3.2) n'est autre que l'homomorphisme de faisceaux correspondant à
l'homomorphisme canonique de $A'$-modules défini dans (7.2.2)
\[
  T^A_p(A)\otimes_A A'
  \longrightarrow
  T^{A'}_p(A')=T^A_p(A')
\]
comme il résulte aisément de (6.7.11) (car dans le cas envisagé, on peut prendre
$\mathfrak U'^{(i)}=u_i^{-1}(\mathfrak U^{(i)})$ dans (6.7.11)).
\end{env}

\begin{env}[7.7.4]
Lorsque $f$ est un morphisme \emph{propre}, $Y=\operatorname{Spec}(A)$ un schéma
affine noethérien et $\mathscr P_\bullet$ un complexe de
$\mathscr O_X$-Modules cohérents et $Y$-plats \emph{limité inférieurement}, on a
vu (6.10.5) que l'on peut écrire à un isomorphisme près,
\[
  \mathscr T_p(\mathscr M)
  =\mathscr H_p
    \bigl(\mathscr L_\bullet\otimes_{\mathscr O_Y}\mathscr M\bigr),
  \qquad
  \mathscr L_\bullet=\widetilde{L_\bullet},
\]
où $L_\bullet$ est un complexe de $A$-modules \emph{libres de type fini limité
inférieurement}; le foncteur $\mathscr T_\bullet$ est donc du type qui a été
étudié en détail dans (7.4) et (7.6). Nous allons traduire les résultats de cette
étude :
\end{env}

\begin{theorem}[7.7.5]
Soient $Y$ un préschéma localement noethérien, $(U_\alpha)$ un recouvrement de
$Y$ formé d'ouverts affines, $f:X\to Y$ un morphisme propre,
$\mathscr P_\bullet$ un complexe de $\mathscr O_X$-Modules cohérents et $Y$-plats
limité inférieurement. Le foncteur homologique $\mathscr T_\bullet(\mathscr M)$
défini par (7.7.1.1) possède alors les propriétés suivantes :

\medskip
\noindent
\textnormal{I) (La propriété de semi-continuité).}\footnote{Un cas particulier de
ce théorème se trouve déjà dans la note [3] de Chow-Igusa. La propriété de
semi-continuité a été découverte, dans le cadre des espaces analytiques (et sous
des hypothèses assez particulières), par Kodaira-Spencer (« On the variations of
almost-complex structures », \emph{Algebraic Geometry and Topology, A Symposium in
honor of S. Lefschetz}, Princeton Series n\textsuperscript{o} 12, p. 139--150,
Princeton, 1957) et la version générale démontrée par Grauert [5].}
\emph{La fonction}
\[
  y\longmapsto
  d_p(y)
  =
  \operatorname{rang}_{\kappa(y)}
    \mathscr T_p^{\kappa(y)}(\kappa(y))
  \tag{7.7.5.1}\label{III.7.7.5.1.fr}
\]
\emph{est semi-continue supérieurement.}

\oldpage[III]{200}
\medskip
\noindent
\textnormal{II) (La propriété d'échange).}
\emph{Pour un entier $p$ donné, les conditions suivantes sont équivalentes :}
\begin{enumerate}
  \item[a)] $\mathscr T_p$ est exact à droite.

  \item[a$'$)] $\mathscr T_p(\mathscr M)$ est isomorphe à un foncteur de la forme
    $\mathscr N\otimes_{\mathscr O_Y}\mathscr M$, $\mathscr N$ étant nécessairement
    isomorphe à
    \[
      \mathscr T_p(\mathscr O_Y)
      =
      \mathscr H^{-p}(f,\mathscr P^\bullet).
    \]

  \item[a$''$)] L'homomorphisme canonique fonctoriel (7.7.3.1) est un
    isomorphisme.

  \item[b)] $\mathscr T_{p-1}$ est exact à gauche.

  \item[b$'$)] Il existe un $\mathscr O_Y$-Module $\mathscr Q$ (nécessairement
    cohérent, et déterminé à un isomorphisme unique près) et un isomorphisme de
    foncteurs
    \[
      \mathscr T_{p-1}(\mathscr M)
      \xrightarrow{\sim}
      \mathscr{H}\!om_{\mathscr O_Y}(\mathscr Q,\mathscr M).
      \tag{7.7.5.2}\label{III.7.7.5.2.fr}
    \]

  \item[c)] En désignant par $A_\alpha$ l'anneau de l'ouvert affine $U_\alpha$,
    pour tout indice $\alpha$ le foncteur de $A_\alpha$-modules
    $T_p^{A_\alpha}$ est exact à droite.

  \item[d)] Pour tout morphisme $g:Y'\to Y$, l'homomorphisme canonique
    \[
      \mathscr T_p(\mathscr O_Y)
      \otimes_{\mathscr O_Y}\mathscr O_{Y'}
      \longrightarrow
      \mathscr T_p^{Y'}(\mathscr O_{Y'})
      \tag{7.7.5.3}\label{III.7.7.5.3.fr}
    \]
    est un isomorphisme.
\end{enumerate}
\end{theorem}

\begin{proof}
La propriété de semi-continuité est locale sur $Y$ et résulte donc de la
remarque (7.7.4) et de (7.6.9). Il est clair que a$''$) entraîne a$'$) et que
a$'$) entraîne a). L'équivalence de a), a$''$), b) et b$'$) a été démontrée
dans (7.3.1) et (7.4.6), compte tenu de la remarque (7.7.4), lorsque $Y$ est
affine. Pour passer au cas général, prouvons d'abord que a) est équivalent à c),
ce qui prouvera le caractère local sur $Y$ de la propriété a) ; la démonstration
s'appliquera également pour prouver le caractère local de a$''$) et b). Comme il
est clair que c) entraîne a), tout revient à prouver la réciproque. Il suffit
évidemment de montrer que pour tout ouvert affine $U$ de $Y$ et toute suite
exacte
\[
  0\longrightarrow\mathscr F'
  \longrightarrow\mathscr F
  \longrightarrow\mathscr F''
  \longrightarrow0
\]
de $(\mathscr O_Y|U)$-Modules quasi-cohérents, il existe une suite exacte
\[
  0\longrightarrow\mathscr G'
  \longrightarrow\mathscr G
  \longrightarrow\mathscr G''
  \longrightarrow0
\]
de $\mathscr O_Y$-Modules quasi-cohérents telle que
\[
  \mathscr F'=\mathscr G'|U,
  \qquad
  \mathscr F=\mathscr G|U,
  \qquad
  \mathscr F''=\mathscr G''|U;
\]
or, cela résulte aussitôt de l'hypothèse que $Y$ est localement noethérien, et
de (I, 9.4.2) : on prolonge en effet $\mathscr F$ en un $\mathscr O_Y$-Module
quasi-cohérent $\mathscr G$, $\mathscr F'$ en un sous-$\mathscr O_Y$-Module
$\mathscr G'$ de $\mathscr G$, et il suffit de prendre
$\mathscr G''=\mathscr G/\mathscr G'$.

Pour démontrer l'équivalence de b) et b$'$) dans le cas général, remarquons que
lorsque $Y$ est affine, on sait que $\mathscr Q$ est déterminé à un isomorphisme
unique près ; si alors $U$ est un ouvert affine du schéma affine $Y$, on en
déduit qu'il existe un isomorphisme fonctoriel
\[
  \mathscr T_{p-1}^{U}(\mathscr M|U)
  \xrightarrow{\sim}
  \mathscr{H}\!om_{\mathscr O_Y|U}
    (\mathscr Q|U,\mathscr M|U).
\]
Dans le cas général, pour tout ouvert affine $U$ de $Y$, il y a un
$(\mathscr O_Y|U)$-Module cohérent $\mathscr Q_U$ et un isomorphisme fonctoriel
\[
  \mathscr T_{p-1}^{U}(\mathscr M|U)
  \xrightarrow{\sim}
  \mathscr{H}\!om_{\mathscr O_Y|U}
    (\mathscr Q_U,\mathscr M|U);
\]
la remarque précédente montre que si $V$ est un ouvert affine contenu dans $U$,
on a $\mathscr Q_U|V=\mathscr Q_V$ ; d'où l'existence et l'unicité du
$\mathscr O_Y$-Module $\mathscr Q$ vérifiant (7.7.5.2).

Reste enfin à montrer l'équivalence de a) et d) ; il est clair que d) est de
caractère local sur $Y$, et l'on a vu ci-dessus qu'il en est de même de a) ; en
outre, d) est aussi local sur $Y'$. Or, lorsque $Y=\operatorname{Spec}(A)$,
$Y'=\operatorname{Spec}(A')$, on a vu que $T_\bullet^{A'}$ est le foncteur obtenu

\oldpage[III]{201}
à partir de $T_\bullet^A$ par extension des scalaires à $A'$, et il est clair
alors que a$'$) entraîne que (7.7.5.3) est un isomorphisme. Inversement, supposons
toujours $Y=\operatorname{Spec}(A)$ affine et soit $A'$ la $A$-algèbre
$A\oplus M$, où $M$ est un $A$-module quelconque, la multiplication dans $A'$
étant donnée par
\[
  (a_1,m_1)(a_2,m_2)
  =
  (a_1a_2,a_1m_2+a_2m_1);
\]
alors
\[
  T_p^{A'}(A')
  =
  T_p(A\oplus M)
  =
  T_p(A)\oplus T_p(M),
\]
et l'hypothèse que (7.7.5.3) soit bijectif entraîne qu'il en est de même de
l'application canonique
\[
  T_p(A)\otimes_A M\longrightarrow T_p(M),
\]
autrement dit d) entraîne a$''$), ce qui achève la démonstration.
\end{proof}

\begin{theorem}[7.7.6]
Soient $Y$ un préschéma localement noethérien, $f:X\to Y$ un morphisme propre,
$\mathscr F$ un $\mathscr O_X$-Module cohérent et $Y$-plat. Il existe alors un
$\mathscr O_Y$-Module cohérent $\mathscr Q$ (déterminé à isomorphisme unique près)
et un isomorphisme de foncteurs en le $\mathscr O_Y$-Module quasi-cohérent
$\mathscr M$ :
\[
  f_*(\mathscr F\otimes_{\mathscr O_Y}\mathscr M)
  \xrightarrow{\sim}
  \mathscr{H}\!om_{\mathscr O_Y}(\mathscr Q,\mathscr M)
  \tag{7.7.6.1}\label{III.7.7.6.1.fr}
\]
(d'où un isomorphisme de foncteurs
\[
  \Gamma(X,\mathscr F\otimes_{\mathscr O_Y}\mathscr M)
  \xrightarrow{\sim}
  \operatorname{Hom}_{\mathscr O_Y}(\mathscr Q,\mathscr M)).
  \tag{7.7.6.2}\label{III.7.7.6.2.fr}
\]
\end{theorem}

\begin{proof}
En effet, comme $\mathscr M\mapsto\mathscr F\otimes_{\mathscr O_Y}\mathscr M$ est
exact (0\textsubscript{I}, 6.7.4) et $f_*$ exact à gauche, le foncteur
$\mathscr M\mapsto f_*(\mathscr F\otimes_{\mathscr O_Y}\mathscr M)$ est exact à
gauche. Il suffit alors d'appliquer l'équivalence de (7.7.5, b)) et (7.7.5,
b$'$)) pour $p=1$.
\end{proof}

\begin{corollary}[7.7.7]
Soient $Y$ un préschéma localement noethérien, $f:X\to Y$ un morphisme propre,
$\mathscr F$, $\mathscr F'$ deux $\mathscr O_X$-Modules cohérents et $Y$-plats,
$u:\mathscr F\to\mathscr F'$ un homomorphisme. Considérons les deux foncteurs en
le $\mathscr O_Y$-Module quasi-cohérent $\mathscr M$ :
\begin{align*}
  \mathscr T(\mathscr M)
  &=
  \operatorname{Ker}
  \bigl(
    f_*(\mathscr F\otimes_{\mathscr O_Y}\mathscr M)
    \longrightarrow
    f_*(\mathscr F'\otimes_{\mathscr O_Y}\mathscr M)
  \bigr),\\
  T(\mathscr M)
  &=
  \Gamma(Y,\mathscr T(\mathscr M))\\
  &=
  \operatorname{Ker}
  \bigl(
    \Gamma(X,\mathscr F\otimes_{\mathscr O_Y}\mathscr M)
    \longrightarrow
    \Gamma(X,\mathscr F'\otimes_{\mathscr O_Y}\mathscr M)
  \bigr).
\end{align*}
Alors il existe un $\mathscr O_Y$-Module cohérent $\mathscr R$ (déterminé à
isomorphisme unique près) et des isomorphismes de foncteurs
\[
  \mathscr T(\mathscr M)
  \xrightarrow{\sim}
  \mathscr{H}\!om_{\mathscr O_Y}(\mathscr R,\mathscr M)
  \tag{7.7.7.1}\label{III.7.7.7.1.fr}
\]
\[
  T(\mathscr M)
  \xrightarrow{\sim}
  \operatorname{Hom}_{\mathscr O_Y}(\mathscr R,\mathscr M).
  \tag{7.7.7.2}\label{III.7.7.7.2.fr}
\]
\end{corollary}

\begin{proof}
On peut se borner à démontrer (7.7.7.2) ; cela prouvera en effet (7.7.7.1)
dans le cas où $Y$ est affine, et on passera de là au cas général en raisonnant
comme dans la démonstration de l'équivalence de (7.7.5, b) et b$'$)), grâce à
l'unicité à isomorphisme unique près d'un représentant d'un foncteur
représentable (0, 8.1.8). Il résulte de (7.7.6) qu'il existe deux
$\mathscr O_Y$-Modules cohérents $\mathscr Q$, $\mathscr Q'$ définissant des
isomorphismes fonctoriels
\[
  \Gamma(X,\mathscr F\otimes_{\mathscr O_Y}\mathscr M)
  \xrightarrow{\sim}
  \operatorname{Hom}_{\mathscr O_Y}(\mathscr Q,\mathscr M),
  \qquad
  \Gamma(X,\mathscr F'\otimes_{\mathscr O_Y}\mathscr M)
  \xrightarrow{\sim}
  \operatorname{Hom}_{\mathscr O_Y}(\mathscr Q',\mathscr M).
\]
Or, $u:\mathscr F\to\mathscr F'$ définit canoniquement un morphisme de foncteurs
\[
  \Gamma(X,\mathscr F\otimes_{\mathscr O_Y}\mathscr M)
  \longrightarrow
  \Gamma(X,\mathscr F'\otimes_{\mathscr O_Y}\mathscr M);
\]

\oldpage[III]{202}
il correspond à ce dernier un homomorphisme unique
$v:\mathscr Q'\to\mathscr Q$ de $\mathscr O_Y$-Modules tel que le diagramme
\[
\xymatrix@C=42pt@R=28pt{
  \Gamma(X,\mathscr F\otimes_{\mathscr O_Y}\mathscr M)
    \ar[r]\ar[d]_-{\sim}
  &
  \Gamma(X,\mathscr F'\otimes_{\mathscr O_Y}\mathscr M)
    \ar[d]^-{\sim}
  \\
  \operatorname{Hom}_{\mathscr O_Y}(\mathscr Q,\mathscr M)
    \ar[r]
  &
  \operatorname{Hom}_{\mathscr O_Y}(\mathscr Q',\mathscr M)
}
\]
soit commutatif (0, 8.1.4). Comme le foncteur contravariant
$\mathscr N\mapsto\operatorname{Hom}_{\mathscr O_Y}(\mathscr N,\mathscr M)$ est
exact à gauche dans la catégorie des $\mathscr O_Y$-Modules, il suffit de prendre
$\mathscr R=\operatorname{Coker}(v)$ pour obtenir l'isomorphisme (7.7.7.2)
cherché.
\end{proof}

\begin{corollary}[7.7.8]
Sous les hypothèses de (7.7.6) relatives à $X$, $Y$ et $f$, soient $\mathscr F$,
$\mathscr G$ deux $\mathscr O_X$-Modules cohérents vérifiant les conditions
suivantes : (i) $\mathscr F$ est $Y$-plat ; (ii) $\mathscr G$ est isomorphe au
conoyau d'un homomorphisme de $\mathscr O_X$-Modules localement libres de type
fini $\mathscr E_1\to\mathscr E_0$. Considérons les deux foncteurs en le
$\mathscr O_Y$-Module quasi-cohérent $\mathscr M$ :
\begin{align*}
  \mathscr T(\mathscr M)
  &=
  f_*
  \bigl(
    \mathscr{H}\!om_{\mathscr O_X}
      (\mathscr G,\mathscr F\otimes_{\mathscr O_Y}\mathscr M)
  \bigr),\\
  T(\mathscr M)
  &=
  \Gamma(Y,\mathscr T(\mathscr M))
  =
  \operatorname{Hom}_{\mathscr O_X}
    (\mathscr G,\mathscr F\otimes_{\mathscr O_Y}\mathscr M).
\end{align*}
Alors il existe un $\mathscr O_Y$-Module cohérent $\mathscr N$ (déterminé à
isomorphisme unique près) et des isomorphismes de foncteurs
\[
  \mathscr T(\mathscr M)
  \xrightarrow{\sim}
  \mathscr{H}\!om_{\mathscr O_Y}(\mathscr N,\mathscr M)
  \tag{7.7.8.1}\label{III.7.7.8.1.fr}
\]
\[
  T(\mathscr M)
  \xrightarrow{\sim}
  \operatorname{Hom}_{\mathscr O_Y}(\mathscr N,\mathscr M).
  \tag{7.7.8.2}\label{III.7.7.8.2.fr}
\]
\end{corollary}

\begin{proof}
En vertu de l'isomorphisme fonctoriel ($0_{\mathrm I}$, 5.4.2.1), on a des
isomorphismes fonctoriels en $\mathscr M$
\begin{align*}
  \mathscr{H}\!om_{\mathscr O_X}
    (\mathscr E_i,\mathscr F\otimes_{\mathscr O_Y}\mathscr M)
  &\xrightarrow{\sim}
  \check{\mathscr E}_i
    \otimes_{\mathscr O_X}
    (\mathscr F\otimes_{\mathscr O_Y}\mathscr M)\\
  &\xrightarrow{\sim}
  (\check{\mathscr E}_i\otimes_{\mathscr O_X}\mathscr F)
    \otimes_{\mathscr O_Y}\mathscr M\\
  &\xrightarrow{\sim}
  \mathscr{H}\!om_{\mathscr O_X}(\mathscr E_i,\mathscr F)
    \otimes_{\mathscr O_Y}\mathscr M
\end{align*}
pour $i=0,1$. Posons
$\mathscr F_i=\mathscr{H}\!om_{\mathscr O_X}(\mathscr E_i,\mathscr F)$ pour
$i=0,1$ ; ce sont des $\mathscr O_X$-Modules cohérents ($0_{\mathrm I}$, 5.3.5)
et $Y$-plats ($0_{\mathrm I}$, 5.4.2) ; soit
$u=\mathscr{H}\!om(v,1_{\mathscr F}):\mathscr F_0\to\mathscr F_1$. En vertu de
l'exactitude à gauche du foncteur
$\mathscr H\mapsto\mathscr{H}\!om_{\mathscr O_X}
(\mathscr H,\mathscr F\otimes_{\mathscr O_Y}\mathscr M)$, on a des isomorphismes
fonctoriels en $\mathscr M$
\begin{align*}
  \mathscr{H}\!om_{\mathscr O_X}
    (\mathscr G,\mathscr F\otimes_{\mathscr O_Y}\mathscr M)
  &\xrightarrow{\sim}
  \operatorname{Ker}
  \bigl(
    \mathscr{H}\!om_{\mathscr O_X}
      (\mathscr E_0,\mathscr F\otimes_{\mathscr O_Y}\mathscr M)
    \longrightarrow
    \mathscr{H}\!om_{\mathscr O_X}
      (\mathscr E_1,\mathscr F\otimes_{\mathscr O_Y}\mathscr M)
  \bigr)\\
  &\xrightarrow{\sim}
  \operatorname{Ker}
  \bigl(
    \mathscr F_0\otimes_{\mathscr O_Y}\mathscr M
    \longrightarrow
    \mathscr F_1\otimes_{\mathscr O_Y}\mathscr M
  \bigr).
\end{align*}
Puisque $f_*$ est exact à gauche, on en déduit un isomorphisme fonctoriel
\[
  f_*
  \bigl(
    \mathscr{H}\!om_{\mathscr O_X}
      (\mathscr G,\mathscr F\otimes_{\mathscr O_Y}\mathscr M)
  \bigr)
  \xrightarrow{\sim}
  \operatorname{Ker}
  \bigl(
    f_*(\mathscr F_0\otimes_{\mathscr O_Y}\mathscr M)
    \longrightarrow
    f_*(\mathscr F_1\otimes_{\mathscr O_Y}\mathscr M)
  \bigr)
\]
et il suffit alors d'appliquer (7.7.7).
\end{proof}

\oldpage[III]{203}
\begin{remarks}[7.7.9]
\textnormal{(i)} Dans (7.7.6), (7.7.7), (7.7.8), la formation des
$\mathscr O_Y$-Modules $\mathscr Q$, $\mathscr R$, $\mathscr N$ \emph{permute aux
changements de base}. Par exemple (gardant les notations de (7.7.2)), dans le
cas (7.7.6), on a, pour tout $\mathscr O_{Y'}$-Module quasi-cohérent
$\mathscr M'$, l'isomorphisme
\[
  f'_*(\mathscr F\otimes_{\mathscr O_Y}\mathscr M')
  \xrightarrow{\sim}
  \mathscr{H}\!om_{\mathscr O_{Y'}}(g^*(\mathscr Q),\mathscr M')
\]
car en vertu de la remarque faite dans (7.7.2), tout revient à voir que l'on a
\[
  \operatorname{Hom}_{\mathscr O_Y}(\mathscr Q,g_*(\mathscr M'))
  =
  \operatorname{Hom}_{\mathscr O_{Y'}}(g^*(\mathscr Q),\mathscr M')
\]
ce qui n'est autre que ($0_{\mathrm I}$, 4.4.3.1). De même, quand dans (7.7.7)
on remplace $Y$, $f$, $\mathscr M$, $\mathscr F$, $\mathscr F'$ par $Y'$, $f'$,
$\mathscr M'$, $\mathscr F\otimes_{\mathscr O_X}\mathscr O_{X'}$,
$\mathscr F'\otimes_{\mathscr O_X}\mathscr O_{X'}$, il faut remplacer
$\mathscr R$ par $g^*(\mathscr R)$. Enfin, dans (7.7.8), lorsqu'on remplace $X$,
$Y$, $f$, $\mathscr M$, $\mathscr F$ par $X'$, $Y'$, $f'$, $\mathscr M'$,
$\mathscr F\otimes_{\mathscr O_X}\mathscr O_{X'}$, et $\mathscr G$ par
$\mathscr G'=\mathscr G\otimes_{\mathscr O_Y}\mathscr O_{X'}$, il faut remplacer
$\mathscr N$ par $g^*(\mathscr N)$ : cela résulte de ce que l'on a encore une
suite exacte
\[
  \mathscr E'_1\longrightarrow
  \mathscr E'_0\longrightarrow
  \mathscr G'\longrightarrow0
  \qquad\text{avec}\qquad
  \mathscr E'_i=\mathscr E_i\otimes_{\mathscr O_Y}\mathscr O_{X'}
\]
et de ce que
\[
  \check{\mathscr E}'_i
  \otimes_{\mathscr O_{X'}}
  (\mathscr F\otimes_{\mathscr O_Y}\mathscr M')
  =
  \check{\mathscr E}_i
  \otimes_{\mathscr O_X}
  (\mathscr F\otimes_{\mathscr O_Y}\mathscr M')
  \qquad(i=0,1).
\]

\medskip
\textnormal{(ii)} La condition (ii) de l'énoncé de (7.7.8) relative à
$\mathscr G$ est toujours satisfaite pour un $\mathscr O_X$-Module cohérent
$\mathscr G$ quelconque lorsqu'il existe un $\mathscr O_X$-Module inversible
$Y$-ample, par exemple lorsque $Y$ est affine et $f:X\to Y$ est un morphisme
projectif. Il suffit alors de noter (compte tenu de (II, 5.5.1)) qu'il existe un
$\mathscr O_X$-Module localement libre de type fini $\mathscr E_0$ tel que
$\mathscr G$ soit isomorphe à un quotient de $\mathscr E_0$ (II, 2.7.10) ; comme
$\mathscr E_0$ et $\mathscr G$ sont cohérents, il en est de même du noyau
$\mathscr G_1$ de $\mathscr E_0\to\mathscr G$, et en appliquant le même résultat
à $\mathscr G_1$, on obtient bien une suite exacte
$\mathscr E_1\to\mathscr E_0\to\mathscr G\to0$ où $\mathscr E_0$ et
$\mathscr E_1$ sont localement libres de type fini.

\medskip
\textnormal{(iii)} Nous prouverons au chap. V que, dans (7.7.8), l'hypothèse
restrictive (ii) est \emph{surabondante}.
\end{remarks}

\begin{proposition}[7.7.10]
\textnormal{(Critères locaux pour la propriété d'échange).} Sous les conditions
générales de (7.7.5), soient $y$ un point de $Y$, $p$ un entier. Les propriétés
suivantes sont équivalentes :
\begin{enumerate}
  \item[a)] Le foncteur $T_p^{\mathscr O_y}$ est exact à droite.

  \item[b)] L'homomorphisme canonique
    \[
      T_p^{\mathscr O_y}(\mathscr O_y)
      \longrightarrow
      T_p^{\mathscr O_y}(\kappa(y))
    \]
    est surjectif.

  \item[c)] Pour tout entier $n$, l'homomorphisme canonique
    \[
      T_p^{\mathscr O_y}
      (\mathscr O_y/\mathfrak m_y^{n+1})
      \longrightarrow
      T_p^{\mathscr O_y}(\kappa(y))
    \]
    est surjectif.
\end{enumerate}
De plus, l'ensemble des $y\in Y$ vérifiant ces conditions est le plus grand
ensemble ouvert $U$ de $Y$ tel que $\mathscr T_p^U$ soit exact à droite.
\end{proposition}

\begin{proof}
Compte tenu de (7.7.4), l'équivalence de a), b) et c) résulte de (7.4.7) et
(7.5.2). Le fait que l'ensemble $U$ où $T_p^{\mathscr O_y}$ est exact à droite
soit ouvert est aussi conséquence de (7.4.4), et inversement si
$\mathscr T_p^V$ est exact à droite, il en est de même de
$T_p^{\mathscr O_y}$ pour tout $y\in V$, par la condition c) de (7.7.5) et
(7.6.1).
\end{proof}

\begin{corollary}[7.7.11]
Si $\mathscr T_p$ est exact à droite (resp. à gauche), alors, pour tout morphisme
$g:Y'\to Y$, $\mathscr T_p^{Y'}$ est exact à droite (resp. à gauche). La
réciproque est vraie lorsque le morphisme $g$ est fidèlement plat.
\end{corollary}

\oldpage[III]{204}
\begin{proof}
La première assertion est conséquence immédiate de (7.6.1) et du fait que la
question est locale sur $Y$ et $Y'$, d'après (7.7.5, c) et b)). Pour démontrer
la seconde assertion, il suffit de voir que pour tout $y\in Y$,
$T_p^{\mathscr O_y}$ est exact à droite (resp. à gauche), en vertu de (7.7.10).
Mais par hypothèse, il existe $y'\in Y'$ tel que $g(y')=y$, et
$\mathscr O_{y'}$ est un $\mathscr O_y$-module fidèlement plat ; la conclusion
résulte alors de l'hypothèse et de (7.6.1).
\end{proof}

\begin{remarks}[7.7.12]
\textnormal{(i)} Sous les hypothèses de (7.7.4), supposons en outre que
$\mathscr P_\bullet$ soit un complexe \emph{fini} ; alors il résulte de (7.7.1)
(puisqu'on peut prendre pour $\mathfrak U$ un recouvrement ouvert affine
\emph{fini} de $X$) que le bicomplexe
$C^\bullet(\mathfrak U,\mathscr P^\bullet\otimes_{\mathscr O_Y}\mathscr M)$ est
aussi fini, et de façon précise qu'il existe un ensemble fini $E$ de couples,
\emph{indépendant de $\mathscr M$}, tel que
\[
  C^h(\mathfrak U,\mathscr P^k\otimes_{\mathscr O_Y}\mathscr M)=0
  \qquad\text{pour tous les couples }(h,k)\notin E.
\]
On en conclut qu'il existe $i_1$ tel que, pour $i\geq i_1$, on ait
$\mathscr T_i(\mathscr M)=0$ pour \emph{tout} $\mathscr O_Y$-Module
quasi-cohérent $\mathscr M$. En particulier, $\mathscr T_i$ est trivialement un
foncteur exact en $\mathscr M$ pour ces valeurs de $i$, et par suite (7.4.1),
$Z'_i(L_\bullet)$ est un $A$-module \emph{plat} de type fini (donc
\emph{projectif} de type fini, puisque $A$ est noethérien) pour ces valeurs de
$i$. Considérons alors le complexe $(L'_\bullet)$ de $A$-modules tel que
$L'_i=L_i$ pour $i<i_1$, $L'_{i_1}=Z'_{i_1}(L_\bullet)$ et $L'_i=0$ pour
$i>i_1$ et soit $\mathscr L'_\bullet=\widetilde{L'_\bullet}$. Il est clair que
\[
  \mathscr H_i(\mathscr L'_\bullet\otimes_{\mathscr O_Y}\mathscr M)
  =
  \mathscr H_i(\mathscr L_\bullet\otimes_{\mathscr O_Y}\mathscr M)
\]
pour $i<i_1-1$ et aussi pour $i\geq i_1$ (les deux membres étant alors nuls) ;
enfin, comme
$\operatorname{Im}(Z'_{i_1}\otimes_A M)=\operatorname{Im}(L_{i_1}\otimes_A M)$
par définition, on a aussi
$\mathscr H_i(\mathscr L'_\bullet\otimes_{\mathscr O_Y}\mathscr M)
=\mathscr H_i(\mathscr L_\bullet\otimes_{\mathscr O_Y}\mathscr M)$ pour
$i=i_1-1$. On voit donc qu'on peut supposer dans (7.7.4) que
$\mathscr L_\bullet$ est aussi un complexe \emph{fini}, à condition d'exiger
seulement que les $\mathscr L_i$ soient des $\mathscr O_Y$-Modules
\emph{localement libres} (associés à des $A$-modules projectifs de type fini).

Ce raisonnement s'applique en particulier au cas où $\mathscr P_\bullet$ est
réduit à \emph{un seul} terme $\mathscr F\ne0$, de degré 0 (auquel cas
$\mathscr T_n(\mathscr M)=\mathrm R^{-n}f_*
(\mathscr F\otimes_{\mathscr O_Y}\mathscr M)$) ; on peut alors supposer que les
$\mathscr L_i$ sont nuls pour $i>0$ ; on utilisera de préférence dans ce cas les
notations cohomologiques, écrivant donc $\mathscr T^{-p}$ au lieu de
$\mathscr T_p$.

\medskip
\textnormal{(ii)} Lorsque dans l'énoncé de (7.7.5), on ne suppose plus que les
$\mathscr P_i$ sont $Y$-plats, les conclusions restent valables à condition de
poser cette fois
\[
  \mathscr T_p(\mathscr M)
  =
  \operatorname{Tor}^Y_p
    (f,1_Y;\mathscr P_\bullet,\mathscr M).
  \tag{7.7.12.1}\label{III.7.7.12.1.fr}
\]
En effet, $\operatorname{Tor}^Y_n(f,1_Y;\mathscr P_\bullet,\mathscr O_Y)$ est
alors un $\mathscr O_Y$-Module \emph{cohérent} en vertu de (6.7.9). La
démonstration de (6.10.5) s'applique sans changement, compte tenu de (6.10.1)
et montre encore que lorsque $Y=\operatorname{Spec}(A)$ est affine, l'on a
\[
  \mathscr T_p(\mathscr M)
  =
  \mathscr H_p(\mathscr L_\bullet\otimes_{\mathscr O_Y}\mathscr M),
  \qquad
  \mathscr L_\bullet=\widetilde{L_\bullet},
\]
où $L_\bullet$ est un complexe de $A$-modules libres de type fini ; cela prouve
notre assertion.
\end{remarks}

\subsection{Application aux morphismes propres : II. Critères de platitude
cohomologique}
\label{subsection:III.7.8.fr}

\begin{definition}[7.8.1]
Soient $X$, $Y$ deux préschémas, $f:X\to Y$ un morphisme quasi-compact et
séparé, $\mathscr P_\bullet$ un complexe de $\mathscr O_X$-Modules
quasi-cohérents et $Y$-plats, $\mathscr T_\bullet$ le foncteur homologique de
$\mathscr O_Y$-Modules quasi-cohérents défini par (7.7.1.1), $y$ un point de $Y$.
On dit que $\mathscr P_\bullet$ est

\oldpage[III]{205}
\emph{homologiquement plat sur $Y$ au point $y$, en dimension $p$} (ou
\emph{cohomologiquement plat sur $Y$ au point $y$, en dimension $-p$}) s'il
existe un voisinage ouvert $U$ de $y$ dans $Y$ tel que
$\mathscr T_p^U=\mathscr T_U^{-p}$ soit exact. On dit que $\mathscr P_\bullet$ est
\emph{homologiquement plat en dimension $p$ sur $Y$} (ou
\emph{cohomologiquement plat en dimension $-p$ sur $Y$}) s'il est
homologiquement plat sur $Y$ en tout point $y\in Y$, en dimension $p$.

Lorsque $\mathscr P_\bullet$ est homologiquement plat sur $Y$ (resp. sur $Y$ au
point $y$) pour \emph{toute dimension $p$}, on dit simplement que
$\mathscr P_\bullet$ est \emph{homologiquement plat sur $Y$} (resp. sur $Y$ au
point $y$) ou \emph{cohomologiquement plat sur $Y$} (resp. sur $Y$ au point
$y$).
\end{definition}

\begin{env}[7.8.2]
Par définition, la notion de platitude homologique sur $Y$ est \emph{locale sur
$Y$}. Si $Y$ est localement noethérien, ou un \emph{schéma}, pour que
$\mathscr P_\bullet$ soit homologiquement plat sur $Y$ en dimension $p$, il faut
et il suffit que le foncteur $\mathscr T_p$ soit \emph{exact} : la démonstration
a été faite dans le cas où $Y$ est localement noethérien au cours de la
démonstration de (7.7.5) ; le raisonnement est le même (basé sur (I, 9.4.2)
appliqué à un ouvert affine dans un schéma quasi-compact) lorsque $Y$ est un
schéma.
\end{env}

\begin{proposition}[7.8.3]
Les notations et hypothèses étant celles de (7.8.1), les conditions suivantes
sont équivalentes :
\begin{enumerate}
  \item[a)] $\mathscr P_\bullet$ est homologiquement plat sur $Y$ au point $y$
    en dimension $p$.
  \item[b)] Il existe un voisinage ouvert $U$ de $y$ dans $Y$ tel que
    $\mathscr T_p^U$ et $\mathscr T_{p+1}^U$ soient exacts à droite.
  \item[c)] Il existe un voisinage ouvert $U$ de $y$ dans $Y$ tel que
    $\mathscr T_p^U$ et $\mathscr T_{p-1}^U$ soient exacts à gauche.
  \item[d)] Il existe un voisinage ouvert $U$ de $y$ dans $Y$ tel que
    $\mathscr T_{p+1}^U$ soit exact à droite et $\mathscr T_{p-1}^U$ exact à
    gauche.
\end{enumerate}
\end{proposition}

\begin{proof}
Compte tenu de l'interprétation de $\mathscr T_p$ lorsque $Y$ est affine, cela
n'est qu'une traduction d'une partie de (7.3.3).
\end{proof}

\begin{proposition}[7.8.4]
Soient $Y$ un préschéma localement noethérien, $f:X\to Y$ un morphisme propre,
$\mathscr P_\bullet$ un complexe de $\mathscr O_X$-Modules cohérents et $Y$-plats
limité inférieurement, $\mathscr T_\bullet$ le foncteur défini par (7.7.1.1).
Pour tout $y\in Y$, les conditions suivantes sont équivalentes :
\begin{enumerate}
  \item[a)] $\mathscr P_\bullet$ est homologiquement plat sur $Y$ en $y$ en
    dimension $p$.
  \item[b)] Le foncteur $T_p^{\mathscr O_y}$ est exact.
  \item[c)] Il existe un entier $n_0$ tel que pour $n\geq n_0$, on ait
    \[
      \operatorname{long}T_p^{\mathscr O_y}
        (\mathscr O_y/\mathfrak m_y^{n+1})
      =
      \operatorname{long}T_p^{\mathscr O_y}(\kappa(y))
      \cdot
      \operatorname{long}(\mathscr O_y/\mathfrak m_y^{n+1})
      \tag{7.8.4.1}\label{III.7.8.4.1.fr}
    \]
    (où il s'agit de longueurs de $\mathscr O_y$-modules).
  \item[d)] Il y a un voisinage ouvert $U$ de $y$ tel que
    $(\mathscr H^{-p}(f,\mathscr P^\bullet))|U$ soit isomorphe à un
    $(\mathscr O_Y|U)$-Module de la forme $(\mathscr O_Y|U)^m$ et que, pour tout
    $(\mathscr O_Y|U)$-Module quasi-cohérent $\mathscr M$, l'homomorphisme
    canonique
    \[
      ((\mathscr H^{-p}(f,\mathscr P^\bullet))|U)
      \otimes_{\mathscr O_Y|U}\mathscr M
      \longrightarrow
      \mathscr H^{-p}
      \bigl(
        f,(\mathscr P^\bullet|U)
          \otimes_{\mathscr O_Y|U}\mathscr M
      \bigr)
      \tag{7.8.4.2}\label{III.7.8.4.2.fr}
    \]
    soit bijectif.
\end{enumerate}
Lorsque ces conditions sont vérifiées, on a en outre la propriété suivante :
\begin{enumerate}
  \item[e)] Il existe un voisinage de $y$ dans lequel la fonction
    $z\mapsto d_p(z)$ (définie dans (7.7.5.1)) est constante.
\end{enumerate}
De plus, si $Y$ est réduit au point $y$ ($0_{\mathrm I}$, 4.1.4), e) est
équivalente aux autres conditions.
\end{proposition}

\oldpage[III]{206}
\begin{proof}
En effet, la condition b) équivaut à dire que $T_p^{\mathscr O_y}$ et
$T_{p+1}^{\mathscr O_y}$ sont exacts à droite (7.3.3). L'équivalence de a) et b)
résulte alors de (7.7.10) et (7.8.3). Comme
$\mathscr O_y/\mathfrak m_y^{n+1}$ est artinien, et que
$T_p^{\mathscr O_y}(\mathscr O_y/\mathfrak m_y^{n+1})$ et
$T_p^{\mathscr O_y}(\kappa(y))$ sont des
$(\mathscr O_y/\mathfrak m_y^{n+1})$-modules de type fini (7.7.4), donc de
longueur finie, l'équivalence de b) et c) résulte encore de (7.7.10) et de
(7.3.5.7). Le fait que a) entraîne e), et lui est équivalente lorsque
$\mathscr O_y$ est réduit, est conséquence de (7.6.9). Enfin, a) entraîne que
(7.8.4.2) est bijectif en vertu de la définition (7.8.1) et de (7.7.5) ; d'autre
part, a) entraîne que $(\mathscr H^{-p}(f,\mathscr P^\bullet))_y$ est un
$\mathscr O_y$-module plat (7.3.3, f)), donc libre ($0_{\mathrm I}$, 10.1.3),
puisqu'il s'agit d'un $\mathscr O_y$-module de type fini en vertu de (7.7.4) ;
puisque $\mathscr H^{-p}(f,\mathscr P^\bullet)$ est un $\mathscr O_Y$-Module
cohérent (7.7.4), il est localement libre dans un voisinage de $y$
($0_{\mathrm I}$, 5.2.7). Inversement, il est clair que d) entraîne a) par
définition du foncteur $\mathscr T_p^U$ (7.7.2.2).
\end{proof}

\begin{proposition}[7.8.5]
Sous les hypothèses de (7.8.4), les conditions suivantes sont équivalentes :
\begin{enumerate}
  \item[a)] $\mathscr P_\bullet$ est homologiquement plat sur $Y$ en toute
    dimension $i\leq p$.
  \item[b)] Pour $i\leq p+1$ les foncteurs $\mathscr T_i$ sont exacts à droite.
  \item[c)] Pour $i\geq-p$, les $\mathscr O_Y$-Modules
    $\mathscr H^i(f,\mathscr P^\bullet)$ sont localement libres.
\end{enumerate}
\end{proposition}

\begin{proof}
L'équivalence de a) et b) est triviale (7.8.3) et a) entraîne c) en vertu de
(7.8.4). Inversement, supposons c) vérifiée ; notons d'autre part qu'on a
$\mathscr L_i=0$ pour $i<i_0$ (7.7.4), donc aussi $\mathscr T_i=0$ pour
$i\leq i_0$. Tout point $y\in Y$ a donc un voisinage affine
$U=\operatorname{Spec}(A)$ tel que $T_i^A(A)$ soit un $A$-module libre pour
$i\leq p$ ; en vertu de (7.3.7), on en conclut que
$\mathscr T_i^U=T_i^A$ est exact pour $i\leq p$.
\end{proof}

Nous allons surtout appliquer les critères de platitude cohomologique au cas où
le complexe $\mathscr P_\bullet$ est réduit à un \emph{seul}
$\mathscr O_X$-Module cohérent $\mathscr F$ \emph{plat} sur $Y$, pris égal à
$\mathscr P_0$ ; rappelons que l'on a alors
$\mathscr T_p(\mathscr M)=\mathrm R^{-p}f_*
(\mathscr F\otimes_{\mathscr O_Y}\mathscr M)$.

\begin{proposition}[7.8.6]
Soient $Y$ un préschéma localement noethérien, $f:X\to Y$ un morphisme propre
et plat, $y$ un point de $Y$ ; désignons par $X_y$ la fibre
$f^{-1}(y)=X\otimes_Y\kappa(y)$. Supposons que
$\Gamma(X_y,\mathscr O_{X_y})=R$ soit une $\kappa(y)$-algèbre séparable
(Bourbaki, \emph{Alg.}, chap. VIII, § 7, n\textsuperscript{o} 5) autrement dit
composée d'un nombre fini d'extensions séparables de degré fini de $\kappa(y)$.
Alors $\mathscr O_X$ est cohomologiquement plat sur $Y$ au point $y$ en
dimension 0.
\end{proposition}

\begin{proof}
En vertu de (7.8.4), on peut se borner au cas où $Y$ est le spectre de l'anneau
local $A=\mathscr O_y$ ; l'hypothèse que $f$ est plat entraîne
$\mathscr T_{-1}=0$, donc on voit déjà que $T_0^A$ est exact à gauche et tout
revient à voir qu'il est exact à droite ; en vertu de (7.7.10, c)), on est même
ramené au cas où $A=\mathscr O_y$ est \emph{artinien}. Soit $k'$ une extension
finie de $\kappa(y)$ qui soit un corps neutralisant de $R$, de sorte que
$R\otimes_{\kappa(y)}k'$ est composée directe d'un nombre fini de corps
isomorphes à $k'$. On sait qu'il existe un homomorphisme local de $A$ dans un
anneau local $A'$, faisant de $A'$ une $A$-algèbre \emph{libre} finie sur $A$,
et tel que le corps résiduel de $A'$ soit isomorphe à $k'$ ($0_{\mathrm I}$,
10.3.2). En vertu de (7.6.1), on est ramené à prouver que $T_0^{A'}$ est exact
à droite, autrement dit on peut supposer que $R$ est composé direct de $m$
corps isomorphes à $\kappa(y)$. Notons maintenant le lemme élémentaire
suivant :

\begin{lemma}[7.8.6.1]
Soit $Z$ un espace annelé en anneaux locaux ; pour que $Z$ soit connexe,

\oldpage[III]{207}
il faut et il suffit que l'anneau $\Gamma(Z,\mathscr O_Z)$ ne soit pas un produit
de deux anneaux non réduits à 0.
\end{lemma}

\begin{proof}
Il est clair en effet que si $Z$ est réunion de deux ouverts non vides disjoints,
$\Gamma(Z,\mathscr O_Z)$ est isomorphe au produit des deux anneaux
$\Gamma(Z_1,\mathscr O_Z)$ et $\Gamma(Z_2,\mathscr O_Z)$ non réduits à 0.
Inversement, dire que $\Gamma(Z,\mathscr O_Z)$ est un tel produit équivaut à
dire qu'il y a dans $\Gamma(Z,\mathscr O_Z)$ un idempotent $s$ distinct de 0 et
de 1 ; pour tout $z\in Z$, $s_z$ est alors un idempotent dans $\mathscr O_z$,
donc égal à 0 ou 1. Mais il est clair que l'ensemble des $z$ tels que $s_z=0$
est ouvert ; d'autre part, si $s_z=1$, on a par définition $s(z)\ne0$, donc
l'ensemble des $z$ où $s_z=1$ est aussi ouvert ($0_{\mathrm I}$, 5.5.2) ; d'où
la conclusion.
\end{proof}

Il résulte de ce lemme que $X_y$ a exactement $m$ composantes connexes $X'_i$
et que $\Gamma(X'_i,\mathscr O_{X'_i})=\kappa(y)$ pour tout $i$. Comme $A$ a été
supposé local et artinien, son spectre est réduit à un point, donc $X$ et $X_y$
ont même espace sous-jacent ; $X$ a donc $m$ composantes connexes $X_i$ telles
que $X'_i=X_i\otimes_Y\kappa(y)$. On est ainsi finalement ramené au cas où
$R=\kappa(y)$ ; en vertu de (7.7.10, b)), on est ramené à prouver que
l'homomorphisme canonique
\[
  \Gamma(X,\mathscr O_X)
  \longrightarrow
  \Gamma(X_y,\mathscr O_{X_y})
\]
est surjectif ; mais cela est trivial, car le composé
\[
  \Gamma(Y,\mathscr O_Y)=A
  \longrightarrow
  \Gamma(X,\mathscr O_X)
  \longrightarrow
  \Gamma(X_y,\mathscr O_{X_y})=\kappa(y)
\]
est déjà surjectif.
\end{proof}

\begin{corollary}[7.8.7]
Sous les hypothèses de (7.8.6), il existe un voisinage ouvert $U$ de $y$ tel
que :
\begin{enumerate}
  \item[(i)] $f_*(\mathscr O_X)|U$ soit isomorphe à un
    $(\mathscr O_Y|U)$-Module de la forme $(\mathscr O_Y|U)^m$.
  \item[(ii)] Pour tout $z\in U$, l'homomorphisme canonique
    \[
      (f_*(\mathscr O_X))_z
      \otimes_{\mathscr O_z}\kappa(z)
      \longrightarrow
      \Gamma(X_z,\mathscr O_{X_z})
    \]
    est bijectif.
\end{enumerate}
\end{corollary}

\begin{proof}
(i) résulte de (7.8.6) et (7.8.4).

(ii) résulte de ce que $\mathscr T_0^U$ est exact (pour $U$ convenablement
choisi), et de (7.7.5.3).
\end{proof}

\begin{corollary}[7.8.8]
Supposons vérifiées les conditions de (7.8.6) et en outre que
$\Gamma(X_y,\mathscr O_{X_y})=\kappa(y)$. Alors il existe un voisinage ouvert
$U$ de $y$ tel que l'homomorphisme canonique
$\mathscr O_Y|U\to f_*(\mathscr O_X)|U$ soit bijectif.
\end{corollary}

\begin{proof}
En effet, il résulte de (7.8.7, (ii)) que l'entier $m$ figurant dans (7.8.7,
(i)) est nécessairement égal à 1.
\end{proof}

\begin{corollary}[7.8.9]
Sous les hypothèses de (7.8.6), il existe un voisinage ouvert $U$ de $y$, un
$\mathscr O_U$-Module cohérent $\mathscr Q$ (déterminé à isomorphisme unique
près) et un isomorphisme de foncteurs en le $\mathscr O_U$-Module quasi-cohérent
$\mathscr M$ :
\[
  \mathrm R^1f_*(f^*(\mathscr M))
  \xrightarrow{\sim}
  \mathscr{H}\!om_{\mathscr O_U}(\mathscr Q,\mathscr M).
  \tag{7.8.9.1}\label{III.7.8.9.1.fr}
\]
\end{corollary}

\begin{proof}
En effet, l'hypothèse entraîne que $\mathscr T_0^U$ est exact pour un $U$
convenable ; il suffit donc d'appliquer l'équivalence de (7.7.5, a)) et
(7.7.5, b$'$)) au cas $p=0$ et en prenant pour $\mathscr P_\bullet$ le complexe
réduit à son terme de degré 0 égal à $\mathscr O_X$.
\end{proof}

\begin{remarks}[7.8.10]
\textnormal{(i)} Sous les conditions de (7.8.6), considérons la factorisation de
Stein de $f$ (4.3.3)
\[
\xymatrix{
  X \ar[r]^{f'} & Y' \ar[r]^{g} & Y .
}
\]

\oldpage[III]{208}
avec $Y'=\operatorname{Spec}(f_*(\mathscr O_X))$ ; le morphisme fini $g$ est alors
tel que $g_*(\mathscr O_{Y'})=f_*(\mathscr O_X)$ soit localement libre au
voisinage de $y$, et sa fibre en $y$ est le spectre d'une algèbre séparable sur
$\kappa(y)$ (II, 1.5.1). Nous en déduirons au chap. IV qu'il y a un voisinage
ouvert $U$ de $y$ dans $Y$ tel que pour la restriction $g^{-1}(U)\to U$ de $g$,
\emph{toute} fibre $g^{-1}(z)$ (où $z\in U$) soit spectre d'une algèbre séparable
sur $\kappa(z)$ (c'est ce que nous appellerons un \emph{revêtement étale} de
$U$) ; il résultera alors de (7.8.7, (ii)) que l'hypothèse faite sur le point
$y$ dans (7.8.6) est vérifiée aussi en tous les points d'un voisinage de $y$.

\medskip
\textnormal{(ii)} Nous verrons au chap. V que, même si $X$ est projectif sur $Y$
(et même s'il est en outre « simple » sur $Y$, propriété qui sera définie au
chap. IV), le $\mathscr O_U$-Module $\mathscr Q$ de (7.8.9) n'est pas
nécessairement localement libre ; en d'autres termes, $\mathscr O_X$ (sous ces
conditions) n'est pas nécessairement cohomologiquement plat en \emph{dimension
1} sur $Y$ au point $y$. Au chap. V, nous interpréterons $\mathscr Q$ comme le
faisceau des 1-différentielles du schéma de Picard de $X$ par rapport à $Y$ le
long de la section unité.
\end{remarks}

\subsection{Application aux morphismes propres : III. Invariance de la
caractéristique d'Euler-Poincaré et du polynôme de Hilbert}
\label{subsection:III.7.9.fr}
\label{III.7.9.fr}

\begin{env}[7.9.1]
\label{III.7.9.1.fr}
Soient $A$ un anneau, $M$ un $A$-module projectif de type fini ; rappelons
(Bourbaki, \emph{Alg. comm.}, chap. II, \S~5, n\textsuperscript{o}~2) qu'il
revient au même de dire que le $\mathscr O_X$-Module associé $\widetilde M$ sur
$X=\operatorname{Spec}(A)$ est \emph{localement libre de type fini}. Pour tout
$\mathfrak p\in\operatorname{Spec}(A)$ on appelle \emph{rang de $M$ en
$\mathfrak p$} et on note $\operatorname{rang}_{\mathfrak p}(M)$ le rang du
$A_{\mathfrak p}$-module libre $M_{\mathfrak p}$ (ou encore le rang en
$\mathfrak p$ du $\mathscr O_X$-Module localement libre $\widetilde M$). On a
donc
\[
  \operatorname{rang}_{\mathfrak p}M
  =\operatorname{rang}_{\mathfrak p}(M_{\mathfrak p})
  =\operatorname{rang}_{\kappa(\mathfrak p)}
  \bigl(M\otimes_A\kappa(\mathfrak p)\bigr).
  \tag{7.9.1.1}\label{III.7.9.1.1.fr}
\]
\end{env}

\begin{proposition}[7.9.2]
\label{III.7.9.2.fr}
Soit $P_\bullet$ un complexe fini de $A$-modules projectifs de type fini, et
pour tout $A$-module $M$, soit
$T_\bullet(M)=H_\bullet(P_\bullet\otimes_A M)$. Alors, pour tout
$\mathfrak p\in\operatorname{Spec}(A)$, on a
\[
  \sum_i(-1)^i\operatorname{rang}_{\kappa(\mathfrak p)}
  T_i\bigl(\kappa(\mathfrak p)\bigr)
  =\sum_i(-1)^i\operatorname{rang}_{\mathfrak p}(P_i).
  \tag{7.9.2.1}\label{III.7.9.2.1.fr}
\]
\end{proposition}

\begin{proof}
En effet, on a par définition
$T_i(\kappa(\mathfrak p))=H_i(P_\bullet\otimes_A\kappa(\mathfrak p))$ et,
compte tenu de (7.9.1.1), la formule (7.9.1.2) n'est autre que l'invariance de
la caractéristique d'Euler-Poincaré d'un complexe fini d'espaces vectoriels de
dimension finie par passage à l'homologie ($0_{\mathrm I}$, 11.10.2).
\end{proof}

\begin{corollary}[7.9.3]
\label{III.7.9.3.fr}
La fonction
\[
  \mathfrak p\longmapsto
  \sum_i(-1)^i\operatorname{rang}_{\kappa(\mathfrak p)}
  T_i\bigl(\kappa(\mathfrak p)\bigr)
\]
est localement constante dans $\operatorname{Spec}(A)$.
\end{corollary}

\begin{theorem}[7.9.4]
\label{III.7.9.4.fr}
Soient $Y$ un préschéma localement noethérien, $f:X\to Y$ un morphisme propre,
$\mathscr P_\bullet$ un complexe fini de $\mathscr O_X$-Modules cohérents et
$Y$-plats. Si l'on pose
\[
  \mathscr T_\bullet(\mathscr M)
  =\mathscr H^{-\bullet}
  \bigl(f,\mathscr P^\bullet\otimes_{\mathscr O_Y}\mathscr M\bigr)
  \qquad\text{(cf. (7.7.1.1))}
\]
la fonction

\oldpage[III]{209}
\[
  y\longmapsto
  \sum_i(-1)^i\operatorname{rang}_{\kappa(y)}
  \mathscr T_i\bigl(\kappa(y)\bigr)
  \tag{7.9.4.1}\label{III.7.9.4.1.fr}
\]
est localement constante dans $Y$.
\end{theorem}

\begin{proof}
On peut se borner au cas où $Y=\operatorname{Spec}(A)$ est affine d'anneau $A$
noethérien. Comme le complexe $\mathscr P_\bullet$ est fini, on sait (7.7.12,
(i)) qu'on a
\[
  \mathscr T_p(\mathscr M)
  =\mathscr H_p(\mathscr L_\bullet\otimes_{\mathscr O_Y}\mathscr M),
  \qquad \mathscr L_\bullet=\widetilde{L_\bullet},
\]
$L_\bullet$ étant un complexe \emph{fini} de $A$-modules projectifs de type
fini. Le théorème résulte alors de (7.9.3).
\end{proof}

\begin{env}[7.9.5]
\label{III.7.9.5.fr}
Sous les conditions de (7.9.4), la fonction (7.9.4.1) est \emph{constante}
lorsque $Y$ est \emph{connexe}. Lorsque $Y$ est connexe et non vide, on désigne
la valeur unique (entière) de (7.9.4.1) par
\[
  \operatorname{EP}(f,\mathscr P_\bullet)
  \quad\text{ou}\quad
  \operatorname{EP}(Y,\mathscr P_\bullet),
  \quad\text{ou simplement}\quad
  \operatorname{EP}(\mathscr P_\bullet)
\]
s'il ne peut en résulter de confusion, et l'on dit que cet entier est la
\emph{caractéristique d'Euler-Poincaré de $\mathscr P_\bullet$} relativement à
$f$ (ou à $Y$). Dans le cas général, on notera aussi
\[
  \operatorname{EP}(f,\mathscr P_\bullet;y)
  \quad\text{ou}\quad
  \operatorname{EP}(Y,\mathscr P_\bullet;y)
  \quad\text{ou}\quad
  \operatorname{EP}(\mathscr P_\bullet;y)
\]
le second membre de (7.9.4.1).
\end{env}

\begin{env}[7.9.6]
\label{III.7.9.6.fr}
Sous les hypothèses de (7.9.4) relativement à $X$, $Y$ et $f$, soit
\[
  0\longrightarrow\mathscr P'_\bullet
  \xrightarrow{u}\mathscr P_\bullet
  \xrightarrow{v}\mathscr P''_\bullet
  \longrightarrow0
\]
une suite exacte de complexes finis de $\mathscr O_X$-Modules cohérents et
$Y$-plats, les homomorphismes $u$ et $v$ étant de degrés \emph{pairs} $2d$,
$2d'$ respectivement. Comme $\mathscr T_\bullet$ est un foncteur homologique
(7.7.1), on a une suite exacte d'homologie
\[
  \cdots\longrightarrow
  \mathscr T_i(\mathscr P'_\bullet,\kappa(y))
  \longrightarrow
  \mathscr T_{i+2d}(\mathscr P_\bullet,\kappa(y))
  \longrightarrow
  \mathscr T_{i+2d+2d'}(\mathscr P''_\bullet,\kappa(y))
  \longrightarrow
  \mathscr T_{i-1}(\mathscr P'_\bullet,\kappa(y))
  \longrightarrow\cdots
\]
n'ayant d'ailleurs qu'un nombre fini de termes. En écrivant que la
caractéristique d'Euler-Poincaré de ce complexe est nulle ($0_{\mathrm I}$,
11.10.1), il vient aussitôt
\[
  \operatorname{EP}(\mathscr P_\bullet;y)
  =\operatorname{EP}(\mathscr P'_\bullet;y)
  +\operatorname{EP}(\mathscr P''_\bullet;y)
  \tag{7.9.6.1}\label{III.7.9.6.1.fr}
\]
pour tout $y\in Y$. Or, si par exemple $\mathscr P_\bullet=(\mathscr P_i)$ avec
$\mathscr P_i=0$ pour $i<0$, on a la suite exacte de complexes
\[
\xymatrix@C=2.4em{
  \cdots \ar[r] & 0 \ar[r]\ar[d] & 0 \ar[r]\ar[d]
    & 0 \ar[r]\ar[d] & 0 \ar[r] & \cdots \\
  \cdots \ar[r] & 0 \ar[r]\ar[d] & 0 \ar[r]\ar[d]
    & \mathscr P_1 \ar[r]\ar[d] & \mathscr P_2 \ar[r] & \cdots \\
  \cdots \ar[r] & 0 \ar[r]\ar[d] & \mathscr P_0 \ar[r]\ar[d]
    & \mathscr P_1 \ar[r]\ar[d] & \mathscr P_2 \ar[r] & \cdots \\
  \cdots \ar[r] & 0 \ar[r]\ar[d] & \mathscr P_0 \ar[r]\ar[d]
    & 0 \ar[r]\ar[d] & 0 \ar[r] & \cdots \\
  \cdots \ar[r] & 0 \ar[r] & 0 \ar[r]
    & 0 \ar[r] & 0 \ar[r] & \cdots
}
\]
les flèches verticales non nulles étant les automorphismes identiques ; on peut
appliquer (7.9.6.1) à cette suite exacte, d'où, par récurrence sur la longueur de
$\mathscr P_\bullet$, la formule
\[
  \operatorname{EP}(\mathscr P_\bullet;y)
  =\sum_i(-1)^i\operatorname{EP}(\mathscr P_i;y).
  \tag{7.9.6.2}\label{III.7.9.6.2.fr}
\]

\oldpage[III]{210}
où, pour tout $\mathscr O_X$-Module cohérent $\mathscr F$, plat sur $Y$, on
désigne par $\operatorname{EP}(\mathscr F;y)$ (ou
$\operatorname{EP}(f,\mathscr F;y)$ ou $\operatorname{EP}(Y,\mathscr F;y)$) la
fonction $\operatorname{EP}(\mathscr Q_\bullet;y)$ correspondant au complexe
$\mathscr Q_\bullet$ dont le seul terme $\ne0$ est de degré $0$ et égal à
$\mathscr F$. On voit donc qu'on peut se ramener à étudier les caractéristiques
d'Euler-Poincaré de complexes réduits à un seul terme.
\end{env}

\begin{proposition}[7.9.7]
\label{III.7.9.7.fr}
Sous les hypothèses de (7.9.4), soient $Y'$ un préschéma localement noethérien,
$g:Y'\to Y$ un morphisme,
\[
  X'=X\times_Y Y',
  \qquad f'=f_{(Y')}:X'\longrightarrow Y',
\]
$\mathscr P'_\bullet$ le complexe fini
\[
  \mathscr P_\bullet\otimes_{\mathscr O_Y}\mathscr O_{Y'}
\]
de $\mathscr O_{X'}$-Modules ; $\mathscr P'_\bullet$ est formé de
$\mathscr O_{X'}$-Modules cohérents et $Y'$-plats, et pour tout $y'\in Y'$, on a
\[
  \operatorname{EP}(\mathscr P'_\bullet;y')
  =\operatorname{EP}(\mathscr P_\bullet;g(y')).
  \tag{7.9.7.1}\label{III.7.9.7.1.fr}
\]
\end{proposition}

\begin{proof}
Les $\mathscr O_{X'}$-Modules $\mathscr P'_i$ étant images réciproques des
$\mathscr P_i$ par la projection $X'\to X$ sont cohérents, ils sont $Y'$-plats en
vertu de ($0_{\mathrm I}$, 6.2.1) et (I, 4.14.5), la question étant locale sur
$X$, $Y$ et $Y'$ ; enfin, on sait que $f'$ est propre (II, 5.4.2), donc le
premier membre de (7.9.7.1) est défini. La formule (7.9.7.1) résulte alors de
(6.10.4.2), (7.7.2) et du lemme (7.6.7), en se ramenant, comme on peut toujours
le faire, au cas où $Y$ et $Y'$ sont affines.
\end{proof}

\begin{proposition}[7.9.8]
\label{III.7.9.8.fr}
Supposons vérifiées les hypothèses de (7.9.4) et en outre qu'il existe un entier
$i_0$ tel que
\[
  \mathscr T_i(\kappa(y))=0
  \qquad\text{pour $i\ne i_0$ et tout $y\in Y$.}
\]
Alors
\[
  \mathscr T_{i_0}(\mathscr O_Y)
  =\mathscr H^{-i_0}(f,\mathscr P_\bullet)
\]
est un $\mathscr O_Y$-Module localement libre, dont le rang en $y\in Y$ est égal
à $(-1)^{i_0}\operatorname{EP}(f,\mathscr P_\bullet;y)$.
\end{proposition}

\begin{proof}
Remarquons d'abord que les hypothèses de (7.4.4) sont vérifiées par les
$\mathscr T_j^{\mathscr O_y}$, donc (7.4.7) leur est applicable, et l'hypothèse
entraîne que $\mathscr T_i^{\mathscr O_y}$ est \emph{nul} pour $i\ne i_0$ en
vertu de (7.5.3) ; en raison de (7.3.3), $\mathscr T_{i_0}$ est donc aussi exact,
et par suite (7.8.4), $\mathscr H^{-i_0}(f,\mathscr P_\bullet)$ est localement
libre et son rang en un point $y\in Y$ est
\[
  \operatorname{rang}_{\kappa(y)}
  \mathscr T_{i_0}(\kappa(y))
  =\operatorname{EP}(f,\mathscr P_\bullet;y)
\]
par définition, puisque $\mathscr T_i(\kappa(y))=0$ pour $i\ne i_0$.
\end{proof}

\begin{corollary}[7.9.9]
\label{III.7.9.9.fr}
Soient $Y$ un préschéma localement noethérien, $f:X\to Y$ un morphisme propre,
$\mathscr F$ un $\mathscr O_X$-Module cohérent et $Y$-plat ; on suppose qu'il
existe un entier $i_0$ tel que
\[
  H^i\bigl(f^{-1}(y),
  \mathscr F\otimes_{\mathscr O_Y}\kappa(y)\bigr)=0
  \qquad\text{pour tout $i\ne i_0$ et tout $y\in Y$.}
\]
Alors $\mathrm R^{i_0}f_*(\mathscr F)$ est un $\mathscr O_Y$-Module localement
libre, dont le rang en $y$ est égal à
$(-1)^{i_0}\operatorname{EP}(f,\mathscr F;y)$.
\end{corollary}

En particulier :

\begin{corollary}[7.9.10]
\label{III.7.9.10.fr}
Sous les conditions préliminaires de (7.9.9) pour $X$, $Y$ et $\mathscr F$,
supposons que l'on ait $\mathrm R^if_*(\mathscr F)=0$ pour tout $i>0$. Alors
$f_*(\mathscr F)$ est un $\mathscr O_Y$-Module localement libre, dont le rang en
$y$ est égal à $\operatorname{EP}(f,\mathscr F;y)$.
\end{corollary}

\begin{proof}
Il suffira, en vertu de (7.9.9) de prouver le lemme suivant :
\end{proof}

\begin{lemma}[7.9.10.1]
\label{III.7.9.10.1.fr}
Sous les hypothèses de (7.9.10), on a
\[
  H^i\bigl(f^{-1}(y),
  \mathscr F\otimes_{\mathscr O_Y}\kappa(y)\bigr)=0
\]
pour tout $i>0$ et tout $y\in Y$.
\end{lemma}

\begin{proof}
En effet, on peut se borner au cas où $Y=\operatorname{Spec}(A)$ est affine.
Avec les notations de (7.9.4), et $\mathscr P_\bullet$ étant réduit à son terme
de degré $0$ égal à $\mathscr F$, on a en effet
$\mathscr T_p(\mathscr O_Y)=0$ pour $p<0$ par hypothèse ; on conclut de (7.3.7)
que $\mathscr T_p$ est exact pour $p<0$, et le lemme résulte alors de
l'équivalence de (7.7.5, a)) et (7.7.5, d)).
\end{proof}

\oldpage[III]{211}
\begin{proposition}[7.9.11]
\label{III.7.9.11.fr}
Les hypothèses étant celles de (7.9.4), soit $\mathscr L$ un
$\mathscr O_X$-Module inversible très ample pour $Y$, et posons
$\mathscr P_\bullet(n)=\mathscr P_\bullet\otimes_{\mathscr O_X}
\mathscr L^{\otimes n}$ pour tout $n\in\mathbf Z$. Alors, pour tout $y\in Y$,
la fonction
\[
  n\longmapsto\operatorname{EP}(f,\mathscr P_\bullet(n);y)
  \tag{7.9.11.1}\label{III.7.9.11.1.fr}
\]
est un polynôme à coefficients dans $\mathbf Q$, qui est le même pour tous les
points d'une même composante connexe de $Y$.
\end{proposition}

\begin{proof}
Il est clair que $\mathscr P_\bullet(n)$ est un complexe de
$\mathscr O_X$-Modules $Y$-plats. En vertu de (7.9.6.2), on peut se borner au cas
où $\mathscr P_\bullet$ est réduit à un seul terme $\mathscr F\ne0$ de degré 0 ;
en outre, comme il s'agit de questions locales sur $Y$, on peut supposer $Y$
affine et $f$ projectif (II, 5.5.3) ; posons $X_y=f^{-1}(y)$, et soit
$\mathscr L_y=\mathscr L\otimes_{\mathscr O_Y}\kappa(y)$, qui est un
$\mathscr O_{X_y}$-Module très ample (II, 4.4.10) ; en vertu de (7.7.2), on a,
pour le foncteur $\mathscr T_\bullet$ relatif au complexe
$\mathscr P_\bullet(n)$,
\[
  \mathscr T_i(\kappa(y))
  =H^{-i}(X_y,\mathscr F_y\otimes\mathscr L_y^{\otimes n})
  \qquad
  \bigl(\text{où }\mathscr F_y=\mathscr F\otimes_{\mathscr O_Y}\kappa(y)\bigr) ;
\]
d'où résulte que $\operatorname{EP}(f,\mathscr F(n);y)$ n'est autre que la
caractéristique d'Euler-Poincaré $\chi_{\kappa(y)}(\mathscr F_y(n))$ définie dans
(2.5.1) ; le fait que (7.9.11.1) soit un polynôme résulte alors de (2.5.3) ; en
outre, pour chaque $n$, sa valeur est constante dans une composante connexe de
$Y$ (7.9.4), ce qui achève la démonstration.
\end{proof}

\begin{env}[7.9.12]
\label{III.7.9.12.fr}
Nous désignerons par $\operatorname{PH}(f,\mathscr P_\bullet;y)$ ou
$\operatorname{PH}(\mathscr P_\bullet;y)$ le polynôme (7.9.11.1), à coefficients
rationnels, et nous dirons que c'est le \emph{polynôme de Hilbert en $y$
relatif à $\mathscr P_\bullet$, $f$ et $\mathscr L$} (ou simplement le
\emph{polynôme de Hilbert en $y$ de $\mathscr P_\bullet$}, ou de $f$, s'il n'en
résulte pas de confusion) ; lorsque $Y$ est connexe non vide, on supprime la
mention de $y$ dans la notation et la terminologie. L'invariant ainsi obtenu
jouera un rôle essentiel au chap. V, dans la théorie des « modules » des
faisceaux cohérents quotients d'un faisceau cohérent donné.

Avec les notations de (7.9.6) et (7.9.11), on a
\[
  \operatorname{PH}(\mathscr P_\bullet;y)
  =\operatorname{PH}(\mathscr P'_\bullet;y)
  +\operatorname{PH}(\mathscr P''_\bullet;y)
  \tag{7.9.12.1}\label{III.7.9.12.1.fr}
\]
et en particulier
\[
  \operatorname{PH}(\mathscr P_\bullet;y)
  =\sum_i(-1)^i\operatorname{PH}(\mathscr P_i;y) ;
  \tag{7.9.12.2}\label{III.7.9.12.2.fr}
\]
cela résulte trivialement de (7.9.6.1) et (7.9.6.2). De même, avec les notations
et hypothèses de (7.9.7), on a
\[
  \operatorname{PH}(\mathscr P'_\bullet;y')
  =\operatorname{PH}(\mathscr P_\bullet;g(y')).
  \tag{7.9.12.3}\label{III.7.9.12.3.fr}
\]
\end{env}

La formule (7.9.12.2) ramène l'étude des polynômes de Hilbert d'un complexe à
celle des polynômes de Hilbert d'un seul $\mathscr O_X$-Module $Y$-plat. Ces
derniers admettent une interprétation remarquable indépendante de considérations
homologiques :

\begin{corollary}[7.9.13]
\label{III.7.9.13.fr}
Soient $Y$ un préschéma noethérien, $f:X\to Y$ un morphisme propre,
$\mathscr L$ un $\mathscr O_X$-Module inversible très ample pour $Y$,
$\mathscr F$ un $\mathscr O_X$-Module cohérent et $Y$-plat. Il existe un entier
$n_0$ tel que pour $n\geq n_0$, $f_*(\mathscr F(n))$ soit un
$\mathscr O_Y$-Module localement libre, de rang en $y\in Y$ égal à
$\operatorname{PH}(f,\mathscr F;y)(n)$.
\end{corollary}

\begin{proof}
Comme le morphisme $f$ est projectif (II, 5.5.3), il existe $n_0$ tel que pour
$n\geq n_0$

\oldpage[III]{212}
on ait $\mathrm R^if_*(\mathscr F(n))=0$ pour tout $i>0$ (2.2.1) ; la conclusion
résulte donc de (7.9.10).
\end{proof}

Le critère de platitude suivant sera important dans la théorie des « modules »
des faisceaux cohérents du chap. V :

\begin{proposition}[7.9.14]
\label{III.7.9.14.fr}
Soient $Y$ un préschéma noethérien, $f:X\to Y$ un morphisme projectif,
$\mathscr L$ un $\mathscr O_X$-Module inversible ample pour $f$, et posons
$\mathscr F(n)=\mathscr F\otimes\mathscr L^{\otimes n}$ pour tout
$\mathscr O_X$-Module $\mathscr F$ et tout $n\in\mathbf Z$. Pour qu'un
$\mathscr O_X$-Module cohérent $\mathscr F$ soit $Y$-plat, il faut et il suffit
qu'il existe un entier $n_0$ tel que, pour tout $n\geq n_0$,
$f_*(\mathscr F(n))$ soit un $\mathscr O_Y$-Module localement libre.
\end{proposition}

\begin{proof}
La nécessité de la condition se démontre comme dans (7.9.13) (le résultat de
(2.2.1) s'appliquant à un faisceau ample $\mathscr L$, puisque $f$ est
projectif). Pour démontrer la réciproque, on peut se borner au cas où $Y$ est
affine d'anneau $A$ ; en vertu de l'hypothèse et de (2.2.2, (i)), les
$A$-modules $\Gamma(X,\mathscr F(n))$ sont de type fini et projectifs (Bourbaki,
\emph{Alg. comm.}, chap. II, \S~5, n\textsuperscript{o}~2, th. 1). Soit $S$
l'anneau gradué
\[
  S=\bigoplus_{n\geq0}\Gamma(X,\mathscr L^{\otimes n}) ;
\]
on sait que $X$ s'identifie canoniquement à $\operatorname{Proj}(S)$ (II,
4.5.2, (b) et 5.4.4). Soit
\[
  M=\bigoplus_{n\geq n_0}\Gamma(X,\mathscr F(n)) ;
\]
remplaçant au besoin $\mathscr L$ par une puissance $\mathscr L^{\otimes d}$,
on peut supposer que $S$ est engendré par un nombre fini d'éléments de degré 1
(2.3.5.1), et il résulte alors de (II, 2.7.5 et 2.7.2) que $\mathscr F$
s'identifie à $\mathscr P\!roj_0(M)$. Pour tout élément homogène $g\in S$ de
degré $>0$, on a donc $\Gamma(X_g,\mathscr F)=M_{(g)}$ ; or, $M$, somme directe
de $A$-modules projectifs, est un $A$-module plat, donc il en est de même de
$M_g$ ($0_{\mathrm I}$, 6.3.2), et par suite aussi de $M_{(g)}$, qui est un
facteur direct de $M_g$ ($0_{\mathrm I}$, 6.1.2). On en conclut (I, 4.14.5) que
$\mathscr F$ est $Y$-plat en tout point de $X_g$, et comme les $X_g$ recouvrent
$X$, la proposition est démontrée.
\end{proof}

\bigskip
\noindent\emph{(A suivre.)}

\clearpage
\oldpage[III]{213}
\section*{INDEX DES NOTATIONS}

\begin{flushleft}
$\mathbf{Tor}_n^A(P_\bullet,Q_\bullet)$ : 6.3.1.

\medskip
$\mathscr{T}\!or_n^{\mathscr O_S}(\mathscr P_\bullet,\mathscr Q_\bullet)$,
$\mathscr{T}\!or_n^S(\mathscr P_\bullet,\mathscr Q_\bullet)$ : 6.4.1 et 6.5.1.

\medskip
$\mathscr{T}\!or_n^S(\mathscr F,\mathscr G)$ : 6.5.1.

\medskip
$\mathscr{T}\!or_q^S(\mathscr P_\bullet^{(1)},\mathscr P_\bullet^{(2)},\ldots,
\mathscr P_\bullet^{(m)})$ : 6.5.15.

\medskip
$\mathbf H^{-n}(\mathfrak U^{(i)},\mathscr P_\bullet^{(i)})$,
$\mathbf H^{-n}(X^{(i)},\mathscr P_\bullet^{(i)})$ : 6.6.1.

\medskip
$\mathbf{Tor}_n^A(L_{\bullet\bullet}^{(1)},L_{\bullet\bullet}^{(2)})$,
$\mathbf{Tor}_n^S(\mathfrak U^{(1)},\mathfrak U^{(2)};
\mathscr P_\bullet^{(1)},\mathscr P_\bullet^{(2)})$ : 6.6.2.

\medskip
$\mathscr{T}\!or_n^S(\mathfrak U^{(1)},\mathfrak U^{(2)};
\mathscr F^{(1)},\mathscr F^{(2)})$,
$\mathscr{T}\!or_n^S(\mathfrak U^{(1)},\mathfrak U^{(2)};
\mathscr P_\bullet^{(1)},\mathscr P_\bullet^{(2)})$ : 6.6.2.

\medskip
$\mathscr{T}\!or_n^S(f_1,f_2;\mathscr P_\bullet^{(1)},
\mathscr P_\bullet^{(2)})$ : 6.6.8, 6.7.1 et 6.7.3.

\medskip
$\mathscr{T}\!or_n^S((f_i)_{i\in I};(\mathscr P_\bullet^{(i)})_{i\in I})$,
$\mathscr{T}\!or_n^S(f_1,\ldots,f_m;\mathscr P_\bullet^{(1)},\ldots,
\mathscr P_\bullet^{(m)})$ : 6.7.11.

\medskip
$\mathbf{Ab}$, $\mathbf{Ab}_A$ : 7.1.1.

\medskip
$T_{(B)}$, $T^{(B)}$, $T\otimes_A B$
($T$ foncteur de $\mathbf{Ab}_A$ dans $\mathbf{Ab}$) : 7.1.3.

\medskip
$t_M$ : 7.2.2.

\medskip
$T_{\mathfrak p}$ : 7.1.4.

\medskip
$\mathscr T_\bullet^{Y'}(\mathscr M')$,
$\mathscr T_\bullet^U(\mathscr M|U)$ : 7.7.2.

\medskip
$\operatorname{EP}(f,\mathscr P_\bullet;y)$,
$\operatorname{EP}(Y,\mathscr P_\bullet;y)$,
$\operatorname{EP}(\mathscr P_\bullet;y)$ : 7.9.5.

\medskip
$\operatorname{PH}(f,\mathscr P_\bullet;y)$,
$\operatorname{PH}(\mathscr P_\bullet;y)$ : 7.9.12.
\end{flushleft}

\clearpage
\oldpage[III]{214}
\section*{INDEX TERMINOLOGIQUE}

\begin{flushleft}
Caractéristique d'Euler-Poincaré d'un complexe de Modules : 7.9.5.

\medskip
Cohomologiquement plat en un point $y\in Y$ en dimension $p$,
cohomologiquement plat sur $Y$ en dimension $p$, cohomologiquement plat sur
$Y$ au point $y$, cohomologiquement plat sur $Y$ (complexe de
$\mathscr O_X$-Modules $Y$-plats) : 7.8.1.

\medskip
Complexes homotopes : 6.1.4.

\medskip
Espace annelé de dimension cohomologique $\leq n$ : 6.5.5.

\medskip
Extension des scalaires dans un foncteur covariant additif : 7.1.3.

\medskip
Faisceau d'anneaux de dimension cohomologique $\leq n$ : 6.5.5.

\medskip
Formule de Künneth : 6.7.8.

\medskip
Homologiquement plat en un point $y\in Y$ en dimension $p$,
homologiquement plat sur $Y$ en dimension $p$, homologiquement plat sur $Y$
au point $y$, homologiquement plat sur $Y$ (complexe de $\mathscr O_X$-Modules
$Y$-plats) : 7.8.1.

\medskip
Homotopisme : 6.1.4.

\medskip
Hypertor de deux complexes de $A$-modules : 6.3.1.

\medskip
Hypertor local de deux complexes de Modules : 6.4.1.

\medskip
Hypertor global de deux complexes de Modules : 6.6.2, 6.6.8 et 6.7.3.

\medskip
Polynôme de Hilbert relatif à un complexe de Modules : 7.9.12.

\medskip
Suites spectrales de Künneth : 6.7.3.
\end{flushleft}

\clearpage
\oldpage[III]{215}
\section*{TABLE DES MATIÈRES}

\begin{flushleft}
\textsc{Chapitre III.} --- \textbf{Étude cohomologique des faisceaux cohérents
(suite)} \dotfill 5

\medskip
\S~6. Foncteurs Tor locaux et globaux ; formule de Künneth \dotfill 5

6.1. Introduction \dotfill 5

6.2. Hypercohomologie des complexes de Modules sur un préschéma \dotfill 6

6.3. Hypertor de deux complexes de modules \dotfill 9

6.4. Foncteurs hypertor locaux de complexes de Modules quasi-cohérents ; cas
des schémas affines \dotfill 14

6.5. Foncteurs hypertor locaux de complexes de Modules quasi-cohérents : cas
général \dotfill 16

6.6. Foncteurs hypertor globaux de complexes de Modules quasi-cohérents et
suites spectrales de Künneth : cas de la base affine \dotfill 21

6.7. Foncteurs hypertor globaux de complexes de Modules quasi-cohérents et
suites spectrales de Künneth : cas général \dotfill 25

6.8. Les suites spectrales d'associativité des hypertor globaux \dotfill 32

6.9. Les suites spectrales de changement de base dans les hypertor globaux
\dotfill 34

6.10. Structure locale de certains foncteurs cohomologiques \dotfill 39

\medskip
\S~7. Étude du changement de base dans les foncteurs homologiques covariants
de Modules \dotfill 43

7.1. Foncteurs de $A$-modules \dotfill 43

7.2. Caractérisation du foncteur produit tensoriel \dotfill 44

7.3. Critères d'exactitude des foncteurs homologiques de modules \dotfill 48

7.4. Critères d'exactitude pour les foncteurs
$\mathrm H_i(P\otimes_A M)$ \dotfill 53

7.5. Cas des anneaux locaux noethériens \dotfill 58

7.6. Descente des propriétés d'exactitude. Théorème de semi-continuité et
critère d'exactitude de Grauert \dotfill 60

7.7. Application aux morphismes propres : I. La propriété d'échange
\dotfill 65

7.8. Application aux morphismes propres : II. Critères de platitude
cohomologique \dotfill 72

7.9. Application aux morphismes propres : III. Invariance de la
caractéristique d'Euler-Poincaré et du polynôme de Hilbert \dotfill 76

\begin{flushright}
\emph{Reçu le 15 décembre 1962.}
\end{flushright}
\end{flushleft}

% La page imprimée 216 est blanche.
\clearpage
\thispagestyle{empty}
\null
\clearpage

\oldpage[III]{217}
\section*{ERRATA ET ADDENDA}
\begin{center}
(Liste 2)

\medskip
A) \emph{Erreurs typographiques}
\end{center}

\noindent\textbf{$(0_{\mathrm I},\,3.2.1)$} Ligne 2 de la p. 27, remplacer $X$
par $U$.

\noindent\textbf{$(0_{\mathrm I},\,3.2.6)$} Ligne 1 du bas de la p. 27 et ligne
1 de la p. 28, remplacer (3 fois) $\mathfrak S_\lambda$ par
$\mathscr H_\lambda$.

\noindent\textbf{$(0_{\mathrm I},\,3.4.5)$} Ajouter une parenthèse devant
3.4.5).

\noindent\textbf{$(0_{\mathrm I},\,4.2.1)$} Ligne 10 de la p. 39, remplacer
$\psi_*(A)$ par $\psi_*(\mathscr A)$.

\noindent\textbf{$(0_{\mathrm I},\,5.1.3)$} Ligne 16 de la p. 45, remplacer
$\mathscr F|V$ par $\mathscr F|U$.

\noindent\textbf{$(0_{\mathrm I},\,5.3.9)$} Ligne 12 de la p. 48, avant
« telle que », ajouter : au-dessus d'un voisinage ouvert $U$ de $x$.

\noindent\textbf{$(0_{\mathrm I},\,7.6.9)$} Ligne 3 du bas de la p. 73,
remplacer $I_\lambda$ par $J_\lambda$.

\noindent\textbf{(I, 1.1.15)} Ligne 3 de la p. 83, remplacer $\mathfrak a$ par
$\mathfrak a\ne A$.

\noindent\textbf{(I, 1.4.1)} Ligne 2 du bas de la p. 90, remplacer
$\widetilde N$ par $N$.

\noindent\textbf{(I, 2.3.1)} Ligne 3 du bas de la p. 101, remplacer (0, 4.1.6)
par (0, 4.1.7).

\noindent\textbf{(I, 2.3.2)} Ligne 11 de la p. 101, remplacer $B$ par $B_s$ et
$C$ par $C_t$.

\noindent\textbf{(I, 2.5.5)} Ligne 19 de la p. 104, remplacer $S$-morphisme
par $S$-préschéma.

\noindent\textbf{(I, 3.3.7)} Ligne 17 de la p. 109, remplacer $Y_{(S)}$ par
$Y_{(S')}$.

\noindent\textbf{(I, 3.4.5)} Ligne 9 de la p. 113, remplacer $\mathrm s$ par
$s$.

\noindent\textbf{(I, 3.4.8)} Ligne 4 de la p. 114, remplacer $p(x)$ par
$f(x)$.

\noindent\textbf{(I, 3.5.1)} Ligne 3 de la p. 115, remplacer $1_X\times g$
par $1_{X'}\times g$.

\noindent\textbf{(I, 3.7.2)} Ligne 8 de la p. 119, remplacer le premier $X$
par $X'$.

\noindent\textbf{(I, 4.2.2)} Ligne 8 de la p. 123, dans la flèche verticale de
gauche du diagramme, remplacer $\alpha_{\psi(y)}$ par $\varphi_{\psi(y)}$.

\noindent\textbf{(I, 4.4.3)} Ligne 16 de la p. 126, remplacer isomorphée par
isomorphes.

\noindent\textbf{(I, 4.5.5)} Ligne 17 du bas de la p. 127, remplacer (4.2.4)
par (4.2.5).

\noindent\textbf{(I, 5.1.4)} Ligne 14 du bas de la p. 128, remplacer (2.1.7)
par (2.1.8).

\noindent\textbf{(I, 5.3.13)} Ligne 20 de la p. 134, remplacer (4.2.4) par
(4.2.5).

\noindent\textbf{(I, 6.4.2)} Ligne 1 du bas de la p. 147, remplacer
recouvrement par recouvrement fini.

\noindent\textbf{(I, 9.1.13)} Ligne 15 de la p. 171, remplacer
$p^{-1}(\mathscr F)$ par $p^*(\mathscr F)$.

\noindent\textbf{(I, 9.5.11)} Ligne 7 de la p. 179, remplacer $Y$ par $Y'$.

\noindent\textbf{(I, 9.6.5)} Ligne 17 du bas de la p. 180, remplacer
$(0,\,4.1.4)$ par $(0,\,4.1.3)$.

\noindent\textbf{(I, 10.12.2)} Ligne 14 du bas de la p. 206, remplacer
$(X_n)_{(S_n)}$ par $(X_n)_{(S_m)}$.

\clearpage
\oldpage[III]{218}
\noindent\textbf{(I), Index terminologique :} Ligne 6 du bas de la p. 219,
remplacer 0, 4.1.4 par 0, 4.1.5. Lignes 16, 17, 18 du bas de la p. 222,
remplacer I, 2.1.7 par I, 2.1.8 ; ligne 19 du bas de la p. 222, remplacer I,
2.1.8 par I, 2.1.7.

\noindent\textbf{(II, 1.5.2)} Ligne 17 de la p. 12, insérer un $f$ à côté de
la flèche verticale de gauche du diagramme.

\noindent\textbf{(II, 4.2.7)} Ligne 16 du bas de la p. 75, remplacer (III,
2.1.14) par (III, 2.1.13).

\noindent\textbf{(II, 5.2.1)} Ligne 16 du bas de la p. 97, remplacer $X-U$ par
$U$. Ligne 8 du bas de la p. 97, remplacer $X'$ par $X_f$.

\noindent\textbf{(II, 5.3.6)} Ligne 16 de la p. 100, remplacer (4.6.18) par
(4.6.17).

\noindent\textbf{(II, 6.2.7)} Ligne 9 du bas de la p. 116, remplacer chap. V
par chap. IV.

\noindent\textbf{(II, 6.6.5)} Ligne 15 du bas de la p. 132, remplacer chap. V
par chap. IV.

\noindent\textbf{(II, 7.4.12)} Ligne 1 de la p. 151, remplacer chap. V par
chap. III.

\noindent\textbf{(II, 8.1.4)} Ligne 15 de la p. 154, remplacer (III, 2.3.8)
par (III, 2.3.7).

\noindent\textbf{(II, 8.10.5)} Ligne 17 du bas de la p. 187, remplacer
$g_{(j)x}$ par $g_{j(x)}$.

\noindent\textbf{$(0_{\mathrm{III}},\,10.2.7)$} Ligne 1 de la p. 20,
remplacer $u$-adique par $\mathfrak m$-adique. Ligne 13 du bas de la p. 20,
remplacer $k$ par $k_i$.

\noindent\textbf{$(0_{\mathrm{III}},\,11.1.1)$} Ligne 18 de la p. 23,
remplacer $\operatorname{Im}(d_r^{p+r,q-r+1})$ par
$\operatorname{Im}(d_r^{p-r,q+r-1})$.

\noindent\textbf{$(0_{\mathrm{III}},\,11.8.6)$} Ligne 14 du bas de la p. 45,
remplacer les flèches $\to$ par $\rightsquigarrow$ (4 fois).

\noindent\textbf{$(0_{\mathrm{III}},\,12.3.4)$} Ligne 4 de la p. 61,
remplacer $\rightsquigarrow$ par $\to$.

\noindent\textbf{$(0_{\mathrm{III}},\,13.2.4)$} Ligne 2 du bas de la p. 67,
remplacer $\varinjlim$ par $\varprojlim$ (2 fois).

\noindent\textbf{(III, 1.1.7)} Ligne 15 de la p. 85, remplacer $A^r$ par
$\bigwedge(A^r)$ (2 fois).

\noindent\textbf{(III, 1.2.4)} Ligne 7 de la p. 87, remplacer
$\Gamma(U,\mathscr F)$ par $\Gamma(U,\mathscr O_X)$ et remplacer $\mathscr G$
par $\mathscr F$. Ligne 8 de la p. 87, remplacer
$\mathrm H^p(\mathfrak U,\mathscr G)$ par $\mathrm H^p(\mathfrak U,\mathscr F)$.

\noindent\textbf{(III, 2.5.3.1)} Ligne 3 de la p. 111, remplacer $n>0$ par
$n\geq0$. Ligne 8 de la p. 111, remplacer $-r\leq n\leq0$ par
$-r\leq n<0$.

\noindent\textbf{(III, 3.1.2)} Ligne 13 de la p. 116, remplacer $\mathscr G$
par $\mathscr G\in\mathbf K'$.

\noindent\textbf{(III, 4.3.6)} Ligne 7 de la p. 133, remplacer
$K'=K\otimes_A\widehat A$ par $K'\supset K\otimes_A\widehat A$. Ligne 16 de
la p. 133, remplacer $R(X)$ par $R(Y)$.

\noindent\textbf{(III, 4.4.12)} Ligne 10 de la p. 138, remplacer chap. V par
chap. IV.

\noindent\textbf{(III, 4.6.6)} Ligne 19 de la p. 142, remplacer « chap. V » par
« dans un paragraphe ultérieur ».

\begin{center}
B) \emph{Modifications de texte}\footnote{Pour faciliter la recherche des
références, les modifications appartenant à la liste d'errata insérée dans le
chapitre N seront désormais désignées par $\mathbf{Err}_{\mathrm N}$ suivi d'un
numéro.}
\end{center}

\noindent\textbf{$(\mathbf{Err}_{\mathrm{III}},\,1)$} Dans
$(0_{\mathrm I},\,5.1.3)$, ligne 18 de la p. 45, remplacer « somme directe »
par « somme directe finie ».

\noindent\textbf{$(\mathbf{Err}_{\mathrm{III}},\,2)$} Dans
$(0_{\mathrm I},\,5.4.3)$, lignes 8 à 4 du bas de la p. 49, remplacer depuis
« Dans le cas général\ldots{} » par le texte suivant :

Dans le cas général où $X$ est un espace annelé tel que $\mathscr O_x$ soit un
anneau \emph{local} pour
\oldpage[III]{219}
tout $x\in X$, si $\mathscr L$ est un $\mathscr O_X$-Module \emph{de type fini}
tel qu'il existe un $\mathscr O_X$-Module $\mathscr F$ pour lequel
$\mathscr L\otimes_{\mathscr O_X}\mathscr F$ soit \emph{isomorphe à}
$\mathscr O_X$, le raisonnement précédent montre que pour tout $x\in X$,
$\mathscr L_x$ est isomorphe à $\mathscr O_x$. On en déduit que $\mathscr L$ est
alors \emph{inversible}. En effet, pour tout $x\in X$, soient $U$ un voisinage
ouvert de $x$ tel que $\mathscr L|U$ soit engendré par $n$ sections
$s_i\ (1\leq i\leq n)$ au-dessus de $U$ (5.2.1). On peut supposer par exemple
que $s=s_1$ est telle que $s_x\ne0$, et comme $\mathscr L_x$ est isomorphe à
$\mathscr O_x$, il existe pour chaque $i$ une section $t_i$ de $\mathscr O_X$
au-dessus d'un voisinage ouvert $V_i\subset U$ de $x$ telle que
$(t_i)_x s_x=(s_i)_x$. Il y a par suite un voisinage ouvert
$V\subset\bigcap_{i=1}^nV_i$ de $x$ tel que $t_i s=s_i$ dans $V$, autrement
dit $\mathscr L|V$ est engendré par \emph{l'unique} section $s|V$. En outre,
si $z$ est une section au-dessus d'un ouvert $W\subset V$ du noyau de
l'homomorphisme $\mathscr O_X|V\to\mathscr L|V$ défini par $s$ (5.1.1), $z_y$
annule $\mathscr L_y$ pour tout $y\in W$, donc $z_y=0$ par hypothèse et par
suite $z=0$, ce qui achève de prouver notre assertion. De plus, la
considération du produit tensoriel
$\mathscr L^{-1}\otimes\mathscr L\otimes\mathscr F$ montre aussitôt que
$\mathscr F$ est isomorphe à $\mathscr L^{-1}$.

\noindent\textbf{$(\mathbf{Err}_{\mathrm{III}},\,3)$} Dans
$(0_{\mathrm I},\,7.2.4)$, il faut imposer la condition que les puissances
$\mathfrak J^n$ sont des idéaux \emph{fermés} dans l'anneau admissible $A$ pour
que la proposition soit exacte. La démonstration donnée est incorrecte, et
l'énoncé rectifié résulte de Bourbaki, \emph{Top. gén.}, chap. III,
3\textsuperscript{e} éd., \S~3, n\textsuperscript{o}~5, cor. 1 de la prop. 9.
Il faut de même supposer les $\mathfrak J^n$ \emph{fermés} dans $A$ dans les
énoncés $(0_{\mathrm I},\,7.2.5)$ et $(0_{\mathrm I},\,7.2.6)$.

\noindent\textbf{$(\mathbf{Err}_{\mathrm{III}},\,4)$} Après la proposition
(I, 2.2.5), ajouter : avec les notations de (I, 2.2.5), si $\mathfrak J$ est un
idéal de $A$, on note $\widetilde{\mathfrak J\mathscr F}$ le sous-$\mathscr O_X$-Module
$\mathfrak J\widetilde{\mathscr F}$ défini dans $(0,\,4.3.5)$.

\noindent\textbf{$(\mathbf{Err}_{\mathrm{III}},\,5)$} Dans (I, 3.4.5), ligne 1
du bas de la p. 112, remplacer « \emph{un corps} $K$ » par « \emph{un corps
algébriquement clos} $K$ ». Ligne 1 de la p. 113, remplacer « \emph{extension} »
par « \emph{extension algébriquement close} ». Lignes 1 à 4 de la p. 113,
remplacer depuis « $K$ sera appelé\ldots{} » par le texte suivant :

Pour tout point de $X$ à valeurs dans un corps $K$, $K$ sera appelé le
\emph{corps des valeurs} du point correspondant, et si $x$ est la localité de
ce point on dit encore que ce dernier est \emph{localisé en} $x$. On définit
ainsi une application $X(K)\to X$, faisant correspondre à un point à valeurs
dans $K$ sa localité.

Lignes 7 et 16 de la p. 113, supprimer « géométrique » ; lignes 10 et 13 de
la p. 113, supprimer « géométriques ». Lignes 10 et 11 du bas de la p. 115,
supprimer « \emph{géométrique} ».

\noindent\textbf{$(\mathbf{Err}_{\mathrm{III}},\,6)$} À la fin de la
démonstration de (I, 3.6.1), ligne 12 du bas de la p. 117, ajouter : Si l'on
pose
$X'=X\times_Y\operatorname{Spec}(\mathscr O_y/\mathfrak a_y)$, alors, pour tout
point $x\in X'$, identifié par $p$ à un point de $X$, on a
$\mathscr O_{X',x}=\mathscr O_{X,x}/\mathfrak a_y\mathscr O_{X,x}$. La question
étant en effet locale sur $X$ et $Y$, on peut supposer que
$X=\operatorname{Spec}(B)$, $Y=\operatorname{Spec}(A)$, et l'on a
$(B/\mathfrak a_yB)_x=B_x/\mathfrak a_yB_x$ par platitude $(0,\,1.3.2)$.

\noindent\textbf{$(\mathbf{Err}_{\mathrm{III}},\,7)$} Dans (I, 5.1.1), ligne 3
du bas de la p. 127, remplacer « $\mathscr O_X$-Module quasi-cohérent » par
« \emph{Idéal quasi-cohérent de} $\mathscr B$ ».

\noindent\textbf{$(\mathbf{Err}_{\mathrm{III}},\,8)$} Après (I, 5.1.10),
ligne 16 de la p. 131, ajouter :

\emph{Remarque (5.1.11).} --- Une légère adaptation du raisonnement fait dans
(5.1.9) montre que la conclusion reste valable \emph{sans supposer a priori
que $X$ soit un préschéma}, mais

\oldpage[III]{220}
en supposant seulement que $(X,\mathscr O_X)$ soit un espace annelé en anneaux
locaux, $\mathscr J$ un Idéal de $\mathscr O_X$ tel que $\mathscr J^n=0$, que
l'espace annelé $X_0=(X,\mathscr O_X/\mathscr J)$ soit un schéma affine et enfin
que les Idéaux $\mathscr J^k/\mathscr J^{k+1}$ soient des
$\mathscr O_{X_0}$-Modules quasi-cohérents. Il suffit en effet d'utiliser
(1.8.1) (cf. $(\mathbf{Err}_{\mathrm{II}})$) au lieu de (2.2.4) dans le
raisonnement.

\noindent\textbf{$(\mathbf{Err}_{\mathrm{III}},\,9)$} Dans (I, 5.3.1), ligne
11 de la p. 132, insérer, après « ou $\Delta_X$ », « ou $\Delta_\varphi$ si
$\varphi:X\to S$ est le morphisme structural ».

\noindent\textbf{$(\mathbf{Err}_{\mathrm{III}},\,10)$} Dans (I, 5.3.9), la
démonstration est insuffisante, car elle ne prouve pas que $\Delta_X(X)$ soit
localement fermé dans $X\times_SX$. Pour avoir une démonstration correcte, il
suffit d'utiliser (4.2.4, $a$)) : pour tout $x\in X$ et tout voisinage affine
$U$ de $x$ dans $X$, $U\times_SU$ est un voisinage affine de $\Delta_X(x)$ ;
compte tenu de (5.3.16) (dont la démonstration n'utilise que la définition
(5.3.1.1)), on est ramené à démontrer (5.3.9) lorsque
$S=\operatorname{Spec}(B)$ et $X=\operatorname{Spec}(A)$ sont des schémas
affines ; il est clair alors en vertu de (5.3.1.1) que $\Delta_X$ correspond à
l'homomorphisme canonique $A\otimes_BA\to A$ qui transforme $x\otimes y$ en
$xy$ ; cet homomorphisme étant surjectif, $\Delta_X$ est dans ce cas une
immersion fermée (4.2.3), ce qui achève la démonstration.

\noindent\textbf{$(\mathbf{Err}_{\mathrm{III}},\,11)$} Dans (I, 7.1.15) et
(I, 7.1.16), lignes 7 et 15 du bas de la p. 158, supprimer
« \emph{géométriques} ».

\noindent\textbf{$(\mathbf{Err}_{\mathrm{III}},\,12)$} Remplacer (I, 7.4.7)
par : On étend (par abus de langage) les définitions de (7.4.1) au cas où
$X$ est un préschéma \emph{réduit} dont tout point admet un voisinage ouvert
n'ayant qu'un nombre \emph{fini} de composantes irréductibles ; il résulte
alors de (7.3.4) et (7.4.6) que, pour un $\mathscr O_X$-Module quasi-cohérent de
type fini $\mathscr F$, dire que $\mathscr F$ est un \emph{faisceau de torsion}
équivaut à dire que $\operatorname{Supp}(\mathscr F)$ \emph{ne contient aucune
composante irréductible de $X$}.

\noindent\textbf{$(\mathbf{Err}_{\mathrm{III}},\,13)$} Dans (I, 10.11.7),
lignes 17 à 23 de la p. 205, la fin de la démonstration depuis « Reste à
prouver\ldots{} » est inutile en vertu de (10.11.6), les faisceaux considérés
étant cohérents par $(0,\,5.3.5)$.

\noindent\textbf{$(\mathbf{Err}_{\mathrm{III}},\,14)$} Dans (II, 1.7.8),
après la ligne 10 du bas de la p. 16, ajouter : lorsque
$\mathscr E=\mathscr O_S^n$, on écrit aussi $\mathbf V_S^n$ au lieu de
$\mathbf V(\mathscr E)$ ; si en outre $S=\operatorname{Spec}(A)$ est affine,
on écrit $\mathbf V_A^n$ au lieu de $\mathbf V_S^n$ ; on a
$\mathbf V_A^n=\operatorname{Spec}(A[T_1,\ldots,T_n])$ où les $T_i$ sont des
indéterminées.

\noindent\textbf{$(\mathbf{Err}_{\mathrm{III}},\,15)$} Dans (II, 1.7.10),
ligne 10 de la p. 17, supprimer « \emph{géométriques} » ; lignes 12 et 16--17
de la p. 17, supprimer « géométrique ».

\noindent\textbf{$(\mathbf{Err}_{\mathrm{III}},\,16)$} Dans (II, 1.7.12),
ajouter : On pose de même :
\[
  \mathbf V(\mathscr O_S^n)=\mathbf V_S^n=S[T_1,\ldots,T_n].
\]

\noindent\textbf{$(\mathbf{Err}_{\mathrm{III}},\,17)$} Dans (II, 4.2.6),
ligne 7 de la p. 75, supprimer « \emph{géométriques} » ; ligne 8 de la p.
75, supprimer « \emph{géométrique} ».

\noindent\textbf{$(\mathbf{Err}_{\mathrm{III}},\,18)$} Dans (II, 4.6.13),
ligne 7 de la p. 92, le raisonnement doit être précisé, car avec les notations
de (4.4.10), il faut ici prouver que l'immersion $\Gamma_f$ est
\emph{quasi-compacte}, afin d'appliquer (i bis). En vertu de (4.6.4), pour
prouver que $\mathscr L$ est ample relativement à $f$, on peut se borner au
cas où $Y$ est affine. Notons d'autre part que dans chacune des hypothèses de
(v), $f$ est quasi-compact (I, 6.6.4). D'autre part, si $g$ est
\oldpage[III]{221}
séparé, $\Gamma_f$ est une immersion fermée (I, 5.4.3) donc quasi-compacte
(I, 6.6.4). Si au contraire $X$ est localement noethérien, comme $Y$ est affine
et $f$ quasi-compact, l'espace sous-jacent à $X$ est quasi-compact, donc
noethérien, et on peut de nouveau appliquer (I, 6.6.4) pour prouver que
$\Gamma_f$ est quasi-compacte.

\noindent\textbf{$(\mathbf{Err}_{\mathrm{III}},\,19)$} L'énoncé de (II,
5.2.2) est incorrect, les conditions $b)$, $c)$ et $c')$ n'étant pas locales
sur $Y$ lorsqu'on ne sait pas si pour un ouvert $U$ de $Y$, un
$(\mathscr O_X|f^{-1}(U))$-Module quasi-cohérent est restriction à $f^{-1}(U)$
d'un $\mathscr O_X$-Module quasi-cohérent. Il faut donc remplacer dans la ligne
16 de la p. 98 « \emph{Soit $f:X\to Y$ un morphisme séparé quasi-compact} »
par : « \emph{Soient $X$, $Y$ deux préschémas tels que $X$ soit un schéma ou
que l'espace sous-jacent à $X$ soit localement noethérien, et soit
$f:X\to Y$ un morphisme quasi-compact.} » Dans la démonstration, on observera
que les hypothèses entraînent que $f$ est \emph{séparé} lorsqu'on suppose que
$X$ est un schéma, en vertu de (I, 5.5.5 et 5.5.8) ; on peut donc appliquer
(I, 9.2.2) dans les deux cas.

\noindent\textbf{$(\mathbf{Err}_{\mathrm{III}},\,20)$} Dans (II, 6.2.3),
ajouter, après la ligne 18 du bas de la p. 115 :

\emph{On dit qu'un morphisme $f:X\to Y$ est quasi-fini en un point $x\in X$
s'il existe un voisinage ouvert affine $V$ de $y=f(x)$ et un voisinage ouvert
affine $U$ de $x$ tels que $f(U)\subset V$ et que le morphisme $U\to V$
restriction de $f$ soit quasi-fini. On dit qu'un morphisme $f:X\to Y$ est
localement quasi-fini s'il est quasi-fini en tout point de $X$.}

\noindent\textbf{$(\mathbf{Err}_{\mathrm{III}},\,21)$} Dans (II, 6.4.3), la
démonstration est incorrecte, les éléments d'un sous-$A$-module de type fini de
$K$ n'étant pas nécessairement entiers sur $A$. Remplacer les 9 dernières lignes
de la p. 121 par le texte suivant : Comme on peut supposer que $E$ est sans
torsion, donc fidèle, $E$ est aussi un $A[u]$-module fidèle ($A$ étant plongé
dans l'anneau des endomorphismes de $E$). Comme $E$ est un $A$-module de type
fini, il en résulte que $u$ est \emph{entier} sur $A$ (Bourbaki,
\emph{Alg. comm.}, chap. V, \S~1, n\textsuperscript{o}~1, lemme 1). On en
conclut immédiatement que les valeurs propres de $u\otimes1$ (dans une clôture
algébrique de $K$) sont des éléments \emph{entiers} sur $A$, et il en est donc
de même des $\sigma_i(u)$.

\noindent\textbf{$(\mathbf{Err}_{\mathrm{III}},\,22)$} Dans
$(0_{\mathrm{III}},\,9.1.1)$, ligne 18 du bas de la p. 12, supprimer
« ouverte ».

\noindent\textbf{$(\mathbf{Err}_{\mathrm{III}},\,23)$} Dans
$(0_{\mathrm{III}},\,11.7.3)$, ligne 10 de la p. 43, après $\mathsf C''$,
ajouter : tel que pour tout objet projectif $P$ de $\mathsf C'$ (resp. tout
objet projectif $P'$ de $\mathsf C''$) le foncteur
$A'\rightsquigarrow T(P,A')$ (resp. $A\rightsquigarrow T(A,P')$) soit exact
dans $\mathsf C'$ (resp. $\mathsf C$).

\noindent\textbf{$(\mathbf{Err}_{\mathrm{III}},\,24)$} L'énoncé de
$(0_{\mathrm{III}},\,13.7.7)$ est inexact, et doit être modifié comme suit :
ligne 6 de la p. 78, supprimer « et
$(\mathrm R^{n+1}T(A_k))_{k\in\mathbf Z}$ » ; ligne 7 de la p. 78, remplacer
« $\mathrm R'^nT(A)$ est un $S$-module de type fini » par
« $\mathrm R'^nT(A)$ et $\mathrm R'^{n+1}T(A)$ sont des $S$-modules de type
fini » ; lignes 12--13 de la p. 78, supprimer « et $p+q=n+1$ ». Dans la
démonstration, ligne 23 de la p. 78, supprimer « et pour $n+1$ ».

\noindent\textbf{$(\mathbf{Err}_{\mathrm{III}},\,25)$} Dans (III, 1.4.15),
ligne 6 du bas de la p. 92, remplacer « \emph{de type fini} » par
« \emph{quasi-compact} ».

\noindent\textbf{$(\mathbf{Err}_{\mathrm{III}},\,26)$} Dans (III, 2.2.4),
ligne 1 de la p. 102, remplacer « \emph{que les supports de $\mathscr F$ et de
$\mathscr H$ soient propres sur $Y$} » par « \emph{que le support de
$\operatorname{Im}(\mathscr F\to\mathscr G)$ soit propre sur $Y$} ». Dans la
démonstration, remplacer le texte des lignes 10 à 16, depuis « Cela
étant\ldots{} », par :

\oldpage[III]{222}
Posons
$\mathscr G_1=\operatorname{Im}(\mathscr F\to\mathscr G)
=\operatorname{Ker}(\mathscr G\to\mathscr H)$, qui est cohérent. Comme on a la
suite exacte
$0\to\mathscr G_1(n)\to\mathscr G(n)\to\mathscr H(n)$ et que le foncteur
$f_*$ est exact à gauche, la suite
$0\to f_*(\mathscr G_1(n))\to f_*(\mathscr G(n))\to f_*(\mathscr H(n))$ est
exacte pour tout $n$, et il suffit donc de montrer que, pour $n$ assez grand,
la suite $f_*(\mathscr F(n))\to f_*(\mathscr G_1(n))\to0$ est exacte. Autrement
dit, on peut se borner au cas où $\mathscr H=0$ et où le support de
$\mathscr G$ est propre sur $Y$, donc \emph{fermé dans $Z$} (II, 5.4.10) ;
$\mathscr G'=i_*(\mathscr G)$ est par suite un $\mathscr O_Z$-Module cohérent tel
que $\mathscr G'|X=\mathscr G$. On sait (I, 9.4.3) qu'il existe un
$\mathscr O_Z$-Module cohérent $\mathscr F'$ tel que
$\mathscr F'|X=\mathscr F$. En outre, comme $X$ est ouvert dans $Z$ et
$\operatorname{Supp}(\mathscr G)\subset X$ fermé dans $Z$, il est immédiat que
l'on définit un homomorphisme surjectif $u':\mathscr F'\to\mathscr G'$ de
faisceaux tel que $u'|X$ soit l'homomorphisme surjectif donné
$u:\mathscr F\to\mathscr G$, en prenant $u'|U=0$ pour tout ouvert $U$ de $Z$
ne rencontrant pas $\operatorname{Supp}(\mathscr G)$ et $u'|U=u|U$ pour tout
ouvert $U\subset X$. Cela étant, on voit comme au début de la démonstration de
(2.2.2) que l'on peut se borner au cas où $Y$ est affine, et il s'agit donc de
montrer que, pour $n$ assez grand, l'homomorphisme
$\Gamma(u):\Gamma(X,\mathscr F(n))\to\Gamma(X,\mathscr G(n))$ est surjectif. Or,
on a le diagramme commutatif
\[
\xymatrix@C=35mm@R=12mm{
\Gamma(Z,\mathscr F'(n)) \ar[r]^{\Gamma(u')} \ar[d]_v &
\Gamma(Z,\mathscr G'(n)) \ar[d]^w \\
\Gamma(X,\mathscr F(n)) \ar[r]_{\Gamma(u)} & \Gamma(X,\mathscr G(n))
}
\]
où $v$ et $w$ sont les homomorphismes de restriction. La définition de
$\mathscr G'$ montre en outre que $w$ est \emph{bijectif} ; par ailleurs, en
vertu de (2.2.3), $\Gamma(u')$ est \emph{surjectif} pour $n$ assez grand, donc
il en est de même de $\Gamma(u)$.

\noindent\textbf{$(\mathbf{Err}_{\mathrm{III}},\,27)$} Dans (III, 2.2.5),
après la ligne 7 du bas de la p. 102, ajouter :

(iii) Sous les hypothèses de (2.2.4) concernant $X$, $Y$, $f$ et $\mathscr L$,
il est immédiat que si $\mathscr H$ est un $\mathscr O_X$-Module cohérent dont
le support est \emph{propre sur $Y$}, un raisonnement analogue à celui de
(2.2.4) montre qu'il existe un entier $N$ tel que pour $n\geq N$, on ait
$\mathrm R^if_*(\mathscr H(n))=0$ pour tout $i>0$. On en conclut, par la suite
exacte de cohomologie, que si $u:\mathscr F\to\mathscr G$ est un homomorphisme
de $\mathscr O_X$-Modules cohérents tel que $\operatorname{Ker}(u)$ et
$\operatorname{Coker}(u)$ aient leurs supports \emph{propres sur $Y$}, alors il
existe $N$ tel que pour $n\geq N$, l'homomorphisme correspondant
$\mathrm R^if_*(\mathscr F(n))\to\mathrm R^if_*(\mathscr G(n))$ soit
\emph{bijectif} pour tout $i>0$.

\noindent\textbf{$(\mathbf{Err}_{\mathrm{III}},\,28)$} Dans (III, 4.4.9),
ligne 17 du bas de la p. 137, remplacer « \emph{Soient $Y$ un préschéma
intègre localement noethérien} » par « \emph{Soient $X$ et $Y$ deux
préschémas intègres localement noethériens} » ; ligne 16 du bas de la p.
137, remplacer « \emph{de type fini} » par « \emph{localement de type
fini} ». La démonstration est essentiellement inchangée,
$f^{-1}(y)$ étant localement de type fini sur $k(y)$, donc encore discret.

\noindent\textbf{$(\mathbf{Err}_{\mathrm{III}},\,29)$} dans (I, 9.3.4), ligne
19 du bas de la p. 173, remplacer « \emph{noethérien} » par
« \emph{quasi-compact} » ; lignes 18 et 19 du bas de la p. 173, remplacer
« \emph{cohérent} » par « \emph{quasi-cohérent de type fini} ». Dans la
démonstration, on se ramène au cas où

\oldpage[III]{223}
$X=\operatorname{Spec}(A)$, $\mathscr F=\widetilde M$, où $M$ est un $A$-module
de type fini, $\mathscr J=\widetilde{\mathfrak J}$, où $\mathfrak J$ est un
idéal de type fini de $A$ ; le reste du raisonnement est alors inchangé.

\noindent\textbf{$(\mathbf{Err}_{\mathrm{III}},\,30)$} Dans (I, 9.3.5),
remplacer les lignes 1 à 7 du bas de la p. 173 par le texte suivant :

\emph{Proposition (9.3.5). --- Soient $X$ un préschéma, $\mathscr F$ un
$\mathscr O_X$-Module quasi-cohérent de type fini. Alors il existe un
sous-préschéma fermé $Y$ de $X$, dont l'espace sous-jacent est égal à
$\operatorname{Supp}(\mathscr F)$, et un $\mathscr O_Y$-Module quasi-cohérent de
type fini $\mathscr G$ tels que, si $j:Y\to X$ est l'injection canonique,
$\mathscr F$ soit isomorphe à $j_*(\mathscr G)$.}

Il suffira de montrer que l'Idéal $\mathscr J$ de $\mathscr O_X$,
\emph{annulateur de $\mathscr F$}, est \emph{quasi-cohérent} ; on prendra alors
pour $Y$ le sous-préschéma fermé de $X$ défini par $\mathscr J$ (I, 4.1.2), et
comme $\mathscr J\mathscr F=0$, $\mathscr F$ est un
$(\mathscr O_X/\mathscr J)$-Module et on répondra à la question en prenant
$\mathscr G=j^*(\mathscr F)$. Pour voir que $\mathscr J$ est quasi-cohérent, on
peut (la question étant locale) se borner au cas où
$X=\operatorname{Spec}(A)$, $\mathscr F=\widetilde M$, où $M$ est un $A$-module
engendré par un nombre fini d'éléments $x_i\ (1\leq i\leq r)$ ; l'Idéal
$\mathscr J$ est alors l'intersection des annulateurs des $x_i$. Mais
l'annulateur de $x_i$ est le noyau de l'homomorphisme
$\mathscr O_X\to\mathscr F$ correspondant à l'homomorphisme $s\mapsto sx_i$ de
$A$ dans $M$ ; c'est donc bien un Idéal quasi-cohérent (I, 4.1.1) et toute
intersection finie de tels Idéaux est aussi un Idéal quasi-cohérent (I,
1.3.10).

% FIN SOURCE ega3/ega3-7-fr.tex

\end{document}
% EGA-COMPONENT-END BASE

% EGA-COMPONENT-BEGIN IV1 244250DFD67FE0C9E0D84F23EF7579F0B8871BC6FEC491B312CE0D6561D6D916
\documentclass[10pt,oneside]{book}

\usepackage[T1]{fontenc}
\usepackage[utf8]{inputenc}
\usepackage[french]{babel}
\usepackage{amsmath,amssymb,mathrsfs}
\usepackage[all]{xy}
\usepackage[a4paper,margin=22mm]{geometry}
\usepackage[
  colorlinks=true,
  linkcolor=blue,
  hypertexnames=false,
  unicode=true,
  pdfencoding=auto
]{hyperref}
\hypersetup{
  pdftitle={Éléments de géométrie algébrique IV, première partie -- édition française},
  pdfauthor={Alexander Grothendieck et Jean Dieudonné},
  pdfsubject={Transcription diplomatique française depuis l'autorité NUMDAM}
}

\renewcommand{\thesection}{\arabic{section}}
\renewcommand{\thesubsection}{\thesection.\arabic{subsection}}
\newcommand{\gr}{\operatorname{gr}}
\newcommand{\sref}[1]{(\ref{#1})}
\makeatletter
\newcommand{\setmanualcurrentlabel}[1]{\def\@currentlabel{#1}}
\makeatother
\newcommand{\oldpage}[2][]{%
  \setcounter{footnote}{0}%
  \marginpar{\scriptsize\textbf{#1}~|~#2}\ignorespaces}
\newenvironment{env}[1][]{\setmanualcurrentlabel{#1}\par\medskip\noindent\textbf{(#1)}\ }{\par}
\newenvironment{definition}[1][]{%
  \setmanualcurrentlabel{#1}\par\medskip\noindent\itshape Définition (#1). ---\ }{\par}
\newenvironment{proposition}[1][]{%
  \setmanualcurrentlabel{#1}\par\medskip\noindent\itshape Proposition (#1). ---\ }{\par}
\newenvironment{theorem}[1][]{%
  \setmanualcurrentlabel{#1}\par\medskip\noindent\itshape Théorème (#1). ---\ }{\par}
\newenvironment{corollary}[1][]{%
  \setmanualcurrentlabel{#1}\par\medskip\noindent\itshape Corollaire (#1). ---\ }{\par}
\newenvironment{remark}[1][]{%
  \setmanualcurrentlabel{#1}\par\medskip\noindent\itshape Remarque (#1). ---\ }{\par}
\newenvironment{remarks}[1][]{%
  \setmanualcurrentlabel{#1}\par\medskip\noindent\itshape Remarques (#1). ---\ }{\par}
\newenvironment{lemma}[1][]{%
  \setmanualcurrentlabel{#1}\par\medskip\noindent\itshape Lemme (#1). ---\ }{\par}
\newenvironment{example}[1][]{%
  \setmanualcurrentlabel{#1}\par\medskip\noindent\textit{Exemple (#1). ---}\ }{\par}
\newenvironment{examples}[1][]{%
  \setmanualcurrentlabel{#1}\par\medskip\noindent\textit{Exemples (#1). ---}\ }{\par}
\newenvironment{proof}{%
  \par\medskip\noindent\textit{Démonstration. ---\ }}{\par}

\makeatletter
\newcommand{\EGAReaderTarget}[2]{\hypertarget{#1}{}\@bsphack\begingroup\def\@currentHref{#1}\label{#2}\endgroup\@esphack}
\makeatother
\begin{document}
% BEGIN RIGHTS AND AI TRANSCRIPTION NOTICE
\clearpage
\begingroup
\small
\noindent\textbf{Droits, disponibilité et provenance}

Les droits sur les œuvres originales et sur toute contribution protégée à ces traductions et transcriptions dérivées restent ceux de leurs titulaires respectifs. Ce projet ne place pas ces œuvres ou contributions dans le domaine public et n'accorde aucune licence sur elles. Les mentions de droits de tiers et les autorisations déjà accordées restent applicables.

Ces éditions sont mises à disposition gratuitement, à des fins non commerciales et exclusivement pédagogiques, dans la présomption que les titulaires de droits ne s'opposeraient pas à la libre disponibilité de ces traductions et transcriptions aux côtés des originaux français sur NUMDAM (www.numdam.org/journals/PMIHES/). Cette présomption ne signifie pas qu'une autorisation a été obtenue ; le libre accès aux originaux n'autorise pas, à lui seul, leur traduction, leur adaptation ou leur republication.

Si vous détenez des droits sur un élément présenté ici et vous opposez à sa disponibilité, veuillez ouvrir une demande sur GitHub (github.com/KokunoYumeto/ega-fr/issues). Le responsable de la publication retirera l'élément concerné à la demande du titulaire des droits.

La transcription moderne, les corrections, la composition et les notes mathématiques sont produites par IA. Les travaux documentés comprennent ceux d'OpenAI Codex, GPT-6 Astra, effort Ultra. Les indications de provenance des contributions antérieures restent applicables. Les textes originaux restent l'œuvre d'Alexandre Grothendieck et Jean Dieudonné ; ils ne sont pas attribués à l'IA. Cette édition n'est approuvée ni par les auteurs, ni par l'IHÉS, ni par NUMDAM.
\par
\endgroup
\clearpage
% END RIGHTS AND AI TRANSCRIPTION NOTICE

%
% DEBUT SOURCE ega4/ega0-14-fr.tex
\clearpage
\thispagestyle{plain}
\markboth{PRÉLIMINAIRES}{PRÉLIMINAIRES}
\oldpage[IV-1]{5}
\phantomsection
\label{chapter:0IV1-fr}

\begin{center}
{\large CHAPITRE 0 (suite)\footnotemark\par}
\vspace{1.5em}
{\Large\bfseries PRÉLIMINAIRES\par}
\vspace{2em}
{\large\bfseries Sommaire\par}
\end{center}

\footnotetext{On rappelle que le symbole $0_N$ (resp. $\operatorname{Err}_N$)
réfère à la partie du chapitre 0 (resp. la liste d'Errata) jointe au chapitre
$N$.}

\begin{center}
\begin{tabular}{@{}r l@{}}
\S~14. & Dimension combinatoire d'un espace topologique.\\
\S~15. & Suites M-régulières et suites $\mathscr F$-régulières.\\
\S~16. & Dimension et profondeur dans les anneaux locaux noethériens.\\
\S~17. & Anneaux réguliers.\\
\S~18. & Compléments sur les extensions d'algèbres.\\
\S~19. & Algèbres formellement lisses et anneaux de Cohen.\\
\S~20. & Dérivations et différentielles.\\
\S~21. & Différentielles dans les anneaux de caractéristique $p$.\\
\S~22. & Critères différentiels de lissité formelle et de régularité.\\
\S~23. & Anneaux japonais.
\end{tabular}
\end{center}

La presque totalité des paragraphes précédents est consacrée à l'exposé de
notions d'Algèbre commutative qui seront utilisées au cours du chapitre IV.
Bien qu'une bonne partie de ces notions figure déjà dans plusieurs ouvrages
([1], [12], [13], [17], [30]), on a pensé qu'il serait plus commode pour le
lecteur d'en avoir une exposition suivie et à peu près indépendante. Joints
aux \S\S~5, 6 et 7 du chapitre IV (où est utilisé le langage des schémas), ces
paragraphes constituent, à l'intérieur de notre Traité, un petit Traité
spécial, à peu près indépendant des chapitres I à III, et qui vise à présenter
sous une forme cohérente les propriétés des anneaux qui «~se comportent bien~»
relativement à des opérations telles que la complétion ou la fermeture
intégrale, en rattachant systématiquement ces propriétés à quelques
conceptions générales\footnote{La plupart des propriétés dont il s'agit ont
été découvertes par Chevalley, Zariski, Nagata et Serre. La méthode utilisée
ici a d'abord été développée en automne 1961, dans un cours professé à Harvard
University, par A.~Grothendieck.}.

\setcounter{section}{13}
\section{Dimension combinatoire d'un espace topologique}
\label{section:0.14-fr}

\setcounter{subsection}{0}
\subsection{Dimension combinatoire d'un espace topologique}
\label{subsection:0.14.1-fr}

\oldpage[0\textsubscript{IV-1}]{6}
\begin{env}[14.1.1]
\label{0.14.1.1-fr}
Soit $I$ un ensemble ordonné~; une \emph{chaîne} d'éléments de $I$ est par
définition une suite finie strictement croissante
$i_0<i_1<\cdots<i_n$ d'éléments de $I$ ($n\geq 0$)~; par définition la
\emph{longueur} de cette chaîne est $n$. Si $X$ est un espace topologique,
l'ensemble de ses parties \emph{fermées irréductibles} est ordonné par
inclusion, d'où la notion de \emph{chaîne} de parties fermées irréductibles de
$X$.
\end{env}

\begin{definition}[14.1.2]
\label{0.14.1.2-fr}
Soit $X$ un espace topologique. On appelle \emph{dimension combinatoire de
$X$} (ou simplement \emph{dimension de $X$} s'il ne peut y avoir de
confusion) et on note $\dim_c(X)$, ou simplement $\dim(X)$, la borne
supérieure des longueurs des chaînes de parties fermées irréductibles de $X$.
Pour tout $x\in X$, on appelle \emph{dimension combinatoire de $X$ en $x$}
(ou simplement \emph{dimension de $X$ en $x$}) et on note $\dim_x(X)$ le
nombre $\inf_U(\dim(U))$, où $U$ parcourt l'ensemble des voisinages ouverts de
$x$ dans $X$.
\end{definition}

Il résulte de cette définition que l'on a
\[
  \dim(\varnothing)=-\infty
\]
(la borne supérieure dans $\overline{\mathbb R}$ de la partie vide étant
$-\infty$). Si $(X_\alpha)$ est la famille des composantes irréductibles de
$X$, on a
\[
  \dim(X)=\sup_\alpha\dim(X_\alpha)
  \tag{14.1.2.1}\label{0.14.1.2.1-fr}
\]
car toute chaîne de parties fermées irréductibles de $X$ est par définition
contenue dans une composante irréductible de $X$, et inversement ces
composantes sont fermées dans $X$, donc toute partie fermée irréductible d'un
$X_\alpha$ est une partie fermée irréductible de $X$.

\begin{definition}[14.1.3]
\label{0.14.1.3-fr}
On dit qu'un espace topologique $X$ est \emph{équidimensionnel} si toutes ses
composantes irréductibles ont même dimension (égale alors à $\dim(X)$ par
(14.1.2.1)).
\end{definition}

\begin{proposition}[14.1.4]
\label{0.14.1.4-fr}
\begin{enumerate}
  \item[\rm(i)] Pour toute partie $Y$ d'un espace topologique $X$, on a
  $\dim(Y)\leq\dim(X)$.
  \item[\rm(ii)] Si un espace topologique $X$ est réunion finie de parties
  fermées $X_i$, on a $\dim(X)=\sup_i\dim(X_i)$.
\end{enumerate}
\end{proposition}

\begin{proof}
Pour toute partie fermée irréductible $Z$ de $Y$, l'adhérence $\overline Z$ de
$Z$ dans $X$ est irréductible ($0_{\mathrm I}$, 2.1.2) et
$\overline Z\cap Y=Z$, d'où (i). D'autre part, si
$X=\bigcup_{i=1}^n X_i$, où les $X_i$ sont fermés, toute partie fermée
irréductible de $X$ est contenue dans un des $X_i$ ($0_{\mathrm I}$, 2.1.1),
et par suite toute chaîne de parties fermées irréductibles de $X$ est contenue
dans un des $X_i$~; d'où (ii).
\end{proof}

On déduit de (14.1.4, (i)) que pour tout $x\in X$, on peut aussi écrire
\[
  \dim_x(X)=\lim_U\dim(U)
  \tag{14.1.4.1}\label{0.14.1.4.1-fr}
\]
la limite étant prise suivant l'ensemble filtrant décroissant des voisinages
ouverts de $x$ dans $X$.

\oldpage[0\textsubscript{IV-1}]{7}
\begin{corollary}[14.1.5]
\label{0.14.1.5-fr}
Soient $X$ un espace topologique, $x$ un point de $X$, $U$ un voisinage de
$x$, $Y_i$ ($1\leq i\leq n$) des parties fermées de $U$ telles que
$x\in Y_i$ pour tout $i$ et que $U$ soit réunion des $Y_i$. Alors on a
\[
  \dim_x(X)=\sup_i\bigl(\dim_x(Y_i)\bigr).
  \tag{14.1.5.1}\label{0.14.1.5.1-fr}
\]
\end{corollary}

\begin{proof}
Il résulte de (14.1.4, (ii)) que l'on a
$\dim_x(X)=\inf_V(\sup_i(\dim(Y_i\cap V)))$ lorsque $V$ parcourt l'ensemble
des voisinages ouverts de $x$ contenus dans $U$~; de même on a
$\dim_x(Y_i)=\inf_V(\dim(Y_i\cap V))$ pour tout $i$. Le corollaire est donc
évident si
\[
  \sup_i\bigl(\dim_x(Y_i)\bigr)=+\infty~;
\]
dans le cas contraire, il y a un voisinage ouvert $V_0\subset U$ de $x$ tel
que $\dim(Y_i\cap V)=\dim_x(Y_i)$ pour $1\leq i\leq n$ et pour tout
$V\subset V_0$, d'où la conclusion.
\end{proof}

\begin{proposition}[14.1.6]
\label{0.14.1.6-fr}
Pour tout espace topologique $X$, on a
$\dim(X)=\sup_{x\in X}\dim_x(X)$.
\end{proposition}

\begin{proof}
Il résulte de la définition (14.1.2) et de la prop. (14.1.4) que
$\dim_x(X)\leq\dim(X)$ pour tout $x\in X$. D'autre part, soit
$Z_0\subset Z_1\subset\cdots\subset Z_n$ une chaîne de parties fermées
irréductibles de $X$, et soit $x\in Z_0$~; pour tout ouvert $U\subset X$
contenant $x$, $U\cap Z_i$ est irréductible ($0_{\mathrm I}$, 2.1.6) et
fermé dans $U$, et comme on a $\overline{U\cap Z_i}=Z_i$ dans $X$, les
$U\cap Z_i$ sont distincts deux à deux~; donc on a $\dim(U)\geq n$, ce qui
achève la démonstration.
\end{proof}

\begin{corollary}[14.1.7]
\label{0.14.1.7-fr}
Si $(X_\alpha)$ est un recouvrement ouvert de $X$, ou un recouvrement fermé
localement fini de $X$, on a
$\dim(X)=\sup_\alpha(\dim(X_\alpha))$.
\end{corollary}

\begin{proof}
Si $X_\alpha$ est un voisinage de $x\in X$, on a
$\dim_x(X)\leq\dim(X_\alpha)$, d'où la première assertion. D'autre part, si
les $X_\alpha$ sont fermés et si $U$ est un voisinage de $x\in X$ qui ne
rencontre qu'un nombre fini d'ensembles $X_\alpha$, on a
\[
  \dim_x(X)\leq\dim(U)=\sup_\alpha\bigl(\dim(U\cap X_\alpha)\bigr)
  \leq\sup_\alpha\bigl(\dim(X_\alpha)\bigr)
\]
d'après (14.1.4), d'où la seconde assertion.
\end{proof}

\begin{corollary}[14.1.8]
\label{0.14.1.8-fr}
Soit $X$ un espace de Kolmogoroff ($0_{\mathrm I}$, 2.1.3) noethérien, et
soit $F$ l'ensemble des points fermés de $X$. Alors
$\dim(X)=\sup_{x\in F}\dim_x(X)$.
\end{corollary}

\begin{proof}
Avec les notations de la démonstration de (14.1.6), il suffit de remarquer
qu'il existe dans $Z_0$ un point fermé ($0_{\mathrm I}$, 2.1.3).
\end{proof}

\begin{proposition}[14.1.9]
\label{0.14.1.9-fr}
Soit $X$ un espace de Kolmogoroff noethérien non vide. Pour que
$\dim(X)=0$, il faut et il suffit que $X$ soit fini et discret.
\end{proposition}

\begin{proof}
Si un espace $X$ est séparé (et \emph{a fortiori} si $X$ est un espace
discret), les seuls ensembles fermés irréductibles de $X$ sont réduits à un
point, donc $\dim(X)=0$. Inversement, supposons que $X$ soit un espace de
Kolmogoroff noethérien tel que $\dim(X)=0$~; comme toute composante
irréductible de $X$ contient un point fermé ($0_{\mathrm I}$, 2.1.3), elle
doit être réduite à ce point. Comme $X$ n'a qu'un nombre fini de composantes
irréductibles, il est fini et discret.
\end{proof}

\begin{corollary}[14.1.10]
\label{0.14.1.10-fr}
Soit $X$ un espace de Kolmogoroff noethérien. Pour qu'un point $x\in X$ soit
isolé, il faut et il suffit que $\dim_x(X)=0$.
\end{corollary}

\begin{proof}
La condition est évidemment nécessaire (sans hypothèse sur $X$). Elle est
suffisante,
\oldpage[0\textsubscript{IV-1}]{8}
car il en résulte que $\dim(U)=0$ pour un voisinage ouvert $U$ de $x$, et
comme $U$ est un espace de Kolmogoroff noethérien, $U$ est fini et discret.
\end{proof}

\begin{proposition}[14.1.11]
\label{0.14.1.11-fr}
La fonction $x\mapsto\dim_x(X)$ est semi-continue supérieurement dans $X$.
\end{proposition}

\begin{proof}
Il est clair que cette fonction est semi-continue supérieurement en tout
point où elle vaut $+\infty$. Supposons donc $\dim_x(X)=n<+\infty$~; alors,
la formule (14.1.4.1) montre qu'il existe un voisinage ouvert $U_0$ de $x$ tel
que $\dim(U)=n$ pour tout voisinage ouvert $U\subset U_0$ de $x$. Cela étant,
pour tout $y\in U_0$ et tout voisinage ouvert $V\subset U_0$ de $y$, on a
$\dim(V)\leq\dim(U_0)=n$ (14.1.4)~; on déduit donc de (14.1.4.1) que
$\dim_y(X)\leq n$.
\end{proof}

\begin{remark}[14.1.12]
\label{0.14.1.12-fr}
Si $X$, $Y$ sont deux espaces topologiques, et $f:X\to Y$ une application
continue, on peut avoir $\dim(f(X))>\dim(X)$~; on en a un exemple en prenant
pour $X$ un espace discret à 2 éléments $a$, $b$, pour $Y$ l'ensemble
$\{a,b\}$ muni de la topologie pour laquelle les ensembles fermés sont
$\varnothing$, $\{a\}$ et $\{a,b\}$~; si $f:X\to Y$ est l'application
identique, on a $\dim(Y)=1$ et $\dim(X)=0$. On notera que $Y$ est le spectre
d'un anneau de valuation discrète $A$, dont $a$ est l'unique point fermé et
$b$ le point générique~; si $K$ et $k$ sont le corps des fractions et le
corps résiduel de $A$, $X$ est le spectre de l'anneau $k\times K$ et $f$
l'application continue correspondant à l'homomorphisme
$(\varphi,\psi):A\to k\times K$, où $\varphi:A\to k$ et $\psi:A\to K$ sont
les homomorphismes canoniques (cf. IV, 5.4.3).
\end{remark}

\subsection{Codimension d'une partie fermée}
\label{subsection:0.14.2-fr}

\begin{definition}[14.2.1]
\label{0.14.2.1-fr}
Étant donnée une partie fermée irréductible $Y$ d'un espace topologique $X$,
on appelle \emph{codimension combinatoire} (ou simplement \emph{codimension})
de $Y$ dans $X$, et on note $\operatorname{codim}(Y,X)$ la borne supérieure
des longueurs des chaînes de parties fermées irréductibles de $X$, dont $Y$
est le plus petit élément. Si $Y$ est une partie fermée quelconque de $X$, on
appelle \emph{codimension de $Y$ dans $X$}, et on note encore
$\operatorname{codim}(Y,X)$ la borne inférieure des codimensions dans $X$ des
composantes irréductibles de $Y$. On dit que $X$ est
\emph{équicodimensionnel} si toutes les parties fermées irréductibles minimales
de $X$ ont même codimension dans $X$.
\end{definition}

Il résulte de cette définition que
$\operatorname{codim}(\varnothing,X)=+\infty$, la borne inférieure de la
partie vide de $\overline{\mathbb R}$ étant $+\infty$. Si $Y$ est fermé dans
$X$ et si $(X_\alpha)$ (resp. $(Y_\beta)$) est la famille des composantes
irréductibles de $X$ (resp. $Y$), tout $Y_\beta$ est contenu dans un
$X_\alpha$, et plus généralement toute chaîne de parties fermées
irréductibles de $X$ dont $Y_\beta$ est le plus petit élément est formée de
parties d'un $X_\alpha$~; on a donc
\[
  \operatorname{codim}(Y,X)=
  \inf_\beta\left(\sup_\alpha
  \bigl(\operatorname{codim}(Y_\beta,X_\alpha)\bigr)\right)
  \tag{14.2.1.1}\label{0.14.2.1.1-fr}
\]
où pour chaque $\beta$, $\alpha$ parcourt l'ensemble des indices tels que
$Y_\beta\subset X_\alpha$.

\begin{proposition}[14.2.2]
\label{0.14.2.2-fr}
Soit $X$ un espace topologique.
\begin{enumerate}
  \item[\rm(i)] Si $\Phi$ est l'ensemble des parties fermées irréductibles de
  $X$, on a
  \[
    \dim(X)=\sup_{Y\in\Phi}\bigl(\operatorname{codim}(Y,X)\bigr).
    \tag{14.2.2.1}\label{0.14.2.2.1-fr}
  \]
  \oldpage[0\textsubscript{IV-1}]{9}
  \item[\rm(ii)] Pour toute partie fermée non vide $Y$ de $X$, on a
  \[
    \dim(Y)+\operatorname{codim}(Y,X)\leq\dim(X).
    \tag{14.2.2.2}\label{0.14.2.2.2-fr}
  \]
  \item[\rm(iii)] Si $Y$, $Z$, $T$ sont trois parties fermées de $X$ telles
  que $Y\subset Z\subset T$, on a
  \[
    \operatorname{codim}(Y,Z)+\operatorname{codim}(Z,T)
    \leq\operatorname{codim}(Y,T).
    \tag{14.2.2.3}\label{0.14.2.2.3-fr}
  \]
  \item[\rm(iv)] Pour qu'une partie fermée $Y$ de $X$ soit telle que
  $\operatorname{codim}(Y,X)=0$, il faut et il suffit que $Y$ contienne une
  composante irréductible de $X$.
\end{enumerate}
\end{proposition}

\begin{proof}
Les assertions (i) et (iv) sont des conséquences immédiates de la définition
(14.2.1). Pour démontrer (ii), on peut se borner au cas où $Y$ est
irréductible, et alors la formule résulte des définitions (14.1.1) et
(14.2.1). Enfin, pour démontrer (iii), on peut, en vertu de la définition
(14.2.1), se borner d'abord au cas où $Y$ est irréductible~; alors
$\operatorname{codim}(Y,Z)=
\sup_\alpha(\operatorname{codim}(Y,Z_\alpha))$ pour les composantes
irréductibles $Z_\alpha$ de $Z$ contenant $Y$~; il est clair que
$\operatorname{codim}(Y,T)\geq\operatorname{codim}(Y,Z)$, donc l'inégalité
est vraie si $\operatorname{codim}(Y,Z)=+\infty$~; sinon, il existe un
$\alpha$ tel que
$\operatorname{codim}(Y,Z)=\operatorname{codim}(Y,Z_\alpha)$ et en vertu de
(14.2.1), on peut se borner au cas où $Z$ est lui aussi irréductible~; mais
alors l'inégalité (14.2.2.3) est conséquence évidente de la définition
(14.2.1).
\end{proof}

\begin{proposition}[14.2.3]
\label{0.14.2.3-fr}
Soient $X$ un espace topologique, $Y$ une partie fermée de $X$. Pour tout
ouvert $U$ dans $X$, on a
\[
  \operatorname{codim}(Y\cap U,U)\geq\operatorname{codim}(Y,X).
  \tag{14.2.3.1}\label{0.14.2.3.1-fr}
\]
En outre, pour que les deux membres de (14.2.3.1) soient égaux, il faut et il
suffit que, si $(Y_\alpha)$ est la famille des composantes irréductibles de
$Y$ rencontrant $U$, on ait
$\operatorname{codim}(Y,X)=
\inf_\alpha(\operatorname{codim}(Y_\alpha,X))$.
\end{proposition}

\begin{proof}
On sait ($0_{\mathrm I}$, 2.1.6) que $Z\mapsto\overline Z$ est une bijection
de l'ensemble des parties fermées irréductibles de $U$ sur l'ensemble des
parties fermées irréductibles de $X$ rencontrant $U$, et en particulier fait
correspondre aux composantes irréductibles de $Y\cap U$ les composantes
irréductibles de $Y$ rencontrant $U$~; si $Y_\alpha$ est une de ces
dernières, on a donc
$\operatorname{codim}(Y_\alpha,X)=
\operatorname{codim}(Y_\alpha\cap U,U)$, et la proposition résulte alors de
la définition (14.2.1).
\end{proof}

\begin{definition}[14.2.4]
\label{0.14.2.4-fr}
Soient $X$ un espace topologique, $Y$ une partie fermée de $X$, $x$ un point
de $X$. On appelle \emph{codimension de $Y$ dans $X$ au point $x$} et on note
$\operatorname{codim}_x(Y,X)$ le nombre
$\sup_U(\operatorname{codim}(Y\cap U,U))$, où $U$ parcourt l'ensemble des
voisinages ouverts de $x$ dans $X$.
\end{definition}

En vertu de (14.2.3), on peut aussi écrire
\[
  \operatorname{codim}_x(Y,X)=
  \lim_U\bigl(\operatorname{codim}(Y\cap U,U)\bigr)
  \tag{14.2.4.1}\label{0.14.2.4.1-fr}
\]
la limite étant prise suivant l'ensemble filtrant décroissant des voisinages
ouverts de $x$ dans $X$. On notera que l'on a
\[
  \operatorname{codim}_x(Y,X)=+\infty\quad\text{si}\quad x\in X-Y.
\]

\oldpage[0\textsubscript{IV-1}]{10}
\begin{proposition}[14.2.5]
\label{0.14.2.5-fr}
Si $(Y_i)_{1\leq i\leq n}$ est une famille finie de parties fermées d'un
espace topologique $X$, et $Y$ la réunion de cette famille, on a
\[
  \operatorname{codim}(Y,X)=
  \inf_i\bigl(\operatorname{codim}(Y_i,X)\bigr).
  \tag{14.2.5.1}\label{0.14.2.5.1-fr}
\]
\end{proposition}

\begin{proof}
En effet, toute composante irréductible de l'un des $Y_i$ est contenue dans
une composante irréductible de $Y$, et inversement toute composante
irréductible de $Y$ est aussi composante irréductible de l'un des $Y_i$
($0_{\mathrm I}$, 2.1.1)~; la conclusion résulte donc de la définition
(14.2.1) et de l'inégalité (14.2.2.3).
\end{proof}

\begin{corollary}[14.2.6]
\label{0.14.2.6-fr}
Soient $X$ un espace topologique, $Y$ un sous-espace fermé localement
noethérien de $X$.
\begin{enumerate}
  \item[\rm(i)] Pour tout $x\in X$, il n'existe qu'un nombre fini de
  composantes irréductibles $Y_i$ ($1\leq i\leq n$) de $Y$ contenant $x$,
  et l'on a
  $\operatorname{codim}_x(Y,X)=
  \inf_i(\operatorname{codim}(Y_i,X))$.
  \item[\rm(ii)] La fonction
  $x\mapsto\operatorname{codim}_x(Y,X)$ est semi-continue inférieurement
  dans $X$.
\end{enumerate}
\end{corollary}

\begin{proof}
En effet, par hypothèse il y a un voisinage ouvert $U_0$ de $x$ dans $X$ tel
que $Y\cap U_0$ soit noethérien, donc n'ait qu'un nombre fini de composantes
irréductibles, qui sont traces sur $U_0$ de composantes irréductibles de $Y$~;
\emph{a fortiori} il n'y a qu'un nombre fini de composantes irréductibles
$Y_i$ ($1\leq i\leq n$) de $Y$ contenant $x$, et on peut, en remplaçant $U_0$
par un voisinage ouvert $U\subset U_0$ de $x$ ne rencontrant aucun des $Y_j$
qui ne contiennent pas $x$, supposer que les $Y_i\cap U$ sont les composantes
irréductibles de $Y\cap U$~; pour tout voisinage ouvert $V\subset U$ de $x$
dans $X$, les $Y_i\cap V$ sont alors les composantes irréductibles de
$Y\cap V$, et (14.2.3) montre alors que
$\operatorname{codim}(Y_i,X)=\operatorname{codim}(Y_i\cap V,V)$, ce qui
prouve (i). En outre, pour tout $x'\in U$, les composantes irréductibles de
$Y$ contenant $x'$ sont certaines des $Y_i$, donc
$\operatorname{codim}_{x'}(Y,X)\geq\operatorname{codim}_x(Y,X)$, ce qui
démontre (ii).
\end{proof}

\subsection{La condition des chaînes}
\label{subsection:0.14.3-fr}

\begin{env}[14.3.1]
\label{0.14.3.1-fr}
Dans un espace topologique $X$, nous dirons qu'une chaîne
$Z_0\subset Z_1\subset\cdots\subset Z_n$ de parties fermées irréductibles est
\emph{saturée} s'il n'existe pas de partie fermée irréductible $Z'$ distincte
des $Z_i$ et telle que $Z_k\subset Z'\subset Z_{k+1}$ pour un indice $k$.
\end{env}

\begin{proposition}[14.3.2]
\label{0.14.3.2-fr}
Soit $X$ un espace topologique tel que, pour deux parties fermées
irréductibles quelconques $Y$, $Z$ de $X$ telles que $Y\subset Z$, on ait
$\operatorname{codim}(Y,Z)<+\infty$. Les deux conditions suivantes sont
équivalentes~:
\begin{enumerate}
  \item[\rm(a)] Deux chaînes saturées de parties fermées irréductibles de $X$,
  ayant mêmes extrémités, ont même longueur.
  \item[\rm(b)] Si $Y$, $Z$, $T$ sont trois parties fermées irréductibles de
  $X$ telles que $Y\subset Z\subset T$, on a
  \[
    \operatorname{codim}(Y,T)=
    \operatorname{codim}(Y,Z)+\operatorname{codim}(Z,T).
    \tag{14.3.2.1}\label{0.14.3.2.1-fr}
  \]
\end{enumerate}
\end{proposition}

\begin{proof}
Il est immédiat que a) entraîne b). Réciproquement, supposons b) vérifiée, et
démontrons que si deux chaînes saturées de mêmes extrémités ont pour longueurs
$m$ et $n\leq m$, on a nécessairement $m=n$. Raisonnons par récurrence sur
$n$, la proposition étant évidente pour $n=1$. Supposons donc $n>1$, $n<m$,
et soit $Z_0\subset Z_1\subset\cdots\subset Z_n$ une
\oldpage[0\textsubscript{IV-1}]{11}
chaîne saturée telle qu'il existe une autre chaîne saturée d'extrémités
$Z_0$, $Z_n$ et de longueur $m$. Comme
$\operatorname{codim}(Z_0,Z_n)\geq m>n$ et que
$\operatorname{codim}(Z_0,Z_1)=1$, il résulte de b) que
$\operatorname{codim}(Z_1,Z_n)=\operatorname{codim}(Z_0,Z_n)-1>n-1$, ce qui
contredit l'hypothèse de récurrence.
\end{proof}

Lorsque les conditions de (14.3.2) sont remplies, on dit que $X$ satisfait à
la \emph{condition des chaînes}, ou encore est un \emph{espace caténaire}. Il
est clair que tout sous-espace fermé d'un espace caténaire est caténaire.

\begin{proposition}[14.3.3]
\label{0.14.3.3-fr}
Soit $X$ un espace de Kolmogoroff noethérien de dimension finie. Les
conditions suivantes sont équivalentes~:
\begin{enumerate}
  \item[\rm(a)] Deux chaînes maximales de parties fermées irréductibles de $X$
  ont même longueur.
  \item[\rm(b)] $X$ est équidimensionnel, équicodimensionnel, et caténaire.
  \item[\rm(c)] $X$ est équidimensionnel, et pour deux parties fermées
  irréductibles $Y$, $Z$ de $X$ telles que $Y\subset Z$, on a
  \[
    \dim(Z)=\dim(Y)+\operatorname{codim}(Y,Z).
    \tag{14.3.3.1}\label{0.14.3.3.1-fr}
  \]
  \item[\rm(d)] $X$ est équicodimensionnel et pour deux parties fermées
  irréductibles $Y$, $Z$ de $X$ telles que $Y\subset Z$, on a
  \[
    \operatorname{codim}(Y,X)=
    \operatorname{codim}(Y,Z)+\operatorname{codim}(Z,X).
    \tag{14.3.3.2}\label{0.14.3.3.2-fr}
  \]
\end{enumerate}
\end{proposition}

\begin{proof}
Les hypothèses sur $X$ entraînent que les extrémités d'une chaîne maximale de
parties fermées irréductibles de $X$ sont nécessairement un point fermé et
une composante irréductible de $X$ ($0_{\mathrm I}$, 2.1.3)~; en outre, toute
chaîne saturée d'extrémités $Y$, $Z$ (avec $Y\subset Z$) est contenue dans une
chaîne maximale dont les éléments autres que ceux de la chaîne donnée, ou
bien sont contenus dans $Y$ ou bien contiennent $Z$. Ces remarques
établissent aussitôt l'équivalence de a) et b), et prouvent aussi que si a)
est vérifiée, on a, pour toute partie fermée irréductible $Y$ de $X$
\[
  \dim(Y)+\operatorname{codim}(Y,X)=\dim(X)~;
  \tag{14.3.3.3}\label{0.14.3.3.3-fr}
\]
comme (14.3.2.1) a lieu, on déduit aussitôt (14.3.3.1) et (14.3.3.2) de
(14.3.3.3). Inversement, (14.3.3.1) entraîne (14.3.2.1), donc (14.3.3.1)
entraîne la condition des chaînes en vertu de (14.3.2)~; en outre, en
appliquant (14.3.3.1) au cas où $Y$ est réduit à un point fermé $x$ de $X$ et
$Z$ à une composante irréductible de $X$, on obtient
$\operatorname{codim}(\{x\},X)=\dim(Z)$~; on en conclut que c) entraîne b).
De même, (14.3.3.2) entraîne (14.3.2.1), donc la condition des chaînes~; en
outre, avec le même choix de $Y$ et $Z$ que ci-dessus, (14.3.3.2) entraîne
encore $\operatorname{codim}(\{x\},X)=\dim(Z)$, donc (toute composante
irréductible de $X$ contenant un point fermé en vertu de ($0_{\mathrm I}$,
2.1.3)), d) entraîne b).
\end{proof}

On dit qu'un espace de Kolmogoroff noethérien est \emph{biéquidimensionnel}
s'il est de \emph{dimension finie} et s'il vérifie les conditions équivalentes
de (14.3.3).

\begin{corollary}[14.3.4]
\label{0.14.3.4-fr}
Soit $X$ un espace de Kolmogoroff noethérien biéquidimensionnel~; alors, pour
tout point fermé $x$ de $X$ et toute composante irréductible $Z$ de $X$, on a
\[
  \dim(X)=\dim(Z)=\operatorname{codim}(\{x\},X)=\dim_x(X).
  \tag{14.3.4.1}\label{0.14.3.4.1-fr}
\]
\end{corollary}

\oldpage[0\textsubscript{IV-1}]{12}
\begin{proof}
La dernière égalité résulte de ce que si
$Y_0=\{x\}\subset Y_1\subset\cdots\subset Y_m$ est une chaîne maximale de
parties fermées irréductibles de $X$ et $U$ un voisinage ouvert de $x$, les
$U\cap Y_i$ sont des parties fermées irréductibles de $U$ deux à deux
distinctes (puisque $\overline{U\cap Y_i}=Y_i$), d'où
$\dim(U)=\dim(X)$ en vertu de (14.1.4).
\end{proof}

\begin{corollary}[14.3.5]
\label{0.14.3.5-fr}
Soit $X$ un espace de Kolmogoroff noethérien~; si $X$ est
biéquidimensionnel, il en est de même de toute réunion de composantes
irréductibles de $X$ et de toute partie fermée irréductible de $X$. En outre,
pour toute partie fermée $Y$ de $X$, on a alors
\[
  \dim(Y)+\operatorname{codim}(Y,X)=\dim(X).
  \tag{14.3.5.1}\label{0.14.3.5.1-fr}
\]
\end{corollary}

\begin{proof}
Toute chaîne de parties fermées irréductibles de $X$ étant contenue dans une
composante irréductible de $X$, la première assertion découle aussitôt de
(14.3.3). En outre, si $X'$ est une partie fermée irréductible de $X$, $X'$
vérifie trivialement les conditions (14.3.3, c)), d'où la seconde assertion.

Enfin, pour démontrer (14.3.5.1), remarquons que l'on a vu dans la
démonstration de (14.3.3) que cette relation est vérifiée lorsque $Y$ est
irréductible~; si $Y_i$ ($1\leq i\leq m$) sont les composantes irréductibles
de $Y$, celle des $Y_i$ pour laquelle $\dim(Y_i)$ est le plus grand est aussi
celle pour laquelle $\operatorname{codim}(Y_i,X)$ est le plus petit~; donc
(14.3.5.1) résulte des définitions de $\dim(Y)$ et de
$\operatorname{codim}(Y,X)$.
\end{proof}

\begin{remark}[14.3.6]
\label{0.14.3.6-fr}
Le lecteur remarquera que la démonstration de (14.3.2) s'applique à un
ensemble ordonné quelconque, le fait qu'il s'agit de parties fermées
irréductibles d'un espace topologique n'intervenant pas dans cette
démonstration. Il en est de même de la démonstration de (14.3.3) dans un
ensemble ordonné $E$ tel que pour tout $x\in E$, il existe un $z\leq x$ qui
soit \emph{minimal} dans $E$, et où la longueur des chaînes d'éléments de $E$
est bornée.
\end{remark}

% FIN SOURCE ega4/ega0-14-fr.tex

%
% DEBUT SOURCE ega4/ega0-15-fr.tex
\section{Suites \texorpdfstring{$M$}{M}-régulières et suites
\texorpdfstring{$\mathscr F$}{F}-régulières}
\label{section:0.15-fr}

\subsection{Suites \texorpdfstring{$M$}{M}-régulières et suites
\texorpdfstring{$M$}{M}-quasi-régulières}
\label{subsection:0.15.1-fr}
\label{0.15.1-fr}

\begin{env}[15.1.1]
\label{0.15.1.1-fr}
\oldpage[0\textsubscript{IV-1}]{12}
Soient $A$ un anneau, $N$ un $A$-module~; dans ce qui suit, nous noterons
$N[T_1,\ldots,T_n]$, pour abréger, le $A$-module gradué
$N\otimes_A A[T_1,\ldots,T_n]$ ($T_i$ indéterminées), où $N$ est gradué par la
graduation triviale et $A[T_1,\ldots,T_n]$ par le degré total.

Soient $M$ un $A$-module, $f_i$ ($1\leq i\leq n$) des éléments de $A$,
$\mathfrak{J}=\sum_{i=1}^n f_iA$ l'idéal de $A$ engendré par les $f_i$, et
munissons $M$ de la filtration $\mathfrak{J}$-préadique~; on définit alors un
\emph{homomorphisme gradué surjectif de $A$-modules gradués}, de degré $0$,
\[
  \varphi : (\gr_0(M))[T_1,\ldots,T_n]\longrightarrow\gr_\bullet(M)
  \tag{15.1.1.1}\label{0.15.1.1.1-fr}
\]
de la façon suivante~: si $\xi\in\gr_0(M)$ est la classe mod.
$\mathfrak{J}M$ d'un élément $x\in M$, au couple formé de $\xi$ et d'un
polynôme $P[T_1,\ldots,T_n]$ homogène de degré $k$, on fait correspondre la
classe mod. $\mathfrak{J}^{k+1}M$ de $P(f_1,\ldots,f_n)x$~; il est immédiat
que cette classe
\oldpage[0\textsubscript{IV-1}]{13}
ne dépend que de $\xi$ et de $P$, et si $S_k$ est l'ensemble des polynômes
homogènes de degré $k$ dans $A[T_1,\ldots,T_n]$, on a ainsi défini une
application $A$-linéaire
\[
  \varphi_k:\gr_0(M)\otimes_A S_k\longrightarrow\gr_k(M),
\]
qui est surjective par définition de $\mathfrak{J}^kM$.

Nous dirons dans ce qui suit qu'un élément $f\in A$ est \emph{non diviseur de
zéro dans un $A$-module $M$} si l'homothétie $f_M:x\mapsto fx$ de $M$ est
\emph{injective}.
\end{env}

\begin{lemma}[15.1.2]
\label{0.15.1.2-fr}
Soient $M$ un $A$-module, $f$ un élément de $A$, et munissons $M$ de la
filtration $(fA)$-préadique. Si $f$ est non diviseur de zéro dans $M$,
l'homomorphisme canonique défini dans (15.1.1)~:
\[
  \varphi:(M/fM)[T]\longrightarrow\gr_\bullet(M)
  \tag{15.1.2.1}\label{0.15.1.2.1-fr}
\]
est bijectif. La réciproque est vraie si $M$ est séparé pour la topologie
$(fA)$-préadique.
\end{lemma}

\begin{proof}
On a ici $\gr_k(M)=f^kM/f^{k+1}M$~; dire que $\varphi$ est injective signifie
que pour tout $x\in M$ et tout $k$, la relation $f^kx\in f^{k+1}M$ implique
$x\in fM$~; il est clair qu'il en est bien ainsi si $f$ est non diviseur de
zéro dans $M$. Inversement, remarquons que l'homothétie
$f_M:x\mapsto fx$ définit par passage aux quotients, des homomorphismes
$\gr_k(M)\to\gr_{k+1}(M)$, donc un homomorphisme gradué $g$ de degré $1$ de
$\gr_\bullet(M)$ dans lui-même~; et dire que $\varphi$ est injectif exprime
encore que $g$ est injectif. Si on suppose la filtration de $M$ séparée, il en
résulte que $f_M$ est injective (Bourbaki, \emph{Alg.~comm.}, chap.~III,
\S~2, no~8, cor.~1 du th.~1).
\end{proof}

\begin{corollary}[15.1.3]
\label{0.15.1.3-fr}
Soient $A$ un anneau noethérien, $M$ un $A$-module de type fini. Pour que $f$
soit non diviseur de zéro dans $M$, il faut et il suffit que $\varphi$ soit
bijectif.
\end{corollary}

\begin{proof}
On doit seulement démontrer que la condition est suffisante~; pour prouver
que $f_M$ est injective, il suffit de montrer que pour tout idéal premier
$\mathfrak p$ de $A$, l'homothétie de rapport
$f'=f/1\in A'=A_{\mathfrak p}$ dans le $A_{\mathfrak p}$-module
$M_{\mathfrak p}=M'$ est injective (Bourbaki, \emph{Alg.~comm.}, chap.~II,
\S~3, no~4, th.~1). Or on a
$f'^kM'/f'^{k+1}M'=(f^kM/f^{k+1}M)\otimes_A A'$, et il est immédiat que
l'homomorphisme canonique
$\varphi':(M'/f'M')[T]\to\gr_\bullet(M')$ est égal à
$\varphi\otimes 1_{A'}$~; si $\varphi$ est injectif, il en est donc de même de
$\varphi'$ ($0_{\mathrm I}$, 1.3.2). On est donc ramené au cas où $A$ est un
anneau noethérien \emph{local}, dont nous désignerons par $\mathfrak m$
l'idéal maximal. Si $f\notin\mathfrak m$, alors $f$ est inversible et il n'y a
rien à démontrer. Si $f\in\mathfrak m$, on sait que $M$ est séparé pour la
topologie $(fA)$-préadique ($0_{\mathrm I}$, 7.3.5), et il suffit donc
d'appliquer (15.1.2).
\end{proof}

\begin{definition}[15.1.4]
\label{0.15.1.4-fr}
Soient $A$ un anneau, $M$ un $A$-module. On dit qu'un élément $f\in A$ est
\emph{$M$-régulier} (resp. \emph{$M$-quasi-régulier}) si $f$ est non diviseur
de zéro dans $M$ (resp. si l'homomorphisme (15.1.2.1) est bijectif).
\end{definition}

Nous venons donc de démontrer que tout élément $M$-régulier est
$M$-quasi-régulier, et que la réciproque est vraie si $A$ est noethérien et
$M$ un $A$-module de type fini. Lorsqu'on ne suppose plus $M$ de type fini,
cette réciproque peut être en défaut~: par exemple, si $A$ est un anneau
principal qui n'est pas un corps, $K$ son corps des fractions et $M$ le
$A$-module $K/A$, on a $fM=M$ pour tout $f\neq0$ dans $A$, donc les deux
membres de (15.1.2.1) sont réduits à $0$, mais $f$ est alors diviseur de zéro
dans $M$.

\begin{env}[15.1.5]
\label{0.15.1.5-fr}
\oldpage[0\textsubscript{IV-1}]{14}
Supposons maintenant $M$ muni d'une \emph{filtration} $(M_k)$ telle que
$M_0=M$, et désignons par
$\gr_\bullet(M)=\bigoplus_{k\geq0}M_k/M_{k+1}$ le $A$-module gradué associé.
Les notations étant celles de (15.1.1), posons, pour tout $k\geq0$,
\[
  M'_k=M_k+\mathfrak{J}M_{k-1}+\cdots+
  \mathfrak{J}^{k-1}M_1+\mathfrak{J}^kM_0.
  \tag{15.1.5.1}\label{0.15.1.5.1-fr}
\]
Il est clair que $(M'_k)$ est une filtration sur $M$. Lorsque
$M_k=\mathfrak{Q}^kM$, où $\mathfrak Q$ est un idéal de $A$, la filtration
$(M'_k)$ n'est autre que la filtration
$(\mathfrak{J}+\mathfrak{Q})$-préadique. Nous noterons
$\gr'_\bullet(M)$ le $A$-module gradué associé à $(M'_k)$. Notons encore
\[
  (\gr_\bullet(M)\otimes_A(A/\mathfrak{J}))[T_1,\ldots,T_n]
\]
le produit tensoriel de $A$-modules gradués ($\mathrm{II}$, 2.1.2)
\[
  (\gr_\bullet(M)\otimes_A(A/\mathfrak{J}))
  \otimes_A A[T_1,\ldots,T_n].
\]
On définit un \emph{homomorphisme gradué surjectif de $A$-modules gradués}, de
degré $0$,
\[
  \psi:(\gr_\bullet(M)\otimes_A(A/\mathfrak{J}))[T_1,\ldots,T_n]
  \longrightarrow\gr'_\bullet(M)
  \tag{15.1.5.2}\label{0.15.1.5.2-fr}
\]
de la façon suivante~: si $\alpha\in A/\mathfrak{J}$ est la classe d'un
élément $a\in A$, $\xi\in M_k/M_{k+1}$ la classe d'un élément $x\in M_k$ et
$P[T_1,\ldots,T_n]$ un polynôme homogène de degré $h$, on fait correspondre au
triplet $(\alpha,\xi,P)$ la classe mod. $M'_{h+k+1}$ de
$P(f_1,\ldots,f_n)ax$~; on vérifie encore aussitôt que cet élément ne dépend
bien que de $\xi$, $\alpha$ et $P$, et cela définit l'homomorphisme gradué
(15.1.5.2), qui est surjectif en vertu de la définition (15.1.5.1). On
retrouve l'homomorphisme (15.1.1.1) en considérant la filtration sur $M$ telle
que $M_1=0$.
\end{env}

\begin{lemma}[15.1.6]
\label{0.15.1.6-fr}
On garde les notations de (15.1.5) et on suppose $n=1$, en posant $f=f_1$. Si
$f$ est $\gr_\bullet(M)$-régulier, l'homomorphisme (15.1.5.2) est bijectif.
\end{lemma}

\begin{proof}
Soit $Q_k$ (resp. $Q'_k$) le sous-$A$-module des termes de degré $k$ au
premier (resp. second) membre de (15.1.5.2). Munissons $Q_k$ de la filtration
\[
  (Q_k)_i=\sum_{j\leq k-i}
  (\gr_{k-j}(M)\otimes_A(A/fA))T^j
  \tag{15.1.6.1}\label{0.15.1.6.1-fr}
\]
de sorte que $(Q_k)_0=Q_k$ et $(Q_k)_{k+1}=0$, et pour cette filtration, on a
\[
  \gr_i(Q_k)=(\gr_i(M)\otimes_A(A/fA))T^{k-i}.
\]
Munissons $Q'_k$ de la filtration formée des
$\psi((Q_k)_i)=(Q'_k)_i$. Comme ces filtrations sont \emph{finies}, il suffit
de prouver que les homomorphismes
$\psi_{ki}:\gr_i(Q_k)\to\gr_i(Q'_k)$ déduits de $\psi_k$ sont injectifs
(Bourbaki, \emph{Alg.~comm.}, chap.~III, \S~3, no~8, cor.~1 du th.~1). Or on
a
\[
  \gr_i(Q_k)=((M_i/M_{i+1})\otimes_A(A/fA))T^{k-i}~;
\]
d'autre part, $(Q'_k)_{i+1}$ est l'image de
\[
  M_k+fM_{k-1}+\cdots+f^{k-i-1}M_{i+1}
\]
\oldpage[0\textsubscript{IV-1}]{15}
dans $M'_k/M'_{k+1}$. On est donc ramené à écrire que, pour $x\in M_i$, la
relation
\[
  f^{k-i}x\in M_k+fM_{k-1}+\cdots+
  f^{k-i-1}M_{i+1}+M'_{k+1}
  \tag{15.1.6.2}\label{0.15.1.6.2-fr}
\]
implique $x\in fM_i+M_{i+1}$. Or, le second membre de (15.1.6.2) est contenu
dans $M_{i+1}+M'_{k+1}$, et en vertu de (15.1.5.1), $M'_{k+1}$ est contenu
dans $M_{i+1}+f^{k+1-i}M$. L'hypothèse sur $f$ entraîne que $f$ est non
diviseur de zéro dans $M/M_h$ pour tout $h$ (Bourbaki, \emph{loc.~cit.}). La
relation $f^{k-i}x\in f^{k+1-i}M+M_{i+1}$ entraîne donc (en raisonnant dans
$M/M_{i+1}$) l'existence d'un $y\in M$ tel que $x-fy\in M_{i+1}$, et comme
$x\in M_i$, on a $fy\in M_i$, d'où (en raisonnant dans $M/M_i$), $y\in M_i$,
et finalement $x\in fM_i+M_{i+1}$. C.Q.F.D.
\end{proof}

\begin{definition}[15.1.7]
\label{0.15.1.7-fr}
Soit $(f_i)_{1\leq i\leq n}$ une suite d'éléments de $A$. On dit qu'elle est
\emph{$M$-régulière} si, pour $1\leq i\leq n$, $f_i$ est
$\left(M/(\sum_{1\leq j\leq i-1}f_jM)\right)$-régulier. On dit que la suite
$(f_i)$ est \emph{$M$-quasi-régulière} si
l'homomorphisme canonique $\varphi$ défini dans (15.1.1.1) est bijectif.

Lorsque $M=A$, on dit simplement «~suite régulière~» et
«~suite quasi-régulière~».
\end{definition}

\begin{proposition}[15.1.8]
\label{0.15.1.8-fr}
Les notations étant celles de (15.1.5), on suppose que la suite
$(f_i)_{1\leq i\leq n}$ est $\gr_\bullet(M)$-régulière. Alors
l'homomorphisme canonique (15.1.5.2) est bijectif.
\end{proposition}

\begin{proof}
La proposition n'est autre que (15.1.6) pour $n=1$~; raisonnons par récurrence
sur $n$. Posons
\[
  \mathfrak{J}''=\sum_{i=1}^{n-1}f_iA,
\]
de sorte que $\mathfrak{J}=\mathfrak{J}''+f_nA$, et soit
\[
  M''_k=M_k+\mathfrak{J}''M_{k-1}+\cdots+
  \mathfrak{J}''^{k-1}M_1+\mathfrak{J}''^kM_0~;
\]
on peut donc écrire
\[
  M'_k=M''_k+f_nM''_{k-1}+\cdots+f_n^{k-1}M''_1+f_n^kM''_0.
\]
Notons $\gr''_\bullet(M)$ le $A$-module gradué associé à la filtration
$(M''_k)$ de $M$. On peut écrire, à des isomorphismes canoniques près
\[
  (\gr_\bullet(M)\otimes(A/\mathfrak{J}))\otimes_A A[T_1,\ldots,T_n]
  =
  (\gr_\bullet(M)\otimes(A/\mathfrak{J}''))[T_1,\ldots,T_{n-1}]
  \otimes(A/f_nA)\otimes_A A[T_n]
\]
et par suite (15.1.5.2) se factorise en
\[
  (\gr''_\bullet(M)\otimes(A/f_nA))[T_n]\longrightarrow\gr'_\bullet(M)
  \tag{15.1.8.1}\label{0.15.1.8.1-fr}
\]
et
\[
  ((\gr_\bullet(M)\otimes(A/\mathfrak{J}''))[T_1,\ldots,T_{n-1}])
  \otimes(A/f_nA)\otimes_A A[T_n]
  \xrightarrow{\psi'\otimes1}
  \gr''_\bullet(M)\otimes(A/f_nA)\otimes_A A[T_n]
\]
$\psi'$ étant l'application (15.1.5.2) où l'on remplace $n$ par $n-1$,
$\mathfrak{J}$ par $\mathfrak{J}''$ et $\gr'_\bullet(M)$ par
$\gr''_\bullet(M)$. L'hypothèse de récurrence entraîne que $\psi'$ est
bijective et il reste donc à voir qu'il en est de même de (15.1.8.1). En
appliquant (15.1.6), il suffit de montrer que $f_n$ est
$\gr''_\bullet(M)$-régulier. Mais $\gr''_\bullet(M)$ s'identifie à
\[
  (\gr_\bullet(M)\otimes(A/\mathfrak{J}''))
  [T_1,\ldots,T_{n-1}]~;
\]
par hypothèse, l'homothétie de rapport $f_n$ dans
$\gr_\bullet(M)\otimes(A/\mathfrak{J}'')$ est bijective~; il en est par suite
de même de l'homothétie de rapport $f_n$ dans $\gr''_\bullet(M)$, puisque
$A[T_1,\ldots,T_{n-1}]$ est un $A$-module libre.
\end{proof}

\begin{proposition}[15.1.9]
\label{0.15.1.9-fr}
Soient $M$ un $A$-module, $(f_i)_{1\leq i\leq n}$ une suite d'éléments de $A$.
Si la suite $(f_i)$ est $M$-régulière, elle est $M$-quasi-régulière. La
réciproque est vraie lorsqu'on suppose en outre que les $A$-modules
\[
  M,\quad M/f_1M,\quad\ldots,\quad
  M/\left(\sum_{1\leq j\leq n-1}f_jM\right)
\]
sont séparés pour la topologie $\mathfrak{J}$-préadique (où
$\mathfrak{J}=\sum_{i=1}^n f_iA$).
\end{proposition}

\begin{proof}
La première assertion résulte de (15.1.8), où l'on prend $M_1=0$, et par suite
l'homomorphisme (15.1.5.2) se réduit à (15.1.1.1). Inversement, supposons que
l'homomorphisme $\varphi$ de (15.1.1.1) soit bijectif. Notons que si on pose
\[
  \mathfrak{J}''=\sum_{i=1}^{n-1}f_iA,
\]
\oldpage[0\textsubscript{IV-1}]{16}
et si on désigne par $\gr''_\bullet(M)$ le $A$-module gradué associé à $M$
filtré par la filtration $\mathfrak{J}''$-préadique, $\varphi$ se factorise en
\[
  (\gr''_\bullet(M)\otimes(A/f_nA))[T_n]\longrightarrow\gr_\bullet(M)
  \tag{15.1.9.1}\label{0.15.1.9.1-fr}
\]
et en $\varphi'\otimes1$, où $\varphi'$ est l'homomorphisme canonique
\[
  (\gr''_0(M))[T_1,\ldots,T_{n-1}]
  \longrightarrow\gr''_\bullet(M).
\]
Comme ces deux homomorphismes sont surjectifs, ils sont bijectifs si
$\varphi$ est supposé bijectif, donc $\varphi'$ est nécessairement injectif,
et par suite bijectif puisqu'il est surjectif. On peut alors raisonner par
récurrence sur $n$ (le cas $n=1$ résultant de (15.1.2)), car sur
\[
  M,\quad M/f_1M,\quad\ldots,\quad
  M/\left(\sum_{j=1}^{n-2}f_jM\right)
\]
la topologie $\mathfrak{J}''$-préadique est plus fine que la topologie
$\mathfrak{J}$-préadique, et par suite est séparée. On voit donc que la suite
$(f_i)_{1\leq i\leq n-1}$ est $M$-régulière. D'autre part, montrons que
l'hypothèse entraîne que si $y\in M/\mathfrak{J}''M$ est tel que
\[
  f_n^ky\in f_n^{k+1}(M/\mathfrak{J}''M),
\]
alors on a $y\in f_n(M/\mathfrak{J}''M)$. En effet, si $x$ est un représentant
de $y$, montrons que l'on a $f_n^kx\in\mathfrak{J}^{k+1}M$~; notre assertion
résultera alors du fait que (15.1.9.1) est injectif. L'hypothèse entraîne
\[
  f_n^kx\in f_n^{k+1}M+\mathfrak{J}''M\subset\mathfrak{J}M,
\]
d'où, puisque $\varphi$ est injectif,
$f_n^{k-1}x\in\mathfrak{J}^kM$~; par récurrence, on obtient ainsi
$x\in\mathfrak{J}M$, d'où finalement
$f_n^kx\in\mathfrak{J}^{k+1}M$. On termine le raisonnement comme dans
(15.1.2)~: sur $M/\mathfrak{J}''M$, la filtration
$\mathfrak{J}$-préadique est identique à la filtration
$(f_nA)$-préadique, donc est séparée par hypothèse, et on conclut que
l'homothétie de rapport $f_n$ dans $M/\mathfrak{J}''M$ est injective, ce qui
achève de montrer que la suite $(f_i)_{1\leq i\leq n}$ est $M$-régulière.
\end{proof}

\begin{corollary}[15.1.10]
\label{0.15.1.10-fr}
Avec les notations de (15.1.9), supposons que pour tout $A$-module $N$ de type
fini, $M\otimes_A N$ soit séparé pour la topologie
$\mathfrak{J}$-préadique. Alors, pour que la suite $(f_i)$ soit
$M$-régulière, il faut et il suffit qu'elle soit $M$-quasi-régulière.
\end{corollary}

On notera que l'hypothèse de (15.1.10) entraîne que, si la suite
$(f_i)_{1\leq i\leq n}$ est $M$-régulière, il en est de même de la suite
$(f_{\sigma(i)})$ pour toute permutation $\sigma$ de l'intervalle $[1,n]$.

\begin{corollary}[15.1.11]
\label{0.15.1.11-fr}
Soient $A$ un anneau noethérien, $M$ un $A$-module de type fini,
$(f_i)_{1\leq i\leq n}$ une suite finie d'éléments contenus dans le radical de
$A$. Alors, pour que la suite $(f_i)$ soit $M$-régulière, il faut et il suffit
qu'elle soit $M$-quasi-régulière.
\end{corollary}

\begin{proof}
C'est un cas particulier de (15.1.10), puisque tout $A$-module de type fini
est alors séparé pour la topologie $\mathfrak{J}$-préadique
($0_{\mathrm I}$, 7.3.5).
\end{proof}

\begin{remarks}[15.1.12]
\label{0.15.1.12-fr}
\mbox{}
\begin{enumerate}
  \item[(i)] Lorsque $A$ est un anneau local noethérien, $M$ un $A$-module de
  type fini, et lorsque les $f_i$ sont contenus dans l'idéal maximal de $A$,
  la notion de suite $M$-régulière a été introduite par J.-P.~Serre sous le
  nom de $M$-\emph{suite} [17].
  \item[(ii)] Même si $A$ est noethérien et $M=A$, une suite $(f,g)$ de deux
  éléments de $A$ qui est quasi-régulière n'est pas nécessairement régulière.
  Par exemple, la suite $(0,1)$ est quasi-régulière, puisqu'alors
  $\mathfrak{J}=fA+gA=A$, donc $\gr_\bullet(A)=0$~; mais si $A\neq0$ elle
  n'est pas régulière~; on notera toutefois que dans ce cas la suite $(1,0)$
  est régulière, car $1$ n'est pas diviseur de $0$ dans $A$, et comme
  $A/A=0$, $0$ n'est pas diviseur de zéro dans ce module (cf. (15.1.1)).
  \item[(iii)]
  \oldpage[0\textsubscript{IV-1}]{17}
  Dans (15.1.6), on ne peut en général remplacer l'hypothèse que $f$ est
  $\gr_\bullet(M)$-régulier par l'hypothèse qu'il est
  $\gr_\bullet(M)$-quasi-régulier. Pour le voir, reprenons l'exemple (15.1.4)
  où $A$ est un anneau principal de corps des fractions $K\neq A$ et prenons
  $M=K$, la filtration $(M_k)$ de $M$ étant définie par
  $M_0=M=K$, $M_1=A$, $M_k=0$ pour $k\geq2$, d'où
  $\gr_0(M)=K/A$, $\gr_1(M)=A$, $\gr_k(M)=0$ pour $k\geq2$. Si $f\neq0$ dans
  $A$, on a vu que $f$ n'est pas $\gr_\bullet(M)$-régulier, mais comme on a
  \[
    f^k\gr_\bullet(M)/f^{k+1}\gr_\bullet(M)=f^kA/f^{k+1}A,
  \]
  il est immédiat que l'homomorphisme (15.1.2.1) (où $M$ est remplacé par
  $\gr_\bullet(M)$) est injectif, donc $f$ est
  $\gr_\bullet(M)$-quasi-régulier. Mais avec les notations de (15.1.6), on a
  alors $M'_k=K$ pour tout $k\geq0$, donc $\gr'_\bullet(M)=0$, tandis que le
  premier membre de (15.1.5.2) n'est pas réduit à $0$.
\end{enumerate}
\end{remarks}

\begin{proposition}[15.1.13]
\label{0.15.1.13-fr}
Soient $A$ un anneau, $M$, $N$ deux $A$-modules, $(f_i)_{1\leq i\leq n}$ une
suite d'éléments de $A$. Si la suite $(f_i)$ est $M$-régulière et si $N$ est
un $A$-module plat, alors la suite $(f_i)$ est
$M\otimes_A N$-régulière. La réciproque est vraie si $N$ est un $A$-module
fidèlement plat.
\end{proposition}

\begin{proof}
En effet (sans hypothèse sur $N$),
\[
  (M\otimes_A N)/
  \left(\sum_{j=1}^{i-1}f_j(M\otimes_A N)\right)
\]
est isomorphe à
\[
  \left(M/\left(\sum_{j=1}^{i-1}f_jM\right)\right)\otimes_A N~;
\]
il suffit alors d'appliquer ($0_{\mathrm I}$, 6.1.1) et
($0_{\mathrm I}$, 6.4.1) à l'homothétie de rapport $f_i$ dans
$M/(\sum_{j=1}^{i-1}f_jM)$.
\end{proof}

\begin{corollary}[15.1.14]
\label{0.15.1.14-fr}
Soient $A$ un anneau, $M$ un $A$-module, $(f_i)_{1\leq i\leq n}$ une suite
finie d'éléments de $A$. Soient $\rho:A\to B$ un homomorphisme d'anneaux,
$M'=M_{(B)}=M\otimes_A B$, $f'_i=\rho(f_i)$ pour $1\leq i\leq n$. Si $B$ est
un $A$-module plat et si la suite $(f_i)$ est $M$-régulière, alors la suite
$(f'_i)$ est $M'$-régulière~; la réciproque est vraie si $B$ est un
$A$-module fidèlement plat.
\end{corollary}

\begin{proof}
En effet,
\[
  M'/\left(\sum_{j=1}^{i-1}f'_jM'\right)
\]
s'identifie à
\[
  \left(M/\left(\sum_{j=1}^{i-1}f_jM\right)\right)\otimes_A B,
\]
et l'homothétie de rapport $f'_i$ dans le $B$-module
$M'/(\sum_{j=1}^{i-1}f'_jM')$ à l'homothétie de rapport $f_i$ dans le
$A$-module
\[
  \left(M/\left(\sum_{j=1}^{i-1}f_jM\right)\right)\otimes_A B~;
\]
l'assertion résulte donc de (15.1.13).
\end{proof}

\begin{proposition}[15.1.15]
\label{0.15.1.15-fr}
Soient $A$ un anneau, $B$ une $A$-algèbre commutative,
$(f_i)_{1\leq i\leq n}$ une suite finie d'éléments de $B$, $M$ un
$B$-module. Supposons que la suite $(f_i)$ soit $M$-régulière et que chacun
des $A$-modules
\[
  M/\left(\sum_{j=1}^{i-1}f_jM\right)\quad(1\leq i\leq n)
\]
soit plat. Alors, pour tout homomorphisme d'anneaux $\rho:A\to A'$, si l'on
pose $B'=B\otimes_A A'$, $M'=M\otimes_A A'$ et $f'_i=f_i\otimes1$
($1\leq i\leq n$), la suite $(f'_i)$ est $M'$-régulière et les
$A'$-modules
\[
  M'/\left(\sum_{j=1}^{i-1}f'_jM'\right)\quad(1\leq i\leq n)
\]
sont plats.
\end{proposition}

\begin{proof}
Considérons la suite (exacte par hypothèse)
\[
  0\longrightarrow M\xrightarrow{f_1}M\longrightarrow M/f_1M
  \longrightarrow0
\]
(la flèche $M\xrightarrow{f_1}M$ étant l'homothétie de rapport $f_1$ dans
$M$). L'hypothèse que $M/f_1M$ est $A$-plat entraîne que la suite
\[
  0\longrightarrow M\otimes_A A'
  \xrightarrow{f_1\otimes1}M\otimes_A A'
  \longrightarrow(M/f_1M)\otimes_A A'\longrightarrow0
\]
\oldpage[0\textsubscript{IV-1}]{18}
est exacte ($0_{\mathrm I}$, 6.1.2). Cela montre que l'homothétie de rapport
$f'_1$ dans $M'$ est injective.

Posons ensuite, pour $1\leq i\leq n-1$,
\[
  M_i=M/\left(\sum_{j=1}^i f_jM\right),\qquad
  M'_i=M'/\left(\sum_{j=1}^i f'_jM'\right)~;
\]
on a $M'_i=M_i\otimes_A A'$, $M_i/f_{i+1}M_i=M_{i+1}$,
$M'_i/f'_{i+1}M'_i=M'_{i+1}$~; il suffit de remplacer dans le raisonnement
précédent $M$ et $f_1$ par $M_i$ et $f_{i+1}$ pour conclure que l'homothétie
de rapport $f'_{i+1}$ dans $M'_i$ est injective, en raison de la platitude de
$M_i/f_{i+1}M_i$. La dernière assertion résulte de ($0_{\mathrm I}$, 6.2.1).
\end{proof}

\begin{proposition}[15.1.16]
\label{0.15.1.16-fr}
Soient $A$, $B$ deux anneaux locaux noethériens, $k$ le corps résiduel de
$A$, $\varphi:A\to B$ un homomorphisme local, $M$ un $B$-module de type fini.
Soit $(f_i)_{1\leq i\leq n}$ une suite d'éléments de l'idéal maximal de $B$,
et pour tout $i$, soit $g_i$ l'image de $f_i$ dans $B\otimes_A k$. Les
conditions suivantes sont équivalentes~:
\begin{enumerate}
  \item[\rm(a)] La suite $(f_i)$ est $M$-régulière et les quotients
  \[
    M_i=M/\left(\sum_{j=1}^i f_jM\right)
  \]
  sont des $A$-modules plats pour $1\leq i\leq n$.
  \item[\rm(b)] La suite $(f_i)$ est $M$-régulière et le $A$-module $M_n$ est
  plat.
  \item[\rm(c)] Le $A$-module $M$ est plat et la suite
  $(g_i)_{1\leq i\leq n}$ est $(M\otimes_A k)$-régulière.
  \item[\rm(d)] Le $A$-module $M$ est plat, et pour tout homomorphisme
  d'anneaux $\rho:A\to A'$, si on pose $B'=B\otimes_A A'$,
  $M'=M\otimes_A A'$, $f'_i=f_i\otimes1$ ($1\leq i\leq n$), la suite
  $(f'_i)$ est $M'$-régulière.
\end{enumerate}
\end{proposition}

\begin{proof}
Il est clair que a) entraîne b)~; pour montrer que b) entraîne d), il suffit,
en vertu de (15.1.15), de montrer que b) entraîne que les $M_i$ sont des
$A$-modules plats pour $0\leq i\leq n-1$ (en posant $M_0=M$). Comme
$M_{i+1}=M_i/f_{i+1}M_i$, on est ramené, par récurrence descendante, à prouver
que si $h$ est un élément de l'idéal maximal de $B$ tel que l'homothétie de
rapport $h$ dans $M$ soit injective et que $M/hM$ soit un $A$-module plat,
alors $M$ est un $A$-module plat~; mais cela n'est autre que
($0_{\mathrm{III}}$, 10.2.7). Il est clair que c) est le cas particulier de
d) obtenu en prenant $A'=k$. Reste donc à prouver que c) entraîne a). Comme
$M$ est un $A$-module plat et un $B$-module de type fini, il résulte de c) et
de ($0_{\mathrm{III}}$, 10.2.4) que l'homothétie de rapport $f_1$ dans $M$
est injective et que son conoyau $M/f_1M=M_1$ est un $A$-module plat.
Raisonnons par récurrence sur $i$ et supposons démontré que $M_i$ est un
$A$-module plat~; comme l'homothétie de rapport $g_{i+1}$ dans
$M_i\otimes_A k$ est injective par hypothèse, il résulte encore de
($0_{\mathrm{III}}$, 10.2.4) que l'homothétie de rapport $f_{i+1}$ dans
$M_i$ est injective et que $M_i/f_{i+1}M_i=M_{i+1}$ est un $A$-module plat,
ce qui achève la démonstration. On notera que la démonstration du fait que c)
entraîne a) n'utilise pas le fait que les $f_i$ appartiennent à l'idéal
maximal de $B$.
\end{proof}

\begin{corollary}[15.1.17]
\label{0.15.1.17-fr}
Les hypothèses sur $A$ et $B$ étant celles de (15.1.16), soient $M$ un
$B$-module de type fini, plat sur $A$, $R$ un sous-$B$-module de $M$ tel que
$N=M/R$ soit un $A$-module plat. Soit $(f_i)_{1\leq i\leq n}$ une suite
d'éléments de l'idéal maximal de $B$ ayant les propriétés suivantes~:
\begin{enumerate}
  \item[(i)] $f_iM\subset R$ pour $1\leq i\leq n$.
  \item[(ii)] Si $g_i$ est l'image de $f_i$ dans $B\otimes_A k$, la suite
  $(g_i)_{1\leq i\leq n}$ est $(M\otimes_A k)$-régulière.
  \item[(iii)] La somme des sous-$(B\otimes_A k)$-modules
  $g_i(M\otimes_A k)$ de $M\otimes_A k$ est l'image canonique de
  $R\otimes_A k$ dans $M\otimes_A k$ (égale au noyau de l'homomorphisme
  canonique $M\otimes_A k\to N\otimes_A k$).
\end{enumerate}
Alors la suite $(f_i)$ est $M$-régulière et l'on a
$R=\sum_{i=1}^n f_iM$.
\end{corollary}

\begin{proof}
\oldpage[0\textsubscript{IV-1}]{19}
Comme $N$ est un $A$-module plat, il résulte de l'exactitude de la suite
$0\to R\to M\to N\to0$ que la suite
\[
  0\longrightarrow R\otimes_A k\longrightarrow M\otimes_A k
  \longrightarrow N\otimes_A k\longrightarrow0
\]
est exacte ($0_{\mathrm I}$, 6.1.2)~; on a donc
$R\otimes_A k=\sum_i g_i(M\otimes_A k)$. Comme $M$ est un $B$-module de type
fini et que $B$ est noethérien, $R$ est aussi un $B$-module de type fini,
d'où $R=\sum_i f_iM$ en vertu du lemme de Nakayama, les $f_i$ appartenant à
l'idéal maximal de $B$.
\end{proof}

\begin{lemma}[15.1.18]
\label{0.15.1.18-fr}
Soient $A$ un anneau, $M$ un $A$-module muni d'une filtration séparée et
exhaustive $(M_k)_{k\geq0}$, et soit $f\in A$ un élément
$\gr_\bullet(M)$-régulier. Alors~:
\begin{enumerate}
  \item[(i)] $f$ est $M$-régulier.
  \item[(ii)] Pour tout $k$, on a $M_k\cap fM=fM_k$, et l'image canonique
  $M'_k$ de $M_k$ dans $M'=M/fM$ est donc canoniquement isomorphe à
  $M_k/fM_k$.
  \item[(iii)] La suite
  \[
    0\longrightarrow M_{k+1}/fM_{k+1}\longrightarrow M_k/fM_k
    \longrightarrow\gr_k(M)/f\cdot\gr_k(M)\longrightarrow0
  \]
  est exacte, et si l'on munit $M'$ de la filtration quotient $(M'_k)$,
  $\gr_\bullet(M')$ est donc canoniquement isomorphe à
  $\gr_\bullet(M)/f\cdot\gr_\bullet(M)$.
\end{enumerate}
\end{lemma}

\begin{proof}
L'assertion (i) résulte de l'hypothèse sur $f$ et de Bourbaki,
\emph{Alg.~comm.}, chap.~III, \S~2, no~8, cor.~1 du th.~1. Comme
$\gr_\bullet(M/M_k)$ est facteur direct de $\gr_\bullet(M)$, $f$ est
$\gr_\bullet(M/M_k)$-régulier pour tout $k$, donc $(M/M_k)$-régulier par
(i), ce qui établit (ii). Pour prouver (iii), il suffit de voir que
l'application $M_{k+1}/fM_{k+1}\to M_k/fM_k$ est injective~: or, si
$x\in M_{k+1}$ est tel que $x\in fM_k$, il résulte de (ii) que
$x\in fM_{k+1}$, ce qui achève de prouver le lemme.
\end{proof}

\begin{proposition}[15.1.19]
\label{0.15.1.19-fr}
Soient $A$ un anneau, $M$ un $A$-module muni d'une filtration
$(M_k)_{k\geq0}$ telle que $M_0=M$, $(f_i)_{1\leq i\leq n}$ une suite
$\gr_\bullet(M)$-régulière d'éléments de $A$. Supposons en outre que pour
$0\leq i\leq n-1$, la filtration quotient de $(M_k)$ sur
\[
  M/\left(\sum_{j\leq i}f_jM\right)
\]
soit séparée. Alors la suite $(f_i)$ est $M$-régulière, et si l'on pose
\[
  M'=M/\left(\sum_{i\leq n}f_iM\right)
\]
et que l'on munit $M'$ de la filtration quotient de $(M_k)$,
$\gr_\bullet(M')$ est canoniquement isomorphe à
\[
  \gr_\bullet(M)/
  \left(\sum_{i\leq n}f_i\gr_\bullet(M)\right).
\]
\end{proposition}

\begin{proof}
Pour $n=1$, la proposition résulte du lemme (15.1.18). Raisonnons par
récurrence sur $n$, en posant
\[
  N=M/\left(\sum_{i\leq n-1}f_iM\right)
\]
et en désignant par $(N_k)$ la filtration quotient de $(M_k)$ sur $N$, qui
est séparée par hypothèse~; en outre $\gr_\bullet(N)$ est isomorphe à
\[
  \gr_\bullet(M)/
  \left(\sum_{i\leq n-1}f_i\gr_\bullet(M)\right).
\]
Par hypothèse $f_n$ est donc $\gr_\bullet(N)$-régulier, et il résulte donc du
lemme (15.1.18) appliqué à $N$ que $f_n$ est $N$-régulier et
$\gr_\bullet(M')$ isomorphe à
$\gr_\bullet(N)/f_n\gr_\bullet(N)$, ce qui démontre la proposition.
\end{proof}

La proposition (15.1.19) s'appliquera en particulier pour les filtrations
vérifiant l'une ou l'autre hypothèse suivante~:
\begin{enumerate}
  \item[$1^\circ$] La filtration $(M_k)$ est finie et séparée (puisque cela
  entraîne $M_k=0$ pour $k$ assez grand)~;
  \item[$2^\circ$] $A$ est un anneau noethérien, $\mathfrak{J}$ un idéal
  contenu dans le radical de $A$, $M$ un $A$-module de type fini et $(M_k)$
  la filtration $\mathfrak{J}$-préadique ($0_{\mathrm I}$, 7.3.5).
\end{enumerate}

\begin{corollary}[15.1.20]
\label{0.15.1.20-fr}
\oldpage[0\textsubscript{IV-1}]{20}
Soient $A$ un anneau, $M$ un $A$-module, $(f_i)_{1\leq i\leq n}$ une suite
$M$-régulière d'éléments de $A$. Alors, quels que soient les entiers
$\nu_i\geq1$, la suite $(f_i^{\nu_i})$ est $M$-régulière et
$M/(\sum_{i\leq n}f_i^{\nu_i}M)$ a une suite de composition dont les
quotients sont isomorphes à $M/(\sum_{i\leq n}f_iM)$.
\end{corollary}

\begin{proof}
Raisonnons par récurrence sur $n$~; pour $n=1$ les assertions sont immédiates
puisque $f^jM/f^{j+1}M$ est isomorphe à $M/fM$, l'homothétie de rapport $f$
dans $M$ étant injective par hypothèse. Posons alors
\[
  N=M/\left(\sum_{i\leq n-1}f_i^{\nu_i}M\right)~;
\]
par l'hypothèse de récurrence, $N$ admet une filtration finie $(N_k)$ pour
laquelle les $\gr_k(N)$ sont isomorphes à
$M/(\sum_{i\leq n-1}f_iM)$~; par hypothèse $f_n$ est donc
$\gr_\bullet(N)$-régulier~; il résulte par suite de (15.1.19) que $f_n$ est
$N$-régulier. Appliquant le cas $n=1$ démontré au début, on en conclut que
$f_n^{\nu_n}$ est $N$-régulier et que $N/f_n^{\nu_n}N$ admet une suite de
composition dont les quotients sont isomorphes à $N/f_nN$~; mais pour la
filtration de $N/f_nN$ quotient de la filtration $(N_k)$,
$\gr_k(N/f_nN)$ est isomorphe à $\gr_k(N)/f_n\gr_k(N)$ par (15.1.19) pour
tout $k$, donc à $M/(\sum_{i\leq n}f_iM)$, ce qui prouve le corollaire.
\end{proof}

\begin{proposition}[15.1.21]
\label{0.15.1.21-fr}
Soient $A$, $B$ deux anneaux locaux noethériens, $\varphi:A\to B$ un
homomorphisme local, $M$ un $B$-module de type fini, $(f_i)_{1\leq i\leq n}$
une suite $A$-régulière formée d'éléments de l'idéal maximal de $A$. Les
propriétés suivantes sont équivalentes~:
\begin{enumerate}
  \item[a)] $M$ est un $A$-module plat.
  \item[b)] Si $\mathfrak{J}$ est l'idéal $\sum_{i\leq n}f_iA$,
  $M/\mathfrak{J}M$ est un $A/\mathfrak{J}$-module plat et la suite $(f_i)$
  est $M$-régulière.
\end{enumerate}
\end{proposition}

\begin{proof}
Le fait que a) entraîne b) résulte de (15.1.13) et de
($0_{\mathrm{III}}$, 10.2.2). Inversement, si la suite $(f_i)$ est
$M$-régulière, elle est $M$-quasi-régulière (15.1.9), donc l'homomorphisme
(15.1.1.1) est bijectif. Mais comme la suite $(f_i)_{1\leq i\leq n}$ est aussi
supposée $A$-régulière, elle est $A$-quasi-régulière (15.1.9), et
l'homomorphisme canonique (15.1.1.1) correspondant
$(A/\mathfrak{J})[T_1,\ldots,T_n]\to\gr_\bullet(A)$ est bijectif. On en
conclut que l'homomorphisme canonique ($0_{\mathrm{III}}$, 10.1.1.2)
\[
  \gr_0(M)\otimes_{A/\mathfrak{J}}\gr_\bullet(A)
  \longrightarrow\gr_\bullet(M)
\]
est bijectif. La conclusion résulte alors de ($0_{\mathrm{III}}$, 10.2.2),
les $f_i$ étant dans l'idéal maximal de $A$ et l'homomorphisme $\varphi$ étant
local.
\end{proof}

\subsection{Suites \texorpdfstring{$\mathscr{F}$}{F}-régulières}
\label{subsection:0.15.2-fr}

\begin{env}[15.2.1]
\label{0.15.2.1-fr}
Soient $X$ un espace annelé, $\mathscr{F}$ un $\mathcal{O}_X$-Module,
$f_i$ ($1\leq i\leq n$) une suite de sections de $\mathcal{O}_X$ au-dessus
de $X$, $\mathscr{J}$ l'Idéal de $\mathcal{O}_X$ engendré par les $f_i$
(image de $\mathcal{O}_X^n$ par l'homomorphisme
$\mathcal{O}_X^n\to\mathcal{O}_X$ défini canoniquement par les $f_i$
($0_{\mathrm I}$, 5.1.1)). Les sous-$\mathcal{O}_X$-Modules
$\mathscr{J}^k\mathscr{F}$ définissent sur $\mathscr{F}$ une filtration, et
on pose encore
$\gr_k(\mathscr{F})=\mathscr{J}^k\mathscr{F}/\mathscr{J}^{k+1}\mathscr{F}$~;
si $S_k$ est l'ensemble des polynômes homogènes de degré total $k$ dans
l'anneau $\mathbb{Z}[T_1,\ldots,T_n]$, $S_k$ peut être considéré comme un
faisceau simple sur $X$, et on définit comme dans (15.1.1) un homomorphisme
canonique
\[
\label{0.15.2.1.1-fr}
  \varphi_k:\gr_0(\mathscr{F})\otimes_{\mathbb Z}S_k
  \longrightarrow\gr_k(\mathscr{F}).
\tag{15.2.1.1}
\]
\oldpage[0\textsubscript{IV-1}]{21}
de la façon suivante~: pour tout ouvert $U$ de $X$ (ou d'une base de la
topologie de $X$), on considère l'homomorphisme (15.1.1.1)
\[
  \varphi_{k,U}:\gr_0(\Gamma(U,\mathscr{F}))\otimes_{\mathbb Z}S_k
  \longrightarrow\gr_k(\Gamma(U,\mathscr{F}))
\]
défini par les $n$ sections $f_i|U\in\Gamma(U,\mathcal{O}_X)$, le
$\Gamma(U,\mathcal{O}_X)$-module $\Gamma(U,\mathscr{F})$ étant filtré par les
sous-modules $(\Gamma(U,\mathscr{J}))^k\Gamma(U,\mathscr{F})$. Le
$\mathcal{O}_X$-Module $\gr_k(\mathscr{F})$ étant associé au préfaisceau
$U\mapsto\gr_k(\Gamma(U,\mathscr{F}))$ pour tout $k$, les $\varphi_{k,U}$
définissent un homomorphisme de $\mathcal{O}_X$-Modules (15.2.1.1). En vertu
des propriétés d'exactitude des limites inductives de modules et de leur
commutation aux produits tensoriels, il résulte de ce qui précède que pour
tout $x\in X$, l'homomorphisme $(\varphi_k)_x$ des fibres des deux membres de
(15.2.1.1) au point $x$ s'identifie à l'homomorphisme de (15.1.1)
\[
\label{0.15.2.1.2-fr}
  \gr_0(\mathscr{F}_x)\otimes_{\mathbb Z}S_k
  \longrightarrow\gr_k(\mathscr{F}_x)
\tag{15.2.1.2}
\]
défini par la suite $((f_i)_x)_{1\leq i\leq n}$ d'éléments de
$\mathcal{O}_x$.
\end{env}

\begin{env}[15.2.2]
\label{0.15.2.2-fr}
Les notations étant celles de (15.2.1), nous dirons que la suite
$(f_i)_{1\leq i\leq n}$ est $\mathscr{F}$-régulière si, en désignant par
$f_j\mathscr{F}$ le sous-$\mathcal{O}_X$-Module de $\mathscr{F}$ image de
l'endomorphisme de $\mathscr{F}$ défini par $f_j$, l'endomorphisme de
$\mathscr{F}/(\sum_{j=1}^{i-1}f_j\mathscr{F})$ défini par $f_i$ est injectif
pour $1\leq i\leq n$~; nous dirons que la suite $(f_i)$ est
$\mathscr{F}$-quasi-régulière si les homomorphismes (15.2.1.1) sont injectifs.
Il résulte de ces définitions et des remarques de (15.2.1) que pour que la
suite $(f_i)$ soit $\mathscr{F}$-régulière (resp.
$\mathscr{F}$-quasi-régulière), il faut et il suffit que, pour tout $x\in X$,
la suite $((f_i)_x)_{1\leq i\leq n}$ soit $\mathscr{F}_x$-régulière (resp.
$\mathscr{F}_x$-quasi-régulière). Toute suite $\mathscr{F}$-régulière est donc
$\mathscr{F}$-quasi-régulière (15.1.9)~; la réciproque est vraie si
$\mathscr{F}$ est un $\mathcal{O}_X$-Module de type fini, si les
$\mathcal{O}_x$ sont des anneaux noethériens pour tout $x\in X$ et si en outre
les $(f_i)_x$ appartiennent au radical de $\mathcal{O}_x$ pour tout $x\in X$
(15.1.11).

Si $X=\operatorname{Spec}(A)$ est un schéma affine, et
$\mathscr{F}=\widetilde{M}$, il
résulte de (I, 1.3) que pour que la suite $(f_i)$ soit
$\mathscr{F}$-régulière (resp. $\mathscr{F}$-quasi-régulière), il faut et il
suffit qu'elle soit $M$-régulière (resp. $M$-quasi-régulière).

Lorsque $\mathscr{F}=\mathcal{O}_X$, on supprimera la mention de
$\mathscr{F}$ dans les définitions précédentes.
\end{env}

\begin{remark}[15.2.3]
\label{0.15.2.3-fr}
Lorsque $X=\operatorname{Spec}(A)$ et
$\mathfrak{J}=\sum_{i=1}^n f_iA$, la suite $(f_j)$ est
quasi-régulière pour tout $A$-module $M$ dans tout ouvert $U$ ne rencontrant
pas $V(\mathfrak{J})$, car en tout point $x\in U$, on a
$\mathfrak{J}_x=\mathcal{O}_x$ et les deux membres de (15.2.1.2) sont nuls. La
notion de suite $\widetilde{M}$-régulière n'a donc d'intérêt qu'au voisinage
des points de $V(\mathfrak{J})$.
\end{remark}

\begin{proposition}[15.2.4]
\label{0.15.2.4-fr}
Soient $X$ un espace annelé, $\mathscr{F}$ un $\mathcal{O}_X$-Module cohérent,
$x$ un point de $X$, $(f_i)_{1\leq i\leq n}$ une suite d'éléments de
$\Gamma(X,\mathcal{O}_X)$. Si la suite $((f_i)_x)_{1\leq i\leq n}$ est
$\mathscr{F}_x$-régulière, il existe un voisinage ouvert $U$ de $x$ tel que la
suite $(f_i|U)_{1\leq i\leq n}$ soit $(\mathscr{F}|U)$-régulière.
\end{proposition}

\begin{proof}
En effet, pour $1\leq i\leq n$, le $\mathcal{O}_X$-Module
$\mathscr{F}/(\sum_{j=1}^{i-1}f_j\mathscr{F})$ est cohérent
($0_{\mathrm I}$, 5.3.4), donc il en est de même du noyau de
l'endomorphisme de ce $\mathcal{O}_X$-Module défini par $f_i$
($0_{\mathrm I}$, 5.3.4)~; comme la fibre au point $x$ de ce noyau est nulle,
il en est de même de sa restriction à un voisinage de $x$
($0_{\mathrm I}$, 5.2.2).
\end{proof}

\begin{proposition}[15.2.5]
\label{0.15.2.5-fr}
\oldpage[0\textsubscript{IV-1}]{22}
Soit $u=(\psi,\theta):X\to Y$ un morphisme plat d'espaces annelés
($0_{\mathrm I}$, 6.7.1)~; soient $\mathscr{F}$ un
$\mathcal{O}_Y$-Module, $\mathscr{G}=u^*(\mathscr{F})$,
$(f_i)_{1\leq i\leq n}$ une suite d'éléments de
$\Gamma(Y,\mathcal{O}_Y)$, $(g_i)_{1\leq i\leq n}$ la suite de leurs images
$g_i=\theta(f_i)\in\Gamma(Y,u_*(\mathcal{O}_X))=\Gamma(X,\mathcal{O}_X)$. Si
la suite $(f_i)$ est $\mathscr{F}$-régulière, la suite $(g_i)$ est
$\mathscr{G}$-régulière. La réciproque est vraie si $u$ est un morphisme
fidèlement plat ($0_{\mathrm I}$, 6.7.8).
\end{proposition}

\begin{proof}
Soient $x$ un point de $X$, $y=u(x)=\psi(x)$~; on a
\[
  \mathscr{G}_x=\psi_x^*(\mathscr{F}_y)
  \otimes_{\psi_x^*(\mathcal{O}_y)}\mathcal{O}_x,
  \qquad\text{et}\qquad
  (g_i)_x=\theta_x^\#(\psi_x^*((f_i)_y))~;
\]
comme par hypothèse
$\theta_x^\#:\psi_x^*(\mathcal{O}_y)\to\mathcal{O}_x$ fait de
$\mathcal{O}_x$ un $\psi_x^*(\mathcal{O}_y)$-module plat, et que le foncteur
$\psi^*$ est exact, la première assertion résulte de l'assertion correspondante
de (15.1.14)~; il en est de même de la seconde, compte tenu de ce que
l'hypothèse que $u$ est fidèlement plat signifie que $\psi$ est surjectif et
que $\theta_x^\#$ fait de $\mathcal{O}_x$ un
$\psi_x^*(\mathcal{O}_y)$-module fidèlement plat.
\end{proof}

% FIN SOURCE ega4/ega0-15-fr.tex

%
% DEBUT SOURCE ega4/ega0-16-fr.tex
\section{Dimension et profondeur dans les anneaux locaux noethériens}
\label{section:0.16-fr}

\subsection{Dimension d'un anneau}
\label{subsection:0.16.1-fr}

\begin{env}[16.1.1]
\label{0.16.1.1-fr}
On appelle dimension (ou dimension de Krull) d'un anneau $A$ et on note
$\dim A$ la dimension (combinatoire) de son spectre
$\operatorname{Spec}(A)$ (14.2.1)~; comme les parties fermées irréductibles de
$\operatorname{Spec}(A)$ sont les
$V(\mathfrak{p})$, où $\mathfrak{p}$ est un idéal premier de $A$ (Bourbaki,
\emph{Alg.~comm.}, chap.~II, \S~4, no~3, cor.~1 de la prop.~14), il revient au
même de dire que $\dim(A)$ est la borne supérieure des longueurs des chaînes
(14.1.1)
\[
  \mathfrak{p}_0\subset\mathfrak{p}_1\subset\cdots\subset\mathfrak{p}_n
\]
d'idéaux premiers de $A$. Si $A\neq0$, on a $\dim A\geq0$. Un anneau artinien
$\neq0$ (en particulier un corps) est de dimension 0 (Bourbaki,
\emph{Alg.~comm.}, chap.~IV, \S~2, no~5, prop.~7)~; un anneau de Dedekind qui
n'est pas un corps (en particulier $\mathbb{Z}$ ou un anneau $k[T]$ de
polynômes à une indéterminée sur un corps $k$) est de dimension 1 (Bourbaki,
\emph{Alg.~comm.}, chap.~VII, \S~2). Un anneau noethérien n'est pas
nécessairement de dimension finie [30].

Si $(\mathfrak{p}_\alpha)$ est la famille des idéaux premiers minimaux de $A$,
on a (14.1.2.1)
\[
\label{0.16.1.1.1-fr}
  \dim(A)=\sup_\alpha(\dim(A/\mathfrak{p}_\alpha)).
\tag{16.1.1.1}
\]
\end{env}

\begin{env}[16.1.2]
\label{0.16.1.2-fr}
Pour tout idéal $\mathfrak{a}$ de $A$, les idéaux premiers de
$A/\mathfrak{a}$ correspondent biunivoquement aux idéaux premiers de $A$
contenant $\mathfrak{a}$, donc
\[
\label{0.16.1.2.1-fr}
  \dim(A/\mathfrak{a})\leq\dim(A).
\tag{16.1.2.1}
\]
Si en outre $\mathfrak{a}$ n'est contenu dans aucun idéal premier minimal
$\mathfrak{p}_\alpha$ de $A$, toute chaîne d'idéaux premiers de $A$ contenant
$\mathfrak{a}$ peut être complétée par un des $\mathfrak{p}_\alpha$, donc on a
\[
\label{0.16.1.2.2-fr}
  1+\dim(A/\mathfrak{a})\leq\dim(A).
\tag{16.1.2.2}
\]
\end{env}

\begin{env}[16.1.3]
\label{0.16.1.3-fr}
\oldpage[0\textsubscript{IV-1}]{23}
Si $S$ est une partie multiplicative de $A$, les idéaux premiers de $S^{-1}A$
correspondent biunivoquement aux idéaux premiers de $A$ ne rencontrant pas
$S$, donc on a
\[
\label{0.16.1.3.1-fr}
  \dim(S^{-1}A)\leq\dim(A).
\tag{16.1.3.1}
\]
Pour tout idéal premier $\mathfrak{p}$ de $A$, on a, par définition (14.2.1)
\[
\label{0.16.1.3.2-fr}
  \dim(A_{\mathfrak{p}})
  =\operatorname{codim}(V(\mathfrak{p}),\operatorname{Spec}(A)).
\tag{16.1.3.2}
\]
Ce nombre est encore appelé hauteur de l'idéal premier $\mathfrak{p}$ et noté
$\operatorname{ht}(\mathfrak{p})$. Plus généralement, pour tout idéal
$\mathfrak{a}$ de $A$, on appelle hauteur de $\mathfrak{a}$ le nombre
\[
\label{0.16.1.3.3-fr}
  \operatorname{ht}(\mathfrak{a})
  =\operatorname{codim}(V(\mathfrak{a}),\operatorname{Spec}(A)).
\tag{16.1.3.3}
\]
Si $(\mathfrak{p}_\lambda)$ est la famille des idéaux premiers minimaux parmi
ceux qui contiennent $\mathfrak{a}$, on a donc (14.2.1)
\[
\label{0.16.1.3.4-fr}
  \operatorname{ht}(\mathfrak{a})
  =\inf_\lambda(\operatorname{ht}(\mathfrak{p}_\lambda)).
\tag{16.1.3.4}
\]
On a évidemment
\[
\label{0.16.1.3.5-fr}
  \dim(A)=\sup_{\mathfrak{m}}(\dim(A_{\mathfrak{m}}))
\tag{16.1.3.5}
\]
où $\mathfrak{m}$ parcourt l'ensemble des idéaux maximaux de $A$.
\end{env}

\begin{env}[16.1.4]
\label{0.16.1.4-fr}
Pour tout idéal premier $\mathfrak{p}$ de $A$, on a (14.2.2.2)
\[
\label{0.16.1.4.1-fr}
  \dim(A_{\mathfrak{p}})+\dim(A/\mathfrak{p})\leq\dim A.
\tag{16.1.4.1}
\]
On dit que $A$ est caténaire (resp. équidimensionnel, équicodimensionnel,
biéquidimensionnel) lorsque $\operatorname{Spec}(A)$ est caténaire (resp.
équidimensionnel,
équicodimensionnel, biéquidimensionnel). Lorsque $A$ est noethérien et
biéquidimensionnel, les deux membres de (16.1.4.1) sont égaux pour tout idéal
premier de $A$ (14.3.5).

Pour tout idéal $\mathfrak{a}$ de $A$, les idéaux premiers de
$A/\mathfrak{a}$ correspondent biunivoquement aux idéaux premiers de $A$
contenant $\mathfrak{a}$, donc, si $A$ est caténaire, il en est de même de
$A/\mathfrak{a}$. De même, pour toute partie multiplicative $S$ de $A$, les
idéaux premiers de $S^{-1}A$ correspondent biunivoquement aux idéaux premiers
de $A$ ne rencontrant pas $S$, donc si $A$ est caténaire, il en est de même de
$S^{-1}A$.

En outre, comme les idéaux premiers de $A$ contenus dans un idéal premier
$\mathfrak{p}$ correspondent biunivoquement aux idéaux premiers de
$A_{\mathfrak{p}}$, pour que $A$ soit caténaire, il faut et il suffit que
$A_{\mathfrak{p}}$ le soit pour tout $\mathfrak{p}$. Dire que $A$ est
caténaire signifie donc ((14.3.2) et (16.1.3.2)) que pour tout couple d'idéaux
premiers $\mathfrak{p}$, $\mathfrak{q}$ de $A$ tels que
$\mathfrak{q}\subset\mathfrak{p}$, on a
\[
\label{0.16.1.4.2-fr}
  \dim(A_{\mathfrak{q}})
  +\dim(A_{\mathfrak{p}}/\mathfrak{q}A_{\mathfrak{p}})
  =\dim(A_{\mathfrak{p}}).
\tag{16.1.4.2}
\]
Si $A$ est biéquidimensionnel, il en est de même de $A_{\mathfrak{p}}$ pour
tout idéal premier $\mathfrak{p}$ de $A$.
\end{env}

\begin{proposition}[16.1.5]
\label{0.16.1.5-fr}
Soit $\rho:A\to B$ un homomorphisme d'anneaux faisant de $B$ une $A$-algèbre
entière. Alors on a $\dim(B)\leq\dim(A)$~; si de plus $\rho$ est injectif, on a
$\dim(B)=\dim(A)$.
\end{proposition}

\begin{proof}
\oldpage[0\textsubscript{IV-1}]{24}
Si $\mathfrak{n}$ est le noyau de $\rho$, $\rho(A)$ est isomorphe à
$A/\mathfrak{n}$, donc $\dim\rho(A)\leq\dim(A)$ par (16.1.2.1), et comme $B$
est entier sur $\rho(A)$, on peut se borner à démontrer la seconde assertion
lorsque $A\subset B$. Si
$\mathfrak{p}'_0\subset\mathfrak{p}'_1\subset\cdots\subset\mathfrak{p}'_n$
est une chaîne d'idéaux premiers de $B$, les $\mathfrak{p}'_i\cap A$ sont
alors des idéaux premiers distincts dans $A$ (Bourbaki, \emph{Alg.~comm.},
chap.~V, \S~2, no~1, cor.~1 de la prop.~1), donc forment une chaîne d'idéaux
premiers. Inversement, pour toute chaîne
$\mathfrak{p}_0\subset\mathfrak{p}_1\subset\cdots\subset\mathfrak{p}_m$
d'idéaux premiers de $A$, on sait qu'il existe pour chaque $i$ un idéal
premier $\mathfrak{p}'_i$ de $B$ tel que
$\mathfrak{p}'_i\cap A=\mathfrak{p}_i$ et que
$\mathfrak{p}'_i\subset\mathfrak{p}'_{i+1}$ pour $0\leq i\leq m-1$
(\emph{loc.~cit.}, cor.~2 du th.~1). D'où la conclusion.
\end{proof}

\begin{proposition}[16.1.6]
\label{0.16.1.6-fr}
Soient $B$ un anneau intègre, $A$ un sous-anneau local de $B$,
$\mathfrak{m}$ son idéal maximal. Supposons que $A$ soit intégralement clos et
que $B$ soit entier sur $A$~; alors, pour tout idéal maximal $\mathfrak{n}$ de
$B$, on a $\dim(B_{\mathfrak{n}})=\dim(A)$.
\end{proposition}

\begin{proof}
La première partie du raisonnement de (16.1.5) montre que
$\dim(B_{\mathfrak{n}})\leq\dim(A)$. Inversement, considérons une chaîne
d'idéaux premiers
$\mathfrak{p}_0\subset\mathfrak{p}_1\subset\cdots\subset
\mathfrak{p}_k=\mathfrak{m}$ de $A$~; on sait que
$\mathfrak{n}\cap A=\mathfrak{m}$ (Bourbaki, \emph{Alg.~comm.}, chap.~V,
\S~2, no~1, prop.~1) et en vertu des hypothèses faites, on peut appliquer le
second théorème de Cohen-Seidenberg (Bourbaki, \emph{Alg.~comm.}, chap.~V,
\S~2, no~4, th.~3), qui prouve l'existence d'une chaîne d'idéaux premiers
$\mathfrak{p}'_0\subset\mathfrak{p}'_1\subset\cdots\subset
\mathfrak{p}'_k=\mathfrak{n}$ de $B$ telle que
$\mathfrak{p}'_i\cap A=\mathfrak{p}_i$ pour tout $i$~; on en déduit aussitôt
$\dim(B_{\mathfrak{n}})\geq\dim(A)$, d'où la proposition.
\end{proof}

\begin{env}[16.1.7]
\label{0.16.1.7-fr}
Soit $M$ un $A$-module~; on appelle dimension de $M$ et on note
$\dim_A(M)$ ou $\dim(M)$ la dimension de l'anneau $A/\operatorname{Ann}(M)$
(isomorphe à l'anneau des homothéties $A_M$ de $M$). On a donc
$\dim(M^k)=\dim(M)$ pour tout $k>0$. Si $M$ est de type fini, les idéaux
premiers de $A$ contenant $\operatorname{Ann}(M)$ sont exactement ceux qui
appartiennent à $\operatorname{Supp}(M)$ ($0_{\mathrm I}$, 1.7.4) et ce
dernier est fermé dans $\operatorname{Spec}(A)$, donc on a alors
\[
\label{0.16.1.7.1-fr}
  \dim(M)=\dim(\operatorname{Supp}(M))\leq\dim(A).
\tag{16.1.7.1}
\]
Si $M$, $N$ sont deux $A$-modules tels que $N\subset M$, on a
\[
\label{0.16.1.7.2-fr}
  \dim(N)\leq\dim(M),\qquad \dim(M/N)\leq\dim(M).
\tag{16.1.7.2}
\]
Si $S$ est une partie multiplicative de $A$ et si $M$ est de type fini, on a
\[
\label{0.16.1.7.3-fr}
  \dim_{S^{-1}A}(S^{-1}M)\leq\dim_A(M).
\tag{16.1.7.3}
\]
En effet, on a
$\operatorname{Ann}(S^{-1}M)=S^{-1}\operatorname{Ann}(M)$ (Bourbaki,
\emph{Alg.~comm.}, chap.~II, \S~2, no~4) et
$S^{-1}A/S^{-1}\operatorname{Ann}(M)=S^{-1}(A/\operatorname{Ann}(M))$.

Par ailleurs, toute chaîne maximale d'idéaux premiers contenant
$\operatorname{Ann}(M)$ se termine par un idéal maximal $\mathfrak{m}$, donc
sa longueur est égale à celle de la chaîne des idéaux premiers correspondants
de $A_{\mathfrak{m}}$ (qui contiennent
$\operatorname{Ann}(M_{\mathfrak{m}})=(\operatorname{Ann}(M))_{\mathfrak{m}}$)~;
de cette remarque et de (16.1.7.3), on tire
\[
\label{0.16.1.7.4-fr}
  \dim_A(M)=\sup_{\mathfrak{m}}
  (\dim_{A_{\mathfrak{m}}}(M_{\mathfrak{m}}))
\tag{16.1.7.4}
\]
où $\mathfrak{m}$ parcourt l'ensemble des idéaux maximaux de $A$.

On dit que $M$ est caténaire (resp. équidimensionnel, équicodimensionnel,
biéquidimensionnel) lorsque $A/\operatorname{Ann}(M)$ l'est.
\end{env}

\begin{proposition}[16.1.8]
\label{0.16.1.8-fr}
\oldpage[0\textsubscript{IV-1}]{25}
Soient $A$ un anneau, $M$ un $A$-module de type fini, $\mathfrak{p}$ un idéal
premier de $A$. On a alors
\[
\label{0.16.1.8.1-fr}
  \dim_{A_{\mathfrak{p}}}(M_{\mathfrak{p}})
  +\dim_A(M/\mathfrak{p}M)\leq\dim_A(M).
\tag{16.1.8.1}
\]
\end{proposition}

\begin{proof}
On a
$\operatorname{Ann}(M_{\mathfrak{p}})=(\operatorname{Ann}(M))_{\mathfrak{p}}$,
donc les idéaux premiers de $A_{\mathfrak{p}}$ contenant
$\operatorname{Ann}(M_{\mathfrak{p}})$ correspondent biunivoquement aux
idéaux premiers de $A$ contenant $\operatorname{Ann}(M)$ et contenus dans
$\mathfrak{p}$. D'autre part, si $\operatorname{Ann}(M)=\mathfrak{n}$,
$V(\operatorname{Ann}(M/\mathfrak{p}M))=V(\mathfrak{p}+\mathfrak{n})$
($0_{\mathrm I}$, 1.7.5), donc les idéaux premiers de $A$ contenant
$\operatorname{Ann}(M/\mathfrak{p}M)$ contiennent
$\mathfrak{p}+\mathfrak{n}$~; d'où la conclusion.

Cette dernière remarque montre en outre que l'on a
\[
\label{0.16.1.8.2-fr}
  \dim_{A/\mathfrak{p}}(M/\mathfrak{p}M)=\dim_A(M/\mathfrak{p}M).
\tag{16.1.8.2}
\]
\end{proof}

\begin{proposition}[16.1.9]
\label{0.16.1.9-fr}
Soient $A$ un anneau noethérien, $\rho:A\to B$ un homomorphisme d'anneaux,
$M$ un $B$-module tel que $M_{[\rho]}$ soit un $A$-module de type fini. Alors
on a
\[
\label{0.16.1.9.1-fr}
  \dim_A(M_{[\rho]})=\dim_B(M).
\tag{16.1.9.1}
\]
\end{proposition}

\begin{proof}
En effet, l'anneau $B_M$ des homothéties du $B$-module $M$ s'identifie à un
sous-anneau de $C=\operatorname{End}_A(M_{[\rho]})$, et $C$ est un $A$-module
de type fini (Bourbaki, \emph{Alg.~comm.}, chap.~III, \S~3, no~1, lemme~2),
donc il en est de même de $B_M$~; d'ailleurs l'anneau $A_M$ des homothéties du
$A$-module $M_{[\rho]}$ est un sous-anneau de $B_M$, image canonique de $A$
dans $C$, donc $B_M$ est fini sur $A_M$~; la conclusion résulte donc de
(16.1.5).
\end{proof}

\begin{env}[16.1.10]
\label{0.16.1.10-fr}
Soit $A$ un anneau noethérien~; si $M$ est un $A$-module de longueur finie, on
sait (Bourbaki, \emph{Alg.~comm.}, chap.~IV, \S~2, no~5, prop.~7 et cor.~1 de
la prop.~7) que $\operatorname{Supp}(M)$ est un espace fini discret, donc
$\dim(M)=0$ si $M\neq0$~; réciproquement, si $M$ est un $A$-module de type
fini tel que $\dim(M)=0$, tout point de $\operatorname{Supp}(M)$ est fermé,
autrement dit est un idéal maximal de $A$, donc (\emph{loc.~cit.}, prop.~7)
$M$ est de longueur finie.
\end{env}

\subsection[Dimension d'un anneau semi-local noethérien]{Dimension d'un
anneau semi-local noethérien\protect\footnote{Ce numéro reproduit l'exposé
donné par Serre dans un cours au Collège de France en 1958.}}
\label{subsection:0.16.2-fr}

\begin{env}[16.2.1]
\label{0.16.2.1-fr}
Soient $A$ un anneau semi-local noethérien, $\mathfrak{r}$ son radical~;
rappelons qu'un idéal de définition $\mathfrak{q}$ de $A$ est un idéal tel que
$\mathfrak{q}\subset\mathfrak{r}$ et que $\mathfrak{q}$ contienne une
puissance de $\mathfrak{r}$~; il revient au même de dire que les seuls idéaux
premiers contenant $\mathfrak{q}$ sont maximaux, ou que $A/\mathfrak{q}$ est
un anneau artinien. Soit $M$ un $A$-module de type fini~; pour tout entier
$n>0$, $M/\mathfrak{q}^nM$ est un $A$-module de longueur finie, égale à
$\sum_{i=0}^{n-1}\operatorname{long}(\gr_i(M))$, où $\gr_\bullet(M)$ est le
$(A/\mathfrak{q})$-module gradué associé à $M$ muni de la filtration
$\mathfrak{q}$-préadique. On sait qu'il existe un polynôme $Q(n)$ à
coefficients rationnels tel que $\operatorname{long}(\gr_n(M))=Q(n)$ pour $n$
assez grand (III, 2.5.4)~; on en déduit aussitôt qu'il existe un polynôme
$P_{\mathfrak{q}}(M,n)$ en $n$, à coefficients rationnels, tel que
\[
  \operatorname{long}(M/\mathfrak{q}^nM)=P_{\mathfrak{q}}(M,n)
\]
pour $n$ assez grand~; il est clair que ce polynôme est unique et que son
coefficient dominant est $>0$ si $M\neq0$~; on dit que c'est le polynôme de
Hilbert-Samuel de $M$ pour la filtration $\mathfrak{q}$-préadique.
\end{env}

\begin{lemma}[16.2.2.1]
\label{0.16.2.2.1-fr}
\oldpage[0\textsubscript{IV-1}]{26}
Soit $(M_n)$ une filtration $\mathfrak{q}$-bonne de $M$ (Bourbaki,
\emph{Alg.~comm.}, chap.~III, \S~3, no~1)~; alors
$\operatorname{long}(M/M_n)$ est égal, pour $n$ grand, à un polynôme en $n$
dont le degré et le coefficient dominant sont les mêmes que ceux de
$P_{\mathfrak{q}}(M,n)$.
\end{lemma}

\begin{proof}
En effet, il existe $n_0$ tel que $M_{n+1}=\mathfrak{q}M_n$ pour $n\geq n_0$,
d'où les inclusions
$\mathfrak{q}^{n+n_0}M\subset M_{n+n_0}=\mathfrak{q}^nM_{n_0}
\subset\mathfrak{q}^nM\subset M_n$ pour $n$ grand, ce qui entraîne
\[
  P_{\mathfrak{q}}(M,n+n_0)\geq
  \operatorname{long}(M/M_{n+n_0})\geq
  P_{\mathfrak{q}}(M,n)\geq\operatorname{long}(M/M_n)
\]
et par suite la conclusion.
\end{proof}

Si $\mathfrak{q}'$ est un second idéal de définition, on a
$\mathfrak{q}'\supset\mathfrak{q}^m$ pour un entier $m$, donc
$P_{\mathfrak{q}'}(M,n)\leq P_{\mathfrak{q}}(M,mn)$~; par suite le degré de
$P_{\mathfrak{q}}(M,n)$ ne dépend pas de l'idéal de définition
$\mathfrak{q}$ considéré~; on le note $d(M)$.

\begin{lemma}[16.2.2.2]
\label{0.16.2.2.2-fr}
Soit $0\to M'\to M\to M''\to0$ une suite exacte de $A$-modules de type
fini. Alors le polynôme
\[
  P_{\mathfrak{q}}(M,n)-P_{\mathfrak{q}}(M',n)
  -P_{\mathfrak{q}}(M'',n)
\]
a un degré $\leq d(M')-1\leq d(M)-1$.
\end{lemma}

\begin{proof}
En effet, la filtration des
$M'_n=M'\cap\mathfrak{q}^nM'$ est $\mathfrak{q}$-bonne (Bourbaki,
\emph{Alg.~comm.}, chap.~III, \S~3, no~1, prop.~1)~; comme
\[
  \operatorname{long}(M/\mathfrak{q}^nM)
  =\operatorname{long}(M''/\mathfrak{q}^nM'')
  +\operatorname{long}(M'/M'_n),
\]
le lemme résulte aussitôt de (16.2.2.1) appliqué à $M'$.
\end{proof}

\begin{theorem}[16.2.3]
\label{0.16.2.3-fr}
\textup{(Krull-Chevalley-Samuel).}
Soient $A$ un anneau semi-local noethérien, $\mathfrak{r}$ son radical,
$M\neq0$ un $A$-module de type fini, $d(M)$ le degré du polynôme de
Hilbert-Samuel de $M$ (pour un idéal de définition quelconque de $A$)~; soit
d'autre part $s(M)$ la borne inférieure des entiers $n$ tels qu'il existe des
éléments $x_1,\ldots,x_n$ de $\mathfrak{r}$ pour lesquels
$M/(x_1M+\cdots+x_nM)$ soit de longueur finie. Alors on a
\[
  d(M)=s(M)=\dim(M).
\]
\end{theorem}

\begin{lemma}[16.2.3.1]
\label{0.16.2.3.1-fr}
Pour tout $x\in\mathfrak{r}$, soit ${}_xM$ le sous-module de $M$ formé des
éléments annulés par $x$. Alors~:
\begin{enumerate}
  \item[\rm{(i)}] On a $s(M)\leq s(M/xM)+1$.
  \item[\rm{(ii)}] Soient $\mathfrak{p}_i$ les idéaux premiers appartenant à
  $\operatorname{Supp}(M)$ et tels que
  \[
    \dim(A/\mathfrak{p}_i)=\dim(M)\quad(1\leq i\leq m).
  \]
\end{enumerate}
Si $x\notin\mathfrak{p}_i$ pour tout $i$, on a
$\dim(M/xM)\leq\dim(M)-1$.
\begin{enumerate}
  \item[\rm{(iii)}] Si $\mathfrak{q}$ est un idéal de définition de $A$, le
  polynôme
  \[
    P_{\mathfrak{q}}({}_xM,n)-P_{\mathfrak{q}}(M/xM,n)
  \]
  est de degré $\leq d(M)-1$.
\end{enumerate}
\end{lemma}

\begin{proof}
L'assertion (i) est triviale, car si $(x_i)_{1\leq i\leq m}$ est tel que pour
le module $N=M/xM$, $N/(x_1N+\cdots+x_mN)$ soit de longueur finie, il suffit
d'observer que ce module est isomorphe à
$M/(xM+x_1M+\cdots+x_mM)$.

Pour prouver (ii), on remarque que, si $\mathfrak{a}$ est l'annulateur de
$M$, les idéaux premiers qui contiennent l'annulateur de $M/xM$ sont ceux qui
contiennent $\mathfrak{a}+Ax$ ($0_{\mathrm I}$, 1.7.5)~; aucun des
$\mathfrak{p}_i$ ne peut donc contenir $\mathfrak{a}+Ax$ et par définition de
$\dim(M)$ (16.1.7) toute chaîne d'idéaux premiers contenant
$\mathfrak{a}+Ax$ a donc une longueur $\leq\dim(M)-1$.

\oldpage[0\textsubscript{IV-1}]{27}

Enfin, pour démontrer (iii), on note que l'on a deux suites exactes
\[
  0\to{}_xM\to M\xrightarrow{x}xM\to0,\qquad
  0\to xM\to M\to M/xM\to0,
\]
et l'assertion découle aussitôt du lemme (16.2.2.2).
\end{proof}

La démonstration de (16.2.3) se fait en trois étapes~:

\begin{env}[16.2.3.2]
\label{0.16.2.3.2-fr}
On a $\dim(M)\leq d(M)$.
\end{env}

Raisonnons par récurrence sur $d(M)$, le cas $d(M)=0$ étant trivial.
Supposons donc $d(M)\geq1$, et soit $\mathfrak{p}_0\in\operatorname{Supp}(M)$
tel que $\dim(A/\mathfrak{p}_0)=\dim(M)$. Cela entraîne que
$\mathfrak{p}_0$ est un élément \emph{minimal} de $\operatorname{Supp}(M)$,
donc il est associé à $M$ (Bourbaki, \emph{Alg.~comm.}, chap.~IV, \S~1,
no~4, th.~2) et $M$ contient par suite un sous-module $N$ isomorphe à
$A/\mathfrak{p}_0$ (\emph{loc.~cit.}, \S~1, no~1)~; comme $d(M)\geq d(N)$
par (16.2.2.2), on est ramené à prouver que $\dim(N)\leq d(N)$. Soit donc
$\mathfrak{p}_0\subset\mathfrak{p}_1\subset\cdots\subset\mathfrak{p}_n$ une
chaîne d'idéaux premiers distincts dans $A$, et montrons que $n\leq d(N)$.
C'est évident si $n=0$. Sinon, il existe
$x\in\mathfrak{p}_1\cap\mathfrak{r}$ tel que $x\notin\mathfrak{p}_0$, car la
relation $\mathfrak{p}_0\supset\mathfrak{p}_1\cap\mathfrak{r}$ entraînerait
$\mathfrak{p}_0\supset\mathfrak{r}$ puisque $\mathfrak{p}_0$ ne contient pas
$\mathfrak{p}_1$, donc $\mathfrak{p}_0$ serait maximal, ce qui est absurde.
On a $\mathfrak{p}_i\in\operatorname{Supp}(N/xN)$, pour $i\geq1$, donc
$n-1\leq\dim(N/xN)$ (16.1.7). Comme ${}_xN=0$ en vertu du choix de $x$, on a
$d(N/xN)\leq d(N)-1$ par (16.2.3.1, (iii)). L'hypothèse de récurrence implique
alors $d(N/xN)=\dim(N/xN)\geq n-1$, donc $n-1\leq d(N)-1$.

\begin{env}[16.2.3.3]
\label{0.16.2.3.3-fr}
On a $d(M)\leq s(M)$.
\end{env}

Soit en effet $(x_i)_{1\leq i\leq n}$ une famille d'éléments de
$\mathfrak{r}$ telle que, si l'on pose
$\mathfrak{a}=x_1A+\cdots+x_nA$, $M/\mathfrak{a}M$ soit de longueur finie.
Les idéaux associés à $M/\mathfrak{a}M$ sont alors maximaux (Bourbaki,
\emph{Alg.~comm.}, chap.~IV, \S~2, no~5, prop.~7) et comme ce sont les seuls
idéaux premiers contenant
$\mathfrak{q}=\mathfrak{a}+(\mathfrak{r}\cap\operatorname{Ann}(M))$,
$A/\mathfrak{q}$ est artinien (\emph{loc.~cit.}, prop.~9) et $\mathfrak{q}$
est un idéal de définition de $A$. Mais on a
$\mathfrak{q}^mM/\mathfrak{q}^{m+1}M=
\mathfrak{a}^mM/\mathfrak{a}^{m+1}M$, et si $M$ est engendré par $r$ éléments
$z_j$, $\mathfrak{a}^mM/\mathfrak{a}^{m+1}M$ est un $(A/\mathfrak{q})$-module
engendré par les images canoniques des
$x_1^{\nu_1}x_2^{\nu_2}\cdots x_n^{\nu_n}z_j$ où
$\sum_i\nu_i=m$, donc par $r\binom{m+n-1}{n-1}$ éléments~; sa longueur est
par suite $\leq ar\binom{m+n-1}{n-1}$, où $a$ est une constante, et on en
déduit aussitôt que le degré de $P_{\mathfrak{q}}(M,n)$ est $\leq n$~; d'où
$d(M)\leq s(M)$.

\begin{env}[16.2.3.4]
\label{0.16.2.3.4-fr}
On a $s(M)\leq\dim(M)$.
\end{env}

Notons d'abord que $n=\dim(M)$ est fini par (16.2.3.2). Raisonnons par
récurrence sur $n$~; la proposition est évidente pour $n=0$, puisque $M$ est
alors de longueur finie (16.1.10). Supposons $n\geq1$, et soient
$\mathfrak{p}_i$ $(1\leq i\leq m)$ les idéaux premiers de
$\operatorname{Supp}(M)$ tels que $\dim(A/\mathfrak{p}_i)=n$~; la définition
des $\mathfrak{p}_i$ montre que ce sont des éléments minimaux de
$\operatorname{Supp}(M)$, donc en nombre fini (Bourbaki, \emph{Alg.~comm.},
chap.~IV, \S~1, no~4, th.~2). Comme $n\geq1$, les $\mathfrak{p}_i$ ne sont pas
maximaux, et il existe donc $x\in\mathfrak{r}$ tel que
$x\notin\mathfrak{p}_i$ pour tout $i$ (Bourbaki, \emph{Alg.~comm.}, chap.~II,
\S~1, no~1, prop.~2). En vertu de (16.2.3.1), on a
$s(M)\leq s(M/xM)+1$ et $\dim(M)\geq\dim(M/xM)+1$~; l'hypothèse de récurrence
entraîne que $s(M/xM)\leq\dim(M/xM)$, d'où la conclusion.

\begin{corollary}[16.2.4]
\label{0.16.2.4-fr}
Sous les hypothèses de (16.2.3), on a
\begin{equation}
\label{0.16.2.4.1-fr}
\tag{16.2.4.1}
  \dim_{\widehat{A}}(\widehat{M})=\dim_A(M).
\end{equation}
\end{corollary}

\oldpage[0\textsubscript{IV-1}]{28}

\begin{proof}
En effet, $M/\mathfrak{r}^nM$ et
$\widehat{M}/\mathfrak{r}^n\widehat{M}$ sont isomorphes, donc
$d(\widehat{M})=d(M)$.
\end{proof}

\begin{corollary}[16.2.5]
\label{0.16.2.5-fr}
La dimension d'un anneau semi-local noethérien $A$ est finie et égale au
nombre minimum d'éléments du radical de $A$ engendrant un idéal de définition.
\end{corollary}

\begin{proof}
C'est l'égalité $s(M)=\dim(M)$ pour $M=A$.
\end{proof}

En particulier~:

\begin{corollary}[16.2.6]
\label{0.16.2.6-fr}
Soient $A$ un anneau local noethérien, $\mathfrak{m}$ son idéal maximal,
$k=A/\mathfrak{m}$ son corps résiduel~; on a
\begin{equation}
\label{0.16.2.6.1-fr}
\tag{16.2.6.1}
  \dim(A)\leq\operatorname{rg}_k(\mathfrak{m}/\mathfrak{m}^2).
\end{equation}
\end{corollary}

\begin{proof}
En effet, on sait que si les $x_i$ $(1\leq i\leq n)$ sont des éléments de
$\mathfrak{m}$ dont les classes mod.~$\mathfrak{m}^2$ forment une base du
$k$-espace vectoriel $\mathfrak{m}/\mathfrak{m}^2$, les $x_i$ engendrent
$\mathfrak{m}$ (Bourbaki, \emph{Alg.~comm.}, chap.~II, \S~3, no~2, cor.~2 de
la prop.~4).
\end{proof}

\begin{proposition}[16.2.7]
\label{0.16.2.7-fr}
Soient $k$ un corps, $B$ une $k$-algèbre graduée de type fini à degrés
positifs engendrée par ses éléments homogènes de degré 1 et telle que $B_0=k$~;
soit $E$ un $B$-module gradué de type fini, de sorte (III, 2.5.4) que, pour $n$
assez grand, $\operatorname{rg}_k(E_n)=r(n)$ est de la forme
$\Phi(n)-\Phi(n-1)$, où $\Phi$ est un polynôme en $n$ de degré $d$. Alors il
existe dans $B$ une famille $(t_i)_{1\leq i\leq d}$ de $d$ éléments homogènes
de degrés $\geq1$ tels que $E/(\sum_i t_iE)$ soit de rang fini sur $k$.
\end{proposition}

\begin{proof}
Soit $\mathfrak{q}$ l'idéal maximal $\bigoplus_{n\geq1}B_n$ de $B$, et, pour
tout $n$, soit $E'_n$ le sous-$B$-module $\bigoplus_{h\geq n+1}E_h$ de $E$.
Tout élément de $B-\mathfrak{q}$ détermine une homothétie injective dans le
$B$-module artinien $E/E'_n$, et cette homothétie est donc bijective
(Bourbaki, \emph{Alg.}, chap.~VIII, \S~2, no~2, lemme~3), ce qui entraîne que
$(E/E'_n)_{\mathfrak{q}}$ s'identifie à $E/E'_n$~; comme le corps résiduel de
$B_{\mathfrak{q}}$ est $k$, on déduit de l'hypothèse que
$\operatorname{long}_{B_{\mathfrak{q}}}((E/E'_n)_{\mathfrak{q}})$ est un
polynôme en $n$ de degré $d$. D'autre part, comme $E$ est un $B$-module de
type fini et que $B$ est engendré par ses éléments homogènes de degré 1, la
filtration $(E'_n)$ sur $E$ est $\mathfrak{q}$-bonne (Bourbaki,
\emph{Alg.~comm.}, chap.~III, \S~1, no~3, lemme~1), donc la filtration
$((E'_n)_{\mathfrak{q}})$ sur $E_{\mathfrak{q}}$ est
$\mathfrak{q}B_{\mathfrak{q}}$-bonne. Puisque $B_{\mathfrak{q}}$ est un anneau
local noethérien, on conclut de (16.2.3) et de (16.2.2.1) que l'on a
$\dim(E_{\mathfrak{q}})=d$. Raisonnons alors par récurrence sur $d$, le lemme
étant évident si $d=0$. Supposons $d\geq1$, et observons que les idéaux premiers
minimaux $\mathfrak{p}_i$ $(1\leq i\leq r)$ dans
$\operatorname{Ass}(E)$ sont \emph{gradués} (Bourbaki, \emph{Alg.~comm.},
chap.~IV, \S~3, no~1, prop.~1), et que $\mathfrak{q}$ est distinct de la
réunion des $\mathfrak{p}_i$ puisque $d\geq1$~; il y a donc un élément
\emph{homogène} $t_1\in\mathfrak{q}$ n'appartenant à aucun des
$\mathfrak{p}_i$ (Bourbaki, \emph{Alg.~comm.}, chap.~III, \S~1, no~4,
prop.~8). Comme les idéaux premiers $(\mathfrak{p}_i)_{\mathfrak{q}}$ sont les
éléments minimaux de $\operatorname{Ass}(E_{\mathfrak{q}})$ (Bourbaki,
\emph{Alg.~comm.}, chap.~IV, \S~1, no~2, prop.~5), on a
$\dim(E_{\mathfrak{q}}/t_1E_{\mathfrak{q}})=d-1$ (16.3.4). Il suffit alors
d'appliquer l'hypothèse de récurrence au $B$-module gradué $E/t_1E$ pour
obtenir la conclusion.
\end{proof}

\subsection{Systèmes de paramètres dans un anneau local noethérien}
\label{subsection:0.16.3-fr}

\begin{proposition}[16.3.1]
\label{0.16.3.1-fr}
Soient $A$ un anneau noethérien, $\mathfrak{p}$ un idéal premier de $A$, $n$ un
entier. Les conditions suivantes sont équivalentes~:
\begin{enumerate}
  \item[$a)$] $\operatorname{ht}(\mathfrak{p})=
  \dim(A_{\mathfrak{p}})\leq n$~;
  \item[$b)$] Il existe $n$ éléments $x_i$ de $A$ $(1\leq i\leq n)$ tels que
  $\mathfrak{p}$ soit minimal dans l'ensemble des idéaux premiers contenant
  l'idéal $\mathfrak{a}$ engendré par les $x_i$ (autrement dit
  $V(\mathfrak{p})$ est une composante irréductible de $V(\mathfrak{a})$).
\end{enumerate}
\end{proposition}

\begin{proof}
La condition $b)$ entraîne que $\mathfrak{a}A_{\mathfrak{p}}$ est un idéal de
définition de $A_{\mathfrak{p}}$, d'où $a)$ en vertu de (16.2.5).
Inversement, si $a)$ est vérifiée, il existe un idéal de définition
$\mathfrak{b}$ de $A_{\mathfrak{p}}$, engendré par $n$ éléments $x_i/s$, avec
$s\in A-\mathfrak{p}$. En vertu de la correspondance biunivoque entre idéaux
premiers de $A_{\mathfrak{p}}$ et idéaux premiers de $A$ contenus dans
$\mathfrak{p}$, l'idéal $\mathfrak{a}$ de $A$ engendré par les $x_i$ vérifie
$b)$.
\end{proof}

\oldpage[0\textsubscript{IV-1}]{29}

Pour $n=1$, on obtient le cas particulier~:

\begin{corollary}[16.3.2]
\label{0.16.3.2-fr}
\textup{(``Hauptidealsatz'').}
Pour qu'un idéal premier dans un anneau noethérien soit de hauteur $\leq1$, il
faut et il suffit qu'il soit minimal dans l'ensemble des idéaux premiers
contenant un idéal principal convenable.
\end{corollary}

\begin{corollary}[16.3.3]
\label{0.16.3.3-fr}
\textup{(Artin-Tate).}
Soit $A$ un anneau intègre noethérien. Les conditions suivantes sont
équivalentes~:
\begin{enumerate}
  \item[$a)$] $A$ est un anneau semi-local de dimension $\leq1$.
  \item[$b)$] Il existe $f\neq0$ dans $A$ tel que $A_f$ soit un corps.
\end{enumerate}
\end{corollary}

\begin{proof}
La condition $a)$ équivaut à dire que $A$ n'a qu'un nombre fini d'idéaux
premiers $\mathfrak{m}_i\neq0$, et que ces idéaux sont tous maximaux ($0$ étant
un idéal premier). Comme le produit des $\mathfrak{m}_i$ n'est pas nul, puisque
$A$ est intègre, un élément $f\neq0$ de ce produit appartient à tous les
$\mathfrak{m}_i$, donc $(0)$ est le seul idéal premier de l'anneau intègre
$A_f$, autrement dit $A_f$ est un corps. (On notera que cette partie de la
démonstration n'utilise pas le fait que $A$ est noethérien.) Prouvons
maintenant que $b)$ entraîne $a)$~; l'hypothèse $b)$ signifie que tout idéal
premier $\mathfrak{p}\neq0$ de $A$ contient $f$. Soient $\mathfrak{p}_i$
$(1\leq i\leq n)$ les idéaux premiers minimaux parmi ceux qui contiennent $f$
(qui sont en nombre fini puisque $A/fA$ est noethérien)~; il suffit de prouver
que les $\mathfrak{p}_i$ sont des idéaux \emph{maximaux}, car alors tout idéal
premier $\neq0$ contient nécessairement un des $\mathfrak{p}_i$, donc lui est
égal. Supposons donc qu'il existe un idéal maximal $\mathfrak{m}$ contenant un
des $\mathfrak{p}_i$ et distinct des $\mathfrak{p}_i$~; $\mathfrak{m}$ n'est
donc pas contenu dans la réunion des $\mathfrak{p}_i$ (Bourbaki,
\emph{Alg.~comm.}, chap.~II, \S~1, no~1, prop.~2), autrement dit il existe
$g\in\mathfrak{m}$ n'appartenant à aucun des $\mathfrak{p}_i$. Soit
$\mathfrak{q}$ un des idéaux premiers minimaux parmi ceux qui contiennent
$g$~; il résulte du ``Hauptidealsatz'' (16.3.2) que $\mathfrak{q}$ est de
hauteur 1, donc $\neq0$~; il contient par suite $f$, donc aussi un des
$\mathfrak{p}_i$, et comme $g\in\mathfrak{q}$ et que $g$ n'appartient à aucun
des $\mathfrak{p}_i$, $\mathfrak{q}$ est distinct des $\mathfrak{p}_i$~; il
serait donc de hauteur $\geq2$, ce qui est contradictoire.
\end{proof}

\begin{proposition}[16.3.4]
\label{0.16.3.4-fr}
Soient $A$ un anneau semi-local noethérien, $\mathfrak{r}$ son radical, $M$ un
$A$-module de type fini, $\mathfrak{p}_i$ les éléments de
$\operatorname{Supp}(M)$ tels que
$\dim(A/\mathfrak{p}_i)=\dim(M)$. Alors les $\mathfrak{p}_i$ sont des éléments
minimaux de $\operatorname{Ass}(M)$, et l'on a
$\dim(M/\mathfrak{p}_iM)=\dim(M)$. Pour tout $x\in\mathfrak{r}$, on a
\begin{equation}
\label{0.16.3.4.1-fr}
\tag{16.3.4.1}
  \dim(M/xM)\geq\dim(M)-1.
\end{equation}
En outre, pour que les deux membres de (16.3.4.1) soient égaux, il faut et il
suffit que $x$ n'appartienne à aucun des $\mathfrak{p}_i$.
\end{proposition}

\begin{proof}
La première assertion est immédiate, car les $\mathfrak{p}_i$ sont par
définition des éléments minimaux de $\operatorname{Supp}(M)$ (16.1.6), et ces
derniers sont aussi éléments minimaux de $\operatorname{Ass}(M)$ (Bourbaki,
\emph{Alg.~comm.}, chap.~IV, \S~1, no~4, th.~2). En outre,
$\operatorname{Supp}(M/\mathfrak{p}_iM)$ est l'ensemble des idéaux premiers
contenant $\mathfrak{p}_i$ ($0_{\mathrm I}$, 1.7.5), donc
$\dim(M/\mathfrak{p}_iM)=\dim(A/\mathfrak{p}_i)=\dim(M)$.

Posons $N=M/xM$. Si $x_1,\ldots,x_n$ sont des éléments de $\mathfrak{r}$ tels
que $N/(x_1N+\cdots+x_nN)$ soit de longueur finie, cela signifie que
$M/(xM+x_1M+\cdots+x_nM)$ est de longueur finie, d'où l'inégalité (16.3.4.1)
en vertu de (16.2.3).

Le fait que les deux membres de (16.3.4.1) soient égaux lorsque $x$
n'appartient à aucun des $\mathfrak{p}_i$ résulte de (16.3.4.1) et de
(16.2.3.1). Inversement, si
\oldpage[0\textsubscript{IV-1}]{30}
$\dim(M/xM)=\dim(M)-1$, aucun des $\mathfrak{p}_i$ ne peut appartenir à
$\operatorname{Supp}(M/xM)$, et les idéaux premiers de ce support sont ceux
qui contiennent $\operatorname{Ann}(M)+Ax$~; comme les $\mathfrak{p}_i$
contiennent $\operatorname{Ann}(M)$, ils ne peuvent contenir $x$.
\end{proof}

\begin{corollary}[16.3.5]
\label{0.16.3.5-fr}
Avec les notations de (16.3.4), pour tout idéal
$\mathfrak{a}\subset\mathfrak{p}_i$, on a
$\dim(M/\mathfrak{a}M)=\dim(M)$ et si l'on pose
$N=M/\mathfrak{a}M$, la relation $\dim(M/xM)=\dim(M)-1$ entraîne
$\dim(N/xN)=\dim(N)-1$.
\end{corollary}

\begin{proof}
En effet, $\operatorname{Supp}(N)$ est l'ensemble des idéaux premiers
contenant $\mathfrak{a}+\operatorname{Ann}(M)$, donc $\mathfrak{p}_i$
appartient à $\operatorname{Supp}(N)$ et comme
$\dim(N)\leq\dim(M)=\dim(A/\mathfrak{p}_i)$, on a
$\dim(A/\mathfrak{p}_i)=\dim(N)$~; en outre les idéaux premiers
$\mathfrak{q}\in\operatorname{Supp}(N)$ tels que
$\dim(A/\mathfrak{q})=\dim(N)$ sont certains des $\mathfrak{p}_j$~; si $x$
n'appartient à aucun des $\mathfrak{p}_j$, on a donc
$\dim(N/xN)=\dim(N)-1$ en vertu de (16.3.4).
\end{proof}

\begin{definition}[16.3.6]
\label{0.16.3.6-fr}
Soient $A$ un anneau semi-local noethérien, $\mathfrak{r}$ son radical, $M$ un
$A$-module de type fini, et posons $n=\dim(M)$. On appelle \emph{système de
paramètres pour $M$} tout système $(x_i)_{1\leq i\leq n}$ de $n$ éléments de
$\mathfrak{r}$ tel que $M/(x_1M+\cdots+x_nM)$ soit de longueur finie.
\end{definition}

On a vu (16.2.3) qu'il existe toujours de tels systèmes.

On a vu au cours de la démonstration de (16.2.3.3) que si
$\mathfrak{J}=\sum_{i=1}^n x_iA$ est tel que $M/\mathfrak{J}M$ soit de
longueur finie, $\mathfrak{q}=\mathfrak{J}+\operatorname{Ann}(M)$ est un
\emph{idéal de définition} de $A$, et réciproquement, s'il en est ainsi,
$M/\mathfrak{J}M$ est un module de type fini sur l'anneau artinien
$A/\mathfrak{q}$, donc est de longueur finie. Il revient donc au même de dire
que $(x_i)_{1\leq i\leq n}$ est un système de paramètres \emph{pour $M$} ou
\emph{pour $A/\operatorname{Ann}(M)$} (ou encore que leurs images dans
$A/\operatorname{Ann}(M)$ forment un système de paramètres pour cet anneau).
Si $(x_i)_{1\leq i\leq n}$ est un système de paramètres pour $M$, il en est de
même de $(x_i^k)$ pour tout entier $k\geq1$, car on peut se borner, en vertu de
ce qui précède, au cas où $M=A$, et si $\mathfrak{J}$ est un idéal de
définition de $A$, l'idéal $\sum_{i=1}^n x_i^kA$ contient
$\mathfrak{J}^{kn}$, donc est aussi un idéal de définition.

\begin{proposition}[16.3.7]
\label{0.16.3.7-fr}
Soient $A$ un anneau semi-local noethérien, $\mathfrak{r}$ son radical,
$x_1,\ldots,x_k$ des éléments de $\mathfrak{r}$, $M$ un $A$-module de type
fini. On a alors
\begin{equation}
\label{0.16.3.7.1-fr}
\tag{16.3.7.1}
  \dim(M/(x_1M+\cdots+x_kM))\geq\dim(M)-k.
\end{equation}
En outre, pour que les deux membres de (16.3.7.1) soient égaux, il faut et il
suffit qu'il existe des éléments $x_i$
$(k+1\leq i\leq n=\dim(M))$ de $\mathfrak{r}$ tels que
$(x_i)_{1\leq i\leq n}$ soit un système de paramètres pour $M$.
\end{proposition}

\begin{proof}
L'inégalité (16.3.7.1) résulte de (16.3.4.1) par récurrence sur $k$. Si les
deux membres de (16.3.7.1) sont égaux, et si $x_{k+1},\ldots,x_n$ est un
système de paramètres pour $M/(x_1M+\cdots+x_kM)$, il est clair que
$M/(x_1M+\cdots+x_kM+x_{k+1}M+\cdots+x_nM)$ est de longueur finie, donc
$(x_i)_{1\leq i\leq n}$ est un système de paramètres pour $M$~;
réciproquement, s'il existe des $x_i$ $(k+1\leq i\leq n)$ ayant cette
propriété et si l'on pose $N=M/(x_1M+\cdots+x_kM)$, il est clair que
$N/(x_{k+1}N+\cdots+x_nN)$ est de longueur finie, donc
$\dim(N)\leq n-k$, et les deux membres sont égaux en vertu de (16.3.7.1).
\end{proof}

\oldpage[0\textsubscript{IV-1}]{31}

\begin{proposition}[16.3.8]
\label{0.16.3.8-fr}
Soient $A$ un anneau semi-local noethérien, $X=\operatorname{Spec}(A)$ son
spectre, $\mathfrak{a}$ un idéal de $A$ distinct de $A$ tel que
$\operatorname{codim}(V(\mathfrak{a}),X)=r>0$. Il existe alors $r$ éléments
$x_1,\ldots,x_r$ de $\mathfrak{a}$, faisant partie d'un système de paramètres
de $A$, tels que
\[
  \operatorname{codim}(V(\mathfrak{a}),
  V(Ax_1+\cdots+Ax_r))=0.
\]
\end{proposition}

\begin{proof}
Soient $\mathfrak{p}_i$ $(1\leq i\leq m)$ les idéaux premiers minimaux de
l'anneau $A$, $\mathfrak{q}_j$ $(1\leq j\leq n)$ les idéaux premiers minimaux
parmi ceux qui contiennent $\mathfrak{a}$~; on a (14.2.1)
\[
  \operatorname{codim}(V(\mathfrak{a}),X)
  =\inf_j(\operatorname{codim}(V(\mathfrak{q}_j),X))
  =\inf_j(\dim(A_{\mathfrak{q}_j}))
\]
(16.1.3.2). L'hypothèse $r>0$ entraîne que $\mathfrak{a}$ n'est contenu dans
aucun des $\mathfrak{p}_i$~; par suite il existe $x\in\mathfrak{a}$ qui
n'appartient à aucun des $\mathfrak{p}_i$ (Bourbaki, \emph{Alg.~comm.},
chap.~II, \S~1, no~1, prop.~2). On a comme ci-dessus
\[
  \operatorname{codim}(V(\mathfrak{a}),V(Ax))
  =\inf_j(\dim((A/Ax)_{\mathfrak{q}_j/Ax}))
  =\inf_j(\dim(A_{\mathfrak{q}_j}/A_{\mathfrak{q}_j}x))~;
\]
mais comme $x/1$ n'est contenu dans aucun des idéaux premiers minimaux de
$A_{\mathfrak{q}_j}$, on a
$\dim(A_{\mathfrak{q}_j}/A_{\mathfrak{q}_j}x)
=\dim(A_{\mathfrak{q}_j})-1$ (16.3.4), d'où
$\operatorname{codim}(V(\mathfrak{a}),V(Ax))=r-1$. Pour la même raison, on a
$\dim(A/Ax)=\dim(A)-1$. Raisonnons maintenant par récurrence sur $r$~; si
$r>1$, il existe $r-1$ éléments $x'_i$ $(2\leq i\leq r)$ de $A'=A/Ax$,
faisant partie d'un système de paramètres de cet anneau, appartenant à
$\mathfrak{a}/Ax$ et tels que
\[
  \operatorname{codim}(V(\mathfrak{a}/Ax),
  V(A'x'_2+\cdots+A'x'_r))=0.
\]
Soit $(x'_i)_{r+1\leq i\leq s}$ un système d'éléments de $A'$ tel que
$(x'_i)_{2\leq i\leq s}$ soit un système de paramètres de $A'$~; pour tout
$i$ tel que $2\leq i\leq s$, soit $x_i$ un élément de la classe $x'_i$ dans
$A$, avec $x_i\in\mathfrak{a}$ pour $2\leq i\leq r$~; il est immédiat que
$A/(Ax+Ax_2+\cdots+Ax_s)$ est de longueur finie, et comme $\dim(A)=s$, on voit
que $x_1=x$ et les $x_i$ d'indices $i\geq2$ répondent aux conditions de
l'énoncé.
\end{proof}

\begin{proposition}[16.3.9]
\label{0.16.3.9-fr}
Soient $A$, $B$ deux anneaux locaux noethériens, $\mathfrak{m}$ l'idéal maximal
de $A$, $k$ son corps résiduel, $\varphi:A\to B$ un homomorphisme local.
\begin{enumerate}
  \item[\rm{(i)}] On a
  \begin{equation}
  \label{0.16.3.9.1-fr}
  \tag{16.3.9.1}
    \dim(B)\leq\dim(A)+\dim(B\otimes_A k).
  \end{equation}
  \item[\rm{(ii)}] Soient $M\neq0$ un $A$-module de type fini, $N\neq0$ un
  $B$-module de type fini. On a
  \begin{equation}
  \label{0.16.3.9.2-fr}
  \tag{16.3.9.2}
    \dim_B(M\otimes_A N)
    \leq\dim_A(M)+\dim_{B\otimes_A k}(N\otimes_A k).
  \end{equation}
\end{enumerate}
\end{proposition}

\begin{proof}
Soient $m$ la dimension de $A$, $(s_i)_{1\leq i\leq m}$ un système de
paramètres de $A$ (16.3.6), de sorte que si
$\mathfrak{a}=\sum_i s_iA$, $A/\mathfrak{a}$ est un $A$-module de longueur
finie. Alors l'anneau $A/\mathfrak{a}$ est artinien, donc l'idéal maximal
$\mathfrak{m}/\mathfrak{a}$ de cet anneau est nilpotent~; l'idéal
$\mathfrak{m}B/\mathfrak{a}B$ de $B/\mathfrak{a}B$, image de
$(\mathfrak{m}/\mathfrak{a})\otimes_A B$, est donc lui aussi nilpotent, et par
suite on a
$\dim(B/\mathfrak{a}B)=\dim(B/\mathfrak{m}B)$ (les spectres de ces deux
anneaux ayant même espace sous-jacent). D'autre part,
$B/\mathfrak{a}B=B/(s_1B+\cdots+s_mB)$, et les images des $s_i$ dans $B$
appartiennent à l'idéal maximal de $B$. Par suite (16.3.7), on a
\[
  \dim(B)\leq m+\dim(B\otimes_A k),
\]
ce qui n'est autre que (16.3.9.1).

Pour prouver (16.3.9.2), notons que si $\mathfrak{r}$ (resp.~$\mathfrak{s}$)
est l'annulateur de $M$ (resp.~$N$)
\oldpage[0\textsubscript{IV-1}]{32}
dans $A$ (resp.~$B$), on a
$\dim_A(M)=\dim(A/\mathfrak{r})$, $\dim_B(N)=\dim(B/\mathfrak{s})$ par
définition (16.1.7)~; d'autre part, $\operatorname{Supp}(M\otimes_A N)$ est un
sous-ensemble fermé de $\operatorname{Spec}(B)$ qui s'identifie (I, 9.1.13) à
\[
  \operatorname{Supp}(M)\times_{\operatorname{Spec}(A)}\operatorname{Supp}(N)
  =\operatorname{Spec}((B/\mathfrak{s})\otimes_A(A/\mathfrak{r}))
  =\operatorname{Spec}(B/(\mathfrak{s}+\mathfrak{r}B)).
\]
Comme $N\otimes_A k$ est un $B$-module de type fini, on a
$\dim_{B\otimes_A k}(N\otimes_A k)=\dim_B(N\otimes_A k)$ (16.1.9), et
$\operatorname{Supp}(N\otimes_A k)$ est égal à
$\operatorname{Spec}(B/(\mathfrak{s}+\mathfrak{m}B))$ par le même raisonnement
que ci-dessus. Appliquant (16.3.9.1) à $A/\mathfrak{r}$ et à
$B/(\mathfrak{s}+\mathfrak{r}B)$ il vient
\[
  \dim(B/(\mathfrak{s}+\mathfrak{r}B))
  \leq\dim(A/\mathfrak{r})+
  \dim(B/(\mathfrak{s}+\mathfrak{m}B)),
\]
ce qui n'est autre que (16.3.9.2).
\end{proof}

\begin{corollary}[16.3.10]
\label{0.16.3.10-fr}
Sous les hypothèses de (16.3.9), si $\mathfrak{m}B$ est un idéal de définition
de $B$, on a $\dim(B)\leq\dim(A)$~; si de plus $A$ est intègre et
$\dim(A)=\dim(B)$, $\varphi$ est injectif.
\end{corollary}

\begin{proof}
La première assertion résulte de (16.3.9) puisque alors
$\dim(B/\mathfrak{m}B)=0$. On peut d'ailleurs remplacer $A$ par
$\varphi(A)$, donc $\dim(B)\leq\dim(\varphi(A))$~; si
$\mathfrak{a}=\operatorname{Ker}(\varphi)\neq0$, et si $A$ est intègre, on a
$\dim(\varphi(A))=\dim(A/\mathfrak{a})<\dim(A)$ (16.1.2.2), donc on ne peut
alors avoir $\dim(A)=\dim(B)$ que si $\mathfrak{a}=0$.
\end{proof}

\subsection[Profondeur et coprofondeur]{Profondeur et
coprofondeur\protect\footnote{Dans la 3e partie du chap.~III, nous développons
la notion de profondeur du point de vue cohomologique et dans un cadre plus
général.}}
\label{subsection:0.16.4-fr}

\begin{proposition}[16.4.1]
\label{0.16.4.1-fr}
Soient $A$ un anneau local noethérien, $\mathfrak{m}$ son idéal maximal, $M$ un
$A$-module de type fini. Alors toute suite $M$-régulière
$(x_i)_{1\leq i\leq r}$ (15.1.7) formée d'éléments de $\mathfrak{m}$, fait
partie d'un système de paramètres pour $M$.
\end{proposition}

\begin{proof}
Raisonnons par récurrence sur $r$~; considérons le $A$-module
\[
  N=M/(x_1M+\cdots+x_{r-1}M)~;
\]
par hypothèse $x_1,\ldots,x_{r-1}$ font partie d'un système de paramètres pour
$M$, donc $\dim(N)=\dim(M)-(r-1)$ (16.3.7). Comme par hypothèse l'homothétie
de rapport $x_r$ dans $N$ est injective, $x_r$ n'appartient à aucun des idéaux
premiers associés à $N$, donc (16.3.4) on a
$\dim(N/x_rN)=\dim(N)-1$, ce qui s'écrit aussi
\[
  \dim(M/(x_1M+\cdots+x_rM))=\dim(M)-r~;
\]
la conclusion résulte donc de (16.3.7).
\end{proof}

Pour un $A$-module $M\neq0$, une suite $M$-régulière d'éléments de
$\mathfrak{m}$ a donc \emph{au plus} $\dim(M)$ éléments.

\begin{lemma}[16.4.2]
\label{0.16.4.2-fr}
Soient $A$ un anneau local noethérien, $\mathfrak{m}$ son idéal maximal,
$k=A/\mathfrak{m}$ son corps résiduel, $M$ un $A$-module. Pour qu'une suite
$M$-régulière $(x_i)_{1\leq i\leq n}$ d'éléments de $\mathfrak{m}$ soit
maximale, il faut et il suffit que l'on ait
\[
  \operatorname{Hom}_A(k,M/(x_1M+\cdots+x_nM))\neq0.
\]
\end{lemma}

\begin{proof}
Dire que $(x_i)$ est maximale signifie que, pour aucun $x\in\mathfrak{m}$,
l'homothétie de rapport $x$ dans $N=M/(x_1M+\cdots+x_nM)$ n'est injective, ou
encore que $\mathfrak{m}$ est contenu dans la réunion des idéaux premiers
associés à $N$ (Bourbaki, \emph{Alg.~comm.}, chap.~IV, \S~1, no~1, cor.~2 de
la prop.~2), donc égal à un de ces derniers, puisque c'est un idéal
\oldpage[0\textsubscript{IV-1}]{33}
maximal (Bourbaki, \emph{Alg.~comm.}, chap.~II, \S~1, no~1, prop.~2)~;
autrement dit, il y a un élément $z\in N$ dont l'annulateur est
$\mathfrak{m}$, et par suite le sous-module $Az$ de $N$ est isomorphe à
$A/\mathfrak{m}=k$~; d'où le lemme.
\end{proof}

\begin{lemma}[16.4.3]
\label{0.16.4.3-fr}
Soient $A$ un anneau local noethérien, $\mathfrak{m}$ son idéal maximal,
$k=A/\mathfrak{m}$ son corps résiduel, $M$ un $A$-module,
$(x_i)_{1\leq i\leq n}$ une suite $M$-régulière d'éléments de
$\mathfrak{m}$. Alors les $A$-modules
\[
  \operatorname{Hom}_A(k,M/(x_1M+\cdots+x_nM))
  \quad\text{et}\quad \operatorname{Ext}^n_A(k,M)
\]
sont isomorphes, et pour $n\geq1$, ils sont aussi isomorphes à
$\operatorname{Ext}^{n-1}_A(k,M/x_1M)$.
\end{lemma}

\begin{proof}
Raisonnons par récurrence sur $n$, la proposition étant évidente pour $n=0$.
Posons $N=M/x_1M$~; la suite $(x_i)_{2\leq i\leq n}$ est $N$-régulière, donc
\[
  \operatorname{Hom}_A(k,M/(x_1M+\cdots+x_nM))
  =\operatorname{Hom}_A(k,N/(x_2N+\cdots+x_nN))
\]
est isomorphe à $\operatorname{Ext}^{n-1}_A(k,N)$ en vertu de l'hypothèse de
récurrence. Considérons alors la suite exacte
$0\to M\xrightarrow{x_1}M\to M/x_1M=N\to0$~; on en déduit la suite exacte des
Ext
\begin{equation}
\label{0.16.4.3.1-fr}
\tag{16.4.3.1}
  \operatorname{Ext}^{n-1}_A(k,M)
  \to\operatorname{Ext}^{n-1}_A(k,N)
  \to\operatorname{Ext}^n_A(k,M)
  \xrightarrow{x_1}\operatorname{Ext}^n_A(k,M).
\end{equation}
En vertu de l'hypothèse de récurrence, $\operatorname{Ext}^{n-1}_A(k,M)$ est
isomorphe à
\[
  \operatorname{Hom}_A(k,M/(x_1M+\cdots+x_{n-1}M)),
\]
qui est nul, puisque $(x_i)_{1\leq i\leq n-1}$ n'est pas une suite
$M$-régulière maximale (16.4.2). D'autre part, comme
$x_1\in\mathfrak{m}$, l'homothétie de rapport $x_1$ dans le $A$-module
$\operatorname{Hom}_A(k,T)$ est nulle pour tout $A$-module $T$, donc il en est
de même de l'homothétie de rapport $x_1$ dans tout $A$-module
$\operatorname{Ext}^n_A(k,T)$~; les assertions du lemme découlent donc
aussitôt de la suite exacte (16.4.3.1).
\end{proof}

\begin{corollary}[16.4.4]
\label{0.16.4.4-fr}
Soient $A$ un anneau local noethérien, $\mathfrak{m}$ son idéal maximal,
$k=A/\mathfrak{m}$ son corps résiduel, $M$ un $A$-module de type fini
$\neq0$. Toutes les suites $M$-régulières maximales d'éléments de
$\mathfrak{m}$ ont le même nombre d'éléments, qui est le plus petit entier $n$
tel que $\operatorname{Ext}^n_A(k,M)\neq0$.
\end{corollary}

\begin{definition}[16.4.5]
\label{0.16.4.5-fr}
Soient $A$ un anneau local noethérien, $\mathfrak{m}$ son idéal maximal, $M$ un
$A$-module de type fini. On appelle \emph{profondeur} de $M$, et l'on note
$\operatorname{prof}_A M$ ou $\operatorname{prof}(M)$ la borne supérieure du
nombre d'éléments d'une suite $M$-régulière formée d'éléments de
$\mathfrak{m}$. On a donc $\operatorname{prof}(M)=+\infty$ si et seulement si
$M=0$, et, par (16.4.1)
\begin{equation}
\label{0.16.4.5.1-fr}
\tag{16.4.5.1}
  \operatorname{prof}(M)\leq\dim(M)\quad\text{si}\quad M\neq0.
\end{equation}
On a évidemment $\operatorname{prof}(M^k)=\operatorname{prof}(M)$ pour tout
$k>0$.
\end{definition}

\begin{proposition}[16.4.6]
\EGAReaderTarget{canon-ref.0.16.4.6-fr}{0.16.4.6-fr}
Soient $A$ un anneau local noethérien, $\mathfrak{m}$ son idéal maximal, $M$ un
$A$-module de type fini.
\begin{enumerate}
  \item[\rm{(i)}] Pour que $\operatorname{prof}(M)=0$, il faut et il suffit que
  $\mathfrak{m}\in\operatorname{Ass}(M)$ (ou encore qu'il existe un sous-module
  de $M$ isomorphe au corps résiduel $k=A/\mathfrak{m}$ de $A$).
  \item[\rm{(ii)}] Pour tout élément $M$-régulier $x$ appartenant à
  $\mathfrak{m}$, on a
  \begin{equation}
  \label{0.16.4.6.1-fr}
  \tag{16.4.6.1}
    \operatorname{prof}(M/xM)=\operatorname{prof}(M)-1.
  \end{equation}
  \item[\rm{(iii)}] Si $M\neq0$, on a
  \begin{equation}
  \label{0.16.4.6.2-fr}
  \tag{16.4.6.2}
    \operatorname{prof}(M)
    \leq\inf_{\mathfrak{p}\in\operatorname{Ass}(M)}\dim(A/\mathfrak{p})
    \leq\dim(M).
  \end{equation}
  \oldpage[0\textsubscript{IV-1}]{34}
  \item[\rm{(iv)}] On a
  $\operatorname{prof}_A(M)=
  \operatorname{prof}_{\widehat{A}}(\widehat{M})$.
\end{enumerate}
\end{proposition}

\begin{proof}
(i) Dire que $\operatorname{prof}(M)=0$ signifie qu'il n'existe aucun
$x\in\mathfrak{m}$ tel que l'homothétie de rapport $x$ dans $M$ soit
injective, autrement dit que tout $x\in\mathfrak{m}$ appartient à un idéal
premier associé à $M$ (Bourbaki, \emph{Alg.~comm.}, chap.~IV, \S~1, no~1,
cor.~2 de la prop.~2)~; mais $\mathfrak{m}$ ne peut être réunion des idéaux
premiers associés à $M$ que s'il est égal à l'un d'eux (Bourbaki,
\emph{Alg.~comm.}, chap.~II, \S~1, no~1, prop.~2).

(ii) Si $(x_i)_{1\leq i\leq m}$ est une suite $(M/xM)$-régulière formée
d'éléments de $\mathfrak{m}$, la suite formée de $x$ et des $x_i$ est
$M$-régulière, et il résulte de la définition qu'elle est maximale lorsque la
suite $(x_i)_{1\leq i\leq m}$ est une suite $(M/xM)$-régulière maximale, d'où
la conclusion.

(iii) Raisonnons par récurrence sur $n=\operatorname{prof}(M)$, l'assertion
étant évidente pour $n=0$. Nous utiliserons le lemme suivant~:
\end{proof}

\begin{lemma}[16.4.6.3]
\label{0.16.4.6.3-fr}
Soient $A$ un anneau noethérien, $M$ un $A$-module de type fini, $t\in A$ un
élément $M$-régulier. Alors, pour tout
$\mathfrak{p}\in\operatorname{Ass}(M)$, tout idéal premier minimal parmi ceux
qui contiennent $\mathfrak{p}+At$ est associé à $M/tM$.
\end{lemma}

\begin{proof}
On sait qu'il existe une suite exacte
\[
  0\to M'\to M\to M''\to0
\]
telle que $\operatorname{Ass}(M')=\{\mathfrak{p}\}$,
$\operatorname{Ass}(M'')=\operatorname{Ass}(M)-\{\mathfrak{p}\}$ (Bourbaki,
\emph{Alg.~comm.}, chap.~IV, \S~1, no~1, prop.~4)~; comme $t$ est
$M$-régulier, il est aussi $M'$-régulier et $M''$-régulier
(\emph{loc.~cit.}, \S~1, no~1, cor.~2 de la prop.~2). On en déduit que la
suite
\[
  0\to M'/tM'\to M/tM\to M''/tM''\to0
\]
est exacte (15.1.18). Par suite, on a
$\operatorname{Ass}(M'/tM')\subset\operatorname{Ass}(M/tM)$. Mais les idéaux
premiers de $A$ contenant $\mathfrak{p}$ sont les points de
$\operatorname{Supp}(M')$ (\emph{loc.~cit.}, \S~1, no~3, prop.~7), donc les
idéaux premiers contenant $\mathfrak{p}+At$ sont les points du support de
$M'/tM'=M'\otimes_A(A/tA)$ ($0_{\mathrm I}$, 1.7.5)~; un élément minimal de
l'ensemble de ces idéaux premiers appartient donc à
$\operatorname{Ass}(M'/tM')$ (Bourbaki, \emph{loc.~cit.}, \S~1, no~3,
cor.~1 de la prop.~7), et \emph{a fortiori} à
$\operatorname{Ass}(M/tM)$.
\end{proof}

\begin{proof}[Suite de la démonstration de (16.4.6)]
Ce lemme étant établi, revenons à la démonstration de (iii). Soit
$x\in\mathfrak{m}$ un élément $M$-régulier, de sorte que
$\operatorname{prof}(M/xM)=n-1$ par (ii)~; pour tout
$\mathfrak{p}\in\operatorname{Ass}(M)$, il résulte du lemme qu'il y a un idéal
$\mathfrak{p}'$ associé à $M/xM$ et contenant $\mathfrak{p}+xA$~; mais comme
$x$ est $M$-régulier, on a $x\notin\mathfrak{p}$, d'où
$\mathfrak{p}'\neq\mathfrak{p}$ et par suite
$\dim(A/\mathfrak{p}')\leq\dim(A/\mathfrak{p})-1$. Or l'hypothèse de
récurrence montre que $n-1\leq\dim(A/\mathfrak{p}')$~; on en conclut que
$n\leq\dim(A/\mathfrak{p})$ pour tout
$\mathfrak{p}\in\operatorname{Ass}(M)$.

(iv) On peut se borner au cas où $M\neq0$. Rappelons que
$\widehat{M}=M\otimes_A\widehat{A}$ à un isomorphisme canonique près
($0_{\mathrm I}$, 7.3.3)~; comme $\widehat{A}$ est un $A$-module plat, on a des
isomorphismes canoniques
\[
  \operatorname{Ext}^i_A(k,M)\otimes_A\widehat{A}
  \xrightarrow{\sim}
  \operatorname{Ext}^i_{\widehat{A}}
  (k,M\otimes_A\widehat{A})
\]
puisque $k=k\otimes_A\widehat{A}$ à un isomorphisme canonique près
($0_{\mathrm{III}}$, 12.3.5). L'assertion (iv) résulte donc de (16.4.4) et du
fait que $\widehat{A}$ est un $A$-module fidèlement plat
($0_{\mathrm I}$, 7.3.5).
\end{proof}

\begin{env}[16.4.6.4]
\label{0.16.4.6.4-fr}
On notera que si $\mathfrak{p}$ est un idéal premier de $A$, on peut avoir
aussi bien
$\operatorname{prof}_{A_{\mathfrak{p}}}(M_{\mathfrak{p}})
<\operatorname{prof}_A(M)$ que
$\operatorname{prof}_{A_{\mathfrak{p}}}(M_{\mathfrak{p}})
>\operatorname{prof}_A(M)$~; pour un exemple du premier cas,
\oldpage[0\textsubscript{IV-1}]{35}
il suffit de prendre pour $A$ un anneau local intègre de dimension $\geq1$~; on a
$\operatorname{prof}(A)\geq1$ mais $\operatorname{prof}(K)=0$ pour le corps
des fractions $K$ de $A$. Pour un exemple du second cas, considérons un
anneau local noethérien intègre $A_0$ de dimension $\geq2$, d'idéal maximal
$\mathfrak{m}_0$ et soit $k=A_0/\mathfrak{m}_0$ son corps résiduel~; soit
$A=A_0\oplus\mathfrak{J}$ l'extension triviale de $A_0$ par le $A_0$-module
$k$ (18.2.3), l'idéal $\mathfrak{J}$ étant $A_0$-isomorphe à $k$ (de sorte
que la multiplication dans $A$ est donnée par
$(x,z)(x',z')=(xx',xz'+x'z)$). Il est clair que $\mathfrak{J}$ est le
nilradical de $A$ et $A/\mathfrak{J}$ est isomorphe à $A_0$, de sorte que
tout idéal premier $\mathfrak{p}$ de $A$ est de la forme
$\mathfrak{p}_0\oplus\mathfrak{J}$, où $\mathfrak{p}_0$ est un idéal premier
de $A_0$~; en particulier $\mathfrak{m}=\mathfrak{m}_0\oplus\mathfrak{J}$
est l'unique idéal maximal de $A$. Comme tout élément de $\mathfrak{m}_0$
annule $\mathfrak{J}$, on a
$\operatorname{prof}(A)=\operatorname{prof}(A_{\mathfrak{m}})=0$~; au
contraire, si $\mathfrak{p}_0\neq\mathfrak{m}_0$, les éléments de
$\mathfrak{J}$ sont annulés par ceux de
$\mathfrak{m}_0-\mathfrak{p}_0$, donc $A_{\mathfrak{p}}$ est canoniquement
isomorphe à $(A_0)_{\mathfrak{p}_0}$, et comme $A_0$ est intègre,
$\operatorname{prof}(A_{\mathfrak{p}})
=\operatorname{prof}((A_0)_{\mathfrak{p}_0})\geq1$ si
$\mathfrak{p}_0\neq0$.
\end{env}

\begin{corollary}[16.4.7]
\label{0.16.4.7-fr}
Pour qu'un anneau local noethérien réduit $A$ soit de profondeur $0$, il faut
et il suffit que $A$ soit un corps.
\end{corollary}

\begin{proof}
En effet, comme l'intersection des idéaux premiers minimaux de $A$ est $0$,
$A$ n'a pas d'idéaux premiers associés immergés~; en vertu de (16.4.6, (i)),
si $\operatorname{prof}(A)=0$, l'idéal maximal $\mathfrak{m}$ de $A$ est
aussi idéal premier minimal, donc c'est le seul idéal premier de $A$, et
comme $A$ est réduit, $\mathfrak{m}=0$.
\end{proof}

\begin{proposition}[16.4.8]
\label{0.16.4.8-fr}
Soient $A$, $B$ deux anneaux locaux noethériens,
$\rho:A\to B$ un homomorphisme local, $M$ un $B$-module tel que
$M_{[\rho]}$ soit un $A$-module de type fini. Alors on a
\begin{equation}
\label{0.16.4.8.1-fr}
\tag{16.4.8.1}
  \operatorname{prof}_A(M_{[\rho]})=\operatorname{prof}_B(M).
\end{equation}
\end{proposition}

\begin{proof}
On peut se borner au cas $M\neq0$~; supposons que
$\operatorname{prof}_A(M_{[\rho]})=n$. Si $(x_i)_{1\leq i\leq n}$ est une
suite $(M_{[\rho]})$-régulière maximale formée d'éléments de l'idéal maximal
de $A$, les $\rho(x_i)$ constituent une suite $M$-régulière formée d'éléments
de l'idéal maximal de $B$, et on a donc, en posant
$N=M/(\rho(x_1)M+\cdots+\rho(x_n)M)$,
$\operatorname{prof}_B(N)=\operatorname{prof}_B(M)-n$ (16.4.6, (ii))~; de
même
$N_{[\rho]}=M_{[\rho]}/(x_1M_{[\rho]}+\cdots+x_nM_{[\rho]})$ et
$\operatorname{prof}_A(N_{[\rho]})=0$~; on est ainsi ramené à démontrer la
proposition lorsque $n=0$. Soit $k$ le corps résiduel de $A$~; comme $k$ est
un quotient de $A$, le $A$-module
$P=\operatorname{Hom}_A(k,M_{[\rho]})$ est un sous-module de
$\operatorname{Hom}_A(A,M_{[\rho]})=M_{[\rho]}$ et l'hypothèse entraîne que
$P\neq0$ (16.4.6, (i))~; en outre, $P$ est un $A$-module de type fini
(puisque $M_{[\rho]}$ est un $A$-module de type fini), donc un $k$-espace
vectoriel de type fini, étant annulé par l'idéal maximal de $A$, et finalement
un $A$-module de longueur finie. D'autre part, $P$ est aussi un sous-module du
$B$-module $M$, et puisqu'il est de longueur finie en tant que $A$-module, il
l'est \emph{a fortiori} en tant que $B$-module. Comme il est $\neq0$, il
contient un sous-$B$-module simple, c'est-à-dire isomorphe au corps résiduel de
$B$~; on conclut alors de (16.4.6, (i)) que l'on a bien
$\operatorname{prof}_B(M)=0$.
\end{proof}

\begin{definition}[16.4.9]
\label{0.16.4.9-fr}
Soient $A$ un anneau local noethérien, $M$ un $A$-module de type fini. Si
$M\neq0$, on appelle \emph{coprofondeur} de $M$ et on note
$\operatorname{coprof}_A(M)$ ou $\operatorname{coprof}(M)$, l'entier fini
$\dim_A(M)-\operatorname{prof}_A(M)\geq0$ (16.4.5.1). Si $M=0$, on pose
$\operatorname{coprof}(M)=0$.

On a $\operatorname{coprof}(M^k)=\operatorname{coprof}(M)$ pour tout $k>0$.
\end{definition}

\begin{proposition}[16.4.10]
\label{0.16.4.10-fr}
Soient $A$ un anneau local noethérien, $M$ un $A$-module de type fini.
\oldpage[0\textsubscript{IV-1}]{36}
\begin{enumerate}
  \item[(i)] Pour tout élément $M$-régulier $x$ appartenant à l'idéal maximal
  de $A$, on a
  $\operatorname{coprof}(M/xM)=\operatorname{coprof}(M)$.
  \item[(ii)] On a
  $\operatorname{coprof}_A(M)=\operatorname{coprof}_{\widehat A}(\widehat M)$.
\end{enumerate}
\end{proposition}

\begin{proof}
En effet, (i) résulte de (16.3.4) et (16.4.6, (ii)), et (ii) de (16.2.4) et
(16.4.6, (iv)).
\end{proof}

Nous montrerons plus tard (IV, 6.11.5) que pour tout idéal premier
$\mathfrak{p}$ de $A$, on a
$\operatorname{coprof}_{A_{\mathfrak{p}}}(M_{\mathfrak{p}})
\leq\operatorname{coprof}_A(M)$.

\begin{proposition}[16.4.11]
\label{0.16.4.11-fr}
Sous les hypothèses de (16.4.8), on a
\begin{equation}
\label{0.16.4.11.1-fr}
\tag{16.4.11.1}
  \operatorname{coprof}_A(M_{[\rho]})=\operatorname{coprof}_B(M).
\end{equation}
\end{proposition}

\begin{proof}
Cela résulte de (16.1.9) et (16.4.8).
\end{proof}

\subsection{Modules de Cohen-Macaulay}
\label{subsection:0.16.5-fr}

\begin{definition}[16.5.1]
\label{0.16.5.1-fr}
Soit $A$ un anneau local noethérien. On dit qu'un $A$-module de type fini $M$
est un \emph{module de Cohen-Macaulay} si
$\operatorname{coprof}(M)=0$, autrement dit si $M=0$ ou si $M\neq0$ et
$\dim(M)=\operatorname{prof}(M)$. On dit que $A$ est un
\emph{anneau de Cohen-Macaulay} si le $A$-module $A$ est un module de
Cohen-Macaulay.
\end{definition}

Un $A$-module de longueur finie est un module de Cohen-Macaulay puisqu'il est
de dimension $0$ (16.1.10). Un anneau local réduit $A$ de dimension $1$ est un
anneau de Cohen-Macaulay, car alors son idéal maximal $\mathfrak{m}$ ne peut
être associé à $A$ (Bourbaki, \emph{Alg.~comm.}, chap.~IV, \S~2, no~5,
prop.~10), donc $\operatorname{prof}(A)=1$ en vertu de (16.4.6, (i)). Si $A$
est un \emph{anneau local intègre et intégralement clos de dimension $\geq2$},
on a $\operatorname{prof}(A)\geq2$~: en effet, si $x\neq0$ est un élément de
$\mathfrak{m}$, on sait (Bourbaki, \emph{Alg.~comm.}, chap.~VII, \S~1, no~4,
prop.~8) que les idéaux premiers associés à $A/xA$ sont de hauteur $1$, donc
distincts de $\mathfrak{m}$, et leur réunion est par suite distincte de
$\mathfrak{m}$ (Bourbaki, \emph{Alg.~comm.}, chap.~II, \S~1, no~1, prop.~2)~;
il existe donc des suites $(x,y)$ qui sont $A$-régulières et formées
d'éléments de $\mathfrak{m}$, d'où notre assertion. On en conclut qu'un anneau
local intégralement clos de dimension $2$ est un anneau de Cohen-Macaulay.

\begin{proposition}[16.5.2]
\label{0.16.5.2-fr}
Soient $A$ un anneau local noethérien, $M$ un $A$-module de type fini. Pour que
$M$ soit un module de Cohen-Macaulay, il faut et il suffit que $\widehat M$
soit un $\widehat A$-module de Cohen-Macaulay.
\end{proposition}

\begin{proof}
Cela résulte aussitôt de (16.4.10, (ii)).
\end{proof}

\begin{proposition}[16.5.3]
\label{0.16.5.3-fr}
Sous les hypothèses de (16.4.8), pour que $M$ soit un $B$-module de
Cohen-Macaulay, il faut et il suffit que $M_{[\rho]}$ soit un $A$-module de
Cohen-Macaulay.
\end{proposition}

\begin{proof}
Cela résulte de (16.4.11).
\end{proof}

\begin{proposition}[16.5.4]
\label{0.16.5.4-fr}
Soient $A$ un anneau local noethérien, $M$ un $A$-module de type fini $\neq0$.
Si $M$ est un $A$-module de Cohen-Macaulay, on a
$\dim(A/\mathfrak{p})=\dim(M)$ pour tout idéal premier
$\mathfrak{p}\in\operatorname{Ass}(M)$~; en particulier aucun idéal premier
associé à $M$ n'est immergé.
\end{proposition}

\begin{proof}
Cela résulte aussitôt des inégalités (16.4.6.2).
\end{proof}

On peut encore dire que $M$ (ou $\operatorname{Supp}(M)$) est
\emph{équidimensionnel} (16.1.7).

\oldpage[0\textsubscript{IV-1}]{37}

\begin{proposition}[16.5.5]
\label{0.16.5.5-fr}
Soient $A$ un anneau local noethérien, $\mathfrak{m}$ son idéal maximal, $M$
un $A$-module de type fini $\neq0$, $x$ un élément de $\mathfrak{m}$.
Supposons que $M$ soit un module de Cohen-Macaulay~; alors les conditions
suivantes sont équivalentes~:
\begin{enumerate}
  \item[a)] $x$ est $M$-régulier~;
  \item[b)] $\dim(M/xM)=\dim(M)-1$~;
  \item[c)] $x$ n'appartient à aucun élément minimal de
  $\operatorname{Supp}(M)$.
\end{enumerate}
En outre, $M/xM$ est alors un module de Cohen-Macaulay.
\end{proposition}

\begin{proof}
On a vu (16.5.4) que tous les éléments $\mathfrak{p}$ de
$\operatorname{Ass}(M)$ sont minimaux dans $\operatorname{Supp}(M)$ et que
$\dim(A/\mathfrak{p})=\dim(M)$~; l'équivalence de a) et c) résulte donc de
Bourbaki, \emph{Alg.~comm.}, chap.~IV, \S~1, no~1, cor.~2 de la prop.~2~;
l'équivalence de b) et c) résulte de (16.3.4). Le fait que $M/xM$ soit un
module de Cohen-Macaulay résulte enfin de (16.4.10, (i)).
\end{proof}

\begin{corollary}[16.5.6]
\label{0.16.5.6-fr}
Sous les hypothèses de (16.5.5), soit $(x_i)_{1\leq i\leq r}$ une suite
d'éléments de $\mathfrak{m}$. Les conditions suivantes sont équivalentes~:
\begin{enumerate}
  \item[a)] La suite $(x_i)$ est $M$-régulière.
  \item[b)] On a
  $\dim(M/(x_1M+\cdots+x_rM))=\dim(M)-r$.
  \item[c)] Les $x_i$ font partie d'un système de paramètres de $M$.
\end{enumerate}
En outre, lorsque ces conditions sont satisfaites,
$N=M/(x_1M+\cdots+x_rM)$ est un module de Cohen-Macaulay, donc, pour tout
idéal $\mathfrak{p}\in\operatorname{Ass}(N)$, on a
$\dim(A/\mathfrak{p})=\dim(M)-r$.
\end{corollary}

\begin{proof}
On sait déjà que b) et c) sont équivalentes (16.3.6) et que a) entraîne c)
(16.4.1) sans hypothèse sur $M$. Pour voir que c) implique a) et prouver la
dernière assertion de l'énoncé, raisonnons par récurrence sur $r$, l'assertion
étant triviale pour $r=0$. Puisque $x_1$ fait partie d'un système de
paramètres, $\dim(M/x_1M)=\dim(M)-1$ (16.3.6), donc $x_1$ est $M$-régulier et
$M/x_1M$ est un module de Cohen-Macaulay par (16.5.5). Comme la suite
$(x_i)_{2\leq i\leq r}$ fait partie d'un système de paramètres pour
$P=M/x_1M$, l'hypothèse de récurrence montre que $(x_i)_{2\leq i\leq r}$ est
une suite $P$-régulière et que
$P/(x_2P+\cdots+x_rP)=N$ est un module de Cohen-Macaulay~; en outre
$(x_i)_{1\leq i\leq r}$ est $M$-régulière.
\end{proof}

\begin{remark}[16.5.7]
\label{0.16.5.7-fr}
Il résulte de (16.5.6) qu'il y a identité entre les \emph{systèmes de
paramètres pour $M$} et les \emph{suites $M$-régulières maximales d'éléments de
$\mathfrak{m}$} lorsque $M$ est un module de Cohen-Macaulay. Inversement, il
résulte de (16.4.1) que si un $A$-module de type fini non nul est tel qu'il y
ait une suite $M$-régulière qui soit un système de paramètres pour $M$, alors
$M$ est un module de Cohen-Macaulay. La dernière propriété énoncée dans
(16.5.6) \emph{caractérise aussi} les modules de Cohen-Macaulay~:
\end{remark}

\begin{proposition}[16.5.8]
\label{0.16.5.8-fr}
Soient $A$ un anneau local noethérien, $M$ un $A$-module de type fini. On
suppose que pour toute suite $(x_i)_{1\leq i\leq r}$ d'éléments de $A$
faisant partie d'un système de paramètres pour $M$, et tout idéal premier
$\mathfrak{p}$ associé à
$N=M/(x_1M+\cdots+x_rM)$, on ait
$\dim(A/\mathfrak{p})=\dim(N)$. Alors $M$ est un module de Cohen-Macaulay.
\end{proposition}

\begin{proof}
Soit $(y_i)_{1\leq i\leq n}$ un système de paramètres pour $M$~; il suffit de
montrer que la suite $(y_i)$ est $M$-régulière. Raisonnons par récurrence sur
$n=\dim(M)$, la proposition étant triviale pour $n=0$. Par hypothèse
(appliquée avec $r=0$) les idéaux premiers $\mathfrak{p}_j$ associés à $M$
sont tous tels que $\dim(A/\mathfrak{p}_j)=\dim(M)$~; comme $y_1$ fait partie
d'un système de paramètres
\oldpage[0\textsubscript{IV-1}]{38}
pour $M$, il résulte de (16.3.7) et (16.3.4) que $y_1$ n'appartient à aucun
des $\mathfrak{p}_j$, donc est $M$-régulier. Si on pose $P=M/y_1M$,
$(y_i)_{2\leq i\leq n}$ est un système de paramètres pour $P$, et il est
immédiat que $P$ vérifie l'hypothèse de l'énoncé~; appliquant l'hypothèse de
récurrence, on voit donc que $(y_i)_{2\leq i\leq n}$ est une suite
$P$-régulière, ce qui démontre la proposition.
\end{proof}

\begin{proposition}[16.5.9]
\label{0.16.5.9-fr}
Soient $A$ un anneau local noethérien, $M$ un $A$-module de type fini~; on
suppose que $M$ soit un module de Cohen-Macaulay. Soit
$\mathfrak{p}\in\operatorname{Supp}(M)$, et posons
$r=\dim(M)-\dim(M/\mathfrak{p}M)$. Alors il existe une suite $M$-régulière
$(x_i)_{1\leq i\leq r}$ formée d'éléments de $\mathfrak{p}$~; pour toute
suite ayant ces propriétés, on a
\[
  \dim(M/(x_1M+\cdots+x_rM))
  =\dim(M/\mathfrak{p}M)=\dim(A/\mathfrak{p}),
\]
et $\mathfrak{p}$ est un élément minimal de
$\operatorname{Ass}(M/(x_1M+\cdots+x_rM))$.
\end{proposition}

\begin{proof}
Prouvons d'abord l'existence des $x_i$~; il n'y a rien à démontrer pour
$r=0$. Pour $r>0$, procédons par récurrence sur $r$. Comme pour tout idéal
premier $\mathfrak{q}_i$ associé à $M$ on a
$\dim(A/\mathfrak{q}_i)=\dim(M)$ (16.5.4), l'hypothèse
$\dim(M/\mathfrak{p}M)<\dim(M)$ entraîne que $\mathfrak{p}$ ne peut être
contenu dans un des $\mathfrak{q}_i$, ces derniers étant les idéaux premiers
minimaux parmi ceux qui contiennent $\operatorname{Ann}(M)$ (16.1.2.2 et
16.1.7)~; on en conclut que $\mathfrak{p}$ n'est pas non plus contenu dans la
réunion des $\mathfrak{q}_i$ (Bourbaki, \emph{Alg.~comm.}, chap.~II, \S~1,
no~1, prop.~2), et par suite il existe un $x_1\in\mathfrak{p}$ qui est
$M$-régulier (Bourbaki, \emph{Alg.~comm.}, chap.~IV, \S~1, no~1, cor.~2 de
la prop.~2). On déduit donc de (16.5.5) que $N=M/x_1M$ est un module de
Cohen-Macaulay, de dimension $\dim(M)-1$~; comme $x_1\in\mathfrak{p}$, on a
d'ailleurs $N/\mathfrak{p}N=M/\mathfrak{p}M$, d'où
$\dim(N)-\dim(N/\mathfrak{p}N)=r-1$~; on peut par suite appliquer à $N$
l'hypothèse de récurrence, et si $(x_i)_{2\leq i\leq r}$ est une suite
$N$-régulière formée d'éléments de $\mathfrak{p}$, il est clair que
$(x_i)_{1\leq i\leq r}$ est une suite $M$-régulière formée d'éléments de
$\mathfrak{p}$.

Si on pose $P=M/(x_1M+\cdots+x_rM)$, on a
$P/\mathfrak{p}P=M/\mathfrak{p}M$ puisque les $x_i$ sont dans
$\mathfrak{p}$, et $\dim(P)=\dim(P/\mathfrak{p}P)$ par (16.5.6)~; mais comme
$P$ est un module de Cohen-Macaulay, on a aussi
$\dim(P)=\dim(A/\mathfrak{p}')$ pour tout
$\mathfrak{p}'\in\operatorname{Ass}(P)$ (16.5.4). Comme
$\mathfrak{p}\in\operatorname{Supp}(P)$ ($0_{\mathrm I}$, 1.7.5),
$\mathfrak{p}$ contient un des idéaux
$\mathfrak{p}'\in\operatorname{Ass}(P)$ (Bourbaki, \emph{Alg.~comm.},
chap.~IV, \S~1, no~4, th.~2), et comme
\[
  \dim(P)=\dim(P/\mathfrak{p}P)
  \leq\dim(A/\mathfrak{p})
  \leq\dim(A/\mathfrak{p}')=\dim(P),
\]
on a nécessairement $\mathfrak{p}=\mathfrak{p}'$ (16.1.2.2), ce qui achève
la démonstration.
\end{proof}

\begin{corollary}[16.5.10]
\label{0.16.5.10-fr}
Sous les mêmes hypothèses que dans (16.5.9)~:
\begin{enumerate}
  \item[(i)] $M_{\mathfrak{p}}$ est un $A_{\mathfrak{p}}$-module de
  Cohen-Macaulay.
  \item[(ii)] On a
  \begin{equation}
  \label{0.16.5.10.1-fr}
  \tag{16.5.10.1}
    \dim(M)=\dim_A(M/\mathfrak{p}M)
    +\dim_{A_{\mathfrak{p}}}(M_{\mathfrak{p}}).
  \end{equation}
\end{enumerate}
\end{corollary}

\begin{proof}
Avec les notations de (16.5.9), soient $x_i'$ les images canoniques des $x_i$
dans $A_{\mathfrak{p}}$ $(1\leq i\leq r)$~; comme les $x_i'$ appartiennent à
l'idéal maximal de $A_{\mathfrak{p}}$ et forment une suite
$M_{\mathfrak{p}}$-régulière par platitude (15.1.14), on a
\[
  \operatorname{prof}_{A_{\mathfrak{p}}}(M_{\mathfrak{p}})
  \geq r=\dim(M)-\dim(M/\mathfrak{p}M)
  \geq\dim_{A_{\mathfrak{p}}}(M_{\mathfrak{p}})
\]
(en vertu de (16.1.8.1)). Mais compte tenu de (16.4.5.1), les trois termes de
cette inégalité sont nécessairement égaux, d'où le corollaire.
\end{proof}

\oldpage[0\textsubscript{IV-1}]{39}

\begin{corollary}[16.5.11]
\label{0.16.5.11-fr}
Soit $A$ un anneau local noethérien~; supposons qu'il existe un $A$-module de
type fini $M$ de support $\operatorname{Spec}(A)$, qui soit un module de
Cohen-Macaulay. Alors, pour tout idéal premier $\mathfrak{p}$ de $A$ on a
\begin{equation}
\label{0.16.5.11.1-fr}
\tag{16.5.11.1}
  \dim(A)=\dim(A/\mathfrak{p})+\dim(A_{\mathfrak{p}}).
\end{equation}
\end{corollary}

\begin{proof}
En effet, $\operatorname{Supp}(M/\mathfrak{p}M)
=\operatorname{Spec}(A/\mathfrak{p})$ ($0_{\mathrm I}$, 1.7.5) et
$\operatorname{Supp}(M_{\mathfrak{p}})
=\operatorname{Spec}(A_{\mathfrak{p}})$~; la relation (16.5.11.1) est donc un
cas particulier de (16.5.10.1).
\end{proof}

\begin{corollary}[16.5.12]
\label{0.16.5.12-fr}
Tout anneau quotient d'un anneau local noethérien $A$ vérifiant les conditions
de (16.5.11) (en particulier tout anneau quotient d'un anneau local de
Cohen-Macaulay) est \emph{caténaire}.
\end{corollary}

\begin{proof}
Il suffit évidemment de prouver que $A$ lui-même est caténaire c'est-à-dire
que, pour deux idéaux premiers $\mathfrak{p}\subset\mathfrak{q}$ de $A$, on a
(16.1.4.2)
\[
  \dim(A_{\mathfrak{q}})
  =\dim(A_{\mathfrak{p}})
  +\dim(A_{\mathfrak{q}}/\mathfrak{p}A_{\mathfrak{q}}).
\]
Or, en vertu de (16.5.10, (i)) $A_{\mathfrak{q}}$ vérifie les mêmes
hypothèses que $A$, et la relation précédente n'est donc autre que
(16.5.11.1) appliquée à $A_{\mathfrak{q}}$ et à l'idéal premier
$\mathfrak{p}A_{\mathfrak{q}}$ de $A_{\mathfrak{q}}$.
\end{proof}

\begin{env}[16.5.13]
\label{0.16.5.13-fr}
Soient $A$ un anneau noethérien, $M$ un $A$-module de type fini. On dit que
$M$ est un \emph{$A$-module de Cohen-Macaulay} si, pour tout idéal premier
$\mathfrak{p}$ de $A$, $M_{\mathfrak{p}}$ est un $A_{\mathfrak{p}}$-module
de Cohen-Macaulay~; en vertu de (16.5.10) cette définition coïncide avec
(16.5.1) lorsque $A$ est \emph{local}. On dit que $A$ est un
\emph{anneau de Cohen-Macaulay} si tous les $A_{\mathfrak{p}}$ le sont. Il est
clair que si $M$ (resp. $A$) est un $A$-module de Cohen-Macaulay (resp. un
anneau de Cohen-Macaulay), $S^{-1}M$ est un $S^{-1}A$-module de Cohen-Macaulay
(resp. $S^{-1}A$ est un anneau de Cohen-Macaulay) pour toute partie
multiplicative $S$ de $A$.
\end{env}

% FIN SOURCE ega4/ega0-16-fr.tex

%
% DEBUT SOURCE ega4/ega0-17-fr.tex
\section{Anneaux réguliers}
\label{section:0.17-fr}

\subsection{Définition des anneaux réguliers}
\label{subsection:0.17.1-fr}

\begin{proposition}[17.1.1]
\label{0.17.1.1-fr}
Soient $A$ un anneau local noethérien de dimension $n$, $\mathfrak{m}$ son
idéal maximal, $k=A/\mathfrak{m}$ son corps résiduel. Les conditions suivantes
sont équivalentes~:
\begin{enumerate}
  \item[\rm a)] L'homomorphisme canonique surjectif de $k$-modules gradués
  (15.1.1.1)
  \[
    \varphi:\mathbf{S}^*_k(\mathfrak{m}/\mathfrak{m}^2)
    \longrightarrow\operatorname{gr}^*_{\mathfrak{m}}(A)
  \]
  (où le second membre est le $k$-module gradué associé à $A$ muni de la
  filtration $\mathfrak{m}$-préadique) est bijectif.
  \item[\rm b)] On a
  $\operatorname{rg}_k(\mathfrak{m}/\mathfrak{m}^2)=n$.
  \item[\rm c)] L'idéal $\mathfrak{m}$ admet un système de générateurs de $n$
  éléments.
  \item[\rm d)] L'idéal $\mathfrak{m}$ admet un système de générateurs qui est
  une suite $A$-régulière.
\end{enumerate}
\end{proposition}


\begin{proof}
En vertu du lemme de Nakayama,
$\operatorname{rg}_k(\mathfrak{m}/\mathfrak{m}^2)$ est le plus petit nombre
d'éléments d'un système de générateurs de $\mathfrak{m}$, donc b) et c) sont
équivalentes. D'autre part, si $(x_i)_{1\leq i\leq r}$ est un système de
générateurs de $\mathfrak{m}$ dont les classes $\overline{x}_i$ mod.
$\mathfrak{m}^2$ forment une base du $k$-espace vectoriel
$\mathfrak{m}/\mathfrak{m}^2$, $\mathbf{S}^*(\mathfrak{m}/\mathfrak{m}^2)$
est isomorphe à l'anneau de
\oldpage[0\textsubscript{IV-1}]{40}
polynômes $k[\mathrm{T}_1,\ldots,\mathrm{T}_r]$~; compte tenu de ce que tout
$A$-module de type fini est séparé pour la filtration $\mathfrak{m}$-préadique
($0_{\mathrm I}$, 7.3.5), il résulte de (15.1.9) que les conditions a) et d)
sont équivalentes~; en outre, comme toute suite $A$-régulière d'éléments de
$\mathfrak{m}$ a au plus $n$ éléments (16.4.1), on voit que d) entraîne c).
Reste à prouver que c) entraîne a).

Posons pour abréger
$\mathrm{S}=\mathbf{S}^*(\mathfrak{m}/\mathfrak{m}^2)
=k[\mathrm{T}_1,\ldots,\mathrm{T}_n]$,
$\mathrm{G}=\operatorname{gr}^*_{\mathfrak{m}}(A)$, et considérons la suite
exacte
$0\longrightarrow\mathfrak{J}\longrightarrow\mathrm{S}
\xrightarrow{\varphi}\mathrm{G}\longrightarrow0$, où le noyau
$\mathfrak{J}$ de $\varphi$ est donc un idéal gradué de $\mathrm{S}$. Pour
tout entier $s>0$, on a donc
\begin{equation}
\label{0.17.1.1.1-fr}
\tag{17.1.1.1}
  \binom{s+n-1}{n-1}
  =\operatorname{long}(\mathrm{S}_s)
  =\operatorname{long}(\mathrm{G}_s)
  +\operatorname{long}(\mathfrak{J}_s).
\end{equation}

Supposons $\mathfrak{J}\neq0$, de sorte qu'il existe un élément homogène
$u\in\mathfrak{J}$ de degré $h>0$~; $\mathfrak{J}_s$ contient
$u\mathrm{S}_{s-h}$ pour $s\geq h$, et comme $u\neq0$ et que $\mathrm{S}$ est
intègre, $u\mathrm{S}_{s-h}$ est isomorphe à $\mathrm{S}_{s-h}$ comme
$k$-espace vectoriel~; on aurait donc
$\operatorname{long}(\mathfrak{J}_s)\geq
\operatorname{long}(\mathrm{S}_{s-h})=\binom{s-h+n-1}{n-1}$, d'où pour
$s\geq h$,
\begin{equation}
\label{0.17.1.1.2-fr}
\tag{17.1.1.2}
  \operatorname{long}(\mathrm{G}_s)
  \leq\binom{s+n-1}{n-1}-\binom{s-h+n-1}{n-1}.
\end{equation}
Or, le second membre de (17.1.1.2) est un polynôme en $s$ de degré
$\leq n-2$~; mais on a
$\mathrm{G}_s=\mathfrak{m}^s/\mathfrak{m}^{s+1}$, et pour $s$ assez grand,
$\operatorname{long}(\mathrm{G}_s)$ est un polynôme en $s$ de degré exactement
égal à $n-1$ dont le coefficient dominant est $>0$ (16.2.1), ce qui contredit
l'inégalité (17.1.1.2). C.Q.F.D.
\end{proof}

\begin{definition}[17.1.2]
\label{0.17.1.2-fr}
On dit qu'un anneau local noethérien qui vérifie les conditions équivalentes de
(17.1.1) est \emph{régulier}.
\end{definition}

\begin{corollary}[17.1.3]
\label{0.17.1.3-fr}
Un anneau local régulier est intègre, intégralement clos, et un anneau de
Cohen-Macaulay.
\end{corollary}

\begin{proof}
La dernière assertion résulte de (17.1.1, d))~; d'autre part, si $A$ est
régulier, $\operatorname{gr}^*_{\mathfrak{m}}(A)$ est intègre et complètement
intégralement clos, étant isomorphe à $k[\mathrm{T}_1,\ldots,\mathrm{T}_n]$~;
on en conclut que $A$ possède aussi ces deux propriétés (Bourbaki,
\emph{Alg.~comm.}, chap.~III, \S~2, no~3, cor. de la prop.~1 et chap.~V,
\S~1, no~5, prop.~15).
\end{proof}

\begin{examples}[17.1.4]
\label{0.17.1.4-fr}
\begin{enumerate}
  \item[\rm (i)] Un anneau local régulier de \emph{dimension} $0$, étant
  intègre par (17.1.3), est nécessairement un \emph{corps}, et réciproquement.
  \item[\rm (ii)] Pour qu'un anneau local noethérien de \emph{dimension} $1$
  soit régulier, il faut et il suffit qu'il soit un \emph{anneau de valuation
  discrète}~: en effet, dire qu'un anneau local noethérien de dimension $1$ est
  régulier signifie que son idéal maximal est \emph{principal} et la conclusion
  résulte de Bourbaki, \emph{Alg.~comm.}, chap.~VI, \S~3, no~6, prop.~9.
  \item[\rm (iii)] Soit $k$ un corps~; l'anneau de séries formelles
  $A=k[[\mathrm{T}_1,\ldots,\mathrm{T}_n]]$ est un \emph{anneau régulier de
  dimension $n$}~; en effet, il est clair que les $\mathrm{T}_i$
  $(1\leq i\leq n)$ engendrent l'idéal maximal $\mathfrak{m}$ de $A$ et
  forment une suite $A$-régulière, car
  $A/(\mathrm{T}_1A+\cdots+\mathrm{T}_iA)$ est isomorphe à
  $k[[\mathrm{T}_{i+1},\ldots,\mathrm{T}_n]]$.
\end{enumerate}
\end{examples}

\begin{proposition}[17.1.5]
\label{0.17.1.5-fr}
Pour qu'un anneau local noethérien $A$ soit régulier, il faut et il suffit que
son complété $\widehat{A}$ le soit.
\end{proposition}

\begin{proof}
\oldpage[0\textsubscript{IV-1}]{41}
En effet, l'idéal maximal de $\widehat{A}$ est $\mathfrak{m}\widehat{A}$ et on
sait que $\mathfrak{m}^h/\mathfrak{m}^{h+1}$ et
$(\mathfrak{m}\widehat{A})^h/(\mathfrak{m}\widehat{A})^{h+1}$ sont des
$k$-espaces vectoriels isomorphes~; la condition a) de (17.1.1) est donc la
même pour $A$ et $\widehat{A}$.
\end{proof}

\begin{definition}[17.1.6]
\label{0.17.1.6-fr}
Étant donné un anneau local noethérien $A$, d'idéal maximal $\mathfrak{m}$, on
appelle \emph{système régulier de paramètres de $A$} un système de paramètres
pour $A$ (16.3.5) qui engendre $\mathfrak{m}$.
\end{definition}

L'existence d'un tel système équivaut au fait que $A$ est régulier (17.1.1)~;
si $(x_i)_{1\leq i\leq n}$ est un système régulier de paramètres, c'est un
système minimal de générateurs de $\mathfrak{m}$, dont les classes mod.
$\mathfrak{m}^2$ forment une base de $\mathfrak{m}/\mathfrak{m}^2$ sur $k$~;
en vertu de (17.1.1, a)) et de (15.1.9), $(x_i)$ est une suite $A$-régulière
maximale (d'où la terminologie).

On prendra garde toutefois que, dans un anneau régulier, un système de
paramètres pour $A$ qui est aussi une suite $A$-régulière \emph{n'est pas
nécessairement} un système régulier de paramètres, comme le montre l'exemple
des puissances $k$-ièmes d'un système régulier de paramètres (15.1.20).

\begin{proposition}[17.1.7]
\label{0.17.1.7-fr}
Soient $A$ un anneau local noethérien, $\mathfrak{m}$ son idéal maximal,
$k=A/\mathfrak{m}$ son corps résiduel, $(x_i)_{1\leq i\leq r}$ une suite
d'éléments de $\mathfrak{m}$, $\mathfrak{J}$ l'idéal
$x_1A+\cdots+x_rA$. Les conditions suivantes sont équivalentes~:
\begin{enumerate}
  \item[\rm a)] $A$ est régulier et les $x_i$ font partie d'un système
  régulier de paramètres de $A$.
  \item[\rm a')] $A$ est régulier et les classes $\overline{x}_i$ des $x_i$
  mod. $\mathfrak{m}^2$ sont linéairement indépendantes sur $k$.
  \item[\rm b)] Les $x_i$ font partie d'un système de paramètres pour $A$ et
  $A/\mathfrak{J}$ est un anneau régulier.
\end{enumerate}
En outre, lorsque ces conditions sont vérifiées, $\mathfrak{J}$ est un idéal
premier.
\end{proposition}

\begin{proof}
La dernière assertion résulte trivialement du fait que $A/\mathfrak{J}$,
étant régulier, est intègre (17.1.3). Posons $n=\dim(A)$. L'équivalence de a)
et a') est immédiate~; en effet, les classes mod. $\mathfrak{m}^2$ des éléments
d'un système régulier de paramètres forment une base de
$\mathfrak{m}/\mathfrak{m}^2$ par rapport à $k$, et réciproquement, si les
$\overline{x}_i$ $(1\leq i\leq r)$ sont linéairement indépendants sur $k$, on
peut trouver $n-r$ éléments $x_i\in\mathfrak{m}$ $(r+1\leq i\leq n)$ tels que
les $\overline{x}_i$ $(1\leq i\leq n)$ forment une base de
$\mathfrak{m}/\mathfrak{m}^2$, donc $(x_i)_{1\leq i\leq n}$ est un système
régulier de paramètres. Pour voir que a) entraîne b), remarquons que puisque
les $x_i$ font partie d'un système de paramètres de $A$, on a
$\dim(A/\mathfrak{J})=n-r$ (16.3.6)~; soit
$\mathfrak{n}=\mathfrak{m}/\mathfrak{J}$ l'idéal maximal de
$A/\mathfrak{J}$. On a une suite exacte
\begin{equation}
\label{0.17.1.7.1-fr}
\tag{17.1.7.1}
  0\longrightarrow(\mathfrak{m}^2+\mathfrak{J})/\mathfrak{m}^2
  \longrightarrow\mathfrak{m}/\mathfrak{m}^2
  \longrightarrow\mathfrak{n}/\mathfrak{n}^2\longrightarrow0
\end{equation}
puisque $\mathfrak{n}^2=(\mathfrak{m}^2+\mathfrak{J})/\mathfrak{J}$~;
l'hypothèse a) entraîne que
$\operatorname{rg}_k((\mathfrak{m}^2+\mathfrak{J})/\mathfrak{m}^2)=r$, donc
$\operatorname{rg}_k(\mathfrak{n}/\mathfrak{n}^2)=n-r
=\dim(A/\mathfrak{J})$, ce qui prouve b) en vertu de (17.1.1, b)).

Inversement, pour prouver que b) entraîne a), notons que le fait que les $x_i$
font partie d'un système de paramètres entraîne que
$\dim(A/\mathfrak{J})=n-r$ (16.3.6)~; l'hypothèse que $A/\mathfrak{J}$ est
régulier implique d'autre part
$\operatorname{rg}_k(\mathfrak{n}/\mathfrak{n}^2)=n-r$ (17.1.1)~; par ailleurs
on a évidemment
$\operatorname{rg}_k((\mathfrak{m}^2+\mathfrak{J})/\mathfrak{m}^2)\leq r$,
donc on tire de (17.1.7.1) que
$\operatorname{rg}_k(\mathfrak{m}/\mathfrak{m}^2)\leq r+(n-r)=n$~; donc $A$
est régulier en vertu de (16.2.6) et (17.1.1). En outre la relation
$\operatorname{rg}_k(\mathfrak{m}/\mathfrak{m}^2)=n$ entraîne alors
$\operatorname{rg}_k((\mathfrak{m}^2+\mathfrak{J})/\mathfrak{m}^2)=r$, et par
suite les $\overline{x}_i$ sont linéairement indépendants sur $k$. C.Q.F.D.
\end{proof}

\oldpage[0\textsubscript{IV-1}]{42}

\begin{corollary}[17.1.8]
\label{0.17.1.8-fr}
Soient $A$ un anneau local noethérien, $\mathfrak{m}$ son idéal maximal, $t$ un
élément de $\mathfrak{m}$. Les conditions suivantes sont équivalentes~:
\begin{enumerate}
  \item[\rm a)] $A/tA$ est régulier et $t$ est non diviseur de zéro dans $A$.
  \item[\rm b)] $A$ est régulier et $t\notin\mathfrak{m}^2$.
\end{enumerate}
\end{corollary}

\begin{proof}
C'est le cas particulier $r=1$ de (17.1.7) (compte tenu de ce qu'un élément de
$\mathfrak{m}$ non diviseur de zéro fait partie d'un système de paramètres
(16.4.1)).
\end{proof}

\begin{corollary}[17.1.9]
\label{0.17.1.9-fr}
Soient $A$ un anneau local régulier, $\mathfrak{m}$ son idéal maximal,
$\mathfrak{J}$ un idéal de $A$ contenu dans $\mathfrak{m}$. Les conditions
suivantes sont équivalentes~:
\begin{enumerate}
  \item[\rm a)] L'anneau $A/\mathfrak{J}$ est régulier.
  \item[\rm b)] $\mathfrak{J}$ est engendré par une suite
  $(x_i)_{1\leq i\leq r}$ faisant partie d'un système régulier de paramètres
  de $A$.
\end{enumerate}
\end{corollary}

\begin{proof}
On a déjà vu (17.1.7) que b) entraîne a). Inversement, supposons que
$A/\mathfrak{J}$ soit régulier (ce qui entraîne que $\mathfrak{J}$ est premier)
et soient $n=\dim(A)$, $n-r=\dim(A/\mathfrak{J})$~; avec les notations de
(17.1.7), la suite exacte (17.1.7.1) donne
$\operatorname{rg}_k((\mathfrak{m}^2+\mathfrak{J})/\mathfrak{m}^2)=r$, donc il
existe une suite $(x_i)_{1\leq i\leq r}$ d'éléments de $\mathfrak{J}$ faisant
partie d'un système régulier de paramètres de $A$. Posons
$\mathfrak{J}'=x_1A+\cdots+x_rA$~; il résulte de (17.1.7) que $\mathfrak{J}'$
est un idéal premier de $A$ et que $A/\mathfrak{J}'$ est régulier et de
dimension $n-r$~; mais comme
$A/\mathfrak{J}=(A/\mathfrak{J}')/(\mathfrak{J}/\mathfrak{J}')$ et que
$\mathfrak{J}/\mathfrak{J}'$ est premier dans l'anneau intègre
$A/\mathfrak{J}'$, les dimensions de $A/\mathfrak{J}$ et de
$A/\mathfrak{J}'$ ne peuvent être égales que si
$\mathfrak{J}=\mathfrak{J}'$ (16.1.2.2).
\end{proof}

\subsection{Rappels sur la dimension projective et la dimension injective des
modules}
\label{subsection:0.17.2-fr}

\begin{env}[17.2.1]
\label{0.17.2.1-fr}
Soient $A$ un anneau, $M$ un $A$-module. Rappelons (M, VI, 2) qu'on appelle
\emph{dimension projective} (resp. \emph{dimension injective}) de $M$ et que
l'on note $\operatorname{dim.proj}(M)$ ou
$\operatorname{dim.proj}_A(M)$ (resp. $\operatorname{dim.inj}(M)$ ou
$\operatorname{dim.inj}_A(M)$) le plus petit $n$ (entier ou égal à $+\infty$)
tel qu'il existe une résolution projective gauche (resp. une résolution
injective droite) de $M$ qui soit \emph{de longueur $n$}. Il revient au même
(\emph{loc. cit.}) de dire que $n$ est le plus petit nombre tel que pour tout
$A$-module $N$, on ait $\operatorname{Ext}^i_A(M,N)=0$ (resp.
$\operatorname{Ext}^i_A(N,M)=0$) pour tout $i>n$, ou seulement pour $i=n+1$.
Pour tout \emph{facteur direct} $M'$ de $M$, on a donc
$\operatorname{dim.proj}(M')\leq\operatorname{dim.proj}(M)$ puisque
$\operatorname{Ext}^i_A(M',N)$ est facteur direct de
$\operatorname{Ext}^i_A(M,N)$~; de même
$\operatorname{dim.inj}(M')\leq\operatorname{dim.inj}(M)$. Une condition
équivalente à $\operatorname{dim.proj}(M)=n$ (resp.
$\operatorname{dim.inj}(M)=n$) est que $n$ est le plus petit nombre tel que,
pour toute suite exacte
\[
  0\longrightarrow R\longrightarrow P_{n-1}\longrightarrow\cdots
  \longrightarrow P_0\longrightarrow M\longrightarrow0
\]
(resp.
$0\longrightarrow M\longrightarrow Q_0\longrightarrow\cdots
\longrightarrow Q_{n-1}\longrightarrow C\longrightarrow0$), où tous les
$P_i$ pour $0\leq i<n$ sont projectifs (resp. tous les $Q_i$ pour $0\leq i<n$
sont injectifs), alors $R$ est projectif (resp. $C$ est injectif).
\end{env}

\begin{remarks}[17.2.2]
\label{0.17.2.2-fr}
\begin{enumerate}
  \item[\rm (i)] Dire que $M$ est un $A$-module \emph{projectif} (resp.
  \emph{injectif}) équivaut donc à dire que
  $\operatorname{dim.proj}(M)=0$ (resp. $\operatorname{dim.inj}(M)=0$).
  \item[\rm (ii)] Soit $T$ un \emph{foncteur covariant} additif de la catégorie
  des $A$-modules dans une catégorie abélienne. Il résulte aussitôt des
  définitions de (17.2.1) et de celle des
  \oldpage[0\textsubscript{IV-1}]{43}
  foncteurs dérivés que si $\operatorname{dim.proj}(M)\leq n$ (resp.
  $\operatorname{dim.inj}(M)\leq n$), on a $L^iT(M)=0$ (resp.
  $R^iT(M)=0$) \emph{pour tout $i>n$}.
  \item[\rm (iii)] Si on suppose $A$ \emph{noethérien} et $M$ \emph{de type
  fini}, la dernière interprétation de la dimension projective donnée dans
  (17.2.1) montre que si $\operatorname{dim.proj}(M)=n$, $M$ admet une
  résolution gauche de longueur $n$ formée de modules projectifs \emph{de type
  fini}. Si de plus $A$ est un anneau \emph{local} noethérien, $M$ admet une
  résolution de longueur $n$ par des modules \emph{libres de type fini}, un
  $A$-module projectif de type fini étant alors libre (Bourbaki,
  \emph{Alg.~comm.}, chap.~II, \S~3, no~2, cor.~2 de la prop.~5).
\end{enumerate}
\end{remarks}

\begin{lemma}[17.2.3]
\label{0.17.2.3-fr}
Soient $A$ un anneau, $M$ un $A$-module. Pour que
$\operatorname{dim.inj}(M)\leq n$, il faut et il suffit que pour tout
$A$-module monogène $N$, on ait $\operatorname{Ext}^{n+1}_A(N,M)=0$.
\end{lemma}

\begin{proof}
Avec les notations de (17.2.1), il s'agit de prouver que $C$ est injectif~;
pour tout $A$-module $N$, $\operatorname{Ext}^{n+1}_A(N,M)$ est isomorphe à
$\operatorname{Ext}^1_A(N,C)$ (M, V, 7), donc on a
$\operatorname{Ext}^1_A(N,C)=0$ pour tout $A$-module $N$ de la forme
$A/\mathfrak{J}$, où $\mathfrak{J}$ est un idéal quelconque de $A$. La suite
exacte des Ext montre alors que l'homomorphisme canonique
$\operatorname{Hom}(A,C)\longrightarrow\operatorname{Hom}(\mathfrak{J},C)$ est
surjectif pour tout idéal $\mathfrak{J}$ de $A$, ce qui implique que $C$ est un
$A$-module injectif (M, I, 3.2).
\end{proof}

\begin{lemma}[17.2.4]
\label{0.17.2.4-fr}
Soient $A$ un anneau noethérien, $M$ un $A$-module de type fini. Pour que
$\operatorname{dim.proj}(M)\leq n$, il faut et il suffit que pour tout
$A$-module monogène $N$, on ait $\operatorname{Ext}^{n+1}_A(M,N)=0$.
\end{lemma}

\begin{proof}
On sait en effet (M, VI, 2.5) que la condition
$\operatorname{dim.proj}(M)\leq n$ est équivalente à
$\operatorname{Ext}^{n+1}_A(M,N)=0$ pour tout $A$-module $N$ \emph{de type
fini}~; pour voir que la condition de l'énoncé est aussi suffisante, on
raisonne par récurrence sur le nombre de générateurs $m$ de $N$~: il y a un
sous-module $N_1$ de $N$ engendré par $m-1$ éléments et tel que
$N_2=N/N_1$ soit monogène~; de la suite exacte
$0\longrightarrow N_1\longrightarrow N\longrightarrow N_2\longrightarrow0$,
on déduit alors la suite exacte
$\operatorname{Ext}^{n+1}_A(M,N_1)\longrightarrow
\operatorname{Ext}^{n+1}_A(M,N)\longrightarrow
\operatorname{Ext}^{n+1}_A(M,N_2)$ et l'hypothèse de récurrence montre que la
condition de l'énoncé entraîne bien $\operatorname{Ext}^{n+1}_A(M,N)=0$.
\end{proof}

\begin{corollary}[17.2.5]
\label{0.17.2.5-fr}
Soient $A$ un anneau noethérien, $M$ un $A$-module. On a
\begin{equation}
\label{0.17.2.5.1-fr}
\tag{17.2.5.1}
  \operatorname{dim.inj}_A(M)
  =\sup_{\mathfrak{m}}
  \bigl(\operatorname{dim.inj}_{A_{\mathfrak{m}}}(M_{\mathfrak{m}})\bigr)
\end{equation}
et si $M$ est \emph{de type fini}
\begin{equation}
\label{0.17.2.5.2-fr}
\tag{17.2.5.2}
  \operatorname{dim.proj}_A(M)
  =\sup_{\mathfrak{m}}
  \bigl(\operatorname{dim.proj}_{A_{\mathfrak{m}}}(M_{\mathfrak{m}})\bigr)
\end{equation}
où $\mathfrak{m}$ parcourt l'ensemble des idéaux premiers (ou l'ensemble des
idéaux maximaux) de $A$.
\end{corollary}

\begin{proof}
En effet, si $N$ est un $A$-module de type fini (donc de présentation finie),
on a, pour toute partie multiplicative $S$ de $A$ et pour tout $i\geq0$
\[
  S^{-1}\operatorname{Ext}^i_A(N,M)
  =\operatorname{Ext}^i_{S^{-1}A}(S^{-1}N,S^{-1}M)
\]
par platitude, en considérant une résolution libre de $M$ et en utilisant le
fait que la relation précédente est vraie pour $i=0$ (Bourbaki,
\emph{Alg.~comm.}, chap.~II, \S~2, no~7, prop.~19). En particulier
$\operatorname{Ext}^i_{A_{\mathfrak{m}}}
(A_{\mathfrak{m}}/\mathfrak{J}A_{\mathfrak{m}},M_{\mathfrak{m}})
=(\operatorname{Ext}^i_A(A/\mathfrak{J},M))_{\mathfrak{m}}$ pour tout idéal
premier $\mathfrak{m}$ et tout idéal $\mathfrak{J}$ de $A$~; si l'on tient
compte de (17.2.3) et de ce que tout
\oldpage[0\textsubscript{IV-1}]{44}
idéal de $A_{\mathfrak{m}}$ est de la forme $\mathfrak{J}A_{\mathfrak{m}}$
pour un idéal convenable $\mathfrak{J}$ de $A$, on déduit la formule
(17.2.5.1) de ce qui précède et de Bourbaki, \emph{Alg.~comm.}, chap.~II,
\S~3, no~3, th.~1. On procède de même pour (17.2.5.2) en utilisant cette fois
(17.2.4) et en échangeant les rôles de $M$ et de $N$.
\end{proof}

Pour les anneaux noethériens et les modules de type fini, l'étude de la
dimension projective ou de la dimension injective est donc ramenée par
(17.2.5) au cas des anneaux \emph{locaux}. On a alors le

\begin{lemma}[17.2.6]
\label{0.17.2.6-fr}
Soient $A$ un anneau local noethérien, $k$ son corps résiduel, $M$ un
$A$-module de type fini. Pour que $\operatorname{dim.proj}(M)\leq n$, il est
nécessaire que $\operatorname{Tor}^A_i(M,k)=0$ pour $i>n$, et il suffit que
$\operatorname{Tor}^A_{n+1}(M,k)=0$.
\end{lemma}

\begin{proof}
La nécessité est un cas particulier de la remarque (17.2.2 (ii)), appliquée au
foncteur covariant $M\mapsto k\otimes_A M$. Pour prouver que la condition est
suffisante, il s'agit, avec les notations de (17.2.1), d'établir que $R$ est
projectif lorsque les $P_i$ sont supposés de type fini~; or
$\operatorname{Tor}^A_{n+1}(M,k)$ est isomorphe à
$\operatorname{Tor}^A_1(R,k)$ (M, V, 7)~; et on sait que, puisque $R$ est de
type fini, la condition $\operatorname{Tor}^A_1(R,k)=0$ entraîne que $R$ est
libre (Bourbaki, \emph{Alg.~comm.}, chap.~II, \S~3, no~2, cor.~2 de la
prop.~5).
\end{proof}

\begin{corollary}[17.2.7]
\label{0.17.2.7-fr}
Sous les hypothèses de (17.2.6), soit $x$ un élément $M$-régulier de $A$
appartenant à l'idéal maximal de $A$~; alors on a
\begin{equation}
\label{0.17.2.7.1-fr}
\tag{17.2.7.1}
  \operatorname{dim.proj}(M/xM)=\operatorname{dim.proj}(M)+1.
\end{equation}
\end{corollary}

\begin{proof}
En effet, de la suite exacte
$0\longrightarrow M\xrightarrow{x}M\longrightarrow M/xM\longrightarrow0$
définie par la multiplication par $x$ dans $M$, on déduit la suite exacte des
Tor~:
\[
  \operatorname{Tor}^A_i(M,k)\xrightarrow{x}
  \operatorname{Tor}^A_i(M,k)\longrightarrow
  \operatorname{Tor}^A_i(M/xM,k)\longrightarrow
  \operatorname{Tor}^A_{i-1}(M,k)\xrightarrow{x}
  \operatorname{Tor}^A_{i-1}(M,k)
\]
pour tout $i\geq1$. Or, pour tout $A$-module $N$, l'homothétie de rapport $x$
dans $M\otimes_A N$ provient aussi bien de l'homothétie de rapport $x$ dans
$M$ que de l'homothétie de rapport $x$ dans $N$~; on en conclut aussitôt,
puisque l'homothétie de rapport $x$ dans $k$ est nulle par définition, que,
pour tout $i$, l'homomorphisme
$\operatorname{Tor}^A_i(M,k)\xrightarrow{x}\operatorname{Tor}^A_i(M,k)$ est
\emph{nul}~; autrement dit, on a la suite exacte
\[
  0\longrightarrow\operatorname{Tor}^A_i(M,k)\longrightarrow
  \operatorname{Tor}^A_i(M/xM,k)\longrightarrow
  \operatorname{Tor}^A_{i-1}(M,k)\longrightarrow0.
\]
Si $\operatorname{dim.proj}(M)=n$, on a
$\operatorname{Tor}^A_n(M,k)\neq0$ et
$\operatorname{Tor}^A_{n+1}(M,k)=\operatorname{Tor}^A_{n+2}(M,k)=0$ en vertu
de (17.2.6). Il résulte donc de ce qui précède que l'on a
$\operatorname{Tor}^A_{n+1}(M/xM,k)\neq0$ et
$\operatorname{Tor}^A_{n+2}(M/xM,k)=0$, donc
$\operatorname{dim.proj}(M/xM)=n+1$ par (17.2.6).
\end{proof}

\begin{proposition}[17.2.8]
\label{0.17.2.8-fr}
\textnormal{(M.~Auslander)} Soit $A$ un anneau. Les conditions suivantes sont
équivalentes~:
\begin{enumerate}
  \item[\rm a)] Tout $A$-module est de dimension projective $\leq n$.
  \item[\rm a')] Tout $A$-module de type fini est de dimension projective
  $\leq n$.
  \item[\rm b)] Tout $A$-module est de dimension injective $\leq n$.
  \item[\rm c)] Pour tout couple de $A$-modules $M$, $N$ on a
  $\operatorname{Ext}^{n+1}_A(M,N)=0$.
  \item[\rm c')] Pour tout couple de $A$-modules $M$, $N$ tel que $M$ soit de
  type fini (ou seulement monogène), on a
  $\operatorname{Ext}^{n+1}_A(M,N)=0$.
\end{enumerate}
\end{proposition}

\begin{proof}
\oldpage[0\textsubscript{IV-1}]{45}
Cela résulte aussitôt de (17.2.1) et (17.2.3).
\end{proof}

Le plus petit nombre $n$ (entier ou $+\infty$) pour lequel les conditions
équivalentes de (17.2.8) sont satisfaites est appelé \emph{dimension
cohomologique globale} (ou simplement \emph{dimension cohomologique}) de $A$
et noté $\operatorname{dim.coh}(A)$.

\begin{proposition}[17.2.9]
\label{0.17.2.9-fr}
Soit $A$ un anneau noethérien. Les conditions suivantes sont équivalentes~:
\begin{enumerate}
  \item[\rm a)] $\operatorname{dim.coh}(A)\leq n$.
  \item[\rm b)] Tout $A$-module de type fini est de dimension injective
  $\leq n$.
  \item[\rm c)] Pour tout couple de $A$-modules de type fini, $M$, $N$, on a
  $\operatorname{Ext}^{n+1}_A(M,N)=0$.
\end{enumerate}
\end{proposition}

\begin{proof}
Cela résulte aussitôt de la définition et de (17.2.4).
\end{proof}

\begin{corollary}[17.2.10]
\label{0.17.2.10-fr}
Si $A$ est un anneau noethérien, on a
\begin{equation}
\label{0.17.2.10.1-fr}
\tag{17.2.10.1}
  \operatorname{dim.coh}(A)
  =\sup_{\mathfrak{m}}
  \bigl(\operatorname{dim.coh}(A_{\mathfrak{m}})\bigr)
\end{equation}
où $\mathfrak{m}$ parcourt le spectre de $A$ (ou l'ensemble des idéaux
maximaux de $A$).
\end{corollary}

\begin{proof}
Cela résulte de (17.2.9) et (17.2.5).
\end{proof}

\begin{proposition}[17.2.11]
\label{0.17.2.11-fr}
Soient $A$ un anneau local noethérien, $k$ son corps résiduel. Pour que
$\operatorname{dim.coh}(A)\leq n$, il est nécessaire que
$\operatorname{Tor}^A_i(k,k)=0$, pour $i>n$, et il suffit que
$\operatorname{Tor}^A_{n+1}(k,k)=0$.
\end{proposition}

\begin{proof}
Compte tenu de (17.2.6), il suffit de prouver que les relations
$\operatorname{dim.coh}(A)\leq n$ et
$\operatorname{dim.proj}_A(k)\leq n$ sont équivalentes. Il est clair que la
première entraîne la seconde par définition. Inversement, si
$\operatorname{dim.proj}_A(k)\leq n$, on a
$\operatorname{Tor}^A_i(M,k)=0$ pour $i>n$ et pour tout $A$-module $M$ en
vertu de (17.2.2, (ii))~; donc $\operatorname{dim.proj}(M)\leq n$, ce qui
prouve la proposition en vertu de (17.2.8).
\end{proof}

\begin{corollary}[17.2.12]
\label{0.17.2.12-fr}
Soit $A$ un anneau noethérien. Pour que l'on ait
$\operatorname{dim.coh}(A)\leq n$, il est nécessaire que, pour tout idéal
maximal $\mathfrak{m}$ de $A$, on ait
$\operatorname{Tor}^{A_{\mathfrak{m}}}_i(A/\mathfrak{m},A/\mathfrak{m})=0$
pour $i>n$, et il suffit que ces relations soient vérifiées pour $i=n+1$.
\end{corollary}

\begin{proof}
Cela résulte aussitôt de (17.2.11) et (17.2.10).
\end{proof}

\begin{proposition}[17.2.13]
\label{0.17.2.13-fr}
Soient $A$, $B$ deux anneaux locaux noethériens,
$\varphi:A\longrightarrow B$ un homomorphisme local faisant de $B$ un
$A$-module plat. Alors on a
\begin{equation}
\label{0.17.2.13.1-fr}
\tag{17.2.13.1}
  \operatorname{dim.coh}(A)\leq\operatorname{dim.coh}(B).
\end{equation}
\end{proposition}

\begin{proof}
Supposons que $\operatorname{dim.coh}(B)=n$ soit finie~; il suffit de prouver
que pour tout couple $(M,N)$ de $A$-modules de type fini, on a
$\operatorname{Tor}^A_i(M,N)=0$ pour $i>n$ (17.2.11). Comme $B$ est un
$A$-module fidèlement plat ($0_{\rm I}$, 6.6.2), il revient au même
($0_{\rm I}$, 6.4.1) de montrer que l'on a
\begin{equation}
\label{0.17.2.13.2-fr}
\tag{17.2.13.2}
  \operatorname{Tor}^A_i(M,N)\otimes_A B=0.
\end{equation}
Or, si $L_\bullet=(L_j)$ est une résolution droite de $M$ par des $A$-modules
libres, il résulte de ce que $B$ est un $A$-module plat que
$L_\bullet\otimes_A B=(L_j\otimes_A B)$ est une résolution droite du
$B$-module $M\otimes_A B$ par des $B$-modules libres~; en outre, on a
$(L_j\otimes_A B)\otimes_B(N\otimes_A B)
=(L_j\otimes_A N)\otimes_A B$, d'où on conclut aussitôt que le premier membre
\oldpage[0\textsubscript{IV-1}]{46}
de (17.2.13.2) est égal à
$\operatorname{Tor}^B_i(M\otimes_A B,N\otimes_A B)$~; l'hypothèse sur $B$
entraîne que ce $B$-module est nul pour $i>n$ (17.2.2, (ii)), d'où la
conclusion.
\end{proof}

\begin{env}[17.2.14]
\label{0.17.2.14-fr}
Soient $(X,\mathscr{O}_X)$ un espace annelé, $\mathscr{F}$ un
$\mathscr{O}_X$-Module~; on appelle \emph{dimension projective} (resp.
\emph{injective}) \emph{ponctuelle} de $\mathscr{F}$ et on note
$\operatorname{dim.proj}(\mathscr{F})$ (resp.
$\operatorname{dim.inj}(\mathscr{F})$) le nombre
$\sup_{x\in X}(\operatorname{dim.proj}(\mathscr{F}_x))$ (resp.
$\sup_{x\in X}(\operatorname{dim.inj}(\mathscr{F}_x))$. On appelle
\emph{dimension cohomologique ponctuelle de $X$} et on note
$\operatorname{dim.coh}(X)$ le nombre
$\sup_{x\in X}(\operatorname{dim.coh}(\mathscr{O}_x))$. Il résulte de
(17.2.5) et (17.2.10) que lorsque $A$ est un anneau noethérien, $M$ un
$A$-module de type fini et $X=\operatorname{Spec}(A)$, on a
$\operatorname{dim.proj}(\widetilde{M})=\operatorname{dim.proj}(M)$,
$\operatorname{dim.inj}(\widetilde{M})=\operatorname{dim.inj}(M)$ et
$\operatorname{dim.coh}(X)=\operatorname{dim.coh}(A)$. On appelle
\emph{dimension projective} (resp. \emph{injective}) \emph{de $\mathscr{F}$ en
un point $x\in X$} la dimension projective (resp. injective) de
$\mathscr{F}_x$, \emph{dimension cohomologique de $X$ en un point $x$} la
dimension cohomologique de $\mathscr{O}_x$.
\end{env}

\begin{proposition}[17.2.15]
\label{0.17.2.15-fr}
Soient $X$, $Y$ deux espaces annelés en anneaux locaux noethériens,
$f:X\longrightarrow Y$ un morphisme plat. Si
$\operatorname{dim.coh}(X)\leq n$, alors $Y$ est de dimension cohomologique
$\leq n$ en tout point de $f(X)$.
\end{proposition}

\begin{proof}
Cela résulte aussitôt de (17.2.13).
\end{proof}

\subsection{Théorie cohomologique des anneaux réguliers}
\label{subsection:0.17.3-fr}

\begin{theorem}[17.3.1]
\label{0.17.3.1-fr}
\textnormal{(Hilbert-Serre)} Soit $A$ un anneau local noethérien. Pour que $A$
soit de dimension cohomologique finie, il faut et il suffit que $A$ soit
régulier~; on a alors
\begin{equation}
\label{0.17.3.1.1-fr}
\tag{17.3.1.1}
  \operatorname{dim.coh}(A)=\dim(A).
\end{equation}
\end{theorem}

\begin{proof}
Supposons $A$ régulier~; soient $\mathfrak{m}$ son idéal maximal,
$\mathbf{x}=(x_i)_{1\leq i\leq n}$ un système régulier de paramètres de $A$
(17.1.6)~; considérons le complexe de l'algèbre extérieure
$K_\bullet(\mathbf{x})$ (\textbf{III}, 1.1.1), qui est formé de $A$-modules
libres, avec $K_i(\mathbf{x})=0$ pour $i>n$~; comme $(x_i)$ est une suite
régulière, on a $H_i(K_\bullet(\mathbf{x}))=0$ pour $i>0$ (\textbf{III},
1.1.4 et 1.1.3.3) et
$H_0(K_\bullet(\mathbf{x}))=A/(x_1A+\cdots+x_nA)=A/\mathfrak{m}=k$~; les
$K_i(\mathbf{x})$ forment par suite une \emph{résolution libre} de $k$ de
longueur $n$. Or, le fait que les $x_i$ appartiennent à $\mathfrak{m}$
entraîne aussitôt que dans le complexe
$K_\bullet(\mathbf{x},k)=K_\bullet(\mathbf{x})\otimes_A k$, l'opérateur bord
est nul en toutes dimensions, si bien que l'on a, par définition,
$\operatorname{Tor}^A_i(k,k)=\bigwedge^i(k^n)$~; l'égalité (17.3.1.1) résulte
donc aussitôt de (17.2.11) (ce résultat est essentiellement le «~théorème des
syzygies~» de Hilbert).

Montrons maintenant que si $A$ est un anneau local noethérien,
$\mathfrak{m}$ son idéal maximal, et si
$\operatorname{dim.proj}_A(\mathfrak{m})$ est \emph{finie}, alors $A$ est
régulier, ce qui achèvera de démontrer (17.3.1). Nous procéderons par
récurrence sur $n=\operatorname{rg}_k(\mathfrak{m}/\mathfrak{m}^2)$. Pour
$n=0$, on a $\mathfrak{m}=0$ et l'assertion est triviale.
\end{proof}

\begin{lemma}[17.3.1.2]
\label{0.17.3.1.2-fr}
\textnormal{(Nagata)} Si tout élément de
$\mathfrak{m}-\mathfrak{m}^2$ est diviseur de zéro dans $A$, il existe
$c\neq0$ dans $A$ tel que $c\mathfrak{m}=0$ (autrement dit on a
$\mathfrak{m}\in\operatorname{Ass}(A)$).
\end{lemma}

\begin{proof}
On peut se borner au cas où $\mathfrak{m}\neq0$, donc
$\mathfrak{m}\neq\mathfrak{m}^2$. L'hypothèse entraîne que
$\mathfrak{m}-\mathfrak{m}^2$ est contenu dans la réunion des idéaux
$\mathfrak{p}_i$ de $\operatorname{Ass}(A)$ (Bourbaki,
\emph{Alg.~comm.}, chap.~IV, \S~1, no~1, cor.~3 de la prop.~2)~; donc
$\mathfrak{m}$ est contenu dans la réunion de $\mathfrak{m}^2$
\oldpage[0\textsubscript{IV-1}]{47}
et des $\mathfrak{p}_i$, donc dans l'un des $\mathfrak{p}_i$, puisque
$\mathfrak{m}\not\subset\mathfrak{m}^2$ (Bourbaki, \emph{Alg.~comm.},
chap.~II, \S~1, no~1, prop.~2)~; $\mathfrak{m}$ étant maximal, cela prouve le
lemme.
\end{proof}

\begin{lemma}[17.3.1.3]
\label{0.17.3.1.3-fr}
Si $a\in\mathfrak{m}-\mathfrak{m}^2$, $\mathfrak{m}/Aa$ est isomorphe à un
facteur direct de $\mathfrak{m}/a\mathfrak{m}$.
\end{lemma}

\begin{proof}
Il existe un système minimal de générateurs de $\mathfrak{m}$ contenant $a$
(dont les images dans $\mathfrak{m}/\mathfrak{m}^2$ forment une base de ce
$k$-espace vectoriel)~; soit $\mathfrak{b}$ l'idéal engendré par les éléments
de ce système autres que $a$. Comme la relation $xa\in\mathfrak{b}$ entraîne
$x\in\mathfrak{m}$ (en considérant l'image de $xa$ dans
$\mathfrak{m}/\mathfrak{m}^2$), on a
$\mathfrak{b}\cap Aa\subset a\mathfrak{m}$ d'où un homomorphisme
$\mathfrak{b}/(\mathfrak{b}\cap Aa)\longrightarrow
\mathfrak{m}/a\mathfrak{m}$ déduit de l'injection
$\mathfrak{b}\longrightarrow\mathfrak{m}$ par passage aux quotients, et qui
est \emph{injectif}. En outre, $\mathfrak{b}+Aa=\mathfrak{m}$, et
l'homomorphisme composé
\[
  \mathfrak{m}/Aa=(\mathfrak{b}+Aa)/Aa
  \xrightarrow{\sim}\mathfrak{b}/(\mathfrak{b}\cap Aa)
  \longrightarrow\mathfrak{m}/a\mathfrak{m}
  \longrightarrow\mathfrak{m}/Aa
\]
est l'identité~; d'où le lemme.
\end{proof}

\begin{lemma}[17.3.1.4]
\label{0.17.3.1.4-fr}
Soient $A$ un anneau local noethérien, $\mathfrak{m}$ son idéal maximal, $E$
un $A$-module de type fini et de dimension projective finie. Si
$a\in\mathfrak{m}$ est $A$-régulier et $E$-régulier, alors $E/aE$ est un
$(A/aA)$-module de dimension projective finie, au plus égale à
$\operatorname{dim.proj}_A(E)$.
\end{lemma}

\begin{proof}
Raisonnons par récurrence sur $h=\operatorname{dim.proj}_A(E)$, le cas $h=0$
étant trivial puisque alors $E$ est un $A$-module projectif, donc $E/aE$ un
$(A/aA)$-module projectif. Il existe une suite exacte
\[
  0\longrightarrow N\longrightarrow L\longrightarrow E\longrightarrow0
\]
où $L$ est libre et $\operatorname{dim.proj}_A(N)=h-1$ (17.2.2, (iii)), $N$
étant de type fini. En outre, la suite
\[
  0\longrightarrow N/aN\longrightarrow L/aL\longrightarrow E/aE
  \longrightarrow0
\]
est exacte (15.1.18). Comme $a$ est $A$-régulier, il est aussi $N$-régulier
puisque $L$ est libre~; l'hypothèse de récurrence entraîne que $N/aN$ est un
$(A/aA)$-module de dimension projective $\leq h-1$, et comme $L/aL$ est un
$(A/aA)$-module libre, $E/aE$ est un $(A/aA)$-module de dimension projective
$\leq h$.
\end{proof}

Examinons maintenant deux cas~:

\medskip\noindent\textbf{I.} --- Supposons d'abord que tout élément de
$\mathfrak{m}-\mathfrak{m}^2$ soit diviseur de zéro dans $A$, auquel cas
(17.3.1.2) il existe $c\neq0$ dans $A$ tel que $c\mathfrak{m}=0$. Montrons
alors que $\mathfrak{m}=0$. S'il n'en était pas ainsi, notons d'abord que
$\mathfrak{m}$ ne pourrait être un $A$-module projectif, car il serait libre
(Bourbaki, \emph{Alg.~comm.}, chap.~II, \S~3, no~2, cor.~2 de la prop.~5),
ce qui contredit la relation $c\mathfrak{m}=0$. On aurait donc
$n=\operatorname{dim.coh}(A)\geq1$. Comme
$\mathfrak{m}\in\operatorname{Ass}(A)$, il existerait une suite exacte de
$A$-homomorphismes
\[
  0\longrightarrow k\longrightarrow A\longrightarrow E\longrightarrow0.
\]
Mais cela est absurde, car en vertu de la relation $n\geq1$, la suite exacte
des Tor donnerait la suite exacte
$0\longrightarrow\operatorname{Tor}^A_{n+1}(E,k)\longrightarrow
\operatorname{Tor}^A_n(k,k)\longrightarrow0$~; or on a
$\operatorname{Tor}^A_{n+1}(E,k)=0$ (17.2.2, (ii)) et
$\operatorname{Tor}^A_n(k,k)\neq0$ (17.2.11) et nous avons abouti à une
contradiction.

\medskip\noindent\textbf{II.} --- On peut donc se borner au cas où il existe
$a\in\mathfrak{m}-\mathfrak{m}^2$ qui est un élément $A$-régulier, et par
suite aussi $\mathfrak{m}$-régulier. Considérons l'anneau $A'=A/aA$ et son
idéal
\oldpage[0\textsubscript{IV-1}]{48}
maximal $\mathfrak{m}'=\mathfrak{m}/aA$~; il est clair que
$\operatorname{rg}_k(\mathfrak{m}'/\mathfrak{m}'^2)=n-1$. En vertu de
(17.3.1.4), $\mathfrak{m}/a\mathfrak{m}$ est un $A'$-module de dimension
projective finie, donc il en est de même de
$\mathfrak{m}'=\mathfrak{m}/aA$, qui en est un facteur direct (17.3.1.3 et
17.2.1). L'hypothèse de récurrence entraîne par suite que $A'=A/aA$ est
régulier, ce qui prouve par (17.1.8) que $A$ est régulier.

\begin{corollary}[17.3.2]
\label{0.17.3.2-fr}
Si $A$ est un anneau local régulier, $A_{\mathfrak{p}}$ est régulier pour tout
idéal premier $\mathfrak{p}$ de $A$.
\end{corollary}

\begin{proof}
En effet, on a vu (17.2.10) que
$\operatorname{dim.coh}(A_{\mathfrak{p}})\leq
\operatorname{dim.coh}(A)$, donc la conclusion résulte aussitôt de (17.3.1).
\end{proof}

\begin{proposition}[17.3.3]
\label{0.17.3.3-fr}
Soient $A$, $B$ deux anneaux locaux noethériens,
$\varphi:A\longrightarrow B$ un homomorphisme local, $\mathfrak{m}$ (resp.
$\mathfrak{n}$) l'idéal maximal de $A$ (resp. $B$),
$k=A/\mathfrak{m}$ (resp. $k'=B/\mathfrak{n}$) le corps résiduel de $A$
(resp. $B$). On a donc un homomorphisme canonique de $k'$-espaces vectoriels
\begin{equation}
\label{0.17.3.3.1-fr}
\tag{17.3.3.1}
  \psi:(\mathfrak{m}/\mathfrak{m}^2)\otimes_k k'
  \longrightarrow\mathfrak{n}/\mathfrak{n}^2.
\end{equation}
\begin{enumerate}
  \item[\rm (i)] Si $B$ est régulier et si $B$ est un $A$-module plat, $A$
  est régulier.
  \item[\rm (ii)] Les conditions suivantes sont équivalentes~:
  \begin{enumerate}
    \item[\rm a)] $B$ est régulier et l'homomorphisme $\psi$ (17.3.3.1) est
    injectif.
    \item[\rm b)] $A$ et $B$ sont réguliers, et pour un système régulier de
    paramètres $(x_i)_{1\leq i\leq m}$ de $A$, les $\varphi(x_i)$ font partie
    d'un système régulier de paramètres de $B$ (auquel cas cette propriété a
    lieu pour \emph{tout} système régulier de paramètres de $A$).
    \item[\rm c)] $B$ et $B\otimes_A k$ sont réguliers, et $B$ est un
    $A$-module plat.
    \item[\rm d)] $A$ et $B\otimes_A k$ sont réguliers, et $B$ est un
    $A$-module plat.
    \item[\rm e)] $A$ et $B\otimes_A k$ sont réguliers, et l'on a
    \begin{equation}
    \label{0.17.3.3.2-fr}
    \tag{17.3.3.2}
      \dim(B)=\dim(A)+\dim(B\otimes_A k).
    \end{equation}
  \end{enumerate}
\end{enumerate}
\end{proposition}

\begin{proof}
\textnormal{(i)} On a
$\operatorname{dim.coh}(A)\leq\operatorname{dim.coh}(B)$ par (17.2.13), donc
il suffit d'appliquer (17.3.1).

\textnormal{(ii)} Lorsque $A$ est supposé régulier et que $(x_i)$ est un
système régulier de paramètres de $A$, dire que $B$ est un $A$-module
\emph{plat} équivaut, en vertu de (15.1.21) à dire que la suite des
$y_i=\varphi(x_i)$ est $B$-régulière (puisque $B\otimes_A k$ est un
$k$-module plat). D'autre part, comme $\dim(A)=m$, et que la suite $(y_i)$
engendre l'idéal $\mathfrak{m}B$, la relation (17.3.3.2), qui s'écrit aussi
$\dim(B/\mathfrak{m}B)=\dim(B)-m$, équivaut à dire que la suite
$(y_i)_{1\leq i\leq m}$ fait partie d'un système de paramètres de $B$
(16.3.7). On voit donc que les conditions b), d) et e) équivalent
respectivement aux suivantes~:
\begin{enumerate}
  \item[\rm b')] $A$ et $B$ sont réguliers et $(y_i)_{1\leq i\leq m}$ fait
  partie d'un système régulier de paramètres de $B$.
  \item[\rm d')] $A$ est régulier,
  $B/(\sum_{i=1}^m y_iB)$ est régulier et la suite $(y_i)$ est
  $B$-régulière.
  \item[\rm e')] $A$ est régulier,
  $B/(\sum_{i=1}^m y_iB)$ est régulier et la suite $(y_i)$ fait partie d'un
  système de paramètres de $B$.
\end{enumerate}
Or, b') et e') sont équivalentes en vertu de (17.1.7), et comme d') entraîne
e') (16.4.1) et est entraînée par b') (17.1.7), elle leur est équivalente. La
conjonction
\oldpage[0\textsubscript{IV-1}]{49}
de b) et d) implique trivialement c), et en vertu de (i), c) entraîne d)~; on
a donc prouvé l'équivalence de b), c), d) et e). En outre, il est clair que b)
entraîne a), les classes des $x_i$ dans
$\mathfrak{m}/\mathfrak{m}^2$ formant alors une base de ce $k$-espace
vectoriel et les classes des $y_i$ dans $\mathfrak{n}/\mathfrak{n}^2$ un
système libre de ce $k'$-espace vectoriel. Reste donc à prouver que a)
entraîne e). Posons $V=\mathfrak{m}/\mathfrak{m}^2$,
$W=\mathfrak{n}/\mathfrak{n}^2$~; il est immédiat que l'on a
$\psi(V\otimes_k k')=(\mathfrak{n}^2+\mathfrak{m}B)/\mathfrak{n}^2$. On a
d'une part, en vertu de (16.3.9)
\begin{equation}
\label{0.17.3.3.3-fr}
\tag{17.3.3.3}
  \dim(B)\leq\dim(A)+\dim(B\otimes_A k)~;
\end{equation}
en second lieu, d'après (16.2.6), on a
\begin{equation}
\label{0.17.3.3.4-fr}
\tag{17.3.3.4}
  \dim(A)\leq\operatorname{rg}_{k'}(V\otimes_k k').
\end{equation}
Enfin, dans l'anneau local $B\otimes_A k=B/\mathfrak{m}B$,
$\mathfrak{n}'=\mathfrak{n}/\mathfrak{m}B$ est l'idéal maximal, $k'$ le corps
résiduel et $\mathfrak{n}'/\mathfrak{n}'^2$ est isomorphe à
$\mathfrak{n}/(\mathfrak{n}^2+\mathfrak{m}B)
=W/\psi(V\otimes_k k')$~; on a donc, d'après (16.2.6)
\begin{equation}
\label{0.17.3.3.5-fr}
\tag{17.3.3.5}
  \dim(B\otimes_A k)
  \leq\operatorname{rg}_{k'}W
  -\operatorname{rg}_{k'}\psi(V\otimes_k k').
\end{equation}
Enfin, puisque $B$ est supposé régulier, on a
$\dim(B)=\operatorname{rg}_{k'}W$ (17.1.1)~; on conclut donc de (17.3.3.3),
(17.3.3.4) et (17.3.3.5) que l'on a
\[
  \operatorname{rg}_{k'}W
  \leq\dim(A)+\dim(B\otimes_A k)
  \leq\operatorname{rg}_{k'}W
  +\operatorname{rg}_{k'}(V\otimes_k k')
  -\operatorname{rg}_{k'}\psi(V\otimes_k k')
\]
et dire que les deux termes extrêmes de cette inégalité sont égaux signifie
que $\psi$ est \emph{injectif}. La condition a) entraîne donc nécessairement
que dans chacune des relations (17.3.3.3), (17.3.3.4) et (17.3.3.5), les deux
membres sont \emph{égaux}~; or, l'égalité dans (17.3.3.4) (resp. (17.3.3.5))
signifie que $A$ (resp. $B\otimes_A k$) est \emph{régulier} (17.1.1)~; on a
donc bien prouvé que a) entraîne e).
\end{proof}

\begin{proposition}[17.3.4]
\label{0.17.3.4-fr}
Soient $A$ un anneau local régulier de dimension $n$, $\mathfrak{m}$ son idéal
maximal. Pour tout $A$-module non nul de type fini $M$, on a
\begin{equation}
\label{0.17.3.4.1-fr}
\tag{17.3.4.1}
  \operatorname{prof}(M)+\operatorname{dim.proj}(M)=n.
\end{equation}
\end{proposition}

\begin{proof}
Raisonnons par récurrence sur $m=\operatorname{prof}(M)$. Si $m=0$, on sait
(16.4.6, (i)) qu'il existe un sous-$A$-module $N$ de $M$ isomorphe à
$k=A/\mathfrak{m}$~; appliquant la suite exacte des Tor à la suite exacte
$0\longrightarrow N\longrightarrow M\longrightarrow M/N\longrightarrow0$,
on obtient une suite exacte
\[
  \operatorname{Tor}^A_{n+1}(M/N,k)\longrightarrow
  \operatorname{Tor}^A_n(N,k)\longrightarrow
  \operatorname{Tor}^A_n(M,k),
\]
et l'hypothèse que $A$ est régulier entraîne
$\operatorname{Tor}^A_{n+1}(M/N,k)=0$ ((17.3.1.1) et (17.2.6)), donc
$\operatorname{Tor}^A_n(k,k)$ est isomorphe à un sous-$A$-module de
$\operatorname{Tor}^A_n(M,k)$~; appliquant de nouveau (17.3.1.1) et (17.2.6),
on voit que $\operatorname{Tor}^A_n(M,k)\neq0$, donc
$\operatorname{dim.proj}(M)\geq n$ (17.2.6)~; mais comme
$\operatorname{dim.proj}(M)\leq n$ par (17.3.1.1), on a bien
$\operatorname{dim.proj}(M)=n$.

Supposons maintenant $m>0$, et soit $x$ un élément $M$-régulier appartenant à
$\mathfrak{m}$~; on sait alors (16.4.6, (i)) que l'on a
$\operatorname{prof}(M/xM)=m-1$, et d'autre part (17.2.7)
$\operatorname{dim.proj}(M/xM)=\operatorname{dim.proj}(M)+1$~; l'hypothèse de
récurrence prouve aussitôt la relation (17.3.4.1).
\end{proof}

\oldpage[0\textsubscript{IV-1}]{50}

\begin{corollary}[17.3.5]
\label{0.17.3.5-fr}
\begin{enumerate}
  \item[\rm(i)] Soit $A$ un anneau local régulier de dimension $n$~; pour
  qu'un $A$-module $M\neq0$ de type fini soit libre, il faut et il suffit que
  $M$ soit un $A$-module de Cohen-Macaulay de dimension $n$.
  \item[\rm(ii)] Soient $A$ un anneau local régulier, $B$ un anneau local,
  $\rho:A\to B$ un homomorphisme local faisant de $B$ un $A$-module de type
  fini. Pour que $B$ soit un $A$-module libre, il faut et il suffit que $B$
  soit un anneau de Cohen-Macaulay et que $\rho$ soit injectif (ou, ce qui
  revient au même (16.1.5), que $\dim(B)=\dim(A)$).
\end{enumerate}
\end{corollary}

\begin{proof}
(i) Cela résulte de (17.3.4.1) et du fait que pour un $A$-module de type fini,
il revient au même de dire que ce module est projectif ou libre (Bourbaki,
\emph{Alg. comm.}, chap.~II, \S~3, n\textsuperscript{o}~2, cor.~2 de la
prop.~5)~; les $A$-modules libres de type fini $M$ sont donc caractérisés par
la relation $\operatorname{dim.proj}(M)=0$ (17.2.2, (i)).

(ii) Dire que $B$ est un anneau de Cohen-Macaulay équivaut à dire que $B$ est
un $A$-module de Cohen-Macaulay (16.5.3), donc il suffit d'appliquer (i),
puisque $\dim(B)=\dim(A)$ (16.1.5).
\end{proof}

\begin{env}[17.3.6]
\label{0.17.3.6-fr}
On dit qu'un anneau noethérien $A$ est \emph{régulier} si pour tout idéal
premier $\mathfrak{p}$ de $A$, l'anneau local $A_{\mathfrak{p}}$ est
régulier~; lorsque $A$ lui-même est un anneau local, il résulte de (17.3.2)
que cette définition est équivalente à celle de (17.1.2). Pour que $A$ soit
régulier, il suffit, en vertu de (17.3.2), que $A_{\mathfrak{m}}$ soit
régulier pour tout idéal maximal $\mathfrak{m}$ de $A$. En outre, il résulte
aussitôt de cette définition que pour toute partie multiplicative $S$ de
$A$, $S^{-1}A$ est régulier.
\end{env}

\begin{proposition}[17.3.7]
\label{0.17.3.7-fr}
Si $A$ est un anneau noethérien régulier, tout anneau de polynômes
$A[T_1,\ldots,T_n]$ est régulier.
\end{proposition}

\begin{proof}
Il suffit évidemment de prouver que l'anneau de polynômes $B=A[T]$ est
régulier~; comme $B$ est un $A$-module libre, pour tout idéal premier
$\mathfrak{q}$ de $B$, $B_{\mathfrak{q}}$ est un
$A_{\mathfrak{p}}$-module plat, où $\mathfrak{p}=\mathfrak{q}\cap A$
($0_{\mathrm I}$, 6.3.1)~; il suffit donc, en vertu de (17.1.10), de prouver
que $B_{\mathfrak{q}}/\mathfrak{p}B_{\mathfrak{q}}$ est régulier, et comme
cet anneau est un anneau local en un idéal premier de
$A_{\mathfrak{p}}[T]/\mathfrak{p}A_{\mathfrak{p}}[T]=k[T]$ (où $k$ est le
corps résiduel $A_{\mathfrak{p}}/\mathfrak{p}A_{\mathfrak{p}}$), il suffit
de prouver que $C=k[T]$ est régulier~; or, $C$ étant principal, les anneaux
locaux en les idéaux premiers de $C$ sont des anneaux de valuation discrète
ou un corps (pour l'idéal $(0)$), donc réguliers (17.1.4), ce qui achève la
démonstration.
\end{proof}

\begin{corollary}[17.3.8]
\label{0.17.3.8-fr}
Si $A$ est un anneau régulier, tout anneau de séries formelles
$A[[T_1,\ldots,T_n]]$ est régulier.
\end{corollary}

\begin{proof}
Soit $\mathfrak{J}$ l'idéal engendré par les $T_i$ dans l'anneau de polynômes
$A[T_1,\ldots,T_n]$~; comme ce dernier est régulier par (17.3.7), et que
$A[[T_1,\ldots,T_n]]$ est le complété de $A[T_1,\ldots,T_n]$ pour la
topologie $\mathfrak{J}$-préadique (Bourbaki, \emph{Alg. comm.}, chap.~III,
\S~2, n\textsuperscript{o}~12, \emph{Exemple 1}), la conclusion résulte du
\end{proof}

\begin{lemma}[17.3.8.1]
\label{0.17.3.8.1-fr}
Soient $A$ un anneau régulier, $\mathfrak{J}$ un idéal de $A$,
$\widehat{A}$ le séparé complété de $A$ pour la topologie
$\mathfrak{J}$-préadique. Alors $\widehat{A}$ est régulier.
\end{lemma}

\begin{proof}
Il suffit en effet (17.3.6) de voir que pour tout idéal maximal
$\mathfrak{n}$ de $\widehat{A}$, l'anneau local
$\widehat{A}_{\mathfrak{n}}$ est régulier~; on sait (Bourbaki,
\emph{Alg. comm.}, chap.~III, \S~3, n\textsuperscript{o}~4, prop.~8) que
$\mathfrak{n}$
\oldpage[0\textsubscript{IV-1}]{51}
est de la forme $\mathfrak{m}\widehat{A}$, où $\mathfrak{m}$ est un idéal
maximal de $A$ contenant $\mathfrak{J}$, et que l'homomorphisme canonique
$A\to\widehat{A}$ donne un homomorphisme \emph{injectif}
$A_{\mathfrak{m}}\to\widehat{A}_{\mathfrak{n}}$, tel que la topologie
$\mathfrak{n}\widehat{A}_{\mathfrak{n}}$-préadique de
$\widehat{A}_{\mathfrak{n}}$ induise sur $A_{\mathfrak{m}}$ la topologie
$\mathfrak{m}A_{\mathfrak{m}}$-préadique, et que $A_{\mathfrak{m}}$ soit
dense dans $\widehat{A}_{\mathfrak{n}}$ pour la topologie
$\mathfrak{n}\widehat{A}_{\mathfrak{n}}$-préadique~; par suite les
\emph{complétés} des anneaux locaux noethériens $A_{\mathfrak{m}}$ et
$\widehat{A}_{\mathfrak{n}}$ sont canoniquement \emph{isomorphes}. Or, par
hypothèse $A_{\mathfrak{m}}$ est régulier, donc il en est de même de son
complété (17.1.5) et comme le complété de $\widehat{A}_{\mathfrak{n}}$ est
régulier, il en est de même de $\widehat{A}_{\mathfrak{n}}$ par (17.1.5).
\end{proof}

\begin{corollary}[17.3.9]
\label{0.17.3.9-fr}
Soit $A$ un anneau noethérien, quotient d'un anneau noethérien régulier $B$.
Si $C$ est une $A$-algèbre de type fini, tout anneau de fractions $S^{-1}C$
de $C$ est quotient d'un anneau régulier.
\end{corollary}

\begin{proof}
En effet, $C$ est quotient d'un anneau de polynômes $A[T_1,\ldots,T_n]$~;
comme $A=B/\mathfrak{J}$, où $\mathfrak{J}$ est un idéal de $B$, on a
$A[T_1,\ldots,T_n]=B[T_1,\ldots,T_n]/\mathfrak{J}B[T_1,\ldots,T_n]$~; en
vertu de (17.3.7), on peut donc se borner au cas où $C$ est quotient d'un
anneau régulier $B'$~; mais si $S'$ est l'image réciproque de $S$ dans $B'$,
$S'$ est une partie multiplicative de $B'$ et $C$ est un anneau quotient de
$S'^{-1}B'$, si bien que finalement tout revient à voir que lorsque $C$ est
\emph{régulier}, il en est de même de $S^{-1}C$, ce qu'on a vu en (17.3.6).
\end{proof}

L'importance des quotients d'anneaux réguliers tient entre autres à la
propriété précédente et au fait que ce sont des anneaux \emph{caténaires}
(16.5.12).

Tous les anneaux importants dans les applications à la Géométrie algébrique
sont des quotients d'anneaux réguliers.

% FIN SOURCE ega4/ega0-17-fr.tex

%
% DEBUT SOURCE ega4/ega0-18-fr.tex
\section{Compléments sur les extensions d'algèbres}
\label{section:0.18-fr}

Ce paragraphe rassemble un certain nombre de constructions fonctorielles sur
les anneaux, qui seront utilisées de façon répétée aux \S\S~19 et 20~; il ne
contient aucun résultat non trivial.

\subsection{Images réciproques d'anneaux augmentés}
\label{subsection:0.18.1-fr}

\begin{env}[18.1.1]
\label{0.18.1.1-fr}
Étant donné un anneau $A$ (non nécessairement commutatif), la catégorie des
\emph{$A$-anneaux} a pour objets les couples $(B,\rho)$ formés d'un anneau
$B$ et d'un homomorphisme d'anneaux $\rho:A\to B$, pour morphismes (dits aussi
\emph{$A$-homomorphismes}) les homomorphismes d'anneaux $u:B\to C$ tels que,
si $\rho:A\to B$ et $\sigma:A\to C$ sont les homomorphismes (dits
\emph{structuraux}) définissant la structure de $A$-anneau sur $B$ et $C$
respectivement, le diagramme
\begin{equation}
\label{0.18.1.1.1-fr}
\tag{18.1.1.1}
  \xymatrix{
    B\ar[r]^u & C\\
    & A\ar[lu]^\rho\ar[u]_\sigma
  }
\end{equation}
soit commutatif. Le \emph{noyau} $\mathfrak{J}$ de $u$ est un idéal bilatère
qui est un $A$-\emph{bimodule} (pour $\rho$).

Si $f:A'\to A$ est un homomorphisme d'anneaux, pour tout $A$-anneau
$(B,\rho)$, $(B,\rho\circ f)$ est un $A'$-anneau, et si $u:B\to C$ est un
$A$-homomorphisme du $A$-anneau $(B,\rho)$
\oldpage[0\textsubscript{IV-1}]{52}
dans le $A$-anneau $(C,\sigma)$, c'est aussi un $A'$-homomorphisme du
$A'$-anneau $(B,\rho\circ f)$ dans le $A'$-anneau $(C,\sigma\circ f)$~; on
définit ainsi un foncteur canonique de la catégorie des $A$-anneaux dans celle
des $A'$-anneaux.
\end{env}

\begin{env}[18.1.2]
\label{0.18.1.2-fr}
Soient $B$, $E$, $F$ trois $A$-anneaux, $f:E\to B$, $g:F\to B$ deux
$A$-homomorphismes. Rappelons que l'on appelle \emph{produit fibré} de $E$ et
$F$ sur $B$ (pour les $A$-homomorphismes $f$ et $g$) et que l'on note
$E\times_BF$ le sous-anneau $G$ de l'anneau produit $E\times F$, formé des
couples $(x,y)$ tels que $f(x)=g(y)$~; les restrictions $p_1:G\to E$,
$p_2:G\to F$ des projections $\operatorname{pr}_1$ et
$\operatorname{pr}_2$ à $G$ sont encore appelées les \emph{projections
canoniques}~; la structure de $A$-anneau de $G$ est définie par
l'homomorphisme $\alpha\mapsto(\rho(\alpha),\sigma(\alpha))$ (où
$\rho:A\to E$ et $\sigma:A\to F$ sont les homomorphismes structuraux), qui
applique effectivement $A$ dans $G$ en vertu de (18.1.1). La propriété
caractéristique de $E\times_BF$ est que, pour tout couple de
$A$-homomorphismes $u:C\to E$, $v:C\to F$ tels que
$f\circ u=g\circ v$, il existe un $A$-homomorphisme et un seul $w:C\to G$
tel que $u=p_1\circ w$, $v=p_2\circ w$. On peut encore dire que
$E\times_BF$ est \emph{limite projective} du système projectif formé de $B$,
$E$, $F$ et des $A$-homomorphismes $f$, $g$, dans la catégorie des
$A$-anneaux ($0_{\mathrm{III}}$, 8.1.9).
\end{env}

\begin{env}[18.1.3]
\label{0.18.1.3-fr}
Soit $\mathfrak{J}$ l'idéal bilatère de $E$, noyau de $f$~; il est immédiat
que le noyau $\mathfrak{J}'$ de $p_2:G\to F$ est l'idéal formé des éléments
$(x,0)$ où $x\in\mathfrak{J}$~; la restriction
$i_1:\mathfrak{J}'\to\mathfrak{J}$ de $p_1$ est donc un
\emph{isomorphisme} d'anneau \emph{sans élément unité}, et aussi un
isomorphisme de $A$-bimodules. De même, si $\mathfrak{R}$ est le noyau de
$g$, le noyau $\mathfrak{R}'$ de $p_1:G\to E$ est l'idéal des éléments
$(0,y)$ où $y\in\mathfrak{R}$, et la restriction
$i_2:\mathfrak{R}'\to\mathfrak{R}$ de $p_2$ est un isomorphisme (au même
sens). Enfin, il est clair que le noyau du $A$-homomorphisme
$q=f\circ p_1=g\circ p_2$ de $E\times_BF$ dans $B$ est
$\mathfrak{J}\times\mathfrak{R}=\mathfrak{J}'\oplus\mathfrak{R}'$, de sorte
que l'on a le diagramme commutatif
\begin{equation}
\label{0.18.1.3.1-fr}
\tag{18.1.3.1}
  \xymatrix@C=2.2pc@R=1.8pc{
    0\ar[dr] & & 0\ar[d] & 0\ar[d]\\
    & \mathfrak{J}'\oplus\mathfrak{R}'\ar[r]\ar[d]\ar[dr]
      & \mathfrak{R}'\ar[r]^{i_2}_{\sim}\ar[d]^j
      & \mathfrak{R}\ar[d]^{p_2}\\
    0\ar[r] & \mathfrak{J}'\ar[r]\ar[d]^{i_1}_{\sim}
      & G\ar[r]\ar[d]^{p_1}\ar[dr]^q & F\ar[d]^g\\
    0\ar[r] & \mathfrak{J}\ar[r] & E\ar[r]_f & B
  }
\end{equation}

Les définitions et résultats de (18.1.2) et (18.1.3) s'étendent aussitôt au
produit fibré d'une \emph{famille} quelconque $(E_\lambda)_{\lambda\in I}$ de
$A$-anneaux défini par une famille de $A$-homomorphismes
$f_\lambda:E_\lambda\to B$. Nous laissons la formulation de ces résultats au
lecteur.
\end{env}

\begin{env}[18.1.4]
\label{0.18.1.4-fr}
Conformément à la terminologie de (M, VIII), nous appellerons
\emph{$A$-anneau augmenté sur $B$} un $A$-anneau $E$ muni d'un
$A$-homomorphisme \emph{surjectif} (dit \emph{augmentation} de $E$),
$f:E\to B$~; le noyau $\mathfrak{J}$ de $f$ est appelé
\emph{l'idéal d'augmentation}. On dit que le $A$-anneau augmenté $E$ est
\emph{trivial} s'il existe un $A$-homomorphisme de $A$-anneaux $s:B\to E$
qui est \emph{inverse à droite} de l'augmentation $f:E\to B$ (autrement dit
$f\circ s=1_B$). La suite exacte de $A$-bimodules
\[
  0\longrightarrow\mathfrak{J}\longrightarrow E
  \xrightarrow{f}B\longrightarrow0
\]
\oldpage[0\textsubscript{IV-1}]{53}
est alors \emph{scindée}~; autrement dit, on peut identifier le $A$-bimodule
$E$ à $B\times\mathfrak{J}$, et la multiplication dans $E$ est alors donnée
par
\[
  (b,z)(b',z')=(bb',bz'+zb'+zz'),
\]
$\mathfrak{J}$ étant considéré comme $B$-bimodule au moyen de $s:B\to E$.
\end{env}

\begin{env}[18.1.5]
\label{0.18.1.5-fr}
Avec les notations de (18.1.2), supposons que le $A$-homomorphisme $f$ soit
\emph{surjectif}, autrement dit fasse de $E$ un $A$-anneau \emph{augmenté sur
$B$}. Alors il est clair que $p_2:G\to F$ est aussi \emph{surjectif},
autrement dit définit sur $G$ une structure de $A$-anneau \emph{augmenté sur
$F$}, que l'on appelle \emph{l'image réciproque} par $g:F\to B$ de l'anneau
augmenté $E$.
\end{env}

\begin{proposition}[18.1.6]
\label{0.18.1.6-fr}
Pour que l'image réciproque $G=E\times_BF$ par $g:F\to B$ du $A$-anneau
augmenté $E$ soit un $A$-anneau augmenté trivial, il faut et il suffit qu'il
existe un $A$-homomorphisme $u:F\to E$ rendant commutatif le diagramme
\[
  \xymatrix{
    & F\ar[ld]_u\ar[d]^g\\
    E\ar[r]_f & B
  }
\]
\end{proposition}

\begin{proof}
La condition est évidemment nécessaire, en prenant $u=p_1\circ s$, où
$s:F\to G$ est un $A$-homomorphisme inverse à droite de $p_2$. Inversement,
s'il existe un $A$-homomorphisme $u$ vérifiant la condition de l'énoncé,
l'existence de l'inverse à droite $s$ de $p_2$ résulte de la propriété
universelle du produit fibré (18.1.2) appliquée aux $A$-homomorphismes
$u:F\to E$ et $1_F:F\to F$.

Ce résultat entraîne en particulier que si $E$ est un anneau augmenté
\emph{trivial}, il en est de même de toutes ses images réciproques.
\end{proof}

\begin{env}[18.1.7]
\label{0.18.1.7-fr}
Reprenons la situation décrite en (18.1.2) et (18.1.3)~; on a évidemment sur
$\mathfrak{R}$ une structure de $G$-bimodule, provenant de sa structure de
$F$-bimodule et de l'homomorphisme d'anneaux $p_2:G\to F$. Comme $i_2$ est
bijectif on a en outre une injection
$\theta=j\circ i_2^{-1}:\mathfrak{R}\to G$, qui est un homomorphisme de
$G$-bimodules. Nous allons voir inversement que, lorsqu'on suppose en outre
que $f$ et $g$ sont \emph{surjectifs}, autrement dit que $E$ et $F$ sont des
$A$-anneaux \emph{augmentés} sur $B$, la donnée d'un tel homomorphisme
$\theta$ permet de reconstituer l'anneau augmenté $E$ (sur $B$) à partir de
l'anneau augmenté $G$ (sur $F$).

De façon précise, donnons-nous un $A$-anneau $F$ augmenté sur $B$ et un
$A$-anneau $G$ augmenté sur $F$, les idéaux d'augmentation étant notés
$\mathfrak{R}$ et $\mathfrak{J}'$ respectivement~:
\begin{equation}
\label{0.18.1.7.1-fr}
\tag{18.1.7.1}
  \xymatrix@C=2pc@R=1.5pc{
    & & 0\ar[d] & & \\
    & & \mathfrak{R}\ar[ld]^\theta\ar[d] & & \\
    0\ar[r] & \mathfrak{J}'\ar[r] & G\ar[r]_h & F\ar[r]\ar[d]^g & 0\\
    & & & B\ar[d] & \\
    & & & 0 &
  }
\end{equation}

\oldpage[0\textsubscript{IV-1}]{54}
L'homomorphisme $h$ définit sur tous les idéaux de $F$ (et en particulier sur
$\mathfrak{R}$) une structure de $G$-bimodule. Soit
$\theta:\mathfrak{R}\to G$ un homomorphisme de $G$-bimodules rendant
commutatif le diagramme (18.1.7.1)~; cela implique que $\theta$ est
\emph{injectif} et que $\mathfrak{J}'\cap\theta(\mathfrak{R})=0$~; comme
$\theta(\mathfrak{R})$ est un idéal bilatère de $G$, et que
$h(\theta(\mathfrak{R}))=\mathfrak{R}$, l'homomorphisme composé
$g\circ h:G\to B$ se factorise en
$G\xrightarrow{q}G/\theta(\mathfrak{R})\xrightarrow{f}B$, où $f$ est
surjectif~; en outre l'image $\mathfrak{J}$ de $\mathfrak{J}'$ dans
$E=G/\theta(\mathfrak{R})$ est le noyau de $f$ (celui de $g\circ h$ étant
$\mathfrak{J}'\oplus\theta(\mathfrak{R})$), et la restriction à
$\mathfrak{J}'$ de l'homomorphisme canonique $q:G\to E$ est injective. On
peut donc former le \emph{produit fibré} $G'=E\times_BF$, et comme les deux
$A$-homomorphismes $q:G\to E$ et $h:G\to F$ sont tels que
$f\circ q=g\circ h$, ils définissent un unique $A$-homomorphisme $u:G'\to G$
par la propriété universelle de $G'$ (18.1.2)~; nous allons voir que $u$ est
\emph{bijectif}. Il suffit de le prouver lorsque $u$ est considéré comme
homomorphisme de $A$-bimodules~; mais on notera alors que $u$ est compatible
avec les filtrations finies sur $G'$ et $G$ formées respectivement par $G'$
et $\mathfrak{J}''=\operatorname{Ker}(G'\to F)$, et par $G$ et
$\mathfrak{J}'$~; en outre, on a vu (18.1.3) que
$\operatorname{gr}_0u:F\to F$ est l'identité et que
$\operatorname{gr}_1u:\mathfrak{J}''\to\mathfrak{J}'$ est bijectif, donc $u$
lui-même est bijectif (Bourbaki, \emph{Alg. comm.}, chap.~III, \S~2,
n\textsuperscript{o}~8, cor.~3 du th.~1).
\end{env}

\subsection{Extensions d'un anneau par un bimodule}
\label{subsection:0.18.2-fr}

\begin{env}[18.2.1]
\label{0.18.2.1-fr}
Soient $E$ un $A$-anneau augmenté sur $B$, $f:E\to B$ l'augmentation,
$\mathfrak{J}=\operatorname{Ker}(f)$ l'idéal d'augmentation. Si l'on a
$\mathfrak{J}^2=0$, $\mathfrak{J}$ est non seulement un $E$-bimodule mais
aussi un $B$-\emph{bimodule} puisque $B$ est isomorphe à
$E/\mathfrak{J}$~; de façon précise, tout $b\in B$ est de la forme $f(x)$
avec $x\in E$, et si $z,z'$ sont dans $\mathfrak{J}$ on a
$(x+z)z'=xz'$ (resp. $z'(x+z)=z'x$) de sorte que la valeur de $xz'$ (resp.
$z'x$) ne dépend pas de l'élément $x$ de $f^{-1}(b)$, et peut s'écrire $bz'$
(resp. $z'b$) ce qui définit la structure de $B$-bimodule considérée.
Inversement, si $\mathfrak{J}$ est muni d'une structure de $B$-bimodule telle
que $xz'=f(x)z'$ et $z'x=z'f(x)$ pour $x\in E$ et
$z'\in\mathfrak{J}$, il est clair que $\mathfrak{J}^2=0$.
\end{env}

\begin{env}[18.2.2]
\label{0.18.2.2-fr}
On appelle \emph{$A$-extension d'un $A$-anneau $B$ par un $B$-bimodule $L$}
une \emph{suite exacte} d'homomorphismes de $A$-bimodules
\[
  0\longrightarrow L\xrightarrow{j}E\xrightarrow{f}B\longrightarrow0
\]
où $E$ est un $A$-anneau, $f$ un $A$-homomorphisme d'anneaux et l'on a, pour
$x\in E$ et $z\in L$,
\[
  j(f(x)z)=xj(z),\qquad j(zf(x))=j(z)x
\]
d'où résulte (18.2.1) que $j(L)$ est un idéal bilatère de carré nul de $E$.
Par abus de langage on dira aussi que $E$ est une extension de $B$ par $L$.
On dit que deux $A$-extensions $E$, $E'$ de $B$ par $L$ sont
\emph{$A$-équivalentes} s'il existe un isomorphisme de $A$-anneaux
$u:E\xrightarrow{\sim}E'$ (dit aussi \emph{$A$-équivalence de
$A$-extensions}) rendant commutatif le diagramme
\begin{equation}
\label{0.18.2.2.1-fr}
\tag{18.2.2.1}
  \xymatrix@C=2.1pc@R=1.2pc{
    & & E\ar[dr]\ar[dd]_u\ar@{}[dd]|{\wr} & & \\
    0\ar[r] & L\ar[ur]\ar[dr] & & B\ar[r] & 0\\
    & & E'\ar[ur] & &
  }
\end{equation}
\end{env}

\oldpage[0\textsubscript{IV-1}]{55}

\begin{env}[18.2.3]
\label{0.18.2.3-fr}
On dit qu'une $A$-extension $E$ de $B$ par un $B$-bimodule $L$ est
\emph{$A$-triviale} si $E$ est un $A$-anneau augmenté (sur $B$) trivial
(18.1.4). On définit sur le $A$-bimodule produit $B\times L$ une structure de
$A$-anneau en posant $(x,s)(y,t)=(xy,xt+sy)$, comme on le vérifie aussitôt,
et il est immédiat que les applications canoniques $j:L\to B\times L$ et
$f:B\times L\to B$ définissent une $A$-extension de $B$ par $L$ qui est
$A$-triviale, l'application canonique $g:B\to B\times L$ étant un
$A$-homomorphisme inverse à droite de $f$. On dit que cette extension est
\emph{l'extension triviale type} de $B$ par $L$ et on la note $D_B(L)$~; il
est immédiat que toute $A$-extension $A$-triviale de $B$ par $L$ est
$A$-équivalente à $D_B(L)$.

On notera que toute extension du $A$-anneau $A$ lui-même par un $A$-bimodule
est nécessairement $A$-triviale.
\end{env}

\begin{env}[18.2.4]
\label{0.18.2.4-fr}
Étant données deux $A$-extensions
$L\xrightarrow{j}E\xrightarrow{f}B$,
$L'\xrightarrow{j'}E'\xrightarrow{f'}B'$, un \emph{morphisme} de la première
dans la seconde est par définition un triplet d'homomorphismes de
$A$-bimodules $(u,v,w)$ tel que le diagramme
\begin{equation}
\label{0.18.2.4.1-fr}
\tag{18.2.4.1}
  \xymatrix{
    0\ar[r] & L\ar[r]^j\ar[d]_w & E\ar[r]^f\ar[d]_u
      & B\ar[r]\ar[d]_v & 0\\
    0\ar[r] & L'\ar[r]_{j'} & E'\ar[r]_{f'} & B'\ar[r] & 0
  }
\end{equation}
soit commutatif, $u$ et $v$ étant des $A$-homomorphismes \emph{d'anneaux} et
$w$ étant tel que $w(bz)=v(b)w(z)$ et $w(zb)=w(z)v(b)$ pour $z\in L$ et
$b\in B$ (en d'autres termes, le couple $(v,w)$ constitue un
\emph{di-homomorphisme} du $B$-bimodule $L$ dans le $B'$-bimodule $L'$)~; il
est clair que si $(u',v',w')$ est un morphisme de
$L'\to E'\to B'$ dans une $A$-extension $L''\to E''\to B''$,
$(u'\circ u,v'\circ v,w'\circ w)$ est encore un morphisme, ce qui justifie la
terminologie.

La considération des deux carrés commutatifs du diagramme (18.2.4.1) va nous
conduire à deux opérations sur les extensions de $A$-anneaux.
\end{env}

\begin{env}[18.2.5]
\label{0.18.2.5-fr}
En premier lieu, considérons une $A$-extension $E'$ de $B'$ par $L'$
\[
  0\longrightarrow L'\xrightarrow{j'}E'\xrightarrow{f'}B'\longrightarrow0
\]
et un $A$-homomorphisme d'anneaux $v:B\to B'$, et soit
$F=E'\times_{B'}B$ \emph{l'image réciproque} par $v$ du $A$-anneau augmenté
$E'$ (18.1.5), de sorte que l'on a un diagramme commutatif
\[
  \xymatrix{
    0\ar[r] & L_0\ar[r]\ar[d]^i\ar@{}[d]|{\wr} & F\ar[r]^{p_2}\ar[d]^{p_1}
      & B\ar[r]\ar[d]^v & 0\\
    0\ar[r] & L'\ar[r]_{j'} & E'\ar[r]_{f'} & B'\ar[r] & 0
  }
\]
dont les lignes sont exactes, $p_1$ et $p_2$ étant les homomorphismes
canoniques~; on a vu (18.1.3) que $i$ est bijectif et il résulte aussi de la
définition (18.1.2) que $L_0^2=0$, de sorte que l'on peut considérer que $F$
est une $A$-extension de $B$ par $L'$, que l'on appelle \emph{image
réciproque par $v$} de l'extension $E'$ de $B'$ par $L'$ ($L'$ étant
naturellement considéré comme $B$-bimodule au moyen de l'homomorphisme
d'anneaux $v$). Le caractère fonctoriel
\oldpage[0\textsubscript{IV-1}]{56}
du produit fibré vis-à-vis de chacun des facteurs montre en outre que si l'on
a un morphisme de deux extensions de $B'$
\[
  \xymatrix{
    0\ar[r] & L'_1\ar[r]\ar[d]^h & E'_1\ar[r]\ar[d]^g
      & B'\ar[r]\ar[d]^{1_{B'}} & 0\\
    0\ar[r] & L'_2\ar[r] & E'_2\ar[r] & B'\ar[r] & 0
  }
\]
on en déduit un morphisme \emph{image réciproque par $v$}
\[
  \xymatrix{
    0\ar[r] & L'_1\ar[r]\ar[d]^h & E'_1\times_{B'}B\ar[r]
      \ar[d]^{g\times_{B'}1_B} & B\ar[r]\ar[d]^{1_B} & 0\\
    0\ar[r] & L'_2\ar[r] & E'_2\times_{B'}B\ar[r] & B\ar[r] & 0.
  }
\]
En particulier, si $E'_1$ et $E'_2$ sont des $A$-extensions
$A$-équivalentes de $B'$ par $L'$, leurs images réciproques par $v$ sont des
$A$-extensions $A$-équivalentes de $B$ par $L'$.

La définition du produit fibré montre que lorsque l'on a un morphisme
(18.2.4.1) de $A$-extensions, il se \emph{factorise à travers} l'image
réciproque de $E'$ par $v$~; de façon précise, il existe un
$A$-homomorphisme unique $u_0:E\to F=E'\times_{B'}B$ rendant commutatif le
diagramme
\[
  \xymatrix{
    0\ar[r] & L\ar[r]^j\ar[d]^{w_0} & E\ar[r]^f\ar[d]^{u_0}
      & B\ar[r]\ar[d]^{1_B} & 0\\
    0\ar[r] & L_0\ar[r]\ar[d]^i\ar@{}[d]|{\wr} & F\ar[r]^{p_2}\ar[d]^{p_1}
      & B\ar[r]\ar[d]^v & 0\\
    0\ar[r] & L'\ar[r]_{j'} & E'\ar[r]_{f'} & B'\ar[r] & 0
  }
\]
où $w_0$ est la restriction de $u_0$ à $L$ et
$p_1\circ u_0=u$, $i\circ w_0=w$.
\end{env}

\begin{env}[18.2.6]
\label{0.18.2.6-fr}
Étudions en particulier les images réciproques d'extensions par des
$A$-homomorphismes \emph{surjectifs}. Considérons un $A$-homomorphisme
surjectif $v:B\to B'$, une $A$-extension $F$ de $B$ par un $B$-bimodule $L$,
et l'idéal bilatère $\mathfrak{R}$ de $B$, noyau de $v$ (qui peut être
considéré comme $F$-bimodule au moyen de l'augmentation $F\to B$)~; on a vu
(18.1.7) que tout homomorphisme de $F$-bimodules
$\theta:\mathfrak{R}\to F$ rendant commutatif le diagramme
\[
  \xymatrix{
    & \mathfrak{R}\ar[ld]^\theta\ar[d]\\
    F\ar[r] & B
  }
\]
détermine une extension $E'$ de $B'=B/\mathfrak{R}$ par $L$ dont l'image
réciproque par $v:B\to B'$ est équivalente à $F$, et que toute $A$-extension
$E'$ de $B'$ par $L$ ayant cette dernière propriété s'obtient ainsi (à une
$A$-équivalence près). En outre, pour que deux homomorphismes
$\theta_1,\theta_2$ de $F$-bimodules de $\mathfrak{R}$ dans $F$ donnent deux
  $A$-extensions $A$-équivalentes de $B$ par $L$, il faut et il suffit qu'il
existe une $A$-équivalence $u$ de la $A$-extension $F$ sur elle-même telle que
\oldpage[0\textsubscript{IV-1}]{57}
$\theta_2=u\circ\theta_1$~; cela résulte aussitôt de ce qu'on a vu dans
(18.2.5) et de la définition de la bijection canonique de $F$ sur le produit
fibré $E'\times_{B'}B$ (18.1.7).
\end{env}

\begin{env}[18.2.7]
\label{0.18.2.7-fr}
Considérons maintenant le carré de gauche de (18.2.4.1), et rappelons d'abord
la notion de \emph{somme amalgamée} dans la catégorie des $A$-bimodules~:
étant donnés trois $A$-bimodules $X$, $Y$, $Z$ et deux $A$-homomorphismes
$f:X\to Y$, $g:X\to Z$, la somme amalgamée $Y\oplus_XZ$ est \emph{limite
inductive} du système inductif formé par $X$, $Y$, $Z$ et les
$A$-homomorphismes $f$, $g$ dans la catégorie des $A$-bimodules
($0_{\mathrm{III}}$, 8.1.11). On définit ce $A$-bimodule comme quotient du
produit $Y\times Z$ par le sous-$A$-bimodule $M$ image de $X$ par
l'homomorphisme $x\mapsto(f(x),-g(x))$. Sa propriété caractéristique est que,
pour tout couple d'homomorphismes de $A$-bimodules $u:Y\to T$, $v:Z\to T$
tels que $u\circ f=v\circ g$, il existe un homomorphisme et un seul
$w:Y\oplus_XZ\to T$ tel que $u=w\circ j_1$ et $v=w\circ j_2$, où
$j_1:Y\to Y\oplus_XZ$ et $j_2:Z\to Y\oplus_XZ$ sont les applications
canoniques.
\end{env}

\begin{env}[18.2.8]
\label{0.18.2.8-fr}
Considérons maintenant une $A$-extension $E$ de $B$ par un $B$-bimodule
$L$~:
\[
  0\longrightarrow L\xrightarrow{j}E\xrightarrow{f}B\longrightarrow0
\]
et soit d'autre part $w:L\to L'$ un homomorphisme de $B$-bimodules. Soit $H$
le $A$-bimodule somme amalgamée $E\oplus_LL'$~; montrons comment on peut munir
ce $A$-bimodule d'une structure de $A$-anneau et définir une $A$-extension
\[
  0\longrightarrow L'\longrightarrow H\longrightarrow B\longrightarrow0.
\]

Notons pour cela que $L'$ est muni d'une structure de $E$-\emph{bimodule} au
moyen de l'homomorphisme d'augmentation $E\to B$~; on peut donc former la
$A$-extension triviale type $G=D_E(L')$ (18.2.3). Considérons alors
l'application $\theta:z\mapsto(j(z),-w(z))$ de $L$ dans $G$~; c'est un
homomorphisme de $G$-bimodules ($L$ étant considéré comme $G$-bimodule au
moyen de l'homomorphisme canonique $\rho:G\to E$). En effet, pour
$(x,z')\in G$ et $z\in L$, on a
$(j(z),-w(z))(x,z')=(j(z)x,j(z)z'-w(z)f(x))$~; or
$j(z)z'=f(j(z))z'=0$ par définition de la structure de $E$-bimodule sur
$L'$, et $w(z)f(x)=w(zf(x))$ par définition de la structure de $B$-bimodule
sur $L'$~; on vérifie de même que $\theta$ est un homomorphisme de
$G$-module à gauche. On peut alors appliquer au diagramme commutatif
\[
  \xymatrix@C=2pc@R=1.4pc{
    & & 0\ar[d] & & \\
    & & L\ar[ld]^\theta\ar[d]^j & & \\
    0\ar[r] & L'\ar[r] & G\ar[r]_p & E\ar[r]\ar[d]^f & 0\\
    & & & B\ar[d] & \\
    & & & 0 &
  }
\]
le résultat de (18.1.7). Comme $H=G/\theta(L)$ par définition, notre
assertion est une conséquence immédiate de (18.1.7).

\oldpage[0\textsubscript{IV-1}]{58}

On dit que la $A$-extension $H$ de $B$ par $L'$ est \emph{déduite de $E$ au
moyen de l'homomorphisme} $w:L\to L'$. Le caractère fonctoriel de la somme
amalgamée en chacun des sommandes montre en outre que si l'on a un morphisme
d'extensions
\[
  \xymatrix{
    0\ar[r] & L\ar[r]\ar[d]^{1_L}\ar@{}[d]|{\wr}
      & E_1\ar[r]\ar[d]^g & B_1\ar[r]\ar[d]^h & 0\\
    0\ar[r] & L\ar[r] & E_2\ar[r] & B_2\ar[r] & 0
  }
\]
on en déduit canoniquement un morphisme d'extensions
\[
  \xymatrix{
    0\ar[r] & L'\ar[r]\ar[d]^{1_{L'}}\ar@{}[d]|{\wr}
      & E_1\oplus_LL'\ar[r]\ar[d]^{g\oplus_L1_{L'}}
      & B_1\ar[r]\ar[d]^h & 0\\
    0\ar[r] & L'\ar[r] & E_2\oplus_LL'\ar[r] & B_2\ar[r] & 0.
  }
\]
En particulier, si $E_1$ et $E_2$ sont des $A$-extensions $A$-équivalentes de
$B$ par $L$, les extensions de $B$ par $L'$ qu'on en déduit au moyen de $w$
sont $A$-équivalentes.

Lorsque l'on a un morphisme (18.2.4.1) de $A$-extensions, il se \emph{factorise
à travers} la $A$-extension $H$ de $B$ par $L'$ déduite de $E$ au moyen de
l'homomorphisme $w:L\to L'$ ($L'$ étant considéré comme $B$-bimodule au moyen
de l'homomorphisme $v:B\to B'$)~: en effet, la définition de la somme amalgamée
montre qu'il existe un $A$-homomorphisme unique $u_0:H\to E'$ de
$A$-bimodules, rendant commutatif le diagramme
\[
  \xymatrix{
    0\ar[r] & L\ar[r]^j\ar[d]^{w_0} & E\ar[r]^f\ar[d]^{j_1}
      & B\ar[r]\ar[d]^{1_B}\ar@{}[d]|{\wr} & 0\\
    0\ar[r] & L'_0\ar[r]^{j_2}\ar[d]^{j_0}\ar@{}[d]|{\wr}
      & H\ar[r]^{f_0}\ar[d]^{u_0} & B\ar[r]\ar[d]^v & 0\\
    0\ar[r] & L'\ar[r]_{j'} & E'\ar[r]_{f'} & B'\ar[r] & 0
  }
\]
avec $u=u_0\circ j_1$, $w=j_0\circ w_0$, $j_1$ et $j_2$ étant les
homomorphismes canoniques~; on vérifie immédiatement que $u_0$ est aussi un
homomorphisme d'anneaux.

Notons enfin les propriétés fonctorielles relatives aux extensions triviales~:
\end{env}

\begin{proposition}[18.2.9]
\label{0.18.2.9-fr}
Soient $B$, $B'$ deux $A$-anneaux, $L$ un $B$-bimodule, $L'$ un
$B'$-bimodule, $v:B\to B'$ un $A$-homomorphisme d'anneaux, $w:L\to L'$ un
homomorphisme de $A$-bimodules tel que $(v,w)$ soit un di-homomorphisme de
bimodules. Alors, il existe un $A$-homomorphisme unique d'anneaux
$u:D_B(L)\to D_{B'}(L')$ rendant commutatif les diagrammes
\[
  \xymatrix@C=2.2pc@R=1.6pc{
    D_B(L)\ar[r]^u & D_{B'}(L') & & D_B(L)\ar[r]^u & D_{B'}(L')\\
    B\ar[u]\ar[r]_v & B'\ar[u] & & L\ar[u]\ar[r]_w & L'\ar[u]
  }
\]
où les flèches verticales sont les injections canoniques.
\end{proposition}

\oldpage[0\textsubscript{IV-1}]{59}

\begin{proof}
En effet, $u$ ne peut être que l'application
$(x,s)\mapsto(v(x),w(s))$ et il reste à vérifier que c'est un
$A$-homomorphisme d'anneaux, ce qui résulte trivialement de la définition
(18.2.3). On notera que $u$ rend aussi commutatif le diagramme
\[
  \xymatrix@C=2.2pc@R=1.5pc{
    D_B(L)\ar[r]^u\ar[d] & D_{B'}(L')\ar[d]\\
    B\ar[r]_v & B'
  }
\]
où cette fois les flèches verticales sont les augmentations.
\end{proof}

\begin{proposition}[18.2.10]
\label{0.18.2.10-fr}
Soient $B$ un $A$-anneau, $L$ un $B$-bimodule, $E$ une $A$-extension de $B$
par $L$. On définit une application bijective de l'ensemble $G$ des
homomorphismes de $B$-anneaux de $B$ dans $E$ (autrement dit, l'ensemble des
$A$-homomorphismes inverses à droite de l'augmentation $E\to B$) sur
l'ensemble $G'$ des $A$-équivalences de $D_B(L)$ sur $E$ en faisant
correspondre à tout $g\in G$ la $A$-équivalence
$g':(x,s)\mapsto g(x)+s$~; l'application réciproque fait correspondre à tout
$g'\in G'$ le $B$-homomorphisme $x\mapsto g'(x,0)$.
\end{proposition}

\begin{proof}
Ceci résulte aussitôt des définitions.
\end{proof}

\subsection{Le groupe des classes de $A$-extensions}
\label{subsection:0.18.3-fr}

\begin{env}[18.3.1]
\label{0.18.3.1-fr}
Considérons un $A$-anneau $B$ et un $B$-bimodule $L$ fixés~; alors la relation
«~$E$ et $E'$ sont $A$-équivalentes~» entre $A$-extensions $E$, $E'$ de $B$
par $L$ est une relation d'équivalence, et pour cette relation on peut parler
de l'\emph{ensemble} des classes de $A$-extensions $A$-équivalentes de $B$ par
$L$. Il suffit pour le voir de remarquer que si, pour tout $x\in B$, $c_x$ est
un élément de $E$ dont l'image dans $B$ est $x$, tout $z\in E$ s'écrit d'une
manière et d'une seule sous la forme $c_x+t$, où $t\in L$, et l'on peut écrire
\[
  c_{x+y}=c_x+c_y+\varphi(x,y),\qquad
  c_{xy}=c_xc_y+\psi(x,y),
\]
où $\varphi(x,y)$ et $\psi(x,y)$ sont des éléments de $L$, les applications
$\varphi$ et $\psi$ de $B\times B$ dans $L$ devant vérifier des conditions
exprimant que $E$ est un $A$-anneau, qu'il est inutile d'écrire ici. Toute
$A$-extension de $B$ par $L$ est donc $A$-équivalente à une $A$-extension dont
$B\times L$ est l'ensemble sous-jacent, d'où on tire aussitôt notre conclusion.
\end{env}

\begin{env}[18.3.2]
\label{0.18.3.2-fr}
Le $A$-anneau $B$ étant fixé, désignons provisoirement par $T(L)$ l'ensemble
des classes de $A$-extensions de $B$ par $L$. Pour tout $B$-homomorphisme
$w:L\to L'$ de $B$-bimodules, on définit canoniquement une application
$T(w):T(L)\to T(L')$ en faisant correspondre à la classe d'une $A$-extension
$E$ de $B$ par $L$ la classe de la $A$-extension $E\oplus_LL'$ déduite de $E$
par $w$, en vertu de (18.2.8). Si $w':L'\to L''$ est un second homomorphisme
de $B$-bimodules, on a en outre
\begin{equation}
\label{0.18.3.2.1-fr}
\tag{18.3.2.1}
  T(w'\circ w)=T(w')\circ T(w).
\end{equation}
En effet, on sait qu'il existe un isomorphisme canonique de $A$-bimodules
\[
  E\oplus_LL''\xrightarrow{\sim}(E\oplus_LL')\oplus_{L'}L''
\]
en vertu des propriétés générales de limites inductives (cf. par exemple
($\mathrm{I}$, 3.3.9)), et il est immédiat de vérifier que c'est bien une
$A$-équivalence de $A$-extensions, d'où

\oldpage[0\textsubscript{IV-1}]{60}

la relation (18.3.2.1). Si $\mathbf C_B$ désigne la catégorie des
$B$-bimodules, on voit qu'on a ainsi défini un \emph{foncteur covariant}
$T:\mathbf C_B\to\mathbf{Ens}$.
\end{env}

\begin{env}[18.3.3]
\label{0.18.3.3-fr}
Considérons maintenant une famille $(L_\alpha)_{\alpha\in I}$ de
$B$-bimodules et leur produit $L=\prod_{\alpha\in I}L_\alpha$~; les
projections $\operatorname{pr}_\alpha:L\to L_\alpha$ définissent des
applications
\[
  T(\operatorname{pr}_\alpha):T(L)\to T(L_\alpha),
\]
d'où une application canonique
\begin{equation}
\label{0.18.3.3.1-fr}
\tag{18.3.3.1}
  \prod_\alpha T(\operatorname{pr}_\alpha):T(L)
  \longrightarrow\prod_{\alpha\in I}T(L_\alpha).
\end{equation}
Nous allons voir que cette application est \emph{bijective}. En effet, pour
tout $\alpha\in I$, soit $E_\alpha$ une $A$-extension de $B$ par $L_\alpha$,
et soit $F$ le \emph{produit fibré} des $E_\alpha$ sur $B$ (18.1.3)~; il est
immédiat que $F$ est une $A$-extension de $B$ par
$L=\prod_{\alpha\in I}L_\alpha$ et que si l'on remplace chaque $E_\alpha$ par
une $A$-extension $A$-équivalente $E'_\alpha$, le produit fibré $F'$ des
$E'_\alpha$ sur $B$ est $A$-équivalente à $F$. On a ainsi défini une
application $\prod_{\alpha\in I}T(L_\alpha)\to T(L)$ et il est clair que cette
application est réciproque de (18.3.3.1), d'où notre assertion. On identifiera
canoniquement $T(L)$ à $\prod_{\alpha\in I}T(L_\alpha)$ par l'application
(18.3.3.1). On vérifie en outre immédiatement que si
$(L'_\alpha)_{\alpha\in I}$ est une seconde famille de $B$-bimodules et, pour
chaque $\alpha$, $w_\alpha:L_\alpha\to L'_\alpha$ un homomorphisme de
$B$-bimodules, alors, en posant
$w=\prod_\alpha w_\alpha:L\to L'=\prod_{\alpha\in I}L'_\alpha$, $T(w)$
s'identifie à $\prod_{\alpha\in I}T(w_\alpha)$ quand on fait l'identification
précédente.
\end{env}

\begin{env}[18.3.4]
\label{0.18.3.4-fr}
Cela étant, pour un $B$-bimodule $L$, l'addition est un homomorphisme
$s:L\times L\to L$ de $B$-bimodules, et il en est de même de la symétrie
$t:L\to L$ de la loi additive de $L$. On en déduit une loi de composition
\[
  T(s):T(L)\times T(L)\longrightarrow T(L)
\]
dans $T(L)$ en vertu de (18.3.3), et cette loi est une loi de \emph{groupe
commutatif} dont $T(t)$ est la symétrie, comme il résulte de la définition d'un
objet en groupes au moyen de diagrammes commutatifs
($\mathbf 0_{\mathrm{III}}$, 8.2.5 et 8.2.6). Nous désignerons par
$\operatorname{Exan}_A(B,L)$ le groupe commutatif ainsi défini et nous dirons
que c'est le \emph{groupe des classes de $A$-extensions de $B$ par $L$}.
\end{env}

\begin{env}[18.3.5]
\label{0.18.3.5-fr}
Désignons par $\mathbf K$ la catégorie dont les objets sont les triplets
$(A,B,L)$ où $A$ est un anneau, $B$ un $A$-anneau et $L$ un $B$-bimodule~;
les \emph{morphismes} de cette catégorie sont les triplets $(u,v,w)$ où
$u:A'\to A$ et $v:B'\to B$ sont deux homomorphismes d'anneaux rendant
commutatif le carré de gauche du diagramme
\[
  \xymatrix@C=2pc@R=1.5pc{
    A\ar[r] & B & L\ar[d]^w\\
    A'\ar[u]^u\ar[r] & B'\ar[u]^v & L'
  }
\]
où les flèches horizontales sont les homomorphismes structuraux~; enfin
$w:L\to L'$ est un homomorphisme de groupes commutatifs tel que l'on ait
$w(v(b')z)=b'w(z)$ et

\oldpage[0\textsubscript{IV-1}]{61}

$w(zv(b'))=w(z)b'$ quels que soient $z\in L$ et $b'\in B'$ (en d'autres
termes, $w$ est un homomorphisme de $B'$-bimodules lorsqu'on munit $L$ de la
structure de $B'$-bimodule définie par $v$). La composition des morphismes est
définie par
$(u',v',w')\circ(u,v,w)=(u\circ u',v\circ v',w'\circ w)$, qui se justifie
aussitôt. Nous nous proposons de montrer que
\begin{equation}
\label{0.18.3.5.1-fr}
\tag{18.3.5.1}
  (A,B,L)\longmapsto\operatorname{Exan}_A(B,L)
\end{equation}
est un \emph{foncteur covariant} de la catégorie $\mathbf K$ dans la catégorie
$\mathbf{Ab}$ des groupes commutatifs. Il s'agit donc pour tout triplet
$(u,v,w)$ comme ci-dessus, de définir un homomorphisme de groupes commutatifs
\[
  (u,v,w)_*:\operatorname{Exan}_A(B,L)
  \longrightarrow\operatorname{Exan}_{A'}(B',L').
\]
En vertu de la définition des morphismes dans $\mathbf K$, on peut écrire
\[
  (u,v,w)=(1_{A'},1_{B'},w)\circ(1_{A'},v,1_L)
  \circ(u,1_B,1_L)
\]
où, dans le premier facteur, $L$ est muni de sa structure de $B'$-bimodule
définie par $v$~; nous définirons donc d'abord $(u,v,w)_*$ lorsque deux des
homomorphismes $u$, $v$, $w$ sont réduits à l'identité.
\end{env}

\begin{env}[18.3.6]
\label{0.18.3.6-fr}
Nous prendrons d'abord pour $(1_A,1_B,w)_*$ l'application
\begin{equation}
\label{0.18.3.6.1-fr}
\tag{18.3.6.1}
  w_*:\operatorname{Exan}_A(B,L)\longrightarrow
  \operatorname{Exan}_A(B,L')
\end{equation}
notée $T(w)$ dans (18.3.2)~; il est immédiat de vérifier que c'est un
homomorphisme de groupes, cette propriété s'exprimant par la commutativité de
diagrammes, transformés par $T$ de diagrammes analogues pour $L$ et $L'$.

L'application $(1_A,v,1_L)_*$ est l'application
\begin{equation}
\label{0.18.3.6.2-fr}
\tag{18.3.6.2}
  v^*:\operatorname{Exan}_A(B,L)\longrightarrow
  \operatorname{Exan}_A(B',L)
\end{equation}
définie de la façon suivante~: si $E$ est une $A$-extension de $B$ par $L$,
on a vu que $E\times_BB'$ est une $A$-extension de $B'$ par $L$ (18.2.5) et
que si l'on remplace $E$ par une $A$-extension $A$-équivalente $E'$,
$E'\times_BB'$ est $A$-équivalente à $E\times_BB'$~; l'image par $v^*$ de la
classe de $E$ est la classe de $E\times_BB'$. On vérifie aussitôt que si
$w:L\to L'$ est un homomorphisme de $B$-bimodules, le diagramme
\begin{equation}
\label{0.18.3.6.3-fr}
\tag{18.3.6.3}
  \xymatrix@C=3pc@R=2.4pc{
    \operatorname{Exan}_A(B,L)\ar[r]^{v^*}\ar[d]_{w_*}
      & \operatorname{Exan}_A(B',L)\ar[d]^{w_*}\\
    \operatorname{Exan}_A(B,L')\ar[r]_{v^*}
      & \operatorname{Exan}_A(B',L')
  }
\end{equation}
est commutatif, $L$ et $L'$ étant considérés comme $B'$-bimodules au moyen
de $v$ dans la colonne de droite. Remplaçant $L$ et $L'$ respectivement par
$L\times L$ et $L$, et $w$ par l'addition $s$ dans $L$, on en conclut que
$v^*$ est bien un \emph{homomorphisme de groupes}.
\end{env}

\oldpage[0\textsubscript{IV-1}]{62}

\begin{env}[18.3.6.4]
\label{0.18.3.6.4-fr}
Enfin, l'application $(u,1_B,1_L)_*$ est l'application
\[
  u^*:\operatorname{Exan}_A(B,L)\longrightarrow
  \operatorname{Exan}_{A'}(B,L)
\]
qui s'obtient en faisant correspondre à une $A$-extension $E$ de $B$ par $L$
l'anneau $E$ considéré comme $A'$-\emph{anneau} au moyen de $u$ (18.1.1), qui
est évidemment une $A'$-extension de $B$ par $L$, $B$ étant aussi considéré
comme $A'$-anneau au moyen de $u$~; il est clair qu'une $A$-équivalence est
aussi une $A'$-équivalence, d'où l'application (18.3.6.4), qui, pour tout
homomorphisme $w:L\to L'$ de $B$-bimodules, rend encore commutatif le diagramme
\begin{equation}
\label{0.18.3.6.5-fr}
\tag{18.3.6.5}
  \xymatrix@C=3pc@R=2.4pc{
    \operatorname{Exan}_A(B,L)\ar[r]^{u^*}\ar[d]_{w_*}
      & \operatorname{Exan}_{A'}(B,L)\ar[d]^{w_*}\\
    \operatorname{Exan}_A(B,L')\ar[r]_{u^*}
      & \operatorname{Exan}_{A'}(B,L')
  }
\end{equation}
d'où l'on conclut comme ci-dessus que $u^*$ est un homomorphisme de groupes.

Cela étant, on pose
$(u,v,w)_*=(1_{A'},1_{B'},w)_*\circ(1_{A'},v,1_L)_*
\circ(u,1_B,1_L)_*$ et l'on vérifie facilement, en raison de la commutativité
des diagrammes (18.3.6.3) et (18.3.6.5), que l'on a
\[
  (u\circ u',v\circ v',w'\circ w)_*
  =(u',v',w')_*\circ(u,v,w)_*,
\]
ce qui achève de prouver que (18.3.4.1) est un foncteur.

L'existence de l'homomorphisme de groupes (18.3.6.1) montre en particulier
que si $E$ est une $A$-extension \emph{triviale} de $B$ par $L$, l'extension
$E\oplus_LL'$ de $B$ par $L'$ définie dans (18.2.8) est aussi triviale, ce que
l'on vérifie d'ailleurs sans peine de façon directe.

D'autre part, pour tout élément $z$ du \emph{centre} $Z$ de $B$, l'homothétie
$h_z:y\mapsto yz=zy$ est un endomorphisme du $B$-bimodule $L$, donc $(h_z)_*$
est un endomorphisme du groupe commutatif $\operatorname{Exan}_A(B,L)$, et par
fonctorialité, ces endomorphismes définissent sur
$\operatorname{Exan}_A(B,L)$ une structure canonique de $Z$-\emph{module}.
\end{env}

\begin{env}[18.3.7]
\label{0.18.3.7-fr}
Soient $A$, $A'$ deux anneaux, $u:A'\to A$ un homomorphisme, $B$ un
$A$-anneau et $L$ un $B$-bimodule. Le \emph{noyau} de l'homomorphisme de
groupes
\[
  u^*:\operatorname{Exan}_A(B,L)\longrightarrow
  \operatorname{Exan}_{A'}(B,L)
\]
est formé par définition des classes de $A$-extensions de $B$ par $L$ qui
sont $A'$-\emph{triviales} quand on les considère comme $A'$-extensions au
moyen de $u$. On désigne ce noyau par la notation
$\operatorname{Exan}_{A/A'}(B,L)$ quand cela ne prête pas à confusion.

Si $\Lambda$ est un anneau, et si on a un diagramme commutatif de
$\Lambda$-\emph{homomorphismes} de $\Lambda$-anneaux
\begin{equation}
\label{0.18.3.7.1-fr}
\tag{18.3.7.1}
  \xymatrix@C=2pc@R=1.5pc{
    B'\ar[r] & B\\
    A'\ar[u]\ar[r] & A\ar[u]
  }
\end{equation}

\oldpage[0\textsubscript{IV-1}]{63}

on en déduit canoniquement des homomorphismes
\begin{equation}
\label{0.18.3.7.2-fr}
\tag{18.3.7.2}
  \operatorname{Exan}_{A/\Lambda}(B,L)\longrightarrow
  \operatorname{Exan}_{A'/\Lambda}(B,L)\longrightarrow
  \operatorname{Exan}_{A'/\Lambda}(B',L)
\end{equation}
qui proviennent de la commutativité du diagramme
\[
  \xymatrix@C=2.8pc@R=2.3pc{
    \operatorname{Exan}_A(B,L)\ar[r]\ar[d]
      & \operatorname{Exan}_{A'}(B,L)\ar[r]\ar[d]
      & \operatorname{Exan}_{A'}(B',L)\ar[d]\\
    \operatorname{Exan}_{\Lambda}(B,L)\ar[r]
      & \operatorname{Exan}_{\Lambda}(B,L)\ar[r]
      & \operatorname{Exan}_{\Lambda}(B',L)
  }
\]
où les flèches sont déduites de celles de (18.3.7.1) par fonctorialité.
\end{env}

\begin{proposition}[18.3.8]
\label{0.18.3.8-fr}
Soient $B$ un anneau, $\mathfrak J$ un idéal bilatère de $B$,
$C=B/\mathfrak J$ l'anneau quotient~; $\mathfrak J/\mathfrak J^2$ est alors
canoniquement muni d'une structure de $C$-bimodule. Pour tout $C$-bimodule
$L$, soit $\operatorname{Hom}_C(\mathfrak J/\mathfrak J^2,L)$ le groupe
additif des homomorphismes de $C$-bimodules de $\mathfrak J/\mathfrak J^2$
dans $L$. On définit alors un isomorphisme canonique de groupes commutatifs
\begin{equation}
\label{0.18.3.8.1-fr}
\tag{18.3.8.1}
  \eta_L:\operatorname{Hom}_C(\mathfrak J/\mathfrak J^2,L)
  \xrightarrow{\sim}\operatorname{Exan}_B(C,L)
\end{equation}
en faisant correspondre à tout $C$-homomorphisme
$w:\mathfrak J/\mathfrak J^2\to L$ (qui est \emph{a fortiori} un
$B$-homomorphisme) la classe de l'extension
$(B/\mathfrak J^2)\oplus_{(\mathfrak J/\mathfrak J^2)}L$ déduite de
l'extension $B/\mathfrak J^2$ de $C$ par $\mathfrak J/\mathfrak J^2$ au moyen
de $w$~; l'isomorphisme réciproque fait correspondre à la classe d'une
$B$-extension $E$ de $C$ par $L$ l'homomorphisme $w$ tel que le composé
$\mathfrak J\to\mathfrak J/\mathfrak J^2\xrightarrow{w}L$ soit la restriction
à $\mathfrak J$ de l'homomorphisme structural $B\to E$.
\end{proposition}

\begin{proof}
Soit $F=B/\mathfrak J^2$, qui est une $B$-extension de $C$ par
$\mathfrak J/\mathfrak J^2$. Pour toute $B$-extension
\[
  0\longrightarrow L\xrightarrow{j}E\xrightarrow{p}C\longrightarrow0
\]
de $C$ par $L$, l'homomorphisme structural $f:B\to E$ est tel que le composé
$B\xrightarrow{f}E\xrightarrow{p}C$ soit l'homomorphisme canonique
$B\to B/\mathfrak J=C$. Comme l'image de $\mathfrak J$ par $p\circ f$ est
nulle, $f(\mathfrak J)$ est contenu dans le noyau de $p$, c'est-à-dire
$j(L)$, et comme $j(L)$ est de carré nul, on a $f(\mathfrak J^2)=0$~; donc
$f$ se factorise en
$B\to B/\mathfrak J^2=F\xrightarrow{u}E$, et si $w$ est la restriction de
$u$ à $\mathfrak J/\mathfrak J^2$, on a un diagramme commutatif
\[
  \xymatrix@C=2pc@R=1.5pc{
    0\ar[r] & \mathfrak J/\mathfrak J^2\ar[r]\ar[d]^w
      & F\ar[r]^g\ar[d]^u
      & C\ar[r]\ar[d]^{1_C}\ar@{}[d]|{\wr} & 0\\
    0\ar[r] & L\ar[r]_j & E\ar[r]_p & C\ar[r] & 0
  }
\]

autrement dit $(u,1_C,w)$ est un morphisme d'extensions (18.2.4). On en
déduit un morphisme d'extensions
\[
  \xymatrix@C=2pc@R=1.5pc{
    0\ar[r] & L\ar[r]\ar[d]^{1_L}\ar@{}[d]|{\wr}
      & E'\ar[r]\ar[d]^{u'}
      & C\ar[r]\ar[d]^{1_C}\ar@{}[d]|{\wr} & 0\\
    0\ar[r] & L\ar[r]_j & E\ar[r]_p & C\ar[r] & 0
  }
\]

\oldpage[0\textsubscript{IV-1}]{64}

où $E'=F\oplus_{(\mathfrak J/\mathfrak J^2)}L$ (18.2.8)~; il s'agit de
prouver que $u'$ est une $B$-\emph{équivalence}, ce qui établira que $\eta_L$
est bijectif. En vertu de la construction faite dans (18.2.8), tout revient à
prouver (compte tenu de (18.1.7)) que l'image réciproque $G$ de l'extension
$E$ de $C$ par l'homomorphisme canonique $g:F\to C$ est $F$-triviale, ce qui
est évident puisque $g$ se factorise en
$F\xrightarrow{u}E\xrightarrow{p}C$ (18.1.6).

Il reste à voir que $\eta_L$ est un homomorphisme de \emph{groupes}~; or, pour
tout $B$-homomorphisme $h:L\to L'$, il est immédiat que le diagramme
\[
  \xymatrix@C=3pc@R=2.4pc{
    \operatorname{Hom}_C(\mathfrak J/\mathfrak J^2,L)
      \ar[r]^{\eta_L}\ar[d]_{\operatorname{Hom}(1,h)}
      & \operatorname{Exan}_B(C,L)\ar[d]^{h_*}\\
    \operatorname{Hom}_C(\mathfrak J/\mathfrak J^2,L')
      \ar[r]_{\eta_{L'}}^{\sim}
      & \operatorname{Exan}_B(C,L')
  }
\]
est commutatif. Il suffit d'appliquer cette remarque à l'homomorphisme
$L\times L\to L$ définissant l'addition pour conclure.
\end{proof}

\subsection{Extensions d'algèbres}
\label{subsection:0.18.4-fr}

\begin{env}[18.4.1]
\label{0.18.4.1-fr}
Soit $A$ un anneau \emph{commutatif}. La catégorie des $A$-\emph{algèbres} est
alors une sous-catégorie pleine de celle des $A$-anneaux, caractérisée par le
fait que les homomorphismes structuraux $\rho:A\to B$ sont \emph{centraux}.

Si $B$ est une $A$-algèbre, $L$ un $B$-bimodule, il est immédiat que la
$A$-extension triviale type $D_B(L)$ est une $A$-algèbre. De même, dans la
construction de (18.2.5) (resp. de (18.2.8)), si $B$, $B'$ et $E'$ (resp.
$B$ et $E$) sont des $A$-algèbres, il en est de même de
$E'\times_{B'}B$ (resp. de $E\oplus_LL'$). Enfin, il est clair que si $B$ est
une $A$-algèbre et $E$ une $A$-extension de $B$ par un $B$-bimodule $L$ qui
est une $A$-algèbre, toute $A$-extension de $B$ par $L$, $A$-\emph{équivalente}
à $E$, est aussi une $A$-algèbre. On déduit alors aussitôt de la définition
(18.3.4) que les classes de $A$-extensions équivalentes de $B$ par $L$ qui
sont des $A$-\emph{algèbres} forment un \emph{sous-groupe}, noté
$\operatorname{Exal}_A(B,L)$, de $\operatorname{Exan}_A(B,L)$. Soit
$\mathbf K'$ la sous-catégorie pleine de la catégorie $\mathbf K$ définie dans
(18.3.5), dont les objets $(A,B,L)$ sont tels que $A$ soit commutatif et $B$
une $A$-algèbre. Alors ce qui précède montre que
\[
  (A,B,L)\longmapsto\operatorname{Exal}_A(B,L)
\]
est un \emph{foncteur covariant} de $\mathbf K'$ dans $\mathbf{Ab}$. Les
résultats de (18.3.7) sont inchangés lorsqu'on remplace partout
$\operatorname{Exan}$ par $\operatorname{Exal}$.
\end{env}

\begin{env}[18.4.2]
\label{0.18.4.2-fr}
Supposons toujours l'anneau $A$ commutatif. Les remarques de (18.4.1) restent
valables lorsque l'on remplace «~algèbre~» par «~algèbre commutative~» et
«~bimodule~» par «~module~». Si $B$ est une $A$-algèbre commutative et $L$ un
$B$-module,

\oldpage[0\textsubscript{IV-1}]{65}

les classes de $A$-extensions équivalentes de $B$ par $L$ qui sont des
$A$-\emph{algèbres commutatives} (ou, ce qui revient au même, des
$A$-\emph{anneaux commutatifs}) forment un \emph{sous-groupe} de
$\operatorname{Exal}_A(B,L)$, noté $\operatorname{Exalcom}_A(B,L)$. Si
$\mathbf K''$ est la sous-catégorie pleine de $\mathbf K'$ formée des triplets
$(A,B,L)$ où $A$ et $B$ sont commutatifs et $L$ un $B$-module,
\[
  (A,B,L)\longmapsto\operatorname{Exalcom}_A(B,L)
\]
est encore un \emph{foncteur covariant} de $\mathbf K''$ dans $\mathbf{Ab}$. On
peut aussi dans (18.3.7) remplacer partout $\operatorname{Exan}$ par
$\operatorname{Exalcom}$. Enfin, si dans (18.3.8), on suppose que $B$ soit un
anneau \emph{commutatif} et que $L$ soit un $C$-\emph{module}, le même
raisonnement donne un isomorphisme canonique
\begin{equation}
\label{0.18.4.2.1-fr}
\tag{18.4.2.1}
  \operatorname{Hom}_C(\mathfrak J/\mathfrak J^2,L)
  \xrightarrow{\sim}\operatorname{Exalcom}_B(C,L)
\end{equation}
où le premier membre est le groupe des homomorphismes de $C$-\emph{module}.
\end{env}

\begin{env}[18.4.3]
\label{0.18.4.3-fr}
Soit $A$ un anneau commutatif. Un cas important d'extensions de
$A$-algèbres est formé des $A$-algèbres $E$, extensions d'une $A$-algèbre $B$
par un $B$-bimodule $L$ telles que $E$, en tant que $A$-\emph{module}, soit une
extension \emph{triviale} du $A$-module $B$ par le $A$-module $L$~; autrement
dit, la suite exacte de $A$-\emph{modules}
$0\to L\to E\to B\to0$ est \emph{scindée}. On dit alors que $E$ est une
\emph{extension de Hochschild} de $B$ par $L$. Ce sera toujours le cas lorsque
$B$ est un $A$-module \emph{projectif}, et en particulier lorsque $A$ est un
\emph{corps} commutatif. En tant que $A$-module, on peut identifier $E$ à
$B\times L$, la multiplication dans $E$ étant donnée par
$(x,s)(y,t)=(xy,xt+sy+f(x,y))$ avec $f(x,y)\in L$. Si l'on écrit que cette
multiplication définit sur $B\times L$ une structure de $A$-algèbre, on trouve
(M, XIV, 2) que $f$ doit être une application $A$-\emph{bilinéaire} de
$B\times B$ dans $L$, telle que
\begin{equation}
\label{0.18.4.3.1-fr}
\tag{18.4.3.1}
  f(xy,z)+f(x,y)z=xf(y,z)+f(x,yz),
\end{equation}
autrement dit, $f$ est un $2$-\emph{cocycle} sur $B$ à valeurs dans $L$, au
sens de Hochschild~; pour que l'extension $E$ soit $A$-triviale, il faut et il
suffit que l'on ait
\begin{equation}
\label{0.18.4.3.2-fr}
\tag{18.4.3.2}
  f(x,y)=xg(y)-g(xy)+g(x)y,
\end{equation}
où $g$ est une application $A$-\emph{linéaire} de $B$ dans $L$, autrement dit
$f$ doit être un $2$-\emph{cobord} au sens de Hochschild. On en déduit aussitôt
que les classes d'extensions de Hochschild de $B$ par $L$ forment un
sous-groupe de $\operatorname{Exal}_A(B,L)$, isomorphe au \emph{groupe de
cohomologie} de Hochschild $\mathrm H_A^2(B,L)$.

Si $B$ est une $A$-algèbre \emph{commutative}, et $L$ un $B$-module, les
extensions de Hochschild \emph{commutatives} de $B$ par $L$ correspondent aux
$2$-cocycles \emph{symétriques}, c'est-à-dire tels que $f(y,x)=f(x,y)$. Les
classes d'extensions de Hochschild commutatives de $B$ par $L$ forment donc un
sous-groupe de $\operatorname{Exalcom}_A(B,L)$, isomorphe au sous-groupe de
$\mathrm H_A^2(B,L)$ image du groupe des $2$-cocycles symétriques, que nous
noterons $\mathrm H_A^2(B,L)^s$.
\end{env}

\begin{env}[18.4.4]
\label{0.18.4.4-fr}
Bornons-nous au cas où $A$ et $B$ sont \emph{commutatifs} et $L$ un
$B$-module, et rappelons dans ce cas la définition équivalente des groupes de
Hochschild $\mathrm H_A^i(B,L)$ (M, IX, 6). On considère un complexe
$\mathbf P_\bullet=(\mathbf P_n)_{n\geqslant0}$ de $B$-modules, où
$\mathbf P_n=B^{\otimes(n+1)}=B\otimes_A B\otimes_A\cdots\otimes_A B$
($n+1$ fois), la structure de $B$-module étant définie

\oldpage[0\textsubscript{IV-1}]{66}

par
\[
  x(y_1\otimes y_2\otimes\cdots\otimes y_{n+1})
  =(xy_1)\otimes y_2\otimes\cdots\otimes y_{n+1}\,;
\]
la dérivation, de degré $-1$, $d_n:\mathbf P_n\to\mathbf P_{n-1}$, est
définie par
\begin{align*}
  d_n(x_1\otimes x_2\otimes\cdots\otimes x_{n+1})
    &=(x_1x_2)\otimes x_3\otimes\cdots\otimes x_{n+1}
      -x_1\otimes(x_2x_3)\otimes\cdots\otimes x_{n+1}+\cdots\\
    &\quad+(-1)^{n-1}x_1\otimes x_2\otimes\cdots\otimes x_{n-1}
      \otimes(x_nx_{n+1})
      +(-1)^n(x_{n+1}x_1)\otimes x_2\otimes\cdots\otimes x_n,
\end{align*}
qui est bien $B$-linéaire puisque $B$ est commutatif.

Un $2$-cocycle de ce complexe à valeurs dans $L$ est une application
$B$-linéaire $h$ de $B\otimes B\otimes B$ dans $L$ telle que
$h(d_3(x\otimes y\otimes z\otimes t))=0$~; mais comme
$h(x\otimes y\otimes z)=xh(1\otimes y\otimes z)$, le cocycle $h$ est déterminé
par l'application $A$-bilinéaire
$(y,z)\longmapsto f(y,z)=h(1\otimes y\otimes z)$ de $B\times B$ dans $L$, et,
en écrivant la condition précédente pour $h$, on retrouve pour $f$ la
condition (18.4.3.1). De même, un $2$-cobord sera une application de la forme
$x\otimes y\otimes z\longmapsto h'(d_2(x\otimes y\otimes z))$, où
$h':B\otimes_A B\to L$ est linéaire, et ici encore $h'$ est déterminé par
l'application $B$-linéaire $g:y\longmapsto h'(1\otimes y)$ de $B$ dans $L$~;
on obtient alors
$h'(d_2(1\otimes x\otimes y))=xg(y)-g(xy)+yg(x)$, ce qui redonne (18.4.3.2).
On procède de même pour tout $i$, et l'on voit ainsi que l'on a
\begin{equation}
\label{0.18.4.4.1-fr}
\tag{18.4.4.1}
  \mathrm H_A^i(B,L)=\mathrm H^i(\operatorname{Hom}_B(\mathbf P_\bullet,L)).
\end{equation}
\end{env}

\begin{env}[18.4.5]
\label{0.18.4.5-fr}
Sous les conditions de (18.4.4), on peut interpréter de la même façon le
groupe $\mathrm H_A^2(B,L)^s$ (18.4.3). Pour cela, modifions le complexe
$\mathbf P_\bullet$ en degré $3$, en considérant un nouveau complexe
\[
  \mathbf P'_\bullet:\mathbf P'_3\xrightarrow{d'_3}\mathbf P_2
  \xrightarrow{d_2}\mathbf P_1\,;
\]
nous prendrons
$\mathbf P'_3=\mathbf P_3\oplus(B\otimes_A B\otimes_A B)$, $d'_3$ coïncidant
avec $d_3$ sur $\mathbf P_3$, et étant donné dans $B\otimes_A B\otimes_A B$
par
\[
  d'_3(x\otimes y\otimes z)=x\otimes y\otimes z-x\otimes z\otimes y.
\]

La relation $d_2d'_3=0$ résulte de la commutativité de $B$. Avec les notations
introduites ci-dessus, un $2$-cocycle de $\mathbf P'_\bullet$ correspond
maintenant à une application $A$-bilinéaire $f:B\times B\to L$ qui est
\emph{symétrique} et vérifie (18.4.3.1)~; par suite, on a
\[
  \mathrm H_A^2(B,L)^s
  =\mathrm H^2(\operatorname{Hom}_B(\mathbf P'_\bullet,L)).
\]
\end{env}

\begin{env}[18.4.6]
\label{0.18.4.6-fr}
Dans le cas particulier où l'on considère un \emph{corps} commutatif $k$, une
\emph{extension} $K$ de $k$, on a $\operatorname{Ext}_K^1(M,L)=0$ quels que
soient les $K$-espaces vectoriels $L$, $M$, et par suite (M, VI, 3.3 $a$)), on
a un isomorphisme canonique
\begin{equation}
\label{0.18.4.6.1-fr}
\tag{18.4.6.1}
  \mathrm H_k^i(K,L)\xrightarrow{\sim}
  \operatorname{Hom}_K(\mathrm H_i(\mathbf P_\bullet),L)
\end{equation}
et de même
\begin{equation}
\label{0.18.4.6.2-fr}
\tag{18.4.6.2}
  \mathrm H_k^2(K,L)^s\xrightarrow{\sim}
  \operatorname{Hom}_K(\mathrm H_2(\mathbf P'_\bullet),L).
\end{equation}
\end{env}

\subsection{Cas des anneaux topologiques}
\label{subsection:0.18.5-fr}

\begin{env}[18.5.1]
\label{0.18.5.1-fr}
Soient $A$, $B$ deux anneaux \emph{topologiques} dont la topologie est
\emph{linéaire}, $\rho:A\to B$ un homomorphisme continu, $L$ un $B$-bimodule
topologique, et supposons qu'il existe un idéal bilatère \emph{ouvert}
$\mathfrak R_0$ de $B$ tel que
$\mathfrak R_0.L=L.\mathfrak R_0=0$, de sorte que $L$ peut être

\oldpage[0\textsubscript{IV-1}]{67}

considéré comme un $(B/\mathfrak R)$-bimodule pour tout idéal ouvert
$\mathfrak R\subset\mathfrak R_0$. Soient alors
$\mathfrak R\subset\mathfrak R_0$ un idéal bilatère ouvert de $B$,
$\mathfrak J$ un idéal bilatère ouvert de $A$ tel que
$\rho(\mathfrak J)\subset\mathfrak R$, de sorte que $B/\mathfrak R$ peut être
considéré comme un $(A/\mathfrak J)$-anneau discret. La remarque précédente
prouve que le groupe
$\operatorname{Exan}_{A/\mathfrak J}(B/\mathfrak R,L)$ est défini~; en outre,
si $\mathfrak R'\subset\mathfrak R$, $\mathfrak J'\subset\mathfrak J$ sont
deux idéaux bilatères ouverts de $A$ et $B$ respectivement tels que
$\rho(\mathfrak J')\subset\mathfrak R'$, on a par (18.3.5.1) un homomorphisme
canonique
\begin{equation}
\label{0.18.5.1.1-fr}
\tag{18.5.1.1}
  \operatorname{Exan}_{A/\mathfrak J}(B/\mathfrak R,L)
  \longrightarrow
  \operatorname{Exan}_{A/\mathfrak J'}(B/\mathfrak R',L).
\end{equation}

L'ensemble des couples d'idéaux ouverts $(\mathfrak J,\mathfrak R)$ tels que
$\rho(\mathfrak J)\subset\mathfrak R\subset\mathfrak R_0$ est évidemment
ordonné filtrant à droite pour la relation «~$\mathfrak J\supset\mathfrak J'$
et $\mathfrak R\supset\mathfrak R'$~» et les applications (18.5.1.1)
définissent un \emph{système inductif} de groupes additifs ayant cet ensemble
pour ensemble d'indices. On pose, par abus de notation (car il ne s'agit plus
d'un groupe en correspondance biunivoque naturelle avec un ensemble
d'extensions)
\begin{equation}
\label{0.18.5.1.2-fr}
\tag{18.5.1.2}
  \operatorname{Exantop}_A(B,L)
  =\varinjlim\operatorname{Exan}_{A/\mathfrak J}(B/\mathfrak R,L).
\end{equation}

Dire que le second membre de (18.5.1.2) est \emph{nul} signifie donc que, pour
tout couple d'idéaux ouverts $\mathfrak J\subset A$, $\mathfrak R\subset B$
tels que $\rho(\mathfrak J)\subset\mathfrak R\subset\mathfrak R_0$ et toute
$(A/\mathfrak J)$-extension $E$ de $B/\mathfrak R$ par $L$, il existe deux
idéaux ouverts $\mathfrak J'\subset\mathfrak J$,
$\mathfrak R'\subset\mathfrak R$ tels que
$\rho(\mathfrak J')\subset\mathfrak R'$ et que l'image réciproque par
l'homomorphisme $B/\mathfrak R'\to B/\mathfrak R$ de $E$ soit \emph{triviale}.

Nous laissons au lecteur la définition analogue des limites inductives
$\operatorname{Exaltop}_A(B,L)$, $\operatorname{Exalcotop}_A(B,L)$ à partir de
$\operatorname{Exal}_{A/\mathfrak J}(B/\mathfrak R,L)$ et de
$\operatorname{Exalcom}_{A/\mathfrak J}(B/\mathfrak R,L)$ pour le cas où $A$
est commutatif et $B$ une $A$-algèbre (resp. une $A$-algèbre commutative)
topologique.
\end{env}

\begin{env}[18.5.2]
\label{0.18.5.2-fr}
Si l'on a un diagramme commutatif d'homomorphismes continus d'anneaux
\[
  \xymatrix@C=3pc@R=2pc{
    B'\ar[r] & B\\
    A'\ar[u]\ar[r] & A\ar[u]
  }
\]
on en déduit canoniquement deux homomorphismes de groupes additifs
\[
  \operatorname{Exantop}_A(B,L)
  \longrightarrow\operatorname{Exantop}_{A'}(B,L)
  \longrightarrow\operatorname{Exantop}_{A'}(B',L)
\]
par passage à la limite inductive à partir de (18.3.5.1).

En vertu de l'exactitude du foncteur $\varinjlim$ dans la catégorie des groupes
commutatifs, le noyau de l'homomorphisme
\[
  \operatorname{Exantop}_A(B,L)
  \longrightarrow\operatorname{Exantop}_{A'}(B,L)
\]
est la limite inductive des noyaux des homomorphismes
\[
  \operatorname{Exan}_{A/\mathfrak J}(B/\mathfrak R,L)
  \longrightarrow
  \operatorname{Exan}_{A'/\mathfrak J'}(B/\mathfrak R,L)
\]
où on a pris pour $\mathfrak J'$ l'image réciproque de $\mathfrak J$~; on note
ce noyau $\operatorname{Exantop}_{A/A'}(B,L)$. On définit de même
$\operatorname{Exaltop}_{A/A'}(B,L)$ et
$\operatorname{Exalcotop}_{A/A'}(B,L)$. Enfin, si on a un homomorphisme

\oldpage[0\textsubscript{IV-1}]{68}

continu de $B$-bimodules $L\to L'$, on en déduit canoniquement un
homomorphisme de groupes additifs
\[
  \operatorname{Exantop}_A(B,L)\longrightarrow
  \operatorname{Exantop}_A(B,L')
\]
par passage à la limite inductive à partir de (18.3.6.1).
\end{env}

\begin{env}[18.5.3]
\label{0.18.5.3-fr}
Étant donnés un anneau topologique $C$ et deux $C$-bimodules topologiques
$M$, $N$,\linebreak on note $\operatorname{Hom.\ cont}_C(M,N)$ le groupe additif des
$C$-homomorphismes \emph{continus} de $M$ dans $N$.
\end{env}

\begin{lemma}[18.5.3.1]
\label{0.18.5.3.1-fr}
Soient $C$ un anneau topologique, $E$, $L$ deux $C$-bimodules topologiques~;
on suppose que les topologies sont linéaires, que $L$ est discret et annulé
par un idéal bilatère ouvert de $C$. Alors on a un isomorphisme canonique
\begin{equation}
\label{0.18.5.3.2-fr}
\tag{18.5.3.2}
  \varinjlim\operatorname{Hom}_{C/\mathfrak R}(E/V,L)
  \xrightarrow{\sim}\operatorname{Hom.\ cont}_C(E,L)
\end{equation}
où dans le premier membre la limite inductive est prise suivant l'ensemble
ordonné filtrant des couples $(\mathfrak R,V)$ tels que $\mathfrak R$ soit un
idéal bilatère ouvert de $C$, $V$ un sous-$C$-bimodule ouvert de $E$, tels que
$\mathfrak R.L=L.\mathfrak R=0$, $\mathfrak R.E\subset V$,
$E.\mathfrak R\subset V$.
\end{lemma}

\begin{proof}
Comme $C/\mathfrak R$ et $E/V$ sont discrets, on a des homomorphismes
canoniques
\[
  w_{\mathfrak R,V}:\operatorname{Hom}_{C/\mathfrak R}(E/V,L)
  \longrightarrow\operatorname{Hom.\ cont}_C(E,L)
\]
formant un système inductif, d'où un homomorphisme (18.5.3.2) par passage à la
limite inductive. Comme l'homomorphisme $E/V'\to E/V$ est surjectif pour
$V'\subset V$, il résulte aussitôt de la définition que l'homomorphisme
\[
  \operatorname{Hom}_{C/\mathfrak R}(E/V,L)
  \longrightarrow\operatorname{Hom}_{C/\mathfrak R}(E/V',L)
\]
(avec $\mathfrak R.L=L.\mathfrak R=0$, $\mathfrak R.E\subset V'$,
$E.\mathfrak R\subset V'$) est injectif, et il en est évidemment de même de
l'homomorphisme
\[
  \operatorname{Hom}_{C/\mathfrak R}(E/V,L)
  \longrightarrow\operatorname{Hom}_{C/\mathfrak R'}(E/V,L)
\]
pour $\mathfrak R'\subset\mathfrak R$~; on en conclut que l'homomorphisme
(18.5.3.2) est injectif. D'autre part, si $u$ est un $C$-homomorphisme continu
de $E$ dans $L$, son noyau est un sous-bimodule ouvert $V_0$ de $E$, et si
$\mathfrak R_0$ est un idéal bilatère ouvert de $C$ tel que
$\mathfrak R_0.L=L.\mathfrak R_0=0$ et $\mathfrak R_0.E\subset V_0$,
$E.\mathfrak R_0\subset V_0$, il est clair que $u$ est l'image canonique d'un
$(C/\mathfrak R_0)$-homomorphisme de $E/V_0$ dans $L$, donc (18.5.3.2) est
surjectif.
\end{proof}

Cela étant, la proposition (18.3.8) se généralise comme suit aux anneaux
topologiques~:

\begin{proposition}[18.5.4]
\label{0.18.5.4-fr}
Soient $B$ un anneau topologique linéairement topologisé, $\mathfrak J$ un
idéal bilatère de $B$, $C=B/\mathfrak J$ l'anneau topologique quotient~;
$\mathfrak J/\mathfrak J^2$ (où $\mathfrak J$ est muni de la topologie induite
par celle de $B$ et $\mathfrak J/\mathfrak J^2$ de la topologie quotient de
celle de $\mathfrak J$) est alors canoniquement muni d'une structure de
$C$-bimodule topologique. Pour tout $C$-bimodule discret $L$ annulé par un
idéal ouvert de $C$, il existe alors un isomorphisme canonique
\begin{equation}
\label{0.18.5.4.1-fr}
\tag{18.5.4.1}
  \operatorname{Hom.\ cont}_C(\mathfrak J/\mathfrak J^2,L)
  \xrightarrow{\sim}\operatorname{Exantop}_B(C,L).
\end{equation}
\end{proposition}

\begin{proof}
En effet, pour tout idéal ouvert $\mathfrak R$ de $B$ tel que
$(\mathfrak J+\mathfrak R)/\mathfrak J$ annule $L$, on a, d'après (18.3.9),
un isomorphisme canonique
\[
  \operatorname{Hom}_{B/(\mathfrak J+\mathfrak R)}
    ((\mathfrak J+\mathfrak R)/(\mathfrak J^2+\mathfrak R),L)
  \xrightarrow{\sim}
  \operatorname{Exan}_{B/\mathfrak R}(B/(\mathfrak J+\mathfrak R),L)
\]
et il suffit de passer à la limite inductive, en tenant compte de (18.5.3.1).
\end{proof}

% FIN SOURCE ega4/ega0-18-fr.tex

%
% DEBUT SOURCE ega4/ega0-19-fr.tex
\section{Algèbres formellement lisses et anneaux de Cohen}
\label{section:0.19-fr}
\label{0.19-fr}

\setcounter{subsection}{-1}
\subsection{Introduction}
\label{subsection:0.19.0-fr}
\label{0.19.0-fr}

\oldpage[0\textsubscript{IV-1}]{69}

\begin{env}[19.0.1]
\label{0.19.0.1-fr}
Au chapitre IV, nous introduirons et étudierons entre autres une classe
importante de morphismes de préschémas, les morphismes \emph{lisses}\footnote{Aussi
appelés morphismes \emph{simples} dans certains travaux récents (en s'inspirant
de la terminologie classique de «~point simple~»)~; cette terminologie prête
cependant à confusion, notamment en théorie des groupes algébriques.}. Une de
leurs propriétés fondamentales (et qui, jointe à une condition de finitude, les
caractérise) est une propriété de \emph{relèvement des morphismes}~: si
$f:X\to Y$ est un morphisme lisse, $g:Y'\to Y$ un morphisme, alors pour tout
morphisme $h:Y''\to Y'$ faisant de $Y''$ un préschéma «~peu différent~» de
$Y'$, \emph{tout $Y$-morphisme $Y''\to X$ se factorise} en
$Y''\xrightarrow{h}Y'\to X$. De façon précise, bornons-nous au cas où
$Y=\operatorname{Spec}(A)$, $X=\operatorname{Spec}(B)$ sont \emph{affines}~;
$B$ est alors appelée une $A$-algèbre \emph{formellement lisse} si, pour toute
$A$-algèbre $C$ et tout idéal \emph{nilpotent} $\mathfrak J$ de $C$, tout
$A$-homomorphisme $B\to C/\mathfrak J$ \emph{se relève} en
$B\to C\to C/\mathfrak J$. En d'autres termes, l'application
\[
  \operatorname{Hom}_A(B,C)\longrightarrow
  \operatorname{Hom}_A(B,C/\mathfrak J)
\]
est \emph{surjective}. Dans beaucoup d'applications, $X$ apparaîtra comme objet
représentant un \emph{foncteur contravariant représentable} $Y'\mapsto F(Y')$
de la catégorie des $Y$-préschémas dans celle des ensembles, de sorte
($\mathbf{0}_{\mathrm{III}}$, 8.1.8) que l'on aura
$F(Y')=\operatorname{Hom}_Y(Y',X)$. Dans le cas affine, si l'on pose
$F^0(C)=F(\operatorname{Spec}(C))$, la vérification du fait que, sous les
conditions ci-dessus, $F^0(C)\to F^0(C/\mathfrak J)$ est \emph{surjectif}
(qui peut se faire même \emph{sans savoir} que $F$ est représentable) établira
que $f$ est lisse, pourvu que la condition additionnelle de finitude soit
remplie.
\end{env}

\begin{env}[19.0.2]
\label{0.19.0.2-fr}
Pour les \emph{algèbres topologiques} sur un anneau topologique $A$, il y a une
notion analogue d'algèbre \emph{formellement lisse} que nous ne préciserons pas
ici (cf.~19.3.1). L'étude de ces notions est faite d'abord d'un point de vue
élémentaire au §~19, puis, à l'aide des propriétés des différentielles qui
seront développées aux §§~20 et 21, des théorèmes plus délicats seront prouvés
au §~22. Résumons ici les principaux résultats sur les algèbres formellement
lisses~:

\medskip\noindent\textbf{I.} --- Soient $A$, $B$ deux anneaux locaux
noethériens, $\varphi:A\to B$ un homomorphisme local, $k$ le corps résiduel de
$A$, et soit $B_0=B\otimes_A k$~; on munit $A$, $B$ et $B_0$ de leurs
topologies préadiques et $k$ de la topologie discrète. Alors, pour que $B$ soit
une $A$-algèbre formellement lisse, il faut et il suffit que $B$ soit un
$A$-module \emph{plat} et que $B_0$ soit une $k$-algèbre \emph{formellement
lisse} (19.7.1). Ce théorème ramène donc la lissité formelle, pour les anneaux
locaux noethériens, à la même question pour les anneaux locaux noethériens qui
sont des algèbres sur un \emph{corps}.

\medskip\noindent\textbf{II.} --- Soient $k$ un corps, $A$ un anneau local
noethérien qui est une $k$-algèbre. Pour que $A$ soit formellement lisse, il
faut et il suffit que $A$ soit \emph{géométriquement régulier} sur $k$,
c'est-à-dire que pour toute extension \emph{finie} $k'$ de $k$, l'anneau
semi-local $A\otimes_k k'$ soit

\oldpage[0\textsubscript{IV-1}]{70}

\noindent régulier (22.5.8)~; le \emph{complété} $\widehat A$ de $A$ est alors
un anneau isomorphe à un anneau de séries formelles
$K[[T_1,\ldots,T_n]]$ (19.6.5). En outre, la structure de $k$-algèbre de $A$,
lorsque $A$ est supposé être \emph{complet} et formellement lisse sur $k$, est
entièrement déterminée par le corps résiduel $K$ de $A$ et la \emph{dimension}
de $A$~; celle-ci peut d'ailleurs être arbitraire à condition de satisfaire à
l'inégalité
\[
  \dim(A)\geq \operatorname{rg}_k(\Gamma_{K/k}),
\]
où $\Gamma_{K/k}$ est le «~module d'imperfection~» de $K$ (22.2.6).

En particulier, pour qu'une \emph{extension} $K$ de $k$ soit formellement
lisse, il faut et il suffit que $K$ soit une extension \emph{séparable} de $k$
(19.6.1).

\medskip\noindent\textbf{III.} --- Soient $A$ un anneau local noethérien,
$\mathfrak J$ un idéal de $A$ distinct de $A$, $A_0=A/\mathfrak J$, $B_0$ un
anneau local noethérien \emph{complet}, $A_0\to B_0$ un homomorphisme local
faisant de $B_0$ une $A_0$-algèbre formellement lisse. Alors il existe un
anneau local noethérien \emph{complet} $B$, un homomorphisme local $A\to B$
faisant de $B$ un $A$-module plat, et un $A_0$-\emph{isomorphisme}
\[
  B\otimes_A A_0\xrightarrow{\sim}B_0
\]
(de sorte que, par I), $B$ est une $A$-algèbre \emph{formellement lisse})~; en
outre $B$ est déterminé par ces propriétés à isomorphisme près (19.7.2). Ce
théorème contient en particulier les \emph{théorèmes de Cohen} sur la structure
des anneaux locaux noethériens \emph{complets} (19.8), qui joueront un rôle
important dans les §§~6 et 7 du chapitre IV.

\medskip\noindent\textbf{IV.} --- L'intérêt de l'étude des anneaux locaux
noethériens formellement lisses sur un autre provient de la caractérisation
«~ponctuelle~» suivante des morphismes lisses~: si $X$ et $Y$ sont des
préschémas localement noethériens, $f:X\to Y$ un morphisme localement de type
fini, alors, pour que $f$ soit lisse, il faut et il suffit que pour tout
$x\in X$, l'anneau $\mathcal O_x$ soit formellement lisse sur
$\mathcal O_{f(x)}$. En particulier, si $Y=\operatorname{Spec}(k)$, où $k$ est
un corps parfait, dire que $f$ est lisse équivaut à dire que $X$ est un
préschéma \emph{régulier}.

\medskip\noindent\textbf{V.} --- Enfin, nous verrons aux §§~20, 21 et 22 que
la notion d'algèbre formellement lisse s'introduit de façon naturelle dans la
théorie des \emph{différentielles de Kähler}, les deux théories s'éclairant
mutuellement.
\end{env}

\begin{env}[19.0.3]
\label{0.19.0.3-fr}
Dans tout ce paragraphe et les suivants, les anneaux et modules topologiques
seront supposés \emph{linéairement topologisés} ($\mathbf{0}_{\mathrm I}$,
7.1.1)~; les anneaux topologiques considérés seront supposés \emph{commutatifs},
sauf mention expresse du contraire. Rappelons que si $A$ et $B$ sont deux
anneaux topologiques, $\rho:A\to B$ un homomorphisme d'anneaux définissant sur
$B$ une structure de $A$-algèbre, on dit que $B$ est une \emph{$A$-algèbre
topologique} si $\rho$ est \emph{continu} pour les topologies envisagées.

Pour abréger, dans un anneau topologique $A$ (resp. un $A$-module topologique
$M$), nous dirons «~\emph{système fondamental d'idéaux} (resp.
\emph{sous-modules}) \emph{ouverts}~» au lieu de «~système fondamental de
voisinages de 0 formé d'idéaux (resp. sous-modules)~».

Étant donnés un anneau topologique $A$ et un $A$-module $M$, les ensembles
$\mathfrak J M$, où $\mathfrak J$ parcourt un système fondamental d'idéaux
ouverts, forment un système fondamental de sous-modules ouverts d'une topologie
sur $M$ faisant de $M$ un $A$-module topologique, et que l'on dit \emph{déduite}
de la topologie de $A$.

Soit $M$ un $A$-module topologique dont la topologie est \emph{moins fine} que
la topologie déduite de celle de $A$~; alors, si $N$ est un sous-module
\emph{ouvert} de $M$, le $A$-module \emph{discret} $M/N$ est \emph{annulé par
un idéal ouvert} de $A$, car par hypothèse il existe un tel idéal $\mathfrak R$
tel que $\mathfrak R.M\subset N$.

\oldpage[0\textsubscript{IV-1}]{71}

Si $M$ et $N$ sont deux $A$-modules topologiques dont les topologies sont toutes
deux déduites de celle de $A$, alors \emph{tout} $A$-homomorphisme $u:M\to N$
est \emph{continu}, car pour tout voisinage $V$ de 0 dans $N$, il existe par
hypothèse un idéal ouvert $\mathfrak J$ de $A$ tel que $\mathfrak JN\subset V$,
et l'on a donc $u(\mathfrak JM)\subset\mathfrak JN\subset V$.

Lorsque $B$ est une $A$-algèbre topologique, la topologie sur $B$ \emph{déduite}
de celle de $A$ est \emph{plus fine} que la topologie donnée, car pour tout
idéal ouvert $\mathfrak R$ de $B$, il y a par hypothèse un idéal ouvert
$\mathfrak J$ de $A$ tel que $\mathfrak JB\subset\mathfrak R$.
\end{env}

\subsection{Épimorphismes et monomorphismes formels}
\label{subsection:0.19.1-fr}

\begin{proposition}[19.1.1]
\label{0.19.1.1-fr}
Soient $A$ un anneau topologique, $M$, $N$ deux $A$-modules topologiques,
$(W_\lambda)$ un système fondamental de sous-modules ouverts dans $N$,
$u:M\to N$ un $A$-homomorphisme continu, $\widehat M$, $\widehat N$ les
séparés complétés de $M$ et $N$, $\widehat u:\widehat M\to\widehat N$ le
prolongement continu de $u$ aux séparés complétés.
\begin{enumerate}
  \item[(i)] Les conditions suivantes sont équivalentes~:
  \begin{enumerate}
    \item[a)] $u(M)$ est \emph{dense} dans $N$.
    \item[a$'$)] $\widehat u(\widehat M)$ est \emph{dense} dans $\widehat N$.
    \item[a$''$)] Pour tout $\lambda$, l'homomorphisme composé
    $M\xrightarrow{u}N\to N/W_\lambda$ est \emph{surjectif}.
  \end{enumerate}
  \item[(ii)] Les conditions suivantes sont équivalentes~:
  \begin{enumerate}
    \item[b)] L'image réciproque par $u$ de la topologie de $N$ est égale à la
    topologie de $M$.
    \item[b$'$)] $\widehat u$ est un \emph{isomorphisme} du $\widehat A$-module
    topologique $\widehat M$ sur le sous-$\widehat A$-module topologique
    $\widehat u(\widehat M)$ de $\widehat N$ (qui est nécessairement
    \emph{fermé}).
    \item[b$''$)] Les $u^{-1}(W_\lambda)$ forment un système fondamental de
    voisinages de 0 dans $M$.
  \end{enumerate}
\end{enumerate}
Cela résulte immédiatement de la définition des séparés complétés, et du fait
que $N/W_\lambda$ est discret.
\end{proposition}

\begin{env}[19.1.2]
\label{0.19.1.2-fr}
Lorsque les conditions équivalentes de (i) (resp. (ii)) dans (19.1.1) sont
remplies, on dit que $u$ est un \emph{épimorphisme formel} (resp. un
\emph{monomorphisme formel}). On dit que $u$ est un \emph{bimorphisme formel}
s'il est à la fois un monomorphisme formel et un épimorphisme formel~; il
revient au même, en vertu de (19.1.1), de dire que $\widehat u$ est un
\emph{isomorphisme} du $\widehat A$-module topologique $\widehat M$ sur le
$\widehat A$-module topologique $\widehat N$.
\end{env}

\begin{proposition}[19.1.3]
\label{0.19.1.3-fr}
Soient $A$ un anneau topologique, $M$, $N$ deux $A$-modules topologiques,
$u:M\to N$ un $A$-homomorphisme continu. On suppose qu'il existe deux systèmes
fondamentaux de sous-modules ouverts, $(V_\lambda)$, $(W_\lambda)$ dans $M$ et
$N$ respectivement, ayant même ensemble d'indices et tels que
$u(V_\lambda)\subset W_\lambda$ pour tout $\lambda$~; soit
$u_\lambda:M/V_\lambda\to N/W_\lambda$ l'homomorphisme déduit de $u$ par
passage aux quotients. Alors~:
\begin{enumerate}
  \item[(i)] Pour que $u$ soit un épimorphisme formel, il faut et il suffit que,
  pour tout $\lambda$, $u_\lambda$ soit surjectif.
  \item[(ii)] Pour que $u$ soit un monomorphisme formel, il faut et il suffit
  que, pour tout $\lambda$, il existe un $\mu$ tel que
  $V_\mu\subset V_\lambda$ et que $\operatorname{Ker}(u_\mu)$ soit contenu
  dans le noyau $V_\lambda/V_\mu$ de l'application canonique
  $M/V_\mu\to M/V_\lambda$.
  \item[(iii)] Pour que $u$ soit un bimorphisme formel, il faut et il suffit
  que, pour tout $\lambda$, $u_\lambda$ soit surjectif et qu'il existe un
  $\mu$ tel que $V_\mu\subset V_\lambda$ et que l'homomorphisme canonique
  $M/V_\mu\to M/V_\lambda$

  \oldpage[0\textsubscript{IV-1}]{72}

  \noindent se factorise en
  $M/V_\mu\xrightarrow{u_\mu}N/W_\mu\to M/V_\lambda$ (où
  l'homomorphisme $N/W_\mu\to M/V_\lambda$ est nécessairement unique).
\end{enumerate}
\end{proposition}

\begin{proof}
L'assertion (i) est immédiate~; d'autre part,
$\operatorname{Ker}(u_\mu)=u^{-1}(W_\mu)/V_\mu$, et le noyau de l'application
canonique $M/V_\mu\to M/V_\lambda$ est $V_\lambda/V_\mu$~; la condition de
(ii) revient donc à dire que $u^{-1}(W_\mu)\subset V_\lambda$, et l'assertion
(ii) en découle aussitôt~; enfin, lorsque $u_\mu$ est surjectif, il revient au
même de dire que $\operatorname{Ker}(u_\mu)$ est contenu dans
$V_\lambda/V_\mu$, ou de dire que $M/V_\mu\to M/V_\lambda$ se factorise en
$v\circ u_\mu$, où $v$ est un homomorphisme
$N/W_\mu\to M/V_\lambda$, puisque $N/W_\mu$ s'identifie alors à
$(M/V_\mu)/\operatorname{Ker}(u_\mu)$.
\end{proof}

\begin{corollary}[19.1.4]
\label{0.19.1.4-fr}
Soient $A$ un anneau topologique, $M$, $N$ deux $A$-modules topologiques dont
les topologies sont déduites de celle de $A$, $u:M\to N$ un épimorphisme
formel. Soit $(\mathfrak J_\lambda)$ un système fondamental d'idéaux ouverts
dans $A$. Pour que $u$ soit un bimorphisme formel, il faut et il suffit que pour
tout $\lambda$, l'homomorphisme
$u_\lambda:M/\mathfrak J_\lambda M\to N/\mathfrak J_\lambda N$ déduit de $u$
par passage aux quotients soit bijectif.
\end{corollary}

\begin{proof}
On a en effet $u(\mathfrak J_\lambda M)\subset\mathfrak J_\lambda N$ et on
peut appliquer le critère (19.1.3, (iii))~; mais si l'on a une factorisation
\[
  M/\mathfrak J_\mu M\xrightarrow{u_\mu}
  N/\mathfrak J_\mu N\xrightarrow{v}M/\mathfrak J_\lambda M,
\]
l'image de $\mathfrak J_\lambda M/\mathfrak J_\mu M$ par $u_\mu$ est
$\mathfrak J_\lambda N/\mathfrak J_\mu N$, donc l'image de
$\mathfrak J_\lambda N/\mathfrak J_\mu N$ par $v$ est 0 et $v$ se factorise en
\[
  N/\mathfrak J_\mu N\longrightarrow N/\mathfrak J_\lambda N
  \xrightarrow{v'}M/\mathfrak J_\lambda M~;
\]
mais alors l'automorphisme identique de $M/\mathfrak J_\lambda M$ se factorise
en
\[
  M/\mathfrak J_\lambda M\xrightarrow{u_\lambda}
  N/\mathfrak J_\lambda N\xrightarrow{v'}M/\mathfrak J_\lambda M,
\]
ce qui montre que $u_\lambda$ est injectif, et comme on sait déjà qu'il est
surjectif, cela prouve le corollaire.
\end{proof}

\begin{proposition}[19.1.5]
\label{0.19.1.5-fr}
Soient $A$ un anneau topologique, $M$, $N$ deux $A$-modules topologiques,
$u:M\to N$ un $A$-homomorphisme continu. Les conditions suivantes sont
équivalentes~:
\begin{enumerate}
  \item[a)] Pour tout $A$-module topologique discret $E$ et tout
  $A$-homomorphisme continu $v:M\to E$, il existe un $A$-homomorphisme continu
  $w:N\to E$ tel que $v=w\circ u$.
  \item[b)] Pour tout sous-module ouvert $V'$ de $M$, il existe un sous-module
  ouvert $W''$ de $N$, un sous-module ouvert
  $V''\subset V'\cap u^{-1}(W'')$ et un homomorphisme
  $h:N/W''\to M/V'$ tels que l'homomorphisme canonique
  $M/V''\to M/V'$ se factorise en
  \[
    M/V''\xrightarrow{u''}N/W''\xrightarrow{h}M/V',
  \]
  où $u''$ est déduit de $u$ par passage aux quotients.
\end{enumerate}
\end{proposition}

\begin{proof}
Pour voir que a) entraîne b), il suffit d'appliquer a) à $E=M/V'$ et à
l'homomorphisme canonique $v:M\to M/V'$~; alors $W''=w^{-1}(0)$ est un
sous-module ouvert de $N$ et $V''=u^{-1}(W'')$ un sous-module ouvert de $M$
contenu dans $V'$~; ces sous-modules et l'homomorphisme
$h:N/W''\to M/V'$ déduit de $w$ par passage au quotient répondent à la
question. Inversement, pour montrer que b) entraîne a), on peut se limiter au
cas où $v$ est surjectif, en remplaçant $E$ par $v(M)$~; alors
$V'=\operatorname{Ker}(v)$ est un sous-module ouvert de $M$, donc $E$ est
isomorphe au $A$-module discret $M/V'$, et en appliquant b), on obtient

\oldpage[0\textsubscript{IV-1}]{73}

le $A$-homomorphisme continu $w$ cherché en prenant le composé
$N\to N/W''\xrightarrow{h}M/V'$, le diagramme
\[
  \xymatrix{
    M\ar[r]^u\ar[d] & N\ar[d]\\
    M/V''\ar[r]_{u''} & N/W''
  }
\]
étant commutatif.
\end{proof}

Lorsque les conditions équivalentes de (19.1.5) sont remplies, nous dirons que
$u$ est \emph{formellement inversible à gauche}~; comme la condition b) de
(19.1.5) implique que $u^{-1}(W'')\subset V'$, $u$ est alors un monomorphisme
formel. La terminologie est en outre justifiée par les corollaires suivants~:

\begin{corollary}[19.1.6]
\label{0.19.1.6-fr}
S'il existe un $\widehat A$-homomorphisme continu
$s:\widehat N\to\widehat M$ tel que
$s\circ\widehat u=1_{\widehat M}$, alors $u$ est formellement inversible à
gauche.
\end{corollary}

On notera que $\widehat u$ est alors un isomorphisme topologique de
$\widehat M$ sur un \emph{facteur direct topologique} de $\widehat N$. Pour
prouver (19.1.6), il suffit de noter qu'avec les notations de (19.1.5, a)),
$v:M\to E$ se prolonge en un $\widehat A$-homomorphisme continu
$\widehat v:\widehat M\to E$ puisque $E$ est discret~; soient
$j:M\to\widehat M$ et $j':N\to\widehat N$ les homomorphismes canoniques~;
alors, en posant $w=\widehat v\circ s\circ j'$, on a
$w\circ u=\widehat v\circ s\circ\widehat u\circ j=\widehat v\circ j=v$.

\begin{corollary}[19.1.7]
\label{0.19.1.7-fr}
Supposons que les topologies de $M$ et de $N$ soient déduites de celle de $A$,
et soit $(\mathfrak J_\lambda)$ un système fondamental d'idéaux ouverts dans
$A$. Pour que $u$ soit formellement inversible à gauche il faut et il suffit
que, pour tout $\lambda$, l'homomorphisme
$u_\lambda:M/\mathfrak J_\lambda M\to N/\mathfrak J_\lambda N$, déduit de
$u$ par passage aux quotients, soit inversible à gauche (autrement dit, soit
un isomorphisme de $M/\mathfrak J_\lambda M$ sur un facteur direct de
$N/\mathfrak J_\lambda N$).
\end{corollary}

En effet, la condition est suffisante, car, pour
$V'=\mathfrak J_\lambda M$ dans la condition b) de (19.1.5), on répond à la
question en prenant $W''=\mathfrak J_\lambda N$, $V''=V'$ et $h$ tel que
$h\circ u_\lambda$ soit l'identité dans $M/\mathfrak J_\lambda M$.
Inversement, si $u$ est formellement inversible à gauche, alors, pour tout
$\lambda$, il y a, en vertu de (19.1.5, b)), un $\mu$ tel que
$\mathfrak J_\mu\subset\mathfrak J_\lambda$ et un homomorphisme
$h:N/\mathfrak J_\mu N\to M/\mathfrak J_\lambda M$ tels que l'homomorphisme
canonique $M/\mathfrak J_\mu M\to M/\mathfrak J_\lambda M$ se factorise en
\[
  M/\mathfrak J_\mu M\xrightarrow{u_\mu}N/\mathfrak J_\mu N
  \xrightarrow{h}M/\mathfrak J_\lambda M~;
\]
mais comme
$h(\mathfrak J_\lambda(N/\mathfrak J_\mu N))
\subset\mathfrak J_\lambda(M/\mathfrak J_\lambda M)=0$, $h$ se factorise
canoniquement en
\[
  N/\mathfrak J_\mu N\longrightarrow N/\mathfrak J_\lambda N
  \xrightarrow{s}M/\mathfrak J_\lambda M,
\]
et il est immédiat que $s$ est inverse à gauche de $u_\lambda$.

\begin{proposition}[19.1.8]
\label{0.19.1.8-fr}
Soit $A$ un anneau topologique admettant un système fondamental dénombrable
décroissant $(\mathfrak J_n)$ d'idéaux ouverts. Soient $M$, $N$ deux
$A$-modules topologiques dont les topologies sont déduites de celle de $A$,
$u:M\to N$ un $A$-homomorphisme. Posons
$M_n=M/\mathfrak J_nM$, $N_n=N/\mathfrak J_nN$ et soit
$u_n:M_n\to N_n$ le $(A/\mathfrak J_n)$-homomorphisme déduit de $u$ par
passage aux quotients. Supposons que, pour tout $n$,
$P_n=\operatorname{Coker}(u_n)$ soit un $(A/\mathfrak J_n)$-module projectif
et que $M$ soit séparé et complet. Alors, pour que $u$ soit formellement
inversible à gauche, il faut et il suffit que $u$ soit inversible à gauche (et
$u$ est alors un isomorphisme topologique de $M$ sur un facteur direct
topologique de $N$).
\end{proposition}

\oldpage[0\textsubscript{IV-1}]{74}

\begin{proof}
En vertu de (19.1.7), on a des diagrammes commutatifs
\[
  \xymatrix{
    0\ar[r] & M_n\ar[r]^{u_n} & N_n\ar[r]^{p_n} & P_n\ar[r] & 0\\
    0\ar[r] & M_{n+1}\ar[u]^f\ar[r]_{u_{n+1}} &
      N_{n+1}\ar[u]^g\ar[r]_{p_{n+1}} & P_{n+1}\ar[u]^h\ar[r] & 0
  }
\]
où les lignes sont des suites exactes \emph{scindées}, autrement dit, il existe
pour chaque $n$ un homomorphisme $s_n:P_n\to N_n$ tel que
$p_n\circ s_n=1_{P_n}$. Nous allons montrer que l'on peut, par récurrence sur
$n$, définir un homomorphisme $s'_n:P_n\to N_n$ tel que
$p_n\circ s'_n=1_{P_n}$ et que les diagrammes
\[
  \xymatrix{
    P_n\ar[r]^{s'_n} & N_n\\
    P_{n+1}\ar[u]^h\ar[r]_{s'_{n+1}} & N_{n+1}\ar[u]^g
  }
\]
soient commutatifs. En effet, $g\circ s_{n+1}-s'_n\circ h$ est un
homomorphisme de $P_{n+1}$ dans
\[
  u_n(M_n)=u_{n+1}(M_{n+1})/
  \mathfrak J_nu_{n+1}(M_{n+1})~;
\]
comme $P_{n+1}$ est un $(A/\mathfrak J_{n+1})$-module \emph{projectif}, cet
homomorphisme se factorise en
\[
  P_{n+1}\xrightarrow{t_{n+1}}u_{n+1}(M_{n+1})
  \longrightarrow
  u_{n+1}(M_{n+1})/\mathfrak J_nu_{n+1}(M_{n+1}),
\]
d'où l'on conclut aussitôt que $s'_{n+1}=s_{n+1}-t_{n+1}$ répond à la
question. Cela étant, de la décomposition de $N_n$ en somme directe de
$u_n(M_n)$ et de $s'_n(P_n)$, on déduit aussitôt un homomorphisme
$w_n:N_n\to M_n$ inverse à gauche de $u_n$ et tel que les diagrammes
\[
  \xymatrix{
    N_n\ar[r]^{w_n} & M_n\\
    N_{n+1}\ar[u]^g\ar[r]_{w_{n+1}} & M_{n+1}\ar[u]^f
  }
\]
soient commutatifs. Le système projectif $(w_n)$ admet alors une limite
projective $\widehat w:\widehat N\to\widehat M=M$, d'où par composition avec
l'homomorphisme canonique $N\to\widehat N$, un homomorphisme $w:N\to M$ tel
que, pour tout $n$, l'endomorphisme $(w\circ u)_n$ de
$M_n=M/\mathfrak J_nM$ déduit de $w\circ u$ par passage aux quotients soit
l'identité~; cela entraîne que $w\circ u=1_M$ puisque $M$ est séparé et
complet.
\end{proof}

\begin{proposition}[19.1.9]
\label{0.19.1.9-fr}
Soient $A$ un anneau topologique préadmissible ($\mathbf{0}_{\mathrm I}$,
7.1.2), $\mathfrak Q$ un idéal de définition de $A$,
$(\mathfrak J_\lambda)$ un système fondamental d'idéaux ouverts de $A$.
Soient $M$, $N$ deux $A$-modules topologiques dont les topologies sont déduites
de celles de $A$, et tels que pour tout $\lambda$,
$N/\mathfrak J_\lambda N$ soit un $(A/\mathfrak J_\lambda)$-module projectif
(cf. (19.2.3)). Soit $u:M\to N$ un $A$-homomorphisme. Alors les conditions
suivantes sont équivalentes~:
\begin{enumerate}
  \item[a)] $u$ est formellement inversible à gauche.

  \oldpage[0\textsubscript{IV-1}]{75}

  \item[b)] L'homomorphisme
  $u_0:M/\mathfrak QM\to N/\mathfrak QN$ déduit de $u$ par passage aux
  quotients est inversible à gauche.
\end{enumerate}
\end{proposition}

\begin{proof}
On a vu (19.1.7) que la condition a) équivaut à dire que $u_\lambda$ est
inversible à gauche pour tout $\lambda$~; comme $\mathfrak Q$ est un idéal
ouvert, donc contient un $\mathfrak J_\lambda$, on en déduit aussitôt que
$u_0$ est inversible à gauche. Pour montrer inversement que b) entraîne a),
notons que pour tout $\lambda$, $\mathfrak Q/\mathfrak J_\lambda$ est par
hypothèse un idéal nilpotent de $A/\mathfrak J_\lambda$. Notre assertion
résultera de la proposition suivante.
\end{proof}

\begin{proposition}[19.1.10]
\label{0.19.1.10-fr}
Soient $A$ un anneau, $M$, $N$ deux $A$-modules, $N$ étant projectif,
$u:M\to N$ un $A$-homomorphisme. Soit $\mathfrak J$ un idéal de $A$~; on
suppose vérifiée l'une des conditions suivantes~:
\begin{enumerate}
  \item[(i)] $\mathfrak J$ est nilpotent.
  \item[(ii)] $\mathfrak J$ est contenu dans le radical de $A$ et $M$ est de
  type fini.
\end{enumerate}
Alors, pour que $u$ soit inversible à gauche, il faut et il suffit que
l'homomorphisme
$u_0:M/\mathfrak JM\to N/\mathfrak JN$ de $(A/\mathfrak J)$-modules, déduit
de $u$ par passage aux quotients, soit inversible à gauche.
\end{proposition}

\begin{proof}
La condition étant évidemment nécessaire, prouvons qu'elle est suffisante.
Soit $v_0$ un inverse à gauche de $u_0$~; l'homomorphisme composé
$N\to N/\mathfrak JN\xrightarrow{v_0}M/\mathfrak JM$ se factorise en
$N\xrightarrow{v}M\to M/\mathfrak JM$ puisque $N$ est projectif~; alors
$w=v\circ u$ est un endomorphisme de $M$ tel que l'endomorphisme $w_0$ de
$M/\mathfrak JM$ déduit par passage aux quotients soit l'identité, et il suffit
de prouver que $w$ est lui-même bijectif (car alors $w^{-1}\circ v$ sera un
inverse à gauche de $u$). Distinguons maintenant les deux cas.

(i) Pour tout $n$ on a le diagramme commutatif
\[
  \xymatrix{
    (\mathfrak J^n/\mathfrak J^{n+1})\otimes_{A/\mathfrak J}
      (M/\mathfrak JM)\ar[r]\ar[d]_{1\otimes\operatorname{gr}^0(w)} &
      \mathfrak J^nM/\mathfrak J^{n+1}M\ar[d]^{\operatorname{gr}^n(w)}\\
    (\mathfrak J^n/\mathfrak J^{n+1})\otimes_{A/\mathfrak J}
      (M/\mathfrak JM)\ar[r] & \mathfrak J^nM/\mathfrak J^{n+1}M
  }
\]
où les flèches horizontales sont \emph{surjectives}, et comme
$\operatorname{gr}^0(w)=w_0$ est l'identité, il en est de même de
$\operatorname{gr}^n(w)$ qui, \emph{a fortiori}, est bijectif. La filtration
$\mathfrak J$-préadique sur $M$ étant finie puisque $\mathfrak J$ est
nilpotent, on en conclut que $w$ est bijectif (Bourbaki, \emph{Alg. comm.},
chap.~III, §~2, n\textsuperscript{o}~8, cor.~3 du th.~1).

(ii) Il suffit de montrer que pour tout idéal maximal $\mathfrak m$ de $A$,
l'endomorphisme $w_{\mathfrak m}$ de $M_{\mathfrak m}$ est bijectif
(Bourbaki, \emph{Alg. comm.}, chap.~II, §~3, n\textsuperscript{o}~3, th.~1)
et comme $\mathfrak JA_{\mathfrak m}\subset\mathfrak mA_{\mathfrak m}$, et
$A_{\mathfrak m}/\mathfrak JA_{\mathfrak m}=(A/\mathfrak J)_{\mathfrak m}$,
on est ramené à prouver la proposition lorsque $A$ est un anneau local. En
outre, on peut supposer que $\mathfrak J$ est l'idéal maximal
$\mathfrak J_0$ de $A$, car si $u_0$ est inversible

\oldpage[0\textsubscript{IV-1}]{76}

à gauche, il en est de même de
$u_{00}:M/\mathfrak J_0M\to N/\mathfrak J_0M$ obtenu par tensorisation de
$u_0$ avec $1_{A/\mathfrak J_0}$, puisque l'on a
\[
  M/\mathfrak J_0M=(M/\mathfrak JM)\otimes_{A/\mathfrak J}(A/\mathfrak J_0).
\]
Supposons donc $\mathfrak J$ maximal, de sorte que $A/\mathfrak J$ est un
corps. Il suffit évidemment de montrer que $M$ est un $A$-module \emph{libre}
sous les conditions de l'énoncé~: en effet, $\det(w_0)$ est alors l'image
canonique de $\det(w)$ dans $A/\mathfrak J$, donc $\det(w)$ n'appartient pas à
l'idéal $\mathfrak J$ et est par suite inversible. Or le
$(A/\mathfrak J)$-espace vectoriel $M/\mathfrak JM$ étant libre de type fini,
il y a un $A$-module de type fini $L$ et un $A$-homomorphisme $f:L\to M$ tel
que l'homomorphisme $f_0:L/\mathfrak JL\to M/\mathfrak JM$ déduit de $f$ par
passage aux quotients soit bijectif. Comme $M$ est de type fini, on en conclut
tout d'abord que $f$ est surjectif par le lemme de Nakayama (Bourbaki,
\emph{Alg. comm.}, chap.~II, §~3, n\textsuperscript{o}~2, cor.~1 de la
prop.~4)~; en outre, si $g=u\circ f$, l'homomorphisme
$g_0:L/\mathfrak JL\to N/\mathfrak JN$ déduit par passage aux quotients est
inversible à gauche, et puisqu'ici $L$ est libre, la remarque du début prouve
que $g$ est lui-même inversible à gauche~; or cela entraîne évidemment que
$f$ est \emph{injectif}, ce qui achève la démonstration.
\end{proof}

Mentionnons en passant les corollaires suivants de (19.1.10)~:

\begin{corollary}[19.1.11]
\label{0.19.1.11-fr}
Soient $A$ un anneau local, $k$ son corps résiduel, $M$ un $A$-module de type
fini, $N$ un $A$-module projectif, $u:M\to N$ un homomorphisme. Pour que $u$
soit inversible à gauche, il faut et il suffit qu'il existe un système de
générateurs $(x_i)_{1\leq i\leq m}$ de $M$ tel que les images par
$u\otimes 1:M\otimes_Ak\to N\otimes_Ak$ des $x_i\otimes1$ soient linéairement
indépendantes dans $N\otimes_Ak$~; les $x_i$ forment alors une base de $M$.
\end{corollary}

La condition est évidemment nécessaire, car si $u$ est inversible à gauche,
$M$ est un $A$-module projectif de type fini, donc libre. Inversement, si la
condition est satisfaite, il est clair que les $x_i\otimes1$ forment une base
de $M\otimes_Ak$ et que $u\otimes1$ est inversible à gauche~; il suffit donc
d'appliquer (19.1.10) à l'idéal maximal $\mathfrak J$ de $A$.

\begin{corollary}[19.1.12]
\label{0.19.1.12-fr}
Soient $A$ un anneau, $M$ un $A$-module de type fini, $N$ un $A$-module
projectif, $u:M\to N$ un homomorphisme. Pour tout idéal premier
$\mathfrak p\in\operatorname{Spec}(A)$, les conditions suivantes sont
équivalentes~:
\begin{enumerate}
  \item[a)] L'homomorphisme
  $u_{\mathfrak p}:M_{\mathfrak p}\to N_{\mathfrak p}$ est inversible à
  gauche.
  \item[b)] L'homomorphisme
  $u\otimes1:M\otimes_Ak(\mathfrak p)\to N\otimes_Ak(\mathfrak p)$ est
  injectif (on rappelle que $k(\mathfrak p)$ désigne le corps résiduel de $A$
  en $\mathfrak p$).
  \item[c)] Il existe un système fini d'éléments $x_i\in M$
  $(1\leq i\leq m)$ tel que les images des $x_i$ dans $M_{\mathfrak p}$
  engendrent $M_{\mathfrak p}$, et un système de $m$ formes linéaires
  $y_i^*$ sur $N$ $(1\leq i\leq m)$ tel que
  $\det(\langle y_i^*,u(x_j)\rangle)\notin\mathfrak p$.
  \item[d)] Il existe $f\in A-\mathfrak p$ tel que l'homomorphisme
  $u_f:M_f\to N_f$ soit inversible à gauche.
\end{enumerate}
En outre, l'ensemble des $\mathfrak p\in\operatorname{Spec}(A)$ vérifiant ces
conditions est ouvert dans $\operatorname{Spec}(A)$.
\end{corollary}

La dernière condition est conséquence triviale de d). Comme $N$ est
projectif, il est facteur direct d'un $A$-module libre $A^{(I)}$~; en outre,
comme $M$ est de type fini, $u(M)$ est contenu dans un sous-module de $A^{(I)}$
de la forme $A^n$ pour $n$ fini~; comme chacun des énoncés a), b), c), d) est
équivalent à l'énoncé correspondant où l'on remplace $N$ par $N\oplus P$ ($u$
étant toujours supposé appliquer $M$ dans $N$), on est ramené au cas où $N$
est \emph{libre de type fini}. Il est trivial que d) entraîne a), et a) et b)
sont équivalentes en vertu de (19.1.11)~; d'ailleurs a) entraîne que
$M_{\mathfrak p}$ est libre (19.1.11), donc a) entraîne c), en prenant pour les
$x_i$ une famille dont les images dans $M_{\mathfrak p}$ forment une base de ce

\oldpage[0\textsubscript{IV-1}]{77}

$A_{\mathfrak p}$-module et en notant que (puisque $N$ est libre de type fini)
toute forme linéaire sur le $A_{\mathfrak p}$-module $N_{\mathfrak p}$ s'écrit
$u=v/f$, où $f\in A-\mathfrak p$, et où $v$ est une forme linéaire sur le
$A$-module $N$. Il est clair que c) entraîne b), et il reste donc à voir que a)
entraîne d). Or, comme $N$ est de présentation finie,
$(\operatorname{Hom}_A(N,M))_{\mathfrak p}$ s'identifie canoniquement à
$\operatorname{Hom}_{A_{\mathfrak p}}(N_{\mathfrak p},M_{\mathfrak p})$
(Bourbaki, \emph{Alg. comm.}, chap.~II, §~2, n\textsuperscript{o}~7,
prop.~19). Si $w_{\mathfrak p}$ est inverse à gauche de $u_{\mathfrak p}$, il
existe donc un homomorphisme $w:N\to M$ et un élément $f\in A-\mathfrak p$
tels que $w=w_{\mathfrak p}\otimes(1/f)1_{A_{\mathfrak p}}$. La relation
$w_{\mathfrak p}\circ u_{\mathfrak p}=1_{M_{\mathfrak p}}$ s'écrit donc aussi
$(w\circ u)\otimes1_{A_{\mathfrak p}}=f.1_{M_{\mathfrak p}}$. Mais comme $M$
est un $A$-module de type fini, il existe $g\in A-\mathfrak p$ tel que
l'endomorphisme $g((w\circ u)-f.1_M)$ s'annule en tous les générateurs de $M$,
donc soit nul. Si on pose $h=fg$, et
$u_h=u\otimes1_{A_h}$, $w_h=w\otimes1_{A_h}$, on a donc
$w_h(u_h(z))=(f/1)z$ pour tout $z\in M_h$~; mais comme $h/1$ est inversible
dans $A_h$, il en est de même de $f/1=f'$, et $f'^{-1}w_h$ est par suite
inverse à gauche de $u_h$, ce qui achève la démonstration.

\begin{proposition}[19.1.13]
\label{0.19.1.13-fr}
Supposons vérifiées les hypothèses de (19.1.9) et supposons en outre que
$(\mathfrak J_\lambda)$ soit une suite décroissante $(\mathfrak J_n)$ et que
$M$ soit séparé et complet. Alors les conditions a) et b) de (19.1.9) sont
aussi équivalentes à~:
\begin{enumerate}
  \item[c)] $u$ est inversible à gauche.
\end{enumerate}
\end{proposition}

On sait déjà que c) entraîne a) (19.1.6). Inversement, si a) est vérifiée, on
sait (avec les notations de (19.1.8)) que $M_n$ est facteur direct du
$(A/\mathfrak J_n)$-module projectif $N_n$, donc $P_n$ est aussi isomorphe à
un facteur direct de $N_n$, et par suite est projectif~; il suffit donc
d'appliquer (19.1.8).

Notons enfin la proposition suivante~:

\begin{proposition}[19.1.14]
\label{0.19.1.14-fr}
Soient $A$ un anneau, $M$ un $A$-module de type fini, $N$ un $A$-module
projectif, $u:M\to N$ un $A$-homomorphisme.
\begin{enumerate}
  \item[(i)] Pour que $u$ soit inversible à gauche, il faut et il suffit que,
  pour tout idéal maximal $\mathfrak m$ de $A$,
  $u_{\mathfrak m}:M_{\mathfrak m}\to N_{\mathfrak m}$ soit inversible à
  gauche.
  \item[(ii)] Soit $A'$ une $A$-algèbre qui soit un $A$-module fidèlement
  plat. Pour que $u$ soit inversible à gauche, il faut et il suffit que
  $u\otimes1:M\otimes_AA'\to N\otimes_AA'$ soit inversible à gauche.
\end{enumerate}
\end{proposition}

Comme dans (19.1.12), on peut se borner au cas où $N$ est libre de type fini~;
dire que $u$ est inversible à gauche signifie alors que $u$ est injectif et
que le module quotient $P=N/u(M)$ est \emph{projectif}, car $u(M)$ sera alors
facteur direct de $N$. Notons en outre que puisque $M$ est de type fini, $P$
est de \emph{présentation finie}. Cela étant~:

(i) La condition est évidemment nécessaire. Inversement, si elle est
satisfaite, on sait que $u$ est injectif (Bourbaki, \emph{Alg. comm.},
chap.~II, §~3, n\textsuperscript{o}~3, th.~1) et comme
$P_{\mathfrak m}=N_{\mathfrak m}/u_{\mathfrak m}(M_{\mathfrak m})$ est
projectif pour tout $\mathfrak m$, on sait que cela entraîne que $P$ est
projectif (\emph{loc. cit.}, §~5, n\textsuperscript{o}~2, th.~1).

(ii) Ici encore, la condition est trivialement nécessaire. Inversement, si
elle est remplie, on sait que $u$ est injectif ($\mathbf{0}_{\mathrm I}$,
6.4.1) et comme $P\otimes_AA'=\operatorname{Coker}(u\otimes1)$ est projectif,
donc plat, on en déduit que $P$ est un $A$-module plat
($\mathbf{0}_{\mathrm I}$, 6.6.3), donc projectif puisqu'il est de présentation
finie (Bourbaki, \emph{Alg. comm.}, chap.~II, §~5,
n\textsuperscript{o}~2, cor.~2 du th.~1).

\oldpage[0\textsubscript{IV-1}]{78}

\begin{remark}[19.1.15]
\label{0.19.1.15-fr}
La notion «~duale~» de celle d'homomorphisme formellement inversible à gauche
est celle d'homomorphisme \emph{formellement inversible à droite}~; un tel
$A$-homomorphisme continu $u:M\to N$ vérifie, par définition, la condition
suivante~: pour tout sous-module ouvert $W'$ de $N$, il existe un sous-module
ouvert $V'\subset u^{-1}(W')$ de $M$, un sous-module ouvert $W''\subset W'$ de
$N$ et un homomorphisme $h:N/W''\to M/V'$ tels que l'homomorphisme canonique
$N/W''\to N/W'$ se factorise en
\[
  N/W''\xrightarrow{h}M/V'\xrightarrow{u'}N/W',
\]
où $u'$ est déduit de $u$ par passage aux quotients. Cela entraîne que $u$ est
un épimorphisme formel~; s'il existe un $\widehat A$-homomorphisme continu
$r:\widehat N\to\widehat M$ tel que $\widehat u\circ r=1_{\widehat N}$, on
vérifie aussitôt que $u$ est formellement inversible à droite. Si les
conditions de (19.1.7) sont remplies, pour que $u$ soit formellement inversible
à droite, il faut et il suffit que les $u_\lambda$ soient inversibles à droite
(c'est-à-dire que le noyau de $u_\lambda$ est facteur direct de
$M/\mathfrak J_\lambda M$ et que $u_\lambda$ est un isomorphisme d'un
supplémentaire de $\operatorname{Ker}(u_\lambda)$ sur
$N/\mathfrak J_\lambda N$). Nous laissons au lecteur le soin d'énoncer et de
démontrer les propositions correspondant à (19.1.8) et (19.1.9) (dans
l'analogue de (19.1.8), il faut supposer $M$ séparé et complet et que $M_n$ est
un $(A/\mathfrak J_n)$-module projectif~; dans l'analogue de (19.1.9), il faut
supprimer l'hypothèse sur les $N/\mathfrak J_\lambda N$, mais supposer par
contre que, pour tout $\lambda$, $M/\mathfrak J_\lambda M$ est un
$(A/\mathfrak J_\lambda)$-module projectif).
\end{remark}

\subsection{Modules formellement projectifs}
\label{subsection:0.19.2-fr}

\begin{definition}[19.2.1]
\label{0.19.2.1-fr}
Soient $A$ un anneau topologique. On dit qu'un $A$-module topologique $M$ est
\emph{formellement projectif} s'il vérifie la condition suivante~: pour tout
idéal ouvert $\mathfrak J$ de $A$, tout couple de $(A/\mathfrak J)$-modules
(discrets) $P$, $Q$, tout $A$-homomorphisme surjectif $u:P\to Q$ et tout
$A$-homomorphisme continu $v:M\to Q$, il existe un $A$-homomorphisme continu
$w:M\to P$ tel que $v=u\circ w$.
\end{definition}

\begin{env}[19.2.2]
\label{0.19.2.2-fr}
Pour vérifier la condition de (19.2.1), on peut évidemment se borner (en
remplaçant $Q$ par $v(M)$ et $P$ par $u^{-1}(v(M))$) au cas où $v$ est
lui-même surjectif~; alors $Q$ est isomorphe à $M/V$, où $V$ est un
sous-module ouvert de $M$ tel que $\mathfrak JM\subset V$~; la condition de
(19.2.1) équivaut alors à dire que pour tout $(A/\mathfrak J)$-module discret
$P$ et tout $A$-homomorphisme surjectif $u:P\to M/V$, il existe un sous-module
ouvert $V'\subset V$ dans $M$ et un $A$-homomorphisme $w:M/V'\to P$, tels que
l'homomorphisme canonique $M/V'\to M/V$ se factorise en
\[
  M/V'\xrightarrow{w}P\xrightarrow{u}M/V.
\]
On notera qu'il suffit de vérifier cette condition lorsque $\mathfrak J$
parcourt un système fondamental de voisinages de 0 dans $A$ formé d'idéaux.
\end{env}

\begin{env}[19.2.3]
\label{0.19.2.3-fr}
Supposons qu'il existe un système fondamental $(\mathfrak J_\lambda)$ d'idéaux
ouverts dans $A$ et un système fondamental $(V_\lambda)$ de sous-modules
ouverts de $M$, ayant \emph{même} ensemble d'indices que
$(\mathfrak J_\lambda)$ et tel que, pour tout $\lambda$, $M/V_\lambda$ soit un
$(A/\mathfrak J_\lambda)$-module \emph{projectif}. Alors $M$ est formellement
projectif~: il suffit en effet, avec les notations de (19.2.2), de prendre
$\lambda$ tel que $\mathfrak J_\lambda\subset\mathfrak J$ et
$V_\lambda\subset V$~; comme $P$ est aussi un
$(A/\mathfrak J_\lambda)$-module, la factorisation de l'homomorphisme canonique
$M/V_\lambda\to M/V$ en $M/V_\lambda\to P\to M/V$ résulte de ce qu'il s'agit
d'un $(A/\mathfrak J_\lambda)$-homomorphisme et de ce que $M/V_\lambda$ est
supposé être un $(A/\mathfrak J_\lambda)$-module projectif.

Lorsque la condition plus stricte de ce numéro est satisfaite, on dit que $M$
est \emph{strictement formellement projectif}.
\end{env}

\begin{proposition}[19.2.4]
\label{0.19.2.4-fr}
Soient $A$ un anneau topologique, $M$ un $A$-module topologique dont la
topologie est déduite de celle de $A$. Pour que $M$ soit formellement
projectif, il faut et il suffit qu'il soit strictement formellement projectif.
\end{proposition}

\oldpage[0\textsubscript{IV-1}]{79}

\begin{proof}
On vient de voir que la condition est suffisante. Inversement, supposons $M$
formellement projectif et soit $(\mathfrak J_\lambda)$ un système fondamental
d'idéaux ouverts dans $A$. Pour tout $\lambda$, soient $P$ un
$(A/\mathfrak J_\lambda)$-module libre et
$u:P\to M/\mathfrak J_\lambda M$ un
$(A/\mathfrak J_\lambda)$-homomorphisme surjectif. Il existe donc un
$\mathfrak J_\mu\subset\mathfrak J_\lambda$ tel que l'homomorphisme canonique
$M/\mathfrak J_\mu M\to M/\mathfrak J_\lambda M$ se factorise en
\[
  M/\mathfrak J_\mu M\xrightarrow{w}P\xrightarrow{u}M/\mathfrak J_\lambda M;
\]
mais comme
$w(\mathfrak J_\lambda(M/\mathfrak J_\mu M))\subset\mathfrak J_\lambda P=0$,
$w$ se factorise en $M/\mathfrak J_\lambda M\xrightarrow{v}P$, et il est
clair que $u\circ v$ est l'identité dans $M/\mathfrak J_\lambda M$, ce qui
prouve que ce $(A/\mathfrak J_\lambda)$-module est projectif.
\end{proof}

\begin{proposition}[19.2.5]
\label{0.19.2.5-fr}
Soient $A$ un anneau topologique, $M$ un $A$-module topologique.
\begin{enumerate}
  \item[(i)] Pour que $M$ soit formellement projectif (resp. strictement
  formellement projectif), il faut et il suffit que le $\hat A$-module
  topologique $\hat M$ le soit.
  \item[(ii)] Soit $A'$ une $A$-algèbre topologique. Si $M$ est formellement
  projectif (resp. strictement formellement projectif), alors
  $M\otimes_A A'$ (muni de la topologie produit tensoriel) est un $A'$-module
  topologique formellement projectif (resp. strictement formellement
  projectif).
\end{enumerate}
\end{proposition}

\begin{proof}
\begin{enumerate}
  \item[(i)] Il suffit de remarquer que lorsque $\mathfrak J$ (resp. $V$)
  parcourt l'ensemble des idéaux ouverts de $A$ (resp. l'ensemble des
  sous-modules ouverts de $M$), le séparé complété $\hat{\mathfrak J}$
  (resp. $\hat V$) parcourt l'ensemble des idéaux ouverts de $\hat A$
  (resp. l'ensemble des sous-modules ouverts de $\hat M$), et
  $\hat A/\hat{\mathfrak J}=A/\mathfrak J$ (resp.
  $\hat M/\hat V=M/V$) à un isomorphisme canonique près~; comme la notion
  de module formellement projectif (resp. strictement formellement projectif)
  ne fait intervenir que les $A/\mathfrak J$ et les $M/V$, on en déduit
  aussitôt (i).
  \item[(ii)] Supposons d'abord $M$ formellement projectif et posons
  $M'=M\otimes_A A'$~; soient $\mathfrak J'$ un idéal ouvert de $A'$, $P'$,
  $Q'$ deux $(A'/\mathfrak J')$-modules discrets, $u':P'\to Q'$ un
  $A'$-homomorphisme surjectif, $v':M'\to Q'$ un $A'$-homomorphisme continu.
  Il y a un idéal ouvert $\mathfrak J$ de $A$ tel que
  $\mathfrak JA'\subset\mathfrak J'$, donc $P'$ et $Q'$ sont aussi des
  $(A/\mathfrak J)$-modules discrets. Si l'on considère le
  $A$-homomorphisme composé $v:M\xrightarrow{j}M'\xrightarrow{v'}Q'$ qui est
  continu, l'hypothèse entraîne qu'il existe un $A$-homomorphisme continu
  $w:M\to P'$ tel que $v=u'\circ w$~; mais comme $P'$ est un $A'$-module
  topologique, $w$ se factorise en
  $M\xrightarrow{j}M'\xrightarrow{w'}P'$, où $w'$ est un $A'$-homomorphisme
  continu, et comme $v'\circ j=(u'\circ w')\circ j$, on en conclut que
  $v'=u'\circ w'$.

  Supposons ensuite que $M$ soit strictement formellement projectif~; soit
  $(\mathfrak J_\lambda)$ (resp. $(V_\lambda)$) un système fondamental
  d'idéaux ouverts (resp. de sous-modules ouverts) dans $A$ (resp. $M$) tel
  que $(M/V_\lambda)$ soit un $(A/\mathfrak J_\lambda)$-module projectif,
  et soit $(\mathfrak J'_\mu)$ un système fondamental d'idéaux ouverts dans
  $A'$. On a un système fondamental de voisinages de $0$ dans $M'$ en prenant
  les sous-modules
  \[
    \operatorname{Im}(M\otimes_A\mathfrak J'_\mu)
    +\operatorname{Im}(V_\lambda\otimes_A A')=W_{\lambda\mu},
  \]
  où $\lambda$ et $\mu$ sont tels que
  $\mathfrak J_\lambda A'\subset\mathfrak J'_\mu$. Comme
  \[
    (M\otimes_A A')/W_{\lambda\mu}
    =(M/V_\lambda)\otimes_{A/\mathfrak J_\lambda}
      (A'/\mathfrak J'_\mu)
  \]
  et que $(M/V_\lambda)$ est un
  $(A/\mathfrak J_\lambda)$-module projectif, $M'/W_{\lambda\mu}$
  est un $(A'/\mathfrak J'_\mu)$-module projectif.
\end{enumerate}
\end{proof}

\subsection{Algèbres formellement lisses}
\label{subsection:0.19.3-fr}

\begin{definition}[19.3.1]
\label{0.19.3.1-fr}
Soient $A$ un anneau topologique, $B$ une $A$-algèbre topologique. On dit que
$B$ est une $A$-algèbre \emph{formellement lisse} si, pour toute
$A$-algèbre topologique discrète $C$ et tout idéal nilpotent $\mathfrak J$ de
$C$, tout $A$-homomorphisme continu $u:B\to C/\mathfrak J$ se factorise en
\[
  B\xrightarrow{v}C\xrightarrow{\varphi}C/\mathfrak J,
\]
où $v$ est un homomorphisme continu et $\varphi$ l'homomorphisme canonique.
\end{definition}

\oldpage[0\textsubscript{IV-1}]{80}

La définition (19.3.1) revient à dire que la propriété suivante a lieu~:
\begin{enumerate}
  \item[(P)] Pour tout idéal ouvert $\mathfrak R$ de la $A$-algèbre $B$ et
  tout $A$-homomorphisme $u':B/\mathfrak R\to C/\mathfrak J$, il y a un
  idéal ouvert $\mathfrak R'\subset\mathfrak R$ de $B$ tel que
  l'homomorphisme
  $B\to B/\mathfrak R\xrightarrow{u'}C/\mathfrak J$ se factorise en
  \[
    B\to B/\mathfrak R'\xrightarrow{v'}C\to C/\mathfrak J,
  \]
  où $v'$ est un $A$-homomorphisme.
\end{enumerate}
En effet, si $u:B\to C/\mathfrak J$ est un $A$-homomorphisme continu, il a
pour noyau $\mathfrak R$ un idéal ouvert de $B$, donc $u$ se factorise en
$B\to B/\mathfrak R\xrightarrow{u'}C/\mathfrak J$, et si (P) est
vérifiée, il suffit de l'appliquer à $u'$ et de prendre pour $v$ la composée
$B\to B/\mathfrak R'\xrightarrow{v'}C$ pour satisfaire aux conditions de
(19.3.1). Inversement, supposons que $B$ soit une $A$-algèbre formellement
lisse~; donnons-nous un idéal ouvert $\mathfrak R$ de $B$ et un
$A$-homomorphisme $u':B/\mathfrak R\to C/\mathfrak J$, et appliquons la
définition (19.3.1) à
$u:B\to B/\mathfrak R\xrightarrow{u'}C/\mathfrak J$~; si
$v:B\to C$ est un $A$-homomorphisme continu tel que $u$ se factorise en
$B\xrightarrow{v}C\to C/\mathfrak J$, l'idéal
$\mathfrak R'=\operatorname{Ker}(v)\cap\mathfrak R$ de $B$ est ouvert et
contenu dans $\mathfrak R$~; par suite $v$ se factorise en
$B\to B/\mathfrak R'\xrightarrow{v'}C$, et $v'$ vérifie bien la condition
(P).

\begin{proposition}[19.3.2]
\label{0.19.3.2-fr}
Soient $A$ un anneau discret, $V$ un $A$-module projectif~; l'algèbre
symétrique $B=\mathbf S_A^\bullet(V)$, munie de la topologie discrète, est une
$A$-algèbre formellement lisse.
\end{proposition}

\begin{proof}
En effet, les notations $C$ et $\mathfrak J$ ayant le même sens que dans
(19.3.1), soit $u:B\to C/\mathfrak J$ un homomorphisme de $A$-algèbres, qui
par restriction à $V=\mathbf S_A^1(V)$, donne un homomorphisme de $A$-modules
$u_1:V\to C/\mathfrak J$~; comme $V$ est projectif, $u_1$ se factorise en
$V\xrightarrow{v_1}C\xrightarrow{\varphi}C/\mathfrak J$, et $v_1$ se
prolonge en un homomorphisme de $A$-algèbres
$v:\mathbf S_A^\bullet(V)\to C$, tel que le composé
$\varphi\circ v$ coïncide avec $u$ dans $V$, donc
$\varphi\circ v=u$, ce qui démontre la proposition.
\end{proof}

\begin{corollary}[19.3.3]
\label{0.19.3.3-fr}
Si $A$ est un anneau discret, toute algèbre de polynômes
$B=A[T_\alpha]_{\alpha\in I}$, munie de la topologie discrète, est une
$A$-algèbre formellement lisse.
\end{corollary}

\begin{proposition}[19.3.4]
\label{0.19.3.4-fr}
Soit $A$ un anneau topologique, et soit
$B=A[[T_\alpha]]_{\alpha\in I}$ un anneau de séries formelles
$\sum_{\lambda\in\mathbb{N}^{(I)}}c_\lambda T^\lambda$ (algèbre large sur
$A$ du monoïde $\mathbb{N}^{(I)}$, identifié en tant que $A$-module au produit
$A^{\mathbb{N}^{(I)}}$). Si l'on munit $B$ de la topologie produit, $B$ est une
$A$-algèbre formellement lisse.
\end{proposition}

Soit $(\mathfrak L_\mu)_{\mu\in M}$ un système fondamental d'idéaux ouverts
dans $A$. Pour toute partie finie $H$ de $I$, tout $\mu\in M$ et tout entier
$n$, désignons par $\mathfrak K_{H,\mu,n}$ le voisinage de $0$ dans $B$ formé
des $(c_\lambda)$ tels que pour tout
$\lambda=(\lambda_\alpha)_{\alpha\in I}$ tel que
$\lambda_\alpha=0$ pour $\alpha\notin H$ et
$\sum_{\alpha\in H}\lambda_\alpha\leq n$, on ait
$c_\lambda\in\mathfrak L_\mu$. On vérifie immédiatement que les
$\mathfrak K_{H,\mu,n}$ forment un système fondamental de voisinages de $0$
dans $B$ et sont des idéaux de $B$, donc la topologie produit est compatible
avec la structure de $A$-algèbre de $B$.

Notons d'abord, avec les mêmes notations, le
\begin{lemma}[19.3.4.1]
\label{0.19.3.4.1-fr}
Soit $E$ une $A$-algèbre discrète.
\begin{enumerate}
  \item[(i)] Si $f:B\to E$ est un $A$-homomorphisme continu, il existe une
  partie finie $H$ de $I$ telle que $f(T_\alpha)=0$ pour $\alpha\notin H$,
  et $f(T_\alpha)$ est nilpotent dans $E$ pour tout $\alpha\in H$.
  \item[(ii)] Inversement, soient $H$ une partie finie de $I$,
  $(z_\alpha)_{\alpha\in H}$ une famille d'éléments nilpotents de $E$.
  Il existe un $A$-homomorphisme continu $g:B\to E$ et un seul tel que
  $g(T_\alpha)=z_\alpha$ pour $\alpha\in H$ et
  $g(T_\alpha)=0$ pour $\alpha\notin H$.
\end{enumerate}
\end{lemma}

\oldpage[0\textsubscript{IV-1}]{81}

\begin{proof}
(i) résulte aussitôt de ce que $f^{-1}(0)$ est un voisinage de $0$ dans le
$A$-module produit $A^{\mathbb{N}^{(I)}}$, d'où $f(T^\lambda)=0$ sauf pour un
nombre fini de valeurs de $\lambda\in\mathbb{N}^{(I)}$. Pour prouver (ii), il
suffit de remarquer que l'anneau de polynômes
$B'=A[T_\alpha]_{\alpha\in I}$ est dense dans $B$~; l'existence et
l'unicité de la restriction $g|B'$ sont triviales et sa continuité résulte de
l'hypothèse que les $f(T_\alpha)$ pour $\alpha\in H$ sont nilpotents,
car si $(f(T_\alpha))^n=0$ pour tout $\alpha\in H$, on a
$g(T^\lambda)=0$ pour tout $\lambda=(\lambda_\alpha)_{\alpha\in I}$ de
support fini sauf ceux tels que $\lambda_\alpha=0$ pour
$\alpha\notin H$ et $\lambda_\alpha\leq n$ pour
$\alpha\in H$, c'est-à-dire sauf pour un nombre fini de valeurs de
$\lambda$.
\end{proof}

\begin{proof}[Démonstration de la proposition~{\rm(19.3.4)}]
Ce lemme étant établi, et les notations $C$ et $\mathfrak J$ ayant le même
sens que dans (19.3.1), on a $u(T_\alpha)=0$ sauf pour
$\alpha\in H$, $H$ étant une partie finie de $I$, et les
$z_\alpha=u(T_\alpha)$ pour $\alpha\in H$ sont nilpotents dans
$C/\mathfrak J$~; comme $\mathfrak J$ est nilpotent, il existe une famille
$(x_\alpha)_{\alpha\in H}$ d'éléments nilpotents de $C$ dont les images
canoniques dans $C/\mathfrak J$ sont les $z_\alpha$~; si $v$ est le
$A$-homomorphisme continu de $B$ dans $C$ tel que
$v(T_\alpha)=0$ pour $\alpha\notin H$, $v(T_\alpha)=x_\alpha$ pour
$\alpha\in H$, il est clair que $u$ se factorise en
$B\xrightarrow{v}C\to C/\mathfrak J$.
\end{proof}

\begin{proposition}[19.3.5]
\label{0.19.3.5-fr}
Soit $A$ un anneau topologique.
\begin{enumerate}
  \item[(i)] $A$ est une $A$-algèbre formellement lisse.
  \item[(ii)] Si $B$ est une $A$-algèbre formellement lisse et $C$ une
  $B$-algèbre formellement lisse, alors $C$ est une $A$-algèbre formellement
  lisse.
  \item[(iii)] Soient $B$ une $A$-algèbre formellement lisse, $A'$ une
  $A$-algèbre topologique~; alors la $A'$-algèbre topologique
  $B\otimes_A A'$ ($\mathbf 0_{\rm I}$, 7.7.5 et 7.7.6) est formellement
  lisse.
  \item[(iv)] Soient $B$ une $A$-algèbre topologique, $S$ (resp. $T$) une
  partie multiplicative de $A$ (resp. $B$) telle que l'image canonique de $S$
  dans $B$ soit contenue dans $T$. Si $B$ est une $A$-algèbre formellement
  lisse, alors $T^{-1}B$ est une $S^{-1}A$-algèbre formellement lisse.
  \item[(v)] Soient $B_i$ ($1\leq i\leq n$) des $A$-algèbres topologiques.
  Pour que $\prod_{i=1}^n B_i$ soit une $A$-algèbre formellement lisse, il
  faut et il suffit que chacune des $B_i$ le soit.
\end{enumerate}
\end{proposition}

\oldpage[0\textsubscript{IV-1}]{82}

\begin{proof}
\begin{enumerate}
  \item[(i)] Si $C$ est une $A$-algèbre discrète,
  $\varphi:C\to C/\mathfrak J$ l'homomorphisme canonique de $C$ sur une
  $A$-algèbre quotient quelconque de $C$, le \emph{seul} $A$-homomorphisme
  de $A$ dans $C/\mathfrak J$ est l'homomorphisme composé
  $A\xrightarrow{\psi}C\xrightarrow{\varphi}C/\mathfrak J$, où
  $\psi$ est l'homomorphisme définissant la structure de $A$-algèbre de
  $C$~; comme $\psi$ est continu, la condition de (19.3.1) est trivialement
  vérifiée.
  \item[(ii)] Soient $\alpha:A\to B$, $\beta:B\to C$ les homomorphismes
  continus définissant respectivement la structure de $A$-algèbre sur $B$ et
  celle de $B$-algèbre sur $C$, de sorte que $\beta\circ\alpha$ définit la
  structure de $A$-algèbre sur $C$. Soient $E$ une $A$-algèbre discrète,
  $\mathfrak L$ un idéal nilpotent de $E$, $u:C\to E/\mathfrak L$ un
  $A$-homomorphisme continu, de sorte que $u\circ\beta\circ\alpha$ est
  l'homomorphisme définissant la structure de $A$-algèbre de
  $E/\mathfrak L$. Comme $B$ est une $A$-algèbre formellement lisse, le
  $A$-homomorphisme continu $u\circ\beta:B\to E/\mathfrak L$ se
  factorise en $B\xrightarrow{v}E\to E/\mathfrak L$, où $v$ est un
  $A$-homomorphisme continu~; $v$ et $u\circ\beta$ définissent alors sur
  $E$ et $E/\mathfrak L$ respectivement des structures de $B$-algèbre
  topologique, pour lesquelles $E/\mathfrak L$ est encore l'algèbre quotient
  (discrète) de la $B$-algèbre $E$. En outre, $u$ est un $B$-homomorphisme
  continu, donc se factorise en $C\xrightarrow{w}E\to E/\mathfrak L$,
  où $w$ est un $B$-homomorphisme continu~; comme
  $v\circ\alpha$ est le $A$-homomorphisme définissant la structure de
  $A$-algèbre sur $E$, $w$ est bien un $A$-homomorphisme continu, d'où notre
  assertion.
  \item[(iii)] Soient $C$ une $A'$-algèbre topologique discrète,
  $\mathfrak J$ un idéal nilpotent de $C$,
  $u:B\otimes_A A'\to C/\mathfrak J$ un $A'$-homomorphisme continu. Si on
  compose $u$ et l'homomorphisme canonique
  $\rho:B\to B\otimes_A A'$, on obtient (puisque
  $A\to B\xrightarrow{\rho}B\otimes_A A'$ est égal au composé
  $A\to A'\xrightarrow{\sigma}B\otimes_A A'$) un $A$-homomorphisme
  continu qui, par hypothèse, se factorise donc en
  $B\xrightarrow{v}C\to C/\mathfrak J$, où $v$ est un $A$-homomorphisme
  continu. L'égalité des composés $A\to B\xrightarrow{v}C$ et
  $A\to A'\xrightarrow{w}C$ entraîne donc l'existence d'un homomorphisme
  continu d'anneaux $f:B\otimes_A A'\to C$ tel que
  $v=f\circ\rho$ et $w=f\circ\sigma$ ($\mathbf 0_{\rm I}$, 7.7.6)~;
  comme les homomorphismes composés
  $B\xrightarrow{\rho}B\otimes_A A'\xrightarrow{f}C\to C/\mathfrak J$
  et
  $B\xrightarrow{\rho}B\otimes_A A'\xrightarrow{u}C/\mathfrak J$
  (resp.
  $A'\xrightarrow{\sigma}B\otimes_A A'\xrightarrow{f}C\to C/\mathfrak J$
  et
  $A'\xrightarrow{\sigma}B\otimes_A A'\xrightarrow{u}C/\mathfrak J$)
  sont égaux, on a bien la factorisation
  $u:B\otimes_A A'\xrightarrow{f}C\to C/\mathfrak J$, ce qui établit
  (iii).
  \item[(iv)] La topologie considérée sur $S^{-1}A$ (resp. $T^{-1}B$) est
  naturellement celle pour laquelle un système fondamental de voisinages de
  $0$ est formé des $S^{-1}\mathfrak J$ (resp. $T^{-1}\mathfrak R$), où
  $\mathfrak J$ (resp. $\mathfrak R$) parcourt un système fondamental
  d'idéaux ouverts dans $A$ (resp. $B$) ($\mathbf 0_{\rm I}$, 7.6.1). Si
  $\alpha:A\to B$ est l'homomorphisme canonique, il est clair que
  l'homomorphisme canonique $\alpha':S^{-1}A\to T^{-1}B$ déduit de
  $\alpha$ est continu. Soient alors $C$ une $S^{-1}A$-algèbre topologique
  discrète, $\mathfrak J$ un idéal nilpotent de cette algèbre,
  $u:T^{-1}B\to C/\mathfrak J$ un $S^{-1}A$-homomorphisme continu~; alors
  le composé $B\to T^{-1}B\xrightarrow{u}C/\mathfrak J$ est un
  $A$-homomorphisme continu qui, par hypothèse, se factorise en
  $B\xrightarrow{v}C\to C/\mathfrak J$, où $v$ est un $A$-homomorphisme
  continu. Comme pour tout $t\in T$, $t/1$ est inversible dans $T^{-1}B$,
  $u(t/1)$ est inversible dans $C/\mathfrak J$. Comme $\mathfrak J$ est
  nilpotent, tout élément de la classe $u(t/1)$ dans $C$, et en particulier
  $v(t)$, est inversible dans $C$ ($\mathbf 0_{\rm I}$, 7.1.12), et par suite
  $v$ se factorise en $B\to T^{-1}B\xrightarrow{w}C$~; comme $v$ est
  continu, il en est de même de $w$ ($\mathbf 0_{\rm I}$, 7.6.6), et c'est
  un $S^{-1}A$-homomorphisme car le composé
  $A\to B\to T^{-1}B\xrightarrow{w}C$ est égal à
  \[
    A\to S^{-1}A\xrightarrow{\alpha'}T^{-1}B\xrightarrow{w}C,
  \]
  donc $S^{-1}A\xrightarrow{\alpha'}T^{-1}B\xrightarrow{w}C$ est
  l'homomorphisme canonique définissant sur $C$ la structure de
  $S^{-1}A$-algèbre. Enfin, les homomorphismes composés
  $B\to T^{-1}B\xrightarrow{w}C\to C/\mathfrak J$ et
  $B\to T^{-1}B\xrightarrow{u}C/\mathfrak J$ étant égaux, le même
  raisonnement montre que $u$ est bien égal au composé
  $T^{-1}B\xrightarrow{w}C\to C/\mathfrak J$, d'où (iv).
  \item[(v)] Enfin, (v) est immédiat, la donnée d'un $A$-homomorphisme
  continu de $B=\prod_{i=1}^n B_i$ dans $C$ (resp. $C/\mathfrak J$)
  équivalant à celle de $n$ $A$-homomorphismes continus $B_i\to C$
  (resp. $B_i\to C/\mathfrak J$) et tout $A$-homomorphisme continu
  $B_j\to C$ (resp. $B_j\to C/\mathfrak J$) donnant par composition un
  $A$-homomorphisme continu $B\to B_j\to C$
  (resp. $B\to B_j\to C/\mathfrak J$).
\end{enumerate}
\end{proof}

\oldpage[0\textsubscript{IV-1}]{83}

\begin{proposition}[19.3.6]
\label{0.19.3.6-fr}
Soient $A$ un anneau topologique, $B$ une $A$-algèbre topologique, $\hat A$
et $\hat B$ les séparés complétés respectifs de $A$ et $B$, de sorte que
$\hat B$ est une $\hat A$-algèbre topologique. Les conditions suivantes
sont équivalentes~:
\begin{enumerate}
  \item[a)] $B$ est une $A$-algèbre formellement lisse.
  \item[b)] $\hat B$ est une $A$-algèbre formellement lisse.
  \item[c)] $\hat B$ est une $\hat A$-algèbre formellement lisse.
\end{enumerate}
\end{proposition}

\begin{proof}
Bien entendu, la structure de $\hat A$-algèbre sur $\hat B$ est définie
par l'homomorphisme $\hat\varphi$, si $\varphi:A\to B$ est
l'homomorphisme définissant la structure de $A$-algèbre sur $B$. Comme toute
$A$-algèbre discrète $C$ est séparée et complète, elle est aussi une
$\hat A$-algèbre (par prolongement canonique de l'homomorphisme de $A$ dans
$C$), et un $A$-homomorphisme continu de $B$ dans $C$ donne par prolongement
canonique un $\hat A$-homomorphisme (et \emph{a fortiori} un
$A$-homomorphisme) de $\hat B$ dans $C$ (autrement dit, tout
$A$-homomorphisme $B\to C$ se factorise en $B\to\hat B\to C$ de façon
unique). Ces remarques et la définition (19.3.1) entraînent aussitôt la
proposition.
\end{proof}

\begin{corollary}[19.3.7]
\label{0.19.3.7-fr}
Sous les hypothèses de (19.3.5), (iv), la $A\{S^{-1}\}$-algèbre topologique
$B\{T^{-1}\}$ est formellement lisse.
\end{corollary}

Cela résulte des définitions ($\mathbf 0_{\rm I}$, 7.6.1 et 7.6.7) et de
(19.3.5), (iv) et (19.3.6).

\begin{proposition}[19.3.8]
\label{0.19.3.8-fr}
Soient $A$ un anneau topologique, $B$ une $A$-algèbre topologique, et supposons
que pour tout idéal ouvert $\mathfrak R$ de $B$, $\mathfrak R^2$ soit aussi
ouvert. Soient $A'$, $B'$ des anneaux topologiques obtenus en munissant $A$ et
$B$ de topologies plus fines que les topologies initiales, et supposons que
l'homomorphisme canonique $\varphi:A\to B$ soit encore un homomorphisme
continu de $A'$ dans $B'$. Alors, si $B'$ est une $A'$-algèbre formellement
lisse, $B$ est une $A$-algèbre formellement lisse.
\end{proposition}

Il suffit d'appliquer le lemme suivant~:

\begin{lemma}[19.3.8.1]
\label{0.19.3.8.1-fr}
Soient $C$ une $A$-algèbre discrète, $\mathfrak J$ un idéal nilpotent de $C$.
Supposons que le carré de tout idéal ouvert de $B$ soit ouvert. Alors, si un
homomorphisme composé
$v:B\xrightarrow{u}C\to C/\mathfrak J$ est continu, l'homomorphisme $u$
est continu.
\end{lemma}

\begin{proof}
En effet, $\mathfrak R=u^{-1}(\mathfrak J)$ est le noyau de $v$ et il est
ouvert par hypothèse~; comme il existe un entier $n$ tel que
$\mathfrak J^n=0$, $\mathfrak R^n$ est contenu dans
$\operatorname{Ker}(u)$~; mais par récurrence sur $h$, il résulte de
l'hypothèse que $\mathfrak R^{2^h}$ est ouvert, donc aussi
$\mathfrak R^n$ pour tout $n$, et par suite $\operatorname{Ker}(u)$ est lui
aussi ouvert, ce qui prouve notre assertion.
\end{proof}

On notera que l'hypothèse sur $B$ signifie que la topologie de $B$ est borne
supérieure des topologies $\mathfrak R$-préadiques, où $\mathfrak R$ parcourt
l'ensemble des idéaux ouverts de $B$. Si $B$ est un anneau préadique
($\mathbf 0_{\rm I}$, 7.1.9), cette condition est toujours vérifiée.

\begin{env}[19.3.9]
\label{0.19.3.9-fr}
Nous allons maintenant voir que la propriété d'être formellement lisse
implique des propriétés de «~relèvement~» d'homomorphismes sous des conditions
plus générales que celles de la définition (19.3.1).
\end{env}

\oldpage[0\textsubscript{IV-1}]{84}

\begin{proposition}[19.3.10]
\label{0.19.3.10-fr}
Soient $A$ un anneau topologique, $B$ une $A$-algèbre formellement lisse.
Soient $C$ une $A$-algèbre topologique, $\mathfrak J$ un idéal de $C$,
vérifiant les conditions suivantes~:
\begin{enumerate}
  \item[$1^\circ$] $C$ est métrisable et complet.
  \item[$2^\circ$] $\mathfrak J$ est fermé et la suite
  $(\mathfrak J^n)$ tend vers $0$.
\end{enumerate}
Alors, tout $A$-homomorphisme continu $u:B\to C/\mathfrak J$ se factorise en
$B\xrightarrow{v}C\to C/\mathfrak J$, où $v$ est un
$A$-homomorphisme continu.
\end{proposition}

\begin{proof}
Soit $(\mathfrak Q_n)$ une suite décroissante d'idéaux de $C$ formant un
système fondamental de voisinages de $0$. Pour tout $n$, considérons la
$A$-algèbre discrète $C/\mathfrak Q_n$ et l'idéal
$(\mathfrak J+\mathfrak Q_n)/\mathfrak Q_n$ de cette algèbre~; comme il
existe $k$ tel que $\mathfrak J^k\subset\mathfrak Q_n$,
$(\mathfrak J+\mathfrak Q_n)/\mathfrak Q_n$ est nilpotent. Pour tout
$n$, soit $u_n$ l'homomorphisme continu
\[
  B\xrightarrow{u}C/\mathfrak J\to C/(\mathfrak J+\mathfrak Q_n)
  =(C/\mathfrak Q_n)/((\mathfrak J+\mathfrak Q_n)/\mathfrak Q_n)~;
\]
par hypothèse $u_n$ se factorise en
\[
  B\xrightarrow{v_n}C/\mathfrak Q_n
  \xrightarrow{\varphi_n}C/(\mathfrak J+\mathfrak Q_n),
\]
où $v_n$ est un $A$-homomorphisme continu. Montrons qu'on peut choisir de
proche en proche les $v_n$ de sorte que $v_n$ se factorise en
\[
  B\xrightarrow{v_{n+1}}C/\mathfrak Q_{n+1}\to C/\mathfrak Q_n
\]
pour tout $n$ (autrement dit, que les $v_n$ forment un \emph{système
projectif} d'homomorphismes). Cela résultera du lemme suivant.
\end{proof}

\begin{lemma}[19.3.10.1]
\label{0.19.3.10.1-fr}
Soient $B$ une $A$-algèbre formellement lisse, $E$, $E'$ deux $A$-algèbres
topologiques discrètes, $\mathfrak R$ (resp. $\mathfrak R'$) un idéal
nilpotent de $E$ (resp. $E'$), $f:E\to E'$ un $A$-homomorphisme surjectif tel
que $f(\mathfrak R)\subset\mathfrak R'$, $f':E/\mathfrak R\to
E'/\mathfrak R'$ l'homomorphisme déduit de $f$ par passage aux quotients.
Soient $u:B\to E/\mathfrak R$ un $A$-homomorphisme continu, $u'$
l'homomorphisme composé
$B\xrightarrow{u}E/\mathfrak R\xrightarrow{f'}E'/\mathfrak R'$, et soit
$v':B\to E'$ un $A$-homomorphisme continu tel que $u'$ se factorise en
$B\xrightarrow{v'}E'\to E'/\mathfrak R'$. Alors il existe un
$A$-homomorphisme continu $v:B\to E$ tel que $v'$ se factorise en
$B\xrightarrow{v}E\xrightarrow{f}E'$.
\end{lemma}

Dans le diagramme commutatif
\[
\xymatrix{
  E \ar[r]^f \ar[d] & E' \ar[d] \\
  E/\mathfrak R \ar[r]_{f'} & E'/\mathfrak R'
}
\]
tous les homomorphismes sont surjectifs~; si $\mathfrak L$ est le noyau de
$f$, $E'$ s'identifie à $E/\mathfrak L$ et $\mathfrak R'$ à
$(\mathfrak R+\mathfrak L)/\mathfrak L$. Utilisons maintenant le lemme
élémentaire suivant.

\begin{lemma}[19.3.10.2]
\label{0.19.3.10.2-fr}
Soient $F$ un anneau (non nécessairement commutatif), $\mathfrak a$,
$\mathfrak b$ deux idéaux bilatères de $F$~; alors le produit fibré
\[
  (F/\mathfrak a)\times_{F/(\mathfrak a+\mathfrak b)}(F/\mathfrak b)
\]
s'identifie canoniquement à $F/(\mathfrak a\cap\mathfrak b)$.
\end{lemma}

\begin{proof}
Ce n'est autre qu'une application particulière de (18.1.7), où, dans le
diagramme (18.1.7.1), on remplace $G$ par
$F/(\mathfrak a\cap\mathfrak b)$, $\mathfrak J'$ par
$\mathfrak b/(\mathfrak a\cap\mathfrak b)$, $F$ par
$F/\mathfrak b$, $\mathfrak R$ par
$(\mathfrak a+\mathfrak b)/\mathfrak b$, $B$ par
$F/(\mathfrak a+\mathfrak b)$ et $\theta$ par le $F$-homomorphisme
bijectif canonique
$(\mathfrak a+\mathfrak b)/\mathfrak b\to
\mathfrak a/(\mathfrak a\cap\mathfrak b)$.
\end{proof}

\oldpage[0\textsubscript{IV-1}]{85}

\begin{proof}[Suite de la démonstration du lemme~{\rm(19.3.10.1)}]
Appliquant ce lemme à la situation de (19.3.10.1), la commutativité du
diagramme
\[
\xymatrix{
  B \ar[r]^{v'} \ar[d]_u & E'=E/\mathfrak L \ar[d] \\
  E/\mathfrak R \ar[r]_{f'} & E'/\mathfrak R'
}
\]
montre l'existence d'un homomorphisme
$w:B\to E/(\mathfrak R\cap\mathfrak L)$ tel que $v'$ et $u$ se
factorisent en
$B\xrightarrow{w}E/(\mathfrak R\cap\mathfrak L)\to E/\mathfrak L$
et
$B\xrightarrow{w}E/(\mathfrak R\cap\mathfrak L)\to E/\mathfrak R$
respectivement~; comme le noyau de $w$ contient l'intersection de ceux de
$v'$ et de $u$, il est ouvert dans $B$ et $w$ est continu. Enfin, il est clair
que $\mathfrak R\cap\mathfrak L$ est nilpotent, donc $w$ se factorise en
$B\xrightarrow{v}E\to E/(\mathfrak R\cap\mathfrak L)$, où $v$ est
un $A$-homomorphisme continu répondant à la question.
\end{proof}

\begin{proof}[Suite de la démonstration de la proposition~{\rm(19.3.10)}]
L'existence de $v_{n+1}$ résulte du lemme (19.3.10.1) appliqué en remplaçant
$E$, $E'$, $\mathfrak R$, $\mathfrak R'$ par
$C/\mathfrak Q_{n+1}$, $C/\mathfrak Q_n$,
$(\mathfrak J+\mathfrak Q_{n+1})/\mathfrak Q_{n+1}$ et
$(\mathfrak J+\mathfrak Q_n)/\mathfrak Q_n$ respectivement, $u$, $u'$
et $v'$ par $u_{n+1}$, $u_n$ et $v_n$ respectivement. Notons que, comme $C$
est séparé et complet, il est égal à $\varprojlim_n(C/\mathfrak Q_n)$, et
$v=\varprojlim_n v_n$ est un $A$-homomorphisme continu de $B$ dans $C$.
Comme $\mathfrak J$ est fermé dans $C$, on a
$\mathfrak J=\bigcap_n(\mathfrak J+\mathfrak Q_n)$~; comme
$C/\mathfrak J$ est métrisable et complet et que les
$(\mathfrak J+\mathfrak Q_n)/\mathfrak J$ forment un système fondamental
de voisinages de $0$ dans $C/\mathfrak J$, on a
\[
  C/\mathfrak J=\varprojlim_n C/(\mathfrak J+\mathfrak Q_n)
\]
et $\varprojlim_n\varphi_n$ est l'application canonique
$\varphi:C\to C/\mathfrak J$. Comme $\varprojlim_n u_n=u$, on a bien
$u=v\circ\varphi$, C.Q.F.D.
\end{proof}

\begin{corollary}[19.3.11]
\label{0.19.3.11-fr}
Soient $A$ un anneau topologique, $B$ une $A$-algèbre formellement lisse,
$C$ un anneau local noethérien complet, $\mathfrak J$ un idéal distinct de
$C$, $\varphi:A\to C$ un homomorphisme continu, faisant de $C$ une
$A$-algèbre topologique. Alors tout $A$-homomorphisme continu
$u_0:B\to C_0=C/\mathfrak J$ se factorise en
$B\xrightarrow{u}C\to C/\mathfrak J$, où $u$ est un $A$-homomorphisme
continu.
\end{corollary}

Toutes les conditions de (19.3.10) sont en effet remplies
($\mathbf 0_{\rm I}$, 7.3.5).

\begin{proposition}[19.3.12]
\label{0.19.3.12-fr}
Soient $A$ un anneau topologique, $B$ une $A$-algèbre formellement lisse,
$C$ une $A$-algèbre topologique, $\mathfrak J$ un idéal de $C$, vérifiant les
conditions suivantes~:
\begin{enumerate}
  \item[$1^\circ$] Il existe un système fondamental d'idéaux ouverts
  $\mathfrak Q_\lambda$ de $C$, tels que les
  $C_\lambda=C/\mathfrak Q_\lambda$ soient des anneaux artiniens et que
  l'homomorphisme canonique $C\to\varprojlim_\lambda C_\lambda$ soit un
  isomorphisme d'anneaux topologiques.
  \item[$2^\circ$] L'idéal $\mathfrak J$ est fermé dans $C$ et
  topologiquement nilpotent.
  \item[$3^\circ$] Le carré de tout idéal ouvert de $B$ est ouvert.
\end{enumerate}
Dans ces conditions, tout $A$-homomorphisme continu
$u:B\to C/\mathfrak J$ se factorise en
$B\xrightarrow{v}C\to C/\mathfrak J$, où $v$ est un $A$-homomorphisme
continu.
\end{proposition}

\begin{proof}
Soit $\mathfrak J_\lambda=(\mathfrak J+\mathfrak Q_\lambda)/\mathfrak Q_\lambda$
l'image canonique de $\mathfrak J$ dans $C_\lambda$, qui est un idéal
nilpotent, et soit $u_\lambda$ l'homomorphisme composé
$B\xrightarrow{u}C/\mathfrak J\to C_\lambda/\mathfrak J_\lambda
=C/(\mathfrak J+\mathfrak Q_\lambda)$. Montrons que tout
$A$-homomorphisme $w:B\to C$ tel que $u$ se factorise en
$B\xrightarrow{w}C\to C/\mathfrak J$ est nécessairement continu~; en
effet, pour tout $\lambda$, il y a un idéal ouvert $\mathfrak R$ de $B$
tel que $u(\mathfrak R)\subset
(\mathfrak J+\mathfrak Q_\lambda)/\mathfrak J$, d'où
$w(\mathfrak R)\subset\mathfrak J+\mathfrak Q_\lambda$, donc il y a
une puissance $\mathfrak R^{2^m}$ de $\mathfrak R$ telle que
$w(\mathfrak R^{2^m})\subset\mathfrak Q_\lambda$, ce qui établit
notre assertion puisque $\mathfrak R^{2^m}$ est ouvert dans $B$. On peut
donc se borner à trouver un $A$-homomorphisme $w$ ayant la propriété précédente
sans s'inquiéter de ses propriétés de continuité. Or, l'ensemble
$\operatorname{Hom}(B,C)$ de tous les $A$-homomorphismes d'algèbres de $B$
dans $C$ est égal à
$\varprojlim_\lambda\operatorname{Hom}(B,C_\lambda)$, et ce dernier est
\emph{fermé} dans le $C_\lambda$-module $C_\lambda^B$, muni de la topologie
produit, pour laquelle il est \emph{linéairement compact} puisque
  $C_\lambda$ est artinien. Pour tout $\lambda$, soit $W_\lambda$

\oldpage[0\textsubscript{IV-1}]{86}

l'ensemble des $w_\lambda\in\operatorname{Hom}(B,C_\lambda)$ tels que
$u_\lambda$ se factorise en
$B\xrightarrow{w_\lambda}C_\lambda\to
C_\lambda/\mathfrak J_\lambda$~; $W_\lambda$ est une variété
linéaire fermée dans le $C_\lambda$-module
  $\operatorname{Hom}(B,C_\lambda)$, non vide puisque $B$ est formellement
  lisse. Dans le produit
$\prod_{\mu\leq\lambda}\operatorname{Hom}(B,C_\mu)=E_\lambda$,
considérons la sous-variété linéaire $U_\lambda$ formée des
$(w_\mu)_{\mu\leq\lambda}$ tels que
$w_\mu=\varphi_{\mu\lambda}\circ w_\lambda$ pour
$\mu\leq\lambda$ et $w_\lambda\in W_\lambda$, où
$\varphi_{\mu\lambda}:C_\lambda\to C_\mu$ est
l'homomorphisme canonique, qui est aussi fermée. Enfin, soit $V_\lambda$
la variété linéaire dans le produit
$\prod_\mu\operatorname{Hom}(B,C_\mu)$, produit de $U_\lambda$
et des $\operatorname{Hom}(B,C_\mu)$ pour
$\mu\not\leq\lambda$, qui est encore fermée. Tout revient à voir que
l'intersection des $V_\lambda$ est \emph{non vide}, car un élément $w$ de
cette intersection appartiendra à
$\varprojlim_\lambda\operatorname{Hom}(B,C_\lambda)$ par définition.
D'ailleurs, comme $\mathfrak J$ est fermé, et
$C=\varprojlim_\lambda C_\lambda$ linéairement compact,
$C/\mathfrak J$ s'identifie à
$\varprojlim_\lambda C_\lambda/\mathfrak J_\lambda$ et
l'application canonique $\psi:C\to C/\mathfrak J$ à
$\varprojlim_\lambda\psi_\lambda$, où $\psi_\lambda$ est
l'application canonique
$C_\lambda\to C_\lambda/\mathfrak J_\lambda$. On conclut donc
comme dans (19.3.10) que $\psi\circ w=u$.
\end{proof}

\subsection{Premiers critères de lissité formelle}
\label{subsection:0.19.4-fr}

\begin{proposition}[19.4.1]
\label{0.19.4.1-fr}
Soient $A$ un anneau topologique, $B$ une $A$-algèbre topologique~; supposons
qu'il existe deux familles filtrantes décroissantes
$(\mathfrak J_\alpha)_{\alpha\in I}$,
$(\mathfrak R_\alpha)_{\alpha\in I}$ d'idéaux de $A$ et $B$
respectivement, telles que~:
\begin{enumerate}
  \item[$1^\circ$] $(\mathfrak J_\alpha)$ tend vers $0$ dans $A$ et
  $(\mathfrak R_\alpha)$ tend vers $0$ dans $B$~;
  \item[$2^\circ$] pour tout $\alpha\in I$ on a
  $\mathfrak J_\alpha B\subset\mathfrak R_\alpha$ (de sorte que
  $B/\mathfrak R_\alpha$ est une
  $(A/\mathfrak J_\alpha)$-algèbre topologique)~;
  \item[$3^\circ$] pour tout $\alpha\in I$,
  $B/\mathfrak R_\alpha$ est une
  $(A/\mathfrak J_\alpha)$-algèbre formellement lisse.
\end{enumerate}
Alors $B$ est une $A$-algèbre formellement lisse.
\end{proposition}

\begin{proof}
En effet, soient $C$ une $A$-algèbre discrète, $\mathfrak L$ un idéal
nilpotent de $C$, $\mathfrak R$ un idéal ouvert de $B$~; par hypothèse il y a
un $\alpha\in I$ tel que
$\mathfrak R_\alpha\subset\mathfrak R$, donc $B/\mathfrak R$ est
quotient de $B/\mathfrak R_\alpha$ par un idéal ouvert. Tout
$A$-homomorphisme $u':B/\mathfrak R\to C/\mathfrak L$ est aussi un
$(A/\mathfrak J_\alpha)$-homomorphisme, donc il existe un idéal ouvert
$\mathfrak R'$ de $B$ tel que
$\mathfrak R_\alpha\subset\mathfrak R'\subset\mathfrak R$ et que
$B/\mathfrak R_\alpha\to B/\mathfrak R\xrightarrow{u'}C/\mathfrak L$
se factorise en
$B/\mathfrak R_\alpha\to B/\mathfrak R'\to C\to C/\mathfrak L$,
où le second homomorphisme est un
$(A/\mathfrak J_\alpha)$-homomorphisme~; d'où la conclusion, en vertu de
la remarque suivant (19.3.1).
\end{proof}

\begin{corollary}[19.4.2]
\label{0.19.4.2-fr}
Soient $A$ un anneau topologique, $B$ une $A$-algèbre topologique,
$(\mathfrak J_\alpha)_{\alpha\in I}$ une famille filtrante décroissante
d'idéaux de $A$ tendant vers $0$. Pour que $B$ soit une $A$-algèbre
formellement lisse, il faut et il suffit que pour tout $\alpha\in I$, si
l'on pose $A_\alpha=A/\mathfrak J_\alpha$,
$B_\alpha=B/\mathfrak J_\alpha B$, $B_\alpha$ soit une
$A_\alpha$-algèbre formellement lisse.
\end{corollary}

La condition est suffisante par (19.4.1), et elle est aussi nécessaire par
(19.3.5), (iii).

\begin{proposition}[19.4.3]
\label{0.19.4.3-fr}
Soient $A$ un anneau topologique, $B$ une $A$-algèbre topologique. Supposons
que pour toute $A$-algèbre discrète $C$ et tout idéal $\mathfrak J$ de $C$
tel que $\mathfrak J^2=0$, tout $A$-homomorphisme continu
$u:B\to C/\mathfrak J$ se factorise en
$B\xrightarrow{v}C\to C/\mathfrak J$, où $v$ est un $A$-homomorphisme
continu. Alors $B$ est une $A$-algèbre formellement lisse.
\end{proposition}

\begin{proof}
En effet, avec les mêmes notations, soit $\mathfrak R$ un idéal nilpotent
quelconque de $C$, et considérons un $A$-homomorphisme continu
$u':B\to C/\mathfrak R$. Supposons que $\mathfrak R^m=0$, et posons
$C_i=C/\mathfrak R^i$ pour $1\leq i\leq m$, de sorte que $C_m=C$, que
$\mathfrak R^{i-1}/\mathfrak R^i$ est un idéal de carré nul
$\mathfrak J_i$ dans $C_i$, et $C_i=C_{i+1}/\mathfrak J_{i+1}$~;
l'hypothèse entraîne alors par récurrence sur $i$ l'existence de
$A$-homomorphismes continus $u_i:B\to C_i$ tels que $u_1=u$ et que $u_i$
se factorise en
$B\xrightarrow{u_{i+1}}C_{i+1}\to C_{i+1}/\mathfrak J_{i+1}=C_i$~;
d'où la conclusion.
\end{proof}

\begin{proposition}[19.4.4]
\label{0.19.4.4-fr}
Soient $A$ un anneau topologique, $B$ une $A$-algèbre topologique
(commutative). Pour que $B$ soit une $A$-algèbre formellement lisse, il faut et
il suffit que pour

\oldpage[0\textsubscript{IV-1}]{87}

tout $B$-module topologique discret $L$, annulé par un idéal ouvert de $B$, on
ait (cf. (18.5.1))
\[
  \operatorname{Exalcotop}_A(B,L)=0.
\]
\end{proposition}

\begin{proof}
Soit $(\mathfrak R_\lambda)$ un système fondamental décroissant d'idéaux
ouverts de $B$ et posons $B_\lambda=B/\mathfrak R_\lambda$.
Considérons une $A$-algèbre topologique discrète $C$ et un idéal
$\mathfrak J$ de $C$ tel que $\mathfrak J^2=0$, de sorte que $C$ est une
$A$-extension de $C/\mathfrak J$ par $\mathfrak J$~; supposons donné un
$A$-homomorphisme $u:B_\lambda\to C/\mathfrak J$ et formons la
$A$-extension $E_\lambda$ de $B_\lambda$ par $\mathfrak J$, image
réciproque de $C$ par l'homomorphisme $u$ (18.2.5)~; c'est une $A$-algèbre
topologique pour la topologie discrète. Si
$\operatorname{Exalcotop}_A(B,L)=0$, il existe par définition (18.5.1) un
$\mu$ tel que $\mathfrak R_\mu\subset\mathfrak R_\lambda$ et que
l'extension image réciproque $E_\mu$ soit $A$-triviale~; mais cela signifie
(18.1.6) qu'il existe un $A$-homomorphisme continu
$v:B_\mu\to E_\lambda$ tel que l'homomorphisme canonique
$B_\mu\to B_\lambda$ se factorise en
$B_\mu\xrightarrow{v}E_\lambda\to B_\lambda$~; on en conclut
aussitôt que $B_\mu\to B_\lambda\xrightarrow{u}C/\mathfrak J$ se
factorise en
$B_\mu\xrightarrow{v}E_\lambda\to C\to C/\mathfrak J$, et cela
prouve, en vertu de (19.3.1) et (19.4.3), que $B$ est une $A$-algèbre
formellement lisse. La réciproque est immédiate, en appliquant (19.3.1) au cas
où $C$ est une $A$-algèbre topologique qui est une $A$-extension de
$B_\lambda$ par $L$, et à l'homomorphisme identique
$B_\lambda\to B_\lambda=C/L$.
\end{proof}

Lorsque $A$ et $B$ sont des anneaux discrets, le critère de lissité formelle
(19.4.4) se réduit à
\begin{equation}
\tag{19.4.4.1}
\label{0.19.4.4.1-fr}
\operatorname{Exalcom}_A(B,L)=0\qquad\text{pour tout $B$-module $L$};
\end{equation}
autrement dit, \emph{toute $A$-extension commutative de $B$ par un
$B$-module doit être $A$-triviale}.

\begin{corollary}[19.4.5]
\label{0.19.4.5-fr}
Soient $A$ un anneau topologique, $B$ une $A$-algèbre topologique
(commutative).
\begin{enumerate}
  \item[(i)] Supposons que $B$ soit une $A$-algèbre formellement lisse.
  Alors, pour tout idéal ouvert $\mathfrak R$ de $B$, tout
  $(B/\mathfrak R)$-module $L$ et tout $2$-cocycle de Hochschild
  $A$-bilinéaire et symétrique $f$ de
  $(B/\mathfrak R)\times(B/\mathfrak R)$ dans $L$ (18.4.3), il existe un
  idéal ouvert $\mathfrak R'\subset\mathfrak R$ tel que, si
  $\varphi:B/\mathfrak R'\to B/\mathfrak R$ est l'homomorphisme
  canonique, $f\circ(\varphi\times\varphi)$ soit un $2$-cobord de
  Hochschild $A$-bilinéaire de
  $(B/\mathfrak R')\times(B/\mathfrak R')$ dans $L$.
  \item[(ii)] Si $B$ est un $A$-module formellement projectif (19.2.1) et si
  la condition de (i) est vérifiée, $B$ est une $A$-algèbre formellement
  lisse.
\end{enumerate}
\end{corollary}

\begin{proof}
\begin{enumerate}
  \item[(i)] Le $2$-cocycle $f$ définit une $A$-extension de Hochschild $C$
  de $B/\mathfrak R$ par $L$ (18.4.3). Appliquant (19.3.1) à $C$, à l'idéal
  de carré nul $L$ de $C$ et à l'homomorphisme identique
  $B/\mathfrak R\to B/\mathfrak R=C/L$, on en déduit la condition (i) en
  vertu de (18.4.3).
  \item[(ii)] Appliquons le critère (19.4.3), en considérant un idéal ouvert
  $\mathfrak R$ de $B$, un idéal ouvert $\mathfrak J$ de $A$ tel que
  $\mathfrak JB\subset\mathfrak R$, et une $(A/\mathfrak J)$-extension $E$
  de $B/\mathfrak R$ par $L$. Comme $B$ est un $A$-module formellement
  projectif, l'application $A$-linéaire continue canonique
  $\rho:B\to B/\mathfrak R$ se factorise en
  $B\xrightarrow{w}E\to B/\mathfrak R$, où $w$ est une application
  $A$-linéaire continue (19.2.1), qui se factorise donc elle-même en
  $B\to B/\mathfrak R'\xrightarrow{w'}E$ où $\mathfrak R'$ est un idéal
  ouvert de $B$ contenu dans $\mathfrak R$~; remplaçant $\mathfrak J$ par
  un idéal plus petit, on peut supposer que
  $\mathfrak JB\subset\mathfrak R'$. Alors l'image réciproque par
  l'homomorphisme canonique $B/\mathfrak R'\to B/\mathfrak R$ de
  l'extension $E$ de $B/\mathfrak R$ par $L$, est une
  $(A/\mathfrak J)$-extension \emph{de Hochschild} $E'$ de
  $B/\mathfrak R'$ par $L$~;

  \oldpage[0\textsubscript{IV-1}]{88}

  appliquant (i) à un cocycle définissant cette extension (18.4.3), on en
  conclut qu'il y a un idéal ouvert
  $\mathfrak R''\subset\mathfrak R'$ de $B$ tel que l'image réciproque
  $E''$ de $E$ par $B/\mathfrak R''\to B/\mathfrak R$ soit
  $A$-triviale, C.Q.F.D.
\end{enumerate}
\end{proof}

\begin{corollary}[19.4.6]
\label{0.19.4.6-fr}
Soient $A$ un anneau topologique, $B$ une $A$-algèbre topologique qui soit un
$A$-module formellement projectif. Soit $A'$ une $A$-algèbre \emph{munie de
la topologie déduite de celle de $A$}. Supposons en outre que $A'$ soit un
$A$-module fidèlement plat, et que l'une des conditions suivantes soit
remplie~:
\begin{enumerate}
  \item[$1^\circ$] Il existe un système fondamental
  $(\mathfrak J_\lambda)$ d'idéaux ouverts de $A$ et un système fondamental
  $(\mathfrak M_\lambda)$ d'idéaux ouverts de $B$, ayant même ensemble
  d'indices et tels que, pour tout $\lambda$, on ait
  $\mathfrak J_\lambda B\subset\mathfrak M_\lambda$ et que
  $B/\mathfrak M_\lambda$ soit un
  $(A/\mathfrak J_\lambda)$-module projectif de type fini.
  \item[$2^\circ$] $A'$ est un $A$-module projectif de type fini.
\end{enumerate}
Alors, pour que $B'=B\otimes_A A'$ (muni de la topologie produit tensoriel)
soit une $A'$-algèbre formellement lisse, il faut et il suffit que $B$ soit une
$A$-algèbre formellement lisse.
\end{corollary}

\begin{proof}
La suffisance de la condition est contenue dans (19.3.5), (iii), sans aucune
hypothèse supplémentaire sur $B$ ni $A'$. Pour démontrer la réciproque, nous
allons appliquer le critère (19.4.5)~; sous l'hypothèse $2^\circ$, nous
désignerons encore par $(\mathfrak M_\lambda)$ un système fondamental
d'idéaux ouverts de $B$, et, pour tout $\lambda$, par
$\mathfrak J_\lambda$ un idéal ouvert de $A$ tel que
$\mathfrak J_\lambda B\subset\mathfrak M_\lambda$~; dans les deux
cas, nous poserons
$A_\lambda=A/\mathfrak J_\lambda$,
$B_\lambda=B/\mathfrak M_\lambda$,
$A'_\lambda=A_\lambda\otimes_A A'$,
$B'_\lambda=B_\lambda\otimes_A A'$. Soit $f$ un $2$-cocycle de
Hochschild $A_\lambda$-bilinéaire symétrique de
$B_\lambda\times B_\lambda$ dans un $B_\lambda$-module $L$~; par
extension des scalaires, on en déduit un $2$-cocycle de Hochschild
$f'=f\otimes1$, $A'_\lambda$-bilinéaire symétrique, de
$B'_\lambda\times B'_\lambda$ dans $L'=L\otimes_A A'$. Comme par
hypothèse $B'$ est une $A'$-algèbre formellement lisse, il existe, par (19.4.5),
(i), un indice $\mu$ tel que
$\mathfrak M_\mu\subset\mathfrak M_\lambda$, et tel que, si
$\varphi':B'_\mu\to B'_\lambda$ est l'application canonique,
$f'\circ(\varphi'\times\varphi')$ soit un $2$-cobord de Hochschild de
$B'_\mu\times B'_\mu$ dans $L'$~; autrement dit, son image $c'$ dans le
groupe de Hochschild
$\mathrm H^2_{A'_\mu}(B'_\mu,L')^s$ est nulle~; il est clair que si
$\varphi:B_\mu\to B_\lambda$ est l'homomorphisme canonique, et $c$ la
classe de $f\circ(\varphi\times\varphi)$ dans le groupe de Hochschild
$\mathrm H^2_{A_\mu}(B_\mu,L)^s$, $c'$ est l'image canonique de $c$.
Or, si $P_\bullet$ est le complexe relatif aux anneaux $A_\mu$ et
$B_\mu$ défini dans (18.4.5), servant au calcul de
$\mathrm H^2_{A_\mu}(B_\mu,L)^s$, le complexe analogue relatif aux
anneaux $A'_\mu$ et $B'_\mu$ est évidemment
$P_\bullet\otimes_A A'$~; dans l'hypothèse $1^\circ$, la construction de
$P_\bullet$ montre que c'est un $A_\mu$-module projectif de type fini.
On conclut donc de Bourbaki, \emph{Alg.}, chap.~II, 3\textsuperscript{e} éd., §~5,
n\textsuperscript{o}~3, prop.~7 que, sous les deux hypothèses, on a
\[
  \operatorname{Hom}_{A'_\mu}(P_\bullet\otimes_A A',L\otimes_A A')
  =\operatorname{Hom}_{A_\mu}(P_\bullet,L)\otimes_A A'
\]
à un isomorphisme canonique près~; comme $A'$ est un $A$-module plat, on a par
suite (18.4.5)
\[
  \mathrm H^2_{A'_\mu}(B'_\mu,L')^s
  =\mathrm H^2_{A_\mu}(B_\mu,L)^s\otimes_A A'
\]
et l'on peut donc écrire $c'=c\otimes1$~; mais comme $A'$ est un $A$-module
fidèlement plat, l'hypothèse $c'=0$ entraîne $c=0$, ce qui achève la
démonstration.
\end{proof}

\begin{proposition}[19.4.7]
\label{0.19.4.7-fr}
Soient $A$ un anneau topologique préadmissible, $\mathfrak J$ un idéal de
définition de $A$ ($\mathbf 0_{\rm I}$, 7.1.2), $B$ une $A$-algèbre
topologique qui soit un $A$-module formellement projectif. Considérons les
anneaux topologiques quotients $A_0=A/\mathfrak J$,
$B_0=B/\mathfrak JB$~; alors, pour que $B$ soit une $A$-algèbre formellement
lisse, il faut et il suffit que $B_0$ soit une $A_0$-algèbre formellement
lisse.
\end{proposition}

\oldpage[0\textsubscript{IV-1}]{89}

\begin{proof}
La nécessité de la condition résulte de (19.4.2). Pour voir qu'elle est
suffisante, notons d'abord qu'en considérant un système fondamental de
voisinages ouverts de $0$ dans $A$ formés d'idéaux
$\mathfrak J_\alpha$ contenus dans $\mathfrak J$, on peut, en vertu de
(19.4.2), se ramener au cas où $A$ est discret puisque
$B/\mathfrak J_\alpha B$ est un
$(A/\mathfrak J_\alpha)$-module formellement projectif (19.2.5)~; par
définition d'un anneau préadmissible, $\mathfrak J$ est alors
\emph{nilpotent}. Il suffit en outre de démontrer la proposition lorsque
$\mathfrak J^2=0$, car si $\mathfrak J^m=0$, on l'appliquera successivement
aux anneaux $A_k=A/\mathfrak J^k$ et $B_k=B/\mathfrak J^kB$ et à l'idéal
$\mathfrak J^{k-1}/\mathfrak J^k$ de $A_k$ pour $k=2,3,\ldots,m$, en
notant (19.2.5) que $B_k=B\otimes_A A_k$ est un $A_k$-module formellement
projectif. Appliquons le critère (19.4.5), (ii) en considérant un idéal ouvert
$\mathfrak R$ de $B$ et un $2$-cocycle de Hochschild $A$-bilinéaire
symétrique $f$ de $(B/\mathfrak R)\times(B/\mathfrak R)$ dans un
$(B/\mathfrak R)$-module $L$. Considérons d'abord le cas particulier où
$\mathfrak JL=0$, de sorte que $L$ peut aussi être considéré comme un
$(B_0/\mathfrak RB_0)$-module, et $f$ se factorise en
\[
  (B/\mathfrak R)\times(B/\mathfrak R)\to
  (B_0/\mathfrak RB_0)\times(B_0/\mathfrak RB_0)
  \xrightarrow{f_0}L,
\]
où $f_0$ est un $2$-cocycle de Hochschild bilinéaire symétrique. En vertu de
l'hypothèse, il y a donc un idéal ouvert
$\mathfrak R'\subset\mathfrak R$ dans $B$ et une application
$A_0$-linéaire $g_0:B_0/\mathfrak R'B_0\to L$ telle que
\[
  f_0(\varphi_0(x),\varphi_0(y))
  =xg_0(y)-g_0(xy)+g_0(x)y
\]
pour $x,y$ dans $B_0/\mathfrak R'B_0$, $\varphi_0$ étant
l'homomorphisme canonique
$B_0/\mathfrak R'B_0\to B_0/\mathfrak RB_0$. On en conclut aussitôt que
l'application $A$-linéaire composée
$g:B/\mathfrak R'\to B_0/\mathfrak R'B_0\xrightarrow{g_0}L$ est telle
que $dg=f\circ(\varphi\times\varphi)$, où
$\varphi:B/\mathfrak R'\to B/\mathfrak R$ est l'homomorphisme
canonique.

Passons maintenant au cas général, et considérons d'abord le
$(B/\mathfrak R)$-module $L/\mathfrak JL=L'$, pour lequel on a
$\mathfrak JL'=0$~; si $f'$ est l'application $A$-bilinéaire composée
$(B/\mathfrak R)\times(B/\mathfrak R)\xrightarrow{f}L\to L'$, $f'$ est
encore un $2$-cocycle de Hochschild symétrique, et en vertu de ce qui précède,
il existe un idéal ouvert $\mathfrak R'\subset\mathfrak R$ dans $B$ et
une application $A$-linéaire $g':B/\mathfrak R'\to L'$ vérifiant
$dg'=f'\circ(\varphi'\times\varphi')$ pour l'application canonique
$\varphi':B/\mathfrak R'\to B/\mathfrak R$. Comme $B$ est un
$A$-module formellement projectif, il existe un idéal ouvert
$\mathfrak R'_1\subset\mathfrak R'$ et une application $A$-linéaire
$g_1:B/\mathfrak R'_1\to L$ telle que l'homomorphisme
$B/\mathfrak R'_1\to B/\mathfrak R'\xrightarrow{g'}L'$ se factorise en
$B/\mathfrak R'_1\xrightarrow{g_1}L\to L'$. Considérons alors le
$2$-cocycle de Hochschild
\[
  f_1(x,y)=f(\varphi_1(x),\varphi_1(y))-xg_1(y)+g_1(xy)-g_1(x)y,
\]
application $A$-bilinéaire symétrique de
$(B/\mathfrak R'_1)\times(B/\mathfrak R'_1)$ dans $L$. Le choix de
$g_1$ entraîne que $f_1$ prend ses valeurs dans $\mathfrak JL$. Comme
$\mathfrak J(\mathfrak JL)=0$, on peut de nouveau appliquer le premier cas,
et il y a donc un idéal ouvert
$\mathfrak R''\subset\mathfrak R'_1$ et une application $A$-linéaire
$g_2:B/\mathfrak R''\to\mathfrak JL$ telle que
\[
  f_1(\varphi_2(x),\varphi_2(y))=xg_2(y)-g_2(xy)+g_2(x)y
\]
dans $(B/\mathfrak R'')\times(B/\mathfrak R'')$,
$\varphi_2:B/\mathfrak R''\to B/\mathfrak R'_1$ étant l'application
canonique~; l'application $A$-linéaire
$g=g_2+g_1\circ\varphi_2:B/\mathfrak R''\to L$ vérifie donc
$dg=f\circ(\varphi\times\varphi)$ pour l'application canonique
$\varphi:B/\mathfrak R''\to B/\mathfrak R$. C.Q.F.D.
\end{proof}

\subsection{Lissité formelle et anneaux gradués associés}
\label{subsection:0.19.5-fr}

\begin{env}[19.5.1]
\label{0.19.5.1-fr}
\oldpage[0\textsubscript{IV-1}]{90}
Soient $C$ un anneau topologique (commutatif), $V$ un $C$-module topologique,
et considérons l'algèbre symétrique
$\mathbf S_C^*(V)=\bigoplus_{n\geq0}\mathbf S_C^n(V)$, que nous allons munir
canoniquement d'une \emph{topologie linéaire}, compatible avec sa structure de
$C$-algèbre. Pour cela, considérons un sous-module ouvert $U$ de $V$, et
l'\emph{idéal gradué} $U.\mathbf S_C^*(V)$ qu'il engendre dans
$\mathbf S_C^*(V)$~; nous prendrons comme système fondamental de voisinages de
$0$ dans $\mathbf S_C^*(V)$ l'ensemble des sommes
$\mathfrak a.\mathbf S_C^*(V)+U.\mathbf S_C^*(V)$, où $\mathfrak a$ (resp.
$U$) parcourt un système fondamental d'idéaux ouverts (resp. de sous-modules
ouverts) de $C$ (resp. $V$). On notera que si la topologie de $V$ est moins
fine que la topologie déduite de celle de $C$, on peut se borner aux couples
$(\mathfrak a,U)$ tels que $\mathfrak aV\subset U$, donc
$\mathfrak a.\mathbf S_C^*(V)+U.\mathbf S_C^*(V)
=\mathfrak a.1_C+U.\mathbf S_C^*(V)$~; dans ce cas la topologie induite sur
chacun des $\mathbf S_C^n(V)$ pour $n\geq1$ a pour système fondamental de
voisinages de $0$ les $U.\mathbf S_C^{n-1}(V)$, où $U$ parcourt les
sous-modules ouverts de $V$~; en particulier, sur $V=\mathbf S_C^1(V)$, elle
coïncide avec la topologie donnée (en général elle est moins fine que cette
dernière). En outre, dans tous les cas, l'\emph{algèbre topologique}
$\mathbf S_C^*(V)$ ainsi définie est, pour les catégories des $C$-modules
topologiques et des $C$-algèbres topologiques, la solution du même problème
universel que pour les catégories de $C$-modules et de $C$-algèbres. En effet,
soit $E$ une $C$-algèbre topologique, $u:V\to E$ un homomorphisme de
$C$-modules, $\mathbf S(u)=\widetilde u$ son extension canonique à
$\mathbf S_C^*(V)$. Supposons $u$ continu~; alors, si $\mathfrak b$ est un
idéal ouvert de $E$, son image réciproque $U=u^{-1}(\mathfrak b)$ est un
sous-$C$-module ouvert de $V$, et l'image par $\widetilde u$ de
$U.\mathbf S_C^*(V)$ est donc contenue dans $\mathfrak b$~; comme par ailleurs
il existe un idéal ouvert $\mathfrak a$ de $C$ tel que
$\mathfrak a.E\subset\mathfrak b$, donc
$\widetilde u(\mathfrak a.\mathbf S_C^*(V))\subset\mathfrak b$, cela prouve
que $\widetilde u$ est continu. Réciproquement, si $\widetilde u$ est continu,
et si $\mathfrak b$ est un idéal ouvert de $E$, il existe un sous-module ouvert
$U$ de $V$ tel que $\widetilde u(U.\mathbf S_C^*(V))\subset\mathfrak b$, et en
particulier $\widetilde u(U.1_C)\subset\mathfrak b$, c'est-à-dire
$u(U)\subset\mathfrak b$, donc $u$ est continu. Rappelons en outre que l'on a
un isomorphisme canonique de $C$-modules topologiques (discrets)
\begin{equation}
\label{0.19.5.1.1-fr}
  \mathbf S_C^n(V)/U.\mathbf S_C^{n-1}(V)
  \xrightarrow{\sim}\mathbf S_C^n(V/U).
  \tag{19.5.1.1}
\end{equation}
\end{env}

\begin{env}[19.5.2]
\label{0.19.5.2-fr}
Soient $A$ un anneau topologique, $B$ une $A$-algèbre topologique
(commutative), $\mathfrak J$ un idéal de $B$ (non nécessairement ouvert ni
fermé)~; dans tout ce qui suit, on munit $\mathfrak J$ et les
$\mathfrak J^n$ ($n\geq2$) de la topologie \emph{induite} par celle de $B$,
les quotients $C=B/\mathfrak J$,
$\mathfrak J^n/\mathfrak J^{n+1}=\gr_{\mathfrak J}^n(B)$ de la topologie
\emph{quotient}, de sorte que les $\mathfrak J^n/\mathfrak J^{n+1}$ sont des
$C$-modules topologiques~; l'injection canonique
$\mathfrak J/\mathfrak J^2\to\gr_{\mathfrak J}^*(B)$ se prolonge en un
homomorphisme (d'abord non topologique) de $C$-algèbres
\[
  \varphi:\mathbf S_C^*(\mathfrak J/\mathfrak J^2)
  \longrightarrow\gr_{\mathfrak J}^*(B)
\]
qui pour chaque $n$ donne un homomorphisme \emph{surjectif} de $C$-modules
\begin{equation}
\label{0.19.5.2.1-fr}
  \varphi_n:\mathbf S_C^n(\mathfrak J/\mathfrak J^2)
  \longrightarrow\mathfrak J^n/\mathfrak J^{n+1}
  =\gr_{\mathfrak J}^n(B).
  \tag{19.5.2.1}
\end{equation}

\oldpage[0\textsubscript{IV-1}]{91}

Lorsque $\mathbf S_C^n(\mathfrak J/\mathfrak J^2)$ est muni de la topologie
définie dans (19.5.1), les homomorphismes $\varphi_n$ sont continus, en vertu
de la propriété universelle (19.5.1) de
$\mathbf S_C^*(\mathfrak J/\mathfrak J^2)$ appliquée à la $C$-algèbre
topologique $E=\gr_{\mathfrak J/\mathfrak J^n}(B/\mathfrak J^n)$ munie de la
topologie produit de celles des $\mathfrak J^k/\mathfrak J^{k+1}$, et à
l'injection canonique $\mathfrak J/\mathfrak J^2\to E$. Notons qu'ici la
topologie de $\mathfrak J/\mathfrak J^2$ est moins fine que la topologie
déduite de celle de $C$ (cette dernière topologie sur
$\mathfrak J/\mathfrak J^2$ étant aussi la topologie déduite de celle de $B$).
\end{env}

\begin{theorem}[19.5.3]
\label{0.19.5.3-fr}
Soient $A$ un anneau topologique, $B$ une $A$-algèbre topologique,
$\mathfrak J$ un idéal de $B$, $C=B/\mathfrak J$ la $A$-algèbre topologique
quotient. On suppose que la $A$-algèbre discrète $C$ est formellement lisse.
Alors~:
\begin{enumerate}
  \item[\rm(i)] Si $B$ est une $A$-algèbre formellement lisse,
  $\mathfrak J/\mathfrak J^2$ est un $C$-module topologique formellement
  projectif (19.2.1).
  \item[\rm(ii)] Si $B$ est une $A$-algèbre formellement lisse et un anneau
  préadmissible ($\mathbf 0_{\rm I}$, 7.1.2), les homomorphismes $\varphi_n$
  (19.5.2) sont des bimorphismes formels (19.1.2).
  \item[\rm(iii)] Réciproquement, supposons que $B$ soit préadmissible, que la
  suite $(\mathfrak J^n)$ tende vers $0$ dans $B$, que
  $\mathfrak J/\mathfrak J^2$ soit un $C$-module formellement projectif et que
  les $\varphi_n$ soient des bimorphismes formels. Alors $B$ est une
  $A$-algèbre formellement lisse.
\end{enumerate}
\end{theorem}

La démonstration de ce théorème est longue et encombrée de détails techniques~;
aussi commencerons-nous par en démontrer un corollaire plus simple, où
apparaissent plus nettement les idées directrices~; ce corollaire est d'ailleurs
le cas particulier du théorème (19.5.3) qui sera le plus fréquemment employé
par la suite.

\begin{corollary}[19.5.4]
\label{0.19.5.4-fr}
Soient $A$ un anneau topologique, $B$ une $A$-algèbre topologique,
$\mathfrak J$ un idéal de $B$ tel que la topologie de $B$ soit la topologie
$\mathfrak J$-préadique. On suppose que la $A$-algèbre discrète
$C=B/\mathfrak J$ est formellement lisse. Alors les trois conditions suivantes
sont équivalentes~:
\begin{enumerate}
  \item[a)] $B$ est une $A$-algèbre formellement lisse.
  \item[b)] $\mathfrak J/\mathfrak J^2$ est un $C$-module projectif et
  l'homomorphisme canonique
  \[
    \varphi:\mathbf S_C^*(\mathfrak J/\mathfrak J^2)
    \longrightarrow\gr_{\mathfrak J}^*(B)
  \]
  est bijectif.
  \item[c)] Le séparé complété $\widehat B$ de $B$ est une
  $\widehat A$-algèbre topologique isomorphe à une $\widehat A$-algèbre de la
  forme $\widehat{B'}$, où $B'=\mathbf S_C^*(V)$, $V$ étant un $C$-module
  projectif et $B'$ étant muni de la topologie ${B'}^+$-préadique, où
  ${B'}^+$ est l'idéal d'augmentation.
\end{enumerate}
\end{corollary}

\begin{env}[19.5.4.1]
\label{0.19.5.4.1-fr}
Prouvons d'abord que $a)$ entraîne que $\mathfrak J/\mathfrak J^2$ est un
$C$-module projectif. Soient donc $P$ et $Q$ deux $C$-modules,
$u:P\to Q$ un $C$-homomorphisme surjectif, et
$v:\mathfrak J/\mathfrak J^2\to Q$ un $C$-homomorphisme. L'anneau
$B/\mathfrak J^2$ étant considéré comme une $B$-extension de $C$ par
$\mathfrak J/\mathfrak J^2$, on en déduit par $v$ une $B$-extension
$E=(B/\mathfrak J^2)\oplus_{(\mathfrak J/\mathfrak J^2)}Q$ de $C$ par $Q$
(18.2.8). Comme $C$ est une $A$-algèbre discrète formellement lisse,
l'extension $E$ est $A$-\emph{triviale} (19.4.4) et peut donc s'identifier à
$\mathrm D_C(Q)$ (18.2.3). L'homomorphisme surjectif $u:P\to Q$ définit alors
canoniquement un $A$-homomorphisme surjectif
$\psi:\mathrm D_C(P)\to\mathrm D_C(Q)$ (18.2.9) dont le noyau est un idéal
$\mathfrak L$ de l'extension $E'=\mathrm D_C(P)$, contenu dans $P$, et
\emph{a fortiori} de carré nul. Soit $f:B\to E=E'/\mathfrak L$
l'homomorphisme définissant la structure de $B$-algèbre de $E$~: comme
$\mathfrak L$ est de carré nul, et que $B$ est une $A$-algèbre formellement
lisse, $f$ se factorise en $B\xrightarrow{g}E'\to E'/\mathfrak L$, où $g$ est
un $A$-homomorphisme continu.

\oldpage[0\textsubscript{IV-1}]{92}

Le diagramme
\[
  \xymatrix{
    E' \ar[r]^{\psi} & E \\
    B \ar[u]^{g} \ar[r]_{\theta} & B/\mathfrak J^2 \ar[u]_{v'}
  }
\]
où $v'$ est déduit de $v$ (18.2.8), est commutatif. D'ailleurs l'image de
$\mathfrak J$ par $\psi\circ g$ est contenue dans $Q$, donc l'image de
$\mathfrak J$ par $g$ est contenue dans $P+\mathfrak L=P$, et l'image de
$\mathfrak J^2$ par $g$ est nulle. Restreignant $g$ et $\theta$ à
$\mathfrak J$, on obtient un diagramme commutatif
\[
  \xymatrix{
    P \ar[r]^{u} & Q \\
    \mathfrak J/\mathfrak J^2 \ar[u]^{w} \ar[ur]_{v}
  }
\]
ce qui prouve que le $C$-module $\mathfrak J/\mathfrak J^2$ est projectif.
\end{env}

\begin{env}[19.5.4.2]
\label{0.19.5.4.2-fr}
Prouvons en second lieu que $a)$ entraîne que $\varphi$ est bijectif, ce qui
achèvera de montrer que $a)$ entraîne $b)$. Posons
$E_n=B/\mathfrak J^{n+1}$, et désignons par $F_n$ le quotient de
$\mathbf S_C^*(\mathfrak J/\mathfrak J^2)$ par la puissance $(n+1)$-ième de
son idéal d'augmentation. Comme $\mathfrak J/\mathfrak J^{n+1}$ est nilpotent
dans $E_n$ et que $C$ est une $A$-algèbre discrète formellement lisse,
l'automorphisme identique de $C$ se factorise en
$C\xrightarrow{f}E_n\to C=E_n/(\mathfrak J/\mathfrak J^{n+1})$ où $f$ est un
$A$-homomorphisme~; d'autre part, $\mathfrak J/\mathfrak J^2$ étant un
$C$-module projectif d'après (19.5.4.1), l'automorphisme identique de
$\mathfrak J/\mathfrak J^2$ se factorise en
$\mathfrak J/\mathfrak J^2\xrightarrow{g}\mathfrak J/\mathfrak J^{n+1}
\to\mathfrak J/\mathfrak J^2$, où $g$ est une application $C$-linéaire~; à
partir de $f$ et $g$ on obtient canoniquement un homomorphisme de $C$-algèbres
$\mathbf S_C^*(\mathfrak J/\mathfrak J^2)\to E_n$, qui d'ailleurs (par
définition de $g$) s'annule sur la puissance $(n+1)$-ième de l'idéal
d'augmentation de $\mathbf S_C^*(\mathfrak J/\mathfrak J^2)$, d'où, en
passant au quotient, un $A$-homomorphisme surjectif d'algèbres
$v:F_n\to E_n$ tel que $\gr^0(v)$ et $\gr^1(v)$ soient les automorphismes
identiques de $C$ et de $\mathfrak J/\mathfrak J^2$. Par définition de
l'homomorphisme canonique $\varphi$, on voit que $\gr^j(v)=\varphi_j$ pour tout
$j\leq n$. Notons maintenant que le noyau $\mathfrak R$ de $v$ est un idéal
nilpotent de $F_n$, de sorte que $E_n$ s'identifie à $F_n/\mathfrak R$. Comme
$B$ est une $A$-algèbre formellement lisse, le $A$-homomorphisme canonique
$p_n:B\to E_n=B/\mathfrak J^{n+1}$, qui est continu, se factorise en
$B\xrightarrow{w}F_n\xrightarrow{v}E_n$, où $w$ est un $A$-homomorphisme
continu~; comme $\gr^0(v)$ est l'identité, $w(\mathfrak J)$ est contenu dans
l'idéal d'augmentation de $F_n$, d'où $w(\mathfrak J^{n+1})=0$, si bien que
$w$ se factorise en
$B\xrightarrow{p_n}B/\mathfrak J^n=E_n\xrightarrow{w'}F_n$, et le composé
$E_n\xrightarrow{w'}F_n\xrightarrow{v}E_n$ est l'identité. En outre, comme
$\gr^0(v)$ et $\gr^1(v)$ sont les automorphismes identiques de $C$ et de
$\mathfrak J/\mathfrak J^2$, il en est de même de $\gr^0(w')$ et de
$\gr^1(w')$. Comme $\gr^*(E_n)$ est engendré par $\gr^1(E_n)$,
l'homomorphisme composé
\[
  \gr^j(F_n)\xrightarrow{\varphi_j}\gr^j(E_n)
  \xrightarrow{\gr^j(w')}\gr^j(F_n)
\]
est l'identité pour tout $j\leq n$, puisque c'est vrai pour $j=0$ et $j=1$~;
faisant $j=n$, on prouve ainsi que $\varphi_n$ est injectif, ce qui achève de
montrer que $a)$ entraîne $b)$.
\end{env}

\begin{env}[19.5.4.3]
\label{0.19.5.4.3-fr}
\oldpage[0\textsubscript{IV-1}]{93}
Prouvons ensuite que $b)$ entraîne $a)$. Le même raisonnement qu'au début de
(19.5.4.2) prouve l'existence d'un $A$-homomorphisme surjectif d'algèbres
$v:F_n\to E_n$ tel que $\gr^j(v)=\varphi_j$ pour tout $j$~; comme $\varphi_j$
est bijectif pour tout $j$ et que les filtrations de $F_n$ et $E_n$ sont
finies, on en conclut que $v$ est \emph{bijectif} (Bourbaki,
\emph{Alg. comm.}, chap.~III, §~2, n\textsuperscript{o}~8, cor.~3 du th.~1).
Soient maintenant $G$ une $A$-algèbre topologique discrète, $\mathfrak R$ un
idéal de carré nul dans $G$, $f:B\to G/\mathfrak R$ un $A$-homomorphisme
continu d'algèbres. Comme $G$ est discrète, il existe un entier $m$ tel que
$f$ s'annule dans $\mathfrak J^m$, donc $f$ se factorise en
$B\to E_n\xrightarrow{f_n}G/\mathfrak R$, où l'on prend $n=2m$. On a donc par
composition un $A$-homomorphisme continu d'algèbres
$r:C\to F_n\xrightarrow{v}E_n\xrightarrow{f_n}G/\mathfrak R$, et comme $C$
est une $A$-algèbre formellement lisse discrète, $r$ se factorise en
$C\xrightarrow{r'}G\to G/\mathfrak R$, de sorte que $G$ est muni par $r'$ d'une
structure de $C$-algèbre. D'autre part, la restriction à
$\mathfrak J/\mathfrak J^2$ de l'homomorphisme
$f'_n:F_n\xrightarrow{v}E_n\xrightarrow{f_n}G/\mathfrak R$ est une application
$C$-linéaire $g:\mathfrak J/\mathfrak J^2\to G/\mathfrak R$. Puisque
$\mathfrak J/\mathfrak J^2$ est un $C$-module projectif, $g$ se factorise en
$\mathfrak J/\mathfrak J^2\xrightarrow{h}G\to G/\mathfrak R$, où $h$ est une
application $C$-linéaire~; par prolongement, on en déduit un homomorphisme de
$C$-algèbres $w:\mathbf S_C^*(\mathfrak J/\mathfrak J^2)\to G$, et par
construction tout élément de degré $m$ a pour image par $w$ un élément de
$\mathfrak R$, donc tout élément de degré $n=2m$ a une image nulle puisque
$\mathfrak R^2=0$~; autrement dit, $w$ se factorise en
$\mathbf S_C^*(\mathfrak J/\mathfrak J^2)\to F_n\xrightarrow{w'}G$. Par
construction, le composé $F_n\xrightarrow{w'}G\xrightarrow{\theta}G/\mathfrak R$
coïncide avec $f'_n$ dans $C$ et dans $\mathfrak J/\mathfrak J^2$, donc est
égal à $f'_n$. Finalement, le composé
$B\to E_n\xrightarrow{v^{-1}}F_n\xrightarrow{w'}G\to G/\mathfrak R$ étant
égal à $f$, on voit que $B$ est une $A$-algèbre formellement lisse.
\end{env}

\begin{env}[19.5.4.4]
\label{0.19.5.4.4-fr}
Il reste à prouver l'équivalence de $a)$ et $c)$. En premier lieu, $c)$ entraîne
$a)$~: en effet, $B'$ est une $C$-algèbre formellement lisse pour la topologie
discrète (19.3.2), donc aussi pour la topologie ${B'}^+$-préadique (19.3.8)~;
comme $C$ est une $A$-algèbre formellement lisse, $B'$ est aussi une
$A$-algèbre formellement lisse (19.3.5, (ii)) et finalement $\widehat{B'}$ est
une $\widehat A$-algèbre formellement lisse (19.3.6), donc $B$ est une
$A$-algèbre formellement lisse (19.3.6). Reste à voir que $a)$ entraîne $c)$.
Notons d'abord que puisque $C$ est une $A$-algèbre formellement lisse,
l'homomorphisme identique $C\to B/\mathfrak J$ se factorise en
$C\to\widehat B\to B/\mathfrak J$, où $C\to\widehat B$ est un
$A$-homomorphisme (19.3.10)~; $\widehat B$, et par suite tous les
$B/\mathfrak J^{n+1}$, sont donc munis de structures de $C$-algèbres. D'autre
part, comme $\mathfrak J/\mathfrak J^2$ est un $C$-module projectif par $b)$,
l'injection canonique $\mathfrak J/\mathfrak J^2\to B/\mathfrak J^2$ permet de
former un système projectif de $C$-homomorphismes
$u_n:\mathfrak J/\mathfrak J^2\to B/\mathfrak J^{n+1}$ pour $n\geq1$, donc
aussi (par la propriété universelle de l'algèbre symétrique) un système
projectif d'homomorphismes de $C$-algèbres
$v_n:\mathbf S_C^*(\mathfrak J/\mathfrak J^2)=B'\to B/\mathfrak J^{n+1}$~;
d'ailleurs il est clair que $v_n$ s'annule dans $({B'}^+)^{n+1}$ et comme
$\gr^0(v_n)=\varphi_0$ et $\gr^1(v_n)=\varphi_1$, on a
$\gr^j(v_n)=\varphi_j$ pour $0\leq j\leq n$. Comme les $\varphi_j$ sont des
isomorphismes par $b)$, et que les filtrations de
$B'/({B'}^+)^{n+1}$ et de $B/\mathfrak J^{n+1}$ sont finies, on en conclut que
$v_n$ est \emph{bijectif} pour

\oldpage[0\textsubscript{IV-1}]{94}

tout $n$ (Bourbaki, \emph{Alg. comm.}, chap.~III, §~2,
n\textsuperscript{o}~8, cor.~3 du th.~1)~; d'où $c)$ par passage à la limite
projective.
\end{env}

\begin{remark}[19.5.5]
\label{0.19.5.5-fr}
Sous les hypothèses générales de (19.5.4), supposons en outre que
$\mathfrak J$ soit un idéal \emph{maximal} de $B$, de sorte que
$C=B/\mathfrak J$ soit un corps, et que $\mathfrak J/\mathfrak J^2$ soit un
$C$-espace vectoriel de dimension \emph{finie}. Alors les conditions $a)$,
$b)$ et $c)$ de (19.5.4) sont encore équivalentes à la suivante~:
\begin{enumerate}
  \item[d)] \emph{Étant donnés une algèbre de polynômes}
  $E=C[T_1,\ldots,T_n]$ \emph{sur le corps} $C$, \emph{et en désignant par}
  $\mathfrak n$ \emph{l'idéal maximal de} $E$ \emph{engendré par les} $T_i$
  $(1\leq i\leq n)$, \emph{alors, pour tout idéal} $\mathfrak b$ \emph{de}
  $E$ \emph{contenant une puissance} $\mathfrak n^k$, \emph{tout
  homomorphisme} $v:B\to E/\mathfrak b$ \emph{de} $A$-\emph{algèbres}
  $C$-\emph{augmentées se factorise en}
  $B\xrightarrow{w}E/\mathfrak n^k\to E/\mathfrak b$, \emph{où} $w$ \emph{est
  un} $A$-\emph{homomorphisme}.
\end{enumerate}
En effet, $\mathfrak J/\mathfrak J^2$ étant ici un $C$-module libre, il suffit
de vérifier que la condition $d)$ entraîne que l'homomorphisme canonique
$\varphi:\mathbf S_C^*(\mathfrak J/\mathfrak J^2)\to\gr_{\mathfrak J}^*(B)$
est bijectif. Or ici $\mathbf S_C^*(\mathfrak J/\mathfrak J^2)$ s'identifie à
l'algèbre de polynômes $E=C[T_1,\ldots,T_n]$, où
$n=\operatorname{rg}_C(\mathfrak J/\mathfrak J^2)$, et l'idéal d'augmentation
de $E$ est l'idéal $\mathfrak n$ engendré par les $T_i$. Pour tout entier
$k\geq0$, l'hypothèse que $C$ est une $A$-algèbre discrète formellement lisse
entraîne, comme dans (19.5.4.2), l'existence d'un $A$-homomorphisme surjectif
$f:E/\mathfrak n^{k+1}\to B/\mathfrak J^{k+1}$, tel que
$\gr^j(f)=\varphi_j$ pour tout $j\leq k$. Si
$\mathfrak b/\mathfrak n^{k+1}=\operatorname{Ker}(f)$,
$B/\mathfrak J^{k+1}$ s'identifie donc à $E/\mathfrak b$, et l'hypothèse $d)$
permet de factoriser l'homomorphisme canonique $B\to B/\mathfrak J^{k+1}$ en
$B\xrightarrow{w}E/\mathfrak n^{k+1}\xrightarrow{f}B/\mathfrak J^{k+1}$~;
comme $\gr^0(f)$ est l'identité, on a
$w(\mathfrak J)\subset\mathfrak n/\mathfrak n^{k+1}$, donc
$w(\mathfrak J^{k+1})=0$~; on conclut comme dans (19.5.4.2) que $\varphi_k$
est injectif pour tout $k$, ce qui achève de prouver notre assertion.
\end{remark}

\begin{env}[19.5.6]
\label{0.19.5.6-fr}
\emph{Démonstration du théorème} (19.5.3). --- Soit
$(\mathfrak b_\alpha)$ un système fondamental décroissant d'idéaux ouverts de
$B$. Nous poserons
\[
  B_\alpha=B/\mathfrak b_\alpha,\qquad
  C_\alpha=B/(\mathfrak b_\alpha+\mathfrak J),\qquad
  \mathfrak J_\alpha=(\mathfrak b_\alpha+\mathfrak J)/\mathfrak b_\alpha,
\]
de sorte que
\[
  C_\alpha=B_\alpha/\mathfrak J_\alpha\qquad\text{et}\qquad
  \mathfrak J_\alpha^n/\mathfrak J_\alpha^{n+1}
  =(\mathfrak b_\alpha+\mathfrak J^n)/(\mathfrak b_\alpha+\mathfrak J^{n+1}).
\]
Les $C$-modules
$((\mathfrak b_\alpha\cap\mathfrak J^n)+\mathfrak J^{n+1})/
\mathfrak J^{n+1}$ de $\mathfrak J^n/\mathfrak J^{n+1}$ forment un système
fondamental de sous-modules ouverts du $C$-module topologique
$\mathfrak J^n/\mathfrak J^{n+1}$~; comme
$(\mathfrak b_\alpha\cap\mathfrak J^n)+\mathfrak J^{n+1}
=\mathfrak J^n\cap(\mathfrak b_\alpha+\mathfrak J^{n+1})$, on a
\[
  \mathfrak J^n/((\mathfrak b_\alpha\cap\mathfrak J^n)+\mathfrak J^{n+1})
  =\mathfrak J^n/(\mathfrak J^n\cap(\mathfrak b_\alpha+\mathfrak J^{n+1}))
  =(\mathfrak b_\alpha+\mathfrak J^n)/(\mathfrak b_\alpha+\mathfrak J^{n+1})
  =\mathfrak J_\alpha^n/\mathfrak J_\alpha^{n+1}.
\]
\end{env}

\begin{env}[19.5.6.1]
\label{0.19.5.6.1-fr}
Soient $P$, $Q$ deux $C_\alpha$-modules discrets,
$u:P\to Q$ un $C_\alpha$-homomorphisme surjectif. L'anneau discret
$B/(\mathfrak b_\alpha+\mathfrak J^2)$ est une $B$-extension de $C_\alpha$ par
l'idéal de carré nul $\mathfrak J_\alpha/\mathfrak J_\alpha^2$. Soit
$v:\mathfrak J/\mathfrak J^2\to Q$ un $C$-homomorphisme continu~; remplaçant
au besoin $\mathfrak b_\alpha$ par un idéal ouvert plus petit de $B$, on peut
supposer que le noyau de $v$ contient le sous-$C$-module ouvert
$((\mathfrak b_\alpha\cap\mathfrak J)+\mathfrak J^2)/\mathfrak J^2$~;
passant au quotient, on en déduit un $C_\alpha$-homomorphisme de modules
discrets $v':\mathfrak J_\alpha/\mathfrak J_\alpha^2\to Q$~; soit
$E_\alpha$ la $B_\alpha$-extension de $C_\alpha$ par $Q$ déduite de
$B/(\mathfrak b_\alpha+\mathfrak J^2)$ au moyen de $v'$, et soit
$v'_*:B/(\mathfrak b_\alpha+\mathfrak J^2)\to E_\alpha$ le
$B_\alpha$-homomorphisme correspondant (18.2.8)~; $E_\alpha$ est une
$A$-algèbre topologique discrète, et l'isomorphisme canonique
$C_\alpha\xrightarrow{\sim}E_\alpha/Q$ donne par composition un
$A$-homomorphisme continu $f:C\to C_\alpha\to E_\alpha/Q$. Mais comme $Q$ est
de carré nul dans $E_\alpha$ et que $C$ est une $A$-algèbre formellement
lisse, $f$ se factorise en $C\xrightarrow{g}E_\alpha\to E_\alpha/Q$ où $g$
est un $A$-homomorphisme continu. Comme $g$ est continu et $E_\alpha$ discret,
il existe un $\beta\geq\alpha$ tel que $g$ se factorise en
$C\to C_\beta\xrightarrow{g'}E_\alpha$. D'autre part, soit $E_\beta$ la
$A$-extension de $C_\beta$ par $Q$, image réciproque de $E_\alpha$ par
l'homomorphisme canonique $C_\beta\to C_\alpha$~; l'existence de $g'$ (tel que
le composé $C_\beta\to E_\alpha\to C_\alpha$ soit l'homomorphisme canonique)
équivaut au fait que $E_\beta$ est une extension $A$-\emph{triviale}, et peut
donc s'identifier à $\mathrm D_{C_\beta}(Q)$. Cela étant, l'homomorphisme
surjectif $u:P\to Q$ définit canoniquement un homomorphisme surjectif
$\mathrm D_{C_\beta}(P)\to\mathrm D_{C_\beta}(Q)$ (18.2.9), dont le noyau est
un idéal $\mathfrak L$ de l'extension $E'_\beta=\mathrm D_{C_\beta}(P)$
contenu dans $P$, et \emph{a fortiori} de carré nul. Soit
$f':B\to E_\beta=E'_\beta/\mathfrak L$ le $A$-homomorphisme continu définissant
la structure de $B$-algèbre topologique de $E_\beta$~; comme $\mathfrak L$ est
de carré nul et que $B$ est une $A$-algèbre formellement lisse, $f'$ se
factorise en $B\xrightarrow{h'}E'_\beta\to E'_\beta/\mathfrak L$, où $h'$ est
un $A$-homomorphisme continu. La construction de $E'_\beta$ donne par
composition un $A$-homomorphisme $\psi:E'_\beta\to E_\beta\to E_\alpha$, et
il est clair que le diagramme
\[
  \xymatrix{
    E'_\beta \ar[r]^{\psi} & E_\alpha \\
    B \ar[u]^{h'} \ar[r]_{\theta} &
    B/(\mathfrak b_\alpha+\mathfrak J^2) \ar[u]_{v'_*}
  }
\]
est commutatif. D'ailleurs, l'image par $\psi\circ h'$ de $\mathfrak J$ est
contenue dans $Q$, donc l'image par $h'$ de $\mathfrak J$ est contenue dans
$P$, et l'image de $\mathfrak J^2$ par $h'$ est nulle. Restreignant $h'$ et
$\theta$ à $\mathfrak J$, on obtient un diagramme commutatif
\[
  \xymatrix{
    P \ar[r]^{u} & Q \\
    \mathfrak J/\mathfrak J^2 \ar[u]^{w} \ar[ur]_{v}
  }
\]
où $w$ est continu, ce qui prouve que
$\mathfrak J/\mathfrak J^2$ est un $C$-module formellement projectif.
\end{env}

\oldpage[0\textsubscript{IV-1}]{95}

\begin{env}[19.5.6.2]
\label{0.19.5.6.2-fr}
Pour tout entier $n\geq0$, nous poserons
\[
  E_n=B/\mathfrak J^{n+1},
  \qquad\text{de sorte que}\qquad E_0=C~;
\]
les idéaux $(\mathfrak b_\alpha+\mathfrak J^{n+1})/\mathfrak J^{n+1}$ forment
un système fondamental d'idéaux ouverts dans $E_n$, et nous poserons
\[
  E_{\alpha,n}=B/(\mathfrak b_\alpha+\mathfrak J^{n+1})
  =B_\alpha/\mathfrak J_\alpha^{n+1},
\]
anneau quotient discret. De même, dans
$\gr_{\mathfrak J}^n(B)=\mathfrak J^n/\mathfrak J^{n+1}$, nous avons vu que les
$((\mathfrak b_\alpha\cap\mathfrak J^n)+\mathfrak J^{n+1})/
\mathfrak J^{n+1}$ forment un système fondamental de voisinages de $0$, le
quotient de $\mathfrak J^n/\mathfrak J^{n+1}$ par ce sous-module s'identifiant
canoniquement à $\mathfrak J_\alpha^n/\mathfrak J_\alpha^{n+1}$. Considérons
l'algèbre symétrique
$\mathbf S_{C_\alpha}^*(\mathfrak J_\alpha/\mathfrak J_\alpha^2)$~; nous
noterons $F_{\alpha,n}$ le quotient de
$\mathbf S_{C_\alpha}^*(\mathfrak J_\alpha/\mathfrak J_\alpha^2)$ par la
puissance $(n+1)$-ième de son idéal d'augmentation. Pour un $n\geq1$ fixé, il
résulte de (19.5.1.1) que les
$\mathbf S_{C_\alpha}^n(\mathfrak J_\alpha/\mathfrak J_\alpha^2)$ sont les
quotients de $\mathbf S_C^n(\mathfrak J/\mathfrak J^2)$ par un système
fondamental de sous-modules ouverts dans ce $C$-module topologique.

Pour abréger le langage, nous dirons que pour $\alpha\leq\beta$, les
homomorphismes canoniques $B_\beta\to B_\alpha$, $C_\beta\to C_\alpha$,
$\mathfrak J_\beta/\mathfrak J_\beta^2\to
\mathfrak J_\alpha/\mathfrak J_\alpha^2$, $E_{\beta,n}\to E_{\alpha,n}$,
$F_{\beta,n}\to F_{\alpha,n}$, etc., sont les \emph{homomorphismes de
transition}.
\end{env}

\begin{lemma}[19.5.6.3]
\label{0.19.5.6.3-fr}
Supposons que la $A$-algèbre $C$ soit formellement lisse, l'anneau $B$
préadmissible et le $C$-module $\mathfrak J/\mathfrak J^2$ formellement
projectif. Alors~:
\begin{enumerate}
  \item[\rm(i)] Pour tout $\alpha$, il existe $\beta\geq\alpha$ et un
  $A$-homomorphisme surjectif d'algèbres
  \[
    v_{\alpha\beta}:F_{\beta,n}\longrightarrow E_{\alpha,n}
  \]
  tel que $\gr^0(v_{\alpha\beta}):C_\beta\to C_\alpha$ et
  $\gr^1(v_{\alpha\beta}):\mathfrak J_\beta/\mathfrak J_\beta^2
  \to\mathfrak J_\alpha/\mathfrak J_\alpha^2$ soient les homomorphismes de
  transition.
  \item[\rm(ii)] Si $\beta$ et $v_{\alpha\beta}$ vérifient les conditions de
  \emph{(i)}, alors, pour tout $\gamma\geq\beta$, il existe
  $\delta\geq\gamma$ et un $A$-homomorphisme surjectif d'algèbres
  $v_{\gamma\delta}:F_{\delta,n}\to E_{\gamma,n}$ vérifiant les conditions de
  \emph{(i)} (pour $\gamma$ et $\delta$) et rendant commutatif le diagramme
  \begin{equation}
  \label{0.19.5.6.4-fr}
  \xymatrix@C=3em{
    F_{\beta,n} \ar[r]^{v_{\alpha\beta}} & E_{\alpha,n} \\
    F_{\delta,n} \ar[u] \ar[r]_{v_{\gamma\delta}} & E_{\gamma,n} \ar[u]
  }
  \tag{19.5.6.4}
  \end{equation}
  où les flèches verticales sont les homomorphismes de transition.
\end{enumerate}
\end{lemma}

\begin{env}[19.5.6.4]
\label{0.19.5.6.4-proof-fr}
\begin{enumerate}
  \item[\rm(i)] Dans la $A$-algèbre topologique discrète $E_{\alpha,n}$,
  l'idéal $\mathfrak J_\alpha/\mathfrak J_\alpha^{n+1}$ est nilpotent, et
  l'isomorphisme identique
  $C_\alpha\xrightarrow{\sim}
  E_{\alpha,n}/(\mathfrak J_\alpha/\mathfrak J_\alpha^{n+1})$ donne par
  composition un $A$-homomorphisme continu
  $C\to C_\alpha\to
  E_{\alpha,n}/(\mathfrak J_\alpha/\mathfrak J_\alpha^{n+1})$~;

\oldpage[0\textsubscript{IV-1}]{96}

  comme $C$ est une $A$-algèbre formellement lisse, cet homomorphisme se
  factorise en
  $C\xrightarrow{f_\alpha}E_{\alpha,n}\to
  E_{\alpha,n}/(\mathfrak J_\alpha/\mathfrak J_\alpha^{n+1})$, où $f_\alpha$
  est un $A$-homomorphisme continu~; par suite,
  $\mathfrak J_\alpha/\mathfrak J_\alpha^{n+1}$ devient, au moyen de
  $f_\alpha$, un $C$-module topologique discret annulé par un idéal ouvert de
  $C$. L'hypothèse que $\mathfrak J/\mathfrak J^2$ est un $C$-module
  formellement projectif entraîne alors l'existence d'une application
  $C$-linéaire continue
  $g_\alpha:\mathfrak J/\mathfrak J^2\to
  \mathfrak J_\alpha/\mathfrak J_\alpha^{n+1}$ rendant commutatif le diagramme
  \[
    \xymatrix{
      \mathfrak J_\alpha/\mathfrak J_\alpha^{n+1} \ar[r] &
      \mathfrak J_\alpha/\mathfrak J_\alpha^2 \\
      & \mathfrak J/\mathfrak J^2 \ar[ul]^{g_\alpha} \ar[u]
    }
  \]
  Comme $f_\alpha$ et $g_\alpha$ sont continus, il y a un
  $\beta\geq\alpha$ tel que ces homomorphismes se factorisent respectivement en
  \[
    C\to C_\beta\xrightarrow{f'_{\alpha\beta}}E_{\alpha,n},
  \]
  \[
    \mathfrak J/\mathfrak J^2\to
    \mathfrak J_\beta/\mathfrak J_\beta^2
    \xrightarrow{g'_{\alpha\beta}}
    \mathfrak J_\alpha/\mathfrak J_\alpha^{n+1}
  \]
  et à partir de $f'_{\alpha\beta}$ et $g'_{\alpha\beta}$, on obtient donc
  canoniquement un homomorphisme de $C_\beta$-algèbres
  \[
    \mathbf S_{C_\beta}^*(\mathfrak J_\beta/\mathfrak J_\beta^2)
    \longrightarrow E_{\alpha,n}
  \]
  qui d'ailleurs (par définition de $g'_{\alpha\beta}$) s'annule dans la
  puissance $(n+1)$-ième de l'idéal d'augmentation de
  $\mathbf S_{C_\beta}^*(\mathfrak J_\beta/\mathfrak J_\beta^2)$~; passant au
  quotient par cette puissance $(n+1)$-ième, on obtient l'homomorphisme
  $v_{\alpha\beta}$ cherché, compte tenu des définitions de $f_\alpha$ et
  $g_\alpha$~; la surjectivité de $v_{\alpha\beta}$ résulte en effet de celle
  des deux homomorphismes $\gr^0(v_{\alpha\beta})$ et
  $\gr^1(v_{\alpha\beta})$, puisque cela entraîne que
  $\gr(v_{\alpha\beta})$ est surjectif (l'algèbre
  $\gr^*(E_{\alpha,n})$ étant engendrée par $\gr^0(E_{\alpha,n})$ et
  $\gr^1(E_{\alpha,n})$), et comme les filtrations considérées sont finies, on
  peut appliquer Bourbaki, \emph{Alg. comm.}, chap.~III, §~2,
  n\textsuperscript{o}~8, cor.~3 du th.~1.

  \item[\rm(ii)] L'hypothèse que $B$ est \emph{préadmissible} signifie que
  l'on peut supposer tous les $\mathfrak b_\alpha$ contenus dans un même
  $\mathfrak b_{\alpha_0}$ dont les puissances \emph{tendent vers} $0$. Cela
  entraîne en particulier que le noyau de l'homomorphisme de transition
  $E_{\gamma,n}\to E_{\alpha,n}$, égal à
  $(\mathfrak b_\alpha+\mathfrak J^{n+1})/
  (\mathfrak b_\gamma+\mathfrak J^{n+1})$, est \emph{nilpotent}~; appliquant le
  lemme (19.3.10.1), on voit qu'on peut supposer $f_\gamma$ choisi de sorte que
  le diagramme
  \[
    \xymatrix{
      & E_{\alpha,n} \\
      C \ar[ur]^{f_\alpha} \ar[r]_{f_\gamma} & E_{\gamma,n} \ar[u]
    }
  \]
  (où la flèche verticale est l'homomorphisme de transition) soit commutatif.
  L'hypothèse que $\mathfrak J/\mathfrak J^2$ est un $C$-module formellement
  projectif permet d'autre part de choisir
  $g_\gamma:\mathfrak J/\mathfrak J^2\to
  \mathfrak J_\gamma/\mathfrak J_\gamma^{n+1}$ de sorte que le diagramme
  \[
    \xymatrix@C=5em@R=3em{
      \mathfrak J_\alpha/\mathfrak J_\alpha^{n+1} \ar[r] &
      \mathfrak J_\alpha/\mathfrak J_\alpha^2 \\
      \mathfrak J_\gamma/\mathfrak J_\gamma^{n+1} \ar[u] \ar[r] &
      \mathfrak J_\gamma/\mathfrak J_\gamma^2 \ar[u] \\
      & \mathfrak J/\mathfrak J^2
        \ar[ul]_{g_\gamma} \ar[u]_{p_\gamma} \ar[uul]^{g_\alpha}
    }
  \]
  soit commutatif (il suffit de remarquer que l'image par $(g_\alpha,p_\gamma)$
  de $\mathfrak J/\mathfrak J^2$ dans le module produit
  $(\mathfrak J_\alpha/\mathfrak J_\alpha^{n+1})
  \times(\mathfrak J_\gamma/\mathfrak J_\gamma^2)$ est contenue (vu la
  définition de $g_\alpha$ et de la relation $\beta\leq\gamma$) dans l'image
  canonique $Q$ dans ce module produit du module
  $P=\mathfrak J_\gamma/\mathfrak J_\gamma^{n+1}$, et d'appliquer la définition
  (19.2.1) à l'homomorphisme \emph{surjectif} $P\to Q$ et à l'homomorphisme
  $(g_\alpha,p_\gamma)$ de $\mathfrak J/\mathfrak J^2$ dans $Q$). Ce choix de
  $f_\gamma$ et de $g_\gamma$ permet alors, en construisant
  $v_{\gamma\delta}$ comme dans \emph{(i)}, d'obtenir en outre la commutativité
  du diagramme (19.5.6.4).
\end{enumerate}
\end{env}

\oldpage[0\textsubscript{IV-1}]{97}

\begin{lemma}[19.5.6.5]
\label{0.19.5.6.5-fr}
Supposons que les $A$-algèbres $B$ et $C$ soient formellement lisses et que
$B$ soit préadmissible (de sorte qu'en vertu de (19.5.6.1), les conditions de
(19.5.6.3) sont satisfaites). Pour tout système de deux indices
$\alpha\leq\beta$ et d'un homomorphisme $v_{\alpha\beta}$ vérifiant les
conditions de (19.5.6.3, (i)), il existe un indice $\lambda\geq\beta$ et un
$A$-homomorphisme d'algèbres
\[
  w_{\beta\lambda}:E_{\lambda,n}\longrightarrow F_{\beta,n}
\]
tel que~: $1^\circ$ $\gr^0(w_{\beta\lambda}):C_\lambda\to C_\beta$ et
$\gr^1(w_{\beta\lambda}):
\mathfrak J_\lambda/\mathfrak J_\lambda^2\to
\mathfrak J_\beta/\mathfrak J_\beta^2$ soient les homomorphismes de
transition~; $2^\circ$ le composé
$E_{\lambda,n}\xrightarrow{w_{\beta\lambda}}F_{\beta,n}
\xrightarrow{v_{\alpha\beta}}E_{\alpha,n}$ soit l'homomorphisme de transition.
\end{lemma}

\begin{proof}
Appliquons le lemme (19.5.6.3, (ii)) avec $\gamma=\beta$, ce qui donne un
$\delta\geq\beta$ et un
$v_{\beta\delta}:F_{\delta,n}\to E_{\beta,n}$. Rappelons que
$v_{\beta\delta}$ est surjectif~; d'autre part, son noyau $\mathfrak R$ est
\emph{nilpotent}~: en effet, l'idéal d'augmentation $F_{\delta,n}^+$ de
$F_{\delta,n}$ est nilpotent par définition, et le noyau de
$\gr^0(v_{\beta\delta}):C_\delta\to C_\beta$ est aussi nilpotent en vertu de
l'hypothèse que $B$ est préadmissible. On voit donc que $E_{\beta,n}$
s'identifie à $F_{\delta,n}/\mathfrak R$. Comme $B$ est une $A$-algèbre
formellement lisse, le $A$-homomorphisme canonique
$p_{\beta,n}:B\to E_{\beta,n}$, qui est continu, se factorise en
$B\xrightarrow{w}F_{\delta,n}\xrightarrow{v_{\beta\delta}}E_{\beta,n}$, où
$w$ est un $A$-homomorphisme continu. Comme $F_{\delta,n}$ est discret, il
existe $\lambda\geq\delta$ tel que $w$ soit nul dans $\mathfrak b_\lambda$,
donc $w$ se factorise en
$B\to B/\mathfrak b_\lambda\xrightarrow{w'_\lambda}F_{\delta,n}$. Considérons
l'homomorphisme composé
$B/\mathfrak b_\lambda\xrightarrow{w'_\lambda}F_{\delta,n}
\xrightarrow{q_{\beta\delta}}F_{\beta,n}$, où $q_{\beta\delta}$ est
l'homomorphisme de transition~; notons que
$\gr^0(q_{\beta\delta})=\gr^0(v_{\beta\delta}):C_\delta\to C_\beta$ est
l'homomorphisme de transition~; comme le composé
$v_{\alpha\beta}\circ q_{\beta\delta}\circ w'_\lambda$ est l'homomorphisme
canonique $B/\mathfrak b_\lambda\to
B/(\mathfrak b_\alpha+\mathfrak J^{n+1})$, cela montre que l'image de
$\mathfrak J_\lambda=(\mathfrak b_\lambda+\mathfrak J)/\mathfrak b_\lambda$
par $q_{\beta\delta}\circ w'_\lambda$ est contenue dans l'idéal d'augmentation
$F_{\beta,n}^+$, et par suite l'image de
$\mathfrak J_\lambda^{n+1}=(\mathfrak b_\lambda+\mathfrak J^{n+1})/
\mathfrak b_\lambda$ est nulle. Autrement dit,
$q_{\beta\delta}\circ w'_\lambda$ se factorise en
\[
  B/\mathfrak b_\lambda\to
  B/(\mathfrak b_\lambda+\mathfrak J^{n+1})=E_{\lambda,n}
  \xrightarrow{w_{\beta\lambda}}F_{\beta,n}
\]
tel que le composé
$E_{\lambda,n}\xrightarrow{w_{\beta\lambda}}F_{\beta,n}
\xrightarrow{v_{\alpha\beta}}E_{\alpha,n}$ soit l'homomorphisme de transition.
Le raisonnement précédent montre aussi que $\gr^0(w_{\beta\lambda})$, qui est
le composé
$B/(\mathfrak b_\lambda+\mathfrak J)\to\gr^0(F_{\delta,n})
\xrightarrow{\gr^0(q_{\beta\delta})}\gr^0(F_{\beta,n})$, est l'homomorphisme
de transition $C_\lambda\to C_\beta$, puisque
$v_{\beta\delta}\circ w$ est l'homomorphisme canonique. En outre, on a aussi
$\gr^1(q_{\beta\delta})=\gr^1(v_{\beta\delta})$ (19.5.6.3), donc le même
raisonnement prouve que $\gr^1(w_{\beta\lambda})$ est l'homomorphisme de
transition.
\end{proof}

\begin{env}[19.5.6.6]
\label{0.19.5.6.6-fr}
\emph{Démonstration de} (19.5.3, (ii)). --- Pour montrer que $\varphi_n$ est un
bimorphisme formel, nous allons appliquer le critère de (19.1.3, (iii)). Les
conditions du lemme (19.5.6.5) étant satisfaites par hypothèse, déterminons,
pour tout indice $\alpha$, $v_{\alpha\beta}$ et $w_{\beta\lambda}$ vérifiant
les conclusions de ce lemme. L'homomorphisme
\[
  \gr^n(v_{\alpha\beta}):\gr^n(F_{\beta,n})
  \longrightarrow\gr^n(E_{\alpha,n})
\]
n'est autre que l'homomorphisme
\[
  \varphi_{\alpha\beta,n}:
  \mathbf S_{C_\beta}^n(\mathfrak J_\beta/\mathfrak J_\beta^2)
  \longrightarrow\mathfrak J_\alpha^n/\mathfrak J_\alpha^{n+1}
\]
déduit de l'homomorphisme canonique $\varphi_n$ (19.5.2.1) par passage aux
quotients~; en effet, il résulte de (19.5.6.3) que $\gr^0(v_{\alpha\beta})$ et
$\gr^1(v_{\alpha\beta})$ coïncident respectivement avec
$\varphi_{\alpha\beta,0}$ et $\varphi_{\alpha\beta,1}$, et la définition de
l'homomorphisme canonique $\varphi$ montre alors, par récurrence sur $j\leq n$,
que $\gr^j(v_{\alpha\beta})$ et $\varphi_{\alpha\beta,j}$ sont égaux pour tout
$j$. Cela étant, comme $\varphi_n$ est surjectif, c'est \emph{a fortiori} un
épimorphisme formel~; en outre, pour $j\leq n$, l'homomorphisme composé
\[
  \gr^j(F_{\lambda,n})\xrightarrow{\varphi_{\lambda\lambda,j}}
  \gr^j(E_{\lambda,n})\xrightarrow{\gr^j(w_{\beta\lambda})}
  \gr^j(F_{\beta,n})
\]
est l'homomorphisme de transition, car c'est vrai pour $j=0$ et $j=1$ en
vertu de (19.5.6.5), et comme $\gr^*(E_{\lambda,n})$ est engendré par
$\gr^1(E_{\lambda,n})$ cela démontre l'assertion par récurrence sur $j$.
Composant avec l'homomorphisme de transition
$\gr^j(F_{\beta,n})\to\gr^j(F_{\alpha,n})$, on voit donc, pour $j=n$, que
l'on a factorisé l'homomorphisme de transition
$\gr^n(F_{\lambda,n})\to\gr^n(F_{\alpha,n})$ en
\[
  \gr^n(F_{\lambda,n})\xrightarrow{\varphi_{\lambda\lambda,n}}
  \gr^n(E_{\lambda,n})\longrightarrow\gr^n(F_{\alpha,n}),
\]
ce qui est la condition du critère (19.1.3) pour que $\varphi_n$ soit un
bimorphisme formel.
\end{env}

\oldpage[0\textsubscript{IV-1}]{98}

\begin{lemma}[19.5.6.7]
\label{0.19.5.6.7-fr}
Supposons que $B$ soit préadmissible, que $\mathfrak J/\mathfrak J^2$ soit un
$C$-module formellement projectif, que $C$ soit une $A$-algèbre formellement
lisse et que les $\varphi_n$ soient des bimorphismes formels. Alors, pour tout
couple d'indices $\alpha$, $\beta$ tel que $\alpha\leq\beta$ et tout
homomorphisme $v_{\alpha\beta}:F_{\beta,n}\to E_{\alpha,n}$ vérifiant les
conditions de (19.5.6.3, (i)), il existe un indice $\gamma\geq\beta$ tel que,
pour tout indice $\mu\geq\gamma$, il y ait un diagramme commutatif de
$A$-homomorphismes
\[
\xymatrix@C=4em@R=3em{
  F_{\beta,n}\ar[r]^{v_{\alpha\beta}} & E_{\alpha,n} \\
  & E_{\mu,n}\ar[ul]_{u_{\beta\mu}}\ar[u]_{p_{\alpha\mu}}
}
\]
où $p_{\alpha\mu}$ est l'homomorphisme de transition.
\end{lemma}

\begin{proof}
Appliquant le critère (19.1.3) à chacun des $\varphi_j$ pour $0\leq j\leq n$,
on voit qu'il existe un indice $\gamma\geq\beta$ et des homomorphismes
uniquement déterminés (et surjectifs) de $C_\gamma$-modules
\[
  w_{\beta\gamma}^j:\gr^j(E_{\gamma,n})\longrightarrow
  \gr^j(F_{\beta,n})
\]
tels que les composés
\[
  \gr^j(F_{\gamma,n})
  \xrightarrow{\varphi_{\gamma\gamma,j}}\gr^j(E_{\gamma,n})
  \xrightarrow{w_{\beta\gamma}^j}\gr^j(F_{\beta,n})
\]
soient les homomorphismes de transition (le fait qu'on puisse choisir le même
indice $\gamma$ pour tous les $\varphi_j$ résulte de ce qu'ils sont en nombre
fini). En outre, l'unicité des $w_{\beta\gamma}^j$ prouve (puisque $\varphi$ est
un homomorphisme d'algèbres graduées) que
$w_{\beta\gamma}=(w_{\beta\gamma}^j)_{0\leq j\leq n}$ est un homomorphisme de
$C_\gamma$-algèbres graduées
\[
  \gr^*(E_{\gamma,n})\longrightarrow
  \gr^*(F_{\beta,n})=F_{\beta,n}.
\]
En outre, comme $\varphi_{\gamma\gamma,0}$ et $\varphi_{\gamma\gamma,1}$ sont
les homomorphismes identiques,
$w_{\beta\gamma}^0:C_\gamma\to C_\beta$ et
$w_{\beta\gamma}^1:\mathfrak J_\gamma/\mathfrak J_\gamma^2
\to\mathfrak J_\beta/\mathfrak J_\beta^2$ sont les homomorphismes de
transition, et il en est donc de même de
$\gr^0(v_{\alpha\beta}\circ w_{\beta\gamma})$ et de
$\gr^1(v_{\alpha\beta}\circ w_{\beta\gamma})$~; comme
$\gr^*(E_{\gamma,n})$ est engendré par $\gr^1(E_{\gamma,n})$, on en conclut que
$\gr^j(v_{\alpha\beta}\circ w_{\beta\gamma})$ est aussi l'homomorphisme de
transition pour $0\leq j\leq n$.

Appliquons maintenant à $\alpha$, $\beta$ et $\gamma$ le lemme (19.5.6.3,
(ii)), ce qui donne le diagramme (19.5.6.4) avec $\delta\geq\gamma$~; puis
répétons le raisonnement du début en remplaçant $\alpha$ et $\beta$ par
$\gamma$ et $\delta$. On obtient ainsi un indice $\lambda\geq\delta$ et un
diagramme commutatif d'homomorphismes
\[
\xymatrix@C=4em@R=4em{
  \gr^*(E_{\gamma,n})\ar[r]^{w_{\beta\gamma}} &
  F_{\beta,n}\ar[r]^{v_{\alpha\beta}} & E_{\alpha,n} \\
  \gr^*(E_{\lambda,n})\ar[u]_{\gr^*(p_{\gamma\lambda})}
    \ar[r]_{w_{\delta\lambda}} &
  F_{\delta,n}\ar[u]_{q_{\beta\delta}}\ar[r]_{v_{\gamma\delta}} &
  E_{\gamma,n}\ar[u]_{p_{\alpha\gamma}}
    \ar@{-->}[ul]_{u_{\beta\gamma}}
}
\]
où les flèches verticales sont les homomorphismes de transition. Tout revient
à prouver l'existence de l'homomorphisme $u_{\beta\gamma}$ laissant le
diagramme commutatif, et pour cela il suffit évidemment de montrer que l'on a
\[
  \operatorname{Ker}(v_{\gamma\delta})
  \subset\operatorname{Ker}(q_{\beta\delta}),
\]
$v_{\gamma\delta}$ et $q_{\beta\delta}$ étant surjectifs. Comme
$w_{\delta\lambda}$ est surjectif, cette dernière relation équivaut à
\[
\begin{split}
  \operatorname{Ker}(v_{\gamma\delta}\circ w_{\delta\lambda})
  &\subset\operatorname{Ker}(q_{\beta\delta}\circ w_{\delta\lambda})\\
  &=\operatorname{Ker}\bigl(w_{\beta\gamma}\circ
  (\gr(p_{\gamma\lambda}))\bigr).
\end{split}
\]
Mais on a vu ci-dessus que
$\gr(p_{\gamma\lambda})=\gr(v_{\gamma\delta}\circ w_{\delta\lambda})$ et il
est clair que
\[
\begin{split}
  \operatorname{Ker}(v_{\gamma\delta}\circ w_{\delta\lambda})
  &\subset\operatorname{Ker}\bigl(\gr(v_{\gamma\delta}\circ
  w_{\delta\lambda})\bigr)\\
  &=\operatorname{Ker}(\gr(p_{\gamma\lambda}))
  \subset\operatorname{Ker}\bigl(w_{\beta\gamma}\circ
  (\gr(p_{\gamma\lambda}))\bigr),
\end{split}
\]
ce qui achève la démonstration du lemme, car pour tout $\mu\geq\gamma$, il
suffira de définir $u_{\beta\mu}$ comme le composé
\[
  E_{\mu,n}\xrightarrow{p_{\gamma\mu}}E_{\gamma,n}
  \xrightarrow{u_{\beta\gamma}}F_{\beta,n}.
\]
\end{proof}

\begin{env}[19.5.6.8]
\label{0.19.5.6.8-fr}
\emph{Démonstration de} (19.5.3, (iii)). --- Soient $G$ une $A$-algèbre
topologique discrète, $\mathfrak R$ un idéal de carré nul dans $G$,
$f:B\to G/\mathfrak R$ un $A$-homomorphisme continu d'algèbres. Comme
$G/\mathfrak R$ est discrète, il y a un indice $\alpha$ tel que $f$ s'annule
dans $\mathfrak b_\alpha$~; par hypothèse, il existe un entier $m$ tel que
$\mathfrak J^m\subset\mathfrak b_\alpha$, donc $f$ se factorise en
\[
  B\xrightarrow{p_\alpha}E_{\alpha,n}
  \xrightarrow{f_{\alpha,n}}G/\mathfrak R,
\]
où l'on prend $n=2m$. Les hypothèses du lemme (19.5.6.3) étant satisfaites, on
a tout d'abord un $\beta\geq\alpha$ et un $A$-homomorphisme composé
\begin{equation}
\label{0.19.5.6.9-fr}
  f_{\alpha,n}\circ v_{\alpha\beta}:F_{\beta,n}
  \xrightarrow{v_{\alpha\beta}}E_{\alpha,n}
  \xrightarrow{f_{\alpha,n}}G/\mathfrak R.
  \tag{19.5.6.9}
\end{equation}
Comme $F_{\beta,n}$ est une $C_\beta$-algèbre, donc \emph{a fortiori} une
$C$-algèbre topologique (discrète), cela donne par composition un
$A$-homomorphisme continu
$r:C\to F_{\beta,n}\to G/\mathfrak R$. Comme d'autre part $C$ est une
$A$-algèbre formellement lisse,

\oldpage[0\textsubscript{IV-1}]{99}

cet homomorphisme se factorise en
$r:C\xrightarrow{r'}G\to G/\mathfrak R$, où $r'$ est un
$A$-homomorphisme continu, si bien que $G$ est ainsi muni d'une structure de
$C$-algèbre topologique (discrète). D'autre part, par composition avec
l'homomorphisme canonique
\[
  \mathfrak J/\mathfrak J^2\longrightarrow
  \mathfrak J_\beta/\mathfrak J_\beta^2\longrightarrow F_{\beta,n},
\]
on déduit de (19.5.6.9) un $C$-homomorphisme continu
$g:\mathfrak J/\mathfrak J^2\to G/\mathfrak R$. Comme
$\mathfrak J/\mathfrak J^2$ est un $C$-module formellement projectif, $g$ se
factorise en
\[
  \mathfrak J/\mathfrak J^2\xrightarrow{h}G\longrightarrow G/\mathfrak R,
\]
où $h$ est un $C$-homomorphisme continu. Comme $G$ est discret, il existe un
indice $\gamma\geq\beta$ tel que l'image par $h$ de
$((\mathfrak b_\gamma\cap\mathfrak J)+\mathfrak J^2)/\mathfrak J^2$ soit
nulle, ainsi que l'image par $r'$ de
$(\mathfrak b_\gamma+\mathfrak J)/\mathfrak J$, de sorte que $h$ se factorise
en
\[
  \mathfrak J/\mathfrak J^2\longrightarrow
  \mathfrak J_\gamma/\mathfrak J_\gamma^2
  \xrightarrow{h'}G,
\]
où $h'$ est un $C_\gamma$-homomorphisme. Par prolongement, on déduit donc de
$h'$ un homomorphisme de $C_\gamma$-algèbres
\[
  w:\mathbf S_{C_\gamma}^*(\mathfrak J_\gamma/\mathfrak J_\gamma^2)
  \longrightarrow G
\]
et par construction, tout élément de degré $m$ a pour image par $w$ un élément
de $\mathfrak R$, donc tout élément de degré $n=2m$ a une image nulle, puisque
$\mathfrak R^2=0$~; autrement dit, $w$ se factorise en
\[
  \mathbf S_{C_\gamma}^*(\mathfrak J_\gamma/\mathfrak J_\gamma^2)
  \longrightarrow F_{\gamma,n}\xrightarrow{w'_\gamma}G.
\]
Appliquons maintenant à
\[
  v_{\alpha\gamma}:F_{\gamma,n}
  \xrightarrow{q_{\beta\gamma}}F_{\beta,n}
  \xrightarrow{v_{\alpha\beta}}E_{\alpha,n}
\]
le lemme (19.5.6.7), dont les hypothèses sont vérifiées~; il existe donc un
$\delta\geq\gamma$ et un homomorphisme
$u_{\gamma\delta}:E_{\delta,n}\to F_{\gamma,n}$ tel que le composé
$v_{\alpha\gamma}\circ u_{\gamma\delta}$ soit l'homomorphisme de transition
$p_{\alpha\delta}:E_{\delta,n}\to E_{\alpha,n}$. On obtient donc finalement un
diagramme commutatif de $A$-homomorphismes continus
\[
\xymatrix@C=3em@R=2.5em{
  && E_{\alpha,n}\ar[r]^{f_{\alpha,n}} & G/\mathfrak R \\
  && F_{\beta,n}\ar[u]_{v_{\alpha\beta}} & \\
  B\ar[r]_{p_\delta} & E_{\delta,n}\ar[uur]^{p_{\alpha\delta}}
    \ar[r]_{u_{\gamma\delta}} &
  F_{\gamma,n}\ar[u]_{q_{\beta\gamma}}\ar[r]^{w'_\gamma} & G\ar[uu]
}
\]
et le composé $f':B\to G$ des homomorphismes de la ligne inférieure est donc
tel que $f$ se factorise en $B\xrightarrow{f'}G\to G/\mathfrak R$, ce qui
prouve que $B$ est une $A$-algèbre formellement lisse.

La démonstration du th.~(19.5.3) est ainsi achevée.
\end{env}

\begin{corollary}[19.5.7]
\label{0.19.5.7-fr}
Soient $A$ un anneau topologique, $B$ une $A$-algèbre topologique,
$(\mathfrak b_\lambda)$ un système fondamental d'idéaux ouverts dans $B$,
$\mathfrak J$ un idéal de $B$, $C=B/\mathfrak J$ la $A$-algèbre topologique
quotient. On pose $C_\lambda=B/(\mathfrak b_\lambda+\mathfrak J)$. On suppose
que~: $1^\circ$ pour tout $n$, la topologie induite sur $\mathfrak J^n$ par
celle de $B$ est aussi la topologie du $B$-module $\mathfrak J^n$ déduite de
la topologie de $B$ (19.0.2) (cette condition sera remplie en particulier si
$B$ est noethérien et sa topologie préadique ($\mathbf 0_{\mathrm I}$,
7.3.2))~; $2^\circ$ $C$ est une $A$-algèbre formellement lisse. Sous ces
conditions~:
\begin{enumerate}
  \item[\rm(i)] Si $B$ est une $A$-algèbre formellement lisse, alors, pour tout
  $\lambda$, $(\mathfrak J/\mathfrak J^2)\otimes_C C_\lambda$ est un
  $C_\lambda$-module projectif.
  \item[\rm(ii)] Si $B$ est une $A$-algèbre formellement lisse et un anneau
  préadmissible, alors pour tout $\lambda$, l'homomorphisme canonique
  \begin{equation}
  \label{0.19.5.7.1-fr}
    \varphi_\lambda=\varphi\otimes 1_{C_\lambda}:
    \mathbf S_{C_\lambda}^*((\mathfrak J/\mathfrak J^2)\otimes_C C_\lambda)
    \longrightarrow\gr_{\mathfrak J}^*(B)\otimes_C C_\lambda
    \tag{19.5.7.1}
  \end{equation}
  est bijectif.

  \oldpage[0\textsubscript{IV-1}]{100}

  \item[\rm(iii)] Réciproquement, supposons que $B$ soit préadmissible, que la
  suite $(\mathfrak J^n)$ tende vers $0$, et que, pour tout $\lambda$,
  $(\mathfrak J/\mathfrak J^2)\otimes_C C_\lambda$ soit un $C_\lambda$-module
  projectif et l'homomorphisme (19.5.4.1) soit bijectif. Alors $B$ est une
  $A$-algèbre formellement lisse.
\end{enumerate}
\end{corollary}

\begin{proof}
En effet, la topologie de $\mathfrak J/\mathfrak J^2$ et celle des
$\mathbf S_C^n(\mathfrak J/\mathfrak J^2)$ sont alors aussi déduites de celle
de $B$ (19.5.1)~; les conclusions sont alors des conséquences immédiates de
(19.5.3) et des critères de (19.1.4) et (19.2.4) spéciaux à ce type de
topologies.
\end{proof}

\begin{remark}[19.5.8]
\label{0.19.5.8-fr}
Supposons que $\mathfrak J/\mathfrak J^2$ soit un $C$-module de type fini et
que, pour la topologie quotient de celle de $B$, $C$ soit un anneau de Zariski~;
soit $\mathfrak r$ un idéal de définition de $C$, de sorte que les
$\mathfrak r^n$ forment un système fondamental de voisinages de $0$ dans $C$.
Alors les conclusions de (i) et (ii) dans (19.5.7) peuvent être remplacées par
les suivantes~:
\begin{enumerate}
  \item[\rm(i')] $\mathfrak J/\mathfrak J^2$ est un $C$-module projectif.
  \item[\rm(ii')] L'homomorphisme canonique
  $\varphi:\mathbf S_C^*(\mathfrak J/\mathfrak J^2)
  \to\gr_{\mathfrak J}^*(B)$ est bijectif.
\end{enumerate}
En effet, il est clair que (i') entraîne la conclusion de (i) dans (19.5.7).
Inversement, si $(\mathfrak J/\mathfrak J^2)\otimes_C C_\lambda$ est un
$C_\lambda$-module projectif pour tout $\lambda$, alors
$(\mathfrak J/\mathfrak J^2)\otimes_C(C/\mathfrak r^n)$ est un
$(C/\mathfrak r^n)$-module projectif (donc plat) pour tout $n$~; on en conclut
que $\mathfrak J/\mathfrak J^2$ est un $C$-module plat
($\mathbf 0_{\mathrm{III}}$, 10.2.2), donc projectif puisqu'il est de
présentation finie (Bourbaki, \emph{Alg. comm.}, chap.~II, \textsection5,
n\textsuperscript{o}~2, cor.~2 du th.~1). D'autre part, les $C$-modules
$\mathbf S_C^n(\mathfrak J/\mathfrak J^2)$ et $\gr_{\mathfrak J}^n(B)$ sont
de type fini, et on sait que lorsque $C$ est un anneau de Zariski, il revient
au même de dire alors que $\varphi_n$ est bijectif ou que
$\widehat\varphi_n$ est bijectif (Bourbaki, \emph{Alg. comm.}, chap.~III,
\textsection3, n\textsuperscript{o}~5, prop.~9), donc (ii) est équivalent à
(ii').
\end{remark}

\subsection{Cas des algèbres sur un corps}
\label{subsection:0.19.6-fr}

\begin{theorem}[19.6.1) (Cohen]
\label{0.19.6.1-fr}
Soient $k$ un corps, $K$ une extension de $k$, $k$ et $K$
étant munis des topologies discrètes. Pour que $K$ soit une $k$-algèbre
formellement lisse, il faut et il suffit que $K$ soit une extension séparable
de $k$.
\end{theorem}

\begin{proof}
La nécessité de la condition sera établie en (19.6.5.1) (et naturellement ne
sera pas utilisée jusque-là)~; nous nous bornerons ici à prouver que la
condition est \emph{suffisante}. Distinguons deux cas~:

\emph{I. --- $K$ est une extension séparable de type fini de $k$.} On sait
alors (Bourbaki, \emph{Alg.}, chap.~V, \textsection9,
n\textsuperscript{o}~3, th.~2) qu'il existe une sous-extension pure
$K'=k(T_1,\ldots,T_n)$ de $K$ telle que $K$ soit une extension algébrique
finie séparable de $K'$. Compte tenu de (19.3.5, (iii)), on peut donc se
borner, soit au cas où $K=K'$, soit au cas où $K$ est algébrique fini sur $k$.
Dans le premier cas, on sait que $A=k[T_1,\ldots,T_n]$ est une $k$-algèbre
formellement lisse (19.3.3), donc il en est de même de
$k(T_1,\ldots,T_n)=S^{-1}A$, où $S=A-\{0\}$ (19.3.5, (iv)). Dans le second
cas, \emph{tous} les groupes de cohomologie de Hochschild
$\mathrm H_k^n(K,L)$, pour un $(K,K)$-bimodule quelconque $L$, sont nuls~: en
effet, si l'on considère la $k$-algèbre produit tensoriel $C=K\otimes_k K$,
$L$ est un $C$-module à gauche et le groupe de cohomologie
$\mathrm H_k^n(K,L)$ est égal à $\operatorname{Ext}_C^n(K,L)$, où $K$ est
aussi considéré comme $(K,K)$-bimodule,

\oldpage[0\textsubscript{IV-1}]{101}

donc comme $C$-module à gauche (M, IX, 4). Or, comme $K$ est une extension
finie séparable de $k$, on sait que $K\otimes_k K$ est composée directe
d'extensions de $k$, dont l'une est $K$ lui-même (Bourbaki, \emph{Alg.},
chap.~VIII, \textsection8, prop.~3)~; cela entraîne donc que $K$ est un
$C$-module \emph{projectif}, d'où notre assertion. Toute $k$-extension de $K$
de noyau $L$ est donc $k$-triviale, et \emph{a fortiori} les $k$-extensions
commutatives le sont, d'où le théorème dans ce cas (19.4.4).

\emph{II. --- Cas général.} --- Avec les notations de (18.4.5), il s'agit de
prouver que $\operatorname{Hom}_K(\mathrm H_2(P'_\bullet),L)=0$ pour
\emph{tout} $K$-espace vectoriel $L$, ce qui signifie évidemment que
$\mathrm H_2(P'_\bullet)=0$. Si $K$ est réunion d'une famille filtrante de
sous-extensions $K^{(\alpha)}$ de $k$, $P'_\bullet$ est \emph{limite
inductive} des complexes correspondants $P'^{(\alpha)}_\bullet$, puisque le
foncteur $\varinjlim$ commute au produit tensoriel dans la catégorie des
$k$-modules~; par l'exactitude du foncteur $\varinjlim$ dans cette catégorie,
on a donc
$\mathrm H_2(P'_\bullet)=\varinjlim\mathrm H_2(P'^{(\alpha)}_\bullet)$.
Comme l'hypothèse que $K$ est séparable entraîne qu'il en est de même de toute
sous-extension de $K$ et que $K$ est réunion des sous-extensions de type fini,
la première partie de la démonstration entraîne bien que $K$ est une
$k$-algèbre formellement lisse. C.Q.F.D.
\end{proof}

\begin{corollary}[19.6.2]
\label{0.19.6.2-fr}
Soient $A$ un anneau local séparé et complet, $K$ son corps résiduel, $k$ un
sous-corps de $K$ tel que $K$ soit une extension séparable de $k$. Alors il
existe un sous-corps $K'$ de $A$ contenant $k$, tel que la restriction à $K'$
de l'homomorphisme canonique $A\to K$ soit un isomorphisme de $K'$ sur $K$.
\end{corollary}

\begin{proof}
Soit $\mathfrak m$ l'idéal maximal de $A$. En vertu de l'hypothèse et de
(19.6.1), $K$ est une $k$-algèbre formellement lisse~; appliquant (19.3.10) en
remplaçant $C$ par $A$ et $\mathfrak J$ par $\mathfrak m$, on voit que
l'automorphisme identique de $K=A/\mathfrak m$ se factorise en
$K\xrightarrow{u}A\to A/\mathfrak m$, où $u$ est un $k$-homomorphisme, donc
nécessairement un $k$-isomorphisme de $K$ sur un sous-corps $K'$ contenant
$k$.
\end{proof}

\begin{corollary}[19.6.3]
\label{0.19.6.3-fr}
Soient $A$ un anneau local noethérien complet, $\mathfrak m$ son idéal maximal,
$k$ son corps résiduel. Les conditions suivantes sont équivalentes~:
\begin{enumerate}
  \item[a)] $A$ contient un sous-corps.
  \item[b)] Si $p$ est la caractéristique de $k$, on a $pA=0$.
  \item[c)] Il existe un corps $K$ tel que $A$ soit isomorphe à un anneau
  quotient d'un anneau de séries formelles $B=K[[T_1,\ldots,T_n]]$.
\end{enumerate}
Lorsqu'il en est ainsi, on peut supposer que $A$ est isomorphe à $B/\mathfrak b$,
où $\mathfrak b$ est contenu dans le carré de l'idéal maximal de $B$.
\end{corollary}

\begin{proof}
Il est immédiat que $c)$ implique $a)$, car $A$ est $\neq0$ puisque c'est un
anneau local~; comme $K$ est un corps et que l'homomorphisme composé
$K\xrightarrow{\varphi}B\to A$ (où $\varphi$ est l'injection canonique) n'est
pas nul, il est injectif. Pour voir que $a)$ entraîne $b)$, il suffit de
remarquer que si $K$ est un sous-corps de $A$, $K$ est isomorphe à un
sous-corps de $k$ et a par suite même caractéristique $p$~; comme $p.1=0$ dans
$K$, donc dans $A$, on a $pA=0$. Réciproquement, $b)$ entraîne $a)$, car le
composé des homomorphismes canoniques
$\mathbf Z\xrightarrow{j}A\to k$ a pour noyau $p\mathbf Z$, et l'on a
$j(p\mathbf Z)=0$, donc $\operatorname{Ker}(j)=p\mathbf Z$ et $A$ contient un
anneau $A'$ isomorphe à $\mathbf Z/p\mathbf Z$ et ne rencontrant pas
$\mathfrak m$~; si $p>0$, $A'$ est un corps~; sinon, tout

\oldpage[0\textsubscript{IV-1}]{102}

élément de $A'$ étant inversible dans $A$, le corps des fractions de $A'$
(isomorphe à $\mathbf Q$) est aussi contenu dans $A$.

Prouvons enfin que $a)$ entraîne $c)$~; la condition $a)$ entraîne que $A$
contient un sous-corps premier $k_0$, isomorphe au sous-corps premier de $k$,
et dont $k$ est par suite extension séparable. Appliquant (19.6.2), on voit
d'abord qu'il existe un isomorphisme $f$ de $k$ sur un sous-corps $K$ de $A$.
D'autre part, soit $(x_i)_{1\leq i\leq n}$ un système d'éléments de
$\mathfrak m$ tel que les classes mod.~$\mathfrak m^2$ des $x_i$ forment une
base (sur $k$) de $\mathfrak m/\mathfrak m^2$. Puisque $A$ est complet, il y a
alors un homomorphisme continu d'anneaux
$u:k[[T_1,\ldots,T_n]]\to A$ tel que $u$ soit égal à $f$ dans $k$ et que
$u(T_i)=x_i$ pour tout $i$, et cet homomorphisme est surjectif en vertu du
choix des $x_i$ (Bourbaki, \emph{Alg. comm.}, chap.~III, \textsection2,
n\textsuperscript{o}~9, prop.~11).
\end{proof}

\begin{theorem}[19.6.4]
\label{0.19.6.4-fr}
Soient $k$ un corps, $A$ une $k$-algèbre locale noethérienne, $\mathfrak m$ son
idéal maximal, $K=A/\mathfrak m$ son corps résiduel, $A$ étant muni de la
topologie $\mathfrak m$-préadique. On suppose que $K$ soit une extension
séparable de $k$. Alors les conditions suivantes sont équivalentes~:
\begin{enumerate}
  \item[a)] $A$ est une $k$-algèbre formellement lisse.
  \item[b)] Le complété $\widehat A$ de $A$ est $k$-isomorphe à un anneau de
  séries formelles $K[[T_1,\ldots,T_n]]$ (dont la structure de $k$-algèbre est
  définie par l'homomorphisme composé
  \[
    k\xrightarrow{\varphi}K\xrightarrow{\psi}K[[T_1,\ldots,T_n]],
  \]
  où $\psi$ est l'injection canonique et $\varphi$ l'homomorphisme définissant
  la structure d'extension de $K$).
  \item[b')] $\widehat A$ est un anneau isomorphe à un anneau de séries
  formelles $K[[T_1,\ldots,T_n]]$.
  \item[c)] $A$ est un anneau local régulier.
\end{enumerate}
\end{theorem}

\begin{proof}
Le fait que $a)$ entraîne $b)$ est le cas particulier de ((19.5.4),
équivalence de $a)$ et $c)$), appliqué en y remplaçant $A$ par $k$, $B$ par
$A$ et $\mathfrak J$ par $\mathfrak m$, compte tenu de (19.6.1). Il est clair
que $b)$ entraîne $b')$ et $b')$ entraîne $c)$ par (17.1.1). Enfin si $c)$ est
vérifiée, il résulte de (17.1.1, $a)$) et de ((19.5.4), équivalence de $b)$ et
$a)$), appliqué à $K=A/\mathfrak m$ remplaçant $C$, que $A$ est une
$k$-algèbre formellement lisse.
\end{proof}

\begin{corollary}[19.6.5]
\label{0.19.6.5-fr}
Soient $k$ un corps, $A$ une $k$-algèbre locale noethérienne, $\mathfrak m$ son
idéal maximal, $A$ étant muni de la topologie $\mathfrak m$-préadique.
Supposons que $A$ soit une $k$-algèbre formellement lisse~; alors $A$ est
géométriquement régulier sur $k$, autrement dit (cf.~$\mathbf{IV}$, 6.7.6),
pour toute extension finie $k'$ de $k$, l'anneau semi-local
$A'=A\otimes_k k'$ est régulier (17.3.6). En outre, si $K$ est le corps
résiduel de $A$, $\widehat A$ est isomorphe à un anneau de séries formelles
$K[[T_1,\ldots,T_n]]$.
\end{corollary}

\begin{proof}
Comme $\widehat A'=\widehat A\otimes_k k'$, il résulte de (19.3.6) et de
(17.1.5) (appliqué aux anneaux locaux en les idéaux maximaux de $A'$) qu'on
peut se borner au cas où $A$ est complet~; alors $A'$ est une $k'$-algèbre
formellement lisse (19.3.5, (iii)) et par ailleurs est composée directe de
$k'$-algèbres locales, qui sont aussi formellement lisses (19.3.5, (v)). On est
donc ramené à prouver que $A$ est \emph{régulier}. Soit $k_0$ le sous-corps
premier de $k$~; comme $k$ est une extension séparable de $k_0$, c'est une
$k_0$-algèbre formellement lisse (19.6.1) et en vertu de l'hypothèse, il en est
de même de $A$ (19.3.5, (ii)). Comme le corps résiduel $K$ de $A$ est une
extension séparable de $k_0$, on peut alors appliquer (19.6.4) à $A$ et $k_0$,
d'où la conclusion.
\end{proof}

\oldpage[0\textsubscript{IV-1}]{103}

\begin{corollary}[19.6.5.1]
\label{0.19.6.5.1-fr}
Soient $A$ un anneau local noethérien, $k$ un sous-corps de $A$ tel que $A$
soit une $k$-algèbre formellement lisse (pour sa topologie préadique). Alors
tout corps $K$ tel que $k\subset K\subset A$ est une extension séparable de
$k$.
\end{corollary}

\begin{proof}
En effet, pour toute extension finie $k'$ de $k$, l'anneau $K\otimes_k k'$ 
s'identifie à un sous-anneau de $A'=A\otimes_k k'$~; comme $A'$ est un anneau
régulier, il est réduit, donc il en est de même de $K$, ce qui prouve que $K$
est une extension séparable de $k$ (Bourbaki, \emph{Alg.}, chap.~VIII,
\textsection7, n\textsuperscript{o}~3, th.~1).
\end{proof}

On notera que cela démontre que la condition de l'énoncé de (19.6.1) est
\emph{nécessaire}.

\begin{remark}[19.6.5.2]
\label{0.19.6.5.2-fr}
Si $K$ est une extension \emph{non séparable} de $k$, l'anneau de séries
formelles $K[[T_1,\ldots,T_n]]$ \emph{n'est donc pas} une $k$-algèbre
formellement lisse (pour la structure de $k$-algèbre usuelle rappelée dans
(19.6.4)). D'autre part, il y a des $k$-algèbres formellement lisses $A$ qui
sont des anneaux locaux noethériens complets dont le corps résiduel $K$ est une
extension \emph{non séparable} de $k$ (par exemple les complétés des localisés
de $B=k[T_1,\ldots,T_n]$ en des idéaux maximaux $\mathfrak n$ tels que
$B/\mathfrak n$ soit une extension algébrique finie non séparable de $k$)~; un
tel anneau ne peut donc être $k$-isomorphe à $K[[T_1,\ldots,T_n]]$ bien qu'il
lui soit isomorphe en vertu de (19.6.5).

On peut dans certains cas préciser le corollaire (19.6.5)~:
\end{remark}

\begin{proposition}[19.6.6]
\label{0.19.6.6-fr}
Soient $k$ un corps, $A$ une $k$-algèbre locale noethérienne, $\mathfrak m$ son
idéal maximal, $A$ étant muni de la topologie $\mathfrak m$-préadique.
Supposons que le corps résiduel $K=A/\mathfrak m$ de $A$ soit tel qu'il existe
une extension radicielle finie $k_0$ de $k$ pour laquelle le corps résiduel de
l'anneau local $K\otimes_k k_0$ soit une extension séparable de $k_0$ (nous
exprimerons plus tard cette propriété en disant que $K$ est de
\emph{multiplicité radicielle finie} sur $k$ ($\mathbf{IV}$, 4.7.4) et nous
verrons que cette condition est toujours satisfaite lorsque $K$ est une
\emph{extension de type fini} de $k$). Alors les conditions suivantes sont
équivalentes~:
\begin{enumerate}
  \item[a)] $A$ est une $k$-algèbre formellement lisse.
  \item[b)] Il existe une extension radicielle finie $k'$ de $k$ telle que
  $\widehat A\otimes_k k'$ soit $k'$-isomorphe à une algèbre de séries
  formelles $K'[[T_1,\ldots,T_n]]$, où $K'$ est une extension séparable de
  $k'$.
  \item[b')] Il existe une extension finie $k'$ de $k$ telle que l'anneau
  semi-local $\widehat A\otimes_k k'$ ait au moins un anneau local composant
  qui soit $k'$-isomorphe à une algèbre de séries formelles
  $K'[[T_1,\ldots,T_n]]$, où $K'$ est une extension séparable de $k'$.
  \item[c)] $A$ est géométriquement régulier sur $k$ (19.6.5).
\end{enumerate}
\end{proposition}

\begin{proof}
Notons d'abord que si $k'$ est une extension radicielle de $k$, il n'y a qu'un
seul idéal de $A'=A\otimes_k k'$ au-dessus de $\mathfrak m$, formé des éléments
dont une puissance $p^h$-ième ($p$ exposant caractéristique de $k$) est dans
$\mathfrak m$ pour un $h$ convenable (Bourbaki, \emph{Alg. comm.}, chap.~V,
\textsection2, n\textsuperscript{o}~3, lemme~4)~; $A'$ est donc un anneau
local, et il en est de même de
$K\otimes_k k'=(A\otimes_k k')/(\mathfrak m\otimes_k k')$~; en outre les corps
résiduels de ces deux anneaux sont identiques. Rappelons d'autre part que si
$K$ est une extension séparable de $k$, alors, pour toute extension finie
$k''$ de $k$, $K\otimes_k k''$ est composé direct de corps (Bourbaki,
\emph{Alg.}, chap.~VIII, \textsection7, n\textsuperscript{o}~3, cor.~1 du
th.~1), et par suite $\mathfrak m\otimes_k k''$ est le radical de
$A\otimes_k k''$, et les corps composants de $K\otimes_k k''$ sont les corps
résiduels aux idéaux maximaux de $A\otimes_k k''$~; en

\oldpage[0\textsubscript{IV-1}]{104}

outre, ces corps sont des extensions séparables de $k''$ puisque pour toute
extension $k_1$ de $k''$,
$K\otimes_k k_1=(K\otimes_k k'')\otimes_{k''}k_1$ est sans radical
(\emph{loc. cit.}, th.~1). Si l'on applique ces remarques au corps $K$ de
multiplicité radicielle finie sur $k$, on voit que pour \emph{toute} extension
finie $k''$ de $k_0$, $(K\otimes_k k'')_{\mathrm{red}}$ est séparable sur
$k''$.

Passons à la démonstration de (19.6.6). Il est trivial que $b)$ implique
$b')$. Montrons que $b')$ entraîne $a)$. Si l'on pose
$B'=K'[[T_1,\ldots,T_n]]$, l'hypothèse que $K'$ est séparable sur $k'$
entraîne, par les remarques du début, que pour toute extension finie $k''$ de
$k'$, les composants de l'anneau semi-local complet $B'\otimes_{k'}k''$
(égal à l'anneau de séries formelles
$(K'\otimes_{k'}k'')[[T_1,\ldots,T_n]]$) sont les anneaux
$K_i''[[T_1,\ldots,T_n]]$, où les $K_i''$ sont les corps composants de
$K'\otimes_{k'}k''$. Comme les composants locaux de $B'\otimes_{k'}k''$ sont
aussi des composants locaux de
\[
  \widehat A\otimes_k k''
  =(\widehat A\otimes_k k')\otimes_{k'}k'',
\]
on voit que l'hypothèse $b')$ est encore vérifiée lorsqu'on remplace $k'$ par
une quelconque de ses extensions finies. En particulier, on peut supposer que
$k'$ est \emph{quasi-galoisienne}, donc extension galoisienne d'une extension
radicielle $k'_0$ de $k$~; $A'=\widehat A\otimes_k k'$ peut alors s'écrire
$A'_0\otimes_{k'_0}k'$, où $A'_0=\widehat A\otimes_k k'_0$ est un anneau
local complet. On sait alors (Bourbaki, \emph{Alg.}, chap.~VIII,
\textsection8, prop.~4) que $A'$ est composé direct d'anneaux locaux
\emph{isomorphes} à $A'_0$ en tant que $k'_0$-algèbres. Or, un de ces anneaux
est par hypothèse $k'$-isomorphe à un anneau de séries formelles
$K'[[T_1,\ldots,T_n]]$, où $K'$ est une extension séparable de $k'$, donc
\emph{aussi} une extension séparable de $k'_0$~; compte tenu de (19.6.4),
$A'_0$ est donc une $k'_0$-algèbre formellement lisse. Mais comme $k'_0$ est
un $k$-module fidèlement plat, projectif et de type fini, on conclut de
(19.4.7) que $\widehat A$ est une $k$-algèbre formellement lisse.

On a déjà vu (19.6.5) que $a)$ implique $c)$~; il reste donc à vérifier que
$c)$ implique $b)$. Or, si $c)$ est vérifiée, alors, pour toute extension
radicielle finie $k'$ de $k$ contenant $k_0$, $A\otimes_k k'$ est un anneau
local régulier, dont le corps résiduel $K'$ est une extension séparable de
$k'$~; il résulte donc de (19.6.4) que son complété
$\widehat A\otimes_k k'$ est $k'$-isomorphe à un anneau de séries formelles
$K'[[T_1,\ldots,T_n]]$.
\end{proof}

\begin{remarks}[19.6.7]
\label{0.19.6.7-fr}
\begin{enumerate}
  \item[\rm(i)] On notera que l'hypothèse que $K$ est de multiplicité
  radicielle finie sur $k$ n'a été utilisée que dans la démonstration de
  l'implication $b')\Rightarrow a)$.
  \item[\rm(ii)] Nous prouverons plus tard (22.5.8) que $a)$ et $c)$ sont
  équivalentes, \emph{sans hypothèse} sur l'extension $K$ de $k$.
\end{enumerate}
\end{remarks}

\subsection{Cas des homomorphismes locaux; théorèmes d'existence et
d'unicité}
\label{subsection:0.19.7-fr}

Dans ce numéro, lorsqu'un anneau semi-local est considéré comme un anneau
topologique, il est toujours sous-entendu qu'il s'agit de sa topologie
$\mathfrak r$-préadique, où $\mathfrak r$ est son radical. Tout homomorphisme
local d'anneaux locaux est donc automatiquement \emph{continu}.

\begin{theorem}[19.7.1]
\label{0.19.7.1-fr}
Soient $A$, $B$ deux anneaux locaux noethériens, $\mathfrak m$, $\mathfrak n$
leurs idéaux maximaux respectifs, $k=A/\mathfrak m$ le corps résiduel de $A$~;
on suppose $A$ et $B$ munis respectivement des topologies
$\mathfrak m$-préadique et $\mathfrak n$-préadique. Soit
$\varphi:A\to B$ un homomorphisme local, et posons $B_0=B\otimes_A k$.
Les propriétés suivantes sont équivalentes~:
\begin{enumerate}
  \item[a)] $B$ est une $A$-algèbre \emph{formellement lisse}.

\oldpage[0\textsubscript{IV-1}]{105}

  \item[b)] $B$ est un $A$-module \emph{plat} et
  $B_0=B/\mathfrak mB$ (muni de la topologie quotient) est une $k$-algèbre
  \emph{formellement lisse}.
\end{enumerate}
\end{theorem}

\begin{proof}
La démonstration se fait en plusieurs étapes.

\medskip
\noindent\phantomsection\label{0.19.7.1.1-fr}\textbf{(19.7.1.1)}\enspace
Démontrons d'abord que $b)$ entraîne $a)$~; nous allons appliquer le critère
(19.4.7), avec $\mathfrak J=\mathfrak m$~; en vertu de la seconde hypothèse
dans $b)$, tout revient à montrer que $B$ est un $A$-module
\emph{formellement projectif}. L'hypothèse entraîne que pour tout $h>0$,
$B/\mathfrak m^hB$ est un $(A/\mathfrak m^h)$-module plat
($\mathbf 0_{\mathrm{III}}$, 10.2.1)~; comme les $\mathfrak m^h$ forment un
système fondamental de voisinages de 0 dans $A$ et que
$(B/\mathfrak m^hB)\otimes_{A/\mathfrak m^h}k=B_0$, on peut remplacer $A$ et
$B$ par $A/\mathfrak m^h$ et $B/\mathfrak m^hB$ respectivement, et par suite
supposer $A$ \emph{artinien} (donc discret). Comme $B_0$ est une $k$-algèbre
formellement lisse, c'est un anneau \emph{régulier} (19.6.5)~; soit
$(x_i^0)_{1\leq i\leq n}$ un système régulier de paramètres pour $B_0$
(17.1.6), et pour tout $i$, soit $x_i\in B$ tel que $x_i^0$ soit son image
dans $B_0=B/\mathfrak mB$~; comme les $x_i^0$ engendrent l'idéal maximal
$\mathfrak n_0=\mathfrak n/\mathfrak mB$ de $B_0$, les idéaux
\[
  \mathfrak Q_m=\sum_{i=1}^n(x_i^0)^mB_0 \qquad (m>0)
\]
forment un système fondamental de voisinages de 0 dans $B_0$, car
$\mathfrak Q_m$ est évidemment contenu dans $\mathfrak n_0^m$, et d'autre
part contient $\mathfrak n_0^{mn}$. Posons
$\mathfrak J_m=\sum_{i=1}^n x_i^mB$ pour tout $m>0$~; il est clair que
$\mathfrak n=\mathfrak J_1+\mathfrak mB$~; comme il existe un $h>0$ tel que
$\mathfrak m^h=0$, on a
$\mathfrak n^m\subset\mathfrak J_1^{m-h}\subset\mathfrak n^{m-h}$ pour
$m>h$, et comme on a vu que $\mathfrak J_m\supset\mathfrak J_1^{mn}$, on voit
que les $\mathfrak J_m$ forment un système fondamental de voisinages de 0 dans
$B$. Tout revient par suite à prouver que les $B/\mathfrak J_m$ sont des
$A$-modules \emph{libres}, et il revient au même de voir que ce sont des
$A$-modules \emph{plats} ($\mathbf 0_{\mathrm{III}}$, 10.1.3). Or,
l'hypothèse que $(x_i^0)$ est une suite $B_0$-régulière d'éléments de l'idéal
maximal de $B_0$ entraîne la même propriété pour la suite des $(x_i^0)^m$
($1\leq i\leq n$) pour tout $m>0$ (15.1.20)~; la conclusion résulte donc de
(15.1.16, $b)$ et $c)$).
\end{proof}

\begin{lemma}[19.7.1.2]
\label{0.19.7.1.2-fr}
Soient $A$ un anneau topologique, $B$, $C$ deux $A$-algèbres topologiques qui
sont des anneaux locaux noethériens. On suppose en outre que $C$ est complet
et que le corps résiduel $B/\mathfrak m$ de $B$ est un $A$-module de type fini.
Soit $E$ le produit tensoriel complété $B\widehat\otimes_A C$. Alors~:
\begin{enumerate}
  \item[(i)] $E$ est un anneau semi-local noethérien complet.
  \item[(ii)] L'idéal $\mathfrak mE$ est contenu dans le radical de $E$, et
  pour tout $h>0$, $E/\mathfrak m^hE$ est isomorphe à
  $(B/\mathfrak m^h)\otimes_A C$.
  \item[(iii)] Si $C$ est un $A$-module plat, $E$ est un $B$-module plat.
\end{enumerate}
\end{lemma}

\begin{proof}
Par définition, $E$ est le séparé complété du produit tensoriel $B\otimes_A C$
pour la topologie définie par les idéaux
$\operatorname{Im}(\mathfrak m^i\otimes_A C)+
\operatorname{Im}(B\otimes_A\mathfrak n'{}^j)$
($\mathbf 0_{\mathrm I}$, 7.7.5). Si l'on pose
\[
  \mathfrak r=\operatorname{Im}(\mathfrak m\otimes_A C)+
  \operatorname{Im}(B\otimes_A\mathfrak n'),
\]
on a
$\mathfrak r^{2i}\subset
\operatorname{Im}(\mathfrak m^i\otimes_A C)+
\operatorname{Im}(B\otimes_A\mathfrak n'{}^i)\subset\mathfrak r^i$, donc $E$
est aussi le séparé complété de $B\otimes_A C$ pour la topologie
$\mathfrak r$-préadique. Par hypothèse,
$(B\otimes_A C)/\mathfrak r=(B/\mathfrak m)\otimes_A(C/\mathfrak n')$ est un
$(C/\mathfrak n')$-module de type fini, donc un anneau artinien~; en outre,
$\mathfrak r/\mathfrak r^2$, étant un quotient de
\[
  ((\mathfrak m/\mathfrak m^2)\otimes_A C)
  \oplus(B\otimes_A(\mathfrak n'/\mathfrak n'{}^2)),
\]
est un $(B\otimes_A C)$-module de type fini~; appliquant
($\mathbf 0_{\mathrm I}$, 7.2.11), on voit donc que $E$ est un anneau
\emph{noethérien}~; en outre, $E/\mathfrak rE$, isomorphe à
$(B\otimes_A C)/\mathfrak r$, étant artinien, $E$ est \emph{semi-local}.
Notons maintenant que $E$, qui est isomorphe à
\[
  \varprojlim_{i,j}((B/\mathfrak m^i)\otimes_A(C/\mathfrak n'{}^j)),
\]
est aussi isomorphe, par le théorème de la double limite projective, à
\[
  \varprojlim_i\left(\varprojlim_j
  ((B/\mathfrak m^i)\otimes_A(C/\mathfrak n'{}^j))\right)~;
\]

\oldpage[0\textsubscript{IV-1}]{106}

mais $\varprojlim_j((B/\mathfrak m^i)\otimes_A(C/\mathfrak n'{}^j))$ est le
séparé complété de $(B/\mathfrak m^i)\otimes_A C$, et comme $C$ est complet,
ce n'est autre que $(B/\mathfrak m^i)\otimes_A C$ lui-même,
$B/\mathfrak m^i$ étant un $A$-module de type fini puisque
$\mathfrak m/\mathfrak m^i$ est un $(B/\mathfrak m)$-module de type fini
($\mathbf 0_{\mathrm I}$, 7.3.6). On a donc
\[
  E=\varprojlim_i((B/\mathfrak m^i)\otimes_A C).
\]
Pour tout entier $h>0$, $\mathfrak m^hE$, étant un idéal de $E$, est fermé
dans $E$ ($\mathbf 0_{\mathrm I}$, 7.3.5), donc complet, et d'autre part il est
évidemment dense dans
\[
  \varprojlim_i
  (\operatorname{Im}(\mathfrak m^h/\mathfrak m^{h+i})\otimes_A C),
\]
donc égal à cette dernière limite projective. En outre, tous les systèmes
projectifs considérés sont définis par des homomorphismes \emph{surjectifs}~;
il résulte donc de ($\mathbf 0_{\mathrm{III}}$, 13.2.2) que
$E/\mathfrak m^hE$ est isomorphe à
\[
  (B/\mathfrak m^h)\otimes_A C
  =(B\otimes_A C)/\operatorname{Im}(\mathfrak m^h\otimes_A C).
\]
En particulier, comme
$\operatorname{Im}(\mathfrak m\otimes_A C)\subset\mathfrak r$, cela montre
que $\mathfrak mE$ est contenu dans $\mathfrak rE$, donc dans le radical de
$E$. Enfin, l'hypothèse que $C$ est un $A$-module plat entraîne que
$(B/\mathfrak m^h)\otimes_A C=E/\mathfrak m^hE$ est un
$(B/\mathfrak m^h)$-module plat pour tout $h>0$~: comme $B$ et $E$ sont
noethériens et que $\mathfrak mE$ est contenu dans le radical de $E$, il
résulte de ($\mathbf 0_{\mathrm{III}}$, 10.2.2) que $E$ est un $B$-module
plat.
\end{proof}

\begin{lemma}[19.7.1.3]
\label{0.19.7.1.3-fr}
Soient $A$ un anneau local noethérien, $\mathfrak m$ son idéal maximal, $k$ son
corps résiduel, $B_0$ une $k$-algèbre~; on suppose que $B_0$ est un anneau
local noethérien, complet et régulier. Alors il existe une $A$-algèbre
topologique $B$ qui est un anneau local noethérien complet, un $A$-module plat,
et tel que $B_0$ soit $k$-isomorphe à
$B\otimes_A k=B/\mathfrak mB$.
\end{lemma}

\begin{proof}
Comme $\widehat A$ est plat sur $A$ et a même corps résiduel, on peut se borner
au cas où $A$ est \emph{complet}.

Soit $K$ le corps résiduel de $B_0$, et distinguons deux cas~:

\medskip
\noindent I) $K$ est une extension \emph{séparable} de $k$. En vertu de
(19.6.4), $B_0$ est $k$-isomorphe à un anneau de séries formelles
$K[[T_1,\ldots,T_n]]$. Lorsque $B_0=K$, le lemme a déjà été démontré
($\mathbf 0_{\mathrm{III}}$, 10.3.1)~; soit $C$ un anneau local noethérien
complet qui est un $A$-module plat et tel que $C\otimes_A k$ soit isomorphe à
$K$. Pour $n\geq1$, il suffit de prendre (avec la notation précédente)
$B=C[[T_1,\ldots,T_n]]$~; on sait en effet (Bourbaki,
\emph{Alg. comm.}, chap.~III, §~3, n\textsuperscript{o}~4, cor.~3 du th.~1)
que $B$ est un $C$-module plat, donc aussi un $A$-module plat, et d'autre part,
il est immédiat que $C[[T_1,\ldots,T_n]]\otimes_A k$ est isomorphe à
$(C/\mathfrak mC)[[T_1,\ldots,T_n]]=B_0$.

\medskip
\noindent II) $K$ est de \emph{caractéristique} $p>0$, et par suite il en est
de même de $k$. Notons $P$ le corps premier $\mathbf F_p$, et $W(P)$ l'anneau
local complet des nombres $p$-adiques $\mathbf Z_p$, qui est un anneau de
valuation discrète (donc \emph{régulier}), et a $P$ pour corps résiduel.
Montrons d'abord qu'il existe un homomorphisme continu d'anneaux $W(P)\to A$
faisant donc de $A$ une $W(P)$-\emph{algèbre topologique}. En effet, si
$j:\mathbf Z\to A$ est l'homomorphisme canonique, on a
$j(p\mathbf Z)\subset\mathfrak m$ par hypothèse, d'où
$j^{-1}(\mathfrak m)=p\mathbf Z$, et par suite $j$ se factorise en
$\mathbf Z\to\mathbf Z_{p\mathbf Z}\xrightarrow{f}A$, où $f$ est un
homomorphisme local, donc continu, qui (puisque $A$ est complet) se prolonge
par continuité en l'homomorphisme $W(P)\to A$ cherché.

Comme $k$ est extension séparable de $P$, le cas I) montre qu'il y a un
homomorphisme local $W(P)\to W(k)$, où $W(k)$ est un anneau local noethérien
complet et un $W(P)$-module plat, tel que $W(k)\otimes_{W(P)}P$ soit isomorphe
à $k$. D'ailleurs, comme l'uniformisante $p$ de $W(P)$ est un élément
$W(k)$-régulier par platitude ($\mathbf 0_{\mathrm I}$, 6.3.4) et comme

\oldpage[0\textsubscript{IV-1}]{107}

$W(k)/pW(k)=k$, $pW(k)$ est l'idéal maximal de $W(k)$, ce qui entraîne que ce
dernier anneau est un \emph{anneau de valuation discrète complet} (Bourbaki,
\emph{Alg. comm.}, chap.~VI, §~3, n\textsuperscript{o}~5, prop.~9). Par
(19.7.1.1) on voit en outre (puisque $k$ est séparable sur $P$, donc une
$P$-algèbre formellement lisse (19.6.1)) que $W(k)$ est une $W(P)$-algèbre
\emph{formellement lisse}. Le $W(P)$-homomorphisme continu $W(k)\to k$ se
factorise donc en $W(k)\to A\to k$ (19.3.11), ce qui permet de considérer $A$
comme une $W(k)$-algèbre topologique. Appliquant maintenant le cas I) à $B_0$
considéré comme $P$-algèbre et à $W(P)$, on voit qu'il existe une
$W(P)$-algèbre $B_P$ qui est un anneau local noethérien complet, un
$W(P)$-module plat, et telle que $B_P\otimes_{W(P)}P$ soit $P$-isomorphe à
$B_0$. Utilisant de nouveau le fait que $W(k)$ est une $W(P)$-algèbre
formellement lisse, on voit par (19.3.11) que l'homomorphisme composé
$W(k)\to k\to B_0$ se factorise en $W(k)\to B_P\to B_0$~; en outre, comme
$k=W(k)/pW(k)$, on a
$B_P\otimes_{W(k)}k=B_P/pB_P=B_P\otimes_{W(P)}P=B_0$. Montrons que $B_P$ est
un $W(k)$-module \emph{plat}~; comme $W(k)$ est un anneau de valuation discrète
dont $p$ est l'uniformisante, il suffit de vérifier que $B_P$ est un
$W(k)$-module \emph{sans torsion} ($\mathbf 0_{\mathrm I}$, 6.3.4), ou encore
que $p$ est un élément $B_P$-régulier, ce qui résulte de ce que $B_P$ est un
$W(P)$-module plat ($\mathbf 0_{\mathrm I}$, 6.3.4).

Posons maintenant
\[
  B=B_P\widehat\otimes_{W(k)}A
\]
et notons que le corps résiduel de $A$ étant égal à celui de $W(k)$, est
\emph{a fortiori} un $W(k)$-module de type fini. Il résulte donc tout d'abord
de (19.7.1.2) que $B$ est un anneau semi-local noethérien complet,
$\mathfrak mB$ étant contenu dans le radical de $B$~; en outre
$B/\mathfrak mB$ est $k$-isomorphe à
$B_P\otimes_{W(k)}k=B_0$, donc $B$ est en fait un anneau local. Comme $B_P$
est un $W(k)$-module plat, (19.7.1.2) montre enfin que $B$ est un $A$-module
plat. C.Q.F.D.
\end{proof}

\begin{lemma}[19.7.1.4]
\label{0.19.7.1.4-fr}
Soient $A$ un anneau, $\mathfrak J$ un idéal de $A$, $M$, $N$ deux $A$-modules
séparés pour la topologie $\mathfrak J$-préadique. On suppose en outre que $N$
est complet pour la topologie $\mathfrak J$-préadique et que $M$ soit un
$A$-module plat. Soit $u:N\to M$ un $A$-homomorphisme~; si
$u\otimes1:N\otimes_A(A/\mathfrak J)\to M\otimes_A(A/\mathfrak J)$ est
bijectif, alors $u$ est bijectif.
\end{lemma}

\begin{proof}
Les modules gradués associés étant pris relatifs aux filtrations
$\mathfrak J$-préadiques, il résulte des hypothèses sur $M$ et $N$ relativement
aux topologies $\mathfrak J$-préadiques qu'il suffit de prouver que
$\gr(u):\gr_\bullet(N)\to\gr_\bullet(M)$ est \emph{bijectif} (Bourbaki,
\emph{Alg. comm.}, chap.~III, §~2, n\textsuperscript{o}~8, cor.~3 du th.~1).
Or, on a un diagramme commutatif
\[
\xymatrix{
  \gr_0(N)\otimes_{A_0}\gr_\bullet(A)
    \ar[r]^-{\gr_0(u)\otimes1}\ar[d]_{\varphi_N} &
  \gr_0(M)\otimes_{A_0}\gr_\bullet(A)\ar[d]^{\varphi_M} \\
  \gr_\bullet(N)\ar[r]_-{\gr(u)} & \gr_\bullet(M)
}
\]

\oldpage[0\textsubscript{IV-1}]{108}

où $A_0=A/\mathfrak J=\gr_0(A)$, et $\varphi_M$ et $\varphi_N$ sont les
applications canoniques ($\mathbf 0_{\mathrm{III}}$, 10.1.1.2). Par
hypothèse, $\gr_0(u)$ est bijectif ainsi que $\varphi_M$
($\mathbf 0_{\mathrm{III}}$, 10.2.1), et $\varphi_N$ est surjectif~; on en
déduit d'abord que $\gr_0(u)\otimes1$ est bijectif, puis que $\varphi_N$ est
injectif, donc bijectif, et enfin que $\gr(u)$ est bijectif.
\end{proof}

\begin{lemma}[19.7.1.5]
\label{0.19.7.1.5-fr}
Soient $A$ un anneau noethérien, $\mathfrak J$ un idéal de $A$, $B$, $B'$ deux
$A$-algèbres qui sont des anneaux locaux noethériens, les homomorphismes
$A\to B$, $A\to B'$ étant continus pour la topologie $\mathfrak J$-préadique
sur $A$. On suppose que~:
\begin{enumerate}
  \item[$1^{\mathrm o}$] $B$ et $B'$ sont complets pour les topologies
  $\mathfrak J$-préadiques~;
  \item[$2^{\mathrm o}$] $B$ est une $A$-algèbre formellement lisse~;
  \item[$3^{\mathrm o}$] $B'$ est un $A$-module plat.
\end{enumerate}
Posons $A_0=A/\mathfrak J$, et soit
$u_0:B\otimes_A A_0\to B'\otimes_A A_0$ un $A_0$-isomorphisme~; alors il
existe un $A$-isomorphisme $u:B\to B'$ tel que $u_0=u\otimes1$ (ce qui
entraîne que $B'$ est une $A$-algèbre formellement lisse et $B$ un $A$-module
plat).
\end{lemma}

\begin{proof}
Posons $B_0=B\otimes_A A_0$, $B'_0=B'\otimes_A A_0$. Notons que si
$\mathfrak m$ et $\mathfrak m'$ sont les idéaux maximaux de $B$ et $B'$, les
topologies $\mathfrak J$-préadiques sur $B$ et $B'$ sont séparées puisque
$\mathfrak JB\subset\mathfrak m$ et $\mathfrak JB'\subset\mathfrak m'$~; en
outre, comme $\mathfrak JB'$ est fermé dans $B'$ pour la topologie
$\mathfrak J$-préadique, l'homomorphisme composé
$B\to B_0\xrightarrow{u_0}B'_0$, qui est continu pour les topologies
$\mathfrak J$-préadiques, se factorise en $B\xrightarrow{u}B'\to B'_0$, où
$u$ est un $A$-homomorphisme continu (19.3.10). On a évidemment
$u_0=u\otimes1$, et l'hypothèse que $u_0$ est bijectif entraîne qu'il en est de
même de $u$ en vertu de (19.7.1.4).
\end{proof}

\medskip
\noindent\textbf{(19.7.1.6)}\enspace
\emph{Fin de la démonstration.} --- Pour achever de prouver (19.7.1), il faut
montrer que $a)$ entraîne $b)$~; on sait déjà que $a)$ entraîne que $B_0$ est
une $k$-algèbre formellement lisse (19.3.5, (iii)), donc tout revient à prouver
que $B$ est un $A$-module plat. Il revient au même d'établir que $\widehat B$
est un $\widehat A$-module plat (Bourbaki, \emph{Alg. comm.}, chap.~III, §~5,
n\textsuperscript{o}~4, prop.~4), et l'on sait que $\widehat B$ est une
$\widehat A$-algèbre formellement lisse (19.3.6)~; on peut donc se borner au
cas où $A$ et $B$ sont complets. Comme $B_0$ est une $k$-algèbre formellement
lisse, c'est un anneau régulier (19.6.5) et complet
($\mathbf 0_{\mathrm I}$, 6.3.5)~; appliquant (19.7.1.3), on voit qu'il existe
une $A$-algèbre $B'$ qui est un anneau local noethérien complet et un
$A$-module plat, un homomorphisme local $A\to B'$ et un $k$-isomorphisme
$B_0\simeq B'_0=B'\otimes_A k$. Il suffit alors d'appliquer (19.7.1.5) en
prenant pour $\mathfrak J$ l'idéal maximal de $A$, pour obtenir que $B$ est
$A$-isomorphe à $B'$, donc est un $A$-module plat. C.Q.F.D.

\begin{theorem}[19.7.2]
\label{0.19.7.2-fr}
Soient $A$ un anneau local noethérien, $\mathfrak J$ un idéal contenu dans
l'idéal maximal de $A$, $A_0=A/\mathfrak J$, $B_0$ un anneau local noethérien
complet, $A_0\to B_0$ un homomorphisme local faisant de $B_0$ une
$A_0$-algèbre formellement lisse. Alors il existe un anneau local noethérien
complet $B$, un homomorphisme local $A\to B$ faisant de $B$ un $A$-module
plat, et un $A_0$-isomorphisme
$u:B\otimes_A A_0\xrightarrow{\sim}B_0$. Si $(B',u')$ est un couple
satisfaisant aux mêmes conditions que $(B,u)$, il existe un $A$-isomorphisme
$v:B\xrightarrow{\sim}B'$ rendant commutatif le diagramme
\[
\xymatrix{
  B\otimes_A A_0\ar[rr]^{v\otimes1}_{\sim}\ar[dr]^{\sim}_u &&
  B'\otimes_A A_0\ar[dl]^{u'}_{\sim} \\
  & B_0
}
\]
\end{theorem}

\oldpage[0\textsubscript{IV-1}]{109}

\begin{proof}
Soit $\mathfrak m$ l'idéal maximal de $A$, de sorte que
$\mathfrak m_0=\mathfrak m/\mathfrak J$ est l'idéal maximal de $A_0$, $A$ et
$A_0$ ayant le même corps résiduel $k$. Posons
$B_{00}=B_0\otimes_{A_0}k$~; comme $B_{00}$ est une $k$-algèbre formellement
lisse (19.3.5, (iii)), c'est un anneau local régulier (19.6.5)~; en appliquant
(19.7.1.3), on voit qu'il existe une $A$-algèbre topologique $B$ qui est un
anneau local noethérien complet, un $A$-module plat, et pour laquelle on a un
$k$-isomorphisme
$u_{00}:B\otimes_A k\xrightarrow{\sim}B_{00}$. Notons qu'en vertu de (19.7.1),
$B$ est une $A$-algèbre formellement lisse, donc
$B\otimes_A A_0=B/\mathfrak JB$ est une $A_0$-algèbre formellement lisse
(19.3.5, (iii)) et un anneau local noethérien complet~; en outre, $B_0$ est un
$A_0$-module plat en vertu de l'hypothèse et de (19.7.1)~; comme on a un
$k$-isomorphisme
\[
  u_{00}:B\otimes_A k=(B\otimes_A A_0)\otimes_{A_0}k
  \xrightarrow{\sim}B_0\otimes_{A_0}k=B_{00},
\]
on déduit de (19.7.1.5), appliqué en y remplaçant $A$ par $A_0$ et
$\mathfrak J$ par $\mathfrak m_0$, qu'il existe un $A_0$-isomorphisme
$u:B\otimes_A A_0\xrightarrow{\sim}B_0$ tel que $u_{00}=u\otimes1$. Quant à
l'assertion d'unicité, notons que les idéaux $\mathfrak J^hB$ (resp.
$\mathfrak J^hB'$) sont fermés dans $B$ (resp. $B'$)
($\mathbf 0_{\mathrm I}$, 7.3.5), donc $B$ et $B'$ sont séparés et complets
pour les topologies $\mathfrak J$-préadiques (Bourbaki, \emph{Top. gén.},
chap.~III, 3\textsuperscript{e} éd., §~3, n\textsuperscript{o}~5, cor.~2 de la
prop.~9)~; on a par hypothèse un $A_0$-isomorphisme
$v_0:B\otimes_A A_0\xrightarrow{\sim}B'\otimes_A A_0$ tel que
$u'\circ v_0=u$~; comme $B$ est une $A$-algèbre formellement lisse et $B'$ un
$A$-module plat, on peut appliquer (19.7.1.5), d'où l'existence du
$A$-isomorphisme $v$ répondant à la question.
\end{proof}

\begin{remarks}[19.7.3]
\label{0.19.7.3-fr}
\begin{enumerate}
  \item[(i)] On notera que l'assertion d'unicité dans (19.7.2) est encore
  valable si l'on suppose seulement que $B$ et $B'$ sont complets \emph{pour
  les topologies $\mathfrak J$-préadiques}. Nous ignorons si l'on peut
  améliorer de même l'assertion d'existence, autrement dit si l'on peut se
  dispenser de supposer l'anneau local $B_0$ \emph{complet} (pour sa topologie
  $\mathfrak n_0$-préadique, en désignant par $\mathfrak n_0$ son idéal
  maximal) en exigeant seulement que $B$ soit complet \emph{pour la topologie
  $\mathfrak J$-préadique}. Lorsque $A_0$ est complet pour la topologie
  $\mathfrak m$-préadique, on peut voir que ce problème revient au suivant~:
  si $B_0$ est un anneau local noethérien régulier (non nécessairement
  complet) contenant le corps premier $\mathbf F_p=\mathbf Z/p\mathbf Z$,
  existe-t-il pour tout $n>1$ une $(\mathbf Z/p^n\mathbf Z)$-algèbre plate $B$
  telle que $B/pB$ soit isomorphe à $B_0$~?
  \item[(ii)] On notera qu'en général, l'isomorphisme $v$ dont l'existence est
  affirmée dans (19.7.2) \emph{n'est pas unique} (cf.~(19.8.7)).
\end{enumerate}
\end{remarks}

\subsection{Algèbres de Cohen et $p$-anneaux de Cohen; application à la
structure des anneaux locaux complets}
\label{subsection:0.19.8-fr}

Les résultats de cette section sont des applications immédiates des théorèmes
de (19.7), mais méritent d'être explicités en raison de leur importance
pratique.

\begin{definition}[19.8.1]
\label{0.19.8.1-fr}
Soient $A$, $B$ deux anneaux locaux noethériens, $\mathfrak m$ l'idéal maximal
de $A$, $k=A/\mathfrak m$ son corps résiduel,
$\varphi:A\to B$ un homomorphisme local, faisant de $B$ une $A$-algèbre. On
dit que $B$ est une \emph{$A$-algèbre de Cohen} si elle vérifie les conditions
suivantes~:
\begin{enumerate}
  \item[(i)] $B$ est un anneau \emph{complet}.
  \item[(ii)] $B$ est un $A$-module \emph{plat}.
  \item[(iii)] $B\otimes_A k$ est un \emph{corps} (autrement dit,
  $\mathfrak mB$ est l'idéal maximal de $B$) qui est une \emph{extension
  séparable} de $k$.
\end{enumerate}
\end{definition}

\oldpage[0\textsubscript{IV-1}]{110}

\begin{theorem}[19.8.2]
\label{0.19.8.2-fr}
Soient $A$ un anneau local noethérien, $k$ son corps résiduel.
\begin{enumerate}
  \item[(i)] Si $B$ est une $A$-algèbre de Cohen, $B$ est une $A$-algèbre
  formellement lisse. Pour tout anneau local noethérien complet $C$, tout
  homomorphisme local $A\to C$ et tout idéal $\mathfrak J\neq C$ dans $C$,
  tout $A$-homomorphisme $B\to C/\mathfrak J$ se factorise donc en
  $B\xrightarrow{v}C\to C/\mathfrak J$, où $v$ est un $A$-homomorphisme
  (nécessairement local).
  \item[(ii)] Pour tout corps $K$, extension séparable de $k$, il existe une
  $A$-algèbre de Cohen $B$ telle que $B\otimes_A k$ soit $k$-isomorphe à $K$,
  et une telle $A$-algèbre est unique à isomorphisme près.
\end{enumerate}
\end{theorem}

\begin{proof}
Comme $K$ est une $k$-algèbre formellement lisse (19.6.1), l'assertion (i)
résulte de (19.7.1). Pour prouver (ii), on peut se borner au cas où $A$ est
\emph{complet}, car il revient au même de dire que $B$ est un $A$-module plat
ou un $\widehat A$-module plat ($\mathbf 0_{\mathrm{III}}$, 10.2.3), on a
$\mathfrak mB=\mathfrak m\widehat AB$ et $k$ est le corps résiduel de
$\widehat A$. Il suffit alors d'appliquer (19.7.2) en prenant
$\mathfrak J=\mathfrak m$ et $B_0=K$ (et utilisant (19.6.1)).
\end{proof}

\begin{definition}[19.8.3]
\label{0.19.8.3-fr}
On appelle \emph{anneau local premier} un anneau local de la forme
$\mathbf Z_{p\mathbf Z}$, où $p\mathbf Z$ est un idéal premier de $\mathbf Z$.
On appelle \emph{anneau local complet premier} le complété d'un anneau local
premier.
\end{definition}

Les anneaux locaux premiers sont donc de deux sortes~:
\begin{enumerate}
  \item[$1^{\mathrm o}$] Ceux qui correspondent aux idéaux maximaux
  $p\mathbf Z$ où $p\neq0$ est un nombre premier~; $\mathbf Z_{p\mathbf Z}$
  est un anneau de valuation discrète, dont le complété est \emph{l'anneau des
  entiers $p$-adiques} noté d'ordinaire $\mathbf Z_p$\footnote{Cette notation,
  actuellement universellement utilisée, est dans ce cas en conflit avec la
  notation $A_f$ adoptée dans ($\mathbf 0_{\mathrm I}$, 1.2.3)~: avec
  $A=\mathbf Z$ et $f=p$, $A_f$ signifie en effet l'anneau des nombres
  rationnels de la forme $k/p^n$ ($k\in\mathbf Z$, $n$ entier $\geq0$)~; nous
  éviterons toujours d'utiliser la notation $\mathbf Z_p$ pour désigner ce
  dernier anneau.}.
  \item[$2^{\mathrm o}$] Pour l'idéal premier $p\mathbf Z=(0)$,
  $\mathbf Z_{p\mathbf Z}$ est le corps des nombres rationnels $\mathbf Q$,
  identique à son complété (la topologie étant naturellement la topologie
  d'anneau local noethérien, donc ici la topologie discrète).
\end{enumerate}

La terminologie de (19.8.3), analogue à celle des «~corps premiers~», se
justifie de la même manière~: pour tout anneau local $A$, considérons
l'homomorphisme canonique $\varphi:\mathbf Z\to A$, et soit $p\mathbf Z$
l'image réciproque par cet homomorphisme de l'idéal maximal $\mathfrak m$ de
$A$~; $p\mathbf Z$ est un idéal premier de $\mathbf Z$ et l'homomorphisme
précédent se factorise donc en
$\mathbf Z\to\mathbf Z_{p\mathbf Z}\xrightarrow{\psi}A$~; d'ailleurs, comme
$\varphi$ est l'\emph{unique} homomorphisme de $\mathbf Z$ dans $A$, $p$ et
$\psi$ sont déterminés de façon \emph{unique}. Autrement dit, pour tout anneau
local $A$, il y a un \emph{unique homomorphisme} $\psi:P\to A$, où $P$ est un
anneau local premier~; si en outre $A$ est séparé et complet, on peut prolonger
par complétion cet homomorphisme, et il y a donc un \emph{unique
homomorphisme} $\widehat\psi:\widehat P\to A$, où $\widehat P$ est un anneau
local complet premier. D'ailleurs, par passage aux quotients, $\psi$ donne un
homomorphisme du corps résiduel $\mathbf F_p$ si $p>0$ (resp. $\mathbf Q$ si
$p=0$) dans le corps résiduel $k$ de $A$, et $p$ est donc la
\emph{caractéristique} de $k$.

Si l'on prend en particulier pour $A$ un anneau local premier (resp. complet
premier), on voit qu'il n'existe dans un tel anneau qu'\emph{un seul
endomorphisme}, savoir l'identité.

% FIN SOURCE ega4/ega0-19-fr.tex

%
% DEBUT SOURCE ega4/ega0-19b-fr.tex
\oldpage[0\textsubscript{IV-1}]{111}

\begin{definition}[19.8.4]
\label{0.19.8.4-fr}
Soient $A$ un anneau local, $P\to A$ l'unique homomorphisme d'un anneau local
premier $P$ dans $A$, $p$ la caractéristique des corps résiduels de $P$ et
$A$. On dit que $A$ est un \emph{anneau de Cohen} si c'est une $P$-algèbre de
Cohen, c'est-à-dire (19.8.1) si~:
\begin{enumerate}
  \item[$1^{\mathrm o}$] $A$ est noethérien et complet.
  \item[$2^{\mathrm o}$] $A$ est un $P$-module plat (ce qui revient aussi à
  dire que $A$ est un $\widehat P$-module plat (Bourbaki, \emph{Alg. comm.},
  chap.~III, §~5, n\textsuperscript{o}~4, prop.~4)).
  \item[$3^{\mathrm o}$] $A/pA$ est un corps (nécessairement séparable sur le
  corps résiduel de $P$, ce corps étant premier).
\end{enumerate}
\end{definition}

Si $p=0$, ces conditions équivalent à dire que $A$ est un \emph{corps de
caractéristique $0$}. Si $p>0$, on a nécessairement $pA\neq0$~; la condition
$3^{\mathrm o}$ signifie que $pA$ est l'idéal maximal $\mathfrak m$ de $A$~;
la condition $2^{\mathrm o}$ signifie que $p$ est $A$-régulier, puisque $P$
est un anneau de valuation discrète ($\mathbf 0_{\mathrm I}$, 6.3.4). Donc $A$
est un anneau régulier (17.1.1, $d)$) de dimension $1$, et par suite un anneau
de valuation discrète, complet en vertu de $1^{\mathrm o}$~; en résumé~:

\begin{proposition}[19.8.5]
\label{0.19.8.5-fr}
Les anneaux de Cohen sont les corps de caractéristique $0$ et les anneaux de
valuation discrète complets, dont le corps résiduel a une caractéristique
$p>0$, et dont l'idéal maximal est engendré par $p\cdot1$ ($1$ étant l'unité de
l'anneau).
\end{proposition}

On notera que dans le second cas, $p\cdot1\neq0$ puisque $p$ est $A$-régulier,
donc on peut identifier $p\cdot1$ à l'entier $p$~; l'homomorphisme canonique
$\mathbf Z_p\to A$ est injectif, et on identifie $p\cdot1$ à l'élément $p$ de
$\mathbf Z_p$~; on dit dans ce cas que $A$ est un \emph{$p$-anneau de Cohen}.

\begin{theorem}[19.8.6]
\label{0.19.8.6-fr}
\textup{(Cohen).} ---
\begin{enumerate}
  \item[(i)] Soient $W$ un anneau de Cohen, $C$ un anneau local noethérien
  complet, $\mathfrak J$ un idéal de $C$ distinct de $C$. Alors tout
  homomorphisme local $u:W\to C/\mathfrak J$ se factorise en
  $W\xrightarrow{v}C\to C/\mathfrak J$, où $v$ est un homomorphisme local.
  \item[(ii)] Soit $K$ un corps. Il existe un anneau de Cohen $W$ dont le corps
  résiduel est isomorphe à $K$. Si $W'$ est un second anneau de Cohen, $K'$ son
  corps résiduel, tout isomorphisme $u:K\xrightarrow{\sim}K'$ provient par
  passage aux quotients d'un isomorphisme $v:W\xrightarrow{\sim}W'$.
\end{enumerate}
\end{theorem}

\begin{proof}
Cela n'est autre que (19.8.2) appliqué au cas où $A$ est un anneau local
premier.
\end{proof}

\begin{remarks}[19.8.7]
\label{0.19.8.7-fr}
\begin{enumerate}
  \item[(i)] Lorsque $K$ est de caractéristique $0$, la partie (ii) de (19.8.6)
  devient triviale.
  \item[(ii)] L'homomorphisme $v$ de (19.8.6, (i)) \emph{n'est pas
  nécessairement déterminé de façon unique par $u$}, comme le montre déjà le
  cas où $W$ est un corps de caractéristique $0$, $\mathfrak J^2=0$ et $u$ est
  un isomorphisme (cf.~(21.5.5)). De même, dans (19.8.6, (ii)) l'isomorphisme
  $v$ \emph{n'est pas nécessairement déterminé de façon unique par $u$}
  (cf.~(21.5.5)).

  Toutefois, lorsque $K$ est \emph{parfait} et de caractéristique $p>0$, on
  verra (21.5.5) que dans (19.8.6, (ii)) l'isomorphisme $v$ est unique. On verra
  aussi plus tard que dans ce cas $W$ s'identifie à l'anneau $W_\infty(K)$ des
  \emph{vecteurs de Witt de longueur infinie} sur $K$.
  \item[(iii)] Dans (19.8.6, (i)), on peut affaiblir les hypothèses sur $C$ en
  utilisant (19.3.10) et (19.3.12).
\end{enumerate}
\end{remarks}

\begin{theorem}[19.8.8]
\label{0.19.8.8-fr}
\textup{(Cohen).} --- Soient $A$ un anneau local noethérien complet, $k$ son
corps résiduel.

\oldpage[0\textsubscript{IV-1}]{112}

\begin{enumerate}
  \item[(i)] Il existe un anneau de Cohen $W$ tel que $A$ soit isomorphe à un
  anneau quotient d'un anneau de séries formelles
  $W[[T_1,\ldots,T_n]]$ (et en particulier $A$ est isomorphe à un quotient d'un
  anneau local régulier complet (17.3.8)). Si $A$ contient un corps, il est
  isomorphe à un anneau quotient de $k[[T_1,\ldots,T_n]]$.
  \item[(ii)] Supposons en outre $A$ intègre. Alors il existe un sous-anneau
  $B$ de $A$, tel que~: $1^{\mathrm o}$ $B$ est isomorphe à un anneau de séries
  formelles sur un anneau $C$ qui est un corps ou un anneau de Cohen (ce qui
  entraîne que $B$ est un anneau local régulier et complet (17.3.8))~;
  $2^{\mathrm o}$ $B$ a même corps résiduel que $A$ et l'injection $B\to A$ est
  un homomorphisme local~; $3^{\mathrm o}$ $A$ est une $B$-algèbre finie.
\end{enumerate}
\end{theorem}

\begin{proof}
Soit $\mathfrak m$ l'idéal maximal de $A$. Il existe un anneau de Cohen $W$ dont
le corps résiduel est isomorphe à $k$ (19.8.6, (ii))~; on a donc un
homomorphisme local $W\to A/\mathfrak m$, qui par suite se factorise en
$W\xrightarrow{u}A\to A/\mathfrak m$, où $u$ est un homomorphisme local
(19.8.6, (i)). Pour toute famille finie $(x_i)_{1\leq i\leq n}$ d'éléments de
$\mathfrak m$, il existe alors un homomorphisme local
$v:W[[T_1,\ldots,T_n]]\to A$ prolongeant $u$ et tel que $v(T_i)=x_i$ pour tout
$i$ (Bourbaki, \emph{Alg. comm.}, chap.~III, §~4, n\textsuperscript{o}~5,
prop.~6). Lorsque $A$ contient un corps, il contient un corps premier $P$, dont
$k$ est une extension (nécessairement séparable), et par suite $A$ contient un
corps isomorphe à $k$ (19.6.2)~; on peut alors remplacer $W$ par $k$ dans la
définition précédente de $v$.

\medskip
\noindent(i)\enspace
Prenons d'abord pour les $x_i$ un système de générateurs de $\mathfrak m$. Comme
$W$ a même corps résiduel que $A$, et que les classes des $x_i$ dans l'anneau
gradué $\gr_\bullet(A)$ engendrent $\gr_\bullet(A)$ en tant que $k$-algèbre,
$\gr(v):\gr_\bullet(W[[T_1,\ldots,T_n]])\to\gr_\bullet(A)$ est surjectif~; on
en déduit que $v$ est lui-même surjectif (Bourbaki, \emph{Alg. comm.},
chap.~III, §~2, n\textsuperscript{o}~8, cor.~2 du th.~1). Rappelons que le cas
où $A$ contient un corps a déjà été vu et ne figure ici que pour mémoire
(19.6.3).

\medskip
\noindent(ii)\enspace
Si $A$ contient un corps, il contient un corps $k'$ isomorphe à $k$ comme on
l'a vu~; on considère alors un système de paramètres
$(y_j)_{1\leq j\leq m}$ de $A$ (16.3.6), on prend
$B=k[[T_1,\ldots,T_m]]$ et on considère l'homomorphisme local $w:B\to A$ qui
coïncide dans $k$ avec un isomorphisme $k\to k'$ et qui est tel que
$w(T_j)=y_j$ pour $1\leq j\leq m$. Si $A$ ne contient pas de corps, l'unique
homomorphisme $\mathbf Z_{p\mathbf Z}\to A$ (19.8.3) est nécessairement
injectif (sans quoi, comme $A$ est intègre, son noyau serait l'idéal maximal de
$\mathbf Z_{p\mathbf Z}$ et son image isomorphe à un corps)~; en outre, on a
alors $p>0$ par hypothèse, $\mathbf Z_{p\mathbf Z}$ étant un corps si $p=0$.
L'élément $p\cdot1$ de $A$ (identifié à $p$) est non diviseur de zéro dans $A$,
et est contenu dans $\mathfrak m$, donc (16.3.4 et 16.3.7) il existe une famille
$(z_i)_{1\leq i\leq m-1}$ qui, avec $p$, forme un système de paramètres de $A$.
L'anneau de Cohen $W$ considéré au début de la démonstration est alors un
anneau de valuation discrète de corps résiduel $k$, dans lequel $p$ engendre
l'idéal maximal (19.8.5), et l'homomorphisme unique $u:W\to A$ défini au début
applique $p$ sur lui-même. On prend alors $B=W[[T_1,\ldots,T_{m-1}]]$ et on
considère l'homomorphisme local $w:B\to A$ qui coïncide avec $u$ dans $W$ et
est tel que $w(T_i)=z_i$ pour $1\leq i\leq m-1$. Dans les deux cas, si
$\mathfrak n$ est l'idéal maximal de $B$, il est clair que $\mathfrak nA$ est
un idéal de définition de $A$~; comme en outre $B/\mathfrak n$ et
$A/\mathfrak m$ sont isomorphes, $A$ est un $B$-module quasi-fini
($\mathbf 0_{\mathrm I}$, 7.4.4), donc un $A$-module de type fini puisque $B$
est complet et $A$ séparé pour les topologies $\mathfrak n$-préadiques
($\mathbf 0_{\mathrm I}$, 7.4.1). D'autre part,

\oldpage[0\textsubscript{IV-1}]{113}

dans les deux cas, on a $\dim(B)=\dim(A)=m$~; dans le premier cas, cela résulte
de (17.1.4, (iii))~; dans le second, on voit directement que $p$ et les $T_i$
($1\leq i\leq m-1$) forment une suite $B$-régulière engendrant $\mathfrak n$,
ou on peut aussi utiliser le fait que ces éléments engendrent $\mathfrak n$ et
que l'on a $\dim(B)\geq\dim(A)$ par (16.3.10). Comme $A$ et $B$ sont intègres,
on tire finalement de (16.3.10) que $w$ est injectif, ce qui achève la
démonstration.
\end{proof}

\begin{corollary}[19.8.9]
\label{0.19.8.9-fr}
Soit $A$ un anneau local noethérien intègre complet contenant un corps $k_0$~;
soit $k$ le corps résiduel de $A$, et supposons que $k$ soit fini sur $k_0$.
Alors, dans la conclusion de (19.8.8, (ii)), on peut remplacer $1^{\mathrm o}$
et $2^{\mathrm o}$ par la condition que $B$ est de la forme
$k_0[[T_1,\ldots,T_n]]$, l'injection canonique $B\to A$ étant un
$k_0$-homomorphisme local (pour la structure usuelle de $k_0$-algèbre de $B$).
\end{corollary}

\begin{proof}
En effet, en reprenant la démonstration de (19.8.8, (ii)), on définit cette fois
$w:k_0[[T_1,\ldots,T_n]]\to A$ comme coïncidant dans $k_0$ avec l'identité et
appliquant $T_j$ sur $y_j$ pour $1\leq j\leq n$. L'hypothèse que $k$ est de
degré fini sur $k_0$ entraîne encore que $A$ est un $B$-module quasi-fini, donc
de type fini par ($\mathbf 0_{\mathrm I}$, 7.4.1), et on conclut comme dans
(19.8.8).
\end{proof}

\begin{corollary}[19.8.10]
\label{0.19.8.10-fr}
Soit $A$ un anneau local artinien dont l'idéal maximal $\mathfrak m$ est de
carré nul~; il existe alors un anneau local noethérien régulier $B$, d'idéal
maximal $\mathfrak n$, tel que $A$ soit isomorphe à $B/\mathfrak n^2$.
\end{corollary}

\begin{proof}
Soient $K$ le corps résiduel $A/\mathfrak m$ de $A$, $n$ le rang de
$\mathfrak m/\mathfrak m^2=\mathfrak m$ sur $K$. Si $A$ contient un corps, il
résulte de (19.6.3) que $A$ est isomorphe à $B/\mathfrak b$, où
$B=K[[T_1,\ldots,T_n]]$ et $\mathfrak b$ est contenu dans le carré de l'idéal
maximal $\mathfrak n$ de $B$~; mais comme
$\operatorname{long}(B/\mathfrak n^2)=n+1=\operatorname{long}(A)$, on a
nécessairement $\mathfrak b=\mathfrak n^2$.

Supposons ensuite que $A$ ne contienne pas de corps~; cela entraîne que $K$ est
de caractéristique $p>0$ et que $p\cdot1\neq0$ dans $A$ (19.6.3)~; donc
$p\cdot1$ est un élément de $\mathfrak m$, et il y a par suite $n-1$ autres
éléments $x_i$ ($2\leq i\leq n$) de $\mathfrak m$ formant avec $p\cdot1$ une
base de $\mathfrak m$ sur $K$. Soit $W$ un anneau de Cohen dont le corps
résiduel est isomorphe à $K$~; $W$ est un anneau de valuation discrète dont $p$
engendre l'idéal maximal~; on a vu dans la démonstration de (19.8.8) qu'il y a
un homomorphisme $u:W\to A$ appliquant $p$ sur lui-même et qui par passage aux
quotients donne l'identité sur $K$. On prend $B=W[[T_2,\ldots,T_n]]$ et l'on
considère l'homomorphisme local $w:B\to A$ qui coïncide avec $u$ dans $W$ et
est tel que $w(T_i)=x_i$ pour $2\leq i\leq n$. Il est clair que $w$ est
surjectif et que son noyau $\mathfrak b$ est contenu dans le carré de l'idéal
maximal $\mathfrak n=pB+BT_2+\cdots+BT_n$ de $B$~; comme
$\operatorname{long}(B/\mathfrak n^2)=n+1=\operatorname{long}(A)$, on a encore
$\mathfrak b=\mathfrak n^2$.
\end{proof}

\begin{proposition}[19.8.11]
\label{0.19.8.11-fr}
Soient $A$ un anneau local artinien, $\mathfrak m$ son idéal maximal, $k$ son
corps résiduel. Pour que $A$ soit isomorphe à un anneau quotient d'un anneau de
Cohen, il faut et il suffit que $\mathfrak m$ soit engendré par $p\cdot1$, où
$p$ est la caractéristique de $k$.
\end{proposition}

\begin{proof}
La condition est évidemment nécessaire (19.8.5). Pour voir qu'elle est
suffisante, on observe, comme au début de la démonstration de (19.8.8), qu'il
existe un anneau de Cohen $W$ dont le corps résiduel est isomorphe à $k$ et un
homomorphisme local $u:W\to A$. En outre, si l'on considère l'homomorphisme
composé $\mathbf Z_{p\mathbf Z}\to W\xrightarrow{u}A$ (qui est nécessairement
l'unique homomorphisme de $\mathbf Z_{p\mathbf Z}$ dans $A$), on voit que
l'image

\oldpage[0\textsubscript{IV-1}]{114}

par $u$ de l'élément $p\cdot1$ de $W$ est l'élément $p\cdot1$ de $A$~; comme
l'élément $p\cdot1$ de $W$ engendre l'idéal maximal de cet anneau, on déduit
aussitôt de l'hypothèse que
$\gr(u):\gr_\bullet(W)\to\gr_\bullet(A)$ est surjectif, et par suite il en est
de même de $u$ (Bourbaki, \emph{Alg. comm.}, chap.~III, §~2,
n\textsuperscript{o}~8, cor.~2 du th.~1).
\end{proof}

\subsection{Algèbres relativement formellement lisses}
\label{subsection:0.19.9-fr}

\begin{definition}[19.9.1]
\label{0.19.9.1-fr}
Soient $\Lambda$ un anneau topologique, $A$ une $\Lambda$-algèbre topologique,
$B$ une $A$-algèbre topologique. On dit que $B$ est une $A$-algèbre
\emph{formellement lisse relativement à $\Lambda$} si, pour toute $A$-algèbre
topologique discrète $C$, et tout idéal nilpotent $\mathfrak J$ de $C$, tout
$A$-homomorphisme continu $u_0:B\to C/\mathfrak J$ qui se factorise en
$B\xrightarrow{u}C\to C/\mathfrak J$, où $u$ est un $\Lambda$-homomorphisme
continu, se factorise aussi en $B\xrightarrow{v}C\to C/\mathfrak J$, où $v$
est un $A$-homomorphisme continu.
\end{definition}

Il résulte de cette définition que si $B$ est une $A$-algèbre \emph{formellement
lisse}, alors $B$ est aussi \emph{formellement lisse relativement à
$\Lambda$}, pour toute structure de $\Lambda$-algèbre topologique définie sur
$A$ (autrement dit, tout homomorphisme continu d'anneaux $\Lambda\to A$).

\begin{proposition}[19.9.2]
\label{0.19.9.2-fr}
Soient $\Lambda$ un anneau topologique, $A$ une $\Lambda$-algèbre topologique.
\begin{enumerate}
  \item[(i)] $A$ est une $\Lambda$-algèbre formellement lisse relativement à
  $\Lambda$.
  \item[(ii)] Si $B$ est une $A$-algèbre formellement lisse relativement à
  $\Lambda$ et $C$ une $B$-algèbre formellement lisse relativement à
  $\Lambda$, alors $C$ est une $A$-algèbre formellement lisse relativement à
  $\Lambda$.
  \item[(iii)] Soient $B$ une $A$-algèbre formellement lisse relativement à
  $\Lambda$, $A'$ une $A$-algèbre topologique, alors la $A'$-algèbre
  topologique $B\otimes_A A'$ est formellement lisse relativement à $\Lambda$.
  \item[(iv)] Soient $B$ une $A$-algèbre topologique, $S$ (resp. $T$) une
  partie multiplicative de $A$ (resp. $B$) telle que l'image canonique de $S$
  dans $B$ soit contenue dans $T$. Si $B$ est une $A$-algèbre formellement
  lisse relativement à $\Lambda$, alors $T^{-1}B$ est une $S^{-1}A$-algèbre
  formellement lisse relativement à $\Lambda$.
  \item[(v)] Soient $B_i$ ($1\leq i\leq n$) des $A$-algèbres topologiques.
  Pour que $\prod_{i=1}^n B_i$ soit une $A$-algèbre formellement lisse
  relativement à $\Lambda$, il faut et il suffit que chacune des $B_i$ le
  soit.
\end{enumerate}
\end{proposition}

\begin{proof}
L'assertion (i) est triviale, et la démonstration des autres est étroitement
calquée sur les démonstrations de (19.3.5)~; elle est donc laissée au lecteur.
\end{proof}

\begin{corollary}[19.9.3]
\label{0.19.9.3-fr}
Soient $A$ un anneau topologique, $A'$ et $B$ deux $A$-algèbres topologiques.
Alors la $A$-algèbre topologique $B'=B\otimes_A A'$ est une $B$-algèbre
formellement lisse relativement à $A$.
\end{corollary}

\begin{proof}
Cela résulte de (19.9.2, (i) et (iii)).
\end{proof}

\begin{proposition}[19.9.4]
\label{0.19.9.4-fr}
Soient $\Lambda$ un anneau topologique, $A$ une $\Lambda$-algèbre topologique,
$B$ une $A$-algèbre topologique. Les conditions suivantes sont équivalentes~:
\begin{enumerate}
  \item[a)] $B$ est une $A$-algèbre formellement lisse relativement à
  $\Lambda$.
  \item[b)] $\widehat B$ est une $A$-algèbre formellement lisse relativement à
  $\Lambda$.
  \item[c)] $\widehat B$ est une $\widehat A$-algèbre formellement lisse
  relativement à $\Lambda$.
  \item[d)] $\widehat B$ est une $\widehat A$-algèbre formellement lisse
  relativement à $\widehat\Lambda$.
\end{enumerate}
\end{proposition}

\begin{proof}
On laisse encore au lecteur la démonstration, calquée sur celle de (19.3.6).
\end{proof}

\oldpage[0\textsubscript{IV-1}]{115}

\begin{env}[19.9.5]
\label{0.19.9.5-fr}
De même, l'énoncé (19.3.8) est encore valable (avec la même démonstration) quand
on y remplace «~formellement lisse~» par «~formellement lisse relativement à
$\Lambda$~». Si dans l'énoncé de (19.3.10) on remplace «~formellement lisse~»
par «~formellement lisse relativement à $\Lambda$~» la conclusion est remplacée
par la suivante (la démonstration restant essentiellement inchangée)~: \emph{tout
$A$-homomorphisme continu $u_0:B\to C/\mathfrak J$ qui se factorise en
$B\xrightarrow{u}C\to C/\mathfrak J$, où $u$ est un $\Lambda$-homomorphisme
continu, se factorise aussi en $B\xrightarrow{v}C\to C/\mathfrak J$, où $v$
est un $A$-homomorphisme continu}.
\end{env}

\begin{env}[19.9.6]
\label{0.19.9.6-fr}
Les critères de lissité formelle (19.4.1) et (19.4.2) sont valables lorsqu'on y
remplace «~formellement lisse~» par «~formellement lisse relativement à
$\Lambda$~», les démonstrations restant pratiquement inchangées.
\end{env}

\begin{proposition}[19.9.7]
\label{0.19.9.7-fr}
Soient $\Lambda$ un anneau topologique, $A$ une $\Lambda$-algèbre topologique,
$B$ une $A$-algèbre topologique. Supposons que pour toute $A$-algèbre discrète
$C$ et tout idéal $\mathfrak J$ de $C$ tel que $\mathfrak J^2=0$, tout
$A$-homomorphisme continu $u_0:B\to C/\mathfrak J$ qui se factorise en
$B\xrightarrow{u}C\to C/\mathfrak J$, où $u$ est un $\Lambda$-homomorphisme
continu, se factorise aussi en $B\xrightarrow{v}C\to C/\mathfrak J$, où $v$
est un $A$-homomorphisme continu. Alors $B$ est une $A$-algèbre formellement
lisse relativement à $\Lambda$.
\end{proposition}

\begin{proof}
La démonstration de (19.4.3) se transcrit aussitôt.
\end{proof}

\begin{proposition}[19.9.8]
\label{0.19.9.8-fr}
Soient $\Lambda$ un anneau topologique, $A$ une $\Lambda$-algèbre topologique,
$B$ une $A$-algèbre topologique. Pour que $B$ soit une $A$-algèbre
formellement lisse relativement à $\Lambda$, il faut et il suffit que pour tout
$B$-module topologique discret $L$, annulé par un idéal ouvert de $B$, on ait
(cf.~18.4.2)
\[
  \operatorname{Exalcotop}_{A/\Lambda}(B,L)=0.
\]
\end{proposition}

\begin{proof}
Avec les notations de la démonstration de (19.4.4), il suffit de noter ici que
l'on peut supposer que l'extension $E_\Lambda$ de $B_\Lambda$ est
$\Lambda$-triviale~; le reste de la démonstration est alors inchangé.
\end{proof}

Lorsque $\Lambda$, $A$ et $B$ sont des anneaux discrets, le critère (19.9.8) se
réduit à
\begin{equation}
\tag{19.9.8.1}
\label{0.19.9.8.1-fr}
  \operatorname{Exalcom}_{A/\Lambda}(B,L)=0
  \quad\text{pour tout $B$-module $L$}~;
\end{equation}
autrement dit, \emph{toute $A$-extension commutative de $B$ par un $B$-module,
qui est $\Lambda$-triviale, est aussi $A$-triviale}.

\subsection{Algèbres formellement non ramifiées et algèbres formellement
étales}
\label{subsection:0.19.10-fr}
\label{0.19.10-fr}

\begin{env}[19.10.1]
\label{0.19.10.1-fr}
Nous aurons aussi à utiliser deux notions voisines de la notion d'algèbre
formellement lisse.
\end{env}

\begin{definition}[19.10.2]
\label{0.19.10.2-fr}
Soient $A$ un anneau topologique, $B$ une $A$-algèbre topologique. On dit que
$B$ est une $A$-algèbre \emph{formellement non ramifiée} si, pour tout
$A$-homomorphisme $p:E\to C$ de $A$-algèbres discrètes, dont le noyau est
nilpotent, et tout $A$-homomorphisme continu $u:B\to C$, il existe \emph{au plus
un} $A$-homomorphisme continu $v:B\to E$ tel que $u$ se factorise en
$B\xrightarrow{v}E\xrightarrow{p}C$. On dit que $B$ est une $A$-algèbre
\emph{formellement étale} si elle est formellement lisse et formellement non
ramifiée, autrement dit, si, pour tout $A$-homomorphisme surjectif $p:E\to C$
de $A$-algèbres discrètes, dont le noyau est nilpotent, et pour tout
$A$-homomorphisme continu $u:B\to C$, il existe \emph{un et un seul}
$A$-homomorphisme continu $v:B\to E$ tel que $u$ se factorise en
$B\xrightarrow{v}E\xrightarrow{p}C$.
\end{definition}

\oldpage[0\textsubscript{IV-1}]{116}

\begin{examples}[19.10.3]
\label{0.19.10.3-fr}
Nous nous bornerons aux anneaux \emph{discrets}.
\begin{enumerate}
  \item[(i)] Si l'homomorphisme structural $\rho:A\to B$ est \emph{surjectif},
  $B$ est une $A$-algèbre \emph{formellement non ramifiée}, car avec les
  notations de (19.10.2), l'homomorphisme composé
  $A\xrightarrow{j}B\xrightarrow{v}E$ doit être l'homomorphisme structural,
  donc $v$ est unique (s'il existe). Par contre, $B$ n'est formellement lisse
  que si $\rho$ est \emph{bijectif}.
  \item[(ii)] Soient $A$ un anneau, $S$ une partie multiplicative de $A$~;
  alors $B=S^{-1}A$ est une $A$-algèbre \emph{formellement étale}. En effet,
  soient $j:A\to B$ l'homomorphisme canonique, $\rho:A\to E$ un homomorphisme
  d'anneaux, $\mathfrak J$ un idéal nilpotent de $E$,
  $p:E\to C=E/\mathfrak J$ l'homomorphisme canonique. Si $u:B\to C$ est un
  $A$-homomorphisme, $u(s/1)$ est inversible dans $C$ pour $s\in S$ puisque
  $s/1$ est inversible dans $B$~; mais $u(s/1)=p(\rho(s))$, et comme
  $\mathfrak J$ est nilpotent, le fait que l'image de $\rho(s)$ dans
  $E/\mathfrak J$ soit inversible entraîne que $\rho(s)$ lui-même est
  inversible dans $E$ ($\mathbf 0_{\mathrm I}$, 7.1.12)~; on en conclut que
  $\rho$ se factorise d'une seule manière en
  $A\xrightarrow{j}B\xrightarrow{v}E$, et l'on a évidemment
  $u\circ j=(p\circ v)\circ j$, d'où $u=p\circ v$, ce qui prouve notre
  assertion.
  \item[(iii)] Une extension finie séparable $k'$ d'un corps $k$ est une
  $k$-algèbre formellement lisse (19.6.1)~; nous verrons plus loin que pour
  qu'une extension de \emph{type fini} $K$ d'un corps $k$ soit une $k$-algèbre
  formellement étale, il faut et il suffit qu'elle soit formellement non
  ramifiée, et cela équivaut aussi au fait que $K$ est une extension finie et
  séparable de $k$ (21.7.4).
  \item[(iv)] Une algèbre de polynômes $A[X_\alpha]_{\alpha\in I}$ sur un
  anneau $A$ est \emph{formellement lisse} sur $A$ (19.3.3) mais \emph{non
  formellement étale} si $I$ n'est pas vide, comme il résulte aussitôt des
  définitions.
\end{enumerate}
\end{examples}

\begin{remark}[19.10.4]
\label{0.19.10.4-fr}
Le même raisonnement que dans (19.4.3) montre que pour vérifier qu'une
$A$-algèbre topologique $B$ est formellement non ramifiée, on peut, dans la
définition (19.10.2), se borner au cas où le noyau de $p:E\to C$ est de
\emph{carré nul}~; on peut en outre se borner au cas où $u(B)\subset p(E)$
(sans quoi il n'y a aucun $v$ factorisant $u$) et, en remplaçant $C$ par $u(B)$
et $E$ par $p^{-1}(u(B))$, supposer $u$ et $p$ surjectifs.
\end{remark}

% FIN SOURCE ega4/ega0-19b-fr.tex

%
% DEBUT SOURCE ega4/ega0-20-fr.tex
\section{Dérivations et différentielles}
\label{section:0.20-fr}
\label{0.20-fr}

Les notions introduites dans ce paragraphe seront reprises sous forme
géométrique au chapitre~IV, §~16, et joueront un rôle important dans l'étude
des préschémas. Leur importance dans ce chapitre tient d'abord à leurs
relations avec la notion d'algèbre formellement lisse, et notamment aux
théorèmes (20.4.9), (20.5.7) et (20.5.12), qui seront traduits en langage
géométrique dans le paragraphe du chapitre~IV consacré aux morphismes lisses,
et sont d'un usage constant dans les applications. D'autre part, les notions
différentielles serviront à démontrer au §~22 d'importants critères de
régularité qui joueront un rôle essentiel dans l'étude approfondie des anneaux
locaux noethériens faite au §~7 du chapitre~IV.

\subsection{Dérivations et extensions d'algèbres}
\label{subsection:0.20.1-fr}

\oldpage[0\textsubscript{IV-1}]{117}
\begin{proposition}[20.1.1]
\label{0.20.1.1-fr}
Soient $A$ un anneau (non nécessairement commutatif), $B$, $E$, $C$ des
$A$-anneaux, $p:E\to C$ un $A$-homomorphisme dont le noyau $\mathfrak{Q}$
est de carré nul, $u:B\to C$ un $A$-homomorphisme. Supposons qu'il existe
un $A$-homomorphisme $v_0:B\to E$ tel que $u$ se factorise en
$B\xrightarrow{v_0}E\xrightarrow{p}C$. Alors l'ensemble des
$A$-homomorphismes $v:B\to E$ tels que $u=p\circ v$ est identique à
l'ensemble des applications $v_0+D$, où $D:B\to\mathfrak{Q}$ est une
application vérifiant les deux conditions suivantes :
\begin{enumerate}
  \item[\rm(i)] $D$ est un homomorphisme de $A$-bimodules.
  \item[\rm(ii)] Quels que soient $f$, $g$ dans $B$, on a
\end{enumerate}
\begin{equation}
\tag{20.1.1.1}
\label{0.20.1.1.1-fr}
  D(fg)=f.D(g)+D(f).g.
\end{equation}
\end{proposition}

\begin{proof}
Dire que $p(v(f))=p(v_0(f))$ pour $f\in B$ signifie que
$D(f)=v(f)-v_0(f)$ appartient à $\mathfrak{Q}$ ; si l'on écrit que
$v(fg)=v(f)v(g)$, on obtient la relation \sref{0.20.1.1.1-fr},
$\mathfrak{Q}$ étant de carré nul, et la condition (i) résulte de ce que $v$
et $v_0$ sont des $A$-homomorphismes.

Si $\rho:A\to B$ est l'homomorphisme structural, on tire de
\sref{0.20.1.1.1-fr} que l'on a
$D(\rho(a)f)=\rho(a)D(f)+D(\rho(a))f$ pour tout $a\in A$ ; mais on doit
avoir $D(\rho(a)f)=\rho(a)D(f)$ d'après (i), donc, en faisant $f=1$, il
vient
\[
  D(\rho(a))=0 \qquad\text{pour}\quad a\in A\,;
\]
réciproquement, si $D$ est nulle dans $\rho(A)$ et vérifie
\sref{0.20.1.1.1-fr}, elle vérifie aussi (i).
\end{proof}

\begin{definition}[20.1.2]
\label{0.20.1.2-fr}
Étant donnés un $A$-anneau $B$ et un $B$-bimodule $L$, on appelle
\emph{$A$-dérivation de $B$ dans $L$} une application $D:B\to L$ vérifiant
les conditions (i) et (ii) de \sref{0.20.1.1-fr}.
\end{definition}

Il résulte de \sref{0.20.1.1-fr} que le \emph{noyau} d'une $A$-dérivation
$D:B\to L$ est un sous-$A$-anneau de $B$.

On appelle parfois \emph{dérivation} de $B$ dans $L$ une
$\mathbf{Z}$-dérivation, c'est-à-dire une application \emph{additive} de $B$
dans $L$ vérifiant \sref{0.20.1.1.1-fr} ; on a donc $D(1)=0$ pour une telle
application. Si $B$ est une algèbre sur un \emph{corps premier} $P$, toute
$\mathbf{Z}$-dérivation de $B$ est une $P$-dérivation : c'est évident d'après
ce qui précède si $P$ est de caractéristique $>0$, et dans le cas contraire,
la relation $n.n^{-1}=1$ pour $n\in\mathbf{Z}$ donne
\[
  nD(n^{-1})=0
\]
donc $D(n^{-1})=0$ puisque $P$ est de caractéristique $0$.

Il résulte aussitôt de la définition \sref{0.20.1.2-fr} que si $D$, $D'$ sont
deux $A$-dérivations de $B$ dans $L$, il en est de même de $D-D'$. Autrement
dit, l'ensemble des $A$-dérivations de $B$ dans $L$ est muni d'une structure
de \emph{groupe additif} ; on note ce groupe
$\operatorname{D\acute{e}r}_A(B,L)$. Si $A$ est \emph{commutatif}, $B$ une
$A$-\emph{algèbre commutative}, et $L$ un $B$-module, alors, pour toute
$A$-dérivation $D$ de $B$ dans $L$ et tout $a\in A$, $aD$ est encore une
$A$-dérivation de $B$ dans $L$, autrement dit $\operatorname{D\acute{e}r}_A(B,L)$
est alors muni d'une structure de $A$-\emph{module}.

\oldpage[0\textsubscript{IV-1}]{118}
La prop. \sref{0.20.1.1-fr} s'interprète de la façon suivante :

\begin{corollary}[20.1.3]
\label{0.20.1.3-fr}
Étant donnés deux $A$-homomorphismes de $A$-anneaux $p:E\to C$,
$u:B\to C$ dont le premier a un noyau $\mathfrak{Q}$ de carré nul,
l'ensemble des $A$-homomorphismes $v:B\to E$ tels que $u=p\circ v$ est vide
ou est un espace homogène principal pour le groupe
$\operatorname{D\acute{e}r}_A(B,\mathfrak{Q})$.
\end{corollary}

En particulier :

\begin{corollary}[20.1.4]
\label{0.20.1.4-fr}
Soient $A$ un anneau, $C$ un $A$-anneau, $L$ un $C$-bimodule, $E$ une
$A$-extension de $C$ par $L$, $p:E\to C$ l'augmentation. L'application qui,
à toute dérivation $D\in\operatorname{D\acute{e}r}_A(C,L)$, associe l'application
$x\mapsto x+D(p(x))$ est un isomorphisme du groupe
$\operatorname{D\acute{e}r}_A(C,L)$ sur le groupe des $A$-équivalences de $E$ sur
elle-même.
\end{corollary}

\begin{proof}
Appliquons \sref{0.20.1.1-fr} pour $B=E$, $v_0=\mathrm{I}_E$ : l'ensemble
des $A$-homomorphismes $v:E\to E$ tels que $p\circ v=p$ est identique à
l'ensemble des applications $v=\mathrm{I}_E+D'$, où
$D'\in\operatorname{D\acute{e}r}_A(E,L)$. Dire qu'un tel $A$-homomorphisme $v$ est
une $A$-équivalence équivaut à dire que $v$ se réduit à l'injection canonique
dans $L$, ou encore que $D'(x)=0$ dans $L$ ; mais cela signifie aussi que
$D'$ se factorise en $E\xrightarrow{p}C\xrightarrow{D}L$, où $D$ est une
$A$-dérivation. D'où le corollaire.
\end{proof}

Pour les extensions triviales, \sref{0.20.1.3-fr} donne :

\begin{corollary}[20.1.5]
\label{0.20.1.5-fr}
Soient $A$ un anneau, $B$, $C$ deux $A$-anneaux, $L$ un $C$-bimodule,
$u:B\to C$ un $A$-homomorphisme ; l'application qui, à toute dérivation
$D\in\operatorname{D\acute{e}r}_A(B,L)$, associe l'application
$\alpha:x\mapsto(u(x),D(x))$ est une bijection sur l'ensemble des
$A$-homomorphismes $v:B\to D_C(L)$ (cf. 18.2.3) tels que $u$ se factorise en
$B\xrightarrow{v}D_C(L)\to C$.
\end{corollary}

Plus particulièrement :

\begin{corollary}[20.1.6]
\label{0.20.1.6-fr}
Soient $A$ un anneau, $B$ un $A$-anneau, $L$ un $B$-bimodule. Si, à toute
dérivation $D\in\operatorname{D\acute{e}r}_A(B,L)$, on associe : $1^\circ$ la
$A$-équivalence $(x,y)\mapsto(x,y+D(x))$ de l'extension $D_B(L)$ sur
elle-même ; $2^\circ$ le $A$-homomorphisme $x\mapsto(x,D(x))$ de $B$ dans
$D_B(L)$, inverse à droite de l'homomorphisme d'augmentation $D_B(L)\to B$,
on définit des correspondances biunivoques canoniques entre :
\begin{enumerate}
  \item[\rm(i)] l'ensemble $\operatorname{D\acute{e}r}_A(B,L)$ ;
  \item[\rm(ii)] l'ensemble des $A$-équivalences de $D_B(L)$ sur elle-même ;
  \item[\rm(iii)] l'ensemble des $A$-homomorphismes inverses à droite de
  l'homomorphisme d'augmentation $D_B(L)\to B$.
\end{enumerate}
\end{corollary}

On notera que la correspondance biunivoque ainsi établie entre (i) et (ii)
respecte les structures de groupe, et que celle qu'on en déduit entre (ii) et
(iii) n'est autre que la correspondance biunivoque déjà définie dans (18.3.8).

\begin{corollary}[20.1.7]
\label{0.20.1.7-fr}
Soient $A$ un anneau, $C$ un $A$-anneau, $L$ un $A$-bimodule, $E$ une
$A$-extension triviale de $C$ par $L$, $p:E\to C$ l'augmentation. L'ensemble
des $A$-dérivations $D$ de $E$ dans $L$ telles que $D(x)=x$ dans $L$ est
identique à l'ensemble des applications $x\mapsto x-w(p(x))$, où $w$
parcourt l'ensemble des $A$-homomorphismes inverses à droite de $p$.
\end{corollary}

\begin{proof}
Il suffit encore d'appliquer \sref{0.20.1.1-fr} pour $B=E$,
$v_0=\mathrm{I}_E$ ; si $v=\mathrm{I}_E-D$, la condition $D(x)=x$ pour
$x\in L$ équivaut à $v(x)=0$ pour $x\in L$, c'est-à-dire à $v=w\circ p$,
où $w:C\to E$ est un $A$-homomorphisme ; en outre la condition
$p\circ v=p$ équivaut à $p\circ w=\mathrm{I}_C$, autrement dit au fait que
$w$ est inverse à droite de $p$.
\end{proof}

\subsection{Propriétés fonctorielles des dérivations}
\label{subsection:0.20.2-fr}

\oldpage[0\textsubscript{IV-1}]{119}
\begin{env}[20.2.1]
\label{0.20.2.1-fr}
Soient $A$ un anneau, $B$ un $A$-anneau, $L$ un $B$-bimodule ; si $L'$ est
un second $B$-bimodule et $w:L\to L'$ un homomorphisme de $B$-bimodules, il
est clair que l'application $D\mapsto w\circ D$ est un homomorphisme de
groupes additifs
\begin{equation}
\tag{20.2.1.1}
\label{0.20.2.1.1-fr}
  w_0:\operatorname{D\acute{e}r}_A(B,L)\to\operatorname{D\acute{e}r}_A(B,L')
\end{equation}
et que si $w':L'\to L''$ est un second homomorphisme de $B$-bimodules, on a
$(w'\circ w)_0=w'_0\circ w_0$. Lorsque $A$ est commutatif, $B$ une
$A$-algèbre commutative et $L$ un $B$-module, \sref{0.20.2.1.1-fr} est un
homomorphisme de $A$-modules.

En second lieu, soient $B'$ un $A$-anneau, $v:B'\to B$ un
$A$-homomorphisme qui fait de $L$ un $B'$-bimodule ; alors l'application
$D\mapsto D\circ v$ est un homomorphisme de groupes additifs,
\begin{equation}
\tag{20.2.1.2}
\label{0.20.2.1.2-fr}
  v^0:\operatorname{D\acute{e}r}_A(B,L)\to\operatorname{D\acute{e}r}_A(B',L)
\end{equation}
comme il résulte de \sref{0.20.1.1.1-fr} ; si $v':B''\to B'$ est un second
$A$-homomorphisme, on a $(v\circ v')^0=v'^0\circ v^0$. Lorsque $A$, $B$ et
$B'$ sont commutatifs et $L$ un $B$-module, \sref{0.20.2.1.2-fr} est un
homomorphisme de $A$-modules.

Enfin, soit $u:A'\to A$ un homomorphisme d'anneaux faisant de $B$ un
$A'$-anneau ; toute $A$-dérivation $D\in\operatorname{D\acute{e}r}_A(B,L)$ est aussi
une $A'$-dérivation, d'où une injection canonique de groupes commutatifs
\begin{equation}
\tag{20.2.1.3}
\label{0.20.2.1.3-fr}
  u^0:\operatorname{D\acute{e}r}_A(B,L)\to\operatorname{D\acute{e}r}_{A'}(B,L)
\end{equation}
et si $u':A''\to A'$ est un second homomorphisme d'anneaux, on a
$(u\circ u')^0=u'^0\circ u^0$ ; lorsque $A$, $A'$ et $B$ sont commutatifs et
$L$ un $B$-module, \sref{0.20.2.1.3-fr} est un di-homomorphisme de modules
(relatif à $u$).

On peut encore dire que
\[
  (A,B,L)\leadsto\operatorname{D\acute{e}r}_A(B,L)
\]
est un \emph{foncteur covariant} de la catégorie $\boldsymbol{K}$ définie dans
(18.3.5) dans la catégorie $\boldsymbol{Ab}$ des groupes commutatifs, en
faisant correspondre, à tout triplet $(u,v,w)$ constituant un morphisme de
$\boldsymbol{K}$ l'homomorphisme $w_0\circ v^0\circ u^0$ ; la vérification de
la fonctorialité résulte de la commutativité des diagrammes
\[
\xymatrix{
  \operatorname{D\acute{e}r}_A(B,L) \ar[r]^{v^0} \ar[d]_{w_0} &
    \operatorname{D\acute{e}r}_A(B',L) \ar[d]^{w_0} \\
  \operatorname{D\acute{e}r}_A(B,L') \ar[r]^{v^0} &
    \operatorname{D\acute{e}r}_A(B',L')
}
\qquad
\xymatrix{
  \operatorname{D\acute{e}r}_A(B,L) \ar[r]^{u^0} \ar[d]_{w_0} &
    \operatorname{D\acute{e}r}_{A'}(B,L) \ar[d]^{w_0} \\
  \operatorname{D\acute{e}r}_A(B,L') \ar[r]^{u^0} &
    \operatorname{D\acute{e}r}_{A'}(B,L')
}
\]
pour tout homomorphisme $w:L\to L'$ de $B$-bimodules.
\end{env}

\oldpage[0\textsubscript{IV-1}]{120}
\begin{theorem}[20.2.2]
\label{0.20.2.2-fr}
Soient $u:A\to B$, $v:B\to C$ deux homomorphismes d'anneaux, $L$ un
$C$-bimodule. On a une suite exacte canonique de groupes commutatifs
\begin{align}
\tag{20.2.2.1}
\label{0.20.2.2.1-fr}
  0\to\operatorname{D\acute{e}r}_B(C,L)\xrightarrow{u^0}
  \operatorname{D\acute{e}r}_A(C,L)&\xrightarrow{v^0}
  \operatorname{D\acute{e}r}_A(B,L)\xrightarrow{\partial} \\
  &\to\operatorname{Exan}_B(C,L)\xrightarrow{u^1}
  \operatorname{Exan}_A(C,L)\xrightarrow{v^1}
  \operatorname{Exan}_A(B,L) \notag
\end{align}
où $u^0$, $v^0$ sont les homomorphismes (20.2.1.3) et (20.2.1.2)
respectivement, $u^1$, $v^1$ les homomorphismes définis dans (18.3.4.1) et
(18.3.3.1) respectivement, et où $\partial$ est défini comme suit : pour
toute $A$-dérivation $D$ de $B$ dans $L$, $\partial(D)$ est la classe de la
$B$-extension de $C$ par $L$ définie sur l'anneau $D_C(L)$ par le
$A$-homomorphisme $\alpha:x\mapsto(v(x),D(x))$ (cf.
\sref{0.20.1.5-fr}). En outre, la suite exacte \sref{0.20.2.2.1-fr} est
fonctorielle en $L$ (pour les homomorphismes définis dans (20.2.1.1) et
(18.3.1.1) respectivement).
\end{theorem}

\begin{proof}
Comme $D$ est une $A$-dérivation (et \emph{a fortiori} une
$\mathbf{Z}$-dérivation) de $B$ dans $L$, le $A$-homomorphisme
$\alpha:x\mapsto(v(x),D(x))$ définit bien sur $D_C(L)$ une structure de
$B$-extension, donc $\partial(D)$ est bien définie \sref{0.20.1.5-fr}. Il
faut vérifier l'exactitude en $5$ endroits :

\medskip
\noindent 1)\enspace
L'exactitude en $\operatorname{D\acute{e}r}_B(C,L)$ est triviale (cf.
\sref{0.20.2.1-fr}).

\medskip
\noindent 2)\enspace
Par définition \sref{0.20.2.1-fr} le noyau de $v^0$ est l'ensemble des
$A$-dérivations de $C$ dans $L$ qui s'annulent dans $v(B)$, c'est-à-dire
celles des $A$-dérivations qui sont aussi des $B$-dérivations
\sref{0.20.1.1-fr} ; d'où l'exactitude en $\operatorname{D\acute{e}r}_A(C,L)$.

\medskip
\noindent 3)\enspace
Le noyau de $\partial$ est formé des dérivations
$D\in\operatorname{D\acute{e}r}_A(B,L)$ pour lesquelles la $B$-extension définie par
$\alpha:x\mapsto(v(x),D(x))$ est $B$-triviale ; cela signifie (18.2.3) qu'il
existe un $B$-homomorphisme $z\mapsto(z,w(z))$ de $C$ dans $D_C(L)$ (la
structure de $B$-anneau de $D_C(L)$ étant définie par $\alpha$) ; mais un tel
homomorphisme, étant \emph{a fortiori} un $A$-homomorphisme, est de la forme
$z\mapsto(z,D'(z))$ où $D'\in\operatorname{D\acute{e}r}_A(C,L)$
\sref{0.20.1.6-fr} ; et en écrivant que c'est un $B$-homomorphisme, il vient
\[
  D(x)z+v(x)D'(z)=D'(v(x)z)=v(x)D'(z)+D'(v(x))z
\]
pour $x\in B$, $z\in C$, ce qui donne $D'(v(x))=D(x)$ ; le noyau de
$\partial$ est donc l'image de $v^0$.

\medskip
\noindent 4)\enspace
Le noyau de $u^1$ est l'ensemble des classes de $B$-extensions de $C$ par
$L$ qui sont $A$-triviales (18.3.7), donc (à équivalence près) de la forme
$D_C(L)$, où la structure de $A$-anneau est définie par l'homomorphisme
$t\mapsto(u(t),0)$. Or, toute structure de $B$-extension sur $D_C(L)$ est
définie par un homomorphisme $\alpha:x\mapsto(v(x),D(x))$, où $D$ est une
$\mathbf{Z}$-dérivation de $B$ dans $L$ \sref{0.20.1.5-fr} ; dire que la
structure de $A$-anneau de cette $B$-extension est déduite de sa structure de
$B$-anneau au moyen de $u:A\to B$ signifie que $D(u(t))=0$ pour $t\in A$,
donc que $D$ est une $A$-dérivation ou encore que la classe de la $B$-extension
considérée est de la forme $\partial(D)$ ; d'où l'exactitude en
$\operatorname{Exan}_B(C,L)$.

\medskip
\noindent 5)\enspace
Le noyau de $v^1$ est l'ensemble des classes des $A$-extensions $E$ de $C$
par $L$ qui se trivialisent sur $v(B)$, c'est-à-dire pour lesquelles il y a un
$A$-homomorphisme $w:B\to E$ tel que $v$ se factorise en
$B\xrightarrow{w}E\to C$ ; or un tel $A$-homomorphisme définit sur $E$ une
structure de $B$-extension dont la classe a pour image par $u^1$ la classe de
la $A$-extension donnée ; la réciproque étant triviale, on a prouvé
l'exactitude en $\operatorname{Exan}_A(C,L)$.

Enfin, la fonctorialité en $L$ découle trivialement des définitions.
\end{proof}

\oldpage[0\textsubscript{IV-1}]{121}
\begin{corollary}[20.2.3]
\label{0.20.2.3-fr}
Soient $A$, $B$, $C$ trois anneaux commutatifs, $u:A\to B$, $v:B\to C$
deux homomorphismes d'anneaux, $L$ un $C$-module. On a une suite exacte
canonique de $A$-modules
\begin{align}
\tag{20.2.3.1}
\label{0.20.2.3.1-fr}
  0\to\operatorname{D\acute{e}r}_B(C,L)\xrightarrow{u^0}
  \operatorname{D\acute{e}r}_A(C,L)&\xrightarrow{v^0}
  \operatorname{D\acute{e}r}_A(B,L)\xrightarrow{\partial} \\
  &\to\operatorname{Exalcom}_B(C,L)\xrightarrow{u^1}
  \operatorname{Exalcom}_A(C,L)\xrightarrow{v^1}
  \operatorname{Exalcom}_A(B,L) \notag
\end{align}
fonctorielle en $L$.
\end{corollary}

\begin{proof}
Le raisonnement est le même que dans \sref{0.20.2.2-fr} une fois qu'on a
vérifié que pour toute dérivation $D\in\operatorname{D\acute{e}r}_A(B,L)$,
$\partial(D)$ est bien la classe d'une $B$-extension de $C$ par $L$ qui est
une $B$-algèbre commutative ; mais cela résulte aussitôt de la commutativité
de $C$ et du fait que $L$ est un $C$-module.
\end{proof}

\begin{corollary}[20.2.4]
\label{0.20.2.4-fr}
Sous les hypothèses de \sref{0.20.2.2-fr} (resp.
\sref{0.20.2.3-fr}), on a une suite exacte canonique, fonctorielle en $L$,
\begin{equation}
\tag{20.2.4.1}
\label{0.20.2.4.1-fr}
  0\to\operatorname{D\acute{e}r}_B(C,L)\xrightarrow{u^0}
  \operatorname{D\acute{e}r}_A(C,L)\xrightarrow{v^0}
  \operatorname{D\acute{e}r}_A(B,L)\xrightarrow{\partial}
  \operatorname{Exan}_{B/A}(C,L)\to0
\end{equation}
(resp.
\begin{equation}
\tag{20.2.4.2}
\label{0.20.2.4.2-fr}
  0\to\operatorname{D\acute{e}r}_B(C,L)\xrightarrow{u^0}
  \operatorname{D\acute{e}r}_A(C,L)\xrightarrow{v^0}
  \operatorname{D\acute{e}r}_A(B,L)\xrightarrow{\partial}
  \operatorname{Exalcom}_{B/A}(C,L)\to0\text{).}
\end{equation}
Cela résulte de la définition de $\operatorname{Exan}_{B/A}(C,L)$ (resp.
$\operatorname{Exalcom}_{B/A}(C,L)$) ((18.3.7) et (18.4.2)).
\end{corollary}

\begin{remark}[20.2.5]
\label{0.20.2.5-fr}
Supposons que l'on ait un diagramme commutatif d'homomorphismes d'anneaux
\[
\xymatrix{
  A \ar[r] & B \ar[r] & C \\
  A' \ar[u] \ar[r] & B' \ar[u] \ar[r] & C' \ar[u]
}
\]
Alors on a un diagramme commutatif
\[
\xymatrix@C=0.5em{
  0 \ar[r] & \operatorname{D\acute{e}r}_B(C,L) \ar[r] \ar[d] &
  \operatorname{D\acute{e}r}_A(C,L) \ar[r] \ar[d] &
  \operatorname{D\acute{e}r}_A(B,L) \ar[r]^-\partial \ar[d] &
  \operatorname{Exan}_B(C,L) \ar[r] \ar[d] &
  \operatorname{Exan}_A(C,L) \ar[r] \ar[d] &
  \operatorname{Exan}_A(B,L) \ar[d] \\
  0 \ar[r] & \operatorname{D\acute{e}r}_{B'}(C',L) \ar[r] &
  \operatorname{D\acute{e}r}_{A'}(C',L) \ar[r] &
  \operatorname{D\acute{e}r}_{A'}(B',L) \ar[r]^-\partial &
  \operatorname{Exan}_{B'}(C',L) \ar[r] &
  \operatorname{Exan}_{A'}(C',L) \ar[r] &
  \operatorname{Exan}_{A'}(B',L)
}
\]
et de même pour les suites exactes \sref{0.20.2.3.1-fr},
\sref{0.20.2.4.1-fr} et \sref{0.20.2.4.2-fr}.
\end{remark}

\subsection{Dérivations continues dans les anneaux topologiques}
\label{subsection:0.20.3-fr}

\begin{env}[20.3.1]
\label{0.20.3.1-fr}
Étant donnés deux anneaux topologiques $A$, $B$ (linéairement topologisés
comme toujours), nous désignerons par $\operatorname{Hom.\ cont}(A,B)$
l'ensemble des homomorphismes continus de $A$ dans $B$. Étant donné un anneau
topologique $A$, la catégorie des $A$-anneaux topologiques se définit comme
celle des $A$-anneaux (18.1.4) en remplaçant partout « anneau » par « anneau
topologique » et « homomorphisme » par « homomorphisme continu » ; si $B$

\oldpage[0\textsubscript{IV-1}]{122}
et $C$ sont deux $A$-anneaux topologiques, on notera
$\operatorname{Hom.\ cont}_A(B,C)$ l'ensemble des $A$-homomorphismes continus
de $B$ dans $C$.

Soient $A$ un anneau topologique, $B$ un $A$-anneau topologique, $L$ un
$B$-bimodule topologique ; on note $\operatorname{D\acute{e}r.\ cont}_A(B,L)$
l'ensemble des $A$-dérivations \emph{continues} de $B$ dans $L$ ; il est clair
que c'est un \emph{sous-groupe} de $\operatorname{D\acute{e}r}_A(B,L)$ (et un
sous-$A$-module lorsque $A$ et $B$ sont commutatifs et $L$ un $B$-module).
\end{env}

\begin{env}[20.3.2]
\label{0.20.3.2-fr}
Il est immédiat que la prop. \sref{0.20.1.1-fr} subsiste quand on y remplace
« anneau » par « anneau topologique » et « homomorphisme » par
« homomorphisme continu ». De même, \sref{0.20.1.3-fr} et
\sref{0.20.1.4-fr} restent valables en remplaçant $\operatorname{D\acute{e}r}$ par
$\operatorname{D\acute{e}r.\ cont}$, « anneau » (resp. « bimodule ») par « anneau
topologique » (resp. « bimodule topologique »), « homomorphisme » par
« homomorphisme continu » ; il faut naturellement supposer dans
\sref{0.20.1.4-fr} que $p$ est continu, et remplacer « $A$-équivalences » par
« $A$-équivalences continues ». On a les mêmes résultats pour
\sref{0.20.1.5-fr} et \sref{0.20.1.6-fr}, à condition de prendre comme
topologie sur $D_C(L)$ (resp. $D_B(L)$) la topologie \emph{produit} (sur
$C\times L$, resp. $B\times L$).
\end{env}

\begin{proposition}[20.3.3]
\label{0.20.3.3-fr}
Soient $A$ un anneau topologique, $B$ un $A$-anneau topologique, $L$ un
$B$-bimodule topologique discret annulé par un idéal bilatère ouvert de $B$.
Si dans $B$ le carré de tout idéal bilatère ouvert est ouvert, on a
$\operatorname{D\acute{e}r.\ cont}_A(B,L)=\operatorname{D\acute{e}r}_A(B,L)$.
\end{proposition}

\begin{proof}
En effet, si $\mathfrak{R}$ est un idéal bilatère ouvert de $B$ annulant $L$,
et $D$ une $A$-dérivation de $B$ dans $L$, on a
$D(\mathfrak{R}^2)\subset\mathfrak{R}.D(\mathfrak{R})+
D(\mathfrak{R}).\mathfrak{R}=0$ \sref{0.20.1.1.1-fr}, donc $D$ est continue.
\end{proof}

\begin{env}[20.3.4]
\label{0.20.3.4-fr}
Tous les résultats de \sref{0.20.2.1-fr} restent valables lorsqu'on y remplace
« anneau » par « anneau topologique », « bimodule » par « bimodule
topologique », « homomorphisme » par « homomorphisme continu » et
$\operatorname{D\acute{e}r}$ par $\operatorname{D\acute{e}r.\ cont}$.
\end{env}

\begin{proposition}[20.3.5]
\label{0.20.3.5-fr}
Soient $A$ un anneau topologique, $B$ un $A$-anneau topologique, $L$ un
$B$-bimodule topologique discret annulé par un idéal bilatère ouvert de $B$.
On a alors un isomorphisme canonique
\begin{equation}
\tag{20.3.5.1}
\label{0.20.3.5.1-fr}
  \varinjlim\operatorname{D\acute{e}r}_{A/\mathfrak{J}}
  (B/\mathfrak{R},L)\xrightarrow{\sim}
  \operatorname{D\acute{e}r.\ cont}_A(B,L)
\end{equation}
où dans le premier membre la limite inductive est prise suivant l'ensemble
ordonné filtrant des couples $(\mathfrak{J},\mathfrak{R})$ d'idéaux bilatères
tels que $\mathfrak{R}.L=L.\mathfrak{R}=0$,
$\mathfrak{J}.B\subset\mathfrak{R}$,
$B.\mathfrak{J}\subset\mathfrak{R}$.
\end{proposition}

\begin{proof}
Comme $A/\mathfrak{J}$ et $B/\mathfrak{R}$ sont discrets, on a des
homomorphismes canoniques
$w_{\mathfrak{J},\mathfrak{R}}:
\operatorname{D\acute{e}r}_{A/\mathfrak{J}}(B/\mathfrak{R},L)\to
\operatorname{D\acute{e}r.\ cont}_A(B,L)$ formant un système inductif
\sref{0.20.3.4-fr}, d'où l'homomorphisme \sref{0.20.3.5.1-fr} par passage à
la limite inductive. Comme l'homomorphisme
$B/\mathfrak{R}'\to B/\mathfrak{R}$ est surjectif pour
$\mathfrak{R}\supset\mathfrak{R}'$, il résulte aussitôt de la définition que
l'homomorphisme
$\operatorname{D\acute{e}r}_{A/\mathfrak{J}}(B/\mathfrak{R},L)\to
\operatorname{D\acute{e}r}_{A/\mathfrak{J}}(B/\mathfrak{R}',L)$ (avec
$\mathfrak{R}.L=L.\mathfrak{R}=0$,
$\mathfrak{J}.B\subset\mathfrak{R}'$,
$B.\mathfrak{J}\subset\mathfrak{R}'$) est injectif, et il en est évidemment
de même de l'homomorphisme
$\operatorname{D\acute{e}r}_{A/\mathfrak{J}}(B/\mathfrak{R},L)\to
\operatorname{D\acute{e}r}_{A/\mathfrak{J}'}(B/\mathfrak{R},L)$ pour
$\mathfrak{J}'\subset\mathfrak{J}$ \sref{0.20.2.1-fr} ; on en conclut que
l'homomorphisme \sref{0.20.3.5.1-fr} est injectif. D'autre part, si $D$ est
une $A$-dérivation continue de $B$ dans $L$, son noyau contient un idéal
bilatère ouvert $\mathfrak{R}_0$ de $B$, et si $\mathfrak{J}_0$ est un idéal
bilatère ouvert de $A$ tel que $\mathfrak{J}_0B\subset\mathfrak{R}_0$ et
$B\mathfrak{J}_0\subset\mathfrak{R}_0$, il est clair que $D$ est l'image
canonique d'une $(A/\mathfrak{J}_0)$-dérivation de
$B/\mathfrak{R}_0$ dans $L$, donc \sref{0.20.3.5.1-fr} est surjectif.
\end{proof}

\oldpage[0\textsubscript{IV-1}]{123}
\begin{proposition}[20.3.6]
\label{0.20.3.6-fr}
Soient $u:A\to B$, $v:B\to C$ deux homomorphismes continus d'anneaux
topologiques, $L$ un $C$-bimodule discret annulé par un idéal bilatère ouvert
de $C$. On a une suite exacte canonique
\begin{align}
\tag{20.3.6.1}
\label{0.20.3.6.1-fr}
  0\to\operatorname{D\acute{e}r.\ cont}_B(C,L)\xrightarrow{u^0}
  \operatorname{D\acute{e}r.\ cont}_A(C,L)&\xrightarrow{v^0}
  \operatorname{D\acute{e}r.\ cont}_A(B,L)\xrightarrow{\partial} \\
  &\to\operatorname{Exantop}_B(C,L)\xrightarrow{u^1}
  \operatorname{Exantop}_A(C,L)\xrightarrow{v^1}
  \operatorname{Exantop}_A(B,L) \notag
\end{align}
où $\partial$ est défini par passage à la limite inductive à partir de
l'homomorphisme $\partial$ de \sref{0.20.2.2.1-fr} ; cette suite exacte est
fonctorielle en $L$ (dans la catégorie des $C$-bimodules discrets annulés par
des idéaux bilatères ouverts).
\end{proposition}

\begin{proof}
Cela résulte de l'exactitude du foncteur $\varinjlim$, à partir de
\sref{0.20.2.2-fr}.
\end{proof}

\begin{corollary}[20.3.7]
\label{0.20.3.7-fr}
Soient $A$, $B$, $C$ trois anneaux topologiques commutatifs, $u:A\to B$,
$v:B\to C$ deux homomorphismes continus, $L$ un $C$-module discret annulé
par un idéal ouvert de $C$. On a une suite exacte canonique de $A$-modules,
fonctorielle en $L$,
\begin{align}
\tag{20.3.7.1}
\label{0.20.3.7.1-fr}
  0\to\operatorname{D\acute{e}r.\ cont}_B(C,L)\xrightarrow{u^0}
  \operatorname{D\acute{e}r.\ cont}_A(C,L)&\xrightarrow{v^0}
  \operatorname{D\acute{e}r.\ cont}_A(B,L)\xrightarrow{\partial} \\
  &\to\operatorname{Exalcotop}_B(C,L)\xrightarrow{u^1}
  \operatorname{Exalcotop}_A(C,L)\xrightarrow{v^1}
  \operatorname{Exalcotop}_A(B,L). \notag
\end{align}
\end{corollary}

\begin{corollary}[20.3.8]
\label{0.20.3.8-fr}
Sous les hypothèses de \sref{0.20.3.5-fr} (resp.
\sref{0.20.3.6-fr}) on a une suite exacte canonique, fonctorielle en $L$
\begin{equation}
\tag{20.3.8.1}
\label{0.20.3.8.1-fr}
  0\to\operatorname{D\acute{e}r.\ cont}_B(C,L)\xrightarrow{u^0}
  \operatorname{D\acute{e}r.\ cont}_A(C,L)\xrightarrow{v^0}
  \operatorname{D\acute{e}r.\ cont}_A(B,L)\to
  \operatorname{Exantop}_{B/A}(C,L)\to0
\end{equation}
(resp.
\begin{equation}
\tag{20.3.8.2}
\label{0.20.3.8.2-fr}
  0\to\operatorname{D\acute{e}r.\ cont}_B(C,L)\xrightarrow{u^0}
  \operatorname{D\acute{e}r.\ cont}_A(C,L)\xrightarrow{v^0}
  \operatorname{D\acute{e}r.\ cont}_A(B,L)\to
  \operatorname{Exalcotop}_{B/A}(C,L)\to0\text{).}
\end{equation}
\end{corollary}

Nous laissons au lecteur le soin d'écrire les diagrammes analogues à ceux de
\sref{0.20.2.5-fr}.

\subsection{Parties principales et différentielles}
\label{subsection:0.20.4-fr}

Dans toute la suite de ce paragraphe et dans les trois suivants, tous les
anneaux sont supposés \emph{commutatifs}.

\begin{env}[20.4.1]
\label{0.20.4.1-fr}
Soient $A$ un anneau topologique, $B$ une $A$-algèbre topologique ; la
$A$-algèbre $B\otimes_A B$ sera munie de la topologie \emph{produit
tensoriel}, qui en fait une $A$-algèbre topologique ; nous désignerons par $p$
(ou $p_{B/A}$) le $A$-homomorphisme canonique surjectif
\begin{equation}
\tag{20.4.1.1}
\label{0.20.4.1.1-fr}
  p:B\otimes_A B\to B
\end{equation}
tel que $p(b\otimes b')=bb'$ ; il est immédiat que $p$ est continu. Le
\emph{noyau} de $p$ sera noté $\mathfrak{J}_{B/A}$ (ou seulement
$\mathfrak{J}$ s'il n'y a pas de confusion). On désignera par
$j_1:B\to B\otimes_A B$ et $j_2:B\to B\otimes_A B$ les deux
$A$-homomorphismes canoniques, tels que
\[
  j_1(b)=b\otimes1, \qquad j_2(b)=1\otimes b
\]
qui sont continus.
\end{env}

\oldpage[0\textsubscript{IV-1}]{124}
\begin{definition}[20.4.2]
\label{0.20.4.2-fr}
On appelle \emph{$B$-algèbre augmentée des parties principales d'ordre $1$ de
$B$ par rapport à $A$} et l'on note $\mathrm{P}^1_{B/A}$ la $A$-algèbre
topologique quotient
\begin{equation}
\tag{20.4.2.1}
\label{0.20.4.2.1-fr}
  \mathrm{P}^1_{B/A}=(B\otimes_A B)/\mathfrak{J}^2
\end{equation}
munie de la structure de $B$-algèbre topologique définie par l'homomorphisme
$\bar{j}_1:B\to\mathrm{P}^1_{B/A}$ (déduit de $j_1$ par composition avec
l'homomorphisme canonique $B\otimes_A B\to\mathrm{P}^1_{B/A}$), et de
l'augmentation de $B$-algèbre $\varepsilon:\mathrm{P}^1_{B/A}\to B$ (aussi
notée $\varepsilon_{B/A}$) déduite de $p$ par passage au quotient.

Comme $p(b\otimes1)=b$ par définition, il est clair que $\varepsilon$ est
bien une augmentation de $B$-algèbre.
\end{definition}

\begin{definition}[20.4.3]
\label{0.20.4.3-fr}
Le noyau de l'augmentation $\varepsilon:\mathrm{P}^1_{B/A}\to B$,
\begin{equation}
\tag{20.4.3.1}
\label{0.20.4.3.1-fr}
  \Omega^1_{B/A}=\mathfrak{J}_{B/A}/(\mathfrak{J}_{B/A})^2
\end{equation}
muni de la topologie induite par celle de $\mathrm{P}^1_{B/A}$, qui en fait
un $B$-module topologique, est appelé le $B$-module des
\emph{$1$-différentielles} (ou simplement des \emph{différentielles}) de $B$
par rapport à $A$.
\end{definition}

On notera que la topologie de $\Omega^1_{B/A}$ est aussi la topologie quotient
de la topologie induite sur $\mathfrak{J}_{B/A}$ par celle de
$B\otimes_A B$ (Bourbaki, \emph{Top. gén.}, chap.~III, 3\textsuperscript{e}
éd., §~2, n\textsuperscript{o}~7, prop.~20). Si $B$ est discret il en est de
même de $\Omega^1_{B/A}$. On note $\widehat{\Omega}^1_{B/A}$ le
\emph{séparé complété} du $B$-module topologique $\Omega^1_{B/A}$.

Tout anneau topologique $B$ pouvant être considéré comme une
$\mathbf{Z}$-algèbre topologique (où $\mathbf{Z}$ est muni de la topologie
discrète), on peut définir le $B$-module topologique $\Omega^1_{B/\mathbf{Z}}$,
qu'on appelle parfois aussi le $B$-module des \emph{différentielles absolues}
sur $B$ et que l'on note $\Omega^1_B$. Si $B$ est une algèbre topologique sur
un corps premier $P$ (discret), on a
$B\otimes_P B=B\otimes_{\mathbf{Z}}B$, donc
$\Omega^1_{B/P}=\Omega^1_B$.

\begin{lemma}[20.4.4]
\label{0.20.4.4-fr}
Soient $A$ un anneau, $B$ une $A$-algèbre. L'idéal $\mathfrak{J}_{B/A}$ de
$B\otimes_A B$ est engendré par les éléments $1\otimes s-s\otimes1$, où $s$
parcourt un ensemble de générateurs de la $A$-algèbre $B$.
\end{lemma}

\begin{proof}
Il est clair que pour tout $x\in B$, on a
$x\otimes1-1\otimes x\in\mathfrak{J}$ ; d'autre part, quels que soient $x$,
$y$ dans $B$, on a
$x\otimes y=xy\otimes1+(x\otimes1)(1\otimes y-y\otimes1)$. Si
$\sum_i(x_i\otimes y_i)\in\mathfrak{J}$, on a par définition
$\sum_i x_i y_i=0$, donc
\begin{equation}
\tag{20.4.4.1}
\label{0.20.4.4.1-fr}
  \sum_i(x_i\otimes y_i)=
  \sum_i(x_i\otimes1)(1\otimes y_i-y_i\otimes1)
\end{equation}
ce qui prouve que $\mathfrak{J}$ est l'idéal engendré par les éléments
$1\otimes x-x\otimes1$. En outre, si $x=st$, on a
\begin{equation}
\tag{20.4.4.2}
\label{0.20.4.4.2-fr}
  x\otimes1-1\otimes x=(s\otimes1)(t\otimes1-1\otimes t)
  +(s\otimes1-1\otimes s)(1\otimes t)
\end{equation}
ce qui termine aussitôt la démonstration par récurrence.
\end{proof}

\begin{proposition}[20.4.5]
\label{0.20.4.5-fr}
Soient $A$ un anneau topologique, $B$ une $A$-algèbre topologique. La
topologie de $\Omega^1_{B/A}$ est moins fine que la topologie déduite de celle
de $B$ (19.0.2) ; si dans $B$ le carré de tout idéal ouvert est ouvert, ces
deux topologies sont identiques.
\end{proposition}

\begin{proof}
La première assertion est triviale, la topologie produit tensoriel sur
$B\otimes_A B$ étant moins fine que la topologie déduite de celle de $B$ ;
\emph{a fortiori} la topologie induite sur
$\mathfrak{J}=\mathfrak{J}_{B/A}$

\oldpage[0\textsubscript{IV-1}]{125}
par celle de $B\otimes_A B$ est moins fine que la topologie sur $\mathfrak{J}$
déduite de celle de $B$. Pour démontrer la seconde assertion, écrivons
$M\otimes N$, par abus de notation, le sous-module
$\operatorname{Im}(M\otimes_A N)$ pour deux sous-$A$-modules $M$, $N$ de $B$.
Utilisant la relation
\[
\begin{aligned}
  (xy)\otimes z-x\otimes(yz)
  &=(x\otimes1)(1\otimes z)(y\otimes1-1\otimes y)\\
  &=xz.(y\otimes1-1\otimes y)
    -x.(z\otimes1-1\otimes z)(y\otimes1-1\otimes y)
\end{aligned}
\]
dans le $B$-module $B\otimes_A B$ (défini par $j_1$), on voit aussitôt,
compte tenu de \sref{0.20.4.4-fr}, que, si $\mathfrak{R}$ est un idéal de
$B$, on a
\[
  ((\mathfrak{R}^2\otimes B)+(B\otimes\mathfrak{R}^2))\cap\mathfrak{J}
  \subset(\mathfrak{R}\otimes\mathfrak{R})\cap\mathfrak{J}
  +\mathfrak{R}\mathfrak{J}+\mathfrak{J}^2
\]
et d'autre part, si $x_i$, $y_i$ sont des éléments de $\mathfrak{R}$ tels
que $\sum_i(x_i\otimes y_i)\in\mathfrak{J}$, il résulte de
\sref{0.20.4.4.1-fr} que l'on a
$\sum_i(x_i\otimes y_i)\in\mathfrak{R}\mathfrak{J}$, si bien que finalement
\begin{equation}
\tag{20.4.5.1}
\label{0.20.4.5.1-fr}
  (\mathfrak{R}^2\otimes B+B\otimes\mathfrak{R}^2)\cap\mathfrak{J}
  \subset\mathfrak{R}\mathfrak{J}+\mathfrak{J}^2.
\end{equation}
Or, on a un système fondamental de voisinages de $0$ dans $\mathfrak{J}$
(pour la topologie induite par celle de $B\otimes_A B$) en prenant pour
voisinages de $0$ les ensembles
$(\mathfrak{R}\otimes B+B\otimes\mathfrak{R})\cap\mathfrak{J}$, où
$\mathfrak{R}$ parcourt l'ensemble des idéaux ouverts. Comme la topologie de
$\Omega^1_{B/A}$ déduite de celle de $B$ est aussi le quotient par
$\mathfrak{J}^2$ de la topologie de $\mathfrak{J}$ déduite de celle de $B$,
l'hypothèse sur les idéaux ouverts de $B$ et la relation
\sref{0.20.4.5.1-fr} achèvent de prouver la seconde assertion.
\end{proof}

\begin{definition}[20.4.6]
\label{0.20.4.6-fr}
Soient $p_1$ et $p_2$ les applications composées
$B\xrightarrow{j_1}B\otimes_A B\to\mathrm{P}^1_{B/A}$,
$B\xrightarrow{j_2}B\otimes_A B\to\mathrm{P}^1_{B/A}$, qui sont des
$A$-homomorphismes continus tels que
$\varepsilon\circ p_1=\varepsilon\circ p_2=\mathrm{I}_B$. Le
$A$-homomorphisme continu de $A$-modules
\begin{equation}
\tag{20.4.6.1}
\label{0.20.4.6.1-fr}
  d_{B/A}=p_1-p_2:B\to\Omega^1_{B/A}
\end{equation}
est appelé la \emph{différentielle extérieure de $B$ relative à $A$} ; pour
tout $x\in B$, $d_{B/A}(x)$ (aussi noté $d(x)$ ou $dx$) est appelé la
\emph{différentielle de $x$} (relativement à $A$).
\end{definition}

Lorsque $A=\mathbf{Z}$, on écrit $d_B$ au lieu de $d_{B/\mathbf{Z}}$ ; si $B$
est une algèbre sur un corps premier $P$, on a $d_B=d_{B/P}$.

\begin{proposition}[20.4.7]
\label{0.20.4.7-fr}
Le $B$-module $\Omega^1_{B/A}$ est engendré par les éléments $d_{B/A}(x)$,
où $x$ parcourt un système de générateurs de la $A$-algèbre $B$.
\end{proposition}

\begin{proof}
Comme $d_{B/A}(x)$ est l'image canonique de $x\otimes1-1\otimes x$ dans
$\Omega^1_{B/A}$, la proposition est une conséquence immédiate de
\sref{0.20.4.4-fr}.
\end{proof}

\begin{theorem}[20.4.8]
\label{0.20.4.8-fr}
Soient $A$ un anneau topologique, $B$ une $A$-algèbre topologique.
\begin{enumerate}
  \item[\rm(i)] Il existe un isomorphisme unique de $B$-algèbres augmentées
  topologiques
\begin{equation}
\tag{20.4.8.1}
\label{0.20.4.8.1-fr}
  \varphi:\mathrm{P}^1_{B/A}\xrightarrow{\sim}D_B(\Omega^1_{B/A})
\end{equation}
qui se réduit à l'identité dans $\Omega^1_{B/A}$.
  \item[\rm(ii)] L'homomorphisme $d_{B/A}$ est une $A$-dérivation de $B$
  dans $\Omega^1_{B/A}$, ayant la propriété universelle suivante : pour tout
  $B$-module topologique $L$, l'application $u\mapsto u\circ d_{B/A}$ est
  un isomorphisme de $A$-modules
\begin{equation}
\tag{20.4.8.2}
\label{0.20.4.8.2-fr}
  \operatorname{Hom.\ cont}_B(\Omega^1_{B/A},L)
  \xrightarrow{\sim}\operatorname{D\acute{e}r.\ cont}_A(B,L).
\end{equation}
\end{enumerate}
\end{theorem}

\begin{proof}
\oldpage[0\textsubscript{IV-1}]{126}
(i) Il est immédiat que $\varphi$ (avec les notations de
\sref{0.20.4.6-fr}) est nécessairement l'application
$z\mapsto(\varepsilon(z),z-p_1(\varepsilon(z)))$, et l'isomorphisme
réciproque l'application $(b,x)\mapsto p_1(b)+x$ ; ces deux applications
étant continues, cela prouve la première assertion. On notera que cela
entraîne que la topologie de $B$ s'identifie (par $p_1$) au \emph{quotient}
de la topologie de $B\otimes_A B$ par $\mathfrak{J}_{B/A}$.

(ii) Le fait que $d_{B/A}$ soit une $A$-dérivation de $B$ résulte de la
définition \sref{0.20.4.6-fr} et de \sref{0.20.1.1-fr}.

Pour prouver la propriété universelle, rappelons que
$\operatorname{D\acute{e}r.\ cont}_A(B,L)$ s'identifie canoniquement à
l'ensemble $G$ des homomorphismes continus de $A$-algèbres
$u:B\to D_B(L)$ tels que le composé
$B\xrightarrow{u}D_B(L)\xrightarrow{q}B$ soit l'identité (où $q(b,0)=b$ est
la projection ; cf. \sref{0.20.1.6-fr} et \sref{0.20.3.2-fr}). D'autre part,
grâce à l'isomorphisme $\varphi$,
$\operatorname{Hom.\ cont}_B(\Omega^1_{B/A},L)$ s'identifie canoniquement à
l'ensemble des homomorphismes continus de $B$-algèbres
$v:\mathrm{P}^1_{B/A}\to D_B(L)$ tels que le composé
$\mathrm{P}^1_{B/A}\xrightarrow{v}D_B(L)\xrightarrow{q}B$ soit
l'augmentation $\varepsilon$. Comme on a $p_2=p_1-d_{B/A}$ par définition,
tout revient à prouver que tout $u\in G$ se factorise en
\[
\xymatrix{
  B \ar[r]^u \ar[d]_{p_2} & D_B(L) \\
  \mathrm{P}^1_{B/A} \ar[ur]_v &
}
\]
où $v$ est un $B$-homomorphisme continu. Or, on a déjà un homomorphisme
continu de $A$-algèbres $j:b\mapsto(b,0)$ de $B$ dans $D_B(L)$, qui
appartient à $G$ ; par définition du produit tensoriel topologique
d'algèbres topologiques ($\mathbf{0_I}$, 7.7.6), il existe donc un
$A$-homomorphisme continu d'algèbres
$w:B\otimes_A B\to D_B(L)$ rendant commutatif le diagramme
\[
\xymatrix{
  B\otimes_A B \ar[dr]^w & B \ar[l]_{j_1} \ar[d]^j \\
  B \ar[u]^{j_2} \ar[r]_u & D_B(L)
}
\]
On a donc par définition
$w(b\otimes1-1\otimes b)=j(b)-u(b)\in L$, et en vertu de
\sref{0.20.4.4-fr}, cela entraîne $w(\mathfrak{J})\subset L$ donc
$w(\mathfrak{J}^2)=0$ ; par suite $w$ se factorise en
\[
  B\otimes_A B\to\mathrm{P}^1_{B/A}\xrightarrow{v}D_B(L)
\]
où $v$ est un homomorphisme continu de $A$-algèbres ; d'ailleurs, comme
$v\circ p_1=j$ est un homomorphisme de $B$-algèbres, il en est de même de
$v$ par définition de la structure de $B$-algèbre de
$\mathrm{P}^1_{B/A}$ ; comme on a par définition $u=v\circ p_2$, cela
termine la démonstration.
\end{proof}

\begin{theorem}[20.4.9]
\label{0.20.4.9-fr}
Supposons que $B$ soit une $A$-algèbre topologique formellement lisse. Alors
le $B$-module topologique $\Omega^1_{B/A}$ est \emph{formellement projectif}.
\end{theorem}

\begin{proof}
En effet, l'hypothèse entraîne que $B\otimes_A B$ (muni de la structure de
$B$-algèbre topologique définie par $j_1$) est une $B$-algèbre topologique
formellement lisse \sref{0.19.3.5-fr}, (iii), et par suite aussi une
$A$-algèbre topologique formellement lisse \sref{0.19.3.5-fr}, (ii) ; comme
$B$ est topologiquement isomorphe à la $A$-algèbre quotient
$(B\otimes_A B)/\mathfrak{J}$ et est une $A$-algèbre formellement lisse, la
conclusion résulte de \sref{0.19.5.3-fr}.
\end{proof}

\oldpage[0\textsubscript{IV-1}]{127}
\begin{corollary}[20.4.10]
\label{0.20.4.10-fr}
Supposons en outre que dans $B$ le carré de tout idéal ouvert soit ouvert ;
alors, pour tout idéal ouvert $\mathfrak{R}$ de $B$,
$\Omega^1_{B/A}\otimes_B(B/\mathfrak{R})
=\Omega^1_{B/A}/\mathfrak{R}.\Omega^1_{B/A}$ est un
$(B/\mathfrak{R})$-module projectif.
\end{corollary}

\begin{proof}
En effet, la topologie de $\Omega^1_{B/A}$ est alors déduite de celle de
$B$ \sref{0.20.4.5-fr}, et il suffit d'appliquer \sref{0.19.2.4-fr}.
\end{proof}

\begin{corollary}[20.4.11]
\label{0.20.4.11-fr}
Soient $A$, $B$ deux anneaux locaux noethériens, $\rho:A\to B$ un
homomorphisme local, faisant de $B$ une $A$-algèbre topologique formellement
lisse (pour les topologies préadiques). Alors, pour tout idéal de définition
$\mathfrak{b}$ de $B$,
$\Omega^1_{B/A}\otimes_B(B/\mathfrak{b})$ est un
$(B/\mathfrak{b})$-module libre.
\end{corollary}

\begin{proof}
En effet, il résulte de \sref{0.20.4.10-fr} que ce module est projectif, et
comme $B/\mathfrak{b}$ est un anneau artinien, tout
$(B/\mathfrak{b})$-module projectif est libre
($\mathbf{0_{III}}$, 10.1.3).
\end{proof}

\begin{proposition}[20.4.12]
\label{0.20.4.12-fr}
Soient $A$ un anneau topologique, $B$ une $A$-algèbre topologique. Si
l'homomorphisme structural $A\to B$ est surjectif, on a
$\Omega^1_{B/A}=0$.
\end{proposition}

\begin{proof}
En effet, on a $\operatorname{D\acute{e}r}_A(B,L)=0$ pour tout $B$-module
$L$ en vertu de \sref{0.20.1.1-fr}, et la proposition découle donc aussitôt
de \sref{0.20.4.8-fr}.
\end{proof}

\begin{examples}[20.4.13]
\label{0.20.4.13-fr}
\begin{enumerate}
  \item[\rm(i)] Soient $A$ un anneau,
  $B=A[X_\alpha]_{\alpha\in I}$ une algèbre de polynômes sur $A$. Alors
  $\Omega^1_{B/A}$ est un $B$-module libre, dont les $dX_\alpha$ forment une
  base. En effet, les $dX_\alpha$ engendrent ce $B$-module
  \sref{0.20.4.7-fr}. D'autre part, si $L$ est un $B$-module libre, ayant
  une base $(z_\alpha)_{\alpha\in I}$ dont $I$ est l'ensemble d'indices, il
  existe un $A$-homomorphisme $u$ de $B$ dans $D_B(L)$ tel que
  $u(X_\alpha)=(X_\alpha,z_\alpha)$, pour tout $\alpha\in I$, donc
  \sref{0.20.1.5-fr} une $A$-dérivation $D$ de $B$ dans $L$ telle que
  $D(X_\alpha)=z_\alpha$ pour tout $\alpha$ ; en vertu de
  \sref{0.20.4.8.1-fr}, cela prouve que les $dX_\alpha$ sont linéairement
  indépendants.
  \item[\rm(ii)] Soient $A$ un anneau, $L$ un $A$-module, $B$ la
  $A$-algèbre $D_A(L)$ ; alors l'homomorphisme canonique
  $x\mapsto d_{B/A}x$ de $L$ dans $\Omega^1_{B/A}$ est \emph{bijectif}, car
  il est immédiat que les $B$-dérivations de $B=D_A(L)$ dans un $B$-module
  $M$ sont les applications de la forme $(b,x)\mapsto u(x)$, où
  $u\in\operatorname{Hom}_A(L,M)$ ; on conclut donc par
  \sref{0.20.4.8-fr}, (ii).
  \item[\rm(iii)] Supposons que $B$ soit le produit de deux $A$-algèbres
  topologiques $B_1$, $B_2$ (identifiées à des idéaux de $B$). Alors les
  images de $B_1\otimes_A B_2$ et de $B_2\otimes_A B_1$ par l'homomorphisme
  $p$ \sref{0.20.4.1.1-fr} sont nulles, d'où il résulte aussitôt que
  $\mathrm{P}^1_{B/A}$ s'identifie au \emph{produit}
  $\mathrm{P}^1_{B_1/A}\times\mathrm{P}^1_{B_2/A}$, et que le $B$-module
  $\Omega^1_{B/A}$ est \emph{somme directe} (topologique) des $B$-modules
  $\Omega^1_{B_1/A}$ et $\Omega^1_{B_2/A}$ (annulés respectivement par
  $B_2$ et $B_1$).
  \item[\rm(iv)] Soient $A$, $B$ deux anneaux intègres tels que $A\subset
  B$, que $A$ soit intégralement clos, $B$ entier sur $A$, et que le corps
  des fractions de $B$ soit une extension séparable de celui de $A$. Alors
  le $B$-module $\Omega^1_{B/A}$ est un \emph{module de torsion}. En effet,
  pour tout $x\in B$, le polynôme minimal de $x$ par rapport au corps des
  fractions $K$ de $A$ est un polynôme $f(T)$ appartenant à $A[T]$ ; comme
  $x$ est séparable sur $K$, on a $f'(x)\neq0$ et d'autre part on déduit de
  la relation $f(x)=0$ que $f'(x)d_{B/A}x=0$, d'où notre assertion en vertu
  de \sref{0.20.4.7-fr}.
\end{enumerate}
\end{examples}

\begin{remarks}[20.4.14]
\label{0.20.4.14-fr}
\begin{enumerate}
  \item[\rm(i)] On notera que le $B$-module $\Omega^1_{B/A}$, privé de sa
  topologie, est indépendant des topologies de $A$ et de $B$.
  \item[\rm(ii)] Nous introduirons plus tard l'« algèbre des parties
  principales d'ordre $n$ » de $B$ par rapport à $A$,
  $\mathrm{P}^n_{B/A}=(B\otimes_A B)/(\mathfrak{J}_{B/A})^{n+1}$, qui est à
  la base du « calcul différentiel d'ordre $n$ ».
\end{enumerate}
\end{remarks}

\subsection{Propriétés fonctorielles fondamentales de
\texorpdfstring{$\Omega^1_{B/A}$}{Omega1(B/A)}}
\label{subsection:0.20.5-fr}

\oldpage[0\textsubscript{IV-1}]{128}
\begin{env}[20.5.1]
\label{0.20.5.1-fr}
Dans tout ce numéro et le suivant, sauf mention expresse du contraire, les
anneaux et modules considérés sont supposés munis de la topologie
\emph{discrète}.
\end{env}

\begin{env}[20.5.2]
\label{0.20.5.2-fr}
Soient $A$ un anneau, $B$, $C$ deux $A$-algèbres, $u:B\to C$ un
$A$-homomorphisme ; on a un diagramme commutatif
\[
\xymatrix{
  B\otimes_A B \ar[r]^{u\otimes u} \ar[d]_{p_{B/A}} &
  C\otimes_A C \ar[d]^{p_{C/A}} \\
  B \ar[r]^u & C
}
\tag{20.5.2.1}
\label{0.20.5.2.1-fr}
\]
d'où par passage aux quotients, un $A$-homomorphisme d'algèbres
\begin{equation}
\tag{20.5.2.2}
\label{0.20.5.2.2-fr}
  u':\mathrm{P}^1_{B/A}\to\mathrm{P}^1_{C/A}
\end{equation}
tel que le diagramme
\[
\xymatrix{
  \mathrm{P}^1_{B/A} \ar[r]^{u'} & \mathrm{P}^1_{C/A} \\
  B \ar[u]^{p_1} \ar[r]^u & C \ar[u]_{p_1}
}
\]
soit commutatif ; comme $u\otimes u$ applique $\mathfrak{J}_{B/A}$ dans
$\mathfrak{J}_{C/A}$, on obtient, en restreignant $u'$ à
$\Omega^1_{B/A}$, une application
\begin{equation}
\tag{20.5.2.3}
\label{0.20.5.2.3-fr}
  u'':\Omega^1_{B/A}\to\Omega^1_{C/A}
\end{equation}
telle que le couple $(u'',u)$ soit un \emph{di-homomorphisme} pour la
structure de $B$-module de $\Omega^1_{B/A}$ et la structure de $C$-module de
$\Omega^1_{C/A}$ ; ce dernier fait permet d'en déduire canoniquement un
\emph{homomorphisme de $C$-modules}
\begin{equation}
\tag{20.5.2.4}
\label{0.20.5.2.4-fr}
  u_{C/B/A}:\Omega^1_{B/A}\otimes_B C\to\Omega^1_{C/A}.
\end{equation}

\oldpage[0\textsubscript{IV-1}]{129}
En outre, comme le diagramme
\[
\xymatrix{
  \mathrm{P}^1_{B/A} \ar[r]^{u'} & \mathrm{P}^1_{C/A} \\
  B \ar[u]^{p_2} \ar[r]^u & C \ar[u]_{p_2}
}
\tag{20.5.2.5}
\label{0.20.5.2.5-fr}
\]
est aussi commutatif, on en déduit que le diagramme
\[
\xymatrix{
  \Omega^1_{B/A} \ar[r]^{u''} & \Omega^1_{C/A} \\
  B \ar[u]^{d_{B/A}} \ar[r]^u & C \ar[u]_{d_{C/A}}
}
\tag{20.5.2.6}
\label{0.20.5.2.6-fr}
\]
est commutatif.

Enfin, si $w:C\to D$ est un second homomorphisme de $A$-algèbres, on a la
propriété de transitivité
\begin{equation}
\tag{20.5.2.7}
\label{0.20.5.2.7-fr}
  (w\circ u)_{D/B/A}
  =w_{D/C/A}\circ(u_{C/B/A}\otimes\mathrm{I}_D)
\end{equation}
comme il résulte de la définition.
\end{env}

\begin{env}[20.5.3]
\label{0.20.5.3-fr}
Soient maintenant $A$, $B$ deux anneaux, $v:A\to B$ un homomorphisme
d'anneaux, $C$ une $B$-algèbre qui devient une $A$-algèbre au moyen de $v$ ;
alors l'application canonique $v_0:C\otimes_A C\to C\otimes_B C$ est un
\emph{di-homomorphisme} surjectif d'algèbres (relatif à $v:A\to B$) tel que
le diagramme
\[
\xymatrix{
  C\otimes_A C \ar[r]^{v_0} \ar[d]_{p_{C/A}} &
  C\otimes_B C \ar[d]^{p_{C/B}} \\
  C \ar[r]^{\mathrm{I}_C} & C
}
\tag{20.5.3.1}
\label{0.20.5.3.1-fr}
\]
soit commutatif ; par passage aux quotients, on en déduit un
di-homomorphisme d'algèbres
\begin{equation}
\tag{20.5.3.2}
\label{0.20.5.3.2-fr}
  v':\mathrm{P}^1_{C/A}\to\mathrm{P}^1_{C/B}
\end{equation}
tel que le diagramme
\[
\xymatrix{
  \mathrm{P}^1_{C/A} \ar[r]^{v'} & \mathrm{P}^1_{C/B} \\
  C \ar[u]^{p_1} \ar[r]^{\mathrm{I}_C} & C \ar[u]_{p_1}
}
\]
soit commutatif. Comme $v_0$ applique $\mathfrak{J}_{C/A}$ dans
$\mathfrak{J}_{C/B}$ on obtient, en restreignant $v'$ à
$\Omega^1_{C/A}$, une application
\begin{equation}
\tag{20.5.3.3}
\label{0.20.5.3.3-fr}
  v_{C/B/A}:\Omega^1_{C/A}\to\Omega^1_{C/B}
\end{equation}
qui est un homomorphisme de $C$-modules.

\oldpage[0\textsubscript{IV-1}]{130}
En outre, comme le diagramme
\[
\xymatrix{
  \mathrm{P}^1_{C/A} \ar[r]^{v'} & \mathrm{P}^1_{C/B} \\
  C \ar[u]^{p_2} \ar[r]^{\mathrm{I}_C} & C \ar[u]_{p_2}
}
\tag{20.5.3.4}
\label{0.20.5.3.4-fr}
\]
est aussi commutatif, on en déduit que le diagramme
\[
\xymatrix{
  \Omega^1_{C/A} \ar[r]^{v_{C/B/A}} & \Omega^1_{C/B} \\
  C \ar[u]^{d_{C/A}} \ar[r]^{\mathrm{I}_C} & C \ar[u]_{d_{C/B}}
}
\tag{20.5.3.5}
\label{0.20.5.3.5-fr}
\]
est commutatif.

Enfin, si $s:A'\to A$ est un second homomorphisme d'anneaux, on a la
propriété de transitivité
\begin{equation}
\tag{20.5.3.6}
\label{0.20.5.3.6-fr}
  (v\circ s)_{C/B/A'}=v_{C/B/A}\circ s_{C/A/A'}.
\end{equation}
\end{env}

\begin{env}[20.5.4]
\label{0.20.5.4-fr}
Si maintenant on a un diagramme commutatif d'homomorphismes d'anneaux
\[
\xymatrix{
  B \ar[r]^v & B' \\
  A \ar[u]^f \ar[r]^u & A' \ar[u]_g
}
\]
on déduit de \sref{0.20.5.2.4-fr} et \sref{0.20.5.3.3-fr}, par composition,
un homomorphisme de $B'$-modules
\begin{equation}
\tag{20.5.4.1}
\label{0.20.5.4.1-fr}
  \Omega^1_{B/A}\otimes_B B'
  \xrightarrow{v_{B'/B/A}}\Omega^1_{B'/A}
  \xrightarrow{u_{B'/A'/A}}\Omega^1_{B'/A'}
\end{equation}
tel que le diagramme de $A'$-homomorphismes
\[
\xymatrix{
  \Omega^1_{B/A}\otimes_B B' \ar[r] & \Omega^1_{B'/A'} \\
  B' \ar[u]^{d_{B/A}\otimes1} \ar[r]^{\mathrm{I}_{B'}} &
  B' \ar[u]_{d_{B'/A'}}
}
\tag{20.5.4.2}
\label{0.20.5.4.2-fr}
\]
soit commutatif.

L'homomorphisme \sref{0.20.5.4.1-fr} correspond d'ailleurs à un
di-homomorphisme de $B$-modules
\begin{equation}
\tag{20.5.4.3}
\label{0.20.5.4.3-fr}
  \Omega^1_{B/A}\to\Omega^1_{B'/A'}.
\end{equation}
\end{env}

\begin{proposition}[20.5.5]
\label{0.20.5.5-fr}
Si $A'$, $B$ sont deux $A$-algèbres et $B'=B\otimes_A A'$,
l'homomorphisme canonique \sref{0.20.5.4.1-fr}
\begin{equation}
\tag{20.5.5.1}
\label{0.20.5.5.1-fr}
  \Omega^1_{B/A}\otimes_B B'\to\Omega^1_{B'/A'}
\end{equation}
est bijectif.
\end{proposition}

\begin{proof}
\oldpage[0\textsubscript{IV-1}]{131}
Le premier membre de \sref{0.20.5.5.1-fr} n'est autre alors que
$\Omega^1_{B/A}\otimes_A A'$. On peut écrire
$B'\otimes_{A'}B'=(B\otimes_A B)\otimes_A A'$ à un isomorphisme canonique
près et $p_{B'/A'}$ s'identifie alors à
$p_{B/A}\otimes\mathrm{I}_{A'}$ ; par suite (comme $p_{B/A}$ est surjectif)
$\mathfrak{J}_{B'/A'}=\operatorname{Im}(\mathfrak{J}_{B/A}\otimes_A A')$
et
$\mathfrak{J}^2_{B'/A'}
=\operatorname{Im}(\mathfrak{J}^2_{B/A}\otimes_A A')$, d'où
$\mathrm{P}^1_{B'/A'}=\mathrm{P}^1_{B/A}\otimes_A A'$ à un isomorphisme
canonique près, qui transforme en eux-mêmes les idéaux d'augmentation ;
comme le $A$-module $\mathrm{P}^1_{B/A}$ s'identifie canoniquement à la
somme directe de $B$ et de $\Omega^1_{B/A}$, on a
\begin{equation}
\tag{20.5.5.2}
\label{0.20.5.5.2-fr}
  \Omega^1_{B/A}\otimes_A A'=\Omega^1_{B'/A'}
\end{equation}
par le même isomorphisme, et on vérifie aussitôt que le composé de cet
isomorphisme et de l'isomorphisme canonique
$\Omega^1_{B/A}\otimes_B B'\xrightarrow{\sim}
\Omega^1_{B/A}\otimes_A A'$ n'est autre que \sref{0.20.5.5.1-fr}.
\end{proof}

\begin{env}[20.5.6]
\label{0.20.5.6-fr}
Les homomorphismes canoniques \sref{0.20.5.2.4-fr} et
\sref{0.20.5.3.3-fr} donnent, par fonctorialité, pour tout $C$-module $L$,
des homomorphismes canoniques
\begin{align}
\tag{20.5.6.1}
\label{0.20.5.6.1-fr}
  &\operatorname{Hom}_C(\Omega^1_{C/A},L)
  \to\operatorname{Hom}_C(\Omega^1_{B/A}\otimes_B C,L)
  =\operatorname{Hom}_B(\Omega^1_{B/A},L) \\
\tag{20.5.6.2}
\label{0.20.5.6.2-fr}
  &\operatorname{Hom}_C(\Omega^1_{C/B},L)
  \to\operatorname{Hom}_C(\Omega^1_{C/A},L).
\end{align}
Compte tenu de \sref{0.20.4.8.2-fr} et des diagrammes commutatifs
\sref{0.20.5.2.6-fr} et \sref{0.20.5.3.5-fr}, ces homomorphismes ne sont
autres (à une identification canonique près) que les homomorphismes
\sref{0.20.2.1.2-fr} et \sref{0.20.2.1.3-fr} respectivement.
\end{env}

\begin{theorem}[20.5.7]
\label{0.20.5.7-fr}
Soient $u:A\to B$, $v:B\to C$ deux homomorphismes d'anneaux.
\begin{enumerate}
  \item[\rm(i)] La suite de $C$-modules
\begin{equation}
\tag{20.5.7.1}
\label{0.20.5.7.1-fr}
  \Omega^1_{B/A}\otimes_B C
  \xrightarrow{v_{C/B/A}}\Omega^1_{C/A}
  \xrightarrow{u_{C/B/A}}\Omega^1_{C/B}\to0
\end{equation}
est exacte.
  \item[\rm(ii)] Pour que $v_{C/B/A}$ soit inversible à gauche, il faut et
  il suffit que $C$ soit une $B$-algèbre formellement lisse relativement à
  $A$ (pour les topologies discrètes (cf. \sref{0.19.9.1-fr})) ; en
  particulier, il suffit pour cela que $C$ soit une $B$-algèbre formellement
  lisse (pour les topologies discrètes).
\end{enumerate}
\end{theorem}

\begin{proof}
(i) L'exactitude de la suite \sref{0.20.2.4.2-fr} montre tout d'abord,
compte tenu de \sref{0.20.5.6-fr}, que la suite
\[
  0\to\operatorname{Hom}_C(\Omega^1_{C/B},L)
  \to\operatorname{Hom}_C(\Omega^1_{C/A},L)
  \to\operatorname{Hom}_C(\Omega^1_{B/A}\otimes_B C,L)
\]
est exacte pour tout $C$-module $L$. On sait que cela implique l'exactitude
de la suite \sref{0.20.5.7.1-fr} (Bourbaki, \emph{Alg.}, chap.~II,
3\textsuperscript{e} éd., §~2, n\textsuperscript{o}~1, th.~1).

(ii) En vertu de l'exactitude de \sref{0.20.5.7.1-fr}, dire que
$v_{C/B/A}$ est inversible à gauche signifie que la suite
\begin{equation}
\tag{20.5.7.2}
\label{0.20.5.7.2-fr}
  0\to\Omega^1_{B/A}\otimes_B C
  \xrightarrow{v_{C/B/A}}\Omega^1_{C/A}
  \xrightarrow{u_{C/B/A}}\Omega^1_{C/B}\to0
\end{equation}
est \emph{exacte et scindée} ; on sait (Bourbaki, \emph{loc. cit.},
n\textsuperscript{o}~1, prop.~1) que cela équivaut à dire que pour tout
$C$-module $L$, la suite
\[
  0\to\operatorname{Hom}_C(\Omega^1_{C/B},L)
  \to\operatorname{Hom}_C(\Omega^1_{C/A},L)
  \to\operatorname{Hom}_C(\Omega^1_{B/A}\otimes_B C,L)\to0
\]
est exacte ; compte tenu de \sref{0.20.5.6-fr} et
\sref{0.20.2.4.2-fr}, cette condition équivaut à
\[
  \operatorname{Exalcom}_{B/A}(C,L)=0
\]
pour tout $C$-module $L$, et la conclusion résulte donc de
\sref{0.19.9.8.1-fr}.
\end{proof}

\oldpage[0\textsubscript{IV-1}]{132}
Notons en outre que si l'on a un diagramme commutatif d'homomorphismes
d'anneaux
\[
\xymatrix{
  A' \ar[r] & B' \ar[r] & C' \\
  A \ar[u] \ar[r] & B \ar[u] \ar[r] & C \ar[u]
}
\]
on en déduit un diagramme commutatif
\[
\xymatrix{
  \Omega^1_{B/A}\otimes_B C \ar[r] \ar[d] &
  \Omega^1_{C/A} \ar[r] \ar[d] &
  \Omega^1_{C/B} \ar[r] \ar[d] & 0 \\
  \Omega^1_{B'/A'}\otimes_{B'} C' \ar[r] &
  \Omega^1_{C'/A'} \ar[r] &
  \Omega^1_{C'/B'} \ar[r] & 0
}
\tag{20.5.7.3}
\label{0.20.5.7.3-fr}
\]
où les flèches verticales proviennent des di-homomorphismes
\sref{0.20.5.4.3-fr}.

\begin{corollary}[20.5.8]
\label{0.20.5.8-fr}
Supposons que l'homomorphisme $v:B\to C$ fasse de $C$ une $B$-algèbre
formellement étale (pour les topologies discrètes
\sref{0.19.10.2-fr}) ; alors l'homomorphisme \sref{0.20.5.3.3-fr}
\[
  v_{C/B/A}:\Omega^1_{B/A}\otimes_B C\to\Omega^1_{C/A}
\]
est bijectif.
\end{corollary}

\begin{proof}
En effet, si $C$ est une $B$-algèbre formellement non ramifiée pour les
topologies discrètes, il résulte de \sref{0.19.10.4-fr},
\sref{0.20.4.8-fr} et \sref{0.20.1.1-fr} que l'on a
$\operatorname{Hom}_C(\Omega^1_{C/B},L)=0$ pour tout $C$-module $L$, donc
$\Omega^1_{C/B}=0$ (cf. (20.7.4)) ; d'autre part, si $C$ est une
$B$-algèbre formellement lisse pour les topologies discrètes, la suite
\sref{0.20.5.7.2-fr} est exacte ; d'où le corollaire.
\end{proof}

\begin{corollary}[20.5.9]
\label{0.20.5.9-fr}
Soient $A$ un anneau, $B$ une $A$-algèbre, $S$ une partie multiplicative de
$B$ ; alors l'homomorphisme canonique
\begin{equation}
\tag{20.5.9.1}
\label{0.20.5.9.1-fr}
  S^{-1}\Omega^1_{B/A}\to\Omega^1_{S^{-1}B/A}
\end{equation}
est bijectif.
\end{corollary}

\begin{proof}
Il suffit d'appliquer \sref{0.20.5.8-fr} à $C=S^{-1}B$, qui est une
$B$-algèbre formellement étale pour les topologies discrètes
\sref{0.19.10.3-fr}, (ii).
\end{proof}

Compte tenu de \sref{0.20.5.5-fr}, on peut donc écrire
\begin{equation}
\tag{20.5.9.2}
\label{0.20.5.9.2-fr}
  \Omega^1_{S^{-1}B/S^{-1}A}
  =S^{-1}\Omega^1_{B/A}
  =\Omega^1_{S^{-1}B/A},
\end{equation}
à des isomorphismes canoniques près.

\begin{corollary}[20.5.10]
\label{0.20.5.10-fr}
Si $k$ est un corps et $K=k(X_\alpha)_{\alpha\in I}$ une extension pure de
$k$, les $dX_\alpha$ forment une base du $K$-espace vectoriel
$\Omega^1_{K/k}$.
\end{corollary}

\begin{proof}
Comme $K$ est le corps des fractions de l'anneau de polynômes
$k[X_\alpha]_{\alpha\in I}$, cela résulte de \sref{0.20.4.13-fr}, (i) et de
\sref{0.20.5.9-fr}.
\end{proof}

\oldpage[0\textsubscript{IV-1}]{133}
\begin{env}[20.5.11]
\label{0.20.5.11-fr}
Soient $A$ un anneau, $B$ une $A$-algèbre, $\mathfrak{R}$ un idéal de $B$,
$C$ la $A$-algèbre quotient $B/\mathfrak{R}$, et considérons
l'homomorphisme composé de $A$-modules
\begin{equation}
\tag{20.5.11.1}
\label{0.20.5.11.1-fr}
  \delta:\mathfrak{R}\longrightarrow B\xrightarrow{d_{B/A}}\Omega^1_{B/A}
\end{equation}
où la première flèche est l'injection canonique; comme
$d(xy)=xdy+ydx$, on voit que l'on a
$\delta(\mathfrak{R}^2)\subset\mathfrak{R}.\Omega^1_{B/A}$, d'où, par
passage aux quotients, un homomorphisme de $A$-modules
\begin{equation}
\tag{20.5.11.2}
\label{0.20.5.11.2-fr}
  \delta_{C/B/A}:\mathfrak{R}/\mathfrak{R}^2\longrightarrow
  \Omega^1_{B/A}\otimes_B C=
  \Omega^1_{B/A}/(\mathfrak{R}.\Omega^1_{B/A}).
\end{equation}
Mais en fait, $\delta_{C/B/A}$ est un homomorphisme de $C$-modules, car
pour $x\in B$ et $y\in\mathfrak{R}$, on a
$ydx\in\mathfrak{R}.\Omega^1_{B/A}$, donc
$d(xy)\equiv xdy\pmod{\mathfrak{R}.\Omega^1_{B/A}}$, ce qui prouve
d'abord que (20.5.11.2) est un homomorphisme de $B$-modules, et comme
$\mathfrak{R}$ annule les deux membres, cela établit notre assertion.

Si $B'$ est une seconde $A$-algèbre, $u:B\to B'$ un $A$-homomorphisme,
$\mathfrak{R}'$ un idéal de $B'$ tel que
$u(\mathfrak{R})\subset\mathfrak{R}'$ et $C'=B'/\mathfrak{R}'$ l'algèbre
quotient, on a un diagramme commutatif
\begin{equation}
\tag{20.5.11.3}
\label{0.20.5.11.3-fr}
\xymatrix{
  \mathfrak{R}/\mathfrak{R}^2 \ar[r]^-{\delta_{C/B/A}} \ar[d] &
  \Omega^1_{B/A}\otimes_B C \ar[d] \\
  \mathfrak{R}'/\mathfrak{R}'^2 \ar[r]_-{\delta_{C'/B'/A}} &
  \Omega^1_{B'/A}\otimes_{B'} C'
}
\end{equation}
où les flèches verticales proviennent de $u$ (20.5.2.4).
\end{env}

\begin{theorem}[20.5.12]
\label{0.20.5.12-fr}
Soient $A$ un anneau, $B$ une $A$-algèbre, $C$ la $A$-algèbre quotient
$B/\mathfrak{R}$, $u:B\to C$ l'homomorphisme canonique.
\begin{enumerate}
  \item[\rm(i)] On a une suite exacte de $C$-modules
\begin{equation}
\tag{20.5.12.1}
\label{0.20.5.12.1-fr}
  \mathfrak{R}/\mathfrak{R}^2
  \xrightarrow{\delta_{C/B/A}}\Omega^1_{B/A}\otimes_B C
  \xrightarrow{u_{C/B/A}}\Omega^1_{C/A}\longrightarrow 0,
\end{equation}
où $u_{C/B/A}$ et $\delta_{C/B/A}$ sont définis par (20.5.2.4) et
(20.5.11.2) respectivement.
  \item[\rm(ii)] Si l'on pose $E=B/\mathfrak{R}^2$, l'homomorphisme
canonique (20.5.2.4)
\[
  \Omega^1_{B/A}\otimes_B C\longrightarrow
  \Omega^1_{E/A}\otimes_E C
\]
est bijectif.
  \item[\rm(iii)] Les trois conditions suivantes sont équivalentes :
  \begin{enumerate}
    \item[\rm a)] $\delta_{C/B/A}$ est inversible à gauche.
    \item[\rm b)] Toute $A$-extension de $C$ par un $C$-module $L$, dont
l'image réciproque par $u:B\to C$ est $A$-triviale, est elle-même
$A$-triviale.
    \item[\rm c)] La $A$-algèbre $E=B/\mathfrak{R}^2$ est une extension
$A$-triviale de $C$ par $\mathfrak{R}/\mathfrak{R}^2$.
  \end{enumerate}
  \item[\rm(iv)] Il y a correspondance biunivoque canonique entre les
inverses à gauche de $\delta_{C/B/A}$ et les inverses à droite de
l'homomorphisme canonique $E\to C$.
\end{enumerate}
\end{theorem}

\begin{proof}
\oldpage[0\textsubscript{IV-1}]{134}
(i) Comme $u$ est surjectif, on a $\operatorname{D\acute{e}r}_B(C,L)=0$ pour tout $C$-module
$L$ en vertu de (20.1.1). La suite exacte (20.2.3.1) devient donc
\begin{equation}
\tag{20.5.12.2}
\label{0.20.5.12.2-fr}
  0\longrightarrow\operatorname{D\acute{e}r}_A(C,L)\xrightarrow{u^0}\operatorname{D\acute{e}r}_A(B,L)
  \xrightarrow{\partial}\operatorname{Exalcom}_B(C,L)
  \xrightarrow{v^1}\operatorname{Exalcom}_A(C,L)
  \xrightarrow{u^1}\operatorname{Exalcom}_A(B,L),
\end{equation}
où $v$ est l'homomorphisme $A\to B$. Rappelons d'autre part (18.3.8) que
$\operatorname{Exalcom}_B(C,L)$ s'identifie canoniquement à
$\operatorname{Hom}_C(\mathfrak{R}/\mathfrak{R}^2,L)$; on déduit donc de (20.5.12.2)
et (20.4.8) la suite exacte
\begin{equation}
\tag{20.5.12.3}
\label{0.20.5.12.3-fr}
\begin{split}
  0\longrightarrow\operatorname{Hom}_C(\Omega^1_{C/A},L)
  \longrightarrow\operatorname{Hom}_C(\Omega^1_{B/A}\otimes_B C,L)
  \xrightarrow{\varphi}\operatorname{Hom}_C(\mathfrak{R}/\mathfrak{R}^2,L)\\
  \xrightarrow{\psi}\operatorname{Ker}(u^1)\longrightarrow 0
\end{split}
\end{equation}
avec $\varphi=\partial\circ\eta^{-1}$ et $\psi=v^1\circ\eta$. Revenant
aux définitions de $\partial$ (20.2.2) et de $\eta$ (18.3.8), on constate
aussitôt que $\varphi$ est précisément l'homomorphisme
$\operatorname{Hom}(\delta_{C/B/A},1_L)$. L'existence de la suite exacte formée des
4 premiers termes de (20.5.12.3) montre donc que la suite (20.5.12.1)
est exacte (Bourbaki, \emph{Alg.}, chap. II, 3\textsuperscript{e} éd.,
\S~2, n\textsuperscript{o}~1, th.~1).

\medskip
(ii) Appliquons à $B$ et à l'idéal $\mathfrak{R}^2$ la suite exacte
(20.5.12.1), ce qui donne
\begin{equation}
\tag{20.5.12.4}
\label{0.20.5.12.4-fr}
  \mathfrak{R}^2/\mathfrak{R}^4\longrightarrow
  \Omega^1_{B/A}\otimes_B E\longrightarrow\Omega^1_{E/A}\longrightarrow 0,
\end{equation}
d'où, en tensorisant par $C$ (considéré comme $E$-algèbre), la suite
exacte
\[
  \mathfrak{R}^2/\mathfrak{R}^4\otimes_E C\longrightarrow
  \Omega^1_{B/A}\otimes_B C\longrightarrow
  \Omega^1_{E/A}\otimes_E C\longrightarrow 0.
\]
Or, si $x,y$ sont deux éléments de $\mathfrak{R}$, et $\zeta$ la classe
de $xy$ mod. $\mathfrak{R}^4$, l'image
$\delta'(\zeta\otimes1)$ est par définition
$d_{B/A}(xy)\otimes1=(xd_{B/A}(y)+yd_{B/A}(x))\otimes1$; mais comme les
images de $x$ et de $y$ dans $C$ sont nulles, on a aussi
$d_{B/A}(xy)\otimes1=0$ dans $\Omega^1_{B/A}\otimes_B C$, ce qui prouve
notre assertion.

\medskip
(iii) Dire que $\delta_{C/B/A}$ est inversible à gauche revient, compte
tenu de l'exactitude de (20.5.12.1), à dire que la suite
\begin{equation}
\tag{20.5.12.5}
\label{0.20.5.12.5-fr}
  0\longrightarrow\mathfrak{R}/\mathfrak{R}^2
  \xrightarrow{\delta_{C/B/A}}\Omega^1_{B/A}\otimes_B C
  \xrightarrow{u_{C/B/A}}\Omega^1_{C/A}\longrightarrow0
\end{equation}
est \emph{exacte et scindée}, et il revient au même (Bourbaki,
\emph{loc. cit.}) de dire que $\operatorname{Ker}(u^1)=0$ dans la suite exacte
(20.5.12.3) pour tout $L$, ce qui montre l'équivalence des conditions
a) et b) (cf. (18.3.6.2)).

Le fait que b) entraîne c) provient de ce que l'image réciproque par
$u:B\to C$ de la $A$-extension $E$ de $C$ par
$\mathfrak{R}/\mathfrak{R}^2$ est $A$-triviale, $u$ étant composée des
homomorphismes canoniques $B\to E\to C$ (18.1.6). Inversement, c)
entraîne b), car toute $B$-extension de $C$ par $L$ est $B$-équivalente
à $E\oplus_{\mathfrak{R}/\mathfrak{R}^2}L$ (18.3.8), autrement dit sa
classe est l'image de la classe de $E$ (considéré comme $B$-extension)
par l'homomorphisme

\oldpage[0\textsubscript{IV-1}]{135}
$w_*:\operatorname{Exan}_B(C,\mathfrak{R}/\mathfrak{R}^2)
\to\operatorname{Exan}_B(C,L)$ correspondant à un $B$-homomorphisme
$w:\mathfrak{R}/\mathfrak{R}^2\to L$. Le fait que c) entraîne b) résulte
alors de la commutativité du diagramme (18.3.6.5)
\[
\xymatrix{
  \operatorname{Exan}_B(C,\mathfrak{R}/\mathfrak{R}^2)
    \ar[r]^{w_*} \ar[d]_{u^*} &
  \operatorname{Exan}_B(C,L) \ar[d]^{u^*} \\
  \operatorname{Exan}_A(C,\mathfrak{R}/\mathfrak{R}^2)
    \ar[r]_{w_*} & \operatorname{Exan}_A(C,L).
}
\]

\medskip
(iv) On a vu (20.1.7) que les inverses à droite de $E\to C$
correspondent biunivoquement de façon canonique à l'ensemble des éléments
$D\in\operatorname{D\acute{e}r}_A(E,\mathfrak{R}/\mathfrak{R}^2)$ tels que $D(x)=x$ dans
$\mathfrak{R}/\mathfrak{R}^2$, donc aussi, par (20.4.8), à l'ensemble des
$E$-homomorphismes
$h:\Omega^1_{E/A}\to\mathfrak{R}/\mathfrak{R}^2$ tels que le composé
\[
  \mathfrak{R}/\mathfrak{R}^2\xrightarrow{d_{E/A}}
  \Omega^1_{E/A}\xrightarrow{h}\mathfrak{R}/\mathfrak{R}^2
\]
soit l'identité. Par tensorisation avec $C$, on en déduit (puisque
$\mathfrak{R}/\mathfrak{R}^2$ est un $C$-module) que le composé
\[
  \mathfrak{R}/\mathfrak{R}^2\xrightarrow{\delta_{C/E/A}}
  \Omega^1_{E/A}\otimes_E C\xrightarrow{h\otimes1}
  \mathfrak{R}/\mathfrak{R}^2
\]
est l'identité; or, puisque $\mathfrak{R}/\mathfrak{R}^2$ est un
$C$-module, $h\mapsto h\otimes1$ est un isomorphisme de l'ensemble
$\operatorname{Hom}_E(\Omega^1_{E/A},\mathfrak{R}/\mathfrak{R}^2)$ sur
$\operatorname{Hom}_C(\Omega^1_{E/A}\otimes_E C,\mathfrak{R}/\mathfrak{R}^2)$; et par
ailleurs (ii) prouve que l'on peut identifier canoniquement
$\Omega^1_{E/A}\otimes_E C$ et $\Omega^1_{B/A}\otimes_B C$,
$\delta_{C/E/A}$ s'identifiant alors à $\delta_{C/B/A}$. C.Q.F.D.
\end{proof}

\begin{example}[20.5.13]
\label{0.20.5.13-fr}
Soient $B=A[X_\alpha]_{\alpha\in I}$ une algèbre de polynômes sur $A$,
$\mathfrak{R}$ un idéal de $B$, $(P_\lambda)$ un système de générateurs
de $\mathfrak{R}$ et $C=B/\mathfrak{R}$; on sait que
$\Omega^1_{B/A}$ est un $B$-module libre dont les $dX_\alpha$ forment une
base (20.4.13, (i)), donc les $dX_\alpha$ forment aussi une base du
$C$-module libre $\Omega^1_{B/A}\otimes_B C$. D'autre part, il résulte
aussitôt de la définition que l'image de
$\mathfrak{R}/\mathfrak{R}^2$ par $\delta_{C/B/A}$ est le sous-$C$-module
engendré par les
\[
  dP_\lambda=\sum_\alpha
  \frac{\partial P_\lambda}{\partial X_\alpha}\,dX_\alpha.
\]
On en conclut que $\Omega^1_{C/A}$ est isomorphe au quotient du
$C$-module libre ayant les $dX_\alpha$ pour base, par le sous-$C$-module
engendré par les $dP_\lambda$, ce qui donne une description d'un module
de différentielles d'une algèbre \emph{quelconque}, toute $A$-algèbre
$C$ pouvant s'obtenir de la façon précédente.
\end{example}

\begin{corollary}[20.5.14]
\label{0.20.5.14-fr}
Si $C$ est une $A$-algèbre formellement lisse (pour les topologies
discrètes), la suite
\begin{equation}
\tag{20.5.14.1}
\label{0.20.5.14.1-fr}
  0\longrightarrow\mathfrak{R}/\mathfrak{R}^2
  \xrightarrow{\delta_{C/B/A}}\Omega^1_{B/A}\otimes_B C
  \xrightarrow{u_{C/B/A}}\Omega^1_{C/A}\longrightarrow0
\end{equation}
est exacte et scindée.
\end{corollary}

\begin{proof}
En effet, toute $A$-extension de $C$ par un $C$-module est alors triviale
(19.4.4.1).
\end{proof}

\oldpage[0\textsubscript{IV-1}]{136}
\begin{remark}[20.5.15]
\label{0.20.5.15-fr}
Soit $u:A\to B$ un homomorphisme \emph{surjectif} d'anneaux; alors, pour
tout homomorphisme d'anneaux $v:B\to C$, l'homomorphisme canonique
\begin{equation}
\tag{20.5.15.1}
\label{0.20.5.15.1-fr}
  u_{C/B/A}:\Omega^1_{C/A}\longrightarrow\Omega^1_{C/B}
\end{equation}
est \emph{bijectif}; cela résulte en effet de la suite exacte (20.5.7.1),
puisque $\Omega^1_{B/A}=0$ (20.4.12).
\end{remark}

\subsection{Modules d'imperfection et homomorphismes caractéristiques}
\label{subsection:0.20.6-fr}
\label{0.20.6-fr}

\begin{definition}[20.6.1]
\label{0.20.6.1-fr}
Étant donnés deux homomorphismes d'anneaux $u:A\to B$, $v:B\to C$, on
appelle \emph{module d'imperfection de la $B$-algèbre $C$ relativement à
$A$}, et on note $\Upsilon_{C/B/A}$, le $C$-module noyau de
l'homomorphisme
$v_{C/B/A}:\Omega^1_{B/A}\otimes_B C\to\Omega^1_{C/A}$.
On a donc par définition (cf. (20.5.7)) la suite exacte
\begin{equation}
\tag{20.6.1.1}
\label{0.20.6.1.1-fr}
  0\longrightarrow\Upsilon_{C/B/A}\longrightarrow
  \Omega^1_{B/A}\otimes_B C\xrightarrow{v_{C/B/A}}
  \Omega^1_{C/A}\xrightarrow{u_{C/B/A}}\Omega^1_{C/B}
  \longrightarrow0.
\end{equation}
\end{definition}

Lorsque $A=\mathbb{Z}$ (les modules $\Omega^1_{B/A}$ et $\Omega^1_{C/A}$
étant donc les modules de différentielles « absolues » $\Omega^1_B$ et
$\Omega^1_C$), on écrit $\Upsilon_{C/B}$ au lieu de
$\Upsilon_{C/B/\mathbb{Z}}$; lorsque $B$ et $C$ sont des algèbres sur un
corps premier $P$, on a $\Upsilon_{C/B/P}=\Upsilon_{C/B}$.

Soient $R$, $S$ des parties multiplicatives de $B$ et $C$ respectivement,
telles que l'image de $R$ soit contenue dans $S$. Il résulte alors de la
suite exacte (20.6.1.1) et de (20.5.9) que l'on a
\begin{equation}
\tag{20.6.1.2}
\label{0.20.6.1.2-fr}
  \Upsilon_{S^{-1}C/R^{-1}B/A}=S^{-1}\Upsilon_{C/B/A}.
\end{equation}

\begin{proposition}[20.6.2]
\label{0.20.6.2-fr}
Si $C$ est une $B$-algèbre formellement lisse relativement à $A$ (et en
particulier si $C$ est une $B$-algèbre formellement lisse), on a
$\Upsilon_{C/B/A}=0$.
\end{proposition}

\begin{proof}
Cela résulte de (20.5.7, (ii)).
\end{proof}

\begin{proposition}[20.6.3]
\label{0.20.6.3-fr}
Soient $k$ un corps, $K$ une extension de $k$. Pour que $K$ soit une
extension séparable de $k$, il faut et il suffit que $\Upsilon_{K/k}=0$
(autrement dit, l'homomorphisme canonique
$\Omega^1_k\otimes_k K\to\Omega^1_K$ est injectif ou encore (20.4.8),
toute dérivation de $k$ dans $K$ se prolonge en une dérivation de $K$
dans lui-même).
\end{proposition}

\begin{proof}
En effet, il est équivalent de dire que $K$ est séparable sur $k$ ou une
$k$-algèbre formellement lisse (19.6.1); d'autre part, si $P$ est le
corps premier de $k$, $K$ est séparable sur $P$, donc une $P$-algèbre
formellement lisse, et il revient donc au même de dire que $K$ est une
$k$-algèbre formellement lisse ou une $k$-algèbre formellement lisse
relativement à $P$. Enfin, dire que $K$ est une $k$-algèbre formellement
lisse relativement à $P$ équivaut d'après (20.5.7, (ii)) à dire que
l'homomorphisme
$v_{K/k/P}:\Omega^1_k\otimes_k K\to\Omega^1_K$ est inversible à gauche;
mais comme $K$ est un corps, cette dernière condition équivaut à dire
que le noyau de $v_{K/k/P}$, c'est-à-dire $\Upsilon_{K/k}$, est nul.
\end{proof}

\begin{env}[20.6.4]
\label{0.20.6.4-fr}
Considérons un diagramme commutatif
\begin{equation}
\tag{20.6.4.1}
\label{0.20.6.4.1-fr}
\xymatrix{
  A' \ar[r]^{u'} & B' \ar[r]^{v'} & C' \\
  A \ar[u]^f \ar[r]_u & B \ar[u]^g \ar[r]_v & C \ar[u]^h
}
\end{equation}

\oldpage[0\textsubscript{IV-1}]{137}
d'homomorphismes d'anneaux commutatifs. La commutativité du diagramme
correspondant (20.5.7.3) entraîne l'existence d'un unique
$C$-homomorphisme
\begin{equation}
\tag{20.6.4.2}
\label{0.20.6.4.2-fr}
  \Upsilon_{C/B/A}\longrightarrow\Upsilon_{C'/B'/A'}
\end{equation}
déduit canoniquement de (20.6.4.1) et rendant commutatif le diagramme
\begin{equation}
\tag{20.6.4.3}
\label{0.20.6.4.3-fr}
\xymatrix{
  0 \ar[r] & \Upsilon_{C/B/A} \ar[r] \ar[d] &
  \Omega^1_{B/A}\otimes_B C \ar[r]^{v_{C/B/A}} \ar[d] &
  \Omega^1_{C/A} \ar[r]^{u_{C/B/A}} \ar[d] &
  \Omega^1_{C/B} \ar[r] \ar[d] & 0 \\
  0 \ar[r] & \Upsilon_{C'/B'/A'} \ar[r] &
  \Omega^1_{B'/A'}\otimes_{B'} C' \ar[r]_{v_{C'/B'/A'}} &
  \Omega^1_{C'/A'} \ar[r]_{u_{C'/B'/A'}} &
  \Omega^1_{C'/B'} \ar[r] & 0
}
\end{equation}

La donnée de l'homomorphisme (20.6.4.2) équivaut d'ailleurs à celle d'un
$C'$-homomorphisme
\begin{equation}
\tag{20.6.4.4}
\label{0.20.6.4.4-fr}
  \Upsilon_{C/B/A}\otimes_C C'\longrightarrow\Upsilon_{C'/B'/A'}
\end{equation}
qui, composé avec l'homomorphisme canonique
$\Upsilon_{C/B/A}\to\Upsilon_{C/B/A}\otimes_C C'$, redonne (20.6.4.2).
Il est clair que (20.6.4.2) jouit d'une propriété évidente de
transitivité, permettant de dire que $\Upsilon_{C/B/A}$ est un
\emph{foncteur} en le triplet $(A,B,C)$.
\end{env}

\begin{env}[20.6.5]
\label{0.20.6.5-fr}
Il nous sera commode pour la suite, sous les conditions de (20.6.1),
d'introduire un \emph{complexe} (de chaînes) \emph{de $C$-modules}
$\mathrm{K}_\bullet(C/B/A)$ dont les termes sont nuls sauf en degrés 0
et 1, où on prend
\begin{equation}
\tag{20.6.5.1}
\label{0.20.6.5.1-fr}
\left\{
\begin{aligned}
  \mathrm{K}_0(C/B/A)&=\Omega^1_{C/A},\\
  \mathrm{K}_1(C/B/A)&=\Omega^1_{B/A}\otimes_B C,
\end{aligned}
\right.
\end{equation}
l'opérateur différentiel $\mathrm{K}_1\to\mathrm{K}_0$ étant
$v_{C/B/A}$. Cela permet d'écrire (à des isomorphismes canoniques près)
$\Omega^1_{C/B}$ et $\Upsilon_{C/B/A}$ comme les \emph{modules
d'homologie} de ce complexe
\begin{equation}
\tag{20.6.5.2}
\label{0.20.6.5.2-fr}
  \mathrm{H}_0(\mathrm{K}_\bullet(C/B/A))=\Omega^1_{C/B},\qquad
  \mathrm{H}_1(\mathrm{K}_\bullet(C/B/A))=\Upsilon_{C/B/A}.
\end{equation}
\end{env}

De même :
\begin{proposition}[20.6.6]
\label{0.20.6.6-fr}
Sous les hypothèses de (20.6.1), pour tout $C$-module $L$, on a des
$C$-isomorphismes canoniques
\begin{align}
\tag{20.6.6.1}
\label{0.20.6.6.1-fr}
  \mathrm{H}^0(\mathrm{K}_\bullet(C/B/A),L)&\simeq\operatorname{D\acute{e}r}_B(C,L),\\
\tag{20.6.6.2}
\label{0.20.6.6.2-fr}
  \mathrm{H}^1(\mathrm{K}_\bullet(C/B/A),L)&\simeq
  \operatorname{Exalcom}_{B/A}(C,L).
\end{align}
\end{proposition}

\begin{proof}
En effet, le complexe de cochaînes
$\operatorname{Hom}_C(\mathrm{K}_\bullet(C/B/A),L)$ n'est autre, en vertu de (20.4.8)
et (20.5.6), que le complexe
\[
  \cdots\longrightarrow0\longrightarrow\operatorname{D\acute{e}r}_A(C,L)
  \longrightarrow\operatorname{D\acute{e}r}_A(B,L)\longrightarrow0\longrightarrow\cdots
\]
\oldpage[0\textsubscript{IV-1}]{138}
où l'opérateur différentiel est $v^0$ (avec les notations de (20.2.1)).
La proposition résulte alors de la suite exacte (20.2.4.2) et de la
définition des modules de cohomologie
\begin{equation}
\tag{20.6.6.3}
\label{0.20.6.6.3-fr}
  \mathrm{H}^i(\mathrm{K}_\bullet,L)=
  \mathrm{H}^i(\operatorname{Hom}_C(\mathrm{K}_\bullet,L)).
\end{equation}
\end{proof}

Si on a un diagramme commutatif d'homomorphismes d'anneaux
\[
\xymatrix{
  A' \ar[r] & B' \ar[r] & C' \\
  A \ar[u] \ar[r] & B \ar[u] \ar[r] & C \ar[u]
}
\]
les di-homomorphismes (20.5.4.3) définissent un di-homomorphisme de
complexes de modules
\begin{equation}
\tag{20.6.6.4}
\label{0.20.6.6.4-fr}
  \mathrm{K}_\bullet(C/B/A)\longrightarrow
  \mathrm{K}_\bullet(C'/B'/A')
\end{equation}
et les di-homomorphismes qu'on en déduit pour l'homologie ou la
cohomologie s'identifient, par les formules (20.6.5.2), (20.6.6.1) et
(20.6.6.2), aux di-homomorphismes déjà définis dans (20.5.4.3),
(20.6.4.2), (20.2.1) et (18.3.7.2) (compte tenu, pour le dernier, de la
définition de l'opérateur $\partial$ dans la suite exacte (20.2.4.2)).

\begin{env}[20.6.7]
\label{0.20.6.7-fr}
On sait que pour un complexe $\mathrm{K}_\bullet$ de $C$-modules et un
$C$-module $L$, on a des homomorphismes canoniques
\[
  \alpha_i:\mathrm{H}^i(\mathrm{K}_\bullet,L)\longrightarrow
  \operatorname{Hom}_C(\mathrm{H}_i(\mathrm{K}_\bullet),L)
\]
(M, IV, 6). Ici, l'homomorphisme canonique
\[
  \alpha_1:\mathrm{H}^1(\mathrm{K}_\bullet(C/B/A),L)
  \longrightarrow
  \operatorname{Hom}_C(\mathrm{H}_1(\mathrm{K}_\bullet(C/B/A)),L)
\]
se définit immédiatement comme obtenu par passage au quotient par l'image
de $\operatorname{Hom}_C(\mathrm{K}_0,L)$ de l'homomorphisme de restriction
\[
  \operatorname{Hom}_C(\mathrm{K}_1(C/B/A),L)\longrightarrow
  \operatorname{Hom}_C(\mathrm{H}_1(\mathrm{K}_\bullet(C/B/A)),L)
\]
puisque $\mathrm{H}_1(\mathrm{K}_\bullet(C/B/A))$ n'est autre que le
noyau de $\mathrm{K}_1\to\mathrm{K}_0$; tenant compte de (20.6.6.2) et
(20.6.5.2), on obtient donc un $C$-homomorphisme canonique
\begin{equation}
\tag{20.6.7.1}
\label{0.20.6.7.1-fr}
  \operatorname{Exalcom}_{B/A}(C,L)\longrightarrow
  \operatorname{Hom}_C(\Upsilon_{C/B/A},L)
\end{equation}
qui s'explicite de la façon suivante : en vertu de (20.2.4.2), toute
$B$-extension de $C$ par $L$ qui est $A$-triviale provient de la donnée
d'une $A$-dérivation $D$ de $B$ dans $L$, donc (20.4.8) d'un
$C$-homomorphisme $f$ de $\Omega^1_{B/A}\otimes_B C$ dans $L$; on fait
correspondre à la classe de cette extension la \emph{restriction} de $f$
à $\Upsilon_{C/B/A}$, qui ne dépend effectivement que de cette classe et
non du choix de $D$.
\end{env}

\begin{definition}[20.6.8]
\label{0.20.6.8-fr}
Soient $u:A\to B$, $v:B\to C$ deux homomorphismes d'anneaux. Pour toute
$B$-extension $E$ de $C$ par un $C$-module $L$, qui est $A$-triviale
(18.3.7), on appelle \emph{homomorphisme caractéristique de $E$}, et on
note $\chi_E$, le $C$-homomorphisme $\Upsilon_{C/B/A}\to L$, image de la
classe de $E$ par l'homomorphisme canonique (20.6.7.1).
\end{definition}

\oldpage[0\textsubscript{IV-1}]{139}
On peut définir l'homomorphisme $\chi_E$ d'une autre manière :
\begin{proposition}[20.6.9]
\label{0.20.6.9-fr}
Soit $E$ une $B$-extension de $C$ par un $C$-module $L$, qui est
$A$-triviale (18.3.7); alors le diagramme
\begin{equation}
\tag{20.6.9.1}
\label{0.20.6.9.1-fr}
\xymatrix{
  0 \ar[r] & \Upsilon_{C/B/A} \ar[r] \ar[d]^{\chi_E} &
  \Omega^1_{B/A}\otimes_B C \ar[r]^{v_{C/B/A}}
    \ar[d]^{q_{E/B/A}\otimes1_C} &
  \Omega^1_{C/A} \ar[d]^{1} \\
  0 \ar[r] & L \ar[r]_{\delta_{C/E/A}} &
  \Omega^1_{E/A}\otimes_E C \ar[r]_{p_{C/E/A}} &
  \Omega^1_{C/A} \ar[r] & 0
}
\end{equation}
où $q:B\to E$ définit la structure de $B$-extension de $E$ et
$p:E\to C$ est l'homomorphisme d'augmentation, est commutatif et ses
lignes sont exactes.
\end{proposition}

\begin{proof}
La ligne inférieure du diagramme est la suite (20.5.12.5) relative aux
deux homomorphismes $A\xrightarrow{q\circ u}E$ et
$p:E\to C=E/L$; comme $p$ est surjectif et que $E$ est une extension
$A$-triviale, cette suite est exacte et scindée en vertu de (20.5.12,
(iii)). La commutativité du carré de droite de (20.6.9.1) résulte de la
relation $v=p\circ q$ (20.5.2.7); l'image par
$q_{E/B/A}\otimes1_C$ du noyau $\Upsilon_{C/B/A}$ de $v_{C/B/A}$ est
donc contenue dans le noyau $L$ de $p_{C/E/A}$. D'autre part soit
$h:C\to E$ un $A$-homomorphisme inverse à droite de $p$, et soit
$j:L\to E$ l'injection canonique, de sorte que l'on a
$q(b)=h(v(b))+j(D(b))$ pour $b\in B$, où $D$ est la $A$-dérivation de
$B$ dans $L$ définissant la $B$-extension $E$; on peut écrire
$D=f\circ d_{B/A}$, où $f:\Omega^1_{B/A}\to L$ est un
$B$-homomorphisme. En vertu de (20.5.2.6), on a
\[
  q_{E/B/A}(d_{B/A}(b))=d_{E/A}(q(b))
  =d_{E/A}(h(v(b)))+d_{E/A}(j(f(d_{B/A}(b))))
\]
pour $b\in B$. Soit alors
$z=\sum_i d_{B/A}(b_i)\otimes c_i$, où $b_i\in B$ et $c_i\in C$, un
élément de $\Omega^1_{B/A}\otimes_B C$; on a
\[
  (q_{E/B/A}\otimes1_C)(z)=
  \sum_i d_{E/A}(h(v(b_i)))\otimes c_i+
  \sum_i d_{E/A}(j(f(d_{B/A}(b_i))))\otimes c_i.
\]
Dans la première somme, on a
$d_{E/A}(h(v(b_i)))=h_{E/C/A}(d_{C/A}(v(b_i)))$, donc cette somme est
$(h_{E/C/A}\otimes1_C)(v_{C/B/A}(z))$ en vertu de (20.5.2.6). Si l'on
prend $z\in\Upsilon_{C/B/A}$, cette somme est donc nulle, et il reste,
par définition de $\delta_{C/E/A}$,
\[
  (q_{E/B/A}\otimes1_C)(z)=
  \sum_i\delta_{C/E/A}(c_i f(d_{B/A}(b_i)))
  =\delta_{C/E/A}(f(z)),
\]
ce qui démontre la commutativité du carré de gauche dans (20.6.9.1).
\end{proof}

On notera que lorsque l'on suppose seulement que $\delta_{C/E/A}$ est
\emph{injectif} (et non nécessairement inversible à gauche), cette
interprétation permettrait encore de définir $\chi_E$ comme la
restriction de $q_{E/B/A}\otimes1_C$ à $\Upsilon_{C/B/A}$.

\begin{corollary}[20.6.10]
\label{0.20.6.10-fr}
Si $\mathfrak{R}$ est un idéal de $B$, $C=B/\mathfrak{R}$,
$E=B/\mathfrak{R}^2$, et si $E$ est une $B$-extension $A$-triviale
(18.3.7) de $C$ par $\mathfrak{R}/\mathfrak{R}^2$, alors
l'homomorphisme caractéristique
$\chi_E:\Upsilon_{C/B/A}\to\mathfrak{R}/\mathfrak{R}^2$ est bijectif.
\end{corollary}

\begin{proof}
\oldpage[0\textsubscript{IV-1}]{140}
En effet, dans le diagramme (20.6.9.1), les deux flèches verticales de
droite sont des homomorphismes bijectifs (20.5.12, (ii)).
\end{proof}

\begin{theorem}[20.6.11]
\label{0.20.6.11-fr}
Soient $u:A\to B$, $v:B\to C$ deux homomorphismes d'anneaux, $L$ un
$C$-module. Supposons vérifiée l'une des conditions suivantes :
\begin{enumerate}
  \item[\rm(i)] $L$ est un $C$-module injectif.
  \item[\rm(ii)] $\Upsilon_{C/B/A}$ est facteur direct du $C$-module
$\Omega^1_{B/A}\otimes_B C$ et
$u_{C/B/A}:\Omega^1_{C/A}\to\Omega^1_{C/B}$ est inversible à droite.
\end{enumerate}
Alors l'homomorphisme canonique (20.6.7.1)
\[
  \operatorname{Exalcom}_{B/A}(C,L)\longrightarrow
  \operatorname{Hom}_C(\Upsilon_{C/B/A},L)
\]
est bijectif.

En particulier, si $\Omega^1_{C/B}$ et $\Omega^1_{C/A}$ sont des
$C$-modules projectifs, l'homomorphisme canonique (20.6.7.1) est
bijectif.
\end{theorem}

\begin{proof}
Le fait que chacune des conditions (i), (ii) entraîne que (20.6.7.1) est
bijectif résulte dans les deux cas de la définition de $\alpha_1$. On
notera d'ailleurs que la condition (ii) est \emph{nécessaire et
suffisante} pour que l'homomorphisme (20.6.7.1) soit bijectif pour
\emph{tout} $C$-module $L$ (Bourbaki, \emph{Alg.}, chap. II,
3\textsuperscript{e} éd., \S~2, n\textsuperscript{o}~1, prop.~1). Si on
suppose que $\Omega^1_{C/B}$ et $\Omega^1_{C/A}$ sont des $C$-modules
projectifs, alors, dans la suite exacte (20.6.1.1),
$\operatorname{Ker}(u_{C/B/A})$ est un $C$-module projectif, car la suite exacte
\[
  0\longrightarrow\operatorname{Ker}(u_{C/B/A})\longrightarrow\Omega^1_{C/A}
  \longrightarrow\Omega^1_{C/B}\longrightarrow0
\]
est scindée, $\Omega^1_{C/B}$ étant projectif; comme
$\operatorname{Ker}(u_{C/B/A})=\operatorname{Im}(v_{C/B/A})$, la suite exacte
\[
  0\longrightarrow\Upsilon_{C/B/A}\longrightarrow
  \Omega^1_{B/A}\otimes_B C\longrightarrow
  \operatorname{Im}(v_{C/B/A})\longrightarrow0
\]
est scindée.
\end{proof}

\begin{corollary}[20.6.12]
\label{0.20.6.12-fr}
Supposons que $C$ soit une $A$-algèbre formellement lisse. Il existe alors
un homomorphisme canonique
\begin{equation}
\tag{20.6.12.1}
\label{0.20.6.12.1-fr}
  \operatorname{Exalcom}_B(C,L)\longrightarrow
  \operatorname{Hom}_C(\Upsilon_{C/B/A},L).
\end{equation}
En outre, cet homomorphisme est bijectif si l'une des conditions (i),
(ii) de (20.6.11) est vérifiée.
\end{corollary}

\begin{proof}
En effet, il résulte de l'hypothèse sur $C$ que
$\operatorname{Exalcom}_B(C,L)=\operatorname{Exalcom}_{B/A}(C,L)$, et
l'homomorphisme (20.6.12.1) n'est autre que (20.6.7.1).
\end{proof}

\begin{corollary}[20.6.13]
\label{0.20.6.13-fr}
Si $C$ est une $A$-algèbre formellement lisse et si $\Omega^1_{C/B}$ est
un $C$-module projectif, l'homomorphisme (20.6.12.1) est bijectif.
\end{corollary}

\begin{proof}
En effet, on sait alors que $\Omega^1_{C/A}$ est un $C$-module projectif
(20.4.9), les topologies étant discrètes.
\end{proof}

\begin{env}[20.6.14]
\label{0.20.6.14-fr}
Les notations restant les mêmes, supposons maintenant en outre que $A$,
$B$, $C$ soient des $\Lambda$-algèbres et $u$, $v$ des
$\Lambda$-homomorphismes, ce qui revient à se donner trois homomorphismes
d'anneaux
\[
  \Lambda\xrightarrow{s}A\xrightarrow{u}B\xrightarrow{v}C.
\]

\oldpage[0\textsubscript{IV-1}]{141}
On a donc, en dehors du module d'imperfection $\Upsilon_{C/B/A}$, les
modules d'imperfection $\Upsilon_{B/A/\Lambda}$,
$\Upsilon_{C/A/\Lambda}$ et $\Upsilon_{C/B/\Lambda}$, et on a déjà défini
des homomorphismes canoniques de $C$-modules (20.6.4.2)
\begin{align}
\tag{20.6.14.1}
\label{0.20.6.14.1-fr}
  u':\Upsilon_{C/A/\Lambda}&\longrightarrow\Upsilon_{C/B/\Lambda},\\
\tag{20.6.14.2}
\label{0.20.6.14.2-fr}
  s':\Upsilon_{C/B/\Lambda}&\longrightarrow\Upsilon_{C/B/A}.
\end{align}

Comme dans le diagramme commutatif (20.5.7.3)
\begin{equation}
\tag{20.6.14.3}
\label{0.20.6.14.3-fr}
\xymatrix{
  \Omega^1_{A/\Lambda}\otimes_A B \ar[r] \ar[d] &
  \Omega^1_{B/\Lambda} \ar[r] \ar[d] &
  \Omega^1_{B/A} \ar[r] \ar[d] & 0 \\
  \Omega^1_{A/\Lambda}\otimes_A C \ar[r] &
  \Omega^1_{C/\Lambda} \ar[r] & \Omega^1_{C/A} \ar[r] & 0
}
\end{equation}
la ligne inférieure est formée de $C$-modules, on en déduit par
tensorisation un diagramme commutatif
\begin{equation}
\tag{20.6.14.4}
\label{0.20.6.14.4-fr}
\xymatrix{
  \Omega^1_{A/\Lambda}\otimes_A C \ar[r] \ar@{=}[d] &
  \Omega^1_{B/\Lambda}\otimes_B C \ar[r] \ar[d] &
  \Omega^1_{B/A}\otimes_B C \ar[r] \ar[d] & 0 \\
  \Omega^1_{A/\Lambda}\otimes_A C \ar[r] &
  \Omega^1_{C/\Lambda} \ar[r] & \Omega^1_{C/A} \ar[r] & 0
}
\end{equation}
où la première ligne est encore exacte et la flèche verticale de gauche
est l'identité; si l'on pose
\begin{equation}
\tag{20.6.14.5}
\label{0.20.6.14.5-fr}
  \Upsilon^C_{B/A/\Lambda}=
  \operatorname{Ker}(u_{B/A/\Lambda}\otimes1_C)=
  \operatorname{Ker}(\Omega^1_{A/\Lambda}\otimes_A C\longrightarrow
       \Omega^1_{B/\Lambda}\otimes_B C)
\end{equation}
on voit, compte tenu de la définition de $\Upsilon_{C/A/\Lambda}$, que
l'on a un unique $C$-homomorphisme
\begin{equation}
\tag{20.6.14.6}
\label{0.20.6.14.6-fr}
  v':\Upsilon^C_{B/A/\Lambda}\longrightarrow\Upsilon_{C/A/\Lambda}
\end{equation}
rendant commutatif le diagramme
\begin{equation}
\tag{20.6.14.7}
\label{0.20.6.14.7-fr}
\xymatrix{
  0 \ar[r] & \Upsilon^C_{B/A/\Lambda} \ar[r] \ar[d]^{v'} &
  \Omega^1_{A/\Lambda}\otimes_A C \ar[r] \ar@{=}[d] &
  \Omega^1_{B/\Lambda}\otimes_B C \ar[r] \ar[d] &
  \Omega^1_{B/A}\otimes_B C \ar[r] \ar[d] & 0 \\
  0 \ar[r] & \Upsilon_{C/A/\Lambda} \ar[r] &
  \Omega^1_{A/\Lambda}\otimes_A C \ar[r] &
  \Omega^1_{C/\Lambda} \ar[r] & \Omega^1_{C/A} \ar[r] & 0
}
\end{equation}
dont les lignes sont exactes.

Lorsque $\Lambda=\mathbb{Z}$, on écrira $\Upsilon^C_{B/A}$ au lieu de
$\Upsilon^C_{B/A/\mathbb{Z}}$. Si $\Lambda$ est un corps premier, on a
$\Upsilon^C_{B/A/\Lambda}=\Upsilon^C_{B/A}$.
\end{env}

\begin{env}[20.6.15]
\label{0.20.6.15-fr}
Pour étudier les relations entre les modules précédents, nous introduirons
d'une part le complexe de $B$-modules
$\mathrm{K}_\bullet(B/A/\Lambda)$, d'autre part, les complexes
\oldpage[0\textsubscript{IV-1}]{142}
de $C$-modules $\mathrm{K}_\bullet(C/A/\Lambda)$ et
$\mathrm{K}_\bullet(C/B/\Lambda)$ (20.6.5), et en outre les complexes de
$C$-modules suivants. On posera tout d'abord
\begin{equation}
\tag{20.6.15.1}
\label{0.20.6.15.1-fr}
  \mathrm{K}^C_\bullet(B/A/\Lambda)=
  \mathrm{K}_\bullet(B/A/\Lambda)\otimes_B C.
\end{equation}
D'autre part, nous désignerons par $\mathrm{T}_\bullet(C/B/\Lambda)$ le
complexe de $C$-modules dont les termes sont nuls sauf en degrés 0 et 1,
et où
\begin{equation}
\tag{20.6.15.2}
\label{0.20.6.15.2-fr}
  \mathrm{T}_0(C/B/\Lambda)=\mathrm{T}_1(C/B/\Lambda)=
  \Omega^1_{B/\Lambda}\otimes_B C
\end{equation}
l'opérateur différentiel étant l'identité, de sorte que ce complexe est
homotope à 0; posons enfin
\begin{equation}
\tag{20.6.15.3}
\label{0.20.6.15.3-fr}
  \mathrm{K}'_\bullet(C/A/\Lambda)=
  \mathrm{K}_\bullet(C/A/\Lambda)\oplus
  \mathrm{T}_\bullet(C/B/\Lambda).
\end{equation}
En vertu du caractère trivial de $\mathrm{T}_\bullet$, il est clair que
l'on a
\begin{equation}
\tag{20.6.15.4}
\label{0.20.6.15.4-fr}
  \mathrm{H}^i(\mathrm{K}'_\bullet,L)=
  \mathrm{H}^i(\mathrm{K}_\bullet,L)
  \quad\text{et}\quad
  \mathrm{H}_i(\mathrm{K}'_\bullet,L)=
  \mathrm{H}_i(\mathrm{K}_\bullet,L)
\end{equation}
pour tout $C$-module $L$ et tout $i$.
\end{env}
\begin{env}[20.6.16]
\label{0.20.6.16-fr}
Définissons maintenant une suite exacte de complexes, \emph{scindée} en chaque
degré
\begin{equation}
\tag{20.6.16.1}
\label{0.20.6.16.1-fr}
  0\to\mathrm{K}^C_\bullet(B/A/\Lambda)
  \xrightarrow{j}\mathrm{K}'_\bullet(C/A/\Lambda)
  \xrightarrow{p}\mathrm{K}_\bullet(C/B/\Lambda)\to0
\end{equation}
de la façon suivante : désignons pour un moment par
\[
  f:\Omega^1_{A/\Lambda}\otimes_A C\to
    \Omega^1_{B/\Lambda}\otimes_B C,
  \qquad
  g:\Omega^1_{B/\Lambda}\otimes_B C\to\Omega^1_{C/\Lambda}
\]
les homomorphismes canoniques $u_{B/A/\Lambda}\otimes1_C$ et
$v_{C/B/\Lambda}$ respectivement, dont le composé est
$g\circ f=(v\circ u)_{C/A/\Lambda}$ (cf.\ \sref{0.20.6.14.4-fr}).
On prendra $j_1(x)=(x,f(x))$, $p_1(y,z)=z-f(y)$,
$j_0(x)=(g(x),x)$, $p_0(y,z)=g(z)-y$ de sorte que
$\operatorname{Im}(j_1)=\operatorname{Ker}(p_1)$ est le graphe de $f$,
supplémentaire de $\{0\}\oplus\mathrm{T}_1$, et
$\operatorname{Im}(j_0)=\operatorname{Ker}(p_0)$ est le graphe de $g$,
supplémentaire de $\mathrm{K}_1(C/A/\Lambda)\oplus\{0\}$ ; la vérification
de la commutativité du diagramme
\[
\xymatrix{
  0 \ar[r] & \mathrm{K}^C_1(B/A/\Lambda) \ar[r]^{j_1} \ar[d] &
  \mathrm{K}'_1(C/A/\Lambda) \ar[r]^{p_1} \ar[d] &
  \mathrm{K}_1(C/B/\Lambda) \ar[r] \ar[d] & 0 \\
  0 \ar[r] & \mathrm{K}^C_0(B/A/\Lambda) \ar[r]_{j_0} &
  \mathrm{K}'_0(C/A/\Lambda) \ar[r]_{p_0} &
  \mathrm{K}_0(C/B/\Lambda) \ar[r] & 0
}
\]
où les flèches verticales sont les opérateurs différentiels, est immédiate.
\end{env}

\begin{theorem}[20.6.17]
\label{0.20.6.17-fr}
On a une suite exacte de $C$-modules
\begin{equation}
\tag{20.6.17.1}
\label{0.20.6.17.1-fr}
  0\to\Upsilon^C_{B/A/\Lambda}\xrightarrow{v'}
  \Upsilon_{C/A/\Lambda}\xrightarrow{u'}
  \Upsilon_{C/B/\Lambda}\xrightarrow{\partial}
  \Omega^1_{B/A}\otimes_B C\xrightarrow{v_{C/B/A}}
  \Omega^1_{C/A}\xrightarrow{u_{C/B/A}}\Omega^1_{C/B}\to0
\end{equation}
\oldpage[0\textsubscript{IV-1}]{143}
où l'opérateur bord $\partial$ est le composé
\begin{equation}
\tag{20.6.17.2}
\label{0.20.6.17.2-fr}
  \Upsilon_{C/B/\Lambda}\xrightarrow{s'}\Upsilon_{C/B/A}
  \longrightarrow\Omega^1_{B/A}\otimes_B C
\end{equation}
la seconde flèche étant l'injection canonique.
\end{theorem}

\begin{proof}
En écrivant la suite exacte d'homologie pour la suite exacte de complexes
\sref{0.20.6.16.1-fr}, on obtient \sref{0.20.6.17.1-fr}, l'homologie étant
nulle en degrés distincts de 0 et de 1 ; le fait que les homomorphismes de
cette suite exacte qui proviennent par fonctorialité de $j$ et de $p$ sont
bien ceux de l'énoncé est immédiat. Reste à vérifier que $\partial$ est égal
à \sref{0.20.6.17.2-fr} ; or un élément $z\in\Upsilon_{C/B/\Lambda}$ est
image par $p_1$ de $(0,z)$, d'où on déduit aussitôt que $\partial(z)$ est
l'image de $z$ par l'homomorphisme canonique
\[
  s_{B/A/\Lambda}\otimes1_C:
  \Omega^1_{B/\Lambda}\otimes_B C\to\Omega^1_{B/A}\otimes_B C.
\]
Notre assertion résulte de la commutativité du diagramme
\begin{equation}
\tag{20.6.17.3}
\label{0.20.6.17.3-fr}
\xymatrix{
  \Upsilon_{C/B/\Lambda} \ar[r]^{s'} \ar[dd] & \Upsilon_{C/B/A} \ar[dd] \\
  & \\
  \Omega^1_{B/\Lambda}\otimes_B C \ar[r] &
  \Omega^1_{B/A}\otimes_B C
}
\end{equation}
(cf.\ \sref{0.20.6.4.3-fr}).
\end{proof}

\begin{corollary}[20.6.18]
\label{0.20.6.18-fr}
\begin{enumerate}
  \item[\rm(i)] La suite de $C$-modules
\begin{equation}
\tag{20.6.18.1}
\label{0.20.6.18.1-fr}
  0\to\Upsilon^C_{B/A/\Lambda}\xrightarrow{v'}
  \Upsilon_{C/A/\Lambda}\xrightarrow{u'}
  \Upsilon_{C/B/\Lambda}\xrightarrow{s'}\Upsilon_{C/B/A}\to0
\end{equation}
est exacte.
  \item[\rm(ii)] Si $B$ est une $A$-algèbre formellement lisse relativement
  à $\Lambda$, on a $\Upsilon^C_{B/A/\Lambda}=0$.
\end{enumerate}
\end{corollary}

\begin{proof}
(i) Dans \sref{0.20.6.17.1-fr}, l'image de $\Upsilon_{C/B/\Lambda}$ par
$\partial$ est le noyau de $v_{C/B/A}$, donc $\Upsilon_{C/B/A}$ par
définition.

(ii) L'hypothèse entraîne que la suite
\[
  0\to\Omega^1_{A/\Lambda}\otimes_A B\to
  \Omega^1_{B/\Lambda}\to\Omega^1_{B/A}\to0
\]
est exacte et scindée \sref{0.20.5.7-fr} ; par tensorisation avec $C$, la
suite
\[
  0\to\Omega^1_{A/\Lambda}\otimes_A C\to
  \Omega^1_{B/\Lambda}\otimes_B C\to
  \Omega^1_{B/A}\otimes_B C\to0
\]
reste donc exacte, d'où notre assertion.
\end{proof}

\begin{corollary}[20.6.19]
\label{0.20.6.19-fr}
Soient $K$ un corps, $E$, $F$ deux extensions de $K$ telles que
$K\subset E\subset F$.
\begin{enumerate}
  \item[\rm(i)] Si $F$ est une extension séparable de $E$, on a
  $\Upsilon_{F/E/K}=0$ (autrement dit, l'homomorphisme canonique
  $\Omega^1_{E/K}\otimes_E F\to\Omega^1_{F/K}$ est injectif).
  \item[\rm(ii)] Inversement, si $F$ est une extension séparable de $K$ et
  si l'on a $\Upsilon_{F/E/K}=0$, alors $F$ est une extension séparable de
  $E$.
\end{enumerate}
\end{corollary}

\begin{proof}
Si $P$ est le corps premier de $K$, on a en effet la suite exacte
\sref{0.20.6.18-fr}
\[
  \Upsilon_{F/K/P}\to\Upsilon_{F/E/P}\to\Upsilon_{F/E/K}\to0.
\]
Si $\Upsilon_{F/E/P}=\Upsilon_{F/E}=0$, on a donc
$\Upsilon_{F/E/K}=0$, d'où (i) en vertu de \sref{0.20.6.3-fr} ;
inversement, si $\Upsilon_{F/E/K}=0$ et
$\Upsilon_{F/K/P}=\Upsilon_{F/K}=0$, on a
$\Upsilon_{F/E/P}=0$, d'où (ii) en vertu de \sref{0.20.6.3-fr}.
\end{proof}

\oldpage[0\textsubscript{IV-1}]{144}
\begin{corollary}[20.6.20]
\label{0.20.6.20-fr}
\begin{enumerate}
  \item[\rm(i)] Si $K$ est une extension algébrique séparable de $k$, on a
  $\Omega^1_{K/k}=0$.
  \item[\rm(ii)] Soient $k$ un corps de caractéristique 0, $K$ une extension
  de $k$. Pour que $\Omega^1_{K/k}=0$, il faut et il suffit que $K$ soit une
  extension algébrique de $k$. En particulier, pour qu'un corps $K$ de
  caractéristique 0 soit tel que $\Omega^1_K=0$, il faut et il suffit que
  $K$ soit une extension algébrique de $\mathbb{Q}$ (cf.\ (21.4.4) et
  (21.7.5)).
\end{enumerate}
\end{corollary}

\begin{proof}
(i) Pour tout $x\in K$, soit $f$ le polynôme minimal de $x$ sur $k$ ; comme
$f'(x)\ne0$ et $d_{K/k}(f(x))=f'(x)d_{K/k}(x)=0$, on a
$d_{K/k}(x)=0$ et notre assertion résulte de \sref{0.20.4.7-fr}.

(ii) Il existe une extension pure $L$ de $k$ telle que
$k\subset L\subset K$ et que $K$ soit une extension algébrique de $L$.
Comme $K$ est séparable sur $L$, il résulte de \sref{0.20.6.19-fr}, (i) que
la suite \sref{0.20.5.7.2-fr}
\[
  0\to\Omega^1_{L/k}\otimes_L K\to\Omega^1_{K/k}
  \to\Omega^1_{K/L}\to0
\]
est exacte et de (i) que $\Omega^1_{K/L}=0$. La relation
$\Omega^1_{K/k}=0$ est donc équivalente à $\Omega^1_{L/k}=0$, et comme
$L$ est une extension pure de $k$, il résulte de \sref{0.20.5.10-fr} que
la relation $\Omega^1_{L/k}=0$ équivaut à $L=k$.
\end{proof}

\begin{remarks}[20.6.21]
\label{0.20.6.21-fr}
\begin{enumerate}
  \item[\rm(i)] Comme $\Upsilon_{B/A/\Lambda}$ est le noyau de
  \[
    u_{B/A/\Lambda}:\Omega^1_{A/\Lambda}\otimes_A B
    \to\Omega^1_{B/\Lambda}
  \]
  et $\Upsilon^C_{B/A/\Lambda}$ le noyau de
  $u_{B/A/\Lambda}\otimes1_C$, on a un homomorphisme canonique
\begin{equation}
\tag{20.6.21.1}
\label{0.20.6.21.1-fr}
  \Upsilon_{B/A/\Lambda}\otimes_B C\to\Upsilon^C_{B/A/\Lambda}.
\end{equation}
Cet homomorphisme est bijectif lorsque la suite
\begin{equation}
\tag{20.6.21.2}
\label{0.20.6.21.2-fr}
  0\to\Upsilon_{B/A/\Lambda}\otimes_B C\to
  \Omega^1_{A/\Lambda}\otimes_A C\to
  \Omega^1_{B/\Lambda}\otimes_B C\to
  \Omega^1_{B/A}\otimes_B C\to0
\end{equation}
est exacte, ce qui se produit dans les cas suivants :
\begin{enumerate}
  \item[$1^\circ$] $C$ est un $B$-module \emph{plat}.
  \item[$2^\circ$] Les $B$-modules $\Omega^1_{B/\Lambda}$ et
  $\Omega^1_{B/A}$ sont \emph{plats} : en effet, il en est alors de même de
  $\operatorname{Ker}(\Omega^1_{B/\Lambda}\to\Omega^1_{B/A})$
  ($\mathbf{0_I}$, 6.1.2), et la suite \sref{0.20.6.21.2-fr} est alors
  exacte en vertu de ($\mathbf{0_I}$, 6.1.2).
\end{enumerate}
  \item[\rm(ii)] Considérons un diagramme commutatif d'homomorphismes
  d'anneaux
\[
\xymatrix{
  \Lambda' \ar[r] & A' \ar[r] & B' \ar[r] & C' \\
  \Lambda \ar[u] \ar[r] & A \ar[u] \ar[r] & B \ar[u] \ar[r] & C \ar[u]
}
\]
Alors les définitions de \sref{0.20.6.16-fr} montrent que l'on a un
diagramme commutatif de complexes (où les flèches verticales proviennent de
\sref{0.20.6.6.4-fr})
\[
\xymatrix{
  0 \ar[r] & \mathrm{K}^C_\bullet(B/A/\Lambda) \ar[r] \ar[d] &
  \mathrm{K}'_\bullet(C/A/\Lambda) \ar[r] \ar[d] &
  \mathrm{K}_\bullet(C/B/\Lambda) \ar[r] \ar[d] & 0 \\
  0 \ar[r] & \mathrm{K}^{C'}_\bullet(B'/A'/\Lambda') \ar[r] &
  \mathrm{K}'_\bullet(C'/A'/\Lambda') \ar[r] &
  \mathrm{K}_\bullet(C'/B'/\Lambda') \ar[r] & 0
}
\]
\oldpage[0\textsubscript{IV-1}]{145}
d'où, par passage à l'homologie, un \emph{diagramme commutatif}
\begin{equation}
\tag{20.6.21.3}
\label{0.20.6.21.3-fr}
\xymatrix@C=0.7em{
  0 \ar[r] & \Upsilon^C_{B/A/\Lambda} \ar[r] \ar[d] &
  \Upsilon_{C/A/\Lambda} \ar[r] \ar[d] &
  \Upsilon_{C/B/\Lambda} \ar[r] \ar[d] &
  \Omega^1_{B/A}\otimes_B C \ar[r] \ar[d] &
  \Omega^1_{C/A} \ar[r] \ar[d] &
  \Omega^1_{C/B} \ar[r] \ar[d] & 0 \\
  0 \ar[r] & \Upsilon^{C'}_{B'/A'/\Lambda'} \ar[r] &
  \Upsilon_{C'/A'/\Lambda'} \ar[r] &
  \Upsilon_{C'/B'/\Lambda'} \ar[r] &
  \Omega^1_{B'/A'}\otimes_{B'} C' \ar[r] &
  \Omega^1_{C'/A'} \ar[r] & \Omega^1_{C'/B'} \ar[r] & 0
}
\end{equation}
On a un diagramme commutatif analogue pour \sref{0.20.6.18.1-fr}.
\end{enumerate}
\end{remarks}

\begin{proposition}[20.6.22]
\label{0.20.6.22-fr}
Soient $s:\Lambda\to A$, $u:A\to B$ deux homomorphismes d'anneaux,
$\mathfrak{R}$ un idéal de $B$, $C$ l'anneau quotient $B/\mathfrak{R}$.
Supposons que $E=B/\mathfrak{R}^2$ soit une $B$-extension
$\Lambda$-triviale de $C$ par $\mathfrak{R}/\mathfrak{R}^2$. On a alors une
suite exacte
\begin{equation}
\tag{20.6.22.1}
\label{0.20.6.22.1-fr}
  0\to\Upsilon^C_{B/A/\Lambda}\xrightarrow{v'}
  \Upsilon_{C/A/\Lambda}\xrightarrow{\chi_E}
  \mathfrak{R}/\mathfrak{R}^2\xrightarrow{\delta_{C/B/A}}
  \Omega^1_{B/A}\otimes_B C\xrightarrow{v_{C/B/A}}
  \Omega^1_{C/A}\to0.
\end{equation}
\end{proposition}

\begin{proof}
En effet, comme $v:B\to C$ est surjectif, on a $\Omega^1_{C/B}=0$
\sref{0.20.4.12-fr}. En outre, il résulte de \sref{0.20.6.10-fr} que
$\Upsilon_{C/B/\Lambda}$ s'identifie canoniquement à
$\mathfrak{R}/\mathfrak{R}^2$. Il suffit alors d'appliquer les suites exactes
\sref{0.20.6.17.1-fr} et \sref{0.20.6.18.1-fr}, en notant que l'on a un
diagramme commutatif
\[
\xymatrix{
  \Upsilon_{C/A/\Lambda} \ar[r] \ar[d]_{\chi_E} &
  \Omega^1_{A/\Lambda}\otimes_A C \ar[d] \\
  \Upsilon_{C/E/\Lambda}=\Upsilon_{C/B/\Lambda}=\mathfrak{R}/\mathfrak{R}^2
    \ar[r]^-{\delta_{C/B/A}} \ar[d]_{\ell} &
  \Omega^1_{B/A}\otimes_B C=\Omega^1_{E/A}\otimes_E C \ar[d] \\
  \mathfrak{R}/\mathfrak{R}^2 \ar[r]^{\delta_{C/B/A}} &
  \Omega^1_{B/A}\otimes_B C
}
\]
et utilisant la commutativité du diagramme \sref{0.20.6.17.3-fr}.
\end{proof}

On a ainsi précisé le \emph{noyau} de $\delta_{C/B/A}$ dans le cas où il
existe un anneau $\Lambda$ tel que $A$ soit une $\Lambda$-algèbre et que
$B/\mathfrak{R}^2$ soit $\Lambda$-triviale ; on observera que c'est le cas
en particulier lorsque $C$ est une \emph{$\Lambda$-algèbre formellement
lisse} (pour la topologie discrète).

Supposons en outre qu'on ait un diagramme commutatif d'homomorphismes
d'anneaux
\[
\xymatrix{
  \Lambda' \ar[r] & A' \ar[r] & B' \\
  \Lambda \ar[u] \ar[r] & A \ar[u] \ar[r] & B \ar[u]^f
}
\]
\oldpage[0\textsubscript{IV-1}]{146}
que $\mathfrak{R}'$ soit un idéal de $B'$ tel que
$f(\mathfrak{R})\subset\mathfrak{R}'$, et que
$E'=B'/\mathfrak{R}'^2$ soit une $B'$-extension $\Lambda'$-triviale de
$C'=B'/\mathfrak{R}'$ par $\mathfrak{R}'/\mathfrak{R}'^2$. On a alors un
diagramme commutatif
\begin{equation}
\tag{20.6.22.2}
\label{0.20.6.22.2-fr}
\xymatrix@C=0.7em{
  0 \ar[r] & \Upsilon^{C'}_{B'/A'/\Lambda'} \ar[r] &
  \Upsilon_{C'/A'/\Lambda'} \ar[r] &
  \mathfrak{R}'/\mathfrak{R}'^2 \ar[r] &
  \Omega^1_{B'/A'}\otimes_{B'} C' \ar[r] &
  \Omega^1_{C'/A'} \ar[r] & 0 \\
  0 \ar[r] & \Upsilon^C_{B/A/\Lambda} \ar[r] \ar[u] &
  \Upsilon_{C/A/\Lambda} \ar[r] \ar[u] &
  \mathfrak{R}/\mathfrak{R}^2 \ar[r] \ar[u] &
  \Omega^1_{B/A}\otimes_B C \ar[r] \ar[u] &
  \Omega^1_{C/A} \ar[r] \ar[u] & 0
}
\end{equation}
comme il résulte de \sref{0.20.6.21.3-fr} et \sref{0.20.5.11.3-fr}.

\begin{corollary}[20.6.23]
\label{0.20.6.23-fr}
Sous les hypothèses de \sref{0.20.6.22-fr}, supposons en outre que $B$ soit
une $A$-algèbre formellement lisse. Alors on a une suite exacte
\begin{equation}
\tag{20.6.23.1}
\label{0.20.6.23.1-fr}
  0\to\Upsilon_{C/A/\Lambda}\xrightarrow{\chi_E}
  \mathfrak{R}/\mathfrak{R}^2\xrightarrow{\delta_{C/B/A}}
  \Omega^1_{B/A}\otimes_B C\to\Omega^1_{C/A}\to0.
\end{equation}
\end{corollary}

\begin{proof}
Cela résulte en effet de \sref{0.20.6.18-fr}, (ii).
\end{proof}

\begin{env}[20.6.24]
\label{0.20.6.24-fr}
Lorsque les hypothèses de \sref{0.20.6.22-fr} sont satisfaites, on dit encore
que l'homomorphisme caractéristique $\chi_E$ est
l'\emph{homomorphisme caractéristique de la $A$-algèbre $B$, relativement à
$\Lambda$ et à l'idéal $\mathfrak{R}$} (en supprimant ces dernières
précisions lorsqu'il n'y a pas de confusion à craindre) ; on le notera
parfois $\chi_B$ ou $\chi_{B/A}$.
\end{env}

\begin{proposition}[20.6.25]
\label{0.20.6.25-fr}
Soient $s:\Lambda\to A$, $u:A\to B$, $v:B\to C$ trois homomorphismes
d'anneaux, $L$ un $C$-module. On a alors une suite exacte
\begin{equation}
\tag{20.6.25.1}
\label{0.20.6.25.1-fr}
\begin{split}
  0\to\operatorname{D\acute{e}r}_B(C,L)\xrightarrow{u^1}
  \operatorname{D\acute{e}r}_A(C,L)\xrightarrow{v^1}
  \operatorname{D\acute{e}r}_A(B,L)\xrightarrow{\partial} \\
  \to\operatorname{Exalcom}_{B/\Lambda}(C,L)\xrightarrow{u^1}
  \operatorname{Exalcom}_{A/\Lambda}(C,L)\xrightarrow{v^1}
  \operatorname{Exalcom}_{A/\Lambda}(B,L)\to0
\end{split}
\end{equation}
où $u^1$, $v^1$ sont les homomorphismes définis dans (18.3.6.4) et
(18.3.6.2), et $\partial$ est défini comme dans \sref{0.20.2.2-fr}.
\end{proposition}

\begin{proof}
En effet, comme la suite exacte \sref{0.20.6.16.1-fr} est scindée, on en
déduit une suite exacte
\[
  0\to\operatorname{Hom}_C(\mathrm{K}_\bullet(C/B/\Lambda),L)\to
  \operatorname{Hom}_C(\mathrm{K}'_\bullet(C/A/\Lambda),L)\to
  \operatorname{Hom}_C(\mathrm{K}^C_\bullet(B/A/\Lambda),L)\to0.
\]
Si on applique à ce complexe la suite exacte de cohomologie, en tenant compte
de \sref{0.20.6.15.4-fr}, et de \sref{0.20.6.6-fr}, on obtient
\sref{0.20.6.25.1-fr} car on a
\[
  \operatorname{Hom}_C(\mathrm{K}^C_\bullet(B/A/\Lambda),L)
  =\operatorname{Hom}_B(\mathrm{K}_\bullet(B/A/\Lambda),L)
\]
par définition ; l'identification de $u^1$ et $v^1$ avec les homomorphismes
de (18.3.4.2) résulte de \sref{0.20.6.6.4-fr}.
\end{proof}

\begin{remark}[20.6.26]
\label{0.20.6.26-fr}
Dans ce numéro, les complexes $\mathrm{K}_\bullet(C/B/A)$ sont apparus comme
un artifice technique destiné à simplifier l'exposition de certains
comportements fonctoriels. En réalité, ces complexes, considérés comme objets
de la catégorie des
\oldpage[0\textsubscript{IV-1}]{147}
complexes de $C$-modules « à homotopie près » (c'est-à-dire où les morphismes
sont les classes d'homomorphismes homotopes) sont des invariants remarquables,
plus fins que le couple formé de $\Omega^1_{C/B}$ et de
$\Upsilon_{C/B/A}$. Lorsque $k$ est un corps premier, et $C$ une $k$-algèbre
formellement lisse (par exemple un anneau régulier de type fini sur une
extension de $k$ (cf.\ ($\mathbf{IV}$, 6.8.6))), on peut montrer que le
complexe $\mathrm{K}_\bullet(C/A/k)$ peut se décrire uniquement en faisant
intervenir $C$ et $A$ (à l'exclusion de $k$) : on exprime $C$ comme quotient
d'une algèbre de polynômes $B$ sur $A$ par un idéal $\mathfrak{Q}$, et on
considère le complexe $\mathrm{F}_\bullet(C/A)$ à 2 termes non nuls
\[
  \cdots\to0\to\mathfrak{Q}/\mathfrak{Q}^2\to
  \Omega^1_{B/A}\otimes_B C\to0\to\cdots
\]
(dont l'homologie coïncide bien avec celle de
$\mathrm{K}_\bullet(C/A/k)$ en vertu de \sref{0.20.6.23.1-fr}). Ces
complexes $\mathrm{F}_\bullet(C/A)$, qui du point de vue de l'algèbre
homologique jouent le rôle d'un fibré conormal pour $\operatorname{Spec}(C)$
au-dessus de $\operatorname{Spec}(A)$, tiendront une place importante dans
les chapitres de cet ouvrage consacrés à la dualité des faisceaux cohérents et
au théorème de Riemann--Roch.
\end{remark}

\subsection{Généralisations aux anneaux topologiques}
\label{subsection:0.20.7-fr}

\begin{env}[20.7.1]
\label{0.20.7.1-fr}
Il résulte aussitôt des définitions que si, dans \sref{0.20.5.2-fr} et
\sref{0.20.5.3-fr}, les anneaux $A$, $B$, $C$ sont supposés être des anneaux
topologiques et les homomorphismes d'anneaux $u$, $v$ continus, alors les
homomorphismes $u_{C/B/A}$ et $v_{C/B/A}$ sont \emph{continus} (il faut
naturellement prendre sur $\Omega^1_{B/A}\otimes_B C$ la topologie produit
tensoriel).

En outre, $u_{C/B/A}:\Omega^1_{C/A}\to\Omega^1_{C/B}$ est un
\emph{morphisme strict surjectif} de $C$-modules topologiques ; en effet, il en
est ainsi de l'homomorphisme canonique $C\otimes_A C\to C\otimes_B C$, compte
tenu de la définition de la topologie produit tensoriel ; par passage aux
quotients on en déduit que l'homomorphisme correspondant
$\mathrm{P}^1_{C/A}\to\mathrm{P}^1_{C/B}$ est un morphisme strict surjectif,
et $u_{C/B/A}$ est la restriction de ce dernier à $\Omega^1_{C/A}$.

Dans \sref{0.20.5.5-fr}, si $A'$ et $B$ sont des $A$-algèbres topologiques,
et si $B'$ et $\Omega^1_{B/A}\otimes_B B'$ sont munis des topologies produits
tensoriels, l'homomorphisme canonique \sref{0.20.5.5.1-fr} est un isomorphisme
topologique, en tenant compte de ce que $\mathrm{P}^1_{B/A}$ est somme directe
topologique de $B$ et de $\Omega^1_{B/A}$ \sref{0.20.4.8-fr}.
\end{env}

\begin{proposition}[20.7.2]
\label{0.20.7.2-fr}
Soient $u:A\to B$, $v:B\to C$ deux homomorphismes continus d'anneaux
topologiques. Pour que l'homomorphisme continu
$v_{C/B/A}:\Omega^1_{B/A}\otimes_B C\to\Omega^1_{C/A}$ soit formellement
inversible à gauche (cf.\ \sref{0.19.1.5-fr}), il faut et il suffit que $C$
soit une $B$-algèbre formellement lisse relativement à $A$
\sref{0.19.9.1-fr} (et \emph{a fortiori} il suffit que $C$ soit une
$B$-algèbre formellement lisse).
\end{proposition}

\begin{proof}
Dire que $v_{C/B/A}$ est formellement inversible à gauche signifie en effet,
vu que les topologies de $\Omega^1_{C/A}$ et $\Omega^1_{B/A}\otimes_B C$
sont moins fines que celles déduites de la topologie de $C$
\sref{0.20.4.5-fr}, que pour tout $C$-module \emph{discret} $L$, annulé par un
idéal ouvert de $C$, l'homomorphisme canonique
\[
  \operatorname{Hom.\ cont}_C(\Omega^1_{C/A},L)\to
  \operatorname{Hom.\ cont}_C(\Omega^1_{B/A}\otimes_B C,L)
\]
\oldpage[0\textsubscript{IV-1}]{148}
est \emph{surjectif} \sref{0.19.1.5-fr} ; comme
$\operatorname{Hom.\ cont}_C(\Omega^1_{B/A}\otimes_B C,L)=
\operatorname{Hom.\ cont}_B(\Omega^1_{B/A},L)$ par définition de la topologie
produit tensoriel, il revient au même, en vertu de \sref{0.20.4.8-fr}, de dire
que l'homomorphisme canonique
\[
  \operatorname{D\acute{e}r.\ cont}_A(C,L)\to
  \operatorname{D\acute{e}r.\ cont}_A(B,L)
\]
est surjectif. Mais la suite exacte \sref{0.20.3.8.2-fr} montre que cette
condition équivaut à $\operatorname{Exalcotop}_{B/A}(C,L)=0$, c'est-à-dire
précisément au fait que $C$ est formellement lisse relativement à $A$
\sref{0.19.9.8-fr}.
\end{proof}

\begin{corollary}[20.7.3]
\label{0.20.7.3-fr}
Supposons que dans $B$ et dans $C$, le carré d'un idéal ouvert soit ouvert.
Pour que $C$ soit une $B$-algèbre formellement lisse relativement à $A$, il
faut et il suffit que si l'on désigne par $(\mathfrak{R}_\lambda)$ un système
fondamental de voisinages de $0$ formé d'idéaux de $C$, alors, pour tout
$\lambda$, l'homomorphisme
\begin{equation}
\tag{20.7.3.1}
\label{0.20.7.3.1-fr}
  v_{C/B/A}\otimes 1_{C/\mathfrak{R}_\lambda}:
  \Omega^1_{B/A}\otimes_B(C/\mathfrak{R}_\lambda)\to
  \Omega^1_{C/A}\otimes_C(C/\mathfrak{R}_\lambda)
\end{equation}
soit inversible à gauche.
\end{corollary}

\begin{proof}
On sait en effet alors que la topologie de $\Omega^1_{B/A}$ (resp.\
$\Omega^1_{C/A}$) est déduite de celle de $B$ (resp.\ de $C$)
\sref{0.20.4.5-fr} ; on en conclut aussitôt que la topologie de
$\Omega^1_{B/A}\otimes_B C$ est aussi déduite de celle de $C$ ; le corollaire
résulte alors de \sref{0.20.7.2-fr} et \sref{0.19.1.7-fr}.
\end{proof}

\begin{proposition}[20.7.4]
\label{0.20.7.4-fr}
Soient $A$ un anneau topologique, $B$ une $A$-algèbre topologique. Pour que
$B$ soit formellement non ramifiée \sref{0.19.10.2-fr}, il faut et il suffit
que le séparé complété $\widehat{\Omega}^1_{B/A}$ soit nul.
\end{proposition}

\begin{proof}
En effet, il résulte aussitôt de \sref{0.19.10.4-fr} et
\sref{0.20.1.1-fr} que, pour que $B$ soit formellement non ramifiée, il faut et
il suffit que pour tout idéal ouvert $\mathfrak{R}$ de $B$, tout idéal ouvert
$\mathfrak{H}$ de $A$ tel que $\mathfrak{H}B\subset\mathfrak{R}$ et tout
$(B/\mathfrak{R})$-module $L$, on ait
$\operatorname{D\acute{e}r}_{A/\mathfrak{H}}(B/\mathfrak{R},L)=0$,
c'est-à-dire $\operatorname{D\acute{e}r.\ cont}_A(B,L)=0$ pour tout
$B$-module discret $L$ annulé par un idéal ouvert de $B$ ; en vertu de
\sref{0.20.4.8-fr}, cela équivaut à
$\operatorname{Hom.\ cont}_B(\Omega^1_{B/A},L)=0$ pour un tel $B$-module,
d'où aussitôt la proposition.

Lorsque $B$ est \emph{discret}, la condition de l'énoncé de
\sref{0.20.7.4-fr} équivaut donc à $\Omega^1_{B/A}=0$.
\end{proof}

\begin{corollary}[20.7.5]
\label{0.20.7.5-fr}
Soient $A$ un anneau, $\mathfrak{m}$ un idéal de $A$, $B$ une $A$-algèbre ;
munissons $A$ de la topologie $\mathfrak{m}$-préadique, $B$ de la topologie
déduite de celle de $A$, et posons $A_0=A/\mathfrak{m}$,
$B_0=B/\mathfrak{m}B=B\otimes_A A_0$. Alors, pour que $B$ soit une
$A$-algèbre formellement non ramifiée, il faut et il suffit que
$\Omega^1_{B_0/A_0}=0$ (ou encore que $B_0$ soit une $A_0$-algèbre
formellement non ramifiée).
\end{corollary}

\begin{proof}
En effet, on a
$\Omega^1_{B/A}\otimes_B B_0=\Omega^1_{B/A}/\mathfrak{m}.\Omega^1_{B/A}
=\Omega^1_{B_0/A_0}$ en vertu de \sref{0.20.5.5-fr} ; écrire que tout
sous-module ouvert de $\Omega^1_{B/A}$ est égal à $\Omega^1_{B/A}$ équivaut
par \sref{0.20.4.5-fr} à écrire que
$\mathfrak{m}.\Omega^1_{B/A}=\Omega^1_{B/A}$, d'où la conclusion.
\end{proof}

\begin{proposition}[20.7.6]
\label{0.20.7.6-fr}
Soient $u:A\to B$, $v:B\to C$ deux homomorphismes continus d'anneaux
topologiques, et supposons que $v$ fasse de $C$ une $B$-algèbre formellement
étale ; alors $v_{C/B/A}:\Omega^1_{B/A}\otimes_B C\to\Omega^1_{C/A}$ est un
bimorphisme formel \sref{0.19.1.2-fr}.
\end{proposition}

\begin{proof}
En effet, pour tout $C$-module discret $L$ annulé par un idéal ouvert de $C$,
on a alors \sref{0.20.7.4-fr}
$\operatorname{D\acute{e}r.\ cont}_B(C,L)=0$ et par suite
\sref{0.20.3.6.1-fr} l'homomorphisme canonique
\oldpage[0\textsubscript{IV-1}]{149}
$\operatorname{D\acute{e}r.\ cont}_A(C,L)\to
\operatorname{D\acute{e}r.\ cont}_A(B,L)$ est injectif ; il est par ailleurs
surjectif en vertu de \sref{0.20.7.2-fr}, donc il est bijectif ; il en résulte
que l'image de $v_{C/B/A}$ est nécessairement dense dans $\Omega^1_{C/A}$,
sans quoi le quotient de $\Omega^1_{C/A}$ par l'adhérence de
$\operatorname{Im}(v_{C/B/A})$ serait séparé et $\neq0$ et aurait donc un
quotient discret $L\neq0$, contrairement à ce qu'on vient de voir (compte tenu
de \sref{0.20.4.8-fr}). Par suite $v_{C/B/A}$, qui est un monomorphisme formel
en vertu de \sref{0.20.7.2-fr}, est aussi un épimorphisme formel
\sref{0.19.1.2-fr}, donc un bimorphisme formel.
\end{proof}

\begin{corollary}[20.7.7]
\label{0.20.7.7-fr}
Supposons que dans $B$ et dans $C$ le carré de tout idéal ouvert soit ouvert.
Si $C$ est une $B$-algèbre formellement étale, alors, pour tout idéal ouvert
$\mathfrak{R}_\lambda$ de $C$, l'homomorphisme \sref{0.20.7.3.1-fr} est
bijectif.
\end{corollary}

\begin{proof}
En effet, en vertu de \sref{0.19.1.1-fr}, cet homomorphisme est surjectif, et
il est injectif par \sref{0.20.7.3-fr}.
\end{proof}

\begin{proposition}[20.7.8]
\label{0.20.7.8-fr}
Soient $u:A\to B$ un homomorphisme continu d'anneaux topologiques,
$\mathfrak{R}$ un idéal de $B$, $C$ l'anneau topologique quotient
$B/\mathfrak{R}$, $v:B\to C$ l'homomorphisme canonique. Alors :
\begin{enumerate}
  \item[\rm(i)] Dans la suite exacte \sref{0.20.5.12.1-fr}
\[
  \mathfrak{R}/\mathfrak{R}^2\xrightarrow{\delta_{C/B/A}}
  \Omega^1_{B/A}\otimes_B C\xrightarrow{v_{C/B/A}}\Omega^1_{C/A}\to0
\]
l'homomorphisme $\delta_{C/B/A}$ est continu et l'homomorphisme $v_{C/B/A}$
est un morphisme strict de $C$-modules topologiques.
  \item[\rm(ii)] Pour que $\delta_{C/B/A}$ soit formellement inversible à
gauche \sref{0.19.1.5-fr}, il faut et il suffit que pour tout $C$-module
discret $L$ annulé par un idéal ouvert de $C$, l'homomorphisme canonique
\begin{equation}
\tag{20.7.8.1}
\label{0.20.7.8.1-fr}
  \operatorname{Exalcotop}_A(C,L)\to\operatorname{Exalcotop}_A(B,L)
\end{equation}
soit injectif.
\end{enumerate}
\end{proposition}

\begin{proof}
(i) La première assertion est évidente. D'autre part, l'homomorphisme
canonique $v\otimes v:B\otimes_A B\to C\otimes_A C$ est un morphisme strict
par définition de la topologie produit tensoriel, et on en déduit aussitôt
(cf.\ \sref{0.20.5.2-fr}) qu'il en est de même de $v_{C/B/A}$.

(ii) Dire que $\delta_{C/B/A}$ est formellement inversible à gauche signifie
que pour tout $C$-module discret $L$ annulé par un idéal ouvert de $C$,
l'homomorphisme canonique
\[
  \operatorname{Hom.\ cont}_C(\Omega^1_{B/A}\otimes_B C,L)\to
  \operatorname{Hom.\ cont}_C(\mathfrak{R}/\mathfrak{R}^2,L)
\]
est surjectif. Or, compte tenu de \sref{0.18.4.3-fr} et de
\sref{0.20.4.8-fr}, cela revient à dire que l'homomorphisme canonique
\[
  \operatorname{D\acute{e}r.\ cont}_A(B,L)\to
  \operatorname{Exalcotop}_B(C,L)
\]
est surjectif, et la conclusion résulte donc de la suite exacte
\sref{0.20.3.6.1-fr}.

Le fait que \sref{0.20.7.8.1-fr} soit injectif s'exprime encore de la manière
suivante, compte tenu de la définition des deux membres
\sref{0.18.4.1-fr} : pour tout idéal ouvert $\mathfrak{M}$ de $A$, tout idéal
ouvert $\mathfrak{N}$ de $B$ tel que $\mathfrak{M}B\subset\mathfrak{N}$, et
toute $(A/\mathfrak{M})$-extension $E$ de
$B/(\mathfrak{R}+\mathfrak{N})$ par un
$(B/(\mathfrak{R}+\mathfrak{N}))$-module $L$, telle que l'image réciproque de
$E$ par l'homomorphisme canonique
$B/\mathfrak{N}\to B/(\mathfrak{R}+\mathfrak{N})$ soit
$(A/\mathfrak{M})$-triviale, il existe un idéal ouvert
$\mathfrak{M}'\subset\mathfrak{M}$ de $A$,
\oldpage[0\textsubscript{IV-1}]{150}
un idéal ouvert $\mathfrak{N}'\subset\mathfrak{N}$ de $B$ tels que
$\mathfrak{M}'B\subset\mathfrak{N}'$ et que l'image réciproque de $E$ par
l'homomorphisme canonique
$B/(\mathfrak{R}+\mathfrak{N}')\to B/(\mathfrak{R}+\mathfrak{N})$ soit
$(A/\mathfrak{M}')$-triviale.
\end{proof}

\begin{corollary}[20.7.9]
\label{0.20.7.9-fr}
Si la $A$-algèbre topologique $C=B/\mathfrak{R}$ est formellement lisse,
l'homomorphisme canonique $\delta_{C/B/A}$ est formellement inversible à
gauche.
\end{corollary}

\begin{env}[20.7.10]
\label{0.20.7.10-fr}
Dans \sref{0.20.6.1-fr}, lorsque $A$, $B$, $C$ sont des anneaux topologiques,
$u$ et $v$ des homomorphismes continus, on munit $\Upsilon_{C/B/A}$ de la
topologie induite par celle de $\Omega^1_{B/A}\otimes_B C$ ; les
homomorphismes \sref{0.20.6.4.2-fr} et \sref{0.20.6.4.4-fr} sont alors
continus pourvu qu'il en soit de même de ceux du diagramme
\sref{0.20.6.4.1-fr}. En outre, si, dans \sref{0.20.6.7-fr}, on suppose que
$L$ est un $C$-module discret annulé par un idéal ouvert de $C$, on déduit,
par passage à la limite inductive, un $C$-homomorphisme canonique
\begin{equation}
\tag{20.7.10.1}
\label{0.20.7.10.1-fr}
  \operatorname{Exalcotop}_{B/A}(C,L)\to
  \operatorname{Hom.\ cont}_C(\Upsilon_{C/B/A},L)
\end{equation}
compte tenu de \sref{0.18.5.3.1-fr} : pour tout idéal ouvert
$\mathfrak{M}$ de $A$, tout idéal ouvert $\mathfrak{N}$ de $B$ tel que
$\mathfrak{M}B\subset\mathfrak{N}$, tout idéal ouvert $\mathfrak{P}$ de $C$
tel que $\mathfrak{N}C\subset\mathfrak{P}$ et que $\mathfrak{P}.L=0$, toute
$(B/\mathfrak{N})$-extension $E$ de $C/\mathfrak{P}$ par $L$ qui est
$(A/\mathfrak{M})$-triviale provient de la donnée d'un $C$-homomorphisme
continu $f$ de $\Omega^1_{B/A}\otimes_B C$ dans $L$, et l'homomorphisme
\sref{0.20.7.10.1-fr} fait correspondre à l'image dans
$\operatorname{Exalcotop}_{B/A}(C,L)$ de la classe de $E$, la restriction de
$f$ à $\Upsilon_{C/B/A}$, dit \emph{homomorphisme caractéristique de $E$} et
noté encore $\chi_E$.
\end{env}

\begin{proposition}[20.7.11]
\label{0.20.7.11-fr}
Supposons que la topologie de $C$ soit telle que le carré d'un idéal ouvert
soit ouvert. Si $C$ est une $A$-algèbre topologique formellement lisse et si
$\Omega^1_{C/B}$ est un $C$-module formellement projectif, on a un
isomorphisme canonique
\begin{equation}
\tag{20.7.11.1}
\label{0.20.7.11.1-fr}
  \operatorname{Exalcotop}_B(C,L)\xrightarrow{\sim}
  \operatorname{Hom.\ cont}_C(\Upsilon_{C/B/A},L)
\end{equation}
pour tout $C$-module discret $L$ annulé par un idéal ouvert de $C$.
\end{proposition}

\begin{proof}
On a en effet, dans la suite exacte \sref{0.20.3.7.1-fr},
$\operatorname{Exalcotop}_A(C,L)=0$ \sref{0.19.4.4-fr}, donc
$\operatorname{Exalcotop}_B(C,L)=\operatorname{Exalcotop}_{B/A}(C,L)$ et
l'homomorphisme \sref{0.20.7.8.1-fr} n'est autre que
\sref{0.20.7.10.1-fr} ; le fait qu'il est bijectif se déduit de
\sref{0.20.6.13-fr} par passage à la limite inductive, compte tenu de ce que
la topologie de $\Omega^1_{C/B}$ est alors déduite de celle de $C$
\sref{0.20.4.5-fr} et de \sref{0.19.2.4-fr}.
\end{proof}

\begin{env}[20.7.12]
\label{0.20.7.12-fr}
Dans \sref{0.20.6.14-fr} on munit encore $\Upsilon^C_{B/A/\Lambda}$ de la
topologie induite par celle de $\Omega^1_{A/\Lambda}\otimes_A C$ et alors
l'homomorphisme \sref{0.20.6.14.6-fr} est continu, lorsque les anneaux
considérés sont topologiques et les homomorphismes d'anneaux continus.
\end{env}

\begin{env}[20.7.13]
\label{0.20.7.13-fr}
Si, dans \sref{0.20.6.23-fr}, on suppose que les anneaux $\Lambda$, $A$, $B$,
$C$ sont topologiques, les homomorphismes $s$, $u$, $v$ continus et $L$ un
$C$-module discret annulé par un idéal ouvert de $C$, alors on peut passer à
la limite inductive comme dans \sref{0.20.3.6-fr}, et on a une suite exacte
\begin{equation}
\tag{20.7.13.1}
\label{0.20.7.13.1-fr}
\begin{split}
  0\to\operatorname{D\acute{e}r.\ cont}_B(C,L)\to
  \operatorname{D\acute{e}r.\ cont}_A(C,L)\to
  \operatorname{D\acute{e}r.\ cont}_A(B,L)\to \\
  \to\operatorname{Exalcotop}_{B/\Lambda}(C,L)\to
  \operatorname{Exalcotop}_{A/\Lambda}(C,L)\to
  \operatorname{Exalcotop}_{A/\Lambda}(B,L)\to0.
\end{split}
\end{equation}
\end{env}

\begin{env}[20.7.14]
\label{0.20.7.14-fr}
Soient $A$ un anneau topologique, $B$ une $A$-algèbre topologique,
$\mathfrak{M}'$ un idéal ouvert de $A$, $\mathfrak{N}'$ un idéal ouvert de
$B$ tels que $\mathfrak{M}'B\subset\mathfrak{N}'$ ; posons
$A'=A/\mathfrak{M}'$,
\oldpage[0\textsubscript{IV-1}]{151}
$B'=B/\mathfrak{N}'$ ; le noyau de l'homomorphisme
$B\otimes_A B\to B'\otimes_{A'}B'$ est
$\mathfrak{U}'=\operatorname{Im}(\mathfrak{N}'\otimes B+B\otimes\mathfrak{N}')$,
d'où résulte aussitôt que le noyau de l'homomorphisme
\begin{equation}
\tag{20.7.14.1}
\label{0.20.7.14.1-fr}
  \varphi_{(\mathfrak{M}',\mathfrak{N}')}:
  \Omega^1_{B/A}\to\Omega^1_{B'/A'}
\end{equation}
est $((\mathfrak{J}\cap\mathfrak{U}')+\mathfrak{J}^2)/\mathfrak{J}^2$ ; par
ailleurs, l'homomorphisme \sref{0.20.7.14.1-fr} est surjectif, comme il
résulte de \sref{0.20.4.7-fr}. Si $\mathfrak{M}''$ (resp.\
$\mathfrak{N}''$) est un second idéal ouvert de $A$ (resp.\ $B$) tel que
$\mathfrak{M}''\subset\mathfrak{M}'$, $\mathfrak{N}''\subset\mathfrak{N}'$
et $\mathfrak{M}''B\subset\mathfrak{N}''$, et si l'on pose
$A''=A/\mathfrak{M}''$, $B''=B/\mathfrak{N}''$, on a de même un
homomorphisme surjectif
\[
  \varphi_{(\mathfrak{M}'',\mathfrak{N}''),(\mathfrak{M}',\mathfrak{N}')}:
  \Omega^1_{B''/A''}\to\Omega^1_{B'/A'}
\]
et ces homomorphismes forment évidemment un \emph{système projectif}. Si
l'on remarque que les
$((\mathfrak{J}\cap\mathfrak{U}')+\mathfrak{J}^2)/\mathfrak{J}^2$ forment
un système fondamental de voisinages de $0$ dans $\Omega^1_{B/A}$, on voit
donc que le \emph{séparé complété} $\widehat{\Omega}^1_{B/A}$ du $B$-module
topologique $\Omega^1_{B/A}$ est donné, à un isomorphisme canonique près, par
\begin{equation}
\tag{20.7.14.2}
\label{0.20.7.14.2-fr}
  \widehat{\Omega}^1_{B/A}=\varprojlim(\Omega^1_{B'/A'}).
\end{equation}

En outre, l'homomorphisme canonique
$j:\Omega^1_{B/A}\to\widehat{\Omega}^1_{B/A}$ est limite projective du
système projectif des $\varphi_{(\mathfrak{M}',\mathfrak{N}')}$, donc
$\varphi_{(\mathfrak{M}',\mathfrak{N}')}$ se factorise en
$\Omega^1_{B/A}\xrightarrow{j}\widehat{\Omega}^1_{B/A}\to
\Omega^1_{B'/A'}$, et comme il est surjectif on en conclut que
l'homomorphisme canonique
\begin{equation}
\tag{20.7.14.3}
\label{0.20.7.14.3-fr}
  \widehat{\Omega}^1_{B/A}\to\Omega^1_{B'/A'}
\end{equation}
est surjectif pour tout couple $(\mathfrak{M}',\mathfrak{N}')$.

On peut d'ailleurs dans ce qui précède remplacer partout $A'$ par $A$.

Enfin, si $L$ est un $B$-module topologique \emph{séparé et complet}, tout
$B$-homomorphisme continu de $\Omega^1_{B/A}$ dans $L$ se prolonge de façon
unique en un $\widehat B$-homomorphisme continu de
$\widehat{\Omega}^1_{B/A}$ dans $L$, et réciproquement un tel homomorphisme
redonne par composition avec $\Omega^1_{B/A}\to\widehat{\Omega}^1_{B/A}$ un
$B$-homomorphisme continu, si bien que l'on a un isomorphisme canonique
\[
  \operatorname{Hom.\ cont}_B(\Omega^1_{B/A},L)\xrightarrow{\sim}
  \operatorname{Hom.\ cont}_{\widehat B}(\widehat{\Omega}^1_{B/A},L).
\]
Plus particulièrement, si $L$ est un $B$-module discret annulé par un idéal
ouvert de $B$, on voit que l'isomorphisme canonique
\sref{0.20.4.8.2-fr} peut aussi s'écrire
\begin{equation}
\tag{20.7.14.4}
\label{0.20.7.14.4-fr}
  \operatorname{Hom.\ cont}_{\widehat B}(\widehat{\Omega}^1_{B/A},L)
  \xrightarrow{\sim}\operatorname{D\acute{e}r.\ cont}_A(B,L).
\end{equation}
\end{env}

\begin{proposition}[20.7.15]
\label{0.20.7.15-fr}
Soient $A$ un anneau discret, $B$ une $A$-algèbre topologique adique
($\mathbf{0}_{\mathrm I}$, 7.1.9), $\mathfrak{n}$ un idéal de définition de
$B$, $B_0=B/\mathfrak{n}$. On suppose que $\Omega^1_{B_0/A}$ et
$\mathfrak{n}/\mathfrak{n}^2$ sont des $B_0$-modules de type fini. Alors
$\widehat{\Omega}^1_{B/A}$ est un $\widehat B$-module de type fini.
\end{proposition}

\begin{proof}
Comme le carré de tout idéal ouvert de $B$ est ouvert, la topologie de
$\Omega^1_{B/A}$ est la topologie $\mathfrak{n}$-préadique
\sref{0.20.4.5-fr}. Compte tenu de l'hypothèse que $B$ est un anneau adique,
il suffit donc, en vertu de ($\mathbf{0}_{\mathrm I}$, 7.2.7 et 7.2.9), de
prouver
\oldpage[0\textsubscript{IV-1}]{152}
que $\Omega^1_{B/A}/\mathfrak{n}.\Omega^1_{B/A}=
\Omega^1_{B/A}\otimes_B B_0$ est un $B_0$-module de type fini. Mais cela
résulte de l'hypothèse et de la suite exacte \sref{0.20.5.12.1-fr}
\[
  \mathfrak{n}/\mathfrak{n}^2\to
  \Omega^1_{B/A}\otimes_B B_0\to\Omega^1_{B_0/A}\to0.
\]
\end{proof}

\begin{env}[20.7.16]
\label{0.20.7.16-fr}
La proposition \sref{0.20.7.15-fr} s'applique par exemple lorsque $A$ est un
corps $k$, $B=k'[[T_1,\ldots,T_n]]$ une algèbre de séries formelles munie de
sa topologie usuelle, $k'$ une extension \emph{finie} de $k$ (cf.\
\sref{0.21.9.2-fr}). On notera que le corps des fractions $K$ de $B$ a un
degré de transcendance infini sur $k$ ; lorsque $k$ est de caractéristique
$0$, on en déduit aussitôt (à l'aide de \sref{0.20.4.13-fr}, (i) et de
\sref{0.20.5.9-fr} notamment, en utilisant aussi le fait que toute dérivation
d'un corps de caractéristique $0$ se prolonge à toute extension) que
$\Omega^1_{B/k}$ \emph{n'est pas un $B$-module de type fini}.
\end{env}

\begin{env}[20.7.17]
\label{0.20.7.17-fr}
Soient $A$, $B$, $C$ trois anneaux topologiques, $u:A\to B$, $v:B\to C$ deux
homomorphismes continus ; remplaçant $A$, $B$, $C$ par des quotients par des
idéaux ouverts $A'=A/\mathfrak{M}'$, $B'=B/\mathfrak{N}'$,
$C'=C/\mathfrak{P}'$, avec $u(\mathfrak{M}')\subset\mathfrak{N}'$,
$v(\mathfrak{N}')\subset\mathfrak{P}'$, de sorte que l'on a des
homomorphismes $u':A'\to B'$, $v':B'\to C'$, on en déduit des homomorphismes
canoniques $u'_{C'/B'/A'}$, $v'_{C'/B'/A'}$ qui, en vertu de
\sref{0.20.5.4-fr}, forment des systèmes projectifs, et donnent par suite,
par passage à la limite, des homomorphismes canoniques, prolongements aux
séparés complétés des homomorphismes de la suite exacte
\sref{0.20.5.7.1-fr}
\[
  \Omega^1_{B/A}\otimes_B C\xrightarrow{v_{C/B/A}}\Omega^1_{C/A}
  \xrightarrow{u_{C/B/A}}\Omega^1_{C/B}\to0:
\]
\begin{align}
  \widehat v_{C/B/A}:\widehat{\Omega}^1_{B/A}
  \widehat\otimes_{\widehat B}\widehat C&\to\widehat{\Omega}^1_{C/A},
  \tag{20.7.17.1}\label{0.20.7.17.1-fr}\\
  \widehat u_{C/B/A}:\widehat{\Omega}^1_{C/A}&\to
  \widehat{\Omega}^1_{C/B}. \tag{20.7.17.2}\label{0.20.7.17.2-fr}
\end{align}
et dans la suite
\begin{equation}
\tag{20.7.17.3}
\label{0.20.7.17.3-fr}
  \widehat{\Omega}^1_{B/A}\widehat\otimes_{\widehat B}\widehat C
  \xrightarrow{\widehat v_{C/B/A}}\widehat{\Omega}^1_{C/A}
  \xrightarrow{\widehat u_{C/B/A}}\widehat{\Omega}^1_{C/B}\to0
\end{equation}
le composé de deux homomorphismes consécutifs est $0$, mais la suite
\emph{n'est pas nécessairement exacte}. Toutefois, si $B$ et $C$ sont
\emph{métrisables}, l'homomorphisme $\widehat u_{C/B/A}$ est \emph{surjectif},
et $\operatorname{Im}(\widehat v_{C/B/A})$ est \emph{dense} dans
$\operatorname{Ker}(\widehat u_{C/B/A})$ : cela résulte aussitôt
(cf.\ ($\mathbf{0}_{\mathrm I}$, 7.3.1)) de ce que, si
$(\mathfrak{M}_k)$ (resp.\ $(\mathfrak{N}_k)$) est une suite décroissante
d'idéaux de $B$ (resp.\ $C$), formant un système fondamental de voisinages de
$0$, et si l'on pose $B_k=B/\mathfrak{M}_k$, $C_k=C/\mathfrak{N}_k$, les
homomorphismes de transition
$\Omega^1_{C_{k+1}/B_{k+1}}\to\Omega^1_{C_k/B_k}$,
$\Omega^1_{C_{k+1}/A}\to\Omega^1_{C_k/A}$ et
$\Omega^1_{B_{k+1}/A}\to\Omega^1_{B_k/A}$ sont surjectifs
\sref{0.20.7.14-fr}.
\end{env}

\begin{proposition}[20.7.18]
\label{0.20.7.18-fr}
Soient $A$, $B$, $C$ trois anneaux topologiques, $u:A\to B$, $v:B\to C$
deux homomorphismes continus. On suppose $B$ et $C$ admissibles
($\mathbf{0}_{\mathrm I}$, 7.1.2) et métrisables. Pour que l'homomorphisme
canonique
\[
  \widehat v_{C/B/A}:\widehat{\Omega}^1_{B/A}
  \widehat\otimes_{\widehat B}\widehat C\to\widehat{\Omega}^1_{C/A}
\]
admette un inverse à gauche qui soit un $C$-homomorphisme continu, il faut et
il suffit que $C$ soit une $B$-algèbre formellement lisse relativement à $A$.
\end{proposition}

\begin{proof}
La condition est nécessaire en vertu de \sref{0.20.7.2-fr} et
\sref{0.19.1.6-fr}. Pour voir qu'elle est suffisante, notons que le $C$-module
topologique
$L=\widehat{\Omega}^1_{B/A}\widehat\otimes_{\widehat B}\widehat C$ est
métrisable et
\oldpage[0\textsubscript{IV-1}]{153}
complet en vertu de l'hypothèse, donc $E=D_C(L)$, muni de la topologie produit,
est métrisable et complet ; il est en outre admissible, car si $\mathfrak{R}$
est un idéal de définition dans $C$, la suite des
$(\mathfrak{R}\oplus L)^n=\mathfrak{R}^n\oplus\mathfrak{R}^{n-1}L$ tend vers
$0$. Comme l'application composée
\[
  D:B\xrightarrow{d_{B/A}}\Omega^1_{B/A}\to
  \widehat{\Omega}^1_{B/A}\widehat\otimes_{\widehat B}\widehat C=L
\]
est une $A$-dérivation continue de $B$ dans $L$, le $A$-homomorphisme continu
\[
  f:x\longmapsto(v(x),D(x))
\]
de $B$ dans $E$ définit sur $E$ une structure de $B$-extension. Comme $L$ est
un idéal fermé dans $E$, il résulte de \sref{0.19.9.5-fr} et de l'hypothèse
que l'application identique $C\to E/L$ (qui est un $B$-homomorphisme) se
factorise en $C\xrightarrow{g}E\to E/L$, où $g$ est un homomorphisme continu
tel que $g\circ v=f$ ; par suite \sref{0.20.1.3-fr}, $g$ est de la forme
$y\mapsto(y,D'(y))$, où $D'$ est une $B$-dérivation continue de $C$ dans $L$,
autrement dit $D'\circ v=D$. Compte tenu de \sref{0.20.7.14.4-fr}, on a
$D'=h\circ\widehat d_{C/A}$, où
$h:\widehat{\Omega}^1_{C/A}\to L$ est un $C$-homomorphisme continu ; mais on
a $\widehat d_{C/A}\circ v=\widehat v_{C/B/A}\circ D$ par définition, et
comme l'image de $B$ par $D$ engendre (topologiquement) le $C$-module $L$
\sref{0.20.4.7-fr}, la relation $D'\circ v=D$ donne bien
$h\circ\widehat v_{C/B/A}=1_L$.
\end{proof}

\begin{corollary}[20.7.19]
\label{0.20.7.19-fr}
Sous les hypothèses de \sref{0.20.7.18-fr}, si l'on suppose en outre que dans
$B$ et $C$ le carré d'un idéal ouvert soit ouvert, alors, pour que $C$ soit
une $B$-algèbre formellement lisse relativement à $A$, il faut et il suffit
que $\widehat v_{C/B/A}$ soit inversible à gauche.
\end{corollary}

\begin{proof}
En effet, les topologies de
$\widehat{\Omega}^1_{B/A}\widehat\otimes_{\widehat B}\widehat C$ et de
$\widehat{\Omega}^1_{C/A}$ sont alors déduites de celle de $C$
\sref{0.20.4.5-fr}, et tout $C$-homomorphisme de l'un dans l'autre est
nécessairement continu.
\end{proof}

\begin{env}[20.7.20]
\label{0.20.7.20-fr}
Soient $A$ un anneau topologique, $B$ une $A$-algèbre topologique
\emph{métrisable et complète}, $\mathfrak{R}$ un idéal \emph{fermé} de $B$,
$C=B/\mathfrak{R}$ l'anneau topologique quotient, qui est métrisable et
complet. Soit $(\mathfrak{M}_k)$ un système fondamental décroissant de
voisinages de $0$ dans $B$ formé d'idéaux, et posons
$B_k=B/\mathfrak{M}_k$,
$\mathfrak{R}_k=(\mathfrak{R}+\mathfrak{M}_k)/\mathfrak{M}_k$,
$C_k=B_k/\mathfrak{R}_k$. On a un système projectif d'homomorphismes
\[
  \delta_{C_k/B_k/A}:\mathfrak{R}_k/\mathfrak{R}_k^2\to
  \Omega^1_{B_k/A}\otimes_{B_k}C_k
\]
\sref{0.20.5.11.3-fr}, d'où l'on déduit en passant à la limite un
homomorphisme canonique
\[
  \widehat\delta_{C/B/A}:\mathfrak{R}/\mathfrak{R}^2\to
  \widehat{\Omega}^1_{B/A}\widehat\otimes_{\widehat B}\widehat C
\]
et en raisonnant comme dans \sref{0.20.7.17-fr}, on voit que l'homomorphisme
canonique
$\widehat u_{C/B/A}:\widehat{\Omega}^1_{B/A}
\widehat\otimes_{\widehat B}\widehat C\to\widehat{\Omega}^1_{C/A}$ est
\emph{surjectif} et que $\operatorname{Im}(\widehat\delta_{C/B/A})$ est
\emph{dense} dans $\operatorname{Ker}(\widehat u_{C/B/A})$.
\end{env}

% FIN SOURCE ega4/ega0-20-fr.tex

%
% DEBUT SOURCE ega4/ega0-21-fr.tex
\section{Différentielles dans les anneaux de caractéristique p}
Les résultats du présent paragraphe, de nature plus spéciale et technique que
ceux des §§19, 20 et 22, ne serviront qu'exceptionnellement dans le cours du
chap. IV. Leur rôle principal est ici dans la démonstration de trois théorèmes
du §22 (22.3.3, 22.5.8 et 22.7.3), dont le premier et le dernier interviennent
de façon essentielle dans la théorie «fine» des anneaux locaux noethériens du
chap. IV, §7.
\oldpage[0\textsubscript{IV-1}]{154}
\subsection{Systèmes de $p$-générateurs et $p$-bases}
\label{0.21.1-fr}

\begin{env}[21.1.1]
\label{0.21.1.1-fr}
Étant donné un nombre $p$ qui est, soit $0$, soit un nombre premier, nous
dirons qu'un anneau $A$ est de caractéristique $p$ s'il existe un homomorphisme
d'anneaux $P\to A$, où $P$ est le corps premier de caractéristique $p$~; on
notera que cet homomorphisme est alors unique, le composé
$\mathbf Z\to P\to A$ étant l'unique homomorphisme de $\mathbf Z$ dans $A$.
Si $A\neq0$, il revient au même de dire que $A$ contient un corps de
caractéristique $p$, l'image de $P$ étant nécessairement un corps isomorphe à
$P$ (et d'ailleurs le seul sous-corps de $A$ isomorphe à $P$).
\end{env}

\begin{env}[21.1.2]
\label{0.21.1.2-fr}
Si $p>0$, dire que $A$ est de caractéristique $p$ équivaut à dire que, dans
$A$, on a $p\cdot1=0$, ou encore $pA=0$. Si $p=0$, dire que $A$ est de
caractéristique $p$ équivaut à dire que pour tout entier $n\neq0$, $n\cdot1$
est inversible dans $A$. Si $A\neq0$, il ne peut exister qu'un seul $p$
(premier ou $0$) tel que $A$ soit de caractéristique $p$~; cela résulte de
l'identité de Bézout $ap+bq=1$ pour deux nombres premiers distincts $p,q$.
Par contre, l'anneau réduit à $0$ est de caractéristique $p$ pour tout $p$.
\end{env}

\begin{env}[21.1.3]
\label{0.21.1.3-fr}
Si $A$ est de caractéristique $p$, il en est de même de toute algèbre sur
$A$. En particulier, pour tout idéal premier $\mathfrak p$ de $A$, le corps
résiduel de $A$ en $\mathfrak p$ est de caractéristique $p$. Réciproquement,
si $p=0$ et si, pour tout idéal maximal $\mathfrak m$ de $A$, le corps
résiduel de $A$ en $\mathfrak m$ est de caractéristique $0$, il en est de même
de $A$, car pour tout entier $n\neq0$, $n\cdot1$ est alors inversible dans
tous les $A_{\mathfrak m}$, donc aussi dans $A$. Par contre, si $p>0$, un
anneau local peut avoir son corps résiduel de caractéristique $p$ sans être
lui-même de caractéristique $p$, comme le montre l'exemple de l'anneau
premier $\mathbf Z_{p\mathbf Z}$ (19.8.3) ou de l'anneau local artinien
$\mathbf Z/p^n\mathbf Z$ pour $n\geq2$, qui ne contiennent pas de corps.

Notons enfin que, pour un anneau (même réduit), les corps résiduels en ses
idéaux premiers peuvent avoir des caractéristiques différentes, comme le
montre l'exemple de $\mathbf Z$.
\end{env}

\begin{env}[21.1.4]
\label{0.21.1.4-fr}
Dans toute la suite de ce paragraphe, nous supposons fixé un nombre premier
$p$ et tous les anneaux seront supposés de caractéristique $p$, sauf mention
expresse du contraire. Pour un tel anneau $A$, l'application $x\mapsto x^p$
est un endomorphisme de $A$, que l'on note $F_A$~; si $A$ est réduit, $F_A$
est injectif. On pose $A^p=F_A(A)$ (sous-anneau de $A$ formé des $x^p$,
$x\in A$)~; on peut naturellement considérer $A$ comme une $A^p$-algèbre.

On peut aussi considérer $A$ comme une $A$-algèbre au moyen de
l'homomorphisme $F_A:x\mapsto x^p$ de $A$ dans $A$~; autrement dit, il
s'agit de l'algèbre $A^{(p)}$ pour laquelle le produit $\lambda\cdot x$ de
$x\in A$ par un scalaire $\lambda\in A$ est le produit $\lambda^p x$ dans
l'anneau $A$. Il est clair que, pour tout homomorphisme d'anneaux $u:A\to B$,
le couple $(u,u)$ est un di-homomorphisme d'algèbres $A^{(p)}\to B^{(p)}$.
Pour tout $A$-module $E$, on pose $E^{(p)}=E\otimes_A A^{(p)}$, où
$A^{(p)}$ est considéré comme un $(A,A)$-bimodule, la structure de
$A$-module à gauche étant celle qu'on vient de définir et la structure de
$A$-module à droite définie par la multiplication dans $A$~; ainsi, pour
$x\in E$, $\alpha,\beta\in A$, on a
\[
  \alpha(x\otimes\beta)=x\otimes(\alpha\beta),\qquad
  (\alpha x)\otimes\beta=x\otimes\alpha^p\beta.
\]
\end{env}

\oldpage[0\textsubscript{IV-1}]{155}
Posant $x^{(p)}=x\otimes1$, on a donc $(\alpha x)^{(p)}=\alpha^p x^{(p)}$.
Lorsque $F_A$ est injectif, on peut identifier $A^{(p)}$ à l'anneau $A$
considéré comme algèbre sur son sous-anneau $A^p$.

\begin{proposition}[21.1.5]
\label{0.21.1.5-fr}
Soient $A$, $B$ deux anneaux, $u:A\to B$ un homomorphisme.
\begin{enumerate}
  \item[(i)] Toute $A$-dérivation de $B$ dans un $B$-module $L$ est nulle
  dans $A[B^p]$ (donc est une $A[B^p]$-dérivation).
  \item[(ii)] Pour toute sous-$A$-algèbre $A'$ de $A[B^p]$, si
  $j:A'\to B$ est l'injection canonique, l'homomorphisme canonique
  \[
  j_{B/A'/A}:\Omega^1_{B/A}\longrightarrow\Omega^1_{B/A'}
  \]
  est bijectif.
  \item[(iii)] Supposons qu'il existe un entier $s\geq0$ tel que
  $u(A)\supset B^{p^s}$. Alors, dans l'anneau $B\otimes_A B$,
  $\mathfrak J_{B/A}$ est un nilidéal.
\end{enumerate}
\end{proposition}

\begin{proof}
(i) Par récurrence à partir de (20.1.1.1), on déduit, pour toute
$A$-dérivation $D:B\to L$, que $D(x^n)=nx^{n-1}D(x)$, d'où en particulier
$D(x^p)=0$.

(ii) Vu (20.4.8), l'assertion (ii) n'est qu'une traduction de (i), cette
dernière s'écrivant $\operatorname{Der}_{A[B^p]}(B,L)=\operatorname{Der}_A(B,L)$
pour tout $B$-module $L$.

(iii) Pour tout $x\in B$, on a
$(x\otimes1-1\otimes x)^{p^s}=x^{p^s}\otimes1-1\otimes x^{p^s}=0$,
puisque $x^{p^s}\in u(A)$. La conclusion résulte de ce que les éléments
$1\otimes x-x\otimes1$ engendrent $\mathfrak J_{B/A}$ (20.4.4).
\end{proof}

Il résulte aussitôt de (21.1.5) que, pour tout couple d'homomorphismes
d'anneaux $A\to B\to C$, on a
\begin{equation}
\tag{21.1.5.1}
\label{0.21.1.5.1-fr}
\Upsilon_{C/B/A}=\Upsilon_{C/B/A'}
\end{equation}
pour toute sous-$A$-algèbre $A'\subset A[B^p]$ de $B$.

D'autre part, (21.1.5) montre en particulier que, pour les modules de
différentielles «absolues»,
\begin{equation}
\tag{21.1.5.2}
\label{0.21.1.5.2-fr}
\Omega^1_A=\Omega^1_{A/A^p}.
\end{equation}

\begin{corollary}[21.1.6]
\label{0.21.1.6-fr}
Supposons que $B$ soit une $A$-algèbre de type fini et qu'il existe un
entier $s\geq0$ tel que $u(A)\supset B^{p^s}$. Si $\Omega^1_{B/A}=0$,
$u$ est surjectif.
\end{corollary}

\begin{proof}
\oldpage[0\textsubscript{IV-1}]{156}
D'après (21.1.5, (iii)), $\mathfrak J_{B/A}$ est un nilidéal~; en outre,
en vertu de (20.4.4) et de l'hypothèse, $\mathfrak J_{B/A}$ est un idéal de
type fini, donc il est nilpotent~; comme la relation $\Omega^1_{B/A}=0$
signifie que $\mathfrak J_{B/A}=\mathfrak J_{B/A}^2$, on en conclut que
$\mathfrak J_{B/A}=0$, ou encore que $\pi:B\otimes_A B\to B$ est bijectif.
Par ailleurs, tout élément de $B$ ayant sa puissance $p^s$-ième dans $u(A)$,
$B$ est entier sur $A$, et étant de type fini, c'est une $A$-algèbre finie.
On est ainsi ramené à prouver le lemme suivant (où l'on ne suppose pas que
les anneaux considérés sont de caractéristique $p$) :
\end{proof}

\begin{lemma}[21.1.6.1]
\label{0.21.1.6.1-fr}
Soient $R$ un anneau, $S$ une $R$-algèbre finie~; si l'homomorphisme
canonique $\pi:S\otimes_R S\to S$ est bijectif, alors l'homomorphisme
structural $u:R\to S$ est surjectif.
\end{lemma}

\begin{proof}
Il suffit de montrer que, pour tout idéal maximal $\mathfrak m$ de $R$, si
l'on pose $T=R-\mathfrak m$, l'homomorphisme
$u_{\mathfrak m}:R_{\mathfrak m}\to T^{-1}S$ est surjectif (Bourbaki,
\emph{Alg. comm.}, chap.~II, §~3, n$^\circ$~3, th.~1)~; or l'hypothèse
entraîne que l'homomorphisme
$(T^{-1}S)\otimes_{R_{\mathfrak m}}(S/\mathfrak mS)\to T^{-1}S$
est bijectif (0$_I$, 1.3.4), et comme $T^{-1}S$ est une
$R_{\mathfrak m}$-algèbre finie, on voit qu'on peut se borner au cas où
$R$ est un anneau local.

Désignant encore par $\mathfrak m$ son idéal maximal, il suffit de prouver que
$u\otimes1:R/\mathfrak m\to S/\mathfrak mS$ est surjectif, en vertu du
lemme de Nakayama ($S$ étant un $R$-module de type fini)~; comme
l'homomorphisme canonique
$(S/\mathfrak mS)\otimes_{R/\mathfrak m}(S/\mathfrak mS)\to S/\mathfrak mS$
est bijectif et que $S/\mathfrak mS$ est une $(R/\mathfrak m)$-algèbre finie,
on est finalement ramené au cas où $R$ est un corps~; mais alors les rangs de
$S\otimes_R S$ et de $S$ sur $R$ ne peuvent être égaux que si $S$ est de rang
$0$ ou $1$, donc si $u:R\to S$ est surjectif.
\end{proof}

\begin{proposition}[21.1.7]
\label{0.21.1.7-fr}
Soient $A$ un anneau, $B$ une $A$-algèbre, $(x_\alpha)_{\alpha\in I}$ une
famille d'éléments de $B$. Considérons les propriétés suivantes~:
\begin{enumerate}
  \item[a)] $B=A[B^p,(x_\alpha)]$, autrement dit la $A[B^p]$-algèbre
  $B$ est engendrée par la famille $(x_\alpha)_{\alpha\in I}$.
  \item[b)] Le $A[B^p]$-module $B$ est engendré par les monômes
  $\prod_\alpha x_\alpha^{n(\alpha)}$, où $(n(\alpha))_{\alpha\in I}$ est
  une famille d'entiers à support fini telle que $0\leq n(\alpha)<p$ pour
  tout $\alpha\in I$.
  \item[c)] Le $B$-module $\Omega^1_{B/A}$ est engendré par les
  $d_{B/A}(x_\alpha)$ $(\alpha\in I)$.
\end{enumerate}
Alors les propriétés a) et b) sont équivalentes et entraînent c)~; si, en
outre, $B$ est une $A[B^p]$-algèbre de type fini, c) est équivalente à a) et b).
\end{proposition}

\begin{proof}
Il est clair que b) entraîne a), et inversement a) entraîne b), car tout
monôme $\prod_\alpha x_\alpha^{m(\alpha)}$, où les $m(\alpha)$ sont des
entiers $\geq0$ (formant une famille à support fini), peut s'écrire
$(\prod_\alpha x_\alpha^{q(\alpha)})^p\prod_\alpha x_\alpha^{r(\alpha)}$,
en divisant chaque $m(\alpha)$ par $p$, ce qui donne
$m(\alpha)=p q(\alpha)+r(\alpha)$ avec $0\leq r(\alpha)<p$.
Le fait que a) entraîne c) résulte de (21.1.5, (ii)) et de (20.4.7).

Supposons inversement c) vérifiée et que $B$ soit une $A[B^p]$-algèbre de
type fini~; soit $B'$ la sous-$A[B^p]$-algèbre de $B$ engendrée par les
$x_\alpha$~; dans la suite exacte (20.5.7.1)
\[
\Omega^1_{B'/A[B^p]}\otimes_{B'}B\longrightarrow
\Omega^1_{B/A[B^p]}\longrightarrow\Omega^1_{B/B'}\longrightarrow0
\]
l'hypothèse entraîne que la flèche de gauche est surjective (compte tenu de
(21.1.5, (ii)))~; on a donc $\Omega^1_{B/B'}=0$, et comme
$B^p\subset B'\subset B$ et que $B$ est une $B'$-algèbre de type fini,
on a nécessairement $B'=B$ par (21.1.6).
\end{proof}

\begin{remark}[21.1.8]
\label{0.21.1.8-fr}
Lorsque $B$ est un corps, nous démontrerons l'équivalence des propriétés a),
b) et c) sans hypothèse de finitude (21.4.5).
\end{remark}

\begin{definition}[21.1.9]
\label{0.21.1.9-fr}
Soient $A$ un anneau (de caractéristique $p$), $B$ une $A$-algèbre. On dit
qu'une famille $(x_\alpha)_{\alpha\in I}$ d'éléments de $B$ est $p$-libre sur
$A$ (resp. un système de $p$-générateurs de $B$ sur $A$, resp. une $p$-base
de $B$ sur $A$) si la famille des monômes
$\prod_\alpha x_\alpha^{n(\alpha)}$ ($0\leq n(\alpha)<p$,
$(n(\alpha))_{\alpha\in I}$ à support fini) est une famille libre (resp. un
système de générateurs, resp. une base) dans le $A[B^p]$-module $B$.
\end{definition}

On a des définitions correspondantes pour un ensemble $M$ d'éléments de $B$,
en considérant la famille définie par l'injection canonique $M\to B$.
Lorsqu'on prend pour $A$ le corps premier $\mathbf F_p$ (auquel cas
$A[B^p]=B^p$), on omet la mention de $A$ dans la définition précédente (ou
l'on dit encore qu'une famille est «absolument» $p$-libre, resp. un système
«absolu» de $p$-générateurs, resp. une $p$-base «absolue»).

\oldpage[0\textsubscript{IV-1}]{157}
\begin{lemma}[21.1.10]
\label{0.21.1.10-fr}
Soient $C$ une sous-$A$-algèbre de $B$ telle que $B^p\subset C$, $M$ une
partie de $C$, $N$ une partie de $B$.
\begin{enumerate}
  \item[(i)] Si $M$ est un système de $p$-générateurs de $C$ sur $A$ et $N$
  un système de $p$-générateurs de $B$ sur $C$, alors $M\cup N$ est un système
  de $p$-générateurs de $B$ sur $A$.
  \item[(ii)] Supposons que $M$ soit une $p$-base de $C$ sur $A$~; alors,
  pour que $N$ soit $p$-libre sur $C$, il faut et il suffit que $M\cup N$
  soit $p$-libre sur $A$.
\end{enumerate}
\end{lemma}

\begin{proof}
On peut se borner au cas où $B^p\subset A\subset C\subset B$. Alors (i)
est un cas particulier du fait que si $P$ (resp. $Q$) est un système de
générateurs du $A$-module $C$ (resp. du $C$-module $B$), l'ensemble des
produits $xy$, où $x\in P$ et $y\in Q$, est un système de générateurs du
$A$-module $B$.

Conservant les mêmes notations, si $P$ est une base du $A$-module $C$,
dire que $Q$ est une famille libre sur $C$ signifie que la relation
$\sum_{\lambda,\mu}a_{\lambda\mu}x_\lambda y_\mu=0$, où
$a_{\lambda\mu}\in A$, $x_\lambda\in P$, $y_\mu\in Q$ (les $x_\lambda$
(resp. les $y_\mu$) étant deux à deux distincts) équivaut à
$a_{\lambda\mu}=0$ pour tout couple $(\lambda,\mu)$, ou encore à
$\sum_\lambda a_{\lambda\mu}x_\lambda=0$ pour tout $\mu$~; d'où
l'assertion (ii).
\end{proof}

\subsection{$p$-bases et lissité formelle}
\label{0.21.2-fr}

\begin{theorem}[21.2.1]
\label{0.21.2.1-fr}
Soient $B$ un anneau, $A$ un sous-anneau de $B$ tel que $B^p\subset A\subset B$,
$(x_\alpha)_{\alpha\in I}$ une $p$-base de $B$ sur $A$.
Soient $E$ une $A$-algèbre, $u:A\to E$ l'homomorphisme structural,
$q:E\to B$ un $A$-homomorphisme surjectif, $\mathfrak R$ son noyau, et
supposons que $\mathfrak R^p=0$. Alors~:
\begin{enumerate}
  \item[(i)] Pour qu'il existe un $A$-homomorphisme $v:B\to E$ inverse à
  droite de l'homomorphisme $q:E\to B$, il faut et il suffit que
  $u(q(z)^p)=z^p$ pour tout $z\in E$.
  \item[(ii)] Lorsque la condition de (i) est satisfaite, pour toute famille
  $(z_\alpha)_{\alpha\in I}$ d'éléments de $E$ telle que
  $q(z_\alpha)=x_\alpha$ pour tout $\alpha\in I$, il existe un
  $A$-homomorphisme $v:B\to E$ et un seul tel que $v(x_\alpha)=z_\alpha$
  pour tout $\alpha\in I$, et $v$ est inverse à droite de $q$.
\end{enumerate}
\end{theorem}

\begin{proof}
S'il existe un $A$-homomorphisme $v:B\to E$, on doit avoir $v(a)=u(a)$ pour
tout $a\in A\subset B$, donc $v(q(z)^p)=u(q(z)^p)$ pour tout $z\in E$,
puisque $B^p\subset A$. Mais on a $v(q(z)^p)=(v(q(z)))^p$ et, par définition,
$v(q(z))\equiv z\pmod{\mathfrak R}$ si $q\circ v=1_B$, donc
$(v(q(z)))^p=z^p$ puisque $\mathfrak R^p=0$~; d'où la nécessité de (i).
La suffisance de la condition (i) résultera de (ii).

Or, sous les hypothèses de (ii), l'unicité de $v$ est évidente puisque les
monômes $\prod_\alpha x_\alpha^{n(\alpha)}$ engendrent le $A$-module $B$~;
comme ces monômes forment en outre une \emph{base} du $A$-module $B$, il existe
une application $A$-\emph{linéaire} $v$ de $B$ dans $E$ telle que
$v(\prod_\alpha x_\alpha^{n(\alpha)})=\prod_\alpha z_\alpha^{n(\alpha)}$
pour toute famille $(n(\alpha))_{\alpha\in I}$ à support fini avec
$0\leq n(\alpha)<p$ pour tout $\alpha$. Reste à voir que $v$ est un
homomorphisme d'anneaux. Or, on peut écrire
$(\prod_\alpha x_\alpha^{m(\alpha)})(\prod_\alpha x_\alpha^{n(\alpha)})
=a\prod_\alpha x_\alpha^{r(\alpha)}$, où
$r(\alpha)=m(\alpha)+n(\alpha)$ si $m(\alpha)+n(\alpha)<p$,
$r(\alpha)=m(\alpha)+n(\alpha)-p$ dans le cas contraire, et $a\in A$
\oldpage[0\textsubscript{IV-1}]{158}
est le produit des $x_\alpha^p$ pour les $\alpha\in I$ tels que
$m(\alpha)+n(\alpha)\geq p$~; il s'agit de voir que
$u(a)=\prod_\alpha z_\alpha^p$~; mais comme $x_\alpha=q(z_\alpha)$, cela
résulte de l'hypothèse.
\end{proof}

\begin{corollary}[21.2.2]
\label{0.21.2.2-fr}
Soient $B$ un anneau, $A$ un sous-anneau de $B$ tel que $B^p\subset A\subset B$.
S'il existe une $p$-base de $B$ sur $A$, $B$ est une $A$-algèbre formellement
lisse relativement à $B^p$ (pour les topologies discrètes).
\end{corollary}

\begin{proof}
Soient $E$ une $A$-extension de $B$ par un $B$-module $L$, $q:E\to B$
l'augmentation, $u:A\to E$ l'homomorphisme structural. Dire que $E$ est
$B^p$-triviale signifie qu'il existe un homomorphisme d'anneaux $v:B\to E$
tel que $q(v(b))=b$ et $v(b^p)=u(b^p)$ pour tout $b\in B$. On en déduit,
pour $z\in E$, que $u(q(z)^p)=v(q(z)^p)=(v(q(z)))^p$~; mais, en vertu de
la relation $L^2=0$, on a aussi $L^p=0$, et comme $v(q(z))-z\in L$, on a
$(v(q(z)))^p=z^p$. La condition de (21.2.1), (i), est donc satisfaite et
$E$ est $B^p$-triviale.
\end{proof}

\begin{corollary}[21.2.3]
\label{0.21.2.3-fr}
Soient $B$ un anneau, $A$ un sous-anneau de $B$ tel que $B^p\subset A\subset B$,
et $\left(x_\alpha\right)_{\alpha\in I}$ une $p$-base de $B$ sur $A$. Soit
$L$ un $B$-module. Alors~:
\begin{enumerate}
  \item[(i)] Pour qu'une dérivation $D$ de $A$ dans $L$ se prolonge en une
  dérivation de $B$ dans $L$, il faut et il suffit que $D$ s'annule dans $B^p$.
  \item[(ii)] Si $D$ s'annule dans $B^p$, alors, pour toute famille
  $(y_\alpha)_{\alpha\in I}$ d'éléments de $L$, il existe une dérivation
  $D'$ de $B$ dans $L$ et une seule prolongeant $D$ et telle que
  $D'(x_\alpha)=y_\alpha$ pour tout $\alpha\in I$.
\end{enumerate}
\end{corollary}

\begin{proof}
Étant donnée une dérivation $D$ de $A$ dans $L$, considérons l'anneau
$E=D_B(L)$ et l'homomorphisme $u:A\to E$ défini par $u(a)=(a,D(a))$~; un
$A$-homomorphisme $v:B\to E$ inverse à droite de l'homomorphisme canonique
$q:E\to B$ est alors de la forme $x\mapsto(x,D'(x))$, où $D'$ est une
dérivation de $B$ dans $L$ prolongeant $D$ (20.1.5)~; comme
$u(q(z)^p)=(q(z)^p,D(q(z)^p))$ pour $z\in E$, le corollaire résulte aussitôt
de (21.2.1) appliqué à $E$.
\end{proof}

Il résulte de (21.2.3) que la suite
\begin{equation}
\tag{21.2.3.1}
\label{0.21.2.3.1-fr}
0\to\operatorname{Der}_A(B,L)\to\operatorname{Der}_{B^p}(B,L)
\to\operatorname{Der}_{B^p}(A,L)\to0
\end{equation}
(cf. (20.2.2.1)) est exacte, et que l'application
\begin{equation}
\tag{21.2.3.2}
\label{0.21.2.3.2-fr}
D'\longmapsto(D'(x_\alpha))_{\alpha\in I}
\end{equation}
est un isomorphisme du $B$-module $\operatorname{Der}_A(B,L)$ sur le
$B$-module produit $L^I$ (en faisant $D=0$ dans (21.2.3)).

\begin{corollary}[21.2.4]
\label{0.21.2.4-fr}
Sous les hypothèses de (21.2.3), la suite
\begin{equation}
\tag{21.2.4.1}
\label{0.21.2.4.1-fr}
0\to\Omega^1_{A/B^p}\otimes_A B\to\Omega^1_{B/B^p}
\to\Omega^1_{B/A}\to0
\end{equation}
est exacte et scindée, et la famille $(d_{B/A}(x_\alpha))_{\alpha\in I}$ est
une base du $B$-module $\Omega^1_{B/A}$.
\end{corollary}

\begin{proof}
Cela résulte aussitôt de (21.2.3) et de la formule
$\operatorname{Der}_A(B,L)=\operatorname{Hom}_B(\Omega^1_{B/A},L)$ (20.4.8).
\end{proof}

\begin{corollary}[21.2.5]
\label{0.21.2.5-fr}
Soient $A$ un anneau, $B$ une $A$-algèbre, $(x_\alpha)_{\alpha\in I}$ une
$p$-base de $B$ sur $A$. Alors $(d_{B/A}(x_\alpha))_{\alpha\in I}$ est une
base du $B$-module $\Omega^1_{B/A}$.
\end{corollary}

\begin{proof}
En utilisant (21.1.5), (ii), on se ramène au cas où $A=A[B^p]$ et il suffit
alors d'appliquer (21.2.4).
\end{proof}

\oldpage[0\textsubscript{IV-1}]{159}
\begin{env}[21.2.6]
\label{0.21.2.6-fr}
Soient $A$, $B$ deux anneaux, $u:A\to B$ un homomorphisme d'anneaux~; on a
(avec les notations de (21.1.4)) le diagramme commutatif
\begin{equation}
\tag{21.2.6.1}
\label{0.21.2.6.1-fr}
\xymatrix{
 A^{(p)} \ar[r]^{u} & B^{(p)}\\
 A \ar[u]^{F_A} \ar[r]^{u} & B \ar[u]_{F_B}
}
\end{equation}
On en déduit canoniquement un homomorphisme de $B$-algèbres
\begin{equation}
\tag{21.2.6.2}
\label{0.21.2.6.2-fr}
u\otimes1:A^{(p)}\otimes_A B\longrightarrow B^{(p)}
\end{equation}
dont l'image est l'anneau $A[B^p]$ (qui est une $B$-algèbre pour
l'homomorphisme $F_B:B\to B^p$).
\end{env}

\begin{theorem}[21.2.7]
\label{0.21.2.7-fr}
Soient $A$ un anneau, $B$ une $A$-algèbre, $u:A\to B$ l'homomorphisme
structural. On suppose vérifiées les conditions suivantes~:
\begin{enumerate}
  \item[(i)] L'homomorphisme (21.2.6.2) déduit de $u$ est injectif.
  \item[(ii)] $B$ admet une $p$-base $(x_\alpha)_{\alpha\in I}$ par rapport à $A$.
\end{enumerate}
Alors $B$ est une $A$-algèbre formellement lisse (pour les topologies
discrètes). De façon plus précise, munissons $A$ et $B$ des topologies
discrètes~; soient $E$ une $A$-algèbre topologique admissible
(0$_I$, 7.1.2), $\mathfrak R$ un idéal de définition de $E$, $C=E/\mathfrak R$,
$v:B\to C$ un $A$-homomorphisme, $q:E\to C$ l'augmentation. Alors, pour
toute famille $(z_\alpha)_{\alpha\in I}$ d'éléments de $E$ telle que
$q(z_\alpha)=v(x_\alpha)$ pour tout $\alpha\in I$, il existe un
$A$-homomorphisme $w:B\to E$ et un seul tel que $v=q\circ w$ et
$w(x_\alpha)=z_\alpha$ pour tout $\alpha\in I$.
\end{theorem}

\begin{proof}
Considérons d'abord le cas où $E$ est discret, donc $\mathfrak R$ nilpotent.
On peut se borner au cas où $v$ est surjectif ; en outre, le raisonnement de
(19.4.3) permet de supposer $\mathfrak R^2=0$ et par suite $\mathfrak R^p=0$.
Enfin, en considérant l'image réciproque par $v$ de l'extension $E$ de $C$
par $\mathfrak R$, on peut se borner au cas où $C=B$ et où $v$ est l'identité
(19.4.4). Comme l'homomorphisme d'anneaux $F_E:z\rightsquigarrow z^p$
s'annule dans $\mathfrak R$, il se factorise en
$E\xrightarrow q B\xrightarrow{F'}E$,
et $F'$, considéré comme homomorphisme de $B$ dans $E^{(p)}$, est un
\emph{$A$-homomorphisme} par définition de la structure de $A$-algèbre de $E^{(p)}$
à l'aide de l'homomorphisme composé
$A\xrightarrow r E\xrightarrow{F_E}E^{(p)}$, où $r:A\to E$ est l'homomorphisme
structural de la $A$-algèbre $E$ ; d'ailleurs
$r:A^{(p)}\to E^{(p)}$ est aussi un $A$-homomorphisme. Il y a donc un unique
$A$-homomorphisme $f:A^{(p)}\otimes_A B\to E^{(p)}$ tel que les composés de $f$
et des homomorphismes canoniques $A^{(p)}\to A^{(p)}\otimes_A B$ et
$B\to A^{(p)}\otimes_A B$ soient respectivement $r$ et $F'$. Or, par hypothèse,
on peut identifier $A^{(p)}\otimes_A B$ à $A[B^p]$, les homomorphismes canoniques
$A^{(p)}\to A^{(p)}\otimes_A B$ et $B\to A^{(p)}\otimes_A B$ s'identifiant
respectivement à $u$ et à $F_B$. On peut donc maintenant considérer $E$ comme
une $A[B^p]$-algèbre au moyen de l'homomorphisme $f$, et, par construction, on a
$f(q(z)^p)=z^p$ pour tout $z\in E$ ; on est donc dans les conditions
d'application de (21.2.1), d'où le théorème dans ce cas.

Pour passer au cas général, considérons un système fondamental
$(\mathfrak J_\lambda)$ d'idéaux ouverts de $E$, et posons
$E_\lambda=E/\mathfrak J_\lambda$,
$\mathfrak R_\lambda=(\mathfrak R+\mathfrak J_\lambda)/\mathfrak J_\lambda$,
$C_\lambda=E_\lambda/\mathfrak R_\lambda=E/(\mathfrak R+\mathfrak J_\lambda)$ ;
comme $E$ est admissible, on a $E=\varprojlim E_\lambda$ (0$_I$, 7.2.2).
Pour tout couple $(\alpha,\lambda)$, désignons par $z_{\alpha\lambda}$
l'image canonique de $z_\alpha$ dans $E_\lambda$, par
$p_\lambda:C\to C_\lambda$ l'homomorphisme canonique,
\oldpage[0\textsubscript{IV-1}]{160}
par $q_\lambda$ l'homomorphisme canonique $E_\lambda\to C_\lambda$.
Comme par hypothèse $\mathfrak R_\lambda$ est nilpotent dans $E_\lambda$,
la première partie de la démonstration prouve l'existence et l'unicité d'un
$A$-homomorphisme $w_\lambda:B\to E_\lambda$ tel que
$p_\lambda\circ v=q_\lambda\circ w_\lambda$ et
$w_\lambda(x_\alpha)=z_{\alpha\lambda}$ pour tout $\alpha$.
L'unicité des $w_\lambda$ montre alors aussitôt que $(w_\lambda)$ est un
système projectif d'homomorphismes, et $w=\varprojlim w_\lambda$ répond
à la question.
\end{proof}

\begin{remark}[21.2.8]
\label{0.21.2.8-fr}
La vérification de l'existence et de l'unicité de l'homomorphisme $w:B\to E$
tel que $w(x_\alpha)=z_\alpha$ pour tout $\alpha$ peut se déduire directement
du fait que $B$ est une $A$-algèbre formellement lisse et que les
$d_{B/A}(x_\alpha)$ forment une base du $B$-module $\Omega^1_{B/A}$ (21.2.4),
sans faire intervenir le fait qu'il s'agit d'une $p$-base (de sorte que le
résultat est valable sans supposer les anneaux de caractéristique $p$).
En effet, on peut se borner au cas où $\mathfrak R^2=0$~; comme
$\operatorname{Der}_A(B,\mathfrak R)=\operatorname{Hom}_B(\Omega^1_{B/A},\mathfrak R)$
(20.4.8), il existe une $A$-dérivation $D$, et une seule, de $B$ dans
$\mathfrak R$ telle que $D(x_\alpha)=z'_\alpha$ pour toute famille
$(z'_\alpha)$ d'éléments de $\mathfrak R$~; la conclusion résulte de (20.1.1).
\end{remark}

\subsection{$p$-bases et modules d'imperfection}
\label{0.21.3-fr}

\begin{env}[21.3.1]
\label{0.21.3.1-fr}
Soient $A$, $B$ deux anneaux (de caractéristique $p$), $i:A\to B$,
$j:B\to A$ deux homomorphismes d'anneaux tels que
\begin{align}
\tag{21.3.1.1}\label{0.21.3.1.1-fr}
j(i(a))&=a^p\quad\text{pour tout }a\in A,\\
\tag{21.3.1.2}\label{0.21.3.1.2-fr}
i(j(b))&=b^p\quad\text{pour tout }b\in B.
\end{align}
Le plus souvent, $i$ sera injectif, de sorte que $A$ s'identifiera par
$i:A\to B$ à un sous-anneau de $B$~; une fois cette identification faite,
l'existence de $j$ implique que $B^p\subset A$, et $j$ s'identifie alors à $F_B$.
\end{env}

\begin{env}[21.3.2]
\label{0.21.3.2-fr}
Si $h:A^p\to A$ est l'injection canonique, on a, d'après (21.3.1.1) et
(21.3.1.2), un diagramme commutatif
\begin{equation}
\tag{21.3.2.1}\label{0.21.3.2.1-fr}
\xymatrix{
 B \ar[r]^{j} & A\\
 A \ar[u]^{i} \ar[r]^{F_A} & A^p \ar[u]_{h}
}
\end{equation}
si bien que le couple $(j,F_A)$ peut être considéré comme un di-homomorphisme
de la $A$-algèbre $B$ (pour $i$) dans la $A^p$-algèbre $A$ (pour $h$).
On en déduit un homomorphisme canonique de $A$-modules
\begin{equation}
\tag{21.3.2.2}\label{0.21.3.2.2-fr}
\pi_{B/A}:\Omega^1_{B/A}\otimes_B A_{[j]}\longrightarrow
\Omega^1_{A/A^p}=\Omega^1_A
\end{equation}
(cf. (20.5.4)~; l'identification de $\Omega^1_{A/A^p}$ et du module des
différentielles absolues $\Omega^1_A$ provient de (21.1.5), (ii)).
\begin{align}
\tag{21.3.2.3}\label{0.21.3.2.3-fr}
\Theta_{B/A}&=\Omega^1_{B/A}\otimes_B A_{[j]},\\
\tag{21.3.2.4}\label{0.21.3.2.4-fr}
\Xi_{B/A}&=\operatorname{Ker}(\pi_{B/A}).
\end{align}
\oldpage[0\textsubscript{IV-1}]{161}
De sorte que l'on a la suite exacte
\begin{equation}
\tag{21.3.2.5}\label{0.21.3.2.5-fr}
0\to\Xi_{B/A}\to\Theta_{B/A}\xrightarrow{\pi_{B/A}}\Omega^1_A
\end{equation}
(on remarquera que ces notations peuvent prêter à confusion puisque
$\Theta_{B/A}$ et $\Xi_{B/A}$ dépendent non seulement de $A$ et de $B$,
mais aussi de $i$ et de $j$).
\end{env}

\begin{env}[21.3.3]
\label{0.21.3.3-fr}
Comme, en vertu de (21.3.1.2), on a $F_B=i\circ j$, on peut écrire pour
tout $B$-module $M$ (cf. (21.1.4))
\begin{equation}
\tag{21.3.3.1}\label{0.21.3.3.1-fr}
M^{(p)}=M\otimes_B B^{(p)}=(M\otimes_B A_{[j]})\otimes_A B_{[i]},
\end{equation}
d'où en particulier
\begin{equation}
\tag{21.3.3.2}\label{0.21.3.3.2-fr}
(\Omega^1_{B/A})^{(p)}=\Theta_{B/A}\otimes_A B_{[i]},
\end{equation}
et l'on déduit donc de (21.3.2.2) un homomorphisme canonique de $B$-modules
\begin{equation}
\tag{21.3.3.3}\label{0.21.3.3.3-fr}
\pi_{B/A}\otimes1_B:(\Omega^1_{B/A})^{(p)}\longrightarrow
\Omega^1_A\otimes_A B_{[i]}.
\end{equation}
Compte tenu de (20.5.4.2), l'image de cet homomorphisme est le $B$-module
engendré par les éléments $d_A(j(b))\otimes1$ pour $b\in B$. Leur image
par l'homomorphisme canonique $i_{B/A}:\Omega^1_A\otimes_A B\to\Omega^1_B$
déduit de $i$ (20.5.2) est donc dans le sous-$B$-module de $\Omega^1_B$
engendré par les $d_B(i(j(b)))$ pour $b\in B$~; en vertu de (21.3.1.2),
cette image est nulle. Autrement dit, on a une suite d'homomorphismes
\begin{equation}
\tag{21.3.3.4}\label{0.21.3.3.4-fr}
0\to\Xi_{B/A}\otimes_A B_{[i]}\to(\Omega^1_{B/A})^{(p)}
\to\Upsilon_{B/A}
\end{equation}
qui n'est pas nécessairement exacte, mais où le composé de deux
homomorphismes consécutifs est nul.
\end{env}

\begin{proposition}[21.3.4]
\label{0.21.3.4-fr}
Si $B$ est un $A$-module plat (pour $i$), la suite (21.3.3.4) est exacte.
\end{proposition}
\begin{proof}
Cela résulte de la définition de la platitude. La proposition s'applique en
particulier lorsque $B^p\subset A\subset B$, $i$ et $j$ étant respectivement
l'injection canonique et $F_B$, et que $B$ admet une $p$-base sur $A$ (de sorte
que $B$ est alors un $A$-module libre). Mais même dans ce cas le noyau
$\Xi_{B/A}$ n'est pas nécessairement nul. Toutefois~:
\end{proof}

\begin{proposition}[21.3.5]
\label{0.21.3.5-fr}
Soient $B$ un anneau réduit, $A$ un sous-anneau de $B$ tel que
$B^p\subset A\subset B$. Supposons qu'il existe une $p$-base
$(x_\lambda)_{\lambda\in L}$ de $B$ sur $A$ et une $p$-base
$(y_\mu)_{\mu\in M}$ de $A$ sur $B^p$~; alors l'homomorphisme canonique
(21.3.3.4)
\begin{equation}
\tag{21.3.5.1}\label{0.21.3.5.1-fr}
(\Omega^1_{B/A})^{(p)}\longrightarrow\Upsilon_{B/A}
\end{equation}
est bijectif, et les éléments $d_A(x_\lambda^p)\otimes1$ forment une base
du $B$-module $\Upsilon_{B/A}$.
\end{proposition}

\begin{proof}
\oldpage[0\textsubscript{IV-1}]{162}
Comme $B$ est réduit, $x\mapsto x^p$ est un isomorphisme de structure de
$B$ sur $B^p$~; par transport de structure au moyen de cet isomorphisme,
on voit que $(x_\lambda^p)_{\lambda\in L}$ et les $y_\mu$ forment une $p$-base
de $B^p$ sur $A^p$ (21.1.10)~; on en conclut par (21.1.5), (ii), et (21.2.5)
que les $d_A(x_\lambda^p)$ et les $d_A(y_\mu)$ forment une base du
$A$-module $\Omega^1_A$. Or l'image de $d_A(x_\lambda^p)\otimes1$ par
$i_{B/A}$ est $d_B(x_\lambda^p)=0$~; d'autre part, l'image de $d_A(y_\mu)$
par $i_{B/A}$ est $d_B(i(y_\mu))$, et comme les $y_\mu$ et les $x_\lambda$
forment une $p$-base de $B$ sur $B^p$ (21.1.10), les $d_B(y_\mu)$ sont une
partie d'une base de $\Omega^1_B$ (21.2.5), donc sont linéairement indépendants
sur $B$. On en déduit que le noyau de $i_{B/A}$ a pour base les
$d_A(x_\lambda^p)\otimes1$~; comme ce sont les images par (21.3.5.1) des
éléments correspondants de $(\Omega^1_{B/A})^{(p)}$, et que ces derniers
forment une base du $B$-module $(\Omega^1_{B/A})^{(p)}$ (21.2.5),
l'homomorphisme (21.3.5.1) de $B$-modules est bijectif.
\end{proof}

\begin{remarks}[21.3.6]
\label{0.21.3.6-fr}
\begin{enumerate}
  \item[(i)] Soient $A'$, $B'$ deux anneaux de caractéristique $p$, et
  supposons donnés deux homomorphismes $i':A'\to B'$, $j':B'\to A'$ vérifiant
  (21.3.1.1) et (21.3.1.2), ainsi que des homomorphismes d'anneaux
  $f:A\to A'$, $g:B\to B'$ rendant commutatif le diagramme
  \[
  \xymatrix{
  A'\ar[r]^{i'}&B'\ar[r]^{j'}&A'\\
  A\ar[u]^f\ar[r]^i&B\ar[u]^g\ar[r]^j&A\ar[u]^f
  }
  \]
  Alors l'homomorphisme canonique $\Omega^1_{B/A}\to\Omega^1_{B'/A'}$
  (20.5.4.3) donne un di-homomorphisme canonique
  $\Theta_{B/A}\to\Theta_{B'/A'}$ rendant commutatif le diagramme
  \[
  \xymatrix{
  0\ar[r]&\Xi_{B/A}\ar[r]\ar[d]&\Theta_{B/A}\ar[r]^{\pi_{B/A}}\ar[d]&\Omega^1_A\ar[d]\\
  0\ar[r]&\Xi_{B'/A'}\ar[r]&\Theta_{B'/A'}\ar[r]^{\pi_{B'/A'}}&\Omega^1_{A'}
  }
  \]
  \item[(ii)] Soit $A'$ une $A$-algèbre quelconque et posons
  $B'=B\otimes_A A'$~; alors $i'=i\otimes1:A'\to B'$ et
  $j'=j\otimes1:B'\to A'$ vérifient (21.3.1.1) et (21.3.1.2), et l'on a
  $\Omega^1_{B'/A'}=\Omega^1_{B/A}\otimes_A A'=\Omega^1_{B/A}\otimes_B B'$
  (20.5.5)~; on en déduit un $A'$-isomorphisme canonique
  \begin{equation}
  \tag{21.3.6.1}\label{0.21.3.6.1-fr}
  \Theta_{B/A}\otimes_A A'\xrightarrow{\sim}\Theta_{B'/A'}
  \end{equation}
  et aussi un $B'$-homomorphisme canonique
  \begin{equation}
  \tag{21.3.6.2}\label{0.21.3.6.2-fr}
  (\Omega^1_{B/A})^{(p)}\otimes_B B'\xrightarrow{\sim}
  (\Omega^1_{B'/A'})^{(p)}.
  \end{equation}
\end{enumerate}
\end{remarks}

\subsection{Cas des extensions de corps}
\label{0.21.4-fr}

\begin{env}[21.4.0]
\label{0.21.4.0-fr}
Soient $K$ un corps de caractéristique $p>0$, $k$ un sous-corps de $K$~;
alors l'anneau $k[K^p]$ est égal au corps $k(K^p)$ puisque $k$ est algébrique
sur $K^p$. On pourra donc appliquer les résultats des numéros précédents en
remplaçant partout $A$, $B$ et $A[B^p]$ par $k$, $K$ et $k(K^p)$.
\end{env}

\begin{lemma}[21.4.1]
\label{0.21.4.1-fr}
Soient $k$ un corps de caractéristique $p>0$, $K$ une extension de $k$.
Pour qu'un élément $x\in K$ soit $p$-libre sur $k$, il faut et il suffit que
$x\notin k(K^p)$.
\end{lemma}

\begin{proof}
En effet, $x$ est racine du polynôme $X^p-x^p$ de $k(K^p)[X]$, et l'on sait
(Bourbaki, \emph{Alg.}, chap.~V, §~8, n$^\circ$~1, prop.~1) que si
$x\notin k(K^p)$ ce polynôme est irréductible, de sorte que les éléments
$1,x,\ldots,x^{p-1}$ forment une base du $k(K^p)$-module $k(K^p)(x)$.
\end{proof}

\begin{theorem}[21.4.2]
\label{0.21.4.2-fr}
Soient $k$ un corps de caractéristique $p>0$, $K$ une extension de $k$,
$S$ un système de $p$-générateurs de $K$ sur $k$, $L\subset S$ une partie
$p$-libre sur $k$. Il existe alors une $p$-base $B$ de $K$ sur $k$ telle que
$L\subset B\subset S$. En particulier, toute extension de $k$ admet une
$p$-base sur $k$.
\end{theorem}

\begin{proof}
On peut se borner au cas où $K^p\subset k$. En vertu du théorème de Zorn,
il existe dans $K$ une partie $B$ telle que $L\subset B\subset S$, $p$-libre
sur $k$ et maximale parmi les parties de $S$ ayant ces propriétés. Il suffit
de voir que le sous-corps $K'$ de $K$ engendré par $k$ et $B$ est égal à $K$.
Dans le cas contraire, il existerait $x\in S$ non dans $K'$~; comme
$K^p\subset k$, on a $K'=K'(K^p)$~; donc $x$ serait $p$-libre sur $K'$
(21.4.1), et par suite $B\cup\{x\}$ serait $p$-libre sur $k$ (21.1.10),
contrairement à la définition de $B$. La dernière assertion de (21.4.2)
s'obtient en prenant $L=\varnothing$, $S=K$.
\end{proof}

\begin{corollary}[21.4.3]
\label{0.21.4.3-fr}
Soient $k$ un corps de caractéristique $p>0$, $K$ une extension de $k$.
Pour qu'une famille $(x_\alpha)$ d'éléments de $K$ soit $p$-libre sur $k$,
il faut et il suffit que, pour tout $\alpha$, $x_\alpha$ n'appartienne pas au
corps $K_\alpha$ engendré par $k(K^p)$ et par les $x_\beta$ d'indice
$\beta\neq\alpha$.
\end{corollary}

\begin{proof}
La condition est nécessaire en vertu de (21.1.10). Inversement, supposons-la
remplie~; on peut se borner au cas où $(x_\alpha)$ est un système de
$p$-générateurs de $K$. Il existe alors une sous-famille de $(x_\alpha)$ qui
est une $p$-base de $K$ (21.4.2), mais cette famille ne peut être distincte de
$(x_\alpha)$, sans quoi, en vertu de l'hypothèse, ce ne serait pas une famille
de $p$-générateurs de $K$.
\end{proof}

\oldpage[0\textsubscript{IV-1}]{163}
\begin{corollary}[21.4.4]
\label{0.21.4.4-fr}
Soient $k$ un corps de caractéristique $p>0$, $K$ une extension de $k$.
Pour que $\Omega^1_{K/k}=0$, il faut et il suffit que $K=k(K^p)$.
En particulier, pour que $\Omega^1_K=0$, il faut et il suffit que $K$ soit
un corps parfait.
\end{corollary}

\begin{proof}
En effet, si $(x_\alpha)$ est une $p$-base de $K$ sur $k$, les
$d_{K/k}(x_\alpha)$ forment une base de l'espace vectoriel $\Omega^1_{K/k}$
(21.2.5).
\end{proof}

\begin{theorem}[21.4.5]
\label{0.21.4.5-fr}
Soient $k$ un corps de caractéristique $p>0$, $K$ une extension de $k$,
$(x_\alpha)$ une famille d'éléments de $K$. Pour que $(x_\alpha)$ soit
$p$-libre sur $k$ (resp. une famille de $p$-générateurs de $K$ sur $k$,
resp. une $p$-base de $K$ sur $k$), il faut et il suffit que la famille
$(d_{K/k}(x_\alpha))$ soit une famille libre (resp. un système de
générateurs, resp. une base) dans l'espace vectoriel $\Omega^1_{K/k}$.
\end{theorem}

\begin{proof}
\oldpage[0\textsubscript{IV-1}]{164}
Soit $K'$ le sous-corps (égal au sous-anneau) de $K$ engendré par
$k(K^p)$ et les $x_\alpha$~; compte tenu de (21.4.3), il suffit de voir que,
pour chaque $\alpha$, $x_\alpha$ n'appartient pas à $K_\alpha$. Or, dans la
suite exacte
\[
\Omega^1_{K_\alpha/k(K^p)}\otimes_{K_\alpha}K\longrightarrow
\Omega^1_{K/k(K^p)}\longrightarrow\Omega^1_{K/K_\alpha}\longrightarrow0,
\]
les images par la flèche de gauche des $d_{K_\alpha/k}(x_\beta)$ pour
$\beta\neq\alpha$ sont les $d_{K/k}(x_\beta)$, donc engendrent un sous-espace
vectoriel de $\Omega^1_{K/k(K^p)}$ ne contenant pas $d_{K/k}(x_\alpha)$~; comme
ce sous-espace est le noyau de la flèche de droite, on voit que
$d_{K/K_\alpha}(x_\alpha)\neq0$, donc $x_\alpha\notin K_\alpha$.

Si maintenant $(x_\alpha)$ est une famille $p$-libre, elle est une sous-famille
d'une $p$-base de $K$ sur $k$ (21.4.2)~; les $d_{K/k}(x_\alpha)$ font donc
partie d'une base de l'espace vectoriel $\Omega^1_{K/k}$ (21.2.5), et sont
par suite linéairement indépendants sur $K$. Réciproquement, si les
$d_{K/k}(x_\alpha)$ sont linéairement indépendants sur $K$, les deux
paragraphes précédents montrent que $(x_\alpha)$ est $p$-libre sur $k$.
\end{proof}

\begin{corollary}[21.4.6]
\label{0.21.4.6-fr}
Soient $k$ un corps de caractéristique $p>0$, $K$ une extension de $k$,
$x$ un élément de $K$. Les trois conditions suivantes sont équivalentes~:
\begin{enumerate}
  \item[a)] $x\notin k(K^p)$.
  \item[b)] $d_{K/k}(x)\neq0$.
  \item[c)] L'élément $x$ est $p$-libre sur $k$.
\end{enumerate}
\end{corollary}

\begin{proposition}[21.4.7]
\label{0.21.4.7-fr}
Soient $L$ un corps de caractéristique $p>0$, $K$ un sous-corps de $L$ tel
que $L^p\subset K\subset L$. Alors la suite d'espaces vectoriels sur $L$
\begin{equation}
\tag{21.4.7.1}\label{0.21.4.7.1-fr}
0\to\Omega^1_{K/L^p}\otimes_KL\to\Omega^1_{L/L^p}
\to\Omega^1_{L/K}\to0
\end{equation}
est exacte~; en d'autres termes, on a $\Upsilon_{L/K/L^p}=0$.
\end{proposition}
\begin{proof}
C'est un cas particulier de (21.2.4), compte tenu de (21.4.2).
\end{proof}

\begin{proposition}[21.4.8]
\label{0.21.4.8-fr}
Sous les hypothèses de (21.4.7), l'homomorphisme canonique
\begin{equation}
\tag{21.4.8.1}\label{0.21.4.8.1-fr}
(\Omega^1_{L/K})^{(p)}\longrightarrow\Upsilon_{L/K}
\end{equation}
est bijectif~; si $(x_\lambda)$ est une $p$-base de $L$ sur $K$, les éléments
$d_K(x_\lambda^p)\otimes1$ forment une base de l'espace vectoriel $\Upsilon_{L/K}$.
\end{proposition}
\begin{proof}
C'est un cas particulier de (21.3.5).
\end{proof}

\subsection{Application~: critères de séparabilité}
\label{0.21.5-fr}

Dans ce numéro et les deux suivants, on ne suppose plus que les anneaux
considérés soient de caractéristique $p>0$.

\begin{env}[21.5.1]
\label{0.21.5.1-fr}
Notons d'abord que le critère (21.2.7) permet de démontrer une partie du
théorème de Cohen sur les extensions séparables (19.6.1), savoir que si $k$
est de caractéristique $p>0$ et si $K$ est une extension séparable de $k$,
alors $K$ est une $k$-algèbre formellement lisse. En effet, $K$ admet une
$p$-base sur $k$ (21.4.2), et d'autre part, il résulte du critère de MacLane
(Bourbaki, \emph{Alg.}, chap.~V, §~8, n$^\circ$~2, prop.~3) que, dans une
clôture algébrique de $K$, $k^{p^{-1}}$ et $K$ sont linéairement disjoints sur
$k$, et par suite que l'homomorphisme canonique
$k^{p^{-1}}\otimes_k K\to k^{p^{-1}}(K)$ est bijectif, ce qui est précisément
la condition (i) de (21.2.7), après transport de structure par
l'isomorphisme $x\mapsto x^p$.
\end{env}
\begin{proposition}[21.5.2]
\label{0.21.5.2-fr}
\textnormal{(MacLane).}
Soient $A$, $B$ deux anneaux de valuation discrète complets, $\mathfrak m$,
$\mathfrak n$ leurs idéaux maximaux respectifs, $K=A/\mathfrak m$ et
$L=B/\mathfrak n$ leurs corps résiduels, $u:A\to B$ un homomorphisme tel que
$Bu(\mathfrak m)=\mathfrak n$, $u_0=u\otimes 1:K\to L$ l'homomorphisme
correspondant pour les corps résiduels. On considère les conditions suivantes~:
\begin{enumerate}
  \item[\rm a)] $L$ est une extension séparable de $K$ (pour $u_0$).
  \item[\rm b)] Pour tout anneau de valuation discrète complet $B'$ d'idéal
  maximal $\mathfrak n'$ et de corps résiduel $L'$,
\oldpage[0\textsubscript{IV-1}]{165}
  tout homomorphisme
  $u':A\to B'$ tel que $B'u'(\mathfrak m)=\mathfrak n'$ et tout
  $K$-isomorphisme $\sigma:L\to L'$ (relatif à $u_0$ et
  $u'_0=u'\otimes1:K\to L'$), il existe un isomorphisme $w:B\to B'$ tel que
  $u'=w\circ u$ et que $w_0=w\otimes1:L\to L'$ soit égal à $\sigma$.
  \item[\rm b$'$)] Pour tout homomorphisme $u':A\to B$ tel que
  $Bu'(\mathfrak m)=\mathfrak n$ et tel que $u'_0=u'\otimes1:K\to L$ soit égal
  à $u_0$, il existe un automorphisme $w$ de $B$ tel que $u'=w\circ u$ et que
  $w_0=w\otimes1:L\to L$ soit l'identité.
  \item[\rm c)] Désignant par $u_1:A/\mathfrak m^2\to B/\mathfrak n^2$
  l'homomorphisme déduit de $u$ par passage aux quotients, alors, pour tout
  homomorphisme local $u'_1:A/\mathfrak m^2\to B/\mathfrak n^2$ tel que
  $\operatorname{gr}_0(u'_1)=\operatorname{gr}_0(u_1)(=u_0)$, il existe un
  automorphisme $w_1$ de $B/\mathfrak n^2$ tel que
  $\operatorname{gr}_0(w_1)$ soit l'identité et que $u'_1=w_1\circ u_1$.
  \item[\rm c$'$)] Pour tout homomorphisme local $u'_1:A/\mathfrak m^2\to
  B/\mathfrak n^2$ tel que $\operatorname{gr}_0(u'_1)=\operatorname{gr}_0(u_1)$
  et $\operatorname{gr}_1(u'_1)=\operatorname{gr}_1(u_1)$
  (homomorphisme de $\mathfrak m/\mathfrak m^2$ dans
  $\mathfrak n/\mathfrak n^2$), il existe un automorphisme $w_1$ de
  $B/\mathfrak n^2$ tel que $\operatorname{gr}_0(w_1)$ soit l'identité et que
  $u'_1=w_1\circ u_1$.
\end{enumerate}
 Alors on a les implications $c)\Leftrightarrow c')\Leftrightarrow a)\Rightarrow b)
 \Rightarrow b')$. Si de plus $A$ est un anneau de Cohen, les cinq conditions
précédentes sont équivalentes.

\end{proposition}
\begin{proof}
Les implications $b)\Rightarrow b')$ et $c)\Rightarrow c')$ sont triviales.
Montrons que $a)$ implique $b)$.
L'homomorphisme $u$ fait de $B$ un $A$-module sans torsion, puisqu'il
transforme par hypothèse une uniformisante de $A$ en une uniformisante de $B$~;
il en résulte que $B$ est un $A$-module plat (0$_{\rm I}$, 6.3.4), donc une
$A$-algèbre de Cohen (19.8.1)~; le fait que $a)$ implique $b)$ est alors
 conséquence de (19.8.2, (i)) appliquée à $G=B'$, $\mathfrak J=\mathfrak n'$,
 $B'$ étant considéré comme $A$-algèbre pour $u'$ et l'homomorphisme
$B\to B'/\mathfrak n'$ étant le composé
$B\to B/\mathfrak n\to B'/\mathfrak n'$, qui est un $A$-homomorphisme en
 vertu de l'hypothèse $u'_0=\sigma\circ u_0$. Le même raisonnement prouve que
$a)$ entraîne $c)$, en prenant cette fois $G=B/\mathfrak n^2$,
$\mathfrak J=\mathfrak n/\mathfrak n^2$, $G$ étant considéré comme
$A$-algèbre pour $u'_1$.
Prouvons en second lieu que $c')$ implique $a)$. Les deux homomorphismes
$u_1$ et $u'_1$ sont tels que $u'_1=u_1+D$, où $D$ est une dérivation de
$A/\mathfrak m^2$ dans $\mathfrak n/\mathfrak n^2$ (20.1.1)~; d'ailleurs,
l'hypothèse $\operatorname{gr}_1(u'_1)=\operatorname{gr}_1(u_1)$ signifie
que $D$ est nulle dans $\mathfrak m/\mathfrak m^2$, et peut par suite être
considérée comme une dérivation de $K=A/\mathfrak m$ dans
$\mathfrak n/\mathfrak n^2$, et $\mathfrak n/\mathfrak n^2$ s'identifie
(par choix d'une uniformisante de $B$) au $K$-module $L$ (pour $u_0$).
D'autre part, les conditions imposées à $w_1$ entraînent que
 $\operatorname{gr}_1(w_1)$ est aussi l'identité (puisque
 $u_1(\mathfrak m/\mathfrak m^2)$ engendre
$\mathfrak n/\mathfrak n^2$)~; $w_1$ est donc de la forme
$x\mapsto x+D'(x)$, où $D'$ est cette fois une dérivation de
$B/\mathfrak n^2$ dans le $(B/\mathfrak n^2)$-module
$\mathfrak n/\mathfrak n^2$~; comme $w_1$ est l'identité dans
$\mathfrak n/\mathfrak n^2$, $D'$ peut encore (par l'identification
précédente de $\mathfrak n/\mathfrak n^2$ et de $L$) être considérée comme
une dérivation de $L$ dans $L$~; enfin, la relation $u'_1=w_1\circ u_1$
signifie que $D'$ prolonge $D$. Notons d'autre part que toute dérivation $D$
de $K$ dans $L$ correspond à un homomorphisme $u'_1$ vérifiant les conditions
de $c')$ (20.1.1)~; la condition $c')$ signifie donc que toute dérivation de
$K$ dans $L$ se prolonge en une dérivation de $L$ dans lui-même, c'est-à-dire
(20.6.5) que $L$ est séparable sur $K$.
Prouvons enfin que, lorsque $A$ est un anneau de Cohen, $b')$ entraîne $c')$.
Prouvons d'abord que sous les hypothèses de $c')$, il existe un homomorphisme
$u':A\to B$ qui vérifie les hypothèses de $b')$ et qui, par passage aux quotients,
donne $u'_1:A/\mathfrak m^2\to B/\mathfrak n^2$. En effet, cela résulte de
(19.8.6, (i)) appliqué à l'homomorphisme composé
$v':A\to A/\mathfrak m^2\xrightarrow{u'_1}B/\mathfrak n^2$~;
\oldpage[0\textsubscript{IV-1}]{166}
ce dernier se factorise donc en $A\xrightarrow{u'}B\to B/\mathfrak n^2$ et
les hypothèses sur $\operatorname{gr}_0(u'_1)$ et
$\operatorname{gr}_1(u'_1)$ entraînent
$\operatorname{gr}_0(u')=\operatorname{gr}_0(u)=u_0$ et
$\operatorname{gr}_1(u')=\operatorname{gr}_1(u)$, donc l'image par $u'$ d'une
uniformisante de $A$ est une uniformisante de $B$, et l'on a par suite
$Bu'(\mathfrak m)=\mathfrak n$. Il suffit alors, pour obtenir un automorphisme
$w_1$ répondant à la question, de prendre l'automorphisme déduit par passage
aux quotients de l'automorphisme $w$ de $B$ fourni par l'application de $b')$
à l'homomorphisme $u'$.
\end{proof}
\begin{remarks}[21.5.3]
\label{0.21.5.3-fr}
\begin{enumerate}
  \item[(i)] Les propriétés différentielles des corps permettent de résoudre la
  question de l'unicité du corps de représentants dans un anneau local
  noethérien complet $C$ (19.8.7, (ii)). Soient $\mathfrak J$ l'idéal maximal
  de $C$, $K=C/\mathfrak J$ le corps résiduel de $C$~; on peut se borner au
  cas $\mathfrak J\ne0$. Supposons qu'il existe un homomorphisme $u:K\to C$
  qui, composé avec l'augmentation $C\to K$, donne l'identité~; alors, pour
  que cet homomorphisme soit unique, il faut et il suffit que
  $\Omega^1_K=0$. En effet, la condition $\Omega^1_K=0$ entraîne que $K$ est
  formellement non ramifié sur son corps premier (20.7.4)~; s'il existait un
  second homomorphisme $v\ne u$ répondant à la question, il y aurait un plus
  grand entier $n$ tel que $\mathfrak J^n$ contienne l'ensemble des
  $u(x)-v(x)$ pour $x\in K$~; par passage au quotient, $u$ et $v$ donneraient
  deux homomorphismes distincts $u_1,v_1$ de $K$ dans
  $C/\mathfrak J^{n+1}$, dont les composés avec
  $C/\mathfrak J^{n+1}\to C/\mathfrak J^n$ seraient égaux, ce qui contredit
  la définition (19.10.2). Inversement, supposons que $\Omega^1_K\ne0$~; il
  existe alors une dérivation $D\ne0$ de $K$ dans
  $\mathfrak J/\mathfrak J^2$ (20.4.8), donc un homomorphisme
  $v_1:K\to C/\mathfrak J^2$ tel que, si $u_1:K\to C/\mathfrak J^2$ est
  obtenu par passage au quotient à partir de $u$, on ait $v_1=u_1+D$
  (20.1.1). Si $k$ est le corps premier de $K$, $C$ est une $k$-algèbre et
  $K$ une $k$-algèbre formellement lisse (19.6.1), et les restrictions de
  $u_1$ et $v_1$ à $k$ coïncident, donc $v_1$ se factorise en
  $K\xrightarrow{v}C\to C/\mathfrak J^2$ et l'on a $v\ne u$.
  Rappelons ((20.6.20) et (21.4.4)) que la condition $\Omega^1_K=0$ signifie
  que $K$ est parfait s'il est de caractéristique $\ne0$, et extension
  algébrique de $\mathbf Q$ s'il est de caractéristique $0$.
  \item[(ii)] De la même manière, soient $W$ un anneau de Cohen, $C$ un anneau
  local noethérien complet, soit $\mathfrak J\ne0$ un idéal de $C$ contenu dans
  l'idéal maximal~; pour que la factorisation $W\to C\to C/\mathfrak J$
  dans (19.8.6, (i)) soit unique, il faut et il suffit que le corps résiduel
  $K$ de l'anneau de Cohen $W$ soit tel que $\Omega^1_K=0$. En effet, si
  $\Omega^1_K\ne0$ et si $\mathfrak m$ désigne l'idéal maximal de $W$, il
  suffit de composer une dérivation $D\ne0$ de $K$ dans
  $\mathfrak J/\mathfrak m\mathfrak J$ avec l'augmentation $W\to K$ pour
  obtenir une dérivation $D'\ne0$ de $W$ dans
  $\mathfrak J/\mathfrak m\mathfrak J$ et l'on termine le raisonnement comme
  dans (i) en formant à l'aide de $D'$ un homomorphisme
  $v_1:W\to C/\mathfrak m\mathfrak J$ distinct de l'homomorphisme
  $u_1:W\to C/\mathfrak m\mathfrak J$ déduit de $u$ par passage au quotient~;
  on achève le raisonnement en invoquant cette fois (19.8.6, (i)). Si au
  contraire $\Omega^1_K=0$, l'unicité de $v$ résulte déjà de (i) lorsque
  $W=K$ est un corps de caractéristique $0$. Dans le cas contraire, on a
  $\widehat{\Omega}^1_W=0$~; en effet, on a alors $\mathfrak m=pW$, $p$ étant
  la caractéristique de $K$ (19.8.5), et l'homomorphisme canonique
  $\mathfrak m/\mathfrak m^2\to\Omega^1_W/\mathfrak m\Omega^1_W$ (20.5.11.2)
  est par suite nul. La suite exacte (20.5.12.1) appliquée à $W$ et à
  $K=W/\mathfrak m$ entraîne alors que $\Omega^1_W=\mathfrak m\Omega^1_W$,
  d'où notre assertion. Mais alors (20.7.4) $W$ est formellement non ramifié
  (pour sa topologie adique) sur $\mathbf Z$, et l'unicité de $v$ se prouve
  comme dans (i).
\end{enumerate}
\end{remarks}

\oldpage[0\textsubscript{IV-1}]{167}
\subsection{Corps admissibles pour une extension}
\label{section:0.21.6-fr}

\begin{env}[21.6.1]
\label{0.21.6.1-fr}
Étant donnés quatre corps $k_0\subset k\subset K\subset L$, il résulte de
(20.6.16) et (20.6.17) que l'on a une suite exacte
\begin{equation}
\tag{21.6.1.1}
\label{0.21.6.1.1-fr}
0\to\Upsilon_{K/k/k_0}\otimes_KL\xrightarrow{v}\Upsilon_{L/k/k_0}
\xrightarrow{u}\Upsilon_{L/K/k_0}\xrightarrow{s}\Upsilon_{L/K/k}\to0.
\end{equation}

Lorsqu'on laisse fixes $k_0$, $K$ et $L$ et qu'on fait «varier» le corps
intermédiaire $k$ entre $k_0$ et $K$, on a évidemment
$\Upsilon_{L/K/k_0}=\Upsilon_{L/K/k}$ lorsque $k=k_0$. Lorsque
l'homomorphisme canonique $s$ de (21.6.1.1) est encore bijectif, on dit que
$k$ est un corps $k_0$-admissible pour l'extension $L$ de $K$. L'intérêt de
l'existence, sous certaines conditions, de tels corps $k$, qui soient pourtant
«suffisamment proches» de $K$ (par exemple tels que $[K:k]$ soit fini), est
qu'ils permettent de remplacer dans certaines questions les modules de
différentielles $\Omega^1_{K/k_0}$ et $\Omega^1_{L/k_0}$ (qui peuvent être
«trop grands», par exemple lorsque $k_0$ est le corps premier) par
$\Omega^1_{K/k}$ et $\Omega^1_{L/k}$, plus aisément maniables.

Lorsque $k_0$ est le \emph{corps premier}, on dira «corps admissible» au lieu
de «corps $k_0$-admissible».
 
On posera
\end{env}

\begin{env}[21.6.1.2]
\label{0.21.6.1.2-fr}
\[
\Delta(L/K,k/k_0)=\operatorname{Coker}
(\Upsilon_{K/k/k_0}\otimes_KL\to\Upsilon_{L/k/k_0})
\simeq\operatorname{Ker}(\Upsilon_{L/K/k_0}\to\Upsilon_{L/K/k})
\]
(espace vectoriel sur $L$)~; son rang sera noté $d(L/K,k/k_0)$ et appelé
\emph{défaut de $k_0$-admissibilité de $k$ pour l'extension $L$ de $K$}
(il est évidemment nul si et seulement si $k$ est $k_0$-admissible pour cette
extension). Lorsque $k_0$ est le corps premier, on écrira
$\Delta(L/K,k)$ et $d(L/K,k)$ au lieu de $\Delta(L/K,k/k_0)$ et
$d(L/K,k/k_0)$.
\end{env}

\begin{proposition}[21.6.2]
\label{0.21.6.2-fr}
Soient $k_0\subset k\subset K\subset L$ quatre corps.
\begin{enumerate}
  \item[\rm (i)] Les conditions suivantes sont équivalentes~:
  \begin{enumerate}
    \item[\rm a)] Le corps $k$ est $k_0$-admissible pour l'extension $L$ de $K$
    (autrement dit, l'homomorphisme $\Upsilon_{L/K/k_0}\to
    \Upsilon_{L/K/k}$ est injectif, donc bijectif).
    \item[\rm b)] L'homomorphisme canonique
    $u:\Upsilon_{K/k/k_0}\to\Upsilon_{L/k/k_0}$ est nul.
    \item[\rm c)] L'homomorphisme canonique
    $v:\Upsilon_{K/k/k_0}\otimes_KL\to\Upsilon_{L/k/k_0}$ est surjectif
    (donc bijectif).
    \item[\rm d)] On a $d(L/K,k/k_0)=0$ (ou
    $\Delta(L/K,k/k_0)=0$).
  \end{enumerate}
  \item[\rm (ii)] Les conditions équivalentes de (i) sont vérifiées lorsque
  l'on est dans l'un des cas suivants~: $\alpha)$ $L$ est séparable sur $k$~;
  $\beta)$ $L$ est séparable sur $K$~; $\gamma)$ on a
  $k\subset k_0(K^p)$, en désignant par $p$ l'exposant caractéristique de $k_0$.
\end{enumerate}
\end{proposition}

\begin{proof}
\begin{enumerate}
  \item[(i)] Les assertions résultent trivialement de l'exactitude de la suite
  (21.6.1.1).
  \item[(ii)] Si $L$ est séparable sur $k$, on a
  $\Upsilon_{L/K/k_0}=0$ (20.6.19), donc la condition b) de (i) est remplie~;
  si $L$ est séparable sur $K$, on a
  $\Upsilon_{L/K/k_0}=0$ (20.6.19), donc la condition a) de (i) est remplie~;
  enfin, si l'on a $k\subset k_0(K^p)$, il en résulte que
  $\Upsilon_{L/K/k_0}=\Upsilon_{L/K/k}$ en vertu de (21.1.5.1), et la
  condition a) de (i) est vérifiée.
\end{enumerate}
\end{proof}
\oldpage[0\textsubscript{IV-1}]{168}
\begin{env}[21.6.3]
\label{0.21.6.3-fr}
Supposons que l'on ait un diagramme commutatif de monomorphismes de corps
\[
\xymatrix{
 k'_0\ar[r]&k'\ar[r]&K'\ar[r]&L'\\
 k_0\ar[u]\ar[r]&k\ar[u]\ar[r]&K\ar[u]\ar[r]&L\ar[u]
}
\]
Il résulte alors de (20.6.17, (ii)) que l'on a un homomorphisme canonique
\begin{equation}
\tag{21.6.3.1}
\label{0.21.6.3.1-fr}
\Delta(L/K,k/k_0)\longrightarrow\Delta(L'/K',k'/k'_0)
\end{equation}
avec une propriété de transitivité évidente, de sorte que l'on peut dire que
$\Delta(L/K,k/k_0)$ est un \emph{foncteur} en le quadruplet
$(k_0,k,K,L)$.
\end{env}

\begin{proposition}[21.6.4]
\label{0.21.6.4-fr}
\begin{enumerate}
  \item[\rm (i)] Soient $k\subset k'\subset k''\subset K\subset L$ cinq corps.
  On a une suite exacte d'homomorphismes canoniques
\begin{equation}
\tag{21.6.4.1}
\label{0.21.6.4.1-fr}
0\to\Delta(L/K,k'/k)\to\Delta(L/K,k''/k)\to\Delta(L/K,k''/k')\to0
\end{equation}
et par suite l'égalité
\begin{equation}
\tag{21.6.4.2}
\label{0.21.6.4.2-fr}
d(L/K,k''/k)=d(L/K,k''/k')+d(L/K,k'/k).
\end{equation}
  \item[\rm (ii)] Soient $k_0\subset k\subset K\subset L\subset M$ cinq corps.
  On a une suite exacte d'homomorphismes canoniques
\begin{equation}
\tag{21.6.4.3}
\label{0.21.6.4.3-fr}
0\to\Delta(L/K,k/k_0)\otimes_LM\to\Delta(M/K,k/k_0)
\to\Delta(M/L,k/k_0)\to0
\end{equation}
et par suite l'égalité
\begin{equation}
\tag{21.6.4.4}
\label{0.21.6.4.4-fr}
d(M/K,k/k_0)=d(M/L,k/k_0)+d(L/K,k/k_0).
\end{equation}
\end{enumerate}
\end{proposition}

\begin{proof}
\begin{enumerate}
  \item[(i)] Considérons le diagramme commutatif
\[
\xymatrix{
0\ar[r]&\Upsilon_{K/k'/k}\otimes_KL\ar[r]\ar[d]&
\Upsilon_{K/k''/k}\otimes_KL\ar[r]\ar[d]&
\Upsilon_{K/k''/k'}\otimes_KL\ar[r]\ar[d]&0\\
0\ar[r]&\Upsilon_{L/k'/k}\ar[r]\ar[d]&
\Upsilon_{L/k''/k}\ar[r]\ar[d]&
\Upsilon_{L/k''/k'}\ar[r]\ar[d]&0\\
0\ar[r]&\Delta(L/K,k'/k)\ar[r]&\Delta(L/K,k''/k)\ar[r]&
\Delta(L/K,k''/k')\ar[r]&0
}
\]
où l'on peut considérer les trois lignes comme des \emph{complexes}
$T_1,T_2,T_3$ respectivement~; la suite exacte (21.6.1.1) et la définition
de $\Delta$ (21.6.1.2) montrent que l'on a une suite exacte de complexes
$0\to T_1\to T_2\to T_3\to0$~; appliquons-lui la suite exacte de
\oldpage[0\textsubscript{IV-1}]{169}
cohomologie, et notons que, en vertu de l'exactitude de (21.6.1.1), la
cohomologie de $T_1$ et celle de $T_2$ sont nulles sauf en un seul et même
degré, pour lequel les modules de cohomologie sont tous deux égaux à
$\Upsilon_{k''/k'/k}\otimes_{k''}L$~; comme $T_1$ et $T_2$ ont ainsi même
cohomologie, celle de $T_3$ est nécessairement nulle, ce qui prouve (i).
  \item[(ii)] Considérons de même le diagramme commutatif
\[
\xymatrix{
0\ar[r]&\Delta(L/K,k/k_0)\otimes_LM\ar[r]\ar[d]&
\Delta(M/K,k/k_0)\ar[r]\ar[d]&\Delta(M/L,k/k_0)\ar[r]&0\\
0\ar[r]&\Upsilon_{L/K/k_0}\otimes_LM\ar[r]\ar[d]&
\Upsilon_{M/K/k_0}\ar[r]\ar[d]&\Upsilon_{M/L/k_0}\ar[r]&0\\
0\ar[r]&\Upsilon_{L/K/k}\otimes_LM\ar[r]&
\Upsilon_{M/K/k}\ar[r]&\Upsilon_{M/L/k}\ar[r]&0
}
\]
où de nouveau on considère les trois lignes comme des \emph{complexes}
$T'_1,T'_2,T'_3$~; la suite exacte (21.6.1.1) et la définition de $\Delta$
(21.6.1.2) donnent ici une suite exacte de complexes
$0\to T'_1\to T'_2\to T'_3\to0$ à laquelle on applique encore la suite exacte
de cohomologie~; cette fois, en vertu de l'exactitude de (21.6.1.1), la
cohomologie de $T'_2$ et celle de $T'_3$ sont nulles sauf en un seul et même
degré, pour lequel les modules de cohomologie sont tous deux égaux à
$\Upsilon_{M/L/k_0}$~; on conclut donc ici que la cohomologie de $T'_1$ est
nulle, ce qui établit (ii).
\end{enumerate}
\end{proof}

\begin{corollary}[21.6.5]
\label{0.21.6.5-fr}
\begin{enumerate}
  \item[(i)] Étant donnés cinq corps $k\subset k'\subset k''\subset K\subset L$,
  pour que $k''$ soit $k$-admissible pour l'extension $L$ de $K$, il faut et
  il suffit que $k'$ soit $k$-admissible et que $k''$ soit $k'$-admissible
  pour l'extension $L$ de $K$.
  \item[(ii)] Étant donnés cinq corps
  $k_0\subset k\subset K\subset L\subset M$, pour que $k$ soit
  $k_0$-admissible pour l'extension $M$ de $K$, il faut et il suffit qu'il le
  soit pour l'extension $L$ de $K$ et pour l'extension $M$ de $L$.
\end{enumerate}

Cela résulte aussitôt des relations (21.6.4.2) et (21.6.4.4), les valeurs de
$d$ étant $\geqslant0$.
\end{corollary}

\begin{corollary}[21.6.6]
\label{0.21.6.6-fr}
Soient $k_0\subset k\subset K\subset L$ quatre corps, et supposons que $k$
soit $k_0$-admissible pour l'extension $L$ de $K$. Alors, si $k'_0$, $k'$,
$K'$, $L'$ sont quatre corps tels que
$k_0\subset k'_0\subset k'\subset k\subset K\subset K'\subset L'\subset L$,
$k'$ est $k'_0$-admissible pour l'extension $L'$ de $K'$.
\end{corollary}

\subsection{L'égalité de Cartier}
\label{subsection:0.21.7-fr}

Le résultat suivant traduit en termes de différentielles un théorème de
MacLane sur les dérivations~:

\oldpage[0\textsubscript{IV-1}]{170}
\begin{theorem}[21.7.1]
\label{0.21.7.1-fr}
\textnormal{(Cartier).}
Soient $K$ un corps, $L$ une extension de type fini de $K$. Alors
$\Omega^1_{L/K}$ et $\Upsilon_{L/K}$ sont des $L$-espaces vectoriels de rang
fini, et l'on a
\begin{equation}
\tag{21.7.1.1}
\label{0.21.7.1.1-fr}
\operatorname{rg}\Omega^1_{L/K}-\operatorname{rg}\Upsilon_{L/K}
=\deg.\operatorname{tr}_K L.
\end{equation}
\end{theorem}

\begin{proof}
Si $L'$ est un corps tel que $K\subset L'\subset L$, on a la suite exacte
(20.6.15.1) (appliquée à $\Lambda=P$, corps premier de $K$, $A=K$,
$B=L'$, $C=L$)
\[
0\to\Upsilon_{L'/K}\otimes_{L'}L\to\Upsilon_{L/K}
\to\Omega^1_{L'/K}\otimes_{L'}L\to\Omega^1_{L/K}
\to\Omega^1_{L/L'}\to0
\]
d'où (0$_{\rm III}$, 11.10.2)
\[
\operatorname{rg}_L\Omega^1_{L/K}-\operatorname{rg}_L\Upsilon_{L/K}
=(\operatorname{rg}_{L'}\Omega^1_{L/L'}-\operatorname{rg}_{L'}\Upsilon_{L/L'})
+(\operatorname{rg}_{L'}\Omega^1_{L'/K}-\operatorname{rg}_{L'}\Upsilon_{L'/K}).
\]
Comme par ailleurs $\deg.\operatorname{tr}_K L=\deg.\operatorname{tr}_K L'
+\deg.\operatorname{tr}_{L'}L$, on voit, par récurrence sur le nombre de
générateurs de l'extension $L$, que l'on est ramené à prouver (21.7.1.1)
lorsque $L=K(x)$. Distinguons alors trois cas~:
\begin{enumerate}
  \item[\rm a)] $x$ est transcendant sur $K$~; comme $L$ est séparable sur
  $K$, on a alors $\Upsilon_{L/K}=0$ (20.6.3)~; d'autre part, (20.5.9) et
  (20.4.13) montrent que $\Omega^1_{L/K}$ est de rang $1$, d'où
  (21.7.1.1) dans ce cas.
  \item[\rm b)] $L$ est une extension algébrique séparable de $K$, de sorte
  que l'on a encore $\Upsilon_{L/K}=0$ (20.6.3). D'autre part, on a
  $\Omega^1_{L/K}=0$ par (20.6.20), d'où encore (21.7.1.1).
  \item[\rm c)] $L$ est une extension algébrique inséparable de $K$~; le
  raisonnement du début montre que l'on peut se borner au cas où $L^p\subset K$~;
  il résulte donc de (21.4.8) que l'on a
  $\operatorname{rg}_L\Upsilon_{L/K}=\operatorname{rg}_L\Omega^1_{L/K}$,
  d'où encore (21.7.1.1).
\end{enumerate}
\end{proof}

\begin{corollary}[21.7.2]
\label{0.21.7.2-fr}
Soient $K$ un corps, $L$ une extension de type fini de $K$, $k$ un sous-corps
de $K$. Alors $\Omega^1_{L/K}$ et $\Upsilon_{L/K/k}$ sont des espaces
vectoriels de rang fini sur $L$, et l'on a
\begin{equation}
\tag{21.7.2.1}
\label{0.21.7.2.1-fr}
\operatorname{rg}\Omega^1_{L/K}-\operatorname{rg}\Upsilon_{L/K/k}
=\deg.\operatorname{tr}_K L+d(L/K,k).
\end{equation}
Par suite, on a l'inégalité
\begin{equation}
\tag{21.7.2.2}
\label{0.21.7.2.2-fr}
\operatorname{rg}\Omega^1_{L/K}-\operatorname{rg}\Upsilon_{L/K/k}
\geqslant\deg.\operatorname{tr}_K L.
\end{equation}
En outre, pour que les deux membres de (21.7.2.2) soient égaux, il faut et
il suffit que $k$ soit un corps admissible pour l'extension $L$ de $K$.
\end{corollary}

\begin{proof}
En effet, comme l'homomorphisme
$\Upsilon_{L/K}\to\Upsilon_{L/K/k}$ est surjectif, on a par définition
(21.6.1.2)
$\operatorname{rg}\Upsilon_{L/K}-\operatorname{rg}\Upsilon_{L/K/k}
=d(L/K,k)$, et le corollaire résulte donc aussitôt de (21.7.1) et de
(21.6.2).
\end{proof}

\begin{corollary}[21.7.3]
\label{0.21.7.3-fr}
Soient $k\subset K\subset L$ trois corps tels que $L$ soit une extension de
type fini de $k$. Alors on a
\begin{equation}
\tag{21.7.3.1}
\label{0.21.7.3.1-fr}
\operatorname{rg}_L\Omega^1_{L/k}-\operatorname{rg}_K\Omega^1_{K/k}
=\deg.\operatorname{tr}_K L+d(L/K,k).
\end{equation}
\oldpage[0\textsubscript{IV-1}]{171}
Par suite, on a l'inégalité
\begin{equation}
\tag{21.7.3.2}
\label{0.21.7.3.2-fr}
\operatorname{rg}_L\Omega^1_{L/k}-\operatorname{rg}_K\Omega^1_{K/k}
\geqslant\deg.\operatorname{tr}_K L.
\end{equation}
L'égalité étant atteinte si et seulement si $k$ est un corps admissible pour
l'extension $L$ de $K$.
\end{corollary}

\begin{proof}
On a en effet la suite exacte (20.6.1.1)
\[
0\to\Upsilon_{L/K/k}\to\Omega^1_{K/k}\otimes_KL\to\Omega^1_{L/k}
\to\Omega^1_{L/K}\to0
\]
donc le premier membre de (21.7.3.1) est égal à
\[
\operatorname{rg}_L\Omega^1_{L/k}-\operatorname{rg}_L\Upsilon_{L/K/k}
\]
et il suffit d'appliquer (21.7.2).
\end{proof}

L'intérêt de ce dernier corollaire est qu'il ne fait intervenir que des
modules de différentielles, à l'exclusion de modules d'imperfection. Nous
verrons d'ailleurs plus bas (21.8.6) que pour toute extension de type fini
$L$ de $K$, il existe un sous-corps $k$ de $K$ tel que $[K:k]$ soit fini et
qui est admissible pour l'extension $L$ de $K$, de sorte que les deux membres
de (21.7.3.2) sont alors égaux.

\begin{corollary}[21.7.4]
\label{0.21.7.4-fr}
Soit $K$ une extension de type fini d'un corps $k$. Les conditions suivantes
sont équivalentes~:
\begin{enumerate}
  \item[\rm a)] $K$ est une extension finie séparable de $k$.
  \item[\rm b)] $K$ est une $k$-algèbre formellement non ramifiée (19.10.2).
  \item[\rm c)] $K$ est une $k$-algèbre formellement étale (19.10.2).
  \item[\rm d)] On a $\Omega^1_{K/k}=0$.
\end{enumerate}
\end{corollary}

\begin{proof}
En effet, on sait (les topologies étant discrètes) que les conditions b) et
d) sont équivalentes (20.7.4) et que a) entraîne que $K$ est une
$k$-algèbre formellement lisse (19.6.1), donc la conjonction de a) et b)
équivaut à c). Tout revient à voir que d) entraîne a). Or, en vertu de
(21.7.1.1), la relation $\Omega^1_{K/k}=0$ entraîne que $K$ est algébrique
(donc fini) sur $k$ et que $\Upsilon_{K/k}=0$, c'est-à-dire (20.6.3) que
$K$ est séparable sur $k$.
\end{proof}

\begin{remark}[21.7.5]
\label{0.21.7.5-fr}
En vertu de (20.6.3.2) et de (0$_{\rm III}$, 11.10.2), le premier membre de
(21.7.1.1) n'est autre que la \emph{caractéristique d'Euler-Poincaré} du
complexe $K_\bullet(C/B/A)$ introduit dans (20.6.3), pour $C=L$, $B=K$ et
$A=P$ (corps premier de $K$). Dans le chapitre consacré à la théorie des
intersections et au théorème de Riemann--Roch, un rôle important sera joué
également par les caractéristiques d'Euler-Poincaré généralisées (à valeurs
dans des groupes de classes de $\mathcal O_X$-Modules) de complexes
généralisant les complexes $F_\bullet(C/A)$ considérés dans (20.6.22).
\end{remark}

\subsection{Critères d'admissibilité}
\label{subsection:0.21.8-fr}

Nous revenons à nos conventions antérieures et supposons donc que tous les
corps considérés dans ce numéro sont de caractéristique $p>0$.

\begin{lemma}[21.8.1]
\label{0.21.8.1-fr}
Soient $K$ un corps, $k$ un sous-corps de $K$, $(k_\alpha)_{\alpha\in I}$ une
famille filtrante décroissante de sous-corps de $K$ telle que
$k=\bigcap_{\alpha\in I}k_\alpha$. Soient $V$ un espace vectoriel sur $K$,
$(a_i)_{1\leqslant i\leqslant n}$ une famille finie de vecteurs de $V$~; si la
famille $(a_i)$ est libre sur $k$, il existe un indice $\gamma$ tel qu'elle
soit aussi libre sur $k_\gamma$.
\end{lemma}

\oldpage[0\textsubscript{IV-1}]{172}
\begin{proof}
Soit $r$ le rang de la famille $(a_i)_{1\leqslant i\leqslant n}$ sur $K$, et
raisonnons par récurrence sur $n-r$~; la proposition est évidente pour $n=r$,
car alors la famille $(a_i)_{1\leqslant i\leqslant r}$ est libre sur $K$, donc
sur tout sous-corps de $K$. Supposons par exemple que
$(a_i)_{1\leqslant i\leqslant r}$ soit libre sur $K$, et écrivons
$a_{r+1}=\sum_{i=1}^r\lambda_i a_i$ avec $\lambda_i\in K$~; la famille
$(a_i)_{1\leqslant i\leqslant r+1}$ étant libre sur $k$, les $\lambda_i$ ne
peuvent tous appartenir à $k$~; supposons par exemple que
$\lambda_{i_0}\notin k$. Alors il existe un indice $\beta$ tel que
$\lambda_{i_0}\notin k_\beta$~; on en conclut que la famille $(a_i)_{1\leqslant i\leqslant r+1}$ est libre sur $k_\beta$~; en effet, comme la
famille $(a_i)_{1\leqslant i\leqslant r}$ est libre sur tout sous-corps de
$K$, si la famille $(a_i)_{1\leqslant i\leqslant r+1}$ n'était pas libre sur
$k_\beta$, $a_{r+1}$ serait égal à une combinaison linéaire des $a_i$ à
coefficients dans $k_\beta$, et comme ces coefficients sont nécessairement
les $\lambda_i$, on aboutit à une contradiction. Il suffit maintenant
d'appliquer l'hypothèse de récurrence en remplaçant $K$ par $k_\beta$ et la
famille $(k_\alpha)_{\alpha\in I}$ par la sous-famille des $k_\alpha$
contenus dans $k_\beta$.
\end{proof}

\begin{lemma}[21.8.2]
\label{0.21.8.2-fr}
Soient $K$ un corps, $p$ son exposant caractéristique, $k_0$ un sous-corps de
$K$, $(k_\alpha)_{\alpha\in I}$ une famille filtrante décroissante de
sous-corps de $K$ telle que $\bigcap_{\alpha\in I}k_\alpha(K^p)=k_0(K^p)$.
Si $(x_i)_{1\leqslant i\leqslant n}$ est une famille finie d'éléments de $K$
qui est $p$-libre sur $k_0$, il existe un indice $\alpha$ tel que $(x_i)$ soit
$p$-libre sur $k_\alpha$.

En effet, dire que la famille $(x_i)$ est $p$-libre sur un sous-corps $k$
de $K$ signifie que la famille finie des monômes
$\prod_{i=0}^n x_i^{m(i)}$ avec $0\leqslant m(i)<p$ est libre sur $k(K^p)$~;
il suffit donc d'appliquer le lemme (21.8.1) à cette famille de monômes dans
l'espace vectoriel $V=K$, et aux sous-corps $k_\alpha(K^p)$ et $k_0(K^p)$
de $K$.
\end{lemma}

\begin{theorem}[21.8.3]
\label{0.21.8.3-fr}
Soient $K$ un corps de caractéristique $p>0$, $k_0$ un sous-corps de $K$,
$(k_\alpha)_{\alpha\in I}$ une famille filtrante décroissante de sous-corps
de $K$ contenant $k_0$. Les conditions suivantes sont équivalentes~:
\begin{enumerate}
  \item[\rm a)] $\bigcap_{\alpha\in I}k_\alpha(K^p)=k_0(K^p)$.
  \item[\rm b)] Pour toute extension $L$ de $K$ telle que
  $\Upsilon_{L/K/k_\alpha}$ soit un $L$-espace vectoriel de rang fini
  (ce qui a lieu en particulier si $L$ est une extension de type fini en
  vertu de (21.7.2)), il existe $\alpha\in I$ tel que $k_\alpha$ soit
  $k_0$-admissible pour l'extension $L$ de $K$.
  \item[\rm b$'$)] Pour toute extension $L=K(x)$ de $K$, avec $x^p\in K$,
  il existe $\alpha\in I$ tel que $k_\alpha$ soit $k_0$-admissible pour
  l'extension $L$ de $K$.
  \item[\rm c)] L'application canonique
\begin{equation}
\tag{21.8.3.1}
\label{0.21.8.3.1-fr}
\Omega^1_{K/k_0}\longrightarrow\varprojlim_{\alpha}\Omega^1_{K/k_\alpha}
\end{equation}
est injective.
\end{enumerate}
\end{theorem}

\begin{proof}
L'application canonique (21.8.3.1) est bien entendu obtenue par passage à la
limite projective dans le système projectif d'homomorphismes
$\Omega^1_{K/k_0}\to\Omega^1_{K/k_\alpha}$ (20.5.3.3). Nous allons prouver
le théorème suivant le schéma logique $c)\Rightarrow b)\Rightarrow b')
\Rightarrow a)\Rightarrow c)$.

Dire que $k_\alpha$ est $k_0$-admissible pour l'extension $L$ de $K$ signifie
que l'homomorphisme
\oldpage[0\textsubscript{IV-1}]{173}
canonique $\Upsilon_{L/K/k_0}\to\Upsilon_{L/K/k_\alpha}$ est injectif~; or,
si $N_\alpha$ est le noyau de cet homomorphisme, les $N_\alpha$ forment un
système filtrant décroissant de sous-espaces vectoriels de
$\Upsilon_{L/K/k_0}$, et comme $\Upsilon_{L/K/k_0}$ est par hypothèse de rang
fini, il revient au même de dire que l'un des $N_\alpha$ est nul ou que leur
intersection est nulle. Mais cette intersection n'est autre que le noyau de
l'homomorphisme limite projective
$\Upsilon_{L/K/k_0}\to\varprojlim_\alpha\Upsilon_{L/K/k_\alpha}$. Or, on a
le diagramme commutatif
\[
\xymatrix{
\Upsilon_{L/K/k_0}\ar[r]\ar[d]&\varprojlim\Upsilon_{L/K/k_\alpha}\ar[d]\\
\Omega^1_{K/k_0}\otimes_KL\ar[r]&
\varprojlim(\Omega^1_{K/k_\alpha}\otimes_KL)
}
\]
où la flèche verticale de gauche est injective par définition. Pour prouver
que $c)$ entraîne $b)$, il suffit donc de montrer que l'application canonique
\begin{equation}
\tag{21.8.3.2}
\label{0.21.8.3.2-fr}
\Omega^1_{K/k_0}\otimes_KV\longrightarrow
\varprojlim_{\alpha}(\Omega^1_{K/k_\alpha}\otimes_KV)
\end{equation}
est injective pour tout espace vectoriel $V$ sur $K$ (et en particulier pour
$V=L$). Or cela est évident si $V=K^n$ puisque alors
$W\otimes_KV=W^n$ pour tout espace vectoriel $W$ sur $K$ et que les produits
et limites projectives commutent. D'autre part, pour tout élément $z$ du
premier membre de (21.8.3.2), il existe un sous-espace $V'$ de $V$ de rang
fini tel que $z\in\Omega^1_{K/k_0}\otimes_KV'$ (identifié canoniquement à un
sous-espace de $\Omega^1_{K/k_0}\otimes_KV$)~; si $z\ne0$, son image dans
$\varprojlim_\alpha(\Omega^1_{K/k_\alpha}\otimes_KV')$ n'est donc pas nulle~;
comme le foncteur $\varprojlim$ est exact à gauche, l'image de $z$ dans le
second membre de (21.8.3.2) est donc aussi $\ne0$, ce qui achève de montrer
que $c)$ implique $b)$.

Il est trivial que $b)$ entraîne $b')$~; montrons que $b')$ entraîne $a)$.
Avec les notations de $b')$, on peut supposer que $L\ne K$. Il résulte alors
de (21.4.7) que l'on a $\operatorname{rg}\Upsilon_{L/K}\leqslant1$. Si
$\Upsilon_{L/K/k_\alpha}=0$, il n'y a rien à démontrer en vertu de
(21.6.1.1). Sinon, $\Upsilon_{L/K/k_\alpha}$ s'identifie canoniquement à
$\Upsilon_{L/K}$ et, si l'on pose $a=x^p$, $\Upsilon_{L/K}$ a une base
formée du seul élément $d_K(a)\otimes1$ (21.4.7)~;
$\Upsilon_{L/K/k_\alpha}$ a donc une base formée du seul élément
$d_{K/k_\alpha}(a)\otimes1$, et dire qu'un sous-corps $k_\alpha$ est
$k_0$-admissible pour l'extension $L$ de $K$ signifie que l'on a
$d_{K/k_\alpha}(a)\ne0$, ou encore (21.4.5) que
$a\notin k_\alpha(K^p)$. Or, pour tout $a\notin k_0(K^p)$, on a
$d_{K/k_0}(a)\ne0$ et $K(a^{1/p})\ne K$, donc on peut appliquer $b')$, qui
prouve l'existence d'un $\alpha$ tel que $a\notin k_\alpha(K^p)$~;
autrement dit $b')$ implique $a)$.

Reste à montrer que $a)$ implique $c)$. Soit $(x_\mu)_{\mu\in M}$ une
$p$-base de $K$ sur $k_0$~; alors, les $d_{K/k_0}(x_\mu)$ forment une base du
$K$-espace vectoriel $\Omega^1_{K/k_0}$ (21.4.5), et la condition $c)$
signifie que pour toute partie finie $J$ de l'ensemble d'indices $M$, il
existe un $\alpha\in I$ tel que les $d_{K/k_\alpha}(x_\mu)$ pour $\mu\in J$
soient linéairement indépendants dans $\Omega^1_{K/k_\alpha}$~; mais cela
signifie aussi (21.4.5) que les $x_\mu$ pour $\mu\in J$ forment une famille
$p$-libre sur $k_\alpha$, et l'existence d'un $\alpha$ ayant cette propriété
découle de (21.8.2) et de l'hypothèse $a)$.
\end{proof}
\oldpage[0\textsubscript{IV-1}]{174}

\begin{corollary}[21.8.4]
\label{0.21.8.4-fr}
Soient $K$ un corps de caractéristique $p>0$, $k_0$ un sous-corps de $K$,
$(k_\alpha)_{\alpha\in I}$ une famille filtrante décroissante de sous-corps
de $K$, contenant $k_0$, telle que
\[
  \bigcap_{\alpha\in I}k_\alpha(K^p)=k_0(K^p).
\]
Soient $L$ une extension de $K$ telle que $\Upsilon_{L/K/k_0}$ soit un
$L$-espace vectoriel de rang fini~; alors, pour tout corps $K'$ tel que
$K\subset K'\subset L$, il existe $\alpha\in I$ tel que $k_\alpha$ soit
$k_0$-admissible pour l'extension $L$ de $K'$.
\end{corollary}

\begin{proof}
Il suffit d'appliquer (21.8.3) et (21.6.6).
\end{proof}

\begin{corollary}[21.8.5]
\label{0.21.8.5-fr}
Soient $K$ un corps de caractéristique $p>0$, $k_0$ un sous-corps de $K$,
$(k_\alpha)_{\alpha\in I}$ une famille filtrante décroissante de sous-corps
de $K$ contenant $k_0$, telle que
\[
  \bigcap_{\alpha\in I}k_\alpha(K^p)=k_0(K^p).
\]
Alors, pour toute extension $L$ de $K$ telle que $\Upsilon_{L/K/k_0}$ soit
un $L$-espace vectoriel de rang fini, on a
\[
  \bigcap_{\alpha\in I}k_\alpha(L^p)=k_0(L^p).
\]
\end{corollary}

\begin{proof}
Supposons en effet que $M$ soit une extension de $L$ telle que
$\Upsilon_{M/L/k_0}$ soit un $M$-espace vectoriel de rang fini. La suite
exacte (21.6.1.1)
\[
  0\longrightarrow \Upsilon_{L/K/k_0}\otimes_L M
  \longrightarrow \Upsilon_{M/K/k_0}
  \longrightarrow \Upsilon_{M/L/k_0}
\]
et l'hypothèse montrent alors que $\Upsilon_{M/K/k_0}$ est aussi un
$M$-espace vectoriel de rang fini. Il existe donc un indice $\alpha$ tel que
$k_\alpha$ soit $k_0$-admissible pour l'extension $M$ de $K$ (21.8.3), donc
aussi pour l'extension $M$ de $L$ (21.6.6)~; cela ayant lieu pour toute
extension $M$ de $L$ telle que $\Upsilon_{M/L/k_0}$ soit de rang fini, le
corollaire résulte de l'équivalence de $a)$ et $b)$ dans (21.8.3).
\end{proof}

\begin{corollary}[21.8.6]
\label{0.21.8.6-fr}
Soient $K$ un corps, $p$ son exposant caractéristique, $k_0$ un sous-corps de
$K$. Si $L$ est une extension de $K$ telle que $\Upsilon_{L/K/k_0}$ soit un
$L$-espace vectoriel de rang fini, il existe un sous-corps $k$ de $K$,
contenant $k_0(K^p)$, tel que $[K:k]$ soit fini, et qui soit
$k_0$-admissible pour l'extension $L$ de $K$.
\end{corollary}

\begin{proof}
Il suffit en effet, en vertu de (21.8.3), de construire une famille filtrante
décroissante $(k_\alpha)_{\alpha\in I}$ de sous-corps de $K$, contenant
$K^p$ et $k_0$, pour lesquels $[K:k_\alpha]<+\infty$ et
$\bigcap_\alpha k_\alpha=k_0(K^p)$. Pour cela on considère une $p$-base
$(x_\lambda)_{\lambda\in J}$ de $K$ sur $k_0$ et, pour toute partie finie
$H$ de $J$, on considère le sous-corps $k_H$ de $K$ engendré par $k_0(K^p)$
et les $x_\lambda$ d'indice $\lambda\in J-H$~; il résulte de cette définition
que $(x_\lambda)_{\lambda\in H}$ est une $p$-base de $K$ sur $k_H$, et on en
conclut aussitôt que les $k_H$ vérifient les conditions voulues.
\end{proof}

\begin{remarks}[21.8.7]
\label{0.21.8.7-fr}
\begin{enumerate}
  \item[(i)] On a déjà vu (21.7.2) que si $L$ est une extension de
  \emph{type fini} de $K$, $\Upsilon_{L/K/k_0}$ est de rang fini pour tout
  sous-corps $k_0$ de $K$. Il en est de même si $L$ est une extension
  \emph{séparable} de $K$, car en vertu de (20.6.19), on a
  $\Upsilon_{L/K/k_0}=0$. Enfin, si $L$ est une \emph{extension de type fini
  d'une extension séparable $L_0$ de $K$}, le même raisonnement que dans
  (21.8.5) montre que $\Upsilon_{L/K/k_0}$ est encore de rang fini (et en fait
  est isomorphe à un sous-espace de $\Upsilon_{L/L_0/k_0}$).

  \oldpage[0\textsubscript{IV-1}]{175}
  \item[(ii)] Dans l'énoncé de (21.8.5), et par suite aussi dans celui de
  (21.8.3, $b)$), on ne peut omettre l'hypothèse que $\Upsilon_{L/K/k_0}$ est
  de rang fini sur $L$. Prenons pour $k_0$ le corps premier $\mathbf F_p$,
  pour $K$ un corps tel que $[K:K^p]$ soit infini dénombrable (par exemple le
  corps de fractions rationnelles $\mathbf F_p(X_1,\ldots,X_n,\ldots)$ à une
  infinité d'indéterminées)~; le procédé de construction de (21.8.6) montre
  aussitôt qu'il existe une suite infinie strictement décroissante $(K_n)$ de
  sous-corps de $K$ telle que $K_0=K$ et $\bigcap_nK_n=K^p$. Nous allons
  construire une suite croissante d'extensions finies $M_n$ de $K$, telle que
  si $M$ est la réunion des $M_n$, l'extension $L=M^{1/p}$ de $K$ mette
  (21.8.5) en défaut. Pour cela, soit $x$ un élément de $K-K^p$, posons
  $M_0=K^p$, et $M_n=K^p(a_1x,a_2x,\ldots,a_nx)$ pour $n\geqslant1$, où les
  $a_n$ sont construits par récurrence de façon que $a_n\in K_n$ et
  $a_{n+1}\notin M_n(x)$ pour tout $n$~: cela est possible, car $M_n(x)$ est
  de degré fini sur $K^p$, tandis qu'il n'en est pas de même de $K_n$. On en
  conclut aussitôt que $x\notin M=L^p$, mais comme
  $x=(a_nx)a_n^{-1}$, on a $x\in K_n(M)=K_n(L^p)$ pour tout $n$.
\end{enumerate}
\end{remarks}

\begin{proposition}[21.8.8]
\label{0.21.8.8-fr}
Soit $k$ un corps de caractéristique $p>0$, et soient
$A=k[[T_1,\ldots,T_r]]$ l'anneau des séries formelles à $r$ indéterminées sur
$k$, $K=k((T_1,\ldots,T_r))$ son corps des fractions. Alors il existe une
famille filtrante décroissante $(A_\alpha)$ de sous-anneaux noethériens de $A$
telle que $A$ soit un $A_\alpha$-module libre de type fini pour tout $\alpha$
et que, si $K_\alpha$ est le corps des fractions de $A_\alpha$, on ait
$\bigcap_\alpha K_\alpha=K^p$.
\end{proposition}

\begin{proof}
On peut écrire $K^p=k^p((T_1^p,\ldots,T_r^p))$~; on a vu dans la démonstration
de (21.8.6) qu'il existe une famille décroissante $(k_\alpha)$ de sous-corps
de $k$ telle que $[k:k_\alpha]$ soit fini pour tout $\alpha$ et
$\bigcap_\alpha k_\alpha=k^p$~; il est clair que si l'on pose
$A_\alpha=k_\alpha[[T_1^p,\ldots,T_r^p]]$, $A$ est un $A_\alpha$-module
libre de type fini~; tout revient donc à prouver la relation
\begin{equation}
\tag{21.8.8.1}
\label{0.21.8.8.1-fr}
  \bigcap_\alpha k_\alpha((T_1^p,\ldots,T_r^p))
  =k^p((T_1^p,\ldots,T_r^p)).
\end{equation}
Comme $\bigcap_\alpha k_\alpha=k^p$, cela va résulter des deux lemmes suivants.
\end{proof}

\begin{lemma}[21.8.8.2]
\label{0.21.8.8.2-fr}
Si $k'$ est une extension d'un corps $k$, on a
\[
  k'[[T_1,\ldots,T_r]]\cap k((T_1,\ldots,T_r))=k[[T_1,\ldots,T_r]].
\]
\end{lemma}

\begin{proof}
En effet, posons $C=k[[T_1,\ldots,T_r]]$, $D=k'[[T_1,\ldots,T_r]]$~; comme
$k((T_1,\ldots,T_r))$ est le corps des fractions de $C$, il suffira de
prouver que $D$ est un $C$-module \emph{fidèlement plat} (Bourbaki,
\emph{Alg. comm.}, chap.~I, \S~3, n\textsuperscript{o}~5, prop.~10). Or, $C$
et $D$ sont des anneaux locaux noethériens, et si $\mathfrak m$ est l'idéal
maximal de $C$, on a
$D/\mathfrak m^jD=(C/\mathfrak m^j)\otimes_k k'$, donc
$D/\mathfrak m^jD$ est un $(C/\mathfrak m^j)$-module plat~; il suffit donc
d'appliquer ($\mathbf 0_{\mathrm{III}}$, 10.2.1) et
($\mathbf 0_{\mathrm I}$, 6.6.2).
\end{proof}

\begin{lemma}[21.8.8.3]
\label{0.21.8.8.3-fr}
Soient $k$ un corps, $(k_\alpha)$ une famille filtrante décroissante de
sous-corps de $k$, et posons $k_0=\bigcap_\alpha k_\alpha$. On suppose qu'il
existe une puissance $q$ de l'exposant caractéristique de $k$ telle que
$k^q\subset k_0$. Alors on a
\begin{equation}
\tag{21.8.8.4}
\label{0.21.8.8.4-fr}
  \bigcap_\alpha k_\alpha((T_1,\ldots,T_r))=k_0((T_1,\ldots,T_r)).
\end{equation}
\end{lemma}

\begin{proof}
\oldpage[0\textsubscript{IV-1}]{176}
Il suffit de prouver qu'un élément $f\ne0$ du premier membre de (21.8.8.4)
appartient au second membre. Soient $\gamma$ un indice,
$g\in k_\gamma[[T_1,\ldots,T_r]]$ tel que
$gf\in k_\gamma[[T_1,\ldots,T_r]]$~; quitte à remplacer $g$ par $g^q$, on
peut supposer que $g\in k_0[[T_1,\ldots,T_r]]$. Alors, pour tout
$\alpha\geqslant\gamma$, on a
\[
  gf\in k_\gamma[[T_1,\ldots,T_r]]\cap k_\alpha((T_1,\ldots,T_r))
  =k_\alpha[[T_1,\ldots,T_r]]
\]
en vertu du lemme (21.8.8.2). Mais il est clair que l'intersection des
anneaux $k_\alpha[[T_1,\ldots,T_r]]$ n'est autre que
$k_0[[T_1,\ldots,T_r]]$, et l'on a donc bien
$f\in k_0((T_1,\ldots,T_r))$.
\end{proof}

\subsection{Modules de différentielles complétés dans les anneaux de séries formelles}
\label{subsection:0.21.9-fr}

Dans ce numéro, les corps ne sont plus nécessairement supposés être de
caractéristique $>0$.

\begin{lemma}[21.9.1]
\label{0.21.9.1-fr}
Soient $k$ un corps, $A$ un anneau local noethérien complet qui est une
$k$-algèbre, $K$ le corps résiduel de $A$. Si $K$ est une extension de type
fini de $k$, le $A$-module $\widehat\Omega^1_{A/k}$ est de type fini.
\end{lemma}

\begin{proof}
En vertu de (20.7.15), il suffit de prouver que $\Omega^1_{K/k}$ est un
$K$-espace vectoriel de rang fini. Or, par hypothèse, $K$ est le corps des
fractions d'une $k$-algèbre de type fini $B$~; comme $\Omega^1_{B/k}$ est un
$B$-module de type fini en vertu de (20.4.7), la conclusion résulte de
(20.5.9).
\end{proof}

\begin{proposition}[21.9.2]
\label{0.21.9.2-fr}
Soient $k_0$ un corps, $A$ un anneau local noethérien complet qui est une
$k_0$-algèbre formellement lisse, $\mathfrak m$ son idéal maximal,
$K=A/\mathfrak m$ son corps résiduel (de sorte qu'en tant qu'anneau $A$ est
isomorphe à un anneau de séries formelles $K[[T_1,\ldots,T_r]]$ (19.6.5)).
Si $K$ est une extension de type fini de $k_0$, alors
$\widehat\Omega^1_{A/k_0}$ est un $A$-module libre de rang égal à
$\dim(A)+\deg.\operatorname{tr}_{k_0}K$. En outre, pour tout sous-corps $k$
de $k_0$ tel que $\Omega^1_{k_0/k}$ soit de rang fini,
$\widehat\Omega^1_{A/k}$ est un $A$-module libre de rang égal à
\[
  \dim(A)+\deg.\operatorname{tr}_{k_0}K
  +\operatorname{rg}_{k_0}(\Omega^1_{k_0/k}).
\]
\end{proposition}

\begin{proof}
Il est clair que, si $P$ est le sous-corps premier de $k_0$, $K$ est une
$P$-algèbre formellement lisse (19.6.1). Notons d'autre part que l'on a
$\Upsilon^K_{A/k_0/P}=0$~: en effet, cela revient à voir que l'homomorphisme
\[
  u_{A/k_0/P}\otimes1_K:
  \Omega^1_{k_0/P}\otimes_{k_0}K
  \longrightarrow \Omega^1_{A/P}\otimes_A K
\]
est injectif (20.6.14.5), $u$ désignant l'homomorphisme structural
$k_0\to A$. Mais comme $A$ est une $k_0$-algèbre formellement lisse par
hypothèse, il résulte de (20.7.2) que $u_{A/k_0/P}$ est formellement
inversible à gauche, donc (19.1.7) $u_{A/k_0/P}\otimes1_K$ est inversible à
gauche et \emph{a fortiori} injectif, ce qui prouve notre assertion.
Appliquant la suite exacte (20.6.22.1), où $\Lambda,A,B,C$ sont remplacés par
$P,k_0,A,K$, il vient la suite exacte
\[
  0\longrightarrow \Upsilon^K_{A/k_0}
  \longrightarrow \mathfrak m/\mathfrak m^2
  \longrightarrow \Omega^1_{A/k_0}\otimes_A K
  \longrightarrow \Omega^1_{K/k_0}\longrightarrow0,
\]
de sorte que l'on a
\[
\begin{split}
  \operatorname{rg}_K(\Omega^1_{A/k_0}\otimes_AK)
  &={\operatorname{rg}}_K(\mathfrak m/\mathfrak m^2)
    +\bigl(\operatorname{rg}_K(\Omega^1_{K/k_0})
    -\operatorname{rg}_K(\Upsilon_{K/k_0})\bigr)\\
  &={\operatorname{rg}}_K(\mathfrak m/\mathfrak m^2)
    +\deg.\operatorname{tr}_{k_0}K
\end{split}
\]
en vertu de (21.7.1). D'autre part, $\widehat\Omega^1_{A/k_0}$ est un
$A$-module formellement projectif (20.4.9), et comme sa topologie est la
topologie $\mathfrak m$-préadique (20.4.5),
$\widehat\Omega^1_{A/k_0}/\mathfrak m^{j+1}\widehat\Omega^1_{A/k_0}$ est un
$(A/\mathfrak m^{j+1})$-module projectif pour tout $j$ (19.2.4), donc libre
($\mathbf 0_{\mathrm{III}}$, 10.1.2) et de rang

\oldpage[0\textsubscript{IV-1}]{177}
$n=\operatorname{rg}_K(\Omega^1_{A/k_0}\otimes_AK)$. En vertu de (21.9.1),
$\widehat\Omega^1_{A/k_0}$ est un $A$-module de type fini, et en outre ce
$A$-module est plat en vertu de ($\mathbf 0_{\mathrm{III}}$, 10.2.1)~; il est
donc libre ($\mathbf 0_{\mathrm{III}}$, 10.1.3) de rang $n$, d'où la première
assertion.

Comme par hypothèse $\Omega^1_{k_0/k}$ est de rang fini sur $k_0$, le produit
tensoriel complété
$\Omega^1_{k_0/k}\widehat\otimes_{k_0}A$ s'identifie à
$\Omega^1_{k_0/k}\otimes_{k_0}A$~; comme $A$ est une $k_0$-algèbre
formellement lisse, il résulte de (20.7.18) que l'homomorphisme
\[
  \widehat u_{A/k_0/k}:
  \Omega^1_{k_0/k}\otimes_{k_0}A
  \longrightarrow\widehat\Omega^1_{A/k}
\]
admet un inverse à gauche qui est un $A$-homomorphisme continu, ce qui
implique en particulier que son image est fermée dans
$\widehat\Omega^1_{A/k}$~; la suite exacte
\[
  0\longrightarrow \Omega^1_{k_0/k}\otimes_{k_0}A
  \xrightarrow{\widehat u_{A/k_0/k}}\widehat\Omega^1_{A/k}
  \longrightarrow\widehat\Omega^1_{A/k_0}\longrightarrow0
\]
est donc exacte et scindée en vertu de (20.7.17) et de ce qui précède~; ce
qui achève de prouver la proposition.
\end{proof}

\begin{corollary}[21.9.3]
\label{0.21.9.3-fr}
Soient $k$ un corps, $A$ l'anneau de séries formelles
$k[[T_1,\ldots,T_r]]$, muni de sa structure usuelle de $k$-algèbre. Alors
$\widehat\Omega^1_{A/k}$ est un $A$-module libre de rang $r=\dim(A)$.
\end{corollary}

\begin{proof}
Il suffit de noter que $A$ est une $k$-algèbre formellement lisse (19.3.4),
et d'appliquer (21.9.2) avec $k_0=k=K$.
\end{proof}

\begin{lemma}[21.9.4]
\label{0.21.9.4-fr}
Soient $k$ un corps de caractéristique $p>0$, $A$ une $k$-algèbre qui est un
anneau local noethérien complet, $B$ une sous-$k$-algèbre de $A$ isomorphe à
une $k$-algèbre topologique de la forme $k[[T_1,\ldots,T_r]]$ et telle que
$A$ soit une $B$-algèbre de type fini. Si $B_1$ est la sous-algèbre
$k[[T_1^p,\ldots,T_r^p]]$ de $B$, $\widehat\Omega^1_{A/k}$ s'identifie
canoniquement à $\Omega^1_{A/B_1}$.
\end{lemma}

\begin{proof}
Toute dérivation continue de $B_1$ dans un $A$-module topologique $P$, qui
est restriction d'une $k$-dérivation continue de $A$ dans $P$, est \emph{nulle},
car elle l'est dans le sous-anneau $k[T_1^p,\ldots,T_r^p]$ des polynômes, et
ce dernier est dense dans $B_1$. La suite exacte (20.3.7.1) montre donc que
l'homomorphisme canonique
\[
  \operatorname{D\acute{e}r.\,cont}_{B_1}(A,P)
  \longrightarrow \operatorname{D\acute{e}r.\,cont}_k(A,P)
\]
est bijectif~; la conclusion résulte de (20.7.14.4), compte tenu de ce que
$\Omega^1_{A/B_1}$ est un $A$-module de type fini (20.4.7), donc séparé et
complet puisque $A$ est un anneau local noethérien complet.
\end{proof}

\begin{proposition}[21.9.5]
\label{0.21.9.5-fr}
Soient $A$ un anneau local noethérien complet intègre, $A_0$ un sous-anneau
de $A$ isomorphe à un anneau de séries formelles $k[[T_1,\ldots,T_r]]$ sur
un corps $k$, tel que $A$ soit une $A_0$-algèbre finie et que le corps des
fractions $E$ de $A$ soit une extension séparable du corps des fractions
$L_0$ de $A_0$. Alors on a
\begin{equation}
\tag{21.9.5.1}
\label{0.21.9.5.1-fr}
  \operatorname{rg}_A\widehat\Omega^1_{A/k}
  =\operatorname{rg}_E(\widehat\Omega^1_{A/k}\otimes_AE)=\dim(A).
\end{equation}
\end{proposition}

\begin{proof}
Notons que si $\mathfrak m_0$ est l'idéal maximal de $A_0$, la topologie de
$A$ est la topologie $\mathfrak m_0$-adique puisque $A$ est une
$A_0$-algèbre finie (Bourbaki, \emph{Alg. comm.}, chap.~IV, \S~2,
n\textsuperscript{o}~5, cor.~3 de la prop.~9) et induit sur $A_0$ la topologie
$\mathfrak m_0$-adique (Bourbaki, \emph{Alg. comm.}, chap.~III, \S~3,
n\textsuperscript{o}~4, th.~3). On sait (21.9.1) que
$\widehat\Omega^1_{A/k}$ est un $A$-module de type fini et,

\oldpage[0\textsubscript{IV-1}]{178}
en vertu de (21.9.3) $\widehat\Omega^1_{A_0/k}$ est un $A_0$-module libre de
rang $r=\dim(A_0)=\dim(A)$ (16.1.5)~; le produit tensoriel
$\widehat\Omega^1_{A_0/k}\otimes_{A_0}A$ est donc un $A$-module libre de rang
$r$ identique à son séparé complété~; enfin $\widehat\Omega^1_{A/A_0}$ est
un $A$-module de type fini (20.4.7), donc identique à son séparé complété
($\mathbf 0_{\mathrm I}$, 7.3.6). Cela étant, dans la suite d'homomorphismes
\begin{equation}
\tag{21.9.5.2}
\label{0.21.9.5.2-fr}
  \widehat\Omega^1_{A_0/k}\otimes_{A_0}A
  \xrightarrow{u}\widehat\Omega^1_{A/k}
  \xrightarrow{v}\widehat\Omega^1_{A/A_0}\longrightarrow0
\end{equation}
on sait que $v$ est surjectif et que $\operatorname{Im}(u)$ est dense dans
$\operatorname{Ker}(v)$ (20.7.17)~; mais tout sous-$A_0$-module de
$\widehat\Omega^1_{A/k}$ est fermé pour la topologie $\mathfrak m_0$-adique
($\mathbf 0_{\mathrm I}$, 7.3.5), donc la suite (21.9.5.2) est exacte. Notons
d'autre part que puisque $A$ est entier sur $A_0$, $A_0$ intégralement clos
et $E$ séparable sur $L_0$, $\widehat\Omega^1_{A/A_0}$ est un $A$-module de
torsion (20.4.13, (iv))~; tensorisant la suite exacte (21.9.5.2) par $E$, il
vient une suite exacte
\[
  \widehat\Omega^1_{A_0/k}\otimes_{A_0}E
  \longrightarrow \widehat\Omega^1_{A/k}\otimes_AE\longrightarrow0
\]
d'où $\operatorname{rg}_E(\widehat\Omega^1_{A/k}\otimes_AE)\leqslant r$. Il
reste donc à montrer qu'il existe, dans $\widehat\Omega^1_{A/k}\otimes_AE$,
$r$ éléments linéairement indépendants sur $E$. Or, soient
$D_i=\partial/\partial T_i$ les dérivations canoniques de $A_0$ dans lui-même
$(1\leqslant i\leqslant r)$~; elles se prolongent de façon unique en des
dérivations (encore notées $D_i$) de $E$ dans lui-même, puisque $E$ est
extension séparable finie de $L_0$~; par restriction à $A$, ces dérivations
donnent des dérivations de $A$ dans $E$, et il est immédiat que ces
dérivations prennent leurs valeurs dans un même sous-$A$-module de type fini
de $E$~; comme elles sont continues dans $A_0$ et que la topologie de $A$ est
la topologie $\mathfrak m_0$-adique, on a ainsi formé $r$ dérivations
continues de $A$ dans un $A$-module topologique, qui sont évidemment
linéairement indépendantes puisque $D_i(T_j)=\delta_{ij}$~; ceci achève de
prouver la proposition.
\end{proof}

La proposition suivante est un analogue «~formel~» de (21.7.3)~:

\begin{proposition}[21.9.6]
\label{0.21.9.6-fr}
Soient $k_0\subset k\subset K_0$ trois corps de caractéristique $p>0$, tels
que $[K_0:k]<+\infty$~; on pose
\[
  L_0=K_0((T_1,\ldots,T_r)),\qquad
  M_0=K_0((T_1^p,\ldots,T_r^p)),
\]
\[
  L=k((T_1,\ldots,T_r)),\qquad
  M=k((T_1^p,\ldots,T_r^p)),\qquad
  N=k((T_1^{p^2},\ldots,T_r^{p^2})).
\]
Soit $E$ une extension de type fini de $L_0$.
\begin{enumerate}
  \item[(i)] Si $\Upsilon_{K_0/k_0}$ est de rang fini, on a
  \begin{equation}
  \tag{21.9.6.1}
  \label{0.21.9.6.1-fr}
    \operatorname{rg}_E\Omega^1_{E/M}
    -\operatorname{rg}_{K_0}\Omega^1_{K_0/k}
    =(r+\deg.\operatorname{tr}_{L_0}E)
      +d(E/K_0,k_0)+d(E/M_0,N/k_0)
  \end{equation}
  (notations de (21.6.1)), et en particulier on a
  \begin{equation}
  \tag{21.9.6.2}
  \label{0.21.9.6.2-fr}
    \operatorname{rg}_E\Omega^1_{E/M}
    -\operatorname{rg}_{K_0}\Omega^1_{K_0/k}
    \geqslant r+\deg.\operatorname{tr}_{L_0}E+d(E/K_0,k_0).
  \end{equation}
  En outre, si les deux membres de (21.9.6.2) sont égaux, ils le restent
  lorsqu'on remplace $k$ par un corps $k'$ tel que $k_0\subset k'\subset k$.

  \item[(ii)] On suppose en outre que $[k_0:k_0^p]<+\infty$. Soit
  $(k_\alpha)$ une famille filtrante décroissante de sous-corps de $K_0$
  contenant $k_0$, tels que $[K_0:k_\alpha]<+\infty$ pour tout $\alpha$ et
  que
  $\bigcap_\alpha k_\alpha(K_0^p)=k_0(K_0^p)$~; alors il existe un $\alpha$
  tel que, pour $k=k_\alpha$, les deux membres de (21.9.6.2) soient égaux.
\end{enumerate}
\end{proposition}

\begin{proof}
\begin{enumerate}
  \oldpage[0\textsubscript{IV-1}]{179}
  \item[(i)] On sait (21.1.5, (ii)) que
  $\operatorname{rg}_E\Omega^1_{E/M}$ ne change pas quand on remplace $M$ par
  $M(L_0^p)$~; comme
  $L_0^p=K_0^p((T_1^p,\ldots,T_r^p))$ et que
  $[k(K_0^p):k]<+\infty$, on a
  \[
    M(L_0^p)=k(K_0^p)((T_1^p,\ldots,T_r^p))~;
  \]
  de même $\Omega^1_{K_0/k}$ ne change pas quand on y remplace $k$ par
  $k(K_0^p)$, si bien que l'on peut \emph{supposer que $K_0^p\subset k$}, ce
  que nous ferons dans toute la suite. Nous introduirons aussi les corps
  \[
    k'_0=k_0(K_0^p),\qquad
    N'=k'_0((T_1^{p^2},\ldots,T_r^{p^2}))
  \]
  de sorte que l'on a le diagramme de corps
  \[
    \xymatrix{
      K_0\ar[r] & M_0\ar[r] & L_0\ar[r] & E\\
      k\ar[u]\ar[r] & N\ar[u]\ar[r] & M\ar[u]\ar[r] & L\ar[u]\\
      k_0\ar[u]\ar[r] & k'_0\ar[u]\ar[r] & N'\ar[u] &
    }
  \]

  Nous nous proposons d'évaluer la différence $\delta$ du premier et du
  second membre de (21.9.6.2).

  Il résulte de (21.7.3) que l'on a
  \[
    \operatorname{rg}_E\Omega^1_{E/M}
    -\operatorname{rg}_{M_0}\Omega^1_{M_0/M}
    =\deg.\operatorname{tr}_{M_0}E+d(E/M_0,M)
  \]
  et comme $L_0$ est une extension finie de $M_0$, on a
  $\deg.\operatorname{tr}_{L_0}E=\deg.\operatorname{tr}_{M_0}E$. D'autre
  part, une $p$-base de $K_0$ sur $k=k(K_0^p)$ est aussi une $p$-base de
  $M_0$ sur $M$, donc (21.4.5), on a
  \[
    \operatorname{rg}_{M_0}\Omega^1_{M_0/M}
    =\operatorname{rg}_{K_0}\Omega^1_{K_0/k}
  \]
  si bien que l'on a
  \begin{equation}
  \tag{21.9.6.3}
  \label{0.21.9.6.3-fr}
    \delta=d(E/M_0,M)-d(E/K_0,k_0)-r.
  \end{equation}

  Notons maintenant le lemme classique~:

  \begin{lemma}[21.9.6.4]
  \label{0.21.9.6.4-fr}
  Pour tout corps $k$, le corps de séries formelles
  $K=k((T_1,\ldots,T_r))$ est une extension séparable de $k$.
  \end{lemma}

  \begin{proof}
  Rappelons rapidement la démonstration de ce lemme pour être complet. Il
  suffit (Bourbaki, \emph{Alg.}, chap.~VIII, \S~7,
  n\textsuperscript{o}~3, démonstration du th.~1) de prouver que pour toute
  extension finie $k'$ de $k$, $K\otimes_k k'$ est sans élément nilpotent~;
  mais si l'on pose $A=k[[T_1,\ldots,T_r]]$ et $S=A-\{0\}$,
  $K\otimes_k k'$ est égal à $S^{-1}(A\otimes_k k')$, et
  $A\otimes_k k'$ s'identifie canoniquement à l'anneau intègre
  $A'=k'[[T_1,\ldots,T_r]]$ et $A$ à un sous-anneau de $A'$, d'où la
  conclusion.
  \end{proof}

  Utilisant ce lemme et (21.6.2, (ii)), on a $d(L_0/K_0,k_0)=0$, et la
  formule (21.6.4.4) donne donc
  \[
    d(E/K_0,k_0)=d(E/L_0,k_0).
  \]

  \oldpage[0\textsubscript{IV-1}]{180}
  Les formules (21.6.4.4) et (21.6.4.2) donnent d'autre part
  \[
    d(E/M_0,M)=d(E/L_0,M)+d(L_0/M_0,M),
  \]
  \[
    d(L_0/M_0,M)=d(L_0/M_0,M/N)+d(L_0/M_0,N),
  \]
  d'où
  \[
    \delta=d(E/L_0,M)+d(L_0/M_0,N)-d(E/L_0,k_0)
    +d(L_0/M_0,M/N)-r.
  \]

  Nous allons montrer que l'on a la relation
  \begin{equation}
  \tag{21.9.6.5}
  \label{0.21.9.6.5-fr}
    d(L_0/M_0,M/N)=r.
  \end{equation}
  Pour cela, remarquons que, par définition, on a
  \[
    d(L_0/M_0,M/N)
    =\operatorname{rg}_{L_0}\Upsilon_{L_0/M/N}
    -\operatorname{rg}_{M_0}\Upsilon_{M_0/M/N}.
  \]
  Compte tenu des deux suites exactes
  \[
  \begin{aligned}
    0&\longrightarrow\Upsilon_{L_0/M/N}
      \longrightarrow\Omega^1_{M/N}\otimes_ML_0
      \longrightarrow\Omega^1_{L_0/N}
      \longrightarrow\Omega^1_{L_0/M}\longrightarrow0,\\
    0&\longrightarrow\Upsilon_{M_0/M/N}
      \longrightarrow\Omega^1_{M/N}\otimes_MM_0
      \longrightarrow\Omega^1_{M_0/N}
      \longrightarrow\Omega^1_{M_0/M}\longrightarrow0,
  \end{aligned}
  \]
  il vient
  \[
  \begin{split}
    d(L_0/M_0,M/N)={}&\operatorname{rg}_{L_0}\Omega^1_{L_0/M}
      -\operatorname{rg}_{L_0}\Omega^1_{L_0/N}\\
      &-\operatorname{rg}_{M_0}\Omega^1_{M_0/M}
      +\operatorname{rg}_{M_0}\Omega^1_{M_0/N}.
  \end{split}
  \]
  Or, on a $N(L_0^p)\supset M$, donc
  $\Omega^1_{L_0/N}=\Omega^1_{L_0/M}$ (21.1.5, (ii)). D'autre part, les
  $T_i^p$ $(1\leqslant i\leqslant r)$ forment une $p$-base de $M$ sur $N$,
  et une $p$-base de $K_0$ sur $k$ est aussi une $p$-base de $M_0$ sur $M$,
  donc, comme $M_0^p\subset N$,
  \[
    \operatorname{rg}_{M_0}\Omega^1_{M_0/N}
    =r+\operatorname{rg}_{K_0}\Omega^1_{K_0/k}
    =r+\operatorname{rg}_{M_0}\Omega^1_{M_0/M}
  \]
  en vertu de (21.4.5) et (21.1.10), ce qui prouve (21.9.6.5).

  On a encore, par (21.6.4.2),
  \[
    d(E/L_0,M)=d(E/L_0,k_0)+d(E/L_0,M/k_0)
  \]
  et
  \[
    d(L_0/M_0,N)=d(L_0/M_0,N/k)+d(L_0/M_0,k).
  \]
  Mais comme $L_0$ est séparable sur $K_0$ (21.9.6.4), on a
  $d(L_0/K_0,k)=0$, et par suite aussi $d(L_0/M_0,k)=0$ ((21.6.2) et
  (21.6.4)). Il vient donc
  \[
    \delta=d(E/L_0,M/k_0)+d(L_0/M_0,N/k).
  \]
  Appliquons encore deux fois (21.6.4.2) au premier terme du second membre de
  cette formule~; on obtient
  \[
    d(E/L_0,M/k_0)=d(E/L_0,M/N)+d(E/L_0,N/k)+d(E/L_0,k/k_0).
  \]
  Comme $M\subset N(L_0^p)$, on a $d(E/L_0,M/N)=0$ (21.6.2), donc (par
  (21.6.4.4))
  \[
    \delta=d(E/L_0,k/k_0)+d(E/M_0,N/k).
  \]
  Mais $M_0$ et $L_0$ sont des extensions séparables de $K_0$ (21.9.6.4),
  donc
  \[
    d(L_0/K_0,k/k_0)=d(M_0/K_0,k/k_0)=0
  \]
  (21.6.2), et appliquant encore deux fois (21.6.4.4),

  \oldpage[0\textsubscript{IV-1}]{181}
  on a
  \[
    d(E/L_0,k/k_0)=d(E/K_0,k/k_0)=d(E/M_0,k/k_0),
  \]
  d'où, par une dernière application de (21.6.4.2),
  \[
    \delta=d(E/M_0,N/k_0),
  \]
  ce qui prouve (21.9.6.1). En outre, lorsque $k$ est remplacé par une
  sous-extension $k'$ de $k_0$, $N$ est remplacé par une sous-extension $N'$,
  donc $d(E/M_0,N'/k_0)\leqslant d(E/M_0,N/k_0)$ (21.6.4.2), ce qui démontre
  la dernière assertion de (i).

  \item[(ii)] Pour tout $\alpha$, posons
  $N_\alpha=k_\alpha((T_1^{p^2},\ldots,T_r^{p^2}))$. En vertu de (21.9.6.1)
  et de (21.8.3), il s'agit de montrer (compte tenu de ce que
  $K_0^p\subset k_\alpha$ et de ce que $E$ est une extension de type fini de
  $M_0$) que l'on a
  \[
    \bigcap_\alpha N_\alpha=k_0(M_0^p).
  \]
  Or, on a
  \begin{equation}
  \tag{21.9.6.6}
  \label{0.21.9.6.6-fr}
    \bigcap_\alpha N_\alpha=N'
    =(k_0(K_0^p))((T_1^{p^2},\ldots,T_r^{p^2}))
  \end{equation}
  en vertu de (21.8.8.3) et de la relation
  $\bigcap_\alpha k_\alpha=k_0(K_0^p)$ (21.8.6). D'autre part, on a
  $k_0(M_0^p)=k_0(K_0^p((T_1^{p^2},\ldots,T_r^{p^2})))$. Pour achever la
  démonstration de (21.9.6, (ii)), il reste donc à prouver le

  \begin{lemma}[21.9.6.7]
  \label{0.21.9.6.7-fr}
  Soient $k_0$ un corps de caractéristique $p>0$ tel que
  $[k_0:k_0^p]<+\infty$, $K_0$ une extension de $k_0$. Alors on a
  \begin{equation}
  \tag{21.9.6.8}
  \label{0.21.9.6.8-fr}
    (k_0(K_0^p))((T_1,\ldots,T_r))
    =k_0\bigl(K_0^p((T_1,\ldots,T_r))\bigr).
  \end{equation}
  \end{lemma}

  \begin{proof}
  Comme $k_0^p\subset K_0^p$, on a
  \[
    k_0^p\bigl(K_0^p((T_1,\ldots,T_r))\bigr)
    =K_0^p((T_1,\ldots,T_r)).
  \]
  Si $(c_i)_{1\leqslant i\leqslant n}$ est une base de $k_0$ sur $k_0^p$,
  il est immédiat que $(c_i)$ est aussi un système de générateurs de chacun
  des deux membres de (21.9.6.8) sur $K_0^p((T_1,\ldots,T_r))$.
  \end{proof}
\end{enumerate}
\end{proof}

\begin{remark}[21.9.7]
\label{0.21.9.7-fr}
L'assertion de (21.9.6, (ii)) \emph{ne s'étend pas} au cas où
$[k_0:k_0^p]$ est infini. Prenons en effet pour $k_0$ le corps
$\mathbf F_p(X_n)_{n\geqslant0}$ des fractions rationnelles à une infinité
d'indéterminées, de sorte que les $X_n$ forment une $p$-base de $k_0$ sur
$\mathbf F_p$, et par suite $[k_0:k_0^p]=+\infty$. Dans l'énoncé de (21.9.6),
prenons $K_0=k_0$ (donc nécessairement $k=k_0$), $L_0=k_0((T))$, pour $E$
l'extension $L_0(y)$, où $y$ est racine du polynôme
\[
  Y^p-\sum_{n=0}^{\infty}X_nT^{np}
\]
de $L_0[Y]$. Alors la différence $\delta$ des deux membres de (21.9.6.2) est
\emph{non nulle}. En effet, $E$ est \emph{séparable} sur $k_0$, sans quoi
$y^p$ serait élément de $k_0(L_0^p)$, et l'on voit aussitôt que cela n'est
pas possible (dû au fait qu'il y a une infinité de $X_n$ algébriquement
indépendants)~; on a donc $d(E/K_0,k_0)=0$, la formule (21.9.6.3) se réduit à
$\delta=d(E/M_0,M_0)-1$, et il est clair que
$d(E/M_0,M_0)=\operatorname{rg}_E\Upsilon_{E/M_0}$. Tout revient à vérifier
que ce rang est $2$. Or on a $E^p\subset M_0$, une $p$-base de $E$ sur $L_0$
est formée du seul élément $y$, et $L_0$ a évidemment une $p$-base sur $M_0$
formée du seul élément $T$~; donc, puisque $E$ est extension algébrique de
$M_0$, notre assertion résulte de (21.7.1) et (21.4.5).
\end{remark}

\begin{proposition}[21.9.8]
\label{0.21.9.8-fr}
Soient $k_0$ un corps (de caractéristique quelconque), $K_0$ une extension
séparable de $k_0$, $A$ un anneau local noethérien complet, qui est une
$k_0$-algèbre et dont le corps résiduel

\oldpage[0\textsubscript{IV-1}]{182}
est une extension finie de $K_0$. Pour tout idéal premier $\mathfrak p$ de
$A$, soit $k(\mathfrak p)$ le corps résiduel de $A$ en $\mathfrak p$. Alors,
pour tout corps $k$ tel que $k_0\subset k\subset K_0$ et tel que
$[K_0:k]<+\infty$, on a
\begin{equation}
\tag{21.9.8.1}
\label{0.21.9.8.1-fr}
\begin{split}
  &\operatorname{rg}_{k(\mathfrak p)}
    \bigl(\widehat\Omega^1_{(A/\mathfrak p)/k}
      \otimes_{A/\mathfrak p}k(\mathfrak p)\bigr)
    -\operatorname{rg}_{K_0}\Omega^1_{K_0/k}\\
  &\hspace{35mm}\geqslant
    \dim(A/\mathfrak p)
    +\operatorname{rg}_{k(\mathfrak p)}\Upsilon_{k(\mathfrak p)/k_0}.
\end{split}
\end{equation}
Supposons en outre que $[k_0:k_0^p]<+\infty$ (où $p$ est l'exposant
caractéristique de $k_0$). Alors on peut trouver un corps $k$ tel que
$k_0(K_0^p)\subset k\subset K_0$ et $[K_0:k]<+\infty$, pour lequel les deux
membres de (21.9.8.1) sont égaux.
\end{proposition}

\begin{proof}
Comme $A/\mathfrak p$ a même corps résiduel $K$ que $A$ et est complet, on
peut se borner au cas où $A$ est intègre et $\mathfrak p=0$~; on posera
alors $E=k(\mathfrak p)$, corps des fractions de $A$. Comme $K_0$ est
formellement lisse sur $k_0$ (19.6.1), il existe un $k_0$-monomorphisme
$K_0\to A$ faisant de $A$ une $K_0$-algèbre~; on sait en outre qu'il existe
une sous-$K_0$-algèbre $A_0$ de $A$, $K_0$-isomorphe à une algèbre de séries
formelles $K_0[[T_1,\ldots,T_r]]$ et telle que $A$ soit une $A_0$-algèbre
\emph{finie} (19.8.9), ce qui entraîne que $E$ est une extension finie du
corps des fractions $L_0=K_0((T_1,\ldots,T_r))$ de $A_0$. On a alors
$\Upsilon_{K_0/k_0}=0$ puisque $K_0$ est séparable sur $k_0$, et
$\deg.\operatorname{tr}_{L_0}E=0$~; en outre, comme $A$ est une
$A_0$-algèbre finie on a $\dim(A)=\dim(A_0)$ (16.1.5).

Si $p=1$, on a
$\operatorname{rg}_E(\widehat\Omega^1_{A/k}\otimes_AE)=r$ par (21.9.5) et
(17.1.4, (iii)), et $\Omega^1_{K_0/k}=0$ puisque $K_0$ est extension finie
séparable de $k$ (21.7.4)~; d'autre part $\Upsilon_{E/k_0}=0$, $E$ étant
extension séparable de $k_0$~; les deux membres de (21.9.8.1) sont donc égaux
dans ce cas.

Supposons maintenant que $p>1$~; si $B=k[[T_1,\ldots,T_r]]$, $A$ est une
$B$-algèbre finie, donc, en posant $B_1=k[[T_1^p,\ldots,T_r^p]]$,
$\widehat\Omega^1_{A/k}$ s'identifie à $\Omega^1_{A/B_1}$ (21.9.4), et en
désignant par $M$ le corps des fractions $k((T_1^p,\ldots,T_r^p))$, il
résulte de (20.5.9) que $\widehat\Omega^1_{A/k}\otimes_AE$ s'identifie à
$\Omega^1_{E/M}$~; l'inégalité (21.9.8.1) n'est autre alors que (21.9.6.2),
et la dernière assertion du corollaire résulte de (21.9.6, (ii)).
\end{proof}

% FIN SOURCE ega4/ega0-21-fr.tex

%
% DEBUT SOURCE ega4/ega0-22-fr.tex
\section{Critères différentiels de lissité formelle et de régularité}
\label{section:0.22-fr}

Les principaux résultats de ce paragraphe sont~:
\begin{enumerate}
  \item[a)] Les critères (22.2.2) et (22.5.8), qui donnent des conditions
  nécessaires et suffisantes pour qu'un anneau local noethérien complet $A$
  contenant un corps $k$ soit formellement lisse sur $k$~; le critère
  (22.5.8) s'énonce sans faire intervenir de notions différentielles, et
  complète (19.6.6)~; l'énoncé (22.2.2), de nature un peu plus technique,
  permet de donner une \emph{classification}, pour $k$ fixé, des anneaux
  locaux $A$ considérés (22.2.5)~: ils sont donnés par leur dimension $n$ et
  leur corps résiduel $K$, sous la seule condition
  $n\geqslant\operatorname{rg}_K\Upsilon_{K/k}$.

  \item[b)] Le théorème (22.3.3), et le corollaire (22.7.6) du critère
  jacobien de Nagata (22.7.3), qui donnent des conditions permettant
  d'affirmer qu'un anneau \emph{localisé} d'un anneau local noethérien complet
  $A$ contenant un corps $k$ est géométriquement régulier sur $k$~; ils
  joueront un rôle important dans l'étude des propriétés «~fines~» des
  anneaux locaux faite au chap.~IV, §~7. On ne possède pour le moment aucune
  démonstration de (22.7.6) indépendante du critère jacobien de Nagata.

  \item[c)] Le critère jacobien de Zariski (22.6.7) et ses variantes, qui
  sont des conséquences faciles de (20.7.8).
\end{enumerate}

\oldpage[0\textsubscript{IV-1}]{183}

\subsection{Relèvement de la lissité formelle}
\label{subsection:0.22.1-fr}

\begin{theorem}[22.1.1]
\label{0.22.1.1-fr}
Soient $s:\Lambda\to A$, $u:A\to B$ deux homomorphismes d'anneaux,
$\mathfrak J$ un idéal de $B$, $C$ l'anneau quotient $B/\mathfrak J$. On
suppose que $C$ est une $\Lambda$-algèbre formellement lisse pour la topologie
discrète.
\begin{enumerate}
  \item[(i)] Supposons vérifiées les conditions suivantes~:
  \begin{enumerate}
    \item[$\alpha$)] $\mathfrak J/\mathfrak J^2$ est un $C$-module projectif.
    \item[$\beta$)] L'homomorphisme canonique
    $\mathbf S_C^*(\mathfrak J/\mathfrak J^2)\to
    \operatorname{gr}_{\mathfrak J}^*(B)$ (19.5.2) est bijectif.
    \item[$\gamma$)] L'homomorphisme caractéristique
    $\chi:\Upsilon_{C/B/\Lambda}\to\mathfrak J/\mathfrak J^2$ (20.6.20)
    est injectif.
    \item[$\delta$)] Le $C$-module
    $\Omega^1_{B/A}\otimes_BC$ est projectif.
  \end{enumerate}
  Alors $B$, muni de la topologie $\mathfrak J$-préadique, est une
  $\Lambda$-algèbre formellement lisse.

  \item[(ii)] Inversement, supposons que $B$ soit une $\Lambda$-algèbre
  formellement lisse pour la topologie $\mathfrak J$-préadique, et que $A$
  soit une $\Lambda$-algèbre formellement lisse pour la topologie discrète.
  Alors les conditions $\alpha$), $\beta$), $\gamma$), $\delta$) de (i) sont
  vérifiées.
\end{enumerate}
\end{theorem}

\begin{proof}
En vertu de la suite exacte (20.6.22.1) (qui est applicable, puisque
$B/\mathfrak J^2$ est une extension $\Lambda$-triviale de $C$ par
$\mathfrak J/\mathfrak J^2$, $C$ étant supposée être une $\Lambda$-algèbre
formellement lisse (19.4.4)), la condition $\gamma$) équivaut à
$\Upsilon^C_{B/A/\Lambda}=0$, ou encore à la suivante~:
\begin{enumerate}
  \item[$\gamma'$)] L'homomorphisme
  \[
    u_{B/A/\Lambda}\otimes 1_C:
    \Omega^1_{A/\Lambda}\otimes_A C
    \longrightarrow \Omega^1_{B/\Lambda}\otimes_B C
  \]
  est injectif.
\end{enumerate}

(i) Les conditions $\alpha$) et $\beta$) entraînent que $B$ est une
$\Lambda$-algèbre formellement lisse pour la topologie
$\mathfrak J$-préadique (19.5.4). Il revient donc au même de dire que $B$ est
(pour la topologie $\mathfrak J$-préadique) une $A$-algèbre formellement
lisse, ou une $A$-algèbre formellement lisse \emph{relativement à $\Lambda$}.
Pour voir que $B$ est une $A$-algèbre formellement lisse, il suffit donc, en
vertu de (20.7.2), de prouver que l'homomorphisme continu de $B$-modules
topologiques
\[
  u_{B/A/\Lambda}:\Omega^1_{A/\Lambda}\otimes_A B
  \longrightarrow\Omega^1_{B/\Lambda}
\]
est \emph{formellement inversible à gauche}. Or, puisque $B$ est une
$\Lambda$-algèbre formellement lisse, le $B$-module topologique
$\Omega^1_{B/\Lambda}$ est formellement projectif (20.4.9) et sa topologie
est déduite de celle de $B$ (20.4.5). En vertu de (19.1.9), il suffit donc,
pour que $u_{B/A/\Lambda}$ soit formellement inversible à gauche, que
$u_{B/A/\Lambda}\otimes 1_C$ soit inversible à gauche. Mais en vertu de
$\gamma'$), cette dernière application est injective, donc on a la suite
exacte (20.6.14.7)
\[
  0\longrightarrow\Omega^1_{A/\Lambda}\otimes_A C
  \longrightarrow\Omega^1_{B/\Lambda}\otimes_B C
  \longrightarrow\Omega^1_{B/A}\otimes_B C\longrightarrow0.
\]
Enfin, en vertu de $\delta$), le $C$-module
$\Omega^1_{B/A}\otimes_BC$ est projectif, donc la suite exacte précédente
est scindée, ce qui achève de prouver (i).

\oldpage[0\textsubscript{IV-1}]{184}

(ii) Si $B$ est une $\Lambda$-algèbre formellement lisse pour la topologie
$\mathfrak J$-préadique, et $A$ une $\Lambda$-algèbre formellement lisse pour
la topologie discrète, $B$ est une $A$-algèbre formellement lisse pour la
topologie $\mathfrak J$-préadique (19.3.5, (ii))~; les conditions $\alpha$) et
$\beta$) résultent donc de (19.5.6). En outre, (20.7.2) montre que
$u_{B/A/\Lambda}$ est formellement inversible à gauche, donc
$u_{B/A/\Lambda}\otimes 1_C$ est inversible à gauche (19.1.7), et \emph{a
fortiori} injectif. Enfin, $\Omega^1_{B/A}$ est un $B$-module formellement
projectif (20.4.9), et par suite $\Omega^1_{B/A}\otimes_BC$ est un $C$-module
projectif (19.2.4).
\end{proof}

\begin{corollary}[22.1.2]
\label{0.22.1.2-fr}
Soient $s:\Lambda\to A$, $u:A\to B$ deux homomorphismes d'anneaux,
$\mathfrak m$ un idéal maximal de $B$, $K=B/\mathfrak m$ le corps quotient.
On suppose que les $\Lambda$-algèbres $A$ et $K$ soient formellement lisses
pour la topologie discrète. Alors les conditions suivantes sont
équivalentes~:
\begin{enumerate}
  \item[a)] $B$ est une $A$-algèbre formellement lisse pour la topologie
  $\mathfrak m$-préadique.
  \item[b)] L'homomorphisme canonique
  \begin{equation}
  \tag{22.1.2.1}
  \label{0.22.1.2.1-fr}
    \mathbf S_K^*(\mathfrak m/\mathfrak m^2)
    \longrightarrow\operatorname{gr}_{\mathfrak m}^*(B)
  \end{equation}
  est bijectif, et l'homomorphisme caractéristique
  \begin{equation}
  \tag{22.1.2.2}
  \label{0.22.1.2.2-fr}
    \chi:\Upsilon_{K/B/\Lambda}\longrightarrow\mathfrak m/\mathfrak m^2
  \end{equation}
  est injectif.
\end{enumerate}
\end{corollary}

\begin{proof}
Cela résulte aussitôt de (22.1.1), les conditions $\alpha$) et $\delta$)
étant ici trivialement vérifiées, puisque $K$ est un corps.
\end{proof}

\begin{remark}[22.1.3]
\label{0.22.1.3-fr}
Supposons vérifiées les hypothèses de (22.1.2) et en outre supposons que
$\mathfrak m/\mathfrak m^2$ soit un $K$-espace vectoriel de rang \emph{fini}.
Alors, pour que $B$ soit une $A$-algèbre formellement lisse pour la topologie
$\mathfrak m$-préadique, il suffit que les conditions suivantes soient
satisfaites~:

\begin{env}[22.1.3.1]
\label{0.22.1.3.1-fr}
Étant donnés une algèbre de polynômes $E=K[T_1,\ldots,T_n]$, l'idéal maximal
$\mathfrak n$ de $E$ engendré par les $T_i$ $(1\leqslant i\leqslant n)$, et
un idéal $\mathfrak b\supset\mathfrak n^k$ dans $E$, tout
$\Lambda$-homomorphisme $v:B\to E/\mathfrak b$ qui, par passage aux
quotients, donne l'identité $K\to K$, se factorise en
\[
  B\xrightarrow{w}E/\mathfrak n^k\longrightarrow E/\mathfrak b,
\]
où $w$ est un $\Lambda$-homomorphisme.
\end{env}

\begin{env}[22.1.3.2]
\label{0.22.1.3.2-fr}
Si $F$ est l'extension triviale $D_K(K)$ («~algèbre des nombres duaux sur
$K$~»), $\rho:A\to F$ un homomorphisme d'anneaux tel que le composé
$A\xrightarrow{\rho}F\to K$ soit l'homomorphisme canonique, alors tout
$\Lambda$-homomorphisme $v:B\to K$ qui, par passage aux quotients, donne
l'identité $K\to K$, se factorise en
\[
  B\xrightarrow{w}F\longrightarrow K,
\]
où $w$ est un $\Lambda$-homomorphisme (pour la structure de $A$-algèbre sur
$F$ définie par $\rho$).
\end{env}

En effet, on a vu (19.5.5) que la condition (22.1.3.1) entraîne que
l'homomorphisme canonique (22.1.2.1) est bijectif. En vertu de (22.1.2), il
suffit ensuite de voir que l'homomorphisme canonique
\[
  u_{B/A/\Lambda}\otimes1_K:
  \Omega^1_{A/\Lambda}\otimes_AK
  \longrightarrow\Omega^1_{B/\Lambda}\otimes_BK
\]
est injectif, ou encore (puisqu'il s'agit d'espaces vectoriels sur $K$) que
pour tout $K$-espace vectoriel $L$, l'homomorphisme canonique
$\operatorname{D\acute{e}r}_\Lambda(B,L)\to\operatorname{D\acute{e}r}_\Lambda(A,L)$ est
\emph{surjectif} (en vertu de (20.4.8))~; il est clair que pour cela, il faut
et il suffit que l'homomorphisme canonique
$\operatorname{D\acute{e}r}_\Lambda(B,K)\to\operatorname{D\acute{e}r}_\Lambda(A,K)$ soit
surjectif, d'où notre assertion, en vertu de (20.1.5).

On notera que les deux conditions (22.1.3.1) et (22.1.3.2) sont entraînées
par la suivante~:

\begin{env}[22.1.3.3]
\label{0.22.1.3.3-fr}
Pour toute $\Lambda$-algèbre locale artinienne $E$ de corps résiduel égal à
$K$, et tout idéal nilpotent $\mathfrak r$ de $E$, tout
$\Lambda$-homomorphisme $B\to E/\mathfrak r$ qui, par passage aux quotients,
donne l'identité $K\to K$, se factorise en
$B\xrightarrow{u}E\to E/\mathfrak r$, où $u$ est un
$\Lambda$-homomorphisme.
\end{env}
\end{remark}

\oldpage[0\textsubscript{IV-1}]{185}

Ce résultat se généralise comme suit~:

\begin{proposition}[22.1.4]
\label{0.22.1.4-fr}
Soit $\varphi:A\to B$ un homomorphisme local d'anneaux locaux noethériens,
$k$ (resp. $K$) le corps résiduel de $A$ (resp. $B$). Pour que $B$ soit une
$A$-algèbre formellement lisse (pour les topologies préadiques), il faut et
il suffit que pour tout anneau local artinien $C$, de corps résiduel égal à
$K$, tout homomorphisme local $A\to C$ (faisant de $C$ une $A$-algèbre
topologique), et tout idéal nilpotent $\mathfrak J$ de $C$, tout
$A$-homomorphisme local $B\to C/\mathfrak J$ tel que l'homomorphisme
$K\to K$ obtenu par passage aux quotients soit l'identité, se factorise en
\[
  B\xrightarrow{f}C\longrightarrow C/\mathfrak J,
\]
où $f$ est un $A$-homomorphisme local.
\end{proposition}

\begin{proof}
La condition étant nécessaire par définition (19.3.1), montrons qu'elle est
suffisante. Comme les $\widehat A$-homomorphismes locaux de $\widehat B$ dans
une $\widehat A$-algèbre locale discrète proviennent par extension des
$A$-homomorphismes locaux de $B$ dans cette algèbre, on peut se borner au cas
où $A$ et $B$ sont \emph{complets}, en vertu de (19.3.6). Lorsque $A$ est un
\emph{corps} $k$, la proposition résulte de (22.1.3.3), où l'on prend pour
$\Lambda$ le sous-corps premier de $k$, et de (19.6.1).

Dans le cas général, posons $B_0=B\otimes_Ak=B/\mathfrak mB$
($\mathfrak m$ idéal maximal de $A$), qui est complet. Si $C_0$ est une
$k$-algèbre locale artinienne de corps résiduel égal à $K$, $\mathfrak J_0$
un idéal nilpotent de $C_0$, tout $k$-homomorphisme local
$g_0:B_0\to C_0/\mathfrak J_0$ donnant l'identité sur $K$ par passage aux
quotients, donne par composition un $A$-homomorphisme local
\[
  g:B\longrightarrow B_0\xrightarrow{g_0}C_0/\mathfrak J_0
\]
qui par hypothèse se factorise donc en
$B\xrightarrow{f}C_0\to C_0/\mathfrak J_0$, où $f$ est un
$A$-homomorphisme local~; mais comme $C_0$ est une $k$-algèbre, $f$ se
factorise en $B\to B_0\xrightarrow{f_0}C_0$, où $f_0$ est un
$k$-homomorphisme local. On voit donc que les hypothèses de l'énoncé sont
encore remplies lorsqu'on y substitue $B_0$ et $k$ à $B$ et $A$
respectivement~; donc $B_0$ est une $k$-algèbre \emph{formellement lisse}
d'après ce que l'on vient de voir.

Appliquant (19.7.2) et (19.7.1), on voit qu'il existe un anneau local
noethérien complet $B'$, un homomorphisme local $A\to B'$ faisant de $B'$ un
$A$-module \emph{plat} et une $A$-algèbre \emph{formellement lisse}, et un
$k$-isomorphisme
\[
  u'_0:B'_0=B'\otimes_Ak\xrightarrow{\sim}B_0.
\]
Soit $v'_0$ l'isomorphisme réciproque de $u'_0$~; nous nous proposons de
montrer qu'il existe un $A$-homomorphisme local $v:B\to B'$ tel que $v'_0$
provienne de $v$ par passage aux quotients. Pour cela, désignons par
$\mathfrak n$ et $\mathfrak n'$ les idéaux maximaux respectifs de $B$ et
$B'$, et, pour tout $j$, soit
\[
  w_j:B/(\mathfrak n^{j+1}+\mathfrak mB)
  \longrightarrow B'/(\mathfrak n'^{j+1}+\mathfrak mB')
\]
l'homomorphisme déduit de $v'_0$ par passage aux quotients~; il suffira de
former pour tout $j$ un $A$-homomorphisme local
$v_j:B/\mathfrak n^{j+1}\to B'/\mathfrak n'^{j+1}$ tel que les diagrammes
\[
\xymatrix@C=4em@R=3em{
  B/\mathfrak n^{j+1} \ar[r]^{v_j} \ar[d] &
  B'/\mathfrak n'^{j+1} \ar[d] \\
  B/\mathfrak n^j \ar[r]_{v_{j-1}} & B'/\mathfrak n'^j
}
\qquad
\xymatrix@C=4em@R=3em{
  B/\mathfrak n^{j+1} \ar[r]^{v_j} \ar[d] &
  B'/\mathfrak n'^{j+1} \ar[d] \\
  B/(\mathfrak n^{j+1}+\mathfrak mB) \ar[r]_{w_j} &
  B'/(\mathfrak n'^{j+1}+\mathfrak mB')
}
\]
soient commutatifs~; $B$ et $B'$ étant complets, l'homomorphisme
$v=\varprojlim v_j$ répondra à la question.

\oldpage[0\textsubscript{IV-1}]{186}

Or, la formation par récurrence de $v_j$ résulte de l'hypothèse sur $B$ et
du même raisonnement que dans (19.3.10.1), en utilisant le lemme
(19.3.10.2) et le fait que
\[
  \mathfrak n'^j+(\mathfrak n'^{j+1}+\mathfrak mB')
  =\mathfrak n'^j+\mathfrak mB'.
\]
Il résulte alors de (19.7.1.4) que $v$ est un $A$-\emph{isomorphisme}, car
$B'$ est un $A$-module plat, et $B$ est complet pour la topologie
$\mathfrak m$-préadique, les idéaux $\mathfrak m^iB$ étant fermés dans $B$
pour la topologie $\mathfrak n$-adique ($\mathbf 0_{\mathrm I}$, 7.3.5).
C.Q.F.D.
\end{proof}

\subsection{Caractérisation différentielle des algèbres locales formellement
lisses sur un corps}
\label{subsection:0.22.2-fr}

\begin{env}[22.2.1]
\label{0.22.2.1-fr}
Soient $k$ un corps, $P$ son sous-corps premier, $A$ une $k$-algèbre qui est
un anneau \emph{local}, $\mathfrak m$ son idéal maximal, $K=A/\mathfrak m$
son corps résiduel. Comme $K$ est séparable sur $P$, $K$ est une
$P$-algèbre formellement lisse pour la topologie discrète (19.6.1), et par
suite la $P$-extension $A/\mathfrak m^2$ de $K$ par
$\mathfrak m/\mathfrak m^2$ est $P$-\emph{triviale}, et l'homomorphisme
caractéristique
\begin{equation}
\tag{22.2.1.1}
\label{0.22.2.1.1-fr}
  \chi_{A/k}:\Upsilon_{K/k}\longrightarrow\mathfrak m/\mathfrak m^2
\end{equation}
est donc défini (20.6.23).
\end{env}

\begin{theorem}[22.2.2]
\label{0.22.2.2-fr}
Sous les conditions de (22.2.1), pour que $A$ soit une $k$-algèbre
formellement lisse pour la topologie $\mathfrak m$-préadique, il faut et il
suffit que les conditions suivantes soient satisfaites~:
\begin{enumerate}
  \item[(i)] L'homomorphisme canonique
  $\mathbf S_K^*(\mathfrak m/\mathfrak m^2)\to
  \operatorname{gr}_{\mathfrak m}^*(A)$ est bijectif.
  \item[(ii)] L'homomorphisme caractéristique
  $\chi_{A/k}:\Upsilon_{K/k}\to\mathfrak m/\mathfrak m^2$ est injectif.
\end{enumerate}
\end{theorem}

\begin{proof}
C'est un cas particulier de (22.1.2), où l'on remplace $\Lambda$ par $P$,
$A$ par $k$ et $B$ par $A$~; $K$ et $k$, étant séparables sur $P$, sont en
effet des $P$-algèbres formellement lisses (19.6.1).
\end{proof}

\begin{remark}[22.2.3]
\label{0.22.2.3-fr}
La condition (ii) de (22.2.2) peut être remplacée par une quelconque des
suivantes~:
\begin{enumerate}
  \item[a)] L'homomorphisme canonique
  $\Omega^1_k\otimes_kK\to\Omega^1_A\otimes_AK$ est injectif.
  \item[b)] Pour tout $n$, l'homomorphisme canonique
  \[
    \Omega^1_k\otimes_k(A/\mathfrak m^{n+1})
    \longrightarrow
    \Omega^1_A\otimes_A(A/\mathfrak m^{n+1})
  \]
  est inversible à gauche.
  \item[c)] L'homomorphisme canonique
  $\Omega^1_k\widehat\otimes_kA\to\widehat\Omega^1_A$ est inversible à
  gauche.
  \item[d)] $A$ est une $k$-algèbre formellement lisse (pour la topologie
  $\mathfrak m$-préadique) \emph{relativement au corps premier}.
\end{enumerate}
En effet, on a déjà remarqué dans (22.1.1) que l'injectivité de $\chi_{A/k}$
est équivalente à $a)$, et que la conjonction de $a)$ et de la condition (i)
de (22.2.2) équivaut à celle de $d)$ et de (i). Enfin, on a aussi vu dans
(22.1.1) que $b)$ et $d)$ sont équivalentes, et l'équivalence de $b)$ et
$c)$ résulte de (19.1.8).
\end{remark}
\begin{env}[22.2.4]
\label{0.22.2.4-fr}
La principale application de (22.2.2) a trait aux $k$-algèbres locales
noethériennes et \emph{complètes} (cas où $\mathfrak m/\mathfrak m^2$ est un
$K$-espace vectoriel de rang \emph{fini}). De façon précise, soient $k$ un
corps, $K$ une extension de $k$, $V$ un $K$-espace vectoriel de

\oldpage[0\textsubscript{IV-1}]{187}

rang fini~; nous considérons les triplets $(A,\alpha,\beta)$, où $A$ est une
$k$-algèbre locale noethérienne complète \emph{formellement lisse} d'idéal
maximal $\mathfrak m$, $\alpha:A/\mathfrak m\xrightarrow{\sim}K$ un
isomorphisme de $k$-algèbres,
$\beta:\mathfrak m/\mathfrak m^2\xrightarrow{\sim}V$ un di-isomorphisme
d'espaces vectoriels (pour l'isomorphisme $\alpha$). Étant donnés deux tels
triplets $(A,\alpha,\beta)$, $(A',\alpha',\beta')$, une \emph{équivalence} du
premier sur le second est un $k$-isomorphisme $u:A\xrightarrow{\sim}A'$ tel
que (si $\mathfrak m'$ est l'idéal maximal de $A'$), les diagrammes
\[
\xymatrix@C=4em@R=3em{
  A/\mathfrak m \ar[r]^{\operatorname{gr}^0(u)} \ar[d]_{\alpha} &
  A'/\mathfrak m' \ar[d]^{\alpha'} \\
  K \ar[r]_{1_K} & K
}
\qquad
\xymatrix@C=4em@R=3em{
  \mathfrak m/\mathfrak m^2 \ar[r]^{\operatorname{gr}^1(u)}
    \ar[d]_{\beta} &
  \mathfrak m'/\mathfrak m'^2 \ar[d]^{\beta'} \\
  V \ar[r]_{1_V} & V
}
\]
sont commutatifs. Il est immédiat que les \emph{classes d'équivalence} des
triplets $(A,\alpha,\beta)$ forment un \emph{ensemble} que nous noterons
$\operatorname{FL}(K/k,V)$. Remarquons maintenant qu'à tout triplet
$(A,\alpha,\beta)$ est associé le $K$-homomorphisme composé d'espaces
vectoriels
\[
  \Upsilon_{K/k}\xrightarrow{\chi_{A/k}}
  \mathfrak m/\mathfrak m^2\xrightarrow{\beta}V
\]
(où $\chi_{A/k}$ est considéré comme un di-homomorphisme relatif à
$\alpha^{-1}$)~; la définition précédente prouve (par (20.6.22.2)) que ce
$K$-homomorphisme \emph{ne dépend que de la classe d'équivalence} de
$(A,\alpha,\beta)$ dans $\operatorname{FL}(K/k,V)$~; autrement dit on a
défini une application canonique
\begin{equation}
\tag{22.2.4.1}
\label{0.22.2.4.1-fr}
  \operatorname{FL}(K/k,V)\longrightarrow
  \operatorname{Hom}_K(\Upsilon_{K/k},V).
\end{equation}
\end{env}

\begin{proposition}[22.2.5]
\label{0.22.2.5-fr}
Soient $k$ un corps, $K$ une extension de $k$, $V$ un $K$-espace vectoriel
de rang fini. L'application canonique (22.2.4.1) est une bijection de
$\operatorname{FL}(K/k,V)$ sur l'ensemble des $K$-homomorphismes injectifs
de $\Upsilon_{K/k}$ dans $V$.
\end{proposition}

\begin{proof}
(i) Pour montrer que (22.2.4.1) est injective, il faut prouver ce qui suit~:
soient $A$, $A'$ deux $k$-algèbres locales noethériennes complètes
formellement lisses, $\mathfrak m$, $\mathfrak m'$ leurs idéaux maximaux
respectifs, $\alpha:A/\mathfrak m\xrightarrow{\sim}K$,
$\alpha':A'/\mathfrak m'\xrightarrow{\sim}K$ deux $k$-isomorphismes~;
supposons donnés deux $k$-isomorphismes
\[
  v^0:A/\mathfrak m\xrightarrow{\sim}A'/\mathfrak m',\qquad
  v^1:\mathfrak m/\mathfrak m^2\xrightarrow{\sim}
      \mathfrak m'/\mathfrak m'^2
\]
tels que les diagrammes
\begin{equation}
\tag{22.2.5.1}
\label{0.22.2.5.1-fr}
\xymatrix@C=4em@R=3em{
  A/\mathfrak m \ar[r]^{v^0} \ar[dr]_{\alpha} &
  A'/\mathfrak m' \ar[dl]^{\alpha'} \\
  & K
}
\qquad
\xymatrix@C=4em@R=3em{
  \mathfrak m/\mathfrak m^2 \ar[r]^{v^1} &
  \mathfrak m'/\mathfrak m'^2 \\
  \Upsilon_{K/k} \ar[u]^{\chi_A} \ar[ur]_{\chi_{A'}} &
}
\end{equation}
soient commutatifs. Il s'agit de prouver qu'il existe un $k$-isomorphisme
$u:A\xrightarrow{\sim}A'$ tel que $\operatorname{gr}^0(u)=v^0$ et
$\operatorname{gr}^1(u)=v^1$. Remarquons en premier lieu que puisque $K$ est
un corps, l'homomorphisme canonique (20.6.7.1) est \emph{bijectif} (20.6.11),
donc il existe déjà un $k$-isomorphisme
$v:A/\mathfrak m^2\xrightarrow{\sim}A'/\mathfrak m'^2$ tel que
$\operatorname{gr}^0(v)=v^0$ et $\operatorname{gr}^1(v)=v^1$. Notons
maintenant que puisque $A'$ est métrisable et complète, et $A$ une
$k$-algèbre formellement lisse, l'homomorphisme composé
$A\to A/\mathfrak m^2\xrightarrow{v}A'/\mathfrak m'^2$ se factorise en
$A\xrightarrow{u}A'\to A'/\mathfrak m'^2$, où $u$ est

\oldpage[0\textsubscript{IV-1}]{188}

un $k$-homomorphisme (19.3.11), et l'on a donc déjà
$\operatorname{gr}^0(u)=v^0$ et $\operatorname{gr}^1(u)=v^1$. Or, on a le
diagramme commutatif
\[
\xymatrix@C=5em@R=3em{
  \mathbf S_K^*(\mathfrak m/\mathfrak m^2) \ar[r] \ar[d]_{w} &
  \operatorname{gr}_{\mathfrak m}^*(A) \ar[d]^{\operatorname{gr}^*(u)} \\
  \mathbf S_K^*(\mathfrak m'/\mathfrak m'^2) \ar[r] &
  \operatorname{gr}_{\mathfrak m'}^*(A')
}
\]
où les flèches horizontales sont les homomorphismes canoniques, et $w$
provient de $v^0$ et $v^1$~; comme $v^0$ et $v^1$ sont bijectifs, il en est
de même de $w$, et puisque $A$ et $A'$ sont des $k$-algèbres formellement
lisses, on déduit de (22.2.2, (i)) que $\operatorname{gr}(u)$ est bijectif~;
mais comme $A$ et $A'$ sont séparés et complets, cela entraîne que $u$
lui-même est \emph{bijectif} (Bourbaki, \emph{Alg. comm.}, chap.~III, §~2,
n\textsuperscript{o}~8, cor.~3 du th.~1).

(ii) Reste à montrer que l'image de (22.2.4.1) est l'ensemble des
homomorphismes \emph{injectifs} de $\Upsilon_{K/k}$ dans $V$~; on sait déjà
par (22.2.2, (ii)) qu'elle est contenue dans cet ensemble~; il suffira alors
de prouver le lemme suivant.
\end{proof}

\begin{lemma}[22.2.5.2]
\label{0.22.2.5.2-fr}
Pour tout $K$-homomorphisme $h:\Upsilon_{K/k}\to V$, il existe une
$k$-algèbre $A$ qui est un anneau local noethérien régulier et complet,
d'idéal maximal $\mathfrak m$, de corps résiduel $K$, et un $K$-isomorphisme
$\beta:\mathfrak m/\mathfrak m^2\xrightarrow{\sim}V$ tels que
$h=\beta\circ\chi_{A/k}$.
\end{lemma}

\begin{proof}
En effet, si ce lemme est prouvé, le fait que $A$ est régulier entraîne que
la condition (i) de (22.2.2) est satisfaite (17.1.1)~; si en outre $h$ est
injectif, (22.2.2) montrera que $A$ est une $k$-algèbre formellement lisse
dont la classe aura $h$ pour image (par contre si $h$ n'est pas injectif, la
$k$-algèbre $A$ dont (22.2.5.2) prouve l'existence n'est pas formellement
lisse).

Pour prouver (22.2.5.2), soient $B$ la $K$-algèbre $\mathbf S_K^*(V)$,
$\mathfrak n$ l'idéal d'augmentation de $B$ et $A=\widehat B$ la
$K$-algèbre complétée de $B$ pour la topologie $\mathfrak n$-préadique (de
sorte que $A$ est isomorphe à l'algèbre de séries formelles
$K[[T_1,\ldots,T_n]]$ si $n=\operatorname{rg}_KV$). En vertu de (20.6.11)
il existe sur $A/\mathfrak m^2$ (où $\mathfrak m=\mathfrak nA$) une structure
de $k$-extension de $K$ par $V$ (\emph{distincte} en général de la structure
$k$-triviale définie par l'injection canonique $k\to K\to A$) telle que $h$
soit l'homomorphisme caractéristique de cette extension. Si
$f:k\to A/\mathfrak m^2$ est l'homomorphisme structural, le fait que $k$
soit séparable sur le corps premier permet (19.6.1) de factoriser $f$ en
$k\xrightarrow{g}A\to A/\mathfrak m^2$, et la $k$-algèbre $A$ ainsi définie
répond à la question (20.6.22.2).
\end{proof}

\begin{theorem}[22.2.6]
\label{0.22.2.6-fr}
Soient $k$ un corps, $K$ une extension de $k$. Pour qu'il existe une
$k$-algèbre locale noethérienne $A$, d'idéal maximal $\mathfrak m$, telle que
$A/\mathfrak m=K$ et que $A$ soit formellement lisse (pour sa topologie
$\mathfrak m$-préadique), il faut et il suffit que $\Upsilon_{K/k}$ soit un
$K$-espace vectoriel de rang fini. La $k$-algèbre $A$ est alors déterminée
(à un $k$-isomorphisme près donnant par passage aux quotients l'identité
$K\to K$) par sa dimension, qui est soumise à la seule condition d'être
$\geqslant\operatorname{rg}_K\Upsilon_{K/k}$.
\end{theorem}

\begin{proof}
Comme $\mathfrak m/\mathfrak m^2$ est de rang fini sur $K$, la nécessité de
la condition résulte de (22.2.2, (ii))~; inversement, (22.2.5) montre que la
condition est suffisante et que l'on peut prendre pour
$\mathfrak m/\mathfrak m^2$ un $K$-espace vectoriel de rang quelconque
$\geqslant\operatorname{rg}_K(\Upsilon_{K/k})$.
\end{proof}

\oldpage[0\textsubscript{IV-1}]{189}

En particulier, l'homomorphisme identique de $\Upsilon_{K/k}$ associe
\emph{canoniquement} à $K$ une $k$-algèbre noethérienne complète,
formellement lisse, pour laquelle $\Upsilon_{K/k}$ est isomorphe à
$\mathfrak m/\mathfrak m^2$.

\begin{remarks}[22.2.7]
\label{0.22.2.7-fr}
\begin{enumerate}
  \item[(i)] Soit $\Lambda$ un anneau local noethérien, de corps résiduel
  $k$~; il résulte de (19.7.1) et (19.7.2) que la détermination des
  $\Lambda$-algèbres $A$ qui sont des anneaux locaux noethériens complets et
  qui sont formellement lisses est \emph{équivalente} au même problème où
  $\Lambda$ est remplacé par $k$, et est donc en principe entièrement
  résolue par (22.2.5), qui ramène la question à la structure des \emph{corps
  résiduels} des $k$-algèbres cherchées, en tant qu'extensions de $k$.

  \item[(ii)] Le fait que $\Upsilon_{K/k}$ soit de rang fini sur $K$
  \emph{n'entraîne pas} que $K$ soit «~de multiplicité radicielle finie~» sur
  $k$ (cf. (19.6.6)). Par exemple, soient $k_0$ un corps parfait de
  caractéristique $p>0$, $k=k_0(T)$ le corps des fractions rationnelles à une
  indéterminée sur $k_0$ et $K=k^{p^{-\infty}}$ la plus petite extension
  parfaite de $k$. On a alors $\Omega^1_K=\Omega^1_{K/k}=0$ (21.4.4), donc
  par définition (20.6.1.1)
  $\Upsilon_{K/k}=\Omega^1_k\otimes_kK$, et comme $k^p=k_0(T^p)$, $T$ est
  une $p$-base de $k$ sur $k^p$, donc (21.4.5), $\Omega^1_k$ est de rang $1$
  sur $k$, et par suite $\Upsilon_{K/k}$ de rang $1$ sur $K$. Mais pour tout
  $n$ le corps résiduel de l'anneau local $K\otimes_kk^{p^{-n}}$ est
  isomorphe à $K$ donc n'est pas séparable sur $k^{p^{-n}}$.
\end{enumerate}
\end{remarks}

Le théorème (22.2.2) se généralise comme suit~:

\begin{proposition}[22.2.8]
\label{0.22.2.8-fr}
Sous les conditions de (22.2.1), supposons que la caractéristique $p$ de $k$
soit $>0$~; soit $(k_\alpha)$ une famille filtrante décroissante de
sous-corps de $k$, telle que
$\bigcap_\alpha k_\alpha(k^p)=k^p$. Alors la condition (ii) de (22.2.2) peut
se remplacer par l'une quelconque des suivantes~:
\begin{enumerate}
  \item[a)] Il existe $\alpha_0$ tel que l'homomorphisme canonique
  \[
    \Omega^1_{k/k_\alpha}\otimes_kK
    \longrightarrow\Omega^1_{A/k_\alpha}\otimes_AK
  \]
  soit injectif pour tout $\alpha\geqslant\alpha_0$.
  \item[b)] Il existe $\alpha_0$ tel que $A$ soit une $k$-algèbre
  formellement lisse (pour la topologie $\mathfrak m$-préadique) relativement
  à $k_\alpha$, pour tout $\alpha\geqslant\alpha_0$.
  \item[c)] Pour tout $\alpha$ et tout entier $n\geqslant0$,
  l'homomorphisme canonique
  \[
    \Omega^1_{k/k_\alpha}\otimes_k(A/\mathfrak m^{n+1})
    \longrightarrow
    \Omega^1_{A/k_\alpha}\otimes_A(A/\mathfrak m^{n+1})
  \]
  est inversible à gauche.
\end{enumerate}
\end{proposition}

\begin{proof}
En effet, si $A$ est une $k$-algèbre formellement lisse, elle l'est \emph{a
fortiori} relativement à tout sous-corps de $k$, ce qui entraîne $c)$ en
vertu de (20.7.2) et (19.1.7)~; il est clair que $c)$ implique $b)$, en vertu
de (20.7.2), et que $b)$ implique $a)$. Enfin, compte tenu de (22.2.3), il
reste à prouver que $a)$ entraîne que l'homomorphisme canonique
\[
  u:\Omega^1_k\otimes_kK\longrightarrow\Omega^1_A\otimes_AK
\]
est injectif. Or, on a pour tout $\alpha$ un diagramme commutatif
\[
\xymatrix@C=5em@R=3em{
  \Omega^1_k\otimes_kK \ar[r]^{u} \ar[d]_{v_\alpha} &
  \Omega^1_A\otimes_AK \ar[d] \\
  \Omega^1_{k/k_\alpha}\otimes_kK \ar[r]_{u_\alpha} &
  \Omega^1_{A/k_\alpha}\otimes_AK
}
\]

\oldpage[0\textsubscript{IV-1}]{190}

et l'hypothèse que $u_\alpha$ est injectif pour
$\alpha\geqslant\alpha_0$ implique que $\operatorname{Ker}(u)$ est contenu
dans l'\emph{intersection} des $\operatorname{Ker}(v_\alpha)$. Mais on a vu
(21.8.3) que cette intersection est nulle en vertu de l'hypothèse
$\bigcap_\alpha k_\alpha(k^p)=k^p$, ce qui achève la démonstration.
\end{proof}

\begin{proposition}[22.2.9]
\label{0.22.2.9-fr}
Sous les hypothèses de (22.2.1), les conditions suivantes sont
équivalentes~:
\begin{enumerate}
  \item[a)] $A/\mathfrak m^2$ est une $k$-extension $k$-triviale de $K$ par
  $\mathfrak m/\mathfrak m^2$.
  \item[b)] $\chi_{A/k}=0$.
  \item[c)] L'homomorphisme canonique (20.5.11.2)
  \[
    \delta_{K/A/k}:\mathfrak m/\mathfrak m^2
    \longrightarrow\Omega^1_{A/k}\otimes_AK
  \]
  est injectif (en d'autres termes, la suite
  \begin{equation}
  \tag{22.2.9.1}
  \label{0.22.2.9.1-fr}
    0\longrightarrow\mathfrak m/\mathfrak m^2
    \longrightarrow\Omega^1_{A/k}\otimes_AK
    \longrightarrow\Omega^1_{K/k}\longrightarrow0
  \end{equation}
  est exacte).
\end{enumerate}
\end{proposition}

\begin{proof}
L'équivalence de $b)$ et $c)$ résulte de la suite exacte (20.6.22.1)~;
l'équivalence de $c)$ et $a)$ résulte de (20.5.12), la suite (22.2.9.1)
étant scindée si elle est exacte, puisqu'il s'agit d'espaces vectoriels.

On notera que si $K$ est \emph{séparable sur $k$}, les conditions
équivalentes de (22.2.9) sont satisfaites, en vertu de la définition de
$\chi_{A/k}$ (22.2.1) et de (20.6.3).
\end{proof}

\begin{corollary}[22.2.10]
\label{0.22.2.10-fr}
Sous les hypothèses de (22.2.1), considérons un sous-corps $k_0$ de $k$. Les
conditions suivantes sont équivalentes~:
\begin{enumerate}
  \item[a)] $A/\mathfrak m^2$ est une $k_0$-extension $k_0$-triviale de $K$
  par $\mathfrak m/\mathfrak m^2$.
  \item[b)] L'homomorphisme composé
  \[
    \Upsilon_{K/k_0}\longrightarrow\Upsilon_{K/k}
    \xrightarrow{\chi_{A/k}}\mathfrak m/\mathfrak m^2
  \]
  est nul.
  \item[c)] L'homomorphisme caractéristique
  $\chi_{A/k}:\Upsilon_{K/k}\to\mathfrak m/\mathfrak m^2$ se factorise en
  \[
    \Upsilon_{K/k}\longrightarrow\Upsilon_{K/k/k_0}
    \xrightarrow{\chi'}\mathfrak m/\mathfrak m^2.
  \]
\end{enumerate}
En outre, lorsque ces conditions sont remplies, l'homomorphisme $\chi'$ est
déterminé de façon unique et coïncide avec l'homomorphisme caractéristique
de la $k_0$-extension $k_0$-triviale $A/\mathfrak m^2$ de $K$ par
$\mathfrak m/\mathfrak m^2$.
\end{corollary}

\begin{proof}
Comme le composé dans $b)$ n'est autre que $\chi_{A/k_0}$, l'équivalence de
$a)$ et $b)$ résulte de (22.2.9)~; d'autre part, $b)$ et $c)$ sont
équivalentes en vertu de la suite exacte (20.6.18.1)
\[
  \Upsilon_{K/k_0}\longrightarrow\Upsilon_{K/k}
  \longrightarrow\Upsilon_{K/k/k_0}\longrightarrow0,
\]
et comme l'homomorphisme canonique
$\Upsilon_{K/k}\to\Upsilon_{K/k/k_0}$ est surjectif, cela implique l'unicité
de $\chi'$~; le fait que $\chi'$ est l'homomorphisme caractéristique de la
$k_0$-extension $A/\mathfrak m^2$ résulte de (20.6.22.2).
\end{proof}

\begin{corollary}[22.2.11]
\label{0.22.2.11-fr}
Sous les hypothèses de (22.2.1), soit $p$ l'exposant caractéristique de $k$,
et soit $(k_\alpha)$ une famille filtrante décroissante de sous-corps de $k$
telle que $\bigcap_\alpha k_\alpha(k^p)=k^p$. Si $\Upsilon_{K/k}$ est de rang
fini sur $K$, il existe un indice $\alpha$ tel que $A/\mathfrak m^2$ soit une
$k_\alpha$-extension $k_\alpha$-triviale de $K$ par
$\mathfrak m/\mathfrak m^2$.
\end{corollary}

\begin{proof}
On sait en effet (21.8.3) qu'il existe alors un indice $\alpha$ tel que
l'homomorphisme canonique
$\Upsilon_{K/k}\to\Upsilon_{K/k/k_\alpha}$ soit bijectif, donc la condition
$b)$ de (22.2.10) (où $k_\alpha$ remplace $k_0$) est satisfaite.
\end{proof}
\oldpage[0\textsubscript{IV-1}]{191}

\subsection{Application aux relations entre certains anneaux locaux et leurs complétés}
\label{subsection:0.22.3-fr}

\begin{lemma}[22.3.1]
\label{0.22.3.1-fr}
Soient $A_0$ un anneau semi-local noethérien, $A$ une $A_0$-algèbre finie
(qui est donc un anneau semi-local noethérien, cf. Bourbaki,
\emph{Alg. comm.}, chap.~IV, §~2, n\textsuperscript{o}~5, cor.~3 de la
prop.~9). Si $\widehat A_0$ et $\widehat A$ sont les complétés de $A_0$ et
$A$ pour leurs topologies préadiques respectives, alors, lorsqu'on munit
$A_0$, $A$ et $\widehat A$ des topologies discrètes, $\widehat A$ est une
$A$-algèbre formellement lisse relativement à $A_0$ (19.9.1).
\end{lemma}

\begin{proof}
On sait en effet que
$\widehat A=A\otimes_{A_0}\widehat A_0$ (Bourbaki,
\emph{Alg. comm.}, chap.~IV, §~2, n\textsuperscript{o}~5, cor.~3 de la
prop.~9 et chap.~III, §~3, n\textsuperscript{o}~4, th.~1), donc la proposition
résulte de (19.9.3).
\end{proof}

\begin{proposition}[22.3.2]
\label{0.22.3.2-fr}
Soient $A$ un anneau local noethérien régulier, $K$ son corps des fractions,
$p$ la caractéristique de $K$. On suppose vérifiée l'une des hypothèses
suivantes~:
\begin{enumerate}
  \item[(i)] $p=0$.
  \item[(ii)] $p>0$ et il existe une famille filtrante décroissante
  $(A_\alpha)$ de sous-anneaux noethériens de $A$, tels que $A$ soit une
  $A_\alpha$-algèbre finie pour tout $\alpha$, que
  $A^{p^{h(\alpha)}}\subset A_\alpha$ pour un entier $h(\alpha)$ convenable,
  et que, si $K_\alpha$ désigne le corps des fractions de $A_\alpha$, on ait
  \[
    \bigcap_\alpha K_\alpha(K^p)=K^p.
  \]
\end{enumerate}
Soient alors $B$ une $A$-algèbre finie intègre et contenant $A$,
$\mathfrak n$ un idéal premier de $B$, $C$ l'anneau local $B_\mathfrak n$,
$\mathfrak q$ un idéal premier du complété $\widehat C$ tel que
$\mathfrak q\cap C=0$, de sorte que l'anneau local
$\widehat C_\mathfrak q$ est une algèbre sur le corps des fractions $L$ de
$B$. Alors $\widehat C_\mathfrak q$ est une $L$-algèbre formellement lisse
pour sa topologie $\mathfrak q$-adique, et par suite un anneau
géométriquement régulier sur $L$ (19.6.5).
\end{proposition}

\begin{proof}
Distinguons deux cas.

I) Supposons que $\mathfrak n$ contienne l'idéal maximal $\mathfrak m$ de
$A$, et par suite $\mathfrak n\cap A=\mathfrak m$. Alors $\mathfrak n$ est
maximal (Bourbaki, \emph{Alg. comm.}, chap.~V, §~2,
n\textsuperscript{o}~1, prop.~1) et si $\mathfrak r$ est le radical de
l'anneau semi-local $B$, $\widehat C$ est le séparé complété de $B$ pour la
topologie $\mathfrak n$-préadique, et un des composants de l'anneau
semi-local $\widehat B$, complété de $B$ pour la topologie
$\mathfrak r$-préadique (Bourbaki, \emph{Alg. comm.}, chap.~III, §~3,
n\textsuperscript{o}~4, prop.~8)~; on peut donc écrire
$\widehat C_\mathfrak q=\widehat B_{\mathfrak q'}$, où $\mathfrak q'$ est
l'idéal premier de $\widehat B$ image réciproque de $\mathfrak q$. Notons
maintenant que $B$ est un sous-anneau de $\widehat B$ et l'hypothèse
$\mathfrak q\cap C=0$ entraîne $\mathfrak q'\cap B=0$, donc
$\widehat B_{\mathfrak q'}$ est aussi localisé de l'anneau
$\widehat B\otimes_B L$ en un de ses idéaux premiers $\mathfrak q''$, dont
$\mathfrak q'$ est l'image réciproque dans $\widehat B$. En outre on a
$\widehat B=\widehat A\otimes_A B$ (Bourbaki, \emph{Alg. comm.}, chap.~IV,
§~2, n\textsuperscript{o}~5, cor. de la prop.~9 et chap.~III, §~3,
n\textsuperscript{o}~4, th.~1), donc
$\widehat B\otimes_B L=\widehat A\otimes_A L$~; enfin, si $\mathfrak p$ est
l'idéal premier de $\widehat A$, image réciproque de $\mathfrak q''$, le
fait que $A\to B$ est injectif et la commutativité du diagramme
\[
  \xymatrix{
    B \ar[r] & \widehat B\\
    A \ar[u] \ar[r] & \widehat A \ar[u]
  }
\]
entraînent que $\mathfrak p\cap A=0$, donc
$(\widehat A\otimes_A L)_{\mathfrak q''}$ est aussi localisé de l'anneau
\[
  \widehat A_\mathfrak p\otimes_A L
  =(\widehat A_\mathfrak p\otimes_A K)\otimes_K L
\]
en un de ses idéaux premiers.

\oldpage[0\textsubscript{IV-1}]{192}

Pour montrer que cet anneau local est une $L$-algèbre formellement lisse,
il suffit donc de prouver que $\widehat A_\mathfrak p\otimes_A K$ est une
$K$-algèbre formellement lisse (19.3.5, (iii) et (iv))~; d'ailleurs
$\widehat A_\mathfrak p\otimes_A K=\widehat A_\mathfrak p$ puisque
$\mathfrak p\cap A=0$. On va appliquer les critères (22.2.2) et (22.2.8).
En premier lieu, comme $A$ est régulier, il en est de même de $\widehat A$
(17.1.5), donc de $\widehat A_\mathfrak p$ (17.3.2), et la condition (i) de
(22.2.2) est vérifiée (17.1.1). Il suffit par suite (lorsque $p>0$) de
vérifier la condition b) de (22.2.8). Or, (22.3.1) montre que, pour tout
$\alpha$, $\widehat A$ est une $A$-algèbre formellement lisse relativement
à $A_\alpha$ (pour les topologies discrètes)~; par suite
\[
  \widehat A\otimes_{A_\alpha}K_\alpha
  =\widehat A\otimes_A(A\otimes_{A_\alpha}K_\alpha)
\]
est une $(A\otimes_{A_\alpha}K_\alpha)$-algèbre formellement lisse
relativement à $K_\alpha$, pour les topologies discrètes (19.9.2, (iii)).
Mais comme $A$ est une $A_\alpha$-algèbre finie, on a
$A\otimes_{A_\alpha}K_\alpha=K$~; comme $\widehat A_\mathfrak p$ est un
anneau de fractions de $\widehat A\otimes_A K$, on en conclut (19.9.2,
(iv)) que $\widehat A_\mathfrak p$ est une $K$-algèbre formellement lisse
relativement à $K_\alpha$, pour la topologie discrète, et \emph{a fortiori}
pour sa topologie $\mathfrak p$-adique ((19.3.8) et (19.9.5)). Mais cela est
précisément la condition b) de (22.2.8), donc la proposition est démontrée
dans ce cas pour $p>0$. Lorsque $p=0$, le corps résiduel de
$\widehat A_\mathfrak p$ est séparable sur $K$, et comme
$\widehat A_\mathfrak p$ est un anneau régulier, c'est une $K$-algèbre
formellement lisse en vertu de (19.6.4).

II) Cas général. Soit $\mathfrak s$ l'idéal premier $\mathfrak n\cap A$ de
$A$, et posons $S=A-\mathfrak s$, de sorte que $S^{-1}A=A_\mathfrak s$, et
l'on a $C=(S^{-1}B)_{S^{-1}\mathfrak n}$, où cette fois
$S^{-1}\mathfrak n$ contient l'idéal maximal
$\mathfrak sA_\mathfrak s$ de $A_\mathfrak s$. Comme $A_\mathfrak s$ est
régulier (17.3.2) et $S^{-1}B$ une $A_\mathfrak s$-algèbre finie intègre
contenant $A_\mathfrak s$, tout revient à vérifier la condition (ii) pour
$A_\mathfrak s$ lorsque $p>0$~: or, posons
$\mathfrak s_\alpha=\mathfrak s\cap A_\alpha$ et considérons l'anneau local
$(A_\alpha)_{\mathfrak s_\alpha}$. Pour tout $t\notin\mathfrak s$, on a par
hypothèse, en posant $q=p^{h(\alpha)}$, $t^q\in A_\alpha$ et
$t^q\notin\mathfrak s$, donc $t^q\notin\mathfrak s_\alpha$~; si l'on pose
$S_\alpha=A_\alpha-\mathfrak s_\alpha$, on voit donc que l'on a
$A_\mathfrak s=S_\alpha^{-1}A$, et l'hypothèse entraîne que
$A_\mathfrak s$ est une $(A_\alpha)_{\mathfrak s_\alpha}$-algèbre finie~;
en outre, $K_\alpha$ est aussi le corps des fractions de
$(A_\alpha)_{\mathfrak s_\alpha}$, ce qui achève la démonstration.
\end{proof}

\begin{theorem}[22.3.3]
\label{0.22.3.3-fr}
Soient $A$ un anneau local noethérien complet, $\mathfrak p$ un idéal
premier de $A$, $B$ l'anneau localisé $A_\mathfrak p$, $\mathfrak q$ un
idéal premier du complété $\widehat B$, $\mathfrak r=\mathfrak q\cap B$.
Alors l'anneau localisé $\widehat B_\mathfrak q$ de $\widehat B$ est une
$B_\mathfrak r$-algèbre formellement lisse pour les topologies préadiques.
\end{theorem}

\begin{proof}
L'idéal premier $\mathfrak r$ de $B$ est de la forme
$\mathfrak n_\mathfrak p$, où $\mathfrak n\subset\mathfrak p$ est un idéal
premier de $A$, et l'on a $B_\mathfrak r=A_\mathfrak n$~; soit $k$ le corps
résiduel de $B_\mathfrak r$, qui est aussi le corps des fractions de l'anneau
local intègre et complet $A'=A/\mathfrak n$. Comme $\widehat B$ est un
$B$-module plat $(\mathbf 0_\mathrm I,~7.3.5)$,
$\widehat B_\mathfrak q$ est un $B_\mathfrak r$-module plat
$(\mathbf 0_\mathrm I,~6.3.2)$, et pour prouver que
$\widehat B_\mathfrak q$ est une $B_\mathfrak r$-algèbre formellement
lisse, il suffit de prouver que
\[
  \widehat B_\mathfrak q\otimes_{B_\mathfrak r}k
  =\widehat B_\mathfrak q/\mathfrak r\widehat B_\mathfrak q
\]
est une $k$-algèbre formellement lisse (19.7.1). Or
$\widehat B_\mathfrak q/\mathfrak r\widehat B_\mathfrak q
=(\widehat B/\mathfrak r\widehat B)_{\mathfrak q'}$, où $\mathfrak q'$ est
l'image canonique de $\mathfrak q$ dans
$\widehat B/\mathfrak r\widehat B$~; on a par définition
$\mathfrak q'\cap(B/\mathfrak r)=0$ et
$B/\mathfrak r=A_{\mathfrak p'}$, où
$\mathfrak p'=\mathfrak p/\mathfrak n$~; en outre
$\widehat B/\mathfrak r\widehat B$ est le complété de $B/\mathfrak r$.

On est donc ramené à démontrer le corollaire suivant~:
\end{proof}

\begin{corollary}[22.3.4]
\label{0.22.3.4-fr}
Soient $A$ un anneau local noethérien complet et intègre, $K$ son corps des
fractions, $\mathfrak p$ un idéal premier de $A$, $B$ l'anneau localisé
$A_\mathfrak p$. Alors, pour tout idéal premier $\mathfrak q$

\oldpage[0\textsubscript{IV-1}]{193}

du complété $\widehat B$, tel que $\mathfrak q\cap B=0$, l'anneau local
$\widehat B_\mathfrak q$ est une $K$-algèbre formellement lisse pour sa
topologie $\mathfrak q$-adique, donc un anneau géométriquement régulier sur
$K$.
\end{corollary}

\begin{proof}
En effet, il résulte de (19.8.9) qu'il existe un sous-anneau $A_0$ de $A$
qui est un anneau de séries formelles sur un corps de caractéristique $p>0$
ou sur un anneau de Cohen (on rappelle que les anneaux de Cohen ont un corps
des fractions de caractéristique $0$), et est tel que $A$ soit une
$A_0$-algèbre finie. Or, on sait que l'anneau $A_0$ est régulier (17.3.8) et
vérifie l'une des hypothèses (i), (ii) de (22.3.2), en vertu de (21.8.8)~;
il suffit donc d'appliquer (22.3.2), en remplaçant $B$ par $A$ et $A$ par
$A_0$.
\end{proof}

\subsection{Résultats préliminaires sur les extensions finies d'anneaux locaux dont l'idéal maximal est de carré nul}
\label{subsection:0.22.4-fr}
\label{0.22.4-fr}

\begin{proposition}[22.4.1]
\label{0.22.4.1-fr}
Soient $A$ un anneau local dont l'idéal maximal $\mathfrak m$ est de carré
nul, $K=A/\mathfrak m$ son corps résiduel~; désignons par $V$ l'idéal
$\mathfrak m$ considéré comme $K$-espace vectoriel. Soient $T_i$
$(1\leqslant i\leqslant r)$ des indéterminées, et pour chaque $i$, soit
$F_i$ un polynôme unitaire de $A[T_i]$ de degré $d_i$~; désignons par $A'$ le
quotient de l'anneau de polynômes $A[T_1,\ldots,T_r]$ par l'idéal
$\mathfrak b$ engendré par les $r$ polynômes $F_i$. Supposons que $A'$ soit
un anneau local, et soient $\mathfrak m'$ son idéal maximal,
$K'=A'/\mathfrak m'$ son corps résiduel. Alors~:
\begin{enumerate}
  \item[(i)] Si l'on pose $d=\prod_{i=1}^r d_i$, $A'$ est un $A$-module
  libre de rang $d$, et l'on a $[K':K]\leqslant d$.
  \item[(ii)] Pour que $[K':K]=d$, il faut et il suffit que
  $A'\otimes_AK=A'/\mathfrak mA'$ soit un corps (nécessairement isomorphe à
  $K'$), autrement dit, que $\mathfrak m'=\mathfrak mA'$. On a dans ce cas
  $\mathfrak m'^2=0$, et si $V'$ désigne l'idéal $\mathfrak m'$ considéré
  comme $K'$-espace vectoriel, l'homomorphisme canonique
  $V\otimes_KK'\to V'$ est bijectif (et par suite
  $\operatorname{rg}_{K'}V'=\operatorname{rg}_K V$).
\end{enumerate}
\end{proposition}

\begin{proof}
Sans supposer $A'$ local, il est évident, par division euclidienne et
récurrence sur $r$, que si $t_i$ est la classe de $T_i$ mod.~$\mathfrak b$,
les monômes
$M_\nu=t_1^{\nu_1}t_2^{\nu_2}\cdots t_r^{\nu_r}$, où
$0\leqslant\nu_i\leqslant d_i-1$, forment une base du $A$-module $A'$. Si
$A'$ est supposé local, $K'$ est un quotient de
$A'/\mathfrak mA'=A'\otimes_AK$, donc
$[K':K]\leqslant[A'\otimes_AK:K]=d$, et il ne peut y avoir égalité que si
$\mathfrak m'=\mathfrak mA'$, ce qui démontre (i) et la première assertion
de (ii). D'autre part, il est clair que lorsque
$\mathfrak m'=\mathfrak mA'$, on a $\mathfrak m'^2=0$, et
$V\otimes_KK'=\mathfrak m\otimes_AA'=\mathfrak mA'=V'$ puisque $A'$ est
un $A$-module libre.
\end{proof}

\begin{remark}[22.4.2]
\label{0.22.4.2-fr}
Lorsque $A$ est artinien (autrement dit $\operatorname{rg}_K(V)$ fini),
l'hypothèse que $A'$ est local entraîne toujours que
$\operatorname{rg}_{K'}(\mathfrak m'/\mathfrak m'^2)
\geqslant\operatorname{rg}_K(V)$, comme nous le verrons plus loin (22.5.2).
On notera que ces deux rangs peuvent être égaux sans que l'on ait
$\mathfrak m'=\mathfrak mA'$.
\end{remark}

\begin{proposition}[22.4.3]
\label{0.22.4.3-fr}
Les hypothèses sur $A$ et les notations étant celles de (22.4.1), supposons
que les $F_i$ soient de la forme
\begin{equation}
\label{0.22.4.3.1-fr}
  F_i(T_i)=T_i^{d_i}+\sum_{1\leqslant k\leqslant d_i}
  c_{ik}T_i^{d_i-k},
  \tag{22.4.3.1}
\end{equation}
où $d_i\geqslant2$ pour tout $i$, et $c_{ik}\in\mathfrak m$
(«~polynômes d'Eisenstein~»). Alors~:
\begin{enumerate}
  \item[(i)] $A'$ est un anneau local, et son corps résiduel
  $K'=A'/\mathfrak m'$ est isomorphe à $K$.

  \oldpage[0\textsubscript{IV-1}]{194}

  \item[(ii)] Le noyau de l'homomorphisme canonique
  $V\to V'=\mathfrak m'/\mathfrak m'^2$ est l'idéal $\mathfrak n$ de $A$
  engendré par les $\xi_i=c_{i,d_i}$ («~termes constants~» des $F_i$), et
  le conoyau de cet homomorphisme a pour base sur $K'$ les images des $T_i$~;
  en particulier, on a
  \begin{equation}
  \label{0.22.4.3.2-fr}
    \operatorname{rg}_{K'}(V')=\operatorname{rg}_K(V)+r-r',
    \tag{22.4.3.2}
  \end{equation}
  où $r'=\operatorname{rg}_K(\mathfrak n)$. Lorsque
  $\operatorname{rg}_K(V)$ est fini (autrement dit lorsque $A$ est
  artinien), pour que
  $\operatorname{rg}_{K'}(V')=\operatorname{rg}_K(V)$, il faut et il suffit
  que les $\xi_i$ soient linéairement indépendants sur $K$.
\end{enumerate}
\end{proposition}

\begin{proof}
(i) Il est clair que $A'\otimes_AK=A'/\mathfrak mA'$ est isomorphe à
$K[T_1,\ldots,T_r]/\mathfrak b'$, où $\mathfrak b'$ est l'idéal engendré
par les $r$ polynômes $T_i^{d_i}$~; c'est donc un anneau local de corps
résiduel $K$, dont l'idéal maximal est engendré par les classes des $T_i$
mod.~$\mathfrak b'$, d'où (i).

(ii) Soit $t_i$ $(1\leqslant i\leqslant r)$ la classe de $T_i$
mod.~$\mathfrak b$~; il résulte de (i) que l'idéal maximal $\mathfrak m'$ de
$A'$ est donné par
\[
  \mathfrak m'=\mathfrak mA'+\sum_{i=1}^r t_iA',
\]
d'où, compte tenu de $\mathfrak m^2=0$,
\[
  \mathfrak m'^2=\sum_{i,j}t_it_jA'+\sum_i t_i\mathfrak mA'.
\]
On en conclut que $A'/\mathfrak m'^2$ est isomorphe à l'anneau
$A[T_1,\ldots,T_r]/\mathfrak c$, où $\mathfrak c$ est l'idéal engendré par
les éléments
\[
  F_i\ (1\leqslant i\leqslant r),\qquad
  T_iT_j\ (1\leqslant i\leqslant r,\ 1\leqslant j\leqslant r),\qquad
  xT_i\ (1\leqslant i\leqslant r,\ x\in\mathfrak m).
\]
Comme $d_i\geqslant2$ et que les $T_i^2$ appartiennent à
$\mathfrak c$, les hypothèses faites entraînent que $\mathfrak c$ est aussi
engendré par les éléments
\[
  \xi_i\ (1\leqslant i\leqslant r),\qquad
  T_iT_j\ (1\leqslant i\leqslant r,\ 1\leqslant j\leqslant r),\qquad
  T_ix\ (1\leqslant i\leqslant r,\ x\in\mathfrak m).
\]
Soient
\[
  \mathfrak c_1=\sum_{i,j}T_iT_jA[T_1,\ldots,T_r],\qquad
  \mathfrak c_2=\mathfrak c_1+\sum_i\mathfrak mT_iA[T_1,\ldots,T_r],
\]
de sorte que
$\mathfrak c=\mathfrak c_2+\sum_i\xi_iA[T_1,\ldots,T_r]$. Il est clair que
$A[T_1,\ldots,T_r]/\mathfrak c_1$ est le $A$-module somme directe de $A$ et
des $At_i^{(1)}$, où $t_i^{(1)}$ est la classe de $T_i$
mod.~$\mathfrak c_1$. On en déduit que
$A[T_1,\ldots,T_r]/\mathfrak c_2$ est somme directe de $A$ et des
$Kt_i^{(2)}$, où $t_i^{(2)}$ est la classe de $T_i$
mod.~$\mathfrak c_2$. Enfin $A[T_1,\ldots,T_r]/\mathfrak c$ est somme
directe de $A/\mathfrak n$ et des $Kt_i'$, où $t_i'$ est la classe de $T_i$
mod.~$\mathfrak c$, et $\mathfrak n$ est l'idéal de $A$ engendré par les
$\xi_i$~; dans l'identification de $A'/\mathfrak m'^2$ à
$A[T_1,\ldots,T_r]/\mathfrak c$, l'image de $V$ s'identifie à
$\mathfrak m/\mathfrak n$, et ce qui précède prouve donc (ii).
\end{proof}

\begin{proposition}[22.4.4]
\label{0.22.4.4-fr}
Les hypothèses sur $A$ et les notations étant celles de (22.4.1), supposons
que $K$ soit de caractéristique $p>0$ (donc $p\cdot1\in\mathfrak m$ dans
$A$) et que les $F_i$ soient de la forme
\begin{equation}
\label{0.22.4.4.1-fr}
  F_i(T_i)=T_i^p+p\sum_{k=1}^{p-1}c_{ik}T_i^{p-k}-f_i
  \tag{22.4.4.1}
\end{equation}
avec $c_{ik}\in A$ et $f_i\in A$ pour tout $i$ et tout $k$ tel que
$1\leqslant k\leqslant p-1$~; en outre si $a_i$ est la classe de $f_i$
mod.~$\mathfrak m$, on suppose que la famille $(a_i)_{1\leqslant i\leqslant r}$
est $p$-libre sur $K^p$ («~absolument $p$-libre~»). Alors~:
\begin{enumerate}
  \item[(i)] $A'$ est un anneau local, d'idéal maximal $\mathfrak mA'$,
  dont le corps résiduel $K'$ est isomorphe à l'extension de $K$ obtenue par
  adjonction des $a_i^{1/p}$ $(1\leqslant i\leqslant r)$.

  \oldpage[0\textsubscript{IV-1}]{195}

  \item[(ii)] Les éléments $d_Af_i\otimes1_{K'}$ forment une base du noyau
  $N=\Upsilon^{K'}_{A'/A}$ de l'homomorphisme canonique
  \[
    \Omega^1_A\otimes_AK'\longrightarrow
    \Omega^1_{A'}\otimes_{A'}K'.
  \]
\end{enumerate}
\end{proposition}

\begin{proof}
(i) Cette fois $A'\otimes_AK=A'/\mathfrak mA'$ est isomorphe à
$K[T_1,\ldots,T_r]/\mathfrak b'$, où $\mathfrak b'$ est l'idéal engendré
par les polynômes $T_i^p-a_i$~; l'hypothèse faite sur les $a_i$ entraîne que
cet anneau quotient est un corps (21.1.8), d'où les assertions de (i).

Avant de démontrer (ii), nous établirons le

\begin{lemma}[22.4.4.2]
\label{0.22.4.4.2-fr}
Soient $\Lambda$ un anneau, $A$ une $\Lambda$-algèbre, $\mathfrak J$ un
idéal de $A$ tel que l'anneau $C=A/\mathfrak J$ soit de caractéristique
$p>0$ (ce qui équivaut à dire que $p\cdot1\in\mathfrak J$ dans $A$). Soit
$\pi$ l'image canonique de $p\cdot1$ dans
$\mathfrak J/\mathfrak J^2$. Si $C$ est une $(\Lambda/p\Lambda)$-algèbre
formellement lisse pour les topologies discrètes, on a une suite exacte
canonique
\begin{equation}
\label{0.22.4.4.3-fr}
  0\longrightarrow(\mathfrak J/\mathfrak J^2)/C\cdot\pi
  \longrightarrow\Omega^1_{A/\Lambda}\otimes_AC
  \longrightarrow\Omega^1_{C/\Lambda}\longrightarrow0.
  \tag{22.4.4.3}
\end{equation}
\end{lemma}

\begin{proof}[Démonstration du lemme]
Posons en effet $A'=A/pA$, $\mathfrak J'=\mathfrak J/pA$, de sorte que
$C=A'/\mathfrak J'$. Comme
$A'=A\otimes_\Lambda(\Lambda/p\Lambda)$, on a
$\Omega^1_{A/\Lambda}\otimes_AA'=\Omega^1_{A'/\Lambda'}$ à un
isomorphisme canonique près, $\Lambda'$ désignant l'anneau
$\Lambda/p\Lambda$ (20.5.5). Appliquant alors la suite exacte (20.5.14.1) à
la $\Lambda'$-algèbre formellement lisse $C=A'/\mathfrak J'$, on obtient la
suite exacte
\[
  0\longrightarrow\mathfrak J'/\mathfrak J'^2
  \longrightarrow\Omega^1_{A'/\Lambda'}\otimes_{A'}C
  \longrightarrow\Omega^1_{C/\Lambda'}\longrightarrow0
\]
et comme l'homomorphisme $\Lambda\to\Lambda'$ est surjectif, on a
$\Omega^1_{C/\Lambda'}=\Omega^1_{C/\Lambda}$ à un isomorphisme canonique
près (20.5.15)~; d'où la suite exacte (22.4.4.3) en vertu de ce qui précède.
\end{proof}

(ii) Comme on a $\mathfrak m'^2=0$, nous désignerons encore par $V'$ l'idéal
$\mathfrak m'$ considéré comme $K'$-espace vectoriel, et nous poserons
\begin{equation}
\label{0.22.4.4.4-fr}
  V_0=V/K\cdot p=\mathfrak m/(\mathfrak m^2+pA),\qquad
  V_0'=V'/K'\cdot p=\mathfrak m'/(\mathfrak m'^2+pA').
  \tag{22.4.4.4}
\end{equation}
Appliquant la suite exacte (22.4.4.3) au cas où $\Lambda=\mathbf Z$, à la
$\mathbf Z$-algèbre $A$ (resp. $A'$) et à son idéal $\mathfrak m$
(resp. $\mathfrak m'$), il vient, puisque $K$ (resp. $K'$) est séparable sur
le corps premier $\mathbf Z/p\mathbf Z$, donc une
$(\mathbf Z/p\mathbf Z)$-algèbre formellement lisse (19.6.1), les suites
exactes
\begin{equation}
\label{0.22.4.4.5-fr}
  0\longrightarrow V_0\longrightarrow\Omega^1_A\otimes_AK
  \longrightarrow\Omega^1_K\longrightarrow0,\qquad
  0\longrightarrow V_0'\longrightarrow\Omega^1_{A'}\otimes_{A'}K'
  \longrightarrow\Omega^1_{K'}\longrightarrow0.
  \tag{22.4.4.5}
\end{equation}
Considérons alors le diagramme (22.4.4.6) où les deux colonnes médianes sont
la seconde suite (22.4.4.5) et la première tensorisée par $K'$~; comme la
première suite exacte (22.4.4.5) est formée d'espaces vectoriels sur $K$,
les deux colonnes médianes de (22.4.4.6) sont exactes~; en outre, le
diagramme formé par ces colonnes et les flèches horizontales qui les
connectent est commutatif, en vertu de la démonstration de (22.4.4.2) et de
la commutativité du diagramme (20.5.11.3). La dernière ligne de (22.4.4.6)
est la suite exacte obtenue à partir de (20.6.1.1) en y remplaçant $A$, $B$,
$C$ par $\mathbf Z$, $K$ et $K'$~; la ligne médiane est obtenue par
tensorisation avec $K'$ de la suite exacte de $K$-modules déduite de
(20.6.1.1) en y remplaçant $A$, $B$, $C$ par $\mathbf Z$, $A$, $K$
respectivement. Le diagramme formé de ces deux lignes et des flèches
verticales qui les connectent est commutatif en vertu de la commutativité de
(20.6.4.3). Notons d'autre part qu'on est sous les conditions d'application
de (22.4.1, (ii)), donc la ligne supérieure du diagramme (22.4.4.6) est
exacte, les colonnes extrêmes du diagramme

\oldpage[0\textsubscript{IV-1}]{196}

\begin{equation}
\label{0.22.4.4.6-fr}
  \xymatrix{
    && 0 \ar[d] & 0 \ar[d] && \\
    & 0 \ar[r] & V_0\otimes_KK' \ar[r] \ar[d] & V_0' \ar[r] \ar[d] & 0 & \\
    0 \ar[r] & N \ar[r] \ar[d] & \Omega^1_A\otimes_AK' \ar[r] \ar[d] &
      \Omega^1_{A'}\otimes_{A'}K' \ar[r] \ar[d] &
      \Omega^1_{A'/A}\otimes_{A'}K' \ar[r] \ar[d] & 0 \\
    0 \ar[r] & \Upsilon_{K'/K} \ar[r] \ar[d] &
      \Omega^1_K\otimes_KK' \ar[r] \ar[d] & \Omega^1_{K'} \ar[r] \ar[d] &
      \Omega^1_{K'/K} \ar[r] \ar[d] & 0 \\
    & 0 & 0 & 0 & 0 &
  }
  \tag{22.4.4.6}
\end{equation}
sont donc formées respectivement des noyaux et des conoyaux des flèches
horizontales médianes, ce qui achève de prouver que le diagramme est
commutatif. En outre~:

\begin{lemma}[22.4.4.7]
\label{0.22.4.4.7-fr}
Les homomorphismes
\[
  N\longrightarrow\Upsilon_{K'/K}\qquad\text{et}\qquad
  \Omega^1_{A'/A}\otimes_{A'}K'\longrightarrow\Omega^1_{K'/K}
\]
du diagramme (22.4.4.6) sont bijectifs.
\end{lemma}

\begin{proof}
En effet, le diagramme du serpent (Bourbaki, \emph{Alg. comm.}, chap.~I,
§~1, n\textsuperscript{o}~4, prop.~2) appliqué aux deux colonnes médianes de
(22.4.4.6) donne une suite exacte
\[
  0\longrightarrow N\longrightarrow\Upsilon_{K'/K}\longrightarrow0
  \longrightarrow\Omega^1_{A'/A}\otimes_{A'}K'
  \longrightarrow\Omega^1_{K'/K}\longrightarrow0.
\]
Or, si $t_i$ est l'image canonique de $T_i$ dans $A'$, la relation
$F_i(t_i)=0$, jointe au fait que dans $A'$ on a
$p\cdot1\in\mathfrak mA'\subset\mathfrak m'$, montre aussitôt que l'on a
$d_{A'}f_i\in\mathfrak m'\Omega^1_{A'}$, donc
$d_{A'}f_i\otimes1_{K'}=0$~; cela prouve que les
$d_Af_i\otimes1_{K'}$ appartiennent à $N$. En outre, leurs images dans
$\Upsilon_{K'/K}$ sont les $d_Ka_i\otimes1_{K'}$, qui forment une base de
$\Upsilon_{K'/K}$ sur $K'$ (21.4.7)~; compte tenu de (22.4.4.7), cela achève
de démontrer (22.4.4).
\end{proof}
\end{proof}
\begin{env}[22.4.5]
\label{0.22.4.5-fr}
Considérons maintenant deux anneaux $A$, $B$ \emph{de caractéristique} $p>0$, et
deux homomorphismes
\[
  A\xrightarrow{i}B\xrightarrow{j}A
\]

\oldpage[0\textsubscript{IV-1}]{197}

vérifiant les conditions de (21.3.1). Soit en outre $\mathfrak m$ un idéal
\emph{de carré nul} dans $A$~; posons
\[
  K=A/\mathfrak m,\qquad
  B_{(K)}=B\otimes_AK=B/Bi(\mathfrak m)
\]
et soient
\[
  \varphi:A\longrightarrow K,\qquad
  \psi:B\longrightarrow B_{(K)}
\]
les homomorphismes canoniques.

Comme $\varphi$ est surjectif, $\Omega^1_{B_{(K)}/A}$ s'identifie
canoniquement à $\Omega^1_{B_{(K)}/K}$ (20.5.15). Notons que, puisque
$p\geqslant2$, pour tout $x\in\mathfrak m$, on a $j(i(x))=x^p=0$, autrement
dit $j(i(\mathfrak m))=0$~; par passage aux quotients on déduit donc de $j$ un
homomorphisme
\[
  j':B_{(K)}\longrightarrow A
\]
tel que, si $i'=\psi\circ i:A\longrightarrow B_{(K)}$, on ait
\[
  j'(i'(x))=x^p\qquad\text{et}\qquad i'(j'(y))=y^p
  \qquad\text{pour }x\in A\text{ et }y\in B_{(K)}.
\]
En posant, suivant les notations de (21.3.2),
\[
  \Theta_{B_{(K)}/A}
  =\Omega^1_{B_{(K)}/A}\otimes_{B_{(K)}}A_{[j']}
  =\Omega^1_{B_{(K)}/K}\otimes_{B_{(K)}}A_{[j']},
\]
on a donc (21.3.2.2) un homomorphisme canonique
\begin{equation}
\label{0.22.4.5.1-fr}
  \pi_{B_{(K)}/A}:\Theta_{B_{(K)}/A}\longrightarrow\Omega^1_A.
  \tag{22.4.5.1}
\end{equation}

D'autre part, on déduit de $i$ et $j$, par passage aux quotients (tenant compte
de ce que $j(i(\mathfrak m))=0$), des homomorphismes
\[
  K\xrightarrow{i_0}B_{(K)}\xrightarrow{j_0}K
\]
de sorte que l'on a le diagramme commutatif
\[
  \xymatrix{
    A \ar[r]^i \ar[d]_\varphi & B \ar[r]^j \ar[d]^\psi & A \ar[d]^\varphi \\
    K \ar[r]_{i_0} & B_{(K)} \ar[r]_{j_0} \ar[ur]_{j'} & K
  }
\]
Il est clair en outre que l'on a
\[
  j_0(i_0(s))=s^p\qquad\text{et}\qquad i_0(j_0(t))=t^p
  \qquad\text{pour }s\in K\text{ et }t\in B_{(K)}.
\]
On définit donc de même
\[
  \Theta_{B_{(K)}/K}
  =\Omega^1_{B_{(K)}/K}\otimes_{B_{(K)}}K_{[j_0]}
  =\Theta_{B_{(K)}/A}\otimes_AK,
\]
donc, en tensorisant (22.4.5.1) par $K$, on obtient un
$K$-\emph{homomorphisme canonique}
\begin{equation}
\label{0.22.4.5.2-fr}
  \pi'_{B_{(K)}/K}=\pi_{B_{(K)}/A}\otimes1_K:
  \Theta_{B_{(K)}/K}\longrightarrow\Omega^1_A\otimes_AK.
  \tag{22.4.5.2}
\end{equation}

\oldpage[0\textsubscript{IV-1}]{198}

Il résulte aussitôt des définitions que le diagramme
\[
  \xymatrix{
    \Theta_{B_{(K)}/K} \ar[r]^{\pi'_{B_{(K)}/K}} \ar[dr]_{\pi_{B_{(K)}/K}}
      & \Omega^1_A\otimes_AK \ar[d]^{\varphi_{K/A}} \\
    & \Omega^1_K
  }
\]
est commutatif, $\pi_{B_{(K)}/K}$ étant l'homomorphisme canonique
correspondant à $i_0$ et $j_0$ (21.3.2.2). Par \emph{restriction} à
$\Xi_{B_{(K)}/K}=\operatorname{Ker}(\pi_{B_{(K)}/K})$, l'homomorphisme
$\pi'_{B_{(K)}/K}$ définit donc un \emph{homomorphisme canonique}
\begin{equation}
\label{0.22.4.5.3-fr}
  \Xi_{B_{(K)}/K}\longrightarrow
  \operatorname{Ker}(\Omega^1_A\otimes_AK\longrightarrow\Omega^1_K)=\Upsilon_{K/A}.
  \tag{22.4.5.3}
\end{equation}
\end{env}

\begin{env}[22.4.6]
\label{0.22.4.6-fr}
Les hypothèses et notations étant celles de (22.4.5), supposons en outre que
$\mathfrak m$ soit un idéal \emph{maximal de carré nul}, donc $K$ un
\emph{corps}, et notons $V$ l'idéal $\mathfrak m$ considéré comme $K$-espace
vectoriel. Comme $K$ est une algèbre formellement lisse sur son corps premier
$\mathbf F_p$ (21.5.2) et que $A$ est une algèbre sur $\mathbf F_p$ (de sorte
que $\Omega^1_A=\Omega^1_{A/\mathbf F_p}$), on déduit de (20.5.14) que l'on a
une \emph{suite exacte}
\begin{equation}
\label{0.22.4.6.1-fr}
  0\longrightarrow V\longrightarrow\Omega^1_A\otimes_AK
  \longrightarrow\Omega^1_K\longrightarrow0.
  \tag{22.4.6.1}
\end{equation}
On a donc par (22.4.5.3) un homomorphisme noté
\begin{equation}
\label{0.22.4.6.2-fr}
  \chi'_{B/A}:\Xi_{B_{(K)}/K}\longrightarrow V
  \tag{22.4.6.2}
\end{equation}
que nous appellerons \emph{homomorphisme caractéristique} de la $A$-algèbre
$B$ (relative à $i$, $j$ et $\mathfrak m$). Il résulte des définitions que pour
que $\chi'_{B/A}$ soit \emph{injectif}, il faut et il suffit que
$\pi'_{B_{(K)}/K}$ soit injectif, le noyau de ce dernier étant évidemment
contenu dans $\Xi_{B_{(K)}/K}$.
\end{env}

\begin{proposition}[22.4.7]
\label{0.22.4.7-fr}
Soient $A$ un anneau local artinien dont l'idéal maximal $\mathfrak m$ est de
carré nul, $K=A/\mathfrak m$ son corps résiduel, $V$ l'idéal $\mathfrak m$
considéré comme $K$-espace vectoriel~; on suppose que $A$ est de
caractéristique $p>0$~; soient
\begin{equation}
\label{0.22.4.7.1-fr}
  F_i(T_i)=T_i^p-f_i\qquad\text{avec }f_i\in A\quad(1\leqslant i\leqslant r)
  \tag{22.4.7.1}
\end{equation}
et désignons par $B$ l'anneau quotient de $A[T_1,\ldots,T_r]$ par l'idéal
engendré par les $r$ polynômes $F_i$. Alors~:
\begin{enumerate}
  \item[\rm(i)] $B$ est un anneau local.
  \item[\rm(ii)] Si $\mathfrak m'$ est l'idéal maximal de $B$,
  $K'=B/\mathfrak m'$ son corps résiduel, $V'$ le $B$-module
  $\mathfrak m'/\mathfrak m'^2$ considéré comme $K'$-espace vectoriel, alors,
  pour que $\operatorname{rg}_K(V)=\operatorname{rg}_{K'}(V')$, il faut et il suffit que l'homomorphisme
  caractéristique de $K$-espaces vectoriels
  \begin{equation}
  \label{0.22.4.7.2-fr}
    \chi'_{B/A}:\Xi_{B_{(K)}/K}\longrightarrow V
    \tag{22.4.7.2}
  \end{equation}
  (cf. (22.4.6.2)) soit injectif.
\end{enumerate}
\end{proposition}

\begin{proof}
Désignons encore par $a_i$ $(1\leqslant i\leqslant r)$ la classe de $f_i$
mod.~$\mathfrak m$~; les $a_i$ ne sont plus nécessairement $p$-indépendants
sur $K^p$, mais on peut toujours supposer que la sous-famille

\oldpage[0\textsubscript{IV-1}]{199}

$(a_i)_{s+1\leqslant i\leqslant r}$ est $p$-libre sur $K^p$ et forme une
$p$-base sur $K^p$ du corps $K^p(a_1,\ldots,a_r)$ de sorte que l'on a
\[
  a_i\in K^p(a_{s+1},\ldots,a_r)\qquad\text{pour }1\leqslant i\leqslant s.
\]
Désignons par $A''$ l'anneau quotient de $A[T_{s+1},\ldots,T_r]$ par l'idéal
engendré par les $F_i$ d'indice $i\geqslant s+1$~; alors, par (22.4.4), $A''$
est un anneau local dont l'idéal maximal est $\mathfrak m''=\mathfrak mA''$
(donc de carré nul) et le corps résiduel
$K''=K(a_{s+1}^{1/p},\ldots,a_r^{1/p})$. On a donc
$a_i\in K''^p$ pour $1\leqslant i\leqslant s$ et par suite, dans $A''$,
l'élément $f_i$, pour $1\leqslant i\leqslant s$, peut s'écrire
\begin{equation}
\label{0.22.4.7.3-fr}
  f_i=g_i^p+h_i
  \tag{22.4.7.3}
\end{equation}
où $g_i\in A''$ et $h_i\in\mathfrak m''$ (comme $\mathfrak m''^2=0$ et
$p\mathfrak m''=0$, il est immédiat que $h_i$ est déterminé de façon
\emph{unique} par ces conditions). Remplaçant $T_i$ par $T_i+g_i$, on voit donc
que si l'on pose
\[
  G_i(T_i)=T_i^p+\sum_{k=1}^{p-1}\binom pk g_i^{p-k}T_i^k-h_i
  \qquad(1\leqslant i\leqslant s),
\]
$B$ est l'anneau quotient de $A''[T_1,\ldots,T_s]$ par l'idéal engendré par les
$G_i$~; or, tous les coefficients de $G_i$ sauf le coefficient dominant
appartiennent à $\mathfrak m''$, donc on est dans la situation de (22.4.3),
$B$ est un anneau local, et son corps résiduel $K'$ est \emph{isomorphe} à
$K''$, ce qui prouve (i) (qui d'ailleurs résulte directement de ce que
$B^p\subset A$, et par suite il n'y a qu'un seul idéal de $B$ au-dessus de
$\mathfrak m$).

Notons maintenant que pour que (22.4.7.2) soit injectif, il faut et il suffit
qu'il en soit ainsi de
\[
  \pi'_{B_{(K)}/K}:\Omega^1_{B_{(K)}/K}\otimes_{B_{(K)}}K
  \longrightarrow\Omega^1_A\otimes_AK
\]
(par (22.4.6)). Or $B_{(K)}=B/\mathfrak mB$ est l'algèbre quotient de
$C=K[T_1,\ldots,T_r]$ par l'idéal engendré par les polynômes
$\overline F_i=T_i^p-a_i$, et comme $d_{C/K}\overline F_i=0$ il résulte de
(20.5.13) que $\Omega^1_{B_{(K)}/K}$ est un $B_{(K)}$-module libre ayant pour
\emph{base} les $d_{B_{(K)}/K}t_i$, où $t_i$ est l'image canonique de $T_i$
$(1\leqslant i\leqslant r)$. Mais par définition (22.4.5) l'image canonique par
$j':B_{(K)}\longrightarrow A$ d'une classe mod.~$\mathfrak mB$ d'un élément
$x\in B$ est l'élément $x^p\in A$~; on en déduit aussitôt, puisque $t_i^p=f_i$,
que l'image de $d_{B_{(K)}/K}t_i\otimes1$ par $\pi'_{B_{(K)}/K}$ est
$d_Af_i\otimes1$ dans $\Omega^1_A\otimes_AK$~; la condition que
$\pi'_{B_{(K)}/K}$ soit injectif est donc équivalente au fait que les
$d_Af_i\otimes1_K$ sont \emph{linéairement indépendants sur $K$}, ou encore
que leurs images $d_Af_i\otimes1_{K'}$ sont \emph{linéairement indépendants
sur $K'$} dans $\Omega^1_A\otimes_AK'$.

Or, si l'on applique (22.4.4) à la $A$-algèbre $A''$, on voit que le noyau de
l'homomorphisme canonique
$\Omega^1_A\otimes_AK'\longrightarrow\Omega^1_{A''}\otimes_{A''}K'$ a pour
base les $d_Af_i\otimes1_{K'}$ pour $s+1\leqslant i\leqslant r$. La condition
précédente équivaut donc encore à dire que les $d_{A''}f_i\otimes1_{K'}$ sont
linéairement indépendants dans $\Omega^1_{A''}\otimes_{A''}K'$ pour
$1\leqslant i\leqslant s$. Or, on a $d_{A''}f_i=d_{A''}h_i$ par (22.4.7.3)~;
d'autre part, si $V''$ est le $K'$-espace vectoriel $\mathfrak m''$, on a
$\operatorname{rg}_{K'}(V'')=\operatorname{rg}_K(V)$ (22.4.1, (ii)), et la condition
$\operatorname{rg}_{K'}(V')=\operatorname{rg}_K(V)$ est donc équivalente à
$\operatorname{rg}_{K'}(V'')=\operatorname{rg}_{K'}(V')$. Mais il résulte de (22.4.3) que cette dernière
relation est équivalente au fait que les classes $h_i$

\oldpage[0\textsubscript{IV-1}]{200}

dans $V''$ sont \emph{linéairement indépendantes sur $K'$}. Or, dans
$\Omega^1_{A''}\otimes_{A''}K'$, les $d_{A''}h_i\otimes1_{K'}$ sont les images
canoniques (20.5.11.2) des $h_i$ par l'injection
$V''\longrightarrow\Omega^1_{A''}\otimes_{A''}K'$ (cf. (22.4.6.1)), et cela
achève de prouver (22.4.7).
\end{proof}

On a en outre prouvé~:

\begin{corollary}[22.4.7.4]
\label{0.22.4.7.4-fr}
Pour que $\operatorname{rg}_K(V)=\operatorname{rg}_{K'}(V')$, il faut et il suffit que les éléments
$d_Af_i\otimes1_K$ $(1\leqslant i\leqslant r)$ soient linéairement
indépendants dans le $K$-espace vectoriel $\Omega^1_A\otimes_AK$.
\end{corollary}

\begin{remark}[22.4.8]
\label{0.22.4.8-fr}
On peut définir les homomorphismes (22.4.5.2) et (22.4.6.2) dans des conditions
un peu différentes. Soient $A$ un anneau, $\mathfrak m$ un idéal \emph{de carré
nul} dans $A$, $K=A/\mathfrak m$, $p$ un nombre premier tel que
$p\cdot1\in\mathfrak m$ (autrement dit $K$ est de caractéristique $p$, mais non
nécessairement $A$). Soit d'autre part $B$ une $A$-algèbre qui est un $A$-module
\emph{fidèlement plat} (de sorte que $A\subset B$), et supposons que l'on ait
\begin{equation}
\label{0.22.4.8.1-fr}
  y^p\in A+pB\subset A+\mathfrak mB
  \tag{22.4.8.1}
\end{equation}
pour tout $y\in B$. Si $y^p=x+pz$ avec $x\in A$, $z\in B$, la \emph{classe}
mod.~$(A\cap pB)$ de $x$ est bien déterminée, et comme $A\cap pB=pA$ puisque
$B$ est un $A$-module fidèlement plat, on a ainsi défini une application
canonique $j:B\longrightarrow A/pA$, telle que, pour $x\in A$, $j(x)$ soit la
classe de $x^p$ mod.~$pA$. D'ailleurs, si $y'\in B$ et
$y'^p=x'+pz'$ avec $x'\in A$, $z'\in B$, il est immédiat que les éléments
$(y+y')^p-x-x'$ et $(yy')^p-xx'$ appartiennent à $pB$ en vertu du fait que les
coefficients binomiaux $\binom pi$ sont des multiples de $p$ pour
$1\leqslant i\leqslant p-1$. L'application $j$ est donc un
\emph{homomorphisme d'anneaux}. Par composition, on déduit de $j$ un
homomorphisme $j':B\xrightarrow{j}A/pA\longrightarrow A/\mathfrak m=K$ et par
suite un homomorphisme $j_0:B_{(K)}\longrightarrow K$ tel que $j'$ soit le
composé $B\longrightarrow B_{(K)}=B\otimes_AK\xrightarrow{j_0}K$~; en outre,
si $i_0:K\longrightarrow B_{(K)}$ est l'application canonique, $i_0$ et $j_0$
vérifient les conditions de (21.3.1) pour les anneaux \emph{de caractéristique
$p$}, $K$ et $B_{(K)}$. Cela étant, pour définir un $K$-homomorphisme
$\pi'_{B_{(K)}/K}:\Omega^1_{B_{(K)}/K}\otimes_{B_{(K)}}K_{[j_0]}
\longrightarrow\Omega^1_A\otimes_AK$, il suffit (20.4.8) de définir une
$K$-dérivation
\begin{equation}
\label{0.22.4.8.2-fr}
  D_0:B_{(K)}\longrightarrow\Omega^1_A\otimes_AK
  \tag{22.4.8.2}
\end{equation}
où le second membre est considéré comme un $B_{(K)}$-module au moyen de $j_0$.
Pour cela, il suffit de définir une $A$-dérivation
\begin{equation}
\label{0.22.4.8.3-fr}
  D:B\longrightarrow\Omega^1_A\otimes_AK
  \tag{22.4.8.3}
\end{equation}
qui \emph{s'annule dans} $\mathfrak mB$ (le second membre étant considéré comme
$B$-module au moyen de $j'$). Or, si l'on a la relation $y^p=x+pz$ avec
$y\in B$, $x\in A$, $z\in B$, l'élément $d_Ax\otimes1_K$ de
$\Omega^1_A\otimes_AK$ ne dépend que de la classe $j(y)$ mod.~$pA$ de $x$
(donc seulement de $y$), car pour tout $t\in A$ on a
$d_A(pt)=p\,d_A(t)$, donc $d_A(pt)\otimes1=0$ dans
$\Omega^1_A\otimes_AK=\Omega^1_A/\mathfrak m\Omega^1_A$, puisque
$p\cdot1\in\mathfrak m$. On peut donc prendre
\[
  D(y)=d_Ax\otimes1_K.
\]
Le fait que $D$ est une dérivation résulte aisément de ce que $j$ est un
homomorphisme d'anneaux~; en outre, pour $t\in A$, on a
$d_A(t^px)=t^p d_A(x)+pt^{p-1}x d_A(t)$ et l'on en tire que
$D(ty)=tD(y)$, donc $D$ est une $A$-dérivation~; enfin, si l'on a de plus
$t\in\mathfrak m$, on a

\oldpage[0\textsubscript{IV-1}]{201}

$D(ty)=tD(y)=0$ puisque $\Omega^1_A\otimes_AK$ est annulé par $\mathfrak m$.
On a ainsi défini l'homomorphisme (22.4.5.2) cherché, et il est immédiat de
vérifier que l'homomorphisme composé
\[
  \Theta_{B_{(K)}/K}\xrightarrow{\pi'_{B_{(K)}/K}}
  \Omega^1_A\otimes_AK\longrightarrow\Omega^1_K
\]
est encore l'homomorphisme canonique $\pi_{B_{(K)}/K}$ (relatif à $i_0$ et
$j_0$). On a donc aussi un homomorphisme canonique analogue à (22.4.5.3).

Si maintenant $K$ est un \emph{corps}, $A/pA$ est une $\mathbf F_p$-algèbre,
et l'on a la suite exacte (20.5.14)
\[
  0\longrightarrow V/pA\longrightarrow
  \Omega^1_{A/pA}\otimes_{A/pA}K\longrightarrow\Omega^1_K\longrightarrow0
\]
et d'autre part on a la suite exacte (20.5.12.1)
\[
  pA\longrightarrow\Omega^1_A\otimes_A(A/pA)
  \longrightarrow\Omega^1_{A/pA}\longrightarrow0
\]
qui montre, en tensorisant avec $K$, que les $K$-espaces vectoriels
$\Omega^1_A\otimes_AK$ et
$\Omega^1_{A/pA}\otimes_{A/pA}K$ sont isomorphes. L'analogue de (22.4.6.2) est
donc ici un homomorphisme
\begin{equation}
\label{0.22.4.8.4-fr}
  \Xi_{B_{(K)}/K}\longrightarrow V/pA=\mathfrak m/(\mathfrak m^2+pA).
  \tag{22.4.8.4}
\end{equation}

Ceci dit, \emph{le critère (22.4.7), où l'on remplace $V$ et $V'$ par $V/pA$
et $V'/pB$ respectivement, reste valable en supposant seulement que $K$ soit
de caractéristique $p>0$} (autrement dit, que $p\cdot1\in\mathfrak m$ dans
$A$), \emph{et que les $F_i(T_i)$ soient de la forme plus générale
(22.4.4.1)}~: en effet, $B$ est un $A$-module libre (donc fidèlement plat), et
il suffit de reprendre le raisonnement de (22.4.7), en y remplaçant (22.4.6.2)
par (22.4.8.4) et $V''$ par $V''/pA''$.
\end{remark}

\subsection{Algèbres géométriquement régulières et algèbres formellement lisses}
\label{subsection:0.22.5-fr}

\begin{proposition}[22.5.1]
\label{0.22.5.1-fr}
Soient $A$ un anneau local régulier, $A'$ un anneau local contenant $A$ et qui
est une $A$-algèbre finie, $\mathfrak m$ (resp. $\mathfrak m'$) l'idéal maximal
de $A$ (resp. $A'$), $K=A/\mathfrak m$ (resp. $K'=A'/\mathfrak m'$) son corps
résiduel. Pour que $A'$ soit un anneau régulier, il faut et il suffit que l'on
ait
\begin{equation}
\label{0.22.5.1.1-fr}
  \operatorname{rg}_K(\mathfrak m/\mathfrak m^2)
  =\operatorname{rg}_{K'}(\mathfrak m'/\mathfrak m'^2).
  \tag{22.5.1.1}
\end{equation}
\end{proposition}

\begin{proof}
En effet, le premier membre de (22.5.1.1) est $\dim(A)$ (17.1.1) et comme
$A\subset A'$ et que $A'$ est une $A$-algèbre finie,
$\dim(A')=\dim(A)$ (16.1.5)~; l'égalité (22.5.1.1) est donc nécessaire et
suffisante pour que $A'$ soit régulier (17.1.1).
\end{proof}

\begin{remark}[22.5.2]
\label{0.22.5.2-fr}
\begin{enumerate}
  \item[\rm(i)] Soient $A_1=A/\mathfrak m^2$,
  $A'_1=A'\otimes_AA_1=A'/\mathfrak m^2A'$~; comme
  $\mathfrak m^2A'\subset\mathfrak m'^2$,
  $\mathfrak m'_1=\mathfrak m'/\mathfrak m^2A'$ est l'idéal maximal de $A'_1$
  et $\mathfrak m'_1/\mathfrak m_1'^2$ est un $K'$-espace vectoriel isomorphe
  à $\mathfrak m'/\mathfrak m'^2$~; la condition de régularité de $A'$ ne
  dépend donc que de la structure de la $A_1$-algèbre $A'_1$, ce qui va nous
  permettre ci-dessous d'appliquer les résultats préliminaires de (22.4) sur
  les algèbres finies sur des anneaux artiniens.
  \item[\rm(ii)] Il résulte de la démonstration de (22.5.1) et de (16.2.6.1)
  que l'on a en tout cas
  \begin{equation}
  \label{0.22.5.2.1-fr}
    \operatorname{rg}_K(\mathfrak m/\mathfrak m^2)
    \leqslant\operatorname{rg}_{K'}(\mathfrak m'/\mathfrak m'^2).
    \tag{22.5.2.1}
  \end{equation}

  \oldpage[0\textsubscript{IV-1}]{202}

  \item[\rm(iii)] Supposons que $A$ et $A'$ vérifient les hypothèses générales
  de (22.4.1). On sait (19.8.10) qu'il existe un anneau local \emph{régulier}
  $B$, d'idéal maximal $\mathfrak n$, tel que $A$ soit isomorphe à
  $B/\mathfrak n^2$~; désignons par $G_i\in B[T_i]$ un polynôme unitaire dont
  l'image canonique dans $A[T_i]$ soit $F_i$ $(1\leqslant i\leqslant r)$~;
  alors, si $B'$ est l'anneau quotient de $B[T_1,\ldots,T_r]$ par l'idéal
  engendré par les $G_i$, il est clair que $B'$ est un $B$-module libre de rang
  $d$ et que $A'=B'\otimes_BA=B'/\mathfrak n^2B'$~; comme $A'$ est supposé être
  un anneau local, il en est de même de $B'$, et si $\mathfrak n'$ est l'idéal
  maximal de $B'$, $B'/\mathfrak n'$ est isomorphe à $K'$ et
  $\mathfrak n'/\mathfrak n'^2$ isomorphe à $\mathfrak m'/\mathfrak m'^2$ en
  tant qu'espace vectoriel sur $K'$~; l'inégalité (22.5.2.1) montre donc que
  l'on a, sous les hypothèses de (22.4.1),
  \begin{equation}
  \label{0.22.5.2.2-fr}
    \operatorname{rg}_K(V)\leqslant\operatorname{rg}_{K'}(\mathfrak m'/\mathfrak m'^2).
    \tag{22.5.2.2}
  \end{equation}
\end{enumerate}
\end{remark}

\begin{corollary}[22.5.3]
\label{0.22.5.3-fr}
Soient $A$ un anneau local régulier, $\mathfrak m$ son idéal maximal,
$T_i$ $(1\leqslant i\leqslant r)$ des indéterminées, et pour chaque $i$, soit
\begin{equation}
\label{0.22.5.3.1-fr}
  F_i(T_i)=T_i^{d_i}+\sum_{1\leqslant k\leqslant d_i}c_{ik}T_i^{d_i-k}
  \tag{22.5.3.1}
\end{equation}
un polynôme unitaire de $A[T_i]$, tel que $d_i\geqslant2$ pour tout $i$ et
$c_{ik}\in\mathfrak m$. Soit $A'$ le quotient de l'anneau de polynômes
$A[T_1,\ldots,T_r]$ par l'idéal engendré par les $r$ polynômes $F_i$. Alors~:
\begin{enumerate}
  \item[\rm(i)] $A'$ est un anneau local, et son corps résiduel $K'$ est
  isomorphe au corps résiduel $K$ de $A$.
  \item[\rm(ii)] Pour que $A'$ soit un anneau régulier, il faut et il suffit
  que les classes $\xi_i$ mod.~$\mathfrak m^2$ des éléments $c_{i,d_i}$
  $(1\leqslant i\leqslant r)$ soient linéairement indépendantes sur $K$.
\end{enumerate}
\end{corollary}

\begin{proof}
En effet, avec les notations de (22.5.2, (i)), si $\overline F_i$ désigne le
polynôme de $A_1[T_i]$ obtenu en appliquant aux coefficients de $F_i$
l'homomorphisme $A\longrightarrow A_1$, $A'_1$ se définit à partir de $A_1$ et
des polynômes $\overline F_i$ comme dans (22.4.1), et la conclusion résulte
donc de (22.5.1) et de (22.4.3).
\end{proof}

\begin{theorem}[22.5.4]
\label{0.22.5.4-fr}
Soient $A$ un anneau local régulier, $K$ son corps résiduel~; on suppose que
$K$ est de caractéristique $p>0$~; soient
\begin{equation}
\label{0.22.5.4.1-fr}
  F_i(T_i)=T_i^p+p\sum_{k=1}^{p-1}c_{ik}T_i^{p-k}-f_i
  \tag{22.5.4.1}
\end{equation}
avec $c_{ik}\in A$ et $f_i\in A$ $(1\leqslant i\leqslant r)$~; soit $B$
l'anneau quotient de $A[T_1,\ldots,T_r]$ par l'idéal engendré par les $r$
polynômes $F_i$. Alors $B$ est un anneau local, et les conditions suivantes
sont équivalentes~:
\begin{enumerate}
  \item[a)] $B$ est régulier~;
  \item[b)] les éléments $d_Af_i\otimes1_K$ $(1\leqslant i\leqslant r)$ sont
  linéairement indépendants dans le $K$-espace vectoriel
  $\Omega^1_A\otimes_AK$~;
  \item[c)] si $B'=B/\mathfrak m^2B$ et si l'on pose
  $\Xi_{B_{(K)}/K}=\Xi_{B'_{(K)}/K}$ (22.4.5), l'homomorphisme
  caractéristique $\Xi_{B_{(K)}/K}\longrightarrow\mathfrak m/\mathfrak m^2$
  (22.4.6.2) est injectif.
\end{enumerate}
\end{theorem}

\begin{proof}
En effet, $B'$ se définit encore à partir de $A_1=A/\mathfrak m^2$ et des
polynômes $\overline F_i$ obtenus en appliquant aux $F_i$ l'homomorphisme
$A\longrightarrow A_1$. Compte tenu de ce que $\Omega^1_A\otimes_AK$ est
canoniquement isomorphe à $\Omega^1_{A_1}\otimes_{A_1}K$ (20.5.12, (ii)), le
théorème résulte de (22.5.1) et de (22.4.7.4), en utilisant la remarque
(22.4.8) lorsque $A_1$ n'est pas de caractéristique $p$.
\end{proof}

\oldpage[0\textsubscript{IV-1}]{203}

\begin{env}[22.5.5]
\label{0.22.5.5-fr}
Soient $k$ un corps de caractéristique $p>0$, $A$ un anneau local, $k'$ une
extension \emph{radicielle finie} de $k$ telle que $k'^p\subset k$,
$\mathfrak m$ l'idéal maximal de $A$, $K=A/\mathfrak m$ son corps résiduel~;
rappelons que $A'=A\otimes_kk'$ est un anneau local dont le corps résiduel est
le même que celui de l'anneau local $K\otimes_kk'$ (Bourbaki,
\emph{Alg. comm.}, chap.~V, §~2, n\textsuperscript{o}~3, lemme~4). Notons que
l'on a $A'_{(K)}=A'\otimes_AK=K\otimes_kk'$, et par suite (21.3.6)
\begin{equation}
\label{0.22.5.5.1-fr}
  \Theta_{A'_{(K)}/K}=\Theta_{k'/k}\otimes_kK
  \tag{22.5.5.1}
\end{equation}
à un isomorphisme canonique près. En outre, on sait (21.4.7) que l'on a
$\Xi_{k'/k}=0$, autrement dit l'homomorphisme canonique (21.3.2.2)
\[
  \pi_{k'/k}:\Theta_{k'/k}\longrightarrow\Omega^1_k
\]
est injectif, si bien que l'on a une \emph{injection canonique}
\begin{equation}
\label{0.22.5.5.2-fr}
  \Theta_{A'_{(K)}/K}\longrightarrow\Omega^1_k\otimes_kK
  \tag{22.5.5.2}
\end{equation}
qui s'explicite de la façon suivante (compte tenu de (22.5.5.1))~: pour tout
$x'\in k'$, à l'élément
$d_{A'_{(K)}/K}(x'\otimes1_K)\otimes1_K$ de
$\Omega^1_{A'_{(K)}/K}\otimes_{A'_{(K)}}K$, on fait correspondre l'élément
$d_k(x'^p)\otimes1_K$ de $\Omega^1_k\otimes_kK$.
\end{env}

\begin{lemma}[22.5.6]
\label{0.22.5.6-fr}
Avec les notations de (22.5.5)~:
\begin{enumerate}
  \item[\rm(i)] Le noyau de l'homomorphisme canonique
  $\pi_{A'_{(K)}/K}$ (21.3.2.2) s'identifie par (22.5.5.2) à
  $(\Theta_{k'/k}\otimes_kK)\cap\Upsilon_{K/k}$.
  \item[\rm(ii)] Soient $A_1=A/\mathfrak m^2$,
  $A'_1=A_1\otimes_kk'=A'/\mathfrak m^2A'$~; alors l'homomorphisme
  caractéristique $\chi'_{A'_1/A_1}$ (22.4.6.2) s'identifie à la restriction de
  l'homomorphisme caractéristique
  $\chi_{A/k}:\Upsilon_{K/k}\longrightarrow V=\mathfrak m/\mathfrak m^2$
  (20.6.24) à $(\Theta_{k'/k}\otimes_kK)\cap\Upsilon_{K/k}$.
\end{enumerate}
\end{lemma}

\begin{proof}
(i) En vertu de (20.6.9), appliqué en remplaçant $A$ par le corps premier, $B$
par $k$, $C$ par $K$ et $E$ par $A_1=A/\mathfrak m^2$ (extension de $K$ par
$V=\mathfrak m/\mathfrak m^2$), il suffit (cf. (22.5.5)) de vérifier que si
$x'\in k'$ et si $z\in A_1$ est un élément dont la classe mod.
$\mathfrak m/\mathfrak m^2$ est $x'^p\in k$, alors l'image de
$d_k(x'^p)\otimes1_K$ par l'homomorphisme canonique
$\Omega^1_k\otimes_kK\longrightarrow\Omega^1_{A_1}\otimes_{A_1}K$ est
$d_{A_1}z\otimes1_K$, ce qui résulte des définitions.

(ii) Comme $A'_{(K)}=K\otimes_kk'=(A'_1)_{(K)}$ (en vertu de la relation
$K=A_1/(\mathfrak m/\mathfrak m^2)$), $\Theta_{A'_{(K)}/K}$ s'identifie à
$\Theta_{(A'_1)_{(K)}/K}$ et la vérification résulte du même calcul que dans
(i) et de la définition (22.4.5.2) de $\pi'_{(A'_1)_{(K)}/K}$.
\end{proof}

\begin{proposition}[22.5.7]
\label{0.22.5.7-fr}
Soient $k$ un corps de caractéristique $p>0$, $A$ une $k$-algèbre qui est un
anneau local régulier, $\mathfrak m$ son idéal maximal, $K=A/\mathfrak m$ son
corps résiduel~; soit $V=\mathfrak m/\mathfrak m^2$, considéré comme
$K$-espace vectoriel. Soient $k'$ une extension finie de $k$ telle que
$k'^p\subset k$, et $A'=A\otimes_kk'$~; pour que l'anneau local $A'$ soit
régulier, il faut et il suffit que la restriction de l'homomorphisme
caractéristique $\chi_{A/k}:\Upsilon_{K/k}\longrightarrow V$ (20.6.24) à
$(\Theta_{k'/k}\otimes_kK)\cap\Upsilon_{K/k}$ soit injective.
\end{proposition}

\begin{proof}
Il s'agit de prouver que la condition de l'énoncé est équivalente à
(22.5.1.1), en désignant par $\mathfrak m'$ l'idéal maximal de $A'$, par
$K'=A'/\mathfrak m'$ son corps résiduel. Avec les notations de la remarque
(22.5.2, (i)), on a $A'_1=A_1\otimes_kk'$. Or, soit
$(a_i)_{1\leqslant i\leqslant r}$ une $p$-base de $k'$ sur $k$, de sorte que
$a_i^p=b_i\in k$, et $k'$ est isomorphe au quotient de l'anneau de

\oldpage[0\textsubscript{IV-1}]{204}

polynômes $k[T_1,\ldots,T_r]$ par l'idéal engendré par les polynômes
$F_i=T_i^p-b_i$ $(1\leqslant i\leqslant r)$~; on en déduit que $A'_1$ est
isomorphe au quotient de l'anneau de polynômes $A_1[T_1,\ldots,T_r]$ par
l'idéal engendré par les $F_i$. On peut alors appliquer le critère (22.4.7), et
la condition de l'énoncé équivaut donc au fait que $\chi'_{A'_1/A_1}$ soit
injectif~; on conclut à l'aide de (22.5.6).
\end{proof}

Nous pouvons maintenant démontrer la réciproque de (19.6.5)~:

\begin{theorem}[22.5.8]
\label{0.22.5.8-fr}
Soient $k$ un corps, $p$ son exposant caractéristique, $A$ une $k$-algèbre
locale noethérienne. Les conditions suivantes sont équivalentes~:
\begin{enumerate}
  \item[a)] $A$ est une $k$-algèbre formellement lisse (pour sa topologie
  préadique).
  \item[b)] $A$ est géométriquement régulier sur $k$.
  \item[b')] Pour toute extension finie $k'$ de $k$ telle que $k'^p\subset k$,
  $A'=A\otimes_kk'$ est un anneau régulier.
\end{enumerate}
\end{theorem}

\begin{proof}
Compte tenu de (19.6.6), on peut se borner au cas $p>1$.

Le fait que a) implique b) n'est autre que (19.6.5), et b) entraîne
trivialement b'). Montrons que b') entraîne que les conditions de (22.2.2) sont
satisfaites~; c'est évident pour la première (17.1.1) puisque b') entraîne
d'abord que $A$ est régulier. D'autre part, soit $(x_\alpha)_{\alpha\in I}$ une
$p$-base de $k^{1/p}$ sur $k$~; on sait que les éléments $d_k(x_\alpha^p)$
forment une base du $k$-espace vectoriel $\Omega^1_k$ (21.4.5), donc les
éléments $d_k(x_\alpha^p)\otimes1_K$ forment une base du $K$-espace vectoriel
$\Omega^1_k\otimes_kK$. On en conclut (cf. (22.5.5)) que lorsque $k'$ parcourt
l'ensemble des sous-extensions de $k^{1/p}$, finies sur $k$, la famille des
sous-espaces $\Theta_{k'/k}\otimes_kK$ de $\Omega^1_k\otimes_kK$ est filtrante
croissante et a pour réunion $\Omega^1_k\otimes_kK$. Il résulte alors de
(22.5.7) que la condition b') entraîne que $\chi_{A/k}$ est injectif, ce qui
n'est autre que la condition (ii) de (22.2.2).
\end{proof}

\begin{corollary}[22.5.9]
\label{0.22.5.9-fr}
Soient $k$ un corps, $A$ une $k$-algèbre locale noethérienne formellement
lisse. Alors, pour tout idéal premier $\mathfrak p$ de $A$, l'anneau local
$A_{\mathfrak p}$ est une $k$-algèbre formellement lisse (pour la topologie
$\mathfrak p$-préadique).
\end{corollary}

\begin{proof}
En effet, avec les notations de (22.5.8, b')), il suffit de montrer que
l'anneau local $A_{\mathfrak p}\otimes_kk'$ est régulier~; mais comme c'est un
anneau de fractions de l'anneau local $A\otimes_kk'$ et que ce dernier est
régulier par hypothèse, il en est de même de $A_{\mathfrak p}\otimes_kk'$
(17.3.6).
\end{proof}

\begin{remark}[22.5.10]
\label{0.22.5.10-fr}
\begin{enumerate}
  \item[\rm(i)] La conclusion de (22.5.8) est en défaut lorsque, dans la
  condition b'), on se borne aux extensions \emph{monogènes} $k'=k(x)$ de $k$
  (avec $x^p\in k$). Il se peut en effet que pour toutes ces extensions $k'$,
  on ait
  $\Upsilon_{K/k}\cap(\Theta_{k'/k}\otimes_kK)=0$, bien que
  $\Upsilon_{K/k}\neq0$~; cela signifie que pour tout $x\in k^{1/p}$ tel que
  $x\notin k$, on doit avoir
  $d_k(x^p)\otimes1_K\notin\Upsilon_{K/k}$, ou encore $d_K(x^p)\neq0$,
  c'est-à-dire $x^p\notin K^p$~; autrement dit, on doit avoir
  $K^p\cap k=k^p$, bien que $K$ soit une extension \emph{inséparable} de $k$.
  Or, on construit aisément des exemples de telles extensions~: partons d'un
  corps parfait $k_0$ de caractéristique $p>0$, soient $X$, $Y$, $Z$ trois
  indéterminées et posons $k=k_0(X,Y)$ et
  $K=k(X^{1/p}+ZY^{1/p})$~; on vérifie aisément que $K$ vérifie les conditions
  précédentes, donc $\Upsilon_{K/k}\neq0$. Appliquons maintenant le lemme
  (22.2.5.2) en prenant $V=K$ et pour $h$ l'homomorphisme nul~; $A$ est alors
  un anneau de valuation discrète ayant $K$ pour corps résiduel, avec
  $\chi_{A/k}=0$, et n'est donc pas une $k$-algèbre formellement lisse par
  (22.2.2)~; toutefois $A\otimes_kk'$ est un anneau régulier pour toute
  extension \emph{monogène} $k'\subset k^{1/p}$ de $k$.

  \oldpage[0\textsubscript{IV-1}]{205}

  \item[\rm(ii)] Le corollaire (22.5.9) amène à considérer la question
  suivante~: si $A$ est une $k$-algèbre noethérienne, à quelles conditions
  l'ensemble des idéaux premiers $\mathfrak p\in\operatorname{Spec}(A)$ tels que
  $A_{\mathfrak p}$ soit une $k$-algèbre formellement lisse est-il
  \emph{ouvert}~? Nous en aborderons plus tard certains cas particuliers.
\end{enumerate}
\end{remark}

\subsection{Critère jacobien de Zariski}
\label{subsection:0.22.6-fr}

\begin{theorem}[22.6.1]
\label{0.22.6.1-fr}
\textnormal{(critère jacobien de lissité formelle).}
Soient $A$, $B$ deux anneaux topologiques,
$u:A\longrightarrow B$ un homomorphisme continu faisant de $B$ une
$A$-algèbre formellement lisse, $\mathfrak J$ un idéal de $B$ (non
nécessairement fermé), $C=B/\mathfrak J$ la $A$-algèbre topologique quotient.
Pour que $C$ soit une $A$-algèbre formellement lisse, il faut et il suffit que
l'homomorphisme canonique (cf. (20.5.11.2))
\begin{equation}
\label{0.22.6.1.1-fr}
  \delta_{C/B/A}:\mathfrak J/\mathfrak J^2
  \longrightarrow\Omega^1_{B/A}\otimes_BC
  \tag{22.6.1.1}
\end{equation}
soit \emph{formellement inversible à gauche} (cf. (19.1.5)).
\end{theorem}

\begin{proof}
En effet, dire que $B$ (resp. $C$) est une $A$-algèbre formellement lisse
signifie (19.4.4) que l'on a
$\operatorname{Exalcotop}_A(B,L)=0$ (resp.
$\operatorname{Exalcotop}_A(C,L)=0$) pour tout $B$-module (resp. $C$-module)
discret $L$ annulé par un idéal ouvert de $B$ (resp. $C$). Comme par hypothèse
$\operatorname{Exalcotop}_A(B,L)=0$ pour tout $C$-module discret $L$ annulé
par un idéal ouvert de $C$, dire que $C$ est une $A$-algèbre formellement lisse
signifie donc que l'homomorphisme canonique
$\operatorname{Exalcotop}_A(C,L)\longrightarrow
\operatorname{Exalcotop}_A(B,L)$ est injectif, et le théorème résulte donc de
(20.7.8).
\end{proof}

Rappelons (19.5.3) que lorsque $B$ et $C$ sont des $A$-algèbres formellement
lisses, $\mathfrak J/\mathfrak J^2$ est un $C$-module topologique
\emph{formellement projectif}~; en outre, les homomorphismes canoniques
$\varphi_n:\mathbf S_C^n(\mathfrak J/\mathfrak J^2)
\longrightarrow\mathfrak J^n/\mathfrak J^{n+1}$ sont des
\emph{bimorphismes formels} lorsque $B$ est supposé être un anneau
\emph{préadmissible}.
\begin{corollary}[22.6.2]
\label{0.22.6.2-fr}
Les hypothèses et notations étant celles de (22.6.1), supposons de plus que
dans $B$ le carré de tout idéal ouvert soit ouvert, et que sur $\mathfrak J$ la
topologie induite par celle de $B$ soit identique à la topologie déduite de
celle de $B$ (19.0) (on notera que ces deux conditions sont vérifiées lorsque
$B$ est noethérien et sa topologie préadique ($0_{\mathrm I}$, 7.3.2), ou si la
topologie de $B$ est la topologie $\mathfrak J$-préadique). Soit
$(\mathfrak b_\lambda)$ un système fondamental d'idéaux ouverts de $B$ et
posons $B_\lambda=B/\mathfrak b_\lambda$ pour tout $\lambda$. Alors~:
\begin{enumerate}
  \item[\rm(i)] Les conditions suivantes sont équivalentes~:
  \begin{enumerate}
    \item[a)] $C$ est une $A$-algèbre formellement lisse.
    \item[b)] Pour tout $\lambda$, l'homomorphisme de $B_\lambda$-modules
    \begin{equation}
    \label{0.22.6.2.1-fr}
      (\mathfrak J/\mathfrak J^2)\otimes_BB_\lambda
      \longrightarrow\Omega^1_{B/A}\otimes_BB_\lambda
      \tag{22.6.2.1}
    \end{equation}
    déduit de $\delta_{C/B/A}$ par tensorisation, est inversible à gauche
    (autrement dit un isomorphisme sur un facteur direct de
    $\Omega^1_{B/A}\otimes_BB_\lambda$).
  \end{enumerate}
  \item[\rm(ii)] Supposons de plus que l'anneau topologique $C$ soit
  préadmissible ($0_{\mathrm I}$, 7.1.2) et

  \oldpage[0\textsubscript{IV-1}]{206}

  soit $\mathfrak L$ un idéal de définition de $C$~; alors les conditions a)
  et b) de (i) équivalent aussi à~:
  \begin{enumerate}
    \item[c)] L'homomorphisme de $(C/\mathfrak L)$-modules
    \[
      (\mathfrak J/\mathfrak J^2)\otimes_C(C/\mathfrak L)
      \longrightarrow\Omega^1_{B/A}\otimes_B(C/\mathfrak L)
    \]
    est inversible à gauche.
  \end{enumerate}
\end{enumerate}
\end{corollary}

\begin{proof}
Les hypothèses entraînent que sur $\Omega^1_{B/A}$ la topologie est déduite de
celle de $B$ (20.4.5), donc (i) résulte de (22.6.1) et de (19.1.7). D'autre
part, rappelons (20.4.9) que puisque $B$ est une $A$-algèbre formellement
lisse, $\Omega^1_{B/A}$ est un $B$-module formellement projectif, donc
$\Omega^1_{B/A}\otimes_BB_\lambda$ est un $B_\lambda$-module projectif pour
tout $\lambda$ (19.2.4)~; l'équivalence de c) avec a) et b) lorsque $C$ est
préadmissible résulte alors de (19.1.9).
\end{proof}

En particulier~:

\begin{corollary}[22.6.3]
\label{0.22.6.3-fr}
Soient $A$ un anneau, $B$ une $A$-algèbre, $\mathfrak J$ un idéal de $B$,
$C=B/\mathfrak J$~; on suppose $A$, $B$, $C$ munis des topologies discrètes
et que $B$ est une $A$-algèbre formellement lisse. Pour que $C$ soit une
$A$-algèbre formellement lisse, il faut et il suffit que l'homomorphisme
canonique (22.6.1.1) soit inversible à gauche.
\end{corollary}

On notera que puisque toute $A$-algèbre $C$ est quotient d'une algèbre de
polynômes $A[T_\alpha]_{\alpha\in I}$, et que cette dernière est formellement
lisse (19.3.3), le critère (22.6.3) permet en principe de reconnaître si $C$
est une $A$-algèbre formellement lisse.

\begin{proposition}[22.6.4]
\label{0.22.6.4-fr}
Soient $A$ un anneau, $B$ une $A$-algèbre formellement lisse (pour les
topologies discrètes), $\mathfrak J$ un idéal de $B$, $C=B/\mathfrak J$~; on
suppose que $\mathfrak J/\mathfrak J^2$ est un $C$-module de type fini. Soient
$\mathfrak p$ un idéal premier de $C$, $k(\mathfrak p)$ le corps résiduel de
$C_{\mathfrak p}$, $\mathfrak p'$ l'image réciproque de $\mathfrak p$ dans
$B$, $\mathfrak q$ l'idéal premier de $A$ image réciproque de $\mathfrak p'$.
Les conditions suivantes sont équivalentes~:
\begin{enumerate}
  \item[a)] $C_{\mathfrak p}$ est une $A_{\mathfrak q}$-algèbre formellement
  lisse (pour les topologies discrètes).
  \item[a')] $C_{\mathfrak p}$ est une $A$-algèbre formellement lisse (pour
  les topologies discrètes).
  \item[b)] L'homomorphisme canonique
  \begin{equation}
  \label{0.22.6.4.1-fr}
    (\mathfrak J/\mathfrak J^2)\otimes_Ck(\mathfrak p)
    \longrightarrow\Omega^1_{B/A}\otimes_Bk(\mathfrak p)
    \tag{22.6.4.1}
  \end{equation}
  est injectif.
  \item[c)] Il existe $f\in C-\mathfrak p$ tel que $C_f$ soit une
  $A$-algèbre formellement lisse.
\end{enumerate}
Si de plus $B$ est noethérien, les conditions précédentes sont encore
équivalentes à~:
\begin{enumerate}
  \item[d)] $C_{\mathfrak p}$ est une $A_{\mathfrak q}$-algèbre formellement
  lisse pour la topologie $\mathfrak p$-préadique sur $C_{\mathfrak p}$ et la
  topologie discrète (ou la topologie $\mathfrak q$-préadique) sur
  $A_{\mathfrak q}$.
\end{enumerate}
\end{proposition}

\begin{proof}
Comme $A_{\mathfrak q}$ est une $A$-algèbre formellement lisse (19.3.5, (iv)
et (ii)), on sait déjà que a) et a') sont équivalentes (19.3.5, (ii)). On a
$C_{\mathfrak p}=B_{\mathfrak p'}/\mathfrak J B_{\mathfrak p'}$, et
$B_{\mathfrak p'}$ est une $A$-algèbre formellement lisse pour la topologie
discrète (19.3.5, (iv))~; en outre
$\Omega^1_{B_{\mathfrak p'}/A}=(\Omega^1_{B/A})_{\mathfrak p'}$ (20.5.9).
L'équivalence de a') et b) résulte alors de l'application de (22.6.3) à
$B_{\mathfrak p'}$ et $C_{\mathfrak p}$, en tenant compte de ce que
$\Omega^1_{B/A}$ est un $B$-module projectif (20.4.9) et (19.2.11), et que
$\mathfrak J/\mathfrak J^2$ est un $B$-module de type fini, et en utilisant
(19.1.12). L'application de (19.1.12) prouve aussi l'équivalence de b) et c).
Enfin, notons que puisque $B_{\mathfrak p'}$ est une
$A_{\mathfrak q}$-algèbre formellement lisse pour les topologies

\oldpage[0\textsubscript{IV-1}]{207}

discrètes (19.3.5, (iv)), $B_{\mathfrak p'}$, muni de la topologie
$\mathfrak p'$-préadique, est encore une $A_{\mathfrak q}$-algèbre
formellement lisse lorsque l'on prend sur $A_{\mathfrak q}$ la topologie
discrète ou la topologie $\mathfrak q$-préadique (19.3.8). Pour montrer
l'équivalence de b) et d) lorsque $B$ est noethérien, appliquons (22.6.2) à
l'anneau $A_{\mathfrak q}$ discret (ou $\mathfrak q$-préadique) et à l'anneau
$B_{\mathfrak p'}$, muni de la topologie $\mathfrak p'$-préadique~; comme
$C_{\mathfrak p}$ est préadmissible et que $\mathfrak pC_{\mathfrak p}$ est
un idéal de définition pour sa topologie, on peut invoquer l'équivalence de
c) et a) dans (22.6.2).
\end{proof}

\begin{corollary}[22.6.5]
\label{0.22.6.5-fr}
Sous les hypothèses générales de (22.6.4), l'ensemble des
$\mathfrak p\in\operatorname{Spec}(C)$ tels que $C_{\mathfrak p}$ soit une $A$-algèbre
formellement lisse (ou une $A_{\mathfrak q}$-algèbre formellement lisse, en
désignant par $\mathfrak q$ l'image réciproque de $\mathfrak p$ dans $A$)
pour les topologies discrètes, est ouvert dans $\operatorname{Spec}(C)$.
\end{corollary}

Cela résulte de la forme c) de (22.6.4). On notera que lorsque $B$ est
noethérien, on peut remplacer les topologies discrètes par les topologies
préadiques.

\begin{corollary}[22.6.6]
\label{0.22.6.6-fr}
Sous les hypothèses générales de (22.6.4), les conditions suivantes sont
équivalentes~:
\begin{enumerate}
  \item[a)] $C$ est une $A$-algèbre formellement lisse (pour les topologies
  discrètes).
  \item[b)] Pour tout $\mathfrak p\in\operatorname{Spec}(C)$ (ou seulement pour tout idéal
  maximal $\mathfrak p$ de $C$), $C_{\mathfrak p}$ est une $A$-algèbre
  formellement lisse (ou une $A_{\mathfrak q}$-algèbre formellement lisse, en
  désignant par $\mathfrak q$ l'image réciproque de $\mathfrak p$ dans $A$)
  pour les topologies discrètes.
\end{enumerate}
En outre, lorsque $B$ est noethérien, on peut remplacer dans b) les topologies
discrètes par les topologies préadiques sur $A_{\mathfrak q}$ et
$C_{\mathfrak p}$.
\end{corollary}

\begin{proof}
La dernière assertion résulte de (22.6.4). En vertu de (22.6.3), la condition
a) équivaut au fait que l'homomorphisme (22.6.1.1) est inversible à gauche~;
l'équivalence de a) et b) résulte donc de (19.1.14), compte tenu du fait que
$\Omega^1_{B/A}$ est un $B$-module projectif et $\mathfrak J/\mathfrak J^2$
un $C$-module de type fini.
\end{proof}

\begin{proposition}[22.6.7]
\label{0.22.6.7-fr}
Soient $k_0$ un corps, $k$ une extension séparable de $k_0$, $C$ une
$k$-algèbre de type fini.
\begin{enumerate}
  \item[\rm(i)] Soit $\mathfrak p$ un idéal premier de $C$. Les conditions
  suivantes sont équivalentes~:
  \begin{enumerate}
    \item[a)] $C_{\mathfrak p}$ est une $k_0$-algèbre formellement lisse pour
    la topologie discrète.
    \item[b)] $C_{\mathfrak p}$ est une $k_0$-algèbre formellement lisse pour
    la topologie $\mathfrak p$-préadique.
    \item[c)] Il existe $f\in C-\mathfrak p$ tel que $C_f$ soit une
    $k_0$-algèbre formellement lisse pour la topologie discrète.
    \item[d)] $C_{\mathfrak p}$ est un anneau géométriquement régulier
    (19.6.5) sur $k_0$.
  \end{enumerate}
  En outre l'ensemble des $\mathfrak p\in\operatorname{Spec}(C)$ ayant l'une quelconque de
  ces propriétés est ouvert dans $\operatorname{Spec}(C)$.
  \item[\rm(ii)] Les conditions suivantes sont équivalentes~:
  \begin{enumerate}
    \item[a)] $C$ est une $k_0$-algèbre formellement lisse pour la topologie
    discrète.
    \item[b)] Tout $\mathfrak p\in\operatorname{Spec}(C)$ vérifie les conditions
    équivalentes de (i).
    \item[c)] Tout idéal maximal $\mathfrak m$ de $C$ vérifie les conditions
    équivalentes de (i).
    \item[d)] Pour toute extension $k'$ de $k_0$, $C\otimes_{k_0}k'$ est un
    anneau régulier (17.3.6)~; on dit encore que $C$ est un anneau
    géométriquement régulier sur $k_0$.
  \end{enumerate}

  \oldpage[0\textsubscript{IV-1}]{208}

  \item[\rm(iii)] Soient $B=k[T_1,\ldots,T_r]$ une algèbre de polynômes,
  $\mathfrak J$ un idéal de $B$ tels que $C$ soit isomorphe à
  $B/\mathfrak J$. Soit $\mathfrak q$ un idéal premier de $B$ contenant
  $\mathfrak J$ et soit $\mathfrak p=\mathfrak q/\mathfrak J$. Les conditions
  de (i) sont alors encore équivalentes à la suivante («~critère jacobien de
  Zariski~»)~:
  \begin{enumerate}
    \item[e)] Il existe un système d'éléments
    $(f_i)_{1\leq i\leq m}$ de $\mathfrak J$, engendrant l'idéal
    $\mathfrak J_{\mathfrak q}$ dans $B_{\mathfrak q}$ et des
    $k_0$-dérivations $D_j$ de $B$ dans lui-même $(1\leq j\leq m)$ tels que
    $\det(D_j(f_i))\notin\mathfrak q$.
  \end{enumerate}
\end{enumerate}
\end{proposition}

\begin{proof}
On sait que $C$ est toujours isomorphe à une $k$-algèbre de la forme
$B/\mathfrak J$. Comme $B$ est une $k$-algèbre formellement lisse (19.3.3) et
que $k$, étant séparable sur $k_0$, est une $k_0$-algèbre formellement lisse
(19.6.1), $B$ est une $k_0$-algèbre formellement lisse (19.3.5, (ii)).
L'équivalence des conditions a), b), c) de (i) résulte donc de (22.6.4), et
celle de a), b), c) dans (ii) résulte de (22.6.6)~; l'équivalence de a) et d)
dans (i) résulte de (22.5.8)~; comme tout localisé de $C\otimes_{k_0}k'$ est
aussi un localisé de $C_{\mathfrak p}\otimes_{k_0}k'$ pour un idéal premier
$\mathfrak p\in\operatorname{Spec}(C)$ convenable, l'équivalence de d) et b) dans (ii)
résulte de l'équivalence de d) et a) dans (i). Enfin, comme
$\Omega^1_{B/k_0}$ est un $B$-module projectif, l'équivalence du critère de
Zariski e) et des autres conditions de (i) découle de (19.1.12) et de
(22.6.3), compte tenu de (20.4.8).
\end{proof}

Zariski s'intéressait en fait à un critère différentiel de régularité pour les
anneaux locaux $C_{\mathfrak p}$. Comme il revient au même de dire qu'un anneau
local noethérien contenant un corps est régulier ou est formellement lisse en
tant qu'algèbre sur son sous-corps premier (19.6.4), on obtient aussitôt un tel
critère en remplaçant $k_0$ par le sous-corps premier de $k$ dans (22.6.7)~; en
particulier, on obtient le résultat suivant, que nous retrouverons plus tard
(IV, 6.12.5) comme cas particulier de résultats plus généraux de Nagata~:

\begin{corollary}[22.6.8]
\label{0.22.6.8-fr}
\textnormal{(Zariski).}
Soit $C$ une algèbre de type fini sur un corps. Alors l'ensemble des
$\mathfrak p\in\operatorname{Spec}(C)$ tels que $C_{\mathfrak p}$ soit un anneau local
régulier est ouvert dans $\operatorname{Spec}(C)$.
\end{corollary}

\begin{remark}[22.6.9]
\label{0.22.6.9-fr}
\begin{enumerate}
  \item[\rm(i)] Les énoncés de (22.6.7) sont encore valables si au lieu de
  supposer $k$ séparable sur $k_0$, on suppose seulement qu'il est de
  \emph{multiplicité radicielle finie}, autrement dit (19.6.6) qu'il existe
  une extension radicielle finie $k'_0$ de $k_0$ telle que le corps résiduel de
  l'anneau artinien local $k\otimes_{k_0}k'_0$ soit une extension séparable
  $k'$ de $k'_0$. Il existe alors un $k'_0$-monomorphisme
  $k'\longrightarrow k\otimes_{k_0}k'_0$ qui, composé avec l'homomorphisme
  canonique $k\otimes_{k_0}k'_0\longrightarrow k'$, donne l'identité, puisque
  $k'$ est une $k'_0$-algèbre formellement lisse (19.6.1)~; on en conclut que
  $C\otimes_{k_0}k'_0$ est muni d'une structure de $k'$-algèbre de type fini,
  et comme $k'$ est séparable sur $k'_0$, on peut appliquer (22.6.7) à cette
  $k'$-algèbre~; on en conclut notre assertion en appliquant (19.4.7). Comme
  nous n'aurons pas à nous servir de cette généralisation, nous en laissons le
  détail au lecteur. Nous ignorons par contre si les résultats de (22.6.7) sont
  valables sans \emph{aucune} hypothèse sur l'extension $k$ de $k_0$.
  \item[\rm(ii)] Sous les hypothèses générales de (22.6.7), soit $(k_\alpha)$
  une famille filtrante décroissante de sous-corps de $k$ contenant $k_0$ et
  telle que
  $\bigcap_\alpha k_\alpha(k^p)=k_0(k^p)$ (où $p$ est l'exposant
  caractéristique de $k_0$). Utilisant une méthode de décompte des dimensions
  due à Nagata, on peut montrer que, dans le critère jacobien (22.6.7, (iii)),
  on peut se borner à des dérivations $D_j$ de $B$ qui soient des
  $k_\alpha$-dérivations pour un $\alpha$ convenable.

  \oldpage[0\textsubscript{IV-1}]{209}

  L'intérêt de ce résultat est qu'il y a toujours de telles familles
  $(k_\alpha)$ pour lesquelles $[k:k_\alpha]$ soit fini (21.8.6). Nous ne
  démontrerons pas ce raffinement du critère de Zariski, dont nous n'aurons pas
  à faire usage. Dans (22.7), on va donner, pour les anneaux locaux complets,
  une variante (due aussi à Nagata) du critère de Zariski, qui se démontre
  essentiellement par la même méthode (avec des difficultés techniques un peu
  plus grandes).
\end{enumerate}
\end{remark}

\subsection{Le critère jacobien de Nagata}
\label{subsection:0.22.7-fr}

\begin{env}[22.7.1]
\label{0.22.7.1-fr}
Le critère jacobien de Nagata est l'analogue du critère jacobien de Zariski,
mais pour des anneaux quotients d'anneaux de \emph{séries formelles} sur un
corps. Nous en donnerons, comme Nagata [31], deux versions, présentées ici
comme critères de lissité formelle.
\end{env}

\begin{proposition}[22.7.2]
\label{0.22.7.2-fr}
Soient $k$ un corps, $A=k[[T_1,\ldots,T_r]]$ un anneau de séries formelles sur
$k$, muni de sa structure usuelle de $k$-algèbre, $\mathfrak q$ un idéal de
$A$, $B=A/\mathfrak q$. Soit $\mathfrak p$ un idéal premier de $A$ contenant
$\mathfrak q$. On suppose qu'il existe une sous-$k$-algèbre $C'$ de
$C=A/\mathfrak p$, isomorphe à une algèbre de séries formelles
$k[[X_1,\ldots,X_s]]$, telle que $C$ soit finie sur $C'$ et que le corps des
fractions $L$ de $C$ soit une extension séparable de
$L'=k((X_1,\ldots,X_s))$ (hypothèse toujours satisfaite si $k$ est de
caractéristique 0). Alors les conditions suivantes sont équivalentes~:
\begin{enumerate}
  \item[a)] $B_{\mathfrak p}=A_{\mathfrak p}/\mathfrak qA_{\mathfrak p}$ est
  une $k$-algèbre formellement lisse pour la topologie
  $\mathfrak p$-préadique.
  \item[b)] Il existe des $k$-dérivations $D_i$ $(1\leq i\leq m)$ de $A$ dans
  lui-même, et des éléments $f_i$ $(1\leq i\leq m)$ de $\mathfrak q$, tels
  que les images des $f_i$ dans $\mathfrak qA_{\mathfrak p}$ engendrent cet
  idéal de $A_{\mathfrak p}$, et que l'on ait
  $\det(D_i(f_j))\notin\mathfrak p$.
  \item[c)] $B_{\mathfrak p}$ est un anneau régulier.
\end{enumerate}
\end{proposition}

\begin{proof}
Le corps résiduel de $B_{\mathfrak p}$ est $k(\mathfrak p)$~; comme
$k((X_1,\ldots,X_s))$ est séparable sur $k$ (21.9.6.4), $k(\mathfrak p)$ est
séparable sur $k$ par hypothèse, et l'on sait déjà que dans ces conditions les
propriétés a) et c) sont équivalentes (19.6.4). D'ailleurs
$A_{\mathfrak p}$ est un anneau régulier ((17.1.4) et (17.3.2)), et comme son
corps résiduel $k(\mathfrak p)$ est séparable sur $k$, $A_{\mathfrak p}$ est
formellement lisse sur $k$ pour la topologie $\mathfrak p$-préadique
(19.6.4). Il résulte donc de (22.6.2, (ii)) que la condition a) équivaut au
fait que l'homomorphisme de $L$-espaces vectoriels
\begin{equation}
\label{0.22.7.2.1-fr}
  j_0:(\mathfrak q/\mathfrak q^2)\otimes_AL
  \longrightarrow\Omega^1_{A_{\mathfrak p}/k}\otimes_{A_{\mathfrak p}}L
  =\widehat\Omega^1_{A/k}\otimes_AL
  \tag{22.7.2.1}
\end{equation}
soit injectif.

Considérons d'autre part l'homomorphisme composé
\begin{equation}
\label{0.22.7.2.2-fr}
  j:(\mathfrak q/\mathfrak q^2)\otimes_AL
  \xrightarrow{j_0}\Omega^1_{A/k}\otimes_AL
  \longrightarrow\widehat\Omega^1_{A/k}\otimes_AL
  \tag{22.7.2.2}
\end{equation}
et montrons que la condition b) équivaut à dire que $j$ est injectif. Notons en
effet que
\[
  (\mathfrak q/\mathfrak q^2)\otimes_AL
  =(\mathfrak qA_{\mathfrak p}/\mathfrak q^2A_{\mathfrak p})
  \otimes_{A_{\mathfrak p}}L~;
\]
la condition b) signifie que si $F_i=j(f'_i\otimes1)$, où $f'_i$ est l'image
de $f_i$ dans
$\mathfrak qA_{\mathfrak p}/\mathfrak q^2A_{\mathfrak p}$, la matrice
$(\langle F_i,D_j\otimes1\rangle)$ est inversible~; donc cela entraîne que les
$F_i$ sont linéairement indépendants, et \emph{a fortiori} qu'il en est de
même des $f'_i\otimes1$~; mais comme ces derniers engendrent
$(\mathfrak qA_{\mathfrak p}/\mathfrak q^2A_{\mathfrak p})
\otimes_{A_{\mathfrak p}}L$, on en conclut que $j$ est injectif. Inversement,
supposons $j$ injectif, et prenons les $f_i$ tels que les $f'_i\otimes1$
forment

\oldpage[0\textsubscript{IV-1}]{210}

une base de
$(\mathfrak qA_{\mathfrak p}/\mathfrak q^2A_{\mathfrak p})
\otimes_{A_{\mathfrak p}}L$~; alors les $F_i$ forment une base de l'image de
$j$~; mais on sait (21.9.3) que $\widehat\Omega^1_{A/k}$ est un $A$-module
libre de rang $r$ et les $k$-dérivations de $A$ dans lui-même engendrent son
dual~; le fait que les $F_i$ soient linéairement indépendants entraîne donc
l'existence de $k$-dérivations $D_j$ de $A$ dans lui-même telles que la matrice
$(\langle F_i,D_j\otimes1\rangle)$ soit inversible, autrement dit la condition
b).

Ces remarques montrent déjà que b) entraîne a). Pour prouver inversement que
la condition a) entraîne que $j$ est injectif, considérons le diagramme
commutatif
\begin{equation}
\label{0.22.7.2.3-fr}
\xymatrix{
  (\mathfrak q/\mathfrak q^2)\otimes_AL
    \ar[r]^j \ar[d]_\alpha &
  \widehat\Omega^1_{A/k}\otimes_AL \ar@{=}[d] \\
  (\mathfrak p/\mathfrak p^2)\otimes_AL
    \ar[r]_i &
  \widehat\Omega^1_{A/k}\otimes_AL
}
\tag{22.7.2.3}
\end{equation}
et notons le lemme suivant~:
\end{proof}

\begin{lemma}[22.7.2.4]
\label{0.22.7.2.4-fr}
Soient $R$ un anneau local régulier, $\mathfrak m$ son idéal maximal,
$K=R/\mathfrak m$ son corps résiduel. Si $\mathfrak n$ est un idéal de $R$ tel
que $R/\mathfrak n$ soit régulier, alors l'homomorphisme canonique
\begin{equation}
\label{0.22.7.2.5-fr}
  \alpha:(\mathfrak n/\mathfrak n^2)\otimes_RK
  \longrightarrow(\mathfrak m/\mathfrak m^2)\otimes_RK
  \tag{22.7.2.5}
\end{equation}
est injectif.
\end{lemma}

\begin{proof}
En effet, on sait (17.1.6) que $\mathfrak n$ est engendré par une suite
$(x_i)_{1\leq i\leq t}$ faisant partie d'un système régulier de paramètres
$(x_i)_{1\leq i\leq n}$ de $R$~; les images des $x_i$ $(1\leq i\leq t)$
forment une base de $(\mathfrak n/\mathfrak n^2)\otimes_RK$, et leurs images
par $\alpha$ font partie de la base de
$(\mathfrak m/\mathfrak m^2)\otimes_RK$ formée des images des $x_i$ pour
$1\leq i\leq n$, d'où le lemme.
\end{proof}

\begin{proof}[Suite de la démonstration de la proposition~{\rm(22.7.2)}]
Ce lemme et le diagramme (22.7.2.3) ramènent donc (en vertu de l'hypothèse a))
à prouver que $i$ est injectif. Or, dans la suite
\begin{equation}
\label{0.22.7.2.6-fr}
  \mathfrak p/\mathfrak p^2\xrightarrow{h}
  \widehat\Omega^1_{A/k}\widehat\otimes_A(A/\mathfrak p)
  \xrightarrow{g}\widehat\Omega^1_{(A/\mathfrak p)/k}\longrightarrow0
  \tag{22.7.2.6}
\end{equation}
on sait (20.7.20) que $g$ est surjectif et que l'image de $h$ est dense dans
le noyau de $g$ pour la topologie $\mathfrak m$-adique ($\mathfrak m$ étant
l'idéal maximal de $A$)~; comme $\widehat\Omega^1_{A/k}$ est un $A$-module de
type fini, tous les sous-modules du $(A/\mathfrak p)$-module
\[
  \widehat\Omega^1_{A/k}\otimes_A(A/\mathfrak p)
  =\widehat\Omega^1_{A/k}\widehat\otimes_A(A/\mathfrak p)
\]
sont fermés pour la topologie $\mathfrak m$-adique ($0_{\mathrm I}$, 7.3.5),
donc la suite (22.7.2.6) est exacte. Tensorisant par $L$, il vient la suite
exacte
\[
  (\mathfrak p/\mathfrak p^2)\otimes_AL\xrightarrow{i}
  \widehat\Omega^1_{A/k}\otimes_AL\longrightarrow
  \widehat\Omega^1_{(A/\mathfrak p)/k}\otimes_{A/\mathfrak p}L
  \longrightarrow0.
\]

Or, l'hypothèse de séparabilité faite sur $L$ entraîne, en vertu de (21.9.5),
que $\widehat\Omega^1_{(A/\mathfrak p)/k}\otimes_{A/\mathfrak p}L$ est de
rang $s$ sur $L$~; comme $\widehat\Omega^1_{A/k}\otimes_AL$ est de rang $r$
sur $L$ (21.9.3), on a $\operatorname{rg}_L(\operatorname{Im}i)=r-s$. Mais comme
$A/\mathfrak p$ est finie sur une sous-algèbre isomorphe à
$k[[T_1,\ldots,T_s]]$, on a $\dim(A/\mathfrak p)=s$ ((17.1.4) et (16.1.5))~;
donc
\[
  r-s=\dim(A)-\dim(A/\mathfrak p)=\dim(A_{\mathfrak p})
\]

\oldpage[0\textsubscript{IV-1}]{211}

par (16.5.11). Or, comme $A_{\mathfrak p}$ est régulier (17.3.2),
$\dim(A_{\mathfrak p})$ est égal au rang sur $L$ de
$\mathfrak pA_{\mathfrak p}/\mathfrak p^2A_{\mathfrak p}$ (17.1.1), donc au
rang sur $L$ de $(\mathfrak p/\mathfrak p^2)\otimes_AL$, ce qui achève de
prouver que $i$ est injectif.
\end{proof}

\begin{theorem}[22.7.3]
\label{0.22.7.3-fr}
Soient $k$ un corps, $A$ un anneau local noethérien complet de corps résiduel
$K$. On suppose que~:
\begin{enumerate}
  \item[$1^\circ$] $[k:k^p]<+\infty$ (où $p$ est l'exposant caractéristique
  de $k$)~;
  \item[$2^\circ$] $K$ est une extension finie d'une extension séparable
  $K_0$ de $k$~;
  \item[$3^\circ$] $A$ est muni d'une structure de $K_0$-algèbre formellement
  lisse (pour la topologie préadique).
\end{enumerate}
Soient $\mathfrak q$ un idéal de $A$, $B=A/\mathfrak q$, $\mathfrak p$ un
idéal premier de $A$ contenant $\mathfrak q$. Les conditions suivantes sont
équivalentes~:
\begin{enumerate}
  \item[a)] $B_{\mathfrak p}$ est une $k$-algèbre formellement lisse (pour la
  topologie $\mathfrak p$-préadique).
  \item[b)] Il existe des $k$-dérivations $D_i$ de $A$ dans lui-même
  $(1\leq i\leq m)$ et des éléments $f_i$ $(1\leq i\leq m)$ de
  $\mathfrak q$, tels que les images des $f_i$ dans
  $\mathfrak qA_{\mathfrak p}$ engendrent cet idéal de $A_{\mathfrak p}$, et
  que l'on ait $\det(D_i(f_j))\notin\mathfrak p$.
  \item[b')] Il existe une sous-extension $k'$ de $K_0$, contenant $K_0^p$,
  telle que $[K_0:k']<+\infty$, des $k'$-dérivations $D_i$ de $A$ dans
  lui-même $(1\leq i\leq m)$ et des éléments $f_i$ de $\mathfrak q$
  $(1\leq i\leq m)$, tels que les images des $f_i$ dans
  $\mathfrak qA_{\mathfrak p}$ engendrent cet idéal de $A_{\mathfrak p}$, et
  que l'on ait $\det(D_i(f_j))\notin\mathfrak p$.
\end{enumerate}
\end{theorem}

\begin{proof}
Distinguons deux cas, suivant que $p=1$ ou $p>1$.

\emph{A)} $k$ \emph{est de caractéristique} $0$. Comme $K$ est alors une
extension séparable de $K_0$, il résulte de (19.6.4) que $A$ est
$K_0$-isomorphe à un anneau de séries formelles $K[[T_1,\ldots,T_r]]$ muni de
sa structure usuelle de $K$-algèbre. Mais alors, compte tenu de (19.8.8, (ii))
appliqué à $A/\mathfrak p$, on voit que les conditions générales de (22.7.2)
sont vérifiées en y remplaçant $k$ par $K$. En outre, en vertu de (19.6.4), il
revient au même de dire que $B_{\mathfrak p}$ est une $k$-algèbre
formellement lisse ou une $K$-algèbre formellement lisse, les deux conditions
étant équivalentes au fait que $B_{\mathfrak p}$ est un anneau régulier. On
peut donc appliquer les conclusions de (22.7.2), et il est immédiat que cela
prouve l'équivalence de a), b) et b') (avec $k'=K_0$ dans b')).

\emph{B)} $k$ \emph{est de caractéristique} $p>0$. Comme $A$ est une
$K_0$-algèbre formellement lisse et que $K_0$ est séparable sur $k$, il
résulte de (19.3.5, (ii)) et (19.6.1) que $A$ est une $k$-algèbre
formellement lisse~; en vertu de (22.5.9), $A_{\mathfrak p}$ est aussi une
$k$-algèbre formellement lisse pour la topologie $\mathfrak p$-préadique. Il
résulte donc encore de (22.6.2, (ii)) que la condition a) équivaut au fait que
l'homomorphisme $j_0$ défini dans (22.7.2.1) (où $L=k(\mathfrak p)$) est
injectif. Donc (19.1.12, c)) montre donc déjà que b) implique a)~; comme par
ailleurs b') implique trivialement b), tout revient à montrer que l'hypothèse
que $j_0$ est injectif entraîne b').

Désignons d'abord par $k'$ une sous-extension quelconque de $K_0$ telle que
$[K_0:k']<+\infty$. Notons que $\widehat\Omega^1_{A/k'}$ est un $A$-module
libre de type fini. Par le raisonnement de (22.7.2), il suffit de montrer que
(pour un choix convenable de $k'$), l'homomorphisme composé
\begin{equation}
\label{0.22.7.3.1-fr}
  j':(\mathfrak q/\mathfrak q^2)\otimes_AL
  \xrightarrow{j_0}\Omega^1_{A/k}\otimes_AL
  \longrightarrow\widehat\Omega^1_{A/k'}\otimes_AL
  \tag{22.7.3.1}
\end{equation}

\oldpage[0\textsubscript{IV-1}]{212}

est injectif. Considérons encore le diagramme commutatif
\begin{equation}
\label{0.22.7.3.2-fr}
\xymatrix{
  (\mathfrak q/\mathfrak q^2)\otimes_AL
    \ar[r]^{j_0} \ar[d]_\alpha &
  \Omega^1_{A/k}\otimes_AL \ar[r] \ar@{=}[d] &
  \widehat\Omega^1_{A/k'}\otimes_AL \ar@{=}[d] \\
  (\mathfrak p/\mathfrak p^2)\otimes_AL
    \ar[r]_{i_0} &
  \Omega^1_{A/k}\otimes_AL \ar[r] &
  \widehat\Omega^1_{A/k'}\otimes_AL
}
\tag{22.7.3.2}
\end{equation}

Comme $j_0$ est injectif, on a $\operatorname{Ker}(i_0\circ\alpha)=0$. Mais comme
$A_{\mathfrak p}$ est un anneau régulier et que l'hypothèse a) implique que
$B_{\mathfrak p}=A_{\mathfrak p}/\mathfrak qA_{\mathfrak p}$ est aussi un
anneau régulier (19.6.5), le lemme (22.7.2.4) montre que $\alpha$ est injectif,
d'où $\alpha^{-1}(\operatorname{Ker} i_0)=0$. Si $i'$ est le composé des homomorphismes de la
seconde ligne de (22.7.3.2), on a de même
$\operatorname{Ker}(j')=\alpha^{-1}(\operatorname{Ker}(i'))$, et tout revient donc à montrer que, pour un
choix convenable de $k'$, on a
\begin{equation}
\label{0.22.7.3.3-fr}
  \operatorname{Ker}(i')=\operatorname{Ker}(i).
  \tag{22.7.3.3}
\end{equation}

Mais le même raisonnement que dans (22.7.2) montre que l'on a une suite exacte
\[
  (\mathfrak p/\mathfrak p^2)\otimes_AL\xrightarrow{i'}
  \widehat\Omega^1_{A/k'}\otimes_AL\longrightarrow
  \widehat\Omega^1_{(A/\mathfrak p)/k'}\otimes_{A/\mathfrak p}L
  \longrightarrow0.
\]
Or, les hypothèses faites permettent d'appliquer (21.9.8), donc il existe une
sous-extension $k'$ de $K_0$ contenant $K_0^p$, telle que
$[K_0:k']<+\infty$, et pour laquelle on a
\[
  \operatorname{rg}_L\bigl(\widehat\Omega^1_{(A/\mathfrak p)/k'}
  \otimes_{A/\mathfrak p}L\bigr)
  =\dim(A/\mathfrak p)+\operatorname{rg}_L\Upsilon_{L/k}
  +\operatorname{rg}_{K_0}\Omega^1_{K_0/k'}.
\]
Mais on a
\[
  \dim(A/\mathfrak p)=\dim(A)-\dim(A_{\mathfrak p})
\]
par (16.5.11),
\[
  \dim(A_{\mathfrak p})
  =\operatorname{rg}_L\bigl((\mathfrak p/\mathfrak p^2)\otimes_AL\bigr)
\]
puisque $A_{\mathfrak p}$ est régulier (17.1.1), et enfin
\[
  \dim(A)+\operatorname{rg}_{K_0}\Omega^1_{K_0/k'}
  =\operatorname{rg}_L\bigl(\widehat\Omega^1_{A/k'}\otimes_AL\bigr)
\]
en vertu de (21.9.2)~; on obtient donc
\begin{equation}
\label{0.22.7.3.4-fr}
  \operatorname{rg}_L(\operatorname{Ker}(i'))=\operatorname{rg}_L\Upsilon_{L/k}.
  \tag{22.7.3.4}
\end{equation}

D'autre part, le corps résiduel de $A_{\mathfrak p}$ étant formellement lisse
sur le corps premier de $k$, on a une suite exacte (20.6.22.1)
\[
  \Upsilon_{L/k}\xrightarrow{\chi_{A_{\mathfrak p}/k}}
  (\mathfrak pA_{\mathfrak p})/(\mathfrak pA_{\mathfrak p})^2
  \xrightarrow{i}\Omega^1_{A_{\mathfrak p}/k}
  \otimes_{A_{\mathfrak p}}L,
\]
et comme $A_{\mathfrak p}$ est une $k$-algèbre formellement lisse pour la
topologie $\mathfrak p$-préadique, il suit de (22.2.2) que
$\chi_{A_{\mathfrak p}/k}$ est un homomorphisme injectif, donc $\operatorname{Ker}(i)$ est
isomorphe à $\Upsilon_{L/k}$, et compte tenu de (22.7.3.4) et de ce que
$\operatorname{Ker}(i)\subset\operatorname{Ker}(i')$, cela achève de prouver (22.7.3.3) et par suite le
théorème.
\end{proof}

\oldpage[0\textsubscript{IV-1}]{213}

\begin{remark}[22.7.4]
\label{0.22.7.4-fr}
On a en fait montré (en utilisant (21.9.6)) que la condition b') est encore
équivalente aux autres conditions de (22.7.3) lorsque l'on assujettit $k'$ à
être un des corps d'une famille filtrante décroissante $(k_\alpha)$ de
sous-corps de $K_0$ contenant $K_0^p$, tels que
$[K_0:k_\alpha]<+\infty$ pour tout $\alpha$ et que
$\bigcap_\alpha k_\alpha=K_0^p$.
\end{remark}

\begin{corollary}[22.7.5]
\label{0.22.7.5-fr}
Sous les hypothèses générales de (22.7.3), l'ensemble des idéaux premiers
$\mathfrak n\in\operatorname{Spec}(B)$ tels que $B_{\mathfrak n}$ soit une
$k$-algèbre formellement lisse (pour la topologie $\mathfrak n$-préadique)
est ouvert dans $\operatorname{Spec}(B)$.
\end{corollary}

\begin{proof}
Avec les notations de (22.7.3), il suffit de voir que si $B_{\mathfrak p}$
est formellement lisse sur $k$, il en est de même de $B_{\mathfrak p'}$, pour
tout $\mathfrak p'\in V(\mathfrak q)$ assez voisin de $\mathfrak p$ dans
$\operatorname{Spec}(A)$. Or, si les $k$-dérivations $D_i$ et les éléments
$f_i$ de $\mathfrak q$ vérifient le critère b) de (22.7.3), on a aussi
$f=\det(D_if_j)\notin\mathfrak p'$ pour $\mathfrak p'\in D(f)$~; d'autre
part, il existe un $g\notin\mathfrak p$ tel que les images des $f_i$ dans
$\mathfrak qA_{\mathfrak p'}$ engendrent $\mathfrak qA_{\mathfrak p'}$ pour
$\mathfrak p'\in D(g)$ (Bourbaki, \emph{Alg. comm.}, chap.~II, §~5,
n\textsuperscript{o}~1, prop.~2), ce qui achève de prouver le corollaire.
\end{proof}

\begin{corollary}[22.7.6]
\label{0.22.7.6-fr}
\textnormal{(Nagata).} Soit $B$ un anneau local noethérien complet contenant
un corps. Alors l'ensemble des $\mathfrak n\in\operatorname{Spec}(B)$ tels
que $B_{\mathfrak n}$ soit régulier est ouvert dans $\operatorname{Spec}(B)$.
\end{corollary}

\begin{proof}
En effet, $B$ est une algèbre sur un corps premier $k$, donc (19.8.8, (ii))
de la forme $A/\mathfrak q$, où
$A=K_0[[T_1,\ldots,T_n]]$ est un anneau de séries formelles et $K_0$ une
extension de $k$~; d'autre part, dire que $B_{\mathfrak n}$ est régulier
équivaut à dire que c'est une $k$-algèbre formellement lisse pour la topologie
$\mathfrak n$-préadique (19.6.4). Toutes les conditions d'application de
(22.7.5) sont donc réalisées.
\end{proof}

\begin{remarks}[22.7.7]
\label{0.22.7.7-fr}
\begin{enumerate}
  \item[(i)] Nous verrons plus loin, avec Nagata, en utilisant (22.7.6), que
  l'on peut étendre la conclusion de (22.7.6) au cas d'un anneau local
  noethérien complet quelconque (IV, 6.12.7).

  \item[(ii)] La conclusion du théorème (22.7.3) n'est plus nécessairement
  exacte lorsqu'on ne suppose plus que $[k:k^p]$ soit fini. Prenons par
  exemple $k=\mathbf F_p(X_n)_{n\geqslant0}$,
  $A=k[[T,U]]$, $\mathfrak q$ égal à l'idéal principal $Af$, où
  \[
    f=U^p-\sum_{n=0}^{\infty}X_nT^{np}~;
  \]
  $A$ est un anneau factoriel dans lequel $f$ est un élément extrémal, car il
  est extrémal dans l'anneau de polynômes $k((T))[U]$, donc aussi dans
  l'anneau de séries formelles $k((T))[[U]]$ (Bourbaki,
  \emph{Alg. comm.}, chap.~VII, §~3), et \emph{a fortiori} dans $A$.
  Prenons $\mathfrak p=\mathfrak q$, de sorte que $B_{\mathfrak p}$ est le
  corps désigné par $E$ dans (21.9.7), qui est séparable sur $k$
  (\emph{loc. cit.}), donc une $k$-algèbre formellement lisse. Mais on a ici,
  avec les notations de (22.7.3), $k=K_0=K$, et les $k$-dérivations de $A$
  dans lui-même sont les combinaisons linéaires de $\partial/\partial T$ et
  $\partial/\partial U$~; comme on a
  $\partial f/\partial T=\partial f/\partial U=0$, les critères b) et b') de
  (22.7.3) ne sont pas vérifiés.
\end{enumerate}
\end{remarks}

% FIN SOURCE ega4/ega0-22-fr.tex

%
% DEBUT SOURCE ega4/ega0-23-fr.tex
\section{Anneaux japonais}
\label{section:0.23-fr}

Les résultats de ce paragraphe seront complétés dans IV, (7.6) et (7.7).

\subsection{Anneaux japonais}
\label{subsection:0.23.1-fr}

\begin{definition}[23.1.1]
\label{0.23.1.1-fr}
On dit qu'un anneau intègre $A$ est un anneau japonais si, pour toute
extension finie $K'$ de son corps des fractions $K$, la fermeture intégrale
$A'$ de $A$ dans $K'$ est

\oldpage[0\textsubscript{IV-1}]{214}

un $A$-module de type fini (autrement dit, une $A$-algèbre finie). On dit
qu'un anneau $A$ est universellement japonais si toute $A$-algèbre intègre de
type fini est un anneau japonais.
\end{definition}

Il est clair que si $A$ est un anneau japonais, tout anneau de fractions
$S^{-1}A$ est alors un anneau japonais, car (avec les notations précédentes)
$S^{-1}A'$ est la fermeture intégrale de $S^{-1}A$ dans $K'$, et est
évidemment un $S^{-1}A$-module de type fini.

Un anneau intègre noethérien, même un \emph{anneau de valuation discrète},
n'est pas nécessairement un anneau japonais [30].

\begin{proposition}[23.1.2]
\label{0.23.1.2-fr}
Soient $A$ un anneau noethérien intègre, $K$ son corps des fractions. Si, pour
toute extension radicielle finie $K'$ de $K$, la fermeture intégrale de $A$
dans $K'$ est un $A$-module de type fini, alors $A$ est un anneau japonais.
\end{proposition}

\begin{proof}
Comme $A$ est noethérien, il suffit, pour vérifier que $A$ est un anneau
japonais, de voir que pour toute extension quasi-galoisienne finie $L$ de
$K$, la fermeture intégrale $B$ de $A$ dans $L$ est un $A$-module de type
fini. Or $L$ est une extension galoisienne de la plus grande extension
radicielle $K'$ de $K$ contenue dans $L$~; et si $A'$ est la fermeture
intégrale de $A$ dans $K'$, $B$ est la fermeture intégrale de $A'$ dans $L$.
Mais on sait que dans une extension séparable de $K'$, la fermeture intégrale
de $A'$ est un $A'$-module de type fini (Bourbaki, \emph{Alg. comm.},
chap.~V, §~1, n\textsuperscript{o}~6, cor.~1 de la prop.~20), d'où la
proposition.
\end{proof}

Il résulte de (23.1.2) que lorsque $K$ est de caractéristique $0$, dire que
$A$ est un anneau japonais signifie que sa \emph{clôture intégrale} $A'$ est
un $A$-module de type fini.

\begin{theorem}[23.1.3]
\label{0.23.1.3-fr}
\textnormal{(Tate).} Soient $A$ un anneau intègre noethérien, $x\neq0$ un
élément de $A$. Supposons vérifiées les conditions suivantes~:
\begin{enumerate}
  \item[(i)] $A$ est intégralement clos.
  \item[(ii)] $\mathfrak p=xA$ est premier et $A$ est séparé et complet pour
  la topologie $\mathfrak p$-préadique.
  \item[(iii)] $A/xA$ est un anneau japonais.
\end{enumerate}
Alors $A$ est un anneau japonais.
\end{theorem}

Nous utiliserons le lemme suivant~:

\begin{lemma}[23.1.3.1]
\label{0.23.1.3.1-fr}
Soient $A$ un anneau, $x$ un élément de $A$ non diviseur de $0$ dans $A$ et
tel que $xA=\mathfrak p$ soit premier~; alors, pour tout entier $n>0$,
l'image réciproque de $x^nA_{\mathfrak p}$ dans $A$ est $x^nA$.
\end{lemma}

\begin{proof}
En effet, supposons que $b\in A$ soit un élément tel que
$b/1=x^na/s$ dans $A_{\mathfrak p}$, où $a\in A$ et
$s\notin\mathfrak p$~; il existe donc $s'\notin\mathfrak p$ tel que
$s'sb=x^nas'$, d'où $s'sb\in\mathfrak p$, et comme
$s's\notin\mathfrak p$, cela entraîne $b\in\mathfrak p$, autrement dit
$b=xb'$ avec $b'\in A$~; puisque $x$ est non diviseur de $0$, on en conclut
$s'sb'=x^{n-1}as$ et il suffit de raisonner par récurrence sur $n$.
\end{proof}

\begin{proof}
Pour démontrer le théorème, on peut se borner au cas où le corps des fractions
$K$ de $A$ est de caractéristique $p>0$, puisque $A$ est intégralement clos.
Soit $K'$ une extension radicielle finie de $K$, de sorte qu'il existe une
puissance $q=p^e$ telle que $K'^q\subset K$~; en remplaçant $K'$ par une
extension radicielle plus grande, on peut même supposer qu'il existe $y\in K'$
tel que $y^q=x$. Soit $A'$ la fermeture intégrale de $A$ dans $K'$~; comme
$A$ est intégralement clos, $A'$ est l'ensemble des $x'\in K'$ tels que
$x'^q\in A'\cap K=A$. Posons $V=A_{\mathfrak p}$~; l'idéal maximal
$\mathfrak m=xV$ de $V$ étant principal, on sait que l'anneau local
noethérien $V$ est un \emph{anneau de valuation discrète} (17.1.4)~; comme
$V$ est intégralement clos, le même

\oldpage[0\textsubscript{IV-1}]{215}

raisonnement que ci-dessus montre que la fermeture intégrale $V'$ de $V$ dans
$K'$ est l'ensemble des $x'\in K'$ tels que $x'^q\in V$~; on sait
(Bourbaki, \emph{Alg. comm.}, chap.~VI, §~8, n\textsuperscript{o}~6,
prop.~6, et chap.~V, §~2, n\textsuperscript{o}~3, lemme~4) que $V'$ est un
\emph{anneau de valuation discrète}, dont l'idéal maximal $\mathfrak m'$ est
l'ensemble des $x'\in K'$ tels que $x'^q\in\mathfrak m$ et en outre le corps
résiduel $V'/\mathfrak m'$ est une extension \emph{finie} de
$V/\mathfrak m$ (\emph{loc. cit.}, chap.~VI, §~8,
n\textsuperscript{o}~1, lemme~2). Montrons que, pour tout entier $n>0$, on a
\begin{equation}
\label{0.23.1.3.2-fr}
  \mathfrak m'^n\cap A'=y^nA'.
  \tag{23.1.3.2}
\end{equation}

En effet, comme $y^q=x$, on a $y\in\mathfrak m'$, donc
$y^nA'\subset\mathfrak m'^n\cap A'$. Inversement, soit
$x'\in\mathfrak m'^n\cap A'$, et posons $x'=z'y^n$ avec $z'\in K'$. On a
$x'^q\in A'$~; d'autre part, $x'$ est somme de produits
$t'_1\cdots t'_n$ avec $t'_i\in\mathfrak m'$, donc
$t_i'^q\in\mathfrak m$, et l'on en conclut que
$x'^q\in\mathfrak m^n\cap A$. Or le lemme (23.1.3.1) prouve que
$\mathfrak m^n\cap A=x^nA$, donc on a $z'^qx^n\in x^nA$, ou encore
$z'^q\in A$ puisque $x\neq0$, ce qui établit (23.1.3.2).

Montrons en second lieu que sur $A'$ la topologie $xA'$-préadique est
\emph{séparée}~; cette topologie est aussi la topologie $yA'$-préadique de
$A'$, et il résulte de (23.1.3.2) que cette topologie est induite par la
topologie $\mathfrak m'$-préadique de $V'$ qui est séparée puisque $V'$ est
un anneau de valuation discrète.

Prouvons ensuite que $A'/xA'$ est un $A$-module de type fini. Comme $x=y^q$
et que $y^kA'/y^{k+1}A'$ est isomorphe à $A'/yA'$, on peut se borner à
montrer que $A'/yA'$ est un $A$-module de type fini~; mais la formule
(23.1.3.2) appliquée pour $n=1$ montre que $A'/yA'$ est un sous-anneau de
$V'/\mathfrak m'$, c'est-à-dire du corps résiduel de $V'$, qui est une
extension finie du corps résiduel $V/\mathfrak m$ de $V$. Or
$V/\mathfrak m$ est le corps des fractions de $A/\mathfrak p$, et comme
$A'$ est entier sur $A$, $A'/yA'$ est entier sur $A/\mathfrak p$, donc
contenu dans la fermeture intégrale de $A/\mathfrak p$ dans
$V'/\mathfrak m'$~; comme par hypothèse $A/\mathfrak p$ est un anneau
japonais noethérien, $A'/yA'$ est un $(A/\mathfrak p)$-module de type fini,
et \emph{a fortiori} un $A$-module de type fini.

Cela étant, la topologie $xA'$-préadique de $A'$ étant séparée, $A'$ est un
sous-anneau de son complété $\widehat{A'}$ pour cette topologie~; mais puisque
$A'/xA'$ est un $A$-module de type fini et que $A$ est séparé et complet pour
la topologie $xA$-préadique, $\widehat{A'}$ est un $A$-module de type fini en
vertu de (0\textsubscript{I}, 7.2.9), donc il en est de même de $A'$, C.Q.F.D.
\end{proof}

\begin{corollary}[23.1.4]
\label{0.23.1.4-fr}
Soit $A$ un anneau japonais noethérien et intégralement clos. Alors tout
anneau de séries formelles $A[[T_1,\ldots,T_r]]$ est un anneau japonais.
\end{corollary}

\begin{proof}
On sait que pour tout anneau noethérien intégralement clos $B$, l'anneau de
séries formelles $B[[T]]$ est noethérien et intégralement clos (Bourbaki,
\emph{Alg. comm.}, chap.~V, §~1, n\textsuperscript{o}~4, prop.~14)~; on peut
donc, par récurrence sur $n$, se borner à prouver que $A[[T]]$ est un anneau
japonais. Or l'élément $x=T$ vérifie toutes les conditions de (23.1.3), d'où
la conclusion.
\end{proof}

\begin{theorem}[23.1.5]
\label{0.23.1.5-fr}
\textnormal{(Nagata).} Tout anneau local noethérien intègre complet $A$ est
un anneau japonais.
\end{theorem}

\begin{proof}
On sait (19.8.8, (ii)) qu'il existe un sous-anneau $B$ de $A$ qui est un
anneau local complet et régulier, tel que $A$ soit une $B$-algèbre finie~; il
suffit évidemment de prouver

\oldpage[0\textsubscript{IV-1}]{216}

que $B$ est un anneau japonais, autrement dit, on peut se borner au cas où
$A$ est en outre \emph{régulier}. Raisonnons par récurrence sur
$n=\dim(A)$, le théorème étant évident pour $n=0$. Supposons donc $n>0$, et,
en désignant par $\mathfrak m$ l'idéal maximal de $A$, soit $x$ un élément de
$\mathfrak m-\mathfrak m^2$~; on sait (17.1.8) que $A/xA$ est un anneau
régulier, et en outre ((17.1.7) et (16.3.4)) que
$\dim(A/xA)=n-1$~; de plus, l'idéal $xA$ est fermé dans $A$
(0\textsubscript{I}, 7.3.5), donc $A/xA$ est complet. En vertu de l'hypothèse
de récurrence, $A/xA$ est donc un anneau japonais, et il résulte alors de
(23.1.3) qu'il en est de même de $A$.
\end{proof}

\begin{corollary}[23.1.6]
\label{0.23.1.6-fr}
Soient $A$ un anneau local noethérien intègre complet, $K$ son corps des
fractions, $K'$ une extension finie de $K$~; alors la fermeture intégrale
$A'$ de $A$ dans $K'$ est un anneau local complet.
\end{corollary}

\begin{proof}
On sait déjà par (23.1.5) que $A'$ est une $A$-algèbre finie, donc un anneau
semi-local noethérien complet, et par suite un \emph{composé direct}
d'anneaux locaux~; mais comme $A'$ est intègre, c'est nécessairement un
anneau local.
\end{proof}

\begin{proposition}[23.1.7]
\label{0.23.1.7-fr}
Soient $A$ un anneau local noethérien intègre, $K$ son corps des fractions,
$K'$ une extension finie de $K$, $A'$ la fermeture intégrale de $A$ dans
$K'$. On suppose que le complété $\widehat A$ est réduit, et l'on désigne par
$R$ son anneau total des fractions.
\begin{enumerate}
  \item[(i)] Si l'anneau $R\otimes_KK'$ est réduit (ce qui aura lieu en
  particulier lorsque $K'$ est une extension séparable de $K$, $R$ étant
  composé direct de corps, extensions de $K$ (Bourbaki, \emph{Alg.},
  chap.~VIII, §~7, n\textsuperscript{o}~3, cor.~1 du th.~1)), alors $A'$ est
  un $A$-module de type fini.
  \item[(ii)] Si en particulier $R$ est une $K$-algèbre séparable, alors $A$
  est un anneau japonais.
\end{enumerate}
\end{proposition}

\begin{proof}
\emph{(i)} Posons pour simplifier $A_1=\widehat A$,
$A'_1=A'\otimes_AA_1$, $K'_1=K'\otimes_AA_1$. Comme $A_1$ est un
$A$-module plat (0\textsubscript{I}, 7.3.5), $A'_1$ s'identifie à un
sous-anneau de $K'_1$ et est évidemment entier sur $A_1$. D'autre part, si
l'on pose $K_1=K\otimes_AA_1$, on peut écrire
$K'_1=K'\otimes_KK_1$~; comme $A$ est intègre et que $A_1$ est un $A$-module
plat, tout élément $\neq0$ de $A$ est $A_1$-régulier
(0\textsubscript{I}, 6.3.4)~; comme $K=S^{-1}A$, où $S=A-\{0\}$, on a
$K_1=S^{-1}A_1$, et la remarque précédente prouve que $K_1$ s'identifie à un
sous-anneau de l'anneau total des fractions $R$ de $A_1$. Puisque tout
$K$-module est plat, $K'_1$ s'identifie à un sous-anneau de $K'\otimes_KR$,
qui par hypothèse est réduit et est une $R$-algèbre finie. Désignons par
$\mathfrak q_i$ ($1\leqslant i\leqslant n$) les idéaux premiers minimaux de
$A_1$, par $B_i$ l'anneau intègre $A_1/\mathfrak q_i$, par $L_i$ le corps des
fractions de $B_i$, de sorte que $R$ est composé direct des $L_i$, et
$K'\otimes_KR$ composé direct des $K'\otimes_KL_i$~; ces dernières
$K$-algèbres, étant réduites, sont des composées directes de corps,
extensions finies des $L_i$. Comme $B_i$ est un anneau local intègre
\emph{complet}, le théorème de Nagata (23.1.5) prouve que la fermeture
intégrale de $B_i$ dans $K'\otimes_KL_i$ est un $B_i$-module de type fini, et
\emph{a fortiori} un $A_1$-module de type fini~; comme tout élément de
$K'\otimes_KL_i$ qui est entier sur $A_1$ est aussi entier sur $B_i$ (étant
annulé par $\mathfrak q_i$), on en conclut finalement que la fermeture
intégrale de $A_1$ dans $K'\otimes_KR$ est un $A_1$-module de type fini. Mais
comme $A'_1$ est contenu dans cette fermeture intégrale et que $A_1$ est
noethérien, $A'_1$ est aussi un $A_1$-module de type fini. Enfin, comme $A_1$
est un $A$-module fidèlement plat (0\textsubscript{I}, 7.3.5), on en conclut
que $A'$ est un $A$-module de type fini (Bourbaki, \emph{Alg. comm.},
chap.~I, §~3, n\textsuperscript{o}~6, prop.~11).

\oldpage[0\textsubscript{IV-1}]{217}

\emph{(ii)} L'hypothèse entraîne que les $L_i$ sont des extensions séparables
de $K$, donc $K'\otimes_KL_i$ est réduit (Bourbaki, \emph{Alg.}, chap.~VIII,
§~7, n\textsuperscript{o}~3, th.~1) et l'on peut appliquer (i) à toute
extension finie $K'$ de $K$, ce qui démontre notre assertion.
\end{proof}

\subsection{Clôture intégrale d'un anneau local noethérien intègre}
\label{subsection:0.23.2-fr}

\begin{env}[23.2.1]
\label{0.23.2.1-fr}
Dans ce qui suit, pour tout anneau $A$, nous noterons $A_{\mathrm{red}}$ pour
simplifier, le quotient de $A$ par son nilradical (de sorte que si
$X=\operatorname{Spec}(A)$, on a
$X_{\mathrm{red}}=\operatorname{Spec}(A_{\mathrm{red}})$).

Nous dirons qu'un anneau \emph{local} $A$ est \emph{unibranche} si l'anneau
$A_{\mathrm{red}}$ est intègre et si la clôture intégrale de
$A_{\mathrm{red}}$ est un anneau \emph{local}, ce qui généralise la définition
donnée dans (III, 4.3.6). Nous dirons que $A$ est \emph{géométriquement
unibranche} s'il est unibranche et si le corps résiduel de l'anneau local,
clôture intégrale de $A_{\mathrm{red}}$, est une extension \emph{radicielle}
de celui de $A$. Il est clair qu'un anneau local intégralement clos est
géométriquement unibranche~; il résulte de (23.1.6) qu'un anneau local
noethérien intègre \emph{complet} est unibranche.
\end{env}

\begin{lemma}[23.2.2]
\label{0.23.2.2-fr}
\textnormal{(*)}\footnote{Dans la suite de ce numéro, nous utilisons des
notions et résultats exposés au chap.~IV, §§~2 et 5~; comme les résultats de
ce numéro ne sont pas utilisés avant le chap.~IV, §~6, cela n'entraîne pas de
cercle vicieux.} Soient $A$ un anneau intègre, $A'$ sa clôture intégrale~; pour
que le morphisme canonique
$f:\operatorname{Spec}(A')\to\operatorname{Spec}(A)$ soit radiciel, il faut et
il suffit que pour tout $\mathfrak p\in\operatorname{Spec}(A)$,
$A_{\mathfrak p}$ soit géométriquement unibranche~; lorsqu'il en est ainsi,
$f$ est un homéomorphisme.
\end{lemma}

\begin{proof}
En effet, pour tout $\mathfrak p\in\operatorname{Spec}(A)$, la clôture
intégrale de $A_{\mathfrak p}$ est $A'_{\mathfrak p}$ (Bourbaki,
\emph{Alg. comm.}, chap.~V, §~1, n\textsuperscript{o}~5, prop.~16), et tous
les idéaux premiers de $A'_{\mathfrak p}$ au-dessus de
$\mathfrak pA_{\mathfrak p}$ sont maximaux (\emph{loc. cit.}, §~2,
n\textsuperscript{o}~1, prop.~1). Dire que $f$ est injectif signifie donc que
pour tout $\mathfrak p\in\operatorname{Spec}(A)$, $A'_{\mathfrak p}$ est un
anneau \emph{local}, c'est-à-dire que les $A_{\mathfrak p}$ sont
unibranches. Dire que tout $A_{\mathfrak p}$ est géométriquement unibranche
signifie alors que $f$ est radiciel en vertu de (I, 3.5.8). Lorsqu'il en est
ainsi, $f$ est surjectif et fermé (II, 6.1.10), donc un homéomorphisme.
\end{proof}

\begin{lemma}[23.2.3]
\label{0.23.2.3-fr}
Soient $A$ un anneau intègre, $K$ son corps des fractions, $B$ un anneau
noethérien, $\varphi:A\to B$ un homomorphisme d'anneau faisant de $B$ un
$A$-module plat. Soit $K'$ une extension de $K$~; posons
$B_1=B_{\mathrm{red}}$, et soit $R$ l'anneau total des fractions de $B_1$.
Alors, pour toute sous-$A$-algèbre $A'$ de $K'$, l'homomorphisme canonique
\begin{equation}
\label{0.23.1.8.1-fr}
  (A'\otimes_AB)_{\mathrm{red}}
  \longrightarrow(K'\otimes_AR)_{\mathrm{red}}
  \tag{23.1.8.1}
\end{equation}
est injectif.
\end{lemma}

\begin{proof}
On peut en effet considérer l'homomorphisme canonique
$h:A'\otimes_AB\to K'\otimes_AR$ comme composé des homomorphismes suivants
\[
  A'\otimes_AB\xrightarrow{u}K'\otimes_AB
  \xrightarrow{v}K'\otimes_AB_1\xrightarrow{w}K'\otimes_AR.
\]
Comme $B$ est un $A$-module plat, $u$ est injectif~; de même, $K'$ étant un
$K$-module plat et $K$ un $A$-module plat, $K'$ est un $A$-module plat, et
$w$ est donc injectif puisque $B_1$ est réduit, donc l'homomorphisme
$B_1\to R$ injectif. Enfin, pour la même raison, si $\mathfrak N$ est

\oldpage[0\textsubscript{IV-1}]{218}

le nilradical de $B$, le noyau de $v$ est $K'\otimes_A\mathfrak N$, donc il est
contenu dans le nilradical de $K'\otimes_AB$~; on en conclut aussitôt que si
l'image par $h=w\circ v\circ u$ d'un élément $x\in A'\otimes_AB$ est
nilpotent, $x$ est nilpotent, ce qui prouve le lemme.
\end{proof}

\begin{proposition}[23.2.4]
\label{0.23.2.4-fr}
Soient $A$ un anneau intègre, $K$ son corps des fractions, $B$ un anneau
noethérien, $\mathfrak q_i$ ($1\leqslant i\leqslant n$) ses idéaux premiers
minimaux, $\varphi:A\to B$ un homomorphisme d'anneaux. On suppose que les
$B/\mathfrak q_i$ sont des anneaux japonais, et que $B$ soit un $A$-module
plat. Soit $K'$ une extension finie de $K$, et soit $(C_\lambda)$ la famille
filtrante croissante des sous-anneaux de $K'$ qui sont des $A$-algèbres
finies et admettent $K'$ pour corps des fractions~; la réunion $A'$ des
$C_\lambda$ est donc la fermeture intégrale de $A$ dans $K'$. Alors~:
\begin{enumerate}
  \item[(i)] Il existe un indice $\alpha$ tel que pour $\lambda\geqslant
  \alpha$, l'homomorphisme canonique
  \begin{equation}
  \label{0.23.2.4.1-fr}
    (C_\alpha\otimes_AB)_{\mathrm{red}}
    \longrightarrow(C_\lambda\otimes_AB)_{\mathrm{red}}
    \tag{23.2.4.1}
  \end{equation}
  soit bijectif.
  \item[(ii)] Si en outre $B$ est un $A$-module fidèlement plat, le morphisme
  $\operatorname{Spec}(A')\to\operatorname{Spec}(C_\alpha)$ est radiciel.
\end{enumerate}
\end{proposition}

\begin{proof}
\emph{(i)} Soit $B_1=B_{\mathrm{red}}$, et soit $R$ l'anneau total des
fractions de $B_1$~; comme $A$ est intègre et $B$ un $A$-module plat, tout
élément $\neq0$ de $A$ est $B$-régulier (0\textsubscript{I}, 6.3.4), donc
l'homomorphisme composé $A\xrightarrow{\varphi}B\to B_1$ se prolonge en un
homomorphisme $K\to R$~; on peut alors écrire
$K'\otimes_AR=(K'\otimes_AK)\otimes_KR$, et comme $K'\otimes_AK$
s'identifie à $K'$, on a $K'\otimes_AR=K'\otimes_KR$. Comme $R$ est composé
direct des corps des fractions $L_i$ des $B/\mathfrak q_i$,
$K'\otimes_AR$ est composé direct des $K'\otimes_KL_i$, et par suite
$(K'\otimes_AR)_{\mathrm{red}}$ est composé direct d'un nombre fini
d'extensions finies des $L_i$. En vertu de l'hypothèse, la fermeture
intégrale de $B/\mathfrak q_i$ dans une extension finie de $L_i$ est un
$(B/\mathfrak q_i)$-module de type fini, donc un $B$-module de type fini~;
on en conclut que la fermeture intégrale $B'$ de $B$ dans
$(K'\otimes_AR)_{\mathrm{red}}$ est un $B$-module de type fini. En vertu de
(23.2.3), les $(C_\lambda\otimes_AB)_{\mathrm{red}}$ s'identifient
canoniquement à des sous-anneaux de $(K'\otimes_AR)_{\mathrm{red}}$ qui sont
des $B$-algèbres finies, donc contenues dans $B'$. Comme $B$ est noethérien
et $B'$ un $B$-module de type fini, la famille filtrante des
$(C_\lambda\otimes_AB)_{\mathrm{red}}$ admet un \emph{plus grand élément}
$(C_\alpha\otimes_AB)_{\mathrm{red}}$, ce qui prouve (i).

\emph{(ii)} Comme $B$ est un $A$-module fidèlement plat, il suffit
(IV, 2.6.1, (v)) de montrer que le morphisme
$\operatorname{Spec}(A'\otimes_AB)\to
\operatorname{Spec}(C_\alpha\otimes_AB)$ est radiciel, ou, ce qui revient au
même (I, 5.1.6), que le morphisme
$\operatorname{Spec}(A'\otimes_AB)_{\mathrm{red}}\to
\operatorname{Spec}(C_\alpha\otimes_AB)_{\mathrm{red}}$ est radiciel~; mais
on a $A'\otimes_AB=\varinjlim_\lambda(C_\lambda\otimes_AB)$, donc
(IV, 5.13.2)
$(A'\otimes_AB)_{\mathrm{red}}=
\varinjlim_\lambda(C_\lambda\otimes_AB)_{\mathrm{red}}$, et la conclusion
résulte de (i).
\end{proof}

\begin{corollary}[23.2.5]
\label{0.23.2.5-fr}
Soient $A$ un anneau local noethérien intègre, $K$ son corps des fractions,
$K'$ une extension finie de $K$. Soit $(C_\lambda)$ la famille filtrante
croissante des sous-anneaux de $K'$ qui sont des $A$-algèbres finies (donc
des anneaux noethériens semi-locaux (Bourbaki, \emph{Alg. comm.}, chap.~IV,
§~2, n\textsuperscript{o}~5, cor.~3 de la prop.~9)) et admettent $K'$ pour
corps des fractions. Alors il existe $\alpha$ tel que l'homomorphisme
$(\widehat C_\alpha)_{\mathrm{red}}\to
(\widehat C_\lambda)_{\mathrm{red}}$ soit un isomorphisme pour
$\lambda\geqslant\alpha$, et si $A'$ est la fermeture intégrale de $A$ dans
$K'$, le morphisme
$\operatorname{Spec}(A')\to\operatorname{Spec}(C_\alpha)$ est radiciel
(cf. (23.2.2)).
\end{corollary}

\oldpage[0\textsubscript{IV-1}]{219}

\begin{proof}
On applique (23.2.4) en prenant $B=\widehat A$~; comme les
$B/\mathfrak q_i$ sont des anneaux locaux noethériens complets, ce sont des
anneaux japonais (23.1.5) et $B$ est un $A$-module fidèlement plat
(0\textsubscript{I}, 7.3.5)~; en outre on a
$C_\lambda\otimes_AB=\widehat C_\lambda$ (Bourbaki, \emph{Alg. comm.},
chap.~III, §~3, n\textsuperscript{o}~4, th.~3 et chap.~IV, §~2,
n\textsuperscript{o}~5, cor.~3 de la prop.~9).
\end{proof}

\begin{corollary}[23.2.6]
\label{0.23.2.6-fr}
Sous les hypothèses de (23.2.5), la fermeture intégrale $A'$ de $A$ dans
$K'$ est un anneau semi-local~; si $\mathfrak m$ est l'idéal maximal de $A$,
alors, pour tout idéal maximal $\mathfrak m'$ de $A'$,
$A'/\mathfrak m'$ est une extension finie de $A/\mathfrak m$.
\end{corollary}

\begin{proof}
Le fait que l'homomorphisme
$(\widehat C_\alpha)_{\mathrm{red}}\to
(\widehat C_\lambda)_{\mathrm{red}}$ soit bijectif pour
$\lambda\geqslant\alpha$ entraîne que le nombre des idéaux maximaux de
$C_\lambda$ est constant pour $\lambda\geqslant\alpha$ et que si
$\mathfrak m_\alpha$ est un idéal maximal de $C_\alpha$ et
$\mathfrak m_\lambda$ l'unique idéal maximal de $C_\lambda$ au-dessus de
$\mathfrak m_\alpha$, les corps $C_\alpha/\mathfrak m_\alpha$ et
$C_\lambda/\mathfrak m_\lambda$ sont canoniquement isomorphes. La conclusion
résulte de ce que $\mathfrak m'=\varinjlim\mathfrak m_\lambda$ et
$A'/\mathfrak m'=\varinjlim(C_\lambda/\mathfrak m_\lambda)$
(0\textsubscript{III}, 10.3.1.1.3).
\end{proof}

\begin{proposition}[23.2.7]
\label{0.23.2.7-fr}
\textnormal{(Y. Mori).} Soient $A$ un anneau local noethérien intègre, $K$ son
corps des fractions, $A'$ la clôture intégrale de $A$. Alors $A'$ est un
anneau de Krull (Bourbaki, \emph{Alg. comm.}, chap.~VII, §~1) semi-local~;
autrement dit, il existe une famille $(V_\lambda)$ d'anneaux de valuation
discrète ayant $K$ pour corps des fractions, tels que
$A'=\bigcap_\lambda V_\lambda$ et que tout $x\in K$ appartienne à tous les
$V_\lambda$ sauf pour un nombre fini d'entre eux. En outre, il existe un
sous-anneau $C$ de $K$ qui est une $A$-algèbre finie, et telle que les
$V_\lambda$ soient les clôtures intégrales des anneaux
$C_{\mathfrak p_\lambda}$, où $(\mathfrak p_\lambda)$ est la famille des
idéaux premiers de hauteur $1$ de l'anneau local $C$.
\end{proposition}

Prouvons d'abord le lemme suivant~:

\begin{lemma}[23.2.7.1]
\label{0.23.2.7.1-fr}
Soient $A$, $B$ deux anneaux, $\varphi:A\to B$ un homomorphisme faisant de $B$
un $A$-module fidèlement plat. On suppose que $A$ est intègre~; alors, si
$C=B_{\mathrm{red}}$, l'homomorphisme composé
$\psi:A\to B\to B_{\mathrm{red}}=C$ est injectif et se prolonge en un
homomorphisme injectif du corps des fractions $K$ de $A$ dans l'anneau total
des fractions $R$ de $C$~; en outre, si $C'$ est la fermeture intégrale de
$C$ dans $R$, $C'\cap K$ est la clôture intégrale de $A$.
\end{lemma}

\begin{proof}
Comme $A$ est intègre et $\varphi$ injectif, l'intersection de $\varphi(A)$
et du nilradical $\mathfrak N$ de $B$ est réduite à $0$, et comme
$C=B/\mathfrak N$, $\psi$ est injectif. Tout $a\neq0$ dans $A$ étant non
diviseur de zéro dans $A$ ne l'est pas non plus dans $B$ par platitude
(0\textsubscript{I}, 6.3.4)~; on en déduit que $a$ n'est pas non plus diviseur
de $0$ dans $C$, car si l'on avait $ax\in\mathfrak N$ pour un
$x\notin\mathfrak N$ dans $B$, on en tirerait $a^nx^n=0$ pour un entier $n$,
ce qui contredit ce qui précède puisque $x^n\neq0$. On peut donc prolonger
$\psi$ en un homomorphisme injectif de $K$ dans $R$. Pour prouver la dernière
assertion, notons qu'il est clair que la clôture intégrale $A'$ de $A$ est
contenue dans $C'\cap K$. Inversement, soit $x\in C'\cap K$~; $x$ est donc
entier sur $C$, et \emph{a fortiori} sur $B$, autrement dit $B[x]$ est une
$B$-algèbre finie~; d'ailleurs, $K$ s'identifie à un sous-anneau de
$B\otimes_AK$, et le sous-anneau $B[x]$ de $B\otimes_AK$ s'identifie à
$B\otimes_AA[x]$ par platitude~; on en conclut que $A[x]$ est un $A$-module
de type fini (Bourbaki, \emph{Alg. comm.}, chap.~I\textsuperscript{er}, §~3,
n\textsuperscript{o}~6, prop.~11), donc que $x\in A'$.
\end{proof}

Ce lemme étant établi, nous allons l'appliquer à l'injection canonique de $A$
dans son complété $\widehat A$, qui est un $A$-module fidèlement plat~;
$B=(\widehat A)_{\mathrm{red}}=\widehat A/\mathfrak N$ (où $\mathfrak N$ est
le nilradical de $\widehat A$) est un anneau local noethérien complet réduit,
dont l'anneau total des

\oldpage[0\textsubscript{IV-1}]{220}

fractions $R$ est donc composé direct d'un nombre fini de corps $L_i$~;
l'image canonique $B_i$ de $B$ dans $L_i$ est un anneau local noethérien
\emph{intègre et complet}, dont $L_i$ est le corps de fractions, et la
fermeture intégrale $B'$ de $B$ dans $R$ est \emph{composée directe} des
$B'_i$, où $B'_i$ est la clôture intégrale de $B_i$. Mais en vertu de
(23.1.5), $B'_i$ est une $B_i$-algèbre \emph{finie}, donc un anneau local
\emph{noethérien} intégralement clos. On sait (Bourbaki,
\emph{Alg. comm.}, chap.~VII, §~1, n\textsuperscript{o}~3, cor. du th.~2) que
$B'_i$ est un anneau de Krull. Or, pour tout $i$, on a un homomorphisme
$\varphi_i:K\to L_i$, qui est injectif, et par suite
$\varphi_i^{-1}(B'_i)$ est un anneau de Krull dans $K$. Mais comme
$A'=B'\cap K$ en vertu du lemme et que
$B'=\bigcap_i\varphi_i^{-1}(B'_i)$, $A'$ est intersection d'un nombre fini
d'anneaux de Krull et est par suite un anneau de Krull.

Pour prouver la dernière assertion de la proposition, notons qu'il existe une
$A$-algèbre finie $C\subset K$ telle que le morphisme
$\operatorname{Spec}(A')\to\operatorname{Spec}(C)$ soit radiciel et un
homéomorphisme (23.2.4)~; pour tout idéal premier
$\mathfrak p'\in\operatorname{Spec}(A')$, $\mathfrak p'$ est donc le seul
idéal premier de $A'$ au-dessus de $\mathfrak p=\mathfrak p'\cap C$ et
$A'_{\mathfrak p'}$ la clôture intégrale de $C_{\mathfrak p}$~; d'autre part
le fait que $\operatorname{Spec}(A')\to\operatorname{Spec}(C)$ est un
homéomorphisme entraîne que l'application
$\mathfrak p'\mapsto\mathfrak p'\cap C$ est une bijection de l'ensemble des
idéaux premiers de hauteur $1$ de $A'$ sur l'ensemble des idéaux premiers de
hauteur $1$ de $C$~; d'où la conclusion, compte tenu de Bourbaki,
\emph{Alg. comm.}, chap.~VII, §~1, n\textsuperscript{o}~6, th.~4.

\begin{remarks}[23.2.8]
\label{0.23.2.8-fr}
\begin{enumerate}
  \item[(i)] On aura soin de noter, dans l'application de (23.2.6) et
  (23.2.7), que l'anneau $A'$ \emph{n'est pas nécessairement un anneau
  noethérien}, comme le montre un exemple de Nagata avec $\dim(A)=3$ [30].

  \item[(ii)] Les conclusions de (23.2.7) sont encore valables si $A'$ est
  la fermeture intégrale de $A$ dans une extension finie $K'$ de $K$~; il
  suffit en effet de considérer une $A$-algèbre finie $B$ dont $K'$ est le
  corps des fractions~; $B$ est un anneau semi-local noethérien, donc
  intersection d'un nombre fini d'anneaux locaux noethériens
  $B_{\mathfrak m_i}$ ($\mathfrak m_i$ idéaux maximaux de $B$), et sa clôture
  intégrale (qui est égale à $A'$) est intersection des
  $A'_{\mathfrak m_i}$ (Bourbaki, \emph{Alg. comm.}, chap.~II, §~3,
  n\textsuperscript{o}~3, cor.~4 du th.~1) qui sont les clôtures intégrales
  des $B_{\mathfrak m_i}$, donc des anneaux de Krull par (23.1.13)~; par
  suite $A'$ est un anneau de Krull.

  \item[(iii)] On peut montrer ([30], 3.3.10) que pour tout anneau noethérien
  intègre $A$ (non nécessairement local), la clôture intégrale de $A$ est un
  anneau de Krull.
\end{enumerate}
\end{remarks}

\begin{corollary}[23.2.9]
\label{0.23.2.9-fr}
Soient $A$ un anneau noethérien intègre, $A'$ sa clôture intégrale. Supposons
que pour tout anneau $C$ tel que $A\subset C\subset A'$ et tel que $C$ soit
une $A$-algèbre finie, et pour tout idéal premier $\mathfrak q$ de hauteur $1$
dans $C$, $\mathfrak q\cap A$ soit un idéal premier de hauteur $1$ dans $A$.
Soit $P$ l'ensemble des idéaux premiers de hauteur $1$ dans $A$~; on a alors
\[
  A'=\bigcap_{\mathfrak p\in P}A'_{\mathfrak p}.
\]
\end{corollary}

\begin{proof}
Comme $A'$ est un $A$-module sans torsion, on a
$A'=\bigcap A'_{\mathfrak m}$, où $\mathfrak m$ parcourt l'ensemble des idéaux
maximaux de $A$ (Bourbaki, \emph{Alg. comm.}, chap.~II, §~3,
n\textsuperscript{o}~3, cor.~4 du th.~1)~;

\oldpage[0\textsubscript{IV-1}]{221}

il suffira donc de prouver que $A'_{\mathfrak m}$ est l'intersection des
$A'_{\mathfrak p}$ pour les idéaux premiers $\mathfrak p\subset\mathfrak m$
de hauteur $1$. Comme $A'_{\mathfrak m}$ est la clôture intégrale de
$A_{\mathfrak m}$, on voit qu'on peut se borner à démontrer le corollaire pour
$A_{\mathfrak m}$. Or, il y a alors, en vertu de (23.2.7), une
$A_{\mathfrak m}$-algèbre finie $B$ contenue dans $A'_{\mathfrak m}$, telle
que $A'_{\mathfrak m}$ soit l'intersection des clôtures intégrales des anneaux
$B_{\mathfrak q}$ où $\mathfrak q$ parcourt l'ensemble des idéaux premiers de
hauteur $1$ de $B$. Mais $B$ est de la forme $C_{\mathfrak m}$, où $C$ est
une $A$-algèbre finie~; donc, pour tout idéal premier $\mathfrak q$ de hauteur
$1$ dans $B$, $\mathfrak q\cap A$ est par hypothèse un idéal premier de
hauteur $1$ dans $A$~; comme la clôture intégrale de $B_{\mathfrak q}$
contient celle de $A_{\mathfrak q\cap A}$ et que cette dernière contient
$A'_{\mathfrak m}$, il en résulte bien que $A'_{\mathfrak m}$ est
l'intersection des $A'_{\mathfrak p}$ pour $\mathfrak p\in P$ et
$\mathfrak p\subset\mathfrak m$, ce qui achève la démonstration.
\end{proof}

\begin{remarks}[23.2.10]
\label{0.23.2.10-fr}
\begin{enumerate}
  \item[(i)] Nous verrons (IV, 5.6.3 et 5.6.10) que l'hypothèse de (23.2.9)
  est vérifiée lorsque $A$ est \emph{universellement caténaire}, en particulier
  (IV, 5.10.17 et 5.11.2) lorsque $A$ est \emph{quotient d'un anneau
  régulier}~; enfin, elle est évidemment vérifiée si
  $\operatorname{Spec}(A')\to\operatorname{Spec}(A)$ est un homéomorphisme,
  en particulier si ce morphisme est radiciel (23.2.2).

  \item[(ii)] La conclusion de (23.2.9) n'est plus nécessairement exacte
  lorsqu'on omet l'hypothèse sur les sous-$A$-algèbres finies de $A'$. En
  effet, on peut définir un anneau local noethérien intègre $A$, de dimension
  $2$, dont la clôture intégrale $A'$ est une $A$-algèbre finie, mais tel qu'il
  y ait un idéal premier $\mathfrak p'$ de $A'$ de hauteur $1$ au-dessus de
  l'idéal maximal $\mathfrak m$ (de hauteur $2$) de $A$, et tel en outre que
  pour tout idéal premier $\mathfrak p$ de hauteur $1$ dans $A$,
  $A_{\mathfrak p}$ soit intégralement clos (IV, 5.6.11)~; l'intersection de
  ces anneaux $A_{\mathfrak p}$ ne peut alors être égale à $A'$ puisque $A'$
  est un anneau de Krull (Bourbaki, \emph{Alg. comm.}, chap.~VII, §~1,
  n\textsuperscript{o}~5, prop.~9).
\end{enumerate}
\end{remarks}

\begin{flushright}
\emph{(A suivre.)}
\end{flushright}

% FIN SOURCE ega4/ega0-23-fr.tex

%
% DEBUT SOURCE ega4/ega4-1-fr.tex
\cleardoublepage
\markboth{CHAPITRE IV}{CHAPITRE IV}
\oldpage[IV-1]{222}
\thispagestyle{plain}

\begin{center}
{\Large CHAPITRE IV}

\vspace{1.5cm}
{\LARGE\bfseries ÉTUDE LOCALE DES SCHÉMAS\\[0.3em]
ET DES MORPHISMES DE SCHÉMAS}

\vspace{1.5cm}
{\large\bfseries Sommaire}
\end{center}

\medskip
\begin{tabular}{@{}r p{0.84\textwidth}}
§ 1. & Conditions de finitude relatives. Ensembles constructibles dans les préschémas.\\
§ 2. & Changement de base et platitude.\\
§ 3. & Cycles premiers associés et décompositions primaires.\\
§ 4. & Changement du corps de base dans les préschémas.\\
§ 5. & Dimension et profondeur dans les préschémas.\\
§ 6. & Morphismes plats de préschémas localement noethériens.\\
§ 7. & Application aux relations entre un anneau local noethérien et son complété.
Anneaux excellents.\\
§ 8. & Limites projectives de préschémas.\\
§ 9. & Propriétés constructibles.\\
§ 10. & Préschémas de Jacobson.\\
§ 11.\footnotemark[1] & Propriétés topologiques des morphismes plats de présentation finie.
Critères locaux de platitude.\\
§ 12. & Étude des fibres des morphismes plats de présentation finie.\\
§ 13. & Morphismes équidimensionnels.\\
§ 14. & Morphismes universellement ouverts.\\
§ 15. & Étude des fibres d'un morphisme universellement ouvert.\\
§ 16. & Invariants différentiels. Morphismes différentiellement lisses.\\
§ 17. & Morphismes lisses, non ramifiés, étales.\\
§ 18. & Compléments sur les morphismes étales. Anneaux locaux henséliens.\\
§ 19. & Immersions régulières et transversalement régulières.\\
§ 20. & Sections hyperplanes~; projections génériques.\\
§ 21. & Prolongements infinitésimaux.\\
\end{tabular}
\footnotetext[1]{L'ordre et le contenu des §§ 11 à 21 ne sont donnés qu'à titre
indicatif et pourront être modifiés avant leur publication.}

\clearpage
\oldpage[IV-1]{223}

Les sujets traités dans ce Chapitre appellent les remarques suivantes~:

\medskip
\noindent
\emph{a)} Leur caractère commun est d'être relatifs à des propriétés
\emph{locales} de préschémas ou de morphismes, i.e. considérées en un point,
ou aux points d'une fibre, ou dans un voisinage (non spécifié) d'un point ou
d'une fibre. Ces propriétés sont généralement de nature \emph{topologique},
\emph{différentielle} ou \emph{dimensionnelle} (i.e. faisant intervenir les
notions de \emph{dimension} et de \emph{profondeur}) et sont liées aux
propriétés des \emph{anneaux locaux} aux points considérés. Un problème type
est de mettre en relation, pour un morphisme donné $f:X\to Y$ et un point
$x\in X$, les propriétés de $X$ en $x$ avec celles de $Y$ en $y=f(x)$ et de
la fibre $X_y=f^{-1}(y)$ en $x$. Un autre est de déterminer la nature
topologique (par exemple la constructibilité, ou le fait d'être ouvert ou
fermé) de l'ensemble des points $x\in X$ en lesquels $X$ possède une certaine
propriété ou pour lesquels la fibre $X_{f(x)}$ passant par $x$ possède une
certaine propriété en $x$. De même, on s'intéressera à la nature topologique
de l'ensemble des points $y\in Y$ tels que $X$ ait une certaine propriété en
tous les points de la fibre $X_y$, ou tels que cette fibre ait une certaine
propriété.

\medskip
\noindent
\emph{b)} Les notions les plus importantes pour les chapitres suivants sont
celle de \emph{morphisme plat de présentation finie} et celles de
\emph{morphisme lisse} et de \emph{morphisme étale}, qui en sont des cas
particuliers. Leur étude détaillée (jointe à celle de questions connexes)
commence à proprement parler au §~11.

\medskip
\noindent
\emph{c)} Les §§~1 à 10 peuvent être considérés comme de nature préliminaire
et développent trois types de techniques, utilisées non seulement dans les
autres paragraphes du chapitre mais aussi bien entendu dans les chapitres
suivants~:

\medskip
\noindent
\emph{c 1)} Aux §§~1 à 4 sont envisagés divers aspects de la notion de
\emph{changement de base}, surtout en relation avec des conditions de
\emph{finitude} ou de \emph{platitude}~; on y amorce la technique de
\emph{descente} sous ses aspects les plus élémentaires (les questions
d'«~effectivité~» liées à cette technique seront étudiées dans le chap.~V).

\medskip
\noindent
\emph{c 2)} Les §§~5 à 7 sont centrés sur ce que l'on peut appeler des
\emph{techniques noethériennes}, les préschémas considérés étant toujours
supposés localement noethériens, tandis qu'au contraire aucune condition de
finitude n'est généralement imposée aux \emph{morphismes}~; cela tient
essentiellement à ce que les notions de dimension et de profondeur ne sont
guère maniables que pour les anneaux locaux noethériens. Rappelons que le §~7
constitue une théorie «~fine~» des anneaux locaux noethériens, assez peu
utilisée dans la suite du chapitre.

\medskip
\noindent
\emph{c 3)} Les §§~8 à 10 donnent entre autres le moyen d'\emph{éliminer les
hypothèses noethériennes} sur les préschémas étudiés, en leur substituant des
hypothèses convenables de \emph{finitude} (la «~présentation finie~») sur les
\emph{morphismes} considérés~: l'avantage de cette substitution est que ces
dernières hypothèses sont \emph{stables par changement de base}, ce qui n'est
pas le cas pour les hypothèses noethériennes sur les préschémas. La technique
permettant ce passage repose, d'une part, sur l'utilisation de la notion de
\emph{limite projective} de préschémas, grâce à laquelle on peut ramener une
question à la même question sous des hypothèses \emph{noethériennes}~; d'autre
part, sur l'emploi systématique des \emph{ensembles constructibles}, qui ont
le double intérêt de se conserver par image réciproque (par un morphisme

\oldpage[IV-1]{224}

quelconque) et par image directe (par un morphisme de présentation finie), et
d'avoir des propriétés topologiques maniables dans les préschémas localement
noethériens. Les mêmes techniques permettent même souvent de se ramener à des
anneaux noethériens plus particuliers, par exemple les
$\mathbf Z$-\emph{algèbres de type fini}, et c'est là que les propriétés des
anneaux «~excellents~» (étudiées au §~7) interviennent de façon décisive.
Indépendamment d'ailleurs de la question de l'élimination des hypothèses
noethériennes, les techniques des §§~8 à 10, de nature très élémentaire, sont
d'usage constant dans presque toutes les applications.

\setcounter{section}{0}

\section{Conditions de finitude relatives. Ensembles constructibles dans les préschémas}
\label{section:IV.1-fr}

Nous allons, dans ce paragraphe, reprendre, en les complétant, l'exposé des
«~conditions de finitude~» pour un morphisme de préschémas $f:X\to Y$, donné
dans (I, 6.3 et 6.6). Il y a essentiellement deux notions de «~finitude~» de
nature \emph{globale} sur $X$, celle de morphisme \emph{quasi-compact}
(définie dans (I, 6.6.1)) et celle de morphisme \emph{quasi-séparé}~; il y a
d'autre part deux notions de «~finitude~» de nature \emph{locale} sur $X$,
celle de morphisme \emph{localement de type fini} (définie dans (I, 6.6.2))
et celle de morphisme \emph{localement de présentation finie}. En combinant
ces notions locales aux notions globales précédentes, on obtient les notions
de morphisme \emph{de type fini} (définie dans (I, 6.3.1)) et de morphisme
\emph{de présentation finie}. Pour la commodité du lecteur, nous donnerons de
nouveau, dans ce paragraphe, les principales propriétés énoncées dans (I, 6.3
et 6.6), en renvoyant bien entendu à ces numéros du chapitre I pour leurs
démonstrations.

Aux n\textsuperscript{os}~(1.8) et (1.9), nous complétons, dans le cadre des
préschémas, et en faisant usage des notions de finitude précédentes, les
résultats sur les ensembles constructibles donnés dans
(0\textsubscript{III}, §~9).

\subsection{Morphismes quasi-compacts}
\label{subsection:IV.1.1-fr}

\begin{definition}[1.1.1]
\label{IV.1.1.1-fr}
On dit qu'un morphisme de préschémas $f:X\to Y$ est \emph{quasi-compact} si
l'application continue $f$ de l'espace topologique $X$ dans l'espace
topologique $Y$ est quasi-compacte (0\textsubscript{III}, 9.1.1), autrement
dit si l'image réciproque $f^{-1}(U)$ de tout ouvert quasi-compact $U$ de $Y$
est quasi-compacte (cf. (I, 6.6.1)).
\end{definition}

Si $\mathfrak B$ est une base de la topologie de $Y$ formée d'ouverts
affines, pour que $f$ soit quasi-compact, il faut et il suffit que pour tout
$V\in\mathfrak B$, $f^{-1}(V)$ soit \emph{réunion finie d'ouverts affines}.
Par exemple, si $Y$ est affine et $X$ quasi-compact, \emph{tout} morphisme
$f:X\to Y$ est quasi-compact (I, 6.6.1).

Si $f:X\to Y$ est un morphisme quasi-compact, il est clair que pour tout
ouvert $V$ de $Y$, la restriction de $f$ à $f^{-1}(V)$ est un morphisme
quasi-compact $f^{-1}(V)\to V$. Inversement, si $(U_\alpha)$ est un
recouvrement ouvert de $Y$ et $f:X\to Y$ un morphisme tel que les restrictions
$f^{-1}(U_\alpha)\to U_\alpha$ soient quasi-compactes, alors $f$ est
quasi-compact. Par

\oldpage[IV-1]{225}

suite, si $f:X\to Y$ est un $S$-morphisme de $S$-préschémas, et s'il existe
un recouvrement ouvert $(S_\lambda)$ de $S$ tel que les restrictions
$g^{-1}(S_\lambda)\to h^{-1}(S_\lambda)$ de $f$ (où $g$ et $h$ sont les
morphismes structuraux) soient quasi-compactes, alors $f$ est quasi-compact.

\begin{proposition}[1.1.2]
\label{IV.1.1.2-fr}
\begin{enumerate}
  \item[(i)] Une immersion $X\to Y$ est quasi-compacte si elle est fermée, ou si
  l'espace sous-jacent à $Y$ est localement noethérien ou si l'espace
  sous-jacent à $X$ est noethérien.
  \item[(ii)] Le composé de deux morphismes quasi-compacts est quasi-compact.
  \item[(iii)] Si $f:X\to Y$ est un $S$-morphisme quasi-compact, il en est de même de
  $f_{(S')}:X_{(S')}\to Y_{(S')}$ pour toute extension $g:S'\to S$ du
  préschéma de base.
  \item[(iv)] Si $f:X\to X'$ et $g:Y\to Y'$ sont deux $S$-morphismes
  quasi-compacts,
  \[
    f\times_S g:X\times_S Y\to X'\times_S Y'
  \]
  est quasi-compact.
  \item[(v)] Si le composé $g\circ f$ de deux morphismes $f:X\to Y$, $g:Y\to Z$
  est quasi-compact, et si $g$ est séparé, ou l'espace sous-jacent à $X$
  localement noethérien, $f$ est quasi-compact.
  \item[(vi)] Pour qu'un morphisme $f$ soit quasi-compact, il faut et il suffit que
  $f_{\mathrm{red}}$ le soit.
\end{enumerate}
\end{proposition}

Pour la démonstration, voir (I, 6.6.4). On notera que l'assertion (vi) est
aussi conséquence de la proposition plus générale~:

\begin{proposition}[1.1.3]
\label{IV.1.1.3-fr}
Soient $f:X\to Y$, $g:Y\to Z$ deux morphismes. Si $g\circ f$ est
quasi-compact et $f$ surjectif, $g$ est quasi-compact.
\end{proposition}

\begin{proof}
En effet, si $V$ est un ouvert quasi-compact dans $Z$,
$f^{-1}(g^{-1}(V))$ est quasi-compact par hypothèse, et l'on a
$g^{-1}(V)=f(f^{-1}(g^{-1}(V)))$, puisque $f$ est surjectif~; donc
$g^{-1}(V)$ est quasi-compact.
\end{proof}

\begin{corollary}[1.1.4]
\label{IV.1.1.4-fr}
Soient $f:X\to Y$, $g:Y'\to Y$ deux morphismes, $X'=X\times_Y Y'$,
$f'=f_{(Y')}:X'\to Y'$~; supposons $g$ quasi-compact et surjectif (resp.
surjectif). Alors, pour que $f$ soit quasi-compact (resp. dominant), il suffit
que $f'$ le soit.
\end{corollary}

\begin{proof}
Si $g':X'\to X$ est la projection canonique, $g'$ est un morphisme surjectif
(I, 3.5.2), et l'on a $f\circ g'=g\circ f'$~; si $f'$ est quasi-compact
(resp. dominant), il en est de même de $g\circ f'$ puisque $g$ est
quasi-compact (1.1.2) (resp. surjectif)~; comme $f\circ g'$ est quasi-compact
(resp. dominant) et que $g'$ est surjectif, on en déduit que $f$ est
quasi-compact (1.1.3) (resp. dominant).
\end{proof}

\oldpage[IV-1]{226}

Pour abréger, nous appellerons \emph{point maximal} d'un préschéma $X$ tout
point générique d'une composante irréductible de $X$~; si
$X=\operatorname{Spec}(A)$ est affine, les points maximaux de $X$ sont les
\emph{idéaux premiers minimaux} de $A$.

\begin{proposition}[1.1.5]
\label{IV.1.1.5-fr}
Soit $f:X\to Y$ un morphisme quasi-compact. Les conditions suivantes sont
équivalentes~:
\begin{enumerate}
  \item[a)] $f$ est dominant~;
  \item[b)] pour tout point maximal $y$ de $Y$, on a
  $f^{-1}(y)\neq\varnothing$~;
  \item[c)] pour tout point maximal $y$ de $Y$, $f^{-1}(y)$ contient un point
  maximal de $X$.
\end{enumerate}
\end{proposition}

L'équivalence de \emph{a)} et \emph{c)} a été démontrée dans (I, 6.6.5). Il
est clair que \emph{c)} entraîne \emph{b)}~; d'autre part, si $X'$ est une
composante irréductible de $X$ rencontrant $f^{-1}(y)$, $f(X')$ est
irréductible et contient $y$, donc est contenu dans la composante irréductible
$Y'=\overline{\{y\}}$ de $Y$ (0\textsubscript{I}, 2.1.6)~; si $x$ est le
point générique de $X'$, on a donc nécessairement $f(x)=y$
(0\textsubscript{I}, 2.1.5)~; donc \emph{b)} entraîne \emph{c)}.

\begin{proposition}[1.1.6]
\label{IV.1.1.6-fr}
Soient $f':X'\to Y$, $f'':X''\to Y$ deux morphismes, $X$ le préschéma somme
$X'\coprod X''$, $f:X\to Y$ le morphisme égal à $f'$ dans $X'$ et à $f''$
dans $X''$. Pour que $f$ soit quasi-compact, il faut et il suffit que $f'$ et
$f''$ le soient.
\end{proposition}

Cela résulte aussitôt de la définition.

\subsection{Morphismes quasi-séparés}
\label{subsection:IV.1.2-fr}

\begin{definition}[1.2.1]
\label{IV.1.2.1-fr}
Soient $X$, $Y$ deux préschémas. On dit qu'un morphisme $f:X\to Y$ est
\emph{quasi-séparé} (ou que $X$ est \emph{quasi-séparé sur} $Y$) si le
morphisme diagonal
$\Delta_f=(1_X,1_X)_Y:X\to X\times_Y X$ est quasi-compact. On dit qu'un
préschéma $X$ est \emph{quasi-séparé} s'il est quasi-séparé sur
$\operatorname{Spec}(\mathbf Z)$.
\end{definition}

Par définition, tout morphisme séparé est quasi-séparé, $\Delta_f$ étant une
immersion fermée (1.1.2, (i))~; un \emph{schéma} $X$ est un préschéma
quasi-séparé, étant séparé sur $\operatorname{Spec}(\mathbf Z)$ (I, 5.4.1).

\begin{proposition}[1.2.2]
\label{IV.1.2.2-fr}
\begin{enumerate}
  \item[(i)] Tout monomorphisme de préschémas $f:X\to Y$ (en particulier toute
  immersion) est quasi-séparé.
  \item[(ii)] Si $f:X\to Y$ et $g:Y\to Z$ sont deux morphismes quasi-séparés,
  $g\circ f:X\to Z$ est quasi-séparé.
  \item[(iii)] Soient $X$, $Y$ deux $S$-préschémas, $f:X\to Y$ un $S$-morphisme
  quasi-séparé. Alors, pour tout changement de base $g:S'\to S$, le morphisme
  $f_{(S')}:X_{(S')}\to Y_{(S')}$ est quasi-séparé.
  \item[(iv)] Si $f:X\to Y$, $f':X'\to Y'$ sont deux $S$-morphismes quasi-séparés,
  alors
  \[
    f\times_S f':X\times_S X'\to Y\times_S Y'
  \]
  est quasi-séparé.
  \item[(v)] Si $f:X\to Y$, $g:Y\to Z$ sont deux morphismes tels que $g\circ f$
  soit quasi-séparé, alors $f$ est quasi-séparé.
  \item[(vi)] Si $f:X\to Y$ est un morphisme quasi-séparé,
  $f_{\mathrm{red}}:X_{\mathrm{red}}\to Y_{\mathrm{red}}$ est quasi-séparé.
\end{enumerate}
\end{proposition}

\begin{proof}
En vertu de (I, 5.5.12), il suffit de démontrer (i), (ii) et (iii).

L'assertion (i) est triviale, tout monomorphisme étant séparé (I, 5.5.1).
Pour démontrer (iii), on se ramène aussitôt au cas où $Y=S$ (I, 3.3.11) et
l'assertion résulte de ce que
$\Delta_{f_{(S')}}=(\Delta_f)_{(S')}$ (I, 5.3.4) et de (1.1.2, (iii)). Pour
démontrer (ii), considérons les projections $p$ et $q$ de $X\times_Y X$ sur
$X$~; si $i=(p,q)_Z$, on sait que le diagramme
\[
\xymatrix{
  X\times_Y X \ar[r]^-{i} \ar[d]_-{\pi} & X\times_Z X
    \ar[d]^-{f\times_Z f} \\
  Y \ar[r]_-{\Delta_g} & Y\times_Z Y
}
\]
\oldpage[IV-1]{227}
(où $\pi$ est le morphisme structural) est commutatif et identifie
$X\times_Y X$ au produit des $(Y\times_Z Y)$-préschémas $Y$ et
$X\times_Z X$ (I, 5.3.5). Si $g$ est quasi-séparé, $\Delta_g$ est
quasi-compact, donc il en est de même de $i$ (1.1.2, (iii))~; si de plus $f$
est quasi-séparé, $\Delta_f$ est quasi-compact, donc il en est de même de
$i\circ\Delta_f$ (1.1.2, (ii)) qui est égal à $\Delta_{g\circ f}$.
\end{proof}

\begin{corollary}[1.2.3]
\label{IV.1.2.3-fr}
\begin{enumerate}
  \item[(i)] Soit $X$ un préschéma quasi-séparé. Alors tout morphisme $f:X\to Y$
  est quasi-séparé.
  \item[(ii)] Soit $Y$ un préschéma quasi-séparé. Pour qu'un morphisme $f:X\to Y$
  soit quasi-séparé, il faut et il suffit que le préschéma $X$ soit
  quasi-séparé.
\end{enumerate}
\end{corollary}

Cela résulte aussitôt de (1.2.2, (ii) et (v)) appliqué à $X$, $Y$ et
$\operatorname{Spec}(\mathbf Z)$.

\begin{proposition}[1.2.4]
\label{IV.1.2.4-fr}
Soient $f:X\to Y$ un morphisme, $g:Y\to Z$ un morphisme quasi-séparé. Si
$g\circ f$ est quasi-compact, il en est de même de $f$.
\end{proposition}

\begin{proof}
On sait que $f$ est le morphisme composé
$X\xrightarrow{\Gamma_f}X\times_ZY\xrightarrow{\mathrm{pr}_2}Y$
(I, 5.3.13)~; d'autre part, $\mathrm{pr}_2$ s'identifie à
$(g\circ f)\times_Z1_Y$ (I, 3.3.4) et si $g\circ f$ est quasi-compact, il en
est de même de $\mathrm{pr}_2$ (1.1.2, (iv)). Enfin, on a le diagramme
commutatif
\[
\xymatrix{
  X \ar[r]^-{\Gamma_f} \ar[d]_-{f} & X\times_ZY
    \ar[d]^-{f\times 1_Y} \\
  Y \ar[r]_-{\Delta_g} & Y\times_ZY
}
\tag{1.2.4.1}
\label{IV.1.2.4.1-fr}
\]
qui identifie $X$ au produit des $(Y\times_ZY)$-préschémas $Y$ et
$X\times_ZY$ (I, 5.3.7). Comme par hypothèse $\Delta_g$ est un morphisme
quasi-compact, il en est de même de $\Gamma_f$ (1.1.2, (iii))~; la conclusion
résulte donc de (1.1.2, (ii)).
\end{proof}

\begin{proposition}[1.2.5]
\label{IV.1.2.5-fr}
Soient $f:X\to Y$, $g:Y'\to Y$ deux morphismes,
$X'=X\times_Y Y'$, $f'=f_{(Y')}:X'\to Y'$. Si $g$ est quasi-compact et
surjectif et $f'$ quasi-séparé, alors $f$ est quasi-séparé.
\end{proposition}

\begin{proof}
On a en effet $(X'\times_{Y'}X')=(X\times_YX)_{(Y')}$ et
$\Delta_{f'}=(\Delta_f)_{(Y')}$ (I, 5.3.4)~; comme la projection
$X'\times_{Y'}X'\to X\times_YX$ est quasi-compacte (1.1.2, (iii)) et
surjective (I, 3.5.2), on peut, en vertu de (I, 3.3.11), appliquer (1.1.4),
qui montre que puisque $\Delta_{f'}$ est un morphisme quasi-compact, il en est
de même de $\Delta_f$.
\end{proof}

\begin{proposition}[1.2.6]
\label{IV.1.2.6-fr}
Soient $f:X\to Y$ un morphisme, $(U_\alpha)$ un recouvrement de $Y$ par des
ouverts tels que les préschémas induits $U_\alpha$ soient quasi-séparés. Pour
que $f$ soit quasi-séparé, il faut et il suffit que chacun des préschémas
induits par $X$ sur les ouverts $f^{-1}(U_\alpha)$ soit quasi-séparé.
\end{proposition}

\begin{proof}
L'image réciproque dans $X\times_YX$ de $U_\alpha$ est
$X_\alpha\times_{U_\alpha}X_\alpha$, en posant
$X_\alpha=f^{-1}(U_\alpha)$, et la restriction
$X_\alpha\to X_\alpha\times_{U_\alpha}X_\alpha$ de $\Delta_f$ n'est autre
que $\Delta_{f_\alpha}$, en désignant par $f_\alpha$ la restriction
$X_\alpha\to U_\alpha$ de $f$~; en vertu de la définition (1.2.1) et du
caractère local de la notion de morphisme quasi-compact (1.1.1), pour que $f$
soit quasi-séparé, il faut et il suffit que chacun des morphismes $f_\alpha$
le soit. Mais comme par hypothèse le morphisme
$U_\alpha\to\operatorname{Spec}(\mathbf Z)$ est quasi-séparé, il revient au
même de dire que $f_\alpha$ est quasi-séparé ou que le composé
$X_\alpha\xrightarrow{f_\alpha}U_\alpha\to\operatorname{Spec}(\mathbf Z)$
l'est (1.2.2, (ii) et (v)), d'où la proposition.
\end{proof}

\oldpage[IV-1]{228}

Pour vérifier qu'un morphisme est quasi-séparé, on est donc ramené à donner
des critères pour qu'un préschéma soit quasi-séparé~:

\begin{proposition}[1.2.7]
\label{IV.1.2.7-fr}
Soient $X$ un préschéma, $(U_\alpha)$ un recouvrement de $X$ formé d'ouverts
quasi-compacts. Les propriétés suivantes sont équivalentes~:
\begin{enumerate}
  \item[a)] $X$ est quasi-séparé.
  \item[b)] Pour tout ouvert quasi-compact $U$ de $X$, l'immersion canonique
  $U\to X$ est quasi-compacte (autrement dit, $U$ est rétrocompact dans $X$).
  \item[b')] L'intersection de deux ouverts quasi-compacts de $X$ est
  quasi-compacte.
  \item[c)] Pour tout couple d'indices $(\alpha,\beta)$, l'intersection
  $U_\alpha\cap U_\beta$ est quasi-compacte.
\end{enumerate}
\end{proposition}

\begin{proof}
Les propriétés \emph{b)} et \emph{b')} sont trivialement équivalentes par
définition d'un morphisme quasi-compact. Comme un ouvert quasi-compact est
réunion finie d'ouverts affines, pour deux ouverts quasi-compacts $U$, $V$ de
$X$, $U\times_{\mathbf Z}V$ est un ouvert quasi-compact de
$X\times_{\mathbf Z}X$ (I, 3.2.7), dont l'image réciproque par $\Delta_X$ est
$U\cap V$, donc \emph{a)} entraîne \emph{b')}. Il est trivial que \emph{b')}
entraîne \emph{c)}~; enfin, si \emph{c)} est vérifiée, les
$U_\alpha\times_{\mathbf Z}U_\beta$ forment un recouvrement de
$X\times_{\mathbf Z}X$ par des ouverts quasi-compacts, et l'on sait que pour
que le morphisme $\Delta_X:X\to X\times_{\mathbf Z}X$ soit quasi-compact, il
suffit que les images réciproques $U_\alpha\cap U_\beta$ des
$U_\alpha\times_{\mathbf Z}U_\beta$ par $\Delta_X$ soient quasi-compactes
(1.1.1)~; donc \emph{c)} entraîne \emph{a)}.
\end{proof}

\begin{corollary}[1.2.8]
\label{IV.1.2.8-fr}
Tout préschéma $X$ dont l'espace sous-jacent est localement noethérien (par
exemple un préschéma localement noethérien) est quasi-séparé~; tout morphisme
$f:X\to Y$ est alors quasi-séparé.
\end{corollary}

\begin{proposition}[1.2.9]
\label{IV.1.2.9-fr}
Soient $X'$, $X''$ deux préschémas, $X=X'\coprod X''$ leur somme,
$f':X'\to Y$, $f'':X''\to Y$ deux morphismes, $f:X\to Y$ le morphisme qui
coïncide avec $f'$ dans $X'$ et avec $f''$ dans $X''$. Pour que $f$ soit
quasi-séparé il faut et il suffit que $f'$ et $f''$ le soient.
\end{proposition}

\begin{proof}
En effet, $X\times_YX$ est somme des quatre préschémas
$X'\times_YX'$, $X'\times_YX''$, $X''\times_YX'$ et $X''\times_YX''$, et
$\Delta_f$ est le morphisme qui coïncide avec $\Delta_{f'}$ dans $X'$ et avec
$\Delta_{f''}$ dans $X''$~; la proposition résulte donc aussitôt des
définitions.
\end{proof}

\subsection{Morphismes localement de type fini}
\label{subsection:IV.1.3-fr}

\begin{env}[1.3.1]
\label{IV.1.3.1-fr}
Rappelons que si $A$ est un anneau, une $A$-algèbre (commutative) $B$ est dite
\emph{de type fini} si elle est engendrée par un nombre fini d'éléments de
$B$, ou, ce qui revient au même, si elle est isomorphe au quotient d'une
algèbre de polynômes $A[T_1,\ldots,T_n]$ par un idéal de cette algèbre. Il est
clair que pour toute $A$-algèbre (commutative) $A'$,
$B\otimes_AA'$ est alors une $A'$-algèbre de type fini. Si $B$ est une
$A$-algèbre de type fini et $C$ une $B$-algèbre de type fini, $C$ est une
$A$-algèbre de type fini, car si $C$ est un quotient de
$B[T_1,\ldots,T_n]$ et $B$ un quotient de $A[S_1,\ldots,S_m]$,
$B[T_1,\ldots,T_n]=B\otimes_AA[T_1,\ldots,T_n]$ est un quotient de
$A[S_1,\ldots,S_m,T_1,\ldots,T_n]$, donc il en est de même de $C$.
\end{env}

\begin{definition}[1.3.2]
\label{IV.1.3.2-fr}
Soient $f:X\to Y$ un morphisme de préschémas, $x$ un point de $X$,
$y=f(x)$. On dit que $f$ est \emph{de type fini en $x$} s'il existe un
voisinage ouvert affine $V$ de $y$ et un voisinage ouvert affine $U$ de $x$
tels que $f(U)\subset V$ et que l'anneau $A(U)$ soit une $A(V)$-algèbre de
type fini. On dit que $f$ est \emph{localement de type fini} s'il est de type
fini en tout point de $X$.
\end{definition}

\oldpage[IV-1]{229}

\begin{proposition}[1.3.3]
\label{IV.1.3.3-fr}
Si $Y$ est un préschéma localement noethérien et $f:X\to Y$ un morphisme
localement de type fini, alors $X$ est localement noethérien.
\end{proposition}

Pour la démonstration, voir (I, 6.3.7).

\begin{proposition}[1.3.4]
\label{IV.1.3.4-fr}
\begin{enumerate}
  \item[(i)] Toute immersion locale est localement de type fini.
  \item[(ii)] Si deux morphismes $f:X\to Y$, $g:Y\to Z$ sont localement de
  type fini, il en est de même de $g\circ f$.
  \item[(iii)] Si $f:X\to Y$ est un $S$-morphisme localement de type fini,
  $f_{(S')}:X_{(S')}\to Y_{(S')}$ est localement de type fini pour toute
  extension $S'\to S$ du préschéma de base.
  \item[(iv)] Si $f:X\to X'$ et $g:Y\to Y'$ sont deux $S$-morphismes
  localement de type fini, $f\times_Sg$ est localement de type fini.
  \item[(v)] Si le composé $g\circ f$ de deux morphismes est localement de
  type fini, $f$ est localement de type fini.
  \item[(vi)] Si un morphisme $f$ est localement de type fini, il en est de
  même de $f_{\mathrm{red}}$.
\end{enumerate}
\end{proposition}

Pour la démonstration, voir (I, 6.6.6).

\begin{corollary}[1.3.5]
\label{IV.1.3.5-fr}
Soit $f:X\to Y$ un morphisme localement de type fini. Pour tout morphisme
$Y'\to Y$ tel que $Y'$ soit localement noethérien, $X\times_YY'$ est
localement noethérien.
\end{corollary}

\begin{proof}
En effet, $f_{(Y')}:X\times_YY'\to Y'$ est localement de type fini par
(1.3.4, (iii)), et il suffit d'appliquer (1.3.3).
\end{proof}

\begin{proposition}[1.3.6]
\label{IV.1.3.6-fr}
Soit $\varphi:A\to B$ un homomorphisme d'anneaux. Pour que le morphisme
correspondant $f:\operatorname{Spec}(B)\to\operatorname{Spec}(A)$ soit localement de type fini, il faut et
il suffit que $B$ soit une $A$-algèbre de type fini.
\end{proposition}

Pour la démonstration, voir (I, 6.3.3), en tenant compte de ce que $f$ est
quasi-compact et de (I, 6.6.3).

\begin{proposition}[1.3.7]
\label{IV.1.3.7-fr}
Pour qu'un morphisme localement de type fini $f:X\to Y$ soit surjectif, il
faut et il suffit que, pour tout corps algébriquement clos $\Omega$,
l'application $X(\Omega)\to Y(\Omega)$ correspondant à $f$ (I, 3.4.1) soit
surjective.
\end{proposition}

Pour la démonstration, voir (I, 6.3.10).

\begin{env}[1.3.8]
\label{IV.1.3.8-fr}
Soit $A$ un anneau. On dit qu'une $A$-algèbre $B$ est \emph{essentiellement de
type fini} si $B$ est $A$-isomorphe à une $A$-algèbre de la forme $S^{-1}C$,
où $C$ est une $A$-algèbre \emph{de type fini} et $S$ une partie multiplicative
de $C$.
\end{env}

\par\medskip\noindent
\begin{proposition}[1.3.9]
\label{IV.1.3.9-fr}
\begin{enumerate}
  \item[(i)] Si $B$ est une $A$-algèbre essentiellement de type fini et $C$
  une $B$-algèbre essentiellement de type fini, alors $C$ est une
  $A$-algèbre essentiellement de type fini.
  \item[(ii)] Soient $B$ une $A$-algèbre essentiellement de type fini, $A'$
  une $A$-algèbre~; alors $B'=B\otimes_AA'$ est une $A'$-algèbre
  essentiellement de type fini.
\end{enumerate}
\end{proposition}

\begin{proof}
\begin{enumerate}
  \item[(i)] Soient $B=S'^{-1}B'$ et $C=T'^{-1}C'$, où $B'$ (resp. $C'$) est
  une $A$-algèbre (resp. une $B$-algèbre) de type fini et $S'$ (resp. $T'$)
  une partie multiplicative de $B'$ (resp. $C'$)~; alors $C'$ est de la forme
  $S'^{-1}C''$, où $C''$ est une $B'$-algèbre de type fini, donc $C$ est aussi
  un anneau de fractions de $C''$~; comme $C''$ est une $A$-algèbre de type
  fini, cela prouve notre assertion.
  \item[(ii)] Si $B$ est anneau de fractions d'une $A$-algèbre de type fini
  $C$, $B'$ est un anneau de fractions de la $A'$-algèbre de type fini
  $C\otimes_AA'$, d'où (ii).
\end{enumerate}
\end{proof}

\oldpage[IV-1]{230}

\begin{env}[1.3.10]
\label{IV.1.3.10-fr}
Si $B$ est une $A$-algèbre \emph{locale} essentiellement de type fini, $B$ est
de la forme $C_{\mathfrak q}$, où $C$ est une $A$-algèbre de type fini et
$\mathfrak q$ un idéal premier de $C$ (Bourbaki, \emph{Alg. comm.}, chap. II,
\S~5, prop. 11). Soit $\mathfrak p$ l'image réciproque dans $A$ de l'idéal
$\mathfrak q$~; si l'on pose $S=A-\mathfrak p$, $C_{\mathfrak q}$ est aussi un
anneau local en un idéal premier de $S^{-1}C$~; comme $S^{-1}C$ est une
algèbre de type fini sur $A_{\mathfrak p}=S^{-1}A$, on voit que $B$ est aussi
une $A_{\mathfrak p}$-algèbre \emph{essentiellement de type fini}, et en outre
l'homomorphisme $A_{\mathfrak p}\to B$ est \emph{local}.
\end{env}

\begin{proposition}[1.3.11]
\label{IV.1.3.11-fr}
Si $B$ est une $A$-algèbre locale essentiellement de type fini, $B$ est
$A$-isomorphe à une $A$-algèbre quotient d'une $A$-algèbre de la forme
$C_{\mathfrak q}$, où $C$ est une algèbre de polynômes
$A[T_1,\ldots,T_n]$ et $\mathfrak q$ un idéal premier de $C$.
\end{proposition}

\begin{proof}
En effet, par définition, $B$ est isomorphe à $C'_{\mathfrak q'}$, où $C'$ est
une $A$-algèbre de type fini et $\mathfrak q'$ un idéal premier de $C'$~; mais
$C'=C/\mathfrak b$, où $C=A[T_1,\ldots,T_n]$ et $\mathfrak b$ est un idéal de
$C$~; donc $\mathfrak q'=\mathfrak q/\mathfrak b$, où $\mathfrak q$ est un
idéal premier de $C$, et $C'_{\mathfrak q'}$ est isomorphe à
$C_{\mathfrak q}/\mathfrak bC_{\mathfrak q}$.
\end{proof}

\subsection{Morphismes localement de présentation finie}
\label{subsection:IV.1.4-fr}

\begin{env}[1.4.1]
\label{IV.1.4.1-fr}
Étant donné un anneau $A$, nous dirons qu'une $A$-algèbre $B$ (commutative)
est \emph{de présentation finie} si elle est isomorphe au quotient
$A[T_1,\ldots,T_n]/\mathfrak a$ d'une algèbre de polynômes sur $A$ par un
idéal \emph{de type fini} $\mathfrak a$ de $A[T_1,\ldots,T_n]$. Pour toute
$A$-algèbre (commutative) $A'$, $B\otimes_AA'$ est alors une $A'$-algèbre de
présentation finie, étant isomorphe au quotient de $A'[T_1,\ldots,T_n]$ par
l'image canonique dans cet anneau de $\mathfrak a\otimes_AA'$, qui est
évidemment un $A'[T_1,\ldots,T_n]$-module de type fini. Si $B$ est une
$A$-algèbre de présentation finie et $C$ une $B$-algèbre de présentation
finie, alors $C$ est une $A$-algèbre de présentation finie. En effet, soient
$B=A[S_1,\ldots,S_m]/\mathfrak a$ et
$C=B[T_1,\ldots,T_n]/\mathfrak b$, où $\mathfrak a$ (resp. $\mathfrak b$) est
un idéal de type fini de $A[S_1,\ldots,S_m]$ (resp.
$B[T_1,\ldots,T_n]$)~; l'anneau $B[T_1,\ldots,T_n]$ est isomorphe à
$A[S_1,\ldots,S_m,T_1,\ldots,T_n]/\mathfrak a'$, où $\mathfrak a'$ est
l'image canonique de
$\mathfrak a\otimes_AA[T_1,\ldots,T_n]$, et est donc un idéal de type fini de
$A[S_1,\ldots,S_m,T_1,\ldots,T_n]$~; l'idéal $\mathfrak b$ de
$B[T_1,\ldots,T_n]$ est de la forme $\mathfrak b'/\mathfrak a'$, où
$\mathfrak b'$ est un idéal de $A[S_1,\ldots,S_m,T_1,\ldots,T_n]$~; comme
$\mathfrak a'$ et $\mathfrak b'/\mathfrak a'$ sont des modules de type fini
sur $A[S_1,\ldots,S_m,T_1,\ldots,T_n]$, il en est de même de $\mathfrak b'$,
et $C$, isomorphe à
$A[S_1,\ldots,S_m,T_1,\ldots,T_n]/\mathfrak b'$, est donc bien une
$A$-algèbre de présentation finie. Notons enfin que si $A$ est
\emph{noethérien}, il en est de même de $A[T_1,\ldots,T_n]$, donc toute
$A$-algèbre de type fini est alors de présentation finie.
\end{env}

\begin{definition}[1.4.2]
\label{IV.1.4.2-fr}
Soient $f:X\to Y$ un morphisme de préschémas, $x$ un point de $X$, $y=f(x)$.
On dit que $f$ est \emph{de présentation finie en $x$} s'il existe un voisinage
ouvert affine $V$ de $y$ et un voisinage ouvert affine $U$ de $x$ tels que
$f(U)\subset V$ et que l'anneau $A(U)$ soit une $A(V)$-algèbre de présentation
finie. On dit que $f$ est \emph{localement de présentation finie} s'il est de
présentation finie en tout point de $X$.
\end{definition}

Si $Y$ est localement noethérien, il revient au même de dire que $f$ est
localement de type fini ou localement de présentation finie.

\begin{proposition}[1.4.3]
\label{IV.1.4.3-fr}
\begin{enumerate}
  \item[(i)] Tout isomorphisme local est localement de présentation finie.
  \item[(ii)] Si deux morphismes $f:X\to Y$, $g:Y\to Z$ sont localement de
  présentation finie, il en est de même de $g\circ f$.

\oldpage[IV-1]{231}

  \item[(iii)] Si $f:X\to Y$ est un $S$-morphisme localement de présentation
  finie, $f_{(S')}:X_{(S')}\to Y_{(S')}$ est localement de présentation finie
  pour toute extension $S'\to S$ du préschéma de base.
  \item[(iv)] Si $f:X\to X'$ et $g:Y\to Y'$ sont deux $S$-morphismes
  localement de présentation finie, $f\times_Sg$ est localement de
  présentation finie.
  \item[(v)] Si le composé $g\circ f$ de deux morphismes est localement de
  présentation finie et si $g$ est localement de type fini, alors $f$ est
  localement de présentation finie.
\end{enumerate}
\end{proposition}

\begin{proof}
L'assertion (i) est triviale. Pour démontrer (iii), on peut se limiter au cas
où $S=Y$ (I, 3.3.11)~; soient $x'$ un point de $X_{(S')}$, $s'$ et $x$ ses
projections sur $S'$ et $X$ respectivement, $s$ la projection commune de $s'$
et $x$ sur $S$. Par hypothèse, il y a un voisinage ouvert affine $V$ (resp.
$U$) de $s$ (resp. $x$) tels que l'image de $U$ soit contenue dans $V$ et que
$A(U)$ soit une $A(V)$-algèbre de présentation finie~; soit $V'$ un voisinage
ouvert affine de $s'$ dont l'image dans $S$ soit contenue dans $V$~; alors
$U\times_VV'$ est un voisinage ouvert affine de $x'$ dont l'image dans $S'$
est contenue dans $V'$ et dont l'anneau $A(U)\otimes_{A(V)}A(V')$ est une
$A(V')$-algèbre de présentation finie (1.4.1).

Pour démontrer (ii), considérons un point $x\in X$, et soient $y=f(x)$,
$z=g(f(x))$. Il y a par hypothèse un voisinage ouvert affine $W=\operatorname{Spec}(C)$
(resp. $V=\operatorname{Spec}(B)$) de $z$ (resp. $y$) tels que $g(V)\subset W$ et que $B$
soit une $C$-algèbre de présentation finie. D'autre part, il y a un voisinage
ouvert affine $V'=\operatorname{Spec}(B')$ (resp. $U=\operatorname{Spec}(A)$) de $y$ (resp. $x$) tels que
$f(U)\subset V'$ et que $A$ soit une $B'$-algèbre de présentation finie. Soit
$s\in B$ tel que l'ouvert $D(s)=\operatorname{Spec}(B_s)$ dans $V$ soit un voisinage de $y$
contenu dans $V'$, et soit $U'=f^{-1}(D(s))\subset U$~; on a
$U'=\operatorname{Spec}(A\otimes_{B'}B_s)$, et $A\otimes_{B'}B_s$ est une $B_s$-algèbre de
présentation finie (1.4.1)~; d'aute part $B_s=B[T]/(1-sT)$ est une
$B$-algèbre de présentation finie par définition~; par suite (1.4.1),
$A\otimes_{B'}B_s$ est une $C$-algèbre de présentation finie, ce qui prouve
(ii).

On sait que (iv) résulte de (i), (ii) et (iii) (I, 3.5.1). Pour prouver (v),
considérons le diagramme commutatif
\[
\xymatrix{
X \ar[r]^{\Gamma_f} \ar[d] & X\times_ZY \ar[d]^{f\times 1} \\
Y \ar[r]_{\Delta_g} & Y\times_ZY
}
\]
qui identifie $X$ au produit des $(Y\times_ZY)$-préschémas $Y$ et
$X\times_ZY$ (I, 5.3.7, $\Delta_g$ étant le morphisme diagonal). Si l'on sait
que $\Delta_g$ est localement de présentation finie, il en résultera par
(iii) qu'il en est de même de $\Gamma_f$. D'autre part, on a la factorisation
de $f$ en
$X\xrightarrow{\Gamma_f}X\times_ZY\xrightarrow{p_2}Y$ (I, 5.3.13), où la
projection $p_2$ est égale à $(g\circ f)\times_Z1_Y$~; comme $g\circ f$ est
supposé localement de présentation finie, il en est de même de $p_2$ par
(iv), et il résultera finalement de (ii) que $f$ est aussi localement de
présentation finie. On voit donc qu'on est ramené à démontrer le
\end{proof}

\begin{corollary}[1.4.3.1]
\label{IV.1.4.3.1-fr}
Soit $g:Y\to Z$ un morphisme localement de type fini~; alors le morphisme
diagonal $\Delta_g:Y\to Y\times_ZY$ est localement de présentation finie.
\end{corollary}

\oldpage[IV-1]{232}

\begin{proof}
On peut se borner au cas où $Z=\operatorname{Spec}(A)$, $Y=\operatorname{Spec}(B)$, $B$ étant une
$A$-algèbre de type fini~; le morphisme $\Delta_g$ correspond alors à
l'\emph{homomorphisme d'augmentation} $\delta:B\otimes_AB\to B$ tel que
$\delta(x\otimes y)=xy$. Compte tenu de la définition (1.4.2), il suffit de
montrer que le noyau $\mathfrak J$ de $\delta$ est un idéal \emph{de type
fini}, ce qui résulte de (0, 20.4.4).
\end{proof}

\begin{proposition}[1.4.4]
\label{IV.1.4.4-fr}
Soient $A$ un anneau, $B$ une $A$-algèbre, $B'$ une algèbre de polynômes
$A[T_1,\ldots,T_n]$, $u:B'\to B$ un homomorphisme surjectif de
$A$-algèbres. Pour que $B$ soit une $A$-algèbre de présentation finie, il faut
et il suffit que le noyau $\mathfrak J$ de $u$ soit un idéal de type fini de
$B'$.
\end{proposition}

\begin{proof}
La condition est suffisante par définition (1.4.1). Pour voir qu'elle est
nécessaire, remarquons que le morphisme $g:\operatorname{Spec}(B')\to\operatorname{Spec}(A)$ est
localement de type fini~; si $B$ est une $A$-algèbre de présentation finie,
le morphisme $f:\operatorname{Spec}(B)\to\operatorname{Spec}(B')$ correspondant à $u$ est tel que
$g\circ f:\operatorname{Spec}(B)\to\operatorname{Spec}(A)$ soit localement de présentation finie, donc il
résulte de (1.4.3, (v)) que $f$ est aussi localement de présentation finie.
Or, pour démontrer que l'idéal $\mathfrak J$ est de type fini, il suffit de
prouver que $\widetilde{\mathfrak J}$ est un $\widetilde{B'}$-Module de type
fini (II, 6.1.4.1). Revenant à la définition (1.4.2), on est finalement ramené
à prouver le
\end{proof}

\begin{lemma}[1.4.4.1]
\label{IV.1.4.4.1-fr}
Soient $\mathfrak a$ un idéal d'un anneau $C$, $s$ un élément de $C$, tel que
$C_s/\mathfrak a_s$ soit une $C$-algèbre de présentation finie~; alors
$\mathfrak a_s$ est un idéal de type fini dans l'anneau $C_s$.
\end{lemma}

\begin{proof}
Par hypothèse, il existe une $C$-algèbre de polynômes
$C'=C[T_1,\ldots,T_n]$ et un $C$-homomorphisme surjectif
$v:C'\to C_s/\mathfrak a_s$ dont le noyau $\mathfrak b$ est de type fini.
Soit $p:C_s\to C_s/\mathfrak a_s$ l'homomorphisme canonique~; pour tout $i$,
il y a un $t_i\in C$ tel que $p(t_i/s^k)=v(T_i)$ (on peut supposer que
l'exposant $k$ est le même pour tous les indices $i$). Considérons alors la
$C_s$-algèbre $D=C_s[T_1,\ldots,T_n]$, et l'homomorphisme de $C_s$-algèbres
$w:D\to C_s$ tel que $w(T_i)=t_i/s^k$~; $w$ est évidemment surjectif, et il
en est par suite de même de $v'=p\circ w:D\to C_s/\mathfrak a_s$. D'autre
part, tout polynôme de $D$ peut s'écrire $P/s^m$, où
$P\in C'=C[T_1,\ldots,T_n]$, et l'on a
$v'(P/s^m)=(1/s^m)v(P)$~; comme $1/s$ est inversible dans $C_s$, la relation
$v'(P/s^m)=0$ dans $C_s/\mathfrak a_s$ équivaut à $v(P)=0$, et par suite le
noyau $\mathfrak b'$ de $v'$ est engendré par l'image canonique de
$\mathfrak b$ dans $D$, et \emph{a fortiori} est de type fini dans l'anneau
$D$. Mais
$\mathfrak b'=v'^{-1}(0)=w^{-1}(p^{-1}(0))=w^{-1}(\mathfrak a_s)$, et comme
$w$ est surjectif, $\mathfrak a_s=w(\mathfrak b')$, ce qui prouve que
$\mathfrak a_s$ est un idéal de type fini de $C_s$.
\end{proof}

\begin{corollary}[1.4.5]
\label{IV.1.4.5-fr}
Soient $X$, $Y$ deux préschémas, $j:X\to Y$ une immersion, $U$ un ouvert de
$Y$ tel que $j(X)$ soit fermé dans $U$, $\mathcal J$ l'idéal quasi-cohérent de
$\mathcal O_U$ qui définit le sous-préschéma fermé de $Y$ associé à $j$. Pour
que $j$ soit localement de présentation finie, il faut et il suffit que
$\mathcal J$ soit un $\mathcal O_U$-Module de type fini.
\end{corollary}

\par\medskip\noindent
\begin{proposition}[1.4.6]
\label{IV.1.4.6-fr}
Soit $\varphi:A\to B$ un homomorphisme d'anneaux. Pour que le morphisme
correspondant $f:\operatorname{Spec}(B)\to\operatorname{Spec}(A)$ soit localement de présentation finie,
il faut et il suffit que $B$ soit une $A$-algèbre de présentation finie.
\end{proposition}

\begin{proof}
La condition étant trivialement suffisante, il reste à prouver qu'elle est
nécessaire. Si $f$ est localement de présentation finie, il résulte déjà de
(1.3.6) que $B$ est une $A$-algèbre de type fini, donc il existe un
homomorphisme surjectif $u:B'=A[T_1,\ldots,T_n]\to B$ de $A$-algèbres~; il
résulte de (1.4.3, (v)) que l'immersion
$j:\operatorname{Spec}(B)\to\operatorname{Spec}(B')$ est localement

\oldpage[IV-1]{233}

de présentation finie~; donc (1.4.5), si $\mathfrak J$ est le noyau de $u$,
le $\widetilde{B'}$-Module $\widetilde{\mathfrak J}$ est de type fini, et par
suite (II, 6.1.4.1) $\mathfrak J$ est un idéal de type fini.
\end{proof}

\begin{proposition}[1.4.7]
\label{IV.1.4.7-fr}
Soit $\rho:A\to B$ un homomorphisme d'anneaux faisant de $B$ une $A$-algèbre
finie. Pour que $B$ soit une $A$-algèbre de présentation finie, il faut et il
suffit que $B$ soit un $A$-module de présentation finie.
\end{proposition}

Nous utiliserons le

\begin{lemma}[1.4.7.1]
\label{IV.1.4.7.1-fr}
Si $B$ est une $A$-algèbre finie, il existe un $A$-homomorphisme surjectif
d'algèbres $u:B'\to B$, où $B'$ est une $A$-algèbre finie et de présentation
finie qui est un $A$-module libre.
\end{lemma}

\begin{proof}
En effet, il y a un système fini $(a_i)_{1\leq i\leq m}$ de générateurs de la
$A$-algèbre $B$ tel que chaque $a_i$ vérifie une relation $F_i(a_i)=0$, où
$F_i\in A[X]$ est un polynôme unitaire de degré $>0$~; la $A$-algèbre quotient
$B'_i=A[X]/(F_i)$ est donc libre de rang fini sur $A$~; soit $c_i$ la classe de
$X$ dans $B'_i$. Il existe un $A$-homomorphisme d'algèbres $u_i:B'_i\to B$
tel que $u_i(c_i)=a_i$~; il suffit alors de prendre pour $B'$ le produit
tensoriel $B'_1\otimes_AB'_2\otimes_A\cdots\otimes_AB'_m$ et de considérer
l'homomorphisme $u_1\otimes u_2\otimes\cdots\otimes u_m:B'\to B$, qui est
surjectif par construction. En outre, si $B''=A[T_1,\ldots,T_m]$ et si $b'_i$
désigne l'image canonique de $c_i$ dans $B'$, le $A$-homomorphisme d'algèbres
$v:B''\to B'$ tel que $v(T_i)=b'_i$ ($1\leq i\leq m$) est surjectif, et son
noyau $\mathfrak b'$ est l'idéal de $B''$ engendré par les $F_i(T_i)$, donc de
type fini~; ceci prouve que $B'$ est une $A$-algèbre de présentation finie.
\end{proof}

\begin{proof}
Ce lemme étant démontré, gardons les mêmes notations, et soit
$w=v\circ u:B''\to B$. Si $\mathfrak b$ (resp. $\mathfrak a$) est le noyau de
$w$ (resp. $u$), on a $\mathfrak a=v(\mathfrak b)$ puisque $v$ est surjectif.
Or (1.4.4), si $B$ est une $A$-algèbre de présentation finie, $\mathfrak b$
est un idéal de type fini de $B''$, donc $\mathfrak a$ est un idéal de type
fini de $B'$, donc un $A$-module de type fini puisque $B'$ est une $A$-algèbre
finie~; comme $B'$ est un $A$-module libre, $B$ est un $A$-module de
présentation finie. Inversement, si $B'$ est un $A$-module de présentation
finie, $\mathfrak a$ est un $A$-module de type fini (Bourbaki,
\emph{Alg. comm.}, chap. I, \S~2, n\textsuperscript{o}~8, lemme 9) et
\emph{a fortiori} un idéal de type fini de $B'$~; par suite, $B$ est par
définition une $B'$-algèbre de présentation finie, et comme $B'$ est une
$A$-algèbre de présentation finie, $B$ est une $A$-algèbre de présentation
finie.
\end{proof}

\subsection{Morphismes de type fini}
\label{subsection:IV.1.5-fr}

\begin{proposition}[1.5.1]
\label{IV.1.5.1-fr}
Soient $f:X\to Y$ un morphisme, $(U_\alpha)$ un recouvrement de $Y$ formé
d'ouverts affines. Les conditions suivantes sont équivalentes~:
\begin{enumerate}
  \item[a)] $f$ est localement de type fini et quasi-compact.
  \item[b)] Pour tout $\alpha$, $f^{-1}(U_\alpha)$ est réunion finie d'ouverts
  affines $V_{\alpha i}$ tels que chaque anneau $A(V_{\alpha i})$ soit une
  $A(U_\alpha)$-algèbre de type fini.
  \item[c)] Pour tout ouvert affine $U$ de $Y$, $f^{-1}(U)$ est réunion finie
  d'ouverts affines $V_j$ tels que les anneaux $A(V_j)$ soient des
  $A(U)$-algèbres de type fini.
\end{enumerate}
\end{proposition}

Pour la démonstration, voir (I, 6.3.2 et 6.6.3).

\begin{definition}[1.5.2]
\label{IV.1.5.2-fr}
On dit qu'un morphisme $f$ est \emph{de type fini} s'il vérifie les conditions
équivalentes de la proposition (1.5.1).
\end{definition}

\oldpage[IV-1]{234}

\begin{proposition}[1.5.3]
\label{IV.1.5.3-fr}
Soit $f:X\to Y$ un morphisme de type fini~; si $Y$ est noethérien, il en est
de même de $X$.
\end{proposition}

Pour la démonstration, voir (I, 6.3.7).

\begin{proposition}[1.5.4]
\label{IV.1.5.4-fr}
\begin{enumerate}
  \item[(i)] Soit $j:X\to Y$ un morphisme d'immersion. Si $j$ est
  quasi-compact (ce qui est le cas si $j$ est une immersion fermée, ou si
  l'espace sous-jacent à $X$ est noethérien, ou si l'espace sous-jacent à $Y$
  est localement noethérien), $j$ est de type fini.
  \item[(ii)] Le composé de deux morphismes de type fini est de type fini.
  \item[(iii)] Si $f:X\to Y$ est un $S$-morphisme de type fini,
  $f_{(S')}:X_{(S')}\to Y_{(S')}$ est de type fini pour toute extension
  $S'\to S$ du préschéma de base.
  \item[(iv)] Si $f:X\to X'$ et $g:Y\to Y'$ sont deux $S$-morphismes de type
  fini, $f\times_Sg$ est de type fini.
  \item[(v)] Si le composé $g\circ f$ de deux morphismes est de type fini et
  si $g$ est quasi-séparé ou si $X$ est noethérien, $f$ est de type fini.
  \item[(vi)] Si un morphisme $f$ est de type fini, il en est de même de
  $f_{\mathrm{red}}$.
\end{enumerate}
\end{proposition}

Cela résulte aussitôt (sauf le cas (v) où $X$ est noethérien, démontré dans
(I, 6.3.6)) des définitions et de (1.1.2), (1.2.4) et (1.3.4).

\begin{corollary}[1.5.5]
\label{IV.1.5.5-fr}
Soit $f:X\to Y$ un morphisme de type fini. Pour tout morphisme $Y'\to Y$ tel
que $Y'$ soit noethérien, $X\times_YY'$ est noethérien.
\end{corollary}

Cela résulte de (1.5.4, (iii)) et de (1.5.3).

\begin{corollary}[1.5.6]
\label{IV.1.5.6-fr}
Soit $X$ un préschéma de type fini sur un préschéma localement noethérien $S$.
Alors tout $S$-morphisme $f:X\to Y$ est de type fini.
\end{corollary}

Pour la démonstration, voir (I, 6.3.9).

\begin{proposition}[1.5.7]
\label{IV.1.5.7-fr}
Soit $\varphi:A\to B$ un homomorphisme d'anneaux. Pour que le morphisme
correspondant $f:\operatorname{Spec}(B)\to\operatorname{Spec}(A)$ soit de type fini, il faut et il suffit
que $\varphi$ fasse de $B$ une $A$-algèbre de type fini.
\end{proposition}

Pour la démonstration, voir (I, 6.3.3).

\subsection{Morphismes de présentation finie}
\label{subsection:IV.1.6-fr}

\begin{definition}[1.6.1]
\label{IV.1.6.1-fr}
Soient $X$, $Y$ deux préschémas. On dit qu'un morphisme $f:X\to Y$ est
\emph{de présentation finie} s'il vérifie les trois conditions suivantes~:
\begin{enumerate}
  \item[(i)] $f$ est localement de présentation finie~;
  \item[(ii)] $f$ est quasi-compact (ce qui, lorsque (i) est vérifiée,
  équivaut à dire que $f$ est de type fini (1.5.2))~;
  \item[(iii)] $f$ est quasi-séparé.
\end{enumerate}
On dit alors aussi que $X$ est \emph{de présentation finie sur $Y$}, ou un
\emph{$Y$-préschéma de présentation finie}.
\end{definition}

La condition (iii) est automatiquement vérifiée si $f$ est séparé, ou si $X$
est localement noethérien (1.2.8)~; lorsque $Y$ est localement noethérien, il
revient au même de dire que $f$ est \emph{de présentation finie} ou que $f$
est \emph{de type fini} (cette dernière condition entraînant que $X$ est
localement noethérien (1.3.3)).

\oldpage[IV-1]{235}

\begin{proposition}[1.6.2]
\label{IV.1.6.2-fr}
\begin{enumerate}
  \item[(i)] Toute immersion quasi-compacte et localement de présentation
  finie (en particulier toute immersion ouverte quasi-compacte) est de
  présentation finie.
  \item[(ii)] Le composé de deux morphismes de présentation finie est de
  présentation finie.
  \item[(iii)] Soient $X$, $Y$ deux $S$-préschémas, $f:X\to Y$ un
  $S$-morphisme de présentation finie. Alors, pour tout morphisme $S'\to S$,
  le morphisme $f_{(S')}:X_{(S')}\to Y_{(S')}$ est de présentation finie.
  \item[(iv)] Soient $f:X\to Y$, $f':X'\to Y'$ deux $S$-morphismes de
  présentation finie~; alors
  $f\times_Sf':X\times_SX'\to Y\times_SY'$ est de présentation finie.
  \item[(v)] Si le composé $g\circ f$ de deux morphismes est de présentation
  finie et si $g$ est quasi-séparé et localement de présentation finie, $f$
  est de présentation finie.
\end{enumerate}
\end{proposition}

Cela résulte immédiatement de (1.1.2), (1.2.2), (1.2.4) et (1.4.3).

Il résulte en particulier de (1.6.2, (iii)) que si $f$ est un morphisme de
présentation finie et $U$ un ouvert de $Y$, la restriction
$f^{-1}(U)\to U$ de $f$ est un morphisme de présentation finie. Inversement,
soit $(U_\alpha)$ un recouvrement de $Y$ par des ouverts affines, et supposons
que les restrictions $f^{-1}(U_\alpha)\to U_\alpha$ de $f$ soient des
morphismes de présentation finie~; il en résulte alors que $f$ est de
présentation finie, car $f$ est évidemment localement de présentation finie,
quasi-compact en vertu de (1.1.1) et quasi-séparé en vertu de (1.2.6). On peut
donc dire que la propriété pour un morphisme $f:X\to Y$ d'être de présentation
finie est \emph{locale sur $Y$}.

Si $X$ est un préschéma quasi-séparé, tout morphisme $f:X\to Y$ est
quasi-séparé (1.2.3). Donc, si $f$ est quasi-compact et localement de
présentation finie, $f$ est alors de présentation finie.

En particulier~:

\begin{corollary}[1.6.3]
\label{IV.1.6.3-fr}
Soit $\varphi:A\to B$ un homomorphisme d'anneaux. Pour que le morphisme
correspondant $f:\operatorname{Spec}(B)\to\operatorname{Spec}(A)$ soit de présentation finie, il faut et
il suffit que $\varphi$ fasse de $B$ une $A$-algèbre de présentation finie.
\end{corollary}

La nécessité de la condition résulte en effet de (1.4.6).

\begin{remark}[1.6.4]
\label{IV.1.6.4-fr}
Dans la définition (1.6.1), la condition (iii) n'est pas conséquence des deux
autres. Soit par exemple $Y$ un schéma affine dont l'espace sous-jacent ne
soit pas noethérien, et soit $U$ un ouvert non quasi-compact dans $Y$. Soit
$X$ le préschéma obtenu en recollant deux préschémas $Y_1$, $Y_2$ isomorphes à
$Y$ le long des ouverts $U_1$, $U_2$ correspondant à $U$ (0$_{\mathrm I}$,
4.1.7), de sorte que $X$ est réunion de deux ouverts isomorphes respectivement
à $Y_1$ et $Y_2$ (et que nous identifions à ces derniers), avec
$Y_1\cap Y_2=U$. En outre, il y a un morphisme $f:X\to Y$ qui coïncide dans
$Y_i$ avec l'isomorphisme donné $Y_i\to Y$ ($i=1,2$). Il est clair que $f$
vérifie la condition (i) de (1.6.1), étant un isomorphisme local~; il vérifie
aussi (ii), l'image réciproque d'un ouvert quasi-compact de $Y$ étant la
réunion de ses images dans $Y_1$ et $Y_2$~; mais comme $Y_1\cap Y_2$ n'est pas
quasi-compact, $f$ n'est pas quasi-séparé (1.2.7).
\end{remark}

\begin{proposition}[1.6.5]
\label{IV.1.6.5-fr}
Soient $X'$, $X''$ deux préschémas, $X=X'\amalg X''$ leur somme,
$f':X'\to Y$, $f'':X''\to Y$ deux morphismes, $f:X\to Y$ le morphisme qui
coïncide avec $f'$ dans $X'$ et avec $f''$ dans $X''$. Pour que $f$ soit de
présentation finie, il faut et il suffit que $f'$ et $f''$ le soient.
\end{proposition}

\par\medskip\noindent
\begin{proof}
Il suffit de montrer que pour que $f$ possède une des trois propriétés de la
définition

\oldpage[IV-1]{236}

(1.6.1), il faut et il suffit que $f'$ et $f''$ la possèdent~; c'est évident
pour la propriété d'être localement de présentation finie, qui est locale sur
$X$~; pour la propriété d'être quasi-compact, cela a été vu en (1.1.6) et pour
la propriété d'être quasi-séparé, en (1.2.9).
\end{proof}

\subsection{Améliorations de résultats antérieurs}
\label{subsection:IV.1.7-fr}

Nous allons donner dans ce numéro une liste de propositions démontrées dans
les chapitres précédents, dont l'énoncé peut être amélioré à l'aide des
nouvelles conditions de finitude introduites ci-dessus.

\begin{env}[1.7.1]
\label{IV.1.7.1-fr}
Dans les énoncés de (I, 6.4.2, 6.4.3 et 6.4.9), on peut, au lieu de supposer
que $X$ et $Y$ sont de type fini sur le corps $K$, supposer seulement qu'ils
sont \emph{localement de type fini} sur $K$~; pour (I, 6.4.2), il suffit
d'observer que tout point d'un préschéma $X$ qui est fermé dans tout ouvert
affine de $X$ est fermé dans $X$, et un schéma affine localement de type fini
sur $K$ est \emph{ipso facto} de type fini (1.5.1). Pour (I, 6.4.9), on utilise
(1.3.7) et (1.3.4, (v)).
\end{env}

\begin{env}[1.7.2]
\label{IV.1.7.2-fr}
Dans (I, 6.5.1, (ii)), on peut au lieu de supposer que $S$ est localement
noethérien et $Y$ de type fini sur $S$, supposer seulement que $Y$ est
\emph{localement de présentation finie} sur $S$, comme le montre aussitôt la
démonstration (en appliquant la définition (1.4.2)). De même dans (I, 6.5.4,
(ii)) et (I, 6.5.5, (ii)) il suffit de supposer que $f$ est un $S$-morphisme,
$Y$ étant localement de présentation finie sur $S$.
\end{env}

\begin{env}[1.7.3]
\label{IV.1.7.3-fr}
Dans (I, 7.1.11), pour la deuxième assertion, au lieu de supposer $S$
localement noethérien et $Y$ de type fini sur $S$, on peut seulement supposer
$Y$ localement de présentation finie sur $S$ (la démonstration dépendant de
(I, 6.5.1, (ii)))~; de même dans (I, 7.1.12 à 7.1.15).
\end{env}

\begin{env}[1.7.4]
\EGAReaderTarget{ega4.c04.statement.1-7-4}{IV.1.7.4-fr}
Dans (I, 9.2.2), on peut remplacer les trois hypothèses par la seule hypothèse
(entraînée par chacune des autres) que $f$ est \emph{quasi-compact et
quasi-séparé}, en vertu de (1.2.6) et (1.2.7). On notera que dans (I, 9.2.1),
l'hypothèse signifie exactement que $f$ est quasi-compact et quasi-séparé.
\end{env}

\begin{env}[1.7.5]
\label{IV.1.7.5-fr}
Dans (I, 9.3.1, 9.3.2 et 9.3.3), on peut remplacer les deux hypothèses sur $X$
par l'hypothèse (moins forte) que $X$ est un préschéma \emph{quasi-compact et
quasi-séparé}, la démonstration utilisant exactement ces deux propriétés (en
vertu de (1.2.7)).
\end{env}

\begin{env}[1.7.6]
\label{IV.1.7.6-fr}
Dans (I, 9.3.4 et 9.3.5), il suffit de supposer $X$ quasi-compact,
$\mathcal F$ et $\mathcal J$ quasi-cohérents et de type fini.
\end{env}

\begin{env}[1.7.7]
\label{IV.1.7.7-fr}
Dans (I, 9.4.7), on peut affaiblir l'hypothèse b) en supposant seulement $X$
\emph{quasi-séparé}, puisqu'il suffit seulement que l'immersion canonique
$U\to X$ soit quasi-compacte (1.2.7). On peut faire la même modification dans
(I, 9.4.8, 9.4.9 et 9.4.10).
\end{env}

\begin{env}[1.7.8]
\label{IV.1.7.8-fr}
Dans (I, 9.5.1), les conditions indiquées entre parenthèses dans l'énoncé
peuvent être remplacées par la condition que $f$ soit quasi-compact et
quasi-séparé.
\end{env}

\begin{env}[1.7.9]
\label{IV.1.7.9-fr}
Dans (I, 9.6.5), on peut remplacer les hypothèses sur $X$ par l'hypothèse
moins forte que $X$ est quasi-compact et quasi-séparé, comme le montre la
démonstration,

\oldpage[IV-1]{237}

qui n'utilise que (I, 9.4.7). Dans (I, 9.6.6), l'hypothèse «~$X$ est un
schéma quasi-compact~» peut être remplacée par «~$X$ est un préschéma
quasi-compact et quasi-séparé~».
\end{env}

\begin{env}[1.7.10]
\label{IV.1.7.10-fr}
Dans (II, 1.7.15), on peut supposer seulement que $Y$ est quasi-compact et
quasi-séparé, la démonstration n'utilisant que (I, 9.6.5).
\end{env}

\begin{env}[1.7.11]
\label{IV.1.7.11-fr}
Dans la seconde assertion de (II, 3.4.5), on peut substituer aux deux
hypothèses sur $Y$ l'hypothèse plus faible que $Y$ est quasi-compact et
quasi-séparé. De même dans (II, 3.4.8).
\end{env}

\begin{env}[1.7.12]
\label{IV.1.7.12-fr}
Dans (II, 3.8.5), on peut de même supposer seulement que $Y$ est quasi-compact
et quasi-séparé, cette hypothèse intervenant par (I, 9.4.7).
\end{env}

\begin{env}[1.7.13]
\label{IV.1.7.13-fr}
Dans (II, 4.4.1, 4.4.6 et 4.4.7), il suffit de supposer que $Y$ est
\emph{quasi-compact et quasi-séparé}, compte tenu de (1.2.3) dans la
démonstration de (II, 4.4.7). De même, dans (II, 4.4.10.1), il suffit de
supposer $Z$ quasi-compact et quasi-séparé.
\end{env}

\begin{env}[1.7.14]
\label{IV.1.7.14-fr}
Dans (II, 4.5.2) et (II, 4.5.5), il suffit de supposer $X$
\emph{quasi-compact et quasi-séparé}, cette hypothèse intervenant par (I,
9.3.1 et 9.3.2).
\end{env}

\begin{env}[1.7.15]
\label{IV.1.7.15-fr}
Dans (II, 4.6.8), il suffit encore de supposer $X$ \emph{quasi-compact et
quasi-séparé}, cette hypothèse intervenant par (I, 9.4.7).
\end{env}

\begin{env}[1.7.16]
\label{IV.1.7.16-fr}
Dans (II, 5.1.2), il suffit de supposer que $X$ est quasi-compact et
quasi-séparé (l'hypothèse intervenant par (II, 4.5.2 et 4.5.5)). De même dans
(II, 5.1.9), où on utilise (I, 9.6.5).
\end{env}

\begin{env}[1.7.17]
\label{IV.1.7.17-fr}
Dans (II, 5.2.1), il suffit toujours de supposer $X$ quasi-compact et
quasi-séparé (l'hypothèse intervenant par (II, 4.5.2)).
\end{env}

\begin{env}[1.7.18]
\label{IV.1.7.18-fr}
Dans (II, 5.2.2), il faut supposer $f$ quasi-compact et $X$ quasi-séparé~; la
démonstration actuelle est en effet insuffisante (avec les hypothèses faites),
car du fait que pour tout $\mathcal O_X$-Module quasi-cohérent $\mathcal F$ on
a $\mathrm R^1f_*(\mathcal F)=0$, il ne s'ensuit pas nécessairement, pour un
ouvert affine $U$ de $Y$, que la même propriété soit valable pour les
$(\mathcal O_X|f^{-1}(U))$-Modules quasi-cohérents, si on ne sait pas qu'un tel
Module est la \emph{restriction} d'un $\mathcal O_X$-Module quasi-cohérent (la
même remarque s'appliquant aux conditions b) et c'))~; lorsque $f$ est
quasi-compact et $X$ quasi-séparé, ce prolongement est possible en vertu de
(I, 9.4.7), modifié plus haut (1.7.7).
\end{env}

\begin{env}[1.7.19]
\label{IV.1.7.19-fr}
Dans (II, 5.3.2), il suffit de supposer $Y$ \emph{quasi-compact et
quasi-séparé} (ce qui intervient par (II, 4.4.6 et 4.4.7)). De même dans (II,
5.5.3, (ii)).
\end{env}

\begin{env}[1.7.20]
\label{IV.1.7.20-fr}
Dans (III, 1.2.2, 1.2.3 et 1.2.4), on peut remplacer les deux hypothèses sur
$X$ par l'hypothèse moins forte que $X$ est un préschéma \emph{quasi-compact
et quasi-séparé}, la démonstration n'utilisant que (I, 9.3.3).
\end{env}

\begin{env}[1.7.21]
\label{IV.1.7.21-fr}
Dans (III, 1.4.10 à 1.4.14), on peut remplacer «~séparé~» par
«~quasi-séparé~», cette hypothèse servant uniquement à pouvoir appliquer (I,
9.2.2) (voir ci-dessus (1.7.4)). De même, dans (III, 1.4.15), il suffit de
supposer que le morphisme $u$ est \emph{quasi-séparé et quasi-compact}.
\end{env}

\begin{env}[1.7.22]
\label{IV.1.7.22-fr}
Dans (III, 2.1.3), on peut encore remplacer les hypothèses par l'hypothèse
moins forte que $X$ est un préschéma quasi-compact et quasi-séparé, le
raisonnement étant le même que dans (III, 1.2.2).
\end{env}

\subsection{Morphismes de présentation finie et ensembles constructibles}
\label{subsection:IV.1.8-fr}

\oldpage[IV-1]{238}

\begin{env}[1.8.1]
\label{IV.1.8.1-fr}
Soit $X$ un préschéma~; on sait que tout ouvert quasi-compact dans $X$ est
réunion finie d'ouverts affines, et réciproquement. Pour qu'un ouvert $U$ de
$X$ soit \emph{rétrocompact} dans $X$ (0\textsubscript{III}, 9.1.1), il faut
et il suffit donc que pour tout \emph{ouvert affine} $V$ de $X$, $U\cap V$
soit quasi-compact. Lorsque $X$ est \emph{quasi-séparé} (1.2.1), tout ouvert
quasi-compact dans $X$ est rétrocompact dans $X$ (1.2.7), donc toute partie
localement constructible de $X$ est rétrocompacte dans $X$
(0\textsubscript{III}, 9.1.13)~; si \emph{de plus} $X$ est
\emph{quasi-compact}, il y a identité entre parties constructibles et parties
localement constructibles dans $X$ (0\textsubscript{III}, 9.1.12), et entre
parties ouvertes constructibles et parties ouvertes quasi-compactes
(0\textsubscript{III}, 9.1.5).
\end{env}

\begin{proposition}[1.8.2]
\label{IV.1.8.2-fr}
Soient $X$, $Y$ deux préschémas, $f:X\to Y$ un morphisme. Pour toute partie
constructible (resp. localement constructible) $Z$ de $Y$, $f^{-1}(Z)$ est
une partie constructible (resp. localement constructible) de $X$.
\end{proposition}

\begin{proof}
Supposons $Z$ localement constructible~; pour tout $x\in X$, il y aura un
voisinage ouvert affine $V$ de $f(x)$ tel que $Z\cap V$ soit constructible
dans $V$~; si la proposition est démontrée pour les parties constructibles,
$f^{-1}(Z)\cap f^{-1}(V)$ sera constructible dans $f^{-1}(V)$, donc
$f^{-1}(Z)$ sera localement constructible. Il suffit par suite de considérer
le cas où $Z$ est constructible, et compte tenu de (0\textsubscript{III},
9.1.3), on est ramené au cas où $Z$ est ouvert et \emph{rétrocompact} dans
$Y$, autrement dit tel que l'injection canonique $Z\to Y$ soit un morphisme
quasi-compact~; il suffit alors de voir que $f^{-1}(Z)$ est
\emph{rétrocompact dans $X$}, autrement dit tel que l'injection canonique
$f^{-1}(Z)\to X$ soit un morphisme quasi-compact~; mais cela résulte de
(1.1.2, (iii)), puisque $f^{-1}(Z)=Z\times_YX$ (I, 4.4.1).
\end{proof}

\begin{lemma}[1.8.3]
\label{IV.1.8.3-fr}
Soient $X$ un préschéma quasi-compact et quasi-séparé, $Z$ une partie
constructible de $X$. Il existe alors un schéma affine $X'$ et un morphisme
de présentation finie $f:X'\to X$ tels que $f(X')=Z$.
\end{lemma}

\begin{proof}
Comme $X$ est quasi-compact, il est réunion finie d'ouverts affines $X_i$
($i\in I$), et comme $X$ est quasi-séparé, il résulte de (1.2.7) que chacune
des immersions canoniques $g_i:X_i\to X$ est de présentation finie~; on en
conclut que si $f_i:X'_i\to X_i$ est un morphisme de présentation finie, il
en est de même de $g_i\circ f_i:X'_i\to X$ (1.6.2), et si $X'$ est le
préschéma somme des $X'_i$, le morphisme $h:X'\to X$ qui coïncide avec
$g_i\circ f_i$ dans chaque $X'_i$ est lui aussi de présentation finie
(1.6.5). Comme $Z\cap X_i$ est constructible dans $X_i$
(0\textsubscript{III}, 9.1.8), on voit qu'on est ramené à prouver le lemme
lorsque $X_i$ est \emph{affine} d'anneau $A$. Comme $Z$ est alors réunion
finie d'ensembles de la forme $U\cap\complement V$, où $U$ et $V$ sont des
ouverts quasi-compacts (0\textsubscript{III}, 9.1.3), on voit, en considérant
encore une somme convenable de préschémas, qu'on peut se ramener au cas où
$Z=U\cap\complement V$~; comme d'ailleurs $U$ est réunion finie d'ouverts
affines, on peut même se borner au cas où $U$ est \emph{affine}. Comme $V$
est quasi-compact, il est réunion finie d'ouverts de la forme $X_{f_i}$, où
$f_i\in A$~; soit $\mathfrak J$ l'idéal de $A$ engendré par les $f_i$, et
soit $Z'$ le sous-schéma fermé affine de $X$ qu'il définit (et qui est par
construction de présentation finie sur $X$)~; on a par définition
$V=\complement Z'$, donc $U\cap\complement V=U\cap Z'$. Considérons le
schéma affine $X'=Z'\times_XU$

\oldpage[IV-1]{239}

et soit $f:X'\to X$ le morphisme structural~; on vient de voir que
l'immersion canonique $Z'\to X$ est de présentation finie, et il en est de
même de l'immersion ouverte $U\to X$, qui est quasi-compacte puisque $U$ et
$X$ sont affines~; on conclut donc de (1.6.2, (iv)) que $f$ est de
présentation finie, et l'on a $f(X')=Z'\cap U$ (I, 3.4.8), ce qui achève la
démonstration.
\end{proof}

\begin{theorem}[1.8.4]
\label{IV.1.8.4-fr}
\textup{(Chevalley).} --- Soit $f:X\to Y$ un morphisme de présentation
finie. Pour toute partie localement constructible $Z$ de $X$, $f(Z)$ est
localement constructible dans $Y$.
\end{theorem}

\begin{proof}
Soient $y\in Y$ et $V$ un voisinage ouvert affine de $y$~; comme $f$ est
quasi-compact et quasi-séparé, il en est de même de sa restriction
$f^{-1}(V)\to V$, donc $f^{-1}(V)$ est un préschéma quasi-compact et
quasi-séparé (1.2.3)~; comme $Z\cap f^{-1}(V)$ est constructible dans
$f^{-1}(V)$ (1.8.1), on voit qu'on peut se borner au cas où $Y$ est affine
et $Z$ constructible~; $X$ est alors quasi-compact et quasi-séparé, donc
(1.8.3), il existe un morphisme de présentation finie $g:X'\to X$ tel que
$Z=g(X')$~; comme $f\circ g$ est de présentation finie, on voit qu'on est
ramené au cas où $Z=X$~; autrement dit, il suffira de prouver le
\end{proof}

\begin{lemma}[1.8.4.1]
\label{IV.1.8.4.1-fr}
Soient $Y$ un schéma affine, $f:X\to Y$ un morphisme quasi-compact et
localement de présentation finie~; alors $f(X)$ est une partie constructible
de $Y$.
\end{lemma}

\begin{proof}
Comme $X$ est quasi-compact, il est réunion finie d'ouverts affines~; on peut
donc se borner au cas où $Y=\operatorname{Spec}(A)$,
$X=\operatorname{Spec}(B)$, $B$ étant une $A$-algèbre de présentation finie
(1.4.6). Or, on a le
\end{proof}

\begin{lemma}[1.8.4.2]
\label{IV.1.8.4.2-fr}
Soient $A_0$ un anneau, $(A_\lambda)_{\lambda\in L}$ un système inductif de
$A_0$-algèbres, $A=\varinjlim A_\lambda$~; si $B$ est une $A$-algèbre de
présentation finie, il existe un $\lambda$ et une $A_\lambda$-algèbre de
présentation finie $B_\lambda$ tels que $B$ soit isomorphe à
$B_\lambda\otimes_{A_\lambda}A$.
\end{lemma}

\begin{proof}
Par hypothèse, $B$ est isomorphe à une algèbre de la forme
$A[T_1,\ldots,T_n]/\mathfrak J$, où les $T_i$ sont des indéterminées et
$\mathfrak J$ un idéal de type fini~; soit $(F_j)$ ($1\leq j\leq m$) un
système de générateurs de $\mathfrak J$. Soit
$\varphi_\lambda:A_\lambda\to A$ l'homomorphisme canonique. Comme $L$ est
filtrant, il existe $\lambda$ tel que chacun des coefficients de chacun des
polynômes $F_j$ soit l'image par $\varphi_\lambda$ d'un élément de
$A_\lambda$. Autrement dit, il existe un système de polynômes
$(F_{j\lambda})_{1\leq j\leq m}$ de $A_\lambda[T_1,\ldots,T_n]$ tels que
$F_j=\varphi_\lambda(F_{j\lambda})$ pour $1\leq j\leq m$. Si
$\mathfrak J_\lambda$ est l'idéal de $A_\lambda[T_1,\ldots,T_n]$ engendré
par les $F_{j\lambda}$, $\mathfrak J$ est l'image de
$\mathfrak J_\lambda\otimes_{A_\lambda}A$ dans
$A[T_1,\ldots,T_n]=A_\lambda[T_1,\ldots,T_n]\otimes_{A_\lambda}A$~;
l'anneau $B_\lambda=A_\lambda[T_1,\ldots,T_n]/\mathfrak J_\lambda$ répond à
la question.
\end{proof}

On appliquera ici ce lemme en considérant $A$ comme limite inductive de ses
sous-$\mathbf Z$-algèbres \emph{de type fini}. On voit donc qu'il existe une
telle sous-$\mathbf Z$-algèbre $A_0$ de $A$, et une $A_0$-algèbre de
présentation finie $B_0$ telles que $B$ soit isomorphe à
$B_0\otimes_{A_0}A$~; posons $X_0=\operatorname{Spec}(B_0)$,
$Y_0=\operatorname{Spec}(A_0)$, de sorte que l'on a $X=X_0\times_{Y_0}Y$,
$f$ étant la projection $X\to Y$~; soient $f_0:X_0\to Y_0$,
$g_0:Y\to Y_0$ les morphismes structuraux~; il résulte de (I, 3.4.8) que
l'on a $f(X)=g_0^{-1}(f_0(X_0))$~; compte tenu de (1.8.2), il suffit donc de
montrer que $f_0(X_0)$ est constructible. Autrement dit, on est finalement
ramené à démontrer le

\begin{corollary}[1.8.5]
\label{IV.1.8.5-fr}
Soient $Y$ un préschéma noethérien, $f:X\to Y$ un morphisme de type fini.
Pour toute partie constructible $Z$ de $X$, $f(Z)$ est une partie
constructible de $Y$.
\end{corollary}

\oldpage[IV-1]{240}

\begin{proof}
On se ramène comme ci-dessus à prouver que $f(X)$ est constructible.
Appliquons le critère (0\textsubscript{III}, 9.2.3)~: il faut prouver que pour
toute partie fermée irréductible $T$ de $Y$, $T\cap f(X)=f(f^{-1}(T))$ est
rare dans $T$ ou contient un ouvert non vide de $T$~; en prenant un
sous-préschéma de $Y$ ayant $T$ pour espace sous-jacent (I, 5.2.1) et en
considérant son image réciproque par $f$, on voit finalement (compte tenu de
(1.5.4, (iii)) qu'on est ramené à prouver que si l'on suppose $Y$
\emph{irréductible et $f$ dominant}, $f(X)$ contient un ouvert non vide de
$Y$. On peut en outre supposer $Y$ affine~; alors $X$ est réunion finie
d'ouverts affines $V_i$, et comme $f(X)$ est dense dans $Y$, il en est de
même d'un au moins des $f(V_i)$. On peut donc aussi supposer $X$ affine~; si
$(X_j)$ est la famille (finie) des composantes irréductibles de $X$, on voit
comme ci-dessus qu'un au moins des $f(X_j)$ est dense dans $Y$~; on peut donc
supposer $X$ irréductible. Enfin, en remplaçant $f$ par $f_{\mathrm{red}}$
(1.5.4, (vi)) on peut supposer $X$ et $Y$ \emph{intègres}. Alors
$X=\operatorname{Spec}(B)$, $Y=\operatorname{Spec}(A)$, où $A$ est un
anneau intègre noethérien, $B$ une $A$-algèbre intègre de type fini contenant
$A$ (I, 1.2.7). Il suffit de montrer qu'il existe $g\in A$ tel que, pour tout
$y\in D(g)$ (c'est-à-dire tel que $g\notin\mathfrak p_y$) l'idéal
$\mathfrak p_y$ soit l'intersection de $A$ et d'un idéal premier de $B$, car
cela montrera que $D(g)\subset f(X)$. Finalement il suffit de prouver le
\end{proof}

\begin{lemma}[1.8.5.1]
\label{IV.1.8.5.1-fr}
Soient $A$ un anneau intègre, $B$ une $A$-algèbre intègre de type fini. Il
existe $g\in A$ tel que tout homomorphisme de $A$ dans un corps
algébriquement clos $\Omega$, non nul en $g$, se prolonge en un homomorphisme
de $B$ dans $\Omega$.
\end{lemma}

\begin{proof}
Or c'est là un résultat classique d'algèbre commutative (Bourbaki,
\emph{Alg. comm.}, chap. V, \S~3, n\textsuperscript{o}~1, cor. 3 du th. 1).
\end{proof}

\begin{corollary}[1.8.6]
\label{IV.1.8.6-fr}
Soient $Y$ un préschéma irréductible, $\eta$ son point générique,
$f:X\to Y$ un morphisme localement de type fini. Si $f^{-1}(\eta)$ n'est pas
vide, il existe un voisinage ouvert $U$ de $\eta$ dans $Y$ tel que
$f^{-1}(y)$ soit non vide pour tout $y\in U$. Si $\mathscr F$ est un
$\mathcal O_X$-Module quasi-cohérent de type fini, et si
$\mathscr F_\eta=\mathscr F\otimes_{\mathcal O_Y}k(\eta)$ n'est pas nul,
alors il existe un voisinage ouvert $U'$ de $\eta$ dans $Y$ tel que
$\mathscr F_y=\mathscr F\otimes_{\mathcal O_Y}k(y)$ ne soit pas nul pour
tout $y\in U'$.
\end{corollary}

\begin{proof}
Si $p_y$ est la projection canonique de la fibre
$f^{-1}(y)=X\times_Y\operatorname{Spec}(k(y))$ dans $X$, on a
$\mathscr F_y=p_y^*(\mathscr F)$, donc
$\operatorname{Supp}(\mathscr F_y)=p_y^{-1}(\operatorname{Supp}(\mathscr F))
=\operatorname{Supp}(\mathscr F)\cap f^{-1}(y)$ en vertu de (I, 9.1.13.1) et
(I, 3.6.1), puisque $\mathscr F$ est de type fini~; en outre
$\operatorname{Supp}(\mathscr F)$ est fermé dans $X$ (0\textsubscript{I},
5.2.2), et si $Z$ est un sous-préschéma fermé de $X$ ayant pour espace
sous-jacent $\operatorname{Supp}(\mathscr F)$ (I, 5.2.1), et $j$ l'immersion
canonique $Z\to X$, $f\circ j$ est localement de type fini (1.3.4)~; cela
prouve que la première assertion entraîne la seconde. Pour prouver la
première assertion, remarquons d'abord que l'on peut supposer $X$ et $Y$
réduits (1.3.4, (vi)), autrement dit $Y$ intègre. Soit $\xi$ un point de
$f^{-1}(\eta)$~; on peut remplacer $X$ et $Y$ par des voisinages ouverts
affines de $\xi$ et $\eta$ respectivement, et par suite supposer que
$X=\operatorname{Spec}(B)$, $Y=\operatorname{Spec}(A)$, $B$ étant une
$A$-algèbre de type fini. En outre, si $Z$ est le sous-préschéma fermé réduit
de $X$ ayant $\overline{\{\xi\}}$ pour ensemble sous-jacent (I, 5.2.1), on
peut, comme ci-dessus, remplacer $X$ par $Z$, autrement dit supposer $B$
\emph{intègre} (I, 5.1.4). Comme le morphisme $f$ est alors dominant,
l'homomorphisme correspondant $\varphi:A\to B$ est injectif (I, 1.2.7)~;
donc le corollaire résulte finalement du lemme (1.8.5.1).
\end{proof}

\oldpage[IV-1]{241}

\begin{proposition}[1.8.7]
\label{IV.1.8.7-fr}
Soient $S$ un préschéma, $g:X\to S$, $h:Y\to S$ deux morphismes de
présentation finie, $f:X\to Y$ un $S$-morphisme. Pour tout $s\in S$, posons
$X_s=g^{-1}(s)=X\times_S\operatorname{Spec}(k(s))$,
$Y_s=h^{-1}(s)=Y\times_S\operatorname{Spec}(k(s))$,
$f_s=f\times 1:X_s\to Y_s$. Alors l'ensemble des $s\in S$ tels que $f_s$
soit surjectif (resp. radiciel) est localement constructible.
\end{proposition}

\begin{proof}
Soit $E$ l'ensemble des $s\in S$ tels que $f_s$ soit surjectif~; l'ensemble
$S-E$ est égal à $h(Y-f(X))$~; or $f$ est de présentation finie (1.6.2, (v))
donc $f(X)$ est localement constructible dans $Y$ (1.8.4)~; comme $h$ est de
présentation finie, $h(Y-f(X))$ est localement constructible dans $S$, donc
aussi $E$.

Pour montrer que l'ensemble $E'$ des $s\in S$ tels que $f_s$ soit radiciel
est localement constructible, nous utiliserons le
\end{proof}

\begin{lemma}[1.8.7.1]
\EGAReaderTarget{ega4.c04.lem.1-8-7-1}{IV.1.8.7.1-fr}
Soient $U$, $V$ deux préschémas~; pour qu'un morphisme $f:U\to V$ soit
radiciel, il faut et il suffit que le morphisme diagonal
$\Delta_f:U\to U\times_VU$ soit surjectif. Par suite, tout morphisme radiciel
est séparé.
\end{lemma}

\begin{proof}
En effet, $\Delta_f$ étant une immersion (I, 5.3.9), est un morphisme
localement de type fini (1.3.4)~; dire que $\Delta_f$ est surjectif signifie
donc que pour tout corps algébriquement clos $K$, l'application
correspondante $U(K)\to(U\times_VU)(K)=U(K)\times_{V(K)}U(K)$ est surjective
(1.3.7 et I, 3.4.2.1)~; mais par définition d'un produit fibré d'ensembles,
cela signifie que l'application $U(K)\to V(K)$ correspondant à $f$ est
injective, et cela équivaut à dire que $f$ est radiciel (I, 3.5.5).

Cela étant, $\Delta_{f_s}$ se déduit de $\Delta_f:X\to X\times_YX$ par le
changement de base $\operatorname{Spec}(k(s))\to S$ (I, 5.3.4). Comme $f$
est de présentation finie, il en est de même du morphisme structural
$X\times_YX\to Y$ (1.6.2, (iv)), donc aussi de $\Delta_f$ (1.6.2, (v))~; il
suffit donc d'appliquer la première partie de la proposition à $\Delta_f$,
en utilisant le lemme (1.8.7.1).
\end{proof}

\subsection{Ensembles pro-constructibles et ind-constructibles}
\label{subsection:IV.1.9-fr}

\begin{lemma}[1.9.1]
\label{IV.1.9.1-fr}
Soit $(A_\alpha)_{\alpha\in I}$ un système inductif d'anneaux dont l'ensemble
d'indices $I$ est filtrant à droite, et soit $A=\varinjlim A_\alpha$. Pour que
$A=0$, il faut et il suffit qu'il existe un indice $\gamma$ tel que
$A_\gamma=0$ (et l'on a alors $A_\beta=0$ pour tout $\beta\geq\gamma$).
\end{lemma}

\begin{proof}
En effet, pour tout $\alpha\in I$, l'homomorphisme canonique
$\varphi_\alpha:A_\alpha\to A$ transforme l'élément unité en élément unité~;
dire que $A=0$ signifie que $\varphi_\alpha(1)=0$, ou encore
$\varphi_\alpha(1)=\varphi_\alpha(0)$~; on sait que cela équivaut à dire
qu'il existe un indice $\gamma\geq\alpha$ tel que l'homomorphisme
$\varphi_{\gamma\alpha}:A_\alpha\to A_\gamma$ soit tel que
$\varphi_{\gamma\alpha}(1)=\varphi_{\gamma\alpha}(0)$, et cette dernière
relation équivaut à $A_\gamma=0$, d'où le lemme.
\end{proof}

\begin{proposition}[1.9.2]
\label{IV.1.9.2-fr}
Soient $B$ un anneau, $(A_\alpha)_{\alpha\in I}$ un système inductif de
$B$-algèbres dont l'ensemble d'indices $I$ est filtrant à droite, et soit
$A=\varinjlim A_\alpha$~; posons $Y=\operatorname{Spec}(B)$,
$X_\alpha=\operatorname{Spec}(A_\alpha)$, $X=\operatorname{Spec}(A)$, et
soient $u:X\to Y$, $u_\alpha:X_\alpha\to Y$ les morphismes correspondant aux
homomorphismes structuraux $B\to A$, $B\to A_\alpha$. Alors~:
\end{proposition}

\begin{env}[1.9.2.i]
\label{IV.1.9.2.i-fr}
\emph{Pour que $X=\varnothing$, il faut et il suffit qu'il existe un
$\gamma\in I$ tel que $X_\gamma=\varnothing$, et l'on a alors
$X_\alpha=\varnothing$ pour $\alpha\geq\gamma$.}
\end{env}

\oldpage[IV-1]{242}

\begin{env}[1.9.2.ii]
\label{IV.1.9.2.ii-fr}
\emph{On a}
\[
  u(X)=\bigcap_{\alpha\in I}u_\alpha(X_\alpha). \tag{1.9.2.1}
  \label{IV.1.9.2.1-fr}
\]
\end{env}

L'assertion (i) n'est autre que la traduction de (1.9.1). Pour démontrer
(ii), notons que, puisque $u$ se factorise en
$X\to X_\alpha\xrightarrow{u_\alpha}Y$, le premier membre de (1.9.2.1) est
contenu dans le second. Inversement, soit $y\in Y-u(X)$~; on a
$u^{-1}(y)=\varnothing$, et $u^{-1}(y)$ est l'espace sous-jacent au
$k(y)$-préschéma $\operatorname{Spec}(A\otimes_Bk(y))$~; comme dans la
catégorie des $B$-modules le foncteur $\varinjlim$ commute au produit
tensoriel, $A\otimes_Bk(y)$ est limite inductive du système inductif d'anneaux
$A_\alpha\otimes_Bk(y)$~; le lemme (1.9.1) montre qu'il existe $\alpha$ tel
que $A_\alpha\otimes_Bk(y)=0$, c'est-à-dire
$u_\alpha^{-1}(y)=\varnothing$, donc $y\notin u_\alpha(X_\alpha)$. C.Q.F.D.

\begin{proposition}[1.9.3]
\label{IV.1.9.3-fr}
Soient $X$ un préschéma quasi-compact et quasi-séparé, $E$ une partie de $X$,
$(X_\alpha)_{\alpha\in I}$ un recouvrement ouvert affine de $X$, et pour
tout $\alpha\in I$, posons $E_\alpha=E\cap X_\alpha$. Les conditions
suivantes sont équivalentes~:
\begin{enumerate}
  \item[a)] Il existe un morphisme quasi-compact $f:X'\to X$ tel que
  $E=f(X')$.
  \item[a$'$)] Pour tout $\alpha$, il existe un morphisme quasi-compact
  $f_\alpha:X'_\alpha\to X_\alpha$ tel que
  $E_\alpha=f_\alpha(X'_\alpha)$.
  \item[a$''$)] Pour tout $\alpha$, il existe un schéma affine $X''_\alpha$
  et un morphisme $f'_\alpha:X''_\alpha\to X_\alpha$ tel que
  $f'_\alpha(X''_\alpha)=E_\alpha$.
  \item[b)] $E$ est intersection de parties constructibles de $X$.
  \item[b$'$)] $F=X-E$ est réunion de parties constructibles de $X$.
\end{enumerate}
\end{proposition}

\begin{proof}
Il est clair que $b)$ et $b')$ sont équivalentes. La condition $a)$ entraîne
$a')$ en prenant pour $X'_\alpha$ le préschéma induit sur l'ouvert
$f^{-1}(X_\alpha)$ de $X'$ et pour $f_\alpha$ la restriction de $f$. Pour
voir que $a')$ entraîne $a'')$, il suffit de remarquer que $X'_\alpha$ est
réunion d'ouverts affines $X''_{\alpha\lambda}$ ($\lambda\in L$)~; soit
$X''_\alpha$ le préschéma \emph{somme} des $X''_{\alpha\lambda}$
($\lambda\in L$), qui est un schéma affine~; si
$g_\alpha:X''_\alpha\to X'_\alpha$ est le morphisme coïncidant avec
l'identité dans chacun des $X''_{\alpha\lambda}$, il est clair que si l'on
pose $f'_\alpha=f_\alpha\circ g_\alpha$, on a
$E_\alpha=f'_\alpha(X''_\alpha)$ puisque $g_\alpha$ est surjectif. Pour
montrer que $a'')$ entraîne $a)$, notons qu'il y a une partie finie $J$ de
$I$ telle que les $X_\alpha$ d'indices $\alpha\in J$ recouvrent $X$~; soit
$X'$ le préschéma somme des $X''_\alpha$ pour $\alpha\in J$, et soit
$f:X'\to X$ le morphisme coïncidant avec $j_\alpha\circ f'_\alpha$ pour tout
$\alpha\in J$, où $j_\alpha:X_\alpha\to X$ est l'injection canonique. Comme
$E$ est réunion des $E\cap X_\alpha$ pour $\alpha\in J$, on a bien
$E=f(X')$ et il reste à voir que $f$ est un morphisme quasi-compact~; mais
l'hypothèse que $X$ est quasi-séparé entraîne que $j_\alpha$ est
quasi-compact (1.2.7), donc il en est de même de
$j_\alpha\circ f'_\alpha$ (1.1.1 et 1.1.2) et finalement de $f$ (1.1.6).

Reste donc à prouver l'équivalence de $a'')$ et de $b)$. Prouvons en premier
lieu que $a'')$ entraîne $b)$~: considérons une partie finie $J$ de $I$ telle
que les $X_\alpha$ pour $\alpha\in J$ recouvrent $X$~; il suffira de montrer
que, pour tout $\alpha\in J$, $E_\alpha$ est intersection de parties
constructibles $F_{\alpha\lambda}$ de $X_\alpha$ ($\lambda\in L_\alpha$)~;
en effet, tout $F_{\alpha\lambda}$ est aussi une partie constructible de $X$
en vertu de (0\textsubscript{III}, 9.1.8, (ii)), car l'hypothèse que $X$ est
quasi-séparé entraîne que tout $X_\alpha$ est rétrocompact dans $X$ (1.2.7)~;
$E$ étant réunion des $E_\alpha$ pour $\alpha\in J$, est intersection des
réunions (finies) $\bigcup_{\alpha\in J}F_{\alpha,\lambda(\alpha)}$, pour
tous les choix des $\lambda(\alpha)\in L_\alpha$~; ces réunions étant
constructibles dans $X$, cela établit notre assertion. On peut donc se borner
au cas

\oldpage[IV-1]{243}

où $X=\operatorname{Spec}(A)$ et où $E=f(X')$, où
$X'=\operatorname{Spec}(A')$ est affine, $f:X'\to X$ un morphisme
correspondant à un homomorphisme d'anneaux $A\to A'$. Notons maintenant que
$A'$ est limite inductive de la famille $(A'_\lambda)$ de ses
sous-$A$-algèbres \emph{de type fini}, ordonnée par inclusion~; si
$X'_\lambda=\operatorname{Spec}(A'_\lambda)$, $f$ se factorise en
$X'\to X'_\lambda\xrightarrow{f_\lambda}X$ et il résulte de (1.9.2, (ii))
que l'on a $f(X')=\bigcap_\lambda f_\lambda(X'_\lambda)$~; si l'on démontre
que $f_\lambda(X'_\lambda)$ est intersection de parties constructibles de
$X$, il en sera de même de $f(X')$. On peut donc se borner au cas où $A'$ est
une $A$-algèbre \emph{de type fini}. Or, on a le lemme suivant~:
\end{proof}

\begin{lemma}[1.9.3.1]
\label{IV.1.9.3.1-fr}
Soient $A$ un anneau, $A'$ une $A$-algèbre de type fini. Alors $A'$ est
$A$-isomorphe à une limite inductive filtrante $\varinjlim A''_\mu$ où les
$A''_\mu$ sont des $A$-algèbres de présentation finie (1.4.1).
\end{lemma}

\begin{proof}
En effet, on peut écrire $A'=B/\mathfrak J$, avec
$B=A[T_1,\ldots,T_n]$ et $\mathfrak J$ un idéal de $B$. Mais $\mathfrak J$
est limite inductive des idéaux \emph{de type fini} $\mathfrak J'_\mu$ de
$B$, contenus dans $\mathfrak J$, ordonnés par inclusion~; comme dans la
catégorie des $B$-modules le foncteur $\varinjlim$ est exact, on en conclut
que $A'=\varinjlim B/\mathfrak J'_\mu$ à un $A$-isomorphisme près.
\end{proof}

Le même raisonnement que ci-dessus permet alors de se ramener au cas où $A'$
est une $A$-algèbre de présentation finie~; mais alors $f(X')$ est
constructible dans $X$ en vertu du théorème de Chevalley (1.8.5), ce qui
prouve $b)$.

Démontrons enfin que $b)$ entraîne $a'')$. Si $E$ est intersection d'une
famille $(G_\lambda)_{\lambda\in L}$ de parties constructibles de $X$,
chaque $E_\alpha$ est intersection des $G_\lambda\cap X_\alpha$, qui sont
constructibles dans $X_\alpha$ (0\textsubscript{III}, 9.1.8, (i)), donc on
est ramené au cas où $X=\operatorname{Spec}(A)$ est affine.

On sait alors (1.8.3) que pour tout $\lambda\in L$ il existe un morphisme
$f_\lambda:X'_\lambda\to X$, où
$X'_\lambda=\operatorname{Spec}(A'_\lambda)$ est affine, tel que
$f_\lambda(X'_\lambda)=G_\lambda$. Il suffit donc de prouver le lemme
suivant~:

\begin{lemma}[1.9.3.2]
\label{IV.1.9.3.2-fr}
Soient $(A'_\lambda)_{\lambda\in L}$ une famille de $A$-algèbres,
$X=\operatorname{Spec}(A)$, $X'_\lambda=\operatorname{Spec}(A'_\lambda)$,
et soit $f_\lambda:X'_\lambda\to X$ le morphisme structural. Pour toute
partie finie $J$ de $L$, posons
$A'_J=\bigotimes_{\lambda\in J}A'_\lambda$ (produit tensoriel de
$A$-algèbres), $X'_J=\operatorname{Spec}(A'_J)$, et soit $f_J:X'_J\to X$ le
morphisme structural. On a alors
$f_J(X'_J)=\bigcap_{\lambda\in J}f_\lambda(X'_\lambda)$. Si
$A'=\varinjlim A'_J$, $J$ parcourant l'ensemble filtrant des parties finies
de $L$, $X'=\operatorname{Spec}(A')$, et si $f:X'\to X$ est le morphisme
structural, on a $f(X')=\bigcap_{\lambda\in L}f_\lambda(X'_\lambda)$.
\end{lemma}

\begin{proof}
La première assertion résulte de (I, 3.4.7)~; la deuxième résulte de (1.9.2,
(ii)), qui donne la relation $f(X')=\bigcap_Jf_J(X'_J)$.
\end{proof}

\begin{definition}[1.9.4]
\label{IV.1.9.4-fr}
Soit $X$ un espace topologique. On dit qu'une partie $E$ de $X$ est
\emph{pro-constructible} (resp. \emph{ind-constructible}) dans $X$ si, pour
tout $x\in X$, il existe un voisinage ouvert $U$ de $x$ dans $X$ tel que
$E\cap U$ soit intersection (resp. réunion) de parties localement
constructibles de $U$.
\end{definition}

Tenant compte de (1.2.7) et de (0\textsubscript{III}, 9.1.8 et 9.1.12), les
conditions équivalentes de la proposition (1.9.3) s'expriment donc en disant
que $E$ est \emph{pro-constructible} dans $X$, les ensembles localement
constructibles dans $X$ étant identiques aux ensembles constructibles dans
$X$ sous les hypothèses de (1.9.3).

\oldpage[IV-1]{244}

\begin{proposition}[1.9.5]
\label{IV.1.9.5-fr}
Soit $X$ un préschéma.
\begin{enumerate}
  \item[(i)] Pour qu'une partie $E$ de $X$ soit pro-constructible, il faut et
  il suffit que $X-E$ soit ind-constructible.
  \item[(ii)] Toute réunion finie et toute intersection de parties
  pro-constructibles de $X$ sont pro-constructibles. Toute intersection
  finie et toute réunion de parties ind-constructibles de $X$ sont
  ind-constructibles.
  \item[(iii)] Toute intersection (resp. réunion) de parties localement
  constructibles de $X$ est pro-constructible (resp. ind-constructible).
  Inversement, si $X$ est quasi-compact et quasi-séparé, toute partie
  pro-constructible (resp. ind-constructible) de $X$ est intersection (resp.
  réunion) de parties constructibles de $X$.
  \item[(iv)] Soit $(U_\alpha)$ un recouvrement ouvert de $X$. Pour qu'une
  partie $E$ de $X$ soit pro-constructible (resp. ind-constructible) dans
  $X$, il faut et il suffit que pour tout $\alpha$, $E\cap U_\alpha$ soit
  pro-constructible (rep. ind-constructible) dans $U_\alpha$.
  \item[(v)] Toute partie pro-constructible $E$ de $X$ est rétrocompacte
  dans $X$. Pour qu'une partie localement fermée $E$ de $X$ soit
  rétrocompacte dans $X$, il faut et il suffit qu'elle soit
  pro-constructible dans $X$.
  \item[(vi)] Soit $f:X'\to X$ un morphisme de préschémas~; pour toute partie
  pro-constructible (resp. ind-constructible) $E$ de $X$, $f^{-1}(E)$ est
  pro-constructible (resp. ind-constructible) dans $X'$.
  \item[(vii)] Soit $f:X'\to X$ un morphisme quasi-compact~; pour toute
  partie pro-constructible $E'$ de $X'$, $f(E')$ est pro-constructible dans
  $X$~; en particulier $f(X')$ est pro-constructible dans $X$.
  \item[(viii)] Soit $f:X'\to X$ un morphisme localement de présentation
  finie~; pour toute partie ind-constructible $E'$ de $X'$, $f(E')$ est
  ind-constructible dans $X$~; en particulier $f(X')$ est ind-constructible
  dans $X$.
  \item[(ix)] Supposons que $X$ soit quasi-compact et quasi-séparé~; alors,
  pour qu'une partie $E$ de $X$ soit pro-constructible, il faut et il suffit
  qu'il existe un schéma affine $X'$ et un morphisme (nécessairement
  quasi-compact) $f:X'\to X$ tels que $E=f(X')$.
\end{enumerate}
\end{proposition}

\begin{proof}
Les propriétés (i), (ii), (iv) et la première assertion de (iii) résultent
de la définition (1.9.4) et sont valables pour un espace topologique
quelconque, en utilisant la distributivité mutuelle de l'intersection et de
la réunion et le fait que si $T$ est localement constructible dans $X$,
$T\cap U$ est localement constructible dans $U$ pour tout ouvert $U$ de $X$.
Si $X$ est quasi-compact et quasi-séparé et $E$ est pro-constructible dans
$X$, il y a un recouvrement fini $(U_i)$ de $X$ formé d'ouverts affines tels
que $E\cap U_i$ soit intersection de parties constructibles $M_{i\lambda}$
de $U_i$ ($\lambda\in L_i$)~; les $M_{i\lambda}$ sont aussi constructibles
dans $X$ en vertu de (1.2.7) et (0\textsubscript{III}, 9.1.8, (ii)) et $E$ est
intersection des réunions finies $\bigcup_iM_{i,\lambda(i)}$ (où
$\lambda(i)\in L_i$ pour tout $i$), qui sont des parties constructibles dans
$X$~; cela démontre la seconde assertion de (iii). L'assertion (ix) résulte
de (iii) et de (1.9.3). Pour démontrer la première assertion de (v), on peut
se limiter au cas où $X$ est affine, et prouver alors que $E$ est
quasi-compact (0\textsubscript{III}, 9.1.1)~; mais comme $X$ est alors
quasi-séparé, il existe en vertu de (ix) un morphisme quasi-compact
$f:X'\to X$ tel que $E=f(X')$~; comme $X'$ est quasi-compact, il en est de
même de son image $E$ par une application continue.

Pour prouver (vi), on peut se borner, en vertu de (iv), au cas où $X$ est
affine~; la conclusion résulte alors de (iii) et de (1.8.2).

\oldpage[IV-1]{245}

De même, pour démontrer (vii), on peut se borner au cas où $X$ est affine~;
alors $X'$ est réunion finie d'ouverts affines $X'_i$ et $f(E')$ est réunion
des $f(E'\cap X'_i)$, si bien qu'on peut aussi supposer que $X'$ est affine,
en vertu de (ii)~; mais alors $E'=g(X'')$, où $g:X''\to X'$ est un morphisme
quasi-compact, en vertu de (ix), et l'on a $f(E')=f(g(X''))$, donc $f(E')$ est
pro-constructible puisque $g\circ f$ est quasi-compact.

On déduit de (vii) la démonstration de la seconde assertion de (v)~: en
effet, soit $E$ une partie localement fermée et rétrocompacte dans $X$, et
considérons un sous-préschéma de $X$ ayant $E$ pour espace sous-jacent (I,
5.2.1)~; l'injection canonique $j:E\to X$ est alors par hypothèse un
morphisme quasi-compact, et il résulte de (vii) que $E=j(E)$ est
pro-constructible dans $X$.

Enfin, pour prouver (viii), on peut encore supposer $X$ affine~; en outre, si
$(X'_\alpha)$ est un recouvrement ouvert affine de $X'$ tel que
$E'\cap X'_\alpha$ soit réunion de parties constructibles de $X'_\alpha$
((iii) et (iv)), $f(E')$ est réunion des $f(E'\cap X'_\alpha)$ et, en vertu
du théorème de Chevalley (1.8.4), chacune des $f(E'\cap X'_\alpha)$ est
réunion de parties constructibles de $X$, d'où la conclusion.
\end{proof}

\begin{env}[1.9.6]
\label{IV.1.9.6-fr}
\textbf{Exemples} (1.9.6). --- Toute partie \emph{finie} d'un préschéma $X$
est pro-constructible~: en effet, il suffit de considérer les parties
$\{x\}$ réduites à un seul point (1.9.5, (ii)) et $\{x\}$ est l'image de
$\operatorname{Spec}(k(x))$ par le morphisme canonique
$\operatorname{Spec}(k(x))\to X$, qui est quasi-compact~; d'où la conclusion
par (1.9.5, (vii)). Par contre, une partie $\{x\}$ réduite à un point n'est
pas nécessairement ind-constructible~; par exemple, soit $A$ un anneau
noethérien intègre ayant une infinité d'idéaux maximaux, et soit $x$ le point
générique de $X=\operatorname{Spec}(A)$~; si $\{x\}$ était ind-constructible
dans $X$, il serait constructible en vertu de (1.9.5, (iii)) et par suite
contiendrait un ouvert non vide de $X$ (0\textsubscript{III}, 9.2.2)~; mais
cela contredit l'hypothèse que $A$ possède une infinité d'idéaux maximaux, en
vertu du théorème d'Artin-Tate (0, 16.3.3).

Toute partie \emph{fermée} $Y$ d'un préschéma $X$ est pro-constructible, par
(1.9.5, (v)). Toute partie ouverte de $X$ est donc ind-constructible, mais
une partie ouverte $U$ de $X$ n'est pro-constructible que si elle est
\emph{rétrocompacte}, en vertu encore de (1.9.5, (v)).

Enfin, si $X$ est un préschéma \emph{noethérien}, donc quasi-séparé (1.2.8),
il résulte de (1.9.5, (iii)) et de (0\textsubscript{III}, 9.1.7) que les
parties ind-constructibles de $X$ sont les \emph{réunions de parties
localement fermées de $X$}.
\end{env}

\begin{theorem}[1.9.7]
\label{IV.1.9.7-fr}
Soient $X$ un préschéma quasi-compact, $E$ une partie ind-constructible de
$X$, $(F_\lambda)_{\lambda\in L}$ une famille de parties pro-constructibles
de $X$ telle que $\bigcap_{\lambda\in L}F_\lambda\subset E$~; alors il existe
une partie finie $J$ de $L$ telle que
$\bigcap_{\lambda\in J}F_\lambda\subset E$.
\end{theorem}

\begin{proof}
Notons que les ensembles $F=X-E$ et $F\cap F_\lambda$ sont
pro-constructibles, donc on est ramené au
\end{proof}

\begin{corollary}[1.9.8]
\label{IV.1.9.8-fr}
Soient $X$ un préschéma quasi-compact, $(F_\lambda)_{\lambda\in L}$ une
famille de parties pro-constructibles de $X$ dont l'intersection est vide~;
alors il existe une sous-famille finie $(F_\lambda)_{\lambda\in J}$ dont
l'intersection est vide.
\end{corollary}

\begin{proof}
Recouvrant $X$ par un nombre fini d'ouverts affines, et utilisant (1.9.5,
(iv)), on est ramené au cas où $X$ est affine~; en vertu de (1.9.5, (ix)), on
a $F_\lambda=f_\lambda(X'_\lambda)$,

\oldpage[IV-1]{246}

où $X'_\lambda=\operatorname{Spec}(A'_\lambda)$ est affine, et $f_\lambda$
est un morphisme $X'_\lambda\to X$~; alors, avec les notations de (1.9.3.2),
on a par hypothèse $f(X')=\varnothing$, donc $X'=\varnothing$~; mais cela
entraîne par (1.9.2, (i)) qu'il existe une partie finie $J$ de $L$ telle que
$X'_J=\varnothing$, donc
$f_J(X'_J)=\bigcap_{\lambda\in J}F_\lambda=\varnothing$.
\end{proof}

\begin{corollary}[1.9.9]
\label{IV.1.9.9-fr}
Soient $X$ un préschéma quasi-compact, $F$ une partie pro-constructible de
$X$, $(E_\lambda)_{\lambda\in L}$ une famille de parties ind-constructibles
de $X$ telle que $F\subset\bigcup_{\lambda\in L}E_\lambda$~; alors il existe
une partie finie $J$ de $L$ telle que
$F\subset\bigcup_{\lambda\in J}E_\lambda$. En particulier, de tout
recouvrement de $X$ par des parties ind-constructibles, on peut extraire un
recouvrement fini de $X$.
\end{corollary}

\begin{proof}
Cela résulte de (1.9.7) par passage aux complémentaires.
\end{proof}

\begin{proposition}[1.9.10]
\label{IV.1.9.10-fr}
Soit $X$ un préschéma. Pour qu'une partie $E$ de $X$ soit
ind-constructible, il est nécessaire que, pour tout $x\in X$, l'intersection
$E\cap\overline{\{x\}}$ soit un voisinage de $x$ dans
$\overline{\{x\}}$. Si $X$ est localement noethérien, cette condition est
suffisante.
\end{proposition}

\begin{proof}
Supposons $E$ ind-constructible, et soit $Y$ l'intersection de
$\overline{\{x\}}$ et d'un ouvert affine dans $X$, contenant $x$~; il y a
donc un sous-préschéma de $X$ ayant $Y$ pour espace sous-jacent (I, 5.2.1),
et si $j:Y\to X$ est l'injection canonique,
$E\cap Y=j^{-1}(E)$ est ind-constructible dans $Y$ (1.9.5, (vi)). Comme $Y$
est quasi-compact et séparé, $E\cap Y$ est donc réunion de parties
constructibles de $Y$ (1.9.5, (iii))~; par suite, il y a deux ouverts $U$,
$V$ dans $Y$ tels que
$x\in U\cap\complement V\subset E\cap Y$ (0\textsubscript{III}, 9.1.3).
Mais comme $x$ est point générique de l'espace irréductible $Y$, $V$ est
nécessairement vide et l'on a $U\subset E\cap Y$.

Supposons maintenant $X$ localement noethérien, et soit $E$ une partie de
$X$ vérifiant la condition de l'énoncé. En vertu de la définition (1.9.4),
on peut se borner au cas où $X$ est noethérien. Alors, pour tout $x\in E$,
$E\cap\overline{\{x\}}$ contient une partie non vide de la forme
$U\cap\overline{\{x\}}$, où $U$ est ouvert dans $X$~; cela montre que $E$
est réunion de parties localement fermées de $X$, donc est
ind-constructible (1.9.6).
\end{proof}

\begin{proposition}[1.9.11]
\label{IV.1.9.11-fr}
Soient $X$ un préschéma, $E$ une partie de $X$. Les propriétés suivantes sont
équivalentes~:
\begin{enumerate}
  \item[a)] $E$ est localement constructible.
  \item[b)] $E$ est ind-constructible et pro-constructible.
  \item[c)] $E$ et $X-E$ sont pro-constructibles.
  \item[d)] $E$ et $X-E$ sont ind-constructibles.
\end{enumerate}
\end{proposition}

\begin{proof}
Il suffit évidemment de prouver que $d)$ entraîne $a)$, et on peut se borner
au cas où $X$ est affine. Alors (1.9.5, (iii)), $E$ (resp. $X-E$) est réunion
d'une famille $(E_\lambda)$ (resp. $(E'_\mu)$) de parties constructibles de
$X$~; comme les $E_\lambda$ et les $E'_\mu$ forment un recouvrement de $X$,
il résulte de (1.9.9) qu'il y a des indices en nombre fini $\lambda_i$,
$\mu_j$ tels que les $E_{\lambda_i}$ et les $E'_{\mu_j}$ forment un
recouvrement de $X$~; cela implique que $E$ est réunion des
$E_{\lambda_i}$, donc est constructible.
\end{proof}

\begin{corollary}[1.9.12]
\label{IV.1.9.12-fr}
Soit $f:X\to Y$ un morphisme surjectif, qui est, soit quasi-compact, soit
localement de présentation finie. Pour qu'une partie $E$ de $Y$ soit
localement constructible (resp. pro-constructible, resp. ind-constructible),
il faut et il suffit que $f^{-1}(E)$ le soit dans $X$.
\end{corollary}

\oldpage[IV-1]{247}

\begin{proof}
On sait que la condition est nécessaire ((1.8.2) et (1.9.5, (vi)))~; pour
montrer qu'elle est suffisante, on est ramené au cas où $X$ est affine, en
vertu de (1.9.5, (iv))~; en outre, en vertu de (1.9.11), on peut se borner au
cas où $f^{-1}(E)$ est pro-constructible, \emph{ou} au cas où $f^{-1}(E)$ est
ind-constructible. Si $f$ est surjectif et quasi-compact, et $f^{-1}(E)$
pro-constructible, il résulte de (1.9.5, (vii)) que
$E=f(f^{-1}(E))$ est pro-constructible. Si $f$ est surjectif et localement de
présentation finie, et $f^{-1}(E)$ ind-constructible,
$E=f(f^{-1}(E))$ est ind-constructible par (1.9.5, (viii)), ce qui achève la
démonstration.
\end{proof}

\begin{env}[1.9.13]
\label{IV.1.9.13-fr}
Soient $X$ un préschéma, $\mathscr T$ sa topologie~; il résulte de (1.9.5,
(i) et (ii)) que les parties ind-constructibles (resp. pro-constructibles) de
$X$ sont les parties ouvertes (resp. fermées) pour une topologie sur $X$,
appelée \emph{topologie constructible} et que nous noterons
$\mathscr T^{\mathrm{cons}}$~; nous désignerons par $X^{\mathrm{cons}}$
l'ensemble $X$ muni de la topologie $\mathscr T^{\mathrm{cons}}$.
\end{env}

\begin{proposition}[1.9.14]
\label{IV.1.9.14-fr}
Soient $X$ un préschéma, $\mathscr T$ sa topologie,
$\mathscr T^{\mathrm{cons}}$ la topologie constructible sur $X$.
\begin{enumerate}
  \item[(i)] La topologie $\mathscr T^{\mathrm{cons}}$ est plus fine que
  $\mathscr T$.
  \item[(ii)] Les parties localement constructibles de $X$ sont identiques
  aux parties à la fois ouvertes et fermées de l'espace $X^{\mathrm{cons}}$.
  \item[(iii)] Pour tout morphisme $f:X\to Y$, l'application sous-jacente de
  $X^{\mathrm{cons}}$ dans $Y^{\mathrm{cons}}$ est continue~; on la note
  $f^{\mathrm{cons}}$.
  \item[(iv)] Si le morphisme $f:X\to Y$ est quasi-compact,
  $f^{\mathrm{cons}}$ est une application fermée~; en particulier, si $f$
  est quasi-compact et bijectif, $f^{\mathrm{cons}}$ est un homéomorphisme.
  \item[(v)] Si un morphisme $f:X\to Y$ est localement de présentation
  finie, $f^{\mathrm{cons}}$ est une application ouverte~; en particulier,
  si $f$ est bijectif et localement de présentation finie,
  $f^{\mathrm{cons}}$ est un homéomorphisme.
  \item[(vi)] Pour tout ouvert $U$ de $X$, la topologie induite par
  $\mathscr T^{\mathrm{cons}}$ sur $U$ est identique à la topologie de
  $U^{\mathrm{cons}}$.
\end{enumerate}
\end{proposition}

\begin{proof}
En effet, (i) résulte de ce que tout ouvert pour $\mathscr T$ est un ouvert
pour $\mathscr T^{\mathrm{cons}}$ (1.9.6), et (ii) est une traduction de
(1.9.11)~; (iii), (iv) et (v) traduisent respectivement (1.9.5), (vi), (vii)
et (viii). Enfin pour démontrer (vi), il suffit de remarquer que l'injection
canonique $j:U\to X$ est localement de présentation finie (1.4.3), donc
$j^{\mathrm{cons}}:U^{\mathrm{cons}}\to X^{\mathrm{cons}}$ est une
application continue ouverte.
\end{proof}

\begin{proposition}[1.9.15]
\label{IV.1.9.15-fr}
Soit $X$ un préschéma.
\begin{enumerate}
  \item[(i)] Pour que $X^{\mathrm{cons}}$ soit quasi-compact, il faut et il
  suffit que $X$ soit quasi-compact.
  \item[(ii)] Pour que $X^{\mathrm{cons}}$ soit séparé, il faut et il suffit
  que $X$ soit quasi-séparé, et alors $X^{\mathrm{cons}}$ est localement
  compact et totalement discontinu.
  \item[(iii)] Pour que $X^{\mathrm{cons}}$ soit compact, il faut et il suffit
  que $X$ soit quasi-compact et quasi-séparé.
  \item[(iv)] Tout point de $X$ admet, pour la topologie
  $\mathscr T^{\mathrm{cons}}$, un voisinage ouvert et compact.
  \item[(v)] Pour qu'un morphisme $f:X\to Y$ soit quasi-compact, il faut et il
  suffit que l'application continue
  $f^{\mathrm{cons}}:X^{\mathrm{cons}}\to Y^{\mathrm{cons}}$ soit propre.
\end{enumerate}
\end{proposition}

\begin{proof}
(i) Comme $\mathscr T^{\mathrm{cons}}$ est plus fine que $\mathscr T$, il est
clair que si $X^{\mathrm{cons}}$ est quasi-compact, il en est de même de $X$~;
la réciproque résulte de (1.9.9).

(ii) Supposons $X$ quasi-séparé, et montrons que $X^{\mathrm{cons}}$ est
totalement discontinu~: en effet, si $x$, $y$ sont deux points distincts de
$X$ et si par exemple $x\notin\overline{\{y\}}$, il existe un
\oldpage[IV-1]{248}
voisinage ouvert affine $U$ de $X$ ne contenant pas $y$~; comme $X$ est
quasi-séparé, $U$ est rétrocompact dans $X$ (1.2.7), donc $U$ et $X-U$ sont
localement constructibles dans $X$, et par suite ouverts dans
$X^{\mathrm{cons}}$ en vertu de (1.9.11), d'où notre assertion. Comme sur tout
ouvert affine $U$ de $X$, la topologie induite par $\mathscr T^{\mathrm{cons}}$
est celle de $U^{\mathrm{cons}}$ en vertu de (1.9.14, (vi)), il résulte de ce
qui précède que $X^{\mathrm{cons}}$ est localement compact, puisque $U$ est
ouvert dans $X^{\mathrm{cons}}$ et que $U^{\mathrm{cons}}$ est compact. Reste à
prouver que si $X^{\mathrm{cons}}$ est séparé, alors $X$ est quasi-séparé~;
considérons en effet un ouvert affine $U$ de $X$~; la topologie induite sur
$U$ par $\mathscr T^{\mathrm{cons}}$ est celle de $U^{\mathrm{cons}}$
(1.9.14, (vi)), donc $U$ est compact pour cette topologie induite, puisque
$U^{\mathrm{cons}}$ est compact par la première partie du raisonnement. Si
$V$ est un second ouvert affine dans $X$, $U\cap V$ est donc une partie
compacte de l'espace séparé $X^{\mathrm{cons}}$, étant intersection de deux
parties compactes de cet espace~; comme la topologie induite par
$\mathscr T^{\mathrm{cons}}$ sur $U\cap V$ est celle de
$(U\cap V)^{\mathrm{cons}}$ (1.9.14, (vi)) et que l'application identique
$(U\cap V)^{\mathrm{cons}}\to U\cap V$ est continue, on en conclut que
$U\cap V$ est quasi-compact pour la topologie induite par $\mathscr T$~; $X$
est par suite quasi-séparé en vertu de (1.2.7).

(iii) est un corollaire immédiat de (i) et (ii).

(iv) Pour tout $x\in X$, un voisinage ouvert affine $U$ de $x$ pour
$\mathscr T$ est aussi un voisinage de $x$ pour $\mathscr T^{\mathrm{cons}}$
et il est compact en vertu de (iii) et de (1.9.14, (vi)), la topologie induite
sur $U$ par $\mathscr T^{\mathrm{cons}}$ étant identique à celle de
$U^{\mathrm{cons}}$.

(v) Supposons $f$ quasi-compact~; alors on sait déjà (1.9.14, (iv)) que
$f^{\mathrm{cons}}$ est une application fermée. Soit d'autre part $y$ un
point de $Y$, et posons
\[
  Z=f^{-1}(y)=X\times_Y\operatorname{Spec}(k(y));
\]
$Z$ est quasi-compact (1.1.2, (iii)) et comme le morphisme canonique
$p:Z\to X$ est injectif, il résulte de (i) et de ce que l'application
$p^{\mathrm{cons}}$ est continue, que la topologie induite sur $f^{-1}(y)$
par celle de $X^{\mathrm{cons}}$ fait de $f^{-1}(y)$ un espace quasi-compact.
Cela prouve que $f^{\mathrm{cons}}$ est une application propre (Bourbaki,
\emph{Top. gén.}, chap. I\textsuperscript{er}, 3\textsuperscript{e} éd.,
\S~10, n\textsuperscript{o}~2, th.~1). Inversement, supposons que
l'application continue $f^{\mathrm{cons}}$ soit propre, et soit $V$ un ouvert
quasi-compact de $Y$~; si $h:V\to Y$ est l'injection canonique,
$h^{\mathrm{cons}}:V^{\mathrm{cons}}\to Y^{\mathrm{cons}}$ est continue et
injective et $V^{\mathrm{cons}}$ est quasi-compact par (i), donc la topologie
induite sur $V$ par celle de $Y^{\mathrm{cons}}$ fait de $V$ un espace
quasi-compact. L'hypothèse que $f^{\mathrm{cons}}$ est propre entraîne alors
que la topologie induite sur $f^{-1}(V)$ par celle de $X^{\mathrm{cons}}$ fait
de $f^{-1}(V)$ un espace quasi-compact (\emph{loc. cit.}, prop.~6), donc
$f^{-1}(V)$ est aussi un sous-espace quasi-compact de $X$, ce qui montre que
le morphisme $f$ est quasi-compact.
\end{proof}

\begin{env}[1.9.16]
\label{IV.1.9.16-fr}
Nous montrerons plus tard comment on peut, pour tout préschéma $X$, munir
l'espace $X^{\mathrm{cons}}$ d'un faisceau d'anneaux qui en fait un préschéma
dont les anneaux locaux sont tous des corps, identiques aux corps résiduels aux
points de $X$. De tels préschémas s'introduisent par exemple de façon naturelle
dans la traduction, en langage des schémas, des constructions de
Néron [32] dans sa théorie de la réduction des variétés abéliennes.
\end{env}

\oldpage[IV-1]{249}

\subsection{Application aux morphismes ouverts}
\label{subsection:IV.1.10-fr}

\begin{theorem}[1.10.1]
\label{IV.1.10.1-fr}
Soient $X$ un préschéma, $E$ une partie ind-constructible de $X$ (1.9.4), $x$
un point de $X$. Pour que $x$ soit intérieur à $E$, il faut et il suffit que
toute générisation ($0_{\mathrm I}$, 2.1.2) $x'$ de $x$ appartienne à $E$,
autrement dit (I, 2.4.2) que l'on ait
$\operatorname{Spec}(\mathcal O_{X,x})\subset E$.
\end{theorem}

\begin{proof}
La condition est évidemment nécessaire, tout voisinage de $x$ contenant les
générisations de $x$. Pour voir qu'elle est suffisante, on peut évidemment se
borner au cas où $X=\operatorname{Spec}(A)$ est affine, $x$ étant un idéal premier
$\mathfrak p$ de $A$. On sait alors (1.9.5, (ix)) qu'il existe un schéma affine
$X'=\operatorname{Spec}(A')$ et un morphisme $f:X'\to X$ tels que
$f(X')=X-E$. Posons $Y=\operatorname{Spec}(\mathcal O_{X,x})$ et
$Y'=X'\times_XY$~; l'hypothèse signifie (en
vertu de (I, 3.6.5)) que l'on a $Y'=\varnothing$, autrement dit
$A'\otimes_AA_{\mathfrak p}=0$. Comme
$A_{\mathfrak p}=\varinjlim A_t$, où $t$ parcourt $A-\mathfrak p$
($0_{\mathrm I}$, 1.4.5), on a $\varinjlim A_t'=0$, le foncteur
$\varinjlim$ commutant au produit tensoriel. Par suite (1.9.1), il existe un
$t\in A-\mathfrak p$ tel que $A_t'=0$, et $\mathrm D(t)$ est alors un
voisinage ouvert de $x$ dans $X$ contenu dans $E$.
\end{proof}

\begin{corollary}[1.10.2]
\label{IV.1.10.2-fr}
Soient $f:X\to Y$ un morphisme localement de présentation finie, et soit
$y\in Y$. Pour que $y$ soit intérieur à $f(X)$ il faut et il suffit que toute
générisation $y'$ de $y$ appartienne à $f(X)$.
\end{corollary}

\begin{proof}
On sait en effet (1.9.5, (viii)) que $f(X)$ est ind-constructible dans $Y$ et
il suffit d'appliquer (1.10.1).
\end{proof}

Nous dirons qu'une application continue $f:X\to Y$ est \emph{ouverte en un
point} $x\in X$ si l'image par $f$ de \emph{tout} voisinage de $x$ dans $X$
est un voisinage de $f(x)$ dans $Y$.

\begin{proposition}[1.10.3]
\EGAReaderTarget{ega4.c04.prop.1-10-3}{IV.1.10.3-fr}
Soient $f:X\to Y$ un morphisme, $x$ un point de $X$, $y=f(x)$. Considérons
les conditions suivantes~:
\begin{enumerate}
  \item[a)] $f$ est ouvert au point $x$.
  \item[b)] Pour toute générisation $y'$ de $y$, il existe une générisation
  $x'$ de $x$ telle que $f(x')=y'$.
  \item[b')] L'image par $f$ de $\operatorname{Spec}(\mathcal O_{X,x})$ est
  $\operatorname{Spec}(\mathcal O_{Y,y})$.
  \item[c)] Pour toute partie fermée irréductible $Y'$ de $Y$ contenant $y$,
  il existe une composante irréductible de $X'=f^{-1}(Y')$ contenant $x$ et
  dominant $Y'$.
\end{enumerate}
Alors on a les implications
$\mathrm{a)}\Rightarrow\mathrm{b)}\Leftrightarrow\mathrm{b')}
\Leftrightarrow\mathrm{c)}$. Si en outre $f$ est localement de présentation
finie, les quatre conditions sont équivalentes.
\end{proposition}

\begin{proof}
Par définition d'une générisation, l'image par $f$ de
$\operatorname{Spec}(\mathcal O_{X,x})$ est toujours contenue dans
$\operatorname{Spec}(\mathcal O_{Y,y})$, donc b) et b') sont équivalentes. Il
est immédiat
que c) entraîne b), car si $y'$ est une générisation de $y$,
$Y'=\overline{\{y'\}}$ est une partie irréductible de $Y$ contenant $y$~; si
$x'$ est le point générique d'une composante irréductible de $f^{-1}(Y')$ qui
contient $x$ et domine $Y'$, on a $f(x')=y'$ ($0_{\mathrm I}$, 2.1.5) et $x'$
est une générisation de $x$. Inversement, b) entraîne c)~: en effet, soit
$y'$ le point générique de $Y'$ et soit $x'$ une générisation de $x$ telle
que $f(x')=y'$~; soit $Z'$ une composante irréductible de $f^{-1}(Y')$
contenant $x'$ (donc $x$)~; comme $f(x')$ est déjà dense dans $Y'$, il en est
ainsi \emph{a fortiori} de $f(Z')$.

Pour voir que a) entraîne b'), on peut évidemment se borner au cas où
$X=\operatorname{Spec}(B)$ et $Y=\operatorname{Spec}(A)$ sont affines, $x$
étant un idéal premier
$\mathfrak q$ de $B$, de sorte que
$\mathcal O_{X,x}=B_{\mathfrak q}=\varinjlim B_t$,
\oldpage[IV-1]{250}
où $t$ parcourt $B-\mathfrak q$. On a, par suite (1.9.2),
\[
  f(\operatorname{Spec}(\mathcal O_{X,x}))
  =\bigcap_t f(\operatorname{Spec}(B_t))=\bigcap_t f(\mathrm D(t)),
\]
et par hypothèse, pour tout $t\in B-\mathfrak q$, $y$ est intérieur à
$f(\mathrm D(t))$, donc $f(\mathrm D(t))$ contient
$\operatorname{Spec}(\mathcal O_{Y,y})$~; d'où
$\operatorname{Spec}(\mathcal O_{Y,y})
\subset f(\operatorname{Spec}(\mathcal O_{X,x}))$, ce qui établit
notre assertion. Enfin, lorsque $f$ est localement de présentation finie, b)
implique a) en vertu de (1.10.2) appliqué à la restriction $U\to Y$ de $f$ à
un voisinage ouvert arbitraire $U$ de $x$.
\end{proof}

\begin{corollary}[1.10.4]
\label{IV.1.10.4-fr}
Soit $f:X\to Y$ un morphisme. Considérons les conditions suivantes~:
\begin{enumerate}
  \item[a)] $f$ est ouvert.
  \item[b)] Pour tout $x\in X$ et toute générisation $y'$ de $y=f(x)$, il
  existe une générisation $x'$ de $x$ telle que $f(x')=y'$.
  \item[b')] Pour tout $x\in X$, l'image par $f$ de
  $\operatorname{Spec}(\mathcal O_{X,x})$ est
  $\operatorname{Spec}(\mathcal O_{Y,f(x)})$.
  \item[c)] Pour toute partie fermée irréductible $Y'$ de $Y$, toute
  composante irréductible de $f^{-1}(Y')$ domine $Y'$.
\end{enumerate}
Alors on a les implications
$\mathrm{a)}\Rightarrow\mathrm{b)}\Leftrightarrow\mathrm{b')}
\Leftrightarrow\mathrm{c)}$. Si en outre $f$ est localement de présentation
finie, les quatre conditions sont équivalentes.
\end{corollary}

\begin{proof}
Dire que $f$ est ouvert signifie que $f$ est ouvert en tout point $x\in X$,
donc les implications
$\mathrm{a)}\Rightarrow\mathrm{b)}\Leftrightarrow\mathrm{b')}$ résultent des
implications analogues dans (1.10.3), ainsi que l'implication
$\mathrm{b)}\Rightarrow\mathrm{a)}$ lorsque $f$ est localement de présentation
finie. L'implication $\mathrm{c)}\Rightarrow\mathrm{b)}$ résulte aussi de
l'implication analogue dans (1.10.3)~; prouvons enfin que b) entraîne c). En
effet, soit $x'$ le point générique d'une composante
irréductible de $f^{-1}(Y')$, et montrons que $y'=f(x')$ est le point générique
de $Y'$. Soit $y''$ une générisation de $y'$~; il existe par hypothèse une
générisation $x''$ de $x'$ telle que $f(x'')=y''$, et comme
$x''\in f^{-1}(Y')$, on a nécessairement $x''=x'$, donc $y''=y'$, ce qui
achève la démonstration.
\end{proof}

\begin{center}
\emph{(A suivre.)}
\end{center}

\clearpage
\oldpage[IV-1]{251}
\phantomsection
\addcontentsline{toc}{section}{Bibliographie (suite)}
\section*{BIBLIOGRAPHIE (\emph{suite})}

\begingroup
\small
\setlength{\parindent}{0pt}
\setlength{\parskip}{3pt}
[30] M.~Nagata, \emph{Local rings}, \emph{Interscience Tracts}, 13 (1962),
Interscience, New York.

[31] M.~Nagata, \emph{A Jacobian criterion of simple points},
\emph{Ill. Journ. of Math.}, 1 (1957), p.~427--432.

[32] A.~Néron, \emph{Modèles minimaux des variétés abéliennes sur les corps
locaux et globaux}, \emph{Publ. Inst. Hautes Études Sci.}, n\textsuperscript{o}
21 (1964).

\clearpage
\oldpage[IV-1]{252}
\phantomsection
\addcontentsline{toc}{section}{Index des notations}
\section*{INDEX DES NOTATIONS}

$\dim_c(X)$, $\dim(X)$ ($X$ espace topologique) : 0, 14.1.2.\par
$\dim_x(X)$ ($X$ espace topologique, $x$ point de $X$) : 0, 14.1.2.\par
$\operatorname{codim}(Y,X)$ ($X$ espace topologique, $Y$ partie fermée de $X$) : 0, 14.2.1.\par
$\operatorname{codim}_x(Y,X)$ ($X$ espace topologique, $Y$ partie fermée de $X$, $x$ point de $X$) : 0, 14.2.4.\par
$N[T_1,\ldots,T_n]$ ($N$ $A$-module) : 0, 15.1.1.\par
$\dim(A)$ ($A$ anneau) : 0, 16.1.1.\par
$\operatorname{ht}(\mathfrak a)$ ($\mathfrak a$ idéal) : 0, 16.1.3.\par
$\dim_A(M)$, $\dim(M)$ ($M$ $A$-module) : 0, 16.1.7.\par
$P_{\mathfrak q}(M,n)$ ($M$ module de type fini sur un anneau semi-local noethérien $A$, $\mathfrak q$ idéal de définition de $A$) : 0, 16.2.1.\par
$\operatorname{prof}_A(M)$, $\operatorname{prof}(M)$ ($M$ $A$-module) : 0, 16.4.5.\par
$\operatorname{coprof}_A(M)$, $\operatorname{coprof}(M)$ ($M$ $A$-module) : 0, 16.4.9.\par
$\operatorname{dim.proj}_A(M)$, $\operatorname{dim.proj}(M)$,
$\operatorname{dim.inj}_A(M)$, $\operatorname{dim.inj}(M)$ ($M$ $A$-module) : 0, 17.2.1.\par
$\operatorname{dim.coh}(A)$ ($A$ anneau) : 0, 17.2.8.\par
$\operatorname{dim.proj}(\mathscr F)$, $\operatorname{dim.inj}(\mathscr F)$
($\mathscr F$ Module sur un espace annelé) : 0, 17.2.14.\par
$\operatorname{dim.coh}(X)$ ($X$ espace annelé) : 0, 17.2.14.\par
$E\times_B F$ ($B$, $E$, $F$ $A$-anneaux) : 0, 18.1.2.\par
$D_B(L)$ ($B$ anneau, $L$ $B$-bimodule) : 0, 18.2.3.\par
$Y\oplus_X Z$ ($X$, $Y$, $Z$ $A$-bimodules) : 0, 18.2.7.\par
$\operatorname{Exan}_A(B,L)$ ($B$ $A$-anneau, $L$ $B$-bimodule) : 0, 18.3.4.\par
$\operatorname{Exan}_{A/A'}(B,L)$ : 0, 18.3.7.\par
$\operatorname{Exal}_A(B,L)$, $\operatorname{Exal}_{A/A'}(B,L)$ ($A$ anneau commutatif, $B$ $A$-algèbre, $L$ $B$-bimodule) : 0, 18.4.1.\par
$\operatorname{Exalcom}_A(B,L)$, $\operatorname{Exalcom}_{A/A'}(B,L)$ ($A$ anneau commutatif, $B$ $A$-algèbre commutative, $L$ $B$-module) : 0, 18.4.2.\par
$\mathrm H_A^2(B,L)$, $\mathrm H_A^2(B,L)^s$ : 0, 18.4.3.\par
$\operatorname{Exantop}_A(B,L)$, $\operatorname{Exaltop}_A(B,L)$,
$\operatorname{Exalcotop}_A(B,L)$ : 0, 18.5.1.\par
$\operatorname{Exantop}_{A/A'}(B,L)$, $\operatorname{Exaltop}_{A/A'}(B,L)$,
$\operatorname{Exalcotop}_{A/A'}(B,L)$ : 0, 18.5.2.\par
$\operatorname{Hom.\ cont}_C(M,N)$ ($M$, $N$ $C$-bimodules topologiques) : 0, 18.5.3.\par
$\operatorname{D\acute{e}r}_A(B,L)$ ($B$ $A$-anneau, $L$ $B$-bimodule) : 0, 20.1.2.\par
$\operatorname{D\acute{e}r.\ cont}_A(B,L)$ : 0, 20.3.1.\par
$p_{B/A}$, $\mathfrak J_{B/A}$ ($B$ $A$-algèbre topologique commutative) : 0, 20.4.1.\par
$P^1_{B/A}$, $\varepsilon_{B/A}$ ($B$ $A$-algèbre topologique commutative) : 0, 20.4.2.\par
$\Omega^1_{B/A}$, $\widehat{\Omega}^1_{B/A}$ ($B$ $A$-algèbre topologique commutative) : 0, 20.4.3.\par
$\Omega^1_B$ ($B$ anneau topologique) : 0, 20.4.3.\par
$d_{B/A}$, $d(x)$, $dx$, $d_B$ : 0, 20.4.6.\par
$\Upsilon_{C/B/A}$, $\Upsilon_{C/B}$ ($A$, $B$, $C$ anneaux) : 0, 20.6.1.\par
$\chi_E$ ($E$ $B$-extension $A$-triviale de $C$ par $L$) : 0, 20.6.8.\par

\clearpage
\oldpage[IV-1]{253}
$\Upsilon^C_{B/A/\Lambda}$, $\Upsilon^C_{B/A}$ ($\Lambda$, $A$, $B$, $C$ anneaux) : 0, 20.6.14.\par
$\chi_{B/A}$, $\chi_B$ ($B$ $A$-algèbre) : 0, 20.6.24.\par
$\chi_E$ ($E$ $B$-algèbre topologique) : 0, 20.7.10.\par
$F_A$, $A^p$, $A^{(p)}$, $E^{(p)}$ ($A$ anneau de caractéristique $p>0$, $E$ $A$-module) : 0, 21.1.4.\par
$\pi_{B/A}$, $\Theta_{B/A}$, $\Xi_{B/A}$ ($A$, $B$ anneaux de caractéristique $p>0$) : 0, 21.3.2.\par
$\Delta(L/K,k/k_0)$, $d(L/K,k/k_0)$, $\Delta(L/K,k)$, $d(L/K,k)$ ($k_0$, $k$, $K$, $L$ corps) : 0, 21.6.1.\par
$\chi'_{B/A}$ ($A$, $B$ anneaux de caractéristique $p>0$) : 0, 22.4.6.\par
$\mathscr T^{\mathrm{cons}}$, $X^{\mathrm{cons}}$ ($X$ préschéma, $\mathscr T$ topologie de $X$) : IV, 1.9.12.\par
$f^{\mathrm{cons}}$ ($f$ morphisme de préschémas) : IV, 1.9.13.\par

\clearpage
\oldpage[IV-1]{254}
\phantomsection
\addcontentsline{toc}{section}{Index terminologique}
\section*{INDEX TERMINOLOGIQUE}

Algèbre augmentée des parties principales d'ordre 1 : 0, 20.4.2.\par
$A$-algèbre de Cohen : 0, 19.8.1.\par
Algèbre de présentation finie : IV, 1.4.1.\par
$A$-algèbre essentiellement de type fini : IV, 1.3.8.\par
$A$-anneau : 0, 18.1.1.\par
$A$-anneau augmenté sur $B$ : 0, 18.1.4.\par
Anneau biéquidimensionnel : 0, 16.1.4.\par
Anneau caténaire : 0, 16.1.4.\par
Anneau de caractéristique $p$ : 0, 21.1.1.\par
Anneau de Cohen : 0, 19.8.4.\par
Anneau de Cohen-Macaulay : 0, 16.5.1 et 16.5.13.\par
Anneau équicodimensionnel, équidimensionnel : 0, 16.1.4.\par
Anneau local géométriquement unibranche : 0, 23.1.7.\par
Anneau japonais : 0, 23.1.1.\par
Anneau local premier, anneau local complet premier : 0, 19.8.3.\par
Anneau régulier : 0, 17.1.1 et 17.3.6.\par
Anneau local unibranche : 0, 23.1.7.\par
Anneau universellement japonais : 0, 23.1.1.\par
Application ouverte en un point : IV, 1.10.2.\par
Augmentation : 0, 18.1.4.\par

Bimorphisme formel : 0, 19.1.2.\par

Chaîne dans un ensemble ordonné, chaîne de parties fermées irréductibles d'un espace topologique : 0, 14.1.1.\par
Chaîne saturée : 0, 14.3.1.\par
Codimension combinatoire, codimension d'une partie fermée d'un espace topologique : 0, 14.2.1.\par
Codimension en un point d'une partie fermée d'un espace topologique : 0, 14.2.4.\par
Condition des chaînes : 0, 14.3.2.\par
Coprofondeur d'un module sur un anneau local noethérien : 0, 16.4.9.\par
Corps $k_0$-admissible, corps admissible pour une extension : 0, 21.6.1.\par
Corps de multiplicité radicielle finie sur un sous-corps : 0, 19.6.6.\par

Défaut de $k_0$-admissibilité d'une extension : 0, 21.6.1.\par
Dérivation, $A$-dérivation : 0, 20.1.2.\par
Différentielles, 1-différentielles, différentielles absolues : 0, 20.4.3.\par
Différentielle extérieure : 0, 20.4.6.\par
Dimension combinatoire, dimension d'un espace topologique : 0, 14.1.2.\par
Dimension d'un espace topologique en un point : 0, 14.1.2.\par
Dimension de Krull, dimension d'un anneau : 0, 16.1.1.\par
Dimension d'un module : 0, 16.1.7.\par
Dimension cohomologique, dimension cohomologique globale d'un anneau : 0, 17.2.8.\par
Dimension cohomologique d'un espace annelé en un point, dimension cohomologique d'un espace annelé : 0, 17.2.14.\par

\clearpage
\oldpage[IV-1]{255}
Dimension injective, dimension projective d'un module : 0, 17.2.1.\par
Dimension injective en un point, dimension injective d'un faisceau : 0, 17.2.14.\par
Dimension projective en un point, dimension projective d'un faisceau : 0, 17.2.14.\par

Élément de $A$ non diviseur de 0 dans un $A$-module : 0, 15.1.1.\par
Élément $M$-régulier, élément $M$-quasi-régulier : 0, 15.1.4.\par
Ensemble ind-constructible, ensemble pro-constructible : IV, 1.9.4.\par
Épimorphisme formel : 0, 19.1.2.\par
$A$-équivalence de $A$-extensions : 0, 18.2.2.\par
Espace biéquidimensionnel : 0, 14.3.3.\par
Espace caténaire : 0, 14.3.2.\par
Espace équicodimensionnel : 0, 14.2.1.\par
Espace équidimensionnel : 0, 14.1.3.\par
$A$-extension d'un $A$-anneau $B$ par un $B$-bimodule : 0, 18.2.2.\par
Extension déduite d'une autre par un homomorphisme de bimodules : 0, 18.2.8.\par
Extension de Hochschild : 0, 18.4.3.\par
$B$-extension $A$-triviale : 0, 18.3.7.\par
Extension triviale type : 0, 18.2.3.\par
$A$-extensions $A$-équivalentes : 0, 18.2.2.\par

Famille $p$-libre, famille absolument $p$-libre : 0, 21.1.9.\par
Formellement étale (algèbre) : 0, 19.10.2.\par
Formellement inversible à droite (homomorphisme) : 0, 19.1.15.\par
Formellement inversible à gauche (homomorphisme) : 0, 19.1.5.\par
Formellement non ramifiée (algèbre) : 0, 19.10.2.\par
Formellement projectif (module) : 0, 19.2.1.\par
Formellement lisse (algèbre) : 0, 19.3.1.\par
Formellement lisse relativement à un anneau (algèbre) : 0, 19.9.1.\par

Géométriquement régulier sur un corps (anneau local) : 0, 19.6.5.\par
Groupe des classes de $A$-extensions : 0, 18.3.4.\par

Hauteur d'un idéal : 0, 16.1.3.\par
$A$-homomorphisme de $A$-anneaux : 0, 18.1.1.\par
Homomorphisme caractéristique d'une $B$-extension $A$-triviale : 0, 20.6.8.\par
Homomorphisme caractéristique d'une $A$-algèbre relativement à un anneau et à un idéal : 0, 20.6.24.\par
Homomorphisme caractéristique d'une algèbre de caractéristique $p>0$ : 0, 22.4.6.\par

Idéal d'augmentation : 0, 18.1.4.\par
Image réciproque d'un anneau augmenté : 0, 18.1.5.\par
Image réciproque d'une $A$-extension : 0, 18.2.5.\par

Longueur d'une chaîne : 0, 14.1.1.\par

Module biéquidimensionnel, caténaire, équicodimensionnel, équidimensionnel : 0, 16.1.7.\par
Module de Cohen-Macaulay : 0, 16.5.1 et 16.5.13.\par
Module d'imperfection : 0, 20.6.1.\par
Monomorphisme formel : 0, 19.1.2.\par
Morphisme de $A$-extensions : 0, 18.2.4.\par
Morphisme de présentation finie en un point, morphisme localement de présentation finie : IV, 1.4.2.\par
Morphisme de présentation finie : IV, 1.6.1.\par
Morphisme quasi-séparé : IV, 1.2.1.\par

$p$-anneau de Cohen : 0, 19.8.4.\par
$p$-base, $p$-base absolue : 0, 21.1.9.\par
Point maximal d'un préschéma : IV, 1.1.4.\par

\clearpage
\oldpage[IV-1]{256}
Polynôme de Hilbert-Samuel : 0, 16.2.1.\par
Préschéma de présentation finie sur un autre : IV, 1.6.1.\par
Préschéma quasi-séparé : IV, 1.2.1.\par
Produit fibré de $A$-anneaux : 0, 18.1.2.\par
Profondeur d'un module : 0, 16.4.5.\par

Somme amalgamée de bimodules : 0, 18.2.7.\par
Strictement formellement projectif (module) : 0, 19.2.3.\par
Suite $M$-régulière, suite $M$-quasi-régulière : 0, 15.1.7.\par
Suite $\mathscr F$-régulière, suite $\mathscr F$-quasi-régulière : 0, 15.2.2.\par
Système de paramètres pour un module sur un anneau semi-local noethérien : 0, 16.3.6.\par
Système de $p$-générateurs, système absolu de $p$-générateurs : 0, 21.1.9.\par
Système fondamental d'idéaux ouverts, de sous-modules ouverts : 0, 19.0.2.\par
Système régulier de paramètres d'un anneau local régulier : 0, 17.1.6.\par

Topologie constructible sur un préschéma : IV, 1.9.12.\par
Topologie sur un $A$-module déduite de la topologie de $A$ : 0, 19.0.2.\par
Trivial ($A$-anneau augmenté) : 0, 18.1.4.\par
$A$-triviale (extension) : 0, 18.2.3 et 18.3.7.\par

\clearpage
\oldpage[IV-1]{257}
\phantomsection
\addcontentsline{toc}{section}{Table des matières}
\section*{TABLE DES MATIÈRES}

\newcommand{\ivonecontentsline}[3]{\noindent\makebox[17mm][l]{#1}\parbox[t]{0.77\linewidth}{#2\dotfill}\makebox[12mm][r]{#3}\par}
\ivonecontentsline{Chapitre 0.}{\textbf{Préliminaires} (\emph{suite})}{5}
\ivonecontentsline{\S~14.}{Dimension combinatoire d'un espace topologique}{6}
\ivonecontentsline{14.1.}{Dimension combinatoire d'un espace topologique}{6}
\ivonecontentsline{14.2.}{Codimension d'une partie fermée}{8}
\ivonecontentsline{14.3.}{La condition des chaînes}{10}
\ivonecontentsline{\S~15.}{Suites $M$-régulières et suites $\mathscr F$-régulières}{12}
\ivonecontentsline{15.1.}{Suites $M$-régulières et suites $M$-quasi-régulières}{12}
\ivonecontentsline{15.2.}{Suites $\mathscr F$-régulières}{20}
\ivonecontentsline{\S~16.}{Dimension et profondeur dans les anneaux locaux noethériens}{22}
\ivonecontentsline{16.1.}{Dimension d'un anneau}{22}
\ivonecontentsline{16.2.}{Dimension d'un anneau semi-local noethérien}{25}
\ivonecontentsline{16.3.}{Systèmes de paramètres dans un anneau local noethérien}{28}
\ivonecontentsline{16.4.}{Profondeur et coprofondeur}{32}
\ivonecontentsline{16.5.}{Modules de Cohen-Macaulay}{36}
\ivonecontentsline{\S~17.}{Anneaux réguliers}{39}
\ivonecontentsline{17.1.}{Définition des anneaux réguliers}{39}
\ivonecontentsline{17.2.}{Rappels sur la dimension projective et la dimension injective des modules}{42}
\ivonecontentsline{17.3.}{Théorie cohomologique des anneaux réguliers}{46}
\ivonecontentsline{\S~18.}{Compléments sur les extensions d'algèbres}{51}
\ivonecontentsline{18.1.}{Images réciproques d'anneaux augmentés}{51}
\ivonecontentsline{18.2.}{Extensions d'un anneau par un bimodule}{54}
\ivonecontentsline{18.3.}{Le groupe des classes de $A$-extensions}{59}
\ivonecontentsline{18.4.}{Extensions d'algèbres}{64}
\ivonecontentsline{18.5.}{Cas des anneaux topologiques}{66}
\ivonecontentsline{\S~19.}{Algèbres formellement lisses et anneaux de Cohen}{69}
\ivonecontentsline{19.0.}{Introduction}{69}
\ivonecontentsline{19.1.}{Épimorphismes et monomorphismes formels}{71}
\ivonecontentsline{19.2.}{Modules formellement projectifs}{78}

\clearpage
\oldpage[IV-1]{258}
\ivonecontentsline{19.3.}{Algèbres formellement lisses}{79}
\ivonecontentsline{19.4.}{Premiers critères de lissité formelle}{86}
\ivonecontentsline{19.5.}{Lissité formelle et anneaux gradués associés}{90}
\ivonecontentsline{19.6.}{Cas des algèbres sur un corps}{100}
\ivonecontentsline{19.7.}{Cas des homomorphismes locaux : théorèmes d'existence et d'unicité}{104}
\ivonecontentsline{19.8.}{Algèbres de Cohen et $p$-anneaux de Cohen; application à la structure des anneaux locaux complets}{109}
\ivonecontentsline{19.9.}{Algèbres relativement formellement lisses}{114}
\ivonecontentsline{19.10.}{Algèbres formellement non ramifiées et algèbres formellement étales}{115}
\ivonecontentsline{\S~20.}{Dérivations et différentielles}{116}
\ivonecontentsline{20.1.}{Dérivations et extensions d'algèbres}{117}
\ivonecontentsline{20.2.}{Propriétés fonctorielles des dérivations}{119}
\ivonecontentsline{20.3.}{Dérivations continues dans les anneaux topologiques}{121}
\ivonecontentsline{20.4.}{Parties principales et différentielles}{123}
\ivonecontentsline{20.5.}{Propriétés fonctorielles fondamentales de $\Omega^1_{B/A}$}{128}
\ivonecontentsline{20.6.}{Modules d'imperfection et homomorphismes caractéristiques}{136}
\ivonecontentsline{20.7.}{Généralisations aux anneaux topologiques}{147}
\ivonecontentsline{\S~21.}{Différentielles dans les anneaux de caractéristique $p$}{153}
\ivonecontentsline{21.1.}{Systèmes de $p$-générateurs et $p$-bases}{154}
\ivonecontentsline{21.2.}{$p$-bases et lissité formelle}{157}
\ivonecontentsline{21.3.}{$p$-bases et modules d'imperfection}{160}
\ivonecontentsline{21.4.}{Cas des extensions de corps}{162}
\ivonecontentsline{21.5.}{Application : critères de séparabilité}{164}
\ivonecontentsline{21.6.}{Corps admissibles pour une extension}{167}
\ivonecontentsline{21.7.}{L'égalité de Cartier}{169}
\ivonecontentsline{21.8.}{Critères d'admissibilité}{171}
\ivonecontentsline{21.9.}{Modules de différentielles complétés dans les anneaux de séries formelles}{176}
\ivonecontentsline{\S~22.}{Critères différentiels de lissité formelle et de régularité}{182}
\ivonecontentsline{22.1.}{Relèvement de la lissité formelle}{183}
\ivonecontentsline{22.2.}{Caractérisation différentielle des algèbres locales formellement lisses sur un corps}{186}
\ivonecontentsline{22.3.}{Application aux relations entre certains anneaux locaux et leurs complétés}{191}
\ivonecontentsline{22.4.}{Résultats préliminaires sur les extensions finies d'anneaux locaux dont l'idéal maximal est de carré nul}{193}

\clearpage
\oldpage[IV-1]{259}
\ivonecontentsline{22.5.}{Algèbres géométriquement régulières et algèbres formellement lisses}{201}
\ivonecontentsline{22.6.}{Critère jacobien de Zariski}{205}
\ivonecontentsline{22.7.}{Le critère jacobien de Nagata}{209}
\ivonecontentsline{\S~23.}{Anneaux japonais}{213}
\ivonecontentsline{23.1.}{Anneaux japonais}{213}
\ivonecontentsline{23.2.}{Clôture intégrale d'un anneau local noethérien intègre}{217}
\ivonecontentsline{}{\textsc{Chapitre IV.} --- \textbf{Étude locale des schémas et des morphismes de schémas}}{222}
\ivonecontentsline{\S~1.}{Conditions de finitude relatives. Ensembles constructibles dans les préschémas}{224}
\ivonecontentsline{1.1.}{Morphismes quasi-compacts}{224}
\ivonecontentsline{1.2.}{Morphismes quasi-séparés}{226}
\ivonecontentsline{1.3.}{Morphismes localement de type fini}{228}
\ivonecontentsline{1.4.}{Morphismes localement de présentation finie}{230}
\ivonecontentsline{1.5.}{Morphismes de type fini}{233}
\ivonecontentsline{1.6.}{Morphismes de présentation finie}{234}
\ivonecontentsline{1.7.}{Amélioration de résultats antérieurs}{236}
\ivonecontentsline{1.8.}{Morphismes de présentation finie et ensembles constructibles}{238}
\ivonecontentsline{1.9.}{Ensembles pro-constructibles et ind-constructibles}{241}
\ivonecontentsline{1.10.}{Applications aux morphismes ouverts}{249}
\ivonecontentsline{}{Bibliographie (\emph{suite})}{251}
\ivonecontentsline{}{Index des notations}{252}
\ivonecontentsline{}{Index terminologique}{254}

\vspace{2em}
\begin{flushright}
\emph{Reçu le 15 mars 1963.}
\end{flushright}
\let\ivonecontentsline\relax
\endgroup

% FIN SOURCE ega4/ega4-1-fr.tex

\end{document}
% EGA-COMPONENT-END IV1

% EGA-COMPONENT-BEGIN IV2 1C4252DFC75EC33A614492EE2B5A15245DEFF8A7677FE88C4524DB9EE54F1616
\documentclass[10pt,oneside]{book}

\usepackage[T1]{fontenc}
\usepackage[utf8]{inputenc}
\usepackage[french]{babel}
\usepackage{amsmath,amssymb,mathrsfs}
\usepackage[all]{xy}
\usepackage{tikz-cd}
\usepackage[a4paper,margin=22mm]{geometry}
\usepackage[
  colorlinks=true,
  linkcolor=blue,
  hypertexnames=false,
  unicode=true,
  pdfencoding=auto
]{hyperref}
\hypersetup{
  pdftitle={Éléments de géométrie algébrique IV, seconde partie -- édition française diplomatique},
  pdfauthor={Alexandre Grothendieck et Jean Dieudonné},
  pdfsubject={Transcription diplomatique française depuis les exemplaires NUMDAM}
}

\renewcommand{\thesection}{\arabic{section}}
\renewcommand{\thesubsection}{\thesection.\arabic{subsection}}
\newcommand{\oldpage}[2][]{%
  \setcounter{footnote}{0}%
  \marginpar{\scriptsize\textbf{#1}~|~#2}\ignorespaces}
\newenvironment{env}[1][]{\par\medskip\noindent\textbf{(#1)}\ }{\par}
\newenvironment{definition}[1][]{%
  \par\medskip\noindent\itshape Définition (#1). ---\ }{\par}
\newenvironment{proposition}[1][]{%
  \par\medskip\noindent\itshape Proposition (#1). ---\ }{\par}
\newenvironment{theorem}[1][]{%
  \par\medskip\noindent\itshape Théorème (#1). ---\ }{\par}
\newenvironment{corollary}[1][]{%
  \par\medskip\noindent\itshape Corollaire (#1). ---\ }{\par}
\newenvironment{example}[1][]{%
  \par\medskip\noindent\textit{Exemple (#1). ---\ }}{\par}
\newenvironment{examples}[1][]{%
  \par\medskip\noindent\textit{Exemples (#1). ---\ }}{\par}
\newenvironment{remark}[1][]{%
  \par\medskip\noindent\itshape Remarque (#1). ---\ }{\par}
\newenvironment{remarks}[1][]{%
  \par\medskip\noindent\itshape Remarques (#1). ---\ }{\par}
\newenvironment{scholium}[1][]{%
  \par\medskip\noindent\itshape Scholie (#1). ---\ }{\par}
\newenvironment{lemma}[1][]{%
  \par\medskip\noindent\itshape Lemme (#1). ---\ }{\par}
\newenvironment{proof}{%
  \par\medskip\noindent\textit{Démonstration. ---\ }}{\par}

\makeatletter
\newcommand{\EGAReaderTarget}[2]{\hypertarget{#1}{}\@bsphack\begingroup\def\@currentHref{#1}\label{#2}\endgroup\@esphack}
\makeatother
\begin{document}
% BEGIN RIGHTS AND AI TRANSCRIPTION NOTICE
\clearpage
\begingroup
\small
\noindent\textbf{Droits, disponibilité et provenance}

Les droits sur les œuvres originales et sur toute contribution protégée à ces traductions et transcriptions dérivées restent ceux de leurs titulaires respectifs. Ce projet ne place pas ces œuvres ou contributions dans le domaine public et n'accorde aucune licence sur elles. Les mentions de droits de tiers et les autorisations déjà accordées restent applicables.

Ces éditions sont mises à disposition gratuitement, à des fins non commerciales et exclusivement pédagogiques, dans la présomption que les titulaires de droits ne s'opposeraient pas à la libre disponibilité de ces traductions et transcriptions aux côtés des originaux français sur NUMDAM (www.numdam.org/journals/PMIHES/). Cette présomption ne signifie pas qu'une autorisation a été obtenue ; le libre accès aux originaux n'autorise pas, à lui seul, leur traduction, leur adaptation ou leur republication.

Si vous détenez des droits sur un élément présenté ici et vous opposez à sa disponibilité, veuillez ouvrir une demande sur GitHub (github.com/KokunoYumeto/ega-fr/issues). Le responsable de la publication retirera l'élément concerné à la demande du titulaire des droits.

La transcription moderne, les corrections, la composition et les notes mathématiques sont produites par IA. Les travaux documentés comprennent ceux d'OpenAI Codex, GPT-6 Astra, effort Ultra. Les indications de provenance des contributions antérieures restent applicables. Les textes originaux restent l'œuvre d'Alexandre Grothendieck et Jean Dieudonné ; ils ne sont pas attribués à l'IA. Cette édition n'est approuvée ni par les auteurs, ni par l'IHÉS, ni par NUMDAM.
\par
\endgroup
\clearpage
% END RIGHTS AND AI TRANSCRIPTION NOTICE

%
% DEBUT SOURCE ega4/ega4-2-fr.tex
\clearpage
\thispagestyle{plain}
\oldpage[IV-2]{5}
\phantomsection
\label{chapter:IV2-fr}

\begin{center}
{\large CHAPITRE IV (suite)\par}
\vspace{1.5em}
{\Large\bfseries ÉTUDE LOCALE DES SCHÉMAS\\
ET DES MORPHISMES DE SCHÉMAS\par}
\end{center}

\setcounter{section}{1}
\section{Changement de base et platitude}
\label{section:IV.2-fr}

Ce paragraphe (contrairement au \S~6) ne fait appel qu'exceptionnellement à
des techniques noethériennes. Les n\textsuperscript{os}~1 et 2 ne sont guère
que des traductions de propriétés élémentaires de platitude en Algèbre
commutative (cf. Bourbaki, \emph{Alg. comm.}, chap.~I\textsuperscript{er}), et
sont incluses ici pour la commodité des références. Les numéros suivants sont
surtout consacrés à des propriétés de «~descente~» par les morphismes plats ou
fidèlement plats~: si $g:Y'\to Y$ est un tel morphisme, il s'agit de pouvoir
affirmer qu'une partie de $Y$, ou un $\mathcal O_Y$-Module, ou un morphisme
$X\to Y$, a une certaine propriété, lorsque l'on sait que son \emph{image
réciproque} par $g$ possède cette propriété. Nous nous bornons ici aux
propriétés ne faisant pas appel à la technique générale de la «~descente~»,
qui sera développée au chapitre V.

\setcounter{subsection}{0}
\subsection{Modules plats sur les préschémas}
\label{subsection:IV.2.1-fr}

\begin{env}[2.1.1]
\EGAReaderTarget{ega4.c04.statement.2-1-1}{IV.2.1.1-fr}
Soient $f:X\to Y$ un morphisme de préschémas, $\mathscr F$ un
$\mathcal O_X$-Module~; rappelons $(0_{\mathrm I},6.7.1)$ que $\mathscr F$ est
dit $f$-plat (ou $Y$-plat) \emph{en un point} $x\in X$ si $\mathscr F_x$ est
un $\mathcal O_{f(x)}$-module plat, $f$-plat (ou $Y$-plat) s'il est $f$-plat
en tout point $x\in X$, et enfin que le morphisme $f$ est dit plat \emph{au
point} $x\in X$ (resp. plat) si $\mathcal O_X$ est $f$-plat au point $x$
(resp. $f$-plat). Lorsque $f=1_X$, on dira simplement qu'un
$\mathcal O_X$-Module $\mathscr F$ est plat \emph{au point} $x$ (resp. plat)
s'il est $X$-plat en ce point (resp. en tout point $x\in X$), c'est-à-dire si
$\mathscr F_x$ est un $\mathcal O_x$-module plat (resp. s'il en est ainsi pour
tout $x\in X$). Rappelons que nous avons démontré (III, 1.4.15.1) la propriété
suivante~:
\end{env}

\begin{proposition}[2.1.2]
\label{IV.2.1.2-fr}
Soient $A$, $B$ deux anneaux, $\varphi:A\to B$ un homomorphisme d'anneaux,
$X=\operatorname{Spec}(B)$, $Y=\operatorname{Spec}(A)$, $f:X\to Y$ le
morphisme correspondant à $\varphi$, $M$ un $B$-module. Pour que
$\mathscr F=\widetilde M$ soit $f$-plat, il faut et il suffit que $M$ soit un
$A$-module plat.
\end{proposition}

\oldpage[IV-2]{6}
\begin{proposition}[2.1.3]
\label{IV.2.1.3-fr}
Soient $f:X\to Y$ un morphisme de préschémas, $\mathscr F$ un
$\mathcal O_X$-Module quasi-cohérent. Les conditions suivantes sont
équivalentes~:
\begin{enumerate}
\item[\emph{a)}] Pour tout changement de base $g:Y'\to Y$, si l'on pose
$X'=X\times_Y Y'$, le foncteur
$\mathscr G'\mapsto\mathscr F\otimes_Y\mathscr G'$ de la catégorie des
$\mathcal O_{Y'}$-Modules quasi-cohérents dans celle des
$\mathcal O_{X'}$-Modules quasi-cohérents, est exact.
\item[\emph{a$'$)}] La condition \emph{a)} est vérifiée pour tous les
morphismes canoniques $g:\operatorname{Spec}(\mathcal O_y)\to Y$ (I, 2.4.1),
où $y$ parcourt $Y$.
\item[\emph{b)}] $\mathscr F$ est $f$-plat.
\end{enumerate}
\end{proposition}

Les questions étant locales sur $X$ et $Y$, on peut se borner au cas où
$Y=\operatorname{Spec}(B)$, $X=\operatorname{Spec}(A)$,
$\mathscr F=\widetilde M$, où $M$ est un $A$-module. Il est clair que
\emph{a)} entraîne \emph{a$'$)}~; la condition \emph{a$'$)} entraîne que
pour tout $x\in X$, le foncteur $N\mapsto M_{\mathfrak n}\otimes_{B_{\mathfrak n}}N$
est exact dans la catégorie des $B_{\mathfrak n}$-modules, $\mathfrak n$
étant l'idéal $\mathfrak j_{f(x)}$ de $B$~; cela signifie que
$M_{\mathfrak n}$ est un $B_{\mathfrak n}$-module plat, et il résulte de
$(0_{\mathrm I},6.3.3)$ et de (2.1.2) que $\mathscr F$ est $f$-plat. Enfin,
pour voir que \emph{b)} entraîne \emph{a)}, on peut aussi se borner au cas où
$Y'=\operatorname{Spec}(A')$ est affine et $\mathscr G'=\widetilde{N'}$, où
$N'$ est un $A'$-module~; la conclusion provient alors encore de (2.1.2) et
de la définition de la platitude, puisque
$(M\otimes_A N')^{\sim}=\mathscr F\otimes_Y\mathscr G'$.

\begin{proposition}[2.1.4]
\label{IV.2.1.4-fr}
Soient $f:X\to Y$, $g:Y'\to Y$ deux morphismes de préschémas, $\mathscr F$
un $\mathcal O_X$-Module quasi-cohérent~; posons $X'=X\times_Y Y'$,
$f'=f_{(Y')}:X'\to Y'$, $\mathscr F'=\mathscr F\otimes_Y\mathcal O_{Y'}$,
et soit $g'$ la projection canonique $X'\to X$. Soient $x'$ un point de
$X'$, $x=g'(x')$, $y'=f'(x')$, $y=g(y')=f(x)$. Si $\mathscr F$ est $f$-plat
au point $x$, $\mathscr F'$ est $f'$-plat au point $x'$~; en particulier si
$\mathscr F$ est $f$-plat, $\mathscr F'$ est $f'$-plat~; si $f$ est plat,
$f'$ est plat.
\end{proposition}

Il suffit de prouver la première assertion~; appliquant trois fois (I, 3.6.5),
ainsi que (I, 2.4.4), on peut se ramener au cas où
$Y=\operatorname{Spec}(\mathcal O_y)$, $X=\operatorname{Spec}(\mathcal O_x)$,
$Y'=\operatorname{Spec}(\mathcal O_{y'})$, $\mathscr F=\widetilde M$, où
$M=\mathscr F_x$~; l'hypothèse et (2.1.2) entraînent alors que $\mathscr F$
est $f$-plat, autrement dit, on est ramené à prouver un cas particulier de la
seconde assertion, et cette dernière découle aussitôt de (2.1.3).

\begin{proposition}[2.1.5]
\label{IV.2.1.5-fr}
Considérons un diagramme commutatif de morphismes de préschémas
\[
\xymatrix{
  X \ar[d]_f & X' \ar[l]_{g'} \ar[d]^{f'} \\
  Y & Y' \ar[l]^g \ar[d]^h \\
  & Z
}
\]
où $X'=X\times_Y Y'$ et $f'=f_{(Y')}$. Soit $x'$ un point de $X'$, et
posons $x=g'(x')$, $y'=f'(x')$, $y=f(x)=g(y')$, $z=h(y')$. Soient
$\mathscr F$ un $\mathcal O_X$-Module quasi-cohérent $f$-plat au point $x$
(resp. $f$-plat), et $\mathscr G'$ un $\mathcal O_{Y'}$-Module quasi-cohérent
$h$-plat au point $y'$ (resp. $h$-plat)~; alors
$\mathscr F\otimes_{Y'}\mathscr G'$ est un $\mathcal O_{X'}$-Module
quasi-cohérent $(hf')$-plat au point $x'$ (resp. $(hf')$-plat).
\end{proposition}

Comme dans (2.1.4), on se ramène au cas où
$X=\operatorname{Spec}(\mathcal O_x)$, $Y=\operatorname{Spec}(\mathcal O_y)$,
\oldpage[IV-2]{7}
$Y'=\operatorname{Spec}(\mathcal O_{y'})$ et
$Z=\operatorname{Spec}(\mathcal O_z)$, et il suffit alors de prouver que
$\mathscr F\otimes_{Y'}\mathscr G'$ est $(hf')$-plat. Compte tenu de
(2.1.2), la proposition résulte de Bourbaki, \emph{Alg. comm.},
chap.~I\textsuperscript{er}, \S~2, n\textsuperscript{o}~7, prop.~8.

\begin{corollary}[2.1.6]
\label{IV.2.1.6-fr}
Soient $f:X\to Y$, $g:Y\to Z$ deux morphismes de préschémas, $\mathscr F$ un
$\mathcal O_X$-Module. Si $\mathscr F$ est $f$-plat au point $x\in X$ et si
$g$ est plat au point $f(x)$, $\mathscr F$ est $(gf)$-plat au point $x$. En
particulier, si $f$ et $g$ sont des morphismes plats, il en est de même de
$g\circ f$.
\end{corollary}

C'est le cas particulier de (2.1.5) avec $Y'=Y$, $\mathscr G'=\mathcal O_Y$.

\begin{corollary}[2.1.7]
\label{IV.2.1.7-fr}
Si $f:X\to X'$, $g:Y\to Y'$ sont deux $S$-morphismes plats,
$f\times_S g:X\times_S Y\to X'\times_S Y'$ est un morphisme plat.
\end{corollary}

Cela résulte de (2.1.4) et (2.1.6) (cf. I, 3.5.1).

\begin{proposition}[2.1.8]
\label{IV.2.1.8-fr}
Soient $f:X\to Y$ un morphisme de préschémas,
\[
0\to\mathscr F'\to\mathscr F\to\mathscr F''\to0
\]
une suite exacte de $\mathcal O_X$-Modules quasi-cohérents, telle que
$\mathscr F''$ soit $Y$-plat.
\begin{enumerate}
\item[(i)] Pour tout morphisme $g:Y'\to Y$ et tout $\mathcal O_{Y'}$-Module
quasi-cohérent $\mathscr G'$, la suite
\[
0\to\mathscr F'\otimes_Y\mathscr G'\to
\mathscr F\otimes_Y\mathscr G'\to
\mathscr F''\otimes_Y\mathscr G'\to0
\]
de $\mathcal O_{X'}$-Modules (où $X'=X\times_Y Y'$) est exacte.
\item[(ii)] Pour que $\mathscr F$ soit $Y$-plat, il faut et il suffit que
$\mathscr F'$ le soit.
\end{enumerate}
\end{proposition}

On peut évidemment supposer $X$, $Y$, $Y'$ affines~; la conclusion résulte
alors de (2.1.2) et de $(0_{\mathrm I},6.1.2)$.

\begin{corollary}[2.1.9]
\label{IV.2.1.9-fr}
Soit $\mathscr L^\bullet$ un complexe de $\mathcal O_X$-Modules
quasi-cohérents, $i$ un indice tel que si l'on note
$d^i:\mathscr L^i\to\mathscr L^{i+1}$ l'opérateur de dérivation,
$\mathscr B^{i+1}(\mathscr L^\bullet)=\operatorname{Im}(d^i)$ et
$\mathscr Z^{i+1}(\mathscr L^\bullet)=\operatorname{Coker}(d^i)$ soient
$Y$-plats. Alors, avec les notations de (2.1.8), l'homomorphisme canonique
\[
\mathscr H^i(\mathscr L^\bullet)\otimes_Y\mathscr G'
\longrightarrow
\mathscr H^i(\mathscr L^\bullet\otimes_Y\mathscr G')
\]
est bijectif.
\end{corollary}

Comme le produit tensoriel est exact à droite, on a
\[
\mathscr Z^{i+1}(\mathscr L^\bullet)\otimes_Y\mathscr G'
=\operatorname{Coker}(d^i\otimes1)
=\mathscr Z^{i+1}(\mathscr L^\bullet\otimes_Y\mathscr G')
\]
et
$\mathscr Z^i(\mathscr L^\bullet)\otimes_Y\mathscr G'
=\mathscr Z^i(\mathscr L^\bullet\otimes_Y\mathscr G')$. En outre, dans la
suite exacte
\[
0\to\mathscr B^{i+1}(\mathscr L^\bullet)\to\mathscr L^{i+1}
\to\mathscr Z^{i+1}(\mathscr L^\bullet)\to0
\]
$\mathscr Z^{i+1}(\mathscr L^\bullet)$ est $Y$-plat, donc il résulte de
(2.1.8, (i)) que l'on a la suite exacte
\[
0\to\mathscr B^{i+1}(\mathscr L^\bullet)\otimes_Y\mathscr G'
\to\mathscr L^{i+1}\otimes_Y\mathscr G'
\to\mathscr Z^{i+1}(\mathscr L^\bullet\otimes_Y\mathscr G')\to0
\]
d'où
$\mathscr B^{i+1}(\mathscr L^\bullet)\otimes_Y\mathscr G'
=\operatorname{Im}(d^i\otimes1)
=\mathscr B^{i+1}(\mathscr L^\bullet\otimes_Y\mathscr G')$. Alors, comme
dans la suite exacte
\[
0\to\mathscr H^i(\mathscr L^\bullet)\to\mathscr Z^i(\mathscr L^\bullet)
\to\mathscr B^{i+1}(\mathscr L^\bullet)\to0
\]
$\mathscr B^{i+1}(\mathscr L^\bullet)$ est $Y$-plat, il résulte de
(2.1.8, (i)) et de ce qui précède que l'on a la suite exacte
\[
0\to\mathscr H^i(\mathscr L^\bullet)\otimes_Y\mathscr G'
\to\mathscr Z^i(\mathscr L^\bullet\otimes_Y\mathscr G')
\to\mathscr B^{i+1}(\mathscr L^\bullet\otimes_Y\mathscr G')\to0,
\]
ce qui démontre le corollaire.

\oldpage[IV-2]{8}
\begin{corollary}[2.1.10]
\label{IV.2.1.10-fr}
Soient $f:X\to Y$ un morphisme de préschémas, $\mathscr F$ un
$\mathcal O_X$-Module quasi-cohérent et $Y$-plat,
$\mathscr L_\bullet=(\mathscr L_i)$ une résolution gauche de $\mathscr F$
formée de $\mathcal O_X$-Modules quasi-cohérents et $Y$-plats. Alors, pour
tout morphisme $g:Y'\to Y$ et tout $\mathcal O_{Y'}$-Module quasi-cohérent
$\mathscr G'$, le complexe
$\mathscr L_\bullet\otimes_Y\mathscr G'=
(\mathscr L_i\otimes_Y\mathscr G')_{i\geq0}$ est une résolution gauche de
$\mathscr F\otimes_Y\mathscr G'$.

En outre, si
$\mathscr Z_i(\mathscr L_\bullet)=\operatorname{Ker}
(\mathscr L_i\to\mathscr L_{i-1})$, les $\mathscr Z_i(\mathscr L_\bullet)$
sont $Y$-plats, et l'on a
\[
\mathscr Z_i(\mathscr L_\bullet)\otimes_Y\mathscr G'
=\mathscr Z_i(\mathscr L_\bullet\otimes_Y\mathscr G')
=\operatorname{Ker}(\mathscr L_i\otimes_Y\mathscr G'
\to\mathscr L_{i-1}\otimes_Y\mathscr G').
\]
Posons
$\mathscr R_i=\operatorname{Im}(\mathscr L_{i+1}\to\mathscr L_i)
=\mathscr Z_i(\mathscr L_\bullet)$~; on a alors les suites exactes
\[
0\leftarrow\mathscr F\leftarrow\mathscr L_0\leftarrow\mathscr R_0\leftarrow0
\]
\[
0\leftarrow\mathscr R_i\leftarrow\mathscr L_{i+1}
\leftarrow\mathscr R_{i+1}\leftarrow0\qquad(i\geq0)
\]
et comme $\mathscr F$ et les $\mathscr L_i$ sont $Y$-plats, on déduit de
(2.1.8, (ii)) par récurrence que tous les $\mathscr R_i$ sont aussi
$Y$-plats~; utilisant (2.1.8, (i)), on a donc les suites exactes
\[
0\leftarrow\mathscr F\otimes_Y\mathscr G'
\leftarrow\mathscr L_0\otimes_Y\mathscr G'
\leftarrow\mathscr R_0\otimes_Y\mathscr G'\leftarrow0,
\]
\[
0\leftarrow\mathscr R_i\otimes_Y\mathscr G'
\leftarrow\mathscr L_{i+1}\otimes_Y\mathscr G'
\leftarrow\mathscr R_{i+1}\otimes_Y\mathscr G'\leftarrow0
\qquad(i\geq0),
\]
qui prouvent le corollaire.
\end{corollary}

\begin{proposition}[2.1.11]
\label{IV.2.1.11-fr}
Soient $f:X\to Y$ un morphisme plat, $\mathscr F$ un $\mathcal O_Y$-Module
quasi-cohérent de présentation finie. Si $\mathscr J$ est l'Idéal de
$\mathcal O_Y$ annulateur de $\mathscr F$, $f^*(\mathscr J)$ est l'Idéal de
$\mathcal O_X$ annulateur de $f^*(\mathscr F)$.
\end{proposition}

On a en effet par définition une suite exacte $(0_{\mathrm I},5.3.7)$
\[
0\to\mathscr J\to\mathcal O_Y\to
\mathscr H\!om_{\mathcal O_Y}(\mathscr F,\mathscr F),
\]
d'où, puisque $f$ est plat, une suite exacte
\[
0\to f^*(\mathscr J)\to\mathcal O_X\to
f^*(\mathscr H\!om_{\mathcal O_Y}(\mathscr F,\mathscr F)),
\]
et comme par hypothèse $\mathscr F$ est un $\mathcal O_Y$-Module de
présentation finie,
$f^*(\mathscr H\!om_{\mathcal O_Y}(\mathscr F,\mathscr F))$ s'identifie
canoniquement à
$\mathscr H\!om_{\mathcal O_X}(f^*(\mathscr F),f^*(\mathscr F))$
$(0_{\mathrm I},6.7.6)$, d'où la conclusion.

\begin{proposition}[2.1.12]
\label{IV.2.1.12-fr}
Soit $X$ un préschéma, $\mathscr F$ un $\mathcal O_X$-Module de présentation
finie, $x$ un point de $X$. Les conditions suivantes sont équivalentes~:
\begin{enumerate}
\item[\emph{a)}] $\mathscr F_x$ est un $\mathcal O_x$-module plat.
\item[\emph{b)}] Il existe un voisinage ouvert $U$ de $x$ tel que
$\mathscr F|U$ soit un $(\mathcal O_X|U)$-Module localement libre.
\end{enumerate}
\end{proposition}

En effet, $\mathscr F_x$ est un $\mathcal O_x$-module de présentation finie
et $\mathcal O_x$ un anneau local~; il revient donc au même de dire que
$\mathscr F_x$ est un $\mathcal O_x$-module plat ou un $\mathcal O_x$-module
libre (Bourbaki, \emph{Alg. comm.}, chap.~II, \S~3,
n\textsuperscript{o}~2, cor.~2 de la prop.~5)~; d'où la conclusion, compte
tenu de $(0_{\mathrm I},5.2.7)$. On notera que la proposition est valable pour
un \emph{espace annelé en anneaux locaux} quelconque.

\begin{proposition}[2.1.13]
\label{IV.2.1.13-fr}
Soit $f:X\to Y$ un morphisme de préschémas. Si $f$ est plat en un point
$x\in X$ et si l'anneau $\mathcal O_x$ est réduit (resp. intègre et
intégralement clos), alors l'anneau $\mathcal O_{f(x)}$
\oldpage[IV-2]{9}
est réduit (resp. intègre et intégralement clos). Si $f$ est fidèlement plat
et si $X$ est réduit (resp. normal), alors $Y$ est réduit (resp. normal).
\end{proposition}

Posons $\mathcal O_{f(x)}=A$, $\mathcal O_x=B$. Si $B$ est un $A$-module
plat, c'est aussi un $A$-module fidèlement plat $(0_{\mathrm I},6.6.2)$, donc
$A$ s'identifie à un sous-anneau de $B$~; si $B$ est réduit, il en est donc de
même de $A$. Supposons maintenant $B$ intègre et intégralement clos, et soit
$L$ son corps des fractions~; alors $A\subset B$ est intègre~; désignons par
$K\subset L$ son corps des fractions. L'hypothèse entraîne que
$B\cap K=A$ (Bourbaki, \emph{Alg. comm.}, chap.~I, \S~3,
n\textsuperscript{o}~5, prop.~10). Si alors $t\in K$ est entier sur $A$, il
est aussi entier sur $B$, donc appartient à $B$ par hypothèse, et par suite
$t\in A$, ce qui prouve que $A$ est intégralement clos.

\begin{proposition}[2.1.14]
\label{IV.2.1.14-fr}
Soit $f:X\to Y$ un morphisme fidèlement plat. Si $X$ est localement intègre
et si l'espace sous-jacent à $Y$ est localement noethérien, alors $Y$ est
localement intègre.
\end{proposition}

En effet, tout $y\in Y$ est de la forme $f(x)$ pour un $x\in X$, et par
hypothèse $\mathcal O_y$ s'identifie à un sous-anneau de $\mathcal O_x$
$(0_{\mathrm I},6.6.1)$~; comme $\mathcal O_x$ est intègre, il en est de même
de $\mathcal O_y$ et cela prouve la proposition (I, 5.1.4).

\subsection{Modules fidèlement plats sur les préschémas}
\label{subsection:IV.2.2-fr}

\begin{proposition}[2.2.1]
\label{IV.2.2.1-fr}
Soient $f:X\to Y$ un morphisme de préschémas, $\mathscr F$ un
$\mathcal O_X$-Module quasi-cohérent. Les conditions suivantes sont
équivalentes~:
\begin{enumerate}
\item[\emph{a)}] Pour tout changement de base $g:Y'\to Y$, si l'on pose
$X'=X\times_Y Y'$, le foncteur
$\mathscr G'\mapsto\mathscr F\otimes_Y\mathscr G'$, de la catégorie des
$\mathcal O_{Y'}$-Modules quasi-cohérents dans celle des
$\mathcal O_{X'}$-Modules quasi-cohérents, est exact et fidèle.
\item[\emph{a$'$)}] La condition \emph{a)} est vérifiée pour tous les
morphismes canoniques $g:\operatorname{Spec}(\mathcal O_y)\to Y$ (I, 2.4.1),
où $y$ parcourt $Y$.
\item[\emph{a$''$)}] La condition \emph{a)} est vérifiée pour toutes les
immersions canoniques $Y'\to Y$, où $Y'$ parcourt l'ensemble des
sous-schémas ouverts affines de $Y$.
\item[\emph{b)}] $\mathscr F$ est $Y$-plat et, pour tout $y\in Y$, si l'on
désigne par $X_y$ la fibre $f^{-1}(y)$, le $\mathcal O_{X_y}$-Module
$\mathscr F_y=\mathscr F\otimes_Y k(y)$ est $\neq0$.
\end{enumerate}
\end{proposition}

Il est clair que \emph{a)} implique \emph{a$'$)} et \emph{a$''$)}~; la
condition \emph{a$'$)} implique d'abord que $\mathscr F$ est $Y$-plat
(2.1.3)~; d'autre part elle implique que pour tout $y\in Y$, le foncteur
$N\mapsto\mathscr F\otimes_{\mathcal O_y}\widetilde N$ est fidèle dans la
catégorie des $\mathcal O_y$-modules~; prenant en particulier $N=k(y)$, on
obtient la seconde assertion de \emph{b)}. Pour montrer que \emph{b)}
implique \emph{a)}, on peut se limiter au cas où $Y$ est \emph{affine}, la
question étant locale sur $Y$. De même, pour prouver que \emph{a$''$)}
implique \emph{a)}, on est ramené à prouver que lorsque $Y$ est affine, le
fait que $\mathscr G\mapsto\mathscr F\otimes_Y\mathscr G$ soit un foncteur
exact et fidèle entraîne la condition \emph{a)}. Autrement dit, on est ramené
à prouver la proposition plus précise suivante~:

\begin{proposition}[2.2.2]
\label{IV.2.2.2-fr}
Soient $Y=\operatorname{Spec}(A)$ un schéma affine, $f:X\to Y$ un morphisme
de préschémas, $\mathscr F$ un $\mathcal O_X$-Module quasi-cohérent. La
condition \emph{a)} de (2.2.1) est équivalente à chacune des suivantes~:
\begin{enumerate}
\item[\emph{b$'$)}] $\mathscr F$ est $Y$-plat et, pour tout point fermé $y$
de $Y$, on a $\mathscr F_y\neq0$.
\oldpage[IV-2]{10}
\item[\emph{c)}] Le foncteur
$\mathscr G\mapsto\mathscr F\otimes_Y\mathscr G$ de la catégorie des
$\mathcal O_Y$-Modules quasi-cohérents dans celle des
$\mathcal O_X$-Modules quasi-cohérents, est exact et fidèle.
\end{enumerate}
\end{proposition}

Si \emph{b$'$)} est vérifié, il y a au moins un $x\in f^{-1}(y)$ tel que
$(\mathscr F_y)_x\neq0$~; soit $U=\operatorname{Spec}(B)$ un voisinage ouvert
affine de $x$, et soit $\mathscr F|U=\widetilde M$, où $M$ est un $B$-module.
Alors \emph{b$'$)} implique que $M/\mathfrak j_yM$ est $\neq0$, et par suite
(puisque $M$ est un $A$-module plat (2.1.2)) que
$M\otimes_A A_y$ est un $A_y$-module fidèlement plat
$(0_{\mathrm I},6.4.5)$. La relation
$(\mathscr F\otimes_Y\mathscr G)\otimes_Y\mathcal O_y=0$ pour un point fermé
$y$ de $Y$ implique
\[
(\mathscr F\otimes_Y\mathcal O_y)\otimes_{\mathcal O_y}\mathscr G_y=0,
\]
donc $\mathscr G_y=0$. Mais si l'on a $\mathscr G_y=0$ pour tout point fermé
de $Y$, on en conclut $\mathscr G=0$, car si $\mathscr G=\widetilde N$,
l'annulateur d'un élément de $N$ n'est contenu dans aucun idéal maximal de
$A$, donc est $A$ tout entier. La relation
$\mathscr F\otimes_Y\mathscr G=0$ implique donc $\mathscr G=0$, autrement dit
le foncteur $\mathscr G\mapsto\mathscr F\otimes_Y\mathscr G$ est fidèle~; on
sait par ailleurs que ce foncteur est exact (2.1.3), ce qui montre bien que
\emph{b$'$)} entraîne \emph{c)}.

Enfin, pour voir que \emph{c)} entraîne \emph{a)}, on peut se borner au cas
où $Y'$ est aussi affine~; comme $g:Y'\to Y$ est alors un morphisme affine,
il en est de même de la projection $g':X'\to X$ (II, 1.5.5)~; de plus, le
foncteur $\mathscr H'\mapsto g'_*(\mathscr H')$ est alors exact dans la
catégorie des $\mathcal O_{X'}$-Modules quasi-cohérents (I, 1.6.4), et si
$g'_*(\mathscr H')=\mathscr H$, on a $\mathscr H'=\widetilde{\mathscr H}$,
donc le foncteur précédent est aussi fidèle~; pour voir que
$\mathscr G'\mapsto\mathscr F\otimes_Y\mathscr G'$ est exact et fidèle, il
suffit par suite de voir que le foncteur
\[
\mathscr G'\mapsto g'_*(\mathscr F\otimes_Y\mathscr G')
\]
l'est. Or, si $f'=f_{(Y')}:X'\to Y'$, on a
$g'_*(\mathscr F\otimes_Y\mathscr G')
=g'_*(g'^*(\mathscr F)\otimes_{X'}f'^*(\mathscr G'))$~; le fait que $g$ est
affine entraîne que l'on a un isomorphisme canonique
\begin{equation}
\tag{2.2.2.1}
\label{IV.2.2.2.1-fr}
\mathscr F\otimes_{\mathcal O_X}f^*(g_*(\mathscr G'))
\xrightarrow{\sim}
g'_*(g'^*(\mathscr F)\otimes_{\mathcal O_{X'}}f'^*(\mathscr G')).
\end{equation}
En effet, on sait (II, 1.5.2) qu'on a un isomorphisme canonique
\[
f^*(g_*(\mathscr G'))\xrightarrow{\sim}g'_*(f'^*(\mathscr G')),
\]
et d'autre part on a un homomorphisme canonique
$\mathscr F\to g'_*(g'^*(\mathscr F))$ $(0_{\mathrm I},4.4.3.2)$~;
composant l'homomorphisme
$\mathscr F\otimes_{\mathcal O_X}f^*(g_*(\mathscr G'))
\to g'_*(g'^*(\mathscr F))\otimes_{\mathcal O_X}g'_*(f'^*(\mathscr G'))$
avec l'homomorphisme canonique $(0_{\mathrm I},4.2.2.1)$, on en déduit
l'homomorphisme (2.2.2.1)~; la vérification du fait qu'il s'agit d'un
isomorphisme est immédiate en se ramenant au cas où $X$ est affine. Cela
étant, le foncteur $\mathscr G'\mapsto g_*(\mathscr G')$ est exact et fidèle
et par hypothèse il en est de même de
$\mathscr G\mapsto\mathscr F\otimes_Y\mathscr G
=\mathscr F\otimes_{\mathcal O_X}f^*(\mathscr G)$~; leur composé est donc
exact et fidèle, ce qui achève de prouver (2.2.1) et (2.2.2).

\begin{corollary}[2.2.3]
\label{IV.2.2.3-fr}
Soient $X=\operatorname{Spec}(B)$, $Y=\operatorname{Spec}(A)$ deux schémas
affines, $f:X\to Y$ un morphisme, $\mathscr F=\widetilde M$ un
$\mathcal O_X$-Module quasi-cohérent. Pour que $\mathscr F$ vérifie les
conditions équivalentes de (2.2.1) (ou (2.2.2)), il faut et il suffit que le
$A$-module $M$ soit fidèlement plat.
\end{corollary}

En effet, la condition \emph{c)} de (2.2.2) signifie alors que le foncteur
$N\mapsto M\otimes_A N$ de la catégorie des $A$-modules dans celle des
$B$-modules est exact et fidèle, et la conclusion résulte de
$(0_{\mathrm I},6.4.1)$.

\begin{definition}[2.2.4]
\label{IV.2.2.4-fr}
Lorsque les conditions équivalentes de (2.2.1) sont satisfaites, on dit que
le $\mathcal O_X$-Module quasi-cohérent $\mathscr F$ est \emph{fidèlement plat
relativement à} $f$ (ou à $Y$).
\end{definition}

\oldpage[IV-2]{11}
On notera que cette notion est locale sur $Y$, mais \emph{non sur} $X$~; en
particulier on peut avoir $\mathscr F_x=0$ pour certains $x\in X$, autrement
dit $\operatorname{Supp}(\mathscr F)$ n'est pas nécessairement égal à $X$.
Toutefois, il résulte aussitôt du critère (2.2.1, \emph{b)}) que pour tout
$y\in Y$, il existe au moins un $x\in f^{-1}(y)$ tel que
$(\mathscr F_y)_x=\mathscr F_x\otimes_{\mathcal O_y}k(y)\neq0$, et
\emph{a fortiori} $\mathscr F_x\neq0$~; en d'autres termes~:

\begin{corollary}[2.2.5]
\label{IV.2.2.5-fr}
Si $\mathscr F$ est un $\mathcal O_X$-Module quasi-cohérent, fidèlement plat
relativement à $f$, on a $f(\operatorname{Supp}(\mathscr F))=Y$, et
\emph{a fortiori} $f$ est un morphisme surjectif.
\end{corollary}

Ce résultat admet une réciproque partielle~:

\begin{corollary}[2.2.6]
\label{IV.2.2.6-fr}
Soit $\mathscr F$ un $\mathcal O_X$-Module quasi-cohérent de type fini. Pour
que $\mathscr F$ soit fidèlement plat relativement à $f$, il faut et il
suffit que $\mathscr F$ soit $f$-plat et que
$f(\operatorname{Supp}(\mathscr F))=Y$.
\end{corollary}

En effet (I, 9.1.13 et 3.6.1) on a
$\operatorname{Supp}(\mathscr F_y)=f^{-1}(y)\cap\operatorname{Supp}(\mathscr F)$,
donc dans ce cas le critère (2.2.1, \emph{b)}) n'est autre que la condition du
corollaire. En particulier, le $\mathcal O_X$-Module $\mathcal O_X$ est
fidèlement plat relativement à $f$ si et seulement s'il est $f$-plat et si
$f$ est surjectif, autrement dit si et seulement si \emph{le morphisme $f$
est fidèlement plat} $(0_{\mathrm I},6.7.8)$.

Explicitons les propriétés intervenant dans la définition (2.2.4)~:

\begin{proposition}[2.2.7]
\label{IV.2.2.7-fr}
Soient $f:X\to Y$ un morphisme de préschémas, $\mathscr F$ un
$\mathcal O_X$-Module quasi-cohérent fidèlement plat relativement à $f$.
Alors, pour qu'une suite $\mathscr G'\to\mathscr G\to\mathscr G''$ de
$\mathcal O_Y$-Modules quasi-cohérents soit exacte, il faut et il suffit que
la suite correspondante
$\mathscr F\otimes_Y\mathscr G'\to\mathscr F\otimes_Y\mathscr G
\to\mathscr F\otimes_Y\mathscr G''$ le soit. En particulier, pour qu'un
homomorphisme $u:\mathscr G\to\mathscr G'$ de $\mathcal O_Y$-Modules
quasi-cohérents soit injectif (resp. surjectif, bijectif, nul), il faut et il
suffit que $1_{\mathscr F}\otimes f^*(u):
\mathscr F\otimes_Y\mathscr G\to\mathscr F\otimes_Y\mathscr G'$ le soit.
Pour qu'un $\mathcal O_Y$-Module quasi-cohérent $\mathscr G$ soit nul, il faut
et il suffit que $\mathscr F\otimes_Y\mathscr G=0$. Pour tout
$\mathcal O_Y$-Module quasi-cohérent $\mathscr G$, l'application
$\mathscr G'\mapsto\mathscr F\otimes_Y\mathscr G'$ de l'ensemble des
sous-$\mathcal O_Y$-Modules quasi-cohérents de $\mathscr G$ dans l'ensemble
des sous-$\mathcal O_X$-Modules quasi-cohérents de
$\mathscr F\otimes_Y\mathscr G$ est injective.
\end{proposition}

Pour démontrer la dernière assertion, c'est-à-dire que pour deux
sous-$\mathcal O_Y$-Modules quasi-cohérents $\mathscr G'$, $\mathscr G''$ de
$\mathscr G$, la relation
$\mathscr F\otimes_Y\mathscr G'=\mathscr F\otimes_Y\mathscr G''$ entraîne
$\mathscr G'=\mathscr G''$, on peut (en remplaçant $\mathscr G''$ par
$\mathscr G'+\mathscr G''$) se ramener au cas où
$\mathscr G'\subset\mathscr G''$, et il suffit alors d'appliquer la seconde
assertion de l'énoncé à l'injection $u:\mathscr G'\to\mathscr G''$.

\begin{corollary}[2.2.8]
\label{IV.2.2.8-fr}
Soit $f:X\to Y$ un morphisme fidèlement plat. Pour tout
$\mathcal O_Y$-Module quasi-cohérent $\mathscr G$, l'application canonique
\begin{equation}
\tag{2.2.8.1}
\label{IV.2.2.8.1-fr}
\Gamma(Y,\mathscr G)\longrightarrow\Gamma(X,f^*(\mathscr G))
\end{equation}
est injective.
\end{corollary}

En effet, $\Gamma(Y,\mathscr G)$ s'identifie canoniquement à
$\operatorname{Hom}_{\mathcal O_Y}(\mathcal O_Y,\mathscr G)$
$(0_{\mathrm I},5.1.1)$, et de même $\Gamma(X,f^*(\mathscr G))$ à
$\operatorname{Hom}_{\mathcal O_X}(\mathcal O_X,f^*(\mathscr G))$. En vertu
de (2.2.1) et (2.2.4) l'hypothèse entraîne que le foncteur
$\mathscr G\mapsto f^*(\mathscr G)$ est exact et fidèle dans la catégorie des
$\mathcal O_Y$-Modules quasi-cohérents, et par suite un homomorphisme
$u:\mathcal O_Y\to\mathscr G$ est nul si et seulement si
$f^*(u):f^*(\mathcal O_Y)=\mathcal O_X\to f^*(\mathscr G)$ est nul.

\begin{remark}[2.2.9]
\label{IV.2.2.9-fr}
Les résultats de (2.2.7) et (2.2.8) sont encore vrais lorsque les
$\mathcal O_Y$-Modules $\mathscr G'$, $\mathscr G$, $\mathscr G''$ qui y
figurent sont quelconques (non nécessairement quasi-cohérents). En effet,
pour tout $y\in Y$, il existe $x\in f^{-1}(y)$ tel que $\mathscr F_x$ soit un
$\mathcal O_y$-module
\oldpage[IV-2]{12}
fidèlement plat, et par suite le foncteur
$\mathscr G_y\mapsto\mathscr G_y\otimes_{\mathcal O_y}\mathscr F_x$ est
fidèle~; comme par ailleurs pour \emph{tout} $x\in f^{-1}(y)$ le foncteur
$\mathscr G_y\mapsto\mathscr G_y\otimes_{\mathcal O_y}\mathscr F_x$ est
exact, on en déduit aussitôt la conclusion.
\end{remark}

\begin{proposition}[2.2.10]
\label{IV.2.2.10-fr}
Soient $f:X\to Y$, $g:Y'\to Y$, $h:Y'\to Z$ trois morphismes de préschémas,
$X'=X\times_Y Y'$, $\mathscr F$ un $\mathcal O_X$-Module quasi-cohérent,
fidèlement plat relativement à $Y$, $\mathscr G'$ un
$\mathcal O_{Y'}$-Module quasi-cohérent. Pour que $\mathscr G'$ soit plat
(resp. fidèlement plat) relativement à $Z$, il faut et il suffit que
$\mathscr F\otimes_Y\mathscr G'$ soit un $\mathcal O_{X'}$-Module plat
(resp. fidèlement plat) relativement à $Z$.
\end{proposition}

On sait déjà que si $\mathscr G'$ est $Z$-plat, il en est de même de
$\mathscr F\otimes_Y\mathscr G'$ (2.1.5). Considérons un changement de base
quelconque $Z''\to Z$ et posons $X''=X'\times_Z Z''$~; si $\mathscr G'$ est
fidèlement plat relativement à $Z$, le foncteur
\begin{equation}
\tag{2.2.10.1}
\label{IV.2.2.10.1-fr}
\mathscr H''\mapsto\mathscr H''\otimes_Z\mathscr G'
\mapsto(\mathscr H''\otimes_Z\mathscr G')\otimes_Y\mathscr F
=\mathscr H''\otimes_Z(\mathscr G'\otimes_Y\mathscr F)
\end{equation}
de la catégorie des $\mathcal O_{Z''}$-Modules quasi-cohérents dans celle des
$\mathcal O_{X''}$-Modules est composé de deux foncteurs exacts et fidèles,
donc est exact et fidèle. Inversement, si ce foncteur composé est exact (resp.
exact et fidèle), il en est de même de
$\mathscr H''\mapsto\mathscr H''\otimes_Z\mathscr G'$, puisque le foncteur
$\mathscr M'\mapsto\mathscr M'\otimes_Y\mathscr F$ (de la catégorie des
$\mathcal O_{Y'}$-Modules quasi-cohérents dans celle des
$\mathcal O_{X'}$-Modules quasi-cohérents) est exact et fidèle par hypothèse.

\begin{corollary}[2.2.11]
\label{IV.2.2.11-fr}
\begin{enumerate}
\item[(i)] Soient $f:X\to Y$, $g:Y'\to Y$ deux morphismes,
$X'=X\times_Y Y'$, $\mathscr F$ un $\mathcal O_X$-Module quasi-cohérent. Si
$\mathscr F$ est fidèlement plat relativement à $Y$,
$\mathscr F'=\mathscr F\otimes_Y\mathcal O_{Y'}$ est fidèlement plat
relativement à $Y'$.
\item[(ii)] Soient $f:X\to Y$, $g:Y\to Z$ deux morphismes, $\mathscr F$ un
$\mathcal O_X$-Module quasi-cohérent, fidèlement plat relativement à $Y$.
Pour que $g$ soit un morphisme fidèlement plat, il faut et il suffit que
$\mathscr F$ soit fidèlement plat relativement à $Z$.
\item[(iii)] Soient $f:X\to Y$, $g:Y\to Z$ deux morphismes, $\mathscr G$ un
$\mathcal O_Y$-Module quasi-cohérent. On suppose le morphisme $f$ fidèlement
plat. Pour que $\mathscr G$ soit plat (resp. fidèlement plat) relativement à
$Z$, il faut et il suffit que $f^*(\mathscr G)$ soit plat (resp. fidèlement
plat) relativement à $Z$.
\item[(iv)] Soient $f:X\to Y$, $g:Y\to Z$ deux morphismes, $x$ un point de
$X$. On suppose $f$ plat au point $x$. Pour que $g\circ f$ soit plat au point
$x$, il faut et il suffit que $g$ soit plat au point $f(x)$.
\end{enumerate}
\end{corollary}

Pour démontrer (i), on applique (2.2.10) en remplaçant $Z$ par $Y'$ et
$\mathscr G'$ par $\mathcal O_{Y'}$. Pour démontrer (ii), on applique
(2.2.10) en prenant pour morphisme $Y'\to Y$ l'identité et remplaçant
$\mathscr G'$ par $\mathcal O_Y$. Pour démontrer (iii), on applique (2.2.10)
en prenant encore l'identité pour $Y'\to Y$ et remplaçant $\mathscr F$ par
$\mathcal O_X$ et $\mathscr G'$ par $\mathscr G$. Enfin (iv) se déduit de
(ii) en remplaçant $X$ par $\operatorname{Spec}(\mathcal O_x)$,
$\mathscr F$ par $\mathcal O_X$, $Y$ par
$\operatorname{Spec}(\mathcal O_{f(x)})$ et $Z$ par
$\operatorname{Spec}(\mathcal O_{g(f(x))})$, et tenant compte de
$(0_{\mathrm I},6.6.2)$.

\begin{corollary}[2.2.12]
\label{IV.2.2.12-fr}
Soient $Y$ un schéma affine, $f:X\to Y$ un morphisme quasi-compact,
$\mathscr F$ un $\mathcal O_X$-Module quasi-cohérent. Si $\mathscr F$ est
fidèlement plat relativement à $f$, il existe un schéma affine $X'$ et un
isomorphisme local surjectif $g:X'\to X$ tels que $g^*(\mathscr F)$ soit
fidèlement plat relativement à $f\circ g$.
\end{corollary}

En effet, $X$ est quasi-compact, donc réunion finie d'ouverts affines $X_i$~;
il suffit de prendre pour $X'$ le schéma affine \emph{somme} des préschémas
induits sur les ouverts $X_i$ et
\oldpage[IV-2]{13}
pour $g:X'\to X$ le morphisme canonique~; il est clair que $g$ est fidèlement
plat, et l'hypothèse entraîne que $g^*(\mathscr F)$ est fidèlement plat
relativement à $f\circ g$, en vertu de (2.2.11, (iii)).

\begin{corollary}[2.2.13]
\label{IV.2.2.13-fr}
\begin{enumerate}
\item[(i)] Soient $f:X\to Y$, $g:Y'\to Y$ deux morphismes,
$X'=X\times_Y Y'$, $f'=f_{(Y')}:X'\to Y'$. Si $f$ est un morphisme
fidèlement plat, il en est de même de $f'$.
\item[(ii)] Si $f:X\to X'$, $g:Y\to Y'$ sont deux $S$-morphismes fidèlement
plats,
\[
f\times_S g:X\times_S Y\longrightarrow X'\times_S Y'
\]
est fidèlement plat.
\item[(iii)] Soient $f:X\to Y$, $g:Y\to Z$ deux morphismes tels que $f$ soit
fidèlement plat. Pour que $g$ soit un morphisme plat (resp. fidèlement plat),
il faut et il suffit que $g\circ f$ soit un morphisme plat (resp. fidèlement
plat).
\end{enumerate}
\end{corollary}

\begin{proposition}[2.2.14]
\label{IV.2.2.14-fr}
Soit $f:X\to Y$ un morphisme fidèlement plat et quasi-compact. Si $X$ est
localement noethérien, il en est de même de $Y$.
\end{proposition}

La question étant locale sur $Y$, on peut supposer $Y=\operatorname{Spec}(A)$
affine~; comme $f$ est quasi-compact, il résulte de (2.2.12) que l'on peut
aussi se borner au cas où $X=\operatorname{Spec}(B)$ est affine. Alors $B$
est un anneau noethérien par hypothèse, et un $A$-module fidèlement plat
(2.2.3)~; donc $A$ est noethérien $(0_{\mathrm I},6.5.2)$.

\begin{proposition}[2.2.15]
\label{IV.2.2.15-fr}
Soit $f:X\to Y$ un morphisme fidèlement plat. Si $\mathfrak S(X)$ et
$\mathfrak S(Y)$ désignent respectivement l'ensemble des sous-préschémas de
$X$ et de $Y$, l'application $Z\mapsto f^{-1}(Z)$ de $\mathfrak S(Y)$ dans
$\mathfrak S(X)$ est injective.
\end{proposition}

Comme $f$ est surjectif, on a, pour l'ensemble sous-jacent à un
sous-préschéma $Z$ de $Y$, $f(f^{-1}(Z))=Z$. D'autre part, si $U$ est un
ouvert de $Y$ contenant $Z$ et dans lequel $Z$ est fermé, $f^{-1}(U)$ est
ouvert dans $X$, $f^{-1}(Z)$ est fermé dans $f^{-1}(U)$, et la restriction
$f^{-1}(U)\to U$ de $f$ est un morphisme fidèlement plat. On peut donc se
borner à ne considérer que des sous-préschémas \emph{fermés} de $Y$. Or, si
$Z$ est un sous-préschéma fermé de $Y$ correspondant à un Idéal
quasi-cohérent $\mathscr J$ de $\mathcal O_Y$, on sait que $f^{-1}(Z)$
correspond à l'Idéal quasi-cohérent $f^*(\mathscr J)\mathcal O_X$ de
$\mathcal O_X$ (I, 4.4.5), et comme $f$ est plat,
$f^*(\mathscr J)\mathcal O_X$ s'identifie à $f^*(\mathscr J)$. Mais
l'application $\mathscr J\mapsto f^*(\mathscr J)$ de l'ensemble des Idéaux
quasi-cohérents de $\mathcal O_Y$ dans l'ensemble des Idéaux quasi-cohérents
de $\mathcal O_X$ est injective (2.2.7), d'où la conclusion.

\begin{corollary}[2.2.16]
\label{IV.2.2.16-fr}
Soient $X$, $Y$ deux $S$-préschémas~; si $S'\to S$ est un morphisme
fidèlement plat, l'application $f\mapsto f_{(S')}$ de
$\operatorname{Hom}_S(X,Y)$ dans
$\operatorname{Hom}_{S'}(X_{(S')},Y_{(S')})$ est injective.
\end{corollary}

On a $X_{(S')}\times_{S'}Y_{(S')}=(X\times_S Y)_{(S')}$ (I, 3.3.10), donc le
morphisme projection
$X_{(S')}\times_{S'}Y_{(S')}\to X\times_S Y$ est fidèlement plat (2.2.13).
Les éléments de $\operatorname{Hom}_S(X,Y)$ correspondent biunivoquement aux
sous-préschémas de $X\times_S Y$ qui sont les \emph{graphes} de ces
$S$-morphismes (I, 5.3.11), et si $Z_f$ est le graphe de
$f\in\operatorname{Hom}_S(X,Y)$, on a $Z_{f_{(S')}}=g^{-1}(Z_f)$ (I, 5.3.12).
Il suffit donc d'appliquer à $g$ la proposition (2.2.15).

\begin{proposition}[2.2.17]
\label{IV.2.2.17-fr}
Soient $A$ un anneau, $B$ une $A$-algèbre telle que $B$ soit un $A$-module
fidèlement plat et de présentation finie. Alors l'homomorphisme structural
$\varphi:A\to B$ est un isomorphisme du $A$-module $A$ sur un \emph{facteur
direct} du $A$-module $B$. Si $A$ est un anneau local, $B$ est un $A$-module
libre et il existe une base de ce module contenant l'élément unité de $B$.
\end{proposition}

\oldpage[IV-2]{14}
En vertu de Bourbaki, \emph{Alg. comm.}, chap.~II, \S~3,
n\textsuperscript{o}~3, prop.~12, il suffit de prouver la proposition lorsque
$A$ est un anneau local~; on sait alors (\emph{loc. cit.},
n\textsuperscript{o}~2, cor.~2 de la prop.~5) que $B$ est un $A$-module libre
de type fini, et la conclusion résulte de \emph{loc. cit.}, prop.~5.

\subsection{Propriétés topologiques des morphismes plats}
\label{subsection:IV.2.3-fr}

\begin{lemma}[2.3.1]
\EGAReaderTarget{ega4.c04.lem.2-3-1}{IV.2.3.1-fr}
Soient $f:X\to Y$ un morphisme quasi-compact et quasi-séparé,
$g:Y'\to Y$ un morphisme plat~; posons $X'=X\times_Y Y'$,
$f'=f_{(Y')}:X'\to Y'$. Alors, pour tout $\mathcal O_X$-Module quasi-cohérent
$\mathscr F$, l'homomorphisme canonique
\begin{equation}
\tag{2.3.1.1}
\label{IV.2.3.1.1-fr}
  g^*(f_*(\mathscr F))\to
  f'_*(\mathscr F\otimes_{\mathcal O_Y}\mathcal O_{Y'})
\end{equation}
est bijectif.
\end{lemma}

C'est un cas particulier de (III, 1.4.15) (amélioré en (1.7.21)).

\begin{proposition}[2.3.2]
\label{IV.2.3.2-fr}
Soient $S$ un préschéma, $f:X\to Y$ un $S$-morphisme quasi-compact et
quasi-séparé~; soient $Z$ le sous-préschéma de $Y$, image fermée de $X$ par
$f$ ((I, 9.5.3) et (1.7.8)), et soit $j:Z\to Y$ l'injection canonique, de
sorte que l'on a $f=j\circ g$, où $g:X\to Z$ est un morphisme
(\emph{loc. cit.}). Soit $h:S'\to S$ un morphisme plat, et posons
$f'=f_{(S')}:X_{(S')}\to Y_{(S')}$~; alors
$j'=j_{(S')}:Z_{(S')}\to Y_{(S')}$ s'identifie à l'injection canonique du
sous-préschéma de $Y_{(S')}$, image fermée de $X_{(S')}$ par $f'$.
\end{proposition}

Comme le morphisme $Y_{(S')}\to Y$ est plat (2.1.4), on peut se borner au cas
où $S=Y$ (I, 3.3.11)~; si $f=(\psi,\theta)$, on sait que $Z$ est le
sous-préschéma fermé de $S$ défini par l'Idéal (quasi-cohérent) $\mathscr J$
de $\mathcal O_S$, noyau de l'homomorphisme
$\theta:\mathcal O_S\to f_*(\mathcal O_X)$ (I, 9.5.2). Comme $h$ est un
morphisme plat, l'Idéal quasi-cohérent $\mathscr J\mathcal O_{S'}$ de
$\mathcal O_{S'}$ s'identifie au noyau de
$h^*(\theta):\mathcal O_{S'}\to h^*(f_*(\mathcal O_X))$. Or, si
$f'=(\psi',\theta')$, on vérifie facilement (par exemple en se ramenant au
cas où $Y$ et $Y'$ sont affines et utilisant (I, 2.2.4)) que
$\theta':\mathcal O_{S'}\to f'_*(\mathcal O_{X'})$ est composé de
l'homomorphisme canonique (2.3.1.1)
$h^*(f_*(\mathcal O_X))\to f'_*(\mathcal O_{X'})$ et de $h^*(\theta)$~; la
conclusion résulte donc de (2.3.1) et de (I, 9.5.2), puisque $f'$ est
quasi-compact et quasi-séparé (1.1.2 et 1.2.2).

\begin{env}[2.3.3]
\label{IV.2.3.3-fr}
Nous dirons qu'un morphisme $f:X\to Y$ est \emph{quasi-plat} s'il existe un
$\mathcal O_X$-Module quasi-cohérent $\mathscr F$ \emph{de type fini} qui soit
$f$-plat et dont le support soit égal à $X$. Nous dirons que $f$ est
\emph{quasi-fidèlement plat} s'il est quasi-plat et surjectif. Tout morphisme
plat (resp. fidèlement plat) est quasi-plat (resp. quasi-fidèlement plat), car
alors $\mathscr F=\mathcal O_X$ vérifie les conditions précédentes.
\end{env}

Il résulte aussitôt de (2.1.4) et de (I, 9.1.13), que si $f$ est quasi-plat,
alors, pour tout morphisme $Y'\to Y$, le morphisme
$f_{(Y')}:X\times_Y Y'\to Y'$ est quasi-plat. De même (compte tenu de
(I, 3.5.2)), si $f$ est quasi-fidèlement plat, il en est de même de
$f_{(Y')}$.

\begin{proposition}[2.3.4]
\EGAReaderTarget{ega4.c04.prop.2-3-4}{IV.2.3.4-fr}
Soit $f:X\to Y$ un morphisme quasi-plat (2.3.3). Alors $f$ possède les
propriétés suivantes (qui sont équivalentes en vertu de (1.10.4))~:
\begin{enumerate}
  \item[(i)] Pour tout $x\in X$ et toute générisation $y'$ de $y=f(x)$, il
  existe une générisation $x'$ de $x$ telle que $f(x')=y'$.
  \item[(ii)] Pour tout $x\in X$, l'image par $f$ de
  $\operatorname{Spec}(\mathcal O_{X,x})$ est
  $\operatorname{Spec}(\mathcal O_{Y,y})$.
\oldpage[IV-2]{15}
  \item[(iii)] Pour toute partie fermée irréductible $Y'$ de $Y$, toute
  composante irréductible de $f^{-1}(Y')$ domine $Y'$.
\end{enumerate}
\end{proposition}

Il suffit par exemple de prouver (ii). Par hypothèse, il y a un
$\mathcal O_X$-Module quasi-cohérent de type fini $\mathscr F$, qui est
$f$-plat et tel que $\operatorname{Supp}(\mathscr F)=X$. Pour tout $x\in X$,
$\mathscr F_x$ est alors un $\mathcal O_x$-module de type fini, non réduit à
0, et en outre $\mathscr F_x$ est un $\mathcal O_y$-module \emph{plat}, pour
l'homomorphisme $\rho:\mathcal O_y\to\mathcal O_x$. Comme ce dernier est
local et que $\mathscr F_x\neq0$, le lemme de Nakayama montre que
$\mathscr F_x\otimes_{\mathcal O_y}k(y)\neq0$, donc $\mathscr F_x$ est un
$\mathcal O_y$-module \emph{fidèlement plat} $(0_{\mathrm I},6.4.1)$. Il
en résulte que pour tout idéal premier $\mathfrak q$ de $\mathcal O_y$, il
existe un idéal premier $\mathfrak p$ de $\mathcal O_x$ tel que
$\mathfrak q=\rho^{-1}(\mathfrak p)$ $(0_{\mathrm I},6.5.1)$, ce qui
prouve (ii).

\begin{corollary}[2.3.5]
\label{IV.2.3.5-fr}
Soit $f:X\to Y$ un morphisme vérifiant les conditions équivalentes (i), (ii),
(iii) de (2.3.4) (en particulier un morphisme quasi-plat ou un morphisme
ouvert (1.10.4)).
\begin{enumerate}
  \item[(i)] Soient $Z$, $Z'$ deux parties fermées irréductibles de $Y$ telles
  que $Z\subset Z'$, et soit $T$ une composante irréductible de $f^{-1}(Z)$~;
  alors il existe une composante irréductible $T'$ de $f^{-1}(Z')$ contenant
  $T$ (et dominant $Z'$).
  \item[(ii)] Pour toute composante irréductible $T$ de $X$,
  $\overline{f(T)}$ est une composante irréductible de $Y$.
  \item[(iii)] Supposons $Y$ irréductible, et, si $y$ est son point générique,
  supposons que $f^{-1}(y)$ soit irréductible. Alors $X$ est irréductible.
\end{enumerate}
\end{corollary}

(i) Il suffit d'appliquer (2.3.4, (i)) en prenant pour $x$ le point générique
de $T$ ($y=f(x)$ étant alors le point générique de $Z$) et pour $y'$ le point
générique de $Z'$.

(ii) Il est clair que $\overline{f(T)}=Z$ est irréductible, et en vertu de (i),
$Z$ ne peut être strictement contenue dans une partie fermée irréductible $Z'$
de $Y$.

(iii) En vertu de (ii), toute composante irréductible de $X$ domine $Y$, donc
rencontre $f^{-1}(y)$~; la conclusion résulte alors de
$(0_{\mathrm I},2.1.8)$.

\begin{proposition}[2.3.6]
\label{IV.2.3.6-fr}
Soit $Y$ un préschéma dont l'ensemble des composantes irréductibles est
\emph{localement fini} (ce qui est le cas si l'espace sous-jacent à $Y$ est
\emph{localement noethérien} (cf. I, 6.1.9)).
\begin{enumerate}
  \item[(i)] Toute partie fermée $W$ de $Y$ \emph{stable par générisation}
  $(0_{\mathrm I},2.1.2)$ est ouverte. En particulier, toute composante
  connexe de $Y$ est ouverte.
  \item[(ii)] Soit $f:X\to Y$ un morphisme fermé vérifiant en outre les
  conditions équivalentes (i), (ii), (iii) de (2.3.4) (ce qui sera le cas si
  $f$ est quasi-plat ou ouvert). Alors l'image par $f$ de toute composante
  connexe $C$ de $X$ est une composante connexe de $Y$.
  \item[(iii)] Soit $f:X\to Y$ un morphisme vérifiant les conditions
  équivalentes (i), (ii), (iii) de (2.3.4) et tel en outre que pour tout
  $y\in Y$, l'ensemble $f^{-1}(y)$ soit fini (ce qui sera le cas si $f$ est
  quasi-fini ou radiciel). Alors l'ensemble des composantes irréductibles de
  $X$ est localement fini.
\end{enumerate}
\end{proposition}

\oldpage[IV-2]{16}
(i) Si $y\in W$, les points génériques $\eta_i$ ($1\leq i\leq m$) des
composantes irréductibles $Y_i$ de $Y$ contenant $y$ appartiennent par
hypothèse à $W$, donc, puisque $W$ est fermé, ces composantes elles-mêmes sont
contenues dans $W$~; comme il y a par hypothèse un voisinage $U$ de $y$ tel que
$U$ soit réunion des $U\cap Y_i$ $(0_{\mathrm I},2.1.6)$, on a
$U\subset W$, donc $W$ est ouvert. Comme pour tout $y\in Y$,
$\overline{\{y\}}$ est connexe, une composante connexe de $Y$ est stable par
générisation, donc la seconde assertion résulte aussitôt de la première.

(ii) Comme $C$ est fermé dans $X$, $f(C)$ est fermé dans $Y$ par hypothèse. En
outre, comme $C$ est stable par générisation, l'hypothèse sur $f$ entraîne que
$f(C)$ est stable par générisation~; donc, en vertu de (i), $f(C)$ est ouvert
et fermé, et comme il est connexe c'est une composante connexe de $Y$.

(iii) Soit $x\in X$~; par hypothèse, il y a un voisinage ouvert $U$ de
$y=f(x)$ ne rencontrant qu'un nombre fini de composantes irréductibles $Y_i$
de $Y$ ($1\leq i\leq n$)~; soit $y_i$ le point générique de $Y_i$
($1\leq i\leq n$). Pour toute composante irréductible $Z$ de $X$ rencontrant
$f^{-1}(U)$, le point générique $z$ de $Z$ est nécessairement dans l'un des
ensembles $f^{-1}(y_i)$ (2.3.4). Comme chacun de ces ensembles est fini par
hypothèse, cela prouve notre assertion.

\begin{proposition}[2.3.7]
\label{IV.2.3.7-fr}
Soient $f:X\to Y$ un morphisme, $g:Y'\to Y$ un morphisme quasi-plat,
$X'=X\times_Y Y'$, $f'=f_{(Y')}:X'\to Y'$.
\begin{enumerate}
  \item[(i)] Si $f$ est quasi-compact et dominant, il en est de même de $f'$.
  \item[(ii)] Si toute composante irréductible de $X$ domine une composante
  irréductible de $Y$, alors toute composante irréductible de $X'$ domine une
  composante irréductible de $Y'$.
\end{enumerate}
\end{proposition}

Désignons par $g':X'\to X$ la projection canonique, qui est un morphisme
quasi-plat (2.3.3).

(i) On sait déjà (1.1.2) que $f'$ est quasi-compact~; en outre, si $y'$ est un
point maximal (1.1.4) de $Y'$, $y=g(y')$ est point maximal de $Y$, comme il
résulte de (2.3.4, (iii)). Par hypothèse (1.1.5), il existe $x\in X$ tel que
$f(x)=y$, donc (I, 3.4.7), il existe $x'\in X'$ tel que $f'(x')=y'$.

(ii) Soit $x'$ un point maximal de $X'$, et soit $x=g'(x')$~; il résulte de
(2.3.4, (ii)) que $x$ est point maximal de $X$, et par hypothèse $y=f(x)$ est
point maximal de $Y$. Posons $y'=f'(x')$, de sorte que $g(y')=y$~; il s'agit
de montrer que $y'$ est point maximal de $Y'$. En vertu de (I, 5.1.7) et
(2.3.3), on peut se borner au cas où $X$ et $Y$ sont \emph{réduits}, donc
$\mathcal O_x$ et $\mathcal O_y$ sont des \emph{corps}~; en outre, en vertu de
(I, 3.6.5) appliqué deux fois et de (I, 2.4.4), on peut se borner au cas où
$Y=\operatorname{Spec}(\mathcal O_y)$ et
$Y'=\operatorname{Spec}(\mathcal O_{y'})$. Alors $f$ est un morphisme
\emph{plat} puisque $\mathcal O_y$ est un corps (2.1.2), et il en est de même
de $f'$ (2.1.4)~; donc il résulte de (2.3.4, (ii)) que $y'=f'(x')$ est point
maximal de $Y'$.

\begin{corollary}[2.3.8]
\label{IV.2.3.8-fr}
\emph{(Zariski).} Soient $A$, $B$ deux anneaux locaux noethériens,
$\mathfrak m$, $\mathfrak n$ leurs idéaux maximaux respectifs,
$\varphi:A\to B$ un homomorphisme local. On suppose vérifiées les hypothèses
suivantes~:
\begin{enumerate}
  \item[$1^\circ$] $B$ est une $A$-algèbre essentiellement de type fini
  (1.3.8).
  \item[$2^\circ$] Le complété $\widehat A$ de $A$ pour la topologie
  $\mathfrak m$-adique est intègre.
  \item[$3^\circ$] $\varphi$ est injectif.
\end{enumerate}
Alors la topologie $\mathfrak m$-adique de $A$ est induite par la topologie
$\mathfrak n$-adique de $B$.
\end{corollary}

\oldpage[IV-2]{17}
Posons $B'=B\otimes_A\widehat A$~; en vertu de $1^\circ$, $B'$ est de la forme
$S^{-1}(C\otimes_A\widehat A)$, où $C$ est une $A$-algèbre de type fini et
$S$ est une partie multiplicative de $C$, donc $B'$ est un anneau noethérien.
Comme $A$ s'identifie à un sous-anneau de $\widehat A$
$(0_{\mathrm I},7.3.5)$, $A$ est intègre par $2^\circ$. L'hypothèse
$3^\circ$ entraîne alors qu'il existe un idéal premier $\mathfrak q$ de $B$
induisant l'idéal 0 de $A$ $(0_{\mathrm I},1.5.8)$, et par suite
l'homomorphisme local $A\to B/\mathfrak q$ est injectif. On peut donc se borner
à prouver la conclusion de (2.3.8) en ajoutant l'hypothèse que $B$ est un
anneau local \emph{intègre}. Appliquons (2.3.7, (ii)) à
$Y=\operatorname{Spec}(A)$, $X=\operatorname{Spec}(B)$,
$Y'=\operatorname{Spec}(\widehat A)$ et $X'=\operatorname{Spec}(B')$~; comme
le morphisme $Y'\to Y$ est plat et que $X$ (qui est intègre) domine le schéma
intègre $Y$, toute composante irréductible de $X'$ domine le schéma (intègre
par hypothèse) $Y'$. Si $y$, $x$, $y'$ sont les points fermés de $Y$, $X$ et
$Y'$ respectivement, il y a un point $x'\in X'$ (d'ailleurs unique) au-dessus
de $x$ et $y'$ (I, 3.4.9) et $\operatorname{Spec}(\mathcal O_{x'})$ domine
donc $\operatorname{Spec}(\mathcal O_{y'})$~; on a par suite un diagramme
commutatif d'homomorphismes locaux d'anneaux locaux noethériens
\[
\xymatrix{
  B=\mathcal O_x \ar[r] & \mathcal O_{x'}\\
  A=\mathcal O_y \ar[u]^{\varphi} \ar[r]_{u} &
  \mathcal O_{y'}=\widehat A \ar[u]_{v}
}
\]
tel que $u$ et $v$ soient injectifs (I, 1.2.7)~; identifiant $A$ et
$\widehat A$ à des sous-anneaux de $\mathcal O_{x'}$, et désignant par
$\mathfrak r$ l'idéal maximal de $\mathcal O_{x'}$, l'intersection des idéaux
$\mathfrak r^k\cap\widehat A$ est donc nulle $(0_{\mathrm I},7.3.5)$~;
comme $\widehat A$ est complet et que ces idéaux sont ouverts dans
$\widehat A$, cela entraîne (Bourbaki, \emph{Alg. comm.}, chap.~III, \S~2,
n\textsuperscript{o}~7, prop.~8) que la topologie de $\widehat A$ est induite
par la topologie $\mathfrak r$-préadique de $\mathcal O_{x'}$~; \emph{a
fortiori} il en est de même de la topologie de $A$
$(0_{\mathrm I},7.3.5)$. D'ailleurs on a
$\mathfrak n^k\cap A\subset\mathfrak r^k\cap A$, donc la topologie
$\mathfrak n$-préadique de $B$ induit sur $A$ une topologie plus fine que la
topologie $\mathfrak m$-préadique~; mais comme
$\mathfrak m^k\subset\mathfrak n^k\cap A$, ces deux topologies sont
identiques. C.Q.F.D.

\begin{remark}[2.3.9]
\label{IV.2.3.9-fr}
Nous verrons plus loin (7.8.3, (vii)) que pour les anneaux locaux noethériens
$A$ les plus usuels en Géométrie algébrique, l'hypothèse que $A$ est intègre
et intégralement clos implique qu'il en est de même de $\widehat A$. C'est
pourquoi, en Géométrie algébrique sur un corps de base, on énonce généralement
(2.3.8) sous l'hypothèse que $A$ est \emph{intègre et intégralement clos}.
\end{remark}

\begin{theorem}[2.3.10]
\label{IV.2.3.10-fr}
Soit $f:X\to Y$ un morphisme quasi-plat (2.3.3). Alors, pour toute partie
pro-constructible (1.9.4) $Z$ de $Y$, on a
$\overline{f^{-1}(Z)}=f^{-1}(\overline Z)$.
\end{theorem}

Comme $f$ est continue, on a
$\overline{f^{-1}(Z)}\subset f^{-1}(\overline Z)$ et tout revient à prouver
que pour tout $x\in X$ tel que $f(x)\in\overline Z$, $x$ est adhérent à
$f^{-1}(Z)$~; il est clair que la question est locale sur $Y$, donc on peut
supposer $Y$ affine. En vertu de l'hypothèse, il existe un schéma affine $Y'$
et un morphisme $g:Y'\to Y$ tels que $g(Y')=Z$ (1.9.5, (ix)). Soit $Y_1$
l'image fermée de $Y'$ par $g$ (I, 9.5.3), et soit $X_1$ le sous-préschéma
fermé $f^{-1}(Y_1)$ de $X$~; si $f_1:X_1\to Y_1$ est la restriction de $f$ à
$X_1$, on sait que $f_1$ est quasi-plat (2.3.3)~; on peut par suite remplacer
$X$, $Y$ par $X_1$, $Y_1$ respectivement, autrement dit supposer que $g$ est
\emph{dominant}. Posons alors $X'=X\times_Y Y'$, et soient $f'$ et $g'$ les
projections de $X'$ dans $Y'$ et $X$ respectivement, de sorte que l'on a le
diagramme commutatif
\[
\xymatrix{
  X \ar[d]_f & X' \ar[l]_{g'} \ar[d]^{f'} \\
  Y & Y' \ar[l]^g
}
\]
Comme $f$ est quasi-plat, $g$ quasi-compact et dominant, il résulte de (2.3.7)
(où les rôles de $f$ et $g$ sont intervertis) que $g'$ est un morphisme
dominant, ce qui démontre le théorème.

\begin{corollary}[2.3.11]
\label{IV.2.3.11-fr}
Soient $f$ un morphisme quasi-plat et quasi-compact, $F$ une partie fermée de
$X$ telle que $F=f^{-1}(f(F))$~; alors on a
$F=f^{-1}(\overline{f(F)})$.
\end{corollary}

Soient $Y'$ le sous-préschéma réduit de $X$ ayant $F$ comme espace sous-jacent
(I, 5.2.1) et soit $j:Y'\to X$ l'injection canonique~; alors $f\circ j$ est
quasi-compact (1.1.2),
\oldpage[IV-2]{18}
donc $Z=f(F)$ est pro-constructible dans $Y$ (1.9.5, (vii)), et le corollaire
résulte de ce que $F$ est fermé.

On peut encore écrire le résultat de (2.3.11) sous la forme
$F=f^{-1}(f(X)\cap\overline{f(F)})$, autrement dit, les ensembles fermés du
sous-espace $f(X)$ de $Y$ sont les parties $Z\subset f(X)$ telles que
$f^{-1}(Z)$ soit fermé dans $X$~; cela signifie aussi que \emph{la topologie
induite par $Y$ sur $f(X)$ est quotient de celle de $X$ par la relation
d'équivalence définie par $f$}. En particulier~:

\begin{corollary}[2.3.12]
\label{IV.2.3.12-fr}
Soit $f:X\to Y$ un morphisme quasi-fidèlement plat (2.3.3) et quasi-compact.
Alors la topologie de $Y$ est quotient de celle de $X$ par la relation
d'équivalence définie par $f$ (autrement dit, pour que $Z\subset Y$ soit
ouvert (resp. fermé) dans $Y$, il faut et il suffit que $f^{-1}(Z)$ soit
ouvert (resp. fermé) dans $X$).
\end{corollary}

En effet, on a alors $f(X)=Y$.

\begin{corollary}[2.3.13]
\label{IV.2.3.13-fr}
Soient $X$, $Y$ deux $S$-préschémas, $f:X\to Y$ un $S$-morphisme fidèlement
plat et quasi-compact. Pour que $Y$ soit séparé au-dessus de $S$, il faut et il
suffit que l'immersion canonique
$j:X\times_Y X\to X\times_S X$ (I, 5.3.10) soit \emph{fermée}.
\end{corollary}

Notons pour cela que l'on a le diagramme commutatif (I, 5.3.5)
\[
\xymatrix{
  X\times_Y X \ar[r]^j \ar[d]_{\pi} &
  X\times_S X \ar[d]^{f\times_S f}\\
  Y \ar[r]_{\Delta_Y} & Y\times_S Y
}
\]
identifiant $X\times_Y X$ au produit des $(Y\times_S Y)$-préschémas $Y$ et
$X\times_S X$. Comme $f$ est surjectif, il en est de même de $\pi$ et de
$f\times_S f$, donc la diagonale $\Delta_Y(Y)$ a pour image réciproque par
$f\times_S f$ l'image $j(X\times_Y X)$ (I, 3.4.8). Comme $f\times_S f$ est
fidèlement plat et quasi-compact (1.1.2 et 2.2.13), il suffit d'appliquer
(2.3.12) à ce morphisme.

\begin{corollary}[2.3.14]
\label{IV.2.3.14-fr}
Soit $f:X\to Y$ un morphisme quasi-fidèlement plat et quasi-compact, et soit
$Z$ une partie de $Y$. Pour que $Z$ soit une partie localement fermée
pro-constructible de $Y$, il faut et il suffit que $f^{-1}(Z)$ soit une partie
localement fermée pro-constructible de $X$.
\end{corollary}

On sait déjà (1.9.12) que, pour que $Z$ soit une partie pro-constructible de
$Y$, il faut et il suffit que $f^{-1}(Z)$ soit une partie pro-constructible de
$X$. La condition est évidemment nécessaire. Pour prouver qu'elle est
suffisante, considérons le sous-préschéma fermé réduit $Y_1$ de $Y$ ayant
$\overline Z$ pour espace sous-jacent et, soit $X_1$ son image réciproque par
$f$, qui a pour espace sous-jacent
$f^{-1}(\overline Z)=\overline{f^{-1}(Z)}$ en vertu de (2.3.10). Comme
$f^{-1}(Z)$ est localement fermé dans $X$, il est ouvert dans
$\overline{f^{-1}(Z)}$, donc dans $X_1$~; or, le morphisme $f_1:X_1\to Y_1$
déduit de $f$ par restriction est quasi-fidèlement plat et quasi-compact
(1.1.2 et 2.3.3)~; on déduit donc de (2.3.12) que $Z$ est ouvert dans
$\overline Z$, ce qui montre que $Z$ est localement fermé dans $Y$.

\begin{remark}[2.3.15]
\label{IV.2.3.15-fr}
Il ne suffit pas, dans (2.3.12), de supposer seulement que $f$ est fidèlement
plat. Par exemple, prenons pour $Y$ une courbe algébrique réduite irréductible
\oldpage[IV-2]{19}
(II, 7.4.2), pour $X$ le préschéma somme des schémas
$\operatorname{Spec}(\mathcal O_y)$, où $y$ parcourt $Y$ (I, 3.1), pour $f$ le
morphisme canonique, qui est fidèlement plat (I, 2.4.2)~; si $\eta$ désigne le
point générique de $Y$, $Z=\{\eta\}$ n'est pas ouvert dans $Y$ (II, 7.4.3),
mais $f^{-1}(Z)=\operatorname{Spec}(\mathcal O_\eta)$ puisque
$\mathcal O_\eta$ est un corps, et par suite $f^{-1}(Z)$ est ouvert dans $X$.
\end{remark}

\subsection{Morphismes universellement ouverts et morphismes plats}

\begin{env}[2.4.1]
\label{IV.2.4.1-fr}
Nous avons déjà défini (II, 5.4.9) la notion de morphisme
\emph{universellement fermé}~; de la même manière, on pose les définitions
suivantes~:
\end{env}

\begin{definition}[2.4.2]
\label{IV.2.4.2-fr}
On dit qu'un morphisme de préschémas $f:X\to Y$ est \emph{universellement
ouvert} (resp. \emph{universellement bicontinu}, resp. un
\emph{homéomorphisme universel}) si, pour tout morphisme $g:Y'\to Y$, le
morphisme $f_{(Y')}:X\times_Y Y'\to Y'$ est ouvert (resp. un homéomorphisme
sur son image, resp. un homéomorphisme sur $Y'$).
\end{definition}

Nous verrons plus loin (14.3.2) que lorsque $Y$ est localement noethérien, la
définition de morphisme universellement ouvert donnée ici est équivalente à la
définition (III, 4.3.9) pour les morphismes de type fini~; le lecteur pourra
vérifier que nous n'utilisons pas cette dernière définition avant le \S~14.

\begin{proposition}[2.4.3]
\EGAReaderTarget{ega4.c04.prop.2-4-3}{IV.2.4.3-fr}
\begin{enumerate}
\item[(i)] Une immersion (resp. une immersion ouverte, resp. une immersion
fermée) est universellement bicontinue (resp. universellement ouverte, resp.
universellement fermée).
\item[(ii)] Le composé de deux morphismes universellement ouverts (resp.
universellement fermés, resp. universellement bicontinus, resp. deux
homéomorphismes universels) l'est aussi.
\item[(iii)] Si $f:X\to Y$ est un $S$-morphisme universellement ouvert (resp.
universellement fermé, resp. universellement bicontinu, resp. un
homéomorphisme universel), il en est de même de
$f_{(S')}:X_{(S')}\to Y_{(S')}$ pour tout changement de base $S'\to S$.
\item[(iv)] Si $f:X\to X'$, $g:Y\to Y'$ sont deux $S$-morphismes
universellement ouverts (resp. universellement fermés, resp. universellement
bicontinus, resp. deux homéomorphismes universels), il en est de même de
\[
f\times_S g:X\times_S Y\longrightarrow X'\times_S Y'.
\]
\item[(v)] Soient $f:X\to Y$, $g:Y\to Z$ deux morphismes tels que $f$ soit
surjectif~; si $g\circ f$ est universellement ouvert (resp. universellement
fermé, resp. universellement bicontinu, resp. un homéomorphisme universel), il
en est de même de $g$.
\item[(vi)] Pour que $f$ soit un morphisme universellement ouvert (resp.
universellement fermé, resp. universellement bicontinu, resp. un
homéomorphisme universel), il faut et il suffit que $f_{\mathrm{red}}$ le soit.
\item[(vii)] Soit $(U_\alpha)$ un recouvrement ouvert de $Y$. Pour que
$f:X\to Y$ soit universellement ouvert (resp. universellement fermé, resp.
universellement bicontinu, resp. un homéomorphisme universel), il faut et il
suffit que pour tout $\alpha$, sa restriction
$f^{-1}(U_\alpha)\to U_\alpha$ le soit.
\end{enumerate}
\end{proposition}

L'assertion (i) résulte de (I, 4.3.2). L'assertion (ii) résulte aussitôt des
définitions, et il en est de même de (iii) en se ramenant au cas où $Y=S$,
$Y'=S'$, ce que l'on peut faire grâce à (I, 3.3.11)~; on sait que (iv)
découle de (ii) et (iii) (I, 3.5.1). Pour prouver (v), notons que pour tout
morphisme $Z'\to Z$, $f_{(Z')}:X_{(Z')}\to Y_{(Z')}$ est surjectif
(I, 3.5.2)~; on peut donc se borner à prouver que si $g\circ f$ est ouvert
(resp. fermé, resp. un
\oldpage[IV-2]{20}
homéomorphisme sur son image, resp. un homéomorphisme bijectif), il en est de
même de $g$, et il s'agit donc d'une question purement topologique. Pour le
cas où $g\circ f$ est ouvert (resp. fermé), le fait que $g$ est alors ouvert
(resp. fermé) résulte de Bourbaki, \emph{Top. gén.}, chap.~I\textsuperscript{er},
3\textsuperscript{e} éd., \S~5, n\textsuperscript{o}~1, prop.~1~; pour les
deux autres cas, on peut se borner à supposer que
$g(f(X))=g(Y)=Z$, autrement dit au cas où $g\circ f$ est un homéomorphisme de
$X$ sur $Z$~; comme $f$ est surjectif, $g$ est nécessairement bijectif, et
comme $g$ est une application continue ouverte d'après ce qui précède, $g$
est bien un homéomorphisme de $Y$ sur $Z$.

Pour prouver (vi), notons que dire qu'un morphisme $g$ est ouvert (resp.
fermé, resp. un homéomorphisme sur son image, resp. un homéomorphisme
bijectif) équivaut à dire que $g_{\mathrm{red}}$ possède la même propriété.
D'autre part (I, 5.1.8), pour tout morphisme $Y'\to Y$, on a
\[
\bigl(X_{\mathrm{red}}\times_{Y_{\mathrm{red}}}Y'_{\mathrm{red}}\bigr)_{\mathrm{red}}
=(X\times_Y Y')_{\mathrm{red}},
\]
donc la remarque précédente montre que si $f_{\mathrm{red}}$ est
universellement ouvert (resp. universellement fermé, resp. universellement
bicontinu, resp. un homéomorphisme universel), il en est de même de $f$. La
réciproque se prouve de même, en notant ici que pour tout morphisme
$Y''\to Y_{\mathrm{red}}$, on a
\[
\bigl(X_{\mathrm{red}}\times_{Y_{\mathrm{red}}}Y''\bigr)_{\mathrm{red}}
=(X\times_Y Y'')_{\mathrm{red}}
\quad\text{(I, 5.1.8).}
\]

Enfin, la nécessité de (vii) résulte aussitôt de (iii). Inversement, supposons
remplie la condition (vii), et soit $g:Y'\to Y$ un morphisme~; alors les
$g^{-1}(U_\alpha)=U'_\alpha$ forment un recouvrement ouvert de $Y'$ et si
l'on note $f_\alpha$ la restriction $f^{-1}(U_\alpha)\to U_\alpha$ de $f$,
et $f'$ le morphisme $f_{(Y')}$, la restriction
$f'^{-1}(U'_\alpha)\to U'_\alpha$ n'est autre que
$(f_\alpha)_{(U'_\alpha)}$. On peut donc se borner à prouver que $f'$ est
ouvert (resp. fermé, resp. un homéomorphisme sur son image, resp. un
homéomorphisme sur $Y'$), ce qui est immédiat.

\begin{proposition}[2.4.4]
\label{IV.2.4.4-fr}
Un morphisme universellement bicontinu $f:X\to Y$ est radiciel (donc séparé
(I, 7.7.1)).
\end{proposition}

En effet, $f$ est universellement injectif par hypothèse (I, 3.5.11).

\begin{proposition}[2.4.5]
\EGAReaderTarget{ega4.c04.prop.2-4-5}{IV.2.4.5-fr}
\begin{enumerate}
\item[(i)] Un morphisme $f:X\to Y$ qui est entier, surjectif et radiciel est
un homéomorphisme universel.
\item[(ii)] Inversement, supposons $Y$ localement noethérien. Alors, si un
morphisme de type fini $f:X\to Y$ est un homéomorphisme universel, il est
fini, surjectif et radiciel.
\end{enumerate}
\end{proposition}

(i) Il suffit d'observer que les trois propriétés de $f$ se conservent par
changement de base (I, 3.5.2, I, 3.5.7 et II, 6.1.5), et comme un morphisme
entier est fermé (II, 6.1.10), il est clair que $f$ est un homéomorphisme de
$X$ sur $Y$.

(ii) Comme $f$ est de type fini et universellement fermé par hypothèse et
qu'il est séparé par (2.4.4), il est propre (II, 5.4.1), et pour tout
$y\in Y$, $f^{-1}(y)$ est réduit à un seul élément~; donc (III, 4.4.2) $f$
est fini~; il est clair que $f$ est surjectif, et il est radiciel puisqu'il
est universellement injectif (I, 3.5.11).

\begin{theorem}[2.4.6]
\label{IV.2.4.6-fr}
Soit $f:X\to Y$ un morphisme quasi-plat (2.3.3) et localement de présentation
finie. Alors $f$ est universellement ouvert. En particulier, un morphisme plat
localement de présentation finie est universellement ouvert.
\end{theorem}

On sait que pour tout changement de base $Y'\to Y$, $f_{(Y')}$ est
quasi-plat (2.3.3) et localement de présentation finie (1.4.3, (iii)). Il
suffit donc de prouver que $f$ est un
\oldpage[IV-2]{21}
morphisme ouvert. Mais cela résulte du critère (1.10.4) pour les morphismes
localement de présentation finie, les conditions \emph{b)}, \emph{b$'$)}, et
\emph{c)} de (1.10.4) n'étant autre que les conditions (i), (ii), (iii) de
(2.3.4).

\begin{corollary}[2.4.7]
\label{IV.2.4.7-fr}
Pour tout préschéma $Y$, le morphisme structural $\mathbf V_Y^n\to Y$ (où
\[
\mathbf V_Y^n
=Y\otimes_{\mathbf Z}\operatorname{Spec}
  \bigl(\mathbf Z[T_1,\ldots,T_n]\bigr),
\]
aussi noté $Y[T_1,\ldots,T_n]$) est universellement ouvert.
\end{corollary}

En effet, pour $Y=\operatorname{Spec}(A)$,
$Y[T_1,\ldots,T_n]=\operatorname{Spec}(A[T_1,\ldots,T_n])$, et
$A[T_1,\ldots,T_n]$ est une $A$-algèbre libre de présentation finie.

\begin{remark}[2.4.8]
\label{IV.2.4.8-fr}
\begin{enumerate}
\item[(i)] On notera qu'un morphisme $f:X\to Y$ fidèlement plat et
quasi-compact \emph{n'est pas nécessairement ouvert}, même lorsque $X$ et $Y$
sont noethériens. Prenons par exemple pour $Y$ une courbe algébrique réduite
irréductible (II, 7.4) de point générique $y$, et soit $X$ le préschéma somme
$Y\amalg\operatorname{Spec}(k(y))$, $f$ étant le morphisme canonique~; il est
clair que $f$ est plat et surjectif, donc fidèlement plat, et quasi-compact,
mais l'image par $f$ de la partie ouverte $\operatorname{Spec}(k(y))$ de $X$
est l'ensemble $\{y\}$, qui n'est pas ouvert dans $Y$ (II, 7.4.3).
\item[(ii)] Pour tout préschéma $X$, le morphisme canonique
$f:X_{\mathrm{red}}\to X$ est une immersion fermée et un
\emph{homéomorphisme universel} (2.4.4, (vi))~; mais lorsque $X$ est
localement noethérien, $f$ n'est plat que si $X$ est réduit, donc
$f=1_X$ (2.2.17).
\end{enumerate}
\end{remark}

\begin{proposition}[2.4.9]
\label{IV.2.4.9-fr}
Soit $Y$ un préschéma discret. Alors tout morphisme $f:X\to Y$ est
universellement ouvert.
\end{proposition}

La question étant locale sur $Y$ (2.4.3, (vii)), on peut se borner au cas où
l'espace sous-jacent $Y$ est réduit à un point~; remplaçant $f$ par
$f_{\mathrm{red}}$ (2.4.3, (vi)), on peut en outre supposer que $Y$ est le
spectre d'un \emph{corps} $k$~; d'autre part, pour tout changement de base
$Y'\to Y$, les ouverts de $X'=X\times_Y Y'$ images réciproques des ouverts
affines de $X$ recouvrent $X'$, donc on peut supposer
$X=\operatorname{Spec}(B)$ affine, $B$ étant une $k$-algèbre. Il s'agit donc
de prouver que pour toute $k$-algèbre $A'$, si l'on pose
$B'=B\otimes_k A'$, l'image par
$f':\operatorname{Spec}(B')\to\operatorname{Spec}(A')$ de tout ouvert de la
forme $U=D(t)$ ($t\in B'$) est ouverte dans $\operatorname{Spec}(A')$. Or,
$B$ est limite inductive de la famille filtrante croissante $(B_\alpha)$ de
ses sous-$k$-algèbres \emph{de type fini}, donc (le foncteur
$\varinjlim$ commutant au produit tensoriel), $B'$ est limite inductive des
$A'$-algèbres $B_\alpha\otimes_k A'=B'_\alpha$~; il existe $\alpha$ tel que
$t$ soit image dans $B'$ d'un élément $t_\alpha\in B'_\alpha$, donc $D(t)$
est l'image réciproque par le morphisme canonique
$u_\alpha:\operatorname{Spec}(B')\to\operatorname{Spec}(B'_\alpha)$ de
l'ouvert $U_\alpha=D(t_\alpha)$ (I, 1.2.2.2). Mais comme $k$ est un corps,
$B_\alpha$ est une $k$-algèbre de présentation finie et un $k$-module plat,
donc $B'_\alpha$ est une $A'$-algèbre de présentation finie et un $A'$-module
plat~; on conclut donc de (2.4.6) que l'image de $U_\alpha$ par
$f'_\alpha:\operatorname{Spec}(B'_\alpha)\to\operatorname{Spec}(A')$ est
\emph{ouverte} dans $\operatorname{Spec}(A')$. Tout revient par suite à voir
que $f'(U)=f'_\alpha(U_\alpha)$~; il est clair que
$f'(U)\subset f'_\alpha(U_\alpha)$ et il reste donc à voir que pour tout point
$y'\in f'_\alpha(U_\alpha)$, l'intersection
$V=U\cap f'^{-1}(y')$ n'est pas vide. Or on a
$V=u_\alpha^{-1}(V_\alpha)$, où
$V_\alpha=U_\alpha\cap f_\alpha'^{-1}(y')$~; autrement dit, $V_\alpha$ est
un ouvert non vide (par définition de $y'$) du préschéma
\[
\operatorname{Spec}(B'_\alpha\otimes_{A'}k(y'))
=\operatorname{Spec}(B_\alpha\otimes_k k(y'))
\]
et $V$ est son image réciproque par le morphisme
\[
v:\operatorname{Spec}(B\otimes_k k(y'))
\longrightarrow\operatorname{Spec}(B_\alpha\otimes_k k(y')).
\]
Comme $B_\alpha$ est une sous-algèbre de $B$ et que $k$ est un corps,
l'homomorphisme
$B_\alpha\otimes_k k(y')\to B\otimes_k k(y')$
\oldpage[IV-2]{22}
est injectif par platitude, donc le morphisme $v$ est \emph{dominant}
(I, 1.2.7), ce qui achève la démonstration.

\begin{corollary}[2.4.10]
\label{IV.2.4.10-fr}
Soient $k$ un corps, $X$, $Y$ des $k$-préschémas~; alors le morphisme
projection $X\times_k Y\to X$ est universellement ouvert. En particulier,
pour toute extension $K$ de $k$, le morphisme projection $X_{(K)}\to X$ est
universellement ouvert.
\end{corollary}

Il suffit d'appliquer (2.4.9) au morphisme structural
$Y\to\operatorname{Spec}(k)$.

\begin{remark}[2.4.11]
\label{IV.2.4.11-fr}
Si $f:X\to Y$ est un morphisme \emph{ouvert}, on sait que, pour toute partie
$E$ de $Y$, on a
$f^{-1}(\overline E)=\overline{f^{-1}(E)}$
(Bourbaki, \emph{Top. gén.}, chap.~I, 3\textsuperscript{e} éd., \S~5,
n\textsuperscript{o}~4, prop.~7). Cette remarque s'applique par exemple
lorsque $f$ est un morphisme \emph{plat localement de présentation finie}
(2.4.6), ou un morphisme projection $X\times_k Y\to X$ où $X$, $Y$ sont des
préschémas sur un corps $k$ (2.4.10), et généralise alors (2.3.10).
\end{remark}

\subsection{Permanence des propriétés des Modules par descente fidèlement plate}

\begin{proposition}[2.5.1]
\label{IV.2.5.1-fr}
Soient $f:X\to Y$ un morphisme, $\mathscr F$ un $\mathcal O_X$-Module
quasi-cohérent, $g:Y'\to Y$ un morphisme fidèlement plat,
$X'=X\times_Y Y'$, $f'=f_{(Y')}:X'\to Y'$,
$\mathscr F'=\mathscr F\otimes_{\mathcal O_Y}\mathcal O_{Y'}$. Pour que
$\mathscr F$ soit plat (resp. fidèlement plat) relativement à $f$, il faut et
il suffit que $\mathscr F'$ soit plat (resp. fidèlement plat) relativement à
$f'$.
\end{proposition}

Il suffit d'appliquer (2.2.10) en remplaçant $X$, $Y'$, $Z$, $\mathscr F$,
$\mathscr G'$ par $Y'$, $X$, $Y$, $\mathcal O_{Y'}$, $\mathscr F$
respectivement~; on en conclut que pour que $\mathscr F$ soit plat (resp.
fidèlement plat) relativement à $f$, il faut et il suffit que $\mathscr F'$
soit plat (resp. fidèlement plat) relativement à $g\circ f'$. Mais si
$g':X'\to X$ est la projection canonique, $g'$ est fidèlement plat (2.2.13)
et on a $g\circ f'=f\circ g'$~; donc pour que $\mathscr F'$ soit plat (resp.
fidèlement plat) relativement à $f\circ g'$, il faut et il suffit que
$\mathscr F$ le soit relativement à $f$ (2.2.11, (iii)).

\begin{proposition}[2.5.2]
\EGAReaderTarget{ega4.c04.prop.2-5-2}{IV.2.5.2-fr}
Soient $f:X'\to X$ un morphisme fidèlement plat et quasi-compact,
$\mathscr F$ un $\mathcal O_X$-Module quasi-cohérent,
$\mathscr F'=f^*(\mathscr F)$. Considérons, pour un Module quasi-cohérent, la
propriété d'être~:
\begin{enumerate}
\item[(i)] de type fini~;
\item[(ii)] de présentation finie~;
\item[(iii)] localement libre de type fini~;
\item[(iv)] localement libre de rang $n$.
\end{enumerate}
Alors, si $\mathbf P$ désigne une des propriétés précédentes, pour que
$\mathscr F$ possède la propriété $\mathbf P$, il faut et il suffit que
$\mathscr F'$ la possède.
\end{proposition}

Pour qu'un $\mathcal O_X$-Module quasi-cohérent soit localement libre de type
fini, il faut et il suffit qu'il soit plat sur $X$ et de présentation finie
(Bourbaki, \emph{Alg. comm.}, chap.~II, \S~5, n\textsuperscript{o}~2,
cor.~2 du th.~1, compte tenu de (2.1.2))~; comme $\mathscr F$ est plat sur
$X$ si et seulement si $\mathscr F'$ est plat sur $X'$ en vertu de (2.5.1)
(appliqué en prenant pour $f$ l'identité), on voit que pour prouver la
proposition dans le cas (iii), il suffit de l'avoir fait dans les cas (i) et
(ii)~; il en est de même pour prouver (iv), car
$f^*(\mathcal O_X^n)=\mathcal O_{X'}^n$, donc si $\mathscr F$ et
$\mathscr F'$ sont localement libres de type fini et $x=f(x')$, le rang de
$\mathscr F$ en $x$ est égal à celui de $\mathscr F'$ en $x'$, et notre
assertion résulte de ce que $f$ est surjectif. Pour traiter les cas (i) et
(ii),
\oldpage[IV-2]{23}
notons que la question est locale sur $X$, et l'on peut donc supposer $X$
affine~; on sait alors (2.2.12) qu'il y a un schéma affine $X''$ et un
morphisme fidèlement plat $g:X''\to X'$ qui soit un isomorphisme local~; par
suite, il revient au même de dire que $f^*(\mathscr F)$ possède la propriété
$\mathbf P$, ou que $g^*(f^*(\mathscr F))$ la possède. On est ainsi ramené au
cas où $X=\operatorname{Spec}(A)$, $X'=\operatorname{Spec}(A')$~; compte tenu
de (2.2.3) et (II, 6.1.4.1), il suffit donc de prouver le

\begin{lemma}[2.5.3]
\label{IV.2.5.3-fr}
Soient $A$ un anneau, $A'$ une $A$-algèbre fidèlement plate, $M$ un
$A$-module, $M'=M\otimes_A A'$. Pour que $M$ soit de type fini (resp. de
présentation finie), il faut et il suffit que $M'$ le soit.
\end{lemma}

Pour la démonstration, voir Bourbaki, \emph{Alg. comm.},
chap.~I\textsuperscript{er}, \S~3, n\textsuperscript{o}~6, prop.~11.

\begin{remark}[2.5.4]
\label{IV.2.5.4-fr}
Les assertions de (2.5.2) pour les propriétés (i) et (ii) sont encore
valables si l'on suppose seulement $f$ quasi-fidèlement plat (2.3.3) et
quasi-compact. On est en effet ramené (Bourbaki, \emph{loc. cit.}) à prouver
le
\end{remark}

\begin{lemma}[2.5.4.1]
\label{IV.2.5.4.1-fr}
Soit $\rho:A\to A'$ un homomorphisme d'anneaux, tel que le morphisme
correspondant $f:\operatorname{Spec}(A')\to\operatorname{Spec}(A)$ soit
surjectif~; supposons qu'il existe un $A'$-module $N'$ de type fini, qui soit
$A$-plat et ait pour support $\operatorname{Spec}(A')$. Alors, si
$u:P\to Q$ est un homomorphisme de $A$-modules tel que
\[
u\otimes 1:P\otimes_A A'\longrightarrow Q\otimes_A A'
\]
soit surjectif, $u$ est surjectif.
\end{lemma}

En effet, on en déduit d'abord que l'homomorphisme
\[
u\otimes 1_{N'}:(P\otimes_A A')\otimes_{A'}N'
\longrightarrow(Q\otimes_A A')\otimes_{A'}N'
\]
est surjectif. Soit $\mathfrak q$ un idéal premier de $A'$, et soit
$\mathfrak p=\rho^{-1}(\mathfrak q)$~; l'homomorphisme
$(P\otimes_A N')_{\mathfrak q}\to(Q\otimes_A N')_{\mathfrak q}$
correspondant est surjectif, et il s'écrit
\[
u_{\mathfrak p}\otimes 1:
P_{\mathfrak p}\otimes_{A_{\mathfrak p}}N'_{\mathfrak q}
\longrightarrow
Q_{\mathfrak p}\otimes_{A_{\mathfrak p}}N'_{\mathfrak q}
\quad(0_{\mathrm I},1.5.4).
\]
Par hypothèse on a $N'_{\mathfrak q}\ne0$ et $N'_{\mathfrak q}$ est un
$A_{\mathfrak p}$-module plat $(0_{\mathrm I},6.3.1)$~; en vertu du lemme de
Nakayama, on a
$\mathfrak qN'_{\mathfrak q}\ne N'_{\mathfrak q}$, et \emph{a fortiori}
$\mathfrak pN'_{\mathfrak q}\ne N'_{\mathfrak q}$, donc
$N'_{\mathfrak q}$ est un $A_{\mathfrak p}$-module fidèlement plat
$(0_{\mathrm I},6.4.1)$. Il en résulte que $u_{\mathfrak p}$ est surjectif
$(0_{\mathrm I},6.4.1)$, et ceci ayant lieu pour tout
$\mathfrak p\in\operatorname{Spec}(A)$, puisque $f$ est surjectif, on en
conclut finalement que $u$ est surjectif (Bourbaki, \emph{Alg. comm.},
chap.~II, \S~3, n\textsuperscript{o}~3, th.~1).

\begin{proposition}[2.5.5]
\label{IV.2.5.5-fr}
Soient $f:X\to Y$ un morphisme, $\mathscr F$ un $\mathcal O_X$-Module
quasi-cohérent de type fini et $f$-plat, $\mathscr E$ un
$\mathcal O_Y$-Module quasi-cohérent de type fini~; pour tout $y\in Y$, on
pose $\mathscr F_y=\mathscr F\otimes_{\mathcal O_Y}k(y)$. Pour qu'un point
$x\in X$ soit point maximal de
$\operatorname{Supp}(\mathscr E\otimes_Y\mathscr F)$, il faut et il suffit que
$y=f(x)$ soit point maximal de $\operatorname{Supp}(\mathscr E)$, et que $x$
soit point maximal de $\operatorname{Supp}(\mathscr F_y)$ dans $f^{-1}(y)$.
S'il en est ainsi, on a
\end{proposition}

\begin{env}[2.5.5.1]
\label{IV.2.5.5.1-fr}
\[
\operatorname{long}\bigl((\mathscr E\otimes_Y\mathscr F)_x\bigr)
=\operatorname{long}(\mathscr E_y)\,
 \operatorname{long}\bigl((\mathscr F_y)_x\bigr).
\]
\end{env}

Il est clair que
$f(\operatorname{Supp}(\mathscr E\otimes_Y\mathscr F))
\subset\operatorname{Supp}(\mathscr E)$ $(0_{\mathrm I},5.2.2)$~; l'image
par $f$ de toute composante irréductible de
$\operatorname{Supp}(\mathscr E\otimes_Y\mathscr F)$ est donc contenue dans
une composante irréductible de $\operatorname{Supp}(\mathscr E)$. On peut se
borner au cas où $\operatorname{Supp}(\mathscr F)=X$. En effet
(\textbf{Err}\textsubscript{\textbf{III}}, 30), puisque $\mathscr F$ est de
type fini, il existe un sous-préschéma fermé $X'$ de $X$ ayant
$\operatorname{Supp}(\mathscr F)$ pour espace sous-jacent et un
$\mathcal O_{X'}$-Module quasi-cohérent de type fini $\mathscr G$ tels que,
si $j:X'\to X$ est l'injection canonique, on ait
$\mathscr F=j_*(\mathscr G)$. Si l'on pose $f'=f\circ j$, il est clair que
$\mathscr G$ est $f'$-plat et que l'on a
\[
\operatorname{Supp}(\mathscr E\otimes_Y\mathscr F)
=\operatorname{Supp}(\mathscr E\otimes_Y\mathscr G).
\]
Supposons donc $\operatorname{Supp}(\mathscr F)=X$~; si $Z$ est une
composante irréductible de $\operatorname{Supp}(\mathscr E)$, on a
\[
f^{-1}(Z)\subset\operatorname{Supp}(\mathscr E\otimes_Y\mathscr F)
\quad\text{(I, 9.1.13),}
\]
et il résulte de (2.3.4) que toute composante irréductible de $f^{-1}(Z)$
domine $Z$~; autrement dit, si $x$ est point générique d'une composante
irréductible de $\operatorname{Supp}(\mathscr E\otimes_Y\mathscr F)$
contenue dans $f^{-1}(Z)$, $y=f(x)$ est le point générique de $Z$~; en outre
$(0_{\mathrm I},2.1.8)$, $x$ est point générique d'une composante
irréductible de
$f^{-1}(y)=\operatorname{Supp}(\mathscr F_y)$ (I, 9.1.13), et
réciproquement tout point générique d'une de ces composantes est point
générique d'une composante irréductible de $f^{-1}(Z)$.

\oldpage[IV-2]{24}
Reste à prouver (2.5.5.1)~; on a
\[
(\mathscr E\otimes_Y\mathscr F)_x
=\mathscr E_y\otimes_{\mathcal O_{Y,y}}\mathscr F_x
\quad\text{(I, 9.1.12)}
\]
et
$(\mathscr F_y)_x=\mathscr F_x/\mathfrak m_y\mathscr F_x$~; on est donc
ramené à prouver le

\begin{lemma}[2.5.5.2]
\label{IV.2.5.5.2-fr}
Soient $A$, $B$ deux anneaux locaux, $\rho:A\to B$ un homomorphisme local,
$\mathfrak m$ l'idéal maximal de $A$. Soient $M$ un $A$-module, $N$ un
$B$-module qui est un $A$-module fidèlement plat et tel que
$N/\mathfrak mN$ soit un $B$-module de longueur finie~; alors on a
\end{lemma}

\begin{env}[2.5.5.3]
\label{IV.2.5.5.3-fr}
\[
\operatorname{long}_B(M\otimes_A N)
=\operatorname{long}_A(M)\,
 \operatorname{long}_B(N/\mathfrak mN).
\]
\end{env}

Si $M$ est de longueur infinie, il en est de même de $M\otimes_A N$, car
pour toute suite strictement croissante de $n$ sous-modules $M_i$ de $M$
($1\leq i\leq n$), les $M_i\otimes_A N$ s'identifient à des sous-$B$-modules
de $M\otimes_A N$ deux à deux distincts~; comme
$N\ne\mathfrak mN$ puisque $N$ est un $A$-module fidèlement plat, la formule
(2.5.5.3) est vraie dans ce cas. Supposons donc $M$ de longueur finie. Si
$M=0$ les deux membres de (2.5.5.3) sont nuls, donc on peut supposer
$M\ne0$. Les $\mathfrak m^kM$ sont des sous-modules de $M$, donc de longueur
finie, et si $\mathfrak m^kM\ne0$, il résulte du lemme de Nakayama que
$\mathfrak m^{k+1}M\ne\mathfrak m^kM$~; donc il y a nécessairement un entier
$r$ tel que $\mathfrak m^rM=0$. Les
$\mathfrak m^kM\otimes_A N$ s'identifient alors à des sous-$B$-modules de
$M\otimes_A N$,
\[
\frac{\mathfrak m^kM\otimes_A N}{\mathfrak m^{k+1}M\otimes_A N}
\]
étant isomorphe à
\[
(\mathfrak m^kM/\mathfrak m^{k+1}M)\otimes_A N,
\]
donc aussi à
\[
(\mathfrak m^kM/\mathfrak m^{k+1}M)
\otimes_{A/\mathfrak m}(N/\mathfrak mN)
\quad(0\leq k\leq r-1)~;
\]
la longueur de ce dernier, considéré comme $B$-module, est donc le produit de
$\operatorname{long}_B(N/\mathfrak mN)$ et du rang du
$(A/\mathfrak m)$-espace vectoriel
$\mathfrak m^kM/\mathfrak m^{k+1}M$, rang qui est aussi la longueur du
$A$-module $\mathfrak m^kM/\mathfrak m^{k+1}M$. Par addition (pour
$0\leq k\leq r-1$), on en déduit aussitôt la formule (2.5.5.3).

\begin{remark}[2.5.5.4]
\label{IV.2.5.5.4-fr}
On notera que lorsque $N$ est un $A$-module \emph{de type fini} dire que $N$
est un $A$-module fidèlement plat équivaut à dire que $N\ne0$ et que $N$ est
un $A$-module plat~; en effet le lemme de Nakayama montre alors que
$\mathfrak mN\ne N$.
\end{remark}

\begin{lemma}[2.5.6]
\label{IV.2.5.6-fr}
Soient $B$ un anneau non nécessairement commutatif, $V$, $W$ deux
$B$-modules à gauche isomorphes, $C$ l'anneau
$\operatorname{End}_B(V)$, et soit
$M=\operatorname{Hom}_B(V,W)$, muni de sa structure canonique de
$C$-module à droite. Alors $M$ est un $C$-module isomorphe à $C_d$~; en
outre, pour tout $u\in M$, les conditions suivantes sont équivalentes~:
\begin{enumerate}
\item[\emph{a)}] $\{u\}$ est une base du $C$-module $M$.
\item[\emph{b)}] $u$ est un isomorphisme de $V$ sur $W$.
\end{enumerate}
\end{lemma}

Si $u$ est un isomorphisme de $V$ sur $W$, l'application
$v\mapsto u\circ v$ de $C$ dans $M$ est évidemment une bijection, donc
\emph{b)} entraîne \emph{a)}. Inversement, supposons que $\{u\}$ soit une
base du $C$-module $M$. Par hypothèse, il existe un isomorphisme $u'$ de $V$
sur $W$, et $\{u'\}$ est par suite une base de $M$~; il y a donc un élément
inversible $w$ de $C$ (c'est-à-dire un automorphisme de $V$) tel que
$u=u'\circ w$, ce qui entraîne que $u$ est un isomorphisme de $V$ sur $W$.

\begin{corollary}[2.5.7]
\label{IV.2.5.7-fr}
Les hypothèses sur $B$, $V$ et $W$ étant celles de (2.5.6), supposons en
outre vérifiée l'une des conditions suivantes~:
\begin{enumerate}
\item[(i)] $V$ et $W$ sont des $B$-modules noethériens~;
\item[(ii)] $V$ et $W$ sont des modules de présentation finie sur un
sous-anneau commutatif de $B$.
\end{enumerate}
Alors les conditions \emph{a)} et \emph{b)} de (2.5.6) sont aussi
équivalentes aux suivantes~:
\begin{enumerate}
\item[\emph{a}$'$)] $u$ est un générateur du $C$-module $M$.
\item[\emph{b}$'$)] $u$ est un épimorphisme de $V$ sur $W$.
\end{enumerate}
\end{corollary}

On sait en effet qu'un épimorphisme d'un $A$-module $E$ sur lui-même est
bijectif dans les deux cas suivants~: $1^\circ$ $E$ est un $A$-module
noethérien (Bourbaki, \emph{Alg.}, chap.~VIII, \S~2,
n\textsuperscript{o}~2, lemme~3)~; $2^\circ$ $A$ est commutatif et $E$ est
un $A$-module de présentation finie (8.9.3)\footnote{Le lecteur pourra
vérifier que (2.5.7) et (2.5.8) ne sont pas utilisés avant le \S~9.}~; donc
\emph{b)} et \emph{b}$'$) sont équivalentes. D'autre part, si $u$ engendre
$M$
\oldpage[IV-2]{25}
et si $\{u'\}$ est une base de $M$, il existe $v\in C$ tel que
$u'=u\circ v$, ce qui prouve que $u$ est surjectif~; donc \emph{a}$'$)
entraîne \emph{b}$'$), et comme \emph{a)} entraîne évidemment
\emph{a}$'$), cela achève de prouver le corollaire.

\begin{proposition}[2.5.8]
\label{IV.2.5.8-fr}
Soient $A$ un anneau commutatif semi-local, $B$ une $A$-algèbre (non
nécessairement commutative), $V$ et $W$ deux $B$-modules. Soit $A'$ une
$A$-algèbre commutative qui est un $A$-module fidèlement plat~; on pose
\[
B'=B\otimes_A A',\qquad V'=V\otimes_A A',\qquad W'=W\otimes_A A',
\]
de sorte que $B'$ est une $A'$-algèbre, $V'$ et $W'$ des $B'$-modules. On
suppose en outre vérifiée l'une des conditions suivantes~:
\begin{enumerate}
\item[(i)] $A$ et $A'$ sont noethériens, $V$ et $W$ sont des $A$-modules de
type fini.
\item[(ii)] $B$ est un $A$-module de type fini, $V$ est un $B$-module
projectif de type fini et un $A$-module de présentation finie.
\end{enumerate}
Alors, si $V'$ et $W'$ sont des $B'$-modules isomorphes, $V$ et $W$ sont des
$B$-modules isomorphes.
\end{proposition}

Notons que dans le cas (ii), $W'$ étant $A'$-isomorphe à $V'$, est un
$A'$-module de type fini, d'où résulte que $W$ est un $A$-module de type fini
(Bourbaki, \emph{Alg. comm.}, chap.~I\textsuperscript{er}, \S~3,
n\textsuperscript{o}~6, prop.~11)~; donc dans tous les cas $V$ et $W$ sont
des $A$-modules de type fini. En outre~:

\begin{env}[2.5.8.1]
\label{IV.2.5.8.1-fr}
Sous l'une des hypothèses (i), (ii), $\operatorname{Hom}_B(V,W)$ est un
$A$-module de type fini.
\end{env}

C'est évident dans le cas (i), car $\operatorname{Hom}_A(V,W)$ est alors un
$A$-module de type fini, et $\operatorname{Hom}_B(V,W)$ est un sous-$A$-module
de $\operatorname{Hom}_A(V,W)$. Dans le cas (ii), $V$ est facteur direct d'un
$B$-module libre $B_s$, donc $\operatorname{Hom}_B(V,W)$ est facteur direct de
$\operatorname{Hom}_B(B_s,W)=W^s$, et comme $W$ est un $A$-module de type
fini, il en est de même de $\operatorname{Hom}_B(V,W)$. Posons
\[
C=\operatorname{End}_B(V),\qquad M=\operatorname{Hom}_B(V,W)
\]
qui sont des $A$-modules de type fini dans les cas (i) et (ii). On sait que
sous l'une des conditions (i), (ii), l'homomorphisme canonique

\begin{env}[2.5.8.2]
\label{IV.2.5.8.2-fr}
\[
\operatorname{Hom}_A(V,W)\otimes_A A'\longrightarrow
\operatorname{Hom}_{A'}(V',W')
\]
\end{env}

est bijectif (Bourbaki, \emph{Alg. comm.}, chap.~II, \S~2,
n\textsuperscript{o}~10, prop.~11). Comme $A'$ est un $A$-module plat,
$\operatorname{Hom}_B(V,W)\otimes_A A'$ s'identifie canoniquement à un
sous-$A'$-module de $\operatorname{Hom}_A(V,W)\otimes_A A'$. L'image de ce
sous-module par l'homomorphisme (2.5.8.2) est contenue dans
$\operatorname{Hom}_{B'}(V',W')$, car si $u\in\operatorname{Hom}_B(V,W)$ et
$a'\in A'$, l'image de $u\otimes a'$ par (2.5.8.2) est l'homomorphisme
$u':V'\to W'$ tel que $u'(x\otimes1)=u(x)\otimes a'$~; pour tout $b\in B$,
on a donc
\[
u'((b\otimes1)(x\otimes1))=u'(bx\otimes1)=u(bx)\otimes a'
=bu(x)\otimes a'=(b\otimes1)(u(x)\otimes a'),
\]
d'où notre assertion. Cela étant~:

\begin{env}[2.5.8.3]
\label{IV.2.5.8.3-fr}
Sous l'une des hypothèses (i), (ii), l'homomorphisme
\end{env}

\begin{env}[2.5.8.4]
\label{IV.2.5.8.4-fr}
\[
\operatorname{Hom}_B(V,W)\otimes_A A'\longrightarrow
\operatorname{Hom}_{B'}(V',W')
\]
est bijectif.
\end{env}

Pour tout $b\in B$, notons $h(b)$ (resp. $h'(b)$) l'homothétie $x\mapsto bx$
dans $V$ (resp. $W$), qui est un $A$-endomorphisme. Soit
$(b_\alpha)_{\alpha\in I}$ un système de générateurs de la $A$-algèbre $B$~;
l'application
\[
u\longmapsto(h'(b_\alpha)\circ u-u\circ h(b_\alpha))_{\alpha\in I}
\]
de $\operatorname{Hom}_A(V,W)$ dans $(\operatorname{Hom}_A(V,W))^I$ est
$A$-linéaire, et par définition, son noyau n'est autre que
$\operatorname{Hom}_B(V,W)$~; autrement dit, on a une suite exacte
\[
0\longrightarrow\operatorname{Hom}_B(V,W)
\longrightarrow\operatorname{Hom}_A(V,W)
\longrightarrow(\operatorname{Hom}_A(V,W))^I.
\]
Le même raisonnement s'applique en remplaçant $A$, $B$, $V$, $W$ par $A'$,
$B'$, $V'$, $W'$~; de plus, on a un diagramme

\begin{env}[2.5.8.5]
\label{IV.2.5.8.5-fr}
\[
\xymatrix{
0 \ar[r] & \operatorname{Hom}_B(V,W)\otimes_A A' \ar[r] \ar[d]^{r} &
\operatorname{Hom}_A(V,W)\otimes_A A' \ar[r] \ar[d]^{s} &
(\operatorname{Hom}_A(V,W))^I\otimes_A A' \ar[d]^{t}\\
0 \ar[r] & \operatorname{Hom}_{B'}(V',W') \ar[r] &
\operatorname{Hom}_{A'}(V',W') \ar[r] &
(\operatorname{Hom}_{A'}(V',W'))^I
}
\]
\end{env}

où $r$ est l'homomorphisme (2.5.8.4), $s$ l'homomorphisme (2.5.8.2) et
$t$ l'homomorphisme composé
\[
(\operatorname{Hom}_A(V,W))^I\otimes_A A'
\xrightarrow{w}(\operatorname{Hom}_A(V,W)\otimes_A A')^I
\xrightarrow{s^I}(\operatorname{Hom}_{A'}(V',W'))^I
\]
$w$ étant l'homomorphisme canonique (Bourbaki, \emph{Alg.}, chap.~II,
3\textsuperscript{e} éd., \S~3, n\textsuperscript{o}~7). On vérifie aussitôt
que le diagramme (2.5.8.5) est commutatif, et comme $A'$ est un $A$-module
plat, ses lignes sont exactes. Enfin, on a vu

\oldpage[IV-2]{26}
que $s$ est un isomorphisme, et il en est donc de même de $s^I$~; dans le cas
(ii), on peut prendre $I$ fini, et on sait alors que $w$ est bijectif
(Bourbaki, \emph{loc. cit.}, prop.~7)~; dans le cas (i), on note que si $B'$
(resp. $B''$) est l'image de $B$ par $h$ (resp. $h'$) dans
$\operatorname{End}_A(V)$ (resp. $\operatorname{End}_A(W)$), $B'$ et $B''$
sont des $A$-modules de type fini, donc on peut encore prendre $I$ fini.
Ainsi, dans tous les cas, $t$ est bijectif, et on en conclut donc qu'il en est
de même de $r$.

Il résulte donc de (2.5.8.4) que si l'on pose
\[
C'=C\otimes_A A',\qquad M'=M\otimes_A A'
\]
on a des bijections canoniques

\begin{env}[2.5.8.6]
\label{IV.2.5.8.6-fr}
\[
C'\xrightarrow{\sim}\operatorname{End}_{B'}(V'),\qquad
M'\xrightarrow{\sim}\operatorname{Hom}_{B'}(V',W')
\]
{\emergencystretch=2em
dont la première est un isomorphisme de $A'$-algèbres, la seconde constituant
avec la première un di-isomorphisme de $C'$-modules à droite.\par}
\end{env}

\begin{env}[2.5.8.7]
\label{IV.2.5.8.7-fr}
\emph{Réduction au cas où $V=B_s$.} L'hypothèse que $V'$ et $W'$ sont des
$B'$-modules isomorphes, entraîne que $C'_d$ et $M'$ sont des $C'$-modules à
droite isomorphes (2.5.6). Montrons que pour prouver (2.5.8), il suffit de
trouver un élément $u\in M$ qui soit un générateur de $M$ en tant que
$C$-module à droite. En effet, $u'=u\otimes1$ sera un générateur de $M'$ en
tant que $C'$-module à droite~; or, dans le cas (i), $V'$ et $W'$ sont des
$A'$-modules de type fini, donc noethériens puisque $A'$ est noethérien~;
dans le cas (ii), $V'$ et $W'$ sont des $A'$-modules (isomorphes) de
présentation finie. On peut donc appliquer (2.5.7) à $A'$, $B'$, $V'$ et
$W'$ et on en conclut que $u'$ est un $B'$-isomorphisme de $V'$ sur $W'$.
Mais comme $A'$ est fidèlement plat sur $A$, cela entraîne que $u$ est
bijectif $(0_{\mathrm I},6.4.1)$, c'est-à-dire la conclusion de (2.5.8). Si
l'on note que dans le cas (i), $C$ et $M$ sont des $A$-modules de type fini,
et que dans le cas (ii), $C$ est (comme on l'a vu dans (2.5.8.1)) facteur
direct de $V^n$, donc un $A$-module projectif de type fini, on voit qu'on est
ramené (en changeant les notations) à prouver (2.5.8) dans le cas particulier
où $V=B_s$, et qu'il suffit de prouver l'existence d'un générateur du
$B$-module $W$. On notera qu'alors $B$ est un $A$-module de type fini.
\end{env}

\begin{env}[2.5.8.8]
\label{IV.2.5.8.8-fr}
\emph{Réduction au cas où $A$ et $A'$ sont des corps, $B$ une $A$-algèbre
simple de centre $A$.} Soit $\mathfrak r$ le radical de l'anneau semi-local
$A$~; il suffit de prouver que $W/\mathfrak rW$ est un
$(B/\mathfrak rB)$-module monogène, car s'il existe un homomorphisme surjectif
$B_s/\mathfrak rB_s\to W/\mathfrak rW$, il donne par composition un
homomorphisme $B_s\to B_s/\mathfrak rB_s\to W/\mathfrak rW$, qui lui-même
(puisque $B_s$ est un $B$-module libre) peut s'écrire
$B_s\xrightarrow{f}W\to W/\mathfrak rW$, si bien que l'homomorphisme
surjectif considéré est
\[
f\otimes1:B_s\otimes_A(A/\mathfrak r)\longrightarrow
W\otimes_A(A/\mathfrak r).
\]
Comme $W$ est un $A$-module de type fini, le lemme de Nakayama prouve que
$f$ est surjectif (Bourbaki, \emph{Alg. comm.}, chap.~II, \S~3,
n\textsuperscript{o}~2, cor.~1 de la prop.~4). Si l'on pose
\[
A_1=A/\mathfrak r,\quad A'_1=A'\otimes_A A_1=A'/\mathfrak rA',\quad
B_1=B/\mathfrak rB=B\otimes_A A_1,\quad
W_1=W/\mathfrak rW=W\otimes_A A_1,
\]
les hypothèses (i) (resp. (ii)) restent vérifiées lorsque l'on y remplace
$A$, $A'$, $B$, $V=B_s$, $W$ par $A_1$, $A'_1$, $B_1$,
$V_1=(B_1)_s$, $W_1$ respectivement~; en outre,
$V'_1=V'\otimes_A A_1=V_1\otimes_{A_1}A'_1$ et
$W'_1=W'\otimes_A A_1=W_1\otimes_{A_1}A'_1$ sont $B'_1$-isomorphes (avec
$B'_1=B'\otimes_A A_1=B_1\otimes_{A_1}A'_1$) et $A'_1$ est un $A_1$-module
fidèlement plat. On peut donc supposer pour démontrer (2.5.8) que $A$ est un
\emph{produit fini de corps} (commutatifs). Comme $B$ est un $A$-module de
type fini, c'est un anneau artinien~; soit $\mathfrak R$ son radical. Il
suffira maintenant de prouver que $W/\mathfrak RW$ est un
$(B/\mathfrak R)$-module monogène, car on voit comme ci-dessus, en utilisant
le lemme de Nakayama, que cela entraîne que $W$ est un $B$-module monogène~;
d'autre part, $W'/\mathfrak RW'$ est $(B'/\mathfrak RB')$-isomorphe à
$(B'/\mathfrak RB')_s$, et l'on a
\[
B'/\mathfrak RB'=(B\otimes_A A')\otimes_B(B/\mathfrak R)
=(B/\mathfrak R)\otimes_A A'
\]
et de même
\[
W'/\mathfrak RW'=(W\otimes_A A')\otimes_B(B/\mathfrak R)
=(W/\mathfrak RW)\otimes_A A'.
\]
On peut donc supposer en outre que $\mathfrak R=0$, c'est-à-dire que $B$ est
une $A$-algèbre semi-simple.

Notons maintenant que puisque $A$ est un produit fini de corps $k_i$
$(1\leq i\leq n)$, $A'$ est composé direct de $k_i$-algèbres $A'_i$
$(1\leq i\leq n)$, chacune des $A'_i$ étant annulée par les $k_j$ d'indice
$j\ne i$~; l'hypothèse que $A'$ est un $A$-module fidèlement plat entraîne
que les $A'_i$ sont $\ne0$~; il y a par suite un quotient $A''_i$ de $A'_i$
qui est un corps, et $A''$, composé direct des $A''_i$, est un $A$-module
fidèlement plat et un quotient de $A'$. Si l'on considère alors l'anneau
$B''=B'\otimes_{A'}A''=B\otimes_A A''$,
$W''=W'\otimes_{A'}A''=W\otimes_A A''$ est un $B''$-module isomorphe à
$B''_s$~; on peut donc se borner à démontrer (2.5.8) en remplaçant $A'$ par
$A''$, c'est-à-dire que l'on peut aussi supposer que $A'$ est un \emph{produit
fini de corps}.

Soit $Z$ le centre de $B$, qui est un produit fini de corps et un $A$-module
de type fini~; notons que $B$ et $W$ sont des $Z$-modules de type fini,
projectifs comme tout $Z$-module~; en outre, on a
$B'=B\otimes_Z Z'$ et $W'=W\otimes_Z Z'$ en posant
$Z'=Z\otimes_A A'$, et $Z'$ est un $Z$-module fidèlement plat. On peut donc
remplacer $A$ par $Z$ dans l'hypothèse, autrement dit, supposer que $A$ est
le centre de $B$, $B$ étant semi-simple et $A'$ un produit fini de corps. Si
$k_i$ $(1\leq i\leq n)$ sont les corps composants de $A$, $B$ est composé
direct d'anneaux simples $B_i$, $k_i$ étant le centre de $B_i$, et $W$ est
somme directe de sous-modules $W_i$ $(1\leq i\leq n)$, $W_i$ étant annulé
par les $B_j$ d'indice $j\ne i$~; en outre, le raisonnement fait plus haut
montre que l'on peut supposer que $A'$ est un produit de corps $k'_i$
$(1\leq i\leq n)$, $k'_i$ étant une extension de $k_i$ et étant annulé par
les $k_j$ d'indice $j\ne i$. L'hypothèse que $B'_s$ et $W'$ sont des
$B'$-modules isomorphes entraîne alors que, pour
\oldpage[IV-2]{27}
tout $i$,
$(B_i\otimes_{k_i}k'_i)_s$ et $W_i\otimes_{k_i}k'_i$ sont des
$(B_i\otimes_{k_i}k'_i)$-modules isomorphes~; il suffit donc de prouver
(2.5.8) lorsque $n=1$, c'est-à-dire dans le cas où $A$ et $A'$ sont des
corps, $B$ une algèbre simple de centre $A$.
\end{env}

\begin{env}[2.5.8.9]
\label{IV.2.5.8.9-fr}
\emph{Fin de la démonstration.} On sait alors (Bourbaki, \emph{Alg.},
chap.~VIII, \S~5, prop.~6 et 8) que tout $B$-module est somme directe de
modules isomorphes à un idéal minimal de $B$, et deux $B$-modules de rang
fini sur $A$ sont donc isomorphes si et seulement s'ils ont même rang sur
$A$. Par hypothèse, on a $[W':A']=[B'_s:A']$. Mais on a
$[W':A']=[W:A]$ et $[B'_s:A']=[B_s:A]$~; donc $[W:A]=[B_s:A]$, ce qui
termine la démonstration.
\end{env}

\subsection{Permanence de propriétés ensemblistes et topologiques de morphismes
par descente fidèlement plate}
\label{subsection:IV.2.6-fr}

\begin{proposition}[2.6.1]
\EGAReaderTarget{ega4.c04.prop.2-6-1}{IV.2.6.1-fr}
Soient $f:X\to Y$ un $S$-morphisme de $S$-préschémas, $g:S'\to S$ un
morphisme surjectif, $X'=X_{(S')}$, $Y'=Y_{(S')}$,
$f'=f_{(S')}:X'\to Y'$. Considérons, pour un morphisme, la propriété d'être~:
\begin{enumerate}
\item[(i)] surjectif~;
\item[(ii)] injectif~;
\item[(iii)] à fibres finies (en tant qu'ensembles)~;
\item[(iv)] bijectif~;
\item[(v)] radiciel.
\end{enumerate}
Alors, si $P$ désigne une des propriétés précédentes et si $f'$ possède la
propriété $P$, il en est de même de $f$.
\end{proposition}

Comme le morphisme projection $Y'\to Y$ est lui aussi surjectif (I, 3.5.2),
on peut, en vertu de (I, 3.3.11), se borner au cas où $Y=S$, $Y'=S'$. Pour
tout $y\in Y$ (resp. $y'\in Y'$) désignons par $X_y$ (resp. $X'_{y'}$) le
préschéma fibre $f^{-1}(y)$ (resp. $f'^{-1}(y')$) (I, 3.6.2)~; on sait que
pour $y'\in Y'$, $y=g(y')$, on a un isomorphisme canonique
\[
X'_{y'}\xrightarrow{\sim}X_y\otimes_{k(y)}k(y')\quad\text{(I, 3.6.4)}~;
\]
comme le morphisme $\operatorname{Spec}(k(y'))\to\operatorname{Spec}(k(y))$
est surjectif, il en est de même de la projection $X'_{y'}\to X_y$
(I, 3.5.2)~; donc si $X'_{y'}$ est non vide (resp. a au plus un point, resp.
est un ensemble fini), il en est de même de $X_y$~; comme $g:Y'\to Y$ est
surjectif, cela démontre la proposition dans les cas (i), (ii) et (iii), et
(iv) résulte de (i) et (ii)~; enfin, pour prouver (v), il suffit de montrer
que si $f'$ est universellement injectif (I, 3.5.11), il en est de même de
$f$~; or, soit $Y_1\to Y$ un morphisme quelconque~; posons
$X_1=X\times_Y Y_1$ et $f_1=f_{(Y_1)}$. D'autre part, posons
\[
Y'_1=Y_1\times_Y Y',\quad X'_1=X'\times_{Y'}Y'_1=X_1\times_{Y_1}Y'_1,
\quad f'_1=f'_{(Y'_1)}=(f_1)_{(Y'_1)}:X'_1\to Y'_1~;
\]
comme $g'=g_{(Y_1)}:Y'_1\to Y_1$ est surjectif (I, 3.5.2) et que $f'$ est
universellement injectif, $f'_1$ est injectif, et il résulte de (ii) que
$f_1$ est injectif, d'où notre assertion.

\begin{proposition}[2.6.2]
\label{IV.2.6.2-fr}
Les notations étant celles de (2.6.1), supposons le morphisme $g:S'\to S$
fidèlement plat et quasi-compact. Considérons, pour un morphisme, la propriété
d'être~:
\begin{enumerate}
\item[(i)] ouvert~;
\item[(ii)] fermé~;
\item[(iii)] quasi-compact et un homéomorphisme sur le sous-espace image~;
\item[(iv)] un homéomorphisme.
\end{enumerate}
Alors, si $P$ désigne une des propriétés précédentes et si $f'$ possède la
propriété $P$, il en est de même de $f$.
\end{proposition}

\oldpage[IV-2]{28}
Comme le morphisme $Y'\to Y$ est fidèlement plat et quasi-compact (2.2.13 et
1.1.2), on peut, en vertu de (I, 3.3.11), se borner au cas où $Y=S$,
$Y'=S'$. Si $g'$ est la projection $X'\to X$, on sait que l'on a, pour toute
partie $M$ de $X$,
\[
g^{-1}(f(M))=f'(g'^{-1}(M))\quad\text{(I, 3.4.8).}
\]
Si $f'$ est un morphisme ouvert (resp. fermé) alors, pour toute partie
ouverte (resp. fermée) $M$ de $X$, $f'(g'^{-1}(M))$ est ouverte (resp.
fermée) dans $Y'$, et comme $g$ est fidèlement plat et quasi-compact, on en
conclut que $f(M)$ est ouverte (resp. fermée) dans $Y$ en vertu de (2.3.12).
Cela démontre la proposition dans les cas (i) et (ii). Démontrons-la dans le
cas (iii) (ce qui l'entraînera dans le cas (iv), compte tenu de (2.6.1,
(iv)). En vertu de (2.6.1, (ii)) $f$ est injectif, et il reste à prouver que
$f$, considéré comme application de $X$ sur $f(X)$, est une application
quasi-compacte et ouverte. Comme $f'$ est quasi-compact, il en est de même de
$f$ (1.1.4). Il suffit alors de prouver que pour toute partie fermée $Z$ de
$X$, on a
\[
Z=f^{-1}(\overline{f(Z)})~;
\]
comme $g'$ est surjective, cette relation équivaut à
\[
g'^{-1}(Z)=g'^{-1}(f^{-1}(\overline{f(Z)}))
=f'^{-1}(g^{-1}(\overline{f(Z)})).
\]
Or, comme $f$ est quasi-compact, il en est de même de sa composée avec
l'injection canonique $Z\to X$ ($Z$ étant ici le sous-préschéma fermé réduit
de $X$ ayant $Z$ comme espace sous-jacent). Appliquant (2.3.10) à la partie
$f(Z)$ de $Y$ (image du morphisme $f|Z$), on obtient
\[
g^{-1}(\overline{f(Z)})=\overline{g^{-1}(f(Z))}~;
\]
la formule à démontrer équivaut donc à la suivante~:
\[
Z'=f'^{-1}(\overline{f'(Z')}),
\]
où l'on a posé $Z'=g'^{-1}(Z)$~; mais cette formule résulte de l'hypothèse
que $f'$ est un homéomorphisme de $X'$ sur $f'(X')$.

\begin{remark}[2.6.3]
\label{IV.2.6.3-fr}
Dans les cas (i) et (ii), les conclusions de (2.6.2) sont encore valables
quand on suppose seulement $g$ quasi-fidèlement plat (2.3.3) et
quasi-compact~; en effet, en vertu de (2.1.4), (I, 3.5.2) et
(I, 9.1.13.1), on peut encore se ramener au cas où $Y=S$, $Y'=S'$~; la
conclusion résulte alors de (2.3.12). Dans les cas (iii) et (iv), les
conclusions restent valables quand on suppose seulement $g$ quasi-fidèlement
plat, pourvu que l'on suppose déjà $f$ quasi-compact~; en effet, on n'utilise
alors que (2.3.10) et le fait que $g$ est surjectif. Enfin, si $g$ est
fidèlement plat et localement de présentation finie ou si $g$ est surjectif
et $S$ discret, la conclusion de (2.6.2) est valable lorsque $P$ est la
propriété~:
\begin{enumerate}
\item[(iii \emph{bis})] être un homéomorphisme sur le sous-espace image~;
\end{enumerate}
cela résulte en effet de la démonstration donnée dans (2.6.2) et de la
Remarque (2.4.11).
\end{remark}

\begin{corollary}[2.6.4]
\label{IV.2.6.4-fr}
Les notations étant celles de (2.6.1), supposons le morphisme $g:S'\to S$
fidèlement plat et quasi-compact. Considérons pour un morphisme, la propriété
d'être~:
\begin{enumerate}
\item[(i)] universellement ouvert~;
\item[(ii)] universellement fermé~;
\item[(iii)] quasi-compact et universellement bicontinu~;
\item[(iv)] un homéomorphisme universel~;
\item[(v)] quasi-compact~;
\item[(vi)] quasi-compact et dominant.
\end{enumerate}
Alors, si $P$ désigne une des propriétés précédentes, pour que $f$ possède la
propriété $P$, il faut et il suffit que $f'$ la possède.
\end{corollary}

\oldpage[IV-2]{29}
Les propriétés (v) et (vi) ne sont mises que pour mémoire, étant conséquences
de (1.1.4), (1.1.6) et (2.3.7). En ce qui concerne les autres, la condition
est nécessaire en vertu de (2.4.3). Inversement, supposons par exemple que
$f'$ soit universellement ouvert, et soit $Y_1\to Y$ un morphisme
quelconque~; posons $X_1=X\times_Y Y_1$ et $f_1=f_{(Y_1)}$. D'autre part,
posons
\[
Y'_1=Y_1\times_Y Y',\qquad
X'_1=X'\times_{Y'}Y'_1=X_1\times_{Y_1}Y'_1,
\qquad f'_1=f'_{(Y'_1)}=f_{1(Y'_1)}:X'_1\to Y'_1~;
\]
comme $g'=g_{(Y')}:Y'_1\to Y'$ est fidèlement plat et quasi-compact
(2.2.13 et 1.1.2), et que $f'$ est universellement ouvert, $f'_1$ est ouvert
et il résulte de (2.6.2) que $f_1$ est ouvert~; donc $f$ est universellement
ouvert. Même raisonnement pour les autres cas.

On notera encore ici que l'on peut remplacer «~fidèlement plat~» par
«~quasi-fidèlement plat~» et, lorsque $g$ est en outre localement de
présentation finie, ou que $g$ est surjectif et $S$ discret, on peut remplacer
la propriété (iii) par~:
\begin{enumerate}
\item[(iii \emph{bis})] universellement bicontinu.
\end{enumerate}

\subsection{Permanence de diverses propriétés des morphismes par descente
fidèlement plate}
\label{subsection:IV.2.7-fr}

\begin{proposition}[2.7.1]
\label{IV.2.7.1-fr}
Soient $f:X\to Y$ un $S$-morphisme de $S$-préschémas, $g:S'\to S$ un
morphisme fidèlement plat et quasi-compact, $X'=X_{(S')}$, $Y'=Y_{(S')}$,
$f'=f_{(S')}:X'\to Y'$. Considérons, pour un morphisme, la propriété d'être~:
\begin{enumerate}
\item[(i)] séparé~;
\item[(ii)] quasi-séparé~;
\item[(iii)] localement de type fini~;
\item[(iv)] localement de présentation finie~;
\item[(v)] de type fini~;
\item[(vi)] de présentation finie~;
\item[(vii)] propre~;
\item[(viii)] un isomorphisme~;
\item[(ix)] un monomorphisme~;
\item[(x)] une immersion ouverte~;
\item[(xi)] une immersion quasi-compacte~;
\item[(xii)] une immersion fermée~;
\item[(xiii)] affine~;
\item[(xiv)] quasi-affine~;
\item[(xv)] fini~;
\item[(xvi)] quasi-fini~;
\item[(xvii)] entier.
\end{enumerate}
Alors, si $P$ désigne une des propriétés précédentes, pour que $f$ possède la
propriété $P$, il faut et il suffit que $f'$ la possède.
\end{proposition}

Il a été démontré aux chapitres I\textsuperscript{er}, II et dans le chapitre
IV, \S~1, que si $f$ possède une des propriétés $P$ ci-dessus, il en est de
même de $f'$ (sans hypothèse sur le morphisme $g:S'\to S$). Il reste donc à
démontrer la réciproque~; comme la projection $Y'\to Y$ est

\oldpage[IV-2]{30}
un morphisme fidèlement plat et quasi-compact (2.2.13 et 1.1.2), on peut se
borner au cas où $S=Y$, $S'=Y'$ en vertu de (I, 3.3.11).

\emph{(i)} Dire que $f$ est séparé signifie que le morphisme diagonal
$\Delta_f:X\to X\times_Y X$ est fermé~; comme
$\Delta_{f'}=(\Delta_f)_{(Y')}$ (I, 5.3.4), si $\Delta_{f'}$ est fermé, il
en est de même de $\Delta_f$ en vertu de (2.6.2), donc $f$ est séparé.

\emph{(ii)} a déjà été démontré sous des hypothèses plus faibles (1.2.5).

\emph{(iii)} et \emph{(iv)}~: La question est évidemment locale sur $X$ et
$Y$, et, compte tenu de (2.2.12), il suffit donc de démontrer le

\begin{lemma}[2.7.1.1]
\label{IV.2.7.1.1-fr}
Soient $A$ un anneau, $B$ une $A$-algèbre, $A'$ une $A$-algèbre qui soit un
$A$-module fidèlement plat, $B'=B\otimes_A A'$. Pour que $B$ soit une
$A$-algèbre de type fini (resp. de présentation finie), il faut et il suffit
que $B'$ soit une $A'$-algèbre de type fini (resp. de présentation finie).
\end{lemma}

On sait déjà que la condition est nécessaire sans hypothèse sur $A'$
(1.3.4, 1.3.6, 1.4.3 et 1.4.6). Supposons que $B'$ soit une $A'$-algèbre de
type fini~; soit $(B_\alpha)_{\alpha\in I}$ la famille filtrante croissante
des sous-$A$-algèbres de $B$, de sorte que l'on a
$B=\varinjlim B_\alpha$, et par suite aussi
$B'=\varinjlim(B_\alpha\otimes_A A')$, le produit tensoriel permutant aux
limites inductives~; si $(x'_i)$ est un système fini de générateurs de la
$A'$-algèbre $B'$, il existe un indice $\alpha$ tel que tous les $x'_i$
appartiennent à la sous-algèbre $B_\alpha\otimes_A A'$ de $B'$, d'où
$B'=B_\alpha\otimes_A A'$, et comme $A'$ est fidèlement plat,
$B=B_\alpha$ $(0_{\mathrm I},6.4.1)$.

Supposons maintenant que $B'$ soit une $A'$-algèbre de présentation finie~;
on sait déjà d'après ce qui précède que $B$ est une $A$-algèbre de type fini,
donc il existe une $A$-algèbre de polynômes
$C=A[T_1,\ldots,T_m]$ et un $A$-homomorphisme surjectif d'algèbres
$C\to B$~; soit $\mathfrak J$ le noyau de cet homomorphisme, de sorte que
l'on a une suite exacte
\[
0\to\mathfrak J\to C\to B\to0,
\]
et par suite aussi une suite exacte $0\to\mathfrak J'\to C'\to B'\to0$
(puisque $A'$ est $A$-plat), en posant
$C'=C\otimes_A A'=A'[T_1,\ldots,T_m]$, et
$\mathfrak J'=\mathfrak J\otimes_A A'$ (identifié à un idéal de $C'$).
Comme $B'$ est une $A'$-algèbre de présentation finie, $\mathfrak J'$ est un
$C'$-module de type fini (1.4.4)~; mais on a
$\mathfrak J'=\mathfrak J\otimes_C C'$ et $C'$ est un $C$-module fidèlement
plat (2.2.13 et 2.2.3)~; on sait alors que l'hypothèse que $\mathfrak J'$ soit
un $C'$-module de type fini entraîne que $\mathfrak J$ est un $C$-module de
type fini (Bourbaki, \emph{Alg. comm.}, chap.~I\textsuperscript{er}, \S~3,
n\textsuperscript{o}~6, prop.~11), donc $B$ est une $A$-algèbre de
présentation finie.

\emph{(v)} résulte de (iii) et de (2.6.2, (v)) en vertu de (1.5.2).

\emph{(vi)} résulte de même de (iv), (v) et (ii) en vertu de (1.6.1).

\emph{(vii)} résulte de (i), (v) et (2.6.4, (ii)) (II, 5.4.1).

\emph{(viii)} Notons d'abord que comme $f'$ est un isomorphisme, c'est un
homéomorphisme universel, donc il en est de même de $f$ (2.6.4)~; on en
conclut déjà que $f$ est quasi-compact et séparé (2.4.4). Posons
$f=(\psi,\theta)$, où $\psi$ est donc un homéomorphisme~; il s'agit de
prouver que
\[
\theta:\mathcal O_Y\longrightarrow f_*(\mathcal O_X)
\]
est un isomorphisme de $\mathcal O_Y$-Modules. Or, si l'on pose
$f'=(\psi',\theta')$, l'homomorphisme
$\theta':\mathcal O_{Y'}\to f'_*(\mathcal O_{X'})$ est composé de
l'homomorphisme canonique
$g^*(f_*(\mathcal O_X))\to f'_*(\mathcal O_{X'})$ et de $g^*(\theta)$
(2.3.2)~; mais le premier de ces deux homomorphismes est bijectif en vertu de
l'hypothèse sur $g$ (2.3.1), donc si $\theta'$ est bijectif,
\oldpage[IV-2]{31}
il en est de même de $g^*(\theta)$, et comme $g$ est fidèlement plat,
$\theta$ est bijectif (2.2.7), ce qui prouve (viii).

\emph{(ix)} La proposition résulte de (viii), de (I, 5.3.4) et de (I, 5.3.8)
qui ramène les monomorphismes aux isomorphismes.

\emph{(x)} Si $f'$ est une immersion ouverte, $f'(X')$ est ouvert dans $Y'$,
et l'on a $f'(X')=g^{-1}(f(X))$ (I, 3.4.8)~; il résulte de (2.3.12) que
$f(X)$ est ouvert. On peut alors remplacer $Y$ (resp. $Y'$) par le
sous-préschéma induit sur l'ouvert $f(X)$ (resp. $f'(X')$), compte tenu de
(1.1.2) et (2.2.13)~; alors $f'$ devient un isomorphisme, donc il en est de
même de $f$ par (viii), et cela établit (x).

\emph{(xi)} Si $f'$ est une immersion quasi-compacte, $f'$ est un morphisme
quasi-compact et quasi-séparé (1.2.2), donc il en est de même de $f$ par (ii)
et (2.6.2, (v)). Soit $Z$ le sous-préschéma de $Y$ image fermée de $X$ par
$f$ (1.7.8), et posons $f=j\circ q$, où $j:Z\to Y$ est l'injection
canonique~; on a alors $f'=j'\circ q'$ avec $j'=j_{(Y')}$,
$q'=q_{(Y')}$, et l'on sait que $j'$ s'identifie à l'injection canonique
$Z'\to Y'$ du sous-préschéma $Z'$ de $Y'$, image fermée de $X'$ par $f'$
(2.3.2). L'hypothèse sur $f'$ signifie alors que $q'$ est une immersion
ouverte (I, 9.5.10), donc il en est de même de $q$ par (x), et cela montre
que $f$ est une immersion.

\emph{(xii)} Dire que $f$ (resp. $f'$) est une immersion fermée signifie que
$f$ (resp. $f'$) est une immersion quasi-compacte et un morphisme fermé~; on
voit donc que (xii) résulte de (xi) et de (2.6.2, (ii)).

\emph{(xiii)} et \emph{(xiv)} Supposons $f'$ affine (resp. quasi-affine)~;
notons alors que $f'$ est quasi-compact et quasi-séparé (II, 5.1.1), donc il
en est de même de $f$ par (ii) et (2.6.2, (v)). Posons
$\mathcal A=f_*(\mathcal O_X)$, $\mathcal A'=f'_*(\mathcal O_{X'})$~; en
vertu de (2.3.1), l'homomorphisme canonique de $\mathcal O_{Y'}$-Algèbres
$g^*(\mathcal A)\to\mathcal A'$ est bijectif~; par conséquent, si
$h:Z=\operatorname{Spec}(\mathcal A)\to Y$ est le morphisme structural, le
morphisme structural $h':Z'=\operatorname{Spec}(\mathcal A')\to Y'$
s'identifie à $h_{(Y')}$ (II, 1.5.2). Soit alors $u:X\to Z$ (resp.
$u':X'\to Z'$) le $Y$-morphisme (resp. $Y'$-morphisme) canonique
correspondant à l'homomorphisme identique de $\mathcal A$ (resp.
$\mathcal A'$) (II, 1.2.7)~; comme on a le diagramme commutatif
\[
\xymatrix{
X \ar[d]_{u} & X' \ar[l] \ar[d]^{u'}\\
Z \ar[d]_{h} & Z' \ar[l]_{g'} \ar[d]^{h'}\\
Y & Y' \ar[l]^{g}
}
\]
et que $h'\circ u'=f'$, il résulte de (II, 1.2.7) que $u'=u_{(Z')}$. En
outre, $g'$ est fidèlement plat et quasi-compact (1.1.2 et 2.2.13). Cela
étant, l'hypothèse sur $f'$ signifie que $u'$ est un isomorphisme (resp. une
immersion ouverte) (II, 5.1.6)~; il résulte alors de (viii) (resp. (x)) que
$u$ est un isomorphisme (resp. une immersion ouverte), d'où (xiii) (resp.
(xiv)).

\oldpage[IV-2]{32}
\emph{(xv)} Si $f'$ est fini, il est affine, donc aussi $f$ d'après (xiii)~;
en outre, avec les notations de la démonstration de (xiii), $\mathcal A'$ est
un $\mathcal O_{Y'}$-Module de type fini, et $\mathcal A'$ est isomorphe à
$g^*(\mathcal A)$~; il résulte de (2.5.2) que $\mathcal A$ est un
$\mathcal O_Y$-Module de type fini, donc $f$ est un morphisme fini.

\emph{(xvi)} Dire que $f$ est quasi-fini signifie que $f$ est un morphisme de
type fini et que pour tout $y\in Y$, $f^{-1}(y)$ est fini (II, 6.2.2 et I,
6.4.4)~; la conclusion résulte donc de (v) et de (xv).

\emph{(xvii)} On voit comme dans (xv) que $f$ est affine. On peut se borner
au cas où $Y=\operatorname{Spec}(A)$, $Y'=\operatorname{Spec}(A')$, et alors
$X=\operatorname{Spec}(B)$, $X'=\operatorname{Spec}(B')$, où
$B'=B\otimes_A A'$~; $B$ est égale à la limite inductive de ses
sous-$A$-algèbres de type fini $B_\alpha$, donc on a
$B'=\varinjlim B'_\alpha$, où $B'_\alpha=B_\alpha\otimes_A A'$, et
$B'_\alpha$ est une $A'$-algèbre de type fini. Mais par hypothèse $B'$ est
entière sur $A'$, donc $B'_\alpha$ est un $A'$-module de type fini, et
$B_\alpha$ est donc un $A$-module de type fini (2.5.2). C.Q.F.D.

\begin{corollary}[2.7.2]
\label{IV.2.7.2-fr}
Les hypothèses et notations étant celles de (2.7.1), supposons $f$
quasi-compact~; soient $\mathcal L$ un $\mathcal O_X$-Module inversible,
$\mathcal L'=\mathcal L\otimes_{\mathcal O_X}\mathcal O_{X'}$, son image
réciproque. Pour que $\mathcal L$ soit ample (resp. très ample) pour $f$, il
faut et il suffit que $\mathcal L'$ soit ample (resp. très ample) pour $f'$.
\end{corollary}

La condition est nécessaire sans hypothèse sur $g:S'\to S$ (II, 4.4.10 et
4.6.13)~; pour voir qu'elle est suffisante, on peut comme dans (2.7.1) se
limiter au cas où $S=Y$, $S'=Y'$. L'hypothèse sur $\mathcal L'$ implique que
$f'$ est quasi-compact et séparé (II, 4.6.1), donc il en est de même de $f$
((2.6.2, (v)) et (2.7.1, (i))). Posons
$\mathcal E=f_*(\mathcal L)$, $\mathcal E'=f'_*(\mathcal L')$~; il résulte
de (2.3.1) que l'homomorphisme canonique
$u:g^*(\mathcal E)\to\mathcal E'$ est bijectif. Si $\mathcal L'$ est très
ample pour $f'$, l'homomorphisme canonique
$\sigma':f'^*(\mathcal E')\to\mathcal L'$ est surjectif, et le morphisme
\[
r'=r_{\mathcal L',\sigma'}:X'\longrightarrow\mathbf P(\mathcal E')
\]
est une immersion (II, 4.4.4, b)) nécessairement quasi-compacte (1.1.2,
(v)). Le fait que $u:g^*(\mathcal E)\to\mathcal E'$ soit bijectif entraîne
que si $h:\mathbf P(\mathcal E)\to Y$, $h':\mathbf P(\mathcal E')\to Y'$
sont les morphismes structuraux, $h'$ s'identifie à $h_{(Y')}$ (II, 4.1.3)~;
d'autre part, en désignant par $g'$ la projection $X'\to X$, $g'$ est
fidèlement plat (2.2.13), on a $f\circ g'=g\circ f'$, et l'on vérifie
aisément que l'homomorphisme
\[
g'^*(f^*(\mathcal E))\longrightarrow g'^*(\mathcal L)
\]
est identique à l'homomorphisme composé
\[
f'^*(g^*(f_*(\mathcal L)))\xrightarrow{f'^*(u)}
f'^*(f'_*(\mathcal L'))\xrightarrow{\sigma'}\mathcal L'
\]
(par exemple en se ramenant au cas où $Y$ et $Y'$ sont affines). Comme
$\sigma'$ est surjectif et $f'^*(u)$ bijectif, on voit que
$g'^*(\sigma)$ est surjectif, donc il en est de même de $\sigma$ (2.2.7).
On en conclut que le morphisme
$r=r_{\mathcal L,\sigma}:X\to\mathbf P(\mathcal E)$ est partout défini
(II, 3.7.4)~; en outre, si l'on pose $P=\mathbf P(\mathcal E)$,
$P'=\mathbf P(\mathcal E')$ et si $g''$ est la projection $P'\to P$, $r'$
s'identifie à $r_{(P')}$ (II, 4.2.10) et $g''$ est fidèlement plat et
quasi-compact (1.1.2 et 2.2.13). On conclut donc de (2.7.1, (xi)) que $r$ est
une immersion, et par suite $\mathcal L$ est très ample (II, 4.4.4, b)).

Supposons maintenant que $\mathcal L'$ soit ample pour $f'$~; pour démontrer
que $\mathcal L$ est ample pour $f$, on peut se borner au cas où $Y$ est
affine (II, 4.6.4), et en vertu de (2.2.12) et de (II, 4.6.13), on peut aussi
supposer que $Y'$ est affine. Alors $X$ et $X'$ sont des

\oldpage[IV-2]{33}
schémas quasi-compacts, et pour prouver que $\mathcal L$ est $f$-ample, on
peut appliquer le critère de (II, 4.6.8, c)). Soit donc $\mathcal F$ un
$\mathcal O_X$-Module quasi-cohérent de type fini~; si
\[
\sigma:f^*(f_*(\mathcal F\otimes\mathcal L^{\otimes n}))
\longrightarrow\mathcal F\otimes\mathcal L^{\otimes n}
\]
est l'homomorphisme canonique, on voit comme ci-dessus que $g'^*(\sigma)$ est
l'homomorphisme composé
\[
f'^*(g^*(f_*(\mathcal F\otimes\mathcal L^{\otimes n})))
\xrightarrow{f'^*(u)}f'^*(f'_*(\mathcal F'\otimes
\mathcal L'^{\otimes n}))\xrightarrow{\sigma'}
\mathcal F'\otimes\mathcal L'^{\otimes n}
\]
en posant $\mathcal F'=g'^*(\mathcal F)$, tenant compte de
$(0_{\mathrm I},4.3.3.1)$ et désignant par $u$ l'homomorphisme canonique
\[
g^*(f_*(\mathcal F\otimes\mathcal L^{\otimes n}))
\longrightarrow f'_*(\mathcal F'\otimes\mathcal L'^{\otimes n}).
\]
Or, on sait que $u$ est bijectif quel que soit $n$ (2.3.1)~; d'autre part,
comme $\mathcal F'$ est quasi-cohérent et de type fini, l'hypothèse que
$\mathcal L'$ est ample pour $f'$ entraîne l'existence d'un $n_0$ tel que
pour $n\geq n_0$, $\sigma'$ soit surjectif~; on voit donc que
$g'^*(\sigma)$ est surjectif pour $n\geq n_0$, et comme $g'$ est fidèlement
plat, $\sigma$ est surjectif pour ces valeurs de $n$ (2.2.7), ce qui achève
la démonstration.

\begin{remark}[2.7.3]
\label{IV.2.7.3-fr}
\begin{enumerate}
\item[(i)] Il résulte de (2.6.1), (2.6.4) et (2.5.4.1) que les conclusions
de (2.7.1) sont encore valables dans les cas (i), (iii), (v), (vii) et
(xvi) lorsqu'on suppose seulement que $g$ est quasi-compact et
quasi-fidèlement plat (2.3.3)~; on a déjà remarqué que (2.7.1) est valable
dans le cas (ii) en supposant seulement $g$ surjectif et quasi-compact.

\item[(ii)] Avec les notations et hypothèses de (2.7.1), il peut se faire
que $f$ soit propre et $f'$ projectif sans que $f$ soit quasi-projectif. En
effet, Hironaka [34] a donné un exemple de morphisme propre non projectif
$f:X\to Y$, où $X$ et $Y$ sont deux schémas algébriques réguliers
$(0_{\mathrm I},4.1.4)$ sur un même corps $k$, $Y$ étant projectif sur $k$~;
en outre $Y$ est réunion de deux ouverts affines $Y_i$ $(i=1,2)$ tels que
$f_i:X\times_Y Y_i\to Y_i$ soit projectif pour $i=1,2$. Soit alors
$Y'=Y_1\amalg Y_2$ le préschéma somme~; il est clair que le morphisme
canonique $g:Y'\to Y$ coïncidant avec les injections canoniques dans $Y_1$
et $Y_2$, est fidèlement plat, et il est quasi-compact en vertu de
(I, 5.5.10)~; pourtant, bien que $f':X\times_Y Y'\to Y'$ (coïncidant avec
$f_i$ dans chacun des $Y_i$) soit projectif (II, 5.5.6), il n'en est pas de
même de $f$. Il existe donc un $\mathcal O_{X'}$-Module inversible
$\mathcal L'$ qui est $f'$-ample mais qui n'est pas de la forme
$g'^*(\mathcal L)$ pour un $\mathcal O_X$-Module inversible $\mathcal L$, en
vertu de (2.7.2).

\item[(iii)] Sous les hypothèses de (2.7.1), il peut se faire que $f'$ soit
un isomorphisme local sans que $f$ soit une immersion locale. Soient en
effet $k$ un corps, $\overline k$ une clôture algébrique de $k$, $K$ une
extension séparable de degré fini de $k$, distincte de $k$~; alors le
morphisme structural $f:X\to Y$, où $X=\operatorname{Spec}(K)$ et
$Y=\operatorname{Spec}(k)$, n'est pas une immersion locale, mais si on prend
$Y'=\operatorname{Spec}(\overline k)$, le morphisme $Y'\to Y$ est
fidèlement plat et quasi-compact et $f'$ est un isomorphisme local, puisque
$X'=X\times_Y Y'$ est somme d'un nombre fini de schémas isomorphes à $Y'$.
\end{enumerate}
\end{remark}

\subsection{Préschémas sur une base régulière de dimension 1; adhérence d'un
sous-préschéma fermé de la fibre générique}
\label{subsection:IV.2.8-fr}

\begin{proposition}[2.8.1]
\label{IV.2.8.1-fr}
Soient $Y$ un préschéma localement noethérien, régulier, irréductible et de
dimension 1, de point générique $\eta$, $f:X\to Y$ un morphisme,
$X_\eta=f^{-1}(\eta)$ la fibre au point

\oldpage[IV-2]{34}
générique, $i:X_\eta\to X$ le morphisme canonique. Soient $\mathcal F$ un
$\mathcal O_X$-Module quasi-cohérent, $\mathcal F_\eta=i^*(\mathcal F)$,
$\mathcal G'$ un $\mathcal O_{X_\eta}$-Module quotient de $\mathcal F_\eta$,
et soit $\overline{\mathcal G'}$ le $\mathcal O_X$-Module image de $\mathcal F$
par l'homomorphisme composé $(0_{\mathrm I},4.4.3.2)$
\[
\mathcal F\xrightarrow{\rho_{\mathcal F}}i_*(i^*(\mathcal F))
\longrightarrow i_*(\mathcal G').
\]
Alors $\overline{\mathcal G'}$ est un $\mathcal O_X$-Module quotient de
$\mathcal F$, quasi-cohérent et $f$-plat, tel que
$i^*(\overline{\mathcal G'})=\mathcal G'$, et c'est le seul
$\mathcal O_X$-Module quotient de $\mathcal F$ ayant ces propriétés.
\end{proposition}

\begin{proof}
Comme $i$ est quasi-compact et quasi-séparé (1.1.2 et 1.2.2), il résulte de
(1.7.4) que pour tout $\mathcal O_{X_\eta}$-Module quasi-cohérent
$\mathcal H'$, $i_*(\mathcal H')$ est un $\mathcal O_X$-Module
quasi-cohérent; en outre, pour tout ouvert $U$ de $X$, on a
\[
(i_*(\mathcal H'))|U=(i|U\cap X_\eta)_*(\mathcal H'|U\cap X_\eta)
\]
par définition $(0_{\mathrm I},3.4.1)$. Si l'on démontre la proposition
lorsque $X$ et $Y$ sont affines, elle en résultera par recollement dans le cas
général, vu l'assertion d'unicité valable dans le cas affine. Autrement dit,
on est ramené à prouver le

\begin{lemma}[2.8.1.1]
\label{IV.2.8.1.1-fr}
Soient $A$ un anneau noethérien régulier $(0,17.3.6)$, intègre de dimension
1, $K$ son corps des fractions, $M$ un $A$-module, $N'$ un $K$-module
quotient de $M_{(K)}=M\otimes_A K$ par un sous-$K$-module $P'$,
$\overline{N'}$ l'image de $M$ par l'homomorphisme composé
$M\to M_{(K)}\to N'$. Alors $\overline{N'}$ est un $A$-module plat, et est
l'unique module quotient $N$ de $M$ qui soit un $A$-module plat et tel que le
noyau de l'homomorphisme surjectif $M_{(K)}\to N_{(K)}$ soit égal à $P'$.
\end{lemma}

Comme pour tout idéal maximal $\mathfrak m$ de $A$, $A_{\mathfrak m}$ est un
anneau local régulier de dimension 1, donc un anneau de valuation discrète,
il revient au même de dire qu'un $A$-module $N$ est plat ou qu'il est
\emph{sans torsion} $(0_{\mathrm I},6.3.4)$. Comme $N'$ est un $K$-espace
vectoriel, c'est un $A$-module sans torsion, donc il en est de même de
$\overline{N'}$, sous-module de $N'$; en outre, il est immédiat de vérifier
que $\overline{N'}_{(K)}$ s'identifie à $N'$. Inversement, si $N$ est un
$A$-module quotient de $M$ ayant les propriétés de l'énoncé, le fait que $N$
soit un $A$-module plat entraîne que l'homomorphisme canonique
$N\to N_{(K)}=N\otimes_A K$ est injectif. Comme $N_{(K)}$ s'identifie à
$N'$, la conclusion résulte de la commutativité du diagramme
\[
\xymatrix{
M \ar[r] \ar[d] & N \ar[d] \\
M_{(K)} \ar[r] & N_{(K)}.
}
\]
\end{proof}

\begin{corollary}[2.8.2]
\label{IV.2.8.2-fr}
Sous les conditions de (2.8.1), pour que $\mathcal F$ soit $f$-plat, il faut
et il suffit que l'homomorphisme canonique
\[
\mathcal F\longrightarrow i_*(\mathcal F_\eta)=i_*(i^*(\mathcal F))
\]
soit injectif.
\end{corollary}

\begin{env}[2.8.3]
\label{IV.2.8.3-fr}
La formation du $\mathcal O_X$-Module $\overline{\mathcal G'}$ est
\emph{fonctorielle en $\mathcal F$ et $\mathcal G'$} : de façon précise, si
$\mathcal F_1$, $\mathcal F_2$ sont deux $\mathcal O_X$-Modules
quasi-cohérents, $u:\mathcal F_1\to\mathcal F_2$ un
$\mathcal O_X$-homomorphisme, $\mathcal G'_i$ un
$\mathcal O_{X_\eta}$-Module quotient de $(\mathcal F_i)_\eta$ $(i=1,2)$ et
$v:\mathcal G'_1\to\mathcal G'_2$ un homomorphisme rendant commutatif le
diagramme
\[
\xymatrix{
(\mathcal F_1)_\eta \ar[r]^{i^*(u)} \ar[d] &
(\mathcal F_2)_\eta \ar[d] \\
\mathcal G'_1 \ar[r]_v & \mathcal G'_2
}
\]

\oldpage[IV-2]{35}
(homomorphisme déterminé de façon unique (lorsqu'il existe) par cette
propriété), alors le diagramme
\[
\xymatrix{
\mathcal F_1 \ar[r]^u \ar[d] & \mathcal F_2 \ar[d] \\
i_*(\mathcal G'_1) \ar[r]_{i_*(v)} & i_*(\mathcal G'_2)
}
\]
est commutatif, et il y a par suite un seul homomorphisme
$w:\overline{\mathcal G'_1}\to\overline{\mathcal G'_2}$ rendant commutatif
le diagramme
\[
\xymatrix{
\mathcal F_1 \ar[r]^u \ar[d] & \mathcal F_2 \ar[d] \\
\overline{\mathcal G'_1} \ar[r]_w & \overline{\mathcal G'_2}
}
\]
\end{env}

\begin{proposition}[2.8.4]
\label{IV.2.8.4-fr}
Les hypothèses sur $Y$ étant celles de (2.8.1), soient $X_1$, $X_2$ deux
$Y$-préschémas, $\mathcal F_i$ un $\mathcal O_{X_i}$-Module
quasi-cohérent, $\mathcal G'_i$ un $\mathcal O_{(X_i)_\eta}$-Module quotient
de $(\mathcal F_i)_\eta$ $(i=1,2)$. Alors on a
\end{proposition}

\begin{env}[2.8.4.1]
\label{IV.2.8.4.1-fr}
\[
\overline{\mathcal G'_1}\otimes_{\mathcal O_Y}\overline{\mathcal G'_2}
=\overline{\mathcal G'_1\otimes_{k(\eta)}\mathcal G'_2}.
\]
\end{env}

\begin{proof}
En effet, posons $X=X_1\times_Y X_2$; le premier membre de (2.8.3.1) est
un $\mathcal O_X$-Module quasi-cohérent qui est $Y$-plat
$(0_{\mathrm I},6.2.1)$, dont l'image réciproque dans $X_\eta$ est
$\mathcal G'_1\otimes_{k(\eta)}\mathcal G'_2$ (I, 9.1.5) et qui est un
quotient de $\mathcal F_1\otimes_{\mathcal O_Y}\mathcal F_2$; la conclusion
résulte donc de la propriété d'unicité de (2.8.1).
\end{proof}

\begin{proposition}[2.8.5]
\label{IV.2.8.5-fr}
Les hypothèses sur $X$ et $Y$ étant celles de (2.8.1), soit $Z'$ un
sous-préschéma fermé de $X_\eta$. Il existe alors un unique sous-préschéma
fermé $\overline{Z'}$ de $X$ qui soit $Y$-plat et tel que
$i^{-1}(\overline{Z'})=Z'$.
\end{proposition}

\begin{proof}
Si $\mathcal J'$ est l'Idéal quasi-cohérent de $\mathcal O_{X_\eta}$
définissant $Z'$, il suffit d'appliquer (2.8.1) au cas où
$\mathcal F=\mathcal O_X$ et
$\mathcal G'=\mathcal O_{X_\eta}/\mathcal J'$; si
$\overline{\mathcal G'}=\mathcal O_X/\mathcal J$, on a en effet
$\mathcal J'=i^*(\mathcal J)\mathcal O_{X_\eta}$, donc
$i^{-1}(\overline{Z'})=Z'$ (I, 4.4.5).
\end{proof}

On notera que le préschéma $\overline{Z'}$ est l'\emph{image fermée} de $Z'$
par le morphisme composé
\[
Z'\longrightarrow X_\eta\xrightarrow{i}X,
\]
où la première flèche est l'injection canonique (I, 9.5.3); son espace
sous-jacent est l'adhérence \emph{dans $X$} de $Z'$ (I, 9.5.4), ce qui
justifie la notation adoptée. On dit encore que $\overline{Z'}$ est le
sous-préschéma de $X$ \emph{adhérence} de $Z'$.

\begin{corollary}[2.8.6]
\label{IV.2.8.6-fr}
Soient $X_1$, $X_2$ deux $Y$-préschémas, $Z'_i$ un sous-préschéma fermé de
$(X_i)_\eta$ $(i=1,2)$. Alors on a
\end{corollary}

\begin{env}[2.8.6.1]
\label{IV.2.8.6.1-fr}
\[
\overline{Z'_1}\times_Y\overline{Z'_2}
=\overline{Z'_1\times_{k(\eta)}Z'_2}.
\]
\end{env}

\begin{proof}
Cela résulte de (2.8.4) et (2.8.5).
\end{proof}

\oldpage[IV-2]{36}
\section{Cycles premiers associés et décompositions primaires}

Nous donnons surtout dans ce paragraphe la traduction des résultats sur les
modules exposés dans Bourbaki, \emph{Alg. comm.}, chap.~IV, que nous suivons de
très près. Les notions qui suivent ne semblent guère avoir d'intérêt que dans
le cas de préschémas localement noethériens.

\subsection{Cycles premiers associés à un Module}

\begin{definition}[3.1.1]
\label{IV.2.3.1.1-def-fr}
Soient $X$ un préschéma, $\mathcal F$ un $\mathcal O_X$-Module
quasi-cohérent. On dit qu'un point $x\in X$ est associé à $\mathcal F$ si
l'idéal maximal $\mathfrak m_x$ de $\mathcal O_x$ est associé au
$\mathcal O_x$-Module $\mathcal F_x$ (autrement dit, si $\mathfrak m_x$ est
l'annulateur d'un élément de $\mathcal F_x$). On dénote par
$\operatorname{Ass}(\mathcal F)$ l'ensemble des $x\in X$ associés à
$\mathcal F$. On dit qu'une partie fermée irréductible $Z$ de $X$ est un
\emph{cycle premier associé à $\mathcal F$} si son point générique est
associé à $\mathcal F$. Lorsque $\mathcal F=\mathcal O_X$, on dira aussi que
les points (resp. \emph{cycles premiers}) associés à $\mathcal F$ sont
associés au préschéma $X$.

On dit qu'un cycle premier associé à $\mathcal F$ (resp. $X$) est
\emph{immergé} s'il est contenu dans un \emph{autre} cycle premier associé à
$\mathcal F$ (resp. $X$) (autrement dit, s'il n'est pas \emph{maximal} dans
l'ensemble de ces cycles).
\end{definition}

Si $X$ est localement noethérien, les \emph{composantes irréductibles} de $X$
ne sont autres que les cycles premiers \emph{maximaux} (ou \emph{non immergés})
associés à $X$.

Il est clair que si $x\in\operatorname{Ass}(\mathcal F)$, on a
$\mathcal F_x\ne0$, autrement dit

\begin{env}[3.1.1.1]
\label{IV.2.3.1.1.1-fr}
\[
\operatorname{Ass}(\mathcal F)\subset\operatorname{Supp}(\mathcal F).
\]
\end{env}

Si $x\in X$ est associé à $\mathcal F$, il est évidemment associé à
$\mathcal F|U$ pour tout voisinage ouvert $U$ de $x$, et réciproquement, s'il
est associé à $\mathcal F|U$ pour un de ces voisinages, il est associé à
$\mathcal F$.

On notera enfin que pour un $\mathcal O_X$-Module quasi-cohérent $\mathcal F$,
l'existence de cycles premiers associés immergés est une question \emph{locale},
car si $y$ et $z$ sont deux points de $\operatorname{Ass}(\mathcal F)$ tels
que $y\in\overline{\{z\}}$, tout voisinage de $y$ contient $z$.

\begin{proposition}[3.1.2]
\label{IV.2.3.1.2-fr}
Soient $A$ un anneau noethérien, $M$ un $A$-module,
$X=\operatorname{Spec}(A)$, $\mathcal F=\widetilde M$~; pour qu'un point
$x\in X$ soit associé à $\mathcal F$, il faut et il suffit que l'idéal
premier $\mathfrak j_x$ de $A$ soit associé au module $M$ (autrement dit,
soit l'annulateur d'un élément $f\in M$).
\end{proposition}

Cela résulte de la définition (3.1.1) et de Bourbaki, \emph{loc. cit.},
\S~1, n\textsuperscript{o}~2, cor. de la prop.~5, appliqué à
$S=A-\mathfrak j_x$.

\begin{proposition}[3.1.3]
\label{IV.2.3.1.3-fr}
Soient $X$ un préschéma localement noethérien, $\mathcal F$ un
$\mathcal O_X$-Module quasi-cohérent, $x$ un point de $X$, $Z$ le
sous-préschéma fermé réduit de $X$ ayant $\overline{\{x\}}$ pour espace
sous-jacent (I, 5.2.1). Les conditions suivantes sont équivalentes~:
\begin{enumerate}
\item[\emph{a)}] $x\in\operatorname{Ass}(\mathcal F)$.

\oldpage[IV-2]{37}
\item[\emph{b)}] Il existe un voisinage ouvert $U$ de $x$ et une section
$f\in\Gamma(U,\mathcal F)$ tels que $U\cap Z$ soit égal à
$\operatorname{Supp}(\mathcal O_U\cdot f)$.
\item[\emph{b$'$)}] Il existe un voisinage ouvert $U$ de $x$ et une section
$f\in\Gamma(U,\mathcal F)$ tels que $U\cap Z$ soit une composante
irréductible de $\operatorname{Supp}(\mathcal O_U\cdot f)$.
\item[\emph{c)}] Il existe un voisinage ouvert $U$ de $x$ et un sous-Module
de $\mathcal F|U$ isomorphe à $\mathcal O_Z|U$ ($\mathcal O_Z$ étant
identifié à un quotient de $\mathcal O_X$).
\item[\emph{c$'$)}] Il existe un voisinage ouvert $U$ de $x$ et un
sous-Module cohérent $\mathcal G$ de $\mathcal F|U$ tels que $U\cap Z$ soit
une composante irréductible de $\operatorname{Supp}(\mathcal G)$.
\end{enumerate}
\end{proposition}

Il est clair que \emph{c)} implique \emph{b)}, car il suffit de prendre pour
$f$ l'élément de $\Gamma(U,\mathcal F)$ qui correspond à la section unité de
$\mathcal O_Z|U$. Comme $U\cap Z$ est irréductible $(0_{\mathrm I},2.1.6)$,
\emph{b)} implique \emph{b$'$)} et \emph{b$'$)} implique \emph{c$'$)} puisque
$\mathcal O_X$ est cohérent $(0_{\mathrm I},5.3.4)$. Pour voir que
\emph{c$'$)} implique \emph{a)}, on peut se borner au cas où
$U=X=\operatorname{Spec}(A)$ est affine, $A$ étant donc noethérien, et où
$\mathcal F=\widetilde M$, où $M$ est un $A$-module, et
$\mathcal G=\widetilde N$, où $N$ est un sous-module de $M$, de type fini.
On sait alors que les éléments minimaux de $\operatorname{Supp}(\mathcal G)$
sont les points maximaux de $V(\operatorname{Ann}(N))$ $(0_{\mathrm I},1.7.4)$,
et ce sont aussi les éléments minimaux de $\operatorname{Ass}(N)$ (Bourbaki,
\emph{Alg. comm.}, chap.~IV, \S~1, n\textsuperscript{o}~3, cor.~1 de la
prop.~7)~; comme $\operatorname{Ass}(N)\subset\operatorname{Ass}(M)
=\operatorname{Ass}(\mathcal F)$, on voit que \emph{c$'$)} implique
\emph{a)}. Enfin \emph{a)} implique \emph{c)} en vertu de (3.1.2), en
prenant encore $X$ affine, $\mathcal F=\widetilde M$, et $Z$ défini par
l'idéal $fA$ (avec les notations de (3.1.2)).

\begin{corollary}[3.1.4]
\label{IV.2.3.1.4-fr}
Soient $X$ un préschéma localement noethérien, $\mathcal F$ un
$\mathcal O_X$-Module cohérent. Alors les cycles premiers maximaux associés à
$\mathcal F$ sont les composantes irréductibles de
$\operatorname{Supp}(\mathcal F)$, et les points génériques de ces composantes
sont les points $x\in X$ tels que $\mathcal F_x$ soit un
$\mathcal O_x$-Module de longueur finie et $\ne0$.
\end{corollary}

En effet, si $x$ est le point générique d'une des composantes irréductibles
$Z$ de $\operatorname{Supp}(\mathcal F)$, il résulte de l'équivalence de
\emph{a)} et \emph{c$'$)} dans (3.1.3) que $x$ appartient à
$\operatorname{Ass}(\mathcal F)$, et $Z$ est un cycle premier associé à
$\mathcal F$, nécessairement maximal en vertu de (3.1.1.1)~; la réciproque
résulte trivialement de (3.1.1.1)~; enfin, la dernière assertion étant
évidemment locale résulte de Bourbaki, \emph{Alg. comm.}, chap.~IV, \S~2,
n\textsuperscript{o}~5, cor.~2 de la prop.~7.

\begin{corollary}[3.1.5]
\EGAReaderTarget{ega4.c04.cor.3-1-5}{IV.2.3.1.5-fr}
Soient $X$ un préschéma localement noethérien, $\mathcal F$ un
$\mathcal O_X$-Module quasi-cohérent. Pour que $\mathcal F=0$, il faut et il
suffit que $\operatorname{Ass}(\mathcal F)=\varnothing$.
\end{corollary}

La question étant locale, on est ramené au cas où $X$ est affine, et la
conclusion résulte aussitôt de (3.1.2) et de Bourbaki, \emph{Alg. comm.},
chap.~IV, \S~1, n\textsuperscript{o}~1, cor.~1 de la prop.~2.

\begin{proposition}[3.1.6]
\label{IV.2.3.1.6-fr}
Soient $X$ un préschéma localement noethérien, $\mathcal F$ un
$\mathcal O_X$-Module cohérent~; alors $\operatorname{Ass}(\mathcal F)$ est
\emph{localement fini} (c'est-à-dire que tout point de $X$ admet un voisinage
dont l'intersection avec $\operatorname{Ass}(\mathcal F)$ est finie).
\end{proposition}

Il suffit de considérer le cas où $X$ est affine, donc noethérien, et alors la
proposition résulte de (3.1.2) et de Bourbaki, \emph{Alg. comm.}, chap.~IV,
\S~1, n\textsuperscript{o}~4, cor. du th.~2.

\oldpage[IV-2]{38}
\begin{proposition}[3.1.7]
\EGAReaderTarget{ega4.c04.prop.3-1-7}{IV.2.3.1.7-fr}
Soit $X$ un préschéma.
\begin{enumerate}
\item[\emph{(i)}] Soit $0\to\mathcal F'\to\mathcal F\to\mathcal F''\to0$
une suite exacte de $\mathcal O_X$-Modules quasi-cohérents. Alors
\[
\operatorname{Ass}(\mathcal F')\subset\operatorname{Ass}(\mathcal F)
\subset\operatorname{Ass}(\mathcal F')\cup\operatorname{Ass}(\mathcal F'').
\]
\item[\emph{(ii)}] Soient $\mathcal F$ un $\mathcal O_X$-Module
quasi-cohérent, $(\mathcal F_\alpha)$ une famille de sous-$\mathcal O_X$-Modules
quasi-cohérents de $\mathcal F$ telle que $\mathcal F$ soit réunion des
$\mathcal F_\alpha$. Alors
$\operatorname{Ass}(\mathcal F)=\bigcup_\alpha\operatorname{Ass}(\mathcal F_\alpha)$.
\item[\emph{(iii)}] Pour toute famille $(\mathcal F_\alpha)$ de
$\mathcal O_X$-Modules quasi-cohérents, on a
$\operatorname{Ass}(\bigoplus_\alpha\mathcal F_\alpha)
=\bigcup_\alpha\operatorname{Ass}(\mathcal F_\alpha)$.
\end{enumerate}
\end{proposition}

On est aussitôt ramené aux propositions correspondantes pour les modules
(Bourbaki, \emph{loc. cit.}, \S~1, n\textsuperscript{o}~1, formule (1), prop.~3
et cor.~1 de la prop.~3).

\begin{proposition}[3.1.8]
\label{IV.2.3.1.8-fr}
Soient $X$ un préschéma localement noethérien, $\mathcal F$ un
$\mathcal O_X$-Module quasi-cohérent, $U$ un ouvert de $X$, $\mathcal J$ un
Idéal cohérent de $\mathcal O_X$ définissant un sous-préschéma fermé de $X$
ayant $X-U$ pour espace sous-jacent. Les conditions suivantes sont
équivalentes~:
\begin{enumerate}
\item[\emph{a)}] $\operatorname{Ass}(\mathcal F)\subset U$.
\item[\emph{b)}] Pour tout ouvert affine $V$, toute section de $\mathcal F$
au-dessus de $V$, dont la restriction à $V\cap U$ est nulle, est égale à $0$.
\item[\emph{c)}] L'homomorphisme canonique
$\mathcal F\to\mathcal Hom_{\mathcal O_X}(\mathcal J,\mathcal F)$ est
injectif.
\end{enumerate}
\end{proposition}

La question étant locale, on peut supposer $X=\operatorname{Spec}(A)$ affine,
$A$ étant un anneau noethérien, $\mathcal F=\widetilde M$, où $M$ est un
$A$-module, et $\mathcal J=\widetilde{\mathfrak J}$, où $\mathfrak J$ est un
idéal de $A$. L'homomorphisme $M\to\operatorname{Hom}_A(\mathfrak J,M)$ fait
correspondre à $m\in M$ l'homomorphisme $x\mapsto xm$ de $\mathfrak J$ dans
$M$~; dire qu'il n'est pas injectif signifie qu'il existe un $m\ne0$ dans
$M$ tel que $\mathfrak Jm=0$.

Montrons d'abord que \emph{c)} implique \emph{a)}~: en effet, si
$\operatorname{Ass}(\mathcal F)$ rencontrait alors $X-U$, il y aurait un
idéal premier $\mathfrak p\in\operatorname{Ass}(M)$ contenant $\mathfrak J$,
donc un élément $m\ne0$ de $M$ tel que $\mathfrak Jm=0$. En second lieu,
\emph{b)} implique \emph{c)}~: en effet, si l'on avait alors
$\mathfrak Jm=0$ pour un $m\ne0$ dans $M$, pour tout idéal premier
$\mathfrak q\not\supset\mathfrak J$, il existe un $a\in\mathfrak J$ non dans
$\mathfrak q$, donc la relation $am=0$ entraîne que l'image canonique de $m$
dans $M_\mathfrak q$ est $0$, autrement dit $m$ serait une section $\ne0$ de
$\mathcal F$ au-dessus de $X$ dont la restriction à $U$ serait nulle.

Enfin \emph{a)} entraîne \emph{b)}. Pour le voir, notons que l'homomorphisme
canonique
\[
M\longrightarrow\prod_{\mathfrak p\in\operatorname{Ass}(M)}M_\mathfrak p
\]
est \emph{injectif}~: en effet, si $N$ est le noyau de cet homomorphisme, on a
$\operatorname{Ass}(N)\subset\operatorname{Ass}(M)$~; s'il existait un
$\mathfrak p\in\operatorname{Ass}(N)$, il y aurait un élément $n\in N$ dont
$\mathfrak p$ serait l'annulateur~; mais par définition de $N$, il y a un
élément $s\notin\mathfrak p$ tel que $sn=0$, ce qui est absurde~; on en conclut
que $\operatorname{Ass}(N)=\varnothing$, d'où $N=0$. Or la condition
\emph{a)} entraîne que si $m\in M$ est une section de $\mathcal F$ dont la
restriction à $U$ est nulle, l'image canonique de $m$ dans $M_\mathfrak p$ est
nulle pour tout $\mathfrak p\in\operatorname{Ass}(M)$, donc $m=0$. C.Q.F.D.

\begin{corollary}[3.1.9]
\label{IV.2.3.1.9-fr}
Soient $X$ un préschéma localement noethérien, $\mathcal F$ un
$\mathcal O_X$-Module quasi-cohérent, $f$ une section de $\mathcal O_X$
au-dessus de $X$. Pour que $f$ soit $\mathcal F$-régulière
$(0,15.2.2)$, il faut et il suffit que
$\operatorname{Ass}(\mathcal F)\subset X_f$.
\end{corollary}

En effet, il est immédiat que dire que $f$ est $\mathcal F$-régulière signifie
que l'homomorphisme canonique
$\mathcal F\to\mathcal Hom_{\mathcal O_X}(f\mathcal O_X,\mathcal F)$ est
injectif, et il suffit d'appliquer (3.1.8) à l'Idéal
$\mathcal J=f\mathcal O_X$.

\oldpage[IV-2]{39}
\begin{proposition}[3.1.10]
\label{IV.2.3.1.10-fr}
Soient $X$, $Y$ des préschémas localement noethériens, $f:X\to Y$ un
morphisme entier. Alors, pour tout $\mathcal O_X$-Module quasi-cohérent
$\mathcal F$, on a
$f(\operatorname{Ass}(\mathcal F))=\operatorname{Ass}(f_*(\mathcal F))$.
\end{proposition}

La question étant locale sur $Y$ et le morphisme $f$ étant affine, on est
immédiatement ramené au cas où $Y$ est affine, autrement dit au

\begin{lemma}[3.1.10.1]
\label{IV.2.3.1.10.1-fr}
Soient $A$, $B$ deux anneaux noethériens, $\rho:A\to B$ un homomorphisme
d'anneaux faisant de $B$ une $A$-algèbre entière, $M$ un $B$-module. Alors les
idéaux premiers $\mathfrak p\in\operatorname{Ass}(M_{[\rho]})$ sont les images
réciproques par $\rho$ des idéaux premiers
$\mathfrak q\in\operatorname{Ass}(M)$.
\end{lemma}

En effet, si $\mathfrak q\in\operatorname{Ass}(M)$, $\mathfrak q$ est
l'annulateur dans $B$ d'un élément $x\in M$, donc
$\rho^{-1}(\mathfrak q)$ est l'annulateur dans $A$ de $x$. Inversement, soit
$\mathfrak p\in\operatorname{Ass}(M_{[\rho]})$, de sorte que $\mathfrak p$
est l'image réciproque par $\rho$ de l'annulateur $\mathfrak b$ dans $B$ d'un
élément $x\in M$~; il résulte donc du premier théorème de Cohen-Seidenberg
qu'il existe un idéal premier $\mathfrak q$ de $B$ contenant $\mathfrak b$ et
dont l'image réciproque est $\mathfrak p$ (Bourbaki, \emph{Alg. comm.},
chap.~V, \S~2, n\textsuperscript{o}~1, cor.~2 du th.~1)~; en considérant un
des idéaux premiers minimaux parmi ceux contenus dans $\mathfrak q$ et
contenant $\mathfrak b$, on peut évidemment supposer que $\mathfrak q$
lui-même est un de ces idéaux minimaux. Mais comme $B\cdot x\subset M$ est
isomorphe à $B/\mathfrak b$, on sait que l'on a alors
$\mathfrak q\in\operatorname{Ass}(B/\mathfrak b)\subset\operatorname{Ass}(M)$
(Bourbaki, \emph{Alg. comm.}, chap.~IV, \S~1, n\textsuperscript{o}~4,
th.~2).

\begin{corollary}[3.1.11]
\label{IV.2.3.1.11-fr}
Sous les hypothèses de (3.1.10), pour que $\mathcal F$ soit sans cycle premier
associé immergé, il suffit qu'il en soit de même de $f_*(\mathcal F)$.
\end{corollary}

Supposons en effet que $f_*(\mathcal F)$ n'ait pas de cycle premier associé
immergé. Notons que si $A$ est une algèbre entière sur un corps $k$, tous les
idéaux premiers de $A$ sont maximaux (Bourbaki, \emph{Alg. comm.}, chap.~V,
\S~2, n\textsuperscript{o}~1, prop.~1)~; il résulte donc de (I, 6.2.2) que les
fibres de $f$ sont des espaces \emph{discrets}~; si $x$, $x'$ sont deux points
distincts de $\operatorname{Ass}(\mathcal F)$, aucun d'eux ne peut donc être
adhérent à l'autre si $f(x)=f(x')$~; et si $f(x)\ne f(x')$, (3.1.10) et
l'hypothèse entraînent qu'aucun des deux points $f(x)$, $f(x')$ ne peut être
adhérent à l'autre, donc il en est de même de $x$ et $x'$.

\begin{remark}[3.1.12]
\label{IV.2.3.1.12-fr}
Sous les hypothèses de (3.1.10), il peut se faire par contre que $\mathcal F$
soit sans cycle premier associé immergé, mais non $f_*(\mathcal F)$. Prenons
par exemple $Y=\operatorname{Spec}(k[T])$ où $k$ est un corps («~droite
affine~»), et $X$ somme de $X_1=Y$ et de $X_2=\operatorname{Spec}(k)$, le
morphisme $X_2\to Y$ correspondant à l'homomorphisme canonique
$k[T]\to k[T]/\mathfrak m$, où $\mathfrak m$ est l'idéal maximal $(T)$. Il
est clair que le morphisme $f:X\to Y$ est fini~; si l'on prend
$\mathcal F=\mathcal O_X$, $\mathcal F$ est sans cycle premier associé
immergé, mais $f_*(\mathcal F)=\widetilde M$, où $M$ est le $k[T]$-module
somme directe de $k[T]$ et de $k$, donc $\operatorname{Ass}(M)$ est formé du
point générique $(0)$ de $Y$ et du point $\mathfrak m$.
\end{remark}

\begin{proposition}[3.1.13]
\label{IV.2.3.1.13-fr}
Soient $X$ un préschéma localement noethérien, $U$ un ouvert de $X$,
$i:U\to X$ l'injection canonique. Pour tout $\mathcal O_U$-Module
quasi-cohérent $\mathcal F$, on a
$\operatorname{Ass}(i_*(\mathcal F))=\operatorname{Ass}(\mathcal F)$.
\end{proposition}

Rappelons que $i_*(\mathcal F)$ est un $\mathcal O_X$-Module quasi-cohérent
(I, 2.2 et 1.7.4)~; comme $i_*(\mathcal F)|U=\mathcal F$, on a
$(\operatorname{Ass}(i_*(\mathcal F)))\cap U=\operatorname{Ass}(\mathcal F)$,
et il reste donc à prouver que $\operatorname{Ass}(i_*(\mathcal F))\subset U$.
Mais d'après (3.1.8), cette relation signifie que pour tout ouvert affine $V$
de $X$, toute section de $i_*(\mathcal F)$ au-dessus de $V$ qui est nulle
dans $U\cap V$, est nulle, condition trivialement vérifiée puisque
$\Gamma(V,i_*(\mathcal F))=\Gamma(U\cap V,\mathcal F)
=\Gamma(U\cap V,i_*(\mathcal F))$.

\oldpage[IV-2]{40}
\subsection{Décompositions irredondantes}

\begin{proposition}[3.2.1]
\label{IV.2.3.2.1-fr}
Soient $X$ un préschéma localement noethérien, $U$ un ouvert dense dans $X$.
Les conditions suivantes sont équivalentes~:
\begin{enumerate}
\item[\emph{a)}] $X$ est réduit.
\item[\emph{b)}] Le sous-préschéma induit sur $U$ est réduit et $X$ est sans
cycle premier immergé.
\item[\emph{c)}] $X$ est sans cycle premier immergé et pour tout point
générique $x$ d'une composante irréductible de $X$, on a
$\operatorname{long}(\mathcal O_x)=1$.
\end{enumerate}
Les cycles premiers associés à $X$ sont alors identiques aux composantes
irréductibles de $X$.
\end{proposition}

Il est clair que si $X$ est réduit il en est de même du sous-préschéma induit
sur $U$. En outre, l'existence des cycles premiers immergés étant locale, on
peut se borner à considérer le cas où $X=\operatorname{Spec}(A)$ est affine,
$A$ étant noethérien. Si $A$ est réduit, on sait que les idéaux premiers
minimaux de $A$ forment une décomposition primaire réduite de $(0)$
(Bourbaki, \emph{Alg. comm.}, chap.~IV, \S~2, n\textsuperscript{o}~5,
prop.~10) et sont les éléments de $\operatorname{Ass}(A)$, donc il n'existe
pas d'idéaux premiers immergés associés à $A$, ce qui montre que \emph{a)}
implique \emph{b)}. Il est immédiat que \emph{b)} entraîne \emph{c)}, car un
point générique $x$ d'une composante irréductible de $X$ appartient à $U$,
donc $\mathcal O_x$ est un corps. Enfin, \emph{c)} entraîne \emph{a)}~: il
suffit en effet de remarquer que si $\mathcal N$ est le Nilradical de
$\mathcal O_X$, qui est un Idéal cohérent, $\operatorname{Supp}(\mathcal N)$
ne peut contenir aucun des points génériques des composantes irréductibles de
$X$ par hypothèse~; si $\operatorname{Supp}(\mathcal N)$ n'était pas vide et
si $x$ était un des points maximaux de cet ensemble fermé, le critère
(3.1.3, \emph{c$'$)}) montrerait que $x\in\operatorname{Ass}(\mathcal O_X)$,
et $\overline{\{x\}}$ serait donc un cycle premier \emph{immergé} de $X$,
contrairement à l'hypothèse~; donc $\mathcal N=0$.

\begin{definition}[3.2.2]
\label{IV.2.3.2.2-fr}
Soient $X$ un préschéma localement noethérien, $\mathcal F$ un
$\mathcal O_X$-Module cohérent. On dit que $\mathcal F$ est \emph{réduit}
s'il vérifie les deux conditions suivantes~: 1\textsuperscript{o}
$\mathcal F$ est sans cycle premier associé immergé~; 2\textsuperscript{o}
pour tout point maximal $x$ de $\operatorname{Supp}(\mathcal F)$, on a
$\operatorname{long}(\mathcal F_x)=1$.
\end{definition}

La condition 1\textsuperscript{o} signifie que les cycles premiers associés à
$\mathcal F$ sont les composantes irréductibles de
$\operatorname{Supp}(\mathcal F)$ (3.1.4), et la condition
2\textsuperscript{o} signifie que pour tout point générique $x$ d'une telle
composante on a $\operatorname{long}(\mathcal F_x)=1$.

Pour un schéma affine $X$, cette définition donne en particulier la notion de
\emph{module réduit} sur un anneau noethérien $A$~; un $A$-module de type fini
$M$ est dit \emph{réduit}, s'il n'a pas d'idéaux premiers associés immergés et
si, pour tout $\mathfrak p\in\operatorname{Ass}(M)$,
$\operatorname{long}_{A_\mathfrak p}(M_\mathfrak p)=1$. Revenant au cas d'un
préschéma localement noethérien $X$ et d'un $\mathcal O_X$-Module cohérent
$\mathcal F$, on dit que $\mathcal F$ est \emph{réduit en un point} $x\in X$
si $\mathcal F_x$ est un $\mathcal O_x$-module réduit~; cela signifie encore
que, sur le schéma local $\operatorname{Spec}(\mathcal O_x)$,
$\widetilde{\mathcal F_x}$ est réduit~; il revient donc au même de dire
\emph{que $x$ n'appartient à aucun cycle premier associé immergé de
$\mathcal F$ et que $\operatorname{long}(\mathcal F_z)=1$ pour tous les
points maximaux $z$ de $\operatorname{Supp}(\mathcal F)$ tels que
$x\in\overline{\{z\}}$}. Il est clair que si $\mathcal F$ est un
$\mathcal O_X$-Module cohérent réduit, il est réduit en chaque point de $X$~;
inversement, si $\mathcal F$ est \emph{réduit au point} $x$, il existe un
voisinage ouvert $U$ de $x$ tel que $\mathcal F|U$ soit un
$\mathcal O_U$-Module \emph{réduit}~: il suffit en effet de prendre $U$ ne
rencontrant aucun cycle premier associé immergé de $\mathcal F$ (un tel
voisinage existe puisque ces cycles forment un ensemble localement fini de
parties
\oldpage[IV-2]{41}
fermées de $X$). Dire que $\mathcal O_X$ est réduit en un point $x$ équivaut à
dire que $X$ est \emph{réduit au point $x$}.

\begin{proposition}[3.2.3]
\label{IV.2.3.2.3-fr}
Soient $X$ un préschéma localement noethérien, $U$ un ouvert dans $X$,
$\mathcal F$ un $\mathcal O_X$-Module cohérent tel que
$U\cap\operatorname{Supp}(\mathcal F)$ soit dense dans
$\operatorname{Supp}(\mathcal F)$. Les conditions suivantes sont
équivalentes~:
\begin{enumerate}
\item[\emph{a)}] $\mathcal F$ est réduit.
\item[\emph{b)}] $\mathcal F|U$ est réduit, et $\mathcal F$ est sans cycle
premier associé immergé.
\item[\emph{c)}] Il existe un sous-préschéma fermé réduit $X'$ de $X$ et un
$\mathcal O_{X'}$-Module cohérent $\mathcal F'$ sans torsion et de rang $1$
sur toute composante irréductible de $X'$, tels que si $j:X'\to X$ est
l'injection canonique, on ait $j_*(\mathcal F')=\mathcal F$.
\end{enumerate}
En outre, lorsqu'il en est ainsi le sous-préschéma $X'$ est défini par l'Idéal
$\mathcal J$ de $\mathcal O_X$ annulateur de $\mathcal F$.
\end{proposition}

Pour voir que \emph{c)} implique \emph{a)}, on peut, en vertu de (3.1.3), se
limiter au cas où $X'=X$ est intègre, de point générique $x$ et l'on peut en
outre supposer $X=\operatorname{Spec}(A)$ affine et
$\mathcal F=\widetilde M$, où $A$ est donc intègre et $M$ est un $A$-module
sans torsion de rang $1$~; l'annulateur de tout élément de $M$ étant alors
réduit à $0$, on a $\operatorname{Ass}(\mathcal F)=\{x\}$ et
$\mathcal F_x$ est isomorphe à $\mathcal O_x$, corps des fractions de $A$,
donc les conditions de la définition (3.2.2) sont satisfaites. Comme
l'existence de cycles premiers associés immergés est locale, il est clair que
si $\mathcal F$ n'a pas de tels cycles et si
$U\cap\operatorname{Supp}(\mathcal F)$ est dense dans
$\operatorname{Supp}(\mathcal F)$, on a
$\operatorname{Ass}(\mathcal F)=\operatorname{Ass}(\mathcal F|U)$, donc
\emph{a)} et \emph{b)} sont équivalentes. Si \emph{a)} est satisfaite,
prenons pour $X'$ le sous-préschéma fermé réduit de $X$ dont
$\operatorname{Supp}(\mathcal F)$ est l'espace sous-jacent (I, 5.2.1), et soit
$\mathcal F'=j^*(\mathcal F)$~; un point $x$ de
$\operatorname{Ass}(\mathcal F)$ est nécessairement point maximal de $X'$ et
comme $\operatorname{long}(\mathcal F_x)=1$, $\mathcal F'_x$ est isomorphe à
$\mathcal F_x$, donc au corps $k(x)=\mathcal O_{X',x}$, ce qui prouve que
$\mathcal F'$ est sans torsion et de rang $1$ (I, 7.4.6 et 7.4.2). Enfin, la
dernière assertion est triviale, puisque pour tout $y\in X'$, l'annulateur du
$\mathcal O_{X',y}$-Module $\mathcal F'_y$ est nul.

\begin{definition}[3.2.4]
\label{IV.2.3.2.4-fr}
Soient $X$ un préschéma localement noethérien, $\mathcal F$ un
$\mathcal O_X$-Module quasi-cohérent. On dit que $\mathcal F$ est
\emph{irredondant} si $\operatorname{Ass}(\mathcal F)$ est réduit à un seul
élément $x$~; si $\mathcal F$ est de type fini, on dit que $\mathcal F$ est
\emph{intègre} si de plus $\mathcal F$ est réduit (autrement dit si
$\operatorname{long}(\mathcal F_x)=1$). On dit qu'un sous-$\mathcal O_X$-Module
quasi-cohérent $\mathcal G$ d'un $\mathcal O_X$-Module quasi-cohérent
$\mathcal F$ est \emph{primaire dans $\mathcal F$} si
$\mathcal F/\mathcal G$ est irredondant.
\end{definition}

Pour un schéma affine $X$, cette définition donne en particulier la notion de
\emph{module intègre} sur un anneau noethérien $A$~; un $A$-module $M$ est dit
\emph{intègre} s'il est de type fini, si $M$ est \emph{primaire} (c'est-à-dire
que $\operatorname{Ass}(M)$ est réduit à un seul idéal premier
$\mathfrak p$) et si en outre
$\operatorname{long}_{A_\mathfrak p}(M_\mathfrak p)=1$. Revenant au cas d'un
préschéma localement noethérien quelconque $X$ et d'un
$\mathcal O_X$-Module cohérent $\mathcal F$, on dit que $\mathcal F$ est
\emph{intègre en un point} $x\in X$ si $\mathcal F_x$ est un
$\mathcal O_x$-module \emph{intègre}~: cela signifie encore que $x$
\emph{n'appartient qu'à un seul cycle premier associé (nécessairement non
immergé) à $\mathcal F$ et que $\operatorname{long}(\mathcal F_z)=1$ en son
point générique $z$}. Il est clair que si $\mathcal F$ est un
$\mathcal O_X$-Module cohérent intègre, il est intègre en tout point de $X$~;
inversement, si $\mathcal F$ est \emph{intègre en un point} $x$, il existe un
voisinage ouvert $U$ de $x$ tel que $\mathcal F|U$

\oldpage[IV-2]{42}
soit un $\mathcal O_U$-Module \emph{intègre}~: il suffit en effet de prendre
$U$ tel que $\mathcal F|U$ soit un $\mathcal O_U$-Module réduit (3.2.2).

On dit que le préschéma localement noethérien $X$ est \emph{irredondant} si
$\mathcal O_X$ est irredondant (ce qui implique que $X$ est
\emph{irréductible})~; pour que $X$ soit \emph{intègre}, il faut et il suffit
que le $\mathcal O_X$-Module $\mathcal O_X$ soit intègre ((I, 2.1.8) et
(3.2.1)). Si $\mathcal O_X$ est intègre en un point $x$, c'est-à-dire si
l'anneau $\mathcal O_x$ est intègre, on dit que $X$ est \emph{intègre au point
$x$}. On dit qu'un sous-préschéma fermé $Y$ de $X$ est \emph{primaire dans
$X$} si l'Idéal $\mathcal J$ de $\mathcal O_X$ qui définit $Y$ est primaire
dans $\mathcal O_X$.

\begin{definition}[3.2.5]
\label{IV.2.3.2.5-fr}
Soient $X$ un préschéma localement noethérien, $\mathcal F$ un
$\mathcal O_X$-Module cohérent. On appelle \emph{décomposition irredondante de
$\mathcal F$} une famille $(\mathcal F_\alpha)_{\alpha\in I}$ de
$\mathcal O_X$-Modules quotients de $\mathcal F$ telle que les
$\mathcal F_\alpha$ soient irredondants, que la famille
$(\operatorname{Supp}(\mathcal F_\alpha))$ soit localement finie, et que
l'homomorphisme canonique
$\mathcal F\to\bigoplus_{\alpha\in I}\mathcal F_\alpha$ soit injectif. On dit
qu'une telle décomposition est \emph{réduite} si les ensembles
$\operatorname{Ass}(\mathcal F_\alpha)$ sont deux à deux distincts, et s'il
n'existe aucune partie $J\ne I$ telle que la sous-famille
$(\mathcal F_\alpha)_{\alpha\in J}$ soit une décomposition irredondante de
$\mathcal F$.
\end{definition}

Si $(\mathcal F_\alpha)_{\alpha\in I}$ est une décomposition irredondante
(resp. irredondante réduite) de $\mathcal F$ et si l'on pose
$\mathcal F_\alpha=\mathcal F/\mathcal G_\alpha$, on dit encore que la
famille $(\mathcal G_\alpha)_{\alpha\in I}$ de sous-$\mathcal O_X$-Modules de
$\mathcal F$ est une \emph{décomposition primaire de $0$ dans $\mathcal F$}~;
on observera que l'hypothèse d'injectivité de l'homomorphisme canonique
$\mathcal F\to\bigoplus_{\alpha\in I}\mathcal F_\alpha$ équivaut à la
condition $\bigcap_{\alpha\in I}\mathcal G_\alpha=0$.

Si $(\mathcal F_\alpha)$ est une décomposition irredondante de $\mathcal F$,
dire qu'elle est \emph{réduite} équivaut à dire que les
$\operatorname{Ass}(\mathcal F_\alpha)$ sont deux à deux distincts et contenus
dans $\operatorname{Ass}(\mathcal F)$~; si
$\operatorname{Ass}(\mathcal F_\alpha)=\{x_\alpha\}$ pour tout $\alpha\in I$,
$\alpha\mapsto x_\alpha$ est une \emph{bijection} de $I$ sur
$\operatorname{Ass}(\mathcal F)$~: ces propriétés sont en effet locales et
résultent donc de Bourbaki, \emph{Alg. comm.}, chap.~IV, \S~2,
n\textsuperscript{o}~3, prop.~4.

\begin{proposition}[3.2.6]
\label{IV.2.3.2.6-fr}
Soient $X$ un préschéma localement noethérien, $\mathcal F$ un
$\mathcal O_X$-Module cohérent. Alors il existe une décomposition irredondante
réduite $(\mathcal F^{(x)})_{x\in\operatorname{Ass}(\mathcal F)}$, formée de
$\mathcal O_X$-Modules cohérents tels que pour tout
$x\in\operatorname{Ass}(\mathcal F)$, on ait
$\operatorname{Ass}(\mathcal F^{(x)})=\{x\}$. Pour tout
$x\in\operatorname{Ass}(\mathcal F)$ tel que $\overline{\{x\}}$ ne soit pas
immergé, $\mathcal F^{(x)}$ est déterminé de façon unique comme image de
l'homomorphisme canonique $\mathcal F\to i_*(i^*(\mathcal F))$, où $i$ est le
morphisme canonique $\operatorname{Spec}(k(x))\to X$.
\end{proposition}

Pour tout $x\in\operatorname{Ass}(\mathcal F)$, soit $U$ un voisinage ouvert
affine de $x$, d'anneau $A$, et soit $\mathcal F|U=\widetilde M$, où $M$ est
un $A$-module de type fini. On sait (Bourbaki, \emph{Alg. comm.}, chap.~IV,
\S~1, n\textsuperscript{o}~1, prop.~4) qu'il existe un sous-module $N$ de $M$
tel que, si l'on pose $P=M/N$, on ait
$\operatorname{Ass}(P)=\{x\}$ et
$\operatorname{Ass}(N)=\operatorname{Ass}(M)-\{x\}$. Soit
$\mathcal G=\widetilde P$, qui est un $\mathcal O_U$-Module quasi-cohérent, et
soit $j$ l'injection canonique $U\to X$~; soit
$u:j^*(\mathcal F)\to\mathcal G$ l'homomorphisme surjectif correspondant à
l'homomorphisme $M\to P$~; on en déduit un homomorphisme
$j_*(u):j_*(j^*(\mathcal F))\to j_*(\mathcal G)$, d'où par composition un
homomorphisme
\[
v:\mathcal F\xrightarrow{\rho_{\mathcal F}}j_*(j^*(\mathcal F))
\xrightarrow{j_*(u)}j_*(\mathcal G)
\]
dont $u$ est la restriction à $U$~; nous désignerons par
$\mathcal F^{(x)}$ l'image de $\mathcal F$ par cet homomorphisme, qui est un
$\mathcal O_X$-Module cohérent (I, 6.1.1). On a
$\operatorname{Ass}(j_*(\mathcal G))=\{x\}$ en vertu de (3.1.13), et
\emph{a fortiori} (3.1.7)
$\operatorname{Ass}(\mathcal F^{(x)})=\{x\}$, puisque
$\mathcal F^{(x)}\ne0$. En outre, si

\oldpage[IV-2]{43}
$\mathcal N^{(x)}=\operatorname{Ker}(v)$, on a
$\mathcal N^{(x)}|U=\widetilde N$, donc
$x\notin\operatorname{Ass}(\mathcal N^{(x)})$. Il en résulte que
l'homomorphisme
$\mathcal F\to\bigoplus_{x\in\operatorname{Ass}(\mathcal F)}
\mathcal F^{(x)}$ est injectif, car son noyau $\mathcal H$ est contenu dans
tous les $\mathcal N^{(x)}$, donc $\operatorname{Ass}(\mathcal H)$ est contenu
dans l'intersection des $\operatorname{Ass}(\mathcal N^{(x)})$, qui est vide~;
par suite (3.1.5), $\mathcal H=0$. Compte tenu de ce que
$\operatorname{Ass}(\mathcal F)$ est localement fini (3.1.6), il est clair que
$(\mathcal F^{(x)})_{x\in\operatorname{Ass}(\mathcal F)}$ est une
décomposition irredondante réduite de $\mathcal F$ vérifiant les conditions de
l'énoncé. La caractérisation de $\mathcal F^{(x)}$ lorsque
$\overline{\{x\}}$ n'est pas immergé résulte de Bourbaki,
\emph{Alg. comm.}, chap.~IV, \S~2, n\textsuperscript{o}~3, prop.~5, la question
étant locale, et compte tenu de (I, 1.6.7).

\begin{corollary}[3.2.7]
\label{IV.2.3.2.7-fr}
Sous les hypothèses de (3.2.6), si $\mathcal F$ n'a pas de cycle premier
associé immergé, il n'existe qu'une seule décomposition irredondante réduite de
$\mathcal F$.
\end{corollary}

\begin{corollary}[3.2.8]
\label{IV.2.3.2.8-fr}
Soient $X$ un préschéma noethérien, $\mathcal F$ un $\mathcal O_X$-Module
cohérent. Il existe une filtration finie $(\mathcal F_i)_{0\le i\le n}$ de
$\mathcal F$ telle que $\mathcal F_0=\mathcal F$, $\mathcal F_n=0$, formée de
$\mathcal O_X$-Modules cohérents et tels que les quotients
$\mathcal F_i/\mathcal F_{i+1}$ soient nuls ou irredondants, et
$\operatorname{Ass}(\mathcal F_i/\mathcal F_{i+1})
\subset\operatorname{Ass}(\mathcal F)$.
\end{corollary}

En effet, $\mathcal F$ est isomorphe à un sous-$\mathcal O_X$-Module d'une
somme directe finie $\bigoplus_{j=1}^n\mathcal G_j$, où les $\mathcal G_j$
sont irredondants et cohérents (3.2.6)~; comme tout sous-$\mathcal O_X$-Module
quasi-cohérent de $\mathcal G_j$ est nul ou irredondant (3.1.7), les
$\mathcal F_i=\mathcal F\cap(\bigoplus_{j=1}^{n-i}\mathcal G_j)$ répondent à
la question, $\mathcal F_i/\mathcal F_{i+1}$ étant isomorphe à un
sous-$\mathcal O_X$-Module cohérent de $\mathcal G_{n-i}$.

\subsection{Relations avec la platitude}

\begin{proposition}[3.3.1]
\EGAReaderTarget{ega4.c04.prop.3-3-1}{IV.2.3.3.1-fr}
Soient $f:X\to Y$ un morphisme, $\mathcal F$ un $\mathcal O_X$-Module
quasi-cohérent et $f$-plat, $\mathcal E$ un $\mathcal O_Y$-Module
quasi-cohérent. Si, pour tout $y\in Y$, on pose
$\mathcal F_y=\mathcal F\otimes_{\mathcal O_Y}k(y)$, on a
\end{proposition}

\begin{env}[3.3.1.1]
\label{IV.2.3.3.1.1-fr}
\[
\operatorname{Ass}(\mathcal E\otimes_Y\mathcal F)
\supset\bigcup_{y\in\operatorname{Ass}(\mathcal E)}
\operatorname{Ass}(\mathcal F_y),
\]
\end{env}

et les deux membres sont égaux si $Y$ est localement noethérien.

(Bien entendu, $\mathcal F_y$ est un faisceau sur la fibre $f^{-1}(y)$, et on
identifie cette fibre à un sous-espace de $X$ (I, 3.6.1).) La question étant
locale sur $X$ et sur $Y$, on est ramené au cas où $X$ et $Y$ sont affines,
et la proposition est alors démontrée dans Bourbaki, \emph{Alg. comm.},
chap.~IV, \S~2, n\textsuperscript{o}~6, th.~2.

\begin{corollary}[3.3.2]
\label{IV.2.3.3.2-fr}
Soient $Y$ un préschéma localement noethérien sans cycles premiers associés
immergés, $f:X\to Y$ un morphisme, $\mathcal F$ un $\mathcal O_X$-Module
quasi-cohérent et $f$-plat. Alors, pour tout
$x\in\operatorname{Ass}(\mathcal F)$, $f(x)$ est point maximal de $Y$.
\end{corollary}

Il suffit d'appliquer (3.3.1) avec $\mathcal E=\mathcal O_Y$, puisque
$\operatorname{Ass}(\mathcal O_Y)$ est par hypothèse l'ensemble des points
maximaux de $Y$.

\begin{corollary}[3.3.3]
\label{IV.2.3.3.3-fr}
Sous les hypothèses de (3.3.1), supposons en outre que $X$ et $Y$ soient
localement noethériens, $\mathcal E$ et $\mathcal F$ cohérents. Alors les
conditions suivantes sont équivalentes~:
\begin{enumerate}
\item[\emph{a)}] $\mathcal E\otimes_Y\mathcal F$ est sans cycle premier
associé immergé.
\item[\emph{b)}] Pour tout point
$y\in\operatorname{Ass}(\mathcal E)\cap f(\operatorname{Supp}(\mathcal F))$,
$\overline{\{y\}}$ est un cycle premier associé non immergé de $\mathcal E$ et
$\mathcal F_y$ est sans cycle premier associé immergé.
\end{enumerate}
\end{corollary}

\oldpage[IV-2]{44}
Supposons \emph{a)} vérifiée. Les hypothèses entraînent que
$\mathcal E\otimes_Y\mathcal F$ est un $\mathcal O_X$-Module cohérent
$(0_{\mathrm I},5.3.11\text{ et }5.3.5)$~; ses cycles premiers associés sont
donc les composantes irréductibles de
$\operatorname{Supp}(\mathcal E\otimes_Y\mathcal F)$ (3.1.4) et pour tout
point maximal $x$ de $\operatorname{Supp}(\mathcal E\otimes_Y\mathcal F)$,
$f(x)=y$ est point maximal de $\operatorname{Supp}(\mathcal E)$ et $x$ point
maximal de $\operatorname{Supp}(\mathcal F_y)$ (2.5.5). Comme, en vertu de
(3.3.1) et du fait que la relation
$y\in f(\operatorname{Supp}(\mathcal F))$ entraîne $\mathcal F_y\ne0$
(I, 9.1.13), tout point de
$\operatorname{Ass}(\mathcal E)\cap f(\operatorname{Supp}(\mathcal F))$ est
l'image par $f$ d'un point maximal de
$\operatorname{Supp}(\mathcal E\otimes_Y\mathcal F)$, on voit que la condition
\emph{b)} est vérifiée.

Inversement, supposons \emph{b)} vérifiée, et montrons que si $z$, $z'$ sont
deux points distincts de $\operatorname{Ass}(\mathcal E\otimes_Y\mathcal F)$,
aucun d'eux ne peut être adhérent à l'autre. En premier lieu, si
$f(z)=f(z')=y$, on a $y\in\operatorname{Ass}(\mathcal E)$ et $z$ et $z'$
appartiennent à $\operatorname{Ass}(\mathcal F_y)$ par (3.3.1), d'où
$y\in f(\operatorname{Supp}(\mathcal F))$~; comme par hypothèse aucun des deux
points $z$, $z'$ n'est adhérent à l'autre \emph{dans} $f^{-1}(y)$, aucun d'eux
ne peut être adhérent à l'autre dans $X$. Si $y=f(z)$ et $y'=f(z')$ sont
distincts, ils appartiennent à
$\operatorname{Ass}(\mathcal E)\cap f(\operatorname{Supp}(\mathcal F))$, donc
aucun d'eux ne peut être adhérent à l'autre dans $Y$~; il résulte de la
continuité de $f$ qu'aucun des points $z$, $z'$ ne peut être adhérent à
l'autre dans $X$.

\begin{proposition}[3.3.4]
\label{IV.2.3.3.4-fr}
Soient $X$, $Y$ deux préschémas localement noethériens, $f:X\to Y$ un
morphisme, $\mathcal E$ un $\mathcal O_Y$-Module cohérent, $\mathcal F$ un
$\mathcal O_X$-Module cohérent et $f$-plat. Alors les conditions suivantes
sont équivalentes~:
\begin{enumerate}
\item[\emph{a)}] $\mathcal E\otimes_Y\mathcal F$ est réduit (3.2.2).
\item[\emph{b)}] Pour tout point
$y\in\operatorname{Ass}(\mathcal E)\cap f(\operatorname{Supp}(\mathcal F))$,
$\overline{\{y\}}$ est un cycle premier associé non immergé de $\mathcal E$,
$\operatorname{long}(\mathcal E_y)=1$ et $\mathcal F_y$ est réduit.
\end{enumerate}
\end{proposition}

Supposons \emph{a)} vérifiée. On sait déjà (3.3.3) que pour tout
$y\in\operatorname{Ass}(\mathcal E)\cap f(\operatorname{Supp}(\mathcal F))$,
$\overline{\{y\}}$ est un cycle premier associé non immergé de $\mathcal E$ et
$\mathcal F_y$ est sans cycle premier associé immergé. En outre (2.5.5), pour
tout $x\in\operatorname{Ass}(\mathcal E\otimes_Y\mathcal F)\cap f^{-1}(y)$,
on a
\[
1=\operatorname{long}((\mathcal E\otimes_Y\mathcal F)_x)
=\operatorname{long}(\mathcal E_y)\cdot\operatorname{long}((\mathcal F_y)_x),
\]
donc $\operatorname{long}(\mathcal E_y)
=\operatorname{long}((\mathcal F_y)_x)=1$, ce qui prouve \emph{b)}.

Inversement, supposons \emph{b)} vérifiée~; on sait déjà que tout point
$x\in\operatorname{Ass}(\mathcal E\otimes_Y\mathcal F)$ est point maximal de
$\operatorname{Supp}(\mathcal E\otimes_Y\mathcal F)$, que $y=f(x)$ est point
maximal de $\operatorname{Supp}(\mathcal E)$ et $x$ point maximal de
$\operatorname{Supp}(\mathcal F_y)$ (3.3.1 et 3.3.3)~; en outre il résulte de
l'hypothèse et de (2.5.5) que
$\operatorname{long}((\mathcal E\otimes_Y\mathcal F)_x)=1$, ce qui prouve
\emph{a)}.

\begin{corollary}[3.3.5]
\label{IV.2.3.3.5-fr}
Soient $X$, $Y$ deux préschémas localement noethériens, $f:X\to Y$ un
morphisme plat~; si $Y$ est réduit aux points de $f(X)$ et si $f^{-1}(y)$ est
un $k(y)$-préschéma réduit pour tout $y\in f(X)$, alors $X$ est réduit.
\end{corollary}

Comme le Nilradical $\mathcal N_Y$ est cohérent, l'ensemble des points où $Y$
est réduit est \emph{ouvert} $(0_{\mathrm I},5.2.2)$, et on peut se borner au
cas où $Y$ est réduit. Il suffit alors d'appliquer (3.3.4) à
$\mathcal E=\mathcal O_Y$ et $\mathcal F=\mathcal O_X$.

\begin{proposition}[3.3.6]
\label{IV.2.3.3.6-fr}
Soient $f:X\to S$, $g:Y\to S$ deux morphismes, $\mathcal F$ un
$\mathcal O_X$-Module quasi-cohérent, $\mathcal G$ un $\mathcal O_Y$-Module
quasi-cohérent. On suppose que~: 1\textsuperscript{o} $\mathcal G$ est
$g$-plat~; 2\textsuperscript{o} $X$ est localement noethérien, et pour tout
$s\in S$, $g^{-1}(s)$ est localement noethérien (ce qui aura lieu si $Y$ est
aussi localement noethérien). Soit $Z=X\times_S Y$~; pour tout couple $(x,y)$
tel que $x\in X$, $y\in Y$ et

\oldpage[IV-2]{45}
$f(x)=g(y)=s$, soit $T_{x,y}$ le préschéma
$\operatorname{Spec}(k(x)\otimes_{k(s)}k(y))$, et soit $I_{x,y}$ l'image de
$\operatorname{Ass}(\mathcal O_{T_{x,y}})$ par le monomorphisme canonique
$T_{x,y}\to Z$ (I, 3.4.9). On a alors
\end{proposition}

\begin{env}[3.3.6.1]
\label{IV.2.3.3.6.1-fr}
\[
\operatorname{Ass}(\mathcal F\otimes_S\mathcal G)
=\bigcup_{x\in\operatorname{Ass}(\mathcal F)}
\left(\bigcup_{y\in\operatorname{Ass}(\mathcal G_{f(x)})}I_{x,y}\right),
\]
\end{env}

où pour tout $s\in S$, $\mathcal G_s=\mathcal G\otimes_{\mathcal O_S}k(s)$.

Soient $p:Z\to X$, $q:Z\to Y$ les projections canoniques, de sorte que l'on
a le diagramme commutatif
\[
\xymatrix{
X \ar[d]_f & Z \ar[l]_p \ar[d]^q \\
S & Y \ar[l]_g
}
\]
Posons $\mathcal G'=q^*(\mathcal G)$, de sorte que
$\mathcal F\otimes_S\mathcal G=\mathcal F\otimes_X\mathcal G'$~; comme
$\mathcal G'$ est $p$-plat (2.1.4), il résulte de (3.3.1) que l'on a

\begin{env}[3.3.6.2]
\label{IV.2.3.3.6.2-fr}
\[
\operatorname{Ass}(\mathcal F\otimes_S\mathcal G)
=\bigcup_{x\in\operatorname{Ass}(\mathcal F)}
\operatorname{Ass}(\mathcal G'_x)
\]
\end{env}

avec $\mathcal G'_x=\mathcal G'\otimes_{\mathcal O_X}k(x)$. Si $s=f(x)$, on a
$\mathcal G'_x=\mathcal G_s\otimes_{k(s)}k(x)$ et
$p^{-1}(x)=g^{-1}(s)\otimes_{k(s)}k(x)$~; en outre, comme le corps $k(x)$ est
un $k(s)$-module plat, le morphisme
$p^{-1}(x)\to g^{-1}(s)$ est plat (2.1.4)~; appliquant (3.3.1) à ce morphisme,
il vient

\begin{env}[3.3.6.3]
\label{IV.2.3.3.6.3-fr}
\[
\operatorname{Ass}(\mathcal G'_x)
=\bigcup_{y\in\operatorname{Ass}(\mathcal G_s)}
\operatorname{Ass}(\mathcal O_{T_{x,y}}),
\]
\end{env}

d'où la proposition.

On notera que si, dans l'énoncé, on supprime l'hypothèse
2\textsuperscript{o} on peut encore conclure, en vertu de (3.3.1), la relation

\begin{env}[3.3.6.4]
\label{IV.2.3.3.6.4-fr}
\[
\operatorname{Ass}(\mathcal F\otimes_S\mathcal G)
\supset\bigcup_{x\in\operatorname{Ass}(\mathcal F)}
\left(\bigcup_{y\in\operatorname{Ass}(\mathcal G_{f(x)})}I_{x,y}\right).
\]
\end{env}

\begin{corollary}[3.3.7]
\label{IV.2.3.3.7-fr}
Sous les hypothèses de (3.3.6), supposons en outre que $S$ soit localement
noethérien et que
$f(\operatorname{Ass}(\mathcal F))\subset\operatorname{Ass}(\mathcal O_S)$.
Alors on a
\end{corollary}

\begin{env}[3.3.7.1]
\label{IV.2.3.3.7.1-fr}
\[
\operatorname{Ass}(\mathcal F\otimes_S\mathcal G)
=\bigcup_{(x,y)\in C}I_{x,y},
\]
\end{env}

où $C$ est l'ensemble des couples $(x,y)$ tels que
$x\in\operatorname{Ass}(\mathcal F)$, $y\in\operatorname{Ass}(\mathcal G)$ et
$f(x)=g(y)$.

Comme $\mathcal G$ est $g$-plat, il résulte en effet de (3.3.1) que la relation
«~$s\in\operatorname{Ass}(\mathcal O_S)$ et
$y\in\operatorname{Ass}(\mathcal G_s)$~» équivaut à
$y\in\operatorname{Ass}(\mathcal G)$~; la conclusion résulte de (3.3.6.1).

\begin{remark}[3.3.8]
\label{IV.2.3.3.8-fr}
Nous verrons plus loin (4.2.2) que sous les hypothèses de (3.3.6),
$T_{x,y}$ est un préschéma \emph{sans cycle premier associé immergé}~; il en
résultera que si $\mathcal F$ et les $\mathcal G_s$ sont sans cycle premier
associé immergé, il en est de même de $\mathcal F\otimes_S\mathcal G$.
\end{remark}

\begin{corollary}[3.3.9]
\label{IV.2.3.3.9-fr}
Sous les conditions de (3.3.7), on a
\end{corollary}

\begin{env}[3.3.9.1]
\label{IV.2.3.3.9.1-fr}
\[
q(\operatorname{Ass}(\mathcal F\otimes_S\mathcal G))
\subset\operatorname{Ass}(\mathcal G)
\]
\end{env}

(où $q:X\times_S Y\to Y$ est la projection canonique).

En effet, si $(x,y)\in Z$, on a
$q(I_{x,y})=\{y\}\subset\operatorname{Ass}(\mathcal G)$.

\oldpage[IV-2]{46}

\subsection{\texorpdfstring{Propriétés des faisceaux $\mathcal F/t\mathcal F$}{Propriétés des faisceaux F/tF}}

\begin{proposition}[3.4.1]
\label{IV.2.3.4.1-fr}
Soient $X$ un préschéma localement noethérien, $t$ une section de
$\mathcal O_X$ au-dessus de $X$, $Y$ le sous-préschéma fermé de $X$ défini par
l'Idéal $t\mathcal O_X$ de $\mathcal O_X$. Soient $\mathcal F$ un
$\mathcal O_X$-Module cohérent, $S$ le sous-préschéma fermé réduit de $X$ ayant
pour espace sous-jacent $\operatorname{Supp}(\mathcal F)$, $(S_i)$ la famille
des sous-préschémas fermés réduits de $X$ ayant pour espaces sous-jacents les
composantes irréductibles de $S$~; on désigne par $s_i$ le point générique de
$S_i$. Enfin, soient $Z$ une composante irréductible de
$\operatorname{Supp}(\mathcal F/t\mathcal F)=S\cap Y$, $z$ son point
générique.
\begin{enumerate}
\item[(i)] Pour tout $i$ tel que $Z\subset S_i$, $Z$ est une composante
irréductible de $S_i\cap Y$.
\item[(ii)] Si $Z$ n'est pas égal à un des $S_i$, on a
\end{enumerate}
\end{proposition}

\begin{env}[3.4.1.1]
\label{IV.2.3.4.1.1-fr}
\[
\operatorname{long}((\mathcal F/t\mathcal F)_z)
\geqslant\sum_i\operatorname{long}(\mathcal F_{s_i})
\]
\end{env}

où la somme du second membre est étendue à tous les $i$ tels que
$Z\subset S_i$.

\begin{enumerate}
\item[(iii)] Supposons que $Z$ ne soit égal à aucun des $S_i$. Pour que les
deux membres de (3.4.1.1) soient égaux, il faut et il suffit que les deux
conditions suivantes soient satisfaites~:
\begin{enumerate}
\item[$\alpha$)] $t_z$ est $\mathcal F_z$-régulier $(0,15.1.4)$.
\item[$\beta$)] Pour tout $i$ tel que $Z\subset S_i$, l'image canonique du
germe $t_z$ dans $\mathcal O_{S_i,z}$ engendre l'idéal maximal de cet anneau
(ce qui entraîne que $\mathcal O_{S_i,z}$ est un anneau de valuation discrète
et l'image de $t_z$ une uniformisante de cet anneau).
\end{enumerate}
\end{enumerate}

\begin{proof}
(i) Si $j:Y\to X$ est l'injection canonique, on a
\[
\mathcal F/t\mathcal F
=\mathcal F\otimes_{\mathcal O_X}\mathcal O_Y
=j^*(\mathcal F),
\]
donc
\[
\operatorname{Supp}(\mathcal F/t\mathcal F)
=j^{-1}(S)=S\cap Y
\quad\text{(I, 9.1.13),}
\]
d'où l'assertion.

(ii) et (iii). Comme les $s_i$ tels que $Z\subset S_i$ sont ceux qui
appartiennent à $\operatorname{Spec}(\mathcal O_z)$, on peut, pour démontrer
(ii) et (iii), remplacer $X$ par $\operatorname{Spec}(\mathcal O_z)$~; et si
$M=\mathcal F_z$, on peut donc supposer que $\mathcal F=\widetilde M$, d'où
$\mathcal F_{s_i}=M_{\mathfrak p_i}$, en désignant par $\mathfrak p_i$ les
idéaux minimaux de $\mathcal O_z$. D'ailleurs, comme $M$ est un
$\mathcal O_z$-module de type fini, on a $S=S'_{\mathrm{red}}$, avec
$S'=\operatorname{Spec}(\mathcal O_z/\mathfrak a)$, où $\mathfrak a$ est
l'annulateur de $M$ dans $\mathcal O_z$ $(0_{\mathrm I},1.7.4)$, et les deux
membres de (3.4.1.1) gardent les mêmes valeurs, que l'on considère $M$ comme
$\mathcal O_z$-module ou comme un $(\mathcal O_z/\mathfrak a)$-module~; on voit
donc qu'on peut finalement remplacer $X$ par $S'=\operatorname{Spec}(A)$, où
$A$ est un anneau local noethérien, $M$ étant un $A$-module fidèle~; comme $Z$
est fermé dans $S$, l'hypothèse que $Z\neq S_i$ pour tous les $i$ signifie que
$s_i\notin Z$, donc que $\dim(A)>0$~; enfin, dire que $z$ est point générique
de $Z$, composante irréductible de $S_i\cap V(t)$, signifie que $A/tA$ est de
dimension $0$ (autrement dit est un anneau local artinien). On est donc ramené
à prouver l'énoncé suivant~:
\end{proof}

\begin{lemma}[3.4.1.2]
\label{IV.2.3.4.1.2-fr}
Soient $A$ un anneau local noethérien de dimension $>0$, $\mathfrak p_i$ les
idéaux premiers minimaux de $A$, $\mathfrak m$ son idéal maximal, $t$ un
élément de $\mathfrak m$ tel que $A/tA$ soit artinien. Alors, pour tout
$A$-module de type fini $M$, on a
\end{lemma}

\begin{env}[3.4.1.3]
\label{IV.2.3.4.1.3-fr}
\[
\operatorname{long}(M/tM)
\geqslant\sum_i\operatorname{long}(M_{\mathfrak p_i})~;
\]
\end{env}

en outre, pour que les deux membres de (3.4.1.3) soient égaux, il faut et il
suffit que les conditions suivantes soient satisfaites~:

\oldpage[IV-2]{47}

\begin{enumerate}
\item[$\alpha$)] $t$ est $M$-régulier~;
\item[$\beta$)] pour tout $i$ tel que $M_{\mathfrak p_i}\neq0$, l'image de $t$
dans $A/\mathfrak p_i$ engendre l'idéal maximal de cet anneau (ce qui entraîne
que $A/\mathfrak p_i$ est un anneau de valuation discrète).
\end{enumerate}

\begin{proof}
Comme $A$ n'est pas de dimension $0$ et que $A/tA$ est artinien, on a
nécessairement $\dim(A)=1$ $(0,16.3.4)$ et $t\notin\mathfrak p_i$ pour tout
$i$~: l'idéal principal $(t)$ est donc un idéal de définition de $A$, et
contient par suite une puissance de son idéal maximal $\mathfrak m$. Soit $N$
le sous-module des éléments de $M$ annulés par une puissance de $t$ (ou par
une puissance de $\mathfrak m$, ce qui revient au même comme on vient de le
voir)~; si l'on pose $P=M/N$, $t$ est $P$-régulier, car la relation $tx\in N$
pour un $x\in M$ entraîne $t^k(tx)=0$ pour un entier $k$, donc $x\in N$. Cela
étant, on a le

\begin{lemma}[3.4.1.4]
\label{IV.2.3.4.1.4-fr}
Soient $A$ un anneau,
\[
0\longrightarrow M'\longrightarrow M\longrightarrow M''\longrightarrow0
\]
une suite exacte de $A$-modules. Si $t\in A$ est $M''$-régulier, la suite
\[
0\longrightarrow M'/tM'\longrightarrow M/tM
\longrightarrow M''/tM''\longrightarrow0
\]
est exacte.
\end{lemma}

\begin{proof}
Comme $M/tM=M\otimes_A(A/tA)$, il suffit de prouver l'exactitude en $M'/tM'$~;
or, si l'image $x\in M$ d'un élément $x'$ de $M'$ est telle que $x=ty$ avec
$y\in M$, on en déduit, pour les images $x''$, $y''$ de $x$, $y$ dans $M''$,
$x''=ty''$~; mais comme $x''=0$, l'hypothèse entraîne $y''=0$, donc $y$ est
l'image d'un élément $y'\in M'$, et la relation $x=ty$ entraîne $x'=ty'$
puisque $M'\to M$ est injectif.
\end{proof}

Ce lemme étant établi, on en tire la relation
\end{proof}

\begin{env}[3.4.1.5]
\label{IV.2.3.4.1.5-fr}
\[
\operatorname{long}(M/tM)
=\operatorname{long}(N/tN)+\operatorname{long}(P/tP).
\]
\end{env}

D'autre part, pour tout $i$, on a $N_{\mathfrak p_i}=0$ puisque
$t\notin\mathfrak p_i$~; donc $M_{\mathfrak p_i}=P_{\mathfrak p_i}$~; pour
prouver (3.4.1.3), il suffit donc de le faire en y remplaçant $M$ par $P$~;
d'autre part, si les deux membres de (3.4.1.3) sont égaux, il résulte de la
même inégalité pour $P$ et de (3.4.1.5) que l'on a nécessairement
$\operatorname{long}(N/tN)=0$, donc $N/tN=0$ et finalement $N=0$, par le lemme
de Nakayama, $N$ étant de type fini~; or, $N=0$ signifie que $t$ est
$M$-régulier. On peut donc se ramener au cas où $M=P$, c'est-à-dire supposer
déjà que $t$ est $M$-régulier. Notons que cela entraîne que
$\mathfrak m\notin\operatorname{Ass}(M)$, puisque $t$ ne peut annuler un
élément $\neq0$ de $M$. Comme $A$ est de dimension $1$, on a donc
nécessairement
\[
\operatorname{Ass}(M)\subset\bigcup_i\{\mathfrak p_i\}.
\]

Procédons alors par récurrence sur
$n=\sum_i\operatorname{long}(M_{\mathfrak p_i})$. Si $n=0$, on a
nécessairement $M_{\mathfrak p_i}=0$ pour tout $i$, donc $M=0$ puisque aucun
des $\mathfrak p_i$ n'appartient à $\operatorname{Ass}(M)$~; les deux membres
de (3.4.1.3) sont alors nuls, et l'assertion $\beta$) de (3.4.1.2) est
triviale. Si $n>0$, le raisonnement du début de la démonstration de (3.4.1)
permet de supposer en outre que le $A$-module $M$ est fidèle~; cela entraîne
$M_{\mathfrak p_i}\neq0$ pour tout $i$ (Bourbaki, \emph{Alg. comm.}, chap.~II,
\S~2, n\textsuperscript{o}~2, cor.~2 de la prop.~4), et par suite
\[
\operatorname{Ass}(M)=\bigcup_i\{\mathfrak p_i\}.
\]

\oldpage[IV-2]{48}

Supposons d'abord $n=1$~; il n'y a alors qu'un seul idéal premier minimal
$\mathfrak p$ de $A$, et dire que $M_{\mathfrak p}$ est de longueur $1$
signifie que $M_{\mathfrak p}$ est isomorphe au corps résiduel
$k=A_{\mathfrak p}/\mathfrak pA_{\mathfrak p}$ en tant que
$A_{\mathfrak p}$-module. Par suite $M_{\mathfrak p}$ est annulé par
$\mathfrak pA_{\mathfrak p}$, donc $\mathfrak p$ est l'annulateur de $M$
(Bourbaki, \emph{Alg. comm.}, chap.~II, \S~2, n\textsuperscript{o}~4,
formule~(9)), ce qui entraîne $\mathfrak p=0$ puisque $M$ est supposé fidèle~;
l'anneau $A$ est donc intègre. Cela étant, l'hypothèse $M\neq0$ entraîne
$M/tM\neq0$ par le lemme de Nakayama, et par suite
$\operatorname{long}(M/tM)\geqslant1$, ce qui démontre (3.4.1.3) dans ce cas.
En outre, si $\operatorname{long}(M/tM)=1$, $M$ est nécessairement monogène
(Bourbaki, \emph{Alg. comm.}, chap.~II, \S~3, n\textsuperscript{o}~2, cor.~2
de la prop.~4) donc isomorphe à un quotient $A/\mathfrak b$~; comme il est
fidèle, on a nécessairement $\mathfrak b=0$ et $M$ est isomorphe à $A$~;
comme $\operatorname{long}(A/tA)=1$, $tA$ est nécessairement égal à l'idéal
maximal $\mathfrak m$, et comme $A$ est un anneau local intègre noethérien,
cela prouve que $A$ est un anneau de valuation discrète (Bourbaki,
\emph{Alg. comm.}, chap.~VI, \S~3, n\textsuperscript{o}~6, prop.~9), dont $t$
est l'uniformisante. Réciproquement, si $A$ est un anneau de valuation
discrète, $t$ son uniformisante, $\operatorname{long}(M_{\mathfrak p})=1$ et
si $t$ est $M$-régulier, $M$ est sans torsion, donc isomorphe à un sous-module
de $A$ ($M$ étant de type fini), et par suite isomorphe à $A$ lui-même, d'où
$\operatorname{long}(M/tM)=\operatorname{long}(A/tA)=1$.

Supposons maintenant $n\geqslant2$~; il existe alors une suite exacte
\[
0\longrightarrow M'\longrightarrow M\longrightarrow M''\longrightarrow0
\]
avec $M'\neq0$, $M''\neq0$ et
$\operatorname{Ass}(M)=\operatorname{Ass}(M')\cup\operatorname{Ass}(M'')$~;
en effet, si $\operatorname{Ass}(M)$ n'est pas réduit à un seul élément, cela
résulte de Bourbaki, \emph{Alg. comm.}, chap.~IV, \S~1,
n\textsuperscript{o}~1, prop.~4~; si au contraire $\operatorname{Ass}(M)$ est
réduit à un seul idéal premier, ce dernier est nécessairement l'unique idéal
premier minimal $\mathfrak p$ de $A$~; l'hypothèse entraîne alors
$\operatorname{long}(M_{\mathfrak p})\geqslant2$ et il suffit de prendre pour
$M'$ l'image réciproque d'un sous-module de $M_{\mathfrak p}$ distinct de $0$
et de $M_{\mathfrak p}$. Comme $t$ est $M$-régulier, $t$ n'appartient à aucun
des idéaux premiers de $\operatorname{Ass}(M)$ (Bourbaki,
\emph{Alg. comm.}, chap.~IV, \S~1, n\textsuperscript{o}~1, cor.~2 de la
prop.~2), donc, pour la même raison, $t$ est $M'$-régulier et
$M''$-régulier. Cette dernière propriété entraîne par (3.4.1.4) que la suite
\[
0\longrightarrow M'/tM'\longrightarrow M/tM
\longrightarrow M''/tM''\longrightarrow0
\]
est exacte~; comme par ailleurs il en est de même de la suite
\[
0\longrightarrow M'_{\mathfrak p_i}\longrightarrow M_{\mathfrak p_i}
\longrightarrow M''_{\mathfrak p_i}\longrightarrow0
\]
pour tout $i$, on a
\[
\begin{aligned}
\operatorname{long}(M/tM)
&=\operatorname{long}(M'/tM')+\operatorname{long}(M''/tM''),\\
\operatorname{long}(M_{\mathfrak p_i})
&=\operatorname{long}(M'_{\mathfrak p_i})
 +\operatorname{long}(M''_{\mathfrak p_i})
\end{aligned}
\]
et l'hypothèse de récurrence entraîne donc l'inégalité (3.4.1.3). En outre,
les deux membres ne peuvent être égaux que si les inégalités analogues pour
$M'$ et $M''$ sont aussi des égalités. En vertu de l'hypothèse de récurrence,
cela équivaut à la propriété $\beta$) pour les $\mathfrak p_i$ tels que
$M'_{\mathfrak p_i}\neq0$ ou $M''_{\mathfrak p_i}\neq0$~; mais ces idéaux sont
précisément ceux pour lesquels $M_{\mathfrak p_i}\neq0$. C.Q.F.D.

\oldpage[IV-2]{49}

\begin{corollary}[3.4.2]
\label{IV.2.3.4.2-fr}
Sous les hypothèses générales de (3.4.1), supposons que $Z$ ne soit pas égal
à un des $S_i$ et que
$\operatorname{long}((\mathcal F/t\mathcal F)_z)=1$. Alors il existe un seul
des $S_i$ contenant $Z$, et pour cette valeur de $i$, on a
$\operatorname{long}(\mathcal F_{s_i})=1$~; en outre
$\mathcal O_{S_i,z}$ est un anneau de valuation discrète dont $t_z$ est une
uniformisante, et $t_z$ est $\mathcal F_z$-régulier.
\end{corollary}

\begin{proof}
Cela résulte de (3.4.1), les deux membres de (3.4.1.1) étant alors égaux.
\end{proof}

\begin{proposition}[3.4.3]
\label{IV.2.3.4.3-fr}
Soient $X$ un préschéma localement noethérien, $t$ une section de
$\mathcal O_X$ au-dessus de $X$, $Y$ le sous-préschéma fermé de $X$ défini par
l'Idéal $t\mathcal O_X$ de $\mathcal O_X$. Soient $\mathcal F$ un
$\mathcal O_X$-Module cohérent, $T$ un cycle premier associé à $\mathcal F$,
$T'$ une composante irréductible de $T\cap Y$, $x$ le point générique de
$T'$. Supposons que $t_x$ soit $\mathcal F_x$-régulier~; alors on a
$x\in\operatorname{Ass}(\mathcal F/t\mathcal F)$.
\end{proposition}

\begin{proof}
Comme dans la démonstration de (3.4.1), on peut se ramener au cas où
$X=\operatorname{Spec}(\mathcal O_x)$~; la proposition est alors (compte tenu
de (3.1.2)) conséquence de $(0,16.4.6.3)$.
\end{proof}

\begin{proposition}[3.4.4]
\label{IV.2.3.4.4-fr}
Soient $X$ un préschéma localement noethérien, $t$ une section de
$\mathcal O_X$ au-dessus de $X$, $Y$ le sous-préschéma fermé de $X$ défini par
l'Idéal $t\mathcal O_X$ de $\mathcal O_X$. Soient $\mathcal F$ un
$\mathcal O_X$-Module cohérent, $(S_i)$ la famille des composantes
irréductibles de $\operatorname{Supp}(\mathcal F)$. Soit $y$ un point de $Y$
tel que $t_y$ soit $\mathcal F_y$-régulier et qu'aucun des cycles premiers
associés immergés de $\mathcal F/t\mathcal F$ ne contienne $y$. Alors les
composantes irréductibles de $\operatorname{Supp}(\mathcal F/t\mathcal F)$ qui
contiennent $y$ sont exactement les composantes irréductibles des $S_i\cap Y$
qui contiennent $y$, et les cycles premiers associés à $\mathcal F$ qui
contiennent $y$ sont non immergés.
\end{proposition}

\begin{proof}
Prouvons d'abord la dernière assertion. Soient $T\supset T_1$ deux cycles
premiers associés à $\mathcal F$ contenant $y$~; si $x$ est point générique
d'une composante irréductible de $T\cap Y$ contenant $y$, $x$ est
générisation de $y$, donc contenu dans tout voisinage de $y$, et l'hypothèse
que $t_y$ est $\mathcal F_y$-régulier entraîne que $t_x$ est
$\mathcal F_x$-régulier $(0,15.2.4)$, donc, en vertu de (3.4.3), on a
$x\in\operatorname{Ass}(\mathcal F/t\mathcal F)$. Soit $x_1$ le point
générique d'une composante irréductible de $T_1\cap Y$ contenant $y$, et soit
$x$ le point générique d'une composante irréductible de $T\cap Y$ contenant
$x_1$~; il résulte de ce qui précède que $x_1$ et $x$ appartiennent tous deux
à $\operatorname{Ass}(\mathcal F/t\mathcal F)$, et comme
$x_1\in\overline{\{x\}}$, l'hypothèse entraîne que $x_1=x$. Désignons encore
par $T$ et $T_1$ les sous-préschémas fermés intègres de $X$ ayant $T$ et
$T_1$ comme espaces sous-jacents respectifs, et posons
$A=\mathcal O_{T,x}$, $A_1=\mathcal O_{T_1,x}$~; on a donc
$A_1=A/\mathfrak p$, où $\mathfrak p$ est un idéal premier de $A$. Par
définition de $x$ et $x_1$, $A/tA$ et $A_1/tA_1$ sont des anneaux artiniens~;
d'autre part, on a vu plus haut que $t_x$ est $\mathcal F_x$-régulier, donc
$x$ ne peut appartenir à $\operatorname{Ass}(\mathcal F)$, et par suite $A_1$
n'est pas artinien. On a donc $\dim A=\dim A_1=1$~; mais cela entraîne
$\mathfrak p=0$ et $A=A_1$ $(0,16.1.2.2)$~; comme
$\operatorname{Spec}(A)$ et $\operatorname{Spec}(A_1)$ sont respectivement
denses dans $T$ et $T_1$, on a bien $T=T_1$.

Les $S_i$ qui contiennent $y$ sont donc tous les cycles premiers associés à
$\mathcal F$ contenant $y$~; si $x$ est le point générique d'une composante
irréductible de $S_i\cap Y$ contenant $y$, on déduit encore de $(0,15.2.4)$
que $t_x$ est $\mathcal F_x$-régulier, donc, par (3.4.3), que
$x\in\operatorname{Ass}(\mathcal F/t\mathcal F)$~; cela démontre la première
assertion de (3.4.4).
\end{proof}

\begin{proposition}[3.4.5]
\label{IV.2.3.4.5-fr}
Soient $X$ un préschéma localement noethérien, $t$ une section de
$\mathcal O_X$ au-dessus de $X$, $Y$ le sous-préschéma fermé de $X$ défini par
l'Idéal $t\mathcal O_X$ de $\mathcal O_X$. Soient $\mathcal F$

\oldpage[IV-2]{50}

un $\mathcal O_X$-Module cohérent, $y$ un point de $Y$~; supposons que $t_y$
soit $\mathcal F_y$-régulier et que $\mathcal F/t\mathcal F$ soit intègre au
point $y$ (3.2.4). Alors $\mathcal F$ est intègre au point $y$.
\end{proposition}

\begin{proof}
Compte tenu de (3.4.4), il suffit de prouver que $y$ est contenu dans une
\emph{seule} composante irréductible de $\operatorname{Supp}(\mathcal F)$, et
que si $s$ est le point générique de cette composante, on a
$\operatorname{long}(\mathcal F_s)=1$. Or, par hypothèse, $y$ n'appartient
qu'à une seule composante irréductible de
$\operatorname{Supp}(\mathcal F/t\mathcal F)$, et si $z$ est le point
générique de cette composante, on a
$\operatorname{long}((\mathcal F/t\mathcal F)_z)=1$~; la conclusion résulte
donc de (3.4.2).
\end{proof}

\begin{proposition}[3.4.6]
\label{IV.2.3.4.6-fr}
Les hypothèses étant celles de (3.4.1), soit $x$ un point de $Y$. Supposons
que $Y$ ne contienne aucun des $S_i$ contenant $x$, et que
$\mathcal F/t\mathcal F$ soit réduit au point $x$ (3.2.2). Alors $t_x$ est
$\mathcal F_x$-régulier et $\mathcal F$ est réduit au point $x$. En outre, si
$z_j$ est point générique d'une composante irréductible de
$\operatorname{Supp}(\mathcal F/t\mathcal F)$ contenant $x$, $z_j$ est
contenu dans un seul des $S_i$, et $\mathcal O_{S_i,z_j}$ est un anneau de
valuation discrète dont $t_{z_j}$ est une uniformisante.
\end{proposition}

\begin{proof}
Le fait que $t_x$ est $\mathcal F_x$-régulier résulte du lemme suivant appliqué
à l'anneau $\mathcal O_x$~:

\begin{lemma}[3.4.6.1]
\label{IV.2.3.4.6.1-fr}
Soient $A$ un anneau noethérien, $M$ un $A$-module de type fini,
$\mathfrak p_i$ les éléments minimaux de $\operatorname{Supp}(M)$, $t$ un
élément de $A$. On suppose que $t$ n'appartient à aucun des $\mathfrak p_i$,
et que $M/tM$ est un $A$-module réduit (3.2.2). Alors $t$ est $M$-régulier.
\end{lemma}

\begin{proof}
Tout idéal premier $\mathfrak p\in\operatorname{Supp}(M)$ contient un des
$\mathfrak p_i$~; comme $t$ n'appartient à aucun des $\mathfrak p_i$,
l'homothétie de rapport $t$ dans $M_{\mathfrak p}$ n'est pas nilpotente
(Bourbaki, \emph{Alg. comm.}, chap.~IV, \S~1, n\textsuperscript{o}~4, cor. de
la prop.~9). Désignons par $N$ le sous-module de $M$ formé des éléments
annulés par une puissance de $t$, et posons $P=M/N$~; nous allons montrer que
$N=0$. Comme $t$ est $P$-régulier, on a une suite exacte (3.4.1.4)
\[
0\longrightarrow N/tN\longrightarrow M/tM\longrightarrow P/tP
\longrightarrow0.
\]
Comme $N$ est de type fini, il est annulé par une puissance de $t$, et il
suffit donc de montrer que $N/tN=0$. Comme $N/tN$ est un sous-module de
$M/tM$, il suffit de prouver que $(N/tN)_{\mathfrak p}=0$ pour tout
$\mathfrak p\in\operatorname{Ass}(M/tM)$, ou encore que l'homomorphisme
\[
u_{\mathfrak p}:(M/tM)_{\mathfrak p}\longrightarrow(P/tP)_{\mathfrak p}
\]
est bijectif pour tout $\mathfrak p\in\operatorname{Ass}(M/tM)$. Or on a
$(P/tP)_{\mathfrak p}\neq0$~; en effet, comme
$\mathfrak p\in\operatorname{Supp}(M/tM)
=\operatorname{Supp}(M)\cap V(t)$ $(0_{\mathrm I},1.7.5)$, l'image de $t$ dans
$A_{\mathfrak p}$ est contenue dans l'idéal maximal
$\mathfrak pA_{\mathfrak p}$, donc l'hypothèse
$P_{\mathfrak p}=tP_{\mathfrak p}$ entraînerait $P_{\mathfrak p}=0$ par le
lemme de Nakayama~; on aurait donc $M_{\mathfrak p}=N_{\mathfrak p}$ et
l'homothétie de rapport $t$ dans $M_{\mathfrak p}$ serait nilpotente~; mais
cela contredit la remarque faite au début, puisque
$\mathfrak p\in\operatorname{Supp}(M)$. Cela étant, l'hypothèse que $M/tM$ est
réduit entraîne $\operatorname{long}((M/tM)_{\mathfrak p})=1$, et comme
$(P/tP)_{\mathfrak p}\neq0$, $u_{\mathfrak p}$ est nécessairement bijectif,
ce qui prouve le lemme.
\end{proof}

Par hypothèse, aucun des cycles premiers associés immergés de
$\mathcal F/t\mathcal F$ ne contient $x$, donc aucun des cycles premiers
associés immergés de $\mathcal F$ ne contient $x$, en vertu de (3.4.4).
D'autre part, appliquant (3.4.2) à une composante irréductible de
$\operatorname{Supp}(\mathcal F/t\mathcal F)$ contenant $x$, on voit que
$\operatorname{long}(\mathcal F_{s_i})=1$ pour tout $S_i$ contenant $x$, ce
qui achève de montrer que $\mathcal F_x$ est réduit~; enfin, les dernières
assertions sont aussi conséquences de (3.4.2).
\end{proof}

\begin{corollary}[3.4.7]
\label{IV.2.3.4.7-fr}
Soient $A$ un anneau local noethérien, $\mathfrak m$ son idéal maximal, $M$ un
$A$-module de type fini, $(x_i)_{1\leqslant i\leqslant k}$ une famille
d'éléments de $\mathfrak m$ formant une partie d'un système de

\oldpage[IV-2]{51}

paramètres pour $M$ $(0,16.3.6)$. Si le $A$-module
\[
N=M\Big/\Big(\sum_{i=1}^k x_iM\Big)
\]
est intègre (3.2.4), alors $M$ est intègre et la suite
$(x_i)_{1\leqslant i\leqslant k}$ est $M$-régulière.
\end{corollary}

\begin{proof}
Par récurrence sur $k$, on est aussitôt ramené au cas $k=1$~; nous écrirons
$x$ au lieu de $x_1$~; l'hypothèse que $x$ fait partie d'un système de
paramètres pour $M$ entraîne que
$\dim(N)=\dim(M)-1$ $(0,16.3.7)$. Posons $n=\dim(N)$~; il y a donc un élément
minimal $\mathfrak p$ de $\operatorname{Supp}(M)$ tel que
$\dim(M/\mathfrak pM)=n+1$ $(0,16.3.4)$, et pour tout entier $j>0$ on a alors
aussi $\dim(M/\mathfrak p^jM)=n+1$ $(0,16.3.5)$~; en outre $x$ fait partie
d'un système de paramètres pour $M/\mathfrak p^jM$ $(0,16.3.5)$, donc, si
l'on pose $M'=M/\mathfrak p^jM$ et $N'=M'/xM'$, on a $\dim(N')=n$. Il est
clair que l'on a un homomorphisme surjectif $v:N\to N'$~; montrons que $v$ est
bijectif. En effet, si $P=\operatorname{Ker}(v)$, on a
$\operatorname{Ass}(P)\subset\operatorname{Ass}(N)$ et puisque $N$ est
intègre, l'hypothèse $P\neq0$ entraînerait que $\operatorname{Ass}(P)$ et
$\operatorname{Ass}(N)$ seraient tous deux réduits à l'unique point
$\mathfrak q$ de
$\operatorname{Ass}(N)=\operatorname{Supp}(N)$~; mais comme
$\dim(N')=\dim(N)$, on a $N'_{\mathfrak q}\neq0$, et l'hypothèse
$\operatorname{long}(N_{\mathfrak q})=1$ entraîne
$\operatorname{long}(N'_{\mathfrak q})=1$ puisque
$N_{\mathfrak q}\to N'_{\mathfrak q}$ est surjective. On aurait donc
$P_{\mathfrak q}=0$, contrairement à l'hypothèse, d'où notre assertion.
Mais alors $N'$, étant isomorphe à $N$, est intègre~; en outre, le support de
$N'$ (égal à l'intersection de $\operatorname{Supp}(M')$ et de $V(x)$) ne
peut contenir $\operatorname{Supp}(M')$, et ce dernier ensemble est
irréductible par construction. L'hypothèse que $M'/xM'$ est intègre (donc
réduit) entraîne alors que $x$ est $M'$-régulier en vertu de (3.4.6). On en
conclut que le noyau de l'homothétie $z\mapsto xz$ dans $M$ est contenu dans
$\mathfrak p^jM\subset\mathfrak m^jM$ pour tout entier $j$, et ce noyau est
donc réduit à $0$ $(0_{\mathrm I},7.3.5)$, ce qui prouve que $x$ est
$M$-régulier. On peut alors appliquer (3.4.5), ce qui prouve que $M$ est
intègre.
\end{proof}

\begin{remark}[3.4.8]
\label{IV.2.3.4.8-fr}
La proposition analogue à (3.4.7), où on remplace «~intègre~» par
«~réduit~» n'est plus nécessairement exacte. Considérons par exemple l'anneau
de polynômes $C=K[X,Y,Z]$ sur un corps $K$, l'anneau quotient
$B=C/\mathfrak p\mathfrak q$, où $\mathfrak p=CZ$,
$\mathfrak q=CX^2+CY$~; soit $A$ l'anneau local de $B$ correspondant à l'idéal
maximal image de $CX+CY+CZ$ dans $B$. Si $x$, $z$ sont les images canoniques
de $X$, $Z$ dans $A$, il est clair que $xz\neq0$ mais $x^2z^2=0$~; d'autre
part, comme $A/xA$ est isomorphe à $K[Y,Z]/(YZ)$, on a
$\dim(A/xA)=1$ tandis que $\dim(A)=2$, donc $x$ appartient à un système de
paramètres de $A$ $(0,16.3.4)$, $A/xA$ est réduit, mais $A$ ne l'est pas.
\end{remark}

\begin{proposition}[3.4.9]
\label{IV.2.3.4.9-fr}
Soient $A$ un anneau noethérien, $M$ un $A$-module de type fini, $f$ un
élément $M$-régulier de $A$ tel que $M/fM$ n'ait pas d'idéaux premiers
associés immergés. Si $\mathfrak p_i$ $(1\leqslant i\leqslant m)$ sont les
idéaux premiers associés à $M/fM$, alors, pour tout entier $n>0$, $f^nM$ est
l'intersection des images réciproques des $f^nM_{\mathfrak p_i}$ par les
applications canoniques $M\to M_{\mathfrak p_i}$ $(1\leqslant i\leqslant m)$.
\end{proposition}

\begin{proof}
Tout revient à montrer que les $\mathfrak p_i$ sont aussi les idéaux premiers
associés à $M/f^nM$, car alors les saturés de $f^nM$ pour les
$\mathfrak p_i$ sont les sous-modules de la décomposition primaire réduite
(nécessairement unique) de $f^nM$ dans $M$. Or, on a
\[
\operatorname{Ass}(f^{n-1}M/f^nM)
\subset\operatorname{Ass}(M/f^nM)
\subset\operatorname{Ass}(M/f^{n-1}M)
\cup\operatorname{Ass}(f^{n-1}M/f^nM)
\]
par (3.1.7), et comme $f$ est $M$-régulier,
$f^{n-1}M/f^nM$ est isomorphe à $M/fM$~; il suffit alors de raisonner par
récurrence sur $n$.
\end{proof}
\oldpage[IV-2]{52}

\section{Changement du corps de base dans les préschémas algébriques}
\label{IV.2.4-fr}

\subsection{Dimension des préschémas algébriques}
\label{subsection:IV.2.4.1-fr}

Nous développerons au \S~5 la théorie générale de la dimension des préschémas~;
mais la théorie de la dimension des préschémas \emph{algébriques} peut se
développer de façon plus élémentaire, et comme elle présente en outre de
nombreux caractères spéciaux, nous en donnerons ici un exposé rapide,
indépendant de la théorie générale.

Si $K$ est un corps et $L$ une extension de $K$, on notera
$\operatorname{deg.tr.}_K L$ le \emph{degré de transcendance} de $L$ sur $K$.

\begin{definition}[4.1.1]
\label{IV.2.4.1.1-fr}
Soient $k$ un corps, $X$ un préschéma localement de type fini sur $k$. On
appelle \emph{dimension} de $X$ le nombre
\begin{equation}
\tag{4.1.1.1}
\label{IV.2.4.1.1.1-fr}
\dim X=\sup_x\operatorname{deg.tr.}_k k(x)
\end{equation}
où $x$ parcourt l'ensemble des points maximaux de $X$.
\end{definition}

Nous verrons au \S~5 (5.2.2) que la définition (4.1.1) ne dépend qu'en
apparence du corps de base $k$ sur lequel $X$ est localement de type fini, et
que le nombre $\dim(X)$ ainsi défini coïncide avec la dimension de
l'\emph{espace sous-jacent} $X$ (définie dans $(0,14.1.2)$). On a évidemment
$\dim X=\dim X_{\mathrm{red}}$.

On notera que chaque $k(x)$ est une extension \emph{de type fini} de $k$
(I, 6.3.3) donc a un degré de transcendance fini sur $k$. Si $X$ est de type
fini sur $k$ et non vide, il est noethérien (I, 6.3.7), donc n'a qu'un nombre
fini de composantes irréductibles, et par suite $\dim(X)$ est \emph{finie} et
$\geqslant0$~; la définition (4.1.1) donne
\[
\dim(\varnothing)=-\infty.
\]

Il est clair que, si l'on désigne par $(X_\alpha)$ la famille des
sous-préschémas fermés réduits de $X$ ayant pour espaces sous-jacents les
composantes irréductibles de $X$ (I, 5.2.1), on a
\begin{equation}
\tag{4.1.1.2}
\label{IV.2.4.1.1.2-fr}
\dim(X)=\sup_\alpha\dim(X_\alpha).
\end{equation}
Cela ramène donc le calcul de la dimension au cas des préschémas \emph{intègres}
(localement de type fini sur $k$). Enfin, on a évidemment
\begin{equation}
\tag{4.1.1.3}
\label{IV.2.4.1.1.3-fr}
\dim(X)=\dim(U)
\end{equation}
pour tout sous-préschéma induit sur un ouvert \emph{partout dense} $U$ de
$X$~; cela ramène donc finalement la notion de dimension au cas d'un schéma
\emph{affine de type fini sur $k$}.

\begin{theorem}[4.1.2]
\label{IV.2.4.1.2-fr}
Soient $X$, $Y$ deux préschémas localement de type fini sur un corps $k$,
$f:X\to Y$ un $k$-morphisme.

\oldpage[IV-2]{53}

\begin{enumerate}
\item[(i)] Si $f$ est quasi-compact et dominant, on a
$\dim(Y)\leqslant\dim(X)$.
\item[(ii)] Si $f$ est quasi-fini, on a
$\dim(X)\leqslant\dim(Y)$.
\item[(iii)] Supposons que $X$ soit de type fini sur $k$. Pour que
$\dim(X)\geqslant n$ (resp. $\dim(X)\leqslant n$, $\dim(X)=n$), il faut et il
suffit qu'il existe un ouvert $U$ dense dans $X$, et un $k$-morphisme
\[
g:U\longrightarrow\mathbf V(k^n)
=\operatorname{Spec}(k[T_1,\ldots,T_n]),
\]
que nous noterons aussi $\mathbf V_k^n$, qui soit surjectif (resp. fini,
resp. fini et surjectif). Si $X$ est un schéma affine, on peut prendre $U=X$.
\end{enumerate}
\end{theorem}

\begin{proof}
(i) Soit $y$ un point maximal de $Y$~; on sait que $f^{-1}(y)$ contient un
point maximal $x$ de $X$ (1.1.5), donc $k(x)$ est une extension de $k(y)$,
d'où l'inégalité
\[
\operatorname{deg.tr.}_k k(y)
\leqslant\operatorname{deg.tr.}_k k(x)
\leqslant\dim(X),
\]
ce qui démontre (i).

(ii) Si $x$ est point maximal de $X$, $k(x)$ est une extension finie de
$k(f(x))$ (II, 6.2.2), donc a même degré de transcendance sur $k$. Si l'on
considère le sous-préschéma fermé réduit de $Y$ ayant $\overline{\{f(x)\}}$
comme espace sous-jacent, on voit qu'on est ramené à démontrer le corollaire
suivant.
\end{proof}

\begin{corollary}[4.1.2.1]
\label{IV.2.4.1.2.1-fr}
Soient $Y$ un $k$-préschéma localement de type fini~; pour tout sous-préschéma
$Z$ de $Y$, on a $\dim(Z)\leqslant\dim(Y)$. Supposons en outre que les
composantes irréductibles de $Y$ aient toutes même dimension~; alors, pour que
$Z$ soit rare dans $Y$, il faut et il suffit que
$\dim(Z)<\dim(Y)$.
\end{corollary}

\begin{proof}
En effet, soit $z$ un point maximal de $Z$~; en considérant un des
sous-préschémas fermés réduits de $Y$ ayant pour espace sous-jacent une des
composantes irréductibles de $Y$ contenant $z$, on se ramène au cas où $Y$ est
intègre de point générique $y$, puis, en considérant dans $Y$ un ouvert affine
contenant $z$, au cas où $Y$ est affine et de type fini sur $k$, soit
$Y=\operatorname{Spec}(A)$, et $Z$ fermé dans $Y$, soit
$Z=\operatorname{Spec}(A/\mathfrak a)$ où $\mathfrak a$ est un idéal de
l'anneau intègre $A$, distinct de $A$. En vertu du lemme de normalisation
(Bourbaki, \emph{Alg. comm.}, chap.~V, \S~3,
n\textsuperscript{o}~1, th.~1), il existe une suite finie
$(x_i)_{1\leqslant i\leqslant n}$ d'éléments de $A$, algébriquement
indépendants sur $k$, tels que $A$ soit entier sur l'anneau
$B=k[x_1,\ldots,x_n]$ et que $\mathfrak a\cap B$ soit engendré par une
sous-famille $(x_i)_{1\leqslant i\leqslant p}$ (éventuellement vide) de
$(x_i)_{1\leqslant i\leqslant n}$. Soit
$g:Y\to\operatorname{Spec}(B)=\mathbf V_k^n$ le morphisme fini dominant qui
correspond à l'injection canonique $B\to A$~; comme
$C=B/(\mathfrak a\cap B)$ est isomorphe à $k[x_{p+1},\ldots,x_n]$, $g$ induit
sur $Z$ un morphisme fini dominant $Z\to\mathbf V_k^{n-p}$~; on a donc
$\operatorname{deg.tr.}_k k(y)=n$ et
$\operatorname{deg.tr.}_k k(z)=n-p$, puisque $k(y)$ (resp. $k(z)$) est une
extension finie de $k(g(y))$ (resp. $k(g(z))$). Cela démontre la première
assertion de (4.1.2.1). En outre, si $p=0$, on a nécessairement $z=y$, car $y$
est le seul point de $Y$ dont l'image par $g$ soit le point générique de
$\mathbf V_k^n$~; dans ce cas, $Z$ contient donc un ouvert non vide de $Y$. Si
au contraire $p>0$, on a nécessairement $z\neq y$ et
$\overline{\{z\}}$ est donc rare dans $\overline{\{y\}}$, ce qui achève de
prouver (4.1.2.1).
\end{proof}

\begin{proof}
(iii) Le fait que les conditions énoncées sont suffisantes résulte aussitôt de
(i) et (ii) et de (1.5.4, (v)). Pour prouver qu'elles sont nécessaires, on peut
considérer un ouvert dense $U$ de $X$, réunion d'ouverts affines $U_i$ deux à
deux disjoints et irréductibles $(1\leqslant i\leqslant m)$, dont chacun
contient un des points maximaux de $X$ $(0_{\mathrm I},2.1.6)$~; on peut en
outre supposer $X$ réduit et si $X$ est affine, on peut bien entendu prendre
$U=X$. Le même

\oldpage[IV-2]{54}

raisonnement que dans (4.1.2.1) montre que si $\dim(U_i)=n_i$, il existe un
$k$-morphisme $g_i:U_i\to\mathbf V_k^{n_i}$ qui est fini et dominant, donc
surjectif (II, 6.1.10). Soit $n=\dim(X)=\sup_i n_i$~; pour chaque $n_i$ il y a
donc un morphisme $h_i:\mathbf V_k^{n_i}\to\mathbf V_k^n$ qui est une
immersion fermée correspondant à l'homomorphisme canonique
$k[T_1,\ldots,T_n]\to k[T_1,\ldots,T_{n_i}]$, et qui est l'identité pour
$n_i=n$. Comme $U$ est somme des $U_i$, on prendra pour $g$ le morphisme qui
dans chaque $U_i$ coïncide avec $h_i\circ g_i$~; il est évidemment fini et
surjectif. Lorsque $n'\geqslant n$, on aura un morphisme fini
$U\to\mathbf V_k^{n'}$ en composant l'immersion fermée canonique
$\mathbf V_k^n\to\mathbf V_k^{n'}$ avec $g$~; lorsque $n'\leqslant n$, on
compose de même avec $g$ le morphisme canonique
$p:\mathbf V_k^n\to\mathbf V_k^{n'}$ correspondant à l'injection canonique
$k[T_1,\ldots,T_{n'}]\to k[T_1,\ldots,T_n]$, en notant que $p$ est fidèlement
plat, donc \emph{surjectif}.
\end{proof}

\begin{remark}[4.1.3]
\label{IV.2.4.1.3-fr}
Le corollaire (4.1.2.1) montre que dans la formule (4.1.1.1), on peut supposer
que $x$ parcourt l'ensemble de \emph{tous les points} de $X$.
\end{remark}

\begin{corollary}[4.1.4]
\label{IV.2.4.1.4-fr}
Soient $X$ un préschéma localement de type fini sur un corps $k$, $K$ une
extension de $k$~; alors $\dim(X\otimes_k K)=\dim(X)$.
\end{corollary}

\begin{proof}
On peut évidemment se borner au cas où $X$ est de type fini sur $k$~; alors le
morphisme $u:X\otimes_k K\to X$ est fidèlement plat (2.2.13, (i)), donc si
$U$ est un ouvert partout dense dans $X$, $u^{-1}(U)=U\otimes_k K$ est dense
dans $X\otimes_k K$ (2.3.10)~; si $g:U\to\mathbf V_k^n$ est fini et
surjectif, il en est de même de
$g_{(K)}:U\otimes_k K\to\mathbf V_K^n$ (I, 3.5.2 et II, 6.1.5), d'où le
corollaire.
\end{proof}

\begin{corollary}[4.1.5]
\label{IV.2.4.1.5-fr}
Soient $X$ et $Y$ deux préschémas localement de type fini sur un corps $k$~;
alors $\dim(X\times_k Y)=\dim(X)+\dim(Y)$.
\end{corollary}

\begin{proof}
Il suffit de démontrer que si $U$ (resp. $V$) est un voisinage affine d'un
point $x$ (resp. $y$) de $X$ (resp. $Y$) tel que
$\operatorname{deg.tr.}_k k(x)=m=\dim U$ et
$\operatorname{deg.tr.}_k k(y)=n=\dim V$, alors on a
$\dim(U\times_k V)=m+n$, autrement dit, on peut supposer en vertu de (4.1.2)
qu'il existe des $k$-morphismes finis surjectifs
$f:X\to\mathbf V_k^m$, $g:Y\to\mathbf V_k^n$~; alors
$f\times g:X\times_k Y\to\mathbf V_k^{m+n}$ est fini et surjectif
(I, 3.5.2 et II, 6.1.5), d'où le corollaire.
\end{proof}

\subsection{Cycles premiers associés sur les préschémas algébriques}
\label{subsection:IV.2.4.2-fr}

\begin{proposition}[4.2.1]
\EGAReaderTarget{ega4.c04.prop.4-2-1}{IV.2.4.2.1-fr}
Soient $K$ et $L$ deux extensions d'un corps $k$, telles que $K\otimes_k L$
soit noethérien. Alors les idéaux premiers associés à $K\otimes_k L$ sont
minimaux, et si $E$ est le corps résiduel de l'anneau local d'un tel idéal, on
a
\begin{equation}
\tag{4.2.1.1}
\label{IV.2.4.2.1.1-fr}
\operatorname{deg.tr.}_K E=\operatorname{deg.tr.}_k L,
\qquad
\operatorname{deg.tr.}_L E=\operatorname{deg.tr.}_k K,
\end{equation}
d'où
\begin{equation}
\tag{4.2.1.2}
\label{IV.2.4.2.1.2-fr}
\operatorname{deg.tr.}_k E
=\operatorname{deg.tr.}_k K+\operatorname{deg.tr.}_k L.
\end{equation}
\end{proposition}

\begin{proof}
On sait que $K$ est une extension algébrique d'une extension transcendante
pure $K'=k(\mathbf t)$, où $\mathbf t=(t_\alpha)_{\alpha\in I}$ est une famille
d'indéterminées~; l'anneau $k[\mathbf t]\otimes_k L=L[\mathbf t]$ est intègre,
donc il en est de même de $K'\otimes_k L$, qui en est un anneau de fractions,
et le corps des fractions de $K'\otimes_k L$ est $L(\mathbf t)$~; on a le
diagramme commutatif d'homomorphismes canoniques

\oldpage[IV-2]{55}

\begin{equation}
\tag{4.2.1.3}
\label{IV.2.4.2.1.3-fr}
\xymatrix@C=3.5em@R=2.8em{
K \ar[r] & K\otimes_k L \ar[r] & K\otimes_{K'}L(\mathbf t)\\
K' \ar[u] \ar[r] & K'\otimes_k L \ar[u] \ar[r] & L(\mathbf t) \ar[u]\\
k \ar[u] \ar[r] & L \ar[u] &
}
\end{equation}

Comme $K$ est fidèlement plat sur $K'$, $K\otimes_k L
=K\otimes_{K'}(K'\otimes_k L)$ est fidèlement plat sur $K'\otimes_k L$,
donc $K'\otimes_k L$ est noethérien $(0_{\mathrm I},6.5.2)$~; en outre
$K'\otimes_k L$ s'identifie à un sous-anneau de $K\otimes_k L$~; la trace
sur $K'\otimes_k L$ d'un idéal premier $\mathfrak p$ associé à $K\otimes_k L$
est l'idéal $0$, un élément $\neq0$ de $K'\otimes_k L$ n'étant pas diviseur de
zéro dans $K\otimes_k L$ $(0_{\mathrm I},6.3.4$ et Bourbaki,
\emph{Alg. comm.}, chap.~IV, \S~1, n\textsuperscript{o}~1, cor.~3 de la
prop.~2). Comme en outre $K$ est algébrique sur $K'$, $K\otimes_k L$ est une
algèbre entière sur $K'\otimes_k L$, et les idéaux premiers de $K\otimes_k L$
induisant $0$ sur $K'\otimes_k L$ sont nécessairement sans relation
d'inclusion mutuelle (Bourbaki, \emph{Alg. comm.}, chap.~V, \S~2,
n\textsuperscript{o}~1, cor.~1 de la prop.~1)~; cela prouve la première
assertion (Bourbaki, \emph{Alg. comm.}, chap.~IV, \S~1,
n\textsuperscript{o}~3, cor.~1 de la prop.~7). De plus, le corps résiduel $E$
de $\mathfrak p$ est algébrique sur le corps résiduel de l'idéal $(0)$ de
$K'\otimes_k L$, c'est-à-dire $L(\mathbf t)$~; donc
$\operatorname{deg.tr.}_L E=\operatorname{deg.tr.}_L L(\mathbf t)
=\operatorname{Card}(I)=\operatorname{deg.tr.}_k K$, autrement dit on a la
première relation (4.2.1.1)~; échangeant les rôles de $K$ et $L$, on a la
seconde relation (4.2.1.1), d'où (4.2.1.2).
\end{proof}

\begin{corollary}[4.2.2]
\label{IV.2.4.2.2-fr}
Sous les hypothèses de (3.3.6), si les préschémas $T_{x,y}$ sont localement
noethériens, ils n'ont pas de cycle premier associé immergé.
\end{corollary}

\begin{corollary}[4.2.3]
\label{IV.2.4.2.3-fr}
Sous les hypothèses de (3.3.6) (resp. (3.3.7)), si les $T_{x,y}$ sont
localement noethériens et si $\mathcal F$ et les $\mathcal G_s$ (pour $s\in S$)
(resp. $\mathcal F$ et $\mathcal G$) sont sans cycle premier associé immergé,
il en est de même de $\mathcal F\otimes_S\mathcal G$.
\end{corollary}

Cela résulte de (4.2.2), (3.3.2) et de la démonstration de (3.3.6). En
particulier, comme tout préschéma sur un corps $k$ est \emph{plat sur $k$}, on
a ainsi démontré l'assertion (i) de la

\begin{proposition}[4.2.4]
\label{IV.2.4.2.4-fr}
Soient $k$ un corps, $X$ et $Y$ deux $k$-préschémas localement noethériens
tels que $X\times_k Y$ soit localement noethérien. Supposons de plus $X$ et
$Y$ intègres. Alors~:
\begin{enumerate}
\item[(i)] $X\times_k Y$ est sans cycle premier associé immergé~; chacune des
composantes irréductibles de $X\times_k Y$ domine $X$ et $Y$, et l'ensemble de
ces composantes est en correspondance biunivoque avec l'ensemble des
composantes irréductibles de
$\operatorname{Spec}(R(X)\otimes_k R(Y))$ (autrement dit l'ensemble des
idéaux premiers minimaux de $R(X)\otimes_k R(Y)$), où $R(X)$ et $R(Y)$ sont
les corps des fonctions rationnelles de $X$ et $Y$ respectivement.
\item[(ii)] Si un point maximal $z$ de $X\times_k Y$ est identifié à un idéal
premier minimal $\mathfrak p$ de $R(X)\otimes_k R(Y)$, l'anneau local
$\mathcal O_{X\times_k Y,z}$ est isomorphe à l'anneau de fractions
$(R(X)\otimes_k R(Y))_{\mathfrak p}$. En particulier, si $R(X)$ ou $R(Y)$ est
séparable sur $k$, $X\times_k Y$ est réduit.

\oldpage[IV-2]{56}

\item[(iii)] Si en outre $X$ et $Y$ sont localement de type fini sur $k$,
toute composante irréductible de $X\times_k Y$ est de dimension
$\dim(X)+\dim(Y)$.
\end{enumerate}
\end{proposition}

\begin{proof}
L'assertion (iii) résulte de (4.2.1.2), compte tenu de (i) et (ii). Pour
démontrer (ii), on peut évidemment se borner au cas où
$X=\operatorname{Spec}(A)$, $Y=\operatorname{Spec}(B)$ sont affines, $A$ et
$B$ étant donc des anneaux intègres de corps des fractions respectifs
$K=R(X)$ et $L=R(Y)$~; l'assertion (i) montre que tout idéal premier minimal
$\mathfrak q$ de $A\otimes_k B$ est la trace sur $A\otimes_k B$ d'un idéal
premier minimal $\mathfrak p$ de $K\otimes_k L$~; comme $K\otimes_k L$ est un
anneau de fractions de $A\otimes_k B$, l'isomorphie des anneaux
$(A\otimes_k B)_{\mathfrak q}$ et $(K\otimes_k L)_{\mathfrak p}$ résulte de
$(0_{\mathrm I},1.2.5)$. Enfin, si par exemple $R(Y)$ est séparable sur $k$,
on sait que l'anneau $R(X)\otimes_k R(Y)$ est réduit (Bourbaki,
\emph{Alg.}, chap.~VIII, \S~7, n\textsuperscript{o}~3, th.~1)~; il en est donc
de même des anneaux locaux des points maximaux de $X\times_k Y$. On en déduit
que $X\times_k Y$ est réduit (3.2.1).
\end{proof}

\begin{proposition}[4.2.5]
\label{IV.2.4.2.5-fr}
Soient $k$ un corps, $X$ et $Y$ des $k$-préschémas localement noethériens,
$\mathcal F$ (resp. $\mathcal G$) un $\mathcal O_X$-Module (resp. un
$\mathcal O_Y$-Module) quasi-cohérent. Soit $(Z'_\lambda)$ (resp.
$(Z''_\mu)$) la famille des cycles premiers associés à $\mathcal F$ (resp.
$\mathcal G$), et désignons encore par $Z'_\lambda$ (resp. $Z''_\mu$) le
sous-préschéma réduit de $X$ (resp. $Y$) ayant $Z'_\lambda$ (resp.
$Z''_\mu$) comme espace sous-jacent. Alors, si $Z'_\lambda\times_k Z''_\mu$
est localement noethérien, les composantes irréductibles $Z_{\lambda\mu\nu}$
de $Z'_\lambda\times_k Z''_\mu$ dominent $Z'_\lambda$ et $Z''_\mu$, et
$(Z_{\lambda\mu\nu})$ est la famille des cycles premiers distincts associés à
$\mathcal F\otimes_k\mathcal G$.
\end{proposition}

\begin{proof}
Il suffit d'appliquer (4.2.4) au produit
$\operatorname{Spec}(k(x))\times_k\operatorname{Spec}(k(y))$.
\end{proof}

En particulier~:

\begin{corollary}[4.2.6]
\label{IV.2.4.2.6-fr}
Soient $k$ un corps, $X$, $Y$ deux $k$-préschémas localement noethériens tels
que $X\times_k Y$ soit localement noethérien. Soit $(Z'_\lambda)$ (resp.
$(Z''_\mu)$) la famille des sous-préschémas réduits de $X$ (resp. $Y$) ayant
pour espaces sous-jacents les composantes irréductibles de $X$ (resp. $Y$).
Alors les composantes irréductibles $Z_{\lambda\mu\nu}$ de
$Z'_\lambda\times_k Z''_\mu$ dominent $Z'_\lambda$ et $Z''_\mu$, et
$(Z_{\lambda\mu\nu})$ est la famille des composantes irréductibles de
$X\times_k Y$.
\end{corollary}

\begin{proof}
En effet, on peut se borner au cas où $X$ et $Y$ sont réduits (I, 5.1.8)~; les
composantes irréductibles de $X$ (resp. $Y$) sont alors les cycles premiers
associés à $\mathcal O_X$ (resp. $\mathcal O_Y$) (3.2.1). Appliquons (4.2.5)
à $\mathcal F=\mathcal O_X$ et $\mathcal G=\mathcal O_Y$, en notant que par
définition $\mathcal O_X\otimes_k\mathcal O_Y=\mathcal O_{X\times_k Y}$~; le
corollaire résulte de ce que $\mathcal O_X\otimes_k\mathcal O_Y$ n'a pas de
cycles premiers associés immergés puisqu'il en est ainsi de $\mathcal O_X$ et
$\mathcal O_Y$ par hypothèse (4.2.3).
\end{proof}

Appliquons les résultats qui précèdent au cas où $Y$ est le spectre d'une
extension $K$ de $k$~:

\begin{proposition}[4.2.7]
\label{IV.2.4.2.7-fr}
Soient $k$ un corps, $X$ un $k$-préschéma, $K$ une extension de $k$ telle que
$X\otimes_k K$ soit localement noethérien, $\mathcal F$ un
$\mathcal O_X$-Module quasi-cohérent, $x'$ un point de $X\otimes_k K$, $x$
son image dans $X$.
\begin{enumerate}
\item[(i)] Soit $(Z_\lambda)$ la famille des sous-préschémas réduits de $X$
ayant pour espaces sous-jacents les cycles premiers associés à $\mathcal F$~;
alors les composantes irréductibles $Z_{\lambda\mu}$ des
$Z_\lambda\otimes_k K$ sont les cycles premiers associés à
$\mathcal F\otimes_k K$, et $Z_{\lambda\mu}$ domine $Z_\lambda$~; en outre,
pour qu'un $Z_{\lambda\mu}$ soit immergé, il faut et il suffit que $Z_\lambda$
le soit.
\item[(ii)] Pour que $x$ appartienne à un cycle premier associé immergé de
$\mathcal F$, il faut et il suffit que $x'$ appartienne à un cycle premier
associé immergé de $\mathcal F\otimes_k K$~; pour que $\mathcal F$ soit sans
cycle premier associé immergé, il faut et il suffit qu'il en soit ainsi de
$\mathcal F\otimes_k K$.

\oldpage[IV-2]{57}

\item[(iii)] Les indices $\lambda$ tels que $x\in Z_\lambda$ sont les mêmes
que ceux pour lesquels il existe un indice $\mu$ tel que
$x'\in Z_{\lambda\mu}$. En particulier, si $x'$ n'appartient qu'à un seul
cycle premier associé à $\mathcal F\otimes_k K$, $x$ n'appartient qu'à un
seul cycle premier associé à $\mathcal F$.
\item[(iv)] Si $X$ est localement de type fini sur $k$, on a
$\dim(Z_{\lambda\mu})=\dim(Z_\lambda)$.
\end{enumerate}
\end{proposition}

\begin{proof}
On notera que l'hypothèse entraîne que $X$ lui-même est localement noethérien
(2.2.13 et 2.2.14)~; l'assertion (i) résulte de (4.2.5) et de la démonstration
de (4.2.3), pour $\mathcal G=\mathcal O_Y$, avec
$Y=\operatorname{Spec}(K)$~; (ii) et (iii) résultent de (i) et de (2.3.5)~;
enfin, (iv) est un cas particulier de (4.2.4, (iii)).
\end{proof}

\begin{corollary}[4.2.8]
\label{IV.2.4.2.8-fr}
Si $X$ est localement de type fini sur $k$, l'ensemble des dimensions des
cycles premiers associés est le même pour $\mathcal F$ et
$\mathcal F\otimes_k K$~; l'ensemble des dimensions des composantes
irréductibles est le même pour $X$ et $X\otimes_k K$.
\end{corollary}

\begin{proposition}[4.2.9]
\label{IV.2.4.2.9-fr}
Supposons vérifiées les hypothèses de (4.2.5), et supposons en outre que
$\mathcal F$ et $\mathcal G$ soient cohérents. Soient
$(\mathcal F_\lambda)_{\lambda\in L}$ et
$(\mathcal G_\mu)_{\mu\in M}$ des décompositions irredondantes de
$\mathcal F$ et $\mathcal G$ respectivement~; pour tout couple
$(\lambda,\mu)\in L\times M$, soit
$(\mathcal H_{\lambda\mu\nu})_{\nu\in S(\lambda,\mu)}$ une décomposition
irredondante réduite de $\mathcal F_\lambda\otimes_k\mathcal G_\mu$, où
$S(\lambda,\mu)=\operatorname{Ass}(\mathcal F_\lambda\otimes_k\mathcal G_\mu)$
(3.2.5). Alors $(\mathcal H_{\lambda\mu\nu})$ (pour tous les triplets
$(\lambda,\mu,\nu)$) est une décomposition irredondante de
$\mathcal F\otimes_k\mathcal G$~; elle est réduite si
$(\mathcal F_\lambda)$ et $(\mathcal G_\mu)$ le sont.
\end{proposition}

\begin{proof}
Notons que $\mathcal F\otimes_k\mathcal G$ et chacun des
$\mathcal F_\lambda\otimes_k\mathcal G_\mu$ sont cohérents, et chacun des
$\mathcal F_\lambda\otimes_k\mathcal G_\mu$ s'identifie à un quotient de
$\mathcal F\otimes_k\mathcal G$~; chacun des $\mathcal H_{\lambda\mu\nu}$
s'identifie par définition à un quotient cohérent de
$\mathcal F_\lambda\otimes_k\mathcal G_\mu$, donc aussi à un quotient
cohérent de $\mathcal F\otimes_k\mathcal G$. La famille des supports des
$\mathcal H_{\lambda\mu\nu}$ est localement finie, car
$\operatorname{Supp}(\mathcal F_\lambda\otimes_k\mathcal G_\mu)$ est contenu
dans l'espace sous-jacent à
$\operatorname{Supp}(\mathcal F_\lambda)\times_k
\operatorname{Supp}(\mathcal G_\mu)$ (où
$\operatorname{Supp}(\mathcal F_\lambda)$ et
$\operatorname{Supp}(\mathcal G_\mu)$ désignent les sous-préschémas fermés
réduits correspondants), et pour $\lambda$ et $\mu$ donnés, la famille des
supports des $\mathcal H_{\lambda\mu\nu}$ est localement finie par hypothèse~;
notre assertion résulte donc de ce que les familles
$(\operatorname{Supp}(\mathcal F_\lambda))$ et
$(\operatorname{Supp}(\mathcal G_\mu))$ sont localement finies. Pour démontrer
que $(\mathcal H_{\lambda\mu\nu})$ est une décomposition irredondante de
$\mathcal F\otimes_k\mathcal G$, il suffit donc de prouver que
l'homomorphisme canonique
\[
\mathcal F\otimes_k\mathcal G
\longrightarrow\bigoplus_{\lambda,\mu,\nu}\mathcal H_{\lambda\mu\nu}
\]
est injectif~; or, il est composé des homomorphismes
\[
\mathcal F\otimes_k\mathcal G
\longrightarrow
\left(\bigoplus_\lambda\mathcal F_\lambda\right)\otimes_k
\left(\bigoplus_\mu\mathcal G_\mu\right)
\longrightarrow\bigoplus_{\lambda,\mu,\nu}\mathcal H_{\lambda\mu\nu}~;
\]
le second de ces homomorphismes est injectif par définition, et il en est de
même du premier, qui n'est autre que le produit tensoriel des homomorphismes
canoniques injectifs
$\mathcal F\to\bigoplus_\lambda\mathcal F_\lambda$,
$\mathcal G\to\bigoplus_\mu\mathcal G_\mu$ (on se souviendra que
$\mathcal F$ et $\mathcal G$ sont \emph{plats sur $k$}). Enfin, si
$(\mathcal F_\lambda)$ et $(\mathcal G_\mu)$ sont réduites, on peut supposer
que $L=\operatorname{Ass}(\mathcal F)$ et
$M=\operatorname{Ass}(\mathcal G)$~; le fait que les
$\mathcal H_{\lambda\mu\nu}$ forment alors une décomposition réduite résulte
de ce que les $\operatorname{Ass}(\mathcal F_\lambda\otimes_k
\mathcal G_\mu)$ forment une partition de
$\operatorname{Ass}(\mathcal F\otimes\mathcal G)$ en vertu de (4.2.2)
(cf. (3.2.5)).
\end{proof}

\begin{corollary}[4.2.10]
\label{IV.2.4.2.10-fr}
Sous les hypothèses de (4.2.7), supposons en outre $\mathcal F$ cohérent, et
soit $(\mathcal F_\lambda)_{\lambda\in L}$ une décomposition irredondante de
$\mathcal F$~; pour tout $\lambda\in L$, soit
$(\mathcal F_{\lambda\mu})_{\mu\in
\operatorname{Ass}(\mathcal F_\lambda\otimes_k K)}$ une décomposition
irredondante réduite de $\mathcal F_\lambda\otimes_k K$. Alors
$(\mathcal F_{\lambda\mu})$ est une décomposition irredondante de
$\mathcal F\otimes_k K$~; elle est réduite si $(\mathcal F_\lambda)$ l'est.
\end{corollary}

\begin{proof}
On applique (4.2.9) avec $Y=\operatorname{Spec}(K)$,
$\mathcal G=\mathcal O_Y=K$, $M$ réduit à un seul élément.
\end{proof}

\oldpage[IV-2]{58}

\subsection{Rappels sur les produits tensoriels de corps}
\label{subsection:IV.2.4.3-fr}

Pour la commodité du lecteur, nous rappelons ici certaines propriétés des
produits tensoriels de corps dont nous ferons usage dans les numéros suivants~;
pour les démonstrations, nous renvoyons à [1].

\begin{env}[4.3.1]
\label{IV.2.4.3.1-fr}
Rappelons qu'une extension $L$ d'un corps $k$ est dite \emph{primaire} si la
plus grande extension algébrique séparable de $k$ contenue dans $L$ est
$k$ lui-même.
\end{env}

\begin{proposition}[4.3.2]
\label{IV.2.4.3.2-fr}
Soient $K$, $L$ deux extensions d'un corps $k$. Si $L$ est une extension
primaire de $k$, alors $\operatorname{Spec}(L\otimes_k K)$ est irréductible,
et si $\xi$ est son point générique, $k(\xi)$ est une extension primaire de
$K$. Réciproquement, si pour toute extension finie séparable $K$ de $k$,
$\operatorname{Spec}(L\otimes_k K)$ est irréductible, $L$ est une extension
primaire de $k$.
\end{proposition}

\begin{proof}
Voir la démonstration dans [1], p.~14-03 à 14-06.
\end{proof}

\begin{corollary}[4.3.3]
\label{IV.2.4.3.3-fr}
Si $k$ est séparablement clos (i.e. si sa clôture algébrique est radicielle sur
$k$), alors, pour deux extensions quelconques $K$, $L$ de $k$,
$\operatorname{Spec}(L\otimes_k K)$ est irréductible, et réciproquement.
\end{corollary}

\begin{proof}
Cela découle aussitôt de (4.3.2).
\end{proof}

\begin{corollary}[4.3.4]
\label{IV.2.4.3.4-fr}
Soient $L$ une extension d'un corps $k$, $L_s$ la plus grande extension
algébrique séparable de $k$ contenue dans $L$ (parfois appelée la
\emph{fermeture algébrique séparable} de $k$ dans $L$)~; supposons $L_s$ de
degré fini sur $k$, et soit $K$ une extension galoisienne de $k$ contenant
$L_s$. Alors $\operatorname{Spec}(L\otimes_k K)$ a $[L_s:k]$ composantes
irréductibles dont il est la somme et les corps résiduels aux points génériques
de ces composantes sont des extensions primaires de $K$.
\end{corollary}

\begin{proof}
En effet, on a $L\otimes_k K=L\otimes_{L_s}(L_s\otimes_k K)$ et on sait que
$L_s\otimes_k K$ est isomorphe à un produit de $[L_s:k]$ corps isomorphes à
$K$ (Bourbaki, \emph{Alg.}, chap.~VIII, \S~7, n\textsuperscript{o}~3,
cor.~2 du th.~1)~; donc $L\otimes_k K$ est isomorphe à un produit de
$[L_s:k]$ anneaux isomorphes à $L\otimes_{L_s}K$~; comme $L$ est une
extension primaire de $L_s$, la conclusion résulte de (4.3.2).
\end{proof}

\begin{proposition}[4.3.5]
\label{IV.2.4.3.5-fr}
Soient $K$, $L$ deux extensions d'un corps $k$. Si $L$ est une extension
séparable de $k$, l'anneau $L\otimes_k K$ est réduit. Réciproquement, si pour
toute extension finie radicielle $K$ de $k$, l'anneau $L\otimes_k K$ est
réduit, $L$ est une extension séparable de $k$.
\end{proposition}

\begin{proof}
Pour la démonstration, voir Bourbaki, \emph{Alg.}, chap.~VIII, \S~7,
n\textsuperscript{o}~3, th.~1.
\end{proof}

\begin{corollary}[4.3.6]
\label{IV.2.4.3.6-fr}
Si $k$ est parfait, alors, pour deux extensions quelconques $K$, $L$ de $k$,
$K\otimes_k L$ est réduit, et réciproquement.
\end{corollary}

\begin{proof}
Cela résulte aussitôt de (4.3.5).
\end{proof}

\begin{corollary}[4.3.7]
\label{IV.2.4.3.7-fr}
Soit $L$ une extension séparable de $k$ et soit $K$ une extension de $k$ telle
que $K$ ou $L$ soit finie sur $k$, alors les corps résiduels de l'anneau
semi-local $L\otimes_k K$ sont des extensions séparables de $K$.
\end{corollary}

\begin{proof}
On sait en effet que $L\otimes_k K$ est composé direct de ces corps résiduels
$E_i$~; pour toute extension $K'$ de $K$, $L\otimes_k K'$ est donc isomorphe
au composé direct des anneaux $E_i\otimes_K K'$~; comme $L\otimes_k K'$ est
réduit, les $E_i$ sont séparables sur $K$ en vertu de (4.3.5).
\end{proof}

\begin{corollary}[4.3.8]
\label{IV.2.4.3.8-fr}
Si $K$ et $L$ sont deux extensions finies séparables de $k$, $K\otimes_k L$
est produit d'un nombre fini d'extensions séparables de $k$.
\end{corollary}

\oldpage[IV-2]{59}

\begin{proposition}[4.3.9]
\label{IV.2.4.3.9-fr}
Si $k$ est algébriquement clos, alors, pour deux extensions quelconques $K$,
$L$ de $k$, $L\otimes_k K$ est intègre, et réciproquement.
\end{proposition}

\begin{proof}
C'est une conséquence de (4.3.3) et (4.3.6), un corps parfait et séparablement
clos étant algébriquement clos.
\end{proof}

\subsection{Préschémas irréductibles et préschémas connexes sur un corps
algébriquement clos}
\label{subsection:IV.2.4.4-fr}

\begin{env}[4.4.1]
\label{IV.2.4.4.1-fr}
Soient $k$ un corps, $X$ un $k$-préschéma, $K$ une extension de $k$. Comme le
morphisme $\operatorname{Spec}(K)\to\operatorname{Spec}(k)$ est fidèlement
plat, quasi-compact, et universellement ouvert (2.4.9), il en est de même du
morphisme projection $p:X\otimes_k K\to X$ (2.2.13). Il résulte donc de
(2.3.5) que toute composante irréductible de $X\otimes_k K$ \emph{domine} une
composante irréductible de $X$, et (comme $p$ est surjectif) toute composante
irréductible de $X$ est dominée par une composante irréductible de
$X\otimes_k K$~; on déduit ainsi de $p$ une application \emph{surjective} de
l'ensemble des composantes irréductibles de $X\otimes_k K$ dans l'ensemble
des composantes irréductibles de $X$. De même, il résulte du fait que $p$ est
continu et surjectif que toute composante connexe de $X$ est image par $p$
d'une réunion de composantes connexes de $X\otimes_k K$, d'où une application
\emph{surjective} de l'ensemble des composantes connexes de $X\otimes_k K$
dans l'ensemble des composantes connexes de $X$. On en conclut que toute
composante irréductible (resp. connexe) $Z'$ de $X\otimes_k K$ est contenue
dans un \emph{unique} ensemble de la forme $p^{-1}(Z)$, où $Z$ est une
composante irréductible (resp. connexe) de $X$~; pour tout sous-préschéma de
$X$ ayant $Z$ pour espace sous-jacent (et encore noté $Z$), $Z'$ est une
composante irréductible (resp. connexe) de $Z\otimes_k K$. En particulier, si
$X\otimes_k K$ a un nombre fini $n$ (resp. $n'$) de composantes irréductibles
(resp. connexes) (ce qui sera le cas lorsque $X$ est \emph{de type fini sur
$k$}, car alors $X\otimes_k K$ est de type fini sur $K$, donc noethérien), le
nombre de composantes irréductibles (resp. connexes) de $X$ est $\leqslant n$
(resp. $\leqslant n'$)~; dire qu'il est égal à $n$ (resp. $n'$) signifie que
pour tout sous-préschéma réduit $Z$ de $X$ ayant pour espace sous-jacent une
composante irréductible (resp. connexe) de $X$, $Z\otimes_k K$ est
irréductible (resp. connexe) (I, 5.1.8)~; le nombre de composantes
irréductibles (resp. connexes) de $X\otimes_k K$ est donc indépendant de $K$
si et seulement si pour toute composante irréductible (resp. connexe) $Z$ de
$X$, $Z\otimes_k K$ est irréductible (resp. connexe) pour \emph{toute}
extension $K$ de $k$.

En particulier, si $X\otimes_k K$ est irréductible (resp. connexe), il en est
de même de $X$ (ce qui résulte déjà du fait que $p$ est surjectif). De même, si
$X\otimes_k K$ est \emph{réduit} (resp. \emph{intègre}), il en est de même de
$X$ puisque $p$ est fidèlement plat (2.1.13).

Dans ce numéro et le suivant, nous allons examiner de plus près ce qu'on peut
dire des composantes irréductibles (resp. connexes) de $X\otimes_k K$ lorsque
$K$ varie~; ce qui précède nous montre déjà que le nombre de ces composantes
est fonction \emph{croissante} de $K$.
\end{env}

\begin{lemma}[4.4.2]
\EGAReaderTarget{ega4.c04.lem.4-4-2}{IV.2.4.4.2-fr}
Soient $X'$, $X$ deux espaces topologiques, $f:X'\to X$ une application
continue. On suppose que les conditions suivantes sont vérifiées~:
\begin{enumerate}
\item[(i)] $f$ est surjective ouverte (resp. telle que l'espace topologique
$X$ s'identifie canoniquement au quotient de $X'$ par la relation d'équivalence
définie par $f$).

\oldpage[IV-2]{60}

\item[(ii)] Pour tout $x\in X$, $f^{-1}(x)$ est irréductible (resp. connexe).
\end{enumerate}
Alors pour que $X'$ soit irréductible (resp. connexe), il faut et il suffit
que $X$ le soit.
\end{lemma}

\begin{proof}
L'assertion relative à la connexité n'est autre que Bourbaki,
\emph{Top. gén.}, chap.~I\textsuperscript{er}, 3\textsuperscript{e} éd.,
\S~11, n\textsuperscript{o}~3, prop.~7. Prouvons l'assertion relative à
l'irréductibilité~; la condition étant trivialement nécessaire puisque $f$
est surjective, prouvons qu'elle est suffisante. Soient donc $X'_1$, $X'_2$
deux parties fermées de $X'$ telles que $X'=X'_1\cup X'_2$, et désignons par
$X_i$ $(i=1,2)$ l'ensemble des $x\in X$ tels que $f^{-1}(x)\subset X'_i$~; on
a donc $X_i=X-f(X'-X'_i)$, et comme $f$ est ouverte, on en conclut que $X_i$
est une partie \emph{fermée} de $X$ $(i=1,2)$. D'autre part, pour tout
$x\in X$, $f^{-1}(x)$ est irréductible par hypothèse, et est réunion des deux
parties fermées $X'_1\cap f^{-1}(x)$ et $X'_2\cap f^{-1}(x)$~; l'une de ces
deux parties doit donc être égale à $f^{-1}(x)$, ce qui signifie que l'on doit
avoir $x\in X_1$ ou $x\in X_2$~; on a donc $X=X_1\cup X_2$, et comme $X$ est
irréductible par hypothèse, on en déduit $X=X_1$ ou $X=X_2$, donc
$X'=X'_1$ ou $X'=X'_2$.
\end{proof}

\begin{remark}[4.4.3]
\label{IV.2.4.4.3-fr}
Si la topologie de $X$ n'est pas quotient de celle de $X'$ par la relation
d'équivalence définie par $f$, il peut se faire que $X$ et toutes les fibres
$f^{-1}(x)$ soient irréductibles, sans que $X'$ soit connexe~: on en a un
exemple en prenant pour $X$ un schéma irréductible (par exemple la droite
affine $\operatorname{Spec}(k[T])$ sur un corps $k$ algébriquement clos), en
considérant un point fermé $x_0$ de $X$ et en prenant pour $X'$ l'espace
\emph{somme} de $\{x_0\}$ et du sous-espace ouvert $X-\{x_0\}$ de $X$. De
même, si l'on remplace l'hypothèse «~$f$ ouverte~» par «~$f$ fermée~» (et même
«~$f$ propre~») (ce qui entraîne pourtant que la topologie de $X$ est alors
quotient de celle de $X'$ par la relation définie par $f$), il peut se faire
que $X$ et toutes les fibres $f^{-1}(x)$ soient irréductibles sans que $X'$ le
soit. Prenons encore pour $X$ la droite affine sur $k$, désignons par $S$ le
produit $X\times_k\mathbf P$, où $\mathbf P$ est la droite projective sur
$k$~; si $x_0$ est un point fermé de $X$, $t_0$ un point fermé de $\mathbf P$,
$p:S\to X$, $q:S\to\mathbf P$ les projections, désignons par $X'$ le
sous-préschéma réduit ayant pour espace sous-jacent l'ensemble fermé
$p^{-1}(x_0)\cup q^{-1}(t_0)$, et prenons pour $f$ la restriction à $X'$ de
$p$~; $f$ est propre mais $X'$ n'est pas irréductible.
\end{remark}

\begin{theorem}[4.4.4]
\label{IV.2.4.4.4-fr}
Soient $k$ un corps algébriquement clos, $X$ un $k$-préschéma. Si $X$ est
irréductible (resp. connexe), il en est de même de $X\otimes_k K$ pour toute
extension $K$ de $k$.
\end{theorem}

\begin{proof}
On a vu en effet (4.4.1) que le morphisme $p:X\otimes_k K\to X$ est
fidèlement plat et ouvert~; pour pouvoir appliquer le lemme (4.4.2), il suffit
donc de vérifier que les fibres $p^{-1}(x)$ sont irréductibles pour tout
$x\in X$. Mais le $k(x)$-préschéma $p^{-1}(x)$ est isomorphe à
$\operatorname{Spec}(k(x)\otimes_k K)$ (I, 3.6.2), donc intègre en vertu de
l'hypothèse sur $k$ et de (4.3.9).
\end{proof}

\begin{corollary}[4.4.5]
\label{IV.2.4.4.5-fr}
Soient $k$ un corps algébriquement clos, $X$ un $k$-préschéma, $K$ une
extension de $k$, $p:X\otimes_k K\to X$ la projection canonique. Si $Z$ est
une partie irréductible (resp. connexe) de $X$, $p^{-1}(Z)$ est irréductible
(resp. connexe)~; en particulier, si $X_0$ est une composante irréductible
(resp. la composante connexe) de $X$ contenant $Z$, $p^{-1}(X_0)$ est une
composante irréductible (resp. connexe) de $X\otimes_k K$ contenant
$p^{-1}(Z)$.
\end{corollary}

\begin{proof}
La seconde assertion résulte de (4.4.1) et de la première appliquée en
remplaçant $Z$ par $X_0$~; pour prouver la première, soient $X'$ l'espace
sous-jacent à $X\otimes_k K$, $Z'=p^{-1}(Z)$~; la

\oldpage[IV-2]{61}

relation d'équivalence $R$ définie par $p$ sur $X'$ est ouverte, donc il en est
de même de la relation $R_{Z'}$ qu'elle induit sur la partie saturée $Z'$, et
$Z$ s'identifie à l'espace quotient $Z'/R_{Z'}$ (Bourbaki,
\emph{Top. gén.}, chap.~I\textsuperscript{er}, 3\textsuperscript{e} éd.,
\S~5, n\textsuperscript{o}~2, prop.~4)~; on peut donc appliquer à $Z$ et $Z'$
le lemme (4.4.2), d'où la première assertion de (4.4.4).
\end{proof}

\begin{corollary}[4.4.6]
\label{IV.2.4.4.6-fr}
Les hypothèses et notations étant celles de (4.4.5), l'application
$Z\mapsto p^{-1}(Z)$ est une bijection de l'ensemble des composantes
irréductibles (resp. connexes) de $X$ sur l'ensemble des composantes
irréductibles (resp. connexes) de $X\otimes_k K$, dont la bijection réciproque
est $Z'\mapsto p(Z')$.
\end{corollary}

\subsection{Préschémas géométriquement irréductibles et géométriquement
connexes}
\label{subsection:IV.2.4.5-fr}

\begin{proposition}[4.5.1]
\label{IV.2.4.5.1-fr}
Soient $k$ un corps, $X$ un $k$-préschéma, $\Omega$ une extension
algébriquement close de $k$. Soit $n$ (resp. $n'$) le cardinal de l'ensemble
des composantes irréductibles (resp. des composantes connexes) de
$X\otimes_k\Omega$~; ce nombre est indépendant de l'extension algébriquement
close $\Omega$ choisie. En outre, pour toute extension $K$ de $k$, le cardinal
de l'ensemble des composantes irréductibles (resp. des composantes connexes)
de $X\otimes_k K$ est $\leqslant n$ (resp. $\leqslant n'$).
\end{proposition}

\begin{proof}
En effet, deux extensions algébriquement closes $\Omega$, $\Omega'$ de $k$
peuvent toujours être considérées comme des sous-extensions d'une même
extension $E$ de $k$~; il suffit alors d'appliquer (4.4.6) à
$X\otimes_k\Omega$ et $X\otimes_k\Omega'$ et à l'extension commune $E$ de
$\Omega$ et $\Omega'$ pour démontrer la première assertion. La seconde
s'obtient en prenant pour $\Omega$ une extension algébriquement close de $K$
et en utilisant (4.4.1).
\end{proof}

\begin{definition}[4.5.2]
\label{IV.2.4.5.2-fr}
Le cardinal $n$ (resp. $n'$) de (4.5.1) s'appelle le \emph{nombre géométrique
de composantes irréductibles} (resp. \emph{de composantes connexes}) de $X$
(relativement à $k$). Si $n=1$ (resp. $n'\leqslant1$) on dit que $X$ est un
$k$-préschéma \emph{géométriquement irréductible} (resp. \emph{géométriquement
connexe}).
\end{definition}

On retrouve ainsi une définition donnée pour un cas particulier dans
(III, 4.3.4). En vertu de (4.5.1), les propriétés suivantes sont
équivalentes~:
\begin{enumerate}
\item[$a)$] $X$ est géométriquement irréductible (resp. géométriquement
connexe).
\item[$b)$] Pour une extension algébriquement close $\Omega$ de $k$,
$X\otimes_k\Omega$ est irréductible (resp. connexe).
\item[$c)$] Pour toute extension $K$ de $k$, $X\otimes_k K$ est irréductible
(resp. connexe).
\end{enumerate}

\begin{env}[4.5.3]
\label{IV.2.4.5.3-fr}
Soient $X$ un $k$-préschéma, $Z$ une partie localement fermée de $X$. Si l'on
désigne encore par $Z$ un quelconque des sous-préschémas de $X$ ayant $Z$ pour
espace sous-jacent (I, 5.2.1), il résulte de (I, 5.1.8) que pour une extension
$K$ de $k$, le nombre de composantes irréductibles (resp. connexes) de
$Z\otimes_k K$ est indépendant du sous-préschéma d'espace sous-jacent $Z$ que
l'on a choisi~; aussi définit-on le nombre géométrique de composantes
irréductibles (resp. de composantes connexes) de $Z$ comme celui de tout
sous-préschéma de $X$ ayant $Z$ pour espace sous-jacent.
\end{env}

\begin{proposition}[4.5.4]
\label{IV.2.4.5.4-fr}
Soient $X$, $Y$ deux $k$-préschémas, $f:X\to Y$ un $k$-morphisme surjectif.
Si $X$ est géométriquement irréductible (resp. géométriquement connexe), il en
est de même de $Y$.
\end{proposition}

\oldpage[IV-2]{62}

\begin{proof}
En effet, pour toute extension $K$ de $k$, $f_{(K)}:X_{(K)}\to Y_{(K)}$ est
surjectif (I, 3.5.2), et $X_{(K)}$ étant irréductible (resp. connexe) par
hypothèse, il en est de même de $Y_{(K)}$.
\end{proof}

\begin{definition}[4.5.5]
\label{IV.2.4.5.5-fr}
On dit qu'un morphisme $f:X\to Y$ de préschémas est \emph{irréductible}
(resp. \emph{connexe}), si, pour tout $y\in Y$, le $k(y)$-préschéma
$f^{-1}(y)$ est géométriquement irréductible (resp. géométriquement connexe).
\end{definition}

\begin{proposition}[4.5.6]
\label{IV.2.4.5.6-fr}
\begin{enumerate}
\item[(i)] Soient $X$ un $k$-préschéma, $K$ une extension de $k$. Alors le
nombre géométrique de composantes irréductibles (resp. connexes) de $X_{(K)}$
relativement à $K$ est égal au nombre géométrique de composantes irréductibles
(resp. connexes) de $X$ relativement à $k$. En particulier, pour que le
$K$-préschéma $X_{(K)}$ soit géométriquement irréductible (resp.
géométriquement connexe), il faut et il suffit que le $k$-préschéma $X$ le
soit.
\item[(ii)] Soient $f:X\to Y$, $g:Y'\to Y$ deux morphismes, et posons
$X'=X_{(Y')}=X\times_Y Y'$, $f'=f_{(Y')}:X'\to Y'$. Si $f$ est irréductible
(resp. connexe), il en est de même de $f'$. La réciproque est vraie si $g$ est
surjectif.
\end{enumerate}
\end{proposition}

\begin{proof}
Pour (i), si $\Omega$ est une extension algébriquement close de $K$, on a
$X\otimes_k\Omega=X_{(K)}\otimes_K\Omega$, d'où l'assertion.

Pour (ii), pour tout $y'\in Y'$, si l'on pose $y=g(y')$, on a
$f'^{-1}(y')=f^{-1}(y)\otimes_{k(y)}k(y')$ (I, 3.6.4), donc les deux assertions
résultent de (i).
\end{proof}

\begin{proposition}[4.5.7]
\label{IV.2.4.5.7-fr}
Soient $f:X\to Y$ un morphisme surjectif irréductible (resp. connexe),
$g:Y'\to Y$ un morphisme, et posons $X'=X\times_Y Y'$. Si en outre $Y'$ est
irréductible (resp. connexe), et $f$ universellement ouvert (resp. plat et
quasi-compact, ou universellement ouvert, ou universellement fermé), alors
$X'$ est irréductible (resp. connexe).
\end{proposition}

\begin{proof}
Les propriétés de $f$ que l'on considère étant toutes stables par changement
de base, on peut se borner au cas où $Y'=Y$~; il suffit alors d'appliquer
(4.4.2) (compte tenu de (2.3.12)).
\end{proof}

\begin{corollary}[4.5.8]
\label{IV.2.4.5.8-fr}
Soient $X$, $Y$ deux $k$-préschémas.
\begin{enumerate}
\item[(i)] Si $X$ est géométriquement irréductible (resp. géométriquement
connexe) et $Y$ irréductible (resp. connexe), alors $X\times_k Y$ est
irréductible (resp. connexe).
\item[(ii)] Si $X$ et $Y$ sont géométriquement irréductibles (resp.
géométriquement connexes), il en est de même de $X\times_k Y$.
\end{enumerate}
\end{corollary}

\begin{proof}
Pour (i), le morphisme structural $X\to\operatorname{Spec}(k)$ est surjectif et
universellement ouvert (2.4.9) et en outre irréductible (resp. connexe)~; il
suffit donc d'appliquer (4.5.7).

Pour (ii), si $\Omega$ est une extension algébriquement close de $k$, on a
$(X\times_k Y)_{(\Omega)}=X_{(\Omega)}\times_\Omega Y_{(\Omega)}$
(I, 3.3.10)~; par hypothèse $X_{(\Omega)}$ et $Y_{(\Omega)}$ sont des
$\Omega$-préschémas géométriquement irréductibles (resp. géométriquement
connexes) (4.5.6) et il suffit donc d'appliquer (i).
\end{proof}

\begin{proposition}[4.5.9]
\label{IV.2.4.5.9-fr}
Soit $X$ un $k$-préschéma. Les conditions suivantes sont équivalentes~:
\begin{enumerate}
\item[$a)$] $X$ est géométriquement irréductible (autrement dit, pour toute
extension $K$ de $k$, $X\otimes_k K$ est irréductible).
\item[$b)$] Pour toute extension finie séparable $K$ de $k$, $X\otimes_k K$
est irréductible.
\item[$c)$] $X$ est irréductible, et si $x$ est son point générique, $k(x)$
est une extension primaire de $k$.
\end{enumerate}
\end{proposition}

\begin{proof}
Il est clair que $a)$ entraîne $b)$. Pour voir que $b)$ entraîne $c)$,
considérons une extension finie séparable $K$ de $k$~; si
$p:X\otimes_k K\to X$ est la projection canonique, les

\oldpage[IV-2]{63}

points maximaux de $X\otimes_k K$ sont ceux de la fibre
$p^{-1}(x)=\operatorname{Spec}(k(x)\otimes_k K)$ ((2.3.4) et
$(0_{\mathrm I},2.1.8)$)~; dire que cette fibre est irréductible pour toute
extension finie séparable $K$ de $k$ revient à dire que $k(y)$ est une
extension primaire de $k$, en vertu de (4.3.2). Inversement, si $c)$ est
vérifiée, le même raisonnement (compte tenu de (4.3.2)) montre que pour toute
extension $K$ de $k$, $X\otimes_k K$ est irréductible, donc $c)$ implique
$a)$.
\end{proof}

\begin{corollary}[4.5.10]
\label{IV.2.4.5.10-fr}
Soient $X$ un $k$-préschéma irréductible, $x$ son point générique, $k'$ la
fermeture algébrique séparable de $k$ dans $k(x)$, $k''$ une extension
galoisienne de $k$ (de degré fini ou non) contenant $k'$. On suppose $k'$ fini
sur $k$~; alors les composantes irréductibles de $X\otimes_k k''$ sont
géométriquement irréductibles~; leur nombre est égal à $[k':k]$ et est aussi le
nombre géométrique de composantes irréductibles de $X$.
\end{corollary}

\begin{proof}
Comme dans (4.5.9), on est ramené à considérer les points maximaux de la fibre
$\operatorname{Spec}(k(x)\otimes_k k'')$~; en vertu de (4.3.4), ils sont au
nombre de $[k':k]$ et les corps résiduels en ces points (qui sont aussi ceux de
$X\otimes_k k''$) sont des extensions primaires de $k''$, ce qui, compte tenu
de (4.5.9), prouve le corollaire.
\end{proof}

\begin{corollary}[4.5.11]
\EGAReaderTarget{ega4.c04.cor.4-5-11}{IV.2.4.5.11-fr}
Soit $X$ un $k$-préschéma~; supposons que $X$ n'ait qu'un nombre fini de
points maximaux $x_i$ $(1\leqslant i\leqslant r)$ et que pour chaque $i$, la
fermeture algébrique séparable $k'_i$ de $k$ dans $k(x_i)$ soit de degré fini
sur $k$. Alors il existe une extension finie séparable $L$ de $k$ telle que les
composantes irréductibles de $X\otimes_k L$ soient géométriquement
irréductibles~; leur nombre, égal à $\sum_i[k'_i:k]$, est le nombre géométrique
des composantes irréductibles de $X$.
\end{corollary}

\begin{proof}
Les points maximaux de $X\otimes_k L$ sont ceux des diverses fibres
$\operatorname{Spec}(k(x_i)\otimes_k L)$ en vertu de (2.3.4) et
$(0_{\mathrm I},2.1.8)$ appliqué aux composantes irréductibles de $X$ et à
leurs images réciproques dans $X\otimes_k L$. Il suffit donc de prendre pour
$L$ une extension galoisienne finie de $k$ contenant tous les $k'_i$ et
d'appliquer (4.5.10).
\end{proof}

On notera que (4.5.11) s'applique en particulier à tout $k$-préschéma $X$ de
type fini sur $k$, car $X$ est alors noethérien, donc n'a qu'un nombre fini de
composantes irréductibles, et chacun des $k(x_i)$ est une extension de type
fini de $k$, donc il en est de même de la fermeture algébrique de $k$ dans
$k(x_i)$ ([1], p.~6-06, lemme~5).

\begin{remark}[4.5.12]
\label{IV.2.4.5.12-fr}
\begin{enumerate}
\item[(i)] Les notions définies dans (4.5.2) dépendent du corps de base $k$.
Pour abréger, et lorsqu'il sera nécessaire de mentionner le corps de base, on
pourra dire «~$k$-irréductible~» (resp. «~$k$-connexe~») au lieu de
«~géométriquement irréductible (resp. connexe) relativement à $k$~». On notera
que cette terminologie, tout à fait naturelle dans le contexte des schémas,
est exactement à l'opposé de celle de Weil~: chez cet auteur,
«~$k$-irréductible~» (resp. «~$k$-connexe~», «~$k$-normal~», etc.) désigne une
propriété \emph{intrinsèque} d'un préschéma $X$, indépendante du corps $k$ pris
comme «~corps de base~» (autrement dit, du morphisme
$X\to\operatorname{Spec}(k)$ choisi)~; nous exprimons ces propriétés en disant
simplement que $X$ est irréductible (resp. connexe, normal, etc.). Weil dit
d'autre part «~absolument irréductible~», «~absolument connexe~», «~absolument
normal~», etc., pour les notions correspondantes relatives au préschéma
$X\otimes_k\Omega$, où $\Omega$ est une extension algébriquement close
convenable de $k$. Cette terminologie s'explique par le point de vue différent
auquel se place cet auteur~: pour lui, une «~variété algébrique~» $V$ est
donnée d'abord comme un objet

\oldpage[IV-2]{64}

géométrique au-dessus d'un corps algébriquement clos $\Omega$~; la donnée d'un
«~corps de définition~» $k$ plus petit (i.e., d'un schéma sur $k$ définissant
$V$ par extension du corps de base à $\Omega$) est pour lui une structure
supplémentaire et, dans une certaine mesure, secondaire, de sorte que sa
qualification d'«~absolu~» ou «~intrinsèque~» d'un côté, «~relatif à $k$~» de
l'autre, est à l'opposé de la nôtre\footnote{Le point de vue de Zariski est
déjà plus proche du nôtre, et sa terminologie n'est généralement pas en conflit
avec celle introduite ici.}. Nous éviterons donc par la suite d'utiliser le
qualificatif «~absolu~», qui peut prêter à confusion, et lorsque nous
utiliserons la terminologie abrégée introduite plus haut (en conflit avec la
terminologie reçue), nous renverrons à (4.5.12) pour éviter toute ambiguïté.

\item[(ii)] La proposition (4.5.9) donne un critère «~birationnel~» (autrement
dit, ne dépendant que du corps résiduel au point générique) pour qu'un
$k$-préschéma $X$ soit géométriquement irréductible. Il n'y a pas de critère
analogue pour que $X$ soit géométriquement connexe. Par exemple, prenons
$X=\operatorname{Spec}(\mathbf R[[T,U]]/(T^2+U^2))$ ($T,U$ indéterminées)~;
$X$ est un $\mathbf R$-schéma intègre, et son corps de fonctions rationnelles
$k(x)$ est $\mathbf C((T))$, comme on le vérifie aisément~; $X\otimes_{\mathbf
R}\mathbf C$ se décompose en deux composantes irréductibles (les deux
«~droites isotropes~») qui ont un point commun (l'idéal maximal image dans
$\mathbf C[[T,U]]/(T^2+U^2)$ de l'idéal maximal $(T)+(U)$ de
$\mathbf C[[T,U]]$)~; donc $X\otimes_{\mathbf R}\mathbf C$ est connexe. Soit
$X'$ le complémentaire (ouvert) du point fermé de $X$ correspondant à l'idéal
maximal $(T)+(U)$ de $\mathbf R[[T,U]]$~; alors $X'\otimes_{\mathbf R}\mathbf C$
n'est pas connexe, bien que $X$ et $X'$ aient même corps de fonctions
rationnelles.
\end{enumerate}
\end{remark}

\begin{proposition}[4.5.13]
\label{IV.2.4.5.13-fr}
Soient $X$, $Y$ deux $k$-préschémas, $f:Y\to X$ un $k$-morphisme. Supposons
$Y$ géométriquement connexe et non vide, et $X$ connexe. Alors $X$ est
géométriquement connexe.
\end{proposition}

\begin{proof}
Soient $\Omega$ une clôture algébrique de $k$, $X'=X_{(\Omega)}$,
$Y'=Y_{(\Omega)}$, $f'=f_{(\Omega)}:Y'\to X'$~; il faut montrer que $X'$ est
connexe. Soit $U'$ une partie ouverte et fermée non vide de $X'$~; notons que
les morphismes $p:X'\to X$ et $q:Y'\to Y$ sont \emph{ouverts} (2.4.10) et
\emph{fermés} puisque $\Omega$ est extension algébrique de $k$ (II, 6.1.10).
Donc $U=p(U')$ est ouvert et fermé non vide dans $X$, et par suite égal à $X$.
Comme $Y$ est non vide, on en conclut (I, 3.4.7) que
$V'=f'^{-1}(U')$ n'est pas vide~; d'ailleurs, c'est une partie ouverte et
fermée de $Y'$, et ce dernier espace est connexe par hypothèse~; donc $V'=Y'$.
On en déduit que $U'=X'$, sans quoi le même raisonnement appliqué à l'ensemble
ouvert et fermé $X'-U'$, montrerait que $f'^{-1}(X'-U')=Y'$, ce qui est absurde
puisque $Y'$ n'est pas vide.
\end{proof}

\begin{corollary}[4.5.13.1]
\label{IV.2.4.5.13.1-fr}
\begin{enumerate}
\item[(i)] Soient $X$, $Y$ deux $k$-préschémas, $f:Y\to X$ un $k$-morphisme.
Si $Y$ est géométriquement connexe et non vide, la composante connexe $X_0$ de
$X$ contenant $f(Y)$ est géométriquement connexe.
\item[(ii)] Soit $X$ un $k$-préschéma. Si $Y$ est une composante irréductible
de $X$ qui est géométriquement irréductible, alors la composante connexe $X_0$
de $X$ contenant $Y$ est géométriquement connexe.
\end{enumerate}
\end{corollary}

\begin{proof}
Pour (i), on peut supposer $Y$ réduit, de sorte que $f$ se factorise en
$Y\xrightarrow{g}X_0\xrightarrow{j}X$, $X_0$

\oldpage[IV-2]{65}

désignant le sous-préschéma réduit de $X$ ayant $X_0$ pour espace sous-jacent
(I, 5.2.2)~; il suffit d'appliquer (4.5.13) à $g$.

Pour (ii), en considérant un préschéma ayant $Y$ pour espace sous-jacent, et
remarquant que $Y$ est \emph{a fortiori} géométriquement connexe, il suffit
d'appliquer (i) à l'injection canonique $Y\to X$.
\end{proof}

\begin{corollary}[4.5.14]
\label{IV.2.4.5.14-fr}
Soit $X$ un $k$-préschéma. Si $x$ est un point de $X$ tel que $k(x)$ soit une
extension primaire de $k$ (en particulier si $x$ est un point rationnel de
$X$), la composante connexe de $x$ dans $X$ est géométriquement connexe.
\end{corollary}

\begin{proof}
Il suffit d'appliquer (4.5.13.1) à $Y=\operatorname{Spec}(k(x))$ et au
morphisme canonique $Y\to X$, en tenant compte de (4.5.9), qui entraîne que
$Y$ est géométriquement irréductible.
\end{proof}

\begin{proposition}[4.5.15]
\label{IV.2.4.5.15-fr}
Soient $X$ un $k$-préschéma, $x$ un point de $X$, $k'$ la fermeture algébrique
séparable de $k$ dans $k(x)$. On suppose vérifiées les conditions suivantes~:
\begin{enumerate}
\item[(i)] $X$ est connexe.
\item[(ii)] $k'$ est une extension finie de $k$.
\end{enumerate}
Alors le nombre géométrique des composantes connexes de $X$ est
$\leqslant[k':k]$, et si $k''$ est une extension galoisienne finie de $k$
contenant $k'$, les composantes connexes de $X\otimes_k k''$ sont
géométriquement connexes.
\end{proposition}

On notera que la condition (ii) est vérifiée si $X$ est \emph{localement de
type fini} sur $k$.

\begin{proof}
Il existe des extensions galoisiennes finies $k''$ de $k$ contenant $k'$ en
vertu de la condition (ii). Posons $X''=X\otimes_k k''$~; le morphisme
projection $p:X''\to X$ est fidèlement plat et fini, donc fermé (II, 6.1.10) et
par suite (2.3.6, (ii)) l'image par $p$ de toute composante connexe $X''_\alpha$
de $X''$ est égale à $X$~; $X''_\alpha$ contient donc un point
$x''_\alpha\in p^{-1}(x)$. Mais $p^{-1}(x)$ a un nombre de points égal au
nombre des composantes irréductibles de
$\operatorname{Spec}(k(x)\otimes_k k'')$ (I, 3.4.9), c'est-à-dire à $[k':k]$
(4.3.4), et \emph{a fortiori} le nombre des composantes connexes de $X''$ est
$\leqslant[k':k]$. Pour tout $x''_\alpha\in p^{-1}(x)$, $k(x''_\alpha)$ est une
extension primaire de $k''$ en vertu de (4.3.5) et de (I, 3.4.9). On peut donc
appliquer à $x''_\alpha$ et $X''$ le corollaire (4.5.14), qui prouve que toutes
les composantes connexes de $X''$ sont géométriquement connexes.
\end{proof}

\begin{corollary}[4.5.16]
\label{IV.2.4.5.16-fr}
Soient $X$ un $k$-préschéma, $(X_\alpha)$ la famille de ses composantes
connexes, et pour tout $\alpha$, soit $x_\alpha\in X_\alpha$. On suppose
vérifiées les conditions suivantes~:
\begin{enumerate}
\item[(i)] La famille $(X_\alpha)$ est finie.
\item[(ii)] Pour tout $\alpha$, la fermeture algébrique séparable $k'_\alpha$
de $k$ dans $k(x_\alpha)$ est une extension finie de $k$.
\end{enumerate}
Alors le nombre géométrique de composantes connexes de $X$ est au plus
$\sum_\alpha[k'_\alpha:k]$, et il existe une extension finie séparable $k''$ de
$k$ telle que toutes les composantes connexes de $X\otimes_k k''$ soient
géométriquement connexes.
\end{corollary}

\begin{proof}
Notant que les composantes connexes de $X$ sont ouvertes en vertu de la
condition (i), il suffit d'appliquer à chacun des sous-préschémas de $X$ induits
sur les ouverts $X_\alpha$ le résultat de (4.5.15)~; pour avoir une extension
$k''$ répondant à la question, il suffit de prendre une extension galoisienne
finie de $k$ contenant tous les $k'_\alpha$.
\end{proof}

\oldpage[IV-2]{66}

\begin{corollary}[4.5.17]
\label{IV.2.4.5.17-fr}
Supposons que le $k$-préschéma $X$ contienne un point $x$ tel que la fermeture
algébrique séparable de $k$ dans $k(x)$ soit finie sur $k$. Pour que $X$ soit
géométriquement connexe, il faut et il suffit que pour toute extension finie
séparable $K$ de $k$, $X\otimes_k K$ soit connexe.
\end{corollary}

\begin{remark}[4.5.18]
\label{IV.2.4.5.18-fr}
Nous verrons au \S~8 (8.4.5) que la conclusion de (4.5.17) est encore valable
si, au lieu de supposer vérifiée la condition de l'énoncé, on suppose que $X$
est \emph{quasi-compact}.
\end{remark}

\begin{proposition}[4.5.19]
\EGAReaderTarget{ega4.c04.prop.4-5-19}{IV.2.4.5.19-fr}
Soient $X$ un $k$-préschéma, $Z$ une partie de $X$, $k'$ une extension
algébriquement close de $k$, $X'=X\otimes_k k'$, $p:X'\to X$ la projection
canonique. On suppose que $Z'=p^{-1}(Z)$ est contenu dans une seule composante
irréductible $X'_0$ de $X'$. Alors on a $X'_0=p^{-1}(X_0)$, où
$X_0=p(X'_0)$ est une composante irréductible de $X$ contenant $Z$, et en outre
$X_0$ est géométriquement irréductible. Si de plus $X'_0$ est la seule
composante irréductible de $X'$ qui rencontre $Z'$, alors $X_0$ est la seule
composante irréductible de $X$ qui rencontre $Z$.
\end{proposition}

Pour montrer que $X'_0=p^{-1}(X_0)$, nous allons appliquer le

\begin{lemma}[4.5.19.1]
\label{IV.2.4.5.19.1-fr}
Soient $T$, $T'$ deux préschémas, $p:T'\to T$ un morphisme~; considérons le
préschéma $T''=T'\times_T T'$, et soient $p_1:T''\to T'$, $p_2:T''\to T'$ les
projections canoniques. Pour qu'une partie $U'$ de $T'$ soit de la forme
$p^{-1}(U)$, où $U\subset T$, il faut et il suffit que l'on ait
$p_1^{-1}(U')=p_2^{-1}(U')$.
\end{lemma}

\begin{proof}
Le morphisme structural $q:T''\to T$ est égal à $p\circ p_1$ et à
$p\circ p_2$ par définition, donc la relation est nécessaire. Inversement,
supposons-la vérifiée~; soient $u'$ un point de $U'$, $t'$ un point de
$p^{-1}(p(u'))$~; comme $p(t')=p(u')$, il existe un point $t''\in T''$ tel que
$p_1(t'')=u'$, $p_2(t'')=t'$ (I, 3.4.7)~; donc
$t''\in p_1^{-1}(U')=p_2^{-1}(U')$ par hypothèse, et par suite
$t'=p_2(t'')\in U'$, ce qui prouve le lemme.
\end{proof}

\begin{proof}
Ce lemme étant établi, pour démontrer que $X'_0=p^{-1}(X_0)$, formons donc le
produit $X''=X'\times_X X'$ (relatif au morphisme $p:X'\to X$)~; si $p_1$ et
$p_2$ sont les projections canoniques de $X''$ dans $X'$, il s'agit donc de
montrer que $p_1^{-1}(X'_0)=p_2^{-1}(X'_0)$. Or, si l'on pose
$S=\operatorname{Spec}(k)$, $S'=\operatorname{Spec}(k')$, observons que si
$S''=S'\times_S S'=\operatorname{Spec}(k'\otimes_k k')$, on peut aussi écrire
$X''=X'\times_{S'}S''$. Si $\varphi'':X''\to S''$ est le morphisme structural,
il suffit de prouver que pour tout $s''\in S''$, on a
\[
p_1^{-1}(X'_0)\cap\varphi''^{-1}(s'')
=p_2^{-1}(X'_0)\cap\varphi''^{-1}(s'').
\]
Si l'on pose $T=\operatorname{Spec}(k(s''))$,
$U=X''\times_{S''}T=\varphi''^{-1}(s'')$, et si $v:U\to X''$ est la projection
canonique, il s'agit donc de prouver que
$U_1=v^{-1}(p_1^{-1}(X'_0))$ et $U_2=v^{-1}(p_2^{-1}(X'_0))$, qui sont des
composantes irréductibles de $U$ (4.4.5) sont identiques. On a le diagramme
\begin{equation}
\tag{4.5.19.2}
\label{IV.2.4.5.19.2-fr}
\xymatrix@C=3.2em@R=3.0em{
X\ar[d]
  & X'\ar[l]_p\ar[d]
  & X''\ar@<.55ex>[l]^{p_2}\ar@<-.55ex>[l]_{p_1}
       \ar[d]^{\varphi''}
  & U\ar[l]_v\ar[d]
\\
S
  & S'\ar[l]
  & S''\ar@<.55ex>[l]\ar@<-.55ex>[l]
  & T\ar[l].
}
\end{equation}
Par construction, $U_1$ et $U_2$ contiennent toutes deux l'ensemble
$w^{-1}(Z)$, où
\[
w=p\circ p_1\circ v=p\circ p_2\circ v.
\]

\oldpage[IV-2]{67}

Or, il existe un corps $K$, extension commune de $k'$ et de $k(s'')$, de sorte
que si l'on pose $P=\operatorname{Spec}(K)$, le diagramme
\[
\xymatrix@C=3.5em@R=3em{
S'\ar[d] & P\ar[l]\ar[d]\\
S & T\ar[l]
}
\]
soit commutatif~; si l'on pose $Q=X\otimes_k K$, le diagramme
\[
\xymatrix@C=3.5em@R=3em{
X'\ar[d]_p & Q\ar[l]_q\ar[d]^r\\
X & U\ar[l]^w
}
\]
est donc lui aussi commutatif, et les deux composantes irréductibles
$r^{-1}(U_1)$ et $r^{-1}(U_2)$ de $Q$ (4.4.5) contiennent par suite
$q^{-1}(Z')$. Or, en vertu de (4.4.5), les images
$q(r^{-1}(U_1))$ et $q(r^{-1}(U_2))$ sont des composantes irréductibles de $X'$
contenant $Z'$~; elles sont donc toutes deux égales à $X'_0$, donc
$r^{-1}(U_1)=r^{-1}(U_2)$ en vertu de (4.4.5), et finalement $U_1=U_2$.

Comme $p$ est surjectif et que $X'_0$ domine une composante irréductible de
$X$, $X_0=p(X'_0)$ est une composante irréductible de $X$ contenant $Z$, et
comme $p^{-1}(X_0)$ est une composante irréductible de $X'$, $X_0$ est
géométriquement irréductible.

Pour prouver la dernière assertion, soit $X_1$ une seconde composante
irréductible de $X$ rencontrant $Z$. Prenant $z\in Z\cap X_1$, on voit qu'il
existe une composante irréductible de $X'$ rencontrant $Z'$ et dominant $X_1$
(2.3.5)~; mais comme $X'_0$ est par hypothèse la seule composante irréductible
de $X'$ rencontrant $Z'$, on a nécessairement $X_1=X_0$.
\end{proof}

\begin{remark}[4.5.20]
\label{IV.2.4.5.20-fr}
Le fait que $Z'$ soit contenu dans une seule composante irréductible $X'_0$ de
$X'$ \emph{n'entraîne pas} que $X_0=p(X'_0)$ soit la seule composante
irréductible de $X$ contenant $Z$. Prenons en effet $k=\mathbf R$,
$k'=\mathbf C$, et pour $X$ le $k$-schéma obtenu par recollement\footnote{La
technique de «~recollement~» des préschémas sera développée en détail au
chap.~V~; pour le cas considéré ici (où il s'agit de courbes algébriques), on
peut consulter [38], p.~68--71.} en un point (non générique) de
$X_1=\operatorname{Spec}(\mathbf R[S,T]/(S^2+T^2+1))$ et de
$X_2=\operatorname{Spec}(\mathbf C[T])$ (les anneaux locaux de ces deux courbes
en chacun de leurs points fermés ayant même corps résiduel isomorphe à
$\mathbf C$)~; $X_1$ et $X_2$ s'identifient aux deux composantes irréductibles
de $X$~; si $x$ est leur point commun, $p^{-1}(x)$ est formé de deux points
distincts $y'$, $z'$, $p^{-1}(X_1)$ est irréductible, et $p^{-1}(X_2)$ est somme
de deux composantes irréductibles $Y'$, $Z'$ de $X'$ telles que $y'\in Y'$ et
$z'\in Z'$, de sorte que $p^{-1}(x)$ n'est contenu que dans une seule
composante irréductible de $X'$.
\end{remark}

\begin{proposition}[4.5.21]
\label{IV.2.4.5.21-fr}
Soit $k$ un corps séparablement clos, et soit $k'$ une clôture algébrique de
$k$. Pour tout $k$-préschéma $X$, la projection canonique
$X\otimes_k k'\to X$ est un homéomorphisme universel.

\oldpage[IV-2]{68}

En particulier les composantes irréductibles (resp. connexes) de $X$ sont
géométriquement irréductibles (resp. géométriquement connexes).
\end{proposition}

\begin{proof}
Par définition, $k'$ est une extension radicielle de $k$~; le morphisme
$\operatorname{Spec}(k')\to\operatorname{Spec}(k)$ est entier, surjectif et
radiciel, donc la première assertion résulte de (2.4.5, (i))~; la seconde en
résulte, vu (4.5.1).
\end{proof}

\subsection{Préschémas algébriques géométriquement réduits}
\label{subsection:IV.2.4.6-fr}

\begin{proposition}[4.6.1]
\label{IV.2.4.6.1-fr}
Soient $k$ un corps, $X$ un $k$-préschéma, $\Omega$ une extension parfaite de
$k$. Les conditions suivantes sont équivalentes~:
\begin{enumerate}
\item[a)] Pour tout $k$-préschéma réduit $S$, $X\times_k S$ est réduit.
\item[b)] Pour toute extension $K$ de $k$, $X\otimes_k K$ est réduit.
\item[c)] $X\otimes_k\Omega$ est réduit.
\item[d)] Pour toute extension radicielle finie $k'$ de $k$, $X\otimes_k k'$ est
réduit.
\item[e)] $X$ est réduite, et pour toute composante irréductible $X_\alpha$ de
$X$, de point générique $x_\alpha$, $k(x_\alpha)$ est une extension séparable de
$k$.
\end{enumerate}
\end{proposition}

\begin{proof}
Il est trivial que a) $\Rightarrow$ b) $\Rightarrow$ c)~; comme $k'$ peut être
considérée comme sous-extension de $\Omega$, on a vu (4.4.1) que c) implique
d). Montrons que d) implique e). En prenant $k'=k$, on voit d'abord que $X$
est réduit~; pour prouver que $k(x_\alpha)$ est séparable sur $k$, il suffit de
montrer que $k(x_\alpha)\otimes_k k'$ est un anneau réduit pour toute extension
finie radicielle $k'$ de $k$ (4.3.5). On peut se borner au cas où
$X=\operatorname{Spec}(A)$ est affine, $A$ étant une $k$-algèbre réduite~; alors
par hypothèse l'anneau $A\otimes_k k'$ est réduit (I, 5.1.4) et
$k(x_\alpha)\otimes_k k'$ est un anneau de fractions de cet anneau, donc il est
réduit $(0_{\mathrm I},1.2.8)$. Enfin, pour prouver que e) entraîne a), on peut
se borner au cas où $X=\operatorname{Spec}(A)$, $S=\operatorname{Spec}(B)$ sont
affines, $A$ et $B$ étant des $k$-algèbres. Les $k(x_\alpha)$ sont les corps des
fractions des anneaux quotients $A/\mathfrak p_\alpha$, où
$\mathfrak p_\alpha=\mathfrak j_{x_\alpha}$ sont les idéaux premiers minimaux de
$A$~; comme par hypothèse $\bigcap_\alpha\mathfrak p_\alpha=0$, $A$ est contenu
dans le produit $\prod_\alpha A/\mathfrak p_\alpha$, donc dans
$\prod_\alpha k(x_\alpha)$~; on peut de même considérer $B$ comme une
sous-algèbre d'un produit $\prod_\beta L_\beta$ d'extensions de $k$. Il en
résulte que $A\otimes_k B$ s'identifie à une sous-algèbre de
\[
\left(\prod_\alpha k(x_\alpha)\right)\otimes_k
\left(\prod_\beta L_\beta\right),
\]
et ce produit tensoriel s'identifie lui-même à une sous-algèbre de
\[
\prod_{\alpha,\beta}\bigl(k(x_\alpha)\otimes_k L_\beta\bigr)
\]
(Bourbaki, \emph{Alg.}, chap.~II, 3\textsuperscript{e} éd., \S~7,
n\textsuperscript{o}~7, prop.~15). Or par hypothèse les
$k(x_\alpha)\otimes_k L_\beta$ sont réduits (4.3.5), donc il en est de même de
$A\otimes_k B$.
\end{proof}

\begin{definition}[4.6.2]
\label{IV.2.4.6.2-fr}
Si les conditions équivalentes a) à e) de la proposition (4.6.1) sont
remplies, on dit que $X$ est \emph{séparable} (ou \emph{géométriquement réduit},
ou \emph{universellement réduit}) sur $k$. On dit qu'un $k$-préschéma $X$ est
\emph{géométriquement intègre} sur $k$ si, pour toute extension $K$ de $k$,
$X\otimes_k K$ est intègre~; il revient au même (en vertu de (4.6.1)) de dire
que $X$ est séparable et géométriquement irréductible sur $k$.
\end{definition}

On dira qu'une $k$-algèbre $A$ (commutative) est \emph{séparable} si
$\operatorname{Spec}(A)$ est séparable sur $k$~; cela signifie donc que pour
toute extension $K$ de $k$, l'anneau $A_{(K)}=A\otimes_k K$ est réduit. On
notera que cette définition coïncide avec celle de Bourbaki, \emph{Alg.},
chap.~VIII,

\oldpage[IV-2]{69}

\S~7, n\textsuperscript{o}~5, déf.~1, lorsque $A$ est de rang \emph{fini} sur
$k$ (\emph{loc. cit.}, cor. de la prop.~7), mais non en général, une $k$-algèbre
pouvant avoir un radical $\neq0$ même si elle est intègre.

\begin{corollary}[4.6.3]
\label{IV.2.4.6.3-fr}
Soit $X$ un $k$-préschéma intègre~; pour que $X$ soit géométriquement réduit
(resp. géométriquement intègre) sur $k$, il faut et il suffit que son corps de
fonctions rationnelles $\mathrm R(X)$ soit une extension séparable (resp.
séparable et primaire) de $k$.
\end{corollary}

\begin{proof}
Cela résulte aussitôt de (4.5.9) et (4.6.1).
\end{proof}

\begin{corollary}[4.6.4]
\label{IV.2.4.6.4-fr}
Soit $X$ un $k$-préschéma réduit. Alors, pour toute extension séparable $k'$ de
$k$, $X'=X\otimes_k k'$ est réduit.
\end{corollary}

\begin{proof}
Il suffit d'appliquer l'équivalence de a) et e) dans (4.6.1), en remplaçant
$X$ par $\operatorname{Spec}(k')$ et $S$ par $X$.
\end{proof}

\begin{proposition}[4.6.5]
\label{IV.2.4.6.5-fr}
\begin{enumerate}
\item[(i)] Soient $X$ un $k$-préschéma, $K$ une extension de $k$. Pour que $X$
soit géométriquement réduit (resp. géométriquement intègre) sur $k$, il faut et
il suffit que $X\otimes_k K$ soit géométriquement réduit (resp. géométriquement
intègre) sur $K$.
\item[(ii)] Soient $X$, $Y$ deux $k$-préschémas. Si $X$ et $Y$ sont
géométriquement réduits (resp. géométriquement intègres) sur $k$, il en est de
même de $X\times_k Y$.
\end{enumerate}
\end{proposition}

\begin{proof}
L'assertion (i) est conséquence triviale des définitions et de (4.4.1). Pour
démontrer la partie de l'assertion (ii) concernant la séparabilité, on observe
que si $\Omega$ est une extension algébriquement close de $k$, on a
\[
(X\times_k Y)\otimes_k\Omega=X\times_k(Y\otimes_k\Omega)~;
\]
comme $Y\otimes_k\Omega$ est réduit, il en est de même de
$X\times_k(Y\otimes_k\Omega)$ en vertu de (4.6.1, a)). Le reste de l'assertion
(ii) résulte de ce qui précède et de (4.5.8, (ii)).
\end{proof}

\begin{proposition}[4.6.6]
\label{IV.2.4.6.6-fr}
Soit $X$ un $k$-préschéma de type fini. Il existe une extension radicielle finie
$k'$ de $k$ telle que $(X_{(k')})_{\mathrm{red}}$ soit géométriquement réduit
sur $k'$.
\end{proposition}

\begin{proof}
Comme $X$ peut être recouvert par un nombre fini d'ouverts affines $U_i$, on
peut se borner au cas où $X$ est affine~: en effet, si pour chaque $i$, $k_i$
est une extension radicielle finie de $k$ telle que
$(U_i\otimes_k k_i)_{\mathrm{red}}$ soit géométriquement réduit sur $k_i$, on
peut supposer que les $k_i$ sont contenus dans une même extension radicielle
finie $k'$ de $k$~; $(U_i\otimes_k k')_{\mathrm{red}}$ s'identifie à
\[
(U_i\otimes_k k_i)_{\mathrm{red}}\otimes_{k_i}k'
\]
en vertu de (I, 5.1.8), donc est géométriquement réduit sur $k'$ par hypothèse,
et il en est donc de même de $(X\otimes_k k')_{\mathrm{red}}$. Supposons donc
que $X=\operatorname{Spec}(A)$, où $A$ est une $k$-algèbre de type fini, et
soit $\Omega$ une clôture algébrique de $k$. Désignons par $\mathfrak N'$ le
nilradical de $A\otimes_k\Omega$~; comme $A\otimes_k\Omega$ est un anneau
noethérien, $\mathfrak N'$ est un idéal engendré par un nombre fini d'éléments
de la forme
\[
y_i=\sum_j x_{ij}\otimes\xi_{ij},
\qquad x_{ij}\in A,\quad \xi_{ij}\in\Omega.
\]
Soient $K$ une sous-extension finie de $\Omega$ contenant les $\xi_{ij}$, et
$\mathfrak N$ l'idéal de $A\otimes_k K$ engendré par les $y_i$
($A\otimes_k K$ étant identifié à un sous-anneau de $A\otimes_k\Omega$)~; il
est clair que les $y_i$ sont nilpotents dans $A\otimes_k K$~; d'autre part,
comme $\mathfrak N\otimes_K\Omega=\mathfrak N'$ (et par suite
$\mathfrak N'\cap(A\otimes_k K)=\mathfrak N$), $\mathfrak N$ contient le
nilradical de $A\otimes_k K$, donc il est égal à ce nilradical. Comme
\[
(A\otimes_k\Omega)/\mathfrak N'
=\bigl((A\otimes_k K)/\mathfrak N\bigr)\otimes_K\Omega,
\]
on voit que $(X\otimes_k K)_{\mathrm{red}}\otimes_K\Omega$ est réduit, et par
suite (4.6.1), $(X_{(K)})_{\mathrm{red}}$ est géométriquement réduit sur $K$.
En remplaçant $K$ par l'extension quasi-galoisienne sur $k$ qu'il engendre dans
$\Omega$, on voit par (I, 5.1.8) que l'on peut supposer en outre $K$
quasi-galoisienne sur $k$, donc extension séparable d'une extension radicielle
finie $k'$ de $k$ (Bourbaki, \emph{Alg.}, chap.~V, \S~10,
n\textsuperscript{o}~9, prop.~14). En vertu

\oldpage[IV-2]{70}

de (I, 5.1.8), $(X_{(k')})_{\mathrm{red}}\otimes_{k'}K$ est isomorphe à
$(X_{(K)})_{\mathrm{red}}$, donc est géométriquement réduit~; il en est alors de
même de $(X_{(k')})_{\mathrm{red}}$ en vertu de (4.6.5, (i)).
\end{proof}

\begin{corollary}[4.6.7]
\label{IV.2.4.6.7-fr}
Si $K$ est une extension de type fini de $k$, il existe une extension radicielle
finie $k'$ de $k$ telle que les corps résiduels de l'anneau semi-local
$K\otimes_k k'$ soient séparables sur $k'$.
\end{corollary}

\begin{proof}
Il suffit d'appliquer (4.6.6) à $X=\operatorname{Spec}(A)$, où $A$ est une
$k$-algèbre de type fini dont $K$ est le corps des fractions.
\end{proof}

\begin{corollary}[4.6.8]
\label{IV.2.4.6.8-fr}
Soit $X$ un $k$-préschéma de type fini. Il existe une extension finie $k'$ de
$k$ telle que $(X_{(k')})_{\mathrm{red}}$ soit géométriquement réduit sur $k'$
et que les composantes irréductibles de $X_{(k')}$ soient géométriquement
irréductibles et les composantes connexes de $X_{(k')}$ géométriquement
connexes.
\end{corollary}

\begin{proof}
Cela résulte aussitôt de (4.6.6), de (4.5.10) et (4.5.15).
\end{proof}

\begin{definition}[4.6.9]
\label{IV.2.4.6.9-fr}
Soient $k$ un corps, $X$ un $k$-préschéma, $x$ un point de $X$. On dit que $X$
est \emph{géométriquement réduit} (ou \emph{séparable}) (resp.
\emph{géométriquement ponctuellement intègre}) au point $x$ sur $k$, si, pour
toute extension $k'$ de $k$ et tout point $x'$ de
$X'=X\otimes_k k'$ au-dessus de $x$, $X'$ est réduit (resp. intègre) en $x'$
(i.e., $\mathcal O_{X',x'}$ est réduit (resp. intègre)). On dit que $X$ est
\emph{géométriquement ponctuellement intègre} si $X$ est géométriquement
ponctuellement intègre en tous ses points, autrement dit si, pour toute
extension $k'$ de $k$, tous les anneaux locaux de $X\otimes_k k'$ sont intègres
(auquel cas on dit aussi que $X\otimes_k k'$ est ponctuellement intègre).
\end{definition}

Notons que, pour que $X$ soit géométriquement réduit sur $k$ (4.6.2), il faut
et il suffit qu'il soit géométriquement réduit sur $k$ en tout point $x\in X$.

\begin{proposition}[4.6.10]
\label{IV.2.4.6.10-fr}
Soient $k$ un corps, $X$ un $k$-préschéma, $k'$ une extension de $k$, $x'$ un
point de $X'=X\otimes_k k'$, $x$ son image dans $X$. Pour que $X$ soit
géométriquement réduit (resp. géométriquement ponctuellement intègre) en $x$
sur $k$, il faut et il suffit que $X'$ le soit en $x'$ sur $k'$. En particulier,
pour que $X$ soit géométriquement ponctuellement intègre sur $k$, il faut et il
suffit que $X'$ soit géométriquement ponctuellement intègre sur $k'$.
\end{proposition}

\begin{proof}
Il suffit de prouver la première assertion. La condition étant évidemment
nécessaire, prouvons qu'elle est suffisante. Supposons donc $X'$
géométriquement réduit (resp. géométriquement ponctuellement intègre) au point
$x'$ sur $k'$ et prouvons que $X$ l'est au point $x$ sur $k$, autrement dit que
pour toute extension $k''$ de $k$ et tout point $x''$ de
$X''=X\otimes_k k''$ au-dessus de $x$, $\mathcal O_{X'',x''}$ est réduit (resp.
intègre). Pour cela, notons qu'il existe un point $z$ de $X'\times_X X''$ qui
se projette en $x'$ et $x''$ (I, 3.4.7). Soit $s$ le point de
$\operatorname{Spec}(k'\otimes_k k'')$ image de $z$ (I, 3.4.9)~; posons
$K=k(s)$, $Z=X\otimes_k K$, de sorte que $K$ peut être considéré comme
extension composée de $k'$ et $k''$, et $z$ comme un point de $Z$ dont l'image
dans $X'$ (resp. $X''$) est $x'$ (resp. $x''$). Comme $X'$ est
géométriquement réduit (resp. géométriquement ponctuellement intègre) au point
$x'$ sur $k'$, l'anneau $\mathcal O_{Z,z}$ est réduit (resp. intègre), et comme
$\mathcal O_{Z,z}$ est un $\mathcal O_{X'',x''}$-module fidèlement plat, il
s'ensuit que $\mathcal O_{X'',x''}$ est réduit (resp. intègre)
$(0_{\mathrm I},6.5.1)$.
\end{proof}

\begin{corollary}[4.6.11]
\label{IV.2.4.6.11-fr}
Supposons $k$ parfait (resp. algébriquement clos). Pour que $X$ soit
géométriquement réduit (resp. géométriquement ponctuellement intègre) au point
$x$ sur $k$, il faut et il suffit que $\mathcal O_{X,x}$ soit réduit (resp.
intègre).
\end{corollary}

\begin{proof}
C'est évidemment nécessaire. Inversement, si cette condition est remplie,
alors, pour toute extension $k'$ de $k$,
$\operatorname{Spec}(\mathcal O_{X,x})\otimes_k k'$ est réduit (4.6.4) (resp.
intègre (4.4.4))~;

\oldpage[IV-2]{71}

donc, pour tout point $x'$ de $X'=X\otimes_k k'$ au-dessus de $x$,
$\mathcal O_{X',x'}$, qui est isomorphe à un anneau local de
$\operatorname{Spec}(\mathcal O_{X,x})\otimes_k k'$, est réduit (resp.
intègre).
\end{proof}

\begin{proposition}[4.6.12]
\label{IV.2.4.6.12-fr}
Soient $k$ un corps~; $X$ un $k$-préschéma, $x$ un point de $X$, $\Omega$ une
extension parfaite de $k$. Les conditions suivantes sont équivalentes~:
\begin{enumerate}
\item[a)] $X$ est géométriquement réduit sur $k$ au point $x$, autrement dit,
pour toute extension $k'$ de $k$ et tout point $x'$ de
$X'=X\otimes_k k'$ au-dessus de $x$, $\mathcal O_{X',x'}$ est réduit.
\item[b)] Le préschéma $X\otimes_k\Omega$ est réduit en un point au-dessus de
$x$.
\item[c)] Pour toute extension finie radicielle $k'$ de $k$,
$X'=X\otimes_k k'$ est réduit en l'unique point au-dessus de $x$.
\item[d)] $\operatorname{Spec}(\mathcal O_{X,x})$ est géométriquement réduit
sur $k$.
\item[e)] $\mathcal O_{X,x}$ est réduit, et pour toute composante irréductible
$Z$ de $X$ contenant $x$, de point générique $z$, $k(z)$ est une extension
séparable de $k$.
\end{enumerate}
\end{proposition}

\begin{proof}
Les implications d) $\Rightarrow$ a) $\Rightarrow$ b) $\Rightarrow$ c) sont
immédiates, en tenant compte pour la dernière de (4.6.10), (4.6.11) et de ce
que toute extension radicielle de $k$ est isomorphe à une sous-extension de
$\Omega$. Prouvons que c) implique d). Il suffit de montrer que pour toute
extension radicielle finie $k'$ de $k$, $\mathcal O_{X,x}\otimes_k k'$ est
réduit (4.6.1)~; or cet anneau est l'anneau local de $X'=X\otimes_k k'$ en
l'unique point $x'$ de $X'$ au-dessus de $x$ (I, 3.5.8 et 3.5.7)~; il est donc
réduit en vertu de l'hypothèse c). Enfin, l'équivalence de d) et de e)
résulte de l'équivalence de b) et e) dans (4.6.1), puisque les composantes
irréductibles de $X$ contenant $x$ correspondent biunivoquement aux composantes
irréductibles de $\operatorname{Spec}(\mathcal O_{X,x})$ (I, 2.4.2).
\end{proof}

\begin{corollary}[4.6.13]
\label{IV.2.4.6.13-fr}
Sous les hypothèses de (4.6.12), supposons de plus que $X$ soit localement
noethérien. Alors les conditions a) à e) de (4.6.12) équivalent aussi à la
suivante~:
\begin{enumerate}
\item[f)] Il existe un voisinage ouvert $U$ de $x$ qui est géométriquement
réduit sur $k$.
\end{enumerate}
\end{corollary}

\begin{proof}
En effet, cette condition implique trivialement que $X$ est géométriquement
réduit au point $x$ sur $k$. Inversement, s'il en est ainsi, puisque l'anneau
$\mathcal O_{X,x}$ est réduit, il existe un voisinage ouvert $U$ de $x$ dans
$X$ qui est réduit (I, 6.1.13)~; prenant $U$ assez petit, on peut en outre
supposer que $U$ ne rencontre aucune des composantes irréductibles de $X$ ne
contenant pas $x$. Alors le critère (4.6.12, e)) prouve que $U$ est
géométriquement réduit sur $k$, compte tenu du critère (4.6.1, e)).
\end{proof}

Il résulte de (4.6.13) que lorsque $X$ est localement noethérien, l'ensemble
des $x\in X$ où $X$ est géométriquement réduit sur $k$ est \emph{ouvert}, et
c'est le plus grand ouvert de $X$ qui soit géométriquement réduit sur $k$.

\begin{env}[4.6.14]
\label{IV.2.4.6.14-fr}
Il résulte de (4.6.10) et (4.6.11) que si $\Omega$ est une extension
algébriquement close de $k$, il revient au même de dire que $X$ est
géométriquement ponctuellement intègre au point $x$, ou que $X\otimes_k\Omega$
est intègre en un point au-dessus de $x$.

Pour que $X$ soit géométriquement ponctuellement intègre au point $x$, il est
nécessaire que $X$ soit intègre au point $x$, et que si $z$ est le point
générique de l'unique composante irréductible de $X$ contenant $x$, $k(z)$ soit
une extension séparable de $k$~; cela résulte de (4.6.10). Mais ces conditions
\emph{ne sont pas suffisantes}, comme il résulte de l'exemple (4.5.12, (ii)).
Une condition suffisante mais non nécessaire est que $X$ soit géométriquement
réduit au point $x$ et n'appartienne qu'à une seule composante irréductible de
$X$, que

\oldpage[IV-2]{72}

l'on suppose en outre \emph{géométriquement irréductible}~; si de plus $X$ est
localement noethérien, il y a alors un voisinage ouvert de $x$ dans $X$ qui est
\emph{géométriquement intègre}.
\end{env}

Rappelons que si un préschéma localement noethérien est ponctuellement intègre,
il est localement intègre (I, 6.1.13). Si un $k$-préschéma $X$, localement de
type fini sur $k$, est géométriquement ponctuellement intègre sur $k$, il
résulte donc de cette remarque que, pour toute extension $k'$ de $k$,
$X\otimes_k k'$ est \emph{localement intègre} (on dit encore dans ce cas que
$X$ est \emph{géométriquement localement intègre}).

\begin{proposition}[4.6.15]
\label{IV.2.4.6.15-fr}
\begin{enumerate}
\item[(i)] Soient $k$ un corps~; $X$ un $k$-préschéma de type fini. Pour que $X$
soit géométriquement ponctuellement intègre, il faut et il suffit qu'il soit
géométriquement réduit et que le nombre géométrique de ses composantes
irréductibles soit égal au nombre géométrique de ses composantes connexes.
\item[(ii)] Soit $X$ un préschéma localement de type fini sur $k$. Si $X$ est
géométriquement ponctuellement intègre en un point $x$, il existe un voisinage
ouvert $U$ de $x$ qui est géométriquement ponctuellement intègre. En d'autres
termes, l'ensemble des $x\in X$ où $X$ est géométriquement ponctuellement
intègre est ouvert dans $X$.
\end{enumerate}
\end{proposition}

\begin{proof}
\begin{enumerate}
\item[(i)] Étant donnée une extension algébriquement close $\Omega$ de $k$,
$X_{(\Omega)}$ est ponctuellement intègre (ou localement intègre, ce qui
revient au même puisque $X_{(\Omega)}$ est noethérien) si et seulement s'il est
réduit et si le nombre de ses composantes connexes est égal au nombre de ses
composantes irréductibles (I, 6.1.10). Il suffit alors d'appliquer (4.5.1),
(4.6.1) et la première remarque de (4.6.14).

\item[(ii)] La question étant locale sur $X$, on peut se borner au cas où $X$
est un préschéma de type fini sur $k$. En outre, on sait qu'il y a un voisinage
ouvert de $x$ dans $X$ qui est géométriquement réduit (4.6.13), donc on peut
supposer $X$ géométriquement réduit. Il existe alors une extension finie $k'$
de $k$ tel que $X'=X\otimes_k k'$ soit réduit et que ses composantes
irréductibles soient géométriquement irréductibles (4.6.8). Soit
$p:X'\to X$ la projection canonique, qui est un morphisme fini surjectif, et
soient $x'_j$ ($1\leqslant j\leqslant r$) les points de $p^{-1}(x)$~; par
hypothèse, $X'$ est intègre en chacun des $x'_j$, donc chaque $x'_j$ a un
voisinage ouvert $V'_j$ dans $X'$ qui est intègre. Comme $p$ est fermé
(II, 6.1.10), chacun des $p(V'_j)$ est un voisinage de $x$, donc il existe un
voisinage ouvert $U$ de $x$ tel que $p^{-1}(U)$ soit contenu dans la réunion
des $V'_j$, et soit par suite intègre en chacun de ses points, donc localement
intègre. Les composantes irréductibles de $p^{-1}(U)$ sont donc ouvertes et
deux à deux disjointes, et comme elles sont géométriquement irréductibles
(4.5.9), $p^{-1}(U)=U\otimes_k k'$ est géométriquement ponctuellement intègre
en vertu de (i), et il en est donc de même de $U$ par définition.
\end{enumerate}
\end{proof}

\begin{proposition}[4.6.16]
\label{IV.2.4.6.16-fr}
Soient $k$ un corps, $X$ un $k$-préschéma localement noethérien, $\mathscr F$
un $\mathcal O_X$-Module cohérent. Les conditions suivantes sont
équivalentes~:
\begin{enumerate}
\item[a)] Pour toute extension de type fini $k'$ de $k$, si l'on pose
$X'=X\otimes_k k'$, alors $\mathscr F'=\mathscr F\otimes_k k'$ est un
$\mathcal O_{X'}$-Module réduit (3.2.2).
\item[b)] Pour toute extension finie radicielle $k'$ de $k$, $\mathscr F'$ est
un $\mathcal O_{X'}$-Module réduit.
\item[c)] $\mathscr F$ est réduit, et si $\mathscr J$ est l'Idéal de
$\mathcal O_X$ annulateur de $\mathscr F$, le sous-préschéma fermé défini par
$\mathscr J$ est géométriquement réduit sur $k$.
\end{enumerate}

\oldpage[IV-2]{73}

En outre, si $X$ est localement de type fini sur $k$, ces conditions sont aussi
équivalentes à la suivante~:
\begin{enumerate}
\item[d)] Pour toute extension $k'$ de $k$ (ou pour une extension
algébriquement close $k'$ de $k$), $\mathscr F'$ est un
$\mathcal O_{X'}$-Module réduit.
\end{enumerate}
\end{proposition}

\begin{proof}
Il est clair que a) entraîne b). La condition b) entraîne d'abord que
$\mathscr F$ est réduit~: cela signifie (3.2.3) que si $Y$ est le
sous-préschéma fermé de $X$ défini par l'Idéal cohérent $\mathscr J$ de
$\mathcal O_X$, annulateur de $\mathscr F$, $Y$ est réduit, et il y a un
$\mathcal O_Y$-Module cohérent sans torsion $\mathscr G$ de rang~1 sur toute
composante irréductible de $Y$, tel que $j_*(\mathscr G)=\mathscr F$,
$j:Y\to X$ étant l'injection canonique. Or, pour toute extension $k'$ de $k$,
l'annulateur de $\mathscr F'$ est $\mathscr J'=\mathscr J\otimes_k k'$, le
morphisme projection $X'\to X$ étant plat (2.1.11)~; le sous-préschéma fermé de
$X'$ défini par $\mathscr J'$ est $Y'=Y\otimes_k k'$, et si
$j':Y'\to X'$ est l'injection canonique, et
$\mathscr G'=\mathscr G\otimes_k k'$, on a $\mathscr F'=j'_*(\mathscr G')$.
Notons aussi que, pour toute extension $k'$ de $k$ telle que $X'$ soit
localement noethérien, $\mathscr F'$ est sans cycle premier associé immergé
(4.2.7). De plus, tout point maximal $y'$ de $Y'$ est au-dessus d'un point
maximal $y$ de $Y$~; comme $\mathscr G_y$ est isomorphe à
$\mathcal O_{Y,y}$, $\mathscr G'_{y'}$ est isomorphe à
$\mathcal O_{Y',y'}$. On voit donc (par (3.2.3)) qu'il revient au même de dire
que $\mathscr F'$ est réduit ou que le préschéma $Y'$ est réduit, lorsque $X'$
est localement noethérien. Le fait que b) entraîne c) et que c) entraîne
a) est donc conséquence de (4.6.1, d) et b))~; dans a), l'hypothèse que
$k'$ est une extension de type fini de $k$ n'est faite que pour assurer que
$X'$ est localement noethérien. Lorsque $X$ est localement de type fini, $X'$
est localement noethérien pour toute extension $k'$ de $k$, ce qui prouve dans
ce cas l'équivalence de d) et des autres conditions.
\end{proof}

\begin{definition}[4.6.17]
\label{IV.2.4.6.17-fr}
Lorsque $\mathscr F$ vérifie les conditions équivalentes de (4.6.16), on dit
que $\mathscr F$ est \emph{géométriquement réduit} sur $k$, ou \emph{séparable}
sur $k$.
\end{definition}

\begin{proposition}[4.6.18]
\label{IV.2.4.6.18-fr}
Soient $k$ un corps~; $X$ un $k$-préschéma localement noethérien, $\mathscr F$
un $\mathcal O_X$-Module cohérent. Les conditions suivantes sont
équivalentes~:
\begin{enumerate}
\item[a)] Pour toute extension de type fini $k'$ de $k$, si l'on pose
$X'=X\otimes_k k'$, alors $\mathscr F'=\mathscr F\otimes_k k'$ est un
$\mathcal O_{X'}$-Module intègre (3.2.4).
\item[b)] Pour toute extension finie $k'$ de $k$, $\mathscr F'$ est un
$\mathcal O_{X'}$-Module intègre.
\item[c)] $\mathscr F$ est réduit (ou intègre), et si $\mathscr J$ est
l'Idéal de $\mathcal O_X$ annulateur de $\mathscr F$, le sous-préschéma fermé
défini par $\mathscr J$ est géométriquement intègre.
\end{enumerate}
En outre, si $X$ est localement de type fini sur $k$, ces conditions sont aussi
équivalentes à la suivante~:
\begin{enumerate}
\item[d)] Pour toute extension $k'$ de $k$ (ou pour une extension
algébriquement close $k'$ de $k$), $\mathscr F'$ est un
$\mathcal O_{X'}$-Module intègre.
\end{enumerate}
\end{proposition}

\begin{proof}
Les conditions a), b) et c) entraînent que $\mathscr F$ est intègre, donc
$\operatorname{Ass}(\mathscr F)$ réduit à un seul point $x$, qui est le point
générique de $\operatorname{Supp}(\mathscr F)$. Pour toute extension $k'$ de
$k$ pour laquelle $X'$ est localement noethérien, les cycles premiers associés
à $\mathscr F'$ sont donc non immergés, et sont les points maximaux de
$\operatorname{Supp}(\mathscr F')$, qui sont tous au-dessus de $x$. D'autre
part, en vertu de (4.6.16), chacune des conditions a), b), c) entraîne que
$\mathscr F$ est géométriquement réduit~; il est trivial que a) entraîne b),
et b) entraîne que $\operatorname{Supp}(\mathscr F)$ est géométriquement
irréductible (4.5.9), donc c), puisque le sous-préschéma $Y$ défini

\oldpage[IV-2]{74}

par $\mathscr J$ est alors géométriquement réduit et géométriquement
irréductible, donc géométriquement intègre. Inversement, si $Y$ est
géométriquement intègre, pour toute extension $k'$ de $k$,
$\operatorname{Supp}(\mathscr F')$ est irréductible, donc, si $X'$ est
localement noethérien, $\mathscr F'$ est intègre~; cela montre que c) entraîne
a), et que c) entraîne d) lorsque $X$ est localement de type fini sur $k$.
\end{proof}

\begin{definition}[4.6.19]
\label{IV.2.4.6.19-fr}
Lorsque $\mathscr F$ vérifie les conditions équivalentes de (4.6.18), on dit
que $\mathscr F$ est \emph{géométriquement intègre} sur $k$.
\end{definition}

\begin{proposition}[4.6.20]
\label{IV.2.4.6.20-fr}
Soient $k$ un corps, $X$ un $k$-préschéma localement noethérien, $\mathscr F$
un $\mathcal O_X$-Module cohérent. Soit $K$ une extension de $k$ telle que
$X\otimes_k K$ soit localement noethérien. Alors, si $\mathscr F$ est
géométriquement réduit (resp. géométriquement intègre) sur $k$,
$\mathscr F\otimes_k K$ est géométriquement réduit (resp. géométriquement
intègre) sur $K$.
\end{proposition}

\begin{proof}
En effet, pour toute extension finie $K'$ de $K$,
\[
(X\otimes_k K)\otimes_K K'=X\otimes_k K'
\]
est localement noethérien, et on a vu dans les démonstrations de (4.6.16) et
(4.6.18) que l'hypothèse sur $\mathscr F$ entraîne que
$\mathscr F\otimes_k K'$ est réduit (resp. intègre), d'où la conclusion.
\end{proof}

\begin{proposition}[4.6.21]
\label{IV.2.4.6.21-fr}
Soient $k$ un corps, $X$, $Y$ deux $k$-préschémas localement noethériens tels
que $X\times_k Y$ soit localement noethérien. Soient $\mathscr F$ un
$\mathcal O_X$-Module cohérent, $\mathscr G$ un $\mathcal O_Y$-Module cohérent.
\begin{enumerate}
\item[(i)] Si $\mathscr F$ est géométriquement réduit (resp. géométriquement
intègre) et $\mathscr G$ réduit (resp. intègre), alors
$\mathscr F\otimes_k\mathscr G$ est réduit (resp. intègre).
\item[(ii)] Si $\mathscr F$ et $\mathscr G$ sont géométriquement réduits (resp.
géométriquement intègres), il en est de même de
$\mathscr F\otimes_k\mathscr G$.
\end{enumerate}
\end{proposition}

\begin{proof}
\begin{enumerate}
\item[(i)] Utilisant (3.2.3), (4.6.16) et (4.6.18), on peut se borner au cas où
$\operatorname{Supp}(\mathscr F)=X$, $\operatorname{Supp}(\mathscr G)=Y$, $X$
étant géométriquement réduit (resp. géométriquement intègre), $Y$ réduit
(resp. intègre), $\mathscr F$ et $\mathscr G$ sans cycles premiers associés
immergés, et $\mathscr F_x$ étant isomorphe à $\mathcal O_x$ en tout point
maximal $x$ de $X$, $\mathscr G_y$ isomorphe à $\mathcal O_y$ en tout point
maximal $y$ de $Y$. On sait alors (4.2.3) que
$\mathscr F\otimes_k\mathscr G$ est sans cycle premier associé immergé, et que
les cycles premiers associés à $\mathscr F\otimes_k\mathscr G$ sont exactement
les composantes irréductibles de $X\times_k Y$ (4.2.5). On peut donc se borner
au cas où $X$ et $Y$ sont irréductibles~; alors, si $X$ est géométriquement
réduit et $Y$ réduit, on sait (4.2.4, (ii)) que $X\times_k Y$ est réduit~; si
$X$ est géométriquement intègre et $Y$ intègre, on sait que $X\times_k Y$ est
en outre irréductible (4.5.8, (i)), donc intègre. Enfin, tout point maximal $z$
de $X\times_k Y$ est au-dessus de $x$ et de $y$, donc les hypothèses sur
$\mathscr F$ et $\mathscr G$ entraînent, en vertu de (I, 9.1.12), que
$(\mathscr F\otimes_k\mathscr G)_z$ est isomorphe à $\mathcal O_z$, ce qui
achève la démonstration dans ce cas.

\item[(ii)] Le raisonnement est analogue, mais ici (après réduction au cas où
$X=\operatorname{Supp}(\mathscr F)$, $Y=\operatorname{Supp}(\mathscr G)$),
$X\times_k Y$ est géométriquement réduit (resp. géométriquement intègre) en
vertu de (4.6.5, (ii)).
\end{enumerate}
\end{proof}

Les définitions (4.6.18) et (4.6.19) se localisent~:

\begin{definition}[4.6.22]
\label{IV.2.4.6.22-fr}
Soient $k$ un corps, $X$ un $k$-préschéma localement noethérien, $\mathscr F$
un $\mathcal O_X$-Module cohérent. On dit que $\mathscr F$ est
\emph{géométriquement réduit} ou \emph{séparable} (resp. géométriquement
ponctuellement intègre) en un point $x\in X$ si, pour toute extension finie
radicielle (resp. finie) $k'$ de $k$, $\mathscr F\otimes_k k'$ est réduit
(resp. intègre) en chacun des points $x'$ de $X\otimes_k k'$ au-dessus de $x$
(cf. (3.2.2) et (3.2.4)).
\end{definition}

\oldpage[IV-2]{75}

Si $\mathscr F$ est réduit en $x$, il y a un voisinage ouvert $U$ de $x$ tel
que $\mathscr F|U$ soit réduit (3.2.2). En raisonnant comme dans (4.6.16), on
voit que dire que $\mathscr F$ est géométriquement réduit au point $x$
équivaut à dire que $\mathscr F$ est réduit au point $x$ et que le
sous-préschéma fermé $Y$ défini par l'annulateur $\mathscr J$ de $\mathscr F$
est géométriquement réduit au point $x$; il en résulte qu'il y a un voisinage
ouvert $U$ de $x$ tel que $\mathscr F|U$ soit géométriquement réduit (4.6.11).
Les raisonnements de (4.6.16) et (4.6.18) montrent aussi que si $X$ est
localement de type fini sur $k$, dire que $\mathscr F$ est géométriquement
ponctuellement intègre au point $x$ équivaut à dire que $\mathscr F$ est
réduit au point $x$ et que le sous-préschéma $Y$ est géométriquement
ponctuellement intègre au point $x$.

\subsection{Multiplicités dans la décomposition primaire sur un préschéma
algébrique}
\label{subsection:IV.4.7-fr}

\begin{lemma}[4.7.1]
\label{IV.2.4.7.1-fr}
Soient $A$, $B$ deux anneaux locaux, $\mathfrak m$ l'idéal maximal de $A$,
$\rho:A\to B$ un homomorphisme local. On suppose que $B$ soit un $A$-module
plat, et $B/\mathfrak mB$ un $B$-module de longueur finie. Alors, pour qu'un
$A$-module $M$ soit de longueur finie, il faut et il suffit que $M_{(B)}$ soit
un $B$-module de longueur finie, et on a
\begin{equation}
\operatorname{long}_B(M_{(B)})
=\operatorname{long}_A(M)\cdot\operatorname{long}_B(B/\mathfrak mB).
\tag{4.7.1.1}
\label{IV.2.4.7.1.1-fr}
\end{equation}
\end{lemma}

\begin{proof}
C'est un cas particulier de (2.5.5.2).
\end{proof}


\begin{corollary}[4.7.2]
\label{IV.2.4.7.2-fr}
Soient $A$, $B$ deux anneaux noethériens, $\rho:A\to B$ un homomorphisme
d'anneaux tel que $B$ soit un $A$-module plat. Soient $\mathfrak q$ un idéal
premier minimal de $B$, $\mathfrak p=\rho^{-1}(\mathfrak q)$, $M$ un
$A$-module de type fini; alors $\mathfrak p$ est un idéal premier minimal de
$A$, $M_{\mathfrak p}$ est un $A_{\mathfrak p}$-module de longueur finie et
l'on a
\begin{equation}
\operatorname{long}_{B_{\mathfrak q}}((M_{(B)})_{\mathfrak q})
=\operatorname{long}_{A_{\mathfrak p}}(M_{\mathfrak p})
 \operatorname{long}_{B_{\mathfrak q}}
 (B_{\mathfrak q}/\mathfrak pB_{\mathfrak q}).
\tag{4.7.2.1}
\label{IV.2.4.7.2.1-fr}
\end{equation}
\end{corollary}

\begin{proof}
En effet, $B_{\mathfrak q}$ est un $A_{\mathfrak p}$-module plat
$(0_{\mathrm I},6.3.2)$ et l'homomorphisme
$A_{\mathfrak p}\to B_{\mathfrak q}$ est local, donc $B_{\mathfrak q}$ est
un $A_{\mathfrak p}$-module fidèlement plat $(0_{\mathrm I},6.6.2)$; le fait
que $\mathfrak p$ est minimal dans $A$ résulte de (2.3.5). En outre,
$A_{\mathfrak p}$ est alors un anneau artinien (Bourbaki, \emph{Alg. comm.},
chap.~IV, \S~2, n\textsuperscript{o}~5, prop.~9) et $M_{\mathfrak p}$ est donc
un $A_{\mathfrak p}$-module de longueur finie; la formule (4.7.2.1) est donc
un cas particulier de (4.7.1.1).
\end{proof}

\begin{proposition}[4.7.3]
\EGAReaderTarget{ega4.c04.prop.4-7-3}{IV.2.4.7.3-fr}
Soient $k$ un corps d'exposant caractéristique $p$, $X$ un préschéma intègre
localement noethérien, $K=\mathrm R(X)$ le corps des fonctions rationnelles
sur $X$. Soient $k'$ une extension de $k$, $(X'_\alpha)$ la famille des
composantes irréductibles de $X'=X_{(k')}$, $x'_\alpha$ le point générique de
$X'_\alpha$.
\begin{enumerate}
\item[(i)] On suppose que $X'$ soit localement noethérien (ce qui a lieu si
$X$ est localement de type fini sur $k$, ou si $k'$ est une extension de type
fini de $k$). Soient alors $k''$ une extension quelconque de $k'$,
$X''=X\otimes_k k''$, $(X''_\beta)$ la famille des composantes irréductibles
de $X''$, $x''_\beta$ le point générique de $X''_\beta$. Si $x''_\beta$ est
au-dessus de $x'_\alpha$, on a
\begin{equation}
\operatorname{long}(\mathcal O_{x''_\beta})
=\operatorname{long}(\mathcal O_{x'_\alpha})
 \operatorname{long}(\mathcal O_{Z'',x''_\beta})
\tag{4.7.3.1}
\label{IV.2.4.7.3.1-fr}
\end{equation}
où l'on a posé $Z''=(X'_{\mathrm{red}})\otimes_{k'}k''$. En particulier, si
$k''$ est une extension séparable de $k'$ telle que $Z''$ soit localement
noethérien, on a
\[
\operatorname{long}(\mathcal O_{x''_\beta})
=\operatorname{long}(\mathcal O_{x'_\alpha}).
\]

\oldpage[IV-2]{76}

\item[(ii)] Si $X$ est localement de type fini sur $k$, ou si $k'$ est une
extension de type fini de $k$, les nombres
$\operatorname{long}(\mathcal O_{x'_\alpha})$ sont des puissances de $p$.

\item[(iii)] On suppose que $X$ est localement de type fini sur $k$. Alors,
si $k'$ est parfait, tous les nombres
$\operatorname{long}(\mathcal O_{x'_\alpha})$ sont égaux et ne dépendent que
de l'extension $K$ de $k$ (et non de l'extension parfaite particulière $k'$
de $k$ considérée). En outre, il existe une extension radicielle finie $k_1$
de $k$ telle que, si $x_1$ est le point générique de
$X_1=X\otimes_k k_1$ (qui est irréductible), on ait
\[
\operatorname{long}(\mathcal O_{x'_\alpha})
=\operatorname{long}(\mathcal O_{x_1})
\quad\text{pour tout }\alpha.
\]
\end{enumerate}
\end{proposition}

\begin{proof}
\begin{enumerate}
\item[(i)] Rappelons (4.2.4) que les composantes irréductibles $X'_\alpha$,
qui sont en nombre fini, correspondent biunivoquement aux idéaux premiers
minimaux $\mathfrak p_\alpha$ de $K\otimes_k k'$, et que
$\mathcal O_{x'_\alpha}$ est un anneau artinien, isomorphe à
$(K\otimes_k k')_{\mathfrak p_\alpha}$; ces anneaux ne dépendent donc que des
extensions $K$ et $k'$ de $k$ (et sont par suite les mêmes ($k'$ étant fixé)
pour tous les $k$-préschémas intègres ayant même corps $K$ de fonctions
rationnelles). La formule (4.7.3.1) résulte de (4.7.2.1) appliquée à l'anneau
noethérien $A$ d'un voisinage affine de $x'_\alpha$ dans $X'$, et à l'anneau
$B$ d'un voisinage affine assez petit de $x''_\beta$ dans $X''$, $\mathfrak p$
et $\mathfrak q$ étant les idéaux premiers minimaux correspondant à
$x'_\alpha$ et $x''_\beta$ respectivement, et en prenant $M=A$: on notera que
$A/\mathfrak p$ est l'anneau d'un voisinage affine de $x'_\alpha$ dans
$X'_{\mathrm{red}}$, et que $B=A\otimes_{k'}k''$; si
$\mathfrak q'=\mathfrak q/\mathfrak pB$, idéal premier minimal de
$B/\mathfrak pB$, on a
\[
\mathcal O_{Z'',x''_\beta}
=(B/\mathfrak pB)_{\mathfrak q'}
=B_{\mathfrak q}/\mathfrak pB_{\mathfrak q},
\]
d'où (4.7.3.1). La dernière assertion de (i) résulte de ce que dans ce cas
$Z''$ est réduit (4.2.4), donc
$\mathfrak pB_{\mathfrak q}=\mathfrak qB_{\mathfrak q}$.

\item[(ii)] Si $k'$ est une extension de type fini de $k$, il y a une
extension transcendante pure $k'_0$ de $k$ telle que $k'$ soit une extension
finie de $k'_0$; comme $k'_0$ est séparable sur $k$, il résulte de (i)
appliqué en remplaçant $k'$ et $k''$ par $k$ et $k'_0$ respectivement, que
l'on peut remplacer $X$ par $X\otimes_k k'_0$, autrement dit, supposer que
$k'$ est une extension finie de $k$. Il y a alors une extension finie
quasi-galoisienne $k''$ de $k$ contenant $k'$, et la formule (4.7.3.1) montre
qu'il suffit de prouver l'assertion pour $k''$. Mais $k''$ est une extension
séparable d'une extension radicielle finie $k_1$ de $k$, donc, en vertu de
(i), on est ramené dans ce cas à prouver l'assertion lorsque $k'=k_1$ est en
outre radicielle. On sait alors que $K\otimes_k k_1$ est un anneau artinien
n'ayant qu'un seul idéal premier minimal (4.3.2); si $\overline K$ est une
clôture algébrique de $K$, la longueur de $K\otimes_k k_1$ divise celle de
l'anneau artinien $\overline K\otimes_k k_1$ en vertu de (4.7.1.1). Mais
$\overline K\otimes_k k_1$ est une $\overline K$-algèbre artinienne n'ayant
qu'un seul idéal premier (nécessairement maximal), et son quotient par cet
idéal est une extension finie de $\overline K$, donc nécessairement égale à
$\overline K$; la longueur de $\overline K\otimes_k k_1$ est donc égale à son
rang sur $\overline K$, c'est-à-dire à $[k_1:k]$, qui est une puissance de
$p$.

Supposons ensuite que $X$ soit localement de type fini sur $k$; alors, $X''$
est localement noethérien pour toute extension $k''$ de $k'$. Prenant pour
$k''$ une extension algébriquement close de $k'$, on est ramené, en vertu de
(4.7.3.1), à prouver l'assertion lorsque $k'$ est algébriquement clos. Comme
$k'$ est alors une extension séparable de la clôture algébrique $\overline k$
de $k$, il résulte encore de (i) qu'on est ramené au cas où
$k'=\overline k$. Or, on sait (en remplaçant $X$ par un ouvert affine de type
fini sur $k$) qu'il y a une extension radicielle

\oldpage[IV-2]{77}

finie $k_1$ de $k$ telle que
$(X\otimes_k k_1)_{\mathrm{red}}$ soit séparable sur $k_1$ (4.6.6); donc
$(X\otimes_k k_1)_{\mathrm{red}}\otimes_{k_1}\overline k$ est réduit
(4.6.4). En vertu de (4.7.3.1), on est donc de nouveau ramené au cas où $k'$
est une extension finie radicielle de $k$, et on conclut comme plus haut.

\item[(iii)] Notons que deux extensions $k'$, $k''$ de $k$ sont contenues
dans une même extension algébriquement close de $k$; pour montrer que
l'ensemble des nombres $\operatorname{long}(\mathcal O_{x'_\alpha})$, pour
$k'$ parfaite, ne dépend pas du choix de $k'$, on peut donc se borner à
comparer ces nombres pour $k'$ et pour une extension algébriquement close
$k''$ de $k'$; dans ce cas, l'assertion résulte de (i), puisque $k''$ est
séparable sur $k'$. Pour montrer que tous les nombres
$\operatorname{long}(\mathcal O_{x'_\alpha})$ sont égaux, on peut donc se
borner au cas où $k'=\overline k$, clôture algébrique de $k$. Or, l'extension
radicielle finie $k_1$ de $k$ étant alors déterminée comme dans (ii), les
$x'_\alpha$ sont tous au-dessus de l'unique point générique de
$X_1=X\otimes_k k_1$, et la relation
$\operatorname{long}(\mathcal O_{x'_\alpha})=
\operatorname{long}(\mathcal O_{x_1})$ pour tout $\alpha$ résulte encore de
(4.7.3.1), vu le choix de $k_1$.
\end{enumerate}
\end{proof}

\begin{definition}[4.7.4]
\label{IV.2.4.7.4-fr}
Soient $k$ un corps, $X$ un $k$-préschéma intègre, $K$ le corps des fonctions
rationnelles sur $X$. On appelle \emph{multiplicité radicielle} de $K$ (ou de
$X$) sur $k$ la borne supérieure des longueurs des anneaux artiniens
$K\otimes_k k_1$, lorsque $k_1$ parcourt l'ensemble des extensions
radicielles finies de $k$. On appelle \emph{multiplicité séparable} de $K$ (ou
de $X$) sur $k$ le nombre géométrique des composantes irréductibles de $X$ si
ce nombre est fini, et $+\infty$ dans le cas contraire. Enfin, on appelle
\emph{multiplicité totale} de $K$ (ou de $X$) sur $k$ le produit de la
multiplicité radicielle et de la multiplicité séparable de $K$ sur $k$.
\end{definition}

On notera que si $p$ est l'exposant caractéristique de $k$ et si la
multiplicité radicielle de $K$ sur $k$ est finie, c'est une puissance $q=p^f$
de $p$, comme il résulte de (4.7.3, (ii)); lorsque $p\ne1$, $f$ est bien
déterminé et est appelé l'\emph{exposant d'inséparabilité} de $K$ sur $k$;
lorsque $p=1$, la multiplicité radicielle est toujours égale à $1$.

Dire que la multiplicité radicielle (resp. séparable, resp. totale) de $X$ sur
$k$ est $1$ signifie que $X$ est géométriquement réduit (resp.
géométriquement irréductible, resp. géométriquement intègre) sur $k$.

Lorsque $X$ est localement de type fini, la multiplicité radicielle de $X$ sur
$k$ est encore la longueur commune des anneaux locaux $A_i$ aux idéaux
premiers minimaux de l'anneau total des fractions $A$ de $K\otimes_k k'$,
pour toute extension parfaite $k'$ de $k$; la multiplicité séparable de $X$
sur $k$ est le nombre d'idéaux premiers minimaux de $A=K\otimes_k\Omega$ pour
toute extension algébriquement close $\Omega$ de $k$, et enfin la multiplicité
totale est la somme des longueurs des anneaux locaux $A_i$ de
$A=K\otimes_k\Omega$ en ses idéaux premiers minimaux; ces nombres sont tous
trois finis dans ce cas.

\begin{definition}[4.7.5]
\label{IV.2.4.7.5-fr}
Soient $k$ un corps, $X$ un $k$-préschéma localement de type fini sur $k$,
$\mathscr F$ un $\mathcal O_X$-Module cohérent, $x$ un point de $X$ tel que
$\mathscr F_x$ soit un $\mathcal O_x$-module de longueur finie. On appelle
\emph{longueur géométrique} de $\mathscr F$ en $x$ (relative à $k$), ou
\emph{multiplicité radicielle de $x$ pour $\mathscr F$} le produit
\[
\lambda_x(\mathscr F)
=\operatorname{long}_{\mathcal O_x}(\mathscr F_x)\mu_i(k(x)\mid k),
\]
où $\mu_i(k(x)\mid k)$ désigne la multiplicité radicielle de $k(x)$ sur $k$.
\end{definition}

Lorsque $X$ est intègre, et $x$ le point générique de $X$, on a
$\operatorname{long}(\mathcal O_x)=1$, donc la multiplicité radicielle de $x$
pour $\mathcal O_X$ n'est autre que la multiplicité radicielle de $X$ sur $k$,
définie dans (4.7.4).

\oldpage[IV-2]{78}

\begin{proposition}[4.7.6]
\label{IV.2.4.7.6-fr}
Soit
\[
0\longrightarrow\mathscr F'\longrightarrow\mathscr F
\longrightarrow\mathscr F''\longrightarrow0
\]
une suite exacte de $\mathcal O_X$-Modules cohérents. Si en un point $x\in X$
la longueur géométrique $\lambda_x(\mathscr F)$ de $\mathscr F$ en ce point
est définie, il en est de même des longueurs géométriques
$\lambda_x(\mathscr F')$ et $\lambda_x(\mathscr F'')$, de $\mathscr F'$ et
$\mathscr F''$ au point $x$, et réciproquement; en outre, on a
\begin{equation}
\lambda_x(\mathscr F)=\lambda_x(\mathscr F')+\lambda_x(\mathscr F'').
\tag{4.7.6.1}
\label{IV.2.4.7.6.1-fr}
\end{equation}
\end{proposition}

\begin{proof}
La proposition résulte aussitôt de la définition, car
$\operatorname{long}(\mathscr F_x)$ est finie si et seulement si
$\operatorname{long}(\mathscr F'_x)$ et
$\operatorname{long}(\mathscr F''_x)$ sont finies, et l'on a
\[
\operatorname{long}(\mathscr F_x)
=\operatorname{long}(\mathscr F'_x)+\operatorname{long}(\mathscr F''_x).
\]
\end{proof}

\begin{env}[4.7.7]
\label{IV.2.4.7.7-fr}
Sous les hypothèses de (4.7.5), les seuls points $x\in X$ tels que
$\mathscr F_x$ soit un $\mathcal O_x$-module de longueur finie non nulle sont
(en vertu de (3.1.2)) les points maximaux de
$\operatorname{Supp}(\mathscr F)$, comme il résulte de Bourbaki,
\emph{Alg. comm.}, chap.~IV, \S~2, n\textsuperscript{o}~5, cor.~2 de la
prop.~7, la question étant locale sur $X$. On sait qu'il existe un
sous-préschéma fermé $Y$ ayant $\operatorname{Supp}(\mathscr F)$ pour espace
sous-jacent, et un $\mathcal O_Y$-Module cohérent $\mathscr G$ tel que
$\mathscr F=j_*(\mathscr G)$, $j:Y\to X$ étant l'injection canonique. Si $x$
est point maximal de $Y$, on voit que les longueurs géométriques de
$\mathscr F$ et de $\mathscr G$ en $x$ sont égales, $\mathcal O_{X,x}$ et
$\mathcal O_{Y,x}$ ayant même corps résiduels. On en déduit la caractérisation
suivante de la longueur géométrique:
\end{env}

\begin{proposition}[4.7.8]
\label{IV.2.4.7.8-fr}
Sous les hypothèses de (4.7.5), soit $x$ un point maximal de
$\operatorname{Supp}(\mathscr F)$. Soit $k'$ une extension parfaite de $k$;
posons $X'=X\otimes_k k'$, $\mathscr F'=\mathscr F\otimes_k k'$; alors la
longueur géométrique $\lambda_x(\mathscr F)$ de $\mathscr F$ en $x$ est égale
à la longueur du $\mathcal O_{x'}$-module $\mathscr F'_{x'}$, en tout point
maximal $x'$ de $\operatorname{Supp}(\mathscr F')$ au-dessus de $x$. En outre,
il existe une extension radicielle finie $k_1$ de $k$ telle que, en posant
$X_1=X\otimes_k k_1$, $\mathscr F_1=\mathscr F\otimes_k k_1$,
$\lambda_x(\mathscr F)$ soit égale à la longueur du
$\mathcal O_{x_1}$-module $(\mathscr F_1)_{x_1}$ en tout point maximal $x_1$
de $\operatorname{Supp}(\mathscr F_1)$ au-dessus de $x$.
\end{proposition}

\begin{proof}
La remarque de (4.7.7) montre qu'on peut se borner au cas où $x$ est point
maximal de $X$. En outre, la question est locale sur $X$ et l'on peut donc
supposer $X=\operatorname{Spec}(A)$ affine noethérien,
$\mathscr F=\widetilde M$, où $M$ est un $A$-module de type fini; soient
$B=A\otimes_k k'$, $\mathfrak p$ (resp. $\mathfrak q$) l'idéal premier
minimal de $A$ (resp. $B$) correspondant à $x$ (resp. $x'$). On a alors, en
appliquant la formule (4.7.2.1)
\[
\operatorname{long}_{B_{\mathfrak q}}(\mathscr F'_{x'})
=\operatorname{long}_{A_{\mathfrak p}}(\mathscr F_x)
 \operatorname{long}_{B_{\mathfrak q}}
 (B_{\mathfrak q}/\mathfrak pB_{\mathfrak q}).
\]
Mais comme $k(x)=A_{\mathfrak p}/\mathfrak pA_{\mathfrak p}$, le même
raisonnement que dans (4.7.3) montre que
$\operatorname{long}_{B_{\mathfrak q}}
(B_{\mathfrak q}/\mathfrak pB_{\mathfrak q})$ est égal à la longueur de
l'anneau local $(B/\mathfrak pB)_{\mathfrak q'}$, où
$\mathfrak q'=\mathfrak q/\mathfrak pB$. L'anneau
$(B/\mathfrak pB)_{\mathfrak q'}$ est un anneau local de
$(A/\mathfrak p)\otimes_k k'$ pour un idéal premier minimal de cet anneau,
donc aussi un anneau local de l'anneau total des fractions de
$k(x)\otimes_k k'$, et la longueur d'un tel anneau est par définition
$\mu_i(k(x)\mid k)$, d'où la première assertion. La seconde se démontre de
même, en remplaçant $k'$ par $k_1$ et utilisant (4.7.3, (iii)).
\end{proof}

On dit encore, dans le cas où $x$ est point maximal de
$\operatorname{Supp}(\mathscr F)$, que la longueur géométrique de $\mathscr F$
en $x$ est la multiplicité radicielle du cycle premier maximal
$\overline{\{x\}}$ associé à $\mathscr F$.

\begin{corollary}[4.7.9]
\label{IV.2.4.7.9-fr}
Sous les hypothèses de (4.7.5), soient $k'$ une extension de $k$,
$X'=X\otimes_k k'$, $\mathscr F'=\mathscr F\otimes_k k'$; soient $Z$ une
composante irréductible de $\operatorname{Supp}(\mathscr F)$, et $Z'$ une
composante irréductible de $\operatorname{Supp}(\mathscr F')$ qui domine $Z$
(4.2.7). Alors les multiplicités radicielles de $Z$ par rapport à
$\mathscr F$ et de $Z'$ par rapport à $\mathscr F'$ sont les mêmes.
\end{corollary}

\oldpage[IV-2]{79}

\begin{proof}
Il suffit de considérer une extension parfaite $k''$ de $k'$ et en posant
$X''=X\otimes_k k''$, $\mathscr F''=\mathscr F\otimes_k k''$, un point
maximal $x''$ de $\operatorname{Supp}(\mathscr F'')$ qui est au-dessus du
point générique de $Z'$; alors il résulte de (4.7.8) que
$\operatorname{long}(\mathscr F''_{x''})$ est égal aux multiplicités de $Z$ et
de $Z'$.
\end{proof}

\begin{proposition}[4.7.10]
\label{IV.2.4.7.10-fr}
Soient $k$ un corps, $X$ un $k$-préschéma localement de type fini sur $k$,
$\mathscr F$ un $\mathcal O_X$-Module cohérent, $x$ un point de $X$, $z_i$
les points génériques des composantes irréductibles de
$\operatorname{Supp}(\mathscr F)$ qui contiennent $x$. Pour que $\mathscr F$
soit géométriquement réduit au point $x$ (4.6.22), il faut et il suffit que
$x$ n'appartienne à aucun cycle premier associé immergé de $\mathscr F$ et
que, pour tout $i$, la longueur géométrique $\lambda_{z_i}(\mathscr F)$ de
$\mathscr F$ au point $z_i$ soit égale à $1$.
\end{proposition}

\begin{proof}
Montrons d'abord que les conditions de l'énoncé sont suffisantes. En effet, si
$k'$ est une extension finie de $k$, $x'$ un point de
$X'=X\otimes_k k'$ au-dessus de $x$, et
$\mathscr F'=\mathscr F\otimes_k k'$, $x'$ n'appartient à aucun cycle premier
associé immergé de $\mathscr F'$ en vertu de (4.2.7); en outre, en vertu de
(4.7.8), la longueur géométrique de $\mathscr F'$ est égale à $1$ en tout
point générique $z'_j$ d'une composante irréductible de
$\operatorname{Supp}(\mathscr F')$ contenant $x'$, un tel point étant
au-dessus d'un des $z_i$ en vertu de (2.3.4); \emph{a fortiori} on a
$\operatorname{long}(\mathscr F'_{z'_j})=1$, ce qui prouve que $\mathscr F'$
est réduit au point $x'$.

Inversement, il résulte d'abord de (4.2.7) que si, pour une extension finie
$k'$ de $k$, $\mathscr F'$ est réduit en tous les points $x'$ de $X'$
au-dessus de $x$, alors $x$ n'appartient à aucun cycle premier associé immergé
de $\mathscr F$. En outre, en prenant pour $k'$ le corps $k_1$ de l'énoncé de
(4.7.8) (où $x$ est remplacé par un des $z_i$), l'hypothèse que $\mathscr F$
est géométriquement réduit au point $x$ entraîne que
$\lambda_{z_i}(\mathscr F)=1$.
\end{proof}

\begin{proposition}[4.7.11]
\label{IV.2.4.7.11-fr}
Soient $k$ un corps, $X$ un préschéma localement de type fini sur $k$,
$\mathscr F$ un $\mathcal O_X$-Module cohérent, $k'$ une extension de $k$,
$X'=X\otimes_k k'$, $\mathscr F'=\mathscr F\otimes_k k'$, $x$ un point de
$X$, $x'$ un point de $X'$ au-dessus de $x$. Pour que $\mathscr F$ soit
géométriquement réduit (resp. géométriquement ponctuellement intègre) au point
$x$, il faut et il suffit que $\mathscr F'$ soit géométriquement réduit (resp.
géométriquement ponctuellement intègre) au point $x'$.
\end{proposition}

\begin{proof}
Les composantes irréductibles de $\operatorname{Supp}(\mathscr F')$ contenant
$x'$ dominent chacune une composante irréductible de
$\operatorname{Supp}(\mathscr F)$ contenant $x$, et inversement chacune de
ces dernières est dominée par une composante irréductible de
$\operatorname{Supp}(\mathscr F')$ contenant $x'$ ((4.2.7, (i)) et (2.3.4)).
L'assertion concernant la propriété de séparabilité résulte donc de (4.2.7,
(ii)), (4.7.9) et (4.7.10). Soient $Y$ le sous-préschéma fermé de $X$ défini
par l'Idéal $\mathscr J$ annulateur de $\mathscr F$, de sorte que
$Y'=Y\otimes_k k'$ est le sous-préschéma fermé de $X'$ défini par l'Idéal
$\mathscr J'=\mathscr J\otimes_k k'$ annulateur de $\mathscr F'$; dire que
$\mathscr F$ (resp. $\mathscr F'$) est géométriquement ponctuellement intègre
au point $x$ (resp. $x'$) revient à dire que $\mathscr F$ (resp.
$\mathscr F'$) est géométriquement réduit au point $x$ (resp. $x'$) et que
$Y$ (resp. $Y'$) est géométriquement ponctuellement intègre au point $x$
(resp. $x'$). En vertu de l'assertion de l'énoncé concernant la séparabilité,
l'hypothèse que $\mathscr F$ est géométriquement ponctuellement intègre au
point $x$, ou l'hypothèse que $\mathscr F'$ est géométriquement ponctuellement
intègre au point $x'$, entraînent toutes deux qu'il existe dans $Y$ un
voisinage de $x$ qui est séparable, et la conclusion résulte de (4.6.10).
\end{proof}

\begin{definition}[4.7.12]
\label{IV.2.4.7.12-fr}
Soient $k$ un corps, $X$ un $k$-préschéma localement de type fini sur $k$,
$\mathscr F$ un $\mathcal O_X$-Module cohérent, $x$ un point maximal du
support de $\mathscr F$. On appelle \emph{multiplicité}

\oldpage[IV-2]{80}

\emph{totale de $x$} (ou du
cycle premier $\overline{\{x\}}$) pour $\mathscr F$ (relative à $k$) le
produit de la multiplicité radicielle de $x$ pour $\mathscr F$ par la
multiplicité séparable de $k(x)$ sur $k$.
\end{definition}

Il revient au même de dire que la multiplicité totale de $x$ pour $\mathscr F$
est le produit de la longueur
$\operatorname{long}_{\mathcal O_x}(\mathscr F_x)$ par la multiplicité totale
de $k(x)$ sur $k$ (4.7.4). Avec les notations de (4.7.11), la multiplicité
totale de $x$ pour $\mathscr F$ est égale à la somme des multiplicités totales
pour $\mathscr F'$ des points maximaux $x'_i$ de
$\operatorname{Supp}(\mathscr F')$ qui sont au-dessus de $x$; en outre, il
existe une extension finie $k'$ de $k$, telle que la multiplicité totale de
$x$ pour $\mathscr F$ soit égale à
$\sum_i\operatorname{long}(\mathscr F'_{x'_i})$.

\begin{proposition}[4.7.13]
\label{IV.2.4.7.13-fr}
Les hypothèses et les notations étant celles de (4.7.10), pour que
$\mathscr F$ soit géométriquement ponctuellement intègre au point $x$, il
suffit que $x$ n'appartienne qu'à un seul cycle premier associé
(nécessairement non immergé) de $\mathscr F$ et que la multiplicité totale
pour $\mathscr F$ du point générique $z$ de ce cycle soit égale à $1$.
\end{proposition}

\begin{proof}
En effet, si $k'$ est une extension finie de $k$, $x'$ un point de
$X'=X\otimes_k k'$ au-dessus de $x$ et
$\mathscr F'=\mathscr F\otimes_k k'$, $x'$ ne peut appartenir qu'à un seul
cycle premier associé à $\mathscr F'$, puisqu'il n'y a qu'un seul point
maximal de $\operatorname{Supp}(\mathscr F')$ au-dessus de $z$ par hypothèse;
en outre, $\mathscr F$ est géométriquement réduit d'après (4.7.10), d'où la
conclusion.
\end{proof}
\subsection{Corps de définition}
\label{subsection:IV.4.8-fr}

\begin{env}[4.8.1]
\label{IV.2.4.8.1-fr}
Étant donné un préschéma $X$, nous désignerons, dans ce numéro, par
$\mathfrak S(X)$ l'ensemble des sous-préschémas de $X$, et comme d'ordinaire
par $\mathfrak P(X)$ l'ensemble des parties de l'espace sous-jacent à $X$;
pour tout $\mathcal O_X$-Module quasi-cohérent $\mathscr F$, nous désignerons
par $\Phi(\mathscr F)$ l'ensemble des sous-$\mathcal O_X$-Modules
quasi-cohérents de $\mathscr F$.
\end{env}

\begin{env}[4.8.2]
\label{IV.2.4.8.2-fr}
Soient $k$ un corps, $X$, $Y$ deux $k$-préschémas, $\mathscr F$, $\mathscr G$
deux $\mathcal O_X$-Modules quasi-cohérents, $K$, $K'$ deux extensions de $k$
telles que $K'\subset K$. On a alors, correspondant au morphisme
$\operatorname{Spec}(K)\to\operatorname{Spec}(K')$, des applications
canoniques
\begin{align}
\Phi(\mathscr F\otimes_kK')
&\longrightarrow\Phi(\mathscr F\otimes_kK),
\tag{4.8.2.1}
\label{IV.2.4.8.2.1-fr}\\
\operatorname{Hom}(\mathscr F\otimes_kK',\mathscr G\otimes_kK')
&\longrightarrow
\operatorname{Hom}(\mathscr F\otimes_kK,\mathscr G\otimes_kK),
\tag{4.8.2.2}
\label{IV.2.4.8.2.2-fr}\\
\mathfrak S(X_{(K')})&\longrightarrow\mathfrak S(X_{(K)}),
\tag{4.8.2.3}
\label{IV.2.4.8.2.3-fr}\\
\operatorname{Hom}_{K'}(X_{(K')},Y_{(K')})
&\longrightarrow\operatorname{Hom}_K(X_{(K)},Y_{(K)}),
\tag{4.8.2.4}
\label{IV.2.4.8.2.4-fr}\\
\mathfrak P(X_{(K')})&\longrightarrow\mathfrak P(X_{(K)}).
\tag{4.8.2.5}
\label{IV.2.4.8.2.5-fr}
\end{align}
Si $p_X:X_{(K)}\to X_{(K')}$ est la projection canonique, l'application
(4.8.2.5) est $M\mapsto p_X^{-1}(M)$, et de même (4.8.2.3) est
$Z\mapsto p_X^{-1}(Z)=Z\otimes_{K'}K$ (I, 4.4.1). L'application (4.8.2.1)
est
\[
\mathscr H\longmapsto p_X^*(\mathscr H)
=\mathscr H\otimes_{\mathcal O_{X_{(K')}}}\mathcal O_{X_{(K)}};
\]
(4.8.2.2) est $u\mapsto p_X^*(u)$ et enfin (4.8.2.4) est l'application
$f\mapsto f_{(K)}$.

Si l'on désigne par $\varepsilon_{K,K'}$ une quelconque de ces applications
canoniques, il est clair que si $K''$, $K'$, $K$ sont trois extensions de $k$
telles que $K''\subset K'\subset K$, on a
\begin{equation}
\varepsilon_{K,K''}=\varepsilon_{K,K'}\circ\varepsilon_{K',K''}.
\tag{4.8.2.6}
\label{IV.2.4.8.2.6-fr}
\end{equation}
\end{env}

\oldpage[IV-2]{81}

\begin{proposition}[4.8.3]
\label{IV.2.4.8.3-fr}
Les applications canoniques (4.8.2.1) à (4.8.2.5) sont injectives.
\end{proposition}

\begin{proof}
On sait en effet que $p_X$ est fidèlement plat et quasi-compact; le fait que
(4.8.2.5) soit injectif résulte donc de ce que $p_X$ est surjectif.
L'injectivité de (4.8.2.1) résulte de (2.2.2), celle de (4.8.2.2) de (2.2.7),
celle de (4.8.2.3) de (2.2.15) et enfin celle de (4.8.2.4) de (2.2.16).
\end{proof}

\begin{definition}[4.8.4]
\label{IV.2.4.8.4-fr}
Les notations étant celles de (4.8.2), on dit qu'un
sous-$\mathcal O_{X_{(K)}}$-Module quasi-cohérent de
$\mathscr F\otimes_kK$ (resp. un homomorphisme
$\mathscr F\otimes_kK\to\mathscr G\otimes_kK$, resp. un sous-préschéma de
$X_{(K)}$, resp. un $K$-morphisme $X_{(K)}\to Y_{(K)}$, resp. un
sous-ensemble de $X_{(K)}$) est \emph{défini sur $K'$} s'il appartient à
l'image de l'application (4.8.2.1) (resp. (4.8.2.2), resp. (4.8.2.3), resp.
(4.8.2.4), resp. (4.8.2.5)). On dit alors que $K'$ est un \emph{corps de
définition} de l'objet considéré.
\end{definition}

Il est clair que $K$ lui-même est corps de définition d'un élément de l'un
quelconque des cinq ensembles d'arrivée des applications (4.8.2.1) à
(4.8.2.5). On notera toutefois que pour un $K$-préschéma $Z$, dire par exemple
qu'un $\mathcal O_Z$-Module quasi-cohérent $\mathscr H$ est «~défini sur un
sous-corps $K'$ de $K$~» n'a de sens que si $Z$ \emph{a été donné} sous la
forme $X_{(K)}$ pour un préschéma $X$ sur un sous-corps $k$ de $K$, et si
$\mathscr H$ est un sous-$\mathcal O_Z$-Module d'un $\mathcal O_Z$-Module qui
\emph{a été donné} sous la forme $\mathscr F\otimes_kK$, où $\mathscr F$ est
un $\mathcal O_X$-Module quasi-cohérent; on a des remarques analogues pour les
autres notions. Il résulte de (4.8.2.6) que si un élément de l'un des ensembles
d'arrivée des applications de (4.8.2) est défini sur $K'$, il est aussi défini
sur tout corps $K''$ tel que $K'\subset K''\subset K$. Enfin, avec les
notations de (4.8.2.6), si $K_1$ est une extension de $K$, pour qu'un élément
de l'ensemble d'arrivée de $\varepsilon_{K,K'}$ soit défini sur $K'$, il
faut et il suffit que son image par $\varepsilon_{K_1,K}$ soit définie sur
$K'$, en vertu de la relation
$\varepsilon_{K_1,K'}=\varepsilon_{K_1,K}\circ\varepsilon_{K,K'}$.

\begin{proposition}[4.8.5]
\label{IV.2.4.8.5-fr}
Avec les notations de (4.8.2), soit $(X_\alpha)$ un recouvrement ouvert de
$X$. Pour qu'un élément $\mathscr H\in\Phi(\mathscr F\otimes_kK)$ (resp.
$u\in\operatorname{Hom}(\mathscr F\otimes_kK,\mathscr G\otimes_kK)$, resp.
$Z\in\mathfrak S(X_{(K)})$, resp. $M\in\mathfrak P(X_{(K)})$) soit défini sur
$K'$, il faut et il suffit que pour tout $\alpha$,
$\mathscr H|(X_\alpha)_{(K)}$ (resp. $u|(X_\alpha)_{(K)}$,
$Z\cap(X_\alpha)_{(K)}$, $M\cap(X_\alpha)_{(K)}$) soit défini sur $K'$.
\end{proposition}

\begin{proof}
Cela résulte aussitôt de l'injectivité des applications (4.8.2.1), (4.8.2.2),
(4.8.2.3) et (4.8.2.5): par exemple, si pour tout $\alpha$, il y a un
sous-$\mathcal O_{(X_\alpha)_{(K')}}$-Module quasi-cohérent
$\mathscr H'_\alpha$ de $(\mathscr F|X_\alpha)\otimes_kK'$ tel que
$\mathscr H|(X_\alpha)_{(K)}=\mathscr H'_\alpha\otimes_{K'}K$, pour deux
indices $\alpha$, $\beta$ quelconques, on aura
\[
\bigl(\mathscr H'_\alpha|(X_\alpha\cap X_\beta)_{(K')}\bigr)
 \otimes_{K'}K
=\bigl(\mathscr H'_\beta|(X_\alpha\cap X_\beta)_{(K')}\bigr)
 \otimes_{K'}K,
\]
donc nécessairement
$\mathscr H'_\alpha|(X_\alpha\cap X_\beta)_{(K')}
=\mathscr H'_\beta|(X_\alpha\cap X_\beta)_{(K')}$ quels que soient
$\alpha$ et $\beta$ et il y a par suite un sous-$\mathcal O_{X_{(K')}}$-Module
quasi-cohérent $\mathscr H'$ de $\mathscr F\otimes_kK'$ tel que
$\mathscr H'|(X_\alpha)_{(K')}=\mathscr H'_\alpha$ pour tout $\alpha$, donc
$\mathscr H=\mathscr H'\otimes_{K'}K$. On raisonne de même dans les autres
cas.
\end{proof}

\begin{lemma}[4.8.6]
\label{IV.2.4.8.6-fr}
Soient $k$ un corps, $K$ une extension de $k$, $A$ une $k$-algèbre, $M$ un
$A$-module, $\overline N$ un sous-$A_{(K)}$-module de $M_{(K)}$, $K'$ une
sous-extension de $K$. Les conditions suivantes sont équivalentes:
\begin{enumerate}
\item[a)] $\overline N$ est de la forme $N'\otimes_{K'}K$, où $N'$ est un
sous-$A_{(K')}$-module de $M_{(K')}$. 
\item[b)] Si $X=\operatorname{Spec}(A)$, le
$\mathcal O_{X_{(K)}}$-Module quasi-cohérent $\widetilde{\overline N}$ sur
$X_{(K)}=\operatorname{Spec}(A_{(K)})$ est défini sur $K'$.
\item[c)] $\overline N$ est de la forme $N'\otimes_{K'}K$, où $N'$ est un
sous-$K'$-espace vectoriel de $M_{(K')}$. 
\end{enumerate}
\end{lemma}

\oldpage[IV-2]{82}

\begin{proof}
L'équivalence de a) et b) résulte trivialement de (I, 1.6.5); il est clair
d'autre part que a) entraîne c). Inversement, si c) est vérifiée, on sait que
l'on peut écrire (à identification canonique près)
$N'=\overline N\cap M_{(K')}$. Comme $\overline N$ et $M_{(K')}$ sont par
hypothèse des $A_{(K')}$-modules, il en est de même de $N'$.
\end{proof}

\begin{corollary}[4.8.7]
\label{IV.2.4.8.7-fr}
Sous les hypothèses de (4.8.6), il existe un plus petit sous-corps $K'$ de $K$
contenant $k$ tel que les conditions équivalentes de (4.8.6) soient vérifiées.
\end{corollary}

\begin{proof}
Il suffit de le voir pour la condition c), où cela résulte de Bourbaki,
\emph{Alg.}, chap.~II, 3\textsuperscript{e} éd., \S~8,
n\textsuperscript{o}~6, prop.~6.
\end{proof}

\begin{lemma}[4.8.8]
\label{IV.2.4.8.8-fr}
Les notations étant celles de (4.8.2):
\begin{enumerate}
\item[(i)] Soit
$\overline u:\mathscr F\otimes_kK\to\mathscr G\otimes_kK$ un
$\mathcal O_{X_{(K)}}$-homomorphisme, et soit
$\overline{\mathscr H}\subset
(\mathscr F\otimes_kK)\oplus(\mathscr G\otimes_kK)$ son graphe. Pour que
$\overline u$ soit défini sur $K'$, il faut et il suffit que
$\overline{\mathscr H}$ le soit.
\item[(ii)] Soit $\overline Z$ un sous-préschéma fermé de $X_{(K)}$, et soit
$\overline{\mathscr J}$ l'Idéal quasi-cohérent de $\mathcal O_{X_{(K)}}$
définissant $\overline Z$. Pour que $\overline Z$ soit défini sur $K'$, il
faut et il suffit que $\overline{\mathscr J}$ le soit.
\item[(iii)] Soit $\overline f$ un $K$-morphisme:
$X_{(K)}\to Y_{(K)}$ et soit $\overline Z$ le sous-préschéma de
$X_{(K)}\times_KY_{(K)}$, graphe de $\overline f$ (I, 5.3.11); pour que
$\overline f$ soit défini sur $K'$, il faut et il suffit que $\overline Z$ le
soit.
\end{enumerate}
\end{lemma}

\begin{proof}
L'assertion (ii) résulte aussitôt de (I, 4.4.5). La nécessité dans l'assertion
(iii) est évidente; supposons réciproquement que l'on ait
$\overline Z=Z'_{(K)}$, où $Z'$ est un sous-préschéma de
$(X\times_kY)_{(K')}$; si $p':Z'\to X_{(K')}$ est la restriction à $Z'$ de la
première projection, on sait que $p'_{(K)}$ est un isomorphisme de préschémas,
donc il en est de même de $p'$ (2.7.1, (viii)); mais cela signifie (I, 5.3.11)
que $Z'$ est le graphe d'un $K'$-morphisme
$f':X_{(K')}\to Y_{(K')}$, et l'on a alors
$\overline f=f'_{(K)}$ (I, 5.3.12).

Enfin, pour prouver (i), on peut, en vertu de (4.8.5), supposer que
$X=\operatorname{Spec}(A)$ est affine, $\mathscr F=\widetilde M$,
$\mathscr G=\widetilde N$, où $M$ et $N$ sont des $A$-modules, et
$\overline u=(\overline v)^{\sim}$, où
$\overline v:M_{(K)}\to N_{(K)}$ est un $A_{(K)}$-homomorphisme. Supposons
que le graphe de $\overline v$ soit de la forme $P'_{(K)}$, où $P'$ est un
sous-$A_{(K')}$-module de $M_{(K')}\oplus N_{(K')}$; si
$p':P'\to M_{(K')}$ est la restriction à $P'$ de la première projection, on
sait que $p'\otimes1_K$ est un isomorphisme de $A_{(K)}$-modules, donc, par
fidèle platitude $(0_{\mathrm I},6.4.1)$, $p'$ est un isomorphisme de
$A_{(K')}$-modules, ce qui prouve que $P'$ est le graphe d'un
$A_{(K')}$-homomorphisme $v':M_{(K')}\to N_{(K')}$ et l'on a évidemment
$\overline v=v'\otimes1_{K'}$.
\end{proof}

\begin{proposition}[4.8.9]
\label{IV.2.4.8.9-fr}
Soient $k$ un corps, $K$ une extension de $k$, $X$ un $k$-préschéma,
$\mathscr F$ un $\mathcal O_X$-Module quasi-cohérent. Soit $\overline{\mathscr H}$ un
sous-$\mathcal O_{X_{(K)}}$-Module quasi-cohérent de
$\mathscr F\otimes_kK$. Alors il existe un plus petit corps de définition de
$\overline{\mathscr H}$.
\end{proposition}

\begin{proof}
Lorsque $X$ est affine, l'assertion n'est autre que (4.8.7). Dans le cas
général, considérons un recouvrement $(X_\alpha)$ de $X$ par des ouverts
affines; en vertu de ce qui précède, il existe pour tout $\alpha$ un plus petit
sous-corps $K'_\alpha$ de $K$ contenant $k$ et tel que
$\overline{\mathscr H}|(X_\alpha)_{(K)}$ soit défini sur $K'_\alpha$. En vertu de
(4.8.5), le sous-corps de $K$ engendré par les $K'_\alpha$ est le plus petit
corps de définition de $\overline{\mathscr H}$.
\end{proof}

\begin{corollary}[4.8.10]
\label{IV.2.4.8.10-fr}
Les notations étant celles de (4.8.9), soit $\mathscr G$ un second
$\mathcal O_X$-Module quasi-cohérent, et soit
$\overline u:\mathscr F\otimes_kK\to\mathscr G\otimes_kK$ un
$\mathcal O_{X_{(K)}}$-homomorphisme; alors il existe un plus petit corps de
définition de $\overline u$.
\end{corollary}

\oldpage[IV-2]{83}

\begin{proof}
En effet, il résulte de (4.8.8) qu'un tel corps est aussi un plus petit corps
de définition du graphe de $\overline u$, et l'existence de ce dernier résulte
de (4.8.9).
\end{proof}

\begin{corollary}[4.8.11]
\label{IV.2.4.8.11-fr}
Soient $k$ un corps, $X$ un $k$-préschéma, $K$ une extension de $k$,
$\overline Z$ un sous-préschéma fermé de $X_{(K)}$; alors il existe un plus
petit corps de définition de $\overline Z$.
\end{corollary}

\begin{proof}
En effet, si $\overline{\mathscr J}$ est l'Idéal quasi-cohérent de
$\mathcal O_{X_{(K)}}$ définissant $\overline Z$, il résulte de (4.8.8) que
cela équivaut à l'existence d'un plus petit corps de définition de
$\overline{\mathscr J}$, qui résulte de (4.8.9).
\end{proof}

\begin{corollary}[4.8.12]
\label{IV.2.4.8.12-fr}
Soient $k$ un corps, $X$, $Y$ deux $k$-préschémas, $Y$ étant supposé séparé,
$K$ une extension de $k$, $\overline f:X_{(K)}\to Y_{(K)}$ un $K$-morphisme.
Il existe alors un plus petit corps de définition de $\overline f$.
\end{corollary}

\begin{proof}
En effet, l'existence d'un tel corps, en vertu de (4.8.8), équivaut à celle
d'un plus petit corps de définition du graphe $\overline Z$ de $\overline f$;
mais comme $\overline Z$ est un sous-préschéma fermé de
$(X\times_kY)_{(K)}$ (I, 5.4.3), l'existence d'un plus petit corps de
définition de $\overline Z$ résulte de (4.8.11).
\end{proof}

\begin{proposition}[4.8.13]
\label{IV.2.4.8.13-fr}
Sous les hypothèses de (4.8.11) (resp. (4.8.9), (4.8.10), (4.8.12))
supposons en outre que $X$ soit de type fini sur $k$ (resp. que $X$ soit de
type fini sur $k$ et $\mathscr F$ cohérent, resp. que $X$ soit de type fini
sur $k$ et $\mathscr F$ et $\mathscr G$ cohérents, resp. que $X$ et $Y$
soient de type fini sur $k$). Alors le plus petit corps de définition de
$\overline Z$ (resp. de $\overline{\mathscr H}$, de $\overline u$, de $\overline f$) est
une extension de type fini de $k$.
\end{proposition}

\begin{proof}
On peut se borner à traiter le cas de (4.8.9), et supposer $X$ affine puisque
$X$ est réunion finie d'ouverts affines de type fini sur $k$. Avec les
notations de (4.8.6), tout revient à montrer que si $A$ est une algèbre de
type fini sur $k$ et $M$ un $A$-module de type fini, il y a une sous-extension
$K'$ de $K$ vérifiant les conditions a), b) ou c) et qui soit de type fini sur
$k$ (toute sous-extension d'une extension de type fini étant de type fini).
Or, soit $(m_i)_{1\leq i\leq n}$ un système de générateurs du $A$-module $M$;
comme $A$ est noethérien, il en est de même de $A_{(K)}$, donc $\overline N$
est un $A_{(K)}$-module de type fini, admettant un système fini de générateurs
$n_j$ ($1\leq j\leq r$) tels que
\[
n_j=\sum_i(m_i\otimes1)b_{ij},
\quad\text{où }b_{ij}\in A_{(K)};
\]
chaque $b_{ij}$ s'écrit d'ailleurs
\[
b_{ij}=\sum_h a_{ijh}\otimes c_{ijh},
\quad\text{où }a_{ijh}\in A\text{ et }c_{ijh}\in K;
\]
il est clair que l'extension $K'$ de $k$ engendrée par les $c_{ijh}$ répond à
la question.
\end{proof}

\begin{proposition}[4.8.14]
\label{IV.2.4.8.14-fr}
Soient $k$ un corps, $X$ un $k$-préschéma de type fini, $K$ une extension de
$k$ contenant une extension parfaite de $k$. Alors le plus petit corps de
définition du sous-préschéma fermé $(X_{(K)})_{\mathrm{red}}$ de $X_{(K)}$
est une extension finie radicielle de $k$.
\end{proposition}

\begin{proof}
Si $p$ est l'exposant caractéristique de $k$, l'hypothèse entraîne que
$k'=k^{p^{-\infty}}$ est contenu dans $K$; on sait que
$(X_{(k')})_{\mathrm{red}}$ est géométriquement réduit sur $k'$ (4.6.1), donc
$(X_{(k')})_{\mathrm{red}}\otimes_{k'}K$ est réduit, et par suite égal à
$(X_{(K)})_{\mathrm{red}}$ (I, 5.1.8); en d'autres termes, $k'$ est un corps
de définition de $(X_{(K)})_{\mathrm{red}}$. Comme $k'$ est une extension
algébrique de $k$, la conclusion résulte de (4.8.13).
\end{proof}

\begin{remarks}[4.8.15]
\label{IV.2.4.8.15-fr}
\begin{enumerate}
\item[(i)] La conclusion de (4.8.14) ne subsiste plus nécessairement si l'on
ne suppose pas que $K$ contient une extension parfaite de $k$. Par exemple,
soient $k_0$

\oldpage[IV-2]{84}

un corps de caractéristique $p>0$, $k$ le corps des fractions
rationnelles $k_0(s,t)$ à $2$ indéterminées, $L$ le corps
$k(s^{1/p},t^{1/p})$ et $X=\operatorname{Spec}(L)$. Soient d'autre part $u$
une troisième indéterminée et soit $K$ le corps
$k(u,s^{1/p}+ut^{1/p})$; on vérifie aisément que $k$ est algébriquement fermé
dans $K$ et que $L\otimes_kK$ a un nilradical non nul, donc $X_{(K)}$ n'est
pas réduit; mais comme $(X_{(K)})_{\mathrm{red}}$ n'est pas défini sur $k$,
et que $K$ ne contient aucune sous-extension de degré fini sur $k$ autre que
$k$, le plus petit corps de définition de $(X_{(K)})_{\mathrm{red}}$ ne peut
être fini sur $k$.

\item[(ii)] Étant donné un élément de l'un des ensembles d'arrivée des
applications (4.8.2), et une extension $K_1$ de $K$, il est équivalent de dire
que l'objet considéré admet un plus petit corps de définition ou que son image
par $\varepsilon_{K_1,K}$ admet un plus petit sous-corps de définition (les
deux sous-corps étant nécessairement les mêmes), comme on l'a vu dans
(4.8.4).
\end{enumerate}
\end{remarks}
\subsection{Corps de définition d'une partie d'un préschéma}

\begin{proposition}[4.9.1]
\label{IV.2.4.9.1-fr}
Soient $k$ un corps, $X$ un $k$-préschéma, $K$ une extension de $k$,
$\overline T$ une partie fermée de $X_{(K)}$, $K'$ une sous-extension de $K$.
Les conditions suivantes sont équivalentes:
\begin{enumerate}
\item[a)] $\overline T$ est définie sur $K'$.
\item[b)] La partie ouverte $\overline U=X_{(K)}-\overline T$ de $X_{(K)}$
est définie sur $K'$.
\item[b')] Le sous-préschéma de $X_{(K)}$ induit sur l'ouvert $\overline U$
est défini sur $K'$.
\item[c)] Il existe un sous-préschéma fermé $\overline Z$ de $X_{(K)}$ ayant
$\overline T$ pour espace sous-jacent, et défini sur $K'$.
\end{enumerate}
\end{proposition}

\begin{proof}
Si $g:X_{(K)}\to X_{(K')}$ est la projection canonique, dire que
$\overline T$ est définie sur $K'$ signifie qu'il existe une partie fermée
$T'$ de $X_{(K')}$ telle que $\overline T=g^{-1}(T')$; comme
$g^{-1}(X_{(K')}-T')=X_{(K)}-g^{-1}(T')$ pour toute partie $T'$ de
$X_{(K')}$, a), b) et b') sont équivalentes. Si $\overline Z$, ayant
$\overline T$ pour espace sous-jacent, est défini sur $K'$, on a
$\overline Z=Z'\otimes_{K'}K$, où $Z'$ est un sous-préschéma fermé de
$X_{(K')}$, et si $T'$ est le sous-espace sous-jacent à $Z'$, on sait que
$\overline T=g^{-1}(T')$ (I, 4.4.1), donc $\overline T$ est définie sur
$K'$. Inversement, pour voir que a) entraîne c), il suffit de considérer un
sous-préschéma fermé $Z'$ de $X_{(K')}$ ayant pour espace sous-jacent $T'$
(I, 5.2.1) et de prendre $\overline Z=Z'\otimes_{K'}K$.
\end{proof}

\begin{corollary}[4.9.2]
\label{IV.2.4.9.2-fr}
Avec les notations de (4.9.1), supposons que $K'$ soit un corps de définition
de $\overline T$; toute extension $K''\subset K'$ de $k$, telle que $K'$ soit
extension radicielle de $K''$, est aussi un corps de définition de
$\overline T$.
\end{corollary}

\begin{proof}
Il suffit d'observer que la projection canonique
$g:X_{(K')}\to X_{(K'')}$ est un morphisme entier, surjectif et radiciel,
donc un homéomorphisme (2.4.5).
\end{proof}

\begin{remark}[4.9.3]
\label{IV.2.4.9.3-fr}
Il résulte de (4.9.2) qu'une partie fermée $\overline T$ de $X_{(K)}$ n'admet
pas nécessairement de \emph{plus petit} corps de définition. Par exemple, si
$K=k(t)$, où $t$ est une indéterminée et $k$ un corps parfait de
caractéristique $p>0$, une partie fermée $\overline F$ de $X_{(K)}$ ne peut
avoir de plus petit corps de définition que si elle est déjà définie sur $k$;
en effet, pour tout sous-corps $K'$ de $K$ contenant $k$ et $\ne k$, on a
$K'^p\ne K'$, $K'^p$ contient $k$ et $K'$

\oldpage[IV-2]{85}

est une extension radicielle de $K'^p$, donc si $K'$ est corps de définition
de $\overline F$, il en est de même de $K'^p$. Toutefois:
\end{remark}

\begin{proposition}[4.9.4]
\label{IV.2.4.9.4-fr}
Avec les notations de (4.9.1), supposons que $K$ soit une extension séparable
de $K'$; pour que $K'$ soit un corps de définition de $\overline T$, il faut
et il suffit que $K'$ soit un corps de définition du sous-préschéma réduit de
$X_{(K)}$ ayant $\overline T$ pour espace sous-jacent.
\end{proposition}

\begin{proof}
La condition est évidemment suffisante en vertu de (4.9.1); pour voir qu'elle
est nécessaire, notons que si $\overline T=g^{-1}(T')$, où $T'$ est une
partie fermée de $X_{(K')}$, et si $Z'$ est le sous-préschéma réduit de
$X_{(K')}$ ayant $T'$ pour espace sous-jacent, alors
$\overline Z=Z'\otimes_{K'}K=\operatorname{Spec}(K)\times_{K'}Z'$ est réduit
en vertu du critère (4.6.1, e)) et de l'hypothèse sur $K$ et a pour espace
sous-jacent $\overline T$.
\end{proof}

\begin{corollary}[4.9.5]
\label{IV.2.4.9.5-fr}
Supposons vérifiée l'une des hypothèses suivantes:
\begin{enumerate}
\item[a)] $k$ est de caractéristique $0$.
\item[b)] $K$ est une extension algébrique de $k$ et $X$ est de type fini sur
$k$.
\end{enumerate}
Alors toute partie fermée $\overline T$ de $X_{(K)}$ possède un plus petit
corps de définition, qui est une extension séparable de $k$ (et finie sur $k$
dans l'hypothèse b)).
\end{corollary}

\begin{proof}
Dans l'hypothèse a), $K$ est séparable sur chacun de ses sous-corps, et il
résulte donc de (4.9.4) que les corps de définition de $\overline T$ sont les
mêmes que ceux du sous-préschéma réduit de $X_{(K)}$ ayant $\overline T$ pour
espace sous-jacent. Le corollaire résulte alors de (4.8.11). Dans l'hypothèse
b) on sait que $K$ est une extension radicielle d'une extension séparable
$K_1$ de $k$. Alors, pour toute sous-extension $K'$ de $K$, $K'$ est
radicielle sur $K'\cap K_1$, donc, en vertu de (4.9.2), $K'$ est un corps de
définition de $\overline T$ si et seulement si $K'\cap K_1$ l'est. Cela nous
ramène au cas où $K'=K_1$, donc au cas où $K'$ est une extension algébrique
séparable de $k$, et on achève le raisonnement comme dans le cas a).
\end{proof}

\begin{proposition}[4.9.6]
\label{IV.2.4.9.6-fr}
Soient $k$ un corps, $X$ un $k$-préschéma, $K$ une extension de $k$,
$\overline U$ une partie à la fois ouverte et fermée de $X_{(K)}$,
$\overline V$ son complémentaire. Soit $K'$ une sous-extension de $K$;
alors les conditions suivantes sont équivalentes:
\begin{enumerate}
\item[a)] $K'$ est un corps de définition de la partie $\overline U$ de
$X_{(K)}$.
\item[a')] $K'$ est un corps de définition du sous-préschéma fermé de
$X_{(K)}$ induit sur l'ouvert $\overline U$.
\item[b)] $K'$ est un corps de définition de la partie $\overline V$ de
$X_{(K)}$.
\item[b')] $K'$ est un corps de définition du sous-préschéma fermé de
$X_{(K)}$ induit sur l'ouvert $\overline V$.
\end{enumerate}
\end{proposition}

\begin{proof}
Il est clair que a') entraîne a), et a) entraîne b') en vertu de (4.9.1).
Comme $\overline U$ et $\overline V$ jouent des rôles symétriques, b')
entraîne b) et b) entraîne a'), ce qui achève la démonstration.
\end{proof}

\begin{corollary}[4.9.7]
\label{IV.2.4.9.7-fr}
Sous les hypothèses de (4.9.6), il existe un plus petit corps de définition
$K'$ pour la partie $\overline U$ de $X_{(K)}$, et $K'$ est aussi le plus
petit corps de définition de la partie $\overline V$ de $X_{(K)}$ et des
sous-préschémas fermés de $X_{(K)}$ induits sur les ouverts $\overline U$ et
$\overline V$. Si en outre $X$ est de type fini sur $k$, $K'$ est une
extension finie séparable de $k$.
\end{corollary}

\begin{proof}
Compte tenu de (4.9.6) et (4.8.11), il n'y a à démontrer que la dernière
assertion. En vertu de (4.8.15, (ii)), on peut se borner au cas où $K$ est
algébriquement clos. Soit alors $\overline k$ la clôture algébrique de $k$
contenue dans $K$, et désignons par $W_i$ les composantes

\oldpage[IV-2]{86}

connexes de $X_{(\overline k)}$ ($1\leq i\leq m$); si
$p:X_{(K)}\to X_{(\overline k)}$ est la projection canonique, on sait que les
$p^{-1}(W_i)$ sont les composantes connexes de $X_{(K)}$ (4.4.6), donc
$\overline U$ est une réunion d'un certain nombre de ces composantes, et est
par suite défini sur $\overline k$; la conclusion résulte par suite de
(4.9.5, b)).
\end{proof}
\section{Dimension, profondeur, régularité dans les préschémas localement
noethériens}

Ce paragraphe se borne à reprendre, dans le langage géométrique, et avec
divers compléments de nature technique, les notions et résultats d'Algèbre
commutative exposés dans le chapitre 0, \S\S~16 et 17.

\subsection{Dimension des préschémas}

\begin{env}[5.1.1]
\label{IV.2.5.1.1-fr}
Soient $A$ un anneau, $\mathfrak J$ un idéal de $A$. Rappelons
$(0,16.1)$ que les définitions de la théorie de la dimension d'un anneau
$(0,16.1.1)$ et celles de la théorie de la dimension combinatoire des espaces
topologiques $(0,14.1.2)$ donnent les relations suivantes:
\begin{align}
\dim(\operatorname{Spec}(A))&=\dim(A),
\tag{5.1.1.1}\label{IV.2.5.1.1.1-fr}\\
\dim(V(\mathfrak J))&=\dim(A/\mathfrak J),
\tag{5.1.1.2}\label{IV.2.5.1.1.2-fr}\\
\operatorname{codim}(V(\mathfrak J),\operatorname{Spec}(A))
&=\operatorname{ht}(\mathfrak J).
\tag{5.1.1.3}\label{IV.2.5.1.1.3-fr}
\end{align}
\par\noindent\textup{(5.1.1.4)} $\operatorname{Spec}(A)$ est un espace
caténaire $\Longleftrightarrow A$ est un anneau caténaire.
\par\noindent\textup{(5.1.1.5)} $\operatorname{Spec}(A)$ est
équidimensionnel $\Longleftrightarrow A$ est équidimensionnel
$\Longleftrightarrow$ les idéaux premiers minimaux $\mathfrak p_\alpha$ de
$A$ sont tels que les dimensions des anneaux $A/\mathfrak p_\alpha$ soient
toutes égales.
\par\noindent\textup{(5.1.1.6)} $\operatorname{Spec}(A)$ est
équicodimensionnel $\Longleftrightarrow A$ est équicodimensionnel
$\Longleftrightarrow$ tous les idéaux maximaux de $A$ ont même hauteur.
On rappelle qu'un anneau noethérien $A$ est dit \emph{biéquidimensionnel} si
$\operatorname{Spec}(A)$ est biéquidimensionnel, c'est-à-dire si
$\operatorname{Spec}(A)$ est noethérien et $A$ est équidimensionnel,
équicodimensionnel, caténaire et de dimension finie.
\end{env}

\begin{proposition}[5.1.2]
\label{IV.2.5.1.2-fr}
Soient $X$ un préschéma, $Y$ une partie fermée irréductible de $X$, $y$ le
point générique de $Y$. On a
\begin{equation}
\operatorname{codim}(Y,X)=\dim(\mathcal O_{X,y}).
\tag{5.1.2.1}\label{IV.2.5.1.2.1-fr}
\end{equation}
\end{proposition}

\begin{proof}
En effet (I, 2.4.2) les parties fermées irréductibles de $X$ contenant $y$
sont canoniquement en correspondance biunivoque avec les parties fermées
irréductibles de $\operatorname{Spec}(\mathcal O_{X,y})$, donc sont en
correspondance biunivoque avec les idéaux premiers de $\mathcal O_{X,y}$, et
(5.1.2.1) résulte des définitions.

En particulier, on retrouve ainsi le fait que les points génériques des
composantes irréductibles de $X$ sont les \emph{seuls} points $x\in X$ tels
que $\dim(\mathcal O_x)=0$ (I, 1.1.14).
\end{proof}

\oldpage[IV-2]{87}

\begin{corollary}[5.1.3]
\label{IV.2.5.1.3-fr}
Pour toute partie fermée $Y$ d'un préschéma $X$, on a
\begin{equation}
\operatorname{codim}(Y,X)=\inf_{y\in Y}\dim(\mathcal O_{X,y}).
\tag{5.1.3.1}\label{IV.2.5.1.3.1-fr}
\end{equation}
En outre, si $X$ est localement noethérien, on a, pour tout $x\in Y$
\begin{equation}
\operatorname{codim}_x(Y,X)=
\inf_{\substack{y\in Y,\,x\in\overline{\{y\}}}}
\dim(\mathcal O_{X,y}).
\tag{5.1.3.2}\label{IV.2.5.1.3.2-fr}
\end{equation}
\end{corollary}

\begin{proof}
La relation (5.1.3.2) résulte en effet de (5.1.3.1) et de $(0,14.2.6)$.

Ce corollaire permet de \emph{définir}, pour une partie \emph{quelconque}
$Y$ d'un préschéma $X$, la \emph{codimension} $\operatorname{codim}(Y,X)$ de
$Y$ dans $X$ comme égale au second membre de (5.1.3.1).
\end{proof}

\begin{proposition}[5.1.4]
\label{IV.2.5.1.4-fr}
Pour tout préschéma $X$, on a
\begin{equation}
\dim(X)=\sup_{x\in X}(\dim(\mathcal O_x)).
\tag{5.1.4.1}\label{IV.2.5.1.4.1-fr}
\end{equation}
Si toute partie fermée irréductible de $X$ contient un point fermé, on a aussi
\begin{equation}
\dim(X)=\sup_{x\in F}(\dim(\mathcal O_x))
\tag{5.1.4.2}\label{IV.2.5.1.4.2-fr}
\end{equation}
où $F$ est l'ensemble des points fermés de $X$.
\end{proposition}

\begin{proof}
Il suffit de remarquer (en vertu de (I, 2.4.2)) que les chaînes de parties
fermées irréductibles de $X$ correspondent biunivoquement aux chaînes de
parties fermées irréductibles d'un schéma local de $X$; lorsque toute partie
fermée irréductible de $X$ contient un point fermé, on peut évidemment se
borner aux schémas locaux aux points fermés de $X$.

On notera que toute partie fermée irréductible de $X$ contient un point fermé
lorsque $X$ est \emph{quasi-compact} $(0_{\mathrm I},2.1.3)$; nous verrons un
peu plus loin (5.1.11) qu'il en est de même lorsque $X$ est \emph{localement
noethérien}.
\end{proof}

\begin{corollary}[5.1.5]
\label{IV.2.5.1.5-fr}
Pour qu'un préschéma $X$ soit caténaire, il faut et il suffit que, pour tout
$x\in X$, l'anneau local $\mathcal O_x$ soit caténaire. Si de plus toute
partie fermée irréductible de $X$ contient un point fermé, il suffit, pour que
$X$ soit caténaire, que $\mathcal O_x$ soit caténaire pour tout point fermé
$x$ de $X$.
\end{corollary}

\begin{proof}
Le raisonnement est le même que dans (5.1.4), puisqu'il s'agit de comparer des
chaînes de parties fermées irréductibles ayant mêmes extrémités.
\end{proof}

\begin{env}[5.1.6]
\label{IV.2.5.1.6-fr}
Nous allons maintenant examiner les propriétés plus spéciales liées à des
hypothèses noethériennes.

Rappelons $(0,16.2.3)$ qu'un anneau \emph{local noethérien} $A\ne0$ est de
dimension \emph{finie}, égale au nombre minimum de générateurs d'un idéal de
définition de $A$; pour tout idéal premier $\mathfrak p$ de $A$, la hauteur de
$\mathfrak p$, égale à $\dim(A_{\mathfrak p})$, est donc aussi finie. Ces
propriétés, jointes à (5.1.2), (5.1.3) et (5.1.4), montrent que:
\end{env}

\begin{proposition}[5.1.7]
\label{IV.2.5.1.7-fr}
Pour toute partie fermée non vide $Y$ d'un préschéma localement noethérien
$X$, $\operatorname{codim}(Y,X)$ est finie. Si $X$ est noethérien et affine
et $Y$ une partie fermée irréductible de $X$, $\operatorname{codim}(Y,X)$ est
égal au nombre minimum des sections $s_i$ de $\mathcal O_X$ au-dessus de $X$
telles que $Y$ soit une composante irréductible de l'ensemble des $x\in X$
tels que $s_i(x)=0$ pour tout $i$.
\end{proposition}

\begin{corollary}[5.1.8]
\label{IV.2.5.1.8-fr}
Soient $X$ un préschéma localement noethérien, $\mathcal L$ un
$\mathcal O_X$-Module inversible, $f$ une section de $\mathcal L$ au-dessus
de $X$. Alors toute composante irréductible de l'ensemble $Z$

\oldpage[IV-2]{88}

des $x\in X$ tels que $f(x)=0$ est de codimension $\leq1$ dans $X$; elle est
de codimension $1$ si $Z$ ne contient aucune composante irréductible de $X$.
\end{corollary}

\begin{proof}
On peut se borner au cas où $\mathcal L=\mathcal O_X$. Si $y$ est un point
générique d'une composante irréductible $Y$ de $Z$, l'idéal $(f_y)$ de
$\mathcal O_{X,y}$ doit être tel que $\mathcal O_{X,y}/(f_y)$ n'ait qu'un
seul idéal premier, ce qui signifie que $f_y$ engendre un idéal de définition
de l'anneau local noethérien $\mathcal O_{X,y}$; on a donc
$\operatorname{codim}(Y,X)\leq1$ (5.1.7); si $Z$ ne contient aucune
composante irréductible de $X$, on ne peut avoir
$\operatorname{codim}(Y,X)=0$ en vertu de $(0,14.2.1)$.
\end{proof}

\begin{proposition}[5.1.9]
\label{IV.2.5.1.9-fr}
Soient $X$ un préschéma localement noethérien, $Y$ une partie fermée de $X$.
Si $x\in Y$ est tel que $\mathcal O_{X,x}$ soit un anneau caténaire, on a
\begin{align}
\operatorname{codim}_x(Y,X)
&=\dim(\mathcal O_{X,x})-\operatorname{codim}(\overline{\{x\}},Y)\notag\\
&=\dim(\mathcal O_{X,x})-\dim(\mathcal O_{Y,x}).
\tag{5.1.9.1}\label{IV.2.5.1.9.1-fr}
\end{align}
\end{proposition}

\begin{proof}
Soient $Y_i$ ($1\leq i\leq n$) les composantes irréductibles de $Y$
contenant $x$ (qui sont en nombre fini puisque $Y$ est localement
noethérien), et soit $y_i$ le point générique de $Y_i$. Si l'on pose
$A=\mathcal O_{X,x}$, et si $\mathfrak p_i$ est l'idéal premier de $A$
correspondant à $Y_i\cap\operatorname{Spec}(A)$, les parties fermées
irréductibles de $Y$ contenant $x$ correspondent biunivoquement aux idéaux
premiers de $A$ contenant un des $\mathfrak p_i$, donc
$\dim(\mathcal O_{Y,x})=\sup_i\dim(A/\mathfrak p_i)$; d'autre part, on a
$\mathcal O_{X,y_i}=A_{\mathfrak p_i}$, et l'hypothèse sur $A$ entraîne
$\dim(A)=\dim(A/\mathfrak p_i)+\dim(A_{\mathfrak p_i})$ $(0,16.1.4)$; la
conclusion résulte donc de la relation
$\operatorname{codim}_x(Y,X)=\inf_i(\dim(\mathcal O_{X,y_i}))$ (5.1.2 et
$(0,14.2.6)$).
\end{proof}

\begin{proposition}[5.1.10]
\label{IV.2.5.1.10-fr}
\begin{enumerate}
\item[(i)] Dans tout préschéma $X$, toute partie localement constructible non
vide $E$ contient un point $x$ tel que $\{x\}$ soit une partie localement
fermée de $X$ (ou, ce qui revient au même, tel que $x$ soit isolé dans
$\overline{\{x\}}$).
\item[(ii)] Soient $X$ un préschéma localement noethérien, $x$ un point de
$X$ tel que $\{x\}$ soit localement fermé dans $X$; alors on a
$\dim(\overline{\{x\}})\leq1$; tout point $y\ne x$ de
$\overline{\{x\}}$ est par suite fermé dans $X$.
\end{enumerate}
\end{proposition}

\begin{proof}
(i) Le résultat est un cas particulier du lemme suivant:
\end{proof}

\begin{lemma}[5.1.10.1]
\label{IV.2.5.1.10.1-fr}
Soit $X$ un espace topologique ayant la propriété suivante: pour toute partie
localement fermée $Z\ne\varnothing$ de $X$, il existe une partie $Z'$ de $Z$,
localement fermée dans $X$ (ou dans $Z$, ce qui revient au même) et un point
$x\in Z'$, fermé dans $Z'$. Alors toute partie localement fermée
$Z\ne\varnothing$ de $X$ contient un point $x$ isolé dans
$\overline{\{x\}}$.
\end{lemma}

\begin{proof}
En effet, soit $Z'\subset Z$ une partie localement fermée de $Z$ contenant un
point $x$ tel que $x$ soit fermé \emph{dans} $Z'$. Il y a un voisinage ouvert
$U$ de $x$ dans $X$ tel que $Z'\cap U$ soit fermé dans $U$, donc $x$ est aussi
fermé dans $U$; cela signifie que $U\cap\overline{\{x\}}=\{x\}$, autrement
dit que $x$ est isolé dans $\overline{\{x\}}$.

Le lemme s'applique à tout espace sous-jacent à un préschéma $X$, car $Z$ est
alors aussi espace sous-jacent à un préschéma (I, 5.2.1) et il suffit de
prendre pour $Z'$ un ouvert affine dans $Z$, qui est un espace de Kolmogoroff
quasi-compact, donc contient un point fermé $(0_{\mathrm I},2.1.3)$.
\end{proof}

\oldpage[IV-2]{89}

\begin{proof}
\textup{(ii)} Soit $Z$ le sous-préschéma réduit de $X$ ayant
$\overline{\{x\}}$ pour espace
sous-jacent; l'hypothèse entraîne que $\{x\}$ est ouvert \emph{dans} $Z$;
donc, pour tout $z\in Z$, le point générique $x$ de
$\operatorname{Spec}(\mathcal O_{Z,z})$ est isolé dans
$\operatorname{Spec}(\mathcal O_{Z,z})$; mais l'anneau
$A=\mathcal O_{Z,z}$ est un anneau local noethérien intègre, et l'hypothèse
entraîne qu'il existe un $f\in A$ tel que $A_f$ soit un corps; on sait
$(0,16.3.3)$ que cela entraîne $\dim(A)\leq1$. Comme
$\dim(\mathcal O_{Z,z})\leq1$ pour tout $z\in Z$, on a bien $\dim(Z)\leq1$.
\end{proof}

\begin{corollary}[5.1.11]
\label{IV.2.5.1.11-fr}
Si $X$ est un préschéma localement noethérien, toute partie fermée non vide de
$X$ contient un point fermé.
\end{corollary}

\begin{proof}
En effet, toute partie fermée de $X$ est constructible
$(0_{\mathrm{III}},9.1.1\text{ et }9.1.5)$ et il suffit d'appliquer (5.1.10).
\end{proof}

\begin{env}[5.1.12]
\label{IV.2.5.1.12-fr}
Soient $X$ un préschéma, $\mathcal F$ un $\mathcal O_X$-Module
quasi-cohérent \emph{de type fini}, $S=\operatorname{Supp}(\mathcal F)$ son
support, qui est \emph{fermé} dans $X$ $(0_{\mathrm I},5.2.2)$. Si, pour tout
$x\in X$, on considère $\operatorname{Supp}(\mathcal F_x)$ comme une partie
fermée du schéma local $\operatorname{Spec}(\mathcal O_x)$, on a, par
définition $(0,16.1.7)$,
$\dim(\mathcal F_x)=\dim(\operatorname{Supp}(\mathcal F_x))$; mais on a
\[
\operatorname{Supp}(\mathcal F_x)
=S\cap\operatorname{Spec}(\mathcal O_{X,x})
=\operatorname{Spec}(\mathcal O_{S,x})
\]
où, dans le dernier terme, $S$ désigne un sous-préschéma fermé quelconque de
$X$ ayant $S$ pour espace sous-jacent. Si l'on remarque que les composantes
irréductibles de $\operatorname{Spec}(\mathcal O_{S,x})$ sont les
intersections de $\operatorname{Spec}(\mathcal O_{X,x})$ et des composantes
irréductibles de $S$ contenant $x$, et correspondent biunivoquement à ces
dernières, on voit que
\begin{equation}
\dim(\mathcal F_x)=\dim(\mathcal O_{S,x})
\tag{5.1.12.1}\label{IV.2.5.1.12.1-fr}
\end{equation}
en vertu de (5.1.1.1); il résulte aussi de (5.1.2) que l'on a
\begin{equation}
\dim(\mathcal F_x)=\operatorname{codim}(\overline{\{x\}},S)
\tag{5.1.12.2}\label{IV.2.5.1.12.2-fr}
\end{equation}
si $X$ est localement noethérien.

On dit que $\mathcal F$ est \emph{équidimensionnel au point} $x\in X$ si
$\mathcal F_x$ est un $\mathcal O_{X,x}$-module équidimensionnel,
c'est-à-dire $(0,16.1.7)$ si $\operatorname{Supp}(\mathcal F_x)$ est
équidimensionnel en tant que partie fermée de
$\operatorname{Spec}(\mathcal O_{X,x})$; il revient au même de dire que
l'anneau $\mathcal O_{S,x}$ est équidimensionnel.

On appelle \emph{dimension de} $\mathcal F$ et l'on note $\dim(\mathcal F)$
la dimension du \emph{support} $\operatorname{Supp}(\mathcal F)$; il résulte
de (5.1.4) et (5.1.12.1) que l'on a
\begin{equation}
\dim(\mathcal F)=\sup_{x\in X}\dim(\mathcal F_x).
\tag{5.1.12.3}\label{IV.2.5.1.12.3-fr}
\end{equation}
Si $X=\operatorname{Spec}(A)$ est un schéma affine et si
$\mathcal F=\widetilde M$, où $M$ est un $A$-module de type fini, on a
$\dim(\mathcal F)=\dim(M)$ en vertu de $(0,16.1.7)$ et (5.1.4).
\end{env}

\begin{proposition}[5.1.13]
\label{IV.2.5.1.13-fr}
Soient $X$ un préschéma, $\mathcal F$ un $\mathcal O_X$-Module
quasi-cohérent de type fini, $x$ un point de $X$, $x'$ une générisation de
$x$ dans $X$; on a alors
\begin{equation}
\dim(\mathcal F_{x'})\leq\dim(\mathcal F_x).
\tag{5.1.13.1}\label{IV.2.5.1.13.1-fr}
\end{equation}
\end{proposition}

\begin{proof}
Cela résulte aussitôt de (5.1.12.1) et des définitions.
\end{proof}

\subsection{Dimension d'un préschéma algébrique}

\oldpage[IV-2]{90}

\begin{proposition}[5.2.1]
\label{IV.2.5.2.1-fr}
Soient $k$ un corps, $X$ un préschéma irréductible localement de type fini sur
$k$, $\xi$ son point générique. Alors $X$ est biéquidimensionnel et l'on a
$\dim(X)=\operatorname{deg.tr}_k k(\xi)$.
\end{proposition}

\begin{proof}
Tout anneau local $\mathcal O_x$ est anneau local d'une $k$-algèbre de type
fini, donc quotient d'un anneau local régulier $(0,17.3.9)$, et l'on sait
$(0,16.5.12)$ que ces anneaux sont caténaires; par suite $X$ est un espace
caténaire (5.1.5) et comme $X$ est irréductible, il suffit (5.1.1) de montrer
que pour \emph{tout point fermé} $x\in X$, on a
\begin{equation}
\dim(\mathcal O_x)=\operatorname{deg.tr}_k k(\xi).
\tag{5.2.1.1}\label{IV.2.5.2.1.1-fr}
\end{equation}
On peut évidemment supposer $X$ réduit et affine, donc \emph{intègre}
d'anneau $A$, algèbre de type fini sur $k$. Soit
$n=\operatorname{deg.tr}_k k(\xi)$, $k(\xi)=K$ étant le corps des fractions
de $A$. On sait (Bourbaki, \emph{Alg. comm.}, chap.~V, \S~3,
n\textsuperscript{o}~1, th.~1) qu'il existe une sous-$k$-algèbre
$B=k[t_1,\ldots,t_n]$ de $A$, où les $t_i$ sont algébriquement indépendants
sur $k$, telle que $A$ soit une $B$-algèbre \emph{finie}. Soit
$\mathfrak m=\mathfrak j_x$, qui est par hypothèse un idéal maximal de $A$;
$\mathfrak n=B\cap\mathfrak m$ est donc un idéal maximal de $B$ (Bourbaki,
\emph{Alg. comm.}, chap.~V, \S~2, n\textsuperscript{o}~1, prop.~1), et
$A_{\mathfrak m}$ est un anneau local de la $B_{\mathfrak n}$-algèbre
\emph{finie} $S^{-1}A$, où $S=B-\mathfrak n$; comme $B_{\mathfrak n}$ est
intégralement clos et $S^{-1}A$ intègre, on a
$\dim(A_{\mathfrak m})=\dim(B_{\mathfrak n})$ $(0,16.1.6)$. On peut donc se
borner au cas où $A=k[t_1,\ldots,t_n]$; on sait alors que
$k'=k(x)$ est une extension \emph{finie} de $k$ (I, 6.4.2); soit
$A'=k'[t_1,\ldots,t_n]$; compte tenu de (I, 2.4.6 et 3.3.14), il existe un
idéal maximal $\mathfrak m'$ de $A'$ au-dessus de $\mathfrak m$ et tel que
$k'$ soit le corps résiduel de $A'_{\mathfrak m'}$. Comme $A'$ est intègre et
est une $A$-algèbre finie, le même raisonnement que ci-dessus montre que
$\dim(A'_{\mathfrak m'})=\dim(A_{\mathfrak m})$, si bien qu'on est ramené au
cas où $A/\mathfrak m=k$. Or dans ce cas $\mathfrak m$ est engendré par des
polynômes $t_i-a_i$ (où $a_i\in k$, $1\leq i\leq n$, les $a_i$ étant les
images canoniques des $t_i$ dans $A/\mathfrak m$). Remplaçant $t_i$ par
$t_i-a_i$, on voit finalement qu'on est ramené au cas où
$\mathfrak m=\sum_{i=1}^nAt_i$. Le complété de l'anneau local
$A_{\mathfrak m}$ est alors l'anneau de séries formelles
$k[[t_1,\ldots,t_n]]$; on sait qu'il a même dimension que $A_{\mathfrak m}$
$(0,16.2.4)$, et d'autre part la dimension de $k[[t_1,\ldots,t_n]]$ est égale
à $n$ $(0,17.1.4)$; d'où la conclusion.
\end{proof}

\begin{corollary}[5.2.2]
\label{IV.2.5.2.2-fr}
Pour un préschéma $X$ localement de type fini sur un corps $k$, $\dim(X)$
coïncide avec le nombre défini dans (4.1.1).
\end{corollary}

\begin{proof}
En effet, si $X_\alpha$ sont les sous-préschémas réduits ayant pour espaces
sous-jacents les composantes irréductibles de $X$, on a
$\dim(X)=\sup_\alpha(\dim(X_\alpha))$ $(0,14.1.2.1)$ et il suffit d'appliquer
(5.2.1) aux $X_\alpha$.
\end{proof}

\begin{corollary}[5.2.3]
\EGAReaderTarget{ega4.c04.cor.5-2-3}{IV.2.5.2.3-fr}
Soient $X$ un préschéma localement de type fini sur un corps $k$, $x$ un
point de $X$. On a
\begin{equation}
\dim_x(X)=\dim(\mathcal O_x)+\operatorname{deg.tr}_k k(x).
\tag{5.2.3.1}\label{IV.2.5.2.3.1-fr}
\end{equation}
\end{corollary}

\begin{proof}
On sait qu'il existe un voisinage ouvert $U$ de $x$ dans $X$ tel que
$\dim_x(X)=\dim(U)$ $(0,14.1.4.1)$, et l'on peut supposer que les composantes
irréductibles de $U$ soient les $X_i\cap U$, où les $X_i$ sont les
composantes irréductibles de $X$ contenant $x$; comme $U\cap X_i$

\oldpage[IV-2]{91}

est dense dans $X_i$, il résulte de (4.1.1.3) que
$\dim(X_i)=\dim(U\cap X_i)$, donc on a
$\dim_x(X)=\sup_i(\dim(X_i))$. En outre, les idéaux premiers minimaux de
$\mathcal O_x$ correspondent aux points génériques des $X_i$, donc
$(0,16.1.1.1)$, on a
$\dim(\mathcal O_{X,x})=\sup_i(\dim(\mathcal O_{X_i,x}))$. On est donc ramené
au cas où $X$ est \emph{irréductible}; comme $X$ est biéquidimensionnel par
(5.2.1), on a
$\dim(X)=\dim(\overline{\{x\}})+\operatorname{codim}(\overline{\{x\}},X)$
$(0,14.3.5.1)$, et l'on sait que
$\dim(\overline{\{x\}})=\operatorname{deg.tr}_k k(x)$ par (5.2.1) et
$\operatorname{codim}(\overline{\{x\}},X)=\dim(\mathcal O_x)$ par (5.1.2).
\end{proof}

\begin{corollary}[5.2.4]
\label{IV.2.5.2.4-fr}
Soient $X$ un préschéma localement de type fini sur un corps $k$, $\mathcal L$
un $\mathcal O_X$-Module inversible, $f$ une section de $\mathcal L$ au-dessus
de $X$, telle que l'ensemble $Y$ des $x\in X$ tels que $f(x)=0$ soit rare
dans $X$. Alors on a $\dim(Y)\leq\dim(X)-1$; les deux membres de cette
inégalité sont égaux si $Y$ rencontre toute composante irréductible de
dimension maxima de $X$.
\end{corollary}

\begin{proof}
Si $(X_\alpha)$ est la famille des composantes irréductibles de $X$, on a
\[
\dim(Y)=\sup_\alpha(\dim(Y\cap X_\alpha))
\]
et l'on est donc ramené au cas où $X$ est irréductible ($Y\cap X_\alpha$
étant rare dans $X_\alpha$ puisque tout $X_\alpha$ a un intérieur non vide
dans $X$). On peut se borner au cas où $Y\ne\varnothing$; alors, pour tout
point maximal $x$ de $Y$, on a (puisque $X$ est biéquidimensionnel)
$\dim(\overline{\{x\}})=\dim(X)-\operatorname{codim}(\overline{\{x\}},X)$,
et puisque $Y$ est rare dans $X$, on a
$\operatorname{codim}(\overline{\{x\}},X)=1$ par (5.1.8); d'où le corollaire.
\end{proof}

\begin{remarks}[5.2.5]
\label{IV.2.5.2.5-fr}
\begin{enumerate}
\item[(i)] Contrairement à ce qui a lieu pour les $k$-préschémas algébriques,
un préschéma localement noethérien $X$ n'est pas nécessairement caténaire
(cf. (5.6.11)); il l'est cependant si tous ses anneaux locaux $\mathcal O_x$
sont des quotients d'anneaux réguliers $(0,16.5.12)$ et en particulier si $X$
est \emph{régulier} (I, 4.1.4). Toutefois, même un schéma (intègre)
$X=\operatorname{Spec}(A)$, où $A$ est un anneau \emph{régulier}, n'est pas
nécessairement \emph{biéquidimensionnel}, autrement dit $(0,14.3.3)$ les
codimensions dans $X$ des divers points fermés de $X$ ne sont pas
nécessairement les mêmes. Par exemple, soient $B$ un anneau de valuation
discrète, $\mathfrak m=B\pi$ son idéal maximal, $k=B/\mathfrak m$ son corps
résiduel, $K$ le corps des fractions de $B$; soit $A$ l'anneau de polynômes
$B[T]$. Il y a dans $A$ des idéaux maximaux \emph{de hauteur} $2$, par exemple
$(\pi)+(T)$; mais il y a aussi des idéaux maximaux \emph{de hauteur} $1$, par
exemple l'idéal principal $(\pi T-1)$: en effet, un idéal premier principal
$\ne0$ est de hauteur $1$ (5.1.8); d'autre part
$A/(\pi T-1)$ est isomorphe à l'anneau de fractions $B_\pi$
$(0_{\mathrm I},1.2.3)$, qui n'est autre ici que $K$, ce qui prouve que
l'idéal $(\pi T-1)$ est maximal et de hauteur $1$.

\item[(ii)] Lorsque $X$ est un préschéma localement de type fini sur un corps,
on a vu (4.1.1.3) que $\dim(U)=\dim(X)$ pour tout ouvert partout dense $U$
dans $X$. Ce résultat \emph{ne s'étend pas} aux préschémas noethériens
réguliers, même s'ils sont biéquidimensionnels. Par exemple, soient $B$ un
anneau de valuation discrète, $\mathfrak m$ son idéal maximal;
$X=\operatorname{Spec}(A)$ a deux points, l'idéal $(0)$ et $\mathfrak m$, ce
dernier étant le seul point fermé; on a $\dim(X)=1$, mais l'ensemble ouvert
$U=\{(0)\}$ est de dimension $0$ (cf. \S~10).
\end{enumerate}
\end{remarks}

\subsection{Dimension du support d'un Module et polynôme de Hilbert}
\label{IV.2.s5.3-fr}

\oldpage[IV-2]{92}

Ce numéro utilise les résultats du chapitre III; il ne sera pas utilisé dans la
suite de ce chapitre.

\begin{proposition}[5.3.1]
\label{IV.2.s5.3.1-fr}
Soient $A$ un anneau local artinien, $X$ un schéma projectif sur
$\operatorname{Spec}(A)$, $\mathcal L$ un $\mathcal O_X$-Module inversible
très ample relativement à $A$, $\mathcal F$ un $\mathcal O_X$-Module cohérent
$\ne0$; on pose
$\mathcal F(n)=\mathcal F\otimes_{\mathcal O_X}\mathcal L^{\otimes n}$ pour
$n\in\mathbf Z$. Alors le degré du polynôme de Hilbert
$P(n)=\chi_A(\mathcal F(n))$ de $\mathcal F$ par rapport à $A$ (III, 2.5.3)
est égal à la dimension de $\operatorname{Supp}(\mathcal F)$.
\end{proposition}

\begin{proof}
Raisonnons par récurrence sur
$d=\dim(\operatorname{Supp}(\mathcal F))$. On sait qu'il existe un
sous-préschéma fermé $Y$ de $X$ dont $\operatorname{Supp}(\mathcal F)$ est
l'espace sous-jacent, et un $\mathcal O_Y$-Module cohérent $\mathcal G$ tel
que $\mathcal F=j_*(\mathcal G)$, où $j:Y\to X$ est l'injection canonique
$(\mathbf{Err}_{\mathrm{III}},30)$. Il est immédiat que les polynômes de
Hilbert de $\mathcal F$ et de $\mathcal G$ sont les mêmes, donc on peut se
borner au cas où $X=\operatorname{Supp}(\mathcal F)$. Supposons d'abord
$d=0$; tous les points de $X$ étant fermés, $X$ est un schéma artinien
(I, 6.2.2), donc $\mathcal F(n)=\mathcal F$ pour tout entier $n$, et l'on a
par suite (III, 2.5.3)
\[
\chi_A(\mathcal F(n))=\operatorname{long}_A(\Gamma(X,\mathcal F))
\]
pour $n$ assez grand, et le degré du polynôme de Hilbert est donc $0$.
Supposons maintenant $d>0$, et posons
$Z=\operatorname{Ass}(\mathcal F)$, qui est un ensemble fini (3.1.6); il
existe un entier $m>0$ et une section
$f\in\Gamma(X,\mathcal L^{\otimes m})$ tels que $X_f$ soit un voisinage de
$Z$ (II, 4.5.4). La multiplication par $f$ est un homomorphisme
\[
\mu_f:\mathcal F\longrightarrow\mathcal F(m)
\]
qui est injectif (3.1.8); on a donc une suite exacte
\begin{equation}
0\longrightarrow\mathcal F\xrightarrow{\mu_f}\mathcal F(m)
\longrightarrow\mathcal G\longrightarrow0
\tag{5.3.1.1}\label{IV.2.s5.3.1.1-fr}
\end{equation}
où $\mathcal G$ est cohérent. En vertu du lemme de Nakayama, les points
$x\in\operatorname{Supp}(\mathcal G)$ sont exactement ceux pour lesquels
$f(x)=0$. Nous allons en déduire que l'on a
\begin{equation}
\dim(\operatorname{Supp}(\mathcal G))=d-1.
\tag{5.3.1.2}\label{IV.2.s5.3.1.2-fr}
\end{equation}
Cela résultera du lemme suivant:
\end{proof}

\begin{lemma}[5.3.1.3]
\label{IV.2.s5.3.1.3-fr}
Soient $A$ un anneau local artinien, $X$ un schéma projectif sur
$\operatorname{Spec}(A)$, $\mathcal L$ un $\mathcal O_X$-Module inversible
ample relativement à $A$; alors, pour toute section $g$ de $\mathcal L$
au-dessus de $X$, l'ensemble $X_g$ des $x\in X$ tels que $g(x)=0$ rencontre
toute composante irréductible de $X$ de dimension $\ne0$.
\end{lemma}

\begin{proof}
En effet (II, 5.5.7), l'ensemble $X_g$ est un ouvert affine de $X$, et s'il
contient une composante irréductible $X'$ de $X$, le sous-préschéma fermé
réduit de $X$ ayant $X'$ pour espace sous-jacent est à la fois projectif et
affine sur $A$, donc fini sur $A$ (III, 4.4.2), et par suite un schéma
artinien, donc de dimension $0$.
\end{proof}

\begin{proof}
Ce lemme étant établi, notons que puisque $Z$ contient les points maximaux de
$X$, $X_f$ est dense, donc $\operatorname{Supp}(\mathcal G)$ est rare dans
$X$, et la relation (5.3.1.2) résulte du lemme et de (5.2.4).

\oldpage[IV-2]{93}

Cela étant, de la suite exacte (5.3.1.1), on déduit, pour tout $n$, la suite
exacte
\[
0\longrightarrow\mathcal F(n)\longrightarrow\mathcal F(n+m)
\longrightarrow\mathcal G(n)\longrightarrow0
\]
et pour $n$ assez grand, on a donc aussi la suite exacte (III, 2.2.3)
\[
0\longrightarrow\Gamma(X,\mathcal F(n))
\longrightarrow\Gamma(X,\mathcal F(n+m))
\longrightarrow\Gamma(X,\mathcal G(n))\longrightarrow0,
\]
d'où, compte tenu de (III, 2.5.3), pour $n$ assez grand,
\[
\chi_A(\mathcal G(n))
=\chi_A(\mathcal F(n+m))-\chi_A(\mathcal F(n)).
\]
Comme en vertu de (5.3.1.2) et de l'hypothèse de récurrence, le degré du
polynôme $\chi_A(\mathcal G(n))$ est $d-1$, la relation précédente entraîne
que le degré de $\chi_A(\mathcal F(n))$ est $d$, C.Q.F.D.
\end{proof}

\subsection{Dimension de l'image d'un morphisme}
\label{IV.2.s5.4-fr}

\begin{proposition}[5.4.1]
\label{IV.2.s5.4.1-fr}
Soient $X$, $Y$ deux préschémas localement noethériens, $f:X\to Y$ un
morphisme.
\begin{enumerate}
\item[(i)] Si $f$ est quasi-fini, on a
$\dim(X)\leq\dim(\overline{f(X)})\leq\dim(Y)$.
\item[(ii)] Si $f$ est surjectif et ouvert (resp. surjectif et fermé), on a
$\dim(X)\geq\dim(Y)$.
\end{enumerate}
\end{proposition}

\begin{proof}
\begin{enumerate}
\item[(i)] On peut remplacer $f$ par $f_{\mathrm{red}}$ (II, 6.2.4), donc
supposer $X$ et $Y$ réduits; si $Z$ est le sous-préschéma fermé réduit de $Y$
ayant pour espace sous-jacent $\overline{f(X)}$, on a alors $f=j\circ g$, où
$j:Z\to Y$ est l'injection canonique et $g:X\to Z$ est un morphisme
quasi-fini (I, 5.2.2 et II, 6.2.4). On peut donc se borner au cas où $f(X)$
est dense dans $Y$. Pour tout $x\in X$, $\mathcal O_x$ est alors un
$\mathcal O_{f(x)}$-module quasi-fini (II, 6.2.2), et par suite
$\mathfrak m_{f(x)}\mathcal O_x$ est un idéal de définition de $\mathcal O_x$
$(0_{\mathrm I},7.4.4)$; mais on sait (0, 16.3.10) que si $A\to B$ est un
homomorphisme local d'anneaux locaux noethériens, tel que, si $\mathfrak m$
est l'idéal maximal de $A$, $\mathfrak mB$ soit un idéal de définition de
$B$, alors on a $\dim(B)\leq\dim(A)$; cela achève de démontrer (i) en vertu
de (5.1.4).

\item[(ii)] La définition de la dimension (0, 14.1.2) montre qu'il suffit de
prouver que pour toute suite $(y_i)_{0\leq i\leq n}$ d'éléments distincts de
$Y$ telle que $y_i\in\overline{\{y_{i+1}\}}$ pour $0\leq i\leq n-1$, il
existe une suite $(x_i)_{0\leq i\leq n}$ de points de $X$ telle que
$x_i\in\overline{\{x_{i+1}\}}$ pour $0\leq i\leq n-1$ et $f(x_i)=y_i$
pour tout $i$. Supposons d'abord $f$ surjectif et ouvert et démontrons
l'existence de $x_i$ par récurrence sur $i$; l'existence de $x_0\in X$ tel
que $f(x_0)=y_0$ résulte de ce que $f$ est surjectif. Si les $x_i$ ont été
déterminés pour $i\leq m$ de façon à satisfaire à $f(x_i)=y_i$ pour
$i\leq m$ et $x_i\in\overline{\{x_{i+1}\}}$ pour $i<m$, on note que puisque
$f$ est ouvert et $y_{m+1}$ une générisation de $y_m$, il existe
$x_{m+1}\in f^{-1}(y_{m+1})$ qui est une générisation de $x_m$ (1.10.3) et
la récurrence peut se poursuivre.

\oldpage[IV-2]{94}

Supposons maintenant $f$ surjectif et fermé et démontrons l'existence de
$x_i$ par récurrence descendante sur $i$, l'existence de $x_n$ tel que
$f(x_n)=y_n$ résultant encore de ce que $f$ est surjectif. Si les $x_i$ ont
été déterminés de façon à satisfaire aux conditions voulues pour $i>m$, on
note que $f(\overline{\{x_{m+1}\}})$ est l'adhérence de
$\{f(x_{m+1})\}=\{y_{m+1}\}$ puisque $f$ est fermé (Bourbaki,
\emph{Top. gén.}, chap. I\textsuperscript{er}, 3\textsuperscript{e} éd.,
\S~5, n\textsuperscript{o}~4, prop.~9); il existe donc
$x_m\in\overline{\{x_{m+1}\}}$ tel que $f(x_m)=y_m$ et la récurrence
descendante se poursuit encore.
\end{enumerate}
\end{proof}

\begin{corollary}[5.4.2]
\label{IV.2.s5.4.2-fr}
Si $X$, $Y$ sont deux préschémas localement noethériens, $f:X\to Y$ un
morphisme fini (donc fermé), on a
$\dim(X)=\dim(\overline{f(X)})$. Si en outre $f$ est surjectif, on a
$\dim(X)=\dim(Y)$.
\end{corollary}

\begin{remarks}[5.4.3]
\label{IV.2.s5.4.3-fr}
\begin{enumerate}
\item[(i)] On a vu dans (4.1.2) que si $X$ et $Y$ sont des préschémas
localement de type fini sur un corps $k$ et $f$ un $k$-morphisme, l'inégalité
$\dim(Y)\leq\dim(X)$ est déjà vérifiée lorsque $f$ est quasi-compact et
dominant. Au contraire, on peut avoir $\dim(Y)>\dim(X)$ même lorsque $f$ est
de type fini, bijectif et une immersion locale, si l'on suppose seulement
$X$ et $Y$ localement noethériens. Par exemple soient $A$ un anneau de
valuation discrète, $K$ son corps des fractions, $k$ son corps résiduel; si
$Y=\operatorname{Spec}(A)$, et si $a$ est le point générique de $Y$ et $b$
son point fermé, on a donc $k(a)=K$, $k(b)=k$. Soit alors $X$ le préschéma
somme de $\operatorname{Spec}(K)$ et de $\operatorname{Spec}(k)$,
$f:X\to Y$ le morphisme qui coïncide dans $\operatorname{Spec}(K)$ et
$\operatorname{Spec}(k)$ avec les morphismes canoniques respectifs
(I, 2.4.5); il est clair que $f$ est une immersion locale bijective, $\{a\}$
étant ouvert dans $Y$; d'autre part, $f$ est de type fini, car
$K=A[\pi^{-1}]$, où $\pi$ est une uniformisante de $A$, donc $K$ est une
$A$-algèbre de type fini. Cependant on a $\dim(X)=0$ et $\dim(Y)=1$.

\item[(ii)] Si $A$ et $B$ sont deux anneaux noethériens tels que $A\subset B$
et que $B$ soit une $A$-algèbre finie, le corollaire (5.4.2) montre à nouveau
que $\dim(B)=\dim(A)$ (0, 16.1.5). Supposons en outre que $A$ soit un anneau
local noethérien; alors $B$ est un anneau noethérien semi-local (Bourbaki,
\emph{Alg. comm.}, chap. IV, \S~2, n\textsuperscript{o}~5, cor.~3 de la
prop.~9); si $\mathfrak n_i$ $(1\leq i\leq r)$ sont les idéaux maximaux de
$B$, on a donc
\begin{equation}
\dim(A)=\sup_i\dim(B_{\mathfrak n_i}).
\tag{5.4.3.1}\label{IV.2.s5.4.3.1-fr}
\end{equation}
On observera que, dans ces conditions, on n'a pas nécessairement
$\dim(B_{\mathfrak n_i})=\dim(A)$ pour tout $i$: il suffit pour le voir de
remplacer $B$ par la composée directe de $B$ et du corps résiduel $k$ de
$A$. Mais on a même des exemples où $A$ et $B$ sont en outre des anneaux
intègres de dimension $2$ et où certains des $B_{\mathfrak n_i}$ ont une
dimension $<\dim(A)$ (5.6.11). Nous verrons toutefois plus loin ((5.6.4) et
(5.6.10)) que ce dernier phénomène ne peut se présenter lorsqu'on suppose
que $A$ est un quotient d'anneau local régulier.
\end{enumerate}
\end{remarks}

\subsection{Formule des dimensions pour un morphisme de type fini}
\label{IV.2.s5.5-fr}

\begin{env}[5.5.1]
\label{IV.2.s5.5.1-fr}
Rappelons (0, 16.3.9) que si $A$, $B$ sont deux anneaux locaux noethériens,
$k$ le corps résiduel de $A$, $\varphi:A\to B$ un homomorphisme local, on a
\begin{equation}
\dim(B)\leq\dim(A)+\dim(B\otimes_A k).
\tag{5.5.1.1}\label{IV.2.s5.5.1.1-fr}
\end{equation}
\end{env}

On en déduit:

\begin{proposition}[5.5.2]
\label{IV.2.s5.5.2-fr}
Soient $X$, $Y$ deux préschémas localement noethériens, $f:X\to Y$ un
morphisme, $x$ un point de $X$, $y=f(x)$. Alors on a
\begin{equation}
\dim(\mathcal O_x)\leq\dim(\mathcal O_y)
+\dim(\mathcal O_x\otimes_{\mathcal O_y}k(y)).
\tag{5.5.2.1}\label{IV.2.s5.5.2.1-fr}
\end{equation}
En particulier, si $x$ est point maximal de la fibre $f^{-1}(y)$, on a
\begin{equation}
\dim(\mathcal O_x)\leq\dim(\mathcal O_y)
\tag{5.5.2.2}\label{IV.2.s5.5.2.2-fr}
\end{equation}
puisque $\mathcal O_x\otimes_{\mathcal O_y}k(y)
=\mathcal O_x/\mathfrak m_y\mathcal O_x$ est l'anneau local de $x$ dans le
préschéma $f^{-1}(y)$, et est donc de dimension $0$ par hypothèse.
\end{proposition}

\oldpage[IV-2]{95}

Nous allons obtenir une formule plus précise que (5.5.2.1) lorsque l'on
suppose que $f$ est un morphisme de type fini.

\begin{proposition}[5.5.3]
\label{IV.2.s5.5.3-fr}
Soient $A$ un anneau noethérien, $T$ une indéterminée, $\mathfrak p$ un idéal
premier de $A$; alors $\mathfrak p'=\mathfrak pA[T]$ est premier dans
$B=A[T]$ et $\mathfrak p'\cap A=\mathfrak p$. Il existe une infinité
d'idéaux premiers de $B$ distincts de $\mathfrak p'$, dont l'intersection
avec $A$ est $\mathfrak p$; ces idéaux sont deux à deux sans relation
d'inclusion. En outre, si $\mathfrak q$ est un tel idéal, on a
\begin{equation}
\dim(B_{\mathfrak q})=\dim(B_{\mathfrak p'})+1
=\dim(A_{\mathfrak p})+1.
\tag{5.5.3.1}\label{IV.2.s5.5.3.1-fr}
\end{equation}
\end{proposition}

\begin{proof}
Dans les premières assertions, on se ramène aussitôt, en remplaçant $A$ par
$A/\mathfrak p$ et observant que
$A[T]/\mathfrak pA[T]=(A/\mathfrak p)[T]$, au cas $\mathfrak p=0$; elles
résultent alors du fait que les idéaux premiers de $B=A[T]$ dont
l'intersection avec $A$ se réduit à $0$ sont exactement ceux qui ne
rencontrent pas la partie multiplicative $S=A-\{0\}$ de l'anneau intègre
$A$; or on sait qu'il y a une bijection croissante de l'ensemble de ces
idéaux sur l'ensemble des idéaux premiers de $S^{-1}A[T]=K[T]$, où $K$ est
le corps des fractions de $A$ (Bourbaki, \emph{Alg. comm.}, chap. II,
\S~2, n\textsuperscript{o}~5, prop.~11). En outre, on a, d'après (5.5.1.2),
\[
\dim(B_{\mathfrak q})\leq\dim(A_{\mathfrak p})
+\dim(B_{\mathfrak q}/\mathfrak pB_{\mathfrak q}),
\]
et si $k$ est le corps des fractions de $A/\mathfrak p$,
$B_{\mathfrak q}/\mathfrak pB_{\mathfrak q}$ s'identifie canoniquement à
$(k[T])_{\mathfrak q}$, donc est un anneau de valuation discrète, donc de
dimension $1$. Enfin, si
$\mathfrak p=\mathfrak p_0\supset\mathfrak p_1\supset\cdots
\supset\mathfrak p_m$ est une chaîne d'idéaux premiers de $A$ de longueur
maxima, les idéaux $\mathfrak p_jB$ $(0\leq j\leq m)$ sont premiers dans
$B$, deux à deux distincts et contenus dans $\mathfrak q$; donc
$\dim(B_{\mathfrak q})\geq m+1=\dim(A_{\mathfrak p})+1$ et par suite
$\dim(B_{\mathfrak q})=\dim(A_{\mathfrak p})+1$. Cette relation s'écrit
encore $\operatorname{ht}(\mathfrak q)=\operatorname{ht}(\mathfrak p)+1$;
comme $\mathfrak q\ne\mathfrak p'$, on a d'ailleurs
$\operatorname{ht}(\mathfrak q)\geq\operatorname{ht}(\mathfrak p')+1
\geq\operatorname{ht}(\mathfrak p)+1$ par définition de la hauteur d'un
idéal premier; cela achève de prouver (5.5.3.1).
\end{proof}

\begin{corollary}[5.5.4]
\EGAReaderTarget{ega4.c04.cor.5-5-4}{IV.2.s5.5.4-fr}
Pour tout anneau noethérien $A$, on a
\begin{equation}
\dim(A[T_1,\ldots,T_r])=\dim(A)+r
\tag{5.5.4.1}\label{IV.2.s5.5.4.1-fr}
\end{equation}
($T_i$ indéterminées).
\end{corollary}

On notera par contre qu'il y a des exemples d'anneaux locaux non noethériens
$A$ tels que $\dim(A)=1$ et $\dim(A[T])=3$ [30].

\begin{corollary}[5.5.5]
\label{IV.2.s5.5.5-fr}
Pour tout préschéma localement noethérien $X$, la dimension de
$X[T_1,\ldots,T_r]=X\otimes_{\mathbf Z}\mathbf Z[T_1,\ldots,T_r]$
($T_i$ indéterminées) est $\dim(X)+r$.
\end{corollary}

Cela résulte de (5.5.4) et de (0, 14.1.7).

\begin{corollary}[5.5.6]
\label{IV.2.s5.5.6-fr}
Sous les hypothèses de (5.5.3), soit $\mathfrak q$ un idéal premier de
$B=A[T]$ tel que $\mathfrak q\cap A=\mathfrak p$; si $k$ et $k'$ sont les
corps résiduels de $A_{\mathfrak p}$ et $B_{\mathfrak q}$ respectivement,
on a
\begin{equation}
\dim(A_{\mathfrak p})+1
=\dim(B_{\mathfrak q})+\operatorname{deg.tr.}_k k'.
\tag{5.5.6.1}\label{IV.2.s5.5.6.1-fr}
\end{equation}
\end{corollary}

\begin{proof}
\oldpage[IV-2]{96}

Si $\mathfrak q=\mathfrak pB$, on a, d'après (5.5.3.1),
$\dim(B_{\mathfrak q})=\dim(A_{\mathfrak p})$, et dans ce cas $k'=k(T)$,
donc la formule (5.5.6.1) est bien vérifiée. Dans le cas contraire,
$\dim(B_{\mathfrak q})=\dim(A_{\mathfrak p})+1$, et comme $\mathfrak q$
correspond à un idéal premier $\mathfrak q'\ne0$ de $k[T]$, $k'$ est une
extension algébrique de $k$, donc on a encore la formule (5.5.6.1).
\end{proof}

\begin{lemma}[5.5.7]
\label{IV.2.s5.5.7-fr}
Soient $A$ un anneau local intègre noethérien, $\mathfrak m$ son idéal
maximal, $k$ son corps résiduel.
\begin{enumerate}
\item[(i)] Pour tout anneau intègre $B$ contenant $A$, tel que $B=A[x]$
(pour un $x\in B$) et tout idéal premier $\mathfrak q$ de $B$ tel que
$\mathfrak q\cap A=\mathfrak m$, on a
\begin{equation}
\dim(A)+\operatorname{deg.tr.}_A B
\geq\dim(B_{\mathfrak q})+\operatorname{deg.tr.}_k k'
\tag{5.5.7.1}\label{IV.2.s5.5.7.1-fr}
\end{equation}
en désignant par $k'$ le corps résiduel de $B_{\mathfrak q}$ et par
$\operatorname{deg.tr.}_A B$ le degré de transcendance du corps des
fractions de $B$ sur celui de $A$.
\item[(ii)] Supposons que pour tout idéal maximal $\mathfrak n$ de $A[T]$
tel que $\mathfrak n\cap A=\mathfrak m$, l'anneau $(A[T])_{\mathfrak n}$
soit caténaire; alors les deux membres de (5.5.7.1) sont égaux.
\end{enumerate}
\end{lemma}

\begin{proof}
\begin{enumerate}
\item[(i)] Si $x$ est transcendant sur le corps des fractions de $A$, on a
$\operatorname{deg.tr.}_A B=1$ et les deux membres de (5.5.7) sont égaux en
vertu de (5.5.6). Dans le cas contraire, on a $B=A[T]/\mathfrak p$, où
$\mathfrak p$ est un idéal premier $\ne0$ de $A[T]$, tel que
$\mathfrak p\cap A=0$ puisque $B$ contient $A$; on a donc
$\operatorname{ht}(\mathfrak p)=1$ en vertu de (5.5.3). L'idéal
$\mathfrak q$ de $B$ est de la forme $\mathfrak n/\mathfrak p$, où
$\mathfrak n\supset\mathfrak p$ est un idéal premier de $A[T]$ tel que
$\mathfrak n\cap A=\mathfrak m$, et l'on a
$B_{\mathfrak q}=(A[T])_{\mathfrak n}/\mathfrak p(A[T])_{\mathfrak n}$;
la formule (0, 16.1.4.1), appliquée à
$X=\operatorname{Spec}((A[T])_{\mathfrak n})$ et à
$Y=\operatorname{Spec}(B_{\mathfrak q})$, donne
\begin{equation}
\dim((A[T])_{\mathfrak n})
\geq\operatorname{ht}(\mathfrak p(A[T])_{\mathfrak n})
+\dim(B_{\mathfrak q})=1+\dim(B_{\mathfrak q})
\tag{5.5.7.2}\label{IV.2.s5.5.7.2-fr}
\end{equation}
car $\operatorname{ht}(\mathfrak p(A[T])_{\mathfrak n})
=\operatorname{ht}(\mathfrak p)=1$ en vertu de la correspondance biunivoque
entre idéaux premiers de $A[T]$ contenus dans $\mathfrak n$ et idéaux
premiers de $(A[T])_{\mathfrak n}$. Enfin, la formule (5.5.6.1) donne
\begin{equation}
\dim((A[T])_{\mathfrak n})
=\dim(A)+1-\operatorname{deg.tr.}_k k'
\tag{5.5.7.3}\label{IV.2.s5.5.7.3-fr}
\end{equation}
puisque les corps résiduels de $B_{\mathfrak q}$ et de
$(A[T])_{\mathfrak n}$ sont les mêmes; par ailleurs, on a alors
$\operatorname{deg.tr.}_A B=0$, ce qui achève de prouver (5.5.7.1).
\item[(ii)] Si $(A[T])_{\mathfrak n}$ est caténaire, les deux membres de
(5.5.7.2) sont égaux (0, 16.1.4), donc aussi les deux membres de (5.5.7.1).
\end{enumerate}
\end{proof}

\begin{theorem}[5.5.8 (formule des dimensions)]
\label{IV.2.s5.5.8-fr}
Soient $A$ un anneau local noethérien intègre, $B$ un anneau intègre contenant
$A$ et qui est une $A$-algèbre de type fini, $\mathfrak q$ un idéal premier
de $B$ tel que $\mathfrak q\cap A$ soit l'idéal maximal $\mathfrak m$ de
$A$, $k$ et $k'$ les corps résiduels de $A$ et $B_{\mathfrak q}$
respectivement. On a alors l'inégalité
\begin{equation}
\dim(A)+\operatorname{deg.tr.}_A B
\geq\dim(B_{\mathfrak q})+\operatorname{deg.tr.}_k k'.
\tag{5.5.8.1}\label{IV.2.s5.5.8.1-fr}
\end{equation}
En outre les deux membres sont égaux si, pour toute sous-$A$-algèbre $A'$ de
type fini de $B[T]$ et tout idéal maximal $\mathfrak m'$ de $A'$ tel que
$\mathfrak m'\cap A=\mathfrak m$, $A'_{\mathfrak m'}$ est caténaire.
\end{theorem}

\begin{proof}
Soit $B=A[x_1,\ldots,x_n]$ et raisonnons par récurrence sur $n$. Posons
$C=A[x_1,\ldots,x_{n-1}]$ et $\mathfrak r=\mathfrak q\cap C$;
$C_{\mathfrak r}$ est un anneau local intègre noethérien, et si on pose
$S=C-\mathfrak r$, $B'=S^{-1}B$, $\mathfrak q'=S^{-1}\mathfrak q$, on a
$B_{\mathfrak q}=B'_{\mathfrak q'}$, $B'=C_{\mathfrak r}[x_n]$ et
$\mathfrak q'\cap C_{\mathfrak r}=\mathfrak rC_{\mathfrak r}$; en outre les
corps des fractions de $B'$ et $C_{\mathfrak r}$ sont ceux de $B$ et $C$
respectivement. Si $k_1$ est le corps résiduel de $C_{\mathfrak r}$, on a
donc, d'après (5.5.7.1)

\oldpage[IV-2]{97}

\begin{equation}
\dim(C_{\mathfrak r})+\operatorname{deg.tr.}_C B
\geq\dim(B_{\mathfrak q})+\operatorname{deg.tr.}_{k_1}k'.
\tag{5.5.8.2}\label{IV.2.s5.5.8.2-fr}
\end{equation}
D'autre part l'hypothèse de récurrence donne
\begin{equation}
\dim(A)+\operatorname{deg.tr.}_A C
\geq\dim(C_{\mathfrak r})+\operatorname{deg.tr.}_k k_1
\tag{5.5.8.3}\label{IV.2.s5.5.8.3-fr}
\end{equation}
d'où (5.5.8.1) en ajoutant (5.5.8.2) et (5.5.8.3) membre à membre. Pour
démontrer la seconde assertion, notons d'abord que toute sous-$A$-algèbre de
type fini de $C[T]$ est aussi une sous-$A$-algèbre de type fini de $B[T]$;
par récurrence sur $n$, on peut donc supposer que les deux membres de
(5.5.8.3) sont égaux. D'autre part, pour voir que les deux membres de
(5.5.8.2) sont égaux, il suffit, en vertu de (5.5.7), de vérifier que si
$\mathfrak n$ est un idéal maximal de $C_{\mathfrak r}[T]$ tel que
$\mathfrak n\cap C_{\mathfrak r}=\mathfrak rC_{\mathfrak r}$, alors
$(C_{\mathfrak r}[T])_{\mathfrak n}$ est caténaire; mais on a
$C_{\mathfrak r}[T]=S'^{-1}C[T]$ où $S'=C-\mathfrak r$; l'idéal
$\mathfrak n$ est donc de la forme $S'^{-1}\mathfrak n'$, où $\mathfrak n'$
est un idéal premier de $C[T]$ tel que $\mathfrak n'\cap C=\mathfrak r$,
d'où $(C_{\mathfrak r}[T])_{\mathfrak n}=(C[T])_{\mathfrak n'}$; si
$\mathfrak m'$ est un idéal maximal de $C[T]$ contenant $\mathfrak n'$,
$(C[T])_{\mathfrak n'}$ est donc un anneau local de l'anneau
$(C[T])_{\mathfrak m'}$, et comme par hypothèse ce dernier est caténaire, il
en est de même de $(C[T])_{\mathfrak n'}$ (0, 16.1.4).
\end{proof}

\subsection{Formule des dimensions et anneaux universellement caténaires}
\label{IV.2.s5.6-fr}

\begin{proposition}[5.6.1]
\EGAReaderTarget{ega4.c04.prop.5-6-1}{IV.2.s5.6.1-fr}
Soit $A$ un anneau noethérien. Les propriétés suivantes sont équivalentes:
\begin{enumerate}
\item[a)] Tout anneau de polynômes $A[T_1,\ldots,T_n]$ est caténaire.
\item[b)] Toute algèbre de type fini sur $A$ est caténaire.
\item[c)] $A$ est caténaire, et pour toute $A$-algèbre intègre locale,
essentiellement de type fini (1.3.8) $B'$, on a, en désignant par
$\mathfrak p$ l'image réciproque dans $A$ de l'idéal maximal $\mathfrak q$
de $B'$, par $A'$ l'image de $A_{\mathfrak p}$ dans $B'$, par $k$ et $k'$
les corps résiduels de $A'$ et $B'$ respectivement,
\begin{equation}
\dim(A')+\operatorname{deg.tr.}_{A'}(B')
=\dim(B')+\operatorname{deg.tr.}_k k'.
\tag{5.6.1.1}\label{IV.2.s5.6.1.1-fr}
\end{equation}
\end{enumerate}
\end{proposition}

Le fait que b) entraîne c) résulte de (5.5.8); en effet, $B'$ est une
$A'$-algèbre locale essentiellement de type fini (1.3.10), donc de la forme
$B''_{\mathfrak q''}$, où $B''$ est une sous-$A'$-algèbre de type fini de
$B'$, $\mathfrak q''$ un idéal premier de $B''$ au-dessus de l'idéal maximal
$\mathfrak p'$ de $A'$; en outre les corps des fractions de $B'$ et $B''$
sont les mêmes, donc
$\operatorname{deg.tr.}_{A'}(B')=\operatorname{deg.tr.}_{A'}(B'')$.
Pour démontrer (5.6.1.1), il suffit de montrer que (sous l'hypothèse b)) toute
sous-$A'$-algèbre $A_1$ de type fini de $B'[T]$ est caténaire; en effet les
deux membres de (5.5.8.1), où on remplace $A$, $B$ et $\mathfrak q$ par
$A'$, $B''$ et $\mathfrak q''$, seront alors égaux, d'où l'égalité (5.6.1.1).
Or l'hypothèse b) entraîne que toute $A$-algèbre essentiellement de type fini
est caténaire (0, 16.1.4), et $A_1$ est une telle $A$-algèbre (1.3.9).

Il est trivial que b) entraîne a); inversement, a) entraîne b), toute
$A$-algèbre de type fini étant quotient d'une algèbre de polynômes
(0, 16.1.4).

Reste à prouver que c) entraîne b). Comme tout quotient par un idéal premier
d'une $A$-algèbre de type fini est une $A$-algèbre intègre de type fini, il
s'agit de voir que, si $B$ est une $A$-algèbre intègre de type fini,
$\mathfrak q$ et $\mathfrak q'$ deux idéaux premiers de $B$ tels que
$\mathfrak q'\subset\mathfrak q$, on a (0, 16.1.4.2)
\begin{equation}
\dim(B_{\mathfrak q}/\mathfrak q'B_{\mathfrak q})
+\dim(B_{\mathfrak q'})=\dim(B_{\mathfrak q}).
\tag{5.6.1.2}\label{IV.2.s5.6.1.2-fr}
\end{equation}

\oldpage[IV-2]{98}

Soient $\mathfrak p$, $\mathfrak p'$ les images réciproques de
$\mathfrak q$, $\mathfrak q'$ respectivement dans $A$, $\mathfrak n$ le
noyau de l'homomorphisme $A_{\mathfrak p}\to B_{\mathfrak q}$.

L'image $A'$ de $A_{\mathfrak p}$ dans $B_{\mathfrak q}$ étant isomorphe à
$A_{\mathfrak p}/\mathfrak n$, la formule (5.6.1.1) appliquée à $A'$ et à
$B'=B_{\mathfrak q}$ s'écrit
\begin{equation}
\dim(A_{\mathfrak p}/\mathfrak n)
+\operatorname{deg.tr.}_{A_{\mathfrak p}/\mathfrak n}(B_{\mathfrak q})
=\dim(B_{\mathfrak q})
+\operatorname{deg.tr.}_{k(\mathfrak p)}k(\mathfrak q).
\tag{5.6.1.3}\label{IV.2.s5.6.1.3-fr}
\end{equation}

D'autre part, le noyau de $A_{\mathfrak p'}\to B_{\mathfrak q'}$ est
$\mathfrak nA_{\mathfrak p'}$, donc, en appliquant la formule (5.6.1.1) où
l'on remplace $A'$ par $A_{\mathfrak p'}/\mathfrak nA_{\mathfrak p'}$ et
$B'$ par $B_{\mathfrak q'}$, on a
\begin{equation}
\dim(A_{\mathfrak p'}/\mathfrak nA_{\mathfrak p'})
+\operatorname{deg.tr.}_{A_{\mathfrak p'}/\mathfrak nA_{\mathfrak p'}}
(B_{\mathfrak q'})
=\dim(B_{\mathfrak q'})
+\operatorname{deg.tr.}_{k(\mathfrak p')}k(\mathfrak q').
\tag{5.6.1.4}\label{IV.2.s5.6.1.4-fr}
\end{equation}

Enfin, $B/\mathfrak q'$ est une $A$-algèbre intègre de type fini, l'image
réciproque de $\mathfrak q/\mathfrak q'$ dans $A$ est $\mathfrak p$, et le
noyau de l'homomorphisme
$A_{\mathfrak p}\to B_{\mathfrak q}/\mathfrak q'B_{\mathfrak q}$ est
$\mathfrak p'A_{\mathfrak p}$, donc, en appliquant (5.6.1.1) où l'on
remplace $A'$ par $A_{\mathfrak p}/\mathfrak p'A_{\mathfrak p}$ et $B'$ par
$B_{\mathfrak q}/\mathfrak q'B_{\mathfrak q}$, on a
\begin{equation}
\dim(A_{\mathfrak p}/\mathfrak p'A_{\mathfrak p})
+\operatorname{deg.tr.}_{A_{\mathfrak p}/\mathfrak p'A_{\mathfrak p}}
(B_{\mathfrak q}/\mathfrak q'B_{\mathfrak q})
=\dim(B_{\mathfrak q}/\mathfrak q'B_{\mathfrak q})
+\operatorname{deg.tr.}_{k(\mathfrak p')}k(\mathfrak q).
\tag{5.6.1.5}\label{IV.2.s5.6.1.5-fr}
\end{equation}

Ajoutons maintenant (5.6.1.4) et (5.6.1.5) membre à membre, et notons que
$k(\mathfrak p')$ (resp. $k(\mathfrak q')$) est le corps des fractions de
$A_{\mathfrak p}/\mathfrak p'A_{\mathfrak p}$ (resp.
$B_{\mathfrak q}/\mathfrak q'B_{\mathfrak q}$); par ailleurs
$A_{\mathfrak p'}/\mathfrak nA_{\mathfrak p'}$ et
$A_{\mathfrak p}/\mathfrak n$ ont même corps des fractions, et
$B_{\mathfrak q'}$ et $B_{\mathfrak q}$ ont même corps des fractions;
enfin, puisque $A$ est supposé caténaire, on a (0, 16.1.4.2)
\[
\dim(A_{\mathfrak p}/\mathfrak p'A_{\mathfrak p})
+\dim(A_{\mathfrak p'}/\mathfrak nA_{\mathfrak p'})
=\dim(A_{\mathfrak p}/\mathfrak n).
\]
Tenant compte de ces remarques et de (5.6.1.3), il vient la relation
(5.6.1.2). C.Q.F.D.

\begin{definition}[5.6.2]
\label{IV.2.s5.6.2-fr}
On dit qu'un anneau noethérien $A$ est \emph{universellement caténaire} s'il
vérifie les conditions équivalentes a), b), c) de (5.6.1).
\end{definition}

\begin{remarks}[5.6.3]
\label{IV.2.s5.6.3-fr}
\begin{enumerate}
\item[(i)] Si $A$ est universellement caténaire, il en est de même de
$S^{-1}A$ pour toute partie multiplicative $S$ de $A$, comme il résulte
aussitôt de (5.6.1, a)) et du fait que tout anneau de fractions d'un anneau
caténaire est caténaire. Inversement, si, pour tout idéal premier
$\mathfrak p$ de $A$, l'anneau $A_{\mathfrak p}$ est universellement
caténaire, alors $A$ est universellement caténaire: cela résulte de
(5.6.1, c)), si l'on remarque qu'en posant $S=A-\mathfrak p$, $S^{-1}B$
est une $A_{\mathfrak p}$-algèbre de type fini, et
$B_{\mathfrak q}=(S^{-1}B)_{S^{-1}\mathfrak q}$.
\item[(ii)] On dit qu'un préschéma localement noethérien $X$ est
\emph{universellement caténaire} si, pour tout $x\in X$, l'anneau
$\mathcal O_x$ est universellement caténaire. Il résulte de (i) que pour que
$A$ soit un anneau universellement caténaire, il faut et il suffit que le
schéma $\operatorname{Spec}(A)$ le soit.
\item[(iii)] Le critère (5.6.1, c)) ne fait intervenir que les images de $A$
dans des $A$-algèbres intègres de type fini, donc des anneaux quotients
intègres de $A$. On en conclut que si

\oldpage[IV-2]{99}

$\mathfrak p_i$ ($1\leq i\leq n$) sont les idéaux premiers minimaux de $A$,
il revient au même de dire que $A$ est universellement caténaire ou que
chacun des $A/\mathfrak p_i$ l'est. Cela montre aussi qu'il revient au même
de dire qu'un préschéma localement noethérien $X$ est universellement
caténaire, ou que $X_{\mathrm{red}}$ l'est.
\item[(iv)] Le critère (5.6.1, b)) et la remarque (i) montrent que, si $A$
est un anneau universellement caténaire, il en est de même de toute
$A$-algèbre essentiellement de type fini (1.3.8).
\end{enumerate}
\end{remarks}

\begin{proposition}[5.6.4]
\label{IV.2.s5.6.4-fr}
Tout anneau quotient d'un anneau régulier est universellement caténaire.
\end{proposition}

Tout revient à voir qu'un anneau régulier $A$ est universellement caténaire
(5.6.3, (iv)); or, on sait que $A[T_1,\ldots,T_n]$ est régulier (0, 17.3.7),
donc caténaire (0, 16.5.12), et l'on conclut par (5.6.1, a)).

En particulier, il résulte du théorème de Cohen (0, 19.8.8) que tout anneau
local noethérien complet est universellement caténaire. De même, toute
algèbre de type fini sur un corps étant quotient d'un anneau régulier, tout
préschéma localement de type fini sur un corps est universellement caténaire.

Nous verrons plus loin (6.3.7) que dans (5.6.4), on peut remplacer « anneau
régulier » par « anneau de Cohen-Macaulay ».

\begin{proposition}[5.6.5]
\EGAReaderTarget{ega4.c04.prop.5-6-5}{IV.2.s5.6.5-fr}
Soient $Y$ un préschéma irréductible localement noethérien, $X$ un préschéma
irréductible, $f:X\to Y$ un morphisme dominant localement de type fini. Soit
$\xi$ (resp. $\eta$) le point générique de $X$ (resp. $Y$), et soit
$e=\dim_\eta(f^{-1}(\eta))=\operatorname{deg.tr.}_{k(\eta)}k(\xi))$
(« dimension de la fibre générique », cf. $(0_{\mathrm I},2.1.8)$ et
(4.1.1)). Pour tout $x\in X$, on a alors, en posant $y=f(x)$,
\begin{equation}
e+\dim(\mathcal O_y)
\geq\operatorname{deg.tr.}_{k(y)}k(x)+\dim(\mathcal O_x)
\tag{5.6.5.1}\label{IV.2.s5.6.5.1-fr}
\end{equation}
\begin{equation}
\dim(\mathcal O_x)
\leq\dim(\mathcal O_y)
+\dim(\mathcal O_x\otimes_{\mathcal O_y}k(y))-\delta(x)
\tag{5.6.5.2}\label{IV.2.s5.6.5.2-fr}
\end{equation}
en posant $\delta(x)=\dim_x(f^{-1}(y))-e$.

En outre, si $Y$ est universellement caténaire, les deux membres de
(5.6.5.1) (resp. (5.6.5.2)) sont égaux. Si de plus $x$ est fermé dans
$f^{-1}(y)$, on a
\begin{equation}
\dim(\mathcal O_x)=\dim(\mathcal O_y)+e.
\tag{5.6.5.3}\label{IV.2.s5.6.5.3-fr}
\end{equation}
\end{proposition}

On peut évidemment se borner au cas où $X$ et $Y$ sont affines (donc $f$ de
type fini) et réduits (1.5.4), donc intègres. Comme $f$ est dominant,
$\mathcal O_y$ s'identifie alors à un sous-anneau de $\mathcal O_x$, et les
corps des fractions respectifs de $\mathcal O_y$ et de $\mathcal O_x$ à
$k(\eta)$ et $k(\xi)$ (I, 8.2.7); en outre, $\mathcal O_x$ est un anneau
local d'une $\mathcal O_y$-algèbre intègre de type fini $B$ en un idéal
premier $\mathfrak q$ (I, 3.6.5); en appliquant (5.5.8.1) à
$A=\mathcal O_y$, $B$ et $\mathfrak q$, on obtient (5.6.5.1). En outre,
$\mathcal O_x\otimes_{\mathcal O_y}k(y)=\mathcal O_x/\mathfrak m_y\mathcal O_x$
n'est autre que l'anneau local de la fibre $f^{-1}(y)$ au point $x$
(I, 3.6.1); comme $f^{-1}(y)$ est un préschéma de type fini sur $k(y)$, il
résulte de (5.2.3) que l'on a
\[
\dim(\mathcal O_x\otimes_{\mathcal O_y}k(y))
=\dim_x(f^{-1}(y))-\operatorname{deg.tr.}_{k(y)}k(x);
\]
l'inégalité (5.6.5.2) n'est donc qu'une autre forme de (5.6.5.1).

Le fait que les deux membres de (5.6.5.1) sont égaux lorsque $Y$ est
universellement caténaire n'est autre que l'égalité (5.6.1.1), appliquée à
$A=\mathcal O_y$ et $B'=\mathcal O_x$. Pour démontrer que l'on a de plus
(5.6.5.3) lorsque $x$ est fermé dans $f^{-1}(y)$, il suffit de remarquer que,
de façon générale, $\operatorname{deg.tr.}_{k(y)}k(x)$ est la dimension de
$\overline{\{x\}}\cap f^{-1}(y)$, comme il résulte de (5.2.1) appliqué au
sous-préschéma fermé réduit de $f^{-1}(y)$ ayant ce sous-espace comme espace
sous-jacent.

\oldpage[IV-2]{100}

Nous démontrerons plus loin (13.1.1) que, sous les conditions de (5.6.5), on
a toujours $\delta(x)\geq0$, et (5.6.5.2) précise donc dans ce cas (5.5.2.1).

\begin{corollary}[5.6.6]
\label{IV.2.s5.6.6-fr}
Sous les hypothèses générales de (5.6.5), on a
\begin{equation}
\dim(X)\leq\dim(Y)+e.
\tag{5.6.6.1}\label{IV.2.s5.6.6.1-fr}
\end{equation}
Si l'on suppose en outre $Y$ universellement caténaire, alors, pour que les
deux membres de (5.6.6.1) soient égaux, il faut et il suffit que l'on ait
\begin{equation}
\dim(Y)=\sup_{y\in f(X)}\dim(\mathcal O_y).
\tag{5.6.6.2}\label{IV.2.s5.6.6.2-fr}
\end{equation}
En particulier, les deux membres de (5.6.6.1) sont égaux lorsque $Y$ est
localement de type fini sur un corps.
\end{corollary}

L'inégalité (5.6.6.1) résulte en effet de (5.6.5.1) appliqué en un point
fermé $x$ de $X$, compte tenu de (5.1.4.2) et de (5.1.11). D'autre part,
toute fibre non vide $f^{-1}(y)$ contient un point $x$ fermé dans cette
fibre (5.1.11), et si l'on suppose $Y$ universellement caténaire, on a donc
en ce point la relation (5.6.5.3). En prenant les bornes supérieures des deux
membres de (5.6.5.3) lorsque $x$ parcourt $X$, et en notant que tout point
$x$ fermé dans $X$ est \emph{a fortiori} fermé dans $f^{-1}(f(x))$, la
seconde assertion du corollaire résulte de (5.1.4.1) et (5.1.4.2).

Pour prouver la dernière assertion, notons qu'il y a un voisinage ouvert
affine $U$ du point générique de $X$ tel que $f|U$ soit de type fini, donc
$f(U)$ est une partie constructible de $Y$, dense dans $Y$ (1.8.5), et par
suite contient une partie ouverte non vide (donc partout dense) de $Y$
$(0_{\mathrm{III}},9.2.2)$. L'hypothèse que $Y$ est localement de type fini
sur un corps assure alors, d'une part que $Y$ est universellement caténaire
(5.6.4), et d'autre part que les deux membres de (5.6.6.2) sont égaux
(5.2.2 et 4.1.1.3).

On notera que la condition (5.6.6.2) est trivialement vérifiée lorsque $f$
est surjectif.

\begin{corollary}[5.6.7]
\label{IV.2.s5.6.7-fr}
Soient $Y$ un préschéma localement noethérien, $f:X\to Y$ un morphisme
localement de type fini, $n$ un entier. Si, pour tout $y\in Y$,
$\dim(f^{-1}(y))\leq n$, alors on a
\begin{equation}
\dim(X)\leq\dim(Y)+n.
\tag{5.6.7.1}\label{IV.2.s5.6.7.1-fr}
\end{equation}
\end{corollary}

En vertu de (0, 14.1.7) et de (1.5.4), on peut se borner au cas où $X$ et
$Y$ sont affines, donc $f$ de type fini, $X$ et $Y$ réduits; soient $X_i$
($1\leq i\leq m$) les sous-préschémas fermés (intègres) de $X$ ayant pour
espaces sous-jacents les composantes irréductibles de $X$; on a
$\dim(X)=\sup_i(\dim(X_i))$. Si $Z_i$ est le sous-préschéma fermé réduit de
$Y$ ayant pour espace sous-jacent $\overline{f(X_i)}$, $Z_i$ est intègre
$(0_{\mathrm I},2.1.5)$; la restriction $X_i\to Y$ de $f$ se factorise en
$X_i\xrightarrow{g_i}Z_i\xrightarrow{j_i}Y$, où $j_i$ est l'injection
canonique (I, 5.2.2), et $g_i$ est dominant et de type fini (1.5.4).

\oldpage[IV-2]{101}

On a $\dim(Z_i)\leq\dim(Y)$ pour tout $i$; d'autre part, si $\zeta_i$ est le
point générique de $Z_i$, on a $\dim(g_i^{-1}(\zeta_i))\leq n$ pour tout $i$
par hypothèse; on voit ainsi que pour prouver (5.6.7.1), il suffit de
démontrer que $\dim(X_i)\leq\dim(Z_i)+n$ pour tout $i$, ce qui résulte de
(5.6.6.1).

\begin{corollary}[5.6.8]
\label{IV.2.s5.6.8-fr}
Soient $Y$ un préschéma localement noethérien, $X$ un préschéma
irréductible, $f:X\to Y$ un morphisme localement de type fini, $n$ un entier
$\geq0$. On suppose que $Y$ est universellement caténaire, et que pour tout
$y\in Y$, on ait $\dim(f^{-1}(y))\geq n$ (resp.
$\dim(f^{-1}(y))=n$). Alors on a
\begin{equation}
\dim(X)\geq\dim(Y)+n
\tag{5.6.8.1}\label{IV.2.s5.6.8.1-fr}
\end{equation}
(resp.
\begin{equation}
\dim(X)=\dim(Y)+n).
\tag{5.6.8.2}\label{IV.2.s5.6.8.2-fr}
\end{equation}
\end{corollary}

Comme $n\geq0$, l'hypothèse entraîne que $f$ est surjectif, donc $Y$
irréductible; en outre, si $\eta$ est le point générique de $Y$, on a
$\dim(f^{-1}(\eta))\geq n$ (resp. $\dim(f^{-1}(\eta))=n$). La conclusion
résulte donc de (5.6.6).

\begin{remarks}[5.6.9]
\label{IV.2.s5.6.9-fr}
\begin{enumerate}
\item[(i)] Même si $Y$ est régulier, $X$ irréductible, $f:X\to Y$ dominant
et de type fini, les deux membres de (5.6.6.1) ne sont pas nécessairement
égaux, comme le montre l'exemple où $Y=\operatorname{Spec}(A)$ où $A$ est un
anneau de valuation discrète, $X=\operatorname{Spec}(K)$ où $K$ est le corps
des fractions de $A$, $f$ étant le morphisme canonique.
\item[(ii)] L'exemple (5.4.3, (i)) montre que dans (5.6.6) et (5.6.8), on ne
peut supprimer l'hypothèse que $X$ est irréductible, les autres hypothèses
étant vérifiées. Nous verrons toutefois (10.6.1) que moyennant des hypothèses
supplémentaires sur $Y$, vérifiées par exemple lorsque $Y$ est un préschéma
localement de type fini sur un corps, ou sur un anneau de Dedekind ayant une
infinité d'idéaux premiers, de tels phénomènes ne peuvent se présenter.
\end{enumerate}
\end{remarks}

\begin{proposition}[5.6.10]
\label{IV.2.s5.6.10-fr}
Soient $A$ un anneau local noethérien intègre universellement caténaire, $B$
un anneau intègre contenant $A$ et qui est une $A$-algèbre finie. Alors, pour
tout idéal maximal $\mathfrak n$ de $B$, on a
$\dim(B_{\mathfrak n})=\dim(A)$.
\end{proposition}

En effet, on a $\operatorname{deg.tr.}_A B=0$ et le corps résiduel $k'$ de
$B_{\mathfrak n}$ est une extension algébrique du corps des fractions de
$A$. La conclusion résulte de la formule (5.6.1.1), tout idéal maximal de
$B$ étant au-dessus de celui de $A$.

\begin{example}[5.6.11]
\label{IV.2.s5.6.11-fr}
Soit $A$ un anneau local noethérien intègre et intégralement clos; alors, on
a vu (0, 16.1.6) que la conclusion de (5.6.10) est valable pour toute
$A$-algèbre finie intègre $B$ contenant $A$. Par contre, nous allons
construire un exemple d'anneau local noethérien intègre caténaire $A$ pour
lequel la conclusion de (5.6.10) sera en défaut. Nous utiliserons la
construction de (5.2.5, (i)). Soient $k_0$ un corps, $k$ une extension
transcendante pure $k_0(S_i)_{i\in\mathbb N}$ de degré de transcendance
infini, $V$ l'anneau de valuation discrète $k[S]_{(S)}$, anneau local de
l'anneau de polynômes $k[S]$ en l'idéal premier $(S)$; enfin soit $E$
l'anneau de polynômes $V[T]$; on a vu que dans $E$ l'idéal maximal
$\mathfrak m=(S)+(T)$ est de hauteur 2 et l'idéal maximal
$\mathfrak m'=(ST-1)$ de hauteur 1;

\oldpage[IV-2]{102}

les corps résiduels correspondants sont $k(\mathfrak m)=k$ et
$k(\mathfrak m')=k(S)$, corps des fractions de $V$; en vertu du choix de
$k$, ces corps sont isomorphes. Désignons par $\varepsilon$ et
$\varepsilon'$ les homomorphismes canoniques de $E$ sur
$E/\mathfrak m=k(\mathfrak m)$ et $E/\mathfrak m'=k(\mathfrak m')$
respectivement; soit d'autre part $\sigma$ un isomorphisme de
$k(\mathfrak m)$ sur $k(\mathfrak m')$, et considérons le sous-anneau $C$
de $E$ formé des $x\in E$ tels que
$\varepsilon'(x)=\sigma(\varepsilon(x))$ (cette construction est un cas
particulier des « procédés de recollement » qui seront étudiés
systématiquement au chap. V). Il est immédiat que
$\mathfrak n=\mathfrak m\mathfrak m'=\mathfrak m\cap\mathfrak m'$ est un
idéal maximal de $C$, $C/\mathfrak m\mathfrak m'$ s'identifiant au
sous-corps de $(E/\mathfrak m)\times(E/\mathfrak m')$ formé des couples
$(z,\sigma(z))$. On a $E=C+C(ST-1)$, autrement dit $E$ est une
$C$-algèbre finie et est évidemment la clôture intégrale de $C$; nous allons
voir en outre que $C$ est noethérien: cela résultera du lemme suivant:

\begin{lemma}[5.6.11.1]
\EGAReaderTarget{ega4.c04.lem.5-6-11-1}{IV.2.s5.6.11.1-fr}
Soient $R$ un anneau, $S$ un sous-anneau,
$\mathfrak K=\operatorname{Ann}_S(R/S)$ le conducteur de $S$ dans $R$
(plus grand idéal de $S$ qui est aussi un idéal de $R$).
\begin{enumerate}
\item[(i)] Pour tout idéal $\mathfrak J\subset\mathfrak K$ de $R$, il existe
une application strictement croissante de l'ensemble des idéaux
$\mathfrak a$ de $S$ tels que $R\mathfrak a=\mathfrak J$ sur l'ensemble des
sous-$(S/\mathfrak K)$-modules de
$\mathfrak J/\mathfrak K\mathfrak J$.
\item[(ii)] Si $S/\mathfrak K$ et $R$ sont noethériens et si $R$ est un
$S$-module de type fini, alors toute suite croissante d'idéaux de $S$
contenus dans $K$ est stationnaire.
\end{enumerate}
\end{lemma}

(i) On a en effet,
$\mathfrak K\mathfrak J=\mathfrak K(R\mathfrak a)
=\mathfrak K\mathfrak a\subset\mathfrak a$
($\mathfrak a$ étant un idéal de $S$), d'où la conclusion.

(ii) Soit $(\mathfrak L_n)$ une suite croissante d'idéaux de $S$ contenus
dans $\mathfrak K$. La suite des idéaux $R\mathfrak L_n$ de $R$ est
stationnaire puisque $R$ est noethérien; on peut donc supposer que tous les
idéaux $R\mathfrak L_n$ sont égaux à un même idéal $\mathfrak J$ de $R$.
Comme $R$ est noethérien, $\mathfrak J/\mathfrak K\mathfrak J$ est un
$R$-module de type fini, donc aussi un $S$-module de type fini, et par suite
un $(S/\mathfrak K)$-module de type fini. Mais comme $S/\mathfrak K$ est
noethérien, $\mathfrak J/\mathfrak K\mathfrak J$ est un
$(S/\mathfrak K)$-module noethérien, et la conclusion résulte de (i).

Nous allons appliquer ce lemme en prenant $R=E$, $S=C$; il est clair que
$\mathfrak n$ est le conducteur de $C$ dans $E$; en outre, pour tout idéal
$\mathfrak a$ de $C$,
$\mathfrak a/(\mathfrak n\cap\mathfrak a)$ est isomorphe à
$(\mathfrak a+\mathfrak n)/\mathfrak n$, donc un sous-$C$-module de
$C/\mathfrak n$, qui est un $C$-module simple; il suffit donc de démontrer
que tout idéal $\mathfrak a\cap\mathfrak n$ de $C$ est de type fini, et cela
résulte de (5.6.11.1).

Il résulte de (0, 16.1.5) que l'on a $\dim(C)=\dim(E)=2$. Prenons
$A=C_{\mathfrak n}$; si $U=C-\mathfrak n$, l'anneau de fractions
$B=U^{-1}E$ est donc la clôture intégrale de $A$, et c'est un $A$-module de
type fini; d'ailleurs, comme $\mathfrak m$ et $\mathfrak m'$ sont les seuls
idéaux premiers de $E$ contenant $\mathfrak n$, $B$ est un anneau semi-local
dont les anneaux locaux en les deux idéaux maximaux sont isomorphes à
$E_{\mathfrak m}$ et $E_{\mathfrak m'}$ respectivement, donc sont
respectivement de dimension 2 et 1. Cela montre que $\dim(B)=2$, donc aussi
$\dim(A)=2$ (0, 16.1.5). Comme $A$ est un anneau local intègre, il est
nécessairement caténaire (tout idéal premier distinct de 0 et de l'idéal
maximal étant nécessairement de hauteur 1); mais il ne vérifie pas la
conclusion de (5.6.10), et \emph{a fortiori} n'est pas universellement
caténaire.

Notons encore que $\mathfrak n$ est un idéal de hauteur 2 dans $C$, et que
pour tout idéal premier $\mathfrak p$ de hauteur 1 dans $C$, il existe un
seul idéal premier $\mathfrak p'$ de $E$ au-dessus de $\mathfrak p$,
nécessairement de hauteur 1, et tel que
$C_{\mathfrak p}=E_{\mathfrak p'}$. En effet, il y a au moins un idéal
premier $\mathfrak p'$ de $E$ au-dessus de $\mathfrak p$ et il résulte du
théorème de Cohen-Seidenberg qu'un

\oldpage[IV-2]{103}

tel idéal est nécessairement de hauteur 1; en outre, comme $E$ est une
$C$-algèbre finie et que $\mathfrak p\not\supset\mathfrak n$,
$C_{\mathfrak p}$ est intégralement clos (Bourbaki, \emph{Alg. comm.},
chap. V, § 1, n° 5, cor. 5 de la prop. 16), donc
$E_{\mathfrak p'}=C_{\mathfrak p}$, ce qui prouve nos assertions.

Il serait intéressant de savoir si tout anneau local intègre noethérien
vérifiant la conclusion de (5.6.10) est universellement caténaire; c'est ce
qu'a affirmé Nagata [33], mais sa démonstration ne paraît pas complète.
\end{example}

\subsection{\texorpdfstring{Profondeur et propriété $(S_k)$}{Profondeur et propriété (S\string_k)}}

\begin{definition}[5.7.1]
\label{IV.2.s5.7.1-fr}
Soient $X$ un préschéma localement noethérien, $\mathscr F$ un
$\mathcal O_X$-Module cohérent. On appelle \emph{profondeur} (resp.
\emph{coprofondeur}) de $\mathscr F$ en un point $x\in X$ le nombre
$\operatorname{prof}(\mathscr F_x)$ (resp.
$\operatorname{coprof}(\mathscr F_x)$) (0, 16.4.5 et 16.4.9). On appelle
\emph{coprofondeur de $\mathscr F$} le nombre
\begin{equation}
\operatorname{coprof}(\mathscr F)
=\sup_{x\in X}\operatorname{coprof}(\mathscr F_x).
\tag{5.7.1.1}\label{IV.2.s5.7.1.1-fr}
\end{equation}

On dit que $\mathscr F$ est un $\mathcal O_X$-Module de Cohen-Macaulay au
point $x$ si $\mathscr F_x$ est un $\mathcal O_x$-module de Cohen-Macaulay,
c'est-à-dire (0, 16.5.1) si $\operatorname{coprof}(\mathscr F_x)=0$. On dit
que $\mathscr F$ est un $\mathcal O_X$-Module de Cohen-Macaulay s'il l'est
en tout point, autrement dit si $\operatorname{coprof}(\mathscr F)=0$. Un
point $x\in X$ tel que $\mathcal O_x$ soit un anneau de Cohen-Macaulay est
encore appelé \emph{point de Cohen-Macaulay} de $X$.

On appelle \emph{coprofondeur de $X$} et on note $\operatorname{coprof}(X)$
le nombre $\operatorname{coprof}(\mathcal O_X)$. On dit que $X$ est un
\emph{préschéma de Cohen-Macaulay} si $\mathcal O_X$ est un
$\mathcal O_X$-Module de Cohen-Macaulay, autrement dit si
$\operatorname{coprof}(X)=0$. Tout préschéma localement noethérien de
dimension 0 est évidemment un préschéma de Cohen-Macaulay. Dire que
$\operatorname{Spec}(A)$ est un schéma de Cohen-Macaulay signifie que $A$
est un anneau de Cohen-Macaulay (0, 16.5.13).
\end{definition}

\begin{definition}[5.7.2]
\label{IV.2.s5.7.2-fr}
Soient $X$ un préschéma localement noethérien, $\mathscr F$ un
$\mathcal O_X$-Module cohérent, $k$ un entier positif ou négatif. On dit que
$\mathscr F$ possède la propriété $(S_k)$ si, pour tout $x\in X$, on a
\begin{equation}
\operatorname{prof}(\mathscr F_x)\geq
\inf(k,\dim(\mathscr F_x)).
\tag{5.7.2.1}\label{IV.2.s5.7.2.1-fr}
\end{equation}
On dit que $\mathscr F$ possède la propriété $(S_k)$ en un point $x\in X$
si, pour toute générisation $x'$ de $x$ dans $X$, on a
\begin{equation}
\operatorname{prof}(\mathscr F_{x'})\geq
\inf(k,\dim(\mathscr F_{x'})).
\tag{5.7.2.2}\label{IV.2.s5.7.2.2-fr}
\end{equation}

On dit que $X$ vérifie la propriété $(S_k)$ (resp. vérifie la propriété
$(S_k)$ en un point $x$) si $\mathcal O_X$ vérifie la propriété $(S_k)$
(resp. vérifie la propriété $(S_k)$ au point $x$).

Pour que $\mathscr F$ vérifie la propriété $(S_k)$, il faut et il suffit
évidemment qu'il la vérifie en tout point de $X$. Si $U$ est un ouvert de
$X$ et si $\mathscr F$ vérifie $(S_k)$, il en est de même de
$\mathscr F|U$; réciproquement, si $(U_\alpha)$ est un recouvrement ouvert
de $X$ et si $\mathscr F|U_\alpha$ vérifie $(S_k)$ pour tout $\alpha$,
$\mathscr F$ vérifie $(S_k)$.
\end{definition}

\begin{remarks}[5.7.3]
\label{IV.2.s5.7.3-fr}
\begin{enumerate}
\item[(i)] Rappelons que l'on a toujours
$\operatorname{prof}(\mathscr F_x)\leq\dim(\mathscr F_x)$ (0, 16.4.5.1)
si $\mathscr F_x\neq0$. Dire que $\mathscr F$ possède la propriété $(S_k)$
signifie donc que l'on a $\operatorname{prof}(\mathscr F_x)\geq k$ sauf aux
points $x\in X$ tels que $\dim(\mathscr F_x)<k$ et qu'en ces derniers points
on a $\dim(\mathscr F_x)=\operatorname{prof}(\mathscr F_x)$, c'est-à-dire
(0, 16.5.1) que $\mathscr F$ est un $\mathcal O_X$-Module de
Cohen-Macaulay en ces points; on notera qu'aux points où
$\dim(\mathscr F_x)=k$, on a $\operatorname{prof}(\mathscr F_x)=k$,

\oldpage[IV-2]{104}

donc $\mathscr F$ est encore un $\mathcal O_X$-Module de Cohen-Macaulay en
ces points. Dire que $\mathscr F$ possède la propriété $(S_k)$ pour tout
$k$ signifie donc que $\mathscr F$ est un $\mathcal O_X$-Module de
Cohen-Macaulay. Il est clair que pour $k'\geq k$, la propriété $(S_{k'})$
implique $(S_k)$; pour $k\leq0$, tout $\mathcal O_X$-Module cohérent a la
propriété $(S_k)$.

\item[(ii)] Pour vérifier la condition (5.7.2.1), on peut se borner au cas
où $x\in\operatorname{Supp}(\mathscr F)$; dans le cas contraire on a en
effet $\dim(\mathscr F_x)=-\infty$ (0, 14.1.2).

\item[(iii)] Si $X=\operatorname{Spec}(A)$, où $A$ est un anneau
noethérien, et $\mathscr F=\widetilde M$, où $M$ est un $A$-module de type
fini, on dit que $M$ possède la propriété $(S_k)$ si $\mathscr F$ possède
cette propriété. Pour un préschéma localement noethérien quelconque $X$,
dire que $\mathscr F$ possède la propriété $(S_k)$ en un point $x\in X$
signifie donc que si on pose $Y=\operatorname{Spec}(\mathcal O_x)$, le
$\mathcal O_Y$-Module $\widetilde{\mathscr F_x}$ possède la propriété
$(S_k)$; on notera que la condition (5.7.2.1) en général n'entraîne pas
(5.7.2.2) pour toute générisation $x'$ de $x$, vu qu'on n'a aucune relation
d'inégalité entre $\operatorname{prof}(\mathscr F_x)$ et
$\operatorname{prof}(\mathscr F_{x'})$ (0, 16.4.6).

\item[(iv)] Il résulte aussitôt de la définition que si $\mathscr F$ vérifie
$(S_k)$ en un point $x$, il vérifie aussi $(S_k)$ en tout point $x'$
générisation de $x$.

\item[(v)] La propriété $(S_k)$ est surtout importante pour $k=1$ et
$k=2$; elle a été introduite pour $k=2$ par Serre, pour exprimer son critère
de normalité (cf. (5.8.5)).

\item[(vi)] Soient $X$ un préschéma localement noethérien, $Y$ un
sous-préschéma fermé de $X$, $j:Y\to X$ l'injection canonique, $\mathscr G$
un $\mathcal O_Y$-Module cohérent. Il est clair que pour tout $x\in Y$, on a
$\dim(\mathscr G_x)=\dim((j_*(\mathscr G))_x)$ et
$\operatorname{prof}(\mathscr G_x)=\operatorname{prof}((j_*(\mathscr G))_x)$,
d'où
$\operatorname{coprof}(\mathscr G_x)=\operatorname{coprof}((j_*(\mathscr G))_x)$.
Pour que $\mathscr G$ vérifie $(S_k)$, il faut et il suffit que
$j_*(\mathscr G)$ vérifie $(S_k)$.
\end{enumerate}
\end{remarks}

\begin{proposition}[5.7.4]
\label{IV.2.s5.7.4-fr}
Soient $X$ un préschéma localement noethérien, $\mathscr F$ un
$\mathcal O_X$-Module cohérent. Posons $S=\operatorname{Supp}(\mathscr F)$,
et pour tout $n\geq0$, désignons par $Z_n$ l'ensemble des $z\in X$ tels que
$\operatorname{coprof}(\mathscr F_x)>n$, de sorte que $Z_n\subset S$.
\begin{enumerate}
\item[(i)] Pour que $\mathscr F$ vérifie la propriété $(S_k)$, il faut et il
suffit que, pour tout $n\geq0$, on ait
\begin{equation}
\operatorname{codim}(Z_n,S)>n+k.
\tag{5.7.4.1}\label{IV.2.s5.7.4.1-fr}
\end{equation}
\item[(ii)] Supposons en outre que les $Z_n$ soient fermés dans $X$. Alors,
pour que $\mathscr F$ vérifie la propriété $(S_k)$ en un point $x\in S$, il
faut et il suffit que, pour tout $n\geq0$, on ait
\begin{equation}
\operatorname{codim}_x(Z_n,S)>n+k.
\tag{5.7.4.2}\label{IV.2.s5.7.4.2-fr}
\end{equation}
\end{enumerate}
\end{proposition}

(i) On a en effet, par définition (5.1.3),
$\operatorname{codim}(Z_n,S)=\inf_{z\in Z_n}(\dim(\mathcal O_{S,z}))$ et
l'inégalité (5.7.4.1) signifie donc (5.1.12.2) que, pour tout $z\in X$, et
tout $n\geq0$, la relation
\[
\dim(\mathscr F_z)-\operatorname{prof}(\mathscr F_z)\geq n+1
\]
entraîne la relation
\[
\dim(\mathscr F_z)\geq n+k+1.
\]
Mais si l'on pose $a=\dim(\mathscr F_z)$,
$b=\operatorname{prof}(\mathscr F_z)$, on a $b\leq a$, et dire que pour tout
$n\geq0$, la relation $b\leq a-n-1$ implique $k\leq a-n-1$ équivaut à dire
que $b\geq\inf(k,a)$, d'où la proposition.

\oldpage[IV-2]{105}

(ii) Le raisonnement est le même que dans (i), à cela près qu'il faut se
limiter aux $z\in X$ qui sont des générisations de $x$, et tenir compte de
(5.1.3.2).

\begin{proposition}[5.7.5]
\label{IV.2.s5.7.5-fr}
Soient $X$ un préschéma localement noethérien, $\mathscr F$ un
$\mathcal O_X$-Module cohérent. Pour que $\mathscr F$ vérifie $(S_k)$
($k\geq1$), il faut et il suffit que pour tout entier $r$ tel que
$0\leq r<k$, tout ouvert $U$ de $X$ et toute suite
$(\mathscr F|U)$-régulière $(f_i)_{1\leq i\leq r}$ de sections de
$\mathcal O_X$ au-dessus de $U$,
\[
(\mathscr F|U)\big/\Big(\sum_{i=1}^r f_i(\mathscr F|U)\Big)
\]
soit sans cycle premier associé immergé.
\end{proposition}

Démontrons d'abord la proposition pour $k=1$; elle s'énonce alors encore en
disant que pour que $\mathscr F$ vérifie $(S_1)$, il faut et il suffit que
$\mathscr F$ soit sans cycle premier associé immergé. En effet, dire que
$\mathscr F$ vérifie $(S_1)$ signifie qu'aux points
$x\in\operatorname{Supp}(\mathscr F)$ tels que $\dim(\mathscr F_x)>0$, on a
$\operatorname{prof}(\mathscr F_x)\geq1$ (puisque aux autres points de
$\operatorname{Supp}(\mathscr F)$ on a $\dim(\mathscr F_x)=0$, donc
$\operatorname{prof}(\mathscr F_x)=0$); mais dire que
$\operatorname{prof}(\mathscr F_x)\geq1$ signifie que $\mathfrak m_x$
n'est pas associé à $\mathscr F_x$ (0, 16.4.6, (i)), ou encore que $x$
n'est pas associé à $\mathscr F$ (3.1.1); d'autre part, si $S$ est un
sous-préschéma de $X$ ayant pour espace sous-jacent
$\operatorname{Supp}(\mathscr F)$, on a (5.1.12.1)
$\dim(\mathscr F_x)=\dim(\mathcal O_{S,x})$; dire que
$\dim(\mathscr F_x)>0$ signifie donc que $x$ n'est pas point maximal de
$\operatorname{Supp}(\mathscr F)$ (5.1.2), d'où la conclusion.

En second lieu, montrons que, pour $k$ quelconque $>1$, la condition de
l'énoncé est nécessaire. On peut se borner à considérer le cas où $U=X$,
et notre assertion sera prouvée (par récurrence sur $k$ et en vertu de la
première partie du raisonnement) si nous montrons que lorsque $\mathscr F$
vérifie $(S_k)$ ($k>1$) et que $f$ est une section $\mathscr F$-régulière
de $\mathcal O_X$ au-dessus de $X$, alors $\mathscr F/f\mathscr F$ vérifie
$(S_{k-1})$ (i.e. vérifie (5.7.2.1) où $k$ est remplacé par $k-1$). Or,
pour tout $x\in X$, $f_x$ est un élément $\mathscr F_x$-régulier; s'il est
inversible, on a $\mathscr F_x/f_x\mathscr F_x=0$ et la conclusion est
triviale. Si au contraire $f_x\in\mathfrak m_x$, on sait que l'on a
$\operatorname{prof}(\mathscr F_x/f_x\mathscr F_x)
=\operatorname{prof}(\mathscr F_x)-1$ (0, 16.4.6, (i)), et
$\dim(\mathscr F_x/f_x\mathscr F_x)=\dim(\mathscr F_x)-1$ (0, 16.3.4), et
notre assertion en résulte.

Prouvons enfin que pour $k>1$, la condition de l'énoncé est suffisante. Nous
allons procéder par récurrence sur $k$; soit $x$ un point de $X$, et
supposons d'abord que $\dim(\mathscr F_x)\geq k$. L'hypothèse de récurrence
entraîne que $\mathscr F$ vérifie $(S_{k-1})$, donc
$\operatorname{prof}(\mathscr F_x)\geq k-1$; compte tenu de (0, 15.2.4), il
y a donc un voisinage ouvert $U$ de $x$ et une suite
$(\mathscr F|U)$-régulière $(f_i)_{1\leq i\leq k-1}$ de sections de
$\mathcal O_X$ au-dessus de $U$, telles que $(f_i)_x\in\mathfrak m_x$ pour
$1\leq i\leq k-1$. L'hypothèse entraîne que
\[
\mathscr G=(\mathscr F|U)\big/\Big(\sum_{i=1}^{k-1}
f_i(\mathscr F|U)\Big)
\]
est sans cycle premier associé immergé; mais on a
$\dim(\mathscr G_x)=\dim(\mathscr F_x)-(k-1)\geq1$ (0, 16.3.4), donc $x$
n'est pas associé à $\mathscr G$; on a donc
$\operatorname{prof}(\mathscr G_x)\geq1$ (0, 16.4.6), et comme
$\operatorname{prof}(\mathscr G_x)=\operatorname{prof}(\mathscr F_x)-(k-1)$
(0, 16.4.6), on a $\operatorname{prof}(\mathscr F_x)\geq k$. Supposons en
second lieu que $\dim(\mathscr F_x)=r<k$; comme $\mathscr F$ vérifie
$(S_{k-1})$, on a
$\operatorname{prof}(\mathscr F_x)\geq\inf(k-1,r)=r$, et cela achève la
démonstration.

\begin{corollary}[5.7.6]
\label{IV.2.s5.7.6-fr}
Supposons que $\mathscr F$ vérifie $(S_k)$ ($k\geq1$); si
$(f_i)$ ($1\leq i\leq r$) est une suite $\mathscr F$-régulière de sections
de $\mathcal O_X$ au-dessus de $X$ et si $r<k$,
$\mathscr F/(\sum_{i=1}^r f_i\mathscr F)$ vérifie $(S_{k-r})$.
\end{corollary}

Cela résulte immédiatement de (5.7.5).

\begin{corollary}[5.7.7]
\label{IV.2.s5.7.7-fr}
Pour que $\mathscr F$ vérifie $(S_2)$, il faut et il suffit que $\mathscr F$
soit sans cycle premier associé immergé et que pour tout ouvert $U$ de $X$
et toute section $(\mathscr F|U)$-régulière $f$ de $\mathcal O_X$ au-dessus
de $U$, $(\mathscr F|U)/f(\mathscr F|U)$ soit sans cycle premier associé
immergé.
\end{corollary}

\oldpage[IV-2]{106}

\begin{remarks}[5.7.8]
\label{IV.2.s5.7.8-fr}
Soit $X$ un préschéma localement noethérien de dimension 1; il revient alors
au même de dire que $X$ est un préschéma de Cohen-Macaulay, ou qu'il vérifie
$(S_1)$, ou qu'il vérifie une des propriétés $(S_n)$ pour $n\geq1$, en vertu
des définitions (5.7.1) et (5.7.2). En vertu de (5.7.5), il revient donc
encore au même, pour les préschémas de dimension 1, de dire que $X$ est un
préschéma de Cohen-Macaulay ou qu'il n'a pas de cycle premier associé
immergé. Par exemple, un préschéma localement noethérien réduit de dimension
1 est un préschéma de Cohen-Macaulay.
\end{remarks}

\begin{proposition}[5.7.9]
\label{IV.2.s5.7.9-fr}
Soient $A$, $B$ deux anneaux noethériens, $\rho:A\to B$ un homomorphisme
d'anneaux, $M$ un $B$-module tel que $M_{[\rho]}$ soit un $A$-module de type
fini. Soit $\mathfrak p$ un idéal premier de $A$; les idéaux premiers de $B$
au-dessus de $\mathfrak p$ et appartenant à $\operatorname{Supp}(M)$ sont en
nombre fini, et si $(\mathfrak q_i)_{1\leq i\leq n}$ est la famille de ces
idéaux, on a
\begin{align}
\dim_{A_{\mathfrak p}}(M_{\mathfrak p})
&=\sup_i\dim_{B_{\mathfrak q_i}}(M_{\mathfrak q_i}),
\tag{5.7.9.1}\label{IV.2.s5.7.9.1-fr}\\
\operatorname{prof}_{A_{\mathfrak p}}(M_{\mathfrak p})
&=\inf_i\operatorname{prof}_{B_{\mathfrak q_i}}(M_{\mathfrak q_i}).
\tag{5.7.9.2}\label{IV.2.s5.7.9.2-fr}
\end{align}
\end{proposition}

Si $\mathfrak b$ est l'annulateur de $M$, $B/\mathfrak b$ s'identifie à un
sous-$A$-module de $\operatorname{End}_A(M_{[\rho]})$, donc est de type fini
puisque $A$ est noethérien; il n'y a donc qu'un nombre fini d'idéaux premiers
de $B$ au-dessus de $\mathfrak p$ et contenant $\mathfrak b$, et ce sont
précisément ceux qui appartiennent à $\operatorname{Supp}(M)$. Remplaçant
$B$ par $B/\mathfrak b$, ce qui ne change pas les seconds membres de
(5.7.9.1) et (5.7.9.2) (par (0, 16.1.9) et (0, 16.4.8)), on peut donc
supposer que $B$ est une $A$-algèbre finie. Posons $S=A-\mathfrak p$, et
$B'=S^{-1}B$; $B'$ est un anneau semi-local noethérien dont les idéaux
maximaux sont $S^{-1}\mathfrak q_i$ ($1\leq i\leq n$) et comme
$M_{\mathfrak p}$ est un $A_{\mathfrak p}$-module de type fini, on a
$\dim_{A_{\mathfrak p}}(M_{\mathfrak p})=\dim_{B'}(M_{\mathfrak p})$
(0, 16.1.9). Cela étant, la relation (5.7.9.1) devient un cas particulier de
(0, 16.1.7.4). Pour démontrer la relation (5.7.9.2), on se ramène aussitôt
comme dans (0, 16.4.8) au cas où
$\operatorname{prof}_{A_{\mathfrak p}}(M_{\mathfrak p})=0$, et le même
raisonnement que dans (0, 16.4.8) montre que $M_{\mathfrak p}$ contient un
sous-$B'$-module de longueur finie, et par suite aussi un sous-$B'$-module
simple; mais un tel sous-module est nécessairement isomorphe au corps
résiduel de l'un des $B_{\mathfrak q_i}$, donc il y a au moins un indice $i$
tel que
$\operatorname{prof}_{B_{\mathfrak q_i}}(M_{\mathfrak q_i})=0$
(0, 16.4.6), ce qui termine la démonstration.

\begin{corollary}[5.7.10]
\label{IV.2.s5.7.10-fr}
Supposons vérifiées les hypothèses de (5.7.9), et supposons en outre que $A$
soit un anneau local; alors, pour que $M$ soit un $A$-module de
Cohen-Macaulay, il faut et il suffit que, pour tous les idéaux premiers
$\mathfrak q_i$ de $B$ au-dessus de l'idéal maximal de $A$,
$M_{\mathfrak q_i}$ soit un $B_{\mathfrak q_i}$-module de Cohen-Macaulay, et
de plus que tous les nombres
$\dim_{B_{\mathfrak q_i}}(M_{\mathfrak q_i})$ soient égaux.
\end{corollary}

Il résulte en effet de (5.7.9.1) et de (5.7.9.2) que ces conditions
équivalent à la relation
$\dim_A(M)=\operatorname{prof}_A(M)$.

\begin{corollary}[5.7.11]
\label{IV.2.s5.7.11-fr}
Supposons vérifiées les hypothèses de (5.7.9).
\begin{enumerate}
\item[(i)] Si $M_{[\rho]}$ vérifie la propriété $(S_k)$, il en est de même
de $M$.
\item[(ii)] Supposons que pour tout couple d'idéaux premiers
$\mathfrak q$, $\mathfrak q'$ de $B$ tels que
$\rho^{-1}(\mathfrak q)=\rho^{-1}(\mathfrak q')$, on ait
$\dim_{B_{\mathfrak q}}(M_{\mathfrak q})
=\dim_{B_{\mathfrak q'}}(M_{\mathfrak q'})$. Alors, si $M$ vérifie la
propriété $(S_k)$, il en est de même de $M_{[\rho]}$.
\end{enumerate}
\end{corollary}

\oldpage[IV-2]{107}

Cela résulte aussitôt des relations (5.7.9.1) et (5.7.9.2) et de la
définition de la propriété $(S_k)$.

\begin{env}[5.7.12]
\label{IV.2.s5.7.12-fr}
Conformément aux définitions de (5.7.1), étant donnés un anneau noethérien
quelconque $A$ et un $A$-module $M$ de type fini, on définit
$\operatorname{coprof}_A(M)$ comme égal à
\[
\operatorname{coprof}(M)=\sup_{x\in X}
\operatorname{coprof}_{A_x}(M_x),\qquad X=\operatorname{Spec}(A);
\]
nous verrons plus loin (6.11.5) que cette définition coïncide avec celle de
(0, 16.4.9) lorsque $A$ est un anneau noethérien local.
\end{env}

\begin{corollary}[5.7.13]
\label{IV.2.s5.7.13-fr}
Soient $A$, $B$ deux anneaux noethériens, $\rho:A\to B$ un homomorphisme
d'anneaux, $M$ un $B$-module tel que $M_{[\rho]}$ soit un $A$-module de type
fini. Alors on a
\begin{equation}
\operatorname{coprof}_A(M_{[\rho]})\geq\operatorname{coprof}_B(M).
\tag{5.7.13.1}\label{IV.2.s5.7.13.1-fr}
\end{equation}
\end{corollary}

Cela résulte de la définition précédente et des relations (5.7.9.1) et
(5.7.9.2).

\subsection{\texorpdfstring{Préschémas réguliers et propriété $(R_k)$. Critère de normalité de Serre}{Préschémas réguliers et propriété (R\string_k). Critère de normalité de Serre}}

\begin{env}[5.8.1]
\label{IV.2.s5.8.1-fr}
Rappelons $(0_{\mathrm I},4.1.4)$ qu'un espace annelé
$(X,\mathcal O_X)$ est dit \emph{régulier} en un point $x\in X$, si
$\mathcal O_x$ est un anneau régulier. Quand il s'agira de préschémas, nous
n'utiliserons cette terminologie dans ce chapitre que lorsque $X$ est
localement noethérien.
\end{env}

\begin{definition}[5.8.2]
\label{IV.2.s5.8.2-fr}
On dit qu'un préschéma localement noethérien $X$ est régulier en codimension
$\leq k$, ou possède la propriété $(R_k)$ si l'ensemble des points où $X$
n'est pas régulier est de codimension $>k$ (autrement dit (5.1.3), si, pour
tout $x\in X$, la relation $\dim(\mathcal O_x)\leq k$ entraîne que
$\mathcal O_x$ est régulier).
\end{definition}

Dire que $X$ est régulier signifie que $X$ possède la propriété $(R_k)$ pour
tout $k$.

Si $X=\operatorname{Spec}(A)$, où $A$ est un anneau noethérien, on dira que
$A$ possède la propriété $(R_k)$ si $X$ possède cette propriété; dire que
$X$ est régulier signifie que l'anneau $A$ est régulier (0, 17.3.6). Pour un
préschéma localement noethérien quelconque $X$, on dira que $X$ possède la
propriété $(R_k)$ en un point $x\in X$ si l'anneau local $\mathcal O_x$
possède la propriété $(R_k)$; cela veut donc dire que pour toute générisation
$x'$ de $x$ dans $X$, la relation $\dim(\mathcal O_{x'})\leq k$ entraîne que
$\mathcal O_{x'}$ est un anneau local régulier. Dire que $X$ est régulier en
un point $x$ équivaut à dire que $X$ vérifie la propriété $(R_n)$ pour tout
$n\geq0$ au point $x$, en vertu de (0, 17.3.6).

\begin{proposition}[5.8.3]
\label{IV.2.s5.8.3-fr}
Si $k$ est un corps, $X$ un préschéma localement de type fini sur $k$, alors,
pour tout $x\in X$, il existe un voisinage ouvert de $x$ dans $X$ isomorphe à
un sous-schéma d'un $k$-schéma régulier.
\end{proposition}

En effet, il y a un voisinage ouvert affine $U$ de $x$ isomorphe à un
$k$-schéma de la forme $\operatorname{Spec}(A)$, où $A$ est une $k$-algèbre
de type fini; $A$ est par suite isomorphe à un quotient d'une algèbre de
polynômes $k[T_1,\ldots,T_n]$, et on sait que cette dernière est un anneau
régulier (0, 17.3.7).

\begin{env}[5.8.4]
\label{IV.2.s5.8.4-fr}
En vertu de (5.8.2), dire que $X$ possède la propriété $(R_0)$ signifie que
pour tout point maximal $x$ de $X$, l'anneau $\mathcal O_x$ est un corps
(0, 17.1.4), autrement dit que $X$ est réduit en ce point. Comme l'ensemble
$U$ des $x\in X$ où $X$ est réduit est ouvert (le

\oldpage[IV-2]{108}

nilradical de $X$ étant un $\mathcal O_X$-Module cohérent (I, 6.1.1 et
$0_{\mathrm I}$, 5.2.2)), il revient au même de dire que $X$ possède la
propriété $(R_0)$ ou que l'ensemble $U$ est partout dense. Par suite:
\end{env}

\begin{proposition}[5.8.5]
\label{IV.2.s5.8.5-fr}
Pour qu'un préschéma localement noethérien $X$ soit réduit, il faut et il
suffit qu'il vérifie les propriétés $(S_1)$ et $(R_0)$.
\end{proposition}

Compte tenu de (5.7.5), cela résulte de la remarque précédente et de (3.2.1).

\begin{theorem}[5.8.6]
\EGAReaderTarget{ega4.c04.thm.5-8-6}{IV.2.s5.8.6-fr}
(critère de Serre). Soit $X$ un préschéma localement noethérien. Pour que
$X$ soit normal, il faut et il suffit que $X$ vérifie les propriétés $(S_2)$
et $(R_1)$, autrement dit, que pour tout $x\in X$, on ait les propriétés
suivantes:
\begin{enumerate}
\item[(i)] Si $\dim(\mathcal O_x)\leq1$, $\mathcal O_x$ est régulier
(c'est-à-dire est un corps ou un anneau de valuation discrète (0, 17.1.4)).
\item[(ii)] Si $\dim(\mathcal O_x)\geq2$, alors
$\operatorname{prof}(\mathcal O_x)\geq2$.
\end{enumerate}
\end{theorem}

Les conditions sont nécessaires. En effet, dire que $X$ est normal signifie
que pour tout $x\in X$, $\mathcal O_x$ est un anneau local noethérien
intégralement clos. Si $\dim(\mathcal O_x)=0$ (resp.
$\dim(\mathcal O_x)=1$), on en conclut que $\mathcal O_x$ est un corps
puisque $\mathcal O_x$ est intègre (resp. que $\mathcal O_x$ est un anneau
de valuation discrète, en vertu de (II, 7.1.6)). D'autre part, pour tout
élément $f_x\neq0$ de $\mathcal O_x$, on sait (Bourbaki, \emph{Alg. comm.},
chap. VII, § 1, n° 4, prop. 8) que les idéaux premiers associés à
$\mathcal O_x/f_x\mathcal O_x$ sont non immergés, donc $\mathcal O_x$
vérifie $(S_2)$ (5.7.7).

Les conditions sont suffisantes. En effet, il résulte d'abord de (5.8.5) que
$X$ est réduit. La question étant locale, on peut en outre supposer que
$X=\operatorname{Spec}(A)$, où $A$ est un anneau noethérien réduit
(I, 5.1.4); si $R$ est l'anneau total des fractions de $A$, $R$ est composé
direct d'un nombre fini de corps, et (compte tenu de (II, 6.3.6)), il suffira
de prouver que $A$ est intégralement fermé dans $R$. Soit donc $h=f/g$ un
élément de $R$ entier sur $A$, $g$ et $f$ étant des éléments de $A$ tels que
$g$ soit non diviseur de 0. On a une relation de la forme
\begin{equation}
f^n+\sum_{i=1}^n a_i f^{n-i}g^i=0
\qquad\text{avec }a_i\in A\quad\text{pour }1\leq i\leq n.
\tag{5.8.6.1}\label{IV.2.s5.8.6.1-fr}
\end{equation}

Soit $\mathfrak p$ un idéal premier de $A$ tel que
$\dim(A_{\mathfrak p})=1$; si $f_{\mathfrak p}$ et $g_{\mathfrak p}$ sont
les images de $f$ et $g$ dans $A_{\mathfrak p}$, il résulte de (5.8.6.1) que
$f_{\mathfrak p}/g_{\mathfrak p}$ (qui appartient à l'anneau total des
fractions de $A_{\mathfrak p}$, puisque $g_{\mathfrak p}$ est non diviseur
de 0 dans $A_{\mathfrak p}$ par platitude $(0_{\mathrm I},5.3.1)$) est entier
sur $A_{\mathfrak p}$; mais comme $A_{\mathfrak p}$ est régulier, donc
intégralement clos, on a $f_{\mathfrak p}/g_{\mathfrak p}\in A_{\mathfrak p}$.
En d'autres termes, on a
$(fA)_{\mathfrak p}\subset(gA)_{\mathfrak p}$. Mais l'hypothèse $(S_2)$
entraîne (5.7.7), puisque $g$ est non diviseur de 0 dans $A$, que $A/gA$
n'a que des idéaux premiers associés non immergés $\mathfrak p_i$
($1\leq i\leq n$); or, $gA$ est l'intersection d'idéaux primaires
$\mathfrak q_i$ correspondant aux $\mathfrak p_i$, et d'après ce qu'on vient
de voir, les $\mathfrak q_i$ sont les images réciproques dans $A$, par les
homomorphismes $A\to A_{\mathfrak p_i}$, des idéaux
$(gA)_{\mathfrak p_i}$ (Bourbaki, \emph{Alg. comm.}, chap. IV, § 2, n° 3,
prop. 5). Mais en vertu du Hauptidealsatz (0, 16.3.2) on a
$\dim(A_{\mathfrak p_i})=1$ pour $1\leq i\leq n$, donc
$(fA)_{\mathfrak p_i}\subset(gA)_{\mathfrak p_i}$ pour tout $i$ d'après ce
qui précède; comme $fA$ est contenu dans l'intersection des images
réciproques des $(fA)_{\mathfrak p_i}$ ($1\leq i\leq n$), on a
$fA\subset gA$, c'est-à-dire $f/g\in A$. C.Q.F.D.

\oldpage[IV-2]{109}

\subsection{\texorpdfstring{Modules $Z$-purs et $Z$-clos}{Modules Z-purs et Z-clos}}

Une partie des notions et résultats de cette section et de la suivante sont des
cas particuliers de notions et résultats développés au chapitre III dans la
théorie de la cohomologie locale. Pour la commodité du lecteur, nous en donnons
ici un exposé indépendant.

\begin{env}[5.9.1]
\label{IV.2.s5.9.1-fr}
Soient $X$ un préschéma localement noethérien, $Z$ une partie de $X$
\emph{stable par spécialisation}: cela signifie que pour toute partie finie
$M$ de $Z$, l'adhérence de $M$ est contenue dans $Z$, et par suite $Z$ est
réunion d'une famille filtrante croissante $(Z_\alpha)$ de parties fermées de
$X$; inversement, il est clair qu'une telle réunion est stable par
spécialisation.

Posons $U_\alpha=X-Z_\alpha$, de sorte que $X-Z$ est intersection de la
famille filtrante décroissante d'ouverts $U_\alpha$; soit
$i_\alpha:U_\alpha\to X$ l'injection canonique et pour
$U_\alpha\supset U_\beta$, soit $i_{\alpha\beta}:U_\beta\to U_\alpha$
l'injection canonique, de sorte que l'on a
$i_\beta=i_\alpha\circ i_{\alpha\beta}$. Soit $\mathcal F$ un
$\mathcal O_X$-Module (non nécessairement quasi-cohérent); on a donc
\[
(i_\beta)_*(\mathcal F|U_\beta)
=(i_\alpha)_*((i_{\alpha\beta})_*(\mathcal F|U_\beta));
\]
de l'homomorphisme canonique $(0_{\mathrm I},4.4.3.2)$
\[
\mathcal F|U_\alpha\longrightarrow
(i_{\alpha\beta})_*(\mathcal F|U_\beta)
\]
on déduit donc, par application du foncteur $(i_\alpha)_*$, un homomorphisme
\[
\rho_{\beta\alpha}:(i_\alpha)_*(\mathcal F|U_\alpha)
\longrightarrow(i_\beta)_*(\mathcal F|U_\beta)
\]
et l'on vérifie aussitôt que l'on a
$\rho_{\gamma\alpha}=\rho_{\gamma\beta}\circ\rho_{\beta\alpha}$ pour
$U_\alpha\supset U_\beta\supset U_\gamma$; autrement dit, les
$\mathcal O_X$-Modules $(i_\alpha)_*(\mathcal F|U_\alpha)$ forment un
\emph{système inductif} pour les homomorphismes $\rho_{\beta\alpha}$.
On pose
\begin{equation}
\mathcal H^0_{X/Z}(\mathcal F)
=\varinjlim_\alpha(i_\alpha)_*(\mathcal F|U_\alpha).
\tag{5.9.1.1}\label{IV.2.s5.9.1.1-fr}
\end{equation}

Ce $\mathcal O_X$-Module ne dépend pas de la famille croissante
$(Z_\alpha)$ de fermés dont $Z$ est la réunion: en effet, soit $V$ un ouvert
noethérien de $X$; on sait $(0_{\mathrm{II}},3.10.1)$ que dans la catégorie
des $\mathcal O_V$-Modules, le foncteur
$\mathcal G\mapsto\Gamma(V,\mathcal G)$ commute aux limites inductives; on a
donc en vertu de (5.9.1.1)
\begin{equation}
\Gamma(V,\mathcal H^0_{X/Z}(\mathcal F))
=\varinjlim_\alpha\Gamma(V\cap U_\alpha,\mathcal F).
\tag{5.9.1.2}\label{IV.2.s5.9.1.2-fr}
\end{equation}

Soit alors $(Z'_\lambda)$ une seconde famille filtrante croissante de fermés
de $X$, de réunion $Z$; $V\cap Z_\alpha$ est donc réunion des
$V\cap Z_\alpha\cap Z'_\lambda$; mais $V\cap Z_\alpha$ est localement fermé
dans $X$, donc toute partie fermée irréductible de $V\cap Z_\alpha$ admet un
point générique; comme les $V\cap Z_\alpha\cap Z'_\lambda$ sont fermés dans
$V\cap Z_\alpha$ et forment (pour $\alpha$ fixe) une famille filtrante
croissante, il existe un indice $\lambda$ tel que
$V\cap Z_\alpha\cap Z'_\lambda=V\cap Z_\alpha$
$(0_{\mathrm{III}},9.2.4)$, autrement dit
$V\cap Z_\alpha\subset V\cap Z'_\lambda$. Ceci prouve que les familles
filtrantes décroissantes $V\cap U_\alpha$, $V\cap U'_\lambda$ (où
$U'_\lambda=X-Z'_\lambda$) sont cofinales l'une de l'autre, d'où notre
assertion, en vertu de (5.9.1.2).
\end{env}

\oldpage[IV-2]{110}

On notera que l'ensemble $Z$ n'est pas nécessairement constructible: on a un
exemple de ce fait en prenant $X=\operatorname{Spec}(A)$, où $A$ est un
anneau intègre noethérien ayant une infinité d'idéaux maximaux, et $Z$ le
complémentaire du point générique de $X$.

Si $Z$ est fermé et si $i:X-Z\to X$ est l'injection canonique, on a
\begin{equation}
\mathcal H^0_{X/Z}(\mathcal F)=i_*(\mathcal F|X-Z)
\tag{5.9.1.3}\label{IV.2.s5.9.1.3-fr}
\end{equation}
et en particulier, pour $Z=\varnothing$,
$\mathcal H^0_{X/Z}(\mathcal F)=\mathcal F$.

\begin{proposition}[5.9.2]
\label{IV.2.s5.9.2-fr}
\begin{enumerate}
\item[(i)] Le foncteur
$\mathcal F\mapsto\mathcal H^0_{X/Z}(\mathcal F)$ est exact à gauche.
\item[(ii)] Si $\mathcal F$ est quasi-cohérent, il en est de même de
$\mathcal H^0_{X/Z}(\mathcal F)$.
\end{enumerate}
\end{proposition}

L'assertion (i) résulte de la définition (5.9.1.1), du fait que
$(i_\alpha)_*$ est un foncteur exact à gauche et de ce que la limite
inductive préserve l'exactitude dans la catégorie des $\mathcal O_X$-Modules.
L'assertion (ii) résulte de (I, 9.2.2) et du fait qu'une limite inductive de
$\mathcal O_X$-Modules quasi-cohérents est quasi-cohérente (I, 1.3.9).

\begin{remark}[5.9.3]
\label{IV.2.s5.9.3-fr}
Si $\mathcal F$ est une $\mathcal O_X$-Algèbre, il en est de même de
$\mathcal H^0_{X/Z}(\mathcal F)$ $(0_{\mathrm I},4.2.4)$; en particulier
$\mathcal H^0_{X/Z}(\mathcal O_X)$ est une $\mathcal O_X$-Algèbre
quasi-cohérente, et pour tout $\mathcal O_X$-Module $\mathcal F$,
$\mathcal H^0_{X/Z}(\mathcal F)$ est un
$\mathcal H^0_{X/Z}(\mathcal O_X)$-Module, qui est quasi-cohérent si
$\mathcal F$ est quasi-cohérent (I, 9.6.1). Plus particulièrement,
supposons que $X=\operatorname{Spec}(A)$, où $A$ est intègre et noethérien;
alors $\mathcal H^0_{X/Z}(\mathcal O_X)$ est la $\mathcal O_X$-Algèbre
$\widetilde B$, où
\begin{equation}
B=\bigcap_{\mathfrak p\in X-Z}A_{\mathfrak p}.
\tag{5.9.3.1}\label{IV.2.s5.9.3.1-fr}
\end{equation}
Cela résulte en effet de (5.9.1.2) et de (I, 8.2.1.1).
\end{remark}

\begin{proposition}[5.9.4]
\label{IV.2.s5.9.4-fr}
Soient $X$ un préschéma localement noethérien, $Z$ une partie de $X$ stable
par spécialisation, $X'$ un préschéma localement noethérien,
$f:X'\to X$ un morphisme plat. Alors $Z'=f^{-1}(Z)$ est stable par
spécialisation et pour tout $\mathcal O_X$-Module quasi-cohérent $\mathcal F$,
on a un isomorphisme canonique
\begin{equation}
f^*(\mathcal H^0_{X/Z}(\mathcal F))
\xrightarrow{\sim}\mathcal H^0_{X'/Z'}(f^*(\mathcal F)).
\tag{5.9.4.1}\label{IV.2.s5.9.4.1-fr}
\end{equation}
\end{proposition}

En effet, avec les notations de (5.9.1), $Z'_\alpha=f^{-1}(Z_\alpha)$ est
fermé dans $X'$ et $Z'$ est réunion des $Z'_\alpha$; en outre,
$(i_\alpha)_{(X')}$ est l'injection canonique
$i'_\alpha:U'_\alpha\to X'$, si
$U'_\alpha=X'-Z'_\alpha=f^{-1}(U_\alpha)$. Comme $f$ est plat, on sait
(2.3.1) que l'homomorphisme canonique
\[
f^*((i_\alpha)_*(\mathcal F|U_\alpha))
\longrightarrow(i'_\alpha)_*(f^*(\mathcal F)|U'_\alpha)
\]
est bijectif; comme, pour $\alpha\leq\beta$, le diagramme
\[
\xymatrix{
f^*((i_\alpha)_*(\mathcal F|U_\alpha))
  \ar[r]^{\sim}\ar[d] &
(i'_\alpha)_*(f^*(\mathcal F)|U'_\alpha)\ar[d]\\
f^*((i_\beta)_*(\mathcal F|U_\beta))
  \ar[r]^{\sim} &
(i'_\beta)_*(f^*(\mathcal F)|U'_\beta)
}
\]
est commutatif, on a donc, en passant à la limite, un isomorphisme canonique
\[
\varinjlim_\alpha
f^*((i_\alpha)_*(\mathcal F|U_\alpha))
\xrightarrow{\sim}\mathcal H^0_{X'/Z'}(f^*(\mathcal F)).
\]

\oldpage[IV-2]{111}

Mais comme le foncteur $f^*$ commute aux limites inductives
$(0_{\mathrm I},4.3.2)$, cela donne par définition l'isomorphisme
(5.9.1.1) cherché.

\begin{corollary}[5.9.5]
\label{IV.2.s5.9.5-fr}
Sous les hypothèses de (5.9.4), si
$\mathcal H^0_{X/Z}(\mathcal F)$ est cohérent, il en est de même de
$\mathcal H^0_{X'/Z'}(f^*(\mathcal F))$. La réciproque est vraie lorsque
$f$ est un morphisme fidèlement plat et quasi-compact.
\end{corollary}

La première assertion résulte de (5.9.4.1) et de $(0_{\mathrm I},5.3.11)$;
la seconde équivaut à dire que si
$\mathcal H^0_{X'/Z'}(f^*(\mathcal F))$ est de type fini, il en est de même
de $\mathcal H^0_{X/Z}(\mathcal F)$; cela résulte de (5.9.4.1) et de
(2.5.2).

\begin{corollary}[5.9.6]
\label{IV.2.s5.9.6-fr}
Soient $X$ un préschéma localement noethérien, $Z$ une partie de $X$ stable
par spécialisation, $\mathcal F$ un $\mathcal O_X$-Module quasi-cohérent.
Pour tout $x\in X$, posons $X_x=\operatorname{Spec}(\mathcal O_x)$,
$Z_x=Z\cap X_x$; on a un isomorphisme canonique fonctoriel
\begin{equation}
(\mathcal H^0_{X/Z}(\mathcal F))_x
\xrightarrow{\sim}\mathcal H^0_{X_x/Z_x}(\widetilde{\mathcal F}_x).
\tag{5.9.6.1}\label{IV.2.s5.9.6.1-fr}
\end{equation}
\end{corollary}

Il suffit d'appliquer (5.9.4) au morphisme canonique $X_x\to X$, qui est
plat, et de tenir compte de (I, 1.6.5).

\begin{env}[5.9.7]
\label{IV.2.s5.9.7-fr}
Avec les notations de (5.9.1), on a pour tout $\alpha$ un homomorphisme
canonique fonctoriel
\[
\mathcal F\longrightarrow(i_\alpha)_*(\mathcal F|U_\alpha)
\quad(0_{\mathrm I},4.4.3.2),
\]
et ces homomorphismes forment un système inductif; par passage à la limite
inductive, on en déduit donc un homomorphisme canonique fonctoriel
\begin{equation}
\rho_{X/Z}:\mathcal F\longrightarrow\mathcal H^0_{X/Z}(\mathcal F).
\tag{5.9.7.1}\label{IV.2.s5.9.7.1-fr}
\end{equation}
\end{env}

\begin{proposition}[5.9.8]
\EGAReaderTarget{ega4.c04.prop.5-9-8}{IV.2.s5.9.8-fr}
Soient $X$ un préschéma localement noethérien, $Z$ une partie de $X$ stable
par spécialisation, $\mathcal F$ un $\mathcal O_X$-Module. Les propriétés
suivantes sont équivalentes:
\begin{enumerate}
\item[a)] L'homomorphisme $\rho_{X/Z}$ (5.9.7.1) est injectif
(resp. bijectif).
\item[b)] Pour tout ouvert noethérien $V$ de $X$, l'homomorphisme
\[
(\rho_{X/Z})_V:\Gamma(V,\mathcal F)\longrightarrow
\varinjlim_\alpha\Gamma(V\cap U_\alpha,\mathcal F)
\]
est injectif (resp. bijectif).
\item[a')] Pour toute partie fermée $T\subset Z$ de $X$, l'homomorphisme
canonique $(0_{\mathrm I},4.4.3.2)$
\[
\mathcal F\longrightarrow i_*(\mathcal F|X-T)
\]
(où $i:X-T\to X$ est l'injection canonique) est injectif
(resp. bijectif).
\item[b')] Pour toute partie fermée $T\subset Z$ de $X$ et tout ouvert
noethérien $V$ de $X$, l'homomorphisme de restriction
\[
\Gamma(V,\mathcal F)\longrightarrow
\Gamma(V\cap(X-T),\mathcal F)
\]
est injectif (resp. bijectif).
\end{enumerate}
\end{proposition}

Compte tenu de (5.9.1.2), l'équivalence de a) et b) (resp. a') et b'))
résulte de
\oldpage[IV-2]{112}

la définition du foncteur $\Gamma$ et du fait qu'il est exact à gauche.
Comme l'homomorphisme $(\rho_{X/Z})_V$ est le composé
\begin{equation}
\Gamma(V,\mathcal F)\longrightarrow
\Gamma(V\cap(X-T),\mathcal F)\longrightarrow
\varinjlim_\alpha\Gamma(V\cap U_\alpha,\mathcal F)
\tag{5.9.8.1}\label{IV.2.s5.9.8.1-fr}
\end{equation}
pour toute partie fermée $T\subset Z$, si $(\rho_{X/Z})_V$ est injectif, il
en est de même de
$\Gamma(V,\mathcal F)\to\Gamma(V\cap(X-T),\mathcal F)$; d'autre part, le
fait que b') implique b) résulte de la définition d'une limite inductive. Il
reste à montrer que si $\rho_{X/Z}$ est bijectif, il en est de même de
$\Gamma(V,\mathcal F)\to\Gamma(V\cap(X-T),\mathcal F)$, et pour cela il
suffit, en vertu de (5.9.8.1), de voir que si $U'\subset U$ sont deux
ouverts contenus dans $V$ et contenant $V\cap Z$, l'homomorphisme de
restriction $\Gamma(U,\mathcal F)\to\Gamma(U',\mathcal F)$ est injectif;
mais cela résulte de ce que $\rho_{X/Z}$ est injectif, en remplaçant dans ce
qui précède $V$ par $U$ et $V\cap(X-T)$ par $U'$.

\begin{definition}[5.9.9]
\label{IV.2.s5.9.9-fr}
Sous les hypothèses de (5.9.8), on dit que $\mathcal F$ est $Z$-pur
(resp. $Z$-clos) si l'homomorphisme $\rho_{X/Z}$ est injectif
(resp. bijectif).
\end{definition}

Si $X=\operatorname{Spec}(A)$ est affine, $\mathcal F=\widetilde M$, où
$M$ est un $A$-module, on dira que $M$ est $Z$-pur (resp. $Z$-clos) lorsque
$\mathcal F$ est $Z$-pur (resp. $Z$-clos).

On dit que $\mathcal F$ est $Z$-pur (resp. $Z$-clos) \emph{en un point}
$x\in X$ si (avec les notations de (5.9.6)) $\widetilde{\mathcal F}_x$ est
$Z_x$-pur (resp. $Z_x$-clos); il revient au même, en vertu de (5.9.6), de
dire que l'homomorphisme canonique
\[
\mathcal F_x\longrightarrow(\mathcal H^0_{X/Z}(\mathcal F))_x
\]
est injectif (resp. bijectif).

On notera que pour tout $x\in X-Z$, $\mathcal F$ est $Z$-clos au point $x$,
en vertu de (5.9.8).

\begin{corollary}[5.9.10]
\label{IV.2.s5.9.10-fr}
\begin{enumerate}
\item[(i)] Soit $(V_\lambda)$ un recouvrement ouvert de $X$. Pour que
$\mathcal F$ soit $Z$-pur (resp. $Z$-clos), il faut et il suffit que pour
tout $\lambda$, $\mathcal F|V_\lambda$ soit $(Z\cap V_\lambda)$-pur
(resp. $(Z\cap V_\lambda)$-clos).
\item[(ii)] Soit $Z'$ une partie de $Z$ stable par spécialisation. Si
$\mathcal F$ est $Z$-pur (resp. $Z$-clos), il est $Z'$-pur
(resp. $Z'$-clos).
\end{enumerate}
\end{corollary}

Cela résulte aussitôt de (5.9.8, b')).

\begin{proposition}[5.9.11]
\label{IV.2.s5.9.11-fr}
Sous les hypothèses de (5.9.8), les $\mathcal O_X$-Modules
$\operatorname{Ker}(\rho_{X/Z})$ et $\operatorname{Coker}(\rho_{X/Z})$ ont
leur support contenu dans $Z$, et le $\mathcal O_X$-Module
$\mathcal H^0_{X/Z}(\mathcal F)$ est $Z$-clos. En outre, si
$u:\mathcal F\to\mathcal F'$ est un homomorphisme de
$\mathcal O_X$-Modules tel que $\mathcal F'$ soit $Z$-clos, $u$ se factorise
d'une seule manière en
\[
\mathcal F\xrightarrow{\rho_{X/Z}}\mathcal H^0_{X/Z}(\mathcal F)
\xrightarrow{v}\mathcal F'.
\]
Si en outre les supports de $\operatorname{Ker}(u)$ et de
$\operatorname{Coker}(u)$ sont contenus dans $Z$, $v$ est un isomorphisme.
\end{proposition}

La première assertion signifie que, pour tout $x\in X-Z$, on a
\[
\operatorname{Ker}(\rho_{X/Z})_x
=\operatorname{Coker}(\rho_{X/Z})_x=0,
\]
i.e. que $(\rho_{X/Z})_x$ est bijectif, ou encore que $\mathcal F$ est
$Z$-pur en $x$, comme il a été signalé plus haut (5.9.9).

Pour montrer que $\mathcal H^0_{X/Z}(\mathcal F)$ est $Z$-clos, il s'agit de
voir que pour un ouvert noethérien $V$ de $X$,
\[
\varinjlim_\beta
\Gamma(V\cap U_\beta,\mathcal H^0_{X/Z}(\mathcal F))
\]
est égale à $\Gamma(V,\mathcal H^0_{X/Z}(\mathcal F))$; mais par définition
\[
\Gamma(V\cap U_\beta,\mathcal H^0_{X/Z}(\mathcal F))
=\varinjlim_\alpha\Gamma(V\cap U_\alpha\cap U_\beta,\mathcal F).
\]
Or, la famille double $(U_\alpha\cap U_\beta)$ est filtrante décroissante,
et notre assertion résulte de (5.9.1) et du théorème de la double limite
inductive.

\oldpage[IV-2]{113}

Passons à la seconde partie de la proposition. L'existence et l'unicité de
$v$ résultent de ce que pour tout $\alpha$, $(i_\alpha)_*(u)$ est l'unique
homomorphisme rendant commutatif le diagramme
\[
\xymatrix{
\mathcal F\ar[r]^u\ar[d] & \mathcal F'\ar[d]\\
(i_\alpha)_*(\mathcal F|U_\alpha)
\ar[r]_{(i_\alpha)_*(u)} &
(i_\alpha)_*(\mathcal F'|U_\alpha)
}
\]
de ce qu'il y a un unique homomorphisme $w$ rendant commutatifs tous les
diagrammes
\[
\xymatrix{
(i_\alpha)_*(\mathcal F|U_\alpha)
\ar[r]^{(i_\alpha)_*(u)}\ar[d] &
(i_\alpha)_*(\mathcal F'|U_\alpha)\ar[d]\\
\mathcal H^0_{X/Z}(\mathcal F)\ar[r]_w &
\mathcal H^0_{X/Z}(\mathcal F')
}
\]
et enfin de ce que $\mathcal F'$ et
$\mathcal H^0_{X/Z}(\mathcal F')$ s'identifient canoniquement par hypothèse.

Reste à voir que si les supports de $\operatorname{Ker}(u)$ et de
$\operatorname{Coker}(u)$ sont contenus dans $Z$, $v$ est un isomorphisme.
Il suffit de voir que pour tout ouvert noethérien $V$, l'homomorphisme
correspondant
\[
\Gamma(V,\mathcal H^0_{X/Z}(\mathcal F))
\longrightarrow\Gamma(V,\mathcal F')
\]
est alors un isomorphisme. Or, si une section
$t\in\Gamma(V,\mathcal H^0_{X/Z}(\mathcal F))$ a pour image $0$ dans
$\Gamma(V,\mathcal F')$, notons que pour un indice $\alpha$, on a
$t\in\Gamma(V\cap U_\alpha,\mathcal F)$, et en vertu de l'hypothèse sur
$u$, on a $t_y=0$ pour tout $y\in V\cap(X-Z)$; il y a par suite un ouvert
contenant $V\cap(X-Z)$ tel que la restriction de $t$ à cet ouvert soit
nulle, donc par définition $t$ est l'élément $0$ de
$\Gamma(V,\mathcal H^0_{X/Z}(\mathcal F))$.

Prouvons maintenant que toute section $s'\in\Gamma(V,\mathcal F')$ est
l'image d'une section de $\mathcal H^0_{X/Z}(\mathcal F)$ au-dessus de $V$.
Par hypothèse, pour tout $x\in V\cap(X-Z)$ il existe une section $s^{(x)}$
de $\mathcal F$ au-dessus d'un voisinage ouvert $W^{(x)}$ de $x$ dans $X$,
dont l'image par $u$ est $s'|W^{(x)}$; $s'|W^{(x)}$ est donc aussi l'image
par $v$ de la section $t^{(x)}$ de $\mathcal H^0_{X/Z}(\mathcal F)$, image
canonique de $s^{(x)}$. En outre, comme on a vu que $v$ est injectif, les
restrictions de $t^{(x)}$ et $t^{(x')}$ à $W^{(x)}\cap W^{(x')}$ sont
identiques pour deux points quelconques $x$, $x'$ de $V\cap(X-Z)$; les
$t^{(x)}$ sont par suite restrictions d'une même section $t$ de
$\mathcal H^0_{X/Z}(\mathcal F)$ au-dessus d'un voisinage ouvert $U$ de
$(X-Z)\cap V$. Mais comme $\mathcal H^0_{X/Z}(\mathcal F)$ est $Z$-clos,
$t$ se prolonge d'une seule manière en une section de
$\mathcal H^0_{X/Z}(\mathcal F)$ au-dessus de $V$, dont l'image par $v$ a
même restriction que $s'$ à $U$, et coïncide donc avec $s'$ pour la même
raison.

On dit que $\mathcal H^0_{X/Z}(\mathcal F)$ est la $Z$-clôture de
$\mathcal F$.

\oldpage[IV-2]{114}

\begin{remarks}[5.9.12]
\label{IV.2.s5.9.12-fr}
\begin{enumerate}
\item[(i)] Soit $\mathcal C(X)$ la catégorie des $\mathcal O_X$-Modules, et
soit $\mathcal C_Z(X)$ la sous-catégorie de $\mathcal C(X)$ formée des
$\mathcal O_X$-Modules de support contenu dans $Z$; cette sous-catégorie est
localisante au sens de Gabriel, et le foncteur $\mathcal H^0_{X/Z}$ n'est
autre que le foncteur localisation de Gabriel (cf. [27]; cela fournirait une
autre démonstration de (5.9.11)). Lorsque $Z$ est fermé, le foncteur
$i^*:\mathcal C(X)\to\mathcal C(X-Z)$ (où $i:X-Z\to X$ est l'injection
canonique) définit une équivalence de catégories
\[
\mathcal C(X)/\mathcal C_Z(X)\approx\mathcal C(X-Z).
\]
\item[(ii)] Il résulte de (5.9.11) que la condition
$\mathcal H^0_{X/Z}(\mathcal F)=0$ est équivalente à
$\operatorname{Supp}(\mathcal F)\subset Z$. Elle entraîne en effet cette
dernière puisque le noyau de $\rho_{X/Z}$ est alors égal à $\mathcal F$.
Inversement, si $\operatorname{Supp}(\mathcal F)\subset Z$, il suffit
d'appliquer la seconde partie de (5.9.11) à l'unique homomorphisme
$u:\mathcal F\to0$ pour en conclure que l'homomorphisme correspondant
$v:\mathcal H^0_{X/Z}(\mathcal F)\to0$ est un isomorphisme.
\item[(iii)] Les développements précédents gardent un sens pour tout espace
annelé localement noethérien dont toute partie fermée irréductible admet
exactement un point générique. En particulier ils s'appliquent, sur un
préschéma localement noethérien, à des \emph{faisceaux de groupes abéliens}
quelconques (considérés comme Modules sur le faisceau simple associé au
préfaisceau constant $\mathbf Z$). On a encore pour tout $x\in X$
l'isomorphisme canonique (5.9.6.1), dans lequel
$\widetilde{\mathcal F}_x$ désigne le faisceau induit sur le sous-espace
$X_x$ de $X$ par le faisceau $\mathcal F$; la démonstration directe résulte
aussitôt de la définition (5.9.1.2) et du théorème de la double limite
inductive.
\end{enumerate}
\end{remarks}

\subsection{\texorpdfstring{Propriété $(S_2)$ et $Z$-clôture}{Propriété (S\string_2) et Z-clôture}}

\begin{env}[5.10.1]
\label{IV.2.s5.10.1-fr}
Soient $X$ un préschéma localement noethérien, $\mathcal F$ un
$\mathcal O_X$-Module cohérent; pour toute partie $T$ de $X$, nous poserons
\begin{equation}
\operatorname{prof}_T(\mathcal F)
=\inf_{x\in T}\operatorname{prof}(\mathcal F_x).
\tag{5.10.1.1}\label{IV.2.s5.10.1.1-fr}
\end{equation}
\end{env}

\begin{proposition}[5.10.2]
\EGAReaderTarget{ega4.c04.prop.5-10-2}{IV.2.s5.10.2-fr}
Soient $X$ un préschéma localement noethérien, $Z$ une partie de $X$ stable
par spécialisation, $\mathcal F$ un $\mathcal O_X$-Module quasi-cohérent.
Les conditions suivantes sont équivalentes:
\begin{enumerate}
\item[a)] $\mathcal F$ est $Z$-pur.
\item[b)] $\operatorname{Ass}(\mathcal F)$ ne rencontre pas $Z$.
\end{enumerate}
Si en outre $\mathcal F$ est cohérent, ces conditions sont aussi équivalentes
à la suivante:
\begin{enumerate}
\item[c)] $\operatorname{prof}_Z(\mathcal F)\geq1$.
\end{enumerate}
\end{proposition}

Dire que $\mathcal F$ est $Z$-pur signifie que pour tout ouvert noethérien
$V$ de $X$, et tout ouvert $U\supset X-Z$, l'homomorphisme de restriction
\[
\Gamma(V,\mathcal F)\longrightarrow\Gamma(V\cap U,\mathcal F)
\]
est injectif (5.9.8); mais d'après (3.1.8) cela équivaut à
$V\cap\operatorname{Ass}(\mathcal F)\subset U$, d'où l'équivalence de a) et
b). En outre, dire que $x\in\operatorname{Ass}(\mathcal F)$ signifie
qu'aucun élément de $\mathfrak m_x$ n'est $\mathcal F_x$-régulier (3.1.2),
donc, lorsque $\mathcal F$ est cohérent, cela s'écrit encore
$\operatorname{prof}(\mathcal F_x)=0$; on en déduit aussitôt dans ce cas
l'équivalence de b) et c).

\begin{corollary}[5.10.3]
\label{IV.2.s5.10.3-fr}
Soit
\[
0\longrightarrow\mathcal F'\longrightarrow\mathcal F
\longrightarrow\mathcal F''\longrightarrow0
\]
une suite exacte de $\mathcal O_X$-Modules quasi-cohérents. Si $\mathcal F$
est $Z$-pur, il en est de même de $\mathcal F'$; inversement, si
$\mathcal F'$ et $\mathcal F''$ sont $Z$-purs, il en est de même de
$\mathcal F$.
\end{corollary}

\oldpage[IV-2]{115}

Cela résulte de la forme (5.10.2, b)) de la condition pour qu'un
$\mathcal O_X$-Module quasi-cohérent soit $Z$-pur, et de (3.1.7).

\begin{corollary}[5.10.4]
\label{IV.2.s5.10.4-fr}
Supposons $\mathcal F$ cohérent. Pour que $\mathcal F$ soit $Z$-pur en un
point $x\in X$, il faut et il suffit que
\[
\operatorname{prof}_{Z_x}(\widetilde{\mathcal F}_x)\geq1
\]
(avec les notations de (5.9.6)).
\end{corollary}

Cela résulte aussitôt de (5.10.2) et (5.9.6).

\begin{theorem}[5.10.5]
\EGAReaderTarget{ega4.c04.thm.5-10-5}{IV.2.s5.10.5-fr}
Soient $X$ un préschéma localement noethérien, $Z$ une partie de $X$ stable
par spécialisation, $\mathcal F$ un $\mathcal O_X$-Module cohérent. Pour que
$\mathcal F$ soit $Z$-clos, il faut et il suffit que l'on ait
\[
\operatorname{prof}_Z(\mathcal F)\geq2.
\]
\end{theorem}

En vertu de (5.10.2), on peut se borner au cas où $\mathcal F$ est $Z$-pur
et $\operatorname{prof}_Z(\mathcal F)\geq1$. De plus, dire que
$\operatorname{prof}_Z(\mathcal F)\geq2$ équivaut à dire que pour toute partie
fermée $Z_\alpha$ de $Z$, $\operatorname{prof}_{Z_\alpha}(\mathcal F)\geq2$;
et de même, il résulte de (5.9.8) que dire que $\mathcal F$ est $Z$-clos
équivaut à dire que $\mathcal F$ est $Z_\alpha$-clos pour tout $\alpha$. On
peut donc déjà se borner au cas où $Z$ est fermé. La question étant locale,
il suffit, pour tout $x\in Z$, de prouver le théorème pour $\mathcal F|U$,
$U$ étant un voisinage ouvert affine de $x$, et l'on peut donc se borner au
cas où $X=U$ est affine. On sait alors que $\operatorname{Ass}(\mathcal F)$
est fini (3.1.6), et comme
$\operatorname{Ass}(\mathcal F)\subset X-Z$, il y a une section $f$ de
$\mathcal O_X$ au-dessus de $X$ telle que
$\operatorname{Ass}(\mathcal F)\subset X_f\subset X-Z$ (II, 4.5.4); on en
déduit que $f$ est $\mathcal F$-régulière (3.1.9) et que pour tout $y\in Z$,
on a $f_y\in\mathfrak m_y$, donc
\[
\operatorname{prof}(\mathcal F_y)
=1+\operatorname{prof}(\mathcal F_y/f_y\mathcal F_y)
\quad(0, 16.4.6).
\]
La condition $\operatorname{prof}_Z(\mathcal F)\geq2$ équivaut donc à
$\operatorname{prof}_Z(\mathcal F/f\mathcal F)\geq1$, ou encore (5.10.2) au
fait que $\mathcal F/f\mathcal F$ est $Z$-pur, et il suffit de voir que cette
dernière propriété équivaut au fait que $\mathcal F$ est $Z$-clos.

Considérons la suite exacte
\[
0\longrightarrow\mathcal F\xrightarrow{f}\mathcal F
\longrightarrow\mathcal F/f\mathcal F\longrightarrow0
\]
(l'homothétie de rapport $f:\mathcal F\to\mathcal F$ étant par hypothèse
injective); si on pose $W=X-Z$, on a le diagramme commutatif
\[
\xymatrix{
0\ar[r] & \Gamma(X,\mathcal F)\ar[r]^f\ar[d] &
\Gamma(X,\mathcal F)\ar[r]\ar[d] &
\Gamma(X,\mathcal F/f\mathcal F)\ar[r]\ar[d] & 0\\
0\ar[r] & \Gamma(W,\mathcal F)\ar[r]_f &
\Gamma(W,\mathcal F)\ar[r] & \Gamma(W,\mathcal F/f\mathcal F)
}
\]
où les deux lignes sont exactes ($X$ étant affine). Si l'homomorphisme de
restriction $\Gamma(X,\mathcal F)\to\Gamma(W,\mathcal F)$ est bijectif, on
déduit de ce diagramme que
\[
\Gamma(X,\mathcal F/f\mathcal F)
\longrightarrow\Gamma(W,\mathcal F/f\mathcal F)
\]
est injectif, et cela montre (5.9.8) que si $\mathcal F$ est $Z$-clos,
$\mathcal F/f\mathcal F$ est $Z$-pur. Inversement, supposons que
$\mathcal F/f\mathcal F$ soit $Z$-pur, et soit $s$ une section de $\mathcal F$
au-dessus de $W$; comme $X_f\subset W$, il existe un entier $n>0$ tel que
$f^n(s|X_f)$ se prolonge en une section $t$ de $\mathcal F$ au-dessus de $X$
(I, 1.4.1); d'ailleurs, les restrictions de $t$ et de $f^ns$ à $X_f$ étant
les mêmes, il en résulte que la restriction de $t$ à $W$ est égale à $f^ns$
en vertu de la relation $\operatorname{Ass}(\mathcal F)\subset X_f$
(5.10.2); comme $f$ est $\mathcal F$-régulier, il suffira de voir que $t$ est
de la forme $f^nt'$, où $t'\in\Gamma(X,\mathcal F)$

\oldpage[IV-2]{116}

pour montrer que l'homomorphisme
$\Gamma(X,\mathcal F)\to\Gamma(W,\mathcal F)$ est surjectif, donc bijectif.
Or, dire que $t=f^nt'$ signifie que l'image de $t$ dans
$\Gamma(X,\mathcal F/f^n\mathcal F)$ est nulle. Mais comme
$f^k\mathcal F/f^{k+1}\mathcal F$ est isomorphe à
$\mathcal F/f\mathcal F$, donc $Z$-pur par hypothèse, on déduit de (5.10.3),
par récurrence sur $n$, que $\mathcal F/f^n\mathcal F$ est $Z$-pur. Mais par
définition l'image de $t|W=f^ns$ dans
$\Gamma(W,\mathcal F/f^n\mathcal F)$ est nulle, d'où la conclusion.

\begin{corollary}[5.10.6]
\label{IV.2.s5.10.6-fr}
Soit $\mathcal F$ un $\mathcal O_X$-Module cohérent. Pour que $\mathcal F$
soit $Z$-clos en un point $x$, il faut et il suffit que
$\operatorname{prof}_{Z_x}(\widetilde{\mathcal F}_x)\geq2$.
\end{corollary}

Cela résulte de (5.9.6) et (5.10.5).

\begin{theorem}[5.10.7]
\label{IV.2.s5.10.7-fr}
\emph{(Hartshorne).} Soient $X$ un préschéma localement noethérien, $Y$ une
partie fermée de $X$. Supposons que pour tout $y\in Y$, on ait
$\operatorname{prof}(\mathcal O_{X,y})\geq2$; alors pour toute composante
connexe $C$ de $X$, $C-(C\cap Y)$ est connexe.
\end{theorem}

On peut se borner au cas où $X$ est connexe; il résulte alors de (5.10.5) que
l'homomorphisme canonique
$\mathcal O_X\to i_*(\mathcal O_X|X-Y)$ (où $i:X-Y\to X$ est l'injection
canonique) est bijectif. Par suite l'homomorphisme de restriction
$\Gamma(X,\mathcal O_X)\to\Gamma(X-Y,\mathcal O_X)$ est aussi bijectif. Il
suffit maintenant d'appliquer le lemme (III, 7.8.6.1).

\begin{corollary}[5.10.8]
\label{IV.2.s5.10.8-fr}
Soient $X$ un préschéma localement noethérien, $d$ un entier tel que pour tout
$x\in X$, la relation $\dim(\mathcal O_x)\geq d$ entraîne
$\operatorname{prof}(\mathcal O_x)\geq2$. Supposons $X$ connexe; alors, si
$X'$, $X''$ sont deux composantes irréductibles distinctes de $X$, il existe
une suite $(X_i)_{0\leq i\leq n}$ de composantes irréductibles de $X$ telle
que $X_0=X'$, $X_n=X''$, et que, pour $1\leq i\leq n$, on ait
\[
\operatorname{codim}(X_{i-1}\cap X_i,X)\leq d-1
\]
(on dit alors que $X$ est connexe en codimension $\leq d-1$).
\end{corollary}

Si $Y$ est une partie fermée de $X$ telle que
$\operatorname{codim}(Y,X)\geq d$, on a $\dim(\mathcal O_{X,y})\geq d$ pour
tout $y\in Y$ (5.1.3), donc $\operatorname{prof}(\mathcal O_{X,y})\geq2$ pour
tout $y\in Y$, et il résulte de (5.10.7) que $X-Y$ est connexe. D'autre part,
pour que $\operatorname{codim}(Y,X)\geq d$, il faut et il suffit que pour tout
$y\in Y$, il existe un voisinage ouvert $V$ de $y$ dans $X$ tel que
$\operatorname{codim}(Y\cap V,V)\geq d$ (0, 14.2.3). Notons enfin que si
$\mathfrak F$ désigne l'ensemble des parties fermées $Y$ de $X$ de
codimension $\geq d$, la réunion de deux ensembles de $\mathfrak F$ appartient
à $\mathfrak F$ (0, 14.2.5), et tout ensemble fermé contenu dans un ensemble
de $\mathfrak F$ appartient à $\mathfrak F$, propriétés qu'on exprime encore
en disant que $\mathfrak F$ est un \emph{antifiltre} de parties fermées de
$X$. Le corollaire résulte alors du lemme topologique suivant:

\begin{lemma}[5.10.8.1]
\label{IV.2.s5.10.8.1-fr}
Soient $X$ un espace topologique connexe localement noethérien,
$\mathfrak F$ un antifiltre de parties fermées de $X$. On suppose que si $Y$
est une partie fermée de $X$ telle que pour tout $y\in Y$, il y ait un
voisinage ouvert $V$ de $y$ dans $X$ et un $Y_y\in\mathfrak F$ tels que
$V\cap Y=V\cap Y_y$, alors $Y\in\mathfrak F$. Les conditions suivantes sont
alors équivalentes:
\begin{enumerate}
\item[a)] Pour tout $Y\in\mathfrak F$, $X-Y$ est connexe.
\item[b)] Si $X'$ et $X''$ sont deux composantes irréductibles distinctes de
$X$, il existe une suite $(X_i)_{0\leq i\leq n}$ de composantes irréductibles
de $X$ telle que $X_0=X'$, $X_n=X''$ et que, pour $1\leq i\leq n$, on ait
$X_{i-1}\cap X_i\notin\mathfrak F$.
\end{enumerate}
\end{lemma}

Supposons b) vérifiée et prouvons que $U=X-Y$ est connexe pour tout
$Y\in\mathfrak F$. Si $U'$ et $U''$ sont deux composantes irréductibles
distinctes de $U$, il existe deux composantes irréductibles $X'$, $X''$ de
$X$ telles que $X'\cap U=U'$, $X''\cap U=U''$ (0$_\mathrm I$, 2.1.6);

\oldpage[IV-2]{117}

formons pour ces deux composantes une suite $(X_i)$ ayant la propriété énoncée
dans b) et posons $U_i=X_i\cap U$ pour $1\leq i\leq n$; alors $U_i$ est une
composante irréductible de $U$ (0$_\mathrm I$, 2.1.6) et de plus
$U_i\cap U_{i-1}\neq\varnothing$ pour $1\leq i\leq n$, sinon on aurait
$X_i\cap X_{i-1}\subset Y$, donc $X_i\cap X_{i-1}\in\mathfrak F$,
contrairement à la définition des $X_i$. Cela entraîne que $U$ est connexe.

Montrons maintenant que a) entraîne b). Désignons par $Y$ la réunion de la
famille $(X_\alpha\cap X_\beta)$, où $(X_\alpha,X_\beta)$ parcourt l'ensemble
des couples de composantes irréductibles distinctes de $X$ telles que
$X_\alpha\cap X_\beta\in\mathfrak F$. Pour tout point $y\in Y$, il y a un
voisinage ouvert $V$ de $y$ dans $X$ ne rencontrant qu'un nombre fini de
composantes irréductibles de $X$; cela montre d'une part que $Y$ est fermé et
d'autre part que $V\cap Y$ est intersection de $V$ et d'un ensemble de
$\mathfrak F$; en vertu de l'hypothèse faite sur $\mathfrak F$, on a
$Y\in\mathfrak F$, donc $U=X-Y$ est connexe; en outre $Y$ est rare dans $X$.
Soient alors $X'$, $X''$ deux composantes irréductibles distinctes de $X$,
$U'$, $U''$ leurs traces respectives sur $U$; ce sont des composantes
irréductibles distinctes de $U$ (0$_\mathrm I$, 2.1.6). Or la réunion des
composantes irréductibles $W$ de $U$ telles qu'il existe une suite
$(U_i)_{0\leq i\leq n}$ de composantes irréductibles de $U$ pour laquelle
$U_0=U'$, $U_{i-1}\neq U_i$ et $U_{i-1}\cap U_i\neq\varnothing$ pour
$1\leq i\leq n$ et $U_n=W$ est un ensemble ouvert et fermé dans $U$, puisque
$U$ est localement noethérien et par suite ses composantes irréductibles
forment une famille localement finie d'ensembles fermés. Il y a donc une telle
suite $(U_i)$ pour laquelle $U_n=U''$; soit $X_i$ ($0\leq i\leq n$) la
composante irréductible de $X$ telle que $X_i\cap U=U_i$
(0$_\mathrm I$, 2.1.6); comme $U_{i-1}\neq U_i$, on a
$X_{i-1}\neq X_i$ pour $1\leq i\leq n$; si l'on avait
$X_{i-1}\cap X_i\in\mathfrak F$ pour un $i$ tel que $1\leq i\leq n$, on en
déduirait $X_{i-1}\cap X_i\subset Y$ par définition de $Y$, d'où
$U_{i-1}\cap U_i=\varnothing$, contrairement à l'hypothèse. Ceci achève la
démonstration du lemme.

On notera que l'hypothèse faite sur $X$ dans (5.10.8) est vérifiée lorsque $X$
est un préschéma de Cohen-Macaulay et $d\geq2$.

\begin{corollary}[5.10.9]
\label{IV.2.s5.10.9-fr}
Si un anneau local noethérien $A$ vérifie $(S_2)$ et est caténaire, il est
équidimensionnel.
\end{corollary}

L'hypothèse de (5.10.8) est alors vérifiée par $X=\operatorname{Spec}(A)$,
avec $d=2$. Pour montrer que toutes les composantes irréductibles de $X$ ont
même dimension, il suffit alors, en vertu de (5.10.8), de montrer que deux
telles composantes $X'$, $X''$ ont même dimension lorsque l'on suppose en
outre que $\operatorname{codim}(X'\cap X'',X)=1$. Il y a alors une composante
irréductible $Z$ de $X'\cap X''$ telle que
$\operatorname{codim}(Z,X)=1$, donc $\operatorname{codim}(Z,X')=1$, puisque
$\operatorname{codim}(Z,X)\geq\operatorname{codim}(Z,X')\geq1$; de même
$\operatorname{codim}(Z,X'')=1$, et comme $X$ est caténaire, cela entraîne
$\dim(X')=\dim(X'')$.

\begin{proposition}[5.10.10]
\label{IV.2.s5.10.10-fr}
Soient $X$ un préschéma localement noethérien, $Z$ une partie de $X$ stable
par spécialisation, $\mathcal F$ un $\mathcal O_X$-Module cohérent et supposons
que le $\mathcal O_X$-Module $\mathcal H^0_{X/Z}(\mathcal F)$ soit cohérent.
Alors:
\begin{enumerate}
\item[(i)] On a
$\operatorname{prof}_Z(\mathcal H^0_{X/Z}(\mathcal F))\geq2$.
\item[(ii)] Pour tout point
$x\in\operatorname{Ass}(\mathcal F)\cap(X-Z)$, on a
\[
\operatorname{codim}(Z\cap\overline{\{x\}},\overline{\{x\}})\geq2.
\]
\item[(iii)] L'ensemble $U$ des $x\in X$ tels que
$\operatorname{prof}_{Z_x}(\widetilde{\mathcal F}_x)\geq2$ (notations de
(5.9.6)) est ouvert dans $X$; on a $X-U\subset Z$, et $U$ est le plus grand
ouvert de $X$ tel que $\mathcal F|U$ soit $(Z\cap U)$-clos.
\end{enumerate}
\end{proposition}

\oldpage[IV-2]{118}

Posons pour abréger $\mathcal F'=\mathcal H^0_{X/Z}(\mathcal F)$. On sait
(5.9.11) que $\mathcal F'$ est $Z$-clos et l'assertion (i) résulte donc de
(5.10.5) appliqué à $\mathcal F'$. Soit
$x\in\operatorname{Ass}(\mathcal F)\cap(X-Z)$; comme les restrictions de
$\mathcal F$ et $\mathcal F'$ à $X-Z$ sont canoniquement isomorphes (5.9.11),
on a $x\in\operatorname{Ass}(\mathcal F')$. Considérons un point
$y\in Z\cap\overline{\{x\}}$; l'idéal premier $\mathfrak p$ de
$\mathcal O_y$ correspondant à $x$ est associé au $\mathcal O_y$-module
$\mathcal F'_y$, donc on a, en vertu de (i) et de (0, 16.4.6.2),
\[
2\leq\operatorname{prof}(\mathcal F'_y)
\leq\dim(\mathcal O_y/\mathfrak p)
=\operatorname{codim}(\overline{\{y\}},\overline{\{x\}}),
\]
d'où (ii). Enfin, pour prouver (iii), notons que $U$ est l'ensemble des
$x\in X$ tels que $\widetilde{\mathcal F}_x$ soit $Z_x$-clos (5.10.5), ou
encore, en vertu de (5.9.6), l'ensemble des points où l'homomorphisme
canonique $\mathcal F_x\to\mathcal F'_x$ est bijectif; c'est donc le
complémentaire de la réunion des supports de $\operatorname{Ker}(\rho_{X/Z})$
et $\operatorname{Coker}(\rho_{X/Z})$, et ces derniers sont des
$\mathcal O_X$-Modules cohérents en vertu de l'hypothèse (0$_\mathrm I$,
5.3.4) donc ont un support fermé (0$_\mathrm I$, 5.2.2); cela montre que $U$
est ouvert, et $U$ est évidemment le plus grand ouvert tel que $\mathcal F|U$
soit $(Z\cap U)$-clos; enfin l'inclusion $X-U\subset Z$ résulte de (5.9.11).

Nous verrons plus loin (5.11.1) que dans les cas les plus importants
l'assertion (ii) entraîne réciproquement que
$\mathcal H^0_{X/Z}(\mathcal F)$ est cohérent.

\begin{env}[5.10.11]
\label{IV.2.s5.10.11-fr}
Soient $X$ un préschéma localement noethérien, $Z$ une partie de $X$ stable
par spécialisation; on a vu que
$\mathcal A=\mathcal H^0_{X/Z}(\mathcal O_X)$ est une
$\mathcal O_X$-Algèbre quasi-cohérente (5.9.3); le $X$-schéma
$X'=\operatorname{Spec}(\mathcal H^0_{X/Z}(\mathcal O_X))$ (II, 1.3.1) est
appelé la $Z$-clôture de $X$. En outre pour tout $\mathcal O_X$-Module
$\mathcal F$, $\mathcal H^0_{X/Z}(\mathcal F)$ est un $\mathcal A$-Module,
qui est quasi-cohérent si $\mathcal F$ est un $\mathcal O_X$-Module
quasi-cohérent; dans ce dernier cas, il y a donc un unique
$\mathcal O_{X'}$-Module $\mathcal F'$ tel que l'on ait, en désignant par
$g:X'\to X$ le morphisme structural (II, 1.4.3)
\begin{equation}
\mathcal H^0_{X/Z}(\mathcal F)=g_*(\mathcal F').
\tag{5.10.11.1}\label{IV.2.s5.10.11.1-fr}
\end{equation}
\end{env}

\begin{proposition}[5.10.12]
\label{IV.2.s5.10.12-fr}
Les notations étant celles de (5.10.11):
\begin{enumerate}
\item[(i)] Soit $x$ un point de $X$. Pour que le morphisme
\[
X'\times_X X_x\longrightarrow X_x(=\operatorname{Spec}(\mathcal O_{X,x}))
\]
déduit de $g$ par localisation soit un isomorphisme, il faut et il suffit que
$\mathcal O_X$ soit $Z$-clos au point $x$ (ce qui a lieu pour tout $x\in X-Z$).
\item[(ii)] Posons $Z'=g^{-1}(Z)$, et supposons $X'$ localement noethérien.
Alors $\mathcal F'$ est $Z'$-clos; si de plus $\mathcal F'$ est un
$\mathcal O_{X'}$-Module cohérent, on a
$\operatorname{prof}_{Z'}(\mathcal F')\geq2$.
\item[(iii)] Supposons que $\mathcal H^0_{X/Z}(\mathcal O_X)$ et
$\mathcal H^0_{X/Z}(\mathcal F)$ soient cohérents. Alors le morphisme
$g:X'\to X$ est fini; l'ensemble $U$ des $x\in X$ tels que l'on ait
$\operatorname{prof}_{Z_x}(\widetilde{\mathcal O}_x)\geq2$ et
$\operatorname{prof}_{Z_x}(\widetilde{\mathcal F}_x)\geq2$ est ouvert dans
$X$ et tel que $X-U\subset Z$; en outre $U$ est le plus grand ensemble ouvert
de $X$ tel que la restriction $g^{-1}(U)\to U$ de $g$ soit un isomorphisme et
que la restriction $\mathcal F|U\to\mathcal F'|g^{-1}(U)$ du $g$-morphisme
canonique $\mathcal F\to\mathcal F'$ soit un isomorphisme.
\end{enumerate}
\end{proposition}

L'assertion (i) résulte des définitions, et (iii) est une conséquence immédiate
de (5.10.10, (iii)). Pour prouver (ii), considérons un ouvert $V$ de $X$
contenant $Z-X$ et son image réciproque $V'=g^{-1}(V)$; si $i:V\to X$ et
$i':V'\to X'$ sont les injections

\oldpage[IV-2]{119}

canoniques, l'homomorphisme canonique
$\rho_{X'/Z'}:\mathcal F'\to i'_*(\mathcal F'|V')$ est tel que
$g_*(\rho_{X'/Z'})$ soit l'homomorphisme canonique
\[
\rho_{X/Z}:\mathcal H^0_{X/Z}(\mathcal F)
\longrightarrow i_*(\mathcal H^0_{X/Z}(\mathcal F)|V)
\]
(II, 1.4.2), compte tenu de (5.10.11.1). Comme
$\mathcal H^0_{X/Z}(\mathcal F)$ est $Z$-clos (5.9.11), $\rho_{X/Z}$ est un
isomorphisme, donc il en est de même de $\rho_{X'/Z'}$. Comme $X'-Z'$ est
intersection de la famille filtrante des $V'_\alpha=g^{-1}(V_\alpha)$, où
$V_\alpha$ parcourt la famille filtrante des ouverts contenant $X-Z$, on en
déduit que $\mathcal F'$ est $Z'$-clos lorsque $X'$ est localement noethérien,
en vertu de (5.9.1).

\begin{env}[5.10.13]
\label{IV.2.s5.10.13-fr}
Nous allons maintenant appliquer les résultats qui précèdent au cas où $Z$ est
un des ensembles $Z^{(n)}(X)$ (ou simplement $Z^{(n)}$), défini comme
l'ensemble des $x\in X$ tels que $\dim(\mathcal O_x)\geq n$; il est clair que
$Z^{(n)}$ est stable par spécialisation; pour qu'une partie fermée $T$ de $X$
soit contenue dans $Z^{(n)}$, il faut et il suffit que
$\operatorname{codim}(T,X)\geq n$. Nous nous intéresserons ici au cas $n=2$.
\end{env}

\begin{proposition}[5.10.14]
\label{IV.2.s5.10.14-fr}
Soient $X$ un préschéma localement noethérien, $\mathcal F$ un
$\mathcal O_X$-Module cohérent de support égal à $X$.
\begin{enumerate}
\item[(i)] Pour que $\mathcal F$ vérifie $(S_1)$, il faut et il suffit qu'il
soit $Z^{(1)}$-pur.
\item[(ii)] Pour que $\mathcal F$ vérifie $(S_2)$, il faut et il suffit qu'il
soit $Z^{(2)}$-clos et $Z^{(1)}$-pur, ou encore qu'il soit $Z^{(2)}$-clos et
n'ait pas de cycle premier associé de codimension 1.
\end{enumerate}
\end{proposition}

(i) Dire que $\mathcal F$ possède la propriété $(S_1)$ signifie que
$\mathcal F$ n'a pas de cycle premier associé immergé (5.7.5), ou encore que
pour tout $x\in\operatorname{Ass}(\mathcal F)$, on a
$\dim(\mathcal F_x)=0$ (3.1.4), autrement dit (5.1.12.1)
$\dim(\mathcal O_x)=0$; mais cela équivaut à dire que
$\operatorname{Ass}(\mathcal F)$ ne rencontre pas $Z^{(1)}$, et la conclusion
résulte de (5.10.2).

(ii) Dire que $\mathcal F$ est $Z^{(2)}$-clos signifie que
$\operatorname{prof}_{Z^{(2)}}(\mathcal F)\geq2$, ou encore que, pour tout
$x\in X$, la relation $\dim(\mathcal O_x)\geq2$ entraîne
$\operatorname{prof}(\mathcal F_x)\geq2$; cela montre que la propriété $(S_2)$
entraîne que $\mathcal F$ est $Z^{(2)}$-clos; elle entraîne en outre que
$\mathcal F$ vérifie $(S_1)$, donc n'a pas de cycle premier associé immergé
(5.7.5), et comme $\operatorname{Supp}(\mathcal F)=X$, cela signifie encore
que tous les cycles premiers associés à $\mathcal F$ sont de codimension 0.
Inversement, supposons que $\mathcal F$ soit $Z^{(2)}$-clos et n'ait pas de
cycle premier associé de codimension 1; pour voir que $\mathcal F$ vérifie
$(S_2)$, il reste à montrer que si $x\in X$ est tel que
$\dim(\mathcal F_x)=1$ (ou, ce qui revient au même,
$\dim(\mathcal O_x)=1$), alors on a $\operatorname{prof}(\mathcal F_x)=1$;
mais par hypothèse la relation $\dim(\mathcal O_x)=1$ entraîne
$x\notin\operatorname{Ass}(\mathcal F)$, et cette dernière relation équivaut à
$\operatorname{prof}(\mathcal F_x)\neq0$, c'est-à-dire ici à
$\operatorname{prof}(\mathcal F_x)=1$. Si $\mathcal F$ est $Z^{(1)}$-pur,
donc vérifie $(S_1)$, on a remarqué plus haut qu'en vertu de la relation
$\operatorname{Supp}(\mathcal F)=X$, tous les cycles premiers associés à
$\mathcal F$ sont de codimension 0, donc ce qui précède s'applique.

On notera qu'il peut se faire que $\mathcal F$ soit $Z^{(2)}$-clos et ne
vérifie pas $(S_1)$: c'est le cas par exemple lorsque $X$ est de dimension 1
(car alors $Z^{(2)}=\varnothing$, et tout $\mathcal O_X$-Module est
$Z^{(2)}$-clos) et a des cycles premiers associés immergés.

Rappelons qu'au chapitre III, dans l'étude de la cohomologie locale, on donne
une caractérisation cohomologique de la propriété $(S_n)$ pour tout $n\geq1$,
généralisant (5.10.14).

\begin{corollary}[5.10.15]
\label{IV.2.s5.10.15-fr}
Soient $X$ un préschéma localement noethérien, $\mathcal F$ un
$\mathcal O_X$-Module cohérent de support $X$. On suppose que $\mathcal F$
n'a pas de cycles premiers associés de codimension 1 et que
$\mathcal F'=\mathcal H^0_{X/Z^{(2)}}(\mathcal F)$ soit cohérent. Alors:
\begin{enumerate}
\item[(i)] $\mathcal F'$ vérifie la propriété $(S_2)$.

\oldpage[IV-2]{120}

\item[(ii)] L'ensemble $U$ des $x\in X$ tels que $\mathcal F$ vérifie $(S_2)$
au point $x$ (5.7.2) est ouvert dans $X$ et on a
$\operatorname{codim}(X-U,X)\geq2$.
\end{enumerate}
\end{corollary}

(i) On sait (5.9.11) que $\mathcal F'$ est $Z^{(2)}$-clos, et en outre
$\operatorname{Supp}(\mathcal F')=X$, car les points maximaux de $X$
appartiennent à $X-Z^{(2)}$ et en ces points
$\mathcal F'_x=\mathcal F_x\neq0$, donc le support de $\mathcal F'$ est dense
dans $X$, et comme $\mathcal F'$ est cohérent,
$\operatorname{Supp}(\mathcal F')$ est fermé, donc égal à $X$. Il reste à voir
que $\mathcal F'$ n'a pas de cycles premiers associés de codimension 1. Mais
si $x\in\operatorname{Ass}(\mathcal F')$ et
$\dim(\mathcal F'_x)=\dim(\mathcal O_x)=1$, on a $x\in X-Z^{(2)}$, donc,
puisque $\mathcal F'_x=\mathcal F_x$, on aurait
$x\in\operatorname{Ass}(\mathcal F)$, contrairement à l'hypothèse, ce qui
achève de prouver (i).

(ii) On a $Z^{(n)}(X_x)=Z^{(n)}(X)\cap X_x$, avec les notations de (5.9.6),
compte tenu de (I, 2.4.2); d'autre part, l'hypothèse que $\mathcal F$ n'a pas
de cycles premiers associés de codimension 1 entraîne la même hypothèse pour
$\widetilde{\mathcal F}_x$; pour que $\widetilde{\mathcal F}_x$ vérifie
$(S_2)$, il faut et il suffit donc, en vertu de (5.10.14), que
$\widetilde{\mathcal F}_x$ soit $Z^{(2)}(X_x)$-clos; l'assertion (ii) résulte
donc de (5.10.6) et de (5.10.10, (iii)).

\begin{proposition}[5.10.16]
\label{IV.2.s5.10.16-fr}
Soient $X$ un préschéma localement noethérien,
\[
X'=\operatorname{Spec}(\mathcal H^0_{X/Z^{(2)}}(\mathcal O_X))
\]
sa $Z^{(2)}$-clôture, $g:X'\to X$ le morphisme structural. Supposons que $X$
n'ait pas de cycle premier associé de codimension 1.
\begin{enumerate}
\item[(i)] Pour qu'en un point $x\in X$, $X$ vérifie $(S_2)$, il faut et il
suffit que le morphisme $X'_x\to X_x$ déduit de $g$ (notations de (5.10.12))
soit un isomorphisme. Cette condition est toujours vérifiée si
$\operatorname{codim}(\overline{\{x\}},X)\leq1$.
\item[(ii)] Supposons de plus que $g$ soit un morphisme fini (voir dans
(5.11.2) des conditions suffisantes pour qu'il en soit ainsi). Alors
l'ensemble $U$ des points où $X$ vérifie $(S_2)$ est ouvert et
$\operatorname{codim}(X-U,X)\geq2$; en outre $U$ est le plus grand ouvert de
$X$ tel que la restriction $g^{-1}(U)\to U$ de $g$ soit un isomorphisme.
\item[(iii)] Sous les mêmes hypothèses que dans (ii), $X'$ satisfait à $(S_2)$
et pour tout $x'\in X'$ tel que\linebreak
$\operatorname{codim}(\overline{\{x'\}},X')\leq1$, le point $x=g(x')$ est tel
que
\[
\operatorname{codim}(\overline{\{x\}},X)
=\operatorname{codim}(\overline{\{x'\}},X').
\]
\item[(iv)] Les hypothèses étant celles de (ii), soit $\mathcal F$ un
$\mathcal O_X$-Module cohérent de support $X$, tel que
$\mathcal H^0_{X/Z^{(2)}}(\mathcal F)$ soit cohérent; alors le
$\mathcal O_{X'}$-Module $\mathcal F'$ tel que
$g_*(\mathcal F')=\mathcal H^0_{X/Z^{(2)}}(\mathcal F)$ est cohérent et
vérifie $(S_2)$, et son support est une réunion de composantes irréductibles de
$X'$.
\end{enumerate}
\end{proposition}

Les assertions (i) et (ii) sont insérées pour mémoire, ayant déjà été
démontrées en substance: (i) résulte en effet de (5.10.12, (i)) et de
(5.10.14), et (ii) est un cas particulier de (5.10.15, (ii)).

Démontrons (iii); posons $x=g(x')$; comme $g$ est fini, il en est de même du
morphisme $X'_x\to X_x$, donc
\[
\dim(\mathcal O_{X',x'})\leq\dim(X'_x)\leq\dim(X_x)
=\dim(\mathcal O_{X,x})
\]
en vertu de (5.4.1). Supposons d'abord que
$\dim(\mathcal O_{X',x'})\leq1$ et montrons qu'alors
$\dim(\mathcal O_{X,x})\leq1$. Sinon, on aurait $x\in Z^{(2)}$, donc en vertu
de (5.10.12, (ii)) appliqué à $\mathcal O_X$ on aurait
$\operatorname{prof}(\mathcal O_{X',x'})\geq2$, ce qui est absurde. On a donc
$x\in X-Z^{(2)}$, et par suite $\mathcal O_{X',x'}$ est isomorphe à
$\mathcal O_{X,x}$ (5.10.12, (i)), d'où
$\dim(\mathcal O_{X',x'})=\dim(\mathcal O_{X,x})$. En outre, comme $X$ n'a
pas de cycles

\oldpage[IV-2]{121}

premiers associés de codimension 1, l'hypothèse
$\dim(\mathcal O_{X,x})=1$ entraîne $x\notin\operatorname{Ass}(\mathcal O_X)$,
donc $\operatorname{prof}(\mathcal O_{X,x})=1$, et par suite aussi
$\operatorname{prof}(\mathcal O_{X',x'})=1$. Supposons maintenant
$\dim(\mathcal O_{X',x'})\geq2$, donc $\dim(\mathcal O_{X,x})\geq2$,
c'est-à-dire $x\in Z^{(2)}$; on en déduit que
$\operatorname{prof}(\mathcal O_{X',x'})\geq2$ par (5.10.12, (ii)). Ceci
établit les assertions de (iii).

Pour prouver (iv) il suffit de remplacer $\mathcal O_X$ par $\mathcal F$ dans
le raisonnement précédent, qui établit que $\mathcal F'$ vérifie $(S_2)$ et
que si $\dim(\mathcal F'_{x'})\leq1$, $\mathcal F_x$ et
$\mathcal F'_{x'}$ sont di-isomorphes; en particulier si
$\dim(\mathcal F'_{x'})=0$, on a $\dim(\mathcal F_x)=0$, donc
$\dim(\mathcal O_{X,x})=0$ puisque $\mathcal F$ a pour support $X$, et enfin
$\dim(\mathcal O_{X',x'})=0$; toute composante irréductible de
$\operatorname{Supp}(\mathcal F')$ est donc une composante irréductible de
$X'$, puisque $\mathcal F'$ est cohérent, donc
$\operatorname{Supp}(\mathcal F')$ fermé.

\begin{proposition}[5.10.17]
\label{IV.2.s5.10.17-fr}
Soient $A$ un anneau intègre noethérien, et désignons par $A^{(1)}$
l'intersection des anneaux locaux $A_{\mathfrak p}$, où $\mathfrak p$
parcourt l'ensemble des idéaux premiers de $A$ de hauteur 1. Supposons que
$A^{(1)}$ soit une $A$-algèbre finie. Alors:
\begin{enumerate}
\item[(i)] L'anneau $A^{(1)}$ vérifie la condition $(S_2)$.
\item[(ii)] L'ensemble $U$ des $\mathfrak p\in\operatorname{Spec}(A)$ tels
que l'homomorphisme canonique
$A_{\mathfrak p}\to(A^{(1)})_{\mathfrak p}$ soit bijectif est égal à
l'ensemble des $\mathfrak p$ tels que $A_{\mathfrak p}$ vérifie $(S_2)$; $U$
est ouvert dans $X=\operatorname{Spec}(A)$ et l'on a
$\operatorname{codim}(X-U,X)\geq2$.
\item[(iii)] Pour toute partie multiplicative $S$ de $A$, $(S^{-1}A)^{(1)}$
est une $(S^{-1}A)$-algèbre finie.
\item[(iv)] Soit $B$ une $A$-algèbre finie, intègre et contenant $A$. Alors
$B^{(1)}$ est une $B$-algèbre finie. En outre, pour tout idéal premier
$\mathfrak q$ de $B$, de hauteur 1, l'idéal premier $\mathfrak q\cap A$ de
$A$ est de hauteur 1.
\end{enumerate}
\end{proposition}

Si l'on tient compte de la formule (5.9.3.1), on voit que
$X'=\operatorname{Spec}(A^{(1)})$ est la $Z^{(2)}$-clôture de
$X=\operatorname{Spec}(A)$; comme $A$ n'a pas d'idéaux premiers associés
immergés, les propriétés (i) et (ii) sont des cas particuliers de (5.10.16,
(i), (ii) et (iii)). Pour prouver (iii), il suffit de remarquer que l'on a
$(S^{-1}A)^{(1)}=S^{-1}A^{(1)}$, ce qui est un cas particulier de (5.9.4): en
effet $S^{-1}A$ est un $A$-module plat, les idéaux premiers de $S^{-1}A$ sont
les idéaux $S^{-1}\mathfrak p$, où $\mathfrak p\in\operatorname{Spec}(A)$ ne
rencontre pas $S$, et l'on a
$\operatorname{ht}(S^{-1}\mathfrak p)=\operatorname{ht}(\mathfrak p)$. Comme
$A^{(1)}$ est une $A$-algèbre finie, $S^{-1}A^{(1)}$ est une
$S^{-1}A$-algèbre finie, d'où (iii).

Pour démontrer (iv), posons $Y=\operatorname{Spec}(B)$, et soit $f:Y\to X$ le
morphisme structural; comme il est fini, il résulte de (5.4.1) que pour tout
$y\in Y$, on a
$\dim(\mathcal O_{Y,y})\leq\dim(\mathcal O_{X,f(y)})$; donc, si
$T=f^{-1}(Z^{(2)}(X))$, on a $T\supset Z^{(2)}(Y)$. Montrons que
\[
\mathcal G=\mathcal H^0_{X/Z^{(2)}(X)}(f_*(\mathcal O_Y))
\]
est cohérent; en effet $f_*(\mathcal O_Y)=\widetilde B$, $B$ étant considéré
comme $A$-module; mais comme $B$ est une $A$-algèbre intègre finie, son corps
des fractions est fini sur le corps des fractions de $A$, donc $B$ est contenu
dans un $A$-module libre de type fini, et par suite (5.9.2, (i)) $\mathcal G$
est un sous-$\mathcal O_X$-Module quasi-cohérent de
$(\mathcal H^0_{X/Z^{(2)}(X)}(\mathcal O_X))^n$ pour un $n$ convenable; ce
dernier étant cohérent par hypothèse, il en est de même de $\mathcal G$. Or,
il résulte de la définition (5.9.1.2) que $\mathcal G$ est isomorphe à
$f_*(\mathcal H^0_{Y/T}(\mathcal O_Y))$; cela prouve a fortiori que
$\mathcal H^0_{Y/T}(\mathcal O_Y)$ est un $\mathcal O_Y$-Module cohérent. Il
résulte alors de (5.10.10, (ii)), appliqué à $\mathcal O_Y$ et au point
générique de $Y$, que l'on a $\operatorname{codim}(T,Y)\geq2$, c'est-à-dire
$T\subset Z^{(2)}(Y)$, et finalement $T=Z^{(2)}(Y)$. Cela prouve les deux
assertions de (iv).

\oldpage[IV-2]{122}

\subsection{\texorpdfstring{Critères de cohérence pour les Modules
$\mathcal H^0_{X/Z}(\mathcal F)$}{Critères de cohérence pour les Modules H\string^0\string_\string{X/Z\string}(F)}}

\begin{proposition}[5.11.1]
\label{IV.2.s5.11.1-fr}
Soient $X$ un préschéma localement noethérien, $Z$ une partie de $X$ stable par
spécialisation, $\mathcal F$ un $\mathcal O_X$-Module cohérent. Désignons par
$(x_\alpha)$ la famille des points de
$\operatorname{Ass}(\mathcal F)\cap(X-Z)$ et, pour tout $\alpha$, soient
$Y_\alpha$ le sous-préschéma fermé réduit de $X$ ayant
$\overline{\{x_\alpha\}}$ pour espace sous-jacent,
$Z_\alpha=Z\cap\overline{\{x_\alpha\}}$. Alors les deux conditions suivantes
sont équivalentes:
\begin{enumerate}
\item[a)] $\mathcal H^0_{X/Z}(\mathcal F)$ est un $\mathcal O_X$-Module
cohérent.
\item[b)] Pour tout $\alpha$,
$\mathcal H^0_{Y_\alpha/Z_\alpha}(\mathcal O_{Y_\alpha})$ est un
$\mathcal O_{Y_\alpha}$-Module cohérent.
\end{enumerate}
Ces deux conditions entraînent la suivante:
\begin{enumerate}
\item[c)] Pour tout $\alpha$, on a
$\operatorname{codim}(Z_\alpha,Y_\alpha)\geq2$.
\end{enumerate}
De plus, les trois conditions a), b), c) sont équivalentes lorsqu'en outre une
des propriétés suivantes est vérifiée:
\begin{enumerate}
\item[(i)] Tout point de $X$ admet un voisinage ouvert isomorphe à un
sous-schéma d'un schéma régulier (auquel cas on dit aussi que $X$ est
localement immersible dans un schéma régulier).
\item[(ii)] Pour tout $\alpha$, $Y_\alpha$ est universellement caténaire
(5.6.2) et son normalisé (II, 6.3.8) $Y'_\alpha$ est fini sur $Y_\alpha$.
\end{enumerate}
\end{proposition}

Toutes les propriétés envisagées sont locales sur $X$, donc on peut se borner
au cas où $X=\operatorname{Spec}(A)$ est affine, $A$ étant un anneau
noethérien, et $\mathcal F=\widetilde M$, où $M$ est un $A$-module de type
fini. Alors, pour tout $\alpha$, si $h_\alpha$ est l'injection canonique
$Y_\alpha\to X$, $(h_\alpha)_*(\mathcal O_{Y_\alpha})=\mathcal G_\alpha$ est
le $\mathcal O_X$-Module correspondant au $A$-module quotient
$A/\mathfrak j_{x_\alpha}$, et, par définition de
$\operatorname{Ass}(\mathcal F)$, ce $A$-module est isomorphe à un
sous-$A$-module de $M$. Comme $\mathcal H^0_{X/Z}(\mathcal G_\alpha)$ est un
sous-$\mathcal O_X$-Module quasi-cohérent de
$\mathcal H^0_{X/Z}(\mathcal F)$ (5.9.2), l'hypothèse que
$\mathcal H^0_{X/Z}(\mathcal F)$ est cohérent entraîne qu'il en est de même de
$\mathcal H^0_{X/Z}(\mathcal G_\alpha)$. D'autre part, il résulte de la
définition (5.9.1.2) que $\mathcal H^0_{X/Z}(\mathcal G_\alpha)$ est
isomorphe à
$(h_\alpha)_*(\mathcal H^0_{Y_\alpha/Z_\alpha}(\mathcal O_{Y_\alpha}))$;
cela prouve que a) entraîne b). Pour voir que b) entraîne a), il suffit de
montrer qu'il y a une filtration finie $(\mathcal F_i)_{0\leq i\leq n}$ de
$\mathcal F$ formée de $\mathcal O_X$-Modules cohérents, avec
$\mathcal F_0=\mathcal F$, $\mathcal F_n=0$, et telle que
$\mathcal H^0_{X/Z}(\mathcal F_i/\mathcal F_{i+1})$ soit cohérent pour tout
$i$: cela résulte, par récurrence descendante sur $i$, des suites exactes
\[
0\longrightarrow\mathcal F_{i+1}\longrightarrow\mathcal F_i
\longrightarrow\mathcal F_i/\mathcal F_{i+1},
\]
du fait que $\mathcal H^0_{X/Z}$ est un foncteur exact à gauche (5.9.2), et
enfin de (0$_\mathrm I$, 5.3.3) et (I, 6.1.1). En vertu de (3.2.8), il suffit
donc de prouver que $\mathcal H^0_{X/Z}(\mathcal F)$ est cohérent lorsque
$\mathcal F$ est irredondant; autrement dit,
$\operatorname{Ass}(M)=\{\mathfrak p\}$ est réduit à un seul élément. Notons
maintenant le

\begin{lemma}[5.11.1.1]
\label{IV.2.s5.11.1.1-fr}
Soient $A$ un anneau noethérien, $M$ un $A$-module de type fini, tel que
$\operatorname{Ass}(M)=\{\mathfrak p\}$. Il existe une filtration finie
$(M_h)_{0\leq h\leq m}$ de $M$ telle que $M_0=M$, $M_m=0$ et que
$M_h/M_{h+1}$ soit isomorphe à un sous-module de $A/\mathfrak p$.
\end{lemma}

Notons d'abord que l'homomorphisme canonique $M\to M_\mathfrak p=N$ est
injectif (Bourbaki, \emph{Alg. comm.}, chap. IV, \S~1, n\textsuperscript{o}~2,
prop. 6). Posons $B=A_\mathfrak p$, $\mathfrak m=\mathfrak pA_\mathfrak p$,
idéal maximal de $B$; on a $\operatorname{Ass}(N)=\{\mathfrak m\}$
(\emph{loc. cit.}, prop. 5), et comme $N$ est un $B$-module de type fini,

\oldpage[IV-2]{123}

il existe un entier $r$ tel que $\mathfrak m^rN=0$; si on pose
$N'_j=\mathfrak m^jN$ pour $0\leq j\leq r$, $N'_j/N'_{j+1}$ est un
$(B/\mathfrak m)$-module de type fini, donc somme directe d'un nombre fini de
sous-modules isomorphes à $B/\mathfrak m$, puisque $B/\mathfrak m$ est un
corps; autrement dit, il y a une filtration finie $(N_h)_{0\leq h\leq m}$ de
$N$ telle que $N_m=0$ et que $N_h/N_{h+1}$ soit isomorphe à
$B/\mathfrak m$, i.e. au corps des fractions de $A/\mathfrak p$; la
filtration des $M_h=M\cap N_h$ répond à la question, car $M_h/M_{h+1}$ est
isomorphe à un sous-$(A/\mathfrak p)$-module de type fini de
$N_h/N_{h+1}=B/\mathfrak m$; mais on sait qu'un tel sous-module est isomorphe
à un sous-module de $A/\mathfrak p$.

L'existence de la filtration $(M_h)$ montre alors, par le même raisonnement que
plus haut, que pour prouver (moyennant b)) que
$\mathcal H^0_{X/Z}(\mathcal F)$ est cohérent, on peut se borner à montrer que
$\mathcal H^0_{X/Z}(\mathcal G)$ est cohérent, où
$\mathcal G=(A/\mathfrak p)^\sim$, $\mathfrak p$ étant un idéal associé à
$M$. Mais si $\mathfrak p=\mathfrak j_y$ avec $y\in Z$, le support de
$\mathcal G$ est un ensemble fermé contenu dans $Z$, puisque $Z$ est stable
par spécialisation; la définition (5.9.1.2) montre alors que
$\mathcal H^0_{X/Z}(\mathcal G)=0$. Si au contraire
$\mathfrak p=\mathfrak j_{x_\alpha}$ pour un $\alpha$, on a
$\mathcal G=(h_\alpha)_*(\mathcal O_{Y_\alpha})$ par définition, et l'on a vu
plus haut que $\mathcal H^0_{X/Z}(\mathcal G)$ est isomorphe à
$(h_\alpha)_*(\mathcal H^0_{Y_\alpha/Z_\alpha}(\mathcal O_{Y_\alpha}))$,
donc est cohérent en vertu de l'hypothèse b).

On a déjà vu (5.10.10, (ii)) que a) entraîne c). Il reste à prouver que, sous
l'une ou l'autre des hypothèses (i), (ii), c) entraîne b). Notons que si $X$
vérifie (i), il en est de même de chaque $Y_\alpha$. Il suffit donc (5.9.3.1)
de prouver le

\begin{corollary}[5.11.2]
\EGAReaderTarget{ega4.c04.cor.5-11-2}{IV.2.s5.11.2-fr}
Soient $A$ un anneau noethérien intègre, vérifiant l'une des deux hypothèses
suivantes:
\begin{enumerate}
\item[(i)] $A$ est quotient d'un anneau noethérien régulier.
\item[(ii)] $A$ est universellement caténaire, et sa clôture intégrale $A'$ est
une $A$-algèbre finie.
\end{enumerate}
Alors l'anneau $A^{(1)}$, intersection des $A_\mathfrak p$ où $\mathfrak p$
parcourt l'ensemble des idéaux premiers de $A$ de hauteur 1, est une
$A$-algèbre finie.
\end{corollary}

(i) Posons $X=\operatorname{Spec}(A)$. L'ensemble $U$ des points $x\in X$ où
$X$ vérifie $(S_2)$ est ouvert dans $X$ moyennant l'hypothèse (i) (6.11.2)
\footnote{Le lecteur pourra vérifier que (5.11.2) n'est pas utilisé dans la
démonstration de (6.11.2).}. En outre, si l'on pose $Z=X-U$, on a
$\operatorname{codim}(Z,X)\geq2$: en effet, pour tout $x\in X$ tel que
$\dim(A_x)\leq1$, on a $\dim(A_{x'})\leq1$ pour tout $x'$ générisation de
$x$, et comme $A_{x'}$ est intègre, $\operatorname{prof}(A_{x'})\geq1$, donc
$X$ vérifie $(S_2)$ au point $x$. On a donc $Z\subset Z^{(2)}$, et
$\mathcal H^0_{X/Z^{(2)}}(\mathcal O_X)
=j_*(\mathcal H^0_{U/Z^{(2)}}(\mathcal O_U))$, où $j:U\to X$ est l'injection
canonique (5.9.1.2). Mais comme le préschéma $U$ vérifie $(S_2)$, il résulte de
(5.10.14) que $\mathcal H^0_{U/Z^{(2)}}(\mathcal O_U)$ est isomorphe à
$\mathcal O_U$; d'autre part, comme $\operatorname{codim}(Z,X)\geq2$, on sait,
d'après le chapitre III, \S~9, que $j_*(\mathcal O_U)$ est un
$\mathcal O_X$-Module cohérent, ce qui démontre la proposition dans le cas
(i).

(ii) L'anneau $A'$ est, en vertu de l'hypothèse (ii), un anneau noethérien
intègre et intégralement clos, donc (Bourbaki, \emph{Alg. comm.}, chap. VII,
\S~1, n\textsuperscript{o}~6, th. 4) intersection de ses anneaux locaux
$A'_{\mathfrak p'}$, où $\mathfrak p'$ parcourt l'ensemble des idéaux premiers
de hauteur 1 de $A'$. Or, pour un tel idéal premier $\mathfrak p'$, si l'on
pose $\mathfrak p=\mathfrak p'\cap A$ et $S=A-\mathfrak p$,
$A'_{\mathfrak p'}$ est un anneau local en l'idéal premier
$S^{-1}\mathfrak p'$ de $S^{-1}A'$, et $S^{-1}A'$ est par hypothèse

\oldpage[IV-2]{124}

une $A_\mathfrak p$-algèbre finie; comme $S^{-1}\mathfrak p'$ est au-dessus
de l'idéal maximal $\mathfrak pA_\mathfrak p$ de $A_\mathfrak p$, c'est un
idéal maximal de $S^{-1}A'$; mais en vertu de l'hypothèse (ii) et de (5.6.3,
(i)), $A_\mathfrak p$ est universellement caténaire. On déduit donc de
(5.6.10) que
$\dim(A_\mathfrak p)=\dim(A'_{\mathfrak p'})=1$. Il résulte de là que l'on a
$A^{(1)}\subset A'$, et comme $A'$ est un $A$-module de type fini par
hypothèse, il en est de même de $A^{(1)}$ puisque $A$ est noethérien; la
proposition est donc démontrée dans le cas (ii).

\begin{remark}[5.11.3]
\label{IV.2.s5.11.3-fr}
Le fait que $A^{(1)}$ est une $A$-algèbre finie n'est plus nécessairement exact
si, dans l'hypothèse (ii) de (5.11.2), on suppose seulement que $A$ est
caténaire. Un exemple est donné par l'anneau local caténaire $A$ construit dans
(5.6.11), dont la clôture intégrale $A'$ (notée $B$ dans (5.6.11)) est une
$A$-algèbre finie; si $A^{(1)}$ était aussi une $A$-algèbre finie, comme il est
contenu dans le corps des fractions de $A$, il serait contenu dans $A'$. Mais
d'autre part, avec les notations de (5.6.11), tout idéal premier de hauteur 1
dans $A$ est de la forme $\mathfrak pA$, où $\mathfrak p$ est un idéal premier
de hauteur 1 dans $C$, et l'on a
$A_{\mathfrak pA}=C_\mathfrak p$; on sait (5.6.11) que
$C_\mathfrak p=E_{\mathfrak p'}$, où $\mathfrak p'$ est l'unique idéal
premier de $E$ au-dessus de $\mathfrak p$, donc $A_{\mathfrak pA}$ est
intégralement clos et contient par suite $A'$; par définition, on a donc
$A'\subset A^{(1)}$, et l'hypothèse que $A^{(1)}$ est une $A$-algèbre finie
entraînerait finalement $A^{(1)}=A'$. Mais cette conclusion est absurde, car un
des deux idéaux premiers de $A'$ au-dessus de l'idéal maximal $\mathfrak nA$ de
$A$ est de hauteur 1, alors que $\mathfrak nA$ est de hauteur 2, ce qui
contredirait (5.10.17, (iv)).
\end{remark}

\begin{corollary}[5.11.4]
\EGAReaderTarget{ega4.c04.cor.5-11-4}{IV.2.s5.11.4-fr}
Soient $X$ un préschéma localement noethérien, $Z$ une partie fermée de $X$,
$U=X-Z$, $i:U\to X$ l'injection canonique, $\mathcal F$ un
$\mathcal O_U$-Module cohérent. Pour que $i_*(\mathcal F)$ soit un
$\mathcal O_X$-Module cohérent, il est nécessaire que pour tout
$x\in\operatorname{Ass}(\mathcal F)$, on ait
\[
\operatorname{codim}(\overline{\{x\}}\cap Z,\overline{\{x\}})\geq2.
\]
Cette condition est suffisante dans chacun des deux cas suivants:
\begin{enumerate}
\item[(i)] Le préschéma $X$ est localement immersible dans un schéma régulier.
\item[(ii)] Le préschéma $X$ est universellement caténaire et universellement
japonais (ce qui signifie, par définition, que tout point $x\in X$ admet un
voisinage ouvert affine dont l'anneau est universellement japonais (0,
23.1.1)).
\end{enumerate}
\end{corollary}

On sait (I, 9.4.7) qu'il existe un sous-$\mathcal O_X$-Module cohérent
$\mathcal G$ de $i_*(\mathcal F)$ tel que $\mathcal G\vert U=\mathcal F$. On a
évidemment
$\operatorname{Ass}(\mathcal F)\subset\operatorname{Ass}(\mathcal G)
\subset\operatorname{Ass}(i_*(\mathcal F))$, et comme
\[
\operatorname{Ass}(i_*(\mathcal F))=\operatorname{Ass}(\mathcal F)
\]
(3.1.13), on a $\operatorname{Ass}(\mathcal G)=\operatorname{Ass}(\mathcal F)$;
il suffit alors d'appliquer (5.11.1) au $\mathcal O_X$-Module cohérent
$\mathcal G$, en notant que $i_*(\mathcal F)=\mathcal H^0_{X/Z}(\mathcal G)$
et que, lorsque $X$ est universellement caténaire et universellement japonais,
l'hypothèse (ii) de (5.11.1) est vérifiée par définition.

\begin{corollary}[5.11.5]
\label{IV.2.s5.11.5-fr}
Soient $X$ un préschéma localement noethérien, $Z$ une partie de $X$ stable par
spécialisation, $\mathcal F$ un $\mathcal O_X$-Module cohérent tel que l'on ait
$\operatorname{Ass}(\mathcal F)\subset X-Z$. Alors la condition:
\begin{enumerate}
\item[a)] $\mathcal H^0_{X/Z}(\mathcal F)$ est un $\mathcal O_X$-Module
cohérent
\end{enumerate}
implique la suivante:
\begin{enumerate}
\item[d)] Pour toute partie $T$ de $Z$ fermée dans $X$ (ou seulement pour une
famille filtrante

\oldpage[IV-2]{125}

croissante de parties fermées $T$ de $Z$, de réunion $Z$),
$i_*(\mathcal F\vert X-T)$ (où $i:X-T\to X$ est l'injection canonique) est
cohérent.
\end{enumerate}
Lorsque $X$ vérifie l'une des hypothèses (i), (ii) de (5.11.1), a) et d) sont
équivalentes.
\end{corollary}

Notons que l'on a
$\operatorname{Ass}(i_*(\mathcal F\vert X-T))=\operatorname{Ass}(\mathcal F)$
en vertu de l'hypothèse et de (3.1.13); il résulte donc de (5.10.2) que les
applications canoniques
\[
i_*(\mathcal F\vert X-T)\longrightarrow\mathcal H^0_{X/Z}(\mathcal F)
\]
sont injectives; le fait que a) implique d) est donc conséquence de cette
remarque. Inversement, la condition d) implique, en vertu de (5.11.1), que
$\operatorname{codim}(T\cap Y_\alpha,Y_\alpha)\geq2$ avec les notations de
(5.11.1); par suite on a
$\operatorname{codim}(Z\cap Y_\alpha,Y_\alpha)\geq2$ puisque $Z$ est réunion
de ses parties qui sont fermées dans $X$, et la dernière assertion du
corollaire découle de (5.11.1).

\begin{corollary}[5.11.6]
\label{IV.2.s5.11.6-fr}
Soient $A$ un anneau noethérien, $X=\operatorname{Spec}(A)$. Considérons les
propriétés suivantes:
\begin{enumerate}
\item[a)] Pour tout anneau quotient intègre $B$ de $A$, l'anneau $B^{(1)}$
(notation de (5.11.2)) est une $B$-algèbre finie.
\item[b)] Pour tout $\mathcal O_X$-Module cohérent $\mathcal F$ et toute partie
$Z$ de $X$, stable par spécialisation, et telle que pour tout
$x\in\operatorname{Ass}(\mathcal F)\cap(X-Z)$ on ait
$\operatorname{codim}(\overline{\{x\}}\cap Z,\overline{\{x\}})\geq2$, le
$\mathcal O_X$-Module $\mathcal H^0_{X/Z}(\mathcal F)$ est cohérent.
\item[c)] Pour toute partie fermée $T$ de $X$ et tout $\mathcal O_U$-Module
cohérent $\mathcal G$ (où $U=X-T$) tel que pour tout
$x\in\operatorname{Ass}(\mathcal G)$ on ait
$\operatorname{codim}(\overline{\{x\}}\cap T,\overline{\{x\}})\geq2$,
$i_*(\mathcal G)$ (où $i:U\to X$ est l'injection canonique) est un
$\mathcal O_X$-Module cohérent.
\item[d)] Pour tout anneau quotient intègre $B$ de $A$ et tout idéal
$\mathfrak J$ de hauteur $\geq2$ dans $B$, l'anneau
\[
\bigcap_{\mathfrak p\not\supset\mathfrak J}B_\mathfrak p
\]
est une $B$-algèbre finie.
\end{enumerate}
On a alors les implications
\[
\mathrm a)\Longleftrightarrow\mathrm b)\Longrightarrow\mathrm c)
\Longleftrightarrow\mathrm d).
\]
En outre, les conditions a), b), c) et d) sont vérifiées dans chacun des deux
cas suivants:
\begin{enumerate}
\item[(i)] $A$ est quotient d'un anneau régulier.
\item[(ii)] $A$ est universellement caténaire et universellement japonais.
\end{enumerate}
\end{corollary}

Soit $B=A/\mathfrak q$, où $\mathfrak q$ est un idéal premier de $A$, de sorte
que $Y=\operatorname{Spec}(B)$ est la partie fermée $V(\mathfrak q)$ de $X$;
posons $Z=Z^{(2)}(Y)$, qui est une partie de $X$ stable par spécialisation;
comme $B$ est intègre, $\operatorname{Ass}(A/\mathfrak q)$ est réduit au point
générique $\mathfrak q$ de $Y$. Si la condition b) est vérifiée, on peut
l'appliquer au $\mathcal O_X$-Module cohérent
$\mathcal F=(A/\mathfrak q)^\sim$ et à $Z$, et en vertu de (5.9.3.1), cela
montre que $B^{(1)}$ est un $A$-module de type fini, et a fortiori un
$B$-module de type fini. Inversement, supposons a) vérifiée; alors, si
$\mathcal F$ est un $\mathcal O_X$-Module cohérent tel que pour tout
$x\in\operatorname{Ass}(\mathcal F)\cap(X-Z)$ on ait
$\operatorname{codim}(\overline{\{x\}}\cap Z,\overline{\{x\}})\geq2$, on peut
appliquer (avec les notations de (5.11.1)) à chacun des schémas affines
$Y_\alpha=\operatorname{Spec}(B_\alpha)$, où $B_\alpha$ est un anneau quotient
intègre de $A$, le résultat de a); comme par hypothèse $Z_\alpha$ est contenu
dans $Z^{(2)}(Y_\alpha)$, la condition a) (compte tenu de (5.10.2)

\oldpage[IV-2]{126}

et du fait que $\operatorname{Ass}(\mathcal O_{Y_\alpha})$ est réduit au point
générique de $Y_\alpha$) entraîne que
$\mathcal H^0_{Y_\alpha/Z_\alpha}(\mathcal O_{Y_\alpha})$ est un
$\mathcal O_{Y_\alpha}$-Module cohérent, et b) découle alors de (5.11.1).

Pour voir que c) entraîne d), on raisonne comme ci-dessus en appliquant c) au
cas où $\mathcal F=(A/\mathfrak q)^\sim$ et $Z=V(\mathfrak J)$;
inversement, on prouve que d) entraîne c) en utilisant encore l'équivalence de
a) et b) dans (5.11.1). Il est évident que c) est un cas particulier de b).
Enfin, le fait que a) (et par suite chacune des autres conditions) est vérifiée
lorsque $A$ vérifie l'une des hypothèses (i), (ii) résulte de (5.11.2), vu la
définition des anneaux universellement caténaires et des anneaux
universellement japonais.

\begin{remarks}[5.11.7]
\label{IV.2.s5.11.7-fr}
\begin{enumerate}
\item[(i)] On ignore si, dans (5.11.5) la condition d) implique a) sans
hypothèse supplémentaire sur $X$; nous verrons plus loin (7.2.4) qu'il en est
bien ainsi lorsque $X$ est un schéma local. De même, nous prouverons que
lorsque $A$ est un anneau local noethérien les quatre propriétés a), b), c),
d) de (5.11.6) sont équivalentes (7.2.4). On ignore si ce résultat s'étend à
tous les anneaux noethériens.
\item[(ii)] Si $A$ vérifie la propriété a) de (5.11.6), il en est de même de
tout anneau de fractions $S^{-1}A$ et de toute $A$-algèbre finie $C$. En effet
tout anneau quotient intègre de $S^{-1}A$ est de la forme
$S^{-1}(A/\mathfrak q)$, où $\mathfrak q$ est un idéal premier de $A$ ne
rencontrant pas $S$; et d'autre part, si $\mathfrak r$ est un idéal premier de
$C$, $\mathfrak p$ son image réciproque dans $A$, $C/\mathfrak r$ est une
$(A/\mathfrak p)$-algèbre finie intègre contenant $A/\mathfrak p$; nos
assertions sont donc des conséquences de (5.10.17, (iii) et (iv)).
\end{enumerate}
\end{remarks}

\subsection{\texorpdfstring{Relations entre les propriétés d'un anneau local noethérien $A$ et
d'un anneau quotient $A/tA$}{Relations entre les propriétés d'un anneau local noethérien A et d'un anneau quotient A/tA}}

On a déjà vu dans (3.4) des relations entre les propriétés de $A$ et de $A/tA$
concernant les idéaux premiers associés, ainsi que les propriétés d'être
intègre ou réduit. On donne dans cette section d'autres relations entre les
propriétés de ces anneaux, liées aux notions de dimension et de profondeur.

\begin{proposition}[5.12.1]
\label{IV.2.s5.12.1-fr}
Soient $A$ un anneau local noethérien, $X=\operatorname{Spec}(A)$, $t$ un
élément de $A$ faisant partie d'un système de paramètres (0, 16.3.6), $X_0$ le
sous-espace fermé $V(t)$ de $X$, $X_1$ l'ouvert complémentaire $X-X_0$. Soit
$Z$ une partie de $X$ telle que toute spécialisation d'un point de $Z$
appartienne à $Z$. On suppose que $X$ est biéquidimensionnel (ce qui sera le
cas si $A$ est équidimensionnel et quotient d'un anneau local régulier (0,
16.5.12)). Alors si l'on pose $Z_0=Z\cap X_0$, $Z_1=Z\cap X_1$, on a
\begin{equation}
\operatorname{codim}(Z_0,X_0)\leq\operatorname{codim}(Z,X)
\leq\operatorname{codim}(Z_1,X_1).
\tag{5.12.1.1}\label{IV.2.s5.12.1.1-fr}
\end{equation}
\end{proposition}

La seconde inégalité résulte de la définition (5.1.3); pour démontrer la
première, on peut se borner à la démontrer lorsque $Z$ est fermé: en effet, par
hypothèse $Z$ est réunion de parties fermées $Z_\alpha$, et si l'on a
$\operatorname{codim}(X_0\cap Z_\alpha,X_0)
\leq\operatorname{codim}(Z_\alpha,X)$ pour tout $\alpha$, on aura aussi
\[
\operatorname{codim}(Z_0,X_0)
=\inf_\alpha\operatorname{codim}(X_0\cap Z_\alpha,X_0)
\leq\inf_\alpha\operatorname{codim}(Z_\alpha,X)
=\operatorname{codim}(Z,X).
\]
Supposons donc $Z$ fermé, de sorte que l'on a
\[
\operatorname{codim}(Z,X)=\dim(X)-\dim(Z)
\]

\oldpage[IV-2]{127}

en vertu de (0, 14.3.5.1). D'autre part, $X_0$ est évidemment caténaire,
équidimensionnel et de dimension $\dim(X)-1$ en vertu de (0, 16.3.4); on a
donc aussi
\[
\operatorname{codim}(Z_0,X_0)=\dim(X_0)-\dim(Z_0)
=\dim(X)-1-\dim(Z_0).
\]
Mais on a $\dim(Z_0)\geq\dim(Z)-1$ (0, 16.3.4), ce qui achève de prouver la
première inégalité (5.12.1.1).

\begin{proposition}[5.12.2]
\label{IV.2.s5.12.2-fr}
Soient $A$ un anneau local noethérien, $M$ un $A$-module de type fini, $t$ un
élément $M$-régulier appartenant à l'idéal maximal $\mathfrak m$ de $A$, $k$
un entier $\geq1$. On suppose que $A$ soit un anneau caténaire. Alors, si
$M/tM$ est équidimensionnel et satisfait à $(S_k)$, il en est de même de $M$.
\end{proposition}

Compte tenu de $(\mathbf{Err}_{\mathrm{III}},30)$ et de (5.7.3, (vi)), on peut
se borner au cas où $\operatorname{Supp}(M)=\operatorname{Spec}(A)=X$. Posons
$X_0=V(t)=\operatorname{Spec}(A/tA)$; comme $M/tM$ vérifie $(S_1)$ par
hypothèse, il n'a pas d'idéaux premiers associés immergés (5.7.5). Appliquant
(3.4.4) au point fermé de $X$ (et de tout sous-préschéma fermé de $X$), on
voit que les composantes irréductibles de $\operatorname{Supp}(M/tM)$ sont
exactement celles de $Y_i\cap X_0$, où $Y_i$ ($1\leq i\leq r$) sont les
composantes irréductibles de $X$. Comme $t$ est par hypothèse $M$-régulier,
$V(t)$ ne contient aucun point maximal de $X$ (3.1.8); chacune des composantes
irréductibles de $Y_i\cap X_0$ est donc de codimension 1 dans $Y_i$ (5.1.8);
comme $X$ est caténaire et que toutes les composantes irréductibles de
$\operatorname{Supp}(M/tM)$ ont par hypothèse même dimension, on en conclut
que les $Y_i$ ont tous même dimension, autrement dit $X$ est équidimensionnel,
donc biéquidimensionnel puisque $A$ est local et caténaire. Pour voir que $M$
vérifie $(S_k)$, appliquons le critère (5.7.4): il faut vérifier (avec les
notations de (5.7.4)) que, pour tout entier $n\geq0$, on a
$\operatorname{codim}(Z_n,X)>n+k$. Or, l'hypothèse sur $M/tM$ et le critère
(5.7.4) montrent que l'on a
$\operatorname{codim}(Z_n\cap X_0,X_0)>n+k$ pour tout entier $n\geq0$. Mais
toute spécialisation d'un point de $Z_n$ appartient encore à $Z_n$, comme nous
le verrons dans (6.11.5)
\footnote{Le lecteur vérifiera que (5.12.2) n'est pas utilisé pour la
démonstration de (6.11.5).}; comme $X$ est biéquidimensionnel et que $t$
appartient à un système de paramètres pour $M$ (0, 16.4.1), donc aussi pour
$A$ (l'annulateur de $M$ étant nilpotent en vertu de l'hypothèse
$\operatorname{Supp}(M)=\operatorname{Spec}(A)$), la conclusion résulte de
(5.12.1).

\begin{remark}[5.12.3]
\label{IV.2.s5.12.3-fr}
Si, dans l'énoncé de (5.12.2), on suppose seulement que $M/tM$ est
équidimensionnel, il n'en résulte pas nécessairement que $M$ soit
équidimensionnel. Par exemple, soient $k$ un corps, $B$ l'anneau de polynômes
$k[T,U,V]$ à 3 indéterminées, $C$ l'anneau local de $B$ en l'idéal maximal
$BT+BU+BV$ de $B$, et soient $\mathfrak p=CU$,
$\mathfrak q=CV+C(T+U)$; considérons alors l'anneau local
$A=C/\mathfrak p\mathfrak q$ (géométriquement, si $X$ est le sous-préschéma
fermé de l'espace affine à 3 dimensions sur $k$, formé d'un plan et d'une
droite coupant ce plan en un point $x$, $A$ est l'anneau local de $X$ au point
$x$). Soient $t,u,v$ les images canoniques de $T,U,V$ dans $A$; il est
immédiat que $t$ n'est pas diviseur de zéro dans $A$, et $A/tA$ est isomorphe
à $C_0/\mathfrak p_0\mathfrak q_0$, où $C_0$ est l'anneau local de
$B_0=k[U,V]$ en l'idéal maximal $B_0U+B_0V$, et
$\mathfrak p_0=C_0U$, $\mathfrak q_0=C_0U+C_0V$ ($\mathfrak q_0$ étant donc
l'idéal maximal de $C_0$). Il est immédiat que $\operatorname{Spec}(A/tA)$

\oldpage[IV-2]{128}

est irréductible et de dimension 1, $A/tA$ ne vérifie pas $(S_1)$ et
$\operatorname{Spec}(A)$ n'est pas équidimensionnel.
\end{remark}

\begin{corollary}[5.12.4]
\label{IV.2.s5.12.4-fr}
Soient $A$ un anneau local noethérien caténaire, $t$ un élément régulier
appartenant à l'idéal maximal de $A$, $k$ un entier $\geq1$. Si $A/tA$
vérifie $(S_k)$, il en est de même de $A$.
\end{corollary}

Si $k=1$, le corollaire résulte de (3.4.4), vu l'interprétation de la condition
$(S_1)$ (5.7.5). Si $k\geq2$, il résulte du critère de Hartshorne (5.10.9) que
$A/tA$ est équidimensionnel, et on peut appliquer (5.12.2).

\begin{proposition}[5.12.5]
\label{IV.2.s5.12.5-fr}
Soient $A$ un anneau local noethérien caténaire, $t$ un élément régulier de $A$
appartenant à l'idéal maximal de $A$, $k$ un entier $\geq0$. Si l'anneau
$A/tA$ est réduit, équidimensionnel et vérifie la propriété $(R_k)$, l'anneau
$A$ possède aussi ces trois propriétés.
\end{proposition}

On sait déjà (3.4.6) que $A$ est réduit, et il résulte de (5.12.2), appliqué
pour $k=1$, que $A$ est équidimensionnel. Posons
$X=\operatorname{Spec}(A)$, $X_0=V(t)=\operatorname{Spec}(A/tA)$, et soit $Z$
l'ensemble des points $x\in X$ où $X$ n'est pas régulier; en vertu de (0,
17.3.2) toute spécialisation d'un point de $Z$ appartient à $Z$. Il résulte
d'autre part de (0, 17.1.8) que pour tout point $x\in X_0$ où $X_0$ est
régulier, $X$ est aussi régulier, donc l'ensemble $Z'_0$ des points où $X_0$
n'est pas régulier contient $Z_0=Z\cap X_0$; l'hypothèse entraîne donc que
\[
k<\operatorname{codim}(Z'_0,X_0)\leq\operatorname{codim}(Z_0,X_0).
\]
Mais, puisque $A$ est équidimensionnel et caténaire, on a, d'après (5.12.1),
\[
\operatorname{codim}(Z_0,X_0)\leq\operatorname{codim}(Z,X)
\]
ce qui prouve que $X$ vérifie $(R_k)$.

\begin{remark}[5.12.6]
\label{IV.2.s5.12.6-fr}
Si, dans l'énoncé de (5.12.5), on suppose seulement que $A/tA$ soit réduit et
possède la propriété $(R_k)$, il n'en résulte pas nécessairement que $A$
possède la propriété $(R_k)$. Par exemple, soient $k$ un corps, $P(U,V)$ un
polynôme irréductible de $k[U,V]$ tel que la courbe $\Gamma$ définie par
l'idéal principal $(P)$ dans le plan affine $\operatorname{Spec}(k[U,V])$ ait
un point singulier correspondant à l'idéal maximal $(U)+(V)$ (par exemple
$P(U,V)=U(U^2+V^2)+(U^2-V^2)$, qui donne une cubique à point double). Soient
alors $B$ l'anneau de polynômes à 4 indéterminées $k[T,U,V,W]$, $C$ l'anneau
local de $B$ en l'idéal maximal $BT+BU+BV+BW$, et soient
\[
\mathfrak p=CW+CP(U-T,V),\qquad \mathfrak q=CU.
\]
Considérons alors l'anneau local $A=C/\mathfrak p\mathfrak q$
(géométriquement, si $X$ est le sous-préschéma fermé de l'espace affine à 4
dimensions, formé d'un hyperplan $H$ et d'un « cylindre » $L$ à 2 dimensions
ayant pour « base » la courbe $\Gamma$ et dont la « droite singulière » $Y$
n'est pas contenue dans l'hyperplan $H$, alors $A$ est l'anneau local de $X$ au
point $x$ où $Y$ rencontre $H$). On voit alors aussitôt que
$\operatorname{Spec}(A/tA)$ (où $t$ est l'image canonique de $T$ dans $A$) a
pour seul point singulier $x$, dont l'anneau local est $A/tA$ lui-même, qui est
réduit et de dimension 2; au contraire, le point générique $y$ de la « droite
singulière » $Y$ (définie par l'image dans $A$ de l'idéal
$C(U-T)+CV+CW$) est un point singulier de $X$, et
$\mathcal O_{X,y}$ est de dimension 1; autrement dit, $A/tA$ est réduit et
vérifie $(R_1)$, mais $A$ ne vérifie pas $(R_1)$.
\end{remark}

\oldpage[IV-2]{129}

\begin{corollary}[5.12.7]
\label{IV.2.s5.12.7-fr}
Soient $A$ un anneau local noethérien caténaire, $t$ un élément régulier de
$A$ appartenant à l'idéal maximal de $A$. Si $A/tA$ est intègre et
intégralement clos, il en est de même de $A$.
\end{corollary}

En vertu du critère de Serre (5.8.6) et du fait que l'anneau $A$ est local, il
suffit de montrer que $\operatorname{Spec}(A)$ vérifie $(S_2)$ et $(R_1)$.
Mais l'hypothèse sur $A/tA$ entraîne que $A$ vérifie $(R_1)$ par (5.12.5), et
que $A$ vérifie $(S_2)$ par (5.12.4).

\begin{proposition}[5.12.8]
\label{IV.2.s5.12.8-fr}
(lemme de Hironaka). Soit $A$ un anneau local noethérien réduit,
équidimensionnel et caténaire. Supposons de plus que pour tout idéal premier
minimal $\mathfrak q_i$ de $A$, l'anneau $B_i=A/\mathfrak q_i$ soit tel que
$B_i^{(1)}$ (notation de (5.10.17)) soit une $B_i$-algèbre finie (ce qui aura
lieu par exemple lorsque $A$ possède l'une des propriétés (i), (ii) de
(5.11.6)). Soit $t$ un élément de $A$, non diviseur de zéro dans $A$, et tel
que:
\begin{enumerate}
\item[(i)] Le $A$-module $A/tA$ n'ait qu'un seul idéal premier associé non
immergé $\mathfrak p$ (nécessairement de hauteur 1 en vertu du
Hauptidealsatz (0, 16.3.2), mais on ne suppose pas que ce soit le seul idéal
premier associé à $A/tA$).
\item[(ii)] L'idéal $\mathfrak pA_{\mathfrak p}$ de $A_{\mathfrak p}$ est
engendré par $t/1$.
\item[(iii)] L'anneau $A/\mathfrak p$ est intégralement clos.
\end{enumerate}
Sous ces conditions, l'anneau $A$ est intègre et intégralement clos, et l'on a
$\mathfrak p=tA$.
\end{proposition}

Si $Z=V(tA)$, l'hypothèse (i) entraîne que $Z$ est une partie fermée
irréductible de $X=\operatorname{Spec}(A)$, dont le point générique $z$ est tel
que $\mathfrak i_z=\mathfrak p$. L'hypothèse (ii) montre que l'idéal maximal
$\mathfrak pA_{\mathfrak p}$ de l'anneau local noethérien $A_{\mathfrak p}$
est engendré par un seul élément, donc $A_{\mathfrak p}$ est un anneau de
valuation discrète dont $t/1$ est une uniformisante (Bourbaki, \emph{Alg.
comm.}, chap. VI, § 3, n\textsuperscript{o} 6, prop. 9);
$A_{\mathfrak p}/tA_{\mathfrak p}$ est le corps résiduel de
$A_{\mathfrak p}$, donc un $A_{\mathfrak p}$-module de longueur 1. En vertu de
(3.4.2), cela entraîne que $Z$ n'est contenu que dans une seule des composantes
irréductibles $X_i$ de $X$; pour toute autre composante irréductible $X_j$ de
$X$, on a par suite $\dim(X_j\cap Z)<\dim(Z)$. D'autre part, comme $t$ n'est
pas diviseur de zéro dans $A$, il n'est contenu dans aucun des
$\mathfrak q_i$, et l'on a par suite (5.1.8)
\[
\operatorname{codim}(X_i\cap Z,X_i)
=\operatorname{codim}(X_j\cap Z,X_j)=1,
\]
quel que soit $j\neq i$. Comme $A$ est supposé biéquidimensionnel, la relation
$\dim(X_i\cap Z)\neq\dim(X_j\cap Z)$ entraînerait donc
$\dim(X_i)\neq\dim(X_j)$ (0, 14.3.3.1), ce qui est absurde. On voit donc qu'il
n'y a qu'un seul idéal premier minimal dans $A$; comme $A$ est par hypothèse
réduit, cela entraîne qu'il est intègre. Remarquons maintenant que puisque
$A^{(1)}$ est une $A$-algèbre finie, c'est un anneau semi-local noethérien, et
tout idéal premier non immergé associé à $A^{(1)}/tA^{(1)}$ est donc de
hauteur 1 en vertu du Hauptidealsatz (0, 16.3.2); or, pour un tel idéal
$\mathfrak r$, $\mathfrak r\cap A$ est de hauteur 1 en vertu de (5.10.17,
(iv)) et comme il contient $tA$, ce ne peut être que $\mathfrak p$ en vertu de
l'hypothèse (i). D'autre part, $A^{(1)}$ est contenu dans la clôture intégrale
$A'$ de $A$; si l'on pose $S=A-\mathfrak p$, la clôture intégrale de
$A_{\mathfrak p}$ est $S^{-1}A'$ (Bourbaki, \emph{Alg. comm.}, chap. V, § 1,
n\textsuperscript{o} 5, prop. 17), donc $S^{-1}A'=A_{\mathfrak p}$ puisque
$A_{\mathfrak p}$ est un anneau de valuation discrète; il en résulte
(\emph{loc. cit.}, § 2, n\textsuperscript{o} 1, lemme 1) qu'il n'existe qu'un
seul idéal premier $\mathfrak p'$ de $A'$ au-dessus de $\mathfrak p$, et a
fortiori un seul idéal premier $\mathfrak p^{(1)}$ de $A^{(1)}$ au-dessus de
$\mathfrak p$, et l'on a
\[
A_{\mathfrak p}=A'_{\mathfrak p'}
=A^{(1)}_{\mathfrak p^{(1)}}.
\]
Notons maintenant que $A^{(1)}$ vérifie $(S_2)$ (5.10.17, (i)) et comme $t$
n'est pas diviseur de zéro dans $A^{(1)}$, $A^{(1)}/tA^{(1)}$ n'a pas
d'idéaux premiers associés

\oldpage[IV-2]{130}

immergés (5.7.7). L'idéal $\mathfrak p^{(1)}$ est donc le seul idéal premier
associé à $A^{(1)}/tA^{(1)}$, autrement dit, $tA^{(1)}$ est un idéal
$\mathfrak p^{(1)}$-primaire de $A^{(1)}$; cela veut dire encore que
$tA^{(1)}$ est l'image réciproque dans $A^{(1)}$ de
$tA^{(1)}_{\mathfrak p^{(1)}}=tA_{\mathfrak p}$, et ce dernier est l'idéal
maximal de $A_{\mathfrak p}$ d'après (ii); donc
$tA^{(1)}=\mathfrak p^{(1)}$ est premier dans $A^{(1)}$. D'autre part,
$A^{(1)}/\mathfrak p^{(1)}$ est fini sur $A/\mathfrak p$ et a même corps des
fractions (savoir le corps résiduel de
$A^{(1)}_{\mathfrak p^{(1)}}=A_{\mathfrak p}$), donc il est identique à
$A/\mathfrak p$ en vertu de l'hypothèse (iii). On peut donc écrire
$A^{(1)}=A+\mathfrak p^{(1)}=A+tA^{(1)}$, et comme $t$ est contenu dans le
radical de l'anneau semi-local noethérien $A^{(1)}$, le lemme de Nakayama
entraîne $A^{(1)}=A$, d'où $\mathfrak p^{(1)}=\mathfrak p=tA$. Mais comme $A$
est caténaire et $A/tA$ intègre et intégralement clos en vertu de l'hypothèse
(iii), on conclut de (5.12.7) que $A$ est intégralement clos. C.Q.F.D.

\begin{remark}[5.12.9]
\label{IV.2.s5.12.9-fr}
La conclusion de (5.12.8) tombe en défaut si l'on ne suppose plus $A$
équidimensionnel. Par exemple, soient $k$ un corps, $B$ l'anneau de polynômes
$k[X,Y,Z]$ à trois indéterminées, $C$ l'anneau
$B/\mathfrak r_1\mathfrak r_2$, où $\mathfrak r_1=BZ$ et
$\mathfrak r_2=BX+BY$; soient $\mathfrak m$ l'idéal maximal $BX+BY+BZ$,
$\mathfrak n=\mathfrak m/\mathfrak r_1\mathfrak r_2$ son image dans $C$, $A$
l'anneau local $C_{\mathfrak n}$. Enfin, soit $t_0$ l'image dans $C$ de
l'élément $Y+Z$ de $B$, $t$ son image dans $A$
($\operatorname{Spec}(C)$ est donc formé d'un plan $P$ et d'une droite $D$ non
contenue dans ce plan, $\operatorname{Spec}(C/t_0C)$ d'une droite $D'$ contenue
dans $P$ et passant par le point $D\cap P$). L'anneau $A$ est réduit et
caténaire (étant quotient d'un anneau régulier) et ses idéaux premiers minimaux
$\mathfrak q_1$, $\mathfrak q_2$ sont les images de $\mathfrak r_1$,
$\mathfrak r_2$. Le $A$-module $A/tA$ n'a qu'un seul idéal premier associé non
immergé $\mathfrak p$, image de $BY+BZ$, $\mathfrak pA_{\mathfrak p}$ est
engendré par $t/1$ et $A/\mathfrak p$ isomorphe à $k[X]$; mais $A$ n'est pas
intègre.
\end{remark}

\begin{corollary}[5.12.10]
\label{IV.2.s5.12.10-fr}
Soit $A$ un anneau local intègre noethérien, vérifiant l'une des hypothèses
suivantes:
\begin{enumerate}
\item[a)] $A$ est quotient d'un anneau régulier.
\item[b)] $A$ est universellement caténaire et universellement japonais.
\end{enumerate}
Soit $(x_i)_{1\leq i\leq n}$ une suite de $n$ éléments de $A$; posons
$\mathfrak J=x_1A+\cdots+x_nA$, et supposons vérifiées les conditions
suivantes:
\begin{enumerate}
\item[(i)] Il n'existe qu'un seul idéal premier non immergé $\mathfrak p$
associé à $A/\mathfrak J$ et $\mathfrak p$ est de hauteur $n$.
\item[(ii)] L'idéal maximal $\mathfrak pA_{\mathfrak p}$ de $A_{\mathfrak p}$
est engendré par les $x_i/1$ (donc $A_{\mathfrak p}$ est un anneau local
régulier de dimension $n$).
\item[(iii)] L'anneau $A/\mathfrak p$ est intégralement clos.
\end{enumerate}
Sous ces conditions, pour tout entier $i$ tel que $0\leq i\leq n$, l'anneau
quotient $A/(x_1A+\cdots+x_iA)$ est intégralement clos et de dimension
$\dim(A)-i$. En particulier, $\mathfrak J$ est premier, égal à $\mathfrak p$,
$A$ est intégralement clos et $(x_i)_{1\leq i\leq n}$ est une suite
$A$-régulière.
\end{corollary}

Procédons par récurrence sur $n$, la proposition étant triviale pour $n=0$ et
se réduisant au lemme de Hironaka (5.12.8) pour $n=1$. On peut donc supposer
$n\geq2$. Soit $\mathfrak J'=x_1A+\cdots+x_{n-1}A$, et soit $\mathfrak q$ un
idéal premier minimal parmi ceux qui contiennent $\mathfrak J'$; on a
$\operatorname{ht}(\mathfrak q)\leq n-1$ (0, 16.3.1); montrons que l'on a
$\mathfrak q\subset\mathfrak p$. En effet, si $\mathfrak p'$ est minimal parmi
les idéaux premiers de $A$ contenant $\mathfrak q+x_nA$,
$\mathfrak p'/\mathfrak q$ est de hauteur $\leq1$ dans $A/\mathfrak q$ par le
Hauptidealsatz (0, 16.3.2), donc $\mathfrak p'$ est de hauteur $\leq n$ puisque
$A$ est caténaire (0, 16.1.4); mais comme il contient $\mathfrak J$, il est
nécessairement égal à $\mathfrak p$, d'où
$\mathfrak q\subset\mathfrak p$,

\oldpage[IV-2]{131}

et $\mathfrak q$ est induit sur $A$ par un idéal premier de
$A_{\mathfrak p}$, minimal parmi ceux qui contiennent
$\mathfrak J'A_{\mathfrak p}$. Mais en vertu de l'hypothèse (ii) et de (0,
17.1.7), $\mathfrak J'A_{\mathfrak p}$ est premier dans $A_{\mathfrak p}$,
donc
\[
\mathfrak q=A\cap\mathfrak J'A_{\mathfrak p}
\]
est l'unique idéal premier non immergé associé à $A/\mathfrak J'$ et il est de
hauteur $n-1$, puisque $\mathfrak J'A_{\mathfrak p}$ est de hauteur $n-1$ et
$\mathfrak q$ de hauteur $\leq n-1$. En outre
$\mathfrak qA_{\mathfrak q}=\mathfrak J'A_{\mathfrak q}$ et comme
$A_{\mathfrak q}=(A_{\mathfrak p})_{\mathfrak qA_{\mathfrak p}}$, l'idéal
maximal $\mathfrak qA_{\mathfrak q}$ est engendré par $\mathfrak J'$. On voit
donc que la suite $(x_i)_{1\leq i\leq n-1}$ vérifie déjà les conditions (i) et
(ii) de l'énoncé. Montrons qu'elle vérifie aussi (iii). Considérons pour cela
l'anneau intègre $B=A/\mathfrak q$, et soit $t$ l'image canonique de $x_n$ dans
$B$. Il est immédiat (cf. (5.6.3, (iv))) que si $A$ vérifie l'hypothèse a)
(resp. b)), il en est de même de $B$; comme $x_n\notin\mathfrak q$, on a
$t\neq0$; montrons que $B$ et $t$ vérifient les hypothèses du lemme de
Hironaka. En effet, un idéal premier $\mathfrak n$ de $B$ minimal parmi ceux
contenant $t$ est l'image d'un idéal premier de $A$ minimal parmi ceux
contenant $\mathfrak q+x_nA$, et l'on a vu que $\mathfrak p$ est cet unique
idéal; comme
\[
B_{\mathfrak n}=A_{\mathfrak p}/\mathfrak qA_{\mathfrak p},\qquad
\mathfrak nB_{\mathfrak n}
=\mathfrak pA_{\mathfrak p}/\mathfrak qA_{\mathfrak p},
\]
le fait que $\mathfrak qA_{\mathfrak p}=\mathfrak J'A_{\mathfrak p}$ et que
$\mathfrak pA_{\mathfrak p}$ soit engendré par $\mathfrak J$ entraîne que
$\mathfrak nB_{\mathfrak n}$ est engendré par $t$. Enfin,
$B/\mathfrak n=A/\mathfrak p$ est intégralement clos. L'application de
(5.12.8) prouve donc que $B=A/\mathfrak q$ est intégralement clos et que
$\mathfrak p=\mathfrak q+x_nA$. Utilisons maintenant l'hypothèse de récurrence,
qui montre que $A/(x_1A+\cdots+x_iA)$ est intégralement clos et de dimension
$\dim(A)-i$ pour $0\leq i\leq n-1$ et que $\mathfrak q=\mathfrak J'$, d'où
$\mathfrak p=\mathfrak J'+x_nA=\mathfrak J$. On en conclut que
$A/\mathfrak J=A/\mathfrak p$ est intégralement clos; comme
$\dim(A_{\mathfrak p})=n$ et que
$\dim(A/\mathfrak p)=\dim(A)-\dim(A_{\mathfrak p})$ puisque $A$ est
biéquidimensionnel, cela achève la démonstration.

\subsection{Propriétés de permanence par passage à la limite inductive}

Dans tout ce numéro, $(A_\alpha,\varphi_{\beta\alpha})$ désigne un système
inductif d'anneaux dont l'ensemble d'indices est préordonné filtrant; on pose
$A=\varinjlim A_\alpha$ et on désigne par $\varphi_\alpha$ l'homomorphisme
canonique $A_\alpha\to A$.

\begin{proposition}[5.13.1]
\label{IV.2.s5.13.1-fr}
\begin{enumerate}
\item[(i)] Pour tout indice $\alpha$, soit $\mathfrak a_\alpha$ un idéal de
$A_\alpha$ tel que, pour $\alpha\leq\beta$, on ait
$\varphi_{\beta\alpha}(\mathfrak a_\alpha)\subset\mathfrak a_\beta$. Alors
les $\mathfrak a_\alpha$ forment un système inductif pour les restrictions des
$\varphi_{\beta\alpha}$, $\mathfrak a=\varinjlim\mathfrak a_\alpha$
s'identifie canoniquement à un idéal de $A$, $A/\mathfrak a$ à
$\varinjlim(A_\alpha/\mathfrak a_\alpha)$. Si de plus on a
$\mathfrak a_\alpha=\varphi_{\beta\alpha}^{-1}(\mathfrak a_\beta)$ pour
$\alpha\leq\beta$, alors
$\mathfrak a_\alpha=\varphi_\alpha^{-1}(\mathfrak a)$ pour tout $\alpha$.
\item[(ii)] Inversement, pour tout idéal $\mathfrak a$ de $A$, si l'on pose
$\mathfrak a_\alpha=\varphi_\alpha^{-1}(\mathfrak a)$ pour tout $\alpha$, les
$\mathfrak a_\alpha$ forment un système inductif tel que
$\mathfrak a=\varinjlim\mathfrak a_\alpha$.
\end{enumerate}
\end{proposition}

Il est clair que (ii) est un cas particulier de (i). Dans (i), le fait que
$\mathfrak a$ s'identifie canoniquement à un idéal de $A$ et $A/\mathfrak a$ à
$\varinjlim(A_\alpha/\mathfrak a_\alpha)$ provient de l'exactitude du foncteur
$\varinjlim$ dans la catégorie des groupes abéliens. En outre, $\mathfrak a$
est réunion de la famille croissante des $\varphi_\alpha(\mathfrak a_\alpha)$;
donc, si $x_\alpha\in A_\alpha$ est tel que
$\varphi_\alpha(x_\alpha)\in\mathfrak a$, il existe $\beta\geq\alpha$ tel que
$\varphi_\alpha(x_\alpha)=\varphi_\beta(y_\beta)$ pour un
$y_\beta\in\mathfrak a_\beta$, donc
$\varphi_\beta(\varphi_{\beta\alpha}(x_\alpha))=\varphi_\beta(y_\beta)$; mais
cela entraîne l'existence d'un $\gamma\geq\beta$ tel que
\[
\varphi_{\gamma\alpha}(x_\alpha)=\varphi_{\gamma\beta}(y_\beta),
\]
et par suite $\varphi_{\gamma\alpha}(x_\alpha)\in\mathfrak a_\gamma$, d'où
$x_\alpha\in\mathfrak a_\alpha$ si l'on a
$\mathfrak a_\alpha=\varphi_{\gamma\alpha}^{-1}(\mathfrak a_\gamma)$.

\begin{corollary}[5.13.2]
\label{IV.2.s5.13.2-fr}
Si $\mathfrak N_\alpha$ est le nilradical de $A_\alpha$, le nilradical de $A$
s'identifie à $\varinjlim\mathfrak N_\alpha$, et $A_{\mathrm{red}}$ à
$\varinjlim(A_\alpha)_{\mathrm{red}}$. En particulier, si les $A_\alpha$ sont
réduits, il en est de même de $A$.
\end{corollary}

\oldpage[IV-2]{132}

Il est clair en effet que
$\varphi_{\beta\alpha}(\mathfrak N_\alpha)\subset\mathfrak N_\beta$ pour
$\alpha\leq\beta$, et que $\varinjlim\mathfrak N_\alpha$ est un idéal
nilpotent de $A$, donc contenu dans le nilradical $\mathfrak N$ de $A$.
Inversement, si $x\in\mathfrak N$, on a $x^h=0$ pour un entier $h$; on peut
écrire $x=\varphi_\alpha(x_\alpha)$ pour un indice $\alpha$ et un
$x_\alpha\in A_\alpha$; la relation $\varphi_\alpha(x_\alpha^h)=0$ entraîne
l'existence d'un $\beta\geq\alpha$ tel que
$\varphi_{\beta\alpha}(x_\alpha^h)=0$, ou encore $x_\beta^h=0$ en posant
$x_\beta=\varphi_{\beta\alpha}(x_\alpha)$; par suite
$x_\beta\in\mathfrak N_\beta$ et comme $x=\varphi_\beta(x_\beta)$ cela achève
de prouver le corollaire, compte tenu de (5.13.1).

\begin{proposition}[5.13.3]
\label{IV.2.s5.13.3-fr}
\begin{enumerate}
\item[(i)] Pour tout indice $\alpha$, soit $\mathfrak p_\alpha$ un idéal
premier de $A_\alpha$ tel que, pour $\alpha\leq\beta$, on ait
$\varphi_{\beta\alpha}(\mathfrak p_\alpha)\subset\mathfrak p_\beta$. Alors
$\mathfrak p=\varinjlim\mathfrak p_\alpha$ est un idéal premier de $A$. En
particulier, si les $A_\alpha$ sont intègres il en est de même de $A$.
\item[(ii)] Si de plus on a
$\mathfrak p_\alpha=\varphi_{\beta\alpha}^{-1}(\mathfrak p_\beta)$ pour
$\alpha\leq\beta$, alors $A_{\mathfrak p}$ s'identifie canoniquement à
$\varinjlim(A_\alpha)_{\mathfrak p_\alpha}$.
\item[(iii)] Si $A$ est un anneau local, $\mathfrak p$ son idéal maximal, et
si les homomorphismes $\varphi_\alpha$ sont injectifs, alors, en posant
$\mathfrak p_\alpha=\varphi_\alpha^{-1}(\mathfrak p)$, on a
$A=\varinjlim(A_\alpha)_{\mathfrak p_\alpha}$ et les homomorphismes
$(A_\alpha)_{\mathfrak p_\alpha}\to A$ sont injectifs.
\end{enumerate}
\end{proposition}

Les deux assertions de (i) sont en fait équivalentes, puisque
$A/\mathfrak p=\varinjlim(A_\alpha/\mathfrak p_\alpha)$. Supposons donc tous
les $A_\alpha$ intègres, et soient $x$, $y$ deux éléments de $A$ tels que
$xy=0$. Il existe donc un indice $\alpha$ tel que
$x=\varphi_\alpha(x_\alpha)$, $y=\varphi_\alpha(y_\alpha)$ avec
$x_\alpha,y_\alpha$ dans $A_\alpha$, et
$\varphi_\alpha(x_\alpha y_\alpha)=0$; mais cela entraîne l'existence d'un
indice $\beta\geq\alpha$ tel que, si l'on pose
$x_\beta=\varphi_{\beta\alpha}(x_\alpha)$,
$y_\beta=\varphi_{\beta\alpha}(y_\alpha)$, on ait $x_\beta y_\beta=0$; comme
$A_\beta$ est intègre, on a $x_\beta=0$ ou $y_\beta=0$ et comme
$x=\varphi_\beta(x_\beta)$ et $y=\varphi_\beta(y_\beta)$, cela achève de
montrer que $A$ est intègre.

Prouvons maintenant (ii). L'hypothèse entraîne que les
$A_\alpha-\mathfrak p_\alpha$ forment un système inductif de parties
multiplicatives des $A_\alpha$, dont $A-\mathfrak p$ est la limite inductive;
par suite les anneaux locaux $(A_\alpha)_{\mathfrak p_\alpha}$ forment un
système inductif (où les homomorphismes de transition sont locaux), et les
homomorphismes canoniques
$(A_\alpha)_{\mathfrak p_\alpha}\to A_{\mathfrak p}$ un système inductif
d'homomorphismes, dont la limite inductive
\[
\psi:\varinjlim (A_\alpha)_{\mathfrak p_\alpha}
\longrightarrow A_{\mathfrak p}
\]
est un homomorphisme surjectif. Reste à voir que $\psi$ est injectif. Or, si
$x_\alpha\in A_\alpha$ et
$s_\alpha\in A_\alpha-\mathfrak p_\alpha$ sont tels que
$\varphi_\alpha(x_\alpha)/\varphi_\alpha(s_\alpha)=0$ dans $A_{\mathfrak p}$,
il existe $s'\in A-\mathfrak p$ tel que
$s'\varphi_\alpha(x_\alpha)=0$, et, en remplaçant au besoin $\alpha$ par un
indice plus grand, on peut supposer que $s'=\varphi_\alpha(s'_\alpha)$ avec
$s'_\alpha\in A_\alpha-\mathfrak p_\alpha$, d'où
$\varphi_\alpha(s'_\alpha x_\alpha)=0$; il y a donc un $\beta\geq\alpha$ tel
que $\varphi_{\beta\alpha}(s'_\alpha x_\alpha)=0$ dans $A_\beta$, ce qui
entraîne
$\varphi_{\beta\alpha}(x_\alpha)/\varphi_{\beta\alpha}(s_\alpha)=0$ dans
$(A_\beta)_{\mathfrak p_\beta}$ et prouve donc que l'image canonique de
$x_\alpha/s_\alpha$ dans
$\varinjlim(A_\alpha)_{\mathfrak p_\alpha}$ est nulle.

Enfin, la première assertion de (iii) résulte de (ii) et de ce que
$A_{\mathfrak p}=A$. En outre, les éléments de
$A_\alpha-\mathfrak p_\alpha$ sont inversibles dans $A$, donc on a aussi
$A_{\mathfrak p_\alpha}=A$, et le fait que
$(A_\alpha)_{\mathfrak p_\alpha}\to A$ soit injectif se déduit par platitude
de l'hypothèse que $\varphi_\alpha$ est injectif.

\begin{corollary}[5.13.4]
\label{IV.2.s5.13.4-fr}
Supposons que les $A_\alpha$ soient tous intègres, et que les
$\varphi_{\beta\alpha}$ soient injectives. Alors, si $A'_\alpha$ est la clôture
intégrale de $A_\alpha$, $A'=\varinjlim A'_\alpha$ est la clôture intégrale de
$A$. En particulier, si $A_\alpha$ est intégralement clos pour tout $\alpha$,
il en est de même de $A$. Si les $A_\alpha$ sont des anneaux locaux unibranches
(resp. géométriquement unibranches) (0, 23.2.1) et les
$\varphi_{\beta\alpha}$ des homomorphismes locaux injectifs, $A$ est unibranche
(resp. géométriquement unibranche).
\end{corollary}

Si $K_\alpha$ est le corps des fractions de $A_\alpha$, il résulte de (5.13.3,
(ii)) appliqué aux $\mathfrak p_\alpha=0$, que
$K=\varinjlim K_\alpha$ est le corps des fractions de $A$, et
$A'=\varinjlim A'_\alpha$ s'identifie à

\oldpage[IV-2]{133}

un sous-anneau de $K$, évidemment entier sur $A$. Inversement, soit $z\in K$
un élément vérifiant une équation de dépendance intégrale
\[
z^n+\sum_{i=1}^{n}c_i z^{n-i}=0,
\]
où les $c_i\in A$. Il existe un indice $\alpha$, un $z_\alpha\in K_\alpha$ et
des $c_{i\alpha}\in A_\alpha$ tels que $z$ (resp. les $c_i$) soient les images
de $z_\alpha$ (resp. des $c_{i\alpha}$), donc on a
$z_\alpha\in A'_\alpha$, ce qui prouve que $z\in A'$. Il en résulte aussitôt
que si les $A_\alpha$ sont intégralement clos, il en est de même de $A$. Si les
$A_\alpha$ sont locaux et les $\varphi_{\beta\alpha}$ des homomorphismes
locaux, on sait $(0_{\mathrm{III}}, 10.3.1.3)$ que $A$ est un anneau local et
que si $k_\alpha$ est le corps résiduel de $A_\alpha$,
$k=\varinjlim k_\alpha$ est celui de $A$; si les $A'_\alpha$ sont locaux, $A'$
est donc local et $A$ est unibranche; si de plus le corps résiduel
$k'_\alpha$ de $A'_\alpha$ est une extension radicielle de $k_\alpha$ pour
tout $\alpha$, $k'=\varinjlim k'_\alpha$, qui est le corps résiduel de $A'$,
est une extension radicielle de $k$, et $A$ est donc géométriquement
unibranche.

\begin{env}[5.13.5]
\label{IV.2.s5.13.5-fr}
Nous dirons qu'un anneau $A$ est \emph{normal} si
$\operatorname{Spec}(A)$ est un schéma normal, autrement dit si, pour tout
idéal premier $\mathfrak p$ de $A$, $A_{\mathfrak p}$ est intègre et
intégralement clos. On notera que cela entraîne que $A$ est réduit et que
$\mathfrak p$ ne peut contenir qu'un seul idéal premier minimal $\mathfrak q$
de $A$, autrement dit $\mathfrak p$ n'appartient qu'à une seule composante
irréductible de $\operatorname{Spec}(A)$; en outre $\mathfrak q$ est l'image
réciproque de $(0)$ par l'homomorphisme canonique
$A\to A_{\mathfrak p}$, et $A_{\mathfrak p}$ est donc isomorphe à
$(A/\mathfrak q)_{\mathfrak p/\mathfrak q}$, ce qui prouve que l'anneau
intègre $A/\mathfrak q$ est intégralement clos comme intersection d'anneaux
intégralement clos (I, 8.2.1.1). Si $A$ est noethérien et normal, $A$ est
composé direct d'un nombre fini d'anneaux intègres et intégralement clos (I,
5.1.4).
\end{env}

\begin{corollary}[5.13.6]
\label{IV.2.s5.13.6-fr}
Supposons que les $A_\alpha$ soient normaux et que pour $\alpha\leq\beta$,
toute composante irréductible de $\operatorname{Spec}(A_\beta)$ domine une
composante irréductible de $\operatorname{Spec}(A_\alpha)$. Alors
$A=\varinjlim A_\alpha$ est normal.
\end{corollary}

En effet, soit $\mathfrak p$ un idéal premier de $A$ et posons, pour tout
$\alpha$, $\mathfrak p_\alpha=\varphi_\alpha^{-1}(\mathfrak p)$. Par
hypothèse, pour tout $\alpha$, $\mathfrak p_\alpha$ contient un seul idéal
minimal $\mathfrak q_\alpha$ de $A_\alpha$, et pour $\alpha\leq\beta$,
$\varphi_{\beta\alpha}^{-1}(\mathfrak q_\beta)$ est un idéal minimal de
$A_\alpha$ contenu dans $\mathfrak p_\alpha$, donc égal à
$\mathfrak q_\alpha$. Si l'on pose
$B_\alpha=A_\alpha/\mathfrak q_\alpha$, les anneaux intègres $B_\alpha$
forment donc un système inductif pour les homomorphismes
$\varphi'_{\beta\alpha}:B_\alpha\to B_\beta$ déduits des
$\varphi_{\beta\alpha}$ par passage aux quotients, et qui sont injectifs par
définition; en outre les anneaux $B_\alpha$ sont intégralement clos, donc il en
est de même de $B=\varinjlim B_\alpha$ (5.13.4). Si l'on pose
$\mathfrak p'_\alpha=\mathfrak p_\alpha/\mathfrak q_\alpha$, on a de plus
$(B_\alpha)_{\mathfrak p'_\alpha}=(A_\alpha)_{\mathfrak p_\alpha}$; comme
\[
\varinjlim(B_\alpha)_{\mathfrak p'_\alpha}
=\varinjlim(A_\alpha)_{\mathfrak p_\alpha}=A_{\mathfrak p}
\]
est un anneau local en un idéal premier de $B$ (5.13.3), c'est un anneau
intègre et intégralement clos, ce qui démontre le corollaire.

\begin{proposition}[5.13.7]
\label{IV.2.s5.13.7-fr}
Supposons que les $A_\alpha$ soient réguliers et que, pour
$\alpha\leq\beta$, $\varphi_{\beta\alpha}$ fasse de $A_\beta$ un
$A_\alpha$-module plat; supposons en outre que
$A=\varinjlim A_\alpha$ soit noethérien. Alors $A$ est régulier.
\end{proposition}

Par définition, il suffit de prouver que pour tout idéal premier $\mathfrak p$
de $A$, l'anneau local noethérien $A_{\mathfrak p}$ est régulier; or, si
$\mathfrak p_\alpha=\varphi_\alpha^{-1}(\mathfrak p)$, on a
$\mathfrak p=\varinjlim\mathfrak p_\alpha$ et
$A_{\mathfrak p}=\varinjlim(A_\alpha)_{\mathfrak p_\alpha}$ par (5.13.3);
comme par hypothèse les $(A_\alpha)_{\mathfrak p_\alpha}$ sont réguliers et
que $(A_\beta)_{\mathfrak p_\beta}$ est un
$(A_\alpha)_{\mathfrak p_\alpha}$-module plat pour $\alpha\leq\beta$
$(0_{\mathrm I}, 6.3.2)$, on voit qu'on est ramené au cas où les $A_\alpha$ et
$A$ sont en outre des anneaux locaux. Pour prouver que $A$ est régulier, il
suffit alors de montrer que si $M$, $N$ sont deux $A$-modules de type fini, il
existe $i_0$ tel que, pour $i\geq i_0$,

\oldpage[IV-2]{134}

on ait $\operatorname{Tor}^A_i(M,N)=0$ (0, 17.3.1 et 17.2.6). Mais on a le
lemme suivant (où les $A_\alpha$ sont de nouveau quelconques):

\begin{lemma}[5.13.7.1]
\label{IV.2.s5.13.7.1-fr}
Pour tout $A$-module $M$ de présentation finie, il existe un indice $\alpha$
et un $A_\alpha$-module de présentation finie $M_\alpha$ tel que $M$ soit
isomorphe à $M_\alpha\otimes_{A_\alpha}A$.
\end{lemma}

En effet, on peut supposer que $M=A^r/P$, où $P$ est un sous-$A$-module de
type fini de $A^r$. Comme $A^r=\varinjlim(A_\alpha^r)$ et que $P$ est de type
fini, il existe un $\alpha$ et un sous-$A_\alpha$-module de type fini
$P_\alpha$ de $A_\alpha^r$ tel que $P$ soit l'image canonique de
$P_\alpha\otimes_{A_\alpha}A$; l'exactitude à droite du produit tensoriel
montre donc que $M$ est isomorphe à
$(A_\alpha^r/P_\alpha)\otimes_{A_\alpha}A$.

Ce lemme étant établi, $M$ et $N$ sont de présentation finie, puisque $A$ est
noethérien; écrivons donc
$M=M_\alpha\otimes_{A_\alpha}A$,
$N=N_\alpha\otimes_{A_\alpha}A$; en vertu de l'hypothèse, $A$ est un
$A_\alpha$-module plat $(0_{\mathrm I}, 6.1.2)$, donc on a
\[
\operatorname{Tor}^A_i(M,N)
=\operatorname{Tor}^{A_\alpha}_i(M_\alpha,N_\alpha)
 \otimes_{A_\alpha}A
\]
à un isomorphisme près (III, 6.3.9); comme par hypothèse $A_\alpha$ est
régulier, il existe $i_0$ tel que
$\operatorname{Tor}^{A_\alpha}_i(M_\alpha,N_\alpha)=0$ pour $i\geq i_0$.
C.Q.F.D.

\begin{corollary}[5.13.8]
\label{IV.2.s5.13.8-fr}
Supposons que les $A_\alpha$ soient noethériens et vérifient la propriété
$(R_k)$ pour un entier $k\geq0$. Supposons en outre que pour
$\alpha\leq\beta$, $\varphi_{\beta\alpha}$ fasse de $A_\beta$ un
$A_\alpha$-module plat, que $A=\varinjlim A_\alpha$ soit noethérien et que,
pour tout $\alpha$, ${}^a\varphi_\alpha$ soit un homéomorphisme de
$\operatorname{Spec}(A)$ sur $\operatorname{Spec}(A_\alpha)$. Alors $A$
vérifie $(R_k)$.
\end{corollary}

Soit en effet $\mathfrak p$ un idéal premier de $A$ et, pour tout $\alpha$,
posons $\mathfrak p_\alpha=\varphi_\alpha^{-1}(\mathfrak p)$, de sorte que
$\mathfrak p=\varinjlim\mathfrak p_\alpha$ et
$A_{\mathfrak p}=\varinjlim(A_\alpha)_{\mathfrak p_\alpha}$ (5.13.3). Il
résulte de (I, 2.4.2) et de l'hypothèse que le morphisme
\[
\operatorname{Spec}(A_{\mathfrak p})
\longrightarrow\operatorname{Spec}((A_\alpha)_{\mathfrak p_\alpha})
\]
est un homéomorphisme surjectif pour tout $\alpha$; on a donc
$\dim(A_{\mathfrak p})=\dim((A_\alpha)_{\mathfrak p_\alpha})$ pour tout
$\alpha$. Supposons alors que $\dim(A_{\mathfrak p})\leq k$; on a donc
$\dim((A_\alpha)_{\mathfrak p_\alpha})\leq k$, et l'hypothèse entraîne que
$(A_\alpha)_{\mathfrak p_\alpha}$ est un anneau régulier; comme en outre, pour
$\alpha\leq\beta$, $(A_\beta)_{\mathfrak p_\beta}$ est un
$(A_\alpha)_{\mathfrak p_\alpha}$-module plat $(0_{\mathrm I}, 6.3.2)$, on
conclut de (5.13.7) que $A_{\mathfrak p}$ est régulier, ce qui démontre le
corollaire.

\section{Morphismes plats de préschémas localement noethériens}

Soient $X$, $Y$ deux préschémas localement noethériens, $f:X\to Y$ un
morphisme. Pour tout $y\in Y$, la fibre
$f^{-1}(y)=X\times_Y\operatorname{Spec}(k(y))$ est aussi un préschéma
localement noethérien, car il suffit de le voir lorsque
$Y=\operatorname{Spec}(A)$ et $X=\operatorname{Spec}(B)$ sont des spectres
d'anneaux noethériens, et alors $B\otimes_A k(y)$ est anneau de fractions de
l'anneau quotient $B/\mathfrak j_yB$, donc est noethérien. Nous nous proposons
en premier lieu dans ce paragraphe de mettre en relation les propriétés de
$X$, de $Y$ et des fibres $f^{-1}(y)$ (où $y$ parcourt $Y$), sous l'hypothèse
que le morphisme $f$ est plat. Les questions traitées se ramèneront à l'étude
des relations entre les anneaux locaux noethériens
$\mathcal O_y$, $\mathcal O_x$ et
$\mathcal O_x\otimes_{\mathcal O_y}k(y)$, l'homomorphisme
$\mathcal O_y\to\mathcal O_x$ étant local et faisant de $\mathcal O_x$ un
$\mathcal O_y$-module plat. Dans les n\textsuperscript{os} 6.11 à 6.13 (assez
séparés du reste du paragraphe, de par leur nature « absolue » plutôt que
« relative »), nous appliquons certains des résultats précédents pour trouver
des

\oldpage[IV-2]{135}

critères permettant d'affirmer que le \emph{lieu singulier} (ou certains
ensembles analogues) de certains préschémas sont des ensembles fermés; ces
critères joueront un rôle important au §~7.

\subsection{Platitude et dimension}

\begin{proposition}[6.1.1]
\label{IV.2.s6.1.1-fr}
Soient $A$, $B$ deux anneaux locaux noethériens, $\mathfrak m$ l'idéal maximal
de $A$, $k=A/\mathfrak m$ son corps résiduel,
$\varphi:A\to B$ un homomorphisme local. On suppose que pour tout idéal
premier $\mathfrak p$ de $A$, distinct de $\mathfrak m$, et pour tout élément
minimal $\mathfrak q$ de l'ensemble des idéaux premiers de $B$ contenant
$\mathfrak pB$, $\varphi^{-1}(\mathfrak q)$ soit distinct de $\mathfrak m$
(autrement dit, qu'aucune composante irréductible de
$\operatorname{Spec}(B/\mathfrak pB)$ ne soit contenue dans l'image réciproque
du point fermé $\mathfrak m$ de $\operatorname{Spec}(A)$ dans
$\operatorname{Spec}(B)$). Alors on a
\begin{equation}
\dim(B)=\dim(A)+\dim(B\otimes_A k).
\tag{6.1.1.1}\label{IV.2.s6.1.1.1-fr}
\end{equation}
\end{proposition}

Raisonnons par récurrence sur $n=\dim(A)$; l'assertion est triviale pour
$n=0$, $\mathfrak mB$ étant alors contenu dans le nilradical de $B$ puisque
$\mathfrak m$ est le nilradical de $A$. On peut donc supposer $n>0$. Soient
$\mathfrak q_i$ $(1\leq i\leq m)$ les idéaux premiers minimaux de $B$, et soit
$\mathfrak p_i=\varphi^{-1}(\mathfrak q_i)$; pour tout $i$, on a
$\mathfrak p_i\neq\mathfrak m$; sinon il existerait (puisque $n>0$) un idéal
premier $\mathfrak p\neq\mathfrak m$ contenu dans $\mathfrak p_i$ et comme
$\mathfrak q_i$ est minimal parmi les idéaux premiers de $B$ contenant
$\mathfrak pB$, on aboutirait à une contradiction avec l'hypothèse. Par suite
$\mathfrak m$ est distinct de la réunion des $\mathfrak p_i$ et des idéaux
minimaux $\mathfrak p'_j$ $(1\leq j\leq r)$ de $A$ (Bourbaki,
\emph{Alg. comm.}, chap.~II, §~1, n\textsuperscript{o}~1, prop.~2) et il
existe un $a\in\mathfrak m$ n'appartenant à aucun des $\mathfrak p_i$ ni des
$\mathfrak p'_j$. Posons $A'=A/aA$, $B'=B/aB$; on a (0, 16.3.4)
\[
\dim(A')=\dim(A)-1,\qquad \dim(B')=\dim(B)-1
\]
par construction de $a$, et d'autre part
$B'\otimes_{A'}k=B\otimes_Ak=B/\mathfrak mB$, donc
\[
\dim(B'\otimes_{A'}k)=\dim(B\otimes_Ak)
\]
et il suffit par suite de démontrer (6.1.1.1) pour $A'$ et $B'$; or, en vertu
de la correspondance biunivoque entre idéaux de $A$ (resp. $B$) contenant $a$
et idéaux de $A'$ (resp. $B'$), les hypothèses de l'énoncé sont aussi
valables pour $A'$ et $B'$; on peut donc appliquer l'hypothèse de récurrence,
ce qui achève la démonstration.

\begin{corollary}[6.1.2]
\label{IV.2.s6.1.2-fr}
Soient $A$, $B$ deux anneaux locaux noethériens, $k$ le corps résiduel de
$A$, $\varphi:A\to B$ un homomorphisme local, $M\neq0$ un $A$-module de type
fini, $N\neq0$ un $B$-module de type fini. Si $N$ est un $A$-module plat
(resp. si $B$ est un $A$-module plat), on a
\begin{equation}
\dim_B(M\otimes_A N)=\dim_A(M)+
\dim_{B\otimes_Ak}(N\otimes_Ak)
\tag{6.1.2.1}\label{IV.2.s6.1.2.1-fr}
\end{equation}
(resp. (6.1.1.1)).
\end{corollary}

Il suffit de prouver l'assertion concernant $N$; d'autre part, si
$\mathfrak b$ est l'annulateur de $N$, on peut remplacer $B$ par
$B/\mathfrak b$, et par suite supposer que
$\operatorname{Supp}(N)=\operatorname{Spec}(B)$; l'hypothèse signifie alors
que le morphisme $f:\operatorname{Spec}(B)\to\operatorname{Spec}(A)$
correspondant à $\varphi$ est \emph{quasi-plat} (2.3.3); on en conclut
(2.3.4) que pour tout idéal premier $\mathfrak p\neq\mathfrak m$ de $A$,

\oldpage[IV-2]{136}

toute composante irréductible de $f^{-1}(V(\mathfrak p))$ domine
$V(\mathfrak p)$, et par suite la condition de (6.1.1) est satisfaite, d'où la
conclusion.

\begin{corollary}[6.1.3]
\EGAReaderTarget{ega4.c04.cor.6-1-3}{IV.2.s6.1.3-fr}
Soient $A$, $B$ deux anneaux locaux noethériens, $\mathfrak m$ l'idéal maximal
de $A$, $k$ son corps résiduel, $\varphi:A\to B$ un homomorphisme local.
Supposons que $\dim(B\otimes_Ak)=0$ (c'est-à-dire (0, 16.2.1) que
$\mathfrak mB$ soit un idéal de définition de $B$). Alors on a
$\dim(B)\leq\dim(A)$. Si en outre il existe un $B$-module de type fini $N$ qui
soit un $A$-module plat et ait un support égal à $\operatorname{Spec}(B)$
(ce qui a lieu en particulier si $B$ est un $A$-module plat), on a
$\dim(B)=\dim(A)$.
\end{corollary}

La première assertion résulte de (5.5.1) et la seconde de (6.1.2).

\begin{corollary}[6.1.4]
\label{IV.2.s6.1.4-fr}
Soient $X$, $Y$ deux préschémas localement noethériens, $f:X\to Y$ un
morphisme surjectif. Pour toute partie fermée $Z$ de $Y$, on a alors
\begin{equation}
\operatorname{codim}(f^{-1}(Z),X)\leq\operatorname{codim}(Z,Y).
\tag{6.1.4.1}\label{IV.2.s6.1.4.1-fr}
\end{equation}
En outre, si $f$ est quasi-plat (2.3.3) les deux membres de (6.1.4.1) sont
égaux.
\end{corollary}

En effet, si $y$ est un point maximal de $Z$ et $x$ un point maximal de la
fibre $f^{-1}(y)$, on a
$\dim(\mathcal O_{X,x})\leq\dim(\mathcal O_{Y,y})$ en vertu de (5.5.2);
l'inégalité (6.1.4.1) résulte alors de (5.1.2.1) et de (5.1.3.1) et du fait
que $f$ est surjectif. Si en outre $f$ est quasi-plat, on sait (2.3.4) que
toute composante irréductible de $f^{-1}(Z)$ domine une composante
irréductible de $Z$; les points maximaux de $f^{-1}(Z)$ sont donc les points
maximaux des fibres $f^{-1}(y)$, où $y$ parcourt l'ensemble des points maximaux
de $Z$ $(0_{\mathrm I}, 2.1.8)$, et en chacun de ces points maximaux $x$ on a
$\dim(\mathcal O_x\otimes_{\mathcal O_y}k(x))=0$; comme en outre
$\mathcal O_x$ est un $\mathcal O_y$-module plat, on a
$\dim(\mathcal O_x)=\dim(\mathcal O_y)$ en vertu de (6.1.1), d'où l'égalité
dans (6.1.4.1), compte tenu de ce que $f$ est surjectif.

La proposition (6.1.1) admet la réciproque partielle suivante:

\begin{proposition}[6.1.5]
\EGAReaderTarget{ega4.c04.prop.6-1-5}{IV.2.s6.1.5-fr}
Soient $A$, $B$ deux anneaux locaux noethériens, $k$ le corps résiduel de
$A$, $\varphi:A\to B$ un homomorphisme local, $M$ un $B$-module de type fini.
On suppose que:
\begin{enumerate}
\item[1\textsuperscript{o}] $A$ est un anneau régulier.
\item[2\textsuperscript{o}] $M$ est un $B$-module de Cohen-Macaulay.
\item[3\textsuperscript{o}] On a
$\dim_B(M)=\dim(A)+\dim_{B\otimes_Ak}(M\otimes_Ak)$.
\end{enumerate}
Alors $M$ est un $A$-module plat.
\end{proposition}

Procédons par récurrence sur $n=\dim(A)$. Si $n=0$, $A$ est un corps
puisqu'il est régulier (0, 17.1.4) et l'assertion est triviale. Supposons
$n>0$; soit $\mathfrak m$ l'idéal maximal de $A$, et soit
$x\in\mathfrak m$ un élément n'appartenant pas à $\mathfrak m^2$; on sait
alors (0, 17.1.8) que $A'=A/xA$ est régulier et
$\dim(A')=\dim(A)-1$. Posons
\[
B'=B/xB,\qquad M'=M/xM=M\otimes_AA';
\]
on a donc
\[
B'\otimes_{A'}k=B\otimes_Ak,\qquad
M'\otimes_{A'}k=M\otimes_Ak
\]
et en vertu de (5.5.1.2) on a
\[
\dim_{B'}(M')\leq\dim(A')+
\dim_{B'\otimes_{A'}k}(M'\otimes_{A'}k);
\]

\oldpage[IV-2]{137}

on conclut donc de (0, 16.3.4) que l'on a
\[
\dim_B(M)\leq\dim_{B'}(M')+1
\leq\bigl(\dim(A')+
\dim_{B\otimes_Ak}(M\otimes_Ak)\bigr)+1
=\dim(A)+\dim_{B\otimes_Ak}(M\otimes_Ak).
\]

Comme par hypothèse les membres extrêmes de ces inégalités sont égaux, on
a nécessairement: (i) $\dim_{B'}(M')=\dim_B(M)-1$, et comme $M$ est par
hypothèse un $B$-module de Cohen-Macaulay, cela signifie que $x$ est
$M$-régulier (0, 16.1.9 et 16.5.5); (ii)
$\dim_{B'}(M')=\dim(A')+
\dim_{B'\otimes_{A'}k}(M'\otimes_{A'}k)$; comme $M'$ est un $B'$-module de
Cohen-Macaulay en vertu de (i) et de (0, 16.1.9 et 16.5.5) et que $A'$ est
régulier, l'hypothèse de récurrence prouve que $M'$ est un $A'$-module plat;
on déduit alors de $(0_{\mathrm{III}}, 10.2.7)$ que $M$ est un $A$-module
plat, puisque $x$ est $M$-régulier par (i).

\subsection{Platitude et dimension projective}

\begin{proposition}[6.2.1]
\label{IV.2.s6.2.1-fr}
\begin{enumerate}
\item[(i)] Soient $A$, $B$ deux anneaux, $\varphi:A\to B$ un homomorphisme tel
que $B$ soit un $A$-module plat. Alors, pour tout $A$-module $M$, on a
\begin{equation}
\dim.\operatorname{proj}_B(M\otimes_AB)
\leq\dim.\operatorname{proj}_A(M).
\tag{6.2.1.1}\label{IV.2.s6.2.1.1-fr}
\end{equation}
\item[(ii)] Supposons en outre que $A$ soit un anneau noethérien, $B$ un
$A$-module fidèlement plat et $M$ un $A$-module de type fini; alors
\begin{equation}
\dim.\operatorname{proj}_B(M\otimes_AB)
=\dim.\operatorname{proj}_A(M).
\tag{6.2.1.2}\label{IV.2.s6.2.1.2-fr}
\end{equation}
\end{enumerate}
\end{proposition}

(i) On peut se borner au cas où
$n=\dim.\operatorname{proj}_A(M)$ est finie; il existe donc une résolution
gauche
\[
0\to P_n\to P_{n-1}\to\cdots\to P_0\to M\to0
\]
de $M$ par des $A$-modules projectifs; comme $B$ est un $A$-module plat, la
suite
\[
0\to P_n\otimes_AB\to P_{n-1}\otimes_AB\to\cdots\to
P_0\otimes_AB\to M\otimes_AB\to0
\]
est exacte; d'autre part les $P_i\otimes_AB$ sont des $B$-modules projectifs;
d'où la conclusion.

(ii) Supposons que $\dim.\operatorname{proj}_B(M\otimes_AB)=m$, et considérons
une suite exacte
\[
0\to R\to P_{m-1}\to P_{m-2}\to\cdots\to P_0\to M\to0
\]
où les $P_i$ $(0\leq i\leq m-1)$ sont des $A$-modules projectifs de type fini;
comme $A$ est noethérien, $R$ est aussi un $A$-module de type fini. Comme $B$
est un $A$-module plat, on a aussi une suite exacte
\[
0\to R\otimes_AB\to P_{m-1}\otimes_AB\to\cdots\to
P_0\otimes_AB\to M\otimes_AB\to0
\]
et l'hypothèse sur $M\otimes_AB$ entraîne que $R\otimes_AB$ est un $B$-module
projectif (0, 17.2.1). Comme $B$ est un $A$-module fidèlement plat, on en
conclut que $R$ est un $A$-module projectif (Bourbaki, \emph{Alg. comm.},
chap.~I\textsuperscript{er}, §~3, n\textsuperscript{o}~6, prop.~12); donc
$\dim.\operatorname{proj}_A(M)\leq m$, ce qui achève la démonstration.

\begin{corollary}[6.2.2]
\label{IV.2.s6.2.2-fr}
Soient $f:X\to Y$ un morphisme plat de préschémas, $\mathscr E$ un
$\mathcal O_Y$-Module quasi-cohérent. Si
$\dim.\operatorname{proj}(\mathscr E)\leq n$ (0, 17.2.14), on a
$\dim.\operatorname{proj}(f^*(\mathscr E))\leq n$.
\end{corollary}

\oldpage[IV-2]{138}

\begin{proposition}[6.2.3]
\EGAReaderTarget{ega4.c04.prop.6-2-3}{IV.2.s6.2.3-fr}
Soient $A$, $B$ deux anneaux locaux noethériens, $k$ le corps résiduel de
$A$, $\varphi:A\to B$ un homomorphisme local, $N$ un $B$-module de type fini.
On suppose que $B$ et $N$ sont des $A$-modules plats. Alors on a
\begin{equation}
\dim.\operatorname{proj}_B(N)
=\dim.\operatorname{proj}_{B\otimes_Ak}(N\otimes_Ak).
\tag{6.2.3.1}\label{IV.2.s6.2.3.1-fr}
\end{equation}
\end{proposition}

Considérons en effet une résolution gauche de $N$ par des $B$-modules libres
de type fini
\[
\cdots\to L_i\to L_{i-1}\to\cdots\to L_0\to N\to0.
\]
Comme les $L_i$ et $N$ sont des $A$-modules plats ($B$ étant un $A$-module
plat) il résulte de (2.1.10) que les $Z_i(L_\bullet)$ sont des $A$-modules
plats, que $L_\bullet\otimes_Ak$ est une résolution gauche du
$(B\otimes_Ak)$-module $N\otimes_Ak$ et que l'on a
$Z_i(L_\bullet\otimes_Ak)=Z_i(L_\bullet)\otimes_Ak$ pour tout $i$. Notons que
puisque $B$ est noethérien, les $Z_i(L_\bullet)$ sont des $B$-modules de type
fini, donc les $Z_i(L_\bullet\otimes_Ak)$ sont des $(B\otimes_Ak)$-modules de
type fini. Or, dire qu'un $B$-module (resp. un $(B\otimes_Ak)$-module) de type
fini est plat équivaut à dire qu'il est libre ou encore projectif
$(0_{\mathrm{III}}, 10.1.3)$. Comme les $Z_i(L_\bullet)$ sont des $A$-modules
plats, il résulte d'autre part de $(0_{\mathrm{III}}, 10.2.5)$ (où l'on fait
$C=B$) que, pour que $Z_i(L_\bullet)$ soit un $B$-module plat, il faut et il
suffit que $Z_i(L_\bullet)\otimes_Ak$ soit un $(B\otimes_Ak)$-module plat. Le
plus petit entier $n$ tel que $Z_{n-1}(L_\bullet)$ soit un $B$-module libre est
donc aussi le plus petit entier tel que $Z_{n-1}(L_\bullet\otimes_Ak)$ soit un
$(B\otimes_Ak)$-module libre, ce qui démontre la proposition (0, 17.2.1).

\subsection{Platitude et profondeur}

\begin{proposition}[6.3.1]
\EGAReaderTarget{ega4.c04.prop.6-3-1}{IV.2.s6.3.1-fr}
Soient $A$, $B$ deux anneaux locaux noethériens, $k$ le corps résiduel de
$A$, $\varphi:A\to B$ un homomorphisme local, $M$ un $A$-module de type fini,
$N$ un $B$-module de type fini. Si $N$ est un $A$-module plat $\neq0$, on a
\begin{equation}
\operatorname{prof}_B(M\otimes_AN)=\operatorname{prof}_A(M)
+\operatorname{prof}_{B\otimes_Ak}(N\otimes_Ak).
\tag{6.3.1.1}\label{IV.2.s6.3.1.1-fr}
\end{equation}
\end{proposition}

On peut se borner au cas où $M\neq0$, sans quoi les deux membres de (6.3.1.1)
sont tous deux égaux à $+\infty$. Nous procéderons par récurrence sur l'entier
$n$ égal au second membre de (6.3.1.1) (qui est fini en vertu des hypothèses,
de (0, 16.4.6.2) et du lemme de Nakayama appliqué à $N$). Supposons d'abord
$n=0$; alors, si $\mathfrak m$ et $\mathfrak n$ désignent les idéaux maximaux
de $A$ et $B$ respectivement, $\mathfrak m$ est un idéal premier associé à $M$
et $\mathfrak n/\mathfrak mB$ un idéal premier de
$B/\mathfrak mB=B\otimes_Ak$ associé à $N\otimes_Ak$ (0, 16.4.6, (i)). On
conclut alors de (3.3.1) que $\mathfrak n$ est un idéal premier de $B$ associé
à $M\otimes_AN$, donc (0, 16.4.6, (i)) que
$\operatorname{prof}_B(M\otimes_AN)=0$. Supposons donc $n>0$, et distinguons
deux cas:

$a)$ Supposons $\operatorname{prof}_A(M)>0$. Soit $x\in\mathfrak m$ un
élément $M$-régulier, et posons
\[
A'=A/xA,\qquad B'=B/xB,\qquad M'=M/xM,\qquad N'=N/xN;
\]
on a
\[
B'\otimes_{A'}k=B\otimes_Ak,\qquad N'\otimes_{A'}k=N\otimes_Ak
\]
puisque $x\in\mathfrak m$, et
\[
M'\otimes_{A'}N'=(M\otimes_AN)/x(M\otimes_AN);
\]

\oldpage[IV-2]{139}

en outre, comme $N$ est un $A$-module plat, l'hypothèse que $x$ est
$M$-régulier entraîne que $x$ est aussi $(M\otimes_AN)$-régulier
$(0_{\mathrm I}, 6.1.1)$. On a par suite (0, 16.4.6, (ii) et 16.4.8)
\[
\operatorname{prof}_{A'}(M')=\operatorname{prof}_A(M)-1,
\qquad
\operatorname{prof}_{B'}(M'\otimes_{A'}N')
=\operatorname{prof}_B(M\otimes_AN)-1
\]
et
\[
\operatorname{prof}_{B'\otimes_{A'}k}(N'\otimes_{A'}k)
=\operatorname{prof}_{B\otimes_Ak}(N\otimes_Ak).
\]

L'égalité (6.3.1.1) est donc conséquence de la même relation pour $A'$,
$B'$, $M'$ et $N'$; mais comme $N$ est un $A$-module plat,
$N'=N\otimes_AA'$ est un $A'$-module plat $(0_{\mathrm I}, 6.2.1)$; on peut
par suite appliquer l'hypothèse de récurrence, ce qui démontre (6.3.1.1) dans
ce cas.

$b)$ Supposons $\operatorname{prof}_{B\otimes_Ak}(N\otimes_Ak)>0$.
Considérons un élément $y\in\mathfrak m$ qui soit
$(N\otimes_Ak)$-régulier; il résulte de $(0_{\mathrm{III}}, 10.2.4)$ que $y$
est alors $N$-régulier et que
\[
N'=N/yN
\]
est un $A$-module plat, puisque $N$ est supposé être un $A$-module plat.
Appliquant $(0_{\mathrm I}, 6.1.2)$ à la suite exacte de $A$-modules
\[
0\to N\xrightarrow{y}N\to N'\to0
\]
on en conclut des isomorphismes
\[
(M\otimes_AN)/y(M\otimes_AN)\xrightarrow{\sim}M\otimes_AN'
\quad\text{et}\quad
N'\otimes_Ak\xrightarrow{\sim}(N\otimes_Ak)/y(N\otimes_Ak)
\]
et en outre que $y$ est $(M\otimes_AN)$-régulier. On a par suite
\[
\operatorname{prof}_{B\otimes_Ak}(N'\otimes_Ak)
=\operatorname{prof}_{B\otimes_Ak}(N\otimes_Ak)-1
\]
et
\[
\operatorname{prof}_B(M\otimes_AN')
=\operatorname{prof}_B(M\otimes_AN)-1;
\]
l'hypothèse de récurrence montre que la relation (6.3.1.1) est valable pour
$A$, $B$, $M$ et $N'$, et d'après ce qui précède, on en déduit (6.3.1.1)
pour $A$, $B$, $M$ et $N$.

\begin{corollary}[6.3.2]
\label{IV.2.s6.3.2-fr}
Sous les hypothèses de (6.3.1), supposons en outre $M\neq0$ et $N\neq0$;
alors on a
\begin{equation}
\operatorname{coprof}_B(M\otimes_AN)=\operatorname{coprof}_A(M)
+\operatorname{coprof}_{B\otimes_Ak}(N\otimes_Ak).
\tag{6.3.2.1}\label{IV.2.s6.3.2.1-fr}
\end{equation}
\end{corollary}

Cela résulte aussitôt de (6.3.1.1) et (6.1.2) et de la définition de la
coprofondeur (0, 16.4.9).

\begin{corollary}[6.3.3]
\label{IV.2.s6.3.3-fr}
Sous les hypothèses de (6.3.1), supposons en outre $M\neq0$ et $N\neq0$; pour
que $M\otimes_AN$ soit un $B$-module de Cohen-Macaulay, il faut et il suffit
que $M$ soit un $A$-module de Cohen-Macaulay et que $N\otimes_Ak$ soit un
$(B\otimes_Ak)$-module de Cohen-Macaulay.
\end{corollary}

Cela résulte du corollaire (6.3.2) et de la définition des modules de
Cohen-Macaulay par le fait que leur coprofondeur est nulle (0, 16.5.1), tenant
compte de ce que la condition $N\neq0$ est équivalente à $N\otimes_Ak\neq0$ en
vertu du lemme de Nakayama.

\begin{corollary}[6.3.4]
\label{IV.2.s6.3.4-fr}
Sous les hypothèses de (6.3.1), supposons en outre que $N\otimes_Ak$ soit un
$B$-module de longueur finie; alors on a
\begin{equation}
\operatorname{prof}_B(M\otimes_AN)=\operatorname{prof}_A(M).
\tag{6.3.4.1}\label{IV.2.s6.3.4.1-fr}
\end{equation}

\oldpage[IV-2]{140}

Si de plus $M\neq0$ et $N\neq0$, on a
\begin{equation}
\operatorname{coprof}_B(M\otimes_AN)=\operatorname{coprof}_A(M).
\tag{6.3.4.2}\label{IV.2.s6.3.4.2-fr}
\end{equation}
\end{corollary}

Il revient en effet au même de dire que $N\otimes_Ak$ est un $B$-module de
longueur finie ou un $(B\otimes_Ak)$-module de longueur finie, et l'on sait
(0, 16.2.3) que les modules de longueur finie $\neq0$ sont de dimension $0$ et
de profondeur $0$.

Le corollaire (6.3.4) sera applicable en particulier lorsque $B\otimes_Ak$ est
un anneau de longueur finie, c'est-à-dire lorsque $\mathfrak mB$ est un idéal
de définition de l'anneau $B$.

\begin{corollary}[6.3.5]
\label{IV.2.s6.3.5-fr}
Soient $X$, $Y$ deux préschémas localement noethériens, $f:X\to Y$ un
morphisme plat.
\begin{enumerate}
\item[(i)] Si, en un point $x\in X$, on a
$\operatorname{coprof}(\mathcal O_x)\leq n$, alors, en posant $y=f(x)$, on a
$\operatorname{coprof}(\mathcal O_y)\leq n$ et
$\operatorname{coprof}(\mathcal O_x\otimes_{\mathcal O_y}k(y))\leq n$; en
particulier, si $\mathcal O_x$ est un anneau de Cohen-Macaulay, il en est de
même de $\mathcal O_y$ et de
$\mathcal O_x\otimes_{\mathcal O_y}k(y)$.
\item[(ii)] Supposons inversement que
$\mathcal O_x\otimes_{\mathcal O_y}k(y)$ soit un anneau de Cohen-Macaulay.
Alors, si $\operatorname{coprof}(\mathcal O_y)\leq n$ (resp. si
$\mathcal O_y$ est un anneau de Cohen-Macaulay), on a
$\operatorname{coprof}(\mathcal O_x)\leq n$ (resp. $\mathcal O_x$ est un
anneau de Cohen-Macaulay).
\end{enumerate}
\end{corollary}

Il suffit d'appliquer (6.3.2) pour $M=A=\mathcal O_y$ et
$N=B=\mathcal O_x$.

\begin{corollary}[6.3.6]
\label{IV.2.s6.3.6-fr}
Soit $A$ un anneau de Cohen-Macaulay. Alors
$A'=A[T_1,\ldots,T_n]$ ($T_i$ indéterminées) est un anneau de
Cohen-Macaulay.
\end{corollary}

En effet, si l'on pose $Y=\operatorname{Spec}(A)$,
$X=\operatorname{Spec}(A')$ et si $f:X\to Y$ est le morphisme canonique, $f$
est plat (puisque $A'$ est un $A$-module libre (2.1.2)) et pour tout $y\in Y$,
$k(y)[T_1,\ldots,T_n]$ est un anneau régulier (0, 17.3.7), donc \emph{a
fortiori} de Cohen-Macaulay; il suffit donc d'appliquer (6.3.5).

\begin{corollary}[6.3.7]
\label{IV.2.s6.3.7-fr}
Tout quotient d'un anneau de Cohen-Macaulay est universellement caténaire.
\end{corollary}

Cela résulte de (6.3.6), de (5.6.1) et de (0, 16.5.12).

\begin{proposition}[6.3.8]
\label{IV.2.s6.3.8-fr}
Soit $A$ un anneau local noethérien; supposons qu'il existe un $A$-module $M$
de type fini qui soit un $A$-module de Cohen-Macaulay et dont le support soit
égal à $\operatorname{Spec}(A)$. Soit $B$ un anneau quotient de $A$, et soit
$f:\operatorname{Spec}(\widehat B)\to\operatorname{Spec}(B)$ le morphisme
canonique; alors, pour tout $x\in\operatorname{Spec}(B)$, $f^{-1}(x)$ est un
schéma de Cohen-Macaulay.
\end{proposition}

Comme $B$ est un $A$-module de type fini, on a
$\widehat B=B\otimes_A\widehat A$ $(0_{\mathrm I}, 7.3.3)$, donc $f$ est la
restriction à $\operatorname{Spec}(\widehat B)$ du morphisme canonique
$\operatorname{Spec}(\widehat A)\to\operatorname{Spec}(A)$, et les fibres de
ces deux morphismes en un point de $\operatorname{Spec}(B)$ sont donc les
mêmes. On est par suite ramené à démontrer la proposition lorsque $B=A$.
Soit donc $\mathfrak p=\mathfrak j_x$ un idéal premier de $A$; comme par
hypothèse $\mathfrak p\in\operatorname{Supp}(M)$, on sait (0, 16.5.6 et
16.5.9) qu'il existe une suite $M$-régulière $(t_i)_{1\leq i\leq r}$
d'éléments de $\mathfrak p$ telle que
\[
N=M/\Bigl(\sum_{i=1}^{r}t_iM\Bigr)
\]
soit un $A$-module de Cohen-Macaulay, $\mathfrak p$ un idéal premier associé
minimal de $N$ et
$\dim(N)=\dim(M/\mathfrak pM)=\dim(A/\mathfrak p)$. Le même raisonnement
qu'au début montre que l'on peut remplacer $M$ par $N$ et $A$ par
$A/(\sum_{i=1}^{r}t_iA)$, et par suite supposer que $\mathfrak p$ est un

\oldpage[IV-2]{141}

idéal premier minimal de $A$. Posons pour simplifier
$A'=\widehat A$, $M'=\widehat M=M\otimes_AA'$ $(0_{\mathrm I}, 7.3.3)$; on
sait (0, 16.5.2) que $M'$ est un $A'$-module de Cohen-Macaulay, ce qui entraîne
que pour tout idéal premier $\mathfrak p'$ de $A'$,
$M'_{\mathfrak p'}$ est un $A'_{\mathfrak p'}$-module de Cohen-Macaulay
(0, 16.5.10). Compte tenu de (I, 3.6.5), on voit que si l'on pose
$A''=A'\otimes_AA_{\mathfrak p}$ et
$M''=M'\otimes_AA_{\mathfrak p}$, $M''$ est un $A''$-module de
Cohen-Macaulay. Pour tout idéal premier $\mathfrak p''$ de $A''$ au-dessus
de $\mathfrak p$,
$M''_{\mathfrak p''}=M_{\mathfrak p}\otimes_{A_{\mathfrak p}}
A''_{\mathfrak p''}$ est donc un $A''_{\mathfrak p''}$-module de
Cohen-Macaulay (0, 16.5.10); d'autre part, $M_{\mathfrak p}$ est par hypothèse
un $A_{\mathfrak p}$-module de Cohen-Macaulay et
$A''_{\mathfrak p''}$ est un $A_{\mathfrak p}$-module plat puisque $A'$ est
un $A$-module plat $(0_{\mathrm I}, 7.3.3\text{ et }6.3.2)$. On conclut de (6.3.3)
que, si $k$ est le corps résiduel de $A_{\mathfrak p}$,
$k\otimes_{A_{\mathfrak p}}A''_{\mathfrak p''}$ est un anneau de
Cohen-Macaulay. Mais puisque $A_{\mathfrak p}$ est de dimension $0$, les
idéaux premiers de $A''$ correspondent biunivoquement à ceux de
$k\otimes_{A_{\mathfrak p}}A''$ (I, 3.5.7), et si $\mathfrak q''$ est l'idéal
premier de $k\otimes_{A_{\mathfrak p}}A''$ correspondant à
$\mathfrak p''$, les anneaux locaux
$(k\otimes_{A_{\mathfrak p}}A'')_{\mathfrak q''}$ et
$k\otimes_{A_{\mathfrak p}}A''_{\mathfrak p''}$ sont isomorphes. Par suite
(0, 16.5.13), l'anneau $k\otimes_{A_{\mathfrak p}}A''$ est un anneau de
Cohen-Macaulay. C.Q.F.D.

\subsection{\texorpdfstring{Platitude et propriété $(S_k)$}{Platitude et propriété (S\string_k)}}

\begin{proposition}[6.4.1]
\label{IV.2.s6.4.1-fr}
Soient $X$, $Y$ deux préschémas localement noethériens, $f:X\to Y$ un
morphisme, $\mathscr E$ un $\mathcal O_Y$-Module cohérent,
$\mathscr F$ un $\mathcal O_X$-Module cohérent et $f$-plat.
\begin{enumerate}
\item[(i)] Soit $x\in\operatorname{Supp}(\mathscr F)$ tel que
$\mathscr E\otimes_Y\mathscr F$ ait la propriété $(S_k)$ au point $x$; alors
$\mathscr E$ a la propriété $(S_k)$ au point $y=f(x)$.
\item[(ii)] Supposons que pour tout
$y\in f(\operatorname{Supp}(\mathscr F))$,
$\mathscr F_y=\mathscr F\otimes_{\mathcal O_y}k(y)$ ait la propriété
$(S_k)$; alors, si pour un point
$y\in f(\operatorname{Supp}(\mathscr F))$, $\mathscr E$ a la propriété
$(S_k)$ au point $y$, $\mathscr E\otimes_Y\mathscr F$ a la propriété
$(S_k)$ en tout point de $f^{-1}(y)$.
\end{enumerate}
\end{proposition}

(i) On sait $(\mathbf{Err}_{\mathrm{III}},30)$ qu'il existe un sous-préschéma
fermé $X'$ de $X$ ayant pour espace sous-jacent
$\operatorname{Supp}(\mathscr F)$ et un $\mathcal O_{X'}$-Module cohérent
$\mathscr F'$ tels que $\mathscr F=j_*(\mathscr F')$, où $j$ est l'injection
canonique $X'\to X$; on peut remplacer $X$ par $X'$ et $\mathscr F$ par
$\mathscr F'$, autrement dit supposer que
$\operatorname{Supp}(\mathscr F)=X$. Soit alors $y'$ une générisation de $y$
dans $Y$; on sait (2.3.4) qu'il y a une générisation $x'$ de $x$ dans $X$
telle que $y'=f(x')$; on peut en outre supposer que $x'$ est point générique
d'une composante irréductible de $f^{-1}(y')$; on a par suite
$\dim(\mathscr F_{x'}\otimes_{\mathcal O_{y'}}k(y'))=0$. En vertu de (6.1.2),
on a donc, en posant $\mathscr G=\mathscr E\otimes_Y\mathscr F$,
\[
\dim(\mathscr G_{x'})=\dim(\mathscr E_{y'})
\]
et en vertu de (6.3.1)
\[
\operatorname{prof}(\mathscr G_{x'})=\operatorname{prof}(\mathscr E_{y'})
\]
compte tenu de ce que la profondeur est toujours au plus égale à la
dimension. Par hypothèse, on a
\[
\operatorname{prof}(\mathscr G_{x'})
\geq\inf(k,\dim(\mathscr G_{x'}))
\]
donc
\[
\operatorname{prof}(\mathscr E_{y'})
\geq\inf(k,\dim(\mathscr E_{y'}))
\]
ce qui démontre la première assertion.

\oldpage[IV-2]{142}

(ii) Comme pour toute générisation $x'$ de $x\in f^{-1}(y)$,
$y'=f(x')$ est une générisation de $y$, on peut se borner à vérifier que si
$x\in\operatorname{Supp}(\mathscr G)$ et $f(x)=y$, on a
$\operatorname{prof}(\mathscr G_x)\geq\inf(k,\dim(\mathscr G_x))$; il résulte
de (6.1.2) et (6.3.1) que l'on a
\[
\dim(\mathscr G_x)=\dim(\mathscr E_y)+
\dim(\mathscr F_x\otimes_{\mathcal O_y}k(y))
\]
\[
\operatorname{prof}(\mathscr G_x)=\operatorname{prof}(\mathscr E_y)+
\operatorname{prof}(\mathscr F_x\otimes_{\mathcal O_y}k(y)).
\]

Par hypothèse, on a
\[
\operatorname{prof}(\mathscr E_y)\geq\inf(k,\dim(\mathscr E_y))
\]
\[
\operatorname{prof}(\mathscr F_x\otimes_{\mathcal O_y}k(y))
\geq\inf(k,\dim(\mathscr F_x\otimes_{\mathcal O_y}k(y)))
\]
d'où, en ajoutant membre à membre,
\begin{align*}
\operatorname{prof}(\mathscr G_x)
&\geq\inf(k,\dim(\mathscr E_y))
+\inf(k,\dim(\mathscr F_x\otimes_{\mathcal O_y}k(y)))\geq{}\\
&\geq\inf\bigl(k,\dim(\mathscr E_y)
+\dim(\mathscr F_x\otimes_{\mathcal O_y}k(y))\bigr)
=\inf(k,\dim(\mathscr G_x))
\end{align*}
ce qui prouve (ii).

\begin{corollary}[6.4.2]
\label{IV.2.s6.4.2-fr}
Soient $X$, $Y$ deux préschémas localement noethériens, $f:X\to Y$ un
morphisme plat, $\mathscr E$ un $\mathcal O_Y$-Module cohérent. Supposons que
pour tout $y\in Y$, le préschéma $f^{-1}(y)$ possède la propriété $(S_k)$.
Alors, pour que $f^*(\mathscr E)$ vérifie $(S_k)$ en un point $x$, il faut et
il suffit que $\mathscr E$ possède la propriété $(S_k)$ au point $f(x)$.
\end{corollary}

\begin{remark}[6.4.3]
\label{IV.2.s6.4.3-fr}
Soient $A$, $B$ deux anneaux locaux noethériens, $k$ le corps résiduel de
$A$, $\varphi:A\to B$ un homomorphisme local, $M$ un $B$-module de type fini
qui soit un $A$-module plat, et supposons en outre que le $A$-module $A$ et le
$(B\otimes_Ak)$-module $M\otimes_Ak$ possèdent la propriété $(S_n)$; peut-on
alors conclure que le $B$-module $M$ possède la propriété $(S_n)$? Nous
l'ignorons même lorsque $n=1$, $M=B$ et $B=\widehat A$, autrement dit nous
ignorons si en général le \emph{complété} d'un anneau local noethérien
vérifiant $(S_n)$ vérifie lui aussi $(S_n)$, même pour $n=1$. Posant
$Y=\operatorname{Spec}(A)$, $X=\operatorname{Spec}(B)$ et
$\mathscr F=\widetilde M$, il suffirait, en vertu de (6.4.1, (ii)), de montrer
que pour \emph{tout} $y\in Y$, $\mathscr F_y$ possède la propriété $(S_n)$
(et non seulement lorsque $y$ est le point fermé de $Y$); il suffirait aussi
de montrer que l'ensemble $U$ des $x\in X$ tels que
$\mathscr F_x\otimes_{\mathcal O_{f(x)}}k(f(x))$ vérifie $(S_k)$ est ouvert
dans $X$ (puisqu'il contient le point fermé de $X$ par hypothèse). Nous
démontrerons plus loin (12.1.4) que $U$ est ouvert quand $B$ est un anneau
local d'une $A$-algèbre de type fini et $M=B$. D'autre part, pour
$B=\widehat A$ et $M=B$, nous avons vu dans (6.3.8) que les préschémas
$f^{-1}(y)$ sont des préschémas de Cohen-Macaulay (autrement dit vérifient
$(S_n)$ pour \emph{tout} $n$) lorsque l'on suppose que $A$ est quotient d'un
anneau local de Cohen-Macaulay. On en conclut que lorsque $A$ est un quotient
d'un anneau local de Cohen-Macaulay (ou plus généralement d'un anneau local
vérifiant les hypothèses de (6.3.8)), \emph{pour que $A$ vérifie $(S_n)$, il
faut et il suffit que son complété $\widehat A$ vérifie $(S_n)$}. Il resterait
à voir si cette propriété subsiste sans restriction sur $A$.
\end{remark}

\oldpage[IV-2]{143}

\subsection{\texorpdfstring{Platitude et propriété $(R_k)$}{Platitude et propriété (R\string_k)}}

\begin{proposition}[6.5.1]
\EGAReaderTarget{ega4.c04.prop.6-5-1}{IV.2.s6.5.1-fr}
Soient $A$, $B$ deux anneaux locaux noethériens, $k$ le corps résiduel
de $A$, $\varphi:A\to B$ un homomorphisme local, pour lequel $B$ est un
$A$-module plat. Alors:
\begin{enumerate}
\item[(i)] Si $B$ est régulier, $A$ est régulier.
\item[(ii)] Si $A$ et $B\otimes_Ak$ sont réguliers, $B$ est régulier.
\end{enumerate}
Cette proposition est donnée pour mémoire, ayant déjà été démontrée dans
(0, 17.3.3).
\end{proposition}

\begin{corollary}[6.5.2]
\label{IV.2.s6.5.2-fr}
Soient $X$, $Y$ deux préschémas localement noethériens, $f:X\to Y$
un morphisme plat.
\begin{enumerate}
\item[(i)] Si $X$ est régulier en un point $x$, $Y$ est régulier au point
$f(x)$.
\item[(ii)] Si pour un $y\in f(X)$, $Y$ est régulier au point $y$ et si
$f^{-1}(y)$ est un préschéma régulier en un point $x$, alors $X$ est
régulier au point $x$.
\end{enumerate}
\end{corollary}

\begin{proposition}[6.5.3]
\label{IV.2.s6.5.3-fr}
Soient $X$, $Y$ deux préschémas localement noethériens, $f:X\to Y$
un morphisme plat.
\begin{enumerate}
\item[(i)] Si $X$ possède la propriété $(R_k)$ en un point $x$
(resp. si $X$ possède la propriété $(R_k)$), $Y$ possède la propriété
$(R_k)$ au point $f(x)$ (resp. $Y$ possède la propriété $(R_k)$ en tous
les points de $f(X)$).
\item[(ii)] Supposons que pour tout $y\in f(X)$, le préschéma
$f^{-1}(y)$ possède la propriété $(R_k)$; alors, si pour un point
$y\in f(X)$, $Y$ possède la propriété $(R_k)$ au point $y$, $X$ possède
la propriété $(R_k)$ en tout point de $f^{-1}(y)$.
\end{enumerate}
\end{proposition}

\begin{proof}
\medskip\noindent
\begin{enumerate}
\item[(i)] Posons $y=f(x)$ et soit $y'$ une générisation de $y$; comme
dans la démonstration de (6.4.1), il y a un point générique $x'$ d'une
composante irréductible de $f^{-1}(y')$ qui est une générisation de $x$;
donc on a
\[
\dim(\mathcal O_{x'}\otimes_{\mathcal O_{y'}}k(y'))=0
\]
et, en vertu de l'hypothèse et de (6.1.2),
\[
\dim(\mathcal O_{x'})=\dim(\mathcal O_{y'}).
\]
L'hypothèse entraîne, soit que
$\dim(\mathcal O_{x'})>k$, auquel cas
$\dim(\mathcal O_{y'})>k$, soit que $\mathcal O_{x'}$ est un anneau
régulier, et alors $\mathcal O_{y'}$ est un anneau régulier en vertu de
(6.5.1).

\item[(ii)] Comme, pour toute générisation $x'$ de $x\in f^{-1}(y)$,
$y'=f(x')$ est une générisation de $y$, on peut se borner à vérifier que
si $\dim(\mathcal O_x)\leqslant k$, alors $\mathcal O_x$ est un anneau
régulier. Or, on a, en vertu de (6.1.2),
\[
\dim(\mathcal O_x)=\dim(\mathcal O_y)
+\dim(\mathcal O_x\otimes_{\mathcal O_y}k(y)).
\]
Donc si $\dim(\mathcal O_x)\leqslant k$, on a a fortiori
$\dim(\mathcal O_y)\leqslant k$ et
$\dim(\mathcal O_x\otimes_{\mathcal O_y}k(y))\leqslant k$, et l'hypothèse
entraîne que $\mathcal O_y$ et
$\mathcal O_x\otimes_{\mathcal O_y}k(y)$ sont des anneaux réguliers. On
conclut alors de (6.5.1) que $\mathcal O_x$ est un anneau régulier.
\end{enumerate}
\end{proof}

\begin{corollary}[6.5.4]
\label{IV.2.s6.5.4-fr}
Soient $X$, $Y$ deux préschémas localement noethériens, $f:X\to Y$
un morphisme plat.
\begin{enumerate}
\item[(i)] Si $X$ est normal en un point $x$, $Y$ est normal au point
$f(x)$.
\item[(ii)] Si, pour tout $y\in f(X)$, $f^{-1}(y)$ est un préschéma
normal et si $Y$ est normal en un point $y\in f(X)$, alors $X$ est un
préschéma normal en tout point de $f^{-1}(y)$.
\end{enumerate}
\end{corollary}

Cela résulte aussitôt du critère de normalité de Serre (5.8.6) et de
(6.4.1) appliqué pour $k=2$ et de (6.5.3) appliqué pour $k=1$.

\oldpage[IV-2]{144}

\begin{remarks}[6.5.5]
\EGAReaderTarget{ega4.c04.item.6-5-5}{IV.2.s6.5.5-fr}
\medskip\noindent
\begin{enumerate}
\item[(i)] Soient $A$, $B$ deux anneaux locaux noethériens,
$\varphi:A\to B$ un homomorphisme local tel que $B$ soit un $A$-module
plat. Soit $k$ le corps résiduel de $A$, et supposons que les deux anneaux
$A$ et $B\otimes_Ak$ vérifient la propriété $(R_i)$ (5.8.2); alors il ne
s'ensuit pas nécessairement que $B$ vérifie $(R_i)$, même dans le cas
particulier où $i=0$ ou $i=1$ et où $B$ est le complété $\widehat A$ de
$A$. Nagata a en effet donné un exemple où $A$ est normal (donc vérifie
$(R_1)$) mais où $\widehat A$ n'est même pas réduit (donc ne vérifie pas
$(R_0)$) [30]. On ne peut appliquer à ce cas la proposition (6.5.3)
parce que les fibres $f^{-1}(y)$ ne vérifient pas nécessairement la
propriété $(R_i)$ aux points distincts du point fermé de $Y=\operatorname{Spec}(A)$.
Nous montrerons toutefois plus loin (7.8.3, (v)) que de tels phénomènes
ne se présentent pas pour les anneaux locaux noethériens que l'on rencontre
le plus souvent dans les applications.

\item[(ii)] La propriété pour un préschéma d'être \emph{intègre} ne se
comporte pas du tout comme les propriétés que nous venons d'examiner dans
ce numéro et les précédents: il peut se faire que
$f:X\to Y$ soit un morphisme plat de type fini, que toutes les fibres
$f^{-1}(y)$ soient régulières (et même géométriquement régulières (6.7.6))
et que $Y$ soit intègre sans que $X$ soit même localement intègre.
Par exemple, soit $k$ un corps algébriquement clos de caractéristique $0$,
et soit $A=k[U,V]/(UV-(U+V)^3)$, l'anneau intègre (dont le spectre est
donc une «~cubique à point double~»); $A$ n'est pas intégralement clos,
et si $u$, $v$ sont les images canoniques de $U$, $V$ dans $A$, la clôture
intégrale de $A$ est l'anneau $C=A[t]$, avec
\[
t=(u-v)/(u+v),\qquad t^2=1+u+v,
\]
d'où l'on tire
\[
u=\frac{1}{2}(t^3+t^2-t),\qquad
v=\frac{1}{2}(-t^3+t^2+t),
\]
et par suite $C=k[t]$, isomorphe à l'anneau de polynômes à une
indéterminée sur $k$. Si $\mathfrak m=Au+Av$, idéal maximal de $A$
(correspondant au «~point double~» de la cubique), il y a deux idéaux
maximaux $\mathfrak n_1$, $\mathfrak n_2$ de $C$, engendrés respectivement
par $t-1$ et $t+1$. Soit alors $B$ le sous-anneau du produit $C\times C$
formé des couples de polynômes $(f,g)$ tels que
$f(1)=g(-1)$ et $f(-1)=g(1)$
($\operatorname{Spec}(B)$ est le schéma obtenu en «~recollant~» deux
exemplaires de $\operatorname{Spec}(C)$, le point $\mathfrak n_1$
(resp. $\mathfrak n_2$) de l'un des deux exemplaires étant «~recollé~»
au point $\mathfrak n_2$ (resp. $\mathfrak n_1$) de l'autre; cf. chap.~V,
où cette opération sera discutée de façon générale). Il y a donc deux
idéaux maximaux $\mathfrak r_1$, $\mathfrak r_2$ de $B$ au-dessus de
l'idéal maximal $\mathfrak m$ de $A$. En outre, comme le processus de
«~recollement~» permute à la localisation et à la complétion, on vérifie
aisément que les homomorphismes canoniques
\[
\widehat A_{\mathfrak m}\to\widehat B_{\mathfrak r_1},
\qquad
\widehat A_{\mathfrak m}\to\widehat B_{\mathfrak r_2}
\]
sont bijectifs, et par suite (Bourbaki, \emph{Alg. comm.}, chap.~III,
\S~3, n\textsuperscript{o}~5, prop.~10) $B_{\mathfrak r_1}$ et
$B_{\mathfrak r_2}$ sont des $A_{\mathfrak m}$-modules plats, ayant
d'ailleurs même corps résiduel que $A_{\mathfrak m}$. Pour tout autre idéal
maximal $\mathfrak p$ de $A$, il est immédiat qu'il y a deux idéaux
maximaux $\mathfrak q_1$, $\mathfrak q_2$ de $B$ au-dessus de $\mathfrak p$,
et que les homomorphismes
$A_{\mathfrak p}\to B_{\mathfrak q_1}$ et
$A_{\mathfrak p}\to B_{\mathfrak q_2}$ sont bijectifs. On voit ainsi que
le morphisme $\operatorname{Spec}(B)\to\operatorname{Spec}(A)$ est plat
et fini, et que toutes ses fibres sont géométriquement régulières (il est
même \emph{\'etale}, comme nous le verrons plus tard (17.6.3)); cependant
il est immédiat que $B$ n'est pas intègre.
\end{enumerate}
\end{remarks}

\oldpage[IV-2]{145}

\subsection{Propriétés de transitivité}

\begin{proposition}[6.6.1]
\label{IV.2.s6.6.1-fr}
Pour un préschéma localement noethérien $Z$, désignons par $P(Z)$ l'une
quelconque des propriétés suivantes:
\begin{enumerate}
\item[(a)] $Z$ est un préschéma de Cohen--Macaulay.
\item[(b)] $Z$ vérifie $(S_n)$.
\item[(c)] $Z$ est régulier.
\item[(d)] $Z$ vérifie $(R_n)$.
\item[(e)] $Z$ est normal.
\item[(f)] $Z$ est réduit.
\end{enumerate}
Soient alors $X$, $Y$, $Z$ trois préschémas localement noethériens,
$f:X\to Y$, $g:Y\to Z$ deux morphismes, et posons $h=g\circ f$.
\begin{enumerate}
\item[(i)] Supposons que $f$ et $g$ soient plats et que, pour tout
$y\in f(X)$ (resp. tout $z\in g(Y)$), le préschéma $f^{-1}(y)$
(resp. $g^{-1}(z)$) vérifie $P$. Alors $h$ est plat et, pour tout
$z\in h(X)$, $h^{-1}(z)$ vérifie $P$.
\item[(ii)] Supposons vérifiées les conditions suivantes: $f$ est
fidèlement plat, $h=g\circ f$ est plat, pour tout $y\in Y$, le
préschéma $f^{-1}(y)$ vérifie $P$, et, pour tout $z\in h(X)$, le
préschéma $h^{-1}(z)$ vérifie $P$. Alors $g$ est plat et, pour tout
$z\in g(Y)$, $g^{-1}(z)$ vérifie $P$.
\end{enumerate}
\begin{enumerate}
\item[(i)] On sait déjà (2.1.6) que $h$ est plat; d'autre part, pour
      tout $z\in Z$,
\[
 f_z=f\otimes 1_{k(z)}:h^{-1}(z)\longrightarrow g^{-1}(z)
\]
est plat (2.1.4), et, pour tout $y\in g^{-1}(z)$, $f_z^{-1}(y)$ est
isomorphe à $f^{-1}(y)$ par transitivité des fibres (I, 3.6.4). La
conclusion résulte alors respectivement de (6.3.5, (ii)), (6.4.2),
(6.5.2, (ii)), (6.5.3, (ii)), (6.5.4, (ii)) et (3.3.5, (ii)).
\item[(ii)] On sait déjà que $g$ est plat (2.2.13); en outre, pour
      tout $z\in Z$, $f_z$ est fidèlement plat (2.2.13). La conclusion
      résulte alors respectivement de (6.3.5, (i)), (6.4.2),
(6.5.2, (i)), (6.5.3, (i)), (6.5.4, (ii)) et (2.1.13).
\end{enumerate}
\end{proposition}

\begin{remark}[6.6.2]
\label{IV.2.s6.6.2-fr}
Supposons $h$ plat, $f$ fidèlement plat, et supposons que, pour tout
$z\in Z$, $h^{-1}(z)$ soit de coprofondeur $\leqslant n$ (5.7.1); alors
il résulte du raisonnement de (6.6.1, (ii)) et de (6.3.2) que
$g^{-1}(z)$ est de coprofondeur $\leqslant n$ pour tout $z\in g(Y)$ et
que, pour tout $y\in Y$, $f^{-1}(y)$ est de coprofondeur $\leqslant n$.
\end{remark}

\oldpage[IV-2]{146}

\subsection{Application aux changements de base dans les préschémas algébriques}

\begin{proposition}[6.7.1]
\label{IV.2.s6.7.1-fr}
Soient $k$ un corps, $X$ un $k$-préschéma localement noethérien,
$\mathscr F$ un $\mathcal O_X$-Module cohérent. Soit $k'$ une extension
de $k$; posons $X'=X\otimes_k k'$, $\mathscr F'=\mathscr F\otimes_k k'$
et soit $p:X'\to X$ la projection canonique. On suppose, soit que
$X$ est localement de type fini sur $k$, soit que $k'$ est une extension
de type fini de $k$, de sorte que dans les deux cas $X'$ est localement
noethérien. Soient $x'$ un point de $X'$, $x=p(x')$. Alors:
\begin{enumerate}
\item[(i)] On a $\operatorname{coprof}(\mathscr F_x)=
\operatorname{coprof}(\mathscr F'_{x'})$; en particulier, pour que
$\mathscr F_x$ soit un $\mathcal O_x$-module de Cohen--Macaulay, il
faut et il suffit que $\mathscr F'_{x'}$ soit un
$\mathcal O_{x'}$-module de Cohen--Macaulay.
\item[(ii)] Pour que $\mathscr F$ vérifie la propriété $(S_n)$ au point
$x$, il faut et il suffit que $\mathscr F'$ vérifie $(S_n)$ au point
$x'$.
\end{enumerate}
\end{proposition}

\begin{proof}
On sait que $p$ est fidèlement plat (2.2.13); les deux assertions sont
donc conséquences respectives de (6.3.2) et (6.4.2), pourvu que l'on
prouve que $p^{-1}(x)=\operatorname{Spec}(k(x)\otimes_k k')$ est un
préschéma de Cohen--Macaulay; comme l'un des deux corps $k(x)$, $k'$ est
une extension de type fini de $k$ par hypothèse, on est ramené à prouver
le lemme suivant.
\end{proof}

\begin{lemma}[6.7.1.1]
\label{IV.2.s6.7.1.1-fr}
Soient $K$, $L$ deux extensions d'un corps $k$, dont l'une est de type
fini (de sorte que l'anneau $K\otimes_k L$ est noethérien). Alors
$K\otimes_k L$ est un anneau de Cohen--Macaulay.
\end{lemma}

\begin{proof}
Supposons par exemple que $L$ soit une extension de type fini de $k$,
de sorte que $L$ soit extension finie d'une extension pure $k(\boldsymbol t)$
de $k$ ($\boldsymbol t$ étant un système $(t_i)_{1\leqslant i\leqslant m}$
d'indéterminées). Si l'on pose $A=K\otimes_k k(\boldsymbol t)$,
$B=K\otimes_k L$, $B$ est un $A$-module plat (0.6.2.1) et de type fini;
le morphisme $h:\operatorname{Spec}(B)\to\operatorname{Spec}(A)$ est
donc fini, et comme tout préschéma artinien est de Cohen--Macaulay, pour
prouver que $B$ est un anneau de Cohen--Macaulay, il suffit de prouver que
$A$ est un anneau de Cohen--Macaulay (6.3.5). Or, si $S$ est l'ensemble
des éléments non nuls de $k[\boldsymbol t]$, on a $A=S^{-1}A'$, où
$A'=K\otimes_k k[\boldsymbol t]=K[\boldsymbol t]$; en vertu de
(0.16.5.13), il suffit de prouver que $A'$ est un anneau de
Cohen--Macaulay; mais cela résulte de ce que $A'$ est régulier (0.17.3.7).
\end{proof}

\begin{corollary}[6.7.2]
\EGAReaderTarget{ega4.c04.cor.6-7-2}{IV.2.s6.7.2-fr}
Pour que $\mathcal O_x$ soit un anneau de Cohen--Macaulay (resp. pour
que $X$ vérifie la propriété $(S_n)$ au point $x$), il faut et il suffit
que $\mathcal O_{x'}$ soit un anneau de Cohen--Macaulay (resp. que $X'$
vérifie $(S_n)$ au point $x'$).
\end{corollary}

\begin{corollary}[6.7.3]
\label{IV.2.s6.7.3-fr}
Soient $X$, $Y$ deux $k$-préschémas localement noethériens, l'un au
moins étant localement de type fini sur $k$. Soit $\mathscr F$ (resp.
$\mathscr G$) un $\mathcal O_X$-Module cohérent (resp. un
$\mathcal O_Y$-Module cohérent). Si $\mathscr F$ et $\mathscr G$
vérifient la propriété $(S_n)$, il en est de même de
$\mathscr F\otimes_k\mathscr G$.
\end{corollary}

\begin{proof}
L'hypothèse entraîne que $X\times_kY$ est localement noethérien; notons
$p:X\times_kY\to X$ et $q:X\times_kY\to Y$ les projections canoniques,
qui sont des morphismes plats. Supposons, par exemple, que $X$ soit
localement de type fini sur $k$; on peut écrire
$\mathscr F\otimes_k\mathscr G=p^*(\mathscr F)\otimes_k\mathscr G$,
et puisque $\mathscr F$ est plat relativement au morphisme structural
$X\to\operatorname{Spec}(k)$, $p^*(\mathscr F)$ est $q$-plat (2.1.4).
Appliquons le critère (6.4.1, (ii)) au morphisme $q$; il suffit de voir
que, pour tout $y\in Y$, $p^*(\mathscr F)\otimes_{\mathcal O_Y}k(y)$
a la propriété $(S_n)$; mais
$p^*(\mathscr F)\otimes_{\mathcal O_Y}k(y)=\mathscr F\otimes_k k(y)$,
et comme $X$ est localement de type fini sur $k$, la conclusion résulte
de (6.7.1, (ii)).
\end{proof}

\begin{proposition}[6.7.4]
\label{IV.2.s6.7.4-fr}
Pour un préschéma localement noethérien $Z$, et un point $z\in Z$,
soit $P(Z,z)$ l'une des propriétés suivantes:
\begin{enumerate}
\item[(a)] $Z$ est régulier au point $z$.
\item[(b)] $Z$ vérifie $(R_n)$ au point $z$.
\item[(c)] $Z$ est normal au point $z$.
\item[(d)] $Z$ est réduit au point $z$.
\end{enumerate}
Cela étant, sous les hypothèses et avec les notations de (6.7.1), si
$P(X',x')$ est vraie, il en est de même de $P(X,x)$. La réciproque est
vraie si $k'$ est une extension séparable de $k$.
\end{proposition}

\begin{proof}
Le cas où $P$ est la propriété (d) n'a été mis que pour mémoire, ayant
déjà été traité dans (2.1.13 et 4.6.1). Le même raisonnement qu'au début
de la démonstration de (6.7.1) montre que la première assertion résulte
respectivement de (6.5.2, (i)), (6.5.3, (i)) et (6.5.4, (i)); de même,
la seconde assertion résulte de (6.5.2, (ii)), (6.5.3, (ii)) et
(6.5.4, (ii)),

\oldpage[IV-2]{147}

pourvu que l'on prouve que $p^{-1}(x)$ est un préschéma régulier; autrement
dit, on est ramené à prouver le lemme suivant.
\end{proof}

\begin{lemma}[6.7.4.1]
\label{IV.2.s6.7.4.1-fr}
Soient $K$, $L$ deux extensions d'un corps $k$, dont l'une est de type
fini. Si $L$ est une extension séparable de $k$, l'anneau $K\otimes_kL$
est régulier.
\end{lemma}

\begin{proof}
Distinguons deux cas.

\begin{enumerate}
\item[A)] $L$ est une extension de type fini de $k$. Alors, avec les
notations de (6.7.1.1), on peut supposer que $L$ est une extension finie
séparable de $k(\boldsymbol t)$. Pour tout $s\in\operatorname{Spec}(A)$,
$k(s)\otimes_{k(\boldsymbol t)}L$ est alors un composé direct d'un nombre
fini de corps, extensions finies séparables de $k(s)$, donc un anneau
régulier; il résulte alors de (6.5.2, (ii)) qu'il suffit de prouver que
l'anneau $A$ est régulier. Comme $A=S^{-1}A'$, il suffit de prouver que
$A'$ est régulier (0.17.3.6); mais on a vu qu'il en était bien ainsi dans
la démonstration de (6.7.1.1).
\item[B)] Soit $(L_\lambda)$ la famille filtrante des sous-extensions de
$L$ qui sont de type fini sur $k$; en vertu de A), chacun des anneaux
$C_\lambda=K\otimes_kL_\lambda$ est régulier; d'autre part, si
$L_\lambda\subset L_\mu$, $L_\mu$ est un $L_\lambda$-module plat, donc
$C_\mu$ est un $C_\lambda$-module plat (0.6.2.1). Comme
$C=K\otimes_kL=\varinjlim(K\otimes_kL_\lambda)$ est noethérien, on peut
appliquer (5.13.7), et $C$ est donc régulier.
\end{enumerate}
\end{proof}

\begin{remarks}[6.7.5]
\label{IV.2.s6.7.5-fr}
\begin{enumerate}
\item[(i)] Dans la démonstration du fait que $P(X,x)$ entraîne
$P(X',x')$, on ne peut se dispenser de l'hypothèse que $k'$ est une
extension séparable de $k$. On l'a déjà vu (4.6.1) lorsque $P$ est la
propriété (d); mais même lorsque $X$ est géométriquement intègre (4.6.2),
il peut se faire que $X$ soit régulier sans que $X'$ soit normal.

Prenons par exemple pour $X$ une courbe algébrique normale sur $k$
(II, 7.4.2); les anneaux locaux de $X$ étant intégralement clos et de
dimension $1$, ce sont des anneaux de valuation discrète, donc réguliers
(II, 7.1.6), et $X$ est donc un $k$-schéma régulier (et a fortiori
vérifie $(R_n)$ pour tout $n\geqslant0$). Dire que $X$ est géométriquement
intègre signifie alors que le corps $K$ des fonctions rationnelles sur
$X$ est une extension séparable et primaire de $k$ (4.6.3).

Or, prenons pour $k$ un corps non parfait de caractéristique $p>2$, et
soit $a\in k$ un élément n'appartenant pas à $k^p$. Soit
$B=k[S,T]$ l'anneau de polynômes en deux indéterminées $S$, $T$; le
polynôme
\[
P(S,T)=T^2-S^p+a
\]
est irréductible dans $B$, car on vérifie aussitôt que $S^p-a$ n'est pas
un carré dans $k[S]$; donc $X=\operatorname{Spec}(A)$, où $A=B/PB$,
est une courbe affine irréductible sur $k$. Pour voir que le schéma $X$
est régulier, il suffit de montrer qu'il est normal (II, 7.4.5); or
$A=k[S][t]$, où $t$ est une racine du polynôme $P$ considéré comme
polynôme en $T$ sur $k[S]$, de sorte que le corps $K=R(X)$ des fonctions
rationnelles sur $X$ est le corps des fractions $k(S)[T]/(P)$ de $A$,
extension quadratique de $k(S)$, donc séparable sur $k(S)$, et a fortiori
sur $k$. Comme $2$ est inversible dans $k$, on vérifie aussitôt que $A$
est la fermeture intégrale de $k[S]$ dans $K$, donc est intégralement clos,
ce qui montre que $X$ est régulier. Par ailleurs, si un élément $f+tg$
de $K$ (avec $f,g\in k(S)$) est algébrique sur $k$, il en est de même de
sa norme et de sa trace sur $k(S)$, et comme $k(S)$ est une extension pure
de $k$, on en conclut aisément que l'on doit avoir $f\in k$ et $g=0,
$autrement dit $k$ est algébriquement fermé dans $K$, et a fortiori $K$
est une extension primaire de $k$ (4.3.1). Cependant, si
$k'=k(a^{1/p})$, $X'=X\otimes_k k'$ n'est pas normal, car dans $k'[S]$
on peut écrire $S^p-a=(S-a^{1/p})^p$, et $X'$ est donc isomorphe à
$\operatorname{Spec}(A')$, où $A'=k'[S][t']$ et $t'$ est racine du
polynôme $T^2-S^p$. Or, $A'$ n'est pas intégralement clos: l'élément
$t''=t'/S$ du corps des fractions $K'=k'(S)[t']$ de $A'$ vérifie
l'équation de dépendance intégrale $(t'')^2-S^{p-2}=0$ sur $A'$, mais
n'appartient pas à $A'$. Dans la théorie classique, cela s'exprime en
disant que le «genre» sur $k'$ du corps des fonctions rationnelles $K'$ de
$X'$ est strictement inférieur à celui de $K$ sur $k$.

\oldpage[IV-2]{148}

\item[(ii)] Comme on l'a rappelé en (i), il résulte de (4.6.1) que,
lorsque $X$ est un préschéma algébrique sur $k$ qui n'est pas
géométriquement réduit, il y a des extensions radicielles finies $k'$ de
$k$ telles que $X'=X\otimes_k k'$ soit non réduit. Il est intéressant de
donner un exemple de ce fait où $X$ est un schéma régulier sur $k$, tel
que $k$ soit algébriquement fermé dans le corps $K$ des fonctions
rationnelles de $X$. Soit $k$ un corps de caractéristique $p>0$, dans
lequel il existe deux éléments $a$, $b$ formant une famille $p$-libre sur
$k^p$. Désignant encore par $B$ l'anneau $k[S,T]$, considérons cette fois
le polynôme $P(S,T)=T^p-aS^p-b$; comme $aS^p+b$ n'est pas une puissance
$p$-ième dans $k(S)$, $P$ est irréductible comme polynôme de $k(S)[T]$,
et le schéma $X_0=\operatorname{Spec}(B/PB)$ est donc une courbe affine
intègre sur $k$, dont le corps des fonctions rationnelles est
$K=k(S)[t]$, où $t$ est une racine de $P$; montrons que $k$ est
algébriquement fermé dans $K$. Supposons en effet que $K$ contienne un
élément $z$ algébrique sur $k$ et non dans $k$; on aurait alors aussi
$z\notin k(S)$, donc $[k(S)[z]:k(S)]>1$; comme $[K:k(S)]=p$, on aurait
$z^p\in k(S)$, donc $c=z^p\in k$ puisque $z$ est algébrique sur $k$;
mais on aurait $t^p=aS^p+b\in k^p(S^p)(c)$, donc $a$ et $b$
appartiendraient à $k^p(c)$, ce qui est absurde. Soit alors $X$ la
normalisée de la courbe $X_0$ dans $K$, qui est donc une courbe normale
(et par suite régulière) sur $k$. Si $k'=k(a^{1/p},b^{1/p})$, il est
clair que $aS^p+b$ est une puissance $p$-ième dans $k'(S)$, donc
$K\otimes_k k'$ n'est pas réduit, ni a fortiori le schéma
$X\otimes_k k'$.
\end{enumerate}
\end{remarks}

\begin{definition}[6.7.6]
\label{IV.2.s6.7.6-fr}
Soient $k$ un corps, $X$ un $k$-préschéma localement noethérien. On dit
que $X$ est géométriquement régulier en un point $x$ (resp. possède la
propriété $(R_n)$ géométrique en un point $x$, resp. est géométriquement
normal en un point $x$) si, pour toute extension finie $k'$ de $k$,
$X'=X_{(k')}=X\otimes_k k'$ est régulier (resp. possède la propriété
$(R_n)$, resp. est normal) en tout point $x'$ dont la projection dans
$X$ est $x$. On dit que $X$ est géométriquement régulier (resp. possède
la propriété $(R_n)$ géométrique (ou encore est géométriquement régulier
en codimension $\leqslant n$), resp. est géométriquement normal) si $X$
est géométriquement régulier (resp. possède la propriété $(R_n)$
géométrique, resp. est géométriquement normal) en tout point.

On dit qu'une algèbre $A$ sur $k$ est un anneau géométriquement régulier
(resp. géométriquement normal, resp. géométriquement réduit, resp.
possédant la propriété $(R_n)$ géométrique) si $\operatorname{Spec}(A)$
a la même propriété.

Si $X=\operatorname{Spec}(K)$, où $K$ est une extension de $k$, il revient
au même de dire que $X$ est géométriquement régulier, ou géométriquement
normal, ou géométriquement réduit, ou que $K$ est une extension séparable
de $k$: cela résulte de (4.6.1) et du fait que si $K$ est une extension
séparable de $k$ et $k'$ une extension finie de $k$, $K\otimes_k k'$ est
composé direct d'un nombre fini de corps (Bourbaki, Alg., chap. VIII,
§~7, n°~3, cor. 1 du th. 1).
\end{definition}

\oldpage[IV-2]{149}

\begin{proposition}[6.7.7]
\label{IV.2.s6.7.7-fr}
Soient $k$ un corps, $X$ un $k$-préschéma localement noethérien, $x$ un
point de $X$; désignons par $Q(k')$ la relation «$k'$ est une extension de
$k$, et $P(X_{(k')},x')$ est vraie pour tout point $x'$ de $X_{(k')}$ dont
$x$ est la projection dans $X$», où $P(Z,z)$ est une des propriétés a),
b), c) de l'énoncé (6.7.4). Alors les propriétés suivantes sont équivalentes:
\begin{enumerate}
\item[a)] $Q(k')$ est vraie pour toute extension finie $k'$ de $k$.
\item[b)] $Q(k')$ est vraie pour toute extension finie radicielle $k'$ de $k$.
\item[c)] $Q(k')$ est vraie pour toute extension de type fini $k'$ de $k$.
\end{enumerate}
Supposons en outre que $X$ soit localement de type fini sur $k$; alors les
trois propriétés précédentes sont aussi équivalentes aux suivantes:
\begin{enumerate}
\item[d)] $Q(k')$ est vraie pour toute extension $k'$ de $k$.
\item[e)] $Q(k')$ est vraie pour une extension parfaite $k'$ de $k$.
\item[f)] $Q(k')$ est vraie pour toute extension $k'=k^{p^{-s}}$, où $p$ est
l'exposant caractéristique de $k$, et $s>0$.
\end{enumerate}
\end{proposition}

\begin{proof}
Pour prouver l'équivalence de a), b) et c), il suffit évidemment d'établir
que b) entraîne c). Soit donc $K$ une extension de type fini de $k$, qui est
par suite le corps des fractions d'une $k$-algèbre de type fini $A$; posons
$Y=\operatorname{Spec}(A)$. On sait (4.6.6) qu'il existe une extension finie
et radicielle $k'$ de $k$ telle que $Y'=(Y\otimes_k k')_{\mathrm{red}}$ soit
séparable sur $k'$; si $\eta'$ est le point générique de $Y'$, $K'=k(\eta')$
est une extension séparable de $k'$ par (4.6.1). Cela étant,
$P(X_{(k')},x')$ est vraie par hypothèse pour tout point $x'$ de $X_{(k')}$
au-dessus de $x$, donc il résulte de (6.7.4) que $P(X_{(K')},x'')$ est vraie
pour tout point $x''$ de $X_{(K')}$ au-dessus de $x$, puisque
$X_{(K')}=(X_{(k')})_{(K')}$, que $x''$ est au-dessus d'un point $x'$ de
$X_{(k')}$, lui-même au-dessus de $x$, et que $K'$ est séparable sur $k'$.
Mais on a aussi
$X_{(K')}=(X_{(K)})_{(K')}$ et, pour tout $x_0\in X_{(K)}$ au-dessus de
$x$, il existe $x''\in X_{(K')}$ au-dessus de $x_0$; il résulte donc de
(6.7.4) que $P(X_{(K)},x_0)$ est vraie.

Supposons maintenant $X$ localement de type fini sur $k$, de sorte que, pour
toute extension $K$ de $k$, $X_{(K)}$ est localement noethérien. Comme une
extension radicielle de $k$ est isomorphe à une sous-extension d'une extension
parfaite quelconque de $k$, il résulte aussitôt de (6.7.4) que e) implique f);
de même, toute extension finie radicielle de $k$ est contenue dans une
extension $k^{p^{-s}}\subset k^{p^{-\infty}}$, donc f) entraîne b); d) entraîne
trivialement e), et enfin e) entraîne d). En effet, si $K$ est une extension
quelconque de $k$,
il existe une extension $K'$ de $k$ contenant (à des $k$-isomorphismes près)
$k'$ et $K$; comme $K'$ est extension séparable de $k'$, on déduit de (6.7.4)
que $Q(k')$ est vraie, puis que $Q(K)$ est vraie, en raisonnant comme dans la
première partie de la démonstration. Il suffit donc de prouver que b) entraîne
e). Nous prendrons $k'=k^{p^{-\infty}}$ dans e).

La question étant locale sur $X$, on peut en outre se borner au cas où $X$
est affine et de type fini sur $k$, en remplaçant $X$ par un voisinage de $x$;
soit donc $X=\operatorname{Spec}(B)$, où $B$ est une $k$-algèbre de type fini;
en outre, par définition de la propriété $P$ dans les cas considérés, on peut
aussi remplacer $X$ par le schéma local de $X$ au point $x$, donc supposer que
$B$ est un anneau local noethérien. Posons $B'=B\otimes_k k'$; $k'$ est
limite inductive de ses sous-extensions finies $k'_\lambda$ et, si l'on pose
$B'_\lambda=B\otimes_k k'_\lambda$, le morphisme
$f_\lambda:\operatorname{Spec}(B')\to\operatorname{Spec}(B'_\lambda)$ est
un homéomorphisme de $\operatorname{Spec}(B')$ sur
$\operatorname{Spec}(B'_\lambda)$ pour tout $\lambda$.

\oldpage[IV-2]{150}

En effet $f_\lambda$ est bijectif en vertu de (I, 3.5.2, 3.5.7 et 3.5.8);
d'autre part, il est fermé par (II, 6.1.10). On peut maintenant appliquer
(5.13.6) qui prouve (en vertu de l'hypothèse b)) que b) entraîne e) lorsque
$P(Z,z)$ est la propriété $(R_n)$ au point $z$. Le cas où $P(Z,z)$ est la
propriété d'être régulier (i.e. de posséder la propriété $(R_n)$ pour tout $n$)
en résulte trivialement. Enfin, compte tenu du critère de normalité de Serre
(5.8.6), b) entraîne encore e) lorsque $P(Z,z)$ est la propriété d'être normal
au point $z$, car cela équivaut à dire que $Z$ possède au point $z$ les
propriétés $(S_2)$ et $(R_1)$: on applique alors ce qui précède pour $n=1$,
et (6.7.2, (ii)) pour $n=2$.
\end{proof}

\begin{corollary}[6.7.8]
\label{IV.2.s6.7.8-fr}
Soit $k$ un corps; pour un $k$-préschéma localement noethérien $X$ et un
point $x\in X$, on désigne par $P(X,x)$ une des propriétés suivantes:
\begin{enumerate}
\item[(i)] $\operatorname{coprof}(\mathcal O_x)\leq n$.
\item[(ii)] $\mathcal O_x$ est un anneau de Cohen--Macaulay.
\item[(iii)] $X$ vérifie la propriété $(S_n)$ au point $x$.
\item[(iv)] $X$ est géométriquement régulier au point $x$.
\item[(v)] $X$ possède la propriété $(R_n)$ géométrique au point $x$.
\item[(vi)] $X$ est géométriquement normal au point $x$.
\item[(vii)] $X$ est géométriquement réduit (i.e. séparable) au point $x$.
\end{enumerate}
Soit $k'$ une extension de $k$; on suppose, soit que $k'$ est de type fini,
soit que $X$ est localement de type fini sur $k$, de sorte que
$X'=X\otimes_k k'$ est localement noethérien. Soit $x'$ un point de $X'$ dont
$x$ est la projection dans $X$. Alors, pour que $P(X,x)$ soit vraie, il faut
et il suffit que $P(X',x')$ le soit.
\end{corollary}

\begin{proof}
Cela a déjà été vu pour la propriété (vii) (4.6.11), et pour les propriétés
(i), (ii) et (iii) (6.7.1). Pour (iv), (v) et (vi), cela résulte de (6.7.7):
en effet, le fait que la condition est nécessaire résulte de l'équivalence des
critères a) et c), et de l'équivalence de c) et d) lorsque $X$ est localement
de type fini sur $k$. Pour voir que la condition est suffisante, soit $k''$ une
extension finie radicielle de $k$; on peut toujours considérer $k'$ et $k''$
comme sous-extensions d'une extension $K$ de $k$; posons
$X''=X\otimes_k k''$, $X_0=X\otimes_k K$, et notons que puisque $k''$ est une
extension radicielle de $k$, il n'y a qu'un seul point $x''$ de $X''$ au-dessus
de $x$ (I, 3.5.7 et 3.5.8). Soit alors $x_0$ un point quelconque de $X_0$
au-dessus de $x'$; si $P(X',x')$ est vraie, il en est de même de $P(X_0,x_0)$
en vertu de (6.7.7, c) et d)) (car si $k'$ est une extension de $k$ de type
fini, on peut supposer qu'il en est de même de $K$, et si $X$ est localement
de type fini sur $k$, $X'$ est localement de type fini sur $k'$). On déduit
alors de (6.7.4) que la propriété $Q(k'')$ est vraie (avec les notations de
(6.7.7)), donc $P(X,x)$ est vraie par (6.7.7, b)).
\end{proof}

\subsection{Morphismes réguliers, normaux, réduits, lisses}
\label{subsection:IV.6.8}

\oldpage[IV-2]{150}

\begin{definition}[6.8.1]
\label{IV.2.s6.8.1-fr}
Soient $X$, $Y$ deux préschémas, $f:X\to Y$ un morphisme tel que la fibre
$f^{-1}(y)$ soit un préschéma localement noethérien pour tout $y\in Y$, $x$
un point de $X$. On dit respectivement que $f$ est un morphisme:
\begin{enumerate}
\item[(i)] de coprofondeur $\leq n$ au point $x$;
\item[(ii)] de Cohen--Macaulay au point $x$;
\oldpage[IV-2]{151}
\item[(iii)] $(S_n)$ au point $x$;
\item[(iv)] régulier au point $x$;
\item[(v)] $(R_n)$ au point $x$;
\item[(vi)] normal au point $x$;
\item[(vii)] réduit au point $x$;
\end{enumerate}
si $f$ est plat au point $x$, et si en outre la propriété correspondante
$P(f^{-1}(f(x)),x)$ (notation de (6.7.8)) est vraie.

On dit que $f$ est lisse au point $x$ s'il est localement de présentation finie
dans un voisinage de $x$ dans $X$, et s'il est régulier au point $x$. On dit
respectivement que $f$ est un morphisme: de coprofondeur $\leq n$, de
Cohen--Macaulay, $(S_n)$, régulier, $(R_n)$, normal, réduit, lisse, s'il a la
propriété correspondante en tout point de $X$.
\end{definition}

\begin{proposition}[6.8.2]
\label{IV.2.s6.8.2-fr}
Soient $X$, $Y$ deux préschémas localement noethériens, $f:X\to Y$ un
morphisme, $x$ un point de $X$. Désignons par $M(f,x)$ une des propriétés
(i) à (vii) de la définition (6.8.1), ou la propriété pour $f$ d'être lisse au
point $x$. Soient $Y'$ un préschéma localement noethérien, $g:Y'\to Y$ un
morphisme, $X'=X\times_Y Y'$, $f'=f_{(Y')}:X'\to Y'$. On suppose que $f$
ou $g$ est localement de type fini. Alors, pour tout $x'\in X'$ au-dessus de
$x$, la propriété $M(f,x)$ entraîne $M(f',x')$.
\end{proposition}

\begin{proof}
Posons $y=f(x)$, $y'=f'(x')$; par transitivité des fibres (I, 3.6.4), on a
\[
f'^{-1}(y')=f^{-1}(y)\otimes_{k(y)}k(y').
\]
Comme ou bien $f^{-1}(y)$ est localement de type fini sur $k(y)$, ou bien
$k(y')$ est de type fini sur $k(y)$ (I, 6.4.11), il résulte de (6.7.8) que
les propriétés $P(f^{-1}(y),x)$ et $P(f'^{-1}(y'),x')$ sont équivalentes; en
outre, si $f$ est plat au point $x$, $f'$ est plat au point $x'$ (2.1.4), ce
qui prouve la proposition, compte tenu de (1.4.3, (iii)).
\end{proof}

\begin{proposition}[6.8.3]
\label{IV.2.s6.8.3-fr}
Pour un morphisme $f$ de préschémas localement noethériens, soit $M(f)$
l'une des propriétés suivantes: être de Cohen--Macaulay, $(S_n)$, régulier,
$(R_n)$, normal, réduit.
\begin{enumerate}
\item[(i)] Soient $X$, $Y$, $Z$ trois préschémas localement noethériens,
$f:X\to Y$, $g:Y\to Z$ deux morphismes. Si $M(f)$ et $M(g)$ sont vraies,
$M(g\circ f)$ est vraie.
\item[(ii)] Inversement, si $f$ est surjectif et si $M(f)$ et $M(g\circ f)$
sont vraies, alors $M(g)$ est vraie.
\item[(iii)] Soient $X$, $Y$, $Y'$ trois préschémas localement noethériens,
$f:X\to Y$, $h:Y'\to Y$ deux morphismes; on pose $X'=X\times_Y Y'$,
$f'=f_{(Y')}:X'\to Y'$. On suppose que $f$ ou $h$ est localement de type
fini. Alors, si $M(f)$ est vraie, il en est de même de $M(f')$; la réciproque
est vraie lorsque $h$ est fidèlement plat.
\end{enumerate}
Les conclusions de (i) et (iii) sont encore vraies lorsque $M$ est la propriété
d'être lisse et que (dans (iii)) $h$ est quasi-compact.
\end{proposition}

\begin{proof}
\oldpage[IV-2]{152}
\begin{enumerate}
\item[(i)] On sait déjà que si $f$ et $g$ sont plats, il en est de même de
$u=g\circ f$, et que si $f$ et $g\circ f$ sont plats et $f$ surjectif, $g$ est
plat (2.1.6 et 2.2.13). D'autre part, pour tout $z\in Z$, le morphisme
$f_z=f\otimes\mathbf 1_{k(z)}:u^{-1}(z)\to g^{-1}(z)$ est plat (resp.
fidèlement plat si $f$ l'est) et, pour tout $y\in g^{-1}(z)$,
$f_z^{-1}(y)$ est isomorphe à $f^{-1}(y)$ (I, 3.6.4). Si $M(f)$ et $M(g)$
sont vraies, $M(g\circ f)$ est donc vraie pour les cas où $M$ est la propriété
d'être de Cohen--Macaulay ou $(S_n)$, en vertu de (6.6.1). D'autre part soit
$k'$ une extension finie de $k(z)$; posons
$Y'_z=g^{-1}(z)\otimes_{k(z)}k'$, $X'_z=u^{-1}(z)\otimes_{k(z)}k'$, et
$f'_z=f_z\otimes\mathbf 1_{k'}:X'_z\to Y'_z$. Le morphisme $f'_z$ est plat
(resp. fidèlement plat) et, pour tout $y'\in Y'_z$, la fibre $f_z'^{-1}(y')$
est isomorphe à $f_z^{-1}(y)\otimes_{k(y)}k(y')$, en désignant par $y$ l'image
de $y'$ dans $g^{-1}(z)$ (I, 3.6.4). Lorsque $M$ est la propriété d'être
régulier, $(R_n)$, normal ou réduit, l'hypothèse que $M(f)$ et $M(g)$ sont
vraies entraîne que $Y'_z$ et chacune des fibres $f_z'^{-1}(y')$ possèdent
pour $y'\in Y'_z$ la propriété correspondante parmi les propriétés c), d),
e), f) de (6.6.1); on déduit donc de (6.6.1, (i)) que $X'_z$ possède la même
propriété, donc $M(g\circ f)$ est vraie.
\item[(ii)] Inversement, l'hypothèse que $M(g\circ f)$ et $M(f)$ sont vraies
et que $f$ est surjectif entraîne que $Y'_z$ a la propriété correspondante en
vertu de (6.6.1, (ii)), $f'_z$ étant surjectif pour tout $z$; donc $M(g)$ est
vraie.
\item[(iii)] La première assertion découle aussitôt de (6.8.2). D'autre part,
si $h$ est fidèlement plat et $f'$ plat, $f$ est plat (2.4.1); comme, avec les
notations de (6.8.2), les propriétés $P(f^{-1}(y),x)$ et
$P(f'^{-1}(y'),x')$ sont équivalentes (6.7.8), on voit que $M(f)$ et $M(f')$
sont alors équivalentes.
\end{enumerate}
Enfin, la dernière assertion de la proposition résulte de (1.4.3, (iii)) et
de (2.7.1, (iv)).
\end{proof}

\begin{remarks}[6.8.4]
\label{IV.2.s6.8.4-fr}
\begin{enumerate}
\item[(i)] Si $f$ est fidèlement plat, $g$ plat et si $g\circ f$ est de
coprofondeur $\leq n$, alors $g$ est de coprofondeur $\leq n$, comme il
résulte de (6.6.2).
\item[(ii)] Lorsque, dans (6.8.1), on prend pour $Y$ le spectre d'un corps
$k$, les notions (iv), (v) et (vi) se réduisent à celles définies dans (6.7.5).
Il est clair que ces dernières sont relatives au corps de base $k$. La
définition (6.8.1) conduit alors à dire qu'un préschéma $X$ est «régulier
(resp. $(R_n)$, resp. normal) sur $k$» au lieu de dire qu'il est
«géométriquement régulier (resp. $(R_n)$, resp. normal) relativement à $k$»;
on évitera de confondre cette notion avec la propriété d'être régulier (resp.
$(R_n)$, resp. normal) qui est indépendante de $k$. On peut faire à ce propos
les mêmes remarques que dans (4.5.12).
\end{enumerate}
\end{remarks}

\begin{proposition}[6.8.5]
\label{IV.2.s6.8.5-fr}
Soient $k$ un corps, $X$, $Y$ deux $k$-préschémas localement noethériens,
dont l'un est localement de type fini sur $k$. Pour un $k$-préschéma $Z$,
désignons par $P(Z)$ l'une des propriétés c) à f) de (6.6.1); la propriété
«$P(Z)$ géométrique» est alors définie dans (6.7.6) (resp. (4.6.1)) lorsque
$P(Z)$ est l'une des propriétés c), d), e) (resp. f)) de (6.6.1). Alors:
\begin{enumerate}
\item[(i)] Si $X$ possède la propriété $P$, et $Y$ la propriété $P$
géométrique, $X\times_kY$ possède la propriété $P$.
\item[(ii)] Si $X$ et $Y$ possèdent la propriété $P$ géométrique, il en est
de même de $X\times_kY$.
\end{enumerate}
\end{proposition}

\begin{proof}
Soient en effet $f:X\times_kY\to X$, $g:X\to\operatorname{Spec}(k)$ les
morphismes structuraux, qui sont fidèlement plats (2.2.13). L'hypothèse que
$Y$ possède la propriété $P$ géométrique entraîne en vertu de (6.8.2) que
$M(f)$ est vraie, $M$ étant la propriété de (6.8.3) qui correspond à $P$; sous
l'hypothèse (ii), $M(g)$ est aussi vraie, donc l'assertion (ii) résulte de
(6.8.3, (i)). Quant à l'assertion (i), elle résulte directement de (6.6.1, (i)).
\end{proof}

\begin{theorem}[6.8.6]
\label{IV.2.s6.8.6-fr}
Soient $X$, $Y$ deux préschémas localement noethériens, $f:X\to Y$ un
morphisme localement de type fini, $x$ un point de $X$, $y=f(x)$. Les
propriétés suivantes sont équivalentes:
\begin{enumerate}
\item[a)] $f$ est lisse au point $x$.
\item[b)] $f$ est régulier au point $x$.
\oldpage[IV-2]{153}
\item[c)] $\mathcal O_x$ est une $\mathcal O_y$-algèbre formellement lisse
(0, 19.3.1) pour les topologies préadiques sur $\mathcal O_x$ et $\mathcal O_y$.
\item[c')] $\mathcal O_x$ est une $\mathcal O_y$-algèbre formellement lisse
(0, 19.3.1) pour les topologies discrètes sur $\mathcal O_x$ et $\mathcal O_y$.
\end{enumerate}
\end{theorem}

\begin{proof}
L'équivalence de a) et b) résulte des définitions (6.8.1) et du fait que,
pour les préschémas localement noethériens, les morphismes localement de type
fini sont localement de présentation finie (1.4.2). En second lieu, pour que
$\mathcal O_x$ soit une $\mathcal O_y$-algèbre formellement lisse pour les
topologies préadiques, il faut et il suffit que $\mathcal O_x$ soit un
$\mathcal O_y$-module plat et que $\mathcal O_x\otimes_{\mathcal O_y}k(y)$
soit une $k(y)$-algèbre formellement lisse pour sa topologie préadique
(0, 19.7.1); mais pour que $\mathcal O_x\otimes_{\mathcal O_y}k(y)$ soit une $k(y)$-algèbre formellement lisse pour sa topologie préadique,
il faut et il suffit qu'elle soit une $k(y)$-algèbre géométriquement régulière
(0, 19.6.6); ceci démontre donc l'équivalence de b) et c). Enfin, pour prouver
l'équivalence de c) et c'), on peut se borner au cas où $Y=\operatorname{Spec}(A)$
et $X=\operatorname{Spec}(C)$ sont affines, $A$ étant noethérien et $C$ une
$A$-algèbre de type fini, que l'on peut donc écrire sous la forme
$A[T_1,\ldots,T_n]/\mathfrak J$. Comme ici $\mathfrak J/\mathfrak J^2$ est un
$C$-module de présentation finie, l'équivalence de c) et c') résulte de
(0, 22.6.4) appliqué à $A$, $B=A[T_1,\ldots,T_n]$ et $C=B/\mathfrak J$.
\end{proof}

\begin{corollary}[6.8.7]
\label{IV.2.s6.8.7-fr}
Soient $X$, $Y$ deux préschémas localement noethériens, $f:X\to Y$ un
morphisme localement de type fini. Alors l'ensemble des points $x\in X$ où
$f$ est lisse (ou régulier) est ouvert dans $X$.
\end{corollary}

\begin{proof}
Il résulte en effet de (0, 22.6.5) que l'ensemble des $x\in X$ vérifiant la
condition c') de (6.8.6) est ouvert dans $X$, et on conclut par (6.8.6).
\end{proof}

\begin{remark}[6.8.8]
\label{IV.2.s6.8.8-fr}
Dans (17.5.1), nous montrerons que l'équivalence de b) et c') dans (6.8.6),
ainsi que le corollaire (6.8.7), restent valables sans hypothèse noethérienne
sur $X$ et $Y$, pourvu que l'on se borne aux morphismes localement de
présentation finie.
\end{remark}

\subsection{Le théorème de platitude générique}
\label{subsection:IV.6.9}

\begin{theorem}[6.9.1]
\label{IV.2.s6.9.1-fr}
Soient $Y$ un préschéma localement noethérien et intègre, $u:X\to Y$ un
morphisme de type fini, $\mathcal F$ un $\mathcal O_X$-Module cohérent. Il
existe alors un ouvert non vide $U$ de $Y$ tel que
$\mathcal F|_{u^{-1}(U)}$ soit plat sur $U$.
\end{theorem}

\begin{proof}
On peut évidemment se borner au cas où $Y=\operatorname{Spec}(A)$ est affine;
alors $X$ est réunion finie d'ouverts affines $X_i$ de type fini sur $Y$; si,
pour chaque $i$, il y a un ouvert non vide $U_i$ dans $Y$ tel que
$\mathcal F|_{X_i\cap u^{-1}(U_i)}$ soit plat sur $U_i$, il est clair qu'en
prenant pour $U$ l'intersection des $U_i$, on répondra à la question; on peut
donc supposer que $X=\operatorname{Spec}(B)$, où $B$ est une $A$-algèbre de
type fini. On a alors $\mathcal F=\widetilde M$, où $M$ est un $B$-module de
type fini; le théorème résultera par suite du lemme suivant.
\end{proof}

\begin{lemma}[6.9.2]
\label{IV.2.s6.9.2-fr}
Soient $A$ un anneau intègre noethérien, $B$ une $A$-algèbre de type fini,
$M$ un $B$-module de type fini. Alors il existe $f\ne0$ dans $A$ tel que
$M_f$ soit un $A_f$-module libre.
\end{lemma}

\begin{proof}
\oldpage[IV-2]{154}
Désignons par $K$ le corps des fractions de $A$; alors $B\otimes_AK$ est une
algèbre de type fini sur $K$ et $M\otimes_AK$ un $(B\otimes_AK)$-module de
type fini. Nous raisonnerons par récurrence sur la dimension $n$ du support de
$M\otimes_AK$, qui est $-\infty$ ou finie et $\geq0$. Supposons d'abord
$n=-\infty$, c'est-à-dire $M\otimes_AK=0$; si $(m_i)_{1\leq i\leq r}$ est
un système de générateurs du $B$-module $M$, il existe donc un $f\ne0$ dans
$A$ tel que $fm_i=0$ pour $1\leq i\leq r$; donc $M_f=0$ et le lemme est vrai
dans ce cas. Supposons maintenant $n\geq0$. On sait qu'il existe une suite de
composition $M=M_1\supset M_2\supset\cdots\supset M_q=0$ du $B$-module $M$
telle que chacun des quotients $N_i=M_i/M_{i+1}$ soit isomorphe à un
$B$-module de la forme $B/\mathfrak p_i$, où $\mathfrak p_i$ est un idéal
premier de $B$ (Bourbaki, \emph{Alg. comm.}, chap. IV, \S~1, n°~4, th.~1).
Si le théorème est vrai pour chacun des $N_i$, il y a pour chaque $i$ un
$f_i\ne0$ dans $A$ tel que $(N_i)_{f_i}$ soit libre sur $A_{f_i}$; posant
$f=f_1f_2\cdots f_{q-1}$, il en résulte que $(N_i)_f$ est un $A_f$-module
libre pour $1\leq i\leq q-1$. Mais $(N_i)_f=(M_i)_f/(M_{i+1})_f$
$(0_{\mathrm I},1.3.2)$ et comme une extension de modules libres est libre, on
en déduit alors que $M_f$ est un $A_f$-module libre. Remplaçant $B$ par
$B/\mathfrak p$ ($\mathfrak p$ idéal premier de $B$), qui est encore de type
fini sur $A$, on voit qu'on peut se borner au cas où $M=B$ et $B$ est intègre.
On sait alors (Bourbaki, \emph{Alg. comm.}, chap. V, \S~3, n°~1, cor. 1 du
th. 1) qu'il existe un élément $g\ne0$ dans $A$ et des éléments $t_i$
$(1\leq i\leq m)$ de $B$, algébriquement indépendants sur $A$ et tels que
$B_g$ soit entier sur $A_g[t_1,\ldots,t_m]$. On peut remplacer $A$ par $A_g$,
$B$ par $B_g$ et supposer par suite que $B$ est entier sur
$C=A[t_1,\ldots,t_m]$, donc un $C$-module de type fini et sans torsion. On sait
en outre (4.1.2) que la dimension de $\operatorname{Spec}(B\otimes_AK)$ est
égale à $m$, donc on a $m=n$.

Cela étant, si $h$ est le rang du $C$-module sans torsion $B$, il existe une
suite exacte de $C$-modules
\[
0\longrightarrow C^h\longrightarrow B\longrightarrow M'\longrightarrow0
\]
où $M'$ est un $C$-module de torsion de type fini; le support de $M'$ ne
contient donc pas le point générique de $\operatorname{Spec}(C)$ (I, 7.4.6)
et par suite le support de $M'\otimes_AK$ ne contient pas le point générique
de $\operatorname{Spec}(C\otimes_AK)$ (I, 9.1.13.1); on en conclut (4.1.2.1)
que sa dimension est $<n$. En vertu de l'hypothèse de récurrence, il existe un
$f\ne0$ dans $A$ tel que $M'_f$ soit un $A_f$-module libre; par ailleurs
$C_f^h$ est un $A_f$-module libre; donc il en est de même de $B_f$, qui en
vertu de $(0_{\mathrm I},1.3.2)$ est une extension de modules libres. C.Q.F.D.
\end{proof}

\begin{corollary}[6.9.3]
\label{IV.2.s6.9.3-fr}
Soient $S$ un préschéma noethérien, $u:X\to S$ un morphisme de type fini,
$\mathcal F$ un $\mathcal O_X$-Module cohérent. Il y a alors une partition
$(S_i)_{1\leq i\leq n}$ de $S$ en un nombre fini d'ensembles localement fermés
dans $S$, telle que, si on désigne encore par $S_i$ le sous-préschéma réduit
de $S$ ayant pour espace sous-jacent $S_i$, et si on pose
$X_i=X\times_S S_i$, le $\mathcal O_{X_i}$-Module
$\mathcal F_i=\mathcal F\otimes_{\mathcal O_S}\mathcal O_{S_i}$ soit plat
sur $S_i$.
\end{corollary}

\begin{proof}
Procédons par récurrence noethérienne $(0_{\mathrm I},2.2.2)$ sur l'ensemble
des parties fermées $T$ de $S$ telles que le sous-préschéma réduit de $S$
ayant $T$ pour espace sous-jacent vérifie la conclusion de (6.9.3). On peut
donc se borner à démontrer le corollaire pour $S$ en supposant qu'il soit vrai
pour tout sous-préschéma fermé réduit de $S$ ayant pour espace sous-jacent une
partie fermée $\ne S$. Comme le morphisme $S_{\mathrm{red}}\to S$ est de type
fini et surjectif, on peut remplacer $S$ par $S_{\mathrm{red}}$, donc supposer
$S$ réduit et non vide. Comme $S$ est noethérien, l'intérieur $T$ d'une
composante irréductible de $S$ est non vide, et le préschéma induit sur $T$
est intègre; il y a donc en vertu de (6.9.1) un ouvert non vide $U\subset T$

\oldpage[IV-2]{155}

tel que $\mathcal F|_{u^{-1}(U)}$ soit plat sur $U$. Si $Y$ est alors le
sous-préschéma réduit de $S$ ayant pour espace sous-jacent $S-U$, il y a par
hypothèse une partition $(Y_i)$ de $Y$ en ensembles localement fermés dans
$Y$ (donc dans $S$) tels que
$\mathcal F_i=\mathcal F\otimes_{\mathcal O_S}\mathcal O_{Y_i}$ soit plat
sur $Y_i$ pour tout $i$; il est clair que les $Y_i$ et $U$ forment une
partition répondant à la question.
\end{proof}

\subsection{Dimension et profondeur d'un Module normalement plat le long d'un
sous-préschéma fermé}
\label{subsection:IV.6.10}

\begin{env}[6.10.1]
\EGAReaderTarget{ega4.c04.statement.6-10-1}{IV.2.s6.10.1-fr}
Soient $X$ un préschéma localement noethérien, $\mathcal J$ un Idéal
quasi-cohérent de $\mathcal O_X$, $Y$ le sous-préschéma fermé de $X$ défini
par $\mathcal J$, $j:Y\to X$ l'injection canonique. Pour tout entier $k\geq0$,
le $\mathcal O_X$-Module
$\mathcal J^k\mathcal F/\mathcal J^{k+1}\mathcal F$ est annulé par
$\mathcal J$, donc de la forme $j_*(\mathcal G_k)$, où
\[
\mathcal G_k=j^*(\mathcal J^k\mathcal F/\mathcal J^{k+1}\mathcal F)
=j^*(\mathcal J^k\mathcal F)
\]
est un $\mathcal O_Y$-Module cohérent. Par abus de langage, nous noterons
$\operatorname{gr}_{\mathcal J}^{\bullet}(\mathcal F)$ le
$\mathcal O_Y$-Module gradué égal à la somme directe
\[
\bigoplus_{k=0}^{\infty}\mathcal G_k
=j^*\left(\bigoplus_{k=0}^{\infty}
\mathcal J^k\mathcal F/\mathcal J^{k+1}\mathcal F\right);
\]
en particulier, on a
$\operatorname{gr}_{\mathcal J}^{0}(\mathcal F)=\mathcal G_0
=\mathcal F\otimes_{\mathcal O_X}\mathcal O_Y=j^*(\mathcal F)$. Nous dirons
(avec Hironaka) que $\mathcal F$ est normalement plat le long de $Y$ si
$\operatorname{gr}_{\mathcal J}^{\bullet}(\mathcal F)$ est un
$\mathcal O_Y$-Module plat. Il revient au même ((2.1.12) et
$(0_{\mathrm I},6.1.2)$) de dire que chacun des $\mathcal O_Y$-Modules
$\mathcal G_k=j^*(\mathcal J^k\mathcal F)$ est localement libre.
\end{env}

\begin{proposition}[6.10.2]
\label{IV.2.s6.10.2-fr}
Soient $X$ un préschéma localement noethérien, $Y$ un sous-préschéma fermé
intègre de $X$. Pour tout $\mathcal O_X$-Module cohérent $\mathcal F$, il
existe un ouvert $U$ de $X$ tel que $Y\cap U\ne\varnothing$ et que
$\mathcal F|U$ soit normalement plat le long de $Y\cap U$.
\end{proposition}

\begin{proof}
Soit en effet $\mathcal J$ l'Idéal cohérent de $\mathcal O_X$ définissant
$Y$; la $\mathcal O_Y$-Algèbre
$\mathcal B=\operatorname{gr}_{\mathcal J}^{\bullet}(\mathcal O_X)$ est
quasi-cohérente et de type fini car elle est engendrée par
$\operatorname{gr}_{\mathcal J}^{1}(\mathcal O_X)$, image réciproque de
$\mathcal J/\mathcal J^2$; comme
$\operatorname{gr}_{\mathcal J}^{\bullet}(\mathcal F)$ est un
$\mathcal B$-module quasi-cohérent engendré par
$\operatorname{gr}_{\mathcal J}^{0}(\mathcal F)$, c'est un
$\mathcal B$-Module de type fini. Si l'on pose
$Z=\operatorname{Spec}(\mathcal B)$, le morphisme structural $u:Z\to Y$ est
alors de type fini, et si $\mathcal H$ est le $\mathcal O_Z$-Module cohérent
tel que $u_*(\mathcal H)=\operatorname{gr}_{\mathcal J}^{\bullet}(\mathcal F)$,
il suffit d'appliquer à $u$ et à $\mathcal H$ le théorème de platitude
générique (6.9.1) pour prouver la proposition.
\end{proof}

\begin{proposition}[6.10.3]
\label{IV.2.s6.10.3-fr}
Les notations étant celles de (6.10.1), supposons que $\mathcal F$ soit
normalement plat le long de $Y$. Alors:
\begin{enumerate}
\item[(i)] $Y\cap\operatorname{Supp}(\mathcal F)$ est une partie à la fois
ouverte et fermée de $Y$ (autrement dit, une réunion de composantes connexes
de $Y$).
\item[(ii)] Si $\mathcal J$ est localement nilpotent,
$\operatorname{Supp}(\mathcal F)$ est une partie ouverte et fermée de $X$, et
pour tout $x\in\operatorname{Supp}(\mathcal F)$, on a
\end{enumerate}
\begin{align}
\dim(\mathcal F_x)&=\dim(\mathcal O_{Y,x})
  \tag{6.10.3.1}\label{IV.2.s6.10.3.1-fr}\\
\operatorname{prof}(\mathcal F_x)&=\operatorname{prof}(\mathcal O_{Y,x})
  \tag{6.10.3.2}\label{IV.2.s6.10.3.2-fr}\\
\operatorname{coprof}(\mathcal F_x)&=\operatorname{coprof}(\mathcal O_{Y,x}).
  \tag{6.10.3.3}\label{IV.2.s6.10.3.3-fr}
\end{align}
\end{proposition}

\begin{proof}
\begin{enumerate}
\item[(i)] Avec les notations de (6.10.1), on a
$Y\cap\operatorname{Supp}(\mathcal F)=\operatorname{Supp}(\mathcal G_0)$
(I, 9.1.13), et comme par hypothèse $\mathcal G_0$ est un
$\mathcal O_Y$-Module localement libre de type fini, son support est à la fois
ouvert et fermé dans $Y$.

\oldpage[IV-2]{156}

\item[(ii)] L'hypothèse que $\mathcal J$ est localement nilpotent entraîne que
les espaces sous-jacents à $X$ et à $Y$ sont les mêmes, d'où les deux premières
assertions de (ii), compte tenu de (5.1.12.1); il reste à prouver (6.10.3.2),
car on en déduira aussitôt (6.10.3.3) par différence. Soit
$(f_i)_{1\leq i\leq p}$ une suite finie d'éléments de l'idéal maximal de
$\mathcal O_{X,x}$ dont les images dans
$\mathcal O_{Y,x}=\mathcal O_{X,x}/\mathcal J_x$ forment une suite
$\mathcal O_{Y,x}$-régulière maximale; l'hypothèse sur $\mathcal F$ et
$\mathcal J$ entraîne que le $\mathcal O_{Y,x}$-module
$\operatorname{gr}_{\mathcal J_x}^{\bullet}(\mathcal F_x)$ est libre de type
fini, donc la suite $(f_i)$ est
$\operatorname{gr}_{\mathcal J_x}^{\bullet}(\mathcal F_x)$-régulière; on en
déduit que cette suite est aussi $\mathcal F_x$-régulière (0, 15.1.19). Soit
d'autre part $n$ le plus grand entier tel que
$\mathcal F_x^{(n)}=\mathcal J_x^n\mathcal F_x\ne0$. L'hypothèse entraîne
aussi que
$\operatorname{gr}_{\mathcal J_x}^{\bullet}
(\mathcal F_x/\mathcal F_x^{(n)})$ est libre de type fini, donc la suite
$(f_i)$ est aussi $(\mathcal F_x/\mathcal F_x^{(n)})$-régulière
(\emph{loc. cit.}). Appliquant le lemme (3.4.1.4), on en conclut par récurrence
sur $i$, une suite exacte
\[
0\longrightarrow
\mathcal F_x^{(n)}/\left(\sum_{i=1}^{p}f_i\mathcal F_x^{(n)}\right)
\longrightarrow
\mathcal F_x/\left(\sum_{i=1}^{p}f_i\mathcal F_x\right).
\]
Mais l'hypothèse entraîne que $\mathcal F_x^{(n)}$ est un
$\mathcal O_{Y,x}$-module libre de type fini et $\ne0$, donc
$\mathcal F_x^{(n)}/(\sum_{i=1}^{p}f_i\mathcal F_x^{(n)})$ est isomorphe à
un module de la forme
$(\mathcal O_{Y,x}/(\sum_{i=1}^{p}f_i\mathcal O_{Y,x}))^m$ avec $m>0$; comme
\[
\operatorname{prof}\left(\mathcal O_{Y,x}/
\left(\sum_{i=1}^{p}f_i\mathcal O_{Y,x}\right)\right)=0
\]
(0, 16.4.6), on a aussi, par la caractérisation (0, 16.4.6, (i)) des modules
de profondeur nulle,
\[
\operatorname{prof}\left(\mathcal F_x^{(n)}/
\left(\sum_{i=1}^{p}f_i\mathcal F_x^{(n)}\right)\right)=0,
\]
puis
\[
\operatorname{prof}\left(\mathcal F_x/
\left(\sum_{i=1}^{p}f_i\mathcal F_x\right)\right)=0;
\]
les $(f_i)$ appartenant à l'idéal maximal de $\mathcal O_{Y,x}$ et formant
une suite $\mathcal F_x$-régulière, cela montre que l'on a
$\operatorname{prof}(\mathcal F_x)=p$ (0, 16.4.6, (ii)).
\end{enumerate}
\end{proof}

\begin{corollary}[6.10.4]
\label{IV.2.s6.10.4-fr}
Soit $U_{S_n}(\mathcal F)$ (resp. $U_{C_n}(\mathcal F)$) l'ensemble des
$x\in X$ tels que $\mathcal F$ vérifie $(S_n)$ au point $x$ (resp. tels que
$\operatorname{coprof}(\mathcal F_x)\leq n$). Si $\mathcal F$ est normalement
plat le long de $Y$, si $\operatorname{Supp}(\mathcal F)=X$ et si $\mathcal J$
est localement nilpotent, on a
$U_{S_n}(\mathcal F)=U_{S_n}(\mathcal O_Y)$ et
$U_{C_n}(\mathcal F)=U_{C_n}(\mathcal O_Y)$.
\end{corollary}

\begin{proposition}[6.10.5]
\label{IV.2.s6.10.5-fr}
Les notations étant celles de (6.10.1), supposons que $Y$ soit connexe et non
vide, et que $\mathcal F$ soit normalement plat le long de $Y$. Pour tout
entier $n\geq0$, on pose
\begin{equation}
r(n)=\operatorname{rg}_{\mathcal O_Y}
(\operatorname{gr}_{\mathcal J}^{n}(\mathcal F))
\tag{6.10.5.1}\label{IV.2.s6.10.5.1-fr}
\end{equation}
(le $\mathcal O_Y$-Module localement libre
$\operatorname{gr}_{\mathcal J}^{n}(\mathcal F)$ étant nécessairement de rang
constant). Alors:
\begin{enumerate}
\item[(i)] Il existe un polynôme $P\in\mathbf Q[T]$ tel que $P(n)=r(n)$ pour
tout $n$ assez grand.
\item[(ii)] On a $Y\cap\operatorname{Supp}(\mathcal F)=\varnothing$
(autrement dit $\mathcal F/\mathcal J\mathcal F=0$), ou
$Y\cap\operatorname{Supp}(\mathcal F)=Y$ (autrement dit
$Y\subset\operatorname{Supp}(\mathcal F)$).
\end{enumerate}
Dans le second cas, notons $d-1$ le degré de $P$; pour tout point maximal $y$
de $Y$, on a
\begin{equation}
\dim(\mathcal F_y)=d
\tag{6.10.5.2}\label{IV.2.s6.10.5.2-fr}
\end{equation}
et en particulier
\begin{equation}
\operatorname{codim}(Y,\operatorname{Supp}(\mathcal F))=d.
\tag{6.10.5.3}\label{IV.2.s6.10.5.3-fr}
\end{equation}
\begin{enumerate}
\item[(iii)] Supposons $Y\subset\operatorname{Supp}(\mathcal F)$. Pour tout
$x\in Y$, on a
\end{enumerate}
\begin{equation}
\dim(\mathcal F_x)=\dim(\mathcal O_{Y,x})+d.
\tag{6.10.5.4}\label{IV.2.s6.10.5.4-fr}
\end{equation}
De façon plus précise, il existe dans $\mathcal J_x$ des éléments $f_i$
$(1\leq i\leq d)$ faisant partie d'un système de paramètres pour
$\mathcal F_x$ (0, 16.3.6) et tels que, dans
$\operatorname{Spec}(\mathcal O_{X,x})$, on ait
\[
V(\mathcal J_x)=V\left(\sum_{i=1}^{d}f_i\mathcal O_{X,x}\right)
\cap\operatorname{Supp}(\mathcal F).
\]
\end{proposition}

\begin{proof}
Comme $Y$ est supposé connexe, la première assertion de (ii) découle de
(6.10.3, (i)). Si l'on a
$Y\cap\operatorname{Supp}(\mathcal F)=\varnothing$, l'assertion (i) est
triviale avec $P=0$. Supposons que
$\operatorname{Supp}(\mathcal F)\supset Y$, et soit $y$ un point maximal de
$Y$; comme on a alors $Y=\operatorname{Supp}(\mathcal F/\mathcal J\mathcal F)$,
$\mathcal F_y/\mathcal J_y\mathcal F_y$ est un
$\mathcal O_{X,y}$-module de longueur finie (3.1.4); posant
$s(n)=\operatorname{long}(\operatorname{gr}_{\mathcal J_y}^{n}(\mathcal F_y))$,
on sait qu'il y a un polynôme $R\in\mathbf Q[T]$ tel que $s(n)=R(n)$ pour
$n$ assez grand, savoir le polynôme $H(n+1)-H(n)$, où $H$ est le polynôme
de Hilbert--Samuel de $\mathcal F_y$ pour la filtration
$\mathcal J_y$-préadique (0, 16.2.1). Comme les
$\mathcal O_{Y,y}$-modules
$\operatorname{gr}_{\mathcal J_y}^{n}(\mathcal F_y)$ sont libres par
hypothèse, on a $r(n)=s(n)/m$, en désignant par $m$ la longueur du
$\mathcal O_{X,y}$-module $\mathcal O_{X,y}/\mathcal J_y$; on satisfait donc
à (i) en
prenant $P=R/m$. On sait en outre (0, 16.2.3) que
$\deg(H)=\dim(\mathcal F_y)$, d'où la relation (6.10.5.2); la relation
(6.10.5.3) s'en déduit à l'aide de (5.1.12.2).

\oldpage[IV-2]{157}

Reste à prouver (iii) pour un point quelconque $x\in Y$. Posons
$A=\mathcal O_{X,x}$, $\mathfrak J=\mathcal J_x$, de sorte que
$\mathcal O_{Y,x}=A/\mathfrak J$, et $M=\mathcal F_x$; soit
$S=\operatorname{gr}_{\mathfrak J}^{\bullet}(A)$, qui est une
$(A/\mathfrak J)$-algèbre graduée de type fini, à degrés positifs, telle
que $S_0=A/\mathfrak J$, et engendrée par ses éléments homogènes de degré
$1$; soit enfin
$N=\operatorname{gr}_{\mathfrak J}^{\bullet}(M)$, qui est un $S$-module
gradué de type fini, chaque composante homogène $N_n$ étant par hypothèse
un $(A/\mathfrak J)$-module libre de longueur $r(n)$. Soient $\mathfrak m$
l'idéal maximal de $A$, $k=A/\mathfrak m$ son corps résiduel;
$B=S\otimes_A k$ est une $k$-algèbre graduée de type fini, à degrés
positifs, engendrée par ses éléments homogènes de degré $1$ et telle que
$B_0=k$, de sorte que $\mathfrak q=B_+=\bigoplus_{n\geq1}B_n$ est un idéal
maximal dans $B$; $E=N\otimes_A k$ est un $B$-module gradué de type fini
tel que $\operatorname{rg}_k(E_n)=r(n)$. Appliquons (0, 16.2.7) à l'anneau
gradué $B=S\otimes_A k$ et au $B$-module gradué $E=N\otimes_A k$, et soit
$f_i\in\mathfrak J^{n_i}$ $(1\leq i\leq d)$ un élément dont $t_i$ est
l'image dans $B_{n_i}$. Pour $n\geq\sup(n_i)$, considérons le sous-$A$-module
$\sum_{i=1}^{d}f_i\mathfrak J^{n-n_i}M$ de $M$; comme le composant homogène
de degré $n$ du sous-module $\sum_{i=1}^{d}t_iE$ de $E$ est égal à
$E_n$ dès que $n$ est assez grand, on voit que, pour $n$ assez grand, on a
\[
\mathfrak J^nM=\sum_{i=1}^{d}f_i\mathfrak J^{n-n_i}M+
\mathfrak m\mathfrak J^nM
\]
et comme $\mathfrak J^nM$ est un $A$-module de type fini, cela entraîne,
par le lemme de Nakayama,
\[
\mathfrak J^nM=\sum_{i=1}^{d}f_i\mathfrak J^{n-n_i}M
\subset\sum_{i=1}^{d}f_iM.
\]
Si $\mathfrak a$ est l'annulateur de $M$, on a donc (0$_{\rm I}$, 1.7.5)
dans $\operatorname{Spec}(A)$
\[
V\left(\sum_{i=1}^{d}f_iA\right)\cap V(\mathfrak a)
\subset V(\mathfrak J^n)\cap V(\mathfrak a)
=V(\mathfrak J)\cap V(\mathfrak a)=V(\mathfrak J)
\]
puisque par hypothèse
$Y\cap\operatorname{Spec}(A)=V(\mathfrak J)\cap\operatorname{Supp}(M)
=V(\mathfrak J)$. Comme d'autre part les $f_i$ appartiennent à
$\mathfrak J$, on a
$V(\mathfrak J)\subset V(\sum_{i=1}^{d}f_iA)$, ce qui prouve la dernière
relation de (iii). Il reste à montrer que les $f_i$ font partie d'un système
de paramètres pour $M$. Remplaçant $A$ par $A/\mathfrak a$ et les $f_i$ par
leurs images dans $A/\mathfrak a$, on peut se borner au cas où
$\mathfrak a=0$; on vient de voir que l'on a
\[
\dim(A/\mathfrak J)=
\dim\left(A/\left(\sum_{i=1}^{d}f_iA\right)\right)
\]
et il reste donc à prouver (0, 16.3.7) que l'on a
\[
\dim(A)\geq\dim(A/\mathfrak J)+d.
\]
Or, soit $\mathfrak p$ un idéal premier de $A$ contenant $\mathfrak J$ et
tel que $\dim(A/\mathfrak J)=\dim(A/\mathfrak p)$, de sorte que
$\mathfrak p$ est minimal parmi les idéaux premiers contenant
$\mathfrak J$. On a donc $\mathfrak p=\mathfrak j_y$ où
$y\in\operatorname{Spec}(A)$ est un point maximal de $Y$. Mais en vertu de
(6.10.5.2) et de l'hypothèse
$\operatorname{Spec}(A)=\operatorname{Supp}(M)$, on a
$\dim(A_{\mathfrak p})=d$; l'inégalité
$\dim(A/\mathfrak p)+\dim(A_{\mathfrak p})\leq\dim(A)$ (0, 16.1.4)
achève la démonstration.
\end{proof}

\begin{proposition}[6.10.6]
\label{IV.2.s6.10.6-fr}
Soient $X$ un préschéma localement noethérien, $\mathcal F$ un
$\mathcal O_X$-Module cohérent, $Y$ un sous-préschéma fermé irréductible de
$X$, de point générique $y\in\operatorname{Supp}(\mathcal F)$. Il existe
alors un voisinage ouvert non vide $U$ de $y$ dans $X$ tel que, pour tout
$x\in U\cap Y$, on ait
\begin{align}
\dim(\mathcal F_x)&=\dim(\mathcal F_y)+\dim(\mathcal O_{Y,x})
\tag{6.10.6.1}\label{IV.2.s6.10.6.1-fr}\\
\operatorname{prof}(\mathcal F_x)&=\operatorname{prof}(\mathcal F_y)+
\operatorname{prof}(\mathcal O_{Y,x})
\tag{6.10.6.2}\label{IV.2.s6.10.6.2-fr}\\
\operatorname{coprof}(\mathcal F_x)&=\operatorname{coprof}(\mathcal F_y)+
\operatorname{coprof}(\mathcal O_{Y,x}).
\tag{6.10.6.3}\label{IV.2.s6.10.6.3-fr}
\end{align}
\end{proposition}

\begin{proof}
\oldpage[IV-2]{158}
Soit $Y'=Y_{\mathrm{red}}$, qui est un sous-préschéma fermé intègre de $Y$
ayant même espace sous-jacent, et défini par un Idéal localement nilpotent
$\mathcal J$ de $\mathcal O_Y$ (I, 6.1.6). Il en résulte que
$\operatorname{gr}_{\mathcal J}^{\bullet}(\mathcal O_Y)$ est un
$\mathcal O_{Y'}$-Module cohérent, et comme $Y'$ est intègre, il y a un
voisinage ouvert $V$ de $y$ dans $Y'$ tel que
$\operatorname{gr}_{\mathcal J}^{\bullet}(\mathcal O_Y)|V$ soit localement
libre (0$_{\rm I}$, 5.2.7); autrement dit, quitte à remplacer $X$ par un
voisinage de $y$, on peut supposer que $\mathcal O_Y$ est normalement plat
le long de $Y'$; on déduit donc de (6.10.3) que l'on a
\[
\dim(\mathcal O_{Y,x})=\dim(\mathcal O_{Y',x}),\qquad
\operatorname{prof}(\mathcal O_{Y,x})=
\operatorname{prof}(\mathcal O_{Y',x})
\]
pour tout $x\in Y$; cela permet de se borner au cas où le sous-préschéma
fermé $Y$ est intègre.

Cela étant, en vertu de (6.10.2), on peut, en remplaçant $X$ par un
voisinage ouvert de $y$, supposer que $\mathcal F$ est normalement plat le
long de $Y$. Comme $y\in\operatorname{Supp}(\mathcal F)$, la relation
(6.10.6.1) résulte de (6.10.5.2) et (6.10.5.4). Posons maintenant
$p=\operatorname{prof}(\mathcal F_y)$; remplaçant au besoin $X$ par un
voisinage ouvert de $y$, on peut supposer qu'il existe $p$ sections $f_i$
$(1\leq i\leq p)$ de $\mathcal O_X$ au-dessus de $X$, formant une suite
$\mathcal F$-régulière, tels que les $(f_i)_y$ appartiennent à l'idéal
maximal $\mathfrak m_y$ de $\mathcal O_{X,y}$ et que
\[
\operatorname{prof}\left(\mathcal F_y/
\left(\sum_{i=1}^{p}(f_i)_y\mathcal F_y\right)\right)=0
\]
(0, 16.4.6). Si $\mathcal J$ est l'Idéal de $\mathcal O_X$ qui définit
$Y$, on a $\mathcal J_y=\mathfrak m_y$ et, en remplaçant encore au besoin $X$
par un voisinage de $y$, on peut supposer que
$f_i\in\Gamma(X,\mathcal J)$ pour $1\leq i\leq p$. En outre, si l'on pose
$\mathcal G=\mathcal F/(\sum_{i=1}^{p}f_i\mathcal F)$, l'hypothèse
$\operatorname{prof}(\mathcal G_y)=0$ entraîne que $\mathcal G_y$ contient
un sous-module isomorphe à
$\mathcal O_{X,y}/\mathfrak m_y$ (0, 16.4.6); le même raisonnement (compte
tenu de (0$_{\rm I}$, 5.3.8 et 5.2.7)) montre qu'on peut supposer qu'il
existe un sous-$\mathcal O_X$-Module $\mathcal G'$ de $\mathcal G$
isomorphe à $\mathcal O_X/\mathcal J$ de sorte que
$\mathcal G'_x=\mathcal O_{Y,x}$ pour tout $x\in Y$. Posons
$\mathcal G''=\mathcal G/\mathcal G'$; remplaçant $X$ par un voisinage
ouvert de $y$, on peut, en vertu de (6.10.2), supposer que $\mathcal G'$ et
$\mathcal G''$ sont normalement plats le long de $Y$. Soit alors $x$ un
point quelconque de $Y$, et posons
$q=\operatorname{prof}(\mathcal O_{Y,x})=\operatorname{prof}(\mathcal G'_x)$;
soit $(g_j)_{1\leq j\leq q}$ une suite maximale
$\mathcal G'_x$-régulière d'éléments de l'idéal maximal
$\mathfrak m_x$ de $\mathcal O_{X,x}$. Chacun des composants homogènes de
$\operatorname{gr}_{\mathcal J_x}^{\bullet}(\mathcal G''_x)$ est un
$\mathcal O_{Y,x}$-module plat de type fini par hypothèse, donc est un
$\mathcal O_{Y,x}$-module libre, et il en est par suite de même de
$\operatorname{gr}_{\mathcal J_x}^{\bullet}(\mathcal G''_x)$; comme la
suite $(g_j)$ est $\mathcal O_{Y,x}$-régulière, elle est par suite
$\operatorname{gr}_{\mathcal J_x}^{\bullet}(\mathcal G''_x)$-régulière,
d'où l'on conclut qu'elle est $\mathcal G''_x$-régulière (0, 15.1.19).
Appliquant le lemme (0, 15.1.18) à la suite exacte
\[
0\longrightarrow\mathcal G'_x\longrightarrow\mathcal G_x
\longrightarrow\mathcal G''_x\longrightarrow0
\]
par récurrence sur $j$, on en conclut une suite exacte
\[
0\longrightarrow\mathcal G'_x/
\left(\sum_{j=1}^{q}g_j\mathcal G'_x\right)
\longrightarrow\mathcal G_x/
\left(\sum_{j=1}^{q}g_j\mathcal G_x\right).
\]
Mais par hypothèse
$\operatorname{prof}(\mathcal G'_x/(\sum_{j=1}^{q}g_j\mathcal G'_x))=0$
(0, 16.4.6); par la caractérisation (0, 16.4.6) des modules de profondeur
nulle, on en conclut que
\[
\operatorname{prof}\left(\mathcal G_x/
\left(\sum_{j=1}^{q}g_j\mathcal G_x\right)\right)=0
=\operatorname{prof}\left(\mathcal F_x/
\left(\sum_{i=1}^{p}f_i\mathcal F_x+
\sum_{j=1}^{q}g_j\mathcal F_x\right)\right);
\]
la suite $(g_j)$ étant $\mathcal G'_x$-régulière et
$\mathcal G''_x$-régulière, est aussi $\mathcal G_x$-régulière; enfin, la
suite $((f_i)_x)$ est $\mathcal F_x$-régulière par hypothèse et est formée
d'éléments de l'idéal maximal $\mathfrak m_x$; on en déduit (0, 16.4.6)
que $\operatorname{prof}(\mathcal F_x)=p+q$, ce qui achève la
démonstration.
\end{proof}

\subsection{\texorpdfstring{Critères pour que les ensembles
$U_{S_n}(\mathcal F)$ ou $U_{C_n}(\mathcal F)$ soient ouverts}{Critères pour que les ensembles U\string_\string{S\string_n\string}(F) ou U\string_\string{C\string_n\string}(F) soient ouverts}}
\label{IV.2.s6.11-fr}

\begin{lemma}[6.11.1]
\label{IV.2.s6.11.1-fr}
Soit $X$ un préschéma localement noethérien, $\mathcal F$ un
$\mathcal O_X$-Module cohérent; alors la fonction
$x\mapsto\operatorname{dim.proj}(\mathcal F_x)$ est semi-continue
supérieurement dans $X$.
\end{lemma}

\begin{proof}
On peut se borner au cas où $X=\operatorname{Spec}(A)$ est le spectre d'un
anneau noethérien et $\mathcal F=\widetilde M$, où $M$ est un $A$-module de
type fini. Supposons que pour un $x\in X$, on ait
$\operatorname{dim.proj}(M_x)=n<+\infty$ (si $n=+\infty$ il n'y a rien à
démontrer); il existe une résolution de $M$
\[
L_{n-1}\longrightarrow L_{n-2}\longrightarrow\cdots
\longrightarrow L_0\longrightarrow M\longrightarrow0
\]

\oldpage[IV-2]{159}

où les $L_i$ sont des $A$-modules libres de type fini ($A$ étant
noethérien), d'où une suite exacte
\[
0\longrightarrow R\longrightarrow L_{n-1}\longrightarrow L_{n-2}
\longrightarrow\cdots\longrightarrow L_0\longrightarrow M
\longrightarrow0
\]
où $R$ est un $A$-module de type fini; on en déduit une suite exacte
\[
0\longrightarrow R_x\longrightarrow(L_{n-1})_x\longrightarrow\cdots
\longrightarrow(L_0)_x\longrightarrow M_x\longrightarrow0
\]
où les $(L_i)_x$ sont des $A_x$-modules libres de type fini; comme par
hypothèse $\operatorname{dim.proj}(M_x)=n$, cela entraîne que $R_x$ est un
$A_x$-module projectif de type fini (M, VI, 2.1) et par suite un
$A_x$-module libre de type fini (0$_{\rm III}$, 10.1.3). On en conclut qu'il
existe un voisinage ouvert $U$ de $x$ dans $X$ tel que, pour tout
$x'\in U$, $R_{x'}$ soit un $A_{x'}$-module libre (0$_{\rm I}$, 5.2.7),
donc $M_{x'}$ admet une résolution projective de longueur $n$, autrement
dit $\operatorname{dim.proj}(M_{x'})\leq n$, ce qui prouve le lemme.
\end{proof}

\begin{proposition}[6.11.2]
\label{IV.2.s6.11.2-fr}
Soit $X$ un préschéma localement noethérien et localement immersible dans
un schéma régulier (5.11.1); soit $\mathcal F$ un $\mathcal O_X$-Module
cohérent. Alors:
\begin{enumerate}
\item[(i)] (M. Auslander) La fonction
$x\mapsto\operatorname{coprof}(\mathcal F_x)$ est semi-continue
supérieurement dans $X$ (autrement dit, pour tout entier $n$, l'ensemble
$U_{C_n}(\mathcal F)$ des $x\in X$ tels que
$\operatorname{coprof}(\mathcal F_x)\leq n$ est ouvert).
\item[(ii)] Pour tout entier $n$, l'ensemble $U_{S_n}(\mathcal F)$ des
$x\in X$ tels que $\mathcal F$ possède la propriété $(S_n)$ au point $x$
est ouvert.
\end{enumerate}
\end{proposition}

\begin{proof}
\begin{enumerate}
\item[(i)] La question étant locale sur $X$, on peut, en vertu de
l'hypothèse, se borner au cas où $X$ est un sous-schéma fermé d'un schéma
affine régulier $Y$. Si $j:X\to Y$ est l'injection canonique, et
$\mathcal G=j_*(\mathcal F)$, on sait qu'on a alors, pour tout $x\in X$,
$\dim(\mathcal F_x)=\dim(\mathcal G_x)$ et
$\operatorname{prof}(\mathcal F_x)=\operatorname{prof}(\mathcal G_x)$,
donc $\operatorname{coprof}(\mathcal F_x)=
\operatorname{coprof}(\mathcal G_x)$, et comme
$\operatorname{coprof}(\mathcal G_y)=0$ pour $y\in Y-X$, on est ramené à
prouver la propriété pour $\mathcal G$; autrement dit, on peut se borner au
cas où $X$ est un schéma affine régulier. On sait alors (0, 17.3.4) que
l'on a
\[
\operatorname{prof}(\mathcal F_x)=\dim(\mathcal O_{X,x})-
\operatorname{dim.proj}(\mathcal F_x).
\]
D'autre part, si $S$ est l'unique sous-préschéma fermé de $X$ ayant pour
espace sous-jacent $\operatorname{Supp}(\mathcal F)$, on a
$\dim(\mathcal F_x)=\operatorname{codim}(\{x\},S)$ (5.1.12.1), et, en
vertu de (5.1.9)
\[
\dim(\mathcal F_x)=\dim(\mathcal O_{X,x})-\operatorname{codim}_x(S,X)
\]
puisque $\mathcal O_{X,x}$ est un anneau régulier, et a fortiori
biéquidimensionnel (0, 16.5.12). On peut donc écrire
\begin{equation}
\operatorname{coprof}(\mathcal F_x)=
\operatorname{dim.proj}(\mathcal F_x)-\operatorname{codim}_x(S,X)
\tag{6.11.2.1}\label{IV.2.s6.11.2.1-fr}
\end{equation}
et la proposition résulte alors de (6.11.1) et de (0, 14.2.6).

\oldpage[IV-2]{160}

\item[(ii)] Comme les $U_{C_n}(\mathcal F)$ sont ouverts pour tout $n$ par
(i), les $Z_n=X-U_{C_n}(\mathcal F)$ sont fermés dans $X$; en outre il est
clair que $Z_{n+1}\subset Z_n$, et comme $\dim(\mathcal F_x)$ est finie
pour tout $x\in X$ et
$\operatorname{coprof}(\mathcal F_x)\leq\dim(\mathcal F_x)$,
l'intersection des $Z_n$ est vide; comme on peut se borner au cas où $X$
est affine, donc quasi-compact, on peut supposer qu'il existe un $m$ tel
que $Z_m=\varnothing$. Or, il résulte de (5.7.4) que la relation
$x\in U_{S_k}(\mathcal F)$ équivaut à l'ensemble des relations
\begin{equation}
\operatorname{codim}_x(Z_n,S)>n+k
\tag{6.11.2.2}\label{IV.2.s6.11.2.2-fr}
\end{equation}
pour tout $n\geq0$; mais pour $n\geq m$ cette relation est automatiquement
vérifiée, donc on n'a en fait à considérer les relations (6.11.2.2) que
pour $0\leq n<m$. Or (0, 14.2.6) l'ensemble $V_{n,k}$ des $x$ vérifiant
(6.11.2.2) est ouvert, et $U_{S_k}(\mathcal F)$, intersection des $V_{n,k}$
pour $0\leq n<m$, est aussi ouvert. C.Q.F.D.
\end{enumerate}
\end{proof}

Rappelons que l'hypothèse que $X$ est localement immersible dans un schéma
régulier est toujours remplie lorsque $X$ est un préschéma localement de
type fini sur un corps $k$ (5.8.3).

\begin{corollary}[6.11.3]
\label{IV.2.s6.11.3-fr}
Sous les hypothèses de (6.11.2), l'ensemble $\operatorname{CM}(\mathcal F)$
des points $x\in X$ tels que $\mathcal F_x$ soit un module de
Cohen--Macaulay est ouvert dans $X$.
\end{corollary}

\begin{proof}
En effet, c'est l'ensemble $U_{C_0}(\mathcal F)$.
\end{proof}

\begin{remarks}[6.11.4]
\label{IV.2.s6.11.4-fr}
\begin{enumerate}
\item[(i)] Le raisonnement de (6.11.2, (ii)) prouve que (sans hypothèse sur
$X$) lorsque $U_{C_n}(\mathcal F)$ est ouvert pour tout entier $n$, alors
$U_{S_n}(\mathcal F)$ est ouvert pour tout entier $n$.
\item[(ii)] On a
$\operatorname{CM}(\mathcal F)=\bigcap_nU_{S_n}(\mathcal F)$. Si tout point
$x\in X$ admet un voisinage ouvert $V$ qui soit de dimension finie, la
suite des intersections $V\cap U_{S_n}(\mathcal F)$ est stationnaire
puisqu'il existe un $m$ tel que $\dim(\mathcal F_x)\leq m$ pour tout
$x\in V$; par suite, si les $U_{S_n}(\mathcal F)$ sont ouverts pour tout
entier $n$, $\operatorname{CM}(\mathcal F)$ est alors ouvert dans $X$.
\item[(iii)] On écrira $U_{S_n}(X)$, $U_{C_n}(X)$,
$\operatorname{CM}(X)$ au lieu de $U_{S_n}(\mathcal O_X)$,
$U_{C_n}(\mathcal O_X)$, $\operatorname{CM}(\mathcal O_X)$.
\end{enumerate}
\end{remarks}

\begin{proposition}[6.11.5]
\label{IV.2.s6.11.5-fr}
Soient $A$ un anneau local noethérien, $M$ un $A$-module de type fini. Pour
tout idéal premier $\mathfrak p$ de $A$, on a
\begin{equation}
\operatorname{coprof}_{A_{\mathfrak p}}(M_{\mathfrak p})\leq
\operatorname{coprof}_A(M).
\tag{6.11.5.1}\label{IV.2.s6.11.5.1-fr}
\end{equation}
\end{proposition}

\begin{proof}
On peut se borner au cas où $M_{\mathfrak p}\ne0$. Comme $\widehat A$ est
un $A$-module fidèlement plat (0$_{\rm I}$, 7.3.5), il existe un idéal
premier $\mathfrak q$ de $\widehat A$ au-dessus de $\mathfrak p$
(0$_{\rm I}$, 6.5.1); comme
$\widehat M_{\mathfrak q}=(M\otimes_A\widehat A)_{\mathfrak q}
=M\otimes_A\widehat A_{\mathfrak q}$ et que
$\widehat A_{\mathfrak q}$ est un $\widehat A$-module plat, donc un
$A_{\mathfrak p}$-module plat (0$_{\rm I}$, 6.2.1), il résulte de (6.3.2),
appliqué à l'homomorphisme local
$A_{\mathfrak p}\to\widehat A_{\mathfrak q}$, à $M_{\mathfrak p}$ et
$\widehat M_{\mathfrak q}$, que l'on a
$\operatorname{coprof}_{A_{\mathfrak p}}(M_{\mathfrak p})\leq
\operatorname{coprof}_{\widehat A_{\mathfrak q}}
(\widehat M_{\mathfrak q})$. D'autre part,
$X=\operatorname{Spec}(\widehat A)$ est isomorphe à un sous-schéma d'un
schéma régulier en vertu du théorème de Cohen (0, 19.8.8, (i)). On déduit
donc de (6.11.2) que l'on a
$\operatorname{coprof}_{\widehat A_{\mathfrak q}}
(\widehat M_{\mathfrak q})\leq
\operatorname{coprof}_{\widehat A}(\widehat M)$; enfin, on sait que
$\operatorname{coprof}_{\widehat A}(\widehat M)=
\operatorname{coprof}_A(M)$ (0, 16.4.10), ce qui achève de démontrer
(6.11.5).
\end{proof}

Cette proposition justifie la définition de la coprofondeur d'un
$A$-module lorsque $A$ n'est pas un anneau local, donnée dans (5.7.12).

\begin{proposition}[6.11.6]
\label{IV.2.s6.11.6-fr}
Soient $X$ un préschéma localement noethérien, $\mathcal F$ un
$\mathcal O_X$-Module cohérent, $n$ un entier $>0$. Supposons que pour tout
sous-préschéma fermé intègre $Y$ de $X$, il existe une partie ouverte non
vide $W$ de $Y$ telle que le sous-préschéma de $Y$ induit sur l'ouvert $W$
vérifie $(S_n)$. Alors l'ensemble $U_{S_n}(\mathcal F)$ est ouvert dans
$X$.
\end{proposition}

\begin{proof}
La question étant locale sur $X$, on peut se borner au cas où $X$ est
noethérien. Nous allons raisonner par récurrence sur $n$: pour $n=1$,
l'ensemble $U_{S_1}(\mathcal F)$ est ouvert dans $X$, car l'ensemble des
$x\in X$ où $\mathcal F$ ne vérifie pas $(S_1)$ est l'ensemble des $x$ tels
que $\mathcal F_x$ admette des idéaux premiers associés immergés (5.7.5);
si $(Z_j)$ est la famille finie des cycles premiers associés à $\mathcal F$
qui sont immergés, on a par suite
$U_{S_1}(\mathcal F)=X-\bigcup_jZ_j$, d'où notre assertion puisque les
$Z_j$ sont fermés. Nous supposerons donc désormais que $n>1$. En second
lieu, on peut se borner au cas où
$\operatorname{Supp}(\mathcal F)=X$, car il y a un sous-préschéma fermé
$T$ de $X$ ayant $\operatorname{Supp}(\mathcal F)$ pour espace sous-jacent,
et un $\mathcal O_T$-Module cohérent $\mathcal G$ tels que
$\mathcal F=j_*(\mathcal G)$, $j:T\to X$ étant l'injection canonique
(Err$_{\rm III}$, 30); comme il revient au même de prouver que $\mathcal F$
ou $\mathcal G$ vérifie $(S_n)$ en un point
$x\in\operatorname{Supp}(\mathcal F)$, on peut se borner à considérer le
cas où $T=X$. Notons enfin que par définition (5.7.2),
$U_{S_n}(\mathcal F)$ est stable par générisation. Nous allons utiliser le
lemme suivant:

\begin{lemma}[6.11.6.1]
\label{IV.2.s6.11.6.1-fr}
Soient $X$ un espace noethérien dans lequel toute partie fermée irréductible
admet un point générique, $E$ une partie de $X$. Pour que $E$ soit ouverte
dans $X$, il faut et il suffit que $E$ soit stable par générisation, et que,
pour toute partie ouverte $V$ de $X$ et toute partie irréductible $Y$ fermée
dans $V$, telles que $V-Y\subset E$ et que le point générique de $Y$
appartienne à $E$, $E\cap Y$ contienne une partie ouverte non vide de $Y$.
\end{lemma}

\begin{proof}
On observera que ce critère entraîne celui de (0$_{\rm III}$, 9.2.6)
lorsque toute partie fermée irréductible de $X$ admet un point générique.
Il n'y a évidemment à démontrer que la suffisance des conditions énoncées.

Considérons l'intérieur $U$ de $E$; l'ensemble fermé $X-U$ est réunion de
ses composantes irréductibles, qui sont en nombre fini et fermées dans $X$.
Si on avait $E\ne U$, l'hypothèse que $E$ est stable par générisation
entraînerait que le point générique $z$ d'une des composantes irréductibles
de $X-U$ appartiendrait à $E$. Or, $z$ n'appartient qu'à une seule des
composantes irréductibles de $X-U$; si $T$ est la réunion des autres
composantes irréductibles de $X-U$, $V=X-T$ est ouvert dans $X$, réunion de
$U$ et de l'ensemble $Y=\overline{\{z\}}\cap V$ fermé dans $V$ et
irréductible. Par hypothèse $E\cap Y$ contient une partie $W$ ouverte dans
$Y$; on en conclut que $U\cup W$ est ouvert dans $V$, donc dans $X$, ce qui
est absurde puisque $U$ est supposé être l'intérieur de $E$.
\end{proof}

\oldpage[IV-2]{161}

En vertu de ce lemme, on peut supposer qu'il existe dans $X$ un
sous-préschéma intègre $Y$ de point générique $y$ tel que $\mathcal F$
vérifie $(S_n)$ au point $y$ et en tous les points de $X-Y$, et tout se
borne à vérifier qu'il existe dans $X$ un voisinage ouvert de $y$ tel que
$\mathcal F$ vérifie $(S_n)$ dans ce voisinage. Distinguons alors deux cas:

$1^\circ$ $y$ est point maximal de $X$; comme il existe un voisinage ouvert
de $y$ ne rencontrant aucune composante irréductible de $X$ autre que
$\overline{\{y\}}$, on peut supposer que $X$ est irréductible, donc a même
espace sous-jacent que $Y$, de sorte que $Y$ est défini par le Nilradical de
$\mathcal O_X$, qui est nilpotent. D'autre part, on peut, en remplaçant $X$
par un voisinage ouvert de $y$, supposer que $\mathcal F$ est normalement
plat le long de $Y$ (6.9.1); il résulte alors de (6.10.4) que
$U_{S_n}(\mathcal F)=U_{S_n}(\mathcal O_Y)$, et comme ce dernier est par
hypothèse un voisinage de $y$ dans $X$, cela prouve la proposition dans ce
cas.

$2^\circ$ $y$ n'est pas un point maximal de $X$, autrement dit (puisque
$\operatorname{Supp}(\mathcal F)=X$), $\dim(\mathcal F_y)\geq1$, donc aussi
$\operatorname{prof}(\mathcal F_y)\geq1$ puisque $\mathcal F$ vérifie par
hypothèse $(S_n)$ (et a fortiori $(S_1)$) au point $y$. Remplaçant au besoin
$X$ par un voisinage ouvert de $y$, on peut donc supposer qu'il existe une
section $f$ de $\mathcal O_X$ au-dessus de $X$, $\mathcal F$-régulière et
telle que $f_y\in\mathfrak m_y$, ou encore $f(y)=0$ (0, 15.2.4); on a donc
encore $f(x)=0$ pour tout $x\in Y$ (0$_{\rm I}$, 5.5.2). On sait que
$\mathcal F/f\mathcal F$ vérifie $(S_{n-1})$ au point $y$ (5.7.6).
Appliquant l'hypothèse de récurrence, et remplaçant au besoin $X$ par un
voisinage ouvert de $y$, on peut donc supposer que
$\mathcal F/f\mathcal F$ vérifie $(S_{n-1})$ en tout point de $X$. Mais pour
tout $x\in Y$, la relation $f(x)=0$ entraîne que l'on a
$\operatorname{prof}(\mathcal F_x/f_x\mathcal F_x)=
\operatorname{prof}(\mathcal F_x)-1$ et
$\dim(\mathcal F_x/f_x\mathcal F_x)=\dim(\mathcal F_x)-1$
(0, 16.3.4 et 16.4.6); la relation
\[
\operatorname{prof}(\mathcal F_x/f_x\mathcal F_x)\geq
\inf(n-1,\dim(\mathcal F_x/f_x\mathcal F_x))
\]
est donc équivalente à
\[
\operatorname{prof}(\mathcal F_x)\geq\inf(n,\dim(\mathcal F_x)).
\]
Comme on a supposé que $\mathcal F$ vérifie $(S_n)$ en tout point de $X-Y$,
cela termine la démonstration.
\end{proof}

\begin{corollary}[6.11.7]
\label{IV.2.s6.11.7-fr}
Les notations étant celles de (6.11.6):
\begin{enumerate}
\item[(i)] L'ensemble $U_{S_1}(\mathcal F)$ est ouvert dans $X$.
\item[(ii)] Pour que l'ensemble $U_{S_2}(\mathcal F)$ soit ouvert il suffit
que tout point maximal $x$ de $\operatorname{Supp}(\mathcal F)$,
appartenant à $U_{S_2}(\mathcal F)$, soit intérieur à
$U_{S_2}(\mathcal F)$.
\end{enumerate}
\end{corollary}

\begin{proof}
L'assertion (i) a été prouvée dans le cours de la démonstration de (6.11.6);
d'autre part, pour $n=2$ le cas $2^\circ$ de la démonstration de (6.11.6)
est valable sans aucune hypothèse sur $X$, puisque (avec les mêmes
notations) $U_{S_1}(\mathcal F)$ et
$U_{S_1}(\mathcal F/f\mathcal F)$ sont ouverts dans $X$. Quant au cas
$1^\circ$ de cette démonstration, l'hypothèse assure précisément qu'il est
inutile de le considérer.
\end{proof}

\oldpage[IV-2]{162}

\begin{proposition}[6.11.8]
\label{IV.2.s6.11.8-fr}
Soit $X$ un préschéma localement noethérien vérifiant la propriété suivante:
\begin{description}
\item[(CMU)] Tout sous-préschéma fermé intègre $Y$ de $X$ contient un ouvert
non vide $W$ tel que le préschéma induit par $Y$ sur $W$ soit un préschéma
de Cohen--Macaulay.
\end{description}
Alors, pour tout $\mathcal O_X$-Module cohérent $\mathcal F$, la fonction
$x\mapsto\operatorname{coprof}(\mathcal F_x)$ est localement constructible
et semi-continue supérieurement; les ensembles $U_{C_n}(\mathcal F)$ et
$U_{S_n}(\mathcal F)$ sont ouverts dans $X$.
\end{proposition}

\begin{proof}
En effet, soit $Y$ un sous-préschéma fermé intègre de $X$, de point
générique $y$; en vertu de (6.10.6), il y a dans $Y$ un voisinage ouvert
$V$ de $y$ tel que, pour tout $x\in V\cap Y$, on ait
\begin{equation}
\operatorname{coprof}(\mathcal F_x)=
\operatorname{coprof}(\mathcal F_y)+
\operatorname{coprof}(\mathcal O_{Y,x}).
\tag{6.11.8.1}\label{IV.2.s6.11.8.1-fr}
\end{equation}
Mais par hypothèse il existe un ouvert non vide $W$ de $Y$ tel que, pour
$x\in V\cap W$, on ait $\operatorname{coprof}(\mathcal O_{Y,x})=0$, donc
$\operatorname{coprof}(\mathcal F_x)$ est constante dans un voisinage de
$y$ dans $Y$, ce qui prouve que la fonction
$x\mapsto\operatorname{coprof}(\mathcal F_x)$ est localement constructible
(0$_{\rm III}$, 9.3.2); en outre il résulte alors de (6.11.5) et de
(0$_{\rm III}$, 9.3.4) que cette fonction est semi-continue supérieurement.
La dernière assertion résulte de (6.11.4, (i)), ou encore de (6.11.6).
\end{proof}

\begin{remarks}[6.11.9]
\label{IV.2.s6.11.9-fr}
\begin{enumerate}
\item[(i)] Si $X$ vérifie l'hypothèse (CMU) de (6.11.8), alors, pour tout
morphisme $u:X'\to X$ localement de type fini, $X'$ vérifie aussi (CMU).
Soient en effet $Y'$ un sous-préschéma fermé intègre de $X'$, $y'$ son point
générique, $y=u(y')$, et soit $Y$ le sous-préschéma intègre de $X$ ayant
pour espace sous-jacent $\overline{\{y\}}$; alors $u|Y'$ se factorise en
$Y'\xrightarrow{v}Y\xrightarrow{j}X$, où $j$ est l'injection canonique
(I, 5.2.2), et $v$ est localement de type fini (I, 6.6.6). En se
restreignant à des voisinages ouverts affines de $y$ et $y'$ respectivement,
on peut donc se borner au cas où $X=\operatorname{Spec}(A)$, où $A$ est un
anneau intègre de Cohen--Macaulay, et $X'=\operatorname{Spec}(A')$, où $A'$
est un anneau intègre contenant $A$ et qui est une $A$-algèbre de type fini.

\oldpage[IV-2]{163}

Remplaçant au besoin $A$ par un anneau de fractions $A_f$ (avec $f\ne0$),
on peut en outre supposer que $A'$ contient un anneau de polynômes
$A[T_1,\ldots,T_n]=A''$, et est une $A''$-algèbre finie (Bourbaki,
\emph{Alg. comm.}, chap. V, \S3, n\textsuperscript{o} 1, cor. 1 du th. 1).
Mais $A''$ est un anneau de Cohen--Macaulay (6.3.6); donc on peut se borner
au cas où de plus $A'$ est une $A$-algèbre finie. Il y a alors $g\ne0$ dans
$A$ tel que $A'_g$ soit un $A_g$-module libre de type fini (Bourbaki,
\emph{Alg. comm.}, chap. II, \S5, n\textsuperscript{o} 1, cor. de la
prop. 2), donc on peut encore supposer de plus que $A'$ soit un $A$-module
libre. Mais alors $A'$ est un $A$-module de Cohen--Macaulay (0, 16.5.1), et
comme $A'$ est un $A$-module de type fini, $A'$ est aussi un $A'$-module de
Cohen--Macaulay (0, 16.5.3), donc un anneau de Cohen--Macaulay.
\item[(ii)] Supposons qu'il existe un $\mathcal O_X$-Module cohérent
$\mathcal F$ tel que $\operatorname{Supp}(\mathcal F)=X$ et que
$\mathcal F$ soit un $\mathcal O_X$-Module de Cohen--Macaulay. Alors $X$
vérifie la condition (CMU): en effet, avec les notations de la démonstration
de (6.11.8), la relation (6.11.8.1) montre que l'on a
$\operatorname{coprof}(\mathcal O_{Y,x})=0$ dans un voisinage (par rapport
à $Y$) du point générique $y$ de $Y$.
\end{enumerate}

On ignore s'il existe des préschémas localement noethériens de dimension
$\geq2$ qui ne vérifient pas (CMU) (si $\dim(X)=1$, il est immédiat que tout
point maximal de $X_{\mathrm{red}}$ admet un voisinage ouvert intègre de
dimension $1$, donc de Cohen--Macaulay).
\end{remarks}

\subsection{\texorpdfstring{Critères de Nagata pour que $\operatorname{Reg}(X)$ soit ouvert}{Critères de Nagata pour que Reg(X) soit ouvert}}
\label{IV.2.s6.12-fr}

\begin{env}[6.12.1]
\label{IV.2.s6.12.1-fr}
Étant donné un préschéma localement noethérien $X$, on appelle lieu
singulier de $X$ et on note $\operatorname{Sing}(X)$ l'ensemble des points
$x\in X$ tels que $X$ ne soit pas régulier au point $x$ (autrement dit,
tels que l'anneau local $\mathcal O_{X,x}$ ne soit pas régulier); le
complémentaire $X-\operatorname{Sing}(X)$, ensemble des $x\in X$ où $X$ est
régulier, se note $\operatorname{Reg}(X)$. Nous nous proposons dans ce
numéro de donner des conditions moyennant lesquelles
$\operatorname{Sing}(X)$ est fermé (i.e. $\operatorname{Reg}(X)$ ouvert).
\end{env}

\begin{proposition}[6.12.2]
\label{IV.2.s6.12.2-fr}
Soit $X$ un préschéma localement noethérien. Les conditions suivantes sont
équivalentes:
\begin{enumerate}
\item[a)] $\operatorname{Reg}(X)$ est ouvert dans $X$.
\item[b)] Pour tout $x\in\operatorname{Reg}(X)$, il existe une partie
ouverte non vide dans $\overline{\{x\}}$ contenue dans
$\operatorname{Reg}(X)$.
\end{enumerate}
En outre, ces conditions sont entraînées par la suivante:
\begin{enumerate}
\item[c)] Pour tout $x\in\operatorname{Reg}(X)$, si on désigne par $Y$ le
sous-préschéma fermé réduit de $X$ ayant $\overline{\{x\}}$ pour espace
sous-jacent, $\operatorname{Reg}(Y)$ est un voisinage de $x$ dans $Y$.
\end{enumerate}
\end{proposition}

\begin{proof}
L'équivalence de a) et b) résulte de ce que $\operatorname{Reg}(X)$ est
stable par générisation (0, 17.3.2) et de (0$_{\rm III}$, 9.2.6). Pour
prouver que c) entraîne b), on peut se borner au cas où
$X=\operatorname{Spec}(A)$ est un voisinage ouvert affine de $x$; si
$(t_i)_{1\leq i\leq n}$ est un système régulier de paramètres (0, 17.1.6)
de l'anneau local régulier $A_x$, on peut supposer (en remplaçant au besoin
$X$ par un voisinage ouvert de $x$) que $t_i=(s_i)_x$ pour tout $i$, où
$s_i\in A$ et où la famille $(s_i)_{1\leq i\leq n}$ est $A$-régulière
(0, 15.2.4). On a alors $Y=\operatorname{Spec}(A/\mathfrak p)$, où
$\mathfrak p=\mathfrak j_x$; comme les $t_i$ engendrent l'idéal maximal
$\mathfrak m=\mathfrak pA_x$ de $A_x$, on peut encore supposer (en
remplaçant $X$ par un voisinage plus petit de $x$) que les $s_i$ engendrent
$\mathfrak p$. Pour tout $y\in Y$, les $(s_i)_y$ engendrent donc
$\mathfrak p_y$; comme ils forment une suite $A_y$-régulière, ils font
partie d'un système de paramètres de $A_y$ (0, 16.4.1); donc on déduit de
(0, 17.1.7) que si $(A/\mathfrak p)_y=A_y/\mathfrak p_y$ est régulier, il
en est de même de $A_y$, d'où la conclusion.
\end{proof}

\begin{corollary}[6.12.3]
\label{IV.2.s6.12.3-fr}
Soit $X$ un préschéma localement noethérien. Les conditions suivantes sont
équivalentes:
\begin{enumerate}
\item[a)] Pour tout sous-préschéma $Y$ de $X$,
$\operatorname{Reg}(Y)$ est ouvert dans $Y$.

\oldpage[IV-2]{164}

\item[b)] Pour tout sous-préschéma fermé intègre $Y$ de $X$,
$\operatorname{Reg}(Y)$ contient une partie ouverte non vide de $Y$.
\end{enumerate}
\end{corollary}

\begin{proof}
Il est clair que a) entraîne b), car si $Y$ est intègre et $y$ est son point
générique, $\mathcal O_{Y,y}$ est un corps, donc $y\in\operatorname{Reg}(Y)$.
Inversement, pour voir que b) entraîne a), considérons un sous-préschéma
fermé intègre $Z$ de $Y$ de point générique $z$; si $Y'$ est le
sous-préschéma intègre de $X$ ayant pour espace sous-jacent l'adhérence
$\overline{\{z\}}$ de $z$ dans $X$, $Z$ est ouvert dans $Y'$ et le
sous-préschéma $Z$ de $Y'$, étant réduit, est induit par $Y'$ sur l'ouvert
$Z$ de l'espace sous-jacent à $Y'$. Or l'hypothèse entraîne que
$\operatorname{Reg}(Y')$ est un voisinage de $z$ dans $Y'$, donc
$\operatorname{Reg}(Z)$ est un voisinage de $z$ dans $Z$; il suffit alors
d'appliquer (6.12.2) en y remplaçant $X$ par $Y$ et $Y$ par $Z$.
\end{proof}

\begin{theorem}[6.12.4]
\label{IV.2.s6.12.4-fr}
(Nagata). Soient $A$ un anneau noethérien,
$X=\operatorname{Spec}(A)$. Les conditions suivantes sont équivalentes:
\begin{enumerate}
\item[a)] Pour tout préschéma $X'$ localement de type fini sur $X$,
$\operatorname{Reg}(X')$ est ouvert dans $X'$.
\item[b)] Pour toute $A$-algèbre $A'$ finie et intègre, il existe une
partie ouverte non vide de $X'=\operatorname{Spec}(A')$ contenue dans
$\operatorname{Reg}(X')$.
\item[c)] Pour tout idéal premier $\mathfrak p$ de $A$ et toute extension
finie radicielle $K'$ du corps des fractions $K$ de $A/\mathfrak p$, il
existe une sous-$A$-algèbre $A'$ de $K'$, finie sur $A$, ayant $K'$ pour
corps des fractions, et telle qu'il existe dans
$X'=\operatorname{Spec}(A')$ un ensemble ouvert non vide contenu dans
$\operatorname{Reg}(X')$.
\end{enumerate}
\end{theorem}

\begin{proof}
Il est clair que a) implique b). Pour voir que b) entraîne c), il suffit de
remarquer qu'on peut prendre pour générateurs de l'extension $K'$ de $K$
des éléments entiers sur $A/\mathfrak p$ (et a fortiori sur $A$), et comme
ces éléments sont en nombre fini, ils engendrent une $A$-algèbre finie
$A'$ dont $K'$ est le corps des fractions; on peut alors appliquer b) à
$A'$. Reste donc à prouver que c) entraîne a). La question étant locale sur
$X'$, on peut d'abord supposer que $X'=\operatorname{Spec}(A')$, où $A'$
est une $A$-algèbre de type fini; compte tenu de (6.12.2), il suffit de
prouver que pour tout sous-préschéma fermé intègre $Y'$ de $X'$,
$\operatorname{Reg}(Y')$ contient une partie ouverte non vide de $Y'$.
Autrement dit, on peut se borner à prouver que si $A'$ est une $A$-algèbre
intègre de type fini et $X'=\operatorname{Spec}(A')$,
$\operatorname{Reg}(X')$ contient une partie ouverte non vide de $X'$.
Soit $K'$ le corps des fractions de $A'$; si $\mathfrak p$ est l'image
canonique dans $X$ du point générique de $X'$, $K'$ est une extension du
corps des fractions $K$ de $A/\mathfrak p$, et $A/\mathfrak p$
s'identifie à un sous-anneau de $A'$, $A'$ étant une
$(A/\mathfrak p)$-algèbre de type fini. On peut évidemment se borner dans
la suite de la démonstration au cas $\mathfrak p=0$. Distinguons maintenant
deux cas:

I) $K'$ est une extension séparable de $K$. --- On est alors ramené à
prouver le

\begin{lemma}[6.12.4.1]
\label{IV.2.s6.12.4.1-fr}
Soient $A$ un anneau intègre noethérien, $A'$ une $A$-algèbre intègre de
type fini, contenant $A$ et telle que le corps des fractions $K'$ de $A'$
soit une extension séparable du corps des fractions $K$ de $A$. Si
$\operatorname{Spec}(A)$ contient un ouvert non vide formé de points
réguliers, il en est de même de $\operatorname{Spec}(A')$.
\end{lemma}

\begin{proof}
Remplaçant $A$ par un anneau de fractions $A_f$ de sorte que
$D(f)\subset\operatorname{Reg}(\operatorname{Spec}(A))$, on peut déjà
supposer l'anneau $A$ régulier. Par hypothèse il existe dans $K'$ un système
$(t_i)_{1\leq i\leq n}$ d'éléments algébriquement indépendants sur $K$ et
tels que $K'$ soit une extension algébrique séparable finie de
$K(t_1,\ldots,t_n)$; en considérant un système fini
\oldpage[IV-2]{165}
de générateurs $t'_j$ $(1\leq j\leq m)$ de $K'$ sur
$K(t_1,\ldots,t_n)$, que l'on peut supposer entiers sur
$A_1=A[t_1,\ldots,t_n]$, on voit que
$A'_1=A_1[t'_1,\ldots,t'_m]$ est fini sur $A_1$ et a pour corps de fractions
$K'$. Si on pose $X'=\operatorname{Spec}(A')$,
$X'_1=\operatorname{Spec}(A'_1)$, il résulte de ce que les corps de fonctions
rationnelles $R(X')$ et $R(X'_1)$ sont tous deux isomorphes à $K'$ et de ce
que $X'$ et $X'_1$ sont des $A$-préschémas de type fini, qu'il existe un
ouvert $U'\subset X'$ et un ouvert $U'_1\subset X'_1$ qui sont
$A$-isomorphes (I, 6.5.5). On est donc ramené à prouver que
$\operatorname{Reg}(X'_1)$ contient un ouvert non vide, autrement dit on
peut supposer que $A'$ est une $A_1$-algèbre finie. Or on sait (0, 17.3.7)
que $A_1$ est un anneau régulier, et l'on peut donc se borner au cas où $A'$
est une $A$-algèbre finie et $K'$ une extension séparable et finie de $K$.
Si $\xi$ est le point générique de $X$, $A'_\xi=K'$ est alors un module
libre sur $A_\xi=K$, donc (Bourbaki, \emph{Alg. comm.}, chap. II, \S5,
n\textsuperscript{o} 1,
cor. de la prop. 2) on peut, en remplaçant au besoin $A$ par un $A_f$,
supposer que $A'$ est un $A$-module libre de type fini. Soit alors
$(x_h)_{1\leq h\leq r}$ une base de ce $A$-module, et posons
\[
d=\det(\operatorname{Tr}_{A'/A}(x_hx_k))=
\det(\operatorname{Tr}_{K'/K}(x_hx_k))\in A.
\]
Comme $K'$ est séparable sur $K$, on sait (Bourbaki, \emph{Alg.}, chap. IX, \S2,
prop. 5) que $d\neq0$; remplaçant au besoin $A$ par l'anneau de fractions
$A_d$, on peut supposer $d$ inversible dans $A$. Mais alors, pour tout
$z\in\operatorname{Spec}(A)$, si l'on désigne par $\bar x_h$
$(1\leq h\leq r)$ l'image canonique de $x_h$ dans
$A'(z)=A'\otimes_A k(z)$, on a
$\det(\operatorname{Tr}_{A'(z)/k(z)}(\bar x_h\bar x_k))=\bar d$, image
canonique de $d$ dans $k(z)=A_z/\mathfrak m_z$ et comme $\bar d$ est
inversible (donc $\neq0$) dans $k(z)$, on sait (loc. cit.) que $A'(z)$ est
une $k(z)$-algèbre séparable, donc composée directe de corps extensions
finies et séparables de $k(z)$. Une telle algèbre étant un anneau régulier,
on voit que le morphisme $g:X'\to X$ est plat et que ses fibres $g^{-1}(z)$
sont régulières pour tout $z\in X$; on conclut alors de (6.5.2, (ii)) que
$X'$ est régulier, ce qui termine dans ce cas la démonstration.
\end{proof}

II) Cas général. --- Comme $A'$ est un $A$-module sans torsion,
$A'\otimes_AK$ s'identifie à un sous-anneau de $K'$, donc
$X''=\operatorname{Spec}(A'\otimes_AK)$ est un $K$-préschéma intègre, dont
$K'$ est le corps des fonctions rationnelles. On sait (4.6.6) qu'il existe
une extension radicielle finie $K_1$ de $K$ telle que si
$X'_1=\operatorname{Spec}(A'\otimes_AK_1)=X''\otimes_KK_1$,
$(X'_1)_{\mathrm{red}}$ soit un $K_1$-préschéma géométriquement réduit de
type fini; en outre, le morphisme $\operatorname{Spec}(K_1)\to
\operatorname{Spec}(K)$ étant radiciel, fini et surjectif, est un
homéomorphisme universel (2.4.5), donc $X'_1$ est homéomorphe à $X''$, et
par suite $(X'_1)_{\mathrm{red}}$ est intègre; en outre, son corps de
fonctions rationnelles $K'_1$ est une extension radicielle finie de $K'$ et
une extension séparable de type fini de $K_1$ (4.6.1). En vertu de
l'hypothèse c) de l'énoncé, il y a une sous-$A$-algèbre finie $A_1$ de
$K_1$, ayant $K_1$ pour corps des fractions et telle que si l'on pose
$X_1=\operatorname{Spec}(A_1)$, $\operatorname{Reg}(X_1)$ contienne un
ouvert non vide de $X_1$. Soit $A'_1$ l'image de l'homomorphisme canonique
$A'\otimes_AA_1\to K'_1$ et soit $X'_1=\operatorname{Spec}(A'_1)$; $A'_1$
est un anneau intègre qui est une $A'$-algèbre finie et dont le corps des
fractions est $K'_1$ par construction; en outre, comme l'homomorphisme
composé $A'\to A'\otimes_AA_1\to A'_1\to K'_1$ est identique à
$A'\to K'\to K'_1$, donc injectif, l'homomorphisme $A'\to A'_1$ est
injectif; le morphisme $g:X'_1\to X'$ est donc fini et surjectif
(II, 6.1.10). Cela étant, l'hypothèse sur $A_1$ et la partie I) de la
démonstration entraînent que $\operatorname{Reg}(X'_1)$
\oldpage[IV-2]{166}
contient un ouvert non vide $V'_1$; comme $g$ est fermé (II, 6.1.10), on
peut supposer que $V'_1=g^{-1}(V')$, où $V'$ est un ouvert affine de $X'$;
remplaçant $A'$ par l'anneau de $V'$ et $A'_1$ par celui de $V'_1$, on peut
donc supposer que $X'_1$ est régulier. En outre, le même raisonnement que
dans I) appliqué au point générique $\xi'$ de $X'$ permet (en remplaçant au
besoin $X'$ par un voisinage de $\xi'$) de supposer que $A'_1$ est un
$A'$-module libre, et par suite que le morphisme $g$ est plat. Mais il
résulte alors de (6.5.2, (i)) que $X'$ est régulier. C.Q.F.D.
\end{proof}

\begin{corollary}[6.12.5]
\label{IV.2.s6.12.5-fr}
\textnormal{(Zariski).} --- Soit $k$ un corps; pour tout $k$-préschéma $X$
localement de type fini sur $k$, l'ensemble des $x\in X$ où $X$ est
régulier (resp. géométriquement régulier sur $k$) est ouvert dans $X$.
\end{corollary}

\begin{proof}
L'assertion du corollaire relative à la propriété d'être régulier résulte de
(6.12.4) en prenant $A=k$. L'assertion relative à la propriété d'être
géométriquement régulier résulte déjà de (6.8.7); on peut aussi la déduire de
(6.12.4) de la façon suivante: posons $k'=k^{p^{-\infty}}$ ($p$ exposant
caractéristique de $k$); comme le morphisme
$\operatorname{Spec}(k')\to\operatorname{Spec}(k)$ est radiciel, entier et
surjectif, c'est un homéomorphisme universel (2.4.5), donc le morphisme
projection $X\otimes_kk'\to X$ est un homéomorphisme. La projection dans
$X$ de $\operatorname{Reg}(X\otimes_kk')$ est l'ensemble des points de $X$
où $X$ est géométriquement régulier sur $k$, en vertu de (6.7.7, e)), donc
cet ensemble est ouvert d'après ce qui précède.
\end{proof}

\begin{corollary}[6.12.6]
\label{IV.2.s6.12.6-fr}
Soit $A$ un anneau ayant l'une des propriétés suivantes:
\begin{enumerate}
\item[(i)] $A$ est un anneau de Dedekind et son corps des fractions $K$ est
de caractéristique 0.
\item[(ii)] $A$ est un anneau semi-local noethérien de dimension $\leq1$.
\end{enumerate}
Alors, pour tout préschéma $X$ localement de type fini sur $A$,
$\operatorname{Reg}(X)$ est ouvert dans $X$.
\end{corollary}

\begin{proof}
Vérifions dans les deux cas la condition c) de (6.12.4). Dans les deux cas,
un idéal premier $\mathfrak p$ de $A$ est maximal ou minimal; si
$\mathfrak p$ est maximal, une $(A/\mathfrak p)$-algèbre finie et intègre
est un corps et la condition c) de (6.12.4) est trivialement vérifiée.
Supposons donc $\mathfrak p$ non maximal, et distinguons les deux cas de
l'énoncé.

(i) Si $K$ est de caractéristique 0, il n'y a pas d'autre extension
radicielle de $K$ que $K$ lui-même; comme un anneau de Dedekind est régulier
(0, 17.1.4), la condition c) de (6.12.4) est trivialement vérifiée.

(ii) On peut alors supposer $A$ intègre (6.12.2), soit $K$ son corps des
fractions; si $K'$ est une extension radicielle finie de $K$, $A'$ une
sous-$A$-algèbre de $K'$ engendrée par un système fini de générateurs de
$K'$ sur $K$, entiers sur $A$, $A'$ est un anneau semi-local intègre de
dimension 1 (0, 16.1.5), et par suite, dans
$X'=\operatorname{Spec}(A')$, l'ensemble réduit au point générique est
ouvert et évidemment contenu dans $\operatorname{Reg}(X')$, ce qui prouve
la condition c) de (6.12.4) dans ce cas.
\end{proof}

Ce corollaire s'applique notamment lorsque $A=\mathbf Z$.

\begin{theorem}[6.12.7]
\label{IV.2.s6.12.7-fr}
\textnormal{(Nagata).} --- Soient $A$ un anneau local noethérien complet,
$X=\operatorname{Spec}(A)$. Alors $\operatorname{Reg}(X)$ est ouvert dans
$X$.
\end{theorem}

\begin{proof}
Compte tenu de (6.12.2), on est ramené au cas où de plus $A$ est intègre, et
à prouver que dans ce cas $\operatorname{Reg}(X)$ contient une partie
ouverte non vide de $X$. Distinguons deux cas:

\oldpage[IV-2]{167}
I) Le corps des fractions $K$ de $A$ est de caractéristique 0. --- On sait
alors (0, 19.8.8) qu'il existe un anneau de valuation discrète complet $C$
et un sous-anneau $B$ de $A$ tel que $A$ soit une $B$-algèbre finie et que
$B$ soit isomorphe à un anneau de séries formelles
$C[[T_1,\ldots,T_n]]$. Comme $C$ est régulier (II, 7.1.6), il en est de
même de $B$ (0, 17.3.8); en outre, le corps des fractions $L$ de $B$ étant
de caractéristique 0, $K$ est une extension séparable finie de $L$, donc on
peut appliquer (6.12.4.1) à $B$, et cela prouve alors la proposition.

II) Le corps des fractions $K$ de $A$ est de caractéristique $p>0$. ---
Alors $A$ contient le corps premier $\mathbf F_p$, et le théorème a été
prouvé dans ce cas (0, 22.7.6).
\end{proof}

\begin{corollary}[6.12.8]
\label{IV.2.s6.12.8-fr}
Soit $A$ un anneau local noethérien complet. Alors, pour tout préschéma $X$
localement de type fini sur $A$, $\operatorname{Reg}(X)$ est ouvert dans
$X$.
\end{corollary}

\begin{proof}
Vérifions la condition c) de (6.12.4). Si $\mathfrak p$ est premier dans
$A$, $A/\mathfrak p$ est encore un anneau local noethérien complet; si $K'$
est une extension finie du corps des fractions $K$ de $A/\mathfrak p$, $K'$
est corps des fractions d'une sous-$A$-algèbre finie $A'$ de $K'$, engendrée
par un système de générateurs de $K'$ sur $K$, entiers sur $A$. On sait
alors que $A'$ est un anneau semi-local complet (Bourbaki, Alg. comm.,
chap. IV, § 2, n$^o$ 5, cor. 3 de la prop. 9), donc produit d'anneaux locaux
complets, et comme $A'$ est intègre, c'est un anneau local complet; en vertu
de (6.12.7), si $X'=\operatorname{Spec}(A')$,
$\operatorname{Reg}(X')$ est ouvert et non vide, d'où la conclusion.
\end{proof}

\begin{proposition}[6.12.9]
\label{IV.2.s6.12.9-fr}
Soit $X$ un préschéma localement noethérien tel que $\operatorname{Reg}(X)$
soit ouvert; alors, pour tout $n\geq0$, l'ensemble $U_{R_n}(X)$ des points
$x\in X$ où $X$ vérifie $(R_n)$ est ouvert.
\end{proposition}

\begin{proof}
En effet (5.8.2), $U_{R_n}(X)$ est réunion de $\operatorname{Reg}(X)$ et de
l'ensemble des $x\in\operatorname{Sing}(X)$ tels que les points génériques
$z_i$ des composantes irréductibles de $\operatorname{Sing}(X)$ contenant
$x$ vérifient la relation $\dim(\mathcal O_{X,z_i})\geq n+1$ (puisque l'on a
$z_i\notin\operatorname{Reg}(X)$); autrement dit, ces points
$x\in\operatorname{Sing}(X)$ sont ceux pour lesquels on a
$\operatorname{codim}_x(\operatorname{Sing}(X),X)\geq n+1$ (5.1.2); comme
la fonction $x\mapsto\operatorname{codim}_x(\operatorname{Sing}(X),X)$ est
semi-continue intérieurement (0, 14.2.6), l'ensemble de ces points est
ouvert dans $\operatorname{Sing}(X)$, ce qui achève la démonstration.
\end{proof}

Notons aussi le résultat élémentaire suivant:

\begin{proposition}[6.12.10]
\label{IV.2.s6.12.10-fr}
Soit $X$ un préschéma localement noethérien. L'ensemble $U_{R_0}(X)$ est
ouvert dans $X$. Pour que $U_{R_1}(X)$ soit ouvert dans $X$, il faut et il
suffit que tout point maximal $x\in X$ tel que $x\in U_{R_1}(X)$ (ce qui
signifie que $X$ est réduit en $x$) soit intérieur à $U_{R_1}(X)$.
\end{proposition}

\begin{proof}
Par définition (5.8.2), $U_{R_0}(X)$ est l'ensemble
$X-\bigcup_\alpha X'_\alpha$, où $X'_\alpha$ parcourt l'ensemble des
composantes irréductibles de $X$ telles que $X$ ne soit pas réduit en leur
point générique; comme l'ensemble des composantes irréductibles de $X$ est
localement fini, $U_{R_0}(X)$ est ouvert. En ce qui concerne $U_{R_1}(X)$,
la condition de l'énoncé étant trivialement nécessaire, prouvons qu'elle est
suffisante. Soient $X''_\beta$ les composantes irréductibles de $X$ telles
qu'en leur point générique $x''_\beta$ le préschéma $X$ soit réduit; par
hypothèse il existe un ouvert $U\subset U_{R_1}(X)$ contenant tous les
$x''_\beta$. Désignons alors par $Z_\lambda$ les composantes irréductibles
de l'ensemble fermé $Z=X-U$ dont le point générique $z_\lambda$ est tel que
$\dim(\mathcal O_{X,z_\lambda})\leq1$ et que
$\mathcal O_{X,z_\lambda}$ soit non régulier. Aucun point appartenant à
l'un des $Z_\lambda$ ne peut appartenir
\oldpage[IV-2]{168}
à $U_{R_1}(X)$; mais inversement, si $x\in Z$ n'appartient à aucun des
$Z_\lambda$, alors, pour toute générisation $x'$ de $x$, ou bien $x'$
appartient à $U$, ou bien on a $\dim(\mathcal O_{X,x'})\geq2$, ou bien on a
$\dim(\mathcal O_{X,x'})=1$, et comme $x'$ n'appartient à aucun des
$Z_\lambda$, $\overline{\{x'\}}$ est nécessairement une composante
irréductible de $Z$, de codimension 1 dans $X$, distincte des $Z_\lambda$,
donc $\mathcal O_{X,x'}$ est régulier par définition. On en conclut que
$U_{R_1}(X)=X-\bigcup_\lambda Z_\lambda$; comme l'ensemble des $Z_\lambda$
est localement fini dans $Z$, $\bigcup_\lambda Z_\lambda$ est fermé, ce qui
achève de prouver que $U_{R_1}(X)$ est ouvert dans $X$.
\end{proof}

\subsection{\texorpdfstring{Critères pour que $\operatorname{Nor}(X)$ soit ouvert}{Critères pour que Nor(X) soit ouvert}}
\label{IV.2.s6.13-fr}

\begin{env}[6.13.1]
\label{IV.2.s6.13.1-fr}
Étant donné un préschéma localement noethérien $X$, nous noterons
$\operatorname{Nor}(X)$ l'ensemble des $x\in X$ où $X$ est normal; cet
ensemble contient $\operatorname{Reg}(X)$ et est contenu dans l'ensemble
(ouvert) des points $x$ tels que $\mathcal O_{X,x}$ soit intègre (i.e.
l'ensemble des points $x$ où $X$ est réduit, et qui n'appartiennent qu'à
une seule composante irréductible de $X$).
\end{env}

\begin{proposition}[6.13.2]
\label{IV.2.s6.13.2-fr}
Soient $X$ un préschéma localement noethérien réduit, $X'$ son normalisé
(II, 6.3.8). Si le morphisme canonique $f:X'\to X$ est fini,
$\operatorname{Nor}(X)$ est ouvert dans $X$.
\end{proposition}

\begin{proof}
En effet, on a $X'=\operatorname{Spec}(f_*(\mathcal O_{X'}))$ et
$f_*(\mathcal O_{X'})$ est une $\mathcal O_X$-Algèbre cohérente par
hypothèse (II, 6.1.3). Dire que $X$ est normal en un point $x$ signifie que
l'homomorphisme canonique
$\mathcal O_X\to(f_*(\mathcal O_{X'}))_x$ est bijectif (II, 6.3.4); mais
l'ensemble de ces points est ouvert puisque $\mathcal O_X$ et
$f_*(\mathcal O_{X'})$ sont cohérents (0$_{\rm I}$, 5.2.7).
\end{proof}

\begin{corollary}[6.13.3]
\label{IV.2.s6.13.3-fr}
Si $A$ est un anneau intègre noethérien japonais,
$\operatorname{Nor}(\operatorname{Spec}(A))$ est ouvert dans
$\operatorname{Spec}(A)$.
\end{corollary}

\begin{proposition}[6.13.4]
\label{IV.2.s6.13.4-fr}
Soit $X$ un préschéma localement noethérien. Pour que
$\operatorname{Nor}(X)$ soit ouvert, il faut et il suffit que tout point
maximal $x$ de $X$ tel que $x\in\operatorname{Nor}(X)$ (ce qui signifie que
$X$ est réduit au point $x$) soit intérieur à $\operatorname{Nor}(X)$.
\end{proposition}

\begin{proof}
Il n'y a à prouver que la suffisance de la condition. En vertu de la
remarque faite dans (6.13.1), on peut se borner (en remplaçant $X$ par un
ouvert de $X$ dans lequel $X$ est réduit) au cas où $X$ est réduit, i.e.
supposer que les points maximaux de $X$ appartiennent à $U_{S_1}(X)$ et à
$U_{R_1}(X)$; comme
$\operatorname{Nor}(X)=U_{S_1}(X)\cap U_{R_1}(X)$ en vertu du critère de
Serre (5.8.6), on déduit donc de (6.11.7) et (6.12.10) que ces deux ensembles
sont ouverts; par suite $\operatorname{Nor}(X)$ est ouvert dans $X$.
\end{proof}

\begin{corollary}[6.13.5]
\label{IV.2.s6.13.5-fr}
Si $X$ est tel que $\operatorname{Reg}(X)$ soit ouvert dans $X$, alors
$\operatorname{Nor}(X)$ est ouvert dans $X$.
\end{corollary}

\begin{proof}
En effet, tout point maximal $x$ où $X$ est réduit (donc
$\mathcal O_{X,x}$ un corps) appartient à $\operatorname{Reg}(X)$, donc est
intérieur à $\operatorname{Reg}(X)$ par hypothèse, et a fortiori est
intérieur à $\operatorname{Nor}(X)$.
\end{proof}

\begin{proposition}[6.13.6]
\label{IV.2.s6.13.6-fr}
\textnormal{(Nagata).} --- Soient $A$ un anneau intègre noethérien, $K$ son
corps des fractions, $K'$ une extension finie de $K$, $A'$ la fermeture
intégrale de $A$ dans $K'$. Pour que $A'$ soit une $A$-algèbre finie, il
faut et il suffit que les deux conditions suivantes soient satisfaites:
\begin{enumerate}
\item[(i)] Il existe $f\neq0$ dans $A$ tel que la fermeture intégrale $A'_f$
de l'anneau de fractions $A_f$ dans $K'$ soit une $A_f$-algèbre finie;
\item[(ii)] Pour tout idéal premier $\mathfrak p\in\operatorname{Spec}(A)$,
la fermeture intégrale $A'_\mathfrak p$ de l'anneau local $A_\mathfrak p$
dans $K'$ est une $A_\mathfrak p$-algèbre finie.
\end{enumerate}
\end{proposition}

\begin{proof}
Les conditions sont nécessaires, car pour toute partie multiplicative $S$
de $A$, $S^{-1}A'$ est la fermeture intégrale de $S^{-1}A$ dans $K'$ et est
par hypothèse un $S^{-1}A$-module de type fini. Pour voir que les conditions
sont suffisantes, considérons la famille filtrante croissante $(B_\lambda)$
des sous-$A$-algèbres de $K'$ qui sont des $A$-algèbres finies et ont $K'$
pour corps des fractions. Posons $Y=\operatorname{Spec}(A)$,
$X_\lambda=\operatorname{Spec}(B_\lambda)$, désignons par $u_\lambda$ le
morphisme $X_\lambda\to Y$ et soient
$S_\lambda=X_\lambda-\operatorname{Nor}(X_\lambda)$,
$T_\lambda=u_\lambda(S_\lambda)$. On peut se borner aux $B_\lambda$ qui
contiennent un ensemble fini dont l'image dans $A'_f$ est un système de
générateurs du $A_f$-module $A'_f$; cela entraîne, en vertu de l'hypothèse
(i), que $(B_\lambda)_f=A'_f$ pour tout $\lambda$, ou encore que
$u_\lambda^{-1}(D(f))$ est contenu dans $\operatorname{Nor}(X_\lambda)$. En
vertu de (6.13.2), $S_\lambda$ est donc fermé dans $X_\lambda$, et comme
$u_\lambda$ est un morphisme fini, $T_\lambda$ est fermé dans $Y$. Mais
pour tout $\mathfrak p\in\operatorname{Spec}(A)$, il existe en vertu de
(ii) un $\lambda$ tel que $(B_\lambda)_\mathfrak p=A'_\mathfrak p$, et par
suite tous les points
\oldpage[IV-2]{169}
de $X_\lambda$ au-dessus de $\mathfrak p$ appartiennent à
$\operatorname{Nor}(X_\lambda)$; en d'autres termes on a
$\bigcap_\lambda T_\lambda=\varnothing$; comme $Y$ est noethérien et les
$T_\lambda$ fermés, il existe un $\lambda$ tel que
$T_\lambda=\varnothing$; donc $B_\lambda$ est intégralement clos, et comme
son corps des fractions est $K'$, on a $B_\lambda=A'$. C.Q.F.D.
\end{proof}

\begin{proposition}[6.13.7]
\label{IV.2.s6.13.7-fr}
Soient $A$ un anneau noethérien, $X=\operatorname{Spec}(A)$. Les conditions
suivantes sont équivalentes:
\begin{enumerate}
\item[a)] Pour tout préschéma $X'$ localement de type fini sur $X$,
$\operatorname{Nor}(X')$ est ouvert dans $X'$;
\item[b)] Pour toute $A$-algèbre finie et intègre $A'$,
$\operatorname{Nor}(\operatorname{Spec}(A'))$ est ouvert dans
$\operatorname{Spec}(A')$;
\item[c)] Pour tout idéal premier $\mathfrak p$ de $A$ et toute extension
finie radicielle $K'$ du corps des fractions $K$ de $A/\mathfrak p$, il
existe une sous-$A$-algèbre $A'$ de $K'$, finie sur $A$, ayant $K'$ pour
corps des fractions et telle que $\operatorname{Nor}(\operatorname{Spec}(A'))$
soit ouvert dans $\operatorname{Spec}(A')$.
\end{enumerate}
\end{proposition}

\begin{proof}
Le fait que a) implique b) et que b) implique c) se prouve comme dans
(6.12.4). Pour montrer que c) entraîne a), compte tenu de (6.13.2), on se
ramène comme dans (6.12.4) à prouver que si $A'$ est une $A$-algèbre intègre
et de type fini, le point générique de $\operatorname{Spec}(A')$ est
intérieur à $\operatorname{Nor}(\operatorname{Spec}(A'))$. On distingue
alors deux cas comme dans (6.12.4) en prouvant d'abord le

\begin{lemma}[6.13.7.1]
\label{IV.2.s6.13.7.1-fr}
Soient $A$ un anneau intègre noethérien, $A'$ une $A$-algèbre intègre de
type fini, contenant $A$ et telle que le corps des fractions $K'$ de $A'$
soit une extension séparable du corps des fractions $K$ de $A$. Si
$\operatorname{Nor}(\operatorname{Spec}(A))$ est ouvert dans
$\operatorname{Spec}(A)$, $\operatorname{Nor}(\operatorname{Spec}(A'))$ est
ouvert dans $\operatorname{Spec}(A')$.
\end{lemma}

\begin{proof}
Il s'agit de prouver (compte tenu de (6.13.2)) que le point générique de
$\operatorname{Spec}(A')$ est intérieur à
$\operatorname{Nor}(\operatorname{Spec}(A'))$. La démonstration suit la
même marche que celle de (6.12.4.1), dont nous conservons les notations. On
remarque d'abord que l'on peut supposer que $A$ est intégralement clos, et
alors on sait que $A_1=A[t_1,\ldots,t_n]$ est intégralement clos (Bourbaki,
Alg. comm., chap. V, § 1, n$^o$ 3, cor. 2 de la prop. 13); on se ramène
ensuite au cas où $A'$ est un $A$-module libre de type fini; le raisonnement
de (6.12.4.1) prouve alors (en remplaçant au besoin $A$ par un anneau $A_f$
avec $f\neq0$) que les fibres $g^{-1}(z)$ du morphisme $g:X'\to X$ sont
régulières et a fortiori normales. En outre $g$ est plat et $X$ est normal,
donc (6.5.4, (ii)) $X'$ est normal.
\end{proof}

Ce lemme étant démontré, on passe au cas général comme dans (6.12.4, II)),
dont nous conservons encore les notations; appliquant l'hypothèse c), on
voit cette fois que $\operatorname{Nor}(X'_1)$ est ouvert et on se ramène
ainsi au cas où $X'_1$ est normal et $g:X'_1\to X'$ plat et surjectif; on
conclut cette fois que $X'$ est normal à l'aide de (6.5.4, (i)).
\end{proof}

\subsection{Changement de base et clôture intégrale}
\label{IV.2.s6.14-fr}

\begin{proposition}[6.14.1]
\EGAReaderTarget{ega4.c04.prop.6-14-1}{IV.2.s6.14.1-fr}
Soient $Y$, $Y'$ deux préschémas localement noethériens,
$g:Y'\to Y$ un morphisme normal (6.8.1). Alors, pour tout $Y$-préschéma
normal $X$, le préschéma $X'=X\times_Y Y'$ est normal.
\end{proposition}

Notons que dans cet énoncé on ne suppose pas $X$ localement noethérien.

\begin{lemma}[6.14.1.1]
\label{IV.2.s6.14.1.1-fr}
Soit $R$ un anneau composé direct d'un nombre fini de corps.
\begin{enumerate}
\item[(i)] Pour qu'un sous-anneau $A$ de $R$ ayant $R$ pour anneau total
des fractions soit normal, il faut et il suffit qu'il soit intégralement
fermé dans $R$.
\item[(ii)] Soit $(A_\lambda)$ une famille de sous-anneaux normaux de $R$;
si $A=\bigcap_\lambda A_\lambda$ admet $R$ pour anneau total des fractions,
$A$ est normal.
\end{enumerate}
\end{lemma}

\begin{proof}
(i) Comme $A$ est un sous-anneau de $R$, $A_\mathfrak p$ est un sous-anneau
de $R_\mathfrak p$ pour tout idéal premier $\mathfrak p$ de $A$, et
$R_\mathfrak p$ est un anneau de fractions de $A_\mathfrak p$; en outre
$R_\mathfrak p$ est nécessairement composé direct d'un nombre fini de
corps, donc tout élément non diviseur de $0$ dans $R_\mathfrak p$ est
inversible; cela prouve que $R_\mathfrak p$ est l'anneau total des fractions
de $A_\mathfrak p$. Si $A$ est intégralement fermé dans $R$,
$A_\mathfrak p$ est donc intégralement fermé dans $R_\mathfrak p$; mais si
$R_\mathfrak p$ est composé direct de 2 corps au moins, la fermeture
intégrale de $A_\mathfrak p$ dans $R_\mathfrak p$ est composée directe de
2 anneaux au moins non réduits à $0$, ce qui est absurde puisque
$A_\mathfrak p$ est un anneau local; donc $R_\mathfrak p$ est
nécessairement un corps et $A_\mathfrak p$ est intègre et intégralement clos,
ce qui par définition signifie que $A$ est normal. Inversement, si $A$ est
normal, $R_\mathfrak p$, anneau total des fractions d'un anneau intègre
$A_\mathfrak p$, est un corps et $A_\mathfrak p$ est intégralement fermé
dans $R_\mathfrak p$; si $x\in R$ est un élément entier sur $A$, son image
dans chaque $R_\mathfrak p$ est entière sur $A_\mathfrak p$, donc appartient
à $A_\mathfrak p$; on en conclut que $x\in A$ (Bourbaki, Alg. comm., chap.
II, § 3, n$^o$ 3, cor. 1 du th. 1), et par suite $A$ est intégralement fermé
dans $R$.

\oldpage[IV-2]{170}
(ii) Comme $A\subset A_\lambda\subset R$ pour tout $\lambda$, $R$ est aussi
l'anneau total des fractions de chaque $A_\lambda$. Vu la caractérisation
des anneaux normaux ayant $R$ pour anneau total des fractions, donnée dans
(i), l'assertion (ii) résulte de Bourbaki, Alg. comm., chap. V, § 1, n$^o$
3, prop. 12.
\end{proof}

Ce lemme étant démontré, la démonstration de (6.14.1) procède en plusieurs
étapes.
I) \emph{Réduction au cas où $Y=\operatorname{Spec}(A)$,
$Y'=\operatorname{Spec}(A')$, $X=\operatorname{Spec}(B)$, $A$, $A'$ étant
des anneaux locaux noethériens, $A$ intègre, $B$ la clôture intégrale de
$A$.} --- Il s'agit de prouver que pour tout $x'\in X'$, l'anneau local
$\mathcal O_{X',x'}$ est intègre et intégralement clos; soient $x$, $y$,
$y'$ les images canoniques de $x'$ dans $X$, $Y$, $Y'$ respectivement; si
l'on pose $A=\mathcal O_{Y,y}$, $A'=\mathcal O_{Y',y'}$,
$B=\mathcal O_{X,x}$, les anneaux $\mathcal O_{X',x'}$ pour $x$, $y$, $y'$
fixés, sont les anneaux locaux aux idéaux premiers de l'anneau
$B'=B\otimes_A A'$ (I, 3.6.5), et il s'agit donc de prouver que $B'$ est un
anneau normal. Notons d'autre part que
$\operatorname{Spec}(A')\to\operatorname{Spec}(A)$ est un morphisme normal
par définition (6.8.1 et I, 3.6.5); on est donc ramené au cas où $Y$ et $Y'$
sont des schémas locaux, $X$ le spectre d'un anneau local intègre et
intégralement clos $B$. Désignons par $(B_\alpha)$ la famille des clôtures
intégrales des sous-$A$-algèbres de type fini de $B$; il est clair que $B$
est réunion de la famille filtrante croissante des $B_\alpha$. Comme le
foncteur $\varinjlim$ commute au produit tensoriel, $B'$ est donc isomorphe
à $\varinjlim B'_\alpha$, où l'on a posé
$B'_\alpha=B_\alpha\otimes_A A'$. Pour prouver que $B'$ est normal, il
suffira, en vertu de (5.13.6), de montrer que les anneaux $B'_\alpha$ sont
normaux et que, pour $\alpha\leqslant\beta$, toute composante irréductible
de $\operatorname{Spec}(B'_\beta)$ domine une composante irréductible de
$\operatorname{Spec}(B'_\alpha)$. Mais cette dernière propriété résulte de
l'hypothèse que $A'$ est un $A$-module plat et de (2.3.7, (ii)), puisque
$B_\alpha$ et $B_\beta$ sont intègres et que $B_\alpha\subset B_\beta$.

On peut donc se borner à démontrer que $B'=B\otimes_A A'$ est normal lorsque
$B$ est la clôture intégrale d'une $A$-algèbre intègre de type fini $C$;
dans ce cas, si l'on pose $C'=C\otimes_A A'$, le morphisme
$\operatorname{Spec}(C')\to\operatorname{Spec}(C)$ est normal (6.8.2);
comme $B'=B\otimes_C C'$, on voit qu'on peut remplacer $A$ et $A'$ par $C$
et $C'$ respectivement, donc supposer que $A$ est intègre et que $B$ est la
clôture intégrale de $A$. Enfin, le procédé du début permet de se borner au
cas où $A$ est local (tenant compte de ce que, si $B$ est la clôture
intégrale de $A$, $B_\mathfrak p$ est la clôture intégrale de
$A_\mathfrak p$ pour tout idéal premier $\mathfrak p$ de $A$).

\oldpage[IV-2]{171}
II) \emph{Réduction au cas où $A$ est un anneau local intègre de dimension
1, $B$ un anneau de valuation discrète, clôture intégrale de $A$ et le
morphisme $\operatorname{Spec}(B)\to\operatorname{Spec}(A)$ radiciel.} ---
Soit $K$ le corps des fractions de $A$. On sait (0, 23.2.7) que $B$ est
intersection d'une famille $(V_\lambda)$ d'anneaux de valuation discrète
telle que, pour tout $x\in K$, on ait $x\in V_\lambda$ sauf pour un nombre
fini d'indices $\lambda$. On a donc une suite exacte de $A$-modules
\[
0\longrightarrow B\longrightarrow K\longrightarrow
\bigoplus_\lambda(K/V_\lambda).
\]
Posons $K'=K\otimes_A A'$, $V'_\lambda=V_\lambda\otimes_A A'$; par
platitude, on déduit de la suite exacte précédente une nouvelle suite
exacte
\[
0\longrightarrow B'\longrightarrow K'\longrightarrow
\bigoplus_\lambda(K'/V'_\lambda)
\]
et par suite $B'=\bigcap_\lambda V'_\lambda$. En outre
$\operatorname{Spec}(K')$ est la fibre du morphisme
$\operatorname{Spec}(A')\to\operatorname{Spec}(A)$ au point générique de
$\operatorname{Spec}(A)$, donc $K'$ est un anneau noethérien normal, et par
suite composé direct d'un nombre fini d'anneaux intègres et intégralement
clos; le composé direct $L'$ des corps des fractions de ces derniers est
l'anneau total des fractions de $A'$, donc aussi celui de $B'$. On peut donc
appliquer le lemme (6.14.1.1), et si l'on prouve que chacun des
$V'_\lambda$ est un anneau normal, il en sera de même de $B'$.

On sait d'autre part (0, 23.2.7) qu'on peut prendre les $V_\lambda$ tels
qu'il y ait une sous-$A$-algèbre finie $C$ de $K$ telle que $V_\lambda$
soit la clôture intégrale de $C_{\mathfrak p_\lambda}$, où
$\mathfrak p_\lambda$ est un idéal premier de hauteur 1 dans $C$. Si l'on
pose $C'_{\mathfrak p_\lambda}=C_{\mathfrak p_\lambda}\otimes_A A'$, le
morphisme $\operatorname{Spec}(C'_{\mathfrak p_\lambda})\to
\operatorname{Spec}(C_{\mathfrak p_\lambda})$ est normal (6.8.2) et l'on a
$V'_\lambda=V_\lambda\otimes_{C_{\mathfrak p_\lambda}}
C'_{\mathfrak p_\lambda}$; on peut donc remplacer $B$ par $V_\lambda$ et
$A$ par $C_{\mathfrak p_\lambda}$, donc supposer que $A$ est local, intègre
et de dimension 1, $B$ sa clôture intégrale et un anneau de valuation
discrète. Il y a une sous-$A$-algèbre finie $A_1$ de $B$ telle que le
morphisme $\operatorname{Spec}(B)\to\operatorname{Spec}(A_1)$ soit radiciel
(0, 23.2.5), ce qui entraîne en particulier que $A_1$ est aussi un anneau
local (évidemment de dimension 1); en outre on peut supposer que $B$ et
$A_1$ ont même corps résiduel (loc. cit.). On peut donc par la même méthode
remplacer $A$ par $A_1$. Appliquant si nécessaire le procédé du début de
I), on peut enfin supposer que $A'$ est lui aussi un anneau local et que
l'homomorphisme $A\to A'$ est local.

III) \emph{Fin de la démonstration.} --- Nous établirons d'abord le lemme
suivant:

\begin{lemma}[6.14.1.2]
\label{IV.2.s6.14.1.2-fr}
Soient $A$ un anneau local noethérien intègre de dimension 1, $A'$ un
anneau local noethérien, $A\to A'$ un homomorphisme local tel que le
morphisme correspondant $\operatorname{Spec}(A')\to\operatorname{Spec}(A)$
soit normal. Soient $K$ le corps des fractions de $A$, $\mathfrak m$
l'idéal maximal de $A$, $k=A/\mathfrak m$ son corps résiduel.
\begin{enumerate}
\item[(i)] Si $\mathfrak q'_j$ $(1\leqslant j\leqslant r)$ sont les idéaux
minimaux de $A'$, $K'=K\otimes_A A'$ est composé direct d'anneaux intègres
et intégralement clos $K'_j\supset A'/\mathfrak q'_j$
$(1\leqslant j\leqslant r)$, $K'_j$ ayant pour corps des fractions le corps
des fractions $L'_j$ de $A'/\mathfrak q'_j$, de sorte que l'anneau total des
fractions $L'$ de $A'$ s'identifie au composé direct des $L'_j$ et $A'$ à
un sous-anneau de $K'$.
\item[(ii)] L'idéal $\mathfrak p'=\mathfrak mA'$ de $A'$ est premier;
l'anneau $A'_{\mathfrak p'}$ est de dimension 1 et s'identifie à un
sous-anneau du produit des corps $L'_j$ pour les indices $j$ tels que
$\mathfrak q'_j\subset\mathfrak p'$.
\item[(iii)] Si $\overline{A'_{\mathfrak p'}}$ désigne le sous-anneau de
$L'$, produit de $A'_{\mathfrak p'}$ et des $L'_j$ tels que
$\mathfrak q'_j\not\subset\mathfrak p'$, on a
\begin{equation}
\tag{6.14.1.3}
\label{IV.2.s6.14.1.3-fr}
A'=K'\cap\overline{A'_{\mathfrak p'}}.
\end{equation}
\end{enumerate}
\end{lemma}

\begin{proof}
\oldpage[IV-2]{172}
L'assertion (i) a déjà été vue au cours de la démonstration de II) et est
indépendante de l'hypothèse dimensionnelle sur $A$. L'hypothèse que le
morphisme $\operatorname{Spec}(A')\to\operatorname{Spec}(A)$ est normal
entraîne que $A'\otimes_A k=A'/\mathfrak mA'$ est un anneau normal local et
noethérien, donc intègre, ce qui montre déjà que
$\mathfrak p'=\mathfrak mA'$ est premier. On a, d'après (6.1.2),
$\dim(A'_{\mathfrak p'})=\dim(A)+
\dim(A'_{\mathfrak p'}/\mathfrak mA'_{\mathfrak p'})$. Mais
$A'_{\mathfrak p'}/\mathfrak mA'_{\mathfrak p'}=
A'_{\mathfrak p'}/\mathfrak p'A'_{\mathfrak p'}$ est le corps résiduel de
$A'_{\mathfrak p'}$, donc $\dim(A'_{\mathfrak p'})=1$. Le fait que
$A'_{\mathfrak p'}$ est contenu dans le composé direct des $L'_j$ tels que
$\mathfrak q'_j\subset\mathfrak p'$ résulte de ce que les idéaux
$\mathfrak q'_jA'_{\mathfrak p'}$ pour ces indices $j$ sont les idéaux
premiers minimaux de $A'_{\mathfrak p'}$.

Reste à prouver (6.14.1.3). On a évidemment
$A'\subset K'\cap\overline{A'_{\mathfrak p'}}$. Inversement, soit $y'$ un
élément de cette intersection; soit $a$ un «paramètre» pour $A$, donc tel
que $Aa$ soit $\mathfrak m$-primaire, et soit $a'$ l'image de $a$ dans
$A'$; tout élément de $K$ peut s'écrire $x/a^n$ pour $x\in A$ et un $n>0$
puisque $Aa$ contient une puissance de $\mathfrak m$; donc on peut écrire
$y'=x'/a'^n$ avec $x'\in A'$. Notons maintenant que $\mathfrak p'$ est le
seul idéal premier associé à
$A'/a'^nA'=(A/a^nA)\otimes_A A'$: comme $A'$ est un $A$-module plat, cela
résulte de (3.3.1), $\mathfrak m$ étant le seul idéal premier associé à
$A/a^nA$ et $k\otimes_A A'$ étant intègre. Par suite, $a'^nA'$ est l'image
réciproque dans $A'$ de $a'^nA'_{\mathfrak p'}$; comme par hypothèse
$y'\in\overline{A'_{\mathfrak p'}}$, l'image de $x'$ dans
$A'_{\mathfrak p'}$ appartient à $a'^nA'_{\mathfrak p'}$; d'où
$x'\in a'^nA'$ et $y'\in A'$, ce qui achève la démonstration.

Ce lemme étant établi, dans le cas où nous sommes ramenés à la fin de II),
$B'$ est radiciel sur $A'$ (I, 3.5.7) et par suite est aussi un anneau
local; en outre $B'$ est entier sur $A'$, donc, si $\mathfrak q'$ est
l'unique idéal premier de $B'$ au-dessus de $\mathfrak p'$,
$\mathfrak q'B'_{\mathfrak p'}$ est le seul idéal maximal de
$B'_{\mathfrak p'}$ (Bourbaki, Alg. comm., chap. V, § 2, n$^o$ 1, prop.
1), donc $B'_{\mathfrak q'}=B'_{\mathfrak p'}=B\otimes_A
A'_{\mathfrak p'}$. Nous allons d'abord démontrer que
$B'_{\mathfrak p'}$ est un anneau noethérien et normal. Or, comme $B$
contient $A$ et est contenu dans $K$ et que $A'_{\mathfrak p'}$ est un
$A$-module plat, $B'_{\mathfrak p'}$ contient $A'_{\mathfrak p'}$ et est
contenu dans $K'_{\mathfrak p'}$, donc dans le produit $L_1$ des $L'_j$
tels que $\mathfrak q'_j\subset\mathfrak p'$. Pour tout indice $j$ tel que
$\mathfrak q'_j\subset\mathfrak p'$, soit $\mathfrak a'_j$ le produit des
$L'_h$ tels que $h\neq j$, de sorte que
$\mathfrak a'_j\cap A'=\mathfrak q'_j$; comme tout élément de
$A'-\mathfrak p'$ est régulier dans $L_1$, on a aussi
\[
\mathfrak a'_j\cap A'_{\mathfrak p'}=
\mathfrak q'_jA'_{\mathfrak p'}.
\]
Soit $\mathfrak r'_j=B'_{\mathfrak p'}\cap\mathfrak a'_j$, de sorte que
$B'_{\mathfrak p'}/\mathfrak r'_j$ est isomorphe à la projection de
$B'_{\mathfrak p'}$ dans $L'_j$; donc $B'_{\mathfrak p'}/\mathfrak r'_j$
contient l'anneau local intègre de dimension 1,
$A'_{\mathfrak p'}/\mathfrak q'_jA'_{\mathfrak p'}$, et est contenu dans
son corps des fractions $L'_j$; il est par suite noethérien en vertu du
théorème de Krull-Akizuki (Bourbaki, Alg. comm., chap. VII, § 2, n$^o$ 5,
prop. 5).

\oldpage[IV-2]{173}
Comme l'intersection des $\mathfrak r'_j$ est réduite à 0, on en déduit que
$B'_{\mathfrak p'}$ lui-même est noethérien, en raison du lemme classique
suivant:

\begin{lemma}[6.14.1.4]
\label{IV.2.s6.14.1.4-fr}
Soient $R$ un anneau, $\mathfrak a$ et $\mathfrak b$ deux idéaux de $R$; si
$R/\mathfrak a$ et $R/\mathfrak b$ sont noethériens, il en est de même de
$R/(\mathfrak a\cap\mathfrak b)$.
\end{lemma}

\begin{proof}
En effet, soit $\mathfrak c$ un idéal de $R$ tel que
$\mathfrak a\cap\mathfrak b\subset\mathfrak c$; on a donc
$\mathfrak c\cap\mathfrak b\subset\mathfrak a\cap\mathfrak c\subset
\mathfrak c$; or $\mathfrak c/(\mathfrak a\cap\mathfrak c)$ est un
$R$-module isomorphe à $(\mathfrak a+\mathfrak c)/\mathfrak a$, donc à un
idéal de $R/\mathfrak a$, et par suite est de type fini; d'autre part
$(\mathfrak a\cap\mathfrak c)/(\mathfrak a\cap\mathfrak b)$ est un
sous-$R$-module de $\mathfrak a/(\mathfrak a\cap\mathfrak b)$, et ce
dernier est isomorphe à $(\mathfrak a+\mathfrak b)/\mathfrak b$, donc est
de type fini comme idéal de $R/\mathfrak b$; donc
$(\mathfrak a\cap\mathfrak c)/(\mathfrak a\cap\mathfrak b)$ est aussi un
$R$-module de type fini, et il en est finalement de même de
$\mathfrak c/(\mathfrak a\cap\mathfrak b)$.
\end{proof}

Notons d'autre part que $\operatorname{Spec}(A)$ se compose du point fermé
$\mathfrak m$ et du point générique $(0)$, de corps résiduels $k$ et $K$
respectivement; comme dans le cas où nous nous sommes placés, le corps des
fractions de $B$ et son corps résiduel sont respectivement isomorphes à $K$
et $k$, les fibres du morphisme
$u:\operatorname{Spec}(B'_{\mathfrak p'})\to\operatorname{Spec}(B)$ au
point fermé et au point générique de $\operatorname{Spec}(B)$ sont
respectivement isomorphes aux fibres du morphisme
$\operatorname{Spec}(A'_{\mathfrak p'})\to\operatorname{Spec}(A)$ au point
fermé et au point générique de $\operatorname{Spec}(A)$ (I, 3.6.4), donc
sont géométriquement normales par hypothèse; comme d'autre part le
morphisme $u$ est plat, on en conclut qu'il est normal (6.8.1). Mais comme
$B$ et $B'_{\mathfrak p'}$ sont noethériens et que $B$ est normal, on
déduit de (6.5.4) que $B'_{\mathfrak p'}$ est normal.

Cela étant, $B$ est réunion de la famille filtrante croissante de ses
sous-$A$-algèbres finies $A_\alpha$; par platitude, $B'$ est réunion de la
famille filtrante croissante des
$A'_\alpha=A_\alpha\otimes_A A'$; si l'on pose
$\mathfrak p'_\alpha=\mathfrak q'\cap A'_\alpha$,
$B'_{\mathfrak q'}$ est aussi réunion de la famille filtrante croissante
des $(A'_\alpha)_{\mathfrak p'_\alpha}$ (5.13.3). Désignons par $L''$ le
composé direct des corps $L'_j$ tels que
$\mathfrak q'_j\not\subset\mathfrak p'$; alors pour tout $\alpha$,
$(A'_\alpha)_{\mathfrak p'_\alpha}$ est contenu dans
$K'_{\mathfrak p'}$, et l'anneau
$\overline{B'_{\mathfrak q'}}=L''\times B'_{\mathfrak q'}$ est donc réunion
des anneaux
$L''\times(A'_\alpha)_{\mathfrak p'_\alpha}=
(\overline{A'_\alpha})_{\mathfrak p'_\alpha}$. Mais chacun des $A_\alpha$
est local, noethérien, intègre et de dimension 1, et le morphisme
$\operatorname{Spec}(A'_\alpha)\to\operatorname{Spec}(A_\alpha)$ est
normal (6.8.2), donc on peut lui appliquer le lemme (6.14.1.2) et l'on a
\[
A'_\alpha=K'\cap(\overline{A'_\alpha})_{\mathfrak p'_\alpha}
\]
pour tout $\alpha$; prenant la limite inductive de chacun des deux membres,
il vient
\[
B'=K'\cap\overline{B'_{\mathfrak q'}}.
\]
Mais $\overline{B'_{\mathfrak q'}}$, composé direct des anneaux normaux
$L''$ et $B'_{\mathfrak q'}$, est normal, et comme il en est de même de
$K'$, le lemme (6.14.1.1) montre que $B'$ est normal. C.Q.F.D.
\end{proof}

\begin{corollary}[6.14.2]
\label{IV.2.s6.14.2-fr}
Soient $k$ un corps, $X$ un $k$-préschéma normal (non nécessairement
localement noethérien). Alors, pour toute extension séparable $k'$ de $k$,
$X_{(k')}=X\otimes_k k'$ est normal.
\end{corollary}

\begin{proof}
On sait en effet (6.7.6) que le morphisme
$\operatorname{Spec}(k')\to\operatorname{Spec}(k)$ est alors normal.
\end{proof}

\begin{corollary}[6.14.3]
\label{IV.2.s6.14.3-fr}
Soient $k$ un corps; $X$, $Y$ deux $k$-préschémas intègres et normaux, dont
les corps de fonctions rationnelles $K=R(X)$, $L=R(Y)$ sont des extensions
séparables de $k$. Alors le préschéma $X\times_k Y$ est normal.
\end{corollary}

\begin{proof}
\oldpage[IV-2]{174}
La question étant locale sur $X$ et $Y$, on peut supposer que
$X=\operatorname{Spec}(A)$, $Y=\operatorname{Spec}(B)$, où $A$ et $B$ sont
deux anneaux intègres et intégralement clos, de corps des fractions $K$ et
$L$ respectivement; supposons d'abord que $A$ soit une $k$-algèbre de type
fini, donc $K$ une extension de type fini de $k$. Par platitude,
$A\otimes_k B$ est une sous-$k$-algèbre de $K\otimes_k L$; or
$K\otimes_k L$ est un anneau noethérien normal (en vertu de (6.7.4.1) ou
(6.14.2)), donc composé direct d'anneaux intègres $C_i$ en nombre fini, de
sorte que si $E_i$ est le corps des fractions de $C_i$, le composé direct
$E$ des $E_i$ est l'anneau total des fractions de $K\otimes_k L$; c'est
d'ailleurs aussi l'anneau total des fractions de $A\otimes_k B$, puisque
$K\otimes_k L$ est un anneau de fractions de $A\otimes_k B$. Cela étant,
on peut écrire
$A\otimes_k B=(A\otimes_k L)\cap(K\otimes_k B)$, et il résulte de (6.14.2)
que $A\otimes_k L$ et $K\otimes_k B$ sont normaux, donc $A\otimes_k B$ est
normal en vertu de (6.14.1.1).

Considérons maintenant le cas général; $A$ est alors réunion de la famille
filtrante croissante de ses sous-$k$-algèbres de type fini $A_\alpha$,
donc $A\otimes_k B$ est limite inductive des anneaux normaux
$A_\alpha\otimes_k B$. Pour prouver que $A\otimes_k B$ est normal, il
suffit donc (5.13.6) de prouver que toute composante irréductible de
$\operatorname{Spec}(A_\beta\otimes_k B)$ domine une composante
irréductible de $\operatorname{Spec}(A_\alpha\otimes_k B)$ pour
$A_\alpha\subset A_\beta$, ce qui résulte aussitôt de ce que $B$ est un
$k$-module plat, de (2.3.7, (ii)) et du fait que $A_\alpha$ et $A_\beta$
sont des anneaux intègres.
\end{proof}

\begin{proposition}[6.14.4]
\EGAReaderTarget{ega4.c04.prop.6-14-4}{IV.2.s6.14.4-fr}
Soient $A$ un anneau noethérien, $A'$ une $A$-algèbre noethérienne telle
que le morphisme $\operatorname{Spec}(A')\to\operatorname{Spec}(A)$ soit
normal. Soient $B$ une $A$-algèbre, $C$ la fermeture intégrale de $A$ dans
$B$; on pose $B'=B\otimes_A A'$, $C'=C\otimes_A A'$, $C'$ s'identifiant à
un sous-anneau de $B'$; alors $C'$ est la fermeture intégrale de $A'$ dans
$B'$.
\end{proposition}

\begin{proof}
Nous procéderons en plusieurs étapes.

I) \emph{Réduction au cas où l'anneau $B$ est réduit.} --- Posons
$B_0=B_{\mathrm{red}}=B/\mathfrak N$, où $\mathfrak N$ est le nilradical
de $B$, et soit $C_0$ la fermeture intégrale de $A$ dans $B_0$. On a le
lemme suivant:

\begin{lemma}[6.14.4.1]
\label{IV.2.s6.14.4.1-fr}
Soient $A$ un anneau, $B$ une $A$-algèbre, $B_0=B/\mathfrak n$ le quotient
de $B$ par un nilidéal. Si $C_0$ est la fermeture intégrale de $A$ dans
$B_0$, l'image réciproque $C$ de $C_0$ par l'homomorphisme canonique
$\varphi:B\to B_0$ est la fermeture intégrale de $A$ dans $B$.
\end{lemma}

\begin{proof}
En effet, si $x\in B$ est tel que $\varphi(x)$ vérifie une équation de
dépendance intégrale à coefficients dans $A$, on en déduit que $x$ vérifie
une relation de la forme
$x^n+a_1x^{n-1}+\cdots+a_n\in\mathfrak n$ avec $a_k\in A$, d'où, en
élevant à une puissance assez grande, une équation de dépendance intégrale
pour $x$, à coefficients dans $A$.
\end{proof}

On a donc la suite exacte de $A$-modules
\[
0\longrightarrow C\longrightarrow B\longrightarrow B_0/C_0
\]
d'où, par platitude, la suite exacte
\[
0\longrightarrow C'\longrightarrow B'\longrightarrow B'_0/C'_0
\]
où l'on a posé $B'_0=B_0\otimes_A A'$, $C'_0=C_0\otimes_A A'$. Si l'on
prouve que $C'_0$ est la fermeture intégrale de $A'$ dans $B'_0$, le lemme
(6.14.4.1) prouvera que $C'$ est la fermeture intégrale de $A'$ dans $B'$.

\oldpage[IV-2]{175}
II) \emph{Réduction au cas où $B$ est une $A$-algèbre intègre de type fini
contenant $A$.} --- Soit $(B_\alpha)$ la famille filtrante croissante des
sous-$A$-algèbres de type fini de $B$, et soit $C_\alpha$ la fermeture
intégrale de $A$ dans $B_\alpha$. Il résulte aussitôt de la définition de
la fermeture intégrale que $C$ est réunion des $C_\alpha$; si l'on pose
$B'_\alpha=B_\alpha\otimes_A A'$, $C'_\alpha=C_\alpha\otimes_A A'$,
$C'_\alpha$ est contenu dans $B'_\alpha$ par platitude, et pour la même
raison $B'$ est réunion de la famille filtrante croissante des
$B'_\alpha$ et $C'$ réunion de la famille filtrante croissante des
$C'_\alpha$. Si l'on prouve que $C'_\alpha$ est la fermeture intégrale de
$A'$ dans $B'_\alpha$, il en résultera aussitôt que $C'$ est la fermeture
intégrale de $A'$ dans $B'$. On peut donc se borner au cas où $B$ est une
$A$-algèbre de type fini, donc noethérienne; soient $\mathfrak q_i$
$(1\leqslant i\leqslant n)$ ses idéaux premiers minimaux; comme $B$ est
supposé réduit, il s'identifie à un sous-anneau du produit $B_0$ des
$B/\mathfrak q_i$; si $C_0$ est la fermeture intégrale de $A$ dans $B_0$,
on a $C=B\cap C_0$; si l'on pose $B'_0=B_0\otimes_A A'$,
$C'_0=C_0\otimes_A A'$, on a, par platitude,
$C'=B'\cap C'_0$ (0$_{\mathrm I}$, 6.1.3); il suffit donc de prouver que
$C'_0$ est la fermeture intégrale de $A'$ dans $B'_0$. Mais $C_0$ est
composé direct des $C_i$, fermetures intégrales de $A$ dans les
$B_i=B/\mathfrak q_i$; par suite $C'_0$ est composé direct des
$C'_i=C_i\otimes_A A'$ et il suffit de montrer que $C'_i$ est la fermeture
intégrale de $A'$ dans $B'_i=B_i\otimes_A A'$. On est ainsi ramené au cas
où $B$ est intègre et une $A$-algèbre de type fini; si $\mathfrak p$ est le
noyau de l'homomorphisme $A\to B$, on a aussi
$B'=B\otimes_{A/\mathfrak p}(A'/\mathfrak pA')$; comme le morphisme
$\operatorname{Spec}(A'/\mathfrak pA')\to
\operatorname{Spec}(A/\mathfrak p)$ est normal, on peut remplacer $A$ par
$A/\mathfrak p$ et $A'$ par $A'/\mathfrak pA'$, et supposer par suite que
$A\subset B$.

III) \emph{Cas où $A$ est un corps, $B$ un corps, extension de type fini de
$A$ et tel que $A$ soit algébriquement fermé dans $B$.} --- On a alors
$C=A$, et $A'$ est une $A$-algèbre noethérienne géométriquement normale,
donc composée directe d'anneaux intègres et intégralement clos $A'_i$, les
corps des fractions $K'_i$ des $A'_i$ étant des extensions séparables de
$A$ (4.6.1). La $A'$-algèbre $B'$ est donc composée directe des
$B\otimes_A A'_i=B'_i$ et l'on peut donc se borner au cas où $A'$ est
intègre, son corps des fractions $K'$ étant une extension séparable de $A$.
Or, comme $B$ est un $A$-module plat, $B\otimes_A A'=B'$ s'identifie à un
sous-anneau de $B\otimes_A K'$; comme $K'$ est extension séparable de $A$,
$B\otimes_A K'$ est réduit (4.3.7), et comme de plus $B$ est une extension
primaire de $A$, $B\otimes_A K'$ est intègre (4.3.2); il en est par suite
de même de $B'$. Soit $L'$ le corps des fractions de $B\otimes_A K'$ (qui
est aussi celui de $B'$); c'est évidemment un corps composé $B(K')$ de ses
sous-corps $B$ et $K'$; comme $A$ est algébriquement fermé dans $B$, que
$K'$ est une extension séparable de $A$, et que $B$ et $K'$ sont
linéairement disjoints sur $A$, $K'$ est algébriquement fermé dans
$B(K')=L'$ (Bourbaki, Alg., chap. V, § 9, exerc. 2); a fortiori $A'$, étant
intégralement clos, est intégralement fermé dans $L'$, donc dans $B'$, ce
qui prouve la proposition dans ce cas.

IV) \emph{Cas où $A$ est intègre, $B$ un corps, extension finie du corps
des fractions $K$ de $A$.} --- Alors
$B'=(A'\otimes_A K)\otimes_K B$ est une $B$-algèbre noethérienne
géométriquement normale, et par suite $B'$ est composé direct d'anneaux
intègres; l'anneau total des fractions $L'$ de $B'$ est alors composé
direct d'un nombre fini de corps. Cela étant, comme $B$ est le corps des
fractions de $C$, $B'$ est un anneau de fractions de $C'$ et $L'$ est donc
aussi l'anneau total des fractions de $C'$. Or (6.14.1) comme $C$ est un
anneau normal et $\operatorname{Spec}(A')\to\operatorname{Spec}(A)$ un
morphisme normal d'anneaux noethériens, $C'$ est un anneau normal; mais il
résulte de (6.14.1.1) que $C'$ est alors intégralement fermé dans $L'$, et
a fortiori dans $B'$, et par suite est la fermeture intégrale de $A'$ dans
$B'$.

\oldpage[IV-2]{176}
V) \emph{Fin de la démonstration.} --- D'après II), on peut supposer que
$B$ est une $A$-algèbre intègre de type fini contenant $A$. Soient $K$ le
corps des fractions de $A$, $L$ celui de $B$, qui est une extension de type
fini de $K$. Soit $M$ la fermeture algébrique de $K$ dans $L$, qui est une
extension algébrique finie de $K$; soit $C_0$ la fermeture intégrale de $A$
dans $M$, qui est aussi la fermeture intégrale de $A$ dans $L$; on a donc
$C=B\cap C_0$; si l'on pose $C'_0=C_0\otimes_A A'$, on a par suite
$C'=B'\cap C'_0$ par platitude (0$_{\mathrm I}$, 6.1.3). Or, il résulte de
IV) que $C'_0$ est la fermeture intégrale de $A'$ dans
$M'=M\otimes_A A'$; en outre, $M'$ est un anneau noethérien et le morphisme
$\operatorname{Spec}(M')\to\operatorname{Spec}(M)$ est normal (comme on
l'a vu dans IV)); appliquant III) à $M$ et $M'$ au lieu de $A$ et $A'$ et
à $L$ au lieu de $B$, on en déduit que $M'$ est intégralement fermé dans
$L'=L\otimes_A A'$; donc $C'_0$ est la fermeture intégrale de $A'$ dans
$L'$, et $C'=C'_0\cap B'$ est la fermeture intégrale de $A'$ dans $B'$.
C.Q.F.D.
\end{proof}

\begin{corollary}[6.14.5]
\label{IV.2.s6.14.5-fr}
Soient $A$ un anneau noethérien, $A'$ une $A$-algèbre noethérienne telle
que le morphisme $\operatorname{Spec}(A')\to\operatorname{Spec}(A)$ soit
normal. Soient $B$, $C$ deux $A$-algèbres,
$\varphi:B\to C$ un $A$-homomorphisme faisant de $C$ une $B$-algèbre, $D$
la fermeture intégrale de $B$ dans $C$. Si l'on pose
$B'=B\otimes_A A'$, $C'=C\otimes_A A'$, $D'=D\otimes_A A'$, $D'$ est la
fermeture intégrale de $B'$ dans $C'$.
\end{corollary}

\begin{proof}
Soit $(B_\lambda)$ la famille filtrante croissante des sous-$A$-algèbres
de type fini de $B$, et posons $B'_\lambda=B_\lambda\otimes_A A'$, de
sorte que $B'$ est réunion de la famille filtrante croissante des
sous-$A'$-algèbres de type fini $B'_\lambda$. Soit $D_\lambda$ la
fermeture intégrale de $B_\lambda$ dans $C$, et posons
$D'_\lambda=D_\lambda\otimes_A A'$, de sorte que $D$ est réunion des
$D_\lambda$ et $D'$ réunion des $D'_\lambda$; si nous prouvons que
$D'_\lambda$ est la fermeture intégrale de $B'_\lambda$ dans $C'$, il en
résultera que $D'$ est la fermeture intégrale de $B'$ dans $C'$. Or
$B_\lambda$ et $B'_\lambda$ sont noethériens et le morphisme
$\operatorname{Spec}(B'_\lambda)\to\operatorname{Spec}(B_\lambda)$ est
normal (6.8.2); on peut donc appliquer (6.14.4) en remplaçant $A$, $A'$ et
$B$ par $B_\lambda$, $B'_\lambda$ et $C$ respectivement, d'où le
corollaire.
\end{proof}

On peut par exemple appliquer (6.14.5) lorsque $A$ est un anneau local
excellent et $A'$ son complété $\widehat A$, puisque dans ce cas
$\operatorname{Spec}(A')\to\operatorname{Spec}(A)$ est un morphisme
régulier (7.8.2).

\subsection{Préschémas géométriquement unibranches}
\label{IV.2.s6.15-fr}

\begin{env}[6.15.1]
\label{IV.2.s6.15.1-fr}
Nous dirons qu'un préschéma $X$ est \emph{unibranche} (resp.
\emph{géométriquement unibranche}) en un point $x$, ou que le point $x$
est unibranche (resp. géométriquement unibranche) dans $X$ si l'anneau
local $\mathcal O_{X,x}$ est unibranche (resp. géométriquement
unibranche) (0, 23.2.1); nous dirons que $X$ est unibranche (resp.
géométriquement unibranche) s'il l'est en tout point. Il suit de cette
définition que, pour que $X$ soit unibranche (resp. géométriquement
unibranche) en un point, il faut et il suffit que $X_{\mathrm{red}}$ le soit
en ce point.
\end{env}

\begin{env}[6.15.2]
\EGAReaderTarget{ega4.c04.statement.6-15-2}{IV.2.s6.15.2-fr}
Il faut prendre garde qu'avec les définitions de (6.15.1), un anneau local
$A$ peut être unibranche (resp. géométriquement unibranche) sans que
$\operatorname{Spec}(A)$ le soit; en d'autres termes, il peut se faire qu'il
y ait des idéaux premiers $\mathfrak p$ de $A$ tels que $A_\mathfrak p$ ne
soit pas unibranche; il revient aussi au même de dire que sur un préschéma
la notion de point géométriquement unibranche n'est pas stable par
générisation. Nous allons le voir sur l'exemple suivant.
\end{env}

\oldpage[IV-2]{177}
Soient $K$ un corps algébriquement clos de caractéristique 0, $B$ l'anneau
intègre
\[
K[U,V,W]/(U^2(U-W)-V^2(U+W))
\]
($U$, $V$, $W$ indéterminées) de sorte que $Y=\operatorname{Spec}(B)$ est
un « cône de sommet l'origine, ayant une génératrice double ». Nous
désignerons par $u$, $v$, $w$ les images de $U$, $V$, $W$ dans $B$. Soit
$R$ le corps des fractions de $B$, et considérons dans $R$ l'élément
$t=v(u+w)/u$, qui n'appartient pas à $B$; montrons que $C=B[t]$ est la
clôture intégrale de $B$. On a en effet $t^2=u^2-w^2$, donc $t$ est entier
sur $B$, et $v=tu/(u+w)$; l'anneau $C_1=K[t,u,w]$ est intégralement clos,
car il est isomorphe à
$K[T,U,W]/(T^2-U^2+W^2)$ et est donc la fermeture intégrale de l'anneau
intégralement clos $K[U,W]$ dans l'extension quadratique
$K(U,W)(\sqrt{U^2+W^2})$ de son corps des fractions (Bourbaki, Alg. comm.,
chap. V, § 1, n$^o$ 6, prop. 18). L'anneau de fractions
$K[t,u,w,1/(u+w)]$ de $C_1$ est donc aussi intégralement clos. De la même
manière, on voit que l'anneau $C_2=K[t,v,w]$ est intégralement clos, car
$t$ est racine de l'équation
$t(t-v)-w^2(2v-t)=0$, et par suite
$K[t,v,w,1/(t-v)]=K[t,tu/(u+w),w,(u+w)/(tw)]$ est intégralement clos. Finalement, compte tenu de ce
que $u$ et $w$ sont algébriquement indépendants sur $K$, on prouve aisément
que
\[
C=K[t,u,v,w]=K[t,u,w,1/(u+w)]\cap K[t,v,w,1/(t-v)],
\]
ce qui achève de montrer que $C$ est la clôture intégrale de $B$. Il est
immédiat que si $\mathfrak m_0$ est l'idéal maximal de $B$ engendré par
$u$, $v$, $w$ (« sommet du cône »), il existe un seul idéal maximal
$\mathfrak n_0$ de $C$ au-dessus de $\mathfrak m_0$, savoir l'idéal
engendré par $t$, $u$, $v$, $w$. Si l'on pose $A=B_{\mathfrak m_0}$, on en
déduit aisément que $A'=C_{\mathfrak n_0}$ est la clôture intégrale de
$A$, qui est donc unibranche, et par suite aussi géométriquement unibranche
puisque son corps résiduel $K$ est algébriquement clos. Mais dans $A$
l'idéal premier $\mathfrak p$ engendré par $u$ et $v$ est tel que la
clôture intégrale $A_\mathfrak p[t]$ de $A_\mathfrak p$ ne soit pas un
anneau local.

Nous verrons toutefois plus loin (9.7.10) que lorsque $X$ est un préschéma
localement noethérien tel que, si $X'$ est le normalisé de
$X_{\mathrm{red}}$, le morphisme canonique $X'\to X$ soit fini (ce qui sera
le cas si $X$ est tel que les anneaux de ses ouverts affines soient
universellement japonais (0, 23.1.1)), alors l'ensemble des points $x\in X$
où $X$ est géométriquement unibranche est localement constructible.

\begin{env}[6.15.3]
\label{IV.2.s6.15.3-fr}
Nous dirons pour abréger qu'un morphisme $f:X\to Y$ est \emph{radiciel en
un point} $y\in Y$, si $f^{-1}(y)$ est vide ou réduit à un seul point $x$
et si $k(x)$ est une extension radicielle de $k(y)$. Il revient donc au
même de dire que $f$ est radiciel en tous les points de $Y$ ou que $f$ est
radiciel (I, 3.5.8). Si
$g:f^{-1}(y)\to\operatorname{Spec}(k(y))$ est le morphisme image réciproque
de $f$ par le changement de base
$\operatorname{Spec}(k(y))\to Y$, il revient au même de dire que $f$ est
radiciel au point $y\in Y$, ou que $g$ est radiciel.
\end{env}

\begin{lemma}[6.15.3.1]
\label{IV.2.s6.15.3.1-fr}
\begin{enumerate}
\item[(i)] Soient $f:X\to Y$, $g:Y\to Z$ deux morphismes. Si $f$ est
radiciel en un point $y$ et $g$ radiciel au point $z=g(y)$, alors $g\circ f$
est radiciel au point $g(y)$. La réciproque est vraie si $f$ est surjectif.
\item[(ii)] Soient $f:X\to Y$, $h:Y'\to Y$ deux morphismes, et soit
$f'=f_{(Y')}:X\times_Y Y'\to Y'$. Soient $y'\in Y'$ et $y=h(y')$; alors
pour que $f$ soit radiciel au point $y$, il faut et il suffit que $f'$ soit
radiciel au point $y'$.
\end{enumerate}
\end{lemma}

\begin{proof}
\oldpage[IV-2]{178}
(i) Si $f$ est radiciel au point $y$ et $g$ radiciel au point $z$,
$g^{-1}(z)$ est réduit à $y$ et
$f^{-1}(g^{-1}(z))=f^{-1}(y)$ vide ou réduit à un seul point $x$; en
outre $k(x)$ est extension radicielle de $k(y)$ et $k(y)$ extension
radicielle de $k(z)$, donc $k(x)$ est extension radicielle de $k(z)$.
Inversement, supposons $f$ surjectif; si $g\circ f$ est radiciel au point
$z=g(y)$, $g^{-1}(z)$ se réduit au seul point $y$, sans quoi
$f^{-1}(g^{-1}(z))$ aurait au moins deux points distincts; en outre,
$f^{-1}(y)=f^{-1}(g^{-1}(z))$ se réduit à un seul point $x$, et par
hypothèse on a $k(z)\subset k(y)\subset k(x)$, et $k(x)$ est radiciel sur
$k(z)$, donc $k(y)$ est radiciel sur $k(z)$ et $k(x)$ radiciel sur $k(y)$.

(ii) Soient $g:f^{-1}(y)\to\operatorname{Spec}(k(y))$,
$g':f'^{-1}(y')\to\operatorname{Spec}(k(y'))$ les morphismes images
réciproques de $f$ et $f'$ respectivement; comme
$f'^{-1}(y')=f^{-1}(y)\otimes_{k(y)}k(y')$ (I, 3.6.4), $g'$ est image
réciproque de $g$ et l'assertion se réduit à (2.6.1, (v)).
\end{proof}

Soit alors $X$ un préschéma réduit n'ayant qu'un nombre fini de
composantes irréductibles, et soit $X'$ son normalisé (II, 6.3.8); on sait
(loc. cit.) que le morphisme canonique $f:X'\to X$ est surjectif. La
définition (6.15.1) montre alors que pour que $X$ soit géométriquement
unibranche en un point $x$, il faut et il suffit que $f$ soit radiciel en ce
point; pour que $X$ soit géométriquement unibranche, il faut et il suffit
donc que $f$ soit radiciel.

\begin{env}[6.15.4]
\EGAReaderTarget{ega4.c04.statement.6-15-4}{IV.2.s6.15.4-fr}
Généralisant la définition donnée dans (I, 2.2.9), nous dirons que, pour
deux préschémas réduits $X$, $Y$, un morphisme $f:X\to Y$ est
\emph{birationnel} si la restriction de $f$ à l'ensemble des points
maximaux de $X$ est une bijection de cet ensemble sur l'ensemble des points
maximaux de $Y$, et si, pour tout point maximal $y$ de $Y$, le morphisme
$X\times_Y\operatorname{Spec}(k(y))\to\operatorname{Spec}(k(y))$ déduit de
$f$ est un isomorphisme (autrement dit, si la fibre $f^{-1}(y)$ est réduite
à un seul point $x$ (maximal dans $X$) et si l'homomorphisme
$k(y)\to k(x)$ déduit de $f$ est bijectif); cette notion se réduit à celle
de (I, 2.2.9) lorsque $X$ et $Y$ n'ont qu'un nombre fini de composantes
irréductibles.
\end{env}

\begin{lemma}[6.15.4.1]
\label{IV.2.s6.15.4.1-fr}
Soit $f:X\to Y$ un morphisme, $g:Y'\to Y$ un morphisme plat,
$X'=X\times_Y Y'$, $f'=f_{(Y')}:X'\to Y'$. Si $f_{\mathrm{red}}$ est
birationnel, il en est de même de $f'_{\mathrm{red}}$.
\end{lemma}

\begin{proof}
En effet, si $y'$ est un point maximal de $Y'$, on sait (2.3.4, (ii)) que
$y=g(y')$ est un point maximal de $Y$; il existe un seul point $x$ de $X$
dans $f^{-1}(y)$ par hypothèse; de plus $x$ est maximal dans $X$ et tel que
$k(x)=k(y)$. Il résulte donc de (I, 3.4.9) que $f'^{-1}(y')$ est réduit à
un seul point $x'$ et que $k(x')=k(y')$. On conclut le raisonnement en
remarquant que d'après (2.3.7, (ii)), toute composante irréductible de $X'$
domine une composante irréductible de $Y'$.
\end{proof}

\begin{proposition}[6.15.5]
\label{IV.2.s6.15.5-fr}
Soit $f:X\to Y$ un morphisme tel que $f_{\mathrm{red}}$ soit entier et
birationnel (6.15.4), donc surjectif.
\begin{enumerate}
\item[(i)] Pour que $Y$ soit géométriquement unibranche en un point $y$,
il faut et il suffit que $f$ soit radiciel au point $y$ et que $X$ soit
géométriquement unibranche en l'unique point $x$ de $f^{-1}(y)$.
\item[(ii)] Pour que $Y$ soit géométriquement unibranche, il faut et il
suffit que $X$ soit géométriquement unibranche et que $f$ soit radiciel.
\end{enumerate}
\end{proposition}

\begin{proof}
On peut évidemment se borner au cas où $X$ et $Y$ sont réduits, le fait
que $f$ soit surjectif résulte de ce que $f$ est un morphisme fermé (II,
6.1.10) et que $f(X)$ contient par hypothèse tous les points maximaux de
$Y$. Cela montre aussitôt que (i) entraîne (ii). Utilisant (I, 3.6.5) et
(II, 6.1.5), on peut supposer, pour prouver (i), que
$Y=\operatorname{Spec}(\mathcal O_{Y,y})$, autrement dit
$Y=\operatorname{Spec}(A)$, où $A$ est local. Comme $f$ est affine, on a
aussi $X=\operatorname{Spec}(B)$. Si $A$ est géométriquement unibranche,
il est intègre, donc il en est de même de $B$ puisque $f$ est birationnel;
inversement, si $f$ est radiciel au point $y$ et $X$ géométriquement
unibranche au point $x$, $B$ n'a qu'un seul idéal maximal (puisqu'il est
entier sur $A$ et que tout idéal maximal de $B$ est donc au-dessus de celui
de $A$), autrement dit $B$ est un anneau local, et dire que $X$ est
géométriquement unibranche au point $x$ signifie que $B$ est
géométriquement unibranche, donc intègre. Comme $f$ est dominant par
hypothèse et que $A$ est réduit, on en conclut que $A\subset B$ (I,
1.2.7), donc $A$ est aussi intègre. Ainsi, pour prouver (i), on peut se
borner au cas où $A$ et $B$ sont locaux, intègres, $A$ étant contenu dans
$B$ et ayant même corps des fractions, et par suite (puisque $B$ est entier
sur $A$) même clôture intégrale $C$. Mais alors l'assertion résulte de
(6.15.3.1, (i)) appliqué aux morphismes
$\operatorname{Spec}(C)\to\operatorname{Spec}(B)$ et
$\operatorname{Spec}(B)\to\operatorname{Spec}(A)$, vu l'interprétation
donnée dans (6.15.3) de la propriété d'être géométriquement unibranche en
un point.
\end{proof}

\oldpage[IV-2]{179}
\begin{proposition}[6.15.6]
\label{IV.2.s6.15.6-fr}
Soient $k$ un corps, $X$ un $k$-préschéma. Si $X$ est normal, alors, pour
toute extension $k'$ de $k$, $X'=X\otimes_k k'$ est géométriquement
unibranche.
\end{proposition}

\begin{proof}
On sait que $k'$ est extension algébrique d'une extension pure $k_0$ de
$k$, et si $k''$ est la plus grande extension séparable de $k_0$ contenue
dans $k'$, $k'$ est une extension radicielle de $k''$ et $k''$ extension
séparable de $k$. On sait (6.14.2) que $X''=X\otimes_k k''$ est normal;
comme $X\otimes_k k'=X''\otimes_{k''}k'$, on voit qu'on peut se borner au
cas où l'extension $k'$ de $k$ est radicielle. En outre (I, 3.6.5), on peut
supposer que $X=\operatorname{Spec}(A)$, où $A$ est un anneau local intègre
et intégralement clos (puisque $X$ est normal). Le morphisme projection
$f:X'\to X$ est un homéomorphisme, puisque
$\operatorname{Spec}(k')\to\operatorname{Spec}(k)$ est un homéomorphisme
universel (2.4.5); comme $X'=\operatorname{Spec}(A')$ où
$A'=A\otimes_k k'$, on voit donc que $A'$ est un anneau local dont le
nilradical $\mathfrak N'$ est l'unique idéal premier minimal, d'où
$X'_{\mathrm{red}}=\operatorname{Spec}(A'_0)$, où
$A'_0=A'/\mathfrak N'$ est un anneau local intègre; en outre, si $K$ est
le corps des fractions de $A$, le corps des fractions $K'_0$ de $A'_0$ est
radiciel sur $K$, puisque le morphisme $f$ est radiciel. Comme $A'_0$ est
entier sur $A$, sa clôture intégrale $B$ est aussi la fermeture intégrale
de $A$ dans $K'_0$. Mais comme $A$ est intégralement clos, on sait
(Bourbaki, Alg. comm., chap. V, § 2, n$^o$ 3, lemme 3) que $B$ est
l'ensemble des $x\in K'_0$ dont une puissance $p^m$-ème (pour un $m$ assez
grand) appartient à $A$ ($p$ étant l'exposant caractéristique de $K$);
en outre il n'existe qu'un seul idéal premier de $B$ au-dessus de tout idéal
premier de $A$; en particulier $B$ est un anneau local et son corps résiduel
est une extension radicielle de celui de $A$, et a fortiori de celui de
$A'_0$, ce qui prouve que $A'_0$ est géométriquement unibranche, et par
suite il en est de même de $X'$.
\end{proof}

\begin{proposition}[6.15.7]
\EGAReaderTarget{ega4.c04.prop.6-15-7}{IV.2.s6.15.7-fr}
Soient $k$ un corps, $X$ un $k$-préschéma, $k'$ une extension de $k$,
$X'=X\otimes_k k'$. Soient $x'$ un point de $X'$, $x$ sa projection dans
$X$. Pour que $X$ soit géométriquement unibranche au point $x$, il faut et
il suffit que $X'$ soit géométriquement unibranche au point $x'$. Pour que
$X$ soit géométriquement unibranche, il faut et il suffit que $X'$ le soit.
\end{proposition}

\begin{proof}
La seconde assertion résulte de la première et du fait que la projection
$f:X'\to X$ est un morphisme surjectif. Pour démontrer la première
assertion, on peut, en vertu de (I, 5.1.8), se borner au cas où $X$ est
réduit, et en vertu de (I, 3.6.5), au cas
\oldpage[IV-2]{180}
où $X=\operatorname{Spec}(A)$ (avec $A=\mathcal O_{X,x}$) est un schéma local.
Notons que $A'=\mathcal O_{X',x'}$ est par hypothèse un $A$-module
fidèlement plat, donc contenant $A$, et par suite l'hypothèse que $X$ est
géométriquement unibranche au point $x$, ou l'hypothèse que $X'$ est
géométriquement unibranche au point $x'$, entraînent toutes deux que $A$
est un anneau local intègre (puisque $A$, étant réduit, est aussi isomorphe
à un sous-anneau de $A'_{\mathrm{red}}$). Soient $K$ le corps des fractions
de $A$, $B$ la clôture intégrale de $A$, et posons
$Y=\operatorname{Spec}(B)$; il revient au même de dire que $X$ est
géométriquement unibranche au point $x$ ou que le morphisme $g:Y\to X$
est radiciel au point $x$ (6.15.3). Soit
$Y'=Y\otimes_k k'=Y\times_X X'$ de sorte que l'on a le diagramme
commutatif
\[
\xymatrix{
Y \ar[d]_g & Y' \ar[l]^h \ar[d]_{g'}\\
X & X' \ar[l]_f
}
\]
Notons que $f$ est un morphisme plat; donc (6.15.4.1) le morphisme entier
$g'_{\mathrm{red}}$ est birationnel. D'autre part, comme $Y$ est normal,
$Y'$ est géométriquement unibranche (6.15.6). Pour que $X'$ soit
géométriquement unibranche en $x'$, il faut et il suffit donc que $g'$
soit radiciel au point $x'$ (6.15.5). Or, $x$ est la projection de $x'$ dans
$X$ et il est équivalent de dire que $g$ est radiciel au point $x$ ou que
$g'$ est radiciel au point $x'$ (6.15.3.1); enfin, pour que $X$ soit
géométriquement unibranche au point $x$, il faut et il suffit que $g$ soit
radiciel en ce point, d'où la proposition.
\end{proof}

\begin{lemma}[6.15.8]
\label{IV.2.s6.15.8-fr}
Soient $k$ un corps séparablement clos (autrement dit, tel que la clôture
algébrique de $k$ soit radicielle sur $k$), $X$ un $k$-préschéma localement
de type fini sur $k$, $x$ un point fermé de $X$. Si $X$ est unibranche au
point $x$, il est géométriquement unibranche en ce point.
\end{lemma}

\begin{proof}
En effet, on sait (I, 6.4.2) que $k(x)$ est une extension algébrique de
$k$; comme le corps résiduel de la clôture intégrale de
$(\mathcal O_{X,x})_{\mathrm{red}}$ (clôture intégrale qui est par
hypothèse un anneau local) est une extension algébrique de $k(x)$, c'est
une extension radicielle de $k$ par hypothèse, donc aussi de $k(x)$.
\end{proof}

\begin{corollary}[6.15.9]
\label{IV.2.s6.15.9-fr}
Soient $k$ un corps, $X$ un $k$-préschéma localement de type fini sur $k$,
$k'$ une extension séparablement close de $k$ (autrement dit, telle que la
clôture algébrique de $k'$ soit radicielle sur $k'$); alors, pour que $X$
soit géométriquement unibranche, il faut et il suffit que
$X'=X\otimes_k k'$ soit unibranche.
\end{corollary}

\begin{proof}
Compte tenu de (6.15.7), on est ramené à prouver que si $X$ est unibranche
et $k$ séparablement clos, alors $X$ est géométriquement unibranche. Or, $X$
est géométriquement unibranche en ses points fermés (6.15.8); nous en
conclurons dans (10.4.9) que $X$ est géométriquement unibranche en tous ses
points, en nous appuyant sur le fait (signalé à la fin de (6.15.2)) que
l'ensemble des points où $X$ est géométriquement unibranche est
constructible (bien entendu, le cor. (6.15.9) ne sera pas utilisé
jusque-là).
\end{proof}

Ce résultat justifie, en une certaine mesure, la terminologie de
« géométriquement unibranche ».

\oldpage[IV-2]{181}
\begin{proposition}[6.15.10]
\label{IV.2.s6.15.10-fr}
Soient $Y$, $Y'$ deux préschémas localement noethériens, $g:Y'\to Y$ un
morphisme normal (6.8.1), $f:X\to Y$ un morphisme; posons
$X'=X\times_Y Y'$ et soient $p:X'\to X$ la projection canonique, $x'$ un
point de $X'$. Alors, si $X$ est réduit (resp. géométriquement unibranche,
resp. intègre et géométriquement unibranche) au point $x=p(x')$, $X'$ est
réduit (resp. géométriquement unibranche, resp. intègre et
géométriquement unibranche) au point $x'$.
\end{proposition}

\begin{proof}
En vertu de (I, 3.6.5), on peut se borner au cas où
$Y=\operatorname{Spec}(A)$, $Y'=\operatorname{Spec}(A')$,
$X=\operatorname{Spec}(B)$, où $A$, $A'$ et $B$ sont des anneaux locaux,
les homomorphismes $A\to A'$ et $A\to B$ étant locaux, $A$, $A'$ des
anneaux noethériens, et $x$ étant le point fermé de $X$; il s'agit de
prouver que si $B$ est réduit (resp. géométriquement unibranche, resp.
intègre et géométriquement unibranche) alors $B'=B\otimes_A A'$ est réduit
(resp. $\operatorname{Spec}(B')$ est géométriquement unibranche aux points
de $p^{-1}(x)$, resp. $\operatorname{Spec}(B')$ est intègre et
géométriquement unibranche en ces points). Supposons d'abord seulement
$B$ réduit; $B$ est réunion de la famille croissante de ses
sous-$A$-algèbres de type fini $B_\lambda$, et par platitude, $B'$ est
réunion de la famille filtrante croissante de ses sous-$A'$-algèbres
$B'_\lambda=B_\lambda\otimes_A A'$. Or, le morphisme
$\operatorname{Spec}(B'_\lambda)\to\operatorname{Spec}(B_\lambda)$ est
normal (6.8.2), et a fortiori réduit; le fait que $B'_\lambda$ soit réduit
est donc conséquence de (3.3.5); on conclut que $B'$ lui-même est réduit par
(5.13.2).

Supposons maintenant $B$ géométriquement unibranche; compte tenu de (I,
5.1.8), on peut en outre supposer $B$ réduit, donc intègre. Soit $C$ sa
clôture intégrale, qui est par hypothèse un anneau local; posons
$Z=\operatorname{Spec}(C)$, $C'=C\otimes_A A'$,
$Z'=\operatorname{Spec}(C')=Z\times_X X'$, de sorte que l'on a le diagramme
commutatif
\[
\xymatrix{
Z \ar[d]_l & Z' \ar[l] \ar[d]_{l'}\\
X & X' \ar[l]_p
}
\]
En vertu de la première partie du raisonnement, $X'$ est réduit; d'autre
part, comme $Z'=Z\times_Y Y'$, il résulte de (6.14.1) que $Z'$ est normal
(et a fortiori géométriquement unibranche). Comme $f$ est entier et
birationnel, il en est de même de $f'$ par (6.15.4.1), puisque $p$ est
plat. Enfin, l'hypothèse que $X$ est géométriquement unibranche au point
$x$ entraîne que $f$ est radiciel en ce point (6.15.5); on conclut donc de
(6.15.3.1) que $f'$ est radiciel en tout point $x'\in p^{-1}(x)$, et il
résulte alors de (6.15.5) que $X'$ est géométriquement unibranche (donc
intègre puisqu'il est réduit) en chacun de ces points.
\end{proof}

\begin{remarks}[6.15.11]
\label{IV.2.s6.15.11-fr}
(i) Dans la démonstration de (6.15.10), on ne peut se dispenser de faire
appel à (6.14.1), même lorsque $X=Y$, car on fait intervenir la clôture
intégrale de l'anneau $B$, qui n'est pas nécessairement un anneau
noethérien même si $B$ l'est.

(ii) L'exemple (6.5.5, (ii)) montre que dans (6.15.10), on ne peut remplacer
« géométriquement unibranche » par « intègre » dans l'énoncé, même si
le corps résiduel $k(x)$ est algébriquement clos et le morphisme $g$ étale.

\oldpage[IV-2]{182}
On ne peut pas non plus dans cet énoncé remplacer « géométriquement
unibranche » par « unibranche ». Soit en effet $A$ l'anneau local intègre
complet
\[
\mathbb R[[U,V]]/(U^2+V^2);
\]
si $u$, $v$ sont les images de $U$, $V$ dans $A$, l'idéal maximal de $A$
est $Au+Av$. On vérifie aisément que la clôture intégrale de $A$ est
l'anneau $A[t]$, où $t=u/v$ vérifie la relation $t^2=-1$, si bien que
$A[t]$ est isomorphe à l'anneau local $\mathbb C[[U]]$; comme le corps
résiduel de $\mathbb C[[U]]$ est $\mathbb C$, on voit que $A$ est
unibranche mais non géométriquement unibranche. Mais
$A\otimes_{\mathbb R}\mathbb C$ n'est pas un anneau intègre, étant
isomorphe à
\[
\mathbb C[[U,V]]/(U-iV)(U+iV).
\]
\end{remarks}

\section{Relations entre un anneau local noethérien et son complété. Anneaux excellents}
\label{section:IV.7-fr}

Le présent paragraphe est surtout consacré à l'exposé de certaines propriétés
d'anneaux noethériens, généralement stables par passage à une algèbre de type
fini ou par localisation, qui sont vraies pour les anneaux que l'on rencontre
le plus souvent (algèbres de type fini sur $\mathbb Z$, ou sur un corps, ou sur
un anneau local noethérien complet), sans l'être pour tous les anneaux
noethériens. Une première série de propriétés, liées à la théorie de la
dimension, et notamment à la condition des chaînes, fait l'objet des
n\textsuperscript{os} 1 et
2; toutes les propriétés envisagées sont vraies pour les quotients d'anneaux
noethériens réguliers et leurs démonstrations dans ce cas sont pour la plupart
faciles et bien connues [30]; aussi pour ces anneaux les développements des
\S\S~1 et 2 sont sans intérêt.

Il n'en est plus de même des propriétés étudiées dans les
n\textsuperscript{os} 3, 4 et 5 (qui ne dépendent pas des
n\textsuperscript{os} 1 et 2); leurs démonstrations, même dans le cas des
localisés d'algèbres de type fini sur un corps ou sur $\mathbb Z$ (où elles
sont dues pour la plupart à Zariski et à Nagata) sont le plus souvent
délicates; des exemples classiques sont par exemple les équivalences
$A$ normal $\Longleftrightarrow \widehat A$ normal, et $A$ réduit
$\Longleftrightarrow \widehat A$ réduit. Nous développons une méthode
systématique pour formuler et prouver de telles propriétés, à l'aide des
propriétés des \emph{fibres} du morphisme canonique
$\operatorname{Spec}(\widehat A)\to\operatorname{Spec}(A)$, et en utilisant
les notions et résultats exposés au § 6; le succès de cette méthode repose en
dernière analyse sur les excellentes propriétés des \emph{anneaux locaux
complets} et
de ceux qui s'en déduisent par localisation ou passage à une algèbre de type
fini. A cet égard, le théorème de Nagata (6.12.7) disant que, pour un tel
anneau, le lieu singulier dans $\operatorname{Spec}(A)$ est fermé, joue un rôle
crucial; il en est de même (en ce qui concerne les propriétés de permanence)
du critère de régularité (0, 22.3.4), de nature plus technique.

Enfin, dans les n\textsuperscript{os} 6 et 7, nous appliquons les résultats
obtenus à l'étude de
la finitude de la clôture intégrale des anneaux intègres noethériens.

Le lecteur notera que les résultats les plus délicats du § 7 ne serviront
qu'assez peu dans la suite du chapitre IV, et même dans les chapitres
ultérieurs.

\subsection{Équidimensionalité formelle et anneaux formellement caténaires}
\label{subsection:IV.7.1-fr}

\oldpage[IV-2]{183}

\begin{definition}[7.1.1]
\label{IV.2.s7.1.1-fr}
On dit qu'un anneau local noethérien $A$ est \emph{formellement
équidimensionnel} si son complété $\widehat A$ est équidimensionnel (0,
16.1.4).
\end{definition}

\begin{proposition}[7.1.2]
\label{IV.2.s7.1.2-fr}
Soient $A$ un anneau local noethérien, $\mathfrak p_i$ ($1\leqslant i\leqslant
n$) ses idéaux premiers minimaux. Pour que $A$ soit formellement
équidimensionnel, il faut et il suffit que les $A/\mathfrak p_i$ le soient et
que $A$ soit équidimensionnel (autrement dit que les $A/\mathfrak p_i$ aient
même dimension).
\end{proposition}

\begin{proof}
En effet, cela résulte du fait que pour tout idéal premier $\mathfrak p'$ de
$\widehat A$, $\mathfrak p'\cap A$ contient l'un des $\mathfrak p_i$, donc
$\mathfrak p'$ contient l'un des $\mathfrak p_i\widehat A$, et
$\widehat A/\mathfrak p_i\widehat A=(A/\mathfrak p_i)\otimes_A\widehat A$
est le complété de l'anneau local $A/\mathfrak p_i$ (0$_{\mathrm I}$,
7.3.3), donc a même dimension. Toute chaîne maximale d'idéaux premiers de
$\widehat A$ s'identifie donc canoniquement à une chaîne maximale d'idéaux
premiers de l'un des $\widehat A/\mathfrak p_i\widehat A$, et réciproquement,
d'où aussitôt la conclusion.
\end{proof}

\begin{proposition}[7.1.3]
\label{IV.2.s7.1.3-fr}
Soient $A$, $A'$ deux anneaux locaux noethériens, $\mathfrak m$ l'idéal
maximal de $A$, $\varphi:A\to A'$ un homomorphisme local; on suppose que $A'$
est un $A$-module plat, et que $A'$ soit équidimensionnel et caténaire. Alors:
\begin{enumerate}
\item[(i)] $A$ est équidimensionnel et caténaire.
\item[(ii)] Supposons en outre que $\mathfrak mA'$ soit un idéal de définition de
$A'$. Alors, si $\mathfrak a$ est un idéal de $A$, pour que
$A'/\mathfrak aA'$ soit équidimensionnel, il faut et il suffit que
$A/\mathfrak aA$ le soit. En particulier, pour tout idéal premier
$\mathfrak p$ de $A$, $A'/\mathfrak pA'$ est équidimensionnel.
\end{enumerate}
\end{proposition}

\begin{proof}
Prouvons d'abord la seconde assertion de (ii). Posons
$X=\operatorname{Spec}(A)$, $X'=\operatorname{Spec}(A')$,
$Y=V(\mathfrak p)=\operatorname{Spec}(A/\mathfrak p)$,
$Y'=V(\mathfrak pA')=\operatorname{Spec}(A'/\mathfrak pA')$; si
$f:X'\to X$ est le morphisme correspondant à $\varphi$, $Y'$ est aussi le
sous-préschéma fermé $f^{-1}(Y)$ de $X'$. Comme $\mathfrak mA'$ est un idéal
de définition de $A'$, on a
\begin{equation}
\dim(X)=\dim(X')
\tag{7.1.3.1}
\label{IV.2.eq.7.1.3.1-fr}
\end{equation}
en vertu de (6.1.3); de même $\mathfrak mA'/\mathfrak pA'$ est un idéal de
définition de $A'/\mathfrak pA'$ et $Y'$ est plat sur $Y$ (2.1.4), donc
(6.1.3)
\begin{equation}
\dim(Y)=\dim(Y').
\tag{7.1.3.2}
\label{IV.2.eq.7.1.3.2-fr}
\end{equation}
Soient $y$ le point générique de $Y$, $Y'_i$ ($1\leqslant i\leqslant r$) les
composantes irréductibles de $Y'$, $y'_i$ le point générique de $Y'_i$
($1\leqslant i\leqslant r$); on a $f(y'_i)=y$ pour tout $i$ (2.3.4).
D'autre part, $y'_i$ est point maximal de la fibre $f^{-1}(y)$ (0$_{\mathrm
I}$, 2.1.8), donc
$\mathscr O_{X',y'_i}\otimes_{\mathscr O_{X,y}}k(y)=
\mathscr O_{X',y'_i}/\mathfrak m_y\mathscr O_{X',y'_i}$ est de dimension
0, autrement dit $\mathfrak m_y\mathscr O_{X',y'_i}$ est un idéal de
définition de $\mathscr O_{X',y'_i}$; on peut de nouveau appliquer (6.1.3),
qui donne
\begin{equation}
\dim(\mathscr O_{X',y'_i})=\dim(\mathscr O_{X,y}),
\tag{7.1.3.3}
\label{IV.2.eq.7.1.3.3-fr}
\end{equation}
c'est-à-dire (5.1.2)
\begin{equation}
\operatorname{codim}(Y'_i,X')=\operatorname{codim}(Y,X)
\tag{7.1.3.4}
\label{IV.2.eq.7.1.3.4-fr}
\end{equation}
pour tout $i$. Comme $A'$ est équidimensionnel et caténaire, on a, en vertu
de (0, 14.3.5)
\begin{equation}
\dim(X')=\dim(Y'_i)+\operatorname{codim}(Y'_i,X')
\tag{7.1.3.5}
\label{IV.2.eq.7.1.3.5-fr}
\end{equation}

\oldpage[IV-2]{184}
pour tout $i$; en vertu de (7.1.3.1), (7.1.3.2) et (7.1.3.4) et de
l'inégalité $\dim(Y'_i)\leqslant\dim(Y')$, on en déduit
\begin{equation}
\dim(X)\leqslant\dim(Y)+\operatorname{codim}(Y,X)
\tag{7.1.3.6}
\label{IV.2.eq.7.1.3.6-fr}
\end{equation}
donc les deux membres de cette inégalité sont égaux d'après (0, 14.2.2.2);
en outre cette égalité entraîne
\begin{equation}
\dim(Y'_i)=\dim(Y')=\dim(Y)
\tag{7.1.3.7}
\label{IV.2.eq.7.1.3.7-fr}
\end{equation}
pour tout $i$, ce qui prouve la seconde assertion de (ii).

Démontrons maintenant la première assertion de (ii). Soient $\mathfrak p_j$
($1\leqslant j\leqslant n$) les idéaux premiers de $A$ minimaux parmi ceux
qui contiennent $\mathfrak a$, et soient $\mathfrak p'_{jh}$ les idéaux
premiers de $A'$ minimaux parmi ceux contenant $\mathfrak p_jA'$
($1\leqslant h\leqslant m_j$); on sait alors (2.3.4) que les
$\mathfrak p'_{jh}$ ($1\leqslant j\leqslant n$, $1\leqslant h\leqslant
m_j$ pour tout $j$) sont aussi les idéaux premiers de $A'$ minimaux parmi
ceux qui contiennent $\mathfrak aA'$. D'après ce qu'on a vu plus haut, pour
tout $j$, les dimensions des anneaux $A'/\mathfrak p'_{jh}$
($1\leqslant h\leqslant m_j$) sont toutes égales et sont donc égales à
$\dim(A'/\mathfrak p_jA')$, elle-même égale à $\dim(A/\mathfrak p_j)$ par
(7.1.3.2). Dire que $A'/\mathfrak aA'$ est équidimensionnel signifie donc
que les dimensions des $A/\mathfrak p_j$ sont toutes égales, c'est-à-dire
que $A/\mathfrak a$ est équidimensionnel.

Prouvons enfin (i). Soit $\mathfrak r'$ un idéal premier de $A'$ minimal
parmi ceux contenant $\mathfrak mA'$, et soit $A''=A'_{\mathfrak r'}$;
comme $A''$ est un $A'$-module plat, c'est aussi un $A$-module plat; en
outre $A''$ est équidimensionnel et caténaire (0, 16.1.4), et
$\mathfrak mA''$ est un idéal de définition de $A''$. On peut donc se
borner au cas où $\mathfrak mA'$ est un idéal de définition de $A'$. Si
alors $\mathfrak p$ et $\mathfrak q\subset\mathfrak p$ sont deux idéaux
premiers de $A$, on peut appliquer (ii) à $A/\mathfrak q$ et à
$A'/\mathfrak qA'$, qui est équidimensionnel et plat sur
$A/\mathfrak q$; si $Z=\operatorname{Spec}(A/\mathfrak q)$ et
$Y=\operatorname{Spec}(A/\mathfrak p)$, on a donc
$\dim(Z)=\dim(Y)+\operatorname{codim}(Y,Z)$; en outre on peut appliquer
(7.1.3.6), qui montre que
$\operatorname{codim}(Y,X)=\dim X-\dim Y$; comme $X$ est
équicodimensionnel, ces relations montrent qu'il est biéquidimensionnel en
vertu de (0, 14.3.3).
\end{proof}

\begin{corollary}[7.1.4]
\label{IV.2.s7.1.4-fr}
Soit $A$ un anneau local noethérien formellement équidimensionnel. Alors:
\begin{enumerate}
\item[(i)] $A$ est équidimensionnel et caténaire (autrement dit,
biéquidimensionnel).
\item[(ii)] Soit $\mathfrak a$ un idéal de $A$; pour que $A/\mathfrak a$ soit
équidimensionnel, il faut et il suffit que $A/\mathfrak a$ soit formellement
équidimensionnel. En particulier, pour tout idéal premier $\mathfrak p$ de
$A$, $A/\mathfrak p$ est formellement équidimensionnel.
\end{enumerate}
\end{corollary}

\begin{proof}
Si $\mathfrak m$ est l'idéal maximal de $A$, $\mathfrak m\widehat A$ est
idéal de définition de $\widehat A$. Par hypothèse $\widehat A$ est
équidimensionnel, et on sait d'autre part qu'il est caténaire (5.6.4); il
suffit alors d'appliquer (7.1.3) à $A'=\widehat A$.
\end{proof}

\begin{corollary}[7.1.5]
\label{IV.2.s7.1.5-fr}
Soit $A$ un anneau local noethérien tel qu'il existe un $A$-module de type
fini $M$ qui soit un $A$-module de Cohen--Macaulay et pour lequel
$\operatorname{Supp}(M)=\operatorname{Spec}(A)$ (ce qui aura lieu par
exemple si $A$ est un anneau de Cohen--Macaulay). Alors $A$ est formellement
équidimensionnel, donc (7.1.4) pour qu'un anneau quotient $B$ de $A$ soit
formellement équidimensionnel, il faut et il suffit qu'il soit
équidimensionnel.
\end{corollary}

\begin{proof}
\oldpage[IV-2]{185}
En effet, $\widehat M=M\otimes_A\widehat A$ est un $\widehat A$-module de
Cohen--Macaulay (0, 16.5.2), de support égal à
$\operatorname{Spec}(\widehat A)$; par suite (0, 16.5.4), $\widehat A$ est
équidimensionnel.
\end{proof}

\begin{remark}[7.1.6]
\label{IV.2.s7.1.6-fr}
On a vu (6.3.8) que sous les hypothèses de (7.1.5), pour tout anneau quotient
$B$ de $A$, les fibres du morphisme canonique
$\operatorname{Spec}(\widehat B)\to\operatorname{Spec}(B)$ sont des schémas
de Cohen--Macaulay; par suite ((6.4.3) et (5.7.5)), \emph{pour que $B$ n'ait
pas d'idéaux premiers associés immergés, il faut et il suffit qu'il en soit
de même de $\widehat B$}.
\end{remark}

\begin{lemma}[7.1.7]
\label{IV.2.s7.1.7-fr}
Soient $A$ un anneau local noethérien, $\mathfrak p_i$
($1\leqslant i\leqslant r$) ses idéaux premiers minimaux. Supposons que
chacun des anneaux $A/\mathfrak p_i$ soit formellement équidimensionnel.
Alors, pour tout idéal premier $\mathfrak q$ de $A$ et tout $i$ tel que
$\mathfrak p_i\subset\mathfrak q$, l'anneau
$A_{\mathfrak q}/\mathfrak p_iA_{\mathfrak q}$ est formellement
équidimensionnel.
\end{lemma}

\begin{proof}
Comme $A_{\mathfrak q}/\mathfrak p_iA_{\mathfrak q}$ est l'anneau local de
$A/\mathfrak p_i$ en l'idéal premier $\mathfrak q/\mathfrak p_i$, on peut
se borner au cas où $A$ est intègre et formellement équidimensionnel. Posons
$A'=\widehat A$, et soit $\mathfrak q'$ un des idéaux premiers de $A'$
minimaux parmi ceux contenant $\mathfrak qA'$; si l'on pose
$B=A_{\mathfrak q}$, $B'=A'_{\mathfrak q'}$, $B'$ est un $B$-module plat
(0$_{\mathrm I}$, 6.3.2). Posons $C=\widehat B$, $C'=\widehat {B'}$; comme
$B'$ est un $B$-module plat, on sait que $C'$ est un $C$-module plat
(Bourbaki, \emph{Alg. comm.}, chap. III, § 5, n\textsuperscript{o} 4,
prop. 4). Comme $A'$ est caténaire (5.6.4) et équidimensionnel par
hypothèse, il en est de même de $B'=A'_{\mathfrak q'}$ (0, 16.1.4); en
outre, comme $A'$ est isomorphe à un quotient d'anneau régulier en vertu du
théorème de Cohen (0, 19.8.8), il en est de même de $B'$ (0, 17.3.9); on
conclut donc de (7.1.5) que $C'$ est équidimensionnel. D'autre part $C'$ est
caténaire (5.6.4), donc $C$ est équidimensionnel en vertu de (7.1.3, (ii)).
\end{proof}

\begin{theorem}[7.1.8]
\label{IV.2.s7.1.8-fr}
Soit $A$ un anneau local noethérien. Les conditions suivantes sont
équivalentes:
\begin{enumerate}
\item[a)] Tout anneau quotient intègre de $A$ est formellement équidimensionnel.
\item[b)] Les anneaux quotients de $A$ par ses idéaux premiers minimaux sont
formellement équidimensionnels.
\item[c)] Tout anneau local $B$ qui est une $A$-algèbre essentiellement de type
fini (1.3.8) et qui est équidimensionnel, est formellement
équidimensionnel.
\end{enumerate}
\end{theorem}

\begin{proof}
Il est trivial que c) entraîne a) et que a) entraîne b). D'autre part,
comme tout idéal premier de $A$ contient un idéal premier minimal de $A$,
b) entraîne a) en vertu de (7.1.4, (ii)). Il reste à voir que a) implique
c).

Il suffit de prouver que les quotients de $B$ par ses idéaux premiers
minimaux sont formellement équidimensionnels (7.1.2); comme tout anneau
quotient de $B$ est encore une $A$-algèbre essentiellement de type fini
(1.3.9), on peut en premier lieu supposer $B$ \emph{intègre}. Si
$\mathfrak q$ est
l'image réciproque dans $A$ de l'idéal maximal de $B$, on sait que $B$ est
aussi une $A_{\mathfrak q}$-algèbre essentiellement de type fini (1.3.10),
et il résulte de a) et de (7.1.7) que tout anneau quotient intègre de
$A_{\mathfrak q}$ est formellement équidimensionnel; on peut donc supposer
que l'homomorphisme $A\to B$ est un homomorphisme \emph{local}. Le noyau
$\mathfrak r$ de cet homomorphisme est alors un idéal premier de $A$, et en
vertu de a), tout anneau quotient intègre de $A/\mathfrak r$ est
formellement équidimensionnel; comme $B$ est une
$(A/\mathfrak r)$-algèbre essentiellement de type fini, on voit qu'on peut
aussi supposer

\oldpage[IV-2]{186}
que $A$ est \emph{intègre} et est un sous-anneau de $B$. On sait (1.3.11)
que $B$
est quotient d'un anneau local de la forme $C_{\mathfrak p}$, où
$C=A[T_1,\ldots,T_n]$ est une algèbre de polynômes et $\mathfrak p$ un
idéal premier de $C$ au-dessus de l'idéal maximal $\mathfrak m$ de $A$. En
vertu de (7.1.7), il suffit de prouver que $C_{\mathfrak p}$ est
formellement équidimensionnel, autrement dit on peut se borner au cas où
$B=C_{\mathfrak p}$.

Posons $A'=\widehat A$, et
$C'=A'[T_1,\ldots,T_n]=C\otimes_AA'$; il y a un idéal premier unique
$\mathfrak p'$ de $C'$ au-dessus de l'idéal maximal $\mathfrak mA'$ de
$A'$, donc au-dessus de $\mathfrak p$; posons $B'=C'_{\mathfrak p'}$;
l'homomorphisme $B\to B'$ est local et fait de $B'$ un $B$-module plat. On
sait alors que $\widehat {B'}$ est un $\widehat B$-module plat (Bourbaki,
\emph{Alg. comm.}, chap. III, § 5, n\textsuperscript{o} 4, prop. 4); comme
$\widehat {B'}$ est caténaire (5.6.4), il suffira de prouver que
$\widehat {B'}$ est \emph{équidimensionnel} pour en déduire que $\widehat B$ l'est
aussi (7.1.3, (i)), ce qui achèvera la démonstration.

Or, $A'$ est quotient d'un anneau régulier par le théorème de Cohen (0,
19.8.8), donc il en est de même de $B'$ (0, 17.3.9); en vertu de (7.1.5),
pour montrer que $B'$ est formellement équidimensionnel, il suffit de
prouver que $B'$ est équidimensionnel; et pour cela, il suffit de montrer
que $C'$ est équidimensionnel, puisque $C'$ est quotient d'anneau régulier,
donc caténaire (5.6.4). Or les idéaux premiers minimaux de $C'$ sont les
idéaux $\mathfrak q'C'$, où $\mathfrak q'$ parcourt les idéaux premiers
minimaux de $A'$ (5.5.3), et on a
$C'/\mathfrak q'C'=(A'/\mathfrak q')[T_1,\ldots,T_n]$. Comme $A'$ est
équidimensionnel par hypothèse (puisque $A$ est intègre), il en est donc de
même de $C'$ par (5.5.4). C.Q.F.D.
\end{proof}

\begin{definition}[7.1.9]
\label{IV.2.s7.1.9-fr}
Lorsque les conditions équivalentes de (7.1.8) sont vérifiées, on dit que
$A$ est un anneau \emph{formellement caténaire}.
\end{definition}

Pour un anneau local noethérien, il revient donc au même de dire qu'il est
formellement équidimensionnel ou formellement caténaire et
équidimensionnel, en vertu de (7.1.2).

\begin{proposition}[7.1.10]
\label{IV.2.s7.1.10-fr}
Tout anneau local $A$ quotient d'un anneau local de Cohen--Macaulay est
formellement caténaire.
\end{proposition}

\begin{proof}
Cela résulte aussitôt de (7.1.8) et de (7.1.5) appliqué aux anneaux
quotients intègres de $A$.
\end{proof}

\begin{proposition}[7.1.11]
\label{IV.2.s7.1.11-fr}
\begin{enumerate}
\item[(i)] Un anneau local noethérien formellement caténaire est universellement
caténaire.
\item[(ii)] Si $A$ est un anneau local noethérien formellement caténaire, tout
anneau local $B$ qui est une $A$-algèbre essentiellement de type fini est
formellement caténaire.
\end{enumerate}
\end{proposition}

\begin{proof}
L'assertion (ii) résulte aussitôt de la condition c) de (7.1.8) et du fait
que si un anneau local $C$ est une $B$-algèbre essentiellement de type
fini, $C$ est aussi une $A$-algèbre essentiellement de type fini (1.3.9).
Pour prouver (i), notons d'abord qu'en vertu de la condition b) de (7.1.8),
un anneau local noethérien formellement caténaire $A$ est caténaire,
puisque les quotients de $A$ par ses idéaux premiers minimaux le sont
(7.1.4). Si maintenant $E$ est une $A$-algèbre de type fini, il résulte de
ce qui précède et de (ii) que tout anneau local de $E$ en un idéal premier
est caténaire, donc que $E$ est caténaire. Donc $A$ est universellement
caténaire.
\end{proof}

\oldpage[IV-2]{187}

\begin{remarks}[7.1.12]
\label{IV.2.s7.1.12-fr}
\begin{enumerate}
\item[(i)] On ignore si, réciproquement, un anneau local noethérien
universellement caténaire est toujours formellement caténaire.
\item[(ii)] Un anneau local noethérien $A$ \emph{de dimension 1} est nécessairement
équidimensionnel, son idéal maximal ne pouvant être minimal (car dans ce
cas $A$ serait de dimension 0). Comme $\dim(\widehat A)=1$, cela montre que
$A$ est même \emph{formellement équidimensionnel}, et \emph{a fortiori}
formellement caténaire (donc \emph{universellement caténaire}). Par contre,
l'anneau local de dimension 2 défini dans (5.6.11), qui est caténaire mais
non universellement caténaire, est \emph{a fortiori} non formellement
caténaire.
\end{enumerate}
\end{remarks}

\begin{corollary}[7.1.13]
\label{IV.2.s7.1.13-fr}
Soient $Y$ un préschéma localement noethérien irréductible de dimension 1,
$X$ un préschéma irréductible, $f:X\to Y$ un morphisme dominant de type
fini, $\xi$, $\eta$ les points génériques de $X$ et $Y$ respectivement.
Alors, pour tout $y\in f(X)$, les dimensions des composantes irréductibles
de $f^{-1}(y)$ sont toutes égales à
$\deg.\operatorname{tr}_{k(\eta)}k(\xi)$.
\end{corollary}

\begin{proof}
Par hypothèse, $f^{-1}(\eta)$ est irréductible de point générique $\xi$ et
de dimension $\deg.\operatorname{tr}_{k(\eta)}k(\xi)$ (5.2.1). Comme $Y$
est irréductible et de dimension 1, on a $\dim(\mathscr O_y)=1$ pour tout
$y\neq\eta$, et $\mathscr O_y$ est universellement caténaire en vertu de
(7.1.12, (ii)). Si $y\in f(X)$ et si $z$ est un point générique d'une
composante irréductible $Z$ de $f^{-1}(y)$, on a donc, par (5.6.5)
\[
\dim(Z)=1+\dim(f^{-1}(\eta))-\dim(\mathscr O_{X,z}).
\]
Mais on a $\dim(\mathscr O_{X,z})\leqslant\dim(\mathscr O_y)$ (0,
16.3.9), et d'autre part, comme $z$ n'est pas le point générique de $X$,
$\dim(\mathscr O_{X,z})>0$; donc $\dim(\mathscr O_{X,z})=1$ et
$\dim(Z)=\dim(f^{-1}(\eta))$.
\end{proof}

\subsection{Anneaux strictement formellement caténaires}
\label{subsection:IV.7.2-fr}

\begin{env}[7.2.1]
\label{IV.2.s7.2.1-fr}
\emph{Notations.} --- Pour un anneau intègre noethérien $A$, nous noterons
$A^{(1)}$ l'intersection des anneaux locaux $A_{\mathfrak p}$, où
$\mathfrak p$ parcourt l'ensemble des idéaux premiers de $A$ de hauteur 1
(5.10.17). Si $A$ est un anneau local noethérien intègre, nous noterons
$A^{(\omega)}$ l'intersection des $A_{\mathfrak p}$ où cette fois
$\mathfrak p$ parcourt l'ensemble de tous les idéaux premiers de $A$
distincts de l'idéal maximal.

Nous dirons qu'un anneau local noethérien $A$ est \emph{strictement
équidimensionnel} si, pour tout idéal premier
$\mathfrak p\in\operatorname{Ass}(A)$, on a
$\dim(A/\mathfrak p)=\dim(A)$; il revient au même de dire que $A$ est
équidimensionnel et \emph{sans idéaux premiers associés immergés}.
\end{env}

\begin{example}[7.2.1.1]
\label{IV.2.s7.2.1.1-fr}
Soit $A$ un anneau local noethérien \emph{de dimension 1}. Alors les
conditions suivantes sont équivalentes:
\begin{enumerate}
\item[a)] $A$ est sans idéaux premiers associés immergés.
\item[$a'$)] $\widehat A$ est sans idéaux premiers associés immergés.
\item[b)] $A$ est strictement équidimensionnel.
\item[$b'$)] $\widehat A$ est strictement équidimensionnel.
\item[c)] $A$ est un anneau de Cohen--Macaulay.
\end{enumerate}
\end{example}

En effet a) signifie que l'idéal maximal $\mathfrak m$ de $A$ n'est pas
associé à $A$, donc les idéaux premiers associés à $A$ sont les idéaux
minimaux de $A$, tous distincts de $\mathfrak m$, et

\oldpage[IV-2]{188}

pour un tel idéal $\mathfrak p$, on a nécessairement
$\dim(A/\mathfrak p)=1$; donc a) implique b). Inversement, b) implique que
tout idéal premier $\mathfrak p\in\operatorname{Ass}(A)$ est distinct de
$\mathfrak m$, donc minimal, et par suite b) implique a). On sait déjà que
a) et c) sont équivalentes pour un anneau local de dimension 1 (5.7.8).
Comme il revient au même de dire que $A$ est un anneau de Cohen--Macaulay ou
que $\widehat A$ est un anneau de Cohen--Macaulay (0, 16.5.2), et que l'on
a $\dim(\widehat A)=1$, on voit enfin que $a'$) et $b'$) sont équivalentes à
c).

\begin{proposition}[7.2.2]
\label{IV.2.s7.2.2-fr}
Soient $A$ un anneau local noethérien intègre; posons $A'=\widehat A$,
$X=\operatorname{Spec}(A)$, $X'=\operatorname{Spec}(A')$ et désignons par
$f:X'\to X$ le morphisme canonique. Soient $a$ le point fermé de $X$,
$j:X-\{a\}\to X$ l'injection canonique. Soient $\mathscr F$ un
$\mathscr O_X$-Module cohérent,
$\mathscr F'=f^*(\mathscr F)=\mathscr F\otimes_{\mathscr O_X}\mathscr O_{X'}$;
alors les conditions suivantes sont équivalentes:
\begin{enumerate}
\item[a)] Le $\mathscr O_X$-Module
$j_*(\mathscr F|_{X-\{a\}})$ est cohérent.
\item[b)] Pour tout $x'\in\operatorname{Ass}(\mathscr F')$, on a
$\dim(\overline{\{x'\}})\geqslant2$.
\end{enumerate}
\end{proposition}

\begin{proof}
Soit en effet $a'$ le point fermé de $X'$, qui est l'unique point de la
fibre $f^{-1}(a)$, et soit $j':X'-\{a'\}\to X'$ l'injection canonique.
Comme le morphisme $f$ est fidèlement plat et quasi-compact, il est
équivalent de dire que $j_*(\mathscr F|_{X-\{a\}})$ est cohérent ou que
$j'_*(\mathscr F'|_{X'-\{a'\}})$ est cohérent (5.9.5). D'autre part, comme
$A'$ est isomorphe au quotient d'un anneau local régulier par le théorème de
Cohen (0, 19.8.8), il résulte de (5.11.4) que la condition b) de l'énoncé
équivaut au fait que $j'_*(\mathscr F'|_{X'-\{a'\}})$ est cohérent. D'où la
proposition.

Au lieu d'appliquer (5.11.4) en utilisant le th. de Cohen, ce qui (par
(5.11.2)) fait implicitement appel aux résultats cohomologiques du chap.
III, on peut aussi utiliser le fait que $A'$ est universellement caténaire
(5.6.4) et universellement japonais en vertu de (7.6.5) (la démonstration de
ce dernier résultat n'utilisant pas (7.2.2)).
\end{proof}

\begin{proposition}[7.2.3]
\label{IV.2.s7.2.3-fr}
Soit $A$ un anneau local noethérien intègre. Les notations étant celles de
(7.2.2), les conditions suivantes sont équivalentes:
\begin{enumerate}
\item[a)] $A^{(1)}$ est une $A$-algèbre finie.
\item[b)] Pour toute partie fermée $T$ de $X$ de codimension $\geqslant2$,
si l'on désigne par $i:X-T\to X$ l'injection canonique,
$i_*(\mathscr O_{X-T})$ est un $\mathscr O_X$-Module cohérent.
\item[c)] Pour tout $x'\in\operatorname{Ass}(\mathscr O_{X'})$ et toute
partie fermée $T$ de $X$ de codimension $\geqslant2$, on a
$\operatorname{codim}(f^{-1}(T)\cap\overline{\{x'\}},
\overline{\{x'\}})\geqslant2$.
\end{enumerate}
\end{proposition}

\begin{proof}
Posons $Z=Z^{(2)}(X)$ (5.10.13) et $Z'=f^{-1}(Z)$, qui sont des parties
stables par spécialisation; les conditions a) et b) équivalent
respectivement aux propriétés suivantes:

$a_1)$ $\mathscr H^0_{X/Z}(\mathscr O_X)$ est cohérent;

$b_1)$ $\mathscr H^0_{X/T}(\mathscr O_X)$ est cohérent pour toute partie
fermée $T$ de $X$ de codimension $\geqslant2$.

Compte tenu de (5.9.5), ces deux dernières propriétés équivalent
respectivement à:

$a'_1)$ $\mathscr H^0_{X'/Z'}(\mathscr O_{X'})$ est cohérent;

$b'_1)$ posant $T'=f^{-1}(T)$,
$\mathscr H^0_{X'/T'}(\mathscr O_{X'})$ est cohérent pour toute partie
fermée $T$ de $X$ de codimension $\geqslant2$.

Or, tout point de $\operatorname{Ass}(\mathscr O_{X'})$ se projette dans
$X$ sur le point générique de $X$ puisque $f$ est un morphisme plat
(3.3.2); comme par définition ce point n'appartient pas à $Z$, on voit que
$\operatorname{Ass}(\mathscr O_{X'})$ ne rencontre pas $Z'$;
l'équivalence de $a'_1)$ et $b'_1)$ résulte donc de (5.11.5), puisque $A'$
est isomorphe à un quotient d'anneau régulier en vertu du théorème de Cohen
(0, 19.8.8) (ou encore

\oldpage[IV-2]{189}

parce que $A'$ est universellement japonais (7.6.5)). Pour cette raison,
les conditions $b'_1)$ et c) sont aussi équivalentes en vertu de (5.11.4).
\end{proof}

\begin{proposition}[7.2.4]
\label{IV.2.s7.2.4-fr}
Soit $A$ un anneau local noethérien et posons
$X=\operatorname{Spec}(A)$. Les conditions suivantes sont équivalentes:
\begin{enumerate}
\item[a)] Pour tout anneau quotient intègre $B$ de $A$, l'anneau $B^{(1)}$
est une $B$-algèbre finie.
\item[b)] Pour tout $\mathscr O_X$-Module cohérent $\mathscr F$ et toute
partie $Z$, stable par spécialisation, et telle que pour tout
$x\in\operatorname{Ass}(\mathscr F)\cap(X-Z)$, on ait
$\operatorname{codim}(\overline{\{x\}}\cap Z,\overline{\{x\}})\geqslant2$,
le $\mathscr O_X$-Module $\mathscr H^0_{X/Z}(\mathscr F)$ est cohérent.
\item[c)] Pour toute partie fermée $T$ de $X$ et tout
$\mathscr O_U$-Module cohérent $\mathscr G$ (où $U=X-T$) tel que, pour tout
$x\in\operatorname{Ass}(\mathscr G)$, on ait
$\operatorname{codim}(\overline{\{x\}}\cap T,\overline{\{x\}})\geqslant2$,
$i_*(\mathscr G)$ (où $i:U\to X$ est l'injection canonique) est un
$\mathscr O_X$-Module cohérent.
\item[d)] Pour tout anneau quotient intègre $B$ de $A$, et tout idéal
$\mathfrak J$ de hauteur $\geqslant2$ dans $B$, l'anneau
$\displaystyle\bigcap_{\mathfrak p\not\supset\mathfrak J}B_{\mathfrak p}$
est une $B$-algèbre finie.
\item[e)] Pour tout anneau quotient intègre $B$ de $A$ et tout anneau local
$C=B_{\mathfrak q}$ en un idéal premier $\mathfrak q$ de $B$, tel que
$\dim(C)\geqslant2$, l'anneau $C^{(\omega)}$ est une $C$-algèbre finie (ou,
ce qui revient au même, si $Y=\operatorname{Spec}(C)$ et si $U$ est le
complémentaire dans $Y$ du point fermé de $C$ et $j:U\to Y$ l'injection
canonique, $j_*(\mathscr O_U)$ est un $\mathscr O_Y$-Module cohérent).
\item[f)] Pour tout anneau quotient intègre $B$ de $A$, tout anneau local
$C=B_{\mathfrak q}$ en un idéal premier $\mathfrak q$ de $B$, tel que
$\dim(C)\geqslant2$, et pour tout idéal
$\mathfrak r\in\operatorname{Ass}(\widehat C)$, on a
$\dim(\widehat C/\mathfrak r)\geqslant2$.
\end{enumerate}
\end{proposition}

\begin{proof}
On sait déjà (5.11.6) que a) et b) sont équivalentes, ainsi que c) et d),
et que a) entraîne d). L'équivalence de a) et d) dans le cas actuel résulte
de (7.2.3) appliqué à un anneau quotient intègre $B$ de $A$, la condition
d) étant une formulation équivalente de la condition (7.2.3, b)) en vertu
de (5.9.3.1). L'équivalence de e) et f) est une conséquence de (7.2.2)
appliqué au $\mathscr O_Y$-Module cohérent $\mathscr O_Y$ lui-même. On a
déjà remarqué (5.11.7, (ii)) que si $A$ vérifie a), il en est de même de
toute $A$-algèbre finie et de tout anneau de fractions de $A$. La condition
a) pour $A$ entraîne donc (avec les notations de e)) que l'anneau $C$
vérifie a), donc aussi c); mais comme $C$ est intègre et $\dim(C)\geqslant
2$, on peut appliquer à $C$ la condition c) en prenant pour $T$ le point
fermé de $Y=\operatorname{Spec}(C)$; cela prouve que a) entraîne e). Reste
à prouver que e) entraîne a); on peut évidemment se borner au cas où $A$
est intègre et $B=A$, et il s'agit alors de montrer que la condition e)
entraîne la condition (7.2.3, c)). Avec les notations de (7.2.3, c)),
considérons donc un point
$y'\in f^{-1}(T)\cap\overline{\{x'\}}$, et posons $y=f(y')$ et
$C=\mathscr O_{X,y}$; comme $y\in T$, on a par hypothèse
$\dim(C)\geqslant2$, et par suite (avec les notations de e)),
$j_*(\mathscr O_U)$ est un $\mathscr O_Y$-Module cohérent. Or, posons
$Y'=X'\times_XY$; le morphisme $f:X'\to X$ étant plat, il en est de même de
$g=f_{(Y)}:Y'\to Y$. D'ailleurs l'espace sous-jacent à $Y'$ s'identifie à
$f^{-1}(Y)$ (I, 3.6.5), et comme $A$ est intègre et $f$ plat,
$\operatorname{Ass}(\mathscr O_{X'})$ est contenu dans la fibre du point
générique de $X$ (3.3.2); ce dernier étant contenu dans $Y$, on a
$x'\in Y'$. Soit $U'=g^{-1}(U)$, et soit $j'$ l'injection canonique
$U'\to Y'$; comme $j_*(\mathscr O_U)$ est un $\mathscr O_Y$-Module
cohérent et que $g$ est plat, il résulte de (5.9.4) que
$j'_*(\mathscr O_{U'})$ est un $\mathscr O_{Y'}$-Module cohérent. Or on a
$x'\in\operatorname{Ass}(\mathscr O_{U'})$, donc on conclut de (5.10.10)
que l'on a, dans $Y'$,

\oldpage[IV-2]{190}

$\operatorname{codim}(\overline{\{y'\}}\cap\overline{\{x'\}},
\overline{\{x'\}})\geqslant2$.
Comme $y'$ est arbitraire dans
$f^{-1}(T)\cap\overline{\{x'\}}$, on a bien
$\operatorname{codim}(f^{-1}(T)\cap\overline{\{x'\}},
\overline{\{x'\}})\geqslant2$ dans $X'$. C.Q.F.D.
\end{proof}

\begin{theorem}[7.2.5]
\label{IV.2.s7.2.5-fr}
Soit $A$ un anneau local noethérien. Les conditions suivantes sont
équivalentes:
\begin{enumerate}
\item[a)] Pour tout anneau quotient intègre $B$ de $A$, le complété
$\widehat B$ est strictement équidimensionnel (7.2.1).
\item[b)] $A$ est formellement caténaire (7.1.9) et les fibres du morphisme
canonique $\operatorname{Spec}(\widehat A)\to\operatorname{Spec}(A)$
vérifient $(\mathrm S_1)$ (autrement dit, n'ont pas de cycle premier associé
immergé).
\item[c)] $A$ est universellement caténaire (5.6.2) et pour tout anneau
quotient intègre $B$ de $A$, l'anneau $B^{(1)}$ est une $B$-algèbre finie
(cf. (7.2.4)).
\item[d)] $A$ est universellement caténaire et pour tout anneau quotient
intègre $B$ de $A$ et tout anneau local $C=B_{\mathfrak q}$ en un idéal
premier $\mathfrak q$ de $B$, tel que $\dim(C)\geqslant2$, l'anneau
$C^{(\omega)}$ est une $C$-algèbre finie.
\item[e)] $A$ est universellement caténaire et pour tout anneau quotient
intègre $B$ de $A$ et tout anneau local $C=B_{\mathfrak q}$ en un idéal
premier $\mathfrak q$ de $B$, tel que $\dim(C)\geqslant2$, l'anneau complété
$\widehat C$ est tel que $\operatorname{Spec}(\widehat C)$ n'ait pas de
cycle premier associé de dimension 1.
\end{enumerate}
De plus, lorsque ces conditions sont remplies, alors, pour tout anneau
quotient $B$ de $A$ qui est strictement équidimensionnel, le complété
$\widehat B$ est strictement équidimensionnel.
\end{theorem}

\begin{proof}
L'équivalence de c), d) et e) a été démontrée dans (7.2.4). Pour montrer
que a) et b) sont équivalentes, rappelons (7.1.9) que dire que $A$ est
formellement caténaire signifie que pour tout anneau quotient intègre $B$
de $A$, $\widehat B$ est équidimensionnel. D'autre part, toute fibre du
morphisme $\operatorname{Spec}(\widehat A)\to\operatorname{Spec}(A)$ en un
point $x\in\operatorname{Spec}(A)$ est la fibre du morphisme
$\operatorname{Spec}(\widehat B)\to\operatorname{Spec}(B)$ au point
générique $x$ de $\operatorname{Spec}(B)$, où
$B=A/\mathfrak J_x$; dire que cette fibre est sans cycle premier associé
immergé équivaut à dire que $\widehat B$ n'a pas d'idéaux premiers associés
immergés, en vertu de (3.3.3); ceci prouve l'équivalence de a) et b). Le
même raisonnement montre que si a) est vérifiée, alors, pour tout anneau
quotient $B$ de $A$ qui est sans idéaux premiers associés immergés, le
complété $\widehat B$ est aussi sans idéaux premiers associés immergés;
d'autre part, l'hypothèse que $A$ est formellement caténaire entraîne que
si $B$ est équidimensionnel, il en est de même de $\widehat B$ (7.1.8);
ceci établit la dernière assertion du théorème.

Montrons maintenant que a) implique c). La condition a) implique que $A$
est universellement caténaire (7.1.11); montrons d'autre part que pour tout
anneau quotient intègre $B$ de $A$, $B^{(1)}$ est alors une $B$-algèbre
finie. Compte tenu de (7.2.3) appliqué à $B$, il s'agit de montrer que si
$X=\operatorname{Spec}(B)$, si $T$ est une partie fermée de $X$, de
codimension $\geqslant2$, $X'=\operatorname{Spec}(\widehat B)$, et
$g:X'\to X$ le morphisme canonique, alors on a, pour tout
$x'\in\operatorname{Ass}(\mathscr O_{X'})$,
$\operatorname{codim}(g^{-1}(T)\cap\overline{\{x'\}},
\overline{\{x'\}})\geqslant2$.
Mais par hypothèse $X'$ n'a pas de cycle premier associé immergé, donc
$\inf(\operatorname{codim}(g^{-1}(T)\cap\overline{\{x'\}},
\overline{\{x'\}}))$, lorsque $x'$ parcourt
$\operatorname{Ass}(\mathscr O_{X'})$, n'est autre que
$\operatorname{codim}(g^{-1}(T),X')$.
Or, comme $g$ est un morphisme fidèlement plat, on a (6.1.4)
\[
\operatorname{codim}(g^{-1}(T),X')=\operatorname{codim}(T,X)\geqslant2.
\]

\oldpage[IV-2]{191}

Reste à prouver que c) entraîne a). Raisonnons par récurrence sur
$n=\dim(A)$, le théorème étant trivial pour $n=0$; on peut en outre se
borner au cas où $A=B$ est intègre. Procédons en deux étapes, $n$ étant
$\geqslant1$.

\medskip\noindent
I) \emph{Supposons d'abord que $A$ vérifie $(\mathrm S_2)$.} --- Posons
$X=\operatorname{Spec}(A)$, $A'=\widehat A$,
$X'=\operatorname{Spec}(A')$ et soit $u:X'\to X$ le morphisme canonique;
il s'agit de montrer que pour tout élément
$x'\in\operatorname{Ass}(\mathscr O_{X'})$, on a
$\dim(\overline{\{x'\}})=n$. Soit $f\neq0$ un élément de l'idéal maximal
de $A$, et posons $C=A/fA$; on sait (5.7.6) que $A/fA$ vérifie
$(\mathrm S_1)$; en outre, comme les idéaux premiers minimaux de $A$ parmi
ceux qui contiennent $f$ sont de hauteur 1 en vertu du Hauptidealsatz, et
que $C$ est caténaire, $C=A/fA$ est \emph{strictement équidimensionnel} et
$\dim(C)=n-1$. On a
$\widehat C=C\otimes_AA'=A'/fA'$, et $f$ est $A'$-régulier par platitude
(0$_{\mathrm I}$, 6.3.4); si l'on pose
$Y'=V(fA')=\operatorname{Spec}(\widehat C)$, alors, pour tout point maximal
$y'$ de $Y'\cap\overline{\{x'\}}$, on a
$y'\in\operatorname{Ass}(\mathscr O_{Y'})$ en vertu de (3.4.3); d'autre
part, on a $y'\neq x'$ car $u(x')$ est le point générique de $X$ (3.3.2)
et l'on a $u(y')\in V(fA)$. On en conclut (5.1.8) que
\begin{equation}
\dim(\overline{\{y'\}})=\dim(\overline{\{x'\}})-1.
\tag{7.2.5.1}
\label{IV.2.eq.7.2.5.1-fr}
\end{equation}

Mais l'anneau quotient $C=A/fA$ vérifie aussi l'hypothèse c) de l'énoncé
(5.6.1 et 5.11.7, (ii)); en vertu de l'hypothèse de récurrence, on a
$\dim(\overline{\{y'\}})=\dim(C)=n-1$; la conclusion
$\dim(\overline{\{x'\}})=n$ résulte donc dans ce cas de (7.2.5.1).

\medskip\noindent
II) \emph{Cas général.} --- Comme par hypothèse l'anneau $A^{(1)}$ est une
$A$-algèbre finie, il vérifie $(\mathrm S_2)$ en vertu de (5.10.17, (i));
il en est de même de chacun de ses anneaux locaux
$(A^{(1)})_{\mathfrak n}$ en un idéal maximal $\mathfrak n$ et de plus,
comme $A$ est universellement caténaire, les anneaux
$(A^{(1)})_{\mathfrak n}$ (en nombre fini) sont tous de \emph{dimension}
$n$ (5.6.10); en outre, ces anneaux vérifient l'hypothèse c) de l'énoncé
(5.6.1 et 5.11.7, (ii)). On sait que le complété de $A^{(1)}$, égal à
$\widehat A\otimes_AA^{(1)}$, est composé direct des complétés des anneaux
locaux $(A^{(1)})_{\mathfrak n}$. Posons
$X_1=\operatorname{Spec}(A^{(1)})$,
$X'_1=\operatorname{Spec}(\widehat{A^{(1)}})=X'\times_XX_1$; pour tout
$x'_1\in\operatorname{Ass}(\mathscr O_{X'_1})$, il résulte de ce qui
précède et du cas I) que l'on a $\dim(\overline{\{x'_1\}})=n$. Soit
$u_1=u_{(X_1)}:X'_1\to X_1$ le morphisme canonique; comme $A$ et $A^{(1)}$
ont même corps des fractions, l'image réciproque du point générique $x$ de
$X$ par la projection $X_1\to X$ se réduit au point générique $x_1$ de
$X_1$, l'image réciproque par la projection $X'_1\to X'$ de la fibre
$u^{-1}(x)$ est la fibre $u_1^{-1}(x_1)$ et cette projection induit un
isomorphisme du préschéma $u_1^{-1}(x_1)$ sur $u^{-1}(x)$. Cela étant, les
points de $\operatorname{Ass}(\mathscr O_{X'})$ (resp.
$\operatorname{Ass}(\mathscr O_{X'_1})$) sont les points génériques des
cycles premiers associés à $u^{-1}(x)$ (resp. $u_1^{-1}(x_1)$) (3.3.1).
Pour tout $x'\in\operatorname{Ass}(\mathscr O_{X'})$, il y a donc un
$x'_1\in\operatorname{Ass}(\mathscr O_{X'_1})$ au-dessus de $x'$, et si
$Z'$ (resp. $Z'_1$) est le sous-préschéma fermé réduit de $X'$ (resp.
$X'_1$) ayant $\overline{\{x'\}}$ (resp. $\overline{\{x'_1\}}$) pour
espace sous-jacent, la projection $Z'_1\to Z'$ est un morphisme fini et
surjectif; on en conclut (5.4.2) que
$\dim(\overline{\{x'\}})=\dim(\overline{\{x'_1\}})=n$. C.Q.F.D.
\end{proof}

\begin{definition}[7.2.6]
\label{IV.2.s7.2.6-fr}
Lorsqu'un anneau local noethérien $A$ vérifie les conditions équivalentes de
(7.2.5), on dit qu'il est \emph{strictement formellement caténaire}.
\end{definition}

\oldpage[IV-2]{192}

\begin{corollary}[7.2.7]
\label{IV.2.s7.2.7-fr}
Soit $A$ un anneau local noethérien tel qu'il existe un $A$-module $M$ de
type fini qui soit un $A$-module de Cohen--Macaulay et pour lequel
$\operatorname{Supp}(M)=\operatorname{Spec}(A)$; alors $A$ est strictement
formellement caténaire.
\end{corollary}

\begin{proof}
En effet, $A$ est formellement caténaire (7.1.5 et 7.1.9), et d'autre part
les fibres du morphisme canonique
$\operatorname{Spec}(\widehat A)\to\operatorname{Spec}(A)$ sont des
préschémas de Cohen--Macaulay (6.3.8), donc \emph{a fortiori} vérifient
$(\mathrm S_1)$.
\end{proof}

\begin{corollary}[7.2.8]
\label{IV.2.s7.2.8-fr}
Si $A$ est un anneau local noethérien strictement formellement caténaire,
tout anneau local $B$ qui est une $A$-algèbre essentiellement de type fini
est strictement formellement caténaire.
\end{corollary}

\begin{proof}
En effet, $B$ est formellement caténaire (7.1.11) et il résulte de
(7.4.4)\footnote{Le cor. (7.2.8) n'est pas utilisé pour démontrer (7.4.4).}
que les fibres de
$\operatorname{Spec}(\widehat B)\to\operatorname{Spec}(B)$ vérifient
$(\mathrm S_1)$, d'où la conclusion.
\end{proof}

\begin{corollary}[7.2.9]
\EGAReaderTarget{ega4.c04.cor.7-2-9}{IV.2.s7.2.9-fr}
Tout anneau local noethérien de dimension 1 est strictement formellement
caténaire.
\end{corollary}

\begin{proof}
En effet, tout anneau quotient intègre $B$ d'un tel anneau $A$ est de
dimension 0 ou 1, donc son complété est strictement équidimensionnel
(7.2.1.1).
\end{proof}

\begin{remarks}[7.2.10]
\label{IV.2.s7.2.10-fr}
\begin{enumerate}
\item[(i)] Il résulte de (7.2.7) et (7.2.8) que tout anneau quotient d'un
anneau local de Cohen--Macaulay est strictement formellement caténaire.

Un anneau local noethérien de dimension 2 vérifiant $(\mathrm S_2)$ est
strictement formellement caténaire, puisque c'est un anneau de
Cohen--Macaulay. Rappelons par contre qu'il y a des anneaux locaux
noethériens intègres de dimension 2 qui ne sont pas universellement
caténaires, ni \emph{a fortiori} strictement formellement caténaires
(5.6.11).
\item[(ii)] On ignore si un anneau formellement caténaire (7.1.9) est
strictement formellement caténaire; cela tient à ce qu'on ne sait pas si
pour un anneau local noethérien intègre $B$, $\widehat B$ est sans idéaux
premiers associés immergés (6.4.3). Nous ignorons également si, dans les
conditions équivalentes c), d), e) de (7.2.5), on peut remplacer
l'hypothèse que $A$ est universellement caténaire par l'hypothèse que $A$
est caténaire; on peut montrer qu'il en est ainsi lorsque $A$ est hensélien
(18.9.6).
\end{enumerate}
\end{remarks}

\subsection{Fibres formelles des anneaux locaux noethériens}
\label{subsection:IV.7.3-fr}

\begin{env}[7.3.1]
\label{IV.2.s7.3.1-fr}
Nous allons considérer dans ce numéro et les deux suivants des propriétés
$\boldsymbol P(Z,k)$ de la forme suivante :
\begin{quote}
« $Z$ est un préschéma localement noethérien sur un corps $k$,

et pour tout $z\in Z$, on a $\boldsymbol Q(\mathscr O_z,k)$ »
\end{quote}
où $\boldsymbol Q(A,k)$ est une propriété d'un anneau local noethérien $A$
qui est une $k$-algèbre; on suppose que si $k'$ est un corps isomorphe à
$k$, et si $A'$ est un anneau local noethérien qui est une $k'$-algèbre
di-isomorphe à $A$, les propriétés $\boldsymbol Q(A,k)$ et
$\boldsymbol Q(A',k')$ sont équivalentes.
\end{env}

\oldpage[IV-2]{193}

Si $X$, $Y$ sont deux préschémas localement noethériens, nous dirons qu'un
morphisme $f:X\to Y$ est un \emph{$\boldsymbol P$-morphisme} si :
\begin{enumerate}
\item[1\textsuperscript{o}] $f$ est \emph{plat};
\item[2\textsuperscript{o}] pour tout $y\in Y$, la propriété
$\boldsymbol P(f^{-1}(y),k(y))$ est vraie.
\end{enumerate}

\begin{lemma}[7.3.2]
\label{IV.2.s7.3.2-fr}
Soient $X$, $Y$ deux préschémas localement noethériens,
$f:X\to Y$ un morphisme. Les propriétés suivantes sont équivalentes :
\begin{enumerate}
\item[\emph{a)}] $f$ est un $\boldsymbol P$-morphisme.
\item[\emph{b)}] Pour tout $x\in X$, si l'on pose $y=f(x)$, alors $\mathscr O_x$ est
un $\mathscr O_y$-module plat et la propriété
$\boldsymbol Q(\mathscr O_x\otimes_{\mathscr O_y}k(y),k(y))$ est vraie.
\item[\emph{c)}] Pour tout $x\in X$, si l'on pose $y=f(x)$, le morphisme
$\operatorname{Spec}(\mathscr O_x)\to\operatorname{Spec}(\mathscr O_y)$
correspondant à l'homomorphisme $\mathscr O_y\to\mathscr O_x$ est un
$\boldsymbol P$-morphisme.
\item[\emph{c$'$)}] Pour tout point fermé $x\in X$, le morphisme
$\operatorname{Spec}(\mathscr O_x)\to\operatorname{Spec}(\mathscr O_y)$
est un $\boldsymbol P$-morphisme.
\end{enumerate}
\end{lemma}

\begin{proof}
L'équivalence de a) et b) résulte en effet des définitions; il en est de
même de l'équivalence de b) et c), compte tenu de (I, 2.4.2); enfin
l'équivalence de b) et c$'$) résulte de ce que pour tout $x\in X$,
$\overline{\{x\}}$ contient un point fermé (5.1.11).
\end{proof}

\begin{corollary}[7.3.3]
\label{IV.2.s7.3.3-fr}
Soient $A$, $B$ deux anneaux noethériens, $\varphi:A\to B$ un homomorphisme
tel que le morphisme correspondant
$\operatorname{Spec}(B)\to\operatorname{Spec}(A)$ soit un
$\boldsymbol P$-morphisme. Alors, pour toute partie multiplicative $S$ de
$A$, le morphisme
$\operatorname{Spec}(S^{-1}B)\to\operatorname{Spec}(S^{-1}A)$ est un
$\boldsymbol P$-morphisme.
\end{corollary}

\begin{proof}
Cela résulte aussitôt de (7.3.2) et de (I, 1.6.2).
\end{proof}

\begin{env}[7.3.4]
\label{IV.2.s7.3.4-fr}
Nous supposerons toujours dans ce qui suit que la propriété
$\boldsymbol Q$ est telle que les trois conditions suivantes soient
vérifiées :
\begin{enumerate}
\item[$(\mathrm P_{\mathrm I})$] \emph{(transitivité).} --- Si
$f:X\to Y$ est un morphisme régulier (6.8.1) et $g:Y\to Z$ un
$\boldsymbol P$-morphisme, alors $g\circ f$ est un
$\boldsymbol P$-morphisme.
\item[$(\mathrm P_{\mathrm{II}})$] \emph{(descente).} --- Si
$f:X\to Y$ et $g:Y\to Z$ sont des morphismes de préschémas localement
noethériens tels que $f$ soit fidèlement plat et que $g\circ f$ soit un
$\boldsymbol P$-morphisme, alors $g$ est un $\boldsymbol P$-morphisme.
\item[$(\mathrm P_{\mathrm{III}})$] --- Pour tout corps $k$, la propriété
$\boldsymbol P(\operatorname{Spec}(k),k)$ est vraie.
\end{enumerate}
\end{env}

\begin{remarks}[7.3.5]
\label{IV.2.s7.3.5-fr}
\begin{enumerate}
\item[\emph{i)}] Les conditions $(\mathrm P_{\mathrm I})$ et
$(\mathrm P_{\mathrm{III}})$ entraînent que \emph{tout morphisme régulier
est un $\boldsymbol P$-morphisme}.
\item[\emph{ii)}] Notons que les hypothèses de $(\mathrm P_{\mathrm I})$
(resp. $(\mathrm P_{\mathrm{II}})$) entraînent que $h=g\circ f$ (resp.
$g$) est plat (2.2.13); les hypothèses de $(\mathrm P_{\mathrm I})$ ou de
$(\mathrm P_{\mathrm{II}})$ entraînent d'autre part que pour tout
$z\in Z$, $f_z:h^{-1}(z)\to g^{-1}(z)$ est plat (2.1.4); comme, pour tout
$y\in g^{-1}(z)$, $f_z^{-1}(y)$ est isomorphe à $f^{-1}(y)$ (I, 3.6.4),
cela montre qu'il suffit, pour vérifier les conditions
$(\mathrm P_{\mathrm I})$ et $(\mathrm P_{\mathrm{II}})$, de le faire
seulement lorsque $Z$ est le \emph{spectre d'un corps}.
\item[\emph{iii)}] Dans certains cas, la propriété $\boldsymbol Q$ sera telle que la
condition suivante soit vérifiée :
\begin{enumerate}
\item[$(\mathrm P'_{\mathrm I})$] --- Si $f:X\to Y$ et $g:Y\to Z$ sont
deux $\boldsymbol P$-morphismes, alors $g\circ f$ est un
$\boldsymbol P$-morphisme.
\end{enumerate}
\end{enumerate}
\end{remarks}

\oldpage[IV-2]{194}

\begin{env}[7.3.6]
\label{IV.2.s7.3.6-fr}
Nous dirons que la propriété $\boldsymbol P$ est \emph{géométrique} si elle
satisfait (en outre de $(\mathrm P_{\mathrm I})$,
$(\mathrm P_{\mathrm{II}})$ et $(\mathrm P_{\mathrm{III}})$) à la
condition :
\begin{enumerate}
\item[$(\mathrm P_{\mathrm{IV}})$] \emph{(extension de type fini du corps
de base).} --- Si $\boldsymbol P(Z,k)$ est vraie, alors, pour toute
extension $k'$ de type fini de $k$,
$\boldsymbol P(Z\otimes_k k',k')$ est vraie.
\end{enumerate}
\end{env}

\begin{lemma}[7.3.7]
\label{IV.2.s7.3.7-fr}
Soient $f:X\to Y$ un $\boldsymbol P$-morphisme de préschémas localement
noethériens, $g:Y'\to Y$ un morphisme localement de type fini. Alors, si
$\boldsymbol P$ vérifie la condition $(\mathrm P_{\mathrm{IV}})$, le
morphisme $f'=f_{(Y')}:X\times_Y Y'\to Y'$ est un
$\boldsymbol P$-morphisme.
\end{lemma}

\begin{proof}
En effet, pour tout $y'\in Y'$, si l'on pose $y=g(y')$, $k(y')$ est une
extension de type fini de $k(y)$ (I, 6.4.11) et
$f'^{-1}(y')=f^{-1}(y)\otimes_{k(y)}k(y')$ (I, 3.6.4); il suffit donc
d'appliquer $(\mathrm P_{\mathrm{IV}})$.
\end{proof}

\begin{env}[7.3.8]
\label{IV.2.s7.3.8-fr}
\emph{Exemples.} --- Les propriétés $\boldsymbol P(Z,k)$ suivantes
vérifient les conditions $(\mathrm P_{\mathrm I})$,
$(\mathrm P_{\mathrm{II}})$ et $(\mathrm P_{\mathrm{III}})$ :
\begin{enumerate}
\item[(i)] (aussi notée $(\mathrm i_n)$) $Z$ est de coprofondeur
$\leqslant n$.
\item[(ii)] $Z$ est un préschéma de Cohen--Macaulay.
\item[(iii)] (aussi notée $(\mathrm{iii}_n)$) $Z$ vérifie $(\mathrm S_n)$.
\item[(iv)] $Z$ est régulier.
\item[(v)] (aussi notée $(\mathrm v_n)$) $Z$ vérifie $(\mathrm R_n)$.
\item[(vi)] $Z$ est réduit.
\item[(vii)] $Z$ est normal.
\end{enumerate}
Pour les propriétés (ii) à (vii), cela résulte de (6.6.1) qui, en fait,
démontre la condition plus forte $(\mathrm P'_{\mathrm I})$. Pour (i), la
propriété $(\mathrm P_{\mathrm{II}})$ résulte de (6.6.2); la propriété
$(\mathrm P_{\mathrm I})$ résulte, par le même raisonnement que dans
(6.6.1, (i)), de (6.3.2) et du fait qu'un préschéma régulier est de
coprofondeur 0.

En outre, il résulte de (6.7.1) que les propriétés (i), (ii) et (iii) sont
géométriques; en vertu de (6.7.8), il en est de même des suivantes :
\begin{enumerate}
\item[(iv$'$)] $Z$ est géométriquement régulier.
\item[(v$'$)] (aussi notée $(\mathrm v'_n)$) $Z$ possède la propriété
$(\mathrm R_n)$ géométrique.
\item[(vi$'$)] $Z$ est géométriquement réduit.
\item[(vii$'$)] $Z$ est géométriquement normal.
\end{enumerate}
\end{env}

\begin{remarks}[7.3.9]
\label{IV.2.s7.3.9-fr}
\begin{enumerate}
\item[\emph{i)}] Chacune des propriétés (iv$'$), $(\mathrm v'_n)$, (vi$'$),
(vii$'$) entraîne respectivement la propriété correspondante (iv),
$(\mathrm v_n)$, (vi), (vii). La propriété (iv) implique toutes les
propriétés de (i) à (vii), et la propriété (iv$'$) implique toutes les
propriétés énumérées dans (7.3.8). Rappelons aussi que $(\mathrm i_0)$ est
équivalente à (ii); la conjonction de $(\mathrm{iii}_1)$ et de
$(\mathrm v_0)$ (resp. de $(\mathrm{iii}_1)$ et de
$(\mathrm v'_0)$) équivaut à (vi) (resp. (vi$'$)) (5.8.5); enfin la
conjonction de $(\mathrm{iii}_2)$ et $(\mathrm v_1)$ (resp. de
$(\mathrm{iii}_2)$ et $(\mathrm v'_1)$) équivaut à (vii) (resp.
(vii$'$)) (5.8.6).
\item[\emph{ii)}] Dans tous les exemples de (7.3.8), la propriété
$\boldsymbol Q(\mathscr O_z,k)$ qui sert à définir $\boldsymbol P$ est
telle que, pour toute \emph{générisation} $z'$ de $z$ dans $Z$,
$\boldsymbol Q(\mathscr O_z,k)$ \emph{entraîne}
$\boldsymbol Q(\mathscr O_{z'},k)$ : en vertu de (2.3.4), il suffit de le
vérifier pour les propriétés (i) à (vii) (et même (i) à (v) d'après la
remarque (7.3.9, (i)) : or c'est inclus dans la définition pour (iii) et
(v) (5.7.2 et 5.8.2), cela résulte de (6.3.9) pour (i), et enfin de
(0, 16.5.10 et 17.3.2) pour (ii) et (iv).
\end{enumerate}
\end{remarks}

\oldpage[IV-2]{195}

\begin{env}[7.3.10]
\label{IV.2.s7.3.10-fr}
Nous dirons qu'une propriété $\boldsymbol P(Z,k)$ est du \emph{premier
type} si la propriété $\boldsymbol Q(A,k)$ qui sert à définir
$\boldsymbol P$ est une propriété $\boldsymbol R(A)$ \emph{ne faisant pas
intervenir $k$} et est vraie lorsque $A$ est un \emph{corps}; nous dirons
que $\boldsymbol P(Z,k)$ est du \emph{second type} si la propriété
$\boldsymbol Q(A,k)$ est de la forme
\begin{quote}
« Pour toute extension $k'$ de type fini de $k$ et tout point $z'$ de
$\operatorname{Spec}(A\otimes_k k')$ au-dessus du point fermé de
$\operatorname{Spec}(A)$, on a $\boldsymbol R(\mathscr O_{z'})$ »
\end{quote}
$\boldsymbol R(A)$ étant encore supposée vraie lorsque $A$ est un corps.
Il est clair que dans les exemples de (7.3.8), les propriétés (i) à (vii)
sont du premier type, les propriétés (iv$'$) à (vii$'$) du second type.

Cela étant, si l'on reprend les raisonnements de (6.6.1) et (6.8.3), on
constate que lorsque $\boldsymbol P$ est du premier ou du second type pour
une propriété $\boldsymbol R(A)$, \emph{les conditions
$(\mathrm P_{\mathrm I})$ et $(\mathrm P_{\mathrm{II}})$ sont
conséquences des conditions suivantes sur $\boldsymbol R$} :

$(\mathrm R_{\mathrm I})$ Soient $X$, $Y$ deux préschémas localement
noethériens, $f:X\to Y$ un morphisme \emph{régulier}; alors, pour tout
$x\in X$, la propriété $\boldsymbol R(\mathscr O_{f(x)})$ implique la
propriété $\boldsymbol R(\mathscr O_x)$.

$(\mathrm R_{\mathrm{II}})$ Soit $\varphi:A\to B$ un homomorphisme local
d'anneaux locaux noethériens faisant de $B$ un $A$-module \emph{plat};
alors $\boldsymbol R(B)$ implique $\boldsymbol R(A)$.

En outre, la propriété $(\mathrm P_{\mathrm{III}})$ résulte de ce que
$\boldsymbol R(k)$ est vraie par hypothèse pour tout corps $k$; enfin,
lorsque $\boldsymbol P$ est du second type, la condition
$(\mathrm P_{\mathrm{IV}})$ est conséquence de la définition de
$\boldsymbol P$ et de la transitivité des extensions de type fini.
\end{env}

\begin{remark}[7.3.11]
\label{IV.2.s7.3.11-fr}
Nous laissons au lecteur le soin de formuler la propriété $\boldsymbol R$
correspondant à chacun des exemples de (7.3.8), et de vérifier les
conditions $(\mathrm R_{\mathrm I})$ et $(\mathrm R_{\mathrm{II}})$ dans
chacun des cas, à l'aide des résultats du \S\,6. En fait, sauf pour
l'exemple (i) de (7.3.8), la propriété $\boldsymbol R$ vérifie dans les
autres cas la condition suivante :
\begin{enumerate}
\item[$(\mathrm R'_{\mathrm I})$] Soient $X$, $Y$ deux préschémas
localement noethériens, $f:X\to Y$ un $\boldsymbol P$-morphisme (pour la
propriété $\boldsymbol P$ du premier ou du second type définie par
$\boldsymbol R$); alors, pour tout $x\in X$, la propriété
$\boldsymbol R(\mathscr O_{f(x)})$ implique la propriété
$\boldsymbol R(\mathscr O_x)$.
\end{enumerate}
Les raisonnements de (6.6.1) et (6.8.3) prouvent que si
$\boldsymbol R$ vérifie les conditions $(\mathrm R'_{\mathrm I})$ et
$(\mathrm R_{\mathrm{II}})$, alors $\boldsymbol P$ vérifie les conditions
$(\mathrm P'_{\mathrm I})$ (7.3.4, (iii)) et
$(\mathrm P_{\mathrm{II}})$.
\end{remark}

\begin{proposition}[7.3.12]
\label{IV.2.s7.3.12-fr}
Soit $\boldsymbol P$ une propriété du premier ou du second type définie à
partir d'une propriété $\boldsymbol R$ vérifiant les conditions
$(\mathrm R_{\mathrm I})$ et $(\mathrm R_{\mathrm{II}})$ (resp.
$(\mathrm R'_{\mathrm I})$ et $(\mathrm R_{\mathrm{II}})$). Si pour tout
préschéma localement noethérien $X$, $U_{\boldsymbol R}(X)$ désigne
l'ensemble des $x\in X$ tels que $\boldsymbol R(\mathscr O_x)$ soit vraie,
alors, pour tout morphisme régulier (resp. tout
$\boldsymbol P$-morphisme) $f:X\to Y$ de préschémas localement
noethériens, on a
\begin{equation}
U_{\boldsymbol R}(X)=f^{-1}(U_{\boldsymbol R}(Y)).
\tag{7.3.12.1}
\label{IV.2.eq.7.3.12.1-fr}
\end{equation}
\end{proposition}

\begin{proof}
C'est une conséquence immédiate des définitions.
\end{proof}

\begin{env}[7.3.13]
\label{IV.2.s7.3.13-fr}
Étant donné un anneau semi-local noethérien $A$, nous appellerons
\emph{fibres formelles de $A$} les fibres du morphisme canonique
$f:\operatorname{Spec}(\widehat A)\to\operatorname{Spec}(A)$; pour tout
idéal premier $\mathfrak p$ de $A$, la fibre formelle en $\mathfrak p$ est
donc le schéma
$\operatorname{Spec}(\widehat A\otimes_A k(\mathfrak p))$; comme le
complété de l'anneau local $A/\mathfrak p$ est
$\widehat A/\mathfrak p\widehat A=(A/\mathfrak p)\otimes_A\widehat A$, on
voit que la fibre formelle de $A$ en $\mathfrak p$ est aussi \emph{la
fibre formelle de $A/\mathfrak p$ au point générique (0) de
$\operatorname{Spec}(A/\mathfrak p)$}.
\end{env}

\oldpage[IV-2]{196}

La propriété $\boldsymbol P$ étant définie comme dans (7.3.1), nous dirons
que \emph{les fibres formelles de $A$ ont la propriété $\boldsymbol P$},
ou que $A$ est un \emph{$\boldsymbol P$-anneau}, si, pour tout
$x\in\operatorname{Spec}(A)$,
$\boldsymbol P(f^{-1}(x),k(x))$ est vraie. Comme $f$ est \emph{plat}, il
revient au même de dire que $f$ est un $\boldsymbol P$-morphisme.

\begin{proposition}[7.3.14]
\label{IV.2.s7.3.14-fr}
Soient $A$ un anneau semi-local noethérien,
$\mathfrak m_i$ $(1\leqslant i\leqslant r)$ ses idéaux maximaux, et posons
$A_i=A_{\mathfrak m_i}$; pour que $A$ soit un $\boldsymbol P$-anneau, il
faut et il suffit que chacun des $A_i$ le soit.
\end{proposition}

\begin{proof}
En effet, $\widehat A$ est composé direct des $\widehat A_i$, donc la
fibre formelle de $A$ en un point $x\in\operatorname{Spec}(A)$ est la
\emph{somme} des fibres formelles en $x$ de ceux des $A_i$ tels que
$x\in\operatorname{Spec}(A_i)$.
\end{proof}

\begin{proposition}[7.3.15]
\label{IV.2.s7.3.15-fr}
Soit $A$ un anneau semi-local noethérien.
\begin{enumerate}
\item[\emph{i)}] Si $A$ est un $\boldsymbol P$-anneau, tout anneau quotient de $A$
est un $\boldsymbol P$-anneau.
\item[\emph{ii)}] Si de plus $\boldsymbol P$ satisfait à la condition
$(\mathrm P_{\mathrm{IV}})$ (7.3.6), alors toute $A$-algèbre finie est un
$\boldsymbol P$-anneau.
\end{enumerate}
\end{proposition}

\begin{proof}
Pour tout idéal $\mathfrak a$ de $A$,
$\widehat A\otimes_A(A/\mathfrak a)$ est le complété de
$A/\mathfrak a$, et les fibres formelles de $A/\mathfrak a$ sont donc les
fibres formelles de $A$ aux points de $V(\mathfrak a)$, d'où (i). D'autre
part, si $B$ est une $A$-algèbre finie, donc un anneau semi-local, on a
$\widehat B=B\otimes_A\widehat A$, et (ii) résulte donc de (7.3.7).
\end{proof}

\begin{proposition}[7.3.16]
\label{IV.2.s7.3.16-fr}
On suppose que la propriété $\boldsymbol P$ est du premier (resp. du
second) type, définie à partir d'une propriété $\boldsymbol R$ (7.3.10).
Lorsque $\boldsymbol P$ est du second type, on suppose en outre que la
relation
\begin{quote}
« pour toute extension finie $k'$ de $k$, et tout
$z'\in Z\otimes_k k'$, $\boldsymbol R(\mathscr O_{z'})$ est vraie »
\end{quote}
entraîne $\boldsymbol P(Z,k)$ (ce qui est vérifié pour les exemples
(iv$'$) à (vii$'$) de (7.3.8), en vertu de (6.7.7) et (4.6.1)).

Soit $A$ un anneau semi-local noethérien. Si la propriété
$\boldsymbol R$ vérifie les conditions $(\mathrm R_{\mathrm I})$ et
$(\mathrm R_{\mathrm{II}})$, les propriétés suivantes sont équivalentes :
\begin{enumerate}
\item[\emph{a)}] $A$ est un $\boldsymbol P$-anneau.
\item[\emph{b)}] Pour tout anneau quotient intègre $B$ de $A$ (resp. toute
$A$-algèbre finie intègre $B$) et tout idéal premier $\mathfrak q$ de
$\widehat B$ dont l'image réciproque dans $B$ est 0,
$\boldsymbol R(\widehat B_{\mathfrak q})$ est vraie.
\end{enumerate}
Si de plus $\boldsymbol R$ satisfait à la condition
$(\mathrm R'_{\mathrm I})$, les propriétés a) et b) sont aussi
équivalentes à :
\begin{enumerate}
\item[\emph{c)}] Pour tout anneau quotient $B$ de $A$ (resp. toute $A$-algèbre finie
$B$) si l'on pose $Y=\operatorname{Spec}(B)$,
$Y'=\operatorname{Spec}(\widehat B)$, et si $g:Y'\to Y$ est le morphisme
canonique, on a (avec les notations de (7.3.12))
\begin{equation}
U_{\boldsymbol R}(Y')=g^{-1}(U_{\boldsymbol R}(Y)).
\tag{7.3.16.1}
\label{IV.2.eq.7.3.16.1-fr}
\end{equation}
\end{enumerate}
\end{proposition}

\begin{proof}
L'équivalence de a) et b) est immédiate lorsque $\boldsymbol P$ est du
premier type : en effet, pour tout idéal premier $\mathfrak p$ de $A$, on
a vu que la fibre formelle de $B=A/\mathfrak p$ au point générique (0) de
$\operatorname{Spec}(B)$ n'est autre que la fibre formelle de $A$ au
point $\mathfrak p\in\operatorname{Spec}(A)$ (7.3.13).

Lorsque $\boldsymbol P$ est du second type, l'équivalence de a) et b)
résulte du lemme plus général suivant :
\end{proof}

\begin{lemma}[7.3.16.2]
\label{IV.2.s7.3.16.2-fr}
Soient $A$, $A'$ deux anneaux, $\varphi:A\to A'$ un homomorphisme,
$f:\operatorname{Spec}(A')\to\operatorname{Spec}(A)$ le morphisme
correspondant. Afin que pour tout $x\in\operatorname{Spec}(A)$, toute

\oldpage[IV-2]{197}

extension finie $k$ de $k(x)$ et tout point $z$ de
$f^{-1}(x)\otimes_{k(x)}k=Z$, la propriété
$\boldsymbol R(\mathscr O_{Z,z})$ soit vraie, il faut et il suffit que la
condition suivante soit satisfaite : pour toute $A$-algèbre intègre finie
$B$, si
$g:\operatorname{Spec}(A'\otimes_A B)\to\operatorname{Spec}(B)$ est le
morphisme déduit de $f$ par extension à $B$ de l'anneau de base, et si
$T=g^{-1}(\xi)$ est la fibre du point générique $\xi$ de
$\operatorname{Spec}(B)$, alors, pour tout $t\in T$,
$\boldsymbol R(\mathscr O_{T,t})$ est vraie.
\end{lemma}

\begin{proof}
La condition est trivialement nécessaire puisque si $x$ est le point de
$\operatorname{Spec}(A)$ au-dessus duquel est $\xi$, $k(\xi)$ est une
extension finie de $k(x)$. Inversement, considérons un point quelconque
$x\in\operatorname{Spec}(A)$ et soit $k$ une extension finie de $k(x)$.
Si l'on pose $\mathfrak p=\mathfrak J_x$, $k(x)$ est le corps des
fractions de $A/\mathfrak p$, et il y a une base de $k$ sur $k(x)$ formée
d'éléments entiers sur $A/\mathfrak p$; si $B$ est le sous-anneau de $k$
engendré par ces éléments, $B$ est une $A$-algèbre intègre finie et $k$
est le corps $k(\xi)$ au point générique $\xi$ de
$\operatorname{Spec}(B)$; comme $x$ est l'image de $\xi$ dans
$\operatorname{Spec}(A)$, la fibre $g^{-1}(\xi)$ n'est autre que
$f^{-1}(x)\otimes_{k(x)}k$, ce qui achève de démontrer le lemme.
\end{proof}

\begin{proof}[Fin de la démonstration de (7.3.16)]
Le fait que a) entraîne c) est une conséquence immédiate de (7.3.15) et
(7.3.12). D'autre part, si on spécialise c) au cas où $B$ est intègre, et
si $y$ désigne le point générique de $Y=\operatorname{Spec}(B)$, on a
$\mathscr O_{Y,y}=k(y)$, donc $y\in U_{\boldsymbol R}(Y)$, puisque
$\boldsymbol R(k)$ est vraie pour tout corps $k$; en exprimant que tout
point $y'\in g^{-1}(y)$ appartient à $U_{\boldsymbol R}(Y')$, on obtient
l'énoncé b), donc c) entraîne b). C.Q.F.D.
\end{proof}

\begin{proposition}[7.3.17]
\label{IV.2.s7.3.17-fr}
Supposons que la propriété $\boldsymbol R$ vérifie les conditions
$(\mathrm R'_{\mathrm I})$ et $(\mathrm R_{\mathrm{II}})$ et que
$\boldsymbol P$ soit la propriété du premier (resp. du second) type
définie à partir de $\boldsymbol R$ (7.3.10). Alors, si $A$ est un
$\boldsymbol P$-anneau, les propriétés $\boldsymbol R(A)$ et
$\boldsymbol R(\widehat A)$ sont équivalentes.
\end{proposition}

\begin{proof}
Cela résulte de (7.3.12.1) appliqué à
$Y=\operatorname{Spec}(A)$ et $X=\operatorname{Spec}(\widehat A)$.
\end{proof}

\begin{proposition}[7.3.18]
\label{IV.2.s7.3.18-fr}
Supposons que $\boldsymbol P$ soit une propriété du premier (resp. du
second) type définie à partir d'une propriété $\boldsymbol R$ vérifiant
les conditions $(\mathrm R'_{\mathrm I})$ et
$(\mathrm R_{\mathrm{II}})$ ainsi que la condition suivante :

$(\mathrm R_{\mathrm{III}})$ Pour tout anneau local noethérien complet
$C$, si l'on pose $Z=\operatorname{Spec}(C)$, l'ensemble
$U_{\boldsymbol R}(Z)$ (7.3.12) est ouvert dans $Z$.

Alors, si $A$ est un $\boldsymbol P$-anneau et si l'on pose
$X=\operatorname{Spec}(A)$, l'ensemble $U_{\boldsymbol R}(X)$ est ouvert
dans $X$.
\end{proposition}

\begin{proof}
En effet, si $X'=\operatorname{Spec}(\widehat A)$ et si $f:X'\to X$ est
le morphisme canonique, on a
$U_{\boldsymbol R}(X')=f^{-1}(U_{\boldsymbol R}(X))$ (7.3.12). Comme $f$
est fidèlement plat et quasi-compact et que par hypothèse
$U_{\boldsymbol R}(X')$ est ouvert dans $X'$, la conclusion découle de
(2.3.12).
\end{proof}

\begin{remarks}[7.3.19]
\EGAReaderTarget{ega4.c04.statement.7-3-19}{IV.2.s7.3.19-fr}
\begin{enumerate}
\item[\emph{i)}] La propriété $\boldsymbol R$ qui sert à définir $\boldsymbol P$
vérifie la condition $(\mathrm R_{\mathrm{III}})$ dans tous les exemples
énumérés dans (7.3.8). Pour (i), (ii) et (iii), cela résulte de (6.11.2)
et du théorème de Cohen (0, 19.8.8); lorsque $\boldsymbol P$ est une des
propriétés (iv), (iv$'$), (v), (v$'$), cela découle de (6.12.7) et
(6.12.9), et lorsque $\boldsymbol P$ est l'une des propriétés (vii),
(vii$'$), de (6.12.7) et (6.13.4); enfin, pour (vi) et (vi$'$)
l'assertion de (7.3.18) est triviale, étant vraie pour tout préschéma
localement noethérien.

\oldpage[IV-2]{198}

\item[\emph{ii)}] Nous avons déjà signalé (6.4.3) que lorsque $\boldsymbol P$ est
l'une des propriétés (ii) ou $(\mathrm{iii}_1)$ de (7.3.8), nous ignorons
si tout anneau local noethérien est un $\boldsymbol P$-anneau; rappelons
toutefois que lorsque $A$ est un \emph{anneau quotient d'un anneau de
Cohen--Macaulay}, les fibres formelles de $A$ sont des \emph{schémas de
Cohen--Macaulay} (6.3.8).
\item[\emph{iii)}] Le cas le plus intéressant de la notion de $\boldsymbol P$-anneau
est celui qui correspond à la propriété la plus forte (iv$'$) de
(7.3.8), c'est-à-dire les anneaux dont \emph{les fibres formelles sont
géométriquement régulières}. Les corps, et plus généralement les anneaux
locaux noethériens \emph{complets}, vérifient trivialement cette
propriété.
\item[\emph{iv)}] Soit $A$ un anneau local noethérien de \emph{dimension 1};
$\operatorname{Spec}(A)$ est alors formé du point fermé $a$,
correspondant à l'idéal maximal $\mathfrak m$, et des points maximaux
$b_i$ $(1\leqslant i\leqslant r)$ correspondant aux idéaux premiers
minimaux de $A$. On a $\dim(\widehat A)=1$ et l'idéal maximal de
$\widehat A$ est $\mathfrak m\widehat A$; la fibre formelle de $A$ au
point $a$ est donc $\operatorname{Spec}(k)$, où $k=A/\mathfrak m$ est le
corps résiduel de $A$; la fibre formelle en chacun des $b_i$ est le
spectre d'un anneau artinien dont les corps résiduels sont les corps
résiduels $L_{ij}$ de $\operatorname{Spec}(\widehat A)$ aux points
maximaux $b_{ij}$ au-dessus de $b_i$. Comme un anneau artinien est un
anneau de Cohen--Macaulay, on voit que $A$ est un
$\boldsymbol P$-anneau lorsque $\boldsymbol P$ est la propriété (ii) de
(7.3.8). En outre, comme un anneau artinien réduit est une somme directe
de corps, les propriétés suivantes sont équivalentes (6.7.7) :
\begin{enumerate}
\item[\emph{a)}] les fibres formelles de $A$ sont géométriquement réduites;
\item[\emph{b)}] les fibres formelles de $A$ sont géométriquement normales;
\item[\emph{c)}] les fibres formelles de $A$ sont géométriquement régulières;
\end{enumerate}
en outre, lorsque $A$ est \emph{réduit}, elles sont aussi équivalentes à :
\begin{enumerate}
\item[\emph{d)}] $\widehat A$ est réduit et $L_{ij}$ est une extension séparable de
$K_i$ pour tout couple $(i,j)$ (4.6.1).
\end{enumerate}
En particulier, si $A$ est un \emph{anneau de valuation discrète}, $K$
son corps des fractions, et si $\widehat K$ est le complété de $K$ pour
une valuation correspondant à $A$ (corps des fractions de $\widehat A$),
pour que les fibres formelles de $A$ soient géométriquement régulières,
il faut et il suffit que $\widehat K$ soit une \emph{extension séparable}
de $K$. Ce sera toujours le cas lorsque $K$ est de \emph{caractéristique
0}.
\end{enumerate}
\end{remarks}

\subsection{Permanence des propriétés des fibres formelles}
\label{subsection:IV.7.4-fr}

\begin{env}[7.4.1]
\label{IV.2.s7.4.1-fr}
Dans toute la suite du numéro, nous supposons que la propriété
$\boldsymbol P$ est de la forme définie dans (7.3.1), et \emph{vérifie
les conditions $(\mathrm P_{\mathrm I})$, $(\mathrm P_{\mathrm{II}})$ et
$(\mathrm P_{\mathrm{III}})$} de (7.3.4). En outre, nous supposons que la
propriété $\boldsymbol Q$ qui sert à définir $\boldsymbol P$ est telle
que, pour toute générisation $z'$ de $z$ dans $Z$,
$\boldsymbol Q(\mathscr O_z,k)$ entraîne
$\boldsymbol Q(\mathscr O_{z'},k)$.
\end{env}

\begin{lemma}[7.4.2]
\label{IV.2.s7.4.2-fr}
Soient $A$, $A'$ des anneaux locaux noethériens,
$\varphi:A\to A'$ un homomorphisme local tel que
$f={}^a\varphi:\operatorname{Spec}(A')\to\operatorname{Spec}(A)$ soit un
$\boldsymbol P$-morphisme. Alors, si les fibres formelles de $A'$ sont
géométriquement régulières, $A$ est un $\boldsymbol P$-anneau.
\end{lemma}

\oldpage[IV-2]{199}

\begin{proof}
Considérons l'homomorphisme complété
$\widehat\varphi:\widehat A\to\widehat{A'}$ et le morphisme correspondant
$\widehat f={}^a\widehat\varphi$; on a le diagramme commutatif
\[
\begin{tikzpicture}[>=latex]
\node (tl) at (0,1.5) {$\operatorname{Spec}(\widehat A)$};
\node (tr) at (4,1.5) {$\operatorname{Spec}(\widehat{A'})$};
\node (bl) at (0,0) {$\operatorname{Spec}(A)$};
\node (br) at (4,0) {$\operatorname{Spec}(A')$};
\draw[->] (tr) -- node[above] {$\widehat f$} (tl);
\draw[->] (tl) -- node[left] {$g$} (bl);
\draw[->] (tr) -- node[right] {$g'$} (br);
\draw[->] (br) -- node[below] {$f$} (bl);
\end{tikzpicture}
\]
où $g$ et $g'$ sont les morphismes canoniques. Comme par hypothèse $f$ est
un $\boldsymbol P$-morphisme et $g'$ un morphisme régulier, il résulte de
$(\mathrm P_{\mathrm I})$ que $f\circ g'=g\circ\widehat f$ est un
$\boldsymbol P$-morphisme. D'autre part, l'hypothèse que $f$ est un
$\boldsymbol P$-morphisme implique que $f$ est plat, donc il en est de même
de $\widehat f$ (Bourbaki, \emph{Alg. comm.}, chap. III, \S~5,
n\textsuperscript{o}~4, cor. de la prop. 3), qui est d'ailleurs un
homomorphisme local, donc fidèlement plat (0$_{\mathrm I}$, 6.6.2); il
résulte alors de $(\mathrm P_{\mathrm{II}})$ que $g$ est un
$\boldsymbol P$-morphisme.
\end{proof}

\begin{corollary}[7.4.3]
\label{IV.2.s7.4.3-fr}
\begin{enumerate}
\item[\emph{(i)}] Soient $A$ un $\boldsymbol P$-anneau local noethérien,
$A'=\widehat A$ son complété. Supposons que pour tout idéal premier
$\mathfrak p'$ de $A'$, les fibres formelles de $A'_{\mathfrak p'}$ soient
géométriquement régulières; alors, pour tout idéal premier $\mathfrak p$ de
$A$, $A_{\mathfrak p}$ est un $\boldsymbol P$-anneau.
\item[\emph{(ii)}] Supposons que $\boldsymbol P$ vérifie la condition
$(\mathrm P_{\mathrm{IV}})$ (7.3.6). Soient $A$ un $\boldsymbol P$-anneau
local noethérien, $B$ une $A$-algèbre locale essentiellement de type fini,
et telle que l'homomorphisme $A\to B$ soit local. Posons $A'=\widehat A$,
et soit $\mathfrak n'$ l'unique idéal premier de
$B'=B\otimes_A A'$ au-dessus de l'idéal maximal de $B$ et de l'idéal
maximal de $A'$. Si les fibres formelles de $B'_{\mathfrak n'}$ sont
géométriquement régulières, alors $B$ est un $\boldsymbol P$-anneau.
\end{enumerate}
\end{corollary}

\begin{proof}
\begin{enumerate}
\item[\emph{(i)}] Comme par hypothèse
$\operatorname{Spec}(A')\to\operatorname{Spec}(A)$ est un
$\boldsymbol P$-morphisme, il en est de même de
$\operatorname{Spec}(A'_{\mathfrak p'})\to
\operatorname{Spec}(A_{\mathfrak p})$ en vertu de (7.3.2, b)), pour tout
idéal premier $\mathfrak p$ de $A$ et tout idéal premier $\mathfrak p'$ de
$A'$ au-dessus de $\mathfrak p$. Il suffit alors d'appliquer le lemme
(7.4.2) à ce morphisme, en remarquant que le morphisme
$\operatorname{Spec}(A')\to\operatorname{Spec}(A)$ est surjectif.
\item[\emph{(ii)}] En vertu de (7.4.2), il suffit de prouver que le morphisme
$\operatorname{Spec}(B'_{\mathfrak n'})\to\operatorname{Spec}(B)$ est un
$\boldsymbol P$-morphisme. Or on a $B=C_{\mathfrak n}$, où $C$ est une
$A$-algèbre de type fini et $\mathfrak n$ un idéal premier de $C$ au-dessus
de l'idéal maximal de $A$. Si on pose $C'=C\otimes_A A'$, il résulte des
hypothèses et de (7.3.7) que
$\operatorname{Spec}(C')\to\operatorname{Spec}(C)$ est un
$\boldsymbol P$-morphisme; comme $B'_{\mathfrak n'}$ est un anneau local en
un idéal premier de $C'$ au-dessus de $\mathfrak n$, la conclusion résulte
de (7.3.2, b)).
\end{enumerate}
\end{proof}

\begin{theorem}[7.4.4]
\label{IV.2.s7.4.4-fr}
Les hypothèses sur $\boldsymbol P$ étant celles de (7.4.1), soit $A$ un
$\boldsymbol P$-anneau local noethérien.
\begin{enumerate}
\item[\emph{(i)}] Pour tout idéal premier $\mathfrak p$ de $A$,
$A_{\mathfrak p}$ est un $\boldsymbol P$-anneau.
\item[\emph{(ii)}] Supposons en outre que $\boldsymbol P$ vérifie la condition
$(\mathrm P_{\mathrm{IV}})$. Alors tout anneau local qui est une
$A$-algèbre essentiellement de type fini est un $\boldsymbol P$-anneau.
\end{enumerate}
\end{theorem}

\begin{proof}
\begin{enumerate}
\item[\emph{(i)}] Appliquant (7.4.3, (ii)), tout revient à voir que pour tout
idéal premier $\mathfrak p'$ de $A'=\widehat A$,
$A'_{\mathfrak p'}$ a ses fibres formelles géométriquement régulières. Or,
cela a été démontré dans (0, 22.3.3 et 22.5.8).

\oldpage[IV-2]{200}

\item[\emph{(ii)}] Soit $B$ une $A$-algèbre essentiellement de type fini qui
soit un anneau local; si $\mathfrak p$ est l'idéal premier de $A$ au-dessus
duquel est l'idéal maximal de $B$, $B$ est aussi une
$A_{\mathfrak p}$-algèbre essentiellement de type fini (1.3.8); en vertu de
(i), on peut donc supposer que $\mathfrak p$ est l'idéal maximal $\mathfrak m$
de $A$. On a donc $B=C_{\mathfrak q}$, où $C$ est une $A$-algèbre de type
fini, $\mathfrak q$ un idéal premier de $C$ au-dessus de $\mathfrak m$; soit
$\mathfrak n\supset\mathfrak q$ un idéal maximal de $C$ (nécessairement
au-dessus de $\mathfrak m$); si l'on pose $k=A/\mathfrak m$,
$C/\mathfrak mC$ est une $k$-algèbre de type fini, donc le corps résiduel
$k'$ de $C$ en l'idéal maximal $\mathfrak n$ est fini sur $k$ (I, 6.4.2);
en vertu de (i), il suffira de prouver que $C_{\mathfrak n}$ est un
$\boldsymbol P$-anneau, puisque $C_{\mathfrak q}$ est un anneau local en un
idéal premier de $C_{\mathfrak n}$. On est ainsi ramené à démontrer le
lemme suivant.
\end{enumerate}
\end{proof}

\begin{lemma}[7.4.4.1]
\label{IV.2.s7.4.4.1-fr}
Soient $A$ un $\boldsymbol P$-anneau local noethérien, $k$ son corps
résiduel, $C$ une $A$-algèbre de type fini, $B$ un anneau local en un idéal
premier $\mathfrak n$ de $C$, tel que : 1\textsuperscript{o}
l'homomorphisme $A\to B$ soit local; 2\textsuperscript{o} le corps résiduel
$k'$ de $B$ soit une extension finie de $k$. Si $\boldsymbol P$ satisfait à
$(\mathrm P_{\mathrm{IV}})$, $B$ est un $\boldsymbol P$-anneau.
\end{lemma}

\begin{proof}
Soit $(x_i)_{1\leqslant i\leqslant m}$ un système de générateurs de la
$A$-algèbre $C$; montrons d'abord que l'on peut raisonner par récurrence sur
$m$. Soit $C'$ la sous-algèbre de $C$ engendrée par
$x_1,\ldots,x_{m-1}$, et soit $\mathfrak n'=\mathfrak n\cap C'$.
L'homomorphisme $A\to C_{\mathfrak n}$ se factorise en
$A\to C'_{\mathfrak n'}\to C_{\mathfrak n}$ et il est clair que
$A\to C'_{\mathfrak n'}$ et $C'_{\mathfrak n'}\to C_{\mathfrak n}$ sont
des homomorphismes locaux; si $k''$ est le corps résiduel de
$C'_{\mathfrak n'}$, $k\to k'$ se factorise de même en
$k\to k''\to k'$, donc $k''$ est extension finie de $k$ et $k'$ extension
finie de $k''$. L'hypothèse de récurrence entraîne que
$C'_{\mathfrak n'}$ est un $\boldsymbol P$-anneau; en outre, si
$S'=C'-\mathfrak n'$, $C_{\mathfrak n}$ est un anneau local de
$S'^{-1}C$; comme $C=C'[x_m]$, on a
$S'^{-1}C=C'_{\mathfrak n'}[x_m/1]$ et l'hypothèse de récurrence montre
encore que $C_{\mathfrak n}$ est un $\boldsymbol P$-anneau. On est ainsi
ramené au cas où $C$ est une $A$-algèbre engendrée par un \emph{seul}
élément $t$.

Appliquons (7.4.3, (ii)); en posant $A'=\widehat A$,
$B'=B\otimes_A A'$ est un anneau local de $C\otimes_A A'$ en un idéal
premier au-dessus de $\mathfrak n$; comme $A$ et $A'$ ont même corps
résiduel $k$, le corps résiduel de $C\otimes_A A'$ en cet idéal premier est
égal à $k'$. D'ailleurs $C\otimes_A A'$ est une $A'$-algèbre engendrée par
un seul élément. Il résulte donc de (7.4.3, (ii)) qu'on peut se borner à
démontrer (7.4.4.1) lorsque $\boldsymbol P$ est la propriété (iv$'$) de
(7.3.8), que $A$ est \emph{complet} et $C$ engendrée par un seul élément
$t$.

Pour montrer que les fibres formelles de $B=C_{\mathfrak n}$ sont alors
géométriquement régulières, appliquons le critère (7.3.16, b)). Soit $B_1$
une $B$-algèbre finie \emph{intègre}, donc engendrée par un nombre fini
d'éléments entiers sur $B$. En multipliant ces éléments par un élément de
$S=C-\mathfrak n$, on peut supposer qu'ils sont entiers sur $C$, et on peut
donc écrire $B_1=S^{-1}C_1$, où $C_1$ est une sous-$C$-algèbre de $B_1$
engendrée par un nombre fini d'éléments entiers sur $C$, donc une
$C$-algèbre \emph{finie et intègre}. D'autre part, $B_1$ est un anneau
semi-local, et tout anneau local $B_2$ de $B_1$ en un idéal maximal est un
anneau local de $C_1$ en un idéal premier, tel que $A\to B_2$ soit un
homomorphisme local; en outre, le corps résiduel de $B_2$ est une extension
finie de $k'$, donc de $k$. On voit donc (compte tenu de (7.3.14) et de (i))
qu'on est ramené à la question suivante : soit $C$ un anneau noethérien
\emph{intègre}, contenant un sous-anneau $C_0$ qui est une $A$-algèbre
engendrée par un seul élément $t$ et telle que $C$ soit une $C_0$-algèbre
\emph{finie}; si $\mathfrak n$ est un idéal maximal de $C$

\oldpage[IV-2]{201}

au-dessus de l'idéal maximal de $A$, et $B=C_{\mathfrak n}$, il s'agit de
montrer que pour tout idéal premier $\mathfrak q$ de $\widehat B$, dont
l'image réciproque dans $B$ est 0, l'anneau
$\widehat B_{\mathfrak q}$ est \emph{régulier}. On peut d'ailleurs remplacer
$A$ par son image dans $C$, qui est un anneau local complet (comme quotient
de $A$) et intègre. La conclusion à prouver est alors conséquence du lemme
suivant, plus général en apparence.
\end{proof}

\begin{lemma}[7.4.4.2]
\label{IV.2.s7.4.4.2-fr}
Soient $A$ un anneau local noethérien intègre et complet, $k$ son corps
résiduel, $C$ un anneau intègre contenant $A$, tel qu'il existe $t\in C$
pour lequel $C$ soit une $A[t]$-algèbre finie. Soit $\mathfrak n$ un idéal
maximal de $C$ au-dessus de l'idéal maximal $\mathfrak m$ de $A$; posons
$B=C_{\mathfrak n}$, $X=\operatorname{Spec}(B)$,
$B'=\widehat B$, $X'=\operatorname{Spec}(\widehat B)$; si
$U=\operatorname{Reg}(X)$, $U'=\operatorname{Reg}(X')$, et si
$f:X'\to X$ est le morphisme canonique, on a alors
$f^{-1}(U)\subset U'$.
\end{lemma}

\begin{proof}
L'assertion à démontrer pour obtenir (7.4.4.1) se déduira de ce lemme en
observant que, puisque $B$ est intègre, le point générique de $X$ appartient
à $U$.

On notera que puisque $C$ est une $A$-algèbre de type fini, et $\mathfrak n$
un idéal maximal de $C$, le corps résiduel $k'$ de
$C_{\mathfrak n}=B$ (donc aussi de $B'$) est une extension finie de $k$
(I, 6.4.11 et 6.4.2).

Si l'on pose $Y=\operatorname{Spec}(C)$, il résulte de (6.12.8) que
$\operatorname{Reg}(Y)$ est ouvert dans $Y$; comme les anneaux locaux de
$X$ sont des anneaux locaux de $Y$ (I, 2.4.2), on a
$U=X\cap\operatorname{Reg}(Y)$, donc $U$ est ouvert dans $X$; d'autre part
(6.12.7) $U'$ est ouvert dans $X'$, donc $S'=X'-U'$ est fermé; par suite
$S'\cap f^{-1}(U)$ est localement fermé dans $X'$, et tout revient à
prouver que cet ensemble est vide. On sait (5.1.10) que dans le cas
contraire, il existerait dans $S'\cap f^{-1}(U)$ un idéal premier
$\mathfrak p'$ tel que $\dim(B'/\mathfrak p')\leqslant1$. Remarquons d'abord
que $\mathfrak p'$ ne peut être l'idéal maximal $\mathfrak mB'$ de $B'$,
où $\mathfrak m$ est l'idéal maximal de $B$. En effet, cela signifierait que
$B=B_{\mathfrak m}$ serait régulier, donc aussi $B'=\widehat B$ (0,
17.1.5), et on aurait par suite $\mathfrak mB'\in U'$ contrairement à
l'hypothèse. On doit donc avoir $\dim(B'/\mathfrak p')=1$. Posons
$\mathfrak p=B\cap\mathfrak p'$; par hypothèse $B_{\mathfrak p}$ est
régulier, mais $B'_{\mathfrak p'}$ ne l'est pas; comme
$B'_{\mathfrak p'}$ est un $B_{\mathfrak p}$-module plat, il résulte de
(6.5.2) que la fibre $Z$ de $f$ au point $\mathfrak p$ n'est pas régulière
au point $\mathfrak p'$. Montrons qu'on peut se ramener au cas où
$\mathfrak p=0$. En effet, dans le cas général, si l'on pose
$\mathfrak q=\mathfrak p\cap C$, $\mathfrak r=\mathfrak p\cap A$,
$C/\mathfrak q$ est une $(A/\mathfrak r)[\bar t]$-algèbre finie (où
$\bar t$ est la classe de $t$ mod. $\mathfrak q$); comme
$\mathfrak p=\mathfrak qB$, $B/\mathfrak p$ est égal à
$(C/\mathfrak q)_{\mathfrak n/\mathfrak q}$, et
$\mathfrak n/\mathfrak q$ est un idéal maximal de $C/\mathfrak q$
au-dessus de l'idéal maximal $\mathfrak m/\mathfrak r$ de
$A/\mathfrak r$; on voit donc que les hypothèses de (7.4.4.2) sont encore
remplies par $A/\mathfrak r$, $C/\mathfrak q$ et $B/\mathfrak p$, et comme
le complété de $B/\mathfrak p$ est $B'/\mathfrak pB'$, cela prouve notre
assertion. Supposons donc $\mathfrak p=0$, de sorte que $Z$ est la fibre
générique et que l'homomorphisme $B\to B'/\mathfrak p'$ est injectif.
Posons $V=B'/\mathfrak p'$, et distinguons deux cas :

\emph{I) $V$ est une $A$-algèbre finie.} --- Comme $B\subset V$, $B$ est
\emph{a fortiori} une $A$-algèbre finie, et comme $A$ est complet, il en est
de même de $B$ (Bourbaki, \emph{Alg. comm.}, chap. IV, \S~2,
n\textsuperscript{o}~5, cor. 3 de la prop. 9), d'où $B'=B$,
$\mathfrak p'=0$, donc $B'_{\mathfrak p'}$ est un corps, et par suite un
anneau régulier, contrairement à l'hypothèse.

\emph{II) $V$ n'est pas une $A$-algèbre finie.} --- Comme l'anneau local
$A$ est complet, cela implique que $V$ n'est pas une $A$-algèbre
quasi-finie (0$_{\mathrm I}$, 7.4.1 et 7.4.2); mais par hypothèse le corps
résiduel $k'$ de $V$ est extension finie du corps résiduel $k$ de $A$,

\oldpage[IV-2]{202}

donc (0$_{\mathrm I}$, 7.4.4) l'idéal $\mathfrak mV$ n'est pas un idéal de
définition de $V$. Comme $V$ est un anneau local intègre noethérien de
dimension 1, 0 est le seul idéal de $V$ qui ne soit pas un idéal de
définition, donc $\mathfrak mV=0$. Mais on a $A\subset V$ et $V$ est
intègre, d'où $\mathfrak m=0$ et $A=k$ est un corps. On en déduit tout
d'abord $\dim(C)\leqslant1$ (0, 16.1.5); mais comme $\dim(V)=1$, les
relations
\[
\dim(V)\leqslant\dim(B')=\dim(B)\leqslant\dim(C)
\]
montrent que cela entraîne
$\dim(C)=\dim(B)=\dim(B')=\dim(B'/\mathfrak p')=1$, et par suite
$\mathfrak p'$ est nécessairement un idéal premier minimal de $B'$. Nous
arriverons donc encore à une contradiction si nous prouvons que
$B'_{\mathfrak p'}$ est un corps, ou encore que l'anneau $B'$ est réduit.
Or, comme $C$ est une $k$-algèbre de type fini, la clôture intégrale $C_1$
de $C$ est une $C$-algèbre finie (Bourbaki, \emph{Alg. comm.}, chap. V,
\S~3, n\textsuperscript{o}~2, th. 2); si l'on pose
$S=C-\mathfrak n$, $B_1=S^{-1}C_1$ est la clôture intégrale de $B$, donc
$B_1$ est une $B$-algèbre finie, et par suite un anneau semi-local
noethérien, intègre et intégralement clos et de dimension 1 (0, 16.1.5); si
$\mathfrak m_j$ $(1\leqslant j\leqslant h)$ sont ses idéaux maximaux, les
$(B_1)_{\mathfrak m_j}$ sont donc des anneaux de valuation discrète (II,
7.1.6), et le complété $B'_1$ de $B_1$ est le composé direct des anneaux de
valuation discrète complétés des $(B_1)_{\mathfrak m_j}$ (Bourbaki,
\emph{Alg. comm.}, chap. III, \S~2, n\textsuperscript{o}~13, prop. 18);
$B'_1$ est donc réduit, et comme le complété $B'$ de $B$ est un
sous-anneau de $B'_1$ (Bourbaki, \emph{Alg. comm.}, chap. IV, \S~2,
n\textsuperscript{o}~5, cor. 3 de la prop. 9), c'est aussi un anneau réduit.
C.Q.F.D.
\end{proof}

\begin{corollary}[7.4.5]
\label{IV.2.s7.4.5-fr}
Supposons que la propriété $\boldsymbol P$ satisfasse aux conditions
$(\mathrm P_{\mathrm I})$, $(\mathrm P_{\mathrm{II}})$,
$(\mathrm P_{\mathrm{III}})$. Soit $A$ un anneau noethérien. Les conditions
suivantes sont équivalentes :
\begin{enumerate}
\item[\emph{a)}] Pour tout idéal premier $\mathfrak p$ de $A$,
$A_{\mathfrak p}$ est un $\boldsymbol P$-anneau.
\item[\emph{b)}] Pour tout idéal maximal $\mathfrak m$ de $A$,
$A_{\mathfrak m}$ est un $\boldsymbol P$-anneau.
\end{enumerate}
Si de plus $\boldsymbol P$ satisfait à $(\mathrm P_{\mathrm{IV}})$, les
propriétés a) et b) sont aussi équivalentes à :
\begin{enumerate}
\item[\emph{c)}] Pour toute $A$-algèbre de type fini $B$ et tout idéal
premier $\mathfrak q$ de $B$, $B_{\mathfrak q}$ est un
$\boldsymbol P$-anneau.
\end{enumerate}
\end{corollary}

\begin{proof}
L'équivalence de a) et b) résulte de (7.4.4, (ii)), et celle de a) et c) de
(7.4.4, (ii)).

Lorsque la condition a) de (7.4.5) est remplie, on dit que $A$ est un
$\boldsymbol P$-anneau; pour les anneaux semi-locaux noethériens, cette
définition coïncide avec celle de (7.3.13), lorsque les conditions
$(\mathrm P_{\mathrm I})$, $(\mathrm P_{\mathrm{II}})$ et
$(\mathrm P_{\mathrm{III}})$ sont satisfaites. L'anneau $\mathbf Z$ est un
$\boldsymbol P$-anneau (7.3.19, (iv)); tout anneau local complet est un
$\boldsymbol P$-anneau.
\end{proof}

\begin{proposition}[7.4.6]
\label{IV.2.s7.4.6-fr}
Supposons que la propriété $\boldsymbol P$ satisfasse aux conditions
$(\mathrm P_{\mathrm I})$, $(\mathrm P_{\mathrm{II}})$ et
$(\mathrm P_{\mathrm{III}})$. Soient $A$ un anneau noethérien,
$\mathfrak J$ un idéal de $A$, $\widehat A$ le séparé complété de $A$ pour
la topologie $\mathfrak J$-préadique. Alors, si $A$ est un
$\boldsymbol P$-anneau (7.4.5), le morphisme canonique
$\operatorname{Spec}(\widehat A)\to\operatorname{Spec}(A)$ est un
$\boldsymbol P$-morphisme.
\end{proposition}

\begin{proof}
En utilisant (7.3.2, c$'$)), il suffit de prouver que pour tout idéal
maximal $\mathfrak n$ de $\widehat A$, d'image réciproque $\mathfrak m$
dans $A$, le morphisme
$\operatorname{Spec}((\widehat A)_{\mathfrak n})\to
\operatorname{Spec}(A_{\mathfrak m})$ est un $\boldsymbol P$-morphisme.
On sait (Bourbaki, \emph{Alg. comm.}, chap. III, \S~3,
n\textsuperscript{o}~4, prop. 8) que l'homomorphisme canonique
$A_{\mathfrak m}\to(\widehat A)_{\mathfrak n}$ est injectif, que la
topologie $\mathfrak mA_{\mathfrak m}$-préadique sur $A_{\mathfrak m}$ est
induite par la topologie
$\mathfrak n(\widehat A)_{\mathfrak n}$-préadique et que
$A_{\mathfrak m}$ est dense dans $(\widehat A)_{\mathfrak n}$ pour cette
topologie, de sorte que le complété de $A_{\mathfrak m}$ pour la topologie
$\mathfrak mA_{\mathfrak m}$-préadique s'identifie à celui de
$(\widehat A)_{\mathfrak n}$ pour la topologie
$\mathfrak n(\widehat A)_{\mathfrak n}$-préadique. On a donc deux morphismes
\[
\operatorname{Spec}(\widehat{A_{\mathfrak m}})
\xrightarrow{f}\operatorname{Spec}((\widehat A)_{\mathfrak n})
\xrightarrow{g}\operatorname{Spec}(A_{\mathfrak m})
\]

\oldpage[IV-2]{203}

tels que $f$ soit fidèlement plat; comme par hypothèse $g\circ f$ est un
$\boldsymbol P$-morphisme, il en est de même de $g$ en vertu de
$(\mathrm P_{\mathrm{II}})$.
\end{proof}

\begin{corollary}[7.4.7]
\label{IV.2.s7.4.7-fr}
Supposons que $\boldsymbol P$ vérifie les conditions
$(\mathrm P_{\mathrm I})$, $(\mathrm P'_{\mathrm I})$,
$(\mathrm P_{\mathrm{II}})$, $(\mathrm P_{\mathrm{III}})$ et
$(\mathrm P_{\mathrm{IV}})$. Alors, si $A$ est un
$\boldsymbol P$-anneau (7.4.5), le morphisme canonique
$\operatorname{Spec}(A[[T_1,\ldots,T_r]])\to\operatorname{Spec}(A)$ est un
$\boldsymbol P$-morphisme. En particulier, si de plus $A$ est intègre, $K$
son corps des fractions, et si $\mathfrak p$ est un idéal premier de
$B=A[[T_1,\ldots,T_r]]$ tel que $\mathfrak p\cap A=0$, alors la propriété
$\boldsymbol P(B_{\mathfrak p},K)$ est vraie.
\end{corollary}

\begin{proof}
En effet, le morphisme canonique
$\operatorname{Spec}(B)\to\operatorname{Spec}(A)$ se factorise en
\[
\operatorname{Spec}(A[[T_1,\ldots,T_r]])
\xrightarrow{f}\operatorname{Spec}(A[T_1,\ldots,T_r])
\xrightarrow{g}\operatorname{Spec}(A).
\]
Il est clair que le morphisme $g$ est régulier (0, 17.3.7); en vertu de
(7.4.5), $A[T_1,\ldots,T_r]$ est un $\boldsymbol P$-anneau, donc il résulte
de (7.4.6) que $f$ est un $\boldsymbol P$-morphisme; comme $g$ est régulier,
donc est aussi un $\boldsymbol P$-morphisme (7.3.5, (i)), il en est de même
de $g\circ f$ en vertu de $(\mathrm P'_{\mathrm I})$.

On notera que la conclusion est encore valable si au lieu de supposer que
$\boldsymbol P$ vérifie $(\mathrm P'_{\mathrm I})$ et que tout morphisme
régulier est un $\boldsymbol P$-morphisme, on suppose seulement qu'un
morphisme composé $g\circ f$ est un $\boldsymbol P$-morphisme lorsque $g$
est régulier et $f$ un $\boldsymbol P$-morphisme (condition symétrique en
quelque sorte de $(\mathrm P'_{\mathrm I})$).
\end{proof}

\begin{remark}[7.4.8]
\label{IV.2.s7.4.8-fr}
Les résultats précédents posent les problèmes suivants :
\begin{enumerate}
\item[\emph{A)}] Soit $A$ un anneau de Zariski \emph{complet}, et soit
$\mathfrak J$ un idéal de définition de $A$; si l'anneau $A/\mathfrak J$
est un $\boldsymbol P$-anneau, en est-il de même de $A$? Il en résulterait
que pour tout $\boldsymbol P$-anneau noethérien $A$ et tout idéal
$\mathfrak J$ de $A$, le séparé complété $\widehat A$ pour la topologie
$\mathfrak J$-préadique serait aussi un $\boldsymbol P$-anneau.
\item[\emph{B)}] Soit $k$ un corps valué complet non discret; on appelle
encore anneau des \emph{séries formelles restreintes}
$k\{T_1,\ldots,T_n\}$ le sous-anneau de l'anneau de séries formelles
$k[[T_1,\ldots,T_n]]$ formées des séries dont les coefficients \emph{tendent
vers 0}. Un tel anneau est-il un $\boldsymbol P$-anneau?
\item[\emph{C)}] Si $A$ est un $\boldsymbol P$-anneau linéairement
topologisé, $S$ une partie multiplicative de $A$, les anneaux
$A\{S^{-1}\}$ (0$_{\mathrm I}$, 7.6.1) et $A_{\{S\}}$
(0$_{\mathrm I}$, 7.6.15) sont-ils des $\boldsymbol P$-anneaux?
\end{enumerate}
\end{remark}

\subsection{\texorpdfstring{Un critère pour les $\boldsymbol P$-morphismes}{Un critère pour les P-morphismes}}
\label{subsection:IV.7.5-fr}

\begin{env}[7.5.0]
\label{IV.2.s7.5.0-fr}
Ce numéro ne sera pas utilisé dans la suite du chapitre IV, et peut donc être
omis en première lecture. Nous verrons d'ailleurs plus loin (7.9.8) que les
résultats du présent numéro peuvent être considérablement améliorés quand on
dispose de la « résolution des singularités ».

Dans la suite de ce numéro, nous considérerons une propriété
$\boldsymbol R(A)$, et nous désignerons par $\boldsymbol P(Z,k)$ la propriété
suivante :
\begin{quote}
« $Z$ est un préschéma localement noethérien sur un corps $k$, et, pour toute
extension finie $k'$ de $k$, si l'on pose $Z'=Z\otimes_k k'$, la propriété
$\boldsymbol R(\mathscr O_{Z',z'})$ est vraie pour tout $z'\in Z'$. »
\end{quote}
\end{env}

\begin{theorem}[7.5.1]
\label{IV.2.s7.5.1-fr}
Soient $A$, $B$ deux anneaux locaux noethériens complets, $\mathfrak m$
l'idéal maximal de $A$, $k=A/\mathfrak m$ son corps résiduel,
$\varphi:A\to B$ un homomorphisme local tel que :
\begin{enumerate}
\item[\emph{(i)}] Le corps résiduel de $B$ est une extension finie de $k$.
\item[\emph{(ii)}] $B$ est un $A$-module plat.
\end{enumerate}
Soit d'autre part $\boldsymbol R(C)$ une propriété vérifiant la condition
$(\mathrm R_{\mathrm{III}})$ (7.3.18) et la condition suivante :
\begin{quote}
$(\mathrm R_{\mathrm{IV}})$ Pour tout anneau local $C$ en un idéal premier
d'un anneau local noethérien complet et tout élément $C$-régulier $t$ dans
l'idéal maximal de $C$, la propriété $\boldsymbol R(C/tC)$ implique
$\boldsymbol R(C)$.
\end{quote}

\oldpage[IV-2]{204}

Cela étant, soit $f={}^a\varphi:\operatorname{Spec}(B)\to
\operatorname{Spec}(A)$, et supposons que la propriété
$\boldsymbol P(\operatorname{Spec}(B\otimes_A k),k)$ soit vraie. Alors, pour
tout $x\in\operatorname{Spec}(A)$ la propriété
$\boldsymbol P(f^{-1}(x),k(x))$ est vraie.
\end{theorem}

\begin{proof}
Autrement dit, la propriété $\boldsymbol P$ pour la fibre du point fermé de
$\operatorname{Spec}(A)$ entraîne cette même propriété pour toutes les
fibres du morphisme $f$ (c'est-à-dire que $f$ est un
$\boldsymbol P$-morphisme (7.3.1)). Nous procéderons en plusieurs étapes :

\emph{I) Réduction à l'étude des anneaux locaux de la fibre générique.} ---
Appliquons le lemme (7.3.16.2) : il suffit de voir que pour toute
$A$-algèbre finie et intègre $A'$, de corps des fractions $K'$, si l'on
pose $B'=B\otimes_A A'$, tous les anneaux locaux de
$\operatorname{Spec}(B'\otimes_{A'}K')$ vérifient la propriété
$\boldsymbol R$. L'anneau $A'$ (resp. $B'$) est semi-local complet ($B'$
étant une $B$-algèbre finie), donc produit d'anneaux locaux complets; comme
$A'$ est supposé intègre, il est local; tout idéal maximal $\mathfrak n'$ de
$B'$ est au-dessus de l'idéal maximal $\mathfrak n$ de $B$, donc au-dessus
de $\mathfrak m$, et par suite son image réciproque dans $A'$ est l'idéal
maximal $\mathfrak m'$ de $A'$. Comme tout anneau local de
$\operatorname{Spec}(B'\otimes_{A'}K')$ est un anneau local de
$\operatorname{Spec}(B')$, donc de l'un des
$\operatorname{Spec}(B'_{\mathfrak n'})$ (au-dessus du point générique de
$A'$), on voit qu'il suffira de prouver que les anneaux locaux de
\[
\operatorname{Spec}(B'_{\mathfrak n'}\otimes_{A'}K')
\]
possèdent la propriété $\boldsymbol R$. Or $A'$ et
$B'_{\mathfrak n'}$ sont des anneaux locaux noethériens complets; le corps
résiduel de $B'_{\mathfrak n'}$ est extension finie de celui de $B$, donc
aussi de $k$, et \emph{a fortiori} du corps résiduel $k'$ de $A'$; cette
remarque et (2.1.4) montrent que $A'$ et $B'_{\mathfrak n'}$ vérifient les
conditions (i) et (ii) de l'énoncé. D'autre part, si $k''$ est une extension
finie de $k'$, $k''$ est aussi extension finie de $k$, et les anneaux locaux
de
\[
\operatorname{Spec}(B'_{\mathfrak n'}\otimes_{A'}k'')
\]
sont aussi des anneaux locaux de $\operatorname{Spec}(B\otimes_A k'')$;
donc l'hypothèse que
$\boldsymbol P(\operatorname{Spec}(B\otimes_A k),k)$ est vraie entraîne que
\[
\boldsymbol P(\operatorname{Spec}(B'_{\mathfrak n'}\otimes_{A'}k'),k')
\]
est vraie.

On est ainsi ramené au cas où $A$ est de plus supposé intègre, de corps des
fractions $K$, et à prouver que les anneaux locaux de
$\operatorname{Spec}(B\otimes_A K)$ possèdent la propriété
$\boldsymbol R$.

\emph{II) Cas où $\dim(A)=1$.} --- Soit $A'$ la clôture intégrale de $A$;
on sait (0, 23.1.6) que $A'$ est une $A$-algèbre finie et un anneau local
complet; si l'on pose $B'=B\otimes_A A'$, on a
$B\otimes_A K=B'\otimes_{A'}K$. Le même raisonnement que dans I) montre
qu'il suffit, pour tout idéal maximal $\mathfrak n'$ de $B'$, de prouver que
les anneaux locaux de
$B'_{\mathfrak n'}\otimes_{A'}K$ vérifient $\boldsymbol R$; en outre, ce
raisonnement montre aussi que $A'$ et $B'_{\mathfrak n'}$ vérifient les
hypothèses (i) et (ii) de l'énoncé et que
$\boldsymbol P(\operatorname{Spec}(B'_{\mathfrak n'}\otimes_{A'}k'),k')$
est vraie ($k'$ étant le corps résiduel de $A'$). En outre, comme
$\dim(A')=1$ (0, 16.1.5) et que $A'$ est intègre et intégralement clos,
c'est un anneau de valuation discrète complet. On peut donc se borner au
cas où $A$ est déjà un anneau de valuation discrète complet; si alors $u$
est une uniformisante de $A$, le fait que $u$ soit $A$-régulier et que $B$
soit un $A$-module plat entraîne que $t=\varphi(u)$ est un élément
$B$-régulier, qui appartient à l'idéal maximal de $B$. L'hypothèse que
$\boldsymbol P(\operatorname{Spec}(B\otimes_A k),k)$ est vraie entraîne en
particulier que $\boldsymbol R(B/tB)$ est vraie; en vertu de
$(\mathrm R_{\mathrm{IV}})$, $\boldsymbol R(B)$ est donc vraie. Autrement
dit, $U_{\boldsymbol R}(\operatorname{Spec}(B))$ contient le point fermé de
$\operatorname{Spec}(B)$. Mais comme
$U_{\boldsymbol R}(\operatorname{Spec}(B))$ est ouvert en vertu de
$(\mathrm R_{\mathrm{III}})$ et que $\operatorname{Spec}(B)$ est le seul
ensemble ouvert de $\operatorname{Spec}(B)$ contenant le point fermé, tous
les anneaux locaux de $\operatorname{Spec}(B)$ possèdent la propriété
$\boldsymbol R$, et en particulier ceux de
$\operatorname{Spec}(B\otimes_A K)$.

\emph{III) Cas général.} --- Le cas $\dim(A)=0$ est trivial, puisqu'alors
$A=k$; on peut donc se borner au cas où $\dim(A)\geqslant1$. On sait
(6.12.7) que $\operatorname{Spec}(A)$ contient un ouvert non vide $V$ dont
tous les points sont réguliers; comme $\dim(A)\geqslant1$, l'intersection
$V'$ de $V$ et du complémentaire du point fermé de
$\operatorname{Spec}(A)$ est un ouvert non vide, contenant donc le point
générique de $\operatorname{Spec}(A)$. Si nous prouvons que
\[
f^{-1}(V')\subset U_{\boldsymbol R}(\operatorname{Spec}(B)),
\]
la proposition sera démontrée \emph{a fortiori}. Autrement dit, il suffit de
voir que l'ensemble $Z$, intersection de $f^{-1}(V')$ et du complémentaire
de $U_{\boldsymbol R}(\operatorname{Spec}(B))$, est vide. Raisonnons par
l'absurde : comme, en vertu de $(\mathrm R_{\mathrm{III}})$,
$U_{\boldsymbol R}(\operatorname{Spec}(B))$ est ouvert dans
$\operatorname{Spec}(B)$, $Z$ est localement fermé dans
$\operatorname{Spec}(B)$; s'il n'était pas vide, il contiendrait un point
$x$ tel que $\dim(\overline{\{x\}})\leqslant1$ (5.1.10); comme
$f(x)\in V'$ est distinct du point fermé de $\operatorname{Spec}(A)$, $x$
est distinct du point fermé de $\operatorname{Spec}(B)$, autrement dit
$\dim(\overline{\{x\}})=1$. Nous allons prouver que cela est impossible,
autrement dit que, si $\dim(\overline{\{x\}})=1$ et $f(x)\in V'$ alors on a
$\boldsymbol R(\mathscr O_x)$. En d'autres termes, il s'agit de voir que,
si $\mathfrak q$ est un idéal premier de $B$ tel que
$\dim(B/\mathfrak q)=1$ et si $\mathfrak p=\varphi^{-1}(\mathfrak q)$ est
distinct de $\mathfrak m$ et tel que $A_{\mathfrak p}$ soit régulier, alors
on a $\boldsymbol R(B_{\mathfrak q})$. Comme
$\mathfrak p\ne\mathfrak m$, $\mathfrak m(B/\mathfrak q)$ n'est pas réduit
à 0, donc est un idéal de définition de l'anneau local intègre
$B/\mathfrak q$ de dimension 1; d'autre part, en vertu de (i), le corps
résiduel de $B/\mathfrak q$ est une extension finie du corps résiduel $k$
de $A/\mathfrak p$, donc $B/\mathfrak q$ est une
$(A/\mathfrak p)$-algèbre quasi-finie (0$_{\mathrm I}$, 7.4.4). Mais
$A/\mathfrak p$ est complet et $B/\mathfrak q$ est séparé pour la topologie
$\mathfrak m$-préadique (qui est identique à sa topologie d'anneau local);
donc (0$_{\mathrm I}$, 7.4.1), $B/\mathfrak q$ est une
$(A/\mathfrak p)$-algèbre finie. En outre, par définition,
l'homomorphisme $A/\mathfrak p\to B/\mathfrak q$ est injectif, donc (0,
16.1.5) on a
\[
\dim(A/\mathfrak p)=\dim(B/\mathfrak q)=1.
\]
On peut alors appliquer aux anneaux $A/\mathfrak p$ et
$B/\mathfrak pB$ le résultat de II), car les corps résiduels de ces anneaux
locaux sont respectivement ceux de $A$ et de $B$, et $B/\mathfrak pB$ est
un $(A/\mathfrak p)$-module plat; en outre, on a
\[
(B/\mathfrak pB)\otimes_{A/\mathfrak p}k=B\otimes_A k=B/\mathfrak mB.
\]
Par suite, les anneaux locaux de
\[
\operatorname{Spec}((B/\mathfrak pB)
\otimes_{A/\mathfrak p}k(\mathfrak p))
\]
vérifient $\boldsymbol R$. Or, $B_{\mathfrak q}/\mathfrak pB_{\mathfrak q}$
est un des anneaux locaux de
$\operatorname{Spec}((B/\mathfrak pB)
\otimes_{A/\mathfrak p}k(\mathfrak p))$. En outre,
$B_{\mathfrak q}$ est un $A_{\mathfrak p}$-module plat, et
$A_{\mathfrak p}$ est régulier. Or on a le lemme suivant.

\oldpage[IV-2]{205}

\begin{lemma}[7.5.1.1]
\label{IV.2.s7.5.1.1-fr}
Soit $\mathcal C$ une sous-catégorie pleine de la catégorie des anneaux
locaux noethériens, telle que tout anneau quotient d'un anneau de
$\mathcal C$ appartienne encore à $\mathcal C$. Soit $\boldsymbol R(C)$ une
propriété telle que si $C\in\mathcal C$ et si $t$ est un élément régulier de
l'idéal maximal de $C$ tel que $\boldsymbol R(C/tC)$ soit vraie, alors
$\boldsymbol R(C)$ est vraie.

Cela étant, soient $C$ un anneau local régulier, $k$ son corps résiduel,
$D$ un anneau local appartenant à $\mathcal C$,
$\varphi:C\to D$ un homomorphisme local faisant de $D$ un $C$-module plat.
Alors, si $\boldsymbol R(D\otimes_C k)$ est vraie, il en est de même de
$\boldsymbol R(D)$.
\end{lemma}

\begin{proof}
Raisonnons par récurrence sur $n=\dim(C)$, le lemme étant vrai par hypothèse
si $n=0$, puisqu'alors $C=k$. Soit $t$ un élément de l'idéal maximal
$\mathfrak m$ de $C$ n'appartenant pas à $\mathfrak m^2$; on sait alors que
$C/tC$ est régulier et que $\dim(C/tC)=n-1$ ((0, 17.1.8) et (0, 16.3.4)).
Comme
\[
(D/tD)\otimes_{C/tC}k=D\otimes_C k,
\]
et que $D/tD\in\mathcal C$, l'hypothèse de récurrence montre que
$\boldsymbol R(D/tD)$ est vraie; en outre, $t$ est $C$-régulier, donc aussi
$D$-régulier par platitude (0$_{\mathrm I}$, 6.3.4); l'hypothèse faite sur
$\boldsymbol R$ montre donc que $\boldsymbol R(D)$ est vraie.
\end{proof}

Pour terminer la démonstration de (7.5.1), il suffit ici d'appliquer le
lemme (7.5.1.1) en prenant pour $\mathcal C$ la catégorie des anneaux locaux
des spectres d'anneaux locaux complets, tenant compte de l'hypothèse
$(\mathrm R_{\mathrm{IV}})$, et prenant pour $C$ l'anneau
$A_{\mathfrak p}$, pour $D$ l'anneau $B_{\mathfrak q}$.
\end{proof}

\begin{corollary}[7.5.2]
\label{IV.2.s7.5.2-fr}
Soient $A$ un anneau local noethérien, $\mathfrak m$ l'idéal maximal de $A$,
$k=A/\mathfrak m$ son corps résiduel, $B$ un anneau local noethérien,
$\varphi:A\to B$ un homomorphisme local tel que :
\begin{enumerate}
\item[\emph{(i)}] Le corps résiduel de $B$ est une extension finie de $k$.
\item[\emph{(ii)}] $B$ est un $A$-module plat.
\end{enumerate}
Soit d'autre part $\boldsymbol R(C)$ une propriété vérifiant les conditions
$(\mathrm R_{\mathrm{II}})$ (7.3.10), $(\mathrm R_{\mathrm{III}})$
(7.3.18) et $(\mathrm R_{\mathrm{IV}})$ (7.5.1). Enfin, soit
$\boldsymbol R'(C)$ une propriété vérifiant la condition suivante :
\begin{quote}
$(\mathrm R_{\mathrm V})$ Si $C$, $D$ sont deux anneaux locaux noethériens,
$\psi:C\to D$ un homomorphisme local tel que
\[
g={}^a\psi:\operatorname{Spec}(D)\to\operatorname{Spec}(C)
\]
soit un $\boldsymbol P$-morphisme, alors la propriété
$\boldsymbol R'(C)$ implique $\boldsymbol R(D)$.
\end{quote}
Supposons que le morphisme canonique
$\operatorname{Spec}(\widehat A)\to\operatorname{Spec}(A)$ soit un
$\boldsymbol P'$-morphisme, où $\boldsymbol P'$ se définit à partir de
$\boldsymbol R'$ comme $\boldsymbol P$ à partir de $\boldsymbol R$ dans
(7.5.0).

Cela étant, soit
$f={}^a\varphi:\operatorname{Spec}(B)\to\operatorname{Spec}(A)$, et
supposons que la propriété
\[
\boldsymbol P(\operatorname{Spec}(\widehat B
\otimes_{\widehat A}k),k)
\]
soit vraie. Alors, pour tout $x\in\operatorname{Spec}(A)$, la propriété
$\boldsymbol P(f^{-1}(x),k(x))$ est vraie (autrement dit, $f$ est un
$\boldsymbol P$-morphisme).
\end{corollary}

\begin{proof}
Procédons encore en deux étapes :

\emph{I) Réduction à l'étude des anneaux locaux de la fibre générique.} ---
Appliquons encore le lemme (7.3.16.2); la seule différence avec le
raisonnement de I) dans (7.5.1) est qu'ici l'anneau $A'$ est seulement un
anneau semi-local intègre, mais non nécessairement local; tout idéal maximal
$\mathfrak n'$ de $B'$ est alors au-dessus d'un idéal maximal $\mathfrak m'$
de $A'$; comme tout anneau local de
$\operatorname{Spec}(B'\otimes_{A'}K')$ est un anneau local de
$\operatorname{Spec}(B')$, donc de l'un des
$\operatorname{Spec}(B'_{\mathfrak n'})$ (au-dessus du point générique de
$\operatorname{Spec}(A'_{\mathfrak m'}))$ on est ramené à prouver que les
anneaux locaux de
\[
\operatorname{Spec}(B'_{\mathfrak n'}
\otimes_{A'_{\mathfrak m'}}K')
\]
possèdent la propriété $\boldsymbol R$. Or, la définition de
$\boldsymbol P'$ implique que si $\boldsymbol P'(Z,k)$ est vraie, il en est
de même de $\boldsymbol P'(Z\otimes_k k',k')$ pour toute extension finie
$k'$ de $k$; le même raisonnement que dans (7.3.7) prouve que les fibres
formelles de $A'_{\mathfrak m'}$ possèdent la propriété $\boldsymbol P'$.
On voit comme dans la partie I) de la démonstration de (7.5.1) que
$A'_{\mathfrak m'}$ et $B'_{\mathfrak n'}$ vérifient les conditions (i) et
(ii) de (7.5.2); d'autre part, si $k'$ est le corps résiduel de
$A'_{\mathfrak m'}$, la propriété
\[
\boldsymbol P(\operatorname{Spec}(\widehat{B'_{\mathfrak n'}}
\otimes_{\widehat{A'_{\mathfrak m'}}}k'),k')
\]
est vraie : en effet, le complété de $A'_{\mathfrak m'}$ est l'un des
anneaux locaux composants de
$\widehat{A'}=\widehat A\otimes_A A'$ et le complété de
$B'_{\mathfrak n'}$ l'un des anneaux locaux composants de
\[
\widehat{B'}=\widehat B\otimes_B B'
=\widehat B\otimes_A A'
=\widehat B\otimes_{\widehat A}\widehat{A'};
\]
le raisonnement de I), dans (7.5.1), prouve donc notre assertion.

On peut donc se borner au cas où $A$ est en outre supposé intègre, et à
prouver que, si $K$ est le corps des fractions de $A$, les anneaux locaux de
$\operatorname{Spec}(B\otimes_A K)$ possèdent la propriété
$\boldsymbol R$.

\emph{II) Réduction au cas où $A$ et $B$ sont complets.} --- Soient $\xi$
le point générique de $\operatorname{Spec}(A)$, $y$ un point quelconque de
$f^{-1}(\xi)$; il s'agit de prouver que
$\boldsymbol R(\mathscr O_y)$ est vraie; compte tenu de la condition
$(\mathrm R_{\mathrm{II}})$, il suffira de voir qu'il existe dans
$\operatorname{Spec}(\widehat B)$ un point $z$ au-dessus de $y$ et tel que
$\boldsymbol R(\mathscr O_z)$ soit vraie. Or, il résulte de l'hypothèse
faite sur $\operatorname{Spec}(\widehat B\otimes_{\widehat A}k)$ et de
(7.5.1) que le morphisme
\enlargethispage{5\baselineskip}
\[
{}^a\widehat\varphi:\operatorname{Spec}(\widehat B)
\to\operatorname{Spec}(\widehat A)
\]
est un $\boldsymbol P$-morphisme. D'autre part, pour tout point
$z\in\operatorname{Spec}(\widehat B)$ au-dessus de $y$ (il en existe,
puisque $\widehat B$ est un $B$-module fidèlement plat), l'image $x$ de $z$
dans $\operatorname{Spec}(\widehat A)$ appartient à la fibre
$h^{-1}(\xi)$ pour le morphisme canonique
$h:\operatorname{Spec}(\widehat A)\to\operatorname{Spec}(A)$; en vertu de
l'hypothèse faite sur $A$, la propriété
$\boldsymbol R'(\mathscr O_x)$ est vraie; donc en vertu de
$(\mathrm R_{\mathrm V})$, $\boldsymbol R(\mathscr O_z)$ est vraie.
C.Q.F.D.
\end{proof}

\begin{examples}[7.5.3]
\label{IV.2.s7.5.3-fr}
Considérons les propriétés $\boldsymbol R(C)$ suivantes ($C$ étant un anneau
local noethérien) :
\begin{enumerate}
\item[\emph{(i)}] (aussi notée $(\mathrm i_n)$) $\operatorname{coprof}(C)\leq n$.
\oldpage[IV-2]{206}
\item[\emph{(ii)}] (aussi notée $(\mathrm{ii}_n)$) $C$ vérifie $(\mathrm S_n)$.
\item[\emph{(iii)}] $C$ est un anneau de Cohen--Macaulay.
\item[\emph{(iv)}] $C$ est un anneau régulier.
\item[\emph{(v)}] (aussi notée $(\mathrm v_n)$) $C$ est intègre, intégralement
clos et vérifie $(\mathrm R_n)$.
\item[\emph{(vi)}] $C$ est intègre et intégralement clos.
\item[\emph{(vii)}] $C$ est intègre.
\item[\emph{(viii)}] $C$ est réduit.
\end{enumerate}
Toutes ces propriétés vérifient $(\mathrm R_{\mathrm{III}})$, comme on l'a vu
dans (7.3.19, (i)). Elles vérifient aussi $(\mathrm R_{\mathrm{IV}})$ : pour
(i), cela résulte de (0, 16.4.10, (ii)), et pour (ii) et (iii), cela résulte de
(5.12.4) et du fait qu'un anneau local noethérien complet est caténaire
(5.6.4); pour (iv), c'est un cas particulier de (0, 17.1.8); pour (vii), cela
découle de (3.4.5) et pour (viii) de (3.4.6); pour (vi) cela résulte de
(5.12.7) et de ce qu'un anneau local noethérien complet est caténaire (5.6.4);
enfin, pour (v), cela résulte de ce qui précède et de (5.12.5).

On peut donc appliquer le théorème (7.5.1) lorsque $\boldsymbol R$ est l'une
quelconque des propriétés ci-dessus. D'autre part, on a déjà remarqué (7.3.11)
que les propriétés (i) à (viii) vérifient $(\mathrm R_{\mathrm{II}})$ (sauf
pour (vii), où la vérification de $(\mathrm R_{\mathrm{II}})$ est conséquence
de (2.1.14)). En ce qui concerne $(\mathrm R_{\mathrm V})$, on notera que si
l'on y fait $\boldsymbol R'=\boldsymbol R$, la condition
$(\mathrm R_{\mathrm V})$ se réduit à la condition
$(\mathrm R'_{\mathrm I})$ de (7.3.11), et on a noté que les propriétés (ii) à
(vi), ainsi que (viii), vérifient $(\mathrm R'_{\mathrm I})$ (7.3.11); pour la
propriété (i), on peut prendre pour $\boldsymbol R'$ la propriété d'être un
anneau de Cohen--Macaulay, en vertu de (6.3.2). Enfin, pour la propriété (vii),
la condition $(\mathrm R'_{\mathrm I})$ n'est plus vérifiée (ni d'ailleurs
$(\mathrm R_{\mathrm I})$), comme le montre l'exemple de (6.15.11, (ii)) ou de
(6.5.5, (ii)). Par contre, $(\mathrm R_{\mathrm V})$ est alors vraie lorsque
l'on prend pour $\boldsymbol R'$ la propriété d'être régulier : il suffit en
effet d'appliquer le lemme (7.5.1.1) en prenant pour $\mathcal C$ la catégorie
de tous les anneaux locaux noethériens, et pour $\boldsymbol R(C)$ la propriété
d'être intègre; cela est possible en vertu de (3.4.5).

On voit ainsi en particulier que lorsque $\boldsymbol R$ est l'une quelconque
des propriétés (i) à (viii), la conclusion du corollaire (7.5.2) est applicable
lorsque l'on suppose que les fibres formelles de $A$ sont géométriquement
régulières (cf. (7.8.2)).
\end{examples}

\begin{remarks}[7.5.4]
\EGAReaderTarget{ega4.c04.item.7-5-4}{IV.2.s7.5.4-fr}
\begin{enumerate}
\item[\emph{(i)}] Il serait intéressant de savoir si le théorème (7.5.1)
subsiste sans l'hypothèse de finitude (i) sur le corps résiduel de $B$. La
réponse est affirmative lorsque le corps résiduel $k$ de $A$ est de
caractéristique 0, comme on le voit en utilisant les résultats de Hironaka sur
la résolution des singularités (cf. (7.9.8)). Signalons le cas particulier
suivant de la question soulevée ici : Soient $A$ et $B$ deux anneaux
noethériens complets, $k$ le corps résiduel de $A$,
$\varphi:A\to B$ un homomorphisme local faisant de $B$ une $A$-algèbre
formellement lisse (0, 19.3.1) (ce qui équivaut à dire que $B$ est un
$A$-module plat et que $B\otimes_A k$ est géométriquement régulier sur $k$
(0, 19.7.1)). Alors les fibres du morphisme
${}^a\varphi:\operatorname{Spec}(B)\to\operatorname{Spec}(A)$ sont-elles
géométriquement régulières ? Il en est ainsi lorsque le corps résiduel $k'$ de
$B$ est une extension finie de $k$, d'après (7.5.1); on peut démontrer que la
réponse est encore affirmative lorsque $k'$ est une extension de type fini de
$k$. Nous ignorons la réponse lorsque $B\otimes_A k$ est un corps égal à la
clôture séparable de $k$.

\item[\emph{(ii)}] On pourrait énoncer un résultat analogue à (7.5.1) relatif
à un $A$-module $M$, un $B$-module $N$, tous deux de type fini, concernant les
propriétés de
\[
(M\otimes_A N)_y\otimes_{\mathscr O_x}k(y),
\]
où $y$ parcourt $\operatorname{Spec}(B)$ et $x=f(y)$, découlant de propriétés
de même type de $M_x$ et $N_y$.

\item[\emph{(iii)}] Lorsque la propriété $\boldsymbol R$ vérifie les conditions
$(\mathrm R'_{\mathrm I})$, $(\mathrm R_{\mathrm{II}})$,
$(\mathrm R_{\mathrm{III}})$ et $(\mathrm R_{\mathrm{IV}})$ et que les
hypothèses (i) et (ii) de (7.5.2) sont remplies, alors, des propriétés
$\boldsymbol R(A)$ et
\[
\boldsymbol P(\operatorname{Spec}(\widehat B
\otimes_{\widehat A}k),k)
\]
on déduit $\boldsymbol R(B)$. C'est le cas, comme on l'a vu (7.5.3), pour les
propriétés données en exemples dans (7.5.3), sauf pour (i) et (vii). En ce qui
concerne (vii), il est cependant plausible que la réponse à la question
suivante soit affirmative : soit $\varphi:A\to B$ un homomorphisme local
d'anneaux locaux noethériens, faisant de $B$ un $A$-module plat. Supposons que
$A$ soit complet, intègre et géométriquement unibranche et que la fibre
$\operatorname{Spec}(B\otimes_A k)$ (où $k$ est le corps résiduel de $A$) soit
géométriquement localement intègre; alors est-il vrai que $B$ soit intègre ?
On peut le démontrer lorsque $k$ est de caractéristique 0, en utilisant la
résolution des singularités de Hironaka (7.9.8); on peut montrer aussi que la
réponse est affirmative lorsque $B$ est une $A$-algèbre essentiellement de
type fini (1.3.8) (cf. (11.3.10) et (11.3.11)) ou lorsque
$\operatorname{Spec}(B\otimes_A k)$ est géométriquement normale (car il résulte
alors de (7.5.3) que le morphisme
$\operatorname{Spec}(B)\to\operatorname{Spec}(A)$ est normal, donc on conclut
par (6.15.10)). Mais même en supposant que le corps résiduel de $B$ soit égal
à celui de $A$ et que $B$ est aussi complet, nous ignorons si la réponse est
affirmative dans le cas général.

\item[\emph{(iv)}] Considérons la propriété $\boldsymbol R(C)$ : « $C$ est
réduit, équidimensionnel et vérifie $(\mathrm R_n)$ »; elle vérifie
$(\mathrm R_{\mathrm{IV}})$, en vertu de (5.12.5), mais on ignore par contre si
elle vérifie $(\mathrm R'_{\mathrm I})$, $(\mathrm R_{\mathrm{II}})$ ou
$(\mathrm R_{\mathrm{III}})$, les difficultés provenant de la vérification de
l'équidimensionalité.
\end{enumerate}
\end{remarks}

\oldpage[IV-2]{207}

Dans la suite de ce numéro, nous allons appliquer (7.5.1) à l'étude des
produits tensoriels complétés $A\widehat\otimes_k B$ d'anneaux locaux
noethériens qui sont des algèbres sur un corps $k$.

\begin{lemma}[7.5.5]
\label{IV.2.s7.5.5-fr}
Soient $k$ un corps, $A$, $B$ deux anneaux locaux noethériens complets
contenant $k$, le corps résiduel de $A$ étant une extension finie de $k$. Soit
$C$ le produit tensoriel complété $A\widehat\otimes_k B$
(0$_{\mathrm I}$, 7.7.5). Alors :
\begin{enumerate}
\item[\emph{(i)}] $C$ est un anneau semi-local noethérien complet.
\item[\emph{(ii)}] $C$ est un $A$-module plat et un $B$-module plat.
\item[\emph{(iii)}] Si $\mathfrak m$ est l'idéal maximal de $A$,
$\mathfrak mC$ est contenu dans le radical de $C$, et $C/\mathfrak mC$ est
isomorphe à $(A/\mathfrak m)\otimes_k B$.
\item[\emph{(iv)}] Les corps résiduels des composants locaux de $C$ sont des
extensions finies du corps résiduel de $B$.
\end{enumerate}
Les propriétés (i), (iii) et la première assertion de (ii) sont des cas
particuliers de (0, 19.7.1.2). Pour prouver que $C$ est un $B$-module plat,
notons que pour tout $h>0$,
\[
C/\mathfrak m^hC=(A/\mathfrak m^h)\otimes_k B
\]
est un $B$-module plat, puisque $k$ est un corps; il suffit donc d'appliquer
(0$_{\mathrm{III}}$, 10.2.6) à chacun des composants locaux de $C$ (dont $C$
est le composé direct). Enfin, si $\mathfrak n$ est l'idéal maximal de $B$, les
corps résiduels des composants locaux de $C$ sont aussi ceux des composants
locaux de l'anneau artinien
$(A/\mathfrak m)\otimes_k(B/\mathfrak n)$, qui, par hypothèse, sont des
extensions finies de $B/\mathfrak n$.
\end{lemma}

\begin{proposition}[7.5.6]
\label{IV.2.s7.5.6-fr}
Soient $k$ un corps, $A$, $B$ deux anneaux locaux noethériens complets
contenant $k$, et dont les corps résiduels sont des extensions finies de $k$.
Soit d'autre part $\boldsymbol R$ une propriété vérifiant les conditions
$(\mathrm R'_{\mathrm I})$, $(\mathrm R_{\mathrm{III}})$ et
$(\mathrm R_{\mathrm{IV}})$ (la propriété $\boldsymbol P$ étant définie à
partir de $\boldsymbol R$ par (7.5.0)). Supposons que $\boldsymbol R(A)$ soit
vraie, ainsi que $\boldsymbol P(\operatorname{Spec}(B),k)$ (ce qui signifie que
pour toute extension finie $k'$ de $k$ et pour chacun des anneaux locaux
complets $B_i$ dont le composé direct est $B\otimes_k k'$,
$\boldsymbol R(B_i)$ est vraie). Alors, pour chacun des anneaux locaux complets
$C_j$ dont $A\widehat\otimes_k B$ est le composé direct,
$\boldsymbol R(C_j)$ est vraie.
\end{proposition}

\begin{proof}
Il résulte de (7.5.5, (iii)) que chacun des homomorphismes $A\to C_j$ est local
et de (7.5.5, (ii)) que chacun des $C_j$ est un $A$-module plat; enfin, par
(7.5.5, (iv)), les corps résiduels des $C_j$ sont des extensions finies de
$A/\mathfrak m$. D'autre part, la propriété
\[
\boldsymbol P(\operatorname{Spec}(C\otimes_A(A/\mathfrak m)),A/\mathfrak m)
\]
est vraie, car $C/\mathfrak mC=(A/\mathfrak m)\otimes_kB$, et l'assertion
résulte de la définition de $\boldsymbol P$ et de l'hypothèse que
$A/\mathfrak m$ est une extension finie de $k$. Toutes les conditions de
(7.5.1) sont donc vérifiées pour les homomorphismes locaux $A\to C_j$ et les
morphismes correspondants
$\operatorname{Spec}(C_j)\to\operatorname{Spec}(A)$ sont donc des
$\boldsymbol P$-morphismes; la conclusion résulte alors de
$(\mathrm R'_{\mathrm I})$ et de l'hypothèse que $\boldsymbol R(A)$ est vraie.
\end{proof}

\begin{corollary}[7.5.7]
\label{IV.2.s7.5.7-fr}
(Chevalley). --- Soient $k$ un corps parfait (resp. algébriquement clos), $A$,
$B$ deux anneaux locaux noethériens complets contenant $k$ et dont les corps
résiduels sont des extensions finies de $k$. Alors, si $A$ et $B$ sont réduits
(resp. intègres) le produit tensoriel complété $A\widehat\otimes_k B$ est
réduit (resp. intègre).
\end{corollary}

\begin{proof}
I) Supposons $k$ parfait, $A$ et $B$ réduits; soient $A_i$ (resp. $B_j$) les
quotients de $A$ (resp. $B$) par ses idéaux premiers minimaux
($1\leq i\leq r$, $1\leq j\leq s$), qui sont complets; l'hypothèse entraîne
que $A$ (resp. $B$) s'identifie à un sous-anneau du composé direct des $A_i$
(resp. $B_j$); les produits tensoriels étant pris sur un corps,
$A\otimes_kB$ s'identifie à un sous-anneau du composé direct des
$A_i\otimes_kB_j$, et on vérifie aussitôt que la topologie produit tensoriel
de $A\otimes_kB$ est induite par le produit des topologies des
$A_i\otimes_kB_j$. Il en résulte que $A\widehat\otimes_kB$ s'identifie à un
sous-anneau du composé direct des $A_i\widehat\otimes_kB_j$, et on est donc
ramené au cas où $A$ et $B$ sont intègres. Soient alors $A'$ et $B'$ les
clôtures intégrales de $A$ et $B$ respectivement; on sait, par le théorème de
finitude de Nagata (0, 23.1.6), que $A'$ (resp. $B'$) est un $A$-module (resp.
$B$-module) de type fini et un anneau local complet. En outre,
$A\widehat\otimes_kB$ s'identifie à un sous-anneau de
$A'\widehat\otimes_kB'$ : en effet, il suffira de prouver que
$A\widehat\otimes_kB$ s'identifie à un sous-anneau de
$A'\widehat\otimes_kB$, et ce dernier à un sous-anneau de
$A'\widehat\otimes_kB'$. Cela résulte du lemme suivant :

\begin{lemma}[7.5.7.1]
\label{IV.2.s7.5.7.1-fr}
Soient $A$, $B$ deux anneaux locaux noethériens complets contenant un corps
$k$ et dont les corps résiduels sont des extensions finies de $k$. Soit
$\mathfrak m$ l'idéal maximal de $A$. Pour tout $A$-module $M$ de type fini
(muni de la topologie $\mathfrak m$-adique), le produit tensoriel complété
$M\widehat\otimes_kB$ (0$_{\mathrm I}$, 7.7.1) s'identifie canoniquement à
$M\otimes_A(A\widehat\otimes_kB)$.
\end{lemma}

\begin{proof}
En effet, on a un isomorphisme canonique
$M\otimes_kB\xrightarrow{\sim}M\otimes_A(A\otimes_kB)$ donc un homomorphisme
canonique composé
\[
\varphi:M\otimes_kB\xrightarrow{\sim}M\otimes_A(A\otimes_kB)
\longrightarrow M\otimes_A(A\widehat\otimes_kB)
\]
et il est immédiat que cet homomorphisme est continu pour les topologies
produits tensoriels; en outre, $M\otimes_A(A\widehat\otimes_kB)$ est séparé et
complet ((7.5.5) et (0$_{\mathrm I}$, 7.7.8)), donc on a aussi par complétion
un homomorphisme continu
\[
\widehat\varphi:M\widehat\otimes_kB\longrightarrow
M\otimes_A(A\widehat\otimes_kB).
\]

\oldpage[IV-2]{208}

Il est immédiat que cet homomorphisme est bijectif lorsque $M$ est libre de
type fini; comme on a une suite exacte $L'\to L\to M\to0$, où $L$ et $L'$ sont
des $A$-modules libres de type fini, on en déduit un diagramme commutatif
\[
\xymatrix{
L'\widehat\otimes_kB \ar[r] \ar[d] &
L\widehat\otimes_kB \ar[r] \ar[d] &
M\widehat\otimes_kB \ar[r] \ar[d] & 0 \\
L'\otimes_A(A\widehat\otimes_kB) \ar[r] &
L\otimes_A(A\widehat\otimes_kB) \ar[r] &
M\otimes_A(A\widehat\otimes_kB) \ar[r] & 0
}
\]
où les lignes sont exactes (la première en vertu de la définition du produit
tensoriel complété et de (0$_{\mathrm{III}}$, 13.2.2)); les deux premières
flèches verticales étant des bijections, il en est de même de la troisième.
\end{proof}

Le fait que $A\widehat\otimes_kB$ s'identifie à un sous-anneau de
$A'\otimes_A(A\widehat\otimes_kB)$ résulte alors de ce que $A$ est un
sous-anneau de $A'$ et de ce que $A\widehat\otimes_kB$ est un $A$-module plat
(7.5.5). On peut ainsi supposer en outre que $A$ et $B$ sont intégralement
clos; comme en outre $k$ est supposé parfait, $\operatorname{Spec}(A)$ et
$\operatorname{Spec}(B)$ sont géométriquement normaux sur $k$ en vertu de
(6.7.7, b)); on peut alors appliquer (7.5.6) en prenant pour
$\boldsymbol R$ la propriété (vi) de (7.5.3).

II) Supposons $k$ algébriquement clos, $A$ et $B$ intègres. Le raisonnement de
I) ramène encore au cas où $A$ et $B$ sont intégralement clos; il résulte alors
de (7.5.6) que $\operatorname{Spec}(A\widehat\otimes_kB)$ est normal, donc
$A\widehat\otimes_kB$ est composé direct d'anneaux locaux complets
intégralement clos, et tout revient à voir que $A\widehat\otimes_kB$ est un
anneau local; il suffit pour cela de vérifier que si $\mathfrak m$ et
$\mathfrak n$ sont les idéaux maximaux de $A$ et $B$,
$(A/\mathfrak m)\otimes_k(B/\mathfrak n)$ est un anneau local (voir
démonstration de (0, 19.7.1.2)); mais comme $A/\mathfrak m$ et
$B/\mathfrak n$ sont des extensions finies de $k$, elles sont identiques à
$k$, d'où la conclusion.
\end{proof}

\begin{remark}[7.5.8]
\label{IV.2.s7.5.8-fr}
Il serait intéressant de déterminer si, dans (7.5.7), on peut remplacer
l'hypothèse que $k$ est parfait ou algébriquement clos par l'hypothèse que
$\operatorname{Spec}(A)$ ou $\operatorname{Spec}(B)$ est géométriquement
intègre sur $k$, ou tout au moins par l'hypothèse que $A$ (par exemple) soit
intègre et contienne un sous-anneau $A_0$ isomorphe à un anneau de séries
formelles $k[[T_1,\ldots,T_n]]$, et que le corps des fractions $K$ de $A$ soit
séparable sur le corps des fractions $K_0$ de $A_0$ (on peut montrer que cette
condition entraîne que $A$ est géométriquement intègre sur $k$, et que les
deux propriétés sont équivalentes lorsque $[k:k^p]$ est fini ($p$ étant
l'exposant caractéristique de $k$)). De même il serait désirable de développer
des variantes de (7.5.6) et (7.5.7) où on affaiblirait l'hypothèse de finitude
des corps résiduels de $A$ et $B$, en supposant par exemple qu'un seul d'entre
eux est extension finie de $k$, l'autre étant quelconque.
\end{remark}

\subsection{Applications : I. Anneaux japonais locaux}
\label{IV.2.ss7.6-fr}

\begin{proposition}[7.6.1]
\label{IV.2.s7.6.1-fr}
Soit $A$ un anneau local noethérien réduit dont les fibres formelles sont
géométriquement normales. Alors le complété $\widehat A$ est réduit, la
fermeture intégrale $A'$ de $A$ dans son anneau total des fractions est une
$A$-algèbre finie (donc un anneau noethérien semi-local) et son complété
$\widehat{A'}$ est isomorphe à la fermeture intégrale de $\widehat A$ dans son
anneau total des fractions.
\end{proposition}

\begin{proof}
Les fibres formelles de $A$ sont a fortiori géométriquement réduites, donc
l'hypothèse que $A$ est réduit entraîne qu'il en est de même de $\widehat A$,
en vertu de (7.3.17). Soient $\mathfrak p_i$ ($1\leq i\leq n$) les idéaux
premiers minimaux de $A$, et soit $B_i=A/\mathfrak p_i$; les fibres formelles
des anneaux locaux $B_i$ sont alors aussi géométriquement normales (7.3.15),
donc les $\widehat B_i$ sont aussi réduits et il résulte de (0, 23.1.7, (i))
que la clôture intégrale $B'_i$ de $B_i$ dans son corps des fractions est un
$B_i$-module de type fini, donc un $A$-module de type fini; comme $A'$ est
composé direct des $B'_i$ (II, 6.3.8), on voit que $A'$ est un $A$-module de
type fini. Soient $\mathfrak m_j$ ($1\leq j\leq r$) les idéaux maximaux de
l'anneau semi-local $A'$; on sait que le complété $\widehat{A'}$ de $A'$
s'identifie au composé direct des complétés
$\widehat{A'_{\mathfrak m_j}}$ des $A'_{\mathfrak m_j}$ (Bourbaki,

\oldpage[IV-2]{209}

\emph{Alg. comm.}, chap. III, § 2, n\textsuperscript{o} 13, cor. de la prop. 19).
Or, il résulte de l'hypothèse et de (7.3.15) que les fibres formelles des
$A'_{\mathfrak m_j}$ sont géométriquement normales; comme
$\operatorname{Spec}(A')$ est normal par définition, il en est de même des
$\operatorname{Spec}(A'_{\mathfrak m_j})$, et on déduit donc de (7.3.17) que
$\operatorname{Spec}(\widehat{A'_{\mathfrak m_j}})$ est normal pour tout $j$,
donc aussi $\operatorname{Spec}(\widehat{A'})$. D'autre part (Bourbaki,
\emph{Alg. comm.}, chap. IV, § 2, n\textsuperscript{o} 5, cor. 3 de la prop. 9
et chap. III, § 3, n\textsuperscript{o} 4, th. 3), $\widehat{A'}$ s'identifie
à $A'\otimes_A\widehat A$ puisque $A'$ est un $A$-module de type fini; comme
$A'$ contient $A$ et est contenu dans l'anneau total des fractions $R$ de $A$,
et que $\widehat A$ est un $A$-module plat, $\widehat{A'}$ contient
$\widehat A$ et est contenu dans $R'=R\otimes_A\widehat A$; enfin, comme
$\widehat A$ est un $A$-module plat, tout élément régulier de $A$ est aussi
$\widehat A$-régulier (0$_{\mathrm I}$, 6.3.4); donc $R'$ s'identifie
canoniquement à un sous-anneau de l'anneau total des fractions $R''$ de
$\widehat A$ (Bourbaki, \emph{Alg. comm.}, chap. II, § 2,
n\textsuperscript{o} 1, Remarque 7). Comme
$\operatorname{Spec}(\widehat{A'})$ est normal et que $\widehat{A'}$ est un
$\widehat A$-module de type fini, $\widehat{A'}$ est bien la fermeture
intégrale de $\widehat A$ dans $R''$.
\end{proof}

\begin{corollary}[7.6.2]
\label{IV.2.s7.6.2-fr}
Sous les hypothèses de (7.6.1), il y a une correspondance biunivoque canonique
entre l'ensemble des idéaux maximaux $\mathfrak m_j$ de $A'$ (autrement dit,
l'ensemble des points de $\operatorname{Spec}(A')$ au-dessus du point fermé de
$A$) et l'ensemble des idéaux premiers minimaux $\mathfrak q_j$ de
$\widehat A$ (autrement dit, l'ensemble des points maximaux de
$\operatorname{Spec}(\widehat A)$); dans cette correspondance, le complété
$\widehat{A'_{\mathfrak m_j}}$ de $A'_{\mathfrak m_j}$ est isomorphe à la
clôture intégrale de $\widehat A/\mathfrak q_j$.
\end{corollary}

\begin{proof}
On sait en effet que la fermeture intégrale de $\widehat A$ dans son anneau
total des fractions est composée directe des clôtures intégrales des
$\widehat A/\mathfrak q_j$, qui sont des anneaux locaux complets (0, 23.1.6).
\end{proof}

\begin{corollary}[7.6.3]
\label{IV.2.s7.6.3-fr}
Sous les hypothèses de (7.6.1), pour que $\widehat A$ soit intègre, il faut et
il suffit que $A$ soit unibranche; pour que $\widehat A$ soit géométriquement
unibranche, il faut et il suffit que $A$ le soit.
\end{corollary}

\begin{proof}
C'est un cas particulier de (7.6.2).
\end{proof}

\begin{theorem}[7.6.4]
\EGAReaderTarget{ega4.c04.thm.7-6-4}{IV.2.s7.6.4-fr}
(Zariski--Nagata). --- Soit $A$ un anneau semi-local noethérien. Les conditions
suivantes sont équivalentes :
\begin{enumerate}
\item[\emph{a)}] Pour toute $A$-algèbre finie réduite $C$, le complété
$\widehat C$ est un anneau réduit.
\item[\emph{a$'$)}] Pour tout anneau quotient intègre $B$ de $A$, de corps des
fractions $K$, le complété $\widehat B$ est réduit, et les corps composants
$L_i$ de l'anneau total des fractions de $\widehat B$ sont des extensions
séparables de $K$.
\item[\emph{a$''$)}] Les fibres formelles de $A$ sont géométriquement réduites
(autrement dit, pour tout anneau quotient intègre $B$ de $A$, de corps des
fractions $K$, la $K$-algèbre $\widehat B\otimes_BK$ est séparable (4.6.2)).
\item[\emph{b)}] Tout anneau quotient intègre $B$ de $A$ est un anneau
japonais.
\end{enumerate}
\end{theorem}

\begin{proof}
Pour montrer que $a)$ entraîne $a')$, il suffit de vérifier que les $L_i$ sont
des extensions séparables de $K$, ou encore (4.6.1) que pour toute extension
finie $K'$ de $K$, l'anneau $\widehat B\otimes_BK'$ est réduit; or $K'$ est
engendré par un nombre fini d'éléments entiers sur $B$, et ces derniers
engendrent une sous-$B$-algèbre finie $B'$ de $K'$, dont $K'$ est corps des
fractions. On a $\widehat{B'}=\widehat B\otimes_BB'$
((0$_{\mathrm I}$, 7.3.3) et Bourbaki, \emph{Alg. comm.}, chap. IV, § 2,
n\textsuperscript{o} 5, cor. 3 de

\oldpage[IV-2]{210}

la prop. 9), donc
$\widehat B\otimes_BK'=\widehat{B'}\otimes_{B'}K'$ par associativité du produit
tensoriel; mais $B'$ est une $A$-algèbre intègre finie, donc $\widehat{B'}$ est
un anneau réduit en vertu de $a)$, et comme $\widehat{B'}$ est un $B'$-module
plat, les éléments $\neq0$ de $B'$ sont $\widehat{B'}$-réguliers, ce qui
entraîne que $\widehat{B'}\otimes_{B'}K'$ s'identifie à un sous-anneau de
l'anneau total des fractions de $\widehat{B'}$, et a fortiori est réduit.

Pour voir que $a')$ entraîne $a'')$, considérons un point quelconque $x$ de
$X=\operatorname{Spec}(A)$, et soit $Y$ le sous-schéma fermé intègre de $X$
ayant $\overline{\{x\}}$ pour espace sous-jacent; on a
$Y=\operatorname{Spec}(B)$, où $B$ est un anneau quotient intègre de $A$, et
$\operatorname{Spec}(\widehat B)=Y\times_X\operatorname{Spec}(\widehat A)$,
de sorte que la fibre formelle de $A$ au point $x$ est la même que celle de
$B$, ou encore est égale à $\operatorname{Spec}(\widehat B\otimes_BK)$. Comme
les anneaux locaux de $\operatorname{Spec}(\widehat B\otimes_BK)$ sont ceux de
$\operatorname{Spec}(\widehat B)$ aux points de la fibre de $x$, l'hypothèse
entraîne que $\widehat B\otimes_BK$ est réduit, et les conditions de $a')$
entraînent donc que $\widehat B\otimes_BK$ est une $K$-algèbre séparable
(4.6.1).

La condition $a'')$ entraîne $a)$, car il résulte de (7.3.15) que les fibres
formelles de $C$ sont alors géométriquement réduites, et si $C$ est réduit, il
en est donc de même de $\widehat C$ (cas particulier de (7.3.17)).

Le fait que $a')$ implique $b)$ est un cas particulier de (0, 23.1.7). Reste
donc à prouver que $b)$ entraîne $a)$. Notons que si $C$ est une $A$-algèbre
finie, pour tout idéal premier $\mathfrak q$ de $C$, l'image réciproque
$\mathfrak p$ de $\mathfrak q$ dans $A$ est un idéal premier et
$C/\mathfrak q$ est une $(A/\mathfrak p)$-algèbre finie, donc l'hypothèse $b)$
entraîne que tout anneau quotient intègre de $C$ est un anneau japonais. On est
ainsi ramené à prouver que sous l'hypothèse $b)$, le complété de tout quotient
intègre de $A$ est réduit. Raisonnons par récurrence sur $n=\dim(A)$,
l'assertion étant triviale pour $n=0$; en remplaçant $A$ par les quotients de
$A$ par ses idéaux premiers minimaux $\mathfrak p_i$, on peut se borner au cas
où $A$ est intègre (tout quotient intègre de $A$ étant quotient de l'un des
$A/\mathfrak p_i$). Pour tout idéal premier $\mathfrak p\neq0$, l'hypothèse de
récurrence montre déjà que le complété de $A/\mathfrak p$ est réduit, et il
suffit donc de prouver que $\widehat A$ est réduit. En outre, la clôture
intégrale $A'$ de $A$ est par hypothèse un $A$-module de type fini, donc un
anneau semi-local noethérien et $\widehat A$ s'identifie à un sous-anneau de
$\widehat{A'}$ ((0$_{\mathrm I}$, 7.3.3) et Bourbaki, \emph{Alg. comm.}, chap.
IV, § 2, n\textsuperscript{o} 5, cor. 3 de la prop. 9); il suffira donc de
prouver que $\widehat{A'}$ est réduit; on a vu ci-dessus que l'hypothèse $b)$
est aussi vérifiée par $A'$, qui par ailleurs est aussi de dimension $n$
(0, 16.1.5); on peut donc se borner au cas où $A$ est intégralement clos. Soit
$t\neq0$ un élément du radical de $A$, et soient $\mathfrak q_j$
($1\leq j\leq n$) les idéaux premiers minimaux parmi ceux contenant $tA$; on a
les propriétés suivantes :
\begin{enumerate}
\item[\emph{(i)}] $t$ est régulier.
\item[\emph{(ii)}] $A/tA$ n'a pas d'idéaux premiers associés immergés.
\item[\emph{(iii)}] Les $A_{\mathfrak q_j}$ sont des anneaux de valuation
discrète.
\item[\emph{(iv)}] Les complétés des $A/\mathfrak q_j$ sont réduits.
\end{enumerate}
En effet, (i) est triviale puisque $A$ est intègre et $t\neq0$. Comme $A$ est
intégralement clos, $A/tA$ vérifie $(\mathrm S_1)$, c'est-à-dire (5.7.7) est
sans idéaux premiers associés immergés (Bourbaki, \emph{Alg. comm.}, chap. VII,
§ 1, n\textsuperscript{o} 4, prop. 8). Toujours parce que $A$ est
intégralement

\oldpage[IV-2]{211}

clos, les $A_{\mathfrak q_j}$ le sont et on sait (\emph{loc. cit.}) que ces
anneaux sont de dimension 1, donc ce sont des anneaux de valuation discrète,
d'où (iii). Enfin, (iv) provient de l'hypothèse de récurrence. La démonstration
sera achevée si nous prouvons le

\begin{lemma}[7.6.4.1]
\label{IV.2.s7.6.4.1-fr}
(Zariski). --- Soient $A$ un anneau semi-local noethérien, $t$ un élément de
son radical vérifiant les conditions (i) à (iv) précédentes; alors
$\widehat A$ est réduit.
\end{lemma}

\begin{proof}
Posons pour simplifier $A'=\widehat A$, et considérons $t$ comme un élément de
$A'$; on a $A'/tA'=(A/tA)\otimes_AA'$, et comme $A'$ est un $A$-module plat,
il résulte de (3.3.1) que les idéaux premiers de $A'$ associés au $A'$-module
$A'/tA'$ sont les idéaux premiers $\mathfrak q'_{jh}$ tels que
$\mathfrak q'_{jh}$ soit au-dessus de $\mathfrak q_j$ et associé au
$k_j$-module $(A'/tA')\otimes_Ak_j$, où $k_j$ est le corps des fractions de
$A/\mathfrak q_j$. Or,
$\operatorname{Spec}((A'/tA')\otimes_Ak_j)$ est la fibre formelle de $A$ au
point $\mathfrak q_j$, ou aussi celle de $A/\mathfrak q_j$ au point
$\mathfrak q_j$ (point générique de
$\operatorname{Spec}(A/\mathfrak q_j)$). Comme en vertu de (iv) le complété
$A'/\mathfrak q_jA'$ de $A/\mathfrak q_j$ est réduit, il en est de même de la
fibre formelle de $A/\mathfrak q_j$ au point générique de
$\operatorname{Spec}(A/\mathfrak q_j)$; cette fibre formelle n'a donc pas de
cycle premier associé immergé et ses anneaux locaux aux points génériques
$\mathfrak q'_{jh}$ de ses composantes irréductibles sont des corps (3.2.1).
On voit donc par (3.3.3) et l'hypothèse (ii) que $A'/tA'$ n'a pas d'idéaux
premiers associés immergés, et d'autre part que
$\mathfrak q_jA'_{\mathfrak q'_{jh}}$ est l'idéal maximal de
$A'_{\mathfrak q'_{jh}}$ pour tout $h$; comme $A_{\mathfrak q_j}$ est un
anneau de valuation discrète, son idéal maximal
$\mathfrak q_jA_{\mathfrak q_j}$ est principal, donc l'idéal maximal de
l'anneau local noethérien $A'_{\mathfrak q'_{jh}}$ est principal, ce qui
entraîne que cet anneau est un anneau de valuation discrète (Bourbaki,
\emph{Alg. comm.}, chap. VI, § 3, n\textsuperscript{o} 6, prop. 9). Nous venons
ainsi de vérifier les hypothèses (i) à (iv) pour l'anneau local complet $A'$.
Il suffit donc de montrer que si $A$ est complet et vérifie les hypothèses (i)
à (iii) ((iv) l'étant automatiquement dans ce cas), alors $A$ est réduit. Or,
les hypothèses (i) et (ii) impliquent que, si
\[
\varphi:A\longrightarrow\prod_{j=1}^nA_{\mathfrak q_j}
\]
est l'homomorphisme canonique, on a
$t^mA=\varphi^{-1}(\varphi(t^mA))$ (3.4.9); comme l'image canonique de $tA$
dans l'anneau de valuation discrète $A_{\mathfrak q_j}$ est une puissance
$\mathfrak q_j^hA_{\mathfrak q_j}$ de l'idéal maximal, on voit que sur $A$ la
topologie $(tA)$-préadique est l'image réciproque par $\varphi$ de la topologie
produit sur $\prod_{j=1}^nA_{\mathfrak q_j}$. Or, comme $t$ appartient au
radical de $A$, la topologie $(tA)$-préadique est séparée, donc $\varphi$ est
injectif. Mais en vertu de l'hypothèse (iii), les $A_{\mathfrak q_j}$ sont
intègres, donc $\prod_{j=1}^nA_{\mathfrak q_j}$ est réduit, et a fortiori il en
est de même de $A$. C.Q.F.D.
\end{proof}
\end{proof}

\begin{corollary}[7.6.5]
\label{IV.2.s7.6.5-fr}
Soit $A$ un anneau semi-local noethérien vérifiant les conditions équivalentes
de (7.6.4); toute $A$-algèbre semi-locale essentiellement de type fini vérifie
aussi les conditions de (7.6.4).
\end{corollary}

\begin{proof}
La condition $a'')$ de (7.6.4) signifie en effet que $A$ est un
$\boldsymbol P$-anneau, où $\boldsymbol P(Z,k)$ est la propriété « $Z$ est
géométriquement réduit sur $k$ »; le corollaire est donc conséquence du
théorème général (7.4.4).
\end{proof}

\begin{corollary}[7.6.6]
\label{IV.2.s7.6.6-fr}
Soient $A$ un anneau semi-local noethérien intègre de dimension 1, $K$ son
corps des fractions. Pour que $A$ soit un anneau japonais, il faut et il suffit
que le complété $\widehat A$

\oldpage[IV-2]{212}

de $A$ soit réduit et que, si $\mathfrak q_j$ ($1\leq j\leq n$) sont les
idéaux premiers minimaux de $\widehat A$, les corps des fractions des anneaux
intègres $\widehat A/\mathfrak q_j$ soient des extensions séparables de $K$.
\end{corollary}

\begin{proof}
Les quotients intègres de $A$ sont en effet $A$ lui-même et des corps, donc la
condition $b)$ de (7.6.4) équivaut à l'hypothèse que $A$ est un anneau
japonais, et l'hypothèse $a')$ aux conditions de l'énoncé sur $\widehat A$,
d'où la conclusion.
\end{proof}

\begin{remarks}[7.6.7]
\label{IV.2.s7.6.7-fr}
\begin{enumerate}
\item[\emph{(i)}] Les conditions équivalentes de (7.6.6) signifient encore que
les fibres formelles de $A$ sont géométriquement régulières (7.3.19, (iv)). On
a déjà observé (\emph{loc. cit.}) que cette condition est vérifiée lorsque $A$
est un anneau de valuation discrète dont le corps des fractions est de
caractéristique 0.

\item[\emph{(ii)}] Le même raisonnement que dans la première partie de la
démonstration de (7.6.4) montre que pour un anneau semi-local noethérien $A$,
les deux conditions suivantes sont équivalentes :
\begin{enumerate}
\item[\emph{a)}] Pour tout anneau quotient intègre $B$ de $A$, le complété
$\widehat B$ est réduit.
\item[\emph{a$'$)}] Les fibres formelles de $A$ sont réduites.
\end{enumerate}
Compte tenu de (0, 23.1.7, (i)), ces deux conditions impliquent la suivante :
\begin{enumerate}
\item[\emph{b)}] Pour tout anneau quotient intègre $B$ de $A$, la clôture
intégrale de $B$ est une $B$-algèbre finie.
\end{enumerate}
Lorsque $A$ est un anneau universellement caténaire, on peut prouver que la
condition $b)$ entraîne inversement $a)$; nous n'aurons pas à utiliser ce
résultat.
\end{enumerate}
\end{remarks}

\subsection{Applications : II. Anneaux universellement japonais}
\label{IV.2.ss7.7-fr}

\begin{env}[7.7.1]
\label{IV.2.s7.7.1-fr}
Rappelons (0, 23.1.1) que l'on appelle anneau universellement japonais un
anneau $A$ tel que toute $A$-algèbre intègre de type fini soit un anneau
japonais. Il revient au même de dire que pour toute $A$-algèbre intègre de type
fini $B$, la clôture intégrale de $B$ est une $B$-algèbre finie.
\end{env}

\begin{theorem}[7.7.2]
\EGAReaderTarget{ega4.c04.thm.7-7-2}{IV.2.s7.7.2-fr}
(Nagata). --- Soit $A$ un anneau noethérien. Les conditions suivantes sont
équivalentes :
\begin{enumerate}
\item[\emph{a)}] $A$ est un anneau universellement japonais.
\item[\emph{b)}] Tout anneau quotient intègre de $A$ est un anneau japonais.
\item[\emph{c)}] Pour tout idéal maximal $\mathfrak m$ de $A$, tout anneau
quotient intègre de $A_{\mathfrak m}$ est un anneau japonais, et pour tout
anneau quotient intègre $B$ de $A$, de corps des fractions $K$, toute
extension radicielle finie $K'$ de $K$ et toute sous-$B$-algèbre finie $B'$ de
$K'$, de corps des fractions $K'$, il existe $f'\neq0$ dans $B'$ tel que
$B'_{f'}$ soit intégralement clos (cf. (6.13.7)).
\end{enumerate}
De plus, si $A$ vérifie ces conditions, il en est de même de tout anneau de
fractions $S^{-1}A$ de $A$ et de toute $A$-algèbre de type fini.
\end{theorem}

\begin{proof}
Montrons d'abord que $b)$ entraîne $c)$; tout anneau quotient intègre de
$A_{\mathfrak m}$ est un anneau de fractions d'un anneau quotient intègre de
$A$, donc un anneau japonais (0, 23.1.1). D'autre part, tout quotient $B$ de
$A$ par un de ses idéaux premiers est un anneau japonais, donc il en est de
même de $B'$ (0, 23.1.1). On en déduit que l'ensemble
$\operatorname{Nor}(\operatorname{Spec}(B'))$ est ouvert (6.13.3), donc
contient un ensemble ouvert non vide
$D(f')=\operatorname{Spec}(B'_{f'})$.

\oldpage[IV-2]{213}

Montrons en second lieu que si $A$ vérifie $c)$, il en est de même de toute
$A$-algèbre de type fini $B$. En premier lieu, pour tout idéal premier
$\mathfrak q$ de $B$, $B_{\mathfrak q}$ est anneau local en un idéal premier
de $S^{-1}B$, où $S=A-\mathfrak p$, $\mathfrak p$ étant l'image réciproque de
$\mathfrak q$ dans $A$; comme par hypothèse tout anneau quotient intègre de
$A_{\mathfrak p}$ est un anneau japonais, il en est de même de tout anneau
quotient intègre de $B_{\mathfrak q}$, puisque $S^{-1}B$ est une
$A_{\mathfrak p}$-algèbre de type fini (7.6.5). D'autre part, si $C$ est un
quotient intègre de $B$, de corps des fractions $K$, $K'$ une extension
radicielle finie de $K$ et $C'$ une sous-$C$-algèbre finie de $K'$, de corps
des fractions $K'$, $C'$ est une $A$-algèbre intègre de type fini, et en vertu
de (6.13.7), l'hypothèse $c)$ sur $A$ entraîne que le spectre de $C'$ contient
une partie ouverte non vide dont tous les points sont normaux; autrement dit,
il y a un $g'\neq0$ dans $C'$ tel que $C'_{g'}$ soit intégralement clos, ce qui
prouve notre assertion.

Ce résultat montre que, pour prouver l'équivalence de $a)$, $b)$ et $c)$, il
suffit de montrer que $c)$ entraîne $b)$, et même de prouver que pour tout
anneau quotient intègre $B$ de $A$, de corps des fractions $K$, la clôture
intégrale $B'$ est un $B$-module de type fini. Or, la condition $c)$ entraîne
qu'il existe $f\neq0$ dans $B$ tel que $B_f$ soit intégralement clos, et on a
ainsi vérifié pour $B$ la condition (i) de (6.13.6). D'autre part, la condition
(ii) de (6.13.6) est aussi vérifiée par $B$; en effet, pour tout idéal maximal
$\mathfrak q$ de $B$, $B_{\mathfrak q}$ est quotient d'un anneau local
$A_{\mathfrak p}$, où $\mathfrak p$ est maximal dans $A$; comme
$A_{\mathfrak p}$ est par hypothèse un anneau japonais, il en est de même de
$B_{\mathfrak q}$; en outre, pour tout idéal premier $\mathfrak r$ de $B$,
$B_{\mathfrak r}$ est anneau de fractions d'un $B_{\mathfrak q}$, où
$\mathfrak q$ est un idéal maximal de $B$, donc est encore un anneau japonais,
ce qui prouve notre assertion.

On a en outre vu plus haut que si $A$ vérifie les conditions équivalentes
$a)$, $b)$, $c)$, il en est de même de toute $A$-algèbre de type fini. D'autre
part, si $A$ vérifie $b)$, il en est de même de tout anneau de fractions
$S^{-1}A$, car tout quotient intègre de $S^{-1}A$ est de la forme
$S^{-1}A/S^{-1}\mathfrak p$, où $\mathfrak p$ est un idéal premier de $A$, et
comme $A/\mathfrak p$ est par hypothèse un anneau japonais, il en est de même
de $S^{-1}(A/\mathfrak p)$ (0, 23.1.1). Cela achève de prouver la dernière
assertion de l'énoncé.
\end{proof}

\begin{corollary}[7.7.3]
\label{IV.2.s7.7.3-fr}
Soient $A$ un anneau noethérien vérifiant les conditions suivantes :
\begin{enumerate}
\item[\emph{(i)}] Pour tout idéal maximal $\mathfrak m$ de $A$, les fibres
formelles de $A_{\mathfrak m}$ sont géométriquement réduites.
\item[\emph{(ii)}] Les conditions équivalentes de (6.12.4) sont vérifiées.
\end{enumerate}
Alors $A$ est universellement japonais.
\end{corollary}

\begin{proof}
En effet, il résulte alors de (7.6.4) que tout anneau quotient intègre de
$A_{\mathfrak m}$ est un anneau japonais; d'où la conclusion, en vertu de
(7.7.2).
\end{proof}

\begin{corollary}[7.7.4]
\label{IV.2.s7.7.4-fr}
Un anneau de Dedekind dont le corps des fractions est de caractéristique 0 (en
particulier $\mathbf Z$) est un anneau universellement japonais. Un anneau
local noethérien dont les fibres formelles sont géométriquement réduites (en
particulier un anneau local noethérien complet) est universellement japonais.
\end{corollary}

\begin{proof}
La première assertion résulte de (7.7.3), (7.6.7, (i)) et (6.12.6). La seconde
résulte de (7.6.4) et (7.7.2).
\end{proof}

\subsection{Anneaux excellents}
\label{IV.2.ss7.8-fr}

\begin{env}[7.8.1]
\label{IV.2.s7.8.1-fr}
Les numéros précédents de ce paragraphe (ainsi que (6.11), (6.12) et (6.13))
ont été consacrés à l'étude de problèmes que l'on rencontre fréquemment dans
l'utilisation des anneaux et des préschémas noethériens, et qui peuvent se
grouper dans les types suivants :

\oldpage[IV-2]{214}

A) Pour un anneau \emph{local} noethérien $A$, les propriétés de $A$ se
« transmettent-elles » à son \emph{complété} $\widehat A$ ? Par exemple, si
$A$ est réduit (resp. intègre, resp. intègre et intégralement clos), en est-il
de même de $\widehat A$ ? La plupart de ces questions sont liées aux propriétés
locales des \emph{fibres formelles} de $A$ (rappelons que ce sont les fibres du
morphisme canonique
$\operatorname{Spec}(\widehat A)\to\operatorname{Spec}(A)$). Une exception est
formée par les propriétés liées à la notion de dimension, par exemple la
propriété d'être équidimensionnel; c'est alors la \emph{condition des chaînes}
et ses divers raffinements qui jouent le rôle essentiel.

B) Pour un préschéma localement noethérien $X$ (en particulier pour un schéma
affine $\operatorname{Spec}(A)$), l'ensemble des $x\in X$ où l'anneau local
$\mathcal O_x$ possède une certaine propriété (par exemple être intégralement
clos, ou de Cohen-Macaulay, ou régulier) est-il \emph{ouvert} ?

C) Pour un anneau \emph{intègre} $A$, la fermeture intégrale de $A$ dans une
extension finie de son corps des fractions est-elle un $A$-module \emph{de type
fini} ? Bien entendu, cette question peut se traduire pour les préschémas
noethériens (\textbf{II}, 6.3).

Les problèmes de type B) ou C) peuvent aussi être posés pour des anneaux
\emph{locaux}, mais il ne suffit pas en général qu'ils soient résolus
affirmativement pour tout anneau local $A_{\mathfrak p}$ d'un anneau $A$ pour
qu'ils le soient pour $A$ (cf. (6.13.6)).

Soulignons en outre que dans l'étude de ces problèmes, nous nous sommes
systématiquement préoccupés de savoir si la réponse affirmative à l'un d'eux
est \emph{stable} pour les deux opérations les plus importantes de l'Algèbre
commutative : la localisation et le passage à une algèbre de type fini.

Les résultats obtenus dans l'étude de ces problèmes amènent à dégager la
définition d'une classe d'anneaux noethériens dont le comportement à cet égard
est le meilleur possible :
\end{env}

\begin{definition}[7.8.2]
\EGAReaderTarget{ega4.c04.def.7-8-2}{IV.2.s7.8.2-fr}
On dit qu'un anneau $A$ est excellent s'il est noethérien et s'il
vérifie les conditions suivantes :
\begin{enumerate}
\item[\emph{(i)}] $A$ est universellement caténaire (ou, ce qui revient au même
(5.6.3, (i)), pour tout idéal premier $\mathfrak p$ de $A$,
$A_{\mathfrak p}$ est universellement caténaire).
\item[\emph{(ii)}] Pour tout idéal premier $\mathfrak p$ de $A$, les fibres
formelles de $A_{\mathfrak p}$ sont géométriquement régulières.
\item[\emph{(iii)}] Pour tout anneau quotient intègre $B$ de $A$ et toute
extension finie radicielle $K'$ du corps des fractions $K$ de $B$, il existe
une sous-$B$-algèbre finie $B'$ de $K'$, contenant $B$, ayant $K'$ pour corps
des fractions, et telle que l'ensemble des points réguliers de
$\operatorname{Spec}(B')$ contienne un ouvert non vide.
\end{enumerate}
\upshape
Avec cette terminologie, l'essentiel des résultats du \S~7 et d'une partie du
\S~6 se résume alors comme suit :
\end{definition}

\begin{scholium}[7.8.3]
\EGAReaderTarget{ega4.c04.statement.7-8-3}{IV.2.s7.8.3-fr}
(i) Dans les conditions (i) et (ii) de la définition (7.8.2) on peut se borner
aux idéaux $\mathfrak p$ maximaux dans $A$. Pour qu'un anneau local noethérien
soit excellent, il faut et

\oldpage[IV-2]{215}

il suffit qu'il soit universellement caténaire et que ses fibres formelles
soient géométriquement régulières.

(ii) Si $A$ est un anneau excellent, il en est de même de tout anneau de
fractions $S^{-1}A$ et de toute $A$-algèbre de type fini.

(iii) Un anneau local complet (en particulier un corps) est excellent. Un
anneau de Dedekind dont le corps des fractions est de caractéristique $0$ (en
particulier $\mathbf Z$) est excellent.

(iv) Soient $A$ un anneau excellent, $X=\operatorname{Spec}(A)$; l'ensemble
$\operatorname{Reg}(X)$ (resp. $\operatorname{Nor}(X)$,
$\mathrm U_{\mathrm R_n}(X)$) des points où $X$ est régulier (resp. normal,
resp. vérifie $(\mathrm R_n)$) est ouvert dans $X$. Pour tout
$\mathcal O_X$-Module cohérent $\mathcal F$, l'ensemble
$\mathrm U_{\mathrm C_n}(\mathcal F)$ (resp.
$\mathrm U_{\mathrm S_n}(\mathcal F)$, $\operatorname{CM}(\mathcal F)$) des
$x\in X$ où $\operatorname{coprof}(\mathcal F_x)\leq n$ (resp. des points où
$\mathcal F$ vérifie $(\mathrm S_n)$, resp. des points où $\mathcal F_x$ est un
$\mathcal O_x$-Module de Cohen-Macaulay) est ouvert dans $X$.

(v) Soient $A$ un anneau excellent, $\mathfrak J$ un idéal de $A$,
$\widehat A$ le séparé complété de $A$ pour la topologie $\mathfrak J$-préadique.
Alors le morphisme canonique
$f:\operatorname{Spec}(\widehat A)\to\operatorname{Spec}(A)$ est régulier
(autrement dit, plat et à fibres géométriquement régulières). Si l'on pose
$X=\operatorname{Spec}(A)$, $X'=\operatorname{Spec}(\widehat A)$, on a (avec
les notations de (iv) et en posant
$\mathcal F'=\mathcal F\otimes_{\mathcal O_X}\mathcal O_{X'}$),
\begin{env}[7.8.3.1]
\label{IV.2.s7.8.3.1-fr}
\begin{align*}
\operatorname{Reg}(X')&=f^{-1}(\operatorname{Reg}(X)),&
\operatorname{Nor}(X')&=f^{-1}(\operatorname{Nor}(X)),&
\mathrm U_{\mathrm R_n}(X')&=f^{-1}(\mathrm U_{\mathrm R_n}(X)),\\
\mathrm U_{\mathrm C_n}(\mathcal F')&=f^{-1}(\mathrm U_{\mathrm C_n}(\mathcal F)),&
\mathrm U_{\mathrm S_n}(\mathcal F')&=f^{-1}(\mathrm U_{\mathrm S_n}(\mathcal F)),&
\operatorname{CM}(\mathcal F')&=f^{-1}(\operatorname{CM}(\mathcal F))
\end{align*}
\end{env}

En particulier, si $\mathfrak J$ est contenu dans le radical de $A$ (par
exemple si $A$ est local et $\mathfrak J$ son idéal maximal), pour que $A$ soit
régulier (resp. normal, resp. réduit, resp. vérifie $(\mathrm R_n)$, resp. soit
de coprofondeur $\leq n$, resp. vérifie $(\mathrm S_n)$, resp. soit un anneau
de Cohen-Macaulay), il faut et il suffit qu'il en soit de même de $\widehat A$;
en particulier, pour que $A$ n'ait pas d'idéaux premiers associés immergés, il
faut et il suffit qu'il en soit de même de $\widehat A$.

(vi) Un anneau excellent $A$ est universellement japonais; en particulier, si
$A$ est intègre, sa fermeture intégrale dans toute extension finie de son corps
des fractions est une $A$-algèbre finie.

(vii) Soit $A$ un anneau local excellent réduit. Alors le complété $\widehat A$
est réduit, la fermeture intégrale $A'$ de $A$ dans son anneau total des
fractions est une $A$-algèbre finie (donc un anneau semi-local), et la fermeture
intégrale de $\widehat A$ dans son anneau total des fractions est isomorphe au
complété $\widehat {A'}$ de $A'$. En outre, il y a une correspondance biunivoque
canonique entre l'ensemble des points maximaux de
$\operatorname{Spec}(\widehat A)$ (autrement dit, l'ensemble des idéaux premiers
minimaux de $\widehat A$) et l'ensemble des points fermés de
$\operatorname{Spec}(A')$ (autrement dit, l'ensemble des idéaux maximaux de
$A'$). Pour que l'anneau local $A$ soit unibranche, il faut et il suffit que
$\widehat A$ soit intègre; pour que $A$ soit géométriquement unibranche, il faut
et il suffit que $\widehat A$ le soit.

(viii) Si $A$ est un anneau excellent intègre, la clôture intégrale de $A$ est
l'intersection des clôtures intégrales des anneaux $A_{\mathfrak p}$, où
$\mathfrak p$ parcourt l'ensemble des idéaux premiers de $A$ de hauteur $1$.

(ix) Soient $A$ un anneau excellent, $X=\operatorname{Spec}(A)$, $Z$ une partie
fermée de $X$, $U=X-Z$, $i:U\to X$ l'injection canonique, $\mathcal F$ un
$\mathcal O_U$-Module cohérent. Pour que le $\mathcal O_X$-Module
$i_*(\mathcal F)$ soit cohérent, il faut et il suffit que, pour tout
$x\in\operatorname{Ass}(\mathcal F)$, on ait
$\operatorname{codim}(\overline{\{x\}}\cap Z,\overline{\{x\}})\geq2$. En
particulier, si $A$ est intègre et $\mathcal F$ sans torsion, pour que
$i_*(\mathcal F)$ soit cohérent, il faut et il suffit que
$\operatorname{codim}(Z,X)\geq2$.

\oldpage[IV-2]{216}

(x) Soit $A$ un anneau local excellent. Pour tout anneau quotient intègre $B$
de $A$, $\widehat B$ est équidimensionnel. Pour que $A$ soit équidimensionnel,
il faut et il suffit que $\widehat A$ le soit.
\end{scholium}

\begin{proof}
La plupart de ces résultats ont déjà été démontrés.

(i) La première assertion résulte de (5.6.3, (i)) et (7.4.4, (ii)); la seconde,
de (7.4.4), (7.3.18), (6.12.7) et (6.12.4).

L'assertion (ii) résulte de (5.6.1), (5.6.3, (i)), (7.4.4) et (6.12.4).

(iii) La première assertion a été vue dans (5.6.4) et (7.4.4), compte tenu de
(i). La seconde a été vue dans (5.6.4) et (7.3.19, (iv)), compte tenu de (i).

L'assertion (iv) résulte de (6.12.4), (6.13.5), (6.12.9) et (6.11.8) (compte
tenu de (ii)).

La première assertion de (v) est conséquence de (7.4.6); les relations
(7.8.3.1) en découlent, compte tenu respectivement de (6.5.1), (6.5.4),
(6.5.3), (6.3.5) et (6.4.2). La dernière assertion de (v) est le cas
particulier correspondant à la propriété $(\mathrm S_n)$ pour $n=1$ (5.7.5).

L'assertion (vi) résulte de (7.7.3), et l'assertion (vii) de (7.6.1), (7.6.2)
et (7.6.3). Pour prouver (viii), notons que l'anneau $A^{(1)}$, intersection
des $A_{\mathfrak p}$ pour tous les idéaux premiers $\mathfrak p$ de hauteur
$1$ de $A$, est une $A$-algèbre finie, comme il résulte de (vi), de (7.8.2,
(i)) et de (5.11.2). Comme la clôture intégrale $A'$ de $A$ est une $A$-algèbre
finie, les idéaux premiers de hauteur $1$ de $A'$ sont exactement ceux qui sont
au-dessus des idéaux premiers de hauteur $1$ de $A$ (5.10.17, (iv)); la
conclusion résulte donc de (0, 23.2.9).

La première assertion de (ix) est conséquence de (vi) et de (5.11.4); la
seconde en est un cas particulier, si l'on observe que lorsque $\mathcal F$ est
sans torsion, $\operatorname{Ass}(\mathcal F)$ est réduit au point générique de
$\operatorname{Spec}(A)$. Enfin, pour prouver (x), il suffit d'observer que $B$
est un anneau local excellent par (ii); comme $B$ est universellement caténaire
et universellement japonais en vertu de (vi), l'anneau $B^{(1)}$ est une
$B$-algèbre finie (5.11.2), ce qui montre que $A$ est strictement formellement
caténaire par (7.2.5, $c)$) et achève la démonstration de (x).
\end{proof}

\begin{remarks}[7.8.4]
\EGAReaderTarget{ega4.c04.item.7-8-4}{IV.2.s7.8.4-fr}
\upshape
(i) Si on considère les anneaux noethériens $A$ qui vérifient seulement les
conditions (ii) et (iii) de la définition (7.8.2), on vérifie par inspection
des démonstrations précédentes que les assertions (ii), (iii), (iv), (v), (vi)
et (vii) de (7.8.3) sont encore valables, en remplaçant « excellent » par
« vérifiant » (7.8.2, (ii) et (iii)).

(ii) L'anneau local caténaire $A$ construit dans (5.6.11) vérifie les
conditions (ii) et (iii) (mais non la condition (i)) de la définition (7.8.2).
En effet, la fibre formelle en l'idéal maximal $\mathfrak n C_{\mathfrak n}$ de
$A$ est trivialement géométriquement régulière, et celle aux autres idéaux
premiers $\mathfrak p$ de $A$ coïncide avec celle de l'anneau $E$ en ces idéaux
premiers; or $E$ est une $V$-algèbre de type fini et $V$ est un anneau de
valuation discrète, qui est par suite un anneau excellent si l'on prend le
corps $k_0$ de caractéristique $0$ (7.8.3, (iii)); par suite $E$ est un anneau
excellent (7.8.3, (ii)) et cela prouve donc que $A$ vérifie la condition (7.8.2,
(ii)). D'autre part, le complémentaire du point fermé dans
$\operatorname{Spec}(A)$ est isomorphe à un ouvert de $\operatorname{Spec}(E)$,
dont un ouvert affine est le spectre d'un anneau

\oldpage[IV-2]{217}

excellent, donc vérifie la condition (6.12.4, $a)$) en vertu de (7.8.3, (v));
cela prouve que $A$ lui-même vérifie la condition de (6.12.4, $a)$), donc aussi
(6.12.4, $c)$), qui n'est autre que la condition (7.8.2, (iii)). Pourtant $A$
ne satisfait pas à la condition (7.8.2, (i)), puisque nous avons vu (5.6.11)
qu'il n'est pas universellement caténaire.

(iii) On peut remplacer la condition (i) de la définition (7.8.2) par la
suivante (plus faible en apparence, compte tenu de (5.6.10)) :

(i$'$) \emph{Pour tout idéal premier minimal} $\mathfrak p_i$ \emph{de} $A$
$(1\leq i\leq r)$, \emph{soit} $A'_i$ \emph{la clôture intégrale de l'anneau
intègre} $A_i=A/\mathfrak p_i$; \emph{alors, pour tout idéal maximal}
$\mathfrak m'$ \emph{de} $A'_i$, \emph{on a, en désignant par} $\mathfrak m$
\emph{l'image réciproque de} $\mathfrak m'$ \emph{dans} $A$,
\begin{env}[7.8.4.1]
\label{IV.2.s7.8.4.1-fr}
\[
\dim(A_{\mathfrak m})=\dim((A'_i)_{\mathfrak m'})
\]
\end{env}

En effet, on a vu (Remarque (i)) que sous les seules hypothèses (ii) et (iii)
de (7.8.2), les \emph{analogues} des propriétés (ii), (v), (vi) et (vii) de
(7.8.3) sont valables. On déduit donc d'abord de (vi) et (ii) que les $A'_i$
vérifient (7.8.2, (ii) et (iii)); il résulte alors de (v) que les complétés
$((A'_i)_{\mathfrak m'})^\wedge$ sont intègres (et intégralement clos). Soit
maintenant $\mathfrak m$ un idéal maximal quelconque de $A$; pour tout $i$ tel
que $\mathfrak p_i\subset\mathfrak m$, la clôture intégrale de
$A_{\mathfrak m}/\mathfrak p_iA_{\mathfrak m}$ est une
$(A_{\mathfrak m}/\mathfrak p_iA_{\mathfrak m})$-algèbre finie par (vi), donc
semi-locale, et ses composants locaux sont de la forme
$(A'_i)_{\mathfrak m'}$, où $\mathfrak m'$ est un idéal premier
(nécessairement maximal) de $A'_i$ au-dessus de
$\mathfrak m/\mathfrak p_i$; il résulte alors de (vii), de ce qui précède et de
(7.8.4.1) que les quotients du complété
$(A_{\mathfrak m}/\mathfrak p_iA_{\mathfrak m})^\wedge$ par ses idéaux
premiers minimaux ont tous même dimension égale à $\dim(A_{\mathfrak m})$. Par
suite (7.1.9 et 7.1.8, $b)$), $A_{\mathfrak m}$ est formellement caténaire, et
\emph{a fortiori} (7.1.11) universellement caténaire; il en est donc de même de
$A_{\mathfrak p}$ pour tout idéal premier $\mathfrak p$ de $A$ (5.6.3, (i)),
donc aussi de $A$.

En particulier, si $A$ est normal, ou plus généralement si $A_{\mathfrak m}$
est unibranche pour tout idéal maximal $\mathfrak m$ de $A$, la condition
(i$'$) est toujours remplie, puisque $\mathfrak m$ ne peut contenir qu'un seul
idéal premier $\mathfrak p_i$ (0, 16.1.5); on peut dans ce cas omettre (i) dans
la définition (7.8.2). Plus particulièrement, si $A$ est un anneau noethérien
\emph{local unibranche}, il est excellent si et seulement si ses fibres
formelles sont géométriquement régulières.

(iv) Rappelons qu'un anneau quotient d'un anneau local de Cohen-Macaulay (et
\emph{a fortiori} un anneau quotient d'un anneau local régulier) vérifie aussi
les propriétés (7.8.3, (ix) et (x)) en vertu de (7.2.7); mais l'intérêt des
quotients d'anneaux réguliers réside surtout dans leurs propriétés
cohomologiques (chap. III, 3\textsuperscript{e} Partie).

(v) Nous verrons plus loin (\S~18) que si $k$ est un corps valué complet non
discret, l'anneau $k\{\!\{T_1,\ldots,T_n\}\!\}$ des \emph{séries entières
convergentes} (dans un voisinage de $0$ dans $k^n$) est un \emph{anneau
excellent}.
\end{remarks}

\begin{env}[7.8.5]
\label{IV.2.s7.8.5-fr}
On dit qu'un préschéma localement noethérien $X$ est \emph{excellent} si pour
un recouvrement $(U_\alpha)$ de $X$ formé d'ouverts affines, chacun des anneaux
$A_\alpha$ des $U_\alpha$ est excellent; cette propriété est alors vraie pour
\emph{tout} recouvrement $(U_\alpha)$ de $X$ formé d'ouverts affines.
\end{env}

\begin{proposition}[7.8.6]
\EGAReaderTarget{ega4.c04.prop.7-8-6}{IV.2.s7.8.6-fr}
Soit $X$ un préschéma excellent.
\begin{enumerate}
\item[\emph{(i)}] Si $f:X'\to X$ est un morphisme localement de type fini, $X'$ est
excellent.

\oldpage[IV-2]{218}

\item[\emph{(ii)}] Si $X$ est réduit, son normalisé $X'$ (\textbf{II}, 6.3.8) est fini
sur $X$.
\item[\emph{(iii)}] Les ensembles $\operatorname{Reg}(X)$, $\operatorname{Nor}(X)$,
$\mathrm U_{\mathrm R_n}(X)$ sont ouverts dans $X$, ainsi que
$\mathrm U_{\mathrm C_n}(\mathcal F)$, $\mathrm U_{\mathrm S_n}(\mathcal F)$
et $\operatorname{CM}(\mathcal F)$ pour tout $\mathcal O_X$-Module cohérent
$\mathcal F$.
\end{enumerate}
\end{proposition}

\begin{proof}
Cela résulte aussitôt de (7.8.3, (ii), (vi) et (iv)). Notons aussi que (7.8.3,
(ix)) est valable sans modification d'énoncé lorsque $X$ est un préschéma
excellent.
\end{proof}

\subsection{Anneaux excellents et résolution des singularités}
\label{subsection:IV.2.7.9-fr}

\begin{env}[7.9.1]
\label{IV.2.s7.9.1-fr}
Étant donné un préschéma \emph{localement noethérien réduit} $X$, on appelle
\emph{morphisme résolvant} pour $X$ un morphisme $f:X'\to X$ tel que $X'$ soit
\emph{régulier} et que $f$ soit \emph{propre} et \emph{birationnel}. Lorsqu'il
existe un tel morphisme, on dit qu'on peut \emph{résoudre les singularités} de
$X$ (ou plus simplement que l'on peut \emph{résoudre} $X$). Par exemple, si
$A$ est un \emph{anneau japonais} de dimension $1$, on peut résoudre
$X=\operatorname{Spec}(A)$ puisque le morphisme $X'\to X$ (où $X'$ est le
\emph{normalisé} de $X$) est fini (donc propre) et birationnel, et que les
anneaux locaux de $X'$ sont des anneaux intégralement clos de dimension $1$,
donc des anneaux de valuation discrète, et par suite sont réguliers.
\end{env}

\begin{env}[7.9.2]
\label{IV.2.s7.9.2-fr}
Il est clair que si l'on peut résoudre $X$, on peut aussi résoudre tout
préschéma induit sur un \emph{ouvert} de $X$ et tout préschéma local
$\operatorname{Spec}(\mathcal O_x)$ de $X$ (\textbf{II}, 5.4.2). D'autre part,
si l'on peut résoudre $X$, il est clair qu'il existe dans $X$ un ouvert
\emph{partout} dense contenu dans $\operatorname{Reg}(X)$. Ces remarques
montrent aussitôt que si, pour tout sous-préschéma \emph{fermé intègre} $Y$ de
$X$ et tout $Y$-préschéma $Y'$ \emph{intègre, fini} et \emph{radiciel} sur
$Y$, on peut résoudre $Y'$, alors les ouverts affines de $X$ \emph{vérifient
la condition} (iii) de (7.8.2).

D'autre part, pour les schémas locaux, on a la proposition suivante :
\end{env}

\begin{proposition}[7.9.3]
\label{IV.2.s7.9.3-fr}
Soit $A$ un anneau local noethérien réduit, et supposons que l'on puisse
résoudre $\operatorname{Spec}(A)$. Alors les fibres du morphisme canonique
$\operatorname{Spec}(\widehat A)\to\operatorname{Spec}(A)$ aux points maximaux
de $\operatorname{Spec}(A)$ sont régulières.
\end{proposition}

\begin{proof}
Posons $X=\operatorname{Spec}(A)$, $X'=\operatorname{Spec}(\widehat A)$, et
soit $g:X'\to X$ le morphisme canonique. Soit $f:Y\to X$ un morphisme
résolvant; posons $Y'=X'\times_XY$, et soient $f':Y'\to X'$ et $g':Y'\to Y$
les projections canoniques, de sorte que l'on a un diagramme commutatif
\[
\xymatrix{
X'\ar[d]_g & Y'\ar[l]_{f'}\ar[d]^{g'}\\
X & Y\ar[l]^f
}
\]

Il suffira de prouver que le préschéma $Y'$ est \emph{régulier} : en effet,
comme $f$ est birationnel et de type fini, il y a un ouvert $U$ partout dense
dans $X$ tel que la restriction $f^{-1}(U)\to U$ de $f$ soit un isomorphisme
et que $f^{-1}(U)$ soit partout dense dans $Y$ (\textbf{I}, 6.5.5); les
préschémas induits respectivement sur l'ouvert $g^{-1}(U)$ de $X'$ et l'ouvert
$g'^{-1}(f^{-1}(U))$ de $Y'$ sont donc isomorphes; cela prouve que
$g^{-1}(U)$ est régulier, et \emph{a fortiori} il en est de même des fibres de
$g$ aux points maximaux de $X$, qui sont contenus dans $U$.

\oldpage[IV-2]{219}

Notons que le morphisme $f'$ est propre (\textbf{II}, 5.4.2), et en
particulier de type fini, donc $Y'$ est noethérien, puisqu'il en est ainsi de
$X'$. Soit $a'$ le point fermé de $X'$; pour voir que $Y'$ est régulier, il
suffira de prouver que \emph{pour tout} $y'\in f'^{-1}(a')$, \emph{l'anneau
local $\mathcal O_{Y',y'}$ est régulier}. En effet, on aura alors
$f'^{-1}(a')\subset\operatorname{Reg}(Y')$. Mais comme $X'$ est le spectre
d'un anneau local complet et que $f'$ est de type fini,
$\operatorname{Reg}(Y')$ est \emph{ouvert} dans $Y'$ (6.12.8), et comme le
morphisme $f'$ est \emph{fermé} et que
$f'^{-1}(a')\subset\operatorname{Reg}(Y')$, il existe un voisinage $V'$ de
$a'$ dans $X'$ tel que $f'^{-1}(V')\subset\operatorname{Reg}(Y')$. Comme
$\widehat A$ est un anneau local, on a nécessairement $V'=X'$, donc
$\operatorname{Reg}(Y')=Y'$ et la proposition sera démontrée.

Or, si $a$ est le point fermé de $X$, on a $g^{-1}(a)=\{a'\}$ et les corps
résiduels $k(a)$ et $k(a')$ sont isomorphes; donc
$g'^{-1}(f^{-1}(a))=f'^{-1}(a')$ et la restriction
$f'^{-1}(a')\to f^{-1}(a)$ de $g'$ est un \emph{isomorphisme} de ces deux
fibres. Soit $y=g'(y')$ et posons $B=\mathcal O_{Y,y}$; si
$\mathfrak n=\mathfrak m_y$ est l'idéal maximal de $B$, ce qui précède montre
qu'il existe dans l'anneau (noethérien) $B'=B\otimes_A\widehat A$ un seul
idéal maximal $\mathfrak n'$ au-dessus de $\mathfrak n$ et que l'on a
$\mathcal O_{Y',y'}=B'_{\mathfrak n'}$. Pour montrer que l'anneau local
$B'_{\mathfrak n'}$ est régulier, il suffit d'établir que son \emph{complété}
(0, 17.1.5) l'est; or, par hypothèse $B$ est régulier, donc il en est de même
de son complété $\widehat B$. La conclusion résultera par suite du lemme
suivant :

\begin{lemma}[7.9.3.1]
\label{IV.2.s7.9.3.1-fr}
Soient $A$, $B$ deux anneaux locaux noethériens,
$\rho:A\to B$ un homomorphisme local tel que l'anneau
$B'=B\otimes_A\widehat A$ soit noethérien. Alors, pour tout idéal maximal
$\mathfrak n'$ de $B'$ au-dessus de l'idéal maximal $\mathfrak n$ de $B$ et
de l'idéal maximal de $\widehat A$, les complétés de $B$ et de
$B'_{\mathfrak n'}$ sont isomorphes.
\end{lemma}

\begin{proof}
Soit $\mathfrak m$ l'idéal maximal de $A$ et posons
$B''=B'_{\mathfrak n'}$. Pour tout entier $h>0$, on a
$B'/\mathfrak m^hB'=B\otimes_A(\widehat A/\mathfrak m^h\widehat A)$; mais
$\widehat A/\mathfrak m^h\widehat A$ est isomorphe à
$A/\mathfrak m^h$, donc $B'/\mathfrak m^hB'$ est isomorphe à
$B/\mathfrak m^hB$, et en particulier c'est un anneau local dont l'idéal
maximal est $\mathfrak n'/\mathfrak m^hB'$; par suite
$B''/\mathfrak m^hB''$, isomorphe à
$(B'/\mathfrak m^hB')_{\mathfrak n'}$, est $B$-isomorphe à
$B'/\mathfrak m^hB'$, et finalement à $B/\mathfrak m^hB$; en particulier,
l'idéal maximal de $B''$ est égal à $\mathfrak nB''$, et
$B''/\mathfrak n^hB''$ est isomorphe à $B/\mathfrak n^h$ d'après ce qui
précède. Comme
$\widehat{B''}=\varprojlim_h(B''/\mathfrak n^hB'')$, le lemme est démontré,
ainsi que (7.9.3).
\end{proof}
\end{proof}

\begin{corollary}[7.9.4]
\label{IV.2.s7.9.4-fr}
Soit $A$ un anneau local noethérien.
\begin{enumerate}
\item[\emph{(i)}] Si, pour tout anneau quotient intègre $B$ de $A$, on peut
résoudre $\operatorname{Spec}(B)$, alors les fibres formelles de $A$ sont
régulières.
\item[\emph{(ii)}] Supposons que, pour tout anneau quotient intègre $B$ de
$A$, et tout anneau intègre $B'$ contenant $B$, qui est une $B$-algèbre finie
et dont le corps des fractions est radiciel sur celui de $B$, on puisse
résoudre $\operatorname{Spec}(B')$; alors les fibres formelles de $A$ sont
géométriquement régulières.
\end{enumerate}
\end{corollary}

\begin{proof}
Toute fibre formelle de $A$ étant fibre formelle d'un anneau quotient intègre
de $A$ au point générique de son spectre, il est clair que (i) est conséquence
immédiate de (7.9.3), et (ii) résulte de (i) et de (6.7.7).
\end{proof}

\begin{proposition}[7.9.5]
\label{IV.2.s7.9.5-fr}
Soit $X$ un préschéma localement noethérien, tel que, pour tout préschéma $Y$
intègre et fini sur $X$, on puisse résoudre $Y$. Alors les anneaux des ouverts
affines de $X$ vérifient les conditions (ii) et (iii) de (7.8.2); si en outre
$X$ est universellement caténaire (5.6.3) (et

\oldpage[IV-2]{220}

en particulier si $X$ est \emph{localement immersible dans un préschéma
régulier} (5.6.4)), $X$ est un \emph{préschéma excellent} (7.8.5).
\end{proposition}

\begin{proof}
Cela résulte aussitôt de (7.9.4) et (7.9.2).
\end{proof}

\begin{remark}[7.9.6]
\label{IV.2.s7.9.6-fr}
Il est possible que la réciproque de (7.9.5) soit vraie, c'est-à-dire que
l'hypothèse que $X$ est réduit et que les anneaux des ouverts affines de $X$
vérifient les conditions (ii) et (iii) de (7.8.2) implique la possibilité de
résoudre $X$ (et par suite aussi la possibilité de résoudre tout préschéma
réduit localement de type fini sur $X$, en vertu de (7.4.4) et (6.12.4)).
C'est vrai en tout cas lorsqu'on se borne aux préschémas noethériens réduits
dont les corps résiduels sont de \emph{caractéristique} $0$, comme le montrent
les résultats récents de Hironaka [35] (ce dernier énonce ses résultats sous
des hypothèses trop restrictives, ses raisonnements étant en fait valables
sous les seules hypothèses (ii) et (iii) de (7.8.2) pour les anneaux des
ouverts affines de $X$, jointes au fait que les corps résiduels sont de
caractéristique $0$). Au contraire, pour le cas général, on ne sait jusqu'ici
résoudre que les schémas algébriques de dimension $2$ sur un corps parfait
[36]. Vu l'importance que la résolution des singularités est en train de
prendre dans l'étude topologique des schémas (notamment en ce qui concerne
leurs propriétés homologiques et homotopiques), cela donne un intérêt
particulier à la catégorie des anneaux excellents ou des préschémas
excellents.

On peut aussi noter à ce propos que la partie la moins délicate des
raisonnements de Hironaka (\emph{loc, cit.}) montre qu'indépendamment de toute
hypothèse sur les caractéristiques des corps résiduels, et en utilisant
uniquement les propriétés (ii) et (iii) de (7.8.2) pour les anneaux locaux
d'un préschéma $X$, le problème de la résolution des singularités de $X$ se
ramène au cas où $X=\operatorname{Spec}(A)$, $A$ étant un anneau local
noethérien intègre et \emph{complet}. Par conséquent, si des résultats
ultérieurs devaient mettre en défaut la conjecture avancée plus haut, et
devaient donc amener à formuler des conditions plus restrictives sur les
anneaux locaux de $X$, cela ne pourrait guère être que par l'intermédiaire de
conditions restrictives imposées à des anneaux locaux complets (probablement
des conditions concernant leurs corps résiduels, peut-être la condition
$[k:k^p]<+\infty$ ($p$ exposant caractéristique de $k$)).
\end{remark}

\begin{env}[7.9.7]
\label{IV.2.s7.9.7-fr}
Considérons d'une part une sous-catégorie pleine $\mathcal C$ de la catégorie
des préschémas localement noethériens, d'autre part une propriété
$\boldsymbol R(A)$, soumises aux conditions suivantes :

$1^\circ$ Pour tout $X\in\mathcal C$, tout préschéma localement de type fini
sur $X$ appartient à $\mathcal C$. Pour tout anneau noethérien $A$ tel que
$\operatorname{Spec}(A)\in\mathcal C$ et toute partie multiplicative $S$ de
$A$, on a $\operatorname{Spec}(S^{-1}A)\in\mathcal C$.

$2^\circ$ Pour tout $X\in\mathcal C$, l'ensemble
$U_{\boldsymbol R}(X)$ des $x\in X$ tels que
$\boldsymbol R(\mathcal O_{X,x})$ soit vraie est ouvert dans $X$.

$3^\circ$ Pour tout anneau local noethérien $A$ tel que
$\operatorname{Spec}(A)\in\mathcal C$ et tout élément régulier $t$ de l'idéal
maximal de $A$, $\boldsymbol R(A/tA)$ entraîne $\boldsymbol R(A)$.

Désignons alors comme dans (7.5.0) par $\boldsymbol P(Z,k)$, pour un corps $k$
et un $k$-préschéma $Z\in\mathcal C$, la propriété suivante :

« Pour toute extension finie $k'$ de $k$, tous les anneaux locaux de
$Z\otimes_k k'$ vérifient la propriété $\boldsymbol R$. »

\oldpage[IV-2]{221}

Nous supposerons encore que $\boldsymbol R$ vérifie la condition suivante :

$4^\circ$ Si $\boldsymbol P(Z,k)$ est vraie, alors, pour toute extension
\emph{de type fini} $k''$ de $k$, tous les anneaux locaux de
$Z\otimes_k k''$ vérifient la propriété $\boldsymbol R$.

On notera que ces conditions sont vérifiées lorsque l'on prend pour
$\mathcal C$ la catégorie des préschémas \emph{excellents} et pour
$\boldsymbol R$ une des propriétés (i) à (viii) de (7.5.3); pour la condition
$1^\circ$, cela résulte de (7.8.3, (iii)), et pour les conditions $2^\circ$ et
$3^\circ$, des raisonnements de (7.5.3), compte tenu de ce que tout anneau
excellent est caténaire. Enfin, pour la condition $4^\circ$, elle résulte,
pour (i), (ii) et (iii) de (6.7.1); pour (iv), (v), (vi) de (6.7.7), et pour
(viii) de (4.6.1); en ce qui concerne la propriété (vii) de (7.5.3), la
propriété correspondante $\boldsymbol P(Z,k)$ entraîne que $Z$ est localement
intègre (étant localement noethérien) et que chacun des sous-préschémas
intègres dont $Z$ est somme est géométriquement intègre, en vertu de (4.5.9)
et (4.6.1); on en conclut encore la propriété $4^\circ$ ci-dessus dans ce cas.

Avec ces notations et hypothèses :
\end{env}

\begin{proposition}[7.9.8]
\label{IV.2.s7.9.8-fr}
Soient $A$ un anneau local noethérien, $k$ son corps résiduel, $B$ un anneau
local noethérien tel que $\operatorname{Spec}(B)\in\mathcal C$,
$\varphi:A\to B$ un homomorphisme local faisant de $B$ un $A$-module plat;
posons $Y=\operatorname{Spec}(A)$, $X=\operatorname{Spec}(B)$ et soit
$f:X\to Y$ le morphisme correspondant à $\varphi$. On suppose que :

$1^\circ$ la propriété $\boldsymbol P(B\otimes_A k,k)$ soit vraie;

$2^\circ$ pour tout morphisme fini $Y_1\to Y$, on peut résoudre
$(Y_1)_{\mathrm{red}}$.

Alors, pour tout $y\in Y$, la propriété
$\boldsymbol P(f^{-1}(y),k(y))$ est vraie.
\end{proposition}

\begin{proof}
Notons que la condition $2^\circ$ de (7.9.8) est encore vérifiée quand on
remplace $Y$ par le spectre d'un anneau local en un idéal maximal d'une
$A$-algèbre finie $A'$ (7.9.2); d'autre part, en vertu de la condition
$1^\circ$ de (7.9.7), le spectre d'un anneau de fractions de
$B\otimes_AA'$ appartient aussi à $\mathcal C$. Le lemme (7.3.16.2) montre
alors (comme dans la partie I) du raisonnement de (7.5.2)) qu'il suffit de
prouver que, si $Y$ est intègre et si $y$ est le point générique de $Y$, les
anneaux locaux de la fibre $f^{-1}(y)$ vérifient $\boldsymbol R$.

Cela étant, il existe par hypothèse un préschéma régulier $Y'$ et un morphisme
propre et birationnel $g:Y'\to Y$; comme $Y'$ est noethérien et localement
intègre, l'hypothèse que $g$ est birationnel entraîne que $Y'$ est intègre;
soit $X'=X\times_YY'$, de sorte qu'on a un diagramme commutatif

\begin{env}[7.9.8.1]
\label{IV.2.s7.9.8.1-fr}
\[
\xymatrix{
X\ar[d]_f & X'\ar[l]_{g'}\ar[d]^{f'}\\
Y & Y'\ar[l]^g
}
\]
\end{env}

où $f'$ et $g'$ sont les projections canoniques. Compte tenu de la condition
$2^\circ$ de (7.9.7), le même raisonnement que dans le début de (7.9.3) montre
qu'il suffit de prouver

\oldpage[IV-2]{222}

que l'on a $U_{\boldsymbol R}(X')=X'$, puis, en désignant par $a$ le point
fermé de $X$, que $g'^{-1}(a)\subset U_{\boldsymbol R}(X')$. Or, soit
$x'\in g'^{-1}(a)$, et posons $z'=f'(x')$, de sorte que $b=g(z')$ est le
point fermé de $Y$. On a donc
\[
f'^{-1}(z')=f^{-1}(b)\otimes_k k(z'),
\]
et comme $k(z')$ est une extension de type fini de $k$ (puisque $g$ est
propre, donc de type fini), l'hypothèse que $\boldsymbol P(f^{-1}(b),k)$ est
vraie entraîne en vertu de la condition $4^\circ$ de (7.9.7), qu'il en est de
même de $\boldsymbol P(f'^{-1}(z'),k(z'))$. Mais comme
$\mathcal O_{Y',z'}$ est un anneau régulier et que
$\operatorname{Spec}(\mathcal O_{X',x'})\in\mathcal C$ en vertu de la
condition $1^\circ$ de (7.9.7), le lemme (7.5.1.1) prouve que
$\boldsymbol R(\mathcal O_{X',x'})$ est vraie, ce qui achève la démonstration.
\end{proof}

\begin{corollary}[7.9.9]
\label{IV.2.s7.9.9-fr}
Soient $A$, $B$ deux anneaux locaux noethériens, $\varphi:A\to B$ un
homomorphisme local, faisant de $B$ un $A$-module plat. Supposons que $A$ soit
intègre et géométriquement unibranche, et que pour toute $A$-algèbre finie
intègre $A'$ on puisse résoudre $\operatorname{Spec}(A')$ (ce qui sera par
exemple le cas si les fibres formelles de $A$ sont géométriquement régulières
et le corps résiduel de $A$ de caractéristique $0$, en raison des résultats de
Hironaka [35]). Soit $k$ le corps résiduel de $A$ et supposons que
$\operatorname{Spec}(B\otimes_Ak)$ soit géométriquement ponctuellement intègre
(4.6.9). Alors $B$ est intègre.
\end{corollary}

\begin{proof}
Nous allons appliquer (7.9.8) en prenant pour $\mathcal C$ la catégorie de
tous les préschémas localement noethériens, pour $\boldsymbol R$ la propriété
d'être intègre; les raisonnements de (7.5.3) et (7.9.7) montrent que les
conditions $1^\circ$, $2^\circ$, $3^\circ$ et $4^\circ$ de (7.9.7) sont
satisfaites dans ce cas. L'hypothèse sur $B\otimes_Ak$ et la proposition
(7.9.8) montrent donc (avec les notations de (7.9.8)) que les fibres
$f^{-1}(y)$ sont géométriquement ponctuellement intègres pour $y\in Y$;
\emph{a fortiori} les $f^{-1}(y)$ pour $y\in Y$ sont des préschémas réduits,
et comme $A$ est intègre (et \emph{a fortiori} réduit), on voit déjà que $B$
est réduit (3.3.5). Il reste à prouver que $X=\operatorname{Spec}(B)$ est
irréductible. Or, la démonstration de (7.9.8) prouve que $X'$ est localement
intègre (étant localement noethérien). Comme le morphisme $g':X'\to X$ est
surjectif, il suffit donc de voir que $X'$ est irréductible, ou (puisque $X'$
est localement intègre), que $X'$ est connexe. Mais il suffit pour cela de
prouver que la fibre $g'^{-1}(a)$ est connexe. En effet, si ce point est
établi, $X'$ ne peut être somme de deux préschémas ouverts non vides $X'_1$,
$X'_2$, car on aurait alors par exemple
$g'^{-1}(a)\cap X'_2=\varnothing$; mais la restriction $X'_2\to X$ de $g'$
étant propre (puisque $X'_2$ est un sous-préschéma fermé de $X'$),
$g'(X'_2)$ serait une partie fermée non vide de $X$ ne contenant pas le point
fermé $a$, ce qui est absurde.

Or, on a
$g'^{-1}(a)=g^{-1}(b)\otimes_{k(b)}k(a)$; mais $g:Y'\to Y$ est propre et
birationnel, donc, dans la factorisation de Stein
$Y'\xrightarrow{v}Y''\xrightarrow{u}Y$ du morphisme $g$ (\textbf{III}, 4.3.3),
le morphisme fini $u$ est aussi birationnel et par suite $Y''$ est intègre.
Comme $A$ est supposé géométriquement unibranche, il résulte alors de
(\textbf{III}, 4.3.4) que $g^{-1}(b)$ est géométriquement connexe, ce qui
achève la démonstration.
\end{proof}

\begin{remarks}[7.9.10]
\label{IV.2.s7.9.10-fr}
(i) Nous montrerons plus loin (18.8.11) que pour qu'un anneau local $B$ réduit
soit géométriquement unibranche, il faut et il suffit que pour tout morphisme
étale $h:X'\to X=\operatorname{Spec}(B)$ et tout point $x'\in X'$ au-dessus
du point fermé de $X$, $\mathcal O_{X',x'}$ soit intègre. On en déduit que si,
dans (7.9.9) on suppose, non seulement que
$\operatorname{Spec}(B\otimes_Ak)$ soit géométriquement ponctuellement
intègre, mais géométriquement unibranche, alors on peut en conclure que $B$ est
géométriquement unibranche.

\oldpage[IV-2]{223}

(ii) Soient $X$, $Y$ deux préschémas excellents, $f:X\to Y$ un morphisme
plat, et supposons que pour tout morphisme fini $Y_1\to Y$, on puisse résoudre
$(Y_1)_{\mathrm{red}}$. Il résulte de (7.9.8), en prenant pour
$\boldsymbol R$ la propriété d'être régulier, que l'ensemble $U$ des points
$x\in X$ tels que la fibre au point $f(x)$ du morphisme
\[
\operatorname{Spec}(\mathcal O_x\otimes_{\mathcal O_{f(x)}}k(f(x)))
\longrightarrow\operatorname{Spec}(k(f(x)))
\]
soit géométriquement régulière, est stable par générisation. Nous ignorons si
cet ensemble est ouvert (ou, ce qui revient au même (0$_{\mathrm{III}}$,
9.2.5), s'il est constructible) même dans le cas particulier où $Y$ est le
spectre de $\mathbf Z$ ou d'un anneau de polynômes à une indéterminée $k[T]$
sur un corps $k$, et où par suite la condition de « résolubilité » imposée à
$Y$ est trivialement satisfaite.
\end{remarks}

\begin{center}
\emph{(A suivre.)}
\end{center}

\clearpage
\oldpage[IV-2]{224}
\phantomsection
\addcontentsline{toc}{section}{Bibliographie (suite)}
\section*{BIBLIOGRAPHIE \textnormal{\emph{(suite)}}}

\begingroup
\small
\setlength{\parindent}{0pt}
\setlength{\parskip}{4pt}

\textbf{[33]} M.~\textsc{Nagata}, \emph{Note on a chain condition for prime
ideals}, \emph{Mem. Coll. Sci. Kyoto}, t.~32 (1959), p.~85--90.

\textbf{[34]} H.~\textsc{Hironaka}, article à paraître sur \emph{Les morphismes
projectifs}.

\textbf{[35]} H.~\textsc{Hironaka}, \emph{Resolution of singularities of an
algebraic variety over a field of characteristic zero}, \emph{Ann. of Math.},
t.~79 (1964), p.~109--326.

\textbf{[36]} S.~\textsc{Abhyankhar}, \emph{Local uniformization of algebraic
surfaces over ground fields of characteristic $p\ne0$}, \emph{Amer. J. of
Math.}, t.~63 (1956), p.~491--526.

\textbf{[37]} H.~\textsc{Cartan}, \emph{Séminaire de l'École Normale
Supérieure}, 13\textsuperscript{e} année (1960--61).

\textbf{[38]} J.-P.~\textsc{Serre}, \emph{Groupes algébriques et corps de
classes}, Paris (Hermann), 1959.

\endgroup

\clearpage
\oldpage[IV-2]{225}
\phantomsection
\addcontentsline{toc}{section}{Index des notations}
\section*{INDEX DES NOTATIONS}

\begingroup
\small
\setlength{\parindent}{0pt}
\setlength{\parskip}{2pt}

$\mathfrak S(X)$ ($X$ préschéma) : 2.2.15 et 4.8.1.

$\operatorname{Ass}(\mathscr F)$ ($\mathscr F$ $\mathcal O_X$-Module) : 3.1.1.

$\operatorname{deg.tr.}_K L$ ($L$ extension du corps $K$) : 4.1.

$\lambda_x(\mathscr F)$ ($\mathscr F$ $\mathcal O_X$-Module, $x$ point maximal
de $\operatorname{Supp}(\mathscr F)$) : 4.7.5.

$\Phi(\mathscr F)$ ($\mathscr F$ $\mathcal O_X$-Module) : 4.8.1.

$\operatorname{codim}(Y,X)$ ($Y$ partie quelconque d'un préschéma $X$) : 5.1.3.

$\dim(\mathscr F)$ ($\mathscr F$ $\mathcal O_X$-Module) : 5.1.12.

$\operatorname{coprof}(\mathscr F)$, $\operatorname{coprof}(X)$ ($X$
préschéma, $\mathscr F$ $\mathcal O_X$-Module) : 5.7.1.

$\operatorname{coprof}_A(M)$ ($A$ anneau noethérien, $M$ $A$-module de type
fini) : 5.7.12.

$\mathcal H^0_{X/Z}(\mathscr F)$ ($X$ préschéma, $Z$ partie de $X$ stable par
spécialisation, $\mathscr F$ $\mathcal O_X$-Module) : 5.9.1.

$\rho_{X/Z}$ ($X$ préschéma, $Z$ partie de $X$ stable par
spécialisation) : 5.9.7.

$\operatorname{prof}_T(\mathscr F)$ ($T$ partie d'un préschéma $X$, $\mathscr F$
$\mathcal O_X$-Module) : 5.10.1.

$Z^{(n)}(X)$ ($X$ préschéma, $n$ entier $\geq0$) : 5.10.13.

$A^{(1)}$ ($A$ anneau intègre noethérien) : 5.10.17 et 7.2.1.

$\operatorname{gr}_{\mathscr J}^{\bullet}(\mathscr F)$ ($\mathscr F$
$\mathcal O_X$-Module, $\mathscr J$ Idéal de $\mathcal O_X$) : 6.10.1.

$U_{S_n}(\mathscr F)$, $U_{C_n}(\mathscr F)$ ($\mathscr F$
$\mathcal O_X$-Module) : 6.11.2.

$\operatorname{CM}(\mathscr F)$ ($\mathscr F$ $\mathcal O_X$-Module) : 6.11.3.

$U_{S_n}(X)$, $U_{C_n}(X)$, $\operatorname{CM}(X)$ ($X$ préschéma) : 6.11.4.

$\operatorname{Reg}(X)$, $\operatorname{Sing}(X)$ ($X$ préschéma) : 6.12.1.

$U_{R_n}(X)$ ($X$ préschéma) : 6.12.9.

$\operatorname{Nor}(X)$ ($X$ préschéma) : 6.13.1.

$k\{T_1,\ldots,T_n\}$ : 7.4.8.

\endgroup

\clearpage
\oldpage[IV-2]{226}
\phantomsection
\addcontentsline{toc}{section}{Index terminologique}
\section*{INDEX TERMINOLOGIQUE}

\begingroup
\footnotesize
\setlength{\parindent}{0pt}
\setlength{\parskip}{1.1pt}
\newcommand{\EGAIVtwoTerm}[2]{#1 : #2.\par}

\EGAIVtwoTerm{Algèbre séparable sur un corps}{4.6.2}
\EGAIVtwoTerm{Anneau des séries formelles restreintes sur un corps valué complet}{7.4.8}
\EGAIVtwoTerm{Anneau excellent}{7.8.2}
\EGAIVtwoTerm{Anneau formellement caténaire}{7.1.9}
\EGAIVtwoTerm{Anneau formellement équidimensionnel}{7.1.1}
\EGAIVtwoTerm{Anneau géométriquement normal, géométriquement réduit,
géométriquement régulier, possédant la propriété $(R_k)$ géométrique}{6.7.6}
\EGAIVtwoTerm{Anneau normal}{5.13.5}
\EGAIVtwoTerm{$P$-anneau}{7.3.13}
\EGAIVtwoTerm{Anneau possédant la propriété $(R_k)$}{5.8.2}
\EGAIVtwoTerm{Anneau strictement équidimensionnel}{7.2.1}
\EGAIVtwoTerm{Anneau strictement formellement caténaire}{7.2.6}
\EGAIVtwoTerm{Anneau universellement caténaire}{5.6.2}
\EGAIVtwoTerm{Antifiltre de parties d'un ensemble}{5.10.8}

\medskip
\EGAIVtwoTerm{Codimension d'une partie quelconque d'un préschéma}{5.1.3}
\EGAIVtwoTerm{Coprofondeur d'un Module en un point, coprofondeur d'un Module,
coprofondeur d'un préschéma}{5.7.1}
\EGAIVtwoTerm{Corps de définition}{4.8.4}
\EGAIVtwoTerm{Corps séparablement clos}{4.3.3}
\EGAIVtwoTerm{Cycle premier associé à un Module, à un préschéma}{3.1.1}
\EGAIVtwoTerm{Cycle premier associé immergé}{3.1.1}

\medskip
\EGAIVtwoTerm{Décomposition irredondante, décomposition irredondante réduite
d'un Module}{3.2.5}
\EGAIVtwoTerm{Décomposition primaire de $0$ dans un Module}{3.2.5}
\EGAIVtwoTerm{Dimension d'un Module}{5.1.12}
\EGAIVtwoTerm{Dimension d'un préschéma sur un corps}{4.1.1}

\medskip
\EGAIVtwoTerm{Exposant d'inséparabilité d'une extension d'un corps}{4.7.4}
\EGAIVtwoTerm{Extension primaire d'un corps}{4.3.1}

\medskip
\EGAIVtwoTerm{Fermeture algébrique séparable d'un corps dans une extension}{4.3.4}
\EGAIVtwoTerm{Fibre formelle d'un anneau en un point de son spectre}{7.3.13}

\medskip
\EGAIVtwoTerm{Homéomorphisme universel}{2.4.2}
\EGAIVtwoTerm{Homomorphisme de Modules défini sur un corps}{4.8.4}

\medskip
\EGAIVtwoTerm{Lemme de Hironaka}{5.12.8}
\EGAIVtwoTerm{Longueur géométrique d'un Module en un point}{4.7.5}

\medskip
\EGAIVtwoTerm{Module intègre sur un anneau}{3.2.4}
\EGAIVtwoTerm{$A$-Module possédant la propriété $(S_k)$}{5.7.2}
\EGAIVtwoTerm{Module réduit sur un anneau}{3.2.2}
\EGAIVtwoTerm{$\mathcal O_X$-Module de Cohen-Macaulay en un point,
$\mathcal O_X$-Module de Cohen-Macaulay}{5.7.1}
\EGAIVtwoTerm{$\mathcal O_X$-Module fidèlement plat (relativement à un morphisme)}{2.2.4}
\EGAIVtwoTerm{$\mathcal O_X$-Module géométriquement intègre sur un corps}{4.6.19}

\clearpage
\oldpage[IV-2]{227}

\EGAIVtwoTerm{$\mathcal O_X$-Module géométriquement ponctuellement intègre sur un corps}{4.6.22}
\EGAIVtwoTerm{$\mathcal O_X$-Module géométriquement réduit (ou séparable) sur un corps}{4.6.17}
\EGAIVtwoTerm{$\mathcal O_X$-Module géométriquement réduit (ou séparable) sur un corps en un point}{4.6.22}
\EGAIVtwoTerm{$\mathcal O_X$-Module équidimensionnel en un point}{5.1.12}
\EGAIVtwoTerm{$\mathcal O_X$-Module intègre, intègre en un point}{3.2.4}
\EGAIVtwoTerm{$\mathcal O_X$-Module irredondant}{3.2.4}
\EGAIVtwoTerm{$\mathcal O_X$-Module normalement plat le long d'un sous-préschéma}{6.10.1}
\EGAIVtwoTerm{$\mathcal O_X$-Module possédant la propriété $(S_k)$,
possédant la propriété $(S_k)$ en un point}{5.7.2}
\EGAIVtwoTerm{$\mathcal O_X$-Module réduit, réduit en un point}{3.2.2}
\EGAIVtwoTerm{$\mathcal O_X$-Module $Z$-clos, $Z$-pur}{5.9.9}
\EGAIVtwoTerm{$\mathcal O_X$-Module $Z$-clos, $Z$-pur en un point}{5.9.9}
\EGAIVtwoTerm{Morphisme birationnel de préschémas réduits}{6.15.4}
\EGAIVtwoTerm{Morphisme connexe}{4.5.5}
\EGAIVtwoTerm{Morphisme de Cohen-Macaulay en un point, morphisme de Cohen-Macaulay}{6.8.1}
\EGAIVtwoTerm{Morphisme de coprofondeur $\leq n$ en un point, morphisme de coprofondeur $\leq n$}{6.8.1}
\EGAIVtwoTerm{Morphisme défini sur un corps}{4.8.4}
\EGAIVtwoTerm{Morphisme irréductible}{4.5.5}
\EGAIVtwoTerm{Morphisme lisse en un point, morphisme lisse}{6.8.1}
\EGAIVtwoTerm{Morphisme normal en un point, morphisme normal}{6.8.1}
\EGAIVtwoTerm{$P$-morphisme}{7.3.1}
\EGAIVtwoTerm{Morphisme possédant la propriété $(R_k)$ en un point,
morphisme possédant la propriété $(R_k)$}{6.8.1}
\EGAIVtwoTerm{Morphisme possédant la propriété $(S_k)$ en un point,
morphisme possédant la propriété $(S_k)$}{6.8.1}
\EGAIVtwoTerm{Morphisme quasi-fidèlement plat, morphisme quasi-plat}{2.3.3}
\EGAIVtwoTerm{Morphisme radiciel en un point}{6.15.3}
\EGAIVtwoTerm{Morphisme réduit en un point, morphisme réduit}{6.8.1}
\EGAIVtwoTerm{Morphisme régulier en un point, morphisme régulier}{6.8.1}
\EGAIVtwoTerm{Morphisme résolvant}{7.9.1}
\EGAIVtwoTerm{Morphisme universellement ouvert, universellement bicontinu}{2.4.2}
\EGAIVtwoTerm{Multiplicité radicielle, multiplicité séparable, multiplicité
totale d'une extension d'un corps}{4.7.4}
\EGAIVtwoTerm{Multiplicité radicielle, multiplicité séparable, multiplicité
totale d'un préschéma sur un corps}{4.7.4}
\EGAIVtwoTerm{Multiplicité radicielle d'un point pour un $\mathcal O_X$-Module
($X$ préschéma sur un corps)}{4.7.5}
\EGAIVtwoTerm{Multiplicité radicielle d'un cycle premier maximal associé à un
$\mathcal O_X$-Module ($X$ préschéma sur un corps)}{4.7.8}
\EGAIVtwoTerm{Multiplicité totale d'un point (d'un cycle premier maximal
associé) pour un $\mathcal O_X$-Module ($X$ préschéma sur un corps)}{4.7.12}

\medskip
\EGAIVtwoTerm{Nombre géométrique de composantes connexes, de composantes
irréductibles, pour un préschéma sur un corps}{4.5.2}

\medskip
\EGAIVtwoTerm{Point associé à un $\mathcal O_X$-Module, à un préschéma}{3.1.1}
\EGAIVtwoTerm{Point de Cohen-Macaulay d'un préschéma}{5.7.1}
\EGAIVtwoTerm{Point géométriquement unibranche, point unibranche d'un préschéma}{6.15.1}
\EGAIVtwoTerm{Préschéma connexe en codimension $d-1$}{5.10.8}
\EGAIVtwoTerm{Préschéma de Cohen-Macaulay}{5.7.1}
\EGAIVtwoTerm{Préschéma excellent}{7.8.5}
\EGAIVtwoTerm{Préschéma (sur un corps) géométriquement connexe}{4.5.2}
\EGAIVtwoTerm{Préschéma (sur un corps) géométriquement intègre}{4.6.2}
\EGAIVtwoTerm{Préschéma (sur un corps) géométriquement ponctuellement intègre}{4.6.9}
\EGAIVtwoTerm{Préschéma (sur un corps) géométriquement irréductible}{4.5.2}
\EGAIVtwoTerm{Préschéma (sur un corps) géométriquement normal en un point,
géométriquement normal}{6.7.6}
\EGAIVtwoTerm{Préschéma (sur un corps) géométriquement réduit}{4.6.2}
\EGAIVtwoTerm{Préschéma (sur un corps) géométriquement réduit en un point}{4.6.9}
\EGAIVtwoTerm{Préschéma (sur un corps) géométriquement régulier en un point,
géométriquement régulier}{6.7.6}

\clearpage
\oldpage[IV-2]{228}

\EGAIVtwoTerm{Préschéma géométriquement unibranche en un point,
géométriquement unibranche}{6.15.1}
\EGAIVtwoTerm{Préschéma intègre en un point}{3.2.4}
\EGAIVtwoTerm{Préschéma irredondant}{3.2.4}
\EGAIVtwoTerm{Préschéma localement immersible dans un schéma régulier}{5.11.1}
\EGAIVtwoTerm{Préschéma ponctuellement intègre}{4.6.9}
\EGAIVtwoTerm{Préschéma possédant la propriété $(R_k)$, possédant la propriété
$(R_k)$ en un point}{5.8.2}
\EGAIVtwoTerm{Préschéma (sur un corps) possédant la propriété $(R_k)$
géométrique en un point, possédant la propriété $(R_k)$ géométrique}{6.7.6}
\EGAIVtwoTerm{Préschéma possédant la propriété $(S_k)$, possédant la propriété
$(S_k)$ en un point}{5.7.2}
\EGAIVtwoTerm{Préschéma régulier en codimension $k$}{5.8.2}
\EGAIVtwoTerm{Préschéma (sur un corps) séparable}{4.6.2}
\EGAIVtwoTerm{Préschéma (sur un corps) séparable en un point}{4.6.9}
\EGAIVtwoTerm{Préschéma unibranche en un point, préschéma unibranche}{6.15.1}
\EGAIVtwoTerm{Préschéma universellement caténaire}{5.6.3}
\EGAIVtwoTerm{Préschéma (sur un corps) universellement réduit}{4.6.2}
\EGAIVtwoTerm{Profondeur d'un $\mathcal O_X$-Module en un point}{5.7.1}
\EGAIVtwoTerm{Propriété $P$ du premier type (du second type) associée à une
propriété $R$ d'anneaux locaux noethériens}{7.3.10}
\EGAIVtwoTerm{Propriété $P$ géométrique}{7.3.6}

\medskip
\EGAIVtwoTerm{Résolution des singularités}{7.9.1}

\medskip
\EGAIVtwoTerm{Sous-ensemble (d'un préschéma) défini sur un corps}{4.8.4}
\EGAIVtwoTerm{Sous-$\mathcal O_X$-Module défini sur un corps}{4.8.4}
\EGAIVtwoTerm{Sous-$\mathcal O_X$-Module primaire d'un $\mathcal O_X$-Module}{3.2.4}
\EGAIVtwoTerm{Sous-préschéma défini sur un corps}{4.8.4}
\EGAIVtwoTerm{Sous-préschéma fermé primaire d'un préschéma}{3.2.4}

\medskip
\EGAIVtwoTerm{$Z$-clôture d'un $\mathcal O_X$-Module}{5.9.11}

\endgroup

\clearpage
\oldpage[IV-2]{229}
\phantomsection
\addcontentsline{toc}{section}{Table des matières}
\section*{TABLE DES MATIÈRES}

\begingroup
\small
\setlength{\parindent}{0pt}
\setlength{\parskip}{1.5pt}
\newcommand{\EGAIVtwoContentsLine}[3]{\textbf{#1}\quad #2\dotfill #3\par}

\EGAIVtwoContentsLine{CHAPITRE IV.}{\textbf{Étude locale des schémas et des
morphismes des schémas} \emph{(suite)}}{5}

\EGAIVtwoContentsLine{\S~2.}{Changement de base et platitude}{5}
\EGAIVtwoContentsLine{2.1.}{Modules plats sur les préschémas}{5}
\EGAIVtwoContentsLine{2.2.}{Modules fidèlement plats sur les préschémas}{9}
\EGAIVtwoContentsLine{2.3.}{Propriétés topologiques des morphismes plats}{14}
\EGAIVtwoContentsLine{2.4.}{Morphismes universellement ouverts et morphismes plats}{19}
\EGAIVtwoContentsLine{2.5.}{Permanence des propriétés des Modules par descente fidèlement plate}{22}
\EGAIVtwoContentsLine{2.6.}{Permanence de propriétés ensemblistes et topologiques de morphismes par descente fidèlement plate}{27}
\EGAIVtwoContentsLine{2.7.}{Permanence de diverses propriétés des morphismes par descente fidèlement plate}{29}
\EGAIVtwoContentsLine{2.8.}{Préschémas sur une base régulière de dimension 1; adhérence d'un sous-préschéma fermé de la fibre générique}{33}

\EGAIVtwoContentsLine{\S~3.}{Cycles premiers associés et décompositions primaires}{36}
\EGAIVtwoContentsLine{3.1.}{Cycles premiers associés à un Module}{36}
\EGAIVtwoContentsLine{3.2.}{Décompositions irredondantes}{40}
\EGAIVtwoContentsLine{3.3.}{Relations avec la platitude}{43}
\EGAIVtwoContentsLine{3.4.}{Propriétés des faisceaux $\mathscr F/t\mathscr F$}{46}

\EGAIVtwoContentsLine{\S~4.}{Changement du corps de base dans les préschémas algébriques}{52}
\EGAIVtwoContentsLine{4.1.}{Dimension des préschémas algébriques}{52}
\EGAIVtwoContentsLine{4.2.}{Cycles premiers associés sur les préschémas algébriques}{54}
\EGAIVtwoContentsLine{4.3.}{Rappels sur les produits tensoriels de corps}{58}
\EGAIVtwoContentsLine{4.4.}{Préschémas irréductibles et préschémas connexes sur un corps algébriquement clos}{59}
\EGAIVtwoContentsLine{4.5.}{Préschémas géométriquement irréductibles et géométriquement connexes}{61}
\EGAIVtwoContentsLine{4.6.}{Préschémas algébriques géométriquement réduits}{68}

\clearpage
\oldpage[IV-2]{230}

\EGAIVtwoContentsLine{4.7.}{Multiplicités dans la décomposition primaire sur un préschéma algébrique}{75}
\EGAIVtwoContentsLine{4.8.}{Corps de définition}{80}
\EGAIVtwoContentsLine{4.9.}{Corps de définition d'une partie d'un préschéma}{84}

\EGAIVtwoContentsLine{\S~5.}{Dimension, profondeur, régularité dans les préschémas localement noethériens}{86}
\EGAIVtwoContentsLine{5.1.}{Dimension des préschémas}{86}
\EGAIVtwoContentsLine{5.2.}{Dimension d'un préschéma algébrique}{90}
\EGAIVtwoContentsLine{5.3.}{Dimension du support d'un Module et polynôme de Hilbert}{92}
\EGAIVtwoContentsLine{5.4.}{Dimension de l'image d'un morphisme}{93}
\EGAIVtwoContentsLine{5.5.}{Formule des dimensions pour un morphisme de type fini}{94}
\EGAIVtwoContentsLine{5.6.}{Formule des dimensions et anneaux universellement caténaires}{97}
\EGAIVtwoContentsLine{5.7.}{Profondeur et propriété $(S_k)$}{103}
\EGAIVtwoContentsLine{5.8.}{Préschémas réguliers et propriété $(R_k)$. Critère de normalité de Serre}{107}
\EGAIVtwoContentsLine{5.9.}{Modules $Z$-purs et $Z$-clos}{109}
\EGAIVtwoContentsLine{5.10.}{Propriété $(S_2)$ et $Z$-clôture}{114}
\EGAIVtwoContentsLine{5.11.}{Critères de cohérence pour les Modules $\mathcal H^0_{X/Z}(\mathscr F)$}{122}
\EGAIVtwoContentsLine{5.12.}{Relations entre les propriétés d'un anneau local noethérien $A$ et d'un anneau quotient $A/tA$}{126}
\EGAIVtwoContentsLine{5.13.}{Propriétés de permanence par passage à la limite inductive}{131}

\EGAIVtwoContentsLine{\S~6.}{Morphismes plats de préschémas localement noethériens}{134}
\EGAIVtwoContentsLine{6.1.}{Platitude et dimension}{135}
\EGAIVtwoContentsLine{6.2.}{Platitude et dimension projective}{137}
\EGAIVtwoContentsLine{6.3.}{Platitude et profondeur}{138}
\EGAIVtwoContentsLine{6.4.}{Platitude et propriété $(S_k)$}{141}
\EGAIVtwoContentsLine{6.5.}{Platitude et propriété $(R_k)$}{143}
\EGAIVtwoContentsLine{6.6.}{Propriétés de transitivité}{145}
\EGAIVtwoContentsLine{6.7.}{Application aux changements de base dans les préschémas algébriques}{145}
\EGAIVtwoContentsLine{6.8.}{Morphismes réguliers, normaux, réduits, lisses}{150}
\EGAIVtwoContentsLine{6.9.}{Le théorème de platitude générique}{153}
\EGAIVtwoContentsLine{6.10.}{Dimension et profondeur d'un Module normalement plat le long d'un sous-préschéma fermé}{155}
\EGAIVtwoContentsLine{6.11.}{Critères pour que les ensembles $U_{S_n}(\mathscr F)$ ou $U_{C_n}(\mathscr F)$ soient ouverts}{158}
\EGAIVtwoContentsLine{6.12.}{Critères de Nagata pour que $\operatorname{Reg}(X)$ soit ouvert}{163}
\EGAIVtwoContentsLine{6.13.}{Critères pour que $\operatorname{Nor}(X)$ soit ouvert}{168}
\EGAIVtwoContentsLine{6.14.}{Changement de base et clôture intégrale}{169}
\EGAIVtwoContentsLine{6.15.}{Préschémas géométriquement unibranches}{176}

\clearpage
\oldpage[IV-2]{231}

\EGAIVtwoContentsLine{\S~7.}{Relations entre un anneau local noethérien et son complété. Anneaux excellents}{182}
\EGAIVtwoContentsLine{7.1.}{Équidimensionalité formelle et anneaux formellement caténaires}{183}
\EGAIVtwoContentsLine{7.2.}{Anneaux strictement formellement caténaires}{187}
\EGAIVtwoContentsLine{7.3.}{Fibres formelles des anneaux locaux noethériens}{192}
\EGAIVtwoContentsLine{7.4.}{Permanence des propriétés des fibres formelles}{198}
\EGAIVtwoContentsLine{7.5.}{Un critère pour les $P$-morphismes}{203}
\EGAIVtwoContentsLine{7.6.}{Applications : I. Anneaux japonais locaux}{208}
\EGAIVtwoContentsLine{7.7.}{Applications : II. Anneaux universellement japonais}{212}
\EGAIVtwoContentsLine{7.8.}{Anneaux excellents}{214}
\EGAIVtwoContentsLine{7.9.}{Anneaux excellents et résolution des singularités}{218}
\EGAIVtwoContentsLine{}{BIBLIOGRAPHIE \emph{(suite)}}{224}
\EGAIVtwoContentsLine{}{INDEX DES NOTATIONS}{225}
\EGAIVtwoContentsLine{}{INDEX TERMINOLOGIQUE}{226}

\vspace{2em}
\begin{center}
\emph{Manuscrit reçu le 15 décembre 1963.}
\end{center}

\endgroup

% FIN SOURCE ega4/ega4-2-fr.tex

\end{document}
% EGA-COMPONENT-END IV2

% EGA-COMPONENT-BEGIN IV3 D2D39D659B73DA0550F43F59D7A3E972E90E95615B8DBE0C71A280FDF54984BE
\documentclass[10pt,oneside]{book}

\usepackage[T1]{fontenc}
\usepackage[utf8]{inputenc}
\usepackage[french]{babel}
\usepackage{amsmath,amssymb,mathrsfs}
\usepackage[all]{xy}
\usepackage{tikz-cd}
\usepackage[shortlabels]{enumitem}
\usepackage[a4paper,margin=22mm]{geometry}
\usepackage[
  colorlinks=true,
  linkcolor=blue,
  hypertexnames=false,
  unicode=true,
  pdfencoding=auto
]{hyperref}
\hypersetup{
  pdftitle={Éléments de géométrie algébrique IV, troisième partie -- édition française},
  pdfauthor={Alexander Grothendieck et Jean Dieudonné},
  pdfsubject={Transcription diplomatique française depuis l'autorité NUMDAM}
}

\renewcommand{\thesection}{\arabic{section}}
\renewcommand{\thesubsection}{\thesection.\arabic{subsection}}
\newcommand{\oldpage}[2][]{%
  \setcounter{footnote}{0}%
  \marginpar{\scriptsize\textbf{#1}~|~#2}\ignorespaces}
\newenvironment{env}[1][]{\par\medskip\noindent\textbf{(#1)}\ }{\par}
\newenvironment{definition}[1][]{%
  \par\medskip\noindent\itshape Définition (#1). ---\ }{\par}
\newenvironment{proposition}[1][]{%
  \par\medskip\noindent\itshape Proposition (#1). ---\ }{\par}
\newenvironment{theorem}[1][]{%
  \par\medskip\noindent\itshape Théorème (#1). ---\ }{\par}
\newenvironment{corollary}[1][]{%
  \par\medskip\noindent\itshape Corollaire (#1). ---\ }{\par}
\newenvironment{remark}[1][]{%
  \par\medskip\noindent\itshape Remarque (#1). ---\ }{\par}
\newenvironment{remarks}[1][]{%
  \par\medskip\noindent\itshape Remarques (#1). ---\ }{\par}
\newenvironment{lemma}[1][]{%
  \par\medskip\noindent\itshape Lemme (#1). ---\ }{\par}
\newenvironment{proof}{%
  \par\medskip\noindent\textit{Démonstration. ---\ }}{\par}

\makeatletter
\newcommand{\EGAReaderTarget}[2]{\hypertarget{#1}{}\@bsphack\begingroup\def\@currentHref{#1}\label{#2}\endgroup\@esphack}
\makeatother
\begin{document}
% BEGIN RIGHTS AND AI TRANSCRIPTION NOTICE
\clearpage
\begingroup
\small
\noindent\textbf{Droits, disponibilité et provenance}

Les droits sur les œuvres originales et sur toute contribution protégée à ces traductions et transcriptions dérivées restent ceux de leurs titulaires respectifs. Ce projet ne place pas ces œuvres ou contributions dans le domaine public et n'accorde aucune licence sur elles. Les mentions de droits de tiers et les autorisations déjà accordées restent applicables.

Ces éditions sont mises à disposition gratuitement, à des fins non commerciales et exclusivement pédagogiques, dans la présomption que les titulaires de droits ne s'opposeraient pas à la libre disponibilité de ces traductions et transcriptions aux côtés des originaux français sur NUMDAM (www.numdam.org/journals/PMIHES/). Cette présomption ne signifie pas qu'une autorisation a été obtenue ; le libre accès aux originaux n'autorise pas, à lui seul, leur traduction, leur adaptation ou leur republication.

Si vous détenez des droits sur un élément présenté ici et vous opposez à sa disponibilité, veuillez ouvrir une demande sur GitHub (github.com/KokunoYumeto/ega-fr/issues). Le responsable de la publication retirera l'élément concerné à la demande du titulaire des droits.

La transcription moderne, les corrections, la composition et les notes mathématiques sont produites par IA. Les travaux documentés comprennent ceux d'OpenAI Codex, GPT-6 Astra, effort Ultra. Les indications de provenance des contributions antérieures restent applicables. Les textes originaux restent l'œuvre d'Alexandre Grothendieck et Jean Dieudonné ; ils ne sont pas attribués à l'IA. Cette édition n'est approuvée ni par les auteurs, ni par l'IHÉS, ni par NUMDAM.
\par
\endgroup
\clearpage
% END RIGHTS AND AI TRANSCRIPTION NOTICE

%
% DEBUT SOURCE ega4/ega4-8-fr.tex
\clearpage
\thispagestyle{plain}
\oldpage[IV-3]{5}
\phantomsection
\label{chapter:IV3-fr}

\begin{center}
{\large CHAPITRE IV (suite)\par}
\vspace{1.5em}
{\Large\bfseries ÉTUDE LOCALE DES SCHÉMAS\\
ET DES MORPHISMES DE SCHÉMAS\par}
\end{center}

\setcounter{section}{7}
\section{Limites projectives de préschémas}
\label{section:IV.8-fr}

\setcounter{subsection}{0}
\subsection{Introduction}
\label{subsection:IV.8.1-fr}

\begin{env}[8.1.1]
\label{IV.8.1.1-fr}
Dans ce paragraphe, nous allons étudier systématiquement la situation
suivante. Soient $I$ un ensemble préordonné filtrant croissant,
$(A_\lambda,\varphi_{\mu\lambda})$ un système inductif d'anneaux ayant $I$
pour ensemble d'indices, $A=\varinjlim A_\lambda$ sa limite inductive. Pour
tout $\alpha\in I$ et tout $A_\alpha$-préschéma $X_\alpha$, considérons les
$A_\lambda$-préschémas
$X_\lambda=X_\alpha\otimes_{A_\alpha}A_\lambda$ pour $\lambda\geq\alpha$,
et le $A$-préschéma $X=X_\alpha\otimes_{A_\alpha}A$~; il est clair que les
préschémas $X_\lambda$ (pour $\lambda\geq\alpha$) forment un \emph{système
projectif}, et on verra (8.2.5) que $X$ est une \emph{limite projective} de ce
système dans la catégorie des préschémas. Nous nous proposons de trouver des
conditions sur $X_\alpha$ ou sur les $A_\lambda$ permettant de démontrer des
propriétés du type suivant~: pour que $X$ possède une propriété $P$ (par
exemple, la propriété d'être propre sur $S=\operatorname{Spec}(A)$, ou
irréductible, ou connexe, etc.), il faut et il suffit qu'il existe un indice
$\mu\geq\alpha$ tel que, pour tout $\lambda\geq\mu$, $X_\lambda$ ait (par
rapport à $S_\lambda=\operatorname{Spec}(A_\lambda)$, le cas échéant), la
même propriété $P$. Nous obtiendrons des énoncés analogues pour des propriétés
de $\mathcal O_X$-Modules, de $A$-morphismes de $A$-préschémas, etc. Nous
montrerons également (8.9.1) que la donnée d'un $A$-préschéma de présentation
finie (1.6.1) équivaut essentiellement à la donnée d'un $A_\lambda$-préschéma
de présentation finie $X_\lambda$ pour un $\lambda$ assez grand, $X$ étant
alors \emph{isomorphe à} $X_\lambda\otimes_{A_\lambda}A$. On a aussi des
énoncés analogues pour les $\mathcal O_X$-Modules de présentation finie,
leurs homomorphismes, les $A$-morphismes de $A$-préschémas de présentation
finie, etc.
\end{env}

\begin{env}[8.1.2]
\EGAReaderTarget{ega4.c04.statement.8-1-2}{IV.8.1.2-fr}
L'utilité de tels résultats apparaîtra par exemple dans les questions
suivantes~:
\begin{enumerate}
\item[\emph{a)}] Soient $Y$ un préschéma, $y$ un point de $Y$, $(U_\lambda)$
le système projectif filtrant décroissant des voisinages ouverts affines de
$y$ dans $Y$~; si $A_\lambda$ est l'anneau de $U_\lambda$, les $A_\lambda$
forment un système inductif filtrant croissant, dont la limite inductive $A$
est l'anneau local $\mathcal O_y$~; d'ailleurs, si on désigne par
$\mathfrak p_\lambda$ l'idéal premier $\mathfrak j_y$ dans l'anneau
$A_\lambda$, le système inductif $(A_\lambda)$ est cofinal à tout système
inductif $(A_\alpha)_f$, où $f$ parcourt
$A_\alpha-\mathfrak p_\alpha$ (pour un $\alpha$ \emph{fixé}), car les $D(f)$
forment un système fondamental de voisinages de $y$ dans $U_\alpha$,
\oldpage[IV-3]{6}
donc dans $Y$. Les résultats du présent paragraphe impliqueront que la
géométrie algébrique des $\mathcal O_y$-préschémas de présentation finie (et
la théorie des Modules de présentation finie sur ces préschémas) est
essentiellement équivalente à la géométrie algébrique des préschémas de
présentation finie sur des voisinages ouverts «~assez petits~» de $y$. Ainsi,
l'énoncé (8.10.5, (xii)) implique que si un morphisme $f:X\to Y$ est de
présentation finie, alors, pour que
$X\times_Y\operatorname{Spec}(\mathcal O_y)$ soit propre sur
$\operatorname{Spec}(\mathcal O_y)$, il faut et il suffit qu'il existe un
voisinage ouvert $U$ de $y$ dans $Y$ tel que $f^{-1}(U)$ soit propre sur $U$.

Un cas particulièrement important, et dans une certaine mesure classique,
est celui où $Y$ est intègre et où $y$ est son point générique, de sorte que
$\mathcal O_y$ n'est autre que \emph{le corps $R(Y)=K$ des fonctions
rationnelles sur $Y$}. Les résultats du présent paragraphe reviennent alors à
interpréter la géométrie algébrique sur $K$ en termes de la géométrie
algébrique au-dessus d'ouverts non vides «~assez petits~» de $Y$, c'est-à-dire,
intuitivement, en termes de «~familles~» d'objets géométriques indexées par les
points d'un tel ouvert. Ce point de vue a d'ailleurs été couramment utilisé
depuis longtemps, non seulement en Géométrie algébrique sur des corps
algébriquement clos, mais aussi dans l'étude arithmétique des variétés définies
sur un \emph{corps de nombres} $K$ (extension finie de $\mathbf Q$), en
considérant ce dernier comme le corps des fractions de son anneau des entiers
$A$ («~théorie de la \emph{réduction modulo} $\mathfrak p$~» ($\mathfrak p$
idéal premier de $A$)~; cf. (I, 3.7)). Les résultats des \S\S~8 et 9
fournissent donc entre autres des fondements du langage de la «~réduction
modulo $\mathfrak p$~» en arithmétique.

On notera que dans l'exemple envisagé ici, les morphismes
$S_\mu\to S_\lambda$ (pour $\lambda\leq\mu$) sont les \emph{immersions
ouvertes} canoniques $U_\mu\to U_\lambda$, et \emph{a fortiori} sont des
morphismes \emph{plats} (mais non fidèlement plats en général), ce qui
explique l'intérêt des énoncés qui font appel à une telle restriction.

\item[\emph{b)}] Supposons que les $A_\lambda$ soient des \emph{corps}, de
sorte que $A=\varinjlim A_\lambda$ est aussi un corps. Ce cas se présente
généralement lorsqu'on part de données géométriques au-dessus d'un corps
quelconque $K$, que l'on considère comme extension d'un corps $k$ (par exemple
le sous-corps premier de $K$). Il y a alors intérêt à considérer $K$ comme
limite inductive de ses sous-extensions qui sont \emph{de type fini sur $k$},
ce qui permet dans de nombreuses questions de se ramener au cas où $K$ est
une extension de type fini de $k$. Utilisant aussi la méthode esquissée dans
\emph{a)}, on peut alors se ramener généralement au cas d'un anneau de base
$A$ qui est une \emph{algèbre intègre de type fini sur $k$}.

On notera que dans cet exemple, les morphismes $S_\mu\to S_\lambda$ sont
\emph{fidèlement plats}.

\item[\emph{c)}] Supposons qu'on s'intéresse aux propriétés, locales sur
$Y$, des préschémas de présentation finie au-dessus d'un préschéma quelconque
$Y$, que l'on peut dès lors supposer affine d'anneau $A$. Il y a alors intérêt
à considérer $A$ comme limite inductive de ses sous-anneaux qui sont des
$\mathbf Z$-algèbres \emph{de type fini}, ce qui permet de ramener de
nombreuses questions au cas où $Y$ est le spectre d'une telle algèbre. C'est
là l'explicitation du «~point de vue kroneckerien~» suivant lequel la
Géométrie algébrique se réduit à la Géométrie algébrique des préschémas de
type fini sur $\mathbf Z$ (que l'on qualifie parfois de
\oldpage[IV-3]{7}
«~Géométrie algébrique absolue~»). Cet exemple nous montre en particulier que
dans la plupart des questions «~relatives~» sur un préschéma de base $Y$, on
peut se ramener au cas où $Y$ est \emph{noethérien}.

On notera que dans cet exemple, au contraire des précédents, les morphismes
$S_\mu\to S_\lambda$ n'ont en général aucune propriété particulière de
régularité.
\end{enumerate}

Par la suite, quand nous appliquerons les résultats qui vont suivre à l'une
quelconque des trois situations particulières qui viennent d'être décrites,
nous nous dispenserons de reprendre en détail le procédé qui permet de faire
ces applications, nous contentant de renvoyer à ce qui précède.
\end{env}

\begin{env}[8.1.3]
\label{IV.8.1.3-fr}
Dans l'exemple \emph{a)} de (8.1.2), nous avons vu que si $Y$ est un
préschéma intègre de point générique $y$, $f:X\to Y$ un morphisme de
présentation finie, alors, si la fibre générique $f^{-1}(y)$ est propre sur
$\boldsymbol k(y)$, il y a un voisinage ouvert $U$ de $y$ tel que $f^{-1}(U)$
soit propre sur $U$~; il en résulte \emph{a fortiori} que pour tout $s\in U$,
$f^{-1}(s)$ est propre sur $\boldsymbol k(s)$. Il arrive qu'on ait besoin
d'une réciproque, assurant que si $f^{-1}(s)$ est propre sur
$\boldsymbol k(s)$ pour «~suffisamment~» de points $s\in U$, alors
$f^{-1}(y)$ est propre sur $\boldsymbol k(y)$. Par exemple, supposons que $X$
et $Y$ soient des préschémas algébriques sur un corps $k$ algébriquement clos
(on peut prendre pour $k$ le corps $\mathbf C$ des nombres complexes, pour
fixer les idées)~; on a parfois besoin de savoir que si, pour tout $s\in Y$,
\emph{rationnel sur $k$}, la fibre $f^{-1}(s)$ est propre sur
$\boldsymbol k(s)$, alors $f^{-1}(y)$ est propre sur $\boldsymbol k(y)$, et
par suite $f^{-1}(U)$ est propre sur $U$ pour un voisinage $U$ de
$y$\footnote{On notera qu'un tel énoncé est en fin de compte purement
géométrique, en ce sens qu'il ne fait appel qu'à des points rationnels sur
$k$, et non à des points génériques~; par exemple, lorsque $k=\mathbf C$, cet
énoncé a une signification topologique évidente pour l'analyste, en
interprétant «~propre~» au sens topologique du terme, pour les espaces
sous-jacents aux espaces analytiques formés des points de $X$ et $Y$
rationnels sur $\mathbf C$.}. Or cet énoncé résultera aisément du suivant~:
l'ensemble $E$ des points $s\in Y$ tels que $f^{-1}(s)$ soit propre sur
$\boldsymbol k(s)$ est \emph{constructible} (et par conséquent identique à
$Y$ tout entier s'il contient les points fermés de $Y$, grâce au théorème des
zéros de Hilbert (10.4.8))~; cela équivaut encore à dire que si $f^{-1}(y)$
\emph{n'est pas propre} sur $\boldsymbol k(y)$, alors il existe un voisinage
ouvert $U$ de $y$ tel que $f^{-1}(s)$ ne soit pas propre sur
$\boldsymbol k(s)$ pour tout $s\in U$ (cf. (9.6.1, (iv))). Cet exemple
illustre l'intérêt qu'il y a à développer systématiquement des \emph{critères
de constructibilité} pour les notions les plus importantes~: c'est ce qui
sera fait au \S~9.
\end{env}

\subsection{Limites projectives de préschémas}
\label{subsection:IV.8.2-fr}

\begin{env}[8.2.1]
\label{IV.8.2.1-fr}
Soient $S_0$ un espace annelé, $L$ un ensemble préordonné filtrant croissant,
$(\mathscr A_\lambda,\varphi_{\mu\lambda})$ un système inductif de
$\mathcal O_{S_0}$-Algèbres (non nécessairement commutatives), ayant $L$ pour
ensemble d'indices. On sait que, considéré comme système inductif de
$\mathcal O_{S_0}$-Modules, $(\mathscr A_\lambda,\varphi_{\mu\lambda})$ admet
une limite inductive $\mathscr A$~; désignons par
$\varphi_\lambda:\mathscr A_\lambda\to\mathscr A$ l'homomorphisme canonique
(de $\mathcal O_{S_0}$-Modules). Soit
$m_\lambda:\mathscr A_\lambda\otimes\mathscr A_\lambda\to\mathscr A_\lambda$
l'homomorphisme de $\mathcal O_{S_0}$-Modules qui définit la multiplication
dans la $\mathcal O_{S_0}$-Algèbre $\mathscr A_\lambda$~; l'hypothèse sur les
$\varphi_{\mu\lambda}$ entraîne que les $m_\lambda$ forment un système
inductif d'homomorphismes,
\oldpage[IV-3]{8}
et comme le foncteur $\varinjlim$ commute au produit tensoriel,
$m=\varinjlim m_\lambda$ est un homomorphisme
$\mathscr A\otimes\mathscr A\to\mathscr A$ de
$\mathcal O_{S_0}$-Modules~; par passage à la limite sur les diagrammes
commutatifs exprimant l'associativité de $m_\lambda$ et l'existence de la
section unité dans $\mathscr A_\lambda$, on voit que $m$ définit sur
$\mathscr A$ une structure de $\mathcal O_{S_0}$-Algèbre et que
$\varphi_\lambda$ est un homomorphisme de $\mathcal O_{S_0}$-Algèbres pour
tout $\lambda\in L$. En outre $\mathscr A$ est limite inductive du système
$(\mathscr A_\lambda,\varphi_{\mu\lambda})$ \emph{dans la catégorie des
$\mathcal O_{S_0}$-Algèbres}, autrement dit, pour toute
$\mathcal O_{S_0}$-Algèbre $\mathscr B$, l'application canonique
\begin{equation}
\operatorname{Hom}_{\mathcal O_{S_0}\text{-Alg.}}(\mathscr A,\mathscr B)
\longrightarrow
\varprojlim_\lambda
\operatorname{Hom}_{\mathcal O_{S_0}\text{-Alg.}}
(\mathscr A_\lambda,\mathscr B)
\tag{8.2.1.1}\label{IV.8.2.1.1-fr}
\end{equation}
qui à tout homomorphisme $f:\mathscr A\to\mathscr B$ de
$\mathcal O_{S_0}$-Algèbres, fait correspondre la famille
$(f\circ\varphi_\lambda)$, est \emph{bijective}. En effet, on sait déjà
qu'elle est injective et identifie
$\operatorname{Hom}_{\mathcal O_{S_0}\text{-Alg.}}(\mathscr A,\mathscr B)$
à une partie de
$\varprojlim_\lambda
\operatorname{Hom}_{\mathcal O_{S_0}\text{-Mod.}}
(\mathscr A_\lambda,\mathscr B)$~; tout revient à voir que si $(f_\lambda)$
est un système inductif d'homomorphismes de $\mathcal O_{S_0}$-Algèbres,
$f_\lambda:\mathscr A_\lambda\to\mathscr B$, sa limite inductive
$f:\mathscr A\to\mathscr B$, qui est par définition un homomorphisme de
$\mathcal O_{S_0}$-\emph{Modules}, est aussi un homomorphisme de
$\mathcal O_{S_0}$-Algèbres~; mais cela résulte du passage à la limite
inductive dans le diagramme commutatif d'homomorphismes de
$\mathcal O_{S_0}$-Modules
\[
\xymatrix{
  \mathscr A_\lambda\otimes\mathscr A_\lambda
    \ar[r]^{f_\lambda\otimes f_\lambda} \ar[d]_{m_\lambda}
    & \mathscr B\otimes\mathscr B \ar[d] \\
  \mathscr A_\lambda \ar[r]_{f_\lambda} & \mathscr B
}
\]
et du fait que le foncteur $\varinjlim$ commute aux produits tensoriels.

On notera enfin que si les $\mathscr A_\lambda$ sont des
$\mathcal O_{S_0}$-Algèbres commutatives, il en est de même de $\mathscr A$.
\end{env}

\begin{env}[8.2.2]
\label{IV.8.2.2-fr}
Supposons maintenant que $S_0$ soit un \emph{préschéma}, et que les
$\mathscr A_\lambda$ soient des $\mathcal O_{S_0}$-Algèbres (commutatives)
\emph{quasi-cohérentes}~; on sait alors que
$\mathscr A=\varinjlim\mathscr A_\lambda$ est une
$\mathcal O_{S_0}$-Algèbre quasi-cohérente (I, 4.1.1). Désignons par
$S_\lambda$ (resp. $S$) le spectre de la $\mathcal O_{S_0}$-Algèbre
$\mathscr A_\lambda$ (resp. $\mathscr A$) (II, 1.3.1), et soient
$u_{\lambda\mu}:S_\mu\to S_\lambda$ (pour $\lambda\leq\mu$) et
$u_\lambda:S\to S_\lambda$ les $S_0$-morphismes correspondant aux
homomorphismes $\varphi_{\mu\lambda}$ et $\varphi_\lambda$ respectivement
(II, 1.2.7)~; il est clair que $(S_\lambda,u_{\lambda\mu})$ est un
\emph{système projectif} dans la catégorie des $S_0$-préschémas. On notera
que les $u_{\lambda\mu}$ et $u_\lambda$ sont des morphismes \emph{affines}
(II, 1.6.2), donc \emph{quasi-compacts} et \emph{séparés}.
\end{env}

\begin{proposition}[8.2.3]
\label{IV.8.2.3-fr}
Avec les notations de (8.2.2), les morphismes $u_\lambda:S\to S_\lambda$
font de $S$ une limite projective du système projectif
$(S_\lambda,u_{\lambda\mu})$ dans la catégorie des préschémas. En outre, si
$h:S_0\to T$ est un morphisme, faisant de tout $S_0$-préschéma un
$T$-préschéma, $S$ est aussi limite projective du système
$(S_\lambda,u_{\lambda\mu})$ dans la catégorie des $T$-préschémas.
\end{proposition}

Prouvons d'abord la seconde assertion de l'énoncé dans le cas $T=S_0$.
\oldpage[IV-3]{9}
Tout revient à démontrer que si $X$ est un $S_0$-préschéma quelconque,
l'application canonique
\begin{equation}
\operatorname{Hom}_{S_0}(X,S)\longrightarrow
\varprojlim_\lambda\operatorname{Hom}_{S_0}(X,S_\lambda)
\tag{8.2.3.1}\label{IV.8.2.3.1-fr}
\end{equation}
qui à tout $S_0$-morphisme $v:X\to S$ fait correspondre la famille
$(u_\lambda\circ v)$, est \emph{bijective}. Or, si $g:X\to S_0$ est le
morphisme structural et si on pose $\mathscr B=g_*(\mathcal O_X)$, qui est
une $\mathcal O_{S_0}$-Algèbre, l'application (8.2.3.1) s'identifie
canoniquement à (8.2.1.1) (II, 1.2.7), et la conclusion résulte donc de ce
qu'on a vu dans (8.2.1).

Les autres assertions de (8.2.3) sont conséquences du lemme général suivant~:

\begin{lemma}[8.2.4]
\label{IV.8.2.4-fr}
Soient $\mathbf C$ une catégorie, $T$ un objet de $\mathbf C$,
$\mathbf C_T$ la sous-catégorie des objets de $\mathbf C$ au-dessus de $T$.
Soit $(S_\lambda,u_{\lambda\mu})$ un système projectif dans $\mathbf C_T$~;
alors toute limite projective de ce système dans $\mathbf C_T$ est aussi
limite projective dans $\mathbf C$, et réciproquement.
\end{lemma}

Soit $f_\lambda:S_\lambda\to T$ le morphisme structural. Supposons que $S$
soit limite projective de $(S_\lambda,u_{\lambda\mu})$ dans $\mathbf C$, et
désignons par $u_\lambda:S\to S_\lambda$ les morphismes canoniques
correspondants. Considérons alors un système projectif de $T$-morphismes
$w_\lambda:Y\to S_\lambda$, où $Y\in\mathbf C_T$. Il existe par hypothèse un
morphisme unique (dans $\mathbf C$) $w:Y\to S$ tel que
$w_\lambda=u_\lambda\circ w$ pour tout $\lambda$. L'hypothèse que les
$u_\lambda$ sont des $T$-morphismes entraîne que les morphismes
$f_\lambda\circ u_\lambda:S\to T$ sont tous les mêmes, et ce morphisme $f$
fait donc de $S$ un $T$-objet. Si $g:Y\to T$ est le morphisme structural de
$Y$, on a alors
$f\circ w=f_\lambda\circ u_\lambda\circ w=f_\lambda\circ w_\lambda=g$ pour
tout $\lambda$, ce qui prouve que $w$ est un $T$-morphisme. Inversement,
supposons (avec les mêmes notations) que $S$ soit limite projective de
$(S_\lambda,u_{\lambda\mu})$ dans $\mathbf C_T$, et considérons maintenant un
système projectif de morphismes (de $\mathbf C$) $w_\lambda:Y\to S_\lambda$.
Les morphismes composés $f_\lambda\circ w_\lambda:Y\to T$ sont alors tous
les mêmes~: en effet, pour deux indices quelconques $\lambda$, $\mu$, il y a
un indice $\nu$ tel que $\lambda\leq\nu$ et $\mu\leq\nu$, d'où
$f_\nu=f_\lambda\circ u_{\lambda\nu}=f_\mu\circ u_{\mu\nu}$~; comme
$w_\lambda=u_{\lambda\nu}\circ w_\nu$ et
$w_\mu=u_{\mu\nu}\circ w_\nu$, on a
$f_\lambda\circ w_\lambda=f_\lambda\circ u_{\lambda\nu}\circ w_\nu
=f_\nu\circ w_\nu$, et on voit de même que
$f_\mu\circ w_\mu=f_\nu\circ w_\nu$. Si $g:Y\to T$ est l'unique morphisme
ainsi défini, $g$ fait de $Y$ un $T$-objet, et les $w_\lambda$ sont alors des
$T$-morphismes~; ils ont par suite une limite projective $w:Y\to S$ qui est
un $T$-morphisme, et \emph{a fortiori} un morphisme de $\mathbf C$~;
d'ailleurs la première partie du raisonnement montre que toute limite
projective $w'$ (dans $\mathbf C$) du système projectif $(w_\lambda)$ est
nécessairement aussi un $T$-morphisme, donc égale à $w$, ce qui achève de
prouver le lemme.

\begin{proposition}[8.2.5]
\label{IV.8.2.5-fr}
Sous les conditions de (8.2.2), soient $\alpha$ un élément de $L$,
$X_\alpha$ un $S_\alpha$-préschéma. Pour tout $\lambda\geq\alpha$, posons
$X_\lambda=X_\alpha\times_{S_\alpha}S_\lambda$, et pour
$\alpha\leq\lambda\leq\mu$, posons
$v_{\lambda\mu}=1_{X_\alpha}\times u_{\lambda\mu}$, de sorte que
$(X_\lambda,v_{\lambda\mu})$ est un système projectif de
$X_\alpha$-préschémas, dont l'ensemble d'indices est formé des
$\lambda\geq\alpha$ dans $L$. Posons de même
$X=X_\alpha\times_{S_\alpha}S$ et $v_\lambda=1_{X_\alpha}\times u_\lambda$.
Alors les $X_\alpha$-morphismes $v_\lambda:X\to X_\lambda$ font de $X$ une
limite projective du système projectif $(X_\lambda,v_{\lambda\mu})$ dans la
catégorie des $X_\alpha$-préschémas, ou dans la catégorie de tous les
préschémas.
\end{proposition}

Cela résultera encore du lemme général suivant~:

\begin{lemma}[8.2.6]
\label{IV.8.2.6-fr}
Soient $\mathbf C$ une catégorie dans laquelle les produits fibrés existent,
$q:T'\to T$ un morphisme de $\mathbf C$, $\mathbf C_T$ (resp.
$\mathbf C_{T'}$) la catégorie des objets de $\mathbf C$ au-dessus
\oldpage[IV-3]{10}
de $T$ (resp. $T'$). Soit $(S_\lambda,u_{\lambda\mu})$ un système projectif
(non nécessairement filtrant) dans $\mathbf C_T$, et posons
$S'_\lambda=S_\lambda\times_T T'$,
$u'_{\lambda\mu}=u_{\lambda\mu}\times 1_{T'}$, de sorte que
$(S'_\lambda,u'_{\lambda\mu})$ est un système projectif dans
$\mathbf C_{T'}$. Alors, si $(S,u_\lambda)$ est une limite projective de
$(S_\lambda,u_{\lambda\mu})$ dans $\mathbf C_T$,
$(S\times_T T',u_\lambda\times 1_{T'})$ est une limite projective de
$(S'_\lambda,u'_{\lambda\mu})$ dans $\mathbf C_{T'}$.
\end{lemma}

On a par hypothèse, pour tout $\lambda$, un diagramme commutatif
\[
\xymatrix{
  S' \ar[r]^{u'_\lambda} \ar[d]_p
    & S'_\lambda \ar[r]^{h'_\lambda} \ar[d]_{p_\lambda}
    & T' \ar[d]^q \\
  S \ar[r]_{u_\lambda} & S_\lambda \ar[r]_{h_\lambda} & T
}
\]
où l'on a posé $S'=S\times_T T'$, $u'_\lambda=u_\lambda\times 1_{T'}$,
$h'_\lambda=h_\lambda\times 1_{T'}$. Soient $Y$ un $T'$-objet,
$g':Y\to T'$ le morphisme correspondant, et considérons un système
projectif de $T'$-morphismes $w'_\lambda:Y\to S'_\lambda$. Alors $Y$ est un
$T$-objet pour le morphisme $g=q\circ g'$, et les
$w_\lambda=p_\lambda\circ w'_\lambda$ sont des $T$-morphismes puisque
$h_\lambda\circ w_\lambda=h_\lambda\circ p_\lambda\circ w'_\lambda
=q\circ h'_\lambda\circ w'_\lambda=q\circ g'$ par hypothèse. En outre, on
vérifie aussitôt que $(w_\lambda)$ est un système projectif. Il existe donc
par hypothèse un $T$-morphisme $w:Y\to S$ et un seul tel que
$u_\lambda\circ w=w_\lambda$ pour tout $\lambda$. Par définition du produit
fibré, il y a un unique $T'$-morphisme $w':Y\to S'$ tel que
$p\circ w'=w$. On a alors
$u_\lambda\circ p\circ w'=u_\lambda\circ w=w_\lambda
=p_\lambda\circ w'_\lambda$, ce qui s'écrit aussi
$p_\lambda\circ u'_\lambda\circ w'=p_\lambda\circ w'_\lambda$~; d'autre
part, en écrivant que $w'$ est un $T'$-morphisme, il vient
$h'_\lambda\circ u'_\lambda\circ w'=g'=h'_\lambda\circ w'_\lambda$. La
définition de $S'_\lambda$ comme produit fibré $S_\lambda\times_T T'$ donne
donc $u'_\lambda\circ w'=w'_\lambda$, et il est immédiat que $w'$ est
l'unique $T'$-morphisme vérifiant ces relations, d'où le lemme.

\begin{remark}[8.2.7]
\label{IV.8.2.7-fr}
Étant donné un espace annelé quelconque $S$, les limites inductives par
rapport à un ensemble préordonné \emph{quelconque} $L$ (non nécessairement
filtrant) existent dans la catégorie des $\mathcal O_S$-Algèbres
commutatives, puisque la limite inductive filtrante existe d'après (8.2.1) et
d'autre part, pour deux homomorphismes de $\mathcal O_S$-Algèbres
$\mathscr A\to\mathscr B$, $\mathscr A\to\mathscr C$, le produit tensoriel
$\mathscr B\otimes_{\mathscr A}\mathscr C$ est la «~somme amalgamée~»
correspondante dans cette catégorie.

Lorsque $S$ est un préschéma, on sait que le produit tensoriel
$\mathscr B\otimes_{\mathscr A}\mathscr C$ est une $\mathcal O_S$-Algèbre
quasi-cohérente lorsqu'il en est ainsi de $\mathscr A$, $\mathscr B$ et
$\mathscr C$ (I, 1.3.13)~; on en conclut que, dans la catégorie des
$\mathcal O_S$-Algèbres \emph{quasi-cohérentes}, les limites inductives pour
un ensemble préordonné quelconque d'indices existent toujours. Cela permet de
généraliser la définition d'une limite projective de préschémas et les prop.
(8.2.3) et (8.2.5) au cas où l'ensemble préordonné $L$ n'est pas
nécessairement filtrant.
\end{remark}

\begin{env}[8.2.8]
\label{IV.8.2.8-fr}
Les notations étant celles de (8.2.2), posons
$u_{\lambda\mu}=(\psi_{\lambda\mu},\theta_{\lambda\mu})$ et
$u_\lambda=(\psi_\lambda,\theta_\lambda)$~; $(S_\lambda,\psi_{\lambda\mu})$
est donc un système projectif d'espaces topologiques et $(\psi_\lambda)$ un
système projectif d'applications continues $S\to S_\lambda$, des espaces
sous-jacents respectivement aux préschémas $S$ et $S_\lambda$.
\end{env}

\begin{proposition}[8.2.9]
\label{IV.8.2.9-fr}
Avec les notations de (8.2.8), la limite projective du système projectif
$(\psi_\lambda)$ d'applications continues est un homéomorphisme de l'espace
sous-jacent à $S$ sur la limite projective du système projectif
$(S_\lambda,\psi_{\lambda\mu})$ d'espaces topologiques.
\end{proposition}

\oldpage[IV-3]{11}
Soit $T$ l'espace topologique limite du système
$(S_\lambda,\psi_{\lambda\mu})$ et posons
$\psi=\varprojlim\psi_\lambda:S\to T$. On peut se borner au cas où
$S_0=S_\alpha$ pour un $\alpha\in L$, et $\lambda\geq\alpha$.

\medskip
\noindent
(i) Montrons en premier lieu que la topologie de $S$ est image réciproque par
$\psi$ de la topologie de $T$~; autrement dit, si
$\pi_\lambda:T\to S_\lambda$ est l'application canonique, il faut montrer
que tout ouvert de $S$ est réunion d'ouverts de la forme
$\psi^{-1}(\pi_\lambda^{-1}(U_\lambda))$, où $U_\lambda$ est ouvert dans
$S_\lambda$. Or tout ouvert de $S$ est par définition réunion d'ouverts
obtenus de la façon suivante~: on considère un ouvert affine $U_0$ de $S_0$,
d'anneau $A_0$, de sorte que $\psi_\alpha^{-1}(U_0)$ est l'ouvert affine de
$S$ d'anneau $A=\Gamma(U_0,\mathscr A)$, puis on prend un élément $f\in A$ et
on considère dans $\operatorname{Spec}(A)$, identifié à
$\psi_\alpha^{-1}(U_0)$, l'ouvert $D(f)$. Ce sont ces ouverts $D(f)$ qui
forment une base de la topologie de $S$ (II, 1.3.1). Or, si on pose
$A_\lambda=\Gamma(U_0,\mathscr A_\lambda)$, on a
$A=\varinjlim A_\lambda$ (I, 1.3.9), donc il existe un indice $\lambda$ tel
que $f$ soit l'image canonique d'un élément $f_\lambda\in A_\lambda$~; on a
alors $D(f)=\psi_\lambda^{-1}(D(f_\lambda))$ (I, 1.2.2), et comme
$\psi_\lambda=\pi_\lambda\circ\psi$, notre assertion est démontrée.

\medskip
\noindent
(ii) Prouvons maintenant que $\psi$ est \emph{bijective}, ce qui achèvera la
démonstration. Comme $S$ est un espace de Kolmogoroff, il résulte déjà de (i)
que $\psi$ est injective, et il reste donc à montrer que $\psi$ est
surjective. On peut évidemment remplacer pour cela $T$ par un ouvert
$\pi_\alpha^{-1}(U_0)$, où $U_0$ est ouvert affine dans $S_\alpha=S_0$, donc
on peut se limiter au cas où les $S_\lambda$ et $S$ sont affines, autrement
dit $\mathscr A_\lambda$ est le faisceau associé à une $A_0$-algèbre
$A_\lambda$, et $\mathscr A$ le faisceau d'algèbres associé à
$A=\varinjlim A_\lambda$~; nous noterons encore
$\varphi_{\mu\lambda}:A_\lambda\to A_\mu$ et
$\varphi_\lambda:A_\lambda\to A$ les homomorphismes canoniques. Par
définition, un élément de $T$ est une famille
$(\mathfrak p_\lambda)_{\lambda\in L}$, où $\mathfrak p_\lambda$ est un
idéal premier de $A_\lambda$ et où l'on a
$\mathfrak p_\lambda=\varphi_{\mu\lambda}^{-1}(\mathfrak p_\mu)$ pour
$\lambda\leq\mu$. On sait alors (5.13.3 et 5.13.1) qu'il y a un idéal premier
$\mathfrak p$ de $A$ tel que
$\mathfrak p_\lambda=\varphi_\lambda^{-1}(\mathfrak p)$ pour tout
$\lambda\in L$, ce qui achève la démonstration.

En particulier, on a ainsi démontré le

\begin{corollary}[8.2.10]
\label{IV.8.2.10-fr}
Soit $(A_\lambda)_{\lambda\in L}$ un système inductif filtrant d'anneaux, et
soient $A=\varinjlim A_\lambda$, $\varphi_\lambda:A_\lambda\to A$ les
homomorphismes canoniques. L'application canonique
$\mathfrak p\mapsto(\varphi_\lambda^{-1}(\mathfrak p))$ est un
homéomorphisme de $\operatorname{Spec}(A)$ sur l'espace topologique
$\varprojlim\operatorname{Spec}(A_\lambda)$.
\end{corollary}

\begin{corollary}[8.2.11]
\EGAReaderTarget{ega4.c04.cor.8-2-11}{IV.8.2.11-fr}
Avec les notations de (8.2.8), pour tout ouvert quasi-compact $U$ de $S$, il
existe un indice $\lambda$ et un ouvert quasi-compact $U_\lambda$ de
$S_\lambda$ tel que $U=\psi_\lambda^{-1}(U_\lambda)$.
\end{corollary}

Cela résulte de ce que, par définition de la topologie limite projective, les
$\psi_\lambda^{-1}(U_\lambda)$ ($U_\lambda$ ouvert quasi-compact de
$S_\lambda$) forment une base de la topologie de $S$, et de ce que l'ensemble
d'indices $L$ est filtrant.

\begin{corollary}[8.2.12]
\label{IV.8.2.12-fr}
Avec les notations de (8.2.8), la limite inductive du système inductif
d'homomorphismes
$\theta_\lambda^\sharp:\psi_\lambda^*(\mathcal O_{S_\lambda})\to
\mathcal O_S$ de faisceaux d'anneaux sur $S$ est un isomorphisme
\begin{equation}
\varinjlim_\lambda\psi_\lambda^*(\mathcal O_{S_\lambda})
\xrightarrow{\sim}\mathcal O_S.
\tag{8.2.12.1}\label{IV.8.2.12.1-fr}
\end{equation}
\end{corollary}

On peut évidemment supposer les $S_\lambda$ affines~; avec les notations de
la démonstration de (8.2.9), tout revient à voir que la limite inductive du
système d'applications canoniques
$(A_\lambda)_{\mathfrak p_\lambda}\to A_{\mathfrak p}$ est un isomorphisme,
ce qui n'est autre que (5.13.3, (ii)).

\oldpage[IV-3]{12}
\begin{proposition}[8.2.13]
\label{IV.8.2.13-fr}
Supposons que les morphismes $u_{\lambda\mu}:S_\mu\to S_\lambda$ soient des
immersions ouvertes, de sorte que $S_\mu$ s'identifie à un sous-préschéma
induit sur un ouvert de $S_\lambda$ pour $\lambda\leq\mu$. Alors, pour tout
$\alpha\in L$, l'espace sous-jacent au préschéma $S$ s'identifie au
sous-espace de $S_\alpha$ intersection des $S_\lambda$ pour
$\lambda\geq\alpha$, et le faisceau structural $\mathcal O_S$ au faisceau
induit $(\mathrm G,\mathrm{II},1.5)$ par $\mathcal O_{S_\alpha}$ sur cette
intersection~; plus généralement, pour tout $\mathcal O_{S_\alpha}$-Module
$\mathscr F_\alpha$, $u_\alpha^*(\mathscr F_\alpha)$ s'identifie au
$\mathcal O_S$-Module induit par $\mathscr F_\alpha$ sur $S$.
\end{proposition}

La première assertion résulte de (8.2.9), vu la définition d'une limite
projective d'espaces topologiques~; en outre tous les
$\psi_\lambda^*(\mathcal O_{S_\lambda})$ sont égaux au faisceau induit par
$\mathcal O_{S_\alpha}$ sur $S$ par définition $(0_{\mathrm I},3.7.1)$ et,
avec les notations de la démonstration de (8.2.9), les homomorphismes
$(A_\lambda)_{\mathfrak p_\lambda}\to
(A_\mu)_{\mathfrak p_\mu}$ sont bijectifs pour un système
$(\mathfrak p_\lambda)$ d'idéaux premiers correspondant à un même point de
$S$~; l'assertion relative à $\mathcal O_S$ se déduit donc de (8.2.12). La
dernière assertion résulte alors de la définition de
$u_\alpha^*(\mathscr F_\alpha)$ $(0_{\mathrm I},4.3.1)$.

\begin{remark}[8.2.14]
\label{IV.8.2.14-fr}
Les résultats de (8.2.9) et (8.2.12) montrent que $S$ est limite projective du
système projectif $(S_\lambda,u_{\lambda\mu})$ dans la catégorie de
\emph{tous} les espaces annelés (ou de tous les espaces annelés en anneaux
locaux). En effet, soit $Y$ un espace annelé, et considérons un système
projectif de morphismes d'espaces annelés $w_\lambda:Y\to S_\lambda$. Si
l'on pose $w_\lambda=(\rho_\lambda,\omega_\lambda)$, les $\rho_\lambda$
forment un système projectif d'applications continues et, en vertu de
(8.2.9), leur limite projective $\rho$ s'identifie à une application continue
$Y\to S$ telle que $\rho_\lambda=\psi_\lambda\circ\rho$. D'autre part, les
\[
\omega_\lambda^\sharp:
\rho_\lambda^*(\mathcal O_{S_\lambda})\to\mathcal O_Y
\]
forment un système inductif d'homomorphismes de faisceaux d'anneaux~; comme
on peut écrire
$\rho_\lambda^*(\mathcal O_{S_\lambda})=
\rho^*(\psi_\lambda^*(\mathcal O_{S_\lambda}))$ et que le foncteur $\rho^*$
est exact, la limite inductive des
$\rho_\lambda^*(\mathcal O_{S_\lambda})$ est $\rho^*(\mathcal O_S)$ en
vertu de (8.2.12), et il y a donc un homomorphisme unique
$\omega^\sharp:\rho^*(\mathcal O_S)\to\mathcal O_Y$ tel que
$\omega_\lambda^\sharp=\omega^\sharp\circ
\rho^*(\theta_\lambda^\sharp)$, ce qui démontre notre assertion.
\end{remark}

\subsection{Parties constructibles dans une limite projective de préschémas}
\label{subsection:IV.8.3-fr}

\begin{env}[8.3.1]
\label{IV.8.3.1-fr}
Dans toute la suite de ce paragraphe, nous supposons remplies les conditions
de (8.2.2), dont nous conservons les notations.
\end{env}

\begin{theorem}[8.3.2]
\label{IV.8.3.2-fr}
Pour tout $\lambda$, soient $E_\lambda$, $F_\lambda$ deux parties de
$S_\lambda$. Posons
\begin{equation}
E=\bigcap_\lambda u_\lambda^{-1}(E_\lambda),\qquad
F=\bigcup_\lambda u_\lambda^{-1}(F_\lambda).
\tag{8.3.2.1}\label{IV.8.3.2.1-fr}
\end{equation}
On suppose vérifiées les conditions suivantes~:
\begin{enumerate}
\item[(i)] Pour tout $\lambda$, $E_\lambda$ est pro-constructible et
$F_\lambda$ est ind-constructible (1.9.4).
\item[(ii)] Pour $\lambda\leq\mu$, on a
\[
E_\mu\subset u_{\lambda\mu}^{-1}(E_\lambda),\qquad
F_\mu\supset u_{\lambda\mu}^{-1}(F_\lambda).
\]
\item[(iii)] Il existe $\alpha\in L$ tel que $S_\alpha$ soit quasi-compact
(ce qui entraîne que $S_\lambda$ est quasi-compact pour
$\lambda\geq\alpha$).
\end{enumerate}
Alors, les propriétés suivantes sont équivalentes~:
\begin{enumerate}
\item[\emph{a)}] $E\subset F$.
\oldpage[IV-3]{13}
\item[\emph{b)}] Il existe $\lambda\geq\alpha$ tel que
$u_\lambda^{-1}(E_\lambda)\subset u_\lambda^{-1}(F_\lambda)$ (et alors on a
$u_\mu^{-1}(E_\mu)\subset u_\mu^{-1}(F_\mu)$ pour $\mu\geq\lambda$).
\item[\emph{c)}] Il existe $\lambda\geq\alpha$ tel que
$E_\lambda\subset F_\lambda$ (et alors on a $E_\mu\subset F_\mu$ pour
$\mu\geq\lambda$).
\end{enumerate}
\end{theorem}

Les remarques entre parenthèses dans \emph{b)} et \emph{c)} résultent de (ii).
Posons
\[
G_\lambda=E_\lambda\cap(S_\lambda-F_\lambda),\qquad
G=E\cap(S-F).
\]
Alors $G_\lambda$ est une partie pro-constructible de $S_\lambda$ (1.9.5,
(i)), et en vertu de (8.3.2.1) et de (ii), on a
\[
G_\mu\subset u_{\lambda\mu}^{-1}(G_\lambda)
\quad\text{pour }\lambda\leq\mu,
\qquad
G=\bigcap_\lambda u_\lambda^{-1}(G_\lambda).
\]
On est donc ramené à prouver le cas particulier de (8.3.2) correspondant à
$F_\lambda=\emptyset$ pour tout $\lambda$~:

\begin{corollary}[8.3.3]
\label{IV.8.3.3-fr}
Pour tout $\lambda$, soit $E_\lambda$ une partie pro-constructible de
$S_\lambda$ telle que, pour $\lambda\leq\mu$, on ait
$E_\mu\subset u_{\lambda\mu}^{-1}(E_\lambda)$. Supposons qu'il existe
$\alpha\in L$ tel que $S_\alpha$ soit quasi-compact. Alors les conditions
suivantes sont équivalentes~:
\begin{enumerate}
\item[\emph{a)}]
$E=\bigcap_\lambda u_\lambda^{-1}(E_\lambda)=\emptyset$.
\item[\emph{b)}] Il existe un $\lambda$ tel que
$u_\lambda^{-1}(E_\lambda)=\emptyset$ (et alors
$u_\mu^{-1}(E_\mu)=\emptyset$ pour $\mu\geq\lambda$).
\item[\emph{c)}] Il existe $\lambda$ tel que $E_\lambda=\emptyset$ (et alors
$E_\mu=\emptyset$ pour $\mu\geq\lambda$).
\end{enumerate}
\end{corollary}

Il est clair que \emph{c)} implique \emph{a)}. Prouvons que \emph{a)} entraîne
\emph{b)}~: puisque $S_\lambda$ est quasi-compact, il en est de même de $S$
(8.2.2)~; $E_\lambda$ étant pro-constructible, il en est de même de
$u_\lambda^{-1}(E_\lambda)$ (1.9.5, (vi))~; alors la famille filtrante
décroissante d'ensembles pro-constructibles $u_\lambda^{-1}(E_\lambda)$ a
une intersection vide, donc (1.9.9) l'un d'eux est vide.

Enfin, montrons que \emph{b)} entraîne \emph{c)}. Comme $S_\alpha$ est
quasi-compact et $L$ filtrant, on peut remplacer $S_\alpha$ par un ouvert
affine, donc supposer (8.2.2) que $S$ et les $S_\lambda$ pour
$\lambda\geq\alpha$ sont affines~; on a alors (1.9.2.1), pour
$\lambda\geq\alpha$,
\[
u_\lambda(S)=\bigcap_{\mu\geq\lambda}u_{\lambda\mu}(S_\mu),
\]
d'où
\[
E_\lambda\cap u_\lambda(S)=
\bigcap_{\mu\geq\lambda}
\bigl(E_\lambda\cap u_{\lambda\mu}(S_\mu)\bigr).
\]
Or, comme $u_\lambda$ et les $u_{\lambda\mu}$ sont quasi-compacts,
$u_\lambda(S)$ et $u_{\lambda\mu}(S_\mu)$ sont pro-constructibles dans
$S_\lambda$ (1.9.5, (vii)), donc les ensembles
$E_\lambda\cap u_{\lambda\mu}(S_\mu)$ pour $\mu\geq\lambda$ forment une
famille filtrante décroissante de parties pro-constructibles de $S_\lambda$.
Comme $S_\lambda$ est quasi-compact, l'hypothèse \emph{b)} entraîne que
l'intersection de cette famille est vide, donc (1.9.9) un des ensembles
$E_\lambda\cap u_{\lambda\mu}(S_\mu)$ est vide, donc
$E_\mu\subset u_{\lambda\mu}^{-1}(E_\lambda)$ est vide. C.Q.F.D.

\begin{corollary}[8.3.4]
\label{IV.8.3.4-fr}
Pour tout $\lambda$, soit $F_\lambda$ une partie ind-constructible de
$S_\lambda$ telle que pour $\lambda\leq\mu$ on ait
$u_{\lambda\mu}^{-1}(F_\lambda)\subset F_\mu$. Supposons qu'il existe
$\alpha\in L$ tel que $S_\alpha$ soit quasi-compact. Alors les conditions
suivantes sont équivalentes~:
\begin{enumerate}
\item[\emph{a)}] L'ensemble
$F=\bigcup_\lambda u_\lambda^{-1}(F_\lambda)$ est égal à $S$.
\item[\emph{b)}] Il existe un $\lambda$ tel que
$u_\lambda^{-1}(F_\lambda)=S$ (et alors
$u_\mu^{-1}(F_\mu)=S$ pour $\mu\geq\lambda$).
\oldpage[IV-3]{14}
\item[\emph{c)}] Il existe un $\lambda$ tel que $F_\lambda=S_\lambda$ (et
alors $F_\mu=S_\mu$ pour $\mu\geq\lambda$).
\end{enumerate}
\end{corollary}

Cela se déduit aussitôt de (8.3.3) par passage aux complémentaires.

\begin{corollary}[8.3.5]
\label{IV.8.3.5-fr}
Pour tout $\lambda$, soient $E_\lambda$, $F_\lambda$ deux parties
constructibles de $S_\lambda$ telles que, pour $\mu\geq\lambda$, on ait
$E_\mu\subset u_{\lambda\mu}^{-1}(E_\lambda)$ et
$F_\mu\supset u_{\lambda\mu}^{-1}(F_\lambda)$. Supposons qu'il existe un
indice $\alpha$ tel que $S_\alpha$ soit quasi-compact. Alors, pour que
$E\subset F$ (resp. $E=F$), il faut et il suffit qu'il existe $\lambda$ tel
que $E_\lambda\subset F_\lambda$ (resp. $E_\lambda=F_\lambda$), auquel cas
on a aussi $E_\mu\subset F_\mu$ (resp. $E_\mu=F_\mu$) pour
$\mu\geq\lambda$.
\end{corollary}

C'est un cas particulier de (8.3.2).

En particulier~:

\begin{corollary}[8.3.6]
\label{IV.8.3.6-fr}
Supposons qu'il existe un $\alpha$ tel que $S_\alpha$ soit quasi-compact.
Pour que $S=\emptyset$, il faut et il suffit qu'il existe un $\lambda$ tel
que $S_\lambda=\emptyset$.
\end{corollary}

\begin{corollary}[8.3.7]
\label{IV.8.3.7-fr}
On a, pour tout $\lambda$,
\begin{equation}
u_\lambda(S)=\bigcap_{\mu\geq\lambda}u_{\lambda\mu}(S_\mu).
\tag{8.3.7.1}\label{IV.8.3.7.1-fr}
\end{equation}
\end{corollary}

Il est clair que le premier membre de (8.3.7.1) est contenu dans le second.
Soit $s$ un point de $S_\lambda$ et posons
$X_\lambda=\operatorname{Spec}(k(s))$~; considérons le système projectif
$(X_\mu,v_{\mu\nu})$ où
$X_\mu=X_\lambda\times_{S_\lambda}S_\mu$ et
$v_{\mu\nu}=1\times u_{\mu\nu}$ pour
$\lambda\leq\mu\leq\nu$~; sa limite projective est
$X=X_\lambda\times_{S_\lambda}S$ et $v_\lambda=1\times u_\lambda$ est
l'application canonique $X\to X_\lambda$ (8.2.5). Si
$s\in\bigcap_{\mu\geq\lambda}u_{\lambda\mu}(S_\mu)$, cela entraîne que
$X_\mu\neq\emptyset$ pour tout $\mu\geq\lambda$ (I, 3.4.8)~; il résulte
alors de (8.3.6) que $X\neq\emptyset$, donc $s\in u_\lambda(S)$ par
(I, 3.4.8).

\begin{proposition}[8.3.8]
\label{IV.8.3.8-fr}
\begin{enumerate}
\item[(i)] Pour que le morphisme $u_\lambda:S\to S_\lambda$ soit dominant
(resp. surjectif), il faut et il suffit que pour $\mu\geq\lambda$ le
morphisme $u_{\lambda\mu}:S_\mu\to S_\lambda$ soit dominant (resp.
surjectif).
\item[(ii)] Si, pour un indice $\lambda$, les morphismes
$u_{\lambda\mu}:S_\mu\to S_\lambda$ sont plats (resp. fidèlement plats)
pour tout $\mu\geq\lambda$, alors le morphisme
$u_\lambda:S\to S_\lambda$ est plat (resp. fidèlement plat).
\item[(iii)] Supposons que les morphismes $u_\mu:S\to S_\mu$ soient
surjectifs pour $\mu\geq\lambda$. Pour que $u_\lambda$ soit un morphisme
ouvert (resp. universellement ouvert), il faut et il suffit que, pour tout
$\mu\geq\lambda$, $u_{\lambda\mu}$ soit un morphisme ouvert (resp.
universellement ouvert).
\end{enumerate}
\end{proposition}

\noindent
(i) Comme $u_\lambda(S)\subset u_{\lambda\mu}(S_\mu)$ pour
$\mu\geq\lambda$, la nécessité des conditions est triviale, et il résulte
aussitôt de (8.3.7.1) que si les $u_{\lambda\mu}$ sont surjectifs, il en est
de même de $u_\lambda$. Supposons maintenant les $u_{\lambda\mu}$ dominants
pour $\mu\geq\lambda$, et considérons dans $S_\lambda$ un ouvert
quasi-compact non vide $U_\lambda$~; alors les
$U_\mu=u_{\lambda\mu}^{-1}(U_\lambda)$ pour $\mu\geq\lambda$ forment un
système projectif dont la limite projective est
$U=u_\lambda^{-1}(U_\lambda)$ (8.2.5). Par hypothèse les $U_\mu$ sont tous
non vides, donc il en est de même de $U$ par (8.3.6), ce qui prouve que
$u_\lambda$ est dominant.

\medskip
\noindent
(ii) En vertu de (i), il suffit de considérer le cas où les
$u_{\lambda\mu}$ sont plats~; on peut alors se borner au cas où $S_\lambda$
est affine, donc aussi les $S_\mu$ pour $\mu\geq\lambda$ et $S$, et
l'assertion résulte alors de (2.1.2) et $(0_{\mathrm I},6.2.3)$.

\medskip
\noindent
(iii) En vertu de (8.2.5) et de (I, 3.5.2), il suffit de traiter le cas des
morphismes ouverts. Comme $u_\lambda=u_{\lambda\mu}\circ u_\mu$ et que
$u_\mu$ est surjectif, on sait que si $u_\lambda$ est ouvert il en est de
même de $u_{\lambda\mu}$ pour $\mu\geq\lambda$ (Bourbaki,
\emph{Top. gén.}, chap. $\mathrm I^{\mathrm{er}}$, $3^{\mathrm e}$ éd.,
\S~5, n\textsuperscript{o} 1, prop. 1).

\oldpage[IV-3]{15}
Inversement, pour montrer que $u_\lambda$ est ouvert lorsque tous les
$u_{\lambda\mu}$ le sont pour $\mu\geq\lambda$, il suffit de voir que pour
tout ouvert quasi-compact $V$ de $S$, $u_\lambda(V)$ est ouvert dans
$S_\lambda$~; mais il existe alors $\mu\geq\lambda$ tel que
$V=u_\mu^{-1}(V_\mu)$, où $V_\mu$ est ouvert dans $S_\mu$ (8.2.11)~; comme
$u_\mu$ est surjectif, on a $V_\mu=u_\mu(V)$ et
$u_\lambda(V)=u_{\lambda\mu}(u_\mu(V))$ est donc ouvert par hypothèse.

On notera qu'il peut se faire que tous les $u_{\lambda\mu}$ soient ouverts
sans qu'il en soit ainsi de $u_\lambda$ si les $u_\mu$ ne sont pas
surjectifs. On en a un exemple en considérant un anneau intègre $A$ qui n'est
pas un corps, et son corps des fractions $K$, qui est limite inductive des
$A_f$, où $f$ parcourt $A-\{0\}$~; si l'on pose
$S_1=\operatorname{Spec}(A)$, $S_f=\operatorname{Spec}(A_f)$, on a
$S=\varprojlim S_f=\operatorname{Spec}(K)$, et le morphisme $S\to S_1$ n'est
pas ouvert, bien que les morphismes $S_f\to S_1$ le soient.

\begin{env}[8.3.9]
\label{IV.8.3.9-fr}
Pour tout préschéma $X$, nous noterons comme d'ordinaire $\mathfrak P(X)$
l'ensemble des parties de l'ensemble sous-jacent à $X$, par $\mathfrak C(X)$
(resp. $\mathfrak D_c(X)$, $\mathfrak F_c(X)$, $\mathfrak{LF}_c(X)$)
l'ensemble des parties constructibles (resp. constructibles et ouvertes,
resp. constructibles et fermées, resp. constructibles et localement fermées)
de $X$. Il est clair que
$(\mathfrak P(S_\lambda),u_{\lambda\mu}^{-1})$ est un \emph{système inductif
d'ensembles} et que les applications
$u_\lambda^{-1}:\mathfrak P(S_\lambda)\to\mathfrak P(S)$ forment un système
inductif d'applications, d'où, en passant à la limite inductive, une
application canonique
\begin{equation}
u_{\mathfrak P}:\varinjlim\mathfrak P(S_\lambda)\to\mathfrak P(S).
\tag{8.3.9.1}\label{IV.8.3.9.1-fr}
\end{equation}
En outre, il résulte de (1.8.2) que $u_{\lambda\mu}^{-1}$ applique
$\mathfrak C(S_\lambda)$ (resp. $\mathfrak D_c(S_\lambda)$,
$\mathfrak F_c(S_\lambda)$, $\mathfrak{LF}_c(S_\lambda)$) dans
$\mathfrak C(S_\mu)$ (resp. $\mathfrak D_c(S_\mu)$,
$\mathfrak F_c(S_\mu)$, $\mathfrak{LF}_c(S_\mu)$) et que
$u_\lambda^{-1}$ applique $\mathfrak C(S_\lambda)$ (resp.
$\mathfrak D_c(S_\lambda)$, $\mathfrak F_c(S_\lambda)$,
$\mathfrak{LF}_c(S_\lambda)$) dans $\mathfrak C(S)$ (resp.
$\mathfrak D_c(S)$, $\mathfrak F_c(S)$, $\mathfrak{LF}_c(S)$). On a donc par
restriction de (8.3.9.1), des applications canoniques
\begin{align}
\varinjlim\mathfrak C(S_\lambda)&\longrightarrow\mathfrak C(S)
\tag{8.3.9.2}\label{IV.8.3.9.2-fr}\\
\varinjlim\mathfrak D_c(S_\lambda)&\longrightarrow\mathfrak D_c(S)
\tag{8.3.9.3}\label{IV.8.3.9.3-fr}\\
\varinjlim\mathfrak F_c(S_\lambda)&\longrightarrow\mathfrak F_c(S)
\tag{8.3.9.4}\label{IV.8.3.9.4-fr}\\
\varinjlim\mathfrak{LF}_c(S_\lambda)&\longrightarrow\mathfrak{LF}_c(S).
\tag{8.3.9.5}\label{IV.8.3.9.5-fr}
\end{align}
\end{env}

\begin{env}[8.3.10]
\label{IV.8.3.10-fr}
Soit $g_\alpha:X_\alpha\to S_\alpha$ un morphisme~; avec les notations de
(8.2.5) on a comme ci-dessus une application canonique
$v_{\mathfrak P}:\varinjlim\mathfrak P(X_\lambda)\to\mathfrak P(X)$~;
d'autre part, on a des morphismes de projection
$g_\lambda:X_\lambda\to S_\lambda$ pour tout $\lambda\geq\alpha$ et un
morphisme de projection $g:X\to S$. Il est clair que les
$g_\lambda^{-1}:\mathfrak P(S_\lambda)\to\mathfrak P(X_\lambda)$ forment un
système inductif d'applications, et que les diagrammes
\[
\xymatrix{
  \mathfrak P(S_\lambda) \ar[r]^{u_\lambda^{-1}} \ar[d]_{g_\lambda^{-1}}
    & \mathfrak P(S) \ar[d]^{g^{-1}} \\
  \mathfrak P(X_\lambda) \ar[r]_{v_\lambda^{-1}} & \mathfrak P(X)
}
\]
\oldpage[IV-3]{16}
sont commutatifs~; on en déduit donc par passage à la limite inductive un
diagramme commutatif
\begin{equation}
\xymatrix{
  \varinjlim\mathfrak P(S_\lambda) \ar[r]^{u_{\mathfrak P}}
    \ar[d]_{\varinjlim g_\lambda^{-1}}
    & \mathfrak P(S) \ar[d]^{g^{-1}} \\
  \varinjlim\mathfrak P(X_\lambda) \ar[r]_{v_{\mathfrak P}}
    & \mathfrak P(X)
}
\tag{8.3.10.1}\label{IV.8.3.10.1-fr}
\end{equation}
et il résulte de (1.8.2) que l'on a des diagrammes analogues en remplaçant
$\mathfrak P$ par $\mathfrak C$, $\mathfrak D_c$, $\mathfrak F_c$ ou
$\mathfrak{LF}_c$.
\end{env}

Il résulte de (8.3.5) que sous l'hypothèse que pour un $\alpha\in L$,
$S_\alpha$ est \emph{quasi-compact}, l'application canonique (8.3.9.2) est
\emph{injective}. En outre~:

\begin{theorem}[8.3.11]
\label{IV.8.3.11-fr}
Supposons qu'il existe un $\alpha\in L$ tel que $S_\alpha$ soit quasi-compact
et quasi-séparé. Alors les applications canoniques (8.3.9.2), (8.3.9.3),
(8.3.9.4) et (8.3.9.5) sont \emph{bijectives}.
\end{theorem}

En vertu de la remarque précédente, il reste à prouver que ces applications
sont \emph{surjectives}~; comme toute partie constructible de $S$ est réunion
finie d'ensembles de la forme $U\cap\complement V$, où $U$ et $V$ sont
ouverts et constructibles, il suffira de prouver que (8.3.9.3) est surjective
pour qu'il en soit de même de (8.3.9.2) (et aussi de (8.3.9.4), par passage
aux complémentaires). Or, comme les morphismes $S_\lambda\to S_\alpha$ et
$S\to S_\alpha$ sont affines, donc séparés, les $S_\lambda$ pour
$\lambda\geq\alpha$ et $S$ sont quasi-compacts et quasi-séparés (1.2.2), et
l'on sait que les parties ouvertes constructibles dans un tel préschéma ne
sont autres que les parties ouvertes quasi-compactes (1.8.1). La conclusion
résulte donc de (8.2.11) sauf pour l'application (8.3.9.5). Pour prouver que
cette dernière est surjective, considérons une partie $Z$ localement fermée
et constructible dans $S$~; $Z$ est donc quasi-compact
$(0_{\mathrm{III}},9.1.4)$. Comme tout point $x\in Z$ admet par hypothèse un
voisinage ouvert quasi-compact $V_x$ dans $S$ tel que $Z\cap V_x$ soit fermé
dans $V_x$, on peut couvrir $Z$ par un nombre fini des $V_x$~; autrement dit,
il y a un ouvert quasi-compact $U$ contenant $Z$ et tel que $Z$ soit fermé
dans $U$~; comme $Z$ est constructible dans $S$, il l'est aussi dans $U$
$(0_{\mathrm{III}},9.1.8)$. On sait (8.2.11) qu'il y a un indice $\lambda$ et
un ouvert quasi-compact $U_\lambda$ dans $S_\lambda$ tels que
$U=u_\lambda^{-1}(U_\lambda)$. Appliquant à $U$ (qui est limite projective
des $U_\mu=u_{\lambda\mu}^{-1}(U_\lambda)$ pour $\mu\geq\lambda$) le fait
que l'application (8.3.9.4) est surjective, on voit qu'il existe un
$\mu\geq\lambda$ et un ensemble fermé constructible $Z_\mu$ dans $U_\mu$
tels que $Z=u_\mu^{-1}(Z_\mu)$. Mais comme l'immersion canonique
$U_\mu\to S_\mu$ est quasi-compacte par hypothèse (1.2.7), elle est de
présentation finie (1.6.2, (i)), et $Z_\mu$ est aussi une partie constructible
\emph{de} $S_\mu$ en vertu de (1.8.4) et (1.8.1)~; comme $Z_\mu$ est
évidemment localement fermée dans $S_\mu$, cela achève la démonstration.

\begin{corollary}[8.3.12]
\label{IV.8.3.12-fr}
Supposons qu'il existe un $\alpha$ tel que $S_\alpha$ soit quasi-compact, et
soit, pour tout $\lambda$, $Z_\lambda$ une partie constructible de
$S_\lambda$ telle que
$Z_\mu=u_{\lambda\mu}^{-1}(Z_\lambda)$ pour $\mu\geq\lambda$. Si
$Z=u_\lambda^{-1}(Z_\lambda)$,
\oldpage[IV-3]{17}
alors, pour que $Z$ soit ouverte (resp. fermée, resp. localement fermée) dans
$S$, il faut et il suffit qu'il existe $\lambda\geq\alpha$ tel que
$Z_\lambda$ le soit dans $S_\lambda$.
\end{corollary}

Couvrons $S_\alpha$ par un nombre fini d'ouverts affines
$U_\alpha^{(i)}$~; alors les $U^{(i)}=u_\alpha^{-1}(U_\alpha^{(i)})$ forment
un recouvrement ouvert affine de $S$, et pour que $Z$ soit ouvert (resp.
fermé, resp. localement fermé) dans $S$, il faut et il suffit que chacun des
$Z\cap U^{(i)}$ le soit dans $U^{(i)}$. Comme $L$ est filtrant et que chacun
des $Z\cap U^{(i)}$ est constructible dans $U^{(i)}$
$(0_{\mathrm{III}},9.1.8)$, on peut se borner à démontrer le corollaire
lorsque $S_\alpha$ est affine, donc quasi-compact et quasi-séparé~; mais alors
il résulte de (8.3.11).

\begin{proposition}[8.3.13]
\label{IV.8.3.13-fr}
Supposons que les morphismes $u_{\lambda\mu}:S_\mu\to S_\lambda$ soient
plats pour $\lambda\leq\mu$, et qu'il existe $\alpha$ tel que $S_\alpha$ soit
quasi-compact. Pour tout $\lambda$, soient $Z'_\lambda$, $Z''_\lambda$ deux
parties pro-constructibles de $S_\lambda$, telles que
$Z'_\mu=u_{\lambda\mu}^{-1}(Z'_\lambda)$ et
$Z''_\mu=u_{\lambda\mu}^{-1}(Z''_\lambda)$ pour $\mu\geq\lambda$~; on
suppose en outre que $\overline{Z'_\alpha}$ est constructible dans
$S_\alpha$. Soient $Z'=u_\lambda^{-1}(Z'_\lambda)$,
$Z''=u_\lambda^{-1}(Z''_\lambda)$~; pour que
$Z''\subset\overline{Z'}$, il faut et il suffit qu'il existe
$\lambda\geq\alpha$ tel que
$Z''_\lambda\subset\overline{Z'_\lambda}$.
\end{proposition}

En effet, on sait que $u_\lambda:S\to S_\lambda$ est aussi un morphisme plat
pour tout $\lambda$ (8.3.8)~; comme $Z'_\lambda$ est pro-constructible,
l'adhérence de $Z'_\mu$ pour $\mu\geq\lambda$ (resp. de $Z'$) dans $S_\mu$
(resp. $S$) est égale à
$u_{\lambda\mu}^{-1}(\overline{Z'_\lambda})$ (resp.
$u_\lambda^{-1}(\overline{Z'_\lambda})$) (2.3.10). Comme les
$u_{\lambda\mu}^{-1}(\overline{Z'_\lambda})$ et
$u_\lambda^{-1}(\overline{Z'_\lambda})$ sont constructibles (1.8.2), la
conclusion résulte de (8.3.2).

\subsection{Critères d'irréductibilité et de connexion pour les limites
projectives de préschémas}
\label{subsection:IV.8.4-fr}

\begin{proposition}[8.4.1]
\label{IV.8.4.1-fr}
Supposons qu'il existe un indice $\alpha$ tel que $S_\alpha$ soit
quasi-compact.
\begin{enumerate}
\item[(i)] Si $S$ n'est pas irréductible et si de plus l'espace sous-jacent à
$S$ est noethérien et $S_\alpha$ quasi-séparé, il existe
$\lambda\geq\alpha$ tel que, pour $\mu\geq\lambda$, $S_\mu$ ne soit pas
irréductible.
\item[(ii)] Si $S$ n'est pas connexe, il existe $\lambda\geq\alpha$ tel que,
pour $\mu\geq\lambda$, $S_\mu$ ne soit pas connexe.
\end{enumerate}
\end{proposition}

Supposons que $S$ soit réunion de deux ensembles fermés $S'$, $S''$ distincts
de $S$ (resp. fermés disjoints et non vides). Dans le cas (i), $S'$ et $S''$
sont constructibles puisque l'espace $S$ est noethérien. En vertu de
(8.3.11), il existe donc un $\lambda\geq\alpha$ et deux ensembles fermés
constructibles $S'_\lambda$, $S''_\lambda$ de $S_\lambda$ tels que
$S'=u_\lambda^{-1}(S'_\lambda)$,
$S''=u_\lambda^{-1}(S''_\lambda)$~; comme $S=S'\cup S''$, il résulte aussi
de (8.3.11) que l'on peut supposer que
$S_\lambda=S'_\lambda\cup S''_\lambda$~; comme $S'_\lambda$ et
$S''_\lambda$ sont distincts de $S_\lambda$, cela prouve que $S_\lambda$
n'est pas irréductible.

Dans le cas (ii), $S'$ et $S''$ sont ouverts quasi-compacts, donc, en vertu
de (8.2.11), il existe un $\lambda\geq\alpha$ et deux ensembles ouverts
quasi-compacts $S'_\lambda$, $S''_\lambda$ de $S_\lambda$ tels que
$S'=u_\lambda^{-1}(S'_\lambda)$,
$S''=u_\lambda^{-1}(S''_\lambda)$. D'ailleurs, comme $S'$ et $S''$ sont
ouverts et fermés dans $S$, ils sont à la fois pro-constructibles et
ind-constructibles (1.9.6), donc constructibles (1.9.11), et il résulte donc
de (8.3.5) que l'on peut supposer $\lambda$ pris tel que
$S_\lambda=S'_\lambda\cup S''_\lambda$ et
$S'_\lambda\cap S''_\lambda=\emptyset$, ce qui montre que $S_\lambda$
n'est pas connexe.

\begin{proposition}[8.4.2]
\label{IV.8.4.2-fr}
Supposons que l'espace sous-jacent à $S$ soit noethérien et que l'une des deux
conditions suivantes soit vérifiée~:
\begin{enumerate}
\item[\emph{a)}] Pour $\lambda\leq\mu$,
$u_{\lambda\mu}:S_\mu\to S_\lambda$ est dominant, et il existe $\alpha$ tel
que $S_\alpha$ soit quasi-compact.
\oldpage[IV-3]{18}
\item[\emph{b)}] Il existe $\alpha$ tel que l'espace sous-jacent à
$S_\alpha$ soit noethérien, et pour $\mu\geq\lambda$, $u_{\lambda\mu}$ est
un homéomorphisme de $S_\mu$ sur un sous-espace de $S_\lambda$.
\end{enumerate}
Sous ces conditions, il existe un $\lambda$ tel que, pour tout
$\mu\geq\lambda$~:
\begin{enumerate}
\item[(i)] Pour toute composante irréductible $Y_i$ de $S$
$(1\leq i\leq m)$, $\overline{u_\mu(Y_i)}$ est une composante irréductible
de $S_\mu$, et l'application $Y_i\mapsto\overline{u_\mu(Y_i)}$ est une
bijection de l'ensemble des composantes irréductibles de $S$ sur l'ensemble
des composantes irréductibles de $S_\mu$.
\item[(ii)] Pour toute composante connexe $C_j$ de $S$
$(1\leq j\leq n)$, $\overline{u_\mu(C_j)}$ est une composante connexe de
$S_\mu$, et l'application $C_j\mapsto\overline{u_\mu(C_j)}$ est une
bijection de l'ensemble des composantes connexes de $S$ sur l'ensemble des
composantes connexes de $S_\mu$.
\end{enumerate}
\end{proposition}

Nous établirons d'abord le

\begin{lemma}[8.4.2.1]
\label{IV.8.4.2.1-fr}
Sous la condition \emph{a)} ou \emph{b)} de (8.4.2), il existe $\lambda$ tel
que, pour $\mu\geq\lambda$, $u_\mu:S\to S_\mu$ soit dominant.
\end{lemma}

Dans le cas \emph{a)}, cela a déjà été prouvé sans supposer l'espace $S$
noethérien (8.3.8, (i)). Dans le cas \emph{b)}, posons
$Z_\alpha=\overline{u_\alpha(S)}$~; en tant que partie fermée de l'espace
noethérien $S_\alpha$, $Z_\alpha$ est constructible, et comme
$u_\alpha^{-1}(Z_\alpha)=S$, il résulte de (8.3.5) qu'il existe
$\lambda\geq\alpha$ tel que, pour $\mu\geq\lambda$, on ait
$Z_\mu=u_{\alpha\mu}^{-1}(Z_\alpha)=S_\mu$. Mais comme $u_{\alpha\mu}$ est
un homéomorphisme de $S_\mu$ sur un sous-espace de $S_\alpha$, et comme
l'application composée $S\to Z_\mu\to Z_\alpha$ est dominante, il en est de
même de $S\to Z_\mu=S_\mu$.

Ce lemme étant démontré, on peut supposer que pour tout $\lambda$,
$u_\lambda$ est un morphisme dominant.

\medskip
\noindent
(i) Chacun des $S_\lambda$ est réunion des
$\overline{u_\lambda(Y_i)}$, qui sont irréductibles. D'autre part, si $U_i$
est l'ouvert de $S$ complémentaire de la réunion des $Y_j$ d'indice $j\neq i$
$(1\leq i\leq m)$, les $U_i$ sont deux à deux disjoints et
$Y_i=\overline{U_i}$ $(0_{\mathrm I},2.1.6)$. Comme l'espace sous-jacent à
$S$ est noethérien, les $U_i$ sont quasi-compacts, donc il existe un indice
$\lambda$ et des ouverts $U_{i\lambda}$ de $S_\lambda$ tels que
$U_i=u_\lambda^{-1}(U_{i\lambda})$ pour $1\leq i\leq m$ (8.2.11). On en
conclut que si l'on pose
$U_{i\mu}=u_{\lambda\mu}^{-1}(U_{i\lambda})$ pour $\lambda\leq\mu$, les
$U_{i\mu}$ sont deux à deux disjoints, car les
$U_i=u_\mu^{-1}(U_{i\mu})$ le sont et $u_\mu$ est dominant. Par suite,
aucune des adhérences $\overline{U_{i\mu}}$ n'est contenue dans une autre et
$u_\mu(U_i)$ est dense dans $U_{i\mu}$ puisque $u_\mu$ est dominant~; on a
donc
$\overline{U_{i\mu}}=\overline{u_\mu(Y_i)}$, ce qui prouve que les
$\overline{U_{i\mu}}$ sont les composantes irréductibles de $S_\mu$
$(0_{\mathrm I},2.1.7)$ et achève la démonstration.

\medskip
\noindent
(ii) Comme l'espace $S$ est noethérien, les $C_j$ sont ouverts et fermés dans
$S$ et quasi-compacts~; le même raisonnement que dans (i) montre donc qu'il
existe un $\lambda$ et des ouverts $V_{j\lambda}$ de $S_\lambda$ tels que
$C_j=u_\lambda^{-1}(V_{j\lambda})$ pour $1\leq j\leq n$. On voit aussi,
comme dans (i), que si l'on pose
$V_{j\mu}=u_{\lambda\mu}^{-1}(V_{j\lambda})$ pour $\mu\geq\lambda$, les
$V_{j\mu}$ sont deux à deux disjoints, et $u_\mu(C_j)$ est dense dans
$V_{j\mu}$~; cela entraîne que $V_{j\mu}$ est connexe. En outre, il résulte
de (8.3.4) que pour $\mu$ assez grand, la réunion des $V_{j\mu}$ est
$S_\mu$, puisque tout ouvert dans un préschéma est ind-constructible (1.9.6).
Les $V_{j\mu}$ sont donc les composantes connexes de $S_\mu$, ce qui achève
la démonstration.

On notera que si les morphismes $u_{\lambda\mu}$ sont des \emph{immersions},
ils vérifieront en particulier la condition \emph{b)} de (8.4.2).

\oldpage[IV-3]{19}
\begin{corollary}[8.4.3]
\label{IV.8.4.3-fr}
Supposons vérifiée l'une des conditions \emph{a)}, \emph{b)} de (8.4.2),
l'espace sous-jacent à $S$ étant noethérien~; alors, pour que $S$ soit
irréductible, il faut et il suffit qu'il existe $\lambda$ tel que $S_\mu$ le
soit pour tout $\mu\geq\lambda$.
\end{corollary}

\begin{proposition}[8.4.4]
\label{IV.8.4.4-fr}
Supposons qu'il existe un $\alpha$ tel que $S_\alpha$ soit quasi-compact et
que, pour $\lambda\leq\mu$, $u_{\lambda\mu}:S_\mu\to S_\lambda$ soit dominant.
Alors, pour que $S$ soit connexe, il faut et il suffit qu'il existe $\lambda$
tel que $S_\mu$ le soit pour tout $\mu\geq\lambda$.
\end{proposition}

\begin{proof}
La condition est suffisante en vertu de (8.4.1)~; d'autre part, on a vu
(8.3.8, (i)) que $u_\lambda:S\to S_\lambda$ est dominant pour $\lambda$ assez
grand, donc, si $S$ est connexe, il en est de même de $S_\lambda$, puisque
$u_\lambda(S)$ est dense dans $S_\lambda$ et connexe.
\end{proof}

\begin{corollary}[8.4.5]
\label{IV.8.4.5-fr}
Soient $k$ un corps, $X$ un $k$-préschéma quasi-compact. Pour que $X$ soit
géométriquement connexe (4.5.2), il faut et il suffit que, pour toute extension
finie et séparable $K$ de $k$, $X\otimes_k K$ soit connexe.
\end{corollary}

\begin{proof}
La condition est trivialement nécessaire. Pour voir qu'elle est suffisante,
nous devons prouver que si $\Omega$ est une clôture algébrique de $k$,
$X\otimes_k\Omega$ est connexe (4.5.1). Or, $\Omega$ est limite inductive
filtrante des sous-extensions finies $k'$ de $k$, et pour
$k\subset k'\subset k''\subset\Omega$, le morphisme
$X\otimes_k k''\to X\otimes_k k'$ est surjectif. On est donc ramené, en vertu
de (8.4.4), à prouver que $X\otimes_k k'$ est connexe pour toute extension
finie $k'$ de $k$. Mais si $K$ est la plus grande extension séparable contenue
dans $k'$, le morphisme $X\otimes_k k'\to X\otimes_k K$ est fini, surjectif et
radiciel, donc (2.4.5) un homéomorphisme, et comme $X\otimes_k K$ est connexe
par hypothèse, il en est de même de $X\otimes_k k'$.
\end{proof}

\begin{remark}[8.4.6]
\label{IV.8.4.6-fr}
\begin{enumerate}[label=(\roman*)]
  \item La démonstration de (8.4.2) montre que la conclusion de cette
  proposition est valable si l'on suppose que l'espace sous-jacent à $S$ est
  noethérien, qu'il existe $\alpha$ tel que $S_\alpha$ soit quasi-compact, et
  enfin que les $u_\lambda:S\to S_\lambda$ sont dominants.

  \item Par contre, la conclusion de (8.4.2) peut être en défaut lorsque les
  $u_\lambda$ ne sont pas dominants pour $\lambda$ assez grand, même lorsque
  les $S_\lambda$ et $S$ sont noethériens, ainsi que le montre l'exemple
  suivant. Prenons pour ensemble d'indices $\mathbf N$, tous les $S_n$ étant
  égaux à
  $\operatorname{Spec}(A\times K)=\operatorname{Spec}(A)\amalg
  \operatorname{Spec}(K)$, où $K$ est un corps, $A$ une $K$-algèbre
  quelconque, et tous les morphismes $u_{n,n+1}$ étant égaux au même morphisme
  correspondant à l'homomorphisme $(x,y)\mapsto(j(y),y)$ de $A\times K$ dans
  lui-même, où $j:K\to A$ est l'homomorphisme canonique. On vérifie aisément
  que la limite inductive de ce système d'anneaux est $K$, l'homomorphisme
  canonique $u_n$ correspondant à la seconde projection $A\times K\to K$. On
  voit donc que $S=\operatorname{Spec}(K)$ est irréductible bien qu'aucun des
  $S_n$ ne soit connexe.
\end{enumerate}
\end{remark}

\subsection{Modules de présentation finie sur une limite projective de préschémas}
\label{subsection:IV.8.5-fr}

\begin{env}[8.5.1]
\label{IV.8.5.1-fr}
Nous conservons toujours les notations de (8.2.2)~; nous nous bornerons en
outre au cas où $S_0$ est l'un des $S_\lambda$, auquel on peut toujours se
ramener.

\emph{Lorsque, dans cette section, nous considérerons une famille
$(\mathcal F_\lambda)$, où, pour tout $\lambda\in L$, $\mathcal F_\lambda$ est
un $\mathcal O_{S_\lambda}$-Module, il sera sous-entendu que cette famille
vérifie la condition}
\begin{equation}
  \mathcal F_\mu=u_{\lambda\mu}^*(\mathcal F_\lambda)
  \qquad\text{pour }\lambda\leq\mu.
  \tag{8.5.1.1}\label{IV.8.5.1.1-fr}
\end{equation}

\oldpage[IV-3]{20}
Nous poserons alors
\begin{equation}
  \mathcal F=u_\lambda^*(\mathcal F_\lambda)
  \tag{8.5.1.2}\label{IV.8.5.1.2-fr}
\end{equation}
qui est un $\mathcal O_S$-Module ne dépendant pas de l'indice $\lambda\in L$,
en vertu de l'hypothèse (8.5.1.1).

Soient maintenant $(\mathcal F_\lambda)$, $(\mathcal G_\lambda)$ deux telles
familles de $\mathcal O_{S_\lambda}$-Modules. Il est clair que les applications
$f_\lambda\mapsto u_{\lambda\mu}^*(f_\lambda)$ de
$\operatorname{Hom}_{S_\lambda}(\mathcal F_\lambda,\mathcal G_\lambda)$ dans
$\operatorname{Hom}_{S_\mu}(\mathcal F_\mu,\mathcal G_\mu)$ définissent alors
un \emph{système inductif} de groupes abéliens
$(\operatorname{Hom}_{S_\lambda}(\mathcal F_\lambda,\mathcal G_\lambda),
u_{\lambda\mu}^*)$, et que les applications
$f_\lambda\mapsto u_\lambda^*(f_\lambda)$ forment un \emph{système inductif}
d'homomorphismes de groupes abéliens
\[
  u_\lambda^*:\operatorname{Hom}_{S_\lambda}
  (\mathcal F_\lambda,\mathcal G_\lambda)
  \longrightarrow\operatorname{Hom}_S(\mathcal F,\mathcal G)~;
\]
d'où, en passant à la limite inductive, un homomorphisme canonique de groupes
abéliens
\begin{equation}
  u_{\mathcal F,\mathcal G}:
  \varinjlim\operatorname{Hom}_{S_\lambda}
  (\mathcal F_\lambda,\mathcal G_\lambda)
  \longrightarrow\operatorname{Hom}_S(\mathcal F,\mathcal G).
  \tag{8.5.1.3}\label{IV.8.5.1.3-fr}
\end{equation}

Notons que lorsque $\mathcal F_\lambda=\mathcal O_{S_\lambda}$, la condition
(8.5.1.1) est vérifiée, et que l'on a $\mathcal F=\mathcal O_S$~;
l'homomorphisme (8.5.1.3) donne donc $(0_{\mathrm I},5.1.1)$ un homomorphisme
canonique de groupes abéliens
\begin{equation}
  u_{\mathcal G}:\varinjlim\Gamma(S_\lambda,\mathcal G_\lambda)
  \longrightarrow\Gamma(S,\mathcal G).
  \tag{8.5.1.4}\label{IV.8.5.1.4-fr}
\end{equation}
\end{env}

\begin{theorem}[8.5.2]
\label{IV.8.5.2-fr}
\begin{enumerate}[label=(\roman*)]
  \item Supposons $S_0$ quasi-compact (resp. quasi-compact et quasi-séparé) et
  que, pour un $\lambda\in L$, $\mathcal F_\lambda$ soit quasi-cohérent et de
  type fini (resp. de présentation finie) et $\mathcal G_\lambda$
  quasi-cohérent. Alors l'homomorphisme $u_{\mathcal F,\mathcal G}$ est
  injectif (resp. bijectif).

  \item Supposons $S_0$ quasi-compact et quasi-séparé. Pour tout
  $\mathcal O_S$-Module quasi-cohérent $\mathcal F$ de présentation finie, il
  existe un $\lambda\in L$ et un $\mathcal O_{S_\lambda}$-Module
  quasi-cohérent $\mathcal F_\lambda$ de présentation finie tels que
  $\mathcal F$ soit isomorphe à $u_\lambda^*(\mathcal F_\lambda)$.
\end{enumerate}
\end{theorem}

\begin{proof}
\begin{enumerate}[label=(\roman*)]
  \item On peut évidemment se borner au cas où $S_0=S_\lambda$ puisque les
  morphismes $u_{0\lambda}:S_\lambda\to S_0$ sont affines, donc quasi-compacts
  et séparés. Considérons d'abord le cas où
  $S_0=\operatorname{Spec}(A_0)$ est \emph{affine}. Alors l'assertion (i)
  équivaut au

  \begin{lemma}[8.5.2.1]
  \label{IV.8.5.2.1-fr}
  Soient $A_0$ un anneau, $(A_\lambda)_{\lambda\in L}$ un système inductif de
  $A_0$-algèbres, $A=\varinjlim A_\lambda$~; soient $M_0$, $N_0$ deux
  $A_0$-modules, et posons
  $M_\lambda=M_0\otimes_{A_0}A_\lambda$,
  $N_\lambda=N_0\otimes_{A_0}A_\lambda$,
  $M=M_0\otimes_{A_0}A=\varinjlim M_\lambda$,
  $N=N_0\otimes_{A_0}A=\varinjlim N_\lambda$. Si $M_0$ est de type fini
  (resp. de présentation finie), l'homomorphisme canonique
  \begin{equation}
    \varinjlim\operatorname{Hom}_{A_\lambda}(M_\lambda,N_\lambda)
    \longrightarrow\operatorname{Hom}_A(M,N)
    \tag{8.5.2.2}\label{IV.8.5.2.2-fr}
  \end{equation}
  est injectif (resp. bijectif).
  \end{lemma}

  On sait en effet (Bourbaki, \emph{Alg.}, chap. II, 3\textsuperscript{e} éd.,
  \S~5, n\textsuperscript{o}~1) que l'on a des isomorphismes canoniques
  fonctoriels
  \[
    \operatorname{Hom}_{A_\lambda}(M_\lambda,N_\lambda)
    \xrightarrow{\sim}\operatorname{Hom}_{A_0}(M_0,N_\lambda),
    \qquad
    \operatorname{Hom}_A(M,N)
    \xrightarrow{\sim}\operatorname{Hom}_{A_0}(M_0,N)
  \]

  \oldpage[IV-3]{21}
  si bien que l'homomorphisme (8.5.2.2) n'est autre, à des isomorphismes
  canoniques près, que l'homomorphisme canonique
  \begin{equation}
    \varinjlim\operatorname{Hom}_{A_0}(M_0,N_\lambda)
    \longrightarrow
    \operatorname{Hom}_{A_0}(M_0,\varinjlim N_\lambda),
    \tag{8.5.2.3}\label{IV.8.5.2.3-fr}
  \end{equation}
  qui, à tout système inductif d'homomorphismes de $A_0$-modules
  $\theta_\lambda:M_0\to N_\lambda$, fait correspondre sa limite inductive.

  Or, si $M_0$ est de type fini (resp. de présentation finie), on a une suite
  exacte $A_0^m\to M_0\to0$ (resp. $A_0^n\to A_0^m\to M_0\to0$)~; comme il est
  clair que (8.5.2.3) est bijectif lorsque $M_0$ est de la forme $A_0^q$, il
  suffit d'utiliser l'exactitude à gauche du foncteur
  $M_0\mapsto\operatorname{Hom}_{A_0}(M_0,P)$ et l'exactitude du foncteur
  $\varinjlim$ (dans la catégorie des groupes abéliens) pour conclure.

  Passons au cas où $S_0$ est quasi-compact, et soit $(U_i)$ un recouvrement
  fini de $S_0$ par des ouverts affines~; pour tout $\lambda$, les
  $U_{i\lambda}=u_{0\lambda}^{-1}(U_i)$ forment un recouvrement ouvert affine de
  $S_\lambda$, et les $V_i=u_0^{-1}(U_i)$ un recouvrement ouvert affine de $S$.
  Pour voir que $u_{\mathcal F,\mathcal G}$ est injectif, il faut prouver que
  si $f_\lambda:\mathcal F_\lambda\to\mathcal G_\lambda$ est tel que
  $f=u_\lambda^*(f_\lambda)=0$, alors il existe $\mu\geq\lambda$ tel que
  $f_\mu=u_{\lambda\mu}^*(f_\lambda)=0$. En vertu du lemme (8.5.2.1), pour
  chaque $i$ il existe un $\lambda_i$ tel que
  $f_\mu|U_{i\mu}=0$ pour $\mu\geq\lambda_i$. Il suffit donc de prendre $\mu$
  plus grand que tous les $\lambda_i$.

  Supposons en outre $S_0$ quasi-séparé et $\mathcal F_0$ de présentation finie,
  et soit $f:\mathcal F\to\mathcal G$ un homomorphisme de
  $\mathcal O_S$-Modules. En vertu du lemme (8.5.2.1), pour tout $i$, il existe
  un indice $\lambda_i$ et un homomorphisme
  $f_{\lambda_i}^{(i)}:\mathcal F_{\lambda_i}|U_{i,\lambda_i}
  \to\mathcal G_{\lambda_i}|U_{i,\lambda_i}$ tels que
  $u_{\lambda_i}^*(f_{\lambda_i}^{(i)})=f|V_i$. Comme $L$ est filtrant, on peut
  en outre supposer tous les $\lambda_i$ égaux à un même $\lambda$. Notons
  maintenant que $S_\lambda$ est quasi-séparé (1.2.3) et
  $\mathcal F_\lambda$ est un $\mathcal O_{S_\lambda}$-Module quasi-cohérent de
  présentation finie $(0_{\mathrm I},5.2.5)$~; comme, pour tout couple d'indices
  $i,j$, $U_{ij\lambda}=U_{i\lambda}\cap U_{j\lambda}$ est quasi-compact et que
  l'on a
  \[
    u_\lambda^*(f_\lambda^{(i)}|U_{ij\lambda})
    =u_\lambda^*(f_\lambda^{(j)}|U_{ij\lambda})
    =f|(V_i\cap V_j)
  \]
  par définition, il résulte de ce qu'on a vu plus haut qu'il existe un indice
  $\lambda_{ij}$ tel que
  \[
    u_{\lambda\mu}^*(f_\lambda^{(i)}|U_{ij\lambda})
    =u_{\lambda\mu}^*(f_\lambda^{(j)}|U_{ij\lambda})
    \quad\text{pour }\mu\geq\lambda_{ij}~;
  \]
  prenant $\mu$ plus grand que tous les $\lambda_{ij}$, on voit donc que
  $u_{\lambda\mu}^*(f_\lambda^{(i)})$ et
  $u_{\lambda\mu}^*(f_\lambda^{(j)})$ coïncident dans
  $U_{i\mu}\cap U_{j\mu}$ pour tout couple $(i,j)$, et par suite définissent
  un homomorphisme $f_\mu:\mathcal F_\mu\to\mathcal G_\mu$ tel que
  $f=u_\mu^*(f_\mu)$.
\end{enumerate}
\end{proof}

Avant de passer à la démonstration de (ii), notons les corollaires suivants de
(i)~:

\begin{corollary}[8.5.2.4]
\label{IV.8.5.2.4-fr}
Supposons $S_0$ quasi-compact, $\mathcal F_\lambda$ quasi-cohérent de type fini,
$\mathcal G_\lambda$ quasi-cohérent de présentation finie. Soit
$f_\lambda:\mathcal F_\lambda\to\mathcal G_\lambda$ un homomorphisme. Pour que
$f=u_\lambda^*(f_\lambda):\mathcal F\to\mathcal G$ soit un isomorphisme, il faut
et il suffit qu'il existe $\mu\geq\lambda$ tel que
$f_\mu=u_{\lambda\mu}^*(f_\lambda):\mathcal F_\mu\to\mathcal G_\mu$ soit un
isomorphisme.
\end{corollary}

\begin{proof}
On peut toujours supposer $S_\lambda=S_0$~; la question étant locale sur $S_0$,
on peut en outre ($S_0$ étant quasi-compact et $L$ filtrant) se ramener au cas
où $S_0$ est affine, donc quasi-séparé. La condition étant trivialement
suffisante, il reste à montrer qu'elle est nécessaire~: or, il y a par
hypothèse un $\mathcal O_S$-homomorphisme $g:\mathcal G\to\mathcal F$ tel que
$g\circ f=1_{\mathcal F}$ et $f\circ g=1_{\mathcal G}$. Comme $\mathcal G$ est
de présentation finie, il existe $\nu\geq\lambda$ et un homomorphisme
$g_\nu:\mathcal G_\nu\to\mathcal F_\nu$ tels que $g=u_\nu^*(g_\nu)$ en vertu de
(8.5.2, (i))~; on a par suite
$u_\nu^*(g_\nu\circ f_\nu)=1_{\mathcal F}$

\oldpage[IV-3]{22}
et $u_\nu^*(f_\nu\circ g_\nu)=1_{\mathcal G}$~; compte tenu de ce que
$\mathcal F_\nu$ et $\mathcal G_\nu$ sont de type fini, on en conclut par
(8.5.2, (i)) qu'il existe $\mu\geq\nu$ tel que l'on ait
$g_\mu\circ f_\mu=1_{\mathcal F_\mu}$ et
$f_\mu\circ g_\mu=1_{\mathcal G_\mu}$, d'où le corollaire.
\end{proof}

\begin{corollary}[8.5.2.5]
\label{IV.8.5.2.5-fr}
Supposons $S_0$ quasi-compact et quasi-séparé. Supposons que
$\mathcal F_\lambda$, $\mathcal G_\lambda$ soient des
$\mathcal O_{S_\lambda}$-Modules quasi-cohérents de présentation finie. Pour
que $\mathcal F$ et $\mathcal G$ soient isomorphes, il faut et il suffit qu'il
existe $\mu\geq\lambda$ tel que $\mathcal F_\mu$ et $\mathcal G_\mu$ soient
isomorphes. De plus, pour tout isomorphisme
$f:\mathcal F\xrightarrow{\sim}\mathcal G$, il existe un $\nu\geq\mu$ et un
isomorphisme $f_\nu:\mathcal F_\nu\xrightarrow{\sim}\mathcal G_\nu$ tels que
$f=u_\nu^*(f_\nu)$.
\end{corollary}

\begin{proof}
Cela résulte de (8.5.2.4) et de (8.5.2, (i)) puisque tout homomorphisme
$f:\mathcal F\to\mathcal G$ est de la forme $u_\mu^*(f_\mu)$ pour un
$\mu\geq\lambda$ et un homomorphisme
$f_\mu:\mathcal F_\mu\to\mathcal G_\mu$.
\end{proof}

\par\medskip\noindent\textup{(ii)} Considérons encore en premier lieu le cas où
$S_0=\operatorname{Spec}(A_0)$ est
affine. Alors l'assertion équivaut au lemme (5.13.7.1).

Dans le cas général, en partant d'un recouvrement ouvert affine fini $(U_i)$
de $S_0$, on déduit de (5.13.7.1) que pour tout $i$, il existe un indice
$\lambda(i)$ et un $\mathcal O_{U_{i,\lambda(i)}}$-Module quasi-cohérent de
présentation finie $\mathcal F_i$ tels que
$u_{\lambda(i)}^*(\mathcal F_i)=\mathcal F|V_i$ (avec les notations de (i)). En
outre, comme $L$ est filtrant, on peut supposer que tous les $\lambda_i$ sont
égaux à un même $\lambda$. Comme
$U_{ij\lambda}=U_{i\lambda}\cap U_{j\lambda}$ est quasi-compact et
quasi-séparé (1.2.7), il résulte de (8.5.2.5) que pour tout couple $(i,j)$, il
existe un indice $\lambda_{ij}\geq\lambda$ et un isomorphisme
\[
  \theta_{ij}:u_{\lambda,\lambda_{ij}}^*(\mathcal F_i|U_{ij\lambda})
  \xrightarrow{\sim}
  u_{\lambda,\lambda_{ij}}^*(\mathcal F_j|U_{ij\lambda})
\]
tel que $u_{\lambda_{ij}}^*(\theta_{ij})$ soit l'automorphisme identique de
$\mathcal F|(V_i\cap V_j)$~; on peut encore supposer tous les
$\lambda_{ij}$ égaux à $\lambda$. Changeant les notations, on peut donc
supposer qu'il existe pour tout couple $(i,j)$ un isomorphisme
\[
  \theta_{ij}:\mathcal F_i|U_{ij\lambda}
  \xrightarrow{\sim}\mathcal F_j|U_{ij\lambda},
\]
tels que $u_\lambda^*(\theta_{ij})$ soit l'automorphisme identique de
$\mathcal F|(V_i\cap V_j)$. Enfin, pour trois indices quelconques $i,j,k$, si
l'on pose
$U_{ijk,\lambda}=U_{i\lambda}\cap U_{j\lambda}\cap U_{k\lambda}$,
$U_{ijk,\lambda}$ est quasi-compact, et si $\theta'_{ij}$, $\theta'_{jk}$ et
$\theta'_{ik}$ désignent les restrictions de $\theta_{ij}$, $\theta_{jk}$ et
$\theta_{ik}$ à $U_{ijk,\lambda}$, on a
\[
  u_\lambda^*(\theta'_{ij}\circ\theta'_{jk})=u_\lambda^*(\theta'_{ik}).
\]
Il y a donc, en vertu de (i), un indice $\mu\geq\lambda$ tel que l'on ait
\[
  u_{\lambda\mu}^*(\theta'_{ik})
  =u_{\lambda\mu}^*(\theta'_{ij}\circ\theta'_{jk}),
\]
donc les isomorphismes
\[
  u_{\lambda\mu}^*(\theta'_{ij}):
  u_{\lambda\mu}^*(\mathcal F_i)|U_{ij\mu}
  \xrightarrow{\sim}
  u_{\lambda\mu}^*(\mathcal F_j)|U_{ij\mu}
\]
vérifient la condition de recollement, et définissent par suite sur $S_\mu$ un
$\mathcal O_{S_\mu}$-Module quasi-cohérent $\mathcal F_\mu$ de présentation
finie $(0_{\mathrm I},3.3.1)$ tel que $\mathcal F$ soit isomorphe à
$u_\mu^*(\mathcal F_\mu)$.

\begin{env}[8.5.3]
\label{IV.8.5.3-fr}
Le résultat de (8.5.2) s'exprime encore en disant que si $S_0$ est
quasi-compact et quasi-séparé, la catégorie des $\mathcal O_S$-Modules
quasi-cohérents de présentation finie est déterminée à équivalence près par la
donnée des catégories des $\mathcal O_{S_\lambda}$-Modules quasi-cohérents de
présentation finie, des foncteurs $u_{\lambda\mu}^*$ entre ces catégories, et
des isomorphismes de transition
$u_{\mu\nu}^*\circ u_{\lambda\mu}^*\xrightarrow{\sim}u_{\lambda\nu}^*$. De façon
imagée, on peut dire que la donnée d'un $\mathcal O_S$-Module quasi-cohérent de
présentation finie $\mathcal F$ équivaut «~fonctoriellement~» à la donnée d'un
$\mathcal O_{S_\lambda}$-Module de présentation finie $\mathcal F'_\lambda$
pour $\lambda$ grand et que, si un $\mathcal O_{S_\mu}$-Module quasi-cohérent de
présentation finie $\mathcal F''_\mu$ a aussi $\mathcal F$ comme image
réciproque, alors $\mathcal F'_\lambda$ et $\mathcal F''_\mu$ ont même image
réciproque dans un $S_\nu$ convenable ($\nu\geq\lambda$, $\nu\geq\mu$).

Nous allons interpréter de ce point de vue diverses notions liées aux
$\mathcal O_S$-Modules quasi-cohérents.
\end{env}

\oldpage[IV-3]{23}
\begin{corollary}[8.5.4]
\label{IV.8.5.4-fr}
Supposons $S_0$ quasi-compact et quasi-séparé~; alors, pour tout
$\mathcal O_{S_\lambda}$-Module quasi-cohérent $\mathcal G_\lambda$,
l'homomorphisme canonique (8.5.1.4) est bijectif.
\end{corollary}

\begin{proof}
En effet, il suffit d'appliquer (8.5.2, (i)) au cas où
$\mathcal F_\lambda=\mathcal O_{S_\lambda}$, qui est de présentation finie.
\end{proof}

\begin{proposition}[8.5.5]
\label{IV.8.5.5-fr}
Supposons $S_0$ quasi-compact, et supposons que $\mathcal F_\lambda$ soit un
$\mathcal O_{S_\lambda}$-Module quasi-cohérent de présentation finie. Pour que
$\mathcal F$ soit localement libre (resp. localement libre de rang $n$), il
faut et il suffit qu'il existe $\mu\geq\lambda$ tel que $\mathcal F_\mu$ le
soit.
\end{proposition}

\begin{proof}
La condition étant trivialement suffisante, prouvons qu'elle est nécessaire.
Si $\mathcal F$ est localement libre (resp. localement libre de rang $n$), il
existe un recouvrement ouvert affine fini $(V_i)$ de $S$ tel que
$\mathcal F|V_i$ soit isomorphe à $\mathcal O_S^{n_i}|V_i$ (resp.
$\mathcal O_S^n|V_i$) pour tout $i$. En vertu de (8.2.11), il existe un
$\nu\geq\lambda$ et pour chaque $i$ un ouvert quasi-compact $U_{i\nu}$ de
$S_\nu$ tels que $V_i=u_\nu^{-1}(U_{i\nu})$. Comme $S_\nu$ est quasi-compact,
chaque $U_{i\nu}$ est réunion finie d'ouverts affines~; on est donc ramené au
cas où $S_0$ est affine et $\mathcal F=\mathcal O_S^n$~; on sait alors qu'il
existe $\mu\geq\lambda$ tel que $\mathcal F_\mu$ soit isomorphe à
$\mathcal O_{S_\mu}^n$ (8.5.2.5).
\end{proof}

\begin{proposition}[8.5.6]
\label{IV.8.5.6-fr}
Supposons $S_0$ quasi-compact, et considérons une suite
\[
  \mathcal F_\lambda\longrightarrow\mathcal G_\lambda
  \longrightarrow\mathcal H_\lambda\longrightarrow0
\]
d'homomorphismes de $\mathcal O_{S_\lambda}$-Modules quasi-cohérents, où
$\mathcal F_\lambda$ et $\mathcal G_\lambda$ sont de type fini et
$\mathcal H_\lambda$ de présentation finie. Pour que la suite correspondante
$\mathcal F\to\mathcal G\to\mathcal H\to0$ soit exacte, il faut et il suffit
qu'il existe $\mu\geq\lambda$ tel que la suite
$\mathcal F_\mu\to\mathcal G_\mu\to\mathcal H_\mu\to0$ le soit (auquel cas il
en est de même de la suite
$\mathcal F_\nu\to\mathcal G_\nu\to\mathcal H_\nu\to0$ pour
$\nu\geq\mu$).
\end{proposition}

\begin{proof}
Le fait que la condition soit suffisante et la dernière assertion résultent
de ce que le foncteur $u_\lambda^*$ (resp. $u_{\mu\nu}^*$) est exact à droite.
Pour démontrer que la condition est nécessaire, notons qu'il résulte de
l'hypothèse et de (8.5.2, (i)) qu'il existe $\nu\geq\lambda$ tel que le composé
$\mathcal F_\nu\to\mathcal G_\nu\to\mathcal H_\nu$ soit nul. Si l'on pose
$\mathcal H'_\nu=\operatorname{Coker}(\mathcal F_\nu\to\mathcal G_\nu)$, on a
donc un homomorphisme $f_\nu:\mathcal H'_\nu\to\mathcal H_\nu$~; par
hypothèse, $u_\nu^*(f_\nu)$ est un isomorphisme, et il résulte donc de
(8.5.2.4) qu'il existe $\mu\geq\nu$ tel que $u_{\nu\mu}^*(f_\nu)$ soit un
isomorphisme, ce qui achève la démonstration.
\end{proof}

\begin{corollary}[8.5.7]
\label{IV.8.5.7-fr}
Supposons $S_0$ quasi-compact, $\mathcal F_\lambda$ quasi-cohérent,
$\mathcal G_\lambda$ quasi-cohérent de type fini, et soit
$f_\lambda:\mathcal F_\lambda\to\mathcal G_\lambda$ un homomorphisme. Pour que
$f=u_\lambda^*(f_\lambda):\mathcal F\to\mathcal G$ soit surjectif, il faut et
il suffit qu'il existe $\mu\geq\lambda$ tel que
$f_\mu=u_{\lambda\mu}^*(f_\lambda):\mathcal F_\mu\to\mathcal G_\mu$ le soit.
\end{corollary}

\begin{proof}
C'est le cas particulier de (8.5.6) appliqué à la suite
$0\to\mathcal H_\lambda\to0\to0$, où
$\mathcal H_\lambda=\operatorname{Coker}(f_\lambda)$, qui est quasi-cohérent et
de type fini (compte tenu de ce que l'on a
$\mathcal H=\operatorname{Coker}(f)$ et
$\mathcal H_\mu=\operatorname{Coker}(f_\mu)$, en vertu de l'exactitude à droite
des foncteurs $u_\lambda^*$ et $u_{\lambda\mu}^*$).
\end{proof}

\begin{corollary}[8.5.8]
\EGAReaderTarget{ega4.c04.cor.8-5-8}{IV.8.5.8-fr}
Supposons $S_0$ quasi-compact et les morphismes
$u_{\lambda\mu}:S_\mu\to S_\lambda$ plats. Alors~:
\begin{enumerate}[label=(\roman*)]
  \item Soit
  $\mathcal F_\lambda\xrightarrow{f_\lambda}\mathcal G_\lambda
  \xrightarrow{g_\lambda}\mathcal H_\lambda$ une suite d'homomorphismes de
  $\mathcal O_{S_\lambda}$-Modules quasi-cohérents, tels que
  $\operatorname{Im}f_\lambda$ et $\operatorname{Ker}g_\lambda$ soient de type
  fini. Pour que la suite correspondante
  $\mathcal F\xrightarrow{f}\mathcal G\xrightarrow{g}\mathcal H$ soit

  \oldpage[IV-3]{24}
  exacte, il faut et il suffit qu'il existe $\mu\geq\lambda$ tel que la suite
  $\mathcal F_\mu\xrightarrow{f_\mu}\mathcal G_\mu
  \xrightarrow{g_\mu}\mathcal H_\mu$ soit exacte.

  \item Soit $f_\lambda:\mathcal F_\lambda\to\mathcal G_\lambda$ un
  homomorphisme de $\mathcal O_{S_\lambda}$-Modules quasi-cohérents tel que
  $\operatorname{Ker}f_\lambda$ soit de type fini. Pour que
  $f=u_\lambda^*(f_\lambda):\mathcal F\to\mathcal G$ soit injectif, il faut et
  il suffit qu'il existe $\mu\geq\lambda$ tel que
  $f_\mu=u_{\lambda\mu}^*(f_\lambda):\mathcal F_\mu\to\mathcal G_\mu$ le soit.
\end{enumerate}
\end{corollary}

\begin{proof}
\begin{enumerate}[label=(\roman*)]
  \item Compte tenu de (8.3.8, (ii)), notons que, par platitude,
  $\operatorname{Im}f$ et $\operatorname{Ker}g$ (resp.
  $\operatorname{Im}f_\mu$ et $\operatorname{Ker}g_\mu$ pour
  $\mu\geq\lambda$) sont les images réciproques de
  $\operatorname{Im}f_\lambda$ et $\operatorname{Ker}g_\lambda$
  $(0_{\mathrm I},6.7.2)$. Supposons que la suite
  $\mathcal F\xrightarrow{f}\mathcal G\xrightarrow{g}\mathcal H$ soit exacte.
  Puisque $\operatorname{Im}f_\lambda$ est de type fini, il existe
  $\mu\geq\lambda$ tel que le composé
  $\mathcal F_\mu\xrightarrow{f_\mu}\mathcal G_\mu
  \xrightarrow{g_\mu}\mathcal H_\mu$ soit nul, en vertu de (8.5.2, (i)).
  Changeant les notations, on peut donc déjà supposer que
  $g_\lambda\circ f_\lambda=0$. Alors comme l'homomorphisme
  $\mathcal F\to\operatorname{Ker}g$ est surjectif et que
  $\operatorname{Ker}g_\lambda$ est de type fini, il résulte de (8.5.7) qu'il
  existe $\mu\geq\lambda$ tel que l'homomorphisme
  $\mathcal F_\mu\to\operatorname{Ker}g_\mu$ soit surjectif, ce qui achève de
  prouver (i).

  \item L'assertion est le cas particulier de (i) appliqué à la suite
  $0\to\mathcal F_\lambda\to\mathcal G_\lambda$.
\end{enumerate}
\end{proof}

\begin{lemma}[8.5.9]
\label{IV.8.5.9-fr}
Supposons $S_0$ quasi-compact, $\mathcal F_\lambda$ quasi-cohérent de type
fini~; soient $\mathcal G'_\lambda$ et $\mathcal G''_\lambda$ deux quotients
quasi-cohérents de $\mathcal F_\lambda$, $\mathcal G'_\lambda$ étant en outre
supposé de présentation finie. Pour que $\mathcal G''$ soit quotient de
$\mathcal G'$, il faut et il suffit qu'il existe $\mu\geq\lambda$ tel que
$\mathcal G''_\mu$ soit quotient de $\mathcal G'_\mu$.
\end{lemma}

\begin{proof}
Par hypothèse, il y a deux homomorphismes surjectifs
$p'_\lambda:\mathcal F_\lambda\to\mathcal G'_\lambda$,
$p''_\lambda:\mathcal F_\lambda\to\mathcal G''_\lambda$~; en vertu de
l'exactitude à droite de $u_\lambda^*$ et de $u_{\lambda\mu}^*$,
$p'=u_\lambda^*(p'_\lambda)$, $p''=u_\lambda^*(p''_\lambda)$,
$p'_\mu=u_{\lambda\mu}^*(p'_\lambda)$,
$p''_\mu=u_{\lambda\mu}^*(p''_\lambda)$ sont aussi surjectifs~; en outre, s'il
existe un homomorphisme $f:\mathcal G'\to\mathcal G''$ (resp.
$f_\mu:\mathcal G'_\mu\to\mathcal G''_\mu$) tel que $p''=f\circ p'$ (resp.
$p''_\mu=f_\mu\circ p'_\mu$), cet homomorphisme est nécessairement unique, ce
qui montre que la question est locale sur $S_\lambda$, et qu'on peut par suite
($S_0$ étant quasi-compact et $L$ filtrant) supposer $S_\lambda$ affine, donc
quasi-séparé. Il est clair que la condition de l'énoncé est suffisante.
Inversement, comme $\mathcal G'_\lambda$ est de présentation finie,
$S_\lambda$ quasi-compact et quasi-séparé, il résulte de (8.5.2, (i)) que s'il
existe un homomorphisme $f:\mathcal G'\to\mathcal G''$ tel que
$p''=f\circ p'$, il existe un $\mu\geq\lambda$ et un homomorphisme
$f_\mu:\mathcal G'_\mu\to\mathcal G''_\mu$ tels que
$f=u_\mu^*(f_\mu)$ et $p''_\mu=f_\mu\circ p'_\mu$, d'où le lemme.
\end{proof}

\begin{env}[8.5.10]
\label{IV.8.5.10-fr}
Dans la fin de ce numéro, pour tout Module quasi-cohérent $\mathcal F$ sur un
préschéma, notons $\mathfrak Q(\mathcal F)$ l'ensemble des Modules quotients de
$\mathcal F$ qui sont de présentation finie. Si $\mathcal F_\lambda$ est
quasi-cohérent et $\mathcal G_\lambda\in\mathfrak Q(\mathcal F_\lambda)$, il
résulte du fait que $u_{\lambda\mu}^*$ et $u_\lambda^*$ sont exacts à droite,
que l'on a $\mathcal G_\mu\in\mathfrak Q(\mathcal F_\mu)$ pour
$\mu\geq\lambda$ et $\mathcal G\in\mathfrak Q(\mathcal F)$~; il est clair que
$(\mathfrak Q(\mathcal F_\lambda),u_{\lambda\mu}^*)$ est un système inductif
d'ensembles, et que les
$u_\lambda^*:\mathfrak Q(\mathcal F_\lambda)\to\mathfrak Q(\mathcal F)$
forment un système inductif d'applications, d'où, en passant à la limite
inductive, une application canonique
\begin{equation}
  u_{\mathfrak Q}:\varinjlim\mathfrak Q(\mathcal F_\lambda)
  \longrightarrow\mathfrak Q(\mathcal F).
  \tag{8.5.10.1}\label{IV.8.5.10.1-fr}
\end{equation}

En outre, si $(\mathcal F'_\lambda)$ est une seconde famille de
$\mathcal O_{S_\lambda}$-Modules quasi-cohérents et si, pour tout $\lambda$,
$\mathcal F'_\lambda$ est un quotient de $\mathcal F_\lambda$, alors
$\mathcal F'$ est un quotient de $\mathcal F$ et l'on a un diagramme
commutatif

\oldpage[IV-3]{25}
\[
\begin{tikzcd}
\varinjlim\mathfrak Q(\mathcal F_\lambda)
  \ar[r,"u_{\mathfrak Q}"] \ar[d]
  & \mathfrak Q(\mathcal F) \ar[d] \\
\varinjlim\mathfrak Q(\mathcal F'_\lambda)
  \ar[r,"u_{\mathfrak Q}"']
  & \mathfrak Q(\mathcal F')
\end{tikzcd}
\tag{8.5.10.2}\label{IV.8.5.10.2-fr}
\]
\end{env}

\begin{proposition}[8.5.11]
\label{IV.8.5.11-fr}
Supposons $S_0$ quasi-compact (resp. quasi-compact et quasi-séparé). Supposons
$\mathcal F_\lambda$ quasi-cohérent de type fini (resp. de présentation finie)
pour tout $\lambda$~; alors l'application canonique (8.5.10.1) est injective
(resp. bijective).
\end{proposition}

\begin{proof}
La première assertion résulte du lemme plus précis (8.5.9). Pour démontrer la
seconde, considérons un $\mathcal O_S$-Module quotient $\mathcal G$ de
$\mathcal F$ qui soit de présentation finie. Il résulte de (8.5.2, (ii)) qu'il
existe un $\lambda'\in L$ et un $\mathcal O_{S_{\lambda'}}$-Module
$\mathcal G_{\lambda'}$ quasi-cohérent de présentation finie tels que
$\mathcal G=u_{\lambda'}^*(\mathcal G_{\lambda'})$~; comme $L$ est filtrant, on
peut supposer $\lambda'=\lambda$ (en remplaçant $\lambda$ et $\lambda'$ par un
indice qui les majore). Considérons alors l'homomorphisme canonique
$p:\mathcal F\to\mathcal G$~; il résulte de (8.5.2, (i)) qu'il existe un
$\mu\geq\lambda$ et un homomorphisme
$p_\mu:\mathcal F_\mu\to\mathcal G_\mu$ tels que $p=u_\mu^*(p_\mu)$. En outre,
en vertu de (8.5.7), on peut supposer $\mu$ pris assez grand pour que $p_\mu$
soit surjectif, ce qui termine la démonstration.
\end{proof}

\subsection{Sous-préschémas de présentation finie d'une limite projective de préschémas}
\label{subsection:IV.8.6-fr}

\begin{env}[8.6.1]
\label{IV.8.6.1-fr}
Étant donné un préschéma $Y$, désignons dans ce numéro par
$\mathfrak{Spr}(Y)$ l'ensemble ordonné (I, 4.1.10) des sous-préschémas de $Y$
qui sont de présentation finie sur $Y$ (1.6.1), par $\mathfrak{Spro}(Y)$
(resp. $\mathfrak{Sprf}(Y)$) la partie de $\mathfrak{Spr}(Y)$ formée des
sous-préschémas induits sur des ouverts (resp. des sous-préschémas fermés) de
$Y$, de présentation finie sur $Y$~; il revient au même de dire qu'un
sous-préschéma $Z$ de $Y$ appartient à $\mathfrak{Spro}(Y)$ (resp.
$\mathfrak{Sprf}(Y)$) ou qu'il est induit sur un ouvert et que l'espace
sous-jacent $Z$ est rétrocompact dans $Y$ (resp. qu'il est fermé et que l'idéal
$\mathcal J$ de $\mathcal O_Y$ qui le définit est de type fini, ce qui signifie
aussi que $j_*(\mathcal O_Z)\in\mathfrak Q(\mathcal O_Y)$ si $j:Z\to Y$ est
l'injection canonique) (1.6.1 et 1.4.5).
\end{env}

\begin{env}[8.6.2]
\label{IV.8.6.2-fr}
Nous conservons toujours les notations de (8.2.2) et supposons que $S_0$ est
un des $S_\lambda$. Soit $Y_\lambda$ un sous-préschéma de $S_\lambda$~; alors
$Y_\mu=u_{\lambda\mu}^{-1}(Y_\lambda)$ (resp.
$Y=u_\lambda^{-1}(Y_\lambda)$) est un sous-préschéma de $S_\mu$ pour
$\mu\geq\lambda$ (resp. de $S$)~; il est induit sur un ouvert (resp. il est
fermé) si $Y_\lambda$ l'est (I, 4.3.2) et de présentation finie sur $S_\mu$
(resp. $S$) si $Y_\lambda$ est de présentation finie sur $S_\lambda$
(1.6.2, (iii)). Par suite
$(\mathfrak{Spr}(S_\lambda),u_{\lambda\mu}^{-1})$ (resp.
$(\mathfrak{Spro}(S_\lambda),u_{\lambda\mu}^{-1})$,
$(\mathfrak{Sprf}(S_\lambda),u_{\lambda\mu}^{-1})$) est un système inductif et
les applications
$u_\lambda^{-1}:\mathfrak{Spr}(S_\lambda)\to\mathfrak{Spr}(S)$ (resp.
$\mathfrak{Spro}(S_\lambda)\to\mathfrak{Spro}(S)$,
$\mathfrak{Sprf}(S_\lambda)\to\mathfrak{Sprf}(S)$) forment un système inductif
d'applications~; d'où, en passant à la limite inductive, des applications
canoniques
\begin{align}
  \varinjlim\mathfrak{Spr}(S_\lambda)
    &\longrightarrow\mathfrak{Spr}(S)
    \tag{8.6.2.1}\label{IV.8.6.2.1-fr}\\
  \varinjlim\mathfrak{Spro}(S_\lambda)
    &\longrightarrow\mathfrak{Spro}(S)
    \tag{8.6.2.2}\label{IV.8.6.2.2-fr}\\
  \varinjlim\mathfrak{Sprf}(S_\lambda)
    &\longrightarrow\mathfrak{Sprf}(S).
    \tag{8.6.2.3}\label{IV.8.6.2.3-fr}
\end{align}
\end{env}

\oldpage[IV-3]{26}

Rappelons (I, 4.1.10) que $\mathfrak{Spr}(X)$, pour tout préschéma $X$, est un
ensemble \emph{ordonné} par la relation «~$Z$ est un sous-préschéma du
sous-préschéma $Y$~» qui s'écrit $Z\leq Y$. Les applications
$u_{\lambda\mu}^{-1}$ et $u_\lambda^{-1}$ sont \emph{croissantes} pour les
relations d'ordre correspondantes dans $\mathfrak{Spr}(S_\lambda)$,
$\mathfrak{Spr}(S_\mu)$, $\mathfrak{Spr}(S)$. En outre, on définit une relation
d'ordre dans l'ensemble $\varinjlim\mathfrak{Spr}(S_\lambda)$ en écrivant que
$\eta\leq\eta'$ pour deux éléments de cet ensemble lorsqu'il existe un
$\lambda$ et deux éléments $Y_\lambda$, $Y'_\lambda$ de
$\mathfrak{Spr}(S_\lambda)$, dont $\eta$ et $\eta'$ sont les images canoniques,
et qui sont tels que $Y_\lambda\leq Y'_\lambda$~; on vérifie aisément que cela
ne dépend pas des représentants $Y_\lambda$, $Y'_\lambda$ considérés, et que
l'on a bien ainsi une \emph{relation d'ordre}. Cela étant, le fait que les
$u_\lambda^{-1}$ sont croissantes entraîne aussitôt que l'application
canonique (8.6.2.1) est \emph{croissante}~; il en est de même évidemment de
(8.6.2.2) et (8.6.2.3).

\begin{proposition}[8.6.3]
\label{IV.8.6.3-fr}
Supposons $S_0$ quasi-compact (resp. quasi-compact et quasi-séparé). Alors les
applications (8.6.2.1), (8.6.2.2), (8.6.2.3) sont injectives (resp.
bijectives).
\end{proposition}

\begin{proof}
Compte tenu des remarques de (8.6.1), les assertions relatives à (8.6.2.3)
découlent de (8.5.11) appliqué à $\mathcal F_\lambda=\mathcal O_{S_\lambda}$~;
de même, les assertions relatives à (8.6.2.2) sont des cas particuliers de
(8.3.5) et (8.3.11), compte tenu de ce que les $S_\lambda$ et $S$ sont
quasi-compacts. Reste à considérer l'application (8.6.2.1). Prouvons d'abord
qu'elle est surjective lorsque $S_0$ est quasi-compact et quasi-séparé. Soit
$Z$ un sous-préschéma de $S$, de présentation finie sur $S$~; comme $S$ est
quasi-compact, il en est de même de $Z$, donc il existe un ouvert
quasi-compact $U$ de $S$ tel que $Z$ soit un sous-préschéma fermé de $U$, de
présentation finie sur $U$. Il existe alors un indice $\lambda$ et un ouvert
quasi-compact $U_\lambda$ de $S_\lambda$ tels que
$U=u_\lambda^{-1}(U_\lambda)$ (8.2.11)~; comme $S_\lambda$ est quasi-séparé, il
en est de même de $U_\lambda$ (I, 1.2.7), et par suite on peut se borner au cas
où $U=S$~; mais dans ce cas, on est ramené au fait que (8.6.2.3) est
surjective.

Enfin, pour voir que (8.6.2.1) est injective lorsque $S_0$ est quasi-compact,
il suffira de prouver le résultat plus précis suivant~:
\end{proof}

\begin{corollary}[8.6.3.1]
\label{IV.8.6.3.1-fr}
Supposons $S_0$ quasi-compact et soient $Z'_\lambda$, $Z''_\lambda$ deux
sous-préschémas de $S_\lambda$, de présentation finie sur $S_\lambda$. Pour
que $Z'=u_\lambda^{-1}(Z'_\lambda)$ soit majoré par
$Z''=u_\lambda^{-1}(Z''_\lambda)$ (I, 4.1.10), il faut et il suffit qu'il
existe $\mu\geq\lambda$ tel que
$Z'_\mu=u_{\lambda\mu}^{-1}(Z'_\lambda)$ soit majoré par
$Z''_\mu=u_{\lambda\mu}^{-1}(Z''_\lambda)$.
\end{corollary}

\begin{proof}
Il est trivial que la condition est suffisante. Pour voir qu'elle est
nécessaire, notons d'abord que les ensembles sous-jacents $Z'_\lambda$ et
$Z''_\lambda$ sont localement constructibles dans $S_\lambda$ par hypothèse
(I, 1.8.4), donc l'hypothèse entraîne l'existence d'un $\nu\geq\lambda$ tel
que $Z'_\nu\subset Z''_\nu$ (8.3.5)~; remplaçant $\lambda$ par $\nu$, on peut
donc déjà supposer que l'on a $Z'_\lambda\subset Z''_\lambda$. En outre, par
hypothèse (I, 1.6.1), les sous-espaces $Z'_\lambda$ et $Z''_\lambda$ de
$S_\lambda$ sont quasi-compacts~; pour tout point $x\in Z'_\lambda$, il y a un
voisinage ouvert quasi-compact $V(x)$ dans $S_\lambda$ tels que
$V(x)\cap Z'_\lambda$ et $V(x)\cap Z''_\lambda$ soient fermés dans $V(x)$. En
recouvrant $S_\lambda$ par un nombre fini de voisinages $V(x_i)$ on voit donc
qu'il y a un ouvert quasi-compact $U_\lambda$ de $S_\lambda$ contenant
$Z'_\lambda$ et tel que $Z'_\lambda$ et $Z''_\lambda\cap U_\lambda$ soient
fermés dans $U_\lambda$. Si l'on désigne par $Y_\lambda$ le sous-préschéma
induit par $Z''_\lambda$ sur $U_\lambda\cap Z''_\lambda$, il est clair qu'avec
les notations habituelles, $Y_\mu$ (resp. $Y$) est induit par $Z''_\mu$ sur
$U_\mu\cap Z''_\mu$ pour $\mu\geq\lambda$ (resp. par $Z''$ sur
$U\cap Z''$)~; en outre $Z'$ est majoré par $Y$ (I, 4.4.1), et comme il suffit
de prouver que $Z'_\mu$ est majoré par $Y_\mu$ pour $\mu$
\oldpage[IV-3]{27}
assez grand, on voit finalement qu'on est ramené (en remplaçant $S_\lambda$
par $U_\lambda$) au cas où $Z'_\lambda$ et $Z''_\lambda$ sont des
sous-préschémas fermés de $S_\lambda$. Mais alors, cela a déjà été démontré
puisque (8.6.2.3) est croissante et injective.
\end{proof}

\begin{corollary}[8.6.4]
\label{IV.8.6.4-fr}
Supposons $S_0$ quasi-compact, et soit $Z_\lambda$ un sous-préschéma de
$S_\lambda$, de présentation finie sur $S_\lambda$. Pour que
$Z=u_\lambda^{-1}(Z_\lambda)$ soit un sous-préschéma induit sur un ouvert
(resp. un sous-préschéma fermé) de $S$, il faut et il suffit qu'il existe
$\mu\geq\lambda$ tel que $Z_\mu=u_{\lambda\mu}^{-1}(Z_\lambda)$ soit induit
sur un ouvert (resp. un sous-préschéma fermé) dans $S_\mu$.
\end{corollary}

\begin{proof}
Soit $(U_\lambda^{(i)})$ un recouvrement ouvert affine fini de $S_\lambda$, et
posons $U_\mu^{(i)}=u_{\lambda\mu}^{-1}(U_\lambda^{(i)})$ pour
$\mu\geq\lambda$ et $U^{(i)}=u_\lambda^{-1}(U_\lambda^{(i)})$. Si $Z$ est
ouvert (resp. fermé) dans $S$, $Z\cap U^{(i)}$ l'est dans $U^{(i)}$, et
inversement si chacun des $Z_\mu\cap U_\mu^{(i)}$ est ouvert (resp. fermé)
dans $U_\mu^{(i)}$, $Z_\mu$ l'est dans $S_\mu$. Puisque $L$ est filtrant, il
suffit de démontrer le corollaire lorsque $S_\lambda$ est affine, donc
quasi-séparé. Mais alors le résultat découle de ce que les applications
(8.6.2.1), (8.6.2.2) et (8.6.2.3) sont bijectives.
\end{proof}

\subsection{Critères pour qu'une limite projective de préschémas soit un préschéma réduit (resp. intègre)}
\label{subsection:IV.8.7-fr}

Nous conservons les hypothèses et notations de (8.2.2) et supposons toujours
que $S_0$ est un des $S_\lambda$.

\begin{proposition}[8.7.1]
\label{IV.8.7.1-fr}
Supposons que $S$ soit non réduit. Alors il existe $\lambda_0$ tel que pour
$\lambda\geq\lambda_0$, $S_\lambda$ soit non réduit.
\end{proposition}

\begin{proof}
La question étant locale sur $S_0$, on peut supposer
$S_0=\operatorname{Spec}(A_0)$ affine, d'où $S=\operatorname{Spec}(A)$, où
$A=\varinjlim A_\lambda$ est la limite inductive d'un système inductif de
$A_0$-algèbres $(A_\lambda)$. On sait alors (5.13.2) que le nilradical de $A$
est la limite inductive de ceux des $A_\lambda$~; s'il n'est pas nul, un des
$A_\lambda$ contient donc un élément nilpotent $a_\lambda\neq0$ dont l'image
dans $A$ est un élément nilpotent et $\neq0$, et l'image de $a_\lambda$ dans
les $A_\mu$ pour $\mu\geq\lambda$ est par suite un élément nilpotent et
$\neq0$.
\end{proof}

\begin{proposition}[8.7.2]
\label{IV.8.7.2-fr}
Supposons vérifiée l'une des hypothèses suivantes~:
\begin{enumerate}
\item[\emph{a)}] $S_0$ est quasi-compact, le nilradical $\mathcal N_0$ de
$\mathcal O_{S_0}$ est un Idéal de type fini (ce qui sera par exemple vérifié
lorsque $S_0$ est noethérien), et les morphismes
$u_{\lambda\mu}:S_\mu\to S_\lambda$ sont des immersions ouvertes.
\item[\emph{b)}] Les morphismes $u_{\lambda\mu}:S_\mu\to S_\lambda$ sont
fidèlement plats.
\end{enumerate}
Sous ces conditions, pour que $S$ soit réduit, il faut et il suffit qu'il
existe $\lambda_0$ tel que $S_\lambda$ soit réduit pour
$\lambda\geq\lambda_0$.

En outre, dans le cas \emph{b)}, si $S$ est réduit, tous les $S_\lambda$ le
sont.
\end{proposition}

\begin{proof}
La dernière assertion résulte de ce que le morphisme $S\to S_\lambda$ est
alors fidèlement plat pour tout $\lambda$ (8.3.8) et de (2.1.13). D'autre part,
(8.7.1) prouve que la condition de l'énoncé est suffisante (sans hypothèse sur
$S_0$ ni sur les $u_{\lambda\mu}$). Il reste donc à voir que la condition est
nécessaire dans l'hypothèse \emph{a)}~; alors (8.2.13), l'espace sous-jacent à
$S$ s'identifie à l'intersection des espaces sous-jacents aux $S_\lambda$ (les
$S_\lambda$ étant identifiés à des sous-préschémas induits sur des ouverts de
$S_0$), et le faisceau structural $\mathcal O_S$ s'identifie au
\oldpage[IV-3]{28}
faisceau induit sur $S$ par tous les $\mathcal O_{S_\lambda}$~; en particulier
pour tout $s\in S$, l'anneau local $\mathcal O_s$ est \emph{le même} pour $S$
et pour tous les $S_\lambda$. Si $\mathcal N_\lambda$ est le Nilradical de
$\mathcal O_{S_\lambda}$, le Nilradical $\mathcal N$ de $\mathcal O_S$ a donc
en chaque point de $S$ la même fibre $\mathcal N_s$ (nilradical de
$\mathcal O_s$) que $u_\lambda^*(\mathcal N_\lambda)$ (induit sur $S$ par
$\mathcal N_\lambda$). L'hypothèse que $S$ est réduit entraîne par suite
$u_\lambda^*(\mathcal N_\lambda)=0$~; comme $\mathcal N_0$ est supposé de
type fini, il en est de même des $\mathcal N_\lambda$, et la conclusion résulte
donc de (8.5.7).
\end{proof}

\begin{corollary}[8.7.3]
\label{IV.8.7.3-fr}
Supposons vérifiée l'une des hypothèses suivantes~:
\begin{enumerate}
\item[\emph{a)}] $S_0$ est un préschéma noethérien et les morphismes
$u_{\lambda\mu}:S_\mu\to S_\lambda$ sont des immersions ouvertes.
\item[\emph{b)}] Les morphismes $u_{\lambda\mu}:S_\mu\to S_\lambda$ sont
fidèlement plats.
\end{enumerate}
Alors, pour que $S$ soit intègre, il faut et il suffit qu'il existe
$\lambda_0$ tel que $S_\lambda$ soit intègre pour $\lambda\geq\lambda_0$.
\end{corollary}

\begin{proof}
Dire qu'un préschéma est intègre signifie qu'il est à la fois réduit et
irréductible~; le corollaire résulte donc de (8.7.2) et de (8.4.3).
\end{proof}

\begin{remark}[8.7.4]
\label{IV.8.7.4-fr}
Si l'on ne fait pas d'hypothèse sur les $u_{\lambda\mu}$, il peut se faire que
$S$ soit intègre bien que tous les $S_\lambda$ soient non réduits et non
connexes, comme le montre l'exemple (8.4.4), où on prend l'anneau $A$ non
réduit.
\end{remark}

\subsection{Préschémas de présentation finie sur une limite projective de préschémas}
\label{subsection:IV.8.8-fr}

\begin{env}[8.8.1]
\label{IV.8.8.1-fr}
Conservant toujours les notations et hypothèses de (8.2.2), nous supposerons
donnés dans cette section deux $S_\alpha$-préschémas $X_\alpha$, $Y_\alpha$,
ce qui définit (8.2.5) deux systèmes projectifs de préschémas
$(X_\lambda,v_{\lambda\mu})$ et $(Y_\lambda,w_{\lambda\mu})$ en posant
$X_\lambda=X_\alpha\times_{S_\alpha}S_\lambda$,
$Y_\lambda=Y_\alpha\times_{S_\alpha}S_\lambda$,
$v_{\lambda\mu}=1_{X_\alpha}\times u_{\lambda\mu}$,
$w_{\lambda\mu}=1_{Y_\alpha}\times u_{\lambda\mu}$ (pour
$\alpha\leq\lambda\leq\mu$), dont les limites projectives sont respectivement
$X=X_\alpha\times_{S_\alpha}S$, $Y=Y_\alpha\times_{S_\alpha}S$, les
morphismes canoniques $v_\lambda:X\to X_\lambda$ et
$w_\lambda:Y\to Y_\lambda$ étant respectivement égaux à
$1_{X_\alpha}\times u_\lambda$ et $1_{Y_\alpha}\times u_\lambda$. Pour
$\alpha\leq\lambda\leq\mu$, on a une application canonique
$e_{\mu\lambda}:\operatorname{Hom}_{S_\lambda}(X_\lambda,Y_\lambda)\to
\operatorname{Hom}_{S_\mu}(X_\mu,Y_\mu)$, qui à tout $S_\lambda$-morphisme
$f_\lambda:X_\lambda\to Y_\lambda$ fait correspondre
$f_\mu=f_\lambda\times1_{S_\mu}:X_\lambda\times_{S_\lambda}S_\mu\to
Y_\lambda\times_{S_\lambda}S_\mu$, et il est clair que
$(\operatorname{Hom}_{S_\lambda}(X_\lambda,Y_\lambda),e_{\mu\lambda})$ est un
\emph{système inductif} d'ensembles. De même, on a une application canonique
$e_\lambda:\operatorname{Hom}_{S_\lambda}(X_\lambda,Y_\lambda)\to
\operatorname{Hom}_S(X,Y)$ qui à $f_\lambda$ fait correspondre
$f=f_\lambda\times1_S:X_\lambda\times_{S_\lambda}S\to
Y_\lambda\times_{S_\lambda}S$ et $(e_\lambda)$ est un système inductif
d'applications~; d'où, en passant à la limite inductive, une application
canonique, fonctorielle en $S_\alpha$, $X_\alpha$ et $Y_\alpha$~:
\begin{equation}
e:\varinjlim\operatorname{Hom}_{S_\lambda}(X_\lambda,Y_\lambda)
\longrightarrow\operatorname{Hom}_S(X,Y).
\tag{8.8.1.1}\label{IV.8.8.1.1-fr}
\end{equation}
\end{env}

\begin{theorem}[8.8.2]
\EGAReaderTarget{ega4.c04.thm.8-8-2}{IV.8.8.2-fr}
\begin{enumerate}
\item[(i)] Supposons $X_\alpha$ quasi-compact (resp. quasi-compact et
quasi-séparé), et $Y_\alpha$ localement de type fini (resp. localement de
présentation finie) sur $S_\alpha$. Alors l'application (8.8.1.1) est
injective (resp. bijective).
\item[(ii)] Supposons $S_0$ quasi-compact et quasi-séparé. Pour tout préschéma
$X$ de présentation finie sur $S$, il existe un $\lambda\in L$, un préschéma
$X_\lambda$ de présentation finie sur $S_\lambda$, et un $S$-isomorphisme
$X\xrightarrow{\sim}X_\lambda\times_{S_\lambda}S$.
\end{enumerate}
\end{theorem}

\begin{proof}
\oldpage[IV-3]{29}
(i) Considérons d'abord le cas où $S_0=\operatorname{Spec}(A_0)$,
$X_\alpha=\operatorname{Spec}(B_\alpha)$ et
$Y_\alpha=\operatorname{Spec}(C_\alpha)$ sont affines~; alors les
$S_\lambda=\operatorname{Spec}(A_\lambda)$ et $S=\operatorname{Spec}(A)$ sont
aussi affines, avec $A=\varinjlim A_\lambda$, et les assertions de (i) sont
équivalentes au
\end{proof}

\begin{lemma}[8.8.2.1]
\label{IV.8.8.2.1-fr}
Soient $A_0$ un anneau, $(A_\lambda)_{\lambda\in L}$ un système inductif de
$A_0$-algèbres, $A=\varinjlim A_\lambda$~; soient $B_\alpha$ une
$A_\alpha$-algèbre, $C_\alpha$ une $A_\alpha$-algèbre de type fini (resp. de
présentation finie). Alors l'homomorphisme canonique
\begin{equation}
\varinjlim_\lambda
\operatorname{Hom}_{A_\lambda\text{-alg.}}
  (C_\alpha\otimes_{A_\alpha}A_\lambda,
   B_\alpha\otimes_{A_\alpha}A_\lambda)
\longrightarrow
\operatorname{Hom}_{A\text{-alg.}}
  (C_\alpha\otimes_{A_\alpha}A,
   B_\alpha\otimes_{A_\alpha}A)
\tag{8.8.2.2}\label{IV.8.8.2.2-fr}
\end{equation}
est injectif (resp. bijectif).
\end{lemma}

\begin{proof}
On sait que l'on a des isomorphismes canoniques fonctoriels
\[
\operatorname{Hom}_{A_\lambda\text{-alg.}}
 (C_\alpha\otimes_{A_\alpha}A_\lambda,
  B_\alpha\otimes_{A_\alpha}A_\lambda)
\xrightarrow{\sim}
\operatorname{Hom}_{A_\alpha\text{-alg.}}
 (C_\alpha,B_\alpha\otimes_{A_\alpha}A_\lambda)
\]
\[
\operatorname{Hom}_{A\text{-alg.}}
 (C_\alpha\otimes_{A_\alpha}A,
  B_\alpha\otimes_{A_\alpha}A)
\xrightarrow{\sim}
\operatorname{Hom}_{A_\alpha\text{-alg.}}
 (C_\alpha,B_\alpha\otimes_{A_\alpha}A)
\]
en vertu de la propriété universelle du produit tensoriel de deux algèbres. Il
suffit donc de prouver le
\end{proof}

\begin{lemma}[8.8.2.3]
\label{IV.8.8.2.3-fr}
Soient $E$ un anneau, $G$ une $E$-algèbre de type fini (resp. de présentation
finie), $(F_\lambda)$ un système inductif de $E$-algèbres. Alors
l'homomorphisme canonique
\[
\varinjlim\operatorname{Hom}_{E\text{-alg.}}(G,F_\lambda)
\longrightarrow
\operatorname{Hom}_{E\text{-alg.}}(G,\varinjlim F_\lambda)
\]
qui, à tout système inductif d'homomorphismes
$\theta_\lambda:G\to F_\lambda$ de $E$-algèbres, fait correspondre sa limite
inductive, est injectif (resp. bijectif).
\end{lemma}

\begin{proof}
Supposons d'abord la $E$-algèbre $G$ de type fini, et soit
$(t_i)_{1\leq i\leq n}$ un système de générateurs de cette $E$-algèbre~;
montrons que si $(\theta'_\lambda)$, $(\theta''_\lambda)$ sont deux systèmes
inductifs d'homomorphismes $G\to F_\lambda$ tels que
$\varinjlim\theta'_\lambda=\varinjlim\theta''_\lambda$, il existe un
$\mu\geq\lambda$ tel que $\theta'_\mu=\theta''_\mu$. En effet, si
$\varphi_{\mu\lambda}:F_\lambda\to F_\mu$ et
$\varphi_\lambda:F_\lambda\to F=\varinjlim F_\lambda$ sont les
homomorphismes canoniques du système inductif $(F_\lambda)$, par hypothèse,
pour chaque indice $i$, il existe un indice $\lambda_i$ tel que
$\varphi_{\lambda_i}(\theta'_{\lambda_i}(t_i))=
\varphi_{\lambda_i}(\theta''_{\lambda_i}(t_i))$, et l'on peut supposer tous
les $\lambda_i$ égaux à un même $\lambda$~; il en résulte de la même manière
l'existence d'un $\mu\geq\lambda$ tel que
$\varphi_{\mu\lambda}(\theta'_\lambda(t_i))=
\varphi_{\mu\lambda}(\theta''_\lambda(t_i))$ pour $1\leq i\leq n$,
c'est-à-dire $\theta'_\mu(t_i)=\theta''_\mu(t_i)$ pour $1\leq i\leq n$,
d'où $\theta'_\mu=\theta''_\mu$.

Supposons en second lieu $G$ de présentation finie, de sorte que l'on a
\[
G=E[T_1,\ldots,T_n]/\mathfrak J,
\]
où $\mathfrak J$ est un idéal de type fini, $t_i$ étant la classe de $T_i$
mod. $\mathfrak J$. Soit $(P_j)_{1\leq j\leq m}$ un système de générateurs de
$\mathfrak J$. Supposons donné un homomorphisme de $E$-algèbres
$\theta:G\to F$~; posons $b_i=\theta(t_i)$~; par définition, on a donc
$P_j(b_1,b_2,\ldots,b_n)=\theta(P_j(t_1,\ldots,t_n))=0$ pour
$1\leq j\leq m$. Or, il existe un $\lambda$ et des éléments
$x_1,\ldots,x_n$ dans $F_\lambda$ tels que $b_i=\varphi_\lambda(x_i)$ pour
$1\leq i\leq n$~; on a donc
$\varphi_\lambda(P_j(x_1,\ldots,x_n))=P_j(b_1,\ldots,b_n)=0$, et par suite il
existe $\mu\geq\lambda$ tel que
$\varphi_{\mu\lambda}(P_j(x_1,\ldots,x_n))=
P_j(\varphi_{\mu\lambda}(x_1),\ldots,\varphi_{\mu\lambda}(x_n))=0$ pour
$1\leq j\leq m$~; on en conclut qu'il existe un homomorphisme de
$E$-algèbres $\theta_\mu:G\to F_\mu$ tel que
$\theta_\mu(t_i)=\varphi_{\mu\lambda}(x_i)$
\oldpage[IV-3]{30}
pour $1\leq i\leq n$~; on en déduit pour tout $\nu\geq\mu$ un
homomorphisme de $E$-algèbres
$\theta_\nu=\varphi_{\nu\mu}\circ\theta_\mu$ de $G$ dans $F_\nu$, et il est
clair que $\theta$ est la limite inductive de ce système inductif
d'homomorphismes, ce qui achève de démontrer le lemme.

Passons maintenant au cas où $X_\alpha$ est quasi-compact et $Y_\alpha$
localement de type fini sur $S_\alpha$. Posons
$Z_\alpha=X_\alpha\times_{S_\alpha}Y_\alpha$ et introduisons le système
projectif correspondant des
$Z_\lambda=Z_\alpha\times_{S_\alpha}S_\lambda=
X_\lambda\times_{S_\lambda}Y_\lambda$ et sa limite
$Z=Z_\alpha\times_{S_\alpha}S=X\times_S Y$~; les bijections canoniques
(I, 3.3.14) donnent des diagrammes commutatifs
\[
\begin{tikzcd}[column sep=large,row sep=large]
\operatorname{Hom}_{S_\lambda}(X_\lambda,Y_\lambda)
  \ar[r,"e_{\mu\lambda}"] \ar[d,"\wr"] &
\operatorname{Hom}_{S_\mu}(X_\mu,Y_\mu)
  \ar[r,"e_\mu"] \ar[d,"\wr"] &
\operatorname{Hom}_S(X,Y) \ar[d,"\wr"] \\
\operatorname{Hom}_{X_\lambda}(X_\lambda,Z_\lambda)
  \ar[r,"e_{\mu\lambda}"'] &
\operatorname{Hom}_{X_\mu}(X_\mu,Z_\mu)
  \ar[r,"e_\mu"'] &
\operatorname{Hom}_X(X,Z)
\end{tikzcd}
\]
et par suite on est ramené à prouver que (8.8.1.1) est injective dans le cas
particulier où $S_\alpha=X_\alpha$ (compte tenu de (I, 3.4)). De plus, comme
$X_\alpha$ est quasi-compact, donc réunion finie d'ouverts affines, on peut
supposer $X_\alpha$ affine ($L$ étant filtrant). Supposons donc donnés deux
$X_\alpha$-morphismes $f'_\alpha:X_\alpha\to Y_\alpha$,
$f''_\alpha:X_\alpha\to Y_\alpha$ tels que $f'$ et $f''$ soient des
$X$-morphismes égaux de $X$ dans $Y$~; il s'agit de prouver que
$f'_\mu=f''_\mu$ pour $\mu\geq\alpha$ assez grand. Comme $X_\alpha$ est
quasi-compact, il en est de même de
$f'_\alpha(X_\alpha)\cup f''_\alpha(X_\alpha)$, et comme $Y_\alpha$ est
localement de type fini sur $X_\alpha$,
$f'_\alpha(X_\alpha)\cup f''_\alpha(X_\alpha)$ est contenu dans une réunion
finie d'ouverts affines $V_{i\alpha}$ de $Y_\alpha$, de type fini sur
$X_\alpha$. Posons
$U'_{i\alpha}=f_\alpha^{'-1}(V_{i\alpha})$,
$U''_{i\alpha}=f_\alpha^{''-1}(V_{i\alpha})$,
$U_{i\alpha}=U'_{i\alpha}\cap U''_{i\alpha}$,
$U_\alpha=\bigcup_iU_{i\alpha}$. L'hypothèse $f'=f''$ entraîne
$v_\alpha^{-1}(U'_{i\alpha})=v_\alpha^{-1}(U''_{i\alpha})$, ces deux ensembles
étant respectivement égaux à $f^{'-1}(w_\alpha^{-1}(V_{i\alpha}))$ et à
$f^{''-1}(w_\alpha^{-1}(V_{i\alpha}))$. Comme les $V_{i\alpha}$ recouvrent
$Y_\alpha$, on a $v_\alpha^{-1}(U_\alpha)=f^{'-1}(Y)=X$. Comme $X_\alpha$
est quasi-compact et que toute partie ouverte de $X_\lambda$ est
ind-constructible, il résulte de (8.3.4) qu'il y a un indice
$\lambda\geq\alpha$ tel que les
$U_{i\mu}=v_{\alpha\mu}^{-1}(U_{i\alpha})$ forment un \emph{recouvrement} de
$X_\mu$. Remplaçant $\alpha$ par $\mu$, on peut donc supposer que les
$U_{i\alpha}$ recouvrent $X_\alpha$~; cela entraîne que pour tout
$x\in X_\alpha$, il y a un voisinage ouvert affine $W(x)$ contenu dans un des
$U_{i\alpha}$, autrement dit tel que les restrictions de $f'_\alpha$ et
$f''_\alpha$ à $W(x)$ appliquent $W(x)$ dans un \emph{même} ouvert affine
$V_{i\alpha}$. Comme $X_\alpha$ est quasi-compact, il est recouvert par un
nombre fini d'ouverts affines $W(x_k)$~; en vertu du lemme (8.8.2.1) et du fait
que $L$ est filtrant, il existe un $\lambda\geq\alpha$ tel que les
restrictions de $f'_\lambda$ et $f''_\lambda$ à chacun des ouverts
$v_{\alpha\lambda}^{-1}(W(x_k))$ soient égales, d'où
$f'_\lambda=f''_\lambda$.

Supposons maintenant $X_\alpha$ quasi-compact et quasi-séparé et $Y_\alpha$
localement de présentation finie sur $S_\alpha$, et prouvons que (8.8.1.1) est
surjective. Supposons donc donné un $X$-morphisme $f:X\to Y$. Comme $X$ est
quasi-compact, il en est de même
\oldpage[IV-3]{31}
de $f(X)$, et par suite il y a un ouvert quasi-compact $Y'$ dans $Y$ qui
contient $f(X)$~; il existe par suite un $\lambda\geq\alpha$ et un ouvert
quasi-compact $Y'_\lambda$ dans $Y_\lambda$ tels que
$Y'=w_\lambda^{-1}(Y'_\lambda)$ (8.2.11). Remplaçant au besoin $\alpha$ par
$\lambda$ et $Y_\lambda$ par $Y'_\lambda$, on peut donc se borner au cas où
$Y_\alpha$ est quasi-compact, donc aussi $Y$ et les $Y_\lambda$. Comme $Y$ est
quasi-compact, il est réunion finie d'ouverts affines $V_i$, et par suite $X$
est réunion des ouverts $f^{-1}(V_i)$. Comme tout point de $X$ a un voisinage
ouvert quasi-compact contenu dans un des $f^{-1}(V_i)$ et que $X$ est
quasi-compact, on peut, en tenant compte de (8.2.11) et remplaçant au besoin
$\alpha$ par un indice $\lambda\geq\alpha$, supposer que $Y$ est réunion finie
d'ouverts $V_i=w_\alpha^{-1}(V_{i\alpha})$, où les $V_{i\alpha}$ sont des
ouverts affines de $Y_\alpha$~; par suite $X$ est réunion des ouverts
$f^{-1}(V_i)$. Comme tout point de $X$ a un voisinage ouvert quasi-compact
contenu dans l'un des $f^{-1}(V_i)$ et que $X$ est quasi-compact, on peut
recouvrir $X$ par un nombre fini de tels voisinages, et (en répétant au besoin
certains des $V_i$) supposer que l'on a un recouvrement $(U_i)$ de $X$ par des
ouverts quasi-compacts ayant même ensemble d'indices que $(V_i)$ et tel que
$f(U_i)\subset V_i$ pour tout $i$. En outre, à l'aide de (8.2.11) (et en
remplaçant au besoin $\alpha$ par un indice $\lambda\geq\alpha$), on peut
supposer que l'on a $U_i=v_\alpha^{-1}(U_{i\alpha})$ où les $U_{i\alpha}$ sont
des ouverts quasi-compacts dans $X_\alpha$~; de plus, utilisant (8.3.4) comme
ci-dessus, on peut supposer que $(U_{i\alpha})$ est un \emph{recouvrement} de
$X_\alpha$.

Cela étant, il suffira de montrer qu'il existe un $\lambda\geq\alpha$ et pour
chaque $i$ un morphisme $f_{i\lambda}:U_{i\lambda}\to V_{i\lambda}$ (avec
$U_{i\lambda}=v_{\alpha\lambda}^{-1}(U_{i\alpha})$ et
$V_{i\lambda}=w_{\alpha\lambda}^{-1}(V_{i\alpha})$) tels que le morphisme
correspondant $f_i=e_\lambda(f_{i\lambda}):U_i\to V_i$ soit égal à la
\emph{restriction} de $f$ à $U_i$. En effet, s'il en est ainsi, comme
$X_\lambda$ est quasi-séparé (I, 1.2.3), les
$U_{i\lambda}\cap U_{j\lambda}$ sont \emph{quasi-compacts} et le résultat
d'unicité démontré plus haut (qui s'applique puisque $Y_\lambda$ est
localement de type fini sur $S_\lambda$) prouve qu'il existe
$\mu\geq\lambda$ tel que $f_{i\mu}$ et $f_{j\mu}$ coïncident dans
$U_{i\mu}\cap U_{j\mu}$ pour tout couple $(i,j)$, et par suite définissent un
$S_\mu$-morphisme $f_\mu:X_\mu\to Y_\mu$ tel que $e_\mu(f_\mu)=f$.

On est ainsi ramené au cas où $Y_\alpha$ est \emph{affine}, et comme en outre
on peut supposer que les $V_{i\alpha}$ ont des images dans $S_\alpha$
contenues dans des ouverts affines, on peut aussi supposer que $S_\alpha$ est
affine, soient donc $S_\alpha=\operatorname{Spec}(A_\alpha)$,
$Y_\alpha=\operatorname{Spec}(C_\alpha)$, $C_\alpha$ étant une
$A_\alpha$-algèbre de présentation finie, $S=\operatorname{Spec}(A)$,
$Y=\operatorname{Spec}(C)$, avec $A=\varinjlim A_\lambda$,
$C=\varinjlim(C_\alpha\otimes_{A_\alpha}A_\lambda)=
C_\alpha\otimes_{A_\alpha}A$. On a alors
\[
\operatorname{Hom}_S(X,Y)=
\operatorname{Hom}_{A\text{-alg.}}(C,\Gamma(X,\mathcal O_X))=
\operatorname{Hom}_{A_\alpha\text{-alg.}}(C_\alpha,\Gamma(X,\mathcal O_X))
\]
(I, 2.2.4) et de même
\[
\operatorname{Hom}_{S_\lambda}(X_\lambda,Y_\lambda)=
\operatorname{Hom}_{A_\lambda\text{-alg.}}
(C_\alpha\otimes_{A_\alpha}A_\lambda,\Gamma(X_\lambda,\mathcal O_{X_\lambda}))=
\operatorname{Hom}_{A_\alpha\text{-alg.}}
(C_\alpha,\Gamma(X_\lambda,\mathcal O_{X_\lambda})).
\]
Mais comme $X_\alpha$ est quasi-compact et quasi-séparé, on sait (8.5.4) que
l'on a
$\varinjlim\Gamma(X_\lambda,\mathcal O_{X_\lambda})=
\Gamma(X,\mathcal O_X)$~; puisque $C_\alpha$ est une $A_\alpha$-algèbre de
présentation finie, le fait que (8.8.1.1) est bijectif résulte alors de
(8.8.2.3).

Avant de passer à la démonstration de (ii), notons les corollaires suivants de
(i)~:
\end{proof}

\begin{corollary}[8.8.2.4]
\label{IV.8.8.2.4-fr}
Supposons $S_0$ quasi-compact, $X_\alpha$ de présentation finie sur $S_\alpha$,
$Y_\alpha$ de type fini sur $S_\alpha$ et quasi-séparé sur $S_\alpha$ (ce qui
sera par exemple le cas si $Y_\alpha$ est également
\oldpage[IV-3]{32}
de présentation finie sur $S_\alpha$). Soit
$f_\alpha:X_\alpha\to Y_\alpha$ un $S_\alpha$-morphisme. Pour que
$f:X\to Y$ soit un isomorphisme, il faut et il suffit qu'il existe
$\lambda\geq\alpha$ tel que $f_\lambda:X_\lambda\to Y_\lambda$ soit un
isomorphisme.
\end{corollary}

\begin{proof}
La condition est évidemment suffisante. Pour prouver qu'elle est nécessaire,
notons d'abord que la question étant locale sur $S_0$ (puisque $L$ est
filtrant), on peut toujours supposer $S_0$ affine, donc quasi-séparé. Il y a
par hypothèse un $S$-morphisme $g:Y\to X$ tel que $g\circ f=1_X$ et
$f\circ g=1_Y$. Comme $X_\alpha$ est de présentation finie sur $S_\alpha$,
et $Y_\alpha$ quasi-compact et quasi-séparé (puisque $S_\alpha$ est
quasi-compact et quasi-séparé), il existe un $\mu\geq\alpha$ et un
$S_\mu$-morphisme $g_\mu:Y_\mu\to X_\mu$ tels que
$g=g_\mu\times1_S$ en vertu de (8.8.2, (i)). D'autre part, il résulte aussi de
(8.8.2, (i)) et des relations $g\circ f=1_X$ et $f\circ g=1_Y$ qu'il existe
$\nu\geq\mu$ tel que l'on ait $g_\nu\circ f_\nu=1_{X_\nu}$ et
$f_\nu\circ g_\nu=1_{Y_\nu}$, puisque $X_\alpha$ et $Y_\alpha$ sont de type
fini sur $S_\alpha$ et quasi-compacts. Cela signifie que
$f_\nu:X_\nu\to Y_\nu$ est un isomorphisme, d'où le corollaire.
\end{proof}

\begin{corollary}[8.8.2.5]
\label{IV.8.8.2.5-fr}
Supposons $S_0$ quasi-compact et quasi-séparé, $X_\alpha$ et $Y_\alpha$ de
présentation finie sur $S_\alpha$. Pour que $X$ et $Y$ soient
$S$-isomorphes, il faut et il suffit qu'il existe $\lambda\geq\alpha$ tels que
$X_\lambda$ et $Y_\lambda$ soient $S_\lambda$-isomorphes. De plus, pour tout
$S$-isomorphisme $f:X\xrightarrow{\sim}Y$, il existe $\mu\geq\lambda$ et un
$S_\mu$-isomorphisme $f_\mu:X_\mu\to Y_\mu$ tels que
$f=f_\mu\times1_S$.
\end{corollary}

\begin{proof}
La condition est évidemment suffisante~; inversement, s'il existe un
$S$-isomorphisme $f:X\to Y$, il résulte de (8.8.2, (i)) que $f$ est de la
forme $f_\mu\times1_S$ pour un $\mu\geq\alpha$ et un homomorphisme
$f_\mu:X_\mu\to Y_\mu$~; mais comme $f$ est un isomorphisme, il résulte de
(8.8.2.4) qu'il existe $\nu\geq\mu$ tel que
$f_\nu:X_\nu\to Y_\nu$ soit un isomorphisme.
\end{proof}

\begin{proof}[Démonstration de (8.8.2), (ii)]
Considérons encore en premier lieu le cas où $S_0=\operatorname{Spec}(A_0)$ et
$X=\operatorname{Spec}(B)$ sont affines~; alors l'assertion de (ii) équivaut
au lemme (1.8.4.2).

Pour démontrer (ii) dans le cas général, notons que $S$ est quasi-compact et
quasi-séparé~; comme $X$ est de présentation finie sur $S$ et $S$ affine sur
$S_0$, il existe donc un recouvrement fini $(U_i)$ de $S_0$ et, si $(W_i)$ est
le recouvrement ouvert affine de $S$ formé des images réciproques des $U_i$
par le morphisme $S\to S_0$, un recouvrement fini $(X_r)$ de $X$ par des
ouverts affines tel que l'image par $X\to S$ de chaque $X_r$ soit contenue
dans un $W_{i(r)}$~; l'anneau $A(X_r)$ est alors une algèbre de présentation
finie sur l'anneau $A(W_{i(r)})$ (1.4.6). En vertu du lemme (1.8.4.2) et du
fait que $L$ est filtrant, il existe un indice $\lambda\in L$ et, pour chaque
indice $r$, un schéma affine $Z_{r\lambda}$, de présentation finie sur l'image
réciproque $W_{i(r),\lambda}$ de $U_{i(r)}$ dans $S_\lambda$, et un
$S$-isomorphisme
$g_r:Z_{r\lambda}\times_{S_\lambda}S\xrightarrow{\sim}X_r$. Soit $Z_{rs}$
l'image réciproque par $g_r$ du sous-préschéma induit sur l'ouvert
$X_r\cap X_s$ de $X_r$, qui est quasi-compact puisque $X$ est quasi-séparé,
et désignons par $g'_{rs}$ la restriction
$Z_{rs}\xrightarrow{\sim}X_r\cap X_s$ de l'isomorphisme $g_r$. En vertu de
(8.8.2.5), il existe un $\mu\geq\lambda$ et, pour tout couple $(r,s)$, un
ouvert quasi-compact $Z_{rs\mu}$ de
$Z_{r\mu}=v_{\lambda\mu}^{-1}(Z_{r\lambda})$ tels que $Z_{rs}$ soit l'image
réciproque de $Z_{rs\mu}$~; d'ailleurs, comme $S_\mu$ est quasi-séparé, et
$W_{i(r),\mu}$ ouvert quasi-compact dans $S_\mu$, chacun des $Z_{rs\mu}$ est
de présentation finie sur $S_\mu$ (I, 1.6.2). Considérons alors, pour tout
couple $(r,s)$, l'isomorphisme
$h_{sr}=g_{sr}^{'-1}\circ g'_{rs}$ de $Z_{rs}$ sur $Z_{sr}$~; il résulte de
(8.8.2.4) qu'il existe un $\nu\geq\mu$ et, pour tout couple $(r,s)$, un
isomorphisme $h_{sr\nu}:Z_{rs\nu}\to Z_{sr\nu}$ tels que
$h_{sr}=h_{sr\nu}\times1_S$. Enfin, pour tout triplet $(r,s,t)$ désignons par
$h'_{sr}$ la restriction de $h_{sr}$ à
\oldpage[IV-3]{33}
$Z_{rs}\cap Z_{rt}$~; il résulte aussitôt des définitions que $h'_{sr}$ est un
isomorphisme de $Z_{rs}\cap Z_{rt}$ sur $Z_{sr}\cap Z_{st}$, et que l'on a la
relation $h'_{ts}\circ h'_{sr}=h'_{tr}$. En vertu de (8.8.2, (i)), il existe
donc $\rho\geq\nu$ tel que, pour tout triplet $(r,s,t)$, on ait
$h'_{ts\rho}\circ h'_{sr\rho}=h'_{tr\rho}$. On en conclut que les
isomorphismes $h_{sr\rho}$ vérifient la condition de recollement
$(0_{\mathrm I}, 4.1.7)$ et définissent donc un préschéma $X_\rho$ tel que $X$
soit isomorphe à $X_\rho\times_{S_\rho}S$. En outre, les $Z_{r\rho}$ sont de
présentation finie sur $S_\rho$ et, si on les identifie à des ouverts de
$X_\rho$, les intersections $Z_{r\rho}\cap Z_{s\rho}$, isomorphes aux
$Z_{rs\rho}$, sont quasi-compactes, donc (I, 1.6.3) $X_\rho$ est de
présentation finie sur $S_\rho$.
\end{proof}

\begin{env}[8.8.3]
\label{IV.8.8.3-fr}
Le contenu essentiel de (8.8.2) s'exprime encore en disant que si $S_0$ est
quasi-compact et quasi-séparé, la catégorie des $S$-préschémas de présentation
finie est déterminée à équivalence près par la donnée des catégories des
$S_\lambda$-préschémas de présentation finie, des foncteurs
$\rho_{\mu\lambda}:X_\lambda\rightsquigarrow
X_\lambda\times_{S_\lambda}S_\mu$ (pour $\lambda\leq\mu$) entre ces catégories
et des isomorphismes de transitivité
$\rho_{\nu\mu}\circ\rho_{\mu\lambda}\xrightarrow{\sim}\rho_{\nu\lambda}$. De
façon imagée, on peut dire que la donnée d'un $S$-préschéma de présentation
finie $X$ équivaut «~fonctoriellement~» à la donnée d'un
$S_\lambda$-préschéma de présentation finie $X'_\lambda$~; si un
$S_\mu$-préschéma de présentation finie $X''_\mu$ est tel que $X$ soit
$S$-isomorphe à $X''_\mu\times_{S_\mu}S$, il existe un $\nu$ tel que
$\lambda\leq\nu$, $\mu\leq\nu$ et tel que
$X'_\lambda\times_{S_\lambda}S_\nu$ et
$X''_\mu\times_{S_\mu}S_\nu$ soient $S_\nu$-isomorphes. Le fait que $L$ est
filtrant entraîne en outre que si on se donne une famille \emph{finie}
$(X^{(i)})_{i\in I}$ de $S$-préschémas de présentation finie, et une famille
\emph{finie} $(f^{(j)})_{j\in J}$ de $S$-morphismes entre ces préschémas
($f^{(j)}$ étant donc de la forme
$X^{(\sigma(j))}\to X^{(\tau(j))}$ où $\sigma$ et $\tau$ sont deux
applications de $J$ dans $I$), alors il y a un indice $\lambda\in L$, une
famille $(X_\lambda^{(i)})_{i\in I}$ de $S_\lambda$-préschémas de présentation
finie et une famille $(f_\lambda^{(j)})_{j\in J}$ de $S_\lambda$-morphismes
tels que $X^{(i)}$ s'identifie à
$X_\lambda^{(i)}\times_{S_\lambda}S$ et $f^{(j)}$ à
$f_\lambda^{(j)}\times1_S$ quels que soient $i$ et $j$~; en outre, si on a
une relation $f^{(j)}=f^{(k)}$, on peut supposer $\lambda$ choisi de telle
sorte que $f_\lambda^{(j)}=f_\lambda^{(k)}$.

Considérons en particulier trois $S$-préschémas de présentation finie $X$,
$Y$, $Z$ et deux $S$-morphismes $f:X\to Z$, $g:Y\to Z$, de sorte que le
produit $X\times_ZY$ relatif à ces morphismes est encore un $S$-préschéma de
présentation finie (I, 1.6.2). Alors il résulte de ce qui précède et de
(I, 3.3.11) que l'on peut écrire
$X\times_ZY=(X_\lambda\times_{Z_\lambda}Y_\lambda)\times_{S_\lambda}S$ pour
un $\lambda$ convenable, $X_\lambda$, $Y_\lambda$, $Z_\lambda$ étant des
$S_\lambda$-préschémas de présentation finie~; on peut donc dire que la
détermination des $S$-préschémas de présentation finie par la donnée des
$S_\lambda$-préschémas de présentation finie est «~\emph{compatible avec les
produits fibrés}~». On a vu d'autre part (8.6.3) que si $g$ est une immersion,
on peut supposer qu'il en est de même de $g_\lambda:Y_\lambda\to Z_\lambda$~;
identifiant $Y$ (resp. $Y_\lambda$) avec un sous-préschéma de $Z$ (resp.
$Z_\lambda$), on voit donc que l'on peut écrire, pour un $\lambda$ convenable,
$f^{-1}(Y)=f_\lambda^{-1}(Y_\lambda)\times_{S_\lambda}S$ (I, 4.4.1)~; il y a
donc aussi «~\emph{compatibilité avec la formation des images réciproques de
sous-préschémas}~». Plus particulièrement, si $f_1$, $f_2$ sont deux
$S$-morphismes de $X$ dans $Y$, on appelle \emph{noyau} de ce couple de
morphismes l'image réciproque $N$ du sous-préschéma diagonal de $Y\times_SY$
par le $S$-morphisme $(f_1,f_2)_S$~; on déduit de ce qui précède que l'on aura
pour $\lambda$ convenable $N=N_\lambda\times_{S_\lambda}S$, où $N_\lambda$
est le noyau du couple de $S_\lambda$-morphismes
$(f_{1\lambda},f_{2\lambda})$ de $X_\lambda$ dans $Y_\lambda$. Ces remarques
s'étendent à des produits finis quelconques et aux «~noyaux~» de systèmes
finis quelconques de $S$-morphismes d'un $S$-préschéma de présentation finie
\oldpage[IV-3]{34}
dans un autre~; on en conclut que d'une façon générale la formation des
$S$-préschémas de présentation finie par la donnée des
$S_\lambda$-préschémas de présentation finie est «~\emph{compatible avec les
limites projectives finies}~», une telle limite étant par définition le noyau
d'un système fini de morphismes d'un même $S$-préschéma dans un produit de
$S$-préschémas (T, 1.8).

Nous nous permettrons couramment par la suite de faire les traductions
impliquées par les propriétés précédentes (ainsi que par (8.3.11), (8.5.2) et
(8.6.3)) sans toujours nous astreindre à les justifier pas à pas comme
ci-dessus. Par exemple, la donnée d'un «~\emph{préschéma en groupes}~» $G$ de
présentation finie sur $S$ équivaut à la donnée d'un préschéma en groupes
$G_\lambda$ de présentation finie sur un $S_\lambda$ pour un $\lambda$ assez
grand~; écrire en effet la condition d'associativité pour le $S$-morphisme
«~loi de composition~» $G\times_SG\to G$ du préschéma en groupes revient à
écrire que le noyau de deux $S$-morphismes de la forme
$G\times_SG\times_SG\to G$ est égal à $G\times_SG\times_SG$ (II, 8.3.9), et
les autres conditions intervenant dans la définition d'un préschéma en
groupes s'interprètent de même.
\end{env}

\subsection{Premières applications à l'élimination des hypothèses noethériennes}
\label{subsection:IV.8.9-fr}

\begin{proposition}[8.9.1]
\EGAReaderTarget{ega4.c04.prop.8-9-1}{IV.8.9.1-fr}
Soient $A$ un anneau, $X$ un $A$-préschéma.
\begin{enumerate}
\item[(i)] Les conditions suivantes sont équivalentes~:
\begin{enumerate}
\item[\emph{a)}] $X$ est de présentation finie sur $A$.
\item[\emph{b)}] Il existe un anneau noethérien $A_0$, un préschéma $X_0$ de
type fini sur $A_0$, un homomorphisme d'anneaux $A_0\to A$, et un
$A$-isomorphisme $X_0\otimes_{A_0}A\xrightarrow{\sim}X$.
\item[\emph{c)}] Il existe un sous-anneau $A_0$ de $A$, qui est une
$\mathbf Z$-algèbre de type fini, un préschéma $X_0$ de type fini sur $A_0$ et
un $A$-isomorphisme $X_0\otimes_{A_0}A\xrightarrow{\sim}X$.
\end{enumerate}
\item[(ii)] Si $\mathcal F$ est un $\mathcal O_X$-Module quasi-cohérent de
présentation finie, on peut supposer le sous-anneau $A_0$ tel qu'il existe un
$\mathcal O_{X_0}$-Module cohérent $\mathcal F_0$ tel que $\mathcal F$ soit
isomorphe à $\mathcal F_0\otimes_{A_0}A$~;
$\operatorname{Supp}(\mathcal F)$ est constructible et fermé dans $X$, et il
y a un sous-préschéma fermé $Z$ de $X$, ayant
$\operatorname{Supp}(\mathcal F)$ comme espace sous-jacent, tel que
l'immersion canonique $Z\to X$ soit de présentation finie.
\item[(iii)] Si $Y$ est un second $A$-préschéma de présentation finie, et
$f:X\to Y$ un $A$-morphisme, on peut supposer le sous-anneau $A_0$ de $A$ tel
qu'il existe un préschéma $Y_0$ de type fini sur $A_0$, un $A$-isomorphisme
$Y_0\otimes_{A_0}A\xrightarrow{\sim}Y$ et un $A_0$-morphisme
$f_0:X_0\to Y_0$ (nécessairement de type fini) tels que $f$ s'identifie à
$f_0\otimes1_A$.
\end{enumerate}
\end{proposition}

\begin{proof}
(i) Comme $A$ est limite inductive de ses sous-anneaux qui sont de type fini
sur $\mathbf Z$, \emph{a)} implique \emph{c)} en vertu de (8.8.2, (ii))~;
\emph{c)} implique \emph{b)} car une $\mathbf Z$-algèbre de type fini est un
anneau noethérien~; enfin, si $A_0$ est noethérien, un $A_0$-préschéma de type
fini est de présentation finie (I, 1.6.1), donc \emph{b)} implique \emph{a)}
en vertu de (I, 1.6.2, (iii)).

L'assertion (ii) se déduit immédiatement de (8.5.2, (ii)), (8.3.11) et
(I, 1.8.2) et l'assertion (iii) de (8.8.2, (i)).
\end{proof}

\begin{env}[8.9.2]
\label{IV.8.9.2-fr}
La proposition (8.9.1) et les résultats de (8.5.2) et (8.8.2) permettent
d'étendre dans de nombreux cas à des morphismes de présentation finie
$X\to Y$ des résultats prouvés par des techniques noethériennes en supposant
$Y$ \emph{localement noethérien}.
\oldpage[IV-3]{35}
Nous en verrons de nombreux exemples par la suite et nous nous bornerons ici
à donner quelques résultats de ce type.
\end{env}

\begin{proposition}[8.9.3]
\label{IV.8.9.3-fr}
Soient $A$ un anneau, $M$ un $A$-module de présentation finie. Alors tout
$A$-endomorphisme surjectif de $M$ est bijectif.
\end{proposition}

\begin{proof}
En effet, considérons $A$ comme limite inductive de ses sous-$\mathbf Z$-algèbres
de type fini. Il résulte de (8.5.2.6) qu'il existe une de ces sous-algèbres
$A_0$ et un $A_0$-module $M_0$ de présentation finie tels que $M$ soit
$A$-isomorphe à $M_0\otimes_{A_0}A$~; en outre, si $f:M\to M$ est un
$A$-endomorphisme surjectif, on peut supposer (8.5.2, (i)) qu'il existe un
$A_0$-endomorphisme $f_0:M_0\to M_0$ tel que $f=f_0\otimes1_A$~; enfin
(8.5.7) on peut supposer que $f_0$ est surjectif. Mais comme $A_0$ est
noethérien et $M_0$ un $A_0$-module de type fini, $M_0$ est un $A_0$-module
noethérien, donc (Bourbaki, \emph{Alg.}, chap. VIII, \S~2, n\textsuperscript{o}~2,
lemme 3) $f_0$ est bijectif, et il en est par suite de même de $f$.
\end{proof}

\begin{proposition}[8.9.4]
\label{IV.8.9.4-fr}
(«~théorème de platitude générique~»). --- Soient $Y$ un préschéma intègre,
$u:X\to Y$ un morphisme de type fini et localement de présentation finie,
$\mathcal F$ un $\mathcal O_X$-Module quasi-cohérent de présentation finie. Il
existe alors un ouvert non vide $U$ de $Y$ tel que
$\mathcal F|u^{-1}(U)$ soit plat sur $U$.
\end{proposition}

\begin{proof}
Le raisonnement du début de (6.9.1) ramène à prouver le

\begin{lemma}[8.9.4.1]
\label{IV.8.9.4.1-fr}
Soient $A$ un anneau intègre, $B$ une $A$-algèbre de présentation finie,
$M$ un $B$-module de présentation finie. Alors il existe un $f\neq0$ dans
$A$ tel que $M_f$ soit un $A_f$-module libre.
\end{lemma}

En effet, d'après (8.9.1) il y a une sous-$\mathbf Z$-algèbre de type fini
$A_0$ de $A$, une $A_0$-algèbre de type fini $B_0$ et un $B_0$-module de type
fini $M_0$ tels que $B$ soit $A$-isomorphe à $B_0\otimes_{A_0}A$ et que $M$
soit $B$-isomorphe à $M_0\otimes_{A_0}A$~; d'après (6.9.2), il existe
$f_0\neq0$ dans $A_0$ tel que $(M_0)_{f_0}$ soit un
$(A_0)_{f_0}$-module libre. Comme
$M_{f_0}=(M_0)_{f_0}\otimes_{A_0}A$ et
$A_{f_0}=(A_0)_{f_0}\otimes_{A_0}A$, $M_{f_0}$ est un $A_{f_0}$-module libre.
\end{proof}

\begin{corollary}[8.9.5]
\label{IV.8.9.5-fr}
Soient $S$ un préschéma quasi-compact et quasi-séparé, $u:X\to S$ un
morphisme de présentation finie, $\mathcal F$ un $\mathcal O_X$-Module
quasi-cohérent de présentation finie. Il existe alors une partition
$(S_i)_{1\leq i\leq n}$ de $S$ en un nombre fini d'ensembles localement fermés
dans $S$ telle que, pour $1\leq i\leq n$, il existe un sous-préschéma $S'_i$
de $S$, ayant pour espace sous-jacent $S_i$, de présentation finie sur $S$,
tel que si l'on pose $X_i=X\times_SS'_i$, le $\mathcal O_{X_i}$-Module
$\mathcal F_i:=\mathcal F\otimes_{\mathcal O_S}\mathcal O_{S'_i}$ soit plat
sur $S'_i$.
\end{corollary}

\begin{proof}
Considérons un recouvrement fini $(U_j)_{1\leq j\leq m}$ de $S$ par des
ouverts affines et définissons par récurrence $T_1=U_1$,
$T_j=U_j-\bigcup_{h<j}(U_j\cap U_h)$ pour $j\geq2$~; chaque $T_j$ est fermé
dans l'ouvert affine $U_j$, et les $T_j$ sont deux à deux sans point commun~;
en outre les $U_j\cap U_h$ sont quasi-compacts puisque $S$ est quasi-séparé,
et par suite ($S$ étant aussi quasi-compact), sont constructibles dans $S$
(I, 1.8.1), donc il en est de même des $T_j$. Il suffira évidemment de
démontrer le corollaire pour un sous-préschéma convenable $T'_j$ de $S$ ayant
$T_j$ pour espace sous-jacent, de présentation finie sur $S$ et pour le
morphisme et le Module déduits respectivement de $u$ et $\mathcal F$ par le
changement de base $T'_j\to S$. Notons pour cela que si l'on prend sur $U_j$
la structure de préschéma induite par celle de $S$, l'immersion ouverte
$U_j\to S$ est quasi-compacte puisque $S$ est quasi-séparé (I, 1.2.7), donc
elle est de
\oldpage[IV-3]{36}
présentation finie (I, 1.6.2, (i)). Il suffira donc que $T'_j$ soit de
présentation finie sur $U_j$, autrement dit, on peut déjà se borner au cas où
$U_j=S$ et $T_j=T$ est fermé constructible dans $S$.

Soit $S=\operatorname{Spec}(A)$, et considérons $A$ comme limite inductive de
ses sous-$\mathbf Z$-algèbres de type fini, de sorte que
$S=\varprojlim S_\lambda$, où les $S_\lambda$ sont affines et noethériens. En
vertu de (8.3.11), il existe un $\lambda$ et une partie fermée (nécessairement
constructible) $T_\lambda$ de $S_\lambda$ tels que
$T=u_\lambda^{-1}(T_\lambda)$ ($u_\lambda:S\to S_\lambda$ étant le morphisme
canonique). On munira $T_\lambda$ d'une structure de sous-préschéma de
$S_\lambda$ et on prendra $T'=T_\lambda\times_{S_\lambda}S$~; comme
l'immersion canonique $T_\lambda\to S_\lambda$ est de présentation finie
(I, 1.6.1), il en est de même de l'immersion $T'\to S$. Comme $T'$ est affine,
on voit finalement qu'on peut se borner au cas où $T'=S$ est affine.

Alors (8.9.1), avec les mêmes notations, il existe un $\lambda$, un morphisme
de type fini $u_\lambda:X_\lambda\to S_\lambda$ et un
$\mathcal O_{X_\lambda}$-Module cohérent $\mathcal F_\lambda$ tels que $X$
soit isomorphe à $X_\lambda\times_{S_\lambda}S$, $\mathcal F$ à
$\mathcal F_\lambda\otimes_{\mathcal O_{X_\lambda}}\mathcal O_X$. On peut
alors appliquer à $S_\lambda$, $X_\lambda$ et $\mathcal F_\lambda$ le
résultat de (6.9.3) et il y a des sous-préschémas $S'_{\lambda i}$ de
$S_\lambda$, dont les ensembles sous-jacents forment une partition de
$S_\lambda$ et qui sont tels que si l'on pose
$X_{\lambda i}=X_\lambda\times_{S_\lambda}S'_{\lambda i}$,
$\mathcal F_{\lambda i}=\mathcal F_\lambda\otimes_{\mathcal O_{S_\lambda}}
\mathcal O_{S'_{\lambda i}}$ soit plat sur $S'_{\lambda i}$. Les
$S'_i=S'_{\lambda i}\times_{S_\lambda}S$ sont alors des sous-préschémas de
$S$ répondant à la question en vertu de (2.1.4).
\end{proof}

\subsection{Propriétés de permanence des morphismes par passage à la limite projective}
\label{subsection:IV.8.10-fr}

Dans cette section, nous conservons les hypothèses générales et les notations
de (8.8.1).

\begin{proposition}[8.10.1]
\EGAReaderTarget{ega4.c04.prop.8-10-1}{IV.8.10.1-fr}
S'il existe $\lambda$ tel que, pour $\mu\geq\lambda$, $f_\mu$ soit un morphisme
ouvert (resp. universellement ouvert), alors $f$ est un morphisme ouvert
(resp. universellement ouvert).
\end{proposition}

\begin{proof}
Comme $X=X_\lambda\times_{Y_\lambda}Y$, l'assertion relative aux morphismes
universellement ouverts est un cas particulier de (2.4.3, (iii)). Supposons
donc $f_\mu$ ouvert pour $\mu\geq\lambda$~; il suffit de voir que pour tout
ouvert quasi-compact $U$ de $X$, $f(U)$ est ouvert dans $Y$. Or, il existe
$\mu\geq\lambda$ tel que $U=v_\mu^{-1}(U_\mu)$, où $U_\mu$ est un ouvert
quasi-compact dans $X_\mu$ (2.3.11)~; on a alors
$f(v_\mu^{-1}(U_\mu))=w_\mu^{-1}(f_\mu(U_\mu))$ (I, 3.4.8), donc $f(U)$ est
ouvert dans $Y$.
\end{proof}

\begin{corollary}[8.10.2]
\label{IV.8.10.2-fr}
Soit $f:X\to Y$ un morphisme. Pour que $f$ soit universellement ouvert, il
suffit que, pour tout entier $n>0$, si l'on pose
$Y_n=Y\otimes_{\mathbf Z}\mathbf Z[T_1,\ldots,T_n](=\mathbf V_Y^n)$, et
$X_n=X\times_YY_n$, la projection canonique
$f_n=f\times1_{Y_n}:X_n\to Y_n$ soit un morphisme ouvert.
\end{corollary}

\begin{proof}
Pour démontrer que $f$ est universellement ouvert, il suffit de prouver qu'il
en est ainsi de la restriction $f^{-1}(U)\to U$ de $f$ pour tout ouvert affine
$U$ de $Y$~; comme, par hypothèse, si
$U_n=U\otimes_{\mathbf Z}\mathbf Z[T_1,\ldots,T_n]$ est l'image réciproque de
$U$ dans $Y_n$, le morphisme $f_n^{-1}(U_n)\to U_n$, restriction de $f_n$, est
ouvert, on voit qu'on peut se borner au cas où $Y=\operatorname{Spec}(A)$ est
affine. En outre, il suffit aussi évidemment de montrer que pour tout
morphisme $Y'\to Y$, où $Y'=\operatorname{Spec}(A')$ est lui aussi affine,
$f'=f_{(Y')}$ est ouvert. Supposons d'abord que $A'$ soit une $A$-algèbre de
type fini, donc quotient d'une algèbre de polynômes $A[T_1,\ldots,T_n]$~;
alors $Y'$ est un sous-préschéma fermé de $Y_n$ et $f'$ la
\oldpage[IV-3]{37}
restriction de $f_n$ à $f_n^{-1}(Y')$~; mais pour tout ouvert $V$ de $X_n$,
on a $f_n(V\cap f_n^{-1}(Y'))=f_n(V)\cap Y'$, et comme par hypothèse $f_n(V)$
est ouvert dans $Y_n$, cela montre que $f'$ est aussi un morphisme ouvert.
Lorsque $A'$ est quelconque, il peut être considéré comme limite inductive de
ses sous-$A$-algèbres de type fini $A'_\lambda$, et le fait que $f'$ soit un
morphisme ouvert résulte de ce qui précède et de (8.10.1).
\end{proof}

\begin{proposition}[8.10.3]
\EGAReaderTarget{ega4.c04.prop.8-10-3}{IV.8.10.3-fr}
Supposons qu'il existe $\alpha$ tel que~: 1\textsuperscript{o} $S_\alpha$ soit
quasi-compact~; 2\textsuperscript{o} les morphismes $X_\alpha\to S_\alpha$,
$Y_\alpha\to S_\alpha$ soient quasi-compacts et le morphisme
$Y_\alpha\to S_\alpha$ quasi-séparé~; 3\textsuperscript{o} pour
$\alpha\leq\lambda\leq\mu$, les morphismes
$u_{\lambda\mu}:S_\mu\to S_\lambda$ soient plats~; 4\textsuperscript{o}
$\overline{f_\alpha(X_\alpha)}$ soit constructible dans $Y_\alpha$. Alors,
pour que $f$ soit dominant, il faut et il suffit qu'il existe
$\lambda\geq\alpha$ tel que $f_\lambda$ soit dominant.
\end{proposition}

\begin{proof}
Les hypothèses entraînent que $Y_\alpha$ est quasi-compact et que le morphisme
$f_\alpha$ est quasi-compact (I, 1.2.4)~; par suite
$\overline{f_\alpha(X_\alpha)}=Z_\alpha$ est pro-constructible (I, 1.9.5,
(vii)) dans $Y_\alpha$. Si l'on pose
$Z_\lambda=w_{\alpha\lambda}^{-1}(Z_\alpha)$ pour $\lambda\geq\alpha$ et
$Z=w_\alpha^{-1}(Z_\alpha)$, on a donc
$Z_\lambda=\overline{f_\lambda(X_\lambda)}$ et
$Z=\overline{f(X)}$ (I, 3.4.8), et $Z_\lambda$ est pro-constructible dans
$Y_\lambda$ (I, 1.9.5, (vii)). Il suffit alors d'appliquer (8.3.13) en
remplaçant $S_\lambda$, $Z'_\lambda$ et $Z''_\lambda$ par $Y_\lambda$,
$Y_\lambda$ et $Z_\lambda$ respectivement.
\end{proof}

\begin{proposition}[8.10.4]
\label{IV.8.10.4-fr}
Supposons qu'il existe $\alpha$ tel que $Y_\alpha$ soit quasi-compact et
$f_\alpha$ de type fini et quasi-séparé. Pour que le morphisme $f$ soit
séparé, il faut et il suffit qu'il existe $\lambda\geq\alpha$ tel que
$f_\lambda$ soit séparé.
\end{proposition}

\begin{proof}
La question étant locale sur $Y_\alpha$ (puisque $Y_\alpha$ est quasi-compact
et $L$ filtrant), on peut se borner au cas où $Y_\alpha$ est affine, donc
quasi-séparé, et l'hypothèse entraîne que $X_\alpha$ (donc les $X_\lambda$ et
$X$) sont quasi-compacts et quasi-séparés. Posons
$X'_\lambda=X_\lambda\times_{Y_\lambda}X_\lambda$ pour
$\lambda\geq\alpha$ et $X'=X\times_YX$~; on a
$X'_\lambda=X'_\alpha\times_{Y_\alpha}Y_\lambda$ et
$X'=X'_\alpha\times_{Y_\alpha}Y$~; le morphisme première projection
$X'_\alpha\to X_\alpha$ est quasi-compact et quasi-séparé par hypothèse
(I, 1.2.3, (iii)), donc $X'_\lambda$ est quasi-compact et quasi-séparé.
Notons maintenant que si l'on désigne par $\Delta_\lambda$ (resp. $\Delta$)
la diagonale de $X_\lambda\times_{Y_\lambda}X_\lambda$ (resp. de
$X\times_YX$), il résulte de (I, 5.3.4 et 3.4.8) que $\Delta_\mu$ (resp.
$\Delta$) est l'image réciproque de $\Delta_\lambda$ par l'application
$v'_{\lambda\mu}:X'_\mu\to X'_\lambda$ (resp.
$v'_\lambda:X'\to X'_\lambda$). D'autre part $\Delta_\lambda$ est
constructible dans $X'_\lambda$~: en effet, puisque $f_\lambda$ est
quasi-séparé, l'immersion diagonale $X_\lambda\to X'_\lambda$ est
quasi-compacte, et localement de présentation finie puisque $f_\lambda$ est
de type fini (I, 4.3 et I, 5.4, (iii))~; il résulte alors de (1.8.4.1) que
$\Delta_\lambda$ est localement constructible, donc constructible puisque
$X'_\lambda$ est quasi-compact. On peut maintenant appliquer (8.3.12) en
remplaçant $S_\lambda$ et $Z_\lambda$ par $X'_\lambda$ et $\Delta_\lambda$
respectivement.
\end{proof}

\begin{theorem}[8.10.5]
\EGAReaderTarget{ega4.c04.thm.8-10-5}{IV.8.10.5-fr}
Supposons $S_0$ quasi-compact, $X_\alpha$ et $Y_\alpha$ de présentation finie
sur $S_\alpha$, et soit $f_\alpha:X_\alpha\to Y_\alpha$ un
$S_\alpha$-morphisme. Considérons, pour un morphisme, la propriété d'être~:
\begin{enumerate}
\item[(i)] un isomorphisme~;
\item[(i \emph{bis})] un monomorphisme~;
\item[(ii)] une immersion~;
\item[(iii)] une immersion ouverte~;
\item[(iv)] une immersion fermée~;
\item[(v)] séparé~;
\item[(vi)] surjectif~;
\oldpage[IV-3]{38}
\item[(vii)] radiciel~;
\item[(viii)] affine~;
\item[(ix)] quasi-affine~;
\item[(x)] fini~;
\item[(xi)] quasi-fini~;
\item[(xii)] propre.
\end{enumerate}
Alors, si $\mathbf P$ désigne une des propriétés précédentes, pour que $f$
possède la propriété $\mathbf P$, il faut et il suffit qu'il existe
$\lambda\geq\alpha$ tel que $f_\lambda$ possède la propriété $\mathbf P$
(auquel cas $f_\mu$ la possède aussi pour $\mu\geq\lambda$).

Si $S_0$ est en outre supposé quasi-séparé, la même conclusion est valable
lorsque $\mathbf P$ est la propriété d'être~:
\begin{enumerate}
\item[(xiii)] projectif~;
\item[(xiv)] quasi-projectif.
\end{enumerate}
\end{theorem}

\begin{proof}
Le cas où $\mathbf P$ est l'une des propriétés (i) ou (v) n'a été inséré dans
l'énoncé que pour mémoire~; dans ces cas, le théorème découle de ce qui a été
démontré respectivement dans (8.8.2.4) et (8.10.4). D'ailleurs, compte tenu de
(I, 5.4.1 et 5.3.4), (v) résulte aussi de (iv). Le cas (i \emph{bis}) se
déduit aussitôt de (i), en utilisant (I, 5.3.8) et en notant (comme on l'a
déjà utilisé dans (8.10.4)) que la diagonale $\Delta_f$ se déduit de
$\Delta_{f_\lambda}$ par le changement de base $S\to S_\lambda$.

On notera d'autre part que (vi), (vii) et (xi) sont en fait des conditions
sur les fibres $f^{-1}(y)$ des morphismes envisagés, compte tenu de la
transitivité des fibres par changement de base (I, 3.6.4)~: la condition (vi)
signifie en effet que toutes les fibres doivent être non vides, la condition
(vii) qu'elles doivent être radicielles (I, 3.5.8) et la condition (xi)
qu'elles doivent être finies (III, 6.2.2 et 6.2.3 et I, 6.4.4, compte tenu de
ce que $f$ et les $f_\lambda$ sont des morphismes de type fini par (I, 5.4,
(v)). Le théorème dans ces trois cas sera donc encore conséquence d'un
résultat général sur ce type de propriétés ne concernant que les fibres, qui
sera établi dans (9.3.3)~; nous repoussons donc jusqu'à ce moment-là la
démonstration du théorème dans le cas (xi) (bien entendu, le lecteur pourra
vérifier que, sauf aux n\textsuperscript{os}~8.11 et 8.12, nous ne ferons pas
usage du théorème dans ce cas jusqu'à (9.3.3) et que (8.11) et (8.12) ne
seront pas utilisés avant (9.3.3)).

Pour les cas qui restent à démontrer, on peut se borner à prouver que la
condition de l'énoncé est nécessaire, toutes les propriétés $\mathbf P$
envisagées étant invariantes par changement de base (voir chap. I\textsuperscript{er}
et II dans les numéros concernant chacune de ces propriétés). On peut en
outre supposer que $S_\alpha=S_0$ et que $Y_\alpha=S_\alpha$, donc
$Y_\lambda=S_\lambda$ pour tout $\lambda\geq\alpha$. Enfin, les propriétés
(i) à (xii) sont locales sur $S_0$, donc, comme $S_0$ est réunion finie
d'ouverts affines et $L$ filtrant, on peut se borner pour les démontrer au cas
où $S_0=\operatorname{Spec}(A_0)$ est affine (donc quasi-séparé). On désignera
par $A_\lambda$ (resp. $A$) l'anneau de $S_\lambda$ (resp. $S$).

\emph{Cas (ii), (iii), (iv)}~: Supposons que $f$ soit une immersion (resp.
une immersion ouverte, resp. une immersion fermée), et soit $X'$ le
sous-préschéma (resp. induit sur un ouvert, resp. fermé) de $S$ associé à $f$,
qui est donc un $S$-préschéma de présentation finie.
\oldpage[IV-3]{39}
En vertu de (8.6.3), il existe donc un $\lambda\geq\alpha$ et un
sous-préschéma (resp. induit sur un ouvert, resp. fermé) $X'_\lambda$ de
$S_\lambda$, de présentation finie sur $S_\lambda$, tels que $X'$ soit
isomorphe à $X'_\lambda\times_{S_\lambda}S$. Pour tout $\mu\geq\lambda$,
$X'_\mu=X'_\lambda\times_{S_\lambda}S_\mu$ est donc un sous-préschéma (resp.
induit sur un ouvert, resp. fermé) de $S_\mu$, de présentation finie sur
$S_\mu$, et il résulte donc de (8.8.2.4) et (8.8.2.5) qu'il existe un
$\mu\geq\lambda$ et un isomorphisme $g_\mu:X_\mu\to X'_\mu$ tel que
$g=g_\mu\times1_S$ soit l'isomorphisme $X\to X'$ associé à $f$~; d'où la
conclusion.

\emph{Cas (vi) et (vii)}~: On sait (1.8.4.1) que
$Z_\alpha=f_\alpha(X_\alpha)$ est constructible dans $S_\alpha$~; si l'on
pose $Z_\lambda=u_{\alpha\lambda}^{-1}(Z_\alpha)$ pour $\lambda\geq\alpha$
et $Z=u_\alpha^{-1}(Z_\alpha)$, on a $Z_\lambda=f_\lambda(X_\lambda)$ et
$Z=f(X)$ (I, 3.4.8). Comme en vertu de (8.3.11) l'application canonique
$\varinjlim\mathfrak C(S_\lambda)\to\mathfrak C(S)$ est injective, la
relation $f(X)=S$ implique l'existence d'un $\lambda\geq\alpha$ tel que
$f_\lambda(X_\lambda)=S_\lambda$, ce qui prouve le théorème dans le cas (vi).
Pour le prouver dans le cas (vii), il suffit de noter que le morphisme
structural $X_\alpha\times_{S_\alpha}X_\alpha\to S_\alpha$ est de présentation
finie puisqu'il en est ainsi de la première projection
$X_\alpha\times_{S_\alpha}X_\alpha\to X_\alpha$ (1.6.2)~; il suffit donc en
vertu de (1.8.7.1) d'appliquer le cas (vi) du théorème au morphisme diagonal
$\Delta_{f_\alpha}:X_\alpha\to X_\alpha\times_{S_\alpha}X_\alpha$ en notant
que l'on a $\Delta_{f_\lambda}=\Delta_{f_\alpha}\times1_{S_\lambda}$ et
$\Delta_f=\Delta_{f_\alpha}\times1_S$ (I, 5.3.4 et 3.3.11).

\emph{Cas (viii) et (ix)}~: Comme $S=\operatorname{Spec}(A)$ est affine, dire
que $f:X\to S$ est affine (resp. quasi-affine) signifie qu'il existe un entier
$r$ et une immersion fermée (resp. une immersion)
$j:X\to\mathbf V_S^r=\operatorname{Spec}(A[T_1,\ldots,T_r])$ de
$S$-préschémas, puisque $f$ est de type fini et $S$ quasi-compact
(II, 5.1.9). Comme
$\mathbf V_S^r=\mathbf V_{S_0}^r\times_{S_0}S$, et que
$\mathbf V_{S_0}^r$ est un $S_0$-préschéma de présentation finie, il résulte
de (8.8.2, (i)) appliqué au $S$-morphisme $j$, qu'il existe un $\lambda$ et un
$S_\lambda$-morphisme $j_\lambda:X_\lambda\to\mathbf V_{S_\lambda}^r$ tel que
$j=j_\lambda\times1_S$~; appliquant alors (iv) (resp. (ii)) à $j$, on en
déduit qu'il existe $\mu\geq\lambda$ tel que $j_\mu$ soit une immersion
fermée (resp. une immersion)~; par suite $f_\mu$ est affine (resp.
quasi-affine).

\emph{Cas (x)}~: Par hypothèse, on a $X=\operatorname{Spec}(B)$, où $B$ est
une $A$-algèbre qui est un $A$-module de présentation finie (1.4.7)~; il
résulte donc de (8.5.2, (ii)) qu'il y a un $\lambda$ et un $A_\lambda$-module
de présentation finie $B_\lambda$ tels que le $A$-module $B$ soit isomorphe à
$B_\lambda\otimes_{A_\lambda}A$. La structure de $A$-algèbre de $B$ est
définie par un $A$-homomorphisme $m:B\otimes_AB\to B$~; comme on a
$B\otimes_AB=(B_\lambda\otimes_{A_\lambda}B_\lambda)\otimes_AA$, il existe
d'après (8.5.2, (i)) un $\mu\geq\lambda$ et un $A_\mu$-homomorphisme
$m_\mu:B_\mu\otimes_{A_\mu}B_\mu\to B_\mu$ tel que $m=m_\mu\otimes1$. En
considérant les diagrammes usuels exprimant l'associativité et la
commutativité de $m$, on voit en appliquant encore (8.5.2, (i)) qu'il existe
$\nu\geq\mu$ tel que $m_\nu$ définisse sur $B_\nu$ une multiplication
associative et commutative~; de la même manière on voit qu'on peut supposer
$\nu$ pris assez grand pour que l'anneau $B_\nu$ ainsi défini admette un
élément unité. Si $X'_\nu=\operatorname{Spec}(B_\nu)$, il est clair alors que
$X$ est $S$-isomorphe à $X'_\nu\times_{S_\nu}S$, donc, en vertu de (i), il
existe $\rho\geq\nu$ tel que $X_\rho$ et $X'_\rho$ soient
$S_\rho$-isomorphes, ce qui achève la démonstration dans ce cas.

Pour démontrer le théorème dans le cas (xii), nous démontrerons d'abord la
proposition suivante~:
\end{proof}

\begin{proposition}[8.10.5.1]
\label{IV.8.10.5.1-fr}
(lemme de Chow pour les morphismes de présentation finie). --- Soient $A$ un
anneau, $X$, $Y$ deux $A$-préschémas de présentation finie,
$f:X\to Y$
\oldpage[IV-3]{40}
un $A$-morphisme séparé. Alors il existe deux $A$-préschémas $X'$, $P$ de
présentation finie, et des $A$-morphismes $p:P\to Y$, $j:X'\to P$,
$g:X'\to X$, tels que le diagramme
\[
\begin{tikzcd}[column sep=large,row sep=large]
X' \arrow[r,"j"] \arrow[d,"g"'] & P \arrow[d,"p"] \\
X \arrow[r,"f"'] & Y
\end{tikzcd}
\]
soit commutatif, et que~: 1\textsuperscript{o} $p$ soit projectif~;
2\textsuperscript{o} $g$ soit projectif et surjectif~;
3\textsuperscript{o} $j$ soit une immersion ouverte.
\end{proposition}

\begin{proof}
En effet, soient $A_0\subset A$, $X_0$, $Y_0$ et $f_0$ déterminés comme dans
(8.9.1) de sorte que $Y_0$ est noethérien et $f_0$ de type fini~; on peut en
outre supposer $f_0$ séparé d'après (8.10.4). Le lemme de Chow (II, 5.6.1)
montre alors l'existence de trois morphismes $p_0:P_0\to Y_0$,
$g_0:X'_0\to X_0$ et $j_0:X'_0\to P_0$, de type fini, tels que le diagramme
\[
\begin{tikzcd}[column sep=large,row sep=large]
X'_0 \arrow[r,"j_0"] \arrow[d,"g_0"'] & P_0 \arrow[d,"p_0"] \\
X_0 \arrow[r,"f_0"'] & Y_0
\end{tikzcd}
\]
soit commutatif, et que $p_0$ soit projectif, $g_0$ projectif et surjectif et
$j_0$ une immersion ouverte. Les propriétés de l'énoncé résultent alors de
l'invariance des propriétés précédentes par changement de base (II, 5.5.5,
(iii) et I, 3.5.2 et 4.3.2).
\end{proof}

\begin{proof}[Suite de la démonstration de (8.10.5)]
\emph{Cas (xii)}~: Appliquons au morphisme $f_0:X_0\to S_0$ la prop.
(8.10.5.1)~: on a donc un diagramme commutatif
\[
\begin{tikzcd}[column sep=large,row sep=large]
X'_0 \arrow[r,"j_0"] \arrow[d,"g_0"'] & P_0 \arrow[d,"p_0"] \\
X_0 \arrow[r,"f_0"'] & S_0
\end{tikzcd}
\]
où $p_0$ est projectif, $g_0$ projectif et surjectif, et $j_0$ une immersion
ouverte~; on en déduit pour chaque $\lambda$ un diagramme analogue où les
morphismes $p_\lambda=p_0\times1_{S_\lambda}$,
$g_\lambda=g_0\times1_{S_\lambda}$ et $j_\lambda=j_0\times1_{S_\lambda}$ ont
respectivement les mêmes propriétés, et de même pour les morphismes limites
projectives $p=p_0\times1_S$, $g=g_0\times1_S$, $j=j_0\times1_S$. Comme $g$
est propre (II, 5.5.3), il en est de même de $p\circ j=f\circ g$ (II, 5.4.2)
et comme $p$ est séparé, $j$ est propre, donc une immersion fermée~;
appliquant le cas (iv) au morphisme $j$ (en notant que $X'_0$ et $P_0$ sont
des $S_0$-préschémas de présentation finie (8.10.5.1 et 1.6.2)), on voit
qu'il existe $\lambda$ tel que $j_\lambda$ soit une immersion fermée, donc
soit propre (II, 5.4.2). Mais alors
$f_\lambda\circ g_\lambda=j_\lambda\circ p_\lambda$ est propre (II, 5.5.3 et
5.4.2) et comme $g_\lambda$ est surjectif, et qu'on peut supposer $f_\lambda$
séparé en vertu de l'hypothèse sur $f$ et de (v), il résulte de (II, 5.4.3)
que $f_\lambda$ est propre.

\emph{Cas (xiii) et (xiv)}~: En vertu de (xii) et de (II, 5.5.3) (qui est
applicable puisque les $S_\lambda$ sont quasi-compacts et quasi-séparés,
compte tenu de (1.7.19)), il suffit de
\oldpage[IV-3]{41}
considérer le cas où $f$ est quasi-projectif. Supposons donc qu'il existe un
$\mathcal O_X$-Module inversible $\mathcal L$ qui soit $f$-ample~; comme $S_0$
est quasi-compact et quasi-séparé il en est de même de $X_0$ (1.2.3), et il y
a donc un $\lambda$ et un $\mathcal O_{X_\lambda}$-Module quasi-cohérent
$\mathcal L_\lambda$ de présentation finie tel que
$\mathcal L=v_\lambda^*(\mathcal L_\lambda)$ (8.5.2, (ii))~; en outre, en
vertu de (8.5.5) on peut supposer $\mathcal L_\lambda$ inversible. Le théorème
dans ce cas est alors conséquence du lemme plus précis suivant.

\begin{lemma}[8.10.5.2]
\label{IV.8.10.5.2-fr}
Supposons $S_0$ quasi-compact, et soit $\mathcal L_\lambda$ un
$\mathcal O_{X_\lambda}$-Module inversible. Pour que $\mathcal L$ soit un
$\mathcal O_X$-Module ample pour $f$ (resp. très ample pour $f$), il faut et
il suffit qu'il existe $\mu\geq\lambda$ tel que $\mathcal L_\mu$ soit ample
pour $f_\mu$ (resp. très ample pour $f_\mu$).
\end{lemma}

La condition étant évidemment suffisante (II, 4.4.10 et 4.6.13), montrons
qu'elle est nécessaire~; les $S_\lambda$ étant quasi-compacts et les
$f_\lambda$ de type fini, on peut, en remplaçant $\mathcal L$ par une puissance
tensorielle convenable, se borner à considérer le cas où $\mathcal L$ est très
ample (II, 4.6.11). En outre, la question étant ici locale sur $S_0$ (vu
(II, 4.4.5) et le fait que $L$ est filtrant), on peut se borner au cas où
$S_0$ (et par suite $S$) est affine. Alors, en vertu de (II, 4.4.1, (ii) et
II, 4.1.2), il existe une $S$-immersion
$j:X\to\mathbf P_S^r=P$ telle que $\mathcal L$ soit isomorphe à
$j^*(\mathcal O_P(1))$. Compte tenu de (8.8.2, (i)), de (ii) et de
(II, 4.1.3), il existe donc un $\mu\geq\lambda$ et une immersion
$j_\mu:X_\mu\to\mathbf P_{S_\mu}^r=P_\mu$ tels que
$j=j_\mu\times1_S$~; utilisant ensuite (II, 4.1.3.2) et (8.5.2.5), on voit
qu'il existe $\nu\geq\mu$ tel que $\mathcal L_\nu$ soit isomorphe à
$j_\nu^*(\mathcal O_{P_\nu}(1))$, ce qui montre que $\mathcal L_\nu$ est très
ample pour $f_\nu$ (II, 4.4.2).
\end{proof}

\subsection{Application aux morphismes quasi-finis}
\label{subsection:IV.8.11-fr}
\label{IV.8.11-fr}

Nous nous proposons dans cette section de démontrer les deux théorèmes
suivants~:

\begin{theorem}[8.11.1]
\EGAReaderTarget{ega4.c04.thm.8-11-1}{IV.8.11.1-fr}
Soit $f:X\to Y$ un morphisme propre, localement de présentation finie, et
quasi-fini. Alors le morphisme $f$ est fini.
\end{theorem}

\begin{theorem}[8.11.2]
\label{IV.8.11.2-fr}
Soit $f:X\to Y$ un morphisme localement de présentation finie, quasi-fini et
séparé. Alors le morphisme $f$ est quasi-affine, et \emph{a fortiori}
quasi-projectif.
\end{theorem}

\begin{remark}[8.11.3]
\label{IV.8.11.3-fr}
On va voir ci-dessous que l'on peut se ramener, pour la démonstration de
(8.11.1) et (8.11.2), au cas où $Y$ est localement noethérien~; on notera que
dans ce cas on obtiendra ainsi une autre démonstration du théorème de
Chevalley (III, 4.4.2).
\end{remark}

\begin{env}[8.11.4]
\label{IV.8.11.4-fr}
Les hypothèses et les conclusions de (8.11.1) et (8.11.2) sont toutes locales
sur $Y$ (I, 6.1, I, 2.6, (II, 5.1.1), (II, 5.4.1) et (II, 6.2.2)), donc on
peut supposer $Y=\operatorname{Spec}(A)$ affine. On sait qu'il existe alors un
sous-anneau $A_0$ de $A$, qui est une $\mathbf Z$-algèbre de type fini, et un
morphisme de type fini $f_0:X_0\to\operatorname{Spec}(A_0)$ tel que $X$
s'identifie à $X_0\otimes_{A_0}A$ et $f$ à $f_0\times1$ (8.9.1). En outre,
$A$ peut être considéré comme limite inductive de ses sous-anneaux contenant
$A_0$ et qui sont des $\mathbf Z$-algèbres de type fini~; utilisant la méthode
de (8.1.2, c)) ainsi que (8.10.5, (v), (xi) et (xii)), on voit qu'il suffit de
démontrer les théorèmes (8.11.1) et (8.11.2) pour $f_0$. Supposons donc
désormais $Y$ noethérien~; utilisant maintenant la méthode de (8.1.2, a))
ainsi que (8.10.5, (v), (ix), (x), (xi) et (xii)) on peut remplacer $Y$ par
$\operatorname{Spec}(\mathcal O_y)$, où $y$ est un
\oldpage[IV-3]{42}
point quelconque de $Y$, donc on voit finalement que l'on peut supposer que
$Y=\operatorname{Spec}(A)$, où $A$ est un \emph{anneau local noethérien}.
Soient $\mathfrak m$ l'idéal maximal de $A$, $\widehat A$ le complété de $A$
pour la topologie $\mathfrak m$-préadique~; on sait que $\widehat A$ est un
anneau local noethérien et que le morphisme
$\operatorname{Spec}(\widehat A)\to\operatorname{Spec}(A)$ est fidèlement
plat et quasi-compact (0\textsubscript{I}, 7.3.5)~; quant à (2.7.1, (i),
(vii), (xiv), (xv) et (xvi)), on voit en outre qu'on peut se borner au cas où
$A$ est \emph{complet}. Il résulte alors de (II, 6.2.6) que
$X=X'\amalg X''$, où $X'$ est un $Y$-schéma fini et $X''$ un $Y$-schéma
quasi-fini tel que $X''\cap f^{-1}(y)=\varnothing$.
\end{env}

Plaçons-nous d'abord dans les hypothèses de (8.11.1)~; comme $f$ est propre,
$X''$, qui est fermé dans $X$, est propre sur $Y$ (II, 5.4.10), donc
$f(X'')$ est fermé dans $Y$~; mais $y$ n'est pas contenu dans $f(X'')$, et
par ailleurs est adhérent à tout point de $Y$, donc $f(X'')=\varnothing$, et
par suite $X''$ est vide et $f$ est \emph{fini}.

Plaçons-nous maintenant dans les hypothèses de (8.11.2) et, en nous ramenant
encore (comme on peut le faire d'après ce qui précède) au cas où
$Y=\operatorname{Spec}(A)$ est affine et noethérien de dimension finie,
raisonnons en outre par récurrence sur la dimension de $Y$. En se ramenant
comme ci-dessus au cas où $A$ est en outre local et complet, on a
$\dim(\mathcal O_y)=\dim(A)=\dim(Y)$ et pour tout $z\neq y$,
$\dim(\mathcal O_z)<\dim(\mathcal O_y)$, donc
$\dim(Y-\{y\})<\dim(Y)$. Or, par hypothèse on a
$f(X'')\subset Y-\{y\}$ et la restriction de $f$ à $X''$ est évidemment un
morphisme quasi-fini et séparé~; appliquant à $Y-\{y\}$ et à $X''$
l'hypothèse de récurrence, on voit que $X''$ est quasi-affine sur
$Y-\{y\}$~; mais $Y-\{y\}$ étant quasi-affine sur $Y$ puisque $Y$ est
noethérien, $X''$ est aussi quasi-affine sur $Y$ (II, 5.1.10, (ii))~; comme
par ailleurs $X'$ est fini (et \emph{a fortiori} affine) sur $Y$, $X$ est
quasi-affine sur $Y$ (II, 4.6.17 et 5.1.2, c')).

\begin{proposition}[8.11.5]
\label{IV.8.11.5-fr}
Soit $f:X\to Y$ un morphisme de présentation finie. Les conditions suivantes
sont équivalentes~:
\begin{enumerate}
\item[\emph{a)}] $f$ est une immersion fermée.
\item[\emph{b)}] $f$ est un monomorphisme propre.
\item[\emph{c)}] $f$ est propre et pour tout $y\in Y$, $f^{-1}(y)$ est
radiciel et géométriquement réduit sur $k(y)$ (c'est-à-dire vide ou
$k(y)$-isomorphe à $\operatorname{Spec}(k(y))$).
\end{enumerate}
\end{proposition}

Il est clair que \emph{a)} entraîne \emph{b)}. Pour voir que \emph{b)}
entraîne \emph{c)}, notons (I, 3.3.12) que pour tout $y\in Y$, le morphisme
$f^{-1}(y)\to\operatorname{Spec}(k(y))$ déduit de $f$ par extension de la base
est un monomorphisme, donc est injectif, et par suite $f^{-1}(y)$ est vide ou
réduit à un point, et en tout cas affine. De plus, si $A$ est l'anneau de
$f^{-1}(y)$, l'homomorphisme canonique $A\otimes_kA\to A$ est bijectif
(I, 5.3.8). Cela entraîne évidemment que $A=k$, sans quoi il y aurait un
élément $a\in A$ non dans $k$ et les images de $a\otimes1$ et $1\otimes a$
dans $A$ seraient toutes deux égales à $a$, alors que $a\otimes1\neq1\otimes a$
puisque $1$ et $a$ forment un système libre sur $k$.

Reste à prouver que \emph{c)} entraîne \emph{a)}. Il résulte tout d'abord de
(8.11.1) que $f$ est un morphisme \emph{fini}, donc
$X=\operatorname{Spec}(\mathcal A)$, où $\mathcal A$ est une
$\mathcal O_Y$-Algèbre \emph{finie}. Il suffit donc de prouver que
l'homomorphisme canonique $\mathcal O_Y\to\mathcal A$ est \emph{surjectif}
(II, 1.4.10), ou encore que pour tout $y\in Y$, l'homomorphisme
$\mathcal O_y\to\mathcal A_y$ est surjectif. Mais par hypothèse
\oldpage[IV-3]{43}
$f^{-1}(y)=\operatorname{Spec}(\mathcal A_y/\mathfrak m_y\mathcal A_y)$
(II, 1.5.5) est tel que l'homomorphisme correspondant
$k(y)=\mathcal O_y/\mathfrak m_y\to
\mathcal A_y/\mathfrak m_y\mathcal A_y$ soit bijectif~; comme $\mathcal A_y$
est un $\mathcal O_y$-module de type fini, le lemme de Nakayama montre que
$\mathcal O_y\to\mathcal A_y$ est surjectif, ce qui achève la démonstration.

\begin{remark}[8.11.5.1]
\label{IV.8.11.5.1-fr}
On notera que le raisonnement précédent prouve que si $f$ est un
\emph{monomorphisme}, alors, pour tout $y\in Y$, $f^{-1}(y)$ est vide ou
$k(y)$-isomorphe à $\operatorname{Spec}(k(y))$.
\end{remark}

\begin{proposition}[8.11.6]
\label{IV.8.11.6-fr}
Si un morphisme $f:X\to Y$ de présentation finie est un homéomorphisme
universel, il est fini, surjectif et radiciel (la réciproque étant vraie
d'après (2.4.5, (i))).
\end{proposition}

\begin{proof}
En effet, $f$ étant de type fini, universellement fermé, et séparé en vertu de
(2.4.4), est propre par définition (II, 5.4.1). Comme il est évidemment
quasi-fini (II, 6.2.3), il est fini d'après (8.11.1). On sait par ailleurs
qu'il est radiciel (2.4.4), et évidemment surjectif.
\end{proof}

\subsection{Nouvelle démonstration et généralisation du «~Main Theorem~» de Zariski}
\label{subsection:IV.8.12-fr}
\label{IV.8.12-fr}

\begin{lemma}[8.12.1]
\label{IV.8.12.1-fr}
Soient $f:X\to Y$ un morphisme quasi-compact et quasi-séparé, $\mathcal C$ une
$\mathcal O_Y$-Algèbre quasi-cohérente, $Z=\operatorname{Spec}(\mathcal C)$,
qui est un $Y$-préschéma affine au-dessus de $Y$. Soient $g:X\to Z$ un
$Y$-morphisme, $\varphi=\mathcal A(g):\mathcal C\to f_*(\mathcal O_X)$ le
$\mathcal O_Y$-homomorphisme correspondant de $\mathcal O_Y$-Algèbres
(II, 1.2.7). Supposons que $g$ soit une immersion. Alors, pour que l'image
fermée de $X$ par $g$ (I, 9.5.3) soit égale à $Z$, il faut et il suffit que
$\varphi$ soit injectif~; $g$ est alors une immersion ouverte.
\end{lemma}

\begin{proof}
L'hypothèse entraîne que $f_*(\mathcal O_X)$ est une $\mathcal O_Y$-Algèbre
quasi-cohérente (1.7.5)~; en outre, comme le morphisme canonique $h:Z\to Y$
est affine, donc quasi-compact et séparé, $g$ est un
morphisme quasi-compact et quasi-séparé (1.2.2 et 1.1.2), donc $g_*(\mathcal O_X)$ est une
$\mathcal O_Z$-Algèbre quasi-cohérente (1.7.5). Cela étant, dire que l'image
fermée de $X$ par $g$ est égale à $Z$ signifie (I, 9.5.2) que
l'homomorphisme canonique $\theta:\mathcal O_Z\to g_*(\mathcal O_X)$ est
injectif. Mais on a $h_*(\mathcal O_Z)=\mathcal C$ par définition de $Z$
(II, 1.3.1), et $h_*(g_*(\mathcal O_X))=f_*(\mathcal O_X)$. Comme $Z$ est
affine au-dessus de $Y$, il revient au même de dire que l'homomorphisme
$\theta:\mathcal O_Z\to g_*(\mathcal O_X)$ est injectif ou que
l'homomorphisme correspondant
$\varphi=h_*(\theta):\mathcal C\to f_*(\mathcal O_X)$ est injectif
(I, 1.3.9). Le fait que $g$ est alors une immersion ouverte résulte de
(I, 9.5.10) et de l'hypothèse que $g$ est une immersion.
\end{proof}

\begin{lemma}[8.12.2]
\label{IV.8.12.2-fr}
Soient $Y$ un préschéma quasi-compact et quasi-séparé, $f:X\to Y$ un
morphisme quasi-séparé de type fini, $\mathcal C$ une $\mathcal O_Y$-Algèbre
quasi-cohérente, $Z=\operatorname{Spec}(\mathcal C)$. Soient $g:X\to Z$ un
$Y$-morphisme, $\varphi:\mathcal C\to f_*(\mathcal O_X)$ le
$\mathcal O_Y$-homomorphisme correspondant de $\mathcal O_Y$-Algèbres. Soit
$(\mathcal C_\lambda)$ la famille filtrante croissante des sous-$\mathcal O_Y$-Algèbres
quasi-cohérentes de type fini de $\mathcal C$ (dont $\mathcal C$ est la
réunion ((I, 9.6.6) et (1.7.9)))~; posons
$Z_\lambda=\operatorname{Spec}(\mathcal C_\lambda)$ et soit $g_\lambda$ le
morphisme composé $X\xrightarrow{g}Z\to Z_\lambda$. Alors les conditions
suivantes sont équivalentes~:
\begin{enumerate}
\item[\emph{a)}] $g$ est une immersion.
\item[\emph{b)}] Il existe $\lambda$ tel que $g_\lambda$ soit une immersion.
\end{enumerate}
\oldpage[IV-3]{44}
De plus, lorsque $g_\lambda$ est une immersion, il en est de même de $g_\mu$
pour $\mu\geq\lambda$.

Il suffit d'appliquer (II, 3.8.4) en remplaçant $\mathcal L$ par
$\mathcal O_X$ et $\mathcal S$ par $\mathcal C[\mathrm T]$, et en tenant
compte de (II, 3.1.7).
\end{lemma}

\begin{proposition}[8.12.3]
\EGAReaderTarget{ega4.c04.prop.8-12-3}{IV.8.12.3-fr}
Soient $Y$ un préschéma quasi-compact et quasi-séparé, $f:X\to Y$ un
morphisme séparé de type fini. Soit $\mathcal B=f_*(\mathcal O_X)$, qui est
une $\mathcal O_Y$-Algèbre quasi-cohérente (I, 9.2.2)~; soit $\mathcal C$ la
fermeture intégrale de $\mathcal O_Y$ dans $\mathcal B$, qui est une
$\mathcal O_Y$-Algèbre quasi-cohérente (II, 6.3.4)~; posons
$Z=\operatorname{Spec}(\mathcal C)$, et soit $g:X\to Z$ le $Y$-morphisme
correspondant à l'injection canonique
$\varphi:\mathcal C\to\mathcal B=f_*(\mathcal O_X)$ (II, 1.2.7). Soit
$(\mathcal C_\lambda)$ la famille filtrante croissante des
sous-$\mathcal O_Y$-Algèbres quasi-cohérentes de type fini de $\mathcal C$
(dont $\mathcal C$ est la réunion ((I, 9.6.6) et (1.7.9))), et, pour tout
$\lambda$, posons $Z_\lambda=\operatorname{Spec}(\mathcal C_\lambda)$, et
soit $g_\lambda:X\to Z_\lambda$ le $Y$-morphisme correspondant à l'injection
$\mathcal C_\lambda\to\mathcal B=f_*(\mathcal O_X)$~; les conditions
suivantes sont équivalentes~:
\begin{enumerate}
\item[\emph{a)}] Il existe une factorisation de $f$ en
\[
X\xrightarrow{f'}Y'\xrightarrow{u}Y,
\]
où $f'$ est une immersion et $u$ un morphisme fini.
\item[\emph{a$'$)}] Il existe une factorisation
$X\xrightarrow{f'}Y'\xrightarrow{u}Y$ de $f$, où $f'$ est une immersion
ouverte et $u$ un morphisme fini.
\item[\emph{b)}] Le morphisme $g:X\to Z$ est une immersion.
\item[\emph{c)}] Il existe $\lambda$ tel que
$g_\lambda:X\to Z_\lambda$ soit une immersion.
\end{enumerate}
De plus, lorsqu'il en est ainsi, $g$ est une immersion ouverte, $g(X)$ est
dense dans $Z$, et il existe $\lambda$ tel que, pour $\mu\geq\lambda$,
$g_\mu$ soit une immersion ouverte.
\end{proposition}

Comme l'homomorphisme $\varphi:\mathcal C\to f_*(\mathcal O_X)$ est
injectif, il résulte de (8.12.1) que si $g$ est une immersion, c'est une
immersion ouverte et $g(X)$ est dense dans $Z$, et il en est de même pour
$g_\lambda$. Le fait que \emph{a)} entraîne \emph{a$'$)} résulte aussi de
(8.12.1), appliqué en remplaçant $Z$ par $Y'$ et $g$ par $f'$ ($Y'$ étant
fini, donc affine au-dessus de $Y$)~: en effet si $Y''$ est l'image fermée de
$X$ par $f'$, $Y''$ est fini au-dessus de $Y$ et $f'$ se factorise en
$X\xrightarrow{f''}Y''\xrightarrow{j}Y'$, où $j$ est l'injection canonique,
et $f''$ est une immersion (I, 4.1.10)~; il résulte donc de (8.12.1) que
$f''$ est une immersion ouverte.

L'équivalence de \emph{b)} et \emph{c)} résulte de (8.12.2) ainsi que le fait
que $g_\lambda$ est alors une immersion pour $\lambda$ assez grand. Il est
clair que \emph{c)} implique \emph{a)}, puisque $Z_\lambda$ est fini
au-dessus de $Y$ (II, 6.3.4 et 6.1.2). Montrons enfin que \emph{a)} implique
\emph{c)}. On a vu ci-dessus qu'on peut supposer que $Y'$ est l'image fermée
de $X$ par $f'$, et alors il résulte de (8.12.1) que si l'on pose
$\mathcal B'=u_*(\mathcal O_{Y'})$, de sorte que $Y'$ s'identifie à
$\operatorname{Spec}(\mathcal B')$, l'homomorphisme
$\varphi':\mathcal B'\to\mathcal B=f_*(\mathcal O_X)$ est \emph{injectif}.
Mais comme par hypothèse $\mathcal B'$ est une $\mathcal O_Y$-Algèbre finie,
elle s'identifie par définition de $\mathcal B$ à une des
sous-$\mathcal O_Y$-Algèbres $\mathcal C_\lambda$, ce qui démontre
\emph{c)}.

Nous dirons qu'un morphisme $f:X\to Y$, où $Y$ est quasi-compact et
quasi-séparé, est \emph{pseudo-fini}, s'il est de type fini et vérifie la
condition \emph{a)} de (8.12.3) (auquel cas il est nécessairement séparé).

\begin{corollary}[8.12.4]
\label{IV.8.12.4-fr}
Soient $Y$ un préschéma quasi-compact et quasi-séparé, $f:X\to Y$ un
morphisme.
\oldpage[IV-3]{45}
\begin{enumerate}
\item[(i)] Supposons $f$ pseudo-fini. Alors, pour tout morphisme $Y'\to Y$,
où $Y'$ est quasi-compact et quasi-séparé,
$f_{(Y')}:X'=X_{(Y')}\to Y'$ est pseudo-fini.
\item[(ii)] Soit $(U_\lambda)$ un recouvrement de $Y$ formé d'ouverts
quasi-compacts. Pour que $f$ soit pseudo-fini, il faut et il suffit que pour
tout $\lambda$, la restriction
$f_\lambda:f^{-1}(U_\lambda)\to U_\lambda$ de $f$ soit un morphisme
pseudo-fini.
\item[(iii)] Supposons en outre $Y$ noethérien, et $f$ de type fini. Alors,
pour que $f$ soit pseudo-fini, il faut et il suffit que, pour tout $y\in Y$,
le morphisme
$f_y=f\times 1:X_y=X\times_Y\operatorname{Spec}(\mathcal O_y)
\to\operatorname{Spec}(\mathcal O_y)$ le soit.
\end{enumerate}
\end{corollary}

\begin{proof}
(i) Il est clair que $f_{(Y')}$ est de type fini (1.5.4)~; en outre, une
factorisation
$X\xrightarrow{g}Z\xrightarrow{u}Y$, où $g$ est une immersion et $u$ est
fini, donne une factorisation
$X'\xrightarrow{g'}Z'\xrightarrow{u'}Y'$ de $f_{(Y')}$, où
$Z'=Z_{(Y')}$, $g'=g_{(Y')}$ et $u'=u_{(Y')}$~; $g'$ est une immersion
(I, 4.3.2) et $u'$ est fini (II, 6.1.5)~; donc $f_{(Y')}$ est pseudo-fini.

(ii) La condition est nécessaire en vertu de (i), les $U_\lambda$ étant
quasi-séparés puisque $Y$ l'est. Pour voir qu'elle est suffisante, observons
(avec les notations de (8.12.3)) que si l'on pose
$X_\lambda=f^{-1}(U_\lambda)$, on a
$\mathcal B|U_\lambda=(f_\lambda)_*(\mathcal O_{X_\lambda})$,
$\mathcal C|U_\lambda$ est la fermeture intégrale de $\mathcal O_{U_\lambda}$
dans $\mathcal B|U_\lambda$ et par suite, si $h:Z\to Y$ est le morphisme
canonique, $Z_\lambda=\operatorname{Spec}(\mathcal C|U_\lambda)$ s'identifie à
$h^{-1}(U_\lambda)$. Or, pour que $g:X\to Z$ soit une immersion, il faut et
il suffit que pour tout $\lambda$, la restriction
$g_\lambda:f^{-1}(U_\lambda)\to h^{-1}(U_\lambda)$ de $g$ le soit
(I, 4.2.4). Cela entraîne la conclusion en vertu de (8.12.3).

(iii) Il suffit, en vertu de (ii), de prouver que $y$ admet un voisinage $U$
tel que la restriction $f^{-1}(U)\to U$ de $f$ soit un morphisme pseudo-fini.
Désignons par $(U_\lambda)$ le système projectif filtrant décroissant des
voisinages ouverts affines de $y$, et appliquons la méthode de (8.1.2,
\emph{a)}). Puisque $Y$ est noethérien, les restrictions
$f_\lambda:f^{-1}(U_\lambda)\to U_\lambda$ de $f$ sont de présentation
finie et il en est de même de $f_y$. Par hypothèse $f_y$ se factorise en
$X_y\xrightarrow{g_y}Z_y\xrightarrow{u_y}\operatorname{Spec}(\mathcal O_y)$,
où $u_y$ est fini et $g_y$ est une immersion. Comme $\mathcal O_y$ est
noethérien, il en est de même de $Z_y$, et comme $Z_y$ est de présentation
finie sur $\operatorname{Spec}(\mathcal O_y)$, il existe un $\lambda$ et un
morphisme de présentation finie $u_\lambda:Z_\lambda\to U_\lambda$ tels que
$Z_y$ s'identifie à
$Z_\lambda\times_{U_\lambda}\operatorname{Spec}(\mathcal O_y)$ et $u_y$ à
$u_\lambda\times 1$ (8.8.2, (ii))~; en outre, il existe un morphisme
$g_\lambda:X_\lambda\to Z_\lambda$ tel que $g_y=g_\lambda\times 1$ et que
$f_\lambda=u_\lambda\circ g_\lambda$ (8.8.2, (i)). De plus, on peut supposer
$\lambda$ pris de sorte que $g_\lambda$ soit une immersion et $u_\lambda$ un
morphisme fini (8.10.5, (ii) et (x)), ce qui prouve que $f_\lambda$ est
pseudo-fini.
\end{proof}

\begin{env}[8.12.5]
\label{IV.8.12.5-fr}
Nous pouvons maintenant donner du «~Main theorem~» de Zariski (III, 4.4.3),
une démonstration n'utilisant pas les résultats cohomologiques de nature
«~globale~» du chap.~III, mais faisant appel par contre aux propriétés plus
fines des anneaux locaux noethériens~; nous généraliserons en outre l'énoncé
du théorème en le débarrassant des hypothèses noethériennes~:
\end{env}

\begin{theorem}[8.12.6]
\label{IV.8.12.6-fr}
(«~Main Theorem~» de Zariski). --- Soit $Y$ un préschéma quasi-compact et
quasi-séparé. Si un morphisme $f:X\to Y$ est quasi-fini, séparé et de
présentation finie, il existe une factorisation de $f$
\begin{equation}
\label{IV.8.12.6.1-fr}
X\xrightarrow{f'}Y'\xrightarrow{u}Y. \tag{8.12.6.1}
\end{equation}
où $f'$ est une immersion ouverte et $u$ un morphisme fini.
\end{theorem}

\oldpage[IV-3]{46}
\begin{proof}
En vertu de (8.12.4, (ii)) et du caractère local (sur $Y$) des notions de
morphisme quasi-fini, séparé et de présentation finie, on peut se borner au
cas où $Y=\operatorname{Spec}(A)$ est affine. Appliquant (8.9.1), on peut
supposer qu'il y a un sous-anneau $A_0$ de $A$, qui est une
$\mathbf Z$-algèbre de type fini, et un $A$-isomorphisme
$X_0\otimes_{A_0}A\xrightarrow{\sim}X$, $f$ s'identifiant par cet
isomorphisme à $f_0\times 1$, où
$f_0:X_0\to\operatorname{Spec}(A_0)$ est un morphisme \emph{de type fini}~;
en outre (8.10.5, (v) et (xii)) on peut supposer que $f_0$ est séparé et
quasi-fini~; si l'on prouve que $f_0$ est pseudo-fini, il en sera de même de
$f$ d'après (8.12.4, (i)). Comme $A_0$ est alors noethérien et que les
notions de morphisme de type fini, séparé et quasi-fini se conservent par
changement de base, il résulte de (8.12.4, (iii)) que l'on peut même supposer
que $A$ est un anneau local, \emph{essentiellement de type fini sur}
$\mathbf Z$ (1.3.8). Posons $n=\dim(A)$, et procédons par récurrence sur $n$~;
pour $n=0$, le théorème est évident, $A$ étant un corps et le morphisme $f$
étant déjà fini (II, 6.2.2). Posons $B=\Gamma(X,\mathcal O_X)$~; désignons
par $C$ la fermeture intégrale de $A$ dans $B$, posons
$Z=\operatorname{Spec}(C)$ et soit $g:X\to Z$ le $Y$-morphisme correspondant
au $A$-homomorphisme canonique injectif $C\to B$~; en vertu de (8.12.3), il
s'agit de montrer que $g$ est une immersion ouverte. Soit $a$ le point fermé
de $Y$, et soit $U=Y-\{a\}$~; $U$ est noethérien et tous ses anneaux locaux
sont essentiellement de type fini sur $\mathbf Z$ et de dimension $<n$~;
compte tenu de l'hypothèse de récurrence, et de (8.12.4, (iii)), on voit que
la restriction $f^{-1}(U)\to U$ de $f$ est un morphisme pseudo-fini. On en
conclut (8.12.3) que, si $h:Z\to Y$ est le morphisme structural, la
restriction $f^{-1}(U)\to h^{-1}(U)$ de $g$ est une immersion ouverte.
Posons $A'=\widehat A$, $Y'=\operatorname{Spec}(A')$,
$X'=X_{(Y')}$, $f'=f_{(Y')}:X'\to Y'$. Comme le morphisme canonique
$u:Y'\to Y$ est plat, il résulte de (2.3.1) que
$B'=\Gamma(X',\mathcal O_{X'})$ s'identifie à la $A'$-algèbre
$B\otimes_A A'$.

D'autre part, comme $A$ est un anneau local \emph{excellent} (7.8.3), le
morphisme $u:Y'\to Y$ est régulier, et \emph{a fortiori} normal, et par suite
(6.14.4) \emph{la fermeture intégrale de $A'$ dans $B'$ est égale à}
$C'=C\otimes_A A'$. On voit donc que $Z'=\operatorname{Spec}(C')$ est égal à
$Z_{(Y')}$ et le morphisme $g':X'\to Z'$ provenant de l'injection $C'\to B'$
est égal à $g_{(Y')}$. Comme $u:Y'\to Y$ est fidèlement plat et
quasi-compact, pour prouver que $g$ est une immersion ouverte, il suffit de
prouver que $g'$ est une immersion ouverte (2.7.1, (x)). Notons maintenant
que $u^{-1}(a)$ est réduit au point fermé $a'$ de $Y'$ et par suite
$U'=Y'-\{a'\}=u^{-1}(U)$. Si $h':Z'\to Y'$ est le morphisme canonique, le
fait que la restriction $f^{-1}(U)\to h^{-1}(U)$ de $g$ soit une immersion
ouverte entraîne qu'il en est de même de la restriction
$f'^{-1}(U')\to h'^{-1}(U')$ de $g'$. Notons maintenant que $f'$ est un
morphisme séparé et quasi-fini (II, 6.2.4)~; comme $A'$ est \emph{complet},
on déduit de (II, 6.2.6) que $X'$ est $Y'$-isomorphe à une somme
$X'_1\amalg X'_2$, où la restriction
$f'|X'_1=f'_1:X'_1\to Y'$ est un morphisme \emph{fini}, et
$X'_2\subset f'^{-1}(U')$. Il en résulte que $B'$ est composée directe des
deux $A'$-algèbres
$\Gamma(X'_1,\mathcal O_{X'_1})=B'_1$ et
$\Gamma(X'_2,\mathcal O_{X'_2})=B'_2$~; on en conclut aussitôt que la
fermeture intégrale $C'$ de $A'$ dans $B'$ est composée directe des
fermetures intégrales $C'_1$, $C'_2$ de $A'$ dans $B'_1$, $B'_2$
respectivement, d'où $Z'=Z'_1\amalg Z'_2$, où
$Z'_i=\operatorname{Spec}(C'_i)$ ($i=1,2$)~; et le morphisme canonique
$g':X'\to Z'$ est tel que $g'|X'_i$ soit le morphisme canonique
$g'_i:X'_i\to Z'_i$ ($i=1,2$). Mais comme $B'_1$ est déjà une
$A'$-algèbre finie, on a $C'_1=B'_1$, et $g'_1$ est donc un isomorphisme.
D'autre part, comme $X'_2\subset f'^{-1}(U')$ et est ouvert dans
$f'^{-1}(U')$, on sait
\oldpage[IV-3]{47}
déjà que $g'_2$ est une immersion ouverte. On conclut bien que $g'$ est une
immersion ouverte.
\end{proof}

\begin{remark}[8.12.7]
\label{IV.8.12.7-fr}
Lorsque, dans (8.12.6), on suppose que $Y$ est un \emph{schéma affine}, la
démonstration par réduction au cas noethérien montre que, dans la
factorisation (8.12.6.1), les morphismes $f'$ et $u$ sont aussi des
morphismes \emph{de présentation finie} (1.6.2).
\end{remark}

\begin{corollary}[8.12.8]
\label{IV.8.12.8-fr}
Soient $Y$ un schéma quasi-compact tel qu'il existe un $\mathcal O_Y$-Module
ample (II, 4.5.3), $f:X\to Y$ un morphisme quasi-fini et quasi-projectif.
Alors il existe une factorisation de $f$ en
\[
X\xrightarrow{f'}Y'\xrightarrow{u}Y
\]
où $f'$ est une immersion ouverte et $u$ un morphisme fini.
\end{corollary}

\begin{proof}
L'hypothèse entraîne que $X$ s'identifie à un sous-$Y$-schéma quasi-compact
d'un $Y$-schéma de la forme $Z=\mathbf P_Y^r$ (II, 5.3.3). Il y a par suite
un voisinage ouvert quasi-compact $U$ de $X$ dans $Z$ tel que $X$ soit fermé
dans $U$~; comme $Z$ est un schéma, l'injection canonique $U\to Z$ est un
morphisme de présentation finie ((1.2.7) et (1.6.2)), donc le morphisme
composé $g:U\to Z\to Y$ est aussi un morphisme de présentation finie (le
fait que $\mathbf P_Y^r$ est de présentation finie sur $Y$ résultant
aussitôt de la définition (II, 4.1.1)). Soit $\mathcal J$ l'Idéal
quasi-cohérent de $\mathcal O_U$ définissant le sous-préschéma fermé $X$~;
comme $U$ est un schéma quasi-compact, $\mathcal J$ est limite inductive
filtrante de ses sous-Idéaux quasi-cohérents de type fini $\mathcal J_\lambda$
(I, 9.4.9). Si $X_\lambda$ est le sous-préschéma fermé de $U$ défini par
$\mathcal J_\lambda$, on a par suite $X=\bigcap_\lambda X_\lambda$. Pour tout
$y\in Y$, on a donc
$f^{-1}(y)=\bigcap_\lambda(X_\lambda\cap g^{-1}(y))$, et comme les ensembles
$X_\lambda\cap g^{-1}(y)$ sont fermés dans l'espace noethérien
$g^{-1}(y)$, il existe pour tout $y$ un indice $\lambda(y)$ tel que
$f^{-1}(y)=X_{\lambda(y)}\cap g^{-1}(y)$. Désignons par $E_\lambda$
l'ensemble des $y\in Y$ tels que la fibre $X_\lambda\cap g^{-1}(y)$ de la
restriction de $g$ à $X_\lambda$ soit un $k(y)$-préschéma fini. L'hypothèse
que $f$ est quasi-fini entraîne, en vertu de ce qui précède, que
$Y=\bigcup_\lambda E_\lambda$. Or, chacun des $X_\lambda$ est, par
définition, \emph{de présentation finie} sur $Y$~; il résulte donc de
(9.2.3) et (9.2.6)\footnote{Le lecteur vérifiera que les corollaires (8.12.8)
à (8.12.11) ne sont pas utilisés dans le \S~9.} que les $E_\lambda$ sont des
ensembles \emph{constructibles} dans le schéma $Y$~; comme ils forment un
ensemble filtrant croissant, il existe un indice $\lambda$ tel que
$E_\lambda=Y$ (1.9.9), et pour cet indice $\lambda$, le morphisme
$f_\lambda:X_\lambda\to Y$, restriction de $g$ à $X_\lambda$, est donc
quasi-fini. Comme il est de présentation finie et séparé, on peut lui
appliquer (8.12.6), et $f_\lambda$ se factorise donc en
\[
X_\lambda\xrightarrow{j_\lambda}Y'_\lambda
\xrightarrow{u_\lambda}Y
\]
où $j_\lambda$ est une immersion et $u_\lambda$ un morphisme fini. Comme
$X$ est un sous-préschéma fermé de $X_\lambda$, on a ainsi prouvé que $f$
possède la propriété (8.12.3, \emph{a)}), d'où le corollaire en vertu de
l'équivalence de (8.12.3, \emph{a)}) et (8.12.3, \emph{a$'$)}).
\end{proof}

\oldpage[IV-3]{48}
\begin{corollary}[8.12.9]
\label{IV.8.12.9-fr}
Soit $f:X\to Y$ un morphisme localement quasi-fini
$(\mathbf{Err}_{\mathrm{III}}, 20)$. Pour tout $x\in X$ il existe un
voisinage ouvert $U$ de $x$ dans $X$, un voisinage ouvert $V$ de
$y=f(x)$ dans $Y$, tels que $f(U)\subset V$ et une factorisation
\[
U\xrightarrow{f'}V'\xrightarrow{u}V
\]
de la restriction de $f$ à $U$, où $f'$ est une immersion ouverte et $u$ un
morphisme fini.
\end{corollary}

Il suffit évidemment de prendre pour $V$ un voisinage affine de $y$ dans
$Y$, pour $U$ un voisinage affine de $x$ dans $X$ contenu dans $f^{-1}(U)$
et tel que $f|U$ soit quasi-fini. Le morphisme $U\to V$ restriction de $f$
étant alors affine (donc quasi-projectif), on peut lui appliquer (8.12.8).

\begin{corollary}[8.12.10]
\EGAReaderTarget{ega4.c04.cor.8-12-10}{IV.8.12.10-fr}
Soient $Y$ un préschéma intègre et normal, $X$ un préschéma intègre,
$f:X\to Y$ un morphisme birationnel et localement quasi-fini
$(\mathbf{Err}_{\mathrm{III}}, 20)$. Alors $f$ est un isomorphisme local~;
pour que $f$ soit une immersion ouverte, il faut et il suffit que $f$ soit
en outre séparé.
\end{corollary}

La seconde assertion résulte aussitôt de la première et de (I, 8.2.8). Pour
prouver la première assertion, on peut supposer $X$ et $Y$ affines et $f$
quasi-fini~; considérons la factorisation $f=u\circ f'$ de (8.12.8), qui
permet d'identifier $X$ par $f'$ à un sous-préschéma induit sur un ouvert de
$Y'$. Comme $X$ est intègre, on peut, en vertu de (I, 5.2.3), remplacer
$Y'$ par le sous-préschéma réduit de $Y'$ ayant pour espace sous-jacent
$\overline X$, donc on peut supposer aussi que $Y'$ est \emph{intègre}. En
outre, puisque $f$ est birationnel, $u$ l'est aussi. La conclusion résulte
alors du lemme suivant~:

\begin{lemma}[8.12.10.1]
\label{IV.8.12.10.1-fr}
Soient $Y'$ un préschéma intègre, $Y$ un préschéma intègre et normal~; alors
un morphisme $u:Y'\to Y$ fini et birationnel est un isomorphisme.
\end{lemma}

Posons en effet $\mathcal A=u_*(\mathcal O_{Y'})$, de sorte que $\mathcal A$
est une $\mathcal O_Y$-Algèbre finie, $Y'$ s'identifiant à
$\operatorname{Spec}(\mathcal A)$ (II, 1.3.6). Si $R(Y)$ est le corps des
fonctions rationnelles de $Y$, on a donc, pour tout $y\in Y$,
$\mathcal O_{Y,y}\subset\mathcal A_y\subset R(Y)$~; mais comme l'anneau
$\mathcal O_{Y,y}$ est par hypothèse intégralement clos et a $R(Y)$ pour
corps des fractions, on a nécessairement
$\mathcal A_y=\mathcal O_{Y,y}$, d'où le lemme.

\begin{corollary}[8.12.11]
\EGAReaderTarget{ega4.c04.cor.8-12-11}{IV.8.12.11-fr}
Soient $Y$ un préschéma intègre, $X$ un préschéma intègre et normal,
$f:X\to Y$ un morphisme dominant et localement quasi-fini. Soient $K$ et
$L$ (extension de $K$) les corps de fonctions rationnelles de $Y$ et $X$
respectivement~; et soit $Y'$ la fermeture intégrale de $Y$ relativement à
$L$ (II, 6.3.4)~; alors $f$ se factorise d'une seule manière en
$f:X\xrightarrow{f'}Y'\xrightarrow{u}Y$, où $f'$ est birationnel, et
correspond à l'automorphisme identique de $L$~; $f'$ est alors un
isomorphisme local, et pour que $f'$ soit une immersion ouverte, il faut et
il suffit que $f$ soit séparé.
\end{corollary}

L'existence et l'unicité de la factorisation de $f$ résultent de (II, 6.3.9).
Il résulte de (II, 6.2.4, (v)), en se ramenant au cas affine, que $f'$ est
localement quasi-fini~; en outre, il résulte de (I, 5.5.1) que, pour que
$f'$ soit séparé, il faut et il suffit que $f$ le soit, puisque $u$ est
affine, donc séparé~; les deux dernières assertions sont donc des
conséquences de (8.12.10) appliqué à $f'$.

\oldpage[IV-3]{49}
\subsection{Traduction en termes de pro-objets}
\label{subsection:IV.8.13-fr}

La proposition suivante est essentiellement équivalente à (8.8.2, (i))~:

\begin{proposition}[8.13.1]
\label{IV.8.13.1-fr}
Soient $S$ un préschéma, $(X_\lambda,v_{\lambda\mu})$ un système projectif
filtrant de $S$-préschémas~; on suppose qu'il existe $\alpha$ tel que
$v_{\alpha\lambda}$ soit un morphisme affine pour tout $\lambda\geqslant
\alpha$ (ce qui entraîne (II, 1.6.2) que $v_{\lambda\mu}$ soit affine pour
$\alpha\leqslant\lambda\leqslant\mu$), de sorte que la limite projective
$X=\varprojlim_\lambda X_\lambda$ existe dans la catégorie des $S$-préschémas
(8.2.3). Soit $Y$ un $S$-préschéma, et, pour tout $\lambda\geqslant\alpha$,
soit
\[
e_\lambda:\operatorname{Hom}_S(X_\lambda,Y)
\longrightarrow\operatorname{Hom}_S(X,Y)
\]
l'application qui, à tout $S$-morphisme $f_\lambda:X_\lambda\to Y$, fait
correspondre $f=f_\lambda\circ v_\lambda$, où $v_\lambda:X\to X_\lambda$
est le morphisme canonique. La famille $(e_\lambda)$ est un système inductif
d'applications, qui définit donc une application canonique
\begin{equation}
\label{IV.8.13.1.1-fr}
\varinjlim_\lambda\operatorname{Hom}_S(X_\lambda,Y)
\longrightarrow\operatorname{Hom}_S(X,Y). \tag{8.13.1.1}
\end{equation}
Supposons $X_\alpha$ quasi-compact (resp. quasi-compact et quasi-séparé), et le
morphisme structural $Y\to S$ localement de type fini (resp. localement de
présentation finie). Alors l'application (8.13.1.1) est injective (resp.
bijective).
\end{proposition}

Posons en effet, pour $\lambda\geqslant\alpha$,
$Z_\lambda=Y\times_S X_\lambda$, de sorte que l'on a
$Z_\lambda=Z_\alpha\times_{X_\alpha}X_\lambda$. Posons de même
$Z=Y\times_S X=Z_\alpha\times_{X_\alpha}X$~; on sait alors (8.2.5) que, si
l'on pose aussi
\[
w_{\lambda\mu}=1\times v_{\lambda\mu}:Z_\mu\to Z_\lambda
\quad(\alpha\leqslant\lambda\leqslant\mu),
\qquad
w_\lambda=1\times v_\lambda:Z\to Z_\lambda
\quad(\alpha\leqslant\lambda),
\]
$Z$ est limite projective du système projectif $(Z_\lambda,w_{\lambda\mu})$
et les $w_\lambda$ sont les morphismes canoniques correspondants. Notons
d'autre part que le morphisme $Z_\alpha\to X_\alpha$ est localement de type
fini (resp. localement de présentation finie) (1.3.4 et 1.4.3). Enfin, on
sait que l'on a
\[
\operatorname{Hom}_S(X_\lambda,Y)
=\operatorname{Hom}_{X_\lambda}(X_\lambda,Z_\lambda)
\quad\hbox{et}\quad
\operatorname{Hom}_S(X,Y)=\operatorname{Hom}_X(X,Z)
\]
(I, 3.3.14). Il suffit maintenant d'appliquer (8.8.2, (i)) en y faisant
$X_\lambda=S_\lambda$ et en y remplaçant $Y_\lambda$ par $Z_\lambda$.

\begin{corollary}[8.13.2]
\label{IV.8.13.2-fr}
Avec les notations de (8.13.1), on suppose $X_\alpha$ quasi-compact et
quasi-séparé, et les $v_{\alpha\lambda}$ affines pour
$\alpha\leqslant\lambda$~; on suppose en outre que
$Y=\varprojlim_\rho Y_\rho$, où $(Y_\rho,t_{\rho\sigma})$ est un système
projectif filtrant de $S$-préschémas tel que, pour chaque $\rho$, le morphisme
structural $Y_\rho\to S$ soit localement de présentation finie. On a alors une
bijection canonique
\begin{equation}
\label{IV.8.13.2.1-fr}
\operatorname{Hom}_S(X,Y)
\xrightarrow{\sim}
\varprojlim_\rho\left(\varinjlim_\lambda
\operatorname{Hom}_S(X_\lambda,Y_\rho)\right). \tag{8.13.2.1}
\end{equation}
\end{corollary}

En effet, le fait que $Y$ soit limite projective des $Y_\rho$ entraîne en
particulier que l'application canonique
\[
\operatorname{Hom}_S(X,Y)\longrightarrow
\varprojlim_\rho\operatorname{Hom}_S(X,Y_\rho)
\]
est bijective~; et d'autre part, les hypothèses entraînent, pour chaque
$\rho$, l'existence d'une bijection canonique
\[
\operatorname{Hom}_S(X,Y_\rho)\xrightarrow{\sim}
\varinjlim_\lambda\operatorname{Hom}_S(S_\lambda,Y_\rho)
\]
en vertu de (8.13.1)~; d'où la conclusion.

\begin{env}[8.13.3]
\label{IV.8.13.3-fr}
Les résultats précédents permettent d'interpréter dans la théorie des
préschémas les notions de «~pro-variété~» ou de «~pro-schéma~» qui
s'introduisent dans certaines applications (par exemple dans la théorie du
corps de classes local suivant les idées de Serre [39] ou dans la théorie de
Néron sur la réduction des variétés
\oldpage[IV-3]{50}
abéliennes [32]). Rappelons rapidement ici la notion de \emph{pro-objet} d'une
catégorie, renvoyant au chap.~V pour de plus amples développements (nous
n'utiliserons d'ailleurs pas avant le chap.~V l'interprétation qui va suivre,
et le lecteur peut donc omettre jusque-là la lecture de la fin de ce numéro).
Étant donnée une catégorie $\mathcal C$, la catégorie
$\mathbf{Pro}(\mathcal C)$ des \emph{pro-objets} de $\mathcal C$ a pour objets
les \emph{systèmes projectifs} (dans l'univers où l'on se place)
$\mathbf X=(X_\mu)_{\mu\in M}$ d'objets de $\mathcal C$ dont les ensembles
d'indices (dépendant du système projectif considéré) sont supposés préordonnés
filtrants. Étant donnés deux tels pro-objets
$\mathbf X=(X_\mu)_{\mu\in M}$,
$\mathbf X'=(X'_{\mu'})_{\mu'\in M'}$, les \emph{morphismes} de $\mathbf X$
dans $\mathbf X'$ sont par définition les éléments de l'ensemble
\[
\varprojlim_{\mu'}\left(\varinjlim_\mu
\operatorname{Hom}(X_\mu,X'_{\mu'})\right)~;
\]
la vérification du fait que l'on peut prendre ces ensembles pour ensembles de
morphismes est immédiate, la composition de systèmes de morphismes
$u_{\mu'}^\mu:X_\mu\to X'_{\mu'}$,
$u_{\mu''}^{\prime\mu'}:X'_{\mu'}\to X''_{\mu''}$ qui sont inductifs en
l'indice supérieur et projectifs en l'indice inférieur, se faisant
«~argument par argument~», autrement dit en considérant le système des
$u_{\mu''}^{\prime\mu'}\circ u_{\mu'}^\mu$.
\end{env}

\begin{env}[8.13.4]
\label{IV.8.13.4-fr}
Considérons alors un préschéma quasi-compact et quasi-séparé $S$, et désignons
par $\mathcal C$ la sous-catégorie pleine de la catégorie
$(\mathbf{Sch})_{/S}$ des $S$-préschémas, formée des $S$-préschémas $X$ ayant
la propriété suivante~: le morphisme structural $X\to S$ se factorise en
\[
X\xrightarrow{g}X_0\xrightarrow{f}S,
\]
où $g:X\to X_0$ est affine et $f:X_0\to S$ de présentation finie~; nous
dirons pour abréger que les préschémas de $\mathcal C$ sont
\emph{essentiellement affines au-dessus de $S$}.

Considérons d'autre part la sous-catégorie pleine $\mathcal C'_0$ de
$(\mathbf{Sch})_{/S}$ formée des $S$-préschémas \emph{de présentation finie},
et la catégorie $\mathbf{Pro}(\mathcal C'_0)$ des pro-objets de
$\mathcal C'_0$. Nous dirons qu'un objet
$\mathbf X=(X_\mu)_{\mu\in M}$ de $\mathbf{Pro}(\mathcal C'_0)$ est
\emph{essentiellement affine} s'il existe un $\gamma\in M$ tel que pour tout
$\mu\geqslant\gamma$, le morphisme de transition
$v_{\gamma\mu}:X_\mu\to X_\gamma$ soit affine (ce qui entraîne que pour
$\gamma\leqslant\mu\leqslant\nu$, $v_{\mu\nu}:X_\nu\to X_\mu$ est affine).
On notera qu'un objet de $\mathbf{Pro}(\mathcal C'_0)$ \emph{isomorphe} à un
objet essentiellement affine \emph{n'est pas} nécessairement essentiellement
affine lui-même.

Nous désignerons par $\mathcal C'$ la sous-catégorie pleine de
$\mathbf{Pro}(\mathcal C'_0)$ formée des pro-objets de $\mathcal C'_0$
essentiellement affines.

Cela étant, il résulte de (8.2.2) et (8.2.3) que pour tout objet
$\mathbf X=(X_\mu)_{\mu\in M}$ de $\mathcal C'$, le $S$-préschéma
$X=\varprojlim_\mu X_\mu$ existe~; en outre, comme, pour $\mu$ assez grand,
le morphisme $X\to X_\mu$ est affine (8.2.2), $X$ est
\emph{essentiellement affine au-dessus de $S$} par définition. Posons
$X=L(\mathbf X)$~; montrons qu'on a ainsi défini un \emph{foncteur canonique}
\begin{equation}
\label{IV.8.13.4.1-fr}
L:\mathcal C'\longrightarrow\mathcal C. \tag{8.13.4.1}
\end{equation}

On a en effet, pour deux objets $\mathbf X=(X_\mu)$,
$\mathbf X'=(X'_{\mu'})$ de $\mathcal C'$, une application canonique pour
chaque $\mu'$
\[
\varinjlim_\mu\operatorname{Hom}_S(X_\mu,X'_{\mu'})
\longrightarrow
\operatorname{Hom}_S(\varprojlim_\mu X_\mu,X'_{\mu'})
\]
définie dans (8.13.1.1), et d'autre part, par définition de la limite
projective, une bijection canonique
\[
\varprojlim_{\mu'}
\operatorname{Hom}_S(\varprojlim_\mu X_\mu,X'_{\mu'})
\xrightarrow{\sim}
\operatorname{Hom}_S(\varprojlim_\mu X_\mu,
\varprojlim_{\mu'}X'_{\mu'}),
\]
\oldpage[IV-3]{51}
d'où une application canonique
\begin{equation}
\label{IV.8.13.4.2-fr}
\varprojlim_{\mu'}\left(\varinjlim_\mu
\operatorname{Hom}_S(X_\mu,X'_{\mu'})\right)
\longrightarrow
\operatorname{Hom}_S(\varprojlim_\mu X_\mu,
\varprojlim_{\mu'}X'_{\mu'}), \tag{8.13.4.2}
\end{equation}
évidemment fonctorielle en $\mathbf X$ et $\mathbf X'$, et qui complète la
définition du foncteur $L$.
\end{env}

\begin{proposition}[8.13.5]
\label{IV.8.13.5-fr}
Les hypothèses et notations étant celles de (8.13.4), le foncteur $L$ est
pleinement fidèle. Si de plus $S$ est un préschéma noethérien (ce qui implique
déjà que $S$ est quasi-compact et quasi-séparé (1.2.8)), $L$ est une
\emph{équivalence de catégories}.
\end{proposition}

Dire que $L$ est pleinement fidèle signifie que l'application (8.13.4.2) est
bijective quels que soient $\mathbf X$, $\mathbf X'$ dans $\mathcal C'$, ce
qui est un cas particulier de (8.13.2)~: en effet, les morphismes structuraux
$X_\mu\to S$ étant de présentation finie, sont en particulier quasi-compacts
et quasi-séparés, donc les $X_\mu$ sont quasi-compacts et quasi-séparés.

Pour montrer que lorsque $S$ est noethérien $L$ est une équivalence de
catégories, il suffit, puisque l'on sait déjà que $L$ est pleinement fidèle,
de prouver que tout $S$-préschéma essentiellement affine $X$ est
$S$-\emph{isomorphe} à un objet de la forme $L(\mathbf X)$ où
$\mathbf X\in\mathcal C'$ $(\mathbf 0_{\mathrm{III}},8.1.5)$.

Or, il y a par hypothèse une factorisation
$X\xrightarrow{g}X_0\xrightarrow{f}S$ du morphisme structural, $f$ étant de
présentation finie et $g$ affine. On peut donc écrire
$X=\operatorname{Spec}(\mathcal A)$, où $\mathcal A$ est une
$\mathcal O_{X_0}$-Algèbre quasi-cohérente (II, 1.3.1). Or, comme $X_0$ est
noethérien (puisqu'il en est ainsi de $S$ et que $f$ est de type fini),
$\mathcal A$ est limite inductive filtrante de la famille
$(\mathcal A_\lambda)$ de ses sous-$\mathcal O_{X_0}$-Algèbres quasi-cohérentes
de type fini (I, 9.6.6). Posons
$X_\lambda=\operatorname{Spec}(\mathcal A_\lambda)$~; les morphismes
$X_\lambda\to X_0$ sont de type fini, donc de présentation finie puisque
$X_0$ est noethérien, et par suite il en est de même des morphismes composés
$X_\lambda\to X_0\to S$ (1.6.2)~; autrement dit, les $X_\lambda$
appartiennent à $\mathcal C'_0$ et comme les morphismes $X_\lambda\to X_0$
sont affines, $\mathbf X=(X_\lambda)$ est un objet de $\mathcal C'$ dont la
limite projective existe et est $S$-isomorphe à $X$ en vertu de (8.2.2). Ceci
achève la démonstration.

\begin{remark}[8.13.6]
\label{IV.8.13.6-fr}
Il résulte de (1.6.2) et de (II, 1.6.2) que si $X$ et $Y$ sont
essentiellement affines au-dessus de $S$, alors il en est de même de
$X\times_S Y$. On en conclut par exemple $(\mathbf 0_{\mathrm{III}},8.2.5)$
qu'un $\mathcal C$-groupe n'est autre qu'un $(\mathbf{Sch})_{/S}$-groupe qui
est un préschéma essentiellement affine au-dessus de $S$. D'autre part, les
produits finis existent dans la catégorie $\mathcal C'$~: en effet, si
$\mathbf X=(X_\mu)_{\mu\in M}$, $\mathbf Y=(Y_\rho)_{\rho\in R}$ sont deux
objets de $\mathcal C'$, les produits $X_\mu\times_S Y_\rho$ sont des
$S$-préschémas de présentation finie, et en prenant pour morphismes de
transition
\[
X_\nu\times_S Y_\sigma\longrightarrow X_\mu\times_S Y_\rho
\]
les produits des morphismes de transition $X_\nu\to X_\mu$ et
$Y_\sigma\to Y_\rho$, on voit aussitôt que
$(X_\mu\times_S Y_\rho)$ est produit de $\mathbf X$ et de $\mathbf Y$ dans
$\mathbf{Pro}(\mathcal C'_0)$~; en outre (II, 1.6.2) les morphismes de
transition ainsi définis sont affines pour $\mu$ et $\rho$ assez grands, donc
le produit $\mathbf X\times\mathbf Y$ ainsi défini appartient bien à
$\mathcal C'$. On conclut alors comme plus haut qu'un $\mathcal C'$-groupe est
un $\mathbf{Pro}(\mathcal C'_0)$-groupe qui est essentiellement affine. On
déduit donc de (8.13.5) que les catégories des $\mathcal C'$-groupes et des
$\mathcal C$-groupes sont \emph{équivalentes} lorsque $S$ est
\emph{noethérien}. Il semble plausible que lorsque $S$ est le spectre d'un
corps $k$, la catégorie des $\mathcal C$-groupes soit \emph{équivalente} à
celle des $\boldsymbol K$-groupes, où $\boldsymbol K$ est la catégorie des
$S$-préschémas \emph{quasi-compacts}~; autrement dit, tout préschéma en
groupes sur $k$ qui est quasi-compact serait essentiellement affine. D'autre
part, si l'on désigne
\oldpage[IV-3]{52}
par $\mathcal C'_0$-$\boldsymbol{gr}$ la catégorie des
$\mathcal C'_0$-groupes, il est vraisemblable que la catégorie des
$\mathcal C'$-groupes est équivalente à la sous-catégorie pleine de
$\mathbf{Pro}(\mathcal C'_0\text{-}\boldsymbol{gr})$ formée des «~groupes
proalgébriques essentiellement affines~», c'est-à-dire des systèmes
projectifs $\mathbf G=(G_\mu)_{\mu\in M}$, où les $G_\mu$ sont des groupes
algébriques sur $k$ et les morphismes de transition $G_\nu\to G_\mu$ sont
affines pour $\mu$ assez grand (ce que l'on peut exprimer aussi en disant que
$\mathbf G$ est extension d'un groupe algébrique par un pro-groupe algébrique
affine). La conjonction de ces deux conjectures équivaut d'ailleurs à la
suivante~: tout préschéma en groupes quasi-compacts sur $k$ est «~extension~»
d'un «~groupe algébrique~» (i.e. un préschéma en groupes de type fini sur $k$)
par un préschéma en groupes \emph{affine} sur $k$.

Les seuls pro-groupes algébriques rencontrés en pratique jusqu'à présent étant
en fait essentiellement affines, il y aura donc sans doute avantage à
substituer à l'étude des groupes pro-algébriques généraux (introduits et
étudiés par Serre [40]) celle des schémas en groupes quasi-compacts sur $k$,
dont la définition est conceptuellement plus simple.
\end{remark}

\subsection{Caractérisation d'un préschéma localement de présentation finie sur un autre, en termes du foncteur qu'il représente}
\label{subsection:IV.8.14-fr}

\begin{env}[8.14.1]
\label{IV.8.14.1-fr}
Étant donné un préschéma $S$, nous dirons encore, comme en (8.13.4), qu'un
système projectif filtrant $(X_\lambda,v_{\lambda\mu})$ de $S$-préschémas est
\emph{essentiellement affine} s'il existe $\alpha$ tel que
$v_{\alpha\lambda}$ soit un morphisme affine pour $\lambda\geqslant\alpha$.

L'énoncé suivant, qui sera surtout utile au chap.~V, précise (8.8.2, (i)) en
lui fournissant une réciproque~:
\end{env}

\begin{proposition}[8.14.2]
\label{IV.8.14.2-fr}
Soient $S$ un préschéma, $f:X\to S$ un morphisme. Pour tout $S$-préschéma
$T$, posons
\[
h_X(T)=\operatorname{Hom}_S(T,X)
\]
de sorte que $h_X$ est un foncteur contravariant de la catégorie
$(\mathbf{Sch})_{/S}$ des $S$-préschémas dans la catégorie
$\mathbf{Ens}$ des ensembles $(\mathbf 0_{\mathrm{III}},8.1.1)$, et $X$ un
objet représentant ce foncteur $(\mathbf 0_{\mathrm{III}},8.1.8)$. Les
conditions suivantes sont équivalentes~:
\begin{enumerate}[label=\emph{\alph*)}]
  \item $f$ est localement de présentation finie.
  \item Pour tout système projectif filtrant $(Z_\lambda)$ de
  $S$-préschémas, essentiellement affine (8.13.4) et formé de préschémas
  quasi-compacts et quasi-séparés, l'application canonique (8.13.1.1)
\begin{equation}
\label{IV.8.14.2.1-fr}
\varinjlim_\lambda h_X(Z_\lambda)
\longrightarrow h_X(\varprojlim_\lambda Z_\lambda) \tag{8.14.2.1}
\end{equation}
  est bijective.
  \item Pour tout système projectif filtrant $(Z_\lambda)$ de
  $S$-préschémas, tel que les $Z_\lambda$ soient des schémas affines,
  l'application (8.14.2.1) est bijective.
  \item[\rm c$'$)] Pour tout ouvert affine $U$ de $S$ et tout système
  projectif filtrant $(Z_\lambda)$ de $U$-préschémas, tel que les $Z_\lambda$
  soient des schémas affines, l'application (8.14.2.1) est bijective.
\end{enumerate}
\end{proposition}

Le fait que \emph{a)} implique \emph{b)} n'est autre que (8.13.1)~; il est
trivial que \emph{b)} implique \emph{c)} et que \emph{c)} implique
\emph{c$'$)}. Reste à voir que \emph{c$'$)} entraîne \emph{a)}, et comme la
propriété \emph{a)} est locale sur $S$, on peut se borner au cas où $S$ est
\emph{affine}.

\oldpage[IV-3]{53}
Supposons d'abord que $X$ soit aussi affine~; l'assertion à prouver est alors
équivalente au

\begin{corollary}[8.14.2.2]
\label{IV.8.14.2.2-fr}
Soient $A$ un anneau, $B$ une $A$-algèbre. Afin que, pour tout système
inductif filtrant $(C_\lambda)$ de $A$-algèbres, l'application canonique
\begin{equation}
\label{IV.8.14.2.3-fr}
\varinjlim_\lambda\operatorname{Hom}_{A\text{-}\mathrm{alg.}}(B,C_\lambda)
\longrightarrow
\operatorname{Hom}_{A\text{-}\mathrm{alg.}}
(B,\varinjlim_\lambda C_\lambda) \tag{8.14.2.3}
\end{equation}
soit bijective, il faut et il suffit que $B$ soit une $A$-algèbre de
présentation finie.
\end{corollary}

Il reste seulement à montrer que la condition est nécessaire. Prenons d'abord
pour $(C_\lambda)$ le système inductif filtrant des sous-$A$-algèbres de type
fini de $B$, de sorte que $\varinjlim_\lambda C_\lambda=B$. Le fait que (8.14.2.3)
soit bijective entraîne en particulier que l'application identique
$1_B$ se factorise en $B\to C_\lambda\to B$ pour un $\lambda$ convenable,
ce qui entraîne $C_\lambda=B$, donc $B$ est une $A$-algèbre \emph{de type
fini}. Posons alors $B=C/\mathfrak J$, où
$C=A[T_1,\ldots,T_n]$ et $\mathfrak J$ est un idéal de $C$. Alors
$\mathfrak J$ est limite inductive filtrante des idéaux de type fini
$\mathfrak J_\lambda\subset\mathfrak J$ de $C$~; posant
$C_\lambda=C/\mathfrak J_\lambda$, et utilisant l'exactitude du foncteur
$\varinjlim$, on voit que $B$ est encore isomorphe à la limite inductive du
système inductif filtrant $(C_\lambda)$. Il existe donc un $\lambda$ et un
$A$-homomorphisme $u:B\to C_\lambda$ tels que le composé
\[
B\xrightarrow{u}C_\lambda\xrightarrow{p_\lambda}B
\]
(où $p_\lambda$ est l'homomorphisme canonique) soit l'identité. Soit
$q_\lambda:C\to C_\lambda$ l'homomorphisme canonique, et posons
$t_i=p_\lambda(q_\lambda(T_i))$~; on a donc
$p_\lambda(u(t_i))=p_\lambda(q_\lambda(T_i))$, autrement dit
\[
u(t_i)-q_\lambda(T_i)\in\mathfrak J/\mathfrak J_\lambda.
\]
Il existe donc un $\mu\geqslant\lambda$ tel que les $n$ éléments
$u(t_i)-q_\lambda(T_i)$ appartiennent à
$\mathfrak J_\mu/\mathfrak J_\lambda$ $(1\leqslant i\leqslant n)$~; si
$p_{\mu\lambda}:C_\lambda\to C_\mu$ est l'homomorphisme canonique, on a par
suite
\[
p_{\mu\lambda}(u(t_i))
=p_{\mu\lambda}(q_\lambda(T_i))=q_\mu(T_i).
\]
Remplaçant $\lambda$ par $\mu$ et $u$ par $p_{\mu\lambda}\circ u$, on peut
donc supposer que $u(t_i)=q_\lambda(T_i)$ pour tout $i$, et si
$r=p_\lambda\circ q_\lambda$ est l'homomorphisme canonique
$C\to C/\mathfrak J=B$, on peut donc écrire
$u(r(T_i))=q_\lambda(T_i)$ pour tout $i$, d'où
$q_\lambda=u\circ r$. Mais cela entraîne nécessairement que
$\mathfrak J=\mathfrak J_\lambda$, car si $z\in\mathfrak J$, on a
$r(z)=0$~; donc on a $B=C_\lambda$.

Passons maintenant au cas où $S$ est affine et $X$ quelconque~; tout revient à
prouver qu'un ouvert affine $V$ de $X$ est de présentation finie sur $S$, et
en vertu de ce qui vient d'être démontré, il suffit de prouver que pour tout
système projectif filtrant $(Z_\lambda)$ de $S$-préschémas affines,
l'application
\begin{equation}
\label{IV.8.14.2.4-fr}
\varinjlim_\lambda\operatorname{Hom}_S(Z_\lambda,V)
\longrightarrow
\operatorname{Hom}_S(\varprojlim_\lambda Z_\lambda,V) \tag{8.14.2.4}
\end{equation}
est bijective. Il est immédiat que cette application est injective, car si
$(v_\lambda)$, $(v'_\lambda)$ sont deux systèmes inductifs de
$S$-homomorphismes $v_\lambda:Z_\lambda\to V$,
$v'_\lambda:Z_\lambda\to V$ tels que les morphismes correspondants
\[
Z\xrightarrow{u_\lambda}Z_\lambda\xrightarrow{v_\lambda}V,
\qquad
Z\xrightarrow{u_\lambda}Z_\lambda\xrightarrow{v'_\lambda}V
\]
soient égaux ($u_\lambda$ étant le morphisme canonique), alors les morphismes
\[
Z\xrightarrow{u_\lambda}Z_\lambda\xrightarrow{v_\lambda}V
\xrightarrow{j}X,
\qquad
Z\xrightarrow{u_\lambda}Z_\lambda\xrightarrow{v'_\lambda}V
\xrightarrow{j}X
\]
(où $j$ est l'injection canonique) sont égaux, ce qui entraîne
$j\circ v_\lambda=j\circ v'_\lambda$ par hypothèse pour un $\lambda$
convenable, donc $v_\lambda=v'_\lambda$.

\oldpage[IV-3]{54}
Reste à prouver que (8.14.2.4) est surjective. Soit donc $v:Z\to V$ un
$S$-morphisme~; par hypothèse il existe un $\lambda$ et un $S$-morphisme
$w_\lambda:Z_\lambda\to X$ tels que $j\circ v$ se factorise en
\[
Z\xrightarrow{u_\lambda}Z_\lambda\xrightarrow{w_\lambda}X,
\]
et tout revient à prouver qu'il existe $\mu\geqslant\lambda$ tel que le
morphisme
\[
Z_\mu\xrightarrow{u_{\lambda\mu}}Z_\lambda
\xrightarrow{w_\lambda}X
\]
(où $u_{\lambda\mu}$ est le morphisme de transition) se factorise en
\[
Z_\mu\xrightarrow{v_\mu}V\xrightarrow{j}X.
\]
Posons, pour tout $\mu\geqslant\lambda$,
$U_\mu=u_{\lambda\mu}^{-1}(w_\lambda^{-1}(V))$. On a
\[
u_\mu^{-1}(U_\mu)=u_\lambda^{-1}(U_\lambda)
=u_\lambda^{-1}(w_\lambda^{-1}(V))=v^{-1}(V)=Z.
\]
Comme les $Z_\lambda$ sont quasi-compacts et que les $U_\mu$, étant ouverts,
sont ind-constructibles (1.9.6), on déduit de (8.3.4) qu'il existe
$\mu\geqslant\lambda$ tel que $U_\mu=Z_\mu$. C.Q.F.D.

\begin{remark}[8.14.3]
\label{IV.8.14.3-fr}
Le fait que l'application (8.14.2.1) soit injective lorsque $f$ est localement
de type fini (8.8.2, (i)) amène naturellement à se demander si ce résultat
admet aussi une réciproque. Il n'en est rien, même lorsque $S$ et $X$ sont
affines, car il existe des monomorphismes $X\to S$ qui ne sont pas de type fini
(I, 2.4.2), et qui mettent donc en défaut cette conjecture.
\end{remark}

% FIN SOURCE ega4/ega4-8-fr.tex

%
% DEBUT SOURCE ega4/ega4-9-fr.tex
\section{Propriétés constructibles}
\label{section:IV.9-fr}

Soient $S$ un préschéma, $f:X\to S$ un morphisme de
\emph{présentation finie} (1.6.1), $\mathcal F$ un $\mathcal O_X$-Module
quasi-cohérent de \emph{présentation finie}. On se propose dans ce paragraphe
de donner des critères assurant, par exemple, que l'ensemble des $s\in S$ tels
que le $k(s)$-préschéma
$X_s=f^{-1}(s)=X\times_S\operatorname{Spec}(k(s))$ a une certaine propriété,
ou tels que le $\mathcal O_{X_s}$-Module
$\mathcal F_s=\mathcal F\otimes_{\mathcal O_X}k(s)$ a une certaine propriété,
est \emph{localement constructible}, ou tout au moins
\emph{ind-constructible} (1.9.4)~; on verra que c'est le cas pour la plupart
des propriétés qui s'introduisent en Géométrie algébrique. En supposant
seulement $f$ \emph{localement de présentation finie}, on donnera aussi (9.9)
des critères pour que l'ensemble des points $x\in X$ où la fibre $X_{f(x)}$
(ou le $\mathcal O_{X_{f(x)}}$-Module $\mathcal F_{f(x)}$) a une certaine
propriété, soit \emph{localement constructible}. On verra au \S~12 que ces
résultats, joints à l'hypothèse supplémentaire que $f$ est \emph{plat}
(resp. \emph{propre et plat}), permettent de prouver que les ensembles
considérés dans $X$ (resp. dans $S$) sont même \emph{ouverts} dans de nombreux
cas.

\subsection{Le principe de l'extension finie}
\label{subsection:IV.9.1-fr}

\begin{proposition}[9.1.1]
\label{IV.9.1.1-fr}
\textnormal{(Principe de l'extension finie).} Soient $k$ un corps,
$\mathcal C$ un ensemble d'extensions de $k$. On suppose vérifiées les
conditions suivantes~:
\begin{enumerate}[label=(\roman*)]
  \item Si $K\in\mathcal C$ et s'il existe un $k$-homomorphisme $K\to K'$
  (où $K'$ appartient à l'univers où l'on se place), alors
  $K'\in\mathcal C$.
  \item Si $K\in\mathcal C$, il existe une sous-extension $K'$ de $K$, de
  type fini sur $k$, telle que $K'\in\mathcal C$.
  \item Si $K\in\mathcal C$ est le corps de fractions d'une $k$-algèbre
  intègre de type fini $A$, il existe $f\in A-\{0\}$ tel que pour tout idéal
  maximal $\mathfrak m$ de $A$ on ait $A_f/\mathfrak m\in\mathcal C$.
\end{enumerate}
\oldpage[IV-3]{55}
Soit alors $\Omega$ une extension algébriquement close de $k$ (dans l'univers
considéré). Les conditions suivantes sont équivalentes~:
\begin{enumerate}[label=\emph{\alph*)}]
  \item $\mathcal C$ est non vide.
  \item Il existe $K'\in\mathcal C$ qui est une extension finie de $k$.
  \item On a $\Omega\in\mathcal C$.
\end{enumerate}
\end{proposition}

La condition (i) implique évidemment que \emph{b)} entraîne \emph{c)}, et
\emph{c)} entraîne trivialement \emph{a)}~; prouvons que \emph{a)} entraîne
\emph{b)}. En vertu de (ii) et (iii) il existe une extension
$K'\in\mathcal C$ qui est de la forme $A/\mathfrak m$, où $A$ est une
$k$-algèbre de type fini sur $k$ et $\mathfrak m$ un idéal maximal de $A$.
On sait, par le théorème des zéros de Hilbert (Bourbaki, \emph{Alg. comm.},
chap.~V, \S~3, n\textsuperscript{o}~1, cor.~2 du th.~1) que $K'$ est une
extension finie de $k$.

\begin{corollary}[9.1.2]
\label{IV.9.1.2-fr}
Sous les hypothèses (i), (ii), (iii) de (9.1.1), si $k$ est algébriquement
clos et si $\mathcal C$ n'est pas vide, $\mathcal C$ contient toutes les
extensions de $k$ dans l'univers envisagé.
\end{corollary}

En effet, comme \emph{a)} entraîne \emph{c)}, on a $k\in\mathcal C$, et la
conclusion résulte de l'hypothèse (i).

\begin{remark}[9.1.3]
\label{IV.9.1.3-fr}
En pratique, on vérifiera la condition (ii) de (9.1.1) en notant que $K$ est
limite inductive de ses sous-extensions de type fini $K_\alpha$, et en
appliquant les résultats du \S~8, compte tenu éventuellement du fait que pour
$K_\alpha\subset K_\beta$, $K_\beta$ est fidèlement plat sur $K_\alpha$.
Fréquemment l'ensemble $\mathcal C$ est formé des \emph{corps} appartenant à
un ensemble $\mathcal C'$ de $k$-algèbres qui vérifie la condition suivante~:
\begin{enumerate}[label=(i bis)]
  \item Si $A\in\mathcal C'$ et s'il existe un homomorphisme de $k$-algèbres
  $A\to A'$ (où $A'$ appartient à l'univers où l'on se place), alors
  $A'\in\mathcal C'$.
\end{enumerate}
Lorsqu'il en est ainsi, la condition (i) est trivialement satisfaite, et on
vérifiera généralement la condition (iii) de (9.1.1) en notant que le corps
des fractions $K$ de $A$ est la limite inductive des algèbres $A_f$ (pour la
relation $D(f)\supset D(g)$), et en appliquant les résultats du \S~8, compte
tenu éventuellement du fait que les morphismes $D(g)\to D(f)$ sont des
immersions ouvertes.

Par contre, lorsque (i bis) n'est pas vérifiée, la démonstration de (iii) est
souvent plus délicate, et est liée à des critères de constructibilité qui
seront développés plus loin.
\end{remark}

Voici des exemples typiques d'application du principe de l'extension finie~:

\begin{proposition}[9.1.4]
\label{IV.9.1.4-fr}
Soient $k$ un corps, $\Omega$ une extension algébriquement close de $k$,
$X$ et $Y$ deux préschémas de type fini sur $k$. Les conditions suivantes
sont équivalentes~:
\begin{enumerate}[label=\emph{\alph*)}]
  \item Il existe un $\Omega$-morphisme $X_{(\Omega)}\to Y_{(\Omega)}$
  (resp. un $\Omega$-morphisme possédant l'une (déterminée) des propriétés
  (i) à (xiv) de (8.10.5)).
  \item Il existe une extension finie $k'$ de $k$ et un $k'$-morphisme
  $X_{(k')}\to Y_{(k')}$ (resp. un $k'$-morphisme possédant la propriété
  considérée).
  \item Il existe une extension $K$ de $k$ et un $K$-morphisme
  $X_{(K)}\to Y_{(K)}$ (resp. un $K$-morphisme possédant la propriété
  considérée).
\end{enumerate}
\end{proposition}

On applique la remarque (9.1.3) en prenant pour $\mathcal C'$ l'ensemble de
toutes les $k$-algèbres $A$ (de l'univers où on est placé) telles qu'il existe
un $A$-morphisme $X\otimes_k A\to Y\otimes_k A$ (resp. un morphisme ayant
celle des propriétés de (8.10.5) que l'on considère). La condition (i bis) de
(9.1.3) est alors vérifiée grâce au fait que les

\oldpage[IV-3]{56}
propriétés envisagées dans (8.10.5) sont toutes stables par changement de base.
La condition (iii) de (9.1.1) est donc satisfaite en vertu de (8.8.2, (i))
(resp. (8.10.5)), puisque $\operatorname{Spec}(k)$ est quasi-compact et
quasi-séparé. Reste à vérifier la condition (ii) de (9.1.1) qui découle encore
de (8.8.2, (i)) et de (8.10.5), en considérant $K$ comme limite inductive de
ses sous-extensions de type fini. On conclut donc par (9.1.1).

En particulier, s'il existe une extension $K$ de $k$ et un $K$-isomorphisme
$X_{(K)}\xrightarrow{\sim}Y_{(K)}$, on dit que $X$ et $Y$ sont
\emph{géométriquement isomorphes}.

Le corollaire suivant généralise (II, 6.6.5)~:

\begin{corollary}[9.1.5]
\label{IV.9.1.5-fr}
Soient $k$ un corps, $X$ un $k$-préschéma. S'il existe une extension $K$ de
$k$ telle que $X_{(K)}$ soit projectif (resp. quasi-projectif) sur $K$, alors
$X$ est projectif (resp. quasi-projectif) sur $k$.
\end{corollary}

Le morphisme $\operatorname{Spec}(K)\to\operatorname{Spec}(k)$ étant
fidèlement plat et quasi-compact, il résulte déjà de (2.7.1, (v)) que $X$ est
de type fini sur $k$. L'hypothèse signifie qu'il existe une immersion fermée
(resp. une immersion) $X_{(K)}\to\mathbf P_K^r=\mathbf P_k^r\otimes_k K$
(II, 5.5.4, (ii) et 5.5.3)~; appliquant (9.1.4) pour la propriété (iv)
(resp. (ii)) de (8.10.5), on en déduit qu'il y a une extension finie $k'$ de
$k$ et une immersion fermée (resp. une immersion)
$X_{(k')}\to\mathbf P_{k'}^r$, autrement dit $X_{(k')}$ est projectif
(resp. quasi-projectif) sur $k'$. On conclut alors par (II, 6.6.5).

\begin{proposition}[9.1.6]
\label{IV.9.1.6-fr}
Soient $k$ un corps, $\Omega$ une extension algébriquement close de $k$,
$X$ un préschéma de type fini sur $k$, $\mathcal F$, $\mathcal G$ deux
$\mathcal O_X$-Modules cohérents. Supposons qu'il existe un isomorphisme
$\mathcal F\otimes_k\Omega\xrightarrow{\sim}\mathcal G\otimes_k\Omega$.
Alors il existe une extension finie $k'$ de $k$ et un isomorphisme
$\mathcal F\otimes_k k'\xrightarrow{\sim}\mathcal G\otimes_k k'$.
\end{proposition}

Le raisonnement est le même que dans (9.1.4), en appliquant (8.5.2, (i))
(on utilise ici, dans la démonstration de la propriété (iii) de (9.1.1), le
fait que les morphismes $D(g)\to D(f)$ (avec les notations de (9.1.3)) sont
des immersions ouvertes, et \emph{a fortiori} des morphismes plats).

\subsection{Propriétés constructibles et ind-constructibles}
\label{subsection:IV.9.2-fr}

\begin{definition}[9.2.1]
\label{IV.9.2.1-fr}
Soit $\mathbf P(X,\mathcal F,k)$ une relation. On dit que $\mathbf P$ est une
propriété \emph{constructible} (resp. \emph{ind-constructible}) de préschémas
algébriques si les deux conditions suivantes sont vérifiées~:
\begin{enumerate}[label=(\roman*)]
  \item Si $k$ est un corps, $X$ un préschéma algébrique sur $k$,
  $\mathcal F$ un $\mathcal O_X$-Module cohérent, $k'$ une extension de $k$,
  alors, pour que $\mathbf P(X,\mathcal F,k)$ soit vraie, il faut et il suffit
  que $\mathbf P(X_{(k')},\mathcal F\otimes_k k',k')$ soit vraie.
  \item Soient $S$ un préschéma intègre noethérien, de point générique $\eta$,
  $u:X\to S$ un morphisme de type fini, $\mathcal F$ un $\mathcal O_X$-Module
  cohérent. Pour tout $s\in S$, posons
  $X_s=u^{-1}(s)=X\times_S\operatorname{Spec}(k(s))$,
  $\mathcal F_s=\mathcal F\otimes_{\mathcal O_S}k(s)$. Soit $E$ l'ensemble
  des $s\in S$ tels que $\mathbf P(X_s,\mathcal F_s,k(s))$ soit vraie.
  Alors l'un des ensembles $E$, $S-E$ (resp. l'ensemble $E$) contient un
  ouvert non vide (et par suite est un voisinage de $\eta$) (resp. contient
  un ouvert non vide s'il contient $\eta$).
\end{enumerate}
\end{definition}

\oldpage[IV-3]{57}
\begin{remarks}[9.2.2]
\label{IV.9.2.2-fr}
\begin{enumerate}[label=(\roman*)]
  \item Il s'agit là bien entendu d'une convention de langage de nature
  métamathématique et non d'une définition mathématique proprement dite. On a
  des «~définitions~» analogues pour des relations entre $k$, un ou plusieurs
  $k$-préschémas algébriques, des $k$-morphismes entre ces préschémas, des
  Modules cohérents sur ces préschémas ou des homomorphismes entre ces Modules.

  \item Nous aurons aussi à considérer des relations où figurent des
  \emph{parties constructibles} de préschémas. Par exemple, soit
  $\mathbf P(X,Z,k)$ une relation~; nous dirons (par abus de langage) que
  $\mathbf P$ est une propriété constructible (resp. ind-constructible) de la
  partie constructible $Z$ de $X$ si les deux conditions suivantes sont
  vérifiées~:
  \begin{enumerate}[label=$\arabic*^{\mathrm o}$]
    \item Si $k$ est un corps, $X$ un préschéma algébrique sur $k$, $Z$ une
    \emph{partie constructible de $X$}, $k'$ une extension de $k$, alors, pour
    que $\mathbf P(X,Z,k)$ soit vraie, il faut et il suffit que
    $\mathbf P(X_{(k')},p^{-1}(Z),k')$ soit vraie
    ($p:X_{(k')}\to X$ étant la projection canonique).
    \item Soient $S$ un préschéma intègre noethérien, de point générique
    $\eta$, $u:X\to S$ un morphisme de type fini, $Z$ une
    \emph{partie constructible de $X$}. Pour tout $s\in S$, posons
    $X_s=u^{-1}(s)$, $Z_s=Z\cap X_s$. Soit $E$ l'ensemble des $s\in S$ tels
    que $\mathbf P(X_s,Z_s,k(s))$ soit vraie. Alors l'un des ensembles $E$,
    $S-E$ (resp. l'ensemble $E$) contient un ouvert non vide (resp. contient
    un ouvert non vide s'il contient $\eta$).
  \end{enumerate}

  On notera que dans la condition $2^{\mathrm o}$ il faut supposer que $Z$ est
  une partie constructible de $X$, et non seulement que $Z_s$ est une partie
  constructible de $X_s$ pour tout $s$~; la première de ces deux propriétés
  entraîne la seconde (1.8.2), mais non réciproquement.

  \item Si $\mathbf P$ est une propriété constructible, il est clair qu'il en
  est de même de «~non $\mathbf P$~». Si $\mathbf P$, $\mathbf Q$ sont deux
  propriétés constructibles (resp. ind-constructibles), il en est de même des
  propriétés «~$\mathbf P$ ou $\mathbf Q$~» et «~$\mathbf P$ et
  $\mathbf Q$~»~; en effet, si, sous les hypothèses de (9.2.1, (ii)), $E$,
  $E'$ sont deux parties de $S$ et si $E$ contient un ouvert non vide, il en
  est de même de $E\cup E'$, et si $S-E$ et $S-E'$ contiennent chacun un
  ouvert non vide, il en est de même de
  $S-(E\cup E')=(S-E)\cap(S-E')$.

  \item Soit $\mathbf P(X,\mathcal F,k)$ une relation vérifiant la condition
  (9.2.1, (i))~; soient $S$ un préschéma, $u:X\to S$ un morphisme de type
  fini~; avec les notations de (9.2.1, (ii)), soit $E$ l'ensemble des $s\in S$
  tels que $\mathbf P(X_s,\mathcal F_s,k(s))$ soit vraie. Soit d'autre part
  $g:S'\to S$ un morphisme quelconque, et posons
  $X'=X_{(S')}$, $\mathcal F'=\mathcal F\otimes_{\mathcal O_S}\mathcal O_{S'}$~;
  alors il résulte de la transitivité des fibres (I, 3.6.4) et de la condition
  (9.2.1, (i)) que \emph{l'ensemble des $s'\in S'$ tels que
  $\mathbf P(X'_{s'},\mathcal F'_{s'},k(s'))$ soit vraie est égal à
  $g^{-1}(E)$}. Cela s'étend aussitôt au cas où il figure dans $\mathbf P$
  plusieurs préschémas, Modules, morphismes de préschémas ou homomorphismes de
  Modules, ainsi qu'aux propriétés du type considéré dans (ii).

  \item Comme nous le verrons dans la suite de ce paragraphe et dans le courant
  du reste du chap.~IV, la plupart des propriétés $\mathbf P$ vérifiant la
  condition (9.2.1, (i)) vérifient aussi (9.2.1, (ii)). Comme exceptions
  possibles, notons la propriété d'être \emph{projectif}, ou
  \emph{quasi-projectif}, ou \emph{affine}, ou \emph{quasi-affine} sur le
  corps de base (pour un préschéma algébrique)~; nous verrons (9.6.2) que ces
  propriétés sont \emph{ind-constructibles}, mais nous donnerons plus tard un
  exemple où $S$ est une partie ouverte non vide de $\operatorname{Spec}(\mathbf Z)$
  (ou une partie ouverte d'une courbe elliptique sur un corps fini) et où
  toutes les fibres $X_s$

  \oldpage[IV-3]{58}
  sauf la fibre générique $X_\eta$ sont projectives sur $k(s)$ (tous les
  préschémas $X_s$ étant de dimension 2).
\end{enumerate}
\end{remarks}

\begin{proposition}[9.2.3]
\label{IV.9.2.3-fr}
Soient $\mathbf P$ une propriété constructible (resp. ind-constructible) de
préschémas algébriques, $S$ un préschéma, $X$ un préschéma de présentation
finie sur $S$, $\mathcal F$ un $\mathcal O_X$-Module quasi-cohérent de
présentation finie. Alors l'ensemble $E$ des $s\in S$ tels que
$\mathbf P(X_s,\mathcal F_s,k(s))$ soit vraie est localement constructible
(resp. ind-constructible). En outre, si $S$ est irréductible de point générique
$\eta$, un des deux ensembles $E$, $S-E$ est un voisinage de $\eta$ dans $S$
(resp. $E$ est un voisinage de $\eta$ s'il contient ce point).
\end{proposition}

Pour démontrer ces assertions, on peut se borner au cas où
$S=\operatorname{Spec}(A)$ est affine. On sait alors qu'il existe un
sous-anneau $A_0$ de $A$ qui est une $\mathbf Z$-algèbre de type fini, un
$A_0$-préschéma de type fini $X_0$ et un $\mathcal O_{X_0}$-Module cohérent
$\mathcal F_0$ tels que $X$ soit isomorphe à $X_0\otimes_{A_0}A$ et
$\mathcal F$ à $\mathcal F_0\otimes_{A_0}A$ (8.9.1). Soit
$p:S\to S_0=\operatorname{Spec}(A_0)$ le morphisme correspondant à
l'injection $A_0\to A$, et soit $E_0$ l'ensemble des $s_0\in S_0$ tels que
\[
\mathbf P((X_0)_{s_0},(\mathcal F_0)_{s_0},k(s_0))
\]
soit vraie~; alors, en vertu de (9.2.2, (iv)), on a $E=p^{-1}(E_0)$~; on peut
par suite (1.8.2) se borner au cas où $S$ est le spectre d'une
$\mathbf Z$-algèbre de type fini, donc un schéma noethérien. Utilisons le
critère de constructibilité $(\mathbf 0_{\mathrm{III}},9.2.3)$ (resp. le
critère d'ind-constructibilité (1.9.10))~; on est alors ramené, en utilisant
comme ci-dessus (9.2.2, (iv)) et en remplaçant $S$ par un sous-schéma fermé
intègre de $S$, au cas où $S$ est noethérien et intègre, et où il faut prouver
que $E$ est rare dans $S$ ou contient un ouvert non vide de $S$ (resp. que $E$
contient un ouvert non vide de $S$ s'il contient le point générique)~; mais
cela est garanti en vertu de la condition (9.2.1, (ii)).

On notera que l'on n'a utilisé (9.2.1, (ii)) que lorsque $S$ est le spectre
d'un anneau intègre de type fini sur $\mathbf Z$. Il est clair d'autre part
que l'énoncé de (9.2.3) s'applique aussi lorsqu'il figure dans $\mathbf P$
plusieurs préschémas, Modules sur ces préschémas, morphismes de préschémas ou
homomorphismes de Modules. Il s'applique encore lorsqu'il y figure des parties
(en nombre fini) des préschémas considérés, pourvu qu'on impose à ces parties
la condition d'être \emph{localement constructibles}. En effet, la restriction
au cas où $S$ est affine montre qu'on peut se borner au cas où ces parties sont
\emph{constructibles}~: on applique alors (8.3.11) qui montre (avec les
notations précédentes) qu'une partie constructible de $X$ est l'image
réciproque d'une partie constructible de $X_0$ pour un choix convenable de
$A_0$.

\begin{corollary}[9.2.4]
\label{IV.9.2.4-fr}
Soient $\mathbf P$ une propriété constructible (resp. ind-constructible) de
préschémas algébriques, $X$, $Y$ deux $S$-préschémas de présentation finie,
$f:X\to Y$ un $S$-morphisme. Pour tout $s\in S$, posons
$X_s=X\times_S\operatorname{Spec}(k(s))$,
$Y_s=Y\times_S\operatorname{Spec}(k(s))$,
$f_s=f\times1:X_s\to Y_s$. Alors l'ensemble $E$ des $s\in S$ tels que pour
tout $y\in Y_s$, la propriété $\mathbf P(f_s^{-1}(y),k(y))$ soit vraie, est
localement constructible (resp. ind-constructible).
\end{corollary}

En effet, soit $Z$ l'ensemble des $y\in Y$ tels que
$\mathbf P(f^{-1}(y),k(y))$ soit vraie. Comme les fibres $f^{-1}(y)$ et
$f_s^{-1}(y)$ sont isomorphes, on voit que $E$ est l'ensemble des $s\in S$
tels que $Y_s\subset Z$~; si $g:Y\to S$ est le morphisme structural, on a
donc $E=S-g(Y-Z)$.

\oldpage[IV-3]{59}
Or, $f$ est de présentation finie (1.6.2, (v)), donc il résulte de (9.2.3) que
$Z$ est localement constructible (resp. ind-constructible) dans $Y$, donc
$Y-Z$ est localement constructible (resp. pro-constructible) dans $Y$. Comme
$g$ est de présentation finie, $g(Y-Z)$ est localement constructible
(resp. pro-constructible) dans $S$, en vertu du théorème de Chevalley (1.8.4)
(resp. de (1.9.5, (vii)))~; donc $E$ est localement constructible
(resp. ind-constructible) dans $S$.

\begin{remark}[9.2.5]
\label{IV.9.2.5-fr}
On notera que si $\mathbf P$ est une propriété de préschémas algébriques pour
laquelle la prop. (9.2.3) est vraie, alors $\mathbf P$ satisfait aussi à la
condition (9.2.1, (ii))~: cela résulte en effet du fait que dans un espace
irréductible noethérien, un ensemble constructible est rare ou contient un
ouvert non vide $(\mathbf 0_{\mathrm{III}},9.2.3)$.
\end{remark}

\begin{proposition}[9.2.6]
\label{IV.9.2.6-fr}
Désignons par $\mathbf P$ une des propriétés suivantes d'un $k$-préschéma
algébrique $X$~:
\begin{enumerate}[label=(\roman*)]
  \item $X$ est vide.
  \item $X$ est fini sur $k$.
  \item $X$ est radiciel sur $k$.
  \item $\dim(X)$ appartient à une partie donnée $\Phi$ de l'ensemble
  $\overline{\mathbf Z}=\mathbf Z\cup\{-\infty\}$.
\end{enumerate}
Alors $\mathbf P$ est constructible.
\end{proposition}

Il est clair que (i) et (ii) sont des cas particuliers de (iv), en prenant
respectivement pour $\Phi$ l'ensemble $\{-\infty\}$ et l'ensemble
$\{-\infty,0\}$. On n'a donc qu'à démontrer (iii) et (iv). Dans chacun de ces
deux cas la condition (i) de (9.2.1) est remplie en vertu de (2.7.1, (xv)) et
(4.1.4). D'autre part, dans le cas (iii), la propriété $\mathbf P$ vérifie la
conclusion de (9.2.3) en vertu de (1.8.7)~; il reste donc à voir qu'il en est de
même dans le cas (iv). Cela va résulter de la proposition plus précise
suivante~:

\begin{proposition}[9.2.6.1]
\label{IV.9.2.6.1-fr}
Si $f:X\to S$ est un morphisme de présentation finie, la fonction
$s\mapsto\dim(f^{-1}(s))$ est localement constructible.
\end{proposition}

La question est locale sur $S$, donc on peut supposer que
$S=\operatorname{Spec}(A)$ est affine et prouver que pour tout $n$, l'ensemble
$E$ des $s\in S$ tels que $\dim(X_s)=n$ est constructible. Le même
raisonnement que dans (9.2.3) ramène au cas où $A$ est noethérien et intègre,
et il suffit alors de prouver le

\begin{corollary}[9.2.6.2]
\label{IV.9.2.6.2-fr}
Soient $S$ un préschéma intègre noethérien de point générique $\eta$,
$f:X\to S$ un morphisme de type fini. Alors il existe un voisinage de $\eta$
dans $S$ tel que la fonction $s\mapsto\dim(X_s)$ soit constante dans ce
voisinage.
\end{corollary}

Les images par $f$ des composantes irréductibles (en nombre fini) de $X$ qui
ne rencontrent pas $X_\eta$, sont contenues dans des parties fermées de $S$ ne
contenant pas $\eta$ (puisque $S$ est intègre
$(\mathbf 0_{\mathrm I},2.1.5)$), donc (en remplaçant $S$ par un voisinage
ouvert de $\eta$) on peut se borner au cas où toutes les composantes
irréductibles $X_i$ de $X$ rencontrent $X_\eta$~; désignons encore par $X_i$
le sous-préschéma fermé réduit de $X$ ayant $X_i$ pour espace sous-jacent~;
comme $\dim(X_s)=\sup_i(\dim((X_i)_s))$ (4.1.1), on peut se borner au cas où
$X$ est irréductible. Il existe alors un recouvrement fini $(U_j)$ de $X$ par
des

\oldpage[IV-3]{60}
ouverts affines partout denses, et les nombres $\dim((U_j)_\eta)$ sont tous
égaux à $n=\dim(X_\eta)$ (4.1.1.3)~; on peut par suite se borner au cas où $X$
est affine, donc aussi $X_\eta$. Il existe alors, en vertu de (4.1.2), un
ouvert non vide $W$ de $X$ tel que $W\supset X_\eta$, et un
$\mathbf{k}(\eta)$-morphisme fini surjectif
$h:W_\eta\to\mathbf V_{\mathbf{k}(\eta)}^n
(=\operatorname{Spec}(\mathbf{k}(\eta)[T_1,\ldots,T_n]))$~; en appliquant
(8.8.2, (i)) et la méthode de (8.1.2, \emph{a)}), on en déduit (en remplaçant
au besoin $S$ par un voisinage de $\eta$) que l'on a $h=g_\eta$, où
$g:W\to\mathbf V_A^n(=\operatorname{Spec}(A[T_1,\ldots,T_n]))$ est un
$S$-morphisme, et on peut supposer ce morphisme fini et surjectif en vertu de
(8.10.5, (vi) et (x)). On en conclut que pour tout $s\in S$, le morphisme
$g_s:W_s\to\mathbf V_{\mathbf{k}(s)}^n$ est fini et surjectif, donc que
$\dim(W_s)=n$ (4.1.2).

\subsection{Propriétés constructibles de morphismes de préschémas algébriques}

\begin{proposition}[9.3.1]
\label{IV.9.3.1-fr}
Soit $\mathbf P(X,k)$ une propriété constructible (resp. ind-constructible)
de préschémas algébriques. Désignons par $\mathbf P'(f,X,Y,k)$ la relation
suivante~: $f:X\to Y$ est un $k$-morphisme de $k$-préschémas algébriques tel
que pour tout $y\in Y$, on ait la propriété
$\mathbf P(f^{-1}(y),k(y))$. Alors $\mathbf P'$ est une propriété
constructible (resp. ind-constructible).
\end{proposition}

En effet, comme $\mathbf P$ vérifie la condition (9.2.1, (i)), il en est de
même de $\mathbf P'$ en vertu de la transitivité des fibres (I, 3.6.4)~; par
ailleurs le fait que $\mathbf P'$ vérifie la condition (9.2.1, (ii)) résulte
de (9.2.4), en raison de la remarque (9.2.5).

\begin{proposition}[9.3.2]
\label{IV.9.3.2-fr}
Désignons par $\mathbf P$ une des propriétés suivantes d'un $k$-morphisme
$f:X\to Y$ de $k$-préschémas algébriques~:
\begin{enumerate}[label=(\roman*)]
  \item $f$ est surjectif.
  \item $f$ est quasi-fini.
  \item $f$ est radiciel.
  \item Pour tout $y\in Y$, $\dim(f^{-1}(y))$ appartient à $\Phi$ (notation
  de (9.2.6)).
\end{enumerate}
Alors $\mathbf P$ est constructible.
\end{proposition}

Cela résulte aussitôt de (9.3.1) et (9.2.6) si l'on tient compte de ce que $f$
est de type fini (1.5.4, (v)), de la caractérisation des morphismes radiciels
(I, 3.5.8), et de celle des morphismes quasi-finis (II, 6.2.2).

\begin{proposition}[9.3.3]
\label{IV.9.3.3-fr}
Supposons vérifiées les hypothèses de (8.8.1) dont nous gardons les
notations~; supposons en outre que $S_\alpha$ soit quasi-compact, $X_\alpha$
et $Y_\alpha$ de présentation finie sur $S_\alpha$, et soit
$f_\alpha:X_\alpha\to Y_\alpha$ un $S_\alpha$-morphisme. Soit $\mathbf P$ une
propriété ind-constructible de morphismes de préschémas algébriques. Pour tout
$s\in S$ (resp. $s_\lambda\in S_\lambda$) posons
$X_s=X\times_S\operatorname{Spec}(k(s))$,
$Y_s=Y\times_S\operatorname{Spec}(k(s))$,
$f_s=f\times1:X_s\to Y_s$ (resp.
$X_{\lambda,s_\lambda}=X_\lambda\times_{S_\lambda}
\operatorname{Spec}(k(s_\lambda))$,
$Y_{\lambda,s_\lambda}=Y_\lambda\times_{S_\lambda}
\operatorname{Spec}(k(s_\lambda))$,
$f_{\lambda,s_\lambda}=f_\lambda\times1:
X_{\lambda,s_\lambda}\to Y_{\lambda,s_\lambda}$). Alors, afin que pour
tout $s\in S$ on ait la propriété $\mathbf P(f_s)$, il faut et il suffit
qu'il existe $\lambda\geqslant\alpha$ tel que pour tout
$s_\lambda\in S_\lambda$, on ait $\mathbf P(f_{\lambda,s_\lambda})$.
\end{proposition}

En effet, soit $E$ (resp. $E_\lambda$) l'ensemble des $s\in S$ (resp.
$s_\lambda\in S_\lambda$) tels que la propriété $\mathbf P(f_s)$ (resp.
$\mathbf P(f_{\lambda,s_\lambda})$) soit vraie~; il résulte de
(9.2.2, (iv)) que l'on a
$E_\mu=u_{\lambda\mu}^{-1}(E_\lambda)$ pour $\lambda\leqslant\mu$, et
$E=u_\lambda^{-1}(E_\lambda)$~; en outre, en vertu de (9.2.3), $E$ (resp.
$E_\lambda$) est ind-constructible dans $S$ (resp. $S_\lambda$)~; la
proposition résulte donc de (8.3.4) appliqué à la partie ind-constructible
$E$ de $S$.

\oldpage[IV-3]{61}
On généralise sans peine ce résultat à des propriétés $\mathbf P$ du type
considéré dans (9.2.3, (i) et (ii)).

\begin{remark}[9.3.4]
\label{IV.9.3.4-fr}
La conjonction de (9.3.3) et de (9.3.2, (ii)) démontre l'assertion
(8.10.5, (xi)).
\end{remark}

\begin{proposition}[9.3.5]
\label{IV.9.3.5-fr}
Soit $\mathbf P(X,Y,k)$ la propriété~: «~$X$ et $Y$ sont deux préschémas de
type fini sur le corps $k$, et il existe une extension $k'$ de $k$ et un
$k'$-morphisme $g:X_{(k')}\to Y_{(k')}$ vérifiant $\mathbf Q(g)$~», où
$\mathbf Q$ est une des propriétés (i) à (xiv) de (8.10.5). Alors $\mathbf P$
est une propriété ind-constructible.
\end{proposition}

La définition de $\mathbf P$ montre en effet que la condition (9.2.1, (i))
est satisfaite, compte tenu de ce que la propriété $\mathbf Q(g)$ est stable
par changement du corps de base, et de ce que deux extensions de $k$ peuvent
toujours être considérées comme des sous-extensions d'une troisième extension
de $k$. Pour vérifier (9.2.1, (ii)), on peut se borner au cas où
$S=\operatorname{Spec}(A)$ est affine~; si $K=\mathbf{k}(\eta)$, corps des fractions de
$A$, il existe par hypothèse et en vertu de (9.1.4) une extension finie $K'$ de
$K$ et un $K'$-morphisme
$g':(X_\eta)_{(K')}\to(Y_\eta)_{(K')}$, vérifiant $\mathbf Q(g')$, et $K'$ est
évidemment le corps des fractions d'une $A$-algèbre intègre finie $A'$. Si
l'on pose $S'=\operatorname{Spec}(A')$, il résulte alors de (8.10.5) qu'il y
a un voisinage $U'$ du point générique $\eta'$ de $S'$ tel que, si l'on pose
$X'=X\otimes_A A'$, $Y'=Y\otimes_A A'$, il existe, pour tout $s'\in U'$, un
morphisme $X'_{s'}\to Y'_{s'}$ ayant la propriété $\mathbf Q$. Mais le
morphisme $h:S'\to S$ est fini, donc fermé, et comme
$h^{-1}(\eta)=\{\eta'\}$, $h(U')$ contient un voisinage ouvert $U$ de $\eta$
dans $S$~; pour tout $s\in U$, il y a donc un $s'\in U'$ tel que $h(s')=s$,
et comme
$X'_{s'}=X_s\otimes_{\mathbf{k}(s)}\mathbf{k}(s')$,
$Y'_{s'}=Y_s\otimes_{\mathbf{k}(s)}\mathbf{k}(s')$, la propriété
$\mathbf P(X_s,Y_s,\mathbf{k}(s))$ est vraie pour tout $s\in U$.

\begin{env}[9.3.6]
\label{IV.9.3.6-fr}
\textbf{Exemple.} Prenons par exemple pour $\mathbf Q$ la propriété d'être un
isomorphisme. Alors, en conjuguant (9.3.5) et (9.3.3), on a la propriété
suivante~: les notations et hypothèses étant celles de (8.8.1), $S_\alpha$
étant quasi-compact, $X_\alpha$ et $Y_\alpha$ de présentation finie sur
$S_\alpha$, afin que, pour tout $s\in S$, $X_s$ et $Y_s$ soient
géométriquement isomorphes (9.1.4), il faut et il suffit qu'il existe
$\lambda\geqslant\alpha$ tel que, pour tout $s_\lambda\in S_\lambda$,
$X_{\lambda,s_\lambda}$ et $Y_{\lambda,s_\lambda}$ soient géométriquement
isomorphes.
\end{env}

On a un résultat analogue lorsque les préschémas que l'on considère sont munis
de «~lois de composition~» d'une certaine espèce
$(\mathbf 0_{\mathrm{III}},8.2.1)$, par exemple des «~préschémas en groupes~»,
«~préschémas en anneaux~», etc. $(\mathbf 0_{\mathrm{III}},8.2.3)$. Alors
l'énoncé précédent est encore valable lorsque par «~isomorphisme~» on entend
des isomorphismes de préschémas qui sont des homomorphismes pour les lois de
composition considérées $(\mathbf 0_{\mathrm{III}},8.2.2)$~; il suffit ici
d'utiliser non seulement (8.10.5) mais aussi (8.8.2, (i)) en remarquant que la
notion d'homomorphisme pour une loi de composition s'exprime en écrivant que
des diagrammes de morphismes de préschémas sont commutatifs (il faut bien
entendu que les morphismes de transition $X_\mu\to X_\lambda$ et
$Y_\mu\to Y_\lambda$ pour $\lambda\leqslant\mu$ soient des homomorphismes
pour les lois de composition envisagées).

On peut aussi, au lieu de considérer des morphismes de préschémas comme dans
(9.3.5), considérer des homomorphismes de Modules, en utilisant (9.1.6) au
lieu de (9.1.5).

\oldpage[IV-3]{62}

\subsection{Constructibilité de certaines propriétés des modules}
\label{subsection:IV.9.4-fr}

\begin{env}[9.4.1]
\label{IV.9.4.1-fr}
\textbf{Notations.} Dans ce numéro et les suivants jusqu'à la fin du \S~9,
nous utiliserons systématiquement les notations suivantes~: étant donné un
morphisme $f:X\to S$, nous poserons, pour tout $s\in S$,
$X_s=f^{-1}(s)=X\times_S\operatorname{Spec}(\mathbf{k}(s))$~; pour tout
$\mathcal O_X$-Module quasi-cohérent $\mathcal F$, $\mathcal F_s$ désignera
le $\mathcal O_{X_s}$-Module
$\mathcal F\otimes_{\mathcal O_S}\mathbf{k}(s)$, et pour tout homomorphisme
$u:\mathcal F\to\mathcal G$ de $\mathcal F$ dans un $\mathcal O_X$-Module
quasi-cohérent $\mathcal G$, $u_s:\mathcal F_s\to\mathcal G_s$ sera le
morphisme $p^*(u)$, en désignant par $p$ la projection canonique $X_s\to X$.
Pour toute section $g$ de $\mathcal F$ au-dessus de $X$, on désignera par
$g_s$ l'image de $g$ par l'homomorphisme canonique
$\Gamma(X,\mathcal F)\to\Gamma(X_s,\mathcal F_s)$. Pour toute partie $Z$ de
$X$, on notera $Z_s$ l'image réciproque
$p^{-1}(Z)=Z\cap X_s$ (\textbf{I}, 3.6.1). Enfin, si $Y$ est un second
$S$-préschéma et $h:X\to Y$ un $S$-morphisme, on notera $h_s$ le morphisme
$h\times1:X_s\to Y_s$.
\end{env}

\begin{proposition}[9.4.2]
\label{IV.9.4.2-fr}
Soient $S$ un préschéma intègre de point générique $\eta$, $f:X\to S$ un
morphisme de présentation finie, $\mathcal F$, $\mathcal G$, $\mathcal H$
trois $\mathcal O_X$-Modules quasi-cohérents de présentation finie. Soient
$u:\mathcal F\to\mathcal G$, $v:\mathcal G\to\mathcal H$ deux
homomorphismes de $\mathcal O_X$-Modules, et supposons que la suite
$\mathcal F_\eta\xrightarrow{u_\eta}\mathcal G_\eta
\xrightarrow{v_\eta}\mathcal H_\eta$ soit exacte. Alors il existe un
voisinage ouvert $U$ de $\eta$ dans $S$ tel que, pour tout $s\in U$, la suite
$\mathcal F_s\xrightarrow{u_s}\mathcal G_s
\xrightarrow{v_s}\mathcal H_s$ soit exacte.
\end{proposition}

Avec les notations générales de (9.2.1), il s'agit ici de la relation
$\mathbf P(X,\mathcal F,\mathcal G,\mathcal H,u,v,k)$~: «~$X$ est un
préschéma algébrique sur le corps $k$,
$\mathcal F\xrightarrow{u}\mathcal G\xrightarrow{v}\mathcal H$ une suite
exacte de $\mathcal O_X$-Modules quasi-cohérents~». Comme, pour toute extension
$k'$ de $k$, la projection canonique $X_{(k')}\to X$ est un morphisme
fidèlement plat (2.2.13), la condition (9.2.1, (i)) est satisfaite (2.2.7). En
vertu de (9.2.3), on peut se borner au cas où $S$ est intègre et
\emph{noethérien}, auquel cas $X$ est noethérien, et $\mathcal F$,
$\mathcal G$, $\mathcal H$ sont des $\mathcal O_X$-Modules cohérents.
L'hypothèse implique qu'il existe un voisinage ouvert $U$ de $\eta$ dans $S$
tel que la suite
$\mathcal F|f^{-1}(U)\to\mathcal G|f^{-1}(U)\to\mathcal H|f^{-1}(U)$ soit
\emph{exacte}, en vertu de (8.5.8, (ii)) appliqué suivant la méthode générale
de (8.1.2, \emph{a)}), et l'on peut donc déjà supposer que la suite
$\mathcal F\xrightarrow{u}\mathcal G\xrightarrow{v}\mathcal H$ est exacte~;
il suffit évidemment de prouver que l'on a
$\operatorname{Ker}(v_s)=(\operatorname{Ker}v)_s$ et
$\operatorname{Im}(u_s)=(\operatorname{Im}u)_s$ pour tout $s$ voisin de
$\eta$ dans $S$~; par suite (tenant compte de ce que les
$\mathcal O_X$-Modules $\mathcal F$, $\mathcal G$, $\mathcal H$ sont
\emph{cohérents} et de $(\mathbf 0_{\mathrm I},5.3.4)$) on est ramené à
prouver la proposition dans le cas particulier où la suite
\[
0\to\mathcal F\xrightarrow{u}\mathcal G\xrightarrow{v}\mathcal H\to0
\]
est exacte. Or, il y a alors un ouvert $U$ dans $S$ contenant $\eta$ et tel
que $\mathcal H|f^{-1}(U)$ soit \emph{plat} sur $U$ (6.9.1)~; il résulte donc
de (2.1.8) que pour tout $s\in U$, la suite
$0\to\mathcal F_s\to\mathcal G_s\to\mathcal H_s\to0$ est exacte, ce qui
achève la démonstration.

\begin{corollary}[9.4.3]
\label{IV.9.4.3-fr}
Soient $S$ un préschéma intègre, de point générique $\eta$, $f:X\to S$ un
morphisme de présentation finie. Soit
$\mathcal L^\bullet=(\mathcal L^i)_{i\in\mathbf Z}$ un complexe de
$\mathcal O_X$-Modules quasi-cohérents de présentation finie. Pour tout $i$,
il existe un voisinage ouvert $U$ de $\eta$ dans $S$ tel que les homomorphismes
canoniques
\begin{equation}
\label{IV.9.4.3.1-fr}
(\mathcal H^i(\mathcal L^\bullet))_s\to
\mathcal H^i(\mathcal L_s^\bullet) \tag{9.4.3.1}
\end{equation}
soient bijectifs pour tout $s\in U$.
\end{corollary}

\oldpage[IV-3]{63}

On peut évidemment se borner à un complexe à trois termes de degrés
$-1$, $0$, $+1$~:
\[
\mathcal M\xrightarrow{u}\mathcal N\xrightarrow{v}\mathcal P
\]
et à $i=0$~; l'homomorphisme à considérer est alors l'homomorphisme canonique
$(\operatorname{Ker}v/\operatorname{Im}u)_s\to
\operatorname{Ker}(v_s)/\operatorname{Im}(u_s)$. Utilisant (8.9.1) et
(8.5.2, (i)), on voit qu'on peut se ramener au cas où $S$ (donc aussi $X$)
est \emph{noethérien}, et par suite $\mathcal M$, $\mathcal N$, $\mathcal P$
des $\mathcal O_X$-Modules \emph{cohérents}~; alors
$\operatorname{Im}(u)$ et $\operatorname{Ker}(v)$ sont aussi cohérents
$(\mathbf 0_{\mathrm I},5.3.4)$ et en outre il existe un voisinage $U$ de
$\eta$ tel que pour $s\in U$, on ait
$\operatorname{Ker}(v_s)=(\operatorname{Ker}(v))_s$ et
$\operatorname{Im}(u_s)=(\operatorname{Im}(u))_s$ (9.4.2)~; la conclusion
résulte alors de (9.4.2) appliqué à la suite exacte
\[
0\to\operatorname{Im}(u)\to\operatorname{Ker}(v)\to
\operatorname{Ker}(v)/\operatorname{Im}(u)\to0,
\]
compte tenu de ce que $\mathcal O_\eta=k(\eta)$ (puisque $S$ est intègre) et
par conséquent la suite
\[
0\to(\operatorname{Im}u)_\eta\to(\operatorname{Ker}v)_\eta\to
(\operatorname{Ker}v/\operatorname{Im}u)_\eta\to0
\]
est exacte.

\begin{proposition}[9.4.4]
\label{IV.9.4.4-fr}
Soient $f:X\to S$ un morphisme de présentation finie, $\mathcal F$,
$\mathcal G$, $\mathcal H$ trois $\mathcal O_X$-Modules quasi-cohérents de
présentation finie. Soient $u:\mathcal F\to\mathcal G$,
$v:\mathcal G\to\mathcal H$ deux homomorphismes de $\mathcal O_X$-Modules.
Alors l'ensemble $E$ des $s\in S$ tels que la suite
$\mathcal F_s\xrightarrow{u_s}\mathcal G_s
\xrightarrow{v_s}\mathcal H_s$ soit exacte est localement constructible.
\end{proposition}

Compte tenu de (9.2.3), il faut établir que la propriété $\mathbf P$ considérée
dans (9.4.2) est constructible. On a déjà remarqué dans la démonstration de
(9.4.2) que la condition (9.2.1, (i)) est satisfaite pour cette propriété, et
il reste à vérifier la condition (9.2.1, (ii)). Supposons donc $S$ intègre
noethérien, de point générique $\eta$, et prouvons que $E$ ou $S-E$ est un
voisinage de $\eta$. Si $\eta\in E$, notre assertion résulte de (9.4.2), et
l'on peut donc se borner au cas où $\eta\notin E$, autrement dit, la suite
$\mathcal F_\eta\xrightarrow{u_\eta}\mathcal G_\eta
\xrightarrow{v_\eta}\mathcal H_\eta$ n'est pas exacte. Distinguons alors
deux cas~:

$1^\circ$ Posons $w=v\circ u$, et supposons d'abord que
$w_\eta=v_\eta\circ u_\eta\neq0$. Comme $\mathcal F$, $\mathcal G$,
$\mathcal H$ sont cohérents, il en est de même de
$\mathcal N=\operatorname{Ker}(w)$ $(\mathbf 0_{\mathrm I},5.3.4)$~; il
résulte donc de (9.4.2) appliqué à la suite exacte
$0\to\mathcal N\to\mathcal F\xrightarrow{w}\mathcal H$ qu'il y a un
voisinage $U$ de $\eta$ dans $S$ tel que, pour $s\in S$,
$\operatorname{Ker}(w_s)=\mathcal N_s$~; en restreignant $S$, on peut donc
supposer cette relation vérifiée pour tout $s\in S$. Soit $j$ l'injection
canonique $\mathcal N\to\mathcal F$, et posons
$\mathcal M=\operatorname{Coker}(j)$~; l'exactitude à droite du foncteur
$\mathcal F\mapsto\mathcal F_s$ entraîne que
$\mathcal M_s=\operatorname{Coker}(j_s)$ pour tout $s\in S$. L'hypothèse
$w_\eta\neq0$ signifie que $\mathcal M_\eta\neq0$~; comme $\mathcal M$ est
cohérent $(\mathbf 0_{\mathrm I},5.3.4)$, il résulte de (1.8.6) qu'il y a un
voisinage ouvert $U$ de $\eta$ dans $S$ tel que $\mathcal M_s\neq0$ pour
$s\in U$, donc $w_s\neq0$ pour $s\in U$, et \emph{a fortiori} $S-E$ est un
voisinage de $\eta$.

$2^\circ$ Supposons que $w_\eta=0$~; en vertu de (8.5.2, (i)), appliqué
suivant la méthode générale de (8.1.2, \emph{a)}), il existe un voisinage
ouvert $U$ de $\eta$ tel que $w|f^{-1}(U)=0$~; remplaçant $S$ par $U$, on peut
déjà supposer $w=0$ dans $X$. Alors
$\mathcal F\xrightarrow{u}\mathcal G\xrightarrow{v}\mathcal H$ est un
complexe à trois termes $\mathcal L^\bullet$, auquel on peut appliquer
(9.4.3)~; par hypothèse on a
$(\mathcal H^0(\mathcal L^\bullet))_\eta=
\mathcal H^0(\mathcal L^\bullet_\eta)\neq0$, et
$\mathcal H^0(\mathcal L)$ est cohérent
$(\mathbf 0_{\mathrm I},5.3.4\text{ et }5.3.3)$, donc il résulte de (1.8.6)
qu'il y a un voisinage ouvert $U$ de $\eta$ tel que
$(\mathcal H^0(\mathcal L^\bullet))_s\neq0$ pour tout $s\in U$~; mais comme
on peut supposer que
$\mathcal H^0(\mathcal L_s^\bullet)=(\mathcal H^0(\mathcal L^\bullet))_s$
pour $s\in U$ par (9.4.3), on voit de nouveau que $S-E$ est un voisinage de
$\eta$ dans $S$.

\oldpage[IV-3]{64}

\begin{corollary}[9.4.5]
\label{IV.9.4.5-fr}
Soient $f:X\to S$ un morphisme de présentation finie, $\mathcal F$,
$\mathcal G$ deux $\mathcal O_X$-Modules quasi-cohérents de présentation finie,
$u:\mathcal F\to\mathcal G$ un homomorphisme de $\mathcal O_X$-Modules.
Alors l'ensemble des $s\in S$ tels que $u_s$ soit injectif
\textnormal{(}resp. surjectif, bijectif, nul\textnormal{)} est localement
constructible.
\end{corollary}

Il suffit d'appliquer (9.4.4) aux suites
$0\to\mathcal F\xrightarrow{u}\mathcal G$,
$\mathcal F\xrightarrow{u}\mathcal G\to0$,
$0\to\mathcal F\xrightarrow{u}\mathcal G\to0$,
$\mathcal F\xrightarrow{u}\mathcal G\xrightarrow{1_\mathcal G}\mathcal G$.

\begin{corollary}[9.4.6]
\label{IV.9.4.6-fr}
Soient $f:X\to S$ un morphisme de présentation finie, $\mathcal F$ un
$\mathcal O_X$-Module quasi-cohérent de présentation finie. Soit $h$ une
section de $\mathcal F$ au-dessus de $X$~; pour tout $s\in S$, soit $h_s$ la
section correspondante de $\mathcal F_s$ au-dessus de $X_s$ \textnormal{(}pour
le morphisme projection $X_s\to X$\textnormal{)}. Alors, l'ensemble des
$s\in S$ tels que $h_s=0$ est localement constructible.
\end{corollary}

Il suffit de remarquer que $h$ correspond à un homomorphisme
$u:\mathcal O_X\to\mathcal F$ $(\mathbf 0_{\mathrm I},5.1.1)$ et $h_s$ à
l'homomorphisme $u_s$.

\begin{proposition}[9.4.7]
\label{IV.9.4.7-fr}
Soient $f:X\to S$ un morphisme de présentation finie, $\mathcal F$ un
$\mathcal O_X$-Module quasi-cohérent de présentation finie. L'ensemble $E$
\textnormal{(}resp. $E'$\textnormal{)} des $s\in S$ tels que $\mathcal F_s$
soit un $\mathcal O_{X_s}$-Module localement libre \textnormal{(}resp.
localement libre de rang $n$\textnormal{)} est localement constructible.
\end{proposition}

Si $X$ est un préschéma algébrique sur un corps $k$, $\mathcal F$ un
$\mathcal O_X$-Module cohérent, $k'$ une extension de $k$, alors, pour que
$\mathcal F$ soit localement libre (resp. localement libre de rang $n$), il
faut et il suffit qu'il en soit de même de $\mathcal F\otimes_k k'$, puisque la
projection $X_{(k')}\to X$ est un morphisme fidèlement plat (2.2.7). Autrement
dit, la condition (9.2.1, (i)) est vérifiée pour les propriétés dont on veut
démontrer la constructibilité, et il reste à vérifier (9.2.1, (ii))~; on peut
donc encore supposer que $S$ est affine, noethérien et intègre. Il y a de
nouveau quatre cas à envisager~:

$1^\circ$ $\eta\in E$. Il résulte de (8.5.5), appliqué suivant la méthode
générale de (8.1.2, \emph{a)}), qu'il existe un voisinage ouvert $U$ de $\eta$
dans $S$ tel que $\mathcal F|f^{-1}(U)$ soit localement libre~;
\emph{a fortiori} $\mathcal F_s$ est localement libre pour tout $s\in U$.

$2^\circ$ $\eta\in E'$. Même raisonnement que dans le $1^\circ$.

$3^\circ$ $\eta\in S-E$. Comme $\mathcal F_\eta$ est un
$\mathcal O_{X_\eta}$-Module cohérent, dire qu'il n'est pas localement libre
équivaut à dire qu'il \emph{n'est pas plat} sur $\mathcal O_{X_\eta}$
(Bourbaki, \emph{Alg. comm.}, chap.~II, \S~5, n\textsuperscript{o}~2,
cor.~2 du th.~1). Le fait que $S-E$ est un voisinage de $\eta$ résultera donc
du lemme plus général suivant (appliqué au cas où $g$ est l'identité)~:

\begin{lemma}[9.4.7.1]
\label{IV.9.4.7.1-fr}
Soient $S$ un préschéma intègre noethérien, de point générique $\eta$, $X$,
$Y$ deux $S$-préschémas de type fini sur $S$, $g:X\to Y$ un $S$-morphisme,
$\mathcal F$ un $\mathcal O_X$-Module cohérent. Si $\mathcal F_\eta$ n'est
pas $g_\eta$-plat, alors il existe un voisinage ouvert $U$ de $\eta$ dans $S$
tel que pour tout $s\in U$, $\mathcal F_s$ ne soit pas $g_s$-plat.
\end{lemma}

Compte tenu de (2.1.2) et de Bourbaki, \emph{Alg. comm.}, chap.~I\textsuperscript{er},
\S~2, n\textsuperscript{o}~3, Remarque~1, l'hypothèse signifie qu'il existe un
ouvert non vide $V$ de $Y_\eta$, et un homomorphisme injectif
$v:\mathcal M\to\mathcal N$ de $\mathcal O_V$-Modules \emph{cohérents}, tel
que l'homomorphisme
$1\otimes v:\mathcal F_\eta\otimes_{\mathcal O_V}\mathcal M\to
\mathcal F_\eta\otimes_{\mathcal O_V}\mathcal N$ ne soit pas injectif. On a
$V=Y_\eta\cap W$, où $W$ est ouvert dans $Y$ (\textbf{I}, 3.6.1), et il
résulte de (8.5.2, (i) et (ii)), appliqué suivant la méthode de
(8.1.2, \emph{a)}), qu'il existe un voisinage ouvert $U_0$ de $\eta$ dans $S$,
deux $\mathcal O_Z$-Modules

\oldpage[IV-3]{65}

cohérents $\mathcal G$, $\mathcal H$ (où $Z=W\cap h^{-1}(U_0)$,
$h:Y\to S$ étant le morphisme structural) et un $\mathcal O_Z$-homomorphisme
$u:\mathcal G\to\mathcal H$ tel que $\mathcal M=\mathcal G_\eta$,
$\mathcal N=\mathcal H_\eta$ et $v=u_\eta$~; on peut donc supposer $U_0$ pris
tel que, pour $s\in U_0$, $u_s:\mathcal G_s\to\mathcal H_s$ soit injectif
(9.4.5). Mais pour tout $s\in U_0$, l'homomorphisme
$1\otimes u_s:\mathcal F_s\otimes_{\mathcal O_{Z_s}}\mathcal G_s\to
\mathcal F_s\otimes_{\mathcal O_{Z_s}}\mathcal H_s$ n'est autre que
$(1\otimes u)_s$~; l'hypothèse que $(1\otimes u)_\eta$ est non injectif
entraîne donc (9.4.5) l'existence d'un ouvert non vide $U\subset U_0$ tel que
pour tout $s\in U$, $(1\otimes u)_s$ soit non injectif, et par suite
$\mathcal F_s$ n'est pas $g_s$-plat pour tout $s\in U$.

$4^\circ$ $\eta\in S-E'$. Il est clair que $S-E\subset S-E'$, et si
$\eta\in S-E$, $S-E'$ est \emph{a fortiori} un voisinage de $\eta$ d'après le
$3^\circ$. Supposons donc $\eta\in E$, donc $\mathcal F_\eta$ localement
libre~; ces hypothèses entraînent que $X_\eta$ est non connexe, et que les
rangs du $\mathcal O_{X_\eta}$-Module localement libre $\mathcal F_\eta$ ne
sont pas les mêmes sur les diverses composantes connexes de $X_\eta$. Or, il
résulte de (8.4.2), appliqué suivant la méthode de (8.1.2, \emph{a)}), que l'on
peut supposer (en remplaçant $S$ par un voisinage ouvert de $\eta$) que $X$ et
$X_\eta$ ont le même nombre de composantes connexes, les composantes connexes
de $X_\eta$ étant les intersections de $X_\eta$ et des composantes connexes de
$X$. La conclusion résulte alors du raisonnement fait dans le $2^\circ$,
appliqué à chacune des composantes connexes de $X$ (qui sont en nombre fini).

\begin{remark}[9.4.7.2]
\label{IV.9.4.7.2-fr}
Le lemme (9.4.7.1) sera plus tard généralisé et débarrassé d'hypothèses
noethériennes (11.2.8).
\end{remark}

\begin{proposition}[9.4.8]
\label{IV.9.4.8-fr}
Soient $S$ un préschéma localement noethérien, $f:X\to S$ un morphisme de type
fini, $\mathcal F$ un $\mathcal O_X$-Module cohérent. On suppose que pour tout
$s\in S$, $X_s$ soit un préschéma localement intègre. Alors l'ensemble $E$
\textnormal{(}resp. $E'$\textnormal{)} des $s\in S$ tels que $\mathcal F_s$
soit un $\mathcal O_{X_s}$-Module de torsion \textnormal{(}resp. un
$\mathcal O_{X_s}$-Module sans torsion\textnormal{)} est localement
constructible.
\end{proposition}

On peut évidemment supposer $S=\operatorname{Spec}(A)$ affine et noethérien et
prouver que $E$ (resp. $E'$) est alors constructible en utilisant le critère
$(\mathbf 0_{\mathrm{III}},9.2.3)$~; remplaçant $S$ par le sous-préschéma fermé
réduit de $S$ ayant pour espace sous-jacent une partie fermée irréductible $Y$
de $S$, et notant (\textbf{I}, 3.6.4) que pour $s\in Y$, la fibre
$(X_{(Y)})_s$ s'identifie canoniquement à $X_s$ et le faisceau
$(\mathcal F\otimes_{\mathcal O_S}\mathcal O_Y)_s$ à $\mathcal F_s$, on voit
qu'on est ramené au cas où $S$ est intègre de point générique $\eta$, et à
prouver que $E$ ou $S-E$ (resp. $E'$ ou $S-E'$) est un voisinage de $\eta$
dans $S$. Notons en outre que $X$ est réunion finie d'ouverts affines $U_i$,
de type fini sur $S$, et chacun des $(U_i)_s$ est induit sur un ouvert de
$X_s$, donc localement intègre~; en outre, si $(U_i)_\eta$ est vide (resp. non
vide), on sait que $(U_i)_s$ est aussi vide (resp. non vide) dans un voisinage
de $\eta$ (9.2.6). On peut donc supposer tous les $(U_i)_\eta$ non vides et
intègres, et dire que $\mathcal F_s$ est de torsion (resp. sans torsion)
équivaut à dire que chacun des $\mathcal F_s|(U_i)_s$ est de torsion (resp.
sans torsion). On est donc ramené au cas où $X=\operatorname{Spec}(B)$ est
affine, et $\mathcal F=\widetilde M$, où $M$ est un $B$-module de type fini~;
on posera $B_s=B\otimes_A k(s)$, $M_s=M\otimes_A k(s)$ et l'on peut supposer
$B_\eta$ intègre. Nous avons quatre cas à envisager~:

$1^\circ$ $\eta\in E$~; $M_\eta$ est donc un $B_\eta$-module de torsion de
type fini, et il y a par suite un $h\neq0$ dans $B_\eta$ tel que
$hM_\eta=0$~; en vertu de (8.5.2, (i)), appliqué suivant la méthode

\oldpage[IV-3]{66}

de (8.1.2, \emph{a)}), on peut (en remplaçant au besoin $S$ par un voisinage
de $\eta$) supposer que $h\in\Gamma(X_\eta,\mathcal O_{X_\eta})$ est de la
forme $g_\eta$, où $g\in\Gamma(X,\mathcal O_X)$~; soit
$u:\mathcal F\to\mathcal F$ l'endomorphisme de $\mathcal F$ défini par la
multiplication par $g$~; par hypothèse, on a $u_\eta=0$, donc (9.4.5)
l'endomorphisme $u_s:\mathcal F_s\to\mathcal F_s$ défini par la multiplication
par $g_s$, est nul dans un voisinage de $\eta$. D'autre part, soit
$v:\mathcal O_X\to\mathcal O_X$ l'endomorphisme défini par la multiplication
par $g$~; comme $B_\eta$ est intègre et $h=g_\eta\neq0$, $v_\eta$ est injectif,
et il résulte donc de (9.4.5) que $v_s$ est un endomorphisme injectif de
$\mathcal O_{X_s}$ pour $s$ voisin de $\eta$, autrement dit, $g_s$ est un
élément $\mathcal O_{X_s}$-régulier pour ces valeurs de $s$~; donc
$\mathcal F_s$ est de torsion dans un voisinage de $\eta$.

$2^\circ$ $\eta\in S-E$. Dire qu'un $B_\eta$-module de type fini $M_\eta$
n'est pas un module de torsion signifie que son quotient $M_\eta/T$ par son
sous-module de torsion est $\neq0$, et comme c'est un $B_\eta$-module sans
torsion de type fini, il est isomorphe à un sous-module d'un $B_\eta$-module
$B_\eta^n$~; il existe par suite un homomorphisme
$w:M_\eta\to B_\eta$ qui est $\neq0$. Appliquant (8.5.2, (i)) suivant la
méthode de (8.1.2, \emph{a)}), on en déduit (en remplaçant au besoin $S$ par un
voisinage de $\eta$) qu'il existe un homomorphisme
$v:\mathcal F\to\mathcal O_X$ tel que $v_\eta=\widetilde w$. L'hypothèse
$v_\eta\neq0$ entraîne donc (9.4.5) que $v_s\neq0$ dans un voisinage de
$\eta$, et comme $X_s$ est localement intègre, $\mathcal F_s$ n'est pas de
torsion pour ces valeurs de $s$.

$3^\circ$ $\eta\in E'$. Comme $M_\eta$ est un $B_\eta$-module de type fini
sans torsion, il existe un homomorphisme injectif
$w:M_\eta\to B_\eta^n$. Utilisant (8.5.2, (i)) et (9.4.5) comme dans le
$2^\circ$ (en restreignant au besoin $S$), on en déduit qu'il existe un
homomorphisme $v:\mathcal F\to\mathcal O_X^n$ tel que $v_\eta=\widetilde w$,
et que pour $s$ voisin de $\eta$,
$v_s:\mathcal F_s\to\mathcal O_{X_s}^n$ est injectif~; pour ces valeurs de
$s$, $\mathcal F_s$ est donc sans torsion.

$4^\circ$ $\eta\in S-E'$. Soit $T$ le sous-module de torsion de $M_\eta$~;
par hypothèse $T\neq0$, et $T$ est de type fini puisque $M_\eta$ est
noethérien. Utilisant cette fois (8.5.2, (i) et (ii)) on voit (en restreignant
au besoin $S$) qu'il existe un $\mathcal O_X$-Module cohérent $\mathcal G$ et
un homomorphisme injectif $u:\mathcal G\to\mathcal F$ tels que
$\mathcal G_\eta=\widetilde T$ et que $u_\eta$ soit l'injection canonique
$\widetilde T\to\mathcal F_\eta$. Il résulte alors du $1^\circ$ et de (1.8.6)
que dans un voisinage de $\eta$, $\mathcal G_s$ est un
$\mathcal O_{X_s}$-Module de torsion $\neq0$, et d'autre part il résulte de
(9.4.5) que dans un voisinage de $\eta$, $u_s$ est injectif. On en conclut que
dans un voisinage de $\eta$, le sous-Module de torsion de $\mathcal F_s$ n'est
pas nul. C.Q.F.D.

\begin{remark}[9.4.9]
\label{IV.9.4.9-fr}
La propriété «~$X$ est un $k$-préschéma algébrique localement intègre~» ne
vérifie pas la condition (9.2.1, (i)), et il n'est donc pas certain que
l'énoncé (9.4.8) soit encore valable quand on ne fait aucune hypothèse sur $S$
et qu'on suppose seulement que $f$ est un morphisme de présentation finie et
$\mathcal F$ un $\mathcal O_X$-Module de présentation finie. Considérons
toutefois le cas particulier suivant~: $S_0$ étant un préschéma localement
noethérien, soient $f_0:X_0\to S_0$ un morphisme de type fini, tel que les
fibres $(X_0)_{s_0}$ soient localement intègres (pour tout $s_0\in S_0$), et
$\mathcal F_0$ un $\mathcal O_{X_0}$-Module cohérent~; soit $g:S\to S_0$ un
morphisme quelconque, posons
$X=X_0\times_{S_0}S$, $\mathcal F=\mathcal F_0\otimes_{\mathcal O_{S_0}}
\mathcal O_S$, $f=f_0\times1_S:X\to S$, et \emph{supposons} que pour tout
$s\in S$, la fibre $X_s$ soit encore localement intègre. Alors l'ensemble $E$
(resp. $E'$) des $s\in S$ tels que $\mathcal F_s$ soit de torsion (resp. sans

\oldpage[IV-3]{67}

torsion) est encore \emph{localement constructible}. En effet, soit $s\in S$
et soit $s_0=g(s)$~; il suffira (compte tenu de (1.8.2)) de prouver que, pour
que $\mathcal F_s$ soit de torsion (resp. sans torsion), il faut et il suffit
que $(\mathcal F_0)_{s_0}$ le soit. Or,
$\operatorname{Supp}(\mathcal F_s)$ est l'image réciproque de
$\operatorname{Supp}((\mathcal F_0)_{s_0})$ par la projection
$p:X_s\to(X_0)_{s_0}$ (\textbf{I}, 9.1.13)~; comme $p$ est fidèlement plat et
quasi-compact, dire que $\operatorname{Supp}(\mathcal F_s)$ contient un point
maximal de $X_s$ équivaut à dire que
$\operatorname{Supp}((\mathcal F_0)_{s_0})$ contient un point maximal de
$(X_0)_{s_0}$ (1.1.5 et 2.3.4)~; d'où notre assertion en ce qui concerne
l'ensemble $E$ (\textbf{I}, 7.4.6). Si le sous-Module de torsion $\mathcal T$
de $(\mathcal F_0)_{s_0}$ n'est pas nul,
$\mathcal T\otimes_{k(s_0)}k(s)$ (qui est de torsion d'après ce qui précède)
n'est pas nul et s'identifie à un sous-Module de $\mathcal F_s$ (2.2.7), donc
le sous-Module de torsion de $\mathcal F_s$ n'est pas nul. Enfin, si
$(\mathcal F_0)_{s_0}$ est sans torsion, on peut supposer (en considérant un
ouvert affine de $(X_0)_{s_0}$) que $(\mathcal F_0)_{s_0}$ est isomorphe à un
sous-Module d'un $\mathcal O_{(X_0)_{s_0}}^n$, donc $\mathcal F_s$ est
isomorphe à un sous-Module d'un $\mathcal O_{X_s}^n$ (2.2.7), et ceci établit
notre assertion concernant $E'$.
\end{remark}

\subsection{Constructibilité de propriétés topologiques}
\label{subsection:IV.9.5-fr}

\begin{proposition}[9.5.1]
\label{IV.9.5.1-fr}
Soient $f:X\to S$ un morphisme de présentation finie, $Z$ une partie
localement constructible de $X$. Alors l'ensemble $E$ des $s\in S$ tels que
$Z_s\neq\varnothing$ est localement constructible.
\end{proposition}

En effet, on a $E=f(Z)$, et il suffit d'appliquer le th. de Chevalley (1.8.4).

\begin{corollary}[9.5.2]
\label{IV.9.5.2-fr}
Si $Z'$, $Z''$ sont deux parties localement constructibles de $X$, l'ensemble
des $s\in S$ tels que $Z'_s\subset Z''_s$ \textnormal{(}resp.
$Z'_s=Z''_s$\textnormal{)} est localement constructible.
\end{corollary}

En effet, la relation $Z'_s\subset Z''_s$ équivaut à
$(Z'\cap(\complement Z''))_s=\varnothing$ et
$Z'\cap\complement Z''$ est localement constructible.

\begin{proposition}[9.5.3]
\label{IV.9.5.3-fr}
Soient $f:X\to S$ un morphisme de présentation finie, $Z$, $Z'$ deux parties
localement constructibles de $X$ telles que $Z\subset Z'$. Alors l'ensemble
$E$ des $s\in S$ tels que $Z_s$ soit dense dans $Z'_s$ est localement
constructible dans $S$.
\end{proposition}

Il faut vérifier les deux conditions de (9.2.2, (iii)). En ce qui concerne la
première, considérons un préschéma $X$ algébrique sur un corps $k$, deux parties
constructibles $Z$, $Z'$ de $X$ telles que $Z\subset Z'$, et une extension
$k'$ de $k$. Alors la projection canonique $p:X_{(k')}\to X$ est un morphisme
fidèlement plat et quasi-compact, et l'on a donc
$p^{-1}(\overline Z)=\overline{p^{-1}(Z)}$ et
$p^{-1}(\overline{Z'})=\overline{p^{-1}(Z')}$ en vertu de (2.3.10)~; comme
$p$ est surjectif, la relation $\overline Z=\overline{Z'}$ est équivalente à
$\overline{p^{-1}(Z)}=\overline{p^{-1}(Z')}$.

Vérifions maintenant la seconde condition, et supposons donc $S$ affine,
noethérien et intègre, de point générique $\eta$. Distinguons deux cas~:

$1^\circ$ $\eta\in S-E$, autrement dit, $Z_\eta$ n'est pas dense dans
$Z'_\eta$~; il existe donc dans $X$ un ouvert $V$ tel que
$V\cap Z_\eta=\varnothing$ et $V\cap Z'_\eta\neq\varnothing$. Comme $X$ est
noethérien, $V$ est localement constructible, donc il en est de même de
$V\cap Z$, et en vertu de (9.5.1), il y a un voisinage $U$ de $\eta$ dans $S$
tel que pour tout $s\in U$, on ait $(V\cap Z)_s=\varnothing$ et
$(V\cap Z')_s\neq\varnothing$~; cela entraîne que $Z_s$ n'est pas dense dans
$Z'_s$ pour $s\in U$, autrement dit $U\subset S-E$.

\oldpage[IV-3]{68}

$2^\circ$ $\eta\in E$, donc $Z_\eta$ est dense dans $Z'_\eta$. Montrons
d'abord que l'on peut supposer $Z'$ \emph{fermé}. En effet, $Z_\eta$ est dense
dans $\overline{Z'_\eta}$ (adhérence prise dans $X_\eta$)~; posons
$V_\eta=X_\eta-\overline{Z_\eta}$, qui est ouvert dans $X_\eta$ et ne
rencontre pas $Z'_\eta$~; on peut supposer $V_\eta$ de la forme $V\cap X_\eta$,
où $V$ est ouvert (donc constructible) dans $X$, et l'hypothèse
$V_\eta\cap Z'_\eta=\varnothing$ entraîne alors
$V_s\cap Z'_s=\varnothing$ pour tout $s$ voisin de $\eta$ en vertu de (9.5.1).
Remplaçant $S$ par un voisinage ouvert de $\eta$, on peut donc supposer que
$V\cap Z'=\varnothing$, donc $V\cap\overline{Z'}=\varnothing$ (adhérence prise
dans $X$), et par suite $(\overline{Z'})_\eta=\overline{Z'_\eta}$, d'où notre
assertion. L'ensemble $Z'$ est alors réunion de ses composantes irréductibles
en nombre fini, et en restreignant encore $S$ à un voisinage de $\eta$, on
peut supposer que toutes les composantes irréductibles $Z'_i$ de $Z'$
rencontrent $X_\eta$, d'où résulte que $X_\eta$ contient les points génériques
des $Z'_i$ $(\mathbf 0_{\mathrm I},2.1.8)$. Dire que $Z_s$ est dense dans
$Z'_s$ équivaut alors à dire que chacun des $(Z\cap Z'_i)_s$ est dense dans
$(Z'_i)_s$, et l'on est ainsi ramené au cas où $Z'$ est irréductible.
Remplaçant alors au besoin $X$ par le sous-préschéma réduit ayant $Z'$ pour
espace sous-jacent, on voit qu'on peut supposer que $Z'=X$ et que $X$ est
\emph{intègre} et domine $S$. Enfin, en recouvrant $X$ par un nombre fini
d'ouverts affines $W_j$ et remplaçant $Z$ par $Z\cap W_j$, on peut supposer
que $X=\operatorname{Spec}(A)$, où $A$ est un anneau noethérien \emph{intègre}.
Comme $X_\eta$ est noethérien intègre et que $Z_\eta$ est constructible dans
$X_\eta$ et dense dans $X_\eta$, $Z_\eta$ contient un ouvert non vide de
$X_\eta$ $(\mathbf 0_{\mathrm{III}},9.2.2)$, que l'on peut supposer de la forme
$(D(t))_\eta$, où $t$ est un élément $\neq0$ de $A$. En remplaçant au besoin
$S$ par un voisinage de $\eta$, on peut en outre, en vertu de la relation
$(D(t))_\eta\subset Z_\eta$, supposer que $D(t)\subset Z$ (9.5.2). Enfin,
comme l'homothétie de rapport $t_\eta$ dans $\mathcal O_{X_\eta}$ est
injective, il résulte de (9.4.5) que pour $s$ voisin de $\eta$, $t_s$ est
$\mathcal O_{X_s}$-régulier, donc $(X_s)_{t_s}$ est dense dans $X_s$, et
\emph{a fortiori} il en est de même de $Z_s$ qui contient $(X_s)_{t_s}$.

\begin{corollary}[9.5.4]
\label{IV.9.5.4-fr}
Soient $f:X\to S$ un morphisme de présentation finie, $Z$ une partie
localement constructible de $X$. L'ensemble $E$ des $s\in S$ tels que $Z_s$
soit fermé \textnormal{(}resp. ouvert, resp. localement fermé\textnormal{)}
dans $X_s$ est localement constructible dans $S$.
\end{corollary}

Dire que $Z_s$ est ouvert dans $X_s$ signifie que
$(X-Z)_s=X_s-Z_s$ est fermé dans $X_s$, et comme $X-Z$ est localement
constructible, on peut se borner à considérer l'ensemble des $s\in S$ tels que
$Z_s$ soit fermé et l'ensemble des $s\in S$ tels que $Z_s$ soit localement
fermé.

Vérifions encore dans chaque cas les deux conditions de (9.2.2, (ii)). La
première résulte du fait que $p:X_{(k')}\to X$ est fidèlement plat et
quasi-compact, et de (2.3.12) et (2.3.14). Vérifions donc la seconde condition,
$S$ étant supposé affine, noethérien et intègre, de point générique $\eta$.
Posons $Z'=\overline Z$~; le même raisonnement que dans (9.5.3) montre que
$Z'_\eta$ est égal à l'adhérence de $Z_\eta$ dans $X_\eta$~; en vertu de
(9.5.3), il y a donc un voisinage $U$ de $\eta$ tel que pour $s\in U$, $Z_s$
soit dense dans $Z'_s$, ce dernier étant fermé dans $X_s$. Dire que $Z_s$ est
fermé dans $X_s$ signifie alors que $Z''_s=\varnothing$, où $Z''=Z'-Z$~; il
résulte donc de (9.5.1) que l'ensemble $E$ des $s\in S$ où
$Z''_s=\varnothing$ est tel que $E$ ou $S-E$ contienne un voisinage de
$\eta$. Dire que $Z_s$ est localement fermé dans $X_s$ signifie que $Z''_s$
est fermé dans $X_s$~; il suffit donc d'appliquer le résultat précédent en
remplaçant $Z$ par $Z''$, qui est localement constructible dans $X$.

\oldpage[IV-3]{69}

\begin{proposition}[9.5.5]
\label{IV.9.5.5-fr}
Soient $f:X\to S$ un morphisme de présentation finie, $Z$ une partie
localement constructible de $X$ telle que, pour tout $s\in S$, $Z_s$ soit
localement fermé dans $X_s$. Pour tout $s\in S$, soit
$D(s)\subset\mathbf Z\cup\{-\infty\}$ l'ensemble des dimensions des
composantes irréductibles de $Z_s$. Alors la fonction $s\mapsto D(s)$ est
localement constructible dans $S$.
\end{proposition}

Soit $\Phi$ une partie finie de $\mathbf Z\cup\{-\infty\}$~; il s'agit de
montrer que l'ensemble des $s\in S$ tels que $D(s)=\Phi$ est localement
constructible~; compte tenu de (9.2.3), nous avons encore à vérifier les deux
conditions de (9.2.2, (iii)).

En ce qui concerne la première, il s'agit de voir que si $X$ est un préschéma
algébrique sur un corps $k$, $Z$ une partie localement fermée de $X$, $k'$ une
extension de $k$, $p:X_{(k')}\to X$ la projection canonique, alors l'ensemble
des dimensions des composantes irréductibles de $Z$ est le même que celui des
dimensions des composantes irréductibles de $p^{-1}(Z)$~; compte tenu de
l'existence d'un sous-préschéma de $X$ ayant $Z$ pour espace sous-jacent
(\textbf{I}, 5.2.1), cela résulte de (4.2.8).

Pour la seconde condition de (9.2.2, (iii)), on est dans le cas où $S$ est
noethérien et intègre de point générique $\eta$, et $f:X\to S$ est un
morphisme de type fini, de sorte que $X$ est noethérien. Le sous-espace
$Z_\eta$ de l'espace noethérien $X_\eta$ est par hypothèse localement fermé,
donc a un nombre fini de composantes irréductibles $Z_{i\eta}$, qui sont
localement fermées dans $X_\eta$. Il existe par suite pour chaque indice $i$
une partie localement fermée $Z_i$ de $X$ telle que $(Z_i)_\eta=Z_{i\eta}$,
donc si $Z'=\bigcup_i Z_i$, on a $Z'_\eta=Z_\eta$. Mais comme $Z$ et $Z'$ sont
localement constructibles, on peut, en remplaçant $S$ par un voisinage de
$\eta$, supposer que $Z'=Z$ (9.5.1). En outre, pour $i\neq j$,
$Z_{i\eta}\cap Z_{j\eta}$ est rare et \emph{fermé} dans $Z_{i\eta}$~; donc
(9.5.3 et 9.5.4), on peut encore supposer, en restreignant $S$, que pour
$j\neq i$, $(Z_i)_s\cap(Z_j)_s$ est rare et fermé dans $(Z_i)_s$. Comme $Z_i$
est localement fermé dans $X$, il y a un sous-préschéma de $X$ ayant $Z_i$
pour espace sous-jacent (et encore noté $Z_i$), qui est de type fini sur
$S$ (\textbf{I}, 6.3.5). Posons
$U_i=Z_i-\bigcup_{j\neq i}(Z_i\cap Z_j)$, qui est ouvert dans $Z_i$ et tel que,
pour tout $s\in S$, $(U_i)_s$ soit ouvert partout dense dans $(Z_i)_s$~; en
outre, par construction, les $U_i$ sont deux à deux sans point commun. Comme
$(Z_i)_s$ est un $k(s)$-préschéma algébrique, l'ensemble des dimensions des
composantes irréductibles de $Z_s=\bigcup_i(Z_i)_s$ est le même que l'ensemble
des dimensions des composantes irréductibles de la réunion des $(U_i)_s$
(4.1.1.3), chacune de ces composantes étant déjà une composante irréductible
de l'un des $(U_i)_s$. On peut donc se borner au cas où $Z=U_i$~; en outre,
comme $Z_\eta$ est alors irréductible, il n'y a qu'une seule composante
irréductible de $Z$ qui rencontre $X_\eta$ $(\mathbf 0_{\mathrm I},2.1.8)$ et
l'on peut évidemment, en restreignant $Z$, supposer $Z$ irréductible. On est
finalement ramené à prouver le cas particulier suivant de (9.5.5)~:

\begin{corollary}[9.5.6]
\label{IV.9.5.6-fr}
Soient $S$ un préschéma noethérien et irréductible de point générique $\eta$,
$X$ un préschéma irréductible, $f:X\to S$ un morphisme dominant de type fini.
Alors il existe un voisinage $U$ de $\eta$ dans $S$ tel que, pour tout
$s\in U$, toutes les composantes irréductibles de $X_s$ soient de dimension
$n=\dim(X_\eta)$.
\end{corollary}

On peut évidemment se borner au cas où $S=\operatorname{Spec}(A)$ est affine,
$A$ étant donc noethérien~; remplaçant $f$ par $f_{\mathrm{red}}$, qui est de
type fini (1.5.4, (vii)), on peut

\oldpage[IV-3]{70}

supposer $A$ intègre et $X$ réduit, donc intègre puisqu'il est irréductible. On
sait (4.1.2) qu'il existe un ouvert $W$ dense dans $X_\eta$ et un
$k(\eta)$-morphisme \emph{fini et surjectif}
$h:W\to\mathbf V((k(\eta))^n)=\operatorname{Spec}(B')$, où
$B'=k(\eta)[T_1,\ldots,T_n]$ ($T_i$ indéterminées). Si $V$ est un ouvert de
$X$ tel que $V\cap X_\eta=W$, on sait (9.5.3) que pour $s$ voisin de $\eta$
dans $S$, $V_s$ est un ouvert dense dans $X_s$, et l'on peut par suite (4.1.1.3)
se borner au cas où $V=X$, $W=X_\eta$. Posons
$B=A[T_1,\ldots,T_n]$ et $Y=\operatorname{Spec}(B)=\mathbf V_A^n$, de sorte
que $\operatorname{Spec}(B')=Y_\eta$~; il résulte de (8.8.2, (i)) et
(8.10.5, (vi) et (x)), appliqués suivant la méthode de (8.1.2, \emph{a)}),
qu'en remplaçant au besoin $S$ par un voisinage de $\eta$, on peut supposer que
$h=g_\eta$, où $g:X\to Y$ est un morphisme fini surjectif~; autrement dit, on
a $X=\operatorname{Spec}(C)$, où $C$ est une $B$-algèbre finie et
l'homomorphisme $B\to C$ correspondant à $g$ est \emph{injectif}~; comme $C$
est un anneau intègre, $C$ est donc un $B$-module de type fini \emph{sans
torsion}, et $C_\eta=C\otimes_A k(\eta)$ est donc un module de type fini sans
torsion sur $B_\eta=B'=B\otimes_A k(\eta)=k(\eta)[T_1,\ldots,T_n]$ (étant un
module de fractions dont les dénominateurs sont dans
$A-\{0\}\subset B-\{0\}$). Il résulte donc de (9.4.8) qu'il y a un voisinage
$U$ de $\eta$ dans $S$ tel que pour $s\in S$,
$C_s=C\otimes_A k(s)$ soit un module de type fini \emph{sans torsion} sur
$B_s=B\otimes_A k(s)=k(s)[T_1,\ldots,T_n]$, et en particulier l'homomorphisme
$g_s:B_s\to C_s$ est injectif. Comme aucun élément de $B_s$ n'est diviseur de
0 dans $C_s$, pour tout idéal premier minimal $\mathfrak p_i$ de $C_s$ (dont
les éléments sont diviseurs de 0 dans $C_s$), on a nécessairement
$\mathfrak p_i\cap B_s=0$, donc l'homomorphisme canonique
$B_s\to C_s/\mathfrak p_i$ est injectif. On en déduit que pour chaque
composante irréductible $Z_i=\operatorname{Spec}(C_s/\mathfrak p_i)$ de $X_s$,
la restriction à $Z_i$ de $g$ est un morphisme fini et dominant $Z_i\to Y_s$,
donc \emph{surjectif} (\textbf{II}, 6.1.10). On conclut par suite de (4.1.2)
que $\dim Z_i=n$, ce qui achève la démonstration.

\begin{remark}[9.5.7]
\label{IV.9.5.7-fr}
On aura soin de noter que sous les hypothèses de (9.5.6) il peut se faire que
$X_s$ \emph{ne soit irréductible pour aucun} $s\neq\eta$ dans un voisinage de
$\eta$, autrement dit, la propriété «~$X$ est un $k$-préschéma algébrique
irréductible~» \emph{n'est pas constructible}. Prenons par exemple
$S=\operatorname{Spec}(k[T])$, où $k$ est un corps algébriquement clos, $T$ une
indéterminée~; on a donc $k(\eta)=K=k(T)$. Soit $L$ une extension finie
séparable de $K$ de degré $>1$, et soit $X$ la \emph{fermeture intégrale} de
$S$ dans $L$ (\textbf{II}, 6.3.4)~; on a donc $X=\operatorname{Spec}(B)$, où
$B$ est la fermeture intégrale de $k[T]$ dans $L$. On sait que $B$ est un
anneau de Dedekind, et que tous les idéaux maximaux de $k[T]$, sauf un nombre
fini, sont non ramifiés dans $B$~; comme en outre le corps résiduel de tout
idéal maximal de $B$ est nécessairement $k$ (puisque c'est une extension finie
de $k$), on voit que pour presque tous les idéaux maximaux $\mathfrak j_s$ de
$k[T]$, $B_s=B/\mathfrak j_sB$ est composé direct de $[L:K]$ corps isomorphes
à $k$, autrement dit $X_s$ n'est pas irréductible, bien que
$X_\eta=\operatorname{Spec}(L)$ le soit. Le même exemple montre que la
propriété «~$X$ est un $k$-préschéma algébrique intègre~» n'est pas
constructible. Enfin il en est de même de la propriété «~$X$ est un
$k$-préschéma algébrique réduit~». Il suffit pour le voir de prendre encore
$S=\operatorname{Spec}(k[T])$, où cette fois $k$ est un corps algébriquement
clos de caractéristique $p>0$, et pour $X$ la fermeture intégrale de $S$ dans
$L=K^{1/p}$ (où $K=k(T)$), de sorte que $X=\operatorname{Spec}(B)$ avec
$B=k[T^{1/p}]$ ($k$ étant parfait)~; tout idéal maximal de $k[T]$ est de la
forme $(T-\alpha)$ avec $\alpha\in k$~; l'unique idéal de $B$ au-dessus de
l'idéal $(T-\alpha)$ est l'idéal principal
$(T^{1/p}-\alpha^{1/p})$ et il est immédiat

\oldpage[IV-3]{71}

que l'anneau quotient $B/(T-\alpha)B$ admet par suite des éléments nilpotents~;
autrement dit, $X_s$ n'est réduit pour aucun $s\neq\eta$, tandis que
$X_\eta=\operatorname{Spec}(L)$ est intègre.

Nous verrons un peu plus loin (9.7) que l'on obtient par contre des propriétés
constructibles lorsqu'on considère les notions «~géométriques~»
correspondantes aux notions de préschéma irréductible, réduit ou intègre (4.5
et 4.6).
\end{remark}

\subsection{Constructibilité de certaines propriétés des morphismes}
\label{subsection:IV.9.6-fr}

\begin{proposition}[9.6.1]
\label{IV.9.6.1-fr}
Soient $X$, $Y$ deux $S$-préschémas de présentation finie, $f:X\to Y$ un
$S$-morphisme. Soit $E$ l'ensemble des $s\in S$ pour lesquels $f_s$ a l'une
des propriétés suivantes~: être~:
\begin{enumerate}[label=(\roman*),widest=xii]
  \item surjectif~;
  \item dominant~;
  \item séparé~;
  \item propre~;
  \item radiciel~;
  \item fini~;
  \item quasi-fini~;
  \item une immersion~;
  \item une immersion fermée~;
  \item une immersion ouverte~;
  \item un isomorphisme~;
  \item un monomorphisme.
\end{enumerate}
Alors $E$ est localement constructible dans $S$.
\end{proposition}

Les assertions de (i), (v) et (vii) ne sont insérées que pour mémoire, ayant
déjà été établies dans (9.3.2).

(ii)~: Notons d'abord que $f$ est de présentation finie (1.6.2, (v)), donc, en
vertu du théorème de Chevalley (1.8.4), $f(X)$ est localement constructible dans
$Y$. Par ailleurs, on a $f_s(X_s)=(f(X))_s$ (\textbf{I}, 3.6.1)~; l'ensemble
des $s$ tels que $f_s$ soit dominant est l'ensemble des $s$ tels que
$(f(X))_s$ soit dense dans $Y_s$~; la conclusion résulte donc dans ce cas de
(9.5.3).

(iii)~: Comme $f:X\to Y$ est de présentation finie, l'immersion diagonale
$\Delta_f:X\to X\times_YX$ est de présentation finie (1.6.2, (iv) et (v))~; il
résulte donc de (1.8.4.1) que $\Delta_f(X)=\Delta_X$ est une partie localement
constructible de $X\times_YX$. Dire que $f_s$ est séparé signifie (compte tenu
de (\textbf{I}, 5.3.4)) que $(\Delta_X)_s$ est fermé dans
$X_s\times_{Y_s}X_s=(X\times_YX)_s$~; la conclusion résulte donc ici de
(9.5.4).

(iv)~: Montrons que la propriété «~$X$ et $Y$ sont des préschémas algébriques
sur un corps $k$, et $f:X\to Y$ un $k$-morphisme propre~» est constructible. La
condition (9.2.1, (i)) est vérifiée en vertu de (2.7.1, (vii)). On peut donc se
borner au cas où $S$ est affine, noethérien et intègre, de point générique
$\eta$, et on doit prouver que $E$ ou $S-E$ est un voisinage de $\eta$.
Supposons d'abord que $\eta\in E$~; il résulte de (8.10.5, (xii)) appliqué
suivant la méthode de (8.1.2, \emph{a)}) qu'en remplaçant $S$ par un voisinage
de $\eta$, on peut

\oldpage[IV-3]{72}

supposer que $f$ est lui-même un morphisme propre~; on sait alors qu'il en est
de même de $f_s$ pour tout $s\in S$ (\textbf{II}, 5.4.2, (iii)). Supposons au
contraire que $\eta\in S-E$, et distinguons deux cas~:

\begin{enumerate}[label=\upshape{$\arabic{enumi}^\circ$}]
  \item Supposons que $f_\eta$ soit non séparé~; alors il résulte de (iii) que
  $f_s$ est non séparé (et \emph{a fortiori} non propre) dans un voisinage de
  $\eta$.

  \item Supposons $f_\eta$ séparé~; $Y$ est réunion finie d'ouverts affines
  $V_i$, et pour que $f_s$ soit propre, il faut et il suffit que chacune de ses
  restrictions $f_s^{-1}((V_i)_s)\to(V_i)_s$ le soit (\textbf{II}, 5.4.1)~;
  on peut donc se borner au cas où $Y$ est \emph{affine}, donc un schéma. Dire
  que $f_\eta$ n'est pas propre signifie (\textbf{II}, 5.6.3) qu'il existe un
  morphisme \emph{de type fini} $h:Z\to Y_\eta$ tel que le morphisme
  $(f_\eta)_{(Z)}:X_\eta\times_{Y_\eta}Z\to Z$ ne soit pas fermé. Comme $Y$ est
  un schéma, on déduit de (8.8.2, (i) et (ii)) (en restreignant au besoin $S$ à
  un voisinage de $\eta$) qu'il existe un morphisme de type fini $g:Y'\to Y$
  tel que $Z$ soit isomorphe à $Y'_\eta$, et que $g_\eta=h$~; si l'on pose
  $X'=X\times_YY'$, $f'=f_{(Y')}:X'\to Y'$, on a
  $(f_\eta)_{(Z)}=f'_\eta$, et par hypothèse il existe donc une partie fermée
  $M'$ de $X'_\eta$ telle que $f'_\eta(M')$ ne soit pas fermé dans $Y'_\eta$.
  Or $M'$ est la trace sur $X'_\eta$ d'une partie fermée $N'$ de $X'$~; comme
  $X'$ est noethérien et $f'$ de type fini, $f'(N')$ est constructible dans
  $Y'$ (1.8.4) et par hypothèse
  $(f'(N'))_\eta=f'_\eta(N'_\eta)=f'_\eta(M')$ n'est pas fermé dans
  $Y'_\eta$. On conclut alors de (9.5.4) qu'il existe un voisinage $U$ de
  $\eta$ dans $S$ tel que pour $s\in U$,
  $(f'(N'))_s=f'_s(N'_s)$ ne soit pas fermé dans $Y'_s$~; autrement dit, le
  morphisme $f'_s$ n'est pas fermé, et \emph{a fortiori} le morphisme $f_s$
  n'est pas propre.
\end{enumerate}

(vi) La propriété est conjonction des propriétés (iv) et (vii) (8.11.1), donc
la proposition résulte dans ce cas de ce qui a déjà été prouvé.

(ix) La vérification de la condition (9.2.1, (i)) résulte de (2.7.1, (xi)). On
peut donc se borner au cas où $S$ est affine, noethérien et intègre de point
générique $\eta$, et à prouver que $E$ ou $S-E$ est un voisinage de $\eta$. Si
$\eta\in E$, on peut (en remplaçant $S$ par un voisinage de $\eta$) supposer
que $f$ est une immersion fermée en vertu de (8.10.5, (iv)), et alors il est
clair que $f_s$ est une immersion fermée pour tout $s\in S$ (\textbf{I},
4.3.2). Supposons donc que $\eta\in S-E$ et distinguons deux cas~:

\begin{enumerate}[label=\upshape{$\arabic{enumi}^\circ$}]
  \item $f_\eta$ n'est pas un morphisme fini. Alors il résulte de (vi) que dans
  un voisinage de $\eta$, $f_s$ n'est pas fini, ni \emph{a fortiori} une
  immersion fermée.

  \item $f_\eta$ est fini~; alors, en vertu de (8.10.5, (x)), on peut supposer
  (en restreignant au besoin $S$) que $f$ lui-même est un morphisme \emph{fini}.
  Dans ce cas $\mathcal A=\mathcal A(X)=f_*(\mathcal O_X)$ est un
  $\mathcal O_Y$-Module cohérent (\textbf{II}, 6.1.3) et
  $f=\mathcal A(u)$, où $u:\mathcal O_Y\to\mathcal A$ est un homomorphisme de
  $\mathcal O_Y$-Algèbres (\textbf{II}, 1.1.2)~; comme $f_\eta$ n'est pas une
  immersion fermée par hypothèse, $u_\eta$ n'est pas surjectif (\textbf{II},
  1.4.10), donc (9.4.5) il existe un voisinage de $\eta$ dans lequel $u_s$
  n'est pas surjectif, et par suite (\textbf{I}, 4.2.3) $f_s$ n'est pas une
  immersion fermée.
\end{enumerate}

(viii) La vérification de la condition (9.2.1, (i)) se fait comme dans (ix), en
utilisant cette fois (2.7.1, (x)) et le fait que toute immersion d'un préschéma
noethérien dans un autre est quasi-compacte. On est donc ramené au cas où $S$
est affine, noethérien et intègre de point générique $\eta$, et à prouver que
$E$ ou $S-E$ est un voisinage de $\eta$. Si $\eta\in E$, on conclut comme dans
(ix) à l'aide de (8.10.5, (ii)). Si $\eta\in S-E$, on distingue de nouveau
deux cas~:

\begin{enumerate}[label=\upshape{$\arabic{enumi}^\circ$}]
  \item $f_\eta(X_\eta)$ n'est pas une partie localement fermée de $Y_\eta$.
  Comme $f(X)$ est constructible dans $Y$ (1.8.4) et que
  $f_s(X_s)=(f(X))_s$, on déduit de (9.5.4) que pour $s$ voisin de $\eta$,
  $f_s(X_s)$ n'est pas localement fermé dans $Y_s$, et \emph{a fortiori}
  $f_s$ n'est pas une immersion.

  \item\oldpage[IV-3]{73} $f_\eta(X_\eta)$ est localement fermé dans $Y_\eta$.
  Comme $f(X)$ est constructible dans $Y$ (1.8.4) et qu'il en est de même de
  $\overline{f(X)}$ puisque $Y$ est noethérien, il résulte de (8.3.11) qu'en
  restreignant au besoin $S$, on peut supposer que $f(X)$ est localement fermé
  dans $Y$. Il y a alors un ouvert $V$ de $Y$ contenant $f(X)$ et dans lequel
  $f(X)$ est fermé. Comme $Y$ est noethérien, $V$ est de type fini sur $S$, et
  en remplaçant $Y$ par $V$, on peut donc se ramener au cas où $f(X)$ est
  \emph{fermé} dans $Y$. Mais alors $f_s(X_s)=(f(X))_s$ est fermé dans $Y_s$
  pour tout $s\in S$, et dire que $f_s$ est une immersion équivaut à dire que
  $f_s$ est une immersion fermée~; on est donc ramené à ce qui a été démontré
  dans (ix).
\end{enumerate}

(x) Utilisant cette fois (2.7.1, (ix)) et (8.10.5, (iii)), on est ramené au cas
où $S$ est affine, noethérien, intègre de point générique $\eta$, et où
$\eta\in S-E$. Distinguons trois cas~:

\begin{enumerate}[label=\upshape{$\arabic{enumi}^\circ$}]
  \item $f_\eta(X_\eta)$ n'est pas ouvert dans $Y_\eta$. Comme $f(X)$ est
  constructible dans $Y$ (1.8.4), on déduit de (9.5.4) que $f_s(X_s)$ n'est
  pas ouvert dans $Y_s$ pour $s$ voisin de $\eta$, et \emph{a fortiori} $f_s$
  n'est pas une immersion ouverte.

  \item $f_\eta(X_\eta)$ est ouvert dans $Y_\eta$, mais $f_\eta$ n'est pas une
  immersion. Il résulte alors de (viii) que pour $s$ voisin de $\eta$ dans $S$,
  $f_s$ n'est pas une immersion, ni \emph{a fortiori} une immersion ouverte.

  \item $f_\eta(X_\eta)$ est ouvert dans $Y_\eta$ et $f_\eta$ est une
  immersion. Comme $f(X)$ est constructible dans $Y$, il résulte de (8.3.11)
  qu'en restreignant au besoin $S$, on peut déjà supposer que $f(X)$ est ouvert
  dans $Y$. Comme $Y$ est noethérien, le sous-préschéma induit sur $f(X)$ est
  de type fini sur $S$, donc on peut se ramener au cas où $f$ est \emph{surjectif}
  en remplaçant $Y$ par $f(X)$. Par hypothèse, $f_\eta$ est une immersion
  fermée, donc on peut, comme dans (ix), supposer que $f$ est une immersion
  \emph{fermée}, et par suite que $X$ est un sous-préschéma fermé de $Y$ défini
  par un Idéal cohérent $\mathcal J$ de $\mathcal O_Y$. Par hypothèse
  $f_\eta$ n'est pas un isomorphisme, donc $\mathcal J_\eta\neq0$~; on en
  conclut (9.4.5) que dans un voisinage de $\eta$, on a $\mathcal J_s\neq0$,
  et par suite l'immersion fermée surjective $f_s$ n'est pas ouverte.
\end{enumerate}

(xi) La propriété est conjonction des propriétés (i) et (x) et résulte donc de
ce qui a été démontré.

(xii) En vertu de (\textbf{I}, 5.3.8), dire que $f_s$ est un monomorphisme
signifie que
$\Delta_{f_s}=(\Delta_f)_s:X_s\to X_s\times_{Y_s}X_s=(X\times_YX)_s$ est un
isomorphisme, et comme $X\times_YX$ est un $S$-préschéma de présentation finie
(1.6.2, (iv)), la conclusion résulte de (xi).

\begin{proposition}[9.6.2]
\label{IV.9.6.2-fr}
Soient $X$, $Y$ deux $S$-préschémas de présentation finie, $f:X\to Y$ un
$S$-morphisme.
\begin{enumerate}[label=\upshape{I)}]
  \item Soit $E$ l'ensemble des $s\in S$ pour lesquels $f_s$ a l'une des
  propriétés suivantes~: être~:
  \begin{enumerate}[label=(\roman*)]
    \item affine~;
    \item quasi-affine~;
    \item projectif~;
    \item quasi-projectif.
  \end{enumerate}
  Alors $E$ est ind-constructible dans $S$.
\end{enumerate}

\oldpage[IV-3]{74}

\begin{enumerate}[label=\upshape{II)}]
  \item Soit $\mathcal L$ un $\mathcal O_X$-Module inversible. Alors l'ensemble
  $E'$ des $s\in S$ tels que $\mathcal L_s$ soit un $\mathcal O_{X_s}$-Module
  ample (resp. très ample) relativement à $f_s$, est ind-constructible dans
  $S$.
\end{enumerate}
\end{proposition}

I) Vérifions les conditions (i) et (ii) de (9.2.1). En ce qui concerne la
condition (9.2.1, (i)), elle résulte, pour les propriétés (i) et (ii), de
(2.7.1, (xiii) et (xiv))~; pour les propriétés (iii) et (iv), elle résulte de
(9.1.5). Vérifions ensuite la condition (9.2.1, (ii)), en supposant donc $S$
noethérien, intègre et de point générique $\eta$, et que
$f_\eta:X_\eta\to Y_\eta$ ait l'une des propriétés (i) à (iv) de l'énoncé.
Appliquant (8.10.5, (viii), (ix), (xiii) et (xiv)) suivant la méthode de
(8.1.2, \emph{a)}), on voit d'abord qu'il existe un voisinage ouvert $U$ de
$\eta$ tel que, si $V$ et $W$ sont les images réciproques de $U$ dans $X$ et
$Y$ respectivement par les morphismes structuraux, la restriction $V\to W$ de
$f$ a celle des propriétés (i) à (iv) que l'on considère. La conclusion
résulte alors de ce que ces propriétés sont toutes stables par changement de
base.

II) On procède de la même façon. La condition (9.2.1, (i)) résulte cette fois
de (2.7.2). Pour la condition (9.2.1, (ii)), avec les mêmes notations que dans
I), il résulte de (8.10.5.2) que pour un voisinage $U$ de $\eta$ dans $S$, la
restriction $\mathcal L|V$ est ample (resp. très ample) relativement à la
restriction $V\to W$ de $f$. La conclusion résulte encore de la stabilité des
propriétés considérées par changement de base (\textbf{II}, 4.6.13 et
4.4.10).

On peut améliorer (9.6.2, II) sous certaines conditions~:

\begin{proposition}[9.6.3]
\label{IV.9.6.3-fr}
Soient $X$, $Y$ deux $S$-préschémas de présentation finie, $f:X\to Y$ un
$S$-morphisme propre, $\mathcal L$ un $\mathcal O_X$-Module inversible. Alors
l'ensemble $E$ (resp. $E'$) des $s\in S$ tels que $\mathcal L_s$ soit un
$\mathcal O_{X_s}$-Module ample (resp. très ample) relativement à
$f_s:X_s\to Y_s$ est localement constructible dans $S$.
\end{proposition}

Soient $k$ un corps, $Z$, $T$ deux préschémas algébriques sur $k$, $\mathcal M$
un $\mathcal O_Z$-Module inversible, $g:Z\to T$ un $k$-morphisme de type fini~;
alors, si $\mathbf P(Z,T,\mathcal M,g,k)$ désigne la relation~: «~$\mathcal M$ est ample
(resp. très ample) relativement à $g$~», on a déjà remarqué dans (9.6.2) que
$\mathbf P(Z,T,\mathcal M,g,k)$ vérifie la condition (9.2.1, (i)) en vertu de (2.7.2).
On sait déjà d'autre part que $E$ et $E'$ sont ind-constructibles. Il reste donc
à voir que si $S$ est noethérien, intègre, de point générique $\eta$ et si
$\eta\in S-E$ (resp. $\eta\in S-E'$), alors $S-E$ (resp. $S-E'$) contient un
voisinage de $\eta$. Nous considérerons séparément le cas de $E'$ et celui de
$E$.

I) \emph{Cas de $E'$.} --- Notons que puisque $f$ est séparé et $Y$
quasi-compact, il existe un entier $h$ tel que, pour $q>h$, on ait
$R^qf_*(\mathcal L)=0$ (\textbf{III}, 1.4.12)~; d'autre part, puisque $f$ est
propre et $Y$ noethérien, les $R^qf_*(\mathcal L)$ sont tous des
$\mathcal O_Y$-Modules cohérents (\textbf{III}, 3.2.1)~; comme ils sont nuls
sauf pour un nombre fini de valeurs de $q$, le théorème de platitude générique
(6.9.1) montre qu'en restreignant $S$ à un voisinage de $\eta$, on peut
supposer que $\mathcal L$ et les $R^qf_*(\mathcal L)$ sont tous $S$-plats. On
conclut alors de (\textbf{III}, 6.9.9) que l'homomorphisme canonique
\begin{equation}
\label{IV.9.6.3.1-fr}
 f_*(\mathcal L)\otimes_{\mathcal O_S} k(s)
 \longrightarrow (f_s)_*(\mathcal L_s)
 \tag{9.6.3.1}
\end{equation}
est un \emph{isomorphisme}.

\oldpage[IV-3]{75}

Cela étant, il résulte de (\textbf{II}, 4.4.4) que dire que $\mathcal L_s$
n'est pas très ample relativement à $f_s$ signifie que~: \emph{ou bien}
l'homomorphisme canonique
$(f_s)^*((f_s)_*(\mathcal L_s))\to\mathcal L_s$ n'est pas surjectif~; \emph{ou
bien} l'homomorphisme précédent est surjectif et le morphisme canonique
$r:X_s\to\mathbf P((f_s)_*(\mathcal L_s))$ n'est pas une immersion. Compte
tenu de l'isomorphisme (9.6.3.1), ces conditions s'écrivent respectivement sous
la forme~: $1^\circ$ l'homomorphisme canonique
$(f^*(f_*(\mathcal L)))_s\to\mathcal L_s$ n'est pas surjectif~; $2^\circ$
l'homomorphisme précédent est surjectif et le morphisme canonique
$X_s\to\mathbf P((f_*(\mathcal L))_s)$ n'est pas une immersion.

Supposons d'abord que l'homomorphisme canonique
$(f^*(f_*(\mathcal L)))_\eta\to\mathcal L_\eta$ ne soit pas surjectif. Comme
$f_*(\mathcal L)$ est cohérent, il en est de même de $f^*(f_*(\mathcal L))$,
et alors il résulte de (9.4.5) que pour tout $s$ dans un voisinage de $\eta$,
l'homomorphisme $(f^*(f_*(\mathcal L)))_s\to\mathcal L_s$ n'est pas
surjectif, ce qui prouve dans ce cas que $S-E'$ est un voisinage de $\eta$.

Supposons en second lieu que l'homomorphisme canonique
$(f^*(f_*(\mathcal L)))_\eta\to\mathcal L_\eta$ soit surjectif mais que le
morphisme $X_\eta\to\mathbf P((f_*(\mathcal L))_\eta)$ ne soit pas une
immersion. Alors le même raisonnement que ci-dessus montre d'abord que pour
tout $s$ assez voisin de $\eta$, l'homomorphisme
$(f^*(f_*(\mathcal L)))_s\to\mathcal L_s$ est surjectif~; d'autre part, en
vertu de (9.6.1, (viii)), pour $s$ assez voisin de $\eta$, le morphisme
$X_s\to\mathbf P((f_*(\mathcal L))_s)$ n'est pas une immersion, ce qui achève
la démonstration dans le cas de $E'$.

II) \emph{Cas de $E$.} --- Considérons d'abord le cas particulier où $Y=S$~:

\begin{corollary}[9.6.4]
\label{IV.9.6.4-fr}
Soit $f:X\to S$ un morphisme propre de présentation finie, $\mathcal L$ un
$\mathcal O_X$-Module inversible. Alors, l'ensemble $E$ des $s\in S$ tels que
$\mathcal L_s$ soit ample (relativement à $f_s$) est ouvert dans $S$, et
$\mathcal L|f^{-1}(E)$ est ample relativement à la restriction
$f^{-1}(E)\to E$ de $f$.
\end{corollary}

Comme la condition (9.2.1, (i)) est vérifiée par la propriété $\mathbf P$ définie plus
haut, le résultat de (9.2.2, (iv)) et le raisonnement de (9.2.3) montrent qu'on
peut se borner au cas où $S$ est \emph{noethérien}~; mais alors le résultat
découle de (\textbf{III}, 4.7.1) et de la stabilité de l'amplitude par
changement de base (\textbf{II}, 4.6.13).

\begin{corollary}[9.6.5]
\label{IV.9.6.5-fr}
Sous les hypothèses de (9.6.4), pour que $\mathcal L$ soit ample relativement
à $f$, il faut et il suffit que, pour tout $s\in S$, $\mathcal L_s$ soit ample
relativement à $f_s$.
\end{corollary}

\phantomsection\label{IV.9.6.6-fr}%
\textbf{(9.6.6) Fin de la démonstration de (9.6.3).} --- Revenons au cas
général, $S$ étant noethérien, intègre et de point générique
$\eta\in S-E$. Comme $f$ est propre, il résulte de (9.6.4) que l'ensemble $V$
des $y\in Y$ tels que $\mathcal L_y$ soit un $\mathcal O_{X_y}$-Module ample,
relativement au morphisme $f_y:X_y\to\operatorname{Spec}(k(y))$ est ouvert,
donc $F=Y-V$ est fermé dans $Y$. Cela étant, comme $f_s$ est propre et que,
pour tout $y\in Y$ au-dessus de $s\in S$, on a
$\mathcal L_y=(\mathcal L_s)_y$, il résulte de (9.6.5) que pour que $s$
appartienne à $S-E$, il faut et il suffit que l'on ait $F_s\neq\varnothing$.
Mais comme $F$ est fermé dans $Y$ et que $F_\eta\neq\varnothing$, il résulte
de (9.5.1) que l'ensemble des $s\in S$ tels que $F_s\neq\varnothing$ est un
voisinage de $\eta$ dans $S$. C.Q.F.D.

\begin{remarks}[9.6.7]
\label{IV.9.6.7-fr}
(i) Pour chacune des propriétés $\mathbf P$ considérées dans (9.6.1) la proposition
(9.3.3) est applicable, et ces propriétés (pour les morphismes $f_s$) sont donc
«~stables~» par passage d'une limite projective d'un système projectif
essentiellement affine (8.13.4) de préschémas $X_\lambda$ à l'un convenable
d'entre eux.

(ii) Soit $Z$ une partie \emph{localement constructible} de $X$ telle que, pour
tout $s\in S$, $Z\cap X_s$ soit \emph{ouvert} dans $X_s$, et désignons par
$Z_s$ le sous-préschéma de $X_s$ induit sur l'ouvert $Z\cap X_s$.

\oldpage[IV-3]{76}

Alors, dans les propositions (9.6.1) et (9.6.2, (I)), on peut partout remplacer
$f_s$ par sa restriction $f_s|Z_s:Z_s\to Y_s$ sans changer les conclusions. En
effet, la vérification de (9.2.1, (i)) se fait comme dans (9.6.1) et (9.6.2).
D'autre part, dans la réduction au cas où $S$ est noethérien, faite dans
(9.2.3), si $Z=q^{-1}(Z_0)$, où $q:X\to X_0$ est la projection canonique et
$Z_0$ une partie constructible de $X_0$ (8.3.11), le fait que
$(Z_0)_{s_0}$ soit ouvert dans $(X_0)_{s_0}$, pour $s_0=p(s)$, découle de
(2.4.10) et de ce que la projection $X_s\to(X_0)_{s_0}$ est surjective. On est
donc ramené à vérifier (9.2.1, (ii)) sous les nouvelles hypothèses. Or, comme
$Z_\eta$ est ouvert dans $X_\eta$, il existe un ouvert $Z'\subset X$ tel que
$Z_\eta=Z'\cap X_\eta$~; comme $X$ est alors noethérien, $Z'$ est
constructible, et il en est de même de $Z$ par hypothèse~; on conclut donc de
(9.5.2) et de ($\mathbf 0_{\mathrm{III}}$, 9.2.2) qu'il y a un voisinage $U$
de $\eta$ dans $S$ tel que $Z_s=Z'_s$ pour $s\in U$. Remplaçant $S$ par $U$,
on peut donc se borner au cas où $Z$ est ouvert dans $X$, et alors on est
ramené à ce qui a été prouvé dans (9.6.1) et (9.6.2).
\end{remarks}

\subsection{Constructibilité des propriétés de séparabilité, d'irréductibilité
géométrique et de connexité géométrique}
\label{subsection:IV.9.7-fr}

\begin{lemma}[9.7.1]
\label{IV.9.7.1-fr}
Soient $S$ un préschéma irréductible de point générique $\eta$, $f:X\to S$
un morphisme de présentation finie. Si $X_\eta$ a $n$ composantes irréductibles
(resp. $n'$ composantes connexes), il existe un voisinage ouvert $U$ de $\eta$
dans $S$ tel que pour tout $s\in U$, $X_s$ ait au moins $n$ composantes
irréductibles (resp. $n'$ composantes connexes).
\end{lemma}

On peut se borner au cas où $S$ est affine, donc $X$ quasi-compact et
quasi-séparé. Soient $V_i$ ($1\leqslant i\leqslant n$) les intérieurs des
composantes irréductibles de $X_\eta$, qui sont deux à deux disjoints et
quasi-compacts, et dont la réunion est dense dans $X_\eta$~; en vertu de
(8.2.11), appliqué suivant la méthode de (8.1.2, \emph{a)}), il existe pour
chaque $i$ un ouvert quasi-compact $W_i$ dans $X$ tel que
$W_i\cap X_\eta=V_i$~; comme $X$ est quasi-séparé, les intersections
$W_i\cap W_j$ sont quasi-compactes (1.2.7), donc, en remplaçant $S$
par un voisinage de $\eta$, on peut supposer que $W_i\cap W_j=\varnothing$ pour
$i\neq j$ en vertu de (8.3.3). En outre, comme les $W_i$ sont constructibles,
il y a un voisinage $U$ de $\eta$ dans $S$ tel que pour tout $s\in U$ les
$(W_i)_s$ soient non vides (9.5.1 et 9.2.3) et que la réunion des $(W_i)_s$ soit
\emph{dense} dans $X_s$ (9.5.3 et 9.2.3). Cela étant, les composantes
irréductibles (en nombre fini) des $(W_i)_s$ sont aussi les composantes
irréductibles de la réunion des $(W_i)_s$ $(\mathbf 0_{\mathrm I},2.1.7)$, donc
les adhérences dans $X_s$ de ces composantes sont les composantes
irréductibles de $X_s$ $(\mathbf 0_{\mathrm I},2.1.6)$ et leur nombre est
évidemment $\geqslant n$.

Soient maintenant $C_j$ ($1\leqslant j\leqslant n'$) les composantes connexes
de $X_\eta$, qui sont des ouverts quasi-compacts de $X_\eta$, deux à deux
disjoints. Si on remplace les $V_i$ par les $C_j$ dans le raisonnement
précédent, on voit (utilisant deux fois (8.3.3)) qu'on peut supposer $X$ réunion
de $n'$ ouverts quasi-compacts $M_j$ ($1\leqslant j\leqslant n'$) deux à deux
disjoints et que, dans un voisinage de $\eta$, les $(M_j)_s$ soient non vides.
Comme la réunion des $(M_j)_s$ est $X_s$, les composantes connexes des $(M_j)_s$
sont les composantes connexes de $X_s$, donc leur nombre est $\geqslant n'$.

\oldpage[IV-3]{77}

\begin{lemma}[9.7.2]
\label{IV.9.7.2-fr}
Soient $S$ un préschéma irréductible de point générique $\eta$, $f:X\to S$
un morphisme de présentation finie. Si $X_\eta$ n'est pas réduit, il existe un
voisinage ouvert $U$ de $\eta$ dans $S$ tel que, pour tout $s\in U$, $X_s$ ne
soit pas réduit.
\end{lemma}

En effet, soit $\mathcal N$ le Nilradical de $\mathcal O_X$~; il résulte de
(8.2.13), appliqué suivant la méthode de (8.1.2, \emph{a)}) que $\mathcal N_\eta$
est le nilradical de $\mathcal O_{X_\eta}$, et l'hypothèse est que
$\mathcal N_\eta\neq0$. On conclut de (9.4.5) et (9.2.3) qu'il y a un voisinage
$U$ de $\eta$ tel que, pour tout $s\in U$, $\mathcal N_s$ s'identifie à un Idéal
de $\mathcal O_{X_s}$ et $\mathcal N_s\neq0$~; comme $\mathcal N_s$ est
évidemment contenu dans le Nilradical de $\mathcal O_{X_s}$, on voit que $X_s$
n'est pas réduit pour $s\in U$.

\phantomsection\label{IV.9.7.3-fr}%
\textbf{(9.7.3)} Étant donné un polynôme
$F\in A[\mathrm T_1,\ldots,\mathrm T_n]$, où $A$ est un anneau et les
$\mathrm T_i$ des indéterminées, pour tout homomorphisme d'anneaux
$\rho:A\to B$, nous désignerons par $F_{(\rho)}$ ou $F_{(B)}$ le polynôme de
$B[\mathrm T_1,\ldots,\mathrm T_n]$ obtenu en remplaçant chaque coefficient de
$F$ par son image par $\rho$. Si $k$ est un corps,
$F\in k[\mathrm T_1,\ldots,\mathrm T_n]$ un polynôme non constant et
$X=\operatorname{Spec}(k[\mathrm T_1,\ldots,\mathrm T_n]/(F))$, dire que $X$
est intègre (ou que l'idéal $(F)$ est premier) signifie que $F$ est
\emph{irréductible} (c'est-à-dire que dans toute factorisation $F=F_1F_2$ en
polynômes de $k[\mathrm T_1,\ldots,\mathrm T_n]$, $F_1$ ou $F_2$ est de degré
$0$)~; cela résulte de ce que l'anneau $k[\mathrm T_1,\ldots,\mathrm T_n]$ est
factoriel. On en déduit aussitôt ((4.6.2) et (4.5.2))~:

\begin{lemma}[9.7.4]
\label{IV.9.7.4-fr}
Soient $k$ un corps, $\Omega$ une extension algébriquement close de $k$, $F$ un
polynôme non constant de $k[\mathrm T_1,\ldots,\mathrm T_n]$. Les conditions
suivantes sont équivalentes~:
\begin{enumerate}[label=\emph{\alph*)}]
  \item $X=\operatorname{Spec}(k[\mathrm T_1,\ldots,\mathrm T_n]/(F))$ est
  géométriquement intègre.
  \item $F_{(K)}$ est irréductible pour toute extension $K$ de $k$.
  \item $F_{(\Omega)}$ est irréductible.
\end{enumerate}
Nous dirons dans ce cas que $F$ est \emph{géométriquement irréductible}.
\end{lemma}

\begin{lemma}[9.7.5]
\label{IV.9.7.5-fr}
Soient $A$ un anneau intègre, $K$ son corps des fractions, $F$ un polynôme non
constant de $A[\mathrm T_1,\ldots,\mathrm T_n]$ de degré $d$ tel que $F_{(K)}$
soit géométriquement irréductible. Alors il existe $f\neq0$ dans $A$ tel que
pour tout $x\in D(f)$, $F_{(\boldsymbol{k}(x))}$ soit géométriquement
irréductible.
\end{lemma}

Posons $F=\sum_\alpha c_\alpha\mathrm T^\alpha$ avec comme d'ordinaire
$c_\alpha=c_{\alpha_1\alpha_2\ldots\alpha_n}$,
$\mathrm T^\alpha=\mathrm T_1^{\alpha_1}\mathrm T_2^{\alpha_2}\cdots
\mathrm T_n^{\alpha_n}$,
$|\alpha|=\alpha_1+\alpha_2+\cdots+\alpha_n\leqslant d$ (un au moins des
$c_\alpha$ tels que $|\alpha|=d$ n'étant pas nul). Comme les $c_\alpha$ non
nuls sont inversibles dans $K$, on peut supposer, en remplaçant au besoin $A$
par un anneau $A_g$ ($g\neq0$ dans $A$) que les $c_\alpha$ non nuls sont
inversibles dans $A$. Il en résulte que pour tout $x\in\operatorname{Spec}(A)$,
$F_{(\boldsymbol{k}(x))}$ est \emph{de degré $d$}.

Nous démontrerons d'abord un lemme préliminaire.

\begin{lemma}[9.7.5.1]
\label{IV.9.7.5.1-fr}
Soient $A$ un anneau intègre, $S=\operatorname{Spec}(A)$, $\eta$ le point
générique de $S$, $B$ l'anneau de polynômes
$A[\mathrm T_1,\ldots,\mathrm T_m]$, $(P_i)_{i\in I}$ une famille finie
d'éléments de $B$, $\mathfrak a$ l'idéal de $B$ engendré par les $P_i$. Pour
tout $s\in S$, soit $\mathfrak a_s$ l'idéal de
$\boldsymbol{k}(s)[\mathrm T_1,\ldots,\mathrm T_m]$ engendré par les polynômes
$(P_i)_{(\boldsymbol{k}(s))}$ ($i\in I$). Alors, si
$V(\mathfrak a_\eta)=\varnothing$, il existe un voisinage $U$ de $\eta$ dans
$S$ tel que $V(\mathfrak a_s)=\varnothing$ pour tout $s\in U$.
\end{lemma}

En effet, soit $Y=\operatorname{Spec}(B)$, et soit $Z$ la partie fermée
$V(\mathfrak a)$ de $Y$~; comme $\mathfrak a$ est un idéal de type fini, $Z$
est constructible dans $Y$ $(\mathbf 0_{\mathrm{III}},9.1.5)$ et avec les
notations introduites en (9.4.1), on a $V(\mathfrak a_s)=Z_s$ pour tout
$s\in S$~; comme le morphisme structural $Y\to S$ est de présentation finie,
la conclusion du lemme résulte de (9.5.1) et (9.2.3).

\oldpage[IV-3]{78}

Soit $(p,q)$ un couple d'entiers $>0$ tels que $p+q=d$~; introduisons des
indéterminées $\mathrm T'_\beta$, $\mathrm T''_\gamma$ pour tous les systèmes
d'entiers $\beta$, $\gamma$ tels que $|\beta|\leqslant p$ et
$|\gamma|\leqslant q$~; pour tout système d'entiers $\alpha$ tel que
$|\alpha|\leqslant d$, considérons le polynôme de
\[
B=A[\mathrm T'_\beta,\mathrm T''_\gamma]_{|\beta|\leqslant p,
  |\gamma|\leqslant q}:
\qquad
P_\alpha(\mathrm T'_\beta,\mathrm T''_\gamma)
=\sum_{\beta+\gamma=\alpha}\mathrm T'_\beta\mathrm T''_\gamma-c_\alpha.
\]
Soit $\Omega$ une clôture algébrique de $K$~; dire qu'il existe deux polynômes
$F_1$, $F_2$ de $\Omega[\mathrm T_1,\ldots,\mathrm T_n]$, de degrés respectifs
$p$ et $q$, tels que $F_1F_2=F_{(\Omega)}$, signifie que le système d'équations
$P_\alpha(t'_\beta,t''_\gamma)=0$ ($|\alpha|\leqslant d$) admet une solution
$(t'_\beta,t''_\gamma)$ ($|\beta|\leqslant p$, $|\gamma|\leqslant q$) formée
d'éléments de $\Omega$. Soit $\mathfrak a$ l'idéal de $B$ engendré par les
$P_\alpha$~; l'interprétation précédente, et le théorème des zéros de Hilbert,
montrent que l'hypothèse sur $F_{(K)}$ implique que
$V(\mathfrak a_\eta)=\varnothing$, en désignant par $\eta$ le point générique de
$\operatorname{Spec}(A)$~; le lemme (9.7.5.1) prouve donc que dans un voisinage
de $\eta$, on a $V(\mathfrak a_x)=0$, et par suite pour ces valeurs de $x$,
$F_{(\boldsymbol{k}(x))}$ n'admet pas de factorisation
$F_{(\boldsymbol{k}(x))}=G_1G_2$, où $G_1$, $G_2$ sont des polynômes de degrés
respectifs $p$ et $q$, dont les coefficients sont dans une clôture algébrique
de $\boldsymbol{k}(x)$. Il suffit d'appliquer ce résultat à tous les couples
d'entiers $(p,q)$ tels que $p>0$, $q>0$ et $p+q=d$, pour obtenir la conclusion
du lemme (9.7.5).

\begin{proposition}[9.7.6]
\label{IV.9.7.6-fr}
Soient $S$ un préschéma noethérien intègre de point générique $\eta$, $f:X\to S$
un morphisme de type fini, $\mathcal F$ un $\mathcal O_X$-Module cohérent. Si
$\mathcal F_\eta$ est sans cycle premier associé immergé, il existe un voisinage
$U$ de $\eta$ dans $S$ tel que, pour tout $s\in U$, $\mathcal F_s$ soit sans cycle
premier associé immergé.
\end{proposition}

Comme $X_\eta$ est noethérien, il existe un homomorphisme injectif
$v:\mathcal F_\eta\to\bigoplus_{i=1}^m\mathcal G_i$, où chacun des $\mathcal G_i$ est
irrédondant et $\operatorname{Ass}(\mathcal F_\eta)$ est l'ensemble des $x_i$, où
$\{x_i\}=\operatorname{Ass}(\mathcal G_i)$ (3.2.6)~; si $Z_i$ est l'adhérence de
$\{x_i\}$ dans $X_\eta$, on a $Z_i=\operatorname{Supp}(\mathcal G_i)$ (3.1.4), et
par hypothèse les $Z_i$ sont les composantes irréductibles de
$Z=\operatorname{Supp}(\mathcal F_\eta)$. Il résulte de (8.5.2, (i) et (ii)) (en
restreignant au besoin $S$ à un voisinage de $\eta$) et (8.5.8) qu'il existe des
$\mathcal O_X$-Modules cohérents $\mathcal F_i$ et un homomorphisme injectif
$u:\mathcal F\to\bigoplus_{i=1}^m\mathcal F_i$ tels que $(\mathcal F_i)_\eta=\mathcal G_i$ pour
tout $i$ et $v=u_\eta$. Si $Y_i=\operatorname{Supp}(\mathcal F_i)$, on a
$(Y_i)_s=\operatorname{Supp}((\mathcal F_i)_s)$ pour tout $s\in S$
(\textbf{I}, 9.1.13) et en particulier $(Y_i)_\eta=Z_i$ pour tout $i$.
L'hypothèse entraîne que pour $i\neq j$, $Z_i-(Z_i\cap Z_j)$ est ouvert dense
dans $Z_i$~; on déduit donc de (9.5.3) et (9.5.4) que dans un voisinage de
$\eta$, $(Y_i)_s-((Y_i)_s\cap(Y_j)_s)$ est ouvert et dense dans $(Y_i)_s$.
Supposons que l'on ait prouvé que chacun des $(\mathcal F_i)_s$ est sans cycle
premier associé pour $s\in U$. Alors les éléments de
$\operatorname{Ass}((\mathcal F_i)_s)$ sont les points maximaux de $(Y_i)_s$~;
aucun d'eux ne peut donc appartenir à un $(Y_j)_s$ pour $i\neq j$, et la
proposition sera démontrée. On peut donc supposer que $\mathcal F_\eta$ est
\emph{irrédondant}~; en outre, on peut se borner au cas où
$S=\operatorname{Spec}(A)$ est affine~; $X$ est alors réunion d'un nombre fini
d'ouverts affines $V_j$, et si $V_j\cap\operatorname{Supp}(\mathcal F_\eta)
=\varnothing$, on aura aussi $V_j\cap\operatorname{Supp}(\mathcal F_s)=\varnothing$
pour $s$ voisin de $\eta$ (9.5.1), donc on peut se borner au cas où
$X=\operatorname{Spec}(B)$ est aussi affine et où le morphisme $f:X\to S$ est
dominant~; $A$ est donc un anneau noethérien intègre, \emph{sous-anneau} d'une
$A$-algèbre de type fini $B$, et $\mathcal F=\widetilde M$, où $M$ est un
$B$-module de type fini~; par hypothèse, si $K$ est le corps des fractions de
$A$, le $B_{(K)}$-module $M_{(K)}$ est irrédondant. Soit $\mathfrak q$ l'unique

\oldpage[IV-3]{79}

élément de $\operatorname{Ass}(M_{(K)})$, et soit $\mathfrak p$ l'idéal premier
de $B$ image réciproque de $\mathfrak q$ par l'application canonique
$B\to B_{(K)}$. On sait (5.11.1.1) qu'il existe une filtration finie
$(N_h)_{0\leqslant h\leqslant m}$ de $M_{(K)}$ telle que $N_0=M_{(K)}$,
$N_m=0$, et que $N_h/N_{h+1}$ soit isomorphe à un sous-$B_{(K)}$-module
$\neq0$ de $B_{(K)}/\mathfrak q$. Soit $M_h$ l'image réciproque de $N_h$ par
l'application canonique $M\to M_{(K)}$. Il résulte de (9.4.4) que pour $s$
assez voisin de $\eta$, $(M_h)_{(\boldsymbol{k}(s))}$ s'identifie à un
sous-$B_{(\boldsymbol{k}(s))}$-module de $M_{(\boldsymbol{k}(s))}$, et le
quotient
$(M_h)_{(\boldsymbol{k}(s))}/(M_{h+1})_{(\boldsymbol{k}(s))}$ à un
sous-$B_{(\boldsymbol{k}(s))}$-module $\neq0$ de
$B_{(\boldsymbol{k}(s))}/\mathfrak p_{(\boldsymbol{k}(s))}
=(B/\mathfrak p)_{(\boldsymbol{k}(s))}$. Compte tenu de (3.1.7), on voit qu'on
est ramené à prouver, avec les mêmes notations, que si $B$ est \emph{intègre},
alors $B_{(\boldsymbol{k}(s))}$ n'a pas d'idéaux premiers associés immergés
pour $s$ voisin de $\eta$. Or, en remplaçant au besoin $A$ par $A_g$ et $B$ par
$B_g$ (où $g$ est un élément $\neq0$ de $A$), on peut supposer que $B$ contient
une $A$-algèbre de polynômes $C=A[\mathrm T_1,\ldots,\mathrm T_n]$ telle que
$B$ soit un $C$-module de type fini (Bourbaki, \emph{Alg. comm.}, chap. V,
\S~3, n$^\circ$~1, cor. 1 du th. 1). Comme $B_{(K)}$ est un $C_{(K)}$-module
sans torsion, on peut appliquer de nouveau le raisonnement fait plus haut en
remplaçant $B$, $M$ et $\mathfrak q$ par $C$, $B$ et $(0)$ respectivement, et
il suffit donc de voir que pour $s$ voisin de $\eta$, $C_{(\boldsymbol{k}(s))}$
n'a pas d'idéaux premiers associés immergés. Mais cela est évident puisque
$C_{(\boldsymbol{k}(s))}=\boldsymbol{k}(s)[\mathrm T_1,\ldots,\mathrm T_n]$
est un anneau intègre. C.Q.F.D.

\begin{theorem}[9.7.7]
\label{IV.9.7.7-fr}
Soit $f:X\to S$ un morphisme de présentation finie, et soit $E$ l'ensemble des
$s\in S$ pour lesquels $X_s$ a l'une des propriétés suivantes~: être~:
\begin{enumerate}[label=(\roman*)]
  \item géométriquement irréductible~;
  \item géométriquement connexe~;
  \item géométriquement réduit~;
  \item géométriquement intègre.
\end{enumerate}
Alors $E$ est localement constructible dans $S$.
\end{theorem}

Pour voir que les propriétés envisagées sont constructibles, remarquons d'abord
qu'elles vérifient trivialement la condition (9.2.1, (i)), en vertu de
(4.5.6, (i)) et (4.6.5, (i)). Il reste donc à vérifier (9.2.1, (ii)), donc on
peut supposer $S=\operatorname{Spec}(A)$ affine noethérien et intègre de point
générique $\eta$. En vertu de (4.6.8), il existe une extension finie $K'$ de
$K=\boldsymbol{k}(\eta)$ telle que $(X_\eta)_{(K')}$ soit tel que ses
composantes irréductibles (resp. connexes) soient géométriquement irréductibles
(resp. géométriquement connexes) et que $((X_\eta)_{(K')})_{\mathrm{red}}$ soit
géométriquement réduit. Comme il existe une base de $K'$ sur $K$ formée
d'éléments entiers sur $A$, l'anneau intègre $A'$ engendré par ces éléments est
fini sur $A$ et a $K'$ pour corps des fractions. Posons
$S'=\operatorname{Spec}(A')$~; le morphisme $g:S'\to S$ est fini et dominant,
donc surjectif (\textbf{II}, 6.1.10). Posons $X'=X_{(S')}$, de sorte que si
$\eta'$ est le point générique de $S'$, on a
$X'_{\eta'}=(X_\eta)_{(K')}$~; l'ensemble $E'$ des $s'\in S'$ tels que
$X'_{s'}$ ait l'une des propriétés (i), (ii), (iii) ou (iv) est égal à
$g^{-1}(E)$ (9.2.2, (iv)) ($E$ correspondant bien entendu à la même propriété)~;
$g$ étant surjectif, on a $E=g(E')$ et $S-E=g(S'-E')$~; en outre, $g$ est fermé
et $g^{-1}(\eta)=\{\eta'\}$ puisque $A'$ est intègre et fini sur $A$ (Bourbaki,
\emph{Alg. comm.}, chap. V, \S~2, n$^\circ$~1, cor. 1 de la prop. 1), donc
l'image par $g$ de tout voisinage de $\eta'$ est un voisinage de $\eta$. Le
théorème sera donc démontré si l'on prouve que $E'$ ou $S'-E'$ est un voisinage
de $\eta'$. Autrement dit, on peut désormais supposer que les composantes
irréductibles (resp. connexes) de $X_\eta$ sont géométriquement

\oldpage[IV-3]{80}

irréductibles (resp. géométriquement connexes) et que $(X_\eta)_{\mathrm{red}}$
est géométriquement réduit.

Supposons d'abord que $\eta\in S-E$ pour une des propriétés (i) à (iv). Si
$X_\eta$ n'est pas géométriquement irréductible (resp. géométriquement connexe),
il n'est pas irréductible (resp. connexe) en vertu de l'hypothèse précédente,
donc il en est de même de $X_s$ pour $s$ dans un voisinage de $\eta$ (9.7.1), et
\emph{a fortiori} dans ce voisinage $X_s$ n'est pas géométriquement irréductible
(resp. géométriquement connexe). D'autre part, si $X_\eta$ n'est pas
géométriquement réduit, il n'est pas réduit (sans quoi il serait égal à
$(X_\eta)_{\mathrm{red}}$ qui est géométriquement réduit par hypothèse)~; donc,
dans un voisinage de $\eta$, $X_s$ n'est pas réduit (9.7.2) et \emph{a fortiori}
n'est pas géométriquement réduit. Enfin, si $X_\eta$ n'est pas géométriquement
intègre, ou bien il n'est pas réduit, et alors on vient de voir que $X_s$ n'est
pas réduit (ni \emph{a fortiori} intègre) dans un voisinage de $\eta$~; ou bien
$X_\eta$ est réduit (donc géométriquement réduit par hypothèse), et alors il
n'est pas géométriquement irréductible, et on a vu ci-dessus qu'il en est alors
de même de $X_s$ pour $s$ voisin de $\eta$~; \emph{a fortiori} $X_s$ n'est pas
géométriquement intègre pour ces valeurs de $s$.

Nous allons donc désormais supposer que $\eta\in E$ et examiner séparément
chacune des propriétés considérées.

1$^\circ$ Supposons que $X_\eta$ soit géométriquement intègre. Soit $L$ le corps
des fonctions rationnelles sur $X_\eta$~; l'hypothèse sur $X_\eta$ implique que
$L$ est extension séparable de $K$ (4.6.3), donc extension séparable finie d'une
extension pure $K(\mathrm T_1,\ldots,\mathrm T_n)$ ($\mathrm T_i$
indéterminées)~; il y a par suite un élément $x\in L$, entier sur l'anneau
$K[\mathrm T_1,\ldots,\mathrm T_n]$, et tel que
$L=K(\mathrm T_1,\ldots,\mathrm T_n)(x)$~; soit
$G\in K[\mathrm T_1,\ldots,\mathrm T_n,\mathrm T_{n+1}]$ son polynôme minimal.
Il existe un élément $g\neq0$ de $A$ tel que tous les coefficients $\neq0$ de
$G$ (qui appartiennent à $K$) soient dans l'anneau $A_g$~; remplaçant $A$ par
$A_g$ (ce qui revient à remplacer $S$ par un voisinage de $\eta$), on peut donc
supposer que $G$ a ses coefficients dans $A$~; désignant par $F$ le polynôme
$G$ considéré comme élément de
$A[\mathrm T_1,\ldots,\mathrm T_n,\mathrm T_{n+1}]$, on a donc
$G=F_{(\boldsymbol{k}(\eta))}$. Posons
$Y=\operatorname{Spec}(A[\mathrm T_1,\ldots,\mathrm T_{n+1}]/(F))$, de sorte
que
$Y_\eta=\operatorname{Spec}(K[\mathrm T_1,\ldots,\mathrm T_{n+1}]/(G))$~;
$Y_\eta$ est un schéma intègre ayant $L$ pour corps des fonctions rationnelles.
Comme $X_\eta$ et $Y_\eta$ sont noethériens, il existe un ouvert non vide
$V\subset Y_\eta$ et une immersion ouverte $v:V\to X_\eta$ (nécessairement
dominante) (\textbf{I}, 6.5.1, (ii) et 6.5.4, (ii)). Soit $W$ un ouvert de $Y$
tel que $W\cap Y_\eta=V$~; appliquant (8.8.2, (i)) et (8.10.5, (iii)) suivant
la méthode de (8.1.2, \emph{a)}), on voit qu'en remplaçant au besoin $S$ par un
voisinage de $\eta$, on peut supposer que $v=w_\eta$, où $w:W\to X$ est une
immersion ouverte.

Cela étant, on a vu (4.6.3) que le critère pour qu'un préschéma algébrique
intègre soit géométriquement intègre ne dépend que de son corps de fonctions
rationnelles~; comme $X_\eta$ est géométriquement intègre par hypothèse, il en
est de même de $Y_\eta$, et la définition de $G$ entraîne donc que ce polynôme
est géométriquement irréductible (9.7.4). Appliquant (9.7.5), on voit qu'il y a
un voisinage $U$ de $\eta$ dans $S$ tel que $F_{(\boldsymbol{k}(s))}$ soit
géométriquement irréductible pour tout $s\in U$, donc $Y_s$ est géométriquement
intègre pour $s\in U$ (9.7.4)~; en outre on peut supposer que pour $s\in U$,
$W_s$ n'est pas vide (9.5.1), et par suite est géométriquement intègre (4.6.3)~;
enfin, on peut supposer aussi que

\oldpage[IV-3]{81}

pour $s\in U$, $w_s:W_s\to X_s$ est une immersion ouverte (9.6.1, (x)), et que
son image dans $X_s$ est partout dense (9.5.3). Autrement dit, pour $s\in U$,
il y a dans $X_s$ un ouvert partout dense $W'_s$ qui est géométriquement
intègre~; le critère (4.5.9, \emph{c)}) montre donc déjà que $X_s$ est
géométriquement irréductible pour $s\in U$. Enfin, comme $X_\eta$ est réduit et
par suite n'a pas de cycle premier associé immergé (3.2.1), on peut aussi
supposer que pour $s\in U$, $X_s$ n'a pas de cycle premier associé immergé
(9.7.6)~; soit alors (pour un $s\in U$ fixé) $\Omega$ une extension
algébriquement close de $\boldsymbol{k}(s)$, et soit
$p:(X_s)_{(\Omega)}\to X_s$ la projection canonique~; $p^{-1}(W'_s)$ est un
ouvert dense dans $(X_s)_{(\Omega)}$ (2.3.10) et est intègre par hypothèse~; en
outre $(X_s)_{(\Omega)}$ n'a pas de cycle premier associé immergé (4.2.7), donc
on conclut de (3.2.1) que $(X_s)_{(\Omega)}$ est réduit~; cela achève de prouver
que $X_s$ est géométriquement intègre (4.6.1).

2$^\circ$ Supposons que $X_\eta$ soit géométriquement irréductible~; comme
$X_{\mathrm{red}}$ est aussi de type fini sur $S$ (1.5.4, (vi)) on peut, en
tenant compte de (\textbf{I}, 5.1.8), remplacer $X$ par $X_{\mathrm{red}}$~;
alors $X_\eta$ est aussi intègre, et comme par hypothèse $X_\eta$ est
géométriquement réduit, il est géométriquement intègre. On est alors dans les
conditions de 1$^\circ$ et on en conclut (revenant aux hypothèses initiales) que
$X_s$ est géométriquement irréductible pour $s$ voisin de $\eta$.

3$^\circ$ Supposons $X_\eta$ géométriquement connexe, et soient $Z_i$
($1\leqslant i\leqslant n$) les composantes irréductibles de $X_\eta$~; il
existe (en vertu de (5.10.8.1) appliqué à $\mathfrak T=\{\varnothing\}$) une
application \emph{surjective} $j\mapsto\nu(j)$ d'un intervalle $[1,m]$ de
$\mathbf N$ sur $[1,n]$ telle que
$Z_{\nu(j)}\cap Z_{\nu(j+1)}\neq\varnothing$ pour $1\leqslant j\leqslant m$.
Pour tout $i$, soit $X_i$ l'adhérence de $Z_i$ dans $X$, et soit $Y$ la réunion
des $X_i$~; comme $Y_\eta=X_\eta$ par définition, on peut supposer, en vertu de
(9.5.1), que $Y_s=X_s$ pour tout $s\in S$, donc que $X_s$ est réunion des
$(X_i)_s$. Or, en considérant les sous-préschémas fermés réduits de $X$ ayant
les $X_i$ pour espaces sous-jacents, on voit par 2$^\circ$ qu'il existe un
voisinage $U$ de $\eta$ dans $S$ tel que pour $s\in U$ les $(X_i)_s$ soient
\emph{géométriquement irréductibles} (puisque les $(X_i)_\eta$ peuvent être
supposés géométriquement irréductibles, comme on l'a vu au début). En outre, on
peut aussi supposer que pour $s\in U$, on a
$(X_{\nu(j)})_s\cap(X_{\nu(j+1)})_s\neq\varnothing$ (9.5.1) pour
$1\leqslant j\leqslant m$~; on en conclut aussitôt que $X_s$ est connexe, donc
(4.5.13.1) géométriquement connexe pour $s\in U$.

4$^\circ$ Supposons $X_\eta$ géométriquement réduit~; soient $Z_i$ les
composantes irréductibles de $X_\eta$, $W_i$ l'intérieur de $Z_i$ dans
$X_\eta$~; il y a pour chaque $i$ un ouvert $V_i$ de $X$ tel que
$W_i=V_i\cap X_\eta$ pour tout $i$~; comme les $W_i$ sont ouverts et deux à deux
disjoints et que leur réunion est dense dans $X_\eta$, on peut (9.5.1, 9.5.3 et
9.5.4) supposer que pour $s$ voisin de $\eta$, les $(V_i)_s$ sont des ouverts
deux à deux disjoints dans $X_s$ et que leur réunion est dense dans $X_s$. En
outre, comme les $W_i$ sont géométriquement réduits et ont été supposés au
début géométriquement irréductibles, il résulte du 1$^\circ$ que pour $s$ voisin
de $\eta$, les $(V_i)_s$ sont géométriquement intègres, et \emph{a fortiori}
réduits. D'autre part, on tire de (9.7.6) que pour $s$ voisin de $\eta$, $X_s$
n'a pas de cycle premier associé immergé, puisqu'il en est ainsi de $X_\eta$
qui est réduit (3.2.1)~; on conclut donc de (3.2.1) que $X_s$ est réduit, et de
(4.6.1) qu'il est géométriquement réduit.

\oldpage[IV-3]{82}

Les parties de l'énoncé (9.7.7) concernant les propriétés (i) et (ii) se
généralisent comme suit~:

\begin{proposition}[9.7.8]
\label{IV.9.7.8-fr}
Soient $S$ un préschéma irréductible de point générique $\eta$, et $f:X\to S$
un morphisme de présentation finie. Soit $n$ (resp. $n'$) le nombre géométrique
de composantes irréductibles (resp. connexes) de $X_\eta$ (4.5.2). Alors il
existe un voisinage $U$ de $\eta$ dans $S$ tel que pour tout $s\in U$ le nombre
géométrique de composantes irréductibles (resp. connexes) de $X_s$ soit égal à
$n$ (resp. $n'$).
\end{proposition}

Compte tenu de ce que le nombre géométrique de composantes irréductibles (resp.
connexes) d'un préschéma algébrique est invariant par extension du corps de
base (4.5.6), on voit par la méthode de (9.2.3) que l'on peut se ramener au cas
où $S$ est affine, noethérien et intègre. En outre, en raisonnant comme au début
de la démonstration de (9.7.7), on voit qu'on peut supposer que les composantes
irréductibles (resp. connexes) de $X_\eta$ soient \emph{géométriquement}
irréductibles (resp. \emph{géométriquement connexes}). On sait déjà (9.7.1) que
pour tout $s$ voisin de $\eta$, $X_s$ a au moins $n$ composantes irréductibles
et $n'$ composantes connexes. D'autre part, si $Z_i$ (resp. $Z'_j$) sont les
composantes irréductibles (resp. connexes) de $X_\eta$, $X_i$ (resp. $X'_j$) le
sous-préschéma fermé réduit de $X$ ayant pour espace sous-jacent l'adhérence de
$Z_i$ (resp. $Z'_j$) dans $X$, il résulte de (9.5.1) que dans un voisinage de
$\eta$, les $(X_i)_s$ (resp. $(X'_j)_s$) forment un recouvrement de $X_s$ et que
les $(X'_j)_s$ sont deux à deux disjoints~; comme les $(X_i)_s$ (resp.
$(X'_j)_s$) sont géométriquement irréductibles (resp. géométriquement connexes)
en vertu de (9.7.7) pour $s$ voisin de $\eta$, on voit que $X_s$ a au plus $n$
composantes irréductibles et au plus $n'$ composantes connexes, d'où la
proposition.

\begin{corollary}[9.7.9]
\label{IV.9.7.9-fr}
Soit $f:X\to S$ un morphisme de présentation finie. Pour tout $s\in S$, soit
$n(s)$ (resp. $n'(s)$) le nombre géométrique de composantes irréductibles
(resp. connexes) de $f^{-1}(s)$ (4.5.2). Alors la fonction $s\mapsto n(s)$
(resp. $s\mapsto n'(s)$) est localement constructible dans $S$.
\end{corollary}

Il s'agit de montrer que la propriété «~$X$ est un préschéma algébrique sur un
corps $k$ et le nombre géométrique des composantes irréductibles (resp.
connexes) de $X$ est égal à $n$ (resp. $n'$)~» est constructible. Il résulte de
(4.5.6) que cette propriété vérifie la condition (9.2.1, (i)), et on est donc
ramené au cas où $S$ est affine, noethérien et intègre de point générique
$\eta$~; la conclusion résulte alors de (9.7.8).

\begin{corollary}[9.7.10]
\label{IV.9.7.10-fr}
Soit $X$ un préschéma localement noethérien tel que, si $X'$ est le normalisé de
$X_{\mathrm{red}}$, le morphisme canonique $f:X'\to X$ soit fini. Alors
l'ensemble des points $x\in X$ tels que $X$ soit géométriquement unibranche au
point $x$ est localement constructible dans $X$.
\end{corollary}

En effet, cet ensemble est par définition l'ensemble des points $x\in X$ tels
que le nombre de points géométriques de $f^{-1}(x)$ soit égal à 1. Mais comme
$f$ est fini, ce nombre est aussi le nombre géométrique des composantes
irréductibles de l'espace discret $f^{-1}(x)$ (compte tenu de la définition du
normalisé (\textbf{II}, 6.3.8) et de (4.5.11))~; la conclusion résulte donc de
(9.7.9).

\begin{remark}[9.7.11]
\label{IV.9.7.11-fr}
Soit $Z$ une partie localement constructible de $X$ telle que, pour tout
$s\in S$, $Z\cap X_s$ soit ouvert dans $X_s$, et désignons par $Z_s$ le
sous-préschéma de $X_s$ induit

\oldpage[IV-3]{83}

sur l'ouvert $Z\cap X_s$. Alors, dans le théorème (9.7.7), on peut remplacer
$X_s$ par $Z_s$ sans changer la conclusion~: on le voit en répétant le
raisonnement fait dans (9.6.7, (ii)).
\end{remark}

\begin{proposition}[9.7.12]
\label{IV.9.7.12-fr}
Soient $f:X\to S$ un morphisme de présentation finie, et $g:S\to X$ une
$S$-section de $X$ (\textbf{I}, 2.5.5). Pour tout $s\in S$, soit $X_s^0$ la
composante connexe de $f^{-1}(s)$ qui contient $g(s)$~; alors, la réunion $X^0$
des $X_s^0$ pour $s\in S$ est une partie localement constructible de $X$.
\end{proposition}

Montrons d'abord qu'on peut se ramener au cas où $S$ est affine et
\emph{noethérien}. On peut toujours supposer $S=\operatorname{Spec}(A)$ affine~;
avec les notations de (9.2.3), on a $f=(f_0)_{(S)}$, où $f_0:X_0\to S_0$ est un
morphisme de type fini, et on peut en outre supposer qu'il existe une
$S_0$-section $g_0:S_0\to X_0$ telle que $g=(g_0)_{(S)}$ (8.9.1). Remarquons
maintenant que si $p$ est le morphisme $S\to S_0$, alors, pour tout $s_0\in S_0$,
la composante connexe $(X_0)^0_{s_0}$ de $f_0^{-1}(s_0)$ qui contient $g_0(s_0)$
est \emph{géométriquement connexe} (4.5.13), et par suite, si $s_0=p(s)$, on a
$X_s^0=q^{-1}((X_0)^0_{s_0})$ où $q:X_s\to(X_0)_{s_0}$ est la projection
canonique ((4.5.8) et (4.4.1))~; notre assertion résulte donc de (1.8.2).

Utilisons alors le critère de constructibilité $(\mathbf 0_{\mathrm{III}},9.2.3)$~:
soient $x$ un point de $X$, $Z$ le sous-préschéma réduit de $X$ ayant
$\overline{\{x\}}$ comme espace sous-jacent, $Y$ le sous-préschéma réduit de $S$
ayant $\overline{f(Z)}$ comme espace sous-jacent~; comme la restriction de $f$
à $Z$ se factorise en $Z\to Y\to S$ (\textbf{I}, 5.2.2), on peut remplacer $S$
par $Y$, autrement dit supposer que $f(x)=\eta$, point générique du préschéma
intègre $S$. Par hypothèse, $X_\eta$ est somme de deux sous-préschémas
$X_\eta^0$, $X_\eta^1$, induits sur des ouverts complémentaires de $X_\eta$. En
vertu de (9.5.4) et (9.5.1), on peut donc, en remplaçant au besoin $S$ par un
voisinage de $\eta$, supposer que $X$ est réunion de deux ouverts disjoints
$X^0$, $X^1$ tels que $(X^0)_\eta=X_\eta^0$ et $(X^1)_\eta=X_\eta^1$. Comme
$g:S\to X$ est continu et injectif, $S$ est réunion des deux ouverts disjoints
$g^{-1}(X^0)$ et $g^{-1}(X^1)$~; mais comme $S$ est irréductible, et
\emph{a fortiori} connexe, un de ces deux ouverts est vide, et comme
$g(\eta)\in X_\eta^0$ par définition, on a $g^{-1}(X^1)=\varnothing$. En
d'autres termes, $g$ est une $S$-section de $X^0$~; d'autre part, comme
$(X^0)_\eta$ est géométriquement connexe (4.5.13), il en est de même de
$(X^0)_s$ pour tout $s$ voisin de $\eta$ (9.7.7)~; comme
$g(s)\in(X^0)_s$, on a bien $(X^0)_s=X_s^0$. C.Q.F.D.

\subsection{Décomposition primaire au voisinage d'une fibre générique}
\label{subsection:IV.9.8-fr}

\begin{proposition}[9.8.1]
\label{IV.9.8.1-fr}
Soient $f:X\to S$ un morphisme de présentation finie, $\mathcal F$ un
$\mathcal O_X$-Module quasi-cohérent de présentation finie. Alors l'ensemble
$E$ des $s\in S$ tels que $\mathcal F_s$ soit un $\mathcal O_{X_s}$-Module sans
cycle premier associé immergé est localement constructible.
\end{proposition}

On sait que pour les Modules quasi-cohérents sur les préschémas algébriques, la
propriété d'être sans cycle premier associé immergé est invariante par
changement du corps de base (4.2.7)~; on peut donc se borner au cas où $S$ est
affine, noethérien et intègre de point générique $\eta$, et à prouver que dans
ce cas $E$ ou $S-E$ est un voisinage de $\eta$. On a vu dans (9.7.6) que si
$\eta\in E$, $E$ est un voisinage de $\eta$~; reste donc à considérer le cas où
$\mathcal F_\eta$ a des cycles premiers associés immergés. On peut se borner au
cas où $X$ est affine, et où il y a un sous-$\mathcal O_{X_\eta}$-Module
cohérent $\mathcal H$ de $\mathcal F_\eta$ dont le support $Z'$ est non vide et
\emph{rare} par rapport au support $Z$ de $\mathcal F$ (3.1.3)~; alors, en vertu
de (8.5.2, (i) et (ii))

\oldpage[IV-3]{84}

et (8.5.8), appliqués suivant la méthode de (8.1.2, \emph{a)}), il existe un
sous-$\mathcal O_X$-Module cohérent $\mathcal G$ de $\mathcal F$ tel que
$\mathcal G_\eta=\mathcal H$~; si $Y=\operatorname{Supp}(\mathcal F)$,
$Y'=\operatorname{Supp}(\mathcal G)$, on a par suite
$Y_s=\operatorname{Supp}(\mathcal F_s)$, $Y'_s=\operatorname{Supp}(\mathcal G_s)$
pour tout $s\in S$ (\textbf{I}, 9.1.13) et en particulier $Z=Y_\eta$,
$Z'=Y'_\eta$~; comme $Z-Z'$ est dense dans $Z$ et $Z'\neq\varnothing$, il y a
un voisinage $U$ de $\eta$ tel que pour $s\in U$, $Y_s-Y'_s$ soit dense dans
$Y_s$, et $Y'_s\neq\varnothing$ (9.5.1 et 9.5.3)~; en considérant un point
générique d'une composante irréductible de $Y'_s$ et un voisinage assez petit de
ce point dans $X_s$, on déduit aussitôt de (3.1.3) que $\mathcal F_s$ admet des
cycles premiers associés immergés.

\phantomsection\label{IV.9.8.2-fr}%
\textbf{(9.8.2)} Soient $S$ un préschéma noethérien \emph{intègre}, de point
générique $\eta$, $f:X\to S$ un morphisme de type fini, $\mathcal F$ un
$\mathcal O_X$-Module cohérent. Considérons une décomposition irredondante
réduite $(\mathcal G_i)_{i\in I}$ de $\mathcal F_\eta$ (3.2.6)~; les
$\mathcal G_i$ sont donc des quotients de $\mathcal F_\eta$ et il y a un
homomorphisme injectif
$h:\mathcal F_\eta\to\bigoplus_i\mathcal G_i$~; en outre
$\operatorname{Ass}(\mathcal G_i)$ est réduit à un point $x_i$. Soit $Z_i$
l'adhérence de $\{x_i\}$ dans $X$, de sorte que
$(Z_i)_\eta=\operatorname{Supp}(\mathcal G_i)$. Désignons par $\mathcal J_i$
l'Idéal cohérent de $\mathcal O_X$ définissant le sous-préschéma fermé réduit de
$X$ d'espace sous-jacent $Z_i$ (sous-préschéma encore noté $Z_i$). En vertu de
(8.5.2) et (8.5.8), appliqués suivant la méthode de (8.1.2, \emph{a)}), on peut
(en restreignant au besoin $S$ à un voisinage de $\eta$) supposer qu'il existe
des quotients $\mathcal F_i$ de $\mathcal F$ ($i\in I$) tels que
$\mathcal G_i=(\mathcal F_i)_\eta$ pour tout $i$, et un homomorphisme
$g:\mathcal F\to\bigoplus_i\mathcal F_i$ tel que $h=g_\eta$. En outre
(\textbf{I}, 9.3.5), il existe un entier $m$ tel que
$(\mathcal J_i^m\mathcal F_i)_\eta=(\mathcal J_i)_\eta^m\mathcal G_i=0$ et en
restreignant de nouveau $S$, on peut donc supposer que l'on a aussi
$\mathcal J_i^m\mathcal F_i=0$ (8.5.2.5), donc que le support de $\mathcal F_i$
est contenu dans $Z_i$~; mais comme il est fermé et contient $x_i$, il est
nécessairement égal à $Z_i$.

\begin{proposition}[9.8.3]
\label{IV.9.8.3-fr}
Sous les conditions de (9.8.2), pour tout $s\in S$, et tout $i\in I$, désignons
par $x_{is\alpha}$ ($\alpha\in J_{s,i}$) les points maximaux de $(Z_i)_s$. Il
existe un voisinage $U$ de $\eta$ dans $S$ tel que, pour tout $s\in U$, les
$x_{is\alpha}$ (pour $i\in I$ et, pour chaque $i$, $\alpha\in J_{s,i}$) soient
deux à deux distincts et que $\operatorname{Ass}(\mathcal F_s)$ soit l'ensemble
des $x_{is\alpha}$ (autrement dit, les cycles premiers associés à $\mathcal F_s$
sont les composantes irréductibles des $(Z_i)_s$). En outre, on peut prendre
$U$ tel que, pour que l'adhérence de $\{x_{is\alpha}\}$ dans $X_s$ soit un cycle
premier associé maximal de $\mathcal F_s$, il faut et il suffit que $(Z_i)_\eta$
(adhérence de $\{x_i\}$ dans $X_\eta$) soit un cycle premier associé maximal de
$\mathcal F_\eta$.
\end{proposition}

Il résulte de (3.1.3, \emph{c')}) que pour chaque $i$, il existe un ouvert $W_i$
dans $X_\eta$ tel que $W_i\cap(Z_i)_\eta$ soit non vide, et un
$\mathcal O_{X_\eta}$-Module cohérent $\mathcal H_i$, de support
$W_i\cap(Z_i)_\eta$, tels qu'il y ait un homomorphisme injectif
$v_i:\mathcal H_i\to\mathcal F_\eta|W_i$. Soit $V_i$ un ouvert de $X$ tel que
$V_i\cap X_\eta=W_i$~; appliquant, comme dans (9.8.2), les résultats de (8.5.2)
et (8.5.8), on peut (en restreignant $S$) supposer qu'il existe un Module
cohérent $\mathcal F'_i$, de support $V_i\cap Z_i$ et un homomorphisme
$u_i:\mathcal F'_i\to\mathcal F|V_i$ tels que
$\mathcal H_i=(\mathcal F'_i)_\eta$ et $v_i=(u_i)_\eta$. Nous allons démontrer
qu'il y a un voisinage $U$ de $\eta$ tel que pour $s\in U$, les propriétés
suivantes soient vérifiées~:
\begin{enumerate}[label=\emph{(\roman*)}]
  \item Les $(\mathcal F_i)_s$ sont sans cycle premier associé immergé.
  \item L'homomorphisme
  $g_s:\mathcal F_s\to\bigoplus_i(\mathcal F_i)_s$ est injectif.
  \item $(Z_i)_s\cap(V_i)_s$ est dense dans $(Z_i)_s$,
  $(\mathcal F'_i)_s$ a pour support $(Z_i)_s\cap(V_i)_s$ et
  $(u_i)_s:(\mathcal F'_i)_s\to\mathcal F_s|(V_i)_s$ est injectif.
\end{enumerate}

\oldpage[IV-3]{85}

\begin{enumerate}[label=\emph{(\roman*)}]
  \item[(iv)] Pour $i\neq j$, toute composante irréductible de $(Z_i)_s$ est
  distincte de toute composante irréductible de $(Z_j)_s$.
\end{enumerate}
Or, (i) a déjà été vu (9.7.6)~; (ii) est un cas particulier de (9.4.5)~; (iii)
résulte de même de (9.5.3) et (9.4.5). Enfin, si $i\neq j$,
$(Z_i)_\eta\cap(Z_j)_\eta=(Z_i\cap Z_j)_\eta$ est rare dans $(Z_i)_\eta$ ou
dans $(Z_j)_\eta$~; supposons par exemple que $(Z_i\cap Z_j)_\eta$ soit rare
dans $(Z_j)_\eta$~; alors il résulte de (9.5.3) et (9.5.4) que pour $s$ voisin
de $\eta$, $(Z_i\cap Z_j)_s=(Z_i)_s\cap(Z_j)_s$ est rare dans $(Z_j)_s$, ce qui
montre qu'aucune composante irréductible de $(Z_j)_s$ ne peut être contenue
dans une composante irréductible de $(Z_i)_s$, ni \emph{a fortiori} lui être
égale.

Cela étant, il résulte de (ii) et de (3.1.7) que pour $s\in U$, on a
\[
\operatorname{Ass}(\mathcal F_s)\subset
\bigcup_i\operatorname{Ass}((\mathcal F_i)_s),
\]
de (i) que $\operatorname{Ass}((\mathcal F_i)_s)$ est l'ensemble des
$x_{is\alpha}$ ($\alpha\in J_{s,i}$). D'autre part, en vertu de (iii) et du
critère (3.1.3, \emph{c')}), chacun des $x_{is\alpha}$ ($\alpha\in J_{s,i}$,
$i\in I$) appartient à $\operatorname{Ass}(\mathcal F_s)$. Enfin (iv) signifie
que les $x_{is\alpha}$ sont deux à deux distincts.

Reste à prouver la dernière assertion de l'énoncé. Or, il résulte de (i) que
pour $s$ et pour un $i$ donnés, aucun des ensembles $\{x_{is\alpha}\}$ ne peut
être contenu dans un autre pour $\alpha\in J_{s,i}$. D'autre part, si
$(Z_j)_\eta\subset(Z_i)_\eta$, on peut supposer $U$ pris tel que
$(Z_j)_s\subset(Z_i)_s$ pour $s\in U$ (9.5.1), donc chaque $x_{js\beta}$
appartient à l'adhérence d'un $x_{is\alpha}$~; au contraire, si
$(Z_j)_\eta\cap(Z_i)_\eta$ est rare dans $(Z_j)_\eta$, on a vu en prouvant (iv)
que $(Z_j)_s\cap(Z_i)_s$ est rare dans $(Z_j)_s$, donc aucun des $x_{js\beta}$
n'est adhérent à un $x_{is\alpha}$. En particulier, si $(Z_j)_\eta$ est
maximal, ce qui équivaut à dire que $(Z_j)_\eta\cap(Z_i)_\eta$ est rare dans
$(Z_j)_\eta$ pour \emph{tout} $i\neq j$, on en conclut que
$(Z_j)_s\cap(Z_i)_s$ est rare dans $(Z_j)_s$ pour tout $i\neq j$, donc que tout
$x_{js\alpha}$ ($\alpha\in J_{s,j}$) est \emph{maximal}. C.Q.F.D.

\begin{corollary}[9.8.4]
\label{IV.9.8.4-fr}
Les notations et hypothèses étant celles de (9.8.2), il existe un voisinage $U$
de $\eta$ dans $S$ tel que, pour tout $s\in U$, chacun des $(\mathcal F_i)_s$
soit sans cycle premier associé immergé~; en outre, si
$(\mathcal F_{is\alpha})_{\alpha\in J_{s,i}}$ est l'unique décomposition
irredondante réduite de $(\mathcal F_i)_s$, alors, pour tout $s\in U$, la
famille $(\mathcal F_{is\alpha})_{i\in I,\,\alpha\in J_{s,i}}$ est une
décomposition irredondante réduite de $\mathcal F_s$.
\end{corollary}

La première assertion résulte de (9.8.1) et de la définition des
$\mathcal F_i$. D'autre part, on a vu dans (9.8.3) que l'homomorphisme
$\mathcal F_s\to\bigoplus_i(\mathcal F_i)_s$ est injectif et par définition il
en est de même de chacun des homomorphismes
$(\mathcal F_i)_s\to\bigoplus_\alpha\mathcal F_{is\alpha}$, donc
l'homomorphisme $\mathcal F_s\to\bigoplus_{i,\alpha}\mathcal F_{is\alpha}$ est
injectif. Comme on peut supposer (9.4.5) que pour $s\in U$ chacun des
$(\mathcal F_i)_s$ est un quotient de $\mathcal F_s$, les
$\mathcal F_{is\alpha}$ sont des quotients de $\mathcal F_s$, et il reste à
vérifier (3.2.5) que les $x_{is\alpha}$ sont deux à deux distincts et
appartiennent à $\operatorname{Ass}(\mathcal F_s)$, ce qui a été prouvé dans
(9.8.3).

\begin{proposition}[9.8.5]
\label{IV.9.8.5-fr}
Soient $f:X\to S$ un morphisme de présentation finie, $\mathcal F$ un
$\mathcal O_X$-Module quasi-cohérent de présentation finie. Pour tout $s\in S$,
soit $E(s)$ (resp. $E'(s)$) l'ensemble fini (partie de
$\mathbf Z\cup\{-\infty\}$) des dimensions des cycles premiers associés à
$\mathcal F_s$ (resp. des cycles premiers maximaux associés à $\mathcal F_s$,
c'est-à-dire les composantes irréductibles de $\operatorname{Supp}(\mathcal F_s)$).
Alors les fonctions $s\mapsto E(s)$ et $s\mapsto E'(s)$ sont localement
constructibles dans $S$.
\end{proposition}

\oldpage[IV-3]{86}

Il résulte de (4.2.7) et (4.2.8) que la condition (9.2.1, (i)) est satisfaite
pour les propriétés dont nous voulons démontrer qu'elles sont constructibles.
On est donc ramené au cas où $S$ est affine, noethérien et intègre, et à prouver
que $E$ et $E'$ sont constantes au voisinage du point générique $\eta$ de $S$~;
mais cela résulte de (9.8.3) et (9.5.5).

\begin{proposition}[9.8.6]
\label{IV.9.8.6-fr}
Avec les hypothèses et notations de (9.8.2), soit $i\in I$ tel que $(Z_i)_\eta$
soit maximal (autrement dit, une composante irréductible de
$\operatorname{Supp}(\mathcal F_\eta)$). Alors, il existe un voisinage $U$ de
$\eta$ dans $S$ tel que pour tout $s\in U$ et tout $\alpha\in J_{s,i}$, la
longueur géométrique de $\mathcal F_s$ en $x_{is\alpha}$ (relative à
$\boldsymbol{k}(s)$) (4.7.5) est égale à la longueur géométrique de
$\mathcal F_\eta$ en $x_i$ (relative à $\boldsymbol{k}(\eta)$).
\end{proposition}

On peut évidemment se borner au cas où $S=\operatorname{Spec}(A)$ est affine~;
montrons d'abord qu'on peut se ramener au cas où le sous-préschéma $(Z_i)_\eta$
de $X_\eta$, qui est réduit, est \emph{géométriquement intègre}. Il y a en effet
une extension finie $K'$ de $K=\boldsymbol{k}(\eta)$ telle que
$(((Z_i)_\eta)_{(K')})_{\mathrm{red}}$ soit géométriquement réduit et que les
composantes irréductibles de $((Z_i)_\eta)_{(K')}$ soient géométriquement
irréductibles (4.6.8). Procédons comme dans la démonstration de (9.7.7) en
considérant une sous-$A$-algèbre $A'$ de $K'$ ayant $K'$ pour corps des
fractions et finie sur $A$. On pose $S'=\operatorname{Spec}(A')$ et on considère
le morphisme fini surjectif $g:S'\to S$~; soient ensuite $X'=X_{(S')}$ et
$\mathcal F'=\mathcal F\otimes_{\mathcal O_S}\mathcal O_{S'}$, et soit $\eta'$
le point générique de $S'$. Pour tout $s'\in S'$, soit $s=g(s')$~; si $T$ est
une composante irréductible de $\operatorname{Supp}(\mathcal F_s)$, les
composantes irréductibles $T'_j$ de $T_{(\boldsymbol{k}(s'))}$ sont des
composantes irréductibles de $\operatorname{Supp}(\mathcal F'_{s'})$ et dominent
$T$ (4.2.7) et les multiplicités radicielles de $T$ par rapport à $\mathcal F_s$
et de chaque $T'_j$ par rapport à $\mathcal F'_{s'}$ sont les mêmes (4.7.9). Le
raisonnement de la première partie de (9.7.7) montre donc qu'on peut se borner à
démontrer la proposition pour $X'$ et $\mathcal F'$~; et en vertu du choix de
$K'$, les sous-préschémas réduits ayant pour espaces sous-jacents les composantes
irréductibles de $((Z_i)_\eta)_{(K')}$ sont géométriquement intègres (4.6.1).

Supposons donc désormais $(Z_i)_\eta$ géométriquement intègre~; alors (9.7.7),
il en est de même de $(Z_i)_s$ pour $s$ voisin de $\eta$~; la définition (4.7.5)
montre qu'il suffira donc de prouver que la \emph{longueur} du
$\mathcal O_{X_\eta,x_i}$-module $(\mathcal F_\eta)_{x_i}$ est égale à celle du
$\mathcal O_{X_s,x_{is}}$-module $(\mathcal F_s)_{x_{is}}$ (on a ici supprimé
l'indice $\alpha$, inutile par hypothèse). La question étant évidemment locale
sur $X$, on peut supposer (en se restreignant à un voisinage de $x_i$) que
$\mathcal F=\mathcal F_i$, donc que $\mathcal F_\eta$ est \emph{irredondant} sur
$X=\operatorname{Spec}(B)$ affine, et on écrira $Z$ au lieu de $Z_i$, et $x$
pour le point générique de $Z$ (et de $Z_\eta$). L'anneau noethérien $B$
contient donc $A$ comme sous-anneau, et $\mathcal F=\widetilde M$, où $M$ est un
$B$-module de type fini~; en outre, si $K$ est le corps des fractions de $A$,
le $B_{(K)}$-module $M_{(K)}$ est $\neq0$ et irredondant. Soit $q$ l'unique
élément de $\operatorname{Ass}(M_{(K)})$ et soit $\mathfrak p$ l'idéal premier
de $B$, image réciproque de $q$. En utilisant (5.11.1.1) comme dans la
démonstration de (9.7.6), on se ramène au cas où $B$ est \emph{intègre} et $M$
un sous-module $\neq0$ de $B$~; alors $\mathcal F_\eta$ est un
sous-$\mathcal O_{X_\eta}$-Module non nul de $\mathcal O_{Z_\eta}$, et en vertu
de (9.4.5), pour $s$ voisin de $\eta$, $\mathcal F_s$ est isomorphe à un
sous-$\mathcal O_{X_s}$-Module non nul de $\mathcal O_{Z_s}$~; comme $Z_s$ est
géométriquement intègre, les longueurs de $(\mathcal F_\eta)_x$ et de
$(\mathcal F_s)_{x_s}$ sont toutes deux égales à $1$, ce qui achève la
démonstration.

\oldpage[IV-3]{87}

\begin{corollary}[9.8.7]
\label{IV.9.8.7-fr}
Soient $f:X\to S$ un morphisme de présentation finie, $\mathcal F$ un
$\mathcal O_X$-Module quasi-cohérent de présentation finie. Pour tout $s\in S$,
soit $M(s)$ l'ensemble des couples $(d,m)$ tels qu'il existe une composante
irréductible de $\operatorname{Supp}(\mathcal F_s)$ de dimension $d$ et de
multiplicité radicielle $m$ pour $\mathcal F_s$ (4.7.8). Alors la fonction
$s\mapsto M(s)$ est localement constructible dans $S$.
\end{corollary}

Il résulte de (4.2.7) et (4.7.9) que la condition (9.2.1, (i)) est satisfaite
pour la propriété dont nous voulons démontrer qu'elle est constructible. On est
donc ramené au cas où $S$ est affine, noethérien et intègre et à prouver que $M$
est constante dans un voisinage du point générique de $S$~; mais alors la
proposition résulte de (9.8.3) et (9.8.6).

\begin{proposition}[9.8.8]
\label{IV.9.8.8-fr}
Soient $f:X\to S$ un morphisme de présentation finie, $\mathcal F$ un
$\mathcal O_X$-Module quasi-cohérent de présentation finie. Pour tout $s\in S$,
soit $l(s)$ la somme des multiplicités totales pour $\mathcal F_s$ (relatives à
$\boldsymbol{k}(s)$) des points génériques des composantes irréductibles de
$\operatorname{Supp}(\mathcal F_s)$ (4.7.12). Alors la fonction $s\mapsto l(s)$
est localement constructible dans $S$.
\end{proposition}

Compte tenu de (4.7.12), la condition (9.2.1, (i)) est satisfaite pour la
propriété dont nous voulons démontrer qu'elle est constructible, et on est donc
ramené au cas où $S$ est affine, noethérien et intègre de point générique
$\eta$, et à montrer que $l$ est constante dans un voisinage de $\eta$.
Utilisant les notations de (9.8.2), cela résulte de la définition (4.7.12), du
fait que le nombre géométrique des composantes irréductibles de chaque
$(Z_i)_s$ est constant dans un voisinage de $\eta$ (9.7.8), que la longueur
géométrique de $\mathcal F_s$ en $x_{is\alpha}$ est également à celle de
$\mathcal F_\eta$ en $x_i$ pour chaque $i$ tel que $(Z_i)_\eta$ soit maximal
(9.8.6) et enfin de ce que l'adhérence de $\{x_{is\alpha}\}$ est un cycle
premier associé maximal de $\mathcal F_s$ si et seulement si $(Z_i)_\eta$ est
un cycle premier associé maximal de $\mathcal F_\eta$ (9.8.3).

\begin{remark}[9.8.9]
\label{IV.9.8.9-fr}
On peut préciser de diverses façons les propositions précédentes~;
bornons-nous à un énoncé à titre d'exemple. Disons qu'une partie finie $P$ d'un
$k$-préschéma algébrique $X$ est \emph{saturée} si, pour tout couple de points
$x$, $y$ de $P$, les points génériques des composantes irréductibles de
$\overline{\{x\}}\cap\overline{\{y\}}$ appartiennent aussi à $P$~; pour toute
partie finie $Q$ de $X$, il existe une plus petite partie finie $P$ de $X$
contenant $Q$ et saturée~; nous dirons que $P$ est le \emph{saturé} de $Q$. Pour
tout $\mathcal O_X$-Module cohérent $\mathcal F$, nous appellerons
\emph{squelette primaire} de $\mathcal F$ le système $(P,Q,\omega,d,m)$ où
$Q=\operatorname{Ass}(\mathcal F)$, $P$ est le saturé de $Q$, $\omega$ la
relation d'ordre $\overline{\{x\}}\subset\overline{\{y\}}$ sur $P$, $d$ la
fonction $x\mapsto\dim\overline{\{x\}}$ sur $P$, $m$ la fonction
$x\mapsto\operatorname{long}_{\mathcal O_x}\mathcal F_x$ définie sur l'ensemble
des éléments de $Q$ maximaux pour la relation $\omega$. Appelons d'autre part
\emph{squelette virtuel} tout système $(P,Q,\omega,d,m)$, où $P$ est un
ensemble, $Q$ une partie de $P$, $\omega$ une relation d'ordre sur $P$, $d$ une
application de $P$ dans $\mathbf N$, $m$ une application dans $\mathbf N$ de
l'ensemble des éléments maximaux de $Q$~; on définit de façon évidente la notion
d'\emph{isomorphisme} de deux squelettes virtuels. Enfin, avec les notations
précédentes, appelons \emph{type primaire} de $\mathcal F$ la classe (pour la
relation d'isomorphie des squelettes virtuels) du squelette primaire de
$\mathcal F$. Il résulte de (4.2.6), (4.2.7), (4.2.8), (4.5.1) et (4.7.9) que
si, pour une extension algébriquement close $k'$ de $k$, on pose
$X'=X\otimes_k k'$ et $\mathcal F'=\mathcal F\otimes_k k'$, le
\emph{type primaire} de $\mathcal F'$ est indépendant de

\oldpage[IV-3]{88}

l'extension algébriquement close $k'$ de $k$ considérée~; nous dirons que c'est
le \emph{type primaire géométrique} de $\mathcal F$. Ces définitions posées,
l'énoncé que nous avons en vue est le suivant~:

\phantomsection\label{IV.9.8.9.1-fr}%
\textbf{(9.8.9.1)} Soient $f:X\to S$ un morphisme de présentation finie,
$\mathcal F$ un $\mathcal O_X$-Module quasi-cohérent de présentation finie.
Pour tout $s\in S$, soit $T(s)$ le type primaire géométrique de $\mathcal F_s$.
Alors la fonction $s\mapsto T(s)$ est localement constructible dans $S$.

Compte tenu des remarques qui précèdent, on est ramené comme d'ordinaire à
prouver que si $S$ est affine, noethérien et intègre de point générique $\eta$,
$T(s)$ est constante au voisinage de $\eta$. Raisonnant comme au début de
(9.7.7), on peut supposer que toutes les parties irréductibles de $X_\eta$ qui
interviennent sont géométriquement irréductibles, et alors la proposition
résulte de (9.5.1), (9.5.5), (9.8.3) et (9.8.6)~; nous laissons les détails au
lecteur. On pourrait généraliser en considérant plusieurs Modules cohérents et
en définissant leur «~squelette primaire simultané~», etc. La conclusion
générale de ce qui a été vu depuis le début de ce paragraphe est que pour toutes
les propriétés du type envisagé (et pour un $S$ irréductible) \emph{les
propriétés valables sur la «~fibre générique~» le sont encore sur toutes les
fibres voisines}.
\end{remark}

\subsection{Constructibilité des propriétés locales des fibres}
\label{subsection:IV.9.9-fr}

\begin{proposition}[9.9.1]
\label{IV.9.9.1-fr}
Soient $f:X\to S$ un morphisme localement de présentation finie, $Z$ une
partie localement constructible de $X$ telle que pour tout $s\in S$, $Z_s$ soit
fermé dans $X_s$, $\Phi$ une partie finie de $\mathbf Z\cup\{\pm\infty\}$.
Alors les parties suivantes de $X$ sont localement constructibles~:
\begin{enumerate}
\item[(i)] L'ensemble des $x\in X$ tels que
$\dim_x(Z_{f(x)})$ appartienne à $\Phi$.
\item[(ii)] L'ensemble des $x\in X$ tels que
$\operatorname{codim}_x(Z_{f(x)},X_{f(x)})$ appartienne à $\Phi$.
\item[(iii)] L'ensemble des $x\in X$ tels que l'anneau local
$\mathcal O_{X_{f(x)},x}$ soit équidimensionnel.
\end{enumerate}
\end{proposition}

On notera que les propriétés (i) et (ii) s'expriment encore en disant que les
fonctions $x\mapsto\dim_x(Z_{f(x)})$ et
$x\mapsto\operatorname{codim}_x(Z_{f(x)},X_{f(x)})$ sont
\emph{localement constructibles} dans $X$ ($\mathbf 0_{\mathrm{III}},9.3.1$).

Les questions étant locales sur $X$, on peut se borner au cas où
$S=\operatorname{Spec}(A)$ et $X=\operatorname{Spec}(B)$ sont affines et où
$f$ est un morphisme de présentation finie~; il existe alors un sous-anneau
$A_0$ de $A$ qui est une $\mathbf Z$-algèbre de type fini, un
$A_0$-préschéma de type fini $X_0$ et une partie constructible $Z_0$ de $X_0$
tels que $X=X_0\otimes_{A_0}A$ et $Z=h^{-1}(Z_0)$, où $h:X\to X_0$ est la
projection canonique ((8.9.1) et (8.3.11)). En outre, pour tout $s\in S$, si
$s_0$ est la projection de $s$ dans $S_0=\operatorname{Spec}(A_0)$, on a
\[
X_s=(X_0)_{s_0}\otimes_{k(s_0)}k(s),
\]
et $h_s$ est la projection $X_s\to(X_0)_{s_0}$, on a
$Z_s=h_s^{-1}((Z_0)_{s_0})$. Comme le morphisme $h_s$ est fidèlement plat et
quasi-compact, l'hypothèse que $Z_s$ est fermé dans $X_s$ entraîne que
$(Z_0)_{s_0}$ est fermé dans $(X_0)_{s_0}$ (2.3.12).

Cela étant, la transitivité des fibres (\textbf{I}, 3.6.4) et la prop. (4.2.7)
entraînent que l'ensemble des dimensions des composantes irréductibles de
$Z_s$ contenant $x$ est le même que l'ensemble des dimensions des composantes
irréductibles de $(Z_0)_{s_0}$ contenant $x_0=h_s(x)$. En particulier, on a
$\dim_x(Z_s)=\dim_{x_0}((Z_0)_{s_0})$. D'autre part, si $Z_s^{(\beta)}$ sont les
composantes irréductibles de $Z_s$ contenant $x$ et $X_s^{(\alpha)}$ les
composantes irréductibles de $X_s$

\oldpage[IV-3]{89}

contenant $x$, on a
\[
\operatorname{codim}_x(Z_s,X_s)=
\inf_\beta\left(\sup_\alpha
\bigl(\operatorname{codim}(Z_s^{(\beta)},X_s^{(\alpha)})\bigr)\right),
\]
$(\alpha,\beta)$ variant dans l'ensemble des couples tels que
$x\in Z_s^{(\beta)}\subset X_s^{(\alpha)}$ ($\mathbf 0$, 14.2.6). Comme les
préschémas algébriques irréductibles sont biéquidimensionnels (5.2.1), on peut
écrire, en vertu de ($\mathbf 0$, 14.3.3.1)
\begin{equation}
\tag{9.9.1.1}
\label{IV.9.9.1.1-fr}
\operatorname{codim}_x(Z_s,X_s)=
\inf_\beta\left(\sup_\alpha
\bigl(\dim(X_s^{(\alpha)})-\dim(Z_s^{(\beta)})\bigr)\right)
\end{equation}
avec le même choix des couples $(\alpha,\beta)$. Comme $h_s$ est fidèlement
plat et quasi-compact, pour tout couple formé d'une composante irréductible
$(X_0)_{s_0}^{(\lambda)}$ de $(X_0)_{s_0}$ contenant $x_0$ et d'une composante
irréductible $(Z_0)_{s_0}^{(\mu)}$ de $(Z_0)_{s_0}$ contenant $x_0$ et contenue
dans $(X_0)_{s_0}^{(\lambda)}$, il existe un couple
$(X_s^{(\alpha)},Z_s^{(\beta)})$ du type décrit plus haut et tel que
$Z_s^{(\beta)}$ domine $(Z_0)_{s_0}^{(\mu)}$ et que $X_s^{(\alpha)}$ domine
$(X_0)_{s_0}^{(\lambda)}$ (2.3.5). La formule (9.9.1.1) (et la formule analogue
appliquée à $(X_0)_{s_0}$) montrent alors, en vertu de (4.2.7), que l'on a
\[
\operatorname{codim}_x(Z_s,X_s)=
\operatorname{codim}_{x_0}((Z_0)_{s_0},(X_0)_{s_0}).
\]

On voit ainsi que si $E$ (resp. $E_0$) est l'ensemble des $x\in X$ (resp. des
$x_0\in X_0$) vérifiant l'une des conditions (i), (ii), (iii) de l'énoncé
(resp. la même condition), on a $E=h^{-1}(E_0)$, et en vertu de (1.8.2), on
voit qu'on peut se borner au cas où $A$ est \emph{noethérien}, donc aussi $B$.
Compte tenu de ($\mathbf 0_{\mathrm{III}},9.2.3$), ainsi que de (9.9.1.1), on
est ramené à voir que pour tout $x\in X$, il y a un voisinage $V$ de $x$ dans
$\overline{\{x\}}$ tel que, pour tout $x'\in V$, l'ensemble des dimensions des
composantes irréductibles de $X_{f(x')}$ (resp. $Z_{f(x')}$) contenant $x'$ est
\emph{le même}, et en outre qu'il en est de même de l'ensemble des couples
\[
\bigl(\dim(Z_{f(x')}^{(\beta)}),\dim(X_{f(x')}^{(\alpha)})\bigr)
\]
pour les couples formés d'une composante irréductible $X_{f(x')}^{(\alpha)}$
de $X_{f(x')}$ et d'une composante irréductible $Z_{f(x')}^{(\beta)}$ de
$Z_{f(x')}$ contenue dans $X_{f(x')}^{(\alpha)}$ et contenant $x'$. On peut
évidemment pour cela remplacer $S$ par le sous-préschéma réduit $S'$ de $S$
ayant $\overline{\{f(x)\}}$ pour espace sous-jacent, et $X$ par
$X'=f^{-1}(S')$, les fibres de $X$ et $X'$ aux points de $S'$ étant les mêmes.
Autrement dit, on peut se borner au cas où $S$ est \emph{intègre} et où
$\eta=f(x)$ est son point générique.

Par hypothèse $Z_\eta$ est fermé dans $X_\eta$~; comme $Z$ est constructible,
il résulte de (8.3.11), appliqué suivant la méthode de (8.1.2, a)), que l'on
peut, en remplaçant au besoin $S$ par un voisinage ouvert de $\eta$, supposer
que $Z$ est \emph{fermé} dans $X$. Soient $X_i$ (resp. $Z_j$) les composantes
irréductibles de $X$ (resp. $Z$) contenant $x$~; en vertu de
($\mathbf 0_{\mathrm I},2.1.8$), les $X_i\cap X_\eta$ (resp.
$Z_j\cap X_\eta$) sont les composantes irréductibles de $X_\eta$ (resp.
$Z_\eta$) contenant $x$~; en vertu de (9.5.1), on peut supposer en outre, en
restreignant au besoin $S$ à un voisinage de $\eta$, que les $X_i$ (resp.
$Z_j$) sont exactement les composantes irréductibles de $X$ (resp. $Z$)
rencontrant $\overline{\{x\}}$ et que les
$(X_i)_s\cap\overline{\{x\}}$ et $(Z_j)_s\cap\overline{\{x\}}$ sont non vides
pour tout $s\in S$. Cela étant, il résulte encore de (9.7.1) que l'on peut
supposer, en restreignant $S$, que les couples $(i,j)$ tels que
$(Z_j)_s\subset(X_i)_s$ sont les mêmes pour tout $s\in S$. La conclusion
résulte alors de (9.5.6)~: pour tout $s$ assez voisin de $\eta$, toutes les
composantes irréductibles de $(X_i)_s$ (resp. $(Z_j)_s$) ont une même dimension
égale à celle de $(X_i)_\eta$ (resp. $(Z_j)_\eta$). En outre, si $(i,j)$ est un
couple tel que $Z_j\not\subset X_i$, $(X_i)_\eta$ ne contient pas le

\oldpage[IV-3]{90}

point générique de $(Z_j)_\eta$, donc
$\dim((X_i\cap Z_j)_\eta)<\dim((Z_j)_\eta)$~; par suite, la dimension commune
des composantes irréductibles de $(X_i)_s\cap(Z_j)_s$ est, pour tout $s\in S$,
strictement inférieure à la dimension commune des composantes irréductibles de
$(Z_j)_s$, ce qui prouve qu'aucune des composantes irréductibles de $(Z_j)_s$
n'est contenue dans une composante irréductible de $(X_i)_s$ pour $s\in S$.
C.Q.F.D.

\begin{proposition}[9.9.2]
\label{IV.9.9.2-fr}
Soient $f:X\to S$ un morphisme localement de présentation finie, $\mathcal F$
un $\mathcal O_X$-Module quasi-cohérent de présentation finie, $\Phi$ une
partie finie de $\mathbf Z\cup\{\pm\infty\}$. Les parties suivantes de $X$
sont localement constructibles~:
\begin{enumerate}
\item[(i)] L'ensemble des points $x\in X$ tels que $\mathcal F_{f(x)}$ soit
équidimensionnel au point $x$ (5.1.12).
\item[(ii)] L'ensemble des $x\in X$ tels que les longueurs géométriques de
$\mathcal F_{f(x)}$ relatives à $k(f(x))$, aux points génériques des composantes
irréductibles de $\operatorname{Supp}(\mathcal F_{f(x)})$ contenant $x$
(4.7.5), appartiennent à $\Phi$.
\item[(iii)] L'ensemble des $x\in X$ tels que les dimensions des cycles
premiers associés à $\mathcal F_{f(x)}$ et contenant $x$ appartiennent à
$\Phi$.
\item[(iv)] L'ensemble des $x\in X$ tels que $\mathcal F_{f(x)}$ soit
géométriquement réduit au point $x$ (4.6.17).
\item[(v)] L'ensemble des $x\in X$ tels que $\mathcal F_{f(x)}$ soit
géométriquement intègre au point $x$ (4.6.22).
\item[(vi)] L'ensemble des $x\in X$ tels que
$\operatorname{dim.proj}((\mathcal F_{f(x)})_x)\in\Phi$.
\item[(vii)] L'ensemble des $x\in X$ tels que
$\operatorname{coprof}((\mathcal F_{f(x)})_x)\in\Phi$.
\item[(viii)] L'ensemble des $x\in X$ tels que $\mathcal F_{f(x)}$ possède la
propriété $(S_k)$ au point $x$ (5.7.2).
\item[(ix)] L'ensemble des $x\in X$ tels que $\mathcal F_{f(x)}$ soit un
$\mathcal O_{X_{f(x)}}$-Module de Cohen-Macaulay au point $x$ (5.7.1).
\end{enumerate}
\end{proposition}

(i) Le support $Z$ de $\mathcal F$ est localement constructible et fermé dans
$X$ (8.9.1), et en considérant un sous-préschéma de $X$ ayant $Z$ pour espace
sous-jacent et qui soit de présentation finie sur $S$ (8.9.1), on voit que
l'assertion (i) est un cas particulier de (9.9.1, (iii)).

(ii) Toutes les propriétés envisagées sont locales sur $X$, et on se bornera
donc au cas où $X=\operatorname{Spec}(B)$ et $S=\operatorname{Spec}(A)$ sont
affines et $f$ un morphisme de présentation finie. On garde les notations du
début de la démonstration de (9.9.1), et on suppose en outre $A_0$ choisi de
sorte qu'il existe un $\mathcal O_{X_0}$-Module cohérent $\mathcal F_0$ tel que
$\mathcal F$ soit isomorphe à $\mathcal F_0\otimes_{A_0}A$. Alors ((4.2.7) et
(4.7.9)) l'ensemble des longueurs géométriques de
$(\mathcal F_0)_{f(x_0)}$ aux points génériques des composantes irréductibles de
$\operatorname{Supp}((\mathcal F_0)_{f(x_0)})$ qui contiennent $x_0$ est le
même que l'ensemble analogue pour $x$ et $\mathcal F_{f(x)}$~; autrement dit,
si $E$ (resp. $E_0$) est l'ensemble des $x\in X$ (resp. des $x_0\in X_0$)
vérifiant la condition (ii) de l'énoncé, on a $E=h^{-1}(E_0)$, et en vertu de
(1.8.2), on voit qu'on peut se borner à considérer le cas où $A$ est
\emph{noethérien}. Comme dans la démonstration de (9.9.1), on voit qu'on est
ramené à montrer que, pour tout $x\in X$, il existe un voisinage $V$ de $x$
dans $\overline{\{x\}}$ tel que, pour tout $x'\in V$, l'ensemble des longueurs
géométriques de $\mathcal F_{f(x')}$ aux points génériques des composantes
irréductibles de son support contenant $x'$ soit \emph{le même}. En outre, si
$S'$ est le sous-préschéma réduit de $S$ ayant $\overline{\{f(x)\}}$ pour espace
sous-jacent, et si $X'=f^{-1}(S')$, les fibres de $X$ et de $X'$ aux points de
$S'$ sont les mêmes,

\oldpage[IV-3]{91}

donc on peut remplacer $S$ par $S'$ et $X$ par $X'$, autrement dit supposer que
$S$ est \emph{intègre} et que $\eta=f(x)$ est son point générique.

Or, si l'on pose $Z=\operatorname{Supp}(\mathcal F)$, le même raisonnement que
dans la démonstration de (9.9.1) montre que, si $Z_i$ sont les composantes
irréductibles de $Z$ contenant $x$, on peut supposer que ce sont exactement les
composantes irréductibles de $Z$ rencontrant $\overline{\{x\}}$ et que
$(Z_i)_s\cap\overline{\{x\}}\neq\varnothing$ pour tout $s\in S$. La conclusion
résulte alors de (9.8.3) et (9.8.6).

(iii) On se ramène comme dans (ii) au cas où $S$ est noethérien et intègre et
$\eta=f(x)$ son point générique, en utilisant (4.2.7). Comme dans la
démonstration de (9.9.1), on est ramené à montrer qu'il existe un voisinage $V$
de $x$ dans $\overline{\{x\}}$ tel que, pour tout $x'\in V$, l'ensemble des
dimensions des cycles premiers associés à $\mathcal F_{f(x')}$ qui contiennent
$x'$ soit \emph{le même}. Or, si $Z'_i$ sont les adhérences dans $X$ des cycles
premiers associés à $\mathcal F_\eta$ qui contiennent $x$ (cf. (9.8.2)), il
résulte de (9.8.3) et (9.5.1) que pour tout $s$ assez voisin de $\eta$, les
cycles premiers associés à $\mathcal F_s$ qui rencontrent
$\overline{\{x\}}$ sont des composantes irréductibles des $(Z'_i)_s$ et, en
vertu de (9.5.6), toutes ces composantes irréductibles ont même dimension égale
à $\dim((Z'_i)_\eta)$, d'où la conclusion.

(iv) On raisonne comme dans (iii), en utilisant cette fois (4.7.11). Avec les
mêmes notations que dans (iii), les cycles premiers associés à $\mathcal F_s$
qui sont des composantes irréductibles de $(Z'_i)_s$ sont immergés si et
seulement si $(Z'_i)_\eta$ est un cycle premier associé immergé de
$\mathcal F_\eta$. On conclut déjà que si $x$ appartient à un (resp.
n'appartient à aucun) cycle premier associé immergé de $\mathcal F_\eta$,
l'ensemble des $x'\in\overline{\{x\}}$ tels que $x'$ appartienne à un (resp.
n'appartienne à aucun) cycle premier associé immergé de $\mathcal F_{f(x')}$
est un voisinage de $x$ dans $\overline{\{x\}}$. La conclusion résulte de cette
remarque, de la caractérisation des points où un Module est géométriquement
réduit (4.7.10), et de (ii).

(v) Compte tenu de (4.7.11), on se ramène encore au cas où $S$ est noethérien
et intègre et où $\eta=f(x)$ est son point générique. Soit $Y$ un
sous-préschéma fermé de $X$ ayant $\operatorname{Supp}(\mathcal F)$ pour espace
sous-jacent. Raisonnant comme au début de la démonstration de (9.7.7), on voit
qu'il existe une $A$-algèbre intègre finie $A'$ (de sorte que si
$S'=\operatorname{Spec}(A')$, le morphisme $S'\to S$ est fini et surjectif)
telle que si $Y'=Y_{(S')}$ et si $\eta'$ est le point générique de $S'$, les
composantes irréductibles de $Y'_{\eta'}$ soient géométriquement irréductibles.
D'autre part, si $X'=X_{(S')}$, le morphisme projection $h:X'\to X$ est fini
et surjectif, donc fermé~; par suite, si $x'$ est un point de $X'$ au-dessus de
$x$, on a $h(\overline{\{x'\}})=\overline{\{x\}}$ et l'image par $h$ d'un
voisinage de $x'$ dans $\overline{\{x'\}}$ est un voisinage de $x$ dans
$\overline{\{x\}}$. Compte tenu de (4.7.11) et posant
$\mathcal F'=\mathcal F_{(S')}$, on est donc ramené à prouver que si
$\mathcal F'_{\eta'}$ est (resp. n'est pas) géométriquement intègre au point
$x'$, l'ensemble des $y'\in\overline{\{x'\}}$ tels que
$\mathcal F'_{f'(y')}$ (où $f'=f_{(S')}:X'\to S'$) soit (resp. ne soit pas)
géométriquement intègre au point $y'$ est un voisinage de $x'$ dans
$\overline{\{x'\}}$. On peut donc se borner au cas où $X'=X$ et où les
composantes irréductibles de $Y_\eta$ sont géométriquement irréductibles~; si
l'on désigne par $Z_i$ des parties fermées de $X$ telles que les
$Z_i\cap X_\eta$ soient les

\oldpage[IV-3]{92}

composantes irréductibles de $Y_\eta$, il résulte de (9.7.7), (9.7.8) et
(9.5.3) que pour tout $s$ voisin de $\eta$, les $(Z_i)_s$ sont les composantes
irréductibles de $Y_s$ et qu'elles sont \emph{géométriquement irréductibles}.
Dire qu'en un point $y\in X_s$, $\mathcal F_s$ est géométriquement intègre
signifie alors que $\mathcal F_s$ est géométriquement réduit en ce point et en
outre que $y$ n'appartient qu'à \emph{un seul} des $(Z_i)_s$ (4.6.22). La
conclusion résulte donc d'une part de (iv) et d'autre part de (9.5.1) appliqué à
l'intersection de $\overline{\{x\}}$ et de chaque $Z_i$.

(vi) Gardant les mêmes notations que dans (ii), il résulte de (6.2.1) que l'on
a
\[
\operatorname{dim.proj}((\mathcal F_s)_x)=
\operatorname{dim.proj}(((\mathcal F_0)_{s_0})_{x_0})~;
\]
on peut donc encore se borner au cas où $A$ est \emph{noethérien}. En outre, on
se ramène encore à montrer que, pour tout $x\in X$, il existe un voisinage $V$
de $x$ dans $\overline{\{x\}}$ tel que, pour tout $x'\in V$, on ait
\[
\operatorname{dim.proj}((\mathcal F_{f(x')})_{x'})=
\operatorname{dim.proj}((\mathcal F_{f(x)})_x)~;
\]
et comme ci-dessus, on peut supposer que $S$ est \emph{intègre} et que
$\eta=f(x)$ est son point générique, de sorte que l'on a
$(\mathcal F_\eta)_x=\mathcal F_x$. En vertu du théorème de platitude générique
(6.9.1), on peut, en remplaçant au besoin $S$ par un voisinage ouvert de
$\eta$, supposer que le morphisme $f$ est \emph{plat} et que $\mathcal F$ est
$f$-plat~; on a alors
$\operatorname{dim.proj}((\mathcal F_{f(x')})_{x'})=
\operatorname{dim.proj}(\mathcal F_{x'})$ pour tout $x'\in X$ en vertu de
(6.2.3). Cela étant, en vertu de (6.11.1), on peut (en remplaçant au besoin $X$
par un voisinage ouvert de $x$) supposer que
$\operatorname{dim.proj}(\mathcal F_{x'})\leq
\operatorname{dim.proj}(\mathcal F_x)$ pour tout $x'\in X$. D'autre part, si
$\operatorname{dim.proj}(\mathcal F_x)=n$, il y a par hypothèse un
$\mathcal O_x$-module de type fini $M$ tel que
$\operatorname{Ext}_{\mathcal O_x}^n(\mathcal F_x,M)\neq0$
($\mathbf 0$, 17.2.4). Or, il existe un $\mathcal O_X$-Module cohérent
$\mathcal G$ tel que $M=\mathcal G_x$ (en remplaçant au besoin $X$ par un
voisinage ouvert de $x$ ($\mathbf 0_{\mathrm I}$, 5.3.8))~; en vertu de
(T, 4.2.2), on a donc
$(\mathcal{E}xt_{\mathcal O_X}^n(\mathcal F,\mathcal G))_x\neq0$. Mais
$\mathcal{E}xt_{\mathcal O_X}^n(\mathcal F,\mathcal G)$ est un
$\mathcal O_X$-Module cohérent ($\mathbf 0_{\mathrm{III}}$, 12.3.3), donc son
support $Y$ est fermé ($\mathbf 0_{\mathrm I}$, 5.2.2)~; comme il contient $x$,
il contient aussi $\overline{\{x\}}$, d'où l'on conclut (en appliquant encore
(T, 4.2.2)) que l'on a
$\operatorname{dim.proj}(\mathcal F_{x'})\geq n$ pour tout
$x'\in\overline{\{x\}}$, ce qui achève de prouver l'assertion dans le cas (vi).

(vii) Comme $B$ est une $A$-algèbre de type fini, $X$ est $S$-isomorphe à un
\emph{sous-schéma fermé} d'un $S$-schéma de la forme
$Y=\operatorname{Spec}(A[T_1,\ldots,T_r])$~; si $j:X\to Y$ est l'injection
canonique, et $\mathcal G=j_*(\mathcal F)$, on a
$\mathcal G_s=(j_s)_*(\mathcal F_s)$ pour tout $s\in S$, et, en vertu de
($\mathbf 0$, 16.4.11), on peut se borner à démontrer l'assertion pour $Y$ et
$\mathcal G$. Autrement dit, on peut supposer que $B=A[T_1,\ldots,T_r]$, de
sorte que chacun des schémas $X_s=\operatorname{Spec}(k(s)[T_1,\ldots,T_r])$
est \emph{régulier} ($\mathbf 0$, 17.3.7). Soit alors
$W=\operatorname{Supp}(\mathcal F)$, de sorte que
$W_s=\operatorname{Supp}(\mathcal F_s)$ (\textbf{I}, 9.1.13)~; on a, par
(6.11.2.1)
\begin{equation}
\tag{9.9.2.1}
\label{IV.9.9.2.1-fr}
\operatorname{coprof}((\mathcal F_{f(x)})_x)=
\operatorname{dim.proj}((\mathcal F_{f(x)})_x)-
\operatorname{codim}_x(W_{f(x)},X_{f(x)}).
\end{equation}
Mais comme $W$ est constructible (8.9.1) et chaque $W_s$ fermé, il résulte de
(vi) et de (9.9.1, (ii)) que les deux fonctions du second membre de (9.9.2.1)
sont constructibles~; il en est donc de même de leur différence ce qui achève
de prouver la proposition dans le cas (vii).

(viii) Soit $U_n$ l'ensemble des $x\in X$ tels que
$\operatorname{coprof}((\mathcal F_{f(x)})_x)\leq n$, et posons
$Z_n=X-U_n$~; il résulte de (vii) que les $Z_n$ sont constructibles~; en outre,
comme

\oldpage[IV-3]{93}

la fonction $x\mapsto\dim_x(W_{f(x)})$ est constructible en vertu de
(9.9.1, (i)), elle ne prend qu'un nombre fini de valeurs, donc les nombres
$\dim(W_{f(x)})$ ont une borne supérieure finie $m$ lorsque $x$ parcourt $X$~;
comme
$\operatorname{coprof}((\mathcal F_{f(x)})_x)\leq
\dim((\mathcal F_{f(x)})_x)\leq\dim(W_{f(x)})$, on voit que l'on a
$Z_n=\varnothing$ pour $n\geq m$. Enfin, il résulte de (6.11.2, (i)) que pour
tout $n$ et tout $s\in S$, $(Z_n)_s$ est \emph{fermé} dans $X_s$. D'après
(5.7.4), l'ensemble des $x\in X$ où $(\mathcal F_{f(x)})_x$ possède la
propriété $(S_k)$ est l'ensemble des $x\in X$ vérifiant toutes les relations
\begin{equation}
\tag{9.9.2.2}
\label{IV.9.9.2.2-fr}
\operatorname{codim}_x((Z_n)_{f(x)},W_{f(x)})>n+k
\end{equation}
pour tout $n\geq0$~; comme cette relation est automatiquement vérifiée pour
$n\geq m$, on n'a à considérer les relations (9.9.2.2) que pour $0\leq n<m$.
Mais en vertu de (9.9.1, (ii)), l'ensemble $V_{n,k}$ des $x$ vérifiant (9.9.2.2)
est constructible, et il en est de même de l'intersection de ces ensembles pour
$0\leq n<m$.

(ix) L'assertion résulte ici aussitôt de (vii), l'ensemble des $x\in X$ où
$(\mathcal F_{f(x)})_x$ est un module de Cohen-Macaulay étant défini par la
relation $\operatorname{coprof}((\mathcal F_{f(x)})_x)=0$.

\begin{corollary}[9.9.3]
\label{IV.9.9.3-fr}
Soient $f:X\to S$ un morphisme de présentation finie, $\mathcal F$ un
$\mathcal O_X$-Module de présentation finie, $\mathbf P$ l'une des propriétés
(i) à (ix) de (9.9.2). Alors l'ensemble des $s\in S$ tels que la propriété
$\mathbf P$ soit vraie en tous les points $x\in X_s$, est localement
constructible dans $S$.
\end{corollary}

En effet, son complémentaire est l'image par $f$ du complémentaire de
l'ensemble $E$ des points où $\mathbf P$ est vraie. Comme $E$ est localement
constructible dans $X$, il en est de même de $X-E$, donc $f(X-E)$ est
localement constructible dans $S$ en vertu du théorème de Chevalley (1.8.4).

\begin{proposition}[9.9.4]
\label{IV.9.9.4-fr}
Soit $f:X\to S$ un morphisme localement de présentation finie. L'ensemble des
$x\in X$ tels que $X_{f(x)}$ ait au point $x$ l'une (déterminée) des propriétés
suivantes~:
\begin{enumerate}
\item[(i)] être géométriquement régulier~;
\item[(ii)] posséder la propriété $(R_k)$ géométrique~;
\item[(iii)] être géométriquement normal~;
\item[(iv)] être géométriquement réduit (i.e. séparable)~;
\item[(v)] être géométriquement ponctuellement intègre~;
\end{enumerate}
est une partie localement constructible de $X$.
\end{proposition}

Pour les propriétés (iv) et (v), cela résulte de (9.9.2, (iv) et (v)) appliqué à
$\mathcal F=\mathcal O_X$. Pour les autres propriétés, compte tenu de (6.7.8),
on se ramène, comme au début de (9.9.2) au cas où $S$ est noethérien et intègre
et de point générique $\eta=f(x)$. En outre, en vertu du théorème de platitude
générique (6.9.1), on peut, en remplaçant $S$ par un voisinage ouvert de
$\eta$, supposer que $f$ est un morphisme \emph{plat}. Dire que $X_{f(y)}$ est
géométriquement régulier au point $y$ signifie alors que $f$ est
\emph{régulier} au point $y$ (6.8.1), et on sait que l'ensemble $L$ de ces
points est \emph{ouvert} dans $X$ (6.8.7), ce qui prouve la proposition dans le
cas (i).

Pour prouver le cas (ii), posons $F=X-L$, qui est \emph{fermé} dans $X$. Dire
que $X_s$

\oldpage[IV-3]{94}

vérifie la propriété géométrique $(R_k)$ au point $y\in f^{-1}(s)$ signifie, ou
bien que l'on a $y\in L$, ou bien que les points génériques $z_i$ des
composantes irréductibles de l'ensemble fermé $F_s$ qui contiennent $y$
vérifient la relation $\dim(\mathcal O_{X_s,z_i})\geq k+1$ (compte tenu de
(4.2.7) et (5.2.3))~; autrement dit, les points $y\in F_s$ où $X_s$ vérifie la
propriété $(R_k)$ géométrique sont ceux tels que
$\operatorname{codim}_y(F_s,X_s)\geq k+1$ (5.1.2). La conclusion résulte donc
de (i) et de (9.9.1, (ii)).

Enfin, dans le cas (iii), la conclusion résulte de (ii), de (9.9.2, (viii)), du
fait que la propriété $(S_k)$ est stable par extension du corps de base (6.7.1)
et enfin du critère de Serre (5.8.6).

\begin{corollary}[9.9.5]
\label{IV.9.9.5-fr}
Soit $f:X\to S$ un morphisme de présentation finie, et notons $\mathbf P$
l'une des propriétés (i) à (v) de (9.9.4). Alors l'ensemble des $s\in S$ tels
que la propriété $\mathbf P$ soit vraie en tous les points $x\in X_s$, est
localement constructible dans $S$.
\end{corollary}

La démonstration à partir de (9.9.4) est la même que celle de (9.9.3) à partir
de (9.9.2).

\begin{proposition}[9.9.6]
\label{IV.9.9.6-fr}
Soient $f:X\to S$ un morphisme localement de présentation finie,
$\mathcal L_\bullet$ un complexe formé de $\mathcal O_X$-Modules
quasi-cohérents de présentation finie~; pour tout $s\in S$, soit
$(\mathcal L_\bullet)_s$ le complexe
$\mathcal L_\bullet\otimes_{\mathcal O_S}k(s)$ de
$\mathcal O_{X_s}$-Modules cohérents de type fini. Alors, pour un entier $n$
donné, l'ensemble des $x\in X$ tels que
$(\mathcal H_n((\mathcal L_\bullet)_{f(x)}))_x=0$ est localement constructible
dans $X$.
\end{proposition}

On peut se borner au cas où $\mathcal L_i=0$ sauf pour $i=0$, 1 ou 2, et où
$n=1$. En outre, la question étant locale sur $X$, on peut se borner au cas où
$S=\operatorname{Spec}(A)$ et $X=\operatorname{Spec}(B)$ sont affines, $B$
étant une $A$-algèbre de présentation finie. Il existe alors un sous-anneau
noethérien $A_0$ de $A$, un $A_0$-préschéma de type fini $X_0$ et un complexe
$\mathcal L_\bullet^{(0)}$ de $\mathcal O_{X_0}$-Modules cohérents, nuls sauf en
dimensions 0, 1 et 2 et tels que
$X=X_0\otimes_{A_0}A$ et
$\mathcal L_\bullet=\mathcal L_\bullet^{(0)}\otimes_{A_0}A$. Pour tout
$s\in S$, si $s_0$ est la projection de $s$ dans
$S_0=\operatorname{Spec}(A_0)$, on a
$X_s=(X_0)_{s_0}\otimes_{k(s_0)}k(s)$, et le morphisme projection
$X_s\to(X_0)_{s_0}$ est fidèlement plat~; on en conclut que l'on a
\[
\mathcal H_n((\mathcal L_\bullet)_s)=
\mathcal H_n((\mathcal L_\bullet^{(0)})_{s_0})
\otimes_{k(s_0)}k(s),
\]
et par suite, si $E$ (resp. $E_0$) est l'ensemble des $x\in X$ (resp.
$x_0\in X_0$) tels que
\[
(\mathcal H_n((\mathcal L_\bullet)_{f(x)}))_x=0
\quad\text{(resp.}\quad
(\mathcal H_n((\mathcal L_\bullet^{(0)})_{f_0(x_0)}))_{x_0}=0\text{)},
\]
on a $E=h^{-1}(E_0)$, où $h:X\to X_0$ est la projection canonique. En vertu
de (1.8.2), on peut donc se borner au cas où $A$ est \emph{noethérien}~; il
s'agit de voir ($\mathbf 0_{\mathrm{III}},9.2.3$) que si $x\in X$ est tel que
$(\mathcal H_n((\mathcal L_\bullet)_{f(x)}))_x=0$ (resp. $\neq0$), il existe un
voisinage ouvert $V$ de $x$ dans $\overline{\{x\}}$ tel que, pour tout
$x'\in V$, on ait
$(\mathcal H_n((\mathcal L_\bullet)_{f(x')}))_{x'}=0$ (resp. $\neq0$).
Remplaçant $S$ par le sous-préschéma réduit de $S$ ayant
$\overline{\{f(x)\}}$ pour espace sous-jacent, on peut supposer que $S$ est
intègre et que $f(x)=\eta$ est son point générique. Alors, en restreignant $S$
à un voisinage ouvert de $\eta$, on peut supposer que pour tout $s\in S$, on a
$(\mathcal H_n(\mathcal L_\bullet))_s=
\mathcal H_n((\mathcal L_\bullet)_s)$ (9.4.3), et par suite, si $Z$ est le
support de $\mathcal H_n(\mathcal L_\bullet)$, le support de
$\mathcal H_n((\mathcal L_\bullet)_s)$ est $Z_s=Z\cap X_s$
(\textbf{I}, 9.1.13.1). L'hypothèse est que
$Z_\eta\cap\overline{\{x\}}=\varnothing$ (resp. $\neq\varnothing$). Comme
$Z\cap\overline{\{x\}}$ est fermé dans l'espace noethérien $X$, on conclut de
(9.5.1) qu'il y a un voisinage $U$ de $\eta$ dans $S$ tel que, pour tout
$s\in U$, on ait $Z_s\cap\overline{\{x\}}=\varnothing$ (resp.
$\neq\varnothing$). L'ensemble $V=f^{-1}(U)\cap\overline{\{x\}}$ répond donc à
la question.

% FIN SOURCE ega4/ega4-9-fr.tex

%
% DEBUT SOURCE ega4/ega4-10-fr.tex
\oldpage[IV-3]{95}
\section{Préschémas de Jacobson}
\label{section:IV.10-fr}

On a déjà eu l'occasion d'observer (5.2.5) que même les préschémas
\emph{excellents} (7.8.5) ne se comportent pas toujours comme les
«~variétés~» de la Géométrie algébrique classique, notamment en ce qui
concerne les questions de dimension~; c'est ainsi que si $X$ est le spectre
d'un anneau de valuation discrète complet, l'ensemble des points fermés de
$X$ (réduit à un seul élément) n'est pas partout dense dans $X$, et que son
complémentaire (réduit à un seul élément) est un ouvert partout dense dans
$X$ mais dont la dimension est nulle, donc $<\dim(X)$. On examine dans ce
paragraphe des conditions générales moyennant lesquelles de tels phénomènes
ne se présentent pas~; il en résulte un comportement plus satisfaisant à
certains égards pour les relations entre dimension, codimension, profondeur
et coprofondeur dans ces préschémas (10.6 et 10.8). En outre, dans les
préschémas considérés, le fait que l'ensemble des points fermés soit partout
dense (et même «~très dense~» (10.1.3)) permet de ne considérer que ces points
dans beaucoup de démonstrations~; on rejoint ainsi le point de vue classique
des «~variétés algébriques~» qui, de notre point de vue, sont les ensembles
de points fermés des \emph{préschémas algébriques sur un corps}, et on fait le
lien entre le langage des schémas et celui des «~variétés~» ou «~espaces
algébriques~» de Serre (10.9 et 10.10).

\subsection{Parties très denses d'un espace topologique}
\label{subsection:IV.10.1-fr}

\begin{env}[10.1.1]
\label{IV.10.1.1-fr}
On dit qu'une partie d'un espace topologique $X$ est
\emph{quasi-constructible} si elle est réunion finie de parties localement
fermées de $X$. On dit qu'une partie $T$ de $X$ est \emph{localement
quasi-constructible} si pour tout $x\in X$, il existe un voisinage ouvert $V$
de $x$ tel que $T\cap V$ soit quasi-constructible dans $V$. Il est clair que
toute partie quasi-constructible de $X$ est localement quasi-constructible~:
la réciproque est vraie si $X$ est quasi-compact. Le raisonnement de
($\mathbf 0_{\mathrm{III}}$, 9.1.3), où l'on supprime le mot
«~rétrocompact~», montre que l'ensemble des parties quasi-constructibles
(resp. localement quasi-constructibles) de $X$, que nous désignerons par
$\mathfrak{Qc}(X)$ (resp. $\mathfrak{Lqc}(X)$) est stable par réunion finie,
intersection finie et passage aux complémentaires. Si $f:X\to Y$ est une
application continue, il résulte aussitôt de ces définitions que, pour toute
partie quasi-constructible (resp. localement quasi-constructible) $Z$ de $Y$,
$f^{-1}(Z)$ est quasi-constructible (resp. localement quasi-constructible)
dans $X$.

Les parties constructibles (resp. localement constructibles) de $X$ sont
évidemment quasi-constructibles (resp. localement quasi-constructibles)~; la
réciproque est vraie lorsque $X$ est noethérien (resp. localement
noethérien).

Dans ce qui suit, nous désignerons respectivement par $\mathfrak O(X)$,
$\mathfrak F(X)$, $\mathfrak{Lf}(X)$, $\mathfrak C(X)$, $\mathfrak{Lc}(X)$
l'ensemble des parties de $X$ qui sont respectivement ouvertes, fermées,
localement fermées, constructibles, localement constructibles.
\end{env}

\oldpage[IV-3]{96}
\begin{proposition}[10.1.2]
\label{IV.10.1.2-fr}
Soient $X$ un espace topologique, $X_0$ une partie de $X$. Les conditions
suivantes sont équivalentes~:
\begin{enumerate}[label=\emph{\alph*)}]
  \item Pour toute partie localement fermée $Z\neq\varnothing$ de $X$, on a
  $Z\cap X_0\neq\varnothing$.
  \item[\emph{a$'$)}] Pour toute partie fermée $Z$ de $X$, on a
  $Z=\overline{Z\cap X_0}$.
  \item Pour toute partie $Z\neq\varnothing$ de $X$ localement
  quasi-constructible, on a $Z\cap X_0\neq\varnothing$.
  \item[\emph{b$'$)}] Pour toute partie localement quasi-constructible $Z$ de
  $X$, on a $Z\subset\overline{Z\cap X_0}$ (autrement dit $Z\cap X_0$ est
  \emph{dense} dans $Z$).
  \item L'application $U\mapsto U\cap X_0$ de $\mathfrak O(X)$ dans
  $\mathfrak O(X_0)$ est injective (donc bijective).
  \item[\emph{c$'$)}] L'application $F\mapsto F\cap X_0$ de $\mathfrak F(X)$
  dans $\mathfrak F(X_0)$ est injective (donc bijective).
  \item[\emph{c$''$)}] L'application $Z\mapsto Z\cap X_0$ de
  $\mathfrak{Qc}(X)$ dans $\mathfrak{Qc}(X_0)$ est injective (donc bijective).
  \item[\emph{c$'''$)}] L'application $Z\mapsto Z\cap X_0$ de
  $\mathfrak{Lqc}(X)$ dans $\mathfrak{Lqc}(X_0)$ est injective.
\end{enumerate}
En outre, lorsque ces conditions sont vérifiées, l'application
$Z\mapsto Z\cap X_0$ de $\mathfrak{Lqc}(X)$ dans $\mathfrak{Lqc}(X_0)$ est
bijective.
\end{proposition}

Notons que les assertions de surjectivité dans \emph{c)} et \emph{c$'$)} sont
triviales~: elles entraînent que toute partie localement fermée dans $X_0$ est
trace sur $X_0$ d'une partie localement fermée dans $X$, donc l'application
$\mathfrak{Qc}(X)\to\mathfrak{Qc}(X_0)$ définie dans \emph{c$''$)} est aussi
surjective.

Nous prouverons les implications
\[
c''')\Rightarrow c'')\Rightarrow c')\Rightarrow c)\Rightarrow b)
\Rightarrow b')\Rightarrow a')\Rightarrow a)\Rightarrow c''').
\]
Les trois premières sont triviales. Pour voir que \emph{c)} implique
\emph{b)}, notons d'abord que \emph{c)} entraîne que $X_0$ est \emph{dense}
dans $X$. Remplaçant $X$ et $X_0$ par un ouvert convenable $U$ de $X$ et par
$U\cap X_0$ respectivement, on peut donc se borner au cas où $Z$ est
localement fermée dans $X$, donc $Z=V\cap\complement W$, où $V$ et $W$ sont
ouverts dans $X$~; l'hypothèse $Z\neq\varnothing$ signifie que
$V\not\subset W$, ou encore $V\cup W\neq W$. En vertu de \emph{c)}, on a donc
$(V\cup W)\cap X_0\neq W\cap X_0$, donc $V\cap X_0\not\subset W$, et par
suite $(V\cap\complement W)\cap X_0\neq\varnothing$.

Pour voir que \emph{b)} entraîne \emph{b$'$)}, il suffit d'appliquer
\emph{b)} à $Z\cap U$, où $U$ est un voisinage ouvert arbitraire d'un point de
$Z$. Comme $\overline{Z\cap X_0}\subset\overline{Z}$, il est trivial
que \emph{b$'$)} implique \emph{a$'$)}. Pour montrer que \emph{a$'$)} entraîne
\emph{a)}, remarquons que si $Z$ est localement fermée dans $X$, on peut
écrire $Z=F-F'$ où $F$ et $F'$ sont fermés dans $X$ et $F'\subset F$~; d'où
$Z\cap X_0=(F\cap X_0)-(F'\cap X_0)$. Si l'on avait
$Z\cap X_0=\varnothing$, on en déduirait $F\cap X_0=F'\cap X_0$, donc
$F=F'$ en vertu de \emph{a$'$)}, c'est-à-dire $Z=\varnothing$.

Pour voir que \emph{a)} entraîne \emph{c$'''$)}, il suffit de montrer que si
$Z\neq\varnothing$ est localement quasi-constructible, on a
$Z\cap X_0\neq\varnothing$~: en effet, la relation
$Z'\cap X_0=Z''\cap X_0$ est équivalente à
$((Z'\cup Z'')-(Z'\cap Z''))\cap X_0=\varnothing$. Autrement dit, il suffit de
prouver que \emph{a)} entraîne \emph{b)}~; en outre, remplaçant $X$ et $X_0$
par un ouvert $U$ de $X$ et par $U\cap X_0$ respectivement, on est bien
ramené au cas où $Z$ est localement fermée dans $X$, d'où la conclusion.

Reste à montrer que l'application
$\mathfrak{Lqc}(X)\to\mathfrak{Lqc}(X_0)$ est surjective. Soit donc $Z_0$ une
partie localement quasi-constructible de $X_0$~: il y a un recouvrement
$(V_\alpha)$ de $X_0$ par des ouverts de $X_0$ tel que $Z_0\cap V_\alpha$ soit
quasi-constructible dans $V_\alpha$ (et
\oldpage[IV-3]{97}
donc aussi dans $X_0$). En vertu de \emph{c)}, il y a pour tout $\alpha$ un
seul ouvert $U_\alpha$ de $X$ tel que $X_0\cap U_\alpha=V_\alpha$, et en
vertu de \emph{c$''$)} un seul ensemble quasi-constructible $Z_\alpha$ dans
$X$ tel que $Z_0\cap V_\alpha=X_0\cap Z_\alpha$. Si $\alpha$ et $\beta$ sont
deux indices quelconques, on a
\[
Z_\beta\cap U_\alpha\cap X_0
=U_\alpha\cap Z_0\cap V_\beta
=V_\alpha\cap Z_0\cap V_\beta
=Z_\alpha\cap X_0\cap V_\beta
\subset Z_\alpha\cap X_0~;
\]
comme $Z_\beta\cap U_\alpha$ et $Z_\alpha$ sont quasi-constructibles dans
$X$, il résulte de \emph{c$''$)} que
$Z_\beta\cap U_\alpha\subset Z_\alpha$. Si l'on pose
$Z=\bigcup_\alpha Z_\alpha$, on a donc $Z\cap U_\alpha=Z_\alpha$ pour tout
$\alpha$~; d'ailleurs, comme les $V_\alpha$ recouvrent $X_0$, il résulte de
\emph{c)} que les $U_\alpha$ recouvrent $X$, et l'on voit donc que $Z$ est
localement quasi-constructible dans $X$ et $Z_0=Z\cap X_0$.

\begin{definition}[10.1.3]
\label{IV.10.1.3-fr}
Lorsqu'une partie $X_0$ d'un espace topologique vérifie les conditions
équivalentes de (10.1.2), on dit que $X_0$ est \emph{très dense} dans $X$.
\end{definition}

On a déjà vu au cours de la démonstration de (10.1.2) que $X_0$ est alors
\emph{dense} dans $X$.

\begin{corollary}[10.1.4]
\label{IV.10.1.4-fr}
Si $X_0$ est très dense dans $X$ et $U$ une partie ouverte de $X$,
$U\cap X_0$ est très dense dans $U$. Inversement, si $(U_\alpha)$ est un
recouvrement ouvert de $X$ tel que $U_\alpha\cap X_0$ soit très dense dans
$U_\alpha$ pour tout $\alpha$, alors $X_0$ est très dense dans $X$.
\end{corollary}

Comme toute partie localement fermée dans $U$ est localement fermée dans $X$,
la première assertion résulte du critère \emph{a)} de (10.1.2)~; il en est de
même de la seconde, car si $Z\neq\varnothing$ est localement fermée dans $X$,
$Z\cap U_\alpha$ est localement fermée dans $U_\alpha$ pour tout $\alpha$,
et $Z\cap U_\alpha\neq\varnothing$ pour un $\alpha$ au moins.

\subsection{Quasi-homéomorphismes}
\label{subsection:IV.10.2-fr}

\begin{proposition}[10.2.1]
\EGAReaderTarget{ega4.c04.prop.10-2-1}{IV.10.2.1-fr}
Soient $X_0$, $X$ deux espaces topologiques, $f:X_0\to X$ une application
continue. Les conditions suivantes sont équivalentes~:
\begin{enumerate}[label=\emph{\alph*)}]
  \item L'application $U\mapsto f^{-1}(U)$ de $\mathfrak O(X)$ dans
  $\mathfrak O(X_0)$ est bijective.
  \item[\emph{a$'$)}] L'application $F\mapsto f^{-1}(F)$ de $\mathfrak F(X)$
  dans $\mathfrak F(X_0)$ est bijective.
  \item La topologie de $X_0$ est image réciproque par $f$ de celle de $X$, et
  $f(X_0)$ est très dense (10.1.3) dans $X$.
  \item Le foncteur $\mathcal F\mapsto f^*(\mathcal F)$ de la catégorie des
  faisceaux d'ensembles (resp. des faisceaux de groupes abéliens) sur $X$ dans
  la catégorie des faisceaux d'ensembles (resp. des faisceaux de groupes
  abéliens) sur $X_0$ est une équivalence de catégories.
\end{enumerate}
\end{proposition}

Il est clair que \emph{a)} et \emph{a$'$)} sont équivalentes, et \emph{a)}
implique que la topologie de $X_0$ est image réciproque par $f$ de celle de
$X$. D'autre part, si $f(X_0)$ n'est pas très dense dans $X$, il y a deux
ouverts distincts $U_1$, $U_2$ de $X$ tels que
$U_1\cap f(X_0)=U_2\cap f(X_0)$, et par suite
$f^{-1}(U_1)=f^{-1}(U_2)$, ce qui achève de montrer que \emph{a)} entraîne
\emph{b)}. Inversement, la condition \emph{b)} implique que les applications
$U\mapsto U\cap f(X_0)$ de $\mathfrak O(X)$ dans $\mathfrak O(f(X_0))$ et
$V\mapsto f^{-1}(V)$ de $\mathfrak O(f(X_0))$ dans $\mathfrak O(X_0)$ sont
bijectives, donc il en est de même de leur composée $U\mapsto f^{-1}(U)$.

Pour voir que \emph{a)} entraîne \emph{c)}, il suffit d'appliquer la définition
de $f^*(\mathcal F)$ ($\mathbf 0_{\mathrm I}$, 3.7.1) et les axiomes des
faisceaux~: \emph{a)} entraîne que pour tout ouvert $U$ de $X$, l'application
canonique
$\Gamma(U,\mathcal F)\to\Gamma(f^{-1}(U),f^*(\mathcal F))$ est une bijection
fonctorielle en $\mathcal F$, d'où \emph{c)}. Il reste à montrer que
\emph{c)} entraîne \emph{b)}.

\oldpage[IV-3]{98}
Supposons d'abord que $f(X_0)$ ne soit pas très dense dans $X$~; il existe
alors deux parties fermées distinctes $Y\supset Y'$ de $X$ telles que
$Y\cap f(X_0)=Y'\cap f(X_0)$. Soit $\mathcal F$ (resp. $\mathcal F'$) le
faisceau de groupes abéliens sur $X$ image directe par l'injection canonique
$Y\to X$ (resp. $Y'\to X$) du faisceau simple associé au préfaisceau constant
$\mathbf Z$ sur $Y$ (resp. $Y'$)~; la définition de $f^*$
($\mathbf 0_{\mathrm I}$, 3.7.1) montre que $f^*(\mathcal F)$ est isomorphe à
$f^*(\mathcal F')$ mais $\mathcal F$ n'est pas isomorphe à $\mathcal F'$,
donc la condition \emph{c)} n'est pas vérifiée. (On notera que le foncteur
$f^*$ n'est même pas alors \emph{fidèle}, car si $u$ et $v$ sont
l'automorphisme identique et l'endomorphisme nul de $\mathcal F/\mathcal F'$,
$f^*(u)$ et $f^*(v)$ sont égaux.)

Montrons maintenant que \emph{c)} entraîne que la topologie de $X_0$ est
l'image réciproque de celle de $X$ par $f$. Notons que si la condition
\emph{c)} est vérifiée pour la catégorie des faisceaux d'ensembles, elle l'est
aussi pour la catégorie des faisceaux de groupes abéliens, car il résulte
aussitôt des définitions ($\mathbf 0_{\mathrm{III}}$, 8.2.5) que cette dernière
n'est autre que la catégorie des \emph{objets en groupes abéliens} dans la
catégorie des faisceaux d'ensembles. Il suffira donc de prouver notre assertion
en supposant \emph{c)} vérifiée pour les catégories de faisceaux de groupes
abéliens sur $X$ et $X_0$. Or, si $\mathbf Z_X$ désigne le faisceau simple sur
$X$ associé au préfaisceau constant égal à $\mathbf Z$, on a un isomorphisme
canonique
\[
\Gamma(X,\mathcal F)\xrightarrow{\sim}\operatorname{Hom}(\mathbf Z_X,\mathcal F)
\]
fonctoriel en $\mathcal F$ (même raisonnement que dans
($\mathbf 0_{\mathrm I}$, 5.1.1))~; comme il est clair par définition
($\mathbf 0_{\mathrm I}$, 3.7.1) que $f^*(\mathbf Z_X)=\mathbf Z_{X_0}$,
l'hypothèse que $f^*$ est une équivalence entraîne que l'homomorphisme
canonique $\Gamma(X,\mathcal F)\to\Gamma(X_0,f^*(\mathcal F))$ est
\emph{bijectif} pour tout faisceau de groupes abéliens $\mathcal F$ sur $X$.
Or, soit $Y_0$ une partie fermée de $X_0$~; soit $\mathcal F_0$ le faisceau
sur $X$ image directe par l'injection canonique $Y_0\to X_0$ du faisceau
simple associé au préfaisceau constant $\mathbf Z$ sur $Y_0$. Comme $f^*$ est
une équivalence, $\mathcal F_0$ est isomorphe à un faisceau de la forme
$f^*(\mathcal F)$, où $\mathcal F$ est un faisceau de groupes abéliens sur
$X$. Considérons alors la section $s_0$ de $\mathcal F_0$ au-dessus de $X_0$
telle que $(s_0)_{x_0}=1_{x_0}$ si $x_0\in Y_0$,
$(s_0)_{x_0}=0_{x_0}$ si $x_0\notin Y_0$. Les remarques précédentes montrent
qu'il existe une section $s$ de $\mathcal F$ au-dessus de $X$ telle que
$(s_0)_{x_0}=s_{f(x_0)}$ pour tout $x_0\in X_0$ (les fibres
$(\mathcal F_0)_{x_0}$ et $\mathcal F_{f(x_0)}$ étant canoniquement
identifiées ($\mathbf 0_{\mathrm I}$, 3.7.1))~; cela entraîne que $Y_0$ est
l'image réciproque par $f$ de l'ensemble $Y$ des $x\in X$ tels que
$s_x\neq0_x$, et $y$ est fermé dans $X$. C.Q.F.D.

\begin{definition}[10.2.2]
\label{IV.10.2.2-fr}
Lorsqu'une application continue $f:X_0\to X$ vérifie les conditions
équivalentes de (10.2.1), on dit que $f$ est un \emph{quasi-homéomorphisme} de
$X_0$ dans $X$.
\end{definition}

En vertu de (10.2.1, \emph{b)}), dire qu'une partie $X_0$ d'un espace
topologique $X$ est \emph{très dense} dans $X$ signifie donc que l'injection
canonique $X_0\to X$ est un \emph{quasi-homéomorphisme}.

\begin{corollary}[10.2.3]
\label{IV.10.2.3-fr}
Le composé de deux quasi-homéomorphismes est un quasi-homéomorphisme.
\end{corollary}

Cela résulte aussitôt de (10.2.1, \emph{a)}).

\begin{corollary}[10.2.4]
\label{IV.10.2.4-fr}
Soient $f:X\to Y$ un quasi-homéomorphisme, $Y'$ une partie localement
quasi-constructible de $Y$, $X'=f^{-1}(Y')$~; alors la restriction
$f'=f|X':X'\to Y'$ est un quasi-homéomorphisme.
\end{corollary}

Il est clair en vertu de (10.2.1, \emph{b)}) que la topologie induite sur $X'$
par celle de $X$ est l'image réciproque par $f'$ de la topologie induite sur
$Y'$ par celle de $Y$. D'autre part,
\oldpage[IV-3]{99}
soit $Z\neq\varnothing$ une partie fermée \emph{dans} $Y'$, et $z\in Z$~; il
y a un voisinage ouvert $U$ de $z$ dans $Y$ tel que $U\cap Y'$ soit réunion
d'un nombre fini de parties $Y'_j$ fermées dans $U$~; si $j$ est un indice tel
que $z\in Y'_j$, $Z\cap Y'_j$ est donc fermé dans $U$. Puisque $f(X)\cap U$
est très dense dans $U$ (10.1.4), $Z\cap Y'_j\cap f(X)$ n'est pas vide
(10.1.2, \emph{a)}), ni \emph{a fortiori} $Z\cap f(X)$~; mais comme
$Z\subset Y'$, on a $Z\cap f(X)=Z\cap f'(X')$~; le critère
(10.1.2, \emph{a)}) montre donc que $f'(X')$ est très dense dans $Y'$, et
l'on conclut à l'aide de (10.2.1, \emph{b)}).

\begin{corollary}[10.2.5]
\label{IV.10.2.5-fr}
Soient $f:X\to Y$ une application continue, $(V_\alpha)$ un recouvrement
ouvert de $Y$. Si, pour tout $\alpha$, la restriction
$f^{-1}(V_\alpha)\to V_\alpha$ de $f$ est un quasi-homéomorphisme, alors $f$
est un quasi-homéomorphisme.
\end{corollary}

Cela résulte aussitôt du critère (10.2.1, \emph{b)}) et de (10.1.4).

\begin{corollary}[10.2.6]
\label{IV.10.2.6-fr}
Soient $f:X\to Y$ un quasi-homéomorphisme, $Y'$ une partie localement
quasi-constructible de $Y$, $X'=f^{-1}(Y')$. Pour que $Y'$ soit quasi-compact
(resp. noethérien, resp. rétrocompact dans $Y$), il faut et il suffit que $X'$
soit quasi-compact (resp. noethérien, resp. rétrocompact dans $X$).
\end{corollary}

Démontrons d'abord les deux premières assertions~; en vertu de (10.2.4), on
peut se borner au cas où $Y'=Y$. Dire que $X$ est quasi-compact (resp.
noethérien) signifie que pour toute famille filtrante croissante $(U_\alpha)$
dans $\mathfrak O(X)$ ayant $X$ pour plus grand élément (resp. pour toute
famille filtrante croissante $(U_\alpha)$ dans $\mathfrak O(X)$), il existe
$\gamma$ tel que $U_\alpha=U_\gamma$ pour $\alpha\geq\gamma$. Comme
$U\mapsto f^{-1}(U)$ est une bijection de $\mathfrak O(Y)$ sur
$\mathfrak O(X)$, notre assertion résulte aussitôt de la remarque précédente.

Les ouverts quasi-compacts de $X$ sont donc les ensembles de la forme
$f^{-1}(U)$ où $U$ est ouvert quasi-compact dans $Y$, en vertu de
(10.2.1, \emph{a)}) et de ce qui précède. Pour que $X'$ soit rétrocompact dans
$X$, il faut et il suffit alors que pour tout ouvert quasi-compact $U$ dans
$Y$, $f^{-1}(U)\cap X'=f^{-1}(U\cap Y')$ soit quasi-compact
($\mathbf 0_{\mathrm{III}}$, 9.1.1)~; la première partie de la démonstration
montre que cela équivaut à dire que $U\cap Y'$ est quasi-compact pour tout
ouvert quasi-compact $U$, c'est-à-dire que $Y'$ est rétrocompact dans $Y$.

\begin{proposition}[10.2.7]
\label{IV.10.2.7-fr}
Soit $f:X\to Y$ un quasi-homéomorphisme. Alors l'application
$Z\mapsto f^{-1}(Z)$ de $\mathfrak P(Y)$ dans $\mathfrak P(X)$ définit par
restriction les bijections suivantes (cf. (10.1.1) pour les notations)~:
\[
\begin{gathered}
\mathfrak O(Y)\xrightarrow{\sim}\mathfrak O(X)\\
\mathfrak F(Y)\xrightarrow{\sim}\mathfrak F(X)\\
\mathfrak{Lf}(Y)\xrightarrow{\sim}\mathfrak{Lf}(X)\\
\mathfrak{Qc}(Y)\xrightarrow{\sim}\mathfrak{Qc}(X)\\
\mathfrak{Lqc}(Y)\xrightarrow{\sim}\mathfrak{Lqc}(X)\\
\mathfrak C(Y)\xrightarrow{\sim}\mathfrak C(X)\\
\mathfrak{Lc}(Y)\xrightarrow{\sim}\mathfrak{Lc}(X).
\end{gathered}
\]
\end{proposition}

\oldpage[IV-3]{100}
Pour les deux premières, ce n'est autre que la définition (10.2.2)~; comme la
topologie de $X$ est l'image réciproque par $f$ de celle de $f(X)$ les cinq
dernières applications, où l'on remplace $Y$ par $f(X)$, sont bijectives. On
peut donc (par (10.2.1, \emph{b)})) se borner au cas où $X$ est un sous-espace
très dense de $Y$, et le fait que les applications
$\mathfrak{Lf}(Y)\to\mathfrak{Lf}(X)$,
$\mathfrak{Qc}(Y)\to\mathfrak{Qc}(X)$,
$\mathfrak{Lqc}(Y)\to\mathfrak{Lqc}(X)$ sont bijectives a déjà été prouvé
(10.1.2). Toute partie localement constructible étant localement
quasi-constructible, les applications $\mathfrak C(Y)\to\mathfrak C(X)$ et
$\mathfrak{Lc}(Y)\to\mathfrak{Lc}(X)$ sont injectives~; en outre, pour tout
ouvert $U\subset Y$, la restriction $f^{-1}(U)\to U$ de $f$ est un
quasi-homéomorphisme (10.2.4), donc si l'on montre que
$\mathfrak C(Y)\to\mathfrak C(X)$ est surjective, il en sera de même de
$\mathfrak{Lc}(Y)\to\mathfrak{Lc}(X)$. Mais en vertu de (10.2.6), toute partie
ouverte rétrocompacte $Z$ dans $X$ est de la forme $f^{-1}(T)$, où $T$ est
ouvert rétrocompact dans $Y$~; cela prouve évidemment la surjectivité de
$\mathfrak C(Y)\to\mathfrak C(X)$.

\begin{remarks}[10.2.8]
\label{IV.10.2.8-fr}
\upshape (i) On a vu dans la démonstration de (10.2.1) que si $f$ est un
quasi-homéomorphisme, l'application canonique
\begin{equation}
\label{IV.10.2.8.1-fr}
\tag{10.2.8.1}
\Gamma(Y,\mathcal F)\to\Gamma(X,f^*(\mathcal F))
\end{equation}
est un isomorphisme de groupes abéliens fonctoriel en $\mathcal F$ dans la
catégorie des faisceaux de groupes abéliens sur $Y$. Comme $f^*$ est exact
dans cette catégorie ($\mathbf 0_{\mathrm I}$, 3.7.2), cela implique que les
homomorphismes canoniques (T, 3.2.2)
\[
\mathrm H^i(Y,\mathcal F)\xrightarrow{\sim}
\mathrm H^i(X,f^*(\mathcal F))
\]
sont \emph{bijectifs} pour tout $i$.

(ii) Si $f:X\to Y$ est une application continue et $\mathcal F$,
$\mathcal G$ deux faisceaux de groupes abéliens sur $Y$, on a
\[
f^*(\mathcal F\otimes_{\mathbf Z_Y}\mathcal G)
=f^*(\mathcal F)\otimes_{\mathbf Z_X}f^*(\mathcal G)
\]
à un isomorphisme canonique près ($\mathbf 0_{\mathrm I}$, 4.3.3). On en
conclut que si $f$ est un quasi-homéomorphisme, $f^*$ est encore une
\emph{équivalence} de la catégorie des faisceaux d'anneaux sur $Y$, et de la
catégorie des faisceaux d'anneaux sur $X$~; la donnée d'une structure d'espace
annelé sur $Y$ est donc équivalente à celle d'une structure d'espace annelé
sur $X$.

Étant donnés deux espaces annelés $(X,\mathcal O_X)$, $(Y,\mathcal O_Y)$, nous
dirons qu'un morphisme d'espaces annelés
$f=(\psi,\theta):(X,\mathcal O_X)\to(Y,\mathcal O_Y)$ est un
\emph{quasi-isomorphisme} si $\psi$ est un quasi-homéomorphisme de $X$ dans
$Y$ et $\theta^\sharp:\psi^*(\mathcal O_Y)\to\mathcal O_X$ un isomorphisme de
faisceaux d'anneaux~; lorsqu'il en est ainsi, l'espace annelé
$(X,\mathcal O_X)$ est entièrement déterminé, à isomorphisme près, par
$(Y,\mathcal O_Y)$, l'espace $X$ et le quasi-homéomorphisme $\psi$ (que l'on
peut prendre arbitraire). Lorsque $f$ est un quasi-isomorphisme d'espaces
annelés, le foncteur
\[
\mathcal F\mapsto f^*(\mathcal F)
\]
est une \emph{équivalence} de la catégorie des $\mathcal O_Y$-Modules et de
celle des $\mathcal O_X$-Modules, puisque $f^*(\mathcal F)$ s'identifie
canoniquement ici à $\psi^*(\mathcal F)$. On en conclut par exemple des
isomorphismes de bi-$\partial$-foncteurs
\[
\operatorname{Ext}^i_{\mathcal O_Y}(\mathcal F,\mathcal G)
\xrightarrow{\sim}
\operatorname{Ext}^i_{\mathcal O_X}(f^*(\mathcal F),f^*(\mathcal G)).
\]
De façon générale, on peut dire que les constructions usuelles de la théorie
des faisceaux et de l'algèbre homologique, faites sur l'espace annelé $Y$ ou
l'espace annelé $X$, sont équivalentes.
\end{remarks}

\oldpage[IV-3]{101}
\subsection{Espaces de Jacobson}
\label{subsection:IV.10.3-fr}

\begin{definition}[10.3.1]
\label{IV.10.3.1-fr}
On dit qu'un espace topologique $X$ est un \emph{espace de Jacobson} si
l'ensemble $X_0$ des points fermés de $X$ est très dense dans $X$ (autrement
dit, si l'injection canonique $X_0\to X$ est un
quasi-homéomorphisme).
\end{definition}

Cela signifie donc (10.1.2) que toute partie localement fermée (ou seulement
localement quasi-constructible) $Z\neq\varnothing$ de $X$ contient un point
fermé de $X$, ou encore que \emph{toute partie fermée de $X$ est l'adhérence
de l'ensemble de ses points fermés}.

\begin{proposition}[10.3.2]
\label{IV.10.3.2-fr}
Soient $X$ un espace de Jacobson, $Z$ une partie localement
quasi-constructible de $X$~; alors le sous-espace $Z$ de $X$ est un espace de
Jacobson, et pour qu'un point de $Z$ soit fermé dans $Z$, il faut et il suffit
qu'il soit fermé dans $X$.
\end{proposition}

Si $X_0$ est l'ensemble des points fermés de $X$, $Z\cap X_0$ est très dense
dans $Z$ en vertu de (10.2.4) appliqué à l'injection $i:X_0\to X$~; comme
l'ensemble $Z_0$ des points fermés de $Z$ contient évidemment $Z\cap X_0$,
il est \emph{a fortiori} très dense dans $Z$, donc $Z$ est un espace de
Jacobson. Montrons maintenant qu'en fait on a $Z_0=Z\cap X_0$~; soit $x\in Z$
un point fermé dans $Z$~; soit $\overline{\{x\}}$ son adhérence dans $X$~;
$\overline{\{x\}}\cap Z=\{x\}$ est donc une partie localement
quasi-constructible de $X$, et comme son intersection avec $X_0$ est non vide
(10.1.2), on a $x\in X_0$.

\begin{proposition}[10.3.3]
\label{IV.10.3.3-fr}
Soient $X$ un espace topologique, $(U_\alpha)$ un recouvrement ouvert de $X$.
Pour que $X$ soit un espace de Jacobson, il faut et il suffit que chacun des
sous-espaces $U_\alpha$ le soit.
\end{proposition}

La condition est nécessaire en vertu de (10.3.2). Inversement, montrons
d'abord que l'hypothèse que les $U_\alpha$ sont des espaces de Jacobson
entraîne que pour qu'un point $x\in U_\alpha$ soit fermé dans $X$, il suffit
qu'il soit fermé dans $U_\alpha$. Il suffit en effet de voir que cette
condition entraîne que $x$ est aussi fermé dans chacun des $U_\beta$ qui le
contiennent~; mais $U_\alpha\cap U_\beta$ est ouvert dans $U_\alpha$, donc $x$
est fermé dans $U_\alpha\cap U_\beta$, et en vertu de (10.3.2), $x$ est aussi
fermé dans $U_\beta$, ce qui achève la démonstration.

\subsection{Préschémas de Jacobson et anneaux de Jacobson}
\EGAReaderTarget{ega4.c04.subsec.10-4}{subsection:IV.10.4-fr}

\begin{definition}[10.4.1]
\label{IV.10.4.1-fr}
On dit qu'un préschéma $X$ est un \emph{préschéma de Jacobson} si l'espace
topologique $X$ sous-jacent est un espace de Jacobson. On dit qu'un anneau
$A$ est un \emph{anneau de Jacobson} si $\operatorname{Spec}(A)$ est un
préschéma de Jacobson.
\end{definition}

Toute partie fermée de $X=\operatorname{Spec}(A)$ est de la forme
$Z=V(\mathfrak a)$, où $\mathfrak a$ est un idéal égal à sa racine, et
l'ensemble $X_0$ des points fermés de $X$ est l'ensemble des idéaux maximaux
de $A$~; dire que $Z\cap X_0$ est dense dans $Z$ signifie donc que
$\mathfrak a$ est \emph{intersection d'idéaux maximaux} (\textbf{I}, 1.1.4)~;
comme $\mathfrak a$ est intersection d'idéaux premiers, il revient au même de
dire que tout idéal premier de $A$ est intersection d'idéaux maximaux~; en
vertu de (10.3.1) et (10.1.2), la définition usuelle des anneaux de Jacobson
(Bourbaki, \emph{Alg. comm.}, chap.~V, \S~3, n\textsuperscript{o}~4, déf.~1)
coïncide donc avec la définition (10.4.1).

\oldpage[IV-3]{102}
\begin{proposition}[10.4.2]
\label{IV.10.4.2-fr}
Soient $X$ un préschéma, $(U_\alpha)$ un recouvrement de $X$ formé d'ouverts
affines. Pour que $X$ soit un préschéma de Jacobson, il faut et il suffit que
les anneaux $\Gamma(U_\alpha,\mathcal O_X)$ des $U_\alpha$ soient des anneaux
de Jacobson.
\end{proposition}

Cela résulte de (10.3.3) et de la définition (10.4.1).

\begin{env}[10.4.3]
\label{IV.10.4.3-fr}
Un préschéma discret est un préschéma de Jacobson~; un anneau artinien est donc
un anneau de Jacobson. Tout anneau principal ayant une infinité d'idéaux
maximaux (par exemple $\mathbf Z$) est un anneau de Jacobson~; un anneau local
noethérien $A$ est un anneau de Jacobson si et seulement si son idéal maximal
est son seul idéal premier, c'est-à-dire si $A$ est artinien. Tout
sous-préschéma d'un préschéma de Jacobson est un préschéma de Jacobson en vertu
de (10.3.2).
\end{env}

\begin{proposition}[10.4.4]
\label{IV.10.4.4-fr}
Soit $B$ un anneau intègre. Les conditions suivantes sont équivalentes~:
\begin{enumerate}[label=\emph{\alph*)}]
  \item Il existe $f\neq0$ dans $B$ tel que $B_f$ soit un corps.
  \item Le corps des fractions de $B$ est une $B$-algèbre de type fini.
  \item Il existe un corps $K$ contenant $B$ et qui est une $B$-algèbre de
  type fini.
  \item Le point générique de $\operatorname{Spec}(B)$ est isolé dans
  $\operatorname{Spec}(B)$.
\end{enumerate}
\end{proposition}

Il est clair que \emph{d)} est équivalent à \emph{a)}, puisque \emph{d)}
signifie qu'il existe $f\neq0$ dans $B$ tel que $D(f)$ soit réduit au point
générique de $\operatorname{Spec}(B)$. Il est trivial que \emph{a)} entraîne
\emph{b)} et que \emph{b)} entraîne \emph{c)}. Enfin, \emph{c)} entraîne
\emph{a)}, en vertu de Bourbaki, \emph{Alg. comm.}, chap.~V, \S~3,
n\textsuperscript{o}~1, cor.~2 du th.~1.

\begin{proposition}[10.4.5]
\label{IV.10.4.5-fr}
Soit $A$ un anneau. Les conditions suivantes sont équivalentes~:
\begin{enumerate}[label=\emph{\alph*)}]
  \item $A$ est un anneau de Jacobson.
  \item Pour tout idéal premier non maximal $\mathfrak p$ de $A$ et tout
  $f\neq0$ dans $B=A/\mathfrak p$, $B_f$ n'est pas un corps.
  \item[\emph{b$'$)}] Toute $A$-algèbre de type fini $K$ qui est un corps est
  une $A$-algèbre finie.
\end{enumerate}
\end{proposition}

On sait que \emph{a)} entraîne \emph{b$'$)} (Bourbaki,
\emph{Alg. comm.}, chap.~V, \S~3, n\textsuperscript{o}~4, cor.~3 du th.~3).
En outre, le noyau de l'homomorphisme $A\to K$ est alors un idéal
\emph{maximal} $\mathfrak m$ de $A$, et $K$ est une extension \emph{finie} de
$A/\mathfrak m$ (\emph{loc. cit.}). Il est trivial que \emph{b$'$)} entraîne
\emph{b)}, puisque $A/\mathfrak p=B$ n'est pas un corps, toute $B$-algèbre de
type fini est une $A$-algèbre de type fini et $B_f$ est une $B$-algèbre de
type fini. Reste à voir que \emph{b)} implique \emph{a)}. Nous utiliserons le
lemme suivant~:

\begin{lemma}[10.4.5.1]
\label{IV.10.4.5.1-fr}
Soit $X$ un espace topologique ayant la propriété suivante~: pour toute partie
localement fermée $Z\neq\varnothing$ de $X$, il existe une partie $Z'$ de $Z$,
localement fermée dans $X$ (ou dans $Z$, ce qui revient au même), et un point
$x\in Z'$, fermé dans $Z'$. Alors, pour que $X$ soit un espace de Jacobson, il
faut et il suffit qu'aucun point non fermé $x$ de $X$ ne soit isolé dans
$\overline{\{x\}}$.
\end{lemma}

Si $X$ est un espace de Jacobson, un point non fermé $x$ de $X$ ne peut être
isolé dans $\overline{\{x\}}$, car cela signifierait qu'il existe un ouvert
$U$ de $X$ tel que $U\cap\overline{\{x\}}=\{x\}$~; or
$U\cap\overline{\{x\}}$ est localement fermé dans $X$ et ne contiendrait aucun
point fermé dans $X$, ce qui est contraire à l'hypothèse que $X$ est un espace
de Jacobson. Inversement, supposons vérifiée la condition de l'énoncé~; alors
l'ensemble des points fermés de $X$ est identique à l'ensemble $X_0$ des
$x\in X$ qui sont isolés dans $\overline{\{x\}}$~; mais il résulte de
(5.1.10.1) que

\oldpage[IV-3]{103}
cet ensemble est \emph{très dense} dans $X$, donc $X$ est un espace de
Jacobson par définition.

Rappelons (5.1.10) que l'hypothèse faite sur $X$ dans (10.4.5.1) est toujours
vérifiée lorsque $X$ est l'espace sous-jacent à un \emph{préschéma}.

Revenant alors à la démonstration de (10.4.5), la condition \emph{b)} entraîne
que pour tout point $x$ non fermé de $\operatorname{Spec}(A)$, le point
générique $x$ de
$\operatorname{Spec}(A/\mathfrak j_x)=V(\mathfrak j_x)=\overline{\{x\}}$
n'est pas isolé dans $\overline{\{x\}}$, donc \emph{b)} entraîne \emph{a)} en
vertu du lemme (10.4.5.1) et de (10.4.4).

\begin{corollary}[10.4.6]
\label{IV.10.4.6-fr}
Toute algèbre de type fini $B$ sur un anneau de Jacobson $A$ est un anneau de
Jacobson, et l'image réciproque dans $A$ de tout idéal maximal de $B$ est un
idéal maximal de $A$. En particulier, toute algèbre de type fini sur un corps
ou sur $\mathbf Z$ est un anneau de Jacobson.
\end{corollary}

Une $B$-algèbre de type fini $K$ qui est un corps est aussi une $A$-algèbre de
type fini, donc un $A$-module de type fini et \emph{a fortiori} un $B$-module
de type fini, d'où la première assertion~; la seconde a été démontrée au cours
de la démonstration de (10.4.5), appliqué à $K=B/\mathfrak m$.

\begin{corollary}[10.4.7]
\label{IV.10.4.7-fr}
Si $X$ est un préschéma de Jacobson, $f:Y\to X$ un morphisme localement de type
fini, alors $Y$ est un préschéma de Jacobson et l'image par $f$ de tout point
fermé dans $Y$ est un point fermé dans $X$.
\end{corollary}

La question étant locale sur $X$ et sur $Y$, on est ramené au cas où $X$ et
$Y$ sont affines, et le corollaire résulte alors de (10.4.6).

\begin{corollary}[10.4.8]
\label{IV.10.4.8-fr}
Si $k$ est un corps algébriquement clos, $X$ un $k$-préschéma localement de
type fini, l'ensemble des points de $X$ rationnels sur $k$ est très dense dans
$X$.
\end{corollary}

En effet, $X$ est un préschéma de Jacobson (10.4.7) et les points fermés de
$X$ sont exactement les points rationnels sur $k$ (\textbf{I}, 6.4.2).

\begin{env}[10.4.9]
\label{IV.10.4.9-fr}
Le fait que les préschémas localement de type fini sur un corps ou sur
$\mathbf Z$ sont des préschémas de Jacobson est particulièrement important,
en raison de la possibilité de se ramener à ce cas dans de nombreuses
questions de Géométrie algébrique (8.1.2, \emph{c)}). Nous allons en donner
deux exemples~:
\end{env}

\begin{env}[10.4.10]
\label{IV.10.4.10-fr}
\textbf{Applications (10.4.10)~: I~: Démonstration de (6.15.9).} --- Soient
$k$ un corps séparablement clos, $X$ un $k$-préschéma localement de type fini
sur $k$ et \emph{unibranche}. On sait que la fermeture intégrale d'une
$k$-algèbre intègre de type fini $A$ dans une extension finie de son corps des
fractions est une $A$-algèbre finie (Bourbaki, \emph{Alg. comm.}, chap.~V,
\S~3, n\textsuperscript{o}~2, th.~2), donc toute $k$-algèbre de type fini est
un anneau universellement japonais~; on en conclut que l'ensemble des points
$x\in X$ où $X$ est géométriquement unibranche est localement constructible
(9.7.10). Mais l'hypothèse et le lemme (6.15.8) entraînent que cet ensemble
contient tous les points \emph{fermés} de $X$. La conclusion résulte donc de
(10.4.6), (10.3.1) et de la bijectivité de l'application canonique
$\mathfrak{Lc}(X)\to\mathfrak{Lc}(X_0)$ (où $X_0$ est l'ensemble des points
fermés de $X$) (10.2.7).
\end{env}

\begin{env}[10.4.11]
\EGAReaderTarget{ega4.c04.statement.10-4-11}{IV.10.4.11-fr}
\textbf{Applications~: II~: Proposition (10.4.11).} --- Soient $S$ un
préschéma, $X$ un $S$-préschéma de type fini. Tout $S$-endomorphisme de $X$
qui est radiciel est surjectif (donc bijectif).
\end{env}

\oldpage[IV-3]{104}
Soient $f:X\to S$ le morphisme structural, $g:X\to X$ le $S$-endomorphisme
considéré et, pour tout $s\in S$, soit $g_s$ le morphisme déduit de $g$ par le
changement de base $\operatorname{Spec}(k(s))\to S$, qui est un
$k(s)$-endomorphisme de la fibre
$f^{-1}(s)=X\times_S\operatorname{Spec}(k(s))$.

Pour prouver que $g$ est surjectif, il suffit de prouver que $g_s$ est
surjectif pour tout $s\in S$, donc on peut (en vertu de (\textbf{I}, 3.5.7))
se borner au cas où $S=\operatorname{Spec}(k)$ est le spectre d'un corps $k$,
auquel cas $f$ est un morphisme de présentation finie, puisque $S$ est
noethérien. Appliquant (8.9.1) et (8.10.5, (vii)), on est ramené au cas où
$S=\operatorname{Spec}(A)$, où $A$ est une sous-$\mathbf Z$-algèbre de type
fini de $k$. Or $X$ est alors un préschéma de Jacobson (10.4.7), et $g(X)$ est
constructible dans $X$ (\textbf{I}, 1.8.5), donc, pour prouver que $g(X)=X$,
il suffit de montrer que $g(X)$ contient tous les points \emph{fermés} de $X$
(10.3.1).

\begin{lemma}[10.4.11.1]
\label{IV.10.4.11.1-fr}
Soit $Y$ un $\mathbf Z$-préschéma de type fini.
\begin{enumerate}[label=(\roman*)]
  \item Pour qu'un point $y\in Y$ soit fermé dans $Y$, il faut et il suffit
  que $k(y)$ soit un corps fini.
  \item Pour tout nombre premier $p$ et tout entier $d\geq1$, l'ensemble des
  points $y\in Y$ tels que $k(y)$ soit une extension de $\mathbf F_p$ dont le
  degré divise $d$ est fini.
\end{enumerate}
\end{lemma}

L'assertion (i) résulte de ce que l'image d'un point fermé $y\in Y$ dans
$\operatorname{Spec}(\mathbf Z)$ est un point fermé (10.4.7), autrement dit
un nombre premier, et de (\textbf{I}, 6.4.2). D'autre part, comme $Y$ est
réunion finie d'ouverts affines de type fini sur $\mathbf Z$, on peut se
borner pour prouver (ii) au cas où $Y=\operatorname{Spec}(C)$, où $C$ est une
$\mathbf Z$-algèbre de type fini. Or, les points $y\in Y$ tels que le degré de
$k(y)$ sur $\mathbf F_p$ divise $d$ correspondent biunivoquement aux
homomorphismes $C\to\mathbf F_{p^d}$~; mais si $(t_i)$ est un système de
générateurs de la $\mathbf Z$-algèbre $C$, tout homomorphisme de $C$ est
déterminé par ses valeurs aux éléments $t_i$, et par suite il n'y a qu'un
nombre fini d'homomorphismes de $C$ dans un corps fini.

Cela étant, $X$ est un $\mathbf Z$-préschéma de type fini~; soit $T_{p,d}$
l'ensemble des points fermés $z\in X$ tels que $k(z)$ soit une extension de
$\mathbf F_p$ de degré divisant $d$~; il résulte de (10.4.11.1) que l'ensemble
$T_{p,d}$ est fini et que l'ensemble des points fermés de $X$ est réunion des
$T_{p,d}$. En outre, si $z\in T_{p,d}$ et si $h$ est un endomorphisme
quelconque de $X$, $k(h(z))$ est isomorphe à un sous-corps de $k(z)$, donc
$h(z)\in T_{p,d}$, autrement dit $T_{p,d}$ est \emph{stable} par tout
endomorphisme de $X$. Comme $g$ est par hypothèse injectif, sa restriction à
$T_{p,d}$ est \emph{bijective} puisque $T_{p,d}$ est fini, ce qui achève de
démontrer la proposition.

Nous verrons plus loin (17.9.7) que lorsqu'on suppose de plus, d'une part que
$X$ est un $S$-préschéma de présentation finie, d'autre part que $g$ est un
monomorphisme, alors on peut affirmer que $g$ est un \emph{automorphisme} de
$X$.

\subsection{Préschémas de Jacobson noethériens}
\label{subsection:IV.10.5-fr}

\begin{proposition}[10.5.1]
\label{IV.10.5.1-fr}
Soit $B$ un anneau intègre noethérien. Les conditions équivalentes \emph{a)} à
\emph{d)} de (10.4.4) sont alors aussi équivalentes aux suivantes~:
\begin{enumerate}[label=\emph{\alph*)}]
  \setcounter{enumi}{4}
  \item $\operatorname{Spec}(B)$ est fini.
  \item $B$ est un anneau semi-local de dimension $\leq1$.
\end{enumerate}
\end{proposition}

Il résulte en effet du théorème d'Artin--Tate ($\mathbf 0$, 16.3.3) que les
conditions \emph{a)} et \emph{f)} sont équivalentes. La condition \emph{f)}
implique que $\operatorname{Spec}(B)$ est réunion de l'ensemble

\oldpage[IV-3]{105}
fini de ses points fermés et de son point générique, donc \emph{f)} implique
\emph{e)} sans supposer $A$ noethérien. Enfin \emph{e)} implique \emph{d)} sans
supposer $A$ noethérien, car le point générique $x$ de $X$ est le complémentaire
de la réunion des adhérences $\overline{\{y\}}$, où $y$ parcourt l'ensemble des
points $y\neq x$, et comme ces points sont en nombre fini, $\{x\}$ est
complémentaire d'un ensemble fermé dans $X$.

\begin{corollary}[10.5.2]
\label{IV.10.5.2-fr}
Soit $A$ un anneau noethérien. Pour que $A$ soit un anneau de Jacobson, il faut
et il suffit qu'il n'existe aucun idéal premier $\mathfrak p$ de $A$ tel que
$A/\mathfrak p$ soit un anneau semi-local de dimension $1$.
\end{corollary}

Cela résulte aussitôt de (10.5.1) et de la condition \emph{b)} de (10.4.5),
les idéaux premiers $\mathfrak p$ de $A$ tels que $A/\mathfrak p$ soit
semi-local de dimension $0$ étant les idéaux maximaux de $A$.

\begin{corollary}[10.5.3]
\label{IV.10.5.3-fr}
Soit $X$ un préschéma noethérien irréductible. Les conditions suivantes sont
équivalentes~:
\begin{enumerate}[label=\emph{\alph*)}]
  \item Le point générique de $X$ est isolé.
  \item $X$ est fini.
  \item $X$ est de dimension $\leq1$ et l'ensemble de ses points fermés est
  fini.
\end{enumerate}
\end{corollary}

Il existe par hypothèse un nombre fini d'ouverts affines irréductibles $U_i$
($1\leq i\leq n$) recouvrant $X$, et dont chacun contient donc le point
générique de $X$~; il suffit de prouver l'équivalence de \emph{a)}, \emph{b)}
et \emph{c)} pour chacun des $(U_i)_{\mathrm{red}}$ (compte tenu de
($\mathbf 0$, 14.1.7)). Mais cette équivalence résulte alors de (10.5.1).

\begin{remark}[10.5.4]
\label{IV.10.5.4-fr}
Un préschéma noethérien $X$ vérifiant les conditions équivalentes de (10.5.3)
n'est pas nécessairement un schéma affine~; en fait il peut même être
\emph{non séparé}. On en a un exemple en remplaçant, dans l'exemple
(\textbf{I}, 5.5.11) de «~droite affine avec point dédoublé~», $X_1$ et $X_2$
par le spectre de l'anneau de valuation discrète $(K[s])_{(s)}$, $U_{12}$ et
$U_{21}$ par l'ouvert réduit au point générique dans $X_1$ et $X_2$
respectivement~; le préschéma non séparé $X$ que l'on obtient a exactement
$3$ points.
\end{remark}

\begin{proposition}[10.5.5]
\label{IV.10.5.5-fr}
Soit $X$ un préschéma noethérien.
\begin{enumerate}[label=(\roman*)]
  \item L'ensemble $X_0$ des points $x\in X$ tels que $\overline{\{x\}}$ soit
  fini est très dense dans $X$.
  \item Pour que $X$ soit un préschéma de Jacobson, il faut et il suffit qu'il
  n'existe aucun sous-préschéma de $X$ isomorphe au spectre d'un anneau
  semi-local intègre de dimension $1$.
\end{enumerate}
\end{proposition}

La condition que $\overline{\{x\}}$ soit fini équivaut en effet ici (10.5.3)
au fait que $x$ soit isolé dans $\overline{\{x\}}$, $\overline{\{x\}}$ étant
espace sous-jacent d'un sous-préschéma (noethérien) de $X$~; l'assertion (i)
résulte donc de (5.1.10.1). De même, compte tenu de (10.4.5.1), pour démontrer
l'assertion (ii), remarquons que pour qu'un point non fermé $x$ de $X$
appartienne à $X_0$, il faut et il suffit que le sous-préschéma fermé intègre
$Y$ de $X$ ayant $\overline{\{x\}}$ pour espace sous-jacent soit de dimension
$1$ et fini, donc réunion finie de sous-préschémas ouverts (dans $Y$) affines
$U_i$ qui sont des spectres d'anneaux semi-locaux intègres de dimension $1$.
Inversement, s'il y a un sous-préschéma $Z$ de $X$ qui soit spectre d'un anneau
semi-local intègre de dimension $1$, $Z$ n'est pas un préschéma de Jacobson
(10.5.2), donc il en est de même de $X$ (10.3.2).

\oldpage[IV-3]{106}
\begin{remark}[10.5.6]
\label{IV.10.5.6-fr}
L'assertion (ii) de (10.5.5) est encore valable lorsque $X$ est
\emph{localement noethérien}~: en effet, si $(V_\alpha)$ est un recouvrement de
$X$ formé d'ouverts affines (noethériens), tout sous-préschéma d'un $V_\alpha$
est un sous-préschéma de $X$~; inversement, si un sous-préschéma $Z$ de $X$ est
isomorphe au spectre d'un anneau semi-local intègre de dimension $1$, il y a
un $\alpha$ tel que $V_\alpha\cap Z$ contienne un ouvert affine $U$ de $Z$ non
réduit au point générique de $Z$, et qui est donc aussi spectre d'un anneau
semi-local intègre de dimension $1$. On conclut à l'aide de (10.3.3).
\end{remark}

\begin{proposition}[10.5.7]
\label{IV.10.5.7-fr}
Soient $X$ un préschéma localement noethérien, $Y$ une partie fermée de $X$
telle que toute partie fermée non vide de $X$ rencontre $Y$. Alors le
préschéma induit sur l'ouvert $X-Y$ est un préschéma de Jacobson.
\end{proposition}

Appliquons le critère (10.5.5, (ii)), et supposons qu'il y ait un
sous-préschéma $Z$ de $X-Y$ qui soit spectre d'un anneau semi-local intègre de
dimension $1$, le point générique $z$ de $Z$ étant donc isolé dans $Z$ (ou
dans l'adhérence $\overline Z$ de $Z$ dans $X$), et $Z$ étant distinct de
$\{z\}$. Soit $y\neq z$ un point de $Z$~; comme il n'appartient pas à $Y$, il
n'est pas fermé dans $X$, et son adhérence $\overline{\{y\}}$ dans $X$
rencontre $Y$ en un point $x\neq y$ qui est donc spécialisation de $y$.
L'existence de la chaîne
$\{x\}\subset\overline{\{y\}}\subset\overline{\{z\}}$ montre alors que la
dimension de $\overline{\{z\}}$ serait $\geq2$, et il en serait de même de la
dimension de $U\cap\overline{\{z\}}$, où $U$ est un voisinage ouvert affine
(donc noethérien) de $x$ dans $X$. Or, cela contredit le fait que le point
générique $z$ de $U\cap\overline{\{z\}}$ est isolé dans
$U\cap\overline{\{z\}}$ (10.5.3).

\begin{corollary}[10.5.8]
\label{IV.10.5.8-fr}
Soit $A$ un anneau noethérien~; pour tout élément $f$ du radical $\mathfrak R$
de $A$, l'anneau $A_f$ est un anneau de Jacobson, et l'ouvert
$\operatorname{Spec}(A)-V(\mathfrak R)$ est un schéma de Jacobson.
\end{corollary}

Si $X=\operatorname{Spec}(A)$, $Y=V(\mathfrak J)$, où $\mathfrak J$ est un
idéal de $A$, dire que toute partie fermée non vide de $X$ rencontre $Y$
équivaut ($\mathbf 0_{\mathrm I}$, 2.1.3) à dire que $Y$ contient tous les
\emph{points fermés} de $X$, ou encore que $\mathfrak J$ est \emph{contenu
dans le radical} $\mathfrak R$ de $A$. Si $f\in\mathfrak R$, l'ouvert $D(f)$
ne rencontre donc pas $V(\mathfrak R)$, et est un espace de Jacobson en vertu
de (10.5.7).

\begin{corollary}[10.5.9]
\label{IV.10.5.9-fr}
Soient $A$ un anneau local noethérien, $\mathfrak m$ son idéal maximal,
$X=\operatorname{Spec}(A)$, $Y=X-\{\mathfrak m\}$~; alors $Y$ est un schéma de
Jacobson, dont les points fermés sont les idéaux premiers
$\mathfrak p\in\operatorname{Spec}(A)$ tels que $\dim(A/\mathfrak p)=1$.
\end{corollary}

La première assertion est un cas particulier de (10.5.8)~; d'autre part les
points fermés de $Y$ sont les idéaux premiers $\mathfrak p$ de $A$ qui sont
des éléments maximaux dans l'ensemble des idéaux premiers $\neq\mathfrak m$,
ce qui, par définition de la dimension, signifie que
$\dim(A/\mathfrak p)=1$.

\begin{proposition}[10.5.10]
\label{IV.10.5.10-fr}
Soit $A$ un anneau local noethérien réduit et complet et qui n'est pas un
corps. Alors les intersections finies des noyaux des homomorphismes locaux de
$A$ dans des anneaux de valuation discrète $V$, qui font de $V$ une
$A$-algèbre finie, forment une base de filtre tendant vers $0$ pour la
topologie adique de $A$.
\end{proposition}

Il suffit (Bourbaki, \emph{Alg. comm.}, chap.~III, \S~2,
n\textsuperscript{o}~7, prop.~8) de prouver que l'intersection des noyaux
envisagés dans l'énoncé est réduite à $0$. Supposons d'abord que $A$ soit
intègre et de dimension $1$~; en vertu du théorème de Nagata ($\mathbf 0$,
23.1.5 et 23.1.6), la clôture intégrale $A'$ de $A$ est un anneau local
complet intègre et

\oldpage[IV-3]{107}
intégralement clos et de dimension $1$ et une $A$-algèbre finie~; c'est donc
un anneau de valuation discrète, et la proposition en résulte aussitôt dans
ce cas.

Passons au cas général~; soit $x$ le point fermé de
$X=\operatorname{Spec}(A)$, et posons $Y=X-\{x\}$~; on sait (10.5.9) que $Y$
est un préschéma de Jacobson. Montrons que cela entraîne que l'intersection
des idéaux premiers $\mathfrak p$ de $A$ tels que
$\dim(A/\mathfrak p)=1$ est \emph{réduite à $0$}~; la proposition en résultera
puisque, pour chacun de ces idéaux $\mathfrak p$, l'intersection des noyaux des
homomorphismes locaux de $A/\mathfrak p$ dans des anneaux de valuation
discrète $V$, faisant de $V$ une $(A/\mathfrak p)$-algèbre finie, est réduite
à $0$. Mais dire que l'intersection de ces idéaux premiers est réduite à $0$
signifie que l'ensemble de ces idéaux est dense dans $X$, ou encore dans $Y$
(puisque $Y$ est dense dans $X$), et cela résulte aussitôt de (10.5.9).

\subsection{Dimension dans les préschémas de Jacobson}
\label{subsection:IV.10.6-fr}

Les résultats de ce numéro précisent dans certains cas et généralisent des résultats
du \S~5.

\begin{proposition}[10.6.1]
\label{IV.10.6.1-fr}
Soit $S$ un préschéma localement noethérien, vérifiant en outre les conditions
suivantes~: $1^\circ$ $S$ est un préschéma de Jacobson~; $2^\circ$ pour tout
$s\in S$, $\mathcal O_{S,s}$ est universellement caténaire (5.6.2)~; $3^\circ$ toute
composante irréductible $S'$ de $S$ est équicodimensionnelle (autrement dit, pour
tout point fermé $s$ de $S'$ et tout sous-préschéma de $S$ ayant $S'$ pour espace
sous-jacent, on a $\dim(S')=\dim(\mathcal O_{S',s})$). On a alors les propriétés
suivantes~:
\begin{enumerate}[label=(\roman*)]
  \item Pour tout morphisme localement de type fini $g:X\to S$, $X$ vérifie les
  conditions $1^\circ$, $2^\circ$ et $3^\circ$ précédentes. En particulier, si $X$
  est équidimensionnel (par exemple si $X$ est irréductible), $X$ est
  biéquidimensionnel (autrement dit, $X$ est caténaire et pour tout point fermé
  $x$ de $X$, on a $\dim(\mathcal O_{X,x})=\dim(X)$ ($\mathbf 0$, 14.3.3)).
  \item Soient $X$, $Y$ deux $S$-préschémas localement de type fini sur $S$,
  $f:X\to Y$ un $S$-morphisme~; on suppose $X$ irréductible et $f$ dominant. Si
  $\xi$ (resp. $\eta$) est le point générique de $X$ (resp. $Y$) et
  $e=\dim(f^{-1}(\eta))=\deg.\operatorname{tr}_{k(\eta)}k(\xi)$, on a
  \begin{equation}
  \label{IV.10.6.1.1-fr}
  \tag{10.6.1.1}
  \dim(X)=\dim(Y)+e.
  \end{equation}
  \item Soient $X$, $Y$ deux $S$-préschémas localement de type fini sur $S$,
  $f:X\to Y$ un $S$-morphisme, $n$ un entier $\geq0$ tel que l'on ait
  $\dim(f^{-1}(y))\geq n$ (resp. $\dim(f^{-1}(y))\leq n$) pour tout $y\in Y$.
  Alors on a
  \begin{equation}
  \label{IV.10.6.1.2-fr}
  \tag{10.6.1.2}
  \dim(X)\geq\dim(Y)+n
  \end{equation}
  (resp.
  \begin{equation}
  \label{IV.10.6.1.3-fr}
  \tag{10.6.1.3}
  \dim(X)\leq\dim(Y)+n).
  \end{equation}
\end{enumerate}
\end{proposition}

(i) La propriété $1^\circ$ pour $X$ résulte de (10.4.7). Pour tout $x\in X$,
$\mathcal O_{X,x}$ est anneau local d'une $\mathcal O_{S,g(x)}$-algèbre de type fini en
un idéal premier, et l'homomorphisme
$\mathcal O_{S,g(x)}\to\mathcal O_{X,x}$

\oldpage[IV-3]{108}
est local~; donc (5.6.3, (iv)) $\mathcal O_{X,x}$ est universellement caténaire. Pour
démontrer que $X$ vérifie la condition $3^\circ$, considérons plusieurs cas~:

\emph{a)} $X$ est un sous-préschéma fermé irréductible de $S$~; soient $S'$ une
composante irréductible de $S$ contenant $X$, $\xi$ le point générique de $X$,
$x$ un point fermé de $X$~; pour tout $s\in S'$, $\mathcal O_{S',s}$, quotient de
$\mathcal O_{S,s}$, est caténaire (5.6.1), donc les conditions $2^\circ$ et $3^\circ$
entraînent que $S'$ est biéquidimensionnel ($\mathbf 0$, 14.3.3). En vertu de
(5.1.2) et de ($\mathbf 0$, 14.3.3.2), on a donc
\[
  \dim(\mathcal O_{X,x})=\dim(\mathcal O_{S',x})-\dim(\mathcal O_{S',\xi})
  =\dim(S')-\dim(\mathcal O_{S',\xi})
\]
ce qui prouve que $\dim(\mathcal O_{X,x})$ ne dépend pas du point fermé $x$
considéré, d'où l'assertion dans ce cas (5.1.4).

\emph{b)} $X$ est irréductible et $g$ dominant. Alors, pour tout point fermé $x$
de $X$, $g(x)$ est fermé dans $S$ par (10.4.7)~; comme $\mathcal O_{S,g(x)}$ est
universellement caténaire, il résulte de (5.6.5.3) que l'on a
\[
  \dim(\mathcal O_{X,x})=\dim(\mathcal O_{S,g(x)})+e
\]
où $e=\dim(g^{-1}(\zeta))$, $\zeta$ étant le point générique de $S$. Comme
$\dim(S)=\dim(\mathcal O_{S,g(x)})$ en vertu de la condition $3^\circ$ pour $S$, on
a $\dim(\mathcal O_{X,x})=\dim(S)+e$ pour tout point fermé $x\in X$~; cela prouve
la condition $3^\circ$ pour $X$ (5.1.4), et en même temps la formule
\begin{equation}
\label{IV.10.6.1.4-fr}
\tag{10.6.1.4}
\dim(X)=\dim(S)+e.
\end{equation}

\emph{c)} Cas général. En considérant un sous-préschéma réduit $X'$ de $X$ ayant
pour espace sous-jacent une composante irréductible de $X$, on est ramené au cas
où $X$ est intègre~; utilisant (\textbf{I}, 5.2.2) et le cas \emph{a)} prouvé
ci-dessus, on peut ensuite remplacer $S$ par le sous-préschéma réduit ayant
$\overline{g(X)}$ pour espace sous-jacent~; on est alors ramené au cas \emph{b)},
et cela achève de démontrer (i).

(ii) Le morphisme $f$ étant localement de type fini (\textbf{I}, 1.3.4), on peut
appliquer les résultats de (i) en remplaçant $S$ par $Y$~; en outre, comme $X$
est irréductible et $f$ dominant, on peut aussi remplacer $S$ par $Y$ dans
(10.6.1.4), ce qui donne (10.6.1.1).

(iii) L'assertion relative au cas où $\dim(f^{-1}(y))\leq n$ pour tout $y\in Y$
a déjà été démontrée sous des hypothèses plus générales dans (5.6.7). Supposons
que $\dim(f^{-1}(y))\geq n$ pour tout $y\in Y$, et considérons un point générique
$\eta$ d'une composante irréductible $Y'$ de $Y$~; il existe au moins une
composante irréductible $Z$ de $f^{-1}(\eta)$ de dimension $\geq n$~; si $\xi$
est le point générique de $Z$, $\xi$ est aussi point générique d'une composante
irréductible $X'$ de $X$ telle que $f(X')$ soit dense dans $Y'$ ($\mathbf 0_{\mathrm I}$,
2.1.8). Considérons les sous-préschémas réduits de $X$, $Y$ ayant respectivement
pour espaces sous-jacents $X'$, $Y'$, et la restriction $X'\to Y'$ de $f$
(\textbf{I}, 5.2.2)~; il résulte alors de (ii) que
$\dim(X')\geq\dim(Y')+n$, et a fortiori $\dim(X)\geq\dim(Y')+n$~; ceci étant vrai
pour toute composante irréductible $Y'$ de $Y$, on en conclut l'inégalité
(10.6.1.2).

\begin{corollary}[10.6.2]
\label{IV.10.6.2-fr}
Supposons que $X$ vérifie les conditions $1^\circ$, $2^\circ$ et $3^\circ$ de
(10.6.1). Alors, pour tout ouvert $U$ dense dans $X$, on a
$\dim(U)=\dim(X)$.
\end{corollary}

\oldpage[IV-3]{109}
On sait en effet que $\dim(U)\leq\dim(X)$ ($\mathbf 0$, 14.1.4)~; en outre,
comme $X$ est un préschéma de Jacobson, $U$ contient un point fermé de toute
composante irréductible de $X$, donc $\dim(U)=\dim(X)$ en vertu de (10.6.1) et
($\mathbf 0$, 14.1.2.1).

\begin{proposition}[10.6.3]
\label{IV.10.6.3-fr}
Supposons que $X$ vérifie les conditions $1^\circ$, $2^\circ$ et $3^\circ$ de
(10.6.1), et soit $Y$ une partie fermée de $X$. Alors, pour tout $x\in Y$, et
tout voisinage ouvert $U$ de $x$ dans $X$ ne rencontrant pas les composantes
irréductibles de $Y$ qui ne contiennent pas $x$, on a
\begin{equation}
\label{IV.10.6.3.1-fr}
\tag{10.6.3.1}
\dim_x(Y)=\dim(U\cap Y)=\dim(\overline{\{x\}})
 +\operatorname{codim}(\overline{\{x\}},Y)=\sup_i\dim(Y_i)
\end{equation}
où $Y_i$ ($1\leq i\leq m$) sont les composantes irréductibles de $Y$ contenant
$x$.
\end{proposition}

En considérant un sous-préschéma fermé de $X$ ayant $Y$ pour espace sous-jacent,
on peut se borner, en vertu de (10.6.1, (i)), au cas où $Y=X$. En vertu du choix
de $U$, $U$ est réunion des $U\cap X_i$, donc
$\dim(U)=\sup_i\dim(U\cap X_i)$~; mais d'après (10.6.2) appliqué à un
sous-préschéma fermé de $X$ ayant $X_i$ pour espace sous-jacent, on a (compte
tenu de (10.6.1, (i))) $\dim(U\cap X_i)=\dim(X_i)$~; cela démontre que le second
et le quatrième terme de (10.6.3.1) sont égaux. Comme $U\cap X_i$ est
biéquidimensionnel (10.6.1, (i)) (puisque l'immersion $U\cap X_i\to X_i$ est de
type fini (\textbf{I}, 6.3.5)), on a
\begin{equation}
\label{IV.10.6.3.2-fr}
\tag{10.6.3.2}
\dim(X_i)=\dim(U\cap X_i)=\dim(U\cap\overline{\{x\}})
 +\operatorname{codim}(U\cap\overline{\{x\}},U\cap X_i)
\end{equation}
($\mathbf 0$, 14.3.5). D'après (10.6.2) appliqué à un sous-préschéma fermé de
$X$ ayant $\overline{\{x\}}$ pour espace sous-jacent,
$\dim(U\cap\overline{\{x\}})=\dim(\overline{\{x\}})$~; comme on a aussi, pour les
mêmes raisons,
\begin{equation}
\label{IV.10.6.3.3-fr}
\tag{10.6.3.3}
\dim(X_i)=\dim(\overline{\{x\}})
 +\operatorname{codim}(\overline{\{x\}},X_i)
\end{equation}
on obtient
\[
  \dim(U)=\dim(\overline{\{x\}})
  +\sup_i\operatorname{codim}(\overline{\{x\}},X_i)
  =\dim(\overline{\{x\}})+\operatorname{codim}(\overline{\{x\}},X)
\]
par définition de la codimension ($\mathbf 0$, 14.2.1). Ceci montre que $\dim(U)$
est indépendante du voisinage ouvert $U$ de $x$ vérifiant les conditions de
l'énoncé, donc est égal à $\dim_x(X)$, par ($\mathbf 0$, 14.1.4.1).

\begin{corollary}[10.6.4]
\label{IV.10.6.4-fr}
Sous les hypothèses de (10.6.3), soient $\mathcal F$ un $\mathcal O_X$-Module cohérent,
$Y$ son support. Pour tout $x\in Y$, on a
\begin{equation}
\label{IV.10.6.4.1-fr}
\tag{10.6.4.1}
\dim(\overline{\{x\}})+\dim(\mathcal F_x)=\dim_xY.
\end{equation}
\end{corollary}

En effet, cela résulte de (10.6.3.1) et de la formule
$\dim(\mathcal F_x)=\operatorname{codim}(\overline{\{x\}},Y)$ (5.1.12.2).

\subsection{Exemples et contre-exemples}
\label{subsection:IV.10.7-fr}

\begin{env}[10.7.1]
\label{IV.10.7.1-fr}
Soit $S$ un préschéma localement noethérien de dimension $\leq1$ et supposons
que $S$ soit un \emph{préschéma de Jacobson}~; lorsque $S$ est noethérien, il
revient au même de dire que les composantes irréductibles de $S$ de dimension
$1$ sont infinies, car tout $x\in S$ qui n'est pas fermé est le point générique
d'une telle composante (10.4.5 et 10.5.4). Alors $S$ vérifie aussi les conditions
$2^\circ$ et $3^\circ$ de (10.6.1)~: en effet, tout anneau local $\mathcal O_{S,s}$
est de dimension $0$ ou $1$, et par suite est universellement caténaire (7.2.9)~;
d'autre part, une composante irréductible $S'$ de $S$ est, soit réduite

\oldpage[IV-3]{110}
à un point, soit de dimension $1$, et pour tout point fermé $s\in S'$,
$\mathcal O_{S',s}$ est nécessairement de dimension $1$.

On déduit de ces remarques et de (10.6.1) que tout préschéma localement de type
fini sur $S$ vérifie aussi les propriétés $1^\circ$, $2^\circ$ et $3^\circ$ de
(10.6.1)~: il en est ainsi en particulier des préschémas \emph{localement de type
fini sur un corps ou sur $\mathbf Z$}.
\end{env}

\begin{env}[10.7.2]
\label{IV.10.7.2-fr}
Soit $A$ un anneau local noethérien \emph{universellement caténaire} et soit $S$
le \emph{complémentaire dans $X=\operatorname{Spec}(A)$ du point fermé} $a$.
Alors $S$ vérifie les conditions $1^\circ$, $2^\circ$ et $3^\circ$ de (10.6.1)~:
en effet, on a déjà vu que $S$ est un préschéma de Jacobson (10.5.9)~; comme $A$
est universellement caténaire, il en est de même des anneaux locaux $A_{\mathfrak p}$
aux idéaux premiers de $A$ (5.6.3). D'autre part, une composante irréductible $S'$
de $S$ est le complémentaire de $a$ dans une composante irréductible $X'$ de $X$~;
pour tout point fermé $x$ de $S$, l'adhérence de $x$ dans $X$ est donc
$\{x,a\}$, autrement dit $\dim(\overline{\{x\}})=1$ et par suite, puisque $X'$
est par hypothèse biéquidimensionnel ($\mathbf 0$, 14.3.3), on a
$\dim(\mathcal O_{X',x})=\dim(X')-1$ ($\mathbf 0$, 14.3.2), ce qui prouve que $S'$
est équicodimensionnel (5.1.4).
\end{env}

\begin{env}[10.7.3]
\label{IV.10.7.3-fr}
Soit $A$ un anneau de valuation discrète~; montrons que dans l'anneau
$B=A[T_1,\ldots,T_n]$ des polynômes à $n$ indéterminées, il existe deux idéaux
maximaux $\mathfrak m$, $\mathfrak n$ de \emph{hauteurs respectives} $n$ et
$n+1$. On l'a vu dans (5.2.5, (i)) pour $n=1$~; démontrons-le par récurrence sur
$n$. Comme $A[T_1,\ldots,T_{n+1}]$ est un $A[T_1,\ldots,T_n]$-module libre,
donc fidèlement plat, il y a dans $A[T_1,\ldots,T_{n+1}]$ deux idéaux maximaux
$\mathfrak m'$, $\mathfrak n'$ respectivement au-dessus de $\mathfrak m$ et
$\mathfrak n$ ($\mathbf 0_{\mathrm I}$, 6.5.1)~; en outre, d'après (5.5.3), ces
idéaux sont nécessairement de hauteurs respectives $n+1$ et $n+2$, d'où notre
assertion. Supposons $n\geq2$ dans ce qui suit. Soit
$\mathfrak J=\mathfrak m\cap\mathfrak n=\mathfrak m\mathfrak n$, et
$R=1+\mathfrak J$ qui est une partie multiplicative de $B$~; si l'on pose
$B'=R^{-1}B$, l'idéal $\mathfrak J'=R^{-1}\mathfrak J$ est contenu dans le
radical de $B'$ (Bourbaki, \emph{Alg. comm.}, chap.~III, \S~3,
n\textsuperscript{o}~5, prop.~12)~; on sait que $X'=\operatorname{Spec}(B')$
s'identifie en tant qu'espace topologique à un sous-espace de
$X=\operatorname{Spec}(B)$, et qu'aux points $x$ de $X'$, les anneaux locaux
$\mathcal O_{X,x}$ et $\mathcal O_{X',x}$ sont les mêmes (\textbf{I}, 1.6.2).
Considérons alors dans $X'$ l'ensemble fermé $Y'=V(\mathfrak J')$, et soit
$S=X'-Y'$~; on sait (10.5.7) que $S$ est un préschéma de Jacobson, évidemment
irréductible et noethérien~; en outre les anneaux locaux
$\mathcal O_{S,s}=\mathcal O_{X,s}$ sont universellement caténaires pour tout $s\in S$
en vertu de (5.6.3), puisque $A$ est universellement caténaire (5.6.4). Pourtant,
il y a deux points fermés $a$, $b$ de $S$ tels que $\mathcal O_{S,a}$ et
$\mathcal O_{S,b}$ \emph{n'aient pas même dimension}, autrement dit $S$ ne vérifie
pas la condition $3^\circ$ de (10.6.1). Pour le voir, considérons les deux idéaux
maximaux $\mathfrak m'=R^{-1}\mathfrak m$, $\mathfrak n'=R^{-1}\mathfrak n$ de
$B'$, qui sont de hauteurs $n$ et $n+1$ respectivement~; on a
$\mathfrak J'=\mathfrak m'\cap\mathfrak n'$, et $\mathfrak J'$ n'est donc
contenu dans aucun idéal premier de $B'$ distinct de $\mathfrak m'$ et
$\mathfrak n'$, qui sont par suite les seuls idéaux maximaux de $B'$. Soient
$a'$, $b'$ les seuls points fermés de $X'$, correspondant à $\mathfrak m'$ et
$\mathfrak n'$. Il existe dans $B'$ un idéal premier non maximal
$\mathfrak p'\subset\mathfrak m'$ qui n'est pas contenu dans $\mathfrak n'$~:
il suffit de montrer qu'il y a dans $B$ un idéal premier non maximal contenu dans
$\mathfrak m$ et non dans $\mathfrak n$~; pour cela, on pourra par exemple
considérer la fibre de $\mathfrak m$ pour le morphisme correspondant à l'injection
$A[T_1]\to B=A[T_1,\ldots,T_n]$, et appliquer (6.1.2). En considérant une chaîne
maximale d'idéaux premiers entre $\mathfrak p'$ et $\mathfrak m'$ et remplaçant
$\mathfrak p'$ par l'avant-dernier idéal de cette chaîne, on peut donc supposer
que le point $a$ de $X'$ correspondant à $\mathfrak p'$ est tel que son adhérence
dans $X'$ soit $\{a,a'\}$~; comme $B'_{\mathfrak m'}$ est biéquidimensionnel,
$\mathfrak p'$ est alors de hauteur $n-1$. On construit de la même manière un
idéal premier non maximal $\mathfrak q'$ de $B'$ de hauteur $n$, tel que si $b$
est le point correspondant de $X'$, l'adhérence de $b$ dans $X'$ soit
$\{b,b'\}$. Cela étant, $a$ et $b$ sont dans $S$, donc \emph{fermés dans $S$},
et répondent par suite à la question.
\end{env}

\subsection{Profondeur rectifiée}
\label{subsection:IV.10.8-fr}

\begin{definition}[10.8.1]
\label{IV.10.8.1-fr}
Soient $X$ un préschéma localement noethérien, $\mathcal F$ un $\mathcal O_X$-Module
cohérent. Pour tout $x\in X$, on appelle \emph{profondeur rectifiée de $\mathcal F$
au point $x$} et on note $\operatorname{prof}^{*}_{x}(\mathcal F)$ le nombre (entier
$\geq0$ ou $+\infty$) égal à
\begin{equation}
\label{IV.10.8.1.1-fr}
\tag{10.8.1.1}
\operatorname{prof}^{*}_{x}(\mathcal F)
=\operatorname{prof}(\mathcal F_x)+\dim(\overline{\{x\}})
\end{equation}
où $\overline{\{x\}}$ est l'adhérence du point $x$ dans $X$. Pour toute partie
$Z$ de $X$, on appelle \emph{profondeur rectifiée de $\mathcal F$ le long de $Z$}
et on note $\operatorname{prof}^{*}_{Z}(\mathcal F)$ le nombre
\begin{equation}
\label{IV.10.8.1.2-fr}
\tag{10.8.1.2}
\operatorname{prof}^{*}_{Z}(\mathcal F)
=\inf_{x\in Z}\operatorname{prof}^{*}_{x}(\mathcal F).
\end{equation}
\end{definition}

En d'autres termes pour tout entier $n$, la relation
$\operatorname{prof}^{*}_{Z}(\mathcal F)\geq n$ équivaut à
$\operatorname{prof}^{*}_{x}(\mathcal F)\geq n$ pour tout $x\in Z$. Si $Z=X$, on
écrit $\operatorname{prof}^{*}(\mathcal F)$ au lieu de
$\operatorname{prof}^{*}_{X}(\mathcal F)$.

\oldpage[IV-3]{111}
\begin{remarks}[10.8.2]
\label{IV.10.8.2-fr}
(i) En tout point fermé $x\in X$, la profondeur rectifiée est égale à la
profondeur.

(ii) Soient $Y$ un sous-préschéma fermé de $X$, $j:Y\to X$ l'injection canonique,
$\mathcal G$ un $\mathcal O_Y$-Module cohérent. On sait (5.7.3, (vi)) que l'on a
$\operatorname{prof}(\mathcal G_x)=\operatorname{prof}(j_*(\mathcal G)_x)$ pour tout
$x\in Y$~; on en déduit que l'on a aussi
$\operatorname{prof}^{*}_{x}(\mathcal G)
=\operatorname{prof}^{*}_{x}(j_*(\mathcal G))$ pour tout $x\in Y$.

(iii) La notion de profondeur rectifiée n'a d'intérêt que lorsqu'elle est de
caractère \emph{local}, c'est-à-dire qu'elle ne change pas lorsqu'on remplace $X$
par un voisinage ouvert arbitraire $U$ de $x$. Cela exige évidemment que $x$ ne
soit pas isolé dans $\overline{\{x\}}$ lorsque $x$ n'est pas fermé et par suite
que $X$ soit un \emph{préschéma de Jacobson} (10.4.5.1)~; le plus souvent, il sera
aussi nécessaire de savoir que $\dim(U)=\dim(X)$ pour tout ouvert dense $U$ dans
$X$, et il faudra donc supposer que $X$ vérifie aussi les conditions $2^\circ$ et
$3^\circ$ de (10.6.1).
\end{remarks}

\begin{lemma}[10.8.3]
\label{IV.10.8.3-fr}
Soient $X$ un préschéma régulier et biéquidimensionnel, $\mathcal F$ un
$\mathcal O_X$-Module cohérent. Alors on a, pour tout $x\in X$,
\begin{equation}
\label{IV.10.8.3.1-fr}
\tag{10.8.3.1}
\operatorname{prof}^{*}_{x}(\mathcal F)
=\dim(X)-\dim.\operatorname{proj}(\mathcal F_x).
\end{equation}
\end{lemma}

En effet, comme $X$ est biéquidimensionnel, on a ($\mathbf 0$, 14.3.5.1)
\[
  \dim(\overline{\{x\}})=\dim(X)-\operatorname{codim}(\overline{\{x\}},X)
  =\dim(X)-\dim(\mathcal O_{X,x})
\]
en vertu de (5.1.2). D'autre part, comme $X$ est régulier, on a par
($\mathbf 0$, 17.3.4)
\[
  \operatorname{prof}(\mathcal F_x)
  =\dim(\mathcal O_{X,x})-\dim.\operatorname{proj}(\mathcal F_x)
\]
d'où le lemme.

\begin{corollary}[10.8.4]
\label{IV.10.8.4-fr}
Sous les hypothèses de (10.8.3), la fonction
$x\mapsto\operatorname{prof}^{*}_{x}(\mathcal F)$ est semi-continue inférieurement.
\end{corollary}

Cela résulte de (10.8.3.1), puisque
$x\mapsto\dim.\operatorname{proj}(\mathcal F_x)$ est semi-continue supérieurement
(6.11.1).

\begin{proposition}[10.8.5]
\label{IV.10.8.5-fr}
Soient $S$ un préschéma localement noethérien, $X$ un préschéma localement de
type fini sur $S$, $\mathcal F$ un $\mathcal O_X$-Module cohérent. On suppose que $S$
vérifie les conditions suivantes~: $1^\circ$ $S$ est un préschéma de Jacobson~;
$2^\circ$ $S$ est régulier~; $3^\circ$ les composantes irréductibles de $S$ sont
équicodimensionnelles. Alors la fonction
$x\mapsto\operatorname{prof}^{*}_{x}(\mathcal F)$ est semi-continue inférieurement
dans $X$~; autrement dit, pour tout entier $n$, l'ensemble $U_n$ des $x\in X$
tels que $\operatorname{prof}^{*}_{x}(\mathcal F)\geq n$ est ouvert.
\end{proposition}

Comme les anneaux locaux $\mathcal O_{S,s}$ de $S$ sont réguliers, ils sont
universellement caténaires (5.6.4)~; autrement dit, $S$ vérifie les conditions
$1^\circ$, $2^\circ$ et $3^\circ$ de (10.6.1), donc il en est de même de $X$
(10.6.1, (i)). La notion de profondeur rectifiée étant alors de caractère local
(10.8.2, (iii)), on peut se borner au cas où $S=\operatorname{Spec}(A)$ et
$X=\operatorname{Spec}(B)$ sont affines, $A$ étant un anneau régulier et $B$ une
$A$-algèbre de type fini, donc quotient d'un anneau de polynômes
$C=A[T_1,\ldots,T_n]$, et ce dernier est régulier ($\mathbf 0$, 17.3.7). On peut
donc supposer que $X$ est un sous-préschéma fermé d'un préschéma régulier $Y$
vérifiant également les conditions $1^\circ$, $2^\circ$ et $3^\circ$ de (10.6.1)~;
compte tenu de la remarque (10.8.2, (ii)), on est ainsi ramené au cas où $X$ est
en outre régulier et noethérien. Mais comme les anneaux locaux de $X$ sont alors
intègres, les composantes irréductibles de $X$ sont

\oldpage[IV-3]{112}
ouvertes (\textbf{I}, 6.1.10), et on peut par suite supposer aussi $X$
irréductible. Alors, comme les anneaux locaux de $X$ sont caténaires
($\mathbf 0$, 16.5.12), l'hypothèse $3^\circ$ de (10.6.1) entraîne que $X$ est
biéquidimensionnel ((5.1.5) et ($\mathbf 0$, 14.3.3))~; il suffit donc d'appliquer
(10.8.4).

On notera que si $S$ est le \emph{spectre d'un corps ou de $\mathbf Z$}, il
vérifie les conditions de (10.8.5).

\begin{corollary}[10.8.6]
\label{IV.10.8.6-fr}
Les hypothèses étant celles de (10.8.5), pour tout $x\in X$, le nombre
$\operatorname{prof}^{*}_{x}(\mathcal F)$ est l'unique entier $n$ ayant la propriété
suivante~: il existe un voisinage ouvert $U$ de $x$ dans $X$ tel que pour tout
point $x'\in U\cap\overline{\{x\}}$, fermé dans $U$, on ait
$\operatorname{prof}(\mathcal F_{x'})=n$. En particulier, pour que
$\operatorname{prof}^{*}_{x}(\mathcal F)\geq m$ (resp.
$\operatorname{prof}^{*}_{x}(\mathcal F)\leq m$), il faut et il suffit qu'il existe
un voisinage ouvert $V$ de $x$ dans $X$ tel que, pour tout
$x'\in V\cap\overline{\{x\}}$, fermé dans $V$, on ait
$\operatorname{prof}(\mathcal F_{x'})\geq m$ (resp.
$\operatorname{prof}(\mathcal F_{x'})\leq m$).
\end{corollary}

En effet, on peut se borner au cas où $x$ n'est pas fermé~; si
$\operatorname{prof}^{*}_{x}(\mathcal F)=n$, l'ensemble des $y\in X$ tels que
$\operatorname{prof}^{*}_{y}(\mathcal F)\leq n$ est fermé en vertu de (10.8.5), donc
contient $\overline{\{x\}}$, et en vertu de la semi-continuité inférieure de
$y\mapsto\operatorname{prof}^{*}_{y}(\mathcal F)$, il existe un voisinage ouvert $U$
de $x$ tel que $\operatorname{prof}^{*}_{y}(\mathcal F)=n$ pour tout
$y\in U\cap\overline{\{x\}}$, donc $\operatorname{prof}(\mathcal F_{x'})=n$ si
$x'\in U\cap\overline{\{x\}}$ est fermé dans $U$ (puisque la notion de profondeur
rectifiée est locale).

Pour les préschémas vérifiant les hypothèses de (10.8.5), la notion de profondeur
rectifiée peut donc se définir à l'aide des valeurs de la profondeur aux
\emph{points fermés} de $X$ (ces derniers formant un ensemble \emph{très dense}
dans toute partie fermée de $X$).

\begin{proposition}[10.8.7]
\label{IV.10.8.7-fr}
Soit $S$ un préschéma localement noethérien vérifiant les conditions $1^\circ$,
$2^\circ$ et $3^\circ$ de (10.6.1). Soient $X$ un préschéma localement de type
fini sur $S$, $\mathcal F$ un $\mathcal O_X$-Module cohérent,
$Y=\operatorname{Supp}(\mathcal F)$. Alors, pour tout $x\in Y$, on a
\begin{equation}
\label{IV.10.8.7.1-fr}
\tag{10.8.7.1}
\operatorname{prof}^{*}_{x}(\mathcal F)
=\dim_x(Y)-\operatorname{coprof}(\mathcal F_x).
\end{equation}
\end{proposition}

En effet, par définition, on a
$\operatorname{coprof}(\mathcal F_x)
=\dim(\mathcal F_x)-\operatorname{prof}(\mathcal F_x)$, et il résulte de (10.6.4) que
l'on a $\dim_x(Y)=\dim(\overline{\{x\}})+\dim(\mathcal F_x)$~; d'où (10.8.7.1) par
définition de $\operatorname{prof}^{*}_{x}(\mathcal F)$.

\begin{corollary}[10.8.8]
\label{IV.10.8.8-fr}
Les hypothèses sur $f$ et $\mathcal F$ étant celles de (9.9.1), la fonction
$x\mapsto\operatorname{prof}^{*}_{x}(\mathcal F_{f(x)})$ est constructible.
\end{corollary}

Notons que pour tout $s\in S$, la fibre $X_s$ étant un préschéma localement de
type fini sur $k(s)$, est un préschéma de Jacobson (10.4.7)~; comme en outre
$\operatorname{Spec}(k(s))$ vérifie les conditions de (10.6.1), on a
\[
  \operatorname{prof}^{*}_{x}(\mathcal F_{f(x)})
  =\dim_x(\operatorname{Supp}(\mathcal F_{f(x)}))
  -\operatorname{coprof}((\mathcal F_{f(x)})_x)
\]
(10.8.7). Or, si $Z=\operatorname{Supp}(\mathcal F)$, on a
$Z_{f(x)}=\operatorname{Supp}(\mathcal F_{f(x)})$ (\textbf{I}, 9.1.13) et $Z$ est
localement constructible (8.9.1)~; donc les fonctions
$x\mapsto\dim_x(\operatorname{Supp}(\mathcal F_{f(x)}))$ et
$x\mapsto\operatorname{coprof}((\mathcal F_{f(x)})_x)$ sont localement
constructibles ((9.9.1) et (9.9.3)), ce qui prouve la proposition.

\subsection{Spectres maximaux et ultrapréschémas}
\label{subsection:IV.10.9-fr}

Les résultats de ce numéro ne seront pas utilisés par la suite.

\begin{env}[10.9.1]
\label{IV.10.9.1-fr}
Soit $X$ un préschéma de Jacobson, et soit $S(X)$ l'\emph{espace annelé} dont
l'espace sous-jacent est le sous-espace des points fermés de $X$, et le faisceau
d'anneaux le faisceau induit sur ce sous-espace par $\mathcal O_X$, autrement dit
le faisceau d'anneaux $\psi^*(\mathcal O_X)$, en désignant par
$\psi:S(X)\to X$ l'injection canonique. Comme $\psi$ est un quasi-
\oldpage[IV-3]{113}
homéomorphisme, on a vu (10.2.8, (ii)) que si
$\theta:\mathcal O_X\to\psi_*(\mathcal O_{S(X)})$ est l'homomorphisme de
faisceaux d'anneaux tel que
$\theta^\sharp:\psi^*(\mathcal O_X)\to\mathcal O_{S(X)}$ soit l'identité,
alors $j_X=(\psi,\theta)$ est un \emph{quasi-isomorphisme} d'espaces annelés,
et $\mathcal F\mapsto j_X^*(\mathcal F)$ une équivalence de la catégorie des
$\mathcal O_X$-Modules et de celle des $\mathcal O_{S(X)}$-Modules. Il est clair
que dans cette équivalence, aux $\mathcal O_X$-Modules localement libres (resp.
cohérents) correspondent les $\mathcal O_{S(X)}$-Modules localement libres
(resp. cohérents)~; en outre, si $\mathcal L$ est un $\mathcal O_X$-Module
localement libre et si $U$ est un ouvert de $X$ tel que $\mathcal L|U$ soit
isomorphe à $(\mathcal O_X|U)^n$, $j_X^*(\mathcal L)$ est tel que
$j_X^*(\mathcal L)|(U\cap S(X))$ soit isomorphe à
$(\mathcal O_{S(X)}|(U\cap S(X)))^n$.
\end{env}

\begin{env}[10.9.2]
\label{IV.10.9.2-fr}
Soient $X$, $Y$ deux préschémas de Jacobson et
$f=(\rho,\lambda):X\to Y$ un morphisme \emph{localement de type fini}. On a vu
(10.4.7) que l'on a $\rho(S(X))\subset S(Y)$ et par restriction de $\rho$ à
$S(X)$, on définit donc une application continue $S(\rho):S(X)\to S(Y)$.
D'autre part, pour tout ouvert $V$ de $Y$, on définit par composition un
homomorphisme d'anneaux
\[
  \Gamma(V\cap S(Y),\mathcal O_{S(Y)})\longrightarrow
  \Gamma(V,\mathcal O_Y)\xrightarrow{\lambda_V}
  \Gamma(f^{-1}(V),\mathcal O_X)\longrightarrow
  \Gamma(f^{-1}(V)\cap S(X),\mathcal O_{S(X)})
\]
où les deux isomorphismes extrêmes ont été définis dans (10.9.1)~; il est clair
que cela définit un homomorphisme de faisceaux d'anneaux
$S(\lambda):\mathcal O_{S(Y)}\to\rho_*(\mathcal O_{S(X)})$ (en se rappelant
que les ouverts de $S(X)$ (resp. $S(Y)$) correspondent biunivoquement à ceux de
$X$ (resp. $Y$) (10.2.1))~; on obtient ainsi un morphisme d'espaces annelés
\[
  S(f)=(S(\rho),S(\lambda)):(S(X),\mathcal O_{S(X)})\longrightarrow
  (S(Y),\mathcal O_{S(Y)})
\]
tel que le diagramme
\[
\begin{tikzcd}
S(X) \ar[r,"S(f)"] \ar[d,"j_X"'] & S(Y) \ar[d,"j_Y"] \\
X \ar[r,"f"'] & Y
\end{tikzcd}
\]
soit commutatif~; en outre, si $Z$ est un troisième préschéma de Jacobson et
$g:Y\to Z$ un morphisme localement de type fini, il est clair que
$S(g\circ f)=S(g)\circ S(f)$. On a ainsi défini un \emph{foncteur covariant}
$S:\mathbf C\to\mathbf C'$, où $\mathbf C'$ est la catégorie des \emph{espaces
annelés en anneaux locaux}, et $\mathbf C$ la catégorie dont les objets sont les
\emph{préschémas de Jacobson} et les morphismes sont les morphismes
\emph{localement de type fini} entre préschémas de Jacobson.
\end{env}

\begin{env}[10.9.3]
\label{IV.10.9.3-fr}
Proposons-nous de déterminer la sous-catégorie $\mathbf C''$ de $\mathbf C'$
formée des espaces annelés isomorphes aux $S(X)$ et dont les morphismes
proviennent des $S(f)$. Supposons d'abord que $X=\operatorname{Spec}(A)$, où
$A$ est un anneau de Jacobson~; alors $S(X)$ est l'ensemble des \emph{idéaux
maximaux} de $A$, muni~: $1^\circ$ de la topologie induite par celle de $X$, de
sorte qu'une base de cette topologie est formée des
$D^m(h)=D(h)\cap S(X)$, ensemble des idéaux maximaux $\mathfrak m$ de $A$ tels
que $h\notin\mathfrak m$, où $h$ parcourt $A$~; $2^\circ$ du faisceau d'anneaux
$\mathcal O_{S(X)}$ tel que
$\Gamma(D^m(h),\mathcal O_{S(X)})=A_h$. Nous dirons que cet espace annelé est
le \emph{spectre maximal} de l'anneau de Jacobson $A$ et nous le noterons
$\operatorname{Spm}(A)$.

On notera que si $j:D(h)\to X$ est l'injection canonique, l'espace annelé
induit sur $D^m(h)$ par $S(X)$ est $S(D(h))$ et l'injection canonique
$D^m(h)\to S(X)$ d'espaces annelés est égale à $S(j)$.

Soient $B$ un second anneau de Jacobson, $Y=\operatorname{Spec}(B)$,
$\varphi:B\to A$ un homomorphisme d'anneaux faisant de $A$ une $B$-algèbre de
type fini, $f=({}^a\varphi,\widetilde\varphi):X\to Y$ le morphisme
correspondant de préschémas, et
\[
  S(f):\operatorname{Spm}(A)\longrightarrow\operatorname{Spm}(B)
\]
le morphisme d'espaces annelés correspondant à $f$. Il est clair que
$S(f)=(\psi,\theta)$ est un morphisme d'espaces annelés en anneaux locaux,
c'est-à-dire ($\mathrm{Err}_{\mathrm{II}}$, 1.8.2) que pour tout
$x\in\operatorname{Spm}(A)$, $\theta_x^\sharp$ est un homomorphisme local.
Réciproquement~:
\end{env}

\begin{proposition}[10.9.4]
\label{IV.10.9.4-fr}
Soient $A$, $B$ deux anneaux de Jacobson. Si
$u=(\psi,\theta):\operatorname{Spm}(A)\to\operatorname{Spm}(B)$ est un
morphisme d'espaces annelés en anneaux locaux tel que
$\Gamma(\theta):B\to A$ fasse de $A$ une $B$-algèbre de type fini, il existe
un morphisme de préschémas
$f:\operatorname{Spec}(A)\to\operatorname{Spec}(B)$ et un seul tel que
$u=S(f)$.
\end{proposition}

L'unicité de $f$ est évidente, puisque si
$f=({}^a\varphi,\widetilde\varphi)$, on doit avoir
$\varphi=\Gamma(\theta)$~; il reste à voir que $S(f)$ est défini et qu'on a
bien $u=S(f)$. Or, la première assertion résulte de ce que $\varphi$ est
supposé faire de $A$ une $B$-algèbre de type fini, et par suite
$f(\operatorname{Spm}(A))\subset\operatorname{Spm}(B)$~; le fait que
$\theta_x^\sharp$ est un homomorphisme local pour tout
$x\in\operatorname{Spm}(A)$ permet alors de répéter le raisonnement de
(\textbf I, 1.7.3) en se bornant aux points $x$ de $\operatorname{Spm}(A)$~:
on montre ainsi successivement que $\psi(x)={}^a\varphi(x)$ pour tout
$x\in\operatorname{Spm}(A)$, puis que
$\theta_x^\sharp=\varphi_x:B_{\psi(x)}\to A_x$, ce qui achève de prouver que
$u=S(f)$.

\oldpage[IV-3]{114}

\begin{env}[10.9.5]
\label{IV.10.9.5-fr}
Considérons maintenant un espace annelé $(X,\mathcal O_X)$~; nous dirons qu'une
partie ouverte $U$ de $X$ est \emph{ultra-affine} si l'espace annelé induit
$(U,\mathcal O_X|U)$ est isomorphe à un spectre maximal
$\operatorname{Spm}(A)$, où $A$ est un anneau de Jacobson. Nous dirons que $X$
est un \emph{ultra-préschéma} si tout point de $X$ admet un voisinage ouvert
ultra-affine. On montre, comme dans (\textbf I, 2.1.3 et 2.1.4) que les ouverts
ultra-affines forment une \emph{base} de la topologie de $X$ et que $X$ est un
espace de Kolmogoroff. Si $Y$ est un second ultra-préschéma, nous dirons qu'un
morphisme d'espaces annelés $f:X\to Y$ est un \emph{morphisme
d'ultra-préschémas} s'il vérifie les conditions suivantes~: $1^\circ$ $f$ est
un morphisme d'espaces annelés en anneaux locaux~; $2^\circ$ pour tout $x\in X$,
il y a un voisinage ouvert ultra-affine $V$ de $f(x)$ dans $Y$ et un voisinage
ouvert ultra-affine $U$ de $x$ dans $X$ tels que $f(U)\subset V$ et que
l'homomorphisme
$\Gamma(V,\mathcal O_Y)\to\Gamma(U,\mathcal O_X)$ correspondant à $f$ fasse
de $\Gamma(U,\mathcal O_X)$ une $\Gamma(V,\mathcal O_Y)$-algèbre de type fini.

Il est immédiat qu'on définit bien ainsi des morphismes, le composé de deux
morphismes en étant un troisième grâce à la remarque finale de (10.9.3). Il
est clair que la catégorie $\mathbf C''_0$ ainsi définie est une sous-catégorie
de $\mathbf C'$ qui contient $\mathbf C''$~; nous nous proposons de montrer que
$\mathbf C''=\mathbf C''_0$, autrement dit~:
\end{env}

\begin{proposition}[10.9.6]
\label{IV.10.9.6-fr}
Le foncteur $X\mapsto S(X)$ de $\mathbf C$ dans $\mathbf C''_0$ est une
équivalence de catégories.
\end{proposition}

$1^\circ$ Montrons d'abord que le foncteur $X\mapsto S(X)$ est
\emph{pleinement fidèle}, autrement dit que pour $X$, $Y$ dans $\mathbf C$,
l'application canonique
\[
  \operatorname{Hom}_{\mathbf C}(X,Y)\longrightarrow
  \operatorname{Hom}_{\mathbf C''_0}(S(X),S(Y))
\]
est \emph{bijective}. En premier lieu elle est \emph{injective}~: soient en
effet $f$, $g$ deux morphismes localement de type fini de $X$ dans $Y$ et
supposons que $S(f)=S(g)$. Cela entraîne d'abord que pour tout ouvert $V$ de
$Y$, on a
$f^{-1}(V)\cap S(X)=g^{-1}(V)\cap S(X)$, donc (10.3.1 et 10.2.7)
$f^{-1}(V)=g^{-1}(V)$~; il suffit donc de prouver que pour tout ouvert affine
$V$ de $Y$, $f$ et $g$ coïncident dans
$f^{-1}(V)=g^{-1}(V)$, autrement dit, on est ramené au cas où
$Y=\operatorname{Spec}(B)$ est le spectre d'un anneau de Jacobson $B$. Il
suffit (\textbf I, 2.2.4) de montrer que les homomorphismes d'anneaux
$B\to\Gamma(X,\mathcal O_X)$ correspondant à $f$ et $g$ sont alors les mêmes.
Or, pour tout $s\in B$, les images de $s$ par ces homomorphismes sont deux
sections de $\mathcal O_X$ au-dessus de $X$ qui, par hypothèse, induisent la
même section au-dessus de $S(X)$~; on sait donc (10.2.8) que ces sections sont
identiques, d'où notre assertion. Prouvons en second lieu que tout morphisme
$h:S(X)\to S(Y)$ (pour la catégorie $\mathbf C''_0$) est de la forme $S(f)$,
où $f:X\to Y$ est un morphisme localement de type fini. Par hypothèse, il
existe un recouvrement ouvert ultra-affine $(U'_\lambda)$ (resp.
$(V'_\mu)$) de $S(X)$ (resp. $S(Y)$) tel que pour tout $\lambda$,
$h(U'_\lambda)$ soit contenu dans un $V'_\mu$ et qu'il corresponde au
morphisme $h_\lambda:U'_\lambda\to V'_\mu$, restriction de $h$, un
homomorphisme d'anneaux
$\Gamma(V'_\mu,\mathcal O_{S(Y)})\to
\Gamma(U'_\lambda,\mathcal O_{S(X)})$ faisant du second anneau une algèbre de
type fini sur le premier. On peut supposer que
$U'_\lambda=U_\lambda\cap S(X)$ et $V'_\mu=V_\mu\cap S(Y)$, où
$U_\lambda$ et $V_\mu$ sont des ouverts affines uniquement déterminés dans
$X$ et $Y$ respectivement~; si l'on prouve que, pour tout $\lambda$,
$h_\lambda=S(f_\lambda)$ où $f_\lambda:U_\lambda\to V_\mu$ est un morphisme
de type fini, alors il résulte de la première partie du raisonnement,
appliquée aux restrictions de $f_\alpha$ et $f_\beta$ à
$U_\alpha\cap U_\beta$, que les $f_\lambda$ sont les restrictions d'un même
morphisme $f:X\to Y$, et on aura évidemment $h=S(f)$. On est ainsi ramené au
cas où $X$ et $Y$ sont affines, et la conclusion résulte alors de (10.9.4).

$2^\circ$ Il reste à prouver que tout ultra-préschéma $X'$ est de la forme
$S(X)$ pour un préschéma de Jacobson $X$ (qui sera nécessairement unique à
isomorphisme près, en vertu de $1^\circ$). Il y a un recouvrement
$(U'_\alpha)$ de $X'$ par des ouverts ultra-affines, dont chacun est de la
forme $S(U_\alpha)$, $U_\alpha$ étant le spectre d'un anneau de Jacobson. Pour
tout couple d'indices $\alpha$, $\beta$, considérons l'unique ouvert
$U_{\alpha\beta}$ de $U_\alpha$ dont la trace sur $U'_\alpha$ est
$U'_\alpha\cap U'_\beta$~; en vertu de $1^\circ$, l'automorphisme identique de
$U'_\alpha\cap U'_\beta$ est de la forme $S(\theta_{\alpha\beta})$, où
$\theta_{\alpha\beta}:U_{\alpha\beta}\to U_{\beta\alpha}$ est un
isomorphisme de préschémas. On constate aussitôt (en vertu du $1^\circ$) que la
famille $(\theta_{\alpha\beta})$ vérifie la condition de recollement
($\mathbf 0_{\mathrm I}$, 4.1.7), et que cette famille définit donc un
préschéma $X$, tel que les $U_\alpha$ s'identifient à des ouverts affines de
$X$~; il est clair alors que l'on a $X'=S(X)$, ce qui achève la démonstration.

\subsection{Espaces algébriques de Serre}
\label{subsection:IV.10.10-fr}

\begin{env}[10.10.1]
\label{IV.10.10.1-fr}
Le langage introduit par Serre dans (FAC) est parfois commode, notamment dans
les questions où l'intérêt majeur s'attache aux points rationnels sur le corps
de base $k$ (algébriquement clos par hypothèse) des «~variétés algébriques~»
sur $k$ que l'on considère. Nous allons esquisser ici ce langage en le
rattachant aux considérations qui précèdent, pour permettre au lecteur de
traduire les énoncés de Serre en langage des schémas. Il est d'ailleurs
possible de développer le langage de Serre également pour des préschémas sur
un corps non algébriquement clos (et même sur un anneau artinien)~; mais cela
introduit des complications techniques considérables, et d'ailleurs, sur un
corps de base arbitraire, les avantages (surtout psychologiques) du point de
vue de Serre disparaissent~; aussi
\oldpage[IV-3]{115}
nous tiendrons-nous dans le cadre fixé par Serre. Le présent numéro, tout comme
le précédent, ne sera d'ailleurs pas utilisé dans la suite de ce Traité, et
nous nous bornerons donc à de brèves indications.
\end{env}

\begin{env}[10.10.2]
\label{IV.10.10.2-fr}
Étant donné un ultra-préschéma fixe $R$, on peut naturellement (comme dans
toute catégorie) définir la notion de $R$-\emph{ultra-préschéma}. Considérons
en particulier un corps algébriquement clos $k$~; $\operatorname{Spm}(k)$ est
alors identique à $\operatorname{Spec}(k)$~; nous dirons qu'un
$\operatorname{Spm}(k)$-ultra-préschéma $X'$ est un \emph{espace
préalgébrique sur $k$}~: c'est donc un espace $k$-annelé en anneaux locaux dont
tout point a un voisinage ouvert isomorphe au spectre maximal d'une $k$-algèbre
de type fini~; il revient au même de dire (par (10.9.6)) que $X'=S(X)$, où $X$
est un préschéma localement de type fini sur $k$. Si $X$, $Y$, $Z$ sont trois
$k$-préschémas localement de type fini, il en est de même de
$X\times_ZY$ (\textbf I, 3.4), donc en vertu de (10.9.6), la notion de produit
existe dans la catégorie des espaces préalgébriques sur $k$ (aussi bien le
produit «~sur $k$~» $X'\times_kY'$ que le «~produit fibré~»
$X'\times_{Z'}Y'$, où $X'$, $Y'$, $Z'$ sont trois espaces préalgébriques sur
$k$). On peut donc définir le morphisme diagonal
$\Delta_{X'}:X'\to X'\times_kX'$ (qui n'est autre d'ailleurs que
$S(\Delta_X)$)~; on dit que $X'$ est un \emph{espace algébrique sur $k$} si
$X$ est un schéma, et il revient au même de dire que l'image de $\Delta_{X'}$
est une partie \emph{fermée} de $X'\times_kX'$.
\end{env}

\begin{env}[10.10.3]
\label{IV.10.10.3-fr}
Les simplifications qui proviennent de l'hypothèse que $k$ est algébriquement
clos tiennent d'abord à ce que, pour un $k$-préschéma $X$, localement de type
fini, il y a correspondance biunivoque entre \emph{points fermés de $X$},
\emph{points de $X$ à valeurs dans $k$} (\textbf I, 3.4.4) et \emph{points de
$X$ rationnels sur $k$} (\textbf I, 3.4.5), en vertu de (\textbf I, 6.4.2).
Cela montre en particulier (en vertu de (\textbf I, 3.4.3.1)) que pour deux
espaces $k$-préalgébriques $X'$, $Y'$, l'ensemble sous-jacent au produit
$X'\times_kY'$ est identique à l'\emph{ensemble produit} $X'\times Y'$ des
ensembles sous-jacents (mais bien entendu la topologie de l'espace sous-jacent
à $X'\times_kY'$ n'est pas la \emph{topologie produit} des topologies des
espaces sous-jacents à $X'$ et $Y'$, elle est en général strictement plus fine
que cette dernière).

D'autre part, les anneaux locaux $\mathcal O_x$ aux points d'un espace
préalgébrique $X'$ sur $k$ sont des $k$-algèbres dont on vient de voir que le
corps résiduel est \emph{isomorphe à $k$}~; si $A$ et $B$ sont deux telles
$k$-algèbres locales, tout $k$-homomorphisme $\varphi:A\to B$ est
nécessairement \emph{local}~: en effet, si un élément $x$ de l'idéal maximal de
$A$ était tel que $\varphi(x)$ soit inversible, il existerait $\lambda\in k$
non nul tel que $\varphi(x-\lambda.1)$ appartienne à l'idéal maximal de $B$, ce
qui est absurde puisque $x-\lambda.1$ est inversible dans $A$. On conclut
aussitôt de là que si $X'=S(X)$, $Y'=S(Y)$ sont deux espaces préalgébriques sur
$k$, tout morphisme $X'\to Y'$ d'espaces $k$-annelés est aussi un morphisme
d'espaces $k$-annelés en anneaux locaux (\textbf I, 1.8.2~; cf.
$\mathrm{Err}_{\mathrm{II}}$). D'ailleurs, avec les notations précédentes, si
$A$ et $B$ sont des $k$-algèbres de type fini, $\varphi$ fait de $B$ une
$A$-algèbre de type fini~; donc tout morphisme $X'\to Y'$ de $k$-espaces
annelés est, en vertu de (10.9.6), de la forme $S(f)$, où $f:X\to Y$ est un
morphisme de $k$-préschémas.

Enfin, pour tout ouvert $U$ de $X'$, toute section
$s\in\Gamma(U,\mathcal O_{X'})$, et tout $x\in U$, $s(x)$
($\mathbf 0_{\mathrm I}$, 5.5.1) s'identifie à un élément de $k$, et l'on a
ainsi associé à $s$ une application $\overline s:x\mapsto s(x)$ de $U$ dans
$k$, autrement dit une section au-dessus de $U$ du faisceau $\mathcal A(X')$
des germes d'applications de $X'$ dans $k$~; comme l'application
$h_U:s\mapsto\overline s$ est évidemment un homomorphisme d'anneaux
$\Gamma(U,\mathcal O_{X'})\to\Gamma(U,\mathcal A(X'))$ et commute aux
restrictions à un ouvert $V\subset U$, les $h_U$ définissent un homomorphisme
de faisceaux d'anneaux $h:\mathcal O_{X'}\to\mathcal A(X')$. Si l'on prend
pour $U$ un ouvert ultra-affine $\operatorname{Spm}(A)$, où $A$ est un anneau
de Jacobson, dire que $\overline s=0$ signifie que pour tout idéal maximal
$\mathfrak m$ de $A$, $s$ appartient à $\mathfrak m$, ou encore que $s$ est
dans le \emph{radical} de $A$~; mais comme $A$ est un anneau de Jacobson, son
radical est égal à son nilradical~; pour que $h_U$ soit injective, il faut et
il suffit par suite que $A$ soit réduit.
\end{env}

\begin{env}[10.10.4]
\label{IV.10.10.4-fr}
On dit que l'espace $k$-préalgébrique $X'=S(X)$ est \emph{réduit} s'il en est
ainsi de $X$~; comme l'ensemble des points $x\in X$ où $X$ est réduit est
ouvert ($\mathbf 0_{\mathrm I}$, 5.2.2), son complémentaire contient au moins
un point fermé s'il est non vide (5.1.11), et il revient donc au même de dire
que $X'$ est réduit ou que chacun de ses anneaux locaux $\mathcal O_x$ (pour
$x\in X'$) est réduit. On vient de voir dans (10.10.3) que pour que
l'homomorphisme $h:\mathcal O_{X'}\to\mathcal A(X')$ soit \emph{injectif}, il
faut et il suffit que $X'$ soit réduit. Dans (FAC), Serre se borne en fait aux
espaces préalgébriques \emph{réduits}, ce qui lui permet de définir
$\mathcal O_{X'}$ comme un \emph{sous-faisceau} de $\mathcal A(X')$. On notera
que si $X'$ et $Y'$ sont des $k$-espaces préalgébriques réduits, il en est de
même de $X'\times_kY'$~: en effet, tout revient à voir que si $A$ et $B$ sont
deux $k$-algèbres de type fini réduites, il en est de même de $A\otimes_kB$~;
mais nous avons vu que les radicaux de $A$ et $B$ sont alors réduits à $0$, et
comme $k$ est algébriquement clos, $A$ et $B$ sont des algèbres «~séparables~»
sur $k$ au sens de Bourbaki (Bourbaki, \emph{Alg.}, chap.~VIII, \S~7,
n$^\circ$~5, prop.~5)~; donc $A\otimes_kB$ est sans radical (\emph{loc.
cit.}, n$^\circ$~6, cor.~3 du th.~3), et puisque c'est un anneau de Jacobson,
il est réduit. Toutefois, si $Z'$ est un troisième espace préalgébrique sur
$k$, le «~produit fibré~» $X'\times_{Z'}Y'$ de deux espaces préalgébriques
réduits sur $Z'$ n'est pas réduit en général, ce qui implique que la catégorie
de ces espaces est insuffisante dans de nombreuses questions (notamment dans
la théorie des groupes algébriques). Mais comme on l'a vu ci-dessus, on peut
garder le langage de Serre sans se limiter, comme ce dernier (qui en outre ne
considère que des espaces préalgébriques quasi-compacts), au cas des espaces
préalgébriques réduits.
\end{env}

\oldpage[IV-3]{116}

\begin{env}[10.10.5]
\label{IV.10.10.5-fr}
Enfin, on peut aussi considérer des ultra-préschémas sur un corps $k$
quelconque en conservant un langage qui reste proche de celui de Serre, et en
introduisant, comme chez Weil, une extension algébriquement close fixe $K$ de
$k$ (choisie assez grande, par exemple de degré de transcendance infini sur
$k$, pour avoir suffisamment de «~points génériques~» au sens de Weil). À tout
préschéma $X$ localement de type fini sur $k$, on associe alors l'ensemble
$S_K(X)=S(X\otimes_kK)$ des points de $X$ à valeurs dans $K$~; on a une
application canonique $j:S_K(X)\to X$ que l'on démontre être un
quasi-homéomorphisme quand on munit $S_K(X)$ de la topologie image réciproque
de celle de $X$ par $j$~; on munit $S_K(X)$ du faisceau de $k$-algèbres
$j^*(\mathcal O_X)$, et on obtient ainsi une sous-catégorie de la catégorie des
espaces $k$-annelés en anneaux locaux, que l'on pourrait appeler catégorie des
$(k,K)$-espaces préalgébriques. On peut montrer qu'on y peut encore définir des
produits et généraliser les résultats de (10.10.3) et (10.10.4)
($\mathcal A(X')$ étant ici remplacé par le faisceau $\mathcal A_K(X')$ des
germes d'applications de $X'$ dans $K$). Toutefois ce point de vue fait jouer
un rôle artificiel à un surcorps arbitrairement choisi de $k$, et nous ne le
signalons que pour le rejeter.
\end{env}

% FIN SOURCE ega4/ega4-10-fr.tex

%
% DEBUT SOURCE ega4/ega4-11-fr.tex
\oldpage[IV-3]{116}
\section{Propriétés topologiques des morphismes plats de présentation finie. Critères de platitude}
\label{section:IV.11-fr}

Alors que dans le \S~2 nous avons considéré les énoncés concernant la platitude
qui ne dépendent d'aucune hypothèse de finitude, et que le \S~6 étudie la notion
de platitude dans le cadre des préschémas localement noethériens (mais sans
hypothèse de finitude sur les morphismes), le présent paragraphe est consacré à
la notion de $f$-platitude dans le cas où le morphisme $f:X\to Y$ est
\emph{localement de présentation finie}. L'intérêt de la notion de morphisme
plat de présentation finie tient au fait que c'est elle qui semble exprimer
techniquement de la façon la plus adéquate la notion intuitive de «~famille de
préschémas algébriques paramétrée par un schéma $Y$~», dont l'étude est un des
objets principaux de la Géométrie algébrique. D'ailleurs, alors même qu'on ne
s'intéresserait au départ qu'au cas d'un schéma de base noethérien, il est
indispensable, pour certaines raisons techniques (par exemple pour certaines
applications de la théorie de la «~descente~», qui conduit à introduire des
schémas non nécessairement noethériens) de ne pas se limiter à ce cas, dès
qu'il s'agit de problèmes de nature essentiellement relative liés aux
morphismes localement de présentation finie. Nous suivrons systématiquement ce
principe, déjà étayé par les résultats des \S\S~8 et 9, dans toute la suite de ce
Chapitre, et même dans la suite de notre Traité, quitte à lui sacrifier à
l'occasion la simplicité de certaines démonstrations, que des hypothèses
noethériennes permettent parfois d'alléger\footnote{Ce principe s'inspire
également de la nécessité de donner droit de cité, comme «~espaces de
paramètres~» pour les familles de schémas algébriques, à des espaces annelés
(et même des «~topos~» annelés) quelconques, pour lesquels il ne peut plus être
question en général d'hypothèses noethériennes. Il semble assez clair que l'on
ne pourra plus éluder longtemps cette nouvelle extension de la Géométrie
algébrique, et il convient dès à présent de développer les notions et
techniques de nature «~relative~» de la théorie des schémas de sorte qu'elles
puissent s'adapter pratiquement telles quelles à ce cadre plus général.}. Dans
le présent paragraphe, cela nous conduit à reprendre, dans le contexte de la
«~présentation finie~» (notamment au n$^\circ$~3) certains énoncés de platitude,
déjà obtenus dans le contexte noethérien. L'outil technique essentiel pour
faire la réduction au cas noethérien est le théorème de compatibilité de la
platitude avec les limites projectives de préschémas (11.2.6), complétant les
résultats généraux du \S~8. Nous prouvons aussi en passant (11.3.1) un résultat
souvent utilisé

\oldpage[IV-3]{117}
par la suite, impliquant que l'ensemble des points de platitude d'un morphisme
localement de présentation finie est ouvert.

Dans les n$^\mathrm{os}$~4 à 8, nous étudions la question de la «~descente~» de
la platitude, consistant à trouver des conditions utiles sur un morphisme de
changement de base $Y'\to Y$ (non plat en général) pour pouvoir conclure que si
$X\times_YY'$ est plat sur $Y'$, $X$ est plat sur $Y$. Ces résultats, plus
techniques que ceux des n$^\mathrm{os}$~1 à 3, sont d'une utilisation moins
fréquente dans la suite~; ils joueront cependant un rôle important dans les
techniques de construction non projectives, au chapitre suivant. Le seul
résultat des n$^\mathrm{os}$~4 à 8 utilisé dans la suite du chap.~IV est le
critère valuatif de platitude (11.8), qu'on appliquera dans (15.2).

Enfin, les n$^\mathrm{os}$~9 et 10 sont consacrés à l'étude d'une notion qui
précise, en théorie des schémas, celle de densité au sens topologique, savoir la
notion de famille de sous-préschémas \emph{schématiquement dense} dans un
préschéma donné, et notamment l'étude du comportement de cette notion par
changement de base (plat ou quelconque). Cette notion est surtout utilisée,
pour l'instant, dans l'étude des schémas en groupes.

\subsection{Ensembles de platitude (cas noethérien)}
\label{subsection:IV.11.1-fr}

\begin{theorem}[11.1.1]
\label{IV.11.1.1-fr}
Soient $Y$ un préschéma localement noethérien, $f:X\to Y$ un morphisme
localement de type fini, $\mathcal F$ un $\mathcal O_X$-Module cohérent. Alors
l'ensemble $U$ des $x\in X$ tels que $\mathcal F$ soit $f$-plat au point $x$
est ouvert dans $X$.
\end{theorem}

La question étant évidemment locale sur $X$ et $Y$, on peut supposer
$Y=\operatorname{Spec}(A)$, $X=\operatorname{Spec}(B)$, $A$ étant noethérien et
$B$ une $A$-algèbre de type fini. On a alors $\mathcal F=\widetilde M$, où $M$
est un $B$-module de type fini. Appliquons le critère
($\mathbf 0_{\mathrm{III}}$, 9.2.6)~: il suffit donc de démontrer l'assertion
suivante~:

\begin{env}[11.1.1.1]
\label{IV.11.1.1.1-fr}
\emph{Soient $A$ un anneau noethérien, $B$ une $A$-algèbre de type fini, $M$ un
$B$-module de type fini, $\mathfrak q$ un idéal premier de $B$, $\mathfrak p$
son image réciproque dans $A$. On suppose que $M_{\mathfrak q}$ soit un
$A_{\mathfrak p}$-module plat. Alors il existe $g\in B-\mathfrak q$ tel que,
pour tout idéal premier $\mathfrak q'\supset\mathfrak q$ de $B$ tel que
$g\notin\mathfrak q'$, $M_{\mathfrak q'}$ soit un
$A_{\mathfrak p'}$-module plat, en désignant par $\mathfrak p'$ l'image
réciproque de $\mathfrak q'$ dans $A$ (il revient d'ailleurs au même
($\mathbf 0_{\mathrm I}$, 6.3.1) de dire que $M_{\mathfrak q'}$ est un
$A$-module plat).}
\end{env}

Considérons pour cela $B_{\mathfrak q'}$ comme une $A$-algèbre~; on a
évidemment
$\mathfrak pB_{\mathfrak q'}\subset\mathfrak q'B_{\mathfrak q'}$. On sait
alors ($\mathbf 0_{\mathrm{III}}$, 10.2.2) que pour que $M_{\mathfrak q'}$ soit
un $A$-module plat, il faut et il suffit que
$M_{\mathfrak q'}/\mathfrak pM_{\mathfrak q'}$ soit un
$(A/\mathfrak p)$-module plat et que l'on ait
$\operatorname{Tor}_1^A(M_{\mathfrak q'},A/\mathfrak p)=0$. Or, on a
$M_{\mathfrak q'}=M\otimes_BB_{\mathfrak q'}$~; $B_{\mathfrak q'}$ étant plat
sur $B$, on a
\[
  M_{\mathfrak q'}/\mathfrak pM_{\mathfrak q'}
  =(M/\mathfrak pM)_{\mathfrak q'}
\]
et
\[
  \operatorname{Tor}_1^A(M_{\mathfrak q'},A/\mathfrak p)
  =(\operatorname{Tor}_1^A(M,A/\mathfrak p))_{\mathfrak q'}
\]
(en définissant le Tor à l'aide d'une résolution projective de
$A/\mathfrak p$)~; pour la même raison, comme on doit avoir
$g\notin\mathfrak q'$, $B_{\mathfrak q'}$ est plat sur $B_g$, donc
\[
  (M/\mathfrak pM)_{\mathfrak q'}
  =((M/\mathfrak pM)_g)_{\mathfrak q'B_g}
\]
et
\[
  (\operatorname{Tor}_1^A(M,A/\mathfrak p))_{\mathfrak q'}
  =((\operatorname{Tor}_1^A(M,A/\mathfrak p))_g)_{\mathfrak q'B_g},
\]
où dans ces formules, $M/\mathfrak pM$ et
$\operatorname{Tor}_1^A(M,A/\mathfrak p)$ sont considérés comme des
$B$-modules. Compte tenu de ($\mathbf 0_{\mathrm I}$, 6.3.2), on voit qu'on
est ramené à prouver le

\begin{lemma}[11.1.1.2]
\label{IV.11.1.1.2-fr}
Sous les conditions de (11.1.1.1), il existe $g\in B-\mathfrak q$ tel que~:
\begin{enumerate}[label=(\roman*)]
  \item $(M/\mathfrak pM)_g$ soit un $(A/\mathfrak p)$-module plat~;
  \item $(\operatorname{Tor}_1^A(M,A/\mathfrak p))_g=0$.
\end{enumerate}
\end{lemma}

\oldpage[IV-3]{118}

En vertu du théorème de platitude générique (6.9.1) appliqué à l'anneau
intègre $A/\mathfrak p$, à la $(A/\mathfrak p)$-algèbre de type fini
$B/\mathfrak pB$ et au $(B/\mathfrak pB)$-module de type fini
$M/\mathfrak pM$, il existe un $h\in A-\mathfrak p$ tel que, si $\bar h$ est
son image canonique dans $A/\mathfrak p$,
$(M/\mathfrak pM)_{\bar h}$ soit un $(A/\mathfrak p)$-module plat. D'autre
part, comme $M_{\mathfrak q}$ est un $A_{\mathfrak p}$-module plat, et par
suite un $A$-module plat ($\mathbf 0_{\mathrm I}$, 6.3.1), on a
$\operatorname{Tor}_1^A(M_{\mathfrak q},A/\mathfrak p)=0$, ce qui s'écrit
aussi
$(\operatorname{Tor}_1^A(M,A/\mathfrak p))_{\mathfrak q}=0$. Mais comme $A$
et $B$ sont noethériens, $\operatorname{Tor}_1^A(M,A/\mathfrak p)$ est un
$B$-module de type fini, donc ($\mathbf 0_{\mathrm I}$, 5.2.2) il existe
$g\in B-\mathfrak q$ tel que
$(\operatorname{Tor}_1^A(M,A/\mathfrak p))_g=0$. D'ailleurs, on a
$(M/\mathfrak pM)_{\bar h}=(M/\mathfrak pM)_h$ ($M/\mathfrak pM$ étant
considéré dans le second membre comme un $A$-module)~; de plus,
$(M/\mathfrak pM)_h$ étant un $B$-module,
$(M/\mathfrak pM)_{hg}$ est encore un $(A/\mathfrak p)$-module plat, car il
s'écrit $(M/\mathfrak pM)_{\overline{hg}}$, où $\overline{hg}$ est l'image
canonique de $hg$ dans $B/\mathfrak pB$, et il suffit d'appliquer
($\mathbf 0_{\mathrm I}$, 6.3.2)~; enfin, on a
$(\operatorname{Tor}_1^A(M,A/\mathfrak p))_{hg}=0$ et
$hg\in B-\mathfrak q$ puisque $h\in A-\mathfrak p$. C.Q.F.D.

\begin{corollary}[11.1.2]
\label{IV.11.1.2-fr}
Soient $Y$ un préschéma localement noethérien, $f:X\to Y$ un morphisme
localement de type fini, $\mathcal F$, $\mathcal F'$ deux
$\mathcal O_X$-Modules cohérents,
$u:\mathcal F'\to\mathcal F$ un homomorphisme de $\mathcal O_X$-Modules. On
suppose que $\mathcal F$ est $f$-plat. Alors l'ensemble $U$ des $x\in X$ tels
que, en posant $y=f(x)$, l'homomorphisme
\[
  u_x\otimes1:\mathcal F'_x\otimes_{\mathcal O_y}k(y)\longrightarrow
  \mathcal F_x\otimes_{\mathcal O_y}k(y)
\]
soit injectif, est ouvert dans $X$.
\end{corollary}

En effet, soit $\mathcal N$ (resp. $\mathcal P$) le noyau (resp. le conoyau) de
$u$~; appliquons ($\mathbf 0_{\mathrm{III}}$, 10.2.4)\footnote{Voici une
démonstration de ($\mathbf 0_{\mathrm{III}}$, 10.2.4), qui n'a pas été publiée
dans l'\emph{Algèbre commutative} de N.~Bourbaki. Compte tenu de
($\mathbf 0_{\mathrm I}$, 6.1.2), il suffit de voir que b) entraîne a). Posons
$P=\operatorname{Im}(u)$, $Q=\operatorname{Coker}(u)$,
$R=\operatorname{Ker}(u)$. Le composé
$M\otimes k\to P\otimes k\to N\otimes k$ est injectif, et
$M\otimes k\to P\otimes k$ est surjectif, donc
$P\otimes k\to N\otimes k$ est injectif et
$M\otimes k\to P\otimes k$ est bijectif. La suite exacte
$0\to P\to N\to Q\to0$ donne la suite exacte
\[
0=\operatorname{Tor}_1^A(N,k)\to\operatorname{Tor}_1^A(Q,k)
\to P\otimes k\to N\otimes k\to Q\otimes k\to0
\]
et puisque $P\otimes k\to N\otimes k$ est injectif, on a
$\operatorname{Tor}_1^A(Q,k)=0$~; comme $Q$ est un $B$-module de type fini,
($\mathbf 0_{\mathrm{III}}$, 10.2.2) montre que $Q$ est un $A$-module plat~;
alors $P$ est aussi un $A$-module plat par ($\mathbf 0_{\mathrm I}$, 6.1.2).
La suite $0\to R\to M\to P\to0$ étant exacte, il en est de même de
$0\to R\otimes k\to M\otimes k\to P\otimes k\to0$ par
($\mathbf 0_{\mathrm I}$, 6.1.2)~; puisque $M\otimes k\to P\otimes k$ est
bijectif, on a $R\otimes k=0$~; mais comme $B$ est noethérien, $R$ est un
$B$-module de type fini, donc on a $R=0$ en vertu du lemme de Nakayama.} aux
anneaux locaux $\mathcal O_y$ et $\mathcal O_x$ et aux
$\mathcal O_x$-modules $\mathcal F'_x$ et $\mathcal F_x$~: dire que
$u_x\otimes1$ est injectif équivaut à dire que $\mathcal N_x=0$ et que
$\mathcal P$ est $f$-plat au point $x$. Or comme $\mathcal N$ et $\mathcal P$
sont cohérents ($\mathbf 0_{\mathrm I}$, 5.3.4), l'ensemble des $x$ où
$\mathcal N_x=0$ est ouvert ($\mathbf 0_{\mathrm I}$, 5.2.2), et l'ensemble
des $x$ où $\mathcal P$ est $f$-plat est ouvert par (11.1.1)~; d'où la
conclusion.

En particulier~:

\begin{corollary}[11.1.3]
\label{IV.11.1.3-fr}
Soient $Y$ un préschéma localement noethérien, $f:X\to Y$ un morphisme plat et
localement de type fini, $g$ une section de $\mathcal O_X$ au-dessus de $X$.
L'ensemble des $x\in X$ tels que $g_x\otimes1$ soit non diviseur de zéro dans
$\mathcal O_x\otimes_{\mathcal O_{f(x)}}k(f(x))$ est ouvert dans $X$.
\end{corollary}

Il suffit d'appliquer (11.1.2) à l'endomorphisme de $\mathcal O_X$ défini par
$g$.

\begin{corollary}[11.1.4]
\label{IV.11.1.4-fr}
Soient $Y$ un préschéma localement noethérien, $f:X\to Y$ un morphisme
localement de type fini, $\mathcal F$ un $\mathcal O_X$-Module cohérent et
$f$-plat. Soit $(g_i)_{1\leq i\leq n}$ une suite de sections de $\mathcal O_X$
au-dessus de $X$. Alors l'ensemble $U$ des $x\in X$ tels que la suite
$((g_i)_x\otimes1)$ soit
$(\mathcal F_x\otimes_{\mathcal O_{f(x)}}k(f(x)))$-régulière est ouvert dans
$X$.
\end{corollary}

\oldpage[IV-3]{119}

Comme $\mathcal F$ est $f$-plat, il résulte de ($\mathbf 0$, 15.1.16) que $U$
est aussi l'ensemble des $x\in X$ tels que la suite $((g_i)_x)$ soit
$\mathcal F_x$-régulière et que le $\mathcal O_x$-module $\mathcal G_x$ (où
$\mathcal G=\mathcal F/(\sum_{i=1}^n g_i\mathcal F)$) soit $f$-plat. Mais
$\mathcal F$ et $\mathcal G$ sont cohérents, donc le corollaire résulte de
(11.1.1) et de ($\mathbf 0$, 15.2.4).

\begin{corollary}[11.1.5]
\label{IV.11.1.5-fr}
Soient $X$, $Y$, $Z$ trois préschémas localement noethériens,
$f:X\to Y$, $g:Y\to Z$ deux morphismes de type fini, $\mathcal F$ un
$\mathcal O_X$-Module cohérent. Alors l'ensemble $U$ des $y\in Y$ tels que,
pour toute générisation $y'$ de $y$, $\mathcal F$ soit $(g\circ f)$-plat en
tous les points de $f^{-1}(y')$ (i.e. tels que
$\mathcal F\otimes_Y\operatorname{Spec}(\mathcal O_{Y,y})$ soit plat
relativement au morphisme
$X\times_Y\operatorname{Spec}(\mathcal O_{Y,y})\to Z$) est ouvert dans $Y$.
\end{corollary}

Si $V$ est l'ensemble des $x\in X$ où $\mathcal F$ est $(g\circ f)$-plat, $U$
est l'ensemble des $y\in Y$ tels que pour toute générisation $y'$ de $y$, on
ait $f^{-1}(y')\subset V$. Or $V$ est ouvert (11.1.1) donc localement
constructible dans $X$, et l'ensemble $W$ des $y\in Y$ tels que
$f^{-1}(y)\subset V$ est égal à $Y-f(X-V)$, donc est aussi localement
constructible dans $Y$ en vertu du théorème de Chevalley (1.8.4). Mais il
résulte alors de ($\mathbf 0_{\mathrm{III}}$, 9.2.5) que les points de $U$ sont
les points intérieurs à $W$, d'où la conclusion.

\begin{corollary}[11.1.6]
\label{IV.11.1.6-fr}
Soient $A$ un anneau noethérien, $B$ une $A$-algèbre de type fini, $C$ une
$B$-algèbre de type fini, $M$ un $C$-module de type fini. Alors l'ensemble des
$\mathfrak q\in\operatorname{Spec}(B)$ tels que $M_{\mathfrak q}$ soit un
$A$-module plat est ouvert dans $\operatorname{Spec}(B)$.
\end{corollary}

Compte tenu de (2.1.2), c'est une conséquence de (11.1.5) appliqué à
$Z=\operatorname{Spec}(A)$, $Y=\operatorname{Spec}(B)$,
$X=\operatorname{Spec}(C)$, $\mathcal F=\widetilde M$.

Les résultats de ce numéro seront débarrassés des hypothèses noethériennes
dans (11.3).

\subsection{Platitude d'une limite projective de préschémas}
\label{subsection:IV.11.2-fr}

\begin{env}[11.2.1]
\label{IV.11.2.1-fr}
Soient $A$ un anneau, $M$, $N$ deux $A$-modules, $A'$ une $A$-algèbre~;
posons $M'=M\otimes_AA'$, $N'=N\otimes_AA'$. Rappelons
($\mathbf{III}$, 6.3.8) que l'on définit, pour tout $i$, un homomorphisme
canonique de $A$-modules
\begin{equation}
\psi_i:\operatorname{Tor}_i^A(M,N)\to\operatorname{Tor}_i^{A'}(M',N')
\label{IV.11.2.1.1-fr}
\tag{11.2.1.1}
\end{equation}
de la façon suivante~: on considère une résolution gauche de $M$ par des
$A$-modules libres
\begin{equation}
\cdots\longrightarrow L_{i+1}\xrightarrow{f_i}L_i\longrightarrow\cdots
\longrightarrow L_0\xrightarrow{\varepsilon}M\longrightarrow0
\label{IV.11.2.1.2-fr}
\tag{11.2.1.2}
\end{equation}
d'où l'on déduit par tensorisation avec $A'$ un complexe de $A'$-modules
\begin{equation}
\cdots\longrightarrow L'_{i+1}\xrightarrow{f'_i}L'_i\longrightarrow\cdots
\longrightarrow L'_0\xrightarrow{\varepsilon'}M'\longrightarrow0
\label{IV.11.2.1.3-fr}
\tag{11.2.1.3}
\end{equation}
où l'on a posé $L'_i=L_i\otimes_AA'$, $f'_i=f_i\otimes1_{A'}$,
$\varepsilon'=\varepsilon\otimes1_{A'}$. Considérons d'autre part une
résolution gauche de $M'$ par des $A'$-modules libres
\begin{equation}
\cdots\longrightarrow L''_{i+1}\xrightarrow{f''_i}L''_i\longrightarrow\cdots
\longrightarrow L''_0\xrightarrow{\varepsilon''}M'\longrightarrow0.
\label{IV.11.2.1.4-fr}
\tag{11.2.1.4}
\end{equation}
\end{env}

\oldpage[IV-3]{120}

Comme les $L'_i$ sont des $A'$-modules libres, on sait (M, V, 1.1) qu'il y a
des $A'$-homomorphismes $u_i$ formant un diagramme commutatif
\[
\xymatrix{
\cdots\ar[r]&L'_{i+1}\ar[r]^{f'_i}\ar[d]_{u_{i+1}}&
L'_i\ar[r]\ar[d]_{u_i}&\cdots\ar[r]&
L'_0\ar[r]^{\varepsilon'}\ar[d]_{u_0}&M'\ar[r]\ar[d]_{1_{M'}}&0\\
\cdots\ar[r]&L''_{i+1}\ar[r]_{f''_i}&L''_i\ar[r]&\cdots\ar[r]&
L''_0\ar[r]_{\varepsilon''}&M'\ar[r]&0.
}
\]
Si l'on compose l'homomorphisme $u_\bullet:L'_\bullet\to L''_\bullet$ de
complexes ainsi défini avec l'homomorphisme canonique
$L_\bullet\to L'_\bullet$, on obtient un homomorphisme de complexes de
$A$-modules $L_\bullet\to L''_\bullet$~; en remarquant que l'on a
$(L_i\otimes_AN)\otimes_AA'=L'_i\otimes_{A'}N'$, on en déduit un
homomorphisme de complexes de $A$-modules
$L_\bullet\otimes_AN\to L''_\bullet\otimes_{A'}N'$, d'où, en passant à
l'homologie, les homomorphismes canoniques (11.2.1.1). Comme les $u_i$ sont
bien déterminés à homotopie près (M, V, 1.1), les homomorphismes (11.2.1.1)
ne dépendent pas du choix des $u_i$ ni du choix des résolutions libres
$L_\bullet$ et $L''_\bullet$.

\begin{sloppypar}
Comme les $\operatorname{Tor}_i^{A'}(M',N')$ sont des $A'$-modules, on déduit
canoniquement de (11.2.1.1) des $A'$-homomorphismes
\begin{equation}
\begin{gathered}
\varphi_i:\operatorname{Tor}_i^A(M,N)\otimes_AA'\\[-2pt]
\longrightarrow\operatorname{Tor}_i^{A'}(M',N').
\end{gathered}
\label{IV.11.2.1.5-fr}
\tag{11.2.1.5}
\end{equation}
\end{sloppypar}

Considérons maintenant deux homomorphismes d'anneaux
\[
\rho:A\to A^{(1)},\qquad \sigma:A^{(1)}\to A^{(2)},
\]
et leur composé $\sigma\circ\rho:A\to A^{(2)}$~; posons
$M^{(j)}=M\otimes_AA^{(j)}$, $N^{(j)}=N\otimes_AA^{(j)}$ pour $j=1,2$.
Alors l'homomorphisme canonique composé
\[
\operatorname{Tor}_i^A(M,N)\to
\operatorname{Tor}_i^{A^{(1)}}(M^{(1)},N^{(1)})\to
\operatorname{Tor}_i^{A^{(2)}}(M^{(2)},N^{(2)})
\]
est le même que l'homomorphisme canonique déduit de $\sigma\circ\rho$~;
cela résulte de ce que, si $L_\bullet^{(j)}$ est une résolution libre de
$M^{(j)}$, le diagramme
\[
\xymatrix{
L_\bullet\ar[r]&L_\bullet\otimes_AA^{(1)}\ar[r]\ar[d]&
L_\bullet^{(1)}\ar[d]\\
&L_\bullet\otimes_AA^{(2)}\ar[r]&
L_\bullet^{(1)}\otimes_{A^{(1)}}A^{(2)}\ar[r]&L_\bullet^{(2)}
}
\]
est commutatif.

\begin{env}[11.2.2]
\label{IV.11.2.2-fr}
Les notations étant celles de (11.2.1), considérons maintenant un
\emph{système inductif filtrant} de $A$-algèbres
$(A_\alpha,g_{\beta\alpha})$, et pour tout indice $\alpha$, posons
$M_\alpha=M\otimes_AA_\alpha$, $N_\alpha=N\otimes_AA_\alpha$~; il résulte
alors de (11.2.1) que pour chaque $i$,
$(\operatorname{Tor}_i^{A_\alpha}(M_\alpha,N_\alpha),
\psi_{i\beta\alpha})$, où $\psi_{i\beta\alpha}$ est l'homomorphisme
canonique (11.2.1.1) correspondant à $g_{\beta\alpha}$, est un \emph{système
inductif} de $A_\alpha$-modules. Posons
$A'=\varinjlim A_\alpha$, $M'=\varinjlim M_\alpha=M\otimes_AA'$,
$N'=\varinjlim N_\alpha=N\otimes_AA'$~; si on désigne par
$g_\alpha:A_\alpha\to A'$ l'homomorphisme canonique, on en déduit des
homomorphismes canoniques (11.2.1.1)
$\psi_{i\alpha}:\operatorname{Tor}_i^{A_\alpha}(M_\alpha,N_\alpha)
\to\operatorname{Tor}_i^{A'}(M',N')$ qui (toujours en vertu de (11.2.1))
forment un système inductif d'homomorphismes~; nous nous proposons de
compléter le résultat de (M, V, 9.4$^*$) en montrant que les
\begin{equation}
\psi_i=\varinjlim_\alpha\psi_{i\alpha}:
\varinjlim_\alpha\operatorname{Tor}_i^{A_\alpha}(M_\alpha,N_\alpha)
\to\operatorname{Tor}_i^{A'}(M',N')
\label{IV.11.2.2.1-fr}
\tag{11.2.2.1}
\end{equation}
sont des \emph{isomorphismes de $A'$-modules}. Pour cela, nous procédons
comme dans (M, V, 9.5$^*$), en associant à chaque $M_\alpha$ sa
\emph{résolution libre canonique}. Tout revient (compte tenu de l'exactitude
du foncteur $\varinjlim$) à prouver le
\end{env}

\oldpage[IV-3]{121}

\begin{lemma}[11.2.2.2]
\label{IV.11.2.2.2-fr}
Soient $(A_\alpha,g_{\beta\alpha})$ un système inductif filtrant d'anneaux,
$(M_\alpha,h_{\beta\alpha})$ un système inductif d'ensembles,
$A'=\varinjlim A_\alpha$, $M'=\varinjlim M_\alpha$,
$g_\alpha:A_\alpha\to A'$, $h_\alpha:M_\alpha\to M'$ les applications
canoniques. Pour tout $\alpha$, soit $F(M_\alpha)$ le $A_\alpha$-module des
combinaisons linéaires formelles d'éléments de $M_\alpha$~; soit de même
$F(M')$ le $A'$-module des combinaisons linéaires formelles d'éléments de
$M'$~; si $h^0_{\beta\alpha}:F(M_\alpha)\to F(M_\beta)$ (pour
$\alpha\leq\beta$) et $h'_\alpha:F(M_\alpha)\to F(M')$ sont les
$A_\alpha$-homomorphismes déduits de $h_{\beta\alpha}$ et $h_\alpha$
respectivement, $(F(M_\alpha),h^0_{\beta\alpha})$ est un système inductif
de $A_\alpha$-modules et $(h'_\alpha)$ un système inductif
d'homomorphismes. Alors
\[
h''=\varinjlim h'_\alpha:\varinjlim F(M_\alpha)\to
F(M')=F(\varinjlim M_\alpha)
\]
est un isomorphisme.
\end{lemma}

Pour la démonstration, voir Bourbaki, \emph{Alg.}, chap.~II, 3$^e$ éd.,
\S~6, n$^\circ$~6, cor. de la prop.~10.

\begin{env}[11.2.3]
\label{IV.11.2.3-fr}
Reprenons les notations de (11.2.1) et considérons particulièrement le cas
$i=1$~; posons $T=\operatorname{Tor}_1^A(M,N)$, $T'=T\otimes_AA'$,
$T''=\operatorname{Tor}_1^{A'}(M',N')$~; alors $\psi_1$ est
l'homomorphisme
\[
\operatorname{Ker}(f_0\otimes1_N)/\operatorname{Im}(f_1\otimes1_N)
\to
\operatorname{Ker}(f''_0\otimes1_{N'})/
\operatorname{Im}(f''_1\otimes1_{N'})
\]
qui se déduit par passage aux quotients de la restriction
$\operatorname{Ker}(f_0\otimes1_N)\to
\operatorname{Ker}(f''_0\otimes1_{N'})$ de
\[
u_1\otimes1:L_0\otimes_AN\to L''_0\otimes_{A'}N'=L''_0\otimes_AN.
\]
Posons $R=\operatorname{Ker}(\varepsilon)$,
$R''=\operatorname{Ker}(\varepsilon'')$, de sorte que l'on a les suites
exactes
\[
0\to R\xrightarrow{j}L_0\xrightarrow{\varepsilon}M\to0
\qquad\text{et}\qquad
0\to R''\xrightarrow{j''}L''_0\xrightarrow{\varepsilon''}M'\to0,
\]
d'où on déduit les suites exactes d'homologie
\begin{equation}
0=\operatorname{Tor}_1^A(L_0,N)\to T\xrightarrow{\partial}R\otimes_AN
\xrightarrow{j\otimes1}L_0\otimes_AN\xrightarrow{\varepsilon\otimes1}
M\otimes_AN\to0
\label{IV.11.2.3.1-fr}
\tag{11.2.3.1}
\end{equation}
\begin{equation}
0=\operatorname{Tor}_1^{A'}(L''_0,N')\to T''\xrightarrow{\partial''}
R''\otimes_{A'}N'\xrightarrow{j''\otimes1}L''_0\otimes_{A'}N'
\xrightarrow{\varepsilon''\otimes1}M'\otimes_{A'}N'\to0.
\label{IV.11.2.3.2-fr}
\tag{11.2.3.2}
\end{equation}
On a d'autre part un homomorphisme de $A'$-modules
\[
v:R'=\operatorname{Ker}(\varepsilon)\otimes_AA'
\to\operatorname{Ker}(\varepsilon')\xrightarrow{u_0}
\operatorname{Ker}(\varepsilon'')=R''.
\]
\end{env}

\oldpage[IV-3]{122}

Montrons que le diagramme
\begin{equation}
\xymatrix{
&T'\ar[r]^{\partial\otimes1}\ar[d]_{\varphi_1}&
R'\otimes_{A'}N'\ar[r]\ar[d]_{v\otimes1}&
L'_0\otimes_{A'}N'\ar[r]\ar[d]_{u_0\otimes1}&
M'\otimes_{A'}N'\ar[r]\ar@{=}[d]&0\\
0\ar[r]&T''\ar[r]_{\partial''}&R''\otimes_{A'}N'\ar[r]&
L''_0\otimes_{A'}N'\ar[r]&M'\otimes_{A'}N'\ar[r]&0
}
\label{IV.11.2.3.3-fr}
\tag{11.2.3.3}
\end{equation}
est \emph{commutatif}. Pour cela, on vérifie aussitôt (M, IV, 1) que
l'homomorphisme $\partial:T\to R\otimes_AN$ provient (dans le cas actuel)
par passage au quotient de l'homomorphisme
$w:\operatorname{Ker}(f_0\otimes1_N)\to R\otimes_AN$, restriction de
l'homomorphisme $g_0\otimes1_N:L_1\otimes_AN\to R\otimes_AN$, où
$g_0:L_1\to R$ est tel que $f_0=j\circ g_0$~; de même pour $\partial''$.
Il suffit donc de voir que le diagramme
\[
\xymatrix{
\operatorname{Ker}(f_0\otimes1_N)\ar[r]^w\ar[d]&R\otimes_AN\ar[r]&
R'\otimes_{A'}N'=(R\otimes_AN)\otimes_AA'\ar[d]_{v\otimes1}\\
\operatorname{Ker}(f''_0\otimes1_{N'})\ar[rr]&&R''\otimes_{A'}N'
}
\]
est commutatif, ce qui résulte de la commutativité du diagramme
\[
\xymatrix{
L_1\ar[r]\ar[d]&R\ar[d]\\
L''_1\ar[r]&R''.
}
\]

\begin{lemma}[11.2.4]
\label{IV.11.2.4-fr}
Soient $A$ un anneau, $C$ une $A$-algèbre, $M$ un $A$-module, $N$ un
$C$-module, $A'$ une $A$-algèbre. Posons $C'=C\otimes_AA'$,
$M'=M\otimes_AA'$, $N'=N\otimes_AA'=N\otimes_CC'$. Supposons que
$M\otimes_AN$ soit un $C$-module plat. Alors l'homomorphisme canonique
\[
\varphi_1:\operatorname{Tor}_1^A(M,N)\otimes_AA'
\to\operatorname{Tor}_1^{A'}(M',N')
\]
(cf. (11.2.1.5)) est surjectif.
\end{lemma}

Gardons les notations de (11.2.3)~; l'exactitude à droite du produit
tensoriel montre que la suite
$L'_1\xrightarrow{f'_0}L'_0\xrightarrow{\varepsilon'}M'\to0$ est exacte~;
comme $L'_0$ et $L'_1$ sont des $A'$-modules libres, \emph{on peut supposer
que l'on a pris} $L'_0=L''_0$, $L'_1=L''_1$, $u_0$ et $u_1$ étant les
\emph{applications identiques} et $f''_0=f'_0$. Comme
$R=\operatorname{Im}(f_0)$ et $R''=\operatorname{Im}(f''_0)=
\operatorname{Im}(f'_0)$, l'homomorphisme $v$ est surjectif, et il en est
donc de même de $v\otimes1$. La démonstration sera achevée si l'on prouve que
la première ligne de (11.2.3.3) est \emph{exacte}, $v\otimes1$ étant
surjective et $u_0\otimes1$ injectif (Bourbaki, \emph{Alg. comm.},
chap.~I$^{\mathrm{er}}$, \S~1, n$^\circ$~4, cor.~2 de la prop.~2).
Utilisons pour cela

\oldpage[IV-3]{123}

l'hypothèse que $M\otimes_AN$ est un $C$-module plat. Posant
$Q=\operatorname{Ker}(\varepsilon\otimes1)=\operatorname{Im}(j\otimes1)$
dans la suite exacte (11.2.3.1), on a les deux suites exactes
$0\to Q\to L_0\otimes_AN\to M\otimes_AN\to0$ et
$T\to R\otimes_AN\to Q\to0$, où les homomorphismes sont des homomorphismes
de $C$-modules~; utilisant l'hypothèse de platitude, on en déduit
($\mathbf 0_{\mathrm I}$, 6.1.2) la suite exacte
\[
0\to Q\otimes_CC'\to(L_0\otimes_AN)\otimes_CC'
\to(M\otimes_AN)\otimes_CC'\to0
\]
et d'autre part, le produit tensoriel étant exact à droite, on a la suite
exacte
\[
T\otimes_CC'\to(R\otimes_AN)\otimes_CC'\to Q\otimes_CC'\to0,
\]
d'où finalement la suite exacte
\[
T\otimes_CC'\to(R\otimes_AN)\otimes_CC'\to
(L_0\otimes_AN)\otimes_CC'\to(M\otimes_AN)\otimes_CC'\to0.
\]
Mais par définition, pour tout $C$-module $P$,
$P\otimes_CC'=P\otimes_AA'$, d'où la conclusion.

\begin{lemma}[11.2.5]
\label{IV.11.2.5-fr}
Soient $A$ un anneau, $\mathfrak J$ un idéal de $A$, $B$ une $A$-algèbre,
$M$ un $B$-module, $A\to A'$ un homomorphisme d'anneaux. On pose
$\mathfrak J'=\mathfrak JA'$, $B'=B\otimes_AA'$,
$M'=M\otimes_AA'=M\otimes_BB'$. Soit $\mathfrak p'$ un idéal premier de
$B'$ contenant $\mathfrak J'B'$. On suppose vérifiée l'une des hypothèses
suivantes~:
\begin{enumerate}[label=\emph{\alph*)}]
\item $\mathfrak J$ est nilpotent.
\item $A'$ est noethérien, $B'$ est une $A'$-algèbre de type fini, $M'$ un
$B'$-module de type fini.
\end{enumerate}
Dans ces conditions, supposons vérifiées les deux propriétés suivantes~:
\begin{enumerate}[label=\emph{(\roman*)}]
\item $M\otimes_A(A/\mathfrak J)$ est un $(A/\mathfrak J)$-module plat.
\item L'homomorphisme canonique composé
\[
\operatorname{Tor}_1^A(M,A/\mathfrak J)\xrightarrow{\psi_1}
\operatorname{Tor}_1^{A'}(M',A'/\mathfrak J')\xrightarrow{\theta}
(\operatorname{Tor}_1^{A'}(M',A'/\mathfrak J'))_{\mathfrak p'}
\]
(où $\psi_1$ est l'homomorphisme (11.2.1.1) et $\theta$ l'homomorphisme
canonique d'un $B'$-module dans son localisé en $\mathfrak p'$) est nul.
\end{enumerate}
Alors $M'_{\mathfrak p'}$ est un $A'$-module plat.
\end{lemma}

Notons que dans l'hypothèse b) $M'_{\mathfrak p'}$ est un
$B'_{\mathfrak p'}$-module de type fini, $B'_{\mathfrak p'}$ une
$A'$-algèbre noethérienne, et que $\mathfrak J'B'_{\mathfrak p'}$ est
contenu dans le radical $\mathfrak p'B'_{\mathfrak p'}$ de
$B'_{\mathfrak p'}$~; dans l'hypothèse a), $\mathfrak J'$ est nilpotent~;
on va donc pouvoir appliquer le critère de platitude
($\mathbf 0_{\mathrm{III}}$, 10.2.1) ou ($\mathbf 0_{\mathrm{III}}$, 10.2.2)
suivant que a) ou b) est vérifiée. En premier lieu, on a
\[
M'_{\mathfrak p'}\otimes_{A'}(A'/\mathfrak J')
=(M'\otimes_{A'}(A'/\mathfrak J'))_{\mathfrak p'}
=((M\otimes_A(A/\mathfrak J))\otimes_{A/\mathfrak J}
(A'/\mathfrak J'))_{\mathfrak p'}~;
\]
l'hypothèse (i) entraîne donc que
$M'_{\mathfrak p'}\otimes_{A'}(A'/\mathfrak J')$ est un
$(A'/\mathfrak J')$-module plat, compte tenu de
($\mathbf 0_{\mathrm I}$, 6.2.1 et 6.3.2). Reste donc à voir que
$\operatorname{Tor}_1^{A'}(M'_{\mathfrak p'},A'/\mathfrak J')=0$~; or ce
$B'_{\mathfrak p'}$-module est égal à
$(\operatorname{Tor}_1^{A'}(M',A'/\mathfrak J'))_{\mathfrak p'}$, en vertu
de la platitude de $B'_{\mathfrak p'}$ sur $B'$. Mais en vertu de
l'hypothèse (ii), l'homomorphisme composé
$\theta\circ\varphi_1:\operatorname{Tor}_1^A(M,A/\mathfrak J)\otimes_AA'
\to(\operatorname{Tor}_1^{A'}(M',A'/\mathfrak J'))_{\mathfrak p'}$ est
nul~; par ailleurs, (11.2.4) appliqué à $C=A/\mathfrak J$ et $N=C$ montre
(compte tenu de l'hypothèse (i)) que $\varphi_1$ est \emph{surjectif} (car
$A'/\mathfrak J'=C\otimes_AA'$)~; donc l'homomorphisme $\theta$ est
\emph{nul} et comme l'image par $\theta$ de
$\operatorname{Tor}_1^{A'}(M',A'/\mathfrak J')$ engendre le
$B'_{\mathfrak p'}$-module
$(\operatorname{Tor}_1^{A'}(M',A'/\mathfrak J'))_{\mathfrak p'}$, ce
dernier est nul. C.Q.F.D.

\begin{theorem}[11.2.6]
\label{IV.11.2.6-fr}
Les notations étant celles de (8.5.1) et (8.8.1), on suppose $S_\alpha$
quasi-compact, $X_\alpha$ et $Y_\alpha$ de présentation finie sur
$S_\alpha$~; soient $f_\alpha:X_\alpha\to Y_\alpha$ un
$S_\alpha$-morphisme, $\mathcal F_\alpha$ un
$\mathcal O_{X_\alpha}$-Module quasi-cohérent de présentation finie.

\oldpage[IV-3]{124}

\begin{enumerate}[label=\emph{(\roman*)}]
\item Soient $x$ un point de $X$, $x_\lambda$ sa projection canonique dans
$X_\lambda$. Pour que $\mathcal F$ soit $f$-plat au point $x$, il faut et il
suffit qu'il existe $\lambda\geq\alpha$ tel que $\mathcal F_\lambda$ soit
$f_\lambda$-plat au point $x_\lambda$.
\item Pour que $\mathcal F$ soit $f$-plat, il faut et il suffit qu'il existe
$\lambda\geq\alpha$ tel que $\mathcal F_\lambda$ soit $f_\lambda$-plat.
\end{enumerate}
\end{theorem}

On peut supposer que $S_\alpha=S_0$~; comme $Y_0$ est de présentation finie
sur $S_0$, $Y_0$ est quasi-compact, et $f_0:X_0\to Y_0$ est un morphisme de
présentation finie (1.6.2, (v)), donc on peut aussi se borner au cas où
$S_0=Y_0$. En outre, en vertu de la quasi-compacité de $S_0$ et du fait que
l'ensemble d'indices $L$ est filtrant, on peut se borner au cas où
$S_0=\operatorname{Spec}(A_0)$ est affine. De plus $X_0$ est quasi-compact,
donc le même raisonnement montre qu'on peut aussi supposer
$X_0=\operatorname{Spec}(B_0)$ affine~; on a alors
$\mathcal F_0=\widetilde M_0$, où $M_0$ est un $B_0$-module de présentation
finie, et l'énoncé (11.2.6) est dans ce cas équivalent au suivant (compte tenu
de ($\mathbf 0_{\mathrm I}$, 6.3.1))~:

\begin{corollary}[11.2.6.1]
\label{IV.11.2.6.1-fr}
Soient $A_0$ un anneau, $(A_\lambda)_{\lambda\in L}$ un système inductif
filtrant de $A_0$-algèbres, $B_0$ une $A_0$-algèbre de présentation finie,
$M_0$ un $B_0$-module de présentation finie. On pose
\[
\begin{gathered}
B_\lambda=B_0\otimes_{A_0}A_\lambda,\qquad
M_\lambda=M_0\otimes_{A_0}A_\lambda=M_0\otimes_{B_0}B_\lambda,\\
A=\varinjlim A_\lambda,\qquad
B=\varinjlim B_\lambda=B_0\otimes_{A_0}A,\\
M=\varinjlim M_\lambda=M_0\otimes_{A_0}A=M_0\otimes_{B_0}B.
\end{gathered}
\]
\begin{enumerate}[label=\emph{(\roman*)}]
\item Soit $\mathfrak p$ un idéal premier de $B$, et pour tout $\lambda$,
soit $\mathfrak p_\lambda$ son image réciproque dans $B_\lambda$. Pour que
$M_{\mathfrak p}$ soit un $A$-module plat, il faut et il suffit qu'il existe
$\lambda$ tel que $(M_\lambda)_{\mathfrak p_\lambda}$ soit un
$A_\lambda$-module plat.
\item Pour que $M$ soit un $A$-module plat, il faut et il suffit qu'il existe
$\lambda$ tel que $M_\lambda$ soit un $A_\lambda$-module plat.
\end{enumerate}
\end{corollary}

On doit seulement démontrer que les conditions sont nécessaires (2.1.4).
Nous procéderons en plusieurs étapes.

\medskip
\noindent\textbf{I)} \emph{Réduction au cas où les $A_\lambda$ sont
noethériens.} En vertu de (8.9.1), il existe un sous-anneau $A'_0$ de $A_0$,
qui est une $\mathbf Z$-algèbre de type fini, une $A'_0$-algèbre de type fini
$B'_0$ et un $B'_0$-module de type fini $M'_0$ tels que l'on ait
$B_0=B'_0\otimes_{A'_0}A_0$ et $M_0=M'_0\otimes_{A'_0}A_0$~; comme on a
$B_\lambda=B'_0\otimes_{A'_0}A_\lambda$ et
$M_\lambda=M'_0\otimes_{A'_0}A_\lambda$, on peut remplacer $A_0$, $B_0$,
$M_0$ par $A'_0$, $B'_0$, $M'_0$ dans l'énoncé de (11.2.6.1), en considérant
les $A_\lambda$ comme des $A'_0$-algèbres~; on peut donc déjà supposer que
$A_0$ est noethérien. Soit $H$ l'ensemble des couples
$(\lambda,C_\lambda)$, où $C_\lambda$ est une sous-$A_0$-algèbre de
\emph{type fini} de $A_\lambda$~; ordonnons $H$ en posant
$(\lambda,C_\lambda)\leq(\mu,D_\mu)$ si $\lambda\leq\mu$ et si
l'homomorphisme $\varphi_{\mu\lambda}:A_\lambda\to A_\mu$ est tel que
$\varphi_{\mu\lambda}(C_\lambda)\subset D_\mu$~; pour cet ordre $H$ est
filtrant croissant, car si $(\lambda,C_\lambda)$ et $(\mu,D_\mu)$ sont deux
éléments quelconques de $H$, on les majorera par un $(\nu,E_\nu)$ en prenant
$\nu\geq\lambda$, $\nu\geq\mu$ dans $L$, puis $E_\nu$ égal à la
sous-$A_0$-algèbre de $A_\nu$ engendrée par
$\varphi_{\nu\lambda}(C_\lambda)$ et $\varphi_{\nu\mu}(D_\mu)$. Pour un
élément $\xi=(\lambda,C_\lambda)$ de $H$, on posera $A_\xi=C_\lambda$, et
pour $\xi\leq\eta=(\mu,D_\mu)$ (donc $\lambda\leq\mu$ et
$\varphi_{\mu\lambda}(C_\lambda)\subset D_\mu$),
$\varphi_{\eta\xi}:A_\xi\to A_\eta$ sera la restriction à $C_\lambda$ de
$\varphi_{\mu\lambda}$, considéré comme homomorphisme dans $D_\mu$~; il est
clair que l'on obtient ainsi un système inductif filtrant de $A_0$-algèbres.
On pose $B_\xi=B_0\otimes_{A_0}A_\xi$,
$M_\xi=M_0\otimes_{A_0}A_\xi$~; cette fois les $A_\xi$ sont
noethériens~; en outre la formule de la double limite inductive (Bourbaki,
\emph{Alg.}, chap.~II, 3$^e$ éd., \S~6, n$^\circ$~4, prop.~7) prouve que
l'on a encore
$\varinjlim_HA_\xi=A$, $\varinjlim_HB_\xi=B$,
$\varinjlim_HM_\xi=M$. Supposons (11.2.6.1) prouvé pour le système inductif
$(A_\xi)_{\xi\in H}$~; soit $\mathfrak p$ un idéal premier de $B$,

\oldpage[IV-3]{125}

tel que $M_{\mathfrak p}$ soit un $A$-module plat~; il existe alors
$\xi\in H$ tel que, si $\mathfrak p_\xi$ est l'image réciproque de
$\mathfrak p$ dans $B_\xi$, $(M_\xi)_{\mathfrak p_\xi}$ soit un
$A_\xi$-module plat. Soit $\xi=(\lambda,C_\lambda)$, de sorte que
l'injection $A_\xi=C_\lambda\to A_\lambda$ donne un homomorphisme
$B_\xi=B_0\otimes_{A_0}C_\lambda\to B_\lambda=B_0\otimes_{A_0}A_\lambda$,
et si $\mathfrak p_\lambda$ est l'image réciproque de $\mathfrak p$ dans
$B_\lambda$, $\mathfrak p_\xi$ est l'image réciproque de
$\mathfrak p_\lambda$ dans $B_\xi$~; on a par suite
\[
(M_\lambda)_{\mathfrak p_\lambda}
=(M_\xi\otimes_{C_\lambda}A_\lambda)_{\mathfrak p_\lambda}
=((M_\xi)_{\mathfrak p_\xi}\otimes_{C_\lambda}A_\lambda)_{\mathfrak p_\lambda},
\]
donc $(M_\lambda)_{\mathfrak p_\lambda}$ est un $A_\lambda$-module plat
($\mathbf 0_{\mathrm I}$, 6.2.1 et 6.3.2). On traite de même le cas (ii) de
l'énoncé. Nous pouvons donc par la suite supposer les $A_\lambda$
noethériens pour $\lambda\in L$ (mais non nécessairement $A$ lui-même).

\medskip
\noindent\textbf{II)} \emph{Réduction de l'énoncé global (ii) à l'énoncé
ponctuel (i).} Supposons que $\mathcal F$ soit $f$-plat. Pour tout $\lambda$,
soit $U_\lambda$ l'ensemble des $x_\lambda\in X_\lambda$ tels que
$\mathcal F_\lambda$ soit $f_\lambda$-plat au point $x_\lambda$~; on sait
que $U_\lambda$ est \emph{ouvert} dans $X_\lambda$ puisque $S_\lambda$ est
noethérien et $f_\lambda$ de type fini (11.1.1)~; soit
$V_\lambda=v_\lambda^{-1}(U_\lambda)$ son image réciproque dans $X$. Comme,
par hypothèse, pour tout $x\in X$, il y a un $\lambda$ tel que
$\mathcal F_\lambda$ soit $f_\lambda$-plat au point $x_\lambda$, projection
de $x$ dans $X_\lambda$, cela signifie que $x\in V_\lambda$ pour un
$\lambda$~; autrement dit, $X$ est réunion des $V_\lambda$. D'ailleurs
(2.1.4), pour $\lambda\leq\mu$, on a $V_\lambda\subset V_\mu$, donc, comme
$X$ est quasi-compact, il existe un indice $\mu$ tel que $X=V_\mu$. Comme
les $X_\lambda$ sont quasi-compacts, il résulte de (8.3.4) qu'il existe un
indice $\nu\geq\mu$ tel que $v_{\mu\nu}^{-1}(U_\mu)=X_\nu$~; mais par
(2.1.4), cela entraîne que $\mathcal F_\nu$ est $f_\nu$-plat.

\medskip
\noindent\textbf{III)} \emph{Fin de la démonstration.} Il reste à prouver
(i) en supposant $S_0$ affine et noethérien~; si $y_0$ est la projection de
$x$ dans $S_0$, on peut en outre, en vertu de (2.1.4) et de
($\mathbf I$, 3.6.5), remplacer $S_0$ par
$\operatorname{Spec}(\mathcal O_{y_0})$, autrement dit on peut se borner au
cas où $A_0$ est un anneau \emph{local} noethérien, d'idéal maximal
$\mathfrak m_0=\mathfrak j_{y_0}$~; par définition, $\mathfrak m_0$ est
l'image réciproque de l'idéal premier $\mathfrak p=\mathfrak j_x$ de $B$,
et $M_{\mathfrak p}$ est supposé être un $A$-module plat~; on a donc en
particulier $\operatorname{Tor}_1^A(A/\mathfrak m_0A,M_{\mathfrak p})=0$,
ce qui s'écrit aussi, puisque les $\operatorname{Tor}_1^A(A/\mathfrak m_0A,M)$
sont des $B$-modules et que $B_{\mathfrak p}$ est plat sur $B$,
$(\operatorname{Tor}_1^A(A/\mathfrak m_0A,M))_{\mathfrak p}=0$.
Notons maintenant que $\operatorname{Tor}_1^{A_0}(A_0/\mathfrak m_0,M_0)$
est un $B_0$-module \emph{de type fini}, car on peut le définir en prenant
une résolution de $A_0/\mathfrak m_0$ par des $A_0$-modules libres
\emph{de type fini} ($A_0$ étant noethérien) et en tensorisant par $M_0$, ce
qui donne des $B_0$-modules de type fini~; comme $B_0$ est noethérien,
l'homologie du complexe ainsi obtenu est bien formée de $B_0$-modules de
type fini. Soit $(t_i^0)_{1\leq i\leq m}$ un système de générateurs du
$B_0$-module $\operatorname{Tor}_1^{A_0}(A_0/\mathfrak m_0,M_0)$ et soit
$t_i$ l'image canonique (11.2.1.1) de $t_i^0$ dans
$\operatorname{Tor}_1^A(A/\mathfrak m_0A,M)$. L'hypothèse entraîne qu'il
existe un $h\in B-\mathfrak p$ tel que $ht_i=0$ pour $1\leq i\leq m$. Or,
on a (11.2.2.1)
\[
\operatorname{Tor}_1^A(A/\mathfrak m_0A,M)=
\varinjlim\operatorname{Tor}_1^{A_\lambda}
(A_\lambda/\mathfrak m_0A_\lambda,M_\lambda)~;
\]
il existe donc un $\lambda$ tel que, si les $t_i^\lambda$ sont les images
des $t_i^0$ dans
$\operatorname{Tor}_1^{A_\lambda}(A_\lambda/\mathfrak m_0A_\lambda,
M_\lambda)$, il existe $h_\lambda\in B_\lambda$ d'image $h$, tel que
$h_\lambda t_i^\lambda=0$ pour $1\leq i\leq m$. Soit
$\mathfrak p_\lambda=\mathfrak j_{x_\lambda}$ l'idéal premier de
$B_\lambda$ image réciproque de $\mathfrak p$~; on a
$h_\lambda\in B_\lambda-\mathfrak p_\lambda$, donc les images canoniques
des $t_i^0$ dans
$(\operatorname{Tor}_1^{A_\lambda}(A_\lambda/\mathfrak m_0A_\lambda,
M_\lambda))_{\mathfrak p_\lambda}$ sont nulles et par suite
l'homomorphisme
$\operatorname{Tor}_1^{A_0}(A_0/\mathfrak m_0,M_0)\to
(\operatorname{Tor}_1^{A_\lambda}(A_\lambda/\mathfrak m_0A_\lambda,
M_\lambda))_{\mathfrak p_\lambda}$ est nul. Les conditions du lemme
(11.2.5) sont donc remplies ($A_0/\mathfrak m_0$ étant un corps), et
$(M_\lambda)_{\mathfrak p_\lambda}$ est un $A_\lambda$-module plat, ce qui
achève de démontrer le théorème.

\begin{corollary}[11.2.7]
\label{IV.11.2.7-fr}
Soient $S=\operatorname{Spec}(A)$ un schéma affine, $f:X\to S$ un morphisme,
$\mathcal F$ un $\mathcal O_X$-Module quasi-cohérent, $x$ un point de $X$.
Les conditions suivantes sont équivalentes~:
\begin{enumerate}[label=\emph{\alph*)}]
\item $f$ est un morphisme de présentation finie, $\mathcal F$ est un
$\mathcal O_X$-Module de présentation finie et $\mathcal F$ est $f$-plat au
point $x$ (resp. $f$-plat).

\oldpage[IV-3]{126}

\item Il existe un schéma affine noethérien $S_0=\operatorname{Spec}(A_0)$,
un morphisme de type fini $f_0:X_0\to S_0$, un
$\mathcal O_{X_0}$-Module cohérent $\mathcal F_0$, un morphisme $S\to S_0$
tels que le $S$-préschéma $X_0\otimes_{S_0}S$ soit $S$-isomorphe à $X$ et
que, si on identifie $X$ à $X_0\otimes_{S_0}S$, $\mathcal F$ soit isomorphe
à $\mathcal F_0\otimes_{\mathcal O_{S_0}}\mathcal O_S$, et $\mathcal F_0$
soit $f_0$-plat au point $x_0$ projection de $x$ dans $X_0$ (resp.
$f_0$-plat).
\item Les conditions de b) sont vérifiées et en outre $A_0$ est une
sous-$\mathbf Z$-algèbre de type fini de $A$, le morphisme $S\to S_0$
correspondant à l'injection canonique $A_0\to A$.
\end{enumerate}
\end{corollary}

Il est clair que c) implique b) et b) implique a) en vertu de (2.1.4).
D'autre part, on peut considérer $A$ comme limite inductive de ses
sous-$\mathbf Z$-algèbres de type fini, et on sait par (8.9.1) qu'il y a une
telle sous-algèbre $A_0$ et un morphisme de type fini $f_0:X_0\to S_0$ tels
que $X$ soit $S$-isomorphe à $X_0\times_{S_0}S$ et $\mathcal F$ isomorphe à
$\mathcal F_0\otimes_{\mathcal O_{S_0}}\mathcal O_S$~; on peut donc écrire
$X=\varprojlim X_\lambda$, où $X_\lambda=X_0\otimes_{A_0}A_\lambda$,
$A_\lambda$ parcourant l'ensemble des sous-$\mathbf Z$-algèbres de type
fini de $A$ contenant $A_0$, et $\mathcal F=v_\lambda^*(\mathcal F_\lambda)$,
où $v_\lambda:X\to X_\lambda$ est la projection canonique et
$\mathcal F_\lambda=\mathcal F_0\otimes_{A_0}A_\lambda$. Il résulte de
(11.2.6) qu'il existe $\lambda$ tel que $\mathcal F_\lambda$ soit
$f_\lambda$-plat au point $x_\lambda=v_\lambda(x)$ (resp.
$f_\lambda$-plat)~; alors le sous-anneau $A_\lambda$ de $A$ vérifie les
conditions de c).

\begin{proposition}[11.2.8]
\label{IV.11.2.8-fr}
Soient $g:X\to S$, $h:Y\to S$ deux morphismes de présentation finie,
$f:X\to Y$ un $S$-morphisme, $\mathcal F$ un
$\mathcal O_X$-Module quasi-cohérent de présentation finie~; pour tout
$s\in S$, soient $X_s=g^{-1}(s)=X\times_S\operatorname{Spec}(k(s))$,
$Y_s=h^{-1}(s)=Y\times_S\operatorname{Spec}(k(s))$, $f_s$ le morphisme
$f\times1:X_s\to Y_s$, $\mathcal F_s=\mathcal F\otimes_{\mathcal O_S}k(s)$.
Alors l'ensemble $E$ des $s\in S$ tels que $\mathcal F_s$ soit $f_s$-plat
est localement constructible.
\end{proposition}

La propriété dont nous voulons démontrer qu'elle est constructible vérifie la
condition (9.2.1, (i)), en vertu de (2.2.13) et (2.5.1). Compte tenu de
(9.2.3), on peut donc se limiter au cas où $S$ est affine, noethérien et
intègre de point générique $\eta$ et à prouver que $E$ ou $S-E$ est un
voisinage de $\eta$ dans $S$. Si $\eta\in S-E$, cela résulte aussitôt du
lemme (9.4.7.1). Si au contraire $\eta\in E$, il résulte tout d'abord de
(11.2.6), appliqué suivant la méthode de (8.1.2, a)), qu'il y a un voisinage
ouvert $U$ de $\eta$ dans $S$ tel que
$\mathcal F|g^{-1}(U)$ soit $h^{-1}(U)$-plat~; \emph{a fortiori} (2.1.4)
$\mathcal F_s$ est $f_s$-plat pour tout $s\in U$, ce qui achève la
démonstration.

Le théorème suivant généralise (11.2.6.1, (ii))~:

\begin{theorem}[11.2.9]
\EGAReaderTarget{ega4.c04.thm.11-2-9}{IV.11.2.9-fr}
(Raynaud). --- Soient $A_0$ un anneau, $(A_\lambda)_{\lambda\in L}$ un système
inductif filtrant de $A_0$-algèbres, $B_0$ une $A_0$-algèbre de présentation
finie, $\mathfrak J_0$ un idéal de type fini de $B_0$, $M_0$ un $B_0$-module
de présentation finie. On pose
$B_\lambda=B_0\otimes_{A_0}A_\lambda$,
$\mathfrak J_\lambda=\mathfrak J_0B_\lambda$,
$M_\lambda=M_0\otimes_{A_0}A_\lambda=M_0\otimes_{B_0}B_\lambda$,
$A=\varinjlim A_\lambda$,
$B=\varinjlim B_\lambda=B_0\otimes_{A_0}A$,
$\mathfrak J=\mathfrak J_0B$,
$M=\varinjlim M_\lambda=M_0\otimes_{A_0}A=M_0\otimes_{B_0}B$. Pour que
$\operatorname{gr}_{\mathfrak J}^*(M)$ soit un $A$-module plat, il faut et il
suffit qu'il existe $\lambda$ tel que
$\operatorname{gr}_{\mathfrak J_\lambda}^*(M_\lambda)$ soit un
$A_\lambda$-module plat~; les homomorphismes canoniques
\begin{equation}
\left\{
\begin{aligned}
\operatorname{gr}_{\mathfrak J_\lambda}^*(M_\lambda)
  \otimes_{A_\lambda}A&\longrightarrow
  \operatorname{gr}_{\mathfrak J}^*(M),\\
\operatorname{gr}_{\mathfrak J_\lambda}^*(M_\lambda)
  \otimes_{A_\lambda}A_\mu&\longrightarrow
  \operatorname{gr}_{\mathfrak J_\mu}^*(M_\mu)
\end{aligned}
\right.
\qquad\text{pour }\lambda\leq\mu
\label{IV.11.2.9.1-fr}
\tag{11.2.9.1}
\end{equation}
sont alors bijectifs et
$\operatorname{gr}_{\mathfrak J_\mu}^*(M_\mu)$ est un $A_\lambda$-module plat
pour $\lambda\leq\mu$.
\end{theorem}

\begin{proof}
Montrons d'abord que les conditions sont suffisantes, ce qui revient à prouver
la bijectivité de (11.2.9.1). Cela résulte du lemme suivant~:

\begin{lemma}[11.2.9.2]
\label{IV.11.2.9.2-fr}
Soient $A$ un anneau, $B$ une $A$-algèbre, $M$ un $B$-module, $\mathfrak J$ un
idéal de $B$. Soient $A'$ une $A$-algèbre~; posons
$B'=B\otimes_AA'$, $M'=M\otimes_AA'=M\otimes_BB'$,
$\mathfrak J'=\mathfrak JB'$. Si
$\operatorname{gr}_{\mathfrak J}^*(M)$ est un $A$-module plat,
l'homomorphisme canonique
\[
(\operatorname{gr}_{\mathfrak J}^*(M))\otimes_AA'
\longrightarrow \operatorname{gr}_{\mathfrak J'}^*(M')
\]
est bijectif.
\end{lemma}

\oldpage[IV-3]{127}

En effet, par récurrence sur $k$, l'hypothèse que les
$\mathfrak J^kM/\mathfrak J^{k+1}M$ sont des $A$-modules plats pour $k\geq0$
entraîne d'abord, par ($\mathbf 0_{\mathrm I}$, 6.1.2), que
$M/\mathfrak J^{k+1}M$ est un $A$-module plat~; en outre
($\mathbf 0_{\mathrm I}$, 6.1.2), la suite
\[
0\longrightarrow(\mathfrak J^{k+1}M)\otimes_AA'
\longrightarrow M\otimes_AA'
\longrightarrow(M/\mathfrak J^{k+1}M)\otimes_AA'\longrightarrow0
\]
est alors exacte, autrement dit $(\mathfrak J^{k+1}M)\otimes_AA'$ s'identifie
à son image canonique $\mathfrak J'^{,k+1}M'$ dans $M'=M\otimes_AA'$.
D'autre part, toujours en vertu de ($\mathbf 0_{\mathrm I}$, 6.1.2), la suite
\[
0\longrightarrow(\mathfrak J^{k+1}M)\otimes_AA'
\longrightarrow(\mathfrak J^kM)\otimes_AA'
\longrightarrow(\mathfrak J^kM/\mathfrak J^{k+1}M)\otimes_AA'
\longrightarrow0
\]
est exacte, ce qui prouve le lemme.

Pour démontrer la nécessité des conditions de (11.2.9), nous procéderons en
plusieurs étapes.

\medskip
\noindent\textbf{I)} \emph{Réduction au cas où les $A_\lambda$ sont
noethériens.} --- On procède comme dans la réduction I) de (11.2.6.1), dont
nous gardons les notations~; il faut simplement commencer par remplacer
$A'_0$ par une sous-$A'_0$-algèbre de type fini $A''_0$ de $A_0$ telle que,
si l'on pose $B''_0=B'_0\otimes_{A'_0}A''_0$, il y ait dans $B''_0$ un idéal
de type fini $\mathfrak J''_0$ tel que $\mathfrak J''_0B_0=\mathfrak J_0$.
On considère pour cela les sous-$A'_0$-algèbres $A'_\alpha$ de type fini de
$A_0$, qui forment une famille filtrante, et l'on a
$B_0=\varinjlim B'_\alpha$, où
$B'_\alpha=B'_0\otimes_{A'_0}A'_\alpha$~; il y a donc un indice $\beta$ tel
qu'un système fini de générateurs de $\mathfrak J_0$ soit formé d'images dans
$B_0$ d'éléments de $B'_\beta$~; on prendra alors $A''_0=A'_\beta$,
$B''_0=B'_\beta$ et pour $\mathfrak J''_0$ l'idéal engendré par ces éléments.
On peut donc supposer que $A_0$ (donc aussi $B_0$) est noethérien. On définit
ensuite comme \emph{loc. cit.} l'ensemble filtrant $H$ et les $A_\xi$,
$B_\xi$, $M_\xi$ pour $\xi\in H$~; on posera aussi
$\mathfrak J_\xi=\mathfrak J_0B_\xi$ pour tout $\xi\in H$. Supposons alors
que l'on ait démontré qu'il existe un $\xi=(\lambda,C_\lambda)$ tel que
$\operatorname{gr}_{\mathfrak J_\xi}^*(M_\xi)$ soit un $C_\lambda$-module
plat~; comme $M_\lambda=M_\xi\otimes_{C_\lambda}A_\lambda$, il résulte de
(11.2.9.2) que $\operatorname{gr}_{\mathfrak J_\lambda}^*(M_\lambda)$ est un
$A_\lambda$-module plat.

\medskip
\noindent\textbf{II)} \emph{Lemmes préliminaires.}

\begin{lemma}[11.2.9.3]
\label{IV.11.2.9.3-fr}
Soient $A$ un anneau,
$B=\bigoplus_{i\geq0}B_i$ une $A$-algèbre graduée à degrés positifs,
$M=\bigoplus_{i\in\mathbf Z}M_i$ un $B$-module gradué. On suppose que $B_0$
est un anneau local, et que chacun des $M_i$ est un $B_0$-module de type fini.
Pour que $M$ soit un $B$-module de type fini, il faut et il suffit que, si
$\mathfrak q$ est l'image réciproque dans $A$ de l'idéal maximal $\mathfrak m$
de $B_0$, $M\otimes_Ak(\mathfrak q)$ soit un
$(B\otimes_Ak(\mathfrak q))$-module de type fini.
\end{lemma}

La condition étant évidemment nécessaire, prouvons qu'elle est suffisante.
Comme $B_0$ (et par suite $B$) est une $A_{\mathfrak q}$-algèbre, on peut
remplacer $A$ par $A_{\mathfrak q}$, autrement dit supposer que $A$ est aussi
un anneau local dont $\mathfrak q$ est l'idéal maximal. Par hypothèse, il existe
un entier $N$ tel que l'homomorphisme canonique
\[
\left(\bigoplus_{-N\leq i\leq N}(M_i\otimes_{B_0}B)\right)
\otimes_Ak(\mathfrak q)\longrightarrow M\otimes_Ak(\mathfrak q)
\]

\oldpage[IV-3]{128}

soit surjectif~; cela signifie aussi que, pour tout entier $n$,
l'homomorphisme canonique de $k(\mathfrak q)$-modules
\[
\left(\bigoplus_{-N\leq i\leq N}(M_i\otimes_{B_0}B_{n-i})\right)
\otimes_Ak(\mathfrak q)\longrightarrow M_n\otimes_Ak(\mathfrak q)
\]
est surjectif. Or, $M_n$ est un $B_0$-module de type fini, $B_0$ un anneau
local et $\mathfrak qB_0\subset\mathfrak m$~; donc le lemme de Nakayama prouve
que chacun des homomorphismes canoniques
\[
\bigoplus_{-N\leq i\leq N}(M_i\otimes_{B_0}B_{n-i})\longrightarrow M_n
\]
est surjectif, d'où la conclusion.

\begin{corollary}[11.2.9.4]
\label{IV.11.2.9.4-fr}
Sous les hypothèses de (11.2.9.3), on suppose en outre que chacun des $B_i$ et
chacun des $M_i$ soit un $B_0$-module de présentation finie, et que $M$ soit
un $A$-module plat. Pour que $M$ soit un $B$-module de présentation finie, il
faut et il suffit que $M\otimes_Ak(\mathfrak q)$ soit un
$(B\otimes_Ak(\mathfrak q))$-module de présentation finie.
\end{corollary}

En vertu de (11.2.9.3), si la condition de l'énoncé est vérifiée, $M$ est un
$B$-module de type fini, donc il existe un $B$-module libre \emph{gradué} de
type fini $L$ (ayant donc une base finie formée d'éléments homogènes) et un
homomorphisme gradué surjectif de degré 0, $u:L\to M$. Soit $R=\ker(u)$, qui
est un $B$-module gradué~; il y a donc un nombre fini d'entiers $m_j$
($1\leq j\leq r$) tels que pour chaque entier $i$, $R_i$ soit le noyau d'un
homomorphisme surjectif
\[
\bigoplus_{1\leq j\leq r}B_{i+m_j}\longrightarrow M_i~;
\]
on conclut alors de l'hypothèse sur les $M_i$ et les $B_i$ que $R_i$ est un
$B_0$-module de type fini (Bourbaki, \emph{Alg. comm.}, chap.~I, \S~2,
n$^\circ$~8, lemme 9). Pour prouver que $R$ est un $B$-module de type fini,
notons qu'en vertu de l'hypothèse de platitude et de
($\mathbf 0_{\mathrm I}$, 6.1.2) la suite
\[
0\longrightarrow R\otimes_Ak(\mathfrak q)
\longrightarrow L\otimes_Ak(\mathfrak q)
\longrightarrow M\otimes_Ak(\mathfrak q)\longrightarrow0
\]
est exacte, et l'hypothèse entraîne donc (Bourbaki, \emph{loc. cit.}) que
$R\otimes_Ak(\mathfrak q)$ est un $(B\otimes_Ak(\mathfrak q))$-module de type
fini~; il suffit donc d'appliquer à $R$ le lemme (11.2.9.3).

\medskip
\noindent\textbf{III)} \emph{Réduction au cas où les homomorphismes de
transitions $\varphi_{\mu\lambda}:A_\lambda\to A_\mu$ (pour
$\lambda\leq\mu$) sont injectifs.} --- Soit $A'_\lambda$ l'image de
$A_\lambda$ par l'homomorphisme canonique $\varphi_\lambda:A_\lambda\to A$~;
il est clair que les $A'_\lambda$ forment un système inductif de sous-anneaux
noethériens de $A$, dont la limite inductive est $A$, et les homomorphismes de
transition $A'_\lambda\to A'_\mu$ (pour $\lambda\leq\mu$) sont injectifs.
Posons $B'_\lambda=B_0\otimes_{A_0}A'_\lambda$,
$\mathfrak J'_\lambda=\mathfrak J_0B'_\lambda$,
$M'_\lambda=M_0\otimes_{A_0}A'_\lambda=M_\lambda\otimes_{A_\lambda}A'_\lambda$~;
on a encore $B=\varinjlim B'_\lambda$, $M=\varinjlim M'_\lambda$. Supposons
que l'on ait prouvé qu'il existe un $\lambda$ tel que, pour $\mu\geq\lambda$,
$\operatorname{gr}_{\mathfrak J'_\mu}^*(M'_\mu)$ soit un $A'_\mu$-module
plat. Soit $\mathfrak a_\lambda$ le noyau de l'homomorphisme $\varphi_\lambda$,
qui est donc un idéal de type fini de l'anneau noethérien $A_\lambda$. Il
résulte de la définition de la limite inductive qu'il existe un indice
$\mu\geq\lambda$ tel que $\varphi_{\mu\lambda}(\mathfrak a_\lambda)=0$,
autrement dit l'homomorphisme $\varphi_{\mu\lambda}$ se factorise en
$\varphi_{\mu\lambda}:A_\lambda\to A'_\lambda\to A_\mu$, et l'on peut écrire
$B_\mu=B'_\lambda\otimes_{A'_\lambda}A_\mu$,
$\mathfrak J_\mu=\mathfrak J'_\lambda B_\mu$, et
$M_\mu=M'_\lambda\otimes_{A'_\lambda}A_\mu$~; on déduit donc de (11.2.9.2)
que $\operatorname{gr}_{\mathfrak J_\mu}^*(M_\mu)$ est un $A_\mu$-module
plat.

\medskip
\noindent\textbf{IV)} \emph{Réduction au cas où
$B_0=A_0[T_1,\ldots,T_n]$ et
$\operatorname{gr}_{\mathfrak J_0}^*(B_0)=B_0$.} --- En vertu de l'hypothèse,
il y a un système $(t_i)_{1\leq i\leq n}$ de générateurs de la $A_0$-algèbre
de type fini $B_0$, tel que les $t_i$ pour $i\leq m$ engendrent l'idéal
$\mathfrak J_0$ de $B_0$. Posons $B'_0=A_0[T_1,\ldots,T_n]$,

\oldpage[IV-3]{129}

algèbre de polynômes (donc noethérienne), et soit $\mathfrak J'_0$ l'idéal de
$B'_0$ engendré par les $T_i$ d'indice $i\leq m$~; on a alors un
$A_0$-homomorphisme surjectif $u_0:B'_0\to B_0$ transformant chaque $T_i$ en
$t_i$ ($1\leq i\leq n$), donc tel que $u_0(\mathfrak J'_0)=\mathfrak J_0$~;
cet homomorphisme permet de considérer $M_0$ comme un $B'_0$-module de type
fini. On pose alors
\[
B'_\lambda=B'_0\otimes_{A_0}A_\lambda
=A_\lambda[T_1,\ldots,T_n],\qquad
\mathfrak J'_\lambda=\mathfrak J'_0B'_\lambda,
\]
\[
B'=B'_0\otimes_{A_0}A=A[T_1,\ldots,T_n],\qquad
\mathfrak J'=\mathfrak J'_0B',
\]
de sorte que $\mathfrak J'$ est l'idéal de $B'$ engendré par les $T_i$
d'indice $i\leq m$~; en outre, il est clair que pour tout entier $k\geq0$,
on a $\mathfrak J'^kM=\mathfrak J^kM$ et
$\mathfrak J'_\lambda{}^kM_\lambda=\mathfrak J_\lambda^kM_\lambda$ pour tout
$\lambda$~; donc
$\operatorname{gr}_{\mathfrak J'}^*(M)=\operatorname{gr}_{\mathfrak J}^*(M)$
en tant que $A$-module et
$\operatorname{gr}_{\mathfrak J'_\lambda}^*(M_\lambda)
=\operatorname{gr}_{\mathfrak J_\lambda}^*(M_\lambda)$ en tant que
$A_\lambda$-module~; on peut donc substituer $B'$, $B'_\lambda$,
$\mathfrak J'$, $\mathfrak J'_\lambda$ à $B$, $B_\lambda$, $\mathfrak J$,
$\mathfrak J_\lambda$ respectivement dans la démonstration~; on notera en
outre que par construction $\operatorname{gr}_{\mathfrak J'}^*(B')$
s'identifie à $B'$ et $\operatorname{gr}_{\mathfrak J'_\lambda}^*(B'_\lambda)$
à $B'_\lambda$.

\medskip
\noindent\textbf{V)} \emph{Preuve du fait que
$\operatorname{gr}_{\mathfrak J}^*(M)$ est un
$\operatorname{gr}_{\mathfrak J}^*(B)$-module de présentation finie.} --- Nous
supposons donc désormais $A_0$ et les $A_\lambda$ noethériens, les
homomorphismes de transition $A_\lambda\to A_\mu$ injectifs,
$B_0=A_0[T_1,\ldots,T_n]$, $\mathfrak J_0$ étant engendré par les $T_i$
d'indice $i\leq m$~; l'anneau
$C_0=\operatorname{gr}_{\mathfrak J_0}^*(B_0)$ s'identifie donc à $B_0$ et
$C=\operatorname{gr}_{\mathfrak J}^*(B)$ s'identifie à $B$~; nous utiliserons
seulement tout d'abord dans ce qui suit le fait que
$C\cong C_0\otimes_{A_0}A$.

Nous aurons besoin de la variante suivante de (6.9.3)~:

\begin{lemma}[11.2.9.5]
\label{IV.11.2.9.5-fr}
Soient $A_0$ un anneau noethérien, $B_0$ une $A_0$-algèbre de type fini,
$\mathfrak J_0$ un idéal de $B_0$, $M_0$ un $B_0$-module de type fini. Il
existe alors une suite $(S_{0i})_{1\leq i\leq m}$ de sous-schémas de
$S_0=\operatorname{Spec}(A_0)$ ayant les propriétés suivantes~:
\begin{enumerate}[label=\arabic{enumi}$^\circ$]
\item Les espaces sous-jacents aux $S_{0i}$ sont deux à deux disjoints et
forment un recouvrement de $S_0$.
\item Pour chaque $i$, l'ensemble $S_{0i}$ est ouvert dans
$\bigcup_{j\geq i}S_{0j}$.
\item Chaque schéma $S_{0i}$ est affine.
\item Si $A_{0i}$ est l'anneau de $S_{0i}$ et si l'on pose
$B_{0i}=B_0\otimes_{A_0}A_{0i}$, $\mathfrak J_{0i}=\mathfrak J_0B_{0i}$,
alors
$\operatorname{gr}_{\mathfrak J_{0i}}^*(M_0\otimes_{A_0}A_{0i})$ est un
$A_{0i}$-module plat.
\end{enumerate}
\end{lemma}

On procède comme dans (6.9.3) par récurrence noethérienne, en supposant le
lemme vrai lorsqu'on y remplace $A_0$ par $A'_0=A_0/\mathfrak a_0$, où
l'idéal $\mathfrak a_0$ est tel que $V(\mathfrak a_0)\neq S_0$, $B_0$ par
$B'_0=B_0\otimes_{A_0}A'_0$, $\mathfrak J_0$ par
$\mathfrak J'_0=\mathfrak J_0B'_0$ et $M_0$ par $M_0\otimes_{A_0}A'_0$. On
considère le nilradical $\mathfrak N_0$ de $A_0$, et il suffit évidemment de
prouver le lemme en remplaçant $A_0$ par $A_0/\mathfrak N_0$ et $B_0$,
$\mathfrak J_0$, $M_0$ par les objets correspondants comme ci-dessus.
Autrement dit, on peut supposer que $A_0$ est \emph{réduit}~; d'autre part,
$\operatorname{gr}_{\mathfrak J_0}^*(B_0)$ est une $A_0$-algèbre de type fini
et $\operatorname{gr}_{\mathfrak J_0}^*(M_0)$ un
$\operatorname{gr}_{\mathfrak J_0}^*(B_0)$-module de type fini~; en vertu du
théorème de platitude générique (6.9.1), il existe un ouvert
$D(t_0)\subset S_0$ tel que si l'on pose $A_{01}=(A_0)_{t_0}$,
$\operatorname{gr}_{\mathfrak J_0}^*(M_0)\otimes_{A_0}A_{01}$ soit un
$A_{01}$-module plat~; mais puisque $A_{01}$ est un $A_0$-module plat,
$\operatorname{gr}_{\mathfrak J_0}^*(M_0)\otimes_{A_0}A_{01}$ s'identifie à
$\operatorname{gr}_{\mathfrak J_{01}}^*(M_0\otimes_{A_0}A_{01})$, où
$B_{01}=B_0\otimes_{A_0}A_{01}$ et $\mathfrak J_{01}=\mathfrak J_0B_{01}$.
Le complémentaire de $D(t_0)$ dans $S_0$ est alors de la forme
$V(\mathfrak a_0)$, et on conclut en appliquant l'hypothèse de récurrence
noethérienne.

Appliquons ce lemme à la situation actuelle, en gardant les mêmes notations~;
posons $U_{0i}=\bigcup_{j\leq i}S_{0j}$, qui est un ouvert quasi-compact de
$S_0$. Il y a donc une famille
$(f_{ik})_{1\leq i\leq m,\,1\leq k\leq n}$ d'éléments de $A_0$ tels que pour
chaque $i$, $U_{0i}$ soit réunion des $D(f_{ik})$ ($1\leq k\leq n$).

\oldpage[IV-3]{130}

Le $C_0$-module $\operatorname{gr}_{\mathfrak J_0}^*(M_0)$ est engendré par
un nombre fini d'éléments homogènes de degré 1, de sorte qu'il y a un
épimorphisme de $C_0$-modules gradués
\[
u_0:C_0^r\longrightarrow\operatorname{gr}_{\mathfrak J_0}^*(M_0).
\]
Comme $C=C_0\otimes_{A_0}A$ et que l'homomorphisme canonique
$\operatorname{gr}_{\mathfrak J_0}^*(M_0)\otimes_{A_0}A
\to\operatorname{gr}_{\mathfrak J}^*(M)$ est surjectif, on déduit donc de
$u_0$ un épimorphisme de $C$-modules gradués
\[
u:C^r\xrightarrow{u_0\otimes1_A}
\operatorname{gr}_{\mathfrak J_0}^*(M_0)\otimes_{A_0}A
\longrightarrow\operatorname{gr}_{\mathfrak J}^*(M),
\]
et tout revient à voir que le $C$-module gradué $Q=\ker(u)$ est \emph{de type
fini}.

\begin{lemma}[11.2.9.6]
\label{IV.11.2.9.6-fr}
Sous les hypothèses précédentes, soit $\mathfrak R_0$ un idéal de $A_0$~;
posons
\[
A'_0=A_0/\mathfrak R_0,\qquad
B'_0=B_0\otimes_{A_0}A'_0=B_0/\mathfrak R_0B_0,\qquad
\mathfrak J'_0=\mathfrak J_0B'_0,
\]
et $\mathfrak R=\mathfrak R_0A$, $A'=A/\mathfrak R$, et supposons en outre
que $\operatorname{gr}_{\mathfrak J'_0}^*(M_0\otimes_{A_0}A'_0)$ soit un
$A'_0$-module plat. Alors il existe un nombre fini d'éléments $q_h$
($1\leq h\leq s$) de $Q$ tels que, pour tout idéal premier
$\mathfrak p\supset\mathfrak RB$ de $B$, les images canoniques des $q_h$ dans
$Q_\mathfrak p$ engendrent le $C_\mathfrak p$-module $Q_\mathfrak p$.
\end{lemma}

Supposons ce lemme démontré, et notons qu'on peut encore l'appliquer en
remplaçant $A_0$ par $(A_0)_{t_0}$, $B_0$, $\mathfrak J_0$, $M_0$,
$A_\lambda$ et $A$ par les objets correspondants $(B_0)_{t_0}$,
$(\mathfrak J_0)_{t_0}$, $(M_0)_{t_0}$, $(A_\lambda)_{t_0}$ et $A_{t_0}$ pour
tout $t_0\in A_0$, puisque $(A_0)_{t_0}$ est un $A_0$-module plat. On
appliquera alors ce lemme en remplaçant successivement $A_0$ par chacun des
anneaux $(A_0)_{f_{ik}}$, $\mathfrak R_0$ étant remplacé par l'idéal
$\mathfrak R_{ik}$ définissant le sous-schéma fermé de $D(f_{ik})$ induit par
$S_{0i}$ sur l'ouvert $S_{0i}\cap D(f_{ik})$ de $S_{0i}$. Comme l'hypothèse
de platitude de (11.2.9.6) est vérifiée dans chacun de ces cas en raison de
(11.2.9.5), on obtient ainsi pour chaque couple $(i,k)$ une famille finie
$(a_{ikh})_{1\leq h\leq l}$ d'éléments de $Q_{f_{ik}}$ dont les images dans
$Q_{\mathfrak p_{ik}}$ engendrent ce $C_{\mathfrak p_{ik}}$-module, pour tout
idéal premier $\mathfrak p_{ik}$ de $B_{f_{ik}}$ contenant
$\mathfrak R_{ik}B_{f_{ik}}$. On peut écrire, pour un entier $d>0$ convenable,
$a_{ikh}=b_{ikh}/f_{ik}^d$ pour tous les indices $i$, $k$, $h$, avec
$b_{ikh}\in Q$. Puisque les $S_{0i}$ recouvrent $S_0$, tout idéal premier
$\mathfrak p$ de $B$ est tel que son image dans $S_0$ appartienne à un
ensemble $S_{0i}\cap D(f_{ik})$, donc l'image $\mathfrak p_{ik}$ de
$\mathfrak p$ dans $B_{f_{ik}}$ contient $\mathfrak R_{ik}B_{f_{ik}}$. Comme
les images des $b_{ikh}$ ($1\leq h\leq l$) dans
$Q_{\mathfrak p_{ik}}=Q_\mathfrak p$ engendrent ce $C_\mathfrak p$-module,
on en conclut que les $b_{ikh}$
($1\leq i\leq m$, $1\leq k\leq n$, $1\leq h\leq l$) engendrent le
$C$-module $Q$ (Bourbaki, \emph{Alg. comm.}, chap.~II, \S~3, n$^\circ$~3,
th.~1).

Reste à démontrer le lemme (11.2.9.6). Posons
\[
C'=C\otimes_AA'=C/\mathfrak RC,\qquad
B'=B\otimes_AA',\qquad \mathfrak J'=\mathfrak JB',\qquad
M'=M\otimes_AA'.
\]
Par hypothèse $\operatorname{gr}_{\mathfrak J}^*(M)$ est un $A$-module plat,
donc $\operatorname{gr}_{\mathfrak J'}^*(M')$ s'identifie à
$\operatorname{gr}_{\mathfrak J}^*(M)\otimes_AA'$ (11.2.9.2), et on a la
suite exacte ($\mathbf 0_{\mathrm I}$, 6.1.2)
\begin{equation}
0\longrightarrow Q\otimes_AA'\longrightarrow C'^{r}
\xrightarrow{u\otimes1_{A'}}\operatorname{gr}_{\mathfrak J'}^*(M')
\longrightarrow0.
\label{IV.11.2.9.7-fr}
\tag{11.2.9.7}
\end{equation}
Posons $C'_0=C_0\otimes_{A_0}A'_0=C_0/\mathfrak R_0C_0$ et considérons
l'épimorphisme de $C'_0$-modules gradués
\[
u'_0:C'_0{}^{r}\xrightarrow{u_0\otimes1_{A'_0}}
\operatorname{gr}_{\mathfrak J_0}^*(M_0)\otimes_{A_0}A'_0
\longrightarrow
\operatorname{gr}_{\mathfrak J'_0}^*(M_0\otimes_{A_0}A'_0).
\]
Soit $Q_0=\ker(u'_0)$~; utilisant le fait que
$\operatorname{gr}_{\mathfrak J'_0}^*(M_0\otimes_{A_0}A'_0)$ est un
$A'_0$-module plat, on voit que son produit tensoriel sur $A'_0$ par $A'$
s'identifie à $\operatorname{gr}_{\mathfrak J'}^*(M')$ (11.2.9.2) et on a la
suite exacte ($\mathbf 0_{\mathrm I}$, 6.1.2)
\begin{equation}
0\longrightarrow Q_0\otimes_{A'_0}A'\longrightarrow C'^{r}
\xrightarrow{u\otimes1_{A'}}\operatorname{gr}_{\mathfrak J'}^*(M')
\longrightarrow0.
\label{IV.11.2.9.8-fr}
\tag{11.2.9.8}
\end{equation}

\oldpage[IV-3]{131}

d'où, par comparaison avec (11.2.9.7), un isomorphisme de $C'$-modules
gradués
\[
Q\otimes_AA'\xrightarrow{\sim}Q_0\otimes_{A'_0}A'.
\]
Comme $C'_0$ est noethérien, $Q_0$ est un $C'_0$-module de type fini, donc
$Q_0\otimes_{A'_0}A'$ est un $C'$-module gradué de type fini, et il en est par
suite de même de $Q\otimes_AA'$~; soient $q_h$ ($1\leq h\leq s$) des éléments
homogènes de $Q$ dont les images dans $Q\otimes_AA'$ engendrent ce
$C'$-module.

Considérons maintenant un idéal premier quelconque
$\mathfrak p\supset\mathfrak RB$ de $B$~; on a
$C_\mathfrak p=\operatorname{gr}_{\mathfrak J_\mathfrak p}^*(B_\mathfrak p)$
par platitude, et
$\operatorname{gr}_{\mathfrak J_\mathfrak p}^0(B_\mathfrak p)
=B_\mathfrak p/\mathfrak J_\mathfrak p$ est réduit à 0 ou est un anneau
local~; pour prouver le lemme, on peut évidemment se borner au second cas. Il
est clair alors que chacun des $(B_\mathfrak p/\mathfrak J_\mathfrak p)$-modules
$\operatorname{gr}_{\mathfrak J_\mathfrak p}^i(B_\mathfrak p)$ est de
présentation finie. D'autre part, pour tout indice $i$,
$M/\mathfrak J^iM$ s'identifie à
$(M_0/\mathfrak J_0^iM_0)\otimes_{A_0}A$, et est donc un $B$-module de
présentation finie. Par récurrence sur $i$, l'hypothèse que
$\operatorname{gr}_{\mathfrak J}^*(M)$ est un $A$-module plat entraîne, en
vertu de ($\mathbf 0_{\mathrm I}$, 6.1.2), que $M/\mathfrak J^iM$ est un
$A$-module plat. L'application de (11.2.6.1) où l'on remplace $M_0$ par
$M_0/\mathfrak J_0^kM_0$ pour $k\leq i$ montre qu'il existe un indice
$\lambda_i$ tel que, pour $\mu\geq\lambda_i$, chacun des
$M_\mu/\mathfrak J_\mu^{k+1}M_\mu$ soit un $A_\mu$-module plat, et par suite
aussi ($\mathbf 0_{\mathrm I}$, 6.1.2) chacun des
$\operatorname{gr}_{\mathfrak J_\mu}^k(M_\mu)$ est un $A_\mu$-module plat
pour $k\leq i$. On déduit par suite de (11.2.9.2) que chacun des
$(B/\mathfrak J)$-modules $\operatorname{gr}_{\mathfrak J}^i(M)$ est de
présentation finie, donc
$\operatorname{gr}_{\mathfrak J_\mathfrak p}^i(M_\mathfrak p)$ est un
$(B_\mathfrak p/\mathfrak J_\mathfrak p)$-module de présentation finie. Par
ailleurs les images des $q_h$ dans
\[
Q_\mathfrak p\otimes_{B_\mathfrak p}k(\mathfrak p)
=(Q\otimes_AA')_\mathfrak p\otimes_{A'}k(\mathfrak p)
\]
engendrent ce module sur
$C_\mathfrak p\otimes_{B_\mathfrak p}k(\mathfrak p)$, puisque
\[
C_\mathfrak p\otimes_{B_\mathfrak p}k(\mathfrak p)
=C'_\mathfrak p\otimes_{A'}k(\mathfrak p)~;
\]
comme on a la suite exacte
\[
Q_\mathfrak p\otimes_{B_\mathfrak p}k(\mathfrak p)
\longrightarrow C_\mathfrak p^{r}\otimes_{B_\mathfrak p}k(\mathfrak p)
\longrightarrow
\operatorname{gr}_{\mathfrak J_\mathfrak p}^*(M_\mathfrak p)
\otimes_{B_\mathfrak p}k(\mathfrak p)\longrightarrow0,
\]
on en conclut que
$\operatorname{gr}_{\mathfrak J_\mathfrak p}^*(M_\mathfrak p)
\otimes_{B_\mathfrak p}k(\mathfrak p)$ est un
$(C_\mathfrak p\otimes_{B_\mathfrak p}k(\mathfrak p))$-module de présentation
finie. Appliquant le lemme (11.2.9.4) où $A$, $B$, $M$ sont remplacés
respectivement par $B_\mathfrak p$, $C_\mathfrak p$ et
$\operatorname{gr}_{\mathfrak J_\mathfrak p}^*(M_\mathfrak p)$, on en conclut
que $Q_\mathfrak p$ est un $C_\mathfrak p$-module de type fini. Utilisant
maintenant le lemme de Nakayama (et le fait que $C=B$) on en déduit que les
images des $q_h$ dans $Q_\mathfrak p$ engendrent ce $C_\mathfrak p$-module.
Ceci achève de prouver le lemme (11.2.9.6) et le fait que
$\operatorname{gr}_{\mathfrak J}^*(M)$ est un $C$-module de présentation
finie.

\medskip
\noindent\textbf{VI)} \emph{Fin de la démonstration.} --- Posons pour abréger
$C_\lambda=C_0\otimes_{A_0}A_\lambda$ (égal en fait à $B_\lambda$),
$N_\lambda=\operatorname{gr}_{\mathfrak J_\lambda}^*(M_\lambda)$ et
$N=\operatorname{gr}_{\mathfrak J}^*(M)$. Notons d'abord que pour chaque
entier $k$, $\mathfrak J^kM$ s'identifie à
$\varinjlim\mathfrak J_\lambda^kM_\lambda$ par exactitude du foncteur
$\varinjlim$~; l'image de $\mathfrak J_\lambda^kM_\lambda$ dans $M_\mu$ pour
$\lambda\leq\mu$ (resp. dans $M$) engendrant $\mathfrak J_\mu^kM_\mu$
(resp. $\mathfrak J^kM$)~; utilisant encore l'exactitude de $\varinjlim$, on
en conclut que $N$ s'identifie canoniquement à $\varinjlim N_\lambda$. Faisant
cette identification, nous allons d'abord prouver que~:
\begin{equation}
\text{\emph{Pour $\lambda$ assez grand, l'homomorphisme canonique
$v_\lambda:N_\lambda\otimes_{A_\lambda}A\to N$ est bijectif.}}
\label{IV.11.2.9.9-fr}
\tag{11.2.9.9}
\end{equation}
Comme les $C_\lambda$ sont noethériens et $N_\lambda$ un $C_\lambda$-module
de type fini, les $C$-modules $N_\lambda\otimes_{A_\lambda}A$ sont de
présentation finie et forment un système inductif filtrant, dont la limite
inductive s'identifie canoniquement à $N$ en vertu du fait que $\varinjlim$
commute aux produits tensoriels. En outre, les homomorphismes de transition
$v_{\mu\lambda}:N_\lambda\otimes_{A_\lambda}A
\to N_\mu\otimes_{A_\mu}A$

\oldpage[IV-3]{132}

pour $\lambda\leq\mu$ et les homomorphismes $v_\lambda$ sont \emph{surjectifs}.
Pour un $\lambda$ fixé, considérons les sous-$C$-modules
$\ker(v_{\mu\lambda})=N'_\mu$ de $N_\lambda\otimes_{A_\lambda}A$~; par
définition de la limite inductive, ils forment une famille filtrante croissante
de sous-$C$-modules de $\ker(v_\lambda)$, dont la réunion est
$\ker(v_\lambda)$~; mais on a vu dans V) que $N$ est un $C$-module de
présentation finie, donc (Bourbaki, \emph{Alg. comm.}, chap.~I, \S~2,
n$^\circ$~8, lemme 9) $\ker(v_\lambda)$ est un $C$-module de type fini~; il
existe par suite un indice $\mu\geq\lambda$ tel que
$N'_\mu=\ker(v_\lambda)$, ce qui signifie (vu le fait que $v_{\mu\lambda}$ est
surjectif) que $v_\mu$ est \emph{injectif}~; comme il est surjectif, cela
prouve (11.2.9.9).

Quitte à remplacer $A_0$ et $M_0$ par $A_\lambda$ et $M_\lambda$ pour
$\lambda$ assez grand, on peut donc supposer que l'homomorphisme canonique
$v_\lambda$ est bijectif pour \emph{tout} $\lambda$. Posons alors
$P_\lambda=N_0\otimes_{C_0}C_\lambda=N_0\otimes_{A_0}A_\lambda$, de sorte que
$N=\varinjlim P_\lambda$. Comme $C_0$ est une $A_0$-algèbre de présentation
finie et $P_0$ un $C_0$-module de présentation finie, on peut appliquer à
$C$ et $N$ le résultat de (11.2.6.1) et on voit donc qu'il existe un indice
$\lambda$ tel que, pour $\mu\geq\lambda$, $P_\mu$ soit un $A_\mu$-module
plat. Or, pour $\mu\geq\lambda$, on a un diagramme commutatif
\[
\xymatrix{
P_\mu \ar[r] \ar[d]^{w_\mu} & N \ar@{=}[d] \\
N_\mu \ar[r] & N
}
\]
où il résulte des définitions que $w_\mu$ est \emph{surjectif}. Comme, en
vertu de III), les homomorphismes $A_\mu\to A$ sont injectifs et que $P_\mu$
est un $A_\mu$-module plat, l'homomorphisme canonique
$P_\mu\to N=P_\mu\otimes_{A_\mu}A$ est \emph{injectif}~; on conclut donc du
diagramme commutatif précédent que $w_\mu$ est aussi injectif, donc bijectif,
et par suite $N_\mu$ est un $A_\mu$-module plat pour $\mu\geq\lambda$, ce qui
achève de prouver (11.2.9).
\end{proof}

\begin{remark}[11.2.10]
\label{IV.11.2.10-fr}
Nous ignorons si la généralisation de (11.2.6, (i)) analogue au th. de Raynaud,
est valable.
\end{remark}

\subsection{Application à l'élimination d'hypothèses noethériennes}
\label{subsection:IV.11.3-fr}
\label{IV.11.3-fr}

\begin{theorem}[11.3.1]
\label{IV.11.3.1-fr}
Soient $f:X\to Y$ un morphisme localement de présentation finie,
$\mathcal F$ un $\mathcal O_X$-Module quasi-cohérent de présentation finie.
Alors l'ensemble $U$ des points $x\in X$ tels que $\mathcal F$ soit $f$-plat
au point $x$ est ouvert dans $X$. Si de plus $\operatorname{Supp}(\mathcal F)=X$,
$f|U$ est un morphisme ouvert.
\end{theorem}

\begin{proof}
La seconde assertion a déjà été démontrée (2.4.6) et n'a été insérée que pour
mémoire.

La question étant locale sur $X$ et $Y$, on peut supposer que $X$ et $Y$ sont
affines, et que $f$ est un morphisme de présentation finie. Soit $x\in X$ un
point tel que $\mathcal F$ soit $f$-plat au point $x$. Appliquant (11.2.7), on
peut supposer que
\[
X=X_0\times_{Y_0}Y,\qquad f=f_0\times1,\qquad
\mathcal F=\mathcal F_0\otimes_{\mathcal O_{Y_0}}\mathcal O_Y,
\]
où $Y_0$ est noethérien, $f_0:X_0\to Y_0$ un morphisme de type fini,
$\mathcal F_0$ un $\mathcal O_{X_0}$-Module cohérent~; en outre, si $x_0$ est la
projection canonique de $x$ dans $X_0$, on peut supposer que $\mathcal F_0$ est
$f_0$-plat au point $x_0$. Alors, en vertu de (11.1.1), l'ensemble $U_0$ des
points de $X_0$ où $\mathcal F_0$ est $f_0$-plat est un voisinage de $x_0$~;
donc $\mathcal F$ est $f$-plat aux points de l'image réciproque de $U_0$ dans
$X$ (2.1.4), et cela prouve que $U$ est un voisinage de $x$.
\end{proof}

\oldpage[IV-3]{133}

\begin{corollary}[11.3.2]
\label{IV.11.3.2-fr}
Soient $f:X\to Y$, $g:Y\to Z$ deux morphismes de présentation finie,
$\mathcal F$ un $\mathcal O_X$-Module quasi-cohérent de présentation finie.
Alors l'ensemble $U$ des $y\in Y$ tels que, pour toute générisation $y'$ de
$y$, $\mathcal F$ soit $(g\circ f)$-plat en tous les points de $f^{-1}(y')$
(i.e. tels que $\mathcal F\otimes_Y\operatorname{Spec}(\mathcal O_{Y,y})$
soit plat relativement au morphisme
$X\times_Y\operatorname{Spec}(\mathcal O_{Y,y})\to Z$) est ouvert dans $Y$.
\end{corollary}

Le même raisonnement que dans (11.1.5) montre que si $V$ est l'ensemble des
$x\in X$ où $\mathcal F$ est $(g\circ f)$-plat, $U$ est l'ensemble des $y\in Y$
dont toute générisation appartient à $E=Y-f(X-V)$. Comme $g\circ f$ est de
présentation finie, $V$ est ouvert dans $X$ en vertu de (11.3.1), donc
ind-constructible (1.9.6), et par suite $X-V$ est pro-constructible~; mais comme
$f$ est quasi-compact, $f(X-V)$ est pro-constructible dans $Y$ (1.9.5, (vii)),
et par suite $E$ est ind-constructible dans $Y$. Mais il résulte alors de
(1.10.1) que $U$ est l'intérieur de $E$, donc ouvert dans $Y$.

\begin{corollary}[11.3.3]
\label{IV.11.3.3-fr}
Soient $A$ un anneau, $B$ une $A$-algèbre de présentation finie, $C$ une
$B$-algèbre de présentation finie, $M$ un $C$-module de présentation finie.
Alors l'ensemble des $\mathfrak q\in\operatorname{Spec}(B)$ tels que
$M_\mathfrak q$ soit un $A$-module plat est ouvert dans
$\operatorname{Spec}(B)$.
\end{corollary}

\begin{proposition}[11.3.4]
\EGAReaderTarget{ega4.c04.prop.11-3-4}{IV.11.3.4-fr}
Soient $f:X\to S$ un morphisme localement de présentation finie, $\mathcal J$
un Idéal quasi-cohérent de type fini de $\mathcal O_X$, $Y$ le sous-préschéma
fermé de $X$ défini par $\mathcal J$, $\mathcal F$ un $\mathcal O_X$-Module de
présentation finie. On suppose que $\operatorname{gr}_{\mathcal J}^*(\mathcal F)$
est $f$-plat. Dans ces conditions~:
\begin{enumerate}[label=\emph{(\roman*)}]
\item $\mathcal F$ est $f$-plat aux points de $Y$.
\item Si $\operatorname{gr}_{\mathcal J}^*(\mathcal O_X)$ est $f$-plat,
$\operatorname{gr}_{\mathcal J}^*(\mathcal O_X)$ est une
$(\mathcal O_X/\mathcal J)$-Algèbre de présentation finie et
$\operatorname{gr}_{\mathcal J}^*(\mathcal F)$ est un
$(\operatorname{gr}_{\mathcal J}^*(\mathcal O_X))$-Module de présentation
finie.
\item Supposons $\operatorname{gr}_{\mathcal J}^*(\mathcal O_X)$ $f$-plat
(ce qui entraîne que $\mathcal O_X/\mathcal J$ est $f$-plat). Alors l'ensemble
$W$ des $y\in Y$ tels que $(\operatorname{gr}_{\mathcal J}^*(\mathcal F))_y$
soit un $(\mathcal O_X/\mathcal J)_y$-module plat est ouvert dans $Y$.
\item Supposons $\operatorname{gr}_{\mathcal J}^*(\mathcal O_X)$ $f$-plat.
Soit $S'\to S$ un morphisme quelconque~; soient $X'=X\times_SS'$,
$p:X'\to X$ la projection canonique, $Y'=p^{-1}(Y)=Y\times_SS'$ le
sous-préschéma fermé de $X'$, défini par
$\mathcal J'=p^*(\mathcal J)\mathcal O_{X'}$, $\mathcal F'=p^*(\mathcal F)
=\mathcal F\otimes_{\mathcal O_S}\mathcal O_{S'}$~; alors, si
$W'=p^{-1}(W)$, on a
$\operatorname{gr}_{\mathcal J'}^*(\mathcal F')|W'
\simeq p^*(\operatorname{gr}_{\mathcal J}^*(\mathcal F)|W)$, et
$\operatorname{gr}_{\mathcal J'}^*(\mathcal F')$ est un
$(\mathcal O_{X'}/\mathcal J')$-Module plat aux points de $W'$.
\end{enumerate}
\end{proposition}

\begin{proof}
Les questions étant locales sur $X$ et $S$, on peut supposer que
$S=\operatorname{Spec}(A)$ et $X=\operatorname{Spec}(B)$ sont affines, $B$ étant
une $A$-algèbre de présentation finie, et $\mathcal F=\widetilde M$, où $M$ est
un $B$-module de présentation finie, $\mathcal J=\widetilde{\mathfrak J}$, où
$\mathfrak J$ est un idéal de type fini de $B$~; en vertu de (8.9.1) et
(8.5.11), il existe un sous-anneau noethérien $A_0$ de $A$, une $A_0$-algèbre
de type fini $B_0$ tels que $B=B_0\otimes_{A_0}A$, un idéal $\mathfrak J_0$ de
$B_0$ tel que $\mathfrak J=\mathfrak J_0B$, un $B_0$-module de type fini $M_0$
tel que $M=M_0\otimes_{A_0}A=M_0\otimes_{B_0}B$. En outre, $A$ est limite
inductive de ses sous-$A_0$-algèbres de type fini $A_\lambda$~; on pose
$B_\lambda=B_0\otimes_{A_0}A_\lambda$,
$M_\lambda=M_0\otimes_{A_0}A_\lambda=M_0\otimes_{B_0}B_\lambda$,
$\mathfrak J_\lambda=\mathfrak J_0B_\lambda$, de sorte que
$B=\varinjlim B_\lambda$, $M=\varinjlim M_\lambda$. Cela étant, par hypothèse
$\operatorname{gr}_{\mathfrak J}^*(M)$ est un $A$-module plat~; donc il résulte
de (11.2.9) qu'il existe un $\lambda$ tel que
$\operatorname{gr}_{\mathfrak J_\lambda}^*(M_\lambda)$ soit un
$A_\lambda$-module plat. Pour prouver que $\mathcal F$ est $f$-plat aux points
de $Y$, on peut donc se borner au cas où $S$ est noethérien~; mais alors
$(\mathbf 0_{\mathrm I},6.1.2)$, appliqué par récurrence sur $n$, prouve que les
$\mathcal F/\mathcal J^{n+1}\mathcal F$ sont $f$-plats, et on conclut par
$(\mathbf 0_{\mathrm{III}},10.2.6)$.

\oldpage[IV-3]{134}

La démonstration de (ii) se ramène de la même façon au cas où $S$ est
noethérien, en utilisant (11.2.9)~; la conclusion est alors évidente,
$\operatorname{gr}_{\mathfrak J}^*(B)$ étant dans ce cas une
$(B/\mathfrak J)$-algèbre de type fini et
$\operatorname{gr}_{\mathfrak J}^*(M)$ un
$\operatorname{gr}_{\mathfrak J}^*(B)$-module de type fini.

Pour prouver (iii), on se ramène comme dans (i) au cas où $S$ et $X$ sont
affines~; avec les mêmes notations, on peut supposer, en vertu de (11.2.9) que
$\operatorname{gr}_{\mathfrak J_\lambda}^*(B_\lambda)$ et
$\operatorname{gr}_{\mathfrak J_\lambda}^*(M_\lambda)$ sont des
$A_\lambda$-modules plats et que l'on a
\[
\operatorname{gr}_{\mathfrak J}^*(B)
=\operatorname{gr}_{\mathfrak J_\lambda}^*(B_\lambda)\otimes_{A_\lambda}A
\quad\text{et}\quad
\operatorname{gr}_{\mathfrak J}^*(M)
=\operatorname{gr}_{\mathfrak J_\lambda}^*(M_\lambda)\otimes_{A_\lambda}A.
\]
Par suite $\operatorname{gr}_{\mathfrak J}^*(B)$ est une
$(B/\mathfrak J)$-algèbre de présentation finie et
$\operatorname{gr}_{\mathfrak J}^*(M)$ un
$\operatorname{gr}_{\mathfrak J}^*(B)$-module de présentation finie. Si $W^m$
est l'ensemble des $\mathfrak p\in\operatorname{Spec}(B/\mathfrak J)$ tels que
$\operatorname{gr}_{\mathfrak J}^m(M)_\mathfrak p$ soit un
$(B/\mathfrak J)$-module plat, $W^m$ est ouvert dans
$\operatorname{Spec}(B/\mathfrak J)$ (11.3.1) et l'on a
$W=\bigcap_mW^m$, donc $W$ est stable par générisation. L'assertion (iii)
résulte alors de (11.3.2) appliqué en y prenant
$X=\operatorname{Spec}(\operatorname{gr}_{\mathfrak J}^*(B))$,
$Y=Z=\operatorname{Spec}(B/\mathfrak J)$ et
$\mathcal F=(\operatorname{gr}_{\mathfrak J}^*(M))^{\sim}$.

Enfin, (iv) découle aussitôt de (11.2.9.2).

Généralisant la définition de (6.10.1), on dit que pour un préschéma $X$, un
sous-préschéma fermé $Y$ de $X$ défini par un Idéal quasi-cohérent $\mathcal J$
de $\mathcal O_X$, et un $\mathcal O_X$-Module quasi-cohérent $\mathcal F$,
$\mathcal F$ est \emph{normalement plat le long de $Y$ en un point $y$} si
$(\operatorname{gr}_{\mathcal J}^*(\mathcal F))_y$ est un
$\mathcal O_{Y,y}$-module plat. On dit que $\mathcal F$ est
\emph{normalement plat le long de $Y$} s'il l'est en tous les points de $Y$.
\end{proof}

\begin{corollary}[11.3.5]
\label{IV.11.3.5-fr}
Sous les hypothèses générales de (11.3.4), on suppose que
$\operatorname{gr}_{\mathcal J}^*(\mathcal O_X)$ et
$\operatorname{gr}_{\mathcal J}^*(\mathcal F)$ soient $f$-plats. Alors
l'ensemble $W$ des $y\in Y$ tels que $\mathcal F$ soit normalement plat le long
de $Y$ au point $y$ est ouvert dans $Y$, et (avec les notations de
(11.3.4, (iv))) $\mathcal F'$ est normalement plat le long de $Y'$ en tous les
points de $W'=p^{-1}(W)$.
\end{corollary}

\begin{proposition}[11.3.6]
\label{IV.11.3.6-fr}
Les notations étant celles de (8.5.1) et (8.8.1), on suppose $S_\alpha$
quasi-compact, $X_\alpha$ de présentation finie sur $S_\alpha$, $Y_\alpha$ un
sous-préschéma fermé de $X_\alpha$ défini par un Idéal quasi-cohérent
$\mathcal J_\alpha$ de type fini de $\mathcal O_{X_\alpha}$ tel que
$\operatorname{gr}_{\mathcal J_\alpha}^*(\mathcal O_{X_\alpha})$ soit plat sur
$S_\alpha$~; enfin on suppose que $\mathcal F_\alpha$ est un
$\mathcal O_{X_\alpha}$-Module de présentation finie. Pour que $\mathcal F$
soit normalement plat le long de $Y$, il faut et il suffit qu'il existe
$\lambda$ tel que $\mathcal F_\lambda$ soit normalement plat le long de
$Y_\lambda$.
\end{proposition}

\begin{proof}
Notons que $Y$ (resp. $Y_\lambda$) est le sous-préschéma fermé de $X$ (resp.
$X_\lambda$) défini par $\mathcal J=\mathcal J_\alpha\mathcal O_X$ (resp.
$\mathcal J_\lambda=\mathcal J_\alpha\mathcal O_{X_\lambda}$)~; en vertu de
l'hypothèse et de (11.2.9.2), on a
\[
\operatorname{gr}_{\mathcal J}^*(\mathcal O_X)
=\operatorname{gr}_{\mathcal J_\alpha}^*(\mathcal O_{X_\alpha})
 \otimes_{\mathcal O_{S_\alpha}}\mathcal O_S
\]
et
\[
\operatorname{gr}_{\mathcal J_\lambda}^*(\mathcal O_{X_\lambda})
=\operatorname{gr}_{\mathcal J_\alpha}^*(\mathcal O_{X_\alpha})
 \otimes_{\mathcal O_{S_\alpha}}\mathcal O_{S_\lambda}
\quad\text{pour }\lambda\geq\alpha,
\]
et $\operatorname{gr}_{\mathcal J}^*(\mathcal O_X)$ (resp.
$\operatorname{gr}_{\mathcal J_\lambda}^*(\mathcal O_{X_\lambda})$) est plat
sur $S$ (resp. $S_\lambda$), ce qui entraîne que $Y$ est plat sur $S$ (resp.
$Y_\lambda$ plat sur $S_\lambda$). Si $\mathcal F_\lambda$ est normalement plat
le long de $Y_\lambda$, $\operatorname{gr}_{\mathcal J_\lambda}^*
(\mathcal F_\lambda)|Y_\lambda$ est plat sur $Y_\lambda$, donc aussi sur
$S_\lambda$, et comme
$\operatorname{gr}_{\mathcal J_\lambda}^*(\mathcal F_\lambda)|x_\lambda=0$
pour $x_\lambda\notin Y_\lambda$,
$\operatorname{gr}_{\mathcal J_\lambda}^*(\mathcal F_\lambda)$ est plat sur
$S_\lambda$. On en conclut (11.2.9.2) que
\[
\operatorname{gr}_{\mathcal J}^*(\mathcal F)
=\operatorname{gr}_{\mathcal J_\lambda}^*(\mathcal F_\lambda)
 \otimes_{\mathcal O_{S_\lambda}}\mathcal O_S,
\]
donc $\mathcal F$ est normalement plat le long de $Y$. Pour prouver la
réciproque, on peut supposer que $S_\alpha$ et $X_\alpha$ sont affines et
adopter les notations de (11.2.9)~; puisque
$\operatorname{gr}_{\mathfrak J}^*(B)$ est un $A$-module plat, il en est de même
de $\operatorname{gr}_{\mathfrak J}^*(M)$, donc, en vertu de (11.2.9), il
existe $\lambda\geq\alpha$ tel que
$\operatorname{gr}_{\mathfrak J_\lambda}^*(M_\lambda)$ soit un
$A_\lambda$-module plat, d'où
$\operatorname{gr}_{\mathfrak J}^*(M)
=\operatorname{gr}_{\mathfrak J_\lambda}^*(M_\lambda)\otimes_{A_\lambda}A$.
En outre (11.3.4, (ii)),
$\operatorname{gr}_{\mathfrak J_\lambda}^*(M_\lambda)$ est un
$\operatorname{gr}_{\mathfrak J_\lambda}^*(B_\lambda)$-module de présentation
finie et $\operatorname{gr}_{\mathfrak J_\lambda}^*(B_\lambda)$ une
$(B_\lambda/\mathfrak J_\lambda)$-algèbre de présentation finie. La conclusion
résulte alors de (11.2.6).
\end{proof}

\oldpage[IV-3]{135}

\begin{proposition}[11.3.7]
\label{IV.11.3.7-fr}
Soient $f:X\to Y$ un morphisme localement de présentation finie,
$\mathcal F$, $\mathcal F'$ deux $\mathcal O_X$-Modules quasi-cohérents de
présentation finie, $u:\mathcal F'\to\mathcal F$ un
$\mathcal O_X$-homomorphisme, $x$ un point de $X$ et $y=f(x)$~; on suppose que
$\mathcal F$ est $f$-plat au point $x$. Les conditions suivantes sont alors
équivalentes~:
\begin{enumerate}[label=\emph{\alph*)}]
\item On a $(\ker u)_x=0$ et $\operatorname{Coker}u$ est un
$\mathcal O_X$-Module $f$-plat au point $x$.
\item L'homomorphisme
\[
(u\otimes1)_x:(\mathcal F'\otimes_{\mathcal O_Y}k(y))_x
\longrightarrow(\mathcal F\otimes_{\mathcal O_Y}k(y))_x
\]
est injectif.
\end{enumerate}
Si de plus $\mathcal F$ est $f$-plat, l'ensemble des points $x$ vérifiant les
conditions équivalentes précédentes est ouvert dans $X$.
\end{proposition}

\begin{proof}
La condition a) entraîne b) en vertu de $(\mathbf 0_{\mathrm I},6.1.2)$, sans
hypothèse sur $\mathcal F$. Pour démontrer la réciproque, on peut se borner au
cas où $X$ et $Y$ sont affines, puis, en vertu de (8.9.1) et (11.2.7), supposer
que l'on a
\[
X=X_0\times_{Y_0}Y,\quad f=f_0\times1,\quad
\mathcal F=\mathcal F_0\otimes_{\mathcal O_{Y_0}}\mathcal O_Y,
\]
\[
\mathcal F'=\mathcal F'_0\otimes_{\mathcal O_{Y_0}}\mathcal O_Y,
\qquad u=u_0\otimes1_Y,
\]
où $Y_0$ est noethérien, $f_0:X_0\to Y_0$ un morphisme de type fini,
$\mathcal F_0$, $\mathcal F'_0$ deux $\mathcal O_{X_0}$-Modules cohérents,
$u_0:\mathcal F'_0\to\mathcal F_0$ un homomorphisme~; en outre, si $x_0$ est la
projection de $x$ dans $X_0$, on peut supposer que $\mathcal F_0$ est
$f_0$-plat au point $x_0$. Posons $y_0=f_0(x_0)$~; en vertu de la transitivité
des fibres $(\mathbf I,3.6.4)$, le morphisme projection
$f^{-1}(y)\to f_0^{-1}(y_0)$ est fidèlement plat (2.2.13), et comme
\[
(u\otimes1)_x=((u_0\otimes1)_{x_0})\otimes1_{k(y)},
\]
l'hypothèse que $(u\otimes1)_x$ est injectif entraîne qu'il en est de même de
$(u_0\otimes1)_{x_0}$ (2.2.7). Or, cela entraîne, par
$(\mathbf 0_{\mathrm{III}},10.2.4)$ appliqué aux anneaux locaux noethériens
$\mathcal O_{Y_0,y_0}$ et $\mathcal O_{X_0,x_0}$ et aux
$\mathcal O_{X_0,x_0}$-modules $(\mathcal F'_0)_{x_0}$ et
$(\mathcal F_0)_{x_0}$ (dont le second est plat sur $\mathcal O_{Y_0,y_0}$),
que l'on a $(\ker u_0)_{x_0}=0$ et que $\operatorname{Coker}u_0$ est
$f_0$-plat au point $x_0$. On déduit d'abord de là que
$\operatorname{Coker}u$ est $f$-plat au point $x$ (2.1.4)~; en vertu de
l'exactitude à droite du produit tensoriel, on a d'ailleurs
\[
\operatorname{Coker}u=(\operatorname{Coker}u_0)
 \otimes_{\mathcal O_{Y_0}}\mathcal O_Y~;
\]
appliquant alors $(\mathbf 0_{\mathrm I},6.1.2)$ à la suite (exacte par
hypothèse)
\[
0\longrightarrow(\mathcal F'_0)_{x_0}\longrightarrow
(\mathcal F_0)_{x_0}\longrightarrow(\operatorname{Coker}u_0)_{x_0}
\longrightarrow0,
\]
on en déduit que $u_x$ est un homomorphisme injectif.

Enfin, il résulte de (11.1.2) que l'ensemble $U_0$ des points $z_0\in X_0$ tels
que le morphisme $(u_0\otimes1)_{z_0}$ soit injectif est ouvert~; par platitude
on en déduit que pour tout $z\in X$ au-dessus de $z_0\in U_0$ le morphisme
$(u\otimes1)_z$ est injectif, donc l'ensemble de ces points contient l'image
réciproque $U$ de $U_0$ dans $X$ et est par suite un voisinage de $x$, ce qui
achève la démonstration.
\end{proof}

\begin{theorem}[11.3.8]
\label{IV.11.3.8-fr}
Soient $f:X\to Y$ un morphisme localement de présentation finie,
$\mathcal F$ un $\mathcal O_X$-Module quasi-cohérent de présentation finie,
$(g_i)_{1\leq i\leq n}$ une suite de sections de $\mathcal O_X$ au-dessus de
$X$~; posons
\[
\mathcal F_i=\mathcal F\big/\sum_{j=1}^{i}g_j\mathcal F
\quad\text{pour }0\leq i\leq n
\quad(\text{avec }\mathcal F_0=\mathcal F).
\]
Soient $x$ un point de $X$, $y=f(x)$, et posons
$X_y=f^{-1}(y)=X\times_Y\operatorname{Spec}(k(y))$,
$\mathcal F_y=\mathcal F\otimes_{\mathcal O_Y}k(y)$, qui est un
$\mathcal O_{X_y}$-Module. On suppose que les $(g_i)_x$ appartiennent à l'idéal
maximal $\mathfrak m_x$. Les conditions suivantes sont équivalentes~:
\begin{enumerate}[label=\emph{\alph*)}]
\item La suite $((g_i)_x)$ est $\mathcal F_x$-régulière et les $\mathcal F_i$
$(0\leq i\leq n)$ sont $f$-plats au point $x$.
\item La suite $((g_i)_x)$ est $\mathcal F_x$-régulière et $\mathcal F_n$ est
$f$-plat au point $x$.
\item[\emph{b$'$)}] Il existe un voisinage $U$ de $x$ tel que la suite
$(g_i|U)$ soit $(\mathcal F|U)$-quasi-régulière, et $\mathcal F_n$ est
$f$-plat au point $x$.

\oldpage[IV-3]{136}

\item $\mathcal F$ est $f$-plat au point $x$, et la suite
$((g_i\otimes1)_x)$ d'éléments de $\mathcal O_{X_y,x}$ images des $(g_i)_x$ est
$(\mathcal F_y)_x$-régulière.
\item $\mathcal F$ est $f$-plat au point $x$, et pour tout morphisme
$Y'\to Y$ et tout point $x'\in X'=X\times_YY'$ au-dessus de $x$, si l'on pose
$\mathcal F'=\mathcal F\otimes_YY'$, la suite $(g_i\otimes1)_{x'}$
d'éléments de $\mathcal O_{X',x'}$ est $\mathcal F'_{x'}$-régulière.
\end{enumerate}
En outre l'ensemble des $x\in X$ vérifiant ces conditions est ouvert dans
l'ensemble $Z$ des $x$ tels que $g_i(x)=0$ pour tout $i$.
\end{theorem}

\begin{proof}
Le fait que a) entraîne d) résulte aussitôt de $(\mathbf 0,15.1.15)$, et c) est
un cas particulier de d)~; par ailleurs a) implique b) trivialement. Montrons
que b) ou c) entraîne a). Les $\mathcal F_i$ sont de présentation finie~; le
fait que c) implique a) résulte alors aussitôt de (11.3.7) par récurrence sur
$n$, et cela montre aussi que l'ensemble des $x\in X$ vérifiant c) est ouvert
dans $Z$. Pour montrer que b) entraîne a), on est aussitôt ramené, par
récurrence sur $n$, au cas $n=1$~; nous écrirons $g$ au lieu de $g_1$. La
question étant locale sur $X$ et $Y$, on peut supposer
$Y=\operatorname{Spec}(A)$, $X=\operatorname{Spec}(B)$ affines, $B$ étant une
$A$-algèbre de présentation finie, $\mathcal F=\widetilde M$, où $M$ est un
$B$-module de présentation finie. On peut alors (8.9.1) écrire
$B=B_0\otimes_{A_0}A$, $M=M_0\otimes_{A_0}A$, où $A_0$ est un sous-anneau
noethérien de $A$, $B_0$ une $A_0$-algèbre de type fini et $M_0$ un
$B_0$-module de type fini. En outre $A$ est limite inductive filtrante de ses
sous-anneaux $A_\lambda$ qui sont des $A_0$-algèbres de type fini (donc
noethériens), et si l'on pose
$B_\lambda=B_0\otimes_{A_0}A_\lambda$,
$M_\lambda=M_0\otimes_{A_0}A_\lambda$, $B_\lambda$ est une
$A_\lambda$-algèbre de type fini, $M_\lambda$ un $B_\lambda$-module de type
fini et l'on a $B=\varinjlim B_\lambda$, $M=\varinjlim M_\lambda$. Il existe
donc un indice $\lambda$ et un élément $g_\lambda\in B_\lambda$ tels que
$g=g_\lambda\otimes1$~; revenant aux notations géométriques et posant
$Y_\lambda=\operatorname{Spec}(A_\lambda)$,
$X_\lambda=\operatorname{Spec}(B_\lambda)$ et
$\mathcal F_\lambda=\widetilde{M_\lambda}$, il nous suffira de prouver qu'il y
a un $\mu\geq\lambda$ tel qu'au point $x_\mu\in X_\mu$ projection de $x$,
$(g_\mu)_{x_\mu}$ soit $(\mathcal F_\mu)_{x_\mu}$-régulier et
$\mathcal F_\mu/g_\mu\mathcal F_\mu$ plat sur $Y_\mu$ au point $x_\mu$. On en
déduira en effet, par $(\mathbf 0,15.1.16)$ que $\mathcal F_\mu$ est plat sur
$Y_\mu$ au point $x_\mu$, donc $\mathcal F$ plat sur $Y$ au point $x$.

Le fait que b) entraîne a) résultera donc de la proposition suivante~:

\begin{proposition}[11.3.9]
\label{IV.11.3.9-fr}
Les notations et hypothèses générales étant celles de (8.5.1) et (8.8.1),
supposons que $f_\alpha:X_\alpha\to Y_\alpha$ soit un morphisme localement de
présentation finie. Soient $\mathcal F_\alpha$, $\mathcal G_\alpha$ deux
$\mathcal O_{X_\alpha}$-Modules quasi-cohérents de présentation finie,
$h_\alpha:\mathcal F_\alpha\to\mathcal G_\alpha$ un
$\mathcal O_{X_\alpha}$-homomorphisme, $\mathcal H_\alpha$ son conoyau. Soient
enfin $x$ un point de $X$, $x_\lambda$ sa projection dans $X_\lambda$ pour
$\lambda\geq\alpha$. Pour que $h_x$ soit injectif et que
$\mathcal H=\operatorname{Coker}h$ soit $f$-plat au point $x$, il faut et il
suffit qu'il existe $\lambda\geq\alpha$ tel que $(h_\lambda)_{x_\lambda}$ soit
injectif et que $\mathcal H_\lambda=\operatorname{Coker}h_\lambda$ soit
$f_\lambda$-plat au point $x_\lambda$. En outre, l'ensemble des $x\in X$ ayant
les propriétés précédentes est ouvert dans $X$.
\end{proposition}

Rappelons qu'en vertu de l'exactitude à droite du produit tensoriel, on a
$\mathcal H_\mu=v_{\lambda\mu}^*(\mathcal H_\lambda)$ pour $\lambda\leq\mu$ et
$\mathcal H=v_\lambda^*(\mathcal H_\lambda)$, ce qui justifie les notations.

La suffisance de la condition provient de ce que, si la suite
\[
0\longrightarrow(\mathcal F_\lambda)_{x_\lambda}
\longrightarrow(\mathcal G_\lambda)_{x_\lambda}
\longrightarrow(\mathcal H_\lambda)_{x_\lambda}\longrightarrow0
\]
est exacte et $(\mathcal H_\lambda)_{x_\lambda}$ un
$\mathcal O_{f_\lambda(x_\lambda)}$-module plat, $\mathcal H_x$ est un
$\mathcal O_{f(x)}$-module plat par changement de base (2.1.4) et la suite
$0\to\mathcal F_x\to\mathcal G_x\to\mathcal H_x\to0$ est exacte (2.1.8).
Pour prouver que

\oldpage[IV-3]{137}

la condition est nécessaire, notons que $\mathcal H$ est de présentation
finie~; la question étant locale sur $X$, on peut supposer $X$ et $Y$ affines,
et, en vertu de (11.3.1), supposer que $\mathcal H$ est $f$-plat. Notons
maintenant le

\begin{lemma}[11.3.9.1]
\label{IV.11.3.9.1-fr}
Soient $f:X\to Y$ un morphisme localement de présentation finie,
$\mathcal G$, $\mathcal H$ deux $\mathcal O_X$-Modules quasi-cohérents de
présentation finie tels que $\mathcal H$ soit $f$-plat,
$p:\mathcal G\to\mathcal H$ un homomorphisme surjectif. Alors $\ker(p)$ est un
$\mathcal O_X$-Module de présentation finie.
\end{lemma}

On peut en effet supposer $X$, $Y$ affines, et qu'il existe un morphisme
$Y\to Y_0$ où $Y_0$ est noethérien, un morphisme $f_0:X_0\to Y_0$ de type fini
tels que $X=X_0\times_{Y_0}Y$, $f=(f_0)_{(Y)}$, deux
$\mathcal O_{X_0}$-Modules cohérents $\mathcal G_0$, $\mathcal H_0$ et un
homomorphisme $p_0:\mathcal G_0\to\mathcal H_0$ tels que $\mathcal G$,
$\mathcal H$ et $p$ se déduisent de $\mathcal G_0$, $\mathcal H_0$ et $p_0$ par
changement de base (8.9.1 et 8.5.2). En outre, on peut supposer $p_0$ surjectif
(8.5.7) et $\mathcal H_0$ $f_0$-plat (11.2.7). Alors, si
$\mathcal K_0=\ker(p_0)$, $\mathcal K_0$ est un
$\mathcal O_{X_0}$-Module cohérent $(\mathbf 0_{\mathrm I},5.3.4)$ et en vertu
de $(\mathbf 0_{\mathrm I},6.1.2)$, $\mathcal K=\ker(p)$ se déduit de
$\mathcal K_0$ par changement de base, donc est de présentation finie.

Ce lemme étant démontré, posons $\mathcal K=\ker(\mathcal G\to\mathcal H)$,
qui est donc de présentation finie. On a un homomorphisme canonique
$q:\mathcal F\to\mathcal K$, et par hypothèse
$q_x:\mathcal F_x\to\mathcal K_x$ est un isomorphisme~; par suite
$(\mathbf 0_{\mathrm I},5.2.7)$ il existe un voisinage $U$ de $x$ tel que
$q|U$ soit un isomorphisme, et en restreignant $X$, on peut supposer que la
suite
\[
0\longrightarrow\mathcal F\xrightarrow{h}\mathcal G
\longrightarrow\mathcal H\longrightarrow0
\]
est \emph{exacte}. Cela étant, il résulte de (11.2.6) que pour un $\lambda$
assez grand, $\mathcal H_\lambda$ est un $\mathcal O_{X_\lambda}$-Module plat~;
posons $\mathcal K_\lambda=\ker(\mathcal G_\lambda\to\mathcal H_\lambda)$ et
$\mathcal K_\mu=v_{\lambda\mu}^*(\mathcal K_\lambda)$ pour $\mu\geq\lambda$~;
il résulte de $(\mathbf 0_{\mathrm I},6.1.2)$ que l'on a
$\mathcal K_\mu=\ker(\mathcal G_\mu\to\mathcal H_\mu)$ et
$\mathcal K=v_\lambda^*(\mathcal K_\lambda)
=\ker(\mathcal G\to\mathcal H)=\mathcal F$. On a d'autre part pour tout
$\mu\geq\lambda$ un homomorphisme canonique
$q_\mu:\mathcal F_\mu\to\mathcal K_\mu$ avec
$q_\mu=v_{\lambda\mu}^*(q_\lambda)$ pour $\lambda\leq\mu$, et
$q=v_\mu^*(q_\mu)$. Comme $q$ est un isomorphisme, il résulte de (8.5.2.4) et
(11.3.9.1) qu'il existe $\mu\geq\lambda$ tel que $q_\mu$ soit un isomorphisme,
et par suite $u_\mu:\mathcal F_\mu\to\mathcal G_\mu$ un homomorphisme
injectif.

Revenons à la démonstration de (11.3.8).

Puisque l'ensemble des $x\in X$ vérifiant b) est ouvert dans $X$, il est clair
que b) entraîne b$'$. Montrons enfin que b$'$ entraîne c). En premier lieu,
$\mathcal F_n$ est $f$-plat dans un voisinage de $x$ (11.3.1) et on peut donc
se borner au cas où $U=X$ et où $\mathcal F_n$ est $f$-plat. Comme par
définition $\operatorname{gr}_{\mathcal J}^*(\mathcal F)$ est isomorphe à
$\mathcal F_n\otimes_{\mathcal O_Y}\mathcal O_Y[T_1,\ldots,T_n]$
$(\mathbf 0,15.2.2)$, il est $f$-plat, et on en conclut par (11.3.4, (i)) que
$\mathcal F$ lui-même est $f$-plat au voisinage de $x$. D'autre part, si
$(\mathcal J_y)_x$ est l'idéal de $\mathcal O_{X_y,x}$ engendré par les
$(g_i\otimes1)_x$, il résulte de (11.2.9.2) que, dans le diagramme
\[
\xymatrix{
(\operatorname{gr}_{\mathcal J_x}^0(\mathcal F_x)[T_1,\ldots,T_n])
 \otimes_{\mathcal O_y}k(y) \ar[r] \ar[d] &
\operatorname{gr}_{\mathcal J_x}^*(\mathcal F_x)
 \otimes_{\mathcal O_y}k(y) \ar[d] \\
(\operatorname{gr}_{(\mathcal J_y)_x}^0((\mathcal F_y)_x))
 [T_1,\ldots,T_n] \ar[r] &
\operatorname{gr}_{(\mathcal J_y)_x}^*((\mathcal F_y)_x)
}
\]

\oldpage[IV-3]{138}

les flèches verticales sont des isomorphismes~; comme la première ligne est un
isomorphisme, il en est de même de la seconde, donc la suite
$((g_i\otimes1)_x)$ est $(\mathcal F_y)_x$-quasi-régulière, et par suite aussi
$(\mathcal F_y)_x$-régulière, puisque $X_y$ est localement de type fini sur
$k(y)$. C.Q.F.D.
\end{proof}

\begin{theorem}[11.3.10]
\EGAReaderTarget{ega4.c04.thm.11-3-10}{IV.11.3.10-fr}
(critère de platitude par fibres). --- Soient $S$ un préschéma,
$g:X\to S$, $h:Y\to S$ deux morphismes, $f:X\to Y$ un $S$-morphisme,
$\mathcal F$ un $\mathcal O_X$-Module quasi-cohérent, $x$ un point de $X$,
$y=f(x)$, $s=h(y)=g(x)$. On suppose vérifiée l'une des deux hypothèses
suivantes~:
\begin{enumerate}[label=\arabic{enumi}\textsuperscript{o}]
\item $S$, $X$ et $Y$ sont localement noethériens et $\mathcal F$ est cohérent.
\item $g$ et $h$ sont localement de présentation finie et $\mathcal F$ est de
présentation finie.
\end{enumerate}
Alors, avec les notations de (9.4.1), si $\mathcal F_x\neq0$, les conditions
suivantes sont équivalentes~:
\begin{enumerate}[label=\emph{\alph*)}]
\item $\mathcal F$ est $g$-plat au point $x$ et $\mathcal F_s$ est $f_s$-plat
au point $x$.
\item $h$ est plat au point $y$ et $\mathcal F$ est $f$-plat au point $x$.
\end{enumerate}
En outre, dans l'hypothèse 2\textsuperscript{o}, l'ensemble des $x\in X$
vérifiant la condition b) est ouvert dans $X$.
\end{theorem}

\begin{proof}
La dernière assertion résulte de (11.3.1) appliqué à $\mathcal O_Y$ et au
morphisme $h$ d'une part, à $\mathcal F$ et au morphisme $f$ (qui est
localement de présentation finie) de l'autre.

Comme $g=h\circ f$, il est clair que b) implique a) sans supposer
1\textsuperscript{o} ni 2\textsuperscript{o} (2.1.6 et 2.1.4). Pour prouver que
a) entraîne b), on peut se borner au cas où $S$, $X$ et $Y$ sont affines~;
sous l'hypothèse 2\textsuperscript{o}, en appliquant (11.2.7), on se ramène au
cas où de plus $S$ est noethérien, c'est-à-dire qu'on peut se borner à
considérer le cas où l'hypothèse 1\textsuperscript{o} est satisfaite. Alors
l'assertion équivaut au lemme suivant, qui améliore
$(\mathbf 0_{\mathrm{III}},10.2.5)$~:

\begin{lemma}[11.3.10.1]
\label{IV.11.3.10.1-fr}
Soient $\rho:A\to B$, $\sigma:B\to C$ des homomorphismes locaux d'anneaux
locaux noethériens, $k$ le corps résiduel de $A$, $M$ un $C$-module
$\neq0$ de type fini. Alors les conditions suivantes sont équivalentes~:
\begin{enumerate}[label=\emph{\alph*)}]
\item $M$ est un $A$-module plat et $M\otimes_Ak$ est un
$(B\otimes_Ak)$-module plat.
\item $B$ est un $A$-module plat et $M$ est un $B$-module plat.
\end{enumerate}
\end{lemma}

Nous établirons d'abord le lemme plus général suivant~:

\begin{lemma}[11.3.10.2]
\label{IV.11.3.10.2-fr}
Soient $A$ un anneau commutatif, $B$ une $A$-algèbre commutative,
$\mathfrak J$ un idéal de $A$, $M$ un $B$-module. On considère d'une part les
conditions suivantes~:
\begin{enumerate}[label=\emph{(\roman*)}]
\item $\mathfrak J$ est nilpotent.
\item $B$ est noethérien et $M$ est idéalement séparé pour la topologie
$\mathfrak JB$-préadiques $(\mathbf 0_{\mathrm{III}},10.2.1)$.
\item $\mathfrak JB$ est contenu dans le radical de $B$.
\end{enumerate}
On considère d'autre part les quatre propriétés~:
\begin{enumerate}[label=\emph{\alph*)}]
\item $M$ est un $B$-module plat.
\item $B$ est un $A$-module plat.
\item $M$ est un $A$-module plat et $M/\mathfrak JM$ un
$(B/\mathfrak JB)$-module plat.
\item $B/\mathfrak JB$ est un $(A/\mathfrak J)$-module plat et pour tout idéal
maximal $\mathfrak m\supset\mathfrak JB$ de $B$ on a $\mathfrak mM\neq M$.
\end{enumerate}
Alors~:
\begin{enumerate}[label=\arabic{enumi}\textsuperscript{o}]
\item Si l'une des conditions (i), (ii) est vérifiée, la conjonction de a) et
b) implique c), et c) implique a).

\oldpage[IV-3]{139}

\item Si la condition (i) ou la conjonction de (ii) et (iii) est vérifiée, la
conjonction de c) et d) implique la conjonction de a) et b).
\end{enumerate}
\end{lemma}

\noindent 1\textsuperscript{o} La première assertion est immédiate
$(\mathbf 0_{\mathrm I},6.2.1)$. Supposons donc c) vérifiée, et prouvons a).
Considérons les anneaux gradués $\operatorname{gr}_{\mathfrak J}^*(A)$,
$\operatorname{gr}_{\mathfrak J}^*(B)$ et le module gradué
$\operatorname{gr}_{\mathfrak J}^*(M)$ (à la fois sur
$\operatorname{gr}_{\mathfrak J}^*(A)$ et
$\operatorname{gr}_{\mathfrak J}^*(B)$) relatifs aux filtrations
$\mathfrak J$-préadique, ainsi que les applications canoniques surjectives
$(\mathbf 0_{\mathrm{III}},10.1.1.2)$
\begin{align*}
u&:\operatorname{gr}^0(B)\otimes_{\operatorname{gr}^0(A)}
   \operatorname{gr}^*(A)\longrightarrow\operatorname{gr}^*(B),\\
v&:\operatorname{gr}^0(M)\otimes_{\operatorname{gr}^0(A)}
   \operatorname{gr}^*(A)\longrightarrow\operatorname{gr}^*(M),\\
w&:\operatorname{gr}^0(M)\otimes_{\operatorname{gr}^0(B)}
   \operatorname{gr}^*(B)\longrightarrow\operatorname{gr}^*(M).
\end{align*}
Il est clair qu'on a un diagramme commutatif
\begin{equation}
\xymatrix{
\operatorname{gr}^0(M)\otimes_{\operatorname{gr}^0(A)}\operatorname{gr}^*(A)
 \ar[rr]^v \ar[dr]_{1\otimes u} && \operatorname{gr}^*(M) \\
& \operatorname{gr}^0(M)\otimes_{\operatorname{gr}^0(B)}
  \operatorname{gr}^*(B) \ar[ur]_w &
}
\label{IV.11.3.10.3-fr}
\tag{11.3.10.3}
\end{equation}
L'hypothèse c) entraîne que $v$ est bijective
$(\mathbf 0_{\mathrm{III}},10.2.1)$~; comme les deux autres applications du
diagramme sont surjectives, elles sont aussi bijectives. Mais comme en vertu de
l'hypothèse c), $M/\mathfrak JM$ est un $(B/\mathfrak JB)$-module plat, il
résulte de $(\mathbf 0_{\mathrm{III}},10.2.1)$ que $M$ est un $B$-module plat.

\noindent 2\textsuperscript{o} L'une ou l'autre des conditions (i), (iii)
implique que tout idéal maximal de $B$ contient $\mathfrak JB$. Il résulte donc
de 1\textsuperscript{o} et de la conjonction de c) et d) que $M$ est un
$B$-module \emph{fidèlement plat}, et par suite $\operatorname{gr}^0(M)$ un
$\operatorname{gr}^0(B)$-module \emph{fidèlement plat}
$(\mathbf 0_{\mathrm I},6.2.1)$. On a vu dans 1\textsuperscript{o} que
l'hypothèse c) entraîne que les trois applications du diagramme (11.3.10.3)
sont bijectives~; le fait que $\operatorname{gr}^0(M)$ soit un
$\operatorname{gr}^0(B)$-module fidèlement plat implique donc que $u$ est aussi
bijective $(\mathbf 0_{\mathrm I},6.4.1)$. D'autre part, les conditions (ii) et
(iii) impliquent que $B$ est un $A$-module idéalement séparé pour la filtration
$\mathfrak J$-préadique (Bourbaki, \emph{Alg. comm.}, chap.~III, \S~5,
n$^\circ$~4, prop.~2)~; on déduit donc encore de
$(\mathbf 0_{\mathrm{III}},10.2.1)$ que si la condition (i), ou la conjonction
de (ii) et (iii), est vérifiée, $B$ est un $A$-module plat.

\phantomsection\label{IV.11.3.10.4-fr}
\noindent\textbf{(11.3.10.4)} Le lemme (11.3.10.2) étant établi, on en déduit
(11.3.10.1) en prenant pour $\mathfrak J$ l'idéal maximal de $A$, et en
remarquant qu'alors les conditions (ii) et (iii) de (11.3.10.2) sont
satisfaites (Bourbaki, \emph{Alg. comm.}, chap.~III, \S~5, n$^\circ$~4,
prop.~2). Ceci achève donc aussi la démonstration de (11.3.7).
\end{proof}

\begin{corollary}[11.3.11]
\label{IV.11.3.11-fr}
Soient $g:X\to S$, $h:Y\to S$ deux morphismes localement de présentation
finie, $f:X\to Y$ un $S$-morphisme. Les conditions suivantes sont
équivalentes~:
\begin{enumerate}[label=\emph{\alph*)}]
\item $g$ est plat et $f_s:X_s\to Y_s$ est plat pour tout $s\in S$.
\item $h$ est plat en tous les points de $f(X)$ et $f$ est plat.
\end{enumerate}
\end{corollary}

Il suffit d'appliquer (11.3.10) pour $\mathcal F=\mathcal O_X$.

\begin{remark}[11.3.12]
\label{IV.11.3.12-fr}
Il serait intéressant de pouvoir donner de (11.3.2) et (11.3.10) des
démonstrations n'utilisant pas le passage à la limite inductive~; il suffirait
pour cela de démontrer le critère suivant~:

Soient $A$ un anneau, $B$ une $A$-algèbre de présentation finie, $M$ un
$B$-module de présentation finie, $\mathfrak J$ un

\oldpage[IV-3]{140}

idéal de $A$, $\mathfrak p$ un idéal premier de $B$ contenant $\mathfrak JB$.
Alors les deux conditions suivantes sont équivalentes~:
\begin{enumerate}[label=\emph{\alph*)}]
\item $M_\mathfrak p$ est un $A$-module plat.
\item $M_\mathfrak p/\mathfrak JM_\mathfrak p$ est un
$(A/\mathfrak J)$-module plat et
$\operatorname{Tor}_1^A(A/\mathfrak J,M_\mathfrak p)=0$.
\end{enumerate}
Lorsque $A$ est \emph{noethérien}, c'est une conséquence de
$(\mathbf 0_{\mathrm{III}},10.2.2)$ appliqué à la $A$-algèbre noethérienne
$B_\mathfrak p$~; peut-on en déduire l'énoncé général par un passage à la
limite inductive~?

Il convient de noter à ce propos qu'une telle généralisation n'est certainement
pas possible lorsqu'on remplace la condition b) ci-dessus par l'une des
conditions c), d) de $(\mathbf 0_{\mathrm{III}},10.2.1)$ où $M$ est remplacé
par $M_\mathfrak p$. Prenons par exemple pour $A$ un anneau local dont l'idéal
maximal $\mathfrak J$ est principal et tel que l'intersection
\[
\mathfrak R=\bigcap_{n\geq0}\mathfrak J^{n+1}
\]
ne soit pas réduite à $0$ (par exemple l'anneau des germes de fonctions
numériques indéfiniment différentiables au voisinage de $0$ dans $\mathbf R$).
Prenons $B=A$, $\mathfrak p=\mathfrak J$, et $M=A/\mathfrak N$, où
$\mathfrak N$ est un sous-module monogène $\neq0$ de $\mathfrak R$~; $M$ est
donc de présentation finie. Il est clair que $M$ n'est pas un $A$-module plat,
car étant de présentation finie, il serait libre (Bourbaki, \emph{Alg. comm.},
chap.~II, \S~3, n$^\circ$~2, cor.~2 de la prop.~5), ce qui est absurde puisque
$\mathfrak N\neq0$. Cependant,
\[
\mathfrak J^kM/\mathfrak J^{k+n}M\simeq A/\mathfrak J^n
\]
quels que soient $k$ et $n$ positifs, donc $M$ vérifie bien les conditions c)
et d) de $(\mathbf 0_{\mathrm{III}},10.2.1)$ puisque $A/\mathfrak J$ est un
corps.
\end{remark}
\begin{proposition}[11.3.13]
\label{IV.11.3.13-fr}
Soient $f:X\to Y$ un morphisme de préschémas, $x$ un point de $X$ tel que
$f$ soit plat au point $x$~; posons $y=f(x)$.
\begin{enumerate}[label=\emph{(\roman*)}]
\item Si $X$ est réduit (resp. normal) au point $x$, alors $Y$ est réduit
(resp. normal) au point $y$.
\item On suppose de plus que $f$ soit de présentation finie au point $x$.
Alors, si $Y$ est réduit (resp. normal) au point $y$ et si le morphisme $f$
est réduit (resp. normal) au point $x$ (6.8.1), $X$ est réduit (resp. normal)
au point $x$.
\end{enumerate}
\end{proposition}

\begin{proof}
La première assertion n'est mise que pour mémoire, ayant été démontrée dans
(2.1.13). Pour prouver (ii), on peut se borner au cas où $X=\operatorname{Spec}(B)$,
$Y=\operatorname{Spec}(A)$, où $A$ est un anneau local réduit (resp. intègre et
intégralement clos) et $B$ une $A$-algèbre de présentation finie. On peut alors
(8.9.1) écrire
\[
B=B_0\otimes_{A_0}A,
\]
où $A_0$ est une sous-$\mathbf Z$-algèbre de type fini de $A$ et $B_0$ une
$A_0$-algèbre de type fini. Soit $(C_\alpha)$ la famille filtrante croissante
des sous-$A_0$-algèbres de type fini de $A$, qui sont donc des
$\mathbf Z$-algèbres de type fini~; on a $A=\varinjlim C_\alpha$. Distinguons
maintenant les deux cas~:

\noindent 1\textsuperscript{o} Supposons $A$ réduit et $f$ réduit au point
$x$. Si $\mathfrak m$ est l'idéal maximal de $A$, soit
$\mathfrak p_\alpha$ l'idéal premier $\mathfrak m\cap C_\alpha$, et posons
$A'_\alpha=(C_\alpha)_{\mathfrak p_\alpha}$, de sorte que l'on a aussi
$A=\varinjlim A'_\alpha$ (5.13.3). Posons
$Y_\alpha=\operatorname{Spec}(A'_\alpha)$,
$X_\alpha=\operatorname{Spec}(B_0\otimes_{A_0}A'_\alpha)$ et soient
$f_\alpha:X_\alpha\to Y_\alpha$, $x_\alpha$ (resp. $y_\alpha$) la projection
de $x$ (resp. $y$) dans $X_\alpha$ (resp. $Y_\alpha$). Puisque $f$ est réduit
au point $x$, il existe $\alpha_0$ tel que $f_\alpha$ soit réduit au point
$x_\alpha$ pour $\alpha\geq\alpha_0$ ((6.7.8) et (11.2.6))~; comme
$Y_\alpha$ et $X_\alpha$ sont noethériens et que $A'_\alpha$ est réduit
(puisqu'il en est ainsi de $C_\alpha$), on déduit de (3.3.5) que $X_\alpha$
est réduit au point $x_\alpha$. Mais puisque
$B=\varinjlim(B_0\otimes_{A_0}A'_\alpha)$, on a aussi
$\mathcal O_{X,x}=\varinjlim\mathcal O_{X_\alpha,x_\alpha}$ (5.13.3) et par
suite $\mathcal O_{X,x}$ est réduit (5.13.2).

\noindent 2\textsuperscript{o} Supposons $A$ intégralement clos et $f$ normal
au point $x$. Comme $C_\alpha$ est universellement japonais (7.7.4), sa
clôture intégrale $C'_\alpha$ est une $C_\alpha$-algèbre finie, évidemment
contenue dans $A$. Soit $\mathfrak p'_\alpha$ l'idéal premier
$\mathfrak m\cap C'_\alpha$, et posons
$A''_\alpha=(C'_\alpha)_{\mathfrak p'_\alpha}$, de sorte que $A''_\alpha$
est un anneau local noethérien intègre et intégralement clos, et que l'on a
$A=\varinjlim A''_\alpha$ (5.13.3). Posons
$Y'_\alpha=\operatorname{Spec}(A''_\alpha)$,
$X'_\alpha=\operatorname{Spec}(B_0\otimes_{A_0}A''_\alpha)$ et soient
$f'_\alpha:X'_\alpha\to Y'_\alpha$, $x'_\alpha$ (resp. $y'_\alpha$) la
projection de $x$ (resp. $y$) dans $X'_\alpha$ (resp. $Y'_\alpha$). Puisque
$f$ est normal au point $x$, il existe $\alpha_0$ tel que, pour
$\alpha\geq\alpha_0$, $f'_\alpha$ soit normal au point $x'_\alpha$ ((6.7.8)
et (11.2.6))~;

\oldpage[IV-3]{141}

comme $X'_\alpha$ et $Y'_\alpha$ sont noethériens et que $A''_\alpha$ est
intégralement clos, on déduit de (6.5.4) que $X'_\alpha$ est normal au point
$x'_\alpha$. Mais les morphismes
$\operatorname{Spec}(A''_\beta)\to\operatorname{Spec}(A''_\alpha)$ pour
$\alpha\leq\beta$ sont dominants, donc, comme en vertu de (11.3.1) on peut
supposer que les $f'_\alpha$ sont plats pour $\alpha\geq\alpha_0$, il résulte
de (2.3.7) que toute composante irréductible de $X'_\beta$ domine une
composante irréductible de $X'_\alpha$. On conclut alors de (5.13.4), appliqué
aux préschémas
\[
\operatorname{Spec}(\mathcal O_{X_0,x_0}\otimes_{A_0}A''_\alpha),
\]
que $X$ est normal au point $x$.
\end{proof}

\begin{corollary}[11.3.14]
\label{IV.11.3.14-fr}
Soient $f:X\to Y$ un morphisme localement de présentation finie, $x$ un point
de $X$, $y=f(x)$. Si $Y$ est géométriquement unibranche au point $y$ (6.15.1)
et si le morphisme $f$ est normal au point $x$ (6.8.1), alors $X$ est
géométriquement unibranche au point $x$.
\end{corollary}

\begin{proof}
On peut évidemment se borner au cas où $Y=\operatorname{Spec}(A)$, $A$ étant
un anneau local intègre géométriquement unibranche, $y$ étant le point fermé de
$Y$. Soit $A'$ la clôture intégrale de $A$~; posons
$Y'=\operatorname{Spec}(A')$ et soit $y'$ le point fermé de $Y'$, de sorte que
le morphisme $g:Y'\to Y$ est radiciel au point $y$ (6.15.3), entier et
birationnel. Si l'on pose $X'=X\times_Y Y'$, et si $f':X'\to Y'$ et
$g':X'\to X$ sont les projections, $g'$ est entier et radiciel au point $x$
(6.15.3.1). Soit donc $x'$ l'unique point de $g'^{-1}(x)$, qui est au-dessus
de $y'$. Le morphisme $f'$ est de présentation finie et normal au point $x'$
(6.7.8), et par suite l'anneau local $\mathcal O_{X',x'}$ est intègre et
intégralement clos (11.3.13). D'autre part, on peut supposer $f$ plat (11.3.1),
donc $g'$ est un morphisme birationnel (6.15.4.1)~; par suite
$\operatorname{Spec}(\mathcal O_{X,x})$ est irréductible, et comme
$\mathcal O_{X,x}$ est réduit (11.3.13) il est intègre et géométriquement
unibranche en vertu de (6.15.5).
\end{proof}

\begin{proposition}[11.3.15]
\label{IV.11.3.15-fr}
Soient $A$ un anneau, $B$ une $A$-algèbre de présentation finie, $M$ un
$B$-module de présentation finie, qui est un $A$-module plat. Alors il existe
une suite finie $(f_i)_{1\leq i\leq n}$ d'éléments de $A$, tels que l'idéal
engendré par les $f_i$ soit égal à $A$, et que, pour $0\leq i<n$,
\[
M_{f_{i+1}}\Big/\Big(\sum_{j=1}^{i}f_jM_{f_{i+1}}\Big)
\]
soit un
$A_{f_{i+1}}\big/(\sum_{j=1}^{i}f_jA_{f_{i+1}})$-module libre.
\end{proposition}

On peut encore dire que les $D(f_i)$ forment un recouvrement ouvert de
$X=\operatorname{Spec}(A)$, et que si l'on pose $Z_1=D(f_1)$, puis, par
récurrence,
\[
Z_{i+1}=D(f_{i+1})\cap V(f_1A+\cdots+f_iA),
\]
les $Z_i$ forment une \emph{partition} de $X$ en ensembles localement fermés,
$A_{f_{i+1}}/(\sum_{j=1}^{i}f_jA_{f_{i+1}})$ étant l'anneau d'un sous-schéma
affine de $X$ ayant $Z_{i+1}$ pour espace sous-jacent.

\begin{proof}
Pour démontrer la proposition, on peut d'abord, en vertu de (11.2.7), supposer
qu'il existe un sous-anneau noethérien $A_0$ de $A$, une $A_0$-algèbre de type
fini $B_0$ et un $B_0$-module de type fini $M_0$, plat sur $A_0$ et tels que
$B=B_0\otimes_{A_0}A$, $M=M_0\otimes_{A_0}A$~; il est clair qu'il suffira de
prouver la proposition pour $A_0$, $B_0$ et $M_0$, car si les éléments
$f_i\in A_0$ vérifient dans ce cas les conditions de l'énoncé, ils les
vérifieront aussi pour $A$, $B$, $M$, puisque
\[
M_{f_{i+1}}\Big/\Big(\sum_{j=1}^{i}f_jM_{f_{i+1}}\Big)
=\left((M_0)_{f_{i+1}}\Big/\Big(\sum_{j=1}^{i}f_j(M_0)_{f_{i+1}}\Big)\right)
\otimes_{A_0}A
\]
(Bourbaki, \emph{Alg. comm.}, chap.~II, \S~2, n$^\circ$~7, prop.~18). On
peut donc se borner désormais au cas où $A$ est noethérien.

Notons maintenant que si $C$ est un anneau noethérien, $\mathfrak N$ son
nilradical et $P$ un $C$-module plat, alors il résulte de
$(\mathbf 0_{\mathrm{III}},10.1.2)$ que pour que $P$ soit un $C$-module libre,

\oldpage[IV-3]{142}

il faut et il suffit que $P\otimes_C(C/\mathfrak N)$ soit un
$(C/\mathfrak N)$-module libre. Notons d'autre part que si $C$ est un anneau
noethérien réduit, il existe $g\in C$ tel que $C_g$ soit un anneau intègre.
Utilisons maintenant le lemme (6.9.2)~: on peut définir par récurrence une suite
$(f_i)_{i\geq1}$ d'éléments de $A$ de la façon suivante~:
\begin{enumerate}
\item[$1^\circ$] $f_1$ est tel que $A_1=(A_{\mathrm{red}})_{f_1}$ soit
intègre et $M\otimes_A A_1$ un $A_1$-module libre~;
\item[$2^\circ$] si l'idéal $\mathfrak J_i$ engendré par $f_1,\ldots,f_i$
est $A$, $f_{i+1}=f_i$~;
\item[$3^\circ$] si au contraire $\mathfrak J_i\neq A$, $f_{i+1}$ est un
élément n'appartenant pas à $\mathfrak J_i$ tel que
$A_{i+1}=((A/\mathfrak J_i)_{\mathrm{red}})_{f_{i+1}}$ soit intègre et
$M\otimes_A A_{i+1}$ un $A_{i+1}$-module libre.
\end{enumerate}
Comme $A$ est noethérien, la suite croissante $(\mathfrak J_i)$ est
stationnaire, donc il existe $n$ tel que $f_1,\ldots,f_n$ engendrent l'idéal
$A$, et les $f_i$ répondent à la question puisque
\[
((A/\mathfrak J_i)_{\mathrm{red}})_{f_{i+1}}
=((A/\mathfrak J_i)_{f_{i+1}})_{\mathrm{red}}.
\]
\end{proof}

\begin{proposition}[11.3.16]
\label{IV.11.3.16-fr}
Soient $f:X\to Y$ un morphisme fidèlement plat de présentation finie,
$g:Y\to Z$ un morphisme tel que $g\circ f:X\to Z$ soit un morphisme de type
fini (resp. de présentation finie). Alors $g$ est un morphisme de type fini
(resp. de présentation finie).
\end{proposition}

\begin{proof}
Comme $f$ est surjectif et $g\circ f$ quasi-compact, $g$ est quasi-compact
(1.1.3). Montrons d'abord que si $g\circ f$ est de présentation finie, $g$ est
quasi-séparé. En effet, considérons un ouvert affine $W\subset Z$, et soit
$(V_i)$ un recouvrement affine fini de $g^{-1}(W)$~; il s'agit de voir que les
$V_i\cap V_j$ sont quasi-compacts (1.2.6 et 1.2.7). Pour chaque $i$, soit
$(U_{ih})$ un recouvrement ouvert affine fini de $f^{-1}(V_i)$~; comme $f$ est
de présentation finie, les $U_{ih}\cap U_{jk}$ sont tous quasi-compacts~; or,
puisque $f$ est surjectif, $V_i\cap V_j$ est réunion des images
$f(U_{ih}\cap U_{jk})$ pour $h$, $k$ variables, donc est quasi-compact puisque
$f$ est continue.

La question est donc locale sur $Y$ et on peut supposer que
$Y=\operatorname{Spec}(B)$ et $Z=\operatorname{Spec}(A)$ sont affines, $X$
étant réunion finie d'ouverts affines $X_i=\operatorname{Spec}(C_i)$, où les
$C_i$ sont des $B$-algèbres de type fini (resp. de présentation finie). Si
$X'$ est le préschéma somme des $X_i$, $p:X'\to X$ le morphisme qui coïncide
avec l'injection canonique dans chaque $X_i$, $p$ est un morphisme fidèlement
plat de présentation finie (1.6.5), donc $g\circ f\circ p$ est un morphisme de
type fini (resp. de présentation finie) et $f\circ p$ un morphisme fidèlement
plat de présentation finie. On peut par suite supposer aussi que
$X=\operatorname{Spec}(C)$ est affine, et on est donc ramené à prouver le
corollaire suivant.
\end{proof}

\begin{corollary}[11.3.17]
\label{IV.11.3.17-fr}
Soient $A$ un anneau, $B$ une $A$-algèbre, $C$ une $B$-algèbre de présentation
finie et qui soit un $B$-module fidèlement plat. Alors, si $C$ est une
$A$-algèbre de type fini (resp. de présentation finie), $B$ est une
$A$-algèbre de type fini (resp. de présentation finie).
\end{corollary}

\begin{proof}
Supposons d'abord que $g\circ f$ soit de type fini. Soit $(B_\lambda)$ la
famille filtrante croissante des sous-$A$-algèbres de type fini de $B$~; en
vertu de (8.8.2), il existe un indice $\alpha$ tel que
$C=C_\alpha\otimes_{B_\alpha}B$, où $C_\alpha$ est une
$B_\alpha$-algèbre de présentation finie~; en outre (8.10.5 et 11.2.6) on peut
supposer que $C_\alpha$ est un $B_\alpha$-module fidèlement plat. Pour
$\lambda\geq\alpha$, on a donc
$C=C_\lambda\otimes_{B_\lambda}B$, où $C_\lambda$ est un
$B_\lambda$-module fidèlement plat~; comme l'application $B_\lambda\to B$ est
injective, on en déduit qu'il en est de même de $C_\lambda\to C$. En outre,
comme $C=\varinjlim C_\lambda$, tout élément de $C$ appartient à un
$C_\lambda$, et par suite il existe $\lambda$ tel que l'application
$C_\lambda\to C$

\oldpage[IV-3]{143}

soit bijective, puisque $C$ est une $B$-algèbre de type fini. Mais alors
l'application $B_\lambda\to B$ est bijective par fidèle platitude, donc
$B=B_\lambda$ est une $A$-algèbre de type fini.

Supposons maintenant que $g\circ f$ soit de présentation finie, $C$ étant donc
une $A$-algèbre de présentation finie. D'après la première partie du
raisonnement, on a
\[
B=A[T_1,\ldots,T_n]/\widetilde{\mathfrak J}
\]
pour un idéal $\widetilde{\mathfrak J}$~; soit
$(\widetilde{\mathfrak J}_\lambda)$ la famille filtrante des idéaux de type
fini de $A[T_1,\ldots,T_n]$ contenus dans $\widetilde{\mathfrak J}$, de sorte
que $\widetilde{\mathfrak J}=\varinjlim\widetilde{\mathfrak J}_\lambda$ et
$B=\varinjlim B_\lambda$, avec
$B_\lambda=A[T_1,\ldots,T_n]/\widetilde{\mathfrak J}_\lambda$. Appliquant
comme ci-dessus (8.8.2), (8.10.5) et (11.2.6), on a
$C=C_\alpha\otimes_{B_\alpha}B$, où $C_\alpha$ est une
$B_\alpha$-algèbre de présentation finie et un $B_\alpha$-module fidèlement
plat~; on posera encore $C_\lambda=C_\alpha\otimes_{B_\alpha}B_\lambda$ pour
$\lambda\geq\alpha$, de sorte que
\[
C_\lambda=C_\alpha/\big(C_\alpha\otimes_{B_\alpha}
(\widetilde{\mathfrak J}_\lambda/\widetilde{\mathfrak J}_\alpha)\big)
\]
par platitude, et de même
\[
C=C_\alpha/\big(C_\alpha\otimes_{B_\alpha}
(\widetilde{\mathfrak J}/\widetilde{\mathfrak J}_\alpha)\big).
\]
Comme par hypothèse $C$ est une $A$-algèbre de présentation finie ainsi que
$C_\alpha$, $C_\alpha\otimes_{B_\alpha}
(\widetilde{\mathfrak J}/\widetilde{\mathfrak J}_\alpha)$ est un idéal de
type fini de $C_\alpha$ (1.4.4), et les
$C_\alpha\otimes_{B_\alpha}
(\widetilde{\mathfrak J}_\lambda/\widetilde{\mathfrak J}_\alpha)$
s'identifient à des idéaux de type fini de $C_\alpha$ dont
$C_\alpha\otimes_{B_\alpha}
(\widetilde{\mathfrak J}/\widetilde{\mathfrak J}_\alpha)$ est la réunion. Il
existe par suite un $\lambda\geq\alpha$ tel que
\[
C_\alpha\otimes_{B_\alpha}
(\widetilde{\mathfrak J}/\widetilde{\mathfrak J}_\alpha)
=C_\alpha\otimes_{B_\alpha}
(\widetilde{\mathfrak J}_\lambda/\widetilde{\mathfrak J}_\alpha),
\]
d'où $\widetilde{\mathfrak J}=\widetilde{\mathfrak J}_\lambda$ par fidèle
platitude, ce qui prouve que $B$ est une $A$-algèbre de présentation finie.
\end{proof}

\subsection{Descente de la platitude par des morphismes quelconques~: cas d'un préschéma de base artinien}
\label{subsection:IV.11.4-fr}

\begin{theorem}[11.4.1]
\label{IV.11.4.1-fr}
Soient $A$ un anneau, $\mathfrak J$ un idéal nilpotent de $A$,
$u_\lambda:A\to B_\lambda$ $(\lambda\in L)$ une famille d'homomorphismes
d'anneaux telle que l'intersection des noyaux des $u_\lambda$ soit réduite à
$0$. Soit $M$ un $A$-module tel que pour tout $\lambda\in L$,
$M\otimes_A B_\lambda$ soit un $B_\lambda$-module libre et que
$M\otimes_A(A/\mathfrak J)$ soit un $(A/\mathfrak J)$-module libre. Alors $M$
est un $A$-module libre. Si de plus l'ensemble d'indices $L$ est fini, on peut
remplacer partout «~libre~» par «~plat~» dans l'énoncé précédent.
\end{theorem}

\begin{proof}
Dans les deux cas il suffira de prouver que $M$ est un $A$-module \emph{plat}~:
en effet, lorsque $M\otimes_A(A/\mathfrak J)$ est un
$(A/\mathfrak J)$-module libre, il en résultera que $M$ est un $A$-module libre
en vertu de $(\mathbf 0_{\mathrm{III}},10.1.2)$.

Nous utiliserons le lemme suivant, qui généralise
$(\mathbf 0_{\mathrm{III}},10.2.1)$~:

\begin{lemma}[11.4.1.1]
\label{IV.11.4.1.1-fr}
Soit $A$ un anneau muni d'une filtration finie
$(\mathfrak J_n)_{0\leq n\leq N+1}$, avec $\mathfrak J_0=A$,
$\mathfrak J_{N+1}=0$. Soit $M$ un $A$-module muni de la filtration
$(\mathfrak J_nM)_{0\leq n\leq N+1}$, et désignons par
$\operatorname{gr}(A)$ et $\operatorname{gr}(M)$ l'anneau et le module gradués
correspondants. Supposons que $M\otimes_A(A/\mathfrak J_1)$ soit un
$(A/\mathfrak J_1)$-module plat et que l'homomorphisme canonique
\begin{equation}
\operatorname{gr}_0(M)\otimes_{\operatorname{gr}_0(A)}\operatorname{gr}(A)
\longrightarrow\operatorname{gr}(M)
\label{IV.11.4.1.2-fr}
\tag{11.4.1.2}
\end{equation}
soit injectif. Alors $M$ est un $A$-module plat.
\end{lemma}

L'homomorphisme canonique (11.4.1.2) se définit de la même façon que celui de
$(\mathbf 0_{\mathrm{III}},10.1.1)$ comme étant en degré $n$ l'homomorphisme
\[
(M/\mathfrak J_1M)\otimes(\mathfrak J_n/\mathfrak J_{n+1})
=(M\otimes\mathfrak J_n)/
(\operatorname{Im}(M\otimes\mathfrak J_{n+1})
+\operatorname{Im}(\mathfrak J_1M\otimes\mathfrak J_n))
\longrightarrow\mathfrak J_nM/\mathfrak J_{n+1}M
\]
provenant de l'homomorphisme canonique $M\otimes\mathfrak J_n\to\mathfrak J_nM$
par passage aux quotients. Le lemme se démontre par récurrence sur $N$,
puisqu'il n'y a rien à démontrer pour $N=0$.

\oldpage[IV-3]{144}

Les conditions sur $M$ entraînent, en vertu de l'hypothèse de récurrence, que
$M\otimes_A(A/\mathfrak J_N)$ est un $(A/\mathfrak J_N)$-module plat. Notons
maintenant que l'on a $\mathfrak J_N^2\subset\mathfrak J_{N+1}=0$ si $N\geq1$,
et
\[
(M/\mathfrak J_1M)\otimes(\mathfrak J_N/\mathfrak J_{N+1})
=(M/\mathfrak J_NM)\otimes(\mathfrak J_N/\mathfrak J_{N+1})
=(M/\mathfrak J_NM)\otimes\mathfrak J_N\,;
\]
donc l'homomorphisme canonique
\[
(M/\mathfrak J_NM)\otimes\mathfrak J_N
\longrightarrow\mathfrak J_NM=\mathfrak J_NM/\mathfrak J_N^2M
\]
est injectif. Appliquant $(\mathbf 0_{\mathrm{III}},10.2.1)$ à la filtration
$\mathfrak J_N$-préadique, on en conclut bien que $M$ est un $A$-module plat.

Pour appliquer ce lemme à (11.4.1), nous désignerons par $\mathfrak J_n$
l'idéal de $A$ intersection des images réciproques
$u_\lambda^{-1}(\mathfrak J^nB_\lambda)$~; il est immédiat que l'on a
$\mathfrak J_0=A$, $\mathfrak J_m\mathfrak J_n\subset\mathfrak J_{m+n}$ pour
$m\geq0$, $n\geq0$~; en outre, si $\mathfrak J^{N+1}=0$, on a aussi
$\mathfrak J_{N+1}=0$ puisque l'intersection des noyaux des $u_\lambda$ est
réduite à $0$.

Munissons $A$ de la filtration $(\mathfrak J_n)$, $M$ de la filtration
$(\mathfrak J_nM)$, et, pour chaque $\lambda$, munissons $B_\lambda$ et
$N_\lambda=M\otimes_A B_\lambda$ des filtrations $\mathfrak J$-préadiques~;
considérons pour chaque $\lambda$ le diagramme commutatif
\[
\xymatrix{
\operatorname{gr}^0(M)\otimes_{\operatorname{gr}^0(A)}\operatorname{gr}(A)
 \ar[r]^{\varphi} \ar[d]_{f_\lambda} & \operatorname{gr}(M) \ar[d]^{g_\lambda} \\
\operatorname{gr}_{\mathfrak J}^0(N_\lambda)
 \otimes_{\operatorname{gr}_{\mathfrak J}^0(B_\lambda)}
 \operatorname{gr}_{\mathfrak J}(B_\lambda)
 \ar[r]_{\varphi_\lambda} & \operatorname{gr}_{\mathfrak J}(N_\lambda)
}
\]
où les flèches horizontales sont les homomorphismes canoniques (11.4.1.2) et
les flèches verticales sont déduites des homomorphismes canoniques
$A\to B_\lambda$ et $M\to N_\lambda$. L'hypothèse que $N_\lambda$ est un
$B_\lambda$-module plat entraîne que $\varphi_\lambda$ est injective
$(\mathbf 0_{\mathrm{III}},10.2.1)$, donc
$\operatorname{Ker}(\varphi)\subset\operatorname{Ker}(f_\lambda)$. Posant
$A_0=A/\mathfrak J_1$, $M_0=M\otimes_A A_0$, notons que
\[
\operatorname{gr}_{\mathfrak J}^0(N_\lambda)
=N_\lambda/\mathfrak JN_\lambda
=(M/\mathfrak JM)\otimes_A B_\lambda
=(M/\mathfrak J_1M)\otimes_A(B_\lambda/\mathfrak JB_\lambda),
\]
car ce dernier produit tensoriel est égal à
\[
(M\otimes_A B_\lambda)\Big/
\big(\operatorname{Im}(M\otimes_A\mathfrak JB_\lambda)
 +\operatorname{Im}(\mathfrak J_1M\otimes_A B_\lambda)\big),
\]
et l'on a
$\operatorname{Im}(\mathfrak J_1M\otimes_A B_\lambda)
=\operatorname{Im}(M\otimes_A\mathfrak J_1B_\lambda)
\subset\operatorname{Im}(M\otimes_A\mathfrak JB_\lambda)$ puisque
$\mathfrak J_1B_\lambda\subset\mathfrak JB_\lambda$ par définition~; enfin,
la relation $\mathfrak J_1B_\lambda\subset\mathfrak JB_\lambda$ montre que l'on
a aussi
\[
(M/\mathfrak J_1M)\otimes_A(B_\lambda/\mathfrak JB_\lambda)
=(M/\mathfrak J_1M)\otimes_{A/\mathfrak J_1}
 (B_\lambda/\mathfrak JB_\lambda),
\]
si bien que finalement on a
\[
\operatorname{gr}_{\mathfrak J}^0(N_\lambda)
=M_0\otimes_{A_0}(B_\lambda/\mathfrak JB_\lambda)
\]
et par suite
\[
\operatorname{gr}_{\mathfrak J}^0(N_\lambda)
\otimes_{\operatorname{gr}_{\mathfrak J}^0(B_\lambda)}
\operatorname{gr}_{\mathfrak J}(B_\lambda)
=M_0\otimes_{A_0}\operatorname{gr}_{\mathfrak J}(B_\lambda).
\]
L'homomorphisme $f_\lambda$ peut donc s'écrire
\[
1\otimes\operatorname{gr}(u_\lambda):
M_0\otimes_{A_0}\operatorname{gr}(A)
\longrightarrow M_0\otimes_{A_0}\operatorname{gr}_{\mathfrak J}(B_\lambda),
\]

\oldpage[IV-3]{145}

et comme $M_0$ est par hypothèse un $A_0$-module \emph{plat}, le noyau de
$f_\lambda$ est égal à $M_0\otimes_{A_0}R_\lambda$, où $R_\lambda$ est le
noyau de
$\operatorname{gr}(u_\lambda):\operatorname{gr}(A)
\to\operatorname{gr}_{\mathfrak J}(B_\lambda)$. Tout revient donc à prouver
que
\[
\bigcap_{\lambda\in L}(M_0\otimes_{A_0}R_\lambda)=0.
\]
Or, par définition des $\mathfrak J_n$, l'intersection des noyaux des
homomorphismes
\[
\mathfrak J_nA/\mathfrak J_{n+1}A
\longrightarrow\mathfrak J^nB_\lambda/\mathfrak J^{n+1}B_\lambda,
\]
lorsque $\lambda$ parcourt $L$, est réduite à $0$, autrement dit
$\bigcap_{\lambda\in L}R_\lambda=0$. Cela étant, supposons d'abord que $M_0$
soit un $A_0$-module libre~; prenant une base de $M_0$, on voit aussitôt que
l'on a
\[
M_0\otimes_{A_0}\Big(\bigcap_{\lambda\in L}R_\lambda\Big)
=\bigcap_{\lambda\in L}(M_0\otimes_{A_0}R_\lambda),
\]
d'où la proposition dans ce cas. Lorsque $L$ est fini, la formule précédente
est encore vraie sous la seule hypothèse que $M_0$ est un $A_0$-module plat
$(\mathbf 0_{\mathrm I},6.1.3)$, ce qui termine la démonstration.
\end{proof}

\begin{remark}[11.4.2]
\label{IV.11.4.2-fr}
La conclusion de (11.4.1) peut être inexacte si $L$ est infini et si l'on
suppose seulement que $M\otimes_A(A/\mathfrak J)$ est un
$(A/\mathfrak J)$-module plat. Par exemple, soient $V$ un anneau de valuation
discrète, $K$ son corps des fractions, et soit $A=V[T]/(T^2)$ ($T$
indéterminée)~; prenons pour $\mathfrak J$ l'image de $(T)$ dans $A$, de sorte
que $A/\mathfrak J=V$, et prenons $M=K$, qui est un
$(A/\mathfrak J)$-module, donc égal à $M\otimes_A(A/\mathfrak J)$~; en outre
$M$ est un $(A/\mathfrak J)$-module plat, mais non un $A$-module plat, car
puisque $\mathfrak J$ est nilpotent, il résulterait de
$(\mathbf 0_{\mathrm{III}},10.1.2)$ que $M$ serait un $A$-module libre, ce qui
est absurde puisque $\mathfrak JK=0$. Considérons d'autre part l'idéal maximal
$\mathfrak m$ de l'anneau local noethérien $A$~; on a $\mathfrak mK=K$, donc
$M\otimes_A(A/\mathfrak m^n)=0$ quel que soit l'entier $n$~; les
$(A/\mathfrak m^n)$-modules $M\otimes_A(A/\mathfrak m^n)$ sont donc plats pour
tout $n$, et l'intersection des $\mathfrak m^n$ est réduite à $0$.
\end{remark}

\begin{corollary}[11.4.3]
\label{IV.11.4.3-fr}
Soient $A$ un anneau semi-local dont le radical $\mathfrak J$ est nilpotent
(par exemple un anneau artinien), $u_\lambda:A\to B_\lambda$ $(\lambda\in L)$
une famille d'homomorphismes d'anneaux telle que l'intersection des noyaux des
$u_\lambda$ soit réduite à $0$. Pour qu'un $A$-module $M$ soit plat, il faut et
il suffit que pour tout $\lambda\in L$, $M\otimes_A B_\lambda$ soit un
$B_\lambda$-module plat.
\end{corollary}

\begin{proof}
Comme $A/\mathfrak J$ est composé direct d'un nombre fini de corps (Bourbaki,
\emph{Alg. comm.}, chap.~II, \S~3, n$^\circ$~5, prop.~16) et que $\mathfrak J$
est nilpotent, $A$ est composé direct d'un nombre fini d'anneaux locaux $A_i$
dont le radical est nilpotent (\emph{loc. cit.}, \S~4, n$^\circ$~3, cor.~1 de
la prop.~15) et $M$ est par suite somme directe de $A_i$-modules $M_i$, chaque
$M_i$ étant annulé par les $A_j$ d'indice $j\neq i$~; pour que $M$ soit un
$A$-module plat, il faut et il suffit que chacun des $M_i$ soit un $A_i$-module
plat~; d'ailleurs l'intersection des noyaux des homomorphismes
\[
A_i\longrightarrow A\xrightarrow{u_\lambda}B_\lambda
\]
est réduite à $0$, et $M_i\otimes_{A_i}B_\lambda$ est facteur direct de
$M\otimes_A B_\lambda$. On peut donc se borner au cas où $A$ est en outre
local. Alors $A/\mathfrak J$ est un corps, donc
$M\otimes_A(A/\mathfrak J)$ est un $(A/\mathfrak J)$-module libre, et il
suffit de voir que pour tout $\lambda$, $M\otimes_A B_\lambda$ est un
$B_\lambda$-module libre, en vertu de (11.4.1). Mais si l'on pose
$\mathfrak J_\lambda=\mathfrak JB_\lambda$,
\[
(M\otimes_A B_\lambda)\otimes_{B_\lambda}(B_\lambda/\mathfrak J_\lambda)
=(M\otimes_A(A/\mathfrak J))\otimes_{A/\mathfrak J}
 (B_\lambda/\mathfrak J_\lambda)
\]
est un $(B_\lambda/\mathfrak J_\lambda)$-module libre, et
$\mathfrak J_\lambda$ est nilpotent. La conclusion résulte donc de l'hypothèse
que $M\otimes_A B_\lambda$ est un $B_\lambda$-module plat et de
$(\mathbf 0_{\mathrm{III}},10.1.2)$.
\end{proof}

\begin{corollary}[11.4.4]
\label{IV.11.4.4-fr}
Soient $A$ un anneau, $M$ un $A$-module~; on suppose qu'il existe un idéal
nilpotent $\mathfrak J$ de $A$ tel que $M\otimes_A(A/\mathfrak J)$ soit un
$(A/\mathfrak J)$-module libre. Alors l'ensemble des idéaux $\mathfrak K$ de
$A$ tels que $M\otimes_A(A/\mathfrak K)$ soit un $(A/\mathfrak K)$-module
libre admet un plus petit élément $\mathfrak K_0$ (qui est aussi le plus petit
des idéaux $\mathfrak K$ tels que $M\otimes_A(A/\mathfrak K)$ soit un
$(A/\mathfrak K)$-module plat).
\end{corollary}

\oldpage[IV-3]{146}

\begin{proof}
Pour qu'un homomorphisme $u:A\to A'$ soit tel que $M\otimes_A A'$ soit un
$A'$-module libre (resp. un $A'$-module plat), il faut et il suffit que $u$ se
factorise en $A\to A/\mathfrak K_0\to A'$ (ou encore que
$\mathfrak K_0A'=0$).

Le fait que l'intersection $\mathfrak K_0$ des idéaux $\mathfrak K$ pour
lesquels $M\otimes_A(A/\mathfrak K)$ est un $(A/\mathfrak K)$-module libre
soit le plus petit de ces idéaux résulte de (11.4.1) appliqué à l'anneau
$A/\mathfrak K_0$, à son idéal nilpotent $\mathfrak J/\mathfrak K_0$ et aux
homomorphismes $A/\mathfrak K_0\to A/\mathfrak K$, dont les noyaux ont $0$
pour intersection. Si $A'$ est une $(A/\mathfrak K_0)$-algèbre, on a
\[
M\otimes_A A'=
(M\otimes_A(A/\mathfrak K_0))\otimes_{A/\mathfrak K_0}A',
\]
donc $M\otimes_A A'$ est un $A'$-module libre. Réciproquement, si $A'$ est
une $A$-algèbre telle que $M\otimes_A A'$ soit un $A'$-module libre et si
$\mathfrak K$ est le noyau de l'homomorphisme $A\to A'$, il résulte de
(11.4.1) appliqué à l'anneau $A/\mathfrak K$, au
$(A/\mathfrak K)$-module $M\otimes_A(A/\mathfrak K)$, à l'idéal nilpotent
$(\mathfrak J+\mathfrak K)/\mathfrak K$ et à l'homomorphisme injectif
$A/\mathfrak K\to A'$, que $M\otimes_A(A/\mathfrak K)$ est un
$(A/\mathfrak K)$-module libre, donc que $\mathfrak K\supset\mathfrak K_0$.
Le fait que l'on puisse remplacer «~libre~» par «~plat~» dans ce qui précède
(en conservant naturellement l'hypothèse que $M\otimes_A(A/\mathfrak J)$ est
un $(A/\mathfrak J)$-module libre) résulte de
$(\mathbf 0_{\mathrm{III}},10.1.2)$, comme on l'a vu au début de la
démonstration de (11.4.1).
\end{proof}

\begin{proposition}[11.4.5]
\label{IV.11.4.5-fr}
Soient $Y$ un préschéma irréductible, $f:X\to Y$ un morphisme de présentation
finie, $\mathcal F$ un $\mathcal O_X$-Module de présentation finie. Il existe
alors un ouvert non vide $U$ dans $Y$ et un sous-schéma fermé $Z$ de $U$, de
présentation finie sur $U$, ayant la propriété suivante~: pour tout changement
de base $U'\to U$, si l'on pose $X'=X\times_YU'$, $f'=f_{(U')}$ et
$\mathcal F'=\mathcal F\otimes_{\mathcal O_Y}\mathcal O_{U'}$, alors, pour que
$\mathcal F'$ soit $f'$-plat, il faut et il suffit que le morphisme $U'\to U$
se factorise en $U'\to Z\to U$. Un tel schéma $Z$ a même espace sous-jacent que
$U$. Supposons de plus $Y$ affine, et soit $(W_i)$ un recouvrement fini de $X$
par des ouverts affines~; alors, on peut supposer $U$ choisi de sorte que, si
$U'=\operatorname{Spec}(A')$ est un schéma affine et $U'\to U$ un morphisme qui
se factorise en $U'\to Z\to U$, chacun des
$\Gamma(W_i\times_YU',\mathcal F')$ est un $A'$-module libre.
\end{proposition}

\begin{proof}
On peut évidemment se borner au cas où $Y=\operatorname{Spec}(A)$ est affine.
Utilisant (8.9.1), il y a un sous-anneau noethérien $A_1$ de $A$, un morphisme
de type fini $f_1:X_1\to Y_1=\operatorname{Spec}(A_1)$ et un
$\mathcal O_{X_1}$-Module cohérent $\mathcal F_1$ tels que $f=(f_1)_{(Y)}$ et
$\mathcal F=\mathcal F_1\otimes_{\mathcal O_{Y_1}}\mathcal O_Y$~; on peut en
outre supposer que les $W_i$ sont images réciproques d'ouverts affines de
$X_1$. Notons en outre que $Y_1$ est irréductible, le morphisme $Y\to Y_1$
étant dominant $(\mathbf I,1.2.7)$. Supposons la proposition démontrée pour
$Y_1$, $X_1$ et $\mathcal F_1$, et soient $U_1$ l'ouvert de $Y_1$ et $Z_1$ le
sous-schéma fermé de $U_1$ ayant les propriétés voulues, $U=U_1\times_{Y_1}Y$
et $Z=Z_1\times_{Y_1}Y$ leurs images réciproques. Alors, si $U'\to U$ est un
changement de base tel que
$\mathcal F'=\mathcal F\otimes_{\mathcal O_Y}\mathcal O_{U'}
=\mathcal F_1\otimes_{\mathcal O_{Y_1}}\mathcal O_{U'}$ soit $f'$-plat, le
morphisme $U'\to U_1$ se factorise en $U'\to Z_1\to U_1$~; mais comme
$U'\to U_1$ se factorise aussi en $U'\to U\to U_1$, la définition du produit
de préschémas montre que $U'\to U$ se factorise en $U'\to Z\to U$.

On peut donc se borner au cas où $A$ est \emph{noethérien}~; soit
$\mathfrak N$ son nilradical, qui est nilpotent, et posons
$A_0=A_{\mathrm{red}}=A/\mathfrak N$,
$Y_0=\operatorname{Spec}(A_0)=Y_{\mathrm{red}}$,
$X_0=X\times_YY_0$, $\mathcal F_0=\mathcal F\otimes_{\mathcal O_Y}
\mathcal O_{Y_0}$, $W_{i0}=W_i\times_YY_0$~; si
$M_i=\Gamma(W_i,\mathcal F)$, $M_{i0}=\Gamma(W_{i0},\mathcal F_0)$ est donc
égal à $M_i\otimes_A A_0$. Comme $A_0$ est \emph{intègre}, on peut, en vertu du
théorème de platitude générique (6.9.2), en remplaçant au besoin $Y$ par un
ouvert non vide de $Y$, supposer que les $M_{i0}$ soient des $A_0$-modules
libres. En vertu de (11.4.4), il y a donc pour chaque $i$ un idéal
$\mathfrak J_i\subset\mathfrak N$ tel

\oldpage[IV-3]{147}

que les $A$-algèbres $A'$ pour lesquelles les $M_i\otimes_A A'$ sont des
$A'$-modules libres (ou plats) sont exactement celles pour lesquelles
$\mathfrak J_iA'=0$. Soit $\mathfrak J=\sum_i\mathfrak J_i$, qui est un idéal
contenu dans $\mathfrak N$~; dire que $\mathcal F\otimes_A A'$ est $A'$-plat
équivaut à dire que les $M_i\otimes_A A'$ sont tous des $A'$-modules plats,
donc que $\mathfrak JA'=0$. Il en résulte que si l'on prend
$Z=V(\mathfrak J)$, on répond à la question, car pour qu'un morphisme
$U'\to U$ ait la propriété de l'énoncé, il faut et il suffit évidemment que
pour tout ouvert affine $V'\subset U'$, le morphisme $V'\to U$ ait la même
propriété.
\end{proof}

\begin{corollary}[11.4.6]
\label{IV.11.4.6-fr}
Soient $A$ un anneau tel que $\operatorname{Spec}(A)$ soit irréductible,
$\mathfrak p$ son unique idéal premier minimal, $B$ une $A$-algèbre de
présentation finie, $M$ un $B$-module de présentation finie. Supposons que
$M_\mathfrak p$ soit un $A_\mathfrak p$-module plat~; alors il existe
$t\in A-\mathfrak p$ tel que $M_t$ soit un $A_t$-module libre.
\end{corollary}

\begin{proof}
Appliquant (11.4.5) à $Y=\operatorname{Spec}(A)$, $X=\operatorname{Spec}(B)$,
$\mathcal F=\widetilde M$, on peut (en remplaçant au besoin $A$ par $A_h$, où
$h\in A-\mathfrak p$) supposer qu'il existe un idéal de type fini
$\mathfrak J$ dans $A$ tel que les $A$-algèbres $A'$ pour lesquelles
$M\otimes_A A'$ soit un $A'$-module libre (ou plat) sont exactement celles
telles que $\mathfrak JA'=0$. En particulier, l'hypothèse que
$M_\mathfrak p$ soit un $A_\mathfrak p$-module plat entraîne
$\mathfrak JA_\mathfrak p=0$, ou encore $\mathfrak J_\mathfrak p=0$. Mais
comme $\mathfrak J$ est de type fini, il existe $t\in A-\mathfrak p$ tel que
$\mathfrak J_t=0$, ou encore $\mathfrak JA_t=0$, et par suite $M_t$ est un
$A_t$-module libre.
\end{proof}

\begin{proposition}[11.4.7]
\label{IV.11.4.7-fr}
Soient $A$ un anneau noethérien, $\mathfrak J$ un idéal de $A$, $M$ un
$A$-module idéalement séparé $(\mathbf 0_{\mathrm{III}},10.2.1)$ pour la
topologie $\mathfrak J$-préadique. Soit $(\mathfrak p_i)_{1\leq i\leq r}$ une
famille finie d'idéaux premiers de $A$ contenant $\mathfrak J$~; pour tout
entier $n\geq0$, soit $\mathfrak p_i^{(n)}$ la puissance symbolique $n$-ième de
$\mathfrak p_i$ (noyau de l'homomorphisme
$A\to A_{\mathfrak p_i}/\mathfrak p_i^nA_{\mathfrak p_i}$)~; posons
$\mathfrak J_n=\bigcap_{i=1}^{r}\mathfrak p_i^{(n)}$, de sorte que
$\mathfrak J_n\supset\mathfrak J^n$, et supposons que la topologie définie par
la filtration $(\mathfrak J_n)$ soit identique à la topologie
$\mathfrak J$-préadique (autrement dit que, pour tout $n$, il existe $m$ tel
que $\mathfrak J^n\supset\mathfrak J_m$). Pour que $M$ soit un $A$-module
plat, il faut et il suffit que $M\otimes_A(A/\mathfrak J)$ soit un
$(A/\mathfrak J)$-module plat et que, pour tout $i$,
$M\otimes_AA_{\mathfrak p_i}=M_{\mathfrak p_i}$ soit un
$A_{\mathfrak p_i}$-module plat.
\end{proposition}

\begin{proof}
Comme $M$ est idéalement séparé, il suffit, en vertu de
$(\mathbf 0_{\mathrm{III}},10.2.1)$, de montrer que, pour tout $n\geq0$,
$M\otimes_A(A/\mathfrak J^n)$ est un $(A/\mathfrak J^n)$-module plat~;
puisque tout $\mathfrak J^n$ contient un $\mathfrak J_m$, il revient au même
de prouver que pour tout $n$, $M\otimes_A(A/\mathfrak J_n)$ est un
$(A/\mathfrak J_n)$-module plat. Or, comme
$\mathfrak J_1\supset\mathfrak J$, $M\otimes_A(A/\mathfrak J_1)$ est un
$(A/\mathfrak J_1)$-module plat~; dans l'anneau $A/\mathfrak J_n$, l'idéal
$\mathfrak J_1/\mathfrak J_n$ est nilpotent et enfin l'intersection des noyaux
des homomorphismes
\[
A/\mathfrak J_n\longrightarrow
A_{\mathfrak p_i}/\mathfrak p_i^nA_{\mathfrak p_i}
\]
est nulle dans $A/\mathfrak J_n$, par définition de $\mathfrak J_n$. Il suffit
donc, par (11.4.1), de vérifier que
\[
M\otimes_A(A_{\mathfrak p_i}/\mathfrak p_i^nA_{\mathfrak p_i})
\]
est un $(A_{\mathfrak p_i}/\mathfrak p_i^nA_{\mathfrak p_i})$-module plat, ce
qui résulte de l'hypothèse que $M\otimes_AA_{\mathfrak p_i}$ est un
$A_{\mathfrak p_i}$-module plat.
\end{proof}

\begin{remark}[11.4.8]
\label{IV.11.4.8-fr}
L'hypothèse faite dans (11.4.7) sur la topologie définie par les
$\mathfrak J_n$ est vérifiée si, pour tout $n$ assez grand,
$\operatorname{Ass}(A/\mathfrak J^n)$ est contenu dans l'ensemble des
$\mathfrak p_i$. En effet, $\mathfrak J^n$ est alors intersection d'idéaux
primaires pour les $\mathfrak p_i$, dont chacun contient une puissance
symbolique de $\mathfrak p_i$, d'où la conclusion. En particulier~:
\end{remark}

\begin{corollary}[11.4.9]
\label{IV.11.4.9-fr}
Soient $A$ un anneau noethérien, $\mathfrak J$ un idéal nilpotent de $A$, $M$
un $A$-module. Pour que $M$ soit un $A$-module plat (resp. libre), il faut et
il suffit que $M\otimes_A(A/\mathfrak J)$ soit un $(A/\mathfrak J)$-module
plat (resp. libre) et que pour tout idéal premier
$\mathfrak p\in\operatorname{Ass}(A)$, $M_\mathfrak p$ soit un
$A_\mathfrak p$-module plat.
\end{corollary}
% IV3_B22_APPEND_ANCHOR

\oldpage[IV-3]{148}

L'assertion relative au cas où $M\otimes_A(A/\mathfrak J)$ est libre se déduit
encore de l'assertion relative au cas où $M\otimes_A(A/\mathfrak J)$ est plat
par $(\mathbf 0_{\mathrm{III}},10.1.2)$.

\begin{corollary}[11.4.10]
\label{IV.11.4.10-fr}
Soient $A$ un anneau noethérien, $\mathfrak J$ un idéal de $A$, $M$ un
$A$-module. On suppose que $M$ est idéalement séparé pour la topologie
$\mathfrak J$-préadique et que $\operatorname{gr}_{\mathfrak J}(A)$ est un
$(A/\mathfrak J)$-module plat. Pour que $M$ soit un $A$-module plat, il faut et
il suffit que $M\otimes_A(A/\mathfrak J)$ soit un $(A/\mathfrak J)$-module plat
et que, pour tout $\mathfrak p\in\operatorname{Ass}(A/\mathfrak J)$,
$M_{\mathfrak p}$ soit un $A_{\mathfrak p}$-module plat.
\end{corollary}

\begin{proof}
Compte tenu de (11.4.8), tout revient à montrer que
$\operatorname{Ass}(A/\mathfrak J^n)$ est contenu dans
$\operatorname{Ass}(A/\mathfrak J)$ pour tout $n$. Or, si $a\in A$ n'appartient
à aucun des $\mathfrak p\in\operatorname{Ass}(A/\mathfrak J)$, l'homothétie de
rapport $a$ est injective dans $A/\mathfrak J$~; comme chacun des
$\mathfrak J^k/\mathfrak J^{k+1}$ est un $(A/\mathfrak J)$-module plat, $a$ est
aussi un élément $(\mathfrak J^k/\mathfrak J^{k+1})$-régulier, donc $a$ est
$(A/\mathfrak J^n)$-régulier pour tout $n$ (Bourbaki, \emph{Alg. comm.},
chap.~III, \S~2, n$^\circ$~8, cor.~1 du th.~1), et par suite n'appartient à
aucun idéal premier associé à $A/\mathfrak J^n$, d'où le corollaire (Bourbaki,
\emph{Alg. comm.}, chap.~II, \S~1, n$^\circ$~1, prop.~2).
\end{proof}

\begin{proposition}[11.4.11]
\label{IV.11.4.11-fr}
Soient $A$ un anneau local artinien, d'idéal maximal $\mathfrak m$,
$Y=\operatorname{Spec}(A)$, $y$ l'unique point de $Y$, $f:X\to Y$ un morphisme
localement de type fini, $\mathcal F$ un $\mathcal O_X$-Module cohérent. Soient
$(A'_\alpha)$ une famille d'anneaux locaux, et pour chaque $\alpha$,
$u_\alpha:A\to A'_\alpha$ un homomorphisme d'anneaux (nécessairement local).
Posons $Y'_\alpha=\operatorname{Spec}(A'_\alpha)$,
$X'_\alpha=X\times_Y Y'_\alpha$,
$\mathcal F'_\alpha=\mathcal F\otimes_A A'_\alpha$. Soit $x$ un point de $X$,
et supposons vérifiées les conditions suivantes~:
\begin{enumerate}[label=(\roman*)]
  \item L'intersection des noyaux des $u_\alpha$ est réduite à $0$.
  \item L'extension $k(x)$ du corps résiduel $k$ de $A$ est primaire (4.3.1)
  (condition automatiquement vérifiée si $k$ est séparablement clos).
  \item Pour tout $\alpha$, il existe un point $x'_\alpha\in X'_\alpha$ dont
  les projections respectives dans $X$ et $Y'_\alpha$ sont $x$ et le point
  fermé $y'_\alpha$ de $Y'_\alpha$, et tel que $\mathcal F'_\alpha$ soit
  $Y'_\alpha$-plat au point $x'_\alpha$.
\end{enumerate}
Dans ces conditions, $\mathcal F$ est $f$-plat au point $x$.
\end{proposition}

\begin{proof}
La question étant locale sur $X$, on peut évidemment se borner au cas où $f$
est un morphisme de type fini et supposer $\mathcal F_x\neq0$. Nous procéderons
en plusieurs étapes.

\medskip
\noindent\textbf{I)} \emph{Réduction au cas où $A'_\alpha$ est un anneau local
d'un $Y$-préschéma de type fini et où le corps résiduel $k'_\alpha$ de
$A'_\alpha$ est une extension finie de $k$.}

\medskip
Comme la réduction se fait séparément sur chaque $A'_\alpha$, on peut supprimer
dans cette partie l'indice $\alpha$. Soit $\mathfrak m'$ l'idéal maximal de
$A'$. Considérons $A'$ comme limite inductive de ses sous-$A$-algèbres de type
fini $B_\lambda$, et posons $\mathfrak n_\lambda=\mathfrak m'\cap B_\lambda$~;
$A'$ est aussi limite inductive de ses sous-anneaux locaux
$B'_\lambda=(B_\lambda)_{\mathfrak n_\lambda}$ (5.13.3), et on a évidemment
$\mathfrak m'\cap B'_\lambda=\mathfrak n'_\lambda$, idéal maximal de
$B'_\lambda$. Il existe donc (11.2.6) un indice $\lambda$ tel que, en posant
$Z'_\lambda=\operatorname{Spec}(B'_\lambda)$,
$\mathcal F\otimes_A B'_\lambda$ soit $Z'_\lambda$-plat au point $x'_\lambda$,
projection de $x'$, la projection de $x'_\lambda$ sur $Z'_\lambda$ étant le
point fermé $z'_\lambda$ de $Z'_\lambda$.

On peut donc supposer qu'il existe un schéma affine $Z'$ de type fini sur $Y$
et un point $z'$ de $Z'$ tel que $A'=\mathcal O_{Z',z'}$, et que, si l'on pose
$T'=X\times_Y Z'$, il existe un point $t'\in T'$ dont les projections dans
$Z'$ et $X$ sont $z'$ et $x$, et tel que $\mathcal F\otimes_Y Z'$ soit
$Z'$-plat au point $t'$. Soient $Z'_1$ (resp. $X_1$) un sous-préschéma fermé de
$Z'$ (resp. $X$) ayant pour espace sous-jacent l'adhérence de $\{z'\}$ (resp.
$\{x\}$), et posons $T'_1=X_1\times_Y Z'_1$. L'ensemble $U$ des points de
$T'$ où $\mathcal F\otimes_Y Z'$ est $Z'$-plat est ouvert dans $T'$ (11.1.1) et
contient $t'$, donc

\oldpage[IV-3]{149}

$V=U\cap T'_1$ est ouvert non vide dans $T'_1$. L'anneau $A$, étant artinien,
est un anneau de Jacobson, donc (10.4.6) $Z'_1$ et $T'_1$ sont des préschémas de
Jacobson~; il existe par suite dans $V$ un point $t'_0$ fermé dans $V$, et son
image $z'_0$ dans $Z'_1$ est un point fermé de $Z'_1$ (10.4.7). Soient $f_1$ la
restriction de $f$ à $X_1$ et $f'_1:T'_1\to Z'_1$ la projection canonique~;
$V\cap f_1'^{-1}(z'_0)$ est un ouvert non vide dans
$f_1^{-1}(y)\otimes_{k(y)}k(z'_0)$, et comme ce dernier préschéma est plat sur
$f_1^{-1}(y)$, un point maximal $t'_1$ de $V\cap f_1'^{-1}(z'_0)$ est
nécessairement au-dessus de l'unique point maximal $x$ de $X_1$ (2.3.4).
Enfin, $k(z'_0)$ est une extension finie de $k$ (\textbf I, 6.4.2) et
l'homomorphisme $A=\mathcal O_{Y,y}\to A'=\mathcal O_{Z',z'}$ se factorise en
$\mathcal O_{Y,y}\to\mathcal O_{Z',z'_0}\to\mathcal O_{Z',z'}$, donc son noyau
est contenu dans celui de $\mathcal O_{Y,y}\to\mathcal O_{Z',z'_0}$. Ceci
achève la réduction annoncée.

\medskip
\noindent\textbf{II)} \emph{Réduction au cas où les $A'_\alpha$ sont en nombre
fini et sont des $A$-algèbres finies.} --- Soit $\mathfrak m'_\alpha$ l'idéal
maximal de $A'_\alpha$~; comme $A'_\alpha$ est un anneau local noethérien,
l'intersection des $\mathfrak m'^n_\alpha$ est $0$ (\textbf I, 7.3.5)~;
l'intersection des $u_\alpha^{-1}(\mathfrak m'^n_\alpha)$, pour tous les indices
$n$ et $\alpha$, est donc égale à l'intersection des noyaux des $u_\alpha$, donc
est réduite à $0$ par l'hypothèse (i). Puisque $A$ est artinien, il y a un nombre
fini de ces idéaux dont l'intersection est déjà $0$~; notons-les
$u_i^{-1}(\mathfrak m_i'^{\,n_i})$ $(1\leq i\leq r)$. Comme les $A'_i$ sont
noethériens, les $\mathfrak m_i'^{\,h}/\mathfrak m_i'^{\,h+1}$ sont des
$(A'_i/\mathfrak m'_i)$-modules de longueur finie, et comme
$A'_i/\mathfrak m'_i$ est un $k$-espace vectoriel de rang fini, on voit que
$A''_i=A'_i/\mathfrak m_i'^{\,n_i}$ est une $A$-algèbre finie et un anneau local
artinien. La réduction annoncée résulte donc de (2.1.4).

\medskip
\noindent\textbf{III)} \emph{Fin de la démonstration.} --- Supposons désormais
que les $A'_i$ $(1\leq i\leq r)$ soient en nombre fini et soient des
$A$-algèbres finies. Pour tout $i$, le corps résiduel $k_i$ de $A'_i$ est une
extension finie de $k$~; utilisant l'hypothèse (ii), on en conclut que l'image
réciproque de $x$ dans $X'_i$ est réduite au seul point $x'_i$ (4.3.2). Soient
$Y'$ le préschéma somme des $Y'_i$, $X'=X\times_Y Y'$, somme des $X'_i$,
$\mathcal F'=\mathcal F\otimes_Y Y'$. L'hypothèse entraîne que
$\mathcal F'$ est $Y'$-plat aux points de l'image réciproque de $x$ par la
projection $p:X'\to X$. Comme $p$ est de type fini, il existe par suite un
ouvert $U'\supset p^{-1}(x)$ tel que $\mathcal F'$ soit $Y'$-plat aux points de
$U'$ (11.1.1). En outre, le morphisme $Y'\to Y$ est fini puisque les $A'_i$
sont des $A$-algèbres finies~; donc $p$ est un morphisme fini (\textbf{II},
6.1.5), par suite fermé (\textbf{II}, 6.1.10), et il existe donc un voisinage
affine $U$ de $x$ dans $X$ tel que $p^{-1}(U)\subset U'$. Soient $B$ et $A'$
les anneaux des schémas $U$ et $Y'$ ($A'$ étant le composé direct des $A'_i$)~;
remplaçant $X$ par $U$, on a donc $\mathcal F=\widetilde M$, où $M$ est un
$B$-module et, par hypothèse, $M'=M\otimes_A A'$ est un $A'$-module plat
(2.1.2)~; comme l'homomorphisme $A\to A'$ est injectif par construction, on
peut appliquer (11.4.3), qui prouve que $M$ est un $A$-module plat, C.Q.F.D.
\end{proof}

\begin{proposition}[11.4.12]
\label{IV.11.4.12-fr}
Soient $A$ un anneau local artinien de corps résiduel $k$,
$Y=\operatorname{Spec}(A)$, $f:X\to Y$ un morphisme localement de type fini,
$\mathcal F$ un $\mathcal O_X$-Module cohérent. Soient $(A'_\alpha)$ une famille
d'anneaux locaux, et pour chaque $\alpha$, $u_\alpha:A\to A'_\alpha$ un
homomorphisme d'anneaux. Posons $Y'_\alpha=\operatorname{Spec}(A'_\alpha)$,
$X'_\alpha=X\times_Y Y'_\alpha$, $f'_\alpha=f_{(Y'_\alpha)}$,
$\mathcal F'_\alpha=\mathcal F\otimes_A A'_\alpha$. Soit $x$ un point de $X$,
et supposons vérifiées les conditions suivantes~:
\begin{enumerate}[label=(\roman*)]
  \item L'intersection des noyaux des $u_\alpha$ est réduite à $0$.
  \item Pour tout $\alpha$, $\mathcal F'_\alpha$ est $f'_\alpha$-plat en tous
  les points $x'_\alpha\in X'_\alpha$ dont les projections respectives dans
  $X$ et $Y'_\alpha$ sont $x$ et le point fermé $y'_\alpha$ de $Y'_\alpha$.
\end{enumerate}
Alors $\mathcal F$ est $f$-plat au point $x$.
\end{proposition}

\begin{proof}
% IV3_B22_P149_ANCHOR

\oldpage[IV-3]{150}

Par hypothèse, l'intersection des noyaux des $u_\alpha$ est réduite à $0$~;
comme $A$ est artinien, il y a déjà un nombre fini de ces noyaux dont
l'intersection est $0$, donc on peut se borner au cas où la famille
$(A'_\alpha)$ est finie. Soit $k'$ une clôture algébrique de $k$~; on sait
$(\mathbf 0_{\mathrm{III}},10.3.1)$ qu'il existe un homomorphisme local
$A\to B$ faisant de $B$ un $A$-module plat, tel que $B$ soit un anneau local
artinien, entier sur $A$ et que $B\otimes_A k$ soit isomorphe à $k'$.

Par platitude, les noyaux des homomorphismes
$B\to B'_\alpha=B\otimes_A A'_\alpha$ se déduisent de ceux des $u_\alpha$ par
tensorisation avec $B$, et comme ils sont en nombre fini, leur intersection est
réduite à $0$ $(\mathbf 0_{\mathrm I},6.1.3)$. Considérons les anneaux
$B'_{\alpha\beta}$, localisés de $B'_\alpha$ en ses idéaux maximaux~; on sait
(Bourbaki, \emph{Alg. comm.}, chap.~II, \S~3, n$^\circ$~3, cor.~2 du th.~1)
que l'intersection des noyaux des homomorphismes
$B'_\alpha\to B'_{\alpha\beta}$ (pour un $\alpha$ donné) est réduite à $0$~;
on en conclut que l'intersection des noyaux des homomorphismes composés
$v_{\alpha\beta}:B\to B'_\alpha\to B'_{\alpha\beta}$ ($\alpha$ et $\beta$
variables) est réduite à $0$. D'autre part, comme $B$ est entier sur $A$,
$B'_\alpha$ est entier sur $A'_\alpha$, donc ses idéaux maximaux sont au-dessus
de l'idéal maximal de $A'_\alpha$. Si l'on pose
$Z'_{\alpha\beta}=\operatorname{Spec}(B'_{\alpha\beta})$,
$T'_{\alpha\beta}=X'_\alpha\times_{Y'_\alpha}Z'_{\alpha\beta}$,
$f'_{\alpha\beta}=f_{(Z'_{\alpha\beta})}$,
$\mathcal G=\mathcal F\otimes_A B$ et
$\mathcal G'_{\alpha\beta}=\mathcal G\otimes_B B'_{\alpha\beta}$, on voit donc
que les hypothèses (i) et (ii) sont encore vérifiées lorsqu'on remplace
respectivement $A$, $A'_\alpha$, $u_\alpha$, $\mathcal F$ et $x$ par $B$,
$B'_{\alpha\beta}$, $v_{\alpha\beta}$, $\mathcal G$ et un point $t$ de
$T=X\otimes_A B$ au-dessus de $x$. Comme le corps résiduel de $B$ est
séparablement clos, on déduit de (11.4.11) que $\mathcal G$ est plat sur $B$ au
point $t$. Mais puisque $B$ est un $A$-module fidèlement plat, on conclut de
(2.1.4) que $\mathcal F$ est plat sur $A$ au point $x$, ce qui prouve (11.4.12).
\end{proof}

\begin{corollary}[11.4.13]
\label{IV.11.4.13-fr}
Soient $A$ un anneau artinien local, $Y=\operatorname{Spec}(A)$,
$f:X\to Y$ un morphisme localement de type fini, $\mathcal F$ un
$\mathcal O_X$-Module cohérent. Soit $(Y'_\alpha)$ une famille de
$Y$-préschémas, et pour tout $\alpha$, posons
$X'_\alpha=X\times_Y Y'_\alpha$, $f'_\alpha=f_{(Y'_\alpha)}$,
$\mathcal F'_\alpha=\mathcal F\otimes_Y Y'_\alpha$. Soit $x$ un point de $X$ et
supposons vérifiées les hypothèses suivantes~:
\begin{enumerate}[label=(\roman*)]
  \item L'intersection des noyaux des homomorphismes
  $A\to\Gamma(Y'_\alpha,\mathcal O_{Y'_\alpha})$ correspondant aux morphismes
  structuraux est réduite à $0$.
  \item Pour tout $\alpha$, $\mathcal F'_\alpha$ est $f'_\alpha$-plat en tous
  les points $x'_\alpha\in X'_\alpha$ dont la projection dans $X$ est $x$.
\end{enumerate}
Alors $\mathcal F$ est $f$-plat au point $x$.
\end{corollary}

\begin{proof}
En effet, pour tout $y'_\alpha\in Y'_\alpha$, considérons le schéma local
$Y''_{\alpha,y'_\alpha}=\operatorname{Spec}(\mathcal O_{Y'_\alpha,y'_\alpha})$~;
en vertu de (2.1.4),
$\mathcal F'_\alpha\otimes_{Y'_\alpha}Y''_{\alpha,y'_\alpha}$ est
$Y''_{\alpha,y'_\alpha}$-plat aux points dont les projections sur $X$ et
$Y''_{\alpha,y'_\alpha}$ sont $x$ et le point fermé de
$Y''_{\alpha,y'_\alpha}$. En outre le noyau de l'homomorphisme
$A\to\Gamma(Y'_\alpha,\mathcal O_{Y'_\alpha})$ est l'intersection des noyaux
des homomorphismes
$A\to\Gamma(Y''_{\alpha,y'_\alpha},\mathcal O_{Y''_{\alpha,y'_\alpha}})$, car
on se ramène aussitôt au cas où $Y'_\alpha$ est affine, et il suffit alors
d'appliquer Bourbaki, \emph{Alg. comm.}, chap.~II, \S~3, n$^\circ$~3, cor.~2 du
th.~1. En remplaçant la famille $(Y'_\alpha)$ par la famille des
$Y''_{\alpha,y'_\alpha}$, on est donc ramené à (11.4.12).
\end{proof}

\subsection{Descente de la platitude par des morphismes quelconques : cas général}
\label{subsection:IV.11.5-fr}

\begin{theorem}[11.5.1]
\label{IV.11.5.1-fr}
Soient $Y$ un préschéma localement noethérien, $f:X\to Y$ un morphisme
localement de type fini, $\mathcal F$ un $\mathcal O_X$-Module cohérent, $x$ un
point de $X$, $y=f(x)$. Soit $(Y'_\alpha)$ une famille de $Y$-préschémas locaux
$Y'_\alpha=\operatorname{Spec}(A'_\alpha)$ telle que les images des points
fermés $y'_\alpha$ des $Y'_\alpha$ soient toutes égales à $y$. Pour tout
$\alpha$, soient $\mathfrak m'_\alpha$ l'idéal maximal de $A'_\alpha$, et
$u_\alpha:\mathcal O_y\to A'_\alpha$
\oldpage[IV-3]{151}
% IV3_B22_P151_ANCHOR
l'homomorphisme canonique
(\textbf I, 2.4.4)~; on suppose que les intersections finies des idéaux
$u_\alpha^{-1}(\mathfrak m_\alpha'^{\,n_\alpha})$ forment un système fondamental
de voisinages de $0$ dans $\mathcal O_y$. Posons
$X'_\alpha=X\times_Y Y'_\alpha$, $f'_\alpha=f_{(Y'_\alpha)}$,
$\mathcal F'_\alpha=\mathcal F\otimes_Y Y'_\alpha$, et supposons vérifiée l'une
des hypothèses suivantes~:
\begin{enumerate}[label=(\roman*)]
  \item Pour tout $\alpha$, $\mathcal F'_\alpha$ est $f'_\alpha$-plat en tous
  les points dont la projection dans $X$ est égale à $x$ et la projection dans
  $Y'_\alpha$ égale à $y'_\alpha$.
  \item Pour tout $\alpha$, il existe un $x'_\alpha\in X'_\alpha$ dont la
  projection sur $X$ est $x$ et la projection dans $Y'_\alpha$ égale à
  $y'_\alpha$, tel que $\mathcal F'_\alpha$ soit $f'_\alpha$-plat au point
  $x'_\alpha$, et $k(x)$ est une extension primaire de $k(y)$.
\end{enumerate}
Dans ces conditions, $\mathcal F$ est $f$-plat au point $x$.
\end{theorem}

\begin{proof}
Soit $\mathfrak m_y$ l'idéal maximal de $\mathcal O_y$~; comme $\mathcal O_y$
et $\mathcal O_x$ sont noethériens et $\mathcal O_y\to\mathcal O_x$ un
homomorphisme local, il suffit, en vertu de
$(\mathbf 0_{\mathrm{III}},10.2.2)$, de prouver que pour tout $n>0$,
$\mathcal F_x/\mathfrak m_y^n\mathcal F_x$ est un
$(\mathcal O_y/\mathfrak m_y^n)$-module plat. Désignons par
$(\mathfrak J_\lambda)$ la famille des intersections finies des
$u_\alpha^{-1}(\mathfrak m_\alpha'^{\,n_\alpha})$~; par hypothèse, il existe
$\lambda$ tel que $\mathfrak J_\lambda\subset\mathfrak m_y^n$, et comme
\[
\mathcal F_x\otimes_{\mathcal O_y}(\mathcal O_y/\mathfrak m_y^n)
=\bigl(\mathcal F_x\otimes_{\mathcal O_y}(\mathcal O_y/\mathfrak J_\lambda)\bigr)
 \otimes_{\mathcal O_y/\mathfrak J_\lambda}(\mathcal O_y/\mathfrak m_y^n),
\]
il suffira de prouver que $\mathcal F_x/\mathfrak J_\lambda\mathcal F_x$ est un
$(\mathcal O_y/\mathfrak J_\lambda)$-module plat. Or, pour chaque $\alpha$ tel
que $\mathfrak J_\lambda\subset u_\alpha^{-1}(\mathfrak m_\alpha'^{\,n_\alpha})$,
on a, par passage aux quotients, un homomorphisme
$u'_\alpha:\mathcal O_y/\mathfrak J_\lambda\to
A'_\alpha/\mathfrak m_\alpha'^{\,n_\alpha}$, et l'intersection des noyaux des
$u'_\alpha$ est réduite à $0$. Compte tenu de (\textbf I, 3.6.1), on voit qu'on
est ramené à (11.4.11) dans le cas (ii), et à (11.4.12) dans le cas (i).
\end{proof}

\begin{corollary}[11.5.2]
\label{IV.11.5.2-fr}
Soient $Y$ un préschéma localement noethérien, $f:X\to Y$ un morphisme
localement de type fini, $\mathcal F$ un $\mathcal O_X$-Module cohérent, $x$ un
point de $X$, $y=f(x)$~; posons $A=\mathcal O_y$. Soit $u:A\to A'$ un
homomorphisme de $A$ dans un anneau de Zariski $A'$, tel que l'image réciproque
$u^{-1}(\mathfrak r')$ du radical $\mathfrak r'$ de $A'$ soit l'idéal maximal
$\mathfrak m$ de $A$~; supposons en outre que l'homomorphisme
$\widehat u:\widehat A\to\widehat{A'}$ soit injectif. Posons
$Y'=\operatorname{Spec}(A')$, $X'=X\times_Y Y'$, $f'=f_{(Y')}$,
$\mathcal F'=\mathcal F\otimes_Y Y'$. Pour que $\mathcal F$ soit $f$-plat au
point $x$, il faut et il suffit que $\mathcal F'$ soit $f'$-plat en tout point
dont la projection dans $X$ est égale à $x$ et dont la projection dans $Y'$ est
égale à un point fermé de $Y'$.

Si en outre $A'$ est une $A$-algèbre finie, on peut dans ce qui précède
remplacer l'hypothèse que $\widehat u$ est injectif par l'hypothèse que $u$ est
injectif.
\end{corollary}

\begin{proof}

Comme $A$ (resp. $A'$) s'identifie à un sous-anneau de $\widehat A$ (resp.
$\widehat{A'}$) (Bourbaki, \emph{Alg. comm.}, chap.~III, \S~3, n$^\circ$~3,
prop.~6), on voit d'abord que $u$ lui-même est injectif et que $\widehat u$ est
son prolongement par continuité à $\widehat A$.

Soit $(\mathfrak m'_\alpha)$ la famille des idéaux maximaux de $A'$~; comme on a
\[
\mathfrak m_\alpha'^{\,n}\widehat{A'}
=\widehat{\mathfrak m}_\alpha'^{\,n},
\qquad\text{et}\qquad
\widehat{\mathfrak m}_\alpha'^{\,n}\cap A'
=\mathfrak m_\alpha'^{\,n},
\]
et que les $\mathfrak m_\alpha'^{\,n}$ sont ouverts dans $A'$, on a
$u^{-1}(\mathfrak m_\alpha'^{\,n})
=A\cap\widehat u^{-1}(\widehat{\mathfrak m}_\alpha'^{\,n})$, et il suffira de
montrer que dans $\widehat A$, les intersections finies des
$\widehat u^{-1}(\widehat{\mathfrak m}_\alpha'^{\,n_\alpha})$ forment un système
fondamental de voisinages de $0$, ce qui permettra d'appliquer (11.5.1) aux
homomorphismes composés $v_\alpha\circ u:A\to A'_{\mathfrak m'_\alpha}$,
où $v_\alpha:A'\to A'_{\mathfrak m'_\alpha}$ est l'homomorphisme canonique.
Comme $\widehat A$ est complet, il suffira de montrer que l'intersection des
$\widehat u^{-1}(\widehat{\mathfrak m}_\alpha'^{\,n_\alpha})$ est réduite à
$0$ (Bourbaki, \emph{Alg. comm.}, chap.~III, \S~2, n$^\circ$~7, prop.~8, où on
peut dans la démonstration remplacer la suite décroissante par un ensemble
filtrant quelconque). Or, pour tout $\alpha$
\oldpage[IV-3]{152}
fixé, l'intersection des idéaux
$\widehat{\mathfrak m}_\alpha'^{\,n}\widehat{A'}_{\widehat{\mathfrak m}'_\alpha}$
pour $n>0$ est réduite à $0$ dans l'anneau local noethérien
$\widehat{A'}_{\widehat{\mathfrak m}'_\alpha}$. D'autre part les
$\widehat{\mathfrak m}'_\alpha$ sont les idéaux maximaux de $\widehat{A'}$, donc
l'homomorphisme canonique
$\widehat{A'}\to\prod_\alpha\widehat{A'}_{\widehat{\mathfrak m}'_\alpha}$ est
injectif (Bourbaki, \emph{Alg. comm.}, chap.~II, \S~3, n$^\circ$~3, cor.~2 du
th.~1), et comme par hypothèse $\widehat u:\widehat A\to\widehat{A'}$ est aussi
injectif, cela achève la démonstration dans le cas général. La dernière
assertion résulte de ce que $\widehat A$ est un $A$-module fidèlement plat
($\mathbf 0_{\mathrm I},7.3.5$) et $\widehat{A'}=A'\otimes_A\widehat A$ puisque
$A'$ est par hypothèse un $A$-module de type fini (Bourbaki, \emph{Alg. comm.},
chap.~III, \S~3, n$^\circ$~4, th.~3 et chap.~IV, \S~2, n$^\circ$~5, cor.~3 de la
prop.~9).
\end{proof}

\begin{proposition}[11.5.3]
\label{IV.11.5.3-fr}
Soient $f:X\to Y$ un morphisme localement de présentation finie,
$\mathcal F$ un $\mathcal O_X$-Module quasi-cohérent de présentation finie, $x$
un point de $X$, $y=f(x)$, $g=(\psi,\theta):Y'\to Y$ un morphisme propre de
présentation finie. Supposons que~:
\begin{enumerate}[label=(\roman*)]
  \item l'homomorphisme $\theta_y:\mathcal O_y\to
  (g_*\mathcal O_{Y'})_y$ soit injectif~;
  \item pour tout $x'\in X'=X\times_Y Y'$ dont la projection dans $X$ est $x$,
  $\mathcal F'=\mathcal F\otimes_Y Y'$ soit $Y'$-plat au point $x'$.
\end{enumerate}
Alors $\mathcal F$ est $f$-plat au point $x$.
\end{proposition}

\begin{proof}
La question étant locale sur $X$, on peut supposer $f$ de présentation finie.
Soit $p:X'\to X$ la projection canonique. Comme
$f'=f_{(Y')}:X'\to Y'$ est de présentation finie (1.6.2) et que $\mathcal F'$
est un $\mathcal O_{X'}$-Module de présentation finie
$(\mathbf 0_{\mathrm I},5.2.5)$, il résulte de (11.3.1) que l'ensemble $U'$ des
points de $X'$ où $\mathcal F'$ est $f'$-plat est ouvert. Comme $U'$ contient
$p^{-1}(x)$ par hypothèse, et que $p$ est propre, donc fermé, $U'$ contient un
ensemble de la forme $p^{-1}(U)$, où $U$ est un voisinage de $x$. Remplaçant $X$
par $U$, on peut donc supposer déjà que $\mathcal F'$ est $f'$-plat.

D'autre part, en tenant compte de (\textbf I, 3.6.5), (\textbf{II}, 5.4.2) et
(2.1.4), on peut remplacer $Y$ par $\operatorname{Spec}(\mathcal O_y)$,
autrement dit supposer que $Y=\operatorname{Spec}(A)$, où $A$ est un anneau
local. Sous ces conditions, on va prouver que $\mathcal F$ est $f$-plat. En
vertu de (5.13.3), $A$ est limite inductive de sous-anneaux locaux noethériens
$A_\lambda$ tels que l'injection canonique $A_\lambda\to A$ soit un
homomorphisme local. En vertu de (8.9.1), on peut supposer que
$X=X_\lambda\otimes_{A_\lambda}A$, $f=f_\lambda\otimes1_A$,
$\mathcal F=\mathcal F_\lambda\otimes_{A_\lambda}A$, pour un $\lambda$
convenable, $f_\lambda:X_\lambda\to Y_\lambda=\operatorname{Spec}(A_\lambda)$
étant un morphisme de type fini et $\mathcal F_\lambda$ un
$\mathcal O_{X_\lambda}$-Module cohérent. De même, on peut supposer que
$Y'=Y'_\lambda\otimes_{A_\lambda}A$, $g=g_\lambda\otimes1_A$, où
$g_\lambda:Y'_\lambda\to Y_\lambda$ est un morphisme de présentation finie~;
en outre, (8.10.5, (xii)) permet de supposer $g_\lambda$ propre. Comme par
hypothèse l'homomorphisme $A\to\Gamma(Y',\mathcal O_{Y'})$ est injectif et que
$A_\lambda$ est un sous-anneau de $A$, l'homomorphisme
$A_\lambda\to\Gamma(Y',\mathcal O_{Y'})$ est aussi injectif. Enfin, en vertu de
(11.2.6), on peut supposer $\lambda$ pris tel que
$\mathcal F'_\lambda=\mathcal F_\lambda\otimes_{Y_\lambda}Y'_\lambda$ soit
$f'_\lambda$-plat, puisque $\mathcal F'=\mathcal F'_\lambda\otimes_{Y'_\lambda}Y'$.
Ces remarques prouvent que l'on peut désormais supposer l'anneau $A$
noethérien, les autres hypothèses de (11.5.3) étant vérifiées.

Posons alors $B=\widehat A$, $Z=\operatorname{Spec}(B)$~; comme $B$ est un
$A$-module fidèlement plat $(\mathbf 0_{\mathrm I},7.3.5)$, il revient au même
de dire que $\mathcal F$ est $f$-plat ou que $\mathcal F\otimes_Y Z$ est
$Z$-plat (2.1.4). De même, si l'on pose $Z'=Y'\times_Y Z$, le morphisme
$Z'\to Y'$ est fidèlement plat (2.2.13), donc il revient au même de dire que
$\mathcal F'$ est $f'$-plat ou que $\mathcal F'\otimes_{Y'}Z'$ est $Z'$-plat~;
enfin $h=g_{(Z)}:Z'\to Z$ est propre (\textbf{II}, 5.4.2) et de type fini
(1.5.2), $\widehat A$ est noethérien, et si $z$ est son point
\oldpage[IV-3]{153}
fermé,
l'homomorphisme $\mathcal O_z\to(h_*\mathcal O_{Z'})_z$ est injectif, car il
résulte de (2.3.1) que $h_*\mathcal O_{Z'}=g_*\mathcal O_{Y'}\otimes_Y Z$,
et notre assertion résulte de l'hypothèse (i) et de la définition des modules
plats $(\mathbf 0_{\mathrm I},6.1.1)$.

On peut donc désormais supposer l'anneau local noethérien $A$ complet~; la
démonstration sera achevée si l'on prouve que l'intersection des noyaux des
homomorphismes $A=\mathcal O_y\to\mathcal O_{y'}$ (où $y'$ parcourt
$g^{-1}(y)$) est réduite à $0$. En effet, les $\mathcal O_{y'}$ sont des anneaux
locaux noethériens, donc pour chaque $y'$ l'intersection des
$\mathfrak m_{y'}^n$ $(n\geq0)$ est réduite à $0$~; si $\mathfrak a_{n,y'}$ est
l'image réciproque dans $A$ de $\mathfrak m_{y'}^n$, les intersections finies
des $\mathfrak a_{n,y'}$ sont des voisinages de $0$ dans $A$ et l'intersection de
tous les $\mathfrak a_{n,y'}$ est réduite à $0$~; elles forment donc un système
fondamental de voisinages de $0$ dans $A$ (Bourbaki, \emph{Alg. comm.}, chap.~III,
\S~2, n$^\circ$~7, prop.~8, où l'on peut remplacer la suite décroissante par un
ensemble filtrant quelconque), et l'on pourra par suite appliquer (11.5.1).
Or, soit $s\in A$ un élément appartenant au noyau de chacun des homomorphismes
$A\to\mathcal O_{y'}$~; l'image $s'$ de $s$ dans
$\Gamma(Y',\mathcal O_{Y'})$ est donc une section de $\mathcal O_{Y'}$ au-dessus
de $Y'$ telle que $s'_{y'}=0$ pour tout $y'\in g^{-1}(y)$~; il y a par suite un
voisinage $V'$ de $g^{-1}(y)$ dans $Y'$ tel que $s'|V'=0$. Mais comme $g$ est
fermé, $V'$ contient un ensemble de la forme $g^{-1}(V)$, où $V$ est un
voisinage ouvert de $y$ dans $Y$~; or $Y$ est un schéma local, donc le seul
voisinage du point fermé $y$ de $Y$ est $Y$ tout entier, autrement dit $V'=Y'$,
$s'=0$, et comme $A\to\Gamma(Y',\mathcal O_{Y'})$ est injectif par hypothèse,
$s=0$. C.Q.F.D.
\end{proof}

Le cas particulier suivant de (11.5.3) nous sera utile au chap.~V.

\begin{corollary}[11.5.4]
\label{IV.11.5.4-fr}
Soit $f=(\psi,\theta):X\to Y$ un morphisme propre de présentation finie, et
soit $p:X\times_YX\to X$ la première projection. On suppose que
$\theta:\mathcal O_Y\to f_*(\mathcal O_X)$ est injectif. Alors, pour que $f$
soit plat, il faut et il suffit que $p$ le soit.
\end{corollary}

\begin{proposition}[11.5.5]
\label{IV.11.5.5-fr}
Soient $A$ un anneau, $Y=\operatorname{Spec}(A)$, $f:X\to Y$ un morphisme
localement de présentation finie, $\mathcal F$ un $\mathcal O_X$-Module
quasi-cohérent de présentation finie, $x$ un point de $X$. Soit
$A\to A'$ un homomorphisme injectif faisant de $A'$ une algèbre entière sur $A$.
Posons $Y'=\operatorname{Spec}(A')$, $X'=X\times_Y Y'$, $f'=f\times1$,
$\mathcal F'=\mathcal F\otimes_Y Y'$. Alors, si $\mathcal F'$ est $f'$-plat en
tout point de $X'$ dont la projection dans $X$ est égale à $x$, $\mathcal F$ est
$f$-plat au point $x$.
\end{proposition}

\begin{proof}
Supposons d'abord que $A'$ soit une $A$-algèbre finie et de présentation finie~;
alors le morphisme $Y'\to Y$ est propre (\textbf{II}, 6.1.11) et de présentation
finie, donc les hypothèses de (11.5.3) sont vérifiées, d'où la conclusion. Dans
le cas général, la proposition résultera de ce cas particulier, du fait que $A'$
est limite inductive de ses sous-$A$-algèbres finies $A'_\lambda$ et des deux
lemmes suivants~:
\end{proof}

\begin{lemma}[11.5.5.1]
\label{IV.11.5.5.1-fr}
Toute $A$-algèbre finie $A'$ est limite inductive de $A$-algèbres $A'_\lambda$
qui sont finies et de présentation finie.
\end{lemma}

\begin{proof}
On raisonne comme dans (1.9.3.1). On a en effet $A'=B/\mathfrak J$, où $B$ est
une $A$-algèbre finie qui est un $A$-module libre, et $\mathfrak J$ un idéal de
$B$ (1.4.7.1). Or, $\mathfrak J$ est limite inductive des idéaux
$\mathfrak J_\lambda\subset\mathfrak J$ de $B$ qui sont de type fini (et a
fortiori des $A$-modules de type fini), donc, par l'exactitude du foncteur
$\varinjlim$, $A'$ est limite inductive des $A$-algèbres
$B/\mathfrak J_\lambda$~; or $B/\mathfrak J_\lambda$ est par définition un
$A$-module de présentation finie, donc aussi (1.4.7) une $A$-algèbre de
présentation finie.
\end{proof}
% IV3_B22_P153_END

\oldpage[IV-3]{154}

Pour appliquer ce lemme à la situation de (11.5.5) on notera en outre que si
l'homomorphisme $A\to A'$ est injectif, il en est ainsi \emph{a fortiori} de
$A\to A'_{\lambda}$, pour tout $\lambda$.

\begin{lemma}[11.5.5.2]
\label{IV.11.5.5.2-fr}
Soient $A$ un anneau, $A'$ une $A$-algèbre, $(A'_{\lambda})$ un système inductif
de $A$-algèbres tel que $A'=\varinjlim A'_{\lambda}$~; on pose
$Y=\operatorname{Spec}(A)$, $Y'=\operatorname{Spec}(A')$,
$Y'_{\lambda}=\operatorname{Spec}(A'_{\lambda})$. Soient $f:X\to Y$ un
morphisme de présentation finie, $\mathcal F$ un $\mathcal O_X$-Module
quasi-cohérent de présentation finie~; on pose $X'=X\times_Y Y'$,
$f'=f_{(Y')}$, $\mathcal F'=\mathcal F\otimes_Y Y'$,
$X'_{\lambda}=X\times_Y Y'_{\lambda}$, $f'_{\lambda}=f_{(Y'_{\lambda})}$,
$\mathcal F'_{\lambda}=\mathcal F\otimes_Y Y'_{\lambda}$. Soit $x$ un point de
$X$, tel que $\mathcal F'$ soit $f'$-plat en tous les points $x'\in X'$
au-dessus de $x$~; alors il existe $\lambda$ tel que $\mathcal F'_{\lambda}$
soit $f'_{\lambda}$-plat en tous les points de $X'_{\lambda}$ au-dessus de $x$.
\end{lemma}

\begin{proof}
Soit $U'$ l'ensemble des $x'\in X'$ tels que $\mathcal F'$ soit $f'$-plat au
point $x'$~; on sait (11.3.1) que $U'$ est ouvert dans $X'$ puisque $f'$ est de
présentation finie (1.6.2)~; de même l'ensemble $U'_{\lambda}$ des points de
$X'_{\lambda}$ où $\mathcal F'_{\lambda}$ est $f'_{\lambda}$-plat est ouvert
dans $X'_{\lambda}$, et on sait en outre (11.2.6) que $U'$ est réunion des
$v_{\lambda}^{-1}(U'_{\lambda})$, où
$v_{\lambda}:X'\to X'_{\lambda}$ est la projection canonique. Considérons le
schéma $T=\operatorname{Spec}(k(x))$~; posons $T'=T\times_Y Y'$,
$T'_{\lambda}=T\times_Y Y'_{\lambda}$, et soient
$w_{\lambda}:T'\to T'_{\lambda}$, $g':T'\to X'$,
$g'_{\lambda}:T'_{\lambda}\to X'_{\lambda}$ les projections canoniques. Posons
$V'_{\lambda}=g_{\lambda}'^{-1}(U'_{\lambda})$,
$V'=g'^{-1}(U')=\bigcup_{\lambda}w_{\lambda}^{-1}(V'_{\lambda})$. Par
hypothèse on a (compte tenu de (\textbf I, 3.6.1)) $V'=T'$~; comme $T'$ est
quasi-compact, il existe $\lambda$ tel que
$w_{\lambda}^{-1}(V'_{\lambda})=T'$. On déduit alors de (8.3.3) appliqué aux
parties fermées quasi-compactes $T'_{\lambda}-V'_{\lambda}$ de
$T'_{\lambda}$, qu'il existe $\mu\geq\lambda$ tel que
$T'_{\mu}=V'_{\mu}$~; cela signifie que $\mathcal F'_{\mu}$ est
$f'_{\mu}$-plat en tous les points de $X'_{\mu}$ dont la projection dans $X$
est $x$. C.Q.F.D.
\end{proof}

\subsection{Descente de la platitude par des morphismes quelconques : cas d'un
préschéma de base unibranche}
\label{subsection:IV.11.6-fr}

\begin{theorem}[11.6.1]
\label{IV.11.6.1-fr}
Soient $A$ un anneau local intègre géométriquement unibranche
$(\mathbf 0,23.2.1)$, $Y=\operatorname{Spec}(A)$, $f:X\to Y$ un morphisme
localement de présentation finie, $\mathcal F$ un $\mathcal O_X$-Module
quasi-cohérent de présentation finie. Soient $A'$ un anneau local,
$u:A\to A'$ un homomorphisme local injectif~; on pose
$Y'=\operatorname{Spec}(A')$, $X'=X\times_Y Y'$, $f'=f_{(Y')}$,
$\mathcal F'=\mathcal F\otimes_Y Y'$. Soient $x$ un point de $X$ dont la
projection $f(x)=y$ est le point fermé de $Y$, $x'$ un point de $X'$ dont les
projections dans $X$ et $Y'$ sont respectivement $x$ et le point fermé $y'$ de
$Y'$. Alors, si $\mathcal F'$ est $f'$-plat au point $x'$, $\mathcal F$ est
$f$-plat au point $x$.
\end{theorem}

\begin{proof}
On peut se borner au cas où $f$ est de présentation finie, la question étant
locale sur $X$. Nous procéderons en plusieurs étapes.

\medskip
\noindent\textbf{I)} \emph{Réduction au cas où $A$ et $A'$ sont des anneaux
locaux intégralement clos.}

Comme $u$ est injectif et $A$ intègre, il existe un idéal premier
$\mathfrak p'$ de $A'$ tel que $u^{-1}(\mathfrak p')=0$
$(\mathbf 0_{\mathrm I},1.5.8)$~; l'homomorphisme composé
$A\xrightarrow{u}A'\to A''=A'/\mathfrak p'$ est donc injectif et local, et si
$Y''=\operatorname{Spec}(A'')$, $X''=X'\otimes_{A'}A''=X\times_Y Y''$,
$\mathcal F''=\mathcal F'\otimes_{A'}A''$, $\mathcal F''$ est $Y''$-plat aux
points de $X''$ au-dessus de $x'$ (2.1.4)~; remplaçant au besoin $A'$ par $A''$
et tenant compte de (\textbf I, 3.4.7), on peut donc supposer d'abord que $A'$
est intègre. Si $K'$ est le corps des fractions de $A'$, il y a alors un anneau
de valuation $B'$ dans $K'$ qui domine $A'$~; l'homomorphisme composé
$A\xrightarrow{u}A'\to B'$ étant injectif et local, le même raisonnement que
% IV3_B23_P154_END

\oldpage[IV-3]{155}

précédemment permet de remplacer $A'$ par $B'$~; on peut donc supposer l'anneau
local $A'$ intégralement clos, $A$ étant un sous-anneau local de $A'$ dominé par
$A'$. Soit $A_1$ la clôture intégrale de $A$~; il est clair que
$A\subset A_1\subset A'$, et par hypothèse $A_1$ est un anneau local~; si
$\mathfrak m$, $\mathfrak m_1$, $\mathfrak m'$ sont les idéaux maximaux de
$A$, $A_1$, $A'$, on a $\mathfrak m\subset\mathfrak m_1\subset\mathfrak m'$~;
en effet, $\mathfrak m_1$ est le \emph{seul} idéal premier de $A_1$ au-dessus de
$\mathfrak m$, puisque $A_1$ est un anneau local (Bourbaki,
\emph{Alg. comm.}, chap.~V, \S~2, n$^\circ$~1, prop.~1)~; comme
$\mathfrak m'\cap A=\mathfrak m$, on a
$\mathfrak m'\cap A_1\cap A=\mathfrak m$, donc
$\mathfrak m'\cap A_1=\mathfrak m_1$. Posons
$Y_1=\operatorname{Spec}(A_1)$, $X_1=X\times_Y Y_1$, $f_1=f_{(Y_1)}$,
$\mathcal F_1=\mathcal F\otimes_Y Y_1$, et soit $x_1$ la projection de $x'$ dans
$X_1$~; désignons d'autre part par $y_1$ l'unique point fermé de $Y_1$, de sorte
que $y_1=f_1(x_1)$. Par hypothèse le morphisme
$\operatorname{Spec}(k(y_1))\to\operatorname{Spec}(k(y))$ est radiciel, d'où
l'on conclut, par la transitivité des fibres (\textbf I, 3.6.4) et
(\textbf I, 3.5.7), que le morphisme $f_1^{-1}(y_1)\to f^{-1}(y)$ est radiciel,
et en particulier que $x_1$ est le \emph{seul} point de $X_1$ dont les
projections dans $X$ et $Y_1$ soient $x$ et $y_1$ respectivement~; en outre, on
a vu que $y_1$ est le seul point de $Y_1$ dont la projection dans $Y$ soit $y$,
donc $x_1$ est le seul point de $X_1$ dont la projection dans $X$ soit $x$. Si
l'on prouve que $\mathcal F_1$ est $f_1$-plat au point $x_1$, on pourra donc
appliquer (11.5.5), d'où résultera la conclusion. On est donc ramené au cas où
$A$ lui aussi est \emph{intégralement clos}.

\medskip
\noindent\textbf{II)} \emph{Réduction au cas où $A$ et $A'$ sont des anneaux
locaux de $\mathbf Z$-algèbres de type fini intégralement clos.}

On peut considérer $A'$ comme limite inductive filtrante de ses
sous-$\mathbf Z$-algèbres (intègres) de type fini $B'_{\lambda}$~; en outre,
comme $A'$ est intégralement clos et que la clôture intégrale d'une
$\mathbf Z$-algèbre de type fini est aussi une $\mathbf Z$-algèbre de type fini
(7.8.3), on voit que $A'$ est limite inductive de ses sous-$\mathbf Z$-algèbres
de type fini $B'_{\lambda}$ \emph{intégralement closes}~; si
$\mathfrak m'_{\lambda}=\mathfrak m'\cap B'_{\lambda}$, $A'$ est aussi limite
inductive des sous-anneaux locaux
$(B'_{\lambda})_{\mathfrak m'_{\lambda}}=A'_{\lambda}$ dominés par $A'$
(5.13.3). Pour tout $\lambda$, $B'_{\lambda}\cap A$ est aussi limite inductive
de ses sous-$\mathbf Z$-algèbres de type fini $B_{\alpha\lambda}$, donc
$A=\varinjlim_{\alpha,\lambda}B_{\alpha\lambda}$, et comme plus haut on peut
remplacer $B_{\alpha\lambda}$ dans cette formule par sa clôture intégrale
(contenue par hypothèse dans $B'_{\lambda}\cap A$), puis par l'anneau local
$A_{\alpha\lambda}=(B_{\alpha\lambda})_{\mathfrak m_{\alpha\lambda}}$, où
$\mathfrak m_{\alpha\lambda}=\mathfrak m\cap B_{\alpha\lambda}
=\mathfrak m'_{\lambda}\cap B_{\alpha\lambda}$, de sorte que
$A_{\alpha\lambda}$ est dominé par $A'_{\lambda}$ et par $A$. Posons
$Y_{\alpha\lambda}=\operatorname{Spec}(A_{\alpha\lambda})$~; il résulte de
(8.9.1) qu'il existe un couple $(\alpha,\lambda)$ assez grand, un morphisme
$f_{\alpha\lambda}:X_{\alpha\lambda}\to Y_{\alpha\lambda}$ de type fini et un
$\mathcal O_{X_{\alpha\lambda}}$-Module cohérent $\mathcal F_{\alpha\lambda}$
tels que
\[
X=X_{\alpha\lambda}\times_{Y_{\alpha\lambda}}Y,
\qquad f=f_{\alpha\lambda}\times1_Y,
\qquad \mathcal F=\mathcal F_{\alpha\lambda}\otimes_{Y_{\alpha\lambda}}Y~;
\]
si $x_{\alpha\lambda}$ est la projection de $x$ dans $X_{\alpha\lambda}$, il
suffira de montrer que $\mathcal F_{\alpha\lambda}$ est
$f_{\alpha\lambda}$-plat au point $x_{\alpha\lambda}$. Comme $A'$ est limite
inductive des $A'_{\mu}$ pour $\mu\geq\lambda$, $X'$ est limite projective des
$X_{\alpha\lambda}\times_{Y_{\alpha\lambda}}Y'_{\mu}=X'_{\mu}$, et on a aussi
$\mathcal F'=\mathcal F'_{\alpha\mu}\otimes_{Y'_{\mu}}Y'$, où
$\mathcal F'_{\alpha\mu}
=\mathcal F_{\alpha\lambda}\otimes_{Y_{\alpha\lambda}}Y'_{\mu}$.
Appliquant (11.2.6), on voit qu'on peut prendre $\mu$ assez grand pour que
$\mathcal F'_{\alpha\mu}$ soit $Y'_{\mu}$-plat au point $x'_{\mu}$, projection
de $x'$ dans $X'_{\mu}$, et en outre, par construction des $A'_{\mu}$, la
projection de $x'_{\mu}$ dans $Y'_{\mu}$ est le point fermé de $Y'_{\mu}$.

\medskip
\noindent\textbf{III)} \emph{Réduction au cas où le corps résiduel $k'$ de
$A'$ est une extension finie et radicielle du corps résiduel $k$ de $A$.}

On peut en premier lieu répéter le raisonnement de la partie I) de la
démonstration de (11.4.11), tenant compte du fait que $\mathbf Z$ est un anneau
de
% IV3_B23_P155_END

\oldpage[IV-3]{156}

Jacobson~; on se réduit ainsi au cas où $k'$ est une extension \emph{finie} de
$k$, ce que l'on va supposer dans ce qui suit. Soient $k''$ la plus grande
extension séparable de $k$ contenue dans $k'$, $k_1$ une extension galoisienne
finie de $k$ contenant $k''$, de sorte que $k''\otimes_k k_1$ est composée
directe de corps isomorphes à $k_1$~; comme $k'$ est une extension radicielle de
$k''$, $k'\otimes_k k_1$ est donc composée directe d'extensions radicielles de
$k_1$. Il existe un anneau local $A_1$ qui est une $A$-algèbre et un $A$-module
\emph{libre de type fini}, tel que $A_1/\mathfrak m A_1$ soit $k$-isomorphe à
$k_1$ $(\mathbf 0_{\mathrm{III}},10.3.1.2)$~; de façon précise, on peut supposer
que $k_1=k[T]/(r)$, où $r$ est un polynôme unitaire irréductible et séparable de
$k[T]$ de degré $n$~; si $R$ est un polynôme unitaire de $A[T]$ dont l'image
canonique est $r$ (et qui est donc de degré $n$), on peut prendre
$A_1=A[T]/(R)$. Or, si $K$ est le corps des fractions de $A$, il est clair que
$R$ est un polynôme irréductible et séparable de $K[T]$~; on en déduit donc
d'abord que $A_1$ est un anneau intègre dont le corps des fractions
$K_1=K[T]/(R)$ est une extension séparable de $K$. En outre, si $t$ est l'image
canonique de $T$ dans $A_1$, les $t^j$ $(0\leq j<n)$ forment une base du
$A$-module $A_1$, et leurs images dans $k_1$ une base sur $k$~; on déduit de là
que
\[
d=\det\bigl(\operatorname{Tr}_{A_1/A}(t^{i+j})\bigr)
\]
est un élément de $A$ dont la classe dans $k$ est $\ne0$, et qui est par suite
inversible. Le même raisonnement que dans (6.12.4.1, I)) prouve alors que le
morphisme $\operatorname{Spec}(A_1)\to\operatorname{Spec}(A)$ est plat et a ses
fibres régulières~; on conclut par suite de (6.5.4, (ii)) que $A_1$ est
\emph{intégralement clos}. Considérons alors l'anneau
$A'_1=A'\otimes_A A_1$~; c'est un $A'$-module libre de type fini, donc un anneau
semi-local (Bourbaki, \emph{Alg. comm.}, chap.~IV, \S~2, n$^\circ$~5, cor.~3 de
la prop.~9)~; en outre, les idéaux maximaux de cette $A'$-algèbre finie sont
tous au-dessus de l'idéal maximal $\mathfrak m'$ de $A'$, et \emph{a fortiori}
contiennent $\mathfrak m A'_1$. Mais
\[
A'_1/\mathfrak m A'_1=(A'/\mathfrak m A')\otimes_k k_1,
\]
et comme $k_1$ est une extension séparable et finie de $k$, le radical de
$A'_1/\mathfrak m A'_1$ est égal à
$(\mathfrak m'/\mathfrak m A')\otimes_k k_1$ (Bourbaki, \emph{Alg.},
chap.~VIII, \S~7, n$^\circ$~2, cor.~2 de la prop.~3)~; si $\mathfrak n_i$
$(1\leq i\leq r)$ sont les idéaux maximaux de $A'_1$, les corps
$A'_1/\mathfrak n_i$ sont donc les corps composant l'algèbre
$k'\otimes_k k_1$, autrement dit ce sont des extensions finies
\emph{radicielles} de $k_1$. D'ailleurs, comme $A\to A'$ est un homomorphisme
injectif, il en est de même de $A_1\to A'_1$, $A_1$ étant un $A$-module plat~;
l'homomorphisme canonique
\[
A'_1\longrightarrow\prod_{i=1}^{r}(A'_1)_{\mathfrak n_i}
\]
étant lui aussi injectif (Bourbaki, \emph{Alg. comm.}, chap.~II, \S~3,
n$^\circ$~3, th.~1), il en est de même du composé de $A_1$ dans ce produit.
Mais $A_1$ est intègre, et les noyaux des homomorphismes
$A_1\to(A'_1)_{\mathfrak n_i}$ sont en nombre fini~; comme leur intersection
est nulle, c'est que l'un d'eux déjà est nul. En d'autres termes, il y a un
$B_1=(A'_1)_{\mathfrak n_i}$ tel que l'homomorphisme $A_1\to B_1$ soit
\emph{injectif} et \emph{local}. Posons $Y'_1=\operatorname{Spec}(B_1)$,
$X'_1=X'\times_{Y'}Y'_1$~; $\mathcal F'_1=\mathcal F'\otimes_{Y'}Y'_1$ est
$Y'_1$-plat en tous les points de $X'_1$ au-dessus de $x'$~; d'ailleurs l'idéal
maximal de $B_1$ est le seul au-dessus de $\mathfrak m'$, par suite tous ces
points ont pour projection dans $Y'_1$ le point fermé $y'_1$. Soit $x'_1$ un de
ces points. Posons d'autre part $Y_1=\operatorname{Spec}(A_1)$,
$X_1=X\times_Y Y_1$, $\mathcal F_1=\mathcal F\otimes_Y Y_1$~; si $x_1$ est la
projection de $x'_1$ dans $X_1$, la projection de $x_1$ dans $X$ est $x$ et sa
projection dans $Y_1$ est le point fermé $y_1$. Si l'on prouve que
$(\mathcal F_1)_{x_1}$ est un $\mathcal O_{y_1}$-module plat, il en résultera
que $\mathcal F_x$ est un $\mathcal O_y$-module plat~; en effet
$\mathcal O_{y_1}$ est un $\mathcal O_y$-module plat, donc
$(\mathcal F_1)_{x_1}$ est un $\mathcal O_y$-module plat
$(\mathbf 0_{\mathrm I},6.2.1)$.
% IV3_B23_P156_END

\oldpage[IV-3]{157}

Mais $(\mathcal F_1)_{x_1}=\mathcal F_x\otimes_{\mathcal O_x}\mathcal O_{x_1}$,
et $\mathcal O_{x_1}$ est un $\mathcal O_x$-module \emph{fidèlement plat}~;
donc (2.2.11, (iii)) $\mathcal F_x$ est un $\mathcal O_y$-module plat. Comme
$X'_1=X_1\times_{Y_1}Y'_1$, $\mathcal F'_1=\mathcal F_1\otimes_{Y_1}Y'_1$,
on est bien ramené à la situation de l'énoncé (11.6.1), avec $A$ et $A'$
remplacés par $A_1$ et $B_1$ respectivement.

\medskip
\noindent\textbf{IV)} \emph{Fin de la démonstration.} --- On est finalement
réduit à prouver (11.6.1) sous les hypothèses supplémentaires suivantes~:
\begin{enumerate}[label=(\roman*)]
  \item $A$ et $A'$ sont des anneaux locaux de $\mathbf Z$-algèbres de type fini
  (donc des anneaux \emph{excellents} (7.8.3))~;
  \item $A$ est intégralement clos~;
  \item le corps résiduel $k'$ de $A'$ est une extension finie et radicielle du
  corps résiduel $k$ de $A$.
\end{enumerate}
On sait alors ((7.8.3) et (2.3.8)) que sous les conditions (i) et (ii), si
$\mathfrak m$ et $\mathfrak m'$ sont les idéaux maximaux de $A$ et $A'$
respectivement, la topologie $\mathfrak m$-adique sur $A$ est induite par la
topologie $\mathfrak m'$-adique de $A'$. Le complété
$\widehat u:\widehat A\to\widehat{A'}$ est donc un homomorphisme injectif.
D'autre part, comme le morphisme
$\operatorname{Spec}(k')\to\operatorname{Spec}(k)$ est radiciel, il en est de
même du morphisme $f'^{-1}(y')\to f^{-1}(y)$ (\textbf I, 3.5.7), et il n'y a
donc qu'un seul point $x'$ de $X'$ dont les projections dans $X$ et $Y'$ soient
$x$ et $y'$ respectivement. On peut donc appliquer le résultat de (11.5.2).
C.Q.F.D.
\end{proof}

\begin{corollary}[11.6.2]
\label{IV.11.6.2-fr}
Soient $A$ un anneau local intègre unibranche, $Y=\operatorname{Spec}(A)$,
$f:X\to Y$ un morphisme localement de présentation finie, $\mathcal F$ un
$\mathcal O_X$-Module quasi-cohérent de présentation finie. Soient $A'$ un
anneau local, $A\to A'$ un homomorphisme local injectif~; on pose
$Y'=\operatorname{Spec}(A')$, $X'=X\times_Y Y'$, $f'=f_{(Y')}$,
$\mathcal F'=\mathcal F\otimes_Y Y'$. Soit $x$ un point de $X$ dont la
projection $f(x)=y$ est le point fermé de $Y$~; on suppose que $\mathcal F'$ est
$f'$-plat en tous les points $x'$ de $X'$ dont les projections dans $X$ et $Y'$
sont respectivement $x$ et le point fermé $y'$ de $Y'$. Alors $\mathcal F$ est
$f$-plat au point $x$.
\end{corollary}

\begin{proof}
On peut en effet reprendre la partie I) de la démonstration de (11.6.1), qui
prouve (avec les mêmes notations) que si $\mathcal F_1$ est $f_1$-plat en tous
les points $x_1$ de $X_1$ dont les projections respectives dans $X$ et $Y_1$
sont $x$ et $y_1$, alors $\mathcal F$ est $f$-plat au point $x$~; on est ainsi
ramené au cas où $A$ est intégralement clos, donc géométriquement unibranche, et
la conclusion résulte alors de (11.6.1).
\end{proof}

\subsection{Contre-exemples}
\label{subsection:IV.11.7-fr}

\begin{env}[11.7.1]
\label{IV.11.7.1-fr}
Considérons d'abord le cas où $A$ est un anneau local \emph{artinien}, et où les
hypothèses de (11.4.11) sont remplies \emph{sauf} la condition (ii) concernant
le corps résiduel $k$ de $A$. Nous allons voir que la conclusion de (11.4.11)
peut alors être en défaut. Soit $k$ un corps admettant une extension
\emph{galoisienne} $k'$ de degré $[k':k]>1$, et désignons par $\Gamma$ le groupe
de Galois de $k'$. Soit $A$ une $k$-algèbre ayant une base de 3 éléments
$1,a,b$ avec la table de multiplication $a^2=ab=ba=b^2=0$, de sorte que $A$ est
un anneau local artinien dont l'idéal maximal $\mathfrak m=ka+kb$ est de carré
nul. Soit $A'=A\otimes_k k'$, qui est une $k'$-algèbre de base $1,a,b$, anneau
local artinien d'idéal maximal
$\mathfrak m'=k'a+k'b=\mathfrak m A'$, de carré nul~; $A$ s'identifie
canoniquement à un sous-anneau de $A'$. Soit $\mathfrak J$ le sous-$k'$-espace
vectoriel de $\mathfrak m'$ engendré par $a+\gamma b$, où $\gamma\in k'$
n'appartient pas à $k$~; il est clair que $\mathfrak J$ est un idéal de $A'$.
Posons $B=A'/\mathfrak J$~; c'est un anneau artinien qui est un $A$-module
\emph{non plat}~; sinon (Bourbaki, \emph{Alg. comm.}, chap.~II, \S~3,
n$^\circ$~2, cor.~2 de la prop.~5), $B$ serait un $A$-module \emph{libre}~;
comme $A'$ est lui aussi un $A$-module libre, et que l'homomorphisme canonique
$A'/\mathfrak m A'\to B/\mathfrak m B$ est bijectif, l'homomorphisme canonique
$A'\to B=A'/\mathfrak J$ serait lui
% IV3_B23_P157_END

\oldpage[IV-3]{158}

aussi bijectif (\emph{loc. cit.}, n$^\circ$~2, cor. de la prop.~6) ce qui est
absurde. Autrement dit, si l'on pose $Y=\operatorname{Spec}(A)$,
$X=\operatorname{Spec}(B)$, $\mathcal O_X$ n'est pas $Y$-plat en l'unique point
$x$ de $X$. Mais soit $A_1=B$, et $Y_1=\operatorname{Spec}(A_1)$, et considérons
l'homomorphisme canonique $A\to A_1$, qui est local et injectif puisque
$\mathfrak J\cap A=0$~; si
$B_1=B\otimes_A A_1=B\otimes_A B$, nous allons voir qu'il y a un point de
$X_1=X\times_Y Y_1=\operatorname{Spec}(B_1)$ où $\mathcal O_{X_1}$ est
$Y_1$-plat. Pour cela, remarquons que l'on a
\[
B\otimes_A B
=\frac{A'\otimes_A A'}
{\operatorname{Im}(\mathfrak J\otimes_A A')
+\operatorname{Im}(A'\otimes_A\mathfrak J)}.
\]
Or la structure de $C'=A'\otimes_A A'$ s'obtient aisément~; on considère la
$A$-algèbre produit $C=\prod_{\sigma\in\Gamma}A'_{\sigma}$, où tous les
$A'_{\sigma}$ sont égaux à $A'$, et l'application canonique
$\varphi:A'\otimes_A A'\to C$ telle que
$\varphi(x\otimes y)=(\sigma(x)y)_{\sigma\in\Gamma}$ (le groupe $\Gamma$
opérant canoniquement sur $A'$)~; en passant aux quotients, on en déduit un
homomorphisme $C'/\mathfrak m C'\to C/\mathfrak m C$ qui n'est autre que
l'homomorphisme canonique
$k'\otimes_k k'\to\prod_{\sigma}k'_{\sigma}$ (avec $k'_{\sigma}=k'$ pour tout
$\sigma\in\Gamma$)~; on sait que ce dernier est bijectif (Bourbaki,
\emph{Alg.}, chap.~VIII, \S~8, prop.~4), donc il en est de même de $\varphi$,
puisque $C'$ et $C$ sont des $A$-modules libres (Bourbaki,
\emph{Alg. comm.}, chap.~II, \S~3, n$^\circ$~2, cor. de la prop.~6). De ce qui
précède, il résulte que $B\otimes_A B$ est une $A$-algèbre semi-locale, composée
directe des $A$-algèbres locales $A'/(\sigma(\mathfrak J)+\mathfrak J)$. Celle
de ces algèbres correspondant à l'identité de $\Gamma$ est isomorphe à
$A'/\mathfrak J$, donc est un $A_1$-module plat, mais comme $\gamma\notin k$, il
y a au moins un $\sigma\in\Gamma$ tel que
$\sigma(\mathfrak J)\ne\mathfrak J$~; alors
$\sigma(\mathfrak J)+\mathfrak J=\mathfrak m'$, et $A'/\mathfrak m'$ n'est pas
un $(A'/\mathfrak J)$-module plat, puisque c'est le quotient de
$A'/\mathfrak J$ par un idéal non nul.
\end{env}

\begin{env}[11.7.2]
\label{IV.11.7.2-fr}
Nous allons maintenant montrer que le résultat de (11.5.4) perd sa validité
lorsqu'on ne suppose plus que $f$ soit un morphisme \emph{propre} (et
\emph{a fortiori} (11.5.3) cesse d'être exact lorsqu'on ne suppose plus $g$
propre). Soient $k$ un corps, $A_0$ l'anneau de polynômes $k[S,T]$, $A$
l'anneau quotient $A_0/A_0ST$~; $Y=\operatorname{Spec}(A)$ est donc la courbe
réductible formée des deux «~axes de coordonnées~» dans le plan affine
$V_k^2=\operatorname{Spec}(A_0)$. L'anneau $A$ admet deux idéaux premiers
minimaux $\mathfrak p_1=A_0S/A_0ST$, $\mathfrak p_2=A_0T/A_0ST$, et comme $A$
est réduit, il se plonge canoniquement dans $B=B_1\oplus B_2$, où
$B_1=A/\mathfrak p_1$, $B_2=A/\mathfrak p_2$~; d'ailleurs $B_1$ s'identifie à
$k[T]$ et $B_2$ à $k[S]$, donc ce sont des anneaux intègres intégralement clos
et par suite $Z=\operatorname{Spec}(B)$ n'est autre que le normalisé du
préschéma $Y$ par rapport à $R(Y)$ (\textbf{II}, 6.3.8), somme des deux schémas
$Z_1=\operatorname{Spec}(B_1)$, $Z_2=\operatorname{Spec}(B_2)$. Désignons par
$y$ le «~point double~» de $Y$, correspondant à l'idéal maximal
$\mathfrak p_1+\mathfrak p_2=\mathfrak m$ de $A$, par $z_1$ et $z_2$ les
points de $Z$ qui se projettent en $y$, correspondant respectivement aux idéaux
maximaux $\mathfrak m_1=(T)$ et $\mathfrak m_2=(S)$ de $B_1$ et $B_2$. Nous
désignerons par $X$ le sous-préschéma de $Z$ induit sur le complémentaire de
$z_2$ dans $Z$~; on a donc
$X=\operatorname{Spec}(B_1\oplus(B_2)_S)$~; il est immédiat que
l'homomorphisme $A\to A'=B_1\oplus(B_2)_S$ est injectif, mais le morphisme
correspondant $f:X\to Y$ n'est pas fermé (car $f(Z_2-\{z_2\})$ n'est pas fermé
dans $Y$, bien que $Z_2-\{z_2\}$ soit fermé dans $X$)~; \emph{a fortiori} il
n'est pas propre. Nous allons voir maintenant que $f$ n'est pas plat au point
$z_1$~; il suffira de montrer $(\mathbf 0_{\mathrm I},6.6.2)$ que
$\mathcal O_{z_1}$ n'est pas un $\mathcal O_y$-module fidèlement plat, et pour
cela il suffit de voir que l'homomorphisme canonique
$\mathcal O_y\to\mathcal O_{z_1}$ n'est pas injectif~; mais cela est immédiat
car $\mathcal O_{z_1}$ est un anneau intègre, tandis que $\mathcal O_y$ admet
des diviseurs de zéro. Cependant, la première projection
$p:X\times_Y X\to X$ est un isomorphisme~: en effet, on a
\[
B_1\otimes_A B_1=(A/\mathfrak p_1)\otimes_A(A/\mathfrak p_1)=A/\mathfrak p_1,
\]
$(B_2)_S\otimes_A(B_2)_S=(B_2\otimes_A B_2)_S=(B_2)_S$ pour la même raison,
et enfin $B_1\otimes_A(B_2)_S=0$, car l'image canonique de $S$ dans $B_1$ est
nulle.
\end{env}

\begin{env}[11.7.3]
\label{IV.11.7.3-fr}
L'exemple précédent peut se généraliser~: on considère sur un corps $k$ une
courbe algébrique réduite $Y$ admettant un seul «~point double ordinaire~» $y$
(notion qui sera définie plus tard de façon générale), et sa normalisée $Z$, de
sorte que le morphisme $g:Z\to Y$ est fini, que la restriction de $g$ à
$Z-g^{-1}(y)$ est un isomorphisme sur $Y-\{y\}$, et que $g^{-1}(y)$ se réduit à
deux points «~simples~» $z_1,z_2$~; en outre le préschéma $g^{-1}(y)$ est somme
des deux préschémas $\operatorname{Spec}(k(z_1))$,
$\operatorname{Spec}(k(z_2))$, canoniquement isomorphes à
$\operatorname{Spec}(k(y))$. Soit $X$ le sous-préschéma de $Z$ induit sur
l'ouvert $Z-\{z_2\}$~; le morphisme $f:X\to Y$, restriction de $g$ à $X$, n'est
pas propre, sans quoi (\textbf{II}, 5.4.3) il en serait de même de l'injection
canonique $j:X\to Z$, qui n'est pas fermée. Le morphisme $f$ est radiciel, car
pour tout $y'\in Y$, la fibre $f^{-1}(y')$ ne comprend qu'un seul point $x'$,
$k(y')\to k(x')$ étant un isomorphisme~; on en conclut d'abord que le morphisme
diagonal $\Delta_f:X\to X\times_Y X$ est bijectif (1.7.7.1) et d'autre part,
comme $f$ est non ramifié (17.4.2, d$'$), $\Delta_f$ est une immersion ouverte
(17.4.2, b$'$)~; par suite $\Delta_f$ est un isomorphisme, et la première
projection $p:X\times_Y X\to X$ l'isomorphisme réciproque. Cependant $f$ n'est
pas plat au point $z_1$~; sinon $\mathcal O_{z_1}$ serait un
$\mathcal O_y$-module fidèlement plat $(\mathbf 0_{\mathrm I},6.6.2)$, et comme
$\mathcal O_y$ contient deux idéaux premiers minimaux $\mathfrak p_1$,
$\mathfrak p_2$ distincts (correspondant aux deux «~branches~» de $Y$ au point
$y$) il existerait dans $\mathcal O_{z_1}$ deux idéaux premiers dont les images
réciproques par $u:\mathcal O_y\to\mathcal O_{z_1}$ seraient
$\mathfrak p_1$ et $\mathfrak p_2$ $(\mathbf 0_{\mathrm I},6.5.1)$~; mais cela
est absurde, car $\mathcal O_{z_1}$ n'a que deux idéaux premiers distincts, $0$
et l'idéal maximal $\mathfrak m_1$, et $u^{-1}(\mathfrak m_1)$ est l'idéal
maximal $\mathfrak m$ de $\mathcal O_y$.
\end{env}

\begin{env}[11.7.4]
\label{IV.11.7.4-fr}
On notera que dans l'exemple précédent l'homomorphisme $u$ est \emph{injectif}
lorsque $Y$ est \emph{irréductible} (on peut par exemple prendre
$Y=\operatorname{Spec}(k[S,T]/(S(S^2+T^2)-(S^2-T^2)))$, «~cubique à point
double~»)~; on peut en effet alors (en remplaçant $Y$ par un voisinage affine de
$y$) supposer que $Y=\operatorname{Spec}(A)$, où $A$ est intègre, d'où
$Z=\operatorname{Spec}(B)$, où $B$ est la clôture intégrale de $A$~; comme $B$,
$\mathcal O_y$ et $\mathcal O_{z_1}$ s'identifient alors à des sous-anneaux du
corps
% IV3_B23_P158_END

\oldpage[IV-3]{159}

des fractions de $A$, $u$ est évidemment injectif. On notera par contre que
l'homomorphisme $\widehat u:\widehat{\mathcal O}_y\to
\widehat{\mathcal O}_{z_1}$ n'est pas injectif, car
$\widehat{\mathcal O}_{z_1}$ est un anneau intègre local ($z_1$ étant un point
simple), tandis que $\widehat{\mathcal O}_y$ a deux idéaux premiers minimaux
distincts (correspondant aux deux «~branches~» de $Y$) et admet donc des
diviseurs de zéro. Cela donne un exemple montrant que dans l'énoncé (11.5.2), on
ne peut pas remplacer l'hypothèse que $\widehat u$ est injectif par l'hypothèse
que $u$ lui-même est injectif, même lorsque $A'$ est un anneau local. Il suffit
en effet de prendre (avec les notations précédentes) $A=\mathcal O_y$,
$A'=\mathcal O_{z_1}$, $Y=\operatorname{Spec}(A)$,
$X=\operatorname{Spec}(A')$~; le raisonnement de (11.7.3) prouve encore que la
première projection $p:X\times_Y X\to X$ est un isomorphisme, bien que $f$ ne
soit pas plat au point $z_1$.
\end{env}

\begin{env}[11.7.5]
\EGAReaderTarget{ega4.c04.statement.11-7-5}{IV.11.7.5-fr}
Les exemples de (11.7.2) et (11.7.3) expliquent la restriction aux anneaux
locaux \emph{unibranches} dans (11.6.1) et (11.6.2). Nous allons voir maintenant
que dans (11.6.1), on ne peut affaiblir l'hypothèse sur $A$ en supposant
seulement $A$ unibranche. Considérons en effet l'anneau local intègre complet
$A=\mathbf R[[U,V]]/(U^2+V^2)$ qui est unibranche mais non géométriquement
unibranche (6.5.11). On sait (\emph{loc. cit.}) que si $u,v$ sont les images de
$U$ et $V$ dans $A$, la clôture intégrale de $A$ est $\overline A=A[t]$ avec
$t=v/u$, tel que $t^2=-1$, si bien que $\overline A$ est isomorphe à
$\mathbf C[[U]]$. Posons $Y=\operatorname{Spec}(A)$,
$X=\operatorname{Spec}(\overline A)$ (normalisé de $Y$ (\textbf{II}, 6.3.8)) et
soient $y$ et $x$ les points fermés de $Y$ et $X$ respectivement~; nous allons
montrer que pour une $A$-algèbre locale convenable $A'$, si l'on pose
$Y'=\operatorname{Spec}(A')$, $X'=X\times_Y Y'$, et si $y'$ désigne le point
fermé de $Y'$, $\mathcal O_{X'}$ est $Y'$-plat en un point de $X'$ dont les
projections dans $X$ et $Y'$ sont respectivement $x$ et $y'$, mais n'est pas
$Y'$-plat en tous les points ayant ces projections~; il en résultera évidemment
(2.1.4) que $\mathcal O_X$ n'est pas $Y$-plat au point $x$ (ce qui est d'ailleurs
trivial \emph{a priori}, $\overline A$ n'étant pas un $A$-module libre).

Soit $B=A\otimes_{\mathbf R}\mathbf C$, isomorphe à
$\mathbf C[[U,V]]/(U+iV)(U-iV)$~; $B$ a deux idéaux premiers minimaux
$\mathfrak p'$, $\mathfrak p''$ engendrés respectivement par $u+iv$ et $u-iv$,
et $\mathfrak n=\mathfrak p'+\mathfrak p''$ est l'idéal maximal de l'anneau
local complet $B$. Soit $\overline B=\overline A\otimes_{\mathbf R}\mathbf C$~;
$\overline B$ est composée directe de deux algèbres isomorphes à
$\mathbf C[[U]]$, engendrées par les idempotents
$e'=(1\otimes1+t\otimes i)/2$ et $e''=(1\otimes1-t\otimes i)/2$~; comme
l'homomorphisme $A\to\overline A$ est injectif, il en est de même de
$B\to\overline B$, et les images de $u+iv$ et de $u-iv$ par cette injection sont
respectivement $ue'$ et $ue''$~; on en conclut aussitôt que $\overline B$
s'identifie canoniquement à $(B/\mathfrak p')\oplus(B/\mathfrak p'')$. Cela
étant, prenons pour $A'$ l'anneau local $B/\mathfrak p'$~; alors
$\overline A\otimes_A A'$ s'identifie à $\overline B\otimes_B A'$. Mais on a
$(B/\mathfrak p')\otimes_B(B/\mathfrak p')=B/\mathfrak p'$ et
$(B/\mathfrak p')\otimes_B(B/\mathfrak p'')=B/\mathfrak n$, donc
$\overline A\otimes_A A'$ est isomorphe à $A'\oplus(B/\mathfrak n)$. Cela
établit notre assertion, car $B/\mathfrak n=A'/(\mathfrak n/\mathfrak p')$
n'est pas un $A'$-module plat (sans quoi ce serait un $A'$-module libre
(Bourbaki, \emph{Alg. comm.}, chap.~II, \S~3, n$^\circ$~2, cor.~2 de la
prop.~5), ce qui est absurde).
\end{env}

\subsection{Un critère valuatif de platitude}
\label{subsection:IV.11.8-fr}

\begin{theorem}[11.8.1]
\label{IV.11.8.1-fr}
Soient $f:X\to Y$ un morphisme localement de présentation finie, $\mathcal F$ un
$\mathcal O_X$-Module quasi-cohérent de présentation finie, $x$ un point de
$X$, $y=f(x)$. On suppose l'anneau local $\mathcal O_y$ intègre (resp. réduit et
noethérien). Pour que $\mathcal F$ soit $f$-plat au point $x$, il faut et il
suffit que, pour tout anneau de valuation (resp. tout anneau de valuation
discrète) $A'$ et tout homomorphisme local $\mathcal O_y\to A'$, la condition
suivante soit satisfaite~: en posant $Y'=\operatorname{Spec}(A')$,
$X'=X\times_Y Y'$, $f'=f_{(Y')}$, le $\mathcal O_{X'}$-Module
$\mathcal F'=\mathcal F\otimes_Y Y'$ est $f'$-plat en tous les points $x'$ de
$X'$ dont les projections respectives dans $X$ et $Y'$ sont $x$ et le point
fermé $y'$ de $Y'$.
\end{theorem}

\begin{proof}
La condition étant évidemment nécessaire (2.1.4), reste à prouver qu'elle est
suffisante. On peut évidemment ((\textbf I, 3.6.5) et (\textbf I, 2.4.4)) se
borner au cas où $Y=\operatorname{Spec}(A)$ est le spectre d'un anneau local $A$
et $y$ le point fermé de $Y$.

\medskip
\noindent\textbf{(i)} \emph{Cas où $A$ est intègre.} --- Soient $K$ le corps
des fractions de $A$, $A_1$ la clôture intégrale de $A$~; si l'on pose
$Y_1=\operatorname{Spec}(A_1)$, $X_1=X\times_Y Y_1$, $f_1=f_{(Y_1)}$, il suffit,
en vertu de (11.5.5), de montrer que
$\mathcal F_1=\mathcal F\otimes_Y Y_1$ est $f_1$-plat en \emph{tous} les points
$x_1$ de $X_1$ dont $x$ est la projection dans $X$. Or, si $f_1(x_1)=y_1$, on a
$\mathfrak i_{y_1}=\mathfrak m_1$, où $\mathfrak m_1$ est un idéal premier de
$A_1$ dont $\mathfrak m=\mathfrak i_y$ est la trace sur $A$ ($\mathfrak m_1$
est d'ailleurs nécessairement un idéal maximal). Soit alors $A'$ un anneau de
valuation pour $K$ qui domine $\mathcal O_{y_1}=(A_1)_{\mathfrak m_1}$~;
l'homomorphisme $A\to\mathcal O_{y_1}$ étant local, il en est de même de
$A\to A'$. Il existe alors au moins un
% IV3_B23_P159_END

\oldpage[IV-3]{160}

point $x'\in X'$ dont les projections dans $X_1$ et $Y'$ sont respectivement
$x_1$ et $y'$ (\textbf I, 3.4.7)~; comme par hypothèse $\mathcal F'$ est
$f'$-plat au point $x'$, et que $\mathcal O_{y_1}$ est intégralement clos, on
peut appliquer (11.6.1), et on en déduit que $\mathcal F_1$ est $f_1$-plat au
point $x_1$, d'où le théorème dans ce cas.

\medskip
\noindent\textbf{(ii)} \emph{Cas où $A$ est réduit et noethérien.} --- Soient
$\mathfrak p_i$ $(1\leq i\leq m)$ les idéaux minimaux de $A$~; comme $A$ est
réduit, il s'identifie canoniquement à un sous-anneau du produit des
$A_i=A/\mathfrak p_i$, qui sont des anneaux locaux noethériens~; posant
$Y_i=\operatorname{Spec}(A_i)$, $X_i=X\times_Y Y_i$, $f_i=f_{(Y_i)}$, il
résulte de (11.5.2) qu'il suffit de démontrer que pour chaque $i$,
$\mathcal F_i=\mathcal F\otimes_Y Y_i$ est $f_i$-plat en tout point $x_i$ de
$X_i$ dont les projections dans $X$ et $Y_i$ sont $x$ et le point fermé $y_i$
de $Y_i$ respectivement. Or, comme $A_i$ est un anneau local noethérien intègre,
il existe dans son corps des fractions $K_i$ un sous-anneau $A''_i$ contenant
$A_i$, qui soit une $A_i$-algèbre \emph{finie} (donc un anneau semi-local
noethérien) et dont les anneaux locaux soient \emph{géométriquement
unibranches} ((6.15.5) et $(\mathbf 0,23.2.5)$). Comme les idéaux maximaux de
$A''_i$ sont alors nécessairement au-dessus de l'idéal maximal de $A_i$, on
déduit encore de (\textbf I, 3.4.7) et de (11.5.2) qu'il suffit, en posant
$Y''_i=\operatorname{Spec}(A''_i)$, $X''_i=X\times_Y Y''_i$,
$f''_i=f_{(Y''_i)}$, de prouver que
$\mathcal F''_i=\mathcal F\otimes_Y Y''_i$ est $f''_i$-plat en tout point
$x''_i$ de $X''_i$ dont les projections dans $X$ et $Y''_i$ sont $x$ et le
point fermé $y''_i$ de $Y''_i$ respectivement. Or, soit $A'$ un anneau de
valuation discrète pour $K_i$ dominant $A''_i$, et soit $x'$ un point de $X'$
dont les projections dans $X''_i$ et dans $Y'$ sont $x''_i$ et $y'$
respectivement (\textbf I, 3.4.7)~; comme $\mathcal O_{y''_i}$ est
géométriquement unibranche, on peut encore appliquer (11.6.1) et on en déduit
bien que $\mathcal F''_i$ est $f''_i$-plat au point $x''_i$.
\end{proof}

\begin{remarks}[11.8.2]
\label{IV.11.8.2-fr}
\begin{enumerate}[label=(\roman*)]
  \item Dans l'énoncé de (11.8.1), on peut se borner à supposer que la condition
  sur $\mathcal F'$ est vérifiée pour les anneaux de valuation $A'$ \emph{complets}
  et dont le corps résiduel est \emph{algébriquement clos}. On sait en effet que
  tout anneau de valuation $A'$ est dominé par un tel anneau $A''$
  (\textbf{II}, 7.1.2), et que si $A'$ est un anneau de valuation discrète, on
  peut supposer qu'il en est de même de $A''$
  $(\mathbf 0_{\mathrm{III}},10.3.1)$.
  \item La démonstration de (11.8.1) se simplifie lorsqu'on suppose non seulement
  que $A$ est intègre et noethérien, mais que son complété $\widehat A$ est lui
  aussi \emph{intègre}. Remplaçant $X$ par $X\otimes_A\widehat A$ et raisonnant
  comme dans la démonstration de (11.5.3), on peut en effet se ramener alors à
  prouver (11.8.1) lorsque $A=\mathcal O_y$ est intègre, noethérien et
  \emph{complet}. Or, on sait (\textbf{II}, 7.1.7) qu'un tel anneau $A$ est
  dominé par un anneau de valuation discrète complet~; la conclusion résulte
  donc directement de (11.5.2).
\end{enumerate}
\end{remarks}

\subsection{Familles séparantes et universellement séparantes d'homomorphismes
de faisceaux de modules}
\label{subsection:IV.11.9-fr}

\begin{env}[11.9.1]
\label{IV.11.9.1-fr}
Soient $X$ un préschéma, $(f_{\lambda})_{\lambda\in L}$ une famille de morphismes
$f_{\lambda}:Z_{\lambda}\to X$, $\mathcal F$ un $\mathcal O_X$-Module
quasi-cohérent~; pour tout $\lambda\in L$, supposons donnés un
$\mathcal O_{Z_{\lambda}}$-Module quasi-cohérent $\mathcal G_{\lambda}$ et un
homomorphisme
\begin{equation}
u_{\lambda}:\mathcal F\longrightarrow(f_{\lambda})_*(\mathcal G_{\lambda}).
\label{IV.11.9.1.1-fr}
\tag{11.9.1.1}
\end{equation}
\end{env}
% IV3_B23_P160_END

\oldpage[IV-3]{161}

On dit que la famille $(u_\lambda)$ (ou la famille correspondante des
$u^\#_\lambda:(f_\lambda)^*(\mathcal F)\to\mathcal G_\lambda$) est
\emph{séparante} si l'intersection des noyaux des $u_\lambda$ est
\emph{nulle}. En d'autres termes, cela signifie que pour \emph{tout} ouvert
$U$ de $X$, et toute section $t$ de $\mathcal F$ au-dessus de $U$, telle que,
pour \emph{tout} $\lambda$, la section $u_\lambda(t)$ (qui, par définition est
une section de $\mathcal G_\lambda$ au-dessus de $f_\lambda^{-1}(U)$) soit
nulle, alors $t$ est elle-même nulle.

\begin{env}[11.9.2]
\label{IV.11.9.2-fr}
Avec les notations de (11.9.1), soit $M$ un second ensemble d'indices~; pour
tout $\lambda\in L$, soit $(g_{\lambda\mu})_{\mu\in M}$ une famille de
morphismes $g_{\lambda\mu}:Z'_{\lambda\mu}\to Z_\lambda$~; pour tout couple
$(\lambda,\mu)$, supposons donnés un $\mathcal O_{Z'_{\lambda\mu}}$-Module
quasi-cohérent $\mathcal H_{\lambda\mu}$ et un homomorphisme
$v_{\lambda\mu}:\mathcal G_\lambda\to
(g_{\lambda\mu})_*(\mathcal H_{\lambda\mu})$~; posons
$h_{\lambda\mu}=f_\lambda\circ g_{\lambda\mu}:Z'_{\lambda\mu}\to X$ et
considérons l'homomorphisme composé
\[
w_{\lambda\mu}:\mathcal F\xrightarrow{u_\lambda}
(f_\lambda)_*(\mathcal G_\lambda)
\xrightarrow{(f_\lambda)_*(v_{\lambda\mu})}
(h_{\lambda\mu})_*(\mathcal H_{\lambda\mu}).
\]
Supposons que, pour tout $\lambda\in L$, la famille
$(v_{\lambda\mu})_{\mu\in M}$ soit séparante~: alors, il en est de même de la
famille des $(f_\lambda)_*(v_{\lambda\mu})$ ($\mu\in M$), comme on le voit
aussitôt. On en conclut que, pour que la famille $(u_\lambda)$ soit séparante,
il faut et il suffit que la famille
$(w_{\lambda\mu})_{(\lambda,\mu)\in L\times M}$ le soit.
\end{env}

\begin{env}[11.9.3]
\label{IV.11.9.3-fr}
Pour vérifier que la famille $(u_\lambda)$ est séparante (avec les notations de
(11.9.1)), on peut évidemment se ramener tout d'abord au cas où $X$ est
\emph{affine}, la propriété étant locale sur $X$. On peut en outre supposer que
$Z_\lambda=X$ \emph{pour tout} $\lambda\in L$. En effet, soit $M$ un ensemble
d'indices, somme d'une famille $(M_\lambda)_{\lambda\in L}$, et pour tout
$\lambda\in L$, soit $(Y_{\lambda\mu})_{\mu\in M_\lambda}$ un recouvrement
ouvert affine de $Z_\lambda$~; soit
$j_{\lambda\mu}:Y_{\lambda\mu}\to Z_\lambda$ l'injection canonique et posons
$\mathcal H_{\lambda\mu}=j_{\lambda\mu}^*(\mathcal G_\lambda)
=\mathcal G_\lambda|Y_{\lambda\mu}$. Si l'on considère l'homomorphisme
canonique
$v_{\lambda\mu}=\rho_{\mathcal G_\lambda}:\mathcal G_\lambda\to
(j_{\lambda\mu})_*(j_{\lambda\mu}^*(\mathcal G_\lambda))
=(j_{\lambda\mu})_*(\mathcal H_{\lambda\mu})$ relatif à $j_{\lambda\mu}$
$(\mathbf 0_{\mathrm I},4.4.3.2)$, il est immédiat que pour chaque
$\lambda\in L$, la famille $(v_{\lambda\mu})_{\mu\in M_\lambda}$ est
séparante. En vertu de (11.9.2), on est donc ramené à prouver que la famille des
homomorphismes composés $((f_\lambda)_*(v_{\lambda\mu}))\circ u_\lambda$ est
séparante, autrement dit on est ramené au cas où les $Z_\lambda$ sont affines.
Mais alors les $(f_\lambda)_*(\mathcal G_\lambda)$ sont des
$\mathcal O_X$-Modules quasi-cohérents (\textbf I, 1.6.2) et en vertu de la
définition, on peut remplacer les $Z_\lambda$ par $X$ et les
$\mathcal G_\lambda$ par les $(f_\lambda)_*(\mathcal G_\lambda)$, d'où notre
assertion.

On notera en outre que si $L$ est \emph{fini} et les $f_\lambda$
\emph{quasi-compacts}, on peut, dans la réduction précédente, supposer que les
$M_\lambda$ sont aussi \emph{finis}, donc on est dans ce cas ramené à vérifier
qu'une famille \emph{finie} d'homomorphismes de $\mathcal F$ dans des
$\mathcal O_X$-Modules quasi-cohérents est séparante.
\end{env}

\begin{env}[11.9.4]
\label{IV.11.9.4-fr}
Considérons donc le cas où $Z_\lambda=X$ pour tout $\lambda$, et où
$X=\operatorname{Spec}(A)$ est affine~; alors on a
$\mathcal F=\widetilde M$ et $\mathcal G_\lambda=\widetilde N_\lambda$, où
$M$ et $N_\lambda$ sont des $A$-modules, et $u_\lambda=\widetilde\varphi_\lambda$,
où les $\varphi_\lambda:M\to N_\lambda$ sont des $A$-homomorphismes. Dire que
la famille $(u_\lambda)$ est séparante signifie alors que, pour tout $s\in A$,
l'intersection des noyaux des
$(\varphi_\lambda)_s:M_s\to(N_\lambda)_s$ est réduite à 0. On dit alors aussi
que la famille $(\varphi_\lambda)$ est \emph{séparante}. On notera que si $L$
est \emph{fini}, il revient au même de dire que l'intersection des noyaux des
$\varphi_\lambda$ est 0, car on a alors
\[
\bigcap_{\lambda\in L}\operatorname{Ker}((\varphi_\lambda)_s)
=\left(\bigcap_{\lambda\in L}\operatorname{Ker}(\varphi_\lambda)\right)_s
\quad(\mathbf 0_{\mathrm I},1.3.2).
\]
Mais cette relation n'est plus exacte en général lorsque $L$ est \emph{infini},
et le fait que l'intersection des noyaux des $\varphi_\lambda$ soit 0
\emph{n'entraîne donc pas},
\oldpage[IV-3]{162}
en général, que la famille $(\varphi_\lambda)$ soit séparante. Par exemple,
supposons que $A$ soit un anneau de valuation discrète, d'idéal maximal
$\mathfrak m$, et considérons la famille des homomorphismes
$\varphi_k:A\to A/\mathfrak m^k$, dont l'intersection des noyaux
$\mathfrak m^k$ est réduite à 0~; cette famille n'est toutefois pas séparante,
car les fibres de tous les $\mathfrak m^k$ au point générique $x$ de
$X=\operatorname{Spec}(A)$ (qui est ouvert dans $X$) sont égales à
$\mathcal O_x=k(x)$, corps des fractions de $A$, et leur intersection n'est
donc pas réduite à 0.
\end{env}

\begin{env}[11.9.5]
\label{IV.11.9.5-fr}
Nous allons principalement nous occuper dans ce qui suit du problème du
\emph{changement de base} pour les familles séparantes. Les notations étant
celles de (11.9.1), considérons un morphisme $g:X'\to X$ et posons
$\mathcal F'=\mathcal F\otimes_{\mathcal O_X}\mathcal O_{X'}
=g^*(\mathcal F)$, et, pour tout $\lambda$,
$Z'_\lambda=Z_\lambda\times_X X'$, $f'_\lambda=f_\lambda\times 1$,
$\mathcal G'_\lambda=\mathcal G_\lambda\otimes_{\mathcal O_{Z_\lambda}}
\mathcal O_{Z'_\lambda}=(g'_\lambda)^*(\mathcal G_\lambda)$, où
$g'_\lambda:Z'_\lambda\to Z_\lambda$ est la projection canonique. Pour tout
$\lambda$, on désigne alors par
$u'_\lambda:\mathcal F'\to(f'_\lambda)_*(\mathcal G'_\lambda)$
l'homomorphisme obtenu comme suit~: soit
\[
\psi:(f_\lambda)_*(\mathcal G_\lambda)\longrightarrow
(f_\lambda)_*((g'_\lambda)_*((g'_\lambda)^*(\mathcal G_\lambda)))
=g_*((f'_\lambda)_*(\mathcal G'_\lambda))
\]
l'homomorphisme $(f_\lambda)_*(\rho_{\mathcal G_\lambda})$, où
$\rho_{\mathcal G_\lambda}$ est l'homomorphisme canonique
$(\mathbf 0_{\mathrm I},4.4.3.2)$ correspondant à $g'_\lambda$. Alors
$u'_\lambda$ est défini comme le composé
\[
g^*(\mathcal F)\xrightarrow{g^*(u_\lambda)}
g^*((f_\lambda)_*(\mathcal G_\lambda))
\xrightarrow{g^*(\psi)}g^*g_*((f'_\lambda)_*(\mathcal G'_\lambda))
\xrightarrow{\sigma}(f'_\lambda)_*(\mathcal G'_\lambda)
\]
où $\sigma=\sigma_{(f'_\lambda)_*(\mathcal G'_\lambda)}$ est
l'homomorphisme canonique $(\mathbf 0_{\mathrm I},4.4.3.3)$ correspondant à
$g$. On dira pour abréger que $u'_\lambda$ est \emph{déduit de $u_\lambda$
par le changement de base} $g$. Lorsque $f_\lambda:Z_\lambda\to X$ est un
morphisme \emph{affine}, on a
$(f'_\lambda)_*(\mathcal G'_\lambda)=g^*((f_\lambda)_*(\mathcal G_\lambda))$,
et dans ce cas on a donc simplement $u'_\lambda=g^*(u_\lambda)$.

On peut aussi interpréter $u'_\lambda$ de la façon suivante~: il suffit de
connaître la valeur de $u'_\lambda(t')$, lorsque $t'$ est une section de
$\mathcal F'$ au-dessus d'un ouvert $U'$ de $X'$, du type particulier
suivant~: $t'$ est la restriction à $U'$ de l'image canonique par
$\Gamma(U,\mathcal F)\to\Gamma(g^{-1}(U),\mathcal F')$ d'une section $t$ de
$\mathcal F$ au-dessus d'un ouvert $U$ de $X$ contenant $g(U')$ (ces sections
engendrant en effet le $\mathcal O_{X'}$-Module $\mathcal F'$
$(\mathbf 0_{\mathrm I},3.7.1)$). Considérons la section $u_\lambda(t)$ de
$\mathcal G_\lambda$ au-dessus de $f_\lambda^{-1}(U)$, et son image canonique
$t''$ par
$\Gamma(f_\lambda^{-1}(U),\mathcal G_\lambda)\to
\Gamma((g'_\lambda)^{-1}(f_\lambda^{-1}(U)),\mathcal G'_\lambda)$~; alors
$u'_\lambda(t')$ est la restriction de $t''$ à $(f'_\lambda)^{-1}(U')$.
Considérons en particulier le cas où $Z_\lambda$ est un sous-préschéma induit
sur un ouvert de $X$, $\mathcal G_\lambda=j_\lambda^*(\mathcal F)$, où
$j_\lambda:Z_\lambda\to X$ est l'injection canonique, et où $u_\lambda$ est
l'homomorphisme canonique
$\rho_{\mathcal F}:\mathcal F\to(j_\lambda)_*(j_\lambda^*(\mathcal F))
=(j_\lambda)_*(\mathcal G_\lambda)$ $(\mathbf 0_{\mathrm I},4.4.3.2)$
correspondant à $j_\lambda$. Alors $Z'_\lambda$ est induit sur un ouvert de
$X'$, et l'interprétation précédente montre que $u'_\lambda$ n'est autre que
l'homomorphisme canonique
$\rho_{\mathcal F'}:\mathcal F'\to
(j'_\lambda)_*((j'_\lambda)^*(\mathcal F'))
=(j'_\lambda)_*(\mathcal G'_\lambda)$ correspondant à l'injection canonique
$j'_\lambda:Z'_\lambda\to X'$.
\end{env}

\begin{env}[11.9.6]
\label{IV.11.9.6-fr}
Sous les conditions de (11.9.5), supposons que $X$ et $X'$ soient affines, et
que l'on veuille prouver que pour toute section $t'$ de $\mathcal F'$ au-dessus
de $X'$, dont les images par tous les $u'_\lambda$ sont nulles, alors $t'$ est
elle-même nulle. Alors, on peut aussi se borner au cas où $Z_\lambda=X$ pour
tout $\lambda\in L$. En effet, avec les notations de (11.9.3) et de (11.9.5),
si l'on pose $Y'_{\lambda\mu}=(g'_\lambda)^{-1}(Y_{\lambda\mu})$,
l'homomorphisme $v'_{\lambda\mu}$ déduit de $v_{\lambda\mu}$ par le changement
de base $g$ n'est autre que l'homomorphisme canonique
\oldpage[IV-3]{163}
$\rho_{\mathcal G'_\lambda}:\mathcal G'_\lambda\to
(j'_{\lambda\mu})_*((j'_{\lambda\mu})^*(\mathcal G'_\lambda))$
correspondant à l'injection canonique
$j'_{\lambda\mu}:Y'_{\lambda\mu}\to Z'_\lambda$, comme on l'a vu dans
(11.9.5). L'assertion résulte alors des raisonnements de (11.9.2) et (11.9.3),
$Y_{\lambda\mu}$ et $v_{\lambda\mu}$ étant remplacés par
$Y'_{\lambda\mu}$ et $v'_{\lambda\mu}$.

On a une réduction semblable lorsque l'on veut prouver que la famille
$(u'_\lambda)$ est \emph{séparante} ($X$ et $X'$ étant affines)~: cela résulte
encore de (11.9.2) et (11.9.3).
\end{env}

\begin{proposition}[11.9.7]
\label{IV.11.9.7-fr}
Avec les notations de (11.9.1) et (11.9.5), supposons
$X=\operatorname{Spec}(A)$ et $X'=\operatorname{Spec}(A')$ affines, et
supposons en outre que $A'$ soit un $A$-module projectif. Alors, si la famille
$(u_\lambda)$ est séparante, toute section $t'$ de $\mathcal F'$ au-dessus de
$X'$ dont les images par tous les $u'_\lambda$ sont nulles, est elle-même
nulle.
\end{proposition}

\begin{proof}
On a vu (11.9.6) qu'on peut se borner au cas où tous les $Z_\lambda$ sont égaux
à $X$. La proposition est alors conséquence du lemme suivant~:
\end{proof}

\begin{lemma}[11.9.7.1]
\label{IV.11.9.7.1-fr}
Soient $A$ un anneau, $(M_\lambda)_{\lambda\in L}$ une famille de $A$-modules,
$M$ un $A$-module et pour chaque $\lambda$,
$u_\lambda:M\to M_\lambda$ un homomorphisme. Supposons que l'intersection des
noyaux des $u_\lambda$ soit réduite à 0. Alors, pour tout $A$-module projectif
$N$, l'intersection des noyaux des homomorphismes
$u_\lambda\otimes 1:M\otimes_A N\to M_\lambda\otimes_A N$ est réduite à 0.
\end{lemma}

\begin{proof}
En effet, $N$ est facteur direct d'un $A$-module libre $L$, et il suffit
évidemment de prouver que l'intersection des noyaux des homomorphismes
$v_\lambda:M\otimes_A L\to M_\lambda\otimes_A L$ est réduite à 0, puisque
$u_\lambda\otimes 1:M\otimes_A N\to M_\lambda\otimes_A N$ est la restriction
de $v_\lambda$. Mais l'assertion résulte alors trivialement de l'hypothèse.
\end{proof}

\begin{remark}[11.9.8]
\label{IV.11.9.8-fr}
Nous ignorons si, sous les hypothèses de la proposition (11.9.7), la famille
$(u'_\lambda)$ est \emph{séparante}~: il faudrait en effet (11.9.4) prouver
qu'une section $t'$ de $\mathcal F'$ \emph{au-dessus d'un ouvert}
$D(h')\subset X'$ (où $h'\in A'$) telle que les $u'_\lambda(t')$ soient toutes
nulles, est elle-même nulle. Or, on ne peut appliquer la proposition (11.9.4)
à $D(h')=\operatorname{Spec}(A'_{h'})$, car du fait que $A'$ soit un
$A$-module projectif (même libre), il ne résulte pas que $A'_{h'}$ soit un
$A$-module projectif. Par exemple, on peut prendre pour $A$ un anneau de
valuation discrète, pour $A'$ un anneau de valuation discrète qui soit un
$A$-module libre de rang 2, et pour $A'_{h'}$ le corps des fractions de $A'$.
\end{remark}

On a toutefois le résultat suivant~:

\begin{corollary}[11.9.9]
\label{IV.11.9.9-fr}
Soient $X$ un préschéma artinien, $g:X'\to X$ un morphisme plat (on notera que
ces deux conditions sont satisfaites si $X$ est le spectre d'un \emph{corps}
et $g$ un morphisme \emph{quelconque}). Alors, avec les notations de (11.9.1)
et (11.9.5), si la famille $(u_\lambda)$ est séparante, il en est de même de la
famille $(u'_\lambda)$.
\end{corollary}

\begin{proof}
On peut évidemment se borner au cas où $X=\operatorname{Spec}(A)$ est le
spectre d'un anneau local artinien $A$ (\textbf I, 6.2.2)~; on note ensuite que
pour tout ouvert affine $U'=\operatorname{Spec}(A')$ de $X'$, $A'$ est un
$A$-module plat, donc \emph{projectif} $(\mathbf 0_{\mathrm{III}},10.1.3)$.
Il suffit donc d'appliquer (11.9.7) à tout ouvert affine de $X'$ pour obtenir
le corollaire.
\end{proof}

\begin{theorem}[11.9.10]
\label{IV.11.9.10-fr}
Soient $X$ un préschéma, $(u_\lambda)$ une famille d'homomorphismes
(11.9.1.1), $g:X'\to X$ un morphisme, $(u'_\lambda)$ la famille
d'homomorphismes déduite de $(u_\lambda)$ par le changement de base $g$
(11.9.5).
\begin{enumerate}[label=(\roman*)]
  \item Si $g$ est un morphisme fidèlement plat et si la famille
  $(u'_\lambda)$ est séparante, alors la famille $(u_\lambda)$ est séparante.
  \oldpage[IV-3]{164}
  \item Supposons que $g$ soit un morphisme plat, et en outre que l'une des
  deux conditions suivantes soit vérifiée~:
  \begin{enumerate}[label=\emph{\alph*)}]
    \item $L$ est fini et les $f_\lambda$ sont quasi-compacts.
    \item Le morphisme $g$ est localement de présentation finie.
  \end{enumerate}
  Alors, si la famille $(u_\lambda)$ est séparante, il en est de même de la
  famille $(u'_\lambda)$.
\end{enumerate}
\end{theorem}

\begin{proof}
\begin{enumerate}[label=(\roman*)]
  \item En vertu de (2.2.8), il suffit de montrer que si une section $t$ de
  $\mathcal F$ au-dessus d'un ouvert $U$ de $X$ appartient au noyau de chacun
  des $u_\lambda$, son image $t'$ par l'homomorphisme canonique
  \[
  \Gamma(\rho):\Gamma(U,\mathcal F)\longrightarrow
  \Gamma(g^{-1}(U),g^*(\mathcal F))
  =\Gamma(U,g_*(g^*(\mathcal F)))
  \]
  est nulle. Or les images de $t'$ par les $g^*(u_\lambda)$ sont les images
  des $u_\lambda(t)$ par l'homomorphisme
  \[
  \Gamma(U,(f_\lambda)_*(\mathcal G_\lambda))\longrightarrow
  \Gamma(g^{-1}(U),g^*((f_\lambda)_*(\mathcal G_\lambda))),
  \]
  donc sont nulles, et \emph{a fortiori} on a $u'_\lambda(t')=0$ pour tout
  $\lambda$, donc $t'=0$ par hypothèse, ce qui prouve (i).

  \item La question étant locale sur $X$ et $X'$, on peut se borner au cas où
  $X=\operatorname{Spec}(A)$ et $X'=\operatorname{Spec}(A')$ sont affines,
  $A'$ étant un $A$-module plat, et à prouver que, pour toute section $z'$ de
  $\mathcal F'$ au-dessus de $X'$ dont les images par tous les $u'_\lambda$
  sont nulles, alors $z'$ est elle-même nulle. On a vu en outre (11.9.6) que
  l'on peut alors supposer que $Z_\lambda=X$ pour tout $\lambda$.

  Distinguons maintenant les deux cas.

  \emph{a)} Si $L$ est fini et les $f_\lambda$ quasi-compacts, on a vu (11.9.3)
  qu'on peut encore se ramener au cas où $Z_\lambda=X$ pour tout
  $\lambda\in L$, et où en outre $L$ est fini. Il revient au même alors de dire
  que l'intersection des noyaux des
  $u_\lambda:\mathcal F\to\mathcal G_\lambda$ est nulle, ou que
  l'homomorphisme
  $u=(u_\lambda):\mathcal F\to\mathcal G=\bigoplus_\lambda\mathcal G_\lambda$
  est injectif. Comme $u'=g^*(u)$ est injectif puisque $g$ est plat, la
  proposition est démontrée dans ce cas.

  \emph{b)} Soient $\mathcal F=\widetilde M$,
  $\mathcal G_\lambda=\widetilde N_\lambda$, et posons
  $M'=M\otimes_A A'$, $N'_\lambda=N_\lambda\otimes_A A'$, de sorte que
  $\mathcal F'=\widetilde{M'}$, $\mathcal G'_\lambda=\widetilde{N'_\lambda}$~;
  par abus de langage, nous noterons encore $u_\lambda$ l'homomorphisme
  $M\to N_\lambda$, et $u'_\lambda$ l'homomorphisme
  $u_\lambda\otimes 1:M'\to N'_\lambda$. Donnons-nous un élément $z'\in M'$
  tel que $u'_\lambda(z')=0$ pour tout $\lambda$~; il s'agit de prouver que
  l'on a $z'=0$. Or, l'hypothèse que $g$ est plat et de présentation finie
  entraîne, d'après (11.3.15), qu'il existe une suite finie
  $(s_i)_{1\leq i\leq n}$ d'éléments de $A$, telle que, si l'on pose
  $\mathfrak J_k=s_1A+\cdots+s_kA$, et
  $A_i=A_{s_i}/\mathfrak J_{i-1}A_{s_i}$, l'anneau
  $A'_i=A'\otimes_A A_i$ soit un $A_i$-module \emph{libre} pour
  $1\leq i\leq n$, et $\mathfrak J_n=A$. La proposition sera établie si nous
  prouvons pour $1\leq i\leq n$ l'assertion suivante~:
  \begin{enumerate}[label=$(*_i)$]
    \item \emph{Il existe un entier $m_i>0$ tel que $s_j^{m_i}z'=0$ pour
    $j\leq i$.}
  \end{enumerate}
  En effet, posant alors $k=m_n$ et notant que les $s_i^k$ ($1\leq i\leq n$)
  engendrent aussi l'idéal unité de $A$, l'assertion $(*_n)$ montrera que $z'$,
  combinaison linéaire des $s_i^kz'$, est nul.

  Prouvons $(*_i)$ par récurrence sur $i$, l'assertion étant vide pour $i=0$.
  Supposons donc $i>0$ et soit $m$ un multiple commun des $m_j$ pour
  $j\leq i-1$. Remarquons que (pour $1\leq h\leq n$) si
  $\mathfrak J'_h$ est l'idéal engendré par les $s_j^m$ ($0\leq j\leq h$),
  $\mathfrak J_h/\mathfrak J'_h$ est nilpotent~; remplacer les $s_j$ par
  $s_j^m$ pour $1\leq j\leq n$ revient donc à remplacer, pour
  $1\leq h\leq n$, $A_h$ par
  $A_{s_h}/\mathfrak J'_{h-1}A_{s_h}=B_h$, de sorte que
  $A_h=B_h/(\mathfrak J_{h-1}/\mathfrak J'_{h-1})B_h$~; comme $A'$ est un
  $A$-module plat, il résulte de $(\mathbf 0_{\mathrm{III}},10.1.2)$ que
  $B'_h=A'\otimes_A B_h$ est encore un $B_h$-module libre. On peut donc
  \oldpage[IV-3]{165}
  remplacer tous les $s_j$ ($1\leq j\leq n$) par $s_j^m$ sans changer les
  propriétés des $\mathfrak J_h$ et des $A'_h$, et supposer par la suite que
  $m=1$. Alors $z'$, étant annulé par $\mathfrak J_{i-1}$, s'identifie à un
  élément de
  $\operatorname{Hom}_{A'}(A'/\mathfrak J_{i-1}A',M')$, et comme
  $\mathfrak J_{i-1}A'$ est un idéal de type fini de $A'$, ce module
  d'homomorphismes s'identifie lui-même à
  $\operatorname{Hom}_A(A/\mathfrak J_{i-1},M)\otimes_A A'$
  $(\mathbf 0_{\mathrm I},6.2.2)$. Soit
  \[
  v_\lambda=\operatorname{Hom}(1,u_\lambda):
  \operatorname{Hom}_A(A/\mathfrak J_{i-1},M)\longrightarrow
  \operatorname{Hom}_A(A/\mathfrak J_{i-1},N_\lambda)
  \]
  l'homomorphisme déduit de $u_\lambda$~; la famille $(v_\lambda)$ est elle
  aussi \emph{séparante}. En effet, pour tout $t\in A$, on a
  \[
  (\operatorname{Hom}_A(A/\mathfrak J_{i-1},M))_t
  =\operatorname{Hom}_{A_t}(A_t/(\mathfrak J_{i-1})_t,M_t)
  \]
  et de même en remplaçant $M$ par $N_\lambda$, puisque l'idéal
  $\mathfrak J_{i-1}$ est de type fini $(\mathbf 0_{\mathrm I},1.3.5)$~;
  comme par hypothèse l'intersection des noyaux des $(u_\lambda)_t$ est nulle,
  il en est de même de l'intersection des noyaux des
  $(v_\lambda)_t=\operatorname{Hom}(1,(u_\lambda)_t)$, d'où notre assertion
  (11.9.4). Remplaçant $A$ par $A/\mathfrak J_{i-1}$, $M$ par
  $\operatorname{Hom}_A(A/\mathfrak J_{i-1},M)$, $N_\lambda$ par
  $\operatorname{Hom}_A(A/\mathfrak J_{i-1},N_\lambda)$ (qui sont des
  $(A/\mathfrak J_{i-1})$-modules), $u_\lambda$ par $v_\lambda$ et enfin $A'$
  par $A'/\mathfrak J_{i-1}A'$, on voit qu'on peut se ramener au cas où, dans
  la situation initiale, l'élément $s=s_i\in A$ est tel que $A'_s$ soit un
  $A_s$-module libre. Or la famille des
  $(u_\lambda)_s:M_s\to(N_\lambda)_s$ est séparante par hypothèse~; comme on a
  $(u'_\lambda)_s(z'/1)=0$, il résulte de (11.9.7) que l'on a $z'/1=0$ dans
  $M'_s$~; mais cela signifie qu'il existe un entier $r$ tel que $s^rz'=0$
  dans $M'$, ce qui achève de prouver $(*_i)$ par récurrence.
\end{enumerate}
\end{proof}

\begin{remark}[11.9.11]
\label{IV.11.9.11-fr}
Bornons-nous pour simplifier au cas où $Z_\lambda=X$ pour tout $\lambda$. Il
faut noter alors que si la famille des homomorphismes
$u_\lambda:\mathcal F\to\mathcal G_\lambda$ est séparante, \emph{il ne
s'ensuit pas nécessairement} que, pour tout $x\in X$, l'intersection des
noyaux des homomorphismes
$(u_\lambda)_x:\mathcal F_x\to(\mathcal G_\lambda)_x$ soit réduite à 0. Par
exemple, soient $X$ un préschéma de Jacobson localement noethérien, de
dimension $\geq1$, et $\mathcal F$ un $\mathcal O_X$-Module cohérent~; pour
tout point fermé $x$ de $X$, et tout entier $n\geq0$,
$\mathfrak m_x^n\mathcal F$ est un $\mathcal O_X$-Module cohérent de support
contenu dans $\{x\}$. La famille des homomorphismes canoniques
$\mathcal F\to\mathcal F/\mathfrak m_x^n\mathcal F$ (où $x$ parcourt
l'ensemble $X_0$ des points fermés de $X$ et $n$ l'ensemble des entiers
$\geq0$) est \emph{séparante}~: en effet, si $t$ est une section de
$\mathcal F$ au-dessus d'un ouvert $U$ de $X$ dont les images dans les
$\Gamma(U,\mathcal F/\mathfrak m_x^n\mathcal F)$ sont toutes nulles, il en
résulte aussitôt que pour tout point fermé $x\in U$, on a $t_x=0$~; comme
l'ensemble des points fermés contenus dans $U$ est très dense dans $U$, cela
entraîne bien $t=0$ (10.2.1). Cependant, si l'on prend
$\mathcal F=\mathcal O_X$, et si $y\in X$ est un point \emph{non fermé} de
$X$, on a $(\mathcal O_X/\mathfrak m_x^n\mathcal O_X)_y=0$ pour tout point
fermé $x$ de $X$, mais $\mathcal O_{X,y}\neq0$, et l'intersection des noyaux
des homomorphismes
$\mathcal O_{X,y}\to(\mathcal O_X/\mathfrak m_x^n\mathcal O_X)_y$ est égale
à $\mathcal O_{X,y}$.
\end{remark}

\begin{lemma}[11.9.12]
\label{IV.11.9.12-fr}
Les notations étant celles de (11.9.1) et (11.9.5), supposons la famille
$(u_\lambda)$ séparante~; supposons en outre que $X$ soit un $S$-préschéma, où
$S=\operatorname{Spec}(A)$ est un schéma affine, et que
$X'=X\times_S S'$, où $S'=\operatorname{Spec}(A')$, $A'$ étant une
$A$-algèbre~; supposons enfin que les $\mathcal G_\lambda$ soient $S$-plats.
Soit $(A'_\alpha)_{\alpha\in I}$ la famille filtrante des sous-$A$-algèbres de
type fini de $A'$~; pour tout $\alpha\in I$, posons
$S''_\alpha=\operatorname{Spec}(A'_\alpha)$,
$X''_\alpha=X\times_S S''_\alpha$,
$Z''_{\alpha\lambda}=Z_\lambda\times_X X''_\alpha$, et soient
$f''_{\alpha\lambda}:Z''_{\alpha\lambda}\to X''_\alpha$,
$\mathcal F''_\alpha$, $\mathcal G''_{\alpha\lambda}$ et
$u''_{\alpha\lambda}$ les morphismes, Modules et homomorphismes de Modules
déduits de $f_\lambda$, $\mathcal F$, $\mathcal G_\lambda$ et $u_\lambda$ au
moyen du changement de base $X''_\alpha\to X$. Alors, si, pour tout
$\alpha\in I$, la famille $(u''_{\alpha\lambda})_{\lambda\in L}$ est
séparante, il en est de même de $(u'_\lambda)_{\lambda\in L}$.
\end{lemma}

\begin{proof}
Il s'agit de prouver que si une section $t'$ de $\mathcal F'$ au-dessus d'un
ouvert affine $U'$ de $X'$ est telle que $u'_\lambda(t')=0$ pour tout
$\lambda\in L$, on a $t'=0$. Si $h_\alpha:X'\to X''_\alpha$ est la projection
\oldpage[IV-3]{166}
canonique, il résulte de (8.2.11) qu'il existe un indice $\alpha$ et un ouvert
quasi-compact $U''_\alpha\subset X''_\alpha$ tels que
$U'=h_\alpha^{-1}(U''_\alpha)$~; en outre, en vertu de (8.5.2, (i)), on peut
supposer qu'il y a une section $t''_\alpha$ de $\mathcal F''_\alpha$ au-dessus
de $U''_\alpha$ telle que $t'$ soit l'image canonique de $t''_\alpha$. Quitte à
remplacer $S$, $X$, $Z_\lambda$, $f_\lambda$, $\mathcal F$,
$\mathcal G_\lambda$ et $u_\lambda$ par $S''_\alpha$, $U''_\alpha$,
$(f''_{\alpha\lambda})^{-1}(U''_\alpha)$ et les restrictions correspondantes
de $\mathcal F''_\alpha$, $\mathcal G''_{\alpha\lambda}$ et
$u''_{\alpha\lambda}$, on peut donc supposer que $U'=X'$, que $t'$ est
l'image canonique d'une section $t$ de $\mathcal F$ au-dessus de $X$ et que
l'homomorphisme $A\to A'$ est injectif, ou encore, si $p:S'\to S$ est le
morphisme correspondant, que l'homomorphisme
$\mathcal O_S\to p_*(\mathcal O_{S'})$ intervenant dans la définition de $p$
est injectif. Il en résulte aussitôt en vertu de la platitude de
$\mathcal G_\lambda$ sur $S$ (et en se ramenant par exemple au cas où
$Z_\lambda$ est affine sur $S$ (\textbf I, 1.6.3 et 1.6.5)) que
l'homomorphisme canonique
$\rho:\mathcal G_\lambda\to(g'_\lambda)_*(\mathcal G'_\lambda)$ est lui aussi
injectif. Mais l'homomorphisme composé
\[
\Gamma(X,\mathcal F)\xrightarrow{\Gamma(u_\lambda)}
\Gamma(Z_\lambda,\mathcal G_\lambda)\xrightarrow{\Gamma(\rho)}
\Gamma(Z'_\lambda,\mathcal G'_\lambda)
\]
est égal par définition à l'homomorphisme composé
\[
\Gamma(X,\mathcal F)\xrightarrow{\Gamma(\rho)}
\Gamma(X',\mathcal F')\xrightarrow{\Gamma(u'_\lambda)}
\Gamma(Z'_\lambda,\mathcal G'_\lambda)
\]
donc l'image de $t$ par ces homomorphismes composés est
$u'_\lambda(t')=0$~; en vertu de l'injectivité de l'homomorphisme
$\Gamma(Z_\lambda,\mathcal G_\lambda)\xrightarrow{\Gamma(\rho)}
\Gamma(Z'_\lambda,\mathcal G'_\lambda)$ on en conclut que $u_\lambda(t)=0$
pour tout $\lambda\in L$, d'où $t=0$ par hypothèse, et finalement $t'=0$.
\end{proof}

\begin{proposition}[11.9.13]
\label{IV.11.9.13-fr}
Les notations étant celles de (11.9.1) et (11.9.5), supposons que $X$ soit un
préschéma sur un corps $k$, et que, en posant $S=\operatorname{Spec}(k)$, on
ait $X'=X\times_S S'$, où $S'$ est un $k$-préschéma quelconque. Alors, si la
famille $(u_\lambda)_{\lambda\in L}$ est séparante, il en est de même de
$(u'_\lambda)$.
\end{proposition}

\begin{proof}
On peut se borner au cas où $S'=\operatorname{Spec}(A')$ est affine. Si $A'$
est une $k$-algèbre de type fini, le morphisme $g:X'\to X$ est plat et de
présentation finie, et on est donc dans les conditions d'application de
(11.9.10, (ii), b)). Dans le cas général, on considère $A'$ comme limite
inductive de ses $k$-sous-algèbres $A'_\alpha$ de type fini, et on applique à
chaque $A'_\alpha$ le résultat de (11.9.10, (ii), b))~; on conclut alors à
l'aide du lemme (11.9.12), puisque les $\mathcal G_\lambda$ sont $S$-plats.
\end{proof}

\begin{env}[11.9.14]
\label{IV.11.9.14-fr}
Gardons toujours les notations de (11.9.1) et (11.9.5) et supposons que $X$
soit un $S$-préschéma. Si pour tout changement de base
$g:X\times_S S'\to X$, où $S'\to S$ est un morphisme quelconque, la famille
$(u'_\lambda)$ correspondante est séparante, nous dirons que la famille
$(u_\lambda)$ est \emph{universellement séparante relativement à} $S$.
Lorsque la famille $(u_\lambda)$ est réduite à un seul élément $u$, nous
dirons encore que $u$ est \emph{universellement injectif}, relativement à
$S$. Il est clair alors que pour tout morphisme $h:S'\to S$, la famille
$(u'_\lambda)$ correspondante est universellement séparante relativement à
$S'$~; inversement, si $h$ est \emph{fidèlement plat} et si $(u'_\lambda)$ est
universellement séparante relativement à $S'$, alors $(u_\lambda)$ est
universellement séparante relativement à $S$, comme il résulte aussitôt de
(11.9.10, (i)) et du fait que pour tout morphisme $S''\to S$, le morphisme
correspondant $S''\times_S S'\to S''$ est fidèlement plat.
\end{env}

\oldpage[IV-3]{167}

\begin{proposition}[11.9.15]
\label{IV.11.9.15-fr}
Les notations étant celles de (11.9.1), supposons que $X$ soit un
$S$-préschéma, les $\mathcal G_\lambda$ étant $S$-plats. Soit $S_0$ un
sous-préschéma fermé de $S$ défini par un Idéal quasi-cohérent nilpotent
$\mathcal J$ de $\mathcal O_S$, tel que les
$(\mathcal O_S/\mathcal J)$-Modules $\mathcal J^k/\mathcal J^{k+1}$ soient
tous localement libres. Soit $(u_{\lambda0})$ la famille d'homomorphismes
obtenue à partir du changement de base
$X\leftarrow X_0=X\times_S S_0$. Alors, si la famille $(u_{\lambda0})$ est
séparante (resp. universellement séparante relativement à $S_0$), la famille
$(u_\lambda)$ est séparante (resp. universellement séparante relativement à
$S$).
\end{proposition}

\begin{proof}
Notons que si $S'\to S$ est un changement de base quelconque et
$S'_0=S'\times_S S_0$, $S'_0$ est un sous-préschéma fermé de $S'$ défini par
un idéal quasi-cohérent nilpotent $\mathcal J'$ de $\mathcal O_{S'}$, tel que
$\mathcal J'^k/\mathcal J'^{k+1}$ soit un
$(\mathcal O_{S'}/\mathcal J')$-Module localement libre pour tout $k$
(2.1.8, (i))~; comme en outre les
$\mathcal G'_\lambda=\mathcal G_\lambda\otimes_S S'$ sont $S'$-plats, on voit
que l'assertion relative aux familles universellement séparantes est
conséquence de l'assertion relative aux familles séparantes. Pour prouver
cette dernière, on peut (11.9.3) se ramener au cas où
$S=\operatorname{Spec}(A)$, $X=\operatorname{Spec}(B)$ sont affines,
$Z_\lambda=X$ pour tout $\lambda$, $\mathcal F=\widetilde M$,
$\mathcal G_\lambda=\widetilde N_\lambda$, où $M$ et les $N_\lambda$ sont des
$B$-modules, les $N_\lambda$ étant des $A$-modules plats. En outre, la question
étant locale sur $X$ et $S$, il suffit de voir que si $t\in M$ est tel que
$u_\lambda(t)=0$ pour tout $\lambda$, alors $t=0$. On a
$\mathcal J=\widetilde{\mathfrak J}$, où $\mathfrak J$ est un idéal nilpotent
de $A$, tel que les $\mathfrak J^k/\mathfrak J^{k+1}$ soient des
$(A/\mathfrak J)$-modules libres, et par hypothèse les
$u_{\lambda0}:M/\mathfrak JM\to N_\lambda/\mathfrak JN_\lambda$ forment une
famille séparante. Supposons que $\mathfrak J^{n+1}=0$ ($n$ entier $\geq0$) et
raisonnons par récurrence sur $n$, l'assertion étant triviale pour $n=0$. Si
$\bar t$ est la classe de $t$ dans $M/\mathfrak J^nM$, la classe de
$u_\lambda(t)$ dans $N_\lambda/\mathfrak J^nN_\lambda$ est nulle pour tout
$\lambda$, donc, par l'hypothèse de récurrence, $\bar t=0$, autrement dit on a
$t\in\mathfrak J^nM$. Or
$\mathfrak J^n=\mathfrak J^n/\mathfrak J^{n+1}$ est un
$(A/\mathfrak J)$-module libre~; si $(e_\alpha)$ est une base de ce module, on
peut donc écrire $t=\sum_\alpha e_\alpha t^0_\alpha$, avec
$t^0_\alpha\in M/\mathfrak JM$, nul sauf pour un nombre fini d'indices.
D'autre part, puisque $N_\lambda$ est un $A$-module plat,
$\mathfrak J^nN_\lambda$ s'identifie à
$\mathfrak J^n\otimes_{(A/\mathfrak J)}(N_\lambda/\mathfrak JN_\lambda)$ et
l'on peut par suite écrire
\[
u_\lambda(t)=\sum_\alpha e_\alpha u_{\lambda0}(t^0_\alpha)
=\sum_\alpha e_\alpha\otimes u_{\lambda0}(t^0_\alpha).
\]
Comme par hypothèse $u_\lambda(t)=0$, on en déduit que
$u_{\lambda0}(t^0_\alpha)=0$ pour tout $\alpha$ et tout $\lambda$~; d'où
$t^0_\alpha=0$ pour tout $\alpha$ puisque la famille $(u_{\lambda0})$ est
séparante, et par suite $t=0$. C.Q.F.D.
\end{proof}

\begin{theorem}[11.9.16]
\label{IV.11.9.16-fr}
Les notations étant celles de (11.9.1), supposons que $X$ soit un
$S$-préschéma localement noethérien, $\mathcal F$ un $\mathcal O_X$-Module
cohérent et que les $\mathcal G_\lambda$ soient $S$-plats. Pour tout $s\in S$,
soit $((u_\lambda)_s)_{\lambda\in L}$ la famille obtenue à partir de
$(u_\lambda)$ par le changement de base
$X_s=X\times_S\operatorname{Spec}(k(s))\to X$. Alors, pour que la famille
$(u_\lambda)$ soit universellement séparante relativement à $S$, il faut et il
suffit que pour tout $s\in S$, la famille $((u_\lambda)_s)$ soit séparante.
\end{theorem}

\begin{proof}
La nécessité de la condition découle trivialement des définitions.
Inversement, supposons la condition de l'énoncé vérifiée, et prouvons d'abord
que la famille $(u_\lambda)$ est séparante. On peut (11.9.3) se réduire au cas
où $S=\operatorname{Spec}(A)$, $X=\operatorname{Spec}(B)$ sont affines,
$Z_\lambda=X$ pour tout $\lambda$, $\mathcal F=\widetilde M$,
$\mathcal G_\lambda=\widetilde N_\lambda$, où $B$ est noethérien, $M$ est un
$B$-module de type fini et les $N_\lambda$ des $A$-modules plats et se borner à
prouver que, si $t\in M$ est tel que $u_\lambda(t)=0$ pour tout $\lambda$,
alors $t=0$. Pour montrer que $t=0$, il suffit de prouver que pour tout idéal
maximal $\mathfrak p$ de $B$, l'image $t_\mathfrak p$ de $t$ dans
$M_\mathfrak p$ est nulle (Bourbaki, \emph{Alg. comm.}, chap. II, \S~3,
n\textsuperscript{o}~3, cor. 1 du th. 1). On peut donc se borner à montrer que
l'intersection
\oldpage[IV-3]{168}
des noyaux des
$(u_\lambda)_\mathfrak p:M_\mathfrak p\to(N_\lambda)_\mathfrak p$ déduits
des $(u_\lambda)$ par le changement de base
$\operatorname{Spec}(B_\mathfrak p)\to\operatorname{Spec}(B)$ est réduite à
0. Autrement dit, on est ramené au cas où $B$ est un anneau local noethérien,
et en considérant l'idéal premier de $A$ image réciproque de l'idéal maximal de
$B$, on peut aussi supposer que $A$ est un anneau local, d'idéal maximal
$\mathfrak m$. Alors, comme $\mathfrak mB$ est contenu dans l'idéal maximal de
$B$, et que $M$ est un $B$-module de type fini, l'intersection des
$\mathfrak m^nM$ est réduite à 0 $(\mathbf 0_{\mathrm I},7.3.5)$, donc il
suffit de prouver que pour tout $n$, l'image de $t$ dans
$M/\mathfrak m^{n+1}M$ est nulle. Il suffit donc de prouver que la famille
déduite de $(u_\lambda)$ par le changement de base
$\operatorname{Spec}(B/\mathfrak m^{n+1}B)\to\operatorname{Spec}(B)$ est
séparante, ce qui signifie encore qu'on peut se borner au cas où $A$ est un
anneau local dont l'idéal maximal $\mathfrak m$ est nilpotent. Mais alors les
$\mathfrak m^k/\mathfrak m^{k+1}$ sont des $(A/\mathfrak m)$-modules libres,
et en vertu de l'hypothèse sur les $(u_\lambda)_s$, on est précisément dans les
conditions d'application de (11.9.15), d'où la conclusion annoncée.

Soient maintenant $h:S'\to S$ un morphisme changement de base, et
$(u'_\lambda)$ la famille obtenue à partir de $(u_\lambda)$ par le changement
de base $h':X\times_S S'\to X$~; prouvons que $(u'_\lambda)$ est aussi
séparante. Supposons d'abord que $h$ soit \emph{localement de type fini}~; il
en est alors de même de $h'$, donc $X'=X\times_S S'$ est localement
noethérien~; de plus, si $s'\in S'$ est au-dessus du point $s\in S$, il résulte
de (11.9.13) appliqué à $X_s$ et à
$X_{s'}=X_s\otimes_{k(s)}k(s')$ que, pour tout $s'\in S'$, la famille
$((u'_\lambda)_{s'})$ est séparante~; on peut par suite conclure de la première
partie de la démonstration que dans ce cas $(u'_\lambda)$ est séparante.

Enfin, si $h$ est quelconque, on peut évidemment se limiter au cas où
$S=\operatorname{Spec}(A)$ et $S'=\operatorname{Spec}(A')$ sont affines, et
considérer $A'$ comme limite inductive de ses sous-$A$-algèbres de type fini.
Comme les $\mathcal G_\lambda$ sont $S$-plats, il suffit d'appliquer ce qui
précède et le lemme (11.9.5) pour terminer la démonstration.
\end{proof}

\begin{proposition}[11.9.17]
\label{IV.11.9.17-fr}
Soient $f:X\to S$ un morphisme localement de présentation finie, $\mathcal F$
un $\mathcal O_X$-Module quasi-cohérent de présentation finie et $f$-plat,
$U$ un ouvert de $X$, $j:U\to X$ l'injection canonique,
$u:\mathcal F\to j_*(j^*(\mathcal F))$ l'homomorphisme canonique
$(\mathbf 0_{\mathrm I},4.4.3.2)$. Pour tout $s\in S$, soient $X_s$ la fibre
$X\times_S\operatorname{Spec}(k(s))$, $U_s$ l'ouvert $U\cap X_s$ de $X_s$,
$\mathcal F_s=\mathcal F\otimes_{\mathcal O_S}k(s)$,
$j_s:U_s\to X_s$ l'injection canonique,
$u_s:\mathcal F_s\to(j_s)_*((j_s)^*(\mathcal F_s))$ l'homomorphisme
canonique. Pour que $u$ soit universellement injectif relativement à $S$, il
faut et il suffit que $u_s$ soit injectif pour tout $s\in S$.
\end{proposition}

\begin{proof}
Il n'y a encore à prouver que la suffisance de la condition. Lorsque $X$ est
localement noethérien, la proposition est un corollaire immédiat de (11.9.16).
Nous allons nous ramener à ce cas en deux étapes, en nous bornant, comme on
peut évidemment le faire, au cas où $S=\operatorname{Spec}(A)$ et $X$ sont
affines.

\emph{A) Cas où $U$ est quasi-compact.} --- Nous utiliserons le lemme suivant~:

\begin{lemma}[11.9.17.1]
\label{IV.11.9.17.1-fr}
Sous les hypothèses générales de (11.9.17), et en supposant en outre $S$ et
$X$ affines et $U$ quasi-compact, l'ensemble $E$ des $s\in S$ tels que $u_s$
soit injectif est constructible.
\end{lemma}

\begin{proof}
En effet, les fibres $X_s$ sont des préschémas localement noethériens, donc
$E$ peut aussi, en vertu de (5.10.2), être défini comme l'ensemble des
$s\in S$ tels que $\operatorname{Ass}(\mathcal F_s)\subset U_s$. Notons en
outre que $U$, étant quasi-compact dans un schéma affine, est constructible.
Alors, la vérification de la condition (9.2.1, (i)) découle aussitôt de
(4.2.7)~; d'autre
\oldpage[IV-3]{169}
part, la vérification de (9.2.1, (ii)) se fait aisément en utilisant l'étude des
cycles premiers associés au voisinage de la fibre générique (9.8.3), ainsi que
(9.5.2) et (9.5.3).
\end{proof}

Ce lemme étant établi, on peut, en vertu de (8.9.1) et (8.2.11), supposer qu'il
existe un sous-anneau noethérien $A_0$ de $A$, un préschéma de type fini $X_0$
sur $S_0=\operatorname{Spec}(A_0)$, un ouvert $U_0$ de $X_0$ et un
$\mathcal O_{X_0}$-Module cohérent $\mathcal F_0$ tels que
$X=X_0\times_{S_0}S$ et que, si $p_0:X\to X_0$ est la projection canonique,
on ait $U=p_0^{-1}(U_0)$ et $\mathcal F=p_0^*(\mathcal F_0)$. Soit
$(A_\alpha)$ la famille filtrante des sous-anneaux de $A$ qui sont des
$A_0$-algèbres de type fini, et posons
$S_\alpha=\operatorname{Spec}(A_\alpha)$,
$X_\alpha=X_0\times_{S_0}S_\alpha$,
$U_\alpha=U_0\times_{S_0}S_\alpha$,
$\mathcal F_\alpha=\mathcal F_0\otimes_{S_0}S_\alpha$~; soit $u_\alpha$
l'homomorphisme canonique relatif à $\mathcal F_\alpha$ et $U_\alpha$, défini
comme dans (11.9.17). Pour tout $s\in S$, l'hypothèse que $u_s$ soit injectif
entraîne que $(u_\alpha)_{s_\alpha}$ l'est aussi (11.9.10, (i)), où
$s_\alpha=q_\alpha(s)$ est l'image de $s$ par le morphisme
$q_\alpha:S\to S_\alpha$. Si $E_\alpha$ est l'ensemble des
$s_\alpha\in S_\alpha$ tels que $(u_\alpha)_{s_\alpha}$ soit injectif, on a
donc $S=q_\alpha^{-1}(E_\alpha)$, et les $E_\alpha$ forment un système
projectif d'ensembles. Mais le lemme (11.9.17.1) appliqué à $u_\alpha$ montre
que $E_\alpha$ est constructible dans $S_\alpha$~; on déduit donc de (8.3.4)
qu'il existe un indice $\alpha$ tel que $E_\alpha=S_\alpha$. Mais alors, comme
$X_\alpha$ est noethérien, $u_\alpha$ est universellement injectif en vertu de
(11.9.16), donc il en est de même de $u$.

\emph{B) Cas général.} --- L'ouvert $U$ est réunion filtrante croissante
d'ouverts quasi-compacts $U_\lambda$~; si $j_\lambda:U_\lambda\to X$ est
l'injection canonique et
$u_\lambda:\mathcal F\to(j_\lambda)_*(j_\lambda^*(\mathcal F))$
l'homomorphisme correspondant, il résulte du lemme (11.9.17.1) que l'ensemble
$E_\lambda$ des $s\in S$ tels que $(u_\lambda)_s$ soit injectif est
constructible. D'autre part, pour un $s\in S$, dire que $u_s$ (resp.
$(u_\lambda)_s$) est injectif signifie que
$\operatorname{Ass}(\mathcal F_s)\subset U\cap X_s$ (resp.
$\operatorname{Ass}(\mathcal F_s)\subset U_\lambda\cap X_s$) (5.10.2).
Comme $\operatorname{Ass}(\mathcal F_s)$ est fini (3.1.6), dire que $u_s$ est
injectif signifie donc qu'il existe $\lambda$ tel que $s\in E_\lambda$~;
l'hypothèse signifie par suite que $S=\bigcup_\lambda E_\lambda$. En vertu de
(1.9.9), il existe un indice $\lambda$ tel que $S=E_\lambda$, d'où l'on
conclut par la première partie du raisonnement que $u_\lambda$ est
universellement injectif. Il résulte alors de la factorisation de $u_\lambda$~:
\[
\mathcal F\xrightarrow{u}j_*(\mathcal F|U)\longrightarrow
(j_\lambda)_*(\mathcal F|U_\lambda)
\]
que $u$ est aussi universellement injectif.
\end{proof}

\begin{remark}[11.9.18]
\label{IV.11.9.18-fr}
Soient $X$ un préschéma, $\mathcal F$, $\mathcal G$ deux
$\mathcal O_X$-Modules quasi-cohérents de présentation finie, $\mathcal G$
étant en outre supposé \emph{localement libre},
$u:\mathcal F\to\mathcal G$ un homomorphisme. Les conditions suivantes sont
équivalentes~:
\begin{enumerate}[label=\emph{\alph*)}]
  \item pour tout morphisme $g:X'\to X$, l'homomorphisme
  $g^*(u):g^*(\mathcal F)\to g^*(\mathcal G)$ est injectif~;
  \item pour tout $x\in X$, l'homomorphisme
  $u\otimes 1_{k(x)}:\mathcal F\otimes_{\mathcal O_X}k(x)
  \to\mathcal G\otimes_{\mathcal O_X}k(x)$ est injectif~;
  \item pour tout $x\in X$, il existe un voisinage ouvert $U$ de $x$ tel que
  $u|U:\mathcal F|U\to\mathcal G|U$ soit un isomorphisme de $\mathcal F|U$
  sur un facteur direct de $\mathcal G|U$.
\end{enumerate}

En effet, il est clair que c) entraîne a) et que a) entraîne b). Le fait que
b) entraîne c) résulte de $(\mathbf 0,19.1.12)$,
$(\mathbf 0_{\mathrm I},5.2.5)$ et $(\mathbf 0_{\mathrm I},5.5.5)$.

Lorsque les conditions équivalentes précédentes sont vérifiées, on dit que
$u$ est \emph{universellement injectif}.
\end{remark}
\oldpage[IV-3]{170}
\subsection{Familles schématiquement dominantes de morphismes et familles
schématiquement denses de sous-préschémas}
\label{subsection:IV.11.10-fr}

\begin{proposition}[11.10.1]
\label{IV.11.10.1-fr}
Soient $X$ un préschéma, $(Z_\lambda)_{\lambda\in L}$ une famille de
préschémas, et pour tout $\lambda\in L$, soit
$f_\lambda=(\psi_\lambda,\theta_\lambda):Z_\lambda\to X$ un morphisme. Les
conditions suivantes sont équivalentes~:
\begin{enumerate}[label=\emph{\alph*)}]
  \item La famille des homomorphismes
  $\theta_\lambda:\mathcal O_X\to(f_\lambda)_*(\mathcal O_{Z_\lambda})$ est
  séparante (autrement dit (11.9.1), l'intersection des noyaux des
  $\theta_\lambda$ est nulle).
  \item Pour tout ouvert $U$ de $X$, toute section $t$ de $\mathcal O_X$
  au-dessus de $U$ dont toutes les images par les homomorphismes canoniques
  \[
  (\theta_\lambda)_U:\Gamma(U,\mathcal O_X)
  \longrightarrow\Gamma(f_\lambda^{-1}(U),\mathcal O_{Z_\lambda})
  \tag{11.10.1.1}
  \label{IV.11.10.1.1-fr}
  \]
  sont nulles, est elle-même nulle.
  \item Pour tout ouvert $U$ de $X$, et tout sous-préschéma fermé $Y$ de $U$
  tel que pour tout $\lambda\in L$, il existe une factorisation
  \[
  f_\lambda^{-1}(U)\xrightarrow{g_\lambda}Y\xrightarrow{j}U
  \tag{11.10.1.2}
  \label{IV.11.10.1.2-fr}
  \]
  de la restriction $f_\lambda^{-1}(U)\to U$ de $f_\lambda$ (où $j$ est
  l'injection canonique), on a $Y=U$.
\end{enumerate}

Si de plus $X$ est un $S$-préschéma, ces conditions sont aussi équivalentes à
la suivante~:
\begin{enumerate}[label=\emph{\alph*)}]
\setcounter{enumi}{3}
  \item Pour tout morphisme séparé $g:X'\to S$ et tout couple de
  $S$-morphismes $u$, $v$ d'un ouvert $U$ de $X$ dans $X'$, tel que pour tout
  $\lambda$, les composés de $u$ et $v$ avec le morphisme
  $f_\lambda^{-1}(U)\to U$, restriction de $f_\lambda$, soient égaux, on a
  $u=v$.
\end{enumerate}
\end{proposition}

\begin{proof}
L'équivalence de a) et b) résulte des définitions. Pour voir que b) entraîne
c), il suffit de considérer l'Idéal quasi-cohérent $\mathcal J$ de
$\mathcal O_U$ définissant $Y$, et de noter que, pour tout ouvert $V\subset U$,
l'hypothèse entraîne que le morphisme $(\theta_\lambda)_V$ se factorise en
\[
\Gamma(V,\mathcal O_U)\longrightarrow
\Gamma(V,\mathcal O_U/\mathcal J)\longrightarrow
\Gamma(f_\lambda^{-1}(V),\mathcal O_{Z_\lambda}).
\]
On en conclut que toute section $t$ de $\mathcal J$ au-dessus de $V$ a pour
image 0 dans tous les $\Gamma(f_\lambda^{-1}(V),\mathcal O_{Z_\lambda})$,
donc, en vertu de b), $\mathcal J=0$ et $Y=U$. Inversement, si c) est vérifiée,
il suffit, pour prouver b), d'appliquer c) au sous-préschéma fermé $Y$ de $U$
défini par l'Idéal $t\mathcal O_U$~: l'hypothèse que les images de $t$ par les
$(\theta_\lambda)_U$ sont toutes nulles entraîne que l'on a des factorisations
(11.10.1.2) pour tout $\lambda$ ($\mathbf I$, 4.1.9). Pour prouver que c)
entraîne d), il suffit d'appliquer c) au sous-préschéma fermé $Y$ de $U$, image
réciproque de la diagonale de $X'\times_SX'$ par $(u,v)_S$ et d'utiliser
($\mathbf I$, 4.4.1). Inversement, on déduit b) de d) en considérant le
$S$-schéma $X'=S[T]=S\otimes_{\mathbf Z}\mathbf Z[T]$ ($T$ indéterminée) et en
se rappelant que les sections de $\mathcal O_U$ au-dessus de $U$ correspondent
biunivoquement aux $S$-morphismes $U\to S[T]$ ($\mathbf I$, 3.3.15)~; dire que
deux sections de $\mathcal O_U$ au-dessus de $U$ ont mêmes images par tous les
$(\theta_\lambda)_U$ équivaut à dire que les composés des deux morphismes
correspondants avec tous les morphismes $f_\lambda^{-1}(U)\to U$ sont égaux.
\end{proof}

\oldpage[IV-3]{171}
\begin{definition}[11.10.2]
\label{IV.11.10.2-fr}
Lorsque les conditions équivalentes de (11.10.1) sont vérifiées, on dit que la
famille $(f_\lambda)$ est \emph{schématiquement dominante}. Lorsque les
$f_\lambda$ sont les immersions canoniques dans $X$ d'une famille
$(Z_\lambda)$ de sous-préschémas de $X$, on dit aussi que la famille
$(Z_\lambda)$ est \emph{schématiquement dense}.
\end{definition}

\begin{remarks}[11.10.3]
\label{IV.11.10.3-fr}
\begin{enumerate}[label=(\roman*)]
  \item La notion de famille schématiquement dominante est \emph{locale} sur
  $X$, comme il résulte par exemple de la forme b) dans (11.10.1)~: si
  $(W_\alpha)$ est un recouvrement ouvert de $X$, la famille $(f_\lambda)$ est
  schématiquement dominante si et seulement s'il en est ainsi de chacune des
  familles formées des morphismes
  $f_\lambda^{-1}(W_\alpha)\to W_\alpha$, restrictions des $f_\lambda$.
  \item Si $Z$ est le préschéma somme des $Z_\lambda$, $f:Z\to X$ le morphisme
  coïncidant avec $f_\lambda$ dans chacun des $Z_\lambda$, il revient au même
  de dire que la famille $(f_\lambda)$ est schématiquement dominante, ou que
  la famille réduite au seul élément $f$ l'est.
  \item Soit $M$ un second ensemble d'indices et, pour tout $\lambda\in L$,
  soit $(g_{\lambda\mu})_{\mu\in M}$ une famille de morphismes
  $g_{\lambda\mu}:Z'_{\lambda\mu}\to Z_\lambda$~; si, pour chaque
  $\lambda\in L$, la famille $(g_{\lambda\mu})_{\mu\in M}$ est
  schématiquement dominante, alors, il revient au même de dire que la famille
  $(f_\lambda)_{\lambda\in L}$ est schématiquement dominante, ou que la
  famille
  $(f_\lambda\circ g_{\lambda\mu})_{(\lambda,\mu)\in L\times M}$ l'est
  (11.9.2).
  \item Soit $f:Z\to X$ un morphisme tel que $f_*(\mathcal O_Z)$ soit un
  $\mathcal O_X$-Module quasi-cohérent (par exemple un morphisme quasi-compact
  et quasi-séparé (1.7.4)). Alors, dire que $f$ est schématiquement dominant
  signifie que \emph{l'image fermée de $Z$ par $f$} ($\mathbf I$, 9.5.3) est
  \emph{identique à $X$}.
\end{enumerate}
\end{remarks}

\begin{proposition}[11.10.4]
\label{IV.11.10.4-fr}
Si la famille de morphismes $f_\lambda:Z_\lambda\to X$ est schématiquement
dominante, la réunion des $f_\lambda(Z_\lambda)$ est dense dans $X$.
Inversement, si cette réunion est dense dans $X$ et si de plus $X$ est réduit,
la famille $(f_\lambda)$ est schématiquement dominante.
\end{proposition}

\begin{proof}
La première assertion résulte aussitôt de (11.10.1, b)). D'autre part, si $X$
est réduit, et si la réunion des $f_\lambda(Z_\lambda)$ est dense dans $X$,
alors, si l'on a des factorisations (11.10.1.2) pour tout $\lambda\in L$, on a
aussi des factorisations
\[
(f_\lambda^{-1}(U))_{\mathrm{red}}
\xrightarrow{(g_\lambda)_{\mathrm{red}}}Y_{\mathrm{red}}
\xrightarrow{j_{\mathrm{red}}}U,
\]
et l'hypothèse entraîne que l'espace sous-jacent à $Y$ est identique à $U$,
donc $Y=Y_{\mathrm{red}}=U$ puisque $U$ est réduit.
\end{proof}

Les résultats du n$^\circ$~11.9 sur les familles séparantes se traduisent en
résultats sur les familles schématiquement dominantes~:

\begin{theorem}[11.10.5]
\label{IV.11.10.5-fr}
Soient $(f_\lambda)$ une famille de morphismes
$f_\lambda:Z_\lambda\to X$, $g:X'\to X$ un morphisme, et posons
$Z'_\lambda=Z_\lambda\times_XX'$,
$f'_\lambda=(f_\lambda)_{(X')}:Z'_\lambda\to X'$.
\begin{enumerate}[label=(\roman*)]
  \item Si $g$ est fidèlement plat et si la famille $(f'_\lambda)$ est
  schématiquement dominante, alors la famille $(f_\lambda)$ est
  schématiquement dominante.
  \item Supposons que $g$ soit un morphisme plat et en outre que l'une des
  conditions suivantes soit vérifiée~:
  \begin{enumerate}[label=\emph{\alph*)}]
    \item $L$ est fini et les $f_\lambda$ sont quasi-compacts.
    \item Le morphisme $g$ est localement de présentation finie.
  \end{enumerate}
  Alors, si la famille $(f_\lambda)$ est schématiquement dominante, il en est
  de même de la famille $(f'_\lambda)$.
\end{enumerate}
\end{theorem}

\oldpage[IV-3]{172}
\begin{proposition}[11.10.6]
\label{IV.11.10.6-fr}
Les notations étant celles de (11.10.5), supposons que $X$ soit un préschéma
sur un corps $k$, et qu'en posant $S=\operatorname{Spec}(k)$, on ait
$X'=X\times_SS'$, où $S'$ est un $k$-préschéma quelconque. Alors, si la famille
$(f_\lambda)$ est schématiquement dominante, il en est de même de
$(f'_\lambda)$.
\end{proposition}

\begin{corollary}[11.10.7]
\label{IV.11.10.7-fr}
Soient $X$ un préschéma sur un corps $k$, $(Z_\lambda)_{\lambda\in L}$ une
famille de $k$-préschémas géométriquement réduits sur $k$, et pour chaque
$\lambda\in L$, soit $f_\lambda:Z_\lambda\to X$ un $k$-morphisme. Désignons
par $Y$ le sous-préschéma réduit de $X$ ayant pour espace sous-jacent
l'adhérence de $\bigcup_{\lambda\in L}f_\lambda(Z_\lambda)$. Soient $k'$ une
extension de $k$, $X'$, $Z'_\lambda$,
$f'_\lambda:Z'_\lambda\to X'$ les préschémas et morphismes déduits de $X$,
$Z_\lambda$ et $f_\lambda$ par extension de la base à $k'$. Alors, si $Y'$ est
le sous-préschéma réduit de $X'$ ayant pour espace sous-jacent l'adhérence de
$\bigcup_{\lambda\in L}f'_\lambda(Z'_\lambda)$, on a
$Y'=Y\otimes_k k'$. En particulier, $Y$ est géométriquement réduit sur $k$.
\end{corollary}

\begin{proof}
Comme les $Z_\lambda$ sont réduits, les morphismes $f_\lambda$ se factorisent
en
\[
Z_\lambda\xrightarrow{g_\lambda}Y\xrightarrow{j}X,
\]
où $j$ est l'injection canonique ($\mathbf I$, 5.2.2). Il résulte alors de
(11.10.4) que $(g_\lambda)$ est une famille schématiquement dominante. Posons
$Y'_1=Y\otimes_k k'$, sous-préschéma fermé de $X'$, et soit $g'_\lambda$ le
morphisme déduit de $g_\lambda$ par extension de la base à $k'$, de sorte que
$f'_\lambda$ se factorise en
\[
Z'_\lambda\xrightarrow{g'_\lambda}Y'_1\xrightarrow{j'}X',
\]
où $j'$ est l'injection canonique. Il résulte de (11.10.6) que la famille
$(g'_\lambda)$ est schématiquement dominante. Mais par hypothèse les
$Z'_\lambda$ sont \emph{réduits}, donc ($\mathbf I$, 5.2.2) les
$g'_\lambda$ se factorisent en $Z'_\lambda\xrightarrow{h'_\lambda}Y'\to Y'_1$~; on conclut donc
de (11.10.2) que $Y'=Y'_1$ et $h'_\lambda=g'_\lambda$, ce qui établit le
corollaire.
\end{proof}

\begin{definition}[11.10.8]
\label{IV.11.10.8-fr}
Les notations étant celles de (11.10.5), supposons que $X$ soit un
$S$-préschéma. On dit que $(f_\lambda)$ est \emph{universellement
schématiquement dominante} (relativement à $S$) si, pour tout changement de
base $S'\to S$, la famille $(f'_\lambda)$ correspondant au changement de base
$X'=X\times_SS'\to X$, est schématiquement dominante.
\end{definition}

Lorsque $S$ est le spectre d'un corps, la prop. (11.10.6) signifie donc qu'une
famille schématiquement dominante est universellement schématiquement
dominante (relativement à $S$).

Lorsque les $f_\lambda$ sont des immersions canoniques $Z_\lambda\to X$, on
dira aussi que la famille des sous-préschémas $Z_\lambda$ est
\emph{universellement schématiquement dense dans $X$} (relativement à $S$) au
lieu de dire que la famille $(f_\lambda)$ est universellement schématiquement
dominante (relativement à $S$).

\begin{theorem}[11.10.9]
\label{IV.11.10.9-fr}
Les notations étant celles de (11.10.5), supposons que $X$ soit un
$S$-préschéma localement noethérien et que les $Z_\lambda$ soient tous
$S$-plats. Pour tout $s\in S$, soient
$X_s=X\times_S\operatorname{Spec}(k(s))$,
$(Z_\lambda)_s=Z_\lambda\times_S\operatorname{Spec}(k(s))$,
$(f_\lambda)_s=f_\lambda\times1:(Z_\lambda)_s\to X_s$. Pour que la famille
$(f_\lambda)$ soit universellement schématiquement dominante relativement à
$S$, il faut et il suffit que, pour tout $s\in S$, la famille
$((f_\lambda)_s)$ soit schématiquement dominante.
\end{theorem}

\begin{proposition}[11.10.10]
\label{IV.11.10.10-fr}
Soient $f:X\to S$ un morphisme plat localement de présentation finie, $U$ un
ouvert de $X$. Pour que $U$ soit universellement schématiquement dense dans
$X$ relativement à $S$, il faut et il suffit que, pour tout $s\in S$,
$U_s=U\cap X_s$ soit schématiquement dense dans $X_s$ (ou, ce qui revient au
même, que l'on ait $\operatorname{Ass}(\mathcal O_{X_s})\subset U_s$).
\end{proposition}
% IV3_B25_P172_END

% FIN SOURCE ega4/ega4-11-fr.tex

%
% DEBUT SOURCE ega4/ega4-12-fr.tex
\oldpage[IV-3]{173}
\section{Étude des fibres des morphismes plats de présentation finie}
\label{section:IV.12-fr}

\setcounter{subsection}{-1}
\subsection{Introduction}
\label{subsection:IV.12.0-fr}

\emph{Nous utiliserons dans tout ce paragraphe les notations générales de}
(9.4.1).

\begin{env}[12.0.1]
\label{IV.12.0.1-fr}
Étant donné un morphisme localement de présentation finie $f:X\to Y$, nous
avons vu (9.9) que pour certaines propriétés locales $\mathbf P$ de préschémas
sur des corps ou de Modules sur ces préschémas, l'ensemble $E$ des $x\in X$
tels que la propriété $\mathbf P$ ait lieu, pour la fibre $X_{f(x)}$, au point
$x$ de cette fibre, est localement constructible dans $X$. Nous nous proposons
de montrer que, pour la plupart de ces propriétés, si l'on suppose de plus que
le morphisme $f$ est plat, alors l'ensemble $E$ est même ouvert dans $X$. De
même ((9.2) à (9.8)) nous avons montré que si $f$ est de présentation finie,
et si cette fois $\mathbf P$ désigne certaines propriétés globales de
préschémas sur des corps ou de Modules sur ces préschémas, l'ensemble $F$ des
$y\in Y$ tels que la propriété $\mathbf P$ ait lieu pour la fibre $X_y$ est
localement constructible dans $Y$. Nous montrerons que si l'on suppose de plus
que le morphisme $f$ est propre et plat, $F$ est même ouvert dans $Y$.
\end{env}

\begin{env}[12.0.2]
\label{IV.12.0.2-fr}
La méthode générale de démonstration des propriétés en question comporte trois
étapes. On se ramène d'abord au cas où $Y$ est affine et $X$ de présentation
finie~; puis~:
\begin{enumerate}[label=\emph{\Alph*)}]
  \item À l'aide de (8.9.1) (et éventuellement d'autres résultats du \S~8) et
  de (11.2.6), on se ramène au cas où $X$ et $Y$ sont noethériens.
  \item On applique les résultats du \S~9 rappelés en (12.0.1) prouvant que
  $E$ (resp. $F$) est constructible.
  \item Pour voir que $E$ est ouvert, il suffit, en vertu de
  ($\mathbf 0_{\mathrm{III}}$, 9.2.5) de montrer que si $x\in E$, alors toute
  générisation $x'$ de $x$ appartient aussi à $E$. Utilisant ($\mathbf{II}$,
  7.1.7), on voit, puisque $Y$ est noethérien, qu'il existe un spectre d'anneau
  de valuation discrète $Y_1$ et un morphisme $h:Y_1\to X$ tel que, si $y_1$
  (resp. $y'_1$) est le point fermé (resp. le point générique) de $Y_1$, on ait
  $h(y_1)=x$, $h(y'_1)=x'$. On fait alors le changement de base
  $g=f\circ h:Y_1\to Y$~; compte tenu des résultats des \S\S~4 et 6 sur les
  préschémas localement noethériens sur des corps et les changements du corps
  de base, on est ramené à prouver l'assertion en question pour un point $x_1$
  de $X_1=X\times_YY_1$ au-dessus de $y_1$ et pour une générisation $x'_1$ de
  $x_1$ au-dessus de $y'_1$. Comme $Y_1=\operatorname{Spec}(A)$, où $A$ est
  un anneau de valuation discrète, d'uniformisante $t$, on doit en définitive,
  pour un $A$-module $M$, prouver une propriété de $M$, sachant que $M/tM$ a
  la même propriété et que $t$ est $M$-régulier (ce qui résulte de l'hypothèse
  de platitude)~; on utilise pour cela les résultats de (3.4) et de (5.12). On
  procède de même pour l'ensemble $F\subset Y$.
\end{enumerate}
\end{env}

\oldpage[IV-3]{174}
\subsection{Propriétés locales des fibres d'un morphisme plat localement de
présentation finie}
\label{subsection:IV.12.1-fr}

\begin{theorem}[12.1.1]
\EGAReaderTarget{ega4.c04.thm.12-1-1}{IV.12.1.1-fr}
Soient $f:X\to Y$ un morphisme localement de présentation finie, $\mathcal F$
un $\mathcal O_X$-Module $f$-plat et de présentation finie, $\Phi$ une partie
finie de $\mathbf Z\cup\{\pm\infty\}$, $k$ un entier. Les parties suivantes de
$X$ sont ouvertes~:
\begin{enumerate}[label=(\roman*)]
  \item L'ensemble des $x\in X$ tels que les dimensions des cycles premiers
  associés à $\mathcal F_{f(x)}$ et contenant $x$ soient des éléments de
  $\Phi$.
  \item L'ensemble des $x\in X$ tels que les cycles premiers associés à
  $\mathcal F_{f(x)}$ contenant $x$ aient tous la dimension $k$.
  \item L'ensemble des $x\in X$ tels que $x$ n'appartienne à aucun cycle
  premier immergé associé à $\mathcal F_{f(x)}$.
  \item L'ensemble des $x\in X$ tels que $\mathcal F_{f(x)}$ soit
  équidimensionnel au point $x$ et possède la propriété $(S_k)$ au point $x$
  (5.7.2).
  \item L'ensemble des $x\in X$ tels que
  $\operatorname{coprof}((\mathcal F_{f(x)})_x)\leq k$
  ($\mathbf 0$, 16.4.9).
  \item L'ensemble des $x\in X$ tels que $\mathcal F_{f(x)}$ soit un
  $\mathcal O_{X_{f(x)}}$-Module de Cohen-Macaulay au point $x$ (5.7.1).
  \item L'ensemble des $x\in X$ tels que $\mathcal F_{f(x)}$ soit
  géométriquement réduit au point $x$ (4.6.22).
  \item L'ensemble des $x\in X$ tels que $\mathcal F_{f(x)}$ soit
  géométriquement ponctuellement intègre au point $x$ (4.6.22).
\end{enumerate}
\end{theorem}

\begin{proof}
Les questions étant locales sur $X$, on se ramène d'abord au cas où
$X=\operatorname{Spec}(B)$ et $Y=\operatorname{Spec}(A)$ sont affines, $f$ un
morphisme de présentation finie. Il y a alors un sous-anneau noethérien $A_0$
de $A$, un $A_0$-préschéma de type fini $X_0$, un
$\mathcal O_{X_0}$-Module cohérent $\mathcal F_0$ tel que
$\mathcal F_0\otimes_{A_0}A$ soit isomorphe à $\mathcal F$ (8.9.1)~; en
outre, en vertu de (11.2.6), on peut supposer que $\mathcal F_0$ est
$Y_0$-plat (avec $Y_0=\operatorname{Spec}(A_0)$). Si $h:X\to X_0$ est la
projection canonique, l'ensemble $E$ des points $x\in X$ où l'une des
propriétés (i) à (viii) est vérifiée est égal à $h^{-1}(E_0)$, où $E_0$ est
l'ensemble des $x_0\in X_0$ où la propriété correspondante pour $\mathcal F_0$
est vérifiée~: cela résulte, pour les propriétés (i) à (iii), de (4.2.7)~; pour
les propriétés (iv) à (vi), de (6.7.1)~; pour les propriétés (vii) et (viii),
de (4.7.11).

On est ainsi ramené au cas où $A$ et $B$ sont noethériens, $f$ de type fini,
et $\mathcal F=\widetilde M$, où $M$ est un $B$-module de type fini. En vertu
de (9.9.2), on sait que l'ensemble $E$ est constructible pour toutes les
propriétés envisagées, et il reste dans chaque cas l'étape C) de (12.0.2), où
l'on doit prouver que $E$ est stable par générisation.

\begin{env}[12.1.1.1]
\label{IV.12.1.1.1-fr}
Commençons par le cas (iii) qui est le plus simple. Soit $x'\ne x$ une
générisation de $x$ dans $X$. Posons $y=f(x)$, $y'=f(x')$~; il existe un
spectre d'anneau de valuation discrète $Y_1$ et un morphisme $u:Y_1\to X$ tels
que, si $y_1$ et $y'_1$ sont le point fermé et le point générique de $Y_1$, on
ait $u(y_1)=x$ et $u(y'_1)=x'$ ($\mathbf{II}$, 7.1.7)~; posons
$g=f\circ u:Y_1\to Y$~; si $X_1=X\times_YY_1$, il existe une $Y_1$-section
$u_1:Y_1\to X_1$ telle que $u=g_1\circ u_1$, où $g_1:X_1\to X$ est la
projection canonique. Si $x_1=u_1(y_1)$, $x'_1=u_1(y'_1)$, on a donc
$g_1(x_1)=x$ et $g_1(x'_1)=x'$, et $x'_1$ est une générisation de $x_1$.
Appliquant de nouveau (4.2.7), on voit
\oldpage[IV-3]{175}
qu'on est ramené au cas où $Y=\operatorname{Spec}(A)$ est le spectre d'un
anneau de valuation discrète, $y$ étant le point fermé et $y'$ le point
générique de $Y$. Si $t$ est une uniformisante de $A$, l'hypothèse que
$\mathcal F$ est $f$-plat entraîne que $t$ est $\mathcal F$-régulier
($\mathbf 0_{\mathrm I}$, 6.3.4). Par hypothèse aucun des cycles premiers
associés immergés de $\mathcal F/t\mathcal F$ ne contient $x$. Alors, il
résulte de (3.4.4) que $x$ n'appartient à aucun cycle premier immergé associé à
$\mathcal F$, et il en est donc de même de toute générisation de $x$. En
particulier, $x'$ n'appartient à aucun cycle premier immergé associé à
$\mathcal F$, ni \emph{a fortiori} à aucun de ceux associés à
$\mathcal F_{y'}$ ($X_{y'}$ étant un sous-préschéma induit sur un ouvert de
$X$).
\end{env}

\begin{env}[12.1.1.2]
\label{IV.12.1.1.2-fr}
Envisageons ensuite les cas (iv) et (v) ((vi) n'est qu'un cas particulier de
(v)). On procède comme ci-dessus (en utilisant (6.7.1) au lieu de (4.2.7)) et
l'on est ramené au cas où $Y$ est le spectre d'un anneau de valuation discrète
$A$.

L'anneau $C=\mathcal O_{X,x}$ est alors localisé d'une $A$-algèbre de type
fini, donc est caténaire (5.6.4), et par hypothèse le $C$-module
$M_x/tM_x$ vérifie $(S_k)$ et est équidimensionnel (resp. on a
$\operatorname{coprof}(M_x/tM_x)\leq k$)~; on déduit donc de (5.12.2) (resp.
de ($\mathbf 0$, 16.4.10)) que $M_x$ vérifie $(S_k)$ et est équidimensionnel
(resp. que $\operatorname{coprof}(M_x)\leq k$) puisque $t$ est $M_x$-régulier
et appartient à l'idéal maximal de $C$. D'où les conclusions puisque $x'$ est
une générisation de $x$ et que $(\mathcal F_{y'})_{x'}=\mathcal F_{x'}$.
\end{env}

\begin{env}[12.1.1.3]
\label{IV.12.1.1.3-fr}
Passons aux cas (vii) et (viii). On peut évidemment remplacer $Y$ par le
sous-schéma intègre ayant $\overline{\{y'\}}$ pour espace sous-jacent, et $X$
par $f^{-1}(\overline{\{y'\}})$, sans changer les fibres en $y$ et $y'$, donc
on peut supposer $A$ intègre, de corps des fractions $K=k(y')$. On sait
((4.5.11) et (4.7.8)) qu'il existe une extension finie $K'$ de $K$ telle que,
pour le $(\mathcal O_X\otimes_AK')$-Module
$\mathcal F'=\mathcal F\otimes_AK'$, les cycles premiers associés soient
géométriquement irréductibles et que les longueurs géométriques de
$\mathcal F'$ aux points maximaux $z'_i$ de son support soient respectivement
égales aux longueurs de $\mathcal F'_{z'_i}$ en ces points~; il revient donc
au même de dire que $\mathcal F$ est géométriquement réduit (resp.
géométriquement ponctuellement intègre) en un point $x'$ de
$X_{y'}=X\otimes_AK$, ou de dire que $\mathcal F'$ est réduit (resp. intègre)
aux points de $X\otimes_AK'$ au-dessus de $x'$, compte tenu de (4.2.7). Soit
$A'$ une $A$-algèbre finie dont $K'$ est le corps des fractions~; il résulte
de (4.7.11) qu'en tout point de $X\otimes_AA'$ au-dessus de $x$,
$\mathcal F\otimes_AA'$ a la propriété (vii) (resp. (viii)). Remplaçant $A$
par $A'$ et $X$ par $X\otimes_AA'$, on voit qu'on peut se borner à prouver que
$\mathcal F$ est réduit (resp. intègre) en toute générisation $x'$ de $x$. On
procède alors comme dans (12.1.1.1), en utilisant cette fois (4.7.11), et l'on
est encore ramené au cas où $A$ est un anneau de valuation discrète, $y$ étant
le point fermé, $y'$ le point générique de $Y$, et l'uniformisante $t$ de $A$
étant un élément $M$-régulier. Comme par hypothèse
$\mathcal F/t\mathcal F$ est réduit (resp. intègre) au point $x$, il résulte de
(3.4.6) (resp. (3.4.5)) que $\mathcal F$ est réduit (resp. intègre) au point
$x$, donc aussi dans un voisinage de $x$, et en particulier au point $x'$, ce
qui achève la démonstration dans les cas (vii) et (viii).
\end{env}

\begin{env}[12.1.1.4]
\label{IV.12.1.1.4-fr}
Reste à examiner les cas (i) et (ii). Remplaçant $Y$ par le sous-schéma intègre
ayant $\overline{\{y'\}}$ pour espace sous-jacent, on peut se borner au cas où
$Y$ est irréductible et $y'$ son point générique. Soit $z'$ un point générique
d'un cycle premier associé à $\mathcal F_{y'}$
\oldpage[IV-3]{176}
contenant $x'$, et soit $Z'$ l'adhérence de $z'$ dans $X$, de sorte que l'on a
$x\in Z'$. Pour traiter le cas (i), il va suffire de prouver la

\begin{proposition}[12.1.1.5]
\label{IV.12.1.1.5-fr}
Soient $Y$ un préschéma localement noethérien, $f:X\to Y$ un morphisme
localement de type fini, $\mathcal F$ un $\mathcal O_X$-Module cohérent et
$f$-plat~; pour tout $y\in Y$, posons
$\mathcal F_y=\mathcal F\otimes_{\mathcal O_Y}k(y)$. Soient $z'$ un point de
$X$, $y'=f(z')$, et supposons que $z'\in\operatorname{Ass}(\mathcal F_{y'})$.
Soient $Z'$ (resp. $Y'$) le sous-préschéma réduit de $X$ (resp. $Y$) ayant
pour espace sous-jacent l'adhérence de $\{z'\}$ dans $X$ (resp. de $\{y'\}$
dans $Y$). Alors, pour tout $y\in f(Z')$, les dimensions de toutes les
composantes irréductibles de $Z'_y$ sont égales à $\dim(Z'_{y'})$ (dimension
de l'adhérence de $z'$ dans $X_{y'}$) (ce que nous exprimerons plus tard
(13.2.2) en disant que la restriction $Z'\to Y'$ de $f$ est un morphisme
équidimensionnel)~; en outre, en tout point maximal $z$ de $Z'_y$, on a
$z\in\operatorname{Ass}(\mathcal F_y)$.
\end{proposition}

Pour cela, nous allons nous ramener au cas où $Y$ est spectre d'un anneau de
valuation discrète. Prenons un spectre d'anneau de valuation discrète $Y_1$ et
un morphisme $h:Y_1\to X$ tel que $h(y_1)=z$, $h(y'_1)=z'$, $y_1$ et $y'_1$
étant le point fermé et le point générique de $Y_1$ respectivement
($\mathbf{II}$, 7.1.7). Soient $g=f\circ h:Y_1\to Y$,
$X_1=X\times_YY_1$, et soient $f_1:X_1\to Y_1$, $g_1:X_1\to X$ les projections
canoniques~; il y a une $Y_1$-section $h_1:Y_1\to X_1$ telle que
$h=g_1\circ h_1$ ($\mathbf I$, 3.3.14). Si
$Z'_1=Z'_{y'}\otimes_{k(y')}k(y'_1)$, on sait (4.2.6) que les composantes
irréductibles de $Z'_1$ dominent $Z'_{y'}$~; soit $z'_1$ le point générique de
l'une de ces composantes qui contient $h_1(y'_1)$, de sorte que
$f_1(z'_1)=y'_1$ et $g_1(z'_1)=z'$. Comme $h_1(y_1)$ est spécialisation de
$h_1(y'_1)$, il est \emph{a fortiori} spécialisation de $z'_1$~; soit $z_1$ un
point générique d'une des composantes irréductibles de
$\overline{\{z'_1\}}\cap(X_1)_{y_1}$ contenant $h_1(y_1)$. Alors $g_1(z_1)$
est spécialisation de $g_1(z'_1)=z'$ dans $X$, et
$z=h(y_1)=g_1(h_1(y_1))$ est spécialisation de $g_1(z_1)$~; mais comme
$f(g_1(z_1))=y$ et que $z$ est point générique de $Z'_y$, on a nécessairement
$z=g_1(z_1)$.

Supposons maintenant que l'on ait prouvé que, si l'on pose
$\mathcal F_1=\mathcal F\otimes_{\mathcal O_Y}\mathcal O_{Y_1}$, on a
$z_1\in\operatorname{Ass}((\mathcal F_1)_{y_1})$~; il résultera de (4.2.7) que
$z\in\operatorname{Ass}(\mathcal F_y)$~; en outre, les dimensions de
$\overline{\{z_1\}}\cap(X_1)_{y_1}$ et de $\overline{\{z\}}\cap X_y$ sont
égales, et de même les dimensions de
$\overline{\{z'_1\}}\cap(X_1)_{y'_1}$ et de
$\overline{\{z'\}}\cap X_{y'}$ (4.2.7)~; on s'est donc bien ramené, comme
annoncé, au cas où $Y$ est spectre d'un anneau de valuation discrète $A$, $y$
et $y'$ étant son point fermé et son point générique respectivement.

Cela étant, comme $z'\in\operatorname{Ass}(\mathcal F_{y'})$, on a
\emph{a fortiori} $z'\in\operatorname{Ass}(\mathcal F)$. Si $t$ est une
uniformisante de $A$, $t$ est $\mathcal F$-régulier par platitude, donc (3.4.3)
on a $z\in\operatorname{Ass}(\mathcal F/t\mathcal F)
=\operatorname{Ass}(\mathcal F_y)$. Quant à l'assertion relative aux
dimensions, elle résulte de (7.1.13), appliqué à un sous-préschéma fermé de $X$
ayant $Z'$ pour espace sous-jacent.

Abordons enfin le cas (ii) de (12.1.1). Avec les mêmes notations, il faut
prouver que si tous les cycles premiers associés à $\mathcal F_y$ et contenant
$x$ ont la dimension $k$, il en est de même de tous les cycles premiers
associés à $\mathcal F_{y'}$ et contenant $x'$. Or on vient de voir que tout
cycle premier associé à $\mathcal F_{y'}$ et contenant $x'$ a même dimension
que l'un des cycles premiers associés à $\mathcal F_y$ et contenant $x$~; cela
démontre donc (ii).
\end{env}
\end{proof}

\begin{remarks}[12.1.2]
\EGAReaderTarget{ega4.c04.item.12-1-2}{IV.12.1.2-fr}
\begin{enumerate}[label=(\roman*)]
  \item Sous les conditions de (12.1.1.5) avec $\mathcal F=\mathcal O_X$ (de
  sorte que $f$ est plat), on ne peut en général affirmer que la restriction
  $Z'\to Y'$ de $f$ est un morphisme plat ni même un morphisme ouvert. C'est ce
  que montre l'exemple (6.5.5,
  \oldpage[IV-3]{177}
  (ii)) où l'on prend pour $Z'$ une des deux composantes irréductibles de
  $X=\operatorname{Spec}(B)$. Il est immédiat que ce morphisme restriction
  n'est pas ouvert aux points de $Z'$ au-dessus du «~point double~» de $Y'=Y$.
  \item Dans les hypothèses de (12.1.1.5), avec
  $\mathcal F=\mathcal O_X$, on ne peut affaiblir la condition «~$f$ est
  plat~» en «~$f$ est universellement ouvert~» (cf. (2.4.6)), comme nous le
  verrons plus loin sur un exemple (14.4.10, (i)).
\end{enumerate}
\end{remarks}

\begin{corollary}[12.1.3]
\label{IV.12.1.3-fr}
Sous les hypothèses de (12.1.1), la fonction
$x\mapsto\operatorname{coprof}((\mathcal F_{f(x)})_x)$ est semi-continue
supérieurement dans $X$ et la fonction
$x\mapsto\operatorname{prof}_x^*(\mathcal F_{f(x)})$ (10.8.1) est
semi-continue inférieurement dans $X$.
\end{corollary}

\begin{proof}
La première assertion n'est autre qu'une formulation équivalente de (12.1.1,
(v)). Pour prouver la seconde, on peut d'abord, compte tenu de (10.8.7) et
(10.8.8), se ramener comme dans (12.1.1) au cas où $Y$ est le spectre d'un
anneau de valuation discrète $A$ de point fermé $y$ et de point générique $y'$,
$X=\operatorname{Spec}(B)$ étant affine, $\mathcal F=\widetilde M$, où $M$ est
un $B$-module de type fini. Puisque $\mathcal F$ est $f$-plat, toute composante
irréductible de $\operatorname{Supp}(\mathcal F)$ qui contient $x\in X_y$
domine $Y$ (2.3.4), autrement dit son point générique $z'$ appartient à
$X_{y'}$~; les composantes irréductibles $Z$ de
$\operatorname{Supp}(\mathcal F)$ qui contiennent une générisation $x'$ de
$x$ appartenant à $X_{y'}$ sont donc exactement celles qui contiennent $x$ et
en outre, d'après (12.1.1.5), on a $\dim(Z_{y'})=\dim(Z_y)$, d'où, si
$T=\operatorname{Supp}(\mathcal F)$,
$\dim_x(T_y)=\dim_{x'}(T_{y'})$. De plus, pour une uniformisante $t$ de $A$, on
a vu dans (12.1.1, (v)) que l'on a
$\operatorname{coprof}(M_x)=\operatorname{coprof}(M_x/tM_x)$, et puisque
d'autre part $\operatorname{coprof}(M_{x'})\leq\operatorname{coprof}(M_x)$
(6.11.5), on voit que l'on a
$\operatorname{coprof}((\mathcal F_{y'})_{x'})\leq
\operatorname{coprof}((\mathcal F_y)_x)$, la relation
$\operatorname{prof}_x^*(\mathcal F_y)\leq
\operatorname{prof}_{x'}^*(\mathcal F_{y'})$ découle alors de (10.8.7).
\end{proof}

\begin{remark}[12.1.4]
\label{IV.12.1.4-fr}
On déduit aussi de (12.1.1, (i)) que la fonction
$x\mapsto\dim_x(X_{f(x)})$ est semi-continue supérieurement dans $X$, mais
nous verrons plus loin (13.1.3) que cette propriété est vraie même sans
supposer $f$ plat.
\end{remark}

\begin{corollary}[12.1.5]
\label{IV.12.1.5-fr}
Soient $Y$ un préschéma localement noethérien, $f:X\to Y$ un morphisme
localement de type fini, $\mathcal F$ un $\mathcal O_X$-Module cohérent et
$f$-plat, $\mathcal G$ un $\mathcal O_Y$-Module cohérent, $x$ un point de $X$,
$y=f(x)$. Supposons que $\mathcal G$ ait la propriété $(S_k)$ au point $y$, et
que $\mathcal F_y$ ait la propriété $(S_k)$ au point $x$ et soit
équidimensionnel au point $x$. Alors~:
\begin{enumerate}[label=(\roman*)]
  \item $\mathcal F\otimes_{\mathcal O_Y}\mathcal G$ possède la propriété
  $(S_k)$ au point $x$.
  \item Si de plus $Y$ est localement immersible dans un schéma régulier, ou
  est un préschéma excellent (7.8.5), il existe un voisinage de $x$ dans $X$
  tel que $\mathcal F\otimes_{\mathcal O_Y}\mathcal G$ ait la propriété
  $(S_k)$ dans ce voisinage.
\end{enumerate}
\end{corollary}

\begin{proof}
En effet, par (12.1.1, (iv)), il existe un voisinage ouvert $V$ de $x$ tel que
pour tout $x'\in V$, $\mathcal F_{f(x')}$ vérifie $(S_k)$ au point $x'$.
D'autre part, $\mathcal G$ vérifie $(S_k)$ en tous les points de
$Y'=\operatorname{Spec}(\mathcal O_y)$ (5.7.2). Remplaçant $Y$ par $Y'$ et $X$
par $V\times_YY'$, on est ramené (compte tenu de ($\mathbf I$, 3.6.5)) au cas
où $\mathcal G$ vérifie $(S_k)$ dans $Y$ tout entier et où, pour tout
$y\in f(X)$, $\mathcal F_y$ possède la propriété $(S_k)$ en tous les points de
$f^{-1}(y)$~; comme $\mathcal F$ est $f$-plat, on sait alors (6.4.1) que
$\mathcal F\otimes_{\mathcal O_Y}\mathcal G$ possède la propriété $(S_k)$ dans
$X$, ce qui prouve (i). Pour démontrer (ii), il suffit d'observer que, sous les
hypothèses faites, $\mathcal G$ possède la
\oldpage[IV-3]{178}
propriété $(S_k)$ dans un voisinage $U$ de $y$ dans $Y$ ((6.11.2) et (7.8.3,
(iv)))~; on conclut alors de la même manière en remplaçant cette fois $Y$ par
$U$ et $X$ par $V\times_YU$.
\end{proof}

\begin{theorem}[12.1.6]
\label{IV.12.1.6-fr}
Soient $f:X\to Y$ un morphisme plat et localement de présentation finie, $k$
un entier. Les parties suivantes de $X$ sont ouvertes~:
\begin{enumerate}[label=(\roman*)]
  \item L'ensemble des points $x\in X$ tels que $X_{f(x)}$ vérifie la propriété
  $(S_k)$ au point $x$.
  \item L'ensemble des $x\in X$ tels que $X_{f(x)}$ vérifie la propriété
  $(R_k)$ géométrique au point $x$, soit équidimensionnel au point $x$ et tels
  en outre que $x$ n'appartienne à aucun cycle premier immergé associé à
  $X_{f(x)}$.
  \item L'ensemble des $x\in X$ tels que $X_{f(x)}$ soit géométriquement
  régulier (i.e. lisse) au point $x$.
  \item L'ensemble des $x\in X$ tels que $X_{f(x)}$ soit géométriquement
  normal au point $x$.
\end{enumerate}
\end{theorem}

\begin{proof}
Les étapes A) et B) de (12.0.2) se font ici comme dans (12.1.1)~; pour l'étape
A), on utilise (6.7.8), ainsi que (4.2.7) pour (ii)~; pour l'étape B), on
utilise (9.9.2) et (9.9.4). Reste à examiner l'étape C) dans chaque cas.
\begin{enumerate}[label=(\roman*)]
  \item[(i)] Comme dans (12.1.1.2), on se ramène (en utilisant (6.7.8)) au cas
  où $Y$ est le spectre d'un anneau de valuation discrète $A$,
  $X=\operatorname{Spec}(B)$, où $B$ est une $A$-algèbre de type fini, $x$,
  $x'$ deux points de $X$ tels que $f(x)=y$ soit le point fermé et
  $y'=f(x')$ le point générique de $Y$, $x'$ étant en outre une générisation
  de $x$. Comme il s'agit de prouver que $X$ vérifie $(S_k)$ au point $x'$,
  on peut remplacer $X$ par $\operatorname{Spec}(\mathcal O_{X,x})$,
  c'est-à-dire supposer l'anneau $B$ local, l'homomorphisme $A\to B$ étant
  local et $B$ étant une $A$-algèbre essentiellement de type fini (1.3.8).
  Alors, par hypothèse, si $t$ est une uniformisante de $A$, $t$ est un
  élément $B$-régulier par platitude, et $B/tB$ vérifie $(S_k)$. Mais comme
  $A$ est un anneau universellement caténaire (5.6.4), $B$ est caténaire,
  donc, par (5.12.4), il en résulte que $B$ vérifie $(S_k)$, ce qui achève la
  démonstration dans le cas (i).
  \item[(iii)] Ici l'étape C) est inutile~; comme $f$ est plat, on sait en
  effet (6.8.7) que lorsque $Y$ est localement noethérien et $f$ localement de
  type fini, l'ensemble des $x\in X$ tels que $X_{f(x)}$ soit géométriquement
  régulier au point $x$ est ouvert dans $X$.
  \item[(ii)] En raisonnant comme dans (12.1.1.3), pour prouver que la propriété
  considérée est stable par générisation, on peut d'abord, en considérant une
  extension finie quelconque $K'$ de $k(y')$, et tenant compte de la définition
  (6.7.6) de la propriété $(R_k)$ géométrique, ainsi que de l'invariance par
  changement du corps de base des deux autres propriétés figurant dans (ii),
  remplacer dans (ii) la propriété $(R_k)$ géométrique par la propriété
  $(R_k)$. Procédant alors comme dans (i), on se ramène au cas où $Y$ est
  spectre d'un anneau de valuation discrète $A$, et puisque $A$ est caténaire,
  il suffit d'appliquer (5.12.5) pour conclure.
  \item[(iv)] L'ensemble des points $x\in X$ tels que $X_{f(x)}$ soit
  géométriquement normal au point $x$ est contenu dans celui des $x\in X$ tels
  que $X_{f(x)}$ soit géométriquement ponctuellement intègre et vérifie $(S_2)$
  au point $x$, et ce dernier est ouvert dans $X$ en vertu de (12.1.1,
  (viii)). On peut donc déjà supposer que, pour tout $x\in X$, $X_{f(x)}$ est
  géométriquement ponctuellement intègre et vérifie $(S_2)$ au point $x$, et
  \emph{a fortiori} il est équidimensionnel. D'autre part, en vertu du critère
  de Serre (5.8.6), dire que $X_{f(x)}$ est géométriquement normal au point $x$
  signifie que $X_{f(x)}$ vérifie $(S_2)$ et la propriété $(R_1)$ géométrique
  en $x$~;
  \oldpage[IV-3]{179}
  mais en vertu de (ii) et des remarques précédentes, cet ensemble est
  intersection de deux ouverts de $X$, donc est ouvert dans $X$. C.Q.F.D.
\end{enumerate}
\end{proof}

\begin{corollary}[12.1.7]
\label{IV.12.1.7-fr}
Soit $f:X\to Y$ un morphisme localement de présentation finie. Alors
l'ensemble des points $x\in X$ où $f$ possède l'une des propriétés suivantes
(6.8.1)~:
\begin{enumerate}[label=(\roman*)]
  \item vérifier la propriété $(S_k)$~;
  \item être de coprofondeur $\leq n$~;
  \item être de Cohen-Macaulay~;
  \item être régulier (ou lisse, ce qui revient au même)~;
  \item être normal~;
  \item être réduit~;
\end{enumerate}
est ouvert dans $X$.
\end{corollary}

\begin{proof}
En effet, il résulte de (11.3.1) que l'ensemble des $x\in X$ où $f$ est plat
est ouvert. On peut donc se borner au cas où $f$ est plat, et alors le
corollaire résulte de (12.1.1, (iv), (vi) et (vii)) et de (12.1.6, (i), (ii)
et (iv)).
\end{proof}

\begin{remarks}[12.1.8]
\label{IV.12.1.8-fr}
\begin{enumerate}[label=(\roman*)]
  \item Dans (12.1.6, (ii)), on ne peut supprimer l'hypothèse que $x$
  n'appartienne à aucun cycle premier associé immergé de $X_{f(x)}$. C'est ce
  que montre l'exemple (5.12.3), où l'on prend $X=\operatorname{Spec}(A)$, $Y$
  étant le spectre de l'anneau local $A_0$ de $k[T]$ correspondant à l'idéal
  premier $(T)$ et le morphisme $X\to Y$ correspondant à l'homomorphisme
  injectif $A_0\to A$, déduit par localisation et passage au quotient de
  l'injection $k[T]\to k[T,U,V]$~; ce morphisme est plat puisque $A$ est un
  $A_0$-module sans torsion ($\mathbf 0_{\mathrm I}$, 6.3.4). Ici la fibre
  $X_y$ au point fermé $y$ de $Y$, égale à $\operatorname{Spec}(A/tA)$, est
  irréductible, de dimension 1 et vérifie la condition $(R_0)$ géométrique
  puisque l'anneau local en son point générique est un corps. Par contre, au
  point générique $y'$ de $Y$, la fibre $X_{y'}$ a deux composantes
  irréductibles de dimensions respectives 0 et 1, et celle de dimension 0
  n'est pas réduite, donc $X_{y'}$ ne vérifie pas la condition $(R_0)$.
  \item Dans (12.1.6, (ii)), on ne peut pas non plus supprimer l'hypothèse que
  $X_{f(x)}$ est équidimensionnel au point $x$. On le voit ici sur l'exemple
  (5.12.6) avec $X=\operatorname{Spec}(A)$, $Y=\operatorname{Spec}(A_0)$, où
  $A_0$ est défini comme dans (i), le morphisme $X\to Y$ provenant encore de
  l'injection $k[T]\to k[T,U,V,W]$ par localisation et passage au quotient, et
  étant plat puisque $A$ est un $A_0$-module sans torsion
  ($\mathbf 0_{\mathrm I}$, 6.3.4). La fibre $X_y$ au point fermé $y$ de $Y$
  est réduite (donc vérifie $(S_1)$) et vérifie $(R_1)$, mais a deux composantes
  irréductibles de dimensions 2 et 1, tandis que la fibre $X_{y'}$ au point
  générique $y'$ de $Y$ ne vérifie pas $(R_1)$.
\end{enumerate}
\end{remarks}
\subsection{Propriétés locales et globales des fibres d'un morphisme propre,
plat et de présentation finie}
\label{subsection:IV.12.2-fr}

\begin{theorem}[12.2.1]
\label{IV.12.2.1-fr}
Soient $f:X\to Y$ un morphisme \emph{propre} de présentation finie,
$\mathcal F$ un $\mathcal O_X$-Module $f$-plat et de présentation finie,
$\Phi$ une partie finie de $\mathbf Z\cup\{\pm\infty\}$, $k$ un entier. Les
parties suivantes de $Y$ sont ouvertes~:
\begin{enumerate}[label=(\roman*)]
  \item L'ensemble des $y\in Y$ tels que l'ensemble des dimensions des cycles
  premiers associés à $\mathcal F_y$ soit contenu dans $\Phi$.
  \item L'ensemble des $y\in Y$ tels que l'ensemble des dimensions des
  composantes irréductibles de $\operatorname{Supp}(\mathcal F_y)$ contienne
  $\Phi$.
  \item L'ensemble des $y\in Y$ tels que tous les cycles premiers associés à
  $\mathcal F_y$ aient une même dimension égale à $k$.
  \item L'ensemble des $y\in Y$ tels que $\mathcal F_y$ n'ait pas de cycle
  premier associé immergé et que l'ensemble des dimensions des composantes
  irréductibles de $\operatorname{Supp}(\mathcal F_y)$ soit égal à $\Phi$.
  \oldpage[IV-3]{180}
  \item L'ensemble des $y\in Y$ tels que $\mathcal F_y$ ait la propriété
  $(S_k)$ et soit équidimensionnel en tout point de $X_y$.
  \item L'ensemble des $y\in Y$ tels que
  $\operatorname{coprof}(\mathcal F_y)\leq k$.
  \item L'ensemble des $y\in Y$ tels que $\mathcal F_y$ soit un
  $\mathcal O_{X_y}$-Module de Cohen-Macaulay.
  \item L'ensemble des $y\in Y$ tels que $\mathcal F_y$ soit géométriquement
  réduit.
  \item L'ensemble des $y\in Y$ tels que $\mathcal F_y$ soit géométriquement
  ponctuellement intègre en chaque point de $X_y$.
  \item L'ensemble des $y\in Y$ tels que $\mathcal F_y$ soit géométriquement
  intègre.
  \item L'ensemble des $y\in Y$ tels que $\mathcal F_y$ n'ait pas de cycle
  premier associé immergé et que la somme $l(y)$ des multiplicités totales
  (4.7.12) de $\mathcal F_y$ aux points maximaux de
  $\operatorname{Supp}(\mathcal F_y)$ soit $\leq k$.
\end{enumerate}
\end{theorem}

\begin{proof}
À l'exception de (ii), (iii), (iv), (x) et (xi), les propriétés $P(y)$
considérées sont de la forme suivante~: «~pour tout $x\in X_y$, $\mathcal F_y$
a la propriété $Q(x)$~», où $Q(x)$ est l'une des propriétés (i) à (viii) de
(12.1.1). Il résulte de (12.1.1) que l'ensemble $U$ des $x\in X$ tels que
$Q(x)$ soit vraie est ouvert, et l'ensemble $V$ des $y\in Y$ tels que $P(y)$
soit vraie n'est autre que l'ensemble $Y-f(X-U)$~; dans tous ces cas, le
théorème est donc déjà vrai en supposant seulement le morphisme $f$ \emph{fermé}.
Pour (iii), on applique (i) avec $\Phi$ réduit à un seul élément. Il reste à
examiner les cas (ii), (iv) et (xi) séparément, ((x) se déduisant aussitôt de
(xi) et de (4.7.14)) en appliquant toujours la méthode décrite dans (12.0.2).
\end{proof}

\begin{env}[12.2.1.1]
\label{IV.12.2.1.1-fr}
Cas (ii)~: Les étapes A) et B) de la démonstration procèdent comme dans le
début de (12.1.1)~; pour l'étape A), on utilise (8.9.1), (11.2.6) ainsi que
(4.2.7)~; pour l'étape B), on utilise (9.5.5) appliqué à
$\operatorname{Supp}(\mathcal F)$. Il reste à démontrer que si (en supposant
$Y$ noethérien) l'ensemble des dimensions des composantes irréductibles de
$\operatorname{Supp}(\mathcal F_y)$ contient $\Phi$, alors il en est de même
de l'ensemble des dimensions des composantes irréductibles de
$\operatorname{Supp}(\mathcal F_{y'})$, pour toute \emph{générisation} $y'$ de
$y$ dans $Y$. Soit $Y_1$ un spectre d'anneau de valuation discrète tel que, si
$y_1$ et $y'_1$ sont le point fermé et le point générique de $Y_1$, il y ait un
morphisme $h:Y_1\to Y$ tel que $h(y_1)=y$, $h(y'_1)=y'$ (\textbf{II}, 7.1.7).
Appliquant (4.2.7), on voit qu'on peut remplacer $Y$ par $Y_1$ et $X$ par
$X_1=X\times_YY_1$, autrement dit, se borner au cas où $Y$ est spectre
d'anneau de valuation discrète, $y$ son point fermé et $y'$ son point générique.

Utilisant (\textbf{Err}\textsubscript{\textbf{III}}, 30), on peut se borner au
cas où $\operatorname{Supp}(\mathcal F)=X$, de sorte que $f$ est quasi-plat
(2.3.3)~; les composantes irréductibles $Z_i$ de $X$ dominent alors $Y$
(2.3.4), autrement dit leurs points génériques appartiennent à $X_{y'}$, et
les $Z_i\cap X_{y'}$ sont les composantes irréductibles de $X_{y'}$. Mais toute
composante irréductible de $X_y$ est contenue dans un des $Z_i$, donc est une
composante irréductible de $(Z_i)_y$~; or, on sait (7.1.13) que les dimensions
de toutes les composantes irréductibles de $(Z_i)_y$ sont égales à
$\dim((Z_i)_{y'})$, ce qui achève de prouver (ii).
\end{env}

\begin{env}[12.2.1.2]
\label{IV.12.2.1.2-fr}
Cas (iv)~: Le même raisonnement qu'au début montre, en utilisant (12.1.1,
(iii)) que l'on peut déjà supposer (en remplaçant $Y$ par un ouvert de $Y$) que
pour tout $y\in Y$, $\mathcal F_y$ n'a pas de cycle premier associé immergé.
L'assertion du cas (iv) est alors une conséquence immédiate des assertions des
cas (i) et (ii).
\end{env}

\begin{env}[12.2.1.3]
\label{IV.12.2.1.3-fr}
\oldpage[IV-3]{181}
Cas (xi)~: Pour l'étape A) du raisonnement, on utilise (8.9.1), (11.2.6),
(8.10.5, (xii)) (pour conserver l'hypothèse que $f$ est propre), ainsi que
(4.2.7), (4.5.6) et (4.7.9). Pour l'étape B), on voit, comme dans le cas (iv),
qu'on peut supposer que pour tout $y\in Y$, $\mathcal F_y$ n'a pas de cycle
premier associé immergé, et l'on applique (9.8.8) qui prouve que la fonction
$z\mapsto l(z)$ est constructible. Il reste donc à voir (en supposant $Y$
noethérien et $f$ propre) que pour toute générisation $y'$ d'un point $y\in Y$,
on a $l(y')\leq l(y)$. En raisonnant comme dans (12.1.1.3), on voit (à l'aide
de (4.7.8) et (4.5.11)) qu'on peut supposer que les composantes irréductibles
de $\operatorname{Supp}(\mathcal F_{y'})$ sont géométriquement irréductibles
et que $l(y')$ est la \emph{somme des longueurs} des
$(\mathcal F_{y'})_{z'_j}$ aux points maximaux $z'_j$ de
$\operatorname{Supp}(\mathcal F_{y'})$. Appliquant de nouveau (4.2.7), (4.5.6)
et (4.7.9), on se ramène, comme dans (12.2.1.1), au cas où $Y$ est un spectre
d'anneau de valuation discrète, $y$ son point fermé et $y'$ son point générique.
Le fait que $l(y')\leq l(y)$ sera alors conséquence du

\begin{lemma}[12.2.1.4]
\label{IV.12.2.1.4-fr}
Soient $Y$ un spectre d'anneau de valuation discrète, $y$ son point fermé,
$y'$ son point générique, $f:X\to Y$ un morphisme propre, $\mathcal F$ un
$\mathcal O_X$-Module cohérent et $f$-plat, $z_i$ (resp. $z'_j$) les points
maximaux de $\operatorname{Supp}(\mathcal F_y)$ (resp.
$\operatorname{Supp}(\mathcal F_{y'})$). Supposons que $\mathcal F_y$ n'ait
pas de cycle premier associé immergé. Alors on a
\[
  \sum_j\operatorname{long}((\mathcal F_{y'})_{z'_j})
  \leq\sum_i\operatorname{long}((\mathcal F_y)_{z_i}).
\]
\end{lemma}

On a $Y=\operatorname{Spec}(A)$, où $A$ est un anneau de valuation discrète,
dont nous désignerons par $t$ une uniformisante, de sorte que
$\mathcal F_y=\mathcal F/t\mathcal F$. Comme $\mathcal F$ est $f$-plat, les
$z'_j$ sont aussi les points maximaux de $\operatorname{Supp}(\mathcal F)$
(2.3.4)~; pour tout $z_i$, désignons par $z'_{ik}$ ceux des $z'_j$ qui sont des
générisations de $z_i$~; il résulte de (3.4.1.1) que l'on a
\[
  \operatorname{long}((\mathcal F_y)_{z_i})
  \geq\sum_k\operatorname{long}((\mathcal F_{y'})_{z'_{ik}})
\]
d'où en sommant
\[
  \sum_i\operatorname{long}((\mathcal F_y)_{z_i})
  \geq\sum_i\sum_k\operatorname{long}((\mathcal F_{y'})_{z'_{ik}}).
\]

Le lemme sera donc prouvé si nous établissons que pour tout $z'_j$, il y a au
moins un indice $i$ tel que $z'_j$ soit l'un des $z'_{ik}$. Or, comme $f$ est
propre (donc fermé) et $f(z'_j)=y'$, il existe $x\in X$ qui est une
spécialisation de $z'_j$ et est tel que $f(x)=y$, autrement dit
$\overline{\{z'_j\}}\cap X_y$ n'est pas vide. Comme $t$ est
$\mathcal F$-régulier par platitude, on déduit de (3.4.3) qu'il y a au moins un
point de $\operatorname{Ass}(\mathcal F_y)$ dont $z'_j$ est une
générisation~; mais comme les points de $\operatorname{Ass}(\mathcal F_y)$
sont par hypothèse les $z_i$, cela termine la démonstration.
\end{env}

\begin{corollary}[12.2.2]
\label{IV.12.2.2-fr}
Sous les hypothèses de (12.2.1)~:
\begin{enumerate}[label=(\roman*)]
  \item La fonction $y\mapsto\dim(\operatorname{Supp}(\mathcal F_y))$ est
  continue (donc localement constante) dans $Y$.
  \item La fonction $y\mapsto\operatorname{coprof}(\mathcal F_y)$ (5.7.1) est
  semi-continue supérieurement dans $Y$.
  \item La fonction $y\mapsto l(y)$ (12.2.1, (xi)) est semi-continue
  supérieurement dans $Y$ lorsque les $\mathcal F_y$ n'ont pas de cycle
  premier associé immergé.
  \oldpage[IV-3]{182}
  \item La fonction $y\mapsto\operatorname{prof}^{*}(\mathcal F_y)$ (10.8.1)
  est semi-continue inférieurement dans $Y$.
\end{enumerate}
\end{corollary}

\begin{proof}
(i) Il résulte de (12.2.1, (i)) appliqué à $\Phi$ égal au plus petit intervalle de
$\mathbf N$ contenant les dimensions des cycles premiers associés à
$\mathcal F_y$, que $z\mapsto\dim(\operatorname{Supp}(\mathcal F_z))$ est
semi-continue supérieurement au point $y$~; il résulte d'autre part de
(12.2.1, (ii)) appliqué à $\Phi$ égal à l'ensemble des dimensions des
composantes irréductibles de $\operatorname{Supp}(\mathcal F_y)$, que cette
même fonction est semi-continue inférieurement au point $y$, d'où la
conclusion. Les assertions (ii) et (iii) ne sont que d'autres formulations de
(12.2.1, (vi) et (xi)). Enfin, d'après (12.1.3), l'ensemble des $x\in X$ tels
que $\operatorname{prof}^{*}_{f(x)}(\mathcal F_{f(x)})\leq k$ est ouvert, et la
conclusion de (iv) en résulte par le même raisonnement qu'au début de (12.2.1).
\end{proof}

\begin{remarks}[12.2.3]
\label{IV.12.2.3-fr}
\begin{enumerate}[label=(\roman*)]
  \item On observera que pour (12.2.1, (ii)), on peut se dispenser de
  l'hypothèse que $f$ est \emph{propre}.

  \item Dans (12.2.1, (ii)), on ne peut remplacer la condition que l'ensemble
  $E(y)$ des dimensions des composantes irréductibles de
  $\operatorname{Supp}(\mathcal F_y)$ contienne $\Phi$, par la condition que
  $E(y)$ soit \emph{égal à} $\Phi$. En effet, soit $k$ un corps, et considérons
  l'espace projectif à 3 dimensions $P=\mathbf P^3_k=\operatorname{Proj}(C)$,
  où $C=k[t_0,t_1,t_2,t_3]$ ($t_i$ indéterminées)~; dans $P$, soit $X_0$ le
  sous-préschéma fermé «~réunion de la droite $X_1$ définie par $t_1=t_3=0$ et
  du plan $X_2$ défini par $t_2-t_3=0$~», ce qui décrit en termes géométriques
  le fait que $X_0$ est égal à $\operatorname{Proj}(C/\mathfrak a)$, où l'idéal
  gradué $\mathfrak a$ est égal à $\mathfrak b\mathfrak c$, avec
  $\mathfrak b=Ct_1+Ct_3$, $\mathfrak c=C(t_2-t_3)$~; considérons le morphisme
  $f_0:X_0\to X_1$ qui, en termes géométriques, est la «~projection de $X_0$
  sur $X_1$ à partir de la droite à l'infini $D$ définie par $t_0=t_2=0$~»~;
  sous forme algébrique, $f_0$ correspond à l'homomorphisme d'algèbres graduées
  \[
    k[t_0,t_2]\longrightarrow k[t_0,t_1,t_2,t_3]/\mathfrak a
  \]
  obtenu par passage au quotient à partir de l'injection canonique
  $k[t_0,t_2]\to k[t_0,t_1,t_2,t_3]$. Il est clair que $f_0$ est un morphisme
  projectif, donc \emph{propre}~; il en est donc de même de sa restriction
  $f:X\to Y$, où $Y=D_+(t_0)$ et $X=f_0^{-1}(Y)$~; pour voir que $f$ est
  \emph{plat}, il suffit de remarquer que $k[t_2]$ est un anneau principal et
  l'anneau
  \[
    C'=k[t_1,t_2,t_3]/(t_2-t_3)((t_1)+(t_3))
  \]
  un $k[t_2]$-module sans torsion ($\mathbf 0_{\mathrm I}$, 6.3.4). Pour tout
  $y\in Y$ distinct du point $y_0$ défini par $t_2=0$, $f^{-1}(y)$ a alors
  deux composantes irréductibles, de dimensions 0 et 1, tandis que
  $f^{-1}(y_0)$ n'a qu'une seule composante irréductible de dimension 1.

  \item Dans (12.2.1, (i)), on ne peut pas non plus remplacer la condition que
  l'ensemble $E'(y)\supset E(y)$ des dimensions des cycles premiers associés à
  $\mathcal F_y$ soit contenu dans $\Phi$ par la condition qu'il soit
  \emph{égal à} $\Phi$. Soient en effet $k$ un corps, $A$ l'anneau $k[t]$ de
  polynômes, $B$ le sous-anneau $k[t^4,t^3u,tu^3,u^4]$ de l'anneau de polynômes
  $k[t,u]$ ($t,u$ indéterminées)~; $B$ n'est pas intégralement clos, l'élément
  $t^2u^2$ appartenant à sa clôture intégrale mais n'étant pas dans $B$~; si
  $\mathfrak m$ est l'idéal maximal de $B$, engendré par $t^4,t^3u,tu^3$ et
  $u^4$, $\mathfrak m$ est idéal premier associé immergé de $A/At^4$. Posons
  $Y=\operatorname{Spec}(A)$, $X=\operatorname{Spec}(B)$, et soit $f:X\to Y$
  le morphisme correspondant à l'homomorphisme $A\to B$ qui transforme $t$ en
  $t^4$. Comme cet homomorphisme fait de $B$ un $A$-module sans torsion, $f$
  est plat ($\mathbf 0_{\mathrm I}$, 6.3.4). Si $y$ est le point de $Y$
  correspondant à l'idéal maximal $tA$, le préschéma $f^{-1}(y)$ est
  irréductible et de dimension 1, mais admet un cycle premier associé immergé
  de dimension 0~; au contraire, si $y'$ est le point générique de $Y$,
  $f^{-1}(y')$ n'a pas de cycle premier associé immergé, étant le spectre d'une
  $k(y')$-algèbre intègre de dimension 1. Dans cet exemple, $f$ est un morphisme
  affine non propre~; on peut en déduire aussitôt un exemple analogue où $f$
  est propre et plat en considérant $X$ comme un ouvert partout dense d'un
  schéma projectif sur $Y$ (\textbf{II}, 5.3.4 et 5.3.2), ou en procédant
  directement comme dans la Remarque (ii).

  Les Remarques (ii) et (iii) montrent la nécessité d'inclure dans (12.2.1,
  (iv)) la condition que $\mathcal F_y$ n'ait pas de cycle premier associé
  immergé.

  \item Dans (12.2.1, (xi)), on ne peut supprimer l'hypothèse que $f$ soit
  propre. En effet, soient $k$ un corps, $Y$ la «~droite affine~», spectre de
  l'anneau de polynômes $k[t]$ ($t$ indéterminée), $X$ le schéma \emph{somme}
  de $Y$ et de l'ouvert complémentaire du point fermé $y$ de $Y$ défini par
  $t=0$. Il est clair que le morphisme $f:X\to Y$ qui est égal aux injections
  canoniques dans chacune des composantes de $X$ est plat et que si l'on prend
  $\mathcal F=\mathcal O_X$, $\mathcal F_z$ n'a de cycle premier associé
  immergé pour aucun $z\in Y$. Toutefois, on voit aussitôt que l'on a $l(y)=1$
  et $l(z)=2$ pour tout $z\neq y$.

  \item Dans (12.2.1, (xi)), on ne peut pas non plus supprimer l'hypothèse que
  $\mathcal F_y$ soit sans cycle premier associé immergé. C'est ce que montre
  l'exemple de la Remarque (ii)~: on voit aussitôt en effet que l'on a, avec
  $\mathcal F=\mathcal O_X$, $l(y_0)=1$ et $l(y)=2$ pour $y\neq y_0$ dans $Y$.

  \item Nous ignorons si dans (12.2.1, (xi)) on peut remplacer l'hypothèse que
  $\mathcal F$ est $f$-plat par celle que $\operatorname{Supp}(\mathcal F)=X$
  et que $f$ est universellement ouvert. En reprenant la démonstration de
  (12.2.1.4), on voit (en utilisant aussi (14.3.6)) qu'il faudrait résoudre la
  question suivante~: Soient $A$ un anneau local noethérien, $M$ un $A$-module
  \oldpage[IV-3]{183}
  de type fini, $t$ un élément de l'idéal maximal $\mathfrak m$ de $A$, qui
  n'est contenu dans aucun idéal premier minimal de $A$~; s'il existe un de ces
  idéaux premiers minimaux $\mathfrak p$ tels que $\dim(A/\mathfrak p)=1$,
  est-il vrai que $\mathfrak m$ est un idéal premier associé à $M/tM$ ($t$
  n'étant pas supposé $M$-régulier)~?

  On notera toutefois qu'on ne peut, dans (12.2.1, (xi)), supprimer purement et
  simplement l'hypothèse que $\mathcal F$ est $f$-plat. On prendra en effet ici
  $P$ comme dans la Remarque (ii), et dans $P$ le sous-préschéma fermé $X_0$
  réunion des trois droites $X_1,X_2,X_3$ définies respectivement par
  $t_1=t_3=0$, $t_1-t_0=t_3=0$, $t_2=t_3=0$~; on définit la projection $f_0$ de
  $X_0$ sur $X_1$ comme précédemment et sa restriction $f:X\to Y$, où
  $Y=D_+(t_0)$ est la droite affine et $X=f_0^{-1}(Y)$~; $f$ est propre mais
  non plat, et si l'on prend $\mathcal F=\mathcal O_X$, $\mathcal F_y$ n'a de
  cycles premiers associés immergés pour aucun $y\in Y$. Mais on a $l(y_0)=1$
  et $l(y)=2$ pour $y\neq y_0$ dans $Y$.
\end{enumerate}
\end{remarks}

\begin{theorem}[12.2.4]
\label{IV.12.2.4-fr}
Soient $f:X\to Y$ un morphisme propre, plat et de présentation finie, $k$ un
entier $\geq 1$. Les parties suivantes de $Y$ sont ouvertes~:
\begin{enumerate}[label=(\roman*)]
  \item L'ensemble des $y\in Y$ tels que $X_y$ possède la propriété $(S_k)$.
  \item L'ensemble des $y\in Y$ tels que $X_y$ vérifie la propriété $(R_k)$
  géométrique, soit équidimensionnel en chaque point et n'ait pas de cycle
  premier associé immergé.
  \item L'ensemble des $y\in Y$ tels que $X_y$ soit géométriquement régulier
  (i.e. lisse sur $k(y)$).
  \item L'ensemble des $y\in Y$ tels que $X_y$ soit géométriquement normal.
  \item L'ensemble des $y\in Y$ tels que $X_y$ soit géométriquement réduit.
  \item L'ensemble des $y\in Y$ tels que $X_y$ soit géométriquement réduit et
  que le nombre géométrique de composantes connexes de $X_y$ soit égal à $k$.
  \item L'ensemble des $y\in Y$ tels que $X_y$ soit géométriquement
  ponctuellement intègre (4.6.9).
  \item L'ensemble des $y\in Y$ tels que $X_y$ soit géométriquement intègre.
  \item L'ensemble des $y\in Y$ tels que $X_y$ n'ait pas de cycle premier
  associé immergé et que la multiplicité totale (4.7.4) de $X_y$ sur $k(y)$
  soit $\leq k$.
\end{enumerate}
\end{theorem}

\begin{proof}
Sauf pour (vi), ces assertions sont des cas particuliers d'assertions de
(12.2.1), ou se déduisent d'assertions de (12.1.6) comme au début de la
démonstration de (12.2.1). Pour (vi), on se ramène comme toujours (compte tenu
de l'invariance du nombre géométrique de composantes connexes par changement
de base (4.5.6)) au cas où $Y$ est noethérien. L'ensemble $U$ des $y\in Y$ tels
que $X_y$ soit géométriquement réduit est ouvert dans $Y$ en vertu de (v)~; il
résulte alors de (\textbf{III}, 7.8.7 et 7.8.6) que pour tout $y\in U$, il y a
un voisinage $V\subset U$ de $y$ et un entier $m$ tels que, pour \emph{tout}
$z\in V$, $\Gamma(X_z,\mathcal O_{X_z})$ soit isomorphe à $(k(z))^m$~; mais en
vertu de (\textbf{III}, 4.3.4), $m$ est alors le nombre géométrique de
composantes connexes de $X_z$, d'où la conclusion. On donnera une autre
démonstration de (vi) dans (15.5.9).
\end{proof}

\subsection{Propriétés cohomologiques locales des fibres d'un morphisme plat et
localement de présentation finie}
\label{subsection:IV.12.3-fr}

\begin{lemma}[12.3.1]
\label{IV.12.3.1-fr}
Soient $A$ un anneau, $B$ une $A$-algèbre de présentation finie qui soit un
$A$-module plat, $M$ un $B$-module de présentation finie, qui soit un
$A$-module plat. Alors les conditions suivantes sont équivalentes~:
\begin{enumerate}[label=\emph{\alph*)}]
  \item $M$ est un $B$-module projectif.
  \item $M$ est un $B$-module plat.
  \item Pour tout $y\in\operatorname{Spec}(A)$, $M\otimes_Ak(y)$ est un
  $(B\otimes_Ak(y))$-module projectif.
\end{enumerate}
\end{lemma}

\begin{proof}
\oldpage[IV-3]{184}
L'équivalence de $a)$ et $b)$ résulte de Bourbaki, \emph{Alg. comm.}, chap. II,
\S~5, n\textsuperscript{o} 2, cor. 2 du th. 1. Comme $a)$ implique $c)$
trivialement, il reste à prouver que $c)$ implique $b)$, ce qui résulte du
critère de platitude par fibres (11.3.10), appliqué avec $g=h$, $f=1_X$.
\end{proof}

\begin{proposition}[12.3.2]
\label{IV.12.3.2-fr}
Soient $A$ un anneau, $B$ une $A$-algèbre de présentation finie qui soit un
$A$-module plat, $M$ un $B$-module de présentation finie qui soit un
$A$-module plat. Alors~:
\begin{enumerate}[label=(\roman*)]
  \item Il existe une résolution gauche de $M$ par des $B$-modules libres de
  type fini.
  \item On a
  \begin{equation}
    \tag{12.3.2.1}
    \label{IV.12.3.2.1-fr}
    \operatorname{dim.proj}_B(M)
    =\sup_{y\in\operatorname{Spec}(A)}
    \operatorname{dim.proj}_{B\otimes_Ak(y)}(M\otimes_Ak(y))
    =\operatorname{Tor.dim}_B(M)
  \end{equation}
  où $\operatorname{Tor.dim}_B(M)$ est le \emph{plus petit entier} $i$ tel que
  $\operatorname{Tor}^B_j(M,N)=0$ pour tout $j\geq i$ et tout $B$-module $N$
  (et $+\infty$ si un tel entier $i$ n'existe pas).
\end{enumerate}
\end{proposition}

\begin{proof}
\begin{enumerate}[label=(\roman*)]
  \item En vertu de (8.9.1), il existe un sous-anneau noethérien $A_0$ de $A$,
  une $A_0$-algèbre de type fini $B_0$ et un $B_0$-module de type fini $M_0$
  tels que $B$ soit isomorphe à $B_0\otimes_{A_0}A$ et $M$ à
  $M_0\otimes_{A_0}A$. De plus (11.2.7) on peut supposer que $B_0$ et $M_0$
  sont des $A_0$-modules plats. Il y a alors une résolution gauche
  $(L_0)_\bullet$ de $M_0$ formée de $B_0$-modules libres de type fini, et
  comme $M_0$ et les $L_{0i}$ sont des $A_0$-modules plats,
  $(L_0)_\bullet\otimes_{A_0}A$ est une résolution gauche de $M$ formée de
  $B$-modules libres de type fini (2.1.10).

  \item En vertu de ($\mathbf 0$, 17.2.2, (ii)), on a
  $\operatorname{Tor.dim}_B(M)\leq\operatorname{dim.proj}_B(M)$, et la
  définition de la dimension projective prouve aussitôt que, pour tout
  $y\in\operatorname{Spec}(A)$, on a
  \[
    \operatorname{dim.proj}_{B\otimes_Ak(y)}(M\otimes_Ak(y))
    \leq\operatorname{dim.proj}_B(M).
  \]
  Pour prouver les inégalités inverses, considérons une résolution gauche
  $(L_i)$ de $M$ par des $B$-modules libres de type fini, et supposons que
  $\operatorname{Tor.dim}_B(M)=n$ (resp.
  $\operatorname{dim.proj}_{B\otimes_Ak(y)}(M\otimes_Ak(y))\leq n$ pour tout
  $y\in\operatorname{Spec}(A)$). Alors
  $R=\operatorname{Im}(L_n\to L_{n-1})=\operatorname{Ker}(L_{n-1}\to L_{n-2})$
  est un $B$-module de type fini et aussi un $A$-module plat, en vertu de
  l'hypothèse sur $M$ et $B$ et de (2.1.10). En outre, on a
  $\operatorname{Tor}^B_{n+1}(M,N)=\operatorname{Tor}^B_1(R,N)$ pour tout
  $B$-module $N$ (M, V, 7). L'hypothèse $\operatorname{Tor.dim}_B(M)=n$
  entraîne donc $\operatorname{Tor}^B_1(R,N)=0$ pour tout $B$-module $N$,
  c'est-à-dire que $R$ est un $B$-module plat, donc projectif en vertu de
  (12.3.1)~; cela établit que $\operatorname{dim.proj}_B(M)\leq n$.
  L'hypothèse
  $\operatorname{dim.proj}_{B\otimes_Ak(y)}(M\otimes_Ak(y))\leq n$ pour tout
  $y\in\operatorname{Spec}(A)$ entraîne d'autre part, par tensorisation avec
  $k(y)$, que dans chacune des suites (exactes en vertu de la platitude sur $A$
  de $M$, des $L_i$ et de $R$ (2.1.10))
  \[
    0\to R\otimes_Ak(y)\to L_{n-1}\otimes_Ak(y)\to\cdots
    \to L_0\otimes_Ak(y)\to M\otimes_Ak(y)\to0
  \]
  $R\otimes_Ak(y)$ est un $(B\otimes_Ak(y))$-module projectif (pour
  $y\in\operatorname{Spec}(A)$). On conclut alors encore de (12.3.1) que $R$
  est un $B$-module projectif, donc $\operatorname{dim.proj}_B(M)\leq n$.
\end{enumerate}
\end{proof}

\begin{proposition}[12.3.3]
\label{IV.12.3.3-fr}
Soient $f:X\to S$ un morphisme localement de présentation finie,
$\mathcal L_\bullet$ un complexe formé de $\mathcal O_X$-Modules
quasi-cohérents de présentation finie~; pour tout $s\in S$, soit
$(\mathcal L_\bullet)_s$ le complexe
$\mathcal L_\bullet\otimes_{\mathcal O_S}k(s)$ de $\mathcal O_{X_s}$-Modules de
type fini. Supposons que $\mathcal L_{n-1}$ soit $f$-plat. Alors l'ensemble $U$
des $x\in X$ tels que $(\mathcal H_n((\mathcal L_\bullet)_{f(x)}))_x=0$ est
ouvert dans $X$. Si de plus $\mathcal L_n$ est $f$-plat, alors on a
$\mathcal H_n(\mathcal L_\bullet)|U=0$.
\end{proposition}

\begin{proof}
On peut évidemment se limiter au cas où $n=1$ et où le complexe
$\mathcal L_\bullet$ se réduit à
$0\to\mathcal L_2\xrightarrow{u}\mathcal L_1\xrightarrow{v}\mathcal L_0\to0$.
On peut d'abord se ramener au cas où $S=\operatorname{Spec}(A)$ et $X$
\oldpage[IV-3]{185}
sont affines, puis au cas où $S$ est \emph{noethérien}~; en effet, par (8.9.1)
et (8.5.2), on sait qu'il existe un préschéma noethérien
$S'=\operatorname{Spec}(A')$, où $A'$ est un sous-anneau de $A$, un morphisme
de type fini $f':X'\to S'$ et trois $\mathcal O_{X'}$-Modules cohérents
$\mathcal L'_i$ ($0\leq i\leq2$) tels que $X=X'\times_{S'}S$,
$f=f'_{(S)}$, $\mathcal L_i=\mathcal L'_i\otimes_{S'}S$, ainsi que deux
homomorphismes $u'$, $v'$ tels que $u=u'\otimes1$, $v=v'\otimes1$ et
$v'\circ u'=0$. On peut en outre supposer que $\mathcal L'_0$ est $f'$-plat
(11.2.7), et que $\mathcal L'_1$ est $f'$-plat lorsque $\mathcal L_1$ est
supposé $f$-plat. Par fidèle platitude, l'hypothèse
$(\mathcal H_n((\mathcal L_\bullet)_{f(x)}))_x=0$ équivaut à
$(\mathcal H_n((\mathcal L'_\bullet)_{f'(x')}))_{x'}=0$, si $x'$ est la
projection de $x$ dans $X'$~; si $U'$ est l'ensemble des $x'\in X'$ tels que
$(\mathcal H_n((\mathcal L'_\bullet)_{f'(x')}))_{x'}=0$, $U$ est donc l'image
réciproque de $U'$ par la projection $X\to X'$, ce qui ramène, pour la première
assertion, au cas où $S$ est noethérien. Pour la seconde assertion, on remarque
en outre que $A$ est limite inductive de ses sous-anneaux $A_\lambda$ qui sont
des $A'$-algèbres de type fini~; posons
$X_\lambda=X'\times_{S'}S_\lambda$, où
$S_\lambda=\operatorname{Spec}(A_\lambda)$,
$\mathcal L_i^{(\lambda)}=\mathcal L'_i\otimes_{S'}S_\lambda$, et soient
$u_\lambda=u'\otimes1:\mathcal L_2^{(\lambda)}\to\mathcal L_1^{(\lambda)}$,
$v_\lambda=v'\otimes1:\mathcal L_1^{(\lambda)}\to\mathcal L_0^{(\lambda)}$,
de sorte que l'on a
$\mathcal H_1(\mathcal L_\bullet^{(\lambda)})
=\operatorname{Ker}(v_\lambda)/\operatorname{Im}(u_\lambda)$. Or, comme le
foncteur $\varinjlim$ est exact dans la catégorie des groupes commutatifs, on a
$\operatorname{Ker}(v)=\varinjlim\operatorname{Ker}(v_\lambda)$,
$\operatorname{Im}(u)=\varinjlim\operatorname{Im}(u_\lambda)$, et
$\mathcal H_1(\mathcal L_\bullet)
=\varinjlim\mathcal H_1(\mathcal L_\bullet^{(\lambda)})$. Si l'on a supposé
que $\mathcal L_1$ est $f$-plat et (en se ramenant au cas où $X=U$) que l'on
ait prouvé $\mathcal H_1(\mathcal L_\bullet^{(\lambda)})=0$ pour tout
$\lambda$, on en déduira bien $\mathcal H_1(\mathcal L_\bullet)=0$.

\medskip
I) Supposons désormais $S$ noethérien. On sait (sans hypothèse de platitude sur
$\mathcal L_0$) que l'ensemble $U$ est \emph{constructible} dans $X$ (9.9.6).
Utilisant maintenant ($\mathbf 0_{\mathrm{III}}$, 9.2.5), il reste à montrer
que pour toute \emph{générisation} $x'$ de $x\in U$ dans $X$, on a aussi
$x'\in U$. La méthode exposée dans (12.0.2) s'applique sans changement (compte
tenu du fait que pour un préschéma $Z$ sur un corps $k$ et une extension $k'$
de $k$, la projection $Z\otimes_k k'\to Z$ est fidèlement plate). On peut donc
supposer que $S$ est le spectre d'un anneau de valuation discrète $A$, dont on
désigne par $t$ une uniformisante, $x$ étant au-dessus du point fermé de $S$ et
$x'$ au-dessus du point générique de $S$. L'hypothèse que $\mathcal L_0$ est
$f$-plat entraîne que $t$ est $\mathcal L_0$-régulier
($\mathbf 0_{\mathrm I}$, 6.3.4). On est alors ramené à prouver le lemme suivant
(où $B$, $M$, $N$, $P$ seront remplacés par $\mathcal O_x$,
$(\mathcal L_2)_x$, $(\mathcal L_1)_x$ et $(\mathcal L_0)_x$ respectivement)~:

\begin{lemma}[12.3.3.1]
\label{IV.12.3.3.1-fr}
Soient $B$ un anneau local noethérien, $M,N,P$ trois $B$-modules, $N$ étant
supposé de type fini, $M\xrightarrow{u}N\xrightarrow{v}P$ deux homomorphismes
tels que $v\circ u=0$. Soit $t$ un élément de l'idéal maximal de $B$ tel que
$t$ soit $P$-régulier et que la suite
\begin{equation}
  \tag{12.3.3.2}
  \label{IV.12.3.3.2-fr}
  M/tM\xrightarrow{u\otimes1_{B/tB}}N/tN
  \xrightarrow{v\otimes1_{B/tB}}P/tP
\end{equation}
soit exacte. Alors la suite $M\xrightarrow{u}N\xrightarrow{v}P$ est exacte.
\end{lemma}

Notons d'abord que si l'on remplace $M$ par son image $Q$ dans $N$ et $u$ par
l'injection $j:Q\to N$, l'image de $j\otimes1$ est la même que celle de
$u\otimes1$, donc on peut, sans changer ni l'hypothèse ni la conclusion,
supposer $u$ injectif. D'autre part, si $R$ est l'image de $v$ et
$\rho:N\to R$ la surjection canonique, $t$ est évidemment $R$-régulier, on a
$\rho\circ u=0$, et le noyau de $\rho\otimes1$ est contenu dans celui de
$v\otimes1$~; il lui est par suite égal si la suite (12.3.3.2) est exacte~;
comme d'autre part on a évidemment $\operatorname{Ker}(\rho)
=\operatorname{Ker}(v)$,
\oldpage[IV-3]{186}
on voit qu'on peut, pour prouver le lemme, supposer en
outre $v$ surjectif. Le lemme sera alors conséquence du suivant~:

\begin{lemma}[12.3.3.3]
\label{IV.12.3.3.3-fr}
Soient $B$ un anneau, $M,N,P$ trois $B$-modules,
$u:M\to N$, $v:N\to P$ deux homomorphismes tels que $u$ soit injectif, $v$
surjectif et $v\circ u=0$. Soit $t$ un élément de $B$ tel que $t$ soit
$P$-régulier. Alors on a
\begin{equation}
  \tag{12.3.3.4}
  \label{IV.12.3.3.4-fr}
  \operatorname{Ker}(v\otimes1)/\operatorname{Im}(u\otimes1)
  =(\operatorname{Ker}(v)/\operatorname{Im}(u))\otimes_B(B/tB)
\end{equation}
à un isomorphisme canonique près.
\end{lemma}

En effet, l'hypothèse que la suite (12.3.3.2) est exacte entraînera alors
$(\operatorname{Ker}v/\operatorname{Im}u)\otimes_B(B/tB)=0$, et comme $t$
appartient à l'idéal maximal de $B$ et que $\operatorname{Ker}(v)$ est un
$B$-module de type fini (puisque $B$ est supposé noethérien dans (12.3.3.1)),
le lemme de Nakayama prouvera que $\operatorname{Ker}(v)=\operatorname{Im}(u)$.

Reste donc à démontrer (12.3.3.3). Posons $u'=u\otimes1$, $v'=v\otimes1$,
$Z=\operatorname{Ker}(v)$, $H=Z/\operatorname{Im}(u)$, de sorte qu'on a les
suites exactes
\begin{gather*}
  0\to M\to Z\to H\to0,\\
  0\to Z\to N\xrightarrow{v}P\to0,
\end{gather*}
d'où, en tensorisant par $B/tB$ et utilisant le lemme (3.4.1.4) et le fait que
$t$ est $P$-régulier, les suites exactes
\begin{gather*}
  M/tM\xrightarrow{w'}Z/tZ\to H/tH\to0,\\
  0\to Z/tZ\to N/tN\xrightarrow{v'}P/tP\to0,
\end{gather*}
d'où $\operatorname{Ker}(v')=Z/tZ$, et comme
$\operatorname{Im}(w')=\operatorname{Im}(u')$, on obtient (12.3.3.4).

\medskip
II) Supposons maintenant en outre que $\mathcal L_1$ soit $f$-plat~;
remplaçant en outre éventuellement $X$ par $U$, on peut supposer que $X=U$.
Il s'agit de voir que pour tout $x\in X$, on a
$(\mathcal H_1(\mathcal L_\bullet))_x=0$. Soit $s=f(x)$, et soit
$\mathfrak n$ l'idéal $\mathfrak m_s\mathcal O_{X,x}$ qui est contenu dans
l'idéal maximal de $\mathcal O_{X,x}$~; comme
$(\mathcal H_1(\mathcal L_\bullet))_x$ est un $\mathcal O_{X,x}$-module de type
fini ($X$ étant localement noethérien), il est séparé pour la topologie
$\mathfrak n$-adique ($\mathbf 0_{\mathrm I}$, 7.3.5), donc il suffit de
montrer que son \emph{séparé complété} pour cette topologie est 0. Or, en vertu
de (\textbf{III}, 7.4.7.2), ce séparé complété est égal à
$\varprojlim_k\mathcal H_1(\mathcal L_\bullet\otimes_{\mathcal O_S}
(\mathcal O_{S,s}/\mathfrak m_s^{k+1}))$, et il suffira donc de prouver que
chacun des termes de ce système projectif est nul. Mais cela est vrai par
hypothèse pour $k=0$~; raisonnons donc par récurrence sur $k$. La conclusion
résultera du lemme plus général suivant (qu'on appliquera pour
$A=\mathcal O_{S,s}/\mathfrak m_s^{k+1}$ et $\mathfrak J$ égal à l'idéal
maximal $\mathfrak m_s/\mathfrak m_s^{k+1}$ de $A$)~:

\begin{lemma}[12.3.3.5]
\label{IV.12.3.3.5-fr}
Soient $A$ un anneau, $\mathfrak J$ un idéal nilpotent de $A$,
$L_\bullet:P\xrightarrow{u}Q\xrightarrow{v}R$ un complexe de $A$-modules. On
pose $A_0=A/\mathfrak J$,
$L_\bullet^{(0)}=L_\bullet\otimes_AA_0:
P_0\xrightarrow{u_0}Q_0\xrightarrow{v_0}R_0$, et l'on suppose que
$\operatorname{gr}_{\mathfrak J}^{\bullet}(A)$ est un $A_0$-module plat, et
que $Q$ et $R$ sont des $A$-modules plats. Alors la relation
$\mathrm H_1(L_\bullet^{(0)})=0$ entraîne $\mathrm H_1(L_\bullet)=0$.
\end{lemma}

Posons
$L_\bullet^{(k)}=L_\bullet\otimes_A(A/\mathfrak J^{k+1}):
P_k\xrightarrow{u_k}Q_k\xrightarrow{v_k}R_k$~; comme il existe un entier
$k\geq1$ tel que $L_\bullet^{(k)}=L_\bullet$, il suffira de prouver, par
récurrence sur $k$, que la suite
$P_k\xrightarrow{u_k}Q_k\xrightarrow{v_k}R_k$ est
\oldpage[IV-3]{187}
exacte (cette assertion découlant de l'hypothèse pour $k=0$). Soit donc
$x\in Q_k$ un élément tel que $v_k(x)=0$. Comme par l'hypothèse de récurrence
la suite
$P_{k-1}\xrightarrow{u_{k-1}}Q_{k-1}\xrightarrow{v_{k-1}}R_{k-1}$ est exacte,
l'image canonique de $x$ dans
$Q_{k-1}=Q_k/(\mathfrak J^kQ/\mathfrak J^{k+1}Q)$ appartient à
$\operatorname{Im}(u_{k-1})$, donc il existe $z\in P_k$ tel que
$x=u_k(z)+x'$ avec $x'\in\mathfrak J^kQ/\mathfrak J^{k+1}Q$. Comme on a
$v_k\circ u_k=0$, la relation $v_k(x)=0$ entraîne $v_k(x')=0$, et par ailleurs
on a évidemment $v_k(x')\in\mathfrak J^kR/\mathfrak J^{k+1}R$~; tout revient
donc à prouver que la suite
\[
  \mathfrak J^kP/\mathfrak J^{k+1}P\xrightarrow{u'}
  \mathfrak J^kQ/\mathfrak J^{k+1}Q\xrightarrow{v'}
  \mathfrak J^kR/\mathfrak J^{k+1}R
\]
(où $u'$ et $v'$ proviennent de $u$ et $v$ par restriction et passage aux
quotients) est exacte. Or, par hypothèse,
$\mathfrak J^k/\mathfrak J^{k+1}$ est un $(A/\mathfrak J)$-module plat, donc la
suite
\[
  (P/\mathfrak JP)\otimes_{A/\mathfrak J}
  (\mathfrak J^k/\mathfrak J^{k+1})\xrightarrow{u''}
  (Q/\mathfrak JQ)\otimes_{A/\mathfrak J}
  (\mathfrak J^k/\mathfrak J^{k+1})\xrightarrow{v''}
  (R/\mathfrak JR)\otimes_{A/\mathfrak J}
  (\mathfrak J^k/\mathfrak J^{k+1})
\]
est exacte. Mais
$(Q/\mathfrak JQ)\otimes_{A/\mathfrak J}
(\mathfrak J^k/\mathfrak J^{k+1})
=Q\otimes_A(\mathfrak J^k/\mathfrak J^{k+1})$ s'identifie à
$\mathfrak J^kQ/\mathfrak J^{k+1}Q$ en vertu de la platitude de $Q$ sur $A$,
et de même
$(R/\mathfrak JR)\otimes_{A/\mathfrak J}
(\mathfrak J^k/\mathfrak J^{k+1})$ s'identifie à
$\mathfrak J^kR/\mathfrak J^{k+1}R$. Enfin, l'image de $u'$ s'identifie à
celle de $u''$, et cela termine la démonstration.
\end{proof}

\begin{corollary}[12.3.4]
\label{IV.12.3.4-fr}
Soient $f:X\to S$ un morphisme plat localement de présentation finie,
$\mathcal F$, $\mathcal G$ deux $\mathcal O_X$-Modules de présentation finie et
$f$-plats. Soient $n$ un entier $\geq0$, $U$ l'ensemble des $x\in X$ tels que
\[
  \bigl(\mathscr{E}\!xt^n_{\mathcal O_{X_{f(x)}}}
  (\mathcal F_{f(x)},\mathcal G_{f(x)})\bigr)_x=0
\]
(resp.
$\bigl(\mathscr{T}\!or_n^{\mathcal O_{X_{f(x)}}}
(\mathcal F_{f(x)},\mathcal G_{f(x)})\bigr)_x=0$). Alors $U$ est ouvert, et
l'on a
\[
  \mathscr{E}\!xt^n_{\mathcal O_X}(\mathcal F,\mathcal G)|U=0
\]
(resp.
$\mathscr{T}\!or_n^{\mathcal O_X}(\mathcal F,\mathcal G)|U=0$).
\end{corollary}

\begin{proof}
On peut évidemment se borner au cas où $S$ et $X$ sont affines. Alors (12.3.2),
il existe une résolution gauche $\mathcal L_\bullet=(\mathcal L_i)$ de
$\mathcal F$ formée de $\mathcal O_X$-Modules \emph{libres de type fini}. Si
l'on pose $\mathcal M^\bullet=\mathscr{H}\!om(\mathcal L_\bullet,\mathcal G)$
(resp. $\mathcal N_\bullet=\mathcal L_\bullet\otimes_{\mathcal O_X}\mathcal G$),
on en déduit que chacun des $\mathcal M^i$ (resp. $\mathcal N_i$) est isomorphe
à un $\mathcal O_X$-Module de la forme $\mathcal G^m$, donc les $\mathcal M^i$
et $\mathcal N_i$ sont des $\mathcal O_X$-Modules de présentation finie et
$f$-plats. En outre, pour tout changement de base $S'\to S$, si l'on pose
$X'=X\times_SS'$, $\mathcal F'=\mathcal F\otimes_SS'$,
$\mathcal G'=\mathcal G\otimes_SS'$, $\mathcal L'_\bullet
=\mathcal L_\bullet\otimes_SS'$, $\mathcal L'_\bullet$ est encore une
résolution gauche de $\mathcal F'$ par des $\mathcal O_{X'}$-Modules libres de
type fini (2.1.10), et $\mathcal M'^\bullet=\mathcal M^\bullet\otimes_SS'$ est
égal à $\mathscr{H}\!om(\mathcal L'_\bullet,\mathcal G')$ (resp.
$\mathcal N'_\bullet=\mathcal N_\bullet\otimes_SS'$ est égal à
$\mathcal L'_\bullet\otimes_{\mathcal O_{X'}}\mathcal G'$), d'après ce qui
précède. En particulier, on a, pour tout $s\in S$,
\[
  \mathscr{E}\!xt^n_{\mathcal O_{X_s}}(\mathcal F_s,\mathcal G_s)
  =\mathcal H^n((\mathcal M^\bullet)_s)
\]
(resp.
$\mathscr{T}\!or_n^{\mathcal O_{X_s}}(\mathcal F_s,\mathcal G_s)
=\mathcal H_n((\mathcal N_\bullet)_s)$). Appliquant (12.3.3) aux complexes de
Modules $f$-plats $\mathcal M^\bullet$ et $\mathcal N_\bullet$, on en déduit
aussitôt le corollaire.
\end{proof}

% IV3_B26_P187_END_BEFORE_SECTION_13

% FIN SOURCE ega4/ega4-12-fr.tex

%
% DEBUT SOURCE ega4/ega4-13-fr.tex
\section{Morphismes équidimensionnels}
\label{section:IV.13-fr}

Ce paragraphe est consacré à l'étude de la variation de la \emph{dimension}
des fibres d'un morphisme localement de type fini $f:X\to Y$ (qui est déjà
intervenue à propos de la «~formule des dimensions~» dans 5.5 et 5.6). On
prouve d'abord (13.1.3) que la fonction
$x\mapsto\dim_x(f^{-1}(f(x)))$ est toujours \emph{semi-continue
supérieurement} dans $X$ (théorème de semi-continuité de Chevalley). On étudie
ensuite plus particulièrement les morphismes, dits «~équidimensionnels~», pour
lesquels cette fonction est \emph{localement constante}.

\oldpage[IV-3]{188}
Malheureusement la notion de morphisme équidimensionnel n'est pas stable par
changement de base~; c'est pourquoi dans de nombreuses questions il est plus
commode de travailler avec la notion de morphisme universellement ouvert, dont
l'étude fait l'objet des \S\S~14 et 15.

\subsection{Le théorème de semi-continuité de Chevalley}
\label{subsection:IV.13.1-fr}

\begin{lemma}[13.1.1]
\label{IV.13.1.1-fr}
Soient $Y$ un préschéma irréductible localement noethérien, $X$ un préschéma
irréductible, $f:X\to Y$ un morphisme dominant de type fini. Soit $\xi$
(resp. $\eta$) le point générique de $X$ (resp. $Y$) et soit
$e=\dim(f^{-1}(\eta))$. Alors, pour tout $x\in X$, toutes les composantes
irréductibles de $f^{-1}(f(x))$ sont de dimension $\geq e$.
\end{lemma}

\begin{proof}
La proposition est immédiate lorsque, pour tout $y\in f(X)$,
$\mathcal O_y$ est un \emph{anneau universellement caténaire}~: en effet, si
$x$ est point générique d'une composante irréductible $Z$ de $f^{-1}(y)$, il
résulte de (5.6.5), joint à (5.2.1), que l'on a
$e+\dim(\mathcal O_y)=\dim(Z)+\dim(\mathcal O_x)$~; mais en vertu de
($\mathbf 0$, 16.3.9), on a $\dim(\mathcal O_x)\leq\dim(\mathcal O_y)$,
d'où la conclusion dans ce cas.

Nous allons ramener le cas général à ce cas particulier. La question est
évidemment locale sur $Y$, et, compte tenu de (4.1.1.3), elle est aussi locale
sur $X$~; on peut donc se borner au cas où $Y=\operatorname{Spec}(A)$ et
$X=\operatorname{Spec}(B)$ sont affines et irréductibles, $B$ étant une
$A$-algèbre de type fini. En outre (1.5.4), on peut supposer $X$ et $Y$
réduits, donc $A$ et $B$ \emph{intègres} et, comme $f$ est dominant, $A$ est
alors un \emph{sous-anneau} de $B$. Considérons $A$ comme limite inductive de
ses sous-$\mathbf Z$-algèbres de type fini~; il résulte alors de (8.9.1) qu'il
existe une telle sous-algèbre $A_0$ et une $A_0$-algèbre de type fini $B_0$
telles que $B=B_0\otimes_{A_0}A$. Posons $Y_0=\operatorname{Spec}(A_0)$,
$X_0=\operatorname{Spec}(B_0)$, et soit $f_0:X_0\to Y_0$ le morphisme
correspondant à l'homomorphisme $A_0\to B_0$, de sorte que
$f=f_0\times 1_Y$. Il n'est pas évident \emph{a priori} que le préschéma $X_0$
soit intègre, mais nous allons voir qu'on peut se ramener à ce cas. Soit
$\eta_0$ le point générique de $Y_0$, de sorte que si $g$ est le morphisme
$Y\to Y_0$, on a $g(\eta)=\eta_0$~; par transitivité des fibres
(\textbf{I}, 3.6.4), on a
$f^{-1}(\eta)=f_0^{-1}(\eta_0)\otimes_{k(\eta_0)}k(\eta)$, et comme
$f^{-1}(\eta)$ est irréductible par hypothèse, il en est de même de
$f_0^{-1}(\eta_0)$ (4.4.1). Notre assertion résultera alors du lemme suivant~:

\begin{lemma}[13.1.2]
\label{IV.13.1.2-fr}
Soient $Y_0$, $Y$ deux préschémas intègres de points génériques $\eta_0$,
$\eta$, $g:Y\to Y_0$ un morphisme dominant, $f_0:X_0\to Y_0$ un morphisme
dominant tel que $f_0^{-1}(\eta_0)$ soit irréductible. Soit $X'_0$ l'unique
composante irréductible de $X_0$ rencontrant $f_0^{-1}(\eta_0)$
($\mathbf 0_{\mathrm I}$, 2.1.8), et désignons encore par $X'_0$ le
sous-préschéma fermé réduit de $X_0$ ayant $X'_0$ pour espace sous-jacent.
Supposons que le préschéma $X=X_0\times_{Y_0}Y$ soit intègre~; alors $X$ est
isomorphe à $X'_0\times_{Y_0}Y$.
\end{lemma}

En effet, si $j_0:X'_0\to X_0$ est l'injection canonique, qui est une immersion
fermée, $j=j_0\times 1_Y:X'_0\times_{Y_0}Y\to X_0\times_{Y_0}Y=X$ est une
immersion fermée. D'autre part, $X'_0$ contient la fibre
$f_0^{-1}(\eta_0)$, donc $X'_0\times_{Y_0}Y$ contient $f^{-1}(\eta)$~;
notons d'autre part que $f^{-1}(\eta)$ n'est pas vide (\textbf{I}, 3.4.7),
donc contient le point générique de $X$~; par suite l'image de $j$ est
nécessairement $X$ tout entier. Mais comme $X$ est intègre, le seul
sous-préschéma fermé de $X$ ayant $X$ pour espace sous-jacent est $X$
lui-même, donc $j$ est un isomorphisme.

Ce lemme étant établi, on peut donc supposer que $X_0$ est intègre~; pour tout
$y_0\in Y_0$, $\mathcal O_{y_0}$ est une $\mathbf Z$-algèbre essentiellement
de type fini, donc un anneau universellement

\oldpage[IV-3]{189}
caténaire (5.6.4)~; par suite, pour tout $y_0\in f_0(X_0)$, toutes les
composantes irréductibles de $f_0^{-1}(y)$ ont une dimension $\geq e$, car on
sait que $\dim(f_0^{-1}(\eta_0))=e$ par transitivité des fibres (4.1.4). Pour
tout $y\in f(X)$, on a alors $y_0=g(y)\in f_0(X_0)$ et par transitivité des
fibres et (4.2.7), on achève la démonstration.
\end{proof}

\begin{theorem}[13.1.3]
\label{IV.13.1.3-fr}
\textnormal{(Chevalley).} --- Soit $f:X\to Y$ un morphisme localement de type
fini. Pour tout entier $n$, l'ensemble $F_n(X)$ des $x\in X$ tels que
$\dim_x(f^{-1}(f(x)))\geq n$ est fermé~; autrement dit, la fonction
$x\mapsto\dim_x(f^{-1}(f(x)))$ est semi-continue supérieurement dans $X$.
\end{theorem}

\begin{proof}
I) Supposons d'abord que $f$ soit localement de présentation finie.

La question est évidemment locale sur $X$ et sur $Y$, et on peut donc supposer
$Y=\operatorname{Spec}(A)$, $X=\operatorname{Spec}(B)$ affines, $B$ étant une
$A$-algèbre de présentation finie. On sait alors (8.9.1) qu'il y a un
sous-anneau noethérien $A_0$ de $A$, et une $A_0$-algèbre de type fini $B_0$
tels que si l'on pose $Y_0=\operatorname{Spec}(A_0)$,
$X_0=\operatorname{Spec}(B_0)$, on ait

\[
  X=X_0\times_{Y_0}Y,\qquad f=f_0\times 1_Y,
\]

où $f_0:X_0\to Y_0$ correspond à l'homomorphisme $A_0\to B_0$. Soient
$g:Y\to Y_0$ le morphisme correspondant à l'injection canonique $A_0\to A$,
$y$ un point de $Y$, $y_0=g(y)$~; on sait que
$f^{-1}(y)=f_0^{-1}(y_0)\otimes_{k(y_0)}k(y)$, et il résulte de (4.2.7) que
si $p$ est la projection canonique $f^{-1}(y)\to f_0^{-1}(y_0)$, les
composantes irréductibles de $f^{-1}(y)$ sont les composantes irréductibles des
ensembles $p^{-1}(Z_0)$, où $Z_0$ parcourt l'ensemble des composantes
irréductibles de $f_0^{-1}(y_0)$, et chacune des composantes irréductibles de
$p^{-1}(Z_0)$ domine $Z_0$ et a une dimension égale à $\dim(Z_0)$. Tenant
compte de ($\mathbf 0$, 14.1.5), on voit donc que si $g':X\to X_0$ est la
projection canonique, $x$ un point de $X$ et $x_0=g'(x)$, on a
$\dim_x(f^{-1}(f(x)))=\dim_{x_0}(f_0^{-1}(f_0(x_0)))$, d'où
$F_n(X)=g'^{-1}(F_n(X_0))$, et on est par suite ramené à démontrer le théorème
lorsque $Y$ est \emph{noethérien}, ce que nous supposerons désormais.

On peut évidemment supposer $X$ et $Y$ réduits (1.5.4). En considérant
l'ensemble des parties fermées $Y'$ de $Y$ telles que le théorème soit vrai
pour le sous-préschéma fermé de $Y$ ayant $Y'$ pour espace sous-jacent et pour
$X'=f^{-1}(Y')$, on peut raisonner par récurrence noethérienne
($\mathbf 0_{\mathrm I}$, 2.2.2), et supposer que le théorème est vrai pour
toute partie fermée $Y'\neq Y$ de $Y$. Si $X_i$ ($1\leq i\leq m$) sont les
sous-préschémas fermés réduits de $X$ ayant pour espaces sous-jacents les
composantes irréductibles de $X$, on a $F_n(X)=\bigcup_iF_n(X_i)$ en vertu de
($\mathbf 0$, 14.1.5), et l'on peut donc se borner à démontrer le théorème pour
chacun des $X_i$, autrement dit, on peut supposer $X$ irréductible. Si $Z$ est
le sous-préschéma fermé de $Y$ ayant $\overline{f(X)}$ pour espace sous-jacent,
$f$ se factorise en $X\xrightarrow{g}Z\xrightarrow{j}Y$, où $j$ est
l'injection canonique (\textbf{I}, 5.2.2), et $g$ est de type fini (1.5.4),
donc il suffit de démontrer le théorème pour $Z$ et $g$~; en vertu de
l'hypothèse de récurrence, on est donc ramené à considérer seulement le cas où
$Z=Y$, autrement dit où $Y$ est irréductible et $f$ dominant. Soit alors
$\eta$ le point générique de $Y$ et posons $e=\dim(f^{-1}(\eta))$~; il résulte
de (13.1.1) que pour $n\leq e$ on a $F_n(X)=X$, et par suite on peut se borner
au cas où $n>e$. Mais alors (9.5.6), il y a un voisinage ouvert $U$ de $\eta$
dans $Y$ tel que $F_n(X)\subset f^{-1}(Y-U)$~; comme $Y-U\neq Y$, l'hypothèse
de récurrence entraîne que $F_n(X)$ est fermé.

II) Passons maintenant au cas général, en supposant toujours
$S=\operatorname{Spec}(A)$ et $X=\operatorname{Spec}(B)$ affines, $B$ étant
une $A$-algèbre de type fini, donc de la forme

\oldpage[IV-3]{190}
$B=A[T_1,\ldots,T_n]/\mathfrak J$. Soit $(\mathfrak J_\lambda)$ la famille des
idéaux de type fini de $A[T_1,\ldots,T_n]$ contenus dans $\mathfrak J$, de
sorte que $\mathfrak J$ est la réunion filtrante des $\mathfrak J_\lambda$~;
si $X$ et les
$X_\lambda=\operatorname{Spec}(A[T_1,\ldots,T_n]/\mathfrak J_\lambda)$ sont
considérés comme des sous-préschémas fermés de
$Z=\operatorname{Spec}(A[T_1,\ldots,T_n])$, on a donc, pour les espaces
sous-jacents, $X=\bigcap_\lambda X_\lambda$. Si $f_\lambda:X_\lambda\to S$ est
le morphisme structural, on en déduit que
$f^{-1}(s)=\bigcap_\lambda f_\lambda^{-1}(s)$ pour tout $s\in S$, et comme les
ensembles $f_\lambda^{-1}(s)$ sont fermés dans l'espace noethérien $Z_s$, il
existe un $\lambda$ (dépendant de $s$) tel que
$f^{-1}(s)=f_\lambda^{-1}(s)$. Si alors, pour tout $x\in X$, on pose
$d(x)=\dim_x(f^{-1}(f(x)))$, $d_\lambda(x)=\dim_x(f_\lambda^{-1}(f_\lambda(x)))$,
ce qui précède prouve que l'on a $d(x)=\inf_\lambda d_\lambda(x)$~; les
fonctions $d_\lambda$ étant semi-continues supérieurement d'après la première
partie de la démonstration, il en est de même de $d$. C.Q.F.D.
\end{proof}

\begin{corollary}[13.1.4]
\EGAReaderTarget{ega4.c04.cor.13-1-4}{IV.13.1.4-fr}
Sous les hypothèses de (13.1.3), l'ensemble des $x\in X$ tels que $x$ soit
isolé dans $f^{-1}(f(x))$ est ouvert dans $X$.
\end{corollary}

\begin{proof}
En effet, c'est le complémentaire de $F_1(X)$ ($\mathbf 0$, 14.1.10).
\end{proof}

On notera qu'on retrouve ainsi, sous des hypothèses plus générales, la
conséquence (\textbf{III}, 4.4.10) du «~Main theorem~» de Zariski.

\begin{corollary}[13.1.5]
\label{IV.13.1.5-fr}
Sous les hypothèses de (13.1.3), supposons en outre que $f$ soit un morphisme
fermé. Alors, pour tout entier $n$, l'ensemble des $y\in Y$ tels que
$\dim(f^{-1}(y))\geq n$ est fermé, autrement dit, l'application
$y\mapsto\dim(f^{-1}(y))$ est semi-continue supérieurement~; en particulier,
si $y$ est spécialisation de $y'$, on a
$\dim(f^{-1}(y))\geq\dim(f^{-1}(y'))$.
\end{corollary}

\begin{proof}
En effet, dire que $\dim(f^{-1}(y))\geq n$ signifie que $y\in f(F_n(X))$
($\mathbf 0$, 14.1.6).
\end{proof}

\begin{corollary}[13.1.6]
\label{IV.13.1.6-fr}
Soient $X$, $Y$ deux préschémas irréductibles, $f:X\to Y$ un morphisme
dominant localement de type fini. Soit $\xi$ (resp. $\eta$) le point générique
de $X$ (resp. $Y$) et soit $e=\dim(f^{-1}(\eta))$. Alors, pour tout $x\in X$,
on a $\dim_x(f^{-1}(f(x)))\geq e$.
\end{corollary}

\begin{proof}
En effet, l'ensemble des $x$ tels que $\dim_x(f^{-1}(f(x)))<e$ est ouvert en
vertu de (13.1.3), et comme il ne peut contenir $\xi$, il est vide.
\end{proof}

Notons enfin le résultat plus facile suivant~:

\begin{proposition}[13.1.7]
\label{IV.13.1.7-fr}
Soient $Y$ un préschéma quasi-compact, $f:X\to Y$ un morphisme de type fini.
Il existe un entier $n$ tel que, pour tout $y\in Y$, on ait
$\dim(f^{-1}(y))\leq n$.
\end{proposition}

\begin{proof}
Comme il y a un recouvrement ouvert affine fini $(V_i)$ de $Y$ tel que chaque
$f^{-1}(V_i)$ soit réunion finie d'ouverts affines, on est aussitôt ramené au
cas où $Y=\operatorname{Spec}(A)$ et $X=\operatorname{Spec}(B)$ sont affines,
$B$ étant une $A$-algèbre de type fini. Si $B$ admet un système de $n$
générateurs, alors, pour tout $y\in Y$, $B\otimes_Ak(y)$ est une
$k(y)$-algèbre admettant $n$ générateurs, donc $\dim(f^{-1}(y))\leq n$ en
vertu de (4.1.1).
\end{proof}

\subsection{Morphismes équidimensionnels : cas des morphismes dominants de préschémas irréductibles}
\label{subsection:IV.13.2-fr}

\begin{env}[13.2.1]
\label{IV.13.2.1-fr}
Soient $Y$ un préschéma irréductible, $X$ un préschéma irréductible,
$f:X\to Y$ un morphisme dominant localement de type fini~; soit $\eta$ le
point générique de $Y$. On sait (13.1.6) que pour tout $x\in X$, on a
$\dim_x(f^{-1}(f(x)))\geq\dim(f^{-1}(\eta))$.
\end{env}

\oldpage[IV-3]{191}
\begin{definition}[13.2.2]
\label{IV.13.2.2-fr}
Sous les hypothèses de (13.2.1), on dit que $f$ est équidimensionnel au point
$x$ (ou que $X$ est équidimensionnel sur $Y$ au point $x$) si
\[
  \dim_x(f^{-1}(f(x)))=\dim(f^{-1}(\eta)).
\]
On dit que $f$ est équidimensionnel (ou que $X$ est équidimensionnel sur $Y$)
si $f$ est équidimensionnel en tout point $x\in X$.
\end{definition}

Il résulte du théorème de Chevalley (13.1.3) que l'ensemble des $x\in X$ où
$f$ est équidimensionnel est \emph{ouvert non vide}. En outre, si $f$ est
équidimensionnel au point $x$, toutes les composantes irréductibles de
$f^{-1}(f(x))$ qui contiennent $x$ ont \emph{même dimension}, car chacune
d'elles a une dimension qui est $\geq\dim(f^{-1}(\eta))$ (13.1.6) et
$\leq\dim_x(f^{-1}(f(x)))$ en vertu de ($\mathbf 0$, 14.1.5).

\begin{proposition}[13.2.3]
\label{IV.13.2.3-fr}
Soient $Y$ un préschéma irréductible localement noethérien, $X$ un préschéma
irréductible, $f:X\to Y$ un morphisme dominant localement de type fini~;
soient $\eta$ le point générique de $Y$, $x$ un point de $X$, $y=f(x)$, et
supposons que l'on ait
\[
\tag{13.2.3.1}
\label{IV.13.2.3.1-fr}
  \dim(\mathcal O_x)=\dim(\mathcal O_y)
  +\dim(\mathcal O_x\otimes_{\mathcal O_y}k(y)).
\]
Alors $f$ est équidimensionnel au point $x$. La réciproque est vraie si les
deux membres de l'inégalité (5.6.5.2) sont égaux, en particulier si
$\mathcal O_y$ est universellement caténaire.
\end{proposition}

\begin{proof}
Cela résulte aussitôt de (5.6.5.2) et de l'inégalité $\delta(x)\geq0$
(13.1.1).
\end{proof}

\begin{env}[13.2.4]
\label{IV.13.2.4-fr}
Soient maintenant, de façon générale, $Y$ un préschéma localement noethérien,
$f:X\to Y$ un morphisme de type fini, $x$ un point de $X$, $X'$ une partie
fermée irréductible de $X$ contenant $x$, $y=f(x)$,
$Y'=\overline{f(X')}$, $\eta'$ le point générique de $Y'$. Désignons encore
par $X'$ et $Y'$ les sous-préschémas fermés réduits de $X$ et $Y$
respectivement ayant $X'$ et $Y'$ pour espaces sous-jacents~; la restriction
$X'\to Y$ de $f$ se factorise donc en
$X'\xrightarrow{f'}Y'\xrightarrow{j}Y$, où $j$ est l'injection canonique
(\textbf{I}, 5.2.2), et $f'$ est de type fini (1.5.4). Posons alors pour
abréger
\[
\tag{13.2.4.1}
\label{IV.13.2.4.1-fr}
  A=\mathcal O_{Y,y},\quad B=\mathcal O_{X,x},\quad
  A'=\mathcal O_{Y',y},\quad B'=\mathcal O_{X',x}.
\]
La formule (5.6.5.2) appliquée au morphisme dominant $f'$ et aux préschémas
irréductibles $X'$, $Y'$, donne
\[
\tag{13.2.4.2}
\label{IV.13.2.4.2-fr}
  \dim(B')\leq\dim(A')+\dim(B'\otimes_{A'}k(y))
  -(\dim_x(f'^{-1}(y))-\dim(f'^{-1}(\eta'))).
\]
D'autre part, l'anneau local $A'$ (resp. $B'$, $B'\otimes_{A'}k(y)$) est un
anneau quotient de $A$ (resp. $B$, $B\otimes_Ak(y)$), donc
($\mathbf 0$, 16.1.2.1) on a
\[
\tag{13.2.4.3}
\label{IV.13.2.4.3-fr}
  \dim(A')\leq\dim(A),\quad \dim(B')\leq\dim(B),\quad
  \dim(B'\otimes_{A'}k(y))\leq\dim(B\otimes_Ak(y)).
\]
On déduit donc en premier lieu de (13.2.4.3) et ($\mathbf 0$, 16.3.9)
\[
\tag{13.2.4.4}
\label{IV.13.2.4.4-fr}
  \dim(B')\leq\dim(A')+\dim(B'\otimes_{A'}k(y))
  \leq\dim(A)+\dim(B\otimes_Ak(y)).
\]
En outre, en vertu de ($\mathbf 0$, 16.3.9), on a aussi les inégalités
\[
\tag{13.2.4.5}
\label{IV.13.2.4.5-fr}
  \dim(B')\leq\dim(B)\leq\dim(A)+\dim(B\otimes_Ak(y)).
\]
\end{env}

\oldpage[IV-3]{192}
La comparaison de ces inégalités montre donc que~:

\begin{lemma}[13.2.5]
\label{IV.13.2.5-fr}
Avec les notations de (13.2.4), les conditions suivantes sont équivalentes~:
\begin{enumerate}[label=\emph{\alph*)}]
  \item $\dim(B')=\dim(A)+\dim(B\otimes_Ak(y))$.
  \item On a simultanément les relations suivantes~:
  \begin{enumerate}[label=(\roman*)]
    \item $\dim(A')=\dim(A)$.
    \item $\dim(B'\otimes_{A'}k(y))=\dim(B\otimes_Ak(y))$, autrement dit
    \[
      \dim_x(X'\cap f^{-1}(y))=\dim_x(f^{-1}(y)).
    \]
    \item $\dim_x(f'^{-1}(y))=\dim(f'^{-1}(\eta'))$, autrement dit
    $f':X'\to Y'$ est équidimensionnel (13.2.2) au point $x$.
    \item On a l'égalité
    \[
      \dim(B')=\dim(A')+\dim(B'\otimes_{A'}k(y))
      -(\dim_x(f'^{-1}(y))-\dim(f'^{-1}(\eta')))
    \]
  \end{enumerate}
  (relation qui est toujours vérifiée lorsque $A=\mathcal O_{Y,y}$ est un
  anneau universellement caténaire, en vertu de (5.6.5)).
  \item On a simultanément les relations suivantes~:
  \begin{enumerate}[label=(\roman*)]
    \item $\dim(B')=\dim(B)$.
    \item $\dim(B)=\dim(A)+\dim(B\otimes_Ak(y))$.
  \end{enumerate}
\end{enumerate}
\end{lemma}

\begin{env}[13.2.6]
\label{IV.13.2.6-fr}
Rappelons maintenant que les composantes irréductibles $X_i$ de $X$ contenant
$x$ sont en nombre fini et que l'on a ((5.1.2.1) et ($\mathbf 0$, 14.2.1.1))
\[
\tag{13.2.6.1}
\label{IV.13.2.6.1-fr}
  \dim(\mathcal O_{X,x})=\sup_i\dim(\mathcal O_{X_i,x}).
\]
L'équivalence des conditions $b)$ et $c)$ dans (13.2.5) implique par suite,
compte tenu de ($\mathbf 0$, 16.3.9) et ($\mathbf 0$, 14.2.1)~:
\end{env}

\begin{proposition}[13.2.7]
\label{IV.13.2.7-fr}
Soient $Y$ un préschéma localement noethérien, $f:X\to Y$ un morphisme de type
fini, $x$ un point de $X$, $f(x)=y$~; on a
\[
\tag{13.2.7.1}
\label{IV.13.2.7.1-fr}
  \dim(\mathcal O_{X,x})\leq\dim(\mathcal O_{Y,y})
  +\dim(\mathcal O_{X,x}\otimes_{\mathcal O_{Y,y}}k(y)).
\]
Pour que les deux membres de (13.2.7.1) soient égaux, il faut et il suffit
qu'il existe une partie fermée irréductible $X'$ de $X$ contenant $x$ et
vérifiant simultanément les conditions suivantes~:
\begin{enumerate}[label=(\roman*)]
  \item Si $Y'=\overline{f(X')}$, on a
  $\dim(\mathcal O_{Y',y})=\dim(\mathcal O_{Y,y})$.
  \item $\dim_x(X'\cap f^{-1}(y))=\dim_x(f^{-1}(y))$ (autrement dit, $X'$
  contient une des composantes irréductibles de $f^{-1}(y)$ contenant $x$, de
  dimension maxima parmi toutes ces composantes). Il revient au même de dire
  que
  $\dim(\mathcal O_{X',x}\otimes_{\mathcal O_{Y',y}}k(y))
  =\dim(\mathcal O_{X,x}\otimes_{\mathcal O_{Y,y}}k(y))$.
  \item $X'$ est équidimensionnel sur $Y'$ au point $x$, autrement dit, on a
  \[
    \dim_x(X'\cap f^{-1}(y))=\dim(X'\cap f^{-1}(\eta')),
  \]
  où $\eta'$ est le point générique de $Y'$ (et par suite, toutes les
  composantes irréductibles de $X'\cap f^{-1}(y)$ contenant $x$ ont même
  dimension (13.2.2)).

  \oldpage[IV-3]{193}
  \item On a l'égalité
  \[
    \dim(\mathcal O_{X',x})=\dim(\mathcal O_{Y',y})
    +\dim(\mathcal O_{X',x}\otimes_{\mathcal O_{Y',y}}k(y))
  \]
  (condition toujours impliquée par (iii) lorsque $\mathcal O_{Y,y}$ est un
  anneau universellement caténaire).
\end{enumerate}
De plus, $X'$ est alors une composante irréductible de $X$ et $Y'$ une
composante irréductible de $Y$.
\end{proposition}

Dans l'énoncé, les anneaux locaux $\mathcal O_{X',x}$ et
$\mathcal O_{Y',y}$ sont relatifs aux sous-préschémas fermés réduits de $X$,
$Y$ respectivement ayant $X'$ et $Y'$ pour espaces sous-jacents respectifs.

En outre, cela prouve que

\begin{corollary}[13.2.8]
\label{IV.13.2.8-fr}
Si les deux membres de (13.2.7.1) sont égaux, les parties fermées irréductibles
$X'$ de $X$ contenant $x$ et vérifiant les conditions (i) à (iv) de (13.2.7)
sont exactement celles pour lesquelles on a
$\dim(\mathcal O_{X',x})=\dim(\mathcal O_{X,x})$ (ce qui implique
nécessairement que $X'$ est une composante irréductible de $X$).
\end{corollary}

La proposition (13.2.7) entraîne~:

\begin{corollary}[13.2.9]
\label{IV.13.2.9-fr}
Soient $Y$ un préschéma localement noethérien, $f:X\to Y$ un morphisme de type
fini, $x$ un point de $X$, $f(x)=y$. Supposons que $\mathcal O_{Y,y}$ soit un
anneau universellement caténaire. Alors les conditions suivantes sont
équivalentes~:
\begin{enumerate}[label=\emph{\alph*)}]
  \item Les deux membres de (13.2.7.1) sont égaux.
  \item Il existe une composante irréductible $Z$ de $f^{-1}(y)$ contenant
  $x$, de dimension $\dim_x(f^{-1}(y))$, telle que, pour tout $x'$ dans un
  voisinage de $x$ dans $Z$, on ait
  \[
  \tag{13.2.9.1}
  \label{IV.13.2.9.1-fr}
    \dim(\mathcal O_{X,x'})=\dim(\mathcal O_{Y,y})
    +\dim(\mathcal O_{X,x'}\otimes_{\mathcal O_{Y,y}}k(y)).
  \]
  \item Il existe une composante irréductible $Z$ de $f^{-1}(y)$ contenant
  $x$, de dimension $\dim_x(f^{-1}(y))$, telle que pour le point générique $z$
  de $Z$, on ait
  \[
  \tag{13.2.9.2}
  \label{IV.13.2.9.2-fr}
    \dim(\mathcal O_{X,z})=\dim(\mathcal O_{Y,y}).
  \]
\end{enumerate}
\end{corollary}

Montrons que $a)$ entraîne $b)$. Soit $\dim_x(f^{-1}(y))=n$~; en vertu de
(13.2.7), il existe une composante irréductible $X'$ de $X$ vérifiant les
conditions (i) à (iv) de (13.2.7)~; soit $Z$ une composante irréductible de
dimension $n$ de $f^{-1}(y)$, contenant $x$ et contenue dans $X'$. Comme
$f^{-1}(y)$ est localement noethérien, il existe un voisinage ouvert $U$ de
$x$ dans $Z$ tel que $U$ ne rencontre aucune composante irréductible de
$f^{-1}(y)$ autre que celles qui contiennent $x$, donc (4.1.1.3)
$\dim_{x'}(f^{-1}(y))=n$ pour tout $x'\in U$~; il est clair alors que les
conditions (i) à (iii) de (13.2.7) sont satisfaites quand on y remplace $x$
par un point quelconque $x'\in U$, et il en est de même de la condition (iv)
puisque $\mathcal O_{Y,y}$ est universellement caténaire~; d'où la conclusion
par (13.2.7). La condition $b)$ entraîne trivialement $c)$ en vertu de
(5.1.2). Enfin, si $c)$ est vérifiée et si $X''$ est une composante
irréductible de $X$ contenant $Z$ et telle que les conditions (i) à (iv) de
(13.2.7) soient vérifiées quand on y remplace $X'$ par $X''$ et $x$ par $z$,
il est clair que ces conditions sont aussi vérifiées pour $X''$ et $x$ puisque
$\mathcal O_{Y,y}$ est universellement caténaire, donc $c)$ implique $a)$.

\begin{proposition}[13.2.10]
\label{IV.13.2.10-fr}
Soient $Y$ un préschéma irréductible localement noethérien, $\eta$ son point
générique, $f:X\to Y$ un morphisme de type fini, $y$ un point de $f(X)$.
Soient $Z_i$ les

\oldpage[IV-3]{194}
composantes irréductibles de $f^{-1}(y)$, $z_i$ le point générique de $Z_i$,
et considérons les conditions suivantes~:
\begin{enumerate}[label=\emph{\alph*)}]
  \item Pour tout $x\in f^{-1}(y)$, on a la relation
  \[
  \tag{13.2.10.1}
  \label{IV.13.2.10.1-fr}
    \dim(\mathcal O_{X,x})=\dim(\mathcal O_{Y,y})
    +\dim(\mathcal O_{X,x}\otimes_{\mathcal O_{Y,y}}k(y)).
  \]
  \item Pour tout $i$, on a
  \[
  \tag{13.2.10.2}
  \label{IV.13.2.10.2-fr}
    \dim(\mathcal O_{X,z_i})=\dim(\mathcal O_{Y,y}).
  \]
  \item Pour tout $i$, il existe une composante irréductible $X_i$ de $X$
  contenant $Z_i$ et telle que
  $\dim(X_i\cap f^{-1}(\eta))=\dim(Z_i)$ (autrement dit, telle que le
  \emph{sous-préschéma fermé réduit} $X_i$ de $X$ soit équidimensionnel sur $Y$
  au point $z_i$).
\end{enumerate}
Alors $a)$ entraîne $b)$ et $b)$ entraîne $c)$~; en outre, si
$\mathcal O_{Y,y}$ est universellement caténaire, les trois conditions $a)$,
$b)$, $c)$ sont équivalentes.
\end{proposition}

L'anneau $\mathcal O_{X,z_i}\otimes_{\mathcal O_{Y,y}}k(y)$ étant de dimension
$0$ (5.1.2), $a)$ entraîne évidemment $b)$~; $b)$ entraîne $c)$ en vertu de
(13.2.7) appliqué au point $z_i$. Inversement, supposons que
$\mathcal O_{Y,y}$ soit universellement caténaire~; comme la condition $c)$
implique que $f(X_i)$ est dense dans $Y$, il résulte de $c)$ que les conditions
(i) à (iv) de (13.2.7) sont vérifiées en remplaçant $X'$ par $X_i$ et $x$ par
$z_i$, donc $c)$ implique $b)$~; enfin $b)$ implique $a)$ en vertu de
(13.2.9).

\begin{corollary}[13.2.11]
\label{IV.13.2.11-fr}
Les notations étant celles de (13.2.10), supposons que $\mathcal O_{Y,y}$ soit
universellement caténaire. Pour tout $y'\in Y$, soit $E(y')$ l'ensemble des
dimensions des composantes irréductibles de $f^{-1}(y')$ et posons
$d(y')=\dim(f^{-1}(y'))=\sup(E(y'))$. Alors, si les conditions équivalentes
$a)$, $b)$, $c)$ de (13.2.10) sont vérifiées, on a
$E(y)\subset E(\eta)$, d'où $d(y)\leq d(\eta)$.
\end{corollary}

En effet, avec les notations de (13.2.10), $X_i\cap f^{-1}(\eta)$ n'est pas
vide et est par suite une composante irréductible de $f^{-1}(\eta)$
($\mathbf 0_{\mathrm I}$, 2.1.8).

\begin{remarks}[13.2.12]
\label{IV.13.2.12-fr}
\begin{enumerate}[label=(\roman*)]
  \item Rappelons (6.1.2) que la relation (13.2.10.1) est toujours vérifiée
  lorsque le morphisme $f$ est \emph{plat} au point $x$.
  \item Avec les notations de (13.2.10), supposons que $\mathcal O_{Y,y}$ soit
  universellement caténaire et en outre que le morphisme $f$ soit \emph{propre}~;
  alors, si les conditions équivalentes $a)$, $b)$, $c)$ de (13.2.10) sont
  satisfaites, on a même $d(y)=d(\eta)$, car il résulte de (13.1.5) que l'on a
  $d(y)\geq d(\eta)$.
  \item Le morphisme de (12.2.3, (ii)) est \emph{propre} et \emph{plat} et tous
  les anneaux locaux de $Y$ sont universellement caténaires~; en outre, les deux
  composantes irréductibles $X_1$, $X_2$ de $X$ sont \emph{équidimensionnelles}
  sur $Y$ en tout point~; mais $E(y)$ a \emph{deux} éléments pour tout
  $y\neq y_0$, alors que $E(y_0)$ est réduit à un \emph{seul} élément, donc
  $E(y)$ n'est \emph{pas constant} dans $Y$.
\end{enumerate}
\end{remarks}

\subsection{Morphismes équidimensionnels : cas général}
\label{subsection:IV.13.3-fr}

\begin{proposition}[13.3.1]
\label{IV.13.3.1-fr}
Soient $Y$ un préschéma, $f:X\to Y$ un morphisme localement de type fini,
$x$ un point de $X$, $y=f(x)$. Désignons par $Y_\alpha$ les composantes
irréductibles de $Y$ contenant $y$. Alors les conditions suivantes sont
équivalentes~:
\begin{enumerate}[label=\emph{\alph*)}]
\oldpage[IV-3]{195}
  \item Il existe un entier $e$ et un voisinage ouvert $U$ de $x$ tels que
  l'image par $f$ de toute composante irréductible de $U$ soit dense dans un
  $Y_\alpha$, et que, pour tout $x'\in U$, l'espace
  $U\cap f^{-1}(f(x'))$ soit équidimensionnel et de dimension $e$.
  \item[\emph{a$'$)}] Il existe un entier $e$ et un voisinage ouvert $U$ de
  $x$ tels que l'image par $f$ de toute composante irréductible de $U$ soit
  dense dans un $Y_\alpha$ et que, si on désigne par $y_\alpha$ le point
  générique de $Y_\alpha$, toute composante irréductible des espaces
  $U\cap f^{-1}(y)$, $U\cap f^{-1}(y_\alpha)$ soit de dimension $e$.
  \item[\emph{a$''$)}] Il existe un entier $e$ et un voisinage ouvert $U$ de
  $x$ tels que, pour chacune des composantes irréductibles $U_\lambda$ de
  $U$, $f(U_\lambda)$ soit dense dans un $Y_\alpha$ et que, pour tout
  $x'\in U_\lambda$, les composantes irréductibles de
  $U_\lambda\cap f^{-1}(f(x'))$ soient toutes de dimension $e$.
  \item Il existe un entier $e$, un voisinage ouvert $U$ de $x$ et un
  $Y$-morphisme quasi-fini
  \[
    g:U\longrightarrow Y\otimes_{\mathbf Z}\mathbf Z[T_1,\ldots,T_e]
  \]
  (préschéma que nous noterons aussi $Y[T_1,\ldots,T_e]$ pour abréger) tels
  que l'image par $g$ de toute composante irréductible de $U$ soit dense dans
  une composante irréductible de $Y[T_1,\ldots,T_e]$.
\end{enumerate}
\end{proposition}

\begin{proof}
Il est immédiat que $a''$) entraîne $a$), car pour tout $x\in U$, les
composantes irréductibles de $U\cap f^{-1}(f(x))$ sont chacune composante
irréductible d'un des $U_\lambda\cap f^{-1}(f(x))$, d'où la conclusion par
($\mathbf 0$, 14.1.4). La condition $a$) entraîne trivialement $a'$).
Montrons ensuite que $a'$) entraîne $a''$)~; on peut se borner au cas où $f$
est de type fini et $X$ et $Y$ réduits (1.5.4). Soit $U_\lambda$ une
composante irréductible de $U$, et supposons que $f(U_\lambda)$ soit dense
dans $Y_\alpha$~; alors la restriction de $f$ à $U_\lambda$ se factorise en
$U_\lambda\xrightarrow{f_\alpha}Y_\alpha\to Y$, où $f_\alpha$ est de type
fini et dominant (\textbf{I}, 5.2.2). Soit $x_\lambda$ le point générique de
$U_\lambda$~; en vertu de ($\mathbf 0_{\mathrm I}$, 2.1.8),
$U_\lambda\cap f_\alpha^{-1}(y_\alpha)$ est l'unique composante irréductible
de $U\cap f^{-1}(y_\alpha)$ contenant $x_\lambda$ et est par hypothèse de
dimension $e$, égale aux dimensions de toutes les composantes irréductibles
de $U_\lambda\cap f^{-1}(y)$, en vertu de l'hypothèse et de (13.1.1). Mais en
vertu du théorème de Chevalley (13.1.3) et de (13.1.1), l'ensemble $V_\lambda$
des $x'\in U_\lambda$ tels que
$\dim_{x'}(f_\lambda^{-1}(f_\lambda(x')))=e$ est ouvert et contient $x$, et il
suffit de prendre la réunion des $V_\lambda$ pour obtenir un ensemble ouvert
satisfaisant aux conditions de $a''$).

Prouvons maintenant que $a$) entraîne $b$)~; on peut se limiter au cas où
$Y=\operatorname{Spec}(A)$ et $X=\operatorname{Spec}(B)$ sont affines et où
$U=X$. Prouvons d'abord le lemme suivant~:

\begin{lemma}[13.3.1.1]
\label{IV.13.3.1.1-fr}
Soient $Y=\operatorname{Spec}(A)$, $X=\operatorname{Spec}(B)$ deux schémas
affines, $f:X\to Y$ un morphisme de type fini, $y$ un point de $f(X)$, et
posons $e=\dim(f^{-1}(y))$. Alors il existe un $Y$-morphisme
$g:X\to Y[T_1,\ldots,T_e]$ tel que (si l'on pose
$X_y=X\times_Y\operatorname{Spec}(k(y))$), le morphisme
$g_y:X_y\to\operatorname{Spec}(k(y)[T_1,\ldots,T_e])$ soit fini. En outre,
pour un tel morphisme $g$, $g_y$ est nécessairement surjectif et il existe un
voisinage ouvert $U$ de $f^{-1}(y)$ dans $X$ tel que $g|U$ soit quasi-fini.
\end{lemma}

Soit $\mathfrak p=\mathfrak j_y$~; l'anneau $B\otimes_Ak(\mathfrak p)$ est une
$k(\mathfrak p)$-algèbre de type fini, donc le lemme de normalisation
(Bourbaki, \emph{Alg. comm.}, chap.~V, \S~3, n\textsuperscript{o}~1, th.~1)
prouve qu'il y a dans $B\otimes_Ak(\mathfrak p)$ une suite finie
$(t_i)_{1\leq i\leq r}$ d'éléments algébriquement indépendants sur
$k(\mathfrak p)$ et tels que si l'on pose
$C'=k(\mathfrak p)[t_1,\ldots,t_r]$, $B\otimes_Ak(\mathfrak p)$ soit une
$C'$-algèbre \emph{finie}~; on a donc
$\dim(B\otimes_Ak(\mathfrak p))=\dim(C')$ ($\mathbf 0$, 16.1.5) et comme
$\dim(C')=r$ (5.2.1), on a $r=e$. Comme
$B\otimes_Ak(\mathfrak p)=(B/\mathfrak pB)\otimes_{A/\mathfrak p}k(\mathfrak p)$,
on peut, en multipliant les $t_i$ par un élément $\neq0$

\oldpage[IV-3]{196}
convenable de $A/\mathfrak p$, supposer que chaque $t_i$ est l'image canonique
dans $B\otimes_Ak(\mathfrak p)$ d'un élément $s_i\in B$. Soit alors
$u:A[T_1,\ldots,T_e]\to B$ l'homomorphisme tel que $u(T_i)=s_i$ pour tout
$i$, et soit $g:X\to Y[T_1,\ldots,T_e]$ le morphisme correspondant. Il est
clair qu'en raison du choix des $s_i$, $g_y$ est un morphisme fini. Pour tout
morphisme $g:X\to Y[T_1,\ldots,T_e]$, tel que $g_y$ soit fini, il résulte de
(5.4.2) et (4.1.2.1) que $g_y$ est nécessairement surjectif. D'autre part, en
vertu de (13.1.4), l'ensemble $U$ des $x\in X$ qui sont isolés dans leur fibre
$g^{-1}(g(x))$ est ouvert dans $X$ et contient $f^{-1}(y)$, et par définition
la restriction $g|U$ est un morphisme quasi-fini.

Ce lemme étant établi, pour prouver que $a$) implique $b$), il reste à voir
que si $x_\lambda$ est un point maximal de $U$, $g(x_\lambda)=z_\lambda$ est
un point maximal de $Z=Y[T_1,\ldots,T_e]$. En vertu de l'hypothèse $a$), on
peut (en restreignant au besoin $U$), supposer que $f(x_\lambda)$ est l'un des
points génériques $y_\alpha$ des composantes irréductibles de $Y$ contenant
$y$~; si $p:Z\to Y$ est le morphisme structural, on a donc
$p(z_\lambda)=y_\alpha$, et l'on déduit par suite de $g$ un
$k(y_\alpha)$-morphisme quasi-fini
\[
  h:U\cap f^{-1}(y_\alpha)\longrightarrow p^{-1}(y_\alpha).
\]
Or $p^{-1}(y_\alpha)=\operatorname{Spec}(k(y_\alpha)[T_1,\ldots,T_e])$ est
intègre et de dimension $e$~; si $z_\lambda$ n'était pas point maximal de
$Z$, il ne serait pas point maximal de $p^{-1}(y_\alpha)$
($\mathbf 0_{\mathrm I}$, 2.1.8) et son adhérence $Z'_\lambda$ dans
$p^{-1}(y_\alpha)$ serait donc de dimension $<e$ (4.1.2.1). Mais comme $h$
est quasi-fini, il résulte de (4.1.2) et de l'hypothèse $a$) que l'on a
$\dim(Z'_\lambda)\geq e$ (la restriction de $h$ à
$\overline{\{x_\lambda\}}$ se factorisant en
\[
  \overline{\{x_\lambda\}}\longrightarrow Z'_\lambda
  \longrightarrow p^{-1}(y_\alpha)
\]
en vertu de (\textbf{I}, 5.2.2))~; on aboutit ainsi à une contradiction, ce
qui montre que $a$) entraîne $b$).

Prouvons finalement que $b$) implique $a$). Notons que le morphisme structural
$p:Z=Y[T_1,\ldots,T_e]\to Y$ est fidèlement plat~; donc (2.3.4) les points
maximaux de $Z$ ont pour images par $p$ les points maximaux de $Y$~; ceci
prouve déjà que les $f(x_\lambda)$ sont points génériques des $Y_\alpha$. En
outre, si $f(x_\lambda)=y_\alpha$, le $k(y_\alpha)$-morphisme
$h:U\cap f^{-1}(y_\alpha)\to p^{-1}(y_\alpha)$ déduit de $g$ est dominant et
quasi-fini par hypothèse~; on a donc
$\dim(U_\lambda\cap f^{-1}(y_\alpha))=e$ en vertu de (4.1.2)
($U_\lambda$ étant la composante irréductible de $U$ de point générique
$x_\lambda$). De même, pour tout $x'\in U_\lambda$, le morphisme
$U_\lambda\cap f^{-1}(f(x'))\to p^{-1}(f(x'))$ déduit de $g$ est un
$k(f(x'))$-morphisme quasi-fini, donc (4.1.2) on a
$\dim(f^{-1}(f(x'))\cap U_\lambda)\leq e$~; mais par ailleurs on sait (13.1.6)
que toutes les composantes irréductibles de
$U_\lambda\cap f^{-1}(f(x'))$ sont de dimension $\geq e$~; on voit ainsi que
ces composantes sont exactement de dimension $e$, donc $b$) entraîne $a$).
C.Q.F.D.
\end{proof}

\begin{definition}[13.3.2]
\label{IV.13.3.2-fr}
Soient $Y$ un préschéma, $f:X\to Y$ un morphisme localement de type fini, $x$
un point de $X$. On dit que $f$ est équidimensionnel au point $x$ (ou que $X$
est équidimensionnel sur $Y$ au point $x$) si les conditions équivalentes de
(13.3.1) sont vérifiées. On dit que $f$ est équidimensionnel (ou que $X$ est
équidimensionnel sur $Y$) si $f$ est équidimensionnel en tout point $x\in X$.
\end{definition}

Il est clair que, pour que $f$ soit équidimensionnel en un point $x$, il faut
et il suffit que $f_{\mathrm{red}}$ le soit. En outre, les conditions de
(13.3.1) montrent que l'ensemble des points où $f$ est équidimensionnel est
\emph{ouvert} dans $X$.

On notera que lorsque $X$ et $Y$ sont \emph{irréductibles}, dire que $f$ est
équidimen-
\oldpage[IV-3]{197}
sionnel au point $x$ signifie, en vertu de (13.3.1, $a''$)), que $f$ est
dominant et que
$\dim_x(f^{-1}(f(x)))=\dim(f^{-1}(\eta))$ (où $\eta$ est le point générique de
$Y$)~; la définition (13.3.2) coïncide donc dans ce cas avec la définition
(13.2.2).

\begin{proposition}[13.3.3]
\label{IV.13.3.3-fr}
Soient $Y$ un préschéma localement noethérien, $f:X\to Y$ un morphisme
localement de type fini, $x$ un point de $X$, $X_j$ les composantes
irréductibles de $X$ (en nombre fini) contenant $x$. Pour que $f$ soit
équidimensionnel au point $x$, il faut et il suffit que, pour tout $j$,
$Y_j=\overline{f(X_j)}$ soit une composante irréductible de $Y$ et, en
désignant par $X_j$ et $Y_j$ les sous-préschémas fermés réduits de $X$ et $Y$
d'espaces sous-jacents $X_j$ et $Y_j$, par $f_j:X_j\to Y_j$ le morphisme
déduit de $f$ (\textbf{I}, 5.2.2), par $y_j$ le point générique de $Y_j$, que
$f_j$ soit équidimensionnel au point $x$ et que tous les nombres
$\dim(f_j^{-1}(y_j))$ soient égaux.
\end{proposition}

Cela résulte aussitôt de (13.3.1, $a'$)).

\begin{corollary}[13.3.4]
\label{IV.13.3.4-fr}
Avec les notations de (13.3.3), posons $y=f(x)$. Si $\mathcal O_{Y,y}$ est un
anneau universellement caténaire et si $f$ est équidimensionnel au point $x$,
on a
\[
\tag{13.3.4.1}
\label{IV.13.3.4.1-fr}
  \dim(\mathcal O_{X_j,x})=\dim(\mathcal O_{Y_j,y})+e
  -\deg.\operatorname{tr}_{k(y)}k(x)
\]
où $e$ est la valeur commune des nombres $\dim(f_j^{-1}(y_j))$.
\end{corollary}

Comme chacun des $\mathcal O_{Y_j,y}$, quotient de $\mathcal O_{Y,y}$, est un
anneau universellement caténaire (5.6.1), l'égalité est conséquence de (5.6.5).

\begin{corollary}[13.3.5]
\label{IV.13.3.5-fr}
Avec les notations de (13.3.4), supposons que $f$ soit équidimensionnel au
point $x$ et que $\mathcal O_{Y,y}$ soit un anneau universellement caténaire.
Si l'anneau $\mathcal O_{Y,y}$ est équidimensionnel, il en est de même de
$\mathcal O_{X,x}$~; la réciproque est vraie si l'image par $f$ de la réunion
des $X_j$ est dense dans un voisinage de $y$.
\end{corollary}

En effet, dire que $\mathcal O_{Y,y}$ est équidimensionnel signifie que les
nombres $\dim(\mathcal O_{Y'_i,y})$ sont égaux pour toutes les composantes
irréductibles $Y'_i$ de $Y$ contenant $y$, comme il résulte de (5.1.1.5)
appliqué au schéma local $\operatorname{Spec}(\mathcal O_{Y,y})$~; comme on a
la relation (13.3.4.1), il en résulte par le même raisonnement que
$\mathcal O_{X,x}$ est alors équidimensionnel. Inversement, si
$\mathcal O_{X,x}$ est équidimensionnel, tous les nombres
$\mathcal O_{Y_j,y}$ sont égaux par (13.3.4.1)~; on en déduit que
$\mathcal O_{Y,y}$ est équidimensionnel si les $Y_j$ sont \emph{toutes} les
composantes irréductibles de $Y$ contenant $y$, ce qui résulte de l'hypothèse
supplémentaire.

\begin{proposition}[13.3.6]
\label{IV.13.3.6-fr}
Soient $Y$ un préschéma localement noethérien, $f:X\to Y$ un morphisme
localement de type fini, $x$ un point de $X$, $y=f(x)$. Supposons que l'anneau
$\mathcal O_y$ soit équidimensionnel. Alors, si $\mathcal O_x$ est
équidimensionnel et si l'on a l'égalité
\[
\tag{13.3.6.1}
\label{IV.13.3.6.1-fr}
  \dim(\mathcal O_x)=\dim(\mathcal O_y)
  +\dim(\mathcal O_x\otimes_{\mathcal O_y}k(y))
\]
(cf. (13.2.7.1)) $f$ est équidimensionnel au point $x$, et la réciproque est
vraie si $\mathcal O_y$ est un anneau universellement caténaire.
\end{proposition}

Gardons les notations de (13.3.3)~; il résulte de (13.2.8) et de l'hypothèse
que $\mathcal O_x$ est équidimensionnel, que chacun des $Y_j$ est une
composante irréductible de $Y$ et que chacun des $f_j$ est équidimensionnel au
point $x$~; en outre (13.2.8), on a (compte tenu de (5.6.5))
\[
  \dim(\mathcal O_{X_j,x})=\dim(\mathcal O_{Y_j,y})
  +\dim_x(X_j\cap f^{-1}(y))-\deg.\operatorname{tr}_{k(y)}k(x).
\]

\oldpage[IV-3]{198}
Or, comme $\mathcal O_y$ est supposé équidimensionnel, cette égalité s'écrit
\[
\tag{13.3.6.2}
\label{IV.13.3.6.2-fr}
  \dim_x(X_j\cap f^{-1}(y))=\dim(\mathcal O_{X,x})
  -\dim(\mathcal O_{Y,y})+\deg.\operatorname{tr}_{k(y)}k(x).
\]
Le premier membre de (13.3.6.2) est donc indépendant de $j$~; mais comme $f_j$
est équidimensionnel au point $x$, on a
$\dim_x(X_j\cap f^{-1}(y))=\dim(f_j^{-1}(y_j))$, donc le critère (13.3.3)
montre que $f$ est équidimensionnel au point $x$.

Inversement, supposons que $\mathcal O_{Y,y}$ soit un anneau universellement
caténaire et que $f$ soit équidimensionnel au point $x$~; alors (13.3.5)
$\mathcal O_{X,x}$ est équidimensionnel, et il résulte alors de (13.3.4) que
l'on a la relation
\[
  \dim(\mathcal O_{X,x})=\dim(\mathcal O_{Y,y})+e
  -\deg.\operatorname{tr}_{k(y)}k(x)
\]
où $e$ est la valeur commune des nombres
$\dim(f_j^{-1}(y_j))=\dim_x(X_j\cap f^{-1}(y))$~; mais par définition $e$ est
aussi égal à $\dim_x(f^{-1}(y))$, d'où la relation (13.3.6.1), compte tenu de
(5.6.5.2).

\begin{proposition}[13.3.7]
\label{IV.13.3.7-fr}
Soient $Y$ un préschéma, $f:X\to Y$ un morphisme localement de type fini et
équidimensionnel, $Z$ une partie fermée de $X$. Alors la fonction
\[
\tag{13.3.7.1}
\label{IV.13.3.7.1-fr}
  x\longmapsto\operatorname{codim}_x
  (Z\cap f^{-1}(f(x)),f^{-1}(f(x)))
\]
est semi-continue inférieurement dans $X$.
\end{proposition}

Comme toutes les composantes irréductibles de $f^{-1}(f(x))$ contenant $x$
ont par hypothèse la même dimension (13.3.1), on a, en vertu de (5.2.1),
\[
  \operatorname{codim}_x(Z\cap f^{-1}(f(x)),f^{-1}(f(x)))
  =\dim_x(f^{-1}(f(x)))-\dim_x(Z\cap f^{-1}(f(x))).
\]
Mais par hypothèse le premier terme du second membre est fonction
\emph{continue} de $x$ (13.3.1) et le second est fonction semi-continue
supérieurement de $x$ en vertu de (13.3.3)~; d'où la conclusion.

\begin{proposition}[13.3.8]
\label{IV.13.3.8-fr}
Soient $Y$ un préschéma, $f:X\to Y$ un morphisme localement de type fini, $x$
un point de $X$. Soient $Y'$ un préschéma, $g:Y'\to Y$ un morphisme plat,
$X'=X\times_Y Y'$, $f'=f_{(Y')}:X'\to Y'$~; si $f$ est équidimensionnel au
point $x$, alors $f'$ est équidimensionnel en tout point $x'\in X'$ au-dessus
de $x$.
\end{proposition}

La question étant locale sur $X$ et $Y$, on peut se borner au cas où toute
composante irréductible de $X$ (resp. $Y$) contient $x$ (resp. $y$)~; comme
l'image par $f$ de toute composante irréductible de $X$ est alors dense dans
une composante irréductible de $Y$ (13.3.1), on sait (2.3.5) que l'image par
$f'$ de toute composante irréductible de $X'$ est dense dans une composante
irréductible de $Y'$. D'autre part, par transitivité des fibres
(\textbf{I}, 3.6.4), il résulte de (4.2.8) et de (13.3.1) que, pour tout
$z'\in X'$, l'ensemble des dimensions des composantes irréductibles de
${f'}^{-1}(f'(z'))$ est le même que l'ensemble des dimensions des composantes
irréductibles de $f^{-1}(f(z))$, où $z$ est la projection de $z'$ dans $X$~;
d'où la conclusion en vertu de (13.3.1, $a$)).

\begin{remark}[13.3.9]
\EGAReaderTarget{ega4.c04.statement.13-3-9}{IV.13.3.9-fr}
La propriété d'équidimensionnalité d'un morphisme $f:X\to Y$ n'est \emph{pas}
stable par changement de base quelconque $g:Y'\to Y$, même lorsque $g$ est
l'injection canonique d'une composante irréductible $Y'$ de $Y$ dans $Y$
($Y'$ étant

\oldpage[IV-3]{199}
considérée comme sous-préschéma fermé réduit de $Y$). Par exemple, soient $k$
un corps, $A_0=k[S,T]$ l'anneau de polynômes à deux indéterminées,
$\mathfrak a=\mathfrak p_1\mathfrak p_2$, où $\mathfrak p_1$ et
$\mathfrak p_2$ sont les idéaux premiers $A_0S$ et $A_0T$ de $A_0$~; soient
$A=A_0/\mathfrak a$ et $Y=\operatorname{Spec}(A)$, qui a deux composantes
irréductibles $Y_1=\operatorname{Spec}(A/\mathfrak p_1)$,
$Y_2=\operatorname{Spec}(A/\mathfrak p_2)$~; prenons $X=Y_1$,
$f:X\to Y$ étant l'injection canonique, qui est évidemment un morphisme
équidimensionnel. Prenons d'autre part $Y'=Y_2$, $g:Y'\to Y$ étant l'injection
canonique~; alors on a
$X'=X\times_Y Y'=\operatorname{Spec}((A/\mathfrak p_1)\otimes_A
(A/\mathfrak p_2))=\operatorname{Spec}(A/(\mathfrak p_1+\mathfrak p_2))$,
et $\mathfrak p_1+\mathfrak p_2$ est un idéal maximal de $A$, donc $X'$ est
réduit à un point, et le morphisme $f':X'\to Y'$ n'est \emph{pas dominant},
l'image par $f'$ de l'unique point de $X'$ étant un point fermé de $Y'$~;
donc $f'$ n'est pas équidimensionnel.

On peut aussi donner un contre-exemple où $X$ et $Y$ sont \emph{intègres}, $f$
\emph{fini} et \emph{birationnel} (et \emph{a fortiori} équidimensionnel par
(13.3.1, $b$)), $g:Y'\to Y$ \emph{fini et dominant}. Soient $A$ et
$\overline A$ les anneaux locaux définis dans (11.7.5), et prenons
$Y=\operatorname{Spec}(A)$, $X=\operatorname{Spec}(\overline A)$~; d'autre
part, avec les notations de (11.7.5), prenons
$Y'=\operatorname{Spec}(\overline B)$~; alors
$X'=\operatorname{Spec}(\overline A\otimes_A\overline B)
=\operatorname{Spec}(\overline B\otimes_B\overline B)$~; mais on vérifie
aussitôt que $\overline B\otimes_B\overline B$ est composé direct des anneaux
$B/\mathfrak p'$, $B/\mathfrak p''$ et de deux anneaux isomorphes à
$B/\mathfrak n$, dont les spectres sont donc réduits à un point~; comme les
projections de ces points sont des points fermés de $Y'$, ici encore on voit
que $f'$ ne transforme pas une composante irréductible de $X'$ en une partie
partout dense d'une composante irréductible de $Y'$, donc $f'$ n'est pas
équidimensionnel.

Dans cet exemple, l'anneau $A$ n'est pas géométriquement unibranche~; nous
allons voir au paragraphe suivant (14.4.6) que de tels phénomènes ne peuvent se
produire lorsque les points de $Y$ sont \emph{géométriquement unibranches}. Le
manque de stabilité de la notion de morphisme équidimensionnel restreint
grandement son intérêt, au profit de la notion de morphisme universellement
ouvert, qui va être étudiée en détail au paragraphe suivant.
\end{remark}

% IV3_B28_P199_END_BEFORE_SECTION_14

% FIN SOURCE ega4/ega4-13-fr.tex

%
% DEBUT SOURCE ega4/ega4-14-fr.tex
\section{Morphismes universellement ouverts}
\label{section:IV.14-fr}

Les \S\S~14 et 15 sont consacrés à l'étude de la notion de morphisme
\emph{universellement ouvert} (2.4.2). On a déjà vu (2.4.6) qu'un morphisme
\emph{plat} et localement de présentation finie est universellement ouvert, la
réciproque étant inexacte. Au \S~14, on examine d'abord les propriétés des
dimensions des fibres d'un morphisme universellement ouvert $f:X\to Y$~;
lorsque $X$ et $Y$ sont localement noethériens, $f$ se comporte à cet égard
(14.2.1) comme un morphisme plat (cf. 6.1.2), et est en particulier
\emph{équidimensionnel} lorsqu'il est localement de type fini et dominant et que
$X$ est irréductible (14.2.2). Inversement, un morphisme équidimensionnel
localement de type fini $f:X\to Y$ est universellement ouvert lorsqu'on suppose
en outre que $Y$ est \emph{géométriquement unibranche}, et en particulier lorsque
$Y$ est \emph{normal} (critère de Chevalley, 14.4.4). Nous montrons aussi que les
morphismes universellement ouverts localement de type fini $f:X\to Y$ (lorsque
$X$ et $Y$ sont
\oldpage[IV-3]{200}
localement noethériens et $Y$ irréductible) admettent «~suffisamment~» de
\emph{quasi-sections}, i.e. au voisinage d'un point fermé $x$ d'une fibre
$f^{-1}(y)$, il y a un sous-préschéma fermé $X'$ de $X$ contenant $x$ et tel que
la restriction $X'\to Y$ de $f$ soit un morphisme \emph{quasi-fini} (donc à fibres
\emph{discrètes}) et dominant (14.5.3).

Au \S~15 on étudie diverses propriétés des fibres des morphismes universellement
ouverts $f:X\to Y$, localement de type fini, notamment lorsque $X$ et $Y$ sont
localement noethériens. On obtient ainsi en particulier un critère pour qu'un
point $x\in X$ n'appartienne qu'à \emph{une seule} composante irréductible de $X$,
en termes de propriétés de \emph{la fibre} $f^{-1}(f(x))$ de $x$~: il suffit que
$Y$ soit géométriquement unibranche (par exemple normal) au point $f(x)$ et que
$f^{-1}(f(x))$ soit géométriquement ponctuellement intègre au point $x$ (15.3.3)~;
si de plus $Y$ est localement intègre au point $f(x)$, $f$ est \emph{plat} au point
$x$ et $X$ localement intègre au point $x$. On étudie aussi la variation du
\emph{nombre géométrique de composantes connexes} d'une fibre $f^{-1}(y)$~; par
exemple, si $f$ est universellement ouvert et \emph{propre}, et les fibres de $f$
géométriquement réduites, ce nombre est \emph{localement constant} (15.5.7).
Enfin, lorsque $f$ admet une \emph{section} $g$ (ce qui sera le cas lorsque $X$ est
un $Y$-\emph{schéma en groupes}) et que pour tout $y\in Y$, on désigne par $X_y^0$
la composante connexe de la fibre $X_y=f^{-1}(y)$ au point $g(y)$ («~composante
neutre~» dans le cas des groupes), on étudie la réunion $X^0$ des $X_y^0$ pour
$y\in Y$, et on montre (15.6.4) que si $f$ est universellement ouvert et les fibres
$X_y$ géométriquement réduites, alors $X^0$ est un ensemble \emph{ouvert} dans $X$.

\subsection{Morphismes ouverts}
\label{subsection:IV.14.1-fr}

\begin{env}[14.1.1]
\label{IV.14.1.1-fr}
Rappelons (1.10.2) qu'une application continue $\psi:X\to Y$ est dite
\emph{ouverte en un point} $x\in X$ si l'image par $\psi$ de \emph{tout} voisinage
de $x$ dans $X$ est un voisinage de $\psi(x)$ dans $Y$.

On notera que cela \emph{n'implique pas} qu'il existe un système fondamental de
voisinages de $x$ dont les images soient \emph{ouvertes} dans $Y$.
\end{env}

\begin{proposition}[14.1.2]
\label{IV.14.1.2-fr}
Soient $X$, $Y$ deux espaces topologiques, $\psi:X\to Y$ une application
continue, $x$ un point de $X$, $y=\psi(x)$.
\begin{enumerate}[label=(\roman*)]
  \item Si $\psi$ est ouverte en $x$, alors pour toute partie $Y'$ de $Y$ contenant
  $\psi(x)$ la restriction $\psi^{-1}(Y')\to Y'$ de $\psi$ à $Y'$ est ouverte au
  point $x$.
  \item Supposons que $Y$ soit réunion d'une famille localement finie de parties
  fermées $(Y_i)$ et que pour tout $i$ tel que $\psi(x)\in Y_i$, la restriction
  $\psi^{-1}(Y_i)\to Y_i$ de $\psi$ soit ouverte au point $x$~; alors $\psi$ est
  ouverte au point $x$.
  \item Soient $\gamma:X'\to X$ une application continue, $x'$ un point de $X'$~;
  si l'application composée $\psi\circ\gamma:X'\to Y$ est ouverte au point $x'$,
  $\psi$ est ouverte au point $\gamma(x')$.
\end{enumerate}
\end{proposition}

\begin{proof}
Si $U$ est un voisinage de $x$, on a
$\psi(U\cap\psi^{-1}(Y'))=\psi(U)\cap Y'$~; d'où aussitôt (i). Pour prouver (ii),
notons qu'il y a un voisinage $W$ de $\psi(x)$ dans $Y$ ne rencontrant que les
parties fermées $Y_i$ en nombre \emph{fini} qui \emph{contiennent} $\psi(x)$~; donc
$W$ est réunion des $W\cap Y_i$ pour ces indices. Or, si $U$ est un voisinage de
$x$ tel que $\psi(U)\cap Y_i$ soit un voisinage de $\psi(x)$ dans $Y_i$, il existe
un voisinage $V\subset W$ de $\psi(x)$ dans $Y$ tel que pour tous les $i$ tels que
$\psi(x)\in Y_i$,
\oldpage[IV-3]{201}
on ait $V\cap Y_i\subset\psi(U)\cap Y_i$ et comme la réunion des $V\cap Y_i$ pour
ces indices est $V$, on a $V\subset\psi(U)$, donc $\psi(U)$ est un voisinage de
$\psi(x)$ dans $Y$. L'assertion (iii) est triviale.
\end{proof}

\begin{remarks}[14.1.3]
\label{IV.14.1.3-fr}
\begin{enumerate}[label=(\roman*)]
  \item L'ensemble $Z$ des points $x\in X$ où un morphisme est ouvert
  \emph{n'est pas nécessairement ouvert}. Par exemple, soient $K$ un corps, $A$
  l'anneau de polynômes $K[S,T]$, $V$ le plan affine $\operatorname{Spec}(A)$,
  $X$ le sous-préschéma fermé de $V$ «~réunion de la droite $X_1$ définie par
  $T=0$ et de la droite $X_2$ définie par $S=0$~», c'est-à-dire
  $\operatorname{Spec}(A/\mathfrak a)$, où $\mathfrak a=AST$~; prenons
  $Y=X_2=\operatorname{Spec}(A/AS)$ et pour $f$ la projection correspondant à
  l'injection canonique $K[T]\to A/\mathfrak a$~; alors on a $Z=X_2$, qui n'est
  pas ouvert dans $X$.

  \item Soient $X$, $Y$ deux préschémas noethériens irréductibles de points
  génériques $\xi$, $\eta$ respectivement, $f:X\to Y$ un morphisme
  \emph{dominant localement de type fini}~; alors $f$ \emph{est ouvert au point
  $\xi$}. En effet, on peut évidemment se borner au cas où $X$ et $Y$ sont réduits
  (donc intègres) (1.5.4) et affines~; en vertu du théorème de platitude générique
  (6.9.1), il existe un ouvert non vide $V$ de $Y$ tel que la restriction
  $f^{-1}(V)\to V$ de $f$ soit un morphisme \emph{plat}, et on conclut de (2.4.6)
  que cette restriction est un morphisme ouvert. Toutefois, comme il y a des
  morphismes dominants de type fini $f:X\to Y$ (où $X$ et $Y$ sont
  \emph{irréductibles}) qui ne sont pas ouverts (voir par exemple (\textbf{II},
  8.1)), l'ensemble des points où un tel morphisme est ouvert n'est pas
  nécessairement \emph{fermé} dans $X$.

  \item Nous ignorons si, lorsque $X$ et $Y$ sont localement noethériens et
  $f:X\to Y$ un morphisme localement de type fini, l'ensemble des points de $X$
  où $f$ est ouvert est ou non \emph{localement constructible}.
\end{enumerate}
\end{remarks}

\begin{proposition}[14.1.4]
\label{IV.14.1.4-fr}
Soient $X$, $Y$ deux espaces topologiques, $\psi:X\to Y$ une application
continue. Pour tout $y\in Y$, l'ensemble des $x\in\psi^{-1}(y)$ où $\psi$ est
ouvert est une partie fermée de $\psi^{-1}(y)$.
\end{proposition}

\begin{proof}
En effet, supposons que $\psi$ ne soit pas ouvert en un point
$x\in\psi^{-1}(y)$~; il existe alors un voisinage ouvert $V$ de $x$ dans $X$ tel
que $\psi(V)$ ne soit pas un voisinage de $y$~; il en résulte que pour tout
$x'\in V\cap\psi^{-1}(y)$, $\psi$ n'est pas ouvert au point $x'$.
\end{proof}

\begin{remark}[14.1.5]
\label{IV.14.1.5-fr}
Même si $X$ et $Y$ sont localement noethériens et $f$ un morphisme \emph{fini},
il peut se faire que $f$ soit ouvert en tous les points d'une fibre $f^{-1}(y)$,
sans qu'il existe un \emph{voisinage} de $f^{-1}(y)$ en tous les points duquel $f$
soit ouvert. Soient par exemple $K$ un corps, $A$ l'anneau de polynômes
$K[T_1,T_2,T_3]$, $V=\operatorname{Spec}(A)$ l'espace affine à 3 dimensions sur
$K$, $X=\operatorname{Spec}(A/\mathfrak a)$, où $\mathfrak a=\mathfrak b\mathfrak c$,
avec $\mathfrak b=(T_3)$ et $\mathfrak c=(T_1)+(T_2-T_3)$ dans $A$, de sorte que
$X$ est réunion du plan $X_1=\operatorname{Spec}(A/\mathfrak b)$ («~plan
d'équation $T_3=0$~») et de la droite $X_2=\operatorname{Spec}(A/\mathfrak c)$
(«~droite d'équations $T_1=0$, $T_2=T_3$~») qui sont ses composantes
irréductibles. Prenons $Y=X_1$ et soit $f:X\to Y$ la projection correspondant à
l'injection canonique $K[T_1,T_2]\to A/\mathfrak a$~; si $y$ est le point commun
à $X_1$ et $X_2$, $f^{-1}(y)$ est réduit à $y$ et $f$ est ouvert en ce point mais
n'est ouvert en aucun point de $X_2$ dans un voisinage de $y$ et distinct de $y$.
\end{remark}

\begin{env}[14.1.6]
\label{IV.14.1.6-fr}
Dans ce qui suit, le rôle essentiel sera joué par le critère (1.10.3)
caractérisant les morphismes \emph{localement de présentation finie} $f:X\to Y$
qui sont ouverts en un point $x$ par le «~relèvement des générisations~»~: pour
toute générisation $y'$ de $y=f(x)$, il existe $x'\in X$, générisation de $x$, tel
que $y'=f(x')$.
\end{env}

\oldpage[IV-3]{202}
\subsection{Morphismes ouverts et formule des dimensions}
\label{subsection:IV.14.2-fr}

\begin{theorem}[14.2.1]
\EGAReaderTarget{ega4.c04.thm.14-2-1}{IV.14.2.1-fr}
Soient $X$, $Y$ deux préschémas localement noethériens, $f:X\to Y$ un morphisme,
$x$ un point de $X$, $y=f(x)$. On suppose que $f$ est ouvert aux points génériques
des composantes irréductibles de $f^{-1}(y)$ contenant $x$. Alors on a la relation
\begin{equation}
\tag{14.2.1.1}
\label{IV.14.2.1.1-fr}
  \dim(\mathcal O_x)=\dim(\mathcal O_y)+
  \dim(\mathcal O_x\otimes_{\mathcal O_y}k(y)).
\end{equation}
\end{theorem}

\begin{proof}
Soit $y'$ une générisation quelconque de $y$ distincte de $y$, et considérons le
sous-préschéma fermé réduit $Y'$ de $Y$ ayant pour espace sous-jacent
$\overline{\{y'\}}$~; alors aucune composante irréductible $X'$ de $f^{-1}(Y')$
contenant $x$ ne peut être contenue dans $f^{-1}(y)$~: en effet, si $x'$ est le
point générique de $X'$, on aurait $f(x')=y\ne y'$, et comme $x'$ est sa seule
générisation dans $f^{-1}(Y')$, cela contredirait l'hypothèse que $f$ est ouvert
au point $x'$, en vertu du fait que $a)$ implique $c)$ dans (1.10.3). On peut donc
appliquer (6.1.2) à l'homomorphisme local d'anneaux noethériens
$\mathcal O_y\to\mathcal O_x$, d'où la conclusion.
\end{proof}

\begin{corollary}[14.2.2]
\label{IV.14.2.2-fr}
Soient $Y$ un préschéma localement noethérien, $f:X\to Y$ un morphisme localement
de type fini, $x$ un point de $X$, $y=f(x)$. Supposons $f$ ouvert aux points
génériques des composantes irréductibles de $f^{-1}(y)$ contenant $x$. Alors $f$
est équidimensionnel au point $x$ dans chacun des deux cas suivants~:
\begin{enumerate}[label=(\roman*)]
  \item $X$ est irréductible et $f$ est dominant.
  \item Les anneaux $\mathcal O_x$ et $\mathcal O_y$ sont équidimensionnels.
\end{enumerate}
\end{corollary}

\begin{proof}
Pour (i), cela résulte de (14.2.1) et de (13.2.3). Pour (ii), cela résulte de
(14.2.1) et de (13.3.6).
\end{proof}

\begin{proposition}[14.2.3]
\label{IV.14.2.3-fr}
Soient $Y$ un préschéma irréductible localement noethérien, $\eta$ son point
générique, $f:X\to Y$ un morphisme localement de type fini, $y$ un point de
$f(X)$, $Z$ une composante irréductible de $f^{-1}(y)$ telle que $f$ soit ouvert
au point générique de $Z$. Alors $Z$ est contenue dans une composante irréductible
$X'$ de $X$ dominant $Y$ et telle que
\[
  \dim(Z)=\dim(X'\cap f^{-1}(y))=\dim(X'\cap f^{-1}(\eta)).
\]
\end{proposition}

\begin{proof}
Cela résulte en effet de (14.2.1) appliqué au point générique de $Z$, et de
(13.2.7).
\end{proof}

\begin{corollary}[14.2.4]
\label{IV.14.2.4-fr}
Avec les notations de (14.2.3), supposons que $f$ soit ouvert en tous les points
de $f^{-1}(y)$. Pour tout $z\in Y$, soient $E(z)$ l'ensemble des dimensions des
composantes irréductibles de $f^{-1}(z)$ et
$d(z)=\sup E(z)=\dim(f^{-1}(z))$. On a alors $E(y)\subset E(\eta)$, d'où
$d(y)\leq d(\eta)$.
\end{corollary}

\begin{proof}
Cela résulte aussitôt de (14.2.3).
\end{proof}

\begin{corollary}[14.2.5]
\label{IV.14.2.5-fr}
Soient $Y$ un préschéma localement noethérien, $f:X\to Y$ un morphisme propre,
$y\in Y$ un point tel que $f$ soit ouvert en tous les points de $f^{-1}(y)$. Alors
la fonction $z\mapsto\dim(f^{-1}(z))$ est constante dans un voisinage de $y$.
\end{corollary}

\begin{proof}
Soient $Y_i$ ($1\leq i\leq n$) les sous-préschémas fermés réduits de $Y$ ayant
pour espaces sous-jacents les composantes irréductibles de $Y$ contenant $y$~;
si $f_i:f^{-1}(Y_i)\to Y_i$ est la restriction de $f$, on sait que chacun des
$f_i$ est propre (\textbf{II}, 5.4.5) et ouvert en tous les points de $f^{-1}(y)$
(14.1.2). Comme la réunion des $Y_i$ est un voisinage de $y$ dans $Y$,
\oldpage[IV-3]{203}
on voit qu'on peut se borner à démontrer le corollaire lorsque $Y$ est
irréductible~; soit $\eta$ son point générique. Il résulte de (14.2.4) que
$d(y)\leq d(\eta)$~; d'autre part, comme $f$ est propre, on déduit de (13.1.5) que
$z\mapsto\dim(f^{-1}(z))$ est semi-continue supérieurement~; en particulier
$d(y)\geq d(\eta)$, donc $d(y)=d(\eta)$. De plus, il y a un voisinage $V$ de $y$
tel que $d(z)\leq d(y)$ pour tout $z\in V$~; mais comme $z$ est spécialisation de
$\eta$, on a d'autre part $d(z)\geq d(\eta)$, d'où finalement $d(z)=d(y)$ pour
$z\in V$.
\end{proof}

\begin{corollary}[14.2.6]
\label{IV.14.2.6-fr}
Soient $Y$ un préschéma localement noethérien, $f:X\to Y$ un morphisme ouvert de
type fini, $y$ un point de $Y$. Soient $y_i$ ($1\leq i\leq n$) les points
génériques des composantes irréductibles de $Y$ contenant $y$. Alors, avec les
notations de (14.2.4)~:
\begin{enumerate}[label=(\roman*)]
  \item Si pour $1\leq i\leq n$, on a $E(y)=E(y_i)$ (resp. $d(y)=d(y_i)$), la
  fonction $E$ (resp. $d$) est constante dans un voisinage de $y$.
  \item Il existe un voisinage ouvert $U$ de $f^{-1}(y)$ tel que la fonction
  $z\mapsto\dim(U\cap f^{-1}(z))$ soit constante dans un voisinage de $y$.
\end{enumerate}
\end{corollary}

\begin{proof}
\begin{enumerate}[label=(\roman*)]
  \item Le même raisonnement que dans (14.2.5) montre qu'on peut se borner à
  considérer le cas où $Y$ est irréductible, de point générique $\eta$. Si $y'$
  est une générisation quelconque de $y$, on a alors (en appliquant (14.2.4) à la
  restriction $f^{-1}(Y')\to Y'$ de $f$, où $Y'=\overline{\{y'\}}$, et utilisant
  (14.1.2)) $E(y)\subset E(y')\subset E(\eta)$ et
  $d(y)\leq d(y')\leq d(\eta)$~; cela montre que la relation $E(y)=E(\eta)$
  (resp. $d(y)=d(\eta)$) entraîne $E(y)=E(y')$ (resp. $d(y)=d(y')$) pour toute
  générisation $y'$ de $y$. Or, en vertu de (9.5.5), les fonctions $E$ et $d$ sont
  localement constructibles, donc il résulte de ($\mathbf{0}_{\mathrm{III}}$,
  9.2.5) appliqué à l'ensemble des $z$ tels que $E(z)=E(y)$ (resp.
  $d(z)=d(y)$) que cet ensemble est un voisinage de $y$.

  \item Pour chacun des $y_i$, $f^{-1}(y_i)$ est un espace noethérien puisque $f$
  est de type fini~; soient $x_{ij}$ ($1\leq j\leq m_i$) les points génériques de
  ses composantes irréductibles, et posons $X'_{ij}=\overline{\{x_{ij}\}}$~;
  montrons que le complémentaire $U$ dans $X$ de la réunion des $X'_{ij}$ qui ne
  rencontrent pas $f^{-1}(y)$ (qui est évidemment un voisinage ouvert de
  $f^{-1}(y)$) répond à la question. En effet, la restriction $f|U:U\to Y$ de $f$
  est un morphisme ouvert de type fini (I, 6.3.5)~; pour tout couple $(i,j)$ tel
  que $x_{ij}\in U$, $x_{ij}$ est point maximal de $U\cap f^{-1}(y_i)$, et il
  résulte de (13.1.1) appliqué à la restriction
  $X'_{ij}\to Y_i=\overline{\{y_i\}}$ de $f$ (compte tenu de (I, 5.2.2) et de
  (1.5.4)) que toutes les composantes irréductibles de
  $X'_{ij}\cap f^{-1}(y)$ ont une dimension
  $\geq\dim(X'_{ij}\cap f^{-1}(y_i))$. Compte tenu de (14.2.3) appliqué à la
  restriction $f^{-1}(Y_i)\to Y_i$ de $f$, qui est un morphisme ouvert (14.1.2),
  $f^{-1}(y)$ est, pour chaque $i$, la réunion des composantes irréductibles de
  $X'_{ij}\cap f^{-1}(y)$ ($1\leq j\leq m_i$), donc
  $\dim(f^{-1}(y))\geq\dim(U\cap f^{-1}(y_i))$~; mais comme $f|U$ est ouvert, on a
  aussi $\dim(f^{-1}(y))\leq\dim(f^{-1}(y_i)\cap U)$ en vertu de (14.2.4)~; on voit
  donc qu'on peut appliquer (i) à $f|U$, d'où la conclusion.
\end{enumerate}
\end{proof}

\begin{remark}[14.2.7]
\label{IV.14.2.7-fr}
L'exemple (13.2.12, (iii)) montre que sous les hypothèses de (14.2.5) ou
(14.2.6, (ii)), on ne peut, dans la conclusion, remplacer la fonction
$z\mapsto\dim(f^{-1}(z))$ par la fonction $z\mapsto E(z)$~; en effet, dans cet
exemple, $f$ est propre et plat (donc universellement ouvert (2.4.6)) et tout
voisinage de l'unique point de $f^{-1}(y_0)$ contient les points génériques des
composantes irréductibles de $f^{-1}(\eta)$.
\end{remark}

\oldpage[IV-3]{204}
\subsection{Morphismes universellement ouverts}
\label{subsection:IV.14.3-fr}

\begin{env}[14.3.1]
\label{IV.14.3.1-fr}
Rappelons (2.4.2) que dire qu'un morphisme $f:X\to Y$ est universellement ouvert
signifie que pour tout morphisme $g:Y'\to Y$, $f_{(Y')}:X_{(Y')}\to Y'$ est
ouvert~; il suffit d'ailleurs qu'il en soit ainsi lorsque
$Y'=Y[T_1,\ldots,T_e]$, pour tout $e$ (8.10.2) (et si $Y$ est localement
noethérien, il suffit donc qu'il en soit ainsi pour tout $Y'$ localement
noethérien).
\end{env}

\begin{proposition}[14.3.2]
\label{IV.14.3.2-fr}
Soit $f:X\to Y$ un morphisme de préschémas. Si $f$ est universellement ouvert,
alors, pour tout morphisme $g:Y'\to Y$, où $Y'$ est irréductible, l'image par
$f'=f_{(Y')}:X_{(Y')}\to Y'$ de toute composante irréductible de $X_{(Y')}$ est
dense dans $Y'$. Réciproquement, si cette condition est vérifiée pour tout $Y'$
irréductible et tout morphisme de type fini $g:Y'\to Y$, et si de plus $f$ est
localement de présentation finie, alors $f$ est universellement ouvert.
\end{proposition}

\begin{proof}
En effet, il résulte de (1.10.4) que si $f$ est universellement ouvert, il vérifie
la condition de l'énoncé. Inversement, supposons que $f$ soit localement de
présentation finie, et montrons que pour tout entier $e$, si l'on pose
$Y''=Y[T_1,\ldots,T_e]$, $f_{(Y'')}$ est ouvert. En effet, soit $Y'$ un
sous-préschéma fermé de $Y''$ ayant pour espace sous-jacent une partie fermée
irréductible de $Y''$~; le morphisme composé $Y'\to Y''\to Y$ est de type fini,
donc toute composante irréductible de $X_{(Y')}$ domine $Y'$ par hypothèse~; on
déduit donc de (1.10.4) que $f_{(Y'')}$ est un morphisme ouvert.
\end{proof}

Cette proposition montre que la définition (\textbf{III}, 4.3.9) coïncide dans
le cas considéré avec la définition générale des morphismes universellement
ouverts donnée dans (2.4.2).

A la notion de morphisme ouvert en un point (14.1.1) correspond de même la
suivante~:

\begin{definition}[14.3.3]
\label{IV.14.3.3-fr}
Soient $f:X\to Y$ un morphisme de préschémas, $x$ un point de $X$. On dit que $f$
est \emph{universellement ouvert au point $x$} si, pour tout morphisme $g:Y'\to Y$,
en posant $X'=X\times_Y Y'$, le morphisme $f'=f_{(Y')}:X'\to Y'$ est ouvert en
tout point $x'$ de $X'$ dont la projection dans $X$ est $x$.
\end{definition}

\begin{remarks}[14.3.3.1]
\label{IV.14.3.3.1-fr}
\begin{enumerate}[label=(\roman*)]
  \item Le raisonnement de (8.10.1) montre (avec les mêmes notations) que si $x$
  est un point de $X$, $x_\lambda$ sa projection dans $X_\lambda$, alors, si
  $f_\mu$ est ouvert au point $x_\mu$ pour tout $\mu\geq\lambda$, $f$ est ouvert
  au point $x$~; il suffit de se borner aux ouverts $U$ de $X$ contenant $x$ et de
  remarquer que l'hypothèse implique que $f_\mu(U_\mu)$ est un voisinage de
  $f_\mu(x_\mu)$, donc $f(v_\mu^{-1}(U_\mu))$ est un voisinage de $f(x)$. On en
  déduit que l'énoncé (8.10.2) est encore exact quand on y remplace
  «~universellement ouvert~» par «~universellement ouvert au point $x$~», et
  «~morphisme ouvert~» par «~morphisme ouvert en tout point $x_n$ de $X_n$ dont
  la projection dans $X$ est $x$~»~: il suffit dans la démonstration de se limiter
  aux ouverts $V$ contenant un $x_n$.

  \item Le résultat de (14.1.4) est encore valable pour un morphisme $f:X\to Y$,
  en y remplaçant «~ouvert~» par «~universellement ouvert~». En effet, supposons
  que $f$ ne soit pas universellement ouvert en un point $x\in f^{-1}(y)$~; il y a
  par suite un morphisme $Y'\to Y$ et un point $x'\in X'$ se projetant en $x$, tels
  que $f'$ ne soit pas ouvert au point $x'$.
\end{enumerate}
\end{remarks}

\oldpage[IV-3]{205}

Or, si $y'=f'(x')$, la projection $f'^{-1}(y')\to f^{-1}(y)$ est un morphisme
ouvert (2.4.10), et il y a, en vertu de (14.1.4), un voisinage $V'$ de $x'$ dans
$f'^{-1}(y')$ où $f'$ n'est pas ouvert~; donc $f$ n'est pas universellement ouvert
aux points de l'image de $V'$ dans $f^{-1}(y)$, qui est un voisinage de $x$ dans
$f^{-1}(y)$.

\begin{proposition}[14.3.4]
\label{IV.14.3.4-fr}
\begin{enumerate}[label=(\roman*)]
  \item Soient $f:X\to Y$, $g:Y\to Z$ deux morphismes, $x$ un point de $X$,
  $y=f(x)$. Si $f$ est universellement ouvert au point $x$ et $g$ universellement
  ouvert au point $y$, alors $g\circ f$ est universellement ouvert au point $x$.
  Inversement, si $g\circ f$ est universellement ouvert au point $x$, $g$ est
  universellement ouvert au point $y$.

  \item Si $f:X\to Y$ est un $S$-morphisme universellement ouvert au point
  $x\in X$, alors, pour tout changement de base $S'\to S$,
  $f_{(S')}:X_{(S')}\to Y_{(S')}$ est universellement ouvert en tout point de
  $X_{(S')}$ au-dessus de $x$.

  \item Pour que $f:X\to Y$ soit universellement ouvert au point $x\in X$, il
  faut et il suffit que $f_{\mathrm{red}}$ le soit.

  \item Soient $f:X\to Y$ un morphisme localement de présentation finie, $x$ un
  point de $X$, $y=f(x)$~; posons $Y_1=\operatorname{Spec}(\mathcal O_y)$,
  $X_1=X\times_Y Y_1$, $f_1=f_{(Y_1)}$. Pour que $f$ soit universellement ouvert
  au point $x$, il faut et il suffit que $f_1$ le soit (on rappelle (I, 3.6.5) que
  $X_1$ est canoniquement identifié à un sous-espace de $X$).
\end{enumerate}
\end{proposition}

\begin{proof}
En effet (ii) est une conséquence évidente de la définition (14.3.3)~; il résulte
aussi de la définition que pour démontrer l'assertion (i), il suffit de le faire
lorsqu'on y supprime partout le mot «~universellement~», et cela résulte alors de
(14.1.2). L'assertion (iii) résulte de (ii) et de ce que le morphisme canonique
$X_{\mathrm{red}}\to X$ est surjectif. Enfin, la condition (iv) est trivialement
nécessaire. D'autre part, si elle est vérifiée, et si $g:Y'\to Y$ est un morphisme
quelconque, $X'=X_{(Y')}$, $f'=f_{(Y')}$, $x'$ un point de $X'$ au-dessus de $x$,
alors $f'$ est localement de présentation finie, et pour voir qu'il est ouvert au
point $x'$, il suffit d'appliquer le critère (1.10.3, $b$). Soient
$y'=f'(x')$, $Y'_1=\operatorname{Spec}(\mathcal O_{y'})$,
$X'_1=X'\times_{Y'}Y'_1$, $f'_1=f_{(Y'_1)}$. Comme $g(y')=y$, le morphisme
composé $Y'_1\to Y'\to Y$ se factorise en $Y'_1\to Y_1\to Y$ (I, 2.4.4), donc
$X'_1=X_1\times_{Y_1}Y'_1$ et $f'_1=(f_1)_{(Y'_1)}$~; la conclusion résulte alors
aussitôt de l'hypothèse que $f'_1$ est ouvert au point $x'$ et de (1.10.3).
\end{proof}

\begin{proposition}[14.3.5]
\label{IV.14.3.5-fr}
Soient $X$, $Y$ deux préschémas, $f:X\to Y$ un morphisme, $x$ un point de $X$.
Soit $(Y_i)_{1\leq i\leq n}$ une famille localement finie de sous-préschémas
fermés de $Y$ telle que l'espace $Y$ soit réunion des $Y_i$, et supposons que pour
tout $i$ tel que $f(x)\in Y_i$, la restriction $f^{-1}(Y_i)\to Y_i$ de $f$ soit un
morphisme universellement ouvert au point $x$~; alors $f$ est universellement
ouvert au point $x$.
\end{proposition}

\begin{proof}
Compte tenu de la définition, cela résulte de la proposition analogue (14.1.2,
(ii)) pour les morphismes ouverts en un point.
\end{proof}

\begin{proposition}[14.3.6]
\label{IV.14.3.6-fr}
Soient $Y$ un préschéma localement noethérien, $f:X\to Y$ un morphisme localement
de type fini. Pour que $f$ soit universellement ouvert en un point $x\in X$, il
faut et il suffit que la condition suivante soit vérifiée~: pour tout morphisme
$g:Y'\to Y$, où $Y'=\operatorname{Spec}(A)$ est le spectre d'un anneau de
valuation discrète, tel que l'image $g(y')$ du point fermé $y'$ de $Y'$ soit égale
à $y=f(x)$, et pour tout point $x'\in X'=X\times_Y Y'$ dont les projections sur
$X$ et $Y'$ sont $x$ et $y'$ respectivement, il existe une générisation $z'$ de
$x'$ dans $X'$ dont la projection dans $Y'$ est le point générique de $Y'$
(autrement dit, il existe une composante irréductible de $X'$ contenant $x'$ et
dominant $Y'$).

\oldpage[IV-3]{206}

De plus, on peut, dans la condition précédente, se borner au cas où $A$ est
complet, a un corps résiduel algébriquement clos, et où $x'$ est rationnel sur
$k(y')$.
\end{proposition}

\begin{proof}
Si $Y'$ est comme dans l'énoncé, la nécessité de la condition résulte de ce que
$f'$ doit être ouvert au point $x'$, et du critère (1.10.3). Pour voir que la
condition est suffisante, considérons un morphisme de type fini $g:Y''\to U$, et
soient $X''=X\times_Y Y''$, $f''=f_{(Y'')}:X''\to Y''$, et $x''$ un point de $X''$
au-dessus de $x$. Posons $y''=f''(x'')$, et soit $t$ une générisation de $y''$
dans $Y''$, distincte de $y''$. Puisque $Y''$ est localement noethérien, il
résulte de (\textbf{II}, 7.1.9) et ($\mathbf{0}_{\mathrm{III}}$, 10.3.1) qu'il
existe un schéma $Y'=\operatorname{Spec}(A)$, où $A$ est un anneau de valuation
discrète complet, dont le corps résiduel est une clôture algébrique de $k(x'')$,
et un morphisme $g:Y'\to Y''$ tels que, si $s$ et $y'$ sont le point générique et
le point fermé de $Y'$, on ait $g(s)=t$ et $g(y')=y''$. Il y a alors un point $x'$
de $X'=X\times_Y Y'=X''\times_{Y''}Y'$ dont les projections dans $X''$ et $Y'$
soient $x''$ et $y'$ respectivement, et qui soit rationnel sur $k(y')$ (I, 3.4.9).
L'hypothèse entraîne qu'il y a une générisation $z'$ de $x'$ dans $X'$ dont la
projection dans $Y'$ soit $s$~; si $z''$ est la projection de $z'$ dans $X''$,
$z''$ est une générisation de $x''$ et sa projection dans $Y''$ est $t$~; on
conclut donc de (1.10.3) que $f''$ est ouvert au point $x''$, donc que $f$ est
universellement ouvert au point $x$ (14.3.3.1, (i)).
\end{proof}

\begin{corollary}[14.3.7]
\label{IV.14.3.7-fr}
Les notations étant celles de (14.3.6)~:
\begin{enumerate}[label=(\roman*)]
  \item Étant donné un point $y\in Y$, pour que $f$ soit universellement ouvert en
  tous les points de $f^{-1}(y)$, il faut et il suffit que pour tout morphisme
  $g:Y'\to Y$, où $Y'$ est le spectre d'un anneau de valuation discrète, et où
  l'image $g(y')$ du point fermé $y'$ de $Y'$ est $y$, toute composante
  irréductible de $X'=X\times_Y Y'$ domine $Y'$.
  \item Pour que $f$ soit universellement ouvert, il faut et il suffit que pour
  tout morphisme $g:Y'\to Y$, où $Y'$ est le spectre d'un anneau de valuation
  discrète, toute composante irréductible de $X'=X\times_Y Y'$ domine $Y'$.
\end{enumerate}
\end{corollary}

\begin{proof}
Il est clair qu'il suffit de prouver (i)~; la nécessité de (i) résulte de (14.3.2)
et sa suffisance de (14.3.6).
\end{proof}

\begin{proposition}[14.3.8]
\EGAReaderTarget{ega4.c04.prop.14-3-8}{IV.14.3.8-fr}
Soient $Y$ un préschéma localement noethérien, irréductible, régulier et de
dimension 1 (par exemple le spectre d'un anneau de Dedekind), $f:X\to Y$ un
morphisme localement de type fini, $y$ un point de $Y$. Les conditions suivantes
sont équivalentes~:
\begin{enumerate}[label=\emph{\alph*)}]
  \item $f_{\mathrm{red}}$ est plat en tout point de $f^{-1}(y)$.
  \item $f$ est universellement ouvert dans un voisinage de $f^{-1}(y)$.
  \item $f$ est ouvert dans un voisinage de $f^{-1}(y)$.
  \item Toute composante irréductible de $X$ rencontrant $f^{-1}(y)$ domine $Y$.
\end{enumerate}
\end{proposition}

\begin{proof}
Comme $f_{\mathrm{red}}$ est localement de type fini (1.3.4), $a)$ entraîne que
$f_{\mathrm{red}}$ est plat dans un voisinage de $f^{-1}(y)$ (11.1.1), et il
suffit d'appliquer (2.4.6) dans un tel voisinage pour voir que $a)$ entraîne $b)$.
L'implication $b)\Rightarrow c)$ est triviale, et l'implication
$c)\Rightarrow d)$ résulte de (1.10.4) appliqué à un voisinage de $f^{-1}(y)$.
Reste à voir que $d)$ entraîne $a)$. On peut évidemment, en vertu de (1.3.4), se
borner au cas où $X$ est réduit. La question étant en outre locale sur $X$ et sur
$Y$, on peut supposer $Y=\operatorname{Spec}(A)$ et
$X=\operatorname{Spec}(B)$
\oldpage[IV-3]{207}
affines~; si $X_i$ ($1\leq i\leq n$) sont les sous-préschémas fermés (intègres)
de $X$ ayant pour espaces sous-jacents les composantes irréductibles de $X$,
alors, pour tout $x\in X$, $\mathcal O_{X,x}$, étant réduit, est un sous-anneau du
composé direct des $\mathcal O_{X_i,x}$~; si $y=f(x)$, il suffira de montrer que
chacun des $\mathcal O_{X_i,x}$ est un $\mathcal O_y$-module \emph{sans torsion},
car il en sera alors de même de $\mathcal O_{X,x}$~; comme par hypothèse
$\mathcal O_y$ est un anneau local régulier de dimension 1, c'est-à-dire un anneau
de valuation discrète (\textbf{II}, 7.1.6), il résultera alors de
($\mathbf{0}_{\mathrm I}$, 6.3.4) que $\mathcal O_{X,x}$ est un
$\mathcal O_y$-module \emph{plat}. Mais si $X_i=\operatorname{Spec}(B_i)$, où
$B_i$ est un anneau intègre, l'hypothèse $d)$ entraîne que l'homomorphisme
$A\to B_i$ est \emph{injectif} (I, 1.2.7)~; donc $B_i$ est un $A$-module sans
torsion, et \emph{a fortiori} $\mathcal O_{X_i,x}$ est un $\mathcal O_y$-module
sans torsion.
\end{proof}

\begin{remarks}[14.3.9]
\label{IV.14.3.9-fr}
\begin{enumerate}[label=(\roman*)]
  \item Dans l'énoncé de (14.3.8), on ne peut se dispenser de l'hypothèse que $Y$
  est \emph{régulier}. Avec les notations de (11.7.5), prenons en effet
  $Y=\operatorname{Spec}(A)$, $X=\operatorname{Spec}(\overline A)$, de sorte que
  $f:X\to Y$ est un morphisme \emph{fini surjectif}~; comme $A$ est un anneau
  local intègre de dimension 1, ainsi que $\overline A$, il résulte aussitôt de
  (1.10.4) que le morphisme $f$ est \emph{ouvert}. Cependant $f$ n'est pas
  universellement ouvert (ni \emph{a fortiori} plat), comme le montre (11.7.5).
  On aurait un exemple analogue en prenant pour $Y$ le schéma local au point
  double d'une courbe algébrique ayant un «~point double ordinaire~» et pour $X$
  le normalisé de $Y$.

  \item L'exemple de (14.1.3, (i)) montre que l'ensemble des points de $X$ où un
  morphisme est universellement ouvert n'est pas nécessairement ouvert, le
  morphisme $f$ dans cet exemple étant universellement ouvert en tous les points
  où il est ouvert. L'exemple $f:X\to Y$ vu ci-dessus en (i) montre de même que
  l'ensemble des points où un morphisme est universellement ouvert n'est pas
  nécessairement fermé, car il est immédiat qu'en tous les points de $X$ sauf un
  le morphisme $f$ est universellement ouvert (c'est même un isomorphisme local).
  Il serait intéressant de savoir si l'ensemble des points où un morphisme est
  universellement ouvert est localement constructible.
\end{enumerate}
\end{remarks}

Les deux propositions suivantes nous ont été signalées par M.~Artin~:

\begin{lemma}[14.3.10]
\label{IV.14.3.10-fr}
Soient $A$ un anneau de valuation (non nécessairement discrète), $k$ son corps
résiduel, $K$ son corps des fractions, et posons $S=\operatorname{Spec}(A)$. Soient
$X$ un préschéma irréductible, $f:X\to S$ un morphisme dominant de type fini~; soit
$X_0=X\otimes_A k$ (resp. $X_1=X\otimes_A K$) la fibre de $f$ au point fermé
(resp. au point générique) de $S$. Alors, si $X_0\ne\varnothing$, on a
$\dim(X_0)=\dim(X_1)$.
\end{lemma}

On peut se borner au cas où $X$ est affine, en remplaçant au besoin $X$ par un
ouvert affine contenant un point générique d'une composante irréductible de $X_0$,
de dimension maximale, et utilisant (4.1.1.3). Soit $n=\dim(X_0)\geq 0$~; il résulte
de (13.3.1.1) qu'il existe un voisinage $U$ de $X_0$ dans $X$ et un $S$-morphisme
quasi-fini
\[
  g:U\to S[T_1,\ldots,T_n]=Z
\]
tel que le morphisme restriction
$g_0:U_0=U\cap X_0\to Z_0=\operatorname{Spec}(k[T_1,\ldots,T_n])$ soit fini et
\emph{surjectif}. Par changement de base $\operatorname{Spec}(K)\to S$ et restriction
à l'ouvert $U_1=U\cap X_1$ de $X_1$, on déduit de $g$ un morphisme quasi-fini
$g_1:U_1\to Z_1=\operatorname{Spec}(K[T_1,\ldots,T_n])$. Comme $U_1$ est dense dans
$X_1$, on a $\dim(U_1)=\dim(X_1)$ (4.1.1.3)~; la proposition sera établie, en vertu
de (4.1.2), si nous prouvons que le morphisme $g_1$ est \emph{dominant}. Supposons
le contraire~; il existerait alors un polynôme $F_1\ne 0$ dans
$K[T_1,\ldots,T_n]$ tel que $g_1(U_1)\subset V(F_1)$. Si $\omega$ est une valuation
sur $K$ associée à $A$, et si $(a_\alpha)$ est la famille des coefficients de $F_1$,
on peut, après multiplication de $F_1$ par un élément $\ne 0$ de $K$, supposer que
l'on a $\inf_\alpha(\omega(a_\alpha))=0$~; autrement dit, $F_1$ provient d'un
polynôme $F\in A[T_1,\ldots,T_n]$ (avec lequel il s'identifie) tel que l'image $F_0$
de $F$ dans $k[T_1,\ldots,T_n]$ soit \emph{non nul}. Considérons alors dans $Z$
l'ensemble fermé $V(F)$~; on a $V(F)\cap Z_1=V(F_1)$, et comme $U_1$ est dense
dans $U$ (puisqu'il contient le point générique de $X$), $g(U)\subset V(F)$~; en
particulier, on aurait $g_0(U_0)\subset V(F)\cap Z_0=V(F_0)$~; mais puisque
$F_0\ne 0$, $V(F_0)$ est une partie fermée de $Z_0$ \emph{distincte de $Z_0$} et
l'on aboutit à une contradiction. C.Q.F.D.

\oldpage[IV-3]{208}

\begin{proposition}[14.3.11]
\label{IV.14.3.11-fr}
Soient $f:X\to S$ un morphisme de type fini, $(g_\lambda)_{\lambda\in L}$ une
famille de morphismes universellement ouverts $g_\lambda:Y_\lambda\to S$, et pour
tout $\lambda\in L$, soit $u_\lambda:Y_\lambda\to X$ un $S$-morphisme. Pour tout
$s\in S$, posons $X_s=X\times_S\operatorname{Spec}(k(s))$,
$(Y_\lambda)_s=Y_\lambda\times_S\operatorname{Spec}(k(s))$, et soit
$(u_\lambda)_s:(Y_\lambda)_s\to X_s$ le morphisme $u_\lambda\times 1$ déduit de
$u_\lambda$ par changement de base. Soit $Z(s)$ l'adhérence dans $X_s$ de la
réunion des ensembles $(u_\lambda)_s((Y_\lambda)_s)$, et posons
$d(s)=\dim(Z(s))$. Alors, pour toute générisation $s'$ d'un point $s\in S$, on a
$d(s)\leq d(s')$.
\end{proposition}

\begin{proof}
On sait (\textbf{II}, 7.1.4) qu'il existe un anneau de valuation $A'$ et un
morphisme $h:S'=\operatorname{Spec}(A')\to S$ tels que si $a$ (resp. $b$) est le
point fermé (resp. le point générique) de $S'$, on ait $h(a)=s$, $h(b)=s'$. En
outre, le morphisme projection
$p:X_s\otimes_{k(s)}k(a)\to X_s$ est surjectif et ouvert (2.4.10), donc fait de
$X_s$ un espace quotient de $X_s\otimes_{k(s)}k(a)$ par une relation d'équivalence
ouverte~; pour toute partie $M$ de $X_s$, $p^{-1}(\overline M)$ est donc égal à
l'adhérence $\overline{p^{-1}(M)}$ (Bourbaki, \emph{Top. gén.}, chap.~I\textsuperscript{er},
4\textsuperscript{e} éd., \S~5, n\textsuperscript{o}~3, prop.~7)~; on raisonne de
même pour $X_{s'}$, et compte tenu de (I, 3.4.8), de (4.2.7) et du fait que les
$g_\lambda$ sont universellement ouverts, on voit qu'on peut se ramener à prouver
la proposition sur la situation obtenue après changement de base $S'\to S$.
Supposons donc $S'=S$, $s$ étant le point fermé et $s'$ le point générique de $S$.
L'hypothèse que $g_\lambda$ est ouvert entraîne que toute composante irréductible
de $Y_\lambda$ domine $S$ (1.10.4), donc que son point générique est un point
maximal de $(Y_\lambda)_{s'}$~; si $Z$ désigne l'adhérence dans $X$ de $Z(s')$, on
a donc $u_\lambda(Y_\lambda)\subset Z$, et par suite
$(u_\lambda)_s((Y_\lambda)_s)\subset Z_s=Z\cap X_s$~; autrement dit, on a
$Z(s)\subset Z_s$, d'où $\dim(Z(s))\leq\dim(Z_s)$. Mais en appliquant (14.3.10) à
un sous-préschéma réduit de $X$ ayant pour espace sous-jacent une composante
irréductible de $Z$, on obtient $\dim(Z_s)\leq\dim(Z_{s'})$ (l'inégalité provenant
de ce qu'il peut y avoir des composantes irréductibles de $Z$ ne rencontrant pas
$X_s$). D'autre part, comme $Z(s')$ est par définition fermé dans $X_{s'}$, on a
$Z_{s'}=Z(s')$, ce qui achève la démonstration.
\end{proof}

\begin{remark}[14.3.12]
\label{IV.14.3.12-fr}
Le cas envisagé par M.~Artin était celui où $Y_\lambda=S$ pour tout $\lambda$,
autrement dit le cas où $(u_\lambda)$ est une famille de $S$-sections de $X$. Un
autre cas utile est celui où la famille $(u_\lambda)$ est réduite à un seul
élément~; on peut d'ailleurs toujours se ramener à ce cas en considérant le
préschéma $Y$ somme des $Y_\lambda$ et les morphismes $g:Y\to S$ et $u:Y\to X$
dont les restrictions à chaque $Y_\lambda$ sont respectivement $g_\lambda$ et
$u_\lambda$.
\end{remark}

\begin{proposition}[14.3.13]
\label{IV.14.3.13-fr}
Soient $f:X\to Y$ un morphisme localement de type fini, $y$ un point de $Y$, $x$
un point maximal de la fibre $X_y=X\times_Y\operatorname{Spec}(k(y))$.
Considérons les conditions suivantes~:
\begin{enumerate}[label=\emph{\alph*)}]
  \item $f$ est universellement ouvert au point $x$ (ou encore en tout point de la
  composante irréductible de $X_y$ de point générique $x$ (14.3.3.1, (ii))).
  \item Pour toute composante irréductible $Y_0$ de $Y$ contenant $y$, il existe
  une composante irréductible $Z$ de $X$ contenant $x$, dominant $Y_0$ et telle que
  $\dim_x(X_y)=\dim_x(Z\cap X_y)=\dim(Z\cap X_\eta)$, où $\eta$ est le point
  générique de $Y_0$ (ce qui entraîne que $Z$ est équidimensionnelle sur $Y_0$ au
  point $x$ (13.2.2)).
\end{enumerate}

\oldpage[IV-3]{209}

\begin{enumerate}[label=\emph{\alph*$'$)}]
  \item[b$'$)] Pour tout voisinage ouvert $U$ de $x$ dans $X$ et toute
  générisation $y'$ de $y$, on a
  $\dim(U\cap X_{y'})\geq\dim_x(U\cap X_y)$.
\end{enumerate}
Alors on a les implications $a)\Rightarrow b)\Leftrightarrow b')$.
\end{proposition}

\begin{proof}
Pour montrer que $b)$ implique $b')$, il suffit de remarquer que $y'$ appartient à
une composante irréductible $Y_0$ de $Y$ contenant $y$, de point générique $\eta$~;
prenant $Z$ comme dans $b)$ et notant que le point générique de $Z$ (qui est aussi
celui de $Z\cap X_\eta$ ($\mathbf{0}_{\mathrm I}$, 2.1.8)) est contenu dans $U$, on
a $\dim(U\cap X_{y'})\geq\dim(U\cap Z\cap X_{y'})$, et, en vertu de (13.1.6),
$\dim(U\cap Z\cap X_{y'})\geq\dim(Z\cap X_\eta)$~; d'où l'assertion, puisque
$\dim_x(U\cap X_y)=\dim_x(X_y)$ (4.1.1.3).

Pour prouver que $b')$ implique $b)$, on peut d'abord remplacer $X$ par
$f^{-1}(Y_0)$, donc supposer $Y_0=Y$~; on peut se borner au cas où $X$ et $Y$ sont
affines, puisque $\dim_x(U\cap X_y)=\dim_x(X_y)$ (4.1.1.3). Les composantes
irréductibles $W_i$ du préschéma noethérien $X_\eta$ sont alors en nombre fini, et
le complémentaire de la réunion des $\overline W_i$ (adhérences dans $X$) qui ne
contiennent pas $x$ est un voisinage ouvert $V$ de $x$. En remplaçant $X$ par $V$,
on peut donc supposer que $x\in\overline W_i$ pour tout $i$ (l'hypothèse $b')$
entraîne $\dim(V\cap X_\eta)\geq 0$, donc $x$ appartient à un des $\overline W_i$
pour un $i$ au moins). En outre, les $\overline W_i$ sont exactement les
composantes irréductibles de $X$ qui dominent $Y$ ($\mathbf{0}_{\mathrm I}$,
2.1.8)~; cela étant, si l'on avait, pour chacune de ces composantes $Z$,
$\dim(Z\cap X_\eta)<\dim_x(X_y)$, on en conclurait
$\dim(X_\eta)<\dim_x(X_y)$, contrairement à l'hypothèse $b')$. La relation
$\dim_x(Z\cap X_y)=\dim(Z\cap X_\eta)$ résulte alors de (13.1.6).

Reste à démontrer que $a)$ entraîne $b')$. Tenant compte de (\textbf{II}, 7.1.4)
et de l'invariance des hypothèses et de la conclusion par changement de base (en
vertu de (4.2.7)), on peut se borner au cas où $Y$ est un spectre d'anneau de
valuation, de point fermé $y$ et de point générique $y'$, et où $U=X$.
L'hypothèse que $f$ est ouvert au point $x$ entraîne qu'il existe une composante
irréductible $Z$ de $X$ contenant $x$ et dominant $Y$ (1.10.3). Appliquant
(14.3.10) à un voisinage de $x$ dans $Z$, on en conclut que
$\dim(Z\cap X_{y'})=\dim_x(Z\cap X_y)$~; mais comme $x$ est maximal dans $X_y$,
$Z\cap X_y$ contient la composante irréductible de $X_y$ de point générique $x$,
donc $\dim_x(Z\cap X_y)=\dim_x(X_y)$~; d'autre part, on a
$\dim(X_{y'})\geq\dim(Z\cap X_{y'})$, ce qui achève de prouver $b')$.
\end{proof}

\begin{remark}[14.3.14]
\label{IV.14.3.14-fr}
Nous ignorons si dans (14.3.13), la conclusion subsiste lorsqu'on remplace
l'hypothèse $a)$ par l'hypothèse plus faible que $f$ est ouvert au point $x$. On
peut montrer aisément qu'il suffirait de traiter le cas où $Y$ est spectre d'un
anneau local intègre, dont le point générique est isolé, et où $X$ est un
sous-préschéma fermé du fibré vectoriel $Y[T]$.
\end{remark}

\subsection{Le critère de Chevalley pour les morphismes universellement ouverts}
\label{subsection:IV.14.4-fr}

\begin{theorem}[14.4.1]
\EGAReaderTarget{ega4.c04.thm.14-4-1}{IV.14.4.1-fr}
Soient $f:X\to Y$ un morphisme localement de type fini, $y$ un point de $Y$, $x$
un point maximal de la fibre $X_y=X\times_Y\operatorname{Spec}(k(y))$. Supposons
$y$ géométriquement unibranche. Alors les conditions suivantes sont équivalentes~:
\begin{enumerate}[label=\emph{\alph*)}]
  \item $f$ est universellement ouvert au point $x$ (ou encore en tout point de la
  composante irréductible de $X_y$ de point générique $x$ (14.3.3.1, (iii))).

  \oldpage[IV-3]{210}

  \item Si $Y_0$ est l'unique composante irréductible de $Y$ contenant $y$, $\eta$
  son point générique, il existe une composante irréductible $Z$ de $X$ contenant
  $x$, dominant $Y_0$ et telle que
  $\dim_x(Z\cap X_y)=\dim(Z\cap X_\eta)$ (ce qui signifie que $Z$ est
  équidimensionnelle sur $Y_0$ au point $x$ (13.2.2)).
  \item[b$'$)] Pour tout voisinage ouvert $U$ de $x$ dans $X$ et toute
  générisation $y'$ de $y$, on a
  $\dim(U\cap X_{y'})\geq\dim_x(U\cap X_y)$.
\end{enumerate}
Si de plus $Y$ est localement noethérien, ces conditions sont aussi équivalentes à
la suivante~:
\begin{enumerate}[label=\emph{\alph*)}]
  \item[c)] $f$ est ouvert au point $x$.
\end{enumerate}
\end{theorem}

Notons d'abord que puisque $(\mathcal O_y)_{\mathrm{red}}$ est intègre, $y$
n'appartient qu'à une seule composante irréductible $Y_0$ de $Y$. Le fait que $b)$
et $b')$ sont équivalentes et que $a)$ implique $b')$ résulte de (14.3.13)~;
d'autre part, si $Y$ est localement noethérien, on a vu dans (14.2.3) que $c)$
implique $b)$. Il reste donc à montrer que lorsque $y$ est géométriquement
unibranche, $b)$ entraîne $a)$.

\begin{lemma}[14.4.1.1]
\label{IV.14.4.1.1-fr}
Soient $Y$ un préschéma, $Y'=Y[T_1,\ldots,T_n]$, $y'$ un point de $Y'$, $y$ son
image dans $Y$. Si $Y$ est géométriquement unibranche au point $y$, alors $Y'$ est
géométriquement unibranche au point $y'$.
\end{lemma}

En effet, comme pour tout $y\in Y$,
$\operatorname{Spec}(k(y)[T_1,\ldots,T_n])$ est géométriquement régulier sur $k(y)$
($\mathbf{0}_{\mathrm{IV}}$, 17.3.7), le morphisme structural $Y'\to Y$ est lisse
(6.8.1), et il suffit d'appliquer (11.3.14).

\begin{lemma}[14.4.1.2]
\label{IV.14.4.1.2-fr}
Soient $A$ un anneau local intègre unibranche, $B$ un anneau intègre contenant $A$
et entier sur $A$, $\mathfrak n$ un idéal premier de $B$ au-dessus de l'idéal
maximal $\mathfrak m$ de $A$. Alors le morphisme
$\operatorname{Spec}(B_{\mathfrak n})\to\operatorname{Spec}(A)$ est surjectif,
autrement dit, pour tout idéal premier $\mathfrak p$ de $A$, il existe un idéal
premier $\mathfrak q$ de $B$ tel que $\mathfrak q\subset\mathfrak n$ et
$\mathfrak q\cap A=\mathfrak p$.
\end{lemma}

Soient $K$ (resp. $L$) le corps des fractions de $A$ (resp. $B$), $A'$ la clôture
intégrale de $A$, $B'$ le sous-anneau de $L$ engendré par $A'$ et $B$, de sorte
qu'on a un diagramme commutatif d'injections canoniques
\[
  \xymatrix{
    B \ar[r] & B' \\
    A \ar[u] \ar[r] & A' \ar[u]
  }
\]
Comme $B'$ est entier sur $B$, il existe un idéal premier $\mathfrak n'$ de $B'$
tel que $\mathfrak n'\cap B=\mathfrak n$ (Bourbaki, \emph{Alg. comm.}, chap.~V,
\S~2, n\textsuperscript{o}~1, th.~1), et (pour la même raison)
$\operatorname{Spec}(A')\to\operatorname{Spec}(A)$ est surjectif. D'autre part,
comme $A$ est unibranche, $A'$ est un anneau \emph{local}, donc
$\mathfrak n'\cap A'$, qui est au-dessus de l'idéal maximal $\mathfrak m$ de $A$,
est nécessairement égal à l'unique idéal maximal $\mathfrak m'$ de $A'$. En vertu
du second théorème de Cohen--Seidenberg (\emph{loc. cit.}, \S~2,
n\textsuperscript{o}~4, th.~3) le morphisme
$\operatorname{Spec}(B'_{\mathfrak n'})\to\operatorname{Spec}(A')$ est surjectif,
donc il en est de même du composé
$\operatorname{Spec}(B'_{\mathfrak n'})\to\operatorname{Spec}(A')\to
\operatorname{Spec}(A)$~; mais ce morphisme est aussi le composé
$\operatorname{Spec}(B'_{\mathfrak n'})\to\operatorname{Spec}(B_{\mathfrak n})
\to\operatorname{Spec}(A)$, donc le morphisme
$\operatorname{Spec}(B_{\mathfrak n})\to\operatorname{Spec}(A)$ est surjectif.

Ces lemmes étant établis, revenons à la démonstration de l'implication $b)\Rightarrow
a)$ dans (14.4.1). En vertu de (14.3.3.1, (i)), il suffit de prouver que, pour tout
entier $n\geq 0$ et tout point $x'$ de $X'=X[T_1,\ldots,T_n]$ au-dessus de $x$, le
morphisme

\oldpage[IV-3]{211}

$f':X'\to Y'=Y[T_1,\ldots,T_n]$, déduit de $f$ par changement de base, est ouvert
au point $x'$. Compte tenu du lemme (14.4.1.1), de (2.3.4) et (4.2.7), on est donc
ramené à prouver que $f$ est \emph{ouvert} au point $x$~: en outre, il suffit
évidemment (14.1.2, (iii)) de montrer que la restriction de $f$ à un
sous-préschéma fermé de $X$ ayant $Z$ pour espace sous-jacent est ouverte au point
$x$, si bien qu'on peut se borner au cas où $X=Z$ est irréductible. Quitte à
remplacer $X$ par un voisinage ouvert $V$ de $x$ tel que $V\cap X_y$ soit
irréductible, on peut supposer, en vertu de (13.3.1), que le morphisme $f$ se
factorise en
\[
  X\xrightarrow{g}Y''=Y[T_1,\ldots,T_m]\to Y,
\]
où $g$ est quasi-fini, dominant et localement de type fini. Comme le morphisme
structural $Y''\to Y$ est ouvert (2.4.6), on est ramené à prouver que $g:X\to Y''$
est ouvert au point $x$. D'ailleurs, en vertu de (14.4.1.1), $g(x)$ est un point
géométriquement unibranche de $Y''$. On est donc ramené à prouver le lemme suivant~:

\begin{lemma}[14.4.1.3]
\label{IV.14.4.1.3-fr}
Soient $X$, $Y$ deux préschémas irréductibles, $f:X\to Y$ un morphisme localement
quasi-fini et dominant. Si $x\in X$ est tel que $y=f(x)$ soit unibranche sur $Y$,
alors $f$ est ouvert au point $x$.
\end{lemma}

Il suffit de prouver que
$f(\operatorname{Spec}(\mathcal O_{X,x}))=\operatorname{Spec}(\mathcal O_{Y,y})$
(1.10.3). On peut donc se borner, par le changement de base
$\operatorname{Spec}(\mathcal O_{Y,y})\to Y$, au cas où $Y=\operatorname{Spec}(A)$,
où $A$ est un anneau local et $y$ est le point fermé de $Y$ (tenant compte de
(I, 3.6.5) et ($\mathbf{0}_{\mathrm I}$, 2.1.8), qui prouvent que
$X\times_Y\operatorname{Spec}(\mathcal O_{Y,y})$ est irréductible)~; remplaçant $f$
par $f_{\mathrm{red}}$, on peut supposer $X$ et $Y$ réduits, donc intègres.
Remplaçant éventuellement $X$ et $Y$ par des voisinages affines de $x$ et $y$
respectivement, on peut supposer (8.12.9) que le morphisme $f$ se factorise en
$X\xrightarrow{j}X_1\xrightarrow{g}Y$, où $j$ est une immersion ouverte et $g$ un
morphisme \emph{fini} (évidemment dominant)~; comme $X$ et $X_1$ sont affines, $j$
est affine, donc séparé et quasi-compact, et par suite se factorise en
$X\xrightarrow{u}X_2\xrightarrow{h}X_1$, où $X_2$ est l'image fermée de $X$ par
$j$, $h$ l'injection canonique et $u$ une immersion ouverte (I, 9.5.3). En d'autres
termes, on peut supposer que $X_1$ est intègre, ou encore de la forme
$X_1=\operatorname{Spec}(B)$, où $B$ est une $A$-algèbre intègre et finie, contenant
$A$ puisque $g$ est dominant. Si $\mathfrak n$ est l'idéal premier de $B$
correspondant au point $x$, l'hypothèse que $A$ est unibranche implique alors
(14.4.1.2) que le morphisme
$\operatorname{Spec}(B_{\mathfrak n})\to\operatorname{Spec}(A)$ est surjectif,
c'est-à-dire que
$\operatorname{Spec}(\mathcal O_{X,x})\to\operatorname{Spec}(\mathcal O_{Y,y})$
est surjectif. C.Q.F.D.

\begin{corollary}[14.4.2]
\label{IV.14.4.2-fr}
Soient $f:X\to Y$ un morphisme localement de type fini, $y$ un point
géométriquement unibranche de $Y$, $\eta$ le point générique de l'unique composante
irréductible $Y_0$ de $Y$ contenant $y$. Les conditions suivantes sont
équivalentes~:
\begin{enumerate}[label=\emph{\alph*)}]
  \item $f$ est universellement ouvert en tous les points de $X_y$ (ou, ce qui
  revient au même (14.3.3.1, (ii)), aux points maximaux de $X_y$).
  \item Pour tout $x\in X_y$, il existe une composante irréductible $Z$ de $X$
  contenant $x$ et équidimensionnelle sur $Y$ au point $x$ (13.2.2).
  \item[b$'$)] Pour tout $x\in X_y$ et tout voisinage ouvert $U$ de $x$ dans $X$,
  on a $\dim(U\cap X_\eta)\geq\dim_x(U\cap X_y)$.
  \item[b$''$)] Pour tout ouvert $U$ de $X$, on a
  $\dim(U\cap X_\eta)\geq\dim(U\cap X_y)$.
\end{enumerate}

\oldpage[IV-3]{212}

Lorsque de plus $Y$ est localement noethérien, ces conditions sont encore
équivalentes à la suivante~:
\begin{enumerate}[label=\emph{\alph*)}]
  \item[c)] $f$ est ouvert en tous les points de $X_y$ (ou, ce qui revient au même
  (14.3.3.1, (ii)), aux points maximaux de $X_y$).
\end{enumerate}
\end{corollary}

L'équivalence de $a)$ et $c)$ lorsque $Y$ est localement noethérien résulte de
(14.4.1)~; les conditions $b)$ ou $b')$, appliquées aux \emph{points maximaux} de
$X_y$, entraînent $a)$ en vertu également de (14.3.3.1, (ii))~; enfin, $b')$ et
$b'')$ sont équivalentes, car
\[
  \dim(U\cap X_y)=\sup_{x\in X_y}(\dim_x(U\cap X_y)).
\]
Reste à voir que la condition $a)$ entraîne $b)$ et $b')$ en \emph{tout} point
$x\in X_y$. Posons $d=\dim_x(X_y)$, et soit $x'$ le point générique d'une
composante irréductible de $X_y$ contenant $x$ et de dimension $d$. En vertu de
$a)$ et de (14.4.1), il y a une composante irréductible $Z$ de $X$ contenant $x'$
et équidimensionnelle sur $Y$ au point $x'$, donc telle que
$\dim_{x'}(Z\cap X_y)=\dim(Z\cap X_\eta)$. Mais par construction
$\dim_{x'}(Z\cap X_y)=\dim_{x'}(X_y)=d$, et
$\dim_x(Z\cap X_y)\leq\dim(Z\cap X_y)=d$~; compte tenu de (13.1.6), cela prouve
que $Z$ est équidimensionnel sur $Y$ au point $x$~; donc $a)$ entraîne $b)$. De
plus, on a $\dim(X_\eta)\geq\dim(Z\cap X_\eta)=d=\dim_x(X_y)$. Remplaçant $X$ par
un voisinage ouvert $U$ de $x$, on voit ainsi que $a)$ entraîne $b')$. C.Q.F.D.

\begin{corollary}[14.4.3]
\label{IV.14.4.3-fr}
Soient $Y$ un préschéma géométriquement unibranche, $f:X\to Y$ un morphisme
localement de type fini. Les conditions suivantes sont équivalentes~:
\begin{enumerate}[label=\emph{\alph*)}]
  \item $f$ est universellement ouvert.
  \item Pour tout ouvert $U$ de $X$, tout $y\in Y$ et toute générisation $y'$ de
  $y$, on a $\dim(U\cap X_{y'})\geq\dim(U\cap X_y)$.
\end{enumerate}
Si de plus $Y$ est localement noethérien, ces conditions sont aussi équivalentes à
la suivante~:
\begin{enumerate}[label=\emph{\alph*)}]
  \item[c)] $f$ est ouvert.
\end{enumerate}
\end{corollary}

\begin{corollary}[14.4.4]
\EGAReaderTarget{ega4.c04.cor.14-4-4}{IV.14.4.4-fr}
(critère de Chevalley). Soit $f:X\to Y$ un morphisme localement de type fini.
\begin{enumerate}[label=(\roman*)]
  \item Si $f$ est équidimensionnel en un point $x\in X$ (13.3.2) et si
  $y=f(x)$ est un point géométriquement unibranche de $Y$, $f$ est
  universellement ouvert au point $x$.
  \item Si $Y$ est géométriquement unibranche, $f$ est universellement ouvert en
  tous les points de $X$ où $f$ est équidimensionnel, et l'ensemble de ces points
  est ouvert dans $X$. En particulier, si $f$ est équidimensionnel, il est
  universellement ouvert.
\end{enumerate}
\end{corollary}

L'assertion (ii) est conséquence triviale de (i), puisqu'on sait déjà que l'ensemble
des points où $f$ est équidimensionnel est ouvert (13.3.2). Quant à l'assertion (i),
elle résulte de ce que l'hypothèse implique que la condition $b)$ de (14.4.1) est
vérifiée au point générique d'une composante irréductible de $X_y$ contenant $x$,
compte tenu de (13.3.1)~; il suffit donc d'appliquer (14.4.1).

\begin{remark}[14.4.5]
\label{IV.14.4.5-fr}
On peut prouver que si $Y$ est localement noethérien, et si toutes les générisations
de $y=f(x)$ sont des points géométriquement unibranches de $Y$ (cf. (6.15.2)),
alors, si $f$ est équidimensionnel au point $x$, il est universellement ouvert
\emph{dans un voisinage de $x$}.
\end{remark}

\oldpage[IV-3]{213}

\begin{corollary}[14.4.6]
\label{IV.14.4.6-fr}
Soient $Y$ un préschéma localement noethérien, $f:X\to Y$ un morphisme localement
de type fini. Soit $y$ un point géométriquement unibranche de $Y$, et supposons en
outre que pour tout $x\in f^{-1}(y)$, l'anneau $\mathcal O_{X,x}$ soit
équidimensionnel. Alors les conditions suivantes sont équivalentes~:
\begin{enumerate}[label=\emph{\alph*)}]
  \item $f$ est équidimensionnel (13.3.2) en tous les points de $f^{-1}(y)$.
  \item $f$ est ouvert en tous les points de $f^{-1}(y)$.
  \item $f$ est universellement ouvert en tous les points de $f^{-1}(y)$.
\end{enumerate}
\end{corollary}

En effet, $c)$ implique trivialement $b)$, et $a)$ implique $c)$ en vertu de
(14.4.4)~; enfin, en raison des hypothèses sur $\mathcal O_y$ et $\mathcal O_x$,
$b)$ implique $a)$ par (14.2.2). Plus généralement~:

\begin{proposition}[14.4.7]
\label{IV.14.4.7-fr}
Soient $Y$ un préschéma localement noethérien, $f:X\to Y$ un morphisme localement
de type fini, $x$ un point de $X$, tel que $y=f(x)$ soit géométriquement unibranche.
Les conditions suivantes sont équivalentes~:
\begin{enumerate}[label=\emph{\alph*)}]
  \item $f$ est équidimensionnel (13.3.2) au point $x$.
  \item L'anneau $\mathcal O_x$ est équidimensionnel et $f$ est ouvert aux points
  génériques des composantes irréductibles de $f^{-1}(y)$ contenant $x$ (donc
  aussi en tout point d'une telle composante).
  \item L'anneau $\mathcal O_x$ est équidimensionnel, et $f$ est universellement
  ouvert aux points génériques des composantes irréductibles de $f^{-1}(y)$
  contenant $x$ (donc aussi en tout point d'une telle composante).
\end{enumerate}
En outre, lorsque ces conditions sont satisfaites, pour tout sous-préschéma fermé
réduit $X_i$ de $X$ ayant pour espace sous-jacent une composante irréductible de $X$
contenant $x$, la restriction $f_i:X_i\to Y$ de $f$ est un morphisme
équidimensionnel au point $x$, et universellement ouvert en tous les points des
composantes irréductibles de $f^{-1}(y)\cap X_i$ qui contiennent $x$.
\end{proposition}

La condition $a)$ implique que $f$ est équidimensionnel en tous les points
génériques des composantes irréductibles de $f^{-1}(y)$ contenant $x$ (13.3.1), et
par suite (14.4.4) universellement ouvert en ces points~; le même raisonnement
appliqué à chaque $f_i$ (compte tenu de (13.3.3)) prouve la dernière assertion de
la proposition, compte tenu de (14.3.3.1, (iii)). En outre, d'après (14.2.1), on a
les relations
\begin{equation}
\tag{14.4.7.1}
\label{IV.14.4.7.1-fr}
  \dim(\mathcal O_{X_i,x})=\dim(\mathcal O_y)+
  \dim(\mathcal O_{X_i,x}\otimes_{\mathcal O_y}k(y))
\end{equation}
\begin{equation}
\tag{14.4.7.2}
\label{IV.14.4.7.2-fr}
  \dim(\mathcal O_{X,x})=\dim(\mathcal O_y)+
  \dim(\mathcal O_{X,x}\otimes_{\mathcal O_y}k(y))
\end{equation}
et puisque $f$ est équidimensionnel au point $x$, il résulte de (13.3.1) que l'on
a $\dim_x(X_i\cap f^{-1}(y))=\dim_x(f^{-1}(y))$, donc (5.2.3)
\[
  \dim(\mathcal O_{X_i,x}\otimes_{\mathcal O_y}k(y))=
  \dim(\mathcal O_{X,x}\otimes_{\mathcal O_y}k(y)).
\]
On conclut donc que $\dim(\mathcal O_{X_i,x})=\dim(\mathcal O_{X,x})$ pour tout
$i$, autrement dit $\mathcal O_{X,x}$ est équidimensionnel, et cela achève de
montrer que $a)$ entraîne $c)$. Il est clair que $c)$ entraîne $b)$~; enfin, $b)$
entraîne la relation (14.4.7.2) en vertu de (14.2.1)~; il résulte alors de (13.3.6)
que $b)$ entraîne $a)$.

\begin{proposition}[14.4.8]
\label{IV.14.4.8-fr}
Soient $Y$ un préschéma noethérien, $f:X\to Y$ un morphisme localement de type
fini, $y$ un point de $Y$, $x$ un point maximal de $X_y=f^{-1}(y)$. Les conditions
suivantes sont équivalentes~:
\begin{enumerate}[label=\emph{\alph*)}]
  \item Le morphisme $f$ est universellement ouvert au point $x$, autrement dit,
  pour tout changement de base $g:Y'\to Y$, on a la propriété $P(Y')$~: pour tout
  point $x'$ de $X'=X\times_Y Y'$ au-dessus de $x$, le morphisme
  $f'=f_{(Y')}:X'\to Y'$ est ouvert au point $x'$.

  \oldpage[IV-3]{214}

  \item[a$'$)] La propriété $P(Y')$ est vraie pour tout morphisme fini
  $g:Y'\to Y$.
  \item[a$''$)] La propriété $P(Y'')$ est vraie pour le normalisé $Y''$ de
  $Y_{\mathrm{red}}$ (\textbf{II}, 6.3.8).
  \item Pour tout point $x''$ de $X''=X\times_Y Y''$ au-dessus de $x$, il existe
  une composante irréductible $Z''$ de $X''$ contenant $x''$ et
  équidimensionnelle sur $Y''$ au point $x''$.
\end{enumerate}
\end{proposition}

Il est trivial que $a)$ implique $a')$. Pour montrer que $a')$ implique $a'')$
notons que l'on peut écrire $Y''=\operatorname{Spec}(\mathcal B)$, où $\mathcal B$
est une $\mathcal O_Y$-Algèbre quasi-cohérente entière sur $\mathcal O_Y$~; comme
$Y$ est noethérien, $\mathcal B$ est limite inductive de ses sous-$\mathcal O_Y$-
Algèbres $\mathcal B_\lambda$ qui sont quasi-cohérentes et de type fini (I, 9.6.6)~;
mais alors les $\mathcal B_\lambda$ sont des $\mathcal O_Y$-Algèbres \emph{finies}
(\textbf{II}, 6.1.2)~; on peut donc écrire
$Y''=\varprojlim Y'_\lambda$, où
$Y'_\lambda=\operatorname{Spec}(\mathcal B_\lambda)$, d'où
$X''=X\times_Y Y''=\varprojlim X'_\lambda$, avec
$X'_\lambda=X\times_Y Y'_\lambda$. En vertu de $a')$, les morphismes
$f'_\lambda:X'_\lambda\to Y'_\lambda$ sont ouverts en tous les points de
$X'_\lambda$ au-dessus de $x$~; on en conclut que $f''$ est ouvert en tous les
points de $X''$ au-dessus de $x$, par (8.10.1) et (14.3.3.1, (i)).

Comme le préschéma $Y''$ est \emph{normal} par définition, le fait que $b)$ entraîne
$a'')$ résulte de (14.4.4) appliqué à la composante irréductible
équidimensionnelle de l'énoncé et à la restriction de $f''$ à cette composante. Il
reste donc à montrer que $a'')$ entraîne $a)$ et $b)$. Compte tenu de (1.10.3), on
peut se borner au cas où $Y=\operatorname{Spec}(\mathcal O_y)$, en notant que le
morphisme canonique $\operatorname{Spec}(\mathcal O_y)\to Y$ est universellement
bicontinu (I, 3.6.5), et d'autre part que
$Y''\times_Y\operatorname{Spec}(\mathcal O_y)$ est le normalisé de
$(\operatorname{Spec}(\mathcal O_y))_{\mathrm{red}}$ comme il résulte de la
permutabilité des opérations de fermeture intégrale et de localisation (Bourbaki,
\emph{Alg. comm.}, chap.~V, \S~1, n\textsuperscript{o}~5, prop.~16). Supposant donc
$Y=\operatorname{Spec}(A)$, où $A$ est un anneau local noethérien, et $y$ le point
fermé de $Y$, on sait ($\mathbf{0}_{\mathrm{IV}}$, 23.2.5) qu'il existe une
factorisation
\[
  Y''\xrightarrow{u}Y_1\xrightarrow{v}Y
\]
du morphisme structural, telle que $v$ soit un morphisme \emph{fini surjectif}, $u$
un morphisme entier, \emph{radiciel et dominant}, donc surjectif (\textbf{II},
6.1.10) et par suite un \emph{homéomorphisme universel} (2.4.5). Si l'on pose
$X_1=X\times_Y Y_1$, la projection $X''\to X_1$ est donc un homéomorphisme, et
l'hypothèse $a'')$ entraîne par suite que $f_1=f_{(Y_1)}:X_1\to Y_1$ est ouvert en
tous les points de $X_1$ au-dessus de $x$. En outre, cela montre que pour prouver
la propriété $b)$, il suffit de prouver la même propriété où l'on remplace $Y''$,
$X''$ et $x''$ par $Y_1$, $X_1$ et un point $x_1$ de $X_1$ au-dessus de $x$~: Mais
$Y_1$ est noethérien et en outre il est géométriquement unibranche puisque $Y''$
est normal et $u$ radiciel (6.15.1)~; la propriété à prouver résulte donc de
(14.4.1). Il reste à montrer que $f$ est universellement ouvert au point $x$, ce
qui résultera du lemme suivant~:

\begin{lemma}[14.4.8.1]
\label{IV.14.4.8.1-fr}
Soit $v:Y_1\to Y$ un morphisme fermé (resp. universellement fermé) et surjectif.
Pour qu'un morphisme $f:X\to Y$ soit ouvert (resp. universellement ouvert) en un
point $x\in X$, il suffit que $f_1=f_{(Y_1)}:X_1=X\times_Y Y_1\to Y_1$ soit ouvert
(resp. universellement ouvert) en tous les points de $X_1$ au-dessus de $x$.
\end{lemma}

La seconde assertion résulte trivialement de la première et du fait que pour tout
changement de base $Y'\to Y$, le morphisme
$v_{(Y')}:Y_1\times_Y Y'\to Y'$ est encore surjectif et est fermé si $v$ est
universellement fermé. Pour prouver la première assertion, considérons

\oldpage[IV-3]{215}

un voisinage ouvert $U$ de $x$ dans $X$~; comme $v$ est fermé et surjectif, pour
que $f(U)$ soit un voisinage de $y=f(x)$, il faut et il suffit que
$v^{-1}(f(U))$ soit un voisinage de $v^{-1}(y)$. Mais si $p:X_1\to X$ est la
projection canonique, on a $v^{-1}(f(U))=f_1(p^{-1}(U))$ (I, 3.4.8), et
l'hypothèse entraîne que $f_1(p^{-1}(U))$ est un voisinage de
$f_1(p^{-1}(x))=v^{-1}(y)$ (I, 3.4.8).

\begin{corollary}[14.4.9]
\label{IV.14.4.9-fr}
Soient $Y$ un préschéma noethérien, $f:X\to Y$ un morphisme localement de type
fini. Les conditions suivantes sont équivalentes~:
\begin{enumerate}[label=\emph{\alph*)}]
  \item $f$ est universellement ouvert, autrement dit, pour tout changement de
  base $Y'\to Y$, le morphisme $f_{(Y')}:X\times_Y Y'\to Y'$ est ouvert.
  \item[a$'$)] Pour tout morphisme fini $Y_1\to Y$, $f_{(Y_1)}$ est ouvert.
  \item[a$''$)] Si $Y''$ est le normalisé de $Y_{\mathrm{red}}$, $f_{(Y'')}$ est
  ouvert.
  \item Pour tout point $x''$ de $X''=X\times_Y Y''$, il existe une composante
  irréductible $Z''$ de $X''$ contenant $x''$ et équidimensionnelle sur $Y''$ au
  point $x''$ (cf. (14.4.10, (ii))).
\end{enumerate}
\end{corollary}

Cela résulte aussitôt de (14.4.8) et (14.1.4).

\begin{remarks}[14.4.10]
\label{IV.14.4.10-fr}
\begin{enumerate}[label=(\roman*)]
  \item L'équivalence des conditions $a)$ et $b)$ dans (14.4.8) (resp. (14.4.9))
  est encore valable pour un préschéma $Y$ quelconque et un morphisme $f$
  localement de type fini. En effet, $a)$ entraîne $b)$ en vertu de (14.4.1)~;
  inversement, $b)$ entraîne que $f''$ est universellement ouvert aux points de
  $X''$ au-dessus de $x$ (resp. en tout point de $X''$) en vertu de (14.4.1), et
  on en conclut la propriété $a)$ en appliquant le lemme (14.4.8.1) au morphisme
  entier et surjectif $Y''\to Y$.

  Il se peut que, dans (14.4.1), pour l'équivalence de $a)$ et $c)$, l'hypothèse
  supplémentaire que $Y$ est noethérien soit superflue (cf. (14.3.14)). S'il en
  est ainsi, les hypothèses noethériennes sont également superflues dans (14.4.2),
  (14.4.3), (14.4.8) et (14.4.9).

  \item On peut donner des exemples de morphismes $f:X\to Y$ ayant les propriétés
  suivantes~: $Y$ est noethérien, régulier et de dimension 2, $f$ est
  \emph{universellement ouvert} et de type fini, $X$ a deux composantes
  irréductibles $X_1$, $X_2$, mais la restriction $X_1\to Y$ de $f$ à l'une
  d'elles \emph{n'est pas un morphisme ouvert}. Le principe de la construction
  s'appuie sur la méthode générale de «~recollement~» qui sera exposée au chap.~V,
  et ne peut donc qu'être esquissé ici. On part d'un point fermé $y$ de $Y$, et on
  considère le $Y$-schéma $Y_1$ obtenu en faisant \emph{éclater} $y$
  (\textbf{II}, 8.1.3)~; si $f_1:Y_1\to Y$ est le morphisme structural, on sait que
  la restriction de $f_1$ à $f_1^{-1}(Y-\{y\})$ est un \emph{isomorphisme} sur
  $Y-\{y\}$ (\emph{loc. cit.}), tandis que la fibre $f_1^{-1}(y)$ est isomorphe à
  $\operatorname{Proj}(S)$, où
  $S=\bigoplus_{k=0}^{\infty}\mathfrak m_y^k/\mathfrak m_y^{k+1}$
  (\textbf{II}, 3.5.3), c'est-à-dire ici à $\mathbf P^1_{k(y)}$~; il résulte de
  (14.4.1) que $f_1$ n'est pas ouvert au point générique de $f_1^{-1}(y)$. D'autre
  part, posons $Y_2=\mathbf P^1_Y$, et soit $f_2:Y_2\to Y$ le morphisme structural~;
  il résulte de (\textbf{II}, 8.4.4) que $f_2$ est plat, donc universellement ouvert
  (2.4.6)~; en outre (\textbf{II}, 3.5.3), $f_2^{-1}(y)$ est isomorphe à
  $\mathbf P^1_{k(y)}$~; il suffit alors de «~recoller~» $Y_1$ et $Y_2$ le long des
  fibres isomorphes $f_1^{-1}(y)$ et $f_2^{-1}(y)$, ce qui donne un morphisme
  $f:X\to Y$ où les composantes irréductibles $X_1$, $X_2$ de $X$ s'identifient
  canoniquement à $Y_1$ et $Y_2$ respectivement, et les restrictions de $f$ à ces
  composantes à $f_1$ et $f_2$.

  Rappelons toutefois (12.1.1.5) que si $Y$ est localement noethérien, $f:X\to Y$
  de type fini et \emph{plat}, alors toute composante irréductible de $X$ est
  \emph{équidimensionnelle} sur $Y$ (et par suite la restriction de $f$ à une telle
  composante est \emph{universellement ouverte} si tous les points de $Y$ sont
  \emph{géométriquement unibranches}).

  \item Rappelons (12.1.2, (i)) qu'il y a des morphismes $f:X\to Y$ ayant les
  propriétés suivantes~: $Y$ est noethérien (non géométriquement unibranche), $f$
  est fini et plat (et même étale (17.6.3)), mais la restriction de $f$ à une
  composante irréductible de $X$ \emph{n'est pas un morphisme ouvert} (bien que $f$
  lui-même le soit d'après (2.4.6)).

  \item Le critère de Chevalley (14.4.4) explique l'importance de la notion de
  morphisme universellement ouvert. Cette notion permet en effet, dans de nombreux
  résultats, plus ou moins classiques, de remplacer une hypothèse de \emph{normalité}
  par l'hypothèse qu'un certain morphisme est universellement ouvert~; l'énoncé plus
  général

  \oldpage[IV-3]{216}

  ainsi obtenu s'appliquera en particulier à des morphismes \emph{plats} (2.4.6),
  dont l'importance en géométrie algébrique va croissant. On peut considérer que
  les énoncés faisant intervenir l'hypothèse qu'un morphisme est universellement
  ouvert sont des généralisations communes d'énoncés faisant intervenir une
  hypothèse de normalité et d'énoncés faisant intervenir une hypothèse de platitude.
\end{enumerate}
\end{remarks}

\subsection{Morphismes universellement ouverts et quasi-sections}
\label{subsection:IV.14.5-fr}

\begin{lemma}[14.5.1]
\label{IV.14.5.1-fr}
Soient $Y$ un préschéma irréductible localement noethérien, $f:X\to Y$ un
morphisme localement de type fini, $x$ un point de $X$, tel que $f$ soit
équidimensionnel (13.3.2) au point $x$~; on pose $y=f(x)$,
$e=\dim_x(f^{-1}(y))$. Soit $X'$ une partie fermée irréductible de $X$ contenant
$x$, $n$ un entier tel que l'on ait $\operatorname{codim}(X',X)\leq n$ et
$\dim_x(X'\cap f^{-1}(y))\leq e-n$. Alors on a nécessairement
$\operatorname{codim}(X',X)=n$, $\dim_x(X'\cap f^{-1}(y))=e-n$, et la restriction $X'\to Y$
de $f$ est un morphisme équidimensionnel en $x$ (et a fortiori dominant).
\end{lemma}

La question étant locale sur $X$, on peut supposer que $f$ est de type fini et
\emph{équidimensionnel}, de sorte que pour \emph{tout} $z\in Y$, \emph{toutes} les
composantes irréductibles de $f^{-1}(z)$ sont de dimension $e$ (13.3.1, $a''$).
Soient $x'$ le point générique de $X'$, $y'=f(x')$,
$Y'=\overline{\{y'\}}=\overline{f(X')}$ et posons $Z=f^{-1}(Y')$. En vertu de
($\mathbf 0$, 14.2.2), on a
\begin{equation}
  \operatorname{codim}(X',Z)\leq n.
  \tag{14.5.1.1}\label{IV.14.5.1.1-fr}
\end{equation}
Et si les deux membres sont égaux, on a nécessairement $\operatorname{codim}(X',X)=n$ et $Z$
contient une composante irréductible de $X$, ce qui entraîne (13.3.1) que
$f(X')$ est dense dans $Y$ et par suite $Z=X$~; donc l'égalité
$\operatorname{codim}(X',Z)=n$ équivaut à la conjonction de l'égalité $\operatorname{codim}(X',X)=n$ et de
la relation $Z=X$.

D'autre part, en raisonnant dans les préschémas réduits de $Y$ et $X$ ayant $Y'$
et $Z$ respectivement pour espaces sous-jacents, on déduit de (5.1.2) et de
(\textbf{I}, 3.6.5) que l'on a
\begin{equation}
  \operatorname{codim}(X'\cap f^{-1}(y'),f^{-1}(y'))=\operatorname{codim}(X',Z).
  \tag{14.5.1.2}\label{IV.14.5.1.2-fr}
\end{equation}
En vertu de l'hypothèse, $f^{-1}(y')$ est biéquidimensionnel (5.2.1) et de
dimension $e$, donc ($\mathbf 0$, 14.3.5), on a, en vertu de (14.5.1.2) et
(14.5.1.1)
\begin{equation}
  e'=\dim(X'\cap f^{-1}(y'))=e-\operatorname{codim}(X',Z)\geq e-n,
  \tag{14.5.1.3}\label{IV.14.5.1.3-fr}
\end{equation}
l'égalité ayant lieu si et seulement si $\operatorname{codim}(X',X)=n$ et si $f(X')$ est dense
dans $Y$. Enfin, (13.1.1), on a $\dim_x(X'\cap f^{-1}(y))\geq e'$, d'où, par
(14.5.1.3),
\begin{equation}
  \dim_x(X'\cap f^{-1}(y))\geq e'\geq e-n.
  \tag{14.5.1.4}\label{IV.14.5.1.4-fr}
\end{equation}
Or, par hypothèse, on a aussi $\dim_x(X'\cap f^{-1}(y))\leq e-n$, d'où les
conclusions de la proposition.

\begin{corollary}[14.5.2]
\label{IV.14.5.2-fr}
Les hypothèses sur $f$, $X$, $Y$, $x$ étant celles de (14.5.1), supposons en
outre que $x$ ne soit pas point maximal de $f^{-1}(y)$. Alors il existe un
voisinage ouvert affine $U$ de $x$ dans $X$, et une section
$g\in\Gamma(U,\mathcal O_X)$ telle que l'ensemble $X'$ des $x'\in U$ tels que
$g(x')=0$ contienne $x$ et ne contienne aucun point maximal de $f^{-1}(y)$. Pour
tout $g$ ayant ces propriétés, $X'$ est équidimensionnel sur $Y$ au point $x$, et
l'on a
\begin{equation}
  \dim_x(X'\cap f^{-1}(y))=e-1 \qquad\text{et}\qquad \operatorname{codim}(X',X)=1.
  \tag{14.5.2.1}\label{IV.14.5.2.1-fr}
\end{equation}
\end{corollary}

\oldpage[IV-3]{217}

On peut se borner au cas où $X=U$ est un voisinage ouvert affine de $x$ tel que
toutes les composantes irréductibles de $f^{-1}(y)$ contiennent $x$. Ces
composantes correspondent aux idéaux premiers minimaux de
$\mathcal O_x/\mathfrak m_y\mathcal O_x$, et par hypothèse ces idéaux sont
distincts de $\mathfrak m_x/\mathfrak m_y\mathcal O_x$ (Bourbaki,
\emph{Alg. comm.}, chap.~II, \S~1, n\textsuperscript{o}~1, prop.~2)~; pour obtenir
un $g\in\Gamma(U,\mathcal O_X)$ vérifiant les conditions de l'énoncé, il suffit de
prendre $g\in\mathfrak j_x$ tel que l'image de $g$ dans $\mathfrak m_x$
n'appartienne à aucun des idéaux premiers précédents. D'ailleurs, on a
$\operatorname{codim}(X',X)\leq 1$ (5.1.8), et comme $X'$ ne contient aucune des composantes
irréductibles de $f^{-1}(y)$ et que celles-ci sont de dimension $e$, on a
($\mathbf 0$, 14.2.2.2) $\dim_x(X'\cap f^{-1}(y))\leq e-1$. Il suffit alors
d'appliquer (14.5.1).

\begin{proposition}[14.5.3]
\label{IV.14.5.3-fr}
Soient $Y$ un préschéma irréductible localement noethérien, $f:X\to Y$ un
morphisme localement de type fini, $x$ un point de $X$. On suppose que $f$ est
équidimensionnel au point $x$ et que $x$ est fermé dans $f^{-1}(f(x))$. Alors il
existe une partie irréductible $X'$ de $X$, localement fermée dans $X$, contenant
$x$ et telle que la restriction $X'\to Y$ de $f$ (où $X'$ est le sous-préschéma
réduit de $X$ ayant $X'$ pour espace sous-jacent) soit un morphisme quasi-fini et
dominant.
\end{proposition}

En effet, avec les notations de (14.5.2), l'hypothèse que $x$ est fermé dans
$f^{-1}(y)$ entraîne que $x$ n'est pas point maximal de $X'\cap f^{-1}(y)$ tant
que $e-1\geq 1$. Il suffit donc d'appliquer (14.5.2) en raisonnant par récurrence
descendante sur $e=\dim_x(f^{-1}(f(x)))$ jusqu'à ce que l'on arrive à $e=1$~;
l'application de (14.5.2) dans ce dernier cas donne un $X'$ tel que
$X'\cap f^{-1}(f(x))$ soit noethérien et de dimension 0, donc fini et discret~;
comme alors $X'$ est équidimensionnel sur $Y$, $X'\cap f^{-1}(f(x'))$ est de
dimension 0 pour tout $x'\in X'$, ce qui entraîne que la restriction $X'\to Y$ de
$f$ est un morphisme quasi-fini (\textbf{II}, 6.2.2).

\begin{corollary}[14.5.4]
\label{IV.14.5.4-fr}
Sous les hypothèses de (14.5.3), supposons en outre que $Y=\operatorname{Spec}(A)$, où $A$ est
un anneau local noethérien intègre et complet, et que $y=f(x)$ soit l'unique point
fermé de $Y$. Alors, il existe un anneau local intègre $A'$, contenant $A$, qui est
une $A$-algèbre finie et possède la propriété suivante~: si l'on pose
$Y'=\operatorname{Spec}(A')$ et $X'=X\times_Y Y'$, il existe une $Y'$-section $h:Y'\to X'$ telle
que le morphisme composé $Y'\xrightarrow{h}X'\to X$ soit une immersion dont
l'image contient $x$.
\end{corollary}

Remplaçant au besoin $X$ par un sous-préschéma réduit irréductible de $X$, on peut,
en vertu de (14.5.3), se borner au cas où le morphisme $f$ est déjà quasi-fini et
dominant. Utilisant (\textbf{II}, 6.2.5), on en déduit que $\mathcal O_x=A'$ est un
anneau intègre qui est une $A$-algèbre finie, et que $X$ est somme disjointe du
sous-préschéma fermé $Y'=\operatorname{Spec}(A')$ et d'un sous-préschéma $Y''$~; le schéma $Y'$
répond à la question, le morphisme composé $Y'\to X'\to X$ n'étant autre que le
morphisme canonique $\operatorname{Spec}(\mathcal O_x)\to X$.

\begin{remark}[14.5.5]
\label{IV.14.5.5-fr}
Si on n'exige pas que dans l'énoncé de (14.5.4), le morphisme $Y'\to X$ soit une
immersion, on peut supposer que $A'$ est en outre \emph{intégralement clos}~: il
suffit en effet de remplacer $A'$ par sa \emph{clôture intégrale} $A_1$, car on
sait ($\mathbf 0$, 23.1.5) que $A_1$ est un $A$-module de type fini.
\end{remark}

\begin{proposition}[14.5.6]
\label{IV.14.5.6-fr}
Soient $Y$ un préschéma localement noethérien, irréductible, régulier et de
dimension 1, $f:X\to Y$ un morphisme localement de type fini, $y$ un point de $Y$.
Les conditions suivantes sont équivalentes~:
\begin{enumerate}[label=\emph{\alph*)}]
  \item $f_{\mathrm{red}}$ est plat en tout point de $f^{-1}(y)$.

  \oldpage[IV-3]{218}

  \item $f$ est universellement ouvert dans un voisinage de $f^{-1}(y)$.
  \item $f$ est ouvert dans un voisinage de $f^{-1}(y)$.
  \item Toute composante irréductible de $X$ rencontrant $f^{-1}(y)$ domine $Y$.
  \item Pour tout point $x\in f^{-1}(y)$, fermé dans $f^{-1}(y)$, et toute
  composante irréductible $X_i$ de $X$ contenant $x$, il existe une partie
  irréductible $X'$ de $X$, localement fermée dans $X$, contenant $x$, contenue
  dans $X_i$, et telle que la restriction $X'\to Y$ de $f$ soit un morphisme
  quasi-fini et dominant.
\end{enumerate}
\end{proposition}

L'équivalence de $a)$, $b)$, $c)$ et $d)$ a déjà été démontrée (14.3.8). Il est
clair que $e)$ entraîne $d)$, car pour toute composante irréductible $X_i$ de $X$
rencontrant $f^{-1}(y)$, il existe dans $X_i\cap f^{-1}(y)$ un point fermé dans cet
espace (5.1.11). Enfin, pour prouver que $c)$ entraîne $e)$, on peut se borner au
cas où $f$ est un morphisme de type fini~; considérons un point fermé $x$ de
$f^{-1}(y)$ et soit $X_i$ une composante irréductible de $X$ contenant $x$. Comme
la restriction $f_i:X_i\to Y$ de $f$ à $X_i$ est un morphisme dominant, il résulte
de l'équivalence de $c)$ et $d)$ pour $f_i$ que ce morphisme est ouvert aux points
génériques de $X_i\cap f^{-1}(y)$. Il résulte donc de (14.2.2) que $f_i$ est
équidimensionnel au point $x$, et on conclut alors à l'aide de (14.5.3).

\begin{remark}[14.5.7]
\EGAReaderTarget{ega4.c04.rem.14-5-7}{IV.14.5.7-fr}
Si, dans l'énoncé de (14.5.5), on suppose que $Y=\operatorname{Spec}(A)$, où $A$ est un anneau
de valuation discrète complet, et que $y$ est le point fermé de $Y$, on peut en
outre supposer que $X'=\operatorname{Spec}(A')$, où $A'$ est un anneau de valuation discrète qui
est une $A$-algèbre finie, comme le montre la démonstration de (14.5.4) et le fait
qu'un anneau local intègre régulier et de dimension 1 est un anneau de valuation
discrète (\textbf{II}, 7.1.6).
\end{remark}

\begin{proposition}[14.5.8]
\label{IV.14.5.8-fr}
Soient $Y$ un préschéma localement noethérien, $f:X\to Y$ un morphisme localement
de type fini, $y$ un point de $Y$.

Pour tout $Y$-préschéma $Y'$, posons $X(Y')=\operatorname{Hom}_Y(Y',X)$. Soit $x$ un point de
$f^{-1}(y)$~; afin que $f$ soit universellement ouvert au point $x$, il faut et il
suffit que la condition suivante soit vérifiée~:

Pour tout anneau de valuation discrète complet $A$, de corps résiduel
algébriquement clos, tout morphisme $g:Y'=\operatorname{Spec}(A)\to Y$ tel que l'image par $g$ du
point fermé $y'$ de $Y'$ soit égal à $y$, et tout élément
$u_0\in X(\operatorname{Spec}(k(y')))$ tel que $u_0(y')=x$, il existe un anneau de valuation
discrète $B$, un homomorphisme local $A\to B$ faisant de $B$ une $A$-algèbre finie,
et, si l'on pose $Z'=\operatorname{Spec}(B)$, un élément $u\in X(Z')$ tels que, si $z'$ est le
point fermé de $Z'$, le diagramme
\[
  \xymatrix{
    \operatorname{Spec}(k(z')) \ar[r] \ar[d]^{\ell} & Z' \ar[d]^{u} \\
    \operatorname{Spec}(k(y')) \ar[r]_{u_0} & X
  }
\]
soit commutatif.
\end{proposition}

\oldpage[IV-3]{219}

Notons que si $A$ et $B$ vérifient les conditions de l'énoncé, $B$ est un anneau
de valuation discrète complet (Bourbaki, \emph{Alg. comm.}, chap.~III, \S~3,
n\textsuperscript{o}~3, prop.~7 et chap.~IV, \S~2, n\textsuperscript{o}~2,
cor.~3 de la prop.~9) de corps résiduel isomorphe à celui de $A$, donc
algébriquement clos, et puisque $u(z')=x$, il y a dans $X''=X\times_Y Z'$ un point
$x''$, dont les projections dans $X$ et $Z'$ sont $x$ et $z'$ et qui est rationnel
sur $k(z')$~; en outre, puisqu'il existe une $Z'$-section $v$ de $X''$ telle que
$v(z')=x''$, l'image par $v$ du point générique $s$ de $Z'$ est une générisation
$t$ de $x''$ dont la projection dans $Z'$ est $s$~; en appliquant (14.3.6), on voit
que la condition de l'énoncé est \emph{suffisante}. Prouvons maintenant qu'elle est
\emph{nécessaire}. Posons $X'=X\times_Y Y'$, $f'=f_{(Y')}:X'\to Y'$. Il y a par
hypothèse un point $x'\in X'$ au-dessus de $x$ et de $y'$ et rationnel sur $k(y')$
(\textbf{I}, 3.3.14), donc fermé dans $f'^{-1}(y')$. En vertu de (14.3.6), il y a
une composante irréductible $T'$ de $X'$ contenant $x'$ et dominant $Y'$, donc
((14.3.8) et (14.3.13)) la restriction $T'\to Y'$ de $f'$ est équidimensionnelle au
point $x$. On déduit donc de (14.5.5) qu'il y a une $A$-algèbre finie $B$ qui est
un anneau local intègre et intégralement clos et domine $A$ (donc est un anneau de
valuation discrète) et, en posant $Z'=\operatorname{Spec}(B)$ et
$X''=X\times_Y Z'=X'\times_{Y'}Z'$, une $Z'$-section $v:Z'\to X''$ telle que
l'image du point fermé $z'$ de $Z'$ par le morphisme composé $Z'\to X''\to X'$ soit
$x'$~; le morphisme composé $u:Z'\to X''\to X$ répond donc à la question.

\begin{proposition}[14.5.9]
\label{IV.14.5.9-fr}
Soient $Y$ un préschéma localement noethérien irréductible, $f:X\to Y$ un
morphisme localement de type fini, $y$ un point géométriquement unibranche de $Y$.
Alors les conditions suivantes sont équivalentes~:
\begin{enumerate}[label=\emph{\alph*)}]
  \item $f$ est universellement ouvert en tout point de $f^{-1}(y)$.
  \item $f$ est ouvert en tout point de $f^{-1}(y)$.
  \item Pour toute composante irréductible $Z$ de $f^{-1}(y)$, de point générique
  $z$, il existe une composante irréductible $X_i$ de $X$ contenant $z$ et
  équidimensionnelle sur $Y$ au point $z$.
  \item Pour tout point fermé $x$ de $f^{-1}(y)$, il existe une partie irréductible
  $X'$ de $X$, localement fermée dans $X$, contenant $x$, et telle que la
  restriction $X'\to Y$ de $f$ soit un morphisme quasi-fini et dominant.
\end{enumerate}
\end{proposition}

L'équivalence de $a)$, $b)$ et $c)$ a déjà été démontrée (14.4.2). Pour prouver
que $a)$ entraîne $d)$, notons qu'en vertu de (14.4.2), $a)$ entraîne qu'il existe
une composante irréductible $Z$ de $X$ contenant $x$ et équidimensionnelle sur $Y$
au point $x$~; l'existence de $X'$ provient alors de (14.5.3) appliqué à la
restriction $Z\to Y$ de $f$. Inversement, supposons $d)$ vérifiée~; en vertu du
critère de Chevalley (14.4.4), la restriction $X'\to Y$ de $f$ est un morphisme
universellement ouvert au point $x$, et a fortiori $f$ est ouvert au point $x$~;
$f$ est donc ouvert en tous les points fermés de $X_y=f^{-1}(y)$. Mais $X_y$ est
un $k(y)$-préschéma localement de type fini, donc un préschéma de Jacobson~;
l'ensemble des points fermés de $X_y$ est donc \emph{dense} dans $X_y$ (10.3.1),
et il résulte de (14.1.4) que $f$ est ouvert en \emph{tous} les points de $X_y$, ce
qui achève de prouver que $d)$ entraîne $b)$.

Le résultat suivant a été mis en évidence par D.~Mumford~:

\begin{proposition}[14.5.10]
\label{IV.14.5.10-fr}
Soient $Y$ un préschéma noethérien, $f:X\to Y$ un morphisme universellement ouvert,
surjectif et localement de type fini. Alors il existe un morphisme fini surjectif
$g:Y'\to Y$ tel que, si l'on pose $X'=X\times_Y Y'$ et
$f'=f_{(Y')}:X'\to Y'$, tout point $y'\in Y'$ admette un voisinage ouvert $U'$ tel
qu'il existe une $U'$-section de $f'^{-1}(U')$.
\end{proposition}

Nous démontrerons la proposition en plusieurs étapes.

\oldpage[IV-3]{220}

\emph{I) Réduction au cas où $Y$ est intègre.} -- Si l'on a prouvé la proposition
pour chacun des sous-préschémas réduits $Y_i$ ayant pour espace sous-jacent une
composante irréductible de $Y$, et pour les images réciproques $f^{-1}(Y_i)$, il est
clair que le préschéma $Y'$ \emph{somme} des $Y'_i$ correspondants répondra à la
question. On peut donc supposer $Y$ \emph{intègre} et on va dans ce qui suit se
borner à ce cas. Alors, dans la conclusion, on pourra aussi prendre $Y'$ intègre
(quitte à le remplacer par une composante irréductible convenable).

\emph{II) Caractère local sur $Y$.} -- Nous allons montrer que si l'on peut
recouvrir $Y$ par des ouverts $U_j$ en nombre fini, tels que, pour tout $j$, la
conclusion de la proposition soit vraie pour le morphisme $f^{-1}(U_j)\to U_j$,
restriction de $f$, alors la conclusion est vraie également pour $f$. En effet, on
peut évidemment supposer les $U_j$ affines, de sorte que $U_j=\operatorname{Spec}(A_j)$, où
$A_j$ est un anneau intègre noethérien dont le corps des fractions $K=R(Y)$ est le
corps des fonctions rationnelles sur $Y$. Pour tout $j$, il y a par hypothèse une
$A_j$-algèbre finie intègre $A'_j$ telle que l'homomorphisme $A_j\to A'_j$ soit
injectif (\textbf{I}, 1.2.7) et que le morphisme correspondant
$g_j:U'_j=\operatorname{Spec}(A'_j)\to\operatorname{Spec}(A_j)=U_j$ vérifie les conditions de la proposition
(pour $U_j$ et $f^{-1}(U_j)$). Soit alors $K'$ une extension finie de $K$ contenant
les corps des fractions de tous les $A'_j$ (qui sont des extensions finies de $K$).
Considérons le normalisé $Y''$ de $Y$ dans $K'$ (\textbf{II}, 6.3.8), qui est de la
forme $\operatorname{Spec}(\mathcal B)$, où $\mathcal B$ est une $\mathcal O_Y$-Algèbre entière
quasi-cohérente, fermeture intégrale de $\mathcal O_Y$ dans $K'$ (\textbf{II},
6.3.4). Ces définitions prouvent que pour tout $j$,
\[
  A'_j=\Gamma(U'_j,\mathcal O_{U'_j})
      =\Gamma(U_j,(g_j)_*(\mathcal O_{U'_j}))
\]
s'identifie à une sous-$A_j$-algèbre finie de $\Gamma(U_j,\mathcal B)$~; autrement
dit, $(g_j)_*(\mathcal O_{U'_j})=\mathcal A'_j$ est une
$\mathcal O_{U_j}$-Algèbre cohérente, sous-Algèbre de $\mathcal B|U_j$. Il existe
donc un sous-$\mathcal O_Y$-Module $\mathcal C'_j$ cohérent de $\mathcal B$ tel que
$\mathcal C'_j|U_j=\mathcal A'_j$ (\textbf{I}, 9.4.7). Si l'on pose
$\mathcal C'=\sum_j\mathcal C'_j$, la sous-$\mathcal O_Y$-Algèbre $\mathcal B'$ de
$\mathcal B$ engendrée par $\mathcal C'$ est cohérente puisque $\mathcal B$ est une
$\mathcal O_Y$-Algèbre entière. Montrons alors que $Y'=\operatorname{Spec}(\mathcal B')$ répond à
la question. En effet, il est clair que le morphisme $g:Y'\to Y$ est fini surjectif
et que $Y'$ est intègre~; en outre, pour tout $j$, $\Gamma(U_j,\mathcal B')$ est une
$A'_j$-algèbre finie, autrement dit le morphisme $g^{-1}(U_j)\to U_j$, restriction
de $g$, se factorise en $g^{-1}(U_j)\to U'_j\to U_j$, et comme l'existence locale
de sections est stable par changement de base, cela établit notre assertion, tout
$y\in Y$ appartenant à un $U_j$.

\emph{III) Réduction au cas où $Y$ est intègre, local et géométriquement
unibranche.} -- Supposons d'abord que la proposition ait été prouvée lorsque $Y$
est intègre et \emph{local} (avec $Y'$ intègre), et montrons qu'elle est valable
lorsque $Y$ est intègre (noethérien) affine quelconque. En effet, en vertu de la
réduction II), il suffit de prouver que pour tout point $y\in Y$, la proposition est
vraie pour un voisinage ouvert affine $V$ de $y$ dans $Y$. Soient $Y=\operatorname{Spec}(A)$,
$Y_1=\operatorname{Spec}(A_{\mathfrak p})$, où $\mathfrak p=\mathfrak j_y$, et posons
$X_1=X\times_Y Y_1$~; par hypothèse, il existe un morphisme fini surjectif
$g_1:Y'_1\to Y_1$, où $Y'_1=\operatorname{Spec}(B_1)$, $B_1$ étant une
$A_{\mathfrak p}$-algèbre intègre finie, donc un anneau semi-local, tel que $g_1$
vérifie les conditions de l'énoncé pour $Y_1$ et $X_1$. Si $y'_j$
($1\leq j\leq r$) sont les points fermés de $Y'_1$, il y a donc un recouvrement de
$Y'_1$ par des ouverts $U'_j$ tels que $y'_j\in U'_j$ et qu'il existe une
$U'_j$-section $h'_j$ de $X\times_Y U'_j$ ($1\leq j\leq r$). Le
$A_{\mathfrak p}$-module $B_1$ admet un système fini de générateurs de la forme
$z_i/s$ (avec $s\in A-\mathfrak p$,

\oldpage[IV-3]{221}

$z_i$ entiers sur $A$), que l'on peut supposer (en multipliant au besoin $s$ par un
élément de $A$) être des éléments du corps des fractions de $B_1$, entiers sur
$A_s$, de sorte que si $V$ est l'ouvert affine $D(s)=\operatorname{Spec}(A_s)\subset Y$, $Y'_1$
s'identifie à $Y_1\times_V V'$, où $V'$ est le spectre de la $A_s$-algèbre finie
engendrée par les $z_i/s$~; $g:V'\to V$ est donc un morphisme fini surjectif et
$g_1=g_{(Y_1)}$. De plus, appliquant la méthode de (8.1.2, $a)$), on peut supposer
que chacun des $U'_j$ est l'image réciproque par $p:Y'_1\to V'$ d'un ouvert $W'_j$
de $V'$, tel que les $W'_j$ recouvrent $V'$ (8.3.11), et que chacune des sections
$h'_j$ est de la forme $(v'_j)_{(Y_1)}$, où $v'_j$ est une $W'_j$-section de
$X\times_Y W'_j$ (8.8.2, (i)). On est donc bien ramené à prouver la proposition
lorsque $Y=\operatorname{Spec}(A)$, $A$ étant un anneau local intègre et noethérien. On sait alors
qu'il existe une $A$-algèbre intègre finie $B$, ayant même corps des fractions que
$A$, telle que $A\subset B$ et que $\operatorname{Spec}(B)$ soit géométriquement unibranche
(($\mathbf 0$, 23.2.5) et (6.15.5))~; comme le morphisme
$\operatorname{Spec}(B)\to\operatorname{Spec}(A)$ est surjectif, on peut évidemment remplacer $Y$ par
$\operatorname{Spec}(B)$ et $X$ par $X\otimes_A B$ pour prouver la proposition. Raisonnant comme
au début de la réduction III), on peut donc supposer $A$ local, intègre et
$Y=\operatorname{Spec}(A)$ géométriquement unibranche.

\emph{IV) Réduction au cas où $X$ est intègre, affine, et $f$ quasi-fini,
surjectif, birationnel et universellement ouvert.} -- Supposons donc $Y=\operatorname{Spec}(A)$
intègre, local et géométriquement unibranche. Il existe alors un sous-préschéma
irréductible $X_0$ de $X$ tel que la restriction $f_0:X_0\to Y$ de $f$ soit un
morphisme quasi-fini et dominant et que $f_0(X_0)$ contienne le point fermé $y$ de
$Y$ (14.5.9)~; comme $Y$ est géométriquement unibranche, il résulte de (14.4.1) que
$f_0$ est encore \emph{universellement ouvert}. Comme en outre on peut supposer
$X_0$ réduit, donc intègre, on voit qu'on peut, en remplaçant $X$ par $X_0$,
supposer que $X$ est \emph{intègre} et $f$ \emph{quasi-fini et dominant}, et tel que
$f(X)$ contienne $y$~; comme $f$ est \emph{ouvert} et que tout ouvert de $Y$
contenant $y$ est égal à $Y$, $f$ est surjectif.

Soient $\xi$, $\eta$ les points génériques de $X$ et $Y$ respectivement~;
$k(\xi)$ est donc une extension \emph{finie} de $K=k(\eta)$. Par suite (4.6.8), il
y a une extension \emph{finie} $K'$ de $K$ telle que
$(\operatorname{Spec}(k(\xi)\otimes_K K'))_{\mathrm{red}}$ soit géométriquement réduit sur $K$ et
que ses composantes irréductibles soient géométriquement irréductibles~; il en
résulte ((4.5.9) et (4.6.1)) que les corps résiduels de
$\operatorname{Spec}(k(\xi)\otimes_K K')$ sont des extensions finies, primaires et séparables de
$K'$, donc sont \emph{égaux à} $K'$. Appliquant de nouveau ($\mathbf 0$, 23.2.5) et
(6.15.5), il existe une $A$-algèbre intègre \emph{finie} $B$, ayant $K'$ pour
corps des fractions, contenant $A$ et telle que $\operatorname{Spec}(B)$ soit géométriquement
unibranche~; on peut donc de nouveau remplacer $Y$ par $\operatorname{Spec}(B)$ et $X$ par
$(X\otimes_A B)_{\mathrm{red}}$ pour prouver la proposition~; la réduction III)
permet alors de supposer encore $A$ local, intègre et $Y=\operatorname{Spec}(A)$ géométriquement
unibranche, de point fermé $y$. En outre, $f$ est alors un morphisme \emph{quasi-fini,
surjectif et universellement ouvert}~; chacune des composantes irréductibles $X_i$
de $X$ ($1\leq i\leq m$) domine $Y$ (1.10.4) et est \emph{birationnelle} sur $Y$.
Soit alors $x$ un point de $f^{-1}(y)$ et soit $X_i$ une composante irréductible de
$X$ contenant $x$~; désignons encore par $X_i$ le sous-préschéma réduit (donc
intègre) de $X$ ayant $X_i$ pour espace sous-jacent, et soit $f_i:X_i\to Y$ la
restriction de $f$. Appliquant de nouveau (14.4.1) au morphisme quasi-fini et
dominant $f_i$, on voit que $f_i$ est universellement ouvert~; si $U$ est un ouvert
affine de $X_i$ contenant $x$, $f_i(U)$ est donc

\oldpage[IV-3]{222}

ouvert dans $Y$ et contient $y$, donc est égal à $Y$. On peut ainsi remplacer $X$
par $U$ pour prouver la proposition.

\emph{V) Fin de la démonstration.} -- Nous supposons donc $Y$ local, intègre et
géométriquement unibranche, $X$ intègre et affine, $f$ quasi-fini, surjectif,
birationnel et universellement ouvert~; il en résulte que $f$ est automatiquement
\emph{séparé}. En vertu du Main theorem (8.12.6), il existe une factorisation
$X\xrightarrow{j}Z\xrightarrow{u}Y$, où $j$ est une immersion ouverte et $u$ un
morphisme \emph{fini}. Remplaçant en outre $Z$ par l'image fermée de $j$
(\textbf{I}, 9.5.10), on peut supposer que $Z$ est \emph{intègre}~; comme $u$ est
fini et birationnel et $Y$ géométriquement unibranche, il résulte de
(\textbf{III}, 4.3.5 et 4.3.4) que $u$ (et par suite $f$) est un morphisme
\emph{radiciel}~; comme $f$ est universellement ouvert et surjectif, c'est donc un
\emph{homéomorphisme universel}. Par suite (2.4.5, (ii)) $f$ est un morphisme
\emph{fini}. Mais alors on répond aux conditions de l'énoncé avec $Y'$ \emph{intègre}
en prenant simplement $Y'=X$ et $g=f$, la $Y'$-section étant le morphisme diagonal
$\Delta_f$. C.Q.F.D.

Ce résultat peut être utilisé pour développer des critères de «~descente~» de
diverses propriétés par les morphismes universellement ouverts et surjectifs.
Signalons en particulier le critère suivant, dû à D.~Mumford~:

\begin{corollary}[14.5.11]
\EGAReaderTarget{ega4.c04.cor.14-5-11}{IV.14.5.11-fr}
Soient $Y$ un préschéma localement noethérien, $f:X\to Y$ un morphisme localement
de type fini, $g:Y'\to Y$ un morphisme localement de type fini, universellement
ouvert et surjectif~; posons $X'=X\times_Y Y'$, $f'=f_{(Y')}:X'\to Y'$. Alors,
pour que $f$ soit affine, il faut et il suffit que $f'$ le soit.
\end{corollary}

On doit seulement prouver que la condition est suffisante. La question étant locale
sur $Y$, on peut supposer $Y$ affine, donc noethérien. En vertu de (14.5.10), il
existe un morphisme fini surjectif $h:Y_1\to Y$ tel que, si l'on pose
$Y'_1=Y'\times_Y Y_1$, et $g'=g_{(Y_1)}:Y'_1\to Y_1$, tout point de $Y_1$ admette
un voisinage ouvert $U_1$ tel qu'il existe une $U_1$-section de $g'^{-1}(U_1)$. Si
l'on pose $X_1=X\times_Y Y_1$, la projection canonique $p:X_1\to X$ est un
morphisme fini surjectif, donc, si l'on prouve que $X_1$ est un schéma affine, il en
résultera d'abord que $X$ est quasi-compact, donc noethérien, puis que $X$ est affine
en vertu du théorème de Chevalley (\textbf{II}, 6.7.1). Il suffit par suite de
prouver que le morphisme $f_1=f_{(Y_1)}:X_1\to Y_1$ est affine. Or, si l'on pose
\[
  X'_1=X_1\times_{Y_1}Y'_1=X'\times_{Y'}Y'_1
  \qquad\text{et}\qquad
  f'_1=(f_1)_{(Y'_1)}=f'_{(Y'_1)}:X'_1\to Y'_1,
\]
$f'_1$ est affine en vertu de l'hypothèse. On est donc ramené à prouver le corollaire
en y remplaçant $Y$, $X$, $f$ et $Y'$ par $Y_1$, $X_1$, $f_1$ et $Y'_1$,
autrement dit, il suffit de prouver que $f$ est affine lorsqu'on fait de plus, dans
l'énoncé de (14.5.11), l'hypothèse que tout point de $Y$ admet un voisinage ouvert
$U$ tel qu'il existe une $U$-section de $g^{-1}(U)$. La question étant locale sur
$Y$, on peut même supposer qu'il existe une $Y$-section $s$ de $Y'$. Or, on a le
lemme élémentaire suivant (valable dans toute catégorie admettant des produits
fibrés)~:

\begin{lemma}[14.5.11.1]
\label{IV.14.5.11.1-fr}
Soient $f:X\to S$, $g:Y\to S$ deux morphismes, $p_1:X\times_S Y\to X$,
$p_2:X\times_S Y\to Y$ les projections canoniques. Si $s:S\to Y$ est une section
de $g$, alors $s'=(1,s\circ f)_S$ est une section de $p_1$ et $X$, muni des
morphismes $f:X\to S$ et $s':X\to X\times_S Y$, s'identifie au produit des
$Y$-préschémas $S$ et $X\times_S Y$ pour les morphismes $s:S\to Y$ et
$p_2:X\times_S Y\to Y$.
\end{lemma}

\oldpage[IV-3]{223}

\begin{proof}
C'est un cas particulier de (\textbf{I}, 3.3.11), où l'on remplace le diagramme par
\[
  \xymatrix{
    X \ar[r]^{p_1} \ar[d]_{f} & X\times_S Y \ar[r]^{s'} \ar[d]^{p_2} & X \ar[d]^{f} \\
    S \ar[r] & Y \ar[r]_{s} & S
  }
\]
\end{proof}

Appliquant ce lemme en y remplaçant $S$, $Y$ par $Y$, $Y'$, on voit que l'on peut
écrire $f=(f')_{(Y)}$ pour le changement de base $s:Y\to Y'$, donc $f$ est affine
puisque $f'$ l'est. C.Q.F.D.

Une variante de ce critère est la suivante~:

\begin{corollary}[14.5.12]
\label{IV.14.5.12-fr}
Avec les notations et hypothèses générales de (14.5.11), soit $\mathcal L$ un
$\mathcal O_X$-Module inversible. Supposons qu'il existe une partie fermée $Z$ de
$X$, propre sur $Y$ (\textbf{II}, 5.4.10) et telle que le préschéma induit par $X$
sur l'ouvert $X-Z$ soit normal. Alors, pour que $\mathcal L$ soit ample relativement
à $f$, il faut et il suffit que
$\mathcal L'=\mathcal L\otimes_{\mathcal O_Y}\mathcal O_{Y'}$ soit ample
relativement à $f'$.
\end{corollary}

Gardons les notations de la démonstration de (14.5.11). Posons
$\mathcal L_1=\mathcal L\otimes_{\mathcal O_Y}\mathcal O_{Y_1}$~; il suffira de
prouver que $\mathcal L_1$ est ample relativement à $f\circ p$, en vertu de
(\textbf{III}, 2.6.2)~; mais $f\circ p=h\circ f_1$, et, compte tenu de
(\textbf{II}, 4.6.13, (v)), il suffira de prouver que $\mathcal L_1$ est ample
relativement à $f_1$. La question étant locale sur $Y_1$, on peut de nouveau
supposer que $g'$ admet une section $s$~; si
$\mathcal L'_1=\mathcal L'\otimes_{\mathcal O_{Y'}}\mathcal O_{Y'_1}$, on peut
alors écrire, en vertu du lemme (14.5.11.1),
$\mathcal L_1=\mathcal L'_1\otimes_{\mathcal O_{Y'_1}}\mathcal O_{Y_1}$ pour le
changement de base $s:Y_1\to Y'_1$~; la conclusion résulte donc de deux
applications de (\textbf{II}, 4.6.13, (iii)).

% IV3_B31_P223_END_BEFORE_SECTION_15

% FIN SOURCE ega4/ega4-14-fr.tex

%
% DEBUT SOURCE ega4/ega4-15-fr.tex
\section{Étude des fibres d'un morphisme universellement ouvert}
\label{section:IV.15-fr}

\subsection{Multiplicités des fibres d'un morphisme universellement ouvert}
\label{subsection:IV.15.1-fr}

\begin{proposition}[15.1.1]
\label{IV.15.1.1-fr}
\oldpage[IV-3]{223}
Soient $Y$ un préschéma irréductible localement noethérien, de point générique
$\eta$, $f:X\to Y$ un morphisme localement de type fini, $y$ un point de $Y$,
$\mathcal F$ un $\mathcal O_X$-Module cohérent~; on pose
$\mathcal F_y=\mathcal F\otimes_{\mathcal O_Y}\mathbf{k}(y)$ (faisceau de modules
sur la fibre $f^{-1}(y)$). Soit $z$ le point générique d'une composante
irréductible de $\operatorname{Supp}(\mathcal F_y)$, et soient $X_i$
($1\leqslant i\leqslant n$) ceux des sous-préschémas fermés réduits de $X$ dont
l'espace sous-jacent est une composante irréductible de
$\operatorname{Supp}(\mathcal F)$ contenant $z$, et qui sont tels que la
restriction $X_i\to Y$ de $f$ soit un morphisme universellement ouvert au point
$z$~; soit $z_i$ le point générique de $X_i$ ($1\leqslant i\leqslant n$). Si
$\lambda_x(\mathcal F_y)$ (resp. $\lambda_t(\mathcal F_\eta)$) désigne la longueur
géométrique de $\mathcal F_y$ en un point $x\in f^{-1}(y)$ (resp. celle de
$\mathcal F_\eta$ en un point $t\in f^{-1}(\eta)$) (4.7.5), on a
\[
\tag{15.1.1.1}
\label{IV.15.1.1.1-fr}
  \lambda_z(\mathcal F_y)\geqslant
  \sum_{i=1}^{n}\lambda_{z_i}(\mathcal F_\eta).
\]
Si de plus on suppose l'anneau $\mathcal O_y$ régulier, on a
\[
\tag{15.1.1.2}
\label{IV.15.1.1.2-fr}
  \operatorname{long}((\mathcal F_y)_z)\geqslant
  \sum_{i=1}^{n}\operatorname{long}((\mathcal F_\eta)_{z_i}).
\]
\end{proposition}

\begin{proof}
\oldpage[IV-3]{224}
Les fibres $f^{-1}(y')$ de $f$ (pour tout $y'\in Y$) ne changeant pas lorsque
l'on remplace $f$ par $f_{\mathrm{red}}$, on peut supposer que $X$ et $Y$ sont
réduits (2.4.3, (vi)), donc $Y$ intègre~; on posera
$K=\mathbf{k}(\eta)$. On peut en outre remplacer $f$ par la restriction
$\operatorname{Supp}(\mathcal F)\to Y$ de $f$ au sous-préschéma fermé réduit de
$X$ ayant $\operatorname{Supp}(\mathcal F)$ comme espace sous-jacent, donc se
borner au cas où $X=\operatorname{Supp}(\mathcal F)$. Enfin, si $y=\eta$, il n'y
a qu'une seule composante irréductible $X_i$ de $X$ contenant $z$
($\mathbf{0}_{\mathrm I}$, 2.1.8) et on a $z_i=z$, ce qui rend (15.1.1.1) et
(15.1.1.2) triviales (sans hypothèses sur $f$). On peut donc se borner au cas
$y\ne\eta$~; il résulte alors de l'hypothèse et de (1.10.3) que les $z_i$
appartiennent à $f^{-1}(\eta)$, les seconds membres de (15.1.1.1) et
(15.1.1.2) étant donc définis. Soit $Z_i$ le sous-préschéma fermé réduit de
$f^{-1}(\eta)$ ayant pour espace sous-jacent $X_i\cap f^{-1}(\eta)$~; on sait
(4.6.6) qu'il existe une extension finie radicielle $K'$ de $K$ telle que, pour
tout $i$, le $K'$-préschéma $(Z_i\otimes_K K')_{\mathrm{red}}$ soit
géométriquement réduit sur $K'$. En outre, le morphisme projection
$f^{-1}(\eta)\otimes_K K'\to f^{-1}(\eta)$ est fini, dominant et radiciel
(\textbf{I}, 3.5.7), donc est un homéomorphisme (2.4.5). Si $z_i'$ est l'unique
point de $f^{-1}(\eta)\otimes_K K'$ au-dessus de $z_i$, il résulte de (4.7.9) et
(4.7.5) que l'on a
\[
\tag{15.1.1.3}
\label{IV.15.1.1.3-fr}
  \lambda_{z_i}(\mathcal F_\eta)
  =\lambda_{z_i'}(\mathcal F_\eta\otimes_K K')
  =\operatorname{long}((\mathcal F_\eta\otimes_K K')_{z_i'}).
\]

Soit $V$ un anneau de valuation discrète, de corps de fractions $K'$, dominant
$\mathcal O_y$ (\textbf{II}, 7.1.7)~; posons $Y'=\operatorname{Spec}(V)$,
$X'=X\times_Y Y'$, $\mathcal F'=\mathcal F\otimes_Y Y'$,
$f'=f_{(Y')}:X'\to Y'$~; si $g:Y'\to Y$ est le morphisme structural, $\eta'$ le
point générique de $Y'$ et $y'$ son point fermé, on a $g(y')=y$, $g(\eta')=\eta$~;
le morphisme $\operatorname{Spec}(\mathbf{k}(y'))\to
\operatorname{Spec}(\mathbf{k}(y))$ étant fidèlement plat, il en est de même de
la projection $g':f'^{-1}(y')\to f^{-1}(y)$ (2.2.13), donc (2.3.4), il y a un
point générique $z'$ d'une composante irréductible de $f'^{-1}(y')$ tel que
$g'(z')=z$. Par construction, on a $\mathbf{k}(\eta')=K'$, donc
$f'^{-1}(\eta')=f^{-1}(\eta)\otimes_K K'$~; d'autre part, si l'on pose
$X_i'=X_i\times_Y Y'$, la restriction $X_i'\to Y'$ de $f'$ est un morphisme
universellement ouvert au point $z'$, donc (1.10.3) il existe un point
$z_i''\in f'^{-1}(\eta')$ qui est une générisation de $z'$, et l'on peut
évidemment supposer que $z_i''$ est point générique d'une composante
irréductible de $X_i'\cap f'^{-1}(\eta')$~; comme la projection de
$X_i'\cap f'^{-1}(\eta')$ dans $f^{-1}(\eta)$ est $Z_i$, on a $z_i''=z_i'$.
En vertu de (4.7.9) et (4.7.5), on a
\[
\tag{15.1.1.4}
\label{IV.15.1.1.4-fr}
  \lambda_z(\mathcal F_y)=\lambda_{z'}(\mathcal F'_{y'})
  \geqslant\operatorname{long}((\mathcal F'_{y'})_{z'})
\]
et il résulte donc de cette inégalité et de (15.1.1.3) qu'il suffit, pour établir
(15.1.1.1), de démontrer l'inégalité
\[
\tag{15.1.1.5}
\label{IV.15.1.1.5-fr}
  \operatorname{long}((\mathcal F'_{y'})_{z'})\geqslant
  \sum_{i=1}^{n}\operatorname{long}((\mathcal F'_{\eta'})_{z_i'}).
\]

Considérons maintenant la seconde inégalité (15.1.1.2)~; nous utiliserons le
lemme suivant~:

\begin{lemma}[15.1.1.6]
\label{IV.15.1.1.6-fr}
Soient $A$ un anneau local régulier qui n'est pas un corps, $K$ son corps des
fractions, $k$ son corps résiduel. Il existe un anneau de valuation discrète $V$
dominant $A$, ayant $K$ pour corps des fractions et dont le corps résiduel est
une extension transcendante pure de $k$.
\end{lemma}

\oldpage[IV-3]{225}
Posons $S=\operatorname{Spec}(A)$, et considérons le $S$-schéma $P$ obtenu en
faisant éclater le point fermé $s$ de $S$ (\textbf{II}, 8.1.3)~; si $\mathfrak m$
est l'idéal maximal de $A$, on a donc par définition
$P=\operatorname{Proj}(B)$, où $B$ est la $A$-algèbre graduée
$\bigoplus_{n\geqslant0}\mathfrak m^n$. Si $h:P\to S$ est le morphisme
structural, la fibre $h^{-1}(s)$ est isomorphe à
$\operatorname{Proj}(B\otimes_A k)$ (\textbf{II}, 2.8.10), et par définition
$B\otimes_A k$ est la $k$-algèbre graduée
$\bigoplus_{n\geqslant0}\mathfrak m^n/\mathfrak m^{n+1}$~; mais comme $A$ est
régulier, cette algèbre est isomorphe à une algèbre de polynômes
$k[t_1,\ldots,t_e]$ ($t_i$ indéterminées) ($\mathbf{0}_{\mathrm{IV}}$, 17.1.1)
et par suite $h^{-1}(s)$ est isomorphe à $\mathbf P_k^{e-1}$. Soit $\zeta$ le
point générique de $h^{-1}(s)$~; si $t_i'$ ($1\leqslant i\leqslant e$) sont les
éléments de $\mathfrak m$ dont les classes mod. $\mathfrak m^2$ sont les $t_i$,
$B_{(t_e')}$ est l'anneau d'un voisinage ouvert affine de $\zeta$ dans $P$, et
l'on a $\mathbf{k}(\zeta)=k(t_1,\ldots,t_{e-1})$~; d'autre part, comme $A$ est
intégralement clos, on vérifie aisément qu'il en est de même de $B$ (Bourbaki,
\emph{Alg. comm.}, chap.~V, \S~1, n\textsuperscript{o}~8, cor.~1 de la prop.~21),
donc $B_{(t_e')}$ est aussi intégralement clos, et il en est par suite de même de
l'anneau local $\mathcal O_{P,\zeta}$. Enfin, comme $A=\mathcal O_s$ est
régulier, donc un anneau universellement caténaire (5.6.4), on a, en vertu de
(5.6.5),
$\dim(\mathcal O_{P,\zeta})=\dim(\mathcal O_s)-(e-1)$, puisque $h$ est un
morphisme birationnel, que $\dim(h^{-1}(s))=e-1$ et que l'anneau local de la
fibre $h^{-1}(s)$ en son point générique $\zeta$ est de dimension $0$. Mais
$\dim(\mathcal O_s)=e$, donc $\dim(\mathcal O_{P,\zeta})=1$, et
$\mathcal O_{P,\zeta}=V$ est un anneau de valuation discrète, étant noethérien,
de dimension $1$ et intégralement clos (\textbf{II}, 7.1.6)~; il répond
évidemment à la question.

Ce lemme étant établi, on peut reprendre le raisonnement fait pour l'inégalité
(15.1.1.1) avec $Y'=\operatorname{Spec}(V)$~; cette fois $\mathbf{k}(y')$ est
une extension \emph{séparable} de $\mathbf{k}(y)$, et par suite (4.7.9), on a
$\operatorname{long}((\mathcal F_y)_z)=
\operatorname{long}((\mathcal F'_{y'})_{z'})$~; comme par ailleurs
$f'^{-1}(\eta')=f^{-1}(\eta)$ et $\mathcal F'_{\eta'}=\mathcal F_\eta$, on est
encore ramené à prouver l'inégalité (15.1.1.5). Or, si $t$ est une uniformisante
de $\mathcal O_{y'}$, $f'^{-1}(y')$ est le sous-préschéma fermé de $X'$ défini
par l'idéal $t\mathcal O_{X'}$, et on a
$\mathcal F'_{y'}=\mathcal F'/t\mathcal F'$~; d'autre part, on vérifie aussitôt
(\textbf{I}, 9.1.12) que
$(\mathcal F'_{\eta'})_{z_i'}=\mathcal F'_{z_i'}$~; l'inégalité à démontrer
(15.1.1.5) n'est autre alors qu'un cas particulier de (3.4.1.1).
\end{proof}

\begin{corollary}[15.1.2]
\label{IV.15.1.2-fr}
Les hypothèses étant celles de (15.1.1), supposons de plus que
$\lambda_z(\mathcal F_y)=1$ (resp. que $\mathcal O_y$ soit régulier et
$\operatorname{long}((\mathcal F_y)_z)=1$). Alors il existe au plus une composante
irréductible $X_i$ de $\operatorname{Supp}(\mathcal F)$ contenant $z$ et telle que
la restriction $X_i\to Y$ de $f$ soit universellement ouverte au point $z$. De
plus, si $z_i$ est le point générique de cette composante, on a
$\lambda_{z_i}(\mathcal F_\eta)=1$ (resp.
$\operatorname{long}((\mathcal F_\eta)_{z_i})=1$).
\end{corollary}

Cela résulte aussitôt de (15.1.1) en notant que l'on a nécessairement, par
définition de $X_i$, $(\mathcal F_\eta)_{z_i}\ne0$ (\textbf{I}, 9.1.13).

\begin{corollary}[15.1.3]
\label{IV.15.1.3-fr}
Soient $Y$ un préschéma irréductible localement noethérien de point générique
$\eta$, $f:X\to Y$ un morphisme localement de type fini, $\mathcal F$ un
$\mathcal O_X$-Module cohérent de support $X$, $x$ un point de $X$, $y=f(x)$. On
suppose que~: $1^\circ$ $x$ n'appartient qu'à une seule composante irréductible
$Z$ de $f^{-1}(y)$~; $2^\circ$
$\mathcal F_y=\mathcal F\otimes_{\mathcal O_Y}\mathbf{k}(y)$ est géométriquement
réduit sur $\mathbf{k}(y)$, autrement dit, si $z$ est le point générique de $Z$,
$\lambda_z(\mathcal F_y)=1$ (resp. $\mathcal O_y$ est régulier et
$\operatorname{long}((\mathcal F_y)_z)=1$). Alors, il existe au plus une
composante irréductible $X'$ de $X$ contenant $x$, telle que
$\dim_x(X'\cap f^{-1}(y))=\dim_x(f^{-1}(y))$ et que la restriction $X'\to Y$ de
$f$ soit universellement ouverte aux points génériques des composantes
irréductibles de $X'\cap f^{-1}(y)$ contenant $x$. De plus, s'il existe une telle
composante $X'$ et si $z'$ est son point générique, on a
$\lambda_{z'}(\mathcal F_\eta)=1$ (resp.
$\operatorname{long}((\mathcal F_\eta)_{z'})=1$).
\end{corollary}

\oldpage[IV-3]{226}
En effet, une composante irréductible $X'$ de $X$ contenant $x$ et telle que
$\dim_x(X'\cap f^{-1}(y))=\dim_x(f^{-1}(y))$ contient nécessairement $Z$,
puisqu'une composante irréductible de $X'\cap f^{-1}(y)$ contenant $x$ doit être
de même dimension que la seule composante irréductible $Z$ de $f^{-1}(y)$
contenant $x$. On peut alors appliquer (15.1.2) à $z$.

\begin{remarks}[15.1.4]
\label{IV.15.1.4-fr}
\begin{enumerate}[label=(\roman*)]
  \item Dans l'énoncé de (15.1.1), on ne peut remplacer «~universellement
  ouvert~» par «~équidimensionnel~», comme le montre l'exemple (14.4.10, (ii)) où
  l'on fait $\mathcal F=\mathcal O_X$~; les fibres de $f$ sont alors des schémas
  artiniens réduits, donc (avec les notations introduites \emph{loc. cit.}) on a
  $\lambda_{x'}(\mathcal F_y)=1$, mais il y a deux composantes irréductibles
  $X_1$, $X_2$ de $X$ contenant $x'$, et le second membre de (15.1.1.1) est donc
  égal à $2$, bien que la restriction de $f$ à chacune des composantes $X_1$,
  $X_2$ soit un morphisme équidimensionnel en tout point. Ce même exemple montre
  aussi que dans (15.1.2) et (15.1.3), on ne peut remplacer l'hypothèse
  «~universellement ouvert~» par «~équidimensionnel~».
  \item Il est vraisemblable que pour la validité de l'inégalité (15.1.1.2), on
  ne peut supprimer l'hypothèse que $\mathcal O_y$ est \emph{régulier}.
\end{enumerate}
\end{remarks}

\subsection{Platitude des morphismes universellement ouverts à fibres géométriquement réduites}
\label{subsection:IV.15.2-fr}

\begin{env}[15.2.1]
\label{IV.15.2.1-fr}
On a vu (2.4.6) qu'un morphisme \emph{plat} et localement de présentation finie
est \emph{universellement ouvert}. On a une réciproque partielle de ce résultat
moyennant des hypothèses supplémentaires~:
\end{env}

\begin{theorem}[15.2.2]
\label{IV.15.2.2-fr}
Soient $Y$ un préschéma localement noethérien, $f:X\to Y$ un morphisme localement
de type fini, $\mathcal F$ un $\mathcal O_X$-Module cohérent, $x$ un point de
$X$, $y=f(x)$. On suppose vérifiées les conditions suivantes~:
\begin{enumerate}[label=(\roman*)]
  \item $\operatorname{Supp}(\mathcal F)=X$.
  \item $f$ est universellement ouvert aux points génériques des composantes
  irréductibles de $f^{-1}(y)$ contenant $x$.
  \item $\mathcal F_y=\mathcal F\otimes_{\mathcal O_Y}\mathbf{k}(y)$ est
  géométriquement réduit au point $x$ (4.6.22).
  \item L'anneau $\mathcal O_y$ est réduit.
\end{enumerate}
Alors $\mathcal F$ est $f$-plat au point $x$.
\end{theorem}

\begin{proof}
Comme $\mathcal O_y$ est réduit et noethérien, nous allons appliquer le critère
valuatif de platitude (11.8.1). Considérons donc un homomorphisme local
$\mathcal O_y\to A'$, où $A'$ est un anneau de valuation discrète~; posons
$Y'=\operatorname{Spec}(A')$, $X'=X\times_Y Y'$, $f'=f_{(Y')}:X'\to Y'$, et
soit $x'$ un point de $X'$ dont les projections dans $X$ et $Y'$ sont
respectivement $x$ et le point fermé $y'$ de $Y'$~; si l'on pose
$\mathcal F'=\mathcal F\otimes_Y Y'$, on a
$\operatorname{Supp}(\mathcal F')=X'$ par (\textbf{I}, 9.1.13)~; tout point
générique $z'$ d'une composante irréductible de $f'^{-1}(y')$ contenant $x'$ a
pour projection dans $f^{-1}(y)$ un point générique $z$ d'une composante
irréductible de $f^{-1}(y)$ contenant $x$, en vertu de (4.2.6)~; donc $f'$ est
universellement ouvert au point $z'$. Enfin, il résulte de (4.7.11) que
$\mathcal F'_{y'}$ est géométriquement réduit au point $x'$. On voit donc qu'on
est ramené à démontrer le

\oldpage[IV-3]{227}
\begin{lemma}[15.2.2.1]
\label{IV.15.2.2.1-fr}
Soient $Y=\operatorname{Spec}(A)$ le spectre d'un anneau de valuation discrète,
$y$ son point fermé, $f:X\to Y$ un morphisme localement de type fini, $x$ un
point de $f^{-1}(y)$, $\mathcal F$ un $\mathcal O_X$-Module cohérent. On suppose
vérifiées les conditions suivantes~:
\begin{enumerate}[label=(\roman*)]
  \item $\operatorname{Supp}(\mathcal F)=X$.
  \item $f$ est ouvert aux points génériques des composantes irréductibles de
  $f^{-1}(y)$ contenant $x$.
  \item $\mathcal F_y=\mathcal F\otimes_{\mathcal O_y}\mathbf{k}(y)$ est réduit
  au point $x$ (3.2.2).
\end{enumerate}
Alors $\mathcal F$ est $f$-plat au point $x$.
\end{lemma}

Désignons en effet par $t$ une uniformisante de $A$, posons $B=\mathcal O_x$ et
$\mathcal F_x=M$~; la condition (i) implique que
$\operatorname{Supp}(M)=\operatorname{Spec}(B)$ (\textbf{I}, 9.1.13). La
condition (ii) implique, en vertu de (14.3.2), que toute composante irréductible
de $X$ contenant $x$ domine $Y$~; en d'autres termes, l'image réciproque dans $A$
de tout idéal premier minimal $\mathfrak p_i$ de $B$ est $0$, ce qui signifie
encore que l'on a $t\notin\mathfrak p_i$ pour tout $i$. Enfin, (iii) signifie que
$M/tM$ est un $(B/tB)$-module \emph{réduit} (3.2.2). Il s'agit de montrer que
sous ces conditions $M$ est un $A$-module \emph{sans torsion}
($\mathbf{0}_{\mathrm I}$, 6.3.4), et pour cela il suffit évidemment de montrer
que $t$ est un élément $M$-\emph{régulier}~; mais cela résulte de (3.4.6.1).
\end{proof}

\begin{corollary}[15.2.3]
\label{IV.15.2.3-fr}
Soient $Y$ un préschéma localement noethérien, $f:X\to Y$ un morphisme localement
de type fini, $x$ un point de $X$, $y=f(x)$. On suppose vérifiées les conditions
suivantes~:
\begin{enumerate}[label=(\roman*)]
  \item $f$ est universellement ouvert aux points génériques des composantes
  irréductibles de $f^{-1}(y)$ contenant $x$.
  \item $f^{-1}(y)$ est géométriquement réduit (sur $\mathbf{k}(y)$) au point
  $x$ (4.6.9).
  \item L'anneau $\mathcal O_y$ est réduit.
\end{enumerate}
Sous ces conditions $f$ est plat au point $x$.
\end{corollary}

Il suffit d'appliquer (15.2.2) à $\mathcal F=\mathcal O_X$ (4.6.22).

\begin{remarks}[15.2.4]
\label{IV.15.2.4-fr}
\begin{enumerate}[label=(\roman*)]
  \item Les trois premières des conditions de (15.2.2) ne changent pas lorsqu'on
  remplace $f$ par $f_{\mathrm{red}}$~; mais si $\mathfrak N$ est le nilradical
  d'un anneau local noethérien $A$, $A/\mathfrak N$ est plat sur $A/\mathfrak N$
  mais non sur $A$ lui-même lorsque $\mathfrak N\ne0$ (Bourbaki,
  \emph{Alg. comm.}, chap.~II, \S~3, n\textsuperscript{o}~2, cor.~2 de la
  prop.~5). On voit donc qu'on ne peut dans (15.2.2) supprimer la condition (iv).

  \item Il résulte de (2.4.6) que si la conclusion de (15.2.2) est vraie, ainsi
  que l'hypothèse (i), $f$ est universellement ouvert dans un voisinage de $x$,
  et en particulier aux points génériques des composantes irréductibles de
  $f^{-1}(y)$ contenant $x$. Même lorsque (i), (iii) et (iv) sont vérifiées, on
  ne peut remplacer (ii) par l'hypothèse plus faible que $f$ est universellement
  ouvert \emph{au point $x$ seulement}. C'est ce que montre l'exemple
  (14.1.3, (i)), où $f$ est évidemment universellement ouvert en tout point de
  $X_2$, donc en particulier au point $x$, intersection de $X_1$ et $X_2$~; les
  conditions (ii) et (iii) de (15.2.3) sont aussi vérifiées par cet exemple.
  Toutefois, $\mathcal O_x$ n'est pas un $\mathcal O_y$-module sans torsion,
  $\mathcal O_y$ s'identifiant à un sous-anneau de $\mathcal O_x$ qui contient
  des diviseurs de $0$ dans $\mathcal O_x$~; donc $f$ n'est pas plat au point $x$.

  \item Dans (15.2.3), la conclusion n'est plus nécessairement valable lorsqu'on
  remplace l'hypothèse (ii) par l'hypothèse plus faible que $f^{-1}(y)$,
  considéré comme $\mathbf{k}(y)$-préschéma, n'a pas de cycles premiers associés
  immergés. Un exemple est fourni en prenant pour $Y$ une courbe ayant un point
  de rebroussement $y$, par exemple
  $Y=\operatorname{Spec}(K[S,T]/(S^3-T^2))$, $K$ étant un corps algébriquement
  clos, et pour $X$ sa normalisée (\textbf{II}, 6.3.8). Si $s$ et $t$ sont les
  classes de $S$ et $T$ dans $A=K[S,T]/(S^3-T^2)$, on vérifie aussitôt que
  $X=\operatorname{Spec}(B)$, où $B=K[s,t,u]$, $u$ étant l'élément $t/s$ du corps
  des fractions de $A$, et comme $s=u^2$, $t=u^3$, on a aussi $B=K[u]$,
  isomorphe à l'anneau de polynômes à une indéterminée sur $K$. Le seul point $x$
  de $X$ au-dessus du point $y$ correspondant à l'idéal maximal $(s)+(t)$
  correspond à l'idéal maximal $(u)$~; mais comme la classe $\bar u$ de $u$ dans
  $\mathcal O_x/\mathfrak m_y\mathcal O_x$ est telle que $\bar u^2=0$,
  $f^{-1}(y)$ n'est pas un $\mathbf{k}(y)$-préschéma réduit. Ici $f$ est un
  morphisme fini, surjectif et radiciel, donc un homéomorphisme universel (2.4.5),
  mais n'est pas plat au point $x$, car si $\mathcal O_x$ était un
  $\mathcal O_y$-module plat, il serait un $\mathcal O_y$-module
  \oldpage[IV-3]{228}
  libre de rang $1$ engendré par l'élément $1$ de $\mathcal O_x$ (Bourbaki,
  \emph{Alg. comm.}, chap.~II, \S~3, n\textsuperscript{o}~2, prop.~5), ce qui est
  absurde.

  \item Montrons enfin que dans (15.2.2) ou (15.2.3), on ne peut remplacer
  «~géométriquement réduit~» par «~réduit~». Nous allons en effet définir deux
  anneaux $A$, $A'$ ayant les propriétés suivantes~:
  \begin{enumerate}[label=\arabic*)]
    \item $A$ est un anneau local noethérien d'idéal maximal $\mathfrak m$,
    intègre, complet, de dimension $1$ et géométriquement unibranche.
    \item $A'$ est la clôture intégrale de $A$, une $A$-algèbre finie dont
    l'idéal maximal est $\mathfrak mA'$, et le corps résiduel $k'$ une extension
    radicielle finie non triviale du corps résiduel $k=A/\mathfrak m$ de $A$.
  \end{enumerate}
  On peut alors prendre $Y=\operatorname{Spec}(A)$,
  $X=\operatorname{Spec}(A')$, $\mathcal F=\mathcal O_X$, $y$ et $x$ étant les
  points fermés de $Y$ et $X$ respectivement~; les hypothèses (i) et (iv) de
  (15.2.2) sont trivialement vérifiées~; comme $A$ est de dimension $1$,
  $f:X\to Y$ est radiciel, fini et surjectif, donc un homéomorphisme universel
  (2.4.5), ce qui prouve l'hypothèse (i) de (15.2.3). Enfin, il est clair que
  $f^{-1}(y)$ est réduit. Cependant $f$ n'est pas plat au point $x$, car si $A'$
  était un $A$-module plat, ce serait un $A$-module libre de rang $1$ (puisqu'il
  a même corps des fractions que $A$) engendré par l'élément $1$ de $A$
  (Bourbaki, \emph{Alg. comm.}, chap.~II, \S~3, n\textsuperscript{o}~2,
  prop.~5), ce qui est absurde.

  Pour construire les anneaux $A$ et $A'$, partons d'un corps imparfait $k$ de
  caractéristique $p>0$~; soient $K_0=k((T))$ le corps des séries formelles sur
  $k$, $A_0$ l'anneau de valuation pour $K_0$ correspondant à la valuation
  discrète égale à l'\emph{ordre} des séries formelles~; le corps résiduel de
  $A_0$ est donc $k$. Soit $\alpha$ un élément de $k$ qui n'est pas puissance
  $p$-ième, et prenons $k'=k(\alpha^{1/p})$~; $K=k'((T))$ est alors une extension
  radicielle finie de $K_0$, et $A'=A_0\otimes_k k'$ l'anneau de valuation
  discrète pour $K$ correspondant à l'ordre des séries formelles sur $K$. Si
  $\mathfrak m'$ est l'idéal maximal de $A'$, on répondra aux conditions 1) et 2)
  en prenant $A=A_0+\mathfrak m'$~: en effet, comme $A_0$ est noethérien et $A'$
  un $A_0$-module de type fini, $A$ est une $A_0$-algèbre finie, donc un anneau
  noethérien~; comme $A'$ est fini sur $A$, $\mathfrak m'$ est le seul idéal
  maximal de $A$, qui est donc un anneau local complet, évidemment de dimension
  $1$ (étant fini sur $A_0$) et géométriquement unibranche puisque $A'$ est
  intégralement clos et a même corps de fractions $K$ que $A$.

  \item La démonstration de (15.2.2) montre toutefois que l'on peut affaiblir la
  condition (iii) en y remplaçant «~géométriquement réduit~» par «~réduit~»
  lorsqu'on peut appliquer le critère valuatif de platitude (11.8.1) en utilisant
  seulement des anneaux de valuation discrète $A'$ dont le corps résiduel $k'$
  est séparable sur $\mathbf{k}(y)$~; en effet, si $x'$ et $z'$ ont les mêmes
  significations que dans (15.2.2), les multiplicités radicielles
  $\mu_i(\mathbf{k}(z')\mid k')$ et
  $\mu_i(\mathbf{k}(z)\mid\mathbf{k}(y))$ sont les mêmes, en vertu de (4.7.3),
  donc il résulte de (4.7.9) que $\operatorname{long}((\mathcal F_y)_z)$ et
  $\operatorname{long}((\mathcal F'_{y'})_{z'})$ sont égales, et l'on est encore
  ramené au cas où $Y$ est le spectre d'un anneau de valuation discrète et $y$
  son point fermé. Par exemple, on pourra obtenir ce résultat plus fort lorsque
  $\mathcal O_y$ est unibranche et dominé par un anneau de valuation discrète
  dont le corps résiduel est extension séparable de $\mathbf{k}(y)$, en vertu de
  (11.6.4). C'est le cas lorsque $\mathcal O_y$ est un anneau régulier, d'après
  (15.1.1.6). Mais, comme nous l'a signalé Hironaka, il y a des anneaux locaux
  intégralement clos noethériens, de dimension $2$ (provenant de schémas
  algébriques sur des corps imparfaits) qui ne satisfont pas à la condition
  précédente.
\end{enumerate}
\end{remarks}

\subsection{Applications : critères de réduction et d'irréductibilité}
\label{subsection:IV.15.3-fr}

\begin{proposition}[15.3.1]
\label{IV.15.3.1-fr}
Soient $Y$ un préschéma localement noethérien, $f:X\to Y$ un morphisme localement
de type fini, $\mathcal F$ un $\mathcal O_X$-Module cohérent, $x$ un point de
$X$, $y=f(x)$. Supposons vérifiées les conditions (i), (ii) et (iii) de
(15.2.2). Alors~:
\begin{enumerate}[label=(\roman*)]
  \item Il existe un voisinage $U$ de $x$ dans $X$ tel que $f|U$ soit
  universellement ouvert et que, pour tout $x'\in U$,
  $\mathcal F_{f(x')}$ soit géométriquement réduit sur
  $\mathbf{k}(f(x'))$ au point $x'$.
  \item Si de plus $Y$ est réduit au point $y$, il existe un voisinage
  $U_1\subset U$ de $x$ tel que $\mathcal F|U_1$ soit $f$-plat et réduit.
\end{enumerate}
\end{proposition}

\begin{proof}
Compte tenu de (\textbf{I}, 5.1.8), (4.2.7) et (4.7.11), les hypothèses ne
changent pas, ni la conclusion (i) de l'énoncé, si l'on remplace $Y$ par
$Y_{\mathrm{red}}$ et $X$ par $X\times_Y Y_{\mathrm{red}}$, et l'on peut donc
supposer $Y$ \emph{réduit}. Il résulte alors de (15.2.2) que $\mathcal F$ est
$f$-plat au point $x$, donc aussi (11.1.1) dans un voisinage de $x$, et
\emph{a fortiori} $f$ est universellement ouvert dans ce voisinage (2.4.6). Le
fait que $\mathcal F_{f(x')}$ soit alors géométriquement réduit au point $x'$
\oldpage[IV-3]{229}
pour tous les points d'un voisinage de $x$ résulte de (12.1.1, (vii)).
L'assertion (ii) résulte donc de ce qui précède et de (3.3.4).
\end{proof}

\begin{proposition}[15.3.2]
\label{IV.15.3.2-fr}
Les notations étant celles de (15.3.1), supposons que les conditions (i), (ii),
(iii) de (15.2.2) soient vérifiées, et de plus que $f^{-1}(y)$ soit
équidimensionnel au point $x$. Alors $f$ est équidimensionnel au point $x$.
\end{proposition}

\begin{proof}
On a, pour tout $y'\in Y$, $\operatorname{Supp}(\mathcal F_{y'})=f^{-1}(y')$.
Par hypothèse $\mathcal F_y$ est réduit (donc vérifie $(S_1)$) et est
équidimensionnel au point $x$. En vertu de (12.1.1, (ii)), on peut donc supposer
que $\mathcal F_{f(x')}$ est équidimensionnel et de dimension constante
$e=\dim_x(f^{-1}(y))$ au point $x'$ pour tous les $x'\in U$, ce qui signifie que
$f^{-1}(f(x'))$ est équidimensionnel et de dimension $e$ au point $x'$. Comme
par ailleurs $\mathcal F|U$ est $f$-plat, il résulte de (2.3.4) et de
(13.3.1, $a$)) que $f$ est équidimensionnel au point $x$.
\end{proof}

\begin{corollary}[15.3.3]
\label{IV.15.3.3-fr}
Les notations étant celles de (15.3.1), supposons vérifiées les conditions
suivantes~:
\begin{enumerate}[label=(\roman*)]
  \item $\operatorname{Supp}(\mathcal F)=X$.
  \item $f$ est universellement ouvert aux points génériques des composantes
  irréductibles de $f^{-1}(y)$ contenant $x$.
  \item $\mathcal F_y$ est géométriquement ponctuellement intègre sur
  $\mathbf{k}(y)$ au point $x$ (4.6.22).
  \item $Y$ est géométriquement unibranche au point $y$.
\end{enumerate}
Alors $x$ n'appartient qu'à une seule composante irréductible de $X$.

Si de plus $Y$ est intègre au point $y$ et $\mathcal F=\mathcal O_X$, alors $X$
est intègre au point $x$ et $f$ est plat au point $x$.
\end{corollary}

\begin{proof}
Notons d'abord que (i) et (iii) entraînent que $x$ n'appartient qu'à une seule
composante irréductible $Z$ de $f^{-1}(y)$. Il résulte donc de (14.4.7) et de
(15.3.2) que pour toute composante irréductible $X_i$ de $X$ contenant $x$
(donc $Z$), la restriction $f_i$ de $f$ à $X_i$ est un morphisme
équidimensionnel au point $x$, et universellement ouvert au point générique de
$Z$. La première assertion de l'énoncé résulte alors de (15.1.3) et la seconde
de (15.3.1).
\end{proof}

\begin{remarks}[15.3.4]
\label{IV.15.3.4-fr}
\begin{enumerate}[label=(\roman*)]
  \item L'exemple (14.4.10, (ii)) montre que, dans (15.3.3), on ne peut
  supprimer l'hypothèse que $Y$ est géométriquement unibranche au point $y$.
  \item La remarque (15.2.4, (vi)) montre que la conclusion de (15.3.3) est
  encore vraie sous les hypothèses suivantes~: $\mathcal O_y$ est régulier,
  $\operatorname{Supp}(\mathcal F)=X$ et $\mathcal F_y$ est intègre au point
  $x$. Nous ignorons si dans cet énoncé, on peut remplacer l'hypothèse que
  $\mathcal O_y$ est régulier par celle qu'il est géométriquement unibranche, ou
  même intègre et intégralement clos.
\end{enumerate}
\end{remarks}

\subsection{Compléments sur les morphismes de Cohen-Macaulay}
\label{subsection:IV.15.4-fr}

\begin{proposition}[15.4.1]
\label{IV.15.4.1-fr}
Soient $Y$ un préschéma localement noethérien, $f:X\to Y$ un morphisme localement
de type fini, $\mathcal F$ un $\mathcal O_X$-Module cohérent tel que
$\operatorname{Supp}(\mathcal F)=X$, $x$ un point de $X$, $y=f(x)$. On pose,
pour tout $z\in X$, $X_{f(z)}=f^{-1}(f(z))$,
$\mathcal F_{f(z)}=\mathcal F\otimes_{\mathcal O_{f(z)}}\mathbf{k}(f(z))$.
Supposons que $\mathcal F$ soit $f$-plat au point $x$ et que $\mathcal F_y$ soit
un $\mathcal O_{X_y}$-Module de Cohen-Macaulay au point $x$. Alors $f$ est
\oldpage[IV-3]{230}
équidimensionnel au point $x$.

En effet ((11.1.1) et (12.1.1, (vi))), il existe un voisinage $U$ de $x$ tel que
$\mathcal F|U$ soit $f$-plat et que pour tout $x'\in U$,
$\mathcal F_{f(x')}$ soit un $\mathcal O_{X_{f(x')}}$-Module de Cohen-Macaulay
au point $x'$~; cette dernière propriété entraîne que $\mathcal F_{f(x')}$ est
équidimensionnel au point $x'$ et que $x'$ n'appartient à aucun cycle premier
immergé associé à $\mathcal F_{f(x')}$ ($\mathbf{0}_{\mathrm{IV}}$, 16.5.4). En
outre, en vertu de (12.1.1, (ii)), on peut supposer que les dimensions des
composantes irréductibles de
$\operatorname{Supp}(\mathcal F_{f(x')})|(U\cap X_{f(x')})$ aient une (même)
valeur indépendante de $x'$. Comme
$\operatorname{Supp}(\mathcal F_{f(x')})=X_{f(x')}$ par hypothèse, il résulte de
(2.3.4) et de (13.3.1, $a$)) que $f$ est équidimensionnel au point $x$.
\end{proposition}

\begin{proposition}[15.4.2]
\EGAReaderTarget{ega4.c04.prop.15-4-2}{IV.15.4.2-fr}
Soient $Y$ un préschéma localement noethérien, $f:X\to Y$ un morphisme localement
de type fini, $\mathcal F$ un $\mathcal O_X$-Module cohérent tel que
$\operatorname{Supp}(\mathcal F)=X$, $x$ un point de $X$, $y=f(x)$. On suppose
que $\mathcal O_y$ est régulier et que $\mathcal F$ est un $\mathcal O_X$-Module
de Cohen-Macaulay au point $x$. Alors (avec les notations de (15.4.1)) les
conditions suivantes sont équivalentes~:
\begin{enumerate}[label=\emph{\alph*})]
  \item $\mathcal F$ est $f$-plat au point $x$ et $\mathcal F_y$ est un
  $\mathcal O_{X_y}$-Module de Cohen-Macaulay au point $x$.
  \item $\mathcal F$ est $f$-plat au point $x$.
  \item $f$ est universellement ouvert dans un voisinage de $x$.
  \item $f$ est ouvert aux points génériques des composantes irréductibles de
  $X_y$ contenant $x$.
  \item $\dim(\mathcal F_x)=\dim(\mathcal O_y)+\dim((\mathcal F_y)_x)$.
  \item[\emph{e}$'$)]
  $\dim(\mathcal O_{X,x})=\dim(\mathcal O_y)+
  \dim(\mathcal O_{X,x}\otimes_{\mathcal O_y}\mathbf{k}(y))$.
\end{enumerate}
Comme $\operatorname{Supp}(\mathcal F)=X$, les conditions $e)$ et $e')$ sont
équivalentes par définition (5.1.12). La condition $a)$ implique trivialement
$b)$~; $b)$ entraîne que $\mathcal F$ est $f$-plat aux points d'un voisinage de
$x$ (11.1.1), donc $b)$ entraîne $c)$ (2.4.6)~; $c)$ entraîne trivialement $d)$,
et $d)$ entraîne $e')$ par (14.2.1). Enfin, puisque $\mathcal O_y$ est régulier
et que $\mathcal F_x$ est un $\mathcal O_{X,x}$-module de Cohen-Macaulay, il
résulte de (6.1.4) que $e)$ entraîne $a)$.
\end{proposition}

\begin{proposition}[15.4.3]
\label{IV.15.4.3-fr}
Soient $f:X\to Y$ un morphisme localement de présentation finie, $\mathcal F$ un
$\mathcal O_X$-Module de présentation finie, tel que
$\operatorname{Supp}(\mathcal F)=X$, $x$ un point de $X$, $y=f(x)$. Alors (avec
les notations de (15.4.1)) les conditions suivantes sont équivalentes~:
\begin{enumerate}[label=\emph{\alph*})]
  \item $\mathcal F$ est $f$-plat au point $x$ et $\mathcal F_y$ est un
  $\mathcal O_{X_y}$-Module de Cohen-Macaulay au point $x$.
  \item Il existe un voisinage ouvert $U$ de $x$ dans $X$, et un $Y$-morphisme
  quasi-fini $g:U\to Y'=Y[T_1,\ldots,T_n]$ ($T_i$ indéterminées), tels que
  $\mathcal F|U$ soit $g$-plat au point $x$.
\end{enumerate}
\end{proposition}

\begin{proof}
Montrons d'abord que $b)$ entraîne $a)$. Posons $h:Y'\to Y$~; le
$\mathbf{k}(y)$-préschéma $Y_y'=h^{-1}(y)$ est régulier
($\mathbf{0}_{\mathrm{IV}}$, 17.3.7)~; soient $A$ l'anneau local de ce préschéma
au point $y'=g(x)$, $k$ son corps résiduel, et soit $B$ l'anneau local de $X_y$
au point $x$~; posons d'autre part $(\mathcal F_y)_x=M$, qui est un $B$-module
de type fini. L'hypothèse que le morphisme $g$ est quasi-fini entraîne qu'il en
est de même de $g_y:X_y\to Y_y'$ (\textbf{II}, 6.2.4), donc $M\otimes_A k$ est
un $(B\otimes_A k)$-module de longueur finie (\textbf{II}, 6.2.2)~; comme par
hypothèse $M$ est un $A$-module plat et $A$ un anneau de Cohen-Macaulay, $M$ est
un $B$-module de Cohen-Macaulay (6.3.3). Le fait que $\mathcal F$ est $f$-plat
au point $x$ résulte de ce qu'il est $g$-plat et de ce que le morphisme $h$ est
plat (2.1.6).

Montrons maintenant que $a)$ entraîne $b)$. La question étant locale sur $X$ et
$Y$, on peut supposer $X$ et $Y$ affines~; utilisant (8.9.1) et (11.2.7), on se
ramène au cas où $X$ et $Y$ sont noethériens, compte tenu de (6.7.1).
L'hypothèse entraîne que $f$ est
\oldpage[IV-3]{231}
équidimensionnel au point $x$ (15.4.1), d'où l'existence d'un morphisme
quasi-fini $g:U\to Y'$ tel que toute composante irréductible de $U$ domine $Y'$
(13.3.1, $b$)) et que $g_y:U_y\to Y_y'$ soit fini et surjectif (13.3.1.1).
D'autre part, pour prouver que $\mathcal F|U$ est $g$-plat au point $x$, il
suffit, puisque $h:Y'\to Y$ est plat, de montrer (en vertu du critère de
platitude par fibres (11.3.10)) que $\mathcal F_y$ est $g_y$-plat au point $x$.
Mais par hypothèse $\mathcal F_y$ est un $\mathcal O_{X_y}$-Module de
Cohen-Macaulay au point $x$ et $Y_y'$ un préschéma \emph{régulier}~; d'autre
part, comme $g_y$ est fini et surjectif, il vérifie la condition $e')$ de
(15.4.2) en vertu de (5.6.10)~; il suffit donc pour conclure d'appliquer (15.4.2)
à $g_y$ et à $\mathcal F_y$.
\end{proof}

% IV3_B32_P231_END_BEFORE_SUBSECTION_15.5

\subsection{Rang séparable des fibres d'un morphisme quasi-fini et universellement ouvert. Application aux composantes connexes géométriques des fibres d'un morphisme propre}
\label{subsection:IV.15.5-fr}

\begin{proposition}[15.5.1]
\label{IV.15.5.1-fr}
Soient $Y$ un préschéma localement noethérien, $f:X\to Y$ un morphisme séparé et quasi-fini. Pour tout $z\in Y$, soit $n(z)$ le nombre géométrique des points de $f^{-1}(z)$ (ou rang séparable de $f^{-1}(z)$ sur $\mathbf{k}(z)$, cf. \textbf{I}, 6.4.8). On a les propriétés suivantes~:
\begin{enumerate}[label=\emph{(\roman*)}]
  \item La fonction $z\mapsto n(z)$ est semi-continue inférieurement en tout point $y\in Y$ tel que $f$ soit universellement ouvert aux points de $f^{-1}(y)$.
  \item Soit $y$ un point de $Y$ tel que $f$ soit universellement ouvert aux points de $f^{-1}(y)$~; si la fonction $z\mapsto n(z)$ est constante dans un voisinage de $y$, alors $f$ est propre au point $y$ (15.7.1).
  \item On suppose que $f$ soit universellement ouvert et que tout point de $Y$ soit géométriquement unibranche~; alors, si la fonction $z\mapsto n(z)$ est localement constante dans $Y$, les composantes irréductibles de $X$ sont deux à deux disjointes.
\end{enumerate}
\end{proposition}

\begin{proof}
(i) Notons que puisque $f$ est quasi-fini, le nombre $n(z)$ est le nombre géométrique des composantes connexes de $f^{-1}(z)$~; on sait que la fonction $z\mapsto n(z)$ est localement constructible (9.7.9)~; compte tenu de ($\mathbf{0}_{\mathrm{III}}$, 9.3.4), on est ramené à prouver le

\begin{corollary}[15.5.2]
\EGAReaderTarget{ega4.c04.cor.15-5-2}{IV.15.5.2-fr}
Sous les hypothèses générales de (15.5.1), soit $y$ un point de $Y$ tel que $f$ soit universellement ouvert aux points de $f^{-1}(y)$. Alors, si $y'$ est une générisation de $y$, on a
\begin{equation}
\tag{15.5.2.1}
\label{IV.15.5.2.1-fr}
  n(y)\leqslant n(y').
\end{equation}
\end{corollary}

Pour tout changement de base $g:Z\to Y$, $f_{(Z)}:X_{(Z)}\to Z$ est séparé, quasi-fini et universellement ouvert aux points de $X_{(Z)}$ se projetant dans $f^{-1}(y)$~; en outre, si $z,z'$ sont deux points de $Z$ tels que $g(z)=y$, $g(z')=y'$ et que $z'$ soit une générisation de $z$, on a $n(z)=n(y)$ et $n(z')=n(y')$ (I, 6.4.12)~; il suffit donc de démontrer (15.5.2.1) pour un $g$ convenable. On peut évidemment supposer $y\ne y'$. Nous utiliserons le lemme suivant~:

\begin{lemma}[15.5.2.2]
\label{IV.15.5.2.2-fr}
Sous les hypothèses générales de (15.5.1), si $y$ est un point de $Y$, $y'$ une générisation de $y$ distincte de $y$, il existe un spectre d'anneau de valuation discrète $Z$, de point fermé $z$ et de point générique $z'$, et un morphisme $g:Z\to Y$ tels que~:
\begin{enumerate}[label=\emph{\arabic*\textsuperscript{\textdegree}}]
  \item $g(z)=y$, $g(z')=y'$~;
  \item si l'on pose $f'=f_{(Z)}$, $n(y)$ est le nombre de points de $f'^{-1}(z)$ et $n(y')$ est le nombre de points de $f'^{-1}(z')$.
\end{enumerate}
\end{lemma}

En effet, il existe un anneau de valuation discrète $A_1$ et un morphisme $g_1$ de $Z_1=\operatorname{Spec}(A_1)$ dans $Y$ tels que, si $z_1$ et $z'_1$ sont le point fermé et le point générique de $Z_1$,
\oldpage[IV-3]{232}
on ait $g_1(z_1)=y$ et $g_1(z'_1)=y'$ (\textbf{II}, 7.1.9)~; on peut donc déjà supposer que $Y$ est le spectre d'un anneau de valuation discrète $A$, $y$ son point fermé et $y'$ son point générique. Pour tout $x\in f^{-1}(y)$, $\mathbf{k}(x)$ est une extension finie de $\mathbf{k}(y)$ par hypothèse (\textbf{II}, 6.2.2)~; on déduit de (4.5.11) qu'il existe une extension algébrique finie $k$ de $\mathbf{k}(y)$ telle que le nombre de points de $f^{-1}(y)\otimes_{\mathbf{k}(y)}k$ soit égal à $n(y)$. Comme il y a un anneau de valuation discrète $A_2$ dominant $A$ et dont le corps résiduel est fini sur $\mathbf{k}(y)$ et contient $k$ (\textbf{II}, 7.1.2), on peut en second lieu supposer que $n(y)$ est le nombre de points de $f^{-1}(y)$. Enfin, le même raisonnement montre qu'il y a une extension finie $K$ de $\mathbf{k}(y')$ telle que le nombre de points de $f^{-1}(y')\otimes_{\mathbf{k}(y')}K$ soit égal à $n(y')$~; si $w$ est une valuation de $\mathbf{k}(y')$ associée à $A$, il y a une valuation discrète de $K$ prolongeant $w$, et l'anneau $A_3$ de cette valuation domine $A$, a $K$ pour corps des fractions, et répond aux conditions du lemme.

On peut donc désormais supposer que $Y$ est le spectre d'un anneau de valuation discrète, $y$ son point fermé, $y'$ son point générique, et que $n(y)$ et $n(y')$ sont les nombres de points de $f^{-1}(y)$ et $f^{-1}(y')$, de sorte que pour tout $x\in f^{-1}(y)$ (resp. tout $x'\in f^{-1}(y')$) on a $\mathbf{k}(x)=\mathbf{k}(y)$ (resp. $\mathbf{k}(x')=\mathbf{k}(y')$).

Désignons alors par $X_i$ les sous-préschémas réduits de $X$ ayant pour espaces sous-jacents les composantes irréductibles de $X$ ($1\leqslant i\leqslant m$). En vertu de (1.10.4), la restriction $X_i\to Y$ de $f$ à tout $X_i$ est un morphisme dominant, et comme $f^{-1}(y')$ est discret, chacune des intersections $X_i\cap f^{-1}(y')$ est réduite au point générique de $X_i$~; d'où (en désignant par $n_i(z)$ le nombre géométrique de points de $X_i\cap f^{-1}(z)$ pour tout $z\in Y$), $n_i(y')=1$ pour tout $i$, et $n(y')=\sum_{i=1}^{m}n_i(y')=m$. Par ailleurs, $f^{-1}(y)$ est réunion des $X_i\cap f^{-1}(y)$, d'où
\begin{equation}
\tag{15.5.2.3}
\label{IV.15.5.2.3-fr}
  n(y)\leqslant\sum_{i=1}^{m}n_i(y),
\end{equation}
et, pour prouver (15.5.2.1), il suffit d'établir que $n_i(y)\leqslant 1$ pour tout $i$. Autrement dit, on peut supposer $X$ intègre, et le morphisme $f$ birationnel. Mais alors, pour tout $x\in f^{-1}(y)$, on a $\mathcal O_y\subset\mathcal O_x\subset K$, où $K=\mathbf{k}(y')$ est le corps des fractions de $\mathcal O_y$. Comme $\mathcal O_y$ est un anneau de valuation et que $\mathcal O_x$ domine $\mathcal O_y$, on a nécessairement $\mathcal O_y=\mathcal O_x$ (\textbf{II}, 7.1.1). Il existe alors une $Y$-section $h:Y\to X$ telle que $h(y)=x$ (\textbf{I}, 2.4.4)~; comme $Y$ est réduit et que $X$ est séparé sur $Y$, une telle section est unique (\textbf{I}, 7.2.2), donc l'ensemble $f^{-1}(y)$ ne contient qu'un seul point, ce qui achève de prouver (i).

(ii) Comme au début de (i), on est ramené à montrer que si les deux membres de (15.5.2.1) sont égaux pour toute générisation $y'$ de $y$, $f$ est propre au point $y$. Grâce au critère de propreté locale (15.7.5) et aux remarques du début on peut encore supposer que $Y=\operatorname{Spec}(A)$ est le spectre d'un anneau de valuation discrète $A$, $y$ son point fermé et $y'$ son point générique, et il s'agit alors de montrer que $f$ est propre. Utilisons (15.5.2.2) et notons que si $A'$ est un anneau de valuation discrète dominant $A$, $A'$ est un $A$-module sans torsion, donc le morphisme $\operatorname{Spec}(A')\to\operatorname{Spec}(A)$ est fidèlement plat et quasi-compact, et il revient donc au même de dire que $f$ est propre ou que le morphisme déduit de $f$ par le changement de base $\operatorname{Spec}(A')\to\operatorname{Spec}(A)$ est
\oldpage[IV-3]{233}
propre (2.7.1, (vii))~; on peut donc supposer que $n(y)$ et $n(y')$ sont les nombres de points des fibres $f^{-1}(y)$ et $f'^{-1}(y')$ respectivement. Avec les notations de (i), on a alors par (15.5.2.3) et (15.5.2.1)
\begin{equation}
\tag{15.5.2.4}
\label{IV.15.5.2.4-fr}
  n(y)\leqslant\sum_i n_i(y)\leqslant\sum_i n_i(y')=n(y').
\end{equation}
et, pour que les membres extrêmes soient égaux, il faut que $n_i(y)=n_i(y')=1$ pour tout $i$. Comme on l'a vu dans (i), si $n_i(y)=1$ pour tout $i$, $X_i\cap f^{-1}(y)$ n'est pas vide et la restriction $X_i\to Y$ de $f$ est un isomorphisme~; comme les $X_i$ sont des sous-préschémas fermés de $X$, cela entraîne que $f$ est propre (\textbf{II}, 5.4.5).

(iii) On peut se borner au cas où $Y$ est affine et réduit, et comme $Y$ est par hypothèse localement intègre (\textbf{I}, 5.1.4), on peut supposer $Y$ intègre, et la fonction $y\mapsto n(y)$ constante dans $Y$. Soient encore $X_i$ ($1\leqslant i\leqslant m$) les composantes irréductibles de $X$. Il résulte de (1.10.4) que si $y'$ est le point générique de $Y$, chacun des $X_i\cap f^{-1}(y')$ est réduit à un seul point~; la restriction $f_i:X_i\to Y$ de $f$ étant un morphisme quasi-fini et dominant et tout point de $Y$ étant géométriquement unibranche, il résulte du critère de Chevalley (14.4.4) que $f_i$ est universellement ouvert. Si alors $y$ est un point quelconque de $Y$, on a $n_i(y)\leqslant n_i(y')$~; d'autre part, comme les $X_i\cap f^{-1}(y)$ sont deux à deux disjoints, on a $\sum_i n_i(y')=n(y')$~; on a donc encore les relations (15.5.2.4). Mais par hypothèse les membres extrêmes sont égaux, donc $n(y)=\sum_i n_i(y)$, ce qui implique que les $X_i\cap f^{-1}(y)$ sont deux à deux disjoints, et achève la démonstration.
\end{proof}

Combinant cette proposition avec le théorème de connexion de Zariski et ses conséquences (\textbf{III}, 4.3), on obtient les résultats suivants qui complètent (\textbf{III}, 4.3.7 et 4.3.10)~:

\begin{proposition}[15.5.3]
\label{IV.15.5.3-fr}
Soient $Y$ un préschéma irréductible localement noethérien, $f:X\to Y$ un morphisme propre~; pour tout $y\in Y$, soit $n(y)$ le nombre géométrique de composantes connexes (4.5.2) de $f^{-1}(y)$. On suppose que la restriction de $f$ à toute composante irréductible de $X$ soit un morphisme dominant dans $Y$. Alors, si $y_0$ est un point géométriquement unibranche de $Y$, la fonction $y\mapsto n(y)$ est semi-continue inférieurement au point $y_0$.
\end{proposition}

\begin{proof}
Considérons la factorisation de Stein $f=g\circ f'$ de $f$ (\textbf{III}, 4.3.3), où $g:Y'\to Y$ est un morphisme fini et $f':X\to Y$ est un morphisme surjectif dont les fibres sont géométriquement connexes (\textbf{III}, 4.3.4). Puisque $f'$ est surjectif, l'image réciproque par $f'$ de toute composante irréductible de $Y'$ contient une composante irréductible de $X$ au moins~; donc chacune des composantes irréductibles de $Y'$ domine $Y$. On peut donc appliquer le critère de Chevalley (14.4.4) aux restrictions de $g$ à chacune de ces composantes irréductibles, et on en conclut que $g$ est universellement ouvert en tout point de $g^{-1}(y_0)$. La conclusion résulte donc de (15.5.2) et du fait que le nombre géométrique des composantes connexes de $f^{-1}(y)$ est égal au nombre géométrique des points de $g^{-1}(y)$ (\textbf{III}, 4.3.4).
\end{proof}

\begin{corollary}[15.5.4]
\label{IV.15.5.4-fr}
Soient $Y$ un préschéma localement noethérien, $f:X\to Y$ un morphisme propre~; soit $y$ un point de $Y$ tel que $f$ soit universellement ouvert aux points de $f^{-1}(y)$~;
\oldpage[IV-3]{234}
alors, avec les notations de (15.5.3), la fonction $z\mapsto n(z)$ est semi-continue inférieurement au point $y$.
\end{corollary}

\begin{proof}
Avec les notations de (15.5.3), notons que dans la factorisation de Stein $f=g\circ f'$, $f'$ est surjectif~; comme $f$ est universellement ouvert aux points de $f^{-1}(y)$, $g$ est universellement ouvert aux points de $g^{-1}(y)$ (14.3.4, (i))~; il suffit alors d'appliquer à $g$ la proposition (15.5.1), (i) en tenant compte de (\textbf{III}, 4.3.4).
\end{proof}

\begin{remarks}[15.5.5]
\label{IV.15.5.5-fr}
\begin{enumerate}[label=\emph{(\roman*)}]
  \item Même si $Y$ est intègre et normal, $X$ intègre, $f$ fini et surjectif (donc universellement ouvert en vertu du critère de Chevalley (14.4.4)), la fonction $y\mapsto n(y)$ n'est pas nécessairement localement constante. On en a un exemple en prenant $Y=\operatorname{Spec}(k[T])$ où $k$ est algébriquement clos, et $X=\operatorname{Spec}(A)$, où $A=k[S,T]/(S^2+T^2-1)$~; aux points $y',y''$ de $Y$ correspondant aux idéaux maximaux $(T-1)$ et $(T+1)$, on a $n(y')=n(y'')=1$, mais $n(y)=2$ en tous les autres points de $Y$. On va donner ci-dessous une condition supplémentaire assurant que $y\mapsto n(y)$ est localement constante (15.5.7).
  \item L'exemple (14.4.10, (ii)) montre que dans (15.5.1, (iii)), on ne peut supprimer l'hypothèse que les points de $Y$ sont géométriquement unibranches.
  \item L'exemple (11.7.5) montre que la conclusion de (15.5.1, (i)) n'est plus valable si on suppose seulement que le morphisme $f$ est ouvert (même si, comme c'est le cas dans l'exemple cité, $f$ est fini et surjectif).
  \item Enfin, dans (15.5.1), on ne peut se dispenser de l'hypothèse que le morphisme $f$ est séparé, comme le montre l'exemple où $X$ est la droite affine dont on a «~dédoublé un point~» (\textbf{I}, 5.5.11), et $Y$ la droite affine~: ici $f$ est un isomorphisme local (donc est plat) et est quasi-fini, mais $z\mapsto n(z)$ n'est pas semi-continue inférieurement.
\end{enumerate}
\end{remarks}

\begin{lemma}[15.5.6]
\label{IV.15.5.6-fr}
Soit $Y$ le spectre d'un anneau de valuation discrète, $a$ son point fermé, $b$ son point générique. Soit $f:X\to Y$ un morphisme localement de type fini et ouvert. Supposons que $X$ soit connexe et que la fibre $f^{-1}(a)$ soit un préschéma réduit. Alors $f^{-1}(b)$ est connexe.
\end{lemma}

\begin{proof}
Nous utiliserons le lemme purement topologique suivant (cas particulier d'un résultat plus général du chap. III, 3\textsuperscript{e} Partie)~:
\begin{lemma}[15.5.6.1]
\label{IV.15.5.6.1-fr}
Soit $X$ un espace localement noethérien, dont toute partie fermée irréductible admet un point générique. Soit $Z$ une partie fermée rare de $X$ ayant la propriété suivante~: pour tout $x\in Z$, si l'on désigne par $T_x$ l'ensemble des générisations de $x$ dans $X$, alors $T_x-\{x\}$ est connexe. Dans ces conditions, si $X$ est connexe, $X-Z$ est connexe.
\end{lemma}

\begin{proof}
Donnons-en une démonstration indépendante. Raisonnant par l'absurde, supposons que l'ouvert $U=X-Z$, partout dense dans $X$, soit réunion de deux ouverts non vides sans point commun $U'$ et $U''$, et désignons par $X'$ et $X''$ les adhérences dans $X$ de $U'$ et $U''$. Comme $X=\overline{U}=X'\cup X''$ et que $X$ est connexe, on a $X'\cap X''\ne\varnothing$, et il est clair que $X'\cap X''\subset Z$. Soit $x$ un point générique d'une composante irréductible de $X'\cap X''$. Comme $X$ est localement noethérien, il y a un voisinage ouvert $V$ de $x$ dans $X$ tel que $V\cap U'$ et $V\cap U''$ n'aient qu'un nombre fini de composantes irréductibles~; comme $x$ est adhérent à $U'$ et à $U''$, il est nécessairement dans l'adhérence d'une composante irréductible $Z'$ de $V\cap U'$ et dans l'adhérence d'une composante irréductible $Z''$ de $V\cap U''$. Mais $Z'$ (resp. $Z''$) est fermé et irréductible dans $X$, donc admet un point générique $z'$ (resp. $z''$), et on a nécessairement $z'\in U'$ et $z''\in U''$. Ceci prouve que les intersections de $T_x-\{x\}$ avec $U'$ et $U''$ sont non vides. Mais par définition $(X'\cap X'')\cap(T_x-\{x\})=\varnothing$~; donc les intersections de $X'$ et de $X''$ avec $T_x-\{x\}$ sont deux fermés non vides disjoints dans $T_x-\{x\}$, dont la réunion est $T_x-\{x\}$~; or cela contredit l'hypothèse que $T_x-\{x\}$ est connexe. C.Q.F.D.
\end{proof}
\oldpage[IV-3]{235}
Pour prouver (15.5.6), nous appliquerons le lemme (15.5.6.1) à $X$ et à sa partie fermée $Z=f^{-1}(a)$ qui est rare puisque $f$ est ouvert. Notons que l'hypothèse, compte tenu de (15.2.2.1), implique que $f$ est plat aux points de $f^{-1}(a)$. Pour un tel point $x$, $\mathcal O_x$ est donc un $\mathcal O_a$-module sans torsion ($\mathbf{0}_{\mathrm I}$, 6.3.4), et en particulier, si $t$ est une uniformisante de $\mathcal O_a$, $t$ est $\mathcal O_x$-régulier. En outre $\mathcal O_x/t\mathcal O_x$ est réduit par hypothèse~; s'il est de profondeur $0$, c'est donc un corps ($\mathbf{0}$, 16.4.7), et comme $t\mathcal O_x$ est alors l'idéal maximal de $\mathcal O_x$, $\mathcal O_x$ est un anneau de valuation discrète ($\mathbf{0}$, 17.1.4), donc $\operatorname{Spec}(\mathcal O_x)-\{x\}$ est réduit à un point, et a fortiori est connexe. Si au contraire $\operatorname{prof}(\mathcal O_x/t\mathcal O_x)\geqslant 1$, on a $\operatorname{prof}(\mathcal O_x)\geqslant 2$ puisque $t$ est $\mathcal O_x$-régulier ($\mathbf{0}$, 16.4.6)~; il résulte alors du théorème de Hartshorne (5.10.7) que $\operatorname{Spec}(\mathcal O_x)-\{x\}$ est connexe, ce qui termine la démonstration.
\end{proof}

\begin{proposition}[15.5.7]
\label{IV.15.5.7-fr}
Soient $Y$ un préschéma localement noethérien, $f:X\to Y$ un morphisme propre. Pour tout $z\in Y$, soit $n(z)$ le nombre géométrique de composantes connexes de $f^{-1}(z)$. Soit $y$ un point de $Y$ tel que $f$ soit universellement ouvert aux points de $f^{-1}(y)$ et que $f^{-1}(y)$ soit géométriquement réduit sur $\mathbf{k}(y)$ (4.6.2). Alors la fonction $z\mapsto n(z)$ est constante dans un voisinage de $y$.
\end{proposition}

\begin{proof}
On sait déjà (15.5.4) que sous les conditions de l'énoncé, $z\mapsto n(z)$ est semi-continue inférieurement au point $y$~; tout revient à voir qu'elle est aussi semi-continue supérieurement, et par le même raisonnement qu'au début de la démonstration de (15.5.1, (i)), il suffit de montrer que si $y'$ est une générisation de $y$, on a
\begin{equation}
\tag{15.5.7.1}
\label{IV.15.5.7.1-fr}
  n(y')\leqslant n(y).
\end{equation}

Utilisant le raisonnement du début de la démonstration de (15.5.2) et le lemme (15.5.2.2) appliqué au morphisme fini $g$ de la factorisation de Stein $f=g\circ f'$ de $f$ (en se souvenant que le nombre géométrique de points de $g^{-1}(y)$ est égal à celui des composantes connexes de $f^{-1}(y)$ ($\mathbf{III}$, 4.3.4)), on se ramène au cas où $n(y)$ et $n(y')$ sont respectivement le nombre de composantes connexes de $f^{-1}(y)$ et $f^{-1}(y')$, et où $Y$ est le spectre d'un anneau de valuation discrète, $y$ le point fermé et $y'$ le point générique de $Y$. On peut en outre remplacer $X$ par une quelconque de ses composantes connexes, ces dernières étant ouvertes et fermées dans $X$, et propres sur $Y$. Supposons donc $X$ connexe~; l'hypothèse que $f$ est ouvert aux points de $f^{-1}(y)$ entraîne que si $f^{-1}(y)\ne\varnothing$, on a aussi $f^{-1}(y')\ne\varnothing$ ($\mathbf{0}_{\mathrm I}$, 1.10.3)~; l'hypothèse que $f$ est fermé entraîne que si $f^{-1}(y')\ne\varnothing$, on a aussi $f^{-1}(y)\ne\varnothing$ (\textbf{II}, 7.2.1). Donc, si $X\ne\varnothing$ (ce qu'on peut évidemment supposer, la proposition étant triviale dans le cas contraire), $f^{-1}(y)$ et $f^{-1}(y')$ sont tous deux non vides. En outre, l'hypothèse entraîne que le préschéma $f^{-1}(y)$ est réduit (4.6.1)~; comme $f$ est ouvert aux points de $f^{-1}(y)$, le lemme (15.5.6) implique que $f^{-1}(y')$ est connexe, autrement dit $n(y')=1$. Comme $f^{-1}(y)$ n'est pas vide, cela prouve (15.5.7.1).
\end{proof}

\begin{remark}[15.5.8]
\label{IV.15.5.8-fr}
La relation (15.5.7.1) n'est plus nécessairement valable si $f$ n'est pas supposé propre, même s'il vérifie les autres hypothèses de (15.5.7). Avec les notations de (15.5.5, (i)), il suffit pour le voir de considérer la restriction de $f$ à $X-\{x'\}$, où $x'$ est un des deux points de $X$ au-dessus d'un point $y$ distinct de $y'$ et $y''$~; on a alors $n(y)=1$, tandis que si $\eta$ est le point générique de $Y$, $n(\eta)=2$.
\end{remark}

\oldpage[IV-3]{236}
\begin{proposition}[15.5.9]
\label{IV.15.5.9-fr}
Soit $f:X\to Y$ un morphisme plat et localement de présentation finie~; pour tout $y\in Y$, soit $n(y)$ le nombre géométrique de composantes connexes de $f^{-1}(y)$.
\begin{enumerate}[label=\emph{(\roman*)}]
  \item Si $f$ est séparé et quasi-fini, la fonction $y\mapsto n(y)$ (égale alors au nombre géométrique des points de $f^{-1}(y)$) est semi-continue inférieurement dans $Y$~; si elle est constante dans un voisinage de $y_0$, $f$ est propre au point $y_0$.
  \item Si $f$ est propre, la fonction $y\mapsto n(y)$ est semi-continue inférieurement dans $Y$~; si de plus $f^{-1}(y_0)$ est géométriquement réduit sur $\mathbf{k}(y_0)$, $n(y)$ est constante dans un voisinage de $y_0$.
\end{enumerate}
\end{proposition}

\begin{proof}
Dans tous les cas, les questions sont locales sur $Y$, donc on peut supposer $Y$ affine et $f$ de présentation finie~; utilisant (8.9.1), (8.10.5) et (11.2.7), on se ramène au cas où $Y$ est noethérien (en utilisant en outre (8.2.11) pour la seconde assertion de (i)). Comme de plus $f$ est alors universellement ouvert (2.4.6), il suffit d'appliquer (15.5.1), (15.5.4) et (15.5.7) pour conclure.
\end{proof}

\begin{remark}[15.5.10]
\label{IV.15.5.10-fr}
On notera qu'on a ainsi obtenu une autre démonstration de (12.2.4, (vi)). Inversement, on peut déduire (15.5.7) de (12.2.4, (vi))~: en effet, on peut se borner au cas où $Y$ est réduit (en remplaçant $f$ par $f_{\mathrm{red}}$), et il résulte alors des hypothèses de (15.5.7) et de (15.2.3) que $f$ est plat dans un voisinage ouvert de $f^{-1}(y)$~; comme de plus $f$ est propre, on peut supposer que ce voisinage est de la forme $f^{-1}(U)$, où $U$ est un voisinage ouvert de $y$ dans $Y$. On conclut alors par (12.2.4, (vi)). La démonstration de (15.5.7) donnée plus haut a l'avantage de mettre en évidence le résultat (15.5.6), qui a un intérêt indépendant.
\end{remark}

% IV3_B33_P236_END_BEFORE_SUBSECTION_15.6

\subsection{Composantes connexes des fibres le long d'une section}
\label{subsection:IV.15.6-fr}

\begin{proposition}[15.6.1]
\label{IV.15.6.1-fr}
Soient $Y$ un préschéma localement noethérien, $f:X\to Y$ un morphisme localement de type fini, $g:Y\to X$ une $Y$-section de $X$ (\textbf{I}, 2.5.5). Pour tout $y\in Y$, désignons par $X_y^0$ la composante connexe de $f^{-1}(y)$ contenant le point $g(y)$. Soient $y$ un point de $Y$, $y'$ une générisation de $y$ dans $Y$. Alors :
\begin{enumerate}[label=\emph{(\roman*)}]
  \item S'il existe une composante irréductible $Z$ de $X_{y'}^0$, de dimension égale à $\dim(X_{y'}^0)$ et contenant $g(y')$ (ce qui sera le cas si $X_{y'}^0$ est irréductible), on a
\begin{equation}
\tag{15.6.1.1}
\label{IV.15.6.1.1-fr}
  \dim(X_y^0)\geqslant\dim(X_{y'}^0).
\end{equation}
  \item Si de plus $X_y^0$ est irréductible et si les deux membres de (15.6.1.1) sont égaux, alors $X_y^0\subset\overline{X_{y'}^0}$ (adhérence dans $X$).
\end{enumerate}
\end{proposition}

\begin{proof}
(i) Soit $\overline Z$ l'adhérence de $Z$ dans $X$~; comme $y$ est adhérent à $y'$ et $g$ continue, on a $g(y)\in\overline Z$. Comme $Z=\overline Z\cap f^{-1}(y')$, il résulte du théorème de semi-continuité de Chevalley (13.1.3) appliqué à la restriction de $f$ au sous-préschéma fermé réduit de $X$ ayant $\overline Z$ pour espace sous-jacent, que l'on a
\[
  \dim_{g(y)}(\overline Z\cap f^{-1}(y))\geqslant\dim(Z)~;
\]
mais les composantes irréductibles de $\overline Z\cap f^{-1}(y)$ contenant $g(y)$ sont évidemment contenues dans $X_y^0$, d'où \emph{a fortiori} l'inégalité (15.6.1.1).

(ii) Le raisonnement de (i) montre en outre que si les deux membres de (15.6.1.1) sont égaux, on a nécessairement
\[
  \dim_{g(y)}(\overline Z\cap f^{-1}(y))=\dim_{g(y)}(X_y^0)~;
\]
si $X_y^0$ est irréductible, cela entraîne $\overline Z\cap f^{-1}(y)=X_y^0$, donc $X_y^0$ est contenu dans $\overline Z$, et \emph{a fortiori} dans $\overline{X_{y'}^0}$.
\end{proof}

\oldpage[IV-3]{237}
\begin{corollary}[15.6.2]
\EGAReaderTarget{ega4.c04.cor.15-6-2}{IV.15.6.2-fr}
Avec les notations de (15.6.1), supposons en outre que pour tout $y\in Y$, $X_y^0$ soit irréductible. Alors la fonction $z\mapsto\dim(X_z^0)$ est semi-continue supérieurement dans $Y$. Si en outre cette fonction est constante dans un voisinage d'un point $y\in Y$, alors on a $X_y^0\subset\overline{X_{y'}^0}$ pour toute générisation $y'$ de $y$.
\end{corollary}

\begin{proof}
Soit $y$ un point de $Y$, et soit $U$ un voisinage ouvert affine de $g(y)$ dans $X$~; alors il y a un voisinage ouvert $V$ de $y$ dans $Y$ tel que $g(V)\subset U$, et pour tout $z\in V$, $U\cap X_z^0$ est dense dans $X_z^0$, donc $\dim(X_z^0)=\dim(U\cap X_z^0)$ (4.1.1.3). Remplaçant $X$ par $U$, on peut donc supposer que $f$ est un morphisme de type fini. On sait alors (9.7.10) que la fonction $z\mapsto\dim(X_z^0)$ est localement constructible, et les assertions du corollaire résultent donc de (15.6.1) et de ($\mathbf{0}_{\mathrm{III}}$, 9.3.4).
\end{proof}

\begin{proposition}[15.6.3]
\label{IV.15.6.3-fr}
Avec les notations de (15.6.1), supposons que pour tout $y\in Y$, $X_y^0$ soit géométriquement irréductible (4.5.2). Alors, pour tout $y\in Y$ tel que la fonction $z\mapsto\dim(X_z^0)$ soit constante dans un voisinage de $y$, $f$ est universellement ouvert aux points de $X_y^0$.
\end{proposition}

\begin{proof}
Appliquons le critère (14.3.7). Soit donc $h:Y'\to Y$ un morphisme, $Y'$ étant le spectre d'un anneau de valuation discrète, dont nous désignerons par $t$, $t'$ le point fermé et le point générique respectivement, et supposons que $h(t)=y$, de sorte que $y'=h(t')$ est une générisation de $y$. Posons $X'=X\times_Y Y'$, $f'=f_{(Y')}:X'\to Y'$, $g'=g_{(Y')}:Y'\to X'$. Avec les mêmes notations que dans (15.6.1), l'hypothèse faite sur les $X_y^0$ implique que
\[
  X_t^{\prime 0}=X_y^0\otimes_{\mathbf{k}(y)}\mathbf{k}(t)
  \quad\text{et}\quad
  X_{t'}^{\prime 0}=X_{y'}^0\otimes_{\mathbf{k}(y')}\mathbf{k}(t'),
\]
et que $X_t^{\prime 0}$ et $X_{t'}^{\prime 0}$ sont irréductibles (4.4.1)~; comme $\dim(X_y^0)=\dim(X_t^{\prime 0})$ et $\dim(X_{y'}^0)=\dim(X_{t'}^{\prime 0})$ (4.1.4), on voit qu'on est ramené à prouver la proposition pour $f'$ et $X_t^{\prime 0}$, autrement dit on peut se borner au cas où $Y$ est spectre d'un anneau de valuation discrète, $y$ son point fermé et $y'$ son point générique~; si $x'$ est le point générique de $X_{y'}^0$, il résulte alors de (15.6.2) que tout point de $X_y^0$ est spécialisation de $x'$, d'où la conclusion.
\end{proof}

\begin{proposition}[15.6.4]
\label{IV.15.6.4-fr}
Sous les hypothèses générales de (15.6.1), soit $X^0$ la réunion des $X_y^0$ lorsque $y$ parcourt $Y$. Si $y\in Y$ est tel que $X_y^0$ soit géométriquement réduit sur $\mathbf{k}(y)$ (4.6.2) et si $f$ est universellement ouvert en tout point de $X_y^0$, alors $X^0$ est un voisinage de $X_y^0$ dans $X$. En particulier, si $f$ est universellement ouvert et si, pour tout $y\in Y$, $X_y^0$ est géométriquement réduit sur $\mathbf{k}(y)$, $X^0$ est ouvert dans $X$.
\end{proposition}

\begin{proof}
Montrons d'abord qu'on peut se ramener au cas où $f$ est un morphisme \emph{de type fini}. Soit $x$ un point de $X_y^0$~; il résulte de (5.10.8.1) (avec $\mathfrak F=\{\varnothing\}$) qu'il y a une suite finie $(Z_i)_{0\leqslant i\leqslant n}$ de composantes irréductibles de $f^{-1}(y)$ telle que $g(y)\in Z_0$, $x\in Z_n$ et $Z_i\cap Z_{i+1}\ne\varnothing$ pour $0\leqslant i\leqslant n-1$~; les $Z_i$ étant irréductibles, il y a une suite finie $(U_i)$ ($1\leqslant i\leqslant n$) d'ouverts affines dans $f^{-1}(y)$ tels que $g(y)\in U_1$, $x\in U_n$, $U_i$ étant un voisinage affine d'un point $x_i$ de $Z_i\cap Z_{i-1}$ pour $1\leqslant i\leqslant n$. Il y a pour chaque $i$ un ouvert quasi-compact $W_i$ dans $X$ tel que $U_i=f^{-1}(y)\cap W_i$~; soit $W$ l'ouvert quasi-compact de $X$ réunion des $W_i$. Comme $g$ est continue, il y a un voisinage ouvert affine $V$ de $y$ dans $Y$ tel que $g(V)\subset W$~; si $X'=W\cap f^{-1}(V)$, la restriction de $f$ à $X'$ est de type fini~; en outre, si, pour tout $z\in V$, $X_z^{\prime 0}$ est la composante connexe de $X'\cap f^{-1}(z)$ contenant $g(z)$, on a $X_z^{\prime 0}\subset X_z^0$, et $x$ appartient à $X_y^{\prime 0}$. En effet, chacun des $U_i$ contient les deux ensembles irréductibles $U_i\cap Z_{i-1}$ et $U_i\cap Z_i$ qui se rencontrent, et les ensembles irréductibles $U_i\cap Z_{i-1}$ et $U_{i-1}\cap Z_{i-1}$

\oldpage[IV-3]{238}
se rencontrent pour $2\leqslant i\leqslant n$~; la réunion des $U_i\cap Z_{i-1}$ et des $U_i\cap Z_i$ pour $1\leqslant i\leqslant n$ est donc connexe et contient $x$ et $g(y)$. Comme $f|X'$ est universellement ouvert aux points de $X_y^{\prime 0}$, cela achève de prouver notre assertion, car si la proposition est prouvée lorsque $f$ est de type fini, la réunion $X^{\prime 0}$ des $X_z^{\prime 0}$ pour $z\in V$ sera un voisinage de $x$, et il en sera de même de $X^0$.

Supposons donc $f$ de type fini. Alors $X^0$ est localement constructible (9.7.12)~; il suffit par suite de montrer que $X^0$ contient toute générisation $x'$ de $x$ ($\mathbf{0}_{\mathrm{III}}$, 9.2.5). Posons $y'=f(x')$~; il existe un anneau de valuation discrète $A$ et un morphisme $u$ de $Y'=\operatorname{Spec}(A)$ dans $X$ tels que, si $z$ est le point fermé et $z'$ le point générique de $Y'$, on ait $u(z)=x$ et $u(z')=x'$ (\textbf{II}, 7.1.9)~; si $h=f\circ u$, on a donc $h(z)=y$ et $h(z')=y'$. Posons $X'=X\times_Y Y'$, $f'=f_{(Y')}:X'\to Y'$, qui est un morphisme de type fini universellement ouvert aux points de $f'^{-1}(z)$, et $g'=g_{(Y')}$, qui est une $Y'$-section de $X'$. Si $p_y:f'^{-1}(z)\to f^{-1}(y)$ (resp. $p_{y'}:f'^{-1}(z')\to f^{-1}(y')$) est la projection canonique, $p_y^{-1}(X_y^0)$ (resp. $p_{y'}^{-1}(X_{y'}^0)$) est la composante connexe de $f'^{-1}(z)$ (resp. $f'^{-1}(z')$) contenant $g'(z)$ (resp. $g'(z')$), car les $X_t^0$ sont géométriquement connexes (4.5.13) et il suffit d'appliquer (4.4.1). De plus, le morphisme $u$ correspond à une $Y'$-section $v$ de $X'$ telle que $p_y(v(y))=x$, $p_{y'}(v(y'))=x'$, de sorte que $v(y')$ est une générisation de $v(y)$, et il suffira donc de prouver que $v(y')$ appartient à $p_{y'}^{-1}(X_{y'}^0)$. Enfin, il est clair que $p_y^{-1}(X_y^0)$ est géométriquement réduit sur $\mathbf{k}(z)$.

En d'autres termes, on est ramené au cas où $Y$ est spectre d'un anneau de valuation discrète, $y$ son point fermé, $y'$ son point générique. Comme $f^{-1}(y)-X_y^0$ est alors fermé dans $X$, on peut, en remplaçant $X$ par l'ouvert $X-(f^{-1}(y)-X_y^0)$, supposer que $X_y^0=f^{-1}(y)$~; de même, on peut remplacer $X$ par l'ouvert, composante connexe de $X$, qui contient $g(Y)$, cette composante connexe contenant évidemment $X^0$~; autrement dit, on peut supposer $X$ connexe. La proposition sera alors établie si l'on montre que $X^0=X$, ou encore que $X_{y'}$ est connexe~; mais $f^{-1}(y')$ est géométriquement réduit sur $\mathbf{k}(y')$, et \emph{a fortiori} réduit~; comme en outre $f$ est ouvert aux points de $f^{-1}(y)$, on peut appliquer le lemme (15.5.6), d'où la conclusion.
\end{proof}

\begin{corollary}[15.6.5]
\label{IV.15.6.5-fr}
Soient $f:X\to Y$ un morphisme plat de présentation finie, $g$ une $Y$-section de $X$~; pour tout $y\in Y$, soit $X_y^0$ la composante connexe de $f^{-1}(y)$ contenant $g(y)$, et soit $X^0$ la réunion des $X_y^0$ quand $y$ parcourt $Y$. Alors, pour tout $y\in Y$ tel que $X_y^0$ soit géométriquement réduit sur $\mathbf{k}(y)$, $X^0$ est un voisinage de $X_y^0$ dans $X$. En particulier, si $f$ est réduit (6.8.1), $X^0$ est ouvert dans $X$.
\end{corollary}

\begin{proof}
La question étant locale sur $Y$, on peut supposer $Y=\operatorname{Spec}(A)$ affine. Il existe donc un sous-anneau noethérien $A_0$ et un morphisme plat de type fini $f_0:X_0\to Y_0=\operatorname{Spec}(A_0)$ tels que $X=X_0\times_{Y_0}Y$ et $f=(f_0)_{(Y)}$ (8.9.1 et 11.2.7)~; en outre, on peut supposer qu'il existe une $Y_0$-section $g_0$ de $X_0$ telle que $g=(g_0)_{(Y)}$ (8.8.2). Pour tout $y\in Y$, soit $y_0$ sa projection dans $Y_0$~; il résulte de (4.5.13) que $(X_0)_{y_0}^0$ est géométriquement connexe, donc, si $p:f^{-1}(y)\to f_0^{-1}(y_0)$ est la projection canonique, on a $X_y^0=p^{-1}((X_0)_{y_0}^0)$ (4.4.1)~; de plus, l'hypothèse que $X_y^0$ est géométriquement réduit sur $\mathbf{k}(y)$ entraîne que $(X_0)_{y_0}^0$ est géométriquement réduit sur $\mathbf{k}(y_0)$ (4.6.10). On est ainsi ramené au cas où on suppose

\oldpage[IV-3]{239}
en outre $Y$ noethérien~; comme $f$ est universellement ouvert (11.1.1), il suffit alors d'appliquer (15.6.4).
\end{proof}

\begin{proposition}[15.6.6]
\EGAReaderTarget{ega4.c04.prop.15-6-6}{IV.15.6.6-fr}
Soit $f:X\to Y$ un morphisme tel que $Y$ soit localement noethérien et $f$ localement de type fini, ou que $f$ soit de présentation finie. Soit $g$ une $Y$-section de $X$~; pour tout $y\in Y$, soit $X_y^0$ la composante connexe de $f^{-1}(y)$ contenant $g(y)$~; enfin, soit $X^0$ la réunion des $X_y^0$ lorsque $y$ parcourt $Y$.

Supposons que, pour tout $y\in Y$, $X_y^0$ soit géométriquement irréductible. Alors :
\begin{enumerate}[label=\emph{(\roman*)}]
  \item La fonction $y\mapsto\dim(X_y^0)$ est semi-continue supérieurement dans $Y$.
  \item Soient $x_0\in X^0$, $y_0=f(x_0)$. Si la fonction $y\mapsto\dim(X_y^0)$ est constante dans un voisinage de $y_0$, il existe un voisinage ouvert $U$ de $y_0$ tel que $f$ soit universellement ouvert aux points de $X^0\cap f^{-1}(U)$. Si de plus $Y$ est réduit et $f^{-1}(y_0)$ géométriquement réduit sur $\mathbf{k}(y_0)$ au point $x_0$, $f$ est plat au point $x_0$.
  \item Inversement, supposons que $f$ soit universellement ouvert au point générique de $X_{y_0}^0$ et en outre que l'une des conditions suivantes soit vérifiée :
  \begin{enumerate}
    \item[$\alpha$)] La fibre $f^{-1}(y_0)$ est géométriquement réduite sur $\mathbf{k}(y_0)$ au point $g(y_0)$ et l'anneau $\mathcal O_{Y,y_0}$ est noethérien.
    \item[$\beta$)] Pour tout point générique $y'$ d'une composante irréductible de $Y$ contenant $y_0$, on a $\dim(f^{-1}(y'))=\dim(X_{y'}^0)$.
  \end{enumerate}
  Alors la fonction $y\mapsto\dim(X_y^0)$ est constante dans un voisinage de $y_0$.
  \item L'ensemble $W$ des points $x\in X^0$ tels que la fonction $y\mapsto\dim(X_y^0)$ soit constante dans un voisinage de $f(x)$ et que $X_{f(x)}^0$ soit géométriquement réduit sur $\mathbf{k}(f(x))$, est ouvert dans $X$.
\end{enumerate}
\end{proposition}

\begin{proof}
Les questions étant locales sur $Y$, on peut supposer que $Y=\operatorname{Spec}(A)$ est affine. Lorsque $f$ est de présentation finie, il existe donc un sous-anneau noethérien $A_1$ de $A$ et un morphisme de type fini $f_1:X_1\to Y_1=\operatorname{Spec}(A_1)$ tels que $X=X_1\times_{Y_1}Y$ et $f=(f_1)_{(Y)}$ (8.9.1)~; en outre, on peut supposer qu'il existe une $Y_1$-section $g_1$ de $X_1$ telle que $g=(g_1)_{(Y)}$ (8.8.2). Pour tout $y\in Y$, soit $y_1$ sa projection dans $Y_1$~; il résulte de (4.5.13) que $(X_1)_{y_1}^0$ est géométriquement connexe, donc, si $p:f^{-1}(y)\to f_1^{-1}(y_1)$ est la projection canonique, on a $X_y^0=p^{-1}((X_1)_{y_1}^0)$ (4.4.1). Notons d'autre part que l'on peut, en vertu de (9.7.11) et (9.7.12), appliquer (9.3.3) à la propriété d'être géométriquement irréductible~; on peut donc supposer $A_1$ choisi de telle sorte que $(X_1)_{y_1}^0$ soit géométriquement irréductible pour tout $y_1\in Y_1$. On a $\dim(X_y^0)=\dim((X_1)_{y_1}^0)$ (4.1.4), et en appliquant de nouveau (9.3.3) à la propriété d'avoir une dimension donnée (et utilisant (9.5.5) et (9.7.12)), on peut supposer (en restreignant au besoin $Y$ à un voisinage de $y_0$) que si $\dim(X_y^0)$ est constante dans $Y$, alors $\dim((X_1)_{y_1}^0)$ est constante dans $Y_1$. Si $Y$ est réduit, il en est de même de $Y_1$ puisque $A_1\subset A$~; d'autre part, si $f^{-1}(y_0)$ est géométriquement réduit au point $x_0$, $f_1^{-1}(y_{01})$ l'est au point $x_{01}$ ($x_{01}$ et $y_{01}$ étant les projections respectives dans $X_1$ et $Y_1$ de $x_0$ et $y_0$) (4.6.10).

Ces remarques montrent que pour démontrer (i) et (ii), on peut se borner au cas où $Y$ est noethérien et $f$ localement de type fini. Mais dans ce cas, l'assertion (i) résulte de (15.6.2) et la première assertion de (ii) de (15.6.3). D'autre part, si $Y$ est réduit et $f^{-1}(y_0)$ géométriquement réduit sur $\mathbf{k}(y_0)$ au point $x_0$ ($Y$ étant toujours supposé

\oldpage[IV-3]{240}
noethérien), comme $X_y^0$ est la seule composante irréductible de $f^{-1}(y)$ contenant $x$, on déduit de (15.6.3) et (15.2.3) que si $\dim(X_y^0)$ est constante dans un voisinage de $y_0$, $f$ est plat au point $x_0$.

Pour prouver (iii), remarquons d'abord que l'ensemble $X^0$ est localement constructible dans $X$ (9.7.12), donc, par (9.5.5), l'ensemble $E$ des $y\in Y$ tels que $\dim(X_y^0)=\dim(X_{y_0}^0)$ est localement constructible dans $Y$. Appliquons alors (1.10.1) à $E$~: il suffit de prouver (compte tenu de (i)) que $\dim(X_{y'}^0)=\dim(X_{y_0}^0)$ pour le point générique $y'$ d'une composante irréductible de $Y$ contenant $y_0$~; ceci permet déjà de supposer que $Y=\operatorname{Spec}(\mathcal O_{Y,y_0})$.

Si on est dans le cas $\alpha$), on se ramène aussitôt, en vertu de (\textbf{II}, 7.1.9), et utilisant (4.5.13) et (4.4.1) comme dans (15.6.4), au cas où $Y$ est spectre d'un anneau de valuation discrète, $y_0$ son point fermé et $y'$ son point générique. Mais alors les hypothèses entraînent en vertu de (15.2.2.1), que $f$ est \emph{plat} au point $g(y_0)$, donc aussi dans un voisinage $V$ de ce point dans $X$ (11.1.1). Pour démontrer notre assertion, on peut remplacer $X$ par $V$, car $V\cap X_{y'}^0$ n'est pas vide en vertu de (2.3.4), donc est un ouvert partout dense dans $X_{y'}^0$ et a par suite même dimension (4.1.1.3)~; d'ailleurs, comme $X_{y'}^0$ est aussi la composante connexe de $f^{-1}(y')$ contenant $g(y')$, on peut supposer $V$ pris tel que $V$ ne rencontre aucune autre composante irréductible de $f^{-1}(y_0)$ ni de $f^{-1}(y')$, autrement dit on peut supposer que $X^0=X$~; mais alors la conclusion résulte de (12.1.1, (i)), puisque par hypothèse $X_{y_0}^0$ est intègre.

Supposons maintenant qu'on soit dans le cas $\beta$). Appliquons cette fois (\textbf{II}, 7.1.4) de la même manière que (\textbf{II}, 7.1.9) dans le cas $\alpha$)~: on est alors ramené au cas où $Y$ est spectre d'un anneau de valuation (non nécessairement discrète), $y_0$ son point fermé et $y'$ son point générique. En vertu de (14.3.13), il existe une composante irréductible $Z$ de $X$ contenant $X_{y_0}^0$ et dominant $Y$, et telle que $\dim(Z\cap f^{-1}(y'))=\dim(X_{y_0}^0)$~; mais l'hypothèse $\beta$) et le fait que $\dim(Z\cap f^{-1}(y'))\leqslant\dim(f^{-1}(y'))$ montrent que $\dim(X_{y_0}^0)\leqslant\dim(X_{y'}^0)$, d'où $\dim(X_{y_0}^0)=\dim(X_{y'}^0)$ en vertu de (i).

Reste à prouver (iv). Notons d'abord que les ensembles envisagés ne changent pas lorsqu'on remplace $f$ par $f_{(Y_{\mathrm{red}})}:X\times_Y Y_{\mathrm{red}}\to Y_{\mathrm{red}}$, la projection $X\times_Y Y_{\mathrm{red}}\to X$ étant un homéomorphisme. Autrement dit, on peut supposer $Y$ réduit, et alors il résulte de (ii) que $f$ est \emph{plat} aux points de $W$. Considérons un point $x_0$ de $W$ et prouvons que $W$ est un voisinage de $x_0$~; procédant comme au début de la démonstration, et utilisant aussi (11.2.7), on peut de plus supposer $Y$ noethérien~; il résulte alors de (ii) et de (15.6.4) que $X^0$ est un voisinage de $x_0$ dans $X$, et de (12.1.1, (vii)) que $W$ est aussi un voisinage de $x_0$ dans $X$.
\end{proof}

\begin{corollary}[15.6.7]
\label{IV.15.6.7-fr}
Supposons vérifiées les conditions préliminaires de (15.6.6) sur $f$, et supposons en outre que pour tout $y\in Y$, $X_y^0$ soit géométriquement intègre sur $\mathbf{k}(y)$. Alors les conditions suivantes sont équivalentes :
\begin{enumerate}[label=\emph{\alph*)}]
  \item La fonction $y\mapsto\dim(X_y^0)$ est localement constante dans $Y$.
  \item Le morphisme $f_{(Y_{\mathrm{red}})}:X\times_Y Y_{\mathrm{red}}\to Y_{\mathrm{red}}$ déduit de $f$ par changement de base, est plat aux points de $X^0$.
\end{enumerate}

\oldpage[IV-3]{241}
En outre, ces conditions entraînent que $X^0$ est ouvert dans $X$. Enfin, lorsque $Y$ est localement noethérien, a) et b) sont aussi équivalentes à
\begin{enumerate}
  \item[\rm b$'$)] $f$ est universellement ouvert aux points de $X^0$.
\end{enumerate}
\end{corollary}

\begin{proof}
Le fait que a) entraîne b) résulte de (15.6.6, (ii)), ainsi que le fait que $X^0$ est alors ouvert dans $X$. Si b) est vérifiée, on peut se borner au cas où $Y$ est réduit et $f$ plat~; alors, on se ramène, comme au début de (15.6.6), et en utilisant en outre (11.2.7), au cas où $Y$ est noethérien, cas où la conclusion résulte de (15.6.6, (iii), cas $\alpha$)). L'équivalence de b) et b$'$) a déjà été prouvée lorsque $Y$ est localement noethérien (15.2.2.1), compte tenu de ce que le morphisme $X\times_Y Y_{\mathrm{red}}\to X$ est un homéomorphisme universel.
\end{proof}

\begin{proposition}[15.6.8]
\label{IV.15.6.8-fr}
Soient $f:X\to Y$ un morphisme séparé de présentation finie, $g$ une $Y$-section de $X$~; pour tout $y\in Y$, soit $X_y^0$ la composante connexe de $f^{-1}(y)$ qui contient $g(y)$, et soit $X^0$ la réunion des $X_y^0$ quand $y$ parcourt $Y$. Alors, si $y\in Y$ est tel que $X_y^0$ soit propre sur $\mathbf{k}(y)$, $X^0$ est propre sur $Y$ au point $y$ (i.e. (15.7.1), il existe un voisinage ouvert $U$ de $y$ dans $Y$ tel que $X^0\cap f^{-1}(U)$ soit fermé dans $f^{-1}(U)$ et propre sur $U$).
\end{proposition}

\begin{proof}
Procédant comme au début de (15.6.6) (où on utilise (8.10.5, (v)), et où on remplace (9.7.11) par (9.6.7)), on se ramène au cas où $Y$ est noethérien et $f$ de type fini. On sait alors (9.7.12) que $X^0$ est localement constructible dans $X$~; d'autre part, les $X_z^0$ sont géométriquement connexes (pour tout $z\in Y$) en vertu de (4.5.13)~; enfin, comme $g(Y)\subset X^0$, $X^0$ est universellement submersif sur $Y$ (15.7.8). Il suffit alors d'appliquer à $X^0$ le critère (15.7.9).
\end{proof}

\begin{corollary}[15.6.9]
\label{IV.15.6.9-fr}
Soit $f:X\to Y$ un morphisme séparé tel que $Y$ soit localement noethérien et $f$ localement de type fini, ou que $f$ soit de présentation finie. Soit $g$ une $Y$-section de $X$, et pour tout $y\in Y$, soit $X_y^0$ la composante connexe de $f^{-1}(y)$ contenant $g(y)$. Soit $y$ un point de $Y$ tel que $X_y^0$ soit propre sur $\mathbf{k}(y)$. Alors, pour toute générisation $y'$ de $y$, on a $\dim(X_y^0)\geqslant\dim(X_{y'}^0)$.
\end{corollary}

\begin{proof}
Supposons d'abord $f$ de présentation finie~; alors en remplaçant au besoin $Y$ par un voisinage de $y$, on peut supposer que $f|X^0$ est \emph{propre} (15.6.8), et il suffit d'appliquer (13.1.5) à cette restriction. Si $Y$ est localement noethérien et $f$ localement de type fini, $X$ est localement noethérien~; en outre l'hypothèse que $X_y^0$ soit propre sur $\mathbf{k}(y)$ entraîne que $X_y^0$ est noethérien~; il existe donc un ouvert noethérien $V\subset X$ contenant $X_y^0$~; par suite il y a un voisinage ouvert $W$ de $y$ dans $Y$ tel que $g(W)\subset V$. En outre, si $z$ est point maximal de $X_{y'}^0$, on peut supposer que $z\in V$. La restriction de $f$ à $V\cap f^{-1}(W)$ est alors un morphisme de type fini, et $g|W$ une $W$-section~; en appliquant à ce morphisme le résultat déjà obtenu, on voit que si $Z$ est la composante irréductible de $X_{y'}^0$ de point générique $z$, on a $\dim(Z)\leqslant\dim(X_y^0)$. Comme $z$ est un point maximal quelconque de $X_{y'}^0$, on a bien $\dim(X_{y'}^0)\leqslant\dim(X_y^0)$.
\end{proof}

\begin{remarks}[15.6.10]
\EGAReaderTarget{ega4.c04.item.15-6-10}{IV.15.6.10-fr}
(i) L'hypothèse supplémentaire sur l'existence d'une composante irréductible de $X_{y'}^0$ contenant $g(y')$ et de dimension égale à $\dim(X_{y'}^0)$, faite dans (15.6.1), n'est pas superflue, comme le montre l'exemple suivant. Soient $A$ un anneau de valuation discrète, $\pi$ une uniformisante de $A$, $K$ le corps des fractions de $A$, $k$ son corps résiduel. Prenons $Y=\operatorname{Spec}(A)$, et désignons par $y$ et $y'$ le point fermé et le point générique de $Y$. Considérons les deux $Y$-schémas suivants~: $X_1=\operatorname{Spec}(A[t])$, $X_2=\operatorname{Spec}(K[u,v])$, où $t,u,v$ sont des indéterminées (on notera que la projection de $X_2$ dans $Y$ est donc $\{y'\}$). Dans l'anneau $A[t]$, l'idéal principal $(\pi t-1)$ est \emph{maximal}, et correspond donc à un point fermé $x_1$ de $X_1$, qui se projette au point générique $y'$ de $Y$ et est tel que $\mathbf{k}(x_1)=K$.

\oldpage[IV-3]{242}
Considérons d'autre part le point fermé $x_2$ de $X_2$ correspondant à l'idéal maximal $\left(u-\frac{1}{\pi}\right)+(v)$ de $K[u,v]$~; on a aussi $\mathbf{k}(x_2)=K$. Comme on le verra au chap.~V, on peut donc \emph{recoller} $X_1$ et $X_2$ le long des deux sous-préschémas fermés $Z_1$, $Z_2$ de $X_1$, $X_2$ ayant respectivement $\{x_1\}$ et $\{x_2\}$ pour espaces sous-jacents, suivant l'unique $Y$-isomorphisme de $Z_1$ sur $Z_2$. Désignons par $X$ le $Y$-préschéma ainsi obtenu, par $x$ le point de $X$ image de $x_1$ et $x_2$. La «~section nulle~» $g$ de $X_1$ (\textbf{II}, 1.7.9) est encore alors une $Y$-section de $X$~; $f^{-1}(y)$ est égale à $(X_1)_y$, tandis que $f^{-1}(y')$ a deux composantes irréductibles, respectivement isomorphes à $(X_1)_{y'}^0$ et $(X_2)_{y'}^0$, ayant un unique point commun $x$, et de dimensions respectives 1 et 2. Voir cependant la prop.~(15.6.9).

(ii) Considérons le morphisme $f$ défini dans (12.2.3, (ii)), qui est propre et plat~; en outre, la restriction de $f$ à $X_1\cap f_0^{-1}(Y)$ est un isomorphisme, donc le morphisme réciproque $g:Y\to X$ est une $Y$-section de $X$. On a alors $\dim(X_{y_0})=1$ tandis que $\dim(X_y^0)=0$ pour $y\ne y_0$, bien que $X_y^0$ soit géométriquement irréductible pour tout $y\in Y$ (mais $X_{y_0}^0$ n'est pas réduit)~; on voit donc que dans (15.6.6, (iii)), on ne peut supprimer les hypothèses $\alpha$) et $\beta$). En outre, $X^0$ n'est pas un voisinage de $X_{y_0}^{0}$, ce qui prouve que dans (15.6.4), on ne peut se dispenser de l'hypothèse que $X_{y_0}^0$ est réduit.

(iii) Au chap.~VI, nous appliquerons les résultats précédents aux $Y$-préschémas \emph{en groupes} localement de type fini sur un préschéma localement noethérien $Y$. Si $G$ est un tel préschéma, il existe une $Y$-section canonique $g$, la «~section unité~», et l'on montre que $G_y^0$ est toujours géométriquement irréductible et que l'on a $\dim(G_y^0)=\dim(G_y)$, en posant $G_y=f^{-1}(y)$. Il résultera donc de (15.6.2) et (15.6.3) que la fonction $z\mapsto\dim(G_z)$ est semi-continue supérieurement, et que si elle est constante au voisinage d'un point $y$, $f$ est universellement ouvert aux points de $G_y^0$. Si de plus $G_y$ est géométriquement réduit sur $\mathbf{k}(y)$ en un de ses points, on montre que $G_y$ est \emph{lisse} (6.8.1) sur $\mathbf{k}(y)$ en tous ses points~; il résultera alors de (15.6.6) que si de plus $\mathcal O_y$ est réduit, alors $f$ est lisse en tous les points de $G_y^0$ (mais pas nécessairement en tous les points de $G_y$). Ces résultats s'appliqueront par exemple dans la théorie de schémas de Picard, où nous disposerons d'un théorème général assurant que, sous certaines conditions, la fonction $z\mapsto\dim(G_z)$ est localement constante. Remarquons que c'est en vue d'applications de cette nature que les énoncés tels que (15.6.1) sont donnés pour les morphismes \emph{localement} de type fini, et non seulement pour les morphismes de type fini.

On notera aussi que dans le cas d'un $Y$-préschéma en groupes $G$, l'hypothèse $\beta$) de (15.6.6) est toujours vérifiée.
\end{remarks}

%
% DEBUT SOURCE ega4/ega4-15-terminal-fr.tex
% EGA IV-3 terminal continuation, printed pp. 242--255.
% Direct diplomatic French transcription from the NUMDAM authority.

\subsection{Appendice : Critères valuatifs de propreté locale}
\label{subsection:IV.15.7-fr}

Ce numéro donne des compléments au critère valuatif de propreté démontré dans (\textbf{II}, 7.3.10)~; il est indépendant du reste du \S~15.

\begin{definition}[15.7.1]
\label{IV.15.7.1-fr}
Soient $f:X\to Y$ un morphisme de préschémas, $y$ un point de $Y$. On dit que $f$ est propre au point $y$ s'il existe un voisinage ouvert $U$ de $y$ dans $Y$ tel que la restriction $f^{-1}(U)\to U$ de $f$ soit un morphisme propre. On dit qu'une partie $Z$ de $X$ est propre sur $Y$ au point $y$ s'il existe un voisinage ouvert $U$ de $y$ tel que $Z\cap f^{-1}(U)$ soit fermée dans $f^{-1}(U)$ et propre sur $U$ (\textbf{II}, 5.4.10) (ce qui revient à dire que pour tout sous-préschéma fermé de $X$ ayant $Z$ pour espace sous-jacent, la restriction $Z\to Y$ de $f$ à $Z$ est propre au point $y$).
\end{definition}

Soit $g:Y'\to Y$ un morphisme quelconque~; posons $X'=X\times_Y Y'$, $f'=f_{(Y')}$ et, si $p:X'\to X$ est la projection canonique, $Z'=p^{-1}(Z)$. Alors, si $Z$ est propre sur $Y$ au point $y$, $Z'$ est propre sur $Y'$ en tout point $y'$ au-dessus de $y$ (\textbf{II}, 5.4.2).

\begin{proposition}[15.7.2]
\label{IV.15.7.2-fr}
Soient $Y$ un préschéma intègre localement noethérien, $X$ un préschéma intègre, $f:X\to Y$ un morphisme séparé, dominant et de type fini, $y$ un point de $Y$. Les conditions suivantes sont équivalentes :
\begin{enumerate}[label=\emph{\alph*)}]
  \item $f$ est propre au point $y$.
  \item Si $Y_1=\operatorname{Spec}(\mathcal O_y)$, $X_1=X\times_Y Y_1$, le morphisme $f_1=f_{(Y_1)}:X_1\to Y_1$ est propre.
  \item Tout anneau de valuation discrète ayant $R(X)$ pour corps des fractions et dominant $\mathcal O$ domine un anneau local $\mathcal O_x$ de $X$ (auquel cas on a nécessairement $f(x)=y$).
\end{enumerate}
\end{proposition}

\oldpage[IV-3]{243}
\begin{proof}
Le fait que a) implique b) résulte de (\textbf{II}, 5.4.2), et l'implication b) $\Rightarrow$ c) résulte de (\textbf{II}, 7.3.10). Reste donc à montrer que c) entraîne a). La question étant locale sur $Y$, on peut supposer $Y$ affine, donc noethérien. En vertu du lemme de Chow (\textbf{II}, 5.6.1), il existe un préschéma intègre $X'$, un morphisme projectif $p:P\to Y$, une immersion ouverte dominante $j:X'\to P$, et un morphisme projectif et birationnel (donc surjectif) $g:X'\to X$ tels que le diagramme
\[
\xymatrix{
P\ar[d]_-{p} & X'\ar[l]_-{j}\ar[d]^-{g}\\
Y & X\ar[l]^-{f}
}
\]
soit commutatif. Comme $X'$ est intègre et $j$ dominant, $P$ est irréductible, et l'on peut, en remplaçant $P$ par $P_{\mathrm{red}}$, supposer $P$ intègre (\textbf{I}, 5.2.2 et \textbf{II}, 5.5.5, (vi)). Tout revient à prouver qu'il existe un voisinage ouvert $U$ de $y$ tel que $p^{-1}(U)\subset j(X')$, car alors la restriction de $j$ à
\[
j^{-1}\bigl(p^{-1}(U)\bigr)=g^{-1}\bigl(f^{-1}(U)\bigr)=V
\]
sera un isomorphisme sur $p^{-1}(U)$, donc la restriction $V\to U$ de $f\circ g$ sera propre, et comme la restriction $V\to f^{-1}(U)$ de $g$ est surjective, la restriction $f^{-1}(U)\to U$ de $f$ sera propre (\textbf{II}, 5.4.3). Comme $T=P-j(X')$ est fermé dans $P$, $p(T)$ est fermé dans $Y$, et il suffit de montrer que $y\notin p(T)$, ou encore que tout $z'\in P$ tel que $p(z')=y$ \emph{appartient à} $j(X')$. Or, le corps des fractions rationnelles $R(P)$ est égal à $R(X')$, donc à $R(X)$ par construction; comme on peut se borner au cas où $z'$ n'est pas le point générique de $P$, il existe un anneau de valuation discrète $A$ ayant $R(X)$ pour corps des fractions et dominant $\mathcal O_{z'}$ (\textbf{II}, 7.1.7); comme $p(z')=y$, $A$ domine aussi $\mathcal O_y$, donc par hypothèse il existe un $x\in X$ tel que $f(x)=y$ et que $A$ domine $\mathcal O_x$. Comme $g$ est propre, il existe $x'\in X'$ tel que $g(x')=x$ et que $A$ domine $\mathcal O_{x'}$ (\textbf{II}, 7.3.10); donc $z'$ et $j(x')$ sont deux points de $P$ dont les anneaux locaux $\mathcal O_z$ et $\mathcal O_{j(x')}=\mathcal O_{x'}$ sont apparentés (\textbf{I}, 8.1.4); comme $P$ est un schéma, cela entraîne $z'=j(x')$ (\textbf{I}, 8.2.2), ce qui achève la démonstration.
\end{proof}

\begin{remark}[15.7.3]
\label{IV.15.7.3-fr}
On peut dans (15.7.2) supprimer l'hypothèse que $Y$ est localement noethérien, en remplaçant dans la condition c) les anneaux de valuation discrète par des anneaux de valuation quelconque; la démonstration est alors inchangée, compte tenu de (\textbf{II}, 7.3.10).
\end{remark}

\begin{corollary}[15.7.4]
\label{IV.15.7.4-fr}
Soient $Y$ un préschéma noethérien, $f:X\to Y$ un morphisme séparé de type fini, $Z$ une partie de $X$ tel que les points maximaux de son adhérence $\overline Z$ appartiennent à $Z$ (ce qui est le cas lorsque $Z$ est fini, ou lorsque $Z$ est constructible ($\mathbf{0}_{\mathrm{III}}$, 9.2.2)). Soit $y$ un point de $Y$. Les conditions suivantes sont équivalentes :
\begin{enumerate}
  \item[\rm a)] L'ensemble $\overline Z$ est propre sur $Y$ au point $y$.
  \item[\rm b)] Si $Y_1=\operatorname{Spec}(\mathcal O_y)$, $X_1=X\times_Y Y_1$, et si $g_1:X_1\to X$ est la projection canonique et $Z_1=g_1^{-1}(Z)$, l'adhérence $\overline Z_1$ de $Z_1$ dans $X_1$ est propre sur $Y_1$.
  \item[\rm c)] Pour tout schéma $Y'=\operatorname{Spec}(A')$, où $A'$ est un anneau de valuation discrète, et tout morphisme $g:Y'\to Y$, tel que l'image $g(y')$ du point fermé $y'$ de $Y'$ soit $y$, on a la propriété suivante : posant $X'=X\times_Y Y'$, $f'=f_{(Y')}:X'\to Y'$, $g'=g_{(X)}:X'\to X$, $Z'=g'^{-1}(Z)$, et désignant par $\eta'$
\oldpage[IV-3]{244}
le point générique de $Y'$, alors, pour tout point $x'\in Z'\cap f'^{-1}(\eta')$ rationnel sur $\mathbf k(\eta')$, il existe une $Y'$-section de $X'$ contenant $x'$.
\end{enumerate}
\end{corollary}

\begin{proof}
Le fait que a) entraîne b) résulte de ce que $\overline Z_1\subset g_1^{-1}(\overline Z)$ et de ce que $g_1^{-1}(\overline Z)$ est propre sur $Y_1$ (15.7.1). Le fait que b) entraîne c) résulte de (\textbf{II}, 7.3.3). Reste donc à voir que c) implique a). Il suffit évidemment de prouver que toute composante irréductible de $\overline Z$ est propre sur $Y$ au point $y$, et l'hypothèse permet donc de se borner au cas où $Z$ est réduit à un seul point $z$. On peut évidemment aussi, compte tenu de (\textbf{I}, 5.2.2) et (\textbf{II}, 5.4.6), supposer que $X$ et $Y$ sont réduits, puis (par définition d'une partie propre (\textbf{II}, 5.4.10)) supposer que $\overline Z=X=\overline{\{z\}}$ et (appliquant encore (\textbf{I}, 5.2.2)) que $f$ est dominant, donc $X$ et $Y$ intègres et noethériens; il faut alors prouver que $f$ est propre au point $y$ de $Y$. Montrons pour cela que $f$ vérifie la condition c) de (15.7.2). Soit $A'$ un anneau de valuation discrète ayant $R(X)=\mathbf k(z)$ pour corps des fractions et dominant $\mathcal O_y$; avec les notations de c), on a donc $g(y')=y$, et $g(\eta')$ est le point générique $\eta=f(z)$ de $Y$. Comme par définition $\mathbf k(\eta')=\mathbf k(z)$, il existe un $\mathbf k(\eta)$-morphisme $\operatorname{Spec}(\mathbf k(\eta'))\to\operatorname{Spec}(\mathbf k(z))$ transformant $\eta'$ en $z$, donc une $\mathbf k(\eta')$-section de $f'^{-1}(\eta')$ (\textbf{I}, 3.3.14); si $z'$ est l'image de $\eta'$ par cette section, $z'$ est donc rationnel sur $\mathbf k(\eta')$ et $g'(z')=z$. On conclut donc de l'hypothèse c) qu'il existe une $Y'$-section $h'$ de $X'$ telle que $h'(\eta')=z'$; soit $h=g'\circ h':Y'\to X$ le $Y$-morphisme correspondant, et soit $x=h(y')$; il est clair que $f(x)=y$ et que $A'=\mathcal O_{y'}$ domine $\mathcal O_x$, ce qui achève la démonstration.
\end{proof}

\begin{corollary}[15.7.5]
\label{IV.15.7.5-fr}
Soient $Y$ un préschéma localement noethérien, $f:X\to Y$ un morphisme séparé de type fini, $y$ un point de $Y$. Les conditions suivantes sont équivalentes :
\begin{enumerate}
  \item[\rm a)] $f$ est propre au point $y$.
  \item[\rm b)] Si $Y_1=\operatorname{Spec}(\mathcal O_y)$, $X_1=X\times_Y Y_1$, le morphisme $f_1=f_{(Y_1)}:X_1\to Y_1$ est propre.
  \item[\rm c)] Pour tout schéma $Y'=\operatorname{Spec}(A')$, où $A'$ est un anneau de valuation discrète, et tout morphisme $g:Y'\to Y$, tel que l'image $g(y')$ du point fermé $y'$ de $Y'$ soit $y$, on a la propriété suivante : posant $X'=X\times_Y Y'$, $f'=f_{(Y')}:X'\to Y'$ et désignant par $\eta'$ le point générique de $Y'$, alors, pour tout $x'\in f'^{-1}(\eta')$, rationnel sur $\mathbf k(\eta')$, il existe une $Y'$-section de $X'$ contenant $x'$.
\end{enumerate}
\end{corollary}

\begin{corollary}[15.7.6]
\label{IV.15.7.6-fr}
Soient $Y$ un préschéma localement noethérien, $f:X\to Y$ un morphisme de type fini. Les conditions suivantes sont équivalentes :
\begin{enumerate}
  \item[\rm a)] Il existe un morphisme propre $g:X\to Z$ et une immersion $h:Z\to Y$ tels que $f=h\circ g$.
  \item[\rm b)] Le sous-espace $f(X)$ est localement fermé dans $Y$, et si $Z'$ est le sous-préschéma de $Y$ image fermée de $X$ par $f$ (\textbf{I}, 9.5.3 et 9.5.1), et $Z''$ le préschéma induit par $Z'$ sur l'ouvert $f(X)$ de $Z'$, de sorte que $f$ se factorise de façon unique en $f=j\circ g$ (où $j:Z''\to Y$ est l'injection canonique), alors $g:X\to Z''$ est propre.
  \item[\rm c)] Pour tout $y\in f(X)$, $f$ est propre au point $y$.
  \item[\rm d)] Pour tout schéma $Y'=\operatorname{Spec}(A')$, où $A'$ est un anneau de valuation discrète, et tout morphisme $g:Y'\to Y$ tel que $g(Y')\subset f(X)$, toute $Y'$-section rationnelle (\textbf{I}, 7.1.2) de $X'=X\times_Y Y'$ se prolonge de façon unique en une $Y'$-section de $X'$.
\end{enumerate}
\end{corollary}

\begin{proof}
Il est clair que b) entraîne a). Pour voir que a) entraîne c), remarquons d'abord que a) entraîne que $h(Z)$ est localement fermé dans $Y$ et $g(X)$ fermé dans $Z$, donc $f(X)$ est localement fermé dans $Y$; pour tout $y\in f(X)$, il y a par suite un voisinage ouvert $U$
\oldpage[IV-3]{245}
de $y$ dans $Y$ tel que $f(X)\cap U$ soit fermé dans $U$; alors la restriction $h^{-1}(U)\to U$ de $h$ est une immersion fermée, et la restriction $f^{-1}(U)=g^{-1}(h^{-1}(U))\to h^{-1}(U)$ de $g$ est propre, donc la restriction $f^{-1}(U)\to U$ de $f$ est propre (\textbf{II}, 5.4.2). Pour voir que c) entraîne b), notons d'abord que, pour tout $y\in f(X)$, il existe, en vertu de c), un voisinage ouvert $U$ de $y$ dans $Y$ tel que $f(X)\cap U$ soit fermé dans $U$, donc $f(X)$ est localement fermé, et par suite ouvert dans son adhérence. Comme cette dernière est l'espace sous-jacent de $Z'$, et que $f$ se factorise en $f=j_0\circ g_0$, où $j_0:Z'\to Y$ est l'injection canonique et $g_0$ est un morphisme de $X$ dans $Z'$, le fait que $g_0(X)=Z''$ (qui est ouvert dans $Z'$) entraîne la factorisation de l'énoncé de b); le fait que $g$ est propre résulte alors du caractère local (sur $Z''$) de cette propriété, et de (\textbf{II}, 5.4.3). Il est clair que c) implique d) en vertu de (15.7.5), l'unicité du prolongement d'une $Y'$-section rationnelle résultant de l'hypothèse que $f$ est séparé (\textbf{I}, 7.2.2). Prouvons enfin que d) entraîne c); en premier lieu, il résulte de d), compte tenu de (\textbf{II}, 7.2.3 et 7.2.4), que $f$ est séparé; on peut alors appliquer le critère (15.7.5, c)), qui montre que $f$ est propre en tout point de $f(X)$.
\end{proof}

\begin{remarks}[15.7.7]
\label{IV.15.7.7-fr}
(i) Dans (15.7.4, c)), (15.7.5, c)) et (15.7.6, d)), on peut se restreindre au cas où l'anneau de valuation discrète $A'$ est \emph{complet} et admet un corps résiduel \emph{algébriquement clos}. En effet, dans la démonstration de (15.7.4), si l'on considère un anneau de valuation discrète $A''$ complet dominant $A'$ et ayant un corps résiduel algébriquement clos ($\mathbf{0}_{\mathrm{III}}$, 10.3.1), l'hypothèse c) est alors vérifiée pour $Y''=\operatorname{Spec}(A'')$, $X''=X\times_Y Y''$ et $f''=f_{(Y'')}$; mais $X''=X'\times_{Y'}Y''$, et il résulte alors de la remarque (\textbf{II}, 7.3.9, (i)) que $f'$ est propre; par suite $Y'$, $X'$ et $f'$ vérifient aussi l'hypothèse c) (\textbf{II}, 7.3.8).

(ii) Compte tenu de (\textbf{II}, 7.3.8), la condition c) de (15.7.5) peut se remplacer par la condition que $f'$ est \emph{propre}.
\end{remarks}

\begin{env}[15.7.8]
\label{IV.15.7.8-fr}
On dit qu'un morphisme de préschémas $f:X\to Y$ est \emph{submersif} s'il est surjectif et si la topologie de $Y$ est égale au quotient de la topologie de $X$ par la relation d'équivalence définie par $f$. On dit que $f$ est \emph{universellement submersif} si pour tout changement de base $Y'\to Y$, le morphisme $f'=f_{(Y')}:X\times_Y Y'\to Y'$ est submersif. Tout morphisme surjectif \emph{universellement ouvert}, ou \emph{universellement fermé}, ou \emph{fidèlement plat et quasi-compact} est universellement submersif (2.3.12). Étant donné un morphisme $f:X\to Y$, on dit qu'un sous-ensemble $Z$ de $X$ est \emph{submersif sur $f(Z)$} si la topologie du sous-espace $f(Z)$ de $Y$ est quotient de la topologie du sous-espace $Z$ de $X$ par la relation d'équivalence définie par $f|Z$. On dit que $Z$ est \emph{universellement submersif sur $f(Z)$} si, pour tout changement de base $g:Y'\to Y$, l'image réciproque $Z'=p^{-1}(Z)$ de $Z$ par la projection $p:X\times_Y Y'\to Y'$ est un ensemble submersif sur $f'(Z')=g^{-1}(f(Z))$. On notera que si le morphisme $f$ admet une $Y$-section $h:Y\to X$, tout ensemble $Z$ contenant $h(Y)$ est \emph{universellement submersif sur $Y$}.
\end{env}

\begin{proposition}[15.7.9]
\label{IV.15.7.9-fr}
Soient $Y$ un préschéma localement noethérien, $y$ un point de $Y$, $f:X\to Y$ un morphisme séparé de type fini. Soit $Z$ une partie localement constructible de $X$ ayant les propriétés suivantes :
\begin{enumerate}
  \item[\rm (i)] Pour toute générisation $y'$ de $y$, $Z_{y'}=Z\cap f^{-1}(y')$ est une composante connexe de $f^{-1}(y')$ et est géométriquement connexe sur $\mathbf k(y')$ (4.5.2).
\oldpage[IV-3]{246}
  \item[\rm (ii)] $Z_y$ est propre sur $\operatorname{Spec}(\mathbf k(y))$.
  \item[\rm (iii)] $Z$ est universellement submersif sur $f(Z)$ (15.7.8).
\end{enumerate}
Alors $Z$ est propre sur $Y$ au point $y$ (autrement dit (15.7.1) il existe un voisinage ouvert $U$ de $y$ tel que $Z\cap f^{-1}(U)$ soit fermé dans $f^{-1}(U)$ et propre sur $U$).
\end{proposition}

\begin{proof}
Utilisant la méthode de (8.1.2, a)) et (8.10.5, (xii)), on est ramené à prouver que lorsque $Y=\operatorname{Spec}(A)$ est spectre d'un anneau \emph{local} noethérien, les hypothèses (i), (ii) et (iii) entraînent que $Z$ est \emph{propre} sur $Y$.

Notons d'abord que si $A'$ est un anneau local noethérien, $\varphi:A\to A'$ un homomorphisme local, $Y'=\operatorname{Spec}(A')$ et $g:Y'\to Y$ le morphisme correspondant à $\varphi$, l'image réciproque $Z'$ de $Z$ par la projection canonique $p:X\times_Y Y'\to X$ a vis-à-vis de $Y'$ les propriétés (i), (ii), (iii) de l'énoncé, compte tenu de (1.8.2); utilisant (2.7.1, (vii)), il revient donc au même de démontrer que $Z'$ est propre sur $Y'$, pourvu que $A'$ soit un $A$-module \emph{fidèlement plat}. Nous allons utiliser cette remarque en prenant $A'=\widehat A$ ($\mathbf{0}_{\mathrm{I}}$, 7.3.5), autrement dit nous pouvons nous borner au cas où $A$ est \emph{complet}. Comme $Z_y$ est une composante connexe de $f^{-1}(y)$, propre sur $\mathbf k(y)$, il résulte de (\textbf{III}, 5.5.2) que $X$ est \emph{somme} de deux sous-préschémas $X_0$, $X_1$ induits sur des parties ouvertes et fermées de $X$, telles que $X_0$ soit propre sur $Y$ et que $X_0\cap f^{-1}(y)=Z_y$. Posons $Z_0=X_0\cap Z$, $Z_1=X_1\cap Z$, de sorte que ces ensembles forment une partition de $Z$ en deux parties ouvertes et fermées dans $Z$. Comme les $Z_{y'}$ sont connexes, on a nécessairement $Z_{y'}\subset Z_0$ ou $Z_{y'}\subset Z_1$, donc $f(Z_0)\cap f(Z_1)=\varnothing$. Mais comme $Z_0$ et $Z_1$ sont saturés pour la relation d'équivalence définie par $f|Z$, il résulte de (iii) que $f(Z_0)$ et $f(Z_1)$ sont ouverts dans $f(Z)$. Mais $f(Z_0)$ contient $y$, et tout voisinage de $y$ dans $Y$ est égal à $Y$, donc $f(Z_0)=f(Z)$, et par suite $f(Z_1)=\varnothing$, donc $Z_1=\varnothing$ et $Z\subset X_0$.

On peut donc désormais supposer de plus que $f$ est \emph{propre}, et il reste à prouver que $Z$ est fermé dans $X$. Autrement dit, il suffira de prouver qu'une partie \emph{constructible} $Z$ d'un préschéma $X$ \emph{propre} sur $Y$ qui satisfait aux \emph{deux seules conditions} (i) et (iii), est \emph{fermée} dans $X$. En vertu de ($\mathbf{0}_{\mathrm{III}}$, 9.2.5), il suffit donc de prouver que pour tout $x'\in Z$ et toute spécialisation $x$ de $x'$ dans $X$, on a $x\in Z$. Soient $A_1$ un anneau de valuation discrète \emph{complet}, $u$ un morphisme de $Y_1=\operatorname{Spec}(A_1)$ dans $X$ tel que, si $y_1$ et $y_1'$ sont le point fermé et le point générique de $Y_1$, on ait $u(y_1)=x$, $u(y_1')=x'$ (\textbf{II}, 7.1.7); posons $g=f\circ u:Y_1\to Y$ et $X_1=X\times_Y Y_1$; il existe une $Y_1$-section $u_1:Y_1\to X_1$ de $f_1=f_{(Y_1)}$ telle que $u=p\circ u_1$, où $p:X_1\to X$ est la projection canonique. Si $x_1=u_1(y_1)$, $x_1'=u_1(y_1')$, on a donc $p(x_1)=x$ et $p(x_1')=x'$, et $x_1$ est une spécialisation de $x_1'$. En outre, on a $x_1'\in Z_1=p^{-1}(Z)$; comme nous avons remarqué que les conditions (i) et (iii) sont stables pour \emph{tout} changement de base, ainsi que la propriété d'être constructible, on voit qu'on est ramené à la situation suivante : $Y$ est le spectre d'un anneau de valuation discrète complet, $X$ est propre sur $Y$, $f(x)=y$ est le point fermé de $Y$, $y'=f(x')$ son point générique, il existe une $Y$-section $h$ de $X$ telle que $h(y)=x$, $h(y')=x'$, et enfin on a $x'\in Z$; il faut prouver que $x\in Z$. Or, $Z_y$ est une composante connexe de $f^{-1}(y)$, \emph{propre} sur $\operatorname{Spec}(\mathbf k(y))$ puisque $X$ est propre sur $Y$. Appliquant de nouveau (\textbf{III}, 5.5.2), on voit comme plus haut qu'il y a une partie ouverte et fermée $X'$ de $X$ telle que $X'\cap f^{-1}(y)=Z_y$; comme $Z_y$ est connexe,
\oldpage[IV-3]{247}
on peut supposer $X'$ connexe (en remplaçant au besoin $X'$ par celle de ses composantes connexes contenant $Z_y$). Mais alors le même raisonnement qu'au début prouve qu'on a nécessairement $Z\subset X'$; comme $h(Y)$ est connexe et contient $x'$, il est nécessairement contenu dans $X'$, donc $x=h(y)\in X'\cap f^{-1}(y)=Z_y$. C.Q.F.D.
\end{proof}

\begin{corollary}[15.7.10]
\label{IV.15.7.10-fr}
Soient $Y$ un préschéma localement noethérien, $f:X\to Y$ un morphisme séparé de type fini, universellement submersif (15.7.8) et tel que toutes les fibres $X_y=f^{-1}(y)$ soient géométriquement connexes. Alors, si $y\in Y$ est tel que $X_y$ soit propre sur $\mathbf k(y)$, $f$ est propre au point $y$.
\end{corollary}

\begin{proof}
Compte tenu de (15.7.5), on est ramené au cas où $Y=\operatorname{Spec}(\mathcal O_y)$ puisque les hypothèses sont stables par changement de base. Mais dans ce cas il suffit d'appliquer (15.7.9) à $Z=X$.
\end{proof}

\begin{corollary}[15.7.11]
\label{IV.15.7.11-fr}
Soit $f:X\to Y$ un morphisme séparé de présentation finie; supposons qu'il existe un morphisme surjectif de présentation finie $g:Y'\to Y$, qui soit en outre propre ou plat, et tel que si l'on pose $X'=X\times_Y Y'$, il existe une $Y'$-section de $X'$ (ou encore un $Y$-morphisme $Y'\to X$ (\textbf{I}, 3.3.14)). Supposons enfin que les fibres $X_y=f^{-1}(y)$ soient géométriquement connexes. Alors l'ensemble $U$ des $y\in Y$ tels que $X_y$ soit propre sur $\mathbf k(y)$ est ouvert dans $Y$, et la restriction $f^{-1}(U)\to U$ de $f$ est propre.
\end{corollary}

\begin{proof}
La question étant locale sur $Y$, on peut supposer $Y=\operatorname{Spec}(A)$ affine. Appliquant (8.9.1) à $f$ et à $g$, on voit qu'il existe un sous-anneau noethérien $A_0$ de $A$ et deux morphismes de type fini $f_0:X_0\to Y_0=\operatorname{Spec}(A_0)$, $g_0:Y_0'\to Y_0$ tels que $X=X_0\times_{Y_0}Y$, $Y'=Y_0'\times_{Y_0}Y$, $f=(f_0)_{(Y)}$ et $g=(g_0)_{(Y)}$; en outre, utilisant (8.10.5, (v), (vi) et (xii)) et (11.2.7), on peut supposer $f_0$ séparé, $g_0$ surjectif et propre (resp. plat) si $g$ est propre (resp. plat); utilisant (8.8.2), on peut de plus supposer qu'il existe une $Y_0'$-section de $X_0$. D'autre part, utilisant (9.7.7), on peut appliquer (9.3.3) à la propriété d'être géométriquement connexe, et on peut donc supposer $A_0$ choisi de sorte que les $f_0^{-1}(y_0)$ soient géométriquement connexes pour tout $y_0\in Y_0$. Enfin, utilisant (2.7.1, (vii)), il revient au même de dire que $f^{-1}(y)$ est propre sur $\mathbf k(y)$ ou que $f_0^{-1}(y_0)$ est propre sur $\mathbf k(y_0)$, où $y_0$ est la projection de $y$ dans $Y_0$. Ces remarques montrent qu'on est ramené à démontrer le corollaire lorsque $Y$ est \emph{noethérien}. Notons maintenant le

\begin{lemma}[15.7.11.1]
\label{IV.15.7.11.1-fr}
Soient $f:X\to Y$ un morphisme, $g:Y'\to Y$ un morphisme surjectif, qui est en outre propre ou plat et quasi-compact. Posons $X'=X\times_Y Y'$, $f'=f_{(Y')}:X'\to Y'$. Alors, si $f'$ est submersif (resp. universellement submersif), il en est de même de $f$.
\end{lemma}

\begin{proof}
On peut se borner au cas où $f'$ est submersif. En effet, supposons le lemme prouvé dans ce cas, et montrons que si $f'$ est universellement submersif, $f$ l'est aussi. Soit donc $Y_1\to Y$ un morphisme quelconque, et posons $Y_1'=Y'\times_Y Y_1$; si $X_1'=X'\times_{Y'}Y_1'$ et $f_1'=(f')_{(Y_1')}:X_1'\to Y_1'$, $f_1'$ est submersif par hypothèse; de plus le morphisme $g_1:Y_1'\to Y_1$ est surjectif, et propre (resp. plat et quasi-compact) si $g$ l'est. Donc, si l'on pose $X_1=X\times_Y Y_1$, $f_1=f_{(Y_1)}$, il résulte de l'hypothèse et du fait que $X_1'=X_1\times_{Y_1}Y_1'$, que $f_1$ est submersif, donc $f$ universellement submersif. Supposons donc seulement que $f'$ soit submersif. Il est clair tout d'abord que $f$ est surjectif. Soit $E$ une partie de $Y$ telle que $f^{-1}(E)$ soit fermé dans $X$; alors, si $p:X'\to X$ est la projection canonique,
\oldpage[IV-3]{248}
$p^{-1}(f^{-1}(E))=f'^{-1}(g^{-1}(E))$ est fermé dans $X'$, donc, puisque $f'$ est submersif, $g^{-1}(E)$ est fermé dans $Y'$; mais comme $g$ est aussi universellement submersif (15.7.8), $E$ est fermé dans $Y$, ce qui prouve le lemme.
\end{proof}

Cela étant, comme il existe par hypothèse une $Y'$-section de $X'$, $f'$ est universellement submersif (15.7.8), donc il en est de même de $f$ d'après le lemme (15.7.11.1). Il suffit alors d'appliquer (15.7.10), en remarquant que l'ensemble des $y\in Y$ tels que $f$ soit propre en $y$ est par définition ouvert dans $Y$.
\end{proof}

\begin{center}
\emph{(A suivre.)}
\end{center}

\clearpage
\oldpage[IV-3]{249}
\phantomsection
\addcontentsline{toc}{section}{Bibliographie (suite)}
\section*{BIBLIOGRAPHIE (\emph{suite})}

\begingroup
\small
\setlength{\parindent}{0pt}
\setlength{\parskip}{3pt}
[39] J.-P.~\textsc{Serre}, \emph{Corps locaux}, Paris (Hermann), 1962.

[40] J.-P.~\textsc{Serre}, \emph{Groupes proalgébriques}, \emph{Publ. Inst. Hautes Études Sci.}, n\textsuperscript{o} 7 (1960).

\endgroup

\clearpage
\oldpage[IV-3]{250}
\phantomsection
\addcontentsline{toc}{section}{Index des notations}
\section*{INDEX DES NOTATIONS}

\begingroup
\small
\setlength{\parindent}{0pt}
\setlength{\parskip}{2pt}

$\mathfrak C(X)$, $\mathfrak D_c(X)$, $\mathfrak F_c(X)$,
$\mathfrak{LF}_c(X)$ ($X$ préschéma) : 8.3.9.

$\mathfrak Q(\mathcal F)$ ($\mathcal F$ $\mathcal O_X$-Module) : 8.5.10.

$\mathfrak{Spr}(Y)$, $\mathfrak{Spro}(Y)$, $\mathfrak{Sprf}(Y)$
($Y$ préschéma) : 8.6.1.

$\mathbf{Pro}(\mathcal C)$ ($\mathcal C$ catégorie) : 8.13.3.

$X_s$, $\mathcal F_s$, $u_s$, $g_s$, $Z_s$, $h_s$
($X$ $S$-préschéma, $s\in S$, $\mathcal F$ $\mathcal O_X$-Module,
$u$ homomorphisme de $\mathcal O_X$-Modules, $g$ section de
$\mathcal O_X$-Module, $Z$ partie de $X$, $h:X\to Y$ $S$-morphisme) : 9.4.1.

$\mathfrak{Qc}(X)$, $\mathfrak{Lqc}(X)$, $\mathfrak O(X)$,
$\mathfrak F(X)$, $\mathfrak{Lf}(X)$, $\mathfrak{Lc}(X)$
($X$ espace topologique) : 10.1.1.

$\operatorname{prof}^{*}_{x}(\mathcal F)$,
$\operatorname{prof}^{*}_{Z}(\mathcal F)$,
$\operatorname{prof}^{*}(\mathcal F)$
($\mathcal F$ $\mathcal O_X$-Module) : 10.8.1.

$S(X)$ ($X$ préschéma de Jacobson) : 10.9.1.

$S(f)$ ($f$ morphisme de préschémas de Jacobson) : 10.9.2.

$\operatorname{Spm}(A)$, $D^m(h)$ ($A$ anneau de Jacobson, $h\in A$) : 10.9.3.

$\operatorname{Tor.dim}_B(M)$ ($M$ $B$-module) : 12.3.2.

\endgroup

\clearpage
\oldpage[IV-3]{251}
\phantomsection
\addcontentsline{toc}{section}{Index terminologique}
\section*{INDEX TERMINOLOGIQUE}

\begingroup
\footnotesize
\setlength{\parindent}{0pt}
\setlength{\parskip}{1.1pt}
\newcommand{\EGAIVthreeTerm}[2]{#1 : #2.\par}

\EGAIVthreeTerm{Critère de platitude par fibres}{11.3.10}
\EGAIVthreeTerm{Espace algébrique sur $k$, espace $k$-algébrique, espace
préalgébrique sur $k$, espace $k$-préalgébrique}{10.10.2}
\EGAIVthreeTerm{Espace de Jacobson}{10.3.1}
\EGAIVthreeTerm{Espace $k$-préalgébrique réduit}{10.10.4}
\EGAIVthreeTerm{Famille de morphismes schématiquement dominante}{11.10.2}
\EGAIVthreeTerm{Famille de morphismes universellement schématiquement dominante}{11.10.7}
\EGAIVthreeTerm{Famille de sous-préschémas schématiquement dense}{11.10.2}
\EGAIVthreeTerm{Famille séparante d'homomorphismes de faisceaux de modules}{11.9.1}
\EGAIVthreeTerm{Famille séparante d'homomorphismes de modules}{11.9.4}
\EGAIVthreeTerm{Famille universellement séparante d'homomorphismes de faisceaux de modules}{11.9.14}
\EGAIVthreeTerm{Géométriquement isomorphes ($k$-préschémas)}{9.1.4}
\EGAIVthreeTerm{Homomorphisme de Modules universellement injectif}{11.9.14 et 11.9.18}
\EGAIVthreeTerm{Lemme de Chow pour les morphismes de présentation finie}{8.10.5.1}
\EGAIVthreeTerm{« Main theorem » de Zariski pour les morphismes de présentation finie}{8.12.6}
\EGAIVthreeTerm{Morphisme équidimensionnel en un point, morphisme équidimensionnel}{13.2.2 et 13.3.2}
\EGAIVthreeTerm{Morphisme propre en un point}{15.7.1}
\EGAIVthreeTerm{Morphisme pseudo-fini}{8.12.3}
\EGAIVthreeTerm{Morphisme submersif, morphisme universellement submersif}{15.7.8}
\EGAIVthreeTerm{Morphisme universellement ouvert en un point}{14.3.3}
\EGAIVthreeTerm{Ouvert ultraaffine}{10.9.5}
\EGAIVthreeTerm{Partie finie saturée d'un $k$-préschéma algébrique}{9.8.9}
\EGAIVthreeTerm{Partie propre en un point (d'un $S$-préschéma)}{15.7.1}
\EGAIVthreeTerm{Partie quasi-constructible, partie localement quasi-constructible d'un espace topologique}{10.1.1}
\EGAIVthreeTerm{Partie très dense d'un espace topologique}{10.1.3}
\EGAIVthreeTerm{Polynôme géométriquement irréductible}{9.7.4}
\EGAIVthreeTerm{Préschéma de Jacobson}{10.4.1}
\EGAIVthreeTerm{Préschéma équidimensionnel sur un autre en un point, équidimensionnel sur un autre}{13.2.2 et 13.3.2}
\EGAIVthreeTerm{Préschéma essentiellement affine sur un autre}{8.13.4}
\EGAIVthreeTerm{Principe de l'extension finie}{9.1.1}
\EGAIVthreeTerm{Profondeur rectifiée}{10.8.1}
\EGAIVthreeTerm{Pro-objet, morphisme de pro-objets}{8.13.3}
\EGAIVthreeTerm{Pro-objet essentiellement affine}{8.13.4, 8.14.1}
\EGAIVthreeTerm{Propriété constructible, propriété ind-constructible}{9.2.1 et 9.2.2, (i) et (ii)}
\EGAIVthreeTerm{Quasi-homéomorphisme d'espaces topologiques}{10.2.2}
\EGAIVthreeTerm{Quasi-isomorphisme d'espaces annelés}{10.2.8}
\EGAIVthreeTerm{Saturée d'une partie finie d'un $k$-préschéma algébrique}{9.8.9}
\EGAIVthreeTerm{Sous-ensemble d'un $S$-préschéma submersif, universellement submersif sur son image}{15.7.8}
\EGAIVthreeTerm{Squelette primaire d'un Module, squelette virtuel}{9.8.9}
\EGAIVthreeTerm{Spectre maximal d'un anneau de Jacobson}{10.9.3}
\EGAIVthreeTerm{Type primaire d'un Module, type primaire géométrique d'un Module}{9.8.9}
\EGAIVthreeTerm{Ultrapréschéma, morphisme d'ultrapréschémas}{10.9.5}
\EGAIVthreeTerm{$R$-ultrapréschéma}{10.10.2}

\let\EGAIVthreeTerm\relax
\endgroup

\clearpage
% La page imprimée 252 est blanche.
\thispagestyle{empty}
\null

\clearpage
\oldpage[IV-3]{253}
\phantomsection
\addcontentsline{toc}{section}{Table des matières}
\section*{TABLE DES MATIÈRES}

\begingroup
\small
\setlength{\parindent}{0pt}
\setlength{\parskip}{1.5pt}
\newcommand{\EGAIVthreeContentsLine}[3]{\textbf{#1}\quad #2\unskip\nobreak\hspace{0.4em}\nobreak\dotfill\nobreak\mbox{#3}\par}

\EGAIVthreeContentsLine{CHAPITRE IV.}{\textbf{Étude locale des schémas et des morphismes de schémas} \emph{(suite)}}{5}
\EGAIVthreeContentsLine{\S~8.}{Limites projectives de préschémas}{5}
\EGAIVthreeContentsLine{8.1.}{Introduction}{5}
\EGAIVthreeContentsLine{8.2.}{Limites projectives de préschémas}{7}
\EGAIVthreeContentsLine{8.3.}{Parties constructibles dans une limite projective de préschémas}{12}
\EGAIVthreeContentsLine{8.4.}{Critères d'irréductibilité et de connexion pour les limites projectives de préschémas}{17}
\EGAIVthreeContentsLine{8.5.}{Modules de présentation finie sur une limite projective de préschémas}{19}
\EGAIVthreeContentsLine{8.6.}{Sous-préschémas de présentation finie d'une limite projective de préschémas}{25}
\EGAIVthreeContentsLine{8.7.}{Critères pour qu'une limite projective de préschémas soit un préschéma réduit (resp. intègre)}{27}
\EGAIVthreeContentsLine{8.8.}{Préschémas de présentation finie sur une limite projective de préschémas}{28}
\EGAIVthreeContentsLine{8.9.}{Premières applications à l'élimination des hypothèses noethériennes}{34}
\EGAIVthreeContentsLine{8.10.}{Propriétés de permanence des morphismes par passage à la limite projective}{36}
\EGAIVthreeContentsLine{8.11.}{Application aux morphismes quasi-finis}{41}
\EGAIVthreeContentsLine{8.12.}{Nouvelle démonstration et généralisation du « Main Theorem » de Zariski}{43}
\EGAIVthreeContentsLine{8.13.}{Traduction en termes de pro-objets}{49}
\EGAIVthreeContentsLine{8.14.}{Caractérisation d'un préschéma localement de présentation finie sur un autre, en termes du foncteur qu'il représente}{52}

\EGAIVthreeContentsLine{\S~9.}{Propriétés constructibles}{54}
\EGAIVthreeContentsLine{9.1.}{Le principe de l'extension finie}{54}
\EGAIVthreeContentsLine{9.2.}{Propriétés constructibles et ind-constructibles}{56}
\EGAIVthreeContentsLine{9.3.}{Propriétés constructibles de morphismes de préschémas algébriques}{60}

\clearpage
\oldpage[IV-3]{254}
\EGAIVthreeContentsLine{9.4.}{Constructibilité de certaines propriétés des Modules}{62}
\EGAIVthreeContentsLine{9.5.}{Constructibilité de propriétés topologiques}{67}
\EGAIVthreeContentsLine{9.6.}{Constructibilité de certaines propriétés des morphismes}{71}
\EGAIVthreeContentsLine{9.7.}{Constructibilité des propriétés de séparabilité, d'irréductibilité géométrique et de connexité géométrique}{76}
\EGAIVthreeContentsLine{9.8.}{Décomposition primaire au voisinage d'une fibre générique}{83}
\EGAIVthreeContentsLine{9.9.}{Constructibilité des propriétés locales des fibres}{88}

\EGAIVthreeContentsLine{\S~10.}{Préschémas de Jacobson}{95}
\EGAIVthreeContentsLine{10.1.}{Parties très denses d'un espace topologique}{95}
\EGAIVthreeContentsLine{10.2.}{Quasi-homéomorphismes}{97}
\EGAIVthreeContentsLine{10.3.}{Espaces de Jacobson}{101}
\EGAIVthreeContentsLine{10.4.}{Préschémas de Jacobson et anneaux de Jacobson}{101}
\EGAIVthreeContentsLine{10.5.}{Préschémas de Jacobson noethériens}{104}
\EGAIVthreeContentsLine{10.6.}{Dimension dans les préschémas de Jacobson}{107}
\EGAIVthreeContentsLine{10.7.}{Exemples et contre-exemples}{109}
\EGAIVthreeContentsLine{10.8.}{Profondeur rectifiée}{110}
\EGAIVthreeContentsLine{10.9.}{Spectres maximaux et ultrapréschémas}{112}
\EGAIVthreeContentsLine{10.10.}{Espaces algébriques de Serre}{114}

\EGAIVthreeContentsLine{\S~11.}{Propriétés topologiques des morphismes plats de présentation finie. Critères de platitude}{116}
\EGAIVthreeContentsLine{11.1.}{Ensembles de platitude (cas noethérien)}{117}
\EGAIVthreeContentsLine{11.2.}{Platitude d'une limite projective de préschémas}{119}
\EGAIVthreeContentsLine{11.3.}{Application à l'élimination d'hypothèses noethériennes}{132}
\EGAIVthreeContentsLine{11.4.}{Descente de la platitude par des morphismes quelconques : cas d'un préschéma de base artinien}{143}
\EGAIVthreeContentsLine{11.5.}{Descente de la platitude par des morphismes quelconques : cas général}{150}
\EGAIVthreeContentsLine{11.6.}{Descente de la platitude par des morphismes quelconques : cas d'un préschéma de base unibranche}{154}
\EGAIVthreeContentsLine{11.7.}{Contre-exemples}{157}
\EGAIVthreeContentsLine{11.8.}{Un critère valuatif de platitude}{159}
\EGAIVthreeContentsLine{11.9.}{Familles séparantes et universellement séparantes d'homomorphismes de faisceaux de modules}{160}
\EGAIVthreeContentsLine{11.10.}{Familles schématiquement dominantes de morphismes et familles schématiquement denses de sous-préschémas}{170}

\EGAIVthreeContentsLine{\S~12.}{Étude des fibres des morphismes plats de présentation finie}{173}
\EGAIVthreeContentsLine{12.0.}{Introduction}{173}
\EGAIVthreeContentsLine{12.1.}{Propriétés locales des fibres d'un morphisme plat localement de présentation finie}{174}

\clearpage
\oldpage[IV-3]{255}
\EGAIVthreeContentsLine{12.2.}{Propriétés locales et globales des fibres d'un morphisme propre, plat et de présentation finie}{179}
\EGAIVthreeContentsLine{12.3.}{Propriétés cohomologiques locales des fibres d'un morphisme plat et localement de présentation finie}{183}

\EGAIVthreeContentsLine{\S~13.}{Morphismes équidimensionnels}{187}
\EGAIVthreeContentsLine{13.1.}{Le théorème de semi-continuité de Chevalley}{188}
\EGAIVthreeContentsLine{13.2.}{Morphismes équidimensionnels : cas des morphismes dominants de préschémas irréductibles}{190}
\EGAIVthreeContentsLine{13.3.}{Morphismes équidimensionnels : cas général}{194}

\EGAIVthreeContentsLine{\S~14.}{Morphismes universellement ouverts}{199}
\EGAIVthreeContentsLine{14.1.}{Morphismes ouverts}{200}
\EGAIVthreeContentsLine{14.2.}{Morphismes ouverts et formule des dimensions}{202}
\EGAIVthreeContentsLine{14.3.}{Morphismes universellement ouverts}{204}
\EGAIVthreeContentsLine{14.4.}{Le critère de Chevalley pour les morphismes universellement ouverts}{209}
\EGAIVthreeContentsLine{14.5.}{Morphismes universellement ouverts et quasi-sections}{216}

\EGAIVthreeContentsLine{\S~15.}{Étude des fibres d'un morphisme universellement ouvert}{223}
\EGAIVthreeContentsLine{15.1.}{Multiplicités des fibres d'un morphisme universellement ouvert}{223}
\EGAIVthreeContentsLine{15.2.}{Platitude des morphismes universellement ouverts à fibres géométriquement réduites}{226}
\EGAIVthreeContentsLine{15.3.}{Application : critères de réduction et d'irréductibilité}{228}
\EGAIVthreeContentsLine{15.4.}{Compléments sur les morphismes de Cohen-Macaulay}{229}
\EGAIVthreeContentsLine{15.5.}{Rang séparable des fibres d'un morphisme quasi-fini et universellement ouvert. Application aux composantes connexes géométriques des fibres d'un morphisme propre}{231}
\EGAIVthreeContentsLine{15.6.}{Composantes connexes des fibres le long d'une section}{236}
\EGAIVthreeContentsLine{15.7.}{Appendice : Critères valuatifs de propreté locale}{242}
\EGAIVthreeContentsLine{}{BIBLIOGRAPHIE \emph{(suite)}}{249}
\EGAIVthreeContentsLine{}{INDEX DES NOTATIONS}{250}
\EGAIVthreeContentsLine{}{INDEX TERMINOLOGIQUE}{251}

\vspace{2em}
\begin{center}
\emph{Manuscrit reçu le 15 novembre 1964.}
\end{center}

\let\EGAIVthreeContentsLine\relax
\endgroup

% FIN SOURCE ega4/ega4-15-terminal-fr.tex


% FIN SOURCE ega4/ega4-15-fr.tex

\end{document}
% EGA-COMPONENT-END IV3

% EGA-COMPONENT-BEGIN IV4 0C5850584139062701B56589A5A8D4AAFBA6259672719FF42B452CED4E38BA16
\documentclass[10pt,oneside]{book}

\usepackage[T1]{fontenc}
\usepackage[utf8]{inputenc}
\usepackage[french]{babel}
\usepackage{amsmath,amssymb,mathrsfs}
\usepackage[all]{xy}
\usepackage{tikz-cd}
\usepackage[shortlabels]{enumitem}
\usepackage{xparse}
\usepackage{xstring}
\usepackage[a4paper,margin=22mm]{geometry}
\usepackage[
  colorlinks=true,
  linkcolor=blue,
  filecolor=blue,
  hypertexnames=false,
  unicode=true,
  pdfencoding=auto
]{hyperref}
\hypersetup{
  pdftitle={Éléments de géométrie algébrique IV, quatrième partie -- édition française diplomatique},
  pdfauthor={Alexandre Grothendieck et Jean Dieudonné},
  pdfsubject={Transcription diplomatique française depuis les exemplaires NUMDAM}
}

\renewcommand{\thesection}{\arabic{section}}
\renewcommand{\thesubsection}{\thesection.\arabic{subsection}}
%
% DEBUT SOURCE preamble-base.tex
\def\sh{\mathcal}                   % sheaf font
\def\bb{\mathbf}                    % bold font
\def\cat{\mathtt}                   % category font
\def\leq{\leqslant}                 % <=
\def\geq{\geqslant}                 % >=
\def\setmin{-}                      % set minus
\def\rad{\mathfrak{r}}              % radical
\def\nilrad{\mathfrak{N}}           % nilradical
\def\emp{\varnothing}               % empty set
\def\vphi{\phi}                     % for switching \phi and \varphi, change if needed
\def\HH{\mathrm{H}}                 % cohomology H
\def\CHH{\check{\HH}}               % Čech cohomology H
\def\RR{\mathrm{R}}                 % right derived R
\def\LL{\mathrm{L}}                 % left derived L
\def\dual#1{{#1}^\vee}              % dual
\def\kres{k}                        % residue field k
\def\C{\cat{C}}                     % category C
\def\op{^\cat{op}}                  % opposite category
\def\Set{\cat{Set}}                 % category of sets
\def\CHom{\cat{Hom}}                % functor category
\def\supertilde{{\,\widetilde{\,}\,}}   % use \supertilde instead of ^\sim
\def\GL{\bb{GL}}
\def\bbGamma{\boldsymbol\Gamma}
\def\red{\mathrm{red}}
\def\rg{\operatorname{rg}}
\def\gr{\operatorname{gr}}
\def\Gr{\operatorname{Gr}}
\def\Hom{\operatorname{Hom}}
\def\Proj{\operatorname{Proj}}
\def\Tor{\operatorname{Tor}}
\def\Ext{\operatorname{Ext}}
\def\Supp{\operatorname{Supp}}
\def\Ker{\operatorname{Ker}}
\def\Im{\operatorname{Im}}
\def\Coker{\operatorname{Coker}}
\def\Spec{\operatorname{Spec}}
\def\Spm{\operatorname{Spm}}
\def\Spf{\operatorname{Spf}}
\def\grad{\operatorname{grad}}
\def\dimc{\operatorname{dimc}}
\def\codim{\operatorname{codim}}
\def\id{\operatorname{id}}
\def\Der{\operatorname{Der}}
\def\Diff{\operatorname{Diff}}
\def\Hyp{\operatorname{Hyp}}
\def\Tr{\operatorname{Tr}}

\renewcommand{\to}{\mathchoice{\longrightarrow}{\rightarrow}{\rightarrow}{\rightarrow}}
\newcommand{\from}{\mathchoice{\longleftarrow}{\leftarrow}{\leftarrow}{\leftarrow}}
\let\mapstoo\mapsto
\renewcommand{\mapsto}{\mathchoice{\longmapsto}{\mapstoo}{\mapstoo}{\mapstoo}}
\def\isoto{\simeq} % isomorphism
\def\simto{\xrightarrow{\sim}}

\renewcommand{\hat}[1]{\widehat{#1}}
\renewcommand{\tilde}[1]{\widetilde{#1}}

% FIN SOURCE preamble-base.tex

\makeatletter
\newcommand{\EGAUnavailableHyperref}[2]{%
  \begingroup
  \let\@latex@warning\@gobble
  \let\@latex@warning@no@line\@gobble
  \let\G@refundefinedtrue\relax
  \hyperref[#1]{#2}%
  \endgroup
}
\NewDocumentCommand{\sref}{o m o}{%
  \@ifundefined{r@#2}%
    {\IfNoValueTF{#1}%
      {\IfNoValueTF{#3}%
        {\EGAUnavailableHyperref{#2}{\normalfont{(\StrBehind{#2}{.})}}}%
        {\EGAUnavailableHyperref{#2}{\normalfont{(\StrBehind{#2}{.},~{#3})}}}}%
      {\IfNoValueTF{#3}%
        {\EGAUnavailableHyperref{#2}{\normalfont{(\textbf{#1},~\StrBehind{#2}{.})}}}%
        {\EGAUnavailableHyperref{#2}{\normalfont{(\textbf{#1},~\StrBehind{#2}{.},~{#3})}}}}}%
    {\IfNoValueTF{#1}%
      {\IfNoValueTF{#3}%
        {\hyperref[#2]{\normalfont{(\StrBehind{#2}{.})}}}%
        {\hyperref[#2]{\normalfont{(\StrBehind{#2}{.},~{#3})}}}}%
      {\IfNoValueTF{#3}%
        {\hyperref[#2]{\normalfont{(\textbf{#1},~\StrBehind{#2}{.})}}}%
        {\hyperref[#2]{\normalfont{(\textbf{#1},~\StrBehind{#2}{.},~{#3})}}}}}%
}
\newcommand{\EGAIVIndexExternalLabel}[3]{%
  \expandafter\gdef\csname r@canon-index.#1\endcsname{{#1}{0}{}{ega-fr:#2:#3}{EGA_French_Diplomatic_Cumulative_Reader.pdf}}}
\EGAIVIndexExternalLabel{ega4.c04.cor.13-1-4}{iv3}{ega4.c04.cor.13-1-4}
\EGAIVIndexExternalLabel{ega4.c04.cor.14-4-4}{iv3}{ega4.c04.cor.14-4-4}
\EGAIVIndexExternalLabel{ega4.c04.cor.14-5-11}{iv3}{ega4.c04.cor.14-5-11}
\EGAIVIndexExternalLabel{ega4.c04.cor.15-5-2}{iv3}{ega4.c04.cor.15-5-2}
\EGAIVIndexExternalLabel{ega4.c04.cor.15-6-2}{iv3}{ega4.c04.cor.15-6-2}
\EGAIVIndexExternalLabel{ega4.c04.cor.3-1-5}{iv2}{ega4.c04.cor.3-1-5}
\EGAIVIndexExternalLabel{ega4.c04.cor.4-5-11}{iv2}{ega4.c04.cor.4-5-11}
\EGAIVIndexExternalLabel{ega4.c04.cor.5-11-2}{iv2}{ega4.c04.cor.5-11-2}
\EGAIVIndexExternalLabel{ega4.c04.cor.5-11-4}{iv2}{ega4.c04.cor.5-11-4}
\EGAIVIndexExternalLabel{ega4.c04.cor.5-2-3}{iv2}{ega4.c04.cor.5-2-3}
\EGAIVIndexExternalLabel{ega4.c04.cor.5-5-4}{iv2}{ega4.c04.cor.5-5-4}
\EGAIVIndexExternalLabel{ega4.c04.cor.6-1-3}{iv2}{ega4.c04.cor.6-1-3}
\EGAIVIndexExternalLabel{ega4.c04.cor.6-7-2}{iv2}{ega4.c04.cor.6-7-2}
\EGAIVIndexExternalLabel{ega4.c04.cor.7-2-9}{iv2}{ega4.c04.cor.7-2-9}
\EGAIVIndexExternalLabel{ega4.c04.cor.8-12-10}{iv3}{ega4.c04.cor.8-12-10}
\EGAIVIndexExternalLabel{ega4.c04.cor.8-12-11}{iv3}{ega4.c04.cor.8-12-11}
\EGAIVIndexExternalLabel{ega4.c04.cor.8-2-11}{iv3}{ega4.c04.cor.8-2-11}
\EGAIVIndexExternalLabel{ega4.c04.cor.8-5-8}{iv3}{ega4.c04.cor.8-5-8}
\EGAIVIndexExternalLabel{ega4.c04.def.7-8-2}{iv2}{ega4.c04.def.7-8-2}
\EGAIVIndexExternalLabel{ega4.c04.item.12-1-2}{iv3}{ega4.c04.item.12-1-2}
\EGAIVIndexExternalLabel{ega4.c04.item.15-6-10}{iv3}{ega4.c04.item.15-6-10}
\EGAIVIndexExternalLabel{ega4.c04.item.6-5-5}{iv2}{ega4.c04.item.6-5-5}
\EGAIVIndexExternalLabel{ega4.c04.item.7-5-4}{iv2}{ega4.c04.item.7-5-4}
\EGAIVIndexExternalLabel{ega4.c04.item.7-8-4}{iv2}{ega4.c04.item.7-8-4}
\EGAIVIndexExternalLabel{ega4.c04.lem.1-8-7-1}{iv1}{ega4.c04.lem.1-8-7-1}
\EGAIVIndexExternalLabel{ega4.c04.lem.2-3-1}{iv2}{ega4.c04.lem.2-3-1}
\EGAIVIndexExternalLabel{ega4.c04.lem.4-4-2}{iv2}{ega4.c04.lem.4-4-2}
\EGAIVIndexExternalLabel{ega4.c04.lem.5-6-11-1}{iv2}{ega4.c04.lem.5-6-11-1}
\EGAIVIndexExternalLabel{ega4.c04.prop.1-10-3}{iv1}{ega4.c04.prop.1-10-3}
\EGAIVIndexExternalLabel{ega4.c04.prop.10-2-1}{iv3}{ega4.c04.prop.10-2-1}
\EGAIVIndexExternalLabel{ega4.c04.prop.11-3-4}{iv3}{ega4.c04.prop.11-3-4}
\EGAIVIndexExternalLabel{ega4.c04.prop.14-3-8}{iv3}{ega4.c04.prop.14-3-8}
\EGAIVIndexExternalLabel{ega4.c04.prop.15-4-2}{iv3}{ega4.c04.prop.15-4-2}
\EGAIVIndexExternalLabel{ega4.c04.prop.15-6-6}{iv3}{ega4.c04.prop.15-6-6}
\EGAIVIndexExternalLabel{ega4.c04.prop.2-3-4}{iv2}{ega4.c04.prop.2-3-4}
\EGAIVIndexExternalLabel{ega4.c04.prop.2-4-3}{iv2}{ega4.c04.prop.2-4-3}
\EGAIVIndexExternalLabel{ega4.c04.prop.2-4-5}{iv2}{ega4.c04.prop.2-4-5}
\EGAIVIndexExternalLabel{ega4.c04.prop.2-5-2}{iv2}{ega4.c04.prop.2-5-2}
\EGAIVIndexExternalLabel{ega4.c04.prop.2-6-1}{iv2}{ega4.c04.prop.2-6-1}
\EGAIVIndexExternalLabel{ega4.c04.prop.3-1-7}{iv2}{ega4.c04.prop.3-1-7}
\EGAIVIndexExternalLabel{ega4.c04.prop.3-3-1}{iv2}{ega4.c04.prop.3-3-1}
\EGAIVIndexExternalLabel{ega4.c04.prop.4-2-1}{iv2}{ega4.c04.prop.4-2-1}
\EGAIVIndexExternalLabel{ega4.c04.prop.4-5-19}{iv2}{ega4.c04.prop.4-5-19}
\EGAIVIndexExternalLabel{ega4.c04.prop.4-7-3}{iv2}{ega4.c04.prop.4-7-3}
\EGAIVIndexExternalLabel{ega4.c04.prop.5-10-2}{iv2}{ega4.c04.prop.5-10-2}
\EGAIVIndexExternalLabel{ega4.c04.prop.5-6-1}{iv2}{ega4.c04.prop.5-6-1}
\EGAIVIndexExternalLabel{ega4.c04.prop.5-6-5}{iv2}{ega4.c04.prop.5-6-5}
\EGAIVIndexExternalLabel{ega4.c04.prop.5-9-8}{iv2}{ega4.c04.prop.5-9-8}
\EGAIVIndexExternalLabel{ega4.c04.prop.6-1-5}{iv2}{ega4.c04.prop.6-1-5}
\EGAIVIndexExternalLabel{ega4.c04.prop.6-14-1}{iv2}{ega4.c04.prop.6-14-1}
\EGAIVIndexExternalLabel{ega4.c04.prop.6-14-4}{iv2}{ega4.c04.prop.6-14-4}
\EGAIVIndexExternalLabel{ega4.c04.prop.6-15-7}{iv2}{ega4.c04.prop.6-15-7}
\EGAIVIndexExternalLabel{ega4.c04.prop.6-2-3}{iv2}{ega4.c04.prop.6-2-3}
\EGAIVIndexExternalLabel{ega4.c04.prop.6-3-1}{iv2}{ega4.c04.prop.6-3-1}
\EGAIVIndexExternalLabel{ega4.c04.prop.6-5-1}{iv2}{ega4.c04.prop.6-5-1}
\EGAIVIndexExternalLabel{ega4.c04.prop.7-8-6}{iv2}{ega4.c04.prop.7-8-6}
\EGAIVIndexExternalLabel{ega4.c04.prop.8-10-1}{iv3}{ega4.c04.prop.8-10-1}
\EGAIVIndexExternalLabel{ega4.c04.prop.8-10-3}{iv3}{ega4.c04.prop.8-10-3}
\EGAIVIndexExternalLabel{ega4.c04.prop.8-12-3}{iv3}{ega4.c04.prop.8-12-3}
\EGAIVIndexExternalLabel{ega4.c04.prop.8-9-1}{iv3}{ega4.c04.prop.8-9-1}
\EGAIVIndexExternalLabel{ega4.c04.rem.14-5-7}{iv3}{ega4.c04.rem.14-5-7}
\EGAIVIndexExternalLabel{ega4.c04.statement.1-7-4}{iv1}{ega4.c04.statement.1-7-4}
\EGAIVIndexExternalLabel{ega4.c04.statement.10-4-11}{iv3}{ega4.c04.statement.10-4-11}
\EGAIVIndexExternalLabel{ega4.c04.statement.11-7-5}{iv3}{ega4.c04.statement.11-7-5}
\EGAIVIndexExternalLabel{ega4.c04.statement.13-3-9}{iv3}{ega4.c04.statement.13-3-9}
\EGAIVIndexExternalLabel{ega4.c04.statement.2-1-1}{iv2}{ega4.c04.statement.2-1-1}
\EGAIVIndexExternalLabel{ega4.c04.statement.6-10-1}{iv2}{ega4.c04.statement.6-10-1}
\EGAIVIndexExternalLabel{ega4.c04.statement.6-15-2}{iv2}{ega4.c04.statement.6-15-2}
\EGAIVIndexExternalLabel{ega4.c04.statement.6-15-4}{iv2}{ega4.c04.statement.6-15-4}
\EGAIVIndexExternalLabel{ega4.c04.statement.7-3-19}{iv2}{ega4.c04.statement.7-3-19}
\EGAIVIndexExternalLabel{ega4.c04.statement.7-8-3}{iv2}{ega4.c04.statement.7-8-3}
\EGAIVIndexExternalLabel{ega4.c04.statement.8-1-2}{iv3}{ega4.c04.statement.8-1-2}
\EGAIVIndexExternalLabel{ega4.c04.subsec.10-4}{iv3}{ega4.c04.subsec.10-4}
\EGAIVIndexExternalLabel{ega4.c04.thm.11-2-9}{iv3}{ega4.c04.thm.11-2-9}
\EGAIVIndexExternalLabel{ega4.c04.thm.11-3-10}{iv3}{ega4.c04.thm.11-3-10}
\EGAIVIndexExternalLabel{ega4.c04.thm.12-1-1}{iv3}{ega4.c04.thm.12-1-1}
\EGAIVIndexExternalLabel{ega4.c04.thm.14-2-1}{iv3}{ega4.c04.thm.14-2-1}
\EGAIVIndexExternalLabel{ega4.c04.thm.14-4-1}{iv3}{ega4.c04.thm.14-4-1}
\EGAIVIndexExternalLabel{ega4.c04.thm.5-10-5}{iv2}{ega4.c04.thm.5-10-5}
\EGAIVIndexExternalLabel{ega4.c04.thm.5-8-6}{iv2}{ega4.c04.thm.5-8-6}
\EGAIVIndexExternalLabel{ega4.c04.thm.7-6-4}{iv2}{ega4.c04.thm.7-6-4}
\EGAIVIndexExternalLabel{ega4.c04.thm.7-7-2}{iv2}{ega4.c04.thm.7-7-2}
\EGAIVIndexExternalLabel{ega4.c04.thm.8-10-5}{iv3}{ega4.c04.thm.8-10-5}
\EGAIVIndexExternalLabel{ega4.c04.thm.8-11-1}{iv3}{ega4.c04.thm.8-11-1}
\EGAIVIndexExternalLabel{ega4.c04.thm.8-8-2}{iv3}{ega4.c04.thm.8-8-2}
\newcommand{\EGAIVFourExternalLabel}[3]{%
  \expandafter\gdef\csname r@#1\endcsname{{#1}{0}{}{ega-fr:#2:canon-ref.#3}{EGA_French_Diplomatic_Cumulative_Reader.pdf}}}
\newcommand{\EGAIVFourExternalCitation}[2]{%
  \expandafter\gdef\csname b@#1\endcsname{#2}}
\makeatother
\newcommand{\oldpage}[2][]{%
  \setcounter{footnote}{0}%
  \marginpar{\scriptsize\textbf{#1}~|~#2}\ignorespaces}
\newenvironment{env}[1][]{\par\medskip\noindent\textbf{(#1)}\ }{\par}
\newenvironment{definition}[1][]{%
  \par\medskip\noindent\itshape Définition (#1). ---\ }{\par}
\newenvironment{examples}[1][]{%
  \par\medskip\noindent\itshape Exemples (#1). ---\ }{\par}
\newenvironment{proposition}[1][]{%
  \par\medskip\noindent\itshape Proposition (#1). ---\ }{\par}
\newenvironment{theorem}[1][]{%
  \par\medskip\noindent\itshape Théorème (#1). ---\ }{\par}
\newenvironment{corollary}[1][]{%
  \par\medskip\noindent\itshape Corollaire (#1). ---\ }{\par}
\newenvironment{remark}[1][]{%
  \par\medskip\noindent\itshape Remarque (#1). ---\ }{\par}
\newenvironment{remarks}[1][]{%
  \par\medskip\noindent\itshape Remarques (#1). ---\ }{\par}
\newenvironment{lemma}[1][]{%
  \par\medskip\noindent\itshape Lemme (#1). ---\ }{\par}
\newenvironment{proof}{%
  \par\medskip\noindent\textit{Démonstration. ---\ }}{\par}

\makeatletter
\newcommand{\EGAReaderTarget}[2]{\hypertarget{#1}{}\@bsphack\begingroup\def\@currentHref{#1}\label{#2}\endgroup\@esphack}
\makeatother
\begin{document}
% BEGIN RIGHTS AND AI TRANSCRIPTION NOTICE
\clearpage
\begingroup
\small
\noindent\textbf{Droits, disponibilité et provenance}

Les droits sur les œuvres originales et sur toute contribution protégée à ces traductions et transcriptions dérivées restent ceux de leurs titulaires respectifs. Ce projet ne place pas ces œuvres ou contributions dans le domaine public et n'accorde aucune licence sur elles. Les mentions de droits de tiers et les autorisations déjà accordées restent applicables.

Ces éditions sont mises à disposition gratuitement, à des fins non commerciales et exclusivement pédagogiques, dans la présomption que les titulaires de droits ne s'opposeraient pas à la libre disponibilité de ces traductions et transcriptions aux côtés des originaux français sur NUMDAM (www.numdam.org/journals/PMIHES/). Cette présomption ne signifie pas qu'une autorisation a été obtenue ; le libre accès aux originaux n'autorise pas, à lui seul, leur traduction, leur adaptation ou leur republication.

Si vous détenez des droits sur un élément présenté ici et vous opposez à sa disponibilité, veuillez ouvrir une demande sur GitHub (github.com/KokunoYumeto/ega-fr/issues). Le responsable de la publication retirera l'élément concerné à la demande du titulaire des droits.

La transcription moderne, les corrections, la composition et les notes mathématiques sont produites par IA. Les travaux documentés comprennent ceux d'OpenAI Codex, GPT-6 Astra, effort Ultra. Les indications de provenance des contributions antérieures restent applicables. Les textes originaux restent l'œuvre d'Alexandre Grothendieck et Jean Dieudonné ; ils ne sont pas attribués à l'IA. Cette édition n'est approuvée ni par les auteurs, ni par l'IHÉS, ni par NUMDAM.
\par
\endgroup
\clearpage
% END RIGHTS AND AI TRANSCRIPTION NOTICE

\EGAIVFourExternalLabel{0.4.3.3}{base}{0.4.3.3-fr}
\EGAIVFourExternalLabel{0.5.4.3}{base}{0.5.4.3-fr}
\EGAIVFourExternalLabel{0.5.4.5}{base}{0.5.4.5-fr}
\EGAIVFourExternalLabel{0.5.4.7}{base}{0.5.4.7-fr}
\EGAIVFourExternalLabel{0.6.1.2}{base}{0.6.1.2-fr}
\EGAIVFourExternalLabel{0.16.4.6}{iv1}{0.16.4.6-fr}
\EGAIVFourExternalCitation{IV-37}{37}
\setcounter{section}{15}
%
% DEBUT SOURCE ega4/ega4-16-fr.tex
\section{Invariants différentiels. Morphismes différentiellement lisses}
\label{section:IV.16}

\oldpage[IV-4]{5}
Dans ce paragraphe, nous présentons, sous forme globale, quelques notions de
calcul différentiel particulièrement utiles en Géométrie algébrique. Nous passons sous
silence de nombreux développements, classiques en Géométrie différentielle (connexions,
transformations infinitésimales associées à un champ de vecteurs, jets, etc.), bien que ces
notions s'écrivent de façon particulièrement naturelle dans le cadre des schémas. Nous
passons également sous silence ici les phénomènes spéciaux à la caractéristique $p>0$
(dont certains sont étudiés, dans le cadre affine, dans (0, 21)). Pour certains compléments
sur le formalisme différentiel dans les préschémas, le lecteur pourra consulter les
exposés II et VII de [42], ainsi que les chapitres ultérieurs de ce Traité.

\subsection{Invariants normaux d'une immersion}
\EGAReaderTarget{ega4.c04.subsec.16-1}{IV.16.1}

\begin{env}[16.1.1]
\label{IV.16.1.1}
Soient $(X,\mathcal{O}_X)$, $(Y,\mathcal{O}_Y)$ deux espaces annelés,
$f=(\psi,\theta):Y\to X$ un morphisme d'espaces annelés (0\textsubscript{I}, 4.1.1)
tel que l'homomorphisme
\[
  \theta^\#:\psi^*(\mathcal{O}_X)\longrightarrow\mathcal{O}_Y
\]
soit surjectif, de sorte que $\mathcal{O}_Y$ s'identifie à un faisceau d'anneaux quotient
$\psi^*(\mathcal{O}_X)/\mathscr{I}_f$. On peut alors munir
$\psi^*(\mathcal{O}_X)$ de la filtration $\mathscr{I}_f$-préadique.
\end{env}

\begin{definition}[16.1.2]
\EGAReaderTarget{ega4.c04.def.16-1-2}{IV.16.1.2}
Le faisceau d'anneaux $\mathcal{O}_Y$-augmenté
$\psi^*(\mathcal{O}_X)/\mathscr{I}_f^{n+1}$ est appelé le $n$-ème invariant normal
de $f$; l'espace annelé
$(Y,\psi^*(\mathcal{O}_X)/\mathscr{I}_f^{n+1})$ est appelé le $n$-ème voisinage
infinitésimal de $Y$ pour le morphisme $f$, et noté $Y_f^{(n)}$ ou simplement
$Y^{(n)}$. Le faisceau d'anneaux gradués associé au faisceau d'anneaux filtrés
$\psi^*(\mathcal{O}_X)$
\[
  \mathscr{G}r_\bullet(f)=\bigoplus_{n\geq 0}
  (\mathscr{I}_f^n/\mathscr{I}_f^{n+1})
  \tag{16.1.2.1}
  \label{IV.16.1.2.1}
\]
est appelé le faisceau d'anneaux gradués associé à $f$. Le faisceau
$\mathscr{G}r_1(f)=\mathscr{I}_f/\mathscr{I}_f^2$ est appelé le faisceau conormal
de $f$ (que l'on note aussi $\mathscr{N}_{Y/X}$ s'il n'en résulte pas de confusion).
\end{definition}

Il est clair que les
$\mathcal{O}_{Y^{(n)}}=\psi^*(\mathcal{O}_X)/\mathscr{I}_f^{n+1}$
(qu'on note aussi $\mathcal{O}_{Y_f^{(n)}}$) forment un
\oldpage[IV-4]{6}
système projectif de faisceaux d'anneaux sur $Y$, l'homomorphisme de transition
$\varphi_{nm}:\mathcal{O}_{Y^{(m)}}\to\mathcal{O}_{Y^{(n)}}$ pour $n\leq m$
identifiant $\mathcal{O}_{Y^{(n)}}$ au quotient de $\mathcal{O}_{Y^{(m)}}$ par la
puissance $(\mathscr{I}_f/\mathscr{I}_f^{n+1})^m$ de l'Idéal d'augmentation de
$\mathcal{O}_{Y^{(n)}}$, noyau de
$\varphi_{0n}:\mathcal{O}_{Y^{(n)}}\to\mathcal{O}_Y$. Les $Y^{(n)}$ forment par
suite un système inductif d'espaces annelés, ayant tous pour espace sous-jacent
l'espace $Y$, et on a des morphismes canoniques d'espaces annelés
$h_n:Y^{(n)}\to X$ égaux à $(\psi,\theta_n)$, où $\theta_n^\#$ est le morphisme
canonique
$\psi^*(\mathcal{O}_X)\to\psi^*(\mathcal{O}_X)/\mathscr{I}_f^{n+1}$.
Il est clair que le faisceau $\mathscr{G}r_\bullet(f)$ est un faisceau d'algèbres
graduées sur le faisceau d'anneaux
$\mathcal{O}_Y=\mathscr{G}r_0(f)$, et les $\mathscr{G}r_k(f)$ des
$\mathcal{O}_Y$-Modules.

Comme pour tout faisceau d'anneaux filtré, on a un homomorphisme canonique
surjectif de $\mathcal{O}_Y$-Algèbres graduées
\[
  \mathscr{S}_{\mathcal{O}_Y}^{\bullet}(\mathscr{G}r_1(f))
  \longrightarrow\mathscr{G}r_\bullet(f)
  \tag{16.1.2.2}
  \label{IV.16.1.2.2}
\]
coïncidant en degrés 0 et 1 avec les homomorphismes identiques.

\begin{examples}[16.1.3]
\label{IV.16.1.3}
\begin{enumerate}
\item[(i)] Supposons que $X$ soit un espace annelé en anneaux locaux, que $Y$ soit
réduit à un seul point $y$ (muni d'un anneau $\mathcal{O}_y$) et que, si
$x=\psi(y)$, $\theta^\#:\mathcal{O}_x\to\mathcal{O}_y$ soit un homomorphisme
surjectif d'anneaux ayant pour noyau l'idéal maximal $\mathfrak{m}_x$ de
$\mathcal{O}_x$. Alors les $\mathcal{O}_{Y^{(n)}}$ s'identifient aux anneaux
$\mathcal{O}_x/\mathfrak{m}_x^{n+1}$, et $\mathscr{G}r_\bullet(f)$ à l'anneau
gradué associé à l'anneau local $\mathcal{O}_x$ muni de sa filtration
$\mathfrak{m}_x$-préadique.
\item[(ii)] Supposons que $Y$ soit une partie fermée d'un sous-espace ouvert $U$
de $X$ et que $\mathcal{O}_Y$ soit induit sur $Y$ par un faisceau quotient
$\mathcal{O}_U/\mathscr{I}$, où $\mathscr{I}$ est un Idéal de $\mathcal{O}_U$
tel que $\mathscr{I}_x=\mathcal{O}_x$ pour tout $x\notin Y$; si $X$ est un
espace annelé en anneaux locaux, nous supposerons de plus que
$\mathscr{I}_x\neq\mathcal{O}_x$ pour $x\in Y$, de sorte que
$(Y,\mathcal{O}_Y)$ est encore un espace annelé en anneaux locaux.

Soit $\psi_0:Y\to U$ l'injection canonique, et notons
$\theta_0:\mathcal{O}_U\to(\psi_0)_*(\mathcal{O}_Y)$ l'homomorphisme tel que
$\theta_0^\#$ soit l'homomorphisme canonique
$\psi_0^*(\mathcal{O}_U)=\mathcal{O}_U|Y\to
(\mathcal{O}_U/\mathscr{I})|Y$, de sorte que
$j_0=(\psi_0,\theta_0):Y\to U$ est un morphisme d'espaces annelés (et d'espaces
annelés en anneaux locaux si $X$ est un espace annelé en anneaux locaux); si
$i:U\to X$ est l'injection canonique (morphisme d'espaces annelés),
$j=i\circ j_0$ est le morphisme $(\psi,\theta)$ de $Y$ dans $X$, où
$\psi:Y\to X$ est l'injection canonique, et
$\theta:\mathcal{O}_X\to\psi_*(\mathcal{O}_Y)$ est l'homomorphisme tel que
$\theta^\#=\theta_0^\#$. Comme $\theta^\#$ est surjectif, on peut appliquer les
définitions précédentes; $\mathcal{O}_{Y^{(n)}}$ est égal à
$\psi_0^*(\mathcal{O}_U/\mathscr{I}^{n+1})$, et l'on a
$(\psi_0)_*(\mathcal{O}_{Y^{(n)}})=\mathcal{O}_U/\mathscr{I}^{n+1}$, et
$\mathscr{G}r_n(j)=\mathscr{G}r_n(j_0)=
\psi_0^*(\mathscr{I}^n/\mathscr{I}^{n+1})=
j_0^*(\mathscr{I}^n/\mathscr{I}^{n+1})$.
\end{enumerate}
\end{examples}

\begin{env}[16.1.4]
\label{IV.16.1.4}
L'exemple (16.1.3, (ii)) montre qu'en général les
$\mathcal{O}_{Y^{(n)}}$ ne sont pas munis canoniquement d'une structure de
$\mathcal{O}_Y$-Module, ni a fortiori d'une structure de
$\mathcal{O}_Y$-Algèbre. La donnée d'une telle structure équivaut à celle d'un
homomorphisme de faisceaux d'anneaux
$\lambda_n:\mathcal{O}_Y\to\mathcal{O}_{Y^{(n)}}$, inverse à droite de
l'homomorphisme d'augmentation $\varphi_{0n}$; cela revient aussi à la donnée
d'un morphisme d'espaces annelés
$(1_Y,\lambda_n):Y^{(n)}\to Y$, inverse à gauche du morphisme canonique
$(1_Y,\varphi_{0n}):Y\to Y^{(n)}$.
\end{env}

\begin{proposition}[16.1.5]
\EGAReaderTarget{ega4.c04.prop.16-1-5}{IV.16.1.5}
Soit $f=(\psi,\theta):Y\to X$ une immersion de préschémas. Alors:
\begin{enumerate}
\item[(i)] $\mathscr{G}r_\bullet(f)$ est une $\mathcal{O}_Y$-Algèbre graduée
quasi-cohérente.
\oldpage[IV-4]{7}
\item[(ii)] Les $Y^{(n)}$ sont des préschémas, canoniquement isomorphes à des
sous-préschémas de $X$.
\item[(iii)] Tout homomorphisme de faisceaux d'anneaux
$\lambda_n:\mathcal{O}_Y\to\mathcal{O}_{Y^{(n)}}$, inverse à droite de
l'homomorphisme d'augmentation $\varphi_{0n}$, fait de
$\mathcal{O}_{Y^{(n)}}$ et des $\mathcal{O}_{Y^{(k)}}$ pour $k\leq n$, des
$\mathcal{O}_Y$-Algèbres quasi-cohérentes; les structures de
$\mathcal{O}_Y$-Module déduites des structures précédentes sur les
$\mathscr{G}r_k(f)$ pour $k\leq n$ coïncident avec celles définies dans
(16.1.2).
\end{enumerate}
\end{proposition}

\noindent(i) La question étant locale sur $X$ et sur $Y$, on peut se borner au cas
où $Y$ est un sous-préschéma fermé de $X$ défini par un Idéal quasi-cohérent
$\mathscr{I}$ de $\mathcal{O}_X$; comme $\mathcal{O}_Y$ est la restriction à
$Y$ de $\mathcal{O}_X/\mathscr{I}$, l'assertion (i) est évidente, et $Y^{(n)}$
est le sous-préschéma fermé de $X$ défini par l'Idéal quasi-cohérent
$\mathscr{I}^{n+1}$ de $\mathcal{O}_X$. Enfin, pour prouver (iii), notons que la
donnée de $\lambda_n$ fait de l'Idéal
$\mathscr{I}/\mathscr{I}^{n+1}$ de l'augmentation $\varphi_{0n}$ et de ses
quotients $\mathscr{I}/\mathscr{I}^{k+1}$ ($1\leq k\leq n$) des
$\mathcal{O}_Y$-Modules, et il suffit de prouver par récurrence sur $k$ que les
$\mathscr{I}/\mathscr{I}^{k+1}$ sont des $\mathcal{O}_Y$-Modules
quasi-cohérents et que la structure de $\mathcal{O}_Y$-Module quotient qu'on en
déduit sur $\mathscr{I}^k/\mathscr{I}^{k+1}$ est la même que celle définie en
(16.1.2). La seconde assertion est immédiate,
$\mathscr{I}^k/\mathscr{I}^{k+1}$ étant annulé par
$\mathscr{I}/\mathscr{I}^{n+1}$; la première résulte, par récurrence sur $k$,
de ce qu'elle est triviale pour $k=1$ et de ce que
$\mathscr{I}/\mathscr{I}^{k+1}$ est une extension de
$\mathscr{I}/\mathscr{I}^k$ par
$\mathscr{I}^k/\mathscr{I}^{k+1}$ (III, 1.4.17).

\begin{corollary}[16.1.6]
\label{IV.16.1.6}
Sous les hypothèses générales de (16.1.5), si l'immersion $f$ est localement de
présentation finie, les $\mathscr{G}r_n(f)$ sont des $\mathcal{O}_Y$-Modules
quasi-cohérents de type fini.
\end{corollary}

En effet, avec les notations de la démonstration de (16.1.5), $\mathscr{I}$ est un
Idéal de type fini de $\mathcal{O}_X$ (1.4.7), donc les
$\mathscr{I}^n/\mathscr{I}^{n+1}$ sont des $\mathcal{O}_Y$-Modules de type fini,
d'où la conclusion.

\begin{corollary}[16.1.7]
\label{IV.16.1.7}
Sous les hypothèses générales de (16.1.5), soit $g:X\to Y$ un morphisme de
préschémas inverse à gauche de $f$. Alors, pour tout $n$, le morphisme composé
$(1,\lambda_n):Y^{(n)}\xrightarrow{h_n}X\xrightarrow{g}Y$ définit un
homomorphisme de faisceaux d'anneaux
$\lambda_n:\mathcal{O}_Y\to\mathcal{O}_{Y^{(n)}}$ inverse à droite de
l'augmentation $\varphi_{0n}$, faisant de $\mathcal{O}_{Y^{(n)}}$ une
$\mathcal{O}_Y$-Algèbre quasi-cohérente; pour ces homomorphismes, les
homomorphismes de transition
$\varphi_{nm}:\mathcal{O}_{Y^{(m)}}\to\mathcal{O}_{Y^{(n)}}$ ($n\leq m$) sont
des homomorphismes de $\mathcal{O}_Y$-Algèbres. En outre, si $g$ est localement
de type fini, les $\mathcal{O}_{Y^{(n)}}$ sont des $\mathcal{O}_Y$-Modules
quasi-cohérents de type fini.
\end{corollary}

La première assertion résulte aussitôt des définitions et de (16.1.5). D'autre
part, si $g$ est localement de type fini, alors $f$ est localement de
présentation finie (1.4.3, (v)); les $\mathscr{G}r_n(f)$ étant alors des
$\mathcal{O}_Y$-Modules quasi-cohérents de type fini par (16.1.6), il en est de
même des $\mathcal{O}_Y$-Modules
$\mathscr{I}/\mathscr{I}^{n+1}$, qui sont des extensions d'un nombre fini de
$\mathscr{G}r_k(f)$ (III, 1.4.17).

\begin{proposition}[16.1.8]
\label{IV.16.1.8}
Soient $X$ un préschéma localement noethérien, $j:Y\to X$ une immersion. Alors
les $Y^{(n)}$ sont des préschémas localement noethériens, les
$\mathscr{G}r_n(j)$ sont des $\mathcal{O}_Y$-Modules cohérents, et
$\mathscr{G}r_\bullet(j)$ est un faisceau d'anneaux cohérent sur l'espace $Y$.
\end{proposition}

Tout étant local sur $X$ et $Y$, on est ramené au cas où $X$ est affine et $j$
une immersion fermée, et alors toutes les assertions sont évidentes sauf la
dernière, qui résulte de ce que si $A$ est un anneau noethérien et
$\mathfrak{I}$ un idéal de $A$, $\operatorname{gr}_{\mathfrak{I}}^\bullet(A)$
est un anneau noethérien, compte tenu de l'exactitude du foncteur $\psi^*$ et de
(0\textsubscript{I}, 5.3.7).

\begin{proposition}[16.1.9]
\label{IV.16.1.9}
\oldpage[IV-4]{8}
Soient $X$ un préschéma, $j:Y\to X$ une immersion localement de présentation
finie, $y$ un point de $Y$. Les conditions suivantes sont équivalentes:
\begin{enumerate}
\item[(a)] Il existe un voisinage ouvert $U$ de $y$ dans $Y$ tel que $j|U$ soit
un homéomorphisme de $U$ sur un ouvert de $X$.
\item[(b)] Il existe un entier $n>0$ tel que l'homomorphisme canonique
\[
  (\varphi_{n,n-1})_y:
  \mathcal{O}_{Y^{(n)},y}\longrightarrow\mathcal{O}_{Y^{(n-1)},y}
\]
soit bijectif.
\item[(c)] Il existe un entier $n>0$ tel que
$(\mathscr{G}r_n(j))_y=0$.
\end{enumerate}
De plus, si l'entier $n$ vérifie b) ou c), il existe un voisinage $V$ de $y$
dans $Y$ tel que $\mathscr{G}r_m(j)|V=0$ pour $m\geq n$ et que
$\varphi_{nm}|V:\mathcal{O}_{Y^{(m)}}|V\to\mathcal{O}_{Y^{(n)}}|V$ soit
bijectif pour $m\geq n$.
\end{proposition}

La question étant locale sur $Y$, on peut se borner au cas où $j$ est une
immersion fermée, $Y$ étant défini par un Idéal quasi-cohérent de type fini
$\mathscr{I}$ de $\mathcal{O}_X$. L'équivalence de b) et c), pour un $n$ donné,
est alors immédiate; en outre, comme
$\mathscr{I}^n/\mathscr{I}^{n+1}$ est un $\mathcal{O}_X$-Module de type fini,
il existe un voisinage ouvert $U$ de $y$ dans $X$ tel que
$\mathscr{I}^n|U=\mathscr{I}^{n+1}|U$ (0\textsubscript{I}, 5.2.2), donc aussi
$\mathscr{I}^n|U=\mathscr{I}^m|U$ pour $m\geq n$, ce qui prouve les dernières
assertions. Pour prouver que a) entraîne b), on peut se borner au cas où
l'espace sous-jacent à $Y$ est égal à l'espace sous-jacent à $X$, et où
$\mathscr{I}$ est engendré par un nombre fini de ses sections au-dessus de $X$:
comme $\mathscr{I}$ est alors contenu dans le Nilradical $\mathscr{N}$ de
$\mathcal{O}_X$ (I, 5.1.2), il est nilpotent, ce qui prouve b). Enfin, pour
prouver que b) entraîne a), on peut aussi se borner au cas où
$\mathscr{I}^n=\mathscr{I}^{n+1}$; alors, pour tout $y\in Y$, comme
$\mathscr{I}_y\subset\mathfrak{m}_y$, idéal maximal de
$\mathcal{O}_{X,y}$, on a nécessairement $\mathscr{I}_y^n=0$ en vertu du
lemme de Nakayama, puisque $\mathscr{I}_y$ est un idéal de type fini.
L'ensemble des $x\in X$ tels que $\mathscr{I}_x^n=0$ est donc un ouvert $U$ de
$X$ contenant $Y$ (0\textsubscript{I}, 5.2.2); comme d'autre part
$\mathscr{I}_x\neq 0$ pour $x\notin Y$, on a nécessairement $U=Y$.

\begin{corollary}[16.1.10]
\label{IV.16.1.10}
Pour que la restriction de l'immersion $j$ à un voisinage de $y$ dans $Y$ soit
une immersion ouverte (autrement dit pour que $j$ soit un isomorphisme local au
point $y$), il faut et il suffit que
$(\mathscr{G}r_1(j))_y=(\mathscr{N}_{Y/X})_y=0$.
\end{corollary}

La condition est évidemment nécessaire, et le raisonnement précédent, appliqué
pour $n=1$, prouve qu'elle est suffisante.

\begin{remarks}[16.1.11]
\EGAReaderTarget{ega4.c04.rem.16-1-11}{IV.16.1.11}
\begin{enumerate}
\item[(i)] Sous les conditions de la définition (16.1.1), la limite projective
du système projectif
$(\mathcal{O}_{Y^{(n)}},\varphi_{nm})$ de faisceaux d'anneaux sur $Y$ est
appelé l'invariant normal d'ordre infini de $f$, et parfois noté
$\mathcal{O}_{Y^{(\infty)}}$. Lorsque $X$ est un préschéma localement
noethérien, $j:Y\to X$ une immersion fermée, $Y$ étant donc un sous-préschéma
fermé de $X$ défini par un Idéal cohérent $\mathscr{I}$,
$\mathcal{O}_{Y^{(\infty)}}$ n'est autre que le complété formel de
$\mathcal{O}_X$ le long de $Y$ (I, 10.8.4), et
$Y^{(\infty)}=(Y,\mathcal{O}_{Y^{(\infty)}})$ le préschéma formel complété de
$X$ le long de $Y$ (I, 10.8.5). Dans tous les cas, on pourra dire que
$Y^{(\infty)}$ est le voisinage formel de $Y$ dans $X$ (pour le morphisme
$f$). Dans le cas particulier qu'on vient d'envisager, c'est donc le préschéma
formel limite inductive des voisinages infinitésimaux d'ordre $n$.
\item[(ii)] On notera que pour un morphisme de préschémas
$f=(\psi,\theta):Y\to X$, il peut se faire que l'homomorphisme
$\theta^\#:\psi^*(\mathcal{O}_X)\to\mathcal{O}_Y$ soit surjectif sans que $f$
soit une immersion
\oldpage[IV-4]{9}
locale, et sans que $f$ soit injectif. On en a un exemple en prenant pour $Y$
une somme de préschémas $Y_\lambda$ tous isomorphes à
$\operatorname{Spec}(\mathcal{O}_x)$, où $x\in X$, et pour $f$ le morphisme
égal au morphisme canonique dans chacun des $Y_\lambda$.
\end{enumerate}
\end{remarks}

\subsection{Propriétés fonctorielles des invariants normaux d'une immersion}
\EGAReaderTarget{ega4.c04.subsec.16-2}{IV.16.2}

\begin{env}[16.2.1]
\label{IV.16.2.1}
Soient $f=(\psi,\theta):Y\to X$,
$f'=(\psi',\theta'):Y'\to X'$ deux morphismes d'espaces annelés tels que les
homomorphismes $\theta^\#$ et $\theta'^\#$ soient surjectifs; considérons un
diagramme commutatif de morphismes d'espaces annelés
\[
  \xymatrix{
    Y \ar[r]^f & X \\
    Y' \ar[r]_{f'} \ar[u]^u & X' \ar[u]_v
  }
  \tag{16.2.1.1}
  \label{IV.16.2.1.1}
\]
Posons $u=(\rho,\lambda)$, $v=(\sigma,\mu)$. On a
$\rho^*(\psi^*(\mathcal{O}_X))=
\psi'^*(\sigma^*(\mathcal{O}_X))$ et on a par suite un diagramme commutatif
d'homomorphismes de faisceaux d'anneaux sur $Y'$
\[
  \xymatrix{
    \rho^*(\psi^*(\mathcal{O}_X))
      =\psi'^*(\sigma^*(\mathcal{O}_X))
      \ar[r]^-{\psi'^*(\mu^\#)}
      \ar[d]_{\rho^*(\theta^\#)}
      & \psi'^*(\mathcal{O}_{X'}) \ar[d]^{\theta'^\#} \\
    \rho^*(\mathcal{O}_Y) \ar[r]_-{\lambda^\#}
      & \mathcal{O}_{Y'}
  }
\]
d'où l'on conclut, si $\mathscr{I}$ et $\mathscr{I}'$ sont les noyaux de
$\theta^\#$ et $\theta'^\#$, que l'on a
$\psi'^*(\mu^\#)(\rho^*(\mathscr{I}))\subset\mathscr{I}'$, compte tenu de
l'exactitude du foncteur $\rho^*$. On en déduit aussitôt que pour tout entier
$n$,
$\psi'^*(\mu^\#)(\rho^*(\mathscr{I}^n))\subset\mathscr{I}'^n$, ce qui montre
que $\psi'^*(\mu^\#)$ définit, par passage aux quotients, un homomorphisme de
faisceaux d'anneaux
\[
  \nu_n:
  \rho^*(\psi^*(\mathcal{O}_X)/\mathscr{I}^{n+1})
  \longrightarrow
  \psi'^*(\mathcal{O}_{X'})/\mathscr{I}'^{n+1}
  \tag{16.2.1.2}
  \label{IV.16.2.1.2}
\]
et par suite un morphisme d'espaces annelés
$w_n=(\rho,\nu_n):Y'^{(n)}\to Y^{(n)}$ (qui, pour $n=0$, n'est autre que
$u$). Il résulte aussitôt de cette définition que les diagrammes
\[
  \xymatrix@R=1pc{
    Y^{(n)} \ar[r]^-{h_{mn}} & Y^{(m)} \ar[r]^-{h_m} & X \\
    & & & (n\leq m) \\
    Y'^{(n)} \ar[r]_-{h'_{mn}} \ar[uu]^-{w_n}
      & Y'^{(m)} \ar[r]_-{h'_m} \ar[uu]^-{w_m}
      & X' \ar[uu]_-v
  }
\]
(où les flèches horizontales sont les morphismes canoniques (16.1.2)) sont
commutatifs.

Par passage aux quotients à partir des homomorphismes (16.2.1.2), et en tenant
\oldpage[IV-4]{10}
compte de l'exactitude du foncteur $\rho^*$, on obtient un di-homomorphisme
d'Algèbres graduées (relatif à l'homomorphisme
$\lambda^\#:\rho^*(\mathcal{O}_Y)\to\mathcal{O}_{Y'}$)
\[
  \operatorname{gr}(u):
  \rho^*(\mathscr{G}r_\bullet(f))
  \longrightarrow\mathscr{G}r_\bullet(f')
  \tag{16.2.1.3}
  \label{IV.16.2.1.3}
\]
(ou, si on veut, un $\rho$-morphisme (0\textsubscript{I}, 3.5.1)
$\mathscr{G}r_\bullet(f)\to\mathscr{G}r_\bullet(f')$), et en particulier un
di-homomorphisme des faisceaux conormaux
\[
  \operatorname{gr}_1(u):
  \rho^*(\mathscr{G}r_1(f))\longrightarrow\mathscr{G}r_1(f').
\]
Il est immédiat en outre que ces homomorphismes donnent lieu à un diagramme
commutatif
\[
  \xymatrix{
    \rho^*(\mathscr{S}_{\mathcal{O}_Y}^{\bullet}
      (\mathscr{G}r_1(f)))
      \ar[r] \ar[d]_-{\mathscr{S}(\operatorname{gr}_1(u))}
      & \rho^*(\mathscr{G}r_\bullet(f))
        \ar[d]^-{\operatorname{gr}(u)} \\
    \mathscr{S}_{\mathcal{O}_{Y'}}^{\bullet}
      (\mathscr{G}r_1(f'))
      \ar[r]
      & \mathscr{G}r_\bullet(f')
  }
  \tag{16.2.1.4}
  \label{IV.16.2.1.4}
\]
où les flèches horizontales sont les homomorphismes canoniques (16.1.2.2).

Enfin, si l'on a un diagramme commutatif de morphismes d'espaces annelés
\[
  \xymatrix{
    Y \ar[r]^f & X \\
    Y' \ar[r]_{f'} \ar[u]^u & X' \ar[u]_v \\
    Y'' \ar[r]_{f''} \ar[u]^{u'} & X'' \ar[u]_{v'}
  }
\]
où $f''=(\psi'',\theta'')$ est tel que $\theta''^\#$ soit surjectif, et si
$w'_n$ et $w''_n$ sont définis à partir de $u'$, $v'$ d'une part, et de
$u''=u\circ u'$, $v''=v\circ v'$ de l'autre, alors on a
$w''_n=w_n\circ w'_n$, comme il résulte aussitôt des définitions et de
(0\textsubscript{I}, 3.5.5); on a de même
$\operatorname{gr}(u'')=
\operatorname{gr}(u')\circ\rho'^*(\operatorname{gr}(u))$ si
$u'=(\rho',\lambda')$. On peut donc dire que les $Y^{(n)}$ et
$\mathscr{G}r_\bullet(f)$ dépendent fonctoriellement de $f$.
\end{env}

\begin{proposition}[16.2.2]
\label{IV.16.2.2}
Avec les notations et hypothèses de (16.2.1), supposons de plus que $f$, $f'$,
$u$ et $v$ soient des morphismes de préschémas. Alors:
\begin{enumerate}
\item[(i)] Les morphismes $w_n:Y'^{(n)}\to Y^{(n)}$ sont des morphismes de
préschémas.
\item[(ii)] Si $Y'=Y\times_X X'$, $u$ et $f'$ étant les projections
canoniques, et si $f$ est une immersion ou si $v$ est plat, on a
$Y'^{(n)}=Y^{(n)}\times_X X'$.
\item[(iii)] Si $Y'=Y\times_X X'$ et si $v$ est plat (resp. si $f$ est une
immersion), l'homomorphisme
\[
  \operatorname{Gr}(u)=\operatorname{gr}(u)\otimes 1:
  \mathscr{G}r_\bullet(f)\otimes_{\mathcal{O}_Y}\mathcal{O}_{Y'}
  \longrightarrow\mathscr{G}r_\bullet(f')
\]
est bijectif (resp. surjectif).
\end{enumerate}
\end{proposition}

\noindent(i) Les hypothèses entraînent aussitôt que pour tout $y'\in Y'$,
$\rho_{y'}^*(\theta_{\psi'(y')}^\#)$ est un homomorphisme local (I, 1.6.2),
donc $w_n$ est un morphisme de préschémas (I, 2.2.1).

\oldpage[IV-4]{11}
\noindent(ii) et (iii) Si $f$ est une immersion, on peut se borner au cas où
$f$ est une immersion fermée, $Y$ étant donc défini par l'Idéal
quasi-cohérent $\mathscr{I}$ de $\mathcal{O}_X$, et $Y^{(n)}$ par l'Idéal
$\mathscr{I}^{n+1}$; les assertions résultent alors de (I, 4.4.5).

Supposons en second lieu que $v$ soit plat; on peut se borner au cas où
$X=\operatorname{Spec}(A)$, $Y=\operatorname{Spec}(B)$,
$X'=\operatorname{Spec}(A')$ sont affines, $A'$ étant un $A$-module plat;
alors $Y'=\operatorname{Spec}(B')$ avec $B'=B\otimes_A A'$; en outre, si
$\mathfrak{I}$ est le noyau de l'homomorphisme $A\to B$, le noyau
$\mathfrak{I}'$ de $A'\to B'$ s'identifie à $\mathfrak{I}\otimes_A A'$ par
platitude, et
\[
  \mathfrak{I}'^n/\mathfrak{I}'^{n+1}
  =(\mathfrak{I}^n/\mathfrak{I}^{n+1})\otimes_A A'.
\]
On en déduit aussitôt, compte tenu de (0\textsubscript{I}, 4.3.3), que le
$\mathcal{O}_{Y'}$-Module
$\mathscr{I}'^n/\mathscr{I}'^{n+1}$ est égal à
\begin{align*}
&\psi'^*\!\left(\sigma^*
  \left((\mathfrak{I}^n/\mathfrak{I}^{n+1})^{\sim}\right)\right)
  \otimes_{\sigma^*(\mathcal{O}_X)}\mathcal{O}_{X'} \\
={}&
\psi'^*\!\left(\sigma^*
  \left((\mathfrak{I}^n/\mathfrak{I}^{n+1})^{\sim}\right)\right)
  \otimes_{\psi'^*(\sigma^*(\mathcal{O}_X))}
  \psi'^*(\mathcal{O}_{X'}) \\
={}&
\rho^*(\mathscr{I}^n/\mathscr{I}^{n+1})
  \otimes_{\rho^*(\psi^*(\mathcal{O}_X))}
  \psi'^*(\mathcal{O}_{X'}),
\end{align*}
et comme en particulier pour $n=0$, on a
\[
  \mathcal{O}_{Y'}=
  \rho^*(\mathcal{O}_Y)
  \otimes_{\rho^*(\psi^*(\mathcal{O}_X))}
  \psi'^*(\mathcal{O}_{X'}),
\]
on obtient un isomorphisme canonique de
$\mathscr{I}'^n/\mathscr{I}'^{n+1}$ sur
\[
  \rho^*(\mathscr{I}^n/\mathscr{I}^{n+1})
  \otimes_{\rho^*(\mathcal{O}_Y)}\mathcal{O}_{Y'}
  =
  (\mathscr{I}^n/\mathscr{I}^{n+1})
  \otimes_{\mathcal{O}_Y}\mathcal{O}_{Y'},
\]
ce qui prouve (iii). Posons maintenant
$C_n=\Gamma(Y,\mathcal{O}_{Y^{(n)}})$,
$C'_n=\Gamma(Y',\mathcal{O}_{Y'^{(n)}})$. Puisque $Y^{(n)}$ et
$Y'^{(n)}$ sont des schémas affines (16.1.5), le noyau
$\mathfrak{K}_n$ (resp. $\mathfrak{K}'_n$) de l'homomorphisme
$C_n\to C_{n-1}$ (resp. $C'_n\to C'_{n-1}$) est
$\Gamma(Y,\mathscr{I}^n/\mathscr{I}^{n+1})$ (resp.
$\Gamma(Y',\mathscr{I}'^n/\mathscr{I}'^{n+1})$); on déduit donc de ce qui
précède que $\mathfrak{K}'_n=\mathfrak{K}_n\otimes_A A'$. Or, on a un
diagramme commutatif
\[
  \xymatrix{
    0 \ar[r] &
    \mathfrak{K}_n\otimes_A A' \ar[r]\ar[d]^r &
    C_n\otimes_A A' \ar[r]\ar[d]^{s_n} &
    C_{n-1}\otimes_A A' \ar[r]\ar[d]^{s_{n-1}} & 0 \\
    0 \ar[r] &
    \mathfrak{K}'_n \ar[r] &
    C'_n \ar[r] &
    C'_{n-1} \ar[r] & 0
  }
\]
où la flèche verticale de gauche est bijective et les deux lignes sont exactes
($A'$ étant un $A$-module plat). On en déduit par récurrence que $s_n$ est
bijectif pour tout $n$, car c'est vrai par hypothèse pour $n=0$, et se déduit
par application du lemme des cinq pour $n$ quelconque. Cela prouve la seconde
assertion de (ii).

\begin{corollary}[16.2.3]
\label{IV.16.2.3}
Soient $g:X\to Y$, $u:Y'\to Y$ deux morphismes de préschémas,
$X'=X\times_Y Y'$, $g':X'\to Y'$ et $v:X'\to X$ les projections canoniques.
Soient $f:Y\to X$ une $Y$-section de $X$ (donc une immersion),
$f'=f_{(Y')}:Y'\to X'$ la $Y'$-section de $X'$ déduite de $f$ par le
changement de base $u$. Alors:
\begin{enumerate}
\item[(i)] Le morphisme
$w_n:Y_{f'}'^{(n)}\to Y_f^{(n)}$ correspondant à $f$, $f'$, $u$, $v$
(16.2.1) et le morphisme canonique
$h'_n:Y_{f'}'^{(n)}\to X'$ identifient $Y_{f'}'^{(n)}$ au produit
$Y_f^{(n)}\times_X X'$.
\item[(ii)] Si l'on munit $\mathcal{O}_{Y_f^{(n)}}$ (resp.
$\mathcal{O}_{Y_{f'}'^{(n)}}$) de la structure de $\mathcal{O}_Y$-Algèbre
définie par $g$ (resp. de la structure de $\mathcal{O}_{Y'}$-Algèbre définie
par $g'$) (16.1.6), l'homomorphisme de $\mathcal{O}_{Y'}$-Algèbres
\[
  \rho^*(\mathcal{O}_{Y_f^{(n)}})
  \otimes_{\mathcal{O}_Y}\mathcal{O}_{Y'}
  \longrightarrow
  \mathcal{O}_{Y_{f'}'^{(n)}}
  \tag{16.2.3.1}
  \label{IV.16.2.3.1}
\]
\oldpage[IV-4]{12}
déduit de l'homomorphisme $\nu_n$ (16.2.1.2) est bijectif. En outre,
l'homomorphisme de $\mathcal{O}_{Y'}$-Modules
\[
  \operatorname{Gr}_1(u):
  \mathscr{G}r_1(f)\otimes_{\mathcal{O}_Y}\mathcal{O}_{Y'}
  \longrightarrow\mathscr{G}r_1(f')
  \tag{16.2.3.2}
  \label{IV.16.2.3.2}
\]
est bijectif.
\end{enumerate}
\end{corollary}

\noindent(i) Notons d'abord que les morphismes $f':Y'\to X'$ et $u:Y'\to Y$
identifient $Y'$ au produit $Y\times_X X'$ (pour les morphismes structuraux
$f:Y\to X$ et $v:X'\to X$) (14.5.12.1). La conclusion de (i) résulte alors de
(16.2.2, (ii)), le morphisme $g$ étant une immersion.

\noindent(ii) Le diagramme commutatif
\[
  \xymatrix{
    Y_f^{(n)} \ar[d]^{h_n}
      & Y_{f'}'^{(n)} \ar[d]^-{h'_n} \ar[l]^{w_n} \\
    X \ar[d]^g
      & X' \ar[d]^{g'} \ar[l]^v \\
    Y
      & Y' \ar[l]^u
  }
\]
identifie $Y_{f'}'^{(n)}$ au produit $Y_f^{(n)}\times_X X'$ et $X'$ au produit
$X\times_Y Y'$, donc (I, 3.3.9) il identifie (pour les morphismes
$g'\circ h'_n$ et $w_n$) $Y_{f'}'^{(n)}$ au produit
$Y_f^{(n)}\times_Y Y'$. Comme $Y_f^{(n)}$ (resp. $Y_{f'}'^{(n)}$) est le
préschéma affine au-dessus de $Y$ (resp. $Y'$) associé à la
$\mathcal{O}_Y$-Algèbre $\mathcal{O}_{Y_f^{(n)}}$ (resp. à la
$\mathcal{O}_{Y'}$-Algèbre $\mathcal{O}_{Y_{f'}'^{(n)}}$), le fait que
l'homomorphisme canonique (16.2.3.1) soit bijectif résulte de (II, 1.5.2).
Enfin, l'homomorphisme canonique (16.2.3.1) est compatible avec les
augmentations
$\mathcal{O}_{Y_f^{(n)}}\to\mathcal{O}_Y$ et
$\mathcal{O}_{Y_{f'}'^{(n)}}\to\mathcal{O}_{Y'}$; comme
$\mathcal{O}_{Y_f^{(n)}}$ est somme directe (en tant que
$\mathcal{O}_Y$-Module) de $\mathcal{O}_Y$ et de l'Idéal d'augmentation
$\mathscr{I}/\mathscr{I}^{n+1}$, on voit donc que l'homomorphisme canonique
(16.2.3.1), restreint à
$(\mathscr{I}/\mathscr{I}^{n+1})\otimes_{\mathcal{O}_Y}\mathcal{O}_{Y'}$,
est une bijection de ce dernier sur
$\mathscr{I}'/\mathscr{I}'^{n+1}$. Pour $n=1$, cela montre que
$\operatorname{Gr}_1(u)$ est bijectif.

On notera que, sous les hypothèses de (16.2.3), les homomorphismes
$\operatorname{Gr}_n(u)$ sont surjectifs en vertu de ce qui précède, mais ne
sont pas bijectifs en général pour $n\geq 2$. Toutefois:

\begin{corollary}[16.2.4]
\label{IV.16.2.4}
Sous les hypothèses de (16.2.3), supposons que $u:Y'\to Y$ soit un morphisme
plat (resp. que les $\mathscr{G}r_n(f)$ soient des $\mathcal{O}_Y$-Modules
plats pour $n\leq m$). Alors l'homomorphisme
\[
  \operatorname{Gr}_n(u):
  \mathscr{G}r_n(f)\otimes_{\mathcal{O}_Y}\mathcal{O}_{Y'}
  \longrightarrow\mathscr{G}r_n(f')
\]
est bijectif pour tout $n$ (resp. pour $n\leq m$).
\end{corollary}

Si $u$ est plat, il en est de même de $v:X'\to X$ qui s'en déduit par
changement de base, et on sait déjà dans ce cas que $\operatorname{Gr}(u)$ est
bijectif (16.2.2, (iii)). Si les $\mathscr{G}r_n(f)$ sont plats pour
$n\leq m$, on voit d'abord par récurrence sur $n$ qu'il en est de même des
$\mathscr{I}/\mathscr{I}^{n+1}$ pour $n\leq m$, en vertu des suites exactes
\[
  0\longrightarrow\mathscr{I}^n/\mathscr{I}^{n+1}
  \longrightarrow\mathscr{I}/\mathscr{I}^{n+1}
  \longrightarrow\mathscr{I}/\mathscr{I}^n
  \longrightarrow 0
\]
\oldpage[IV-4]{13}
(0\textsubscript{I}, 6.1.2); en outre, on a alors des diagrammes commutatifs
\[
  \xymatrix{
    0 \ar[r] &
    (\mathscr{I}^n/\mathscr{I}^{n+1})
      \otimes_{\mathcal{O}_Y}\mathcal{O}_{Y'}
      \ar[r]\ar[d] &
    (\mathscr{I}/\mathscr{I}^{n+1})
      \otimes_{\mathcal{O}_Y}\mathcal{O}_{Y'}
      \ar[r]\ar[d] &
    (\mathscr{I}/\mathscr{I}^n)
      \otimes_{\mathcal{O}_Y}\mathcal{O}_{Y'}
      \ar[r]\ar[d] & 0 \\
    0 \ar[r] &
    \mathscr{I}'^n/\mathscr{I}'^{n+1}
      \ar[r] &
    \mathscr{I}'/\mathscr{I}'^{n+1}
      \ar[r] &
    \mathscr{I}'/\mathscr{I}'^n
      \ar[r] & 0
  }
\]
dans lesquels les lignes sont exactes (la première par platitude
(0\textsubscript{I}, 6.1.2)) et les deux dernières flèches verticales sont
bijectives en vertu de (16.2.3, (ii)); d'où la conclusion.

\begin{remarks}[16.2.5]
\label{IV.16.2.5}
\begin{enumerate}
\item[(i)] Le raisonnement de (16.2.2, (i)) s'applique encore lorsque dans
(16.2.1.1) il s'agit de morphismes d'espaces annelés en anneaux locaux
($\mathbf{Err}_{\mathrm{II}}$, (1.8.2)).
\item[(ii)] Dans (16.2.2, (ii)), la conclusion n'est plus nécessairement valable
lorsqu'on suppose seulement que $v$ et $f$ sont des morphismes de préschémas
($f$ vérifiant la condition de (16.1.1)). Par exemple (avec les notations de la
démonstration de (16.2.2, (ii)), il peut se faire que $\mathfrak{I}=0$ mais que
le noyau $\mathfrak{I}'$ de $A'\to B'=B\otimes_A A'$ ne soit pas nul et que
$B'\neq 0$, auquel cas on a $Y^{(n)}=Y$ pour tout $n$, mais
$Y'^{(n)}\neq Y'$. On a un exemple de ce fait en prenant
$A=\mathbb{Z}$, $B=\mathbb{Q}$,
$A'=\prod_{h=1}^{\infty}(\mathbb{Z}/m^h\mathbb{Z})$ où $m>1$.
\end{enumerate}
\end{remarks}

\begin{env}[16.2.6]
\label{IV.16.2.6}
Considérons le cas particulier du diagramme (16.2.1.1) où $X'=X$, $v$ étant
l'identité, $X$ est un préschéma, $Y$ un sous-préschéma de $X$, $Y'$ un
sous-préschéma de $Y$, $f$, $u$ et $f'=f\circ u$ les injections canoniques; le
di-homomorphisme (16.2.1.3) donne, par tensorisation avec
$\mathcal{O}_{Y'}$ au-dessus de $\rho^*(\mathcal{O}_Y)$, un homomorphisme de
$\mathcal{O}_{Y'}$-Algèbres graduées
\[
  u^*(\mathscr{G}r_\bullet(f))
  \longrightarrow\mathscr{G}r_\bullet(f')
  \tag{16.2.6.1}
  \label{IV.16.2.6.1}
\]
D'autre part, on identifie $\mathcal{O}_Y$ à
$\psi^*(\mathcal{O}_X)/\mathscr{I}_f$, $\mathcal{O}_{Y'}$ à
$\rho^*(\mathcal{O}_Y)/\mathscr{I}_u$; comme $\rho^*$ est un foncteur exact,
on a
$\rho^*(\mathcal{O}_Y)=
\rho^*(\psi^*(\mathcal{O}_X))/\rho^*(\mathscr{I}_f)=
\psi'^*(\mathcal{O}_X)/\rho^*(\mathscr{I}_f)$, et comme
$\mathcal{O}_{Y'}$ s'identifie par ailleurs à
$\psi'^*(\mathcal{O}_X)/\mathscr{I}_{f'}$, on voit que l'on a
$\mathscr{I}_u=\mathscr{I}_{f'}/\rho^*(\mathscr{I}_f)$. On en déduit pour
tout entier $n$ un homomorphisme canonique
$\mathscr{I}_{f'}^n/\mathscr{I}_{f'}^{n+1}\to
\mathscr{I}_u^n/\mathscr{I}_u^{n+1}$, d'où un homomorphisme canonique de
$\mathcal{O}_{Y'}$-Algèbres graduées
\[
  \mathscr{G}r_\bullet(f')
  \longrightarrow\mathscr{G}r_\bullet(u).
  \tag{16.2.6.2}
  \label{IV.16.2.6.2}
\]
\end{env}

\begin{proposition}[16.2.7]
\label{IV.16.2.7}
Soient $X$ un préschéma, $Y$ un sous-préschéma de $X$, $Y'$ un
sous-préschéma de $Y$, $j:Y'\to Y$ l'injection canonique. On a alors une suite
exacte de faisceaux conormaux ($\mathcal{O}_{Y'}$-Modules)
\[
  \xymatrix{
    j^*(\mathscr{N}_{Y/X}) \ar[r] &
    \mathscr{N}_{Y'/X} \ar[r] &
    \mathscr{N}_{Y'/Y} \ar[r] & 0
  }
  \tag{16.2.7.1}
  \label{IV.16.2.7.1}
\]
où les flèches sont les composantes de degré 1 des homomorphismes canoniques
(16.2.6.1) et (16.2.6.2).
\end{proposition}

La question étant locale, on peut se borner au cas où
$X=\operatorname{Spec}(A)$, $Y=\operatorname{Spec}(A/\mathfrak{I})$,
$Y'=\operatorname{Spec}(A/\mathfrak{K})$, $\mathfrak{I}$ et $\mathfrak{K}$
étant deux idéaux de $A$ tels que $\mathfrak{I}\subset\mathfrak{K}$; tout
revient à voir alors
\oldpage[IV-4]{14}
que la suite d'homomorphismes canoniques
$\mathfrak{I}/\mathfrak{K}\mathfrak{I}\to
\mathfrak{K}/\mathfrak{K}^2\to
(\mathfrak{K}/\mathfrak{I})/(\mathfrak{K}/\mathfrak{I})^2\to 0$
est exacte, ce qui est immédiat puisque l'image de
$\mathfrak{I}/\mathfrak{K}\mathfrak{I}$ dans
$\mathfrak{K}/\mathfrak{K}^2$ est
$(\mathfrak{I}+\mathfrak{K}^2)/\mathfrak{K}^2$ et que
$(\mathfrak{K}/\mathfrak{I})/(\mathfrak{K}/\mathfrak{I})^2$ s'identifie à
$\mathfrak{K}/(\mathfrak{I}+\mathfrak{K}^2)$.

Il est facile de donner des exemples où la suite (16.2.7.1) prolongée à gauche
par un 0 n'est plus exacte; avec les notations précédentes, il suffit de
prendre $A=k[T]$, $\mathfrak{I}=AT^2$, $\mathfrak{K}=AT$, car on a alors
$(\mathfrak{I}+\mathfrak{K}^2)/\mathfrak{K}^2=0$ et
$\mathfrak{I}/\mathfrak{K}\mathfrak{I}\neq 0$. Voir cependant (16.9.13) et
(19.1.5) pour des cas utiles où la suite prolongée reste encore exacte.

\subsection{Invariants différentiels fondamentaux d'un morphisme de préschémas}
\EGAReaderTarget{ega4.c04.subsec.16-3}{IV.16.3}

\begin{definition}[16.3.1]
\EGAReaderTarget{ega4.c04.def.16-3-1}{IV.16.3.1}
Soient $f:X\to S$ un morphisme de préschémas, $\Delta_f:X\to X\times_S X$
le morphisme diagonal correspondant, qui est une immersion (I, 5.3.9 et
$\mathbf{Err}_{\mathrm{III}}$, 10). On désigne par $\mathscr{P}_f^n$ ou
$\mathscr{P}_{X/S}^n$, et l'on appelle faisceau des parties principales d'ordre $n$
du $S$-préschéma $X$, le faisceau d'anneaux $\mathcal{O}_X$-augmentés, $n$-ème
invariant normal de $\Delta_f$ (16.1.2). On pose
$\mathscr{P}_f^\infty=\mathscr{P}_{X/S}^\infty=
\varprojlim_n\mathscr{P}_{X/S}^n$,
$\mathscr{G}r_n(\mathscr{P}_f)=\mathscr{G}r_n(\mathscr{P}_{X/S})=
\mathscr{G}r_n(\Delta_f)$ (16.1.2); le $\mathcal{O}_X$-Module
$\mathscr{G}r_1(\Delta_f)$, Idéal d'augmentation de $\mathscr{P}_{X/S}^1$, est
aussi noté $\Omega_f^1$ ou $\Omega_{X/S}^1$, et appelé le $\mathcal{O}_X$-Module
des 1-différentielles de $f$, ou de $X$ par rapport à $S$, ou du $S$-préschéma $X$.
\end{definition}

Il résulte de cette définition que $\mathscr{P}_{X/S}^0$ s'identifie
canoniquement à $\mathcal{O}_X$ (16.1.2).

On a (16.1.2.2) un homomorphisme canonique surjectif de
$\mathcal{O}_X$-Algèbres graduées
\[
  \mathscr{S}_{\mathcal{O}_X}^{\bullet}(\Omega_{X/S}^1)
  \longrightarrow\mathscr{G}r_\bullet(\mathscr{P}_{X/S}).
  \tag{16.3.1.1}
  \label{IV.16.3.1.1}
\]
Il résulte aussi de la définition (16.3.1) que pour tout ouvert $U$ de $X$, on a
$\mathscr{P}_{f|U}^n=\mathscr{P}_f^n|U$,
$\mathscr{P}_{f|U}^\infty=\mathscr{P}_f^\infty|U$,
$\mathscr{G}r_n(\mathscr{P}_{f|U})=\mathscr{G}r_n(\mathscr{P}_f)|U$,
$\Omega_{f|U}^1=\Omega_f^1|U$ (autrement dit, les notions introduites sont
locales sur $X$).

\begin{env}[16.3.2]
\label{IV.16.3.2}
Notons $p_1$, $p_2$ les deux projections canoniques du produit $X\times_S X$;
comme $\Delta_f$ est une $X$-section de $X\times_S X$ à la fois pour les
morphismes $p_1$ et $p_2$, chacun de ces morphismes définit, pour tout $n$, un
homomorphisme de faisceaux d'anneaux $\mathcal{O}_X\to\mathscr{P}_{X/S}^n$,
inverse à droite de l'augmentation $\mathscr{P}_{X/S}^n\to\mathcal{O}_X$
(16.1.7); on peut encore dire que l'on définit ainsi sur
$\mathscr{P}_{X/S}^n$ deux structures de $\mathcal{O}_X$-Algèbres augmentées
quasi-cohérentes; les structures correspondantes de $\mathcal{O}_X$-Module sur
$\mathscr{G}r_n(\mathscr{P}_{X/S})$ sont les mêmes. On a de même, par passage à
la limite, deux structures de $\mathcal{O}_X$-Algèbre sur
$\mathscr{P}_{X/S}^\infty$.
\end{env}

\begin{env}[16.3.3]
\EGAReaderTarget{ega4.c04.statement.16-3-3}{IV.16.3.3}
Le morphisme $s=(p_2,p_1)_S:X\times_S X\to X\times_S X$ est un automorphisme
involutif de $X\times_S X$, appelé la symétrie canonique, tel que
\[
  p_1\circ s=p_2,\qquad p_2\circ s=p_1,\qquad s\circ\Delta_f=\Delta_f.
  \tag{16.3.3.1}
  \label{IV.16.3.3.1}
\]
Si l'on pose $s=(\rho,\lambda)$, $p_i=(\pi_i,\mu_i)$ ($i=1,2$),
$\Delta_f=(\delta,\nu)$, $\lambda^\#$ est donc un isomorphisme de
$\rho^*(\pi_1^*(\mathcal{O}_X))$ sur $\pi_2^*(\mathcal{O}_X)$, et
$\delta^*(\lambda^\#)$ laisse invariant $\delta^*(\mathcal{O}_{X\times_S X})$
et le noyau $\mathscr{I}$ de l'homomorphisme
$\nu^\#:\delta^*(\mathcal{O}_{X\times_S X})\to\mathcal{O}_X$. Par suite:
\end{env}

\begin{proposition}[16.3.4]
\EGAReaderTarget{ega4.c04.prop.16-3-4}{IV.16.3.4}
L'homomorphisme $\sigma=\delta^*(\lambda^\#)$ déduit de $s$ (et appelé encore
symétrie canonique) est un automorphisme involutif du système projectif
$(\mathscr{P}_{X/S}^n)$ de faisceaux
\oldpage[IV-4]{15}
d'anneaux $\mathcal{O}_X$-augmentés, et par suite aussi de leur limite projective
$\mathscr{P}_{X/S}^\infty$. Cet automorphisme permute les deux structures de
$\mathcal{O}_X$-Algèbre sur les $\mathscr{P}_{X/S}^n$ et sur
$\mathscr{P}_{X/S}^\infty$.
\end{proposition}

\begin{env}[16.3.5]
\label{IV.16.3.5}
Dans ce qui suit, les deux structures de $\mathcal{O}_X$-Algèbre définies sur les
$\mathscr{P}_{X/S}^n$ et sur $\mathscr{P}_{X/S}^\infty$ joueront des rôles très
différents: nous conviendrons désormais, sauf mention expresse du contraire,
que lorsque $\mathscr{P}_{X/S}^n$ ou $\mathscr{P}_{X/S}^\infty$ sera considéré
comme une $\mathcal{O}_X$-Algèbre, c'est de la structure d'Algèbre définie par
$p_1$ qu'il s'agira.
\end{env}

Pour tout ouvert $U$ de $X$ et toute section $t\in\Gamma(U,\mathcal{O}_X)$, on
notera simplement $t.1$ ou même $t$ l'image de $t$ par l'homomorphisme
structural $\Gamma(U,\mathcal{O}_X)\to\Gamma(U,\mathscr{P}_{X/S}^n)$ (resp.
$\Gamma(U,\mathcal{O}_X)\to\Gamma(U,\mathscr{P}_{X/S}^\infty)$)
(c'est-à-dire l'homomorphisme correspondant à $p_1$).

\begin{definition}[16.3.6]
\EGAReaderTarget{ega4.c04.def.16-3-6}{IV.16.3.6}
On désigne par $d_f^n$, ou $d_{X/S}^n$ (resp. $d_f^\infty$, ou
$d_{X/S}^\infty$), ou simplement $d^n$ (resp. $d^\infty$), l'homomorphisme de
faisceaux d'anneaux $\mathcal{O}_X\to\mathscr{P}_f^n=\mathscr{P}_{X/S}^n$
(resp. $\mathcal{O}_X\to\mathscr{P}_f^\infty=\mathscr{P}_{X/S}^\infty$)
déduit de $p_2$. Pour tout ouvert $U$ de $X$, et tout
$t\in\Gamma(U,\mathcal{O}_X)$, $d^nt$ (resp. $d^\infty t$) est appelé la partie
principale d'ordre $n$ (resp. partie principale d'ordre infini) de $t$. On pose
$dt=d^1t-t$, et on dit que $dt$ est la différentielle de $t$ (élément de
$\Gamma(U,\Omega_{X/S}^1)$, aussi noté $d_{X/S}(t)$).
\end{definition}

Il résulte aussitôt de cette définition que l'on a
\[
  d(t_1t_2)=t_1dt_2+t_2dt_1
  \tag{16.3.6.1}
  \label{IV.16.3.6.1}
\]
quels que soient $t_1,t_2$ dans $\Gamma(U,\mathcal{O}_X)$, autrement dit, $d$
est une dérivation de l'anneau $\Gamma(U,\mathcal{O}_X)$ dans le
$\Gamma(U,\mathcal{O}_X)$-module $\Gamma(U,\Omega_{X/S}^1)$.

Dans toutes les notations introduites dans (16.3.1) et (16.3.6), on remplacera
parfois $S$ par $A$ lorsque $S=\operatorname{Spec}(A)$.

\begin{env}[16.3.7]
\label{IV.16.3.7}
Supposons en particulier que $S=\operatorname{Spec}(A)$ et
$X=\operatorname{Spec}(B)$ soient des schémas affines, $B$ étant donc une
$A$-algèbre. Alors $\Delta_f$ correspond à l'homomorphisme canonique surjectif
$\pi:B\otimes_A B\to B$ tel que $\pi(b\otimes b')=bb'$, de noyau
$\mathfrak{I}=\mathfrak{I}_{B/A}$ (0, 20.4.1); $\mathscr{P}_f^n$ est le faisceau
structural du préschéma $\operatorname{Spec}(P_{B/A}^n)$, où
\[
  P_{B/A}^n=(B\otimes_A B)/\mathfrak{I}^{n+1};
\]
$\mathscr{G}r_\bullet(\mathscr{P}_f)$ est le $\mathcal{O}_X$-Module
quasi-cohérent correspondant au $B$-module gradué
\[
  \operatorname{gr}_{\mathfrak{I}}^\bullet(B\otimes_A B)
  =\bigoplus_{n\geq 0}(\mathfrak{I}^n/\mathfrak{I}^{n+1});
\]
en particulier $\Omega_f^1=\Omega_{X/S}^1$ est le $\mathcal{O}_X$-Module
quasi-cohérent correspondant au $B$-module des 1-différentielles de $B$ par
rapport à $A$, $\Omega_{B/A}^1$ (0, 20.4.3). Les morphismes projections
$p_1:X\times_S X\to X$, $p_2:X\times_S X\to X$ correspondent aux deux
homomorphismes d'anneaux $j_1:B\to B\otimes_A B$, $j_2:B\to B\otimes_A B$
tels que $j_1(b)=b\otimes 1$, $j_2(b)=1\otimes b$, de sorte que (par la
convention de (16.3.5)), $P_{B/A}^n$ est toujours considérée comme une
$B$-algèbre pour l'homomorphisme composé
$B\xrightarrow{j_1}B\otimes_A B\to P_{B/A}^n$; l'homomorphisme d'anneaux
$B\xrightarrow{j_2}B\otimes_A B\to P_{B/A}^n$ se note $d_{B/A}^n$ et correspond
à $d_{X/S}^n$ opérant sur $\Gamma(X,\mathcal{O}_X)$; pour tout $t\in B$, $dt$
est égal à $d_{B/A}t$, défini dans (0, 20.4.6) (cf. $\mathbf{Err}_{\mathrm{IV}}$,
11).

Si $\pi_n:B\otimes_A B\to P_{B/A}^n$ est l'homomorphisme canonique, on a donc,
en vertu des définitions précédentes
\[
  \pi_n(b\otimes b')=b.\pi_n(1\otimes b')=b.d_{B/A}^n(b')
  \quad\text{pour }b\in B,\ b'\in B.
  \tag{16.3.7.1}
  \label{IV.16.3.7.1}
\]
\end{env}

\oldpage[IV-4]{16}
\begin{proposition}[16.3.8]
\label{IV.16.3.8}
L'image de l'homomorphisme $d_{X/S}^n:\mathcal{O}_X\to\mathscr{P}_{X/S}^n$
engendre le $\mathcal{O}_X$-Module $\mathscr{P}_{X/S}^n$.
\end{proposition}

On se ramène aussitôt au cas où $X=\operatorname{Spec}(B)$ et
$S=\operatorname{Spec}(A)$ sont affines et la proposition résulte de (16.3.7.1)
puisque $\pi_n$ est surjectif. On notera qu'en général $d_{X/S}^n$ n'est pas
surjectif (même déjà pour $n=1$).

\begin{proposition}[16.3.9]
\label{IV.16.3.9}
Supposons que $f:X\to S$ soit un morphisme localement de type fini. Alors les
$\mathscr{P}_f^n$ et les $\mathscr{G}r_n(\mathscr{P}_f)$ sont des
$\mathcal{O}_X$-Modules quasi-cohérents de type fini.
\end{proposition}

Cela résulte de (16.1.6) et du fait que $\Delta_f$ est localement de présentation
finie (I, 4.3.1).

\subsection{Propriétés fonctorielles des invariants différentiels}
\EGAReaderTarget{ega4.c04.subsec.16-4}{IV.16.4}

\begin{env}[16.4.1]
\label{IV.16.4.1}
Considérons un diagramme commutatif de morphismes de préschémas
\[
  \xymatrix{
    X \ar[d]_-{f} & X' \ar[l]_-u \ar[d]^-{f'} \\
    S & S' \ar[l]^-w
  }
  \tag{16.4.1.1}
  \label{IV.16.4.1.1}
\]
On en déduit un diagramme commutatif
\[
  \xymatrix{
    X \ar[d]_-{\Delta_f} & X' \ar[l]_-u \ar[d]^-{\Delta_{f'}} \\
    X\times_S X & X'\times_{S'}X' \ar[l]^-v
  }
\]
où $v$ est l'homomorphisme composé (I, 5.3.5 et 5.3.15)
\[
  X'\times_{S'}X'
  \xrightarrow{(p'_1,p'_2)_S}X'\times_S X'
  \xrightarrow{u\times_S u}X\times_S X.
  \tag{16.4.1.2}
  \label{IV.16.4.1.2}
\]
On déduit donc de $u$ et $v$, comme il a été expliqué dans (16.2.1), des
homomorphismes de faisceaux d'anneaux augmentés
\[
  \nu_n:\rho^*(\mathscr{P}_{X/S}^n)\longrightarrow
  \mathscr{P}_{X'/S'}^n
  \tag{16.4.1.3}
  \label{IV.16.4.1.3}
\]
(où l'on a posé $u=(\rho,\lambda)$); ces homomorphismes forment un système
projectif, et donnent donc à la limite un homomorphisme de faisceaux d'anneaux
augmentés
\[
  \nu_\infty:\rho^*(\mathscr{P}_{X/S}^\infty)\longrightarrow
  \mathscr{P}_{X'/S'}^\infty;
  \tag{16.4.1.4}
  \label{IV.16.4.1.4}
\]
d'autre part, par passage aux quotients, les homomorphismes $\nu_n$ donnent un
di-homomorphisme d'Algèbres graduées (relatif à $\lambda^\#$):
\[
  \operatorname{gr}(u):\rho^*(\mathscr{G}r_\bullet(\mathscr{P}_{X/S}))
  \longrightarrow\mathscr{G}r_\bullet(\mathscr{P}_{X'/S'}).
  \tag{16.4.1.5}
  \label{IV.16.4.1.5}
\]
\end{env}

\begin{env}[16.4.2]
\label{IV.16.4.2}
\oldpage[IV-4]{17}
Si l'on a un diagramme commutatif
\[
  \xymatrix{
    X \ar[d]_-{f} & X' \ar[l]_-u \ar[d]^-{f'} &
      X'' \ar[l]_-{u'} \ar[d]^-{f''} \\
    S & S' \ar[l]^-w & S'' \ar[l]^-{w'}
  }
\]
on en déduit un diagramme commutatif
\[
  \xymatrix{
    X \ar[d]_-{\Delta_f} & X' \ar[l]_-u \ar[d]^-{\Delta_{f'}} &
      X'' \ar[l]_-{u'} \ar[d]^-{\Delta_{f''}} \\
    X\times_S X & X'\times_{S'}X' \ar[l]^-v &
      X''\times_{S''}X'' \ar[l]^-{v'}
  }
\]
où $v'$ est défini à partir de $u'$, $w'$, $f'$, $f''$ comme $v$ à partir de
$u$, $w$, $f$, $f'$. On vérifie aussitôt que si $u''=u\circ u'$,
$w''=w\circ w'$, alors le morphisme composé $v\circ v'$ est égal au morphisme
$v''$ défini à partir de $u''$, $w''$, $f$, $f''$ comme $v$ à partir de $u$,
$w$, $f$, $f'$. Si l'on pose $u'=(\rho',\lambda')$,
$u''=(\rho'',\lambda'')$ il résulte alors de (16.2.1) que l'homomorphisme
$\nu_n'':\rho''^*(\mathscr{P}_{X/S}^n)\to\mathscr{P}_{X''/S''}^n$ est égal au
composé
\[
  \rho'^*(\rho^*(\mathscr{P}_{X/S}^n))
  \xrightarrow{\rho'^*(\nu_n^\#)}
  \rho'^*(\mathscr{P}_{X'/S'}^n)
  \xrightarrow{\nu_n'}\mathscr{P}_{X''/S''}^n
\]
et l'on a des propriétés de transitivité analogues pour les homomorphismes
(16.4.1.4) et (16.4.1.5), ce qui permet de dire que les
$\mathscr{P}_{X/S}^n$, $\mathscr{P}_{X/S}^\infty$ et
$\mathscr{G}r_\bullet(\mathscr{P}_{X/S})$ dépendent fonctoriellement de $f$.
\end{env}

\begin{env}[16.4.3]
\label{IV.16.4.3}
On vérifie aussitôt (par exemple en se ramenant au cas affine à l'aide de
(16.3.7)) qu'avec les notations de (16.4.1), le diagramme
\[
  \xymatrix{
    \rho^*(\mathcal{O}_X) \ar[r]^{\lambda^\#} \ar[d] &
      \mathcal{O}_{X'} \ar[d] \\
    \rho^*(\mathscr{P}_{X/S}^n) \ar[r]_{\nu_n} &
      \mathscr{P}_{X'/S'}^n
  }
  \tag{16.4.3.1}
  \label{IV.16.4.3.1}
\]
où les flèches verticales sont celles définissant les structures d'Algèbre
choisies dans (16.3.5) (c'est-à-dire celles provenant des premières projections)
est commutatif; il en est de même du diagramme
\[
  \xymatrix{
    \rho^*(\mathcal{O}_X) \ar[r]^{\lambda^\#}
      \ar[d]_{\rho^*(d_{X/S}^n)} &
      \mathcal{O}_{X'} \ar[d]^{d_{X'/S'}^n} \\
    \rho^*(\mathscr{P}_{X/S}^n) \ar[r]_{\nu_n} &
      \mathscr{P}_{X'/S'}^n
  }
  \tag{16.4.3.2}
  \label{IV.16.4.3.2}
\]
\oldpage[IV-4]{18}
les flèches verticales définissant ici les structures d'Algèbre provenant des
secondes projections; d'ailleurs, si $\sigma$ et $\sigma'$ sont les symétries
canoniques correspondant à $f$ et $f'$ (16.3.4), on a
\[
  \nu_n\circ\rho^*(\sigma)=\sigma'\circ\nu_n
\]
qui fait passer d'un des diagrammes précédents à l'autre. On déduit donc de
(16.4.3.1) un homomorphisme canonique de $\mathcal{O}_{X'}$-Algèbres augmentées
\[
  P^n(u):u^*(\mathscr{P}_{X/S}^n)
  =\mathscr{P}_{X/S}^n\otimes_{\mathcal{O}_X}\mathcal{O}_{X'}
  \longrightarrow\mathscr{P}_{X'/S'}^n
  \tag{16.4.3.3}
  \label{IV.16.4.3.3}
\]
et il résulte de (16.4.3.2) que le diagramme
\[
  \xymatrix{
    \mathcal{O}_{X'} \ar[r]^{\operatorname{id}}
      \ar[d]_{u^*(d_{X/S}^n)} &
      \mathcal{O}_{X'} \ar[d]^{d_{X'/S'}^n} \\
    u^*(\mathscr{P}_{X/S}^n) \ar[r]_{P^n(u)} &
      \mathscr{P}_{X'/S'}^n
  }
  \tag{16.4.3.4}
  \label{IV.16.4.3.4}
\]
est commutatif. On en déduit un homomorphisme de $\mathcal{O}_{X'}$-Algèbres
graduées
\[
  \operatorname{Gr}_\bullet(u):u^*(\mathscr{G}r_\bullet(\mathscr{P}_{X/S}))
  \longrightarrow\mathscr{G}r_\bullet(\mathscr{P}_{X'/S'})
  \tag{16.4.3.5}
  \label{IV.16.4.3.5}
\]
et en particulier un homomorphisme de $\mathcal{O}_{X'}$-Modules
\[
  \operatorname{Gr}_1(u):
  \Omega_{X/S}^1\otimes_{\mathcal{O}_X}\mathcal{O}_{X'}
  \longrightarrow\Omega_{X'/S'}^1
  \tag{16.4.3.6}
  \label{IV.16.4.3.6}
\]
donnant lieu à un diagramme commutatif
\[
  \xymatrix{
    \mathcal{O}_{X'} \ar[r]^{\operatorname{id}}
      \ar[d]_{d_{X/S}\otimes 1} &
      \mathcal{O}_{X'} \ar[d]^{d_{X'/S'}} \\
    \Omega_{X/S}^1\otimes_{\mathcal{O}_X}\mathcal{O}_{X'} \ar[r] &
      \Omega_{X'/S'}^1
  }
  \tag{16.4.3.7}
  \label{IV.16.4.3.7}
\]
\end{env}

\begin{env}[16.4.4]
\label{IV.16.4.4}
Lorsque $S=\operatorname{Spec}(A)$, $S'=\operatorname{Spec}(A')$,
$X=\operatorname{Spec}(B)$, $X'=\operatorname{Spec}(B')$ sont affines, de sorte
que l'on a un diagramme commutatif d'homomorphismes d'anneaux
\[
  \xymatrix{
    B \ar[r] & B' \\
    A \ar[u] \ar[r] & A' \ar[u]
  }
\]
l'image de $\mathfrak{I}_{B/A}$ dans $B'\otimes_{A'}B'$ est contenue dans
$\mathfrak{I}_{B'/A'}$, et l'homomorphisme $\nu_n$ correspond à
l'homomorphisme d'anneaux $P_{B/A}^n\to P_{B'/A'}^n$ déduit de l'homomorphisme
$B\otimes_A B\to B'\otimes_{A'}B'$ par passage aux quotients. L'homomorphisme
(16.4.3.6) correspond à l'homomorphisme défini dans (0, 20.5.4.1), et le
diagramme commutatif (16.4.3.7) au diagramme (0, 20.5.4.2).
\end{env}

\oldpage[IV-4]{19}
\begin{proposition}[16.4.5]
\label{IV.16.4.5}
Supposons que $X'=X\times_S S'$, $f'$ et $u$ étant les projections canoniques.
Alors les homomorphismes canoniques $P^n(u)$ (16.4.3.3) et
$\operatorname{Gr}_1(u)$ (16.4.3.6) sont bijectifs.
\end{proposition}

On a en effet alors $X'\times_{S'}X'=(X\times_S X)\times_S S'$, et on peut donc
appliquer (16.2.3, (ii)) en remplaçant $g$ par la première projection
$p_1:X\times_S X\to X$ et $f$ par la diagonale $\Delta_f$.

On notera aussi que sous les hypothèses de (16.4.5) l'homomorphisme
$\operatorname{Gr}_\bullet(u)$ (16.4.3.5) est surjectif, mais non bijectif en
général. Toutefois (16.2.4):

\begin{corollary}[16.4.6]
\label{IV.16.4.6}
Les hypothèses étant celles de (16.4.5), supposons de plus que $w:S'\to S$ soit
plat (resp. que les $\mathscr{G}r_n(\mathscr{P}_{X/S})$ soient des
$\mathcal{O}_X$-Modules plats pour $n\leq m$); alors l'homomorphisme
\[
  \operatorname{Gr}_n(u):u^*(\mathscr{G}r_n(\mathscr{P}_{X/S}))
  \longrightarrow\mathscr{G}r_n(\mathscr{P}_{X'/S'})
\]
est bijectif pour tout $n$ (resp. pour $n\leq m$).
\end{corollary}

En effet, si $w$ est plat, il en est de même de
$v:X'\times_{S'}X'\to X\times_S X$, donc la conclusion résulte de (16.2.4).

\begin{env}[16.4.7]
\label{IV.16.4.7}
Soient $S$ un préschéma, $\mathscr{E}$ un $\mathcal{O}_S$-Module
quasi-cohérent, et posons $X=\mathbf{V}(\mathscr{E})$ (II, 1.7.8), fibré
vectoriel associé à $\mathscr{E}$, égal à
$\operatorname{Spec}(\mathscr{S}_{\mathcal{O}_S}(\mathscr{E}))$. Soit
$f:X\to S$ le morphisme structural. Pour tout ouvert $U$ de $S$, et toute
section $t\in\Gamma(U,\mathscr{E})$, $t$ s'identifie à une section de
$\mathscr{S}_{\mathcal{O}_S}(\mathscr{E})$ au-dessus de $U$; soit $t'$ son
image dans
$\Gamma(f^{-1}(U),\mathcal{O}_X)=\Gamma(U,f_*(\mathcal{O}_X))=
\Gamma(U,\mathscr{S}_{\mathcal{O}_S}(\mathscr{E}))$, et posons
\[
  \delta(t)=d_{X/S}^n(t')-t'\in
  \Gamma(f^{-1}(U),\mathscr{P}_{X/S}^n);
  \tag{16.4.7.1}
  \label{IV.16.4.7.1}
\]
il est clair que $\delta$ est un di-homomorphisme de modules (correspondant à
l'homomorphisme d'anneaux
$\Gamma(U,\mathcal{O}_S)\to\Gamma(f^{-1}(U),\mathcal{O}_X)$) de
$\Gamma(U,\mathscr{E})$ dans $\Gamma(f^{-1}(U),\mathscr{P}_{X/S}^n)$, dont
l'image appartient d'ailleurs à l'idéal d'augmentation de
$\Gamma(f^{-1}(U),\mathscr{P}_{X/S}^n)$. On en déduit (en faisant varier $U$)
un homomorphisme canonique de $\mathcal{O}_X$-Algèbres
\[
  f^*(\mathscr{S}_{\mathcal{O}_S}(\mathscr{E}))
  \longrightarrow\mathscr{P}_{X/S}^n
  \tag{16.4.7.2}
  \label{IV.16.4.7.2}
\]
et d'après la remarque précédente, si $\mathscr{K}$ est l'Idéal noyau de
l'augmentation $\mathscr{S}_{\mathcal{O}_S}(\mathscr{E})\to\mathcal{O}_S$,
l'image de $\mathscr{K}^{n+1}$ par (16.4.7.2) est nulle, si bien qu'en
factorisant par $\mathscr{K}^{n+1}$, on a finalement un homomorphisme canonique
\[
  \delta_n:f^*(\mathscr{S}_{\mathcal{O}_S}(\mathscr{E})/
  \mathscr{K}^{n+1})\longrightarrow\mathscr{P}_{X/S}^n.
  \tag{16.4.7.3}
  \label{IV.16.4.7.3}
\]
\end{env}

\begin{proposition}[16.4.8]
\label{IV.16.4.8}
Sous les conditions de (16.4.7), les homomorphismes $\delta_n$ sont bijectifs et
forment un système projectif d'isomorphismes; on en déduit un isomorphisme de
$\mathcal{O}_X$-Algèbres graduées
\[
  f^*(\mathscr{S}_{\mathcal{O}_S}^{\bullet}(\mathscr{E}))
  \longrightarrow\mathscr{G}r_\bullet(\mathscr{P}_{X/S}).
  \tag{16.4.8.1}
  \label{IV.16.4.8.1}
\]
\end{proposition}

Le fait que les homomorphismes (16.4.7.3) forment un système projectif résulte
aussitôt de leur définition. Pour prouver que ce sont des isomorphismes, il
suffira de
\oldpage[IV-4]{20}
démontrer que (16.4.8.1) est un isomorphisme, les filtrations des deux membres de
(16.4.7.3) étant finies (Bourbaki, \emph{Alg. comm.}, chap. III, \S~2,
n\textsuperscript{o} 8, cor. 3 du th. 1). Pour cela, considérons la suite exacte
scindée de $\mathcal{O}_S$-Modules
\[
  \xymatrix{
    0 \ar[r] & \mathscr{E} \ar[r]^{u} &
    \mathscr{E}\oplus\mathscr{E} \ar[r]^{v} &
    \mathscr{E} \ar[r] & 0
  }
  \tag{16.4.8.2}
  \label{IV.16.4.8.2}
\]
où, pour tout couple de sections $s$, $t$ de $\mathscr{E}$ au-dessus d'un ouvert
$U$ de $S$, on prend $u(s)=(-s,s)$ et $v(s,t)=s+t$. On a
\[
  X\times_S X=\operatorname{Spec}
  (\mathscr{S}_{\mathcal{O}_S}(\mathscr{E})\otimes_{\mathcal{O}_S}
  \mathscr{S}_{\mathcal{O}_S}(\mathscr{E}))
  =\operatorname{Spec}(\mathscr{S}_{\mathcal{O}_S}
  (\mathscr{E}\oplus\mathscr{E}))
\]
(II, 1.4.6 et 1.7.11), et le morphisme diagonal $X\to X\times_S X$ correspond
(II, 1.2.7) à l'homomorphisme de $\mathcal{O}_S$-Algèbres
$\mathscr{S}(v):\mathscr{S}_{\mathcal{O}_S}(\mathscr{E}\oplus\mathscr{E})
\to\mathscr{S}_{\mathcal{O}_S}(\mathscr{E})$ (II, 1.7.4), de sorte que si
$\mathscr{I}$ est le noyau de cet homomorphisme, on a
\[
  \mathscr{P}_{X/S}^n=f^*(\mathscr{S}_{\mathcal{O}_S}
  (\mathscr{E}\oplus\mathscr{E})/\mathscr{I}^{n+1}).
\]

La proposition sera conséquence du lemme suivant:

\begin{lemma}[16.4.8.3]
\label{IV.16.4.8.3}
Soient $Y$ un espace annelé,
$0\to\mathscr{F}'\xrightarrow{u}\mathscr{F}\xrightarrow{v}\mathscr{F}''\to0$
une suite exacte de $\mathcal{O}_Y$-Modules telle que tout point $y\in Y$
possède un voisinage ouvert $V$ tel que la suite
$0\to\mathscr{F}'|V\to\mathscr{F}|V\to\mathscr{F}''|V\to0$ soit scindée.
Soit $\mathscr{I}$ l'Idéal noyau de
$\mathscr{S}(v):\mathscr{S}_{\mathcal{O}_Y}(\mathscr{F})\to
\mathscr{S}_{\mathcal{O}_Y}(\mathscr{F}'')$, et soit
$\operatorname{gr}_{\mathscr{I}}^\bullet
(\mathscr{S}_{\mathcal{O}_Y}(\mathscr{F}))$ la $\mathcal{O}_Y$-Algèbre graduée
associée à la $\mathcal{O}_Y$-Algèbre
$\mathscr{S}_{\mathcal{O}_Y}(\mathscr{F})$ munie de la filtration
$\mathscr{I}$-préadique. Alors l'homomorphisme de $\mathcal{O}_Y$-Algèbres
graduées
\[
  \mathscr{S}_{\mathcal{O}_Y}^{\bullet}(\mathscr{F}')
  \otimes_{\mathcal{O}_Y}
  \mathscr{S}_{\mathcal{O}_Y}^{\bullet}(\mathscr{F}'')
  \longrightarrow
  \operatorname{gr}_{\mathscr{I}}^\bullet
  (\mathscr{S}_{\mathcal{O}_Y}(\mathscr{F}))
  \tag{16.4.8.4}
  \label{IV.16.4.8.4}
\]
(où le premier membre est le produit tensoriel gradué des
$\mathcal{O}_Y$-Algèbres symétriques munies de leur graduation canonique
(II, 1.7.4 et 2.1.2)), provenant de l'injection canonique
\[
  \mathscr{F}'\longrightarrow\mathscr{I}
  =\operatorname{gr}_{\mathscr{I}}^1
  (\mathscr{S}_{\mathcal{O}_Y}(\mathscr{F})),
\]
est bijectif.
\end{lemma}

L'injection $\mathscr{F}'\to\mathscr{I}$ donne en effet canoniquement un
homomorphisme de $\mathcal{O}_Y$-Algèbres graduées\linebreak
$\mathscr{S}_{\mathcal{O}_Y}^{\bullet}(\mathscr{F}')\to\allowbreak
\operatorname{gr}_{\mathscr{I}}^\bullet
(\mathscr{S}_{\mathcal{O}_Y}(\mathscr{F}))$, et comme le second membre est par
définition une $\mathscr{S}_{\mathcal{O}_Y}^{\bullet}(\mathscr{F}'')$-Algèbre
graduée, on en déduit l'homomorphisme canonique (16.4.8.4) par tensorisation du
précédent avec $\mathscr{S}_{\mathcal{O}_Y}^{\bullet}(\mathscr{F}'')$. Pour
prouver le lemme, on peut, la question étant locale, se borner au cas où
$\mathscr{F}=\mathscr{F}'\oplus\mathscr{F}''$, $u$ et $v$ étant les
homomorphismes canoniques. Alors l'Algèbre graduée
$\mathscr{S}_{\mathcal{O}_Y}^{\bullet}(\mathscr{F})$ s'identifie canoniquement au
produit tensoriel gradué
$\mathscr{S}_{\mathcal{O}_Y}^{\bullet}(\mathscr{F}')
\otimes_{\mathcal{O}_Y}
\mathscr{S}_{\mathcal{O}_Y}^{\bullet}(\mathscr{F}'')$ (II, 1.7.4), et il est
immédiat que $\mathscr{I}$ est alors l'Idéal
$\mathscr{I}'\otimes_{\mathcal{O}_Y}
\mathscr{S}_{\mathcal{O}_Y}^{\bullet}(\mathscr{F}'')$, où $\mathscr{I}'$ est
l'Idéal d'augmentation de
$\mathscr{S}_{\mathcal{O}_Y}^{\bullet}(\mathscr{F}')$, c'est-à-dire la somme
(directe) des $\mathscr{S}_{\mathcal{O}_Y}^m(\mathscr{F}')$ pour $m\geq1$. On en
conclut que $\mathscr{I}^n=\mathscr{I}'^n\otimes_{\mathcal{O}_Y}
\mathscr{S}_{\mathcal{O}_Y}^{\bullet}(\mathscr{F}'')$, où cette fois
$\mathscr{I}'^n$ est la somme (directe) des
$\mathscr{S}_{\mathcal{O}_Y}^m(\mathscr{F}')$ pour $m\geq n$; on a par suite
$\mathscr{I}^n/\mathscr{I}^{n+1}=
\mathscr{S}_{\mathcal{O}_Y}^n(\mathscr{F}')\otimes_{\mathcal{O}_Y}
\mathscr{S}_{\mathcal{O}_Y}^{\bullet}(\mathscr{F}'')$, ce qui prouve que
(16.4.8.4) est bijectif.

\oldpage[IV-4]{21}
Ce lemme étant démontré, il reste à voir que l'homomorphisme (16.4.8.1) est
bien l'image par $f^*$ de l'homomorphisme (16.4.8.4) correspondant à la suite
exacte (16.4.8.2); on constate aisément que cela résulte de la définition de
$u$ (16.4.8.2) et de celle de $\delta$ (16.4.7.1), compte tenu de la définition
de la structure de $\mathcal{O}_X$-Algèbre sur $\mathscr{P}_{X/S}^n$ et de
celle de $d_{X/S}^n$ (16.3.5 et 16.3.6).

En particulier:

\begin{corollary}[16.4.9]
\label{IV.16.4.9}
Sous les conditions de (16.4.7), on a un isomorphisme canonique
\[
  \operatorname{gr}_1(\delta):f^*(\mathscr{E})
  \xrightarrow{\sim}\Omega_{X/S}^1.
  \tag{16.4.9.1}
  \label{IV.16.4.9.1}
\]
\end{corollary}

\begin{corollary}[16.4.10]
\label{IV.16.4.10}
Si $S=\operatorname{Spec}(A)$, $\mathscr{E}=\mathcal{O}_S^m$, de sorte que
\[
  X=\operatorname{Spec}(A[T_1,\ldots,T_m]),
\]
$\mathscr{P}_{X/S}^n$ s'identifie canoniquement à la
$\mathcal{O}_X$-Algèbre correspondant à la $A[T_1,\ldots,T_m]$-algèbre quotient
\[
  A[T_1,\ldots,T_m,U_1,\ldots,U_m]/\mathfrak{K}^{n+1},
\]
où les $U_i$ ($1\leq i\leq m$) sont $m$ nouvelles indéterminées et
$\mathfrak{K}$ est l'idéal engendré par $U_1,\ldots,U_m$.
\end{corollary}

On retrouve ainsi en particulier la structure de $\Omega_{X/S}^1$ dans ce cas
(0, 20.5.13).

Notons en outre que $d_{X/S}^n$ fait alors correspondre à un polynôme
$F(T_1,\ldots,T_m)$ la classe mod.\ $\mathfrak{K}^{n+1}$ de
$F(T_1+U_1,\ldots,T_m+U_m)$, comme il résulte de la définition (16.4.7.1).

\begin{proposition}[16.4.11]
\label{IV.16.4.11}
Soient $f:X\to S$ un morphisme, $g:S\to X$ une $S$-section de $X$,
$S_g^{(n)}$ le $n$-ème voisinage infinitésimal de $S$ pour l'immersion $g$
(16.1.2). Il existe alors un isomorphisme et un seul de
$\mathcal{O}_S$-Algèbres
\[
  \varpi_n:g^*(\mathscr{P}_{X/S}^n)\longrightarrow
  \mathcal{O}_{S_g^{(n)}}
  \tag{16.4.11.1}
  \label{IV.16.4.11.1}
\]
(pour la structure de $\mathcal{O}_S$-Algèbre sur
$\mathcal{O}_{S_g^{(n)}}$ définie par $f$ (16.1.7)), rendant commutatif le
diagramme
\[
  \xymatrix{
    \mathcal{O}_S=g^*(\mathcal{O}_X)
      \ar[rr]^-{\lambda_n}\ar[dr]_-{g^*(d_{X/S}^n)}
      &&\mathcal{O}_{S_g^{(n)}}\\
    &g^*(\mathscr{P}_{X/S}^n)\ar[ur]_-{\varpi_n}
  }
  \tag{16.4.11.2}
  \label{IV.16.4.11.2}
\]
(où $\lambda_n$ est l'homomorphisme structural).
\end{proposition}

En vertu de (I, 5.3.7), où l'on remplace $X$, $Y$, $S$ par $S$, $X$, $S$
respectivement et $f$ par $g$, les diagrammes
\[
  \xymatrix@=3pc{
    S\ar[r]^-g\ar[d]_-g&X\ar[d]^-{\Delta_f}
      &S\ar[r]^-g\ar[d]_-g&X\ar[d]^-{\Delta_f}\\
    X\ar[r]_{(g\circ f,1_X)_S}&X\times_S X
      &X\ar[r]_{(1_X,g\circ f)_S}&X\times_S X
  }
  \tag{16.4.11.3}
  \label{IV.16.4.11.3}
\]
\oldpage[IV-4]{22}
identifient $S$ au produit des $(X\times_S X)$-préschémas $X$ et $X$ pour les
morphismes $\Delta_f$ et $(g\circ f,1_X)_S$ (resp.\ $(1_X,g\circ f)_S$).
D'autre part, les diagrammes
\[
  \xymatrix@=3pc{
    X\ar[r]^{(g\circ f,1_X)_S}\ar[d]_-f&X\times_S X\ar[d]^-{p_1}
      &X\ar[r]^{(1_X,g\circ f)_S}\ar[d]_-f&X\times_S X\ar[d]^-{p_2}\\
    S\ar[r]_{g}&X&S\ar[r]_{g}&X
  }
  \tag{16.4.11.4}
  \label{IV.16.4.11.4}
\]
identifient $X$ au produit des $X$-préschémas $S$ et $X\times_S X$ pour les
morphismes $g$ et $p_1$ (resp.\ $p_2$) (cas particulier de la formule
d'associativité (I, 3.3.9.1)). On peut dire que $\Delta_f$, considéré comme
$X$-section de $X\times_S X$ (relativement à $p_1$ ou $p_2$) joue le rôle
d'une \emph{section universelle} pour les $S$-sections de $X$: chacune de ces
sections $g$ s'en déduit en effet par le \emph{changement de base}
$(g\circ f,1_X)_S:X\to X\times_S X$. La définition de l'homomorphisme
$\varpi_n$ et le fait qu'il est bijectif résultent donc de ces remarques et de
(16.2.3, (ii)) appliqué au premier diagramme (16.4.11.4). La commutativité du
diagramme (16.4.11.2) résulte de même de (16.2.3, (ii)) appliqué cette fois au
second diagramme (16.4.11.4). Pour expliciter $\varpi_n$, on peut se borner au
cas où $g$ est une immersion \emph{fermée}: en effet, pour tout $s\in S$, il y
a un voisinage ouvert $W$ de $s$ dans $S$ tel que $g(W)$ soit fermé dans un
ouvert $U$ de $X$, et il est clair que $g|W$ est une $W$-section du morphisme
$U\cap f^{-1}(W)\to W$, restriction de $f$, et $g(W)$ est \emph{a fortiori}
fermé dans $U\cap f^{-1}(W)$. On peut donc supposer que $S$ est un
sous-préschéma fermé de $X$ défini par un Idéal quasi-cohérent $\mathscr{K}$.
Alors les définitions précédentes montrent que si $W$ est un ouvert de $S$,
$t$ une section de $\mathcal{O}_X$ au-dessus de $f^{-1}(W)$,
$\varpi_n(d^n t|W)$ est égal à l'image canonique de $t$ dans
$\Gamma(W,(\mathcal{O}_X/\mathscr{K}^{n+1})|W)$. L'unicité de $\varpi_n$
résulte donc de ce que l'image de $\mathcal{O}_X$ par $d_{X/S}^n$ engendre le
$\mathcal{O}_X$-Module $\mathscr{P}_{X/S}^n$ (16.3.8).

\begin{corollary}[16.4.12]
\label{IV.16.4.12}
Soient $k$ un corps, $X$ un $k$-préschéma, $x$ un point de $X$ rationnel sur
$k$. Alors $(\mathscr{P}_{X/k}^n)_x\otimes_{\mathcal{O}_x}k(x)$ est
canoniquement isomorphe (en tant que $k(x)$-algèbre augmentée) à
$\mathcal{O}_x/\mathfrak{m}_x^{n+1}$.
\end{corollary}

Il suffit de considérer l'unique $k$-section $g$ de $X$ telle que
$g(\operatorname{Spec}(k))=\{x\}$.

\begin{corollary}[16.4.13]
\label{IV.16.4.13}
Soient $f:X\to S$ un morphisme, $s$ un point de $S$,
$X_s=X\times_S\operatorname{Spec}(k(s))$ la fibre de $f$ en $s$. Si
$x\in X_s$ est rationnel sur $k(s)$,
$(\mathscr{P}_{X/S}^n)_x\otimes_{\mathcal{O}_s}k(s)$ est canoniquement
isomorphe à $\mathcal{O}_{X_s,x}/\mathfrak{m}'_x{}^{\,n+1}$, où
$\mathfrak{m}'_x$ est l'idéal maximal de $\mathcal{O}_{X_s,x}$; de façon
précise, cet isomorphisme fait correspondre à $(d^n t)_x\otimes1$ (où $t$ est
une section de $\mathcal{O}_X$ au-dessus d'un voisinage ouvert de $x$ dans
$X$) la classe de $t_x\otimes1$ modulo $\mathfrak{m}'_x{}^{\,n+1}$.
\end{corollary}

Cela résulte de (16.4.5) et (16.4.12).

Les corollaires précédents justifient la terminologie de «faisceau des parties
principales d'ordre $n$».

\begin{proposition}[16.4.14]
\label{IV.16.4.14}
Soit $\rho:A\to B$ un homomorphisme d'anneaux, $S$ une partie
multiplicativement stable de $B$. Alors les homomorphismes canoniques
\[
  S^{-1}P_{B/A}^n\longrightarrow P_{S^{-1}B/A}^n
  \tag{16.4.14.1}
  \label{IV.16.4.14.1}
\]
\oldpage[IV-4]{23}
déduits des homomorphismes canoniques
$P_{B/A}^n\to P_{S^{-1}B/A}^n$ (16.4.4), forment un système projectif et sont
bijectifs.
\end{proposition}

Il suffit de remarquer que
\[
  S^{-1}((B\otimes_A B)/\mathfrak{I}^{n+1})
  =S^{-1}(B\otimes_A B)/(S^{-1}\mathfrak{I})^{n+1}
\]
par platitude, et que
$S^{-1}(B\otimes_A B)=(S^{-1}B)\otimes_A(S^{-1}B)$ (I, 1.3.4).

\begin{corollary}[16.4.15]
\label{IV.16.4.15}
Les notations étant celles de (16.4.14), soit $R$ une partie multiplicative de
$A$ telle que $\rho(R)\subset S$. Alors on a des isomorphismes canoniques
\[
  S^{-1}P_{B/A}^n\xrightarrow{\sim}P_{S^{-1}B/R^{-1}A}^n
  \tag{16.4.15.1}
  \label{IV.16.4.15.1}
\]
formant un système projectif.
\end{corollary}

Il suffit évidemment de définir des isomorphismes canoniques
\[
  P_{S^{-1}B/A}^n\xrightarrow{\sim}P_{S^{-1}B/R^{-1}A}^n,
  \tag{16.4.15.2}
  \label{IV.16.4.15.2}
\]
c'est-à-dire que l'on est ramené au cas où $\rho(R)$ est formé d'éléments
\emph{inversibles} de $B$. Mais alors l'isomorphisme (16.4.15.2) est simplement
déduit de l'isomorphisme canonique
$B\otimes_A B\to B\otimes_{R^{-1}A}B$ par passage aux quotients
(0$_{\rm I}$, 1.5.3).

\begin{corollary}[16.4.16]
\label{IV.16.4.16}
Soient $f:X\to S$ un morphisme de préschémas, $x$ un point de $X$,
$s=f(x)$. Alors on a des isomorphismes canoniques
\[
  (\mathscr{P}_{X/S}^n)_x\xrightarrow{\sim}
  P_{\mathcal{O}_x/\mathcal{O}_s}^n
  \tag{16.4.16.1}
  \label{IV.16.4.16.1}
\]
formant un système projectif.
\end{corollary}

On déduit de là des isomorphismes pour les gradués associés, et en particulier
un isomorphisme canonique
\[
  (\Omega_{X/S}^1)_x\xrightarrow{\sim}
  \Omega_{\mathcal{O}_x/\mathcal{O}_s}^1.
  \tag{16.4.16.2}
  \label{IV.16.4.16.2}
\]

\begin{corollary}[16.4.17]
\label{IV.16.4.17}
Soient $k$ un corps, $K$ le corps des fractions rationnelles
$k(T_1,\ldots,T_r)$. Alors, pour tout entier $n$, l'homomorphisme de
$K[U_1,\ldots,U_r]$ ($U_i$ indéterminées) dans $P_{K/k}^n$ qui, à tout $U_i$,
fait correspondre $d^nT_i-T_i\cdot1$, est surjectif et définit un isomorphisme
du quotient $K[U_1,\ldots,U_r]/\mathfrak{m}^{n+1}$ (où $\mathfrak{m}$ est
l'idéal engendré par les $U_i$) sur $P_{K/k}^n$.
\end{corollary}

Cela résulte de (16.4.8), (16.4.10) et (16.4.14), où l'on fait $A=k$,
$B=k[T_1,\ldots,T_r]$ et $S=B-\{0\}$.

On retrouve ainsi le fait que les $dT_i$ forment une base du $K$-espace
vectoriel $\Omega_{K/k}^1$ (0, 20.5.10).

\begin{proposition}[16.4.18]
\label{IV.16.4.18}
Soient $f:X\to Y$, $g:Y\to Z$ deux morphismes de préschémas; et considérons
les homomorphismes canoniques de $\mathcal{O}_X$-Algèbres augmentées
(16.4.3.3)
\[
  g_{X/Y/Z}:\mathscr{P}_{X/Z}^n\longrightarrow\mathscr{P}_{X/Y}^n
  \tag{16.4.18.1}
  \label{IV.16.4.18.1}
\]
\[
  f_{X/Y/Z}:f^*(\mathscr{P}_{Y/Z}^n)
  \longrightarrow\mathscr{P}_{X/Z}^n.
  \tag{16.4.18.2}
  \label{IV.16.4.18.2}
\]
Alors $g_{X/Y/Z}$ est surjectif, et son noyau est l'Idéal engendré par l'image
par $f_{X/Y/Z}$ de l'Idéal d'augmentation de
$f^*(\mathscr{P}_{Y/Z}^n)$.
\end{proposition}

\oldpage[IV-4]{24}
Notons d'abord que $g_{X/Y/Z}$ correspond au cas où dans (16.4.3.3) on fait
$X'=X$, $S'=Y$, $S=Z$ et $u=1_X$, et $f_{X/Y/Z}$ au cas où l'on remplace
$X'$, $X$, $S$, $S'$ par $X$, $Y$, $Z$, $Z$ respectivement et $u$, $f$ par
$f$, $g$ respectivement.

On a un diagramme commutatif (I, 5.3.5)
\[
  \xymatrix{
    X\ar[r]^-{\Delta_f}\ar[dr]_-f
      &X\times_Y X\ar[d]^-p\ar[r]^-j
      &X\times_Z X\ar[d]^-{f\times_Z f}\\
    &Y\ar[r]_-{\Delta_g}&Y\times_Z Y
  }
  \tag{16.4.18.3}
  \label{IV.16.4.18.3}
\]
où $j=(1_X,1_X)_Z$ est une immersion, $j\circ\Delta_f=\Delta_{g\circ f}$ et
$p$ est le morphisme structural. Comme on peut se borner au cas où $X$, $Y$,
$Z$ sont affines, on peut supposer les immersions $\Delta_f$, $\Delta_g$ et
$j$ fermées, de sorte que $\mathcal{O}_X$ et
$\mathcal{O}_{X\times_Y X}$ s'identifient respectivement à
$\mathcal{O}_{X\times_Z X}/\mathscr{J}$ et
$\mathcal{O}_{X\times_Z X}/\mathscr{L}$, où
$\mathscr{J}\supset\mathscr{L}$ sont les deux Idéaux quasi-cohérents
correspondant respectivement aux immersions $\Delta_{g\circ f}$ et $j$. La
$\mathcal{O}_X$-Algèbre $\mathscr{P}_{X/Z}^n$ s'identifie donc à
$\mathcal{O}_{X\times_Z X}/\mathscr{J}^{n+1}$, et
$\mathscr{P}_{X/Y}^n$ s'identifie à
$\mathcal{O}_{X\times_Y X}/(\mathscr{J}/\mathscr{L})^{n+1}$, c'est-à-dire à
$\mathcal{O}_{X\times_Z X}/(\mathscr{J}^{n+1}+\mathscr{L})$, et par suite au
quotient de $\mathscr{P}_{X/Z}^n$ par
$(\mathscr{J}^{n+1}+\mathscr{L})/\mathscr{J}^{n+1}$. Mais on sait
(\emph{loc. cit.}) que $p$ et $j$ font de $X\times_Y X$ le produit des
$(Y\times_Z Y)$-préschémas $Y$ et $X\times_Z X$, donc si $\mathcal{O}_Y$ est
identifié à $\mathcal{O}_{Y\times_Z Y}/\mathscr{K}$, où $\mathscr{K}$ est
l'Idéal correspondant à $\Delta_g$, $\mathscr{L}$ est égal à
$(f\times_Z f)^*(\mathscr{K})\cdot\mathcal{O}_{X\times_Z X}$ (I, 4.4.5).
Comme $(\mathscr{J}^{n+1}+\mathscr{L})/\mathscr{J}^{n+1}$ est l'Idéal de
$\mathscr{P}_{X/Z}^n$ engendré par l'image de $\mathscr{L}$, on en déduit la
proposition.

\begin{corollary}[16.4.19]
\label{IV.16.4.19}
Avec les notations de (16.4.18), on a une suite exacte de
$\mathcal{O}_X$-Modules quasi-cohérents
\[
  \xymatrix{
    f^*(\Omega_{Y/Z}^1)\ar[r]^-{f_{X/Y/Z}}&
    \Omega_{X/Z}^1\ar[r]^-{g_{X/Y/Z}}&
    \Omega_{X/Y}^1\ar[r]&0.
  }
  \tag{16.4.19.1}
  \label{IV.16.4.19.1}
\]
\end{corollary}

Lorsque $X$, $Y$, $Z$ sont affines, on retrouve ainsi la suite exacte
(0, 20.5.7.1).

\begin{proposition}[16.4.20]
\label{IV.16.4.20}
Soient $f:Y\to Z$ un morphisme, $j:X\to Y$ une immersion fermée,
$\mathscr{K}$ l'Idéal quasi-cohérent de $\mathcal{O}_Y$ correspondant à $j$.
Alors on a
$\mathscr{P}_{X/Y}^n=\mathcal{O}_X=\mathcal{O}_Y/\mathscr{K}$,
l'homomorphisme canonique
$j_{X/Y/Z}:j^*(\mathscr{P}_{Y/Z}^n)\to\mathscr{P}_{X/Z}^n$ est surjectif, et
son noyau est l'Idéal de $j^*(\mathscr{P}_{Y/Z}^n)$ engendré par
$j^*(\mathcal{O}_Y\cdot d_{Y/Z}^n(\mathscr{K}))$ (il faut noter que
$d_{Y/Z}^n(\mathscr{K})$ est un sous-faisceau de groupes commutatifs de
$\mathscr{P}_{Y/Z}^n$, mais non un $\mathcal{O}_Y$-Module en général).
\end{proposition}

On sait (I, 5.3.8) que la diagonale
$\Delta_j:X\to X\times_Y X$ est un isomorphisme, d'où la première assertion.
Si $\varpi_1$ et $\varpi_2$ sont les deux homomorphismes d'Algèbres
$\mathcal{O}_Y\to\mathscr{P}_{Y/Z}^n$ correspondant respectivement aux deux
projections canoniques $p_1$, $p_2$ de $Y\times_Z Y$ sur $Y$, rappelons que
par définition (16.3.5 et 16.3.6) $\varpi_1$ est l'homomorphisme structural de
la $\mathcal{O}_Y$-Algèbre $\mathscr{P}_{Y/Z}^n$ et
$\varpi_2=d_{Y/Z}^n$. La $\mathcal{O}_X$-Algèbre
$j^*(\mathscr{P}_{Y/Z}^n)$ s'identifie donc à
$\mathscr{P}_{Y/Z}^n/\varpi_1(\mathscr{K})\mathscr{P}_{Y/Z}^n$ et son quotient
par l'Idéal engendré par $j^*(d_{Y/Z}^n(\mathscr{K}))$ à
\[
  \mathscr{P}_{Y/Z}^n/
  (\varpi_1(\mathscr{K})+\varpi_2(\mathscr{K}))\mathscr{P}_{Y/Z}^n.
\]
Notons maintenant qu'on a un diagramme commutatif
\oldpage[IV-4]{25}
\[
  \xymatrix{
    Y\ar[d]_-{\Delta_f}&X\ar[l]_-j\ar[d]^-{\Delta_{f\circ j}}\\
    Y\times_Z Y&X\times_Z X\ar[l]^-{j\times_Z j}
  }
\]
identifiant $X$ au produit des $(Y\times_Z Y)$-préschémas $Y$ et
$X\times_Z X$ (I, 5.3.7). Comme $j\times_Z j$ est une immersion, on déduit
donc de cette remarque et de (16.2.2) que si $\Delta_{Y/Z}^n$ et
$\Delta_{X/Z}^n$ désignent les voisinages infinitésimaux d'ordre $n$ de $Y$
et $X$ pour les immersions canoniques $\Delta_f$ et $\Delta_{f\circ j}$
respectivement, on a un diagramme
\[
  \xymatrix{
    \Delta_{Y/Z}^n\ar[d]&\Delta_{X/Z}^n\ar[l]\ar[d]\\
    Y\times_Z Y&X\times_Z X\ar[l]^-{j\times_Z j}
  }
\]
faisant de $\Delta_{X/Z}^n$ le produit des $(Y\times_Z Y)$-préschémas
$\Delta_{Y/Z}^n$ et $X\times_Z X$. On peut encore dire que
$\mathscr{P}_{X/Z}^n$ s'identifie au faisceau d'anneaux
$\mathscr{P}_{Y/Z}^n\otimes_{\mathcal{O}_{Y\times_Z Y}}
\mathcal{O}_{X\times_Z X}$. Mais on voit aussitôt (par exemple en se ramenant
au cas affine) que
\[
  \mathcal{O}_{X\times_Z X}=
  \mathcal{O}_{Y\times_Z Y}/
  (p_1^*(\mathscr{K})+p_2^*(\mathscr{K}))
  \mathcal{O}_{Y\times_Z Y}.
\]
Donc $\mathscr{P}_{X/Z}^n$ s'identifie au quotient de
$\mathscr{P}_{Y/Z}^n$ par l'Idéal engendré par l'image dans
$\mathscr{P}_{Y/Z}^n$ de $p_1^*(\mathscr{K})+p_2^*(\mathscr{K})$. Mais par
définition cet Idéal est aussi engendré par
$\varpi_1(\mathscr{K})+\varpi_2(\mathscr{K})$. C.Q.F.D.

\begin{corollary}[16.4.21]
\label{IV.16.4.21}
Soient $f:Y\to Z$ un morphisme, $j:X\to Y$ une immersion. On a une suite
exacte de $\mathcal{O}_X$-Modules quasi-cohérents
\[
  \xymatrix{
    \mathscr{N}_{X/Y}\ar[r]&j^*(\Omega_{Y/Z}^1)\ar[r]&
    \Omega_{X/Z}^1\ar[r]&0.
  }
  \tag{16.4.21.1}
  \label{IV.16.4.21.1}
\]
\end{corollary}

Lorsque $X$, $Y$ et $Z$ sont affines, on retrouve ainsi la suite exacte
(0, 20.5.12.1).

\begin{corollary}[16.4.22]
\label{IV.16.4.22}
Si $f:X\to S$ est un morphisme localement de présentation finie,
$\mathscr{P}_{X/S}^n$ et $\Omega_{X/S}^1$ sont des
$\mathcal{O}_X$-Modules quasi-cohérents de présentation finie.
\end{corollary}

On est aussitôt ramené au cas où $S=\operatorname{Spec}(A)$ est affine,
$X=\operatorname{Spec}(B)$, où
$B=A[T_1,\ldots,T_r]/\mathfrak{K}$, $\mathfrak{K}$ étant un idéal de type
fini de $C=A[T_1,\ldots,T_r]$. On applique alors (16.4.20) avec $Z=S$,
$Y=\operatorname{Spec}(C)$ et $\mathscr{K}=\widetilde{\mathfrak{K}}$. Alors
$j^*(\mathscr{P}_{Y/Z}^n)$ est un $\mathcal{O}_X$-Module libre de rang fini
(16.4.10) et l'hypothèse sur $\mathfrak{K}$ entraîne que
$j^*(\mathcal{O}_Y\cdot d_{Y/Z}^n(\mathscr{K}))$ engendre un
$\mathcal{O}_X$-Module quasi-cohérent de type fini; d'où la conclusion.

\begin{proposition}[16.4.23]
\label{IV.16.4.23}
Soient $X$, $Y$ deux $S$-préschémas, $Z=X\times_S Y$ leur produit,
$p:X\times_S Y\to X$ et $q:X\times_S Y\to Y$ les projections canoniques.
Alors l'homomorphisme canonique
\[
  p_{Z/X/S}\oplus q_{Z/Y/S}:
  p^*(\Omega_{X/S}^1)\oplus q^*(\Omega_{Y/S}^1)
  \longrightarrow\Omega_{(X\times_S Y)/S}^1
  \tag{16.4.23.1}
  \label{IV.16.4.23.1}
\]
est bijectif.
\end{proposition}

\oldpage[IV-4]{26}
Le diagramme commutatif
\[
  \xymatrix{
    Y\ar[d]_-g&X\times_S Y\ar[l]_-q\ar[d]_-h&
      X\times_S Y\ar[l]_-{\operatorname{id}}\ar[d]^-p\\
    S&S\ar[l]^-{\operatorname{id}}&X\ar[l]^-f
  }
\]
donne une factorisation de l'\emph{isomorphisme} canonique $P^n(p)$ (16.4.5)
\[
  p^*(\mathscr{P}_{X/S}^n)\longrightarrow
  \mathscr{P}_{Z/S}^n\longrightarrow\mathscr{P}_{Z/Y}^n
\]
et de même, en intervertissant les rôles de $X$ et $Y$, on a une factorisation
de l'\emph{isomorphisme} $P^n(q)$
\[
  q^*(\mathscr{P}_{Y/S}^n)\longrightarrow
  \mathscr{P}_{Z/S}^n\longrightarrow\mathscr{P}_{Z/X}^n.
\]
Ceci prouve que l'homomorphisme canonique (16.4.18.1)
\[
  p_{Z/X/S}:p^*(\mathscr{P}_{X/S}^n)\longrightarrow\mathscr{P}_{Z/S}^n
  \quad\text{(resp.\
  $q_{Z/Y/S}:q^*(\mathscr{P}_{Y/S}^n)\to\mathscr{P}_{Z/S}^n$)}
\]
est \emph{injectif}, et que le noyau de l'homomorphisme canonique surjectif
(16.4.18.2)
\[
  \mathscr{P}_{Z/S}^n\longrightarrow\mathscr{P}_{Z/Y}^n
  \quad\text{(resp.\
  $\mathscr{P}_{Z/S}^n\to\mathscr{P}_{Z/X}^n$)}
\]
est supplémentaire de l'image de $p_{Z/X/S}$ (resp.\ $q_{Z/Y/S}$). Mais
d'autre part, ce noyau est, en vertu de (16.4.18), engendré par l'image par
$q_{Z/Y/S}$ (resp.\ $p_{Z/X/S}$) de l'Idéal d'augmentation de
$q^*(\mathscr{P}_{Y/S}^n)$ (resp.\ $p^*(\mathscr{P}_{X/S}^n)$). On en conclut
la proposition en considérant le cas $n=1$.

On généralise aussitôt (16.4.23) au cas d'un produit d'un nombre fini
quelconque de $S$-préschémas.

\begin{remarks}[16.4.24]
\label{IV.16.4.24}
\begin{enumerate}
\item[(i)] Nous verrons (17.2.3) que lorsque le morphisme $f:X\to Y$ dans
(16.4.18) est \emph{lisse}, l'homomorphisme $f_{X/Y/Z}$ dans (16.4.19.1) est
localement \emph{inversible à gauche} et en particulier injectif. De même,
lorsque le morphisme $f\circ j:X\to Z$ de (16.4.20) est \emph{lisse},
l'homomorphisme de gauche dans (16.4.21.2) est localement
\emph{inversible à gauche} et \emph{a fortiori} injectif (17.2.5). Au chap. V,
nous donnerons aussi une variante, dans le cas des Modules sur les
préschémas, des «modules d'imperfection» étudiés dans (0, 20.6), et des suites
exactes où ils figurent.

\item[(ii)] Soient $X$ un espace topologique, $\mathscr{A}$ un faisceau
d'anneaux sur $X$ et $\mathscr{B}$ une $\mathscr{A}$-Algèbre sur $X$. Alors il
est clair que
\[
  U\longmapsto
  P_{\Gamma(U,\mathscr{B})/\Gamma(U,\mathscr{A})}^n
  \qquad (U\text{ ouvert dans }X)
\]
est un préfaisceau de $\Gamma(U,\mathscr{B})$-algèbres augmentées, donc le
faisceau associé $\mathscr{P}_{\mathscr{B}/\mathscr{A}}^n$ est une
$\mathscr{B}$-Algèbre augmentée. Dans le cas particulier où $X$ est un
préschéma, $f=(\psi,\theta):X\to S$ un morphisme de préschémas, il résulte
aisément de (16.4.16) et de l'exactitude du foncteur $\varinjlim$ que
$\mathscr{P}_{X/S}^n$ est canoniquement isomorphe à
$\mathscr{P}_{\mathcal{O}_X/\psi^*(\mathcal{O}_S)}^n$. Il s'ensuit que le
formalisme développé dans le présent paragraphe pourrait être considéré comme
un
\oldpage[IV-4]{27}
cas particulier d'un formalisme différentiel pour des espaces annelés munis
d'un faisceau d'algèbres sur le faisceau structural. Cependant, nous n'avons
pas voulu partir de ce point de vue, moins intuitif et moins commode pour les
applications. Il semble d'ailleurs que, pour les diverses espèces de
«variétés», la construction «globale» des $\mathscr{P}^n$ analogue à celle que
nous utilisons ici soit également mieux adaptée aux applications.
\end{enumerate}
\end{remarks}

\subsection{Faisceaux et fibrés tangents relatifs; dérivations.}
\EGAReaderTarget{ega4.c04.subsec.16-5}{IV.16.5}

\begin{env}[16.5.1]
\EGAReaderTarget{ega4.c04.statement.16-5-1}{IV.16.5.1}
Soit $f=(\psi,\theta):X\to S$ un morphisme d'espaces annelés. Pour tout
$\mathcal{O}_X$-Module $\mathscr{F}$, on appelle \emph{$S$-dérivation}
(ou \emph{$(X/S)$-dérivation}, ou \emph{$f$-dérivation}) \emph{de
$\mathcal{O}_X$ dans $\mathscr{F}$} tout homomorphisme de \emph{faisceaux de
groupes additifs} $D:\mathcal{O}_X\to\mathscr{F}$, vérifiant les conditions
suivantes:
\begin{enumerate}
\item[(a)] pour tout ouvert $V$ de $X$, et tout couple de sections
$(t_1,t_2)$ de $\mathcal{O}_X$ au-dessus de $V$, on a
\[
  D(t_1t_2)=t_1D(t_2)+D(t_1)t_2;
  \tag{16.5.1.1}
  \label{IV.16.5.1.1}
\]
\item[(b)] pour tout ouvert $V$ de $X$, toute section $t$ de $\mathcal{O}_X$
au-dessus de $V$ et toute section $s$ de $\mathcal{O}_S$ au-dessus d'un ouvert
$U$ de $S$ tel que $V\subset f^{-1}(U)$, on a
\[
  D((s|V)t)=(s|V)D(t).
  \tag{16.5.1.2}
  \label{IV.16.5.1.2}
\]
\end{enumerate}
Il est clair qu'il revient au même de dire que pour tout $x\in X$,
l'homomorphisme de groupes additifs
$D_x:\mathcal{O}_x\to\mathscr{F}_x$ est une
$\mathcal{O}_{f(x)}$-dérivation.

Une autre interprétation consiste à considérer la $\mathcal{O}_X$-Algèbre
$\mathscr{D}_{\mathcal{O}_X}(\mathscr{F})$ égale à
$\mathcal{O}_X\oplus\mathscr{F}$, la structure d'Algèbre étant définie par la
condition que pour tout ouvert $V$ de $X$, le produit de deux sections de
$\mathcal{O}_X$ (resp. d'une section de $\mathcal{O}_X$ et d'une section de
$\mathscr{F}$) au-dessus de $V$ est défini par la structure d'anneau de
$\Gamma(V,\mathcal{O}_X)$ (resp. la structure de
$\Gamma(V,\mathcal{O}_X)$-module sur $\Gamma(V,\mathscr{F})$), et le produit de
deux sections de $\mathscr{F}$ au-dessus de $V$ est pris égal à 0; alors,
$\mathscr{F}$ est un Idéal de
$\mathscr{D}_{\mathcal{O}_X}(\mathscr{F})$, noyau de l'augmentation canonique
$\mathscr{D}_{\mathcal{O}_X}(\mathscr{F})\to\mathcal{O}_X$, et dire que $D$
est une $S$-dérivation de $\mathcal{O}_X$ dans $\mathscr{F}$ signifie que
$1_{\mathcal{O}_X}+D$ est un $\mathcal{O}_S$-homomorphisme d'Algèbres de
$\mathcal{O}_X$ dans $\mathscr{D}_{\mathcal{O}_X}(\mathscr{F})$, qui, composé
avec l'augmentation, donne $1_{\mathcal{O}_X}$.

Les $S$-dérivations de $\mathcal{O}_X$ dans $\mathscr{F}$ forment évidemment
un $\Gamma(X,\mathcal{O}_X)$-module
$\operatorname{D\acute{e}r}_S(\mathcal{O}_X,\mathscr{F})$.

Lorsque $\mathscr{F}=\mathcal{O}_X$, une $S$-dérivation de $\mathcal{O}_X$ dans
lui-même est simplement appelée une \emph{$S$-dérivation de $\mathcal{O}_X$}.
\end{env}

\begin{proposition}[16.5.2]
\label{IV.16.5.2}
Soient $A$ un anneau, $B$ une $A$-algèbre, $L$ un $B$-module; on pose
$S=\operatorname{Spec}(A)$, $X=\operatorname{Spec}(B)$,
$\mathscr{F}=\widetilde{L}$. Alors l'application $D\mapsto\Gamma(D)$ qui, à
toute $S$-dérivation $D$ de $\mathcal{O}_X$ dans $\mathscr{F}$ fait
correspondre l'application $\Gamma(D):t\mapsto D(t)$ de $B$ dans $L$, est un
isomorphisme de $B$-modules de
$\operatorname{D\acute{e}r}_S(\mathcal{O}_X,\mathscr{F})$ sur
$\operatorname{D\acute{e}r}_A(B,L)$ (cf. (0, 20.1.2)).
\end{proposition}

Cela résulte aussitôt de l'interprétation donnée ci-dessus des $S$-dérivations
en termes
\oldpage[IV-4]{28}
d'homomorphismes d'Algèbres, de l'interprétation analogue donnée dans
(0, 20.1.6), et de la correspondance canonique entre homomorphismes de
$\mathcal{O}_X$-Algèbres et homomorphismes de $B$-algèbres (I, 1.3.13 et
1.3.8).

\begin{proposition}[16.5.3]
\label{IV.16.5.3}
Soit $f=(\psi,\theta):X\to S$ un morphisme de préschémas.
\begin{enumerate}
\item[(i)] La différentielle
$d_{X/S}:\mathcal{O}_X\to\Omega_{X/S}^1$ (16.3.6) est une $S$-dérivation.
\item[(ii)] Pour tout $\mathcal{O}_X$-Module $\mathscr{F}$, l'application
$u\mapsto u\circ d_{X/S}$ est un isomorphisme de
$\Gamma(X,\mathcal{O}_X)$-modules
\[
  \operatorname{Hom}_{\mathcal{O}_X}(\Omega_{X/S}^1,\mathscr{F})
  \xrightarrow{\sim}
  \operatorname{D\acute{e}r}_S(\mathcal{O}_X,\mathscr{F}).
  \tag{16.5.3.1}
  \label{IV.16.5.3.1}
\]
\end{enumerate}
\end{proposition}

L'assertion (i) a déjà été notée (16.3.6). D'autre part, il est immédiat (en
vertu de (0, 20.4.8)) que $u\mapsto u\circ d_{X/S}$ est injectif, en
considérant les restrictions à une fibre $\mathcal{O}_x$ des deux membres et
utilisant (16.4.16.2). Pour voir que l'homomorphisme (16.5.3.1) est surjectif,
considérons une $S$-dérivation $D:\mathcal{O}_X\to\mathscr{F}$; pour tout
ouvert affine $V=\operatorname{Spec}(B)$ de $X$, tel que $f(V)$ soit contenu
dans un ouvert affine $U=\operatorname{Spec}(A)$ de $S$,
$D_V:B\to\Gamma(V,\mathscr{F})$ est une $A$-dérivation, donc il existe un
\emph{unique} $B$-homomorphisme
$u_V:\Omega_{B/A}^1\to\Gamma(V,\mathscr{F})$ tel que
$D_V=u_V\circ d_{B/A}$ (0, 20.4.8); en outre l'unicité de $u_V$ montre
aussitôt que pour tout ouvert affine $W\subset V$, on a $u_W=u_V|W$, donc les
$u_V$ définissent un homomorphisme de $\mathcal{O}_X$-Modules
$u:\mathcal{O}_X\to\mathscr{F}$ répondant à la question.

\begin{env}[16.5.4]
\EGAReaderTarget{ega4.c04.statement.16-5-4}{IV.16.5.4}
Avec les notations de (16.5.1), pour tout ouvert $U$ de $X$,
$\operatorname{D\acute{e}r}_S(\mathcal{O}_U,\mathscr{F}|U)$ est un
$\Gamma(U,\mathcal{O}_X)$-module et il est clair que l'application
$U\mapsto\operatorname{D\acute{e}r}_S(\mathcal{O}_U,\mathscr{F}|U)$ est un
préfaisceau; en fait, c'est même un \emph{faisceau} (donc un
$\mathcal{O}_X$-Module), en vertu de la caractérisation ponctuelle des
$S$-dérivations, vue en (16.5.1). Cet $\mathcal{O}_X$-Module se note
$\mathscr{D}\!\mathrm{\acute{e}r}_S(\mathcal{O}_X,\mathscr{F})$ et est appelé
\emph{faisceau des $S$-dérivations de $\mathcal{O}_X$ dans $\mathscr{F}$}, et
ce qu'on vient de voir s'exprime encore par le corollaire suivant:
\end{env}

\begin{corollary}[16.5.5]
\label{IV.16.5.5}
Pour tout $\mathcal{O}_X$-Module $\mathscr{F}$, l'homomorphisme de
$\mathcal{O}_X$-Modules déduit de $u\mapsto u\circ d_{X/S}$
\[
  \mathscr{H}\!om_{\mathcal{O}_X}(\Omega_{X/S}^1,\mathscr{F})
  \longrightarrow
  \mathscr{D}\!\mathrm{\acute{e}r}_S(\mathcal{O}_X,\mathscr{F})
  \tag{16.5.5.1}
  \label{IV.16.5.5.1}
\]
est bijectif.
\end{corollary}

\begin{corollary}[16.5.6]
\label{IV.16.5.6}
\begin{enumerate}
\item[(i)] Si le morphisme $f:X\to S$ est localement de présentation finie et
si $\mathscr{F}$ est un $\mathcal{O}_X$-Module quasi-cohérent, alors
$\mathscr{D}\!\mathrm{\acute{e}r}_S(\mathcal{O}_X,\mathscr{F})$ est un
$\mathcal{O}_X$-Module quasi-cohérent.
\item[(ii)] Si de plus $S$ est localement noethérien et si $\mathscr{F}$ est
cohérent, alors $\mathscr{D}\!\mathrm{\acute{e}r}_S(\mathcal{O}_X,\mathscr{F})$ est un
$\mathcal{O}_X$-Module cohérent.
\end{enumerate}
\end{corollary}

L'assertion (i) résulte de l'isomorphisme (16.5.5.1), de (16.4.22) et
(I, 1.3.12); l'assertion (ii) résulte de (0$_{\rm I}$, 5.3.5).

\begin{env}[16.5.7]
\EGAReaderTarget{ega4.c04.statement.16-5-7}{IV.16.5.7}
On pose
\[
  \mathfrak{G}_{X/S}=
  \mathscr{H}\!om_{\mathcal{O}_X}(\Omega_{X/S}^1,\mathcal{O}_X)
  =\mathscr{D}\!\mathrm{\acute{e}r}_S(\mathcal{O}_X,\mathcal{O}_X)
  \tag{16.5.7.1}
  \label{IV.16.5.7.1}
\]
et on dit que c'est le \emph{faisceau des $S$-dérivations de
$\mathcal{O}_X$}, ou encore le \emph{faisceau tangent de $X$ relativement à
$S$}: c'est donc le \emph{dual} du $\mathcal{O}_X$-Module
$\Omega_{X/S}^1$. Si $f$ est localement de présentation
\oldpage[IV-4]{29}
finie, $\mathfrak{G}_{X/S}$ est un $\mathcal{O}_X$-Module quasi-cohérent; si
de plus $S$ est localement noethérien, $\mathfrak{G}_{X/S}$ est cohérent
(16.5.6).
\end{env}

\begin{env}[16.5.8]
\label{IV.16.5.8}
Supposons plus particulièrement que $\Omega_{X/S}^1$ soit un
$\mathcal{O}_X$-Module \emph{localement libre} (de rang fini) (ce qui sera le
cas lorsque $f$ est \emph{lisse} (17.2.3)); alors $\mathfrak{G}_{X/S}$ est un
$\mathcal{O}_X$-Module localement libre et a même rang que $\Omega_{X/S}^1$
en chaque point. De façon plus précise, supposons que $\Omega_{X/S}^1$ soit de
rang $n$ en un point $x$; il y a alors $n$ sections $s_i$ ($1\leq i\leq n$)
de $\mathcal{O}_X$ au-dessus d'un voisinage affine $U$ de $x$ telles que les
images canoniques des $ds_i$ dans
$\Omega_{X/S}^1\otimes_{\mathcal{O}_X}k(x)$ forment une base de ce
$k(x)$-espace vectoriel; en vertu du lemme de Nakayama, les germes
$(ds_i)_x=d(s_i)_x$ des $ds_i$ au point $x$ forment une base du
$\mathcal{O}_x$-module $(\Omega_{X/S}^1)_x$, donc, en restreignant $U$, on
peut supposer que les $ds_i$ forment une \emph{base} du
$\Gamma(U,\mathcal{O}_X)$-module $\Gamma(U,\Omega_{X/S}^1)$. Alors le
$\Gamma(U,\mathcal{O}_X)$-module $\Gamma(U,\mathfrak{G}_{X/S})$ est dual du
précédent; on note $(D_i)_{1\leq i\leq n}$ ou
$\left(\frac{\partial}{\partial s_i}\right)_{1\leq i\leq n}$ la base duale de
$(ds_i)_{1\leq i\leq n}$, de sorte que, par (16.5.3), on a
\[
  D_is_j=\langle D_i,ds_j\rangle
  =\left\langle\frac{\partial}{\partial s_i},ds_j\right\rangle
  =\delta_{ij}\qquad\text{(indice de Kronecker).}
  \tag{16.5.8.1}
  \label{IV.16.5.8.1}
\]
Toute $\Gamma(S,\mathcal{O}_S)$-dérivation de la
$\Gamma(S,\mathcal{O}_S)$-algèbre $\Gamma(U,\mathcal{O}_X)$ s'écrit donc
d'une seule manière
\[
  D=\sum_{i=1}^n a_iD_i=\sum_{i=1}^n a_i\frac{\partial}{\partial s_i}
\]
où les $a_i$ ($1\leq i\leq n$) sont des sections de $\mathcal{O}_X$ au-dessus
de $U$. Pour toute section $g\in\Gamma(U,\mathcal{O}_X)$, si l'on pose
$dg=\sum_{i=1}^n c_i ds_i$, on a
$c_i=\langle D_i,dg\rangle=D_ig$ en vertu de (16.5.8.1), autrement dit
\[
  dg=\sum_{i=1}^n(D_ig)ds_i
  =\sum_{i=1}^n\frac{\partial g}{\partial s_i}ds_i.
  \tag{16.5.8.2}
  \label{IV.16.5.8.2}
\]
\end{env}

\begin{env}[16.5.9]
\label{IV.16.5.9}
Soient $D_1$, $D_2$ deux $S$-dérivations de $\mathcal{O}_X$. Pour tout ouvert
$U$ de $X$, si $D_1^U$, $D_2^U$ sont les dérivations correspondantes de
l'anneau $\Gamma(U,\mathcal{O}_X)$, le \emph{crochet}
\[
  [D_1^U,D_2^U]=D_1^U\circ D_2^U-D_2^U\circ D_1^U
\]
est aussi une dérivation de cet anneau, donc le
$\psi^*(\mathcal{O}_S)$-endomorphisme de $\mathcal{O}_X$
\[
  [D_1,D_2]=D_1\circ D_2-D_2\circ D_1
  \tag{16.5.9.1}
  \label{IV.16.5.9.1}
\]
est encore une $S$-dérivation; comme on vérifie aussitôt que ce crochet
vérifie l'identité de Jacobi, on voit qu'on a ainsi défini sur
$\operatorname{D\acute{e}r}_S(\mathcal{O}_X,\mathcal{O}_X)$ une structure de
$\Gamma(S,\mathcal{O}_S)$-algèbre de Lie. Comme la définition de cette
structure commute à la restriction à un ouvert de $X$, on voit que
$\mathfrak{G}_{X/S}$ est canoniquement muni d'une structure de
$\psi^*(\mathcal{O}_S)$-Algèbre de Lie. On notera que l'application
$(D_1,D_2)\mapsto[D_1,D_2]$ n'est pas
$\Gamma(X,\mathcal{O}_X)$-bilinéaire.
\end{env}

\begin{env}[16.5.10]
\label{IV.16.5.10}
Pour tout changement de base $g:S'\to S$, si l'on pose $X'=X\times_S S'$, on
a vu (16.4.5) que l'on a un isomorphisme canonique
\[
  \Omega_{X/S}^1\otimes_S S'\xrightarrow{\sim}\Omega_{X'/S'}^1
  \tag{16.5.10.1}
  \label{IV.16.5.10.1}
\]
\oldpage[IV-4]{30}
d'où l'on déduit, en vertu de (16.5.10.1), un homomorphisme canonique
(Bourbaki, \emph{Alg.}, chap. II, 3\textsuperscript{e} éd., \S~5,
n\textsuperscript{o} 3)
\[
  \mathfrak{G}_{X/S}\otimes_{\mathcal{O}_S}\mathcal{O}_{S'}
  \longrightarrow\mathfrak{G}_{X'/S'}
  \tag{16.5.10.2}
  \label{IV.16.5.10.2}
\]
qui en général n'est ni injectif ni surjectif. Toutefois:
\end{env}

\begin{proposition}[16.5.11]
\label{IV.16.5.11}
\begin{enumerate}
\item[(i)] Si $g:S'\to S$ est un morphisme plat et si $f$ est localement de
type fini (resp.\ localement de présentation finie), l'homomorphisme
(16.5.10.2) est injectif (resp.\ bijectif).
\item[(ii)] Si $\Omega_{X/S}^1$ est un $\mathcal{O}_X$-Module localement libre
de type fini, l'homomorphisme (16.5.10.2) est bijectif.
\end{enumerate}
\end{proposition}

En effet, l'assertion (ii) résulte de Bourbaki, \emph{Alg.}, chap. II,
3\textsuperscript{e} éd., \S~5, n\textsuperscript{o} 3, prop. 7. L'assertion
(i) résulte de même de Bourbaki, \emph{Alg. comm.}, chap. I, \S~2,
n\textsuperscript{o} 10, prop. 11 et du fait que si $f$ est localement de type
fini (resp.\ localement de présentation finie), $\Omega_{X/S}^1$ est un
$\mathcal{O}_X$-Module de type fini (resp.\ de présentation finie) ((16.3.9)
et (16.4.22)).

\begin{env}[16.5.12]
\EGAReaderTarget{ega4.c04.statement.16-5-12}{IV.16.5.12}
Puisque $\Omega_{X/S}^1$ est un $\mathcal{O}_X$-Module quasi-cohérent, on peut
considérer le fibré vectoriel sur $X$ défini par $\Omega_{X/S}^1$
(II, 1.7.8)
\[
  T_{X/S}=\mathbf{V}(\Omega_{X/S}^1)
  \tag{16.5.12.1}
  \label{IV.16.5.12.1}
\]
qu'on appelle \emph{fibré tangent de $X$ relativement à $S$}. On a donc une
bijection canonique (II, 1.7.9)
\[
  \Gamma(T_{X/S}/S)\xrightarrow{\sim}
  \operatorname{Hom}_{\mathcal{O}_X}(\Omega_{X/S}^1,\mathcal{O}_X)
  =\Gamma(X,\mathfrak{G}_{X/S})
\]
par définition de $\mathfrak{G}_{X/S}$, et on peut dans cet isomorphisme
remplacer $X$ par un ouvert quelconque $U$ de $X$; on peut donc dire que le
\emph{faisceau tangent} de $X$ relativement à $S$ est isomorphe au
\emph{faisceau des germes de $S$-sections} du fibré tangent de $X$
relativement à $S$. Si $f:X\to Y$ est un $S$-morphisme, on a vu (16.4.19)
qu'on a un homomorphisme canonique
$f_{X/Y/S}:f^*(\Omega_{Y/S}^1)\to\Omega_{X/S}^1$, ce qui donne, compte tenu de
ce que
\[
  \mathbf{V}(f^*(\Omega_{Y/S}^1))
  =\mathbf{V}(\Omega_{Y/S}^1)\times_Y X\qquad\text{(II, 1.7.11),}
\]
un $X$-morphisme
$T_{X/S}(f):T_{X/S}\to T_{Y/S}\times_Y X$. Si $g:Y\to Z$ est un second
$S$-morphisme, on a
$T_{X/S}(g\circ f)=(T_{Y/S}(g)\times1_X)\circ T_{X/S}(f)$ (0, 20.5.4.1).

Il résulte de (16.5.10.1) et de (II, 1.7.11) que pour tout changement de base
$g:S'\to S$, on a un isomorphisme canonique
\[
  T_{X'/S'}\xrightarrow{\sim}T_{X/S}\times_S S'
  =T_{X/S}\times_X X'.
  \tag{16.5.12.2}
  \label{IV.16.5.12.2}
\]
\end{env}

\begin{env}[16.5.13]
\EGAReaderTarget{ega4.c04.statement.16-5-13}{IV.16.5.13}
Pour tout point $x\in X$, on appelle \emph{espace tangent à $X$ au point $x$}
(relativement à $S$) l'ensemble des points de la fibre
$T_{X/S}\times_X\operatorname{Spec}(k(x))$ \emph{rationnels sur $k(x)$}, donc
l'ensemble
\[
  T_{X/S}(x)=
  \operatorname{Hom}_{k(x)}
  (\Omega_{X/S}^1\otimes_{\mathcal{O}_x}k(x),k(x))
  \tag{16.5.13.1}
  \label{IV.16.5.13.1}
\]
qui est le \emph{dual} du $k(x)$-espace vectoriel
$\Omega_{\mathcal{O}_x/\mathcal{O}_s}^1/
\mathfrak{m}_x\cdot\Omega_{\mathcal{O}_x/\mathcal{O}_s}^1$. Lorsque
$\Omega_{X/S}^1$ est un $\mathcal{O}_X$-Module de \emph{type fini},
$T_{X/S}(x)$ est donc un espace vectoriel de rang fini sur $k(x)$, et pour
tout
\oldpage[IV-4]{31}
changement de base $g:S'\to S$, et tout point
$x'\in X'=X\times_S S'$ au-dessus de $x$, on a un isomorphisme canonique
\[
  T_{X'/S'}(x')\xrightarrow{\sim}
  T_{X/S}(x)\otimes_{k(x)}k(x').
  \tag{16.5.13.2}
  \label{IV.16.5.13.2}
\]

Si $x$ est \emph{rationnel sur $k(s)$}, où $s=f(x)$ (de sorte que
$k(s)\to k(x)$ est un isomorphisme), il résulte de (16.4.13) que l'on a un
isomorphisme canonique
\[
  T_{X/S}(x)=T_{X_s/k(s)}(x)
  =\operatorname{Hom}_{k(s)}
  (\mathfrak{m}'_x/\mathfrak{m}'_x{}^2,k(s))
  \tag{16.5.13.3}
  \label{IV.16.5.13.3}
\]
où $\mathfrak{m}'_x$ est l'idéal maximal de
$\mathcal{O}_{X_s,x}=\mathcal{O}_{X,x}/
\mathfrak{m}_s\mathcal{O}_{X,x}$. Dans le cas où $S$ est le spectre d'un corps
$k$, on retrouve ainsi la définition de Zariski de l'espace tangent en un
point $x\in X$ \emph{rationnel sur $k$}, comme dual de
$\mathfrak{m}_x/\mathfrak{m}_x^2$.

Soit $Y$ un second $S$-préschéma et soit $g:Y\to X$ un $S$-morphisme; on a
alors un homomorphisme canonique de $\mathcal{O}_Y$-Modules (16.4.19)
\[
  g_{Y/X/S}:g^*(\Omega_{X/S}^1)\longrightarrow\Omega_{Y/S}^1.
  \tag{16.5.13.4}
  \label{IV.16.5.13.4}
\]
Notons maintenant que si $y\in Y$ et $x=g(y)$, on a
\[
  g^*(\Omega_{X/S}^1)\otimes_{\mathcal{O}_Y}k(y)
  =(\Omega_{X/S}^1\otimes_{\mathcal{O}_X}k(x))
  \otimes_{k(x)}k(y)
\]
et par suite, si $\Omega_{X/S}^1$ est un $\mathcal{O}_X$-Module
\emph{de type fini}, on peut identifier
\[
  \operatorname{Hom}_{k(y)}
  (g^*(\Omega_{X/S}^1)\otimes_{\mathcal{O}_Y}k(y),k(y))
\]
à $T_{X/S}(x)\otimes_{k(x)}k(y)$. On déduit donc alors de l'homomorphisme
(16.5.13.4) un homomorphisme de \emph{$k(y)$-espaces vectoriels}
\[
  T_y(g):T_{Y/S}(y)\longrightarrow
  T_{X/S}(x)\otimes_{k(x)}k(y)
  \tag{16.5.13.5}
  \label{IV.16.5.13.5}
\]
dit \emph{application linéaire tangente à $g$ au point $y$}. Lorsque $y$ est
\emph{rationnel sur $k(s)$}, on peut identifier $k(s)$, $k(y)$ et $k(x)$ et
$T_y(g)$ est alors un homomorphisme de $k(s)$-espaces vectoriels
$T_{Y/S}(y)\to T_{X/S}(x)$; on notera d'ailleurs que dans ce cas,
$g^*(\Omega_{X/S}^1)\otimes_{\mathcal{O}_Y}k(y)$ s'identifie à
$\Omega_{X/S}^1\otimes_{\mathcal{O}_X}k(x)$, et l'homomorphisme précédent est
donc défini sans aucune condition de finitude sur $\Omega_{X/S}^1$ et n'est
autre que l'homomorphisme $T_{Y/S}(g)$ (16.5.12) restreint à la fibre au point
$y$ de $T_{Y/S}$.
\end{env}

\begin{env}[16.5.14]
\label{IV.16.5.14}
L'interprétation des dérivations d'une $A$-algèbre $B$ dans un $B$-module
$L$, donnée dans (0, 20.1.1), se traduit dans le langage des préschémas de la
façon suivante.

Considérons deux morphismes de préschémas $f:X\to S$, $g:Y\to S$, et un
sous-préschéma fermé $Y_0$ de $Y$ défini par un Idéal de carré nul
$\mathscr{J}$ de $\mathcal{O}_Y$ (de sorte que $Y$ et $Y_0$ ont
\emph{même espace topologique sous-jacent}). Supposons donné un $S$-morphisme
$u_0:Y_0\to X$, de sorte qu'on a un diagramme commutatif
\[
  \xymatrix{
    X\ar[d]_-f & Y_0\ar[l]_-{u_0}\ar[d]^j\\
    S & Y\ar@{-->}[ul]_-u\ar[l]^-g
  }
  \tag{16.5.14.1}
  \label{IV.16.5.14.1}
\]
\oldpage[IV-4]{32}
et proposons-nous de chercher s'il existe des \emph{$S$-morphismes}
$u:Y\to X$ tels que $u_0=u\circ j$ (autrement dit, s'il est possible de
compléter le diagramme précédent par la flèche pointillée $u$, de façon à le
laisser \emph{commutatif}).

Pour cela, considérons un ouvert affine $U=\operatorname{Spec}(C)$ de $Y$;
son image réciproque $j^{-1}(U)$ est l'ouvert affine
$U_0=\operatorname{Spec}(C/\mathfrak{L})$, où
$\mathfrak{L}=\Gamma(U,\mathscr{J})$, idéal \emph{de carré nul} dans $C$;
nous supposerons $U$ assez petit pour que $u_0(U_0)$ soit contenu dans un
ouvert affine $V=\operatorname{Spec}(B)$ de $X$, et
$g(U)=f(u_0(U_0))$ contenu dans un ouvert affine
$W=\operatorname{Spec}(A)$ de $S$, de sorte que $B$ et $C$ sont des
$A$-algèbres et que $u_0|U_0$ correspond à un $A$-homomorphisme $\psi$ de
$B$ dans $C/\mathfrak{L}$; soit $P(U_0)$ l'ensemble des restrictions $u|U$
des homomorphismes cherchés, qui correspondent canoniquement aux
$A$-homomorphismes d'algèbres $\varphi:B\to C$ tels que le \emph{composé}
\[
  B\xrightarrow{\varphi}C\longrightarrow C/\mathfrak{L}
\]
\emph{soit égal à $\psi$}. On sait donc (0, 20.1.1) que l'ensemble de ces
homomorphismes est vide ou de la forme
$\varphi_1+\operatorname{D\acute{e}r}_A(B,\mathfrak{L})$; lorsque $P(U_0)$ n'est pas
vide, le groupe additif $\operatorname{D\acute{e}r}_A(B,\mathfrak{L})$ opère par
addition dans $P(U_0)$, qui est donc un \emph{espace affine} pour le groupe
additif $\operatorname{D\acute{e}r}_A(B,\mathfrak{L})$ (ou encore un
\emph{espace homogène principal} (ou \emph{torseur}) \emph{sous}
$\operatorname{D\acute{e}r}_A(B,\mathfrak{L})$).

Remarquons maintenant que, puisque $\mathfrak{L}$ est muni d'une structure de
$B$-module au moyen de $\psi$, on a un \emph{isomorphisme}
$v\mapsto v\circ d_{B/A}$ de
$\operatorname{Hom}_B(\Omega_{B/A}^1,\mathfrak{L})$ sur
$\operatorname{D\acute{e}r}_A(B,\mathfrak{L})$ (0, 20.4.8). D'ailleurs, comme
$\mathfrak{L}$ est de carré nul, donc un $(C/\mathfrak{L})$-module, tout
$B$-homomorphisme $v:\Omega_{B/A}^1\to\mathfrak{L}$ peut être considéré
comme un $(C/\mathfrak{L})$-homomorphisme
$\Omega_{B/A}^1\otimes_B(C/\mathfrak{L})\to\mathfrak{L}$. Comme
$\mathscr{J}$ est de carré nul, il peut être considéré comme un
$\mathcal{O}_{Y_0}$-Module quasi-cohérent; introduisons le
$\mathcal{O}_{Y_0}$-Module
\[
  \mathscr{G}=\mathscr{H}om(u_0^*(\Omega_{X/S}^1),\mathscr{J});
  \tag{16.5.14.2}
  \label{IV.16.5.14.2}
\]
il résulte alors du fait que
$\Omega_{B/A}^1=\Gamma(V,\Omega_{X/S}^1)$ (16.3.7) que l'on peut écrire
$\operatorname{D\acute{e}r}_A(B,\mathfrak{L})=\Gamma(U_0,\mathscr{G})$.

Comme $P(U_0)$ est défini comme ensemble de $S$-morphismes $U\to X$, il est
clair que $U_0\mapsto P(U_0)$ est un \emph{faisceau d'ensembles}
$\mathscr{P}$ sur $Y_0$. Utilisons ce fait pour prouver que l'application
\[
  h:\Gamma(U_0,\mathscr{G})\times P(U_0)\longrightarrow P(U_0)
\]
définissant la structure de torseur sur $P(U_0)$ est aussi indépendante du
choix de $V$ et $W$, et en outre que, si $U'\subset U$ est un second ouvert
affine de $Y$, $U'_0$ son image réciproque dans $Y_0$, le diagramme
\[
  \xymatrix{
    \Gamma(U_0,\mathscr{G})\times P(U_0)\ar[d]\ar[r]^h
      & P(U_0)\ar[d]\\
    \Gamma(U'_0,\mathscr{G})\times P(U'_0)\ar[r]_{h'}
      & P(U'_0)
  }
  \tag{16.5.14.3}
  \label{IV.16.5.14.3}
\]
est commutatif (les flèches verticales étant les opérateurs de restriction).
En vertu de la remarque précédente, on est ramené à prouver la commutativité
du diagramme précédent lorsque $h$ est défini comme ci-dessus à partir des
ouverts affines $V$, $W$ et $h'$ à partir d'ouverts
\oldpage[IV-4]{33}
affines $V'\subset V$ et $W'\subset W$. Mais en vertu de la description
précédente de $h$, cela résulte de la commutativité du diagramme (0, 20.5.4.2).

Les applications
$\Gamma(U_0,\mathscr{G})\times P(U_0)\to P(U_0)$ définissent donc un
homomorphisme de \emph{faisceaux d'ensembles}
$m:\mathscr{G}\times\mathscr{P}\to\mathscr{P}$ tel que, pour tout ouvert
$U_0$ pour lequel $\Gamma(U_0,\mathscr{P})\neq\varnothing$,
\[
  m_{U_0}:\Gamma(U_0,\mathscr{G})\times\Gamma(U_0,\mathscr{P})
  \longrightarrow\Gamma(U_0,\mathscr{P})
\]
soit une loi externe définissant sur $\Gamma(U_0,\mathscr{P})$ une structure
de torseur sous le groupe $\Gamma(U_0,\mathscr{G})$.
\end{env}

\begin{env}[16.5.15]
\EGAReaderTarget{ega4.c04.statement.16-5-15}{IV.16.5.15}
D'une façon générale, lorsqu'on se donne sur un espace topologique $Z$ un
faisceau d'ensembles $\mathscr{P}$, un faisceau de groupes $\mathscr{G}$ (non
nécessairement commutatif) et un homomorphisme de faisceaux d'ensembles
$m:\mathscr{G}\times\mathscr{P}\to\mathscr{P}$ tel que, pour tout ouvert
$U\subset Z$ tel que $\Gamma(U,\mathscr{P})\neq\varnothing$,
\[
  m_U:\Gamma(U,\mathscr{G})\times\Gamma(U,\mathscr{P})
  \longrightarrow\Gamma(U,\mathscr{P})
\]
fasse de $\Gamma(U,\mathscr{P})$ un \emph{torseur} sous le groupe
$\Gamma(U,\mathscr{G})$, on dit que $\mathscr{P}$ est un
\emph{pseudo-torseur} (ou \emph{faisceau formellement principal homogène})
sous le faisceau de groupes $\mathscr{G}$. On dit que $\mathscr{P}$ est un
\emph{torseur} (ou \emph{faisceau principal homogène}) sous $\mathscr{G}$ si
de plus $\Gamma(U,\mathscr{P})\neq\varnothing$ pour tout ouvert
$U\neq\varnothing$ d'une base convenable de la topologie de $Z$.

Pour la théorie générale des torseurs, nous renvoyons à [42]; nous nous
bornerons ici à rappeler la correspondance canonique entre les classes
d'isomorphie de ces torseurs (pour un $\mathscr{G}$ donné) et les éléments de
l'ensemble de cohomologie $H^1(Z,\mathscr{G})$. Considérons en effet un
torseur $\mathscr{P}$ sous $\mathscr{G}$ et un recouvrement ouvert
$(U_\lambda)$ de $Z$ tel que
$\Gamma(U_\lambda,\mathscr{P})\neq\varnothing$ pour tout $\lambda$;
désignons par $p_\lambda$ un élément de
$\Gamma(U_\lambda,\mathscr{P})$. Pour tout couple d'indices $\lambda$, $\mu$
tel que $U_\lambda\cap U_\mu\neq\varnothing$, il existe alors un élément et
un seul $\gamma_{\lambda\mu}$ de
$\Gamma(U_\lambda\cap U_\mu,\mathscr{G})$ tel que
\[
  \gamma_{\lambda\mu}\cdot
  (p_\mu|U_\lambda\cap U_\mu)=p_\lambda|U_\lambda\cap U_\mu;
\]
en outre, si $\lambda$, $\mu$, $\nu$ sont trois indices tels que
$U_\lambda\cap U_\mu\cap U_\nu\neq\varnothing$, les restrictions
$\gamma'_{\lambda\mu}$, $\gamma'_{\mu\nu}$, $\gamma'_{\lambda\nu}$ de
$\gamma_{\lambda\mu}$, $\gamma_{\mu\nu}$, $\gamma_{\lambda\nu}$ à
$U_\lambda\cap U_\mu\cap U_\nu$ vérifient la condition
$\gamma'_{\lambda\nu}=\gamma'_{\lambda\mu}\gamma'_{\mu\nu}$; autrement dit
$(\lambda,\mu)\mapsto\gamma_{\lambda\mu}$ est un \emph{1-cocycle} du
recouvrement $(U_\lambda)$ à valeurs dans $\mathscr{G}$. Si, pour tout
$\lambda$, $p'_\lambda$ est un second élément de
$\Gamma(U_\lambda,\mathscr{P})$, il existe un élément et un seul
$\beta_\lambda\in\Gamma(U_\lambda,\mathscr{G})$ tel que
$p'_\lambda=\beta_\lambda\cdot p_\lambda$ et le 1-cocycle
$(\gamma'_{\lambda\mu})$ correspondant à la famille $(p'_\lambda)$ est donné
par
\[
  \gamma'_{\lambda\mu}=\beta_\lambda\gamma_{\lambda\mu}\beta_\mu^{-1},
\]
autrement dit, est \emph{cohomologue} à $(\gamma_{\lambda\mu})$.
Inversement, la donnée d'un 1-cocycle $(\gamma_{\lambda\mu})$ définit, pour
tout couple $(\lambda,\mu)$, un automorphisme $\theta_{\lambda\mu}$ du
faisceau \emph{d'ensembles}
$\mathscr{G}|U_\lambda\cap U_\mu$, savoir la translation à droite par
$\gamma_{\lambda\mu}$, et le fait qu'il s'agisse d'un cocycle montre que l'on
peut \emph{recoller} les faisceaux d'ensembles $\mathscr{G}|U_\lambda$ à
l'aide des automorphismes $\theta_{\lambda\mu}$ (0, 3.3.1); on obtient ainsi
de façon évidente un torseur sous $\mathscr{G}$, soit $\mathscr{P}$, et si
l'on prend pour $p_\lambda$ la section unité au-dessus de $U_\lambda$, le
1-cocycle correspondant n'est autre que le 1-cocycle donné
$(\gamma_{\lambda\mu})$; en outre, si l'on remplace
$(\gamma_{\lambda\mu})$ par un 1-cocycle
$\gamma'_{\lambda\mu}=\beta_\lambda\gamma_{\lambda\mu}\beta_\mu^{-1}$ qui
lui est cohomologue, on vérifie aussitôt que le torseur obtenu est isomorphe
à $\mathscr{P}$.

En particulier, si $(\gamma_{\lambda\mu})$ est un \emph{1-cobord}, autrement
dit de la forme
$\gamma_{\lambda\mu}=\beta_\lambda\beta_\mu^{-1}$, le torseur obtenu
$\mathscr{P}$ est \emph{isomorphe à $\mathscr{G}$} (considéré comme torseur
sous lui-même pour les translations à gauche); on dit dans ce cas que
$\mathscr{P}$ est \emph{trivial}, et la réciproque est évidente.

Plus particulièrement, il résulte de (III, 1.3.1) que l'on a:
\end{env}

\begin{proposition}[16.5.16]
\label{IV.16.5.16}
Soient $Z$ un schéma affine, $\mathscr{G}$ un $\mathcal{O}_Z$-Module
quasi-cohérent; alors tout torseur sous $\mathscr{G}$ est trivial.
\end{proposition}

\oldpage[IV-4]{34}
Revenant au problème considéré en (16.5.13), on obtient donc:

\begin{proposition}[16.5.17]
\label{IV.16.5.17}
Soient $X$, $Y$ deux $S$-préschémas, $Y_0$ un sous-préschéma fermé de $Y$
défini par un Idéal quasi-cohérent $\mathscr{J}$ de $\mathcal{O}_Y$ tel que
$\mathscr{J}^2=0$, $j:Y_0\to Y$ l'injection canonique. Soit
$u_0:Y_0\to X$ un $S$-morphisme, et $\mathscr{P}$ le faisceau d'ensembles sur
$Y$ tel que, pour tout ouvert $U$ de $Y$, $\Gamma(U,\mathscr{P})$ soit
l'ensemble des $S$-morphismes $u:U\to X$ tels que
$u_0|U_0=u\circ(j|U_0)$, où $U_0=j^{-1}(U)$. Alors il existe sur
$\mathscr{P}$ une structure de pseudo-torseur sous le
$\mathcal{O}_{Y_0}$-Module
\[
  \mathscr{G}=\mathscr{H}om_{\mathcal{O}_{Y_0}}
  (u_0^*(\Omega_{X/S}^1),\mathscr{J}).
\]
\end{proposition}

En particulier:

\begin{corollary}[16.5.18]
\label{IV.16.5.18}
Avec les notations de (16.5.16), supposons $Y$ affine et $\Omega_{X/S}^1$ de
présentation finie; s'il existe un recouvrement ouvert $(U_\alpha)$ de $Y$,
et, pour chaque indice $\alpha$, un $S$-morphisme $v_\alpha:U_\alpha\to X$
tel que, si $U_\alpha^0=j^{-1}(U_\alpha)$, on ait
$v_\alpha\circ(j|U_\alpha^0)=u_0|U_\alpha^0$, alors il existe un
$S$-morphisme $u:Y\to X$ tel que $u\circ j=u_0$.
\end{corollary}

En effet, $\mathscr{G}$ est alors un $\mathcal{O}_{Y_0}$-Module
quasi-cohérent (I, 1.3.12); en vertu de (16.5.16) et du fait que $Y_0$ est
alors affine, le faisceau $\mathscr{P}$, qui est par hypothèse un torseur sous
$\mathscr{G}$, et non seulement un pseudo-torseur, est \emph{trivial}; mais si
$w$ est un isomorphisme de $\mathscr{G}$ sur $\mathscr{P}$ (en tant que
torseurs sous $\mathscr{G}$), l'image par $w$ de la section 0 de
$\mathscr{G}$ est le $S$-morphisme $u$ cherché.

\subsection{\texorpdfstring{Faisceaux de $p$-différentielles et différentielle extérieure.}{Faisceaux de p-différentielles et différentielle extérieure.}}
\EGAReaderTarget{ega4.c04.subsec.16-6}{IV.16.6}

\begin{env}[16.6.1]
\EGAReaderTarget{ega4.c04.statement.16-6-1}{IV.16.6.1}
Soit $f:X\to S$ un morphisme de préschémas. On appelle \emph{faisceau des
$p$-différentielles de $X$ relativement à $S$} ($p$ entier) la
\emph{puissance extérieure $p$-ème} (0, 4.1.5) du $\mathcal{O}_X$-Module
$\Omega_{X/S}^1$, notée
\[
  \Omega_{X/S}^p=\bigwedge^p(\Omega_{X/S}^1).
  \tag{16.6.1.1}
  \label{IV.16.6.1.1}
\]
On a donc $\Omega_{X/S}^0=\mathcal{O}_X$, et $\Omega_{X/S}^p=0$ pour $p<0$;
les $\Omega_{X/S}^p$ sont les composants homogènes de l'Algèbre extérieure
sur $\Omega_{X/S}^1$
\[
  \Omega_{X/S}^{\bullet}=\bigwedge(\Omega_{X/S}^1)
  =\bigoplus_{p\in\mathbf{Z}}\bigwedge^p(\Omega_{X/S}^1),
  \tag{16.6.1.2}
  \label{IV.16.6.1.2}
\]
qui est donc une $\mathcal{O}_X$-Algèbre graduée quasi-cohérente,
anticommutative et dont les éléments de degré 1 sont de carré nul. Pour tout
ouvert affine $U$ de $X$, on a
\[
  \Gamma(U,\Omega_{X/S}^{\bullet})
  =\bigwedge(\Gamma(U,\Omega_{X/S}^1)),
\]
où $\Gamma(U,\Omega_{X/S}^1)$ est considéré comme
$\Gamma(U,\mathcal{O}_X)$-module.

Lorsque $S=\operatorname{Spec}(A)$ et $X=\operatorname{Spec}(B)$ sont
affines, $B$ étant donc une $A$-algèbre, on a (0, 4.1.5)
$\Omega_{X/S}^p=(\Omega_{B/A}^p)^{\sim}$, en posant
$\Omega_{B/A}^p=\bigwedge^p\Omega_{B/A}^1$.
\end{env}

\begin{theorem}[16.6.2]
\label{IV.16.6.2}
Il existe un endomorphisme et un seul $d$ du faisceau de groupes additifs
$\Omega_{X/S}^{\bullet}$, ayant les propriétés suivantes:
\begin{enumerate}
\item[(i)] $d\circ d=0$.
\item[(ii)] Pour tout ouvert $U$ de $X$ et toute section
$f\in\Gamma(U,\mathcal{O}_X)$, on a $df=d_{X/S}f$.
\oldpage[IV-4]{35}
\item[(iii)] Pour tout ouvert $U$ de $X$, tout couple d'entiers $p$, $q$ et
tout couple de sections
$\omega'_p\in\Gamma(U,\Omega_{X/S}^p)$,
$\omega''_q\in\Gamma(U,\Omega_{X/S}^q)$, on a
\[
  d(\omega'_p\wedge\omega''_q)
  =(d\omega'_p)\wedge\omega''_q
  +(-1)^p\omega'_p\wedge d\omega''_q.
  \tag{16.6.2.1}
  \label{IV.16.6.2.1}
\]
\end{enumerate}
En outre, $d$ est un endomorphisme de $\psi^*(\mathcal{O}_X)$-Modules
gradués, de degré $+1$.
\end{theorem}

Supposons prouvée l'existence de l'endomorphisme $d$. Pour tout ouvert affine
$U$ de $X$, toute section de $\Omega_{X/S}^p$ au-dessus de $U$ est (en vertu
de (i)) combinaison linéaire d'un nombre fini d'éléments de la forme
$g(df_1\wedge df_2\wedge\cdots\wedge df_p)$, où $g$ et les $f_i$ sont des
sections de $\mathcal{O}_X$ au-dessus de $U$ (0, 20.4.7). Les conditions (i)
et (iii) montrent alors, par récurrence sur $p$, que l'on a nécessairement
\[
  d(g(df_1\wedge df_2\wedge\cdots\wedge df_p))
  =dg\wedge df_1\wedge df_2\wedge\cdots\wedge df_p.
  \tag{16.6.2.2}
  \label{IV.16.6.2.2}
\]
Cela prouve donc l'\emph{unicité} de $d$ et la dernière assertion du théorème.
En vertu de cette propriété d'unicité, pour démontrer l'existence de $d$, on
peut se borner au cas où $S=\operatorname{Spec}(A)$ et
$X=\operatorname{Spec}(B)$ sont affines. Or (Bourbaki, \emph{Alg.}, chap.
III, 3\textsuperscript{e} éd., \S~10) pour définir une $A$-antidérivation
$D$ de degré $+1$ d'une algèbre extérieure $\bigwedge(M)$ (où $M$ est un
$B$-module et $B$ une $A$-algèbre), cette antidérivation prenant ses valeurs
dans une $A$-algèbre \emph{graduée anticommutative}
$C=\bigoplus_{n=0}^{\infty}C_n$, dont les éléments de degré 1 sont de carré
nul, il suffit de se donner arbitrairement une $A$-dérivation $D_0$ de $B$
dans $C_1$ et un $A$-homomorphisme $D_1$ de $M$ dans $C_2$; il existe alors
une $A$-antidérivation $D$ et une seule de $\bigwedge(M)$ dans $C$
coïncidant avec $D_0$ dans $B$ et avec $D_1$ dans $M$.

Dans le cas actuel, $D_0$ est nécessairement égal à $d_{B/A}$ en vertu de
(ii); tout revient à voir, compte tenu de (16.6.2.2), qu'il y a un
$A$-homomorphisme $u$ de $\Omega_{B/A}^1$ dans $\Omega_{B/A}^2$ tel que
l'on ait
\[
  u(g.df)=dg\wedge df
  \tag{16.6.2.3}
  \label{IV.16.6.2.3}
\]
quels que soient $f$, $g$ dans $A$; il suffira pour cela de montrer qu'il
existe un $A$-homomorphisme
$v:B\otimes_A\Omega_{B/A}^1\to\Omega_{B/A}^2$ tel que
\[
  v(g.\omega)=dg\wedge\omega
  \tag{16.6.2.4}
  \label{IV.16.6.2.4}
\]
pour $g\in B$ et $\omega\in\Omega_{B/A}^1$. Finalement, comme
$\Omega_{B/A}^1=\mathfrak{I}/\mathfrak{I}^2$ (où
$\mathfrak{I}=\mathfrak{I}_{B/A}$ est le noyau de l'homomorphisme canonique
$B\otimes_A B\to B$) et que $\Omega_{B/A}^1$ est engendré par les éléments
de la forme $g.df$, il suffit de définir un $A$-homomorphisme
$w:B\otimes_A(B\otimes_A B)\to\Omega_{B/A}^2$ tel que
\[
  w(g'\otimes g\otimes f)=dg'\wedge(g.df)
  \tag{16.6.2.5}
  \label{IV.16.6.2.5}
\]
et tel que $w$ soit nul dans l'image de $B\otimes_A\mathfrak{I}^2$. Or,
comme le second membre de (16.6.2.5) est $A$-trilinéaire en $g'$, $g$, $f$,
l'existence de $w$ vérifiant (16.6.2.5) est immédiate. Comme, d'autre part,
$\mathfrak{I}$ est engendré par les éléments $1\otimes x-x\otimes1$
($x\in B$), on est ramené à vérifier que lorsque
$z=(1\otimes x-x\otimes1)(1\otimes y-y\otimes1)$, on a
$w(g'\otimes z)=0$. Or, comme
$z=1\otimes(xy)+(xy)\otimes1-x\otimes y-y\otimes x$, la formule (16.6.2.4)
montre qu'il
\oldpage[IV-4]{36}
suffit de voir que l'on a $d(xy)-x.dy-y.dx=0$, ce qui exprime que $d$ est une
dérivation.

Il reste à prouver que $d$ vérifie la condition (i). Or, le carré d'une
antidérivation est une dérivation (Bourbaki, \emph{loc. cit.}), et comme
$\Omega_{B/A}^{\bullet}$ est engendré par $\Omega_{B/A}^1$ comme
$B$-algèbre, il suffit de vérifier que $d(dz)=0$ pour $z\in B$ et
$z\in\Omega_{B/A}^1$; dans le premier cas, cela résulte de la formule
(16.6.2.3) avec $g=1$; dans le second, on peut se borner au cas où $z=g.df$
avec $f$, $g$ dans $B$, et alors, en vertu de (16.6.2.1) et (16.6.2.3), on a
\[
  d(d(g.df))=d(dg\wedge df)
  =(d(dg))\wedge(df)-(dg)\wedge(d(df))=0.
\]
\hfill C.Q.F.D.

\begin{definition}[16.6.3]
\EGAReaderTarget{ega4.c04.def.16-6-3}{IV.16.6.3}
L'antidérivation $d$ définie dans (16.6.2) (aussi notée $d_{X/S}$) est
appelée \emph{différentielle extérieure sur $X$} (relativement à $S$).
\end{definition}

\begin{proposition}[16.6.4]
\label{IV.16.6.4}
Pour tout changement de base $g:S'\to S$, si l'on pose
$X'=X\times_S S'$, l'homomorphisme canonique
\[
  \Omega_{X/S}^{\bullet}\otimes_S S'
  \longrightarrow\Omega_{X'/S'}^{\bullet}
  \tag{16.6.4.1}
  \label{IV.16.6.4.1}
\]
déduit de l'isomorphisme (16.5.9.1) est bijectif. En outre, si $s$ est une
section de $\Omega_{X/S}^{\bullet}$ au-dessus d'un ouvert $U$ de $X$,
$s\otimes1$ son image réciproque, section de $\Omega_{X'/S'}^{\bullet}$
au-dessus de l'image réciproque $U'$ de $U$ dans $X'$, on a
$d_{X'/S'}(s\otimes1)=d_{X/S}(s)\otimes1$.
\end{proposition}

La première assertion est immédiate, la formation de l'algèbre extérieure
d'un module permutant avec toute extension de l'anneau des scalaires. Pour
prouver la seconde, on peut, en vertu de (16.6.2.2), se borner au cas où
$s\in\Gamma(U,\mathcal{O}_X)$, et dans ce cas l'assertion a déjà été prouvée
(16.4.3.7).

\begin{env}[16.6.5]
\label{IV.16.6.5}
Supposons que $\Omega_{X/S}^1$ soit un $\mathcal{O}_X$-Module localement libre
de rang $n$ en un point $x$, de sorte qu'il y a $n$ sections
$s_i\in\Gamma(U,\mathcal{O}_X)$ telles que les $ds_i$ forment une base du
$\Gamma(U,\mathcal{O}_X)$-module $\Gamma(U,\Omega_{X/S}^1)$ (16.5.8).
Alors, pour tout entier $p\geq1$, les $p$-différentielles
$ds_{i_1}\wedge ds_{i_2}\wedge\cdots\wedge ds_{i_p}$ (pour
$i_1<i_2<\cdots<i_p$, éléments de $[1,n]$) forment une base de
$\binom{n}{p}$ éléments de $\Gamma(U,\Omega_{X/S}^p)$ sur
$\Gamma(U,\mathcal{O}_X)$. En outre la formule (16.6.2.2) montre que pour
toute section $g\in\Gamma(U,\mathcal{O}_X)$, on a
\[
  d(g.ds_{i_1}\wedge ds_{i_2}\wedge\cdots\wedge ds_{i_p})
  =\sum_k(-1)^r\frac{\partial g}{\partial s_k}
  ds_{i_1}\wedge\cdots\wedge ds_{i_r}\wedge ds_k\wedge
  ds_{i_{r+1}}\wedge\cdots\wedge ds_{i_p}
  \tag{16.6.5.1}
  \label{IV.16.6.5.1}
\]
où, au second membre, $k$ parcourt l'ensemble des $n-p$ indices distincts des
$i_h$, $i_r$ étant le plus grand indice $<k$.

On notera que la relation $d(dg)=0$ pour toute section
$g\in\Gamma(U,\mathcal{O}_X)$ s'exprime sous la forme
\[
  D_i(D_jg)=D_j(D_ig)\qquad\text{pour $i\neq j$};
\]
autrement dit, les dérivations $D_i$ définies dans (16.5.7) sont deux à deux
permutables.
\end{env}

\subsection{\texorpdfstring{Les $\mathscr{P}_{X/S}^n(\mathscr{F})$.}{Les P\string_\string{X/S\string}\string^n(F).}}
\EGAReaderTarget{ega4.c04.subsec.16-7}{IV.16.7}

\begin{env}[16.7.1]
\label{IV.16.7.1}
Soient $f:X\to S$ un morphisme de préschémas, $\mathscr{F}$ un
$\mathcal{O}_X$-Module. Désignons par $X_{\Delta_f}^{(n)}$ le \emph{$n$-ème
voisinage infinitésimal de $X$} pour le morphisme diagonal
\oldpage[IV-4]{37}
$\Delta_f:X\to X\times_S X$, par
$h_n:X_{\Delta_f}^{(n)}\to X\times_S X$ le morphisme canonique (16.1.2), et
considérons les deux morphismes composés
\[
  p_1^{(n)}:X_{\Delta_f}^{(n)}\xrightarrow{h_n}X\times_S X
  \xrightarrow{p_1}X,qquad
  p_2^{(n)}:X_{\Delta_f}^{(n)}\xrightarrow{h_n}X\times_S X
  \xrightarrow{p_2}X
\]
de sorte que, par définition, $p_1^{(n)}$ correspond à l'homomorphisme de
faisceaux d'anneaux $\mathcal{O}_X\to\mathscr{P}_{X/S}^n$ que nous avons
choisi pour définir la structure de $\mathcal{O}_X$-Algèbre sur
$\mathscr{P}_{X/S}^n$ (16.3.5), et $p_2^{(n)}$ à l'homomorphisme de faisceaux
d'anneaux $d_{X/S}^n:\mathcal{O}_X\to\mathscr{P}_{X/S}^n$ (16.3.6). Comme
$X_{\Delta_f}^{(n)}$ et $X$ ont même espace sous-jacent, on peut écrire
\[
  \mathscr{P}_{X/S}^n
  =(p_1^{(n)})_*((p_2^{(n)})^*(\mathcal{O}_X)).
  \tag{16.7.1.1}
  \label{IV.16.7.1.1}
\]
Plus généralement, nous poserons
\[
  \mathscr{P}_{X/S}^n(\mathscr{F})
  =(p_1^{(n)})_*((p_2^{(n)})^*(\mathscr{F}))
  \tag{16.7.1.2}
  \label{IV.16.7.1.2}
\]
de sorte que
$\mathscr{P}_{X/S}^n=\mathscr{P}_{X/S}^n(\mathcal{O}_X)$; par définition,
$\mathscr{P}_{X/S}^n(\mathscr{F})$ est donc un $\mathcal{O}_X$-Module.
\end{env}

\begin{env}[16.7.2]
\label{IV.16.7.2}
Si l'on revient aux définitions des images réciproques de Modules sur les
espaces annelés (0, 4.3.1) et que l'on tienne compte de ce que
$X_{\Delta_f}^{(n)}$ et $X$ ont même espace sous-jacent, on voit que l'on peut
aussi écrire la définition (16.7.1.2) sous la forme
\[
  \mathscr{P}_{X/S}^n(\mathscr{F})
  =\mathscr{P}_{X/S}^n\otimes_{\mathcal{O}_X}\mathscr{F}
  \tag{16.7.2.1}
  \label{IV.16.7.2.1}
\]
mais où il faut prendre garde que, dans l'interprétation du signe $\otimes$,
$\mathscr{P}_{X/S}^n$ est muni de sa structure de $\mathcal{O}_X$-Module
définie par l'homomorphisme de faisceaux d'anneaux
$d_{X/S}^n:\mathcal{O}_X\to\mathscr{P}_{X/S}^n$. Il résulte aussitôt de cette
formule (ou directement de (16.7.1.2)) que
$\mathscr{P}_{X/S}^n(\mathscr{F})$ est canoniquement muni d'une structure de
$\mathscr{P}_{X/S}^n$-Module.
\end{env}

\begin{proposition}[16.7.3]
\label{IV.16.7.3}
\begin{enumerate}
\item[(i)] Le foncteur
$\mathscr{F}\mapsto\mathscr{P}_{X/S}^n(\mathscr{F})$ de la catégorie des
$\mathcal{O}_X$-Modules dans celle des $\mathscr{P}_{X/S}^n$-Modules est
exact à droite, et commute aux limites inductives quelconques; il est exact
lorsque $\mathscr{P}_{X/S}^n$ est un $\mathcal{O}_X$-Module plat.
\item[(ii)] Si $\mathscr{F}$ est un $\mathcal{O}_X$-Module quasi-cohérent
(resp. de type fini, resp. de présentation finie),
$\mathscr{P}_{X/S}^n(\mathscr{F})$ est un $\mathscr{P}_{X/S}^n$-Module
quasi-cohérent (resp. de type fini, resp. de présentation finie).
\end{enumerate}
\end{proposition}

Les assertions de (i) résultent aussitôt de la formule (16.7.2.1) et de la
considération de la symétrie de $\mathscr{P}_{X/S}^n$ (16.3.4). Les assertions
de (ii) découlent de l'exactitude à droite du foncteur
$\mathscr{F}\mapsto\mathscr{P}_{X/S}^n(\mathscr{F})$.

\begin{env}[16.7.4]
\label{IV.16.7.4}
Les deux structures de $\mathcal{O}_X$-Module de $\mathscr{P}_{X/S}^n$
définissent sur $\mathscr{P}_{X/S}^n(\mathscr{F})$ deux structures de
$\mathcal{O}_X$-Module, d'ailleurs permutables, donc une structure de
$\mathcal{O}_X$-Bimodule. Il est commode de noter \emph{à gauche} celle de
ces structures provenant de l'homomorphisme structural
$\mathcal{O}_X\to\mathscr{P}_{X/S}^n$ (choisi dans (16.3.5)) et
\emph{à droite} celle provenant de l'homomorphisme
$d_{X/S}^n:\mathcal{O}_X\to\mathscr{P}_{X/S}^n$. Autrement dit, pour tout
ouvert $U$ de $X$, et tout triplet d'éléments
$a\in\Gamma(U,\mathcal{O}_X)$, $b\in\Gamma(U,\mathscr{P}_{X/S}^n)$,
$t\in\Gamma(U,\mathscr{F})$, on a par définition
\[
  a(b\otimes t)=(ab)\otimes t,qquad
  (b\otimes t)a=(b.d^na)\otimes t=b\otimes(at)
  =(d^na).(b\otimes t).
  \tag{16.7.4.1}
  \label{IV.16.7.4.1}
\]
La structure de $\mathcal{O}_X$-Module provenant de la définition (16.7.1.2)
est donc, avec ces conventions, la structure de $\mathcal{O}_X$-Module à
gauche.

\oldpage[IV-4]{38}
Si $\mathscr{F}$ est un $\mathcal{O}_X$-Module quasi-cohérent, il en est de
même de $\mathscr{P}_{X/S}^n(\mathscr{F})$ pour l'une ou l'autre de ses
structures de $\mathcal{O}_X$-Module. Si de plus $\mathscr{F}$ est de type
fini (resp. de présentation finie) et $f:X\to S$ localement de type fini
(resp. localement de présentation finie),
$\mathscr{P}_{X/S}^n(\mathscr{F})$ est (pour l'une ou l'autre de ses
structures de $\mathcal{O}_X$-Module) de type fini (resp. de présentation
finie), comme il résulte de (16.3.9) et de (16.4.22).
\end{env}

\begin{env}[16.7.5]
\label{IV.16.7.5}
La définition (16.7.2.1) entraîne l'existence d'un homomorphisme de faisceaux
de groupes commutatifs
\[
  d_{X/S,\mathscr{F}}^n:\mathscr{F}
  \longrightarrow\mathscr{P}_{X/S}^n(\mathscr{F})
  \qquad\text{(aussi noté $d_{X/S}^n$)}
  \tag{16.7.5.1}
  \label{IV.16.7.5.1}
\]
tel qu'avec les notations de (16.7.4), on ait
\[
  d_{X/S,\mathscr{F}}^n(t)=1\otimes t
  \tag{16.7.5.2}
  \label{IV.16.7.5.2}
\]
et par suite, en vertu de (16.7.4.1)
\[
  d_{X/S,\mathscr{F}}^n(at)
  =(1\otimes t)a=(d_{X/S,\mathscr{F}}^n(t)).a
  \tag{16.7.5.3}
  \label{IV.16.7.5.3}
\]
\[
  d_{X/S,\mathscr{F}}^n(at)
  =(d_{X/S}^n(a)).(1\otimes t)
  =(d_{X/S}^n(a))(d_{X/S,\mathscr{F}}^n(t)).
  \tag{16.7.5.4}
  \label{IV.16.7.5.4}
\]
Il est donc $\mathcal{O}_X$-linéaire pour la structure de
$\mathcal{O}_X$-Module à droite sur $\mathscr{P}_{X/S}^n(\mathscr{F})$, et
\emph{semi-linéaire} (relativement à l'automorphisme $\sigma$ (16.3.4)) pour
la structure de $\mathcal{O}_X$-Module à gauche.
\end{env}

\begin{proposition}[16.7.6]
\label{IV.16.7.6}
Le $\mathcal{O}_X$-Module à gauche $\mathscr{P}_{X/S}^n(\mathscr{F})$ est
engendré par l'image de $\mathscr{F}$ par l'homomorphisme canonique
$d_{X/S,\mathscr{F}}^n$.
\end{proposition}

Cela résulte aussitôt de (16.7.5.3) et du cas particulier correspondant à
$\mathscr{F}=\mathcal{O}_X$ (16.3.8).

\begin{env}[16.7.7]
\label{IV.16.7.7}
Les homomorphismes canoniques de faisceaux d'anneaux
\[
  \varphi_{nm}:\mathscr{P}_{X/S}^m\longrightarrow\mathscr{P}_{X/S}^n
\]
pour $n\leq m$ (16.1.2) définissent, en vertu de (16.7.2.1) des
homomorphismes canoniques
\[
  \mathscr{P}_{X/S}^m(\mathscr{F})
  \longrightarrow\mathscr{P}_{X/S}^n(\mathscr{F})qquad(n\leq m)
\]
qui sont des homomorphismes de $\mathcal{O}_X$-Bimodules en vertu de (16.1.6)
et (16.7.4.1); en outre on a des diagrammes commutatifs
\[
  \xymatrix{
    \mathscr{P}_{X/S}^m(\mathscr{F})\ar[rr]
      &&\mathscr{P}_{X/S}^n(\mathscr{F})\\
    &\mathscr{F}\ar[ul]^-{d_{X/S,\mathscr{F}}^m}
      \ar[ur]_-{d_{X/S,\mathscr{F}}^n}
  }
\]
On a donc ainsi un système projectif de $\mathcal{O}_X$-Bimodules
$(\mathscr{P}_{X/S}^n(\mathscr{F}))$, et l'on pose
\[
  \mathscr{P}_{X/S}^{\infty}(\mathscr{F})
  =\varprojlim\mathscr{P}_{X/S}^n(\mathscr{F}).
  \tag{16.7.7.1}
  \label{IV.16.7.7.1}
\]
En outre, ce qui précède montre que les homomorphismes (16.7.5.1) forment un
système projectif d'homomorphismes, et définissent donc un homomorphisme
canonique
\[
  d_{X/S,\mathscr{F}}^{\infty}:\mathscr{F}
  \longrightarrow\mathscr{P}_{X/S}^{\infty}(\mathscr{F}).
  \tag{16.7.7.2}
  \label{IV.16.7.7.2}
\]
\end{env}

\oldpage[IV-4]{39}

\begin{env}[16.7.8]
\label{IV.16.7.8}
Soient $\mathscr{F}$, $\mathscr{G}$ deux $\mathcal{O}_X$-Modules; il résulte
aussitôt de la définition (16.7.2.1) que l'on a un isomorphisme canonique de
$\mathscr{P}_{X/S}^n$-Modules
\[
  \mathscr{P}_{X/S}^n(\mathscr{F}\otimes_{\mathcal{O}_X}\mathscr{G})
  \xrightarrow{\sim}
  \mathscr{P}_{X/S}^n(\mathscr{F})
  \otimes_{\mathscr{P}_{X/S}^n}\mathscr{P}_{X/S}^n(\mathscr{G})
  \tag{16.7.8.1}
  \label{IV.16.7.8.1}
\]
(Bourbaki, \emph{Alg.}, chap.~II, 3\textsuperscript{e} éd., §~5,
n\textsuperscript{o}~1, prop.~3).

On en conclut en particulier (ou l'on voit directement sur la définition
(16.7.2.1)) que si $\mathscr{F}$ est muni d'une structure de
$\mathcal{O}_X$-Algèbre (non nécessairement associative),
$\mathscr{P}_{X/S}^n(\mathscr{F})$ est canoniquement muni d'une structure de
$\mathscr{P}_{X/S}^n$-Algèbre; cette dernière est associative (resp.
commutative, resp. unitaire, resp. une Algèbre de Lie) lorsqu'il en est ainsi
de $\mathscr{F}$. En outre les homomorphismes canoniques
$\mathscr{P}_{X/S}^m(\mathscr{F})\to\mathscr{P}_{X/S}^n(\mathscr{F})$ pour
$n\leq m$ (16.7.7) sont alors des di-homomorphismes d'Algèbres; de même
(16.7.5.1) est alors un homomorphisme de $\mathcal{O}_X$-Algèbres lorsque
$\mathscr{P}_{X/S}^n(\mathscr{F})$ est muni de sa structure de
$\mathcal{O}_X$-Algèbre provenant de sa structure de $\mathcal{O}_X$-Module
\emph{à droite}.

Avec les mêmes notations, on a également un homomorphisme canonique de
$\mathscr{P}_{X/S}^n$-Modules
\[
  \mathscr{P}_{X/S}^n
  (\mathscr{H}\!om_{\mathcal{O}_X}(\mathscr{F},\mathscr{G}))
  \longrightarrow
  \mathscr{H}\!om_{\mathscr{P}_{X/S}^n}
  (\mathscr{P}_{X/S}^n(\mathscr{F}),\mathscr{P}_{X/S}^n(\mathscr{G}))
  \tag{16.7.8.2}
  \label{IV.16.7.8.2}
\]
(Bourbaki, \emph{Alg.}, 3\textsuperscript{e} éd., §~5,
n\textsuperscript{o}~3), qui est bijectif lorsque $\mathscr{P}_{X/S}^n$ est
un $\mathcal{O}_X$-Module \emph{localement libre de type fini}
(\emph{loc. cit.}, prop.~7).
\end{env}

\begin{env}[16.7.9]
\label{IV.16.7.9}
Supposons qu'on soit dans la situation décrite dans (16.4.1); alors, de
l'homomorphisme canonique $P^n(u)$ (16.4.3.3), on déduit aussitôt un
homomorphisme canonique de $\mathcal{O}_{X'}$-Bimodules
\[
  u^*(\mathscr{P}_{X/S}^n(\mathscr{F}))
  \longrightarrow
  \mathscr{P}_{X'/S'}^n(u^*(\mathscr{F})).
  \tag{16.7.9.1}
  \label{IV.16.7.9.1}
\]
Nous laissons au lecteur le soin d'étendre à cet homomorphisme les propriétés
vues dans (16.4) pour le cas $\mathscr{F}=\mathcal{O}_X$.
\end{env}

\begin{remark}[16.7.10]
\label{IV.16.7.10}
La définition de $\mathscr{P}_{X/S}^n(\mathscr{F})$ sous la forme (16.7.1.2)
garde un sens lorsque $\mathscr{F}$ est un faisceau d'ensembles quelconque
(l'image réciproque d'un faisceau d'ensembles par $p_2^{(n)}$ étant définie
dans (0\textsubscript{I}, 3.7.1)); une variante de cette définition permet de
définir le «schéma des jets» (relativement à $S$) d'un $X$-préschéma
quelconque.
\end{remark}

\subsection{Opérateurs différentiels.}
\footnote{Pour un formalisme plus général, voir l'Exposé VII de [42]
(dû à P. Gabriel).}
\EGAReaderTarget{ega4.c04.subsec.16-8}{IV.16.8}

\begin{definition}[16.8.1]
\EGAReaderTarget{ega4.c04.def.16-8-1}{IV.16.8.1}
Soient $f=(\psi,\theta):X\to S$ un morphisme de préschémas,
$\mathscr{F}$, $\mathscr{G}$ deux $\mathcal{O}_X$-Modules, $n$ un entier
$\geq0$. On dit qu'un homomorphisme de faisceaux de groupes additifs
$D:\mathscr{F}\to\mathscr{G}$ est un \emph{opérateur différentiel d'ordre
$\leq n$} (relativement à $S$) s'il existe un homomorphisme de
$\mathcal{O}_X$-Modules
$u:\mathscr{P}_{X/S}^n(\mathscr{F})\to\mathscr{G}$ (où
$\mathscr{P}_{X/S}^n(\mathscr{F})$ est muni de sa structure de
$\mathcal{O}_X$-Module à gauche (16.7.4)) tel que l'on ait
$D=u\circ d_{X/S,\mathscr{F}}^n$.
\end{definition}

Il est clair, en vertu de l'existence des homomorphismes canoniques
\[
  \mathscr{P}_{X/S}^m(\mathscr{F})
  \longrightarrow\mathscr{P}_{X/S}^n(\mathscr{F})
\]
\oldpage[IV-4]{40}
pour $n\leq m$ (16.7.7) qu'un opérateur différentiel d'ordre $\leq n$ est
aussi un opérateur différentiel d'ordre $\leq m$ pour tout $m\geq n$. Si
$D:\mathscr{F}\to\mathscr{G}$ est un opérateur différentiel d'ordre
$\leq n$, alors, pour tout ouvert $U$ de $X$,
$D|U:\mathscr{F}|U\to\mathscr{G}|U$ est aussi un opérateur différentiel
d'ordre $\leq n$.

On dit qu'un homomorphisme $D:\mathscr{F}\to\mathscr{G}$ des faisceaux de
groupes additifs sous-jacents à $\mathscr{F}$ et à $\mathscr{G}$ est un
\emph{opérateur différentiel} (relativement à $S$) si, pour tout $x\in X$, il
existe un voisinage ouvert $U$ de $x$ et un entier $n\geq0$ tels que
$D|U:\mathscr{F}|U\to\mathscr{G}|U$ soit un opérateur différentiel d'ordre
$\leq n$. L'\emph{ordre} d'un opérateur différentiel
$D:\mathscr{F}\to\mathscr{G}$ est la borne inférieure des entiers $n$ tels
que $D$ soit un opérateur différentiel d'ordre $\leq n$ (et donc $+\infty$
s'il n'y a pas de tels entiers); cet ordre est toujours fini si $X$ est
\emph{quasi-compact}. Les opérateurs différentiels d'ordre $0$ ne sont autres
que les homomorphismes de $\mathcal{O}_X$-Modules
$\mathscr{F}\to\mathscr{G}$; on convient que tout opérateur différentiel
d'ordre $<0$ est \emph{nul}. Pour $n\geq0$, un opérateur différentiel n'est
pas en général un homomorphisme de $\mathcal{O}_X$-Modules mais est toujours
un homomorphisme de $\psi^*(\mathcal{O}_S)$-Modules.

Lorsque $\mathscr{F}=\mathcal{O}_X$, un opérateur différentiel d'ordre
$\leq1$ de $\mathcal{O}_X$ dans $\mathscr{G}$ s'écrit d'une seule manière
sous la forme $v+D$, où $v:\mathcal{O}_X\to\mathscr{G}$ est un
$\mathcal{O}_X$-homomorphisme, et $D$ une $S$-\emph{dérivation} (16.5.1) de
$\mathcal{O}_X$ dans $\mathscr{G}$: cela résulte de la structure de
$P_{B/A}^1$ (0, 20.4.8).

\begin{env}[16.8.2]
\label{IV.16.8.2}
Pour décrire de façon plus explicite un opérateur différentiel d'ordre
$\leq n$, $D:\mathscr{F}\to\mathscr{G}$, il suffit, pour tout ouvert affine
$U$ de $X$, dont l'image dans $S$ est contenue dans un ouvert affine $V$, de
caractériser l'homomorphisme
$D=D_U:\Gamma(U,\mathscr{F})\to\Gamma(U,\mathscr{G})$. Si l'on pose
$\Gamma(V,\mathcal{O}_S)=A$, $\Gamma(U,\mathcal{O}_X)=B$, de sorte que $B$
est une $A$-algèbre, on a
\[
  \Gamma(U,\mathscr{P}_{X/S}^n)
  =(B\otimes_A B)/\mathfrak{I}^{n+1},
\]
où l'on a posé pour abréger $\mathfrak{I}=\mathfrak{I}_{B/A}$. Posons en
outre $M=\Gamma(U,\mathscr{F})$, $N=\Gamma(U,\mathscr{G})$; alors la
définition de $D$ signifie que pour chaque couple $(U,V)$ vérifiant les
conditions précédentes, le $A$-homomorphisme $D:M\to N$ se factorise en
\[
  M\longrightarrow
  ((B\otimes_A B)/\mathfrak{I}^{n+1})\otimes_B M
  \xrightarrow{v}N
\]
où la première flèche est l'homomorphisme canonique $t\mapsto1\otimes t$, et
$v$ est un $B$-homomorphisme, la structure de $B$-module de
$((B\otimes_A B)/\mathfrak{I}^{n+1})\otimes_B M$ provenant du premier
facteur $B$ (alors qu'on rappelle que dans la formation du produit tensoriel
sur $B$, la structure de $B$-module de
$(B\otimes_A B)/\mathfrak{I}^{n+1}$ provient du second facteur $B$). Notons
maintenant que le $B$-module
$((B\otimes_A B)/\mathfrak{I}^{n+1})\otimes_B M$ est isomorphe à
$(B\otimes_A M)/\mathfrak{I}^{n+1}(B\otimes_A M)$, où $B\otimes_A M$ est
considéré comme $(B\otimes_A B)$-module et sa structure de $B$-module
provient de l'homomorphisme $b\mapsto b\otimes1$ de $B$ dans $B\otimes_A B$.
Soit alors $D'$ le $B$-homomorphisme de $B\otimes_A M$ dans $N$ tel que
$D'(b\otimes t)=bD(t)$; la condition de factorisation sur $D$ s'exprime encore
en disant que $D'$ doit \emph{être nul dans le $B$-module}
$\mathfrak{I}^{n+1}(B\otimes_A M)$.
\end{env}

\begin{env}[16.8.3]
\label{IV.16.8.3}
Il est clair que l'ensemble des opérateurs différentiels d'ordre $\leq n$ de
$\mathscr{F}$ dans $\mathscr{G}$ est un groupe additif, noté
$\operatorname{Diff}_{X/S}^n(\mathscr{F},\mathscr{G})$; lorsque
$\mathscr{F}=\mathscr{G}=\mathcal{O}_X$, on écrit aussi
$\operatorname{Diff}_{X/S}^n$ au lieu de
$\operatorname{Diff}_{X/S}^n(\mathcal{O}_X,\mathcal{O}_X)$.

On a vu (16.8.1) que l'on a pour deux ouverts $U\supset V$ de $X$, un
homomorphisme canonique de restriction
\[
  \operatorname{Diff}_{U/S}^n(\mathscr{F}|U,\mathscr{G}|U)
  \longrightarrow
  \operatorname{Diff}_{V/S}^n(\mathscr{F}|V,\mathscr{G}|V)
\]
\oldpage[IV-4]{41}
donc $U\mapsto\operatorname{Diff}_{U/S}^n(\mathscr{F}|U,\mathscr{G}|U)$ est
un préfaisceau de groupes additifs; en fait c'est même un \emph{faisceau}, car
pour $U$ ouvert variable dans $X$, les homomorphismes
$u\mapsto u\circ d_{U/S,\mathscr{F}|U}^n$ sont des \emph{isomorphismes} de
groupes additifs
\[
  \mathscr{H}\!om_{\mathcal{O}_U}
  (\mathscr{P}_{U/S}^n(\mathscr{F}|U),\mathscr{G}|U)
  \xrightarrow{\sim}
  \operatorname{Diff}_{U/S}^n(\mathscr{F}|U,\mathscr{G}|U),
  \tag{16.8.3.1}
  \label{IV.16.8.3.1}
\]
en vertu du fait que l'image de $\mathscr{F}$ par
$d_{X/S,\mathscr{F}}^n$ engendre $\mathscr{P}_{X/S}^n(\mathscr{F})$
(16.7.6). On note ce faisceau
$\mathscr{D}\!iff_{X/S}^n(\mathscr{F},\mathscr{G})$, et on a donc:
\end{env}

\begin{proposition}[16.8.4]
\label{IV.16.8.4}
Les isomorphismes (16.8.3.1) définissent un isomorphisme de faisceaux de
groupes additifs
\[
  \mathscr{H}\!om_{\mathcal{O}_X}
  (\mathscr{P}_{X/S}^n(\mathscr{F}),\mathscr{G})
  \xrightarrow{\sim}
  \mathscr{D}\!iff_{X/S}^n(\mathscr{F},\mathscr{G}).
  \tag{16.8.4.1}
  \label{IV.16.8.4.1}
\]
\end{proposition}

Lorsque $\mathscr{F}=\mathscr{G}=\mathcal{O}_X$, on écrit aussi
$\mathscr{D}\!iff_{X/S}^n$ au lieu de
$\mathscr{D}\!iff_{X/S}^n(\mathcal{O}_X,\mathcal{O}_X)$; il résulte de
(16.8.4) que $\mathscr{D}\!iff_{X/S}^n$ s'identifie canoniquement au
\emph{dual} du $\mathcal{O}_X$-Module $\mathscr{P}_{X/S}^n$; aussi écrit-on
$\langle t,D\rangle$ au lieu de $u(t)$ si $t$ est une section de
$\mathscr{P}_{X/S}^n$ au-dessus d'un ouvert et si $u$ est l'homomorphisme de
$\mathscr{P}_{X/S}^n$ dans $\mathcal{O}_X$ correspondant à $D$.

\begin{env}[16.8.5]
\label{IV.16.8.5}
Puisque $\mathscr{P}_{X/S}^n(\mathscr{F})$ est muni d'une structure de
$\mathcal{O}_X$-Bimodule (16.7.4), on en déduit canoniquement une structure de
$\mathcal{O}_X$-Bimodule sur
$\mathscr{H}\!om_{\mathcal{O}_X}
(\mathscr{P}_{X/S}^n(\mathscr{F}),\mathscr{G})$, donc aussi sur
$\mathscr{D}\!iff_{X/S}^n(\mathscr{F},\mathscr{G})$ en vertu de (16.8.4.1).
De façon précise, à la structure de $\mathcal{O}_X$-Module à gauche de
$\mathscr{P}_{X/S}^n(\mathscr{F})$ correspond, en vertu de la définition
(16.8.1), la structure de $\mathcal{O}_X$-Module à gauche sur
$\mathscr{D}\!iff_{X/S}^n(\mathscr{F},\mathscr{G})$ explicitée comme suit:
pour tout ouvert $U$ de $X$, toute section
$a\in\Gamma(U,\mathcal{O}_X)$ et tout opérateur différentiel
$D:\mathscr{F}|U\to\mathscr{G}|U$, $aD$ est l'opérateur différentiel qui, à
toute section $t\in\Gamma(U,\mathscr{F})$, fait correspondre la section
\[
  (aD)(t)=a(D(t))
  \tag{16.8.5.1}
  \label{IV.16.8.5.1}
\]
de $\Gamma(U,\mathscr{G})$. De même, à la structure de
$\mathcal{O}_X$-Module à droite de $\mathscr{P}_{X/S}^n(\mathscr{F})$
correspond la structure de $\mathcal{O}_X$-Module à droite sur
$\mathscr{D}\!iff_{X/S}^n(\mathscr{F},\mathscr{G})$ explicitée comme suit:
avec les mêmes notations que ci-dessus, $Da$ est l'opérateur différentiel qui,
à $t\in\Gamma(U,\mathscr{F})$, fait correspondre la section
\[
  (Da)(t)=D(at).
  \tag{16.8.5.2}
  \label{IV.16.8.5.2}
\]
\end{env}

\begin{proposition}[16.8.6]
\label{IV.16.8.6}
Si $f:X\to S$ est un morphisme localement de présentation finie,
$\mathscr{F}$ un $\mathcal{O}_X$-Module quasi-cohérent de présentation finie
et $\mathscr{G}$ un $\mathcal{O}_X$-Module quasi-cohérent, alors
$\mathscr{D}\!iff_{X/S}^n(\mathscr{F},\mathscr{G})$ est un
$\mathcal{O}_X$-Module quasi-cohérent pour l'une ou l'autre structure définies
dans (16.8.5).
\end{proposition}

La proposition résulte de ce que, sous les hypothèses faites,
$\mathscr{P}_{X/S}^n(\mathscr{F})$ est un $\mathcal{O}_X$-Module
quasi-cohérent de présentation finie (16.7.4), et de (I, 1.3.12).

\begin{env}[16.8.7]
\label{IV.16.8.7}
L'ensemble des opérateurs différentiels de $\mathscr{F}$ dans $\mathscr{G}$
(d'ordre non précisé (16.8.1)) se note
$\operatorname{Diff}_{X/S}(\mathscr{F},\mathscr{G})$; on voit encore comme
dans (16.8.3) que
$U\mapsto\operatorname{Diff}_{U/S}(\mathscr{F}|U,\mathscr{G}|U)$ est un
faisceau de groupes additifs, que nous désignerons par
$\mathscr{D}\!iff_{X/S}(\mathscr{F},\mathscr{G})$. Il est immédiat que
$\mathscr{D}\!iff_{X/S}(\mathscr{F},\mathscr{G})$ est réunion de la famille
filtrante croissante de ses sous-faisceaux
$\mathscr{D}\!iff_{X/S}^n(\mathscr{F},\mathscr{G})$; si $X$ est
quasi-compact, $\operatorname{Diff}_{X/S}(\mathscr{F},\mathscr{G})$
\oldpage[IV-4]{42}
est de même réunion de ses sous-groupes
$\operatorname{Diff}_{X/S}^n(\mathscr{F},\mathscr{G})$ (16.8.1). Les
structures de $\mathcal{O}_X$-Bimodule sur les
$\mathscr{D}\!iff_{X/S}^n(\mathscr{F},\mathscr{G})$ définissent donc une
structure de $\mathcal{O}_X$-Bimodule sur
$\mathscr{D}\!iff_{X/S}(\mathscr{F},\mathscr{G})$, explicitée encore par
(16.8.5.1) et (16.8.5.2).

Notons que, pour $n\leq m$, on a un diagramme commutatif
\[
  \xymatrix{
    \mathscr{H}\!om_{\mathcal{O}_X}
      (\mathscr{P}_{X/S}^n(\mathscr{F}),\mathscr{G})
      \ar[r]^-{\sim}\ar[d]
      &\mathscr{D}\!iff_{X/S}^n(\mathscr{F},\mathscr{G})\ar[d]\\
    \mathscr{H}\!om_{\mathcal{O}_X}
      (\mathscr{P}_{X/S}^m(\mathscr{F}),\mathscr{G})
      \ar[r]^-{\sim}
      &\mathscr{D}\!iff_{X/S}^m(\mathscr{F},\mathscr{G})
  }
  \tag{16.8.7.1}
  \label{IV.16.8.7.1}
\]
où les flèches horizontales sont les isomorphismes (16.8.4.2) et la flèche
verticale de gauche provient de l'homomorphisme canonique
$\mathscr{P}_{X/S}^m(\mathscr{F})\to
\mathscr{P}_{X/S}^n(\mathscr{F})$ (16.7.7). Pour tout ouvert $U$ de $X$,
munissons alors
\[
  \Gamma(U,\mathscr{P}_{X/S}^{\infty}(\mathscr{F}))
  =\varprojlim\Gamma(U,\mathscr{P}_{X/S}^n(\mathscr{F}))
\]
de la topologie limite projective des topologies discrètes sur les
$\Gamma(U,\mathscr{P}_{X/S}^n(\mathscr{F}))$, ce qui définit
$\Gamma(U,\mathscr{P}_{X/S}^{\infty}(\mathscr{F}))$ comme un
$\Gamma(U,\mathcal{O}_X)$-bimodule topologique, de sorte que
$\mathscr{P}_{X/S}^{\infty}(\mathscr{F})$ apparaît comme un faisceau à valeurs
dans la catégorie des groupes commutatifs topologiques
(0\textsubscript{I}, 3.2.6). Alors (G, II, 1.11) la limite du système inductif
de faisceaux de groupes commutatifs
\[
  (\mathscr{H}\!om_{\mathcal{O}_X}
  (\mathscr{P}_{X/S}^n(\mathscr{F}),\mathscr{G}))
\]
n'est autre que le faisceau des germes d'homomorphismes \emph{continus} de
$\mathscr{P}_{X/S}^{\infty}(\mathscr{F})$ dans $\mathscr{G}$ (ce dernier
étant muni de la topologie discrète): les homomorphismes continus de
$\Gamma(U,\mathscr{P}_{X/S}^{\infty}(\mathscr{F}))$ dans le groupe discret
$\Gamma(U,\mathscr{G})$ correspondent en effet biunivoquement aux systèmes
inductifs d'homomorphismes de groupes
$\Gamma(U,\mathscr{P}_{X/S}^n(\mathscr{F}))\to\Gamma(U,\mathscr{G})$. On peut
donc encore exprimer (16.8.4) en disant qu'on a un isomorphisme canonique
\[
  \mathscr{H}\!om.\mathrm{cont}_{\mathcal{O}_X}
  (\mathscr{P}_{X/S}^{\infty}(\mathscr{F}),\mathscr{G})
  \xrightarrow{\sim}
  \mathscr{D}\!iff_{X/S}(\mathscr{F},\mathscr{G})
\]
où le premier membre désigne le faisceau des germes d'homomorphismes continus
de $\mathscr{P}_{X/S}^{\infty}(\mathscr{F})$ dans $\mathscr{G}$.
\end{env}

\begin{proposition}[16.8.8]
\label{IV.16.8.8}
Soient $\mathscr{F}$, $\mathscr{G}$ deux $\mathcal{O}_X$-Modules,
$D:\mathscr{F}\to\mathscr{G}$ un homomorphisme de
$\psi^*(\mathcal{O}_S)$-Modules, $n$ un entier $\geq0$. Les conditions
suivantes sont équivalentes:
\begin{enumerate}
  \item[(a)] $D$ est un opérateur différentiel d'ordre $\leq n$.
  \item[(b)] Pour toute section $a$ de $\mathcal{O}_X$ au-dessus d'un ouvert
  $U$, l'homomorphisme $D_a:\mathscr{F}|U\to\mathscr{G}|U$ tel que, pour
  toute section $t$ de $\mathscr{F}$ au-dessus d'un ouvert $V\subset U$, on
  ait
  \[
    D_a(t)=D(at)-aD(t)
    \tag{16.8.8.1}
    \label{IV.16.8.8.1}
  \]
  est un opérateur différentiel d'ordre $\leq n-1$.
  \item[(c)] Pour tout ouvert $U$ de $X$, toute famille
  $(a_i)_{1\leq i\leq n+1}$ de $n+1$ sections de $\mathcal{O}_X$ au-dessus
  de $U$ et toute section $t$ de $\mathscr{F}$ au-dessus de $U$, on a
  l'identité
  \[
    \sum_{H\subset I_{n+1}}(-1)^{\operatorname{Card}(H)}
    \left(\prod_{i\in H}a_i\right)
    D\!\left(\left(\prod_{i\notin H}a_i\right)t\right)=0
    \tag{16.8.8.2}
    \label{IV.16.8.8.2}
  \]
  (où $I_{n+1}$ est l'intervalle $1\leq i\leq n+1$ de $\mathbb{N}$).
\end{enumerate}
\oldpage[IV-4]{43}
\end{proposition}

Prouvons d'abord l'équivalence de a) et c). Par définition, pour prouver que
$D$ est un opérateur différentiel d'ordre $\leq n$, il suffit de montrer qu'il
en est ainsi pour la restriction
$D|U:\mathscr{F}|U\to\mathscr{G}|U$ à tout ouvert affine $U$ de $X$, et
d'autre part la propriété c) est valable pour tout ouvert $U$ de $X$ si elle
l'est pour tout ouvert affine. On peut donc se borner au cas où
$S=\operatorname{Spec}(A)$ et $X=\operatorname{Spec}(B)$ sont affines. En
vertu de (16.8.2) (dont on garde les notations), la condition a) signifie
alors que le $A$-homomorphisme $D':B\otimes_A M\to N$ tel que
$D'(b\otimes t)=bD(t)$ s'annule dans
$\mathfrak{I}^{n+1}(B\otimes_A M)$, ce qui, en vertu de (0, 20.4.4), équivaut
à dire que $D'$ s'annule pour tous les éléments de la forme
\[
  \left(\prod_{i=1}^{n+1}(a_i\otimes1-1\otimes a_i)\right).(1\otimes t)
\]
où $a_i\in B$ et $t\in M$. Or, cet élément s'écrit
\[
  \sum_{H\subset I_{n+1}}
  \left(\prod_{i\in H}a_i\right)\otimes
  \left(\left(\prod_{i\notin H}a_i\right)t\right),
\]
et la valeur de $D'$ pour cet élément n'est autre que le premier membre de
(16.8.8.2), ce qui prouve l'équivalence de a) et c).

Prouvons maintenant l'équivalence de b) et c). Raisonnons par récurrence sur
$n$, l'assertion étant triviale pour $n=0$. Écrivant $a_{n+1}$ au lieu de $a$
dans la condition b), on voit, en vertu de l'hypothèse de récurrence, que la
condition b) signifie que pour toute famille $(a_i)_{1\leq i\leq n}$ de $n$
sections de $\mathcal{O}_X$ au-dessus de $U$ et toute section $t$ de
$\mathscr{F}$ au-dessus de $U$, on a
\[
  \sum_{H'\subset I_n}(-1)^{\operatorname{Card}(H')}
  \left(\prod_{i\in H'}a_i\right)
  D_{a_{n+1}}\!\left(\left(\prod_{i\notin H'}a_i\right)t\right)=0.
\]
Mais si l'on remplace dans cette relation $D_{a_{n+1}}$ par sa définition
(16.8.8.1), on constate aussitôt qu'on obtient, au signe près, le premier
membre de (16.8.8.2); d'où la conclusion.

\begin{proposition}[16.8.9]
\label{IV.16.8.9}
Si $D:\mathscr{F}\to\mathscr{G}$ est un opérateur différentiel d'ordre
$\leq n$, et $D':\mathscr{G}\to\mathscr{H}$ un opérateur différentiel d'ordre
$\leq n'$, alors $D'\circ D:\mathscr{F}\to\mathscr{H}$ est un opérateur
différentiel d'ordre $\leq n+n'$.
\end{proposition}

Par hypothèse, on peut écrire $D=u\circ d_{X/S,\mathscr{F}}^n$ et
$D'=v\circ d_{X/S,\mathscr{G}}^{n'}$, où
$u:\mathscr{P}_{X/S}^n\otimes_{\mathcal{O}_X}\mathscr{F}\to\mathscr{G}$ et
$v:\mathscr{P}_{X/S}^{n'}\otimes_{\mathcal{O}_X}\mathscr{G}\to\mathscr{H}$
sont des $\mathcal{O}_X$-homomorphismes. Tout revient à montrer que
l'homomorphisme composé de faisceaux de groupes additifs
\[
  \xymatrix{
    \mathscr{F}\ar[r]^-{d_{X/S,\mathscr{F}}^n}
      &\mathscr{P}_{X/S}^n\otimes_{\mathcal{O}_X}\mathscr{F}
        \ar[r]^-u
      &\mathscr{G}\ar[r]^-{d_{X/S,\mathscr{G}}^{n'}}
      &\mathscr{P}_{X/S}^{n'}\otimes_{\mathcal{O}_X}\mathscr{G}
  }
\]
se factorise en
\[
  \xymatrix{
    \mathscr{F}\ar[r]^-{d_{X/S,\mathscr{F}}^{n+n'}}
      &\mathscr{P}_{X/S}^{n+n'}\otimes_{\mathcal{O}_X}\mathscr{F}
        \ar[r]^-w
      &\mathscr{P}_{X/S}^{n'}\otimes_{\mathcal{O}_X}\mathscr{G}
  }
\]
où $w$ est un $\mathcal{O}_X$-homomorphisme. Il suffira de prouver le

\begin{lemma}[16.8.9.1]
\label{IV.16.8.9.1}
Il existe un $\mathcal{O}_X$-homomorphisme et un seul
\[
  \delta:\mathscr{P}_{X/S}^{n+n'}
  \longrightarrow
  \mathscr{P}_{X/S}^{n'}(\mathscr{P}_{X/S}^n)
  =\mathscr{P}_{X/S}^{n'}\otimes_{\mathcal{O}_X}\mathscr{P}_{X/S}^n
  \tag{16.8.9.2}
  \label{IV.16.8.9.2}
\]
\oldpage[IV-4]{44}
rendant commutatif le diagramme
\[
  \xymatrix{
    \mathcal{O}_X\ar[r]^-{d_{X/S}^{n+n'}}\ar[d]_-{d_{X/S}^n}
      &\mathscr{P}_{X/S}^{n+n'}\ar[d]^-\delta\\
    \mathscr{P}_{X/S}^n
      \ar[r]_-{d_{X/S,\mathscr{P}_{X/S}^n}^{n'}}
      &\mathscr{P}_{X/S}^{n'}(\mathscr{P}_{X/S}^n)
  }
  \tag{16.8.9.3}
  \label{IV.16.8.9.3}
\]
\end{lemma}

On aura alors en effet un diagramme commutatif déduit de (16.8.9.3) par
tensorisation avec $\mathscr{F}$
\[
  \xymatrix{
    \mathscr{F}\ar[r]^-{d_{X/S,\mathscr{F}}^{n+n'}}
      \ar[d]_-{d_{X/S,\mathscr{F}}^n}
      &\mathscr{P}_{X/S}^{n+n'}(\mathscr{F})
        \ar[d]^-{\delta\otimes1}\\
    \mathscr{P}_{X/S}^n(\mathscr{F})
      \ar[r]_-{d_{X/S,\mathscr{P}_{X/S}^n(\mathscr{F})}^{n'}}
      &\mathscr{P}_{X/S}^{n'}(\mathscr{P}_{X/S}^n(\mathscr{F}))
  }
\]
et d'autre part, on vérifie aussitôt sur la définition (16.7.5) que le
diagramme
\[
  \xymatrix{
    \mathscr{P}_{X/S}^n(\mathscr{F})\ar[r]^-u
      \ar[d]_-{d_{X/S,\mathscr{P}_{X/S}^n(\mathscr{F})}^{n'}}
      &\mathscr{G}\ar[d]^-{d_{X/S,\mathscr{G}}^{n'}}\\
    \mathscr{P}_{X/S}^{n'}(\mathscr{P}_{X/S}^n(\mathscr{F}))
      \ar[r]_-{1\otimes u}
      &\mathscr{P}_{X/S}^{n'}(\mathscr{G})
  }
\]
est commutatif. On répondra donc à la question en prenant pour $w$ le
$\mathcal{O}_X$-homomorphisme composé
\[
  \xymatrix{
    \mathscr{P}_{X/S}^{n+n'}(\mathscr{F})
      \ar[r]^-{\delta\otimes1}
      &\mathscr{P}_{X/S}^{n'}(\mathscr{P}_{X/S}^n(\mathscr{F}))
        \ar[r]^-{1\otimes u}
      &\mathscr{P}_{X/S}^{n'}(\mathscr{G}).
  }
\]

Reste à prouver le lemme (16.8.9.1). Compte tenu de (16.7.6), qui prouve
l'unicité de $\delta$, on est ramené au cas où $S=\operatorname{Spec}(A)$ et
$X=\operatorname{Spec}(B)$ sont affines; posant
$\mathfrak{I}=\mathfrak{I}_{B/A}$, il s'agit de définir un homomorphisme
canonique de $B$-modules
\[
  \varphi:(B\otimes_A B)/\mathfrak{I}^{n+n'+1}
  \longrightarrow
  ((B\otimes_A B)/\mathfrak{I}^{n'+1})
  \otimes_B((B\otimes_A B)/\mathfrak{I}^{n+1})
\]
les structures de $B$-module des deux membres provenant du premier facteur
$B$; rappelons que dans le produit tensoriel du second membre,
$(B\otimes_A B)/\mathfrak{I}^{n'+1}$ doit être considéré
\oldpage[IV-4]{45}
comme $B$-module à droite par son second facteur $B$ et
$(B\otimes_A B)/\mathfrak{I}^{n+1}$ comme $B$-module à gauche par son premier
facteur $B$ (16.7.2). Il revient au même de définir un homomorphisme de
$B$-modules
\[
  \varphi_0:B\otimes_A B\longrightarrow
  ((B\otimes_A B)/\mathfrak{I}^{n'+1})
  \otimes_B((B\otimes_A B)/\mathfrak{I}^{n+1})
\]
et de prouver qu'il s'annule dans $\mathfrak{I}^{n+n'+1}$. Or, on définit
aussitôt un tel homomorphisme par la condition que
\[
  \varphi_0(b\otimes b')
  =\pi_{n'}(b\otimes1)\otimes\pi_n(1\otimes b')
  \qquad\text{pour $b,b'$ dans $B$}
\]
avec les notations de (16.3.7). En outre, il est immédiat que $\varphi_0$ est
un homomorphisme d'\emph{anneaux}. Or, on peut écrire
\[
  \begin{aligned}
  \varphi_0(b\otimes1-1\otimes b)
  ={}&\pi_{n'}(b\otimes1-1\otimes b)\otimes\pi_n(1\otimes1)\\
  &+\pi_{n'}(1\otimes b)\otimes\pi_n(1\otimes1)
  -\pi_{n'}(1\otimes1)\otimes\pi_n(1\otimes b)
  \end{aligned}
\]
et on a
\[
  \begin{aligned}
  \pi_{n'}(1\otimes b)\otimes\pi_n(1\otimes1)
  &=\pi_{n'}(1\otimes1)b\otimes\pi_n(1\otimes1)\\
  &=\pi_{n'}(1\otimes1)\otimes b\pi_n(1\otimes1)\\
  &=\pi_{n'}(1\otimes1)\otimes\pi_n(b\otimes1)
  \end{aligned}
\]
d'où finalement
\[
  \begin{aligned}
  \varphi_0(b\otimes1-1\otimes b)
  ={}&\pi_{n'}(b\otimes1-1\otimes b)\otimes\pi_n(1\otimes1)\\
  &+\pi_{n'}(1\otimes1)\otimes\pi_n(b\otimes1-1\otimes b).
  \end{aligned}
  \tag{16.8.9.4}
  \label{IV.16.8.9.4}
\]
Un produit de $n+n'+1$ termes de la forme (16.8.9.4) est donc
nécessairement nul, car il en est ainsi d'un produit de $n+1$ termes de la
forme $\pi_n(b\otimes1-1\otimes b)$ et d'un produit de $n'+1$ termes de la
forme $\pi_{n'}(b\otimes1-1\otimes b)$. La conclusion résulte donc de
(0, 20.4.4).

\begin{corollary}[16.8.10]
\label{IV.16.8.10}
Le faisceau $\mathscr{D}\!iff_{X/S}(\mathcal{O}_X,\mathcal{O}_X)$ (aussi noté
$\mathscr{D}\!iff_{X/S}$) est canoniquement muni d'une structure de faisceau
d'anneaux, les $\mathscr{D}\!iff_{X/S}^n$ formant une filtration croissante
compatible avec cette structure.
\end{corollary}

En particulier, $\mathscr{D}\!iff_{X/S}^0$ est un faisceau de sous-anneaux de
$\mathscr{D}\!iff_{X/S}$, qui s'identifie canoniquement à $\mathcal{O}_X$
(16.8.1). Les formules (16.8.5.1) et (16.8.5.2) montrent que la structure de
$\mathcal{O}_X$-Bimodule de $\mathscr{D}\!iff_{X/S}$ provient de la
multiplication à gauche et à droite par les sections de $\mathcal{O}_X$
considéré comme faisceau de sous-anneaux de $\mathscr{D}\!iff_{X/S}$.

\begin{remarks}[16.8.11]
\label{IV.16.8.11}
\begin{enumerate}
  \item[(i)] Supposons que
  $\mathscr{F}=\bigoplus_{\lambda\in L}\mathscr{F}_\lambda$; alors il est
  clair (16.7.2.1) que
  $\mathscr{P}_{X/S}^n(\mathscr{F})=
  \bigoplus_{\lambda\in L}\mathscr{P}_{X/S}^n(\mathscr{F}_\lambda)$; comme le
  foncteur $\mathscr{F}\mapsto\Gamma(U,\mathscr{F})$ commute à la formation
  des sommes directes quelconques, $d_{X/S,\mathscr{F}}^n$ est
  l'homomorphisme dont la restriction à chaque $\mathscr{F}_\lambda$ est
  \[
    d_{X/S,\mathscr{F}_\lambda}^n:\mathscr{F}_\lambda
    \longrightarrow\mathscr{P}_{X/S}^n(\mathscr{F}_\lambda);
  \]
  on en conclut aussitôt que l'on a
  \[
    \operatorname{Diff}_{X/S}^n(\mathscr{F},\mathscr{G})
    =\prod_{\lambda\in L}
    \operatorname{Diff}_{X/S}^n(\mathscr{F}_\lambda,\mathscr{G}),
  \]
  et par suite aussi (0\textsubscript{I}, 3.2.6)
  \[
    \mathscr{D}\!iff_{X/S}^n(\mathscr{F},\mathscr{G})
    =\prod_{\lambda\in L}
    \mathscr{D}\!iff_{X/S}^n(\mathscr{F}_\lambda,\mathscr{G}).
  \]
  Par ailleurs, si $\mathscr{G}=\prod_{\mu\in M}\mathscr{G}_\mu$
  (0\textsubscript{I}, 3.2.6), on a
  \[
    \operatorname{Hom}_{\mathcal{O}_X}
    (\mathscr{P}_{X/S}^n(\mathscr{F}),\mathscr{G})
    =\prod_{\mu\in M}\operatorname{Hom}_{\mathcal{O}_X}
    (\mathscr{P}_{X/S}^n(\mathscr{F}),\mathscr{G}_\mu),
  \]
  \oldpage[IV-4]{46}
  tout homomorphisme $u$ de $\mathscr{P}_{X/S}^n(\mathscr{F})$ dans
  $\mathscr{G}$ correspondant biunivoquement à la famille de ses composés
  \[
    u_\mu:\mathscr{P}_{X/S}^n(\mathscr{F})
    \longrightarrow\mathscr{G}\longrightarrow\mathscr{G}_\mu.
  \]
  On a donc
  \[
    \operatorname{Diff}_{X/S}^n(\mathscr{F},\mathscr{G})
    =\prod_{\mu\in M}
    \operatorname{Diff}_{X/S}^n(\mathscr{F},\mathscr{G}_\mu),
  \]
  et par suite aussi
  \[
    \mathscr{D}\!iff_{X/S}^n(\mathscr{F},\mathscr{G})
    =\prod_{\mu\in M}
    \mathscr{D}\!iff_{X/S}^n(\mathscr{F},\mathscr{G}_\mu).
  \]
  \item[(ii)] Jusqu'ici, on n'a guère rencontré d'opérateurs différentiels
  $\mathscr{F}\to\mathscr{G}$ que lorsque $\mathscr{F}$ et $\mathscr{G}$ sont
  des $\mathcal{O}_X$-Modules localement libres de rang fini, auquel cas, leur
  structure se ramène localement, en vertu de (i), à celle du faisceau
  $\mathscr{D}\!iff_{X/S}$; cette dernière va être étudiée plus loin (16.11)
  dans un cas particulier.
\end{enumerate}
\end{remarks}

\subsection{Immersions régulières et quasi-régulières.}
\EGAReaderTarget{ega4.c04.subsec.16-9}{IV.16.9}

\begin{definition}[16.9.1]
\EGAReaderTarget{ega4.c04.def.16-9-1}{IV.16.9.1}
Soit $X$ un espace annelé. On dit qu'un Idéal $\mathscr{I}$ de
$\mathcal{O}_X$ est \emph{régulier} (resp. \emph{quasi-régulier}) si, pour
tout point $x\in\operatorname{Supp}(\mathcal{O}_X/\mathscr{I})$, il existe un
voisinage ouvert $U$ de $x$ dans $X$ et une suite régulière (0, 15.2.2) (resp.
quasi-régulière (0, 15.2.2)) d'éléments de $\Gamma(U,\mathcal{O}_X)$ qui
engendre $\mathscr{I}|U$.
\end{definition}

On dira qu'une suite régulière (resp. quasi-régulière) de sections de
$\mathcal{O}_X$ au-dessus de $U$ qui engendre $\mathscr{I}|U$ est un
\emph{système régulier} (resp. \emph{quasi-régulier}) \emph{de générateurs}
de $\mathscr{I}|U$.

\begin{definition}[16.9.2]
\EGAReaderTarget{ega4.c04.def.16-9-2}{IV.16.9.2}
Soit $j:Y\to X$ une immersion de préschémas et soit $U$ un ouvert de $X$ tel
que $j(Y)\subset U$ et que $j$ soit une immersion fermée de $Y$ dans $U$. On
dit que $j$ est \emph{régulière} (resp. \emph{quasi-régulière}) si le
sous-préschéma fermé $j(Y)$ de $U$ associé à $j$ est défini par un Idéal
régulier (resp. quasi-régulier) de $\mathcal{O}_U$ (condition indépendante de
l'ouvert $U$ choisi).
\end{definition}

On dit qu'un sous-préschéma $Y$ d'un préschéma $X$ est \emph{régulièrement
immergé} (resp. \emph{quasi-régulièrement immergé}) si l'injection canonique
$j:Y\to X$ est une immersion régulière (resp. quasi-régulière). Si $Y$ est un
sous-préschéma fermé de $X$ et $\mathscr{I}$ l'Idéal de $\mathcal{O}_X$ qui
définit $Y$, il revient au même de dire que $\mathscr{I}$ est \emph{régulier}
(resp. \emph{quasi-régulier}).

Par exemple, si $A$ est un anneau \emph{intègre}, $f$ un élément $\neq0$ de
$A$, le sous-préschéma fermé $V(f)$ de $\operatorname{Spec}(A)$ (isomorphe à
$\operatorname{Spec}(A/fA)$) est \emph{régulièrement immergé} dans
$\operatorname{Spec}(A)$.

Tout Idéal régulier est quasi-régulier (0, 15.2.2); toute immersion régulière
est quasi-régulière (cf. (16.9.11) pour une réciproque).

\begin{proposition}[16.9.3]
\label{IV.16.9.3}
Soient $X$ un espace annelé, $\mathscr{I}$ un Idéal de $\mathcal{O}_X$,
$(f_i)_{1\leq i\leq m}$ une suite finie de sections de $\mathcal{O}_X$
au-dessus de $X$ engendrant $\mathscr{I}$. Pour que $(f_i)$ soit une suite
quasi-régulière (0, 15.2.2), il faut et il suffit que les conditions suivantes
soient vérifiées:
\begin{enumerate}
  \item[(i)] Les images canoniques des $f_i$ dans
  $\mathscr{I}/\mathscr{I}^2$ forment une base de cet
  $\mathcal{O}_X/\mathscr{I}$-Module.
  \item[(ii)] L'homomorphisme surjectif canonique (16.1.2.2)
  \[
    \mathbf{S}_{\mathcal{O}_X/\mathscr{I}}^{\bullet}
    (\mathscr{I}/\mathscr{I}^2)
    \longrightarrow\mathscr{G}r_{\mathscr{I}}^{\bullet}(\mathcal{O}_X)
  \]
  est bijectif.
\end{enumerate}
En outre, s'il en est ainsi, toute suite $(f'_i)_{1\leq i\leq n}$ de $n$
sections de $\mathscr{I}$ au-dessus de $X$ qui engendre $\mathscr{I}$ est
quasi-régulière.
\end{proposition}

\oldpage[IV-4]{47}

Les deux conditions de l'énoncé ne font que traduire la définition donnée dans
(0, 15.2.2), compte tenu de la définition des homomorphismes canoniques
(0, 15.2.1.1). La dernière assertion résulte de ce que, si un module $M$ sur
un anneau commutatif $A$ admet une base de $n$ éléments, tout système de
générateurs de $M$ ayant $n$ éléments est une base de $M$ (Bourbaki,
\emph{Alg. comm.}, chap.~II, §~3, cor.~5 du th.~1).

\begin{corollary}[16.9.4]
\label{IV.16.9.4}
Soient $X$ un espace annelé en anneaux locaux, $\mathscr{I}$ un Idéal de
$\mathcal{O}_X$. Pour que $\mathscr{I}$ soit quasi-régulier, il faut et il
suffit qu'il vérifie les conditions suivantes:
\begin{enumerate}
  \item[(i)] $\mathscr{I}$ est de type fini.
  \item[(ii)] $\mathscr{I}/\mathscr{I}^2$ est un
  $(\mathcal{O}_X/\mathscr{I})$-Module localement libre.
  \item[(iii)] L'homomorphisme canonique
  \[
    \mathbf{S}_{\mathcal{O}_X/\mathscr{I}}^{\bullet}
    (\mathscr{I}/\mathscr{I}^2)
    \longrightarrow\mathscr{G}r_{\mathscr{I}}^{\bullet}(\mathcal{O}_X)
    \tag{16.9.4.1}
    \label{IV.16.9.4.1}
  \]
  est bijectif.
\end{enumerate}
\end{corollary}

La nécessité des conditions résulte aussitôt de (16.9.3). Pour voir que ces
conditions sont suffisantes, il faut montrer, en vertu de (16.9.3), que si, en
un point $x\in\operatorname{Supp}(\mathcal{O}_X/\mathscr{I})$, il existe un
voisinage ouvert $U$ de $x$ dans $X$ et $n$ sections $f_i$ ($1\leq i\leq n$)
de $\mathscr{I}$ au-dessus de $U$ dont les images canoniques dans
$\mathscr{I}/\mathscr{I}^2$ forment une base de
$(\mathscr{I}/\mathscr{I}^2)|U$ sur
$(\mathcal{O}_X/\mathscr{I})|U$, alors il y a un voisinage ouvert
$V\subset U$ de $x$ tel que les $f_i|V$ engendrent $\mathscr{I}|V$. Or, par
hypothèse, on a $\mathscr{I}_x\neq\mathcal{O}_x$, donc $\mathscr{I}_x$ est
contenu dans l'idéal maximal de $\mathcal{O}_x$; comme $\mathscr{I}_x$ est un
$\mathcal{O}_x$-module de type fini et que les classes des $(f_i)_x$ dans
$\mathscr{I}_x/\mathscr{I}_x^2$ engendrent cet
$(\mathcal{O}_x/\mathscr{I}_x)$-module, le lemme de Nakayama montre que les
$(f_i)_x$ engendrent $\mathscr{I}_x$. Comme $\mathscr{I}$ est de type fini,
on conclut par (0\textsubscript{I}, 5.2.2).

\begin{corollary}[16.9.5]
\label{IV.16.9.5}
Soient $X$ un espace annelé en anneaux locaux, $\mathscr{I}$ un Idéal
quasi-régulier de $\mathcal{O}_X$, $(f_i)_{1\leq i\leq n}$ une suite de
sections de $\mathscr{I}$ au-dessus de $X$,
$x$ un point de $\operatorname{Supp}(\mathcal{O}_X/\mathscr{I})$. Les
conditions suivantes sont équivalentes:
\begin{enumerate}
  \item[(a)] Il existe un voisinage ouvert $U$ de $x$ dans $X$ tel que les
  $f_i|U$ forment une suite quasi-régulière d'éléments de
  $\Gamma(U,\mathcal{O}_X)$ engendrant $\mathscr{I}|U$.
  \item[(b)] Les $(f_i)_x$ forment un système de générateurs de
  $\mathscr{I}_x$ dont le nombre d'éléments est le plus petit possible.
  \item[(b')] Les $(f_i)_x$ forment un système minimal de générateurs de
  $\mathscr{I}_x$.
  \item[(c)] Si $\bar f_i$ est l'image canonique de $f_i$ dans
  $\Gamma(X,\mathscr{I}/\mathscr{I}^2)$, les $(\bar f_i)_x$ forment une base
  du $(\mathcal{O}_x/\mathscr{I}_x)$-module
  $\mathscr{I}_x/\mathscr{I}_x^2$.
\end{enumerate}
\end{corollary}

Par hypothèse, $\mathcal{O}_x$ est un anneau local, $\mathscr{I}_x$ un idéal
de type fini de $\mathcal{O}_x$ contenu dans l'idéal maximal de
$\mathcal{O}_x$; l'équivalence de b), b') et c) résulte donc du lemme de
Nakayama (Bourbaki, \emph{Alg. comm.}, chap.~II, §~3,
n\textsuperscript{o}~2, prop.~5). Il est clair que a) entraîne c) en vertu de
(16.9.3); d'autre part, il résulte de (0\textsubscript{I}, 5.2.2) que si la
condition c) est vérifiée (donc aussi b)), il y a un voisinage $U$ de $x$ dans
$X$ tel que $(\mathscr{I}/\mathscr{I}^2)|U$ ait un rang constant égal à $n$,
et que les $f_i|U$ engendrent $\mathscr{I}|U$; il suffit donc d'appliquer dans
$U$ la dernière assertion de (16.9.3).

\begin{remarks}[16.9.6]
\label{IV.16.9.6}
\begin{enumerate}
  \item[(i)] Sous les hypothèses générales de (16.9.5), il ne suffit pas que
  les $(\bar f_i)_y$ forment une base du
  $(\mathcal{O}_y/\mathscr{I}_y)$-module
  $\mathscr{I}_y/\mathscr{I}_y^2$ pour tout $y\in X$ pour que la suite
  $(f_i)$ engendre $\mathscr{I}$. On en a un exemple en
\oldpage[IV-4]{48}
  prenant $X=\operatorname{Spec}(A)$, où $A$ est un anneau de Dedekind, et
  $\mathscr{I}=\widetilde{\mathfrak{I}}$, où $\mathfrak{I}$ est un idéal
  premier non principal de $A$; on a alors en effet
  $\mathscr{I}_y/\mathscr{I}_y^2=0$ en tout point $y$ distinct du point
  $x\in X$ correspondant à $\mathfrak{I}$, et
  $\mathscr{I}_x/\mathscr{I}_x^2$ est de rang $1$ sur le corps
  $\mathcal{O}_x/\mathscr{I}_x$; en outre $\mathscr{I}$ est évidemment un
  Idéal régulier.
  \item[(ii)] Dans (16.9.5), on ne peut remplacer « quasi-régulière » par
  « régulière », même lorsque $X$ est un préschéma (cf. (16.9.12)).
  Désignons en effet par $B$ l'anneau des germes de fonctions indéfiniment
  différentiables au point $0$ dans $\mathbb{R}$; il a un idéal maximal
  $\mathfrak{m}$ engendré par le germe $t$ de l'application identique de
  $\mathbb{R}$ au point $0$, et l'intersection $\mathfrak{n}$ des
  $\mathfrak{m}^k$ pour $k>0$ n'est pas réduite à $0$. Soit maintenant $A$
  l'anneau quotient $B[T]/\mathfrak{n}TB[T]$, et soient $f_1,f_2$ les images
  canoniques dans $A$ des éléments $t$ et $T$ de $B[T]$. La suite
  $(f_1,f_2)$ est régulière dans $A$: en effet, $f_1$ n'est pas diviseur de
  $0$ dans $A$, car la relation $tP[T]\in\mathfrak{n}TB[T]$, pour un
  polynôme $P\in B[T]$, entraîne que les produits de $t$ par les coefficients
  de $P$ appartiennent à l'idéal $\mathfrak{n}$, et il en résulte aussitôt que
  ces coefficients sont eux-mêmes dans $\mathfrak{n}$, donc
  $P[T]\in\mathfrak{n}TB[T]$. Comme $B/tB$ est isomorphe à $\mathbb{R}$,
  $A/f_1A$ est isomorphe à l'anneau de polynômes $\mathbb{R}[T]$, donc
  intègre, et l'image de $f_2$ dans $A/f_1A$, étant égale à $T$, n'est donc
  pas diviseur de $0$, ce qui prouve notre assertion. Cependant $f_2$ est
  diviseur de $0$ dans $A$, car pour tout élément $x\in\mathfrak{n}$ non nul,
  l'image de $x$ dans $A$ est $\ne0$, mais l'image de $xT$ est nulle. On en
  conclut que la suite $(f_2,f_1)$ n'est pas régulière dans $A$; d'autre part
  l'idéal $\mathfrak{I}=f_1A+f_2A$ est distinct de $A$, donc les conditions
  b), b') et c) de (16.9.5) n'entraînent pas la condition a) lorsqu'on y
  remplace « quasi-régulière » par « régulière ».
\end{enumerate}
\end{remarks}

\begin{env}[16.9.7]
\EGAReaderTarget{ega4.c04.statement.16-9-7}{IV.16.9.7}
Si $X=\operatorname{Spec}(A)$ est un schéma affine, nous dirons qu'un idéal
$\mathfrak{I}$ de $A$ est régulier (resp. quasi-régulier) si l'Idéal
$\mathscr{I}=\widetilde{\mathfrak{I}}$ de $\mathcal{O}_X$ est régulier
(resp. quasi-régulier); on notera que cette notion est locale et n'implique
nullement l'existence d'un système de générateurs de $\mathfrak{I}$ formant
dans $A$ une suite régulière (resp. quasi-régulière) comme le montre l'exemple
(16.9.6, (i)); il en est toutefois ainsi lorsque $A$ est local (16.9.5).

La proposition (16.9.4) se traduit en termes d'immersions quasi-régulières de
la façon suivante:
\end{env}

\begin{proposition}[16.9.8]
\label{IV.16.9.8}
Soit $j:Y\to X$ un morphisme de préschémas; pour que $j$ soit une immersion
quasi-régulière, il faut et il suffit que $j$ satisfasse aux conditions
suivantes:
\begin{enumerate}
  \item[(i)] $j$ est une immersion localement de présentation finie.
  \item[(ii)] Le faisceau conormal
  $\operatorname{Gr}^1(j)=\mathcal{N}_{Y/X}$ (16.1.2) est un
  $\mathcal{O}_Y$-Module localement libre.
  \item[(iii)] L'homomorphisme canonique
  \[
    \operatorname{S}_{\mathcal{O}_Y}^{\bullet}
      (\operatorname{Gr}^1(j))\longrightarrow\operatorname{Gr}^{\bullet}(j)
  \]
  (16.1.2.2) est bijectif.
\end{enumerate}
\end{proposition}

La question étant locale sur $Y$, on peut se borner au cas où $j$ est
l'injection canonique d'un sous-préschéma fermé $Y$ de $X$, et alors la
traduction de (16.9.4) en (16.9.8) résulte de l'explicitation de
$\operatorname{Gr}^1(j)$ et $\operatorname{Gr}^{\bullet}(j)$ en termes de
l'Idéal $\mathscr{I}$ de $\mathcal{O}_X$ définissant le sous-préschéma $Y$
(16.1.3, (ii)).

\begin{corollary}[16.9.9]
\label{IV.16.9.9}
Soient $Y$ un préschéma, $X$ un $Y$-préschéma, $j:Y\to X$ une $Y$-section
de $X$, de sorte que le $n$-ème invariant normal $\mathscr{A}^{(n)}$ de $j$
(16.1.2) est une $\mathcal{O}_Y$-Algèbre augmentée (16.1.7); posons
$\mathscr{A}^{(\infty)}=\varprojlim\mathscr{A}^{(n)}$. Pour que $j$ soit une
immersion quasi-régulière, il faut et il suffit que $j$ soit localement de
présentation finie, et que tout $y\in Y$ admette un voisinage ouvert affine
$U$ d'anneau $C$ tel que $\mathscr{A}^{(\infty)}|U$ soit isomorphe comme
$\mathcal{O}_U$-Algèbre topologique augmentée à
$\mathcal{O}_U[[T_1,\ldots,T_n]]$.
\end{corollary}

On peut se borner au cas où $j$ est une immersion fermée en se restreignant
à un voisinage assez petit de $y$ (voir le raisonnement de (16.4.11)), et
alors $\mathcal{O}_Y$ s'identifie à une Algèbre quotient
$\mathcal{O}_X/\mathscr{I}$ et l'homomorphisme surjectif canonique
$\mathcal{O}_X\to\mathcal{O}_Y$ admet un
\oldpage[IV-4]{49}
inverse à droite (16.1.7). On peut donc supposer
$X=\operatorname{Spec}(B)$ et $Y=\operatorname{Spec}(A)$ affines, $B$ étant
une $A$-algèbre augmentée, et l'idéal d'augmentation $\mathfrak{I}$ étant de
type fini. Comme $\mathscr{A}^{(n)}$ s'identifie alors à
$(B/\mathfrak{I}^{n+1})^{\sim}$, le corollaire résulte de l'équivalence de b)
et c) dans (0, 19.5.4) puisque $B/\mathfrak{I}=A$.

On notera que, dans le cas affine envisagé, le fait que $j$ soit une immersion
quasi-régulière équivaut encore, en vertu de (0, 19.5.4), à dire que $B$ est
une $A$-algèbre formellement lisse pour la topologie
$\mathfrak{I}$-pré-adique.

On notera aussi que la condition que $j$ soit une immersion localement de
présentation finie est toujours vérifiée lorsque le morphisme $X\to Y$ est
localement de type fini (1.4.3, (v)).

\begin{proposition}[16.9.10]
\EGAReaderTarget{ega4.c04.prop.16-9-10}{IV.16.9.10}
Soient $X$ un préschéma localement noethérien, $Y$ un sous-préschéma de $X$,
$j:Y\to X$ l'injection canonique, $y$ un point de $Y$.
\begin{enumerate}
  \item[(i)] Pour qu'il existe un voisinage ouvert $U$ de $y$ dans $X$ tel que
  la restriction $Y\cap U\to U$ de $j$ soit une immersion régulière, il faut
  et il suffit que le noyau $\mathscr{I}_y$ de l'homomorphisme surjectif
  $\mathcal{O}_{X,y}\to\mathcal{O}_{Y,y}$ soit engendré par une suite régulière
  d'éléments de $\mathcal{O}_{X,y}$.
  \item[(ii)] Pour que l'immersion $j$ soit régulière, il faut et il suffit
  qu'elle soit quasi-régulière.
\end{enumerate}
\end{proposition}

\begin{enumerate}
  \item[(i)] On peut se borner au cas où $Y$ est un sous-préschéma fermé de
  $X$ défini par un Idéal \emph{cohérent} $\mathscr{I}$ de $\mathcal{O}_X$.
  La condition est évidemment nécessaire. Inversement, si $\mathscr{I}_y$
  est engendré par une suite régulière $(s_i)_y$, où les $s_i$ sont des
  sections de $\mathscr{I}$ au-dessus d'un voisinage ouvert $U$ de $y$ dans
  $X$, on peut supposer que les $s_i$ engendrent $\mathscr{I}|U$
  (0\textsubscript{I}, 5.2.2) et forment une suite régulière (0, 15.2.4),
  d'où l'assertion.
  \item[(ii)] Le fait qu'une immersion quasi-régulière soit régulière résulte
  de (i) et de l'identité des suites quasi-régulières et des suites régulières
  dans $\mathcal{O}_{X,y}$, formées d'éléments de l'idéal maximal
  (0, 15.1.11).
\end{enumerate}

Lorsque (sans hypothèse noethérienne sur $X$) le noyau $\mathscr{I}_y$ de
$\mathcal{O}_{X,y}\to\mathcal{O}_{Y,y}$ est engendré par une suite régulière
d'éléments de $\mathcal{O}_{X,y}$, on dit que l'immersion $j$ est
\emph{régulière au point $y$}.

\begin{corollary}[16.9.11]
\label{IV.16.9.11}
Soit $X$ un préschéma localement noethérien; alors tout Idéal quasi-régulier de
$\mathcal{O}_X$ est régulier.
\end{corollary}

\begin{remarks}[16.9.12]
\label{IV.16.9.12}
\begin{enumerate}
  \item[(i)] On notera qu'une immersion régulière n'est pas en général un
  morphisme plat, ni \emph{a fortiori} un morphisme \emph{régulier} au sens de
  (6.8.1).
  \item[(ii)] Soit $A$ un anneau local noethérien; il résulte aussitôt de
  (16.9.4) et de (0, 17.1.1) que pour que $A$ soit \emph{régulier}, il faut et
  il suffit que son idéal maximal $\mathfrak{m}$ soit \emph{quasi-régulier}
  (ou \emph{régulier}, ce qui revient au même puisque $A$ est noethérien).
  Pour qu'un schéma affine noethérien $X$ soit \emph{régulier}, il faut et il
  suffit que pour tout point fermé $x\in X$, l'injection canonique
  $\operatorname{Spec}(k(x))\to X$ soit une immersion \emph{régulière}.
\end{enumerate}
\end{remarks}

\begin{proposition}[16.9.13]
\label{IV.16.9.13}
Soient $X$ un préschéma localement noethérien, $Y$ un sous-préschéma de $X$,
$Y'$ un sous-préschéma de $Y$, tel que l'injection canonique $j:Y'\to Y$ soit
régulière. Alors la suite de $\mathcal{O}_{Y'}$-Modules
\[
  \xymatrix{
    0\ar[r] & j^*(\mathcal{N}_{Y/X})\ar[r]
      & \mathcal{N}_{Y'/X}\ar[r] & \mathcal{N}_{Y'/Y}\ar[r] & 0
  }
  \tag{16.9.13.1}
  \label{IV.16.9.13.1}
\]
\oldpage[IV-4]{50}
est exacte; en outre, pour tout $x\in X$, il existe un voisinage ouvert $U$
de $x$ tel que les restrictions à $U$ des homomorphismes de (16.9.13.1)
forment une suite exacte et scindée.
\end{proposition}

Prouvons d'abord le lemme suivant:

\begin{lemma}[16.9.13.2]
\label{IV.16.9.13.2}
Soient $A$ un anneau, $\mathfrak{I}$ un idéal de $A$,
$A'=A/\mathfrak{I}$, $(f_i)_{1\leq i\leq r}$ une suite d'éléments de $A$ qui
est $A'$-régulière, $\mathfrak{K}=\sum_i f_iA$,
$\mathfrak{L}=\mathfrak{I}+\mathfrak{K}$,
$\mathfrak{K}'=\sum_i f_iA'$, de sorte que $C=A/\mathfrak{L}$ est isomorphe à
$A'/\mathfrak{K}'$. Alors pour tout entier $n>0$ et tout entier $N\geq n$,
on a la relation
\[
  \mathfrak{I}\cap\mathfrak{K}^n
  =\mathfrak{I}\mathfrak{K}^n+(\mathfrak{I}\cap\mathfrak{K}^N).
  \tag{16.9.13.3}
  \label{IV.16.9.13.3}
\]
\end{lemma}

Il suffit évidemment de prouver que tout élément du premier membre est
contenu dans le second, et, par récurrence sur $n$, on se ramène au cas où
$N=n+1$. Un élément du premier membre de (16.9.13.3), étant dans
$\mathfrak{K}^n$, s'écrit $P(f_1,\ldots,f_r)$, où
$P\in A[T_1,\ldots,T_r]$ est homogène et de degré $n$. Si $f_i'$ est l'image
canonique de $f_i$ dans $A'$, l'hypothèse
$P(f_1,\ldots,f_r)\in\mathfrak{I}$ signifie que
$P(f_1',\ldots,f_r')=0$. Mais
$P(f_1',\ldots,f_r')\in\mathfrak{K}'^n$, donc l'image canonique de
$P(f_1',\ldots,f_r')$ dans
$\mathfrak{K}'^n/\mathfrak{K}'^{n+1}$ est nulle. Or l'hypothèse que la suite
$(f_i)$ est $A'$-régulière implique que l'homomorphisme canonique
\[
  \mathbf{S}_C^n(\mathfrak{K}'/\mathfrak{K}'^2)
  \longrightarrow\mathfrak{K}'^n/\mathfrak{K}'^{n+1}
\]
est bijectif (0, 15.1.9); on conclut que les coefficients de $P$ appartiennent
à $\mathfrak{L}=\mathfrak{I}+\mathfrak{K}$. Il en résulte aussitôt que l'on a
$P(f_1,\ldots,f_r)\in\mathfrak{I}\mathfrak{K}^n+\mathfrak{K}^{n+1}$, et comme
$P(f_1,\ldots,f_r)\in\mathfrak{I}$, on a finalement
\[
  P(f_1,\ldots,f_r)
  \in\mathfrak{I}\mathfrak{K}^n
    +(\mathfrak{I}\cap\mathfrak{K}^{n+1}),
\]
ce qui prouve le lemme.

En prenant les quotients des deux membres de (16.9.13.3) par
$\mathfrak{I}\mathfrak{K}^n$, on voit que les relations (16.9.13.3) pour
$N\geq n$ entraînent
\[
  (\mathfrak{I}\cap\mathfrak{K}^n)/\mathfrak{I}\mathfrak{K}^n
  \subset\bigcap_{N\geq n}\mathfrak{K}^N
    \cdot(A/(\mathfrak{I}\mathfrak{K}^n)).
  \tag{16.9.13.4}
  \label{IV.16.9.13.4}
\]
On en déduit le

\begin{corollary}[16.9.13.5]
\label{IV.16.9.13.5}
Supposons vérifiées les hypothèses de (16.9.13.2) et supposons en outre que
l'anneau $A$ soit noethérien et que $\mathfrak{K}$ soit contenu dans le
radical de $A$. Alors on a pour tout entier $n>0$,
\[
  \mathfrak{I}\cap\mathfrak{K}^n=\mathfrak{I}\mathfrak{K}^n.
  \tag{16.9.13.6}
  \label{IV.16.9.13.6}
\]
\end{corollary}

En effet, le second membre de (16.9.13.4) est alors nul, puisque
$A/\mathfrak{I}\mathfrak{K}^n$ est un $A$-module de type fini (Bourbaki,
\emph{Alg. comm.}, chap.~III, §~3, n\textsuperscript{o}~3, prop.~6).

Prenons en particulier $n=2$ dans (16.9.13.6), et remarquons que l'on a
$\mathfrak{L}^2=\mathfrak{I}^2+\mathfrak{I}\mathfrak{K}+\mathfrak{K}^2
=\mathfrak{I}\mathfrak{L}+\mathfrak{K}^2$; puisque
$\mathfrak{I}\mathfrak{L}\subset\mathfrak{L}^2$, on en déduit que
\[
  \mathfrak{I}\cap\mathfrak{L}^2
  =\mathfrak{I}\mathfrak{L}+(\mathfrak{I}\cap\mathfrak{K}^2)
  =\mathfrak{I}\mathfrak{L}+\mathfrak{I}\mathfrak{K}^2
  =\mathfrak{I}\mathfrak{L},
\]
autrement dit
\[
  \mathfrak{I}\cap\mathfrak{L}^2=\mathfrak{I}\mathfrak{L},
  \tag{16.9.13.7}
  \label{IV.16.9.13.7}
\]
ce qu'on peut encore exprimer en disant que l'homomorphisme canonique
\[
  \mathfrak{I}/\mathfrak{I}\mathfrak{L}
  \longrightarrow(\mathfrak{I}+\mathfrak{L}^2)/\mathfrak{L}^2
\]
est bijectif.

Ces lemmes étant démontrés, prouvons la première assertion de (16.9.13): il
suffit
\oldpage[IV-4]{51}
évidemment de prouver que la suite des fibres des faisceaux figurant dans
(16.9.13.1), en un point $x\in Y'$, est exacte. Or, si l'on pose
$A=\mathcal{O}_{X,x}$, on peut écrire
$\mathcal{O}_{Y,x}=A'=A/\mathfrak{I}$ où $\mathfrak{I}$ est un idéal contenu
dans l'idéal maximal de $A$, puis
$\mathcal{O}_{Y',x}=A'/\mathfrak{K}'$, où $\mathfrak{K}'$ est engendré par une
suite $A'$-régulière d'éléments de $A'$, eux-mêmes images des éléments d'une
suite $A'$-régulière d'éléments de $A$ appartenant à l'idéal maximal de $A$.
Si $\mathfrak{K}$ est l'idéal engendré par ces derniers et
$\mathfrak{L}=\mathfrak{I}+\mathfrak{K}$, on a
$\mathcal{O}_{Y',x}=A/\mathfrak{L}$, et comme on est dans la situation de
(16.9.13.5), l'homomorphisme canonique
\[
  \mathfrak{I}/\mathfrak{I}\mathfrak{L}
  \longrightarrow(\mathfrak{I}+\mathfrak{L}^2)/\mathfrak{L}^2
\]
est bijectif. Mais cela montre que la suite
\[
  0\longrightarrow\mathfrak{I}/\mathfrak{I}\mathfrak{L}
  \longrightarrow\mathfrak{L}/\mathfrak{L}^2
  \longrightarrow(\mathfrak{L}/\mathfrak{I})/
    (\mathfrak{L}/\mathfrak{I})^2\longrightarrow0
\]
est exacte (voir la démonstration de (16.2.7)), et les modules figurant dans
cette suite sont précisément les fibres en $x$ des faisceaux de (16.9.13.1).
La seconde assertion résulte de ce que $\mathcal{N}_{Y'/Y}$ est un
$\mathcal{O}_{Y'}$-Module localement \emph{libre} (16.9.8) et de Bourbaki,
\emph{Alg.}, chap.~II, 3\textsuperscript{e} éd., §~1, n\textsuperscript{o}~11,
prop.~21.

\subsection{Morphismes différentiellement lisses}
\EGAReaderTarget{ega4.c04.subsec.16-10}{IV.16.10}

\begin{definition}[16.10.1]
\EGAReaderTarget{ega4.c04.def.16-10-1}{IV.16.10.1}
On dit qu'un morphisme de préschémas $f:X\to S$ est
\emph{différentiellement lisse} (ou que $X$ est \emph{différentiellement
lisse sur $S$}) s'il vérifie les conditions suivantes:
\begin{enumerate}
  \item[(i)] $\Omega_{X/S}^1$ est un $\mathcal{O}_X$-Module localement
  projectif, c'est-à-dire que tout point de $X$ admet un voisinage ouvert
  affine $U$ tel que $\Gamma(U,\Omega_{X/S}^1)$ soit un
  $\Gamma(U,\mathcal{O}_X)$-module projectif (non nécessairement de type
  fini).
  \item[(ii)] L'homomorphisme canonique (16.3.1.1)
  \[
    \mathbf{S}_{\mathcal{O}_X}^{\bullet}(\Omega_{X/S}^1)
    \longrightarrow\mathscr{G}r_{\bullet}(\mathscr{P}_{X/S})
  \]
  est bijectif.
\end{enumerate}
En particulier, si $\Omega_{X/S}^1$ est localement libre de rang fini, les
$\mathscr{P}_{X/S}^n$ sont des $\mathcal{O}_X$-Modules localement libres de
rang fini (étant des extensions de tels Modules).
\end{definition}

On dit que $f$ est \emph{différentiellement lisse en un point $x\in X$} (ou
que $X$ est différentiellement lisse sur $S$ au point $x$) s'il existe un
voisinage ouvert $U$ de $x$ dans $X$ tel que $f|U$ soit différentiellement
lisse.

Nous verrons plus loin (17.12.4) qu'un morphisme lisse est différentiellement
lisse, ce qui justifie la terminologie; mais la réciproque est inexacte; en
effet, un \emph{monomorphisme} $f:X\to S$ est différentiellement lisse,
puisque $\Omega_{X/S}^1=0$ en vertu de (I, 5.3.8), et par suite
l'homomorphisme surjectif (16.3.1.1) est évidemment bijectif; or un
monomorphisme n'est même pas nécessairement plat, ni \emph{a fortiori} lisse.
Bornons-nous ici à noter la proposition suivante:

\begin{proposition}[16.10.2]
\label{IV.16.10.2}
Soient $A$ un anneau, $B$ une $A$-algèbre formellement lisse pour les
topologies discrètes (0, 19.3.1). Alors $\operatorname{Spec}(B)$ est
différentiellement lisse sur $\operatorname{Spec}(A)$.
\end{proposition}

En effet, $B\otimes_A B$ est alors (pour les topologies discrètes) une
$B$-algèbre formellement lisse (pour l'un ou l'autre des homomorphismes
canoniques $b\mapsto b\otimes1$, $b\mapsto1\otimes b$ de $B$
\oldpage[IV-4]{52}
dans $B\otimes_A B$) (0, 19.3.5, (iii)); donc $B\otimes_A B$ est aussi une
$A$-algèbre formellement lisse pour les topologies discrètes (0, 19.3.5,
(ii)). Posant $\mathfrak{I}=\mathfrak{I}_{B/A}$, il en résulte que
$B\otimes_A B$ est aussi une $A$-algèbre formellement lisse pour la topologie
$\mathfrak{I}$-préadique (0, 19.3.8); comme par hypothèse
$B=(B\otimes_A B)/\mathfrak{I}$ est une $A$-algèbre formellement lisse pour
les topologies discrètes, la proposition résulte de l'équivalence de a) et b)
dans (0, 19.5.4).

\begin{proposition}[16.10.3]
\label{IV.16.10.3}
Pour qu'un morphisme $f:X\to S$ soit différentiellement lisse, il faut et il
suffit que pour tout $x\in X$, il existe un voisinage ouvert affine $U$ de
$x$, d'anneau $A$, tel que $\Gamma(U,\mathscr{P}_{X/S}^{\infty})$ soit une
$A$-algèbre topologique augmentée isomorphe à l'algèbre complétée $\widehat B$,
où $B=\mathbf{S}_A(V)$, $V$ étant un $A$-module projectif et $B$ étant muni de
la topologie $B^+$-préadique (où $B^+$ est l'idéal d'augmentation). Si
$\Omega_{X/S}^1$ est localement libre de rang fini, on peut remplacer
$\widehat B$ par l'algèbre de séries formelles $A[[T_1,\ldots,T_n]]$.
\end{proposition}

La notion de morphisme différentiellement lisse étant évidemment locale sur
$X$, on peut se borner au cas où $S=\operatorname{Spec}(B)$,
$X=\operatorname{Spec}(C)$. Considérons $C\otimes_B C$ comme une $C$-algèbre
(pour le premier facteur); posons $\mathfrak{I}=\mathfrak{I}_{C/B}$ et
munissons $C\otimes_B C$ de la topologie $\mathfrak{I}$-préadique; on peut
appliquer à la $C$-algèbre topologique $C\otimes_B C$ et à l'idéal
$\mathfrak{I}$ de $C\otimes_B C$ l'équivalence de b) et c) dans (0, 19.5.4),
puisque $(C\otimes_B C)/\mathfrak{I}=C$ est évidemment une $C$-algèbre
formellement lisse pour les topologies discrètes. La topologie sur
$\Gamma(U,\mathscr{P}_{X/S}^{\infty})$ est évidemment la topologie de limite
projective de cet anneau (16.1.11).

On notera que l'entier $n$ de l'énoncé de (16.10.3) est le \emph{rang} de
$\Omega_{X/S}^1$ au point $x$. Nous verrons plus loin (17.13.5) que lorsque
$f$ est différentiellement lisse et localement de type fini, $n$ est aussi
égal à la dimension de la fibre $f^{-1}(f(x))$ au point $x$.

\begin{proposition}[16.10.4]
\label{IV.16.10.4}
Soient $f:X\to S$, $g:S'\to S$ deux morphismes, et posons
$X'=X\times_S S'$, $f'=f_{(S')}:X'\to S'$.
\begin{enumerate}
  \item[(i)] Si $f$ est différentiellement lisse, il en est de même de $f'$.
  \item[(ii)] Inversement, si $g$ est fidèlement plat et quasi-compact, et si
  $f'$ est différentiellement lisse et $\Omega_{X'/S'}^1$ un
  $\mathcal{O}_{X'}$-Module de type fini, $f$ est différentiellement lisse et
  $\Omega_{X/S}^1$ est un $\mathcal{O}_X$-Module de type fini.
\end{enumerate}
\end{proposition}

En effet, si $f$ est différentiellement lisse, les
$\mathscr{G}r_n(\mathscr{P}_{X/S})$ sont des $\mathcal{O}_X$-Modules
\emph{plats}; par suite (16.4.6), l'homomorphisme
\[
  \mathscr{G}r_n(\mathscr{P}_{X/S})\otimes_{\mathcal{O}_X}\mathcal{O}_{X'}
  \longrightarrow\mathscr{G}r_n(\mathscr{P}_{X'/S'})
\]
est bijectif pour tout $n$, et en vertu de la commutativité du diagramme
(16.2.1.3), il résulte de la définition (16.10.1) que $f'$ est
différentiellement lisse. D'autre part, si $g$ est fidèlement plat et
quasi-compact, il résulte encore de (16.4.6) que
\[
  \mathscr{G}r_n(\mathscr{P}_{X/S})\otimes_{\mathcal{O}_X}\mathcal{O}_{X'}
  \longrightarrow\mathscr{G}r_n(\mathscr{P}_{X'/S'})
\]
est bijectif pour tout $n$. Supposons alors $f'$ différentiellement lisse et
$\Omega_{X'/S'}^1$ de rang fini. Comme la projection canonique $X'\to X$ est
un morphisme fidèlement plat et quasi-compact, il résulte d'abord de (2.5.2)
que $\Omega_{X/S}^1$ est un $\mathcal{O}_X$-Module localement libre de rang
fini, puis de (2.2.7) que l'homomorphisme canonique (16.3.1.1) est bijectif,
donc $f$ est différentiellement lisse.

\begin{proposition}[16.10.5]
\label{IV.16.10.5}
Pour qu'un morphisme localement de type fini $f:X\to S$ soit
différentiellement lisse, il faut et il suffit que l'immersion diagonale
$\Delta_f:X\to X\times_S X$ soit quasi-régulière.
\end{proposition}

\oldpage[IV-4]{53}

La question étant locale, on peut se borner au cas où $S$ et $X$ sont
affines, et par suite le sous-préschéma diagonal de $X\times_S X$ est fermé.
L'hypothèse que $f$ est localement de type fini entraîne que $\Delta_f$ est
localement de présentation finie (I, 4.3.1), donc le sous-préschéma diagonal
de $X\times_S X$ est défini par un Idéal de type fini $\mathscr{I}$, et
$\Omega_{X/S}^1=\mathscr{I}/\mathscr{I}^2$ est un $\mathcal{O}_X$-Module de
type fini. La proposition résulte alors aussitôt de la comparaison des
conditions dans (16.10.1) et (16.9.4).

\begin{remark}[16.10.6]
\label{IV.16.10.6}
Soit $f:X\to S$ un morphisme tel que le $\mathcal{O}_X$-Module
$\Omega_{X/S}^1$ soit localement libre de rang fini. Il résulte alors de
(0, 20.4.7) que tout $x\in X$ possède un voisinage ouvert $U$ tel qu'il existe
une famille finie $(z_\lambda)_{\lambda\in L}$ de sections de
$\mathcal{O}_X$ au-dessus de $U$ pour lesquelles $(dz_\lambda)_{\lambda\in L}$
forme une \emph{base} du $\Gamma(U,\mathcal{O}_X)$-module
$\Gamma(U,\Omega_{X/S}^1)$.
\end{remark}

\subsection{\texorpdfstring{Opérateurs différentiels sur un $S$-préschéma différentiellement
lisse}{Opérateurs différentiels sur un S-préschéma différentiellement lisse}}
\EGAReaderTarget{ega4.c04.subsec.16-11}{IV.16.11}

\begin{env}[16.11.1]
\label{IV.16.11.1}
Soient $f:X\to S$ un morphisme, $U$ un ouvert de $X$,
$(z_\lambda)_{\lambda\in L}$ une famille de sections de $\mathcal{O}_X$
au-dessus de $U$ telle que les $dz_\lambda$ forment un système de générateurs
de $\Omega_{X/S}^1|U=\Omega_{U/S}^1$. Soit $m$ un entier ou le symbole
$\infty$, et posons, pour tout $\lambda$
\[
  \zeta_\lambda=\delta z_\lambda=d^m z_\lambda-z_\lambda
  \in\Gamma(U,\mathscr{P}_{X/S}^m).
  \tag{16.11.1.1}
  \label{IV.16.11.1.1}
\]

Nous utiliserons d'autre part les notations usuelles de l'analyse; pour tout
$\mathbf{p}=(p_\lambda)\in\mathbb{N}^{(L)}$ (avec $p_\lambda=0$ sauf pour un
nombre fini d'indices), nous poserons
\[
  |\mathbf{p}|=\sum_\lambda p_\lambda,
  \qquad \mathbf{p}!=\prod_\lambda(p_\lambda!)
  \tag{16.11.1.2}
  \label{IV.16.11.1.2}
\]
\[
  \binom{\mathbf{p}}{\mathbf{q}}
  =\mathbf{p}!/(\mathbf{q}!(\mathbf{p}-\mathbf{q})!)
  \quad\text{pour $\mathbf{p},\mathbf{q}$ dans $\mathbb{N}^{(L)}$,
  $\mathbf{q}\leq\mathbf{p}$}
  \tag{16.11.1.3}
  \label{IV.16.11.1.3}
\]
et on convient que $\binom{\mathbf{p}}{\mathbf{q}}=0$ si
$\mathbf{q}\nleq\mathbf{p}$,
\[
  \mathbf{z}^{\mathbf{p}}=\prod_\lambda(z_\lambda)^{p_\lambda},
  \qquad
  \boldsymbol{\zeta}^{\mathbf{p}}=\prod_\lambda(\zeta_\lambda)^{p_\lambda}.
  \tag{16.11.1.4}
  \label{IV.16.11.1.4}
\]

On aura donc, avec ces notations
\[
  d^m(\mathbf{z}^{\mathbf{p}})
  =(d^m(\mathbf{z}))^{\mathbf{p}}
  =(\boldsymbol{\zeta}+\mathbf{z})^{\mathbf{p}}
  =\sum_{\mathbf{q}\leq\mathbf{p}}
    \binom{\mathbf{p}}{\mathbf{q}}
    \mathbf{z}^{\mathbf{p}-\mathbf{q}}\boldsymbol{\zeta}^{\mathbf{q}}
  \tag{16.11.1.5}
  \label{IV.16.11.1.5}
\]
\[
  \boldsymbol{\zeta}^{\mathbf{p}}
  =(d^m\mathbf{z}-\mathbf{z})^{\mathbf{p}}
  =\sum_{\mathbf{q}\leq\mathbf{p}}
    (-1)^{|\mathbf{p}-\mathbf{q}|}
    \binom{\mathbf{p}}{\mathbf{q}}
    \mathbf{z}^{\mathbf{p}-\mathbf{q}}d^m(\mathbf{z}^{\mathbf{q}}).
  \tag{16.11.1.6}
  \label{IV.16.11.1.6}
\]

Comme les $dz_\lambda$ engendrent $\Omega_{X/S}^1$ et sont les images des
$\delta z_\lambda$, et que l'homomorphisme canonique (16.3.1.1) est
surjectif, on en conclut que pour $m$ fini, les $\delta z_\lambda$ engendrent
la $\mathcal{O}_U$-Algèbre $\mathscr{P}_{U/S}^m$ (Bourbaki,
\emph{Alg. comm.}, chap.~III, §~2, n\textsuperscript{o}~8, cor.~2 du th.~1).
Donc les $\boldsymbol{\zeta}^{\mathbf{p}}$ (pour $|\mathbf{p}|\leq m$)
engendrent le $\mathcal{O}_U$-Module $\mathscr{P}_{U/S}^m$. Un opérateur
différentiel $D\in\operatorname{Diff}_{U/S}^m$ est par suite entièrement
déterminé par les valeurs des
$\langle\boldsymbol{\zeta}^{\mathbf{p}},D\rangle$ pour
$|\mathbf{p}|\leq m$, ou, ce qui revient au même par (16.11.1.5) et
(16.11.1.6), par les valeurs
\oldpage[IV-4]{54}
des $\langle d^m(\mathbf{z}^{\mathbf{p}}),D\rangle=D(\mathbf{z}^{\mathbf{p}})$
pour $|\mathbf{p}|\leq m$; de façon précise, il résulte de (16.11.1.5) que
l'on a
\[
  D(\mathbf{z}^{\mathbf{p}})
  =\langle d^m(\mathbf{z}^{\mathbf{p}}),D\rangle
  =\sum_{\mathbf{q}\leq\mathbf{p}}
    \binom{\mathbf{p}}{\mathbf{q}}
    \langle\boldsymbol{\zeta}^{\mathbf{q}},D\rangle
    \mathbf{z}^{\mathbf{p}-\mathbf{q}}.
  \tag{16.11.1.7}
  \label{IV.16.11.1.7}
\]
\end{env}

\begin{theorem}[16.11.2]
\label{IV.16.11.2}
Soient $f:X\to S$ un morphisme, $U$ un ouvert de $X$,
$(z_\lambda)_{\lambda\in L}$ une famille de sections de $\mathcal{O}_X$
au-dessus de $U$ telle que la famille $(dz_\lambda)_{\lambda\in L}$ engendre
$\Omega_{X/S}^1|U=\Omega_{U/S}^1$. Les conditions suivantes sont
équivalentes:
\begin{enumerate}
  \item[a)] $f|U$ est différentiellement lisse et $(dz_\lambda)$ est une base
  du $\mathcal{O}_U$-Module $\Omega_{U/S}^1$.
  \item[b)] Il existe une famille
  $(D_{\mathbf{p}})_{\mathbf{p}\in\mathbb{N}^{(L)}}$ d'opérateurs
  différentiels de $\mathcal{O}_U$ dans lui-même vérifiant les conditions
  \[
    D_{\mathbf{p}}(\mathbf{z}^{\mathbf{q}})
    =\binom{\mathbf{q}}{\mathbf{p}}
      \mathbf{z}^{\mathbf{q}-\mathbf{p}}
    \quad(\mathbf{p},\mathbf{q}\text{ dans }\mathbb{N}^{(L)}).
    \tag{16.11.2.1}
    \label{IV.16.11.2.1}
  \]
\end{enumerate}
De plus, lorsque ces conditions sont vérifiées, la famille
$(D_{\mathbf{p}})$ est déterminée de façon unique par les conditions
(16.11.2.1) et vérifie les relations
\[
  D_{\mathbf{p}}\circ D_{\mathbf{q}}
  =D_{\mathbf{q}}\circ D_{\mathbf{p}}
  =\frac{(\mathbf{p}+\mathbf{q})!}{\mathbf{p}!\mathbf{q}!}
    D_{\mathbf{p}+\mathbf{q}}
  \quad(\mathbf{p},\mathbf{q}\text{ dans }\mathbb{N}^{(L)}).
  \tag{16.11.2.2}
  \label{IV.16.11.2.2}
\]
Enfin, si $L$ est fini, pour tout entier $m$, les $D_{\mathbf{p}}$ tels que
$|\mathbf{p}|\leq m$ forment une base du $\mathcal{O}_U$-Module
$\operatorname{Diff}_{U/S}^m$, autrement dit, tout opérateur différentiel
d'ordre $\leq m$ sur $U$ s'écrit d'une seule façon sous la forme
\[
  D=\sum_{|\mathbf{p}|\leq m}a_{\mathbf{p}}D_{\mathbf{p}}
\]
où les $a_{\mathbf{p}}$ sont des sections de $\mathcal{O}_X$ au-dessus de
$U$.
\end{theorem}

Notons d'abord qu'en vertu de (16.11.1.6) et (16.11.1.5), on vérifie aussitôt
que les conditions (16.11.2.1) sont équivalentes à
\[
  \langle\boldsymbol{\zeta}^{\mathbf{p}},D_{\mathbf{q}}\rangle
  =\delta_{\mathbf{p}\mathbf{q}}
  \quad\text{(indice de Kronecker)}.
  \tag{16.11.2.3}
  \label{IV.16.11.2.3}
\]
L'existence de la famille $(D_{\mathbf{p}})$ vérifiant ces relations entraîne
donc d'abord (en prenant $|\mathbf{p}|=1$) que les $dz_\lambda$ sont
linéairement indépendants, donc forment une \emph{base} du
$\mathcal{O}_U$-Module $\Omega_{U/S}^1$. Puis, pour tout entier $m\geq1$, on
déduit de même de (16.11.2.3) que les $\boldsymbol{\zeta}^{\mathbf{p}}$ tels
que $|\mathbf{p}|\leq m$ sont linéairement indépendants; par suite
l'homomorphisme canonique (16.3.1.1) est injectif, donc bijectif, et cela
prouve que b) entraîne a). La réciproque découle aussitôt de la définition
(16.10.1), le fait que les $\boldsymbol{\zeta}^{\mathbf{p}}$ forment une base
de $\mathscr{P}_{U/S}^m$ pour $|\mathbf{p}|\leq m$ entraînant l'existence et
l'unicité d'une famille d'homomorphismes
$u_{\mathbf{q},m}:\mathscr{P}_{U/S}^m\to\mathcal{O}_U$
($|\mathbf{q}|\leq m$) tels que
$\langle\boldsymbol{\zeta}^{\mathbf{p}},u_{\mathbf{q},m}\rangle
=\delta_{\mathbf{p}\mathbf{q}}$ pour $|\mathbf{p}|\leq m$,
$|\mathbf{q}|\leq m$. Pour une valeur de $\mathbf{q}$ donnée, les opérateurs
différentiels correspondant aux $u_{\mathbf{q},m}$ pour
$m\geq|\mathbf{q}|$ sont identifiés à un même opérateur $D_{\mathbf{q}}$.
Cela prouve donc que a) entraîne b), et en outre que la famille
$(D_{\mathbf{p}})$ est déterminée de façon unique et que, si $L$ est fini,
pour $|\mathbf{p}|\leq m$, les $D_{\mathbf{p}}$ forment une base du dual
$\operatorname{Diff}_{U/S}^m$ de $\mathscr{P}_{U/S}^m$. Enfin, les relations
(16.11.2.2) découlent aussitôt de l'expression des valeurs des trois
opérateurs considérés pour les $\mathbf{z}^{\mathbf{r}}$, et du fait que les
$\boldsymbol{\zeta}^{\mathbf{r}}$ pour $|\mathbf{r}|\leq m$ engendrent
$\mathscr{P}_{U/S}^m$.

\oldpage[IV-4]{55}

\begin{remarks}[16.11.3]
\EGAReaderTarget{ega4.c04.statement.16-11-3}{IV.16.11.3}
\begin{enumerate}
  \item[(i)] Le fait que les $D_{\mathbf{p}}$ sont deux à deux permutables en
  vertu de (16.11.2.2) n'implique naturellement pas que la
  $\mathcal{O}_U$-Algèbre $\operatorname{Diff}_{U/S}$ soit commutative, les
  $D_{\mathbf{p}}$ ne permutant aux produits par les sections de
  $\mathcal{O}_U$ que si $n=0$.
  \item[(ii)] Les indices $\mathbf{p}$ tels que $|\mathbf{p}|=1$ sont les
  $\boldsymbol{\varepsilon}_\lambda
  =(\varepsilon_{\lambda\mu})_{\mu\in L}$ avec
  $\varepsilon_{\lambda\mu}=0$ si $\mu\neq\lambda$ et
  $\varepsilon_{\lambda\lambda}=1$; lorsque $L$ est fini, les opérateurs
  $D_{\boldsymbol{\varepsilon}_\lambda}$ ne sont autres que les
  $S$-dérivations $D_i$ introduites dans (16.5.7). On notera qu'en général
  (et contrairement à ce qui se passe en Analyse classique), il n'est pas
  vrai qu'un opérateur différentiel d'ordre quelconque puisse s'écrire comme
  combinaison linéaire de puissances des $D_i$ (cf. (16.12)).
  \item[(iii)] Pour tout entier $r\geq1$, on peut définir la notion de
  morphisme \emph{différentiellement lisse jusqu'à l'ordre $r$} en remplaçant
  dans (16.10.1) la condition (ii) par la condition que les homomorphismes
  \[
    \mathbf{S}_{\mathcal{O}_X}^m(\Omega_{X/S}^1)
    \longrightarrow\mathscr{G}r_m(\mathscr{P}_{X/S})
  \]
  soient bijectifs \emph{pour tout $m\leq r$}. Le raisonnement de (16.11.2)
  prouve alors que si, dans la condition a), on remplace «différentiellement
  lisse» par «différentiellement lisse jusqu'à l'ordre $r$», cette condition
  est équivalente à la condition b) dans laquelle on se borne aux
  $\mathbf{p}\in\mathbb{N}^{(L)}$, $\mathbf{q}\in\mathbb{N}^{(L)}$ tels que
  $|\mathbf{p}|\leq r$, $|\mathbf{q}|\leq r$.
\end{enumerate}
\end{remarks}

\subsection{Cas de la caractéristique nulle: critère jacobien pour les
morphismes différentiellement lisses}
\EGAReaderTarget{ega4.c04.subsec.16-12}{IV.16.12}

\begin{env}[16.12.1]
\EGAReaderTarget{ega4.c04.statement.16-12-1}{IV.16.12.1}
On dit qu'un préschéma $X$ est \emph{de caractéristique $p$} ($p$ égal à $0$
ou à un nombre premier) si, pour tout ouvert affine $U$ de $X$, l'anneau
$\Gamma(U,\mathcal{O}_X)$ est de caractéristique $p$ (0, 21.1.1). Il résulte
de (0, 21.1.3) que pour que $X$ soit de caractéristique $0$, il faut et il
suffit que pour tout point \emph{fermé} $x$ de $X$, le corps résiduel $k(x)$
soit de caractéristique $0$, ou encore que $X$ puisse être muni d'une
structure de $\mathbb{Q}$-préschéma (nécessairement unique).
\end{env}

\begin{theorem}[16.12.2]
\label{IV.16.12.2}
Soient $X$ un préschéma de caractéristique $0$, $f:X\to S$ un morphisme. Si
$\Omega_{X/S}^1$ est un $\mathcal{O}_X$-Module localement libre (non
nécessairement de type fini), $f$ est différentiellement lisse.
\end{theorem}

La question étant locale sur $X$, on peut supposer qu'il existe une famille
$(z_\lambda)$ de sections de $\mathcal{O}_X$ au-dessus de $X$ telle que
$(dz_\lambda)$ soit une base du $\mathcal{O}_X$-Module $\Omega_{X/S}^1$.
Appliquant le critère (16.11.2), il suffit de remarquer que les opérateurs
\[
  D_{\mathbf{p}}=(\mathbf{p}!)^{-1}
  \prod_\lambda D_\lambda^{p_\lambda}
\]
(où les $D_\lambda$ sont les formes coordonnées correspondant à la base
$(dz_\lambda)$) vérifient les relations (16.11.2.1), ce qui est une
conséquence du fait que les $D_\lambda$ sont des dérivations.

\begin{env}[16.12.3]
\label{IV.16.12.3}
Le théorème précédent n'est plus exact lorsqu'on abandonne l'hypothèse que
$X$ est de caractéristique $0$. Par exemple, si
$S=\operatorname{Spec}(k)$, où $k$ est un corps de caractéristique $p>0$,
$X=\operatorname{Spec}(K)$ où $K=k(\alpha)$ avec
$\alpha\notin k$, $\alpha^p\in k$, on vérifie aussitôt
\oldpage[IV-4]{56}
que $\Omega_{X/S}^1$ est de rang $1$, et que le morphisme $X\to S$ est
différentiellement lisse jusqu'à l'ordre $p-1$ (16.11.3, (iii)), mais non
jusqu'à l'ordre $p$. Cependant, la démonstration de (16.12.2) prouve que si
$\Omega_{X/S}^1$ est localement libre, et si $n!1_{\mathcal{O}_X}$ est
inversible dans $\Gamma(X,\mathcal{O}_X)$, alors $X$ est différentiellement
lisse sur $S$ jusqu'à l'ordre $n$.
\end{env}

% FIN SOURCE ega4/ega4-16-fr.tex

%
% DEBUT SOURCE ega4/ega4-17-fr.tex
\section[Morphismes lisses, morphismes non ramifiés (ou nets) morphismes étales]{Morphismes lisses, morphismes non ramifiés (ou nets)\\morphismes étales}
\label{section:IV.17}

Dans le présent paragraphe, nous reprenons les notions étudiées dans (0, 19), exprimées à l'aide du langage géométrique des schémas et du point de vue global, pour les préschémas localement de présentation finie sur un préschéma de base donné. La plupart des résultats (à l'exception des n\textsuperscript{os} 17.7, 17.8, 17.9, 17.13 et 17.16) se réduisent d'ailleurs à des variantes de propriétés déjà rencontrées dans (0, 19). Pour des résultats plus spéciaux sur les morphismes étales, le lecteur consultera le § 18.

\subsection{Morphismes formellement lisses, morphismes formellement non ramifiés, morphismes formellement étales}
\EGAReaderTarget{ega4.c04.subsec.17-1}{subsection:IV.17.1}

\begin{definition}[17.1.1]
\EGAReaderTarget{ega4.c04.def.17-1-1}{IV.17.1.1}
Soit $f:X\to Y$ un morphisme de préschémas. On dit que $f$ est \emph{formellement lisse} (resp. \emph{formellement non ramifié} ou \emph{formellement net}, resp. \emph{formellement étale}) si, pour tout schéma affine $Y'$, tout sous-schéma fermé $Y'_0$ de $Y'$ défini par un Idéal nilpotent $\sh{J}$ de $\sh{O}_{Y'}$, et tout morphisme $Y'\to Y$, l'application
\[
\label{IV.17.1.1.1}
  \Hom_Y(Y',X)\longrightarrow\Hom_Y(Y'_0,X)
  \tag{17.1.1.1}
\]
déduite de l'injection canonique $Y'_0\to Y'$, est \emph{surjective} (resp. \emph{injective}, resp. \emph{bijective}).
\end{definition}

On dit encore alors que $X$ est \emph{formellement lisse} (resp. \emph{formellement non ramifié} ou \emph{formellement net}, resp. \emph{formellement étale}) sur $Y$.

Il est clair que dire que $f$ est formellement étale signifie qu'il est \emph{à la fois} formellement lisse et formellement non ramifié.

\begin{remarks}[17.1.2]
\label{IV.17.1.2}
\begin{enumerate}
  \item[(i)] Supposons que $Y=\Spec(A)$ et $X=\Spec(B)$ soient affines, de sorte que $f$ provient d'un homomorphisme d'anneaux $\varphi:A\to B$. En vertu de (0, 19.3.1 et 0, 19.10.1), dire que $f$ est formellement lisse (resp. formellement non ramifié, resp. formellement étale) signifie que $\varphi$ fait de $B$ une $A$-algèbre \emph{formellement lisse} (resp. \emph{formellement non ramifiée}, resp. \emph{formellement étale}) pour les topologies \emph{discrètes} sur $A$ et $B$.
  \item[(ii)] Pour vérifier que $f$ est formellement lisse (resp. formellement non ramifié, resp. formellement étale), on peut, dans la définition (17.1.1) se borner au cas où $\sh{J}^2=0$. En effet, si $f$ vérifie la condition correspondante de la définition (17.1.1) dans ce cas particulier, et si l'on a $\sh{J}^n=0$, on considère le sous-schéma fermé $Y'_j$ de $Y'$ défini par l'Idéal $\sh{J}^{j+1}$ pour $0\leq j\leq n-1$, de sorte que $Y'_j$ est un sous-schéma fermé de $Y'_{j+1}$ défini par un Idéal de carré nul; l'hypothèse entraîne que chacune des applications
  \[
    \Hom_Y(Y'_{j+1},X)\longrightarrow\Hom_Y(Y'_j,X)qquad(0\leq j\leq n-1)
  \]
  \oldpage[IV-4]{57}
  est surjective (resp. injective, resp. bijective); par composition, on en conclut qu'il en est de même de (17.1.1.1).
  \item[(iii)] On notera que les propriétés du morphisme $f$ définies dans (17.1.1) sont des propriétés du \emph{foncteur représentable} $(0_{\mathrm{III}}, 8.1.8)$
  \[
    Y'\longmapsto\Hom_Y(Y',X)
  \]
  de la catégorie des $Y$-préschémas dans la catégorie des ensembles; elles gardent un sens pour \emph{n'importe quel} foncteur contravariant ayant mêmes catégories de départ et d'arrivée, représentable ou non.
  \item[(iv)] Supposons que le morphisme $f$ soit formellement non ramifié (resp. formellement étale); considérons un $Y$-préschéma \emph{quelconque} $Z$ et un sous-préschéma fermé $Z_0$ de $Z$ défini par un Idéal \emph{localement nilpotent} $\sh{J}$ de $\sh{O}_Z$. Alors l'application
  \[
  \label{IV.17.1.2.1}
    \Hom_Y(Z,X)\longrightarrow\Hom_Y(Z_0,X)
    \tag{17.1.2.1}
  \]
  déduite de l'injection canonique $Z_0\to Z$, est encore injective (resp. bijective). En effet, soit $(U_\alpha)$ un recouvrement ouvert affine de $Z$ tel que les Idéaux $\sh{J}|U_\alpha$ soient nilpotents, et pour tout $\alpha$, soit $U^0_\alpha$ l'image réciproque de $U_\alpha$ dans $Z_0$, qui est le sous-schéma fermé de $U_\alpha$ défini par $\sh{J}|U_\alpha$. Soit $f_0:Z_0\to X$ un $Y$-morphisme; par hypothèse, pour tout $\alpha$, il y a au plus un (resp. un et un seul) $Y$-morphisme $f_\alpha:U_\alpha\to X$ dont la restriction à $U^0_\alpha$ soit égale à $f_0|U^0_\alpha$. On en conclut aussitôt que si $f_\alpha$ et $f_\beta$ sont définis, alors, pour tout ouvert affine $V\subset U_\alpha\cap U_\beta$, on a $f_\alpha|V=f_\beta|V$, puisque les restrictions de ces morphismes à l'image réciproque $V_0$ de $V$ dans $Z_0$ coïncident. Il y a donc au plus un (resp. un et un seul) $Y$-morphisme $f:Z\to X$ dont la restriction à $Z_0$ coïncide avec $f_0$.
\end{enumerate}
\end{remarks}

\begin{proposition}[17.1.3]
\label{IV.17.1.3}
\begin{enumerate}
  \item[(i)] Un monomorphisme de préschémas est formellement non ramifié; une immersion ouverte est formellement étale.
  \item[(ii)] Le composé de deux morphismes formellement lisses (resp. formellement non ramifiés, resp. formellement étales) est formellement lisse (resp. formellement non ramifié, resp. formellement étale).
  \item[(iii)] Si $f:X\to Y$ est un $S$-morphisme formellement lisse (resp. formellement non ramifié, resp. formellement étale), il en est de même de $f_{(S')}:X_{(S')}\to Y_{(S')}$ pour toute extension $S'\to S$ du préschéma de base.
  \item[(iv)] Si $f:X\to X'$ et $g:Y\to Y'$ sont deux $S$-morphismes formellement lisses (resp. formellement non ramifiés, resp. formellement étales), il en est de même de $f\times_S g:X\times_S Y\to X'\times_S Y'$.
  \item[(v)] Soient $f:X\to Y$, $g:Y\to Z$ deux morphismes; si $g\circ f$ est formellement non ramifié, il en est de même de $f$.
  \item[(vi)] Si $f:X\to Y$ est un morphisme formellement non ramifié, il en est de même de $f_{\mathrm{red}}:X_{\mathrm{red}}\to Y_{\mathrm{red}}$.
\end{enumerate}
\end{proposition}

En vertu de (I, 5.5.12), il suffit de prouver (i), (ii) et (iii). Les deux assertions de (i) sont triviales. Pour prouver (ii), considérons deux morphismes $f:X\to Y$, $g:Y\to Z$, un schéma affine $Z'$, un sous-schéma fermé $Z'_0$ de $Z'$ défini par un Idéal nilpotent et un morphisme $Z'\to Z$. Supposons $f$ et $g$ formellement lisses, et considérons un
\oldpage[IV-4]{58}
$Z$-morphisme $u_0:Z'_0\to X$; l'hypothèse sur $g$ entraîne qu'il existe un $Z$-morphisme $v:Z'\to Y$ tel que $f\circ u_0=v\circ j$ (où $j:Z'_0\to Z'$ est l'injection canonique); l'hypothèse sur $f$ entraîne alors qu'il existe un morphisme $u:Z'\to X$ tel que $f\circ u=v$ et $u\circ j=u_0$, donc $(g\circ f)\circ u$ est égal au morphisme donné $Z'\to Z$ et $u\circ j=u_0$, ce qui prouve que $g\circ f$ est formellement lisse; on raisonne de même lorsque l'on suppose $f$ et $g$ formellement non ramifiés.

Enfin, pour démontrer (iii), posons $X'=X_{(S')}$, $Y'=Y_{(S')}$, $f'=f_{(S')}$; considérons un schéma affine $Y''$, un sous-schéma fermé $Y''_0$ de $Y''$ défini par un Idéal nilpotent et un morphisme $g:Y''\to Y'$ faisant de $Y''$ un $Y'$-préschéma; on sait alors (I, 3.3.8) que $\Hom_{Y'}(Y'',X')$ s'identifie canoniquement à $\Hom_Y(Y'',X)$ et $\Hom_{Y'}(Y''_0,X')$ à $\Hom_Y(Y''_0,X)$, et la conclusion résulte alors immédiatement de la définition (17.1.1).

On notera qu'une \emph{immersion fermée} n'est pas nécessairement un morphisme formellement lisse.

\begin{proposition}[17.1.4]
\label{IV.17.1.4}
Soient $f:X\to Y$, $g:Y\to Z$ deux morphismes, et supposons $g$ formellement non ramifié. Alors, si $g\circ f$ est formellement lisse (resp. formellement étale), il en est de même de $f$.
\end{proposition}

En effet, soient $Y'$ un schéma affine, $Y'_0$ un sous-schéma fermé de $Y'$ défini par un Idéal nilpotent, $h:Y'\to Y$ un morphisme, $j:Y'_0\to Y'$ l'injection canonique, $u_0:Y'_0\to X$ un $Y$-morphisme, donc tel que $f\circ u_0=h\circ j$. Supposons $g\circ f$ formellement lisse; alors il existe un morphisme $u:Y'\to X$ tel que $u\circ j=u_0$ et $(g\circ f)\circ u=g\circ h$. Mais ces relations entraînent que $f\circ u$ et $h$ sont deux $Z$-morphismes de $Y'$ dans $Y$ tels que $(f\circ u)\circ j=h\circ j$; en vertu de l'hypothèse que $g$ est formellement non ramifié, on en tire que $f\circ u=h$, autrement dit $u$ est un $Y$-morphisme; donc $f$ est formellement lisse. Compte tenu de (17.1.3, (v)), cela démontre la proposition.

\begin{corollary}[17.1.5]
\label{IV.17.1.5}
Supposons $g$ formellement étale; alors, pour que $g\circ f$ soit formellement lisse (resp. formellement non ramifié, resp. formellement étale), il faut et il suffit que $f$ le soit.
\end{corollary}

Cela résulte de (17.1.4) et de (17.1.3, (ii) et (v)).

\begin{proposition}[17.1.6]
\label{IV.17.1.6}
Soit $f:X\to Y$ un morphisme de préschémas.
\begin{enumerate}
  \item[(i)] Soit $(U_\alpha)$ un recouvrement ouvert de $X$ et, pour tout $\alpha$, soit $i_\alpha:U_\alpha\to X$ l'injection canonique. Pour que $f$ soit formellement lisse (resp. formellement non ramifié, resp. formellement étale), il faut et il suffit que chacun des morphismes $f\circ i_\alpha$ le soit.
  \item[(ii)] Soit $(V_\lambda)$ un recouvrement ouvert de $Y$. Pour que $f$ soit formellement lisse (resp. formellement non ramifié, resp. formellement étale), il faut et il suffit que chacune des restrictions $f^{-1}(V_\lambda)\to V_\lambda$ de $f$ le soit.
\end{enumerate}
\end{proposition}

Notons d'abord que (ii) est conséquence de (i): en effet, si $j_\lambda:V_\lambda\to Y$ et $i_\lambda:f^{-1}(V_\lambda)\to X$ sont les injections canoniques, la restriction $f_\lambda:f^{-1}(V_\lambda)\to V_\lambda$ de $f$ est telle que $j_\lambda\circ f_\lambda=f\circ i_\lambda$; si $f$ est formellement lisse (resp. formellement non ramifié), il en est de même de $f\circ i_\lambda$ puisque $i_\lambda$ est formellement étale (17.1.3); mais comme $j_\lambda$ est formellement étale, cela entraîne que $f_\lambda$ est formellement lisse (resp. formellement non ramifié) en vertu de (17.1.5). Inversement, si tous les $f_\lambda$ sont formellement lisses (resp. formellement non ramifiés), il en est de même des $j_\lambda\circ f_\lambda$ (17.1.3), donc aussi de $f$ en vertu de (i).

\oldpage[IV-4]{59}
Si l'on tient compte de ce que les $i_\alpha$ sont formellement étales, tout revient donc à prouver que si les $f\circ i_\alpha$ sont formellement lisses (resp. formellement non ramifiés), il en est de même de $f$.

Soient donc $Y'$ un schéma affine, $Y'_0$ un sous-schéma fermé de $Y'$ défini par un Idéal nilpotent $\sh{J}$, que l'on peut supposer tel que $\sh{J}^2=0$ (17.1.2, (ii)), et enfin soit $g:Y'\to Y$ un morphisme. Supposons donné un $Y$-morphisme $u_0:Y'_0\to X$; désignons par $W_\alpha$ (resp. $W^0_\alpha$) le préschéma induit par $Y'$ (resp. $Y'_0$) sur l'ouvert $u_0^{-1}(U_\alpha)$ (on rappelle que $Y'$ et $Y'_0$ ont \emph{même espace topologique sous-jacent}). Supposons d'abord que les $f\circ i_\alpha$ soient \emph{formellement non ramifiés}, et montrons que, si $u'$ et $u''$ sont deux $Y$-morphismes de $Y'$ dans $X$ dont les restrictions à $Y'_0$ coïncident, alors on a $u'=u''$. En effet, compte tenu de (17.1.2, (iv)), l'hypothèse que les $f\circ i_\alpha$ sont non ramifiés entraîne que pour tout $\alpha$, on a $u'|W_\alpha=u''|W_\alpha$, puisque les restrictions de ces deux $Y$-morphismes à $W^0_\alpha$ coïncident. D'où la conclusion dans ce cas.

Supposons maintenant tous les $f\circ i_\alpha$ \emph{formellement lisses} et prouvons qu'il existe un $Y$-morphisme $u:Y'\to X$ dont $u_0$ est la restriction à $Y'_0$. Or, puisque $Y'$ est un \emph{schéma affine}, on peut appliquer (16.5.17) dont les hypothèses sont satisfaites, et dont la conclusion démontre précisément l'existence de $u$.

On peut donc dire que les notions introduites dans (17.1.1) sont \emph{locales} sur $X$ et sur $Y$, ce qui permet toujours, en vertu de (17.1.2, (i)), de se ramener à l'étude des \emph{algèbres} formellement lisses (resp. formellement non ramifiées, resp. formellement étales).

\subsection{Propriétés différentielles générales}
\EGAReaderTarget{ega4.c04.subsec.17-2}{subsection:IV.17.2}

\begin{proposition}[17.2.1]
\label{IV.17.2.1}
Pour qu'un morphisme $f:X\to Y$ soit formellement non ramifié, il faut et il suffit que $\Omega_f^1=0$ (ce qu'on écrit encore $\Omega_{X/Y}^1=0$ (16.3.1)).
\end{proposition}

Compte tenu de (17.1.6), on est ramené au cas où $Y=\Spec(A)$ et $X=\Spec(B)$ sont affines, et la conclusion résulte alors de (0, 20.7.4) et de l'interprétation de $\Omega_{X/Y}^1$ dans ce cas (16.3.7).

\begin{corollary}[17.2.2]
\label{IV.17.2.2}
Soient $f:X\to Y$, $g:Y\to Z$ deux morphismes. Pour que $f$ soit formellement non ramifié, il faut et il suffit que l'homomorphisme canonique (16.4.19)
\[
  f^*(\Omega_{Y/Z}^1)\longrightarrow\Omega_{X/Z}^1
\]
soit surjectif.
\end{corollary}

C'est une conséquence immédiate de (17.2.1) et de la suite exacte (16.4.19.1).

\begin{proposition}[17.2.3]
\label{IV.17.2.3}
Soit $f:X\to Y$ un morphisme formellement lisse.
\begin{enumerate}
  \item[(i)] Le $\sh{O}_X$-Module $\Omega_{X/Y}^1$ est localement projectif (16.10.1). Si $f$ est localement de type fini, $\Omega_{X/Y}^1$ est localement libre de type fini.
  \item[(ii)] Pour tout morphisme $g:Y\to Z$, la suite (16.4.19) de $\sh{O}_X$-Modules
  \[
  \label{IV.17.2.3.1}
    0\longrightarrow f^*(\Omega_{Y/Z}^1)\longrightarrow\Omega_{X/Z}^1\longrightarrow\Omega_{X/Y}^1\longrightarrow0
    \tag{17.2.3.1}
  \]
  est exacte; en outre, pour tout $x\in X$, il existe un voisinage ouvert $U$ de $x$ tel que les restrictions à $U$ des homomorphismes de (17.2.3.1) forment une suite exacte et scindée.
\end{enumerate}
\end{proposition}

\oldpage[IV-4]{60}
\begin{enumerate}
  \item[(i)] On sait (16.3.9) que si $f$ est localement de type fini, $\Omega_f^1$ est un $\sh{O}_X$-Module de type fini. Pour prouver que, dans tous les cas, il est localement projectif, on peut se borner, en vertu de (17.1.6), au cas où $Y=\Spec(A)$ et $X=\Spec(B)$ sont affines, et cela résulte de l'hypothèse sur $f$ et de (0, 20.4.9 et 0, 19.2.1).
  \item[(ii)] Ici encore, on peut se borner au cas où $X$, $Y$ et $Z$ sont affines (17.1.6) et la conclusion résulte dans ce cas de l'interprétation des Modules figurant dans la suite (17.2.3.1) et de (0, 20.5.7).
\end{enumerate}

\begin{corollary}[17.2.4]
\label{IV.17.2.4}
Si $f:X\to Y$ est un morphisme formellement étale, alors, pour tout morphisme $g:Y\to Z$, l'homomorphisme canonique de $\sh{O}_X$-Modules
\[
  f^*(\Omega_{Y/Z}^1)\longrightarrow\Omega_{X/Z}^1
\]
est bijectif.
\end{corollary}

Cela résulte de l'exactitude de la suite (17.2.3.1) et du fait que l'on a alors $\Omega_{X/Y}^1=0$ (17.2.1).

\begin{proposition}[17.2.5]
\label{IV.17.2.5}
Soient $f:X\to Y$ un morphisme, $X'$ un sous-préschéma de $X$ tel que le morphisme composé $X'\xrightarrow{j}X\to Y$ (où $j$ est l'injection canonique) soit formellement lisse. Alors la suite de $\sh{O}_{X'}$-Modules (16.4.21)
\[
\label{IV.17.2.5.1}
  0\longrightarrow\mathcal{N}_{X'/X}\longrightarrow\Omega_{X/Y}^1\otimes_{\sh{O}_X}\sh{O}_{X'}\longrightarrow\Omega_{X'/Y}^1\longrightarrow0
  \tag{17.2.5.1}
\]
est exacte; en outre, pour tout $x\in X$, il existe un voisinage ouvert $U$ de $x$ tel que les restrictions à $U$ des homomorphismes de (17.2.5.1) forment une suite exacte et scindée.
\end{proposition}

Toujours en vertu de (17.1.6), on peut se borner au cas où $Y=\Spec(A)$ et $X=\Spec(B)$ sont affines, et $X'=\Spec(B/\mathfrak{J})$, où $\mathfrak{J}$ est un idéal de $B$. Alors le faisceau conormal $\mathcal{N}_{X'/X}$ correspond au $B$-module $\mathfrak{J}/\mathfrak{J}^2$ (16.1.3), et la conclusion découle de (0, 20.5.14).

\begin{proposition}[17.2.6]
\label{IV.17.2.6}
Soient $X$, $Y$ deux préschémas, $f:X\to Y$ un morphisme localement de type fini. Les conditions suivantes sont équivalentes:
\begin{enumerate}
  \item[a)] $f$ est un monomorphisme.
  \item[b)] $f$ est radiciel et formellement non ramifié.
  \item[c)] Pour tout $y\in Y$, la fibre $f^{-1}(y)$ est vide ou $k(y)$-isomorphe à $\Spec(k(y))$ (autrement dit, est réduite à un seul point $z$ tel que $k(y)\to\sh{O}_z/\mathfrak{m}_y\sh{O}_z$ soit un isomorphisme).
\end{enumerate}
\end{proposition}

Le fait que a) entraîne c) résulte de (8.11.5.1). Il est clair que c) entraîne que $f$ est radiciel; montrons qu'il résulte aussi de c) que $\Omega_{X/Y}^1=0$, ce qui prouvera que c) entraîne b) (17.2.1). Notons que le $\sh{O}_X$-Module $\Omega_{X/Y}^1$ est quasi-cohérent de type fini (16.3.9). Il résulte donc de (I, 9.1.13.1) que, pour que $(\Omega_{X/Y}^1)_x=0$, il faut et il suffit que si l'on pose $Y_1=\Spec(k(y))$, $X_1=f^{-1}(y)=X\times_Y Y_1$, on ait $(\Omega_{X_1/Y_1}^1)_x=0$; mais comme le morphisme $f_1:X_1\to Y_1$ déduit de $f$ est formellement non ramifié en vertu de l'hypothèse c) (17.1.3), la conclusion résulte de (17.2.1). Prouvons enfin que b) entraîne a); pour cela, considérons le morphisme diagonal $g=\Delta_f:X\to X\times_Y X$; puisque $f$ est radiciel, $g$ est surjectif (I, 8.7.1); d'autre part, $\Omega_{X/Y}^1$ est par définition le faisceau conormal $\mathop{Gr}\nolimits_1(g)$ de l'immersion $g$ (16.3.1), et dire que $f$ est formellement non ramifié signifie donc que

\oldpage[IV-4]{61}
$\mathop{Gr}\nolimits_1(g)=0$ (17.2.1). En outre, $g$ est localement de présentation finie (I, 4.3.1); donc l'hypothèse $\mathop{Gr}\nolimits_1(g)=0$ entraîne que $g$ est une immersion ouverte (16.1.10); étant surjective, cette immersion est un isomorphisme, donc $f$ est un monomorphisme (I, 5.3.8).

\subsection{Morphismes lisses, morphismes non ramifiés, morphismes étales}
\EGAReaderTarget{ega4.c04.subsec.17-3}{subsection:IV.17.3}

\begin{definition}[17.3.1]
\EGAReaderTarget{ega4.c04.def.17-3-1}{IV.17.3.1}
On dit qu'un morphisme $f:X\to Y$ est \emph{lisse} (resp. \emph{non ramifié}, ou \emph{net}\footnote{Les mots « net » et « formellement net » paraissent bien préférables à la terminologie consacrée de « non ramifié » (resp. « formellement non ramifié ») et seront employés à peu près exclusivement à partir du chap. V. Dans ce chapitre, nous avons conservé l'ancienne terminologie afin de ne pas entrer en conflit avec (0, 19.10).}, resp. \emph{étale}) s'il est localement de présentation finie et formellement lisse (resp. formellement non ramifié, resp. formellement étale).
\end{definition}

On dit encore alors que $X$ est \emph{lisse} (resp. \emph{non ramifié} ou \emph{net}, resp. \emph{étale}) sur $Y$.

Nous verrons plus loin (17.5.2) que cette définition d'un morphisme lisse coïncide avec celle déjà donnée dans (6.8.1); jusque-là, c'est la définition de (17.3.1) que nous utiliserons exclusivement.

Il est clair que dire que $f$ est étale signifie qu'il est \emph{à la fois} lisse et non ramifié.

\begin{remarks}[17.3.2]
\EGAReaderTarget{ega4.c04.subsec.17-3-2}{IV.17.3.2}
\begin{enumerate}
  \item[(i)] On notera que la définition (17.3.1) peut s'exprimer uniquement à l'aide du foncteur
  \[
    Y'\longmapsto\Hom_Y(Y',X)
  \]
  considéré dans (17.1.2, (iii)), car dire que $f$ est localement de présentation finie équivaut à dire que le foncteur précédent \emph{commute aux limites projectives de schémas affines} (8.14.2).
  \item[(ii)] Soient $A$ un anneau, $B$ une $A$-algèbre. On dit que $B$ est une $A$-algèbre \emph{lisse} (resp. \emph{non ramifiée}, resp. \emph{étale}) si le morphisme correspondant $\Spec(B)\to\Spec(A)$ est lisse (resp. non ramifié, resp. étale). Il revient au même de dire que $B$ est une $A$-algèbre \emph{de présentation finie} (I, 4.4.6) et formellement lisse (resp. formellement non ramifiée, resp. formellement étale) pour les topologies discrètes.
  \item[(iii)] Il résulte de (17.1.6) et de la définition d'un morphisme localement de présentation finie (I, 4.4.2) que les notions de morphisme lisse, non ramifié et étale sont \emph{locales sur $X$ et sur $Y$}.
\end{enumerate}
\end{remarks}

\begin{proposition}[17.3.3]
\label{IV.17.3.3}
\begin{enumerate}
  \item[(i)] Une immersion ouverte est étale. Pour qu'une immersion soit non ramifiée, il faut et il suffit qu'elle soit localement de présentation finie.
  \item[(ii)] Le composé de deux morphismes lisses (resp. non ramifiés, resp. étales) est lisse (resp. non ramifié, resp. étale).
  \item[(iii)] Si $f:X\to Y$ est un $S$-morphisme lisse (resp. non ramifié, resp. étale), il en est de même de $f_{(S')}:X_{(S')}\to Y_{(S')}$ pour toute extension $S'\to S$ du préschéma de base.
  \item[(iv)] Si $f:X\to X'$ et $g:Y\to Y'$ sont deux $S$-morphismes lisses (resp. non ramifiés, resp. étales), il en est de même de $f\times_S g:X\times_S Y\to X'\times_S Y'$.
\end{enumerate}

\oldpage[IV-4]{62}

\begin{enumerate}
  \item[(v)] Soient $f:X\to Y$, $g:Y\to Z$ deux morphismes; si $g$ est localement de type fini et si $g\circ f$ est non ramifié, alors $f$ est non ramifié.
\end{enumerate}
\end{proposition}

Cela résulte aussitôt de (I, 4.3) et de (17.1.3).

\begin{proposition}[17.3.4]
\label{IV.17.3.4}
Soient $f:X\to Y$, $g:Y\to Z$ deux morphismes, et supposons $g$ non ramifié. Alors, si $g\circ f$ est lisse (resp. non ramifié, resp. étale), il en est de même de $f$.
\end{proposition}

En effet, comme $g$ et $g\circ f$ sont localement de présentation finie, il en est de même de $f$ (I, 4.4.3, (v)); la conclusion résulte donc de (17.1.4) et (17.1.3, (v)).

\begin{corollary}[17.3.5]
\label{IV.17.3.5}
Supposons $g$ étale; alors, pour que $f$ soit lisse (resp. non ramifié, resp. étale), il faut et il suffit que $g\circ f$ le soit.
\end{corollary}

Cela résulte de (17.3.4) et de (17.3.3, (iii)).

\begin{proposition}[17.3.6]
\label{IV.17.3.6}
Soient $g:Y\to S$, $h:X\to S$ deux morphismes localement de présentation finie. Pour qu'un $S$-morphisme $f:X\to Y$ soit non ramifié, il faut et il suffit que l'homomorphisme canonique (16.4.19)
\[
  f^*(\Omega_{Y/S}^1)\longrightarrow\Omega_{X/S}^1
\]
soit surjectif.
\end{proposition}

Comme $f$ est alors localement de présentation finie (I, 4.4.3, (v)), la proposition résulte aussitôt de (17.2.2).

\begin{definition}[17.3.7]
\EGAReaderTarget{ega4.c04.def.17-3-7}{IV.17.3.7}
Soit $f:X\to Y$ un morphisme. On dit que $f$ est \emph{lisse} (resp. \emph{non ramifié}, resp. \emph{étale}) \emph{en un point} $x\in X$ s'il existe un voisinage ouvert $U$ de $x$ dans $X$ tel que la restriction $f|U$ soit un morphisme lisse (resp. non ramifié, resp. étale) de $U$ dans $Y$.
\end{definition}

On dit encore alors que $X$ est lisse (resp. non ramifié, resp. étale) sur $Y$ au point $x$.

Compte tenu de la remarque (17.3.2, (iii)), il revient au même de dire que $f$ est un morphisme lisse (resp. non ramifié, resp. étale) ou qu'il est lisse (resp. non ramifié, resp. étale) en chaque point de $X$.

Il est clair que l'ensemble des points de $X$ où un morphisme $f:X\to Y$ est lisse (resp. non ramifié, resp. étale) est \emph{ouvert} dans $X$.

\begin{proposition}[17.3.8]
\label{IV.17.3.8}
Pour tout préschéma $Y$ et tout $\sh{O}_Y$-Module localement libre $\sh{E}$ de type fini, le préschéma fibré vectoriel $\mathbf{V}(\sh{E})$ (II, 1.7.8) est un $Y$-préschéma lisse.
\end{proposition}

En effet (17.3.2, (iii)), on peut se borner au cas où $Y=\Spec(A)$ est affine et $\mathbf{V}(\sh{E})=\Spec(A[T_1,\ldots,T_r])$; comme $A[T_1,\ldots,T_r]$ est une $A$-algèbre formellement lisse pour les topologies discrètes (0, 19.3.2) et de présentation finie, cela démontre la proposition (17.3.2, (ii)).

\begin{corollary}[17.3.9]
\label{IV.17.3.9}
Sous les hypothèses de (17.3.8), le préschéma fibré projectif $\mathbf{P}(\sh{E})$ (II, 4.1.1) est un $Y$-préschéma lisse.
\end{corollary}

On peut encore se borner au cas où $Y=\Spec(A)$ est affine et $\mathbf{P}(\sh{E})=\mathbf{P}^r_Y$. On sait alors (II, 2.3.14) qu'on a un recouvrement ouvert fini de $\mathbf{P}^r_A$ en prenant les $D_+(T_i)$ ($0\leq i\leq r$), égaux respectivement aux spectres des anneaux $S_{(f)}$, où l'on remplace $S$ par $A[T_0,T_1,\ldots,T_r]$ et $f$ par $T_i$; mais il résulte aussitôt de la définition de $S_{(f)}$ (II, 2.2.1) que cet anneau, dans le cas considéré, est isomorphe à $A[T_0,\ldots,T_{i-1},T_{i+1},\ldots,T_r]$; donc le corollaire résulte de (17.3.8).

\oldpage[IV-4]{63}

\subsection{Caractérisations des morphismes non ramifiés}
\EGAReaderTarget{ega4.c04.subsec.17-4}{subsection:IV.17.4}

\begin{theorem}[17.4.1]
\label{IV.17.4.1}
Soient $f:X\to Y$ un morphisme localement de présentation finie, $x$ un point de $X$. Les propriétés suivantes sont équivalentes:
\begin{enumerate}
  \item[a)] $f$ est non ramifié au point $x$.
  \item[b)] Le morphisme diagonal $\Delta_f:X\to X\times_Y X$ est un isomorphisme local au point $x$.
  \item[b')] Si l'on pose $\Delta_f=(\psi,\theta)$, $Z=X\times_Y X$ et $z=\psi(x)$, l'homomorphisme $\theta_z^\#:\sh{O}_{Z,z}\to\sh{O}_{X,x}$ est bijectif.
  \item[b'')] Pour tout morphisme $g:Y'\to Y$, et tout point $y'\in Y'$ au-dessus de $y=f(x)$, toute $Y'$-section $s'$ de $X'=X\times_Y Y'$ telle que $x'=s'(y')$ soit au-dessus de $x$ est un isomorphisme local au point $y'$.
  \item[c)] On a $(\Omega^1_{X/Y})_x=0$.
  \item[d)] Le $k(y)$-préschéma $f^{-1}(y)$ est non ramifié sur $k(y)$ au point $x$.
  \item[d')] Le point $x$ est isolé dans $X_y=f^{-1}(y)$ (autrement dit (Err$_{\mathrm{III}}$, 20), le morphisme $f$ est quasi-fini au point $x$) et l'anneau $\sh{O}_{X_y,x}$ est un corps, extension séparable de $k(y)$.
  \item[d'')] L'anneau $\sh{O}_{X_y,x}=\sh{O}_{X,x}/\mathfrak{m}_y\sh{O}_{X,x}$ est un corps, extension finie séparable de $k(y)$.
  \item[e)] L'anneau $\sh{O}_{X,x}$ est une $\sh{O}_{Y,y}$-algèbre formellement non ramifiée pour les topologies discrètes.
\end{enumerate}
\end{theorem}

Comme $f$ est localement de type fini, le $\sh{O}_X$-Module $\Omega^1_{X/Y}$ est de type fini (16.3.9), donc il revient au même de dire que $(\Omega^1_{X/Y})_x=0$ ou qu'il existe un voisinage ouvert $U$ de $x$ tel que $\Omega^1_{X/Y}|U=0$. Compte tenu de (17.2.1), cela prouve l'équivalence de a) et c). D'autre part, si l'on pose $A=\sh{O}_{Y,y}$, $B=\sh{O}_{X,x}$, on a $(\Omega^1_{X/Y})_x=\Omega^1_{B/A}$ (16.4.15), et l'équivalence de c) et e) résulte donc de (0, 20.7.4).

Comme d') ne fait intervenir que des propriétés du morphisme $X_y\to\Spec(k(y))$, l'équivalence de a) et d') entraînera \emph{ipso facto} celle de d) et d'). D'autre part d') et d'') sont équivalentes, car il revient au même de dire que $\sh{O}_{X_y,x}$ est une $k(y)$-algèbre \emph{finie} ou que $x$ est un point isolé de $X_y$, puisque $X_y$ est un $k(y)$-préschéma localement de type fini (I, 6.4.4).

Prouvons maintenant l'équivalence de b) et b'). On peut se limiter au cas où $Y=\Spec(R)$ et $X=\Spec(S)$ sont affines et $f$ de présentation finie; alors on a $Z=\Spec(S\otimes_R S)$ et $\Delta_f$ correspond à l'homomorphisme canonique surjectif $S\otimes_R S\to S$, dont on sait que le noyau $\mathfrak{J}$ est un idéal de type fini (0, 20.4.4). Si l'on pose $\sh{J}=\widetilde{\mathfrak{J}}$, le $\sh{O}_Z$-Module $\psi_*(\sh{O}_X)=\sh{O}_Z/\sh{J}$ est donc de présentation finie, et l'hypothèse que l'homomorphisme $\theta_z^\#: \sh{O}_{Z,z}\to(\psi_*(\sh{O}_X))_z$ est bijectif entraîne qu'en remplaçant au besoin $X$ par un voisinage ouvert de $x$, l'homomorphisme $\theta:\sh{O}_Z\to\psi_*(\sh{O}_X)$ est lui-même bijectif ($0_{\mathrm I}$, 5.2.7). Cela montre donc que b') entraîne b); la réciproque est évidente.

D'autre part, l'équivalence de b) et b'') résulte de (I, 5.3.7) sans hypothèse de finitude sur $f$: la donnée d'une $Y'$-section $s':Y'\to X'$ équivaut en effet à celle d'un

\oldpage[IV-4]{64}
$Y$-morphisme $h=g'\circ s':Y'\to X$ (où $g':X'\to X$ est la projection canonique), de sorte que $s'=(1_{Y'},h)_X$, et alors le diagramme
\[
\tag{17.4.1.1}
\label{IV.17.4.1.1}
\xymatrix{
Y' \ar[r]^{s'} \ar[d]_h & X'=Y'\times_Y X \ar[d]^{h\times 1_X}\\
X \ar[r]_{\Delta_f} & X\times_Y X
}
\]
identifie $Y'$ au produit des $(X\times_Y X)$-préschémas $X$ et $X'$. Par suite (I, 4.3.2) si $\Delta_f$ est un isomorphisme local au point $x$, $s'$ est un isomorphisme local au point $y'$ (puisque $x=h(y')$), ce qui prouve que b) implique b''). La réciproque s'obtient en appliquant b'') au cas où l'on prend $Y'=X$, $y'=x$, $g=f$ et $s'=\Delta_f$.

Pour achever la démonstration de (17.4.1), il suffit de prouver les implications
\[
d'')\Rightarrow c)\Rightarrow b)\Rightarrow d'').
\]

$d'')\Rightarrow c)$: Comme $\Omega^1_{X/Y}$ est un $\sh{O}_X$-Module de type fini, il résulte du lemme de Nakayama que la condition c) équivaut à $(\Omega^1_{X/Y})_x/\mathfrak{m}_y(\Omega^1_{X/Y})_x=0$, c'est-à-dire (16.4.5) à $(\Omega^1_{X_y/\Spec(k(y))})_x=0$. On est donc ramené au cas où $Y$ est le spectre d'un corps $k$ et $X$ un $k$-préschéma algébrique. L'hypothèse que $\sh{O}_{X,x}$ est un corps $k'$, extension finie de $k$, entraîne d'abord que $x$ est \emph{fermé} dans $X$ (I, 6.4.2), puis que $x$ est point \emph{maximal} du préschéma noethérien $X$, donc est un point isolé de $X$. Remplaçant $X$ par l'ouvert $\{x\}$ de $X$, on peut donc supposer que $X=\Spec(k')$; mais alors l'hypothèse que $k'$ est extension \emph{finie séparable} de $k$ entraîne $\Omega^1_{k'/k}=0$ (0, 20.6.20), ce qui prouve c).

$c)\Rightarrow b)$: On a vu plus haut que l'on a alors $\Omega^1_{X/Y}|U=0$ pour un voisinage ouvert $U$ de $x$ dans $X$; l'assertion b) résulte alors de la définition de $\Omega^1_{X/Y}$ (16.3.1) et de (16.1.9).

$b)\Rightarrow d'')$: Remplaçant $X$ par un voisinage ouvert de $x$, on peut supposer que $\Delta_f$ est une immersion ouverte; si l'on désigne par $f_y:X_y\to\Spec(k(y))$ le morphisme déduit de $f$ par changement de base, $\Delta_{f_y}$ est alors aussi une immersion ouverte (I, 5.3.4), et comme la condition d'') ne concerne que le préschéma $X_y$, on voit qu'on peut se borner au cas où $Y$ est le spectre d'un corps $k$, $X$ le spectre d'une $k$-algèbre $A$ de type fini; la propriété d'') sera établie si l'on prouve que $A$ est une $k$-algèbre \emph{finie et séparable}, une telle algèbre étant composée directe d'extensions finies séparables de $k$. Si $K$ est une extension algébriquement close de $k$, il revient au même de dire que $A\otimes_k K$ est une $K$-algèbre finie et séparable (4.6.1), donc on voit qu'on peut se borner au cas où $k$ est algébriquement clos. Montrons d'abord que $A$ est une $k$-algèbre \emph{finie}: il suffira de montrer que tout point \emph{fermé} $x$ de $X$ est isolé, car alors l'ensemble de ces points est ouvert dans $X$ et discret, donc fini puisqu'il est quasi-compact ($X$ étant noethérien), ce qui établira notre assertion en vertu de (I, 6.4.4). Or, on a alors $k(x)=k$ puisque $k$ est algébriquement clos (I, 6.4.3), donc il y a une $Y$-section $s$ de $X$ telle que $s(Y)=\{x\}$, et en vertu de (17.4.1.1), $\{x\}$ est l'image réciproque de la diagonale $\Delta_X(X)$ par un morphisme $X\to X\times_Y X$, donc $\{x\}$ est ouvert dans $X$ en vertu de l'hypothèse b). On a ainsi montré que $A$ est une

\oldpage[IV-4]{65}
$k$-algèbre \emph{finie}, composée directe de $k$-algèbres finies locales. Pour exprimer que $\Delta_f$ est une immersion ouverte, on peut donc se borner au cas où $A$ est une $k$-algèbre locale finie, $X=\Spec(A)$ étant donc réduit à un seul point; le corps résiduel de $A$, étant une extension finie de $k$, est nécessairement identique à $k$, et par suite (I, 3.4.9) $X\times_k X$ est réduit à un seul point et $\Delta_f$ est donc nécessairement un \emph{isomorphisme}. Or, puisque $A$ est une $k$-algèbre, l'homomorphisme canonique $A\otimes_k A\to A$ ne peut être bijectif que si $A=k$. C.Q.F.D.

\begin{remark}[17.4.1.2]
\label{IV.17.4.1.2}
Supposons seulement que $f$ soit \emph{localement de type fini}. Alors $\Omega^1_{X/Y}$ est encore un $\sh{O}_X$-Module de type fini (16.3.9), et $\Delta_f$ un morphisme localement de présentation finie (I, 4.3.1). Toute la démonstration de (17.4.1) est donc valable, à condition de remplacer a) par: \emph{la restriction de $f$ à un voisinage convenable de $x$ est un morphisme formellement non ramifié}. On voit en outre que dans ce cas la restriction de $f$ à un voisinage convenable de $x$ est un morphisme \emph{localement quasi-fini}.
\end{remark}

\begin{corollary}[17.4.2]
\label{IV.17.4.2}
Soit $f:X\to Y$ un morphisme localement de présentation finie. Les propriétés suivantes sont équivalentes:
\begin{enumerate}
  \item[a)] $f$ est non ramifié.
  \item[b)] Le morphisme diagonal $\Delta_f:X\to X\times_Y X$ est une immersion ouverte.
  \item[b')] Pour tout morphisme $Y'\to Y$, toute $Y'$-section de $X'=X\times_Y Y'$ est une immersion ouverte.
  \item[c)] On a $\Omega^1_{X/Y}=0$.
  \item[d)] Pour tout $y\in Y$, le $k(y)$-préschéma $f^{-1}(y)$ est non ramifié sur $k(y)$.
  \item[d')] Pour tout $y\in Y$, le $k(y)$-préschéma $f^{-1}(y)$ est isomorphe à un préschéma de la forme $\coprod_{\lambda\in L}\Spec(K_\lambda)$, où, pour tout $\lambda\in L$, $K_\lambda$ est une extension finie et séparable de $k(y)$.
  \item[e)] Pour tout $x\in X$, l'anneau $\sh{O}_{X,x}$ est une $\sh{O}_{Y,f(x)}$-algèbre formellement non ramifiée pour les topologies discrètes.
\end{enumerate}
\end{corollary}

\begin{corollary}[17.4.3]
\label{IV.17.4.3}
Si $f:X\to Y$ est non ramifié, alors $f$ est localement quasi-fini (Err$_{\mathrm{III}}$, 20).
\end{corollary}

\begin{proposition}[17.4.4]
\label{IV.17.4.4}
Soient $Y$ un préschéma localement noethérien, $f:X\to Y$ un morphisme localement de type fini, $x$ un point de $X$, $y=f(x)$. Posons $A=\sh{O}_{Y,y}$, $B=\sh{O}_{X,x}$, qui sont des anneaux locaux noethériens, et soit $k$ le corps résiduel de $A$. Alors les conditions équivalentes a) à e) du théorème (17.4.1) sont aussi équivalentes à chacune des suivantes:
\begin{enumerate}
  \item[f)] $\widehat{B}\otimes_{\widehat{A}}k$ est un corps, extension finie séparable de $k$ (ce qui entraîne que $\widehat{B}$ est une $\widehat{A}$-algèbre finie).
  \item[f')] $\widehat{B}$ est une $\widehat{A}$-algèbre formellement non ramifiée pour les topologies adiques.
\end{enumerate}
Si de plus $k(x)=k(y)$, ou si $k$ est séparablement clos, ces conditions équivalent aussi à:
\begin{enumerate}
  \item[f'')] L'homomorphisme $\widehat{A}\to\widehat{B}$ est surjectif.
\end{enumerate}
\end{proposition}

Notons d'abord, par le même raisonnement que dans (0, 19.3.6), qu'il revient au même de dire que $B$ est une $A$-algèbre formellement non ramifiée pour les topologies préadiques, ou que $\widehat{B}$ est une $\widehat{A}$-algèbre formellement non ramifiée pour les topologies adiques. D'autre part, l'hypothèse que $f$ est localement de type fini entraîne que $\Omega^1_{B/A}$

\oldpage[IV-4]{66}
est un $B$-module de type fini (16.3.9), donc \emph{séparé} pour la topologie $\mathfrak n$-préadique (où $\mathfrak n$ est l'idéal maximal de $B$) ($0_{\mathrm I}$, 7.3.5); il revient par suite au même de dire que $\Omega^1_{B/A}=0$ ou que $\widehat{\Omega}^1_{B/A}=0$; donc (0, 20.7.4), il revient au même de dire que $B$ est une $A$-Algèbre formellement non ramifiée pour les topologies \emph{discrètes}, ou pour les topologies \emph{préadiques}. Ceci prouve l'équivalence des conditions e) et f'). Si $\mathfrak m$ est l'idéal maximal de $A$, on a
\[
k=A/\mathfrak m=\widehat A/\mathfrak m\widehat A,
\qquad
\widehat B\otimes_{\widehat A}k=\widehat B/\mathfrak m\widehat B
=\widehat B\otimes_B(B/\mathfrak m B),
\]
et par suite ($0_{\mathrm I}$, 7.3.5) $\widehat B/\mathfrak m\widehat B$ est le complété de $B/\mathfrak m B=B\otimes_A k$ pour la topologie $\mathfrak n$-préadique; ceci prouve l'équivalence de d'') et de f). Enfin, lorsque $k(x)=k(y)$ ou lorsque $k$ est séparablement clos, la condition f) entraîne que l'homomorphisme $\widehat A/\mathfrak m\widehat A\to\widehat B/\mathfrak m\widehat B$ est bijectif; la condition f) implique d'autre part que $\widehat B$ est une $\widehat A$-algèbre \emph{quasi-finie} ($0_{\mathrm I}$, 7.4.4), donc \emph{finie} puisque $\widehat A$ est complet et $\widehat B$ séparé pour la topologie $\mathfrak m$-préadique, $\mathfrak m\widehat B$ étant un idéal de définition de $\widehat B$ ($0_{\mathrm I}$, 7.4.1). L'homomorphisme $\widehat A\to\widehat B$ est donc surjectif en vertu du lemme de Nakayama. Donc f) entraîne alors f''), et la réciproque est évidente.

\begin{env}[17.4.5]
\EGAReaderTarget{ega4.c04.statement.17-4-5}{IV.17.4.5}
Étant donnés un $S$-préschéma $Y$ et deux $S$-morphismes $f:X\to Y$, $g:X\to Y$, on en déduit canoniquement un $S$-morphisme $(f,g)_S:X\to Y\times_S Y$. Nous appellerons \emph{préschéma des coïncidences} de $f$ et $g$ l'image réciproque par $(f,g)_S$ de la diagonale $\Delta_{Y|S}$; c'est donc un sous-préschéma de $X$, qui est \emph{fermé} lorsque $Y$ est un $S$-schéma (I, 5.4.1).
\end{env}

\begin{proposition}[17.4.6]
\label{IV.17.4.6}
Soient $h:Y\to S$ un morphisme non ramifiée, et soient $f:X\to Y$, $g:X\to Y$ deux $S$-morphismes. Alors le préschéma $C$ des coïncidences de $f$ et $g$ est un sous-préschéma induit sur un ouvert de $X$; si de plus $Y$ est un $S$-schéma (I, 5.4.1), $C$ est un sous-préschéma fermé de $X$.
\end{proposition}

En effet, puisque $\Delta_h:Y\to Y\times_S Y$ est une immersion ouverte (17.4.2), l'image réciproque par $(f,g)_S$ de $\Delta_{Y|S}$ est un sous-préschéma induit sur un ouvert de $X$ (I, 4.4.1). La dernière assertion résulte de (17.4.5).

\begin{corollary}[17.4.7]
\label{IV.17.4.7}
Sous les hypothèses de (17.4.6), soit $x$ un point de $X$ tel que les deux morphismes composés $\Spec(k(x))\to X\xrightarrow{f}Y$ et $\Spec(k(x))\to X\xrightarrow{g}Y$ soient égaux. Alors il existe un voisinage ouvert $U$ de $x$ tel que $f|U=g|U$. Si de plus $Y$ est un $S$-schéma, il existe un voisinage ouvert et fermé $X'$ de $x$ dans $X$ tel que $f|X'=g|X'$. Si enfin on suppose en outre que $X$ est connexe, on a $f=g$.
\end{corollary}

Cela résulte de (17.4.6) et de (I, 5.3.17).

\begin{corollary}[17.4.8]
\label{IV.17.4.8}
Sous les hypothèses de (17.4.6), supposons que le morphisme structural $\varphi=h\circ f=h\circ g$ de $X$ dans $S$ soit fermé. Soit $s$ un point de $S$; soit $X_s$ le $k(s)$-préschéma $\varphi^{-1}(s)$ et supposons que les deux morphismes composés $X_s\to X\xrightarrow{f}Y$ et $X_s\to X\xrightarrow{g}Y$ soient égaux. Alors il existe un voisinage ouvert $V$ de $s$ dans $S$ tel que $f|\varphi^{-1}(V)=g|\varphi^{-1}(V)$. Si de plus $Y$ est un $S$-schéma et si $\varphi$ est ouvert, on peut prendre $V$ ouvert et fermé. Si enfin on suppose de plus $S$ connexe, on a $f=g$.
\end{corollary}

Il résulte de (17.4.7) que le préschéma $C$ des coïncidences de $f$ et $g$ est induit sur un ouvert de $X$ et contient $X_s$. Comme $\varphi$ est fermé, il existe un voisinage ouvert $V$ de $s$ tel que $\varphi^{-1}(V)\subset C$. Si de plus $Y$ est un $S$-schéma, $C$ est fermé, donc $\varphi(X-C)$ est à la

\oldpage[IV-4]{67}
fois ouvert et fermé dans $S$, et son complémentaire $V$ dans $S$ est donc un voisinage ouvert et fermé de $s$ tel que $\varphi^{-1}(V)\subset C$.

\begin{proposition}[17.4.9]
\label{IV.17.4.9}
Soient $Y$ un préschéma connexe, $f:X\to Y$ un morphisme non ramifié et séparé. Alors toute $Y$-section $g$ de $X$ est un isomorphisme de $Y$ sur une composante connexe ouverte de $X$, et l'application $g\mapsto g(Y)$ est une bijection de $\Gamma(X/Y)$ sur l'ensemble des composantes connexes $Z$ de $X$ (nécessairement ouvertes dans $X$) telles que la restriction de $f$ à $Z$ soit un isomorphisme de $Z$ sur $Y$. En particulier, si $g'$ et $g''$ sont deux $Y$-sections de $X$ telles que $g'(y)=g''(y)$ pour un $y\in Y$, on a $g'=g''$.
\end{proposition}

Il résulte en effet de (17.4.1, b'')) qu'une $Y$-section $s$ de $X$ est une immersion ouverte, et comme $X$ est un $Y$-schéma, $s$ est aussi une immersion fermée (I, 5.4.7); il en résulte que $s$ est un isomorphisme de $Y$ sur un sous-préschéma de $X$ induit sur une partie ouverte et fermée de $X$, et comme $s(Y)$ est connexe, c'est nécessairement une composante connexe de $X$. Le reste de la proposition est immédiat.

\begin{remark}[17.4.10]
\label{IV.17.4.10}
Compte tenu de la remarque (17.4.1.2), on voit que, dans les énoncés (17.4.6) à (17.4.9), on peut remplacer partout les mots « non ramifié » par « formellement non ramifié et localement de type fini ».
\end{remark}

\subsection{Caractérisations des morphismes lisses}
\EGAReaderTarget{ega4.c04.subsec.17-5}{subsection:IV.17.5}

\begin{theorem}[17.5.1]
\label{IV.17.5.1}
Soient $f:X\to Y$ un morphisme localement de présentation finie, $x$ un point de $X$, $y=f(x)$. Les conditions suivantes sont équivalentes:
\begin{enumerate}
  \item[a)] $f$ est lisse au point $x$.
  \item[b)] $f$ est plat au point $x$ et le $k(y)$-préschéma $f^{-1}(y)$ est lisse sur $k(y)$ au point $x$.
  \item[b')] $f$ est régulier au point $x$ (6.8.1).
  \item[c)] L'anneau $\sh{O}_{X,x}$ est une $\sh{O}_{Y,y}$-algèbre formellement lisse pour les topologies discrètes.
\end{enumerate}
\end{theorem}

On peut se borner au cas où $Y=\Spec(A)$, $X=\Spec(C)$, où $C=B/\mathfrak{J}$, $B=A[T_1,\ldots,T_n]$ étant une algèbre de polynômes et $\mathfrak{J}$ un idéal de $B$ de type fini. L'équivalence de a) et c) résulte alors de l'équivalence de a) et c) dans (0, 22.6.4). D'autre part, appliquant ce résultat au morphisme localement de type fini $f^{-1}(y)\to\Spec(k(y))$, on voit que l'équivalence de b) et b') résulte de l'équivalence de a) et b) dans (6.8.6). Reste donc à démontrer l'équivalence de a) et b).

Montrons d'abord que a) entraîne b); notons $\mathfrak p$ l'idéal premier $\mathfrak i_x$ dans $C$, $\mathfrak r$ l'idéal premier $\mathfrak i_y$ dans $A$; on a $\mathfrak p=\mathfrak q/\mathfrak J$, où $\mathfrak q$ est un idéal premier de $B$ et $\mathfrak r$ est l'image réciproque de $\mathfrak q$ dans $A$. L'hypothèse a) entraîne d'abord que $f^{-1}(y)$ est lisse sur $k(y)$ au point $x$ par (17.3.3), et il s'agit de montrer en outre que $C_{\mathfrak p}=\sh{O}_{X,x}$ est un $A_{\mathfrak r}$-module \emph{plat}. Comme $C_{\mathfrak p}$ est une $A_{\mathfrak r}$-algèbre formellement lisse et que $B$ est une $A$-algèbre formellement lisse (pour les topologies discrètes), le critère jacobien (0, 22.6.4), joint à (0, 19.1.12), entraîne qu'il existe dans $\mathfrak J$ un système de $r$ polynômes $u_i$ ($1\leq i\leq r$) et $r$ indices $j_i$ ($1\leq i\leq r$) tels que les images des $u_i$ dans $\mathfrak J_{\mathfrak q}/\mathfrak J_{\mathfrak q}^2$ engendrent ce $B_{\mathfrak q}$-module et que l'on ait
\[
\label{IV.17.5.1.1}
\det(\partial u_i/\partial T_{j_k})\notin\mathfrak q.
\tag{17.5.1.1}
\]

\oldpage[IV-4]{68}
Notons maintenant que si $g:\Spec(B)\to\Spec(A)$ est le morphisme structural, la fibre $g^{-1}(y)$ est le spectre de l'anneau régulier $k(y)[T_1,\ldots,T_n]$ (0, 17.3.7), donc l'anneau local noethérien $B_{\mathfrak q}/\mathfrak rB_{\mathfrak q}$ en un point de cette fibre est régulier. Or, la condition (17.5.1.1) entraîne que les images canoniques $v_i$ des $u_i$ dans l'idéal maximal $\mathfrak m$ de $B_{\mathfrak q}/\mathfrak rB_{\mathfrak q}$ sont \emph{linéairement indépendantes} mod. $\mathfrak m^2$: sinon, en effet, il existerait des polynômes $w_i\in B$ ($1\leq i\leq r$) n'appartenant pas tous à $\mathfrak q$ et tels que $\sum_{i=1}^r w_i u_i\in\mathfrak q^2$. En dérivant par rapport aux $T_{j_k}$, on en conclurait que $\sum_{i=1}^r w_i(\partial u_i/\partial T_{j_k})\in\mathfrak q$ pour $1\leq k\leq r$, ce qui contredirait (17.5.1.1) puisque $\mathfrak q$ est premier. On conclut donc de (0, 17.1.7) que $(v_i)$ est une \emph{suite régulière} dans $B_{\mathfrak q}/\mathfrak rB_{\mathfrak q}$. Mais comme le morphisme $g$ est localement de présentation finie et que $B$ est un $A$-module plat, il résulte de (11.3.8) que les images canoniques $u'_i$ des $u_i$ dans $B_{\mathfrak q}$ forment aussi une \emph{suite régulière} et que $B_{\mathfrak q}/(\sum_i u'_iB_{\mathfrak q})$ est un $A_{\mathfrak r}$-module \emph{plat}. Comme les images des $u'_i$ dans $\mathfrak J_{\mathfrak q}/\mathfrak J_{\mathfrak q}^2$ engendrent ce $B_{\mathfrak q}$-module, il résulte du lemme de Nakayama que l'on a $\sum_i u'_iB_{\mathfrak q}=\mathfrak J_{\mathfrak q}$, et $C_{\mathfrak p}=B_{\mathfrak q}/\mathfrak J_{\mathfrak q}$ est donc bien un $A_{\mathfrak r}$-module plat.

Prouvons enfin que b) implique a). Avec les mêmes notations, l'hypothèse que $B_{\mathfrak q}/\mathfrak J_{\mathfrak q}$ est un $A_{\mathfrak r}$-module plat entraîne que l'homomorphisme canonique $\mathfrak J_{\mathfrak q}/\mathfrak r\mathfrak J_{\mathfrak q}\to B_{\mathfrak q}/\mathfrak rB_{\mathfrak q}$ est \emph{injectif} ($0_{\mathrm I}$, 6.1.2), de sorte que $\mathfrak J_{\mathfrak q}/\mathfrak r\mathfrak J_{\mathfrak q}$ est identifié à un idéal de $B_{\mathfrak q}/\mathfrak rB_{\mathfrak q}$. Comme $B_{\mathfrak q}/\mathfrak rB_{\mathfrak q}$ est une $A_{\mathfrak r}$-algèbre formellement lisse pour les topologies discrètes, on peut appliquer à $C_{\mathfrak p}=(B_{\mathfrak q}/\mathfrak rB_{\mathfrak q})/(\mathfrak J_{\mathfrak q}/\mathfrak r\mathfrak J_{\mathfrak q})$, le critère jacobien (0, 22.6.4); joint à (0, 19.1.12), ce dernier prouve, en vertu de l'hypothèse b), l'existence de $r$ polynômes $v_i\in k(y)[T_1,\ldots,T_n]$ tels que leurs images dans $(\mathfrak J_{\mathfrak q}/\mathfrak r\mathfrak J_{\mathfrak q})/(\mathfrak J_{\mathfrak q}/\mathfrak r\mathfrak J_{\mathfrak q})^2$ engendrent ce $(B_{\mathfrak q}/\mathfrak rB_{\mathfrak q})$-module et que l'on ait
\[
\label{IV.17.5.1.2}
\det(\partial v_i/\partial T_{j_k})\notin\mathfrak qB_{\mathfrak q}/\mathfrak rB_{\mathfrak q}.
\tag{17.5.1.2}
\]
Si, pour tout $i$, on désigne alors par $u_i$ un élément de $\mathfrak J$ dont $v_i$ est l'image canonique, il résulte de (17.5.1.2) que les $u_i$ vérifient la condition (17.5.1.1); d'autre part, en vertu du lemme de Nakayama, les images $u'_i$ des $u_i$ dans $\mathfrak J_{\mathfrak q}$ engendrent ce $B_{\mathfrak q}$-module. Le critère jacobien (0, 22.6.4) joint à (0, 19.1.12), prouve alors que $C_{\mathfrak p}=B_{\mathfrak q}/\mathfrak J_{\mathfrak q}$ est une $A_{\mathfrak r}$-algèbre formellement lisse pour les topologies discrètes. C.Q.F.D.

\begin{corollary}[17.5.2]
\label{IV.17.5.2}
Soit $f:X\to Y$ un morphisme localement de présentation finie. Pour que $f$ soit lisse (au sens de (17.3.1)), il faut et il suffit que $f$ soit régulier (6.8.1), autrement dit que $f$ soit plat et que pour tout $y\in Y$, $f^{-1}(y)$ soit un $k(y)$-préschéma géométriquement régulier (6.7.6).
\end{corollary}

On a ainsi établi l'équivalence des deux définitions de « morphisme lisse » données dans (6.8.1) et (17.3.1).

\begin{proposition}[17.5.3]
\label{IV.17.5.3}
Soient $Y$ un préschéma localement noethérien, $f:X\to Y$ un morphisme localement de type fini, $x$ un point de $X$, $y=f(x)$. Posons $A=\sh{O}_{Y,y}$, $B=\sh{O}_{X,x}$, qui sont des anneaux locaux noethériens. Alors les conditions équivalentes a) à c) de (17.5.1) sont aussi équivalentes à chacune des suivantes:
\begin{enumerate}
  \item[d)] $B$ est une $A$-algèbre formellement lisse pour les topologies préadiques.
  \item[d')] $\widehat B$ est une $\widehat A$-algèbre formellement lisse pour les topologies adiques.
\end{enumerate}
Si de plus $k(x)=k(y)$, ces conditions équivalent aussi à:
\begin{enumerate}
  \item[d'')] $\widehat B$ est une $\widehat A$-algèbre isomorphe à une algèbre de séries formelles $\widehat A[[T_1,\ldots,T_n]]$.
\end{enumerate}
\end{proposition}

\oldpage[IV-4]{69}
L'équivalence de la condition c) de (17.5.1) et de d) résulte de l'équivalence de a) et d) dans le critère jacobien (0, 22.6.4), et l'équivalence de d) et d') résulte de (0, 19.3.6). D'autre part, d'') implique d') sans hypothèse sur les corps résiduels (0, 19.3.4). Enfin, si $\mathfrak m$ désigne l'idéal maximal de $\widehat A$, l'hypothèse d') entraîne que $\widehat B/\mathfrak m\widehat B$ est une $k(y)$-algèbre locale noethérienne complète, formellement lisse pour sa topologie adique (0, 19.3.5); l'hypothèse $k(y)=k(x)$ entraîne donc que $\widehat B/\mathfrak m\widehat B$ est $k(y)$-isomorphe à une algèbre de séries formelles $k(y)[[T_1,\ldots,T_n]]$ (0, 19.6.4). Comme d'autre part, $\widehat A[[T_1,\ldots,T_n]]$ est un $\widehat A$-module plat et une $\widehat A$-algèbre locale noethérienne complète, on conclut de (0, 19.7.1.5) que cette algèbre est isomorphe à $\widehat B$. Donc d') entraîne d'') sous l'hypothèse additionnelle $k(x)=k(y)$.

\begin{remark}[17.5.4]
\label{IV.17.5.4}
Supposons que $Y$ soit un préschéma localement noethérien, et $f:X\to Y$ un morphisme localement de type fini. Le critère (17.5.3, d)), joint à (0, 22.1.4) montre que pour démontrer que $f$ est lisse, on peut appliquer la définition (17.1.1), en se restreignant au cas où le schéma affine $Y'$ est le spectre d'un anneau local artinien.
\end{remark}

\begin{proposition}[17.5.5]
\label{IV.17.5.5}
Soient $Y$ un préschéma localement noethérien, $f:X\to Y$ un morphisme localement de type fini, $x$ un point de $X$, $y=f(x)$. Supposons que $Y$ soit réduit au point $y$. Alors, pour que $f$ soit lisse au point $x$, il faut et il suffit que $f$ soit universellement ouvert dans un voisinage de $x$ dans $f^{-1}(y)$ et que $f^{-1}(y)$ soit un $k(y)$-préschéma géométriquement régulier au point $x$.
\end{proposition}

Compte tenu de (17.5.1), tout revient à voir que, si $f^{-1}(y)$ est un $k(y)$-préschéma géométriquement régulier au point $x$, il est équivalent de dire que $f$ est plat au point $x$ ou universellement ouvert dans un voisinage de $x$ dans $f^{-1}(y)$. Or, si $f$ est plat au point $x$, il l'est dans un voisinage de $x$ dans $X$ (11.1.1) et par suite est universellement ouvert dans ce voisinage (2.4.6). Réciproquement, l'hypothèse que $f$ est universellement ouvert dans un voisinage de $x$ dans $f^{-1}(y)$ et que $f^{-1}(y)$ est un $k(y)$-préschéma géométriquement régulier au point $x$ entraîne que $f$ est plat au point $x$, puisque $\sh{O}_{Y,y}$ est réduit (15.2.2).

\begin{corollary}[17.5.6]
\label{IV.17.5.6}
Soient $Y$ un préschéma localement noethérien, $f:X\to Y$ un morphisme localement de type fini, $x$ un point de $X$, $y=f(x)$. On suppose que $Y$ soit réduit et géométriquement unibranche (6.15.1) au point $y$. Alors, pour que $f$ soit lisse au point $x$, il faut et il suffit que $f$ soit équidimensionnel au point $x$ et que $f^{-1}(y)$ soit un $k(y)$-préschéma géométriquement régulier au point $x$.
\end{corollary}

En remarquant que l'ensemble des points où $f$ est équidimensionnel est ouvert (13.3.2), on voit que le corollaire résulte de (17.5.5) et du critère de Chevalley (14.4.4).

Le fait que $f$ soit un morphisme lisse en un point entraîne en particulier que $f$ vérifie

\oldpage[IV-4]{70}
en ce point toutes les propriétés définies dans (6.8.1). On a donc les propriétés suivantes, que nous rappelons pour la commodité des références:

\begin{proposition}[17.5.7]
\label{IV.17.5.7}
Soit $f:X\to Y$ un morphisme localement de présentation finie, lisse en un point $x\in X$; posons $y=f(x)$. Alors, pour que l'anneau $\sh{O}_{X,x}$ soit réduit (resp. intégralement clos, resp. géométriquement unibranche), il faut et il suffit que $\sh{O}_{Y,y}$ le soit.
\end{proposition}

Cela a en effet été prouvé dans (11.3.13) et (11.3.14) complété par Err$_{\mathrm{IV}}$, 30.

\begin{proposition}[17.5.8]
\label{IV.17.5.8}
Soient $Y$ un préschéma localement noethérien, $f:X\to Y$ un morphisme localement de type fini, lisse en un point $x\in X$; posons $y=f(x)$. Alors:
\begin{enumerate}
  \item[(i)] On a $\dim(\sh{O}_{X,x})=\dim(\sh{O}_{Y,y})+\dim(\sh{O}_{X,x}\otimes_{\sh{O}_{Y,y}}k(y))$.
  \item[(ii)] On a $\operatorname{coprof}(\sh{O}_{X,x})=\operatorname{coprof}(\sh{O}_{Y,y})$.
  \item[(iii)] Pour que l'anneau $\sh{O}_{X,x}$ possède la propriété $(S_n)$ (5.7.2) (resp. $(R_n)$ (5.8.2)), il faut et il suffit que l'anneau $\sh{O}_{Y,y}$ la possède. En particulier, pour que $\sh{O}_{X,x}$ soit régulier, il faut et il suffit que $\sh{O}_{Y,y}$ le soit.
\end{enumerate}
\end{proposition}

Ce sont des cas particuliers de (6.1.2), (6.3.2), (6.4.1) et (6.5.3).

\subsection{Caractérisations des morphismes étales}
\EGAReaderTarget{ega4.c04.subsec.17-6}{subsection:IV.17.6}

\begin{theorem}[17.6.1]
\label{IV.17.6.1}
Soient $f:X\to Y$ un morphisme localement de présentation finie, $x$ un point de $X$, $y=f(x)$. Les conditions suivantes sont équivalentes:
\begin{enumerate}
  \item[a)] $f$ est étale au point $x$.
  \item[a')] $f$ est lisse au point $x$ et non ramifié au point $x$.
  \item[b)] $f$ est lisse au point $x$ et quasi-fini au point $x$ (Err$_{\mathrm{III}}$, 20).
  \item[c)] $f$ est plat au point $x$ et non ramifié au point $x$.
  \item[c')] $f$ est plat au point $x$ et l'anneau $\sh{O}_{X,x}/\mathfrak m_y\sh{O}_{X,x}$ est un corps, extension finie séparable de $k(y)$.
  \item[d)] L'anneau $\sh{O}_{X,x}$ est une $\sh{O}_{Y,y}$-algèbre formellement étale pour les topologies discrètes.
\end{enumerate}
\end{theorem}

L'équivalence de a) et a') résulte aussitôt des définitions; celle de a) et d) résulte de l'équivalence de a) et e) dans (17.4.1) et de l'équivalence de a) et d) dans (17.5.1). L'équivalence de c) et c') résulte de l'équivalence de a) et d') dans (17.4.1). Le fait que a') implique c') découle de (17.5.1); inversement, si c') est vérifiée, $f$ est régulier (donc lisse par (17.5.1)) au point $x$, car si $K$ est une extension finie séparable d'un corps $k$, alors, pour toute extension $k'$ de $k$, $\Spec(K\otimes_k k')$ est régulier, étant somme d'un nombre fini de spectres de corps. Le fait que a') implique b) résulte de (17.4.1, d')) et de (17.5.1, b)). Il reste donc à voir que b) entraîne c), et comme on sait déjà que $f$ est plat au point $x$, par (17.5.1), il suffit de montrer que $f^{-1}(y)$ est un $k(y)$-préschéma non ramifié sur $k(y)$; autrement dit, on est ramené à prouver que b) entraîne c) lorsque $Y=\Spec(k)$ est le spectre d'un corps $k$. Comme la question est locale sur $X$, on peut se borner au cas où $X=\Spec(A)$, où $A$ est une $k$-algèbre finie et locale ($0_{\mathrm I}$, 7.4.1). En vertu de l'hypothèse b), $A$ est une $k$-algèbre formellement lisse pour les topologies discrètes, qui coïncident ici avec les topologies préadiques; donc (0, 19.6.5), $A$ est un anneau local régulier, donc

\oldpage[IV-4]{71}
un corps puisqu'il est artinien, et il résulte alors de (0, 19.6.5.1) que $A$ doit être une extension finie et séparable de $k$, ce qui achève la démonstration (17.4.1).

\begin{corollary}[17.6.2]
\label{IV.17.6.2}
Soit $f:X\to Y$ un morphisme localement de présentation finie. Les conditions suivantes sont équivalentes:
\begin{enumerate}
  \item[a)] $f$ est étale.
  \item[a')] $f$ est lisse et non ramifié.
  \item[b)] $f$ est lisse et localement quasi-fini (Err$_{\mathrm{III}}$, 20).
  \item[c)] $f$ est plat et non ramifié.
  \item[c')] $f$ est plat, et toute fibre $f^{-1}(y)$ est un $k(y)$-schéma somme de spectres de corps, extensions finies et séparables de $k(y)$.
  \item[c'')] $f$ est plat, et pour tout $y\in Y$ et toute extension algébriquement close $k'$ de $k(y)$, la « fibre géométrique » $f^{-1}(y)\otimes_{k(y)}k'$ est somme de spectres de corps isomorphes à $k'$.
\end{enumerate}
\end{corollary}

Le seul point qui reste à démontrer est l'équivalence de c') et c''). Il est clair que c') entraîne c'') par changement de base (17.3.3). D'autre part, comme le morphisme projection $f^{-1}(y)\otimes_{k(y)}k'\to f^{-1}(y)$ est ouvert (2.4.10), l'hypothèse c'') entraîne que l'espace $f^{-1}(y)$ est discret, donc, pour tout point $x\in X_y=f^{-1}(y)$, l'anneau local $\sh{O}_{X_y,x}=A$ est une $k(y)$-algèbre finie, donc un anneau local artinien; en outre il résulte de c'') que $\Spec(A\otimes_{k(y)}k')$ est somme de spectres de corps isomorphes à $k'$, ce qui n'est possible que si $A$ est un corps, extension finie séparable de $k(y)$ (4.6.1).

\begin{proposition}[17.6.3]
\label{IV.17.6.3}
Soient $Y$ un préschéma localement noethérien, $f:X\to Y$ un morphisme localement de type fini, $x$ un point de $X$, $y=f(x)$. Posons $A=\sh{O}_{Y,y}$, $B=\sh{O}_{X,x}$, qui sont des anneaux locaux noethériens, et soit $k$ le corps résiduel de $A$. Alors les conditions équivalentes a) à d) de (17.6.1) sont aussi équivalentes à chacune des suivantes:
\begin{enumerate}
  \item[e)] $\widehat B$ est une $\widehat A$-algèbre formellement étale pour les topologies adiques.
  \item[e')] $\widehat B$ est un $\widehat A$-module libre et $\widehat B\otimes_{\widehat A}k$ est un corps, extension finie et séparable de $k$ (ce qui entraîne que $\widehat B$ est une $\widehat A$-algèbre finie).
\end{enumerate}
Si de plus $k(x)=k(y)$, ou si $k$ est séparablement clos, ces conditions équivalent aussi à:
\begin{enumerate}
  \item[e'')] L'homomorphisme canonique $\widehat A\to\widehat B$ est bijectif.
\end{enumerate}
\end{proposition}

L'équivalence de e) et de chacune des conditions de (17.6.1) résulte aussitôt de (17.4.4, f')) et (17.5.3, d')). Le fait que e) entraîne e') résulte de (17.4.4, f)) et de (0, 19.7.1), compte tenu de ce que $\widehat B$ est alors une $\widehat A$-algèbre finie (17.4.4) et qu'il revient donc au même de dire que $\widehat B$ est un $\widehat A$-module plat ou un $\widehat A$-module libre ($0_{\mathrm{III}}$, 10.1.3). Inversement, le fait que e') entraîne e) résulte de (17.4.4) et de (0, 19.7.1). Enfin, e') entraîne que l'homomorphisme $\widehat A\to\widehat B$ est injectif, et si $k(x)=k(y)$ ou si $k$ est séparablement clos, cet homomorphisme est surjectif par (17.4.4). La réciproque est immédiate.

\begin{proposition}[17.6.4]
\label{IV.17.6.4}
Sous les hypothèses de (17.6.3), si $f$ est étale au point $x$, on a $\dim(\sh{O}_{X,x})=\dim(\sh{O}_{Y,y})$.
\end{proposition}

C'est un cas particulier de (17.5.8, (i)) puisque $x$ est isolé dans sa fibre $f^{-1}(y)$.

\oldpage[IV-4]{72}

\subsection{Propriétés de descente, de passage à la limite et de constructibilité}
\EGAReaderTarget{ega4.c04.subsec.17-7}{subsection:IV.17.7}
\label{IV.17.7}

\begin{proposition}[17.7.1]
\label{IV.17.7.1}
Soient $f:X\to Y$ un morphisme localement de présentation finie, $g:Y'\to Y$ un morphisme, $X'=X\times_Y Y'$, $f'=f_{(Y')}:X'\to Y'$ et $g':X'\to X$ les projections canoniques. Soit $x'$ un point de $X'$ et posons $x=g'(x')$, $y'=f'(x')$.
\begin{enumerate}
  \item[(i)] Si $f'$ est non ramifié au point $x'$, alors $f$ est non ramifié au point $x$.
  \item[(ii)] Supposons de plus que $g$ soit plat au point $y'$. Alors, si $f'$ est lisse (resp. étale) au point $x'$, $f$ est lisse (resp. étale) au point $x$.
\end{enumerate}
\end{proposition}

Posons $y=f(x)=g(y')$, de sorte que l'on a $f'^{-1}(y')=f^{-1}(y)\otimes_{k(y)}k(y')$; notons que $f'$ est localement de présentation finie.

(i) Comme la propriété pour un morphisme (localement de présentation finie) d'être non ramifié en un point ne fait intervenir que la fibre du morphisme en ce point (17.4.1, $d$)), on peut se borner au cas où $Y$ et $Y'$ sont des spectres de corps. Mais alors il revient au même de dire que $f$ (resp. $f'$) est non ramifié en $x$ (resp. $x'$) ou qu'il est étale en ce point (17.6.1, $c$)), donc (i) est une conséquence de (ii).

(ii) Comme $f'$ est plat au point $x'$ (17.5.1), l'hypothèse que $g$ est plat au point $y'$ entraîne que $f$ est plat au point $x$, car la projection $g':X'\to X$ est un morphisme plat au point $x'$, et $f\circ g'=g\circ f'$ est un morphisme plat au point $X'$, d'où la conclusion (2.2.11, (iv)). Le fait que $f$ soit lisse (resp. étale) au point $x$ ne fait plus alors intervenir que la fibre $f^{-1}(y)$ (17.5.1 et 17.6.1), et on est donc encore ramené au cas où $Y=\Spec(k)$ et $Y'=\Spec(k')$ sont des spectres de corps. Dire que $f'$ est lisse au point $x'$ signifie alors (17.5.1) que $X'$ est un $k'$-préschéma géométriquement régulier au point $x'$, et cela entraîne (6.7.8) que $X$ est un $k$-préschéma géométriquement régulier au point $x$, donc que $f$ est lisse au point $x$. Supposons de plus que $f'$ soit étale au point $x'$, de sorte que $x'$ est isolé dans $f'^{-1}(y')$ (17.6.1); comme la projection $f'^{-1}(y')\to f^{-1}(y)$ est un morphisme ouvert (2.4.10), $x$ est isolé dans $f^{-1}(y)$; comme on sait déjà que $f$ est lisse au point $x$, il est étale en ce point (17.6.1).

\begin{corollary}[17.7.2]
\label{IV.17.7.2}
\begin{enumerate}
  \item[(i)] Avec les notations de (17.7.1), soit $U$ (resp. $U'$) l'ensemble des points de $X$ (resp. $X'$) où $f$ (resp. $f'$) est non ramifié; alors on a $U'=g'^{-1}(U)$.
  \item[(ii)] Supposons de plus que $g$ soit plat, et soit $V$ (resp. $V'$) l'ensemble des points de $X$ où $f$ (resp. $f'$) est lisse; alors $V'=g'^{-1}(V)$.
\end{enumerate}
\end{corollary}

\begin{corollary}[17.7.3]
\label{IV.17.7.3}
\begin{enumerate}
  \item[(i)] Supposons $g$ surjectif; alors, pour que $f$ soit non ramifié, il faut et il suffit que $f'$ le soit.
  \item[(ii)] Soient $f:X\to Y$ un $S$-morphisme, $g:S'\to S$ un morphisme fidèlement plat. Supposons que $f$ soit localement de présentation finie, ou que $g$ soit quasi-compact. Alors, pour que $f$ soit lisse (resp. non ramifié, resp. étale), il faut et il suffit que $f'$ le soit.
\end{enumerate}
\end{corollary}

Dans (ii), le cas où $f$ est localement de présentation finie résulte de (17.7.1). Si $g$ est quasi-compact et $f'$ lisse (resp. non ramifié, resp. étale), $f'$ est localement de présentation finie par (2.7.1, (iv)) et on est ramené au premier cas.

\begin{proposition}[17.7.4]
\label{IV.17.7.4}
Soient $f:X\to Y$ un morphisme, $g:Y'\to Y$ un morphisme plat et localement de présentation finie, $X'=X\times_Y Y'$, $f'=f_{(Y')}:X'\to Y'$ et $g':X'\to X$ les

\oldpage[IV-4]{73}

projections canoniques. Soit $V$ (resp. $V'$) l'ensemble des $x\in X$ (resp. $x'\in X'$) où $f$ possède l'une des propriétés suivantes (resp. où $f'$ possède la même propriété): être:
\begin{enumerate}
  \item[(i)] localement de type fini;
  \item[(ii)] localement de présentation finie;
  \item[(iii)] plat;
  \item[(iv)] non ramifié;
  \item[(v)] lisse;
  \item[(vi)] étale.
\end{enumerate}
Alors on a $V'=g'^{-1}(V)$ (autrement dit, pour que $f'$ ait la propriété considérée en un point $x'$, il faut et il suffit que $f$ ait cette propriété au point $x$).
\end{proposition}

La propriété (iii) n'est mise que pour mémoire, et ne nécessite pas l'hypothèse que $g$ soit localement de présentation finie ((2.2.11, (iv)), compte tenu du fait que la projection $g':X'\to X$ est un morphisme plat). En vertu de (17.7.2), les assertions relatives aux propriétés (iv), (v) et (vi) sont des conséquences de l'assertion relative à (ii). Il suffit donc de considérer les cas (i) et (ii). Il est clair que $V'\supset g'^{-1}(V)$; reste donc à montrer que $V'\subset g'^{-1}(V)$; notons que les ensembles $V$ et $V'$ sont ouverts, et $W=g'(V')$ est aussi ouvert dans $X$ en vertu de (2.4.6). Il s'agit donc de prouver que le morphisme $f|W:W\to Y$ est localement de type fini (resp. localement de présentation finie); comme par hypothèse le composé
\[
V'\xrightarrow{g''}W\xrightarrow{f|W}Y
\]
(où $g''$ est la restriction de $g'$), égal à $V'\xrightarrow{f'}Y'\xrightarrow{g}Y$, est localement de type fini (resp. localement de présentation finie) et que $g''$ est surjectif (donc fidèlement plat), on est ramené à prouver le lemme suivant, qui améliore (11.3.16):

\begin{lemma}[17.7.5]
\label{IV.17.7.5}
Soient $f:X\to Y$ un morphisme fidèlement plat et localement de présentation finie, $g:Y\to Z$ un morphisme tel que $g\circ f:X\to Z$ ait l'une des propriétés suivantes: être:
\begin{enumerate}
  \item[(i)] localement de type fini;
  \item[(ii)] localement de présentation finie;
  \item[(iii)] de type fini.
\end{enumerate}
Alors $g$ a la même propriété. Si en outre $f$ est quasi-compact ou $g$ quasi-séparé, la même conclusion est valable pour la propriété:
\begin{enumerate}
  \item[(iv)] être de présentation finie.
\end{enumerate}
\end{lemma}

Dans les cas (i) et (ii), il s'agit de voir que pour tout $y\in Y$, il y a un voisinage ouvert affine $V$ de $y$ dans $Y$ et un voisinage ouvert affine $W$ de $z=g(y)$ dans $Z$ contenant $g(V)$, tels que le morphisme $V\to W$, restriction de $g$, soit de type fini (resp. de présentation finie). Or, il existe par hypothèse un $x\in X$ tel que $f(x)=y$, un voisinage ouvert affine $U$ de $x$ dans $X$, un voisinage ouvert affine $V'$ de $y$ dans $Y$ contenant $f(U)$ et un voisinage ouvert affine $W$ de $z$ dans $Z$ contenant $g(V')$ tels que le morphisme $f_1:U\to V'$, restriction de $f$ soit plat et de présentation finie, et le morphisme $g_1:V'\to W$ restriction de $g$, tel que $g_1\circ f_1$ soit de type fini (resp. de présentation finie). Alors $f(U)$ est ouvert dans $Y$ (2.4.6), et si $V\subset f(U)$ est un voisinage affine de $y$, le morphisme $f_2:f_1^{-1}(V)\to V$,

\oldpage[IV-4]{74}

restriction de $f_1$, est encore de présentation finie et est en outre fidèlement plat; de plus, si $g_2=g_1|V$, $g_2\circ f_2$ est de type fini (resp. de présentation finie), l'ouvert $f_1^{-1}(V)$ étant quasi-compact (resp. quasi-compact et quasi-séparé). On est alors ramené aux hypothèses de (11.3.16), d'où l'on conclut que $g_2$ est un morphisme de type fini (resp. de présentation finie).

Dans le cas (iii) la question est locale sur $Z$, donc on peut supposer $Z$ affine, et il en résulte alors que $X$ est quasi-compact, donc aussi $Y=f(X)$. Le cas (iii) est donc conséquence de (i).

Dans le cas (iv) (avec les hypothèses supplémentaires sur $f$ ou $g$), on peut aussi supposer $Z$ affine, donc $X$ et $Y$ quasi-compacts; on sait en outre déjà que $g$ est localement de présentation finie et quasi-compact (1.1.3), donc tout revient à voir que $g$ est quasi-séparé, et il suffit donc de montrer que cette propriété est vraie lorsque l'on suppose $f$ quasi-compact. Or, puisque $g\circ f$ est quasi-séparé, il en est de même de $f$ (1.2.2, (v)); comme $f$ est quasi-compact et localement de présentation finie, il est de présentation finie (1.6.1); il suffit alors de répéter le raisonnement du premier alinéa de la démonstration de (11.3.16).

\begin{remark}[17.7.6]
\label{IV.17.7.6}
Dans l'assertion relative à la propriété (iv), on ne peut supprimer l'hypothèse que $f$ est quasi-compact; sans quoi, cela impliquerait que tout morphisme $g:Y\to Z$ quasi-compact et localement de présentation finie serait de présentation finie, conclusion que l'on sait être erronée (1.6.4). En effet, on peut se borner au cas où $Z$ est affine, donc $Y$ quasi-compact; il y a par suite un recouvrement fini $(U_i)$ de $Y$ par des ouverts affines tels que les restrictions $g|U_i$ soient de présentation finie; il suffirait alors de prendre pour $X$ le préschéma \emph{somme} des $U_i$, pour $f:X\to Y$ le morphisme canonique, qui est évidemment fidèlement plat et localement de présentation finie (1.4.3); $g\circ f$ serait de présentation finie en vertu du choix des $U_i$ et de (1.6.5), d'où notre assertion.
\end{remark}

\begin{proposition}[17.7.7]
\label{IV.17.7.7}
Soient $f:Y\to S$ et $h:X\to S$ deux morphismes localement de présentation finie, $g:X\to Y$ un $S$-morphisme, $x$ un point de $X$, $y=g(x)$. Supposons que $g$ soit plat au point $x$. Alors, si $h$ est lisse (resp. non ramifié, resp. étale) au point $x$, $f$ est lisse (resp. non ramifié, resp. étale) au point $y$.
\end{proposition}

Posons $s=h(x)=f(y)$. Dire que $h$ est non ramifié au point $x$ (resp. que $f$ est non ramifié au point $y$) équivaut à dire que $h^{-1}(s)$ est étale sur $k(s)$ au point $x$ (resp. que $f^{-1}(s)$ est étale sur $k(s)$ au point $y$) (17.4.1 et 17.6.1); comme le morphisme $g_s:h^{-1}(s)\to f^{-1}(s)$ déduit de $g$ est plat au point $x$, on voit que l'on peut se borner à prouver la proposition lorsque $h$ est lisse ou étale au point $x$. En outre, comme $h$ est alors plat au point $x$ (17.5.1), $f$ est plat au point $y$, comme il résulte de (2.2.11, (iv)). Il revient donc au même (17.5.1) de dire que $f$ est lisse (resp. étale) au point $y$, ou que $f^{-1}(s)$ est lisse (resp. étale) sur $k(s)$ au point $x$. On est ainsi ramené au cas où $S=\Spec(k)$ est le spectre d'un corps.

(i) \emph{Cas des morphismes lisses}. Comme $g$ est un morphisme localement de présentation finie (1.4.3, (v)), il y a un voisinage ouvert $U$ de $x$ dans $X$ dans lequel $g$ est plat (11.3.1) et $h$ lisse. En outre $g(U)$ est un voisinage ouvert de $y$ dans $Y$ (2.4.6); remplaçant $X$ par $U$ et $Y$ par $g(U)$, on peut donc supposer que $g$ est \emph{fidèlement plat} et que $h$ est lisse, et on est ramené à prouver que $f$ est alors lisse. Si $k'$ est une extension algébriquement

\oldpage[IV-4]{75}

close de $k$, $X\otimes_k k'$ est alors lisse sur $k'$, et en vertu de (17.7.3, (ii)), il suffit de prouver que $Y\otimes_k k'$ est lisse sur $k'$ (puisque $\Spec(k')\to\Spec(k)$ est fidèlement plat); on peut donc se borner au cas où $k$ est \emph{algébriquement clos}. Comme l'ensemble des points de $Y$ rationnels sur $k$ est alors très dense dans $Y$ (10.4.8) et comme l'ensemble des points de $Y$ où $Y$ est lisse sur $k$ est ouvert, on voit qu'il suffit de prouver que $Y$ est lisse sur $k$ en tout point rationnel sur $k$. Mais en un tel point $y$, dire que $Y$ est lisse sur $k$ en ce point équivaut à dire que $Y$ est \emph{régulier} en $y$ (17.5.1 et 6.7.8). Or on a $y=g(x)$ pour un $x\in X$ et par hypothèse (17.5.1) $X$ est régulier au point $x$; comme $X$ et $Y$ sont alors localement noethériens et que $g$ est plat, $Y$ est bien régulier au point $y$ (6.5.1, (i)).

(ii) \emph{Cas des morphismes étales}. Par (i) on sait déjà que $f$ est lisse au point $y$; en vertu de (17.6.1), il suffit donc de montrer que $f$ est quasi-fini au point $y$, ou encore que $\sh{O}_{Y,y}$ est une $k$-algèbre finie. Or, comme $\sh{O}_{X,x}$ est un $\sh{O}_{Y,y}$-module fidèlement plat ($0_{\mathrm I}$, 6.6.2), $\sh{O}_{Y,y}$ s'identifie à un sous-$k$-module de $\sh{O}_{X,x}$ et comme $\sh{O}_{X,x}$ est par hypothèse une $k$-algèbre finie, il en est de même de $\sh{O}_{Y,y}$.

\begin{proposition}[17.7.8]
\label{IV.17.7.8}
Les notations étant celles de (8.8.1), on suppose $X_\alpha$ et $Y_\alpha$ localement de présentation finie sur $S_\alpha$. Soient $f_\alpha:X_\alpha\to Y_\alpha$ un $S_\alpha$-morphisme, $f:X\to Y$ le $S$-morphisme correspondant.
\begin{enumerate}
  \item[(i)] Soient $x$ un point de $X$, $x_\lambda$ sa projection canonique dans $X_\lambda$. Pour que $f$ soit lisse (resp. non ramifié, resp. étale) au point $x$, il faut et il suffit qu'il existe $\lambda\geq\alpha$ tel que $f_\lambda$ soit lisse (resp. non ramifié, resp. étale) au point $x_\lambda$.
  \item[(ii)] Supposons de plus $X_\alpha$ quasi-compact. Pour que $f$ soit lisse (resp. non ramifié, resp. étale), il faut et il suffit qu'il existe $\lambda\geq\alpha$ tel que $f_\lambda$ soit lisse (resp. non ramifié, resp. étale).
\end{enumerate}
\end{proposition}

(i) Si $y=f(x)$, $y_\lambda=f_\lambda(x_\lambda)$ est la projection canonique de $y$ dans $Y_\lambda$, et l'on a $f^{-1}(y)=f_\lambda^{-1}(y_\lambda)\otimes_{k(y_\lambda)}k(y)$; la partie de l'énoncé concernant les morphismes non ramifiés résulte donc de (17.7.1, (i)), et il suffit donc de considérer le cas des morphismes lisses. Comme $f$ et $f_\lambda$ sont localement de présentation finie, il revient au même de dire que $f^{-1}(y)$ est géométriquement régulier au point $x$ ou que $f_\lambda^{-1}(y_\lambda)$ est géométriquement régulier au point $x_\lambda$ (6.7.8). La proposition résulte donc de (17.5.1) et de (11.2.6).

(ii) Pour tout $\lambda$, soit $U_\lambda$ l'ensemble ouvert des $x_\lambda\in X_\lambda$ tels que $f_\lambda$ soit lisse (resp. non ramifié, resp. étale) au point $x_\lambda$; soit $V_\lambda$ son image réciproque dans $X$. Comme, par hypothèse, pour tout $x\in X$ il existe, en vertu de (i), un $\lambda$ tel que $f_\lambda$ soit lisse (resp. non ramifié, resp. étale) au point $x_\lambda$, $X$ est réunion des $V_\lambda$. D'ailleurs (17.3.3), pour $\lambda\leq\mu$ on a $V_\lambda\subset V_\mu$, donc, comme $X$ est quasi-compact, il existe un indice $\mu$ tel que $X=V_\mu$. Comme les $X_\lambda$ sont quasi-compacts, il résulte alors de (8.3.4) qu'il existe un indice $\nu\geq\mu$ tel que l'image réciproque de $U_\mu$ dans $X_\nu$ soit $X_\nu$ tout entier, ce qui signifie que $f_\nu$ est lisse (resp. non ramifié, resp. étale) par (17.3.3).

\begin{corollary}[17.7.9]
\label{IV.17.7.9}
Soient $S=\Spec(A)$ un schéma affine, $f:X\to S$ un morphisme. Les conditions suivantes sont équivalentes:
\begin{enumerate}
  \item[a)] $f$ est un morphisme de présentation finie et lisse (resp. non ramifié, resp. étale).

  \oldpage[IV-4]{76}

  \item[b)] Il existe un schéma affine noethérien $S_0=\Spec(A_0)$, un morphisme de type fini $f_0:X_0\to S_0$, et un morphisme $S\to S_0$ tels que le $S$-préschéma $X_0\otimes_{S_0}S$ soit $S$-isomorphe à $X$ et que $f_0$ soit lisse (resp. non ramifié, resp. étale).
  \item[c)] Les conditions de b) sont vérifiées, et en outre $A_0$ est une sous-$\mathbb{Z}$-algèbre de type fini de $A$, le morphisme $S\to S_0$ correspondant à l'injection canonique $A_0\to A$.
\end{enumerate}
\end{corollary}

La démonstration à partir de (17.7.8) est la même que celle de (11.2.7) à partir de (11.2.6).

\begin{proposition}[17.7.10]
\label{IV.17.7.10}
Soient $f:Y\to S$, $h:X\to S$ deux morphismes localement de présentation finie, $g:X\to Y$ un $S$-morphisme, $x$ un point de $X$, $y=g(x)$. Supposons que $h$ soit plat au point $x$ et que $f$ soit non ramifié au point $y$. Alors $f$ est étale au point $y$ et $g$ est plat au point $x$.
\end{proposition}

(On verra plus loin (18.4.9) que l'on peut en fait se dispenser de l'hypothèse que $h$ est localement de présentation finie).

La question étant locale sur $S$, $X$ et $Y$, on peut supposer que $S$, $X$ et $Y$ sont affines, $f,g,h$ des morphismes de présentation finie, $h$ plat et $f$ non ramifié. Compte tenu de (11.2.7) et (17.7.9), on peut supposer en outre que $S$, $X$ et $Y$ sont noethériens; enfin on peut se borner au cas où $S=\Spec(\sh{O}_{S,s})$ (où $s=f(y)=h(x)$). L'anneau $A=\sh{O}_{S,s}$ étant un anneau local noethérien, il existe un anneau local noethérien $B$ complet, de corps résiduel algébriquement clos et un homomorphisme local $A\to B$ faisant de $B$ un $A$-module fidèlement plat ($0_{\mathrm{III}}$, 10.3.1). Remplaçant $X$ et $Y$ par $X\times_S\Spec(B)$ et $Y\times_S\Spec(B)$, on conclut par (2.5.1) et (17.7.1) qu'on peut se borner au cas où $A$ est complet et a un corps résiduel algébriquement clos. L'hypothèse que $f$ est non ramifié au point $y$ entraîne alors (17.4.4) que $\sh{O}_{Y,y}$ est une $\sh{O}_{S,s}$-algèbre finie et un anneau local noethérien complet ($0_{\mathrm I}$, 7.4.2), et comme le corps résiduel de $\sh{O}_{S,s}$ est algébriquement clos, l'homomorphisme $\sh{O}_{S,s}\to\sh{O}_{Y,y}$ est surjectif (17.4.4). Mais, d'autre part, l'hypothèse que $h$ est plat au point $x$ entraîne que l'homomorphisme composé $\sh{O}_{S,s}\to\sh{O}_{Y,y}\to\sh{O}_{X,x}$ est injectif ($0_{\mathrm I}$, 6.5.1), donc l'homomorphisme $\sh{O}_{S,s}\to\sh{O}_{Y,y}$ est bijectif, ce qui montre que $f$ est étale au point $y$ (17.6.3); en outre, $\sh{O}_{X,x}$ est un $\sh{O}_{Y,y}$-module plat, donc $g$ est plat au point $x$.

\begin{proposition}[17.7.11]
\label{IV.17.7.11}
Soient $S$ un préschéma, $X$, $Y$ deux $S$-préschémas localement de présentation finie sur $S$, $f:X\to Y$ un $S$-morphisme. Pour tout $s\in S$, on note $X_s$, $Y_s$, $f_s$ les préschémas et le morphisme déduits de $X$, $Y$, $f$ par le changement de base $\Spec(k(s))\to S$. Alors:
\begin{enumerate}
  \item[(i)] L'ensemble des $x\in X$ tels que, si $s$ est l'image de $x$ dans $S$, $f_s:X_s\to Y_s$ soit lisse (resp. non ramifié, resp. étale, resp. différentiellement lisse) au point $x$, est localement constructible.
  \item[(ii)] Supposons $f$ de présentation finie. L'ensemble des $y\in Y$ tels que, si $s$ est l'image de $y$ dans $S$, $f_s$ soit lisse (resp. non ramifié, resp. étale, resp. différentiellement lisse) en tous les points de $f_s^{-1}(y)$, est localement constructible.
  \item[(iii)] Supposons $X$ et $Y$ de présentation finie sur $S$. L'ensemble des $s\in S$ tels que $f_s$ soit lisse (resp. non ramifié, resp. étale, resp. différentiellement lisse) est localement constructible.
\end{enumerate}
\end{proposition}

Soit $E$ l'ensemble des $x\in X$ vérifiant la propriété considérée dans (i). Alors l'ensemble $F$ des $y\in Y$ vérifiant la propriété correspondante considérée dans (ii) n'est

\oldpage[IV-4]{77}

autre que $Y-f(X-E)$, donc (ii) résulte de (i) et du théorème de Chevalley lorsque $f$ est de présentation finie (1.8.4). De même, si $h:Y\to S$ est le morphisme structural, l'ensemble des $s\in S$ vérifiant la propriété correspondante considérée dans (iii) est $S-h(Y-F)$, donc (iii) découle encore de (ii) et du théorème de Chevalley lorsque $f$ et $h$ sont de présentation finie. Il suffit donc de prouver les assertions de (i).

Prouvons d'abord (i) lorsqu'il s'agit de la propriété d'être \emph{lisse}. La question étant locale sur $X$, on peut se borner au cas où $S=\Spec(A)$, $X=\Spec(B)$ et $Y=\Spec(C)$ sont affines, $B$ et $C$ étant des $A$-algèbres de présentation finie. Raisonnant comme au début de (9.9.1) et utilisant (17.7.2, (ii)), on se ramène au cas où $A$ est noethérien. En vertu de ($0_{\mathrm{III}}$, 9.2.3), on est ramené à voir que si $x\in E$ (resp. si $x\notin E$), il existe un voisinage $V$ de $x$ dans $\overline{\{x\}}$ contenu dans $E$ (resp. dans $X-E$). Désignant par $g:X\to S$ et $h:Y\to S$ les morphismes structuraux, on peut d'abord remplacer $S$ par le sous-préschéma réduit $S'$ de $S$ ayant $\overline{\{g(x)\}}$ pour espace sous-jacent, $X$ par $g^{-1}(S')$, $Y$ par $h^{-1}(S')$, les fibres de $X$ et $X'$ (resp. $Y$ et $Y'$) aux points de $S'$ étant les mêmes. Autrement dit on peut se borner au cas où $S$ est \emph{intègre} et où $\eta=g(x)=h(y)$ (où $y=f(x)$) est son point générique.

$1^\circ$ Supposons d'abord que $x\in E$. Les anneaux locaux $\sh{O}_{X_\eta,x}$ et $\sh{O}_{Y_\eta,y}$ sont respectivement égaux à $\sh{O}_{X,x}$ et $\sh{O}_{Y,y}$; comme la propriété de lissité d'un morphisme de présentation finie en un point ne dépend que de l'anneau local de ce point et de l'anneau local de son image (17.5.1), on voit que l'hypothèse $x\in E$ revient à dire que le morphisme $f$ est lisse au point $x$; il possède alors encore cette propriété aux points d'un voisinage ouvert de $x$ dans $X$, et il suffit d'appliquer (17.3.3, (iii)) pour obtenir alors la conclusion.

$2^\circ$ Supposons en second lieu que $x\in X-E$, et que le morphisme $f_\eta$ ne soit pas plat au point $x$. La conclusion résulte alors du lemme suivant qui précise (11.2.8):

\begin{lemma}[17.7.11.1]
\label{IV.17.7.11.1}
Soient $g:X\to S$, $h:Y\to S$ deux morphismes localement de présentation finie, $f:X\to Y$ un $S$-morphisme, $\sh{F}$ un $\sh{O}_X$-module quasi-cohérent de présentation finie. Alors l'ensemble $E$ des $x\in X$ tels que $\sh{F}_{g(x)}$ soit $f_{g(x)}$-plat au point $x$ est localement constructible.
\end{lemma}

Raisonnant encore comme au début de (9.9.1) et utilisant (2.5.1), on se ramène au cas où $S$, $X$ et $Y$ sont noethériens; puis on se ramène comme ci-dessus au cas où $S$ est intègre, où $\eta=g(x)=h(f(x))$ est son point générique, et il s'agit de montrer que si $x\in E$ (resp. $x\notin E$) il y a un voisinage $V$ de $x$ dans $\overline{\{x\}}$ contenu dans $E$ (resp. dans $X-E$). Le cas où $x\in E$ est conséquence immédiate de (11.1.1). Pour traiter le cas où $x\notin E$, on raisonne comme dans (9.4.7.1), dont nous conservons les notations, de sorte que l'on peut supposer, en remplaçant éventuellement $Y$ et $X$ par des voisinages de $f(x)$ et $x$ respectivement, qu'il existe deux $\sh{O}_Y$-modules cohérents $\sh{G}$, $\sh{H}$ et un $\sh{O}_Y$-homomorphisme $u:\sh{G}\to\sh{H}$ tels que pour tout $s\in S$, $u_s:\sh{G}_s\to\sh{H}_s$ soit injectif, mais que l'homomorphisme $1\otimes u_\eta:\sh{F}_\eta\otimes_{\sh{O}_{Y_\eta}}\sh{G}_\eta\to\sh{F}_\eta\otimes_{\sh{O}_{Y_\eta}}\sh{H}_\eta$ ne soit pas injectif au point $x$, autrement dit $x\in\operatorname{Supp}(\operatorname{Ker}(1\otimes u_\eta))$; si l'on pose $T=\overline{\{x\}}$, on a donc $\operatorname{Supp}(\operatorname{Ker}(1\otimes u_\eta))\supset T_\eta$. Mais on peut supposer que pour tout $s\in S$, on a $\operatorname{Ker}(1\otimes u_s)=(\operatorname{Ker}(1\otimes u))_s$ (9.4.2), donc

\oldpage[IV-4]{78}

$\operatorname{Supp}(\operatorname{Ker}(1\otimes u_s))=(\operatorname{Supp}(\operatorname{Ker}(1\otimes u)))_s$ (I, 9.1.13.1); il résulte enfin de (9.5.2) que pour $s$ dans un voisinage de $\eta$, on a $(\operatorname{Supp}(\operatorname{Ker}(1\otimes u)))_s\supset T_s$, ce qui établit le lemme.

$3^\circ$ Supposons maintenant que $x\in X-E$, que le morphisme $f_\eta$ soit plat au point $x$, mais que $f_\eta$ ne soit pas lisse au point $x$. Remarquons que dire que $f_\eta$ est plat au point $x$ équivaut à dire que $f$ lui-même est plat au point $x$ et en remplaçant $X$ par un voisinage de $x$, on peut supposer que $f$ est plat (11.1.1); on en conclut qu'il en est de même de $f_s$ pour tout $s\in S$, et comme pour tout $y\in Y$, $f^{-1}(y)=f_{h(y)}^{-1}(y)$, il revient au même de dire que $f_{g(x')}$ est lisse au point $x'$ ou de dire que $f$ est lisse au point $x'$. Mais l'ensemble des $x'\in X$ où $f$ est lisse est \emph{ouvert} dans $X$ (12.1.7), donc l'ensemble des $x'\in X$ où $f$ n'est pas lisse est fermé, et puisqu'il contient $x$ par hypothèse, il contient aussi $\overline{\{x\}}$, ce qui achève la démonstration pour la première propriété considérée dans (i).

Prouvons en second lieu (i) lorsqu'il s'agit de la propriété d'être \emph{étale}. Notons pour cela que cette propriété pour $f_s$ au point $x$ équivaut à dire que $f_s$ est à la fois lisse et quasi-fini au point $x$ (17.6.1). Or, il revient au même de dire que $f_{g(x)}$ est quasi-fini au point $x$ ou que $f$ est lui-même quasi-fini en ce point; il résulte donc de (13.1.4) que l'ensemble des points $x$ tels que $f_{g(x)}$ soit quasi-fini au point $x$ est ouvert dans $X$, et \emph{a fortiori} localement constructible; la conclusion résulte donc de ce que l'ensemble des $x$ tels que $f_{g(x)}$ soit lisse au point $x$ est lui aussi localement constructible.

Passons à la preuve de (i) lorsqu'il s'agit de la propriété d'être \emph{différentiellement lisse}. Soit $p:X\times_Y X\to X$ la seconde projection canonique; la seconde projection $X_s\times_{Y_s}X_s\to X_s$ n'est autre que $p_s$ pour tout $s\in S$, et il résulte donc de (17.12.5.1)\footnote{Le lecteur vérifiera que le résultat de (17.7.11) n'est pas utilisé dans la suite du § 17 et qu'il n'y a donc pas de cercle vicieux.} que pour que $f_{g(x)}$ soit différentiellement lisse au point $x$, il faut et il suffit que $p_{g(x)}$ soit lisse au point $\Delta_f(x)$. Comme $p$ est localement de présentation finie, l'ensemble des points $z\in X\times_Y X$ tels que $p_{g(p(z))}$ soit lisse au point $z$ est localement constructible, et l'intersection de cet ensemble avec l'ensemble localement fermé $\Delta_f(X)$ est donc aussi localement constructible dans $\Delta_f(X)$ (1.8.2); la restriction de $p$ à $\Delta_f(X)$ étant un isomorphisme sur $X$, on voit que l'ensemble des $x\in X$ tels que $f_{g(x)}$ soit différentiellement lisse au point $x$ est localement constructible.

Considérons enfin la propriété d'être non ramifié; notons que le morphisme diagonal $\Delta_f:X\to X\times_Y X$ est une immersion localement de présentation finie (1.4.3.1) et pour tout $s\in S$, le morphisme diagonal $\Delta_{f_s}:X_s\to X_s\times_{Y_s}X_s$ n'est autre que $(\Delta_f)_s$; dire que $f_{g(x)}$ est non ramifié au point $x$ équivaut à dire que $(\Delta_f)_{g(x)}$ est un isomorphisme local au point $x$ (17.4.1), et puisque c'est une immersion localement de présentation finie, il revient au même (17.9.1)\footnote{Le lecteur vérifiera que le résultat de (17.7.11) n'est pas utilisé dans le reste du § 17 et qu'il n'y a donc pas de cercle vicieux.} de dire que $(\Delta_f)_{g(x)}$ est étale au point $x$; il suffit donc d'appliquer ce qui a été vu plus haut pour la propriété d'être étale.

\oldpage[IV-4]{79}

\subsection{Critères de lissité et non-ramification par fibres}
\EGAReaderTarget{ega4.c04.subsec.17-8}{subsection:IV.17.8}
\label{IV.17.8}

\begin{proposition}[17.8.1]
\label{IV.17.8.1}
Soient $g:Y\to S$, $h:X\to S$ deux morphismes localement de présentation finie. Pour qu'un $S$-morphisme $f:X\to Y$ soit non ramifié, il faut et il suffit que pour tout $s\in S$, le morphisme $f_s:h^{-1}(s)\to g^{-1}(s)$ déduit de $f$ par le changement de base $\Spec(k(s))\to S$, soit non ramifié.
\end{proposition}

Cela résulte aussitôt du fait que, pour un morphisme localement de présentation finie, le fait d'être non ramifié est une propriété des \emph{fibres} de ce morphisme (17.4.1, $d$)), et de ce que, pour tout $y\in Y$, on a $f^{-1}(y)=f_s^{-1}(y)$ si $s=g(y)$.

\begin{proposition}[17.8.2]
\label{IV.17.8.2}
Soient $g:Y\to S$, $h:X\to S$ deux morphismes localement de présentation finie; on suppose en outre que $h$ est plat. Pour qu'un $S$-morphisme $f:X\to Y$ soit lisse (resp. étale), il faut et il suffit que, pour tout $s\in S$, le morphisme $f_s:h^{-1}(s)\to g^{-1}(s)$ déduit de $f$ par le changement de base $\Spec(k(s))\to S$ soit lisse (resp. étale). Lorsqu'il en est ainsi, le morphisme $g$ est plat aux points de $f(X)$.
\end{proposition}

On sait en effet (11.3.10) que pour que $f$ soit plat, il faut et il suffit que $f_s$ le soit pour tout $s\in S$, et qu'alors $g$ est plat aux points de $f(X)$. Mais pour un morphisme plat localement de présentation finie, le fait d'être lisse est une propriété des fibres de ce morphisme (17.5.1, $b$)), et pour tout $y\in Y$, on a $f^{-1}(y)=f_s^{-1}(y)$ si $s=g(y)$.

\begin{remark}[17.8.3]
\label{IV.17.8.3}
Les démonstrations précédentes montrent (compte tenu de (11.3.10)) que si les hypothèses sur $g$ et $h$ sont les mêmes que ci-dessus, alors, pour que $f$ soit non ramifié (resp. lisse, resp. étale) \emph{en un point} $x\in X$, il suffit que, si l'on pose $s=h(x)$, $f_s$ soit non ramifié (resp. lisse, resp. étale) au point $x$.
\end{remark}

\subsection{Morphismes étales et immersions ouvertes}
\EGAReaderTarget{ega4.c04.subsec.17-9}{subsection:IV.17.9}
\label{IV.17.9}

\begin{theorem}[17.9.1]
\label{IV.17.9.1}
Soit $f:X\to Y$ un morphisme. Les propriétés suivantes sont équivalentes:
\begin{enumerate}
  \item[a)] $f$ est une immersion ouverte.
  \item[b)] $f$ est un monomorphisme plat localement de présentation finie.
  \item[c)] $f$ est étale et radiciel.
\end{enumerate}
\end{theorem}

Il résulte de (1.4.3, (i)) que a) implique b). La condition b) implique que pour tout $y\in Y$, la fibre $f^{-1}(y)$ est vide ou isomorphe à $\Spec(k(y))$ (8.11.5.1), donc b) entraîne c), en vertu de (17.6.2, $c$)). Reste à voir que c) implique a).

La question étant locale sur $X$ et sur $Y$ (puisque $f$ est injectif), on peut se borner au cas où $Y$ est affine et $f$ de présentation finie. Comme $f$ est plat, c'est un morphisme \emph{ouvert} (2.4.6), donc, en remplaçant $Y$ par $f(X)$, on peut supposer que $f$ est \emph{surjectif}. Pour tout morphisme $Y'\to Y$, $f'=f_{(Y')}:X_{(Y')}\to Y'$ est encore étale, radiciel, surjectif et de présentation finie, donc ouvert, et par suite un homéomorphisme; autrement dit $f$ est un \emph{homéomorphisme universel}, et étant de type fini et séparé (1.8.7.1), $f$ est \emph{propre}. L'hypothèse que $f$ est radiciel et de type fini entraîne alors que $f$ est quasi-fini; donc (8.11.1) $f$ est un morphisme \emph{fini}. Pour prouver que $f$ est un isomorphisme, on peut se borner au cas

\oldpage[IV-4]{80}

où $Y=\Spec(A)$, $A$ étant un anneau local. Comme $f$ est de présentation finie, on a $X=\Spec(B)$, où $B$ est un $A$-module plat de présentation finie (1.4.7), donc \emph{libre} (Bourbaki, \emph{Alg. comm.}, chap. II, § 5, n\textsuperscript{o} 2, cor. 2 du th. 1). En outre, si $\mathfrak m$ est l'idéal maximal de $A$ et $k$ son corps résiduel, $B/\mathfrak mB$ est par hypothèse un corps, à la fois extension radicielle et extension finie séparable de $k$, puisque $f$ est étale et radiciel (17.6.1); donc $B/\mathfrak mB$ est isomorphe à $k$, et comme $B$ est un $A$-module libre, $B$ est isomorphe à $A$. C.Q.F.D.

\begin{corollary}[17.9.2]
\label{IV.17.9.2}
Soit $X$ un préschéma connexe; si $f:X\to Y$ est une immersion fermée étale, alors $f$ est un isomorphisme de $X$ sur une composante connexe ouverte de $Y$.
\end{corollary}

En effet $f$ est étale et radiciel, donc un isomorphisme de $X$ sur un préschéma induit sur une partie ouverte de $Y$; mais par hypothèse $f(X)$ est fermé dans $Y$, donc à la fois ouvert et fermé, et puisque $f(X)$ est connexe, c'est une composante connexe de $Y$.

\begin{corollary}[17.9.3]
\label{IV.17.9.3}
Soit $f:X\to Y$ un morphisme étale (resp. étale et séparé). Alors toute $Y$-section $g:Y\to X$ de $X$ est une immersion ouverte (resp. ouverte et fermée); en outre l'application $g\mapsto g(Y)$ est une bijection de l'ensemble $\Gamma(X/Y)$ des $Y$-sections de $X$ sur l'ensemble des parties ouvertes $Z$ (resp. ouvertes et fermées) de $X$ telles que la restriction de $f$ à $Z$ soit un morphisme surjectif et radiciel de $Z$ sur $Y$.
\end{corollary}

En effet, le fait que $g$ soit une immersion ouverte résulte déjà de ce que $f$ est non ramifié (17.4.1, $b''$)), et la restriction de $f$ à l'ouvert $g(Y)$ de $X$ est un isomorphisme. Inversement, si $Z$ est un ouvert de $X$ tel que $f|Z$ soit un morphisme surjectif et radiciel de $Z$ sur $Y$, $f|Z$ est un isomorphisme en vertu de (17.9.1), puisqu'il est étale. Si $f$ est en outre séparé, on sait que $g$ est une immersion fermée (1, 5.4.6), ce qui achève de prouver le corollaire.

\begin{corollary}[17.9.4]
\label{IV.17.9.4}
Soient $Y$ un préschéma connexe, $f:X\to Y$ un morphisme étale et séparé. Alors toute $Y$-section $g$ de $X$ est un isomorphisme de $Y$ sur une composante connexe ouverte de $X$, et l'application $g\mapsto g(Y)$ est une bijection de $\Gamma(X/Y)$ sur l'ensemble des composantes connexes ouvertes $Z$ de $X$ telles que la restriction de $f$ à $Z$ soit un morphisme surjectif et radiciel de $Z$ sur $Y$.
\end{corollary}

\begin{corollary}[17.9.5]
\label{IV.17.9.5}
Soient $g:Y\to S$, $h:X\to S$ deux morphismes localement de présentation finie; on suppose en outre que $h$ est plat. Pour qu'un $S$-morphisme $f:X\to Y$ soit une immersion ouverte (resp. un isomorphisme), il faut et il suffit que, pour tout $s\in S$, le morphisme $f_s:h^{-1}(s)\to g^{-1}(s)$ déduit de $f$ par le changement de base $\Spec(k(s))\to S$, soit une immersion ouverte (resp. un isomorphisme).
\end{corollary}

En effet, si $f_s$ est une immersion ouverte pour tout $s\in S$, il résulte de (17.8.3) que $f$ est un morphisme étale; comme pour tout $y\in Y$, on a $f^{-1}(y)=f_s^{-1}(y)$ avec $s=g(y)$, $f$ est radiciel; donc $f$ est une immersion ouverte en vertu de (17.9.1). Si de plus $f_s$ est surjectif pour tout $s\in S$, $f$ est surjectif, donc un isomorphisme.

La proposition suivante précise (10.4.11):

\begin{proposition}[17.9.6]
\EGAReaderTarget{ega4.c04.prop.17-9-6}{IV.17.9.6}
Soient $S$ un préschéma, $X$ un $S$-préschéma de présentation finie. Tout $S$-endomorphisme de $X$ qui est un monomorphisme est un automorphisme de $X$.
\end{proposition}

Soient $f:X\to S$ le morphisme structural, $g$ le $S$-endomorphisme considéré. La question est locale sur $S$, et on peut par suite supposer que $S$ est affine. Utilisant (8.9.1) et (8.10.5, (i \emph{bis})), on est ramené au cas où $S=\Spec(A)$, où $A$ est une $\mathbb{Z}$-algèbre de type fini, et par suite $X$ est un $\mathbb{Z}$-préschéma de type fini. Il résulte déjà de (10.4.11) que $g$ est un morphisme bijectif, puisqu'un monomorphisme est radiciel (8.11.5.1); il suffira donc de montrer

\oldpage[IV-4]{81}

que $g$ est une immersion ouverte, et puisque $g$ est radiciel, il suffira, en vertu de (17.9.1), de prouver que $g$ est étale. En outre, comme l'ensemble des points où $g$ est étale est ouvert et que $X$ est un préschéma de Jacobson (10.4.7), il suffira de montrer que $g$ est étale en tout point \emph{fermé} $z$ de $X$ (10.3.1). Posons $z'=g(z)$; il résulte de (10.4.11.1, (i)) que $z'$ est aussi un point fermé de $X$ et puisque $g$ est un monomorphisme, l'application $k(z')\to k(z)$ déduite de $g$ est un \emph{isomorphisme}. Pour que $g$ soit étale au point $z$, il faut et il suffit donc que l'homomorphisme canonique $\widehat{\sh{O}}_{X,z'}\to\widehat{\sh{O}}_{X,z}$ soit bijectif (17.6.3, $e''$)). Nous allons pour cela prouver que pour tout entier $n$, l'homomorphisme canonique $\sh{O}_{X,z'}/\mathfrak m_{z'}^{n+1}\to\sh{O}_{X,z}/\mathfrak m_z^{n+1}$ est bijectif, d'où résultera aussitôt la conclusion. On peut supposer que $z$ appartient à l'ensemble fini $T_{p,d}$ des points fermés $t$ de $X$ tels que $k(t)$ soit une extension de $\mathbb{F}_p$ dont le degré divise $d$; auquel cas il en est de même de $z'$. Désignons par $\sh{J}$ l'idéal cohérent de $\sh{O}_X$ tel que $\sh{J}_x=\sh{O}_x$ pour $x\notin T_{p,d}$, et $\sh{J}_x=\mathfrak m_x^{n+1}$ pour $x\in T_{p,d}$. Soient $T_{p,d,n}$ le sous-préschéma fermé de $X$ défini par $\sh{J}$, $j:T_{p,d,n}\to X$ l'injection canonique; le morphisme composé $g\circ j:T_{p,d,n}\to X$ applique l'ensemble $T_{p,d}$ dans lui-même, et pour tout $x\in T_{p,d}$, l'homomorphisme $\sh{O}_{X,g(x)}\to\sh{O}_{X,x}/\mathfrak m_x^{n+1}$ se factorise en
\[
\sh{O}_{X,g(x)}\to\sh{O}_{X,g(x)}/\mathfrak m_{g(x)}^{n+1}\to\sh{O}_{X,x}/\mathfrak m_x^{n+1}.
\]
Donc (I, 4.9), il existe un unique endomorphisme $g'$ de $T_{p,d,n}$ tel que $j\circ g'=g\circ j$. Comme $j$ est un monomorphisme, et qu'il en est de même de $g$ par hypothèse, on en déduit que $g'$ est aussi un monomorphisme, et la proposition sera prouvée si l'on montre que $g'$ est un \emph{automorphisme} de $T_{p,d,n}$. Or, pour tout $x\in T_{p,d,n}$, on a vu que $k(x)$ est un corps fini, et comme $\sh{O}_{X,x}$ est noethérien, chacun des $\mathfrak m_x^h/\mathfrak m_x^{h+1}$ est un $k(x)$-espace vectoriel de rang fini; d'où l'on conclut aussitôt que l'anneau local $\sh{O}_{X,x}/\mathfrak m_x^{n+1}$ de $T_{p,d,n}$ au point $x$ a un nombre \emph{fini} d'éléments, et \emph{a fortiori} est un $\mathbb{Z}$-module de type fini. Comme $T_{p,d,n}$ est somme des préschémas $\Spec(\sh{O}_{X,x}/\mathfrak m_x^{n+1})$ en nombre fini, c'est un $\mathbb{Z}$-préschéma \emph{fini}; l'endomorphisme $g'$ est donc lui aussi \emph{fini} (II, 6.1.5, (v)), donc propre. Mais un monomorphisme propre est une \emph{immersion fermée} (8.11.5); si $T_{p,d,n}=\Spec(B)$, $g'$ correspond donc à un endomorphisme surjectif $\varphi:B\to B$ de l'anneau $B$. Or, l'ensemble $B$ est fini, donc $\varphi$ est nécessairement bijectif. C.Q.F.D.

\subsection{Dimension relative d'un préschéma lisse sur un autre}
\EGAReaderTarget{ega4.c04.subsec.17-10}{subsection:IV.17.10}

\begin{definition}[17.10.1]
\EGAReaderTarget{ega4.c04.def.17-10-1}{IV.17.10.1}
Soit $f:X\to Y$ un morphisme localement de type fini. On appelle \emph{dimension relative de $f$ au point $x\in X$} (ou \emph{dimension relative de $X$ sur $Y$ au point $x$}) et l'on note $\dim_x f$ l'entier positif $\dim_x(f^{-1}(f(x)))$.
\end{definition}

Dire que $f$ est \emph{quasi-fini} au point $x$ (Err$_{\mathrm{III}}$, 20) équivaut donc à dire que $\dim_x f=0$. On a vu (13.1.3) que la fonction $x\mapsto\dim_x f$ est semi-continue supérieurement. On notera que, même lorsque le morphisme $f$ a la propriété $(S_1)$ (autrement dit (6.8.1) est plat et tel que ses fibres n'aient pas de cycle premier associé immergé), la fonction $x\mapsto\dim_x f$ n'est pas nécessairement continue, comme le montre l'exemple où $Y=\Spec(k)$, où $k$ est un corps, et $X=\Spec(k[U,V,W]/\mathfrak p\mathfrak q)$, où $\mathfrak p=(W)$ et $\mathfrak q=(U)+(V-W)$, idéaux premiers de $k[U,V,W]$ ($X$ étant donc la réunion, dans l'espace à 3 dimensions, d'un plan et d'une droite non parallèle à ce plan).

Dire qu'un morphisme localement de présentation finie $f:X\to Y$ est étale au point $x\in X$ signifie encore que $f$ est lisse au point $x$ et que $\dim_x f=0$ (17.6.1).

\begin{proposition}[17.10.2]
\label{IV.17.10.2}
Soit $f:X\to Y$ un morphisme lisse. Pour tout $x\in X$, le $\sh{O}_X$-Module localement libre $\Omega_f^1$ (17.2.3) a un rang en $x$ égal à $\dim_x f$ (ce qui entraîne que $x\mapsto\dim_x f$ est une fonction continue dans $X$).
\end{proposition}

En effet, si $y=f(x)$, $X_y=f^{-1}(y)$ est lisse sur $k(y)$, et si $f_y:X_y\to\Spec(k(y))$ est le morphisme structural, on a $\Omega_{f_y}^1=\Omega_f^1\otimes_{\sh{O}_X}\sh{O}_{X_y}$ (16.4.5); on peut donc se borner au cas où $Y=\Spec(k)$ est le spectre d'un corps; en outre, en vertu de (16.4.5) et (17.7.1), on peut remplacer $k$ par une extension algébriquement close, autrement dit

\oldpage[IV-4]{82}

supposer $k$ algébriquement clos. L'ensemble des points de $X$ rationnels sur $k$ étant alors dense dans $X$ (10.4.8), on peut se borner au cas où $x$ est rationnel sur $k$. Or (16.4.12) $(\Omega^1_{X/k})_x\otimes_{\sh{O}_{X,x}}k(x)$ est alors $k$-isomorphe à $\mathfrak m_x/\mathfrak m_x^2$, et puisque $\sh{O}_x$ est un anneau local régulier (17.5.1), $\operatorname{rg}_{k(x)}(\mathfrak m_x/\mathfrak m_x^2)=\dim(\sh{O}_x)$ (0, 17.1.1); donc $\operatorname{rg}_{\sh{O}_x}(\Omega^1_{X/k})_x=\dim(\sh{O}_x)$. Mais puisque $k(x)=k$, on a $\dim_x f=\dim(\sh{O}_x)$ (5.2.3). C.Q.F.D.

Nous démontrerons plus loin une réciproque de ce résultat (17.15.5).

\begin{corollary}[17.10.3]
\label{IV.17.10.3}
Soient $f:X\to Y$, $g:Y\to Z$ deux morphismes lisses. Alors, pour tout $x\in X$, on a
\begin{equation}
\label{IV.17.10.3.1}
\tag{17.10.3.1}
\dim_x(g\circ f)=\dim_x f+\dim_{f(x)}g.
\end{equation}
\end{corollary}

En effet, $g\circ f$ est lisse (17.3.3), donc les trois $\sh{O}_X$-Modules $\Omega^1_{X/Y}$, $\Omega^1_{X/Z}$ et $f^*(\Omega^1_{Y/Z})$ sont localement libres (17.2.3 et $0_{\mathrm I}$, 5.4.5); en outre le rang en $x$ de $f^*(\Omega^1_{Y/Z})$ est égal au rang en $y=f(x)$ de $\Omega^1_{Y/Z}$. L'égalité (17.10.3.1) est donc conséquence de (17.10.2) et de l'exactitude de la suite (17.2.3.1).

\begin{corollary}[17.10.4]
\label{IV.17.10.4}
Soient $f:X\to Y$ un morphisme lisse, $X'$ un sous-préschéma de $X$ tel que le morphisme composé $X'\xrightarrow{j}X\xrightarrow{f}Y$ (où $j$ est l'injection canonique) soit lisse. Alors le faisceau conormal $\sh{N}_{X'/X}$ est un $\sh{O}_{X'}$-Module localement libre, et pour tout $x\in X'$, on a
\begin{equation}
\label{IV.17.10.4.1}
\tag{17.10.4.1}
\dim_x f=\dim_x(f\circ j)+\operatorname{rg}_{\sh{O}_{X',x}}(\sh{N}_{X'/X})_x.
\end{equation}
\end{corollary}

En effet, $\Omega^1_{X/Y}\otimes_{\sh{O}_X}\sh{O}_{X'}$ et $\Omega^1_{X'/Y}$ sont alors localement libres et la suite exacte (17.2.5.1) est scindée dans un voisinage convenable de chaque point de $X'$, donc $\sh{N}_{X'/X}$ est localement libre (Bourbaki, \emph{Alg. comm.}, chap. II, § 5, n\textsuperscript{o} 2, th. 1), et la relation (17.10.4.1) résulte aussitôt de l'exactitude de la suite (17.2.5.1).

\subsection{Morphismes lisses de préschémas lisses}
\EGAReaderTarget{ega4.c04.subsec.17-11}{subsection:IV.17.11}

\begin{theorem}[17.11.1]
\label{IV.17.11.1}
Soient $f:Y\to S$ et $h:X\to S$ deux morphismes localement de présentation finie, $g:X\to Y$ un $S$-morphisme, $x$ un point de $X$; posons $y=g(x)$, $s=f(y)=h(x)$. Les conditions suivantes sont équivalentes:
\begin{enumerate}
  \item[a)] $f$ est lisse au point $y$ et $g$ est lisse au point $x$.
  \item[b)] $g$ et $h$ sont lisses au point $x$.
  \item[c)] $h$ est lisse au point $x$, et l'homomorphisme canonique (16.4.18)
  \begin{equation}
  \label{IV.17.11.1.1}
  \tag{17.11.1.1}
  (g^*(\Omega^1_{Y/S}))_x\to(\Omega^1_{X/S})_x
  \end{equation}
  est inversible à gauche (autrement dit est un isomorphisme sur un \emph{facteur direct} de $(\Omega^1_{X/S})_x$).
  \item[c')] $h$ est lisse au point $x$, et l'homomorphisme canonique
  \begin{equation}
  \label{IV.17.11.1.2}
  \tag{17.11.1.2}
  ((\Omega^1_{Y/S})_y\otimes_{\sh{O}_{Y,y}}k(y))\otimes_{k(y)}k(x)\to(\Omega^1_{X/S})_x\otimes_{\sh{O}_{X,x}}k(x)
  \end{equation}
  est injectif.
\end{enumerate}
Supposons en outre que l'homomorphisme $k(y)\to k(x)$ soit bijectif. Alors les conditions précédentes sont aussi équivalentes à la suivante:
\begin{enumerate}
  \item[d)] $h$ est lisse au point $x$, et l'application canonique $\mathrm{T}_{X/S}(x)\to\mathrm{T}_{Y/S}(y)$ de l'espace vectoriel tangent en $x$ à $X$ dans l'espace vectoriel tangent en $y$ à $Y$ (16.5.12) est surjective.
\end{enumerate}
\end{theorem}

Le fait que a) entraîne b) résulte trivialement de (17.3.3, (iii)); b) entraîne c) par application de (17.2.3, (ii)). Pour voir que c) est équivalente à c'), notons qu'en vertu de (17.2.3, (i)), $\Omega^1_{X/S}$ est un $\sh{O}_X$-Module localement libre de type fini; il suffit alors d'appliquer (0, 19.1.12) à l'anneau local $\sh{O}_{X,x}$ et à l'homomorphisme (17.11.1.1) de $\sh{O}_{X,x}$-modules de type fini. Lorsque $k(x)=k(y)$, l'application linéaire tangente $\mathrm{T}_{X/S}(x)\to\mathrm{T}_{Y/S}(y)$ est transposée de (17.11.1.2), par (16.5.12), ce qui établit dans ce cas l'équivalence de c') et d).

\oldpage[IV-4]{83}

Il reste donc à prouver que c) entraîne a). On peut se borner au cas où $S=\Spec(A)$, $Y=\Spec(B)$, $X=\Spec(C)$ sont affines. L'hypothèse implique (17.2.3, (i)) que $(\Omega^1_{X/S})_x=\Omega^1_{C_{\mathfrak m}/A_{\mathfrak s}}$ est un $C_{\mathfrak m}$-module libre: en effet (16.10.6) il existe des éléments $t_i$ ($1\leq i\leq r$) de $B_{\mathfrak p}=\sh{O}_{Y,y}$ tels que les différentielles $d_{B_{\mathfrak p}/A_{\mathfrak s}}(t_i)$ engendrent le $B_{\mathfrak p}$-module $\Omega^1_{B_{\mathfrak p}/A_{\mathfrak s}}$ et que leurs images dans $\Omega^1_{C_{\mathfrak m}/A_{\mathfrak s}}$ fassent partie d'une base de ce $C_{\mathfrak m}$-module libre. Comme $\Omega^1_{Y/S}$ (resp. $\Omega^1_{X/S}$) est un $\sh{O}_Y$-Module (resp. un $\sh{O}_X$-Module) de présentation finie (16.4.22), on peut, en remplaçant au besoin $X$ et $Y$ par des voisinages ouverts affines convenables de $x$ et $y$ respectivement, supposer que $\Omega^1_{C/A}$ est un $C$-module libre et que les $t_i$ sont les images d'éléments $s_i$ ($1\leq i\leq r$) de $B$ tels que les $d_{B/A}(s_i)$ engendrent le $B$-module $\Omega^1_{B/A}$ et que leurs images dans $\Omega^1_{C/A}$ fassent partie d'une base de ce $C$-module (Bourbaki, \emph{Alg. comm.}, chap. II, § 5, n\textsuperscript{o} 1, prop. 2). Soit $\varphi$ le $A$-homomorphisme de $B'=A[T_1,\ldots,T_r]$ dans $B$ tel que $\varphi(T_i)=s_i$ pour tout $i$; le di-homomorphisme correspondant $\Omega^1_{B'/A}\to\Omega^1_{B/A}$ (0, 20.5.2) transforme les $d_{B'/A}(T_i)$, qui forment une base de $\Omega^1_{B'/A}$ (0, 20.4.13), en les $d_{B/A}(s_i)$ et est par suite \emph{surjectif}; si $Y'=\Spec(B')$, et si $u:Y\to Y'$ est le $S$-morphisme correspondant à $\varphi$, on conclut de (17.2.2) que $u$ est \emph{non ramifié}. Si l'on prouve que le morphisme composé $u\circ g:X\to Y'$ est lisse au point $x$, il résultera donc de (17.7.10) que $u$ est étale au point $y$, puis de (17.3.5) que $g$ est lisse au point $x$; enfin, comme le morphisme structural $f':Y'\to S$ est lisse (17.3.8), $f=f'\circ u$ sera lisse au point $y$. Comme les images canoniques des $d_{B'/A}(T_i)$ dans $\Omega^1_{C/A}$ sont celles des $d_{B/A}(s_i)$, elles font partie d'une base de $\Omega^1_{C/A}$. On voit ainsi que pour prouver que c) implique a), on peut se borner au cas où $Y'=Y$, et par suite supposer que $f$ est un morphisme \emph{lisse}.

En vertu de (17.8.2), on peut alors se borner au cas où $S=\Spec(k)$ est le spectre d'un corps; de plus, grâce à (17.7.1, (ii)), on peut supposer que $k$ est algébriquement clos; enfin, remplaçant au besoin $X$ et $Y$ par des voisinages ouverts de $x$ et $y$ respectivement, on peut supposer que $f$ et $h$ sont tous deux lisses, et que l'homomorphisme canonique
\[
g^*(\Omega^1_{Y/S})\to\Omega^1_{X/S}
\]
est inversible à gauche (0, 19.1.12). Alors l'ensemble des points de $X$ rationnels sur $k$ est très dense dans $X$ (10.4.8), donc, pour prouver que $g$ est lisse, il suffit de prouver que $g$ est lisse en tout point $x\in X$ rationnel sur $k$, ou encore qu'en un tel point, $\sh{O}_{X,x}$ est un $\sh{O}_{Y,y}$-module plat et $\sh{O}_{X,x}/\mathfrak m_y\sh{O}_{X,x}$ un anneau régulier (17.5.1). Or, $y=g(x)$ est

\oldpage[IV-4]{84}

\emph{a fortiori} rationnel sur $k$, et comme $\sh{O}_{Y,y}$ est alors une $k$-algèbre formellement lisse pour les topologies discrètes et que $\sh{O}_{Y,y}/\mathfrak m_y=k$, il résulte de (0, 20.5.14) que l'homomorphisme canonique $\mathfrak m_y/\mathfrak m_y^2\to(\Omega^1_{Y/S})_y\otimes_{\sh{O}_{Y,y}}k(y)$ est bijectif; de même, l'homomorphisme canonique $\mathfrak m_x/\mathfrak m_x^2\to(\Omega^1_{X/S})_x\otimes_{\sh{O}_{X,x}}k(x)$ est bijectif. L'hypothèse c'), équivalente à c), signifie donc ici que l'homomorphisme canonique
\[
(\mathfrak m_y/\mathfrak m_y^2)\otimes_{k(y)}k(x)\to\mathfrak m_x/\mathfrak m_x^2
\]
est injectif. Comme l'anneau $\sh{O}_{X,x}$ est régulier, la conclusion résulte de (0, 17.3.3). C.Q.F.D.

\begin{corollary}[17.11.2]
\label{IV.17.11.2}
Sous les hypothèses générales de (17.11.1), les conditions suivantes sont équivalentes:
\begin{enumerate}
  \item[a)] $f$ est lisse au point $y$ et $g$ est étale au point $x$.
  \item[b)] $h$ est lisse au point $x$ et $g$ est étale au point $x$.
  \item[c)] $h$ est lisse au point $x$, et l'homomorphisme canonique
  \[
  (g^*(\Omega^1_{Y/S}))_x\to(\Omega^1_{X/S})_x
  \]
  est bijectif.
  \item[c')] $h$ est lisse au point $x$, et l'homomorphisme canonique
  \[
  ((\Omega^1_{Y/S})_y\otimes_{\sh{O}_{Y,y}}k(y))\otimes_{k(y)}k(x)\to(\Omega^1_{X/S})_x\otimes_{\sh{O}_{X,x}}k(x)
  \]
  est bijectif.
\end{enumerate}
Supposons en outre que l'homomorphisme $k(y)\to k(x)$ soit bijectif. Alors les conditions précédentes sont aussi équivalentes à la suivante:
\begin{enumerate}
  \item[d)] $h$ est lisse au point $x$, et l'application canonique $\mathrm{T}_{X/S}(x)\to\mathrm{T}_{Y/S}(y)$ (16.5.12) est bijective.
\end{enumerate}
\end{corollary}

Chacune des conditions a), b), c) de (17.11.2) est équivalente à la conjonction de la condition correspondante de (17.11.1) et du fait que $g$ est non ramifié au point $x$, compte tenu de (17.2.2) en ce qui concerne la condition c); d'où l'équivalence de a), b) et c). L'équivalence de c) et c') résulte de l'équivalence des conditions correspondantes de (17.11.1) et du lemme de Nakayama; l'équivalence de c') et d) lorsque $k(x)=k(y)$ est immédiate par transposition (16.5.12).

\begin{corollary}[17.11.3]
\label{IV.17.11.3}
Soient $h:X\to S$ un morphisme lisse, $s_i$ ($1\leq i\leq r$) des sections de $\sh{O}_X$ au-dessus de $X$ (qui sont aussi des sections de $h_*(\sh{O}_X)$ au-dessus de $S$), $g:X\to S[T_1,\ldots,T_r]=\mathbf{V}_S^r$ le $S$-morphisme correspondant à l'homomorphisme de $\sh{O}_S$-Modules $\sh{O}_S^r\to h_*(\sh{O}_X)$ défini par ces sections (II, 1.2.7). Pour que $g$ soit lisse (resp. étale) en un point $x\in X$ il faut et il suffit que les $(d_{X/S}(s_i))_x$ fassent partie d'une base (resp. constituent une base) du $\sh{O}_{X,x}$-module $(\Omega^1_{X/S})_x$.
\end{corollary}

Il suffit d'appliquer (17.11.1) (resp. (17.11.2)) en prenant pour $f$ le morphisme structural $S[T_1,\ldots,T_r]\to S$.

\begin{corollary}[17.11.4]
\label{IV.17.11.4}
Pour qu'un morphisme $h:X\to S$ soit lisse en un point $x\in X$, il faut et il suffit qu'il existe un voisinage ouvert $U$ de $x$, un entier $r$, et un $S$-morphisme étale $U\to S[T_1,\ldots,T_r]$.
\end{corollary}

\oldpage[IV-4]{85}

La condition est évidemment suffisante puisque le morphisme structural $S[T_1,\ldots,T_r]\to S$ est lisse (17.3.8). Pour montrer qu'elle est nécessaire, notons que puisque $h$ est lisse au point $x$, il existe un voisinage ouvert $U$ de $x$ tel que $\Omega^1_{X/S}|U$ soit localement libre (17.2.3); il suffit alors d'utiliser (16.10.6) pour obtenir (en restreignant au besoin $U$) des sections $s_i$ de $\sh{O}_X$ au-dessus de $U$ telles que les $d_{X/S}(s_i)$ forment une base de $\Omega^1_{X/S}|U$; on conclut en appliquant (17.11.3).

\begin{proposition}[17.11.5]
\label{IV.17.11.5}
Soient $f:Y\to S$, $h:X\to S$ deux morphismes lisses. Pour qu'un $S$-morphisme $g:X\to Y$ soit une immersion ouverte, il faut et il suffit que $g$ soit un monomorphisme de préschémas et que pour tout $x\in X$, si l'on pose $y=g(x)$, on ait $\dim_y(f)=\dim_x(h)$ (pour une généralisation de cette proposition, voir (18.10.5)).
\end{proposition}

La condition est évidemment nécessaire; prouvons qu'elle est suffisante.

En vertu de (17.9.5), on est aussitôt ramené au cas où $S=\Spec(k)$ est le spectre d'un corps, et en vertu de (2.7.1, (x)), on peut supposer $k$ algébriquement clos. Compte tenu de (17.9.1), il suffit alors de prouver que $g$ est étale, et comme l'ensemble des points où $g$ est étale est ouvert, il suffit de montrer que $g$ est étale aux points \emph{fermés} (ou encore, rationnels sur $k$) de $X$ (10.4.8). Soit donc $x$ un tel point, et posons $y=g(x)$, qui est aussi rationnel sur $k$; par hypothèse, l'anneau $A=\sh{O}_{Y,y}$ est régulier et de corps résiduel $k$; soient $d=\dim(A)$ et $(t_i)_{1\leq i\leq d}$ un système régulier de paramètres pour $A$. Posons $B=\sh{O}_{X,x}$, $C=B/\mathfrak m_yB$. Comme $g$ est un monomorphisme, il en est de même du morphisme $\Spec(C)\to\Spec(k(y))=\Spec(k)$ déduit de $g$ par changement de base (I, 3.3.12); mais cela signifie que l'homomorphisme correspondant $u:k\to C$ est \emph{surjectif} (donc bijectif), car $u$ admet un inverse à gauche $v:C\to k$, et $u\circ v$ et l'identité de $C$, composés avec $u$, donnent le même morphisme $u:k\to C$. Comme par hypothèse $B$ est un anneau régulier de dimension $d$, les images des $t_i$ dans $B$ forment un système régulier de paramètres pour $B$ (0, 17.1.7); la condition (17.6.3, $e''$) est donc vérifiée par $g$ au point $x$ (0, 17.1.1 et Bourbaki, \emph{Alg. comm.}, chap. III, § 2, n\textsuperscript{o} 8, cor. 3 du th. 1), ce qui achève la démonstration.

\subsection{Sous-préschémas lisses d'un préschéma lisse. Morphismes lisses et morphismes différentiellement lisses}
\EGAReaderTarget{ega4.c04.subsec.17-12}{subsection:IV.17.12}

\begin{theorem}[17.12.1]
\label{IV.17.12.1}
Soient $f:X\to S$, $h:Y\to S$ deux morphismes localement de présentation finie, $j:Y\to X$ une immersion, $y$ un point de $Y$, $x=j(y)$. Les conditions suivantes sont équivalentes:
\begin{enumerate}
  \item[a)] $h$ est lisse au point $y$ et $f$ est lisse au point $x$.
  \item[b)] $f$ est lisse au point $x$, et l'homomorphisme canonique (16.4.21)
  \[
  (\sh{N}_{Y/X})_y\to(j^*(\Omega^1_{X/S}))_y
  \]
  est inversible à gauche.
  \item[c)] $h$ est lisse au point $y$ et il existe un voisinage ouvert $U$ de $y$ dans $Y$ tel que $j|U:U\to X$ soit une immersion régulière (16.9.2).

\oldpage[IV-4]{86}

  \item[c')] $h$ est lisse au point $y$ et il existe un voisinage ouvert $U$ de $y$ dans $Y$ tel que $j|U:U\to X$ soit une immersion quasi-régulière (autrement dit (16.9.8), il existe un voisinage $U$ de $y$ dans $Y$ dans lequel $\sh{N}_{Y/X}$ est un $\sh{O}_Y$-Module localement libre et l'homomorphisme canonique
  \[
  \mathbf{S}_{\sh{O}_Y}(\sh{N}_{Y/X})\to\operatorname{gr}^{\bullet}(j)
  \]
  est bijectif).
\end{enumerate}
\end{theorem}

Pour démontrer l'équivalence de a) et b), on peut se borner au cas où $S=\Spec(A)$ et $X=\Spec(B)$ sont affines, avec $Y=\Spec(B/\mathfrak J)$, où $\mathfrak J$ est un idéal de type fini de $B$, et $B$ est une $A$-algèbre lisse (17.3.2, (ii)). Il suffit alors d'appliquer le critère jacobien (0, 22.6.1) ainsi que (0, 19.1.14), compte tenu de ce que $\Omega^1_{X/S}$ est un $\sh{O}_X$-Module localement libre dans un voisinage de $x$ et $\sh{N}_{Y/X}$ un $\sh{O}_Y$-Module de type fini (16.1.6).

En second lieu, prouvons que a) entraîne c'). On peut encore se borner au cas où $S$, $X$, $Y$ sont affines, $B$ et $B/\mathfrak J$ des $A$-algèbres lisses, et il résulte alors de (17.7.9) et (8.10.5, (iv)) qu'il existe un sous-anneau noethérien $A'$ de $A$, une $A'$-algèbre de type fini $B'$ et un idéal $\mathfrak J'$ de $B'$ tels que $B=B'\otimes_{A'}A$, $\mathfrak J=\mathfrak J'B$ et que $B'$ et $B'/\mathfrak J'$ soient des $A'$-algèbres lisses. Notons que, d'après (0, 19.3.8), $B'$ est encore une $A'$-algèbre formellement lisse lorsqu'on munit $A'$ de la topologie discrète et $B'$ de la topologie $\mathfrak J'$-préadique. Par suite (0, 19.5.4), $\mathfrak J'/\mathfrak J'^2$ est un $(B'/\mathfrak J')$-module projectif et l'homomorphisme canonique
\[
\mathbf{S}_{B'/\mathfrak J'}(\mathfrak J'/\mathfrak J'^2)\to\operatorname{gr}_{\mathfrak J'}^{\bullet}(B')
\]
est bijectif. En d'autres termes (16.9.8), l'immersion $\Spec(B'/\mathfrak J')\to\Spec(B')$ est quasi-régulière, donc régulière puisque $B'$ est noethérien (16.9.10). Mais comme par hypothèse $B'/\mathfrak J'$ est un $A'$-module plat (17.5.1) et $B'$ une $A'$-algèbre de présentation finie, on peut appliquer (11.3.8) en remplaçant $\sh{F}$ par $\widetilde{\mathfrak J'}$, et on voit donc (par changement de base) que $j$ est une immersion régulière, et \emph{a fortiori} quasi-régulière. Le fait que $B$ est une $A$-algèbre de présentation finie et un $A$-module plat montre d'ailleurs, en vertu de (11.3.8), que les conditions c') et c) sont équivalentes, et sont stables par changement de base.

Montrons enfin que c) entraîne a). Du fait que, dans (11.3.8), la condition b) entraîne c), on voit déjà que si c) est vérifiée, $f$ est \emph{plat} au point $x$; en vertu de (17.5.1), tout revient donc à voir que sous l'hypothèse c), $f^{-1}(s)$ est lisse sur $k(s)$ au point $x$, en posant $s=h(x)$; comme on a remarqué que la condition c) est stable par changement de base, on voit qu'on est ramené au cas où $S=\Spec(k)$ est le spectre d'un corps, et compte tenu de (17.7.1, (ii)) on peut supposer que $k$ est \emph{algébriquement clos}. Il suffit alors de prouver que $\sh{O}_{X,x}$ est un anneau \emph{régulier} (17.5.1). Or, par hypothèse on a $\sh{O}_{Y,y}=\sh{O}_{X,x}/\mathfrak J_x$, où $\mathfrak J_x$ est un idéal engendré par une suite $\sh{O}_{X,x}$-régulière $(t_i)$, et $\sh{O}_{Y,y}$ est un anneau régulier; comme les $t_i$ font partie d'un système de paramètres pour $\sh{O}_{X,x}$ (0, 16.4.1), la conclusion résulte de (0, 17.1.7).

\begin{corollary}[17.12.2]
\label{IV.17.12.2}
Avec les notations de (17.12.1), supposons que $f$ soit lisse au point $x$, et soit $Y$ un sous-préschéma fermé de $X$ défini par un Idéal $\sh{J}$ de $\sh{O}_X$, et $S$-lisse au point $x$. Soient $(g_i)_{1\leq i\leq r}$ des sections de $\sh{J}$ au-dessus de $X$. Les conditions suivantes sont équivalentes:
\begin{enumerate}
  \item[a)] Les $(g_i)_x$ forment un système de générateurs du $\sh{O}_{X,x}$-module $\sh{J}_x$ dont le nombre d'éléments est le plus petit possible.
  \item[b)] Si $g_i'$ est l'image canonique de $g_i$ dans $\Gamma(X,\sh{J}/\sh{J}^2)$, les $(g_i')_x$ forment une base du $(\sh{O}_{X,x}/\sh{J}_x)$-module $\sh{J}_x/\sh{J}_x^2$.

\oldpage[IV-4]{87}

  \item[c)] Les images des $d_{X/S}g_i$ dans $(\Omega^1_{X/S})_x\otimes_{\sh{O}_{X,x}}k(x)$ sont des éléments linéairement indépendants sur $k(x)$, et les $(g_i)_x$ engendrent $\sh{J}_x$.
  \item[d)] Il existe un voisinage ouvert $U$ de $x$ dans $X$ et des sections $g_j$ ($r+1\leq j\leq n$) de $\sh{O}_X$ au-dessus de $U$ telles que le $S$-morphisme $u:U\to S[T_1,\ldots,T_n]$ correspondant à l'homomorphisme $\sh{O}_S^n\to f_*(\sh{O}_U)$ défini par les sections $g_i|U$ ($1\leq i\leq r$) et $g_j$ ($r+1\leq j\leq n$) (II, 1.2.7) soit un morphisme étale, et que l'image réciproque du sous-préschéma $Y'=S[T_{r+1},\ldots,T_n]$ par ce morphisme soit le préschéma induit par $j(Y)$ sur l'ouvert $U\cap j(Y)$.
\end{enumerate}
\end{corollary}

Comme $(g_i')_x$ est l'image canonique de $(g_i)_x$, l'équivalence de a) et b) résulte du lemme de Nakayama, $\sh{J}_x$ étant de type fini et $\sh{J}_x/\sh{J}_x^2$ un $(\sh{O}_{X,x}/\sh{J}_x)$-module libre (17.10.4) (Bourbaki, \emph{Alg. comm.}, chap. II, § 3, n\textsuperscript{o} 2, prop. 5). En vertu de (17.12.1, b)), $\sh{J}_x/\sh{J}_x^2$ s'identifie canoniquement à un facteur direct du $(\sh{O}_{X,x}/\sh{J}_x)$-module libre de rang $n$, $(\Omega^1_{X/S})_x\otimes_{\sh{O}_{X,x}}(\sh{O}_{X,x}/\sh{J}_x)$, et l'équivalence de b) et c) résulte de Bourbaki, \emph{loc. cit.} En outre, si a) est vérifiée, $(g_i')_x$ est ainsi identifié à $(d_{X/S}g_i)_x\otimes1$ pour $1\leq i\leq r$; le $\sh{O}_{X,x}$-module $(\Omega^1_{X/S})_x$ étant libre de rang $n$, il résulte encore de Bourbaki, \emph{loc. cit.}, qu'en remplaçant au besoin $X$ par un voisinage ouvert de $x$, on peut supposer qu'il existe $n-r$ sections $g_i$ ($r+1\leq i\leq n$) de $\sh{O}_X$ au-dessus de $X$, telles que les $(d_{X/S}g_i)_x$ pour $1\leq i\leq n$ forment une base de $(\Omega^1_{X/S})_x$; le fait que le morphisme correspondant $u$ soit étale au point $x$ résulte alors de (17.11.3); pour la même raison, le morphisme $u^{-1}(Y')\to Y'$, restriction de $u$, est étale au point $x$; en remplaçant $X$ par un voisinage de $x$, on peut supposer ces deux morphismes étales. En outre, il est immédiat que $Y$ (identifié, pour simplifier, au sous-préschéma fermé $j(Y)$ de $X$) est un sous-préschéma fermé de $u^{-1}(Y')$, et en vertu du choix des $g_i$ pour $i>r$ et de (17.2.5) et (17.11.2), la restriction à $Y$ de $u$ peut encore être supposée étale. On en déduit que pour tout $y'\in Y'$, l'immersion $Y\cap u^{-1}(y')\to u^{-1}(y')$ est ouverte, d'où l'on conclut à l'aide de (17.9.6) que $Y\to u^{-1}(Y')$ est une immersion ouverte; cela prouve que a) entraîne d). La réciproque découle aussitôt de (17.11.3).

\begin{corollary}[17.12.3]
\label{IV.17.12.3}
Soient $f:X\to S$ un morphisme localement de présentation finie, $u:S\to X$ une $S$-section de $X$. Pour que $f$ soit lisse en un point $x\in X$, il faut et il suffit que $u$ soit une immersion quasi-régulière dans un voisinage de $s=f(x)$.
\end{corollary}

La conclusion résulte de (17.12.1), puisque $f\circ u=1_S$ est lisse.

\begin{proposition}[17.12.4]
\label{IV.17.12.4}
Tout morphisme lisse $f:X\to S$ est différentiellement lisse (autrement dit (16.10.5) $\Delta_f:X\to X\times_SX$ est une immersion quasi-régulière). \emph{En particulier, les $\sh{P}_f^n=\sh{P}_{X/S}^n$ et les $\operatorname{gr}_n(\sh{P}_{X/S})$ sont des $\sh{O}_X$-Modules localement libres de type fini.}
\end{proposition}

Il suffit de remarquer que le morphisme structural $p_1:X\times_SX\to X$ est lisse (17.3.3, (iii)) et que $\Delta_f$ est une $X$-section pour $p_1$; la conclusion résulte de (17.12.3).

\begin{proposition}[17.12.5]
\label{IV.17.12.5}
Soit $f:X\to Y$ un morphisme localement de présentation finie. Les conditions suivantes sont équivalentes:
\begin{enumerate}
  \item[a)] $f$ est différentiellement lisse.
  \item[b)] Pour tout morphisme $g:Y'\to Y$, si l'on pose $X'=X\times_Y Y'$ et $f'=f_{(Y')}:X'\to Y'$, alors, pour toute $Y'$-section $s'$ de $X'$, $f'$ est lisse en tous les points de $s'(Y')$.
  \item[c)] La seconde projection $p_2:X\times_YX\to X$ est lisse en tous les points de la diagonale $\Delta_f(X)$.
\end{enumerate}
\end{proposition}

\oldpage[IV-4]{88}

La condition c) est un cas particulier de b): il suffit en effet de prendre $Y'=X$ et $g=f$ dans b), car alors $f'$ n'est autre que la seconde projection $p_2$, et $\Delta_f$ est une $X$-section de $X\times_YX$. En second lieu, c) entraîne a), car $p_2$ est un morphisme localement de présentation finie, et si $p_2$ est lisse aux points de $\Delta_f(X)$, $\Delta_f$ est une immersion quasi-régulière (17.12.3), donc $f$ est différentiellement lisse. Enfin, montrons que a) entraîne b); si $\Delta_f$ est une immersion quasi-régulière, $p_2$ est lisse en tous les points de $\Delta_f(X)$ (17.12.3).

Or, si $g':X'\to X$ est la projection canonique, et $v=(g',g')_Y:X'\to X\times_YX$, on vérifie aussitôt que le diagramme
\[
\xymatrix{
  X\times_Y X \ar[d]_{p_2} & X' \ar[l]_-v \ar[d]^{f'}\\
  X & Y' \ar[l]^{h=g'\circ s'}
}
\]
est commutatif et identifie $X'$ au produit $(X\times_YX)\times_XY'$ (I, 3.3.9); tenant compte de ce que le diagramme (17.4.1.1) identifie $Y'$ au produit des $(X\times_YX)$-préschémas $X$ et $X'$, on conclut, de ce que $p_2$ est lisse en tout point de $\Delta_f(X)$, que $f'$ est lisse en tout point de $s'(Y')$ (17.3.3, (iii)).

\begin{corollary}[17.12.6]
\label{IV.17.12.6}
Les notations étant celles de (8.8.1), on suppose $X_\alpha$ quasi-compact, et $f_\alpha:X_\alpha\to S_\alpha$ localement de présentation finie. Pour que $f:X\to S$ soit différentiellement lisse, il faut et il suffit qu'il existe $\lambda\geq\alpha$ tel que $f_\lambda:X_\lambda\to S_\lambda$ soit différentiellement lisse (auquel cas $f_\mu:X_\mu\to S_\mu$ est différentiellement lisse pour $\mu\geq\lambda$).
\end{corollary}

La suffisance de la condition et la dernière assertion résultent de (16.10.4). Pour prouver que la condition est nécessaire, notons que $\Delta_f:X\to X\times_SX$ et $p_2:X\times_SX\to X$ se déduisent par changement de base de $\Delta_{f_\lambda}:X_\lambda\to X_\lambda\times_{S_\lambda}X_\lambda$ et $p_{2,\lambda}:X_\lambda\times_{S_\lambda}X_\lambda\to X_\lambda$. En vertu de (17.12.5), pour tout $z\in\Delta_f(X)$, il existe un voisinage ouvert affine $U(z)$ de $z$ dans $X\times_SX$ tel que $p_2$ soit lisse dans $U(z)$; en vertu de (8.2.11) et (17.7.8), il existe un indice $\lambda(z)$ et un voisinage $U_{\lambda(z)}$ de la projection de $z$ dans $X_{\lambda(z)}\times_{S_{\lambda(z)}}X_{\lambda(z)}$ tels que $U(z)$ soit l'image réciproque de $U_{\lambda(z)}$, et que $p_{2,\lambda(z)}$ soit lisse dans $U_{\lambda(z)}$. Comme $X$ est quasi-compact, on peut recouvrir $\Delta_f(X)$ par un nombre fini de voisinages $U(z_i)$; en vertu de (8.3.4), il existe un $\lambda$ supérieur à tous les $\lambda(z_i)$ tels que les images réciproques des $U_{\lambda(z_i)}$ dans $X_\lambda\times_{S_\lambda}X_\lambda$ forment un recouvrement de $\Delta_{f_\lambda}(X_\lambda)$; il résulte alors de (17.12.5) que cet indice $\lambda$ répond à la question.

\begin{examples}[17.12.7]
\label{IV.17.12.7}
Soient $S$ un préschéma, $G$ un $S$-\emph{préschéma en groupes}, localement de présentation finie sur $S$. Alors, pour que $G$ soit différentiellement lisse, il faut et il suffit que $G$ soit lisse sur $S$ aux points de la « section unité » $e$ de $G$. La condition est en effet nécessaire par (17.12.5). Inversement, supposons-la remplie; pour tout morphisme $S'\to S$, $G'=G\times_SS'$ est alors un $S'$-préschéma en groupes localement de présentation finie sur $S'$; on peut donc, en vertu de (17.12.5), se borner à prouver que pour toute $S$-section $s$ de $G$, $f$ est lisse aux points de $s(S)$. Or, si $m:G\times_SG\to G$ est le morphisme qui définit la structure de préschéma en groupes de $G$, $m\circ(s\times1_G):S\times_SG\to G\times_SG\to G$ est un

\oldpage[IV-4]{89}

$S$-isomorphisme du préschéma $G$, la « translation à gauche » $\tau_s$, transformant les points de la section unité de $G$ en les points de $s(S)$; on en conclut que $G$ est lisse sur $S$ aux points de $s(S)$.
\end{examples}

\begin{remark}[17.12.8]
\label{IV.17.12.8}
Un $S$-préschéma en groupes $G$, de type fini (et même fini) sur un préschéma localement noethérien $S$, peut être différentiellement lisse sur $S$ sans être lisse (ni même \emph{plat}) sur $S$. Prenons par exemple pour $S$ le spectre de l'algèbre des nombres duaux $D=k[T]/T^2k[T]$ sur un corps $k$, pour $G$ le spectre de la $D$-algèbre composée directe $E=D\oplus k$. On définit sur $G$ une structure de $S$-schéma en groupes en définissant une « application diagonale » $\delta$, homomorphisme de $E$ dans $E\otimes_DE$: si $e'$, $e''$ sont les idempotents, images canoniques de l'unité $1$ de $D$ dans les facteurs $D$ et $k$ de $E$, on vérifie sans peine qu'on obtient une telle application diagonale en prenant
\[
\delta(e')=e'\otimes e'+e''\otimes e'',\qquad
\delta(e'')=e'\otimes e''+e''\otimes e'.
\]
La section unité du schéma en groupes $G$ correspond à l'homomorphisme $E\to D$ qui est l'identité dans $D$ et $0$ dans $k$; c'est un isomorphisme de $S$ sur une composante connexe de $G$, et \emph{a fortiori} $G$ est différentiellement lisse sur $S$ (17.12.6), mais il est clair que $G$ n'est pas $S$-plat.
\end{remark}

\subsection{Morphismes transversaux.}
\EGAReaderTarget{ega4.c04.subsec.17-13}{subsection:IV.17.13}
\label{IV.17.13}

\begin{env}[17.13.1]
\label{IV.17.13.1}
Soient $S$ un préschéma, $X$, $Y$, $X'$ trois $S$-préschémas, $i:Y\to X$ une $S$-immersion, $f:X'\to X$ un $S$-morphisme. Posons $Y'=Y\times_XX'$, et soient $g:Y'\to Y$, $j:Y'\to X'$ les projections canoniques, de sorte qu'on a le diagramme commutatif
\[
\label{IV.17.13.1.1}
\tag{17.13.1.1}
\xymatrix{
  Y \ar[d]_i & Y' \ar[l]_g \ar[d]^j\\
  X & X' \ar[l]^f
}
\]
et que $j$ est une $S$-immersion. On a alors un diagramme commutatif de $\sh{O}_{Y'}$-Modules quasi-cohérents
\[
\label{IV.17.13.1.2}
\tag{17.13.1.2}
\begin{tikzcd}
  g^*(\sh{N}_{Y/X}) \ar{r} \ar{d}{\operatorname{gr}_1(g)} & f^*(\Omega^1_{X/S})\otimes_{\sh{O}_{X'}}\sh{O}_{Y'} \ar{r} \ar{d} & g^*(\Omega^1_{Y/S}) \ar{r} \ar{d} & 0\\
  \sh{N}_{Y'/X'} \ar{r} & \Omega^1_{X'/S}\otimes_{\sh{O}_{X'}}\sh{O}_{Y'} \ar{r} & \Omega^1_{Y'/S} \ar{r} & 0
\end{tikzcd}
\]
où la ligne inférieure est la suite exacte (16.4.21) appliquée à $j$, la ligne supérieure provient de la même suite exacte pour $i$, par application du foncteur exact à droite $g^*$ (qui la laisse donc exacte); $\operatorname{gr}_1(g)$ est défini en (16.2.1), et la commutativité des deux carrés résulte de (0, 20.5.7.3) et (0, 20.5.11.3).

\oldpage[IV-4]{90}

On a vu en outre (16.2.2, (iii)) que $\operatorname{gr}_1(g)$ est ici surjectif, donc on déduit de (17.13.1.2) la suite exacte
\[
\label{IV.17.13.1.3}
\tag{17.13.1.3}
g^*(\sh{N}_{Y/X})\xrightarrow{\alpha}\Omega^1_{X'/S}\otimes_{\sh{O}_{X'}}\sh{O}_{Y'}\to\Omega^1_{Y'/S}\to0.
\]
\end{env}

\begin{proposition}[17.13.2]
\label{IV.17.13.2}
Les notations étant celles de (17.13.1), soit $x'$ un point de $Y'$, $x=g(x')$ son image dans $Y$; supposons $X$ et $Y$ lisses sur $S$ au point $x$, $X'$ lisse sur $S$ au point $x'$. Soient $m$, $m-c$ les dimensions relatives de $X$ et $Y$ sur $S$ au point $x$ (17.10.1), $n$ celle de $X'$ sur $S$ au point $x'$. Les conditions suivantes sont équivalentes:
\begin{enumerate}
  \item[a)] $Y'$ est lisse sur $S$ au point $x'$ et de dimension relative $n-c$.
  \item[b)] L'homomorphisme $\alpha\otimes1:g^*(\sh{N}_{Y/X})\otimes_{\sh{O}_{Y'}}k(x')\to\Omega^1_{X'/S}\otimes_{\sh{O}_{X'}}k(x')$ déduit de l'homomorphisme $\alpha$ de (17.13.1.3) est injectif.
\end{enumerate}
Soit $s$ l'image canonique de $x$ dans $S$, et supposons que $x'$ soit rationnel sur $k(s)$. Alors les conditions a) et b) sont aussi équivalentes à la suivante:
\begin{enumerate}
  \item[b')] L'homomorphisme transposé de $\alpha\otimes1$
  \[
  \mathrm{T}_{X'/S}(x')\to(\sh{N}_{Y/X}\otimes_{\sh{O}_Y}k(x'))^\vee=\mathrm{T}_{X/S}(x)/\mathrm{T}_{Y/S}(x)
  \]
  (cf. (16.5.12)) est surjectif.
\end{enumerate}
De plus, lorsque les conditions a) et b) sont vérifiées en $x'$, elles le sont dans un voisinage de $x'$ dans $Y'$; en remplaçant au besoin $X'$ par un voisinage de $x'$, l'homomorphisme
\[
\operatorname{gr}_1(g):g^*(\sh{N}_{Y/X})\to\sh{N}_{Y'/X'}
\]
est bijectif, et la suite
\[
\label{IV.17.13.2.1}
\tag{17.13.2.1}
0\to g^*(\sh{N}_{Y/X})\xrightarrow{\alpha}\Omega^1_{X'/S}\otimes_{\sh{O}_{X'}}\sh{O}_{Y'}\to\Omega^1_{Y'/S}\to0
\]
obtenue en ajoutant un $0$ à (17.13.1.3), est exacte.
\end{proposition}

Le fait que si les conditions équivalentes a) et b) sont vérifiées au point $x'$, elles le sont aussi au voisinage de $x'$ dans $Y'$, résulte de ce que l'ensemble des points où un morphisme est lisse est ouvert (17.3.7) et de (17.10.2).

En vertu de (17.12.1) appliqué à $X'$ et $Y'$, dire que $Y'$ est lisse sur $S$ au point $x'$ équivaut à dire que l'homomorphisme
\[
\delta\otimes1:\sh{N}_{Y'/X'}\otimes_{\sh{O}_{Y'}}k(x')\to\Omega^1_{X'/S}\otimes_{\sh{O}_{X'}}k(x')
\]
est injectif, compte tenu de (0, 19.1.12) et de ce que $\Omega^1_{X'/S}$ est un $\sh{O}_{X'}$-Module localement libre au point $x'$ (17.2.3); comme alors $\Omega^1_{Y'/S}$ est aussi un $\sh{O}_{Y'}$-Module localement libre au voisinage de $x'$ et que la suite
\[
\label{IV.17.13.2.2}
\tag{17.13.2.2}
0\to\sh{N}_{Y'/X'}\to\Omega^1_{X'/S}\otimes_{\sh{O}_{X'}}\sh{O}_{Y'}\to\Omega^1_{Y'/S}\to0
\]
est exacte (17.2.5), dire que $Y'$ est de dimension relative $n-c$ sur $S$ au point $x'$ signifie, par (17.10.2), que $\sh{N}_{Y'/X'}$ (qui est localement libre au voisinage de $x'$) est de rang $c$ au point $x'$. Mais en vertu des hypothèses sur $X$ et $Y$ au point $x$, et par le même raisonnement, $\sh{N}_{Y/X}$ est localement libre et de rang $c$ au point $x$, donc $g^*(\sh{N}_{Y/X})$ est aussi localement libre et de rang $c$ au point $x'$. Comme l'homomorphisme $\operatorname{gr}_1(g):g^*(\sh{N}_{Y/X})\to\sh{N}_{Y'/X'}$ est surjectif, les conditions précédentes équivalent à dire que cet homomorphisme est

\oldpage[IV-4]{91}

bijectif au point $x'$ (Bourbaki, \emph{Alg. comm.}, chap. II, § 3, n\textsuperscript{o} 2, cor. de la prop. 6), donc aussi dans un voisinage de $x'$ (0$_{\mathrm I}$, 5.2.7); cela entraîne évidemment b), ainsi que la dernière assertion de l'énoncé, en vertu de l'exactitude de (17.13.2.2). Inversement, comme $\alpha\otimes1$ se factorise en
\[
g^*(\sh{N}_{Y/X})\otimes k(x')\xrightarrow{\operatorname{gr}_1(g)\otimes1}\sh{N}_{Y'/X'}\otimes k(x')\xrightarrow{\delta\otimes1}\Omega^1_{X'/S}\otimes k(x')
\]
et que $\operatorname{gr}_1(g)$ est surjectif, dire que $\alpha\otimes1$ est injectif entraîne que $\delta\otimes1$ l'est et que $\operatorname{gr}_1(g)\otimes1$ est bijectif. On en conclut (17.12.1) que $Y'$ est lisse sur $S$ au point $x'$, et que $\operatorname{gr}_1(g)$ est bijectif dans un voisinage de $x'$ (Bourbaki, \emph{loc. cit.}). En outre, comme la suite (17.13.2.2) est alors exacte, et que $\sh{N}_{Y'/X'}$, isomorphe à $g^*(\sh{N}_{Y/X})$, est de rang $c$ au point $x'$, $\Omega^1_{Y'/S}$ est de rang $n-c$ en ce point, ce qui achève de prouver l'équivalence de a) et b), en vertu de (17.10.2).

Reste à montrer l'équivalence de b) et b') lorsque $x'$ est rationnel sur $k(s)$ (ce qui implique qu'il en est de même de $x$); alors $g^*(\sh{N}_{Y/X})\otimes_{\sh{O}_{Y'}}k(x')$ s'identifie à $\sh{N}_{Y/X}\otimes_{\sh{O}_Y}k(x')$, et puisque la suite
\[
0\to\sh{N}_{Y/X}\to\Omega^1_{X/S}\otimes_{\sh{O}_X}\sh{O}_Y\to\Omega^1_{Y/S}\to0
\]
est exacte au point $x$ et formée de $\sh{O}_Y$-Modules localement libres, le dual de $\sh{N}_{Y/X}\otimes_{\sh{O}_Y}k(x)$ s'identifie à l'espace quotient $\mathrm{T}_{X/S}(x)/\mathrm{T}_{Y/S}(x)$ (16.5.12); d'où l'équivalence de b) et b').

\begin{definition}[17.13.3]
\EGAReaderTarget{ega4.c04.def.17-13-3}{IV.17.13.3}
Avec les notations de (17.13.1), on dit que le morphisme $f$ est \emph{transversal} à $Y$ au point $x'$, relativement à $S$, si $X$ et $Y$ sont lisses sur $S$ au point $x$, $X'$ lisse sur $S$ au point $x'$, et si les conditions équivalentes a), b) de (17.13.2) sont satisfaites.
\end{definition}

Lorsque aucune confusion n'est à craindre, on supprime la mention du préschéma $S$ et on dit simplement que $f$ est transversal à $Y$ au point $x'$.

\begin{remarks}[17.13.4]
\label{IV.17.13.4}
(i) Supposons que $X$, $Y$ et $X'$ soient \emph{plats et localement de présentation finie} sur $S$. Pour tout $s\in S$, notons $X_s$, $Y_s$, $X'_s$, $Y'_s$ les fibres de $X$, $Y$, $X'$, $Y'$ au point $s$, $f_s:X'_s\to X_s$ le morphisme déduit de $f$ par changement de base. Il résulte alors de (17.5.9) que pour que $f$ soit un morphisme transversal à $Y$ au point $x'$, il faut et il suffit que, si $s$ est l'image de $x$ dans $S$, $f_s$ soit transversal à $Y_s$ au point $x'$, relativement à $k(s)$.

(ii) On a vu dans (17.13.2) que l'ensemble des points $x'\in Y'$ où $f$ est transversal à $Y$ est \emph{ouvert} dans $Y'$; si $Y'$ est de plus \emph{propre} sur $S$, on en déduit que l'ensemble des $s\in S$ tels que $f$ soit transversal à $Y$ en tous les points de $Y'_s$, relativement à $S$, est \emph{ouvert} dans $S$.

Lorsque $X$, $Y$ et $X'$ sont plats et localement de présentation finie sur $S$, il résulte de (i) que l'ensemble des $x'\in Y'$ tels que $f_s$ soit transversal à $Y_s$ au point $x'$ ($s$ image de $x'$ dans $S$), relativement à $k(s)$, est \emph{ouvert} dans $Y'$. Si de plus $Y'$ est propre sur $S$, l'ensemble des $s\in S$ tels que $f_s$ soit transversal à $Y_s$ en tous les points de $Y'_s$ (relativement à $k(s)$) est \emph{ouvert} dans $S$.

(iii) La propriété d'être lisse sur $S$, ainsi que la notion de dimension relative à $S$ en un point, étant stables par tout changement de base $S'\to S$ ((17.3.3) et (4.2.7)), il en est de même de la propriété pour un morphisme $f:X'\to X$ d'être transversal à un sous-préschéma de $X$ en un point, relativement à $S$.

\oldpage[IV-4]{92}

(iv) La condition b) de (17.13.2) exprime encore que l'homomorphisme $\alpha:g^*(\sh{N}_{Y/X})\to\Omega^1_{X'/S}\otimes_{\sh{O}_{X'}}\sh{O}_{Y'}$ est \emph{universellement injectif} relativement à $Y'$ (II, 9.18).
\end{remarks}

\begin{env}[17.13.5]
\label{IV.17.13.5}
Considérons maintenant un préschéma $S$, trois $S$-préschémas $X$, $Y$, $Z$, deux $S$-morphismes $f:Y\to X$, $g:Z\to X$; posons $T=Y\times_XZ$; on sait alors (I, 5.3.5) qu'on a un diagramme commutatif
\[
\label{IV.17.13.5.1}
\tag{17.13.5.1}
\xymatrix{
  X \ar[d]^{\Delta} & T=Y\times_X Z \ar[l] \ar[d]^{j}\\
  X\times_S X & Y\times_S Z \ar[l]^{u}
}
\]
faisant de $T$ le produit des $(X\times_SX)$-préschémas $X$ et $Y\times_SZ$, où $u=f\times_Sg$. Comme $\Delta$ est une $S$-immersion (I, 5.3.9), on est dans la situation du diagramme (17.13.1.1); ce qui correspond à $\Omega^1_{X'/S}$ dans (17.13.1) est alors $(\Omega^1_{Y/S}\otimes_{\sh{O}_Y}\sh{O}_{Y\times_SZ})\oplus(\Omega^1_{Z/S}\otimes_{\sh{O}_Z}\sh{O}_{Y\times_SZ})$ en vertu de (16.4.23). D'autre part, ce qui correspond au $\sh{O}_Y$-Module $\sh{N}_{Y/X}$ dans (17.13.1) est ici par définition $\Omega^1_{X/S}$ (16.3.1); il correspond donc à (17.13.1.3) une suite exacte
\[
\label{IV.17.13.5.2}
\tag{17.13.5.2}
\Omega^1_{X/S}\otimes_{\sh{O}_X}\sh{O}_T\xrightarrow{\rho}
(\Omega^1_{Y/S}\otimes_{\sh{O}_Y}\sh{O}_T)\oplus
(\Omega^1_{Z/S}\otimes_{\sh{O}_Z}\sh{O}_T)
\xrightarrow{\sigma}\Omega^1_{T/S}\to0
\]
où il reste à préciser les homomorphismes $\rho$ et $\sigma$. Tenant compte d'abord de (16.4.23) et de (0, 20.5.2), on voit que si $p:T\to Y$, $q:T\to Z$ sont les projections canoniques, on a, avec les notations de (16.4.19),
\[
\label{IV.17.13.5.3}
\tag{17.13.5.3}
\sigma=p_{T/Y/S}+q_{T/Z/S}.
\]
D'autre part, pour évaluer $\rho$, utilisons la commutativité du carré de gauche dans (17.13.1.2), ce qui, dans le cas présent, ramène d'abord à expliciter l'homomorphisme canonique
\[
\rho':\Omega^1_{X/S}\to\Omega^1_{(X\times_SX)/S}\otimes_{\sh{O}_{X\times_SX}}\sh{O}_X
\]
défini dans (16.4.21) appliqué à l'immersion $\Delta$. On peut se borner au cas où $S=\Spec(A)$, $X=\Spec(B)$ sont affines, et alors $\rho'$ correspond à l'homomorphisme $\delta$ de (0, 20.5.11.2), où il faut remplacer $B$ par $B\otimes_AB$ et $C$ par $B=(B\otimes_AB)/\mathfrak J_{B/A}$. On voit alors que $\delta$ fait correspondre à la classe de $x\otimes1-1\otimes x\pmod{\mathfrak J_{B/A}^2}$ (pour un $x\in B$) l'image de
\[
(x\otimes1-1\otimes x)\otimes(1\otimes1)-(1\otimes1)\otimes(x\otimes1-1\otimes x)
\]
dans $\Omega^1_{(B\otimes_AB)/A}\otimes_{B\otimes_AB}B$; mais l'élément précédent peut s'écrire
\[
\begin{split}
&((x\otimes1)\otimes(1\otimes1)-(1\otimes1)\otimes(x\otimes1))\\
&\quad-((1\otimes x)\otimes(1\otimes1)-(1\otimes1)\otimes(1\otimes x))
\end{split}
\]
et on voit donc que $\rho'$ est la différence des deux homomorphismes $\pi^{(1)}$ et $\pi^{(2)}$ de $\Omega^1_{X/S}$ dans $\Omega^1_{(X\times_SX)/S}\otimes_{\sh{O}_{X\times_SX}}\sh{O}_X$, correspondant respectivement à la première et à la seconde projection de $X\times_SX$ dans $X$, par (16.4.3.3). Pour obtenir $\rho$, il faut d'abord considérer l'homomorphisme
\[
\rho'':\Omega^1_{(X\times_SX)/S}\otimes_{\sh{O}_{X\times_SX}}\sh{O}_{Y\times_SZ}\to\Omega^1_{(Y\times_SZ)/S}
\]
correspondant par (16.4.3.3) au morphisme $u$ de (17.13.5.1), puis, après tensorisation par $\sh{O}_T$, former le composé $(\rho''\otimes1)\circ(\rho'\otimes1)$; il résulte de ce qui précède que l'on a, avec les notations de (16.4.18),
\[
\label{IV.17.13.5.4}
\tag{17.13.5.4}
\rho=(f_{Y/X/S}\otimes1_{\sh{O}_T})\oplus(-g_{Z/X/S}\otimes1_{\sh{O}_T}).
\]

\oldpage[IV-4]{93}

Ceci dit, l'application de (17.13.2) à la situation du diagramme (17.13.5.1) (compte tenu de (17.3.3, (iv)), qui implique que $X\times_SX$ est lisse sur $S$ en $x$ si $X$ l'est) donne le corollaire suivant.
\end{env}

\begin{corollary}[17.13.6]
\EGAReaderTarget{ega4.c04.cor.17-13-6}{IV.17.13.6}
Soient $S$ un préschéma, $X$, $Y$, $Z$ trois $S$-préschémas, $f:Y\to X$, $g:Z\to X$ deux $S$-morphismes; posons $T=Y\times_XZ$, et soient $p:T\to Y$, $q:T\to Z$ les projections canoniques. Soient $t$ un point de $T$, $y=p(t)$, $z=q(t)$, $x=f(y)=g(z)$. Supposons que $X$ soit lisse sur $S$ au point $x$, de dimension relative $m$, $Y$ (resp. $Z$) lisse sur $S$ au point $y$ (resp. $z$), de dimension relative $m+a$ (resp. $m+b$), $a$ et $b$ étant positifs ou négatifs. Alors les conditions suivantes sont équivalentes:
\begin{enumerate}
  \item[a)] $T$ est lisse sur $S$ au point $t$, de dimension relative $m+a+b$.
  \item[b)] L'homomorphisme
  \[
  \rho\otimes1:\Omega^1_{X/S}\otimes_{\sh{O}_X}k(t)\to
  (\Omega^1_{Y/S}\otimes_{\sh{O}_Y}k(t))\oplus
  (\Omega^1_{Z/S}\otimes_{\sh{O}_Z}k(t))
  \]
  où $\rho$ est donné par (17.13.5.4), est injectif.
  \item[c)] Le morphisme $T=Y\times_SZ\to X\times_SX$ est transversal à la diagonale de $X\times_SX$ au point $t$.
\end{enumerate}
Si $t$ est rationnel sur le corps résiduel de son image $s$ dans $S$, ces conditions sont aussi équivalentes à la suivante:
\begin{enumerate}
  \item[b')] L'homomorphisme
  \[
  \label{IV.17.13.6.1}
  \tag{17.13.6.1}
  \mathrm{T}_y(f)-\mathrm{T}_z(g):\mathrm{T}_{Y/S}(y)\oplus\mathrm{T}_{Z/S}(z)\to\mathrm{T}_{X/S}(x)
  \]
  (cf. (16.5.12.5)) est surjectif.
\end{enumerate}
De plus, lorsque les conditions a) et b) sont vérifiées en $t$, elles le sont dans un voisinage de $t$ dans $T$, et en restreignant $T$ à un tel voisinage, la suite
\[
\label{IV.17.13.6.2}
\tag{17.13.6.2}
0\to\Omega^1_{X/S}\otimes_{\sh{O}_X}\sh{O}_T\xrightarrow{\rho}
(\Omega^1_{Y/S}\otimes_{\sh{O}_Y}\sh{O}_T)\oplus
(\Omega^1_{Z/S}\otimes_{\sh{O}_Z}\sh{O}_T)
\xrightarrow{\sigma}\Omega^1_{T/S}\to0
\]
(où $\sigma$ est donné par (17.13.5.3)) est exacte.

Le seul point qui reste à prouver dans (17.13.6) concerne b'), et résulte de ce que, sous l'hypothèse que $t$ est rationnel sur $k(s)$, l'homomorphisme $\rho\otimes1$ est \emph{transposé} de $\mathrm{T}_y(f)-\mathrm{T}_z(g)$, en vertu de (17.13.5.4).

Lorsque les conditions équivalentes de (17.13.6) sont satisfaites, on dit que $f$ et $g$ forment un \emph{couple de morphismes transversaux au point $t$, relativement} à $S$; ici encore, on supprime souvent la mention de $S$.
\end{corollary}

\begin{env}[17.13.7]
\EGAReaderTarget{ega4.c04.statement.17-13-7}{IV.17.13.7}
Considérons en particulier le cas où $Y$ et $Z$ sont des \emph{sous-préschémas} de $X$, $f$ et $g$ étant les injections canoniques; $T$ est alors le sous-préschéma « intersection »

\oldpage[IV-4]{94}

$\inf(Y,Z)$ de $X$ (I, 4.4.3), et l'on a $t=y=z=x$. Au lieu de dire que le couple $(f,g)$ est transversal au point $x$, on dit alors que $Y$ et $Z$ \emph{se coupent transversalement au point $x$} (relativement à $S$). Notant $X_s$, $Y_s$, $Z_s$, $T_s$ les fibres de $X$, $Y$, $Z$, $T$ au point $s$, image de $x$ dans $S$, et tenant compte de (5.2.3), (5.1.9) et (0, 16.5.12), on voit que pour qu'il en soit ainsi (lorsque $X$, $Y$, $Z$ sont lisses sur $S$ au point $x$), il faut et il suffit que $T$ soit lisse sur $S$ au point $x$, et que l'on ait la relation
\[
\label{IV.17.13.7.1}
\tag{17.13.7.1}
\operatorname{codim}_x(T_s,X_s)=\operatorname{codim}_x(Y_s,X_s)+\operatorname{codim}_x(Z_s,X_s).
\]
\end{env}

\begin{proposition}[17.13.8]
\label{IV.17.13.8}
Soient $S$ un préschéma, $X$ un $S$-préschéma localement de présentation finie sur $S$, $Y$, $Z$ deux sous-préschémas de $X$, $T=\inf(Y,Z)$ le sous-préschéma « intersection » de $Y$ et $Z$, $x$ un point de $T$. Les conditions suivantes sont équivalentes:
\begin{enumerate}
  \item[a)] L'injection canonique $i:Y\to X$ est un morphisme transversal à $Z$ au point $x$, relativement à $S$ (17.13.3).
  \item[a')] L'injection canonique $j:Z\to X$ est un morphisme transversal à $Y$ au point $x$, relativement à $S$.
  \item[a'')] $Y$ et $Z$ se coupent transversalement au point $x$, relativement à $S$.
  \item[b)] $X$, $Y$, $Z$ sont lisses sur $S$ au point $x$, et l'homomorphisme $\rho$ (17.13.5.4) est tel que
  \[
  \rho\otimes1:\Omega^1_{X/S}\otimes_{\sh{O}_X}k(x)\to
  (\Omega^1_{Y/S}\otimes_{\sh{O}_Y}k(x))\oplus
  (\Omega^1_{Z/S}\otimes_{\sh{O}_Z}k(x))
  \]
  est injectif.
\end{enumerate}
Lorsque $x$ est rationnel sur $k(s)$ ($s$ image de $x$ dans $S$) ces conditions sont aussi équivalentes à la suivante:
\begin{enumerate}
  \item[b')] L'homomorphisme
  \[
  \mathrm{T}_x(i)-\mathrm{T}_x(j):\mathrm{T}_{Y/S}(x)\oplus\mathrm{T}_{Z/S}(x)\to\mathrm{T}_{X/S}(x)
  \]
  est surjectif.
\end{enumerate}
En outre, lorsque les conditions équivalentes a) à b) sont vérifiées au point $x$, elles le sont dans un voisinage de $x$ dans $T$, et en restreignant $X$ à un voisinage de $x$, la suite (17.13.6.2) est exacte, et l'on a un isomorphisme canonique
\[
\label{IV.17.13.8.1}
\tag{17.13.8.1}
\sh{N}_{T/X}\xrightarrow{\sim}
(\sh{N}_{Y/X}\otimes_{\sh{O}_Y}\sh{O}_T)\oplus
(\sh{N}_{Z/X}\otimes_{\sh{O}_Z}\sh{O}_T).
\]
\end{proposition}

Les conditions a), a'), a'') impliquent toutes que $X$, $Y$, $Z$ sont lisses sur $S$ au point $x$. Soient en outre, $m$, $m-a$, $m-b$ les dimensions relatives de $X$, $Y$, $Z$ sur $S$ au point $x$. Il résulte alors de (17.13.2) appliqué en remplaçant $X'$ par $Z$ et $Y'$ par $T=Y\times_XZ$, que les conditions a), a') équivalent toutes deux à dire que $T$ est lisse sur $S$ au point $x$ et de dimension relative $m-a-b$ en ce point; mais en vertu de (17.13.7), cela signifie précisément que la condition a'') est vérifiée, d'où l'équivalence de a), a') et a''). L'équivalence de a'') et de b) (ou b') lorsque $x$ est rationnel sur $k(s)$) a été prouvée dans (17.13.6), ainsi que le fait que si ces conditions sont satisfaites au point $x$, elles le sont dans un voisinage de $x$, et l'exactitude de la suite (17.13.6.2) dans un tel voisinage. Reste à définir l'isomorphisme canonique (17.13.8.1).

\oldpage[IV-4]{95}

Notons $\alpha$ et $-\beta$ les homomorphismes figurant au second membre de (17.13.5.4), $\gamma$ et $\delta$ ceux figurant au second membre de (17.13.5.2). Dire que la suite (17.13.6.2)
\[
0\to\Omega^1_{X/S}\otimes_{\sh{O}_X}\sh{O}_T
\xrightarrow{\alpha\oplus(-\beta)}
(\Omega^1_{Y/S}\otimes_{\sh{O}_Y}\sh{O}_T)\oplus
(\Omega^1_{Z/S}\otimes_{\sh{O}_Z}\sh{O}_T)
\xrightarrow{\gamma+\delta}\Omega^1_{T/S}\to0
\]
est exacte signifie que, dans la catégorie des $\sh{O}_T$-Modules, $\Omega^1_{X/S}\otimes_{\sh{O}_X}\sh{O}_T$ s'identifie canoniquement au \emph{produit fibré} de $\Omega^1_{Y/S}\otimes_{\sh{O}_Y}\sh{O}_T$ et $\Omega^1_{Z/S}\otimes_{\sh{O}_Z}\sh{O}_T$ sur $\Omega^1_{T/S}$, pour les homomorphismes $\gamma$ et $\delta$. Le même raisonnement que dans (0, 18.1.2 et 18.1.3), où l'on remplace anneaux et idéaux bilatères respectivement par Modules et sous-Modules, fournit un diagramme commutatif
{\scriptsize
\[
\begin{tikzcd}[column sep=small,row sep=large]
0 \ar[dr] & & 0 \ar[d] & 0 \ar[d] & \\
& (\sh{N}_{Y/X}\otimes\sh{O}_T)\oplus(\sh{N}_{Z/X}\otimes\sh{O}_T)
  \ar[r] \ar[d] \ar[dr,"\iota" description]
& \sh{N}_{Z/X}\otimes\sh{O}_T \ar[r,"\sim"] \ar[d]
& \sh{N}_{T/Y} \ar[d] & \\
0 \ar[r] & \sh{N}_{Y/X}\otimes\sh{O}_T \ar[r] \ar[d]
& \Omega^1_{X/S}\otimes\sh{O}_T \ar[r,"\alpha"] \ar[d,"\beta"] \ar[dr,"\varepsilon" description]
& \Omega^1_{Y/S}\otimes\sh{O}_T \ar[r] \ar[d,"\gamma"] & 0\\
0 \ar[r] & \sh{N}_{T/Z} \ar[r]
& \Omega^1_{Z/S}\otimes\sh{O}_T \ar[r,"\delta"] \ar[d]
& \Omega^1_{T/S} \ar[r] \ar[d] \ar[dr] & 0\\
& & 0 & 0 & 0
\end{tikzcd}
\]
}
où les 3\textsuperscript{e} et 4\textsuperscript{e} lignes et colonnes sont exactes en vertu des hypothèses de lissité. Le fait que l'homomorphisme composé $\varepsilon=\gamma\circ\alpha=\delta\circ\beta$ soit l'homomorphisme correspondant à l'injection canonique $T\to X$ résulte de (0, 20.5.2.7). Comme la diagonale du diagramme précédent est exacte, $\operatorname{Ker}(\varepsilon)$ s'identifie canoniquement avec $\sh{N}_{T/X}$, compte tenu de ce que $T$ est lisse sur $S$ (17.2.5); donc on a un isomorphisme canonique (17.13.8.1) réciproque de $\iota$.

\begin{env}[17.13.9]
\label{IV.17.13.9}
On peut généraliser les résultats de (17.13.5) et (17.13.6) à un nombre fini quelconque de $S$-morphismes $f_i:Y_i\to X$ ($i\in I$, $I$ ensemble fini). Pour cela, notons $X^I$ le produit de la famille $(X_i)_{i\in I}$ de $S$-préschémas tous identiques à $X$ (I, 3.3.5), et soient $p_i$ ($i\in I$) les projections canoniques. Le \emph{morphisme diagonal} $\Delta:X\to X^I$ est défini (comme dans (I, 5.3.1)) comme l'unique $S$-morphisme tel que $p_i\circ\Delta=1_X$ pour tout $i\in I$. C'est une $X$-section de $X^I$ pour le morphisme $p_1$, donc (I, 5.3.11) une \emph{immersion}.

\oldpage[IV-4]{96}

Soit alors $Y$ le produit des $S$-préschémas $Y_i$ (pour les morphismes composés $Y_i\xrightarrow{f_i}X\to S$), $T$ le produit des $X$-préschémas $Y_i$ (pour les morphismes $f_i$). On prouve comme dans (I, 5.3.5) que l'on a un diagramme commutatif
\[
\label{IV.17.13.9.1}
\tag{17.13.9.1}
\xymatrix{
  X \ar[d]^{\Delta} & T \ar[l]_-v \ar[d]^{j}\\
  X^I & Y \ar[l]^{u}
}
\]
(où $u$ est le produit (sur $S$) des $f_i$), qui fait de $T$ le produit des $X^I$-préschémas $X$ et $Y$.

Par récurrence sur le nombre d'éléments de $I$, il résulte de (16.4.23) que $\Omega^1_{X^I/S}$ (resp. $\Omega^1_{Y/S}$) s'identifie canoniquement à la somme directe des $p_i^*(\Omega^1_{X/S})$ (resp. des $q_i^*(\Omega^1_{Y_i/S})$, si $q_i:Y\to Y_i$ sont les projections canoniques).

Supposons maintenant que $X$ soit \emph{lisse} sur $S$ en un point $x$ (donc dans un voisinage de ce point). Il en est alors de même de $X^I$ en ce point (17.3.3, (iv)), et restreignant $X$ à un voisinage de $x$, on a donc (17.2.5) une suite exacte de $\sh{O}_X$-Modules localement libres
\[
0\to\operatorname{gr}_1(\Delta)\to(\Omega^1_{X/S})^I\to\Omega^1_{X/S}\to0.
\]
Posons $\sh{J}=\operatorname{gr}_1(\Delta)$; on déduit de la suite précédente une suite exacte de $\sh{O}_T$-Modules localement libres
\[
\label{IV.17.13.9.1-seq}
\tag{17.13.9.1}
0\to\sh{J}\otimes_{\sh{O}_X}\sh{O}_T\to(\Omega^1_{X/S})^I\otimes_{\sh{O}_X}\sh{O}_T\xrightarrow{\sigma}\Omega^1_{X/S}\otimes_{\sh{O}_X}\sh{O}_T\to0
\]
qui correspond à la première ligne du diagramme (17.13.1.2), et où $\sigma$ n'est autre que l'homomorphisme canonique qui, à chaque famille de sections $(t_i)_{i\in I}$ de $\Omega^1_{X/S}\otimes_{\sh{O}_X}\sh{O}_T$ au-dessus d'un ouvert de $T$, fait correspondre sa \emph{somme}.

D'autre part, ce qui correspond ici à la seconde flèche verticale du diagramme (17.13.1.2) est l'homomorphisme
\[
\label{IV.17.13.9.2}
\tag{17.13.9.2}
\tau:(\Omega^1_{X/S})^I\otimes_{\sh{O}_X}\sh{O}_T\to\bigoplus_{i\in I}(\Omega^1_{Y_i/S}\otimes_{\sh{O}_{Y_i}}\sh{O}_T)
\]
qui, à toute famille $(t'_i\otimes1)_{i\in I}$, où ici les $t'_i$ sont des sections au-dessus d'un ouvert de $X$ de $\Omega^1_{X/S}$, associe la somme des $((f_i)_{Y_i/X/S}(t'_i))\otimes1$, avec la notation de (16.4.18). L'homomorphisme $\rho$ qui correspond à l'homomorphisme $\alpha$ de (17.13.1.3) est donc la \emph{restriction} à $\sh{J}\otimes_{\sh{O}_X}\sh{O}_T$ de l'homomorphisme $\tau$ précédent.

Ces hypothèses et notations étant précisées, on peut appliquer à la situation de (17.13.9.1) la prop. (17.13.2), ce qui donne le
\end{env}

\begin{corollary}[17.13.10]
\EGAReaderTarget{ega4.c04.cor.17-13-10}{IV.17.13.10}
Sous les hypothèses générales de (17.13.9), soit $t$ un point de $T$, $y_i$ ($i\in I$) ses projections dans les $Y_i$, $x$ le point de $X$ égal à chacun des $f_i(y_i)$. Supposons que $X$ soit lisse sur $S$ au point $x$ et de dimension relative $d$, et que chacun des $Y_i$ soit lisse sur $S$ au point $y_i$; soit $d+c_i$ ($c_i$ entier positif ou négatif) la dimension relative de $Y_i$ au point $y_i$. Alors les conditions suivantes sont équivalentes:

\oldpage[IV-4]{97}

\begin{enumerate}
  \item[a)] $T$ est lisse sur $S$ au point $t$, de dimension relative $d+\sum_{i\in I}c_i$.
  \item[b)] L'homomorphisme
  \[
  \rho\otimes1:\sh{J}\otimes_{\sh{O}_X}k(t)\to\bigoplus_{i\in I}(\Omega^1_{Y_i/S}\otimes_{\sh{O}_{Y_i}}k(t))
  \]
  est injectif.
\end{enumerate}
Lorsque $t$ est rationnel sur le corps $k(s)$ (où $s$ est l'image de $t$ dans $S$), ces conditions sont aussi équivalentes à la suivante:
\begin{enumerate}
  \item[b')] L'homomorphisme transposé de $\rho\otimes1$
  \[
  \bigoplus_{i\in I}\mathrm{T}_{Y_i/X}(y_i)\to(\mathrm{T}_{X/S}(x))^I/\delta(\mathrm{T}_{X/S}(x))
  \]
  (où $\delta$ est l'application diagonale) est surjectif.
\end{enumerate}
Il suffit de remarquer que $X^I$ est lisse sur $S$ et de dimension relative $d\cdot\operatorname{Card}(I)$, et $Y$ lisse sur $S$ de dimension relative $d\cdot\operatorname{Card}(I)+\sum_{i\in I}c_i$.

Lorsque les conditions équivalentes de (17.13.10) sont vérifiées, on dit encore que les $f_i$ forment une \emph{famille de morphismes transversaux au point $t$, relativement} à $S$. On voit encore que l'ensemble des points $t\in T$ où cela a lieu est \emph{ouvert} dans $T$. La remarque (17.13.4, (iii)) montre alors que si $X$ et les $Y_i$ sont plats et localement de présentation finie sur $S$, et de plus si $T$ est \emph{propre} sur $S$, l'ensemble des $s\in S$ tels que les $f_i$ forment une famille de morphismes transversaux en \emph{tous} les points de $T_s$, relativement à $S$, est \emph{ouvert} dans $S$.
\end{corollary}

\begin{env}[17.13.11]
\EGAReaderTarget{ega4.c04.statement.17-13-11}{IV.17.13.11}
Considérons en particulier le cas, généralisant (17.13.7), où les $Y_i$ ($i\in I$) sont des \emph{sous-préschémas} de $X$, les $f_i$ étant les injections canoniques, de sorte que $T=\inf_{i\in I}(Y_i)$ est encore le sous-préschéma « intersection » des $Y_i$, et $t=y_i=x$ pour tout $i$; au lieu de dire que les $f_i$ forment une famille de morphismes transversaux au point $x$, on dit encore que \emph{les $Y_i$ se coupent transversalement au point $x$} (relativement à $S$). La condition a) de (17.13.10) s'exprime encore en la relation qui généralise (17.13.7.1)
\[
\label{IV.17.13.11.1}
\tag{17.13.11.1}
\operatorname{codim}_x(T_s,X_s)=\sum_{i\in I}\operatorname{codim}_x((Y_i)_s,X_s).
\]
En outre, on a la propriété suivante, qui étend (17.13.8.1), et en donne une autre démonstration quand $I$ a 2 éléments:
\end{env}

\begin{corollary}[17.13.12]
\label{IV.17.13.12}
Lorsque les $Y_i$ se coupent transversalement au point $x$, on a un isomorphisme canonique
\[
\label{IV.17.13.12.1}
\tag{17.13.12.1}
\sh{N}_{T/X}\xrightarrow{\sim}\bigoplus_{i\in I}(\sh{N}_{Y_i/X}\otimes_{\sh{O}_{Y_i}}\sh{O}_T).
\]
\end{corollary}

On peut se borner au cas où les $Y_i$ sont des sous-préschémas fermés de $X$, définis par des Idéaux quasi-cohérents $\sh{J}_i$, de sorte que $T$ est défini par l'Idéal $\sh{J}=\sum_i\sh{J}_i$. Par définition du faisceau conormal à une immersion (16.1.3), l'homomorphisme canonique $\bigoplus_i\sh{J}_i\to\sh{J}$ donne, par passage aux quotients, un homomorphisme \emph{surjectif}
\[
\label{IV.17.13.12.2}
\tag{17.13.12.2}
\bigoplus_{i\in I}(\sh{N}_{Y_i/X}\otimes_{\sh{O}_{Y_i}}\sh{O}_T)\to\sh{N}_{T/X}.
\]

\oldpage[IV-4]{98}

Mais ici les $\sh{O}_T$-Modules des deux membres de (17.13.12.2) sont localement libres et \emph{de même rang} $\sum_i c_i$ (si $c_i$ est le rang de $\sh{N}_{Y_i/X}$), en vertu de (17.2.5) et de la condition a) de (17.13.10); on conclut donc de Bourbaki, \emph{Alg. comm.}, chap. II, § 3, n\textsuperscript{o} 2, cor. de la prop. 6, que (17.13.12.2) est bijectif, et (17.13.12.1) est l'isomorphisme réciproque.

\subsection{Caractérisations locales et infinitésimales des morphismes lisses, des morphismes non ramifiés et des morphismes étales.}
\EGAReaderTarget{ega4.c04.subsec.17-14}{subsection:IV.17.14}

\begin{proposition}[17.14.1]
\label{IV.17.14.1}
Soit $f:X\to Y$ un morphisme localement de présentation finie, $x$ un point de $X$, $y=f(x)$. Pour que $f$ soit lisse (resp. non ramifié, resp. étale) au point $x$, il faut et il suffit que la condition suivante soit vérifiée:

Pour tout schéma local $Y'=\operatorname{Spec}(A')$ de point fermé $y'$, tout morphisme $h:Y'\to Y$ tel que $h(y')=y$, tout sous-schéma fermé $Y'_0=\operatorname{Spec}(A'/\mathfrak{J}')$ de $Y'$, où l'idéal $\mathfrak{J}'$ est de carré nul, et tout $Y$-morphisme $g_0:Y'_0\to X$ tel que $g_0(y')=x$, il existe au moins un (resp. au plus un, resp. un et un seul) $Y$-morphisme $g:Y'\to X$ dont $g_0$ soit la restriction à $Y'_0$.
\end{proposition}

Compte tenu des définitions (17.1.1 et 17.3.1), il s'agit de montrer que la condition de l'énoncé est \emph{suffisante} pour que $f$ soit lisse (resp. non ramifié) au point $x$.

\emph{(i) Cas des morphismes lisses.} On peut se borner au cas où $Y=\operatorname{Spec}(A)$ et $X=\operatorname{Spec}(C)$ sont affines, avec $C=B/\mathfrak{R}$, où $B=A[T_1,\ldots,T_n]$ est une $A$-algèbre de polynômes et $\mathfrak{R}$ un idéal de type fini de $B$. Pour prouver que $f$ est lisse au point $x$, il suffit d'établir que $\sh{O}_{X,x}$ est une $\sh{O}_{Y,y}$-algèbre formellement lisse pour les topologies discrètes (17.5.1). Or on a $\sh{O}_{X,x}=B_{\mathfrak q}/\mathfrak{R}B_{\mathfrak q}$, où $\mathfrak q$ est un idéal premier de $B$, et si $\mathfrak p$ est l'image réciproque de $\mathfrak q$ dans $A$, $B_{\mathfrak q}$ est une $A_{\mathfrak p}$-algèbre formellement lisse pour les topologies discrètes (0, 19.3.2 et 19.3.5). Par application du critère jacobien (0, 22.6.1 et 20.5.12) il suffit donc de voir que $B_{\mathfrak q}/\mathfrak{R}^2B_{\mathfrak q}$ est une extension $A_{\mathfrak p}$-triviale de $B_{\mathfrak q}/\mathfrak{R}B_{\mathfrak q}$ par $\mathfrak{R}B_{\mathfrak q}/\mathfrak{R}^2B_{\mathfrak q}$. Mais cela résulte précisément de l'hypothèse appliquée à $A'=B_{\mathfrak q}/\mathfrak{R}^2B_{\mathfrak q}$ et $\mathfrak{J}'=\mathfrak{R}B_{\mathfrak q}/\mathfrak{R}^2B_{\mathfrak q}$ et au morphisme $g_0$ correspondant à l'automorphisme identique de $A'/\mathfrak{J}'=\sh{O}_{X,x}$.

\emph{(ii) Cas des morphismes non ramifiés.} Posons ici $A=\sh{O}_{Y,y}$, $B=\sh{O}_{X,x}$, et notons qu'en vertu de (17.4.1, c')), il suffit de montrer que $\Omega^1_{B/A}=(\Omega^1_{X/Y})_x=0$. Or, avec les notations de (0, 20.4.1), l'anneau $C=\sh{P}^1_{B/A}=(B\otimes_A B)/\mathfrak{J}^2$ est local puisqu'il en est ainsi de $C/(\mathfrak{J}/\mathfrak{J}^2)$ isomorphe à $B$. L'hypothèse appliquée à $A'=C$ et $\mathfrak{J}'=\mathfrak{J}/\mathfrak{J}^2$ montre par définition que l'application $d_{B/A}:B\to\Omega^1_{B/A}$ est nulle (0, 20.4.6), donc que $\Omega^1_{B/A}=0$ par (0, 20.4.7).

\begin{proposition}[17.14.2]
\label{IV.17.14.2}
Soient $Y$ un préschéma localement noethérien, $f:X\to Y$ un morphisme localement de type fini, $x$ un point de $X$, $y=f(x)$. Pour que $f$ soit lisse (resp. non ramifié, resp. étale) au point $x$, il faut et il suffit qu'il vérifie la condition de (17.14.1) où l'on suppose de plus que l'anneau local $A'$ est artinien, que $\mathfrak m'\mathfrak{J}'=0$, où $\mathfrak m'$ est l'idéal maximal de $A'$, et que le corps résiduel $A'/\mathfrak m'$ de $A'$ est égal à $k(x)$.
\end{proposition}

Il s'agit encore de montrer que ces conditions sont suffisantes dans le cas des morphismes lisses et le cas des morphismes non ramifiés.

\oldpage[IV-4]{99}

\emph{(i) Cas des morphismes lisses.} Si l'on pose $A=\sh{O}_{Y,y}$ et $B=\sh{O}_{X,x}$, il s'agit ici de prouver, compte tenu de (17.5.3), que $B$ est une $A$-algèbre formellement lisse pour les topologies préadiques. Or, en vertu de (0, 22.1.4), cela découle de l'hypothèse lorsqu'on remplace dans celle-ci la condition $\mathfrak m'\mathfrak{J}'=0$ par $\mathfrak{J}'\subset\mathfrak m'$. Mais comme $\mathfrak m'$ est nilpotent, on voit que la condition de l'énoncé est déjà suffisante en considérant les anneaux $A'/\mathfrak m'^{,j}\mathfrak{J}'$ et en raisonnant par récurrence sur $j$ comme dans (0, 19.4.3).

\emph{(ii) Cas des morphismes non ramifiés.} Reprenons les notations de la démonstration de (17.14.1, (ii)); pour prouver que l'idéal $\mathfrak{J}/\mathfrak{J}^2$ de $C$ est nul, notons que $C$ est alors un anneau local noethérien, étant anneau local de l'anneau d'un ouvert affine d'un sous-préschéma de $X\times_YX$, qui est localement de type fini sur $X$, donc aussi sur $Y$. Si $\mathfrak r$ est l'idéal maximal de $C$, il suffit donc de vérifier que l'on a $\mathfrak{J}/\mathfrak{J}^2\subset\mathfrak r^n$ pour tout $n$ (0$_{\mathrm I}$, 7.3.5), ou encore que $\mathfrak r^n=\mathfrak r^n+(\mathfrak{J}/\mathfrak{J}^2)$. Avec les notations de (0, 20.4.6), cela signifie encore que pour tout $n$, les deux applications composées $B\xrightarrow{p_1}C\to C/\mathfrak r^n$ et $B\xrightarrow{p_2}C\to C/\mathfrak r^n$ sont identiques; mais comme $C/\mathfrak r^n$ est artinien, cette identité résulte de l'hypothèse appliquée, par récurrence sur $n$, à $A'=C/\mathfrak r^n$ et $\mathfrak{J}'=(\mathfrak r^n+(\mathfrak{J}/\mathfrak{J}^2))/\mathfrak r^n$.

\begin{remark}[17.14.3]
\label{IV.17.14.3}
On notera que, sous les hypothèses de (17.14.2), si de plus le corps $k(x)$ est une extension \emph{finie} de $k(y)$, alors les $A'$-modules $\mathfrak m'^{,j}/\mathfrak m'^{,j+1}$ sont des $k(y)$-espaces vectoriels de rang fini, donc \emph{a fortiori} $A'$ est une $\sh{O}_{Y,y}$-algèbre finie.
\end{remark}

\subsection{Cas des préschémas sur un corps de base.}
\EGAReaderTarget{ega4.c04.subsec.17-15}{subsection:IV.17.15}

Rappelons d'abord (6.7.7, 6.7.8 et 6.8.1) la

\begin{proposition}[17.15.1]
\label{IV.17.15.1}
Soient $k$ un corps, $X$ un préschéma localement de type fini sur $k$. Pour que $X$ soit lisse en un point $x$, il faut et il suffit que pour toute extension radicielle $k'$ de $k$ (ou seulement pour toute extension finie radicielle de $k$), l'anneau local $\sh{O}_{X,x}\otimes_k k'$ soit régulier. Si $k(x)$ est une extension séparable de $k$ (en particulier si $k$ est parfait), il revient au même de dire que $\sh{O}_{X,x}$ est régulier.
\end{proposition}

\begin{corollary}[17.15.2]
\label{IV.17.15.2}
Sous les conditions de (17.15.1), pour que $X$ soit lisse sur $k$, il faut et il suffit que $X$ soit géométriquement régulier sur $k$ (6.7.6), ce qui entraîne que $X$ est régulier. Si $k$ est parfait, alors $X$ est lisse sur $k$ si et seulement si $X$ est régulier.
\end{corollary}

\begin{proposition}[17.15.3]
\label{IV.17.15.3}
Soient $k$ un corps, $X$ un préschéma localement de type fini sur $k$, $x$ un point de $X$, $n=\dim_x(X)$. Soit $(g_i)_{1\leq i\leq n}$ une famille de $n$ sections de $\sh{O}_X$ au-dessus de $X$, et soit $f:X\to Y=\operatorname{Spec}(k[T_1,\ldots,T_n])$ le morphisme correspondant au $k$-homomorphisme $k[T_1,\ldots,T_n]\to\Gamma(X,\sh{O}_X)$ transformant $T_i$ en $g_i$ (I, 2.2.4). Les conditions suivantes sont équivalentes (et impliquent que $X$ est lisse sur $k$ au point $x$ et par suite régulier au point $x$):
\begin{enumerate}
  \item[a)] $f$ est étale au point $x$.
  \item[b)] Les images des $d_{X/k}g_i$ forment une base du $\sh{O}_{X,x}$-module $(\Omega^1_{X/k})_x$.
  \item[c)] Les images des $d_{X/k}g_i$ engendrent le $\sh{O}_{X,x}$-module $(\Omega^1_{X/k})_x$.
\end{enumerate}
\end{proposition}

Si $f$ est étale au point $x$, $X$ est lisse sur $k$ au point $x$ puisque $Y=\operatorname{Spec}(k[T_1,\ldots,T_n])$ est lisse sur $k$ (17.3.8); le fait que a) entraîne b) est donc un cas particulier de (17.11.3). Comme b) entraîne trivialement c), il reste à voir que c) implique a). Notons d'abord

\oldpage[IV-4]{100}

que l'hypothèse entraîne que l'homomorphisme $(f^*(\Omega^1_{Y/k}))_x\to(\Omega^1_{X/k})_x$ est surjectif, donc, en remplaçant $X$ par un voisinage ouvert de $x$, on peut supposer que l'homomorphisme $f^*(\Omega^1_{Y/k})\to\Omega^1_{X/k}$ est surjectif (Bourbaki, \emph{Alg. comm.}, chap. II, § 5, n\textsuperscript{o} 1, prop. 2); par suite $f$ est \emph{non ramifié} (17.2.2). Nous allons voir d'abord qu'on peut se borner au cas où $x$ est \emph{rationnel sur $k$}. En effet, si l'on pose $k'=k(x)$, et $X'=X\otimes_k k'$, $Y'=Y\otimes_k k'=\operatorname{Spec}(k'[T_1,\ldots,T_n])$, il existe un point $x'\in X'$ au-dessus de $x$, tel que $k(x')=k'$. Pour prouver que $f$ est étale au point $x$, il suffit de montrer que $f'=f_{(k')}:X'\to Y'$ est étale au point $x'$ (17.7.1, (ii)); d'ailleurs $f'$ est non ramifié (17.3.3, (iii)) et l'on a $\dim_{x'}(X')=n$ (4.2.7). De la même manière on peut, en remplaçant $k'$ par une extension algébriquement close de $k'$, supposer que $k$ est \emph{algébriquement clos}. Posons alors $y=f(x)$, $A=\sh{O}_{Y,y}$, $B=\sh{O}_{X,x}$; comme le corps résiduel de $B$ est égal à $k$, il en est de même de celui de $A$, donc $x$ (resp. $y$) est un point fermé de $X$ (resp. $Y$) (I, 6.4.2) et l'on a par suite $\dim(A)=\dim(B)=n$. Puisque $f$ est non ramifié, l'homomorphisme $\widehat A\to\widehat B$ est surjectif (17.4.4, f'')); mais comme $\dim(\widehat A)=\dim(\widehat B)=n$, et que $\widehat A$, étant un anneau régulier, est intègre, l'homomorphisme $\widehat A\to\widehat B$ est aussi injectif (0, 16.3.10). Cet homomorphisme est donc bijectif, ce qui achève de prouver que $f$ est étale au point $x$ (17.6.3, e'')).

\begin{corollary}[17.15.4]
\label{IV.17.15.4}
Sous les hypothèses générales de (17.15.3), supposons en outre que $k(x)$ soit une extension finie et séparable de $k$ et que les germes $(g_i)_x$ appartiennent à $\mathfrak m_x$. Alors les conditions a), b), c) de (17.15.3) équivalent aussi à chacune des suivantes:
\begin{enumerate}
  \item[d)] Les $n$ germes $(g_i)_x$ engendrent l'idéal maximal $\mathfrak m_x$ de $\sh{O}_{X,x}$.
  \item[d')] L'anneau $\sh{O}_{X,x}$ est régulier et les $(g_i)_x$ forment un système régulier de paramètres pour cet anneau (0, 17.1.6).
\end{enumerate}
\end{corollary}

En effet, d') entraîne trivialement d). D'autre part, comme $k(x)$ est une extension finie et séparable de $k$, on a $\Omega^1_{k(x)/k}=0$ (0, 20.6.20) et $k(x)=\sh{O}_{X,x}/\mathfrak m_x$ est une $k$-algèbre formellement lisse pour les topologies discrètes (0, 19.6.1), donc la suite exacte (0, 20.5.14.1) s'applique à $A=k$, $B=\sh{O}_{X,x}$, $\mathfrak{R}=\mathfrak m_x$, et fournit un isomorphisme canonique
\[
\delta:\mathfrak m_x/\mathfrak m_x^2\xrightarrow{\sim}(\Omega^1_{X/k})_x\otimes_{\sh{O}_{X,x}}k(x).
\]
De là on déduit d'abord l'équivalence des conditions c) et d), compte tenu du lemme de Nakayama. D'autre part, si $f$ est étale au point $x$, l'anneau $\sh{O}_{X,x}$ est régulier et de dimension $n$, puisque $x$ est un point fermé de $X$ (I, 6.4.2), et $n$ éléments de $\mathfrak m_x$ qui engendrent cet idéal maximal forment alors nécessairement un système régulier de paramètres pour $\sh{O}_{X,x}$ (0, 17.1.6), ce qui prouve que a) entraîne d').

\begin{proposition}[17.15.5]
\label{IV.17.15.5}
Soient $k$ un corps, $X$ un préschéma localement de type fini sur $k$, $x$ un point de $X$, $n=\dim_x(X)$. Les conditions suivantes sont équivalentes:
\begin{enumerate}
  \item[a)] $X$ est lisse sur $k$ au point $x$.
  \item[b)] $X$ est différentiellement lisse sur $k$ au point $x$.
  \item[c)] $(\Omega^1_{X/k})_x$ est un $\sh{O}_{X,x}$-module libre de rang $n$.
  \item[d)] $(\Omega^1_{X/k})_x$ est un $\sh{O}_{X,x}$-module admettant un système de $n$ générateurs.

\oldpage[IV-4]{101}

  \item[e)] Il existe une extension parfaite $k'$ de $k$ et un voisinage ouvert $U$ de $x$ dans $X$ tels que le préschéma $U\otimes_k k'$ soit régulier.
\end{enumerate}
\end{proposition}

Le fait que $X$ soit lisse sur $k$ au point $x$ entraîne l'existence d'un voisinage ouvert $U$ de $x$ qui est lisse sur $k$, donc $U\otimes_k k'$ est régulier pour toute extension $k'$ de $k$ (17.15.2), ce qui prouve que a) implique e). Inversement e) implique a), car alors $U\otimes_k k'$ est lisse sur $k'$ (17.15.2), donc $U$ est lisse sur $k$ (17.7.1, (ii)).

On a déjà prouvé que a) entraîne c) (17.10.2); c) implique d) trivialement et le fait que d) implique a) résulte de (0, 20.4.7) et (17.15.3, c)).

Enfin, on a déjà vu que a) implique b) (17.12.4). Inversement, pour prouver que b) entraîne a), on peut se borner au cas où $x$ est rationnel sur $k$, en considérant, comme dans la démonstration de (17.15.3), un point $x'$ de $X'=X\otimes_k k'$ au-dessus de $x$ et tel que $k(x')=k'=k(x)$, utilisant encore (17.7.1, (ii)) et le fait que si $X$ est différentiellement lisse au point $x$, $X'$ est différentiellement lisse au point $x'$ (16.10.4). Supposant donc $x$ rationnel sur $k$, le fait que $f$ soit lisse en $x$ résulte alors de l'hypothèse b) et de (17.12.5) appliqué à la $k$-section $u:\operatorname{Spec}(k)\to X$ telle que $u(\operatorname{Spec}(k))=\{x\}$.

\begin{corollary}[17.15.6]
\label{IV.17.15.6}
Soit $X$ un préschéma de type fini sur un corps $k$. Pour que $X$ soit lisse sur $k$, il faut et il suffit que le $\sh{O}_X$-Module $\Omega^1_{X/k}$ soit localement libre, et que les anneaux locaux aux points maximaux de $X$ soient des corps, extensions séparables de $k$ (cette dernière condition étant automatiquement vérifiée si $k$ est un corps parfait et $X$ un préschéma réduit).
\end{corollary}

Les conditions sont nécessaires, car si $X$ est lisse sur $k$, il résulte de (17.10.2) que $\Omega^1_{X/k}$ est localement libre; d'autre part, $X$ est régulier, donc \emph{a fortiori} réduit, donc en tout point maximal $x$ de $X$, $\sh{O}_{X,x}$ est un corps, qui doit être une $k$-algèbre formellement lisse pour les topologies discrètes (17.5.1), donc une extension séparable de $k$ (0, 19.6.1).

Les conditions sont suffisantes. On peut en effet se borner au cas où $X$ est connexe, donc $\Omega^1_{X/k}$ localement libre de rang constant $n$. Pour tout point maximal $x$ de $X$, on a alors $\operatorname{rg}_{k(x)}\Omega^1_{k(x)/k}=n$ puisque $\sh{O}_{X,x}=k(x)$; comme par hypothèse $k(x)$ est une extension séparable de $k$, on a $\Upsilon_{k(x)/k}=0$ (0, 20.6.3), donc $\deg.\operatorname{tr}_k k(x)=n$ par l'égalité de Cartier (0, 21.7.1). Toutes les composantes irréductibles de $X$ ont donc même dimension $n$ (5.2.1), et l'on conclut que $X$ est lisse sur $k$ en tout point en vertu de ce que c) entraîne a) dans (17.15.5).

\begin{corollary}[17.15.7]
\label{IV.17.15.7}
Si $k$ est de caractéristique 0, alors, pour que $X$ soit lisse sur $k$ au point $x$, il faut et il suffit que $(\Omega^1_{X/k})_x$ soit un $\sh{O}_{X,x}$-module libre.
\end{corollary}

Cela résulte de (16.12.2) et de l'équivalence de b) et a) dans (17.15.5).

\begin{proposition}[17.15.8]
\label{IV.17.15.8}
Soient $k$ un corps, $X$ un préschéma localement de type fini sur $k$, $x$ un point de $X$; posons $n=\dim_x(X)$, $r=\dim(\sh{O}_{X,x})$, de sorte que $n-r=\deg.\operatorname{tr}_k k(x)$ (5.2.3). Soit $(g_i)_{1\leq i\leq n}$ une famille de $n$ sections de $\sh{O}_X$ au-dessus de $X$, telles que $(g_i)_x\in\mathfrak m_x$ pour $1\leq i\leq r$; soit $f:X\to Y=\operatorname{Spec}(k[T_1,\ldots,T_n])$ le morphisme correspondant au $k$-homomorphisme $k[T_1,\ldots,T_n]\to\Gamma(X,\sh{O}_X)$ transformant $T_i$ en $g_i$. Les conditions suivantes sont équivalentes (et impliquent que $X$ est lisse sur $k$ au point $x$ et que $k(x)$ est une extension séparable de $k$):
\begin{enumerate}
  \item[a)] $f$ est étale au point $x$.

\oldpage[IV-4]{102}

  \item[b)] Les $(g_i)_x$ tels que $1\leq i\leq r$ engendrent $\mathfrak m_x$ (et par suite forment un système régulier de paramètres pour $\sh{O}_{X,x}$ (0, 17.1.6)) et les images dans $\Omega^1_{k(x)/k}$ des éléments $(d_{X/k}g_i)_x$ pour $r+1\leq i\leq n$ engendrent $\Omega^1_{k(x)/k}$.
\end{enumerate}
\end{proposition}

On a en effet (0, 20.5.12.1) la suite exacte de $k(x)$-modules
\[
\label{IV.17.15.8.1}
\tag{17.15.8.1}
\mathfrak m_x/\mathfrak m_x^2\to(\Omega^1_{X/k})_x\otimes_{\sh{O}_{X,x}}k(x)\to\Omega^1_{k(x)/k}\to0
\]
et la condition b) entraîne par suite que les $(d_{X/k}g_i)_x$ engendrent $(\Omega^1_{X/k})_x$ compte tenu du lemme de Nakayama; le fait que b) implique a) résulte donc de (17.15.3). Inversement, si a) est vérifiée, les $(d_{X/k}g_i)_x$ pour $1\leq i\leq n$ forment une base de $(\Omega^1_{X/k})_x$ en vertu de (17.15.3). Si $t_i$ ($1\leq i\leq n$) est l'image de $(g_i)_x$ dans $k(x)$, on conclut de ce qui précède que les $dt_i$ engendrent $\Omega^1_{k(x)/k}$ pour $1\leq i\leq n$, et comme par hypothèse $t_i=0$ pour $1\leq i\leq r$, les $dt_i$ pour $r+1\leq i\leq n$ engendrent déjà $\Omega^1_{k(x)/k}$. Comme $\deg.\operatorname{tr}_k k(x)=n-r$, il suit de l'égalité de Cartier (0, 21.7.1) que $\Upsilon_{k(x)/k}=0$, donc $k(x)$ est une extension séparable de $k$ (0, 20.6.3), et la suite de $k(x)$-espaces vectoriels
\[
\label{IV.17.15.8.2}
\tag{17.15.8.2}
0\to\mathfrak m_x/\mathfrak m_x^2\to(\Omega^1_{X/k})_x\otimes_{\sh{O}_{X,x}}k(x)\to\Omega^1_{k(x)/k}\to0
\]
est exacte ((0, 20.5.14) et (0, 19.6.1)); en outre les $dt_i$ pour $r+1\leq i\leq n$ forment une base de $\Omega^1_{k(x)/k}$, donc aucun des $t_i$ tels que $r+1\leq i\leq n$ ne peut être nul. Cela montre que les images des $(g_i)_x$ dans $\mathfrak m_x/\mathfrak m_x^2$ pour $1\leq i\leq r$ engendrent nécessairement $\mathfrak m_x/\mathfrak m_x^2$, donc les $(g_i)_x$ pour $1\leq i\leq r$ engendrent $\mathfrak m_x$ en vertu du lemme de Nakayama, ce qui achève de prouver que a) implique b).

\begin{corollary}[17.15.9]
\label{IV.17.15.9}
Soient $X$ un préschéma localement de type fini sur un corps $k$, $x$ un point de $X$ et posons $n=\dim_x(X)$, $r=\dim(\sh{O}_{X,x})$. Les conditions suivantes sont équivalentes:
\begin{enumerate}
  \item[a)] L'anneau $\sh{O}_{X,x}$ est régulier et $k(x)$ est une extension séparable de $k$.
  \item[b)] $X$ est lisse sur $k$ au point $x$, et l'homomorphisme canonique
  \[
  \mathfrak m_x/\mathfrak m_x^2\to(\Omega^1_{X/k})_x\otimes_{\sh{O}_{X,x}}k(x)
  \]
  est injectif.
  \item[c)] Il existe $n$ sections $g_i$ de $\sh{O}_X$ au-dessus d'un voisinage ouvert $U$ de $x$ telles que $(g_i)_x\in\mathfrak m_x$ pour $1\leq i\leq r$ et que le morphisme $U\to\operatorname{Spec}(k[T_1,\ldots,T_n])$ correspondant aux $g_i$ (cf. (17.15.8)) soit étale au point $x$.
  \item[d)] Il existe $n$ sections $g_i$ ($1\leq i\leq n$) de $\sh{O}_X$ au-dessus d'un voisinage ouvert $U$ de $x$, telles que les $(g_i)_x$ pour $1\leq i\leq r$ engendrent $\mathfrak m_x$ et que les images dans $\Omega^1_{k(x)/k}$ des $(d_{U/k}g_i)_x$ pour $r+1\leq i\leq n$ engendrent $\Omega^1_{k(x)/k}$.
\end{enumerate}
\end{corollary}

L'équivalence de c) et d) résulte de (17.15.8). De plus, on a vu dans la démonstration de (17.15.8) que la condition c) entraîne que $X$ est lisse sur $k$ au point $x$ et que la suite (17.15.8.2) est exacte, donc c) entraîne b). La condition b) entraîne que l'anneau $\sh{O}_{X,x}$ est régulier, donc $\mathfrak m_x/\mathfrak m_x^2$ est de rang $r$; en outre, $(\Omega^1_{X/k})_x$ est alors un $\sh{O}_{X,x}$-module libre de rang $n$ (17.10.2); comme la suite (17.15.7.2) est exacte par hypothèse, $\Omega^1_{k(x)/k}$ est de rang $n-r=\deg.\operatorname{tr}_k k(x)$ et l'égalité de Cartier montre que $\Upsilon_{k(x)/k}=0$, donc $k(x)$ est séparable sur $k$ (0, 20.6.3); ainsi b) entraîne a). Enfin, si a)

\oldpage[IV-4]{103}

est vérifiée, on déduit encore de l'égalité de Cartier que $\Omega^1_{k(x)/k}$ est de rang $n-r$. Comme d'autre part l'anneau $\sh{O}_{X,x}$ est régulier, l'existence des $g_i$ vérifiant les conditions de d) est immédiate.

\begin{remarks}[17.15.10]
\label{IV.17.15.10}
\begin{enumerate}
  \item[(i)] Pour qu'un préschéma de type fini $X$ sur $k$ soit lisse sur $k$, il ne suffit pas que $\Omega^1_{X/k}$ soit localement libre, comme le montre l'exemple où $X=\operatorname{Spec}(K)$, $K$ étant une extension finie non séparable de $k$.
  \item[(ii)] Lorsque $k$ n'est pas parfait, il peut se faire que $X$ soit lisse sur $k$ sans que $k(x)$ soit séparable sur $k$. On en a un exemple en prenant $X=\operatorname{Spec}(k[T])$ et pour $x$ le point correspondant à l'idéal premier principal $(T^p-\lambda)$ ($p>0$ caractéristique de $k$, $\lambda\in k-k^p$).
  \item[(iii)] Toutefois, si $X$ est un préschéma lisse sur $k$, l'ensemble des points fermés de $X$ tels que $k(x)$ soit séparable (et fini) sur $k$ est dense dans $X$. En effet, soit $f:X\to\operatorname{Spec}(k)$ le morphisme structural; pour tout $x_0\in X$, il y a un voisinage ouvert $U$ de $x_0$ dans $X$ et une factorisation de $f|U:U\xrightarrow{g}\operatorname{Spec}(k[T_1,\ldots,T_n])\xrightarrow{h}\operatorname{Spec}(k)$ où $g$ est étale (17.11.4). Comme on peut se borner au cas où $k$ est non parfait, donc infini, l'ensemble des points de l'ouvert $g(U)$ qui sont rationnels sur $k$ est non vide; si $y$ est un tel point et $x\in U$ un point au-dessus de $y$, $x$ est fermé dans $X$ et $k(x)$ est séparable sur $k(y)=k$ (17.6.2).
\end{enumerate}
\end{remarks}

\begin{proposition}[17.15.11]
\label{IV.17.15.11}
Soit $X$ un préschéma de type fini sur un corps $k$. Les conditions suivantes sont équivalentes:
\begin{enumerate}
  \item[a)] $X$ est étale sur $k$.
  \item[b)] $X$ est non ramifié sur $k$.
  \item[c)] $X$ est isomorphe à $\operatorname{Spec}(A)$, où $A$ est une $k$-algèbre finie et séparable.
\end{enumerate}
\end{proposition}

Cela résulte de (17.6.2) et (17.4.2), compte tenu de ce que la condition b) implique ici que $X$ est fini sur $k$.

\begin{proposition}[17.15.12]
\label{IV.17.15.12}
Soient $k$ un corps, $X$ un préschéma localement de type fini sur $k$. Pour qu'il existe un ouvert partout dense $U$ de $X$ qui soit lisse sur $k$, il faut et il suffit que pour tout point maximal $x$ de $X$, $X$ soit réduit en ce point et que $k(x)$ soit une extension séparable de $k$. Pour qu'il existe un ouvert partout dense $U$ de $X$ tel que $U_{\mathrm{red}}$ soit lisse sur $k$, il faut et il suffit que pour tout point maximal $x$ de $X$, $k(x)$ soit une extension séparable de $k$.
\end{proposition}

La seconde assertion résulte évidemment de la première. Dire qu'il existe un ouvert partout dense $U$ lisse sur $k$ signifie que $X$ est lisse sur $k$ en chacun de ses points maximaux $x$. Il est nécessaire pour cela que $\sh{O}_{X,x}$ soit régulier (17.15.1) et \emph{a fortiori} réduit, donc un corps, égal à $k(x)$; en outre (17.15.1), $k(x)\otimes_k k'$ doit être régulier pour toute extension radicielle $k'$ de $k$, donc $k(x)$ doit être une extension séparable de $k$ (4.3.5); la réciproque est immédiate, en vertu de (17.15.1).

\begin{corollary}[17.15.13]
\label{IV.17.15.13}
Soit $X$ un préschéma algébrique sur un corps $k$. Alors il existe un ouvert partout dense $U$ dans $X$ et une extension finie radicielle $k'$ de $k$ tels que $(U_{(k')})_{\mathrm{red}}$ soit lisse sur $k'$.
\end{corollary}

En effet, il y a une telle extension $k'$ telle que $(X_{(k')})_{\mathrm{red}}$ soit géométriquement réduit sur $k'$ (4.6.6), ce qui équivaut à dire (4.6.1) que pour tout point maximal $x'$ de $X_{(k')}$, $k(x')$ est une extension séparable de $k'$. On conclut donc de (17.15.12) qu'il y a un ouvert partout dense $U'$ de $X_{(k')}$ tel que $U'_{\mathrm{red}}$ soit lisse sur $k'$. Mais le \oldpage[IV-4]{104}morphisme $X_{(k')}\to X$ est un homéomorphisme (2.4.5), donc $U'$ est de la forme $U_{(k')}$, où $U$ est un ouvert partout dense dans $X$.

\begin{proposition}[17.15.14]
\label{IV.17.15.14}
Soit $X$ un préschéma algébrique sur $k$, de dimension $\leq1$. Alors il existe une extension finie radicielle $k'$ de $k$ telle que le normalisé (II, 6.3.8) de $(X_{(k')})_{\mathrm{red}}$ soit lisse sur $k'$.
\end{proposition}

Démontrons d'abord les deux lemmes suivants:

\begin{lemma}[17.15.14.1]
\label{IV.17.15.14.1}
Soient $X$ un préschéma réduit dont l'ensemble des composantes irréductibles est localement fini, $f:Y\to X$ un morphisme. Pour que $Y$ soit $X$-isomorphe au normalisé $X'$ de $X$ (II, 6.3.8), il faut et il suffit que les conditions suivantes soient vérifiées:
\begin{enumerate}
  \item[(i)] $Y$ est normal;
  \item[(ii)] $f$ est un morphisme entier et birationnel.
\end{enumerate}
Lorsqu'il existe dans $X$ un ouvert dense et normal (ce qui est toujours le cas lorsque $X$ est excellent (7.8.3), en particulier lorsque $X$ est localement de type fini sur un corps), on peut remplacer la condition (ii) par la suivante:
\begin{enumerate}
  \item[(ii bis)] $f$ est entier et il existe un ouvert dense $U$ dans $X$ tel que $f^{-1}(U)$ soit dense dans $Y$ et que la restriction $f^{-1}(U)\to U$ de $f$ soit un isomorphisme.
\end{enumerate}
\end{lemma}

La question étant locale sur $X$, on peut supposer que $X$ n'a qu'un nombre fini de composantes irréductibles; soient $X_i$ ($1\leq i\leq m$) les sous-préschémas réduits de $X$ ayant ces composantes pour espaces sous-jacents. On sait (II, 6.3.6) que $X'$ est somme des normalisés $X'_i$ des $X_i$; chacun des morphismes structuraux $f_i:X'_i\to X_i$ est donc entier et birationnel, donc $f$ est entier et birationnel; il est immédiat que (ii) entraîne (ii bis) lorsqu'il existe un ouvert $V$ dense et normal dans $X$, en prenant $U=V$. Inversement, si les conditions (i) et (ii) sont remplies, $Y$ n'a qu'un nombre fini de composantes irréductibles $Y_i$ ($1\leq i\leq m$), $Y_i$ dominant $X_i$ pour chaque indice $i$ (6.15.4); $Y$ est somme des $Y_i$ puisqu'il est normal, et le morphisme $g_i:Y_i\to X_i$ en lequel se factorise la restriction $Y_i\to X$ de $f$ (I, 5.2.2) est birationnel; comme $X_i$ et $Y_i$ sont intègres, il résulte alors du fait que $g_i$ est entier et $Y_i$ normal que, pour tout ouvert affine $V$ de $X_i$, $\Gamma(g_i^{-1}(V),\sh{O}_{Y_i})$ est la clôture intégrale de $\Gamma(V,\sh{O}_{X_i})$, donc $Y$ s'identifie canoniquement à $X'$ (II, 6.3.4). Enfin, si l'on suppose vérifiée la condition (ii bis), on peut se borner au cas où $U$ est réunion d'ouverts irréductibles disjoints $U_i$ ($1\leq i\leq m$), donc $f^{-1}(U)$ réunion disjointe des $V_i=f^{-1}(U_i)$. Comme les $V_i$ sont irréductibles et $f^{-1}(U)$ dense dans $Y$, les $\overline{V_i}$ sont les composantes irréductibles de $Y$ et par suite $f$ est birationnel.

\begin{lemma}[17.15.14.2]
\label{IV.17.15.14.2}
Soient $k$ un corps, $X$ un préschéma algébrique sur $k$. Alors il existe une extension finie radicielle $k'$ de $k$ telle que $(X_{(k')})_{\mathrm{red}}$ soit géométriquement réduit sur $k'$ et que son normalisé soit géométriquement normal sur $k'$.
\end{lemma}

Compte tenu de (4.6.6) et de (I, 5.1.8), on peut se borner au cas où $X$ est déjà géométriquement réduit sur $k$. Soient $p$ l'exposant caractéristique de $k$ et $k_1=k^{p^{-\infty}}$ la plus petite extension parfaite de $k$; $k_1$ est donc limite inductive des extensions finies radicielles de $k$. Posons $X_1=X_{(k_1)}$, qui est réduit par hypothèse, et soit $Y_1$ son normalisé; si $f_1:Y_1\to X_1$ est le morphisme structural, $f_1$ est donc fini (7.8.3, (vi)) et surjectif, et il existe un ouvert dense $U_1$ dans $X_1$ tel que $f_1^{-1}(U_1)=V_1$ soit dense dans $Y_1$ et que \oldpage[IV-4]{105}le morphisme $V_1\to U_1$ restriction de $f_1$ soit un isomorphisme (17.15.14.1). Appliquant (8.9.1) et (8.10.5, (vi) et (x)), on voit donc d'abord qu'il existe une extension finie radicielle $k''$ de $k$ et un morphisme fini surjectif $Y''\to X''=X_{(k'')}$ tels que $Y_1=Y''\times_{X''}X_1=Y''\otimes_{k''}k_1$; puisque $Y_1$ est normal et $k_1$ parfait, il résulte de (6.7.7) que $Y''$ est géométriquement normal sur $k''$. Comme les projections $X_1\to X''$ et $Y_1\to Y''$ sont des morphismes entiers, surjectifs et radiciels, ce sont des homéomorphismes (2.4.5), et si $U''$ et $V''$ sont les images de $U_1$ et $V_1$ dans $X''$ et $Y''$ respectivement, ce sont des ouverts denses dans $X''$ et $Y''$ respectivement tels que $V''=f''^{-1}(U'')$, où $f'':Y''\to X''$ est le morphisme structural. Comme on a $U_1=U''\otimes_{k''}k_1$, il résulte de (8.10.5, (i)) qu'il existe une extension finie radicielle $k'$ de $k''$ telle que si l'on pose $X'=X\otimes_k k'=X''\otimes_{k''}k'$, $Y'=Y''\otimes_{k''}k'$, $f'=f''_{(k')}:Y'\to X'$, et si $U'$ et $V'$ sont les images de $U''$ et $V''$ dans $X'$ et $Y'$ respectivement, la restriction $V'\to U'$ de $f'$ soit un isomorphisme. Comme $Y'$ est normal et $f'$ entier et birationnel, on conclut de (17.15.14.1) que $Y'$ est isomorphe au normalisé de $X'$, ce qui prouve le lemme puisque $Y'$ est géométriquement normal.

Revenons maintenant à la démonstration de (17.15.14), et appliquons à $k$ et $X$ le lemme (17.15.14.2); puisque $\dim(X)\leq1$, on a aussi $\dim(X_{(k')})\leq1$ (4.1.4), et le normalisé $Y'$ de $X_{(k')}$ est aussi de dimension $\leq1$ (5.4.2 et II, 7.4.6). Dire que $Y'$ est géométriquement normal sur $k'$ revient alors à dire, en vertu des définitions et de (II, 7.4.5), que $Y'$ est géométriquement régulier sur $k'$, donc lisse sur $k'$ (17.5.1), ce qui prouve (17.15.14).

\begin{proposition}[17.15.15]
\label{IV.17.15.15}
Soit $f:X\to Y$ un morphisme localement de présentation finie. Pour que $f$ soit lisse en un point $x\in X$, il faut et il suffit que $f$ soit plat au point $x$ et que $\Omega^1_f$ soit un $\sh{O}_X$-Module localement libre dans un voisinage de $x$, de rang en $x$ égal à $\dim_x f$.
\end{proposition}

La nécessité des conditions résulte de (17.5.1) et (17.10.2). Inversement, si ces conditions sont vérifiées, et si $y=f(x)$, il suffit de montrer (17.5.1) que $f^{-1}(y)$ est lisse sur $k(y)$ au point $x$; mais cela résulte de la définition de $\dim_x f$ (17.10.1), de (16.4.5) et de (17.15.5).

\subsection{Quasi-sections de morphismes plats ou lisses.}
\EGAReaderTarget{ega4.c04.subsec.17-16}{subsection:IV.17.16}

Les énoncés de ce numéro complètent ceux de (14.5), moyennant des hypothèses de platitude.

\begin{proposition}[17.16.1]
\label{IV.17.16.1}
Soit $f:X\to S$ un morphisme plat localement de présentation finie. Soient $s\in S$, $x$ un point fermé de $X_s$ tel que $\sh{O}_{X_s,x}$ soit un anneau de Cohen--Macaulay; alors il existe un voisinage ouvert $U$ de $s$ dans $S$ et un sous-préschéma $X'\subset f^{-1}(U)$ de $X$ tel que $x\in X'$ et que le morphisme $X'\to U$, restriction de $f$, soit plat, quasi-fini et de présentation finie.
\end{proposition}

La question étant locale sur $X$ et $Y$, on peut supposer $f$ de présentation finie. Soit $(t_i)_{1\leq i\leq m}$ un système de paramètres de l'anneau local $\sh{O}_{X_s,x}$; il existe un voisinage ouvert affine $V=\Spec(A)$ de $x$ dans $X$ et $m$ sections $g_i$ de $\sh{O}_X$ au-dessus de $V$ telles que

\oldpage[IV-4]{106}

les images des $(g_i)_x\in\sh{O}_{X,x}$ dans $\sh{O}_{X_s,x}=\sh{O}_{X,x}/\mathfrak m_s\sh{O}_{X,x}$ soient égales aux $t_i$. Soit $X'$ le sous-préschéma fermé de $V$ défini par l'idéal de $A$ engendré par les $g_i$; la suite $(t_i)$ étant par hypothèse régulière (0, 16.5.7) il résulte de (11.3.8) qu'en remplaçant au besoin $V$ par un voisinage ouvert plus petit, on peut supposer que le morphisme $X'\to S$, restriction de $f$, est plat et de présentation finie. D'autre part, puisque les $t_i$ forment un système de paramètres de $\sh{O}_{X_s,x}$, l'anneau $\sh{O}_{X'_s,x}$ est par définition artinien, et comme $x$ est fermé dans $X'_s$, on en conclut qu'il est isolé dans $X'_s$. En remplaçant $V$ au besoin par un voisinage plus petit de $x$, on en conclut, grâce à (13.1.4) que le morphisme $X'\to S$ est quasi-fini.

\begin{corollary}[17.16.2]
\label{IV.17.16.2}
Soit $f:X\to S$ un morphisme fidèlement plat et localement de présentation finie. Alors il existe un morphisme $g:S'\to S$, fidèlement plat, localement de présentation finie et localement quasi-fini, tel qu'il existe un $S$-morphisme $S'\to X$ (autrement dit, tel qu'il existe une $S'$-section de $X'=X\times_S S'$ (I, 3.3.14)). Si $S$ est quasi-compact (resp. quasi-compact et quasi-séparé), on peut supposer $S'$ affine (resp. $S'$ affine et le morphisme $g$ quasi-fini).
\end{corollary}

Pour tout $s\in S$, la fibre $X_s$ est non vide par hypothèse et est un préschéma localement de type fini sur $k(s)$; l'ensemble $U_s$ des points $x$ de $X_s$ où $\sh{O}_{X_s,x}$ est un anneau de Cohen--Macaulay est ouvert dans $X_s$ (6.11.3) et est non vide, puisqu'il contient les points maximaux de $X_s$ (0, 16.5.1); il contient par suite un point $x_s$ fermé dans $X_s$ (10.4.7). Soit $X'(s)$ un sous-préschéma affine de $X$ contenant $x_s$ et vérifiant les conditions de (17.16.1). Pour obtenir un préschéma $S'$ vérifiant les conditions de l'énoncé, il suffit de prendre la somme des $X'(s)$, où $s$ parcourt $S$. Puisque le morphisme $X'(s)\to S$ est plat et localement de présentation finie, l'image $U(s)$ de $X'(s)$ est ouverte dans $S$ (2.4.6); lorsque $S$ est quasi-compact, on peut donc recouvrir $S$ par un nombre fini de $U(s_i)$ et le préschéma $S'$ somme des $X'(s_i)$ répond encore à la question et est affine. Si de plus $S$ est quasi-séparé, on peut supposer les immersions ouvertes $U(s_i)\to S$ quasi-compactes (I, 2.7), donc de présentation finie (I, 6.2), et alors les morphismes $X'(s_i)\to S$ sont de présentation finie (I, 6.2) et il en est par suite de même du morphisme $S'\to S$.

\begin{corollary}[17.16.3]
\label{IV.17.16.3}
\begin{enumerate}
  \item[(i)] Soit $f:X\to S$ un morphisme lisse. Soient $s\in S$, $x$ un point fermé de $X_s$ tel que le corps résiduel $k(x)$ soit séparable sur $k(s)$; alors, dans la conclusion de (17.16.1), on peut prendre $X'$ tel que le morphisme $X'\to U$, restriction de $f$, soit étale.
  \item[(ii)] Soit $f:X\to S$ un morphisme lisse surjectif. Alors, dans les conclusions de (17.16.2), on peut supposer en outre que $g:S'\to S$ soit étale.
\end{enumerate}
\end{corollary}

Il est clair que (ii) se déduit de (i) comme (17.16.2) de (17.16.1), compte tenu de ce que, pour tout $s\in S$, comme $X_s$ est non vide et lisse sur $k(s)$, il existe un point fermé $x\in X_s$ tel que $k(x)$ soit séparable sur $k(s)$ (17.15.10, (iii)). Il suffit donc de prouver (i). On notera que l'anneau $\sh{O}_{X_s,x}$ est alors régulier; si l'on répète la construction faite dans (17.16.1) en prenant pour $(t_i)$ un système régulier de paramètres de $\sh{O}_{X_s,x}$, l'anneau $\sh{O}_{X'_s,x}$ est un corps isomorphe à $k(x)$, donc séparable sur $k(s)$ par hypothèse. La conclusion résulte alors de (17.6.1, $c'$)).

\begin{proposition}[17.16.4]
\label{IV.17.16.4}
Soient $S$ un préschéma quasi-compact et quasi-séparé, $f:X\to S$ un morphisme surjectif localement de présentation finie. Alors il existe une famille finie $(S_i)_{i\in I}$

\oldpage[IV-4]{107}

de sous-préschémas affines de $S$, de présentation finie sur $S$, deux à deux disjoints, de réunion $S$, et ayant la propriété suivante: pour chaque $i\in I$, il existe deux morphismes finis de présentation finie et surjectifs
\[
  S_i''\xrightarrow{g_i}S_i'\xrightarrow{h_i}S_i,
\]
où $h_i$ est étale et $g_i$ plat et radiciel, et un $S$-morphisme $S_i''\to X$ (autrement dit une $S_i''$-section de $X_i''=X\times_S S_i''$).
\end{proposition}

Il y a un recouvrement fini $(U_i)_{1\leq i\leq n}$ de $S$ par des ouverts affines; pour tout $i$, l'intersection
\[
  W_i=U_i\cap\left(\bigcup_{j<i}U_j\right)
\]
est quasi-compacte (I, 2.7) donc il existe un sous-préschéma fermé $V_i$ de $U_i$ ayant pour espace sous-jacent $U_i-W_i$ et défini par un idéal de type fini de l'anneau de $U_i$; donc $V_i$ est un sous-préschéma affine de $S$ qui est de présentation finie sur $S$ (I, 6.2). Il est clair qu'on peut se borner à prouver la proposition en y remplaçant $S$ par $V_i$ et $X$ par $X\times_S V_i$. Autrement dit, on peut déjà supposer $S$ affine. D'autre part, $X$ est réunion d'ouverts affines $W_\alpha$ et les $f(W_\alpha)$ sont constructibles dans $S$ (I, 8.4) et forment un recouvrement de $S$; par suite (I, 9.9) il existe une sous-famille finie $(W_{\alpha_j})_{1\leq j\leq m}$ telle que les $f(W_{\alpha_j})$ forment déjà un recouvrement de $S$. Soit alors $X'$ le préschéma somme des sous-préschémas induits sur les ouverts $W_{\alpha_j}$ de $X$; il est immédiat qu'il suffit de prouver la proposition en remplaçant $X$ par $X'$ puisqu'un $S$-morphisme $S_i''\to X'$ donne par composition un $S$-morphisme $S_i''\to X'\to X$.

On peut donc supposer $X$ affine et $f$ de présentation finie. On peut alors (8.9.1) écrire $X$ sous la forme $X_0\times_{S_0}S$, où $S_0$ est noethérien, $X_0\to S_0$ un morphisme de type fini, qu'on peut supposer surjectif (8.10.5, (vi)). Si la proposition est prouvée pour ce morphisme, elle en résultera aussitôt pour $X$ par changement de base $S\to S_0$. On peut donc supposer de plus $S$ noethérien et $f$ de type fini.

Soient $s_h$ $(1\leq h\leq r)$ les points maximaux de $S$. Montrons qu'il suffit de prouver l'énoncé en remplaçant $S$ par un voisinage ouvert assez petit $U_h$ de chacun des $s_h$, $X$ par l'image réciproque de $U_h$ dans $X$. En effet, supposons la proposition établie dans ce cas, et raisonnons par récurrence noethérienne ($0_{\mathrm I}$, 2.2.2) en supposant l'énoncé établi pour tout sous-préschéma fermé de $S$ ayant un espace sous-jacent $\neq S$. On peut supposer les $U_h$ deux à deux sans point commun; si $T$ est un sous-préschéma fermé ayant pour espace sous-jacent $S-(\bigcup_h U_h)$, l'hypothèse de récurrence entraîne que l'énoncé est vrai pour $T$; comme il est vrai aussi pour chacun des $U_h$, il l'est évidemment pour $S$.

Comme on peut évidemment (en remplaçant $S$ par $S_{\mathrm{red}}$ et $X$ par $X\times_S S_{\mathrm{red}}$) supposer $S$ réduit, chacun des $\sh{O}_{S,s_h}$ est un corps $k(s_h)$. Supposons que l'on ait prouvé la proposition lorsque $S$ est le spectre d'un corps et $f$ est de type fini. Alors l'existence des $U_h$ résulte de la méthode de (8.1.2, $a$)) et de (8.8.2, (i) et (ii)), (8.10.5, (vi), (vii) et (x)), (11.2.6, (ii)) et (17.7.8, (ii)).

Supposons donc $S=\Spec(k)$, où $k$ est un corps, $X$ étant de type fini sur $k$ et $X\neq\varnothing$. Comme $X$ est noethérien, il existe dans $X$ un point fermé $x$ ($0_{\mathrm I}$, 2.1.3), donc $k(x)$ est une extension finie de $k$ (I, 6.4.2). Il y a par suite une extension finie séparable $k'$ de $k$ telle que $k''=k(x)$ soit extension finie radicielle de $k'$. On répond alors à la question en prenant $I$ réduit à un élément, $S_i'=\Spec(k')$, $S_i''=\Spec(k'')$.

\oldpage[IV-4]{108}

\begin{corollary}[17.16.5]
\label{IV.17.16.5}
Soit $f:X\to S$ un morphisme surjectif localement de présentation finie. Alors il existe un morphisme $g:S'\to S$ surjectif, localement de présentation finie et localement quasi-fini, tel qu'il existe un $S$-morphisme $S'\to X$ (autrement dit tel qu'il existe une $S'$-section de $X'=X\times_S S'$). Si $S$ est quasi-compact (resp. quasi-compact et quasi-séparé), on peut supposer $S'$ affine (resp. $S'$ affine et le morphisme $g$ de présentation finie et quasi-fini).
\end{corollary}

Il suffit de prouver le corollaire en supposant $S$ affine: $S$ est en effet réunion d'une famille $(S_\alpha)$ d'ouverts affines, et si pour chaque $\alpha$, $S'_\alpha$ est affine et le morphisme $g_\alpha:S'_\alpha\to S_\alpha$ répond à la question et est de présentation finie et quasi-fini, alors en prenant pour $S'$ le préschéma somme des $S'_\alpha$, $g:S'\to S$ coïncidant avec $g_\alpha$ dans chacun des $S'_\alpha$, ce morphisme répond à la question; en outre, si $S$ est quasi-compact, on peut supposer la famille $(S_\alpha)$ finie, donc $S'$ affine; si de plus $S$ est quasi-séparé, les immersions $S_\alpha\to S$ sont de présentation finie (I, 6.2), donc il en est de même de $g$.

Considérons donc le cas où $S$ est affine: on forme alors les morphismes finis et de présentation finie $S_i''\to S_i$ de (17.16.4); comme les immersions $S_i\to S$ sont de présentation finie (les $S_i$ étant affines), les morphismes $g_i'':S_i''\to S_i\to S$ sont de présentation finie et quasi-finis, et on répond à la question en prenant $S'$ égal à la somme des $S_i''$ et $g$ coïncidant avec $g_i''$ dans chacun des $S_i''$.

\begin{proposition}[17.16.6]
\label{IV.17.16.6}
Soient $S$ un préschéma quasi-compact et quasi-séparé, $f:X\to S$ un morphisme de présentation finie tel que, pour tout $s\in S$, $X_s=f^{-1}(s)$ soit un préschéma propre sur $k(s)$. Alors il existe une famille finie $(S_i)_{i\in I}$ de sous-préschémas affines de $S$, de présentation finie sur $S$, deux à deux disjoints, de réunion $S$ et tels que, pour tout $i$, le morphisme $X_i=X\times_S S_i\to S_i$ soit propre et plat. Si de plus, pour tout $s\in S$, $X_s$ est fini sur $k(s)$, on peut prendre les $S_i$ tels que chacun des morphismes $X_i\to S_i$ se factorise en
\[
  X_i\xrightarrow{u_i}X_i'\xrightarrow{h_i}S_i
\]
où $u_i$ et $h_i$ sont finis et localement libres (18.2.7), $h_i$ est étale, $u_i$ est radiciel et surjectif.
\end{proposition}

On suit une marche analogue à celle de (17.16.4). On se ramène d'abord au cas où $S$ est affine et où $f$ est plat, en utilisant le théorème de platitude générique (8.9.5) (on observera que les sous-préschémas $S_i'$ définis dans la démonstration de (8.9.5) sont affines). Puisque $f$ est de présentation finie, on peut ensuite écrire $X=X_0\times_{S_0}S$, où $S_0$ est noethérien, $X_0\to S_0$ un morphisme de type fini et plat (11.2.7); en outre chaque fibre $(X_0)_{s_0}$ est encore propre sur $k(s_0)$, comme il résulte de (2.7.1, (vii)) et on est donc ramené au cas où $S$ est noethérien. Par récurrence noethérienne et application du procédé de (8.1.2, $a$)) (en utilisant aussi cette fois (8.10.5, (xii))), on est finalement ramené, pour prouver la première assertion, au cas où $S$ est le spectre d'un corps, qui résulte trivialement de l'hypothèse (avec $S_i=S$). On traite de même le cas où $X_s$ est supposé fini sur $k(s)$ pour tout $s\in S$ (il faut cette fois utiliser (2.7.1, (xv)), (8.10.5, (x), (iv), (vi) et (vii)), (17.7.8) et (2.1.12)), et on est ramené à prouver la dernière assertion de l'énoncé lorsque $S=\Spec(k)$ est le spectre d'un corps. Mais alors $X=\Spec(A)$, où $A$ est une $k$-algèbre finie (I, 6.4.4); donc $A$ est composée directe de $k$-algèbres finies $A_j$ qui sont des anneaux locaux $(1\leq j\leq m)$; puisque $A_j$ est un anneau artinien et une

\oldpage[IV-4]{109}

$k$-algèbre, il contient un sous-corps $K'_j$ isomorphe canoniquement au corps résiduel de $A_j$ (0, 19.6.2). On prend alors $X_i=X$, $X'_i$ somme des $\Spec(K'_j)$ et il est clair que l'on répond ainsi à la question puisque pour tout $j$ on a deux homomorphismes $K'_j\to A_j\to k(A_j)$ dont le composé est un isomorphisme.

% FIN SOURCE ega4/ega4-17-fr.tex

%
% DEBUT SOURCE ega4/ega4-18-fr.tex
\section[Compléments sur les morphismes étales. Anneaux locaux henséliens et anneaux strictement locaux]{Compléments sur les morphismes étales\\Anneaux locaux henséliens et anneaux strictement locaux}
\label{section:IV.18}

Dans le présent paragraphe, nous étudions diverses propriétés spéciales aux morphismes étales. De plus, la notion de morphisme étale permet de développer de façon fort naturelle la théorie de Nagata des anneaux henséliens, ainsi que celle des anneaux strictement locaux. Ces anneaux jouent un rôle important dans de nombreux développements récents, en s'introduisant chaque fois que l'on a besoin d'un procédé de « localisation » plus fin que celui fourni par la topologie de Zariski (cf. par exemple [43], en attendant la parution du chapitre de notre Traité consacré à l'étude de la « topologie étale »).

\subsection{Une équivalence remarquable de catégories.}
\EGAReaderTarget{ega4.c04.subsec.18-1}{subsection:IV.18.1}

\begin{proposition}[18.1.1]
\label{IV.18.1.1}
Soient $S$ un préschéma, $S_0$ un sous-préschéma fermé de $S$, $X_0$ un $S_0$-préschéma lisse (resp. étale) sur $S_0$, $x_0$ un point de $X_0$. Il existe alors un voisinage ouvert $U_0$ de $x_0$ dans $X_0$, un $S$-préschéma $U$ lisse (resp. étale) sur $S$ et un $S_0$-isomorphisme
\[
  U\times_S S_0\xrightarrow{\sim}U_0.
\]
\end{proposition}

Notons que si $X_0$ est étale sur $S_0$ au point $x_0$, il est a fortiori non ramifié sur $S_0$ en ce point; si l'on a construit un $S$-préschéma $U$ lisse sur $S$ tel que $U\times_S S_0$ soit isomorphe à $U_0$, comme les fibres des morphismes $U_0\to S_0$ et $U\to S$ contenant $x_0$ sont alors isomorphes, il en résultera que $U$ est non ramifié sur $S$ au point $x_0$ (17.4.1, $d$)), donc aussi dans un voisinage de $x_0$; en remplaçant $U$ par ce voisinage, on en conclut que $U$ sera étale sur $S$. Il suffit donc de prouver la proposition lorsqu'on suppose seulement $X_0$ lisse sur $S_0$.

La question étant locale sur $S$ et sur $X_0$, on peut supposer que $S=\Spec(A)$ et $X_0=\Spec(C_0)$ sont affines, de sorte que $S_0=\Spec(A_0)$, où $A_0$ est un anneau quotient de $A$, $C_0=B_0/\mathfrak J_0$, où $B_0=A_0[T_1,\ldots,T_n]$ et $\mathfrak J_0$ est un idéal de type fini de $B_0$; enfin, $C_0$ est une $A_0$-algèbre formellement lisse pour les topologies discrètes. Soit $\mathfrak p_0$ l'idéal $\mathfrak j_{x_0}$ dans $C_0$; on a $\mathfrak p_0=\mathfrak q_0/\mathfrak J_0$, où $\mathfrak q_0$ est un idéal premier de $B_0$. Le critère jacobien (0, 22.6.4) joint à (0, 19.1.12) montre qu'il existe dans $\mathfrak J_0$ une famille de $r$ polynômes $u_i$ $(1\leq i\leq r)$ et $r$ indices $j_h$ $(1\leq h\leq r)$ tels que les images des $u_i$ dans $(\mathfrak J_0)_{\mathfrak q_0}/(\mathfrak J_0^2)_{\mathfrak q_0}$ engendrent ce $(B_0)_{\mathfrak q_0}$-module et que l'on ait
\begin{equation}
  \det\left(\frac{\partial u_i}{\partial T_{j_h}}\right)\notin\mathfrak q_0.
  \tag{18.1.1.1}\label{IV.18.1.1.1}
\end{equation}

Comme $(B_0)_{\mathfrak q_0}$ est un anneau local, il résulte du lemme de Nakayama que l'on peut supposer que les images des $u_i$ dans $(\mathfrak J_0)_{\mathfrak q_0}$ engendrent ce $(B_0)_{\mathfrak q_0}$-module, puis, en remplaçant au besoin $X_0$ par un voisinage ouvert affine de $x_0$, que les $u_i$ engendrent $\mathfrak J_0$ ($0_{\mathrm I}$, 5.2.2). Posons alors $B=A[T_1,\ldots,T_n]$; $B_0$ est donc un anneau quotient

\oldpage[IV-4]{110}

de $B$, $\mathfrak q_0$ est l'image d'un idéal premier $\mathfrak q$ de $B$ et $\mathfrak q$ l'image réciproque de $\mathfrak q_0$. Pour tout $i$, soit $v_i\in B$ un élément dont $u_i$ est l'image dans $B_0$, et soit $\mathfrak J$ l'idéal de $B$ engendré par les $v_i$, de sorte que $\mathfrak J_0$ est l'image de $\mathfrak J$ dans $B_0$. La proposition sera établie en prenant pour $U$ un voisinage ouvert du point de $\Spec(B/\mathfrak J)$ correspondant à l'idéal premier $\mathfrak p=\mathfrak q/\mathfrak J$, pourvu que l'on prouve que $B_{\mathfrak q}/\mathfrak J_{\mathfrak q}$ est une $A$-algèbre formellement lisse pour les topologies discrètes. Or, cela résulte du critère jacobien, car les images des $v_i$ dans $\mathfrak J_{\mathfrak q}/\mathfrak J_{\mathfrak q}^2$ engendrent ce $B_{\mathfrak q}$-module, et il résulte de (18.1.1.1) que l'on a $\det(\partial v_i/\partial T_{j_h})\notin\mathfrak q$.

\begin{theorem}[18.1.2]
\label{IV.18.1.2}
Soient $S$ un préschéma, $S_0$ un sous-préschéma fermé de $S$ dont l'espace sous-jacent est identique à celui de $S$. Alors le foncteur
\[
  X\longmapsto X\times_S S_0
\]
de la catégorie des $S$-préschémas étales sur $S$, dans la catégorie des $S_0$-préschémas étales sur $S_0$, est une équivalence de catégories.
\end{theorem}

Montrons d'abord que ce foncteur est pleinement fidèle. Soient $X$, $Y$ deux $S$-préschémas étales sur $S$, et posons $X_0=X\times_S S_0$, $Y_0=Y\times_S S_0$. Si $Z=X\times_S Y$, l'ensemble $\Hom_S(X,Y)$ est en correspondance biunivoque canonique avec l'ensemble des $X$-sections $\Gamma(Z/X)$, et de même $\Hom_{S_0}(X_0,Y_0)$ est en correspondance biunivoque canonique avec $\Gamma(Z_0/X_0)$, où $Z_0=Z\times_S S_0=X_0\times_{S_0}Y_0$. Or $Z$ est étale sur $X$, $Z_0$ étale sur $X_0$, et $X_0$ (resp. $Z_0$) est un sous-préschéma fermé de $X$ (resp. $Z$) ayant même espace sous-jacent. Les parties ouvertes de $Z$ telles que la restriction du morphisme $Z\to X$ soit surjective et radicielle sont donc les mêmes que les parties de $Z_0$ ayant les propriétés correspondantes et notre assertion découle par suite de (17.9.3).

Pour achever la démonstration, il suffit de voir que pour tout $S_0$-préschéma $X_0$ étale sur $S_0$, il existe un $S$-préschéma $X$ étale sur $S$ et un $S_0$-isomorphisme $X_0\xrightarrow{\sim}X\times_S S_0$. En vertu de la prop. (18.1.1), il existe un recouvrement ouvert $(U_\alpha)$ de $X_0$ et pour tout $\alpha$, un $S$-préschéma $V_\alpha$ qui est étale sur $S$, et enfin un $S_0$-isomorphisme $\theta_\alpha:U_\alpha\xrightarrow{\sim}V_\alpha\times_S S_0$. En outre, en raison de la première partie de la démonstration, il existe un unique $S_0$-isomorphisme $\varphi_{\alpha\beta}$ de $\theta_\alpha(U_\alpha\cap U_\beta)$ sur $\theta_\beta(U_\alpha\cap U_\beta)$, correspondant à l'automorphisme identique de $U_\alpha\cap U_\beta$, et il est immédiat, pour la même raison, que ces isomorphismes vérifient la condition de recollement ($0_{\mathrm I}$, 4.1.7). Il y a par suite un $S$-préschéma $X$ tel que les $V_\alpha$ s'identifient canoniquement à des sous-préschémas induits sur des ouverts de $X$, les $\theta_\alpha$ s'identifiant à des $S_0$-isomorphismes qui coïncident dans les intersections $U_\alpha\cap U_\beta$, et définissent donc un $S_0$-isomorphisme $X_0\xrightarrow{\sim}X\times_S S_0$. Il est clair que $X$ est étale sur $S$ (17.3.2), ce qui achève la démonstration.

\begin{corollary}[18.1.3]
\label{IV.18.1.3}
Soient $S$ un préschéma, $X$ un $S$-préschéma étale sur $S$, $S'$ un $S$-préschéma, $S'_0$ un sous-préschéma fermé de $S'$ ayant même espace sous-jacent. Alors l'application canonique $X(S')_S\to X(S'_0)_S$ (I, 3.4.3) est bijective.
\end{corollary}

Posons en effet $X'=X\times_S S'$, $X'_0=X\times_S S'_0$, de sorte que $X(S')_S=\Gamma(X'/S')$ et $X(S'_0)_S=\Gamma(X'_0/S'_0)$; le corollaire résulte du fait que le foncteur défini dans (18.1.2) (avec $S$ et $S_0$ remplacés par $S'$ et $S'_0$ respectivement) est pleinement fidèle (ou directement de (17.9.3)).

\subsection{Revêtements étales.}
\EGAReaderTarget{ega4.c04.subsec.18-2}{subsection:IV.18.2}

\oldpage[IV-4]{111}

\begin{env}[18.2.1]
\EGAReaderTarget{ega4.c04.statement.18-2-1}{IV.18.2.1}
Étant donnés un anneau $A$ et une $A$-algèbre commutative $B$ qui est \emph{finie} et est un $A$-module \emph{libre}, rappelons (Bourbaki, \emph{Alg.}, chap. VIII, \textsection~12, n\textsuperscript{o}~2) que l'on définit sur $B$ une \emph{forme $A$-linéaire} $\operatorname{Tr}_{B/A}$, la « forme trace »; on en déduit la définition d'une \emph{forme $A$-bilinéaire symétrique} (dite encore « forme trace »)
\begin{equation}
  (x,y)\longmapsto\operatorname{Tr}_{B/A}(xy)
  \tag{18.2.1.1}\label{IV.18.2.1.1}
\end{equation}
dont la donnée est équivalente à celle de l'application $A$-linéaire associée $\operatorname{astr}_{B/A}:B\to\check B$ du $A$-module $B$ dans son \emph{dual} $\check B$, égale à sa transposée. Lorsque $A$ est un \emph{corps}, il est équivalent de dire que cette forme bilinéaire est \emph{non dégénérée} ou que $B$ est une $A$-algèbre \emph{séparable} (Bourbaki, \emph{Alg.}, chap. IX, \textsection~2, prop. 5).

Soit $f:A\to A'$ un homomorphisme d'anneaux; posons $B'=B\otimes_A A'$, et soit $g:B\to B'$ l'homomorphisme canonique; l'image par $g$ d'une base du $A$-module $B$ est alors une base du $A'$-module $B'$, et il résulte des définitions que l'on a, pour tout $x\in B$,
\begin{equation}
  \operatorname{Tr}_{B'/A'}(g(x))=f(\operatorname{Tr}_{B/A}(x)).
  \tag{18.2.1.2}\label{IV.18.2.1.2}
\end{equation}
\end{env}

\begin{env}[18.2.2]
\EGAReaderTarget{ega4.c04.statement.18-2-2}{IV.18.2.2}
Considérons maintenant un \emph{espace annelé} $(X,\sh{O}_X)$ et soit $\sh{B}$ une $\sh{O}_X$-Algèbre qui, en tant que $\sh{O}_X$-Module, est \emph{localement libre de rang fini}; alors, pour tout ouvert $U\subset X$ tel que $\sh{B}|U$ soit (en tant que $\sh{O}_U$-Module) isomorphe à $\sh{O}_U^n$ (pour un $n$ dépendant de $U$), $\Gamma(U,\sh{B})$ est une $\Gamma(U,\sh{O}_X)$-algèbre qui, en tant que $\Gamma(U,\sh{O}_X)$-module, est libre de rang fini et définit donc une forme $\Gamma(U,\sh{O}_X)$-linéaire $\operatorname{Tr}_{\Gamma(U,\sh{B})/\Gamma(U,\sh{O}_X)}$, que nous noterons aussi $\operatorname{Tr}_{\sh{B}|\sh{O}_X,U}$; on en déduit une application linéaire associée
\[
  \operatorname{astr}_{\sh{B}|\sh{O}_X,U}:\Gamma(U,\sh{B})\longrightarrow\Gamma(U,\sh{B})^\vee=\Gamma(U,\check{\sh{B}}).
\]
En outre, il résulte de (18.2.1.2) que ces applications linéaires sont compatibles avec les opérations de restriction de $U$ à un ouvert plus petit, et définissent donc, d'une part un \emph{homomorphisme de $\sh{O}_X$-Modules}, appelé encore \emph{homomorphisme trace}:
\begin{equation}
  \operatorname{Tr}_{\sh{B}|\sh{O}_X}:\sh{B}\longrightarrow\sh{O}_X
  \tag{18.2.2.1}\label{IV.18.2.2.1}
\end{equation}
et d'autre part un \emph{homomorphisme de $\sh{O}_X$-Modules}
\begin{equation}
  \operatorname{astr}_{\sh{B}|\sh{O}_X}:\sh{B}\longrightarrow\check{\sh{B}}
  \tag{18.2.2.2}\label{IV.18.2.2.2}
\end{equation}
dit \emph{associé à la trace}, et égal à son transposé. Il résulte aussi de (18.2.1.2) que pour tout $x\in X$, on a
\begin{equation}
  (\operatorname{Tr}_{\sh{B}|\sh{O}_X})_x=\operatorname{Tr}_{\sh{B}_x|\sh{O}_{X,x}}
  \tag{18.2.2.3}\label{IV.18.2.2.3}
\end{equation}
\begin{equation}
  (\operatorname{astr}_{\sh{B}|\sh{O}_X})_x=\operatorname{astr}_{\sh{B}_x|\sh{O}_{X,x}}.
  \tag{18.2.2.4}\label{IV.18.2.2.4}
\end{equation}
Enfin, sous les conditions de (18.2.1), si l'on pose $X=\Spec(A)$ et si $\sh{B}=\widetilde B$ est la $\sh{O}_X$-Algèbre correspondant à $B$, la forme $\operatorname{Tr}_{\sh{B}|\sh{O}_X}$ (resp. l'homomorphisme $\operatorname{astr}_{\sh{B}|\sh{O}_X}$) correspond à la forme $\operatorname{Tr}_{B/A}$ (resp. à l'homomorphisme de $A$-modules $\operatorname{astr}_{B/A}$), comme il résulte encore de (18.2.1.2).
\end{env}


\oldpage[IV-4]{112}

\begin{proposition}[18.2.3]
\EGAReaderTarget{ega4.c04.prop.18-2-3}{IV.18.2.3}
Soit $f:X\to Y$ un morphisme fini de préschémas et soit $\sh{B}=f_*(\sh{O}_X)$. Les conditions suivantes sont équivalentes:
\begin{enumerate}
  \item[a)] $f$ est étale.
  \item[a$'$)] $f$ est un morphisme plat de présentation finie, et pour tout $x\in X$, si l'on pose $y=f(x)$, $\sh{O}_{X,x}/\mathfrak m_y\sh{O}_{X,x}$ est un corps, extension finie séparable de $\kappa(y)$.
  \item[b)] $\sh{B}$ est un $\sh{O}_Y$-Module localement libre et pour tout $y\in Y$, $\sh{B}_y\otimes_{\sh{O}_{Y,y}}\kappa(y)$ est une $\kappa(y)$-algèbre finie séparable (donc le composé direct d'un nombre fini de corps, extensions finies séparables de $\kappa(y)$).
  \item[c)] $\sh{B}$ est un $\sh{O}_Y$-Module localement libre, et l'homomorphisme $\operatorname{astr}_{\sh{B}|\sh{O}_Y}:\sh{B}\to\check{\sh{B}}$ (18.2.2) est bijectif.
\end{enumerate}
\end{proposition}

Compte tenu de ce que $f$ est quasi-compact, l'équivalence de $a)$ et $a')$ a déjà été démontrée (17.6.2). Pour prouver le reste de la proposition, on peut se borner au cas où $Y=\Spec(A)$ et $X=\Spec(B)$ sont affines, $B$ étant une $A$-algèbre finie et $\sh{B}=\widetilde B$. Dire que $f$ est un morphisme de présentation finie équivaut alors à dire que $B$ est un $A$-module de présentation finie (1.4.7). Si en outre $f$ est plat, donc $B$ un $A$-module plat, on sait (Bourbaki, \emph{Alg. comm.}, chap. II, \textsection~5, n\textsuperscript{o}~2, cor. 2 du th. 1) que $B$ est un $A$-module projectif, donc $\sh{B}$ un $\sh{O}_Y$-Module localement libre (\emph{loc. cit.}, n\textsuperscript{o}~2, th. 1), et la réciproque est immédiate. D'autre part $f^{-1}(y)$ n'est autre que le spectre de la $\kappa(y)$-algèbre $\sh{B}(y)=\sh{B}_y\otimes_{\sh{O}_{Y,y}}\kappa(y)$, ce qui achève de prouver l'équivalence de $a')$ et $b)$. Pour voir que $b)$ est équivalente à $c)$, notons que la seconde assertion de $b)$ équivaut au fait que l'homomorphisme $\operatorname{astr}_{\sh{B}(y)|\kappa(y)}:\sh{B}(y)\to\check{\sh{B}(y)}$ est bijectif; comme $\sh{B}(y)=\sh{B}_y/\mathfrak m_y\sh{B}_y$ et $\check{\sh{B}(y)}=\check{\sh{B}}_y/\mathfrak m_y\check{\sh{B}}_y$ et que $\sh{B}_y$ et $\check{\sh{B}}_y$ sont des $\sh{O}_{Y,y}$-modules libres, il résulte de (18.2.2.4) et de Bourbaki, \emph{Alg. comm.}, chap. II, \textsection~3, n\textsuperscript{o}~3, cor. de la prop. 6, que l'homomorphisme $\operatorname{astr}_{\sh{B}_y|\sh{O}_{Y,y}}:\sh{B}_y\to\check{\sh{B}}_y$ est lui aussi bijectif; la réciproque étant évidente, cela achève la démonstration.

Lorsqu'une $\sh{O}_X$-Algèbre $\sh{B}$ vérifie les conditions équivalentes $b)$ et $c)$ de (18.2.3), on dit que $\sh{B}$ est une $\sh{O}_X$-Algèbre \emph{finie et étale}. Lorsque $X=\Spec(A)$ est affine et que l'on a donc $\sh{B}=\widetilde B$, où $B$ est une $A$-algèbre, il revient au même, en vertu de (18.2.3), de dire que $\sh{B}$ est une $\sh{O}_X$-Algèbre finie et étale, ou que $B$ est une $A$-algèbre finie et étale (au sens de (17.3.2)).

\begin{corollary}[18.2.4]
\label{IV.18.2.4}
Soit $f:X\to Y$ un morphisme fini et de présentation finie, et posons $\sh{B}=f_*(\sh{O}_X)$. Soit $y$ un point de $Y$. Les conditions suivantes sont équivalentes:
\begin{enumerate}
  \item[a)] Il existe un voisinage ouvert $U$ de $y$ dans $Y$ tel que la restriction $f^{-1}(U)\to U$ de $f$ soit un morphisme étale.
  \item[b)] $\sh{B}_y$ est un $\sh{O}_{Y,y}$-module libre de type fini et $\sh{B}_y\otimes_{\sh{O}_{Y,y}}\kappa(y)$ est une $\kappa(y)$-algèbre séparable.
\end{enumerate}
\end{corollary}

Il est clair que $a)$ entraîne $b)$ en vertu de (18.2.3). D'autre part, $\sh{B}$ est un $\sh{O}_Y$-Module de présentation finie (1.6.3 et 1.4.7), donc, si $\sh{B}_y$ est un $\sh{O}_{Y,y}$-module libre, il existe un voisinage ouvert $U$ de $y$ dans $Y$ tel que $\sh{B}|U$ soit un $\sh{O}_U$-Module localement libre ($0_{\mathrm I}$, 5.2.7); en outre, par hypothèse, l'homomorphisme $\operatorname{astr}_{\sh{B}_y|\sh{O}_{Y,y}}:\sh{B}_y\to\check{\sh{B}}_y$ étant bijectif, il résulte aussi de ($0_{\mathrm I}$, 5.2.7) que l'on peut supposer $U$ choisi de sorte que

\oldpage[IV-4]{113}

l'homomorphisme $\operatorname{astr}_{\sh{B}|U/\sh{O}_U}$ soit bijectif. Le fait que $b)$ implique $a)$ résulte alors de (18.2.3).

\begin{corollary}[18.2.5]
\label{IV.18.2.5}
Soient $Y$ un préschéma quasi-compact ou localement noethérien, $f:X\to Y$ un morphisme fini et de présentation finie; posons $\sh{B}=f_*(\sh{O}_X)$. Supposons que pour tout point fermé $y$ de $Y$, $\sh{B}_y$ soit un $\sh{O}_{Y,y}$-module libre et $\sh{B}_y\otimes_{\sh{O}_{Y,y}}\kappa(y)$ une $\kappa(y)$-algèbre séparable. Alors $f$ est étale.
\end{corollary}

En effet, il résulte de (18.2.4) que tout point fermé $y$ de $Y$ possède un voisinage ouvert $U$ tel que la restriction $f^{-1}(U)\to U$ de $f$ soit étale; la conclusion résulte de ce que dans les deux cas considérés, toute partie fermée non vide de $Y$ contient un point fermé (5.1.11 et $0_{\mathrm I}$, 2.1.3).

\begin{corollary}[18.2.6]
\label{IV.18.2.6}
Si $f:X\to Y$ est un morphisme fini et étale, et si l'on pose $\sh{B}=f_*(\sh{O}_X)$, alors le $\sh{O}_Y$-homomorphisme $\operatorname{Tr}_{\sh{B}|\sh{O}_Y}:\sh{B}\to\sh{O}_Y$ (18.2.2) (aussi noté $\operatorname{Tr}_f$) est surjectif.
\end{corollary}

La question étant locale, on peut, en vertu de (18.2.3) supposer que $Y=\Spec(A)$, $X=\Spec(B)$ avec $\sh{B}=\widetilde B$, $B$ étant un $A$-module libre; comme en vertu de (18.2.3) la forme bilinéaire (18.2.1.1) est non dégénérée, cela entraîne en particulier que la forme linéaire $\operatorname{Tr}_{B/A}$ est surjective.

\begin{remarks}[18.2.7]
\EGAReaderTarget{ega4.c04.item.18-2-7}{IV.18.2.7}
\begin{enumerate}
  \item[(i)] Lorsque $f:X\to Y$ est un morphisme fini tel que $f_*(\sh{O}_X)$ soit un $\sh{O}_Y$-Module localement libre (resp. localement libre de rang $n$), on dit encore que $f$ est un morphisme \emph{fini localement libre} (resp. \emph{localement libre de rang $n$}). Cette condition, en vertu de (18.2.3), est vérifiée si $f$ est un morphisme fini étale, mais n'implique pas à elle seule que $f$ soit étale, comme le montre l'exemple où $X=\Spec(K)$ et $Y=\Spec(k)$ sont des spectres de corps, $K$ étant une extension finie non séparable de $k$. Lorsque $f:X\to Y$ est un morphisme fini étale, on dit encore que $X$ est un \emph{revêtement étale} de $Y$.

On notera que dans ce cas, $f$ est universellement ouvert et universellement fermé, et en particulier $f(X)$ est une partie à la fois ouverte et fermée de $Y$.

On dit qu'un revêtement étale $X$ de $Y$ est \emph{trivial} si $X$ est \emph{somme} d'un nombre fini de préschémas \emph{isomorphes} à $Y$. On dit qu'un revêtement étale $X$ de $Y$ est \emph{localement trivial} si le morphisme $f:X\to Y$ est tel que tout point $y\in Y$ ait un voisinage ouvert $U$ pour lequel le revêtement $f^{-1}(U)$ de $U$ soit trivial.

  \item[(ii)] Soit $f:X\to Y$ un morphisme fini, localement libre de rang $n$; posons $\sh{B}=f_*(\sh{O}_X)$ et soit $u=\operatorname{astr}_{\sh{B}|\sh{O}_Y}:\sh{B}\to\check{\sh{B}}$; on en déduit un homomorphisme puissance extérieure $n$-ème
\[
  \bigwedge^n u:\bigwedge^n\sh{B}\longrightarrow\bigwedge^n\check{\sh{B}}=(\bigwedge^n\sh{B})^\vee
\]
de $\sh{O}_Y$-Modules \emph{inversibles}, et par suite ($0_{\mathrm I}$, 5.4.2) un élément
\begin{equation}
  d_{X/Y}\in\Gamma\left(Y,(\bigwedge^n\check{\sh{B}})\otimes_{\sh{O}_Y}(\bigwedge^n\check{\sh{B}})\right)
  \tag{18.2.7.1}\label{IV.18.2.7.1}
\end{equation}
que l'on appelle le \emph{discriminant de $X$ sur $Y$}. D'ailleurs, comme $(\bigwedge^n\check{\sh{B}})\otimes_{\sh{O}_Y}(\bigwedge^n\check{\sh{B}})$ est le dual de $(\bigwedge^n\sh{B})\otimes_{\sh{O}_Y}(\bigwedge^n\sh{B})$, $d_{X/Y}$ peut aussi être identifié à un homomorphisme
\begin{equation}
  (\bigwedge^n\sh{B})\otimes_{\sh{O}_Y}(\bigwedge^n\sh{B})\longrightarrow\sh{O}_Y
  \tag{18.2.7.2}\label{IV.18.2.7.2}
\end{equation}
et on note $\mathscr{D}_{X/Y}$ l'Idéal quasi-cohérent de type fini de $\sh{O}_Y$, \emph{image} de l'homomorphisme (18.2.7.2), encore appelé \emph{Idéal discriminant de $X$ sur $Y$}.
\end{enumerate}
\end{remarks}

\oldpage[IV-4]{114}

Cela étant, pour que l'homomorphisme $u$ soit bijectif, il faut et il suffit que $\bigwedge^n u$ soit bijectif, ou encore que la section $d_{X/Y}$ ait un germe \emph{inversible} en tout point $y\in Y$, ce qui s'écrit aussi $d_{X/Y}(y)\neq 0$ pour tout $y$ ($0_{\mathrm I}$, 5.5.2). Il revient aussi au même de dire que l'Idéal discriminant $\mathscr{D}_{X/Y}$ est \emph{égal à $\sh{O}_Y$}.

La terminologie de « revêtement » introduite dans (18.2.7, (i)), se justifie par la proposition suivante:

\begin{proposition}[18.2.8]
\label{IV.18.2.8}
Soit $f:X\to Y$ un morphisme étale, séparé et de type fini et pour tout $y\in Y$, soit $n(y)$ le nombre géométrique des points de $f^{-1}(y)$. Alors la fonction $y\mapsto n(y)$ est semi-continue inférieurement dans $Y$. Pour qu'elle soit continue en un point $y$ (donc constante dans un voisinage de $y$), il faut et il suffit qu'il existe un voisinage ouvert $U$ de $y$ tel que la restriction $f^{-1}(U)\to U$ de $f$ soit un morphisme fini (étale).
\end{proposition}

Comme $f$ est quasi-fini (17.6.1) et localement de présentation finie, il revient au même de dire que $f$ est fini ou que $f$ est propre (8.11.1): en outre, chaque fibre $f^{-1}(y)$ est géométriquement réduite sur $\kappa(y)$. Les conclusions résultent donc de (15.5.9, (i) et (ii)) et de ce que $f$ est plat.

\begin{corollary}[18.2.9]
\label{IV.18.2.9}
Soient $Y$ un préschéma connexe, $f:X\to Y$ un morphisme étale, séparé et de type fini. Pour que $f$ soit fini (autrement dit, que $X$ soit un revêtement étale de $Y$), il faut et il suffit que toutes les fibres de $f$ aient même nombre géométrique de points.
\end{corollary}

\begin{remarks}[18.2.10]
\label{IV.18.2.10}
\begin{enumerate}
  \item[(i)] L'exemple de la « droite affine avec un point dédoublé » (I, 5.5.11) montre qu'un morphisme étale de type fini de préschémas noethériens peut être \emph{non séparé}; pour cet exemple, la première assertion de (18.2.8) n'est plus exacte.

  \item[(ii)] Pour qu'un morphisme séparé, de type fini et étale $f:X\to Y$ fasse de $X$ un revêtement \emph{localement trivial}, il faut et il suffit que pour tout $x\in X$ il existe un voisinage ouvert $U$ de $y=f(x)$ et une $U$-section $g$ de $f^{-1}(U)$ telle que $g(y)=x$. En effet, la condition est évidemment nécessaire; le fait qu'elle est suffisante résulte de ce que toute fibre $f^{-1}(y)$ est finie (17.6.1) et (I, 6.2.2), de la caractérisation des sections d'un $Y$-schéma étale (17.9.3) et de la prop. (18.2.8).
\end{enumerate}
\end{remarks}

\subsection{Algèbres finies et étales.}
\EGAReaderTarget{ega4.c04.subsec.18-3}{subsection:IV.18.3}

\begin{proposition}[18.3.1]
\label{IV.18.3.1}
Soient $A$ un anneau, $B$ une $A$-algèbre de présentation finie.
\begin{enumerate}
  \item[(i)] Pour que $B$ soit une $A$-algèbre non ramifiée, il faut et il suffit que $B$ soit un $A$-module de présentation finie et que $B$ soit un $(B\otimes_A B)$-module projectif.
  \item[(ii)] Supposons de plus que $B$ soit une $A$-algèbre finie. Pour que $B$ soit une $A$-algèbre étale, il faut et il suffit que $B$ soit un $A$-module projectif et un $(B\otimes_A B)$-module projectif.
\end{enumerate}
\end{proposition}

Bien entendu, la structure de $(B\otimes_A B)$-module sur $B$ est celle provenant de la structure de $(B\otimes_A B)$-algèbre sur $B$ correspondant au $A$-homomorphisme canonique d'anneaux $B\otimes_A B\to B$, qui est surjectif et de noyau $\mathfrak J_{B/A}$ (0, 20.4.1).

(i) Dire que le morphisme $\Spec(B)\to\Spec(A)$ est localement de présentation finie équivaut à dire que $B$ est une $A$-algèbre de présentation finie (1.4.6). Dire

\oldpage[IV-4]{115}

que $B$ est une $A$-algèbre non ramifiée signifie alors (17.4.2) que $\Spec((B\otimes_A B)/\mathfrak J_{B/A})$ est un sous-schéma induit sur une partie \emph{ouverte et fermée} de $\Spec(B\otimes_A B)$ et on sait que pour qu'il en soit ainsi, il faut et il suffit que $\mathfrak J_{B/A}$ soit un idéal \emph{facteur direct} de $B\otimes_A B$ (Bourbaki, \emph{Alg. comm.}, chap. II, \textsection~4, n\textsuperscript{o}~3, prop. 15); mais il revient au même de dire que le $(B\otimes_A B)$-module quotient $(B\otimes_A B)/\mathfrak J_{B/A}$ est projectif (Bourbaki, \emph{Alg.}, chap. II, 3\textsuperscript{e} éd., \textsection~2, n\textsuperscript{o}~2, prop. 4).

(ii) Si on se rappelle qu'un $A$-module plat de présentation finie est projectif et réciproquement (Bourbaki, \emph{Alg. comm.}, chap. II, \textsection~5, n\textsuperscript{o}~2, cor. 2 du th. 1), l'assertion de (ii) résulte de celle de (i) et de (17.6.2).

\begin{proposition}[18.3.2]
\label{IV.18.3.2}
Soient $A$ un anneau, $\mathfrak J$ un idéal de $A$ tel que, pour la topologie $\mathfrak J$-pré-adique, $A$ soit séparé et complet; on pose $A_0=A/\mathfrak J$. Alors le foncteur
\[
  B\longmapsto B\otimes_A A_0
\]
est une équivalence de la catégorie des $A$-algèbres finies et étales, dans la catégorie des $A_0$-algèbres finies et étales.
\end{proposition}

Nous démontrerons d'abord le lemme suivant:

\begin{lemma}[18.3.2.1]
\label{IV.18.3.2.1}
Soient $A$ un anneau, $\mathfrak J$ un idéal de $A$ tel que, pour la topologie $\mathfrak J$-pré-adique, $A$ soit séparé et complet.
\begin{enumerate}
  \item[(i)] Tout $A$-module projectif de type fini $M$ est séparé et complet pour la topologie $\mathfrak J$-pré-adique, donc limite projective des $(A/\mathfrak J^{n+1})$-modules projectifs $M/\mathfrak J^{n+1}M$.
  \item[(ii)] Inversement, posons $A_n=A/\mathfrak J^{n+1}$, et soit $(M_n)$ un système projectif de $A_n$-modules, tel que, pour tout $n$, l'homomorphisme $M_{n+1}\otimes_{A_{n+1}}A_n\to M_n$ déduit du di-homomorphisme de transition $M_{n+1}\to M_n$ soit bijectif. Supposons en outre les $M_n$ projectifs et $M_0$ de type fini. Alors $M=\varprojlim M_n$ est un $A$-module projectif de type fini tel que l'homomorphisme canonique $M\otimes_A A_0\to M_0$ soit bijectif.
\end{enumerate}
\end{lemma}

(i) Il existe un $A$-module libre de type fini $L$ tel que $M$ soit isomorphe à un facteur direct de $L$; comme $L$ est séparé pour la topologie $\mathfrak J$-pré-adique, il en est de même de tout sous-module $N$ de $L$ puisque $\mathfrak J^{n+1}N\subset\mathfrak J^{n+1}L$; en particulier $M$ est séparé pour cette topologie, et comme l'homomorphisme surjectif $f:L\to M$ est continu pour la topologie $\mathfrak J$-pré-adique, son noyau $N$ est fermé pour la topologie induite par celle de $L$; puisque $L$ est complet et $f$ un morphisme strict, on en conclut que $M$ est complet (Bourbaki, \emph{Top. gén.}, chap. IX, 2\textsuperscript{e} éd., \textsection~3, n\textsuperscript{o}~1, prop. 4).

(ii) Il résulte du lemme de Nakayama que si $M_0$ est engendré par une famille finie $(x_{i0})$ de $r$ éléments et si pour tout $n$, $x_{in}$ est un élément de $M_n$ dont $x_{i,n-1}$ est l'image dans $M_{n-1}$, alors $(x_{in})$ $(1\leq i\leq r)$ est un système de générateurs de $M_n$ (Bourbaki, \emph{Alg. comm.}, chap. II, \textsection~3, cor. 2 de la prop. 4). Cela étant, pour tout $n$, posons $L_n=A_n^r$; si $(e_{in})_{1\leq i\leq r}$ est la base canonique de $L_n$, soit $u_n:L_n\to M_n$ l'application $A_n$-linéaire telle que $u_n(e_{in})=x_{in}$ pour tout $i$. Par hypothèse on a une suite exacte scindée
\[
  0\longrightarrow N_n\xrightarrow{v_n}L_n\xrightarrow{u_n}M_n\longrightarrow 0
\]

\oldpage[IV-4]{116}

et comme $L_n=L_{n+1}/\mathfrak J^{n+1}L_{n+1}$ et $M_n=M_{n+1}/\mathfrak J^{n+1}M_{n+1}$, les flèches verticales dans le diagramme commutatif
\[
\xymatrix{
0\ar[r] & N_{n+1}\ar[r]^{v_{n+1}}\ar[d] & L_{n+1}\ar[r]^{u_{n+1}}\ar[d] & M_{n+1}\ar[r]\ar[d] & 0\\
0\ar[r] & N_n\ar[r]_{v_n} & L_n\ar[r]_{u_n} & M_n\ar[r] & 0
}
\]
sont toutes trois surjectives. Or on a $M=\varprojlim M_n$ et $L=A^r=\varprojlim L_n$; si l'on pose $N=\varprojlim N_n$, $u=\varprojlim u_n$, $v=\varprojlim v_n$, on a, en vertu de ($0_{\mathrm{III}}$, 13.2.2), la suite exacte
\begin{equation}
  0\longrightarrow N\xrightarrow{v}L\xrightarrow{u}M\longrightarrow 0.
  \tag{18.3.2.2}\label{IV.18.3.2.2}
\end{equation}
En outre, comme pour tout $n$, $v_n$ est inversible à gauche et que $M_n$ est un $A_n$-module projectif, il résulte de (0, 19.1.8) que la suite exacte (18.3.2.2) est scindée, ce qui prouve le lemme.

Cela étant, montrons d'abord que le foncteur de l'énoncé de (18.3.2) est \emph{pleinement fidèle}. Posons, comme dans le lemme, $A_n=A/\mathfrak J^{n+1}$; soient $B$, $C$ deux $A$-algèbres finies et étales, et posons, pour tout $n$, $B_n=B\otimes_A A_n$, $C_n=C\otimes_A A_n$; en vertu de (18.3.1) et (18.3.2.1), $B$ et $C$ sont séparés et complets pour la topologie $\mathfrak J$-pré-adique, et l'on a $B=\varprojlim B_n$, $C=\varprojlim C_n$; en outre tout homomorphisme de $A$-algèbres $u:B\to C$ est continu pour les topologies $\mathfrak J$-pré-adiques, et donne donc un système projectif d'homomorphismes de $A_n$-algèbres $u_n=u\otimes 1:B_n\to C_n$, dont il est la limite projective; la réciproque étant évidente, on a donc une bijection canonique
\[
  \Hom_{A\text{-alg.}}(B,C)\xrightarrow{\sim}\varprojlim\Hom_{A_n\text{-alg.}}(B_n,C_n).
\]
Mais comme $B$ et $C$ sont des $A$-algèbres étales, il résulte aussitôt de (18.1.2) que l'application canonique
\[
  \Hom_{A_{n+1}\text{-alg.}}(B_{n+1},C_{n+1})\longrightarrow\Hom_{A_n\text{-alg.}}(B_n,C_n)
\]
est bijective pour $n\geq 0$, ce qui achève de prouver que l'application canonique
\[
  \Hom_{A\text{-alg.}}(B,C)\longrightarrow\Hom_{A_0\text{-alg.}}(B_0,C_0)
\]
est bijective.

Pour achever de démontrer (18.3.2), il suffit de voir que pour toute $A_0$-algèbre finie et étale $B_0$, il existe une $A$-algèbre finie et étale $B$ et un $A_0$-isomorphisme $B_0\xrightarrow{\sim}B\otimes_A A_0$. Or, il résulte de (18.1.2) qu'il y a un système projectif $(B_n)$ tel que $B_n$ soit une $A_n$-algèbre étale et que les homomorphismes $B_{n+1}\otimes_{A_{n+1}}A_n\to B_n$ soient bijectifs. Il résulte de (18.3.1) et (18.3.2.1) que la $A$-algèbre $B=\varprojlim B_n$ est un $A$-module projectif de type fini et que $B_0$ est isomorphe à $B\otimes_A A_0$. Pour prouver que $B$ est une $A$-algèbre étale, il suffit, en vertu de (18.2.5), de montrer que pour tout idéal maximal $\mathfrak m$ de $A$, $B_{\mathfrak m}/\mathfrak mB_{\mathfrak m}$ est une $(A/\mathfrak m)$-algèbre séparable. Or, puisque $\mathfrak J$ est contenu dans le radical de $A$ ($0_{\mathrm I}$, 7.1.10),

\oldpage[IV-4]{117}

on a $\mathfrak J\subset\mathfrak m$ et si $\mathfrak m_0=\mathfrak m/\mathfrak J$, on a $A_0/\mathfrak m_0=A/\mathfrak m$ et $B_{\mathfrak m}/\mathfrak mB_{\mathfrak m}=(B_0)_{\mathfrak m_0}/\mathfrak m_0(B_0)_{\mathfrak m_0}$; la conclusion résulte donc du fait que $B_0$ est une $A_0$-algèbre étale (18.2.5).

\begin{env}[18.3.3]
\label{IV.18.3.3}
\emph{Exemple.} --- La prop. (18.3.2) s'applique en particulier lorsque $A$ est un \emph{anneau local} séparé et complet, $\mathfrak J$ étant l'idéal maximal de $A$, de sorte que $A_0$ est un \emph{corps} et la catégorie des $A_0$-algèbres finies et étales identique à celle des $A_0$-algèbres finies et séparables, donc isomorphes à des composées directes de corps, extensions séparables et finies de $A_0$. En particulier, si le corps $A_0$ est \emph{séparablement clos}, ces extensions sont toutes identiques à $A_0$, et par suite tout revêtement étale de $\Spec(A)$ est \emph{trivial} (18.2.7) en vertu de (18.3.2).
\end{env}

\begin{theorem}[18.3.4]
\label{IV.18.3.4}
Soient $A$ un anneau noethérien, $\mathfrak J$ un idéal de $A$ tel que $A$ soit séparé et complet pour la topologie $\mathfrak J$-pré-adique, $A_0=A/\mathfrak J$. Posons $S=\Spec(A)$, $S_0=\Spec(A_0)$. Soit $X$ un $S$-schéma propre sur $S$, et posons $X_0=X\times_S S_0$. Alors le foncteur
\[
  Z\longmapsto Z\times_X X_0
\]
de la catégorie des $X$-schémas finis et étales sur $X$ dans la catégorie des $X_0$-schémas finis et étales sur $X_0$, est une équivalence de catégories.
\end{theorem}

Montrons d'abord que ce foncteur est \emph{pleinement fidèle}. Soient donc $Z'$ et $Z''$ deux $X$-schémas finis et étales sur $X$. Posons $S_n=\Spec(A/\mathfrak J^{n+1})$, $X_n=X\times_S S_n$, $Z'_n=Z'\times_S S_n$, $Z''_n=Z''\times_S S_n$ pour tout $n\geq 0$. Il résulte de (III, 5.4.1) que l'on a une bijection canonique $\Hom_X(Z',Z'')\xrightarrow{\sim}\varprojlim\Hom_{X_n}(Z'_n,Z''_n)$. Or, en vertu de (18.1.2), l'application canonique $\Hom_{X_{n+1}}(Z'_{n+1},Z''_{n+1})\to\Hom_{X_n}(Z'_n,Z''_n)$ est bijective, ce qui achève de prouver notre assertion.

Il reste à prouver que si $\sh{B}_0$ est une $\sh{O}_{X_0}$-Algèbre finie et étale, il existe une $\sh{O}_X$-Algèbre finie et étale $\sh{B}$ et un isomorphisme $\sh{B}_0\xrightarrow{\sim}\sh{B}\otimes_{\sh{O}_X}\sh{O}_{X_0}$. Il résulte de (18.1.2) qu'il y a un système projectif $(\sh{B}_n)$, où $\sh{B}_n$ est une $\sh{O}_{X_n}$-Algèbre finie et étale, et le deuxième théorème de comparaison (III, 5.1.4) prouve l'existence d'un $\sh{O}_X$-Module cohérent $\sh{B}$ et d'un système projectif d'isomorphismes $\sh{B}_n\xrightarrow{\sim}\sh{B}\otimes_{\sh{O}_X}\sh{O}_{X_n}$. La donnée d'une structure d'Algèbre sur un Module $\sh{F}$ étant équivalente à celle d'un homomorphisme $\sh{F}\otimes\sh{F}\to\sh{F}$ rendant commutatifs des diagrammes où n'interviennent que des puissances tensorielles de $\sh{F}$, il résulte de (III, 5.1.3), (I, 10.11.4) et (I, 10.11.7) que $\sh{B}$ est muni de façon naturelle d'une structure de $\sh{O}_X$-Algèbre pour laquelle l'isomorphisme $\sh{B}_0\xrightarrow{\sim}\sh{B}\otimes_{\sh{O}_X}\sh{O}_{X_0}$ est un isomorphisme d'Algèbres. De plus, $\sh{B}$ est un $\sh{O}_X$-Module \emph{localement libre}; cela résulte encore de (III, 5.1.3), (I, 10.11.4), (I, 10.11.7) et du fait que, dans la catégorie des $\sh{O}_X$-Modules cohérents, les $\sh{O}_X$-Modules localement libres $\sh{F}$ peuvent être définis comme ceux pour lesquels le foncteur $\sh{G}\mapsto\mathscr{H}\!om_{\sh{O}_X}(\sh{F},\sh{G})$ est \emph{exact}. Enfin, pour voir que $\sh{B}$ est une $\sh{O}_X$-Algèbre étale, il suffit (18.2.5) de montrer que pour tout point \emph{fermé} $x\in X$, $\sh{B}_x\otimes_{\sh{O}_{X,x}}\kappa(x)$ est une $\kappa(x)$-algèbre séparable. Mais comme le morphisme structural $f:X\to S$ est \emph{propre}, $f(x)$ est un point fermé de $S$, donc appartient à $S_0$, puisque $\mathfrak J$ est contenu dans le radical de $A$ ($0_{\mathrm I}$, 7.1.10); la conclusion résulte donc de ce que $X_0=f^{-1}(S_0)$ et de ce que $\sh{B}_0$ est une $\sh{O}_{X_0}$-Algèbre finie et étale.

\subsection{Structure locale des morphismes non ramifiés et des morphismes étales.}
\EGAReaderTarget{ega4.c04.subsec.18-4}{subsection:IV.18.4}

\oldpage[IV-4]{118}

\begin{lemma}[18.4.1]
\label{IV.18.4.1}
Soient $A$ un anneau, $B$ une $A$-algèbre finie monogène, $u$ un générateur de la $A$-algèbre $B$, $F\in A[T]$ un polynôme tel que $F(u)=0$, $F'$ le polynôme dérivé; posons $u'=F'(u)$. Alors l'idéal de $B$, annulateur de $\Omega^1_{B/A}$, contient $u'B$; il est égal à $u'B$ si l'idéal $\mathfrak J$ de $A[T]$ formé des polynômes $G$ tels que $G(u)=0$, est engendré par $F$, autrement dit si l'homomorphisme surjectif canonique $\varphi:A[T]/F.A[T]\to B$ transformant l'image de $T$ en $u$, est bijectif.
\end{lemma}

Posons $C=A[T]$, de sorte que $B=C/\mathfrak J$. On a la suite exacte (0, 20.5.12.1)
\[
  \mathfrak J/\mathfrak J^2\longrightarrow\Omega^1_{C/A}\otimes_C B\longrightarrow\Omega^1_{B/A}\longrightarrow 0
\]
et $\Omega^1_{B/A}$ s'identifie donc au quotient $B/\mathfrak J'$, $\mathfrak J'$ étant l'idéal engendré par les éléments $G'(u)$, où $G$ parcourt un système de générateurs de l'idéal $\mathfrak J$ (0, 20.5.13); d'où aussitôt le lemme.

\begin{proposition}[18.4.2]
\label{IV.18.4.2}
Les notations étant celles de (18.4.1), soit $\mathfrak q$ un idéal premier de $B$. Alors:
\begin{enumerate}
  \item[(i)] Si $\mathfrak q$ ne contient pas $u'$, $B_{\mathfrak q}$ est une $A_{\mathfrak p}$-algèbre formellement non ramifiée ($\mathfrak p$ étant l'image réciproque de $\mathfrak q$ dans $A$); en d'autres termes, $\Spec(B_{u'})$ est formellement non ramifié sur $\Spec(A)$.
  \item[(ii)] Supposons de plus que $F$ soit unitaire et engendre $\mathfrak J$. Alors, pour que $\Spec(B)$ soit étale sur $\Spec(A)$ au point $\mathfrak q$, il faut et il suffit que $u'\notin\mathfrak q$.
\end{enumerate}
\end{proposition}

L'hypothèse que $u'\notin\mathfrak q$ entraîne que $\Omega^1_{B_{\mathfrak q}/A_{\mathfrak p}}=0$ (0, 20.5.9), donc (i) résulte de (17.2.1). En outre, sous les hypothèses de (ii), $B$ est un $A$-module libre en vertu de la division euclidienne; comme l'annulateur $\mathfrak J'$ de $\Omega^1_{B/A}$ est alors égal à $u'B$ en vertu de (18.4.1) et que $\Omega^1_{B/A}$ est un $B$-module de présentation finie (16.4.22), l'annulateur de $\Omega^1_{B_{\mathfrak q}/A_{\mathfrak p}}$ est égal à $u'B_{\mathfrak q}$ (Bourbaki, \emph{Alg. comm.}, chap. II, \textsection~2, n\textsuperscript{o}~4, formule (9)), et (ii) résulte donc de (i) et de l'implication $c)\Rightarrow a)$ dans (17.6.1).

\begin{corollary}[18.4.3]
\EGAReaderTarget{ega4.c04.cor.18-4-3}{IV.18.4.3}
Les notations étant celles de (18.4.2), supposons que $F$ soit unitaire et engendre $\mathfrak J$. Alors, pour que $B$ soit une $A$-algèbre étale, il faut et il suffit que $u'$ soit inversible dans $B$ (ou, ce qui revient au même, que l'idéal de $A[T]$ engendré par $F$ et $F'$ soit égal à $A[T]$).
\end{corollary}

Compte tenu de (18.4.2), (ii), dire que $\Spec(B)$ est étale sur $\Spec(A)$ signifie en effet que $u'$ n'appartient à aucun idéal premier de $B$, i.e. qu'il est inversible dans $B$.

On dit qu'un polynôme unitaire $F\in A[T]$ tel que l'idéal de $A[T]$ engendré par $F$ et $F'$ soit égal à $A[T]$ lui-même, est \emph{séparable}; il est immédiat que cette définition coïncide avec la définition usuelle (Bourbaki, \emph{Alg.}, chap. V, \textsection~7, n\textsuperscript{o}~6) lorsque $A$ est un corps.

\begin{lemma}[18.4.4]
\label{IV.18.4.4}
Soient $A$ un anneau local, $B$ une $A$-algèbre finie monogène, $u$ un générateur de la $A$-algèbre $B$. Soient $\mathfrak n_i$ ($1\leq i\leq r$) des idéaux maximaux de $B$ tels que $\Spec(B)$ soit formellement non ramifié sur $\Spec(A)$ aux points $\mathfrak n_i$. Alors il existe un polynôme unitaire $F\in A[T]$ tel que $F(u)=0$ et $F'(u)\notin\mathfrak n_i$ pour tout indice $i$ ($1\leq i\leq r$). En outre, si $k$ est le corps

\oldpage[IV-4]{119}

résiduel de $A$, $f\in k[T]$ le polynôme minimal de l'image de $u$ dans $B\otimes_A k$, il existe un $F\in A[T]$ dont $f$ est l'image canonique et tel que $F(u)=0$; un tel polynôme $F$ vérifie les conditions $F'(u)\notin\mathfrak n_i$ pour $1\leq i\leq r$.
\end{lemma}

L'idéal maximal $\mathfrak m$ de $A$ est l'image réciproque de chacun des $\mathfrak n_i$ (II, 6.1.10); soit $L=B\otimes_A k$, qui est une $k$-algèbre finie. Soit $\xi$ l'image de $u$ dans $L$, et soit $n$ le rang de $L$ sur $k$; le polynôme minimal $f\in k[T]$ de $\xi$ sur $k$ est donc de degré $n$, et $L$ est isomorphe à $k[T]/f.k[T]$. Si $\mathfrak n'_i=\mathfrak n_i/\mathfrak mB$, l'hypothèse entraîne que $\Spec(L)$ est étale sur $k$ aux points $\mathfrak n'_i$ ($1\leq i\leq r$) en vertu de (17.6.1, $c)$), donc $f'(\xi)\notin\mathfrak n'_i$ en vertu de (18.4.2, (ii)). Notons maintenant que $\mathfrak mB$ est contenu dans le radical de $B$; comme $1,\xi,\ldots,\xi^{n-1}$ forment une base de $L$ sur $k$, il résulte du lemme de Nakayama que $1,u,\ldots,u^{n-1}$ engendrent le $A$-module $B$, et par suite il existe un polynôme unitaire $F\in A[T]$ de degré $n$, tel que $F(u)=0$; en outre, comme $\xi$ est racine de l'image canonique de $F$ dans $k[T]$, cette image est nécessairement égale à $f$. Mais alors l'image de $F'(u)$ dans $L$ est $f'(\xi)$, et comme $f'(\xi)\notin\mathfrak n'_i$, on a $F'(u)\notin\mathfrak n_i$ pour tout $i$.

\begin{proposition}[18.4.5]
\label{IV.18.4.5}
Soient $A$ un anneau local, $k$ son corps résiduel, $B$ une $A$-algèbre finie (resp. finie et de présentation finie). On suppose en outre, ou bien que le corps $k$ est infini, ou bien que $B$ est un anneau local. Soit $n$ le rang de $L=B\otimes_A k$ sur $k$. Pour que $B$ soit une $A$-algèbre formellement non ramifiée (resp. étale), il faut et il suffit qu'il existe un polynôme unitaire séparable $F\in A[T]$ (18.4.3) tel que $B$ soit isomorphe à un quotient de $A[T]/F.A[T]$ (resp. isomorphe à $A[T]/F.A[T]$). En outre, on peut supposer que $F$ est de degré $n$ (resp. $F$ est nécessairement de degré $n$).
\end{proposition}

Les conditions sont suffisantes en vertu de (18.4.2), sans supposer $k$ infini ni $B$ local. Pour voir que les conditions sont nécessaires, notons que si $B$ est une $A$-algèbre formellement non ramifiée, $L$ est une algèbre finie et séparable sur $k$, donc composée directe d'un nombre fini d'extensions finies et séparables $k_j$ de $k$ ($1\leq j\leq r$), $k_j$ étant donc engendrée par un élément $\xi_j$, de polynôme minimal $f_j$ ($1\leq j\leq r$) (Bourbaki, \emph{Alg.}, chap. V, \textsection~11, n\textsuperscript{o}~4, prop. 4). Montrons qu'en vertu des hypothèses faites sur $k$ ou $B$, il existe un élément $\xi$ de $L$ engendrant la $k$-algèbre $L$. C'est immédiat si $B$ est local, puisque alors $r=1$. Sinon, $k$ étant supposé infini, on peut supposer que les polynômes irréductibles $f_j\in k[T]$ sont tous \emph{distincts} (en remplaçant au besoin chaque $\xi_j$ par $\xi_j+a_j$, pour un élément convenable $a_j\in k$); si l'on pose $f=f_1f_2\cdots f_r$, il est clair que $L$ est isomorphe à $k[T]/f.k[T]$ dans les deux cas envisagés, donc est engendrée par un élément $\xi$ de polynôme minimal $f\in k[T]$ de degré $n$. Si $u\in B$ est un élément d'image $\xi$ dans $L$, le lemme de Nakayama montre que les éléments $1,u,\ldots,u^{n-1}$ engendrent le $A$-module $B$; ceci montre déjà qu'il existe un polynôme unitaire $F\in A[T]$ de degré $n$ tel que $F(u)=0$, $u$ engendrant la $A$-algèbre $B$, qui est par suite isomorphe à une algèbre quotient de $A[T]/F.A[T]$; en outre $B$ est un anneau semi-local et en chacun de ses idéaux maximaux $\mathfrak n_i$, on a $F'(u)\notin\mathfrak n_i$ d'après (18.4.4), ce qui prouve que $F'(u)$ est inversible dans $B$, donc que $F$ est un polynôme \emph{séparable}. Enfin, si $B$ est une $A$-algèbre étale, $B$ étant un $A$-module plat de présentation finie (1.4.7) est un $A$-module libre, et $1,u,\ldots,u^{n-1}$ forment une base du $A$-module $B$ (Bourbaki, \emph{Alg. comm.}, chap. II, \textsection~3, n\textsuperscript{o}~2, prop. 5), autrement dit la $A$-algèbre $B$ est

\oldpage[IV-4]{120}

isomorphe à $A[T]/F.A[T]$, et pour tout autre polynôme unitaire $G\in A[T]$ tel que $B$ soit isomorphe à $A[T]/G.A[T]$, $G$ est nécessairement de degré $n$.

\begin{theorem}[18.4.6]
\label{IV.18.4.6}
(Chevalley). ---
\begin{enumerate}
  \item[(i)] Soient $f:X\to Y$ un morphisme localement de type fini, $x$ un point de $X$, $y=f(x)$, et posons $A=\sh{O}_{Y,y}$. Pour que $\sh{O}_{X,x}$ soit une $A$-algèbre formellement étale (resp. formellement non ramifiée), il faut et il suffit qu'il existe un polynôme unitaire $F\in A[T]$ et un idéal maximal $\mathfrak n$ de $B=A[T]/F.A[T]$ (resp. d'une algèbre quotient $B$ de $A[T]/F.A[T]$) tels que $\sh{O}_{X,x}$ soit $A$-isomorphe à $B_{\mathfrak n}$ et que, si $u$ est l'image de $T$ dans $B$, on ait $F'(u)\notin\mathfrak n$.
  \item[(ii)] Supposons en outre que $f$ soit localement de présentation finie. Alors, pour que $f$ soit étale au point $x$, il faut et il suffit de plus que l'on puisse prendre $B=A[T]/F.A[T]$.
\end{enumerate}
\end{theorem}

Les conditions sont suffisantes en vertu de (18.4.2). Pour voir qu'elles sont nécessaires, on peut évidemment se borner au cas où $X$ et $Y$ sont affines, et, compte tenu de la remarque (17.4.1.2), au cas où $f$ est formellement non ramifié et quasi-fini. Comme $f$ est affine, il résulte de (8.12.8) qu'il existe une $A$-algèbre finie $C$ et un idéal maximal $\mathfrak r$ de $C$ (nécessairement au-dessus de l'idéal maximal $\mathfrak m$ de $A$) tels que $\sh{O}_{X,x}$ soit $A$-isomorphe à $C_{\mathfrak r}$. En outre (17.4.1.2) le corps résiduel $C/\mathfrak r=C_{\mathfrak r}/\mathfrak rC_{\mathfrak r}$ est une extension finie séparable de $k=A/\mathfrak m$, donc de la forme $k[v]$, où $v$ est \emph{séparable} sur $k$. Soient $\mathfrak r_i$ ($1\leq i\leq h$) les idéaux maximaux de l'anneau semi-local $C$ \emph{autres} que $\mathfrak r$; il existe un élément $u\in C$ appartenant à tous les $\mathfrak r_i$ et tel que son image dans $C/\mathfrak r$ soit égale à $v$ (Bourbaki, \emph{Alg. comm.}, chap. II, \textsection~1, n\textsuperscript{o}~2, prop. 5). Nous allons montrer que la sous-$A$-algèbre $B=A[u]$ de $C$ et l'idéal (nécessairement maximal puisque $C$ est finie sur $B$) $\mathfrak n=\mathfrak r\cap B$ de $B$ répondent à la question.

Pour traiter le cas où $\sh{O}_{X,x}$ est une $A$-algèbre formellement ramifiée, il suffira de prouver que $B_{\mathfrak n}$ est isomorphe à $C_{\mathfrak r}$; en effet, $B_{\mathfrak n}$ sera alors formellement non ramifié sur $A$ et l'existence du polynôme $F\in A[T]$ ayant les propriétés de l'énoncé résultera de (18.4.4). Notons maintenant que, puisque $f$ est formellement non ramifié, $C/\mathfrak r$ est isomorphe à la $A$-algèbre $C_{\mathfrak r}/\mathfrak mC_{\mathfrak r}$ (17.4.1.2). On est ainsi ramené à prouver le lemme suivant:

\begin{lemma}[18.4.6.1]
\label{IV.18.4.6.1}
Soient $A$ un anneau local, $\mathfrak m$ son idéal maximal, $C$ une $A$-algèbre finie, $\mathfrak r$ un idéal maximal de $C$. Soit $u$ un élément de $C$ appartenant à tous les idéaux maximaux de $C$ distincts de $\mathfrak r$, n'appartenant pas à $\mathfrak r$, et tel que $C_{\mathfrak r}/\mathfrak mC_{\mathfrak r}$ soit une algèbre monogène sur $k=A/\mathfrak m$, engendrée par l'image de $u$ dans $C_{\mathfrak r}/\mathfrak mC_{\mathfrak r}$. Posons $B=A[u]$, $\mathfrak n=\mathfrak r\cap B$. Alors l'homomorphisme canonique $B_{\mathfrak n}\to C_{\mathfrak r}$ est un isomorphisme.
\end{lemma}

Posons $R=B-\mathfrak n$, $S=C-\mathfrak r$, de sorte que $B_{\mathfrak n}=R^{-1}B$ et $C_{\mathfrak r}=S^{-1}C$; l'homomorphisme canonique $B_{\mathfrak n}\to C_{\mathfrak r}$ peut s'écrire comme le composé
\[
  R^{-1}B\xrightarrow{g}R^{-1}C\xrightarrow{h}S^{-1}C
\]
et il suffit de montrer que chacun de ces deux homomorphismes est bijectif.

Montrons d'abord que $h:R^{-1}C\to S^{-1}C=C_{\mathfrak r}$ est bijectif; il suffit de voir que les images dans $R^{-1}C$ des éléments de $S$ sont inversibles, ou encore que tout idéal maximal $\mathfrak p$ de $R^{-1}C$ a une image réciproque dans $C$ ne rencontrant pas $S$, donc nécessairement égale à $\mathfrak r$. Or, comme $R^{-1}C$ est une $(R^{-1}B)$-algèbre finie, l'image réciproque de $\mathfrak p$ dans

\oldpage[IV-4]{121}

$R^{-1}B=B_{\mathfrak n}$ est l'unique idéal maximal $\mathfrak nB_{\mathfrak n}$ de cet anneau, donc l'image réciproque de $\mathfrak p$ dans $B$ est égale à $\mathfrak n$. Mais d'autre part, si $\mathfrak q$ est l'image réciproque de $\mathfrak p$ dans $C$, on a $\mathfrak q\cap B=\mathfrak n$, et comme $\mathfrak n$ est maximal dans $B$ et $C$ une $B$-algèbre finie, $\mathfrak q$ est nécessairement un des idéaux maximaux de $C$; en outre, on a $u\notin\mathfrak q$ puisque $u\notin\mathfrak r$ et que $u\in B$, donc par hypothèse on a nécessairement $\mathfrak q=\mathfrak r$.

D'autre part, comme $B\subset C$, l'homomorphisme $g:R^{-1}B\to R^{-1}C$ est injectif ($0_{\mathrm I}$, 1.3.2); pour voir qu'il est surjectif, notons que $R^{-1}C$ est un $(R^{-1}B)$-module de type fini, et d'autre part que $\mathfrak mR^{-1}B$ est contenu dans l'idéal maximal de l'anneau local $R^{-1}B$; en vertu du lemme de Nakayama, il suffit de prouver que l'homomorphisme
\[
  R^{-1}B/\mathfrak mR^{-1}B\longrightarrow R^{-1}C/\mathfrak mR^{-1}C
\]
est surjectif. Mais en vertu de la première partie de la démonstration, $R^{-1}C/\mathfrak mR^{-1}C$ s'identifie à $C_{\mathfrak r}/\mathfrak mC_{\mathfrak r}$; par hypothèse cette $k$-algèbre est engendrée par l'image de $u$, et \emph{a fortiori} elle est égale à l'image de $R^{-1}B/\mathfrak mR^{-1}B$.

Considérons en second lieu le cas où $f$ est étale au point $x$. En remplaçant $X$ par un voisinage de $x$, on peut supposer que $X$ est un voisinage de $\mathfrak n$ dans $\Spec(B)$ (I, 7.2). Posons $B'=\Spec(A[T]/F.A[T])$ et soit $\mathfrak n'$ l'image réciproque de $\mathfrak n$ dans $B'$; comme l'image de $F'(T)$ dans $B'$ n'appartient pas à $\mathfrak n'$ par hypothèse, le morphisme $\Spec(B')\to\Spec(A)$ est étale au point $\mathfrak n'$ par (18.4.2, (ii)). Comme par hypothèse $\Spec(B)\to\Spec(A)$ est étale au point $\mathfrak n$, on en conclut (17.3.4) que $\Spec(B)\to\Spec(B')$ est étale au point $\mathfrak n$; mais comme ce morphisme est une immersion, il ne peut être étale en un point que si c'est un isomorphisme local en ce point (17.9.1), donc $B_{\mathfrak n}$ et $B'_{\mathfrak n'}$ sont isomorphes.

Enfin, supposons que $\sh{O}_{X,x}$ soit une $A$-algèbre formellement étale; avec les notations précédentes, il résulte de (17.1.5) que $B_{\mathfrak n}$ est une $B'_{\mathfrak n'}$-algèbre formellement étale; mais puisque l'homomorphisme $B'_{\mathfrak n'}\to B_{\mathfrak n}$ est surjectif, cela ne peut avoir lieu que si cet homomorphisme est \emph{bijectif} (0, 19.10.3, (i)). Ceci achève la démonstration de (18.4.6).

\begin{corollary}[18.4.7]
\label{IV.18.4.7}
Soient $f:X\to Y$ un morphisme localement de type fini, $x$ un point de $X$. Pour que $f$ soit formellement non ramifié au point $x$, il faut et il suffit qu'il existe un voisinage ouvert $U$ de $x$ tel que $f|U$ se factorise en $U\xrightarrow{j}X'\xrightarrow{h}Y$, où $h$ est un morphisme étale et $j$ une immersion fermée.
\end{corollary}

On peut évidemment se borner au cas où $Y=\Spec(R)$ est affine et $f$ de type fini. Si $A=\sh{O}_{Y,y}$, avec $y=f(x)$, la condition pour que $f$ soit formellement non ramifié au point $x$ équivaut, en vertu de (17.4.1.2), à dire que $\sh{O}_{X,x}$ est une $A$-algèbre formellement non ramifiée. S'il en est ainsi, on peut appliquer (18.4.6, (ii)); remplaçant au besoin $Y$ par un voisinage affine de $y$, on peut supposer (avec les notations de (18.4.6)) que le polynôme $F$ est l'image dans $A[T]$ d'un polynôme unitaire $G\in R[T]$. On pose alors $X'=\Spec(R[T]/G.R[T])$; soit $x'$ l'image du point $\mathfrak n$ de $\Spec(B)$ par le morphisme correspondant à l'homomorphisme composé $R[T]/G.R[T]\to A[T]/F.A[T]\to B$. Il résulte de (18.4.2) que le morphisme $h:X'\to Y$ correspondant à l'homomorphisme canonique $R\to R[T]/G.R[T]$ est étale au point $x'$, donc, en restreignant au besoin $X'$ et $Y$ à des voisinages ouverts de $x'$ et $y$ respectivement, on peut supposer $h$ étale. D'autre part, en vertu de (I, 6.5.1, (ii)) et de (I, 7.2), il résulte du fait que l'on a un homomorphisme local $\varphi:\sh{O}_{X',x'}\to\sh{O}_{X,x}$ que cet homomorphisme correspond à un morphisme

\oldpage[IV-4]{122}

$j:U\to X'$ d'un voisinage ouvert $U$ de $x$ dans $X$; restreignant au besoin $U$, on peut, en appliquant (I, 6.5.1, (i)), supposer que $h\circ j=f|U$, donc $j$ est de type fini (I, 6.3.4, (v)); enfin, l'homomorphisme $\varphi$ est surjectif par (18.4.6, (ii)), donc il résulte de (I, 6.5.4, (i)) qu'on peut, en restreignant encore $U$ et $X'$, supposer que $j$ est une immersion fermée. Cela prouve donc la nécessité de la condition énoncée; sa suffisance est immédiate (17.1.3, (i) et (ii)).

\begin{corollary}[18.4.8]
\label{IV.18.4.8}
Soit $f:X\to Y$ un morphisme localement de présentation finie. Pour que $f$ soit non ramifié (resp. étale), il faut et il suffit qu'il existe une famille de morphismes plats $g_\alpha:Y'_\alpha\to Y$ et, pour chaque $\alpha$, un ouvert $U_\alpha$ dans $X'_\alpha=X\times_Y Y'_\alpha$, tels que, si $g'_\alpha:X'_\alpha\to X$ et $f'_\alpha:X'_\alpha\to Y'_\alpha$ sont les projections canoniques, les $g'_\alpha(U_\alpha)$ forment un recouvrement de $X$, et que chacun des morphismes composés $U_\alpha\to X'_\alpha\xrightarrow{f'_\alpha}Y'_\alpha$ soit une immersion fermée (resp. ouverte). On peut de plus alors prendre les $g_\alpha$ étales.
\end{corollary}

La nécessité de la condition pour les morphismes étales est triviale, en prenant un seul $Y'_\alpha$ égal à $X$, le morphisme $g_\alpha$ correspondant étant égal à $f$, et l'ouvert $U\subset X\times_Y X$ étant la diagonale. Lorsque $f$ est non ramifié, la nécessité de la condition résulte de (18.4.7); on prend un recouvrement ouvert $(V_\alpha)$ de $X$ tel que pour chaque $\alpha$, $f|V_\alpha$ se factorise en $V_\alpha\xrightarrow{j_\alpha}Y'_\alpha\xrightarrow{g_\alpha}Y$, où $j_\alpha$ est une immersion fermée et $g_\alpha$ un morphisme étale. Alors $j_\alpha:V_\alpha\to Y'_\alpha$ se factorise en $V_\alpha\xrightarrow{s_\alpha}X'_\alpha\xrightarrow{f'_\alpha}Y'_\alpha$, où $s_\alpha$ est une $V_\alpha$-section de $X'_\alpha$, et comme le morphisme $g'_\alpha:X'_\alpha\to X$ est étale, $s_\alpha$ est une immersion ouverte (17.4.1), et il suffit de prendre $U_\alpha=s_\alpha(V_\alpha)$ pour répondre à la question.

La suffisance des conditions résulte de (17.7.1); on en conclut que $f$ est non ramifié (resp. étale) en chaque point de $g'(U_\alpha)$, donc dans $X$ tout entier puisque les $g'(U_\alpha)$ recouvrent $X$.

\begin{proposition}[18.4.9]
\label{IV.18.4.9}
Soient $S$ un préschéma, $f:X\to S$ un morphisme, $h:Y\to S$ un morphisme localement de présentation finie, $g:X\to Y$ un $S$-morphisme, $x$ un point de $X$, $y=g(x)$. Les conditions suivantes sont équivalentes:
\begin{enumerate}
  \item[a)] $h$ est étale au point $y$ et $g$ est plat au point $x$.
  \item[b)] $h$ est non ramifié au point $y$ et $f$ est plat au point $x$.
\end{enumerate}
\end{proposition}

Comme $f=h\circ g$, $a)$ entraîne que $f$ est plat au point $x$ (2.1.6), et évidemment que $h$ est non ramifié au point $y$, donc $a)$ implique $b)$.

Pour prouver que $b)$ entraîne $a)$, on peut d'abord supposer que $h$ est non ramifié (en remplaçant $Y$ par un voisinage de $y$); puis, en remplaçant $S$ par un voisinage ouvert de $s=h(y)=f(x)$, on peut supposer qu'il existe un morphisme étale $u:S'\to S$, un point $y'$ de $Y'=Y_{(S')}$ au-dessus de $y$ et un voisinage ouvert $V$ de $y'$ dans $Y'$ tel que, si $h'=h_{(S')}:Y'\to S'$, la restriction de $h'$ à $V$ soit une immersion fermée (18.4.8). Si l'on prouve alors que $h'$ est étale au point $y'$, il en résultera que $h$ est étale au point $y$ (17.7.1, (ii)); d'ailleurs, $f'=f_{(S')}:X_{(S')}\to S'$ est plat en tout point $x'$ au-dessus de $x$. Comme les projections $v:Y'\to Y$, $w:X'\to X$ sont des morphismes étales (donc plats), si l'on prouve que $g'=g_{(S')}:X_{(S')}\to Y_{(S')}$ est plat au point $x'$, il résultera de (2.2.11, (iv)) que $g$ sera plat au point $x$. On peut donc se borner au cas où $h$ est une \emph{immersion fermée}

\oldpage[IV-4]{123}

de présentation finie, $f$ étant supposé plat au point $x$. Soit $\mathscr J$ l'Idéal quasi-cohérent de type fini de $\sh{O}_S$ qui définit le sous-préschéma fermé $Y$ de $S$. L'hypothèse que $f$ est plat au point $x$ entraîne que l'homomorphisme $\sh{O}_{S,s}\to\sh{O}_{X,x}$ est injectif ($0_{\mathrm I}$, 6.5.1); \emph{a fortiori} l'homomorphisme $\sh{O}_{S,s}\to\sh{O}_{Y,y}$ est injectif, ce qui signifie que $\mathscr J_s=0$, puisque c'est le noyau de l'homomorphisme précédent. Puisque $\mathscr J$ est de type fini, il y a un voisinage ouvert $U$ de $s$ dans $S$ tel que $\mathscr J|U=0$ ($0_{\mathrm I}$, 5.2.2). On peut donc supposer que $h$ est une immersion ouverte, et alors il est clair que $g$ est plat au point $x$.

\begin{proposition}[18.4.10]
\label{IV.18.4.10}
Désignons par $\mathbf P(f,x)$ une propriété vérifiant les conditions suivantes:
\begin{enumerate}
  \item[$1^\circ$] Pour tout morphisme $f:X\to Y$ et tout isomorphisme local $h:Y\to Z$, $\mathbf P(f,x)$ est équivalente à $\mathbf P(h\circ f,x)$ pour $x\in X$.
  \item[$2^\circ$] Pour tout morphisme $f:X\to Y$, tout morphisme étale $g:Y'\to Y$, tout point $x\in X$, si l'on pose $X'=X\times_Y Y'$, $f'=f_{(Y')}:X'\to Y'$ et si $x'\in X'$ est au-dessus de $x$, les propriétés $\mathbf P(f,x)$ et $\mathbf P(f',x')$ sont équivalentes (« invariance par changement de base étale »).
\end{enumerate}
Soient alors $S$ un préschéma, $f:X\to S$ et $h:Y\to S$ deux morphismes, $g:X\to Y$ un $S$-morphisme, $x$ un point de $X$, $y=g(x)$, et supposons $h$ étale au point $y$. Alors les propriétés $\mathbf P(f,x)$ et $\mathbf P(g,x)$ sont équivalentes.
\end{proposition}

Avec les mêmes notations que dans la démonstration de (18.4.9) et en remplaçant $S$ (resp. $Y$) par un voisinage de $s=h(y)$ (resp. de $y$), on peut ici, en vertu de (18.4.8), trouver un morphisme étale $u:S'\to S$ tel que $h'$ soit une immersion ouverte; si $x'\in X'$ est au-dessus de $x$, $\mathbf P(f',x')$ (resp. $\mathbf P(g',x')$) est alors par hypothèse équivalente à $\mathbf P(f,x)$ (resp. $\mathbf P(g,x)$). On peut donc se borner au cas où $h$ est une immersion ouverte. Comme l'hypothèse entraîne que $\mathbf P(g,x)$ est alors équivalente à $\mathbf P(h\circ g,x)$, on en déduit la conclusion.

\begin{env}[18.4.11]
\label{IV.18.4.11}
\emph{Exemples.} --- On peut prendre pour propriété $\mathbf P(f,x)$ l'une quelconque des suivantes, compte tenu de (17.7.4, (ii)):
\begin{enumerate}
  \item[(i)] $f$ est plat au point $x$ (2.2.11, (iv));
  \item[(ii)] $f$ est localement de présentation finie et de coprofondeur $\leq n$ au point $x$ (6.8.1 et 6.7.8);
  \item[(iii)] $f$ est localement de présentation finie et de Cohen-Macaulay au point $x$ (6.8.1 et 6.7.8);
  \item[(iv)] $f$ est localement de présentation finie et possède la propriété $(S_n)$ au point $x$ (6.8.1 et 6.7.8);
  \item[(v)] $f$ est localement de présentation finie et possède la propriété $(R_n)$ au point $x$ (6.8.1 et 6.7.8);
  \item[(vi)] $f$ est localement de présentation finie et normal au point $x$ (6.8.1 et 6.7.8);
  \item[(vii)] $f$ est localement de présentation finie et réduit au point $x$ (6.8.1 et 6.7.8);
  \item[(viii)] $f$ est non ramifié au point $x$ (17.7.4);
  \item[(ix)] $f$ est lisse au point $x$ (17.7.4);
  \item[(x)] $f$ est étale au point $x$ (17.7.4).
\end{enumerate}
\end{env}

\oldpage[IV-4]{124}

\begin{corollary}[18.4.12]
\label{IV.18.4.12}
\begin{enumerate}
  \item[(i)] Soient $f:X\to Y$ un morphisme localement de type fini, $x$ un point de $X$. Si $f$ est plat et formellement non ramifié au point $x$, alors, dans toute factorisation $f|U:U\xrightarrow{j}X'\xrightarrow{h}Y$, où $U$ est un voisinage ouvert de $x$, $j$ une immersion fermée et $h$ un morphisme étale (18.4.7), l'homomorphisme $\sh{O}_{X',j(x)}\to\sh{O}_{X,x}$ correspondant à $j$ est bijectif (en particulier $\sh{O}_{X,x}$ est une $\sh{O}_{Y,f(x)}$-algèbre \emph{essentiellement étale} (18.6.1)).
  \item[(ii)] Pour qu'un morphisme $f:X\to Y$ soit étale, il faut et il suffit qu'il soit localement de type fini, formellement non ramifié et plat.
\end{enumerate}
\end{corollary}

(i) Comme l'homomorphisme $\sh{O}_{X',j(x)}\to\sh{O}_{X,x}$ est surjectif, il suffit de prouver qu'il est injectif, et pour cela qu'il fait de $\sh{O}_{X,x}$ un $\sh{O}_{X',j(x)}$-module \emph{fidèlement plat}, ou seulement \emph{plat} ($0_{\mathrm I}$, 6.5.1 et 6.6.2); autrement dit, il s'agit de montrer que $j$ est un morphisme \emph{plat au point} $x$; mais puisque $h\circ j=f$ est par hypothèse plat au point $x$ et que $h$ est étale, cela résulte de (18.4.10) et (18.4.11, (i)).

(ii) Il n'y a à prouver que la suffisance des conditions énoncées. Pour tout $x\in X$, on a donc dans un voisinage ouvert $U$ de $x$ une factorisation de $f|U$ ayant les propriétés considérées dans (i). Comme par hypothèse $f$ est plat et formellement non ramifié en \emph{tous} les points de $U$, le résultat de (i) s'applique non seulement à $x$ mais à tous les points de $U$; cela signifie que si $\mathscr J$ est l'Idéal quasi-cohérent de $\sh{O}_{X'}$ correspondant au sous-préschéma fermé de $X'$ associé à $j$, on a $\mathscr J_{j(z)}=0$ pour tout $z\in U$, donc $j$ est une immersion ouverte, et $f$ est par suite étale en tout point de $U$, donc en tout point de $X$.

\begin{corollary}[18.4.13]
\label{IV.18.4.13}
Soient $f:X\to Y$ un morphisme localement de type fini, $x$ un point de $X$, $y=f(x)$. Supposons que $y$ admette un voisinage ouvert qui soit un préschéma réduit et n'ayant qu'un nombre fini de composantes irréductibles. Alors, pour que $f$ soit étale au point $x$, il faut et il suffit que $f$ soit plat et formellement non ramifié au point $x$.
\end{corollary}

Il n'y a à démontrer que la suffisance de la condition. La question étant locale sur $X$ et $Y$, on peut supposer (18.4.7) que $f$ se factorise en $X\xrightarrow{j}X'\xrightarrow{h}Y$ où $h$ est étale et $j$ une immersion fermée, et en outre que $Y$ est réduit et n'a qu'un nombre fini de composantes irréductibles. Alors $X'$ est réduit (17.5.7) et, en remplaçant au besoin $X'$ par un voisinage ouvert de $j(x)$, on peut supposer que $X'$ n'a qu'un nombre fini de composantes irréductibles: en effet, on peut supposer $h$ quasi-fini (17.6.1) et comme les points maximaux de $X'$ sont au-dessus des points maximaux de $Y$ (2.3.4) leur nombre est fini. Tout revient à montrer, avec les notations de la démonstration de (18.4.12), que l'on a $\mathscr J_{x'}=0$ pour tous les points $x'$ d'un voisinage de $j(x)$ dans $X'$, sachant que $\mathscr J_{j(x)}=0$. Or, en remplaçant $X'$ par un voisinage affine de $j(x)$, on peut supposer que toutes les composantes irréductibles de $X'$ contiennent $j(x)$; si $X'=\Spec(A')$, et si $\mathfrak p'$ est l'idéal premier de $A'$ correspondant au point $j(x)$, le morphisme $\Spec(A'_{\mathfrak p'})\to\Spec(A')$ est dominant, donc l'homomorphisme correspondant $A'\to A'_{\mathfrak p'}$ est \emph{injectif} puisque $A'$ est réduit (I, 1.2.7). Si $\mathscr J=\widetilde{\mathfrak J}$, où $\mathfrak J$ est un idéal de $A'$, $\mathfrak J$ s'identifie donc à une partie de $\mathfrak J_{\mathfrak p'}$, et l'hypothèse $\mathfrak J_{\mathfrak p'}=0$ entraîne donc $\mathfrak J=0$.

\begin{corollary}[18.4.14]
\label{IV.18.4.14}
Soient $A$ un anneau local d'idéal maximal $\mathfrak m$, de corps résiduel $k$, $B$ une $A$-algèbre finie.

\oldpage[IV-4]{125}

\begin{enumerate}
  \item[(i)] Pour que $B$ soit une $A$-algèbre formellement non ramifiée, il faut et il suffit que $B\otimes_A k$ soit une $k$-algèbre étale.
  \item[(ii)] Pour que $B$ soit une $A$-algèbre étale, il faut et il suffit que $B\otimes_A k$ soit une $k$-algèbre étale et que $B$ soit un $A$-module plat (ce qui équivaut à dire ($0_{\mathrm I}$, 6.3.3) que pour tout idéal maximal $\mathfrak n$ de $B$ (nécessairement au-dessus de $\mathfrak m$), $B_{\mathfrak n}$ est un $A$-module plat).
\end{enumerate}
\end{corollary}

On n'a évidemment à prouver que la suffisance des conditions énoncées. Il est clair que (ii) résulte de (i) et de (18.4.12, (ii)), compte tenu de (17.1.2, (i)). Pour prouver (i) remarquons que si $B\otimes_A k$ est une $k$-algèbre étale, on a $\Omega^1_{(B\otimes_A k)/k}=0$ (17.2.1). Or on a $\Omega^1_{(B\otimes_A k)/k}=\Omega^1_{B/A}\otimes_B k$ (0, 20.5.5) et puisque $B$ est une $A$-algèbre de type fini, $\Omega^1_{B/A}$ est un $B$-module de type fini (0, 20.4.7). Mais comme $B$ est une $A$-algèbre finie, $\mathfrak mB$ est contenu dans le radical de $B$ (Bourbaki, \emph{Alg. comm.}, chap. V, \textsection~2, n\textsuperscript{o}~1, prop. 1), donc le lemme de Nakayama prouve que $\Omega^1_{B/A}=0$, et par suite $B$ est une $A$-algèbre formellement non ramifiée (17.2.1).

\subsection{Anneaux locaux henséliens\texorpdfstring{\protect\footnotemark}{}.}
\footnotetext{La notion d'anneau local hensélien est due à Azumaya, celle de hensélisation à Nagata, à qui l'on doit aussi les principaux résultats de cette théorie.}
\EGAReaderTarget{ega4.c04.subsec.18-5}{subsection:IV.18.5}

\begingroup
\newcommand{\shHom}{\mathscr{H}\!om}
\begin{env}[18.5.1]
\label{IV.18.5.1}
Soient $X$ un préschéma, $\sh{E}$ un $\sh{O}_X$-Module \emph{localement libre} de rang fini; le dual $\check{\sh{E}}=\shHom_{\sh{O}_X}(\sh{E},\sh{O}_X)$ est donc un $\sh{O}_X$-Module localement libre dont le rang en tout point de $X$ est égal à celui de $\sh{E}$ en ce point, et l'homomorphisme canonique $\sh{E}\to\shHom_{\sh{O}_X}(\check{\sh{E}},\sh{O}_X)=(\check{\sh{E}})^\vee$ est \emph{bijectif}. Pour tout morphisme $X'\to X$, posons $\sh{E}_{(X')}=\sh{E}\otimes_{\sh{O}_X}\sh{O}_{X'}$ qui est un $\sh{O}_{X'}$-Module localement libre, et considérons l'\emph{ensemble} $\Gamma(X',\sh{E}_{(X')})$ des sections de ce $\sh{O}_{X'}$-Module au-dessus de $X'$. Nous allons voir que l'on définit ainsi un \emph{foncteur contravariant représentable}
\begin{equation}
  {}^t\!\mathbf V_{\sh{E}}:X'\longmapsto\Gamma(X',\sh{E}_{(X')})
  \tag{18.5.1.1}\label{IV.18.5.1.1}
\end{equation}
de la catégorie des $X$-préschémas dans celle des \emph{ensembles} ($0_{\mathrm{III}}$, 8.1.8).

Tout d'abord, on a bien défini un foncteur, car si $f:X''\to X'$ est un $X$-morphisme de $X$-préschémas, on a $\sh{E}_{(X'')}=f^*(\sh{E}_{(X')})$, d'où ($0_{\mathrm I}$, 4.4.3.2) une application $\Gamma(X',\sh{E}_{(X')})\to\Gamma(X'',f_*(\sh{E}_{(X'')}))=\Gamma(X'',\sh{E}_{(X'')})$ qui achève de définir le foncteur ${}^t\!\mathbf V$. Montrons ensuite que le $X$-préschéma $\mathbf V(\check{\sh{E}})$ (II, 1.7.8) \emph{représente} le foncteur ${}^t\!\mathbf V$. Il est immédiat en effet que l'on a $(\sh{E}_{(X')})^\vee=(\check{\sh{E}})_{(X')}$, donc $\mathbf V(\check{\sh{E}})\times_X X'=\mathbf V((\check{\sh{E}})_{(X')})$; compte tenu de (I, 3.3.14), on est ramené à définir une bijection $\Gamma(\mathbf V(\check{\sh{E}})/X)\to\Gamma(X,\sh{E})$, et à vérifier que la bijection $\Gamma(\mathbf V((\check{\sh{E}})_{(X')})/X')\to\Gamma(X',\sh{E}_{(X')})$ est fonctorielle en $X'$. Or, on a une bijection canonique de $\Gamma(\mathbf V(\check{\sh{E}})/X)$ sur $\Hom_{\sh{O}_X}(\check{\sh{E}},\sh{O}_X)$ (II, 1.7.8), et la transposition $u\mapsto{}^t u$ est une bijection canonique de $\Hom_{\sh{O}_X}(\check{\sh{E}},\sh{O}_X)$ sur $\Hom_{\sh{O}_X}(\sh{O}_X,\sh{E})=\Gamma(X,\sh{E})$, en raison de l'identification de $(\check{\sh{E}})^\vee$ et de $\sh{E}$. La vérification de la propriété fonctorielle est immédiate.
\end{env}
\endgroup

\oldpage[IV-4]{126}

Notons encore que, conformément à la théorie générale ($0_{\mathrm{III}}$, 8.1.6), l'automorphisme identique de $\mathbf V(\check{\sh{E}})$ correspond canoniquement à une section $c$ de $\sh{E}_{(\mathbf V(\check{\sh{E}}))}$ au-dessus de $\mathbf V(\check{\sh{E}})$, c'est-à-dire (II, 1.4.1) à un homomorphisme de $\mathbf S_{\sh{O}_X}(\check{\sh{E}})$-Modules
\[
  u:\mathbf S_{\sh{O}_X}(\check{\sh{E}})\longrightarrow
  \mathbf S_{\sh{O}_X}(\check{\sh{E}})\otimes_{\sh{O}_X}\sh{E};
\]
pour tout ouvert affine $W$ de $X$, si l'on pose $\Gamma(W,\sh{O}_X)=A$, si on identifie $\Gamma(W,\mathbf S_{\sh{O}_X}(\check{\sh{E}}))$ à une algèbre de polynômes $C=A[T_1,\ldots,T_n]$, de sorte que les $T_i$ forment une base de $\Gamma(W,\check{\sh{E}})$, et si enfin on désigne par $(e_i)$ la base duale de $(T_i)$ dans $\Gamma(W,\sh{E})=\Gamma(W,\check{\sh{E}})^\vee$, on voit aussitôt que $u$ correspond à l'homomorphisme de $C$-modules tel que
\[
  u(1)=\sum_{i=1}^n T_i\otimes e_i.
\]
Si $X'=\Spec(A)$ est affine et tel que $\sh{E}_{(X')}$ soit isomorphe à $\sh{O}_{X'}^n$, $\Gamma(X',\sh{E}_{(X')})$ est un $A$-module libre de rang $n$; de façon imagée, on peut dire que l'objet $\mathbf V(\check{\sh{E}})$ représente « l'ensemble des points de l'espace affine tordu sur $X$ défini par $\sh{E}$ ».

\begin{env}[18.5.2]
\label{IV.18.5.2}
Rappelons (II, 1.7.8) que l'on a par définition $\mathbf V(\check{\sh{E}})=\Spec(\mathbf S_{\sh{O}_X}(\check{\sh{E}}))$. Soit $\mathscr J$ un Idéal quasi-cohérent de $\mathbf S_{\sh{O}_X}(\check{\sh{E}})$, de sorte que $\Spec(\mathbf S_{\sh{O}_X}(\check{\sh{E}})/\mathscr J)$ est un sous-préschéma fermé de $\mathbf V(\check{\sh{E}})$; nous allons l'interpréter comme \emph{représentant un foncteur} de la catégorie des $X$-préschémas dans celle des ensembles. Notons pour cela qu'une section $u\in\Gamma(X,\sh{E})$ s'identifie canoniquement à un $\sh{O}_X$-homomorphisme $u:\sh{O}_X\to\sh{E}$, auquel correspond par transposition un $\sh{O}_X$-homomorphisme ${}^t u:\check{\sh{E}}\to\sh{O}_X$, et par suite un homomorphisme de $\sh{O}_X$-Algèbres $v:\mathbf S_{\sh{O}_X}(\check{\sh{E}})\to\sh{O}_X$. Soit $\operatorname{Al}(X,\sh{E},\mathscr J)$ l'ensemble des $u\in\Gamma(X,\sh{E})$ tels que $\mathscr J$ soit \emph{contenu dans le noyau} de $v$; il résulte aussitôt de ces définitions que
\[
  X'\longmapsto\operatorname{Al}(X',\sh{E}_{(X')},\mathscr J\otimes_{\sh{O}_X}\sh{O}_{X'})
\]
est un foncteur représenté par $\Spec(\mathbf S_{\sh{O}_X}(\check{\sh{E}})/\mathscr J)$. Si $X'=\Spec(A)$ est affine et tel que $\sh{E}_{(X')}$ soit isomorphe à $\sh{O}_{X'}^n$, $\Gamma(X',\sh{E}_{(X')})$ peut s'identifier à l'ensemble $A^n$, et $\mathscr J\otimes_{\sh{O}_X}\sh{O}_{X'}$ à un $\sh{O}_{X'}$-Module de la forme $\widetilde{\mathfrak J}$, où $\mathfrak J$ est un idéal de l'anneau de polynômes $A[T_1,\ldots,T_n]$; l'ensemble $\operatorname{Al}(X',\sh{E}_{(X')},\mathscr J\otimes_{\sh{O}_X}\sh{O}_{X'})$ s'identifie alors à la partie de $A^n$ formée des points $(t_1,\ldots,t_n)$ tels que $P(t_1,\ldots,t_n)=0$ pour tous les polynômes $P\in\mathfrak J$; de façon imagée on peut donc dire que l'objet $\Spec(\mathbf S_{\sh{O}_X}(\check{\sh{E}})/\mathscr J)$ représente « le sous-ensemble algébrique de l'espace affine tordu $\Gamma(X,\sh{E})$ formé des points annulant l'idéal $\Gamma(X,\mathscr J)$ ». On note encore ce $X$-préschéma $\mathbf{Al}(\sh{E},\mathscr J)$. On notera que si $\mathscr J$ est un Idéal de type fini de $\mathbf S_{\sh{O}_X}(\check{\sh{E}})$, $\mathbf{Al}(\sh{E},\mathscr J)$ est un $X$-préschéma \emph{de présentation finie}, puisque $\mathbf S_{\sh{O}_X}(\check{\sh{E}})$ est une $\sh{O}_X$-Algèbre de présentation finie.
\end{env}

\begin{lemma}[18.5.3]
\label{IV.18.5.3}
Soient $S$ un préschéma, $f:X\to S$ un morphisme fini et localement libre (18.2.7). Considérons le foncteur contravariant de la catégorie des $S$-préschémas dans la catégorie des ensembles
\begin{equation}
  S'\longmapsto\operatorname{Of}(X\times_S S')
  \tag{18.5.3.1}\label{IV.18.5.3.1}
\end{equation}
où $\operatorname{Of}(X\times_S S')$ est l'ensemble des parties à la fois ouvertes et fermées de l'espace sous-jacent à $X\times_S S'$. Alors ce foncteur est représentable par un $S$-préschéma $\mathbf{Of}(X)$, qui est affine, étale et de présentation finie sur $S$.
\end{lemma}

\oldpage[IV-4]{127}

On a par hypothèse $X=\Spec(\mathscr B)$, où $\mathscr B=f_*(\sh{O}_X)$ est une $\sh{O}_S$-Algèbre finie et localement libre. Posons $X'=X\times_S S'$, $f'=f_{(S')}$, $\mathscr B'=\mathscr B\otimes_{\sh{O}_S}\sh{O}_{S'}=f'_*(\sh{O}_{X'})$, de sorte que $X'=\Spec(\mathscr B')$; il y a alors une bijection canonique, fonctorielle en $S'$, de l'ensemble $\operatorname{Of}(X')$ sur l'ensemble $\operatorname{Id}(\mathscr B')$ des \emph{idempotents} de l'anneau $\Gamma(S',\mathscr B')=\Gamma(X',\sh{O}_{X'})$. En effet, en vertu de l'\emph{équivalence} de la catégorie des $S'$-schémas affines sur $S'$ et de la catégorie opposée de la catégorie des $\sh{O}_{S'}$-Algèbres quasi-cohérentes (II, 1.2.7 et 1.3.1), il y a correspondance biunivoque canonique entre les décompositions de $X'$ en \emph{somme} $X'_1\amalg X'_2$ de deux sous-préschémas induits sur des ouverts de $X'$ et les décompositions de $\mathscr B'$ en \emph{composée directe} de deux Idéaux $\mathscr B'_1$, $\mathscr B'_2$; ces dernières à leur tour forment un ensemble en correspondance biunivoque canonique (et fonctorielle en $S'$) avec $\operatorname{Id}(\mathscr B')$. Il suffit donc de prouver le lemme pour le foncteur $S'\longmapsto\operatorname{Id}(\mathscr B')$.

Pour cela, nous allons montrer qu'il existe un Idéal de type fini $\mathscr J$ de $\mathbf S_{\sh{O}_S}(\check{\mathscr B})$ tel que $\operatorname{Id}(\mathscr B')$ soit de la forme $\operatorname{Al}(S',\mathscr B',\mathscr J\otimes_{\sh{O}_S}\sh{O}_{S'})$ (18.5.2). Notons à cet effet que puisque $\mathscr B$ est une $\sh{O}_S$-Algèbre, son image réciproque $\mathscr B_{(\mathbf V(\check{\mathscr B}))}$ est une $\sh{O}_{(\mathbf V(\check{\mathscr B}))}$-Algèbre, et l'on peut donc former le \emph{carré} $c^2$, dans cette Algèbre, de la section canonique $c$ (18.5.1); elle correspond canoniquement à un homomorphisme $u^{(2)}:\mathbf S_{\sh{O}_S}(\check{\mathscr B})\to\mathbf S_{\sh{O}_S}(\check{\mathscr B})\otimes_{\sh{O}_S}\mathscr B$ de $\mathbf S_{\sh{O}_S}(\check{\mathscr B})$-Modules, et on vérifie aussitôt que pour un ouvert affine $W$ de $S$, avec les notations de (18.5.1), $u^{(2)}$ correspond à l'homomorphisme de $C$-modules tel que
\[
  u^{(2)}(1)=\sum_k\Bigl(\sum_{i,j}c_{ijk}T_iT_j\Bigr)\otimes e_k,
\]
où $(c_{ijk})$ est la table de multiplication de l'algèbre $\Gamma(W,\mathscr B)$. Montrons que l'Idéal $\mathscr J$ de $\mathbf S_{\sh{O}_S}(\check{\mathscr B})$ engendré par le \emph{noyau} de l'homomorphisme $u-u^{(2)}$ répond à la question. Il suffit en effet de noter que l'idéal $\Gamma(W,\mathscr J)$ de $C$ est engendré par les polynômes $P_k(T_1,\ldots,T_n)=T_k-\sum_{i,j}c_{ijk}T_iT_j$ et que les idempotents de $\Gamma(W,\mathscr B)$ sont précisément les éléments $\sum_i t_i e_i$ de cette algèbre tels que $(t_1,\ldots,t_n)$ annule tous les polynômes $P_k$ ($1\leq i\leq n$). On déduit donc de (18.5.2) que le $S$-préschéma affine de présentation finie $\mathbf{Of}(X)=\mathbf{Al}(\mathscr B,\mathscr J)$ représente bien le foncteur (18.5.3.1). Il reste à prouver que $\mathbf{Of}(X)$ est étale sur $S$, ou, ce qui revient au même, qu'il est \emph{formellement étale} sur $S$. Mais si $S'$ est un $S$-préschéma, $S'_0$ un sous-préschéma fermé de $S'$ défini par un Idéal localement nilpotent de $\sh{O}_{S'}$ (et ayant donc même espace sous-jacent que $S'$), il est clair que $X'=X\times_S S'$ et $X'_0=X\times_S S'_0$ ont même espace sous-jacent, donc l'application canonique $\operatorname{Of}(X')\to\operatorname{Of}(X'_0)$ est \emph{bijective}, ce qui achève la démonstration (17.1.1).

\begin{proposition}[18.5.4]
\label{IV.18.5.4}
Soient $S$ un préschéma, $S_0$ un sous-préschéma fermé de $S$; considérons les propriétés suivantes:
\begin{enumerate}
  \item[a)] Pour tout morphisme fini $g:S'\to S$, l'application canonique
  \begin{equation}
    \operatorname{Of}(S')\longrightarrow\operatorname{Of}(S'\times_S S_0)
    \tag{18.5.4.1}\label{IV.18.5.4.1}
  \end{equation}
  (cf. (18.5.3)) est bijective.
  \item[a')] Pour tout morphisme fini et localement libre $g:S'\to S$, l'application (18.5.4.1) est bijective.

\oldpage[IV-4]{128}

  \item[b)] Pour tout morphisme étale et séparé $g:S'\to S$, l'application canonique
  \begin{equation}
    \Gamma(S'/S)\longrightarrow\Gamma(S'\times_S S_0/S_0)
    \tag{18.5.4.2}\label{IV.18.5.4.2}
  \end{equation}
  est bijective.
\end{enumerate}
La condition $b)$ entraîne $a')$; si de plus $S$ est quasi-compact et quasi-séparé, la condition $a)$ entraîne $b)$.
\end{proposition}

Prouvons d'abord que $b)$ entraîne $a')$. Supposons donc $b)$ vérifiée, et soit $g:S'\to S$ un morphisme fini et localement libre; posons $S'_0=S'\times_S S_0$, de sorte que $g_0=g_{(S_0)}:S'_0\to S_0$ est fini et localement libre. Alors il résulte de (18.5.3) que $P=\mathbf{Of}(S')$ est un $S$-préschéma \emph{étale et séparé}; en outre, la définition du foncteur $\operatorname{Of}$ montre aussitôt que si l'on pose $P_0=\mathbf{Of}(S'_0)$ (pour la catégorie des $S_0$-préschémas), on a $P_0=P\times_S S_0$. Cela étant, on a par définition le diagramme commutatif
\[
\xymatrix{
  \Gamma(P/S)\ar[r]\ar[d]^{\sim} & \Gamma(P_0/S_0)\ar[d]^{\sim}\\
  \operatorname{Of}(S')\ar[r] & \operatorname{Of}(S'_0)
}
\]
où les flèches verticales sont les bijections canoniques. Comme l'hypothèse $b)$, appliquée au morphisme $P\to S$, entraîne que la première ligne est une bijection, il en est de même de la seconde, ce qui établit notre assertion.

Avant de prouver que $a)$ entraîne $b)$ lorsque $S$ est quasi-compact et quasi-séparé, nous établirons le

\begin{lemma}[18.5.4.3]
\label{IV.18.5.4.3}
Si $S$ et $S_0$ vérifient la condition $a)$ de (18.5.4), alors, pour tout morphisme fini $g:S'\to S$, $S'$ est l'unique voisinage de $S'_0=g^{-1}(S_0)=S'\times_S S_0$ dans $S'$.
\end{lemma}

En effet, il revient au même de dire que si $T'$ est une partie fermée de $S'$ telle que $T'\cap S'_0=\varnothing$, alors $T'=\varnothing$. Or, si l'on désigne encore par $T'$ un sous-préschéma fermé de $S'$ ayant $T'$ pour espace sous-jacent, le morphisme composé $h:T'\to S'\xrightarrow{g}S$ est fini et $h^{-1}(S_0)$ est vide; la condition $a)$ appliquée au morphisme $h$ entraîne que $T'$ est nécessairement vide.

Ce lemme étant démontré, prouvons d'abord que, sous l'hypothèse $a)$, l'application (18.5.4.2) est \emph{injective}. En effet, si $u'$, $u''$ sont deux $S$-sections de $S'$, le fait que le morphisme $S'\to S$ soit non ramifié entraîne que le préschéma des coïncidences de $u'$ et $u''$ est induit sur un ouvert $U$ de $S$ (17.4.6). Si les restrictions à $S_0$ de $u'$ et $u''$ sont les mêmes, le fait que $u'$ et $u''$ sont des immersions ouvertes (17.4.1) entraîne que $U$ contient $S_0$, donc est égal à $S$ en vertu du lemme (18.5.4.3) appliqué au cas où $S'=S$.

\oldpage[IV-4]{129}

Reste à montrer que sous l'hypothèse $a)$, l'application (18.5.4.2) est \emph{surjective} ($S$ étant quasi-compact et quasi-séparé). Soit donc $u_0:S_0\to S'_0$ une $S_0$-section de $S'_0$; $u_0(S_0)$ étant quasi-compact dans $S'$, peut être recouvert par un nombre fini d'ouverts affines $V_i$ tels que la restriction $g|V_i$ soit un morphisme de type fini; si $V$ est la réunion des $V_i$, on en conclut que $g|V:V\to S$ est un morphisme \emph{de présentation finie}, étant séparé et localement de présentation finie par hypothèse (1.6.1). Remplaçant $S'$ par $V$, on peut donc supposer que $g$ est de présentation finie. Comme $S$ est quasi-compact et quasi-séparé et $g$ quasi-fini et séparé (17.6.1), il résulte du « Main theorem » (8.12.6) que $g$ se factorise en $S'\xrightarrow{j}S''\xrightarrow{f}S$, où $j$ est une immersion ouverte et $f$ un morphisme \emph{fini}. Posons $S''_0=S''\times_S S_0$, $j_0=j_{(S_0)}:S'_0\to S''_0$, qui est une immersion ouverte et $f_0=f_{(S_0)}:S''_0\to S_0$, qui est un morphisme fini. Alors $u_0$ est aussi une $S_0$-section de $S''_0$. Comme $g_0:S'_0\to S_0$ est étale, $u_0$ est une immersion ouverte de $S_0$ dans $S'_0$ (17.4.1), donc $u_0(S_0)$ est ouvert dans $S'_0$, et \emph{a fortiori} dans $S''_0$; mais d'autre part, comme $f_0$ est un morphisme fini, donc séparé, $u_0$ est une immersion fermée de $S_0$ dans $S''_0$ (I, 5.4.6), donc $X_0=u_0(S_0)$ est \emph{à la fois ouvert et fermé} dans $S''_0$. En vertu de l'hypothèse $a)$, il existe un sous-ensemble à la fois ouvert et fermé $X$ de $S''$ tel que $X\cap S''_0=X_0$. Montrons d'abord que le morphisme $f:S''\to S$ est \emph{étale aux points de $X$}: en effet, l'ensemble $U$ des points de $X$ où $f|X$ est étale est ouvert et contient par hypothèse $X_0\subset S'_0$. Mais le lemme (18.5.4.3) appliqué au morphisme fini $f|X$ prouve que $U=X$. D'autre part, $S'\cap X$ est ouvert dans $X$ et contient $X_0$ par hypothèse, donc le même raisonnement prouve que $S'\cap X=X$, c'est-à-dire $X\subset S'$. Il reste à montrer que, pour tout $s\in S$, le nombre géométrique $n(s)$ de points de $X\cap f^{-1}(s)$ est égal à $1$, car il en résultera que $f|X$ est radiciel et surjectif, et comme $f|X$ est étale, on aura montré (17.9.1) que $f|X$ est un isomorphisme de l'ouvert $X\subset S'$ sur $S$, dont l'isomorphisme réciproque $u$ sera la $S$-section cherchée prolongeant $u_0$. Mais comme $f|X$ est étale et \emph{fini}, $s\mapsto n(s)$ est continue dans $S$ (18.2.8), et comme $X\cap S''_0=X_0$, on a $n(s)=1$ dans $S_0$; l'ensemble des points $s\in S$ tels que $n(s)=1$ étant ouvert dans $S$ et contenant $S_0$, il est égal à $S$ par (18.5.4.3). C.Q.F.D.

\begin{remark}[18.5.4.4]
\label{IV.18.5.4.4}
On peut montrer que l'énoncé (18.5.4) reste valable lorsque, dans la condition $b)$, on suppose seulement le morphisme $g$ étale (mais non nécessairement séparé) [43, exp. XII].
\end{remark}

\begin{definition}[18.5.5]
\EGAReaderTarget{ega4.c04.def.18-5-5}{IV.18.5.5}
On dit qu'un préschéma $S$ et un sous-préschéma fermé $S_0$ de $S$ forment un \emph{couple hensélien} s'ils vérifient la condition $a)$ de (18.5.4).
\end{definition}

Compte tenu de (I, 5.1.8), il revient au même de dire que $(S,S_0)$ est un couple hensélien ou que $(S_{\mathrm{red}},(S_0)_{\mathrm{red}})$ en est un.

\begin{proposition}[18.5.6]
\label{IV.18.5.6}
\begin{enumerate}
  \item[(i)] Si $(S,S_0)$ est un couple hensélien, alors, pour tout morphisme fini $f:S'\to S$, si $S'_0$ est le sous-préschéma $f^{-1}(S_0)$ de $S'$, le couple $(S',S'_0)$ est hensélien.
  \item[(ii)] Soient $S=\coprod_\alpha S^{(\alpha)}$ une somme de préschémas, $S_0$ un sous-préschéma fermé de $S$, somme des sous-préschémas fermés $S_0^{(\alpha)}$ des $S^{(\alpha)}$. Pour que le couple $(S,S_0)$ soit hensélien, il faut et il suffit que chacun des couples $(S^{(\alpha)},S_0^{(\alpha)})$ le soit.
\end{enumerate}
\end{proposition}

L'assertion (i) est conséquence immédiate de la définition, puisque pour tout morphisme fini $g:S''\to S'$, $f\circ g:S''\to S$ est un morphisme fini. De même, sous les conditions de (ii), pour qu'un morphisme $g:S'\to S$ soit fini, il faut et il suffit que chacune

\oldpage[IV-4]{130}

de ses restrictions $g^{(\alpha)}:S'^{(\alpha)}=g^{-1}(S^{(\alpha)})\to S^{(\alpha)}$ le soit, et si l'on pose $S'_0=g^{-1}(S_0)$, $S'^{(\alpha)}_0=(g^{(\alpha)})^{-1}(S^{(\alpha)}_0)$, il y a correspondance biunivoque entre les parties ouvertes et fermées $U$ (resp. $U_0$) de $S'$ (resp. $S'_0$) et les familles $(U^{(\alpha)})$ (resp. $U^{(\alpha)}_0$) où $U^{(\alpha)}$ (resp. $U^{(\alpha)}_0$) est une partie ouverte et fermée de $S'^{(\alpha)}$ (resp. $S'^{(\alpha)}_0$), d'où l'assertion (ii).

\begin{remark}[18.5.7]
\label{IV.18.5.7}
Soient $S=\Spec(A)$ un schéma affine, $S_0$ un sous-schéma fermé de $S$ défini par un idéal $\mathfrak J$ de $A$. Alors, si le couple $(S,S_0)$ est hensélien, l'idéal $\mathfrak J$ est nécessairement contenu dans le \emph{radical} de $A$. En effet, si $\mathfrak m$ est un idéal maximal de $A$, $\mathfrak m$ doit appartenir à $V(\mathfrak J)=S_0$, en vertu de (18.5.4.3), autrement dit on doit avoir $\mathfrak J\subset\mathfrak m$, d'où la conclusion. En particulier, supposons que $S_0$ soit \emph{réduit à un point}, c'est-à-dire que l'idéal $\mathfrak J$ soit \emph{maximal}; alors $\mathfrak J$ doit être le radical de $A$, autrement dit $A$ doit être un anneau \emph{local}, $S_0$ étant l'unique point fermé de $\Spec(A)$.
\end{remark}

\begin{definition}[18.5.8]
\EGAReaderTarget{ega4.c04.def.18-5-8}{IV.18.5.8}
On dit qu'un anneau $A$ est \emph{hensélien} s'il est semi-local et si, en désignant par $\mathfrak r$ le radical de $A$, le couple $(\Spec(A),\Spec(A/\mathfrak r))$ est hensélien. On appelle \emph{schéma local hensélien} un schéma isomorphe au spectre d'un anneau local hensélien.
\end{definition}

\begin{proposition}[18.5.9]
\label{IV.18.5.9}
\begin{enumerate}
  \item[(i)] Pour qu'un anneau semi-local $A$ soit hensélien, il faut et il suffit qu'il soit composé direct d'anneaux locaux henséliens.
  \item[(ii)] Pour qu'un anneau local $A$ soit hensélien, il faut et il suffit que toute $A$-algèbre finie $B$ soit isomorphe à un produit d'anneaux locaux.
\end{enumerate}
\end{proposition}

(i) En effet, la définition, appliquée à $S=\Spec(A)$ et $S_0=\Spec(A/\mathfrak r)$ montre, puisque $S_0$ est un ensemble fini discret et fermé dans $S$, que $S$ est réunion d'un nombre fini de parties ouvertes et fermées $S_i$ ($1\leq i\leq n$) deux à deux disjointes, dont chacune contient exactement un des idéaux maximaux $\mathfrak m_i$ de $A$; la conclusion résulte de (18.5.6, (ii)) et de la remarque (18.5.7).

(ii) Tout morphisme fini $S'\to\Spec(A)$ est de la forme $\Spec(B)\to\Spec(A)$, où $B$ est une $A$-algèbre finie. Si $k$ est le corps résiduel de $A$, $\Spec(B\otimes_A k)$ est un spectre d'anneau artinien, donc fini et discret. Dire que le couple $(\Spec(A),\Spec(k))$ est hensélien signifie donc que $B$ est composé direct d'anneaux $A_i$ tels que $\Spec(A_i\otimes_A k)$ soit réduit à un point, c'est-à-dire que $A_i$ (qui est une $A$-algèbre \emph{finie}) ne doit avoir qu'un seul idéal maximal (Bourbaki, \emph{Alg. comm.}, chap. V, \textsection~2, n\textsuperscript{o}~1, prop. 1).

L'étude des anneaux henséliens est donc essentiellement ramenée à celle des anneaux \emph{locaux} henséliens.

\begin{proposition}[18.5.10]
\label{IV.18.5.10}
Si $A$ est un anneau hensélien, toute $A$-algèbre finie $B$ est un anneau hensélien (donc composé direct d'anneaux locaux henséliens (18.5.9)).
\end{proposition}

En effet, si $\mathfrak r'$ est le radical de $B$, l'image réciproque de $\mathfrak r'$ dans $A$ est le radical $\mathfrak r$ de $A$, et tout idéal de $B$ au-dessus d'un idéal maximal de $A$ est un idéal maximal de $B$, donc l'ensemble $V(\mathfrak r')$ dans $\Spec(B)$ est l'image réciproque de l'ensemble $V(\mathfrak r)$ dans $\Spec(A)$. La proposition est alors une conséquence de (18.5.6, (i)).

\begin{theorem}[18.5.11]
\label{IV.18.5.11}
Soient $A$ un anneau local, $\mathfrak m$ son idéal maximal. Les conditions suivantes sont équivalentes:
\begin{enumerate}
  \item[a)] $A$ est hensélien, autrement dit, toute $A$-algèbre finie $B$ est isomorphe à un produit d'anneaux locaux.

\oldpage[IV-4]{131}

  \item[a')] La condition $a)$ est satisfaite pour toutes les $A$-algèbres $B$ de la forme $A[T]/F.A[T]$, où $F\in A[T]$ est un polynôme unitaire.
  \item[b)] Soient $S=\Spec(A)$, $S_0=\Spec(k)$. Pour tout morphisme étale $g:S'\to S$, si l'on pose $S'_0=S'\otimes_A k$, toute $S_0$-section $u_0$ de $S'_0$ est la restriction d'une $S$-section $u$ de $S'$.
  \item[c)] Pour tout morphisme $f:X\to S$, séparé et localement de type fini, et tout point $x\in X$ tel que $f(x)$ soit égal au point fermé $s$ de $S$ et que $f$ soit quasi-fini au point $x$ (Err$_{\mathrm{III}}$, 20), $X$ est somme de deux préschémas $X'$, $X''$ tels que $X'=\Spec(\sh{O}_{X,x})$ et que $f|X':X'\to S$ soit un morphisme fini.
  \item[c')] Pour tout morphisme localement de type fini $f:X\to S$, et tout point $x\in X$ tel que $f$ soit quasi-fini au point $x$, et que $f(x)$ soit égal au point fermé $s$ de $S$, $\sh{O}_{X,x}$ est une algèbre finie sur $\sh{O}_{S,s}=A$.
  \item[c'')] Pour tout morphisme localement de présentation finie $f:X\to S$, et tout point $x\in X$ tel que $f$ soit quasi-fini au point $x$, et que $f(x)$ soit égal au point fermé $s$ de $S$, $\sh{O}_{X,x}$ est une algèbre finie et de présentation finie sur $\sh{O}_{S,s}=A$.
\end{enumerate}
\end{theorem}

Notons d'abord que la condition $c')$ (resp. $c'')$ est équivalente à la même condition où on suppose de plus $f$ \emph{séparé}, la question étant locale sur $X$. De même, la condition $b)$ est équivalente à la même condition où l'on suppose de plus $g$ \emph{séparé}: en effet, il suffit d'appliquer cette dernière à la restriction de $g$ à un voisinage ouvert affine du point $u_0(S_0)$ dans $S'$. Bornons-nous donc désormais au cas où, dans $b)$, $c')$ et $c'')$, les morphismes donnés sont \emph{séparés}.

Le fait que $a)$ implique $b)$ et que $b)$ implique $a')$ est un cas particulier de (18.5.4). Montrons en outre que $a')$ implique $a)$. Il s'agit de prouver que si $e_0$ est un idempotent de $C=B\otimes_A k$, il existe un idempotent $e\in B$ dont $e_0$ est l'image canonique. Si $b$ est un élément de $B$ dont l'image dans $C$ est $e_0$, la sous-$A$-algèbre $A[b]=B'$ de $B$ est finie, et l'image canonique $C'$ de $B'\otimes_A k$ dans $C$ contient $e_0$. Or, $B'\otimes_A k$ est une $k$-algèbre finie, donc composée directe de $k$-algèbres locales finies, et par suite $e_0$ est l'image dans $C'$ d'un idempotent $e'_0$ de $B'\otimes_A k$. On est ainsi ramené au cas où $B$ est monogène, et par suite isomorphe à une $A$-algèbre quotient d'une algèbre de la forme $A[T]/F.A[T]$, où $F$ est un polynôme unitaire. Or, en vertu de $a')$, $A[T]/F.A[T]$ est composée directe d'anneaux locaux, donc il en est de même de toutes ses algèbres quotients, ce qui achève de prouver l'existence de l'idempotent $e$.

Il est immédiat que $c)$ entraîne $a)$, comme on le voit en raisonnant par récurrence sur le nombre d'idéaux maximaux de l'anneau semi-local $B$. Pour voir que $a)$ implique $c)$, on peut, en vertu de (13.1.4), se borner au cas où $f$ est un morphisme affine et quasi-fini. Alors, par application du « Main theorem » (8.12.8), $f$ peut s'écrire comme un morphisme composé
\[
  X\xrightarrow{j}Y\xrightarrow{g}S,
\]
où $g$ est un morphisme fini et $j$ une immersion ouverte. Comme $Y=\Spec(B)$, où $B$ est une $A$-algèbre finie, il résulte de $a)$ que $B$ est composée directe d'anneaux locaux, qui sont évidemment des $A$-algèbres finies, et $\sh{O}_{X,x}$ s'identifie à l'un de ces anneaux locaux puisque $f(x)=s$; en outre, tout ouvert de $Y$ contenant $x$ contient nécessairement $\Spec(\sh{O}_{X,x})$.

Il est trivial que $c)$ implique $c')$, en vertu de la remarque du début. Prouvons que $c')$ implique $c'')$. Supposons en effet $c')$ vérifiée, et prouvons que sous les conditions

\oldpage[IV-4]{132}

de $c'')$, l'ensemble $Z=\Spec(\sh{O}_{X,x})$ s'identifie alors à une partie \emph{ouverte et fermée} de $X$, ce qui établira $c'')$ (I, 2.4.2). En premier lieu, le morphisme composé $Z\xrightarrow{j}X\xrightarrow{f}S$, où $j$ est le morphisme canonique (I, 2.4.1), est fini et de présentation finie par hypothèse, et comme $f$ est séparé et localement de présentation finie, $j$ est aussi un morphisme fini et de présentation finie ((II, 6.1.5) et (I, 4.3)); $j$ est par suite un morphisme fermé (II, 6.1.10), ce qui prouve que $Z$ est fermé dans $X$. Il résulte alors de (I, 2.4.2) et (I, 4.2.2) que $j$ est une immersion fermée. Mais si $\mathscr J$ est l'idéal de $\sh{O}_X$ définissant $Z$, on a alors par hypothèse $\mathscr J_x=0$, donc aussi $\mathscr J_z=0$ en tout point $z$ d'un voisinage ouvert $V$ de $x$ dans $X$, puisque par hypothèse $\mathscr J$ est un $\sh{O}_X$-Module quasi-cohérent de type fini ((1.4.7) et ($0_{\mathrm I}$, 5.2.2)). Or, un tel voisinage contient $Z$, donc $Z$ est ouvert dans $X$ puisque $\mathscr J|V=0$.

Enfin, $c'')$ implique $a')$: en effet, si $B=A[T]/F.A[T]$, le morphisme $X=\Spec(B)\to\Spec(A)=Y$ est fini et de présentation finie, et la démonstration précédente montre que $B$ est composée directe des $A$-algèbres finies $\sh{O}_{X,x_i}$, où les $x_i$ sont les points de la fibre du point fermé de $Y$.

\begin{corollary}[18.5.12]
\label{IV.18.5.12}
Soient $A$ un anneau semi-local, $\mathfrak r$ son radical; posons $S=\Spec(A)$, $S_0=\Spec(A/\mathfrak r)$. Pour que $A$ soit hensélien, il faut et il suffit que, pour tout morphisme fini $f:X\to S$ et tout morphisme étale et séparé $g:Y\to S$, si l'on pose $X_0=X\times_S S_0$ et $Y_0=Y\times_S S_0$, l'application canonique
\[
  \Hom_S(X,Y)\longrightarrow\Hom_{S_0}(X_0,Y_0)
\]
soit bijective.
\end{corollary}

La suffisance de la condition résulte de l'équivalence de $a)$ et $b)$ dans (18.5.11), en appliquant cette condition au cas où $f=1_S$. Pour voir que la condition est nécessaire, notons que si $A$ est hensélien, le couple $(X,X_0)$ est hensélien par (18.5.6, (i)); en outre $\Hom_S(X,Y)=\Gamma(X\times_S Y/X)$ et $\Hom_{S_0}(X_0,Y_0)=\Gamma(X_0\times_{S_0}Y_0/X_0)$; comme $X\times_S Y$ est étale et séparé sur $X$, la conclusion résulte de ce que $a)$ implique $b)$ dans (18.5.4).

\begin{remark}[18.5.13]
\label{IV.18.5.13}
Les conditions équivalentes $a)$, $a')$, $b)$ du théorème (18.5.11) sont encore équivalentes à la suivante (« lemme de Hensel »):
\begin{enumerate}
  \item[a'')] Pour tout polynôme unitaire $F\in A[T]$, d'image canonique $F_0\in k[T]$, et toute décomposition $F_0=G_0H_0$ de $F_0$ en un produit de deux polynômes unitaires étrangers $G_0$, $H_0$ de $k[T]$, il existe un couple unique $(G,H)$ de polynômes unitaires de $A[T]$ possédant les propriétés suivantes: $G_0$ et $H_0$ sont les images canoniques respectives de $G$ et $H$, on a $F=GH$, et l'idéal de $A[T]$ engendré par $G$ et $H$ est égal à $A[T]$.
\end{enumerate}
\end{remark}

Nous établirons d'abord le lemme suivant:

\begin{lemma}[18.5.13.1]
\label{IV.18.5.13.1}
Soient $A$ un anneau local de corps résiduel $k$, $F\in A[T]$ un polynôme unitaire, $B$ la $A$-algèbre $A[T]/F.A[T]$. Il existe une correspondance canonique entre les décompositions de $B$ en composée directe de deux $A$-algèbres quotients $B'$, $B''$ et les décompositions $F=GH$ de $F$ en produit de deux polynômes unitaires $G$, $H$ de $A[T]$, tels que l'idéal engendré par $G$ et $H$ soit égal à $A[T]$; les algèbres quotients $B'$, $B''$ correspondant à un tel couple de polynômes $G$, $H$ sont respectivement $A[T]/H.A[T]$ et $A[T]/G.A[T]$.
\end{lemma}

\oldpage[IV-4]{133}

Si $F=GH$ et si $G$ et $H$ engendrent l'idéal $A[T]$, il y a deux polynômes $P$, $Q$ de $A[T]$ tels que $1=PG+QH$. On en déduit que l'intersection des idéaux principaux $\mathfrak a=G.A[T]$ et $\mathfrak b=H.A[T]$ est égale à $\mathfrak c=F.A[T]$: en effet, si $R\in G.A[T]\cap H.A[T]$, on peut écrire $R=PRG+QRH$; or $RH$ (resp. $RG$) est un multiple de $F$ puisque $R$ est un multiple de $G$ (resp. $H$), donc $R\in F.A[T]$. Comme $A[T]=\mathfrak a+\mathfrak b$, $A[T]/\mathfrak c$ est somme directe des idéaux $\mathfrak a/(\mathfrak a\cap\mathfrak b)$ et $\mathfrak b/(\mathfrak a\cap\mathfrak b)$, canoniquement isomorphes respectivement à $A[T]/\mathfrak b$ et $A[T]/\mathfrak a$.

Inversement, supposons donnée une décomposition de $B$ en composée directe de deux $A$-algèbres $B'$, $B''$, qui s'identifient canoniquement à deux idéaux $e'B$, $e''B$ de $B$, correspondant à une décomposition $1=e'+e''$ de $1$ en idempotents orthogonaux $e'$, $e''$ de $B$. Posons encore $B_0=B\otimes_A k=k[T]/F_0.k[T]$, où $F_0$ est l'image canonique de $F$ dans $k[T]$, de même degré $n$ que $F$; si $e'_0$, $e''_0$ sont les images canoniques de $e'$, $e''$ dans $B_0$, ce sont deux idempotents orthogonaux tels que $1=e'_0+e''_0$, et $B_0$ est donc composée directe de $B'_0=e'_0B_0$ et $B''_0=e''_0B_0$. Soient $t$ et $t_0$ les images canoniques de $T$ dans $B$ et $B_0$; comme $B'_0$ (resp. $B''_0$) est une $k$-algèbre finie engendrée par $t'_0=e'_0t_0$ (resp. $t''_0=e''_0t_0$), elle admet une base de la forme $\{e'_0,t'_0,t_0'{}^2,\ldots,t_0'{}^{s-1}\}$ (resp. $\{e''_0,t''_0,t_0''{}^2,\ldots,t_0''{}^{r-1}\}$) avec $r+s=n$. D'autre part, $B$ étant un $A$-module libre (de base $\{1,t,\ldots,t^{n-1}\}$), $B'$ et $B''$ sont des $A$-modules projectifs, donc libres puisque $A$ est un anneau local (Bourbaki, \emph{Alg. comm.}, chap. II, \textsection~5, n\textsuperscript{o}~3, cor. de la prop. 5); il résulte donc de ce qui précède et de Bourbaki, \emph{Alg. comm.}, chap. II, \textsection~3, n\textsuperscript{o}~3, prop. 5, que si l'on pose $t'=e't$, $t''=e''t$, $\{e',t',\ldots,t'^{s-1}\}$ (resp. $\{e'',t'',\ldots,t''^{r-1}\}$) est une base du $A$-module $B'$ (resp. $B''$). Il y a donc un polynôme unitaire $H$ (resp. $G$) de degré $s$ (resp. $r$) de $A[T]$ tel que $e'H(t')=0$ et $e''G(t'')=0$; comme $t^h=t'^h+t''^h$ pour tout entier $h\geq1$, et $t'^h=e't'^h$, $t''^h=e''t''^h$, on a aussi $G(t)=e'G(t')$ et $H(t)=e''H(t'')$, d'où $G(t)H(t)=0$; on en conclut que le polynôme $G(T)H(T)$ est divisible par $F(T)$; mais comme les degrés de ces deux polynômes unitaires sont les mêmes, on a $GH=F$. En outre, $B'$ (resp. $B''$) est isomorphe à $A[T]/H.A[T]$ (resp. $A[T]/G.A[T]$). Enfin, il y a deux polynômes $R$, $S$ de $A[T]$ tels que $e'=R(t)$ et $e''=S(t)$; comme $R(t)=e'R(t')+e''R(t'')$, on a nécessairement $e''R(t'')=0$, et de même $e'S(t')=0$, de sorte que, par définition de $G$ et $H$, $R=QH$ et $S=PG$, où $P$, $Q$ appartiennent à $A[T]$; la relation $1=R(t)+S(t)$ dans $B$ donne donc par définition $PG+QH=1+LF$ pour un certain polynôme $L\in A[T]$, et comme $F=GH$, cela prouve que l'idéal engendré par $G$ et $H$ est $A[T]$, et achève de démontrer le lemme.

Ce lemme étant établi, il suffit de l'appliquer à l'anneau local $A$ d'une part, au corps $k$ de l'autre, pour voir aussitôt que les conditions $a')$ et $a'')$ sont équivalentes.

\begin{proposition}[18.5.14]
\label{IV.18.5.14}
Tout anneau semi-local $A$, séparé et complet pour la topologie $\mathfrak r$-préadique (où $\mathfrak r$ est le radical de $A$) est hensélien.
\end{proposition}

En effet, $A$ est composé direct d'anneaux locaux séparés et complets (Bourbaki, \emph{Alg. comm.}, chap. III, \textsection~2, n\textsuperscript{o}~13, cor. de la prop. 19), donc on est ramené au cas où $A$ est un anneau local. Vérifions le critère $a')$ de (18.5.11). Comme $B$ est un $A$-module libre de type fini, il est évidemment séparé et complet pour la topologie $\mathfrak r$-préadique,

\oldpage[IV-4]{134}

qui est aussi la topologie $\mathfrak s$-préadique, où $\mathfrak s$ est le radical de l'anneau semi-local $B$, car $B/\mathfrak rB$ est un anneau artinien de radical $\mathfrak s/\mathfrak rB$. On sait alors (Bourbaki, \emph{Alg. comm.}, chap. III, \textsection~2, n\textsuperscript{o}~13, cor. de la prop. 19) que $B$ est composé direct d'anneaux locaux.

\begin{proposition}[18.5.15]
\label{IV.18.5.15}
Soient $A$ un anneau hensélien, $\mathfrak r$ son radical. Alors le foncteur $B\longmapsto B/\mathfrak rB$ est une équivalence de la catégorie des $A$-algèbres finies et étales avec la catégorie des $(A/\mathfrak r)$-algèbres finies et étales.
\end{proposition}

Le fait que le foncteur de l'énoncé soit pleinement fidèle est un cas particulier de (18.5.12). Pour montrer que ce foncteur est une équivalence de catégories, on peut se borner au cas où $A$ est local; il suffit alors d'appliquer (18.1.1) (pour les morphismes étales), à $S=\Spec(A)$ et $S_0=\Spec(A/\mathfrak r)$, réduit à un seul point.

\begin{remarks}[18.5.16]
\label{IV.18.5.16}
\begin{enumerate}
  \item[(i)] Nous ignorons si la proposition (18.5.15) se généralise à un couple hensélien $(S,S_0)$, même lorsque $S$ est affine et noethérien.
  \item[(ii)] Soient $A$ un anneau noethérien, $\mathfrak J$ un idéal de $A$ tel que $A$ soit \emph{séparé et complet} pour la topologie $\mathfrak J$-préadique. Alors le couple $(\Spec(A),\Spec(A/\mathfrak J))$ est hensélien: en effet, pour toute $A$-algèbre finie $B$, $B$ est séparé et complet pour la topologie $\mathfrak J$-préadique ($0_{\mathrm I}$, 7.3.6). Remplaçant $B$ par $A$ et $\mathfrak J B$ par $\mathfrak J$, tout revient donc à voir que l'application qui, à tout idempotent de $A$, fait correspondre sa classe mod. $\mathfrak J$, est \emph{bijective}. Or on a $A=\varprojlim (A/\mathfrak J^n)$. Notons $\operatorname{Idem}(A)$ l'ensemble des idempotents de $A$, et pour tout homomorphisme d'anneaux $\varphi:A\to B$, soit $\operatorname{Idem}(\varphi)$ l'application de $\operatorname{Idem}(A)$ dans $\operatorname{Idem}(B)$ restriction de $\varphi$; il résulte de la définition de la limite projective que l'on a
  \[
    \operatorname{Idem}(A)=\varprojlim_n\operatorname{Idem}(A/\mathfrak J^n)
  \]
  pour les applications $\psi_{nm}:\operatorname{Idem}(A/\mathfrak J^m)\to\operatorname{Idem}(A/\mathfrak J^n)$ restrictions des applications canoniques $A/\mathfrak J^m\to A/\mathfrak J^n$. Mais puisque $\Spec(A/\mathfrak J^n)\to\Spec(A/\mathfrak J^m)$ est un homéomorphisme, les $\psi_{nm}$ sont des bijections (comme on l'a vu dans la démonstration de (18.5.3)); cela prouve donc notre assertion.
\end{enumerate}
\end{remarks}

\begin{theorem}[18.5.17]
\label{IV.18.5.17}
Soient $A$ un anneau local hensélien, $S=\Spec(A)$, $s$ le point fermé de $S$. Pour tout morphisme lisse $f:X\to S$, l'application canonique
\[
  \Gamma(X/S)\longrightarrow\Gamma(X_s/k(s))
\]
(où $X_s=f^{-1}(s)$) est surjective.
\end{theorem}

La donnée d'une $k(s)$-section de $X_s$ équivaut à celle d'un point $x\in X$ au-dessus de $s$ rationnel sur $k(s)$, et il s'agit de prouver qu'il existe une $S$-section $u:S\to X$ telle que $u(s)=x$. Compte tenu de (17.16.3, (i)), on peut supposer que $f$ est étale. Alors la conclusion résulte du critère (18.5.11, $b)$). (Le lecteur notera qu'en vertu de ce critère, la validité de (18.5.17) pour un anneau local donné $A$ est nécessaire et suffisante pour que $A$ soit hensélien).

\begin{remark}[18.5.18]
\label{IV.18.5.18}
Les conditions de (18.5.11) sont encore équivalentes à la suivante:
\begin{enumerate}
  \item[d)] Pour tout morphisme localement de type fini $f:X\to S$ et tout point $x\in X$ tel que $f(x)$ soit le point fermé $s$ de $S$ et que $\sh{O}_{X,x}/\mathfrak m_s\sh{O}_{X,x}$ soit un corps canoniquement isomorphe à $k(s)=k$, il existe un voisinage ouvert $U$ de $x$ dans $X$ tel que $f|U$ soit une immersion fermée.
\end{enumerate}
\end{remark}

\oldpage[IV-4]{135}

Il est immédiat que $d)$ implique la condition $b)$ de (18.5.11) car une immersion fermée étale est une immersion ouverte (18.9.1). Inversement, supposons vérifiées les conditions de (18.5.11) et prouvons $d)$. L'hypothèse de $d)$ implique que $f$ est quasi-fini au point $x$ (Err$_{\mathrm{III}}$, 20), en vertu de (13.1.4). Donc en vertu de la condition $c)$ de (18.5.11), $\sh{O}_{X,x}$ est une $A$-algèbre finie et $\Spec(\sh{O}_{X,x})$ un voisinage ouvert $U$ de $x$ dans $X$. En outre, comme $\sh{O}_{X,x}/\mathfrak m_s\sh{O}_{X,x}$ est isomorphe à $k=A/\mathfrak m_s$, le lemme de Nakayama prouve que l'homomorphisme $A\to\sh{O}_{X,x}$ est surjectif, donc que $f|U$ est une immersion fermée.

\begin{proposition}[18.5.19]
\label{IV.18.5.19}
Soient $A$ un anneau local noethérien et hensélien, $S=\Spec(A)$, $s$ le point fermé de $S$, $f:X\to S$ un morphisme propre; posons $X_s=f^{-1}(s)$. Alors l'application $Y\longmapsto Y\cap X_s$ est une bijection de l'ensemble des composantes connexes de $X$ sur l'ensemble des composantes connexes de $X_s$.
\end{proposition}

En considérant un sous-préschéma fermé de $X$ ayant pour espace sous-jacent une composante connexe de $X$, on est ramené à prouver que si $X$ est \emph{connexe et non vide}, alors $X_s$ est \emph{connexe et non vide}. Le fait que $X_s$ soit non vide résulte de (II, 7.2.1); pour prouver que $X_s$ est connexe, raisonnons par l'absurde, en considérant la factorisation de Stein $f:X\xrightarrow{f'}S'\xrightarrow{g}S$ du morphisme propre $f$ (III, 4.3.3); par hypothèse, l'ensemble discret fini $g^{-1}(s)$ contiendrait au moins deux points. Puisque $A$ est hensélien et $g$ séparé et de type fini, il résulterait alors de (18.5.11, $c)$ que $S'$ serait somme de deux préschémas non vides $S'_1$, $S'_2$ (puisque l'intersection de l'un d'eux avec $g^{-1}(s)$ est réduite à un seul point). Puisque $f'$ est surjectif (III, 4.3.1), on en conclurait que $X$ est somme de deux préschémas non vides, contrairement à l'hypothèse.

\subsection{Hensélisation.}
\EGAReaderTarget{ega4.c04.subsec.18-6}{subsection:IV.18.6}

\begin{env}[18.6.1]
\EGAReaderTarget{ega4.c04.statement.18-6-1}{IV.18.6.1}
Étant donné un anneau local $A$, nous dirons qu'une $A$-algèbre locale $B$ est \emph{essentiellement étale} s'il existe une $A$-algèbre étale $C$ et un idéal premier $\mathfrak n$ de $C$ tels que $B$ soit $A$-isomorphe à $C_{\mathfrak n}$ et que l'homomorphisme composé $A\to C\to C_{\mathfrak n}$ soit \emph{local} (autrement dit, que $\mathfrak n$ soit au-dessus de l'idéal maximal $\mathfrak m$ de $A$); cela entraîne (17.6.1) que $\mathfrak n$ est le \emph{seul} idéal premier de $B$ au-dessus de $\mathfrak m$ puisque $B$ est un anneau local; si $B$, $B'$ sont deux $A$-algèbres locales essentiellement étales, tout $A$-homomorphisme $B\to B'$ est donc \emph{local}.

Si $A'$ est une $A$-algèbre locale essentiellement étale et $A''$ une $A'$-algèbre locale essentiellement étale, alors $A''$ est une $A$-algèbre locale essentiellement étale. En effet, on a par hypothèse $A'=B_{\mathfrak n}$, où $B$ est une $A$-algèbre étale et $\mathfrak n$ un idéal premier de $B$ au-dessus de l'idéal maximal de $A$, et $A''=B'_{\mathfrak n'}$, où $B'$ est une $A'$-algèbre étale et $\mathfrak n'$ un idéal premier de $B'$ au-dessus de l'idéal $\mathfrak nB_{\mathfrak n}$. Si l'on pose $S=B-\mathfrak n$, de sorte que $A'=S^{-1}B$, $B'$ est de la forme $S^{-1}C$, où $C$ est une $B$-algèbre de présentation finie. Par suite $C$ est une $A$-algèbre de présentation finie et $A''$ est de la forme $C_{\mathfrak r}$, où $\mathfrak r$ est un idéal premier de $C$ au-dessus de l'idéal maximal de $A$. Comme $A'$ est une $A$-algèbre formellement étale et $A''$ une $A'$-algèbre formellement étale, $A''$ est une $A$-algèbre formellement

\oldpage[IV-4]{136}

étale (pour les topologies discrètes (17.1.3)); le morphisme $\Spec(C)\to\Spec(A)$ est donc étale au point $\mathfrak r$ (17.6.1), et il y a par suite un élément $g\in C$ tel que $C_g$ soit une $A$-algèbre étale, ce qui prouve que $A''$ est une $A$-algèbre essentiellement étale.

Étant donné un anneau local $A$, il existe un \emph{ensemble} $\mathfrak E$ de $A$-algèbres locales essentiellement étales tel que toute $A$-algèbre locale essentiellement étale soit $A$-isomorphe à une algèbre appartenant à $\mathfrak E$. Il suffit évidemment pour le voir d'observer qu'il existe un ensemble $\mathfrak F$ de $A$-algèbres de type fini tel que toute $A$-algèbre de type fini soit isomorphe à une algèbre appartenant à $\mathfrak F$; on peut prendre $\mathfrak F$ égal à l'ensemble des quotients des algèbres de polynômes $A[T_1,\ldots,T_n]$ ($n\in\mathbf N$).

Dans ce qui suit, si $A$ et $B$ sont deux anneaux locaux, nous noterons $\Hom.\operatorname{loc}(A,B)$ l'ensemble des homomorphismes \emph{locaux} de $A$ dans $B$.
\end{env}

\begin{lemma}[18.6.2]
\EGAReaderTarget{ega4.c04.lem.18-6-2}{IV.18.6.2}
Soient $A$, $A'$ deux anneaux locaux, $k$, $k'$ leurs corps résiduels respectifs, $\varphi:A\to A'$ un homomorphisme (local) faisant de $A'$ une $A$-algèbre essentiellement étale (18.6.1) et tel que l'homomorphisme correspondant $k\to k'$ soit bijectif. Alors, pour tout anneau local hensélien $B$, l'application canonique
\[
  \Hom(\varphi,1_B):\Hom.\operatorname{loc}(A',B)\longrightarrow\Hom.\operatorname{loc}(A,B)
\]
est bijective.
\end{lemma}

Posons $S=\Spec(A)$, $Y=\Spec(B)$, et soient $s$ et $y$ les points fermés de $S$ et $Y$ respectivement. Par hypothèse, $A'$ est isomorphe à un anneau local $\sh{O}_{X,x}$ d'un $S$-schéma $X$ étale sur $S$, en un point $x\in X$ au-dessus de $s$. Supposons donné un homomorphisme local $\psi:A\to B$, faisant de $Y$ un $S$-schéma; il s'agit de voir qu'il existe \emph{un et un seul} $S$-morphisme $f:Y\to X$ tel que $f(y)=x$. Posons $X'=X\times_S Y$, et notons que puisque $k(x)=k(s)$, il existe un seul point $x'\in X'$ au-dessus de $x$ et de $y$, et que $k(x')=k(y)$. Il faut prouver qu'il existe une seule $Y$-section $f'$ de $X'$ telle que $f'(y)=x'$. Or, le morphisme $g:X'\to Y$ est étale et séparé, et la fibre $X'_0=g^{-1}(y)$ a pour anneau local au point $x'$ le corps $k(x')=k(y)$. Si l'on pose $Y_0=\Spec(k(y))$, il existe donc une unique $Y_0$-section $f'_0$ de $X'_0$ telle que $f'_0(y)=x'$, et la conclusion résulte de ce que $B$ est supposé hensélien et de (18.5.11, $b)$).

Nous dirons qu'une $A$-algèbre locale $A'$ vérifiant les conditions de (18.6.2) est \emph{strictement essentiellement étale}. Notons que le critère (18.5.11, $b)$) signifie que, \emph{pour que $A$ soit hensélien, il faut et il suffit que toute $A$-algèbre strictement essentiellement étale soit $A$-isomorphe à $A$}.

\begin{lemma}[18.6.3]
\label{IV.18.6.3}
Soient $A$ un anneau local, $A_1$, $A_2$ deux $A$-algèbres locales strictement essentiellement étales.
\begin{enumerate}
  \item[(i)] Il existe au plus un $A$-homomorphisme (nécessairement local) de $A_1$ dans $A_2$.
  \item[(ii)] Il existe une $A$-algèbre locale strictement essentiellement étale $A_3$ et deux $A$-homomorphismes $A_1\to A_3$, $A_2\to A_3$.
\end{enumerate}
\end{lemma}

Posons $S=\Spec(A)$; par hypothèse il y a deux $S$-schémas étales sur $S$, $X_1$, $X_2$, et deux points $x_1\in X_1$, $x_2\in X_2$ au-dessus du point fermé $s$ de $S$ et tels que $A_1=\sh{O}_{X_1,x_1}$, $A_2=\sh{O}_{X_2,x_2}$. Posons $X_3=X_1\times_S X_2$; les hypothèses $k(x_1)=k(x_2)=k(s)$ entraînent qu'il existe un seul point $x_3\in X_3$ au-dessus de $x_1$ et $x_2$ et que $k(x_3)=k(s)$ (I, 3.4.9). En outre, $X_3$ est étale sur $S$ (17.3.3), donc $A_3=\sh{O}_{X_3,x_3}$ vérifie les conditions de (ii). Par ailleurs,

\oldpage[IV-4]{137}

on a vu qu'un $A$-homomorphisme de $A_1$ dans $A_2$ est nécessairement local; il correspond à un $S$-morphisme $f$ de $X'_2=\Spec(A_2)$ dans $X_1$ tel que $f(x_2)=x_1$, ou encore, en posant $X'_3=X_1\times_S X'_2$, à une $X'_2$-section $f'$ de $X'_3$ telle que $f'(x_2)=x_3$. Comme $k(x_3)=k(x_2)$ et que $X'_2$ est connexe, l'unicité de $f'$ résulte de (17.4.9), d'où (i).

\begin{env}[18.6.4]
\label{IV.18.6.4}
Désignons par $\mathfrak S$ le sous-ensemble de l'ensemble $\mathfrak E$ défini dans (18.6.1) formé des $A$-algèbres strictement essentiellement étales appartenant à $\mathfrak E$. Il résulte de (18.6.3) que la relation « il existe un $A$-homomorphisme de $A_1$ dans $A_2$ » est une relation de \emph{préordre} dans $\mathfrak S$ et fait de $\mathfrak S$ un ensemble \emph{filtrant croissant}. En indexant $\mathfrak S$ par lui-même au moyen de l'application identique, il résulte de (18.6.3, (i)) que si $\lambda\leq\mu$ dans $\mathfrak S$, il existe un $A$-homomorphisme \emph{unique} $\varphi_{\mu\lambda}:A_\lambda\to A_\mu$ et $(A_\lambda,\varphi_{\mu\lambda})$ est évidemment un \emph{système inductif} de $A$-algèbres locales, les $\varphi_{\mu\lambda}$ étant des homomorphismes locaux. En outre, il résulte de (17.3.5) que pour $\lambda\leq\mu$, $A_\mu$ est une $A_\lambda$-algèbre strictement essentiellement étale.
\end{env}

\begin{definition}[18.6.5]
\EGAReaderTarget{ega4.c04.def.18-6-5}{IV.18.6.5}
On appelle \emph{hensélisé} d'un anneau local $A$ et on désigne par ${}^h\!A$ la $A$-algèbre limite inductive du système inductif $(A_\lambda,\varphi_{\mu\lambda})$ défini dans (18.6.4).
\end{definition}

Cette définition ne dépend qu'en apparence du choix de $\mathfrak E$; si $\mathfrak E'$ est un autre ensemble de $A$-algèbres locales essentiellement étales ayant la même propriété que $\mathfrak E$, $\mathfrak S'$ le sous-ensemble de $\mathfrak E'$ formé des $A$-algèbres strictement essentiellement étales, il résulte de (18.6.3, (ii)) que, pour les relations de préordre envisagées, $\mathfrak S$ et $\mathfrak S'$ sont \emph{cofinaux} à $\mathfrak S''=\mathfrak S\cup\mathfrak S'$, donc donnent la même limite inductive à un isomorphisme près. Nous allons voir d'ailleurs ci-dessous (18.6.6, (i) et (ii)) que ${}^h\!A$ et l'homomorphisme canonique $A\to{}^h\!A$ forment une solution d'un problème universel, et par suite sont déterminés à isomorphisme unique près.

On remarquera que si $A$ est un corps, on a évidemment ${}^h\!A=A$. Dans le cas général, ${}^h\!A$ est aussi le hensélisé de toutes les $A$-algèbres strictement essentiellement étales $A_\lambda$ ($\lambda\in\mathfrak S$), car si $B$ est une $A_\lambda$-algèbre strictement essentiellement étale, $B$ est aussi une $A$-algèbre strictement essentiellement étale (18.6.1), donc (à un $A$-isomorphisme près) une des $A_\mu$ pour $\mu\geq\lambda$, et les $A_\mu$ telles que $\mu\geq\lambda$ et que $A_\mu$ soit une $A_\lambda$-algèbre strictement essentiellement étale forment un ensemble cofinal dans l'ensemble préordonné des $A_\lambda$, en vertu de (18.6.1) et (18.6.3).

\begin{theorem}[18.6.6]
\label{IV.18.6.6}
Soient $A$ un anneau local, ${}^h\!A$ son hensélisé.
\begin{enumerate}
  \item[(i)] ${}^h\!A$ est un anneau local hensélien et l'homomorphisme structural $A\to{}^h\!A$ est local.
  \item[(ii)] Pour tout anneau local hensélien $B$, l'application canonique
  \[
    \Hom.\operatorname{loc}({}^h\!A,B)\longrightarrow\Hom.\operatorname{loc}(A,B)
  \]
  est bijective.
  \item[(iii)] ${}^h\!A$ est un $A$-module fidèlement plat, et si $\mathfrak m$ est l'idéal maximal de $A$, $\mathfrak m.{}^h\!A$ est l'idéal maximal de ${}^h\!A$, et l'homomorphisme $A/\mathfrak mA\to{}^h\!A/\mathfrak m.{}^h\!A$ des corps résiduels est bijective.
  \item[(iv)] Si $\widehat A$ et $({}^h\!A)\widehat{\phantom{x}}$ sont les séparés complétés des anneaux locaux $A$ et ${}^h\!A$, l'homomorphisme $\widehat A\to({}^h\!A)\widehat{\phantom{x}}$ déduit de l'homomorphisme structural $A\to{}^h\!A$ par complétion est bijectif.
  \item[(v)] Pour que ${}^h\!A$ soit noethérien, il faut et il suffit que $A$ le soit.
  \item[(vi)] Si $A$ est hensélien, l'homomorphisme canonique $A\to{}^h\!A$ est bijectif.
\end{enumerate}
\end{theorem}

\oldpage[IV-4]{138}

Soit $\mathfrak m_\lambda$ l'idéal maximal de $A_\lambda$; il résulte du fait que, dans (17.6.1), $a)$ implique $c')$, que pour $\lambda\leq\mu$, on a $\mathfrak m_\mu=\mathfrak m_\lambda A_\mu$, l'homomorphisme $A_\lambda/\mathfrak m_\lambda\to A_\mu/\mathfrak m_\mu$ est bijectif et $A_\mu$ est un $A_\lambda$-module plat. Le fait que ${}^h\!A$ soit local, que l'homomorphisme $A\to{}^h\!A$ soit local, l'assertion (iii) et la suffisance de (v) résultent donc de ($0_{\mathrm{III}}$, 10.3.1.3); la nécessité de (v) résulte de ce que ${}^h\!A$ est un $A$-module fidèlement plat ($0_{\mathrm I}$, 6.5.2).

Pour prouver que ${}^h\!A$ est hensélien, appliquons le critère (18.5.11, $b)$). Posons donc $S=\Spec({}^h\!A)$, $S_0=\Spec({}^h\!A/\mathfrak m.{}^h\!A)$, et soit $g:S'\to S$ un morphisme étale; posons $S'_0=g^{-1}(S_0)$, et soit $f_0:S_0\to S'_0$ une $S_0$-section de $S'_0$. Raisonnant comme dans (18.5.4), on peut supposer que $g$ est de présentation finie. Il résulte alors de (8.8.2) et (17.7.5) qu'il existe un indice $\lambda\in\mathfrak S$, un morphisme étale $g_\lambda:S'^{(\lambda)}\to S^{(\lambda)}=\Spec(A_\lambda)$ et, si l'on pose $S_0^{(\lambda)}=\Spec(A_\lambda/\mathfrak mA_\lambda)$ et $S_0'^{(\lambda)}=g_\lambda^{-1}(S_0^{(\lambda)})$, une $S_0^{(\lambda)}$-section $f_0^{(\lambda)}:S_0^{(\lambda)}\to S_0'^{(\lambda)}$ tels que
\[
  S'=S'^{(\lambda)}\times_{S^{(\lambda)}}S,\qquad
  g=g_\lambda\times1,\qquad
  f_0=f_0^{(\lambda)}\times1.
\]
Soient $s$ le point fermé de $S$, $x=f_0(s)$, $x_\lambda$ la projection de $x$ dans $S'^{(\lambda)}$; comme $x_\lambda$ est au-dessus du point fermé $s_\lambda$ de $S^{(\lambda)}$, l'anneau local $C_\lambda$ de $S'^{(\lambda)}$ au point $x_\lambda$ est une $A_\lambda$-algèbre essentiellement étale; en outre, comme $f_0^{(\lambda)}(s_\lambda)=x_\lambda$, on a $k(x_\lambda)=k(s_\lambda)$, autrement dit, $C_\lambda$ est une $A_\lambda$-algèbre locale strictement essentiellement étale, et par suite (18.6.1) est $A_\lambda$-isomorphe à une $A$-algèbre $A_\mu$ avec $\mu\geq\lambda$. Il y a donc un $A_\lambda$-homomorphisme $C_\lambda\to{}^h\!A$, c'est-à-dire un $S^{(\lambda)}$-morphisme $h:S\to S'^{(\lambda)}$ tel que $h(s)=x_\lambda$, et par suite il existe bien une $S$-section $f$ de $S'$ telle que $f(s)=x$, ce qui achève de prouver (i).

Pour prouver (ii), il suffit de noter que l'on a
\[
  \Hom.\operatorname{loc}({}^h\!A,B)=\varprojlim_\lambda\Hom.\operatorname{loc}(A_\lambda,B)
\]
et que par (18.6.2) les homomorphismes canoniques $\Hom.\operatorname{loc}(A_\lambda,B)\leftarrow\Hom.\operatorname{loc}(A,B)$ sont bijectifs.

Pour démontrer (iv), notons que pour tout $\lambda$ et tout entier $n>0$, $\mathfrak m_\lambda^n=\mathfrak m^nA_\lambda$, et
\[
  (\mathfrak m.{}^h\!A)^n=\mathfrak m^n.{}^h\!A=\varinjlim_\lambda\mathfrak m_\lambda^n.
\]
On en conclut, par l'exactitude du foncteur $\varinjlim$ dans la catégorie des $A$-modules, que
\[
  {}^h\!A/(\mathfrak m.{}^h\!A)^n=\varinjlim_\lambda(A_\lambda/\mathfrak m_\lambda^n),
\]
et par suite, il suffit de montrer que, pour tout entier $n$, et tout indice $\lambda$, l'homomorphisme $A/\mathfrak m^n\to A_\lambda/\mathfrak m_\lambda^n$ est bijectif. Or, cela est vrai par hypothèse pour $n=1$; d'autre part, puisque $A_\lambda$ est un $A$-module plat, on a
\[
  \mathfrak m_\lambda^n/\mathfrak m_\lambda^{n+1}
  =(\mathfrak m^n/\mathfrak m^{n+1})\otimes_A A_\lambda
  =(\mathfrak m^n/\mathfrak m^{n+1})\otimes_{A/\mathfrak m}(A_\lambda/\mathfrak mA_\lambda)
\]
et comme $A/\mathfrak m\to A_\lambda/\mathfrak mA_\lambda=A_\lambda/\mathfrak m_\lambda$ est bijectif, l'homomorphisme $\mathfrak m^n/\mathfrak m^{n+1}\to\mathfrak m_\lambda^n/\mathfrak m_\lambda^{n+1}$ est lui aussi bijectif; la conclusion résulte donc de Bourbaki, \emph{Alg. comm.}, chap. III, \textsection~2, n\textsuperscript{o}~8, cor. 3 du th. 1.

Enfin, si $A$ est hensélien, il résulte de la remarque qui précède (18.6.3) que les homomorphismes $A\to A_\lambda$ sont bijectifs, ce qui prouve (vi) par définition de ${}^h\!A$.

\begin{env}[18.6.7]
\EGAReaderTarget{ega4.c04.statement.18-6-7}{IV.18.6.7}
Soient maintenant $A$ un anneau semi-local, $\mathfrak m_i$ ($1\leq i\leq r$) ses idéaux maximaux. On appelle \emph{hensélisé} de $A$ et l'on note encore ${}^h\!A$ l'anneau produit $\prod_{i=1}^r{}^h(A_{\mathfrak m_i})$ des hensélisés des anneaux locaux $A_{\mathfrak m_i}$. C'est un $A$-module fidèlement plat et une $A$-algèbre semi-locale, dont les idéaux maximaux sont les $\mathfrak m_i.{}^h\!A$ en vertu de (18.6.6, (iii)); en outre, si $\mathfrak r=\bigcap_i\mathfrak m_i$ est le radical de $A$, il résulte de ce qui précède que $\mathfrak r.{}^h\!A$

\oldpage[IV-4]{139}

est le radical de ${}^h\!A$ et que l'application canonique $A/\mathfrak r\to{}^h\!A/\mathfrak r.{}^h\!A$ est bijective. Comme le séparé complété $\widehat A$ de $A$ pour la topologie $\mathfrak r$-préadique est le produit des séparés complétés $\widehat A_{\mathfrak m_i}$ (Bourbaki, \emph{Alg. comm.}, chap. III, \textsection~2, n\textsuperscript{o}~13, prop. 18) l'homomorphisme canonique $\widehat A\to({}^h\!A)\widehat{\phantom{x}}$ est bijectif par (18.6.6, (iv)), et il est clair par (18.6.6, (v)) que pour que ${}^h\!A$ soit noethérien, il faut et il suffit que $A$ le soit.

Pour obtenir l'analogue de la propriété universelle (18.6.6, (ii)), convenons, lorsque $A$ et $B$ sont deux anneaux semi-locaux de radicaux respectifs $\mathfrak r$, $\mathfrak s$, d'appeler \emph{homomorphisme semi-local} de $A$ dans $B$ tout homomorphisme $\varphi$ tel que $\varphi(\mathfrak r)\subset\mathfrak s$. Dans le cas particulier où $B$ est un \emph{produit} d'anneaux locaux $B_j$ ($1\leq j\leq n$), la donnée d'un tel homomorphisme revient à celle de ses projections $\varphi_j:A\to B_j$, qui sont des homomorphismes soumis à la seule condition que $\varphi_j(\mathfrak r)\subset\mathfrak n_j$, où $\mathfrak n_j$ est l'idéal maximal de $B_j$. D'ailleurs, si $\mathfrak m_i$ ($1\leq i\leq m$) sont les idéaux maximaux de $A$, l'idéal premier $\varphi_j^{-1}(\mathfrak n_j)$ doit contenir l'un des $\mathfrak m_i$ puisqu'il contient leur intersection $\mathfrak r$, donc il est égal à l'un des $\mathfrak m_i$ et $\varphi_j$ se factorise en $A\to A_{\mathfrak m_i}\xrightarrow{\psi_{ij}}B_j$, où $\psi_{ij}$ est un homomorphisme \emph{local}. Si $\Hom.\operatorname{sloc}(A,B)$ désigne l'ensemble des homomorphismes semi-locaux de $A$ dans $B$, on peut donc identifier cet ensemble à
\[
  \prod_j\left(\bigcup_i\Hom.\operatorname{loc}(A_{\mathfrak m_i},B_j)\right)
\]
dans le cas considéré; comme un anneau semi-local hensélien est produit direct d'anneaux locaux henséliens, on voit que la propriété universelle (18.6.6, (ii)) est encore valable pour un anneau semi-local $A$, lorsqu'on y remplace les homomorphismes locaux par les homomorphismes semi-locaux.
\end{env}

\begin{proposition}[18.6.8]
\label{IV.18.6.8}
Soient $A$ un anneau semi-local, $B$ une $A$-algèbre semi-locale et entière sur $A$. Alors $B\otimes_A({}^h\!A)$ est une $B$-algèbre semi-locale isomorphe à ${}^h\!B$.
\end{proposition}

Comme $B\otimes_A({}^h\!A)$ est entier au-dessus de ${}^h\!A$, chacun de ses idéaux maximaux est au-dessus d'un des idéaux maximaux de ${}^h\!A$, donc au-dessus d'un idéal maximal de $A$, et comme $B$ est entière sur $A$, les idéaux premiers de $B$ au-dessus d'un idéal maximal de $A$ sont des idéaux maximaux de $B$, si bien que finalement un idéal maximal de $B\otimes_A({}^h\!A)$ a pour projections dans $\Spec({}^h\!A)$ et $\Spec(B)$ des idéaux maximaux; compte tenu de (18.6.6, (iii)) et de (I, 3.4.9), on conclut que l'anneau $B\otimes_A({}^h\!A)$ est semi-local. En outre, pour tout anneau local hensélien $C$, $\Hom.\operatorname{sloc}(B\otimes_A({}^h\!A),C)$ est en correspondance biunivoque avec l'ensemble des couples $(\varphi,\psi)$ d'homomorphismes semi-locaux $\varphi:B\to C$, $\psi:{}^h\!A\to C$ tels que les composés
\[
  A\longrightarrow B\xrightarrow{\varphi}C,\qquad
  A\longrightarrow{}^h\!A\xrightarrow{\psi}C
\]
soient égaux. Mais en vertu de la correspondance biunivoque entre $\Hom.\operatorname{sloc}(A,C)$ et $\Hom.\operatorname{sloc}({}^h\!A,C)$ on voit que pour tout $\varphi\in\Hom.\operatorname{sloc}(B,C)$ il existe un $\psi$ et un seul ayant la propriété précédente, donc l'application
\[
  \Hom.\operatorname{sloc}(B\otimes_A({}^h\!A),C)\longrightarrow\Hom.\operatorname{sloc}(B,C)
\]
est bijective, ce qui prouve la proposition en vertu de l'unicité de la solution d'un problème universel.

\begin{theorem}[18.6.9]
\label{IV.18.6.9}
Soit $A$ un anneau semi-local.
\begin{enumerate}
  \item[(i)] Pour que ${}^h\!A$ soit réduit (resp. normal), il faut et il suffit que $A$ le soit.
  \item[(ii)] Supposons $A$ noethérien. Alors, pour tout idéal premier $\mathfrak p$ de $A$, l'anneau $({}^h\!A)_{\mathfrak p}/\mathfrak p.({}^h\!A)_{\mathfrak p}$ est composé direct d'un nombre fini d'extensions algébriques séparables de $k(\mathfrak p)$ (ce qui entraîne que les fibres du morphisme canonique $\Spec({}^h\!A)\to\Spec(A)$ sont \emph{géométriquement régulières}).
\end{enumerate}
\end{theorem}

\oldpage[IV-4]{140}

(i) Comme ${}^h\!A$ est un $A$-module fidèlement plat, il résulte de (2.1.13) que si ${}^h\!A$ est réduit (resp. normal), il en est de même de $A$. Pour prouver la réciproque, on peut se borner au cas où l'anneau $A$ est local, de sorte qu'avec les notations de (18.6.4), on a ${}^h\!A=\varinjlim A_\alpha$. Or, il résulte de la définition des $A_\alpha$ et de (17.5.7) que si $A$ est réduit (resp. intègre et intégralement clos), il en est de même des $A_\alpha$; en outre, les homomorphismes $A_\alpha\to A_\beta$ pour $\alpha\leq\beta$ sont injectifs en vertu de ($0_{\mathrm I}$, 6.5.1) puisque $A_\beta$ est un $A_\alpha$-module fidèlement plat; donc les morphismes $\Spec(A_\beta)\to\Spec(A_\alpha)$ sont dominants, et on conclut de (5.13.2) (resp. (5.13.4)) que ${}^h\!A$ est réduit (resp. intègre et intégralement clos).

(ii) On peut se borner au cas où $A$ est local. En vertu du fait que le foncteur $\varinjlim$ permute au produit tensoriel, la fibre du morphisme $\Spec({}^h\!A)\to\Spec(A)$ en un point $\mathfrak p$ est limite inductive des fibres des morphismes $\Spec(A_\alpha)\to\Spec(A)$ en ce point. Comme ${}^h\!A$ est noethérien, on voit, compte tenu de (17.6.2), que l'on est ramené à prouver le lemme suivant:

\begin{lemma}[18.6.9.1]
\label{IV.18.6.9.1}
Soient $(B_\alpha)$ un système inductif d'anneaux artiniens, $B=\varinjlim B_\alpha$ sa limite inductive. Si $B$ est noethérien, alors $B$ est artinien; si de plus, pour $\alpha\leq\beta$, les homomorphismes $\varphi_{\beta\alpha}:B_\alpha\to B_\beta$ sont injectifs, alors il existe $\lambda$ tel que, pour $\alpha\geq\lambda$ le nombre des composants locaux $B_\alpha^{(i)}$ de $B_\alpha$ soit constant, que $\varphi_{\beta\alpha}(B_\alpha^{(i)})\subset B_\beta^{(i)}$ pour tout $i$ et que les $B^{(i)}=\varinjlim B_\alpha^{(i)}$ soient les composants locaux de $B$.
\end{lemma}

Soient en effet $Y_\alpha=\Spec(B_\alpha)$, $Y=\Spec(B)$, qui, en tant qu'espace topologique, est égal à $\varprojlim Y_\alpha$ (8.2.10). Les espaces $Y_\alpha$ sont finis et discrets; en outre, si les $\varphi_{\beta\alpha}$ sont injectifs, les morphismes canoniques $Y_\beta\to Y_\alpha$ sont dominants, donc surjectifs. Le lemme résultera donc du lemme purement topologique suivant:

\begin{lemma}[18.6.9.2]
\label{IV.18.6.9.2}
Soit $(Y_\alpha,\psi_{\alpha\beta})$ un système projectif d'espaces topologiques finis et discrets. Si $Y=\varprojlim Y_\alpha$ est noethérien, $Y$ est fini et discret. Si de plus les $\psi_{\alpha\beta}$ sont surjectifs, il existe $\lambda$ tel que pour $\alpha\geq\lambda$, le nombre d'éléments de $Y_\alpha$ soit constant et égal au nombre d'éléments de $Y$.
\end{lemma}

En effet, $Y$ étant compact et noethérien, toute partie de $Y$ est compacte, donc fermée, ce qui implique que $Y$ est discret, donc fini puisqu'il est compact. Si de plus les $\psi_{\alpha\beta}$ sont surjectifs, il en est de même de $\psi_\alpha:Y\to Y_\alpha$, donc $\operatorname{Card}(Y)\geq\operatorname{Card}(Y_\alpha)$, et comme $\operatorname{Card}(Y_\alpha)$ est fonction croissante de $\alpha$, il existe $\lambda$ tel que pour $\alpha\geq\lambda$, $\operatorname{Card}(Y_\alpha)=\operatorname{Card}(Y_\lambda)$; les $\psi_{\alpha\beta}$ sont alors bijectifs pour $\alpha\geq\lambda$ et on a donc $\operatorname{Card}(Y)=\operatorname{Card}(Y_\lambda)$.

\begin{corollary}[18.6.10]
\label{IV.18.6.10}
Soit $A$ un anneau local noethérien. Pour que ${}^h\!A$ possède une des propriétés suivantes:
\begin{enumerate}
  \item[$a)$] être un anneau de Cohen-Macaulay (0, 16.5.3);
  \item[$b)$] vérifier la propriété $(S_n)$ (5.7.3);
  \item[$c)$] être régulier;
  \item[$d)$] vérifier la propriété $(R_n)$ (5.8.2);
\end{enumerate}
il faut et il suffit que $A$ possède cette même propriété.
\end{corollary}

\oldpage[IV-4]{141}

Cela résulte aussitôt de ce que les fibres du morphisme $\Spec({}^h\!A)\to\Spec(A)$ sont géométriquement régulières (18.6.9) et de (6.4.1), (6.5.1) et (6.5.3).

\begin{corollary}[18.6.11]
\label{IV.18.6.11}
Soit $A$ un anneau semi-local noethérien. Pour qu'un idéal premier $\mathfrak p$ de ${}^h\!A$ appartienne à $\operatorname{Ass}({}^h\!A)$, il faut et il suffit que $\mathfrak p\cap A$ appartienne à $\operatorname{Ass}(A)$.
\end{corollary}

Compte tenu de la description des fibres du morphisme $\Spec({}^h\!A)\to\Spec(A)$, cela résulte aussitôt de (3.3.1).

\begin{proposition}[18.6.12]
\label{IV.18.6.12}
Soit $A$ un anneau local; les propriétés suivantes sont équivalentes:
\begin{enumerate}
  \item[$a)$] $A$ est unibranche (0, 23.2.1) (resp. réduit et unibranché).
  \item[$b)$] Pour toute $A$-algèbre locale strictement essentiellement étale $A'$, $\Spec(A')$ est irréductible (resp. intègre).
  \item[$c)$] $\Spec({}^h\!A)$ est irréductible (resp. intègre).
\end{enumerate}
\end{proposition}

Notons d'abord que ${}^h(A_{\mathrm{red}})=({}^h\!A)_{\mathrm{red}}$, car ${}^h(A_{\mathrm{red}})=({}^h\!A)\otimes_A A_{\mathrm{red}}$ (18.6.8), et comme ${}^h(A_{\mathrm{red}})$ est réduit (18.6.9), il est égal à $({}^h\!A)_{\mathrm{red}}$. Ceci permet, dans la démonstration, de considérer seulement le cas où $A$ est réduit, donc aussi $A'$ (17.5.7) et ${}^h\!A$ (18.6.9).

Avec les notations de (18.6.4), la condition $b)$ signifie que tous les $A_\alpha$ sont intègres; on en conclut que ${}^h\!A=\varinjlim A_\alpha$ est intègre (5.13.3), donc $b)$ entraîne $c)$.

Il est clair que si ${}^h\!A$ est intègre, il en est de même des $A_\alpha\subset{}^h\!A$, donc $c)$ entraîne $b)$.

Montrons que $c)$ entraîne $a)$. Notons d'abord que $A\subset{}^h\!A$ est intègre; soient $K$ et $L$ les corps des fractions de $A$ et ${}^h\!A$ respectivement. Il suffit de voir que pour toute $A$-algèbre \emph{finie} $B\subset K$, $B$ est un anneau local. Or $B$ est un anneau semi-local, dont le hensélisé est donc $B\otimes_A({}^h\!A)$ (18.6.8). Par platitude, $B\otimes_A({}^h\!A)$ s'identifie à un sous-anneau de $K\otimes_A({}^h\!A)$ et ce dernier à un sous-anneau de $L$, donc ${}^h\!B$ est intègre. Mais comme ${}^h\!B$ est un anneau semi-local hensélien, il est composé direct d'anneaux locaux (18.5.9), donc ne peut être intègre que s'il est local, ce qui entraîne bien que $B$ est local.

Prouvons enfin que $a)$ implique $c)$. Soit $C$ la clôture intégrale de l'anneau intègre $A$. Comme $C$ est un anneau local par hypothèse, $C\otimes_A({}^h\!A)$ est le hensélisé de $C$ (18.6.8), donc est intègre et intégralement clos (18.6.9). Comme $A$ est un sous-anneau de $C$, ${}^h\!A$ s'identifie à un sous-anneau de ${}^h\!C$ par platitude, donc est lui aussi intègre, ce qui achève la démonstration.

\begin{corollary}[18.6.13]
\label{IV.18.6.13}
Soit $A$ un anneau local hensélien. Alors, pour que $A$ soit unibranche, il faut et il suffit que $A_{\mathrm{red}}$ soit intègre.
\end{corollary}

\begin{proposition}[18.6.14]
\label{IV.18.6.14}
\begin{enumerate}
  \item[(i)] Toute limite inductive filtrante d'anneaux locaux henséliens (où les homomorphismes de transition sont locaux) est un anneau local hensélien.
  \item[(ii)] Le foncteur $A\leadsto{}^h\!A$ dans la catégorie des anneaux locaux (où les morphismes sont les homomorphismes locaux) commute aux limites inductives filtrantes.
  \item[(iii)] Tout anneau local hensélien $A$ est limite inductive d'un système inductif filtrant d'anneaux locaux henséliens et noethériens $A_\lambda$ (les homomorphismes de transition $A_\lambda\to A_\mu$ étant locaux).
\end{enumerate}
\end{proposition}

(i) Soit $(A_\lambda)$ un système inductif filtrant d'anneaux locaux henséliens, où les homomorphismes de transition sont locaux. Montrons que $A'=\varinjlim A_\lambda$, qui est local ($0_{\mathrm{III}}$, 10.3.1.3) est hensélien. Or, soient $C$ une $A'$-algèbre étale, $\mathfrak n$ un idéal premier

\oldpage[IV-4]{142}

de $C$ au-dessus de l'idéal maximal $\mathfrak m'$ de $A'$, tel que $C_{\mathfrak n}$ soit strictement essentiellement étale sur $A'$ (18.6.1), autrement dit que $k(\mathfrak m')\to k(\mathfrak n)$ soit un isomorphisme. En vertu de (8.8.2, (ii)) et de (17.7.8), il existe un indice $\lambda$ et une $A_\lambda$-algèbre étale $C_\lambda$ tels que $C=C_\lambda\otimes_{A_\lambda}A'$ à isomorphisme près; en outre, si $\mathfrak n_\lambda$ est l'image réciproque de $\mathfrak n$ dans $C_\lambda$, et $\mathfrak m_\lambda$ l'idéal maximal de $A_\lambda$, on peut, en vertu de (8.8.2.4) et de la transitivité des fibres (I, 3.6.4), supposer que $(C_\lambda)_{\mathfrak n_\lambda}$ est strictement essentiellement étale sur $A_\lambda$; puisque $A_\lambda$ est hensélien, $(C_\lambda)_{\mathfrak n_\lambda}$ est nécessairement isomorphe à $A_\lambda$ (18.5.11, $b)$); il existe par suite un voisinage de $\mathfrak n_\lambda$ dans $\Spec(C_\lambda)$ isomorphe à $\Spec(A_\lambda)$ (I, 6.5.4); on en conclut qu'il y a un voisinage de $\mathfrak n$ dans $\Spec(C)$ isomorphe à $\Spec(A')$, donc que $C_{\mathfrak n}$ est isomorphe à $A'$, ce qui prouve que $A'$ est hensélien (18.6.2) et achève la démonstration.

(ii) Supposons que $B=\varinjlim B_\lambda$, où les $B_\lambda$ sont des anneaux locaux, les homomorphismes de transition $B_\lambda\to B_\mu$ ($\lambda\leq\mu$) étant locaux. Posons $A_\lambda={}^h\!B_\lambda$; les $A_\lambda$ sont des anneaux locaux henséliens (18.6.6, (vi)), et en vertu de (18.6.6, (ii)) les homomorphismes de transition $B_\lambda\to B_\mu$ déterminent de façon unique des homomorphismes locaux $A_\lambda\to A_\mu$, de sorte que $(A_\lambda)$ est un système inductif. Tout revient à voir que $A'=\varinjlim A_\lambda$ est canoniquement isomorphe à $A={}^h\!B$. En vertu de (18.6.6, (ii)) et de la définition des limites inductives, on a, pour tout anneau local hensélien $E$,
\[
  \Hom.\operatorname{loc}(A',E)
  =\varprojlim\Hom.\operatorname{loc}(A_\lambda,E)
  =\varprojlim\Hom.\operatorname{loc}(B_\lambda,E)
  =\Hom.\operatorname{loc}(B,E).
\]
Mais comme $A'$ est hensélien par (i), cela prouve notre assertion.

(iii) L'anneau $A$ est limite inductive filtrante de ses sous-anneaux locaux noethériens $B_\lambda$, les homomorphismes de transition étant locaux (5.13.3, (iii)). Comme les ${}^h\!B_\lambda$ sont des anneaux locaux henséliens et noethériens (18.6.6, (v)) et que $A={}^h\!A$ (18.6.6, (vi)), il suffit d'appliquer (ii) au système inductif $(B_\lambda)$.

\begin{corollary}[18.6.15]
\label{IV.18.6.15}
Soient $A$ un anneau local hensélien, $X$ un $A$-préschéma de présentation finie sur $A$. Il existe alors un anneau local noethérien et hensélien $A_0$, un homomorphisme local $A_0\to A$, un préschéma $X_0$ de type fini sur $A_0$ et un $A$-isomorphisme $X_0\otimes_{A_0}A\xrightarrow{\sim}X$.
\end{corollary}

Cela résulte de (18.6.14) et de (8.8.2, (ii)).

\subsection{Hensélisation et anneaux excellents.}
\EGAReaderTarget{ega4.c04.subsec.18-7}{subsection:IV.18.7}

\begin{env}[18.7.1]
\label{IV.18.7.1}
Nous désignerons dans ce numéro par $\mathbf P(Z,k)$ une propriété de la forme considérée dans (7.3.1), où nous supposons en outre que la propriété $\mathbf Q(A,k)$ satisfait à la condition suivante:

Pour toute extension algébrique séparable $k'$ de $k$, et toute $k'$-algèbre locale noethérienne $A$, la propriété $\mathbf Q(A,k)$ est \emph{équivalente} à $\mathbf Q(A,k')$.

Il est immédiat que toutes les propriétés $\mathbf P$ considérées dans (7.3.8) vérifient la condition précédente, compte tenu du fait que si $K$ est une extension finie de $k$, $k'\otimes_k K$ est composé direct d'extensions algébriques séparables de $K$.
\end{env}

\oldpage[IV-4]{143}

\begin{proposition}[18.7.2]
\label{IV.18.7.2}
Soit $A$ un anneau local noethérien. Pour que ${}^h\!A$ soit un $\mathbf P$-anneau (7.3.13), il faut et il suffit que $A$ le soit.
\end{proposition}

En effet, en vertu de (18.6.6, (iv)), le complété $\widehat A$ de $A$ est aussi le complété de ${}^h\!A$ et l'on a donc $A\subset{}^h\!A\subset\widehat A$. D'après (18.6.9), pour tout $x\in\Spec(A)$, la fibre en $x$ du morphisme $\Spec({}^h\!A)\to\Spec(A)$ est discrète et finie, et en chacun des points $y_i$ de cette fibre, $k(y_i)$ est une extension algébrique séparable de $k(x)$. La fibre formelle de $A$ au point $x$ est donc un préschéma \emph{somme} des fibres formelles de ${}^h\!A$ aux points $y_i$. La conclusion résulte alors de l'hypothèse faite sur $\mathbf P$ dans (18.7.1).

\begin{corollary}[18.7.3]
\label{IV.18.7.3}
Soit $A$ un anneau local noethérien. Pour que ${}^h\!A$ soit universellement japonais, il faut et il suffit que $A$ le soit.
\end{corollary}

En effet, pour un anneau local noethérien $B$, il revient au même de dire que $B$ est universellement japonais, ou que les fibres formelles de $B$ sont géométriquement réduites (7.6.4 et 7.7.1); la conclusion résulte donc de (18.7.2).

\begin{corollary}[18.7.4]
\label{IV.18.7.4}
Soit $A$ un anneau local noethérien. Pour que les fibres formelles de ${}^h\!A$ soient géométriquement régulières, il faut et il suffit que celles de $A$ le soient.
\end{corollary}

C'est un cas particulier de (18.7.2).

\begin{proposition}[18.7.5]
\label{IV.18.7.5}
Soit $A$ un anneau local noethérien. Si $A$ est universellement caténaire (5.6.2) (resp. formellement caténaire (7.1.9), resp. strictement formellement caténaire (7.2.6)), il en est de même de ${}^h\!A$.
\end{proposition}

Avec les notations de (18.6.4), si $A$ est universellement caténaire, il en est de même des $A_\alpha$, qui sont des $A$-algèbres essentiellement de type fini (5.6.3). Pour prouver que dans ce cas ${}^h\!A$ est universellement caténaire, il suffira donc d'établir le lemme suivant:

\begin{lemma}[18.7.5.1]
\label{IV.18.7.5.1}
Soit $(B_\alpha,\varphi_{\beta\alpha})$ un système inductif d'anneaux noethériens et supposons que, pour $\alpha\leq\beta$, $\varphi_{\beta\alpha}$ fasse de $B_\beta$ un $B_\alpha$-module fidèlement plat et que $B=\varinjlim B_\alpha$ soit noethérien. Alors, si les $B_\alpha$ sont caténaires (resp. universellement caténaires), il en est de même de $B$.
\end{lemma}

Dire que $B$ est universellement caténaire équivaut à dire que $B[T_1,\ldots,T_n]$ est caténaire pour tout $n$ (5.6.2), et il est immédiat que le système inductif $(B_\alpha[T_1,\ldots,T_n])$ vérifie les mêmes conditions que $(B_\alpha)$; il suffit donc de prouver que si les $B_\alpha$ sont caténaires, il en est de même de $B$. Si $\mathfrak p_0\supset\mathfrak p_1\supset\cdots\supset\mathfrak p_r$ est une chaîne saturée d'idéaux premiers de $B$, les images réciproques $\mathfrak p_{0\alpha}\supset\mathfrak p_{1\alpha}\supset\cdots\supset\mathfrak p_{r\alpha}$ de ces idéaux dans $B_\alpha$ forment une chaîne d'idéaux premiers, et il suffit de montrer que pour $\alpha$ assez grand, cette chaîne est saturée. Or, chaque $\mathfrak p_i$ est un idéal de type fini, donc il existe un $\lambda$ tel que, pour $\alpha\geq\lambda$, on ait $\mathfrak p_i=\mathfrak p_{i\alpha}.B$. On a donc
\[
1=\operatorname{codim}(V(\mathfrak p_i),V(\mathfrak p_{i+1}))
=\operatorname{codim}(V(\mathfrak p_{i\alpha}),V(\mathfrak p_{i+1,\alpha}))
\]
puisque $B$ est un $B_\alpha$-module plat ((6.1.4) et ($0_{\mathrm I}$, 6.2.3)), ce qui achève de prouver le lemme.

Supposons maintenant $A$ formellement caténaire. Soient $\mathfrak p$ un idéal premier de ${}^h\!A$, $\mathfrak q$ son image réciproque dans $A$, de sorte que $\mathfrak q.{}^h\!A\subset\mathfrak p$, et par suite ${}^h\!A/\mathfrak p$ est un anneau quotient de ${}^h\!A/\mathfrak q.{}^h\!A$. Il suffit donc de prouver que ${}^h\!A/\mathfrak q.{}^h\!A$ est formellement caténaire pour tout idéal premier $\mathfrak q$ de $A$ (7.1.9), et comme ${}^h\!A/\mathfrak q.{}^h\!A={}^h(A/\mathfrak q)$ par (18.6.8), il suffit (7.1.11) de voir que si $A$ est \emph{intègre et formellement caténaire}, ${}^h\!A$ est \emph{formellement équidimensionnel}. Mais comme le complété de ${}^h\!A$ est égal à $\widehat A$, cela résulte de l'hypothèse que $A$ est formellement caténaire.

\oldpage[IV-4]{144}

Supposons enfin $A$ strictement formellement caténaire; alors on vient de voir que ${}^h\!A$ est formellement caténaire, et il reste à prouver que les fibres du morphisme $\Spec(\widehat A)\to\Spec({}^h\!A)$ vérifient la propriété $(S_1)$ (7.2.5, $b)$). Mais cela résulte de l'hypothèse sur $A$ et de (18.7.2) appliqué au cas où $\mathbf P(Z,k)$ est la propriété $(S_1)$ pour $Z$.

\begin{corollary}[18.7.6]
\label{IV.18.7.6}
Si un anneau local noethérien $A$ est excellent (7.8.2), il en est de même de ${}^h\!A$.
\end{corollary}

\begin{remark}[18.7.7]
\label{IV.18.7.7}
Il peut se faire que ${}^h\!A$ soit un anneau local excellent sans que $A$ soit universellement caténaire (ni par suite excellent). Pour le voir, reprenons l'exemple de l'anneau $A$ de (5.6.11), avec les mêmes notations. L'anneau $E$ construit dans (5.6.11) est excellent lorsque $k_0$ est de caractéristique $0$ (7.8.3, (ii) et (iii)); considérons deux anneaux $E_1$, $E_2$ isomorphes à $E$, et « recollons » $\Spec(E_1)$ et $\Spec(E_2)$ de sorte que si $\mathfrak m_i$ et $\mathfrak m'_i$ sont les idéaux maximaux de $E_i$ ($i=1,2$) correspondant à $\mathfrak m$ et $\mathfrak m'$, le point $\mathfrak m_1$ soit « recollé » à $\mathfrak m'_2$ et le point $\mathfrak m_2$ à $\mathfrak m'_1$; de façon précise, si $\varepsilon_i$ et $\varepsilon'_i$ sont les homomorphismes canoniques de $E_i$ sur $k(\mathfrak m_i)$ et $k(\mathfrak m'_i)$ ($i=1,2$), le schéma obtenu est $\Spec(R)$, où $R$ est le sous-anneau de $E_1\times E_2$ formé des couples $(x_1,x_2)$ tels que $\varepsilon'_2(x_2)=\sigma(\varepsilon_1(x_1))$ et $\varepsilon'_1(x_1)=\sigma(\varepsilon_2(x_2))$. On vérifie aisément (par exemple à l'aide de (17.6.3)) que $R$ est une $C$-algèbre finie et étale, et qu'il y a deux idéaux maximaux $\mathfrak r_1$, $\mathfrak r_2$ de $R$ au-dessus de l'idéal maximal $\mathfrak n$ de $C$, les complétés de $R_{\mathfrak r_1}$ et $R_{\mathfrak r_2}$ étant canoniquement isomorphes à celui de $A=C_{\mathfrak n}$; $R_{\mathfrak r_1}$ est donc une $A$-algèbre strictement essentiellement étale, et par suite son hensélisé est égal à ${}^h\!A$. En outre, on a vu (7.8.4, (ii)) que les fibres formelles de $A$ sont géométriquement régulières, donc il en est de même de celles de ${}^h\!A$ (18.7.2). Enfin, l'anneau $R$ est universellement caténaire, car ses quotients par ses deux idéaux premiers minimaux sont isomorphes à $E$, d'où la conclusion par (5.6.3, (iii)). L'anneau $R_{\mathfrak r_1}$ est donc excellent, et il en est par suite de même de ${}^h\!A$ (18.7.6), alors que $A$ n'est pas universellement caténaire.
\end{remark}

\subsection{Anneaux strictement locaux et hensélisation stricte.}
\EGAReaderTarget{ega4.c04.subsec.18-8}{subsection:IV.18.8}

\begin{proposition}[18.8.1]
\label{IV.18.8.1}
Soit $A$ un anneau local. Les conditions suivantes sont équivalentes:
\begin{enumerate}
  \item[$a)$] $A$ est un anneau hensélien et son corps résiduel $k$ est séparablement clos.
  \item[$b)$] $A$ est un anneau hensélien et tout revêtement étale de $S=\Spec(A)$ est trivial.
  \item[$c)$] Pour tout morphisme étale $f:X\to S$ et tout point $x\in X$ tel que $f(x)$ soit le point fermé $s$ de $S$, il existe une $S$-section $u:S\to X$ de $X$ telle que $u(s)=x$.
\end{enumerate}
\end{proposition}

L'hypothèse $c)$ implique que pour tout morphisme étale $f:X\to S$, les corps résiduels aux points de $f^{-1}(s)$ sont tous isomorphes à $k$, et en outre que la condition $b)$ de (18.5.11) est vérifiée. Donc $A$ est hensélien, et le fait que tout revêtement étale de $S$ est trivial résulte de (18.2.10, (ii)); donc $c)$ implique $b)$. Comme les revêtements étales de $\Spec(k)$ sont triviaux si et seulement si $k$ est séparablement clos, $b)$ implique $a)$ en vertu de (18.5.11, $b)$). Enfin, il résulte aussi de (18.5.11, $b)$) que $a)$ implique $c)$.

\begin{definition}[18.8.2]
\EGAReaderTarget{ega4.c04.def.18-8-2}{IV.18.8.2}
On dit qu'un anneau local est un anneau \emph{strictement local} s'il vérifie les conditions équivalentes de (18.8.1). On appelle \emph{schéma strictement local} un schéma isomorphe au spectre d'un anneau strictement local.
\end{definition}

\begin{remark}[18.8.3]
\label{IV.18.8.3}
Les conditions de (18.8.1) sont encore équivalentes à la suivante:
\begin{enumerate}
  \item[$d)$] Pour tout morphisme localement de type fini $f:X\to S$ et tout point $x\in X$ tel que $f(x)$ soit le point fermé $s$ de $S$ et que $\sh{O}_{X,x}/\mathfrak m_s\sh{O}_{X,x}$ soit un corps, extension finie séparable de $k(s)=k$, il existe un voisinage ouvert $U$ de $x$ dans $X$ tel que $f|U$ soit une immersion fermée.
\end{enumerate}
\end{remark}

\oldpage[IV-4]{145}

(On notera que l'hypothèse de $d)$ est vérifiée si $f$ est \emph{formellement non ramifié} au point $x$ (17.4.1.2).)

Nous laissons au lecteur la démonstration, qui est substantiellement la même que celle de (18.5.18).

\begin{lemma}[18.8.4]
\label{IV.18.8.4}
Soient $A$, $A'$ deux anneaux locaux, $\varphi:A\to A'$ un homomorphisme local faisant de $A'$ une $A$-algèbre essentiellement étale (18.6.1). Alors, pour tout anneau strictement local $B$, tout homomorphisme local $\psi:A\to B$ et tout homomorphisme de $k(A)$-algèbres $\alpha:k(A')\to k(B)$, il existe un et un seul homomorphisme local $\psi':A'\to B$ tel que $\psi=\psi'\circ\varphi$ et que $\alpha$ soit égal à l'homomorphisme $\overline{\psi'}$ déduit de $\psi'$ par passage aux quotients.
\end{lemma}

Avec les notations de (18.6.2), l'homomorphisme $\alpha$ correspond à un $k(s)$-morphisme $\Spec(k(y))\to\Spec(k(x))$, donc à un point $x'$ bien déterminé dans $X'$ au-dessus de $y$. L'existence de la $Y$-section $f'$ de $X'$ telle que $f'(y)=x'$ résulte donc de (18.8.1, $c)$) puisque $f'$ est étale et $B$ strictement local, et son unicité résulte de ce que $f'$ est séparé et de (17.4.9).

\begin{lemma}[18.8.5]
\label{IV.18.8.5}
Soient $A$ un anneau local, $A_1$, $A_2$ deux $A$-algèbres locales essentiellement étales, $K$ un corps, extension de $k(A)$.
\begin{enumerate}
  \item[(i)] Pour tout $k(A)$-homomorphisme $\gamma:k(A_1)\to k(A_2)$, il existe au plus un $A$-homomorphisme (nécessairement local) $\psi:A_1\to A_2$ tel que $\gamma$ soit l'homomorphisme $\overline\psi$ déduit de $\psi$ par passage aux quotients.
  \item[(ii)] Pour tout couple de $A$-homomorphismes locaux $\beta_1:A_1\to K$, $\beta_2:A_2\to K$ il existe une $A$-algèbre locale essentiellement étale $A_3$, deux $A$-homomorphismes $\varphi_1:A_1\to A_3$, $\varphi_2:A_2\to A_3$ et un $A$-homomorphisme $\beta:A_3\to K$ tels que $\beta_1=\beta\circ\varphi_1$ et $\beta_2=\beta\circ\varphi_2$.
\end{enumerate}
\end{lemma}

Avec les notations de (18.6.3), les homomorphismes $\beta_1$, $\beta_2$ dans (ii) correspondent à deux $S$-morphismes $\Spec(K)\to X_1$, $\Spec(K)\to X_2$, d'images respectives $x_1$, $x_2$; on en déduit un $S$-morphisme bien déterminé $\Spec(K)\to X_3$ (I, 3.2.1) d'image $x_3$ au-dessus de $x_1$ et $x_2$; la $A$-algèbre $A_3=\sh{O}_{X_3,x_3}$ et l'homomorphisme $A_3\to K$ correspondant à ces données répondent aux conditions de (ii). D'autre part, la donnée de $\gamma$ dans (i) correspond à un $S$-morphisme $\Spec(k(x_2))\to\Spec(k(x_1))$ ou encore à un point $x'_3$ de $X'_3$ au-dessus de $x_2$; l'unicité de $\psi$ résulte de l'unicité d'une $X_2$-section de $X'_3$ passant par $x'_3$ (17.4.9).

\begin{env}[18.8.6]
\label{IV.18.8.6}
Soient $A$ un anneau local, $\mathfrak m$ son idéal maximal, $k$ son corps résiduel, $\Omega$ une clôture séparable de $k$. Considérons l'ensemble $\mathfrak E$ de $A$-algèbres essentiellement étales défini dans (18.6.1), et désignons par $L$ l'ensemble des $A$-homomorphismes locaux $A'\to\Omega$, où $A'$ parcourt $\mathfrak E$; on notera qu'un tel homomorphisme se factorise en
\[
  \lambda:A'\longrightarrow k(A')\xrightarrow{\omega_\lambda}\Omega,
\]
où $\omega_\lambda:k(A')\to\Omega$ est un $k$-homomorphisme (puisque $\mathfrak mA'$ est l'idéal maximal de $A'$); inversement la donnée d'un tel homomorphisme $\omega_\lambda$ détermine uniquement un homomorphisme $\lambda\in L$. Pour tout $\lambda\in L$, nous désignerons par $A_\lambda$ la $A$-algèbre essentiellement étale appartenant à $\mathfrak E$ sur laquelle $\lambda$ est défini. L'ensemble $L$ est \emph{préordonné} par la relation « il existe un $A$-homomorphisme $\varphi_{\mu\lambda}:A_\lambda\to A_\mu$ tel que $\lambda=\mu\circ\varphi_{\mu\lambda}$ », et il résulte de (18.8.5, (i)) que cet homomorphisme $\varphi_{\mu\lambda}$ est \emph{unique}. En outre, $L$ est \emph{filtrant croissant} pour la relation de préordre précédente (que l'on note encore $\lambda\leq\mu$), en vertu

\oldpage[IV-4]{146}

de (18.8.5, (ii)). Il est clair que $(A_\lambda,\varphi_{\mu\lambda})$ est un \emph{système inductif} de $A$-algèbres locales, les $\varphi_{\mu\lambda}$ étant des homomorphismes locaux; en outre, il résulte de (17.3.5) que pour $\lambda\leq\mu$, $A_\mu$ est une $A_\lambda$-algèbre essentiellement étale. Si $\overline\varphi_{\mu\lambda}:k(A_\lambda)\to k(A_\mu)$ est le $k$-homomorphisme déduit de $\varphi_{\mu\lambda}$ par passage aux quotients, $(k(A_\lambda),\overline\varphi_{\mu\lambda})$ est un \emph{système inductif} d'extensions algébriques séparables de $k$, et les $\omega_\lambda:k(A_\lambda)\to\Omega$ forment un système inductif de $k$-homomorphismes. En outre, la limite inductive
\begin{equation}
  \omega:\varinjlim k(A_\lambda)\longrightarrow\Omega
  \label{IV.18.8.6.1}
  \tag{18.8.6.1}
\end{equation}
est un $k$-\emph{isomorphisme}. Il suffit en effet de prouver que pour toute extension finie séparable $k'$ de $k$, il existe une $A$-algèbre essentiellement étale $A'$ telle que $k'$ soit le corps résiduel de $A'$ et que l'homomorphisme $k\to k'$ se déduise de $A\to A'$ par passage aux quotients. Mais cela résulte de (18.1.1) appliqué aux morphismes étales, compte tenu de ce que $\Spec(A)$ est le seul voisinage de son point fermé.

Notons maintenant que si $\mathfrak m_\lambda$ est l'idéal maximal de $A_\lambda$, il résulte de (17.6.1) que pour $\lambda\leq\mu$, on a $\mathfrak m_\mu=\mathfrak m_\lambda A_\mu$ et que $A_\mu$ est un $A_\lambda$-module plat. Donc il résulte de ($0_{\mathrm{III}}$, 10.3.1.3) que l'anneau $\varinjlim A_\lambda$ est \emph{local}, que l'homomorphisme canonique $w:A\to\varinjlim A_\lambda$ est local, et que l'on a un $k$-isomorphisme canonique
\[
  \varinjlim k(A_\lambda)\xrightarrow{\sim}k(\varinjlim A_\lambda);
\]
identifiant ces deux corps, on en déduit un $k$-\emph{isomorphisme canonique}
\begin{equation}
  \omega^{-1}:\Omega\xrightarrow{\sim}k(\varinjlim A_\lambda).
  \label{IV.18.8.6.2}
  \tag{18.8.6.2}
\end{equation}
\end{env}

\begin{definition}[18.8.7]
\EGAReaderTarget{ega4.c04.def.18-8-7}{IV.18.8.7}
Soient $A$ un anneau local, $k=k(A)$ son corps résiduel, $i:k\to\Omega$ un homomorphisme de $k$ dans une clôture séparable $\Omega$ de $k$; on appelle \emph{hensélisé strict} de $A$ relativement à $i$, et on note ${}^{hs}\!A_{(i)}$ (ou ${}^{hs}\!A$ lorsque cela ne prête pas à confusion) la limite inductive du système inductif $(A_\lambda,\varphi_{\mu\lambda})$ défini dans (18.8.6).
\end{definition}

Comme on a un $k$-isomorphisme canonique $\omega^{-1}:\Omega\xrightarrow{\sim}k({}^{hs}\!A_{(i)})$, $\Omega$ est muni canoniquement d'une structure de ${}^{hs}\!A_{(i)}$-algèbre par l'homomorphisme local
\[
  {}^{hs}\!A_{(i)}\longrightarrow k({}^{hs}\!A_{(i)})\xrightarrow{\omega}\Omega.
\]

Comme dans (18.6.5) on voit que la définition donnée dans (18.8.7) ne dépend qu'en apparence du choix de l'ensemble $\mathfrak E$; nous allons voir encore ci-dessous (18.8.8, (i) et (ii)) que ${}^{hs}\!A_{(i)}$ est un objet représentant un foncteur entièrement défini par la donnée de $A$ et de $i$ ($0_{\mathrm{III}}$, 8.1.8), et est donc défini comme tel à isomorphisme unique près.

\begin{proposition}[18.8.8]
\label{IV.18.8.8}
Soient $A$ un anneau local, $k$ son corps résiduel, $i:k\to\Omega$ un homomorphisme de $k$ dans une clôture séparable de $k$, ${}^{hs}\!A_{(i)}$ le hensélisé strict de $A$ correspondant à $i$.
\begin{enumerate}
  \item[(i)] ${}^{hs}\!A_{(i)}$ est un anneau strictement local (18.8.2) et l'homomorphisme structural $w:A\to{}^{hs}\!A_{(i)}$ est local.
  \item[(ii)] Pour tout homomorphisme local $u:A\to B$, où $B$ est un anneau strictement local, et tout $k$-homomorphisme $\psi:\Omega\to k(B)$, il existe un et un seul $A$-homomorphisme $v:{}^{hs}\!A_{(i)}\to B$ tel que $\psi$ se factorise en $\Omega\xrightarrow{\omega^{-1}}k({}^{hs}\!A_{(i)})\xrightarrow{\overline v}k(B)$, où $\overline v$ est déduit de $v$ par passage aux quotients.

\oldpage[IV-4]{147}

  \item[(iii)] ${}^{hs}\!A_{(i)}$ est un $A$-module fidèlement plat, et si $\mathfrak m$ est l'idéal maximal de $A$, $\mathfrak m.{}^{hs}\!A_{(i)}$ est l'idéal maximal de ${}^{hs}\!A_{(i)}$, et $k({}^{hs}\!A_{(i)})$ est une clôture séparable de $k$.
  \item[(iv)] Pour que ${}^{hs}\!A_{(i)}$ soit noethérien, il faut et il suffit que $A$ le soit.
  \item[(v)] Si $A$ est strictement local (donc $\Omega=k$), on a ${}^{hs}\!A\xrightarrow{\sim}A$.
  \item[(vi)] Si $i':k\to\Omega'$ est un second homomorphisme de $k$ dans une clôture séparable de $k$, pour tout $k$-isomorphisme $\sigma:\Omega\xrightarrow{\sim}\Omega'$, le $A$-homomorphisme correspondant (par (ii)) $v_\sigma:{}^{hs}\!A_{(i)}\to{}^{hs}\!A_{(i')}$ est un isomorphisme.
\end{enumerate}
\end{proposition}

On a déjà vu ci-dessus que ${}^{hs}\!A_{(i)}$ est un anneau local et que $w$ est un homomorphisme local; le fait que ${}^{hs}\!A_{(i)}$ soit hensélien (et par suite strictement local) se démontre comme dans (18.6.6). Les assertions de (iii) découlent encore de ($0_{\mathrm{III}}$, 10.3.1.3), ainsi que la suffisance de la condition (iv); la nécessité de cette condition résulte de ce que ${}^{hs}\!A_{(i)}$ est un $A$-module fidèlement plat ($0_{\mathrm I}$, 6.5.2).

Pour prouver (ii), remarquons, avec les notations de (18.8.6), que pour tout $\lambda\in L$, il existe un et un seul $A$-homomorphisme local $v_\lambda:A_\lambda\to B$, tel que le composé $k(A_\lambda)\xrightarrow{\omega_\lambda}\Omega\xrightarrow{\psi}k(B)$ soit déduit de $v_\lambda$ par passage aux quotients, en vertu de (18.8.4) et de l'hypothèse que $B$ est strictement local. L'unicité de $v_\lambda$ entraîne en outre que les $v_\lambda$ forment un système inductif, d'où l'existence de l'homomorphisme $v$ répondant aux conditions de (ii); son unicité résulte de (17.4.9). L'unicité de $v$ entraîne aussitôt l'assertion (vi), en considérant les composés $v_\sigma\circ v_{\sigma^{-1}}$ et $v_{\sigma^{-1}}\circ v_\sigma$. Enfin, si $A$ est strictement local, et $A'=B_{\mathfrak n}$ une $A$-algèbre locale essentiellement étale, où $B$ est une $A$-algèbre étale, il résulte de (18.8.1) qu'il existe une $\Spec(A)$-section de $\Spec(B)$ prenant la valeur $\mathfrak n$ au point fermé de $\Spec(A)$, donc on déduit de (17.4.1) que l'homomorphisme $A\to A'$ est bijectif, ce qui prouve (v).

\begin{env}[18.8.8.1]
\label{IV.18.8.8.1}
Soit $\mathcal C$ la catégorie des anneaux strictement locaux, avec pour morphismes les homomorphismes locaux; pour tout $B\in\mathcal C$, désignons par $F(B)$ l'ensemble des couples $(u,\psi)$ formés d'un homomorphisme local $u:A\to B$ et d'un $k$-isomorphisme $\psi:\Omega\to k(B)$ tels que le composé $k(A)=k\xrightarrow{i}\Omega\xrightarrow{\psi}k(B)$ soit déduit de $u$ par passage aux quotients. On peut dire que l'objet ${}^{hs}\!A_{(i)}$ \emph{représente} le foncteur covariant $F:\mathcal C\to\mathbf{Ens}$ ($0_{\mathrm{III}}$, 8.1.11).

On remarquera qu'en vertu de (18.8.8, (vi)) les hensélisés stricts ${}^{hs}\!A_{(i)}$ sont tous $A$-isomorphes (pour les divers $k$-homomorphismes $i$ de $k$ dans des clôtures séparables de $k$), mais pour $i$ et $i'$ donnés, il y a en général une infinité de $A$-isomorphismes ${}^{hs}\!A_{(i)}\xrightarrow{\sim}{}^{hs}\!A_{(i')}$. De façon précise, le groupe des $A$-automorphismes de ${}^{hs}\!A_{(i)}$ est isomorphe au groupe de Galois de $\Omega$ sur $k$.
\end{env}

\begin{env}[18.8.9]
\EGAReaderTarget{ega4.c04.statement.18-8-9}{IV.18.8.9}
Soient maintenant $A$ un anneau \emph{semi-local}, $\mathfrak m_j$ ($1\leq j\leq r$) ses idéaux maximaux, et pour tout indice $j$, considérons un homomorphisme $i_j:k(A_{\mathfrak m_j})\to\Omega_j$ de $k(A_{\mathfrak m_j})$ dans une clôture séparable de ce corps. On appelle \emph{hensélisé strict} de $A$ relatif aux $i_j$ et l'on note ${}^{hs}\!A$ (lorsqu'il n'en résulte pas de confusion) le produit $\prod_{j=1}^r({}^{hs}(A_{\mathfrak m_j}))$ des hensélisés stricts des anneaux locaux $A_{\mathfrak m_j}$ relativement aux homomorphismes $i_j$.

\oldpage[IV-4]{148}

C'est donc, en vertu de (18.8.8), un $A$-module fidèlement plat et une $A$-algèbre semi-locale dont les idéaux maximaux sont les $\mathfrak m_j.{}^{hs}\!A$ et le radical $\mathfrak r.{}^{hs}\!A$ (si $\mathfrak r$ désigne le radical de $A$). Enfin, la propriété universelle (18.8.8, (ii)) subsiste en remplaçant « homomorphisme local » par « homomorphisme semi-local » (18.6.7), et en remplaçant $\Omega$ par $\prod_j\Omega_j$.
\end{env}

\begin{proposition}[18.8.10]
\label{IV.18.8.10}
Soient $A$ un anneau semi-local, $B$ une $A$-algèbre finie. Soient $\mathfrak n_j$ ($1\leq j\leq r$) les idéaux maximaux de $B$. Alors il y a, pour chaque $j$, un idéal maximal $\mathfrak q_j$ de $B\otimes_A({}^{hs}\!A)$ au-dessus de $\mathfrak n_j$, tel que ${}^{hs}\!B$ soit isomorphe au composé direct des anneaux locaux $(B\otimes_A({}^{hs}\!A))_{\mathfrak q_j}$ ($1\leq j\leq r$).
\end{proposition}

On peut évidemment se borner au cas où $A$ est local, en remplaçant $A$ par $A_{\mathfrak m_j}$, où $\mathfrak m_j$ est l'idéal maximal de $A$ image réciproque de $\mathfrak n_j$. Posons d'autre part $C=B_{\mathfrak n_j}$; en vertu de la définition de ${}^{hs}\!B$ (18.8.1), tout revient à voir que ${}^{hs}\!C$ est $C$-isomorphe à $(C\otimes_A({}^{hs}\!A))_{\mathfrak q}$, où $\mathfrak q$ est un des idéaux maximaux de $C\otimes_A({}^{hs}\!A)$. Soient $k$ et $K$ les corps résiduels de $A$ et $C$, $\Omega$ celui de ${}^{hs}\!A$, qui est par définition une clôture séparable de $k$; comme $K$ est une extension finie de $k$, on peut supposer $K$ et $\Omega$ contenus dans une même clôture algébrique de $k$; il est alors immédiat que $K\otimes_k\Omega$ est composé direct d'un nombre fini de corps, tous isomorphes à la clôture séparable $K\Omega$ de $K$. D'autre part, l'anneau ${}^{hs}\!A$ étant hensélien, $C\otimes_A({}^{hs}\!A)$ est composé direct des anneaux locaux $(C\otimes_A({}^{hs}\!A))_{\mathfrak q_i}$, où $\mathfrak q_i$ ($1\leq i\leq s$) parcourt les idéaux maximaux de $C\otimes_A({}^{hs}\!A)$, et ces anneaux locaux sont henséliens (18.5.10) et ont $K\Omega$ pour corps résiduel, donc sont strictement locaux (18.8.1). Posons pour abréger $C'=C\otimes_A({}^{hs}\!A)$ et $C'_i=C'_{\mathfrak q_i}$, qui est donc un anneau quotient de $C'$; soient $f:C\to C'$, $g:{}^{hs}\!A\to C'$, $\varphi_i:C'\to C'_i$ les applications canoniques.

\[
\xymatrix{
  &&& {}^{hs}\!C \ar@<0.8ex>[dl]_{w'} \\
  && C'_i \ar@<0.8ex>[ur]_{w} \\
  C \ar[uurrr]^{v} \ar[r]_{f} & C' \ar[ur] \\
  A \ar[u] \ar[r] & {}^{hs}\!A \ar[u]^{g} \ar[uuurr]_{u}
}
\]

Soit $v:C\to{}^{hs}\!C$ l'homomorphisme canonique; il existe un homomorphisme local et un seul $u:{}^{hs}\!A\to{}^{hs}\!C$ qui, par passage aux quotients, donne l'injection canonique $\Omega\to K\Omega$ et tel que le composé $A\to C\xrightarrow{v}{}^{hs}\!C$ soit égal au composé $A\to{}^{hs}\!A\xrightarrow{u}{}^{hs}\!C$ (18.8.8); donc il existe un unique homomorphisme local $w_0:C'\to{}^{hs}\!C$ tel que $w_0\circ g=u$, $w_0\circ f=v$, et puisque ${}^{hs}\!C$ est un anneau local, cet homomorphisme se factorise en $C'\xrightarrow{\varphi_i}C'_i\xrightarrow{w}{}^{hs}\!C$ pour un indice $i$ bien déterminé, $w$ étant un homomorphisme local. D'autre part, l'anneau $C'_i$ étant strictement local, il existe un homomorphisme local $w':{}^{hs}\!C\to C'_i$ qui, par passage aux quotients, donne l'automorphisme identique de $K\Omega$ et qui est tel que $w'\circ v=\varphi_i\circ f$ (18.8.8). On déduit d'abord de là que $w\circ w'$ est un endomorphisme de ${}^{hs}\!C$ qui, par passage aux quotients, donne l'automorphisme identique de $K\Omega$, donc est

\oldpage[IV-4]{149}

l'identité (18.8.8). D'autre part, on a $w'\circ w\circ\varphi_i\circ f=w'\circ v=\varphi_i\circ f$, et $w'\circ w\circ\varphi_i\circ g=w'\circ u=\varphi_i\circ g$; d'où $w'\circ w\circ\varphi_i=\varphi_i$, autrement dit $w'\circ w$ est l'automorphisme identique de $C'_i$. C.Q.F.D.

\begin{remark}[18.8.11]
\label{IV.18.8.11}
La démonstration de (18.8.10) utilise le fait que $B$ est entière sur $A$, mais n'utilise le fait que $B$ est une $A$-algèbre \emph{finie}, que pour établir que $K\otimes_k\Omega$ est composé direct d'un nombre fini de corps. Or, cette dernière propriété est encore vérifiée lorsque l'on suppose que $B$ est une $A$-algèbre \emph{entière semi-locale} dont les idéaux maximaux $\mathfrak n_j$ ont des corps résiduels de degrés séparables $[k(\mathfrak n_j):k]_s$ finis.
\end{remark}

\begin{proposition}[18.8.12]
\label{IV.18.8.12}
Soit $A$ un anneau semi-local.
\begin{enumerate}
  \item[(i)] Pour que ${}^{hs}\!A$ soit réduit (resp. normal), il faut et il suffit que $A$ le soit.
  \item[(ii)] Supposons $A$ noethérien. Alors, pour tout idéal premier $\mathfrak p$ de $A$, l'anneau $({}^{hs}\!A)_{\mathfrak p}/\mathfrak p.({}^{hs}\!A)_{\mathfrak p}$ est composé direct d'un nombre fini d'extensions algébriques séparables de $k(\mathfrak p)$ (ce qui entraîne que les fibres du morphisme canonique $\Spec({}^{hs}\!A)\to\Spec(A)$ sont \emph{géométriquement régulières} et sont des espaces \emph{discrets}).
\end{enumerate}
\end{proposition}

\begin{corollary}[18.8.13]
\label{IV.18.8.13}
Soit $A$ un anneau local noethérien. Pour que ${}^{hs}\!A$ possède une des propriétés suivantes:
\begin{enumerate}
  \item[$a)$] être un anneau de Cohen-Macaulay (0, 16.5.3);
  \item[$b)$] vérifier la condition $(S_n)$ (5.7.3);
  \item[$c)$] être régulier;
  \item[$d)$] vérifier la condition $(R_n)$ (5.8.2);
\end{enumerate}
il faut et il suffit que $A$ possède cette même propriété.
\end{corollary}

\begin{corollary}[18.8.14]
\label{IV.18.8.14}
Soit $A$ un anneau semi-local noethérien. Pour qu'un idéal premier $\mathfrak p$ de ${}^{hs}\!A$ appartienne à $\operatorname{Ass}({}^{hs}\!A)$, il faut et il suffit que $\mathfrak p\cap A$ appartienne à $\operatorname{Ass}(A)$.
\end{corollary}

Les démonstrations sont les mêmes que celles de (18.6.9), (18.6.10) et (18.6.11) respectivement.

\begin{proposition}[18.8.15]
\label{IV.18.8.15}
Soit $A$ un anneau local; les propriétés suivantes sont équivalentes:
\begin{enumerate}
  \item[$a)$] $A$ est géométriquement unibranche (resp. réduit et géométriquement unibranche).
  \item[$b)$] Pour toute $A$-algèbre locale essentiellement étale $A'$, $\Spec(A')$ est irréductible (resp. intègre).
  \item[$c)$] $\Spec({}^{hs}\!A)$ est irréductible (resp. intègre).
\end{enumerate}
\end{proposition}

On se ramène comme dans (18.6.12) au cas où $A$ est réduit, en utilisant la relation ${}^{hs}(A_{\mathrm{red}})=({}^{hs}\!A)\otimes_A A_{\mathrm{red}}$ (18.8.11), et l'équivalence de $b)$ et $c)$ se prouve comme dans (18.6.12); il en est de même du fait que $a)$ entraîne $c)$, en tenant compte de (18.8.11) et de la définition d'un anneau local intègre géométriquement unibranche. Enfin, pour démontrer que $c)$ entraîne $a)$, gardons les mêmes notations que dans (18.6.12); il s'agit encore de voir que ${}^{hs}\!B$ est intègre. Mais, par (18.8.10), ${}^{hs}\!B$ est un anneau localisé de l'anneau semi-local $B\otimes_A({}^{hs}\!A)$, qui s'identifie encore par platitude à un sous-anneau du corps $L$, donc est intègre; il en est par suite de même de ${}^{hs}\!B$, ce qui achève la démonstration.

\begin{corollary}[18.8.16]
\label{IV.18.8.16}
Soit $A$ un anneau local hensélien (resp. strictement local). Pour que $A$ soit unibranche (resp. géométriquement unibranche), il faut et il suffit que $\Spec(A)$ soit irréductible.
\end{corollary}

\oldpage[IV-4]{150}

\begin{proposition}[18.8.17]
\label{IV.18.8.17}
Soit $A$ un anneau local noethérien. Si $A$ est universellement caténaire, il en est de même de ${}^{hs}\!A$.
\end{proposition}

La démonstration est la même que celle de (18.7.5).

\begin{proposition}[18.8.18]
\label{IV.18.8.18}
\begin{enumerate}
  \item[(i)] Toute limite inductive filtrante d'anneaux strictement locaux (où les homomorphismes de transition sont locaux) est un anneau strictement local.
  \item[(ii)] Le foncteur $A\leadsto{}^{hs}\!A$ dans la catégorie des anneaux locaux (où les morphismes sont les homomorphismes locaux) commute aux limites inductives filtrantes.
  \item[(iii)] Tout anneau strictement local $A$ est limite inductive d'un système inductif filtrant d'anneaux strictement locaux noethériens (les homomorphismes de transition étant locaux).
\end{enumerate}
\end{proposition}

Les démonstrations sont calquées sur celles de (18.6.14).

\subsection{Fibres formelles des anneaux noethériens henséliens.}
\EGAReaderTarget{ega4.c04.subsec.18-9}{subsection:IV.18.9}

\begin{theorem}[18.9.1]
\label{IV.18.9.1}
Soit $A$ un anneau local noethérien hensélien, dont les fibres formelles sont géométriquement normales. Alors les fibres formelles de $A$ sont géométriquement intègres (donc géométriquement connexes).
\end{theorem}

Les deux assertions de l'énoncé sont en fait équivalentes, un préschéma localement noethérien connexe et normal étant intègre. Comme toute $A$-algèbre intègre finie est un anneau local hensélien (18.5.9 et 18.5.10), il résulte de (7.3.16.2) que l'on est ramené à prouver que si l'on suppose de plus $A$ \emph{intègre}, alors la fibre du morphisme $\Spec(\widehat A)\to\Spec(A)$ au point générique de $\Spec(A)$ est \emph{intègre}, ce qui résulte du

\begin{corollary}[18.9.2]
\label{IV.18.9.2}
Soit $A$ un anneau local noethérien hensélien intègre, dont les fibres formelles sont géométriquement normales. Alors $\widehat A$ est intègre.
\end{corollary}

En effet, il résulte de (18.8.16) que $A$ est unibranche, donc le théorème résulte de l'hypothèse faite sur les fibres formelles de $A$ et de (7.6.3).

\begin{corollary}[18.9.3]
\label{IV.18.9.3}
Sous les hypothèses de (18.9.2), le corps des fractions $L$ de $\widehat A$ est une extension séparable du corps des fractions $K$ de $A$, et $K$ est algébriquement fermé dans $L$.
\end{corollary}

Cela résulte de ce que la fibre du morphisme $\Spec(\widehat A)\to\Spec(A)$ au point générique de $\Spec(A)$ est géométriquement intègre et de (4.3.2) et (4.3.5).

\begin{corollary}[18.9.4]
\label{IV.18.9.4}
Soit $A$ un anneau local noethérien, hensélien et dont les fibres formelles sont géométriquement normales (par exemple un anneau local noethérien hensélien et excellent), et posons $A'=\widehat A$. Soient $X$ un $A$-préschéma, $X'=X\otimes_A A'$, $g:X'\to X$ la projection canonique. Alors l'application $U\mapsto g^{-1}(U)$ est une bijection de l'ensemble des parties à la fois ouvertes et fermées de $X$ sur l'ensemble des parties à la fois ouvertes et fermées de $X'$.
\end{corollary}

Comme $g$ est un morphisme fidèlement plat et quasi-compact, on sait (2.3.12) que la topologie de $X$ est \emph{quotient} de celle de $X'$ par la relation d'équivalence définie par $g$. Tout revient donc à voir qu'une partie ouverte et fermée $U'$ de $X'$ est \emph{saturée} pour cette relation. Or, comme le morphisme $\Spec(A')\to\Spec(A)$ a ses fibres géométriquement connexes (18.9.1), les fibres de $g$ sont connexes, d'où aussitôt notre assertion.

En particulier, si $X$ est localement connexe (ce qui sera le cas si $X$ est localement noethérien), alors, pour toute composante connexe $U$ de $X$ (qui est ouverte et fermée

\oldpage[IV-4]{151}

dans $X$), $g^{-1}(U)$ est connexe (Bourbaki, \emph{Top. gén.}, chap.~I\textsuperscript{er}, 3\textsuperscript{e} éd., \textsection~11, n\textsuperscript{o}~3, prop.~7). Si l'on désigne par $\pi_0(X)$ l'ensemble des composantes connexes de $X$, l'application $\pi_0(X')\to\pi_0(X)$ déduite canoniquement de $g$ est donc \emph{bijective}.

\begin{corollary}[18.9.5]
\label{IV.18.9.5}
Sous les hypothèses de (18.9.4), le foncteur
\begin{equation}
  Z\longmapsto Z\otimes_A A'
  \tag{18.9.5.1}
  \label{IV.18.9.5.1}
\end{equation}
de la catégorie des préschémas étales sur $X$, dans celle des préschémas étales sur $X'$, est pleinement fidèle.
\end{corollary}

Soient $Z_1$, $Z_2$ deux préschémas étales sur $X$, et posons $Z'_i=Z_i\otimes_A A'$ ($i=1,2$); il s'agit de prouver que tout $X'$-morphisme $Z'_1\to Z'_2$ provient, par changement de base, d'un $X$-morphisme unique $Z_1\to Z_2$. Supposons d'abord que $Z_2$ soit \emph{séparé} sur $X$; alors, comme $\Hom_X(Z_1,Z_2)$ s'identifie à $\Gamma((Z_1\times_X Z_2)/Z_1)$, et que $Z_1\times_X Z_2$ est étale et séparé sur $Z_1$, $\Hom_X(Z_1,Z_2)$ s'identifie fonctoriellement, en vertu de (17.9.3), à l'ensemble des parties ouvertes et fermées $U$ de $Z_1\times_X Z_2$ telles que la restriction à $U$ de la projection $p:Z_1\times_X Z_2\to Z_1$ soit un morphisme surjectif et radiciel. L'assertion résulte donc dans le cas considéré de (18.9.4) et de ce que $Z'_1\times_{X'}Z'_2=(Z_1\times_X Z_2)\otimes_A A'$.

Passons au cas général: il suffira de prouver que le graphe $\Gamma'$ d'un $X'$-morphisme $Z'_1\to Z'_2$, qui est ouvert (17.9.3) dans $Z'_1\times_{X'}Z'_2$, est de la forme $\Gamma\otimes_A A'$, où $\Gamma$ est induit sur une partie \emph{ouverte} de $Z=Z_1\times_X Z_2$. En effet, la restriction à $\Gamma$ de la projection $p:Z\to Z_1$ est alors un \emph{isomorphisme}, car par changement de base $A\to A'$ cette restriction devient l'isomorphisme restriction à $\Gamma'$ de la projection $p':Z'\to Z'_1$ (avec $Z'=Z\otimes_A A'$), et il suffit d'appliquer (2.7.1, (viii)); la conclusion résulte alors de la caractérisation des graphes de morphismes (I, 5.3).

Si $q:Z'\to Z$ est la projection canonique, pour prouver que $\Gamma'$ est de la forme $q^{-1}(\Gamma)$, où $\Gamma$ est ouvert dans $Z$, il suffit, puisque $q$ est un morphisme fidèlement plat et quasi-compact, de prouver qu'il existe un ensemble $U\subset Z$ tel que $\Gamma'=q^{-1}(U)$ (2.3.12). Posons $Z''=Z'\times_Z Z'=Z\times_X X''$, où $X''=X'\times_X X'$, et soient $q_1:Z''\to Z'$, $q_2:Z''\to Z'$ les projections canoniques. Appliquant (4.5.19.1), il suffit de montrer que $q_1^{-1}(\Gamma')=q_2^{-1}(\Gamma')$. Mais c'est une propriété qui est vraie si et seulement si elle l'est après chaque changement de base $\Spec(k(x))\to X$, où $x$ parcourt $X$. Autrement dit, il suffit de prouver le corollaire lorsque $X$ est le \emph{spectre d'un corps}; mais comme tout $X$-préschéma étale est alors automatiquement \emph{séparé} sur $X$ (17.6.2, $c'$)), on est ramené au cas envisagé au début de la démonstration.

\begin{remarks}[18.9.6]
\label{IV.18.9.6}
\begin{enumerate}
  \item[(i)] Il est possible que, si le corps résiduel de $A$ est de caractéristique 0, le foncteur (18.9.5.1) induise même une \emph{équivalence} de la catégorie des revêtements étales de $X$ et de la catégorie des revêtements étales de $X'$. On peut en tout cas le démontrer lorsque $A$ est \emph{excellent}, en utilisant la résolution des singularités de Hironaka (M.~Artin). Il est plausible que l'énoncé analogue soit encore vrai sans restriction sur la caractéristique du corps résiduel, à condition de se borner aux revêtements « galoisiens principaux » dont le groupe est d'ordre premier à la caractéristique résiduelle (cf. [41]).

\oldpage[IV-4]{152}

  \item[(ii)] On ignore si les fibres formelles de $A$ sont géométriquement connexes (voire géométriquement irréductibles) lorsque $A$ est un anneau local hensélien quelconque. Il suffirait que, pour tout anneau local noethérien, hensélien et \emph{intègre} $A$, $\Spec(\widehat A)$ soit irréductible. On ignore s'il en est toujours ainsi.

  \item[(iii)] On ignore si lorsque $A$ est un anneau local excellent, son hensélisé strict ${}^{hs}\!A$ est excellent. On peut voir que, pour élucider cette question, on peut se ramener au cas où $A=k[[T_1,\ldots,T_n]]$, $k$ étant un corps de caractéristique $p>0$. Compte tenu de (18.8.17), la question est de savoir si les fibres formelles de ${}^{hs}\!A$ sont géométriquement régulières. La réponse est affirmative lorsque $n=1$, mais n'est pas connue pour $n=2$. On peut montrer que la réponse est affirmative chaque fois que, pour tout schéma $Y_1$ fini sur $Y=\Spec(\widehat A)$, on peut résoudre les singularités de $Y_1$ (7.9.1).
\end{enumerate}
\end{remarks}

L'assertion de connexité faite dans (18.9.1) se généralise de la façon suivante:

\begin{theorem}[18.9.7]
\label{IV.18.9.7}
Soient $A$, $B$ deux anneaux locaux noethériens, $\rho:A\to B$ un homomorphisme local, $f:\Spec(B)=X\to\Spec(A)=Y$ le morphisme correspondant. Supposons vérifiées les conditions suivantes:
\begin{enumerate}
  \item[(i)] $f$ est plat et toutes les fibres $X_y=f^{-1}(y)$ ($y\in Y$) sont géométriquement réduites (4.6.2).
  \item[(ii)] Ou bien l'anneau $A$ est géométriquement unibranche (0, 23.2.1), ou bien l'anneau $A$ est unibranche (0, 23.2.1) et $k(B)$ est une extension primaire (4.3.1) de $k(A)$.
\end{enumerate}
Alors, si $\eta$ est le point générique de $Y$, $X_\eta$ est connexe, et pour toute partie fermée $T$ de $X$ telle que $f(T)$ soit rare dans $Y$, $X-T$ est connexe.
\end{theorem}

Les deux conclusions énoncées sont en fait équivalentes; en effet, l'hypothèse que $f(T)$ est rare entraîne que $f(T)$ ne contient pas $\eta$; d'autre part, comme $\{\eta\}$ est proconstructible (I, 9.6) et que $f$ est plat, $X_\eta$ est dense dans $X$ (2.3.10), et par suite si $X_\eta$ est connexe, il en est de même de $X-T$ qui le contient. Inversement, lorsque $U$ parcourt l'ensemble des voisinages ouverts affines de $\eta$ dans $Y$, $X_\eta$ est limite projective des sous-préschémas induits sur les ouverts $f^{-1}(U)$ de $X$ (8.1.2, $a)$); si l'on suppose que $X-T$ est connexe pour toute partie fermée $T$ de $X$ telle que $f(T)$ soit rare dans $Y$, les $f^{-1}(U)$ sont connexes, donc il en est de même de $X_\eta$ (8.4.1).

Montrons d'abord qu'on peut se limiter au cas où $A$ est intègre et $Y=\Spec(A)$ géométriquement unibranche (6.15.1) (ce qui implique que $A$ est géométriquement unibranche, mais n'est pas équivalent à cette condition (6.15.2)). Il est clair qu'on peut d'abord supposer $Y$ réduit, en considérant le morphisme $X\times_Y Y_{\mathrm{red}}\to Y_{\mathrm{red}}$, déduit de $f$ par changement de base, qui est plat, et a mêmes fibres que $f$. On peut donc supposer $A$ intègre et unibranche; si $K$ est le corps des fractions de $A$, il existe alors une sous-$A$-algèbre finie $A''$ de $K$ telle que, si $A'$ est la clôture intégrale de $A$, le morphisme $\Spec(A')\to\Spec(A'')$ soit radiciel (0, 23.2.5); on en conclut (6.15.5) que $\Spec(A'')$ est géométriquement unibranche. Comme $A$ est unibranche, $A'$ est un anneau local, donc il en est de même de $A''$; montrons que l'anneau $B''=B\otimes_A A''$ est aussi un anneau local. En effet, $B''$ est une $B$-algèbre finie, donc un anneau semi-local noethérien; si $k''$ est le corps résiduel de $A''$, les idéaux maximaux de $B''$ sont les points de $\Spec(B'')$ au-dessus du point fermé de $\Spec(B)$, donc les points de $\Spec(k(B)\otimes_{k(A)}k'')$ (I, 3.4.9)

\oldpage[IV-4]{153}

puisque le point fermé de $\Spec(A'')$ est le seul point au-dessus du point fermé de $\Spec(A)$. Or, lorsque $A$ est géométriquement unibranche, $k''$ est une extension radicielle de $k(A)$, donc $\Spec(k(B)\otimes_{k(A)}k'')$ est radiciel sur $\Spec(k(B))$ et ne comprend donc qu'un seul point. Lorsque $A$ est unibranche et $k(B)$ extension primaire de $k(A)$, $\Spec(k(B)\otimes_{k(A)}k'')$ est irréductible (4.3.2) et comme c'est un espace fini et discret (I, 6.4.4) il est encore réduit à un point, ce qui montre que dans les deux cas envisagés $B''$ est local. Si $Y''=\Spec(A'')$, $X''=\Spec(B'')=X\times_Y Y''$, le morphisme $f''=f_{(Y'')}:X''\to Y''$ est plat et a ses fibres géométriquement réduites; en outre, si $\eta''$ est le point générique de $Y''$, $k(\eta'')=k(\eta)=K$ par définition, donc les fibres $X_\eta$ et $X''_{\eta''}$ sont isomorphes. Il suffit donc de prouver que $X''_{\eta''}$ est connexe.

\begin{remark}[18.9.7.1]
\label{IV.18.9.7.1}
Lorsque l'anneau $A$ est \emph{japonais}, on peut, dans ce qui précède, prendre $A''=A'$ par définition (0, 23.1.1); on voit donc que dans ce cas on serait même ramené à prouver le théorème lorsque $A$ est \emph{intégralement clos}.
\end{remark}

\begin{env}[18.9.7.2]
\label{IV.18.9.7.2}
Supposons donc désormais que $Y$ soit intègre et géométriquement unibranche. Notons que si $T$ est une partie fermée de $X$ telle que $f(T)$ soit rare dans $Y$, $T$ ne contient aucun point maximal de $X$ puisque $f$ est plat (2.3.4), donc est rare dans $X$; comme $X$ est un schéma local, donc connexe, pour prouver que $X-T$ est connexe, il suffit, en vertu de (15.5.6.1) de montrer que pour tout $x\in X$ tel que $f(x)\neq\eta$, $\Spec(\sh{O}_{X,x})-\{x\}$ est connexe. Posons $y=f(x)$, $A_1=\sh{O}_{Y,y}$, $B_1=\sh{O}_{X,x}$, $Y_1=\Spec(A_1)$, $X_1=\Spec(B_1)$, $f_1:X_1\to Y_1$ le morphisme correspondant à l'homomorphisme local $A_1\to B_1$ déduit de $\rho$; il est clair que $f_1$ est plat, et il résulte de (4.6.1) que ses fibres sont géométriquement réduites; en outre $A_1$ est intègre et géométriquement unibranche (6.15.1) et $\dim(A_1)\geq 1$. On est ainsi ramené à prouver le
\end{env}

\begin{lemma}[18.9.7.3]
\label{IV.18.9.7.3}
Soient $A$, $B$ deux anneaux locaux noethériens, $Y=\Spec(A)$, $X=\Spec(B)$, $\rho:A\to B$ un homomorphisme local, $f:X\to Y$ le morphisme correspondant. On suppose vérifiées les conditions suivantes:
\begin{enumerate}
  \item[(i)] $f$ est plat et toutes les fibres $X_y$ ($y\in Y$) sont géométriquement réduites.
  \item[(ii)] $A$ est intègre, géométriquement unibranche et $\dim(A)\geq 1$.
\end{enumerate}
Alors, si $b$ est le point fermé de $X$, $X-\{b\}$ est connexe.
\end{lemma}

Notons d'abord qu'en vertu du théorème de Hartshorne (5.10.7), $X-\{b\}$ est connexe si l'on a $\operatorname{prof}(B)\geq 2$. Or on a (6.3.1)
\begin{equation}
  \operatorname{prof}(B)=\operatorname{prof}(A)+\operatorname{prof}(B\otimes_A k)
  \tag{18.9.7.4}
  \label{IV.18.9.7.4}
\end{equation}
en désignant par $k$ le corps résiduel de $A$. D'autre part, puisque $A$ est intègre et $\dim(A)\geq 1$, on a $\operatorname{prof}(A)\geq 1$, donc $\operatorname{prof}(B)\geq 2$ \emph{sauf} lorsque $\operatorname{prof}(B\otimes_A k)=0$; d'ailleurs $B\otimes_A k$ est réduit par (i), donc la relation $\operatorname{prof}(B\otimes_A k)=0$ signifie que $B\otimes_A k$ est un \emph{corps} (0, 16.4.7). On supposera donc désormais que cette dernière condition est vérifiée. On commencera par traiter un cas où la démonstration est très simple.

A) \emph{Cas où $A$ est intégralement clos.} Alors, si $\dim(A)\geq 2$, on a $\operatorname{prof}(A)\geq 2$ (0, 16.5.1), donc $\operatorname{prof}(B)\geq 2$. On n'a donc qu'à considérer le cas où $\dim(A)=1$ et où $B\otimes_A k$ est un corps; $A$ est alors un anneau de valuation discrète, donc régulier, et puisque $B\otimes_A k$

\oldpage[IV-4]{154}

est un corps et que $f$ est plat, $B$ est régulier (0, 17.3.3); mais en outre (6.1.1.1), on a $\dim(B)=\dim(A)+\dim(B\otimes_A k)=\dim(A)=1$, donc $B$ est aussi un anneau de valuation discrète. Mais alors $X-\{b\}$ est réduit à un seul point, d'où le lemme (18.9.7.3) dans ce cas.

On observera que cela prouve l'énoncé du théorème (18.9.7) lorsqu'on suppose en outre que, dans cet énoncé, l'anneau $A$ est \emph{japonais}, car en vertu de la remarque (18.9.7.1), on est ramené au cas où $A$ est \emph{intégralement clos}, et alors, dans la réduction qui précède (18.9.7.3), l'anneau $A_1=\sh{O}_{Y,y}$ est lui aussi intégralement clos et on peut appliquer le résultat qui vient d'être démontré.

B) \emph{Cas où $\dim(A)\geq 2$.} Le résultat de (18.9.7.3) résultera alors des deux lemmes suivants:

\begin{lemma}[18.9.7.5]
\label{IV.18.9.7.5}
Soient $A$ un anneau semi-local noethérien et réduit, $A'$ sa fermeture intégrale dans son anneau total des fractions, $X=\Spec(A)$, $X'=\Spec(A')$; soient $S$ l'ensemble des points fermés de $X$ et $U=X-S$. Supposons que pour tout point $x'\in X'$ au-dessus d'un point de $S$, on ait $\dim(\sh{O}_{X',x'})\geq 2$; alors l'anneau $A^{(1)}=\Gamma(U,\sh{O}_X)$ est entier sur $A$ et est une $A$-algèbre isomorphe à une sous-algèbre de $A'$.
\end{lemma}

Rappelons que si les $\mathfrak p_i$ ($1\leq i\leq m$) sont les idéaux minimaux de $A$, $A'$ est le composé direct des clôtures intégrales $A'_i$ des anneaux intègres $A_i=A/\mathfrak p_i$, de sorte que $X'$ est la somme des schémas $X'_i=\Spec(A'_i)$. Si $X_0$ est la somme des schémas $X_i=\Spec(A_i)$ ($1\leq i\leq m$) et $U_0$ l'image réciproque de $U$ dans $X_0$, l'homomorphisme canonique $\Gamma(U,\sh{O}_X)\to\Gamma(U_0,\sh{O}_{X_0})$ est injectif puisque $X$ est réduit et le morphisme $X_0\to X$ (donc aussi le morphisme $U_0\to U$) surjectif. Comme $\Gamma(U_0,\sh{O}_{X_0})$ est composé direct des $\Gamma(U_i,\sh{O}_{X_i})$ ($1\leq i\leq m$), où $U_i$ est l'image réciproque de $U$ dans $X_i$, tout revient à prouver que pour chaque $i$, $\Gamma(U_i,\sh{O}_{X_i})$ est isomorphe à une sous-$A_i$-algèbre de $A'_i$. Autrement dit, on est ramené au cas où $A$ est un anneau semi-local noethérien intègre. Comme $X$ est réduit et le morphisme $X'\to X$ surjectif, l'homomorphisme $\Gamma(U,\sh{O}_X)\to\Gamma(U',\sh{O}_{X'})$ (où $U'$ est l'image réciproque de $U$ dans $X'$) est injectif, et tout revient à voir que l'homomorphisme canonique $A'=\Gamma(X',\sh{O}_{X'})\to\Gamma(U',\sh{O}_{X'})$ est \emph{bijectif}. Or (I, 8.2.1.1), $\Gamma(U',\sh{O}_{X'})$ est l'intersection des anneaux locaux $A'_{\mathfrak p'}$, où $\mathfrak p'$ parcourt $U'$, et en vertu de l'hypothèse, parmi ces anneaux locaux figurent tous ceux pour lesquels $\mathfrak p'$ est de hauteur 1. Mais le raisonnement de (0, 23.2.7) s'applique aussi à un anneau semi-local intègre noethérien, donc $A'$ est un \emph{anneau de Krull semi-local}, et est donc l'intersection des anneaux locaux $A'_{\mathfrak p'}$, où $\mathfrak p'$ parcourt l'ensemble des idéaux premiers de hauteur 1 de $A'$ (Bourbaki, \emph{Alg. comm.}, chap.~VII, \textsection~1, n\textsuperscript{o}~6, th.~4); \emph{a fortiori}, $A'$ est l'intersection des $A'_{\mathfrak p'}$ pour $\mathfrak p'\in U'$, ce qui achève la démonstration du lemme.

\begin{remark}[18.9.7.6]
\label{IV.18.9.7.6}
On sait (0, 23.2.5) qu'il existe une sous-$A$-algèbre finie $A''$ du corps des fractions $K$ de $A$ telle que le morphisme $\Spec(A')\to\Spec(A'')$ soit \emph{radiciel}; comme ce morphisme est aussi entier et dominant, donc surjectif (II, 6.1.10), c'est un \emph{homéomorphisme} (2.4.5). En vertu de l'interprétation géométrique de $\dim(\sh{O}_{X',x'})$ (5.1.2), on a donc, pour tout point $x''$ de $X''=\Spec(A'')$ au-dessus d'un point $x\in X$, $\dim(\sh{O}_{X'',x''})=\dim(\sh{O}_{X',x'})\geq 2$ si $x'$ est l'unique point de $X'$ au-dessus de $x$. Cela

\oldpage[IV-4]{155}

étant, il résulte d'abord de (0, 16.1.5) et (0, 16.1.4.1) que la relation $\dim(\sh{O}_{X',x'})\geq 2$ entraîne $\dim(\sh{O}_{X,x})\geq 2$. Inversement, supposons que $\dim(\sh{O}_{X,x})\geq 2$ et que $A$ soit universellement caténaire (5.6.2); alors il en est de même de $\sh{O}_{X,x}$ (5.6.3), donc $\dim(\sh{O}_{X'',x''})=\dim(\sh{O}_{X,x})\geq 2$ en vertu de (5.6.10) et on a donc alors $\dim(\sh{O}_{X',x'})\geq 2$. La même conclusion est valable si $\sh{O}_{X,x}$ est \emph{géométriquement unibranche}, car la fibre du morphisme $X'\to X$ au point $x$ est alors réduite à un seul point, et $\sh{O}_{X',x'}$ est une $\sh{O}_{X,x}$-algèbre \emph{entière}; il suffit donc d'appliquer (0, 16.1.5).
\end{remark}

\begin{lemma}[18.9.7.7]
\label{IV.18.9.7.7}
Soient $A$ un anneau local noethérien, de dimension $\geq 2$, intègre et géométriquement unibranche, $Y=\Spec(A)$, $y$ le point fermé de $Y$. Soient $X$ un préschéma localement noethérien et connexe, $f:X\to Y$ un morphisme plat. Alors, pour toute partie fermée $T\subset f^{-1}(y)$, $X-T$ est connexe.
\end{lemma}

Le même raisonnement qu'au début de (18.9.7.2) ramène à prouver que, pour tout point $x\in f^{-1}(y)$, $\Spec(\sh{O}_{X,x})-\{x\}$ est connexe. On peut donc se borner au cas où $X=\Spec(B)$, $B$ étant un anneau local noethérien et l'homomorphisme $A\to B$ correspondant à $f$ étant local; posons $U=Y-\{y\}$, et notons que $V=f^{-1}(U)$ est dense dans $X$ (2.3.10); il suffit donc de prouver que $V$ est \emph{connexe}. Gardons les notations de (18.9.7.5), et posons donc $A^{(1)}=\Gamma(U,\sh{O}_Y)$, $Y_1=\Spec(A^{(1)})$, $X_1=X\times_Y Y_1=\Spec(B\otimes_A A^{(1)})$, de sorte qu'on a le diagramme commutatif
\[
  \xymatrix{
    X \ar[d]_{f} & X_1 \ar[l]^{g_1} \ar[d]^{f_1} \\
    Y & Y_1 \ar[l]_{g}
  }
\]

En vertu de (2.3.1) et de la définition de $A^{(1)}$, on a un isomorphisme canonique $\Gamma(V,\sh{O}_X)\xrightarrow{\sim}B\otimes_A A^{(1)}$, et il suffit donc de prouver que ce dernier anneau est \emph{local}, un tel anneau ne pouvant être le produit de deux anneaux non réduits à 0 (III, 7.8.6.1). Mais on a vu (18.9.7.5) que $A^{(1)}$ est une sous-$A$-algèbre de la clôture intégrale $A'$ de $A$, puisque $\dim(A)\geq 2$. L'hypothèse que $A$ est géométriquement unibranche entraîne que le morphisme $\Spec(A')\to\Spec(A)$ est radiciel au point $y$ (6.15.3); \emph{a fortiori} $g:Y_1\to Y$ est radiciel au point $y$, et par suite $g_1$ est radiciel au point fermé $x$ de $X$; comme en outre $g_1$ est un morphisme entier, les points fermés de $X_1$ sont au-dessus de $x$, donc il n'existe qu'un seul point fermé de $X_1$, ce qui achève la démonstration de (18.9.7.7).

Il est clair que (18.9.7.7) démontre (18.9.7.3) lorsque $\dim(A)\geq 2$ (même sans supposer les fibres $X_y$ géométriquement réduites).

C) \emph{Cas où $\dim(A)=1$.} On a vu au début de la preuve de (18.9.7.3) que l'on peut supposer $\dim(B\otimes_A k)=0$, et par platitude (6.1.1.1), on a donc
\[
  \dim(B)=\dim(A)+\dim(B\otimes_A k)=1.
\]
Puisque $A$ est intègre, $Y$ se compose de deux points, le point générique $\eta$ et le point fermé $a$. En outre, puisque $f$ est plat, toute composante irréductible de $X$ domine $Y$ (2.3.4), donc $X-\{b\}$, qui est l'ensemble des points maximaux $\xi_i$ de $X$, est égal à

\oldpage[IV-4]{156}

l'ensemble sous-jacent à $f^{-1}(\eta)$, et la fibre $f^{-1}(\eta)$ est somme des $\Spec(k(\xi_i))$. L'hypothèse sur les $X_y$ et le fait que ces fibres sont de dimension 0, montrent qu'elles sont \emph{géométriquement régulières} (6.7.6), et \emph{a fortiori} géométriquement normales, autrement dit $f$ est un \emph{morphisme normal}. Mais puisque $A$ est intègre et géométriquement unibranche, il résulte alors de (6.15.10) que $B$ est intègre et géométriquement unibranche, donc $f^{-1}(\eta)=X-\{b\}$ est réduit à un seul point, et \emph{a fortiori} connexe; ceci termine donc la démonstration de (18.9.7.3) et celle de (18.9.7).

\begin{remark}[18.9.7.8]
\label{IV.18.9.7.8}
Dans le cas C) de la démonstration de (18.9.7.3), on peut éviter de faire appel au délicat résultat (6.15.10) en raisonnant de la façon suivante: puisque $A$ est intègre et de dimension 1, il résulte du théorème de Krull-Akizuki (Bourbaki, \emph{Alg. comm.}, chap.~VII, \textsection~2, n\textsuperscript{o}~5, prop.~5) que sa clôture intégrale $A'$ est un anneau \emph{noethérien}. Le même raisonnement qu'au début de la démonstration de (18.9.7) montre alors que $B'=B\otimes_A A'$ est un anneau local; en outre, par platitude, $B'$ est contenu dans l'anneau total des fractions $R$ de l'anneau réduit $B$ (3.3.5). Or, le raisonnement qui prouve le théorème de Krull-Akizuki (Bourbaki, \emph{loc. cit.}) s'applique également à un anneau local noethérien \emph{réduit} de dimension 1, et montre que pour un tel anneau $B$, tout anneau compris entre $B$ et son anneau total des fractions est \emph{noethérien}. L'anneau $B'$ étant un anneau local noethérien et le morphisme $\Spec(B')\to\Spec(A')$ étant normal, $B'$ est un anneau intègre et intégralement clos (6.5.4), et on conclut comme au début de la preuve de (18.9.7).
\end{remark}

\begin{corollary}[18.9.8]
\label{IV.18.9.8}
Avec les notations de (18.9.7), supposons vérifiée la condition (i) de (18.9.7) et l'une des deux conditions suivantes:
\begin{enumerate}
  \item[(ii$'$)] $A$ est un anneau strictement local (18.8.2).
  \item[(ii$''$)] $A$ est un anneau local hensélien et $k(B)$ est une extension primaire de $k(A)$.
\end{enumerate}
Alors toutes les fibres $X_y$ ($y\in Y$) sont géométriquement connexes.
\end{corollary}

Soit $\mathfrak p$ l'idéal premier de $A$ correspondant au point $y$, et posons $A'=A/\mathfrak p$ et $B'=B\otimes_A A'=B/\mathfrak pB$, qui sont évidemment des anneaux locaux noethériens; il est clair que le morphisme $f':\Spec(B')\to\Spec(A')$ vérifie la condition (i) de (18.9.7); comme $y$ est le point générique de $\Spec(A')$ et que $k(A')=k(A)$, $k(B')=k(B)$, on voit qu'il suffit de vérifier que, dans le cas (ii$'$) (resp. (ii$''$)), $A'$ est géométriquement unibranche (resp. unibranche). Mais $A'$ est intègre, et dans le cas (ii$'$) (resp. (ii$''$)) il est strictement local (resp. hensélien) en vertu de (18.5.10). On en conclut que $A'$ est géométriquement unibranche (resp. unibranche) en vertu de (18.8.15) (resp. (18.6.12)).

\begin{corollary}[18.9.9]
\label{IV.18.9.9}
Soit $A$ un anneau local noethérien hensélien, universellement japonais (\emph{i.e.} ((7.6.4) et (7.7.2)), dont les fibres formelles (7.3.13) sont géométriquement réduites); alors les fibres formelles de $A$ sont géométriquement connexes.
\end{corollary}

Il suffit d'appliquer (18.9.8) au cas où $B=\widehat A$, puisque $B$ est un $A$-module plat et que $k(B)=k(A)$.

\begin{remark}[18.9.10]
\label{IV.18.9.10}
Les démonstrations de (18.9.4) et (18.9.5) n'utilisent que le fait que les fibres formelles de $A$ sont géométriquement connexes. Les conclusions

\oldpage[IV-4]{157}

de ces deux propositions sont donc encore valables lorsqu'on suppose l'anneau local $A$ noethérien, hensélien et universellement japonais.
\end{remark}

\begin{corollary}[18.9.11]
\label{IV.18.9.11}
Soient $X$, $Y$ deux préschémas localement noethériens, $Y$ étant supposé géométriquement unibranche et $X$ connexe. Soit $f:X\to Y$ un morphisme réduit (\emph{i.e.} (6.8.1) plat et à fibres géométriquement réduites). Alors, pour toute partie fermée $T$ de $X$ telle que $f(T)$ soit rare dans $Y$, $X-T$ est connexe.
\end{corollary}

Comme $f$ est plat et que $f(T)$, étant rare, ne peut contenir de point maximal de $Y$, $T$ ne peut contenir de point maximal de $X$ (2.3.4), donc est rare dans $X$. Utilisant (15.5.6.1), il suffit de montrer que, pour tout $x\in T$, $\Spec(\sh{O}_{X,x})-\{x\}$ est connexe. Posons $y=f(x)$, qui n'est pas un point maximal de $Y$. Comme l'anneau $\sh{O}_{Y,y}$ est géométriquement unibranche, on peut appliquer (18.9.7) au morphisme
\[
  f_1:\Spec(\sh{O}_{X,x})\longrightarrow\Spec(\sh{O}_{Y,y}),
\]
qui vérifie les conditions (i) et (ii) de ce théorème, et à la partie fermée $\{x\}$ de $\Spec(\sh{O}_{X,x})$, puisque $f_1(\{x\})=\{y\}$ est rare dans $\Spec(\sh{O}_{Y,y})$.

\subsection{Préschémas étales sur un préschéma géométriquement unibranche ou normal}
\EGAReaderTarget{ega4.c04.subsec.18-10}{subsection:IV.18.10}

\begin{theorem}[18.10.1]
\label{IV.18.10.1}
Soient $f:X\to Y$ un morphisme localement de présentation finie, $x$ un point de $X$, $y=f(x)$. Supposons que $Y$ soit intègre et géométriquement unibranche au point $y$. Alors, pour que $f$ soit étale au point $x$, il faut et il suffit que les deux conditions suivantes soient vérifiées:
\begin{enumerate}
  \item[(i)] $f$ est non ramifié au point $x$;
  \item[(ii)] l'homomorphisme $\sh{O}_{Y,y}\to\sh{O}_{X,x}$ est injectif.
\end{enumerate}
De plus, s'il en est ainsi, $X$ est intègre et géométriquement unibranche au point $x$.

Enfin, si $\sh{O}_{Y,y}$ est noethérien, on peut, dans le critère précédent, remplacer la condition (ii) par la condition

\emph{(ii bis)} On a $\dim\sh{O}_{Y,y}\leq\dim\sh{O}_{X,x}$.
\end{theorem}

Le fait que, si $f$ est étale, $X$ soit intègre et géométriquement unibranche au point $x$ est un cas particulier de (17.5.7). Il est clair d'autre part que les conditions (i) et (ii) sont nécessaires pour que $f$ soit étale au point $x$, puisque $\sh{O}_{X,x}$ est alors un $\sh{O}_{Y,y}$-module fidèlement plat (17.6.1). Prouvons que ces conditions sont suffisantes.

Posons $A=\sh{O}_{Y,y}$, $A'={}^{hs}\!A$ (18.8.7), $Y'=\Spec(A')$, $X'=X\times_Y Y'$, $f'=f_{(Y')}:X'\to Y'$. Soient $y'$ le point fermé de $Y'$, $x'$ un point de $X'$ au-dessus de $x$ et de $y'$; puisque le morphisme $g:Y'\to Y$ est plat (18.8.8), tout revient à voir que $f'$ est étale au point $x'$ (17.7.1, (iii)). L'hypothèse (ii) entraîne que le morphisme $\Spec(\sh{O}_{X,x})\to\Spec(\sh{O}_{Y,y})$ est dominant (I, 1.2.7); comme $\sh{O}_{Y,y}$ est intègre, $\Spec(\sh{O}_{Y,y})$ a un seul point générique $y_1$, et il y a donc une générisation $x_1$ de $x$ dans $X$ au-dessus de $y_1$. Le morphisme projection $X'\to X$ étant plat, il existe une générisation $x'_1$ de $x'$ au-dessus de $x_1$ (2.3.4); posons $y'_1=f'(x'_1)$; on a par suite $g(y'_1)=y_1$. Or, l'hypothèse sur $y$ entraîne que $A'$ est intègre (18.8.15), donc $Y'$ a un seul point générique $\eta$, nécessairement au-dessus de $y_1$ puisque $g$

\oldpage[IV-4]{158}

est plat (2.3.4); d'autre part, $\eta$ est aussi point générique de $g^{-1}(y_1)$ ($0_{\mathrm I}$, 2.1.8), et on sait par ailleurs que les fibres de $g$ sont des espaces discrets (18.8.12), donc $g^{-1}(y_1)=\{\eta\}$, et par suite $y'_1=\eta$. Le morphisme $\Spec(\sh{O}_{X',x'})\to\Spec(\sh{O}_{Y',y'})$ restriction de $f'$ est donc dominant, et puisque $\sh{O}_{Y',y'}=A'$ est intègre, l'homomorphisme correspondant $\sh{O}_{Y',y'}\to\sh{O}_{X',x'}$ est injectif (I, 1.2.7). D'autre part, $f'$ est non ramifié au point $x'$, (donc (18.8.3, $d$)), il y a un voisinage $U$ de $x'$ dans $X'$ tel que $f'|U$ soit une immersion fermée de $U$ dans $Y'$; \emph{a fortiori} l'homomorphisme $\sh{O}_{Y',y'}\to\sh{O}_{X',x'}$ est surjectif, et puisqu'il est injectif, il est bijectif. Mais alors $\sh{O}_{X',x'}$ est un $\sh{O}_{Y',y'}$-module plat, et $f'$ est donc étale au point $x'$ (17.6.1).

Enfin, lorsque $\sh{O}_{Y,y}$ est noethérien, on sait que si $f$ est étale au point $x$, on a $\dim\sh{O}_{Y,y}=\dim\sh{O}_{X,x}$ (17.6.4). Réciproquement, supposons vérifiées les conditions (i) et (ii \emph{bis}), et montrons qu'elles entraînent (ii). Soit $\mathfrak J$ le noyau de l'homomorphisme canonique $\sh{O}_{Y,y}\to\sh{O}_{X,x}$; restreignant au besoin $Y$ à un voisinage de $y$, on peut supposer que $\mathfrak J=\mathscr{J}_y$, où $\mathscr{J}$ est un Idéal de $\sh{O}_Y$; si $Y_1$ est le sous-préschéma fermé de $Y$ défini par $\mathscr{J}$, on peut, en restreignant encore au besoin $X$ et $Y$ à des voisinages de $x$ et $y$ respectivement, supposer que $f$ se factorise en $X\xrightarrow{f_1}Y_1\to Y$ (I, 6.5.1). Il est clair que $f_1$ est encore non ramifié au point $x$, donc (17.4.1) quasi-fini en ce point; la démonstration de (5.4.1, (i)) prouve alors que l'on a $\dim(\sh{O}_{X,x})\leq\dim(\sh{O}_{Y,y}/\mathfrak J)$. Mais si l'on avait $\mathfrak J\neq0$, on en conclurait que $\dim(\sh{O}_{Y,y}/\mathfrak J)<\dim(\sh{O}_{Y,y})$ puisque $Y$ est intègre ($0$, 16.1.2.2); on aurait par suite $\dim\sh{O}_{X,x}<\dim\sh{O}_{Y,y}$ contrairement à l'hypothèse, ce qui prouve (ii).

\begin{remark}[18.10.2]
\label{IV.18.10.2}
\begin{enumerate}
  \item[(i)] Lorsque $A=\sh{O}_{Y,y}$ est noethérien, et que l'on sait que son complété $\widehat A$ est intègre (par exemple si $A$ est régulier), on peut, dans la démonstration précédente, remplacer $A'$ par $\widehat A$.
  \item[(ii)] Lorsque $Y$ est localement noethérien, on peut donner de (18.10.1) une démonstration plus rapide, n'utilisant pas la hensélisation stricte, mais faisant intervenir les résultats assez délicats des \textsection\textsection~14 et 15. Comme $f$ est quasi-fini au point $x$ par hypothèse (17.4.1), on peut, en remplaçant $X$ par un voisinage ouvert de $x$, supposer que $f^{-1}(y)=\{x\}$. L'hypothèse (ii) entraîne, comme on l'a vu, l'existence d'une composante irréductible $Z$ de $X$ contenant $x$ et dominant l'unique composante irréductible $Y_0$ de $Y$ contenant $y$. Le morphisme $f|Z$ étant quasi-fini au point $x$, donc équidimensionnel en ce point (13.2.2), il résulte de l'hypothèse sur $Y$ et du critère de Chevalley (14.4.4) que $g=f|Z$ est universellement ouvert au point $x$. Par ailleurs puisque $f$ (donc aussi $g$) est non ramifié au point $x$, $g^{-1}(y)$ est géométriquement réduit sur $k(y)$ (17.4.1); comme $\sh{O}_{Y,y}$ est intègre, on conclut de (15.2.3) que $g$ est plat au point $x$, donc $f$ est étale en ce point (18.4.9).
  \item[(iii)] Les notations et les hypothèses sur $Y$ étant celles de (18.10.1), supposons seulement que $f$ soit localement de type fini (et non nécessairement localement de présentation finie). Alors, pour que $\sh{O}_{X,x}$ soit une $\sh{O}_{Y,y}$-algèbre essentiellement étale (18.6.1), il faut et il suffit que $f$ vérifie les deux conditions suivantes:

\oldpage[IV-4]{159}

\begin{enumerate}
  \item[$1^\circ$] $f$ est formellement non ramifié au point $x$;
  \item[$2^\circ$] l'homomorphisme $\sh{O}_{Y,y}\to\sh{O}_{X,x}$ est injectif.
\end{enumerate}

Il n'y a encore à prouver que la suffisance de ces conditions. Tenant compte de (18.4.12), il suffit de montrer que $f$ est plat au point $x$. Or, en reprenant la démonstration de (18.10.1) et les notations utilisées dans cette démonstration, il suffit de montrer que $f'$ est plat au point $x'$ (2.5.1); par ailleurs $f'$ est formellement non ramifié au point $x'$ (17.1.3); on peut donc se borner à faire la démonstration lorsque $Y=\Spec(A)$, où $A$ est strictement local et $y=f(x)$ le point fermé de $Y$. Utilisant alors (18.8.3) on voit que l'homomorphisme $\sh{O}_{Y,y}\to\sh{O}_{X,x}$ est surjectif, donc bijectif en vertu de l'hypothèse, et cela suffit à montrer la platitude de $f$ au point $x$.

Supposons de plus que $Y$ soit localement intègre au point $y$ (I, 2.1.8). Alors les conditions $1^\circ$ et $2^\circ$ précédentes (jointe au fait que $Y$ est géométriquement unibranche au point $y$ et $f$ localement de type fini) entraînent déjà que $f$ est étale au point $x$. En effet, cela résulte de ce qui précède et de (18.4.13).
\end{enumerate}
\end{remark}

\begin{corollary}[18.10.3]
\label{IV.18.10.3}
Soient $Y$ un préschéma intègre et géométriquement unibranche, $\eta$ son point générique, $X$ un préschéma connexe, $f:X\to Y$ un morphisme localement de type fini et tel que $f^{-1}(\eta)$ soit non vide. Alors, si $f$ est formellement non ramifié, $f$ est étale, et $X$ est intègre et géométriquement unibranche.
\end{corollary}

Il résulte de l'hypothèse et de (18.4.13) que $f$ est étale en tous les points où il est plat, et en particulier aux points de $f^{-1}(\eta)$, puisque $\sh{O}_{Y,\eta}=k(\eta)$ est un corps. L'ensemble ouvert $U$ des points de $X$ où $f$ est étale est donc non vide. Montrons d'autre part que pour tout $x\in\overline U$, l'homomorphisme $\sh{O}_{Y,f(x)}\to\sh{O}_{X,x}$ est injectif. Soit en effet $s$ une section de $\sh{O}_Y$ au-dessus d'un voisinage ouvert de $f(x)$, qu'on peut toujours supposer être $Y$ lui-même (la question étant locale sur $X$ et $Y$), et supposons que son image $t\in\Gamma(Y,f_*(\sh{O}_X))=\Gamma(X,\sh{O}_X)$ ait un germe nul au point $x$, donc s'annule dans un voisinage ouvert $V$ de $x$; alors la restriction de $t$ à l'ouvert $U\cap V\neq\varnothing$ est nulle; mais la restriction $f|(U\cap V)$ est étale, donc $W=f(U\cap V)$ est ouvert dans $Y$ et contient par suite le point générique $\eta$; l'homomorphisme $\sh{O}_{Y,\eta}\to\sh{O}_{X,z}$ étant injectif pour tout point $z\in f^{-1}(\eta)$, l'hypothèse que $t|(U\cap V)=0$ entraîne que $s_\eta=0$, et comme $Y$ est intègre, cela entraîne $s=0$, d'où notre assertion. Appliquant maintenant (18.10.2, (iii)), on conclut que $f$ est plat au point $x$, autrement dit $x\in U$, et $U$ est donc à la fois ouvert et fermé dans $X$; comme il est non vide et que $X$ est connexe, on a $U=X$.

Notons maintenant que puisque $f$ est étale, donc localement quasi-fini et que $Y$ est irréductible, les points maximaux de $X$ appartiennent à $f^{-1}(\eta)$ (2.3.4) et pour tout $x\in X$, il y a un voisinage de $x$ qui ne contient qu'un nombre fini de ces points maximaux; autrement dit l'ensemble des composantes irréductibles de $X$ est localement fini. Mais par (18.10.1) $X$ est intègre et géométriquement unibranche en tout point, donc tout point de $X$ n'appartient qu'à une seule composante irréductible, et puisque ces dernières forment un ensemble localement fini, elles sont ouvertes (et fermées) dans $X$; puisque $X$ est connexe, il est intègre.

\oldpage[IV-4]{160}

\begin{remark}[18.10.3.1]
\label{IV.18.10.3.1}
Si $f:X\to Y$ est localement de présentation finie, dominant et quasi-compact, la fibre $f^{-1}(\eta)$ est non vide en vertu de (I, 1.1.5). Par contre, si l'on ne suppose pas $f$ quasi-compact, $f$ peut être non ramifié et dominant sans que $f$ soit étale (en supposant toujours vérifiées les hypothèses de (18.10.3) pour $X$ et $Y$). C'est ce que montre l'exemple suivant: on prend pour $Y$ le plan affine $\Spec(\mathbf C[T,U])$; on considère d'autre part une famille $(X_j)_{j\in\mathbf Z}$ de préschémas isomorphes à la droite affine $\Spec(\mathbf C[T])$, et on forme le préschéma obtenu en « recollant » $X_{2j}$ et $X_{2j+1}$ au point $-1$ et $X_{2j+1}$ et $X_{2j+2}$ au point $+1$. Définissons d'autre part $f:X\to Y$ comme étant égal sur $X_{2j}$ à l'immersion fermée dont l'image est la droite d'équation $y=2j$, transformant le point $-1$ en $(2j-1,2j)$ et le point $+1$ en $(2j+1,2j)$; sur $X_{2j+1}$, $f_j$ est l'immersion fermée dont l'image est la droite d'équation $x=2j+1$, transformant le point $-1$ en $(2j+1,2j)$ et le point $-1$ en $(2j+1,2j+2)$. Il est clair que $f$ est une immersion locale (donc est non ramifié) et est dominant mais la fibre du point générique de $Y$ est vide et $f$ n'est pas étale.
\end{remark}

\begin{corollary}[18.10.4]
\label{IV.18.10.4}
Soient $g:Y\to S$, $h:X\to S$ deux morphismes localement de présentation finie, $f:X\to Y$ un $S$-morphisme, $x$ un point de $X$, $y=f(x)$, $s=g(y)=h(x)$. On suppose que $h$ est plat au point $x$ et $g^{-1}(s)$ normal au point $y$. Alors, pour que $f$ soit étale au point $x$, il faut et il suffit que $f$ soit non ramifié au point $x$ et que $\dim_xh^{-1}(s)=\dim_yg^{-1}(s)$.
\end{corollary}

Les conditions sont nécessaires en vertu du fait que si $f$ est étale au point $x$, il en est de même de $f_s:h^{-1}(s)\to g^{-1}(s)$, et la conclusion résulte alors de ce que $f_s$ est plat et quasi-fini au point $x$ (6.1.2). Inversement, si les conditions de l'énoncé sont vérifiées, pour voir que $f$ est étale au point $x$, il suffit, en vertu de (17.8.3) de voir que $f_s$ l'est. On peut donc se borner au cas où $S$ est le spectre d'un corps. Il résulte alors de l'hypothèse $\dim_xX=\dim_yY$ et du fait que $f$ est non ramifié au point $x$ que l'on a (par (5.2.3) et (17.4.1)) $\dim\sh{O}_{X,x}=\dim\sh{O}_{Y,y}$; puisque $f$ est quasi-fini au point $x$ et $\sh{O}_{Y,y}$ intègre, on déduit de ($0_{\mathrm I}$, 7.4.4) et ($0$, 16.3.10) que l'homomorphisme $\sh{O}_{Y,y}\to\sh{O}_{X,x}$ est injectif; il suffit alors pour conclure d'appliquer (18.10.1).

\begin{corollary}[18.10.5]
\label{IV.18.10.5}
Avec les notations de (18.10.4), on suppose que $h$ est plat et que $g^{-1}(s)$ est normal pour tout $s\in S$ (ce qui sera par exemple le cas si $g$ est lisse). Pour que $f$ soit une immersion ouverte, il faut et il suffit que $f$ soit un monomorphisme et que, pour tout $x\in X$, on ait $\dim_xh^{-1}(s)=\dim_yg^{-1}(s)$, où $y=f(x)$ et $s=g(y)=h(x)$.
\end{corollary}

On n'a à prouver que la suffisance des conditions; or, il résulte de (17.1.3) qu'un monomorphisme localement de présentation finie est non ramifié, donc les hypothèses entraînent que $f$ est étale (18.10.4); on conclut alors par (17.9.1).

\begin{lemma}[18.10.6]
\label{IV.18.10.6}
Soit $f:X\to Y$ un morphisme plat et localement quasi-fini ($\mathrm{Err}_{\mathrm{III}}$, 20).
\begin{enumerate}
  \item[(i)] L'ensemble $\operatorname{Max}(X)$ des points maximaux de $X$ est en correspondance biunivoque canonique avec l'ensemble $\coprod_{y\in\operatorname{Max}(Y)}f^{-1}(y)$.
  \item[(ii)] Si $Y$ est irréductible et de point générique $\eta$, l'ensemble des composantes irréductibles de $X$ est en correspondance biunivoque canonique avec $f^{-1}(\eta)$, et en particulier est fini si et seulement si $f^{-1}(\eta)$ est fini.

\oldpage[IV-4]{161}

  \item[(iii)] Si l'ensemble des composantes irréductibles de $Y$ est localement fini, il en est de même de l'ensemble des composantes irréductibles de $X$, et en particulier $X$ est localement connexe.
\end{enumerate}
\end{lemma}

L'assertion (i) résulte aussitôt de (2.3.4) et ($0_{\mathrm I}$, 2.1.6) et du fait que les fibres $f^{-1}(y)$ sont des ensembles discrets. Il est clair que (ii) est un cas particulier de (i). Enfin, pour prouver (iii), on peut, en vertu de (i), se borner au cas où $f$ est quasi-fini et où $Y$ est irréductible ($0_{\mathrm I}$, 2.1.6) et la conclusion résulte alors de (ii) et du fait que les fibres de $f$ sont des ensembles finis (II, 6.2.2).

\begin{proposition}[18.10.7]
\label{IV.18.10.7}
Soient $Y$ un préschéma géométriquement unibranche et irréductible (resp. intègre), $f:X\to Y$ un morphisme étale. Alors $X$ est isomorphe à une somme de préschémas irréductibles (resp. intègres) $X_\lambda$, où $\lambda$ parcourt la fibre $f^{-1}(\eta)$ du point générique $\eta$ de $Y$. Les $X_\lambda$ sont géométriquement unibranches; si $Y$ est normal, les $X_\lambda$ sont normaux.
\end{proposition}

En effet, $X$ est géométriquement unibranche (17.5.7), donc tout point de $X$ n'appartient qu'à une seule composante irréductible; en outre, en vertu de (18.10.6) et (17.6.1) l'ensemble des composantes irréductibles de $X$ est localement fini, donc ($0_{\mathrm I}$, 2.1.6) ce sont les composantes connexes $X_\lambda$ de $X$ ($\lambda\in f^{-1}(\eta)$), et $X$ est somme des $X_\lambda$; en vertu de (17.5.7), si $Y$ est intègre (resp. normal), les $X_\lambda$ sont intègres (resp. normaux).

\begin{proposition}[18.10.8]
\label{IV.18.10.8}
Soient $Y$ un préschéma normal et intègre, de point générique $\eta$, $K=k(\eta)$ son corps de fonctions rationnelles, $f:X\to Y$ un morphisme étale. On suppose que $f^{-1}(\eta)$ est finie et non vide, de sorte que le $K$-schéma $f^{-1}(\eta)$ est le spectre d'une $K$-algèbre finie séparable $L$, composée directe d'un nombre fini de corps $K_i$ ($1\leq i\leq n$), extensions finies séparables de $K$ (17.6.2). Soit $Y'$ la fermeture intégrale de $Y$ dans $L$, isomorphe à la somme des fermetures intégrales $Y_i$ de $Y$ dans les $K_i$ (II, 6.3.6). Alors le morphisme $f$ se factorise d'une seule manière en $f:X\xrightarrow{f'}Y'\xrightarrow{g}Y$, tel que la restriction $f'_\eta:f^{-1}(\eta)\to g^{-1}(\eta)$ soit l'isomorphisme canonique. En outre $f'$ est un isomorphisme local; pour que $f'$ soit une immersion ouverte, il faut et il suffit que $f$ soit séparé.
\end{proposition}

En vertu de (18.10.7), on peut se borner au cas où $X$ est intègre et normal, $L$ son corps de fonctions rationnelles, $f$ étant dominant. L'existence et l'unicité de la factorisation considérée de $f$ résultent alors de (II, 6.3.9). Comme $f'$ est birationnel et localement quasi-fini, les dernières assertions résultent de (8.12.10).

\begin{corollary}[18.10.9]
\label{IV.18.10.9}
Sous les hypothèses de (18.10.8), pour que $f$ soit un morphisme fini (autrement dit, pour que $X$ soit un revêtement étale de $Y$), il faut et il suffit que $X$ soit isomorphe à la fermeture intégrale $Y'$ de $Y$ dans $L$.
\end{corollary}

Si $f$ est fini, l'injection canonique $j:X\to Y'$ (qui est une immersion ouverte) est aussi un morphisme fini (II, 6.1.5, (v)), donc un morphisme fermé (II, 6.1.10), et comme $j(X)$ est dense dans $Y'$ par (18.10.8), on a $j(X)=Y'$. Inversement, comme $Y$ est normal et $L$ composé direct d'extensions finies séparables de $K$, $Y'$ est fini sur $Y$ (Bourbaki, \emph{Alg. comm.}, chap.~V, \textsection~1, n\textsuperscript{o}~6, cor.~1 de la prop.~18), d'où le corollaire.

\begin{env}[18.10.10]
\label{IV.18.10.10}
Soient $Y$ un préschéma normal et intègre, $K=R(Y)$ son corps de fonctions rationnelles. Nous dirons qu'une $K$-algèbre de rang fini $L$ est non ramifiée sur $Y$ si: $1^\circ$ $L$ est une $K$-algèbre séparable, donc composée directe d'extensions finies séparables $K_i$ de $K$ ($1\leq i\leq n$); $2^\circ$ la fermeture intégrale $Y'$ de $Y$ dans $L$ (somme des préschémas fermetures intégrales de $Y$ dans les $K_i$) est non ramifiée sur $Y$ (ce qui, par (18.10.3),

\oldpage[IV-4]{162}

équivaut à dire que $Y'$ est étale sur $Y$). Lorsque $Y=\Spec(A)$, où $A$ est un anneau intègre et intégralement clos, $K$ son corps des fractions, on dira encore (par abus de langage, et lorsqu'une confusion avec la terminologie de (17.3.2, (ii)) n'est pas à craindre) que $L$ est non ramifiée sur $A$ au lieu de « non ramifiée sur $Y$ ».
\end{env}

\begin{remark}[18.10.11]
\EGAReaderTarget{ega4.c04.rem.18-10-11}{IV.18.10.11}
Au lieu de dire que $L$ est non ramifiée sur $Y$ certains auteurs disent aussi que $L$ est « non ramifiée sur $K$ »; nous nous garderons de suivre cet usage qui prête à confusions.
\end{remark}

On peut alors exprimer (18.10.9) sous la forme suivante:

\begin{corollary}[18.10.12]
\label{IV.18.10.12}
Soient $Y$ un préschéma normal et intègre, $K=R(Y)$ son corps de fonctions rationnelles. Alors le foncteur $X\rightsquigarrow R(X)$ est une équivalence de la catégorie des revêtements étales de $Y$, et de la catégorie des $K$-algèbres finies étales qui sont non ramifiées sur $Y$. On obtient un foncteur quasi-inverse en faisant correspondre à toute $K$-algèbre finie étale $L$, non ramifiée sur $Y$, la fermeture intégrale $Y'$ de $Y$ dans $L$.
\end{corollary}

\begin{proposition}[18.10.13]
\label{IV.18.10.13}
Soient $Y$ un préschéma normal et intègre, $K=R(Y)$ son corps de fonctions rationnelles.
\begin{enumerate}
  \item[(i)] Le corps $K$ est une $K$-algèbre non ramifiée sur $Y$.
  \item[(ii)] Soit $L$ une extension finie de $K$, non ramifiée sur $Y$, et soit $Z$ la fermeture intégrale de $Y$ dans $L$. Alors, si $M$ est une $L$-algèbre non ramifiée sur $Z$, $M$ est aussi une $K$-algèbre non ramifiée sur $Y$ (« transitivité » de la non-ramification).
  \item[(iii)] Soient $Y'$ un préschéma normal et intègre, $K'$ son corps de fonctions rationnelles, $g:Y'\to Y$ un morphisme dominant. Si $L$ est une $K$-algèbre non ramifiée sur $Y$, alors $L\otimes_K K'$ est une $K'$-algèbre non ramifiée sur $Y'$ (propriété de « translation »). Si de plus $Y=\Spec(A)$ et $Y'=\Spec(A')$ sont affines, et si $C$ est la fermeture intégrale de $A$ dans $L$, alors $C'=C\otimes_A A'$ est la fermeture intégrale de $A'$ dans $L\otimes_K K'$.
\end{enumerate}
\end{proposition}

L'assertion (i) est immédiate, $Y$ étant par hypothèse son propre normalisé dans $K$. Pour prouver (ii), soit $Z'$ la fermeture intégrale de $Z$ dans $M$, qui est par hypothèse un revêtement étale de $Z$. Comme $Z$ est étale et séparé sur $Y$, il en est de même de $Z'$ (17.3.3) qui est évidemment la fermeture intégrale de $Y$ dans $M$. Comme $M$ est une $K$-algèbre finie et séparable, il résulte de (18.10.10) que $M$ est non ramifiée sur $Y$. Enfin, pour démontrer (iii), notons que si $Z$ est la fermeture intégrale de $Y$ dans $L$, $Z\times_Y Y'$ est étale et séparé sur $Y'$ (17.3.3), fini sur $Y'$ puisque $Z$ est fini sur $Y$, et admet pour anneau de fonctions rationnelles $L\otimes_K K'$ (I, 3.4.9); comme $Y'$ est normal, il en est de même de $Z\times_Y Y'$ (17.5.7), qui est donc la fermeture intégrale de $Y'$ dans $L\otimes_K K'$, ce qui achève la démonstration.

\begin{corollary}[18.10.14]
\label{IV.18.10.14}
Soient $Y$ et $Y'$ deux préschémas normaux et intègres, $K$, $K'$ leurs corps de fonctions rationnelles respectifs, $g:Y'\to Y$ un morphisme dominant, de sorte que $K'$ est une extension de $K$. Soient $L$ une extension finie de $K$ non ramifiée sur $Y$, $L_1$ une extension composée (Bourbaki, \emph{Alg.}, chap.~VIII, \textsection~8, déf.~1) de $L$ et de $K'$. Alors $L_1$ est une extension de $K'$ non ramifiée sur $Y'$. Si de plus $Y=\Spec(A)$, $Y'=\Spec(A')$ sont affines, et si $C$ est la fermeture intégrale de $A$ dans $L$, alors le sous-anneau $C_1=A[C,A']$ de $L_1$ est la fermeture intégrale de $A'$ dans $L_1$.
\end{corollary}

En effet, $L\otimes_K K'$ est composée directe d'extensions finies séparables de $K'$, et $L_1$

\oldpage[IV-4]{163}

est l'une de ces extensions; donc la fermeture intégrale de $Y$ dans $L$ est somme de préschémas dont l'un est la fermeture intégrale de $Y$ dans $L_1$. Le corollaire résulte donc aussitôt de (18.10.14).

Lorsque par exemple $A$ est l'anneau des entiers d'un corps de nombres algébriques $K$, $K'$ et $L$ des extensions algébriques de $K$, il y a des exemples classiques montrant que la relation $C_1=A[C,A']$ est en défaut lorsque $L$ n'est pas non ramifiée sur $A$ [\ ].

\begin{env}[18.10.15]
\label{IV.18.10.15}
Supposons que $Y=\Spec(A)$ soit affine, $A$ étant intègre et intégralement clos; soient $K$ son corps des fractions, $L$ une extension finie séparable de $K$, et supposons que la fermeture intégrale $C$ de $A$ dans $L$ soit un $A$-module projectif de type fini (ce qui sera par exemple le cas si $A$ est un anneau de Dedekind (Bourbaki, \emph{Alg. comm.}, chap.~VII, \textsection~4, n\textsuperscript{o}~10, prop.~22)). Dire que $L$ est non ramifiée sur $A$ signifie que $C$ est étale (18.10.10) et en vertu de (18.2.7, (ii)), cela équivaut donc à dire que le discriminant $d_{C/A}$ de $\Spec(C)$ sur $\Spec(A)$ est inversible dans $A$. Dans le cas particulier où $C$ est un $A$-module libre et $(x_i)_{1\leq i\leq n}$ une base de $C$ sur $A$, cela signifie que $\det(\operatorname{Tr}_{C/A}(x_ix_j))$ est inversible dans $A$.
\end{env}

\begin{theorem}[18.10.16]
\label{IV.18.10.16}
Soient $Y$ un préschéma intègre, $\eta$ son point générique, $f:X\to Y$ un morphisme séparé et localement de type fini. On suppose que toute composante irréductible de $X$ domine $Y$ et que la fibre générique $X_\eta=f^{-1}(\eta)$ est un préschéma fini sur $K=k(\eta)$, donc $(X_\eta)_{\mathrm{red}}$ est égal à $\Spec(L)$, où $L=\prod_iL_i$ est un produit d'extensions finies $L_i$ de $K$. On pose
\[
  n=[L:K]=\sum_i[L_i:K],\qquad n_s=\sum_i[L_i:K]_s \quad\text{(somme des degrés séparables des $L_i$)}.
\]
Soit $y$ un point géométriquement unibranche de $Y$, et désignons par $n(y)$ la somme des degrés séparables sur $k(y)$ des corps résiduels des points isolés de $f^{-1}(y)$. Alors:
\begin{enumerate}
  \item[(i)] On a $n(y)\leq n_s\leq n$.
  \item[(ii)] Supposons $X$ réduit et le point $y$ normal dans $Y$. Pour que $n(y)=n$, il faut et il suffit qu'il existe un voisinage ouvert $U$ de $y$ dans $Y$ tel que la restriction $f^{-1}(U)\to U$ de $f$ soit un morphisme étale et fini.
\end{enumerate}
\end{theorem}

\emph{A) Réduction au cas où $f$ est quasi-fini.} --- On sait (13.1.4) que l'ensemble $Z$ des points $x\in X$ isolés dans $f^{-1}(f(x))$ est ouvert dans $X$ et par hypothèse $Z$ contient $f^{-1}(\eta)$. Puisque $f^{-1}(\eta)$ est finie, $Z$ est réunion filtrante croissante d'ouverts quasi-compacts $V_\lambda$ contenant $f^{-1}(\eta)$ (donc denses dans $X$ par hypothèse). Si $n_\lambda(y)$ est la somme des degrés séparables sur $k(y)$ des corps résiduels des points de $V_\lambda\cap f^{-1}(y)$, on a $n(y)=\sup_\lambda n_\lambda(y)$; pour prouver (i), il suffira donc de montrer que $n_\lambda(y)\leq n_s$. D'autre part, si $n(y)=n$, il y a un indice $\lambda$ tel que $n_\lambda(y)=n$. Supposons démontré qu'il existe alors un voisinage ouvert $U$ de $y$ dans $Y$ tel que la restriction $V_\lambda\cap f^{-1}(U)\to U$ de $f$ soit un morphisme étale et fini. Puisque $f$ est séparé, l'injection canonique $V_\lambda\cap f^{-1}(U)\to f^{-1}(U)$ est alors aussi un morphisme fini (II, 6.1.5, (v)), donc $V_\lambda\cap f^{-1}(U)$ est fermé dans $f^{-1}(U)$ (II, 6.1.10); mais étant partout dense dans $f^{-1}(U)$, il lui est nécessairement égal.

\emph{B) Réduction au cas où $X$ est réduit, $f$ fini et $Y=\Spec(\sh{O}_{Y,y})$.} --- On notera que l'hypothèse sur $f$ entraîne qu'elle est aussi vérifiée pour $f_{\mathrm{red}}$ et que la conclusion de (i) ne concerne que $f_{\mathrm{red}}$. On peut donc remplacer partout $f$ par $f_{\mathrm{red}}$ et supposer donc $X$ réduit.

\oldpage[IV-4]{164}

D'autre part, en vertu de (18.12.13)\footnote{Le lecteur vérifiera que (18.10.16) n'est pas utilisé dans la démonstration de (18.12.13). Si $f$ est supposé localement de présentation finie (resp. quasi-projectif), on peut remplacer (18.12.13) par (8.12.8) (resp. (8.12.6)).}, on peut (après avoir remplacé $Y$ par un voisinage ouvert affine $V$ de $y$ et $f$ par sa restriction $f^{-1}(V)\to V$) factoriser $f$ en $X\xrightarrow{j}X'\xrightarrow{g}Y$, où $g$ est fini et $j$ une immersion ouverte. Puisque $X$ est réduit, on peut remplacer $X'$ par le sous-préschéma fermé réduit d'espace sous-jacent $\overline{j(X)}$ (I, 5.2.2), donc supposer $X'$ réduit et tel que toute composante irréductible de $X'$ domine $Y$. En outre, $j(X)$ est un ouvert dense dans $X'$ et n'a qu'un nombre fini de composantes irréductibles en vertu de l'hypothèse; donc ($0_{\mathrm I}$, 1.2.7) les composantes irréductibles de $X'$ sont les adhérences de celles de $j(X)$, et par suite la fibre $g^{-1}(\eta)$ s'identifie à $f^{-1}(\eta)$. Si on suppose le théorème démontré lorsqu'en outre $f$ est un morphisme fini, on pourra l'appliquer à $g$, et comme $f^{-1}(y)$ s'identifie à un sous-préschéma de $g^{-1}(y)$, la relation (i) pour $g$ entraîne la même relation pour $f$. Supposons d'autre part que $n(y)=n$; alors, en vertu de ce qui précède, on peut appliquer la conclusion de (ii) au morphisme $g$, et en remplaçant au besoin $Y$ par un voisinage ouvert de $y$, on peut supposer $g$ étale, et il en est par suite de même de $f$; en outre, on a alors $f^{-1}(y)=g^{-1}(y)$. Comme le morphisme $g$ est fermé et que $j(X)$ est ouvert dans $X'$, il existe un voisinage $U$ de $y$ dans $Y$ tel que $j(X)\cap g^{-1}(U)=j(f^{-1}(U))$, ce qui prouve que (ii) est vrai aussi pour $f$.

On est ainsi ramené au cas où $f$ est en outre un morphisme fini. Il est immédiat que, pour prouver (i), on peut se borner au cas où $Y=\Spec(\sh{O}_{Y,y})$. Montrons qu'il en est de même pour prouver (ii). En effet, supposons prouvé que la condition $n(y)=n$ entraîne que $X\times_Y\Spec(\sh{O}_{Y,y})$ est étale et fini sur $\Spec(\sh{O}_{Y,y})$; comme cela entraîne que $f$ est plat et formellement non ramifié aux points de $f^{-1}(y)$, et que $Y$ est intègre, on conclut de (18.4.13) que $f$ est étale aux points de $f^{-1}(y)$, donc aussi dans un voisinage ouvert $V$ de $f^{-1}(y)$ dans $X$; mais puisque $f$ est un morphisme fini, donc fermé, $V$ contient un ouvert de la forme $f^{-1}(U)$, où $U$ est un voisinage ouvert de $y$ dans $Y$.

\emph{C) Réduction au cas où $Y=\Spec(\sh{O}_{Y,y})$ et $\sh{O}_{Y,y}$ est strictement local.} --- Posons $A=\sh{O}_{Y,y}$, et soit $A'={}^{hs}\!A$ le hensélisé strict de $A$ (18.8.7), qui est un anneau local intègre et géométriquement unibranche (18.8.16). Soit $Y'=\Spec(A')$, et soient $y'$ le point fermé et $\eta'$ le point générique de $Y'$; $y'$ est au-dessus de $y$ et puisque le morphisme $g:Y'\to Y$ est plat (18.8.8), $\eta'$ est au-dessus de $\eta$ (2.3.4). Posons $X'=X\times_Y Y'$, $f'=f_{(Y')}:X'\to Y'$; $f'$ est fini et puisque $g$ est plat, toute composante irréductible de $X'$ domine $Y'$ (2.3.7).

Si $n'(y')$, $n'$ et $n'_s$ sont les nombres définis pour $f'$ de la même façon que $n(y)$, $n$ et $n_s$ pour $f$, on a $n'_s=n_s$ et $n'(y')=n(y)$ (I, 6.4.8), et il revient donc au même de prouver l'assertion (i) de l'énoncé pour $f$ ou pour $f'$. En outre, si $X$ est réduit et $y$ un point normal de $Y$ (donc $A$ intégralement clos), alors $A'$ est intégralement clos (18.8.12). D'autre part il résulte de la définition (18.8.7), de la remarque (8.1.2, $a$)) et du théorème de la double limite inductive, que $A'$ est limite inductive d'anneaux $A_\alpha$ tels que les morphismes $\Spec(A_\alpha)\to\Spec(A)=Y$ soient étales; par suite $X'$ est limite projective des $X_\alpha=X\otimes_A A_\alpha$, qui sont étales sur $X$, donc réduits puisque $X$ l'est (17.5.7); par passage à la limite,

\oldpage[IV-4]{165}

on en déduit que $X'$ est réduit (8.7.1). Cela étant, si $n(y)=n$, on a $n_s=n$ et par suite les corps résiduels aux points de $X_\eta$ sont nécessairement séparables sur $k(\eta)$; les corps résiduels aux points de $X'_{\eta'}$ sont par suite algébriques et séparables sur $k(\eta')$ (4.6.1), autrement dit $n'_s=n'=n$, et on a donc en vertu de ce qui précède $n'(y')=n'$. Le morphisme $f'$ a donc les mêmes propriétés que $f$, et si l'on prouve que $f'$ est étale, il en résultera que $f$ est étale par (17.7.1).

\emph{D) Fin de la démonstration.} --- Supposons donc $A$ strictement local. Puisque $A$ est hensélien et le morphisme $f$ fini, il résulte de (18.5.11) que si $x_j$ ($1\leq j\leq m$) sont les points distincts de $f^{-1}(y)$, $X$ est somme des sous-préschémas ouverts $\Spec(\sh{O}_{X,x_j})=X_j$, qui sont donc les composantes connexes de $X$. Par hypothèse, chacune d'elles domine $Y$, donc rencontre $f^{-1}(\eta)$, et par suite le nombre de points de $f^{-1}(\eta)$ est $\geq m$; \emph{a fortiori} on a $n_s\geq m$. Mais puisque $k(y)$ est séparablement clos et que les $k(x_j)$ sont algébriques sur $k(y)$, ce sont nécessairement des extensions radicielles de $k(y)$, donc $m=n(y)$, et cela prouve (i). Supposons maintenant vérifiées les hypothèses de (ii); les relations $n(y)=n_s=n$ entraînent d'abord, puisque toute composante irréductible de $X$ domine $Y$, que chacun des $X_j$ est irréductible; en outre, comme les $X_j$ sont réduits, ils sont intègres; enfin les relations $m=n_s=n$ entraînent que pour le point générique $\xi_j$ de $X_j$, $k(\xi_j)$ est une extension séparable de $k(\eta)$ de degré 1, donc isomorphe à $k(\eta)$; autrement dit, la restriction $f|X_j:X_j\to Y$ de $f$ est un morphisme fini et birationnel; en outre, $X_j$ est intègre et par hypothèse $Y$ est normal; donc (8.12.10.1) $f|X_j$ est un isomorphisme, ce qui prouve que $f$ est étale, et termine la démonstration de (18.10.16).

\begin{remark}[18.10.17]
\label{IV.18.10.17}
La méthode de «localisation étale» utilisée dans la démonstration de (18.10.16) permet aussi d'améliorer les résultats de (15.5.1) en éliminant l'hypothèse noethérienne. On a tout d'abord le résultat suivant qui généralise (15.5.2):
\end{remark}

\begin{env}[18.10.17.1]
\label{IV.18.10.17.1}
Soient $f:X\to Y$ un morphisme séparé et quasi-fini, $y$ un point de $Y$ tel que $f$ soit universellement ouvert aux points de $f^{-1}(y)$. Alors, pour toute générisation $y'$ de $y$, on a $n(y)\leq n(y')$.
\end{env}

On peut remplacer $Y$ par le sous-préschéma fermé réduit $Z$ ayant pour espace sous-jacent $\overline{\{y'\}}$ et $X$ par $f^{-1}(Z)$, donc supposer $Y$ intègre de point générique $y'$. Utilisant (18.12.13) comme dans (18.10.16, B)), on peut supposer que $X$ est ouvert dense dans un préschéma $X'$ et que $f$ est la restriction à $X$ d'un morphisme fini $g:X'\to Y$. Montrons que l'on a $g^{-1}(y')=f^{-1}(y')$. En effet, soit $i:g^{-1}(y')\to X'$ l'injection canonique, qui est un morphisme plat puisque $y'$ est point générique de $Y$ (I, 3.6.5); on peut écrire $f^{-1}(y')=i^{-1}(X)$. D'autre part, l'injection canonique $j:X\to X'$ est un morphisme de type fini puisque le composé $g\circ j=f$ est de type fini et que $g$ est séparé (1.5.4); donc $X$ est un ouvert rétrocompact dans $X'$, et par suite est pro-constructible dans $X'$ (1.9.5, (v)). On conclut donc de (2.3.10) que l'on a $i^{-1}(\overline{X})=\overline{i^{-1}(X)}$, autrement dit $f^{-1}(y')$ est dense dans $g^{-1}(y')$; mais comme $g^{-1}(y')$ est discret, on a nécessairement $f^{-1}(y')=g^{-1}(y')$.

On est ainsi ramené à prouver que la somme des nombres $[k(x):k(y)]_s$ où $x$ parcourt l'ensemble des points de $g^{-1}(y)$ où $g$ est universellement ouvert, est au plus égale

\oldpage[IV-4]{166}

au nombre $n(y')$ (pour le morphisme $g$). Utilisant ensuite le fait que la propriété d'être universellement ouvert en un point se conserve par changement de base, on montre, comme dans (18.10.16, C)), que l'on peut se ramener au cas où $Y=\Spec(\sh{O}_{Y,y})$, avec $\sh{O}_{Y,y}$ strictement local. Alors, avec les notations de (18.10.16, D)), $X$ est somme des sous-préschémas ouverts $X_j=\Spec(\sh{O}_{X,x_j})$, où $x_j$ parcourt $g^{-1}(y)$; l'hypothèse que $g$ est ouvert en un de ces points $x_j$ entraîne que le sous-préschéma $X_j$ correspondant domine $Y$ (I, 10.3); on termine comme dans (18.10.16, D)).

\emph{On tire ensuite de (18.10.17.1) que les conclusions de (15.5.1) sont encore valables lorsqu'on ne suppose plus $Y$ localement noethérien, mais seulement que $f$ est séparé, quasi-fini et de présentation finie.}

En effet, pour prouver l'assertion (i) de (15.5.1), on remarque que, puisque $f$ est de présentation finie, l'ensemble $E$ des points $y'\in Y$ tels que $n(y')\geq n(y)$ est localement constructible (9.7.9); pour montrer que $y$ est intérieur à $E$, il suffit, en vertu de (I, 10.1), de voir que toute générisation $y'$ de $y$ appartient à $E$, ce qui n'est autre que (18.10.17.1).

Pour prouver l'assertion (ii) de (15.5.1), on peut, puisque $f$ est de présentation finie, utiliser (8.10.5, (xii)) appliqué suivant la méthode de (8.1.2, a)), et on peut donc déjà supposer que $Y=\Spec(\sh{O}_{Y,y})$. Appliquant ensuite (2.7.1, (vii)), on voit qu'on peut par le changement de base fidèlement plat $Y'\to Y$, où $Y'=\Spec({}^{hs}\!\sh{O}_{Y,y})$, supposer que l'anneau $\sh{O}_{Y,y}$ est strictement local. Appliquant (18.5.11, c)), on voit alors que si $x_j$ ($1\leq j\leq n$) sont les points de $f^{-1}(y)$, $X$ est somme de $n$ sous-préschémas ouverts $X_j=\Spec(\sh{O}_{X,x_j})$ qui sont finis sur $Y$, et d'un préschéma ouvert $X''$. Si on prouve que $X''=\varnothing$, on aura démontré que $f$ est fini, donc propre. Or, les corps $k(x_j)$ étant des extensions algébriques de $k(y)$, sont radicielles, et par suite $n(y)=n$; comme $f$ est ouvert en chacun des points $x_j$, $X_j$ domine $Y$ (I, 10.3), donc la restriction de $f$ à $X_j$, étant un morphisme fini, est surjective (II, 6.1.10). L'hypothèse que $z\mapsto n(z)$ est constant dans $Y$ entraîne donc que $X''\cap f^{-1}(y')=\varnothing$ pour tout $y'\in Y$, c'est-à-dire $X''=\varnothing$.

Enfin, la preuve de (iii) reproduit sans changement celle donnée dans (15.5.1).

\begin{env}[18.10.17.2]
\label{IV.18.10.17.2}
Par des méthodes analogues utilisant la hensélisation stricte, on peut, dans la plupart des résultats des \textsection\textsection~14 et 15, se débarrasser des hypothèses noethériennes.
\end{env}

\begin{lemma}[18.10.18]
\label{IV.18.10.18}
Soient $X$, $Y$ deux préschémas, $f:X\to Y$ un morphisme birationnel\footnote{Nous entendons par là la notion définie dans (6.15.4), mais où l'on ne suppose pas $X$ et $Y$ réduits.}, $x$ un point de $X$. Pour que $f$ soit un isomorphisme local au point $x$, il faut et il suffit que $f$ soit étale au point $x$. Pour que $f$ soit une immersion ouverte, il faut et il suffit que $f$ soit étale et séparé.
\end{lemma}

Les conditions énoncées étant trivialement nécessaires, tout revient à voir qu'elles sont suffisantes. Pour la première assertion, la question est locale sur $X$ et $Y$, donc on peut supposer $f$ étale, $X$ et $Y$ affines, $f$ séparé, et on est ramené à prouver la seconde assertion. En vertu de (17.9.1), il s'agit de prouver que $f$ est radiciel. Soit donc $y\in Y$ et prouvons que $f^{-1}(y)$ est radiciel sur $k(y)$. Le changement de base $\Spec(\sh{O}_{Y,y})\to Y$ ne changeant pas le fait que $f$ est étale, séparé et birationnel, on peut se borner au cas où $Y=\Spec(A)$, avec $A=\sh{O}_{Y,y}$.

\oldpage[IV-4]{167}

Soit alors $A'={}^{hs}\!A$, qui est un anneau local et un $A$-module fidèlement plat (18.8.8), l'homomorphisme $A\to A'$ étant local. Posons $Y'=\Spec(A')$, $X'=X\times_Y Y'$, $f'=f_{(Y')}:X'\to Y'$, de sorte que $f'$ est étale et séparé; en outre, comme le morphisme $Y'\to Y$ est plat, le raisonnement de (6.15.4.1) montre que $f'$ est aussi birationnel. Si $y'$ est le point fermé de $Y'$, il suffira, en vertu de (2.6.1, (v)), de prouver que $f'^{-1}(y')$ est radiciel sur $k(y')$; puisque $f'$ est étale, il suffira de prouver que $f'^{-1}(y')$ contient au plus un point (17.6.1, c')).

Supposons donc $Y=\Spec(A)$, $A$ étant strictement local, $f$ séparé, étale et birationnel, et montrons que $f^{-1}(y)$ ne peut contenir plus d'un point. En effet, s'il existait dans $f^{-1}(y)$ deux points distincts $x_1$, $x_2$, il existerait deux $Y$-sections $u_1$, $u_2$ de $f$ telles que $u_1(y)=x_1$ et $u_2(y)=x_2$ (18.8.1); mais puisque $f$ est étale et séparé, $u_1$ et $u_2$ seraient des isomorphismes de $Y$ sur deux composantes connexes (ouvertes) de $X$ (17.9.4) et il y aurait donc deux points maximaux au moins de $X$ au-dessus de tout point maximal de $Y$, ce qui contredit l'hypothèse que $f$ est birationnel.

\begin{proposition}[18.10.19]
\label{IV.18.10.19}
Soient $Y$ un préschéma réduit dont l'ensemble des composantes irréductibles est localement fini, $f:X\to Y$ un morphisme séparé, localement de type fini et formellement non ramifié, $g$ une $Y$-section rationnelle de $f$ (c'est-à-dire une $Y$-application rationnelle de $Y$ dans $X$), $U$ le domaine de définition de $g$ (I, 7.2.1), et soit $Z$ l'adhérence de $g(U)$ dans $X$. Alors, pour tout point $y\in Y-U$ tel que $Y$ soit géométriquement unibranche au point $y$, on a $Z_y=Z\cap f^{-1}(y)=\varnothing$. En particulier, si $Y$ est géométriquement unibranche, $g(U)$ est une partie à la fois ouverte et fermée de $X$. Si, pour toute partie fermée irréductible $T$ de $X$, $f(T)$ est fermé dans $Y$, $g$ est définie en tout point géométriquement unibranche de $Y$.
\end{proposition}

Puisque $Y$ est réduit et $f$ séparé, il résulte de (I, 7.2.2) qu'il existe une $U$-section $u$ de $f^{-1}(U)$ appartenant à la classe $g$; d'ailleurs (17.4.1.2), puisque $f$ est formellement non ramifié, $u$ est un isomorphisme de $U$ sur le préschéma induit sur l'ouvert $u(U)$ de $X$. Si on note encore $Z$ le sous-préschéma réduit de $X$ ayant $Z$ pour espace sous-jacent, $Z|_U=u(U)$, étant réduit, est induit aussi par $Z$, donc la restriction $f'=f|_Z$ de $f$ est un morphisme birationnel de $Z$ dans $Y$, d'ailleurs formellement non ramifié, puisque $f$ l'est (17.1.3). La question étant locale sur $Y$, puisque $Y$ est réduit et géométriquement unibranche (donc intègre) au point $y$, on peut se borner (en remplaçant $Y$ par un voisinage ouvert de $y$) au cas où $Y$ est intègre; alors $u(U)$ est irréductible, donc aussi $Z$, et par suite $Z$ est aussi intègre. Alors (I, 8.2.7), puisque $f'$ est dominant, l'homomorphisme $\sh{O}_{Y,f(x)}\to\sh{O}_{Z,x}$ est injectif pour tout $x\in Z_y$, et on conclut de (18.10.2, (iii)) que $f'$ est un morphisme étale; étant séparé et birationnel, c'est un isomorphisme local au point $x$ en vertu de (18.10.18); mais cela entraînerait que la $U$-section $u$ serait prolongeable à un ouvert distinct de $U$, contrairement à la définition de ce dernier. On a donc nécessairement $Z_y=\varnothing$, ce qui établit la première assertion; la seconde en est une conséquence évidente. En outre, sous les hypothèses de la dernière assertion de l'énoncé, comme $Z$ est fermé et irréductible, $f(Z)$ est fermé dans $Y$, donc égal à $Y$, c'est-à-dire $Z_y\ne\varnothing$ pour tout $y\in Y$, ce qui achève la démonstration.

\begin{remark}[18.10.20]
\EGAReaderTarget{ega4.c04.rem.18-10-20}{IV.18.10.20}
Soit $Y$ un préschéma localement noethérien. On dit qu'un morphisme $f:X\to Y$ est essentiellement propre s'il est localement de type fini et si, pour tout $Y$-schéma de la forme $Y'=\Spec(A)$ où $A$ est un anneau de valuation discrète, l'application canonique (II, 7.3.2.2) relative au $Y'$-préschéma $X\times_Y Y'$ est bijective.

\oldpage[IV-4]{168}

Le raisonnement de (II, 7.3.8) prouve qu'il revient au même de dire que $f$ (supposé localement de type fini) est séparé et que, pour tout changement de base $Y''\to Y$, où $Y''$ est localement noethérien, l'image par $f_{(Y'')}:X_{(Y'')}\to Y''$ de toute partie fermée irréductible de $X''$ est fermée. Dire que $f$ est propre (pour $Y$ localement noethérien) signifie donc que $f$ est essentiellement propre et de type fini (II, 7.3.8); mais on rencontre des exemples importants de morphismes essentiellement propres qui ne sont pas de type fini (par exemple certains «préschémas de Picard» ou certains «préschémas de Néron-Severi» (chap. VI)). Il résulte évidemment de (18.10.19) que si $Y$ est localement noethérien, réduit et géométriquement unibranche, et si le morphisme $f:X\to Y$ est essentiellement propre et non ramifié, alors toute $Y$-section rationnelle de $f$ est partout définie.
\end{remark}

\begin{proposition}[18.10.21]
\label{IV.18.10.21}
Soient $f:X\to Y$ un morphisme localement de présentation finie, $x$ un point de $X$, $y=f(x)$ et supposons l'anneau $\sh{O}_{Y,y}$ intègre et géométriquement unibranche. Soit $r$ le nombre minimum de générateurs du $\sh{O}_{X,x}$-module $(\Omega^1_{X/Y})_x$ (égal au rang du $k(x)$-espace vectoriel $(\Omega^1_{X/Y})_x\otimes_{\sh{O}_{X,x}}k(x))$. Alors, pour que $f$ soit lisse au point $x$, il faut et il suffit que, si $y'$ est le point générique de l'unique composante irréductible de $Y$ contenant $y$, il existe une générisation $x'$ de $x$ telle que $f(x')=y'$ et $\dim_{x'}f^{-1}(y')\geq r$.
\end{proposition}

Si $f$ est lisse au point $x$, il l'est aussi dans un voisinage ouvert $U$ de $x$, donc en toute générisation $x'$ de $x$ et $\dim_{x'}(f^{-1}(f(x')))=r$ en chacun de ces points, en vertu de (17.10.2). Comme $f$ est en outre plat dans $U$, il existe une générisation de $x$ au-dessus de toute générisation de $y$ (2.3.4), ce qui prouve que la condition est nécessaire. Pour montrer qu'elle est suffisante, notons d'abord que par (17.5.1), il suffit de prouver que le morphisme déduit de $f$ par le changement de base $\Spec(\sh{O}_{Y,y})\to Y$ est lisse au point $x$; on peut donc supposer que $Y=\Spec(\sh{O}_{Y,y})$.

La question est locale sur $X$, donc ((16.4.22) et ($0_{\mathrm I}$, 5.2.2)) on peut supposer que $X$ est affine, et qu'il existe $r$ sections $t_i$ ($1\leq i\leq r$) de $\sh{O}_X$ au-dessus de $X$ telles que les sections $d_{X/Y}t_i$ engendrent $\Omega^1_{X/Y}$. Soit
\[
  g:X\to Z=Y[T_1,\ldots,T_r]=\mathbf V^r_Y
\]
le $Y$-morphisme correspondant à l'homomorphisme $\sh{O}_Y^r\to f_*(\sh{O}_X)$ de $\sh{O}_Y$-Modules défini par les $t_i$ (considérées comme sections de $f_*(\sh{O}_X)$ au-dessus de $Y$) (II, 1.2.7). On a la suite exacte (16.4.19)
\[
  g^*(\Omega^1_{Z/Y})\longrightarrow\Omega^1_{X/Y}\longrightarrow\Omega^1_{X/Z}\longrightarrow0.
\]
Comme les $dT_i$ ($1\leq i\leq r$) engendrent $\Omega^1_{Z/Y}$, il résulte de la définition de $g$ que l'homomorphisme $g^*(\Omega^1_{Z/Y})\to\Omega^1_{X/Y}$ est surjectif, donc $\Omega^1_{X/Z}=0$ et $g$ est par suite non ramifié (17.4.1). Nous allons voir que $g$ est étale au point $x$, et compte tenu de (17.3.8), cela prouvera que $f$ est lisse au point $x$. Comme $\sh{O}_{Y,y}$ est supposé intègre et géométriquement

\oldpage[IV-4]{169}

unibranche, il en est de même de $\sh{O}_{Z,z}$ où $z=g(x)$, puisque le morphisme structural $Z\to Y$ est lisse (17.3.8). En vertu de (18.10.1), il suffit de montrer que l'homomorphisme $\sh{O}_{Z,z}\to\sh{O}_{X,x}$ est injectif; il revient au même de voir que le morphisme $\Spec(\sh{O}_{X,x})\to\Spec(\sh{O}_{Z,z})$ correspondant est dominant, puisque $\sh{O}_{Z,z}$ est intègre (I, 1.2.7).

Soit $y'$ le point générique de $Y=\Spec(\sh{O}_{Y,y})$ et soit $z'$ le point générique de $Z=\Spec(\sh{O}_{Y,y}[T_1,\ldots,T_r])$ qui est intègre; si $h:Z\to Y$ est le morphisme structural, on a $h(z')=y'$. Il suffit de prouver qu'il existe dans $f^{-1}(y')$ une générisation $x'$ de $x$ telle que l'image de $x'$ par le morphisme $g_{y'}:f^{-1}(y')\to h^{-1}(y')$ soit égale à $z'$. D'ailleurs $z'$ est aussi le point générique de $h^{-1}(y')=\Spec(k(y')[T_1,\ldots,T_r])$; comme le morphisme $g_{y'}$ est non ramifié, donc quasi-fini (17.4.2), il en est de même de sa restriction à toute composante irréductible de $f^{-1}(y')$; comme $f^{-1}(y')$ et $h^{-1}(y')$ sont localement noethériens et que $\dim(h^{-1}(y'))=r$, il résulte de (5.4.1) que la restriction de $g_{y'}$ à toute composante irréductible de dimension $\geq r$ de $f^{-1}(y')$ est un morphisme dominant. Or, il existe par hypothèse une telle composante contenant une générisation $x'$ de $x$; \emph{a fortiori} son point générique est une générisation de $x$, d'où la conclusion.

\subsection{Application aux algèbres locales noethériennes complètes sur un corps.}
\EGAReaderTarget{ega4.c04.subsec.18-11}{subsection:IV.18.11}

Le lemme suivant généralise (0, 21.9.1) et (0, 21.9.2):

\begin{lemma}[18.11.1]
\label{IV.18.11.1}
Soient $k$ un corps; $A$ un anneau local noethérien complet qui est une $k$-algèbre, $K$ le corps résiduel de $A$.
\begin{enumerate}
  \item[(i)] Si le $K$-espace vectoriel $\Omega^1_{K/k}$ est de rang fini, le $A$-module $\hat{\Omega}^1_{A/k}$ est de type fini.
  \item[(ii)] Si de plus $A$ est une $k$-algèbre formellement lisse (pour la topologie adique), $\hat{\Omega}^1_{A/k}$ est un $A$-module libre de rang égal à
\[
    \dim(A)+\operatorname{rg}_K(\Omega^1_{K/k})-\operatorname{rg}_K(\Upsilon_{K/k}).
\]
\end{enumerate}
En outre, pour tout sous-corps $k_0$ de $k$ tel que $\Omega^1_{k/k_0}$ soit un $k$-espace vectoriel de rang fini, $\hat{\Omega}^1_{A/k_0}$ est un $A$-module libre de rang égal à
\[
  \dim(A)+\operatorname{rg}_K(\Omega^1_{K/k_0})-\operatorname{rg}_K(\Upsilon_{K/k_0})+\operatorname{rg}_k(\Omega^1_{k/k_0}).
\]
\end{lemma}

L'assertion (i) n'est mise que pour mémoire, ayant été prouvée dans (0, 20.7.15). Pour prouver (ii), notons que cette assertion a été en fait établie par la démonstration de (0, 21.9.2), l'énoncé de (0, 21.9.2) faisant seul intervenir la finitude du degré de transcendance de $K$ sur $k$ (par l'intermédiaire de l'égalité de Cartier (0, 21.7.1)).

\begin{lemma}[18.11.2]
\label{IV.18.11.2}
Soient $k$ un corps de caractéristique $p>0$, $C$ une $k$-algèbre locale noethérienne, complète, dont le corps résiduel est une extension finie de $k$, et qui est intègre; soient $n$ sa dimension, $L$ son corps des fractions. Alors $\Omega^1_{L/k}$ et $\Upsilon_{L/k}$ sont des $L$-espaces vectoriels de rang fini, et l'on a
\begin{equation}
\label{IV.18.11.2.1}
\tag{18.11.2.1}
  \operatorname{rg}_L(\Omega^1_{L/k})-\operatorname{rg}_L(\Upsilon_{L/k})\geq n.
\end{equation}

\oldpage[IV-4]{170}

Supposons en outre que $[k:k^p]<+\infty$. Alors on a l'égalité
\begin{equation}
\label{IV.18.11.2.2}
\tag{18.11.2.2}
  \operatorname{rg}_L(\Omega^1_{L/k})-\operatorname{rg}_L(\Upsilon_{L/k})=n.
\end{equation}
En outre, $\Omega^1_{C/k}$ est alors isomorphe à $\hat{\Omega}^1_{C/k}$, donc (18.11.1) est un $C$-module de type fini.
\end{lemma}

On sait en effet (0, 19.8.9) qu'il existe une sous-$k$-algèbre $C_0$ de $C$ isomorphe à une algèbre de séries formelles $k[[T_1,\ldots,T_n]]$ et telle que $C$ soit une $C_0$-algèbre finie. Par suite, $L$ est une extension finie du corps des fractions $L_0=k((T_1,\ldots,T_n))$ de $C_0$. Or, on a la suite exacte de $L$-espaces vectoriels obtenue en appliquant (0, 20.6.17.1) au sous-corps premier de $k$ et aux trois corps $k\subset L_0\subset L$ (et en tenant compte de (0, 20.6.21, (i))
\[
  0\to\Upsilon_{L_0/k}\otimes_{L_0}L\to\Upsilon_{L/k}\to\Upsilon_{L/L_0}
  \to\Omega^1_{L_0/k}\otimes_{L_0}L\to\Omega^1_{L/k}\to\Omega^1_{L/L_0}\to0.
\]
Puisque $L$ est une extension finie de $L_0$, $\Omega^1_{L/L_0}$ et $\Upsilon_{L/L_0}$ sont des $L$-espaces vectoriels de rang fini ayant même rang, en vertu de l'égalité de Cartier (0, 21.7.1). Comme $L_0$ est séparable sur $k$ (0, 21.9.6.4), on a $\Upsilon_{L_0/k}=0$ (0, 20.6.3); on déduit donc déjà de la suite exacte précédente que $\Upsilon_{L/k}$ est de rang fini et que l'on a dans tous les cas ($0_{\mathrm{III}}$, 11.10.2)
\[
  \operatorname{rg}_L(\Omega^1_{L/k})-\operatorname{rg}_L(\Upsilon_{L/k})
  =\operatorname{rg}_{L_0}(\Omega^1_{L_0/k})-\operatorname{rg}_{L_0}(\Upsilon_{L_0/k}).
\]

Pour prouver (18.11.2.1) ou (18.11.2.2), on peut donc se borner à prouver ces relations en remplaçant $C$ et $L$ par $C_0$ et $L_0$. Comme $L_0$ est séparable sur $k$ (0, 21.9.6.4), on a $\Upsilon_{L_0/k}=0$ (0, 20.6.3); d'autre part $\Omega^1_{L_0/k}=\Omega^1_{C_0/k}\otimes_{C_0}L_0$ (0, 20.5.9). On sait (0, 21.9.4) que si $C_1=k[[T_1^p,\ldots,T_n^p]]$, $\hat{\Omega}^1_{C_0/k}$ s'identifie à $\Omega^1_{C_0/C_1}$; d'autre part, on a $\Omega^1_{C_0/k}=\Omega^1_{C_0/k[C_0^p]}$ (0, 21.1.5), et comme $C_0^p=k^p[[T_1^p,\ldots,T_n^p]]$, on a $k[C_0^p]\subset C_1$, avec égalité lorsque $[k:k^p]<+\infty$. Or, on a une suite exacte $\Omega^1_{C_0/k[C_0^p]}\to\Omega^1_{C_0/C_1}\to0$ (0, 20.5.7), d'où une suite exacte $\Omega^1_{L_0/k}\to\hat{\Omega}^1_{C_0/k}\otimes_{C_0}L_0\to0$. Comme $\hat{\Omega}^1_{C_0/k}$ est un $C_0$-module libre de rang $n$ (0, 21.9.3), on voit que l'on a dans tous les cas $\operatorname{rg}_{L_0}(\Omega^1_{L_0/k})\geq n$, avec égalité lorsque $[k:k^p]<+\infty$; ceci prouve déjà (18.11.2.1) et (18.11.2.2).

Enfin, pour voir que $\Omega^1_{C/k}$ est isomorphe à $\hat{\Omega}^1_{C/k}$ lorsque $[k:k^p]<+\infty$, il suffit de prouver que $\Omega^1_{C/k}$ est un $C$-module de type fini, puisque $C$ est un anneau local noethérien complet ($0_{\mathrm I}$, 7.3.3). Puisque $\Omega^1_{C/k}$ est isomorphe à $\Omega^1_{C/k[C^p]}$ et que $k[C_0^p]\subset k[C^p]$, tout revient à prouver que $C$ est un $k[C_0^p]$-module de type fini (0, 20.4.7); mais cela résulte de ce que $C$ est un $C_0$-module de type fini et $C_0$ un $k[C_0^p]$-module de type fini en vertu de l'hypothèse $[k:k^p]<+\infty$.

\begin{lemma}[18.11.3]
\label{IV.18.11.3}
Soient $k$ un corps, $p$ son exposant caractéristique, et supposons que $[k:k^p]<+\infty$. Soit $A$ une $k$-algèbre locale noethérienne complète, dont le corps résiduel est une extension finie de $k$. Soit $\mathfrak p$ un idéal premier de $A$ tel que $A_{\mathfrak p}$ soit géométriquement régulier sur $k$ (6.7.6), de sorte qu'il existe un seul idéal premier minimal $\mathfrak q$ de $A$ contenu dans $\mathfrak p$. Alors $(\hat{\Omega}^1_{A/k})_{\mathfrak p}$ est un $A_{\mathfrak p}$-module libre de rang égal à $\dim(A/\mathfrak q)$.
\end{lemma}

Comme $A$ est noethérien, et $\Spec(A)$ régulier (et \,\emph{a fortiori} intègre) au point $\mathfrak p$, $\mathfrak p$ n'appartient qu'à une seule composante irréductible de $\Spec(A)$, donc ne contient

\oldpage[IV-4]{171}

qu'un seul idéal minimal $\mathfrak q$ de $A$, et en outre on a $\mathfrak q_{\mathfrak p}=0$. Si l'on pose $B=A/\mathfrak q$, la suite (0, 20.7.20)
\[
  \mathfrak q/\mathfrak q^2\longrightarrow
  \hat{\Omega}^1_{A/k}\widehat{\otimes}_A B
  \longrightarrow\hat{\Omega}^1_{B/k}\longrightarrow0
\]
est exacte; en effet, $\hat{\Omega}^1_{A/k}$ est un $A$-module de type fini (18.11.1), donc on a $\hat{\Omega}^1_{A/k}\widehat{\otimes}_A B=\hat{\Omega}^1_{A/k}\otimes_A B=\hat{\Omega}^1_{A/k}/\mathfrak q\hat{\Omega}^1_{A/k}$, et comme ce $B$-module est de type fini, tout sous-$B$-module en est fermé ($0_{\mathrm I}$, 7.3.5), compte tenu de (0, 20.4.5). Comme l'image de $\mathfrak q/\mathfrak q^2$ est dense dans le noyau de l'homomorphisme $\hat{\Omega}^1_{A/k}\widehat{\otimes}_A B\to\hat{\Omega}^1_{B/k}$ (0, 20.7.20), elle est nécessairement égale à ce noyau, d'où notre assertion. En localisant en $\mathfrak p$ la suite exacte précédente, on voit en outre que l'homomorphisme canonique $(\hat{\Omega}^1_{A/k})_{\mathfrak p}\to(\hat{\Omega}^1_{B/k})_{\mathfrak p/\mathfrak q}$ est bijectif. On voit ainsi qu'on peut se borner au cas où $\mathfrak q=0$, autrement dit au cas où $A$ est intègre. Distinguons alors deux cas:

\emph{I) $p>1$.} Alors on peut appliquer (18.11.2) à l'algèbre quotient $C=A/\mathfrak p$, dont le corps des fractions $K$ n'est autre que le corps résiduel de $A_{\mathfrak p}$; on voit donc que l'on a
\begin{equation}
\label{IV.18.11.3.1}
\tag{18.11.3.1}
  \operatorname{rg}_K(\Omega^1_{K/k})-\operatorname{rg}_K(\Upsilon_{K/k})=\dim(A/\mathfrak p).
\end{equation}

Notons maintenant (0, 19.6.6) que $A_{\mathfrak p}$ est une $k$-algèbre formellement lisse pour sa topologie $\mathfrak pA_{\mathfrak p}$-préadique; par suite $\Omega^1_{A_{\mathfrak p}/k}$ est formellement projectif (0, 20.4.9) pour la topologie $\mathfrak p$-préadique (0, 20.4.5); d'autre part, $\Omega^1_{A_{\mathfrak p}/k}=(\Omega^1_{A/k})_{\mathfrak p}$ (0, 20.5.9) est un $A_{\mathfrak p}$-module de type fini, en vertu de (18.11.2) appliqué à $A$. Pour tout entier $j$, $\Omega^1_{A_{\mathfrak p}/k}/\mathfrak p^{j+1}\Omega^1_{A_{\mathfrak p}/k}$ est donc un $(A_{\mathfrak p}/\mathfrak p^{j+1}A_{\mathfrak p})$-module projectif de rang $m=\operatorname{rg}_K(\Omega^1_{A_{\mathfrak p}/k}\otimes_{A_{\mathfrak p}}K)$ (0, 19.2.4); on conclut donc de ($0_{\mathrm{III}}$, 10.2.1 et 10.1.3) que $\hat{\Omega}^1_{A_{\mathfrak p}/k}$ est un $A_{\mathfrak p}$-module libre de rang $m$. Soit $A'=(A_{\mathfrak p})^\wedge$ l'algèbre complétée de $A_{\mathfrak p}$ pour sa topologie $\mathfrak pA_{\mathfrak p}$-préadique; $A'$ est encore une $k$-algèbre formellement lisse (0, 19.3.6) pour sa topologie adique, et il résulte de (0, 20.7.14) et de (0, 20.4.5) que $\hat{\Omega}^1_{A'/k}=\hat{\Omega}^1_{A_{\mathfrak p}/k}$; on en conclut que $\hat{\Omega}^1_{A'/k}$ est un $A'$-module libre de rang $m$. Mais il résulte alors de (18.11.1) et de ce que $\dim(A')=\dim(A_{\mathfrak p})$ (0, 16.2.4) que l'on a
\begin{equation}
\label{IV.18.11.3.2}
\tag{18.11.3.2}
  m=\dim(A_{\mathfrak p})+
  \bigl(\operatorname{rg}_K(\Omega^1_{K/k})-\operatorname{rg}_K(\Upsilon_{K/k})\bigr),
\end{equation}
d'où, en vertu de (18.11.3.1)
\[
  m=\dim(A/\mathfrak p)+\dim(A_{\mathfrak p}).
\]
Mais puisque $A$ est un anneau local noethérien complet, c'est un quotient d'anneau régulier (0, 19.8.8), donc (0, 16.5.12) on a $\dim(A)=\dim(A/\mathfrak p)+\dim(A_{\mathfrak p})$, ce qui achève la démonstration dans ce cas.

\emph{II) $p=1$.} On a, comme ci-dessus $\dim(A)=\dim(A/\mathfrak p)+\dim(A_{\mathfrak p})$; posons $n=\dim(A)$, $r=\dim(A/\mathfrak p)$, $s=\dim(A_{\mathfrak p})$; nous allons voir que $(\hat{\Omega}^1_{A/k})_{\mathfrak p}$ est un $A_{\mathfrak p}$-module libre de rang $n$. Nous allons d'abord prouver le lemme suivant:

\begin{lemma}[18.11.3.3]
\label{IV.18.11.3.3}
Soient $A$ un anneau local noethérien, $\mathfrak p$ un idéal premier de $A$ tel que $A_{\mathfrak p}$ soit un anneau de Cohen-Macaulay. Pour tout système de paramètres $(t_i)_{1\leq i\leq s}$ de $A_{\mathfrak p}$, il existe une suite $(x_i)_{1\leq i\leq s}$ d'éléments de $\mathfrak p$ tels que les $x_i$ fassent partie d'un système de paramètres de $A$

\oldpage[IV-4]{172}

et que, pour tout $i$, $x_i/1$ soit congru à $t_i$ modulo l'idéal $t_1A_{\mathfrak p}+\cdots+t_{i-1}A_{\mathfrak p}$. En particulier, si $A_{\mathfrak p}$ est régulier et si les $t_i$ forment un système régulier de paramètres de $A_{\mathfrak p}$, il en est de même des $x_i/1$.
\end{lemma}

Ce lemme sera lui-même une conséquence du suivant:

\begin{lemma}[18.11.3.4]
\label{IV.18.11.3.4}
Soient $A$ un anneau local noethérien, $\mathfrak p$ un idéal premier de $A$, $x$ un élément de $A$; alors il existe un élément $x'\in A$ tel que $x'/1=x/1$ dans $A_{\mathfrak p}$ et que $x'$ n'appartienne à aucun des idéaux premiers minimaux de $A$ non contenus dans $\mathfrak p$.
\end{lemma}

Montrons d'abord comment (18.11.3.4) entraîne (18.11.3.3). En multipliant au besoin les $t_i$ par des éléments inversibles de $A_{\mathfrak p}$, on peut supposer qu'ils sont de la forme $x_i/1$, avec $x_i\in\mathfrak p$ pour $1\leq i\leq s$. Raisonnons par récurrence sur $s$; comme $t_1$ fait partie d'un système de paramètres de $A_{\mathfrak p}$, il n'appartient à aucun des idéaux premiers minimaux de $A_{\mathfrak p}$ (0, 16.5.5), donc aucun $x'_1\in A$ tel que $x'_1/1=t_1$ ne peut appartenir à un idéal premier minimal de $A$ contenu dans $\mathfrak p$. En vertu de (18.11.3.4), on peut en outre choisir $x'_1$ tel que $x'_1/1=t_1$ (donc $x'_1\in\mathfrak p$) et que $x'_1$ n'appartienne à aucun idéal premier minimal de $A$, donc fasse partie d'un système de paramètres de $A$ (0, 16.3.4 et 16.3.7). On raisonne ensuite par récurrence, en considérant l'anneau $A'=A/x'_1A$ et l'idéal premier $\mathfrak p'=\mathfrak p/x'_1A$ de cet anneau; comme $A'_{\mathfrak p'}=A_{\mathfrak p}/t_1A_{\mathfrak p}$, $A'_{\mathfrak p'}$ est aussi un anneau de Cohen-Macaulay (0, 16.5.5), et les images $t'_i$ ($2\leq i\leq s$) des $t_i$ dans $A'_{\mathfrak p'}$ forment un système de paramètres de cet anneau; il suffit d'appliquer à $A'$ et aux $t'_i$ ($2\leq i\leq s$) l'hypothèse de récurrence. La dernière assertion de (18.11.3.3) résulte de ce que dans $A_{\mathfrak p}$, un système de paramètres qui engendre l'idéal maximal est un système régulier de paramètres (0, 17.1.1).

Prouvons donc (18.11.3.4). Soit $(\mathfrak p'_k)_{1\leq k\leq r}$ la suite des idéaux premiers minimaux de $A$ non contenus dans $\mathfrak p$ et tels que $x\in\mathfrak p'_k$, et soit $(\mathfrak p_j)_{1\leq j\leq n}$ la suite des idéaux premiers minimaux de $A$ autres que les $\mathfrak p'_k$; on peut supposer que $r\geq1$. Comme $\mathfrak p'_k$ ne contient pas $\bigcap_{1\leq j\leq n}\mathfrak p_j$, il existe un élément $y\in\bigcap_{1\leq j\leq n}\mathfrak p_j$ qui n'est contenu dans aucun des $\mathfrak p'_k$ (Bourbaki, \emph{Alg. comm.}, chap. II, \textsection~1, n\textsuperscript{o}~1, prop. 2); d'ailleurs $y/1$ appartient à tous les idéaux premiers minimaux de $A_{\mathfrak p}$, donc est nilpotent; si $(y/1)^h=0$, l'élément $x'=x+y^h$ répondra aux conditions de l'énoncé, car d'une part $y^h\notin\mathfrak p'_k$ pour $1\leq k\leq r$, et par définition des $\mathfrak p'_k$ on a bien $x'\notin\mathfrak p'_k$ pour $1\leq k\leq r$, et d'autre part, si $\mathfrak p_j$ est un idéal premier minimal de $A$ non contenu dans $\mathfrak p$ mais tel que $x\notin\mathfrak p_j$, on a aussi $x'\notin\mathfrak p_j$ puisque $y\in\mathfrak p_j$.

Revenons à la preuve de (18.11.3) lorsque $p=1$. En vertu de (18.11.3.3), il existe un système régulier de paramètres $(t_i)_{1\leq i\leq s}$ de $A_{\mathfrak p}$ tel que $t_i=x_i/1$, où les $x_i$ ($1\leq i\leq s$) appartiennent à $\mathfrak p$ et font partie d'un système de paramètres $(x_j)_{1\leq j\leq n}$ de $A$. Posons $A_0=k[[T_1,\ldots,T_n]]$; comme $A$ est une $k$-algèbre complète et que les $x_j$ appartiennent à l'idéal maximal $\mathfrak m$ de $A$, il existe un $k$-homomorphisme local $u:A_0\to A$ tel que $u(T_j)=x_j$ pour $1\leq j\leq n$ (Bourbaki, \emph{Alg. comm.}, chap. III, \textsection~4, n\textsuperscript{o}~5, prop. 6); si $\mathfrak n$ est l'idéal de $A$ engendré par les $x_j$ ($1\leq j\leq n$), c'est par hypothèse un idéal de définition de $A$; on déduit donc de ($0_{\mathrm I}$, 7.4.4 et 7.4.3) que $u$ fait de $A$ une $A_0$-algèbre finie.

Posons $\mathfrak p_0=\sum_{j=1}^s A_0T_j$, $B_0=(A_0)_{\mathfrak p_0}$ et $B=A\otimes_{A_0}B_0$, de sorte que $u_{\mathfrak p_0}:B_0\to B$ fait

\oldpage[IV-4]{173}

de $B$ une $B_0$-algèbre finie; en outre, si $\mathfrak p'$ est l'idéal de $B$ engendré par $\mathfrak p$, on a $B_{\mathfrak p'}=A_{\mathfrak p}$; comme $\mathfrak p'$ contient $\mathfrak p_0B$ par construction, il est au-dessus de l'idéal maximal $\mathfrak p_0B_0$ de $B_0$. Montrons que le morphisme $\Spec(B)\to\Spec(B_0)$ est non ramifié au point $\mathfrak p'$: cela résulte en effet de ce que $k(\mathfrak p')$ est une extension finie du corps $k(\mathfrak p_0)$ de caractéristique 0, donc est nécessairement séparable, et que l'on a $B_{\mathfrak p'}/\mathfrak p_0B_{\mathfrak p'}=k(\mathfrak p')$ en vertu du choix des $x_j$ pour $1\leq j\leq s$ (17.4.1).

Notons maintenant le lemme suivant:

\begin{lemma}[18.11.3.5]
\label{IV.18.11.3.5}
Soient $k$ un corps, $R$, $S$ deux $k$-algèbres locales noethériennes complètes, telles que, si $K$ est le corps résiduel de $S$, $\Omega^1_{K/k}$ soit de rang fini sur $K$; soit $u:R\to S$ un $k$-homomorphisme faisant de $S$ une $R$-algèbre finie; alors on a $\hat{\Omega}^1_{S/R}=\Omega^1_{S/R}$, et la suite
\begin{equation}
\label{IV.18.11.3.6}
\tag{18.11.3.6}
  \hat{\Omega}^1_{R/k}\otimes_R S\xrightarrow{v}\hat{\Omega}^1_{S/k}\xrightarrow{w}\Omega^1_{S/R}\to0
\end{equation}
(cf. (0, 20.7.17.3)) est exacte.
\end{lemma}

La première assertion résulte de ce que $\Omega^1_{S/R}$ est un $S$-module de type fini (0, 20.4.7) et de ($0_{\mathrm I}$, 7.3.6). D'autre part, on sait (0, 20.7.17.3) que l'image de $v$ est dense dans $\operatorname{Ker}(w)$ et que $w$ est surjectif; mais il résulte de (18.11.1) que $\hat{\Omega}^1_{S/k}$ est un $S$-module de type fini; on en conclut que tout sous-$S$-module de $\hat{\Omega}^1_{S/k}$ est fermé ($0_{\mathrm I}$, 7.3.5), d'où le lemme.

Appliquons ce lemme au cas où $R=A_0$, $S=A$, et notons que $\hat{\Omega}^1_{A_0/k}$ est un $A_0$-module libre de rang $n$ (0, 21.9.3); donc on a $\hat{\Omega}^1_{A_0/k}\widehat{\otimes}_{A_0}A=\hat{\Omega}^1_{A_0/k}\otimes_{A_0}A$ et ce $A$-module est libre de rang $n$. Puisque $\Spec(B)$ est non ramifié sur $\Spec(B_0)$ au point $\mathfrak p'$, on a $(\Omega^1_{A/A_0})_{\mathfrak p}=\Omega^1_{B_{\mathfrak p'}/B_0}=0$ (16.4.15 et 17.4.1). Si on localise la suite exacte (18.11.3.6) (appliquée à $A_0$ et $A$) en $\mathfrak p$, on obtient donc un homomorphisme surjectif
\[
  (\hat{\Omega}^1_{A_0/k}\otimes_{A_0}A)_{\mathfrak p}\longrightarrow(\hat{\Omega}^1_{A/k})_{\mathfrak p},
\]
d'où l'on conclut que le $A_{\mathfrak p}$-module $(\hat{\Omega}^1_{A/k})_{\mathfrak p}$ admet un système de $n$ générateurs. Mais le $A$-module $\hat{\Omega}^1_{A/k}$ est de rang $n$, en vertu de (0, 21.9.5), qui est applicable à l'anneau complet et intègre $A$ en raison du théorème de Cohen (0, 19.8.8, (ii)) et du fait que le corps des fractions de $A$ est de caractéristique 0. Le $A_{\mathfrak p}$-module $(\hat{\Omega}^1_{A/k})_{\mathfrak p}$ est donc aussi de rang $n$, et comme son quotient par son sous-module de torsion admet un système de $n$ générateurs, ce quotient est nécessairement libre; on en déduit aussitôt que les $n$ générateurs de $(\hat{\Omega}^1_{A/k})_{\mathfrak p}$ obtenus ci-dessus forment un système libre, d'où la conclusion.

\begin{lemma}[18.11.4]
\label{IV.18.11.4}
Soient $k$ un corps, $A$ une $k$-algèbre locale noethérienne complète, dont le corps résiduel est une extension finie de $k$. Soit $k'$ une extension quelconque de $k$, et posons $A'=A\widehat{\otimes}_k k'$. Alors:
\begin{enumerate}
  \item[(i)] $A'$ est un anneau semi-local noethérien complet, composé direct d'anneaux locaux complets $A'_i$ ($1\leq i\leq r$) qui sont des $A$-modules fidèlement plats, et dont les corps résiduels sont des extensions finies de $k'$; si $\mathfrak m$ est l'idéal maximal de $A$, $\mathfrak mA'$ est un idéal de définition de $A'$; on a $\dim(A'_i)=\dim(A)$ pour tout $i$.
  \item[(ii)] Pour tout $i$, $\hat{\Omega}^1_{A'_i/k'}$ est canoniquement isomorphe à $\hat{\Omega}^1_{A/k}\otimes_A A'_i$.
\end{enumerate}
\end{lemma}

\oldpage[IV-4]{174}

(i) Les premières assertions résultent aussitôt de (7.5.5) et de ($0_{\mathrm I}$, 6.6.2); le fait que $\mathfrak mA'$ soit un idéal de définition de $A'$ résulte aussi de (7.5.5), car si $K$ est le corps résiduel de $A$, $K\otimes_k k'$ est une $k'$-algèbre finie. Enfin $A'_i/\mathfrak mA'_i$ est un des anneaux locaux artiniens composants directs de $K\otimes_k k'$ (7.5.5), donc est de dimension 0; comme $A'_i$ est un $A$-module plat, l'égalité des dimensions de $A$ et de $A'_i$ résulte de (6.1.2).

(ii) Comme $\hat{\Omega}^1_{A/k}$ est, en vertu de l'hypothèse sur $K$, un $A$-module de type fini (18.11.1), $\hat{\Omega}^1_{A/k}\otimes_A A'$ est complet et s'identifie au produit tensoriel complété $\hat{\Omega}^1_{A/k}\widehat{\otimes}_A A'$; en vertu de l'associativité du produit tensoriel complété ($0_{\mathrm I}$, 7.7.4), $\hat{\Omega}^1_{A/k}\widehat{\otimes}_A A'=\hat{\Omega}^1_{A/k}\widehat{\otimes}_A(A\widehat{\otimes}_k k')$ s'identifie à $\hat{\Omega}^1_{A/k}\widehat{\otimes}_k k'$. Mais $\Omega^1_{A/k}\otimes_k k'$ s'identifie à $\Omega^1_{A''/k'}$, où $A''=A\otimes_k k'$ (0, 20.5.5), et comme $A'$ est par définition le séparé complété de $A''$, le séparé complété de $\Omega^1_{A''/k'}$ s'identifie par construction à celui de $\Omega^1_{A'/k'}$ (0, 20.7.4); autrement dit, on a un isomorphisme canonique
\[
  \hat{\Omega}^1_{A/k}\widehat{\otimes}_k k'\xrightarrow{\sim}\hat{\Omega}^1_{A'/k'}.
\]
La conclusion de (ii) résulte maintenant de ce que $\Omega^1_{A'/k'}$ est somme directe des $\Omega^1_{A'_i/k'}$ (0, 20.4.13), et si $\mathfrak r$ est le radical de $A'$, la topologie $\mathfrak r$-préadique sur $\Omega^1_{A'/k'}$ s'identifie au produit des topologies $\mathfrak m'_i$-préadiques sur les $\Omega^1_{A'_i/k'}$ (où $\mathfrak m'_i$ est l'idéal maximal de $A'_i$); finalement, il en résulte que le séparé complété $\hat{\Omega}^1_{A'/k'}$ pour la topologie $\mathfrak r$-préadique s'identifie au produit des séparés complétés $\hat{\Omega}^1_{A'_i/k'}$ pour les topologies $\mathfrak m'_i$-préadiques, et il suffit d'utiliser (0, 20.4.5).

\begin{proposition}[18.11.5]
\label{IV.18.11.5}
Soient $k$ un corps, $A$ une $k$-algèbre locale noethérienne et complète dont le corps résiduel est une extension finie de $k$, $\mathfrak p$ un idéal premier de $A$, distinct de l'idéal maximal $\mathfrak m$, tel qu'il existe un idéal premier minimal $\mathfrak q\subset\mathfrak p$ pour lequel $\dim(A/\mathfrak q)=\dim(A)$ (ce qui aura lieu en particulier lorsque $A$ est équidimensionnel), $m$ un entier $\geq0$. Les conditions suivantes sont équivalentes:
\begin{enumerate}
  \item[$a)$] Le $A_{\mathfrak p}$-module $(\hat{\Omega}^1_{A/k})_{\mathfrak p}$ admet un système de $m$ générateurs.
  \item[$b)$] Il existe un $k$-homomorphisme local $u:B\to A$, où $B=k[[T_1,\ldots,T_m]]$, faisant de $A$ une $B$-algèbre finie, et tel que le morphisme correspondant $\Spec(A)\to\Spec(B)$ soit non ramifié au point $\mathfrak p$.
\end{enumerate}
\end{proposition}

Pour prouver que $b)$ entraîne $a)$, notons qu'en vertu du lemme (18.11.3.5), on a alors la suite exacte
\[
  \hat{\Omega}^1_{B/k}\otimes_B A\longrightarrow\hat{\Omega}^1_{A/k}\longrightarrow\Omega^1_{A/B}\longrightarrow0
\]
puisque $A$ est une $B$-algèbre finie; localisant en $\mathfrak p$ et notant que par hypothèse on a alors $(\Omega^1_{A/B})_{\mathfrak p}=0$ (17.4.1), on obtient un homomorphisme surjectif $(\hat{\Omega}^1_{B/k}\otimes_B A)_{\mathfrak p}\to(\hat{\Omega}^1_{A/k})_{\mathfrak p}$, et la conclusion résulte de ce que $\hat{\Omega}^1_{B/k}$ est un $B$-module libre de rang $m$ (0, 21.9.3).

Pour prouver que $a)$ implique $b)$, prouvons d'abord les lemmes suivants.

\begin{lemma}[18.11.5.1]
\label{IV.18.11.5.1}
Soient $k$ un corps, $A$ une $k$-algèbre locale noethérienne et complète, dont le corps résiduel est une extension finie de $k$, $\mathfrak p$ un idéal premier de $A$, $\mathfrak q$ un idéal minimal de $A$ contenu dans $\mathfrak p$. Alors le nombre minimum de générateurs du $A_{\mathfrak p}$-module $(\hat{\Omega}^1_{A/k})_{\mathfrak p}$ est au moins égal à $\dim(A/\mathfrak q)$.
\end{lemma}

\oldpage[IV-4]{175}

Posons $n=\dim(A/\mathfrak q)$, et soit $m$ le nombre minimum de générateurs du $A_{\mathfrak p}$-module $(\hat{\Omega}^1_{A/k})_{\mathfrak p}$, qui est égal à $\operatorname{rg}_{k(\mathfrak p)}((\hat{\Omega}^1_{A/k})_{\mathfrak p}\otimes_{A_{\mathfrak p}}k(\mathfrak p))$ (Bourbaki, \emph{Alg. comm.}, chap. II, \textsection~3, n\textsuperscript{o}~2, cor. 2 de la prop. 4). Notons que l'on a $(\hat{\Omega}^1_{A/k})_{\mathfrak q}=(\hat{\Omega}^1_{A/k})_{\mathfrak p}\otimes_{A_{\mathfrak p}}A_{\mathfrak q}$, donc le nombre minimum de générateurs du $A_{\mathfrak q}$-module $(\hat{\Omega}^1_{A/k})_{\mathfrak q}$ est au plus égal à $m$. Il suffit par suite de considérer le cas où $\mathfrak p=\mathfrak q$ est minimal. En second lieu, montrons qu'on peut se borner au cas où $k$ est algébriquement clos. En effet, soit $k'$ une clôture algébrique de $k$, et posons $A'=A\widehat{\otimes}_k k'$, qui est composé direct de $k'$-algèbres locales $A'_i$ ($1\leq i\leq r$), dont les corps résiduels sont isomorphes à $k'$, et qui sont des $A$-modules fidèlement plats (18.11.4). Appliquant (2.3.4) et (6.1.1), on voit que dans un $A'_i$ donné, il existe un idéal minimal $\mathfrak q'_i$ au-dessus de $\mathfrak q$ tel que $n'=\dim(A'_i/\mathfrak q'_i)\geq\dim(A/\mathfrak q)=n$. D'autre part (18.11.4), on a $(\hat{\Omega}^1_{A'_i/k'})_{\mathfrak q'_i}=(\hat{\Omega}^1_{A/k})_{\mathfrak q}\otimes_{A_{\mathfrak q}}(A'_i)_{\mathfrak q'_i}$, donc le nombre minimum $m'$ de générateurs du $(A'_i)_{\mathfrak q'_i}$-module $(\hat{\Omega}^1_{A'_i/k'})_{\mathfrak q'_i}$ est au plus égal à $m$; d'où notre assertion.

Supposons donc $k$ algébriquement clos, et montrons qu'on peut en outre se borner au cas où $\mathfrak q=0$. En effet, on a $\Omega^1_{(A/\mathfrak q)/A}=0$ (0, 20.4.12), donc la suite exacte (18.11.3.6) fournit un homomorphisme surjectif $\hat{\Omega}^1_{A/k}\otimes_A(A/\mathfrak q)\to\hat{\Omega}^1_{(A/\mathfrak q)/k}$, d'où on déduit aussitôt que le nombre minimum de générateurs du $K$-module $\hat{\Omega}^1_{(A/\mathfrak q)/k}\otimes_{A/\mathfrak q}K$, où $K$ est le corps des fractions de $A/\mathfrak q$, est au plus égal à $m$.

Mais si $A$ est intègre et $K$ son corps des fractions, $K$ est séparable sur le corps algébriquement clos $k$, donc une $k$-algèbre géométriquement régulière (6.7.6); on peut donc appliquer (18.11.3) pour $\mathfrak p=0$ en observant que $k=k^p$ puisque $k$ est algébriquement clos, et dans ce cas on a $m=n$.

\begin{lemma}[18.11.5.2]
\label{IV.18.11.5.2}
Soient $k$ un corps, $A$ une $k$-algèbre locale noethérienne complète, dont le corps résiduel est une extension finie de $k$, $\mathfrak m$ son idéal maximal, $\mathfrak p\subset\mathfrak m$ un idéal premier de $A$. Soit $m$ le nombre minimum de générateurs du $A_{\mathfrak p}$-module $(\hat{\Omega}^1_{A/k})_{\mathfrak p}$, et soit $(x_i)_{1\leq i\leq m}$ une famille d'éléments de $\mathfrak m$ n'appartenant pas à $\mathfrak p$. Alors il existe $m$ éléments inversibles $u_i$ ($1\leq i\leq m$) de $A$ tels que si l'on pose $y_i=u_ix_i$, les images canoniques des $dy_i$ dans $(\hat{\Omega}^1_{A/k})_{\mathfrak p}$ engendrent cet $A_{\mathfrak p}$-module.
\end{lemma}

Pour tout $x\in A$, désignons par $\delta(x)$ l'image canonique de $dx$ ($=d_{A/k}x$) dans $(\hat{\Omega}^1_{A/k})_{\mathfrak p}\otimes_{A_{\mathfrak p}}k(\mathfrak p)=\hat{\Omega}^1_{A/k}\otimes_A k(\mathfrak p)$; comme ce $k(\mathfrak p)$-espace vectoriel est par hypothèse de rang $m$, il suffit (en vertu du lemme de Nakayama) de prouver qu'on peut déterminer les $u_i$ tels que les $\delta(u_ix_i)$ forment un système libre. Raisonnons par récurrence et supposons que, pour un entier $r<m$, on ait déterminé les $u_i$ ($1\leq i\leq r$) tels que les $\delta(u_ix_i)$ pour $1\leq i\leq r$ soient linéairement indépendants sur $k(\mathfrak p)$; si $\delta(x_{r+1})$ n'est pas combinaison linéaire des $\delta(u_ix_i)$ pour $1\leq i\leq r$, il suffit de prendre $u_{r+1}=1$ pour poursuivre la récurrence. Dans le cas contraire, notons que pour $u\in A$, $\delta(ux_{r+1})$ est l'image canonique de $(du)x_{r+1}+u(dx_{r+1})$; comme $u(dx_{r+1})$ a une image canonique combinaison linéaire des $\delta(u_ix_i)$ pour $1\leq i\leq r$, et que d'autre part l'image canonique de $x_{r+1}$ dans $k(\mathfrak p)$ est $\neq0$, on voit qu'il suffit de prouver qu'il existe un élément inversible $u\in A$ tel que $\delta(u)$ ne soit pas combinaison linéaire des $\delta(u_ix_i)$ pour $1\leq i\leq r$.

\oldpage[IV-4]{176}

Or, il résulte de (0, 20.7.15) et de ($0_{\mathrm I}$, 7.2.9) que les $\delta(x)$ engendrent le $k(\mathfrak p)$-espace vectoriel $\hat{\Omega}^1_{A/k}\otimes_A k(\mathfrak p)$; comme par hypothèse $r<m$, il existe donc $z\in A$ tel que $\delta(z)$ ne soit pas combinaison linéaire des $\delta(u_ix_i)$ pour $1\leq i\leq r$. Si $z\notin\mathfrak m$, on prendra $u=z$; sinon, $1+z$ est inversible et on a $\delta(1+z)=\delta(z)$ puisque $d(1+z)=dz$; on prendra alors $u=1+z$, ce qui achève de prouver (18.11.5.2).

Revenons maintenant à la démonstration de l'implication $a)\Rightarrow b)$ dans (18.11.5). Notons d'abord que puisque $\mathfrak p\neq\mathfrak m$, il existe un système de paramètres $(x_i)_{1\leq i\leq n}$ de $A$ (avec $n=\dim(A)$) tel que $x_i\notin\mathfrak p$ pour $1\leq i\leq n$: en effet, on ne peut avoir $x_i\in\mathfrak p$ pour tout $i$, sans quoi $A/\mathfrak p$ serait de longueur finie, et $\mathfrak p$ serait maximal, contrairement à l'hypothèse. Mais si par exemple $x_1\notin\mathfrak p$, il suffit, pour chaque indice $i$ tel que $x_i\in\mathfrak p$, de remplacer $x_i$ par $x_1+x_i$, pour avoir un système de paramètres dont aucun élément n'appartient à $\mathfrak p$. L'hypothèse $a)$ et la relation $\dim(A/\mathfrak q)=n$ entraînent, en vertu de (18.11.5.1), que $m\geq n$; on peut donc considérer une famille $(x_i)_{1\leq i\leq m}$ d'éléments $x_i\notin\mathfrak p$, dont les $n$ premiers forment un système de paramètres de $A$. Multipliant en outre les $x_i$ par des éléments inversibles $u_i$ de $A$, on peut, grâce à (18.11.5.2) supposer que les images des $dx_i$ ($1\leq i\leq m$) dans $(\hat{\Omega}^1_{A/k})_{\mathfrak p}$ engendrent cet $A_{\mathfrak p}$-module, et la multiplication par les $u_i$ n'a pas altéré le fait que les $x_i$ pour $1\leq i\leq n$ forment un système de paramètres.

Considérons alors le $k$-homomorphisme local $u:B\to A$ tel que $u(T_i)=x_i$ pour $1\leq i\leq m$ (Bourbaki, \emph{Alg. comm.}, chap. III, \textsection~4, n\textsuperscript{o}~5, prop. 6); comme les $x_i$ engendrent un idéal de définition de $A$, il résulte de ($0_{\mathrm I}$, 7.4.4 et 7.4.3) que $u$ fait de $A$ une $B$-algèbre finie. On a donc (18.11.3.5) la suite exacte
\[
  \hat{\Omega}^1_{B/k}\otimes_B A\xrightarrow{v}\hat{\Omega}^1_{A/k}\longrightarrow\Omega^1_{A/B}\longrightarrow0.
\]
Mais les $d_Ax_i$ sont les images canoniques par $v$ des éléments $d_BT_i\otimes1$ (0, 20.5.2.6). Si on localise la suite exacte précédente en $\mathfrak p$, on voit donc que $v_{\mathfrak p}:(\hat{\Omega}^1_{B/k}\otimes_B A)_{\mathfrak p}\to(\hat{\Omega}^1_{A/k})_{\mathfrak p}$ est surjectif, et par suite on a $0=(\Omega^1_{A/B})_{\mathfrak p}=\Omega^1_{A_{\mathfrak p}/B_{\mathfrak r}}$ (où $\mathfrak r$ est l'image réciproque de $\mathfrak p$ dans $A$, cf. (16.4.15)). En vertu de (17.4.1), cela entraîne la propriété $b)$ de (18.11.5). C.Q.F.D.

\begin{remark}[18.11.6]
\label{IV.18.11.6}
Dans l'énoncé de (18.11.5), on ne peut supprimer l'hypothèse $\mathfrak p\neq\mathfrak m$. En effet, pour tout $k$-homomorphisme local $u:B\to A$, où $B=k[[T_1,\ldots,T_r]]$ ($r$ entier quelconque) faisant de $A$ une $B$-algèbre finie, $\mathfrak m$ est le seul point de $\Spec(A)$ au-dessus de l'idéal maximal $\mathfrak n$ de $B$; le morphisme $\Spec(A)\to\Spec(B)$ ne peut être non ramifié en $\mathfrak m$ que si $A/\mathfrak nA$ est un corps, extension séparable de $k$ (17.4.1), ce qui implique que le corps résiduel $K$ de $A$ est une extension séparable de $k$. Si cette condition n'est pas remplie, la conclusion de (18.11.5) ne peut jamais être vérifiée pour $\mathfrak p=\mathfrak m$, quel que soit l'entier $m$.
\end{remark}

\begin{corollary}[18.11.7]
\label{IV.18.11.7}
Soient $k$ un corps, $A$ une $k$-algèbre locale noethérienne complète équidimensionnelle, dont le corps résiduel est une extension finie de $k$, $\mathfrak p$ un idéal premier de $A$. Alors le nombre minimum de générateurs du $A_{\mathfrak p}$-module $(\hat{\Omega}^1_{A/k})_{\mathfrak p}$ est au moins égal à $\dim(A)$.
\end{corollary}

C'est un cas particulier de (18.11.5.1).

Plus particulièrement:

\oldpage[IV-4]{177}

\begin{corollary}[18.11.8]
\label{IV.18.11.8}
Soient $k$ un corps, $A$ une $k$-algèbre locale noethérienne complète intègre, dont le corps résiduel est une extension finie de $k$. Si $K$ est le corps des fractions de $A$, on a
\[
  \operatorname{rg}_K\bigl((\hat{\Omega}^1_{A/k})\otimes_A K\bigr)\geq\dim(A).
\]
\end{corollary}

Il suffit de faire $\mathfrak p=0$ dans (18.11.7).

\begin{corollary}[18.11.9]
\label{IV.18.11.9}
Soient $k$ un corps, $A$ une $k$-algèbre locale noethérienne complète de dimension $n$, dont le corps résiduel est une extension finie de $k$. Soit $\mathfrak p$ un idéal premier de $A$ distinct de l'idéal maximal, et contenant un idéal premier minimal $\mathfrak q$ tel que $\dim(A/\mathfrak q)=n$. Supposons que le $A_{\mathfrak p}$-module $(\hat{\Omega}^1_{A/k})_{\mathfrak p}$ admette $n$ générateurs. Alors:
\begin{enumerate}
  \item[(i)] Le $A_{\mathfrak p}$-module $(\hat{\Omega}^1_{A/k})_{\mathfrak p}$ est libre de rang $n$.
  \item[(ii)] Il existe un $k$-homomorphisme local $u:B\to A$, où $B=k[[T_1,\ldots,T_n]]$, faisant de $A$ une $B$-algèbre finie, et tel que le morphisme correspondant $\Spec(A)\to\Spec(B)$ soit étale au point $\mathfrak p$.
  \item[(iii)] La $k$-algèbre $A_{\mathfrak p}$ est géométriquement régulière.
\end{enumerate}
\end{corollary}

Prouvons d'abord (ii); en vertu de (18.11.5), il existe un homomorphisme local
\[
  u:B=k[[T_1,\ldots,T_n]]\longrightarrow A
\]
faisant de $A$ une $B$-algèbre finie, et tel que le morphisme correspondant $\Spec(A)\to\Spec(B)$ soit non ramifié au point $\mathfrak p$; posons $\mathfrak r=u^{-1}(\mathfrak p)$. L'hypothèse sur $\mathfrak q\subset\mathfrak p$, et le fait que $A$ soit un quotient d'anneau régulier (0, 19.8.8) impliquent que l'on a $\dim(A_{\mathfrak p})=n-\dim(A/\mathfrak p)$ (0, 16.5.12). On a de même $\dim(B_{\mathfrak r})=n-\dim(B/\mathfrak r)$. Enfin, puisque le morphisme $\Spec(A)\to\Spec(B)$ est non ramifié au point $\mathfrak p$, la fibre de ce morphisme au point $\mathfrak r$ est de dimension 0, donc (0, 16.3.9) on a $\dim(A/\mathfrak p)\leq\dim(B/\mathfrak r)$ et par suite $\dim(B_{\mathfrak r})\leq\dim(A_{\mathfrak p})$. On conclut donc de (18.10.1) que le morphisme $\Spec(A)\to\Spec(B)$ est étale au point $\mathfrak p$.

L'assertion (iii) résulte de ce que $A_{\mathfrak p}$ est formellement lisse sur $B_{\mathfrak r}$ et que $B_{\mathfrak r}$ est formellement lisse sur $k$ (pour les topologies préadiques) (0, 19.3.4 et 19.3.5), donc $A_{\mathfrak p}$ est formellement lisse sur $k$ pour sa topologie préadique, et par suite est géométriquement régulier sur $k$ (0, 22.5.8).

Prouvons enfin (i). Soit $k'$ une extension algébriquement close de $k$, et considérons la $k'$-algèbre semi-locale $A'=A\widehat{\otimes}_k k'$. Si l'on considère $A$ comme $B$-module de type fini au moyen de l'homomorphisme $u$, on peut donc écrire, à un isomorphisme canonique près, $A'=A\otimes_B(B\widehat{\otimes}_k k')$ ((7.5.7.1), dans l'énoncé duquel on rappelle qu'il n'est pas nécessaire de supposer que le corps résiduel de $B$ soit une extension finie de $k$). Posons $B'=B\widehat{\otimes}_k k'$, qui s'identifie canoniquement à l'algèbre de séries formelles $k'[[T_1,\ldots,T_n]]$. Puisque le morphisme $\Spec(A)\to\Spec(B)$ est fini et étale au point $\mathfrak p$, le morphisme $\Spec(A')\to\Spec(B')$ est fini, et étale en tout point $\mathfrak p'$ au-dessus de $\mathfrak p$ (17.3.3); d'ailleurs, $A'$ est composé direct d'anneaux locaux $A'_i$ de dimension $n$ ($1\leq i\leq r$) (18.11.4) et $\mathfrak p'$ s'identifie à un idéal premier $\mathfrak p'_i$ de l'un des $A'_i$. Le même raisonnement que ci-dessus prouve alors que $(A'_i)_{\mathfrak p'_i}$ est géométriquement régulier sur $k'$, et comme $k'$ est parfait, il résulte de (18.11.3) que $(\hat{\Omega}^1_{A'_i/k'})_{\mathfrak p'_i}$ est un $(A'_i)_{\mathfrak p'_i}$-module libre de type fini. Mais puisque $(\hat{\Omega}^1_{A'_i/k'})_{\mathfrak p'_i}$ est isomorphe à $(\hat{\Omega}^1_{A/k})_{\mathfrak p}\otimes_{A_{\mathfrak p}}(A'_i)_{\mathfrak p'_i}$ (18.11.4) et que $(A'_i)_{\mathfrak p'_i}$ est un $A_{\mathfrak p}$-module fidèlement plat, on voit que $(\hat{\Omega}^1_{A/k})_{\mathfrak p}$ est un $A_{\mathfrak p}$-module libre de type fini (2.5.2);

\oldpage[IV-4]{178}

d'ailleurs, en vertu de (18.11.5.1) et de l'hypothèse $\dim(A/\mathfrak q)=n$, on conclut que $(\hat{\Omega}^1_{A/k})_{\mathfrak p}$ est un $A_{\mathfrak p}$-module libre de rang $n$. C.Q.F.D.

\begin{theorem}[18.11.10]
\label{IV.18.11.10}
Soient $k$ un corps d'exposant caractéristique $p$, $A$ une $k$-algèbre locale noethérienne complète, dont le corps résiduel est une extension finie de $k$, $\mathfrak p$ un idéal premier de $A$ distinct de l'idéal maximal. Les conditions suivantes sont équivalentes:
\begin{enumerate}
  \item[$a)$] Pour toute extension $k'$ de $k$, et tout idéal premier $\mathfrak p'$ de $A'=A\widehat{\otimes}_k k'$ au-dessus de $\mathfrak p$, $A'_{\mathfrak p'}$ est un anneau régulier.
  \item[$a')$] Il existe une extension parfaite $k'$ de $k$ et un idéal premier $\mathfrak p'$ de $A'=A\widehat{\otimes}_k k'$ au-dessus de $\mathfrak p$, tels que $A'_{\mathfrak p'}$ soit régulier.
  \item[$b)$] Soit $n$ la plus grande des dimensions des composantes irréductibles de $\Spec(A)$ auxquelles appartient $\mathfrak p$; alors $(\hat{\Omega}^1_{A/k})_{\mathfrak p}$ est un $A_{\mathfrak p}$-module libre de rang $n$.
\end{enumerate}
Supposons en outre que $n=\dim(A)$ (ce qui sera le cas si $A$ est équidimensionnel); alors les conditions précédentes sont aussi équivalentes à:
\begin{enumerate}
  \item[$c)$] Il existe un $k$-homomorphisme local $u:B\to A$, où $B=k[[T_1,\ldots,T_n]]$, faisant de $A$ une $B$-algèbre finie, et tel que le morphisme correspondant $\Spec(A)\to\Spec(B)$ soit étale au point $\mathfrak p$.
\end{enumerate}
Chacune des conditions $a)$, $a')$, $b)$ entraîne la suivante:
\begin{enumerate}
  \item[$d)$] L'anneau $A_{\mathfrak p}$ est géométriquement régulier sur $k$.
\end{enumerate}
Si de plus $[k:k^p]<+\infty$, la condition $d)$ est équivalente à $a)$, $a')$ et $b)$.
\end{theorem}

Le fait que $d)$ entraîne $b)$ lorsque $[k:k^p]<+\infty$ n'est autre que (18.11.3). Il est clair que $a)$ entraîne trivialement $a')$; montrons que $c)$ entraîne $a)$ lorsque $\dim(A)=n$. Avec les notations de $a)$, on a alors
\[
  A'=A\otimes_B B',\qquad B'=B\widehat{\otimes}_k k'=k'[[T_1,\ldots,T_n]]
\]
(7.5.7.1). Le morphisme $\Spec(A')\to\Spec(B')$ est alors fini et étale au point $\mathfrak p'$, et le raisonnement qui prouve (18.11.9, (iii)) montre que $A'_{\mathfrak p'}$ est régulier.

Le fait que $b)$ entraîne $c)$ lorsque $\dim(A)=n$ résulte de (18.11.9). Montrons que $a')$ entraîne $b)$ lorsque $\dim(A)=n$. On sait que $A'$ est un $A$-module plat (18.11.4) et $\dim(A')=\dim(A)=n$; il résulte de (2.3.4) et de (6.1.1) qu'il existe un idéal premier minimal $\mathfrak q'$ de $A'$ contenu dans $\mathfrak p'$, au-dessus de $\mathfrak q$ et tel que $\dim(A'/\mathfrak q')=\dim(A/\mathfrak q)=n$. Par ailleurs, il résulte de (18.11.4) que
\[
  (\hat{\Omega}^1_{A'/k'})_{\mathfrak p'}=(\hat{\Omega}^1_{A/k})_{\mathfrak p}\otimes_{A_{\mathfrak p}}A'_{\mathfrak p'};
\]
comme $k'$ est parfait, l'hypothèse que $A'_{\mathfrak p'}$ est régulier entraîne qu'il est géométriquement régulier sur $k'$ (6.7.7). Puisque $k'^p=k'$, on peut appliquer à $A'$ et $\mathfrak p'$ le fait que $d)$ entraîne $b)$, donc $(\hat{\Omega}^1_{A'/k'})_{\mathfrak p'}$ est un $A'_{\mathfrak p'}$-module libre de rang $n$; par fidèle platitude (18.11.4 et 2.5.2), on en conclut que $(\hat{\Omega}^1_{A/k})_{\mathfrak p}$ est un $A_{\mathfrak p}$-module libre de rang $n$.

Ceci montre l'équivalence de $a)$, $a')$, $b)$ et $c)$ lorsque $\dim(A)=n$. Il reste à prouver que $a)$, $a')$ et $b)$ sont encore équivalentes dans le cas général. Soit $\mathfrak J$ l'idéal de $A$, noyau de l'homomorphisme canonique $A\to A_{\mathfrak p}$, et posons $A_1=A/\mathfrak J$; on a nécessairement $\mathfrak J\subset\mathfrak p$, et si l'on pose $\mathfrak p_1=\mathfrak p/\mathfrak J$, l'homomorphisme canonique $A_{\mathfrak p}\to(A_1)_{\mathfrak p_1}$ est bijectif; on en conclut (I, 6.5.4) que l'injection canonique $\Spec(A_1)\to\Spec(A)$ est un isomorphisme local au point $\mathfrak p_1$, et l'on a $\mathfrak J_{\mathfrak p}=0$. On voit comme dans (18.11.3) que l'on a une suite exacte
\[
  \mathfrak J/\mathfrak J^2\longrightarrow\hat{\Omega}^1_{A/k}\otimes_A A_1\longrightarrow\hat{\Omega}^1_{A_1/k}\longrightarrow0.
\]

\oldpage[IV-4]{179}

Et en localisant en $\mathfrak p$, il vient un isomorphisme
\[
  (\hat{\Omega}^1_{A/k})_{\mathfrak p}\xrightarrow{\sim}(\hat{\Omega}^1_{A_1/k})_{\mathfrak p_1}.
\]
Ceci montre que la condition $b)$ pour l'anneau $A$ et l'idéal $\mathfrak p$ est équivalente à la condition $b)$ pour l'anneau $A_1$ et l'idéal $\mathfrak p_1$. D'autre part, avec les notations de $a)$, on a $A'_1=A_1\widehat{\otimes}_k k'=A'/\mathfrak JA'$ à isomorphisme près ((7.5.7.1), où l'hypothèse sur le corps résiduel de $B$ est superflue); si $\mathfrak p'$ est un idéal premier de $A'$ au-dessus de $\mathfrak p$, tout élément de $\mathfrak JA'$ annule un élément de $A'-\mathfrak p'$, donc $\mathfrak JA'\subset\mathfrak p'$, et si l'on pose $\mathfrak p'_1=\mathfrak p'/\mathfrak JA'$, $\mathfrak p'_1$ est au-dessus de $\mathfrak p_1$ et $(A'_1)_{\mathfrak p'_1}$ s'identifie canoniquement à $A'_{\mathfrak p'}$; ceci montre donc que la condition $a)$ (resp. $a')$) pour l'anneau $A$ et l'idéal $\mathfrak p$ est équivalente à la condition $a)$ (resp. $a')$) pour l'anneau $A_1$ et l'idéal $\mathfrak p_1$. Or, tous les idéaux premiers minimaux de $A$ contenus dans $\mathfrak p$ contiennent $\mathfrak J$ puisque $\Spec(A_1)\to\Spec(A)$ est un isomorphisme local au point $\mathfrak p_1$; d'autre part les idéaux de $\operatorname{Ass}_A(A/\mathfrak J)$ sont les idéaux de $\operatorname{Ass}(A)$ qui sont contenus dans $\mathfrak p$ (Bourbaki, \emph{Alg. comm.}, chap. IV, \textsection~1, n\textsuperscript{o}~2, prop. 6); donc les idéaux premiers minimaux de $A_1$ sont tous contenus dans $\mathfrak p_1$, et on a par suite $\dim(A_1)=n$. Il suffit alors d'appliquer à $A_1$ et à $\mathfrak p_1$ ce qui a été prouvé plus haut.

\begin{remarks}[18.11.11]
\label{IV.18.11.11}
\begin{enumerate}
  \item[(i)] L'équivalence des conditions $d)$ et $b)$ dans (18.11.10) n'est plus valable lorsqu'on ne suppose plus que $[k:k^p]<+\infty$. En effet, dans l'exemple de (0, 22.7.7, (ii)), l'anneau $B=A/\mathfrak q$ est intègre et de dimension 1; d'autre part, la suite
  \[
    (\mathfrak q/\mathfrak q^2)\otimes_A L\xrightarrow{j}\hat{\Omega}^1_{A/k}\otimes_A L\longrightarrow\hat{\Omega}^1_{B/k}\otimes_B L\longrightarrow0
  \]
  est exacte (même démonstration que dans (18.11.3)), et on sait que $j$ n'est pas injectif, donc $j=0$ puisque $(\mathfrak q/\mathfrak q^2)\otimes_A L$ est de rang 1. On en conclut que $\operatorname{rg}_L(\hat{\Omega}^1_{B/k}\otimes_B L)=2$. Comme $B_{\mathfrak p}$ est une $k$-algèbre géométriquement régulière, on voit qu'ici la condition $d)$ n'entraîne pas $b)$.
  \item[(ii)] Les notations étant celles de (18.11.10), supposons que $n=\dim(A)$ et soit $(x_i)_{1\leq i\leq n}$ un système de paramètres de $A$ n'appartenant pas à $\mathfrak p$; on en déduit un $k$-homomorphisme local $u:B\to A$ tel que $u(T_i)=x_i$ pour $1\leq i\leq n$, faisant de $A$ une $B$-algèbre finie. Pour que le morphisme correspondant $\Spec(A)\to\Spec(B)$ soit étale au point $\mathfrak p$, il faut et il suffit que les images des $d_Ax_i$ dans $(\hat{\Omega}^1_{A/k})_{\mathfrak p}$ forment un système de générateurs de cet $A_{\mathfrak p}$-module. En effet, on a vu dans la démonstration de (18.11.5) que si le morphisme $\Spec(A)\to\Spec(B)$ est non ramifié au point $\mathfrak p$, l'homomorphisme canonique $(\hat{\Omega}^1_{B/k}\otimes_B A)_{\mathfrak p}\to(\hat{\Omega}^1_{A/k})_{\mathfrak p}$ est surjectif et les images des éléments $d_BT_i\otimes1$, qui engendrent $(\hat{\Omega}^1_{B/k}\otimes_B A)_{\mathfrak p}$, sont les $d_Ax_i$; d'où la nécessité de la condition. Inversement, le même raisonnement montre que si cette condition est vérifiée, le morphisme $\Spec(A)\to\Spec(B)$ est non ramifié en $\mathfrak p$, et le raisonnement de (18.11.9) prouve en fait que ce morphisme est étale en $\mathfrak p$.
\end{enumerate}
\end{remarks}

\begin{corollary}[18.11.12]
\label{IV.18.11.12}
Soient $k$ un corps d'exposant caractéristique $p$ tel que $[k:k^p]<+\infty$, $A$ une $k$-algèbre locale noethérienne complète et intègre, qui n'est pas un corps, et dont le corps résiduel est une extension finie de $k$. Il existe alors une extension finie radicielle $k'$ de $k$ telle qu'en posant $A'=A\otimes_k k'$, la $k'$-algèbre $A'_{\mathrm{red}}$ et l'idéal premier $0$ de cette algèbre vérifient les conditions équivalentes $a)$, $a')$, $b)$, $c)$ et $d)$ de (18.11.10). En particulier, si $n=\dim(A)=\dim(A')$,

\oldpage[IV-4]{180}

il existe un $k'$-homomorphisme local $B'=k'[[T_1,\ldots,T_n]]\to A'_{\mathrm{red}}$ faisant de $A'_{\mathrm{red}}$ une $B'$-algèbre finie, et tel que le corps des fractions $K'$ de $A'_{\mathrm{red}}$ soit une extension finie séparable du corps des fractions $k'((T_1,\ldots,T_n))$ de $B'$.
\end{corollary}

Le morphisme $\Spec(A')\to\Spec(A)$ étant radiciel et fini, $A'$ est un anneau local complet et son nilradical est le seul idéal premier au-dessus de l'idéal $0$ de $A$; par platitude, $A'$ s'identifie à un sous-anneau de $K\otimes_k k'$, qui est d'ailleurs l'anneau total des fractions de $A'$; on a par suite $K'=(K\otimes_k k')_{\mathrm{red}}$. Il s'agit de prouver, vu l'équivalence $d)\Leftrightarrow c)$ de (18.11.10), qu'il existe une extension radicielle finie $k'$ de $k$ telle que $(K\otimes_k k')_{\mathrm{red}}$ soit une extension séparable de $k'$ (6.7.6). Or on sait (0, 19.8.9) que sous les hypothèses faites, il existe une sous-$k$-algèbre $C=k[[T_1,\ldots,T_m]]$ de $A$ telle que $A$ soit une $C$-algèbre finie; $K$ est donc une extension finie du corps des fractions $K_1=k((T_1,\ldots,T_m))$ de $C$, qui est séparable sur $k$ (0, 21.9.6.4). Si $p$ est l'exposant caractéristique de $k$, on peut donc écrire
\[
  K\otimes_k k^{p^{-\infty}}=K\otimes_{K_1}(K_1\otimes_k k^{p^{-\infty}})
\]
et $K_1\otimes_k k^{p^{-\infty}}$ est un corps, extension radicielle de $K_1$; on en conclut que $K\otimes_k k^{p^{-\infty}}$ est une algèbre finie sur le corps $K_1\otimes_k k^{p^{-\infty}}$, donc un anneau artinien. La conclusion résulte donc du lemme plus général suivant:

\begin{lemma}[18.11.12.1]
\label{IV.18.11.12.1}
Soient $k$ un corps de caractéristique $p>0$, $K$ une extension de $k$. Les conditions suivantes sont équivalentes:
\begin{enumerate}
  \item[$a)$] L'anneau $K\otimes_k k^{p^{-\infty}}$ est artinien.
  \item[$b)$] Il existe une extension $k'$ de $k$ de type fini telle que $(K\otimes_k k')_{\mathrm{red}}$ soit une $k'$-algèbre séparable.
  \item[$c)$] Il existe une extension finie radicielle $k'$ de $k$ telle que $(K\otimes_k k')_{\mathrm{red}}$ soit une extension séparable de $k'$.
\end{enumerate}
\end{lemma}

Montrons d'abord que $c)$ entraîne $a)$. L'anneau $A=K\otimes_k k'$ est alors un anneau artinien local, et si $\mathfrak N$ est son nilradical, le corps résiduel $L=A/\mathfrak N$ est séparable sur $k'$; par suite $L\otimes_{k'}k^{p^{-\infty}}$ est un corps, et comme il est égal à
\[
  (A\otimes_{k'}k^{p^{-\infty}})/(\mathfrak N\otimes_{k'}k^{p^{-\infty}}),
\]
on voit que $\mathfrak N\otimes_{k'}k^{p^{-\infty}}$ est le nilradical de $A\otimes_{k'}k^{p^{-\infty}}=K\otimes_k k^{p^{-\infty}}$; comme $\mathfrak N$ est un idéal de type fini, le nilradical de $K\otimes_k k^{p^{-\infty}}$ est donc de type fini, ce qui entraîne que l'anneau $K\otimes_k k^{p^{-\infty}}$ est artinien.

Inversement, prouvons que $a)$ entraîne $c)$. Soit $\mathfrak N$ le nilradical de l'anneau local artinien $B=K\otimes_k k^{p^{-\infty}}$, qui est par hypothèse engendré par un nombre fini d'éléments de la forme
\[
  y_i=\sum_j x_{ij}\otimes\xi_{ij},\qquad x_{ij}\in K,\quad \xi_{ij}\in k^{p^{-\infty}}.
\]
Soient $k'$ l'extension radicielle finie de $k$ engendrée par les $\xi_{ij}$, $\mathfrak N_0$ l'idéal de $B_0=K\otimes_k k'$ engendré par les $y_i$; il est clair que les $y_i$ sont nilpotents dans $B_0$; d'autre part, on a $\mathfrak N_0\otimes_{k'}k^{p^{-\infty}}=\mathfrak N$, et par suite $\mathfrak N\cap B_0=\mathfrak N_0$, $\mathfrak N_0$ contient le nilradical de $B_0$, donc il lui est égal. Comme $B/\mathfrak N=(B_0/\mathfrak N_0)\otimes_{k'}k^{p^{-\infty}}$ est réduit, on en conclut que $B_0/\mathfrak N_0=(K\otimes_k k')_{\mathrm{red}}$ est séparable sur $k$ (4.6.1).

Il est clair que $c)$ entraîne $b)$. Inversement, supposons $b)$ vérifiée, et notons qu'il existe une extension séparable $k_1$ de $k$ telle que $k'$ soit une extension radicielle finie de $k_1$; posons $K_1=K\otimes_k k_1$, qui est un corps. Appliquant l'équivalence de $a)$ et $c)$ à l'extension $K_1$ de $k_1$, on voit que $K_1\otimes_{k_1}k_1^{p^{-\infty}}$ est un anneau artinien; mais cet anneau est égal

\oldpage[IV-4]{181}

à
\[
  K\otimes_k k_1^{p^{-\infty}}=(K\otimes_k k^{p^{-\infty}})\otimes_{k^{p^{-\infty}}}k_1^{p^{-\infty}},
\]
donc $K\otimes_k k^{p^{-\infty}}$ est aussi artinien (Bourbaki, \emph{Alg. comm.}, chap. I, \textsection~3, n\textsuperscript{o}~5, cor. de la prop. 8); on a ainsi prouvé que $b)$ entraîne $a)$, ce qui achève la démonstration de (18.11.12.1) et de (18.11.12).

\subsection{Applications de la localisation étale aux morphismes quasi-finis (généralisations de résultats antérieurs)}
\EGAReaderTarget{ega4.c04.subsec.18-12}{subsection:IV.18.12}

Les résultats de ce numéro nous ont été communiqués par P.~Deligne.

\begin{theorem}[18.12.1]
\label{IV.18.12.1}
Soient $f:X\to Y$ un morphisme localement de type fini, $x$ un point de $X$, $y=f(x)$. Supposons que $x$ soit un point isolé de l'espace $f^{-1}(y)$. Alors, il existe un morphisme étale $Y'\to Y$, un point $x'$ de $X'=X\times_Y Y'$ au-dessus de $x$ et un voisinage ouvert $V'$ de $x'$ dans $X'$ tel que, si $f'=f_{(Y')}:X'\to Y'$, $f'|V'$ soit un morphisme fini. Si de plus $f$ est séparé, $V'$ est à la fois ouvert et fermé, et $X'$ est donc somme de deux sous-préschémas induits sur des ouverts de $X'$, dont l'un est fini sur $Y'$ et contient $x'$.
\end{theorem}

La dernière assertion résulte de ce que, si $f$ est séparé, il en est de même de $f'$; donc, si $j':V'\to X'$ est l'injection canonique, le fait que le morphisme $f'|V'=f'\circ j'$ soit fini entraîne qu'il en est de même de $j'$ (II, 6.1.5), et par suite $V'=j'(V')$ est fermé dans $X'$ (II, 6.1.10).

La question étant locale sur $X$, on peut supposer $Y=\Spec(A)$ et $X=\Spec(B)$ affines, $B$ étant une $A$-algèbre de type fini, donc de la forme $C/\mathfrak J$, où $C=A[T_1,\ldots,T_r]$ et $\mathfrak J$ est un idéal de $C$; $f$ est donc séparé, et on peut en outre supposer que $f^{-1}(y)=\{x\}$. Soit $(\mathfrak J_\lambda)$ la famille des idéaux de type fini de $C$ contenus dans $\mathfrak J$, de sorte que $\mathfrak J$ est la réunion filtrante des $\mathfrak J_\lambda$. Si $X$ et les $X_\lambda=\Spec(C/\mathfrak J_\lambda)$ sont considérés comme des sous-préschémas fermés de $Z=\Spec(C)$, on a donc, pour les espaces sous-jacents, $X=\bigcap_\lambda X_\lambda$; si $f_\lambda:X_\lambda\to Y$ est le morphisme structural, on en déduit que
\[
  f^{-1}(y)=\bigcap_\lambda f_\lambda^{-1}(y),
\]
et comme les ensembles $f_\lambda^{-1}(y)$ sont fermés dans l'espace noethérien $Z_y=\Spec(k(y)[T_1,\ldots,T_r])$, il existe un indice $\lambda$ tel que $f_\lambda^{-1}(y)=f^{-1}(y)=\{x\}$. On peut donc supposer que les $X_\lambda$ vérifient la même condition que $X$ au point $x$; si l'on prouve la proposition pour un tel $X_\lambda$, elle en résultera pour $X$, car si $x'$ est un point au-dessus de $x$ dans $Z'=Z\times_Y Y'$, $X'_\lambda=p^{-1}(X_\lambda)$, où $p:Z'\to Z$ est la projection (de sorte qu'on a aussi $X'=p^{-1}(X)$), s'il existe un voisinage ouvert $V'_\lambda$ de $x'$ dans $X'_\lambda$ qui est fini sur $Y'$, \emph{a fortiori} la restriction de $f'_\lambda$ au sous-préschéma fermé $X'\cap V'_\lambda$ de $V'_\lambda$ sera un morphisme fini. On est donc ramené au cas où $f$ est de présentation finie.

Notons en outre que l'ensemble des points de $X$ isolés dans leur fibre est ouvert dans $X$ (13.1.4), donc $f$ est quasi-fini au point $x$, et on peut par suite se borner au cas où $f$ est quasi-fini. Soit $A''$ l'anneau local hensélisé de $\sh{O}_{Y,y}$, et posons $Y''=\Spec(A'')$. Si $X''=X\times_Y Y''$, $f''=f_{(Y'')}:X''\to Y''$ est séparé, quasi-fini et de présentation finie; comme $A''$ est hensélien, il y a pour tout $x''\in X''$ au-dessus de $x$ un voisinage ouvert et fermé $V''$ de $x''$ dans $X''$ qui est fini sur $Y''$ (18.5.11, c)); comme l'immersion $V''\to X''$ est ouverte et fermée, elle est quasi-compacte, donc de présentation finie (1.6.2), donc $f''|V''$ est de présentation finie.

Cela étant, par définition, $A''$ est limite inductive d'une famille filtrante $(A_\lambda)$ de $\sh{O}_{Y,y}$-algèbres strictement essentiellement étales (18.6.5); chaque $A_\lambda$ est lui-même limite inductive d'une famille filtrante $(B_{\lambda\mu})$ de $\sh{O}_{Y,y}$-algèbres étales ((18.6.1) et (8.1.2, a)). Enfin, comme $B_{\lambda\mu}$ peut être supposée une $\sh{O}_{Y,y}$-algèbre de présentation finie, il résulte encore de (8.1.2, a)) et de (17.7.8) que $B_{\lambda\mu}$ est limite inductive d'une famille filtrante $(C_{\lambda\mu\nu})$ de $A$-algèbres étales. Finalement, on voit, en utilisant le théorème de la double limite inductive, que $Y''$ est limite projective d'une famille filtrante $(Y'_\alpha)$ de schémas affines étales sur $Y$. Si $x'_\alpha$ est la projection de $x''$ sur $X'_\alpha=X\times_Y Y'_\alpha$, on peut supposer que le voisinage $V''$, qui est quasi-compact, est de la forme $V'_\alpha\times_{Y'_\alpha}Y''$ où $V'_\alpha$ est un voisinage ouvert de $x'_\alpha$ dans $X'_\alpha$ (8.2.11). Enfin, $V''$ étant de présentation finie sur $Y''$, on conclut de (8.10.5, (x)) que pour un $\alpha$ convenable, $V'_\alpha$ est fini sur $Y'_\alpha$. C.Q.F.D.

\begin{remark}[18.12.2]
\label{IV.18.12.2}
Dans la démonstration précédente, on voit (compte tenu de (18.6.2)) que l'on a construit un $Y'$ répondant à la question, et tel en plus que si $y'=f'(X')$, l'homomorphisme $k(y)\to k(y')$ soit bijectif.
\end{remark}

\begin{corollary}[18.12.3]
\label{IV.18.12.3}
Soient $f:X\to Y$ un morphisme localement de type fini et séparé, $y$ un point de $Y$ tel que le sous-espace $f^{-1}(y)$ soit fini et discret. Il existe alors un morphisme étale $Y'\to Y$, un point $y'\in Y'$ au-dessus de $y$, tel que $k(y')=k(y)$, et une décomposition de $X'$ en somme de deux sous-préschémas $X'_1$, $X'_2$, induits sur des ouverts de $X'$, tels que la restriction de $f'=f_{(Y')}:X'\to Y'$ à $X'_1$ soit un morphisme fini et que l'on ait $X'_2\cap f'^{-1}(y')=\varnothing$.
\end{corollary}

Si $n$ est le nombre de points de $f^{-1}(y)$, on raisonne par récurrence sur $n$, le corollaire étant trivial pour $n=0$. Soit $x$ un point de $f^{-1}(y)$; en vertu de (18.12.1) et (18.12.2) il y a un morphisme étale $Y_1\to Y$, un point $y_1$ de $Y_1$ au-dessus de $y$ tel que $k(y_1)=k(y)$ et si l'on pose $S_1=X\times_Y Y_1$, $f_1=f_{(S_1)}:S_1\to Y_1$, il existe un point $x_1$ de $f_1^{-1}(y_1)$ tel que $S_1$ soit somme de deux ouverts $V_1$, $X_1$, $V_1$ étant fini sur $Y_1$ et un voisinage de $x_1$.

\oldpage[IV-4]{182}

En vertu de la relation $k(y_1)=k(y)$, la fibre $f_1^{-1}(y_1)$ dans $S_1$ est isomorphe à $f^{-1}(y)$, donc $X_1\cap f_1^{-1}(y_1)$ est fini, discret et a $n-1$ points. Comme $f_1|X_1$ est localement de type fini et séparé, on applique l'hypothèse de récurrence à ce morphisme: il y a un morphisme étale $Y'\to Y_1$, un point $y'\in Y'$ au-dessus de $y_1$ tel que $k(y')=k(y_1)$, et, si $X'=X\times_Y Y'$ et si $p:X'\to S_1$ est la projection canonique, on a une décomposition de $p^{-1}(X_1)$ en somme de deux sous-préschémas $U'$ et $X'_2$ induits sur des ouverts de $X'$ tels que $X'_2\cap f'^{-1}(y')=\varnothing$ et que $U'$ soit fini sur $Y'$. Par ailleurs $V'=p^{-1}(V_1)$ est fini sur $Y'$ et $X'$ est somme de $U'$, $V'$ et $X'_2$; on répondra donc à la question en prenant $X'_1$ somme de $U'$ et $V'$.

Le corollaire suivant améliore (8.11.1):

\begin{corollary}[18.12.4]
\label{IV.18.12.4}
Pour qu'un morphisme $f:X\to Y$ soit fini, il faut et il suffit qu'il soit séparé, universellement fermé, localement de type fini et que, pour tout $y\in Y$, $f^{-1}(y)$ soit un espace fini discret; en particulier, un morphisme propre et quasi-fini est fini.
\end{corollary}

En effet, on peut appliquer le corollaire (18.12.3) à un point quelconque $y$ de $Y$. Avec les notations de ce corollaire, $f'$ est un morphisme fermé, donc, puisque $X'_2$ est fermé dans $X'$, $f'(X'_2)$ est fermé dans $Y'$ et ne contient pas $y'$; il existe par suite un voisinage ouvert $U'$ de $y'$ dans $Y'$ tel que $U'\to Y$ soit de type fini et que $f'^{-1}(U')=X\times_Y U'$ soit fini sur $U'$. Soit $U$ l'image de $U'$ dans $Y$, qui est un ouvert de $Y$ en vertu de (II, 3.1) puisque le morphisme $Y'\to Y$ est étale, donc plat et localement de présentation finie; on a évidemment encore $X\times_Y U'=X\times_U U'=f^{-1}(U)\times_U U'$. Puisque maintenant le morphisme $U'\to U$ est fidèlement plat et quasi-compact, on déduit de (2.7.1, (xv)) que le morphisme $f^{-1}(U)\to U$, restriction de $f$, est fini, ce qui prouve le corollaire.

\begin{remark}[18.12.5]
\label{IV.18.12.5}
On notera que dans la démonstration de (18.12.4) on n'a pas utilisé l'hypothèse que $f$ est universellement fermé sous sa forme générale, mais seulement que $f_{(Y')}$ est un morphisme fermé pour tout morphisme étale $Y'\to Y$.
\end{remark}

\begin{corollary}[18.12.6]
\label{IV.18.12.6}
Soit $f:X\to Y$ un morphisme de préschémas. Les conditions suivantes sont équivalentes:
\begin{enumerate}
  \item[$a)$] $f$ est une immersion fermée.
  \item[$b)$] $f$ est un monomorphisme propre.
  \item[$c)$] $f$ est propre et, pour tout $y\in Y$, le $k(y)$-préschéma $X_y=f^{-1}(y)$ est radiciel et géométriquement réduit sur $k(y)$ (c'est-à-dire vide ou $k(y)$-isomorphe à $\Spec(k(y))$).
\end{enumerate}
\end{corollary}

Cela se déduit de (18.12.4) comme (8.11.5) se déduit de (8.11.1).

\begin{proposition}[18.12.7]
\label{IV.18.12.7}
Soient $f:X\to Y$ un morphisme localement de type fini, $y$ un point de $Y$. Pour qu'il existe un voisinage ouvert $U$ de $y$ tel que la restriction $f^{-1}(U)\to U$ de $f$ soit une immersion fermée, il faut et il suffit que $X_y$ soit un $k(y)$-préschéma radiciel et géométriquement réduit sur $k(y)$ (c'est-à-dire vide ou $k(y)$-isomorphe à $\Spec(k(y))$) et qu'il existe un voisinage ouvert $V$ de $y$ tel que la restriction $f^{-1}(V)\to V$ de $f$ soit un morphisme universellement fermé.
\end{proposition}

Les conditions étant évidemment nécessaires, prouvons qu'elles sont suffisantes. On peut se borner au cas où $X_y\neq\varnothing$. En vertu de la seconde condition, on peut déjà supposer $f$ universellement fermé. Montrons en outre qu'on peut supposer $f$ affine, donc séparé; cela résulte du lemme suivant:

\begin{lemma}[18.12.7.1]
\label{IV.18.12.7.1}
Soient $f:X\to Y$ un morphisme fermé, $y\in Y$ un point tel que tout voisinage de $X_y$ contienne un voisinage affine de $X_y$ (condition toujours vérifiée lorsque $X_y$ est vide ou réduite à un seul point). Alors il existe un voisinage ouvert affine $U$ de $y$ tel que la restriction $f^{-1}(U)\to U$ de $f$ soit un morphisme affine.
\end{lemma}

En effet, soit $U_0$ un voisinage ouvert affine de $y$ dans $Y$ et soit $V$ un voisinage ouvert affine de $X_y$ contenu dans $f^{-1}(U_0)$. Puisque $f$ est fermé, il existe un voisinage ouvert affine $U\subset U_0$ de $y$ tel que $f^{-1}(U)\subset V$. Comme la restriction $g:V\to U_0$ de $f$ est un morphisme affine, il en est de même de sa restriction $f^{-1}(U)\to U$ (II, 1.2.5).

Supposons donc $f$ affine et universellement fermé, et montrons qu'il existe un voisinage ouvert $U$ de $y$ tel que la restriction $f^{-1}(U)\to U$ soit un morphisme fini; comme on peut supposer $Y$ et $X$ affines, on peut supposer $f$ de type fini, donc propre, et il suffit de prouver qu'il existe un voisinage $U$ de $y$ tel que $f^{-1}(U)\to U$ soit un morphisme quasi-fini (18.12.4). Mais il existe un voisinage $V$ de l'ensemble $X_y$ (réduit à un seul point) tel que $f|V$ soit quasi-fini (13.1.4), et comme $f$ est fermé il y a un voisinage $U$ de $y$ tel que $f^{-1}(U)\subset V$, ce qui achève la démonstration, compte tenu de (18.12.6).

\begin{proposition}[18.12.8]
\label{IV.18.12.8}
Soit $f:X\to Y$ un morphisme de préschémas. Pour que $f$ soit entier, il faut et il suffit que $f$ soit affine et universellement fermé.
\end{proposition}

Les conditions sont nécessaires (II, 6.1.10); prouvons qu'elles sont suffisantes. On peut supposer $Y=\Spec(A)$ affine, donc aussi $X=\Spec(B)$, et il faut prouver que tout élément $b\in B$ est entier sur $A$ (II, 6.1.1). Soit $B'$ la sous-$A$-algèbre de $B$ engendrée par $b$, qui est une $A$-algèbre de type fini; posons $X'=\Spec(B')$, de sorte que $f:X\to Y$ se factorise en
\[
  X\xrightarrow{g}X'\xrightarrow{h}Y,
\]
où $h$ est de type fini, et $g$ est dominant puisque l'homomorphisme $B'\to B$ est injectif (I, 1.2.7). Comme $h$ est séparé et $h\circ g$ universellement fermé, $g$ est universellement fermé (II, 5.4.3 et

\oldpage[IV-4]{183}

5.4.9), donc surjectif puisqu'il est dominant; on en conclut (II, 5.4.3 et 5.4.9) que $h$ est aussi universellement fermé, donc propre puisqu'il est de type fini et séparé. Mais alors, pour tout $y\in Y$, le morphisme $h^{-1}(y)\to\Spec(k(y))$ déduit de $h$ est propre et affine, donc fini (III, 4.4.2), autrement dit $h$ est quasi-fini; on déduit donc de (18.12.4) que $h$ est fini, ce qui prouve que $B'$ est une $A$-algèbre finie, et par suite que $b$ est entier sur $A$. C.Q.F.D.

\begin{remark}[18.12.9]
\label{IV.18.12.9}
Il se peut qu'un morphisme $f:X\to Y$ soit entier lorsqu'on le suppose seulement séparé, universellement fermé et tel que pour tout $y\in Y$, le morphisme $f^{-1}(y)\to\Spec(k(y))$ déduit de $f$ soit entier. Il faudrait pour cela prouver que ces conditions impliquent que $f$ est affine, ou encore que toute fibre $X_y$ est contenue dans un voisinage ouvert affine.
\end{remark}

\begin{corollary}[18.12.10]
\label{IV.18.12.10}
Un morphisme $f:X\to Y$ qui est injectif et universellement fermé est entier.
\end{corollary}

Il suffit, par (18.12.8), de prouver que $f$ est affine, ce qui résulte du lemme (18.12.7.1) et de l'hypothèse.

\begin{corollary}[18.12.11]
\label{IV.18.12.11}
Soit $f:X\to Y$ un morphisme de préschémas (resp. un morphisme localement de type fini). Pour que $f$ soit un homéomorphisme universel (2.4.2), il faut et il suffit que $f$ soit entier (resp. fini), radiciel et surjectif.
\end{corollary}

On sait que les conditions sont suffisantes (2.4.5); elles sont nécessaires en vertu de (18.12.10) et (18.12.4).

La proposition suivante améliore (8.11.2):

\begin{proposition}[18.12.12]
\label{IV.18.12.12}
Tout morphisme $f:X\to Y$ qui est quasi-fini et séparé est quasi-affine.
\end{proposition}

Posons $\mathscr A=f_*(\sh{O}_X)$, qui est une $\sh{O}_Y$-Algèbre quasi-cohérente puisque $f$ est quasi-compact et séparé (1.7.4); soit $Z=\Spec(\mathscr A)$ (II, 1.3.1) de sorte que $f$ se factorise canoniquement en
\[
  X\xrightarrow{g}Z\xrightarrow{h}Y,
\]
$h$ étant affine et $g$ correspondant à l'automorphisme identique de $\mathscr A$ (II, 1.2.7); il suffira de prouver que $g$ est une immersion ouverte, ou, ce qui revient au même (17.9.1) que $g$ est étale et radiciel. Il suffira évidemment de montrer que, pour tout $y\in Y$, le morphisme $g$ est étale et radiciel en chaque point de $f^{-1}(y)$. Or, pour tout $y\in Y$, on peut appliquer à $f$ le résultat de (18.12.3), dont nous conservons les notations; puisque le morphisme $Y'\to Y$ est plat, et que $f$ est quasi-compact et séparé, on a, à un isomorphisme canonique près,
\[
  f'_*(\sh{O}_{X'})=\mathscr A\otimes_{\sh{O}_Y}\sh{O}_{Y'}=\mathscr A'
\]
(2.3.1), donc $Z'=Z\times_Y Y'$ s'identifie à $\Spec(f'_*(\sh{O}_{X'}))$, et dans la factorisation canonique
\[
  X'\xrightarrow{g'}Z'\xrightarrow{h'}Y'
\]
de $f'$ (II, 1.2.7), on a $h'=h_{(Y')}$ et $g'=g_{(Y')}$. Cela étant, la décomposition de $X'$ en somme de deux sous-préschémas $X'_1$, $X'_2$ entraîne la décomposition de $\mathscr A'$ en produit direct des deux $\sh{O}_{Y'}$-Algèbres quasi-cohérentes $\mathscr A'_1$, $\mathscr A'_2$, images directes respectives de $\sh{O}_{X'_1}$ et $\sh{O}_{X'_2}$, de sorte que $Z'$ s'identifie à la somme $Z'_1\amalg Z'_2$, où $Z'_i=\Spec(\mathscr A'_i)$ et $g'(X'_i)\subset Z'_i$ pour $i=1,2$. Comme $X'_1$ est fini sur $Y'$, $g'_1=g'|X'_1$ est un isomorphisme de $X'_1$ sur $Z'_1$, puisque $f'|X'_1$ est affine; comme $X'_2\cap f'^{-1}(y')=\varnothing$, on voit que $g'$ est étale et radiciel en chaque point de $f'^{-1}(y')$. Le morphisme $Y'\to Y$ étant plat et localement de présentation finie, on déduit donc d'abord de (17.7.4) que $g$ est étale en tous les points de $X$ projections de points de $f'^{-1}(y')$, c'est-à-dire en tous les points de $f^{-1}(y)$ (I, 3.5.2). D'autre part, le morphisme
\[
  g'_{y'}:f'^{-1}(y')\longrightarrow h'^{-1}(y')
\]
déduit de $g'$ est radiciel; comme $k(y')=k(y)$, le morphisme $g_y:f^{-1}(y)\to h^{-1}(y)$ est aussi radiciel, autrement dit $g$ est radiciel en chaque point de $f^{-1}(y)$, ce qui achève de prouver la proposition.

L'énoncé suivant améliore de même (8.12.6):

\begin{corollary}[18.12.13]
\label{IV.18.12.13}
(« Main Theorem » de Zariski). Soient $Y$ un préschéma quasi-compact et quasi-séparé, $f:X\to Y$ un morphisme quasi-fini et séparé. Alors il existe une factorisation de $f$
\[
  X\xrightarrow{g}Z\xrightarrow{u}Y
\]
où $g$ est une immersion ouverte (nécessairement quasi-compacte) et $u$ un morphisme fini.
\end{corollary}

En effet, il résulte de (18.12.12) que $f$ est quasi-affine, donc quasi-projectif. On peut alors appliquer (8.12.8), où l'hypothèse de l'existence d'un $\sh{O}_Y$-Module ample peut en fait être remplacée par la seule hypothèse que $Y$ est quasi-compact et quasi-séparé: en effet, il résulte de (8.12.3) que l'existence de la factorisation annoncée dans (8.12.8) est une propriété locale sur $Y$, et qu'il suffit donc de la prouver lorsque $Y$ est affine.

\begin{remark}[18.12.14]
\label{IV.18.12.14}
On peut donner de (18.12.13) une démonstration analogue à celle de (18.12.12) n'utilisant pas (8.12.8) (mais utilisant (18.12.3), donc (18.5.11), qui lui-même se sert du « Main Theorem » sous sa forme locale (8.12.9)). Gardons en effet les notations de la démonstration de (18.12.12), et soit $\mathscr B$ la $\sh{O}_Y$-Algèbre quasi-cohérente, fermeture intégrale de $\sh{O}_Y$ dans $\mathscr A$ (II, 6.3.2). Si l'on pose $T=\Spec(\mathscr B)$, il suffit, en vertu de (8.12.3), de prouver que le $Y$-morphisme $g:X\to T$ correspondant à l'injection canonique $\mathscr B\to\mathscr A$ est une immersion, et pour cela il suffit de montrer (17.9.1) que $g$ est étale et radiciel. Avec les notations de (18.12.12), on peut supposer que $Y=\Spec(C)$ et $Y'=\Spec(C')$ sont affines, $C'$ étant une $C$-algèbre étale (17.3.2); donc $\mathscr A=\widetilde A$, où $A$ est une $C$-algèbre, $\mathscr A'=\widetilde{A'}$ avec $A'=A\otimes_C C'$, et $\mathscr B=\widetilde B$, où $B$ est la fermeture intégrale de $C$ dans $A$. L'algèbre $A'$ est isomorphe au produit $A'_1\times A'_2$, où $A'_1$ est une $C'$-algèbre finie. Il suffira de prouver que $B'=B\otimes_C C'$, qui

\oldpage[IV-4]{184}

s'identifie par platitude à une sous-$C'$-algèbre de $A'$, contient $A'_1$; en effet, $B'$ se décomposera alors en produit $A'_1\times A''_2$, où $A''_2$ est une sous-$C'$-algèbre de $A'_2$, et si l'on pose $T''_2=\Spec(A''_2)$, $T'=T\times_Y Y'$ sera somme de $Z'_1=\Spec(A'_1)$ et de $T''_2$, et $g'|X'_1$ sera un isomorphisme de $X'_1$ sur $Z'_1$, ce qui permettra de conclure comme dans (18.12.12).

Pour prouver que $B'$ contient $A'_1$, il suffit évidemment de prouver la proposition suivante, qui étend partiellement (6.14.4) lorsqu'on ne fait plus d'hypothèses noethériennes:
\end{remark}

\begin{proposition}[18.12.15]
\label{IV.18.12.15}
Soient $C$ un anneau, $C'$ une $C$-algèbre étale, $A$ une $C$-algèbre, $B$ la fermeture intégrale de $C$ dans $A$; on pose $A'=A\otimes_C C'$, $B'=B\otimes_C C'$, $B'$ s'identifiant à une sous-algèbre de $A'$; alors $B'$ est la fermeture intégrale de $C'$ dans $A'$.
\end{proposition}

En considérant la famille filtrante croissante des sous-$C$-algèbres de type fini de $A$ et raisonnant comme dans (6.14.4, II)), on peut d'abord supposer que $A$ est une $C$-algèbre de type fini, donc de la forme $E/\mathfrak J$, où $E=C[T_1,\ldots,T_r]$ et $\mathfrak J$ est un idéal de $E$; on a alors $A'=E'/\mathfrak JE'$, où $E'=C'[T_1,\ldots,T_r]$. Soit $(\mathfrak J_\lambda)$ la famille filtrante croissante des idéaux de type fini de $E$ contenus dans $\mathfrak J$, de sorte que $A$ est limite inductive des $E/\mathfrak J_\lambda$; si $B_\lambda$ est la fermeture intégrale de $C$ dans $E/\mathfrak J_\lambda$, $B$ est limite inductive des $B_\lambda$, comme le montre le raisonnement de (5.13.4). De même la fermeture intégrale de $C'$ dans $E'/\mathfrak JE'$ est limite inductive de la fermeture intégrale de $C'$ dans $E'/\mathfrak J_\lambda E'$; si on prouve que cette dernière est égale à $B_\lambda\otimes_C C'$, il en résultera que
\[
  B'=\varinjlim_\lambda(B_\lambda\otimes_C C')
\]
sera la fermeture intégrale de $C'$ dans $A'$. On est ainsi ramené à prouver la proposition lorsque $A$ est une $C$-algèbre de présentation finie.

Montrons ensuite qu'on peut se ramener au cas où $C$ est noethérien. En effet, $C$ est réunion filtrante de ses sous-$\mathbf Z$-algèbres de type fini $C_\alpha$, donc il résulte de (17.7.8) qu'il existe un indice $\alpha$ et une $C_\alpha$-algèbre étale $C'_\alpha$ tels que $C'=C'_\alpha\otimes_{C_\alpha}C$; en outre $C'$ est limite inductive des $C'_\beta=C'_\alpha\otimes_{C_\alpha}C_\beta$ pour $\beta\geq\alpha$. On peut en outre supposer, puisque $A$ est une $C$-algèbre de présentation finie, que $A=A_\alpha\otimes_{C_\alpha}C$, où $A_\alpha$ est une $C_\alpha$-algèbre de type fini, et $A$ est limite inductive des $A_\beta=A_\alpha\otimes_{C_\alpha}C_\beta$ pour $\beta\geq\alpha$ (8.9.1). Le raisonnement de (5.13.4) montre alors que $B$ est la limite inductive des fermetures intégrales $B_\beta$ des $C_\beta$ dans $A_\beta$. De même $A'$ est limite inductive des
\[
  A'_\beta=A_\beta\otimes_{C_\beta}C'_\beta=A_\alpha\otimes_{C_\alpha}C'_\beta
\]
pour $\beta\geq\alpha$, donc le même raisonnement que ci-dessus montre qu'il suffira de prouver que $B'_\alpha=B_\alpha\otimes_{C_\alpha}C'_\alpha$ est la fermeture intégrale de $C'_\alpha$ dans $A'_\alpha$.

Une fois ramené au cas où $C$ est noethérien, la proposition devient un cas particulier de (6.14.4). Il faut toutefois observer que, lorsque $C'$ est supposé étale sur $C$ noethérien, la démonstration de (6.14.4) n'exige plus le délicat théorème (6.14.1). En effet, les réductions successives de la preuve de (6.14.4) ramènent sans utiliser (6.14.1), au cas où $C$ est intègre, $A$ le corps des fractions de $C$. Comme alors $B$ est un anneau normal et que le morphisme $\Spec(B')\to\Spec(B)$ est étale, il résulte de (17.5.7) que $B'$ est un anneau normal; le raisonnement se termine en faisant appel seulement au lemme élémentaire (6.14.1.1), mais non à la partie difficile de la preuve de (6.14.1).

\begin{proposition}[18.12.16]
\label{IV.18.12.16}
Soient $g:Y\to S$ un morphisme quasi-compact, $f:X\to Y$ un morphisme séparé et quasi-fini. Pour tout $\sh{O}_Y$-Module inversible $\mathscr L$, ample relativement à $S$ (II, 4.6.1), le $\sh{O}_X$-Module $f^*(\mathscr L)$ est ample relativement à $S$.
\end{proposition}

En effet, en vertu de (18.12.12), le morphisme $f$ est quasi-affine, donc (II, 5.1.6) le $\sh{O}_X$-Module $\sh{O}_X$ est ample pour $f$. On en déduit (II, 4.6.13, (ii)) que
\[
  \sh{O}_X\otimes_{\sh{O}_X}f^*(\mathscr L^{\otimes n})=f^*(\mathscr L^{\otimes n})
\]
est ample relativement à $S$ pour un $n$ assez grand. Mais comme $f^*(\mathscr L^{\otimes n})=(f^*(\mathscr L))^{\otimes n}$, cela entraîne que $f^*(\mathscr L)$ est ample relativement à $S$ (II, 4.6.9, (i)).

\begin{corollary}[18.12.17]
\label{IV.18.12.17}
Soient $Z$ un schéma affine, $h:X\to Z$ un morphisme de type fini, $\mathscr L$ un $\sh{O}_X$-Module inversible. Avec les notations de (II, 4.5.2), les conditions de (II, 4.5.2) (qui équivalent au fait que $\mathscr L$ est ample) sont aussi équivalentes à chacune des suivantes:
\begin{enumerate}
  \item[$b'')$] $h$ est séparé, on a $G(\varepsilon)=X$ et le morphisme canonique $u:G(\varepsilon)\to\operatorname{Proj}(S)$ a ses fibres finies et discrètes.
  \item[$b''')$] On a $G(\varepsilon)=X$ et le morphisme canonique $u$ est radiciel.
\end{enumerate}
\end{corollary}

Si $Z=\Spec(A)$, l'anneau $S$ est canoniquement muni d'une structure de $A$-algèbre graduée, donc on a un morphisme structural $g:\operatorname{Proj}(S)\to Z$ qui est séparé (II, 2.4.2), et on a $h=g\circ u$. Comme tout morphisme radiciel est séparé (I, 1.8.7.1), $b''')$ entraîne que $h$ est séparé, donc entraîne $b'')$. Comme la condition $b)$ de (II, 4.5.2) entraîne évidemment $b''')$, il reste à voir que $b'')$ entraîne que $\mathscr L$ est ample. Puisque $h=g\circ u$ est de type fini par hypothèse et que $g$ est séparé, $u$ est aussi de type fini (I, 6.3.4), donc l'hypothèse $b'')$ entraîne que $u$ est quasi-fini et séparé (I, 5.5.1). Pour prouver que $\mathscr L$ est ample, nous allons appliquer (18.12.16). Posons $Y=\operatorname{Proj}(S)$. Comme $X$ est quasi-compact, il existe un entier $n>0$ et un nombre fini d'éléments $f_j\in S_n$ tels que les images réciproques $u^{-1}(D_+(f_j))$ recouvrent $X$; on peut donc considérer $u$ comme un morphisme quasi-fini et séparé de $X$ dans l'ouvert
\[
  Y'=\bigcup_j D_+(f_j)
\]
de $Y$. Considérons alors le $\sh{O}_{Y'}$-Module inversible
\[
  \mathscr L'=\sh{O}_Y(n)|Y'
\]
(II, 2.5.8); il est très ample relativement à $Z$ (II, 4.4.3), et $u^*(\mathscr L')$ n'est autre par définition que $\mathscr L^{\otimes n}$

\oldpage[IV-4]{185}

(II, 3.7.9, (ii)). En vertu de (18.12.16), $\mathscr L^{\otimes n}$ est donc ample relativement à $Z$ (ce qui équivaut à dire qu'il est ample (II, 4.6.6)), et par suite il en est de même de $\mathscr L$ (II, 4.5.6, (i)).

Nous laissons au lecteur le soin d'énoncer de même des conditions équivalentes à celles de (II, 4.6.3) pour les $\sh{O}_X$-Modules relativement amples.

\section{Immersions régulières et platitude normale}
\label{section:IV.19}

Le présent paragraphe est consacré, d'une part à une étude des immersions régulières (16.9.2) $Y\to X$ de préschémas, notamment dans le cas où $X$ et $Y$ sont plats sur un même préschéma $S$; d'autre part, il donne des compléments sur les suites $M$-régulières et la platitude normale, en généralisant notamment des résultats de platitude établis par H. Hironaka dans sa théorie de la résolution des singularités [35].

\subsection{Propriétés des immersions régulières.}
\EGAReaderTarget{ega4.c04.subsec.19-1}{subsection:IV.19.1}

\begin{proposition}[19.1.1]
\label{IV.19.1.1}
Soient $X$ un préschéma localement noethérien, $j:Y\to X$ une immersion, $y$ un point régulier de $Y$. Pour que $j$ soit une immersion régulière au point $y$, il faut et il suffit que $X$ soit régulier au point $j(y)$.
\end{proposition}

Compte tenu de (16.9.10), on est ramené à prouver le

\begin{corollary}[19.1.2]
\EGAReaderTarget{ega4.c04.cor.19-1-2}{IV.19.1.2}
Soient $A$ un anneau local noethérien, $\mathfrak J$ un idéal de $A$; on suppose que l'anneau quotient $B=A/\mathfrak J$ soit régulier. Pour que $A$ soit régulier, il faut et il suffit que l'idéal $\mathfrak J$ soit régulier.
\end{corollary}

En effet, si $A$ est régulier, on sait ($0$, 17.1.9) que $\mathfrak J$ est engendré par une partie d'un système régulier de paramètres de $A$, donc $\mathfrak J$ est régulier. Inversement, si $\mathfrak J$ est régulier, et si $(x_i)_{1\leq i\leq r}$ est une suite $A$-régulière engendrant $\mathfrak J$, $(x_i)$ fait partie d'un système de paramètres de $A$ ($0$, 16.4.1), donc on déduit de ($0$, 17.1.7) que $A$ est régulier.

\begin{definition}[19.1.3]
\EGAReaderTarget{ega4.c04.def.19-1-3}{IV.19.1.3}
Soient $X$ un préschéma, $Y$ un sous-préschéma de $X$, $U$ un ouvert de $X$ contenant $Y$ et tel que $Y$ soit défini par un Idéal $\mathscr J$ de $\sh{O}_U$; on suppose que $\mathscr J/\mathscr J^2$ soit un $(\sh{O}_U/\mathscr J)$-Module localement libre (ce qui est le cas si $Y$ est quasi-régulièrement immergé dans $X$). Pour tout $y\in Y$, on appelle \emph{codimension transversale de $Y$ dans $X$ au point $y$} et on note $\operatorname{codim}_y^*(Y,X)$ le rang du $(\sh{O}_{U,y}/\mathscr J_y)$-module libre $\mathscr J_y/\mathscr J_y^2$. Si $f:Z\to X$ est une immersion d'image $Y$, on appelle de même \emph{codimension transversale de $f$ au point $z\in Z$} la codimension transversale dans $X$ au point $f(z)$ du sous-préschéma $Y$.
\end{definition}

\begin{proposition}[19.1.4]
\label{IV.19.1.4}
Soient $X$ un préschéma localement noethérien, $Y$ un sous-préschéma de $X$ régulièrement immergé; alors, pour tout $y\in Y$, on a
\begin{equation}
  \operatorname{codim}_y^*(Y,X)=\operatorname{codim}_y(Y,X).
  \tag{19.1.4.1}\label{IV.19.1.4.1}
\end{equation}
\end{proposition}

En vertu de (5.1.3.2), $\operatorname{codim}_y(Y,X)$ est égal à la plus petite valeur des dimensions des anneaux locaux $A=\sh{O}_{X,z}$, où $z$ est point générique d'une composante

\oldpage[IV-4]{186}

irréductible de $Y$ contenant $y$; comme un tel point $z$ est contenu dans tout voisinage de $y$ dans $X$, $\mathscr J_z/\mathscr J_z^2$ est un $A$-module libre de rang $n=\operatorname{codim}_y^*(Y,X)$; tout revient à voir que $\dim(A)=n$. Or, puisque $z$ est point maximal du sous-préschéma $Y$ défini par $\mathscr J$, on a $\dim(\sh{O}_{X,z}/\mathscr J_z)=0$, donc $\mathscr J_z$ est un \emph{idéal de définition} de l'anneau local noethérien $A$; en outre, l'hypothèse que $\mathscr J$ est quasi-régulier entraîne que $\mathscr J_z^m/\mathscr J_z^{m+1}$ est un $(A/\mathscr J_z)$-module libre de rang $\binom{m+n-1}{n-1}$; si $r$ est la longueur de l'anneau artinien $A/\mathscr J_z$, la longueur de $\mathscr J_z^m/\mathscr J_z^{m+1}$ est donc $r\binom{m+n-1}{n-1}$; le polynôme de Hilbert-Samuel de $A$ pour la filtration $\mathscr J_z$-préadique est donc bien de degré $n$, ce qui prouve la proposition ($0$, 16.2.3).

En vertu de (19.1.4), on dira désormais « codimension » au lieu de « codimension transversale » lorsqu'il s'agira d'un sous-préschéma $Y$ régulièrement immergé dans $X$ (\emph{même lorsque $X$ n'est pas localement noethérien}).

\begin{proposition}[19.1.5]
\label{IV.19.1.5}
\begin{enumerate}
  \item[(i)] Pour qu'une immersion $j:Y\to X$ soit ouverte, il faut et il suffit qu'elle soit régulière et de codimension $0$ en tout point.
  \item[(ii)] Soient $f:Y\to X$ une immersion régulière (resp. quasi-régulière), $g:X'\to X$ un morphisme plat, $Y'=Y\times_X X'$; alors le morphisme $f'=f_{(X')}:Y'\to X'$ est une immersion régulière (resp. quasi-régulière), et pour tout $y'\in Y'$, si $y$ est la projection de $y'$ dans $Y$, la codimension de $f'$ au point $y'$ est égale à la codimension de $f$ au point $y$. Inversement, si $f$ est une immersion et $g:X'\to X$ un morphisme fidèlement plat et quasi-compact, alors, si $f'$ est une immersion quasi-régulière, il en est de même de $f$.
  \item[(iii)] Si $f:Y\to X$ et $g:Z\to Y$ sont des immersions régulières, alors $f\circ g:Z\to X$ est une immersion régulière, et pour tout $z\in Z$, la codimension de $f\circ g$ au point $z$ est la somme de la codimension de $g$ au point $z$ et de la codimension de $f$ au point $g(z)$. De plus, la suite d'homomorphismes canoniques (16.2.7.1)
\begin{equation}
  0\longrightarrow g^*(\sh{N}_{Y/X})\longrightarrow\sh{N}_{Z/X}\longrightarrow\sh{N}_{Z/Y}\longrightarrow0
  \tag{19.1.5.1}\label{IV.19.1.5.1}
\end{equation}
est exacte, et pour tout $z\in Z$, il existe un voisinage de $z$ dans $Z$ dans lequel les restrictions des homomorphismes de cette suite forment une suite scindée.
  \item[(iv)] Soient $X$ un préschéma localement noethérien, $f:Y\to X$ et $g:Z\to Y$ deux immersions, $z$ un point de $Z$. Les conditions suivantes sont équivalentes:
  \begin{enumerate}
    \item[a)] $g$ est une immersion régulière au point $z$ et $f$ une immersion régulière au point $g(z)$.
    \item[b)] $g$ et $f\circ g$ sont des immersions régulières au point $z$.
    \item[c)] $f\circ g$ est une immersion régulière au point $z$, $f$ est une immersion régulière au point $g(z)$, et la suite d'homomorphismes canoniques (19.1.5.1), restreinte à un voisinage ouvert assez petit de $z$ dans $Z$, est exacte et scindée.
  \end{enumerate}
  \item[(v)] Soient $X$ un préschéma localement noethérien, $f:Y\to X$, $g:Z\to Y$ deux immersions, $z$ un point de $Z$; supposons que $y=g(z)$ soit un point régulier de $Y$ et $x=f(g(z))$ un point régulier

\oldpage[IV-4]{187}

de $X$ (ce qui entraîne, par (19.1.1), que $f$ est une immersion régulière au point $g(z)$); alors, pour que $f\circ g$ soit une immersion régulière au point $z$, il faut et il suffit que $g$ soit une immersion régulière au point $z$.
\end{enumerate}
\end{proposition}

L'assertion (i) est une autre manière d'exprimer (16.1.10). Pour démontrer la première assertion de (ii), on peut se borner au cas où $Y$ est un sous-préschéma fermé de $X$ défini par un Idéal régulier (resp. quasi-régulier) $\mathscr J$; alors $Y'$ est le sous-préschéma fermé de $X'$ défini par l'Idéal $g^*(\mathscr J)\sh{O}_{X'}$ (I, 4.4.5) qui s'identifie ici à $g^*(\mathscr J)$ puisque $g$ est plat; en remplaçant $X$ par un ouvert affine assez petit, on peut supposer que $\mathscr J$ admet une suite régulière (resp. quasi-régulière) de générateurs $f_i$ (sections de $\sh{O}_X$ au-dessus de $X$). Si la suite $(f_i)$ est régulière, il en est de même de la suite $(f_i\otimes1)$ (qui engendre $g^*(\mathscr J)\sh{O}_{X'}$) en vertu de ($0$, 15.2.5); le résultat analogue pour les suites quasi-régulières découle aussitôt de la platitude de $g$ et du critère (16.9.4). De même, si $g$ est fidèlement plat et quasi-compact, et si $g^*(\mathscr J)$ est quasi-régulier, $\mathscr J$ est de type fini en vertu de (2.5.2); par platitude, on a $g^*(\mathscr J)/g^*(\mathscr J)^2=g^*(\mathscr J/\mathscr J^2)$, donc, puisque $g^*(\mathscr J/\mathscr J^2)$ est localement libre par hypothèse (16.9.4), il en est de même de $\mathscr J/\mathscr J^2$ (2.5.2). Enfin, la condition (iii) de (16.9.4) relative à $g^*(\mathscr J)$ entraîne la même condition pour $\mathscr J$ en vertu de la fidèle platitude de $g$ (2.2.7). On conclut donc de (16.9.4) que $\mathscr J$ est quasi-régulier.

Pour prouver (iii), on peut de même supposer que $Y$ et $Z$ sont des sous-préschémas fermés de $X$, définis respectivement par des Idéaux $\mathscr J$, $\mathscr K$ de $\sh{O}_X$ tels que $\mathscr J\subset\mathscr K$; de plus, on peut supposer qu'il existe des sections $f_i$ ($1\leq i\leq m$) de $\mathscr J$ au-dessus de $X$ engendrant $\mathscr J$ et telles que pour $1\leq i\leq m$, l'endomorphisme de $\sh{O}_X/(\sum_{j=1}^{i-1}f_j\sh{O}_X)$ défini par $f_i$ soit injectif, et des sections $\overline f_i$ ($m+1\leq i\leq n$) de $\mathscr K/\mathscr J$ au-dessus de $X$ engendrant $\mathscr K/\mathscr J$ et telles que pour $m+1\leq i\leq n$, l'endomorphisme de
\[
  (\sh{O}_X/\mathscr J)\big/\Bigl(\sum_{j=m+1}^{i-1}\overline f_j(\sh{O}_X/\mathscr J)\Bigr)
\]
défini par $\overline f_i$ soit injectif. Pour tout $z\in Z$, soit $U$ un voisinage ouvert de $z$ dans $X$ tel qu'il existe une section $f_i$ de $\mathscr K$ au-dessus de $U$, dont $\overline f_i|U$ soit la classe dans $\Gamma(U,\mathscr K/\mathscr J)$ pour $m+1\leq i\leq n$; alors
\[
  (\sh{O}_X/\mathscr J)\big/\Bigl(\sum_{j=m+1}^{i-1}\overline f_j(\sh{O}_X/\mathscr J)\Bigr),
\]
restreint à $U$ est isomorphe à
\[
  \sh{O}_U\big/\Bigl((\mathscr J|U)+\sum_{j=m+1}^{i-1}f_j\sh{O}_U\Bigr)
  =\sh{O}_U\big/\Bigl(\sum_{j=1}^{i-1}f_j\sh{O}_U\Bigr),
\]
et on voit donc que la suite $(f_i)_{1\leq i\leq n}$ est régulière dans $\sh{O}_U$; comme elle engendre $\mathscr K|U$, cela prouve la première assertion de (iii). Pour prouver la dernière assertion de (iii), on peut se borner (avec les mêmes notations) au cas où les images $t_i$ des $f_i$ dans $\mathscr K/\mathscr K^2$ (pour $1\leq i\leq n$) forment une base de cet $(\sh{O}_X/\mathscr K)$-Module (16.9.5); pour la même raison, on peut supposer que les $t_i$ tels que $1\leq i\leq m$ sont les images canoniques des éléments d'une base du $(\sh{O}_X/\mathscr K)$-Module $g^*(\mathscr J/\mathscr J^2)$, et que les images canoniques $\overline t_i$ des $t_i$ dans $(\mathscr K/\mathscr J)/(\mathscr K/\mathscr J)^2$ pour $m+1\leq i\leq n$ forment une base de cet $(\sh{O}_X/\mathscr K)$-Module. Cela achève de prouver (iii).

Passons à la démonstration de (iv). Le fait que $a)$ implique $c)$ résulte de (iii). Prouvons que $c)$ implique $a)$. Posons encore $A=\sh{O}_{X,f(g(z))}$, de sorte que $\sh{O}_{Y,g(z)}=A/\mathfrak J$,

\oldpage[IV-4]{188}

$\sh{O}_{Z,z}=A/\mathfrak R$ avec $\mathfrak J\subset\mathfrak R$. Par hypothèse, $\mathfrak J$ et $\mathfrak R$ sont des idéaux réguliers de $A$ et on a une suite exacte scindée
\[
  0\longrightarrow\mathfrak J/\mathfrak J\mathfrak R\longrightarrow\mathfrak R/\mathfrak R^2\longrightarrow(\mathfrak R/\mathfrak J)/(\mathfrak R/\mathfrak J)^2\longrightarrow0
\]
(cf. (16.2.7)). Cela entraîne qu'il existe des éléments $f_j$ ($1\leq j\leq r+s$) de $\mathfrak R$ dont les images $f'_j$ dans $\mathfrak R/\mathfrak R^2$ forment une base de cet $(A/\mathfrak R)$-module, et tels en outre que les $f_j$ d'indice $j\leq r$ appartiennent à $\mathfrak J$ et sont tels que leurs images dans $\mathfrak J/\mathfrak J\mathfrak R$ forment une base de cet $(A/\mathfrak R)$-module. Les $f_j$ pour $1\leq j\leq r+s$ engendrent donc $\mathfrak R$, et les $f_j$ pour $1\leq j\leq r$ engendrent $\mathfrak J$ puisque $A$ est noethérien (Bourbaki, \emph{Alg. comm.}, chap. II, \textsection~3, n\textsuperscript{o}~2, prop. 5); comme $\mathfrak R$ est un idéal régulier, il résulte de (16.9.5) que la suite $(f_j)_{1\leq j\leq r+s}$ est régulière. Par définition, les images de $f_{r+1},\ldots,f_{r+s}$ dans $\mathfrak R/\mathfrak J$ forment donc une suite régulière dans cet $(A/\mathfrak J)$-module, ce qui achève de prouver que $c)$ entraîne $a)$ dans (iv).

Enfin, compte tenu de (16.9.10), il est immédiat que (v), ainsi que l'équivalence de $a)$ et $b)$ dans (iv), résultent du

\begin{corollary}[19.1.6]
\label{IV.19.1.6}
Soient $A$ un anneau local noethérien, $\mathfrak J$, $\mathfrak R$ deux idéaux de $A$ contenus dans son idéal maximal $\mathfrak m$ et tels que $\mathfrak J\subset\mathfrak R$, $A'=A/\mathfrak J$, $\mathfrak R'=\mathfrak R/\mathfrak J$. Alors:
\begin{enumerate}
  \item[(i)] Si $\mathfrak J$ et $\mathfrak R'$ sont des idéaux réguliers, il en est de même de $\mathfrak R$.
  \item[(ii)] Si $\mathfrak R$ et $\mathfrak R'$ sont des idéaux réguliers, il en est de même de $\mathfrak J$.
  \item[(iii)] Supposons les anneaux $A$ et $A'=A/\mathfrak J$ réguliers. Pour que l'idéal $\mathfrak R$ soit régulier, il faut et il suffit que $\mathfrak R'$ le soit.
\end{enumerate}
\end{corollary}

L'assertion (i) est un cas particulier de (19.1.5, (iii)). Pour démontrer (ii), considérons l'homomorphisme surjectif canonique $u:\mathfrak R/\mathfrak R^2\to\mathfrak R'/\mathfrak R'^2=\mathfrak R/(\mathfrak R^2+\mathfrak J)$. Comme par hypothèse $\mathfrak R/\mathfrak R^2$ et $\mathfrak R'/\mathfrak R'^2$ sont des $(A/\mathfrak R)$-modules libres, dont nous désignerons les rangs respectifs par $p+q$ et $q$, le noyau de $u$ est un $(A/\mathfrak R)$-module \emph{projectif} (Bourbaki, \emph{Alg.}, chap. II, 3\textsuperscript{e} éd., \textsection~2, n\textsuperscript{o}~2, prop. 4), donc \emph{libre} de rang $p$ puisque $A/\mathfrak R$ est un anneau local noethérien ($0_{\mathrm{III}}$, 10.1.3); autrement dit, il existe une base de $\mathfrak R/\mathfrak R^2$ dont les $p$ premiers éléments forment une base de $(\mathfrak R^2+\mathfrak J)/\mathfrak R^2$. Cela signifie encore (en vertu du lemme de Nakayama) qu'il existe un système de générateurs $(f_i)_{1\leq i\leq p+q}$ de $\mathfrak R$, formant une suite régulière dans $A$, telle que la suite $(f_i)_{1\leq i\leq p}$ soit formée d'éléments de $\mathfrak J$; il s'agit de montrer que cette dernière suite engendre $\mathfrak J$. Pour cela, désignons par $\mathfrak J'\subset\mathfrak J$ l'idéal qu'elle engendre; en considérant l'anneau $A/\mathfrak J'$, on voit qu'on est ramené à prouver le lemme suivant:

\begin{lemma}[19.1.6.1]
\label{IV.19.1.6.1}
Soient $B$ un anneau local noethérien d'idéal maximal $\mathfrak n$, $(g_j)_{1\leq j\leq q}$ une suite régulière dans $B$ d'éléments de $\mathfrak n$, $\mathfrak b$ l'idéal qu'elle engendre, $\mathfrak c\subset\mathfrak b$ un idéal tel que les images canoniques des $g_j$ dans $B/\mathfrak c$ forment encore une suite régulière dans cet anneau. Alors on a $\mathfrak c=0$.
\end{lemma}

Le lemme étant trivial pour $q=0$, raisonnons par récurrence sur $q$. L'hypothèse de récurrence, appliquée à l'anneau $B/g_1B$ et aux images canoniques des $g_j$ ($2\leq j\leq q$)

\oldpage[IV-4]{189}

et de $\mathfrak c$ dans cet anneau quotient, montre que $\mathfrak c\subset g_1B$. Soit $\mathfrak d$ l'idéal de $B$ formé des $x\in B$ tels que $g_1x\in\mathfrak c$; on a donc $\mathfrak c=g_1\mathfrak d$; mais comme par hypothèse $g_1$ est régulier dans $B/\mathfrak c$, on a nécessairement $\mathfrak d=\mathfrak c$, donc $\mathfrak c=g_1\mathfrak c$. Comme $\mathfrak c$ est de type fini et que $g_1$ est contenu dans le radical de $B$, le lemme de Nakayama montre que $\mathfrak c=0$.

Démontrons enfin l'assertion (iii) de (19.1.6). En vertu de (i) et de (19.1.2), il suffit de voir que si $\mathfrak R$ est régulier, il en est de même de $\mathfrak R'$. Notons que $\mathfrak J$ est régulier (19.1.2), et raisonnons par récurrence sur le rang $q$ du $(A/\mathfrak J)$-module libre $\mathfrak J/\mathfrak J^2$. Si $q=0$, on a $\mathfrak J=0$ par le lemme de Nakayama et l'assertion est triviale. Comme $A$ et $A/\mathfrak J$ sont réguliers, on sait ($0$, 17.1.9) que $\mathfrak J$ est engendré par une partie à $q$ éléments d'un système régulier de paramètres de $A$, donc, si $f$ est un des éléments de ce système de générateurs de $\mathfrak J$, $A_1=A/fA$ est régulier et $f\notin\mathfrak m^2$ ($0$, 17.1.8). Soient $\mathfrak J_1$, $\mathfrak R_1$ les images canoniques de $\mathfrak J$, $\mathfrak R$ dans $A_1$; on a $\mathfrak J_1\subset\mathfrak R_1$, $A_1/\mathfrak J_1=A/\mathfrak J$, donc $A_1/\mathfrak J_1$ est régulier et par suite $\mathfrak J_1$ est un idéal régulier (19.1.2). Montrons que $\mathfrak R_1$ est un idéal régulier dans $A_1$; comme $\mathfrak R$ est un idéal régulier dans $A$, il suffit de montrer que $f$ fait partie d'un système régulier de générateurs de $\mathfrak R$, et pour cela (16.9.5) il suffit de montrer que l'image de $f$ dans $\mathfrak R/(\mathfrak R^2+\mathfrak m\mathfrak R)=\mathfrak R/\mathfrak m\mathfrak R$ fait partie d'une base de cet $(A/\mathfrak m)$-espace vectoriel, autrement dit que $f\notin\mathfrak m\mathfrak R$, ce qui résulte en effet de $f\notin\mathfrak m^2$. Alors, comme $\mathfrak J_1/\mathfrak J_1^2$ est un $(A_1/\mathfrak J_1)$-module libre de rang $q-1$, l'hypothèse de récurrence montre que $\mathfrak R'_1=\mathfrak R_1/\mathfrak J_1$ est un idéal régulier de $A_1/\mathfrak J_1$, et comme $\mathfrak R'$ s'identifie à $\mathfrak R'_1$ dans l'isomorphisme canonique entre $A/\mathfrak J$ et $A_1/\mathfrak J_1$, $\mathfrak R'$ est régulier. C.Q.F.D.

\begin{remarks}[19.1.7]
\label{IV.19.1.7}
\begin{enumerate}
  \item[(i)] Nous ignorons si le composé de deux immersions quasi-régulières est quasi-régulière.
  \item[(ii)] Dans un espace annelé en anneaux locaux $X$, pour tout Idéal $\mathscr J$ de $\sh{O}_X$, $\sh{O}_X/\mathscr J$ est de type fini et a donc un support fermé dans $X$ ($0_{\mathrm I}$, 5.2.2). On peut par suite définir comme dans (I, 4.1.3 et 4.2.1) les notions de sous-espace annelé (localement fermé) de $X$ défini par un Idéal $\mathscr J$ d'un faisceau $\sh{O}_U$, où $U$ est un ouvert de $X$, et la notion d'immersion. Les définitions d'immersions quasi-régulières et régulières, de codimension transversale s'appliquent alors sans modification, et les résultats de (19.1.4, (i) à (iv)) sont encore valables, en supposant dans la seconde assertion de (i) et dans (iv) que le faisceau $\sh{O}_X$ soit cohérent et les anneaux locaux $\sh{O}_{X,x}$ noethériens; quant à (ii), il faut \emph{définir} $Y'$ comme le sous-espace annelé défini par le $\sh{O}_{X'}$-Idéal $g^*(\mathscr J)$; les démonstrations sont inchangées.
\end{enumerate}
\end{remarks}

\begin{env}[19.1.8]
\label{IV.19.1.8}
Nous avons déjà signalé (notamment dans (16.9.6, (ii))) que dans les anneaux non noethériens, la notion de « suite régulière » d'éléments de l'anneau ne possède pas les propriétés qui la rendraient utilisable. De récentes recherches semblent montrer que la notion qui doit dans ce cas se substituer à celle de suite régulière est celle de suite $\mathbf f=(f_1,\ldots,f_n)$ ayant la propriété que le complexe de l'algèbre extérieure $\mathrm K_0(\mathbf f,M)$ (III, 1.1.3) ait \emph{son homologie nulle en degrés $>0$} (cf. (III, 1.1.4)); voir en particulier A. Grothendieck, Séminaire de géométrie algébrique, 1966, à paraître prochainement.
\end{env}

\subsection{Immersions transversalement régulières.}
\EGAReaderTarget{ega4.c04.subsec.19-2}{subsection:IV.19.2}

\oldpage[IV-4]{190}

\begin{definition}[19.2.1]
\EGAReaderTarget{ega4.c04.def.19-2-1}{IV.19.2.1}
Soient $u:X\to S$ un morphisme de préschémas, $\mathscr F$ un $\sh{O}_X$-Module quasi-cohérent. On dit qu'une suite $(f_i)_{1\leq i\leq n}$ de sections de $\mathscr F$ au-dessus de $X$ est \emph{transversalement $\mathscr F$-régulière relativement à $S$} (ou relativement à $u$) si la suite $(f_i)$ est $\mathscr F$-régulière et telle que les $\sh{O}_X$-Modules
\[
  \mathscr F\big/\Bigl(\sum_{j\leq i}f_j\mathscr F\Bigr)
\]
soient $S$-plats pour $0\leq i\leq n$. On dit qu'un Idéal quasi-cohérent $\mathscr J$ de $\sh{O}_X$ est \emph{transversalement régulier relativement à $S$} (ou relativement à $u$) en un point $x\in\operatorname{Supp}(\sh{O}_X/\mathscr J)$ s'il existe un voisinage ouvert $U$ de $x$ et une suite finie $(f_i)$ de sections de $\mathscr J$ au-dessus de $U$ qui est transversalement $\sh{O}_U$-régulière relativement à $S$ et qui engendre $\mathscr J|U$.

Dans une $A$-algèbre $B$, on dit qu'un idéal $\mathfrak J$ est \emph{transversalement régulier relativement à $A$} si $\widetilde{\mathfrak J}$ est transversalement régulier relativement à $S=\Spec(A)$ en tout point de $V(\mathfrak J)$ dans $\Spec(B)$.
\end{definition}

\begin{definition}[19.2.2]
\EGAReaderTarget{ega4.c04.def.19-2-2}{IV.19.2.2}
Soient $S$ un préschéma, $X$, $Y$ deux $S$-préschémas, $u:Y\to X$ un $S$-morphisme qui est une immersion. On dit que $u$ est une \emph{immersion transversalement régulière relativement à $S$ en un point $y\in Y$} s'il existe un voisinage $V$ de $u(y)$ dans $X$ tel que le sous-préschéma induit sur $u(Y)\cap V$ par le sous-préschéma de $X$ associé à $u$ soit défini par un Idéal de $\sh{O}_U$ transversalement régulier relativement à $S$ au point $u(y)$.

Il est clair que l'ensemble des points de $Y$ où une immersion est transversalement régulière relativement à $S$ est \emph{ouvert} dans $Y$; s'il est égal à $Y$, on dit simplement que $u$ est une \emph{immersion transversalement régulière relativement à $S$}.

Si $S\to T$ est un morphisme \emph{plat}, alors, si $\mathscr J$ est un Idéal de $\sh{O}_X$ transversalement régulier relativement à $S$ en $x\in X$, il l'est aussi relativement à $T$ ($0_{\mathrm I}$, 6.2.1). Une $S$-immersion $Y\to X$ transversalement régulière en un point $y\in Y$ relativement à $S$ est donc aussi une $T$-immersion transversalement régulière relativement à $T$ en ce point.
\end{definition}

\begin{proposition}[19.2.3]
\label{IV.19.2.3}
Soient $S$ un préschéma, $X$, $Y$ deux $S$-préschémas, $u:Y\to X$ un $S$-morphisme qui est une immersion transversalement régulière relativement à $S$ en un point $y\in Y$. Alors, pour tout changement de base $S'\to S$, si l'on pose $X'=X\times_S S'$, $Y'=Y\times_S S'$, l'immersion $u'=u_{(S')}:Y'\to X'$ est transversalement régulière relativement à $S'$ en tout point $y'\in Y'$ au-dessus de $y$.
\end{proposition}

Cela résulte aussitôt de la définition et de ($0$, 15.1.15).

\begin{proposition}[19.2.4]
\label{IV.19.2.4}
Soient $S$ un préschéma, $X$, $Y$ deux $S$-préschémas, $u:Y\to X$ un $S$-morphisme qui est une immersion. On suppose, soit que $S$ et $X$ sont localement noethériens, soit que les morphismes structuraux $g:X\to S$, $h:Y\to S$ soient localement de présentation finie. Soient $y$ un point de $Y$, $s$ son image dans $S$. Les conditions suivantes sont équivalentes:
\begin{enumerate}
  \item[a)] $u$ est transversalement régulière relativement à $S$ au point $y$.
  \item[b)] Les morphismes $g$ et $h$ sont plats au voisinage des points $u(y)$ et $y$ respectivement, et si l'on pose $X_s=g^{-1}(s)$, $Y_s=h^{-1}(s)$, l'immersion $u_s:Y_s\to X_s$ est régulière au point $y$.
  \item[b$'$)] Les morphismes $g$ et $h$ sont plats au voisinage des points $u(y)$ et $y$ respectivement, et, pour tout changement de base $S'\to S$, si l'on pose $X'=X\times_S S'$, $Y'=Y\times_S S'$, l'immersion $u'=u_{(S')}:Y'\to X'$ est régulière en tout point de $Y'$ au-dessus de $y$.
  \item[c)] Le morphisme $h$ est plat au voisinage du point $y$, et l'immersion $u$ est régulière au voisinage du point $y$.

\oldpage[IV-4]{191}

  \item[c$'$)] Le morphisme $h$ est plat au voisinage du point $y$, et l'immersion $u$ est quasi-régulière au voisinage du point $y$.
\end{enumerate}
\end{proposition}

On a déjà prouvé dans (19.2.3) que $a)$ entraîne $b')$ sans hypothèse de finitude, et il est trivial que $b')$ entraîne $b)$. Montrons que $b)$ entraîne $c)$. On peut pour cela (la question étant locale sur $X$ et sur $S$) supposer que $Y$ est un sous-préschéma fermé de $X$ défini par un Idéal quasi-cohérent $\mathscr J$ de $\sh{O}_X$, et $u$ l'injection canonique. L'hypothèse que $\sh{O}_X/\mathscr J$ est $g$-plat entraîne que la suite
\[
  0\longrightarrow\mathscr J\otimes_{\sh{O}_s}\mathbf k(s)\longrightarrow\sh{O}_X\otimes_{\sh{O}_s}\mathbf k(s)\longrightarrow(\sh{O}_X/\mathscr J)\otimes_{\sh{O}_s}\mathbf k(s)\longrightarrow0
\]
est exacte ($0_{\mathrm I}$, 6.1.2), donc $Y_s$ est le sous-préschéma fermé de $X_s$ défini par l'Idéal $\mathscr J_s$ de $\sh{O}_{X_s}=\sh{O}_X\otimes_{\sh{O}_s}\mathbf k(s)$ qui s'identifie à $\mathscr J\otimes_{\sh{O}_s}\mathbf k(s)=\mathscr J/\mathfrak m_s\mathscr J$. Considérons une suite finie $(f_i)$ de sections de $\mathscr J$ au-dessus d'un voisinage de $y$ dans $X$, dont les images dans
\[
  \mathscr J_y/(\mathfrak m_s\mathscr J_y+\mathscr J_y^2)
\]
forment une base de ce $(\sh{O}_{X,y}/(\mathfrak m_s\sh{O}_{X,y}+\mathscr J_y))$-module (qui est libre en vertu de l'hypothèse $b)$). Puisque $X_s$ est localement noethérien, il résulte de (16.9.11), (16.9.5) et de l'hypothèse $b)$ que les images $(f_i)_s=f_i\otimes1$, sections de $\mathscr J_s=\mathscr J\otimes_{\sh{O}_s}\mathbf k(s)$ au-dessus d'un voisinage de $y$ dans $X_s$, forment une suite régulière engendrant la restriction de $\mathscr J_s$ à ce voisinage. Comme $\mathfrak m_s\sh{O}_{X,y}$ est contenu dans le radical de $\sh{O}_{X,y}$, et que dans les deux cas envisagés $\mathscr J_y$ est un idéal de type fini de $\sh{O}_{X,y}$, il résulte du lemme de Nakayama que les $(f_i)_y$ engendrent aussi $\mathscr J_y$; comme $\mathscr J$ est un Idéal de type fini de $\sh{O}_X$ dans les deux cas envisagés, il y a un voisinage $U$ de $y$ dans $X$ tel que les $f_i|U$ engendrent $\mathscr J|U$ ($0_{\mathrm I}$, 5.2.2). Le fait que $b)$ entraîne $c)$ est alors conséquence de ($0$, 15.1.16) lorsque $S$ et $X$ sont localement noethériens, et de (11.3.8) lorsque $g$ et $h$ sont localement de présentation finie; dans ce dernier cas, (11.3.8) prouve aussi l'équivalence de $c)$ et $c')$ (qui est triviale lorsque $X$ et $S$ sont localement noethériens ($0$, 15.1.11)). Enfin, si $c)$ est vérifiée, pour prouver $a)$, on peut encore se borner au cas où $Y$ est un sous-préschéma fermé de $X$ défini par $\mathscr J$, et par hypothèse, il y a une suite régulière $(f_i)$ de sections de $\mathscr J$ au-dessus d'un voisinage de $y$ dans $X$ telle que les $f_i$ engendrent $\mathscr J|U$ et que $h$ soit plat dans $U$; pour voir que $c)$ entraîne $a)$, il suffit encore d'appliquer ($0$, 15.1.6) lorsque $S$ et $X$ sont localement noethériens, et (11.3.8) lorsque $g$ et $h$ sont localement de présentation finie.

\begin{remarks}[19.2.5]
\label{IV.19.2.5}
\begin{enumerate}
  \item[(i)] Lorsque dans la condition $b)$, on supprime l'hypothèse que $h$ est plat au voisinage de $y$, la condition ainsi modifiée n'entraîne plus $a)$. Il suffit de prendre pour $S$ le spectre d'un anneau local qui n'est pas un corps, pour $s$ le point fermé de $S$, et $Y=X_s$.
  \item[(ii)] Supposons que $Y$ soit un sous-préschéma fermé de $X$ défini par un Idéal quasi-cohérent $\mathscr J$ de $\sh{O}_X$. On a prouvé au cours de la démonstration de (19.2.4) que, lorsque les conditions équivalentes de cette proposition sont remplies, alors, pour tout système $(f_i)$ de sections de $\mathscr J$ au-dessus d'un voisinage de $y$ dans $X$, tel que les images des $f_i$ dans $\mathscr J_y/(\mathfrak m_s\mathscr J_y+\mathscr J_y^2)$ forment une base de ce $(\sh{O}_{X,y}/(\mathfrak m_s\sh{O}_{X,y}+\mathscr J_y))$-module, il existe un voisinage ouvert $U$ de $y$ dans $X$ tel que les $f_i|U$ vérifient les conditions de la définition (19.2.1).
\end{enumerate}
\end{remarks}

\oldpage[IV-4]{192}

\begin{corollary}[19.2.6]
\label{IV.19.2.6}
Supposons, soit que $S$ et $X$ soient localement noethériens, soit que $g$ et $h$ soient localement de présentation finie. Supposons en outre que les fibres $X_s$ et $Y_s$ soient des préschémas réguliers aux points $u(y)$ et $y$ respectivement, et que les morphismes $g$ et $h$ soient plats au voisinage de $u(y)$ et $y$ respectivement. Alors l'immersion $Y\to X$ est transversalement régulière au point $x$.
\end{corollary}

En effet, il résulte de (19.1.1) que l'immersion $Y_s\to X_s$ est régulière au point $x$, et il suffit d'appliquer (19.2.4, $b)$).

\begin{proposition}[19.2.7]
\label{IV.19.2.7}
\begin{enumerate}
  \item[(i)] Toute $S$-immersion ouverte est transversalement régulière relativement à $S$.
  \item[(ii)] Soient $f:Y\to X$ une $S$-immersion, $g:S'\to S$ un morphisme et posons $X'=X_{(S')}$, $Y'=Y_{(S')}$, $f'=f_{(S')}:Y'\to X'$. Si $f$ est une immersion transversalement régulière relativement à $S$, $f'$ est une immersion transversalement régulière relativement à $S'$. La réciproque est vraie si $g$ est fidèlement plat, et si de plus $X$ et $Y$ sont localement de présentation finie sur $S$.
  \item[(iii)] Si $f:Y\to X$ et $g:Z\to Y$ sont des immersions transversalement régulières relativement à $S$, il en est de même de $f\circ g:Z\to X$.
  \item[(iv)] Soient $X$ un $S$-préschéma, $f:Y\to X$, $g:Z\to Y$ deux $S$-immersions. On suppose, soit que $S$ et $X$ soient localement noethériens, soit que $X$, $Y$ et $Z$ soient localement de présentation finie sur $S$. Alors, si $g$ et $f\circ g$ sont transversalement régulières au point $z\in Z$ relativement à $S$, $f$ est transversalement régulière au point $g(z)$ relativement à $S$.
  \item[(v)] Sous les conditions générales de (iv), supposons que $X$ et $Y$ soient $S$-plats, $Y_s$ régulier au point $g(z)$ et $X_s$ régulier au point $g(z)$; alors, si $f\circ g$ est transversalement régulière au point $z$, il en est de même de $g$.
\end{enumerate}
\end{proposition}

L'assertion (i) est immédiate, et la première assertion de (ii) résulte aussitôt de (19.2.3). Pour démontrer la seconde assertion de (ii), appliquons (19.2.4): $X'$ et $Y'$ étant localement de présentation finie sur $S'$, il résulte d'abord de ce critère que $X'$ et $Y'$ sont plats sur $S'$, donc $X$ et $Y$ sont plats sur $S$ (2.5.1); d'autre part, pour tout $s\in S$, si $s'\in S'$ est un point au-dessus de $s$, $Y'_{s'}\to X'_{s'}$ est par hypothèse une immersion régulière, et on déduit donc de (19.1.5, (ii)) qu'il en est de même de $Y_s\to X_s$, puisque ces préschémas sont localement noethériens et qu'il revient alors au même de dire que l'immersion $Y_s\to X_s$ est régulière ou quasi-régulière. Pour prouver (iii), on se ramène au cas où $Z$ est un sous-préschéma fermé de $Y$ et $Y$ un sous-préschéma fermé de $X$, et on forme comme dans (19.1.5, (iii)) un système de générateurs de l'Idéal de $\sh{O}_X$ définissant $Z$. Pour prouver (iv), on remarque que l'hypothèse jointe à (19.2.4) montre d'abord que $X$, $Y$ et $Z$ sont $S$-plats dans un voisinage de $z$, et il suffit en outre, si $s$ est l'image de $z$ dans $S$, de prouver que l'immersion $f_s:Y_s\to X_s$ est régulière au point $z$, sachant qu'il en est de même de $g_s:Z_s\to Y_s$ et de $f_s\circ g_s$; mais cela résulte de (19.1.5, (iv)). Enfin, pour (v), on raisonne de la même manière que pour (iv); l'hypothèse sur $f\circ g$ implique déjà par (19.2.4, $b)$) que $Z$ est $S$-plat dans un voisinage de $z$, et il en est de même de $X$ et $Y$ par hypothèse; on est alors ramené à prouver que $g_s$ est régulière au point $z$, sachant qu'il en est de même de $f_s\circ g_s$, ce qui découle de (19.1.5, (v)), vu les hypothèses faites sur $X_s$ et $Y_s$.

\oldpage[IV-4]{193}

\begin{proposition}[19.2.8]
\label{IV.19.2.8}
Soient $X$ un $S$-préschéma, $Y$ un sous-préschéma fermé de $X$ défini par un Idéal quasi-cohérent $\mathscr J$ de $\sh{O}_X$, $u:Y\to X$ l'injection canonique. Si $u$ est transversalement régulière relativement à $S$, les $\sh{O}_X$-Modules $\mathscr J^n$ et $\mathscr J^n/\mathscr J^m$ ($0\leq n<m$) sont $S$-plats aux points de $Y$. Pour tout changement de base $S'\to S$, si $X'=X_{(S')}$, $Y'=Y_{(S')}$ et $u'=u_{(S')}$, on a
\[
  \sh{G}r_n(u')=\sh{G}r_n(u)\otimes_{\sh{O}_S}\sh{O}_{S'}
\]
à isomorphisme près pour tout $n>0$.
\end{proposition}

En effet, $\sh{O}_X$ et $\sh{O}_X/\mathscr J$ sont par hypothèse $S$-plats aux points de $Y$, et $\mathscr J^n/\mathscr J^{n+1}$ est un $(\sh{O}_X/\mathscr J)$-module localement libre (16.9.3), donc est $S$-plat aux points de $Y$. La première assertion résulte alors de ($0_{\mathrm I}$, 6.1.2) et la seconde de (11.2.9.2).

La proposition suivante généralise et précise (17.16.1):

\begin{proposition}[19.2.9]
\label{IV.19.2.9}
Soient $f:X\to S$ un morphisme, $x$ un point de $X$, $s=f(x)$. Supposons que $f$ soit de présentation finie au point $x$, et soit un morphisme de Cohen-Macaulay au point $x$ (6.8.1). Alors il existe un sous-préschéma $Y$ de $X$ contenant $x$ et ayant les propriétés suivantes:
\begin{enumerate}
  \item[(i)] L'injection canonique $j:Y\to X$ est transversalement régulière relativement à $S$.
  \item[(ii)] Le morphisme $g=f\circ j:Y\to S$ est de Cohen-Macaulay (ce qui équivaut à dire, compte tenu de (i), que toutes les fibres $g^{-1}(g(y))$ de $g$ sont des préschémas de Cohen-Macaulay).
  \item[(iii)] Le point $x$ est point maximal de $Y_s=g^{-1}(s)$.
\end{enumerate}
\end{proposition}

Posons $X_s=f^{-1}(s)$; par hypothèse, l'anneau $\sh{O}_{X_s,x}$ est de Cohen-Macaulay; soit $(t_i)_{1\leq i\leq n}$ un système de paramètres de cet anneau, qui est donc une suite régulière ($0$, 16.5.7). Il y a un voisinage ouvert $U$ de $x$ dans $X$ et des sections $h_i$ ($1\leq i\leq n$) de $\sh{O}_X$ au-dessus de $U$ telles que les $t_i$ soient les images canoniques des $(h_i)_x$ dans l'anneau quotient $\sh{O}_{X_s,x}$ de $\sh{O}_{X,x}$. Soit $\mathscr J$ l'Idéal de $\sh{O}_U$ engendré par les $h_i$, et prenons pour $Y$ le sous-préschéma fermé de $U$ défini par $\mathscr J$. On peut supposer que $f|U$ est localement de présentation finie, et il en est donc de même de $g$ par définition de $Y$; il résulte alors de l'équivalence de $c)$ et $a)$ dans (11.3.8) que la suite $((h_i)_x)$ est régulière dans $\sh{O}_{X,x}$ et que $\sh{O}_{Y,x}$ est plat sur $\sh{O}_{S,s}$; en vertu de (11.3.1) $f$ et $g$ sont donc plats dans un voisinage de $x$ et il résulte de (19.2.4) que $j$ est une immersion transversalement régulière relativement à $S$ (en remplaçant au besoin $Y$ par $Y\cap V$ pour un voisinage ouvert $V$ de $x$ dans $X$). On a ainsi prouvé (i) et on voit qu'on peut supposer que $g:Y\to S$ est un morphisme plat. Comme $\sh{O}_{Y_s,x}$, quotient de $\sh{O}_{X_s,x}$ par l'idéal engendré par les $t_i$, est artinien par construction ($0$, 16.3.6), $x$ est point maximal de $Y_s$. Enfin, l'anneau $\sh{O}_{Y_s,x}$ étant artinien, est un anneau de Cohen-Macaulay ($0$, 16.5.1), donc, en remplaçant au besoin $Y$ par son intersection avec un voisinage ouvert de $x$ dans $X$, on peut supposer que $g$ est un morphisme de Cohen-Macaulay, en vertu de (12.1.7, (iii)).

En particulier, on retrouve l'existence des « quasi-sections » \emph{plates} $Y$ d'un morphisme $f:X\to S$ en un point $x\in X_s$ qui est \emph{fermé} dans $X_s$ et tel que $f$ soit plat au point $x$ et que $\sh{O}_{X_s,x}$ soit un anneau de Cohen-Macaulay (17.16.1) et on voit en outre que l'immersion $Y\to X$ est \emph{transversalement régulière} relativement à $S$. Ce dernier complément est aussi valable pour la quasi-section étale définie dans (17.16.3, (i)).

\subsection{Intersections complètes relatives (cas plat).}
\EGAReaderTarget{ega4.c04.subsec.19-3}{subsection:IV.19.3}

\oldpage[IV-4]{194}

\begin{definition}[19.3.1]
\EGAReaderTarget{ega4.c04.def.19-3-1}{IV.19.3.1}
Soit $A$ un anneau local noethérien. On dit que $A$ est un \emph{anneau d'intersection complète} (ou encore, \emph{anneau d'intersection complète absolue}, s'il y a danger de confusion) si le complété $\widehat A$ est isomorphe au quotient d'un anneau local noethérien complet régulier $B$ par un idéal régulier (i.e. (16.9.7) un idéal engendré par une suite régulière d'éléments de $B$).

Tout anneau local régulier est un anneau d'intersection complète.

On dit qu'un préschéma localement noethérien $X$ est \emph{intersection complète} (absolue) \emph{au point} $x\in X$ si l'anneau $\sh{O}_{X,x}$ est un anneau d'intersection complète.
\end{definition}

\begin{proposition}[19.3.2]
\label{IV.19.3.2}
Soient $B$ un anneau local noethérien régulier, $\mathfrak J$ un idéal de $B$. Pour que $A=B/\mathfrak J$ soit un anneau d'intersection complète, il faut et il suffit que l'idéal $\mathfrak J$ soit régulier (i.e. (16.9.7) engendré par une suite régulière d'éléments de $B$).
\end{proposition}

On peut écrire $\widehat A=\widehat B/\mathfrak J\widehat B$ et puisque $\widehat B$ est un $B$-module fidèlement plat et $B$ noethérien, $\mathfrak J\widehat B$ est régulier si et seulement si $\mathfrak J$ l'est (19.1.5, (ii)). Ceci montre (en vertu de ($0$, 17.1.5)) d'une part que la condition de l'énoncé est suffisante, et d'autre part que pour prouver qu'elle est nécessaire, on peut se borner au cas où $A$ et $B$ sont complets. Par hypothèse, il existe alors un anneau local noethérien complet régulier $B'$ tel que $A$ soit isomorphe à un quotient $B'/\mathfrak J'$, où $\mathfrak J'$ est un idéal régulier de $B'$. Considérons alors le \emph{produit fibré} $B\times_A B'=B''$ pour les homomorphismes surjectifs $B\to A$, $B'\to A$ ($0$, 18.1.2); nous prouverons d'abord le lemme suivant:

\begin{lemma}[19.3.2.1]
\label{IV.19.3.2.1}
Soient $A$, $E$, $F$ trois anneaux, $f:E\to A$, $g:F\to A$ deux homomorphismes, et posons $G=E\times_A F$ ($0$, 18.1.2).
\begin{enumerate}
  \item[(i)] Si $f$ est surjectif et si $E$ et $F$ sont des anneaux noethériens, il en est de même de $G$.
  \item[(ii)] Si $A$, $E$ et $F$ sont des anneaux locaux et si $f$ et $g$ sont des homomorphismes locaux, $G$ est un anneau local et les homomorphismes canoniques $G\to E$, $G\to F$ sont locaux.
  \item[(iii)] Si les conditions de (i) et (ii) sont vérifiées, si $g$ est surjectif, et si $E$ et $F$ sont des anneaux locaux noethériens complets, il en est de même de $G$.
\end{enumerate}
\end{lemma}

\begin{enumerate}
  \item[(i)] Compte tenu de ($0$, 18.1.3 et 18.1.5), $G$ est un $A$-anneau augmenté sur $F$, dont l'idéal d'augmentation $\mathfrak J'$ s'identifie canoniquement au noyau $\mathfrak J$ de $f$, la structure de $G$-module sur $\mathfrak J'$ s'identifiant à celle de $E$-module sur $\mathfrak J$ au moyen de l'homomorphisme canonique $G\to E$. Si $\mathfrak a$ est un idéal quelconque de $G$, $\mathfrak a/(\mathfrak a\cap\mathfrak J')$ est isomorphe à $(\mathfrak a+\mathfrak J')/\mathfrak J'$, donc à un idéal de $F$, et par suite est de type fini; d'autre part, $\mathfrak a\cap\mathfrak J'$ est isomorphe à un idéal de $E$ contenu dans $\mathfrak J$, donc est aussi de type fini, ce qui prouve que $\mathfrak a$ est de type fini, donc $G$ est noethérien.
  \item[(ii)] Soient $\mathfrak r$, $\mathfrak s$ les idéaux maximaux de $E$ et $F$ respectivement; l'ensemble $\mathfrak m=\mathfrak r\times_A\mathfrak s$ est évidemment un idéal de $G$, image réciproque de $\mathfrak r$ par la projection $G\to E$ et de $\mathfrak s$ par la projection $G\to F$. On en déduit aussitôt que les éléments de $\mathfrak m$ sont non inversibles dans $G$. D'autre part, si $(x,y)\notin\mathfrak m$, et si par exemple $x\notin\mathfrak r$, $f(x)$ n'appartient pas à l'idéal maximal de $A$, donc, puisque $g(y)=f(x)$, et que $g$ est local, on a aussi $y\notin\mathfrak s$; on en

\oldpage[IV-4]{195}

conclut que $x$ et $y$ sont inversibles, donc il en est de même de $(x,y)$, ce qui achève de prouver (ii).
  \item[(iii)] L'image de $\mathfrak m$ par l'homomorphisme surjectif $G\to E$ est l'idéal maximal $\mathfrak r$ de $E$, donc l'image de $\mathfrak m^n\cap\mathfrak J'$ est $\mathfrak r^n\cap\mathfrak J$, et comme $\mathfrak J$ est fermé (donc complet) pour la topologie induite par la topologie $\mathfrak r$-adique de $E$ ($0_{\mathrm I}$, 7.3.5), $\mathfrak J'$ est complet pour la topologie induite par la topologie $\mathfrak m$-préadique de $G$. D'autre part, $G/\mathfrak J'$, muni de la topologie $\mathfrak m$-préadique, est isomorphe à $F$ muni de la topologie $\mathfrak s$-adique, donc est complet pour la topologie quotient de la topologie $\mathfrak m$-préadique de $G$. On en déduit aussitôt que $G$ est complet pour la topologie $\mathfrak m$-préadique.
\end{enumerate}

Ce lemme étant établi, il résulte du théorème de Cohen ($0$, 19.8.8) que $B''$ est isomorphe à un quotient d'un anneau local noethérien complet \emph{régulier} $C$. Il résulte alors de (19.1.2) que l'immersion $\Spec(B')\to\Spec(C)$ est régulière; comme par hypothèse il en est de même de l'immersion $\Spec(A)\to\Spec(B')$, on conclut de (19.1.5, (iii)) que l'immersion $\Spec(A)\to\Spec(C)$ est régulière; mais cette immersion s'écrit aussi comme composée $\Spec(A)\to\Spec(B)\to\Spec(C)$. Or comme $B$ est régulier par hypothèse, l'immersion $\Spec(B)\to\Spec(C)$ est régulière (19.1.2); donc il en est de même de $\Spec(A)\to\Spec(B)$ par (19.1.5, (iv)).

\begin{corollary}[19.3.3]
\label{IV.19.3.3}
Soit $X$ un préschéma localement noethérien localement immersible dans un préschéma régulier. Alors l'ensemble des $x\in X$ où $X$ est une intersection complète est ouvert dans $X$ et contient l'ensemble des points réguliers de $X$.
\end{corollary}

On peut en effet se borner au cas où $X=\Spec(B/\mathfrak J)$, où $B$ est un anneau régulier et $\mathfrak J$ un idéal de $B$. L'ensemble des $\mathfrak p\supset\mathfrak J$ dans $\Spec(B)$ tels que $\mathfrak J_{\mathfrak p}$ soit régulier dans $B_{\mathfrak p}$ est alors ouvert (16.9.5), et c'est l'ensemble des points de $\Spec(B)$ où $X$ est une intersection complète en vertu de (19.3.2).

\begin{corollary}[19.3.4]
\label{IV.19.3.4}
Soient $k$ un corps, $X$ un préschéma localement de type fini sur $k$, $k'$ une extension de $k$, $X'=X\otimes_k k'$. Soient $x'$ un point de $X'$, $x$ sa projection dans $X$. Pour que $X$ soit intersection complète en $x$, il faut et il suffit que $X'$ soit intersection complète en $x'$.
\end{corollary}

La question étant locale sur $X$, on peut supposer que $X=\Spec(A)$, où $A$ est une $k$-algèbre de type fini, donc quotient d'une algèbre de polynômes $k[T_1,\ldots,T_n]$, si bien que $X$ est sous-préschéma fermé du schéma régulier $Y=\Spec(k[T_1,\ldots,T_n])$ ($0$, 17.3.7). Si l'on pose $Y'=Y\otimes_k k'=\Spec(k'[T_1,\ldots,T_n])$, $Y'$ est aussi un schéma régulier. Or, puisque $\Spec(k')$ est fidèlement plat sur $\Spec(k)$, il résulte de (19.1.5, (ii)) et du fait qu'il s'agit de préschémas noethériens que l'immersion $Y\to X$ est régulière au point $x$ si et seulement si l'immersion $Y'\to X'$ l'est au point $x'$. La conclusion résulte donc du critère (19.3.2).

\begin{corollary}[19.3.5]
\label{IV.19.3.5}
Soient $A$, $C$ deux anneaux locaux noethériens, $\rho:A\to C$ un homomorphisme local, $B=C/\mathfrak J$ un anneau quotient de $C$, $k$ le corps résiduel de $A$. On suppose que $\rho$ fait de $C$ un $A$-module plat, et que l'anneau $C\otimes_A k$ est régulier. Alors les conditions suivantes sont équivalentes:
\begin{enumerate}
  \item[a)] L'idéal $\mathfrak J$ est transversalement régulier relativement à $A$.
  \item[b)] $B$ est un $A$-module plat, et $B\otimes_A k$ est un anneau d'intersection complète.
\end{enumerate}
\end{corollary}

\oldpage[IV-4]{196}

En vertu de (19.2.4), la condition $a)$ équivaut à dire que $B$ est un $A$-module plat et que $\mathfrak J\otimes_A k$ (qui s'identifie à un idéal de $C\otimes_A k$ en vertu de la platitude de $B$ sur $A$ ($0_{\mathrm I}$, 6.1.2)) est régulier dans cet anneau. Comme l'anneau $C\otimes_A k$ est régulier, il revient au même, en vertu de (19.3.2), de dire (lorsque $B$ est un $A$-module plat) que $\mathfrak J\otimes_A k$ est un idéal régulier de $C\otimes_A k$, ou que $B\otimes_A k=(C\otimes_A k)/(\mathfrak J\otimes_A k)$ est un anneau d'intersection complète; d'où le corollaire.

\begin{definition}[19.3.6]
\EGAReaderTarget{ega4.c04.def.19-3-6}{IV.19.3.6}
Soit $f:X\to S$ un morphisme plat localement de présentation finie. On dit que $X$ est une \emph{intersection complète relativement à $S$ en un point} $x$, si la fibre $f^{-1}(f(x))$ est une intersection complète (absolue) au point $x$. On dit que $X$ est une \emph{intersection complète relativement à $S$}, ou que le morphisme $f$ est un \emph{morphisme d'intersection complète}, si $X$ est intersection complète relativement à $S$ en chacun de ses points.
\end{definition}

\begin{proposition}[19.3.7]
\label{IV.19.3.7}
Soient $g:X\to S$, $h:Y\to S$ deux morphismes plats localement de présentation finie, $f:Y\to X$ une $S$-immersion, $y$ un point de $Y$, $x=f(y)$, $s=g(x)=h(y)$. On suppose que la fibre $X_s=g^{-1}(s)$ soit régulière au point $x$. Alors, pour que $f$ soit une immersion transversalement régulière au point $y$, il faut et il suffit que $Y$ soit une intersection complète relativement à $S$ au point $y$.
\end{proposition}

Dire que $f$ est transversalement régulière au point $y$ revient, en vertu de (19.2.4), à dire que $f_s:Y_s\to X_s$ est une immersion régulière au point $y$. Mais comme $X_s$ est régulier au point $x=f(y)$, cela équivaut, en vertu de (19.3.2), à dire que $Y_s$ est intersection complète (absolue) au point $y$, d'où la proposition.

\begin{corollary}[19.3.8]
\label{IV.19.3.8}
Soit $f:X\to S$ un morphisme plat localement de présentation finie. L'ensemble $U$ des $x\in X$ en lesquels $X$ est une intersection complète relativement à $S$ est ouvert dans $X$. Si de plus $f$ est propre, l'ensemble $E$ des $s\in S$ tels que $X_s=f^{-1}(s)$ soit intersection complète en chacun de ses points est une partie ouverte de $S$.
\end{corollary}

Comme $E=Y-f(X-U)$, la seconde assertion résulte de la première et du fait que $f$ est une application fermée. La première assertion est locale sur $X$, donc on peut supposer que $X$ est un sous-préschéma fermé d'un préschéma de la forme $Z=S[T_1,\ldots,T_n]$ ($T_i$ indéterminées). Ce dernier étant lisse sur $S$ (17.3.8), ses fibres $Z_s$ sont régulières, et il revient donc au même, en vertu de (19.3.7), de dire que $x\in U$ ou que l'immersion $X\to Z$ est transversalement régulière au point $x$; on en conclut donc le corollaire par (19.2.2).

\begin{proposition}[19.3.9]
\label{IV.19.3.9}
\begin{enumerate}
  \item[(i)] Toute immersion ouverte est un morphisme d'intersection complète.
  \item[(ii)] Soit $f:X\to S$ un morphisme plat et localement de présentation finie qui est un morphisme d'intersection complète. Pour tout changement de base $g:S'\to S$, $f'=f_{(S')}:X_{(S')}\to S'$ est un morphisme d'intersection complète. La réciproque est vraie si $g$ est fidèlement plat et quasi-compact.
  \item[(iii)] Si $f:X\to Y$ et $g:Y\to Z$ sont deux morphismes plats localement de présentation finie qui sont des morphismes d'intersection complète, il en est de même de $g\circ f:X\to Z$.
\end{enumerate}
\end{proposition}

L'assertion (i) découle aussitôt de la définition (19.3.1). Pour prouver (ii), remarquons d'abord que si $f$ est plat et localement de présentation finie, il en est de

\oldpage[IV-4]{197}

même de $f'$ (2.1.4), la réciproque étant vraie si $g$ est fidèlement plat ((2.5.1) et (2.7.1)); en outre pour tout $s'\in S'$, si $s=g(s')$, il est équivalent de dire que $f^{-1}(s)$ est une intersection complète ou que $f'^{-1}(s')$ est une intersection complète (19.3.4); d'où (ii), compte tenu de ce que $g$ est surjectif lorsqu'il est fidèlement plat.

Enfin, sous les hypothèses de (iii), $g\circ f$ est plat et localement de présentation finie. Par définition (19.3.6), on est donc ramené au cas où $Z$ est le spectre d'un corps, et à prouver alors que les anneaux locaux $\sh{O}_{X,x}$ sont des anneaux d'intersection complète, ce qui est contenu dans le résultat plus général suivant:

\begin{corollary}[19.3.10]
\label{IV.19.3.10}
Soient $f:X\to S$ un morphisme plat localement de présentation finie qui soit un morphisme d'intersection complète, $x$ un point de $X$, $s=f(x)$. Si $\sh{O}_{S,s}$ est un anneau noethérien d'intersection complète, il en est de même de $\sh{O}_{X,x}$.
\end{corollary}

On peut évidemment se borner au cas où $S=\Spec(A)$ avec $A=\sh{O}_{S,s}$. Montrons en outre qu'on peut se borner au cas où l'anneau local $A$ est \emph{complet}. En effet, posons $A'=\widehat A$, $X'=X\otimes_A A'$, et soit $s'$ l'unique point fermé de $S'=\Spec(A')$, qui est l'unique point de $S'$ au-dessus de $s$; si $f'=f_{(S')}:X'\to S'$, la fibre $f'^{-1}(s')$ est isomorphe canoniquement à $f^{-1}(s)$, puisque $A$ et $A'$ ont même corps résiduel (I, 3.6.4); il y a donc un seul point $x'\in X'$ au-dessus de $x$ et de $s'$, et on a par suite $\sh{O}_{X',x'}=\sh{O}_{X,x}\otimes_A A'$. On en déduit que les anneaux locaux $\sh{O}_{X,x}$ et $\sh{O}_{X',x'}$ ont des \emph{complétés isomorphes}, car de façon générale, si $E$ est un anneau local qui est une $A$-algèbre (l'homomorphisme $A\to E$ étant local), le séparé complété $(E\otimes_A\widehat A)^{\wedge}$ de l'anneau $E\otimes_A\widehat A$ muni de la topologie produit tensoriel est isomorphe au séparé complété de $\widehat E\widehat\otimes_A\widehat A\cong\widehat E$, donc au séparé complété $\widehat E$ de $E$ ($0_{\mathrm I}$, 7.7.1). En vertu de la définition (19.3.1), il revient donc au même de dire que $\sh{O}_{X,x}$ est un anneau d'intersection complète, ou que $\sh{O}_{X',x'}$ est un anneau d'intersection complète.

Supposons donc $A$ \emph{complet}; on peut en outre supposer que $X=\Spec(B)$, $B$ étant quotient d'une algèbre de polynômes $C=A[T_1,\ldots,T_r]$. Comme les anneaux de polynômes sur un corps sont des anneaux réguliers ($0$, 17.3.7), il résulte de (19.3.7) que l'immersion $X\to Z=\Spec(C)$ peut être supposée transversalement régulière relativement à $S$; donc on peut supposer que $B=C/\mathfrak J$, où $\mathfrak J$ est engendré par une suite régulière $(g_i)_{1\leq i\leq n}$ d'éléments de $C$. D'autre part, puisque $A$ est complet, on peut par hypothèse l'écrire $A'/\mathfrak R$, où $A'$ est un anneau local régulier et $\mathfrak R$ un idéal de $A'$ ($0$, 19.8.8), et l'hypothèse que $A$ est un anneau d'intersection complète entraîne que $\mathfrak R$ est engendré par une suite régulière $(f_j)_{1\leq j\leq m}$ (19.3.2). Dans l'anneau de polynômes $C'=A'[T_1,\ldots,T_r]$, les éléments $f_j$ ($1\leq j\leq m$) forment encore une suite régulière engendrant l'idéal $\mathfrak R'=\mathfrak R[T_1,\ldots,T_r]$ ($0$, 15.1.4); comme $C'/\mathfrak R'=C$, on voit que si, pour chaque $i$, $g'_i$ est un élément de $C'$ d'image $g_i$ dans $C$, la suite formée des $f_j$ ($1\leq j\leq m$) et des $g'_i$ ($1\leq i\leq n$) est régulière dans $C'$; si $\mathfrak J'$ est l'idéal de $C'$ qu'elle engendre, on a $B=C'/\mathfrak J'$, d'où la conclusion en vertu de (19.3.2), puisque $C'$ est un anneau régulier ($0$, 17.3.7).

\subsection{Application: critères de régularité et de lissité pour les préschémas éclatés.}
\EGAReaderTarget{ega4.c04.subsec.19-4}{subsection:IV.19.4}

\oldpage[IV-4]{198}

\begin{env}[19.4.1]
\label{IV.19.4.1}
Soient $X$ un préschéma, $Y$ un sous-préschéma fermé de $X$, défini par un Idéal quasi-cohérent $\mathscr J$ de $\sh{O}_X$. Rappelons (II, 8.1.3) que le $X$-schéma $X'$ obtenu en faisant éclater $Y$ est le préschéma
\[
  X'=\operatorname{Proj}(\mathscr S),\qquad\text{où}\qquad \mathscr S=\bigoplus_{n\geq0}\mathscr J^n.
\]
Si $\mathscr J$ est de type fini, le morphisme structural $f:X'\to X$ est projectif. Sans hypothèse sur $\mathscr J$, le sous-préschéma fermé
\[
  Y'=f^{-1}(Y)
\]
de $X'$ est défini par l'Idéal quasi-cohérent $\mathscr J\sh{O}_{X'}$ de $\sh{O}_{X'}$, canoniquement isomorphe à $\sh{O}_{X'}(1)$; de façon plus précise, on a une suite exacte
\begin{equation}
  0\longrightarrow\sh{O}_{X'}(1)\xrightarrow{\widetilde u_0}\sh{O}_{X'}\longrightarrow\sh{O}_{Y'}\longrightarrow0.
  \tag{19.4.1.1}\label{IV.19.4.1.1}
\end{equation}
où $\mathscr J\sh{O}_{X'}$ est l'image de $\widetilde u_0$ (II, 8.1.7 et 8.1.8). Comme au voisinage d'un point de $X'$, le $\sh{O}_{X'}$-Module $\sh{O}_{X'}(1)$ est libre de rang $1$, $\mathscr J\sh{O}_{X'}$ est engendré dans un tel voisinage par une section inversible de $\sh{O}_{X'}$ et $\mathscr J\sh{O}_{X'}/\mathscr J^2\sh{O}_{X'}$ est par suite un $\sh{O}_{X'}/\mathscr J\sh{O}_{X'}$-Module libre de rang $1$. Ceci prouve la première assertion du lemme suivant:
\end{env}

\begin{lemma}[19.4.2]
\label{IV.19.4.2}
L'immersion canonique $Y'\to X'$ est régulière et de codimension $1$, et $Y'$ est canoniquement isomorphe au $Y$-schéma $\operatorname{Proj}(\operatorname{gr}_{\mathscr J}^*(\sh{O}_X))$.
\end{lemma}

La seconde assertion résulte de (II, 3.3.3) appliqué à l'injection canonique $Y\to X$, compte tenu de ce que $\mathscr S\otimes_{\sh{O}_X}\sh{O}_Y=\operatorname{gr}_{\mathscr J}^*(\sh{O}_X)$ par définition, puisque $\sh{O}_Y$ est isomorphe à $\sh{O}_X/\mathscr J$ et $\mathscr J^n\otimes_{\sh{O}_X}(\sh{O}_X/\mathscr J)$ à $\mathscr J^n/\mathscr J^{n+1}$.

\begin{corollary}[19.4.3]
\label{IV.19.4.3}
Si l'immersion canonique $Y\to X$ est quasi-régulière, la restriction $g:Y'\to Y$ du morphisme $f$ est lisse.
\end{corollary}

En effet, $\sh{N}_{Y/X}=\mathscr J/\mathscr J^2$ est alors un $\sh{O}_Y$-Module localement libre de rang fini et $\operatorname{gr}_{\mathscr J}^*(\sh{O}_X)$ est isomorphe à $\mathbf S_{\sh{O}_Y}^*(\sh{N}_{Y/X})$ (16.9.8), donc $Y'$ est $Y$-isomorphe à $\mathbf P(\sh{N}_{Y/X})$ et par suite est lisse sur $Y$ (17.3.9).

Nous verrons plus tard (chap. V) que $g$ est encore lisse lorsque l'immersion $Y\to X$ « ne présente que des singularités ordinaires ».

\begin{proposition}[19.4.4]
\label{IV.19.4.4}
Avec les notations de (19.4.1), supposons que $X$ soit localement noethérien et que le morphisme $g:Y'\to Y$, restriction de $f$, soit lisse. Soient $x'$ un point de $Y'$, $x=g(x')$. Pour que $Y$ soit régulier au point $x$, il faut et il suffit que $Y'$ soit régulier au point $x'$; lorsqu'il en est ainsi, $X'$ est alors régulier au point $x'$.
\end{proposition}

La première assertion résulte de (17.5.8) et la seconde de (19.1.1) et (19.4.2).

\begin{corollary}[19.4.5]
\label{IV.19.4.5}
Sous les hypothèses générales de (19.4.4), supposons que $Y$ soit régulier au point $x$. Alors, pour toute générisation $x_1$ de $x$ dans $X$ n'appartenant pas à $Y$, $X$ est régulier au point $x_1$. Si $\operatorname{Reg}(X)$ est ouvert ((6.12), par exemple si $X$ est excellent (7.8.6, (iii)))

\oldpage[IV-4]{199}

et si $Y$ est régulier, alors il existe un voisinage ouvert $U$ de $Y$ dans $X$ tel que $X$ soit régulier aux points de $U-Y$.
\end{corollary}

Pour prouver la première assertion, on peut se borner au cas où $X$ est un schéma local de point fermé $x$; l'hypothèse entraîne alors que $Y$ est régulier ($0$, 17.3.2), donc il en est de même de $Y'$ par (19.4.4). Notons maintenant que le morphisme $f:X'\to X$ est projectif, donc propre, et par suite (II, 7.2.1), si l'on considère l'unique point $x'_1$ de $X'$ au-dessus de $x_1$ (II, 8.1.3), il existe une spécialisation de $x'_1$ appartenant à $Y'$; en lui appliquant (19.4.4) et utilisant ($0$, 17.3.2) on en conclut que $X'$ est régulier au point $x'_1$, donc $X$ régulier au point $x_1$ puisque $X'-Y'$ est isomorphe à $X-Y$ (II, 8.1.3). Si $Y$ est régulier, tout point de $X-Y$ qui est générisation d'un point de $Y$ appartient donc à $\operatorname{Reg}(X)$; si $\operatorname{Reg}(X)$ est ouvert, $U=Y\cup\operatorname{Reg}(X)$ est donc localement constructible et stable par générisation, donc ouvert ($0_{\mathrm{III}}$, 9.2.5) et répond à la question.

\begin{proposition}[19.4.6]
\label{IV.19.4.6}
Soient $g:X\to S$ un morphisme, $Y$ un sous-préschéma fermé de $X$ défini par un Idéal quasi-cohérent $\mathscr J$, $X'$ le $X$-schéma obtenu en faisant éclater $Y$, $Y'$ l'image réciproque de $Y$ dans $X'$.
\begin{enumerate}
  \item[(i)] Les conditions suivantes sont équivalentes:
  \begin{enumerate}
    \item[a)] $\operatorname{gr}_{\mathscr J}^*(\sh{O}_X)$ est une $\sh{O}_X$-Algèbre $g$-plate.
    \item[b)] Pour tout $n\geq0$, le $\sh{O}_X$-Module $\sh{O}_X/\mathscr J^{n+1}$ est $g$-plat (en d'autres termes, le $n$-ème voisinage infinitésimal $Y^{(n)}$ de $Y$ dans $X$ (16.1.2) est $S$-plat).
    \item[c)] Pour tout changement de base $S_1\to S$ et tout entier $n\geq1$, si l'on pose $X_1=X\times_S S_1$, l'homomorphisme canonique $\mathscr J^n\otimes_{\sh{O}_S}\sh{O}_{S_1}\to\mathscr J^n\sh{O}_{X_1}$ est bijectif.
  \end{enumerate}
  Lorsqu'il en est ainsi, $Y'$ est $S$-plat, et pour tout changement de base $S_1\to S$, si $Y_1$ est l'image réciproque de $Y$ dans $X_1$, le préschéma $X'_1$ obtenu en faisant éclater $Y_1$ dans $X_1$ est canoniquement isomorphe à $X'\times_S S_1$.
  \item[(ii)] Supposons que les conditions équivalentes de (i) soient satisfaites et de plus que les morphismes $X\to S$ et $Y\to S$ soient localement de présentation finie. Alors $\mathscr S=\bigoplus_{n\geq0}\mathscr J^n$ est une $\sh{O}_X$-Algèbre de présentation finie, le morphisme $X'\to X$ est de présentation finie, l'immersion canonique $Y'\to X'$ est transversalement régulière de codimension $1$ relativement à $S$, et l'ensemble des points de $X$ (resp. $X'$) où $X$ (resp. $X'$) est $S$-plat, est un voisinage de $Y$ (resp. $Y'$).
\end{enumerate}
\end{proposition}

\begin{enumerate}
  \item[(i)] L'équivalence de $a)$ et $b)$ résulte aussitôt de ($0_{\mathrm I}$, 6.1.2); pour montrer l'équivalence de $b)$ et $c)$, on peut se borner au cas où $S=\Spec(A)$ et $X=\Spec(B)$ sont affines, avec $\mathscr J=\widetilde{\mathfrak J}$, où $\mathfrak J$ est un idéal de $B$; la suite exacte des Tor montre alors que la condition $c)$ implique que $\operatorname{Tor}_1^A(B/\mathfrak J^{n+1},A')=0$ pour tout entier $n\geq0$ et toute $A$-algèbre $A'$; en prenant pour $A'$ une algèbre somme directe $A\oplus M$, où $M$ est un $A$-module quelconque et où la multiplication est définie par $(a,m)(a',m')=(aa',am'+a'm)$ on voit que la condition précédente équivaut à $\operatorname{Tor}_1^A(B/\mathfrak J^{n+1},M)=0$, donc au fait que $B/\mathfrak J^{n+1}$ est un $A$-module plat. Si l'on pose $\mathscr J_1=\mathscr J\otimes_{\sh{O}_S}\sh{O}_{S_1}$, on a alors $\bigoplus_{n\geq0}\mathscr J_1^n=(\bigoplus_{n\geq0}\mathscr J^n)\otimes_{\sh{O}_S}\sh{O}_{S_1}$ et l'assertion relative au préschéma éclaté $X'_1$ s'en déduit aussitôt. Enfin, pour prouver que $Y'$ est $S$-plat, on peut encore se borner au cas où $S$ et $X$ sont affines; si l'on pose $C=\operatorname{gr}_{\mathfrak J}^*(B)$, tout point de $Y'$ admet alors, en vertu de (19.4.2) et de (II, 2.3.6), un

\oldpage[IV-4]{200}

voisinage ouvert affine dont l'anneau est de la forme $C_{(t)}$, où $t$ est un élément homogène de degré $1$ de $C$; comme par hypothèse $C$ est un $A$-module plat, il en est de même de son anneau de fractions $C_t$ et de la composante de degré $0$, $C_{(t)}$, de ce $A$-module gradué, d'où notre assertion.
  \item[(ii)] On peut toujours se borner au cas où $S$ et $X$ sont affines, $B$ étant donc une $A$-algèbre de présentation finie, $\mathfrak J$ un idéal de type fini de $B$. En vertu de (8.9.1), (8.6.3) et (11.2.9), il existe un sous-anneau noethérien $A_0$ de $A$, une $A_0$-algèbre de type fini $B_0$, un idéal $\mathfrak J_0$ de $B_0$ tels que $B=B_0\otimes_{A_0}A$, $\mathfrak J=\mathfrak J_0B$, que $\operatorname{gr}_{\mathfrak J_0}^*(B_0)$ soit un $A_0$-module \emph{plat} et que $\operatorname{gr}_{\mathfrak J}^*(B)=\operatorname{gr}_{\mathfrak J_0}^*(B_0)\otimes_{A_0}A$. En vertu de (i), on a alors $\bigoplus_{n\geq0}\mathfrak J^n=(\bigoplus_{n\geq0}\mathfrak J_0^n)\otimes_{A_0}A=(\bigoplus_{n\geq0}\mathfrak J_0^n)\otimes_{B_0}B$; comme $\bigoplus_{n\geq0}\mathfrak J_0^n$ est une $B_0$-algèbre engendrée par l'idéal de type fini $\mathfrak J_0$ de $B_0$, donc une $B_0$-algèbre de type fini, et par suite de présentation finie puisque $B_0$ est noethérien, on en déduit que $\bigoplus_{n\geq0}\mathfrak J^n$ est une $B$-algèbre de présentation finie. De même, le morphisme $\operatorname{Proj}(\bigoplus_{n\geq0}\mathfrak J_0^n)\to\Spec(B_0)$ est de type fini (II, 2.7.1), et par suite de présentation finie; donc le morphisme $X'\to X$, qui s'en déduit par changement de base, est de présentation finie. On sait déjà (19.4.2) que l'immersion canonique $Y'\to X'$ est régulière; comme $Y'$ est $S$-plat, ainsi qu'on l'a vu dans (i), on déduit de (19.2.4) que l'immersion $Y'\to X'$ est transversalement régulière relativement à $S$ et que $X'$ est $S$-plat aux points de $Y'$, donc (11.3.1) dans un voisinage ouvert de $Y'$. D'autre part, le fait que $\operatorname{gr}_{\mathscr J}^*(\sh{O}_X)$ est une $\sh{O}_X$-Algèbre $S$-plate entraîne, par (11.3.4), que $X$ est $S$-plat aux points de $Y$, donc (11.3.1) dans un voisinage de $Y$.
\end{enumerate}

\begin{corollary}[19.4.7]
\label{IV.19.4.7}
Soient $g:X\to S$ un morphisme localement de présentation finie, $Y$ un sous-préschéma fermé de $X$ tel que le morphisme composé $Y\to X\to S$ soit localement de présentation finie; supposons en outre que $Y$ soit $S$-plat et que $\sh{O}_X$ soit normalement plat le long de $Y$ (11.3.4). Alors (avec les notations de (19.4.1)), le morphisme $X'\to X$ est de présentation finie, l'immersion canonique $Y'\to X'$ est transversalement régulière de codimension $1$ relativement à $S$, $Y'$ est $Y$-plat (donc $S$-plat) et l'ensemble des points de $X$ (resp. $X'$) où $X$ (resp. $X'$) est $S$-plat est un voisinage de $Y$ (resp. $Y'$).
\end{corollary}

En effet, par hypothèse $\operatorname{gr}_{\mathscr J}^*(\sh{O}_X)$ est un $(\sh{O}_X/\mathscr J)$-Module plat, donc est $S$-plat puisque $Y$ est supposé être $S$-plat. On peut par suite appliquer les résultats de (19.4.6, (i) et (ii)). Comme $Y'$ est $Y$-isomorphe à $\operatorname{Proj}(\operatorname{gr}_{\mathscr J}^*(\sh{O}_X))$ (19.4.2), le même raisonnement que dans (19.4.6) montre que $Y'$ est $Y$-plat.

\begin{proposition}[19.4.8]
\label{IV.19.4.8}
Supposons vérifiées les hypothèses de (19.4.7) et en outre que le morphisme $Y'\to Y$ soit lisse. Soient $x'$ un point de $Y'$, $x$ son image dans $Y$. Pour que $Y$ soit lisse sur $S$ au point $x$, il faut et il suffit que $Y'$ soit lisse sur $S$ au point $x'$; lorsqu'il en est ainsi, $X'$ est alors lisse sur $S$ au point $x'$.

Les conclusions précédentes sont en particulier vérifiées lorsque le morphisme $X\to S$ est localement de présentation finie et que l'immersion $Y\to X$ est transversalement régulière relativement à $S$.
\end{proposition}

Compte tenu de la compatibilité de la formation d'un préschéma éclaté avec les changements de base dans le cas envisagé (19.4.6, (i)) et de (17.5.1) et (17.7.1),

\oldpage[IV-4]{201}

il suffit de considérer le cas où $S$ est le spectre d'un corps algébriquement clos. Mais alors il revient au même de dire qu'un $S$-préschéma localement de type fini est lisse sur $S$ en un point ou qu'il est régulier en ce point (17.15.1), donc les conclusions résultent de (19.4.4).

Lorsque $X\to S$ est localement de présentation finie, et que l'immersion canonique $Y\to X$ est quasi-régulière, cette immersion est par définition de présentation finie puisque l'Idéal qui définit $Y$ est de type fini, donc le morphisme composé $Y\to X\to S$ est localement de présentation finie. On a déjà vu (19.4.3) que dans ces conditions le morphisme $Y'\to Y$ est lisse; en outre, on sait alors que $\sh{N}_{Y/X}$ est un $\sh{O}_Y$-Module localement libre et que $\operatorname{gr}_{\mathscr J}^*(\sh{O}_X)$ est isomorphe à $\mathbf S_{\sh{O}_Y}^*(\sh{N}_{Y/X})$, donc est un $\sh{O}_Y$-Module plat. Les hypothèses de (19.4.7) sont donc toutes vérifiées et l'on peut appliquer les conclusions de l'énoncé de (19.4.8).

\begin{corollary}[19.4.9]
\label{IV.19.4.9}
Sous les hypothèses générales de (19.4.8), supposons en outre que $Y$ soit lisse sur $S$ au point $x$. Alors, pour toute générisation $x_1$ de $x$ dans $X$ n'appartenant pas à $Y$, $X$ est lisse sur $S$ au point $x_1$. Si $Y$ est lisse sur $S$, alors il existe un voisinage ouvert $U$ de $Y$ dans $X$ tel que $X$ soit lisse sur $S$ aux points de $U-Y$.
\end{corollary}

Pour la première assertion, on peut se borner au cas où $X$ est un schéma local de point fermé $x$, et l'hypothèse entraîne alors que $Y$ est lisse sur $S$, tout voisinage du point fermé de $Y$ dans $Y$ étant nécessairement $Y$ lui-même; on conclut donc par le même raisonnement que dans (19.4.5), utilisant le fait que l'ensemble des points où $X'$ est lisse sur $S$ est ouvert dans $X'$. Pour prouver la seconde assertion, on peut se ramener au cas où $S$ est noethérien, et $X$ de type fini sur $S$, grâce à (8.9.1), (8.6.3), (11.2.9) et (17.7.6); on conclut alors comme dans (19.4.5), en utilisant le fait que l'ensemble des points où $X$ est lisse sur $S$ est ouvert dans $X$.

\begin{remark}[19.4.10]
\label{IV.19.4.10}
Dans (19.4.6, (ii)), remplaçons l'hypothèse que $X$ et $Y$ sont localement de présentation finie sur $S$ par celle que $S$ et $X$ (donc aussi $Y$) sont \emph{localement noethériens}. Alors il est clair que $\bigoplus_{n\geq0}\mathscr J^n$ est une $\sh{O}_X$-Algèbre de type fini, que le morphisme $X'\to X$ est de type fini, et il résulte encore de (19.2.4) appliqué aux préschémas localement noethériens, que l'immersion $Y'\to X'$ est transversalement régulière relativement à $S$ et que $X'$ est $S$-plat aux points de $Y'$.
\end{remark}

\begin{proposition}[19.4.11]
\EGAReaderTarget{ega4.c04.prop.19-4-11}{IV.19.4.11}
Soient $X$ un préschéma, $\mathscr E$ un $\sh{O}_X$-Module quasi-cohérent, $u:\mathscr E\to\sh{O}_X$ un homomorphisme de $\sh{O}_X$-Modules, $\mathscr J=u(\mathscr E)$ l'Idéal quasi-cohérent de $\sh{O}_X$ image de $u$, $Y$ le sous-préschéma fermé de $X$ défini par $\mathscr J$. Soit $X'$ le $X$-schéma obtenu en faisant éclater $Y$. D'autre part, soient $P=\mathbf P(\mathscr E)$ le fibré projectif sur $X$ défini par $\mathscr E$ (II, 4.1.1), $p:P\to X$ le morphisme structural, $\alpha_1^{\#}:p^*(\mathscr E)\to\sh{O}_P(1)$ l'homomorphisme canonique (II, 4.1.5.1) et $\mathscr H$ son noyau, de sorte qu'on a la suite exacte
\[
  0\longrightarrow\mathscr H\longrightarrow p^*(\mathscr E)\xrightarrow{\alpha_1^{\#}}\sh{O}_P(1)\longrightarrow0.
\]
Enfin, soit $\mathscr K$ l'Idéal quasi-cohérent de $\sh{O}_P$ image de la restriction $v:\mathscr H\to\sh{O}_P$ de $p^*(u)$, et soit $Z$ le sous-préschéma fermé de $P$ défini par $\mathscr K$.

\oldpage[IV-4]{202}

\begin{enumerate}
  \item[(i)] L'image de $X'$ par l'immersion fermée $j:X'\to\mathbf P(\mathscr E)$ correspondant à l'homomorphisme surjectif de $\sh{O}_X$-Algèbres graduées $\mathbf S_{\sh{O}_X}(\mathscr E)\to\mathbf S_{\sh{O}_X}(\mathscr J)\to\bigoplus_{n\geq0}\mathscr J^n$ (II, 3.6.2) est un sous-préschéma fermé de $Z$.
  \item[(ii)] Si $\mathscr E$ est un $\sh{O}_X$-Module localement libre de rang $m$, $\mathscr H$ est un $\sh{O}_P$-Module localement libre de rang $m-1$.
  \item[(iii)] On suppose que $X$ soit localement noethérien, que $\mathscr E$ soit localement libre de type fini, que le sous-préschéma $Y$ soit régulier, et qu'en tout point $x\in Y$, l'immersion canonique $i:Y\to X$ soit régulière et de codimension égale à $\operatorname{rg}_{\sh{O}_{X,x}}(\mathscr E_x)$ (ce que nous exprimerons en disant que l'homomorphisme $u$ est \emph{régulier} au point $x$ (cf. chap. V)). Alors l'image de $X'$ par l'immersion fermée $j:X'\to P$ est le sous-préschéma fermé $Z$ de $P$; en tous les points de $X'-Y'$ ($Y'$ image réciproque de $Y$ dans $X'$) l'immersion $j$ est transversalement régulière relativement à $X$; enfin en tous les points $x'$ de $Y'$, $X'$ est régulier, l'immersion $j$ est régulière et de codimension $\operatorname{rg}_{\sh{O}_{P,z}}(\mathscr H_z)$ où $z=j(x')$ (autrement dit l'homomorphisme $v:\mathscr H\to\sh{O}_P$ est régulier en $z$).
\end{enumerate}
\end{proposition}

Toutes les questions étant locales sur $X$, on peut se borner au cas où $X=\Spec(A)$ est affine et $\mathscr E=\widetilde E$, où $E$ est un $A$-module. Pour faire la vérification de (i), on peut en outre se borner à examiner ce qui se passe dans un ouvert affine de $P$ de la forme $D_+(t)$, où $t\in E=\mathbf S_A^1(E)$ (II, 2.3.14). Si l'on pose $S=\mathbf S_A(E)$, on a alors $\Gamma(D_+(t),\sh{O}_P)=S_{(t)}$ (II, 2.4.1), et $p^*(\mathscr E)|D_+(t)=(E\otimes_A S_{(t)})^\sim$. Si l'on se reporte à la définition de $\alpha_1^{\#}$ (II, 2.6.2.2), on voit qu'à tout élément
\[
  a=\sum_i x^{(i)}\otimes\frac{x_1^{(i)}x_2^{(i)}\cdots x_n^{(i)}}{t^n}
\]
de $E\otimes_A S_{(t)}$ (où $x^{(i)}$ et les $x_j^{(i)}$ sont dans $E$) $\alpha_1^{\#}$ fait correspondre l'élément
\[
  \sum_i\frac{x^{(i)}x_1^{(i)}\cdots x_n^{(i)}}{t^n}
\]
de $(S(1))_{(t)}$. Il s'agit (pour prouver (i)) de montrer que si ce dernier élément est nul, alors il en est de même de l'image de
\[
  a'=v(a)=\sum_i u(x^{(i)})\frac{x_1^{(i)}\cdots x_n^{(i)}}{t^n}
\]
par l'homomorphisme canonique $S_{(t)}\to(\bigoplus_{n\geq0}(u(E))^n)_{(u(t))}$ (II, 3.6.2). Or, cette image n'est autre que
\[
  \sum_i\frac{u(x^{(i)})u(x_1^{(i)})\cdots u(x_n^{(i)})}{(u(t))^n}
\]
dans l'anneau $A_{u(t)}$, c'est-à-dire le produit de l'élément $u(t)$ et de l'image canonique de $\sum_i x^{(i)}x_1^{(i)}\cdots x_n^{(i)}/t^{n+1}$ par l'homomorphisme d'algèbres, prolongement canonique à $S_{(t)}$ de l'homomorphisme $u:E\to A$ de $A$-modules. Cela prouve donc (i).

Pour établir (ii), notons que l'on peut supposer que $E$ est un $A$-module libre de rang $m$; puisque $\alpha_1^{\#}:E\otimes_A S_{(t)}\to(S(1))_{(t)}$ est surjectif et que $(S(1))_{(t)}$ est un $S_{(t)}$-module libre de rang $1$, la suite exacte
\[
  0\longrightarrow H_{(t)}\longrightarrow E\otimes_A S_{(t)}\xrightarrow{\alpha_1^{\#}}(S(1))_{(t)}\longrightarrow0
\]
est \emph{scindée}, et $H_{(t)}$ est donc un $S_{(t)}$-module projectif de rang $m-1$, d'où (ii).

Pour prouver (iii), examinons d'abord ce qui se passe au-dessus d'un point $x\in X-Y$. Alors $u:\mathscr E\to\sh{O}_X$, restreint à un voisinage de $x$, est surjectif, donc on peut se borner au cas où $\mathscr J=\sh{O}_X$, cas où $X'$ s'identifie canoniquement à $X$. Pour voir alors que le sous-préschéma fermé de $P$ image de $X'$ s'identifie à $Z$, on peut se borner comme au début au cas où $X=\Spec(A)$ est affine, et en outre $E=A^m$, de sorte que $S=A[T_0,\ldots,T_{m-1}]$. On peut en outre, avec les notations du début, supposer $u(t)\ne0$ (les valeurs de $t\in E$

\oldpage[IV-4]{203}

qui vérifient cette condition engendrant l'algèbre symétrique $\mathbf S_A(E)$); changeant au besoin de base dans $E$, on peut donc supposer que $t=T_0$, et $u(T_i)=0$ pour $1\leq i\leq m-1$, de sorte que $S_{(t)}$ s'identifie canoniquement à $A[T_1,\ldots,T_{m-1}]$. Les éléments de $H$ sont alors les éléments de la forme
\[
  \sum_\alpha L_\alpha(\mathbf T)\otimes\mathbf T^\alpha
\]
de $E\otimes_A A[T_1,\ldots,T_{m-1}]$ où les $L_\alpha(\mathbf T)=\lambda_{\alpha0}+\lambda_{\alpha1}T_1+\cdots+\lambda_{\alpha,m-1}T_{m-1}$ doivent être tels que $\sum_\alpha L_\alpha(\mathbf T)\mathbf T^\alpha=0$ dans $A[T_1,\ldots,T_{m-1}]$. On vérifie aisément que lorsque la condition $\lambda_{\mathbf0,0}=0$ est remplie, les valeurs des $\lambda_{\alpha0}$ peuvent être prises \emph{arbitrairement} pour $|\alpha|\geq1$, et on peut toujours alors déterminer les $\lambda_{\alpha i}$ pour $i\geq1$ de façon à vérifier la condition $\sum_\alpha L_\alpha(\mathbf T)\mathbf T^\alpha=0$. Les éléments de l'idéal $\mathscr K$ correspondants sont alors les polynômes $\sum_\alpha\lambda_{\alpha0}\mathbf T^\alpha$ \emph{sans terme constant}, et comme $(\mathbf S_A(A))_{(t)}$ s'identifie à $A$, ces polynômes sont exactement les éléments du noyau de l'homomorphisme $(\mathbf S_A(E))_{(t)}\to(\mathbf S_A(A))_{(t)}$. On a ainsi prouvé (19.2.1) que l'immersion $j:X'\to P$ est \emph{transversalement régulière} relativement à $X$ en tous les points de $X'-Y'$ et que $j(X')=Z$ en ces points.

Reste à examiner ce qui se passe au-dessus d'un point $x\in Y$. L'hypothèse de « régularité » de $u$ en ce point faite dans (iii) équivaut (en vertu de (19.1.1) et de ($0$, 17.1.7)) à dire que $X$ est \emph{régulier} au point $x$ et que l'homomorphisme
\[
  u_x\otimes1:\mathscr E_x\otimes_{\sh{O}_{X,x}}\mathbf k(x)\longrightarrow\mathfrak m_x/\mathfrak m_x^2
\]
est \emph{injectif}. Cela étant, comme le morphisme $p:P\to X$ est lisse (17.3.9), pour tout $z\in P$ au-dessus de $x$, $P$ est régulier au point $z$ (17.5.8); en outre ($0$, 17.3.3) l'homomorphisme canonique
\begin{equation}
  (\mathfrak m_x/\mathfrak m_x^2)\otimes_{\mathbf k(x)}\mathbf k(z)\longrightarrow\mathfrak m_z/\mathfrak m_z^2
  \tag{19.4.11.1}\label{IV.19.4.11.1}
\end{equation}
est \emph{injectif}. On en déduit que l'homomorphisme
\[
  (\mathscr E_x\otimes_{\sh{O}_{X,x}}\sh{O}_{P,z})\otimes_{\sh{O}_{P,z}}\mathbf k(z)\longrightarrow\mathfrak m_z/\mathfrak m_z^2
\]
déduit de $p^*(u)$ est aussi injectif, car il s'écrit comme le composé
\[
  (\mathscr E_x\otimes_{\sh{O}_{X,x}}\mathbf k(x))\otimes_{\mathbf k(x)}\mathbf k(z)\longrightarrow(\mathfrak m_x/\mathfrak m_x^2)\otimes_{\mathbf k(x)}\mathbf k(z)\longrightarrow\mathfrak m_z/\mathfrak m_z^2
\]
où la première flèche est injective par platitude et la seconde est l'homomorphisme injectif (19.4.11.1). Comme $\mathscr H$ est (dans un voisinage de $z$ dans $P$) facteur direct de $\mathscr E\otimes_{\sh{O}_X}\sh{O}_P=p^*(\mathscr E)$, l'homomorphisme $v_z\otimes1:\mathscr H_z\otimes_{\sh{O}_{P,z}}\mathbf k(z)\to\mathfrak m_z/\mathfrak m_z^2$ est aussi injectif, ce qui établit la « régularité » de l'homomorphisme $v$ au point $z$. Reste à voir que l'image par $j$ de $X'$ est identique au sous-préschéma fermé $Z$. Il suffira pour cela de montrer d'une part que l'image par $j$ des points de $q^{-1}(x)$ (où $q:X'\to X$ est le morphisme structural) sont exactement ceux de $Z\cap p^{-1}(x)$ et d'autre part qu'en chacun de ces points $z=j(x')$, l'homomorphisme surjectif $\sh{O}_{Z,z}\to\sh{O}_{X',x'}$ est bijectif. Mais, compte tenu de l'hypothèse et de (19.4.3), la restriction $Y'\to Y$ de $q$ est lisse, et puisque $\sh{O}_{Y,x}$ est supposé régulier, il en est de même de $\sh{O}_{Y',x'}$ (17.5.8); en outre, l'immersion $Y'\to X'$ est régulière (19.4.2), donc $\sh{O}_{X',x'}$ est aussi un anneau régulier (19.1.1). Pour prouver que l'homomorphisme $\sh{O}_{Z,z}\to\sh{O}_{X',x'}$ est bijectif, il suffit donc de montrer que les

\oldpage[IV-4]{204}

dimensions de ces deux anneaux réguliers sont égales ($0$, 17.1.9). Mais en vertu de ($0$, 17.3.3), cela équivaut aussi à l'égalité des dimensions des deux anneaux $\sh{O}_{Z,z}\otimes_{\sh{O}_{X,x}}\mathbf k(x)$ et $\sh{O}_{X',x'}\otimes_{\sh{O}_{X,x}}\mathbf k(x)$, si bien que finalement on peut remplacer $X'$ et $Z$ par les \emph{fibres} $q^{-1}(x)$ et $Z\cap p^{-1}(x)$ respectivement, et se ramener à prouver que le morphisme $q^{-1}(x)\to Z\cap p^{-1}(x)$ déduit de $j$ est un \emph{isomorphisme}. Or, on peut ici se borner au cas où $X=\Spec(B)$ avec $B=\sh{O}_{X,x}$ et $E=B^m$, d'où $S=B[T_1,\ldots,T_m]$; en outre, l'hypothèse signifie qu'il existe un système régulier de paramètres $(t_i)_{1\leq i\leq n}$ de $B$ tel que $u(T_i)=t_i$ pour $1\leq i\leq m$ ($0$, 17.1.7); soit $\mathfrak J$ l'image de $u$; si $\mathfrak m$ est l'idéal maximal de $B$, on a $\mathfrak J^n\otimes_B B/\mathfrak m=\mathfrak J^n/\mathfrak m\mathfrak J^n=(\mathfrak J^n/\mathfrak J^{n+1})\otimes_B(B/\mathfrak m)$ et par suite (16.9.4.1), $(\bigoplus_{n\geq0}\mathfrak J^n)\otimes_B\mathbf k(x)$ s'identifie à $\mathbf k(x)[T_1,\ldots,T_m]$, d'où résulte que l'immersion $q^{-1}(x)\to p^{-1}(x)$ est ici un \emph{isomorphisme}; ceci implique \emph{a fortiori} que le sous-préschéma fermé $Z\cap p^{-1}(x)$ est identique à la fibre $p^{-1}(x)$, et termine la démonstration de (19.4.11).

\subsection{Critères de M-régularité.}
\EGAReaderTarget{ega4.c04.subsec.19-5}{subsection:IV.19.5}

Nous reprenons ici, en les complétant, les critères pour qu'une suite d'éléments d'un anneau $A$ soit M-régulière ou M-quasi-régulière ($M$ étant un $A$-module), déjà étudiés dans ($0$, 15.1).

\begin{theorem}[19.5.1]
\label{IV.19.5.1}
Soient $A$ un anneau, $\mathbf f=(f_i)_{1\leq i\leq n}$ une suite d'éléments de $A$, $M$ un $A$-module. Considérons les propriétés suivantes:
\begin{enumerate}
  \item[a)] La suite $\mathbf f$ est M-régulière.
  \item[b)] On a $\operatorname{H}_i(\mathbf f,M)=0$ pour tout $i>0$ (III, 1.1.3).
  \item[b$'$)] On a $\operatorname{H}_1(\mathbf f,M)=0$.
  \item[c)] La suite $\mathbf f$ est M-quasi-régulière.
\end{enumerate}
Alors on a les implications
\[
  a)\Rightarrow b)\Rightarrow b')\qquad\text{et}\qquad a)\Rightarrow c).
\]
Posons $\mathfrak J=\sum_{i=1}^n f_iA$. Si les modules $M_i=M/(\sum_{j=1}^i f_jM)$ sont séparés pour la topologie $\mathfrak J$-préadique (resp. si tout module quotient d'un sous-module de $M$ est séparé pour la topologie $\mathfrak J$-préadique), alors on a aussi l'implication $c)\Rightarrow a)$ (resp. $b')\Rightarrow a)$.
\end{theorem}

Les implications $a)\Rightarrow c)$, et $c)\Rightarrow a)$ lorsque les $M_i$ sont séparés pour la topologie $\mathfrak J$-préadique, ont déjà été prouvées ($0$, 15.1.9), et ne sont données que pour mémoire. On a aussi montré (III, 1.1.4 et 1.1.3.3) que $a)$ entraîne $b)$, et $b)$ implique trivialement $b')$. Il reste donc à voir que lorsque tout module quotient d'un sous-module de $M$ est séparé pour la topologie $\mathfrak J$-préadique, $b')$ entraîne $a)$. Raisonnons par récurrence sur $n$. Pour $n=1$, l'assertion résulte trivialement des définitions, car alors $\operatorname{H}_1(\mathbf f,M)$ n'est autre que le noyau de l'homothétie $z\mapsto f_1z$ dans $M$ (III, 1.1.1 et 1.1.2). Supposons donc $n\geq2$, et posons $\mathbf f'=(f_i)_{1\leq i\leq n-1}$; avec les notations de (III, 1.1.2), on a $\operatorname{K}_\bullet(\mathbf f,M)=\operatorname{K}_\bullet(f_n)\otimes_A\operatorname{K}_\bullet(\mathbf f',M)$. On a donc (III, 1.1.4.1) la suite exacte
\begin{equation}
  0\longrightarrow\operatorname{H}_0(f_n,\operatorname{H}_1(\mathbf f',M))\longrightarrow\operatorname{H}_1(\mathbf f,M)\longrightarrow\operatorname{H}_1(f_n,\operatorname{H}_0(\mathbf f',M))\longrightarrow0.
  \tag{19.5.1.1}\label{IV.19.5.1.1}
\end{equation}

\oldpage[IV-4]{205}

La relation $\operatorname{H}_1(\mathbf f,M)=0$ implique donc $\operatorname{H}_0(f_n,\operatorname{H}_1(\mathbf f',M))=0$ et $\operatorname{H}_1(f_n,\operatorname{H}_0(\mathbf f',M))=0$. La première de ces deux relations signifie que l'on a $f_n\operatorname{H}_1(\mathbf f',M)=\operatorname{H}_1(\mathbf f',M)$ (III, 1.1.3.5). Or, par définition (III, 1.1.1), $\operatorname{H}_1(\mathbf f',M)$ est isomorphe à un quotient d'un sous-module $N$ de $M^{n-1}$; prenant sur $M^{n-1}$ la filtration des $M^j$ pour $j\leq n-1$, sur $N$ la filtration induite et sur $\operatorname{H}_1(\mathbf f',M)$ la filtration quotient de celle de $N$, on obtient sur $\operatorname{H}_1(\mathbf f',M)$ une filtration finie dont les quotients sont isomorphes à des quotients de sous-modules de $M$, donc sont séparés pour la topologie $\mathfrak J$-préadique par hypothèse. Comme la relation $f_n\operatorname{H}_1(\mathbf f',M)=\operatorname{H}_1(\mathbf f',M)$ implique \emph{a fortiori} $\mathfrak J\operatorname{H}_1(\mathbf f',M)=\operatorname{H}_1(\mathbf f',M)$, on déduit de l'hypothèse et des remarques précédentes que $\operatorname{H}_1(\mathbf f',M)=0$. L'hypothèse de récurrence prouve alors que la suite $\mathbf f'$ est M-régulière. D'autre part, la relation $\operatorname{H}_1(f_n,\operatorname{H}_0(\mathbf f',M))=0$ signifie que $f_n$ est $(M/(\sum_{i=1}^{n-1}f_iM))$-régulier, donc la suite $\mathbf f$ est M-régulière.

\begin{corollary}[19.5.2]
\label{IV.19.5.2}
Supposons que $A$ soit un anneau noethérien, que les $f_i$ ($1\leq i\leq n$) appartiennent au radical de $A$ et que $M$ soit un $A$-module de type fini. Alors les quatre conditions $a)$, $b)$, $b')$, $c)$ de (19.5.1) sont équivalentes.
\end{corollary}

On sait en effet alors ($0_{\mathrm I}$, 7.3.5) que tout $A$-module de type fini est séparé pour la topologie $\mathfrak J$-préadique.

\begin{proposition}[19.5.3]
\label{IV.19.5.3}
Soient
\[
  0\longrightarrow M'\longrightarrow M\longrightarrow M''\longrightarrow0
\]
une suite exacte de $A$-modules, $\mathbf f=(f_i)_{1\leq i\leq n}$ une suite M-régulière, $\mathfrak J=\sum_{i=1}^n f_iA$. Considérons les propriétés suivantes:
\begin{enumerate}
  \item[a)] La suite $\mathbf f$ est $M''$-régulière.
  \item[b)] La filtration $\mathfrak J$-préadique de $M'$ est induite sur $M'$ par la filtration $\mathfrak J$-préadique de $M$ (autrement dit, l'homomorphisme canonique $\operatorname{gr}_{\mathfrak J}^*(M')\to\operatorname{gr}_{\mathfrak J}^*(M)$ est injectif).
  \item[c)] L'homomorphisme canonique $M'/(\sum_{i=1}^n f_iM')\to M/(\sum_{i=1}^n f_iM)$ est injectif.
\end{enumerate}
Alors on a les implications $a)\Rightarrow b)\Rightarrow c)$, et $a)$ entraîne en outre que la suite $\mathbf f$ est $M'$-régulière. Si tout quotient d'un sous-module de $M''$ est séparé pour la topologie $\mathfrak J$-préadique, les conditions $a)$, $b)$ et $c)$ sont équivalentes.
\end{proposition}

Il est clair que $c)$ est conséquence de $b)$. Prouvons que sous les hypothèses de la dernière assertion, $c)$ entraîne $a)$: puisque $\mathbf f$ est M-régulière, on a, par (19.5.1), $\operatorname{H}_1(\mathbf f,M)=0$, d'où une suite exacte, portion de la suite exacte d'homologie
\[
  0\longrightarrow\operatorname{H}_1(\mathbf f,M'')\longrightarrow\operatorname{H}_0(\mathbf f,M')\longrightarrow\operatorname{H}_0(\mathbf f,M).
\]
Or, la condition $c)$ exprime que l'homomorphisme $\operatorname{H}_0(\mathbf f,M')\to\operatorname{H}_0(\mathbf f,M)$ est injectif (III, 1.1.3.5); elle équivaut donc à $\operatorname{H}_1(\mathbf f,M'')=0$, d'où, en vertu de l'hypothèse de séparation et de (19.5.1), le fait que $\mathbf f$ est $M''$-régulière.

Montrons ensuite que $a)$ entraîne $c)$ et le fait que $\mathbf f$ est $M'$-régulière, en procédant par récurrence sur $n$. Pour $n=1$, $f_1$ étant M-régulier est aussi $M'$-régulier et la propriété $c)$ n'est autre que le lemme (3.4.1.4). Pour $n>1$, l'hypothèse de récurrence montre que

\oldpage[IV-4]{206}

si l'on pose $\mathbf f'=(f_1,\ldots,f_{n-1})$, la suite $\mathbf f'$ est $M'$-régulière et l'on a, en vertu de $c)$ appliqué à $\mathbf f'$, une suite exacte
\[
  0\longrightarrow M'/(\textstyle\sum_{i=1}^{n-1}f_iM')\longrightarrow M/(\textstyle\sum_{i=1}^{n-1}f_iM)\longrightarrow M''/(\textstyle\sum_{i=1}^{n-1}f_iM'')\longrightarrow0.
\]
Par hypothèse, $f_n$ est $(M/(\sum_{i=1}^{n-1}f_iM))$-régulier et $(M''/(\sum_{i=1}^{n-1}f_iM''))$-régulier, donc le même raisonnement montre d'une part que $f_n$ est $(M'/(\sum_{i=1}^{n-1}f_iM'))$-régulier, et d'autre part, en vertu de (3.4.1.4), que la suite
\[
  0\longrightarrow M'/(\textstyle\sum_{i=1}^{n}f_iM')\longrightarrow M/(\textstyle\sum_{i=1}^{n}f_iM)\longrightarrow M''/(\textstyle\sum_{i=1}^{n}f_iM'')\longrightarrow0
\]
est exacte, d'où $c)$.

Montrons enfin que $a)$ entraîne $b)$. Comme la suite $\mathbf f$ est alors M-régulière et $M'$-régulière, donc M-quasi-régulière et $M'$-quasi-régulière, on a des isomorphismes canoniques
\[
  \operatorname{gr}_{\mathfrak J}^*(M')\xrightarrow{\sim}\operatorname{gr}_{\mathfrak J}^0(M')[T_1,\ldots,T_n],
  \qquad
  \operatorname{gr}_{\mathfrak J}^*(M)\xrightarrow{\sim}\operatorname{gr}_{\mathfrak J}^0(M)[T_1,\ldots,T_n]
\]
($0$, 15.1.7), et comme on a vu ci-dessus que $\operatorname{gr}_{\mathfrak J}^0(M')\to\operatorname{gr}_{\mathfrak J}^0(M)$ est injectif, il en est de même de $\operatorname{gr}_{\mathfrak J}^*(M')\to\operatorname{gr}_{\mathfrak J}^*(M)$.

On notera encore que l'équivalence des conditions $a)$, $b)$, $c)$ de (19.5.3) est valable en particulier lorsque $A$ est noethérien, $M''$ un $A$-module de type fini et que les $f_i$ appartiennent au radical de $A$.

\begin{env}[19.5.4]
\label{IV.19.5.4}
Dans la suite de ce numéro, nous conservons les notations $A$, $\mathbf f$, $M$, $\mathfrak J$ de (19.5.1), et nous supposons de plus $M$ muni d'une filtration décroissante $(M_k)$ formée de sous-$A$-modules, telle que $M_0=M$. Rappelons ($0$, 15.1.5) que l'on définit alors sur $M$ une seconde filtration décroissante formée des sous-modules
\[
  M_k'=M_k+\mathfrak JM_{k-1}+\cdots+\mathfrak J^{k-1}M_1+\mathfrak J^kM_0
\]
et que, si $\operatorname{gr}_\bullet(M)$ et $\operatorname{gr}_\bullet'(M)$ sont les $A$-modules gradués associés aux filtrations $(M_k)$ et $(M_k')$ respectivement, on définit un homomorphisme gradué surjectif de degré $0$ ($0$, 15.1.5.2)
\[
  \psi_M:(\operatorname{gr}_\bullet(M)\otimes_A(A/\mathfrak J))[T_1,\ldots,T_n]\longrightarrow\operatorname{gr}_\bullet'(M).
\]
Rappelons aussi que l'on a l'homomorphisme canonique surjectif ($0$, 15.1.1.1)
\[
  \varphi_{\operatorname{gr}_\bullet(M)}:(\operatorname{gr}_\bullet(M)\otimes_A(A/\mathfrak J))[T_1,\ldots,T_n]\longrightarrow\operatorname{gr}_{\mathfrak J}^*(\operatorname{gr}_\bullet(M)).
\]
\end{env}

\begin{theorem}[19.5.5]
\label{IV.19.5.5}
Avec les notations de (19.5.4), considérons les propriétés suivantes:
\begin{enumerate}
  \item[a)] La suite $\mathbf f$ est $\operatorname{gr}_\bullet(M)$-régulière.
  \item[b)] L'homomorphisme canonique $\psi_M$ est bijectif.
  \item[c)] L'homomorphisme canonique $\varphi_{\operatorname{gr}_\bullet(M)}$ est bijectif (autrement dit ($0$, 15.1.7), la suite $\mathbf f$ est $\operatorname{gr}_\bullet(M)$-quasi-régulière).
\end{enumerate}

\oldpage[IV-4]{207}

On a alors les implications\footnote{L'implication $b)\Rightarrow c)$, sans hypothèse de séparation sur les $\operatorname{gr}_k(M)$, est due à P. Deligne, dont nous reproduisons ci-dessous la démonstration.}
\[
  a)\Rightarrow b)\Rightarrow c).
\]
En outre, si, pour tout entier $k\geq0$, les modules quotients des sous-modules de $\operatorname{gr}_k(M)$ sont séparés pour la topologie $\mathfrak J$-préadique (ce qui est le cas lorsque $A$ est noethérien, que les $f_i$ sont dans le radical de $A$ et que les $\operatorname{gr}_k(M)$ sont des $A$-modules de type fini), alors les conditions $a)$, $b)$ et $c)$ sont équivalentes.
\end{theorem}

L'implication $a)\Rightarrow b)$ a été démontrée dans ($0$, 15.1.8); l'implication $c)\Rightarrow a)$ sous l'hypothèse de séparation est un cas particulier de ($0$, 15.1.9). Il reste à prouver que $b)$ implique $c)$, ce qui va se faire en plusieurs étapes.

Nous désignerons par $\mathfrak J$ l'idéal $\sum_{i=1}^n f_iA$ engendré par les $f_i$; pour toute suite $\mathbf q=(q_1,\ldots,q_n)$ de $n$ entiers $\geq0$, on posera $|\mathbf q|=\sum_{i=1}^nq_i$ et $\mathbf f^{\mathbf q}=f_1^{q_1}f_2^{q_2}\cdots f_n^{q_n}$.

\begin{env}[19.5.5.1]
\label{IV.19.5.5.1}
Pour que l'application $\psi_M$ soit injective (et par suite bijective), il faut et il suffit que pour $p\geq0$ et $q\geq0$ la condition suivante soit vérifiée:

$(\mathrm I_{q,p})$ Pour toute famille $(x_{\mathbf q}^p)_{|\mathbf q|=q}$ d'éléments de $M_p$, la relation
\[
  \sum_{|\mathbf q|=q}\mathbf f^{\mathbf q}x_{\mathbf q}^p\in\sum_{i=1}^q\mathfrak J^{q-i}M_{p+i}+M_{p+q+1}'
\]
implique $x_{\mathbf q}^p\in M_{p+1}+\mathfrak JM_p$ pour tout $\mathbf q$ tel que $|\mathbf q|=q$.

On procède comme dans ($0$, 15.1.6) en considérant le sous-$A$-module $Q_k$ (resp. $Q_k'$) des termes de degré $k$ au premier (resp. second) membre de ($0$, 15.1.5.2). On munit $Q_k$ de la filtration
\[
  (Q_k)_i=\sum_{j\leq k-i}\left(\sum_{|\mathbf l|=j}(\operatorname{gr}_{k-j}(M)\otimes_A(A/\mathfrak J))T^{\mathbf l}\right)
\]
et $Q_k'$ de la filtration image formée des $\psi_M((Q_k)_i)=(Q_k')_i$; il suffit encore de prouver que les homomorphismes $(\psi_M)_{ki}:\operatorname{gr}_i(Q_k)\to\operatorname{gr}_i(Q_k')$ sont injectifs. Or, on a
\[
  \operatorname{gr}_i(Q_k)=\sum_{|\mathbf l|=k-i}((M_i/M_{i+1})\otimes_A(A/\mathfrak J))T^{\mathbf l}
  =\sum_{|\mathbf l|=k-i}(M_i/(\mathfrak JM_i+M_{i+1}))T^{\mathbf l};
\]

\oldpage[IV-4]{208}

d'autre part, $(Q_k')_{i+1}$ est l'image de $\sum_{h=0}^{k-i-1}\mathfrak J^{k-i-h-1}M_{i+h+1}$ dans $M_k'/M_{k+1}'$. Écrire que $(\psi_M)_{ki}$ est injectif équivaut donc à écrire la condition $(\mathrm I_{k-i,i})$; d'où notre assertion.
\end{env}

\begin{env}[19.5.5.2]
\label{IV.19.5.5.2}
Pour que la suite $\mathbf f$ soit $\operatorname{gr}_\bullet(M)$-quasi-régulière, il suffit que, pour $p\geq0$ et $q\geq0$, la condition suivante soit vérifiée:

$(\mathrm S_{q,p})$ Pour toute famille $(x_{\mathbf q}^p)_{|\mathbf q|=q}$ d'éléments de $M_p$, la relation
\[
  \sum_{|\mathbf q|=q}\mathbf f^{\mathbf q}x_{\mathbf q}^p\in M_{p+1}
\]
implique $x_{\mathbf q}^p\in M_{p+1}+\mathfrak JM_p$ pour tout $\mathbf q$ tel que $|\mathbf q|=q$.

Par définition, dire que $\mathbf f$ est $\operatorname{gr}_\bullet(M)$-quasi-régulière signifie que, pour tout $p\geq0$ et tout $q\geq0$, une relation
\[
  \sum_{|\mathbf q|=q}\mathbf f^{\mathbf q}x_{\mathbf q}^p\in M_{p+1}+\mathfrak J^{q+1}M_p
\]
pour une famille $(x_{\mathbf q}^p)_{|\mathbf q|=q}$ d'éléments de $M_p$ entraîne $x_{\mathbf q}^p\in M_{p+1}+\mathfrak JM_p$ pour tout $\mathbf q$ tel que $|\mathbf q|=q$. Or l'hypothèse sur la suite $(x_{\mathbf q}^p)$ signifie qu'il existe une famille $(y_{\mathbf q}^p)_{|\mathbf q|=q}$ d'éléments de $\mathfrak JM_p$ telle que l'on ait $\sum_{|\mathbf q|=q}\mathbf f^{\mathbf q}(x_{\mathbf q}^p-y_{\mathbf q}^p)\in M_{p+1}$. La condition $(\mathrm S_{q,p})$ entraîne alors $x_{\mathbf q}^p-y_{\mathbf q}^p\in M_{p+1}+\mathfrak JM_p$ pour tout $\mathbf q$ tel que $|\mathbf q|=q$, d'où $x_{\mathbf q}^p\in M_{p+1}+\mathfrak JM_p$; on voit donc que les conditions $(\mathrm S_{q,p})$ entraînent que la suite $\mathbf f$ est $\operatorname{gr}_\bullet(M)$-quasi-régulière.
\end{env}

Il suffit donc de prouver la proposition suivante:

\begin{env}[19.5.5.3]
\label{IV.19.5.5.3}
Pour tout $r\geq0$, si les conditions $(\mathrm I_{q,p})$ sont vérifiées pour $p+q<r$, alors les conditions $(\mathrm S_{q,p})$ sont vérifiées pour $p+q<r$.

Introduisons les conditions

$(\mathrm A_{q,p})$ Pour toute famille $(x_{\mathbf q}^p)_{|\mathbf q|=q}$ d'éléments de $M_p$, la relation
\[
  \sum_{|\mathbf q|=q}\mathbf f^{\mathbf q}x_{\mathbf q}^p\in M_{p+1}
\]
implique
\[
  \sum_{|\mathbf q|=q}\mathbf f^{\mathbf q}x_{\mathbf q}^p\in\sum_{i=1}^q\mathfrak J^{q-i}M_{p+i}.
\]

Il est clair que la conjonction de $(\mathrm A_{q,p})$ et de $(\mathrm I_{q,p})$ implique $(\mathrm S_{q,p})$. D'autre part, les conditions $(\mathrm A_{1,p})$ et $(\mathrm S_{0,p})$ sont triviales. Par suite, (19.5.5.3) résultera de la proposition suivante:
\end{env}

\begin{env}[19.5.5.4]
\label{IV.19.5.5.4}
Quels que soient $p\geq0$ et $q\geq1$, la condition $(\mathrm A_{q+1,p})$ est entraînée par la conjonction des conditions $(\mathrm A_{q,p})$, $(\mathrm I_{q-1,p+1})$, $(\mathrm I_{q-2,p+2})$, $\ldots$, $(\mathrm I_{0,p+q})$.

En effet, une fois prouvée cette proposition, les conditions $(\mathrm I_{q,p})$ supposées vérifiées pour $p+q<r$ entraîneront, pour chaque $p\leq r$, la condition $(\mathrm A_{q,p})$ pour $1\leq q<r-p$, par récurrence sur $q$.

Prouvons donc (19.5.5.4). Notons que la condition $(\mathrm A_{q,p})$ équivaut aussi à
\end{env}

\begin{env}[19.5.5.5]
\label{IV.19.5.5.5}
\[
  \mathfrak J^qM_p\cap M_{p+1}\subset\sum_{i=1}^q\mathfrak J^{q-i}M_{p+i}.
\]

Soit donc $m\in\mathfrak J^{q+1}M_p\cap M_{p+1}\subset\mathfrak J^qM_p\cap M_{p+1}$. Appliquant la condition $(\mathrm A_{q,p})$, on voit qu'il existe, pour chaque $i$ tel que $1\leq i\leq q$, une famille $(y_{\mathbf q}^{p+i})_{|\mathbf q|=q-i}$ d'éléments de $M_{p+i}$ telle que
\[
  m=\sum_{i=1}^q\left(\sum_{|\mathbf q|=q-i}\mathbf f^{\mathbf q}y_{\mathbf q}^{p+i}\right).
\]
Supposons que, pour $1\leq i<j$ ($j\geq1$), on ait prouvé que
\[
  y_{\mathbf q}^{p+i}\in\mathfrak JM_{p+i}+M_{p+i+1}\qquad(\text{pour }|\mathbf q|=q-i).
\]
On en déduit par définition (19.5.4)
\[
  m-\sum_{i<j}\left(\sum_{|\mathbf q|=q-i}\mathbf f^{\mathbf q}y_{\mathbf q}^{p+i}\right)\in M_{p+q+1}'.
\]
Puisque par hypothèse $(\mathrm I_{q-j,p+j})$ est vérifiée, on en déduit
\[
  y_{\mathbf q}^{p+j}\in\mathfrak JM_{p+j}+M_{p+j+1}
\]

\oldpage[IV-4]{209}

et, par récurrence sur $j$, on voit donc que cette condition est vérifiée pour tout $j$ tel que $1\leq j\leq q$. On a donc
\[
  m\in\sum_{i=1}^q\mathfrak J^{q-i}(\mathfrak JM_{p+i}+M_{p+i+1})=\sum_{i=1}^{q+1}\mathfrak J^{q-i+1}M_{p+i}
\]
et on a ainsi prouvé $(\mathrm A_{q+1,p})$ et terminé la démonstration de (19.5.5).
\end{env}

\begin{remarks}[19.5.6]
\label{IV.19.5.6}
\begin{enumerate}
  \item[(i)] Les résultats de ce numéro gardent un sens et sont valables lorsque, au lieu d'un $A$-module $M$ et d'éléments $f_i\in A$, on prend un objet $M$ d'une catégorie abélienne quelconque, et $n$ endomorphismes $f_i$ de $M$ qui commutent deux à deux; les démonstrations s'adaptent sans peine à ce cas plus général. L'hypothèse de séparation pour un $A$-module $N$ relativement à la topologie $\mathfrak J$-préadique doit ici être remplacée par l'hypothèse que tout sous-objet de $N$ contenu dans tous les $\sum_i f_i^r(N)$ ($r$ entier arbitraire) est nécessairement nul. La même remarque s'applique aux résultats de 19.7.
  \item[(ii)] Supposons que $A$ soit un anneau \emph{gradué} à degrés $\geq0$. Si $M$ (resp. chaque $\operatorname{gr}_p(M)$) est un $A$-module gradué à degrés bornés inférieurement et si les $f_i$ ne contiennent pas de composant homogène de degré $0$, l'hypothèse de séparation de ($0$, 15.1.9) est \emph{ipso facto} vérifiée pour $M$ (resp. chaque $\operatorname{gr}_p(M)$), et on voit donc que dans ce cas on a l'implication $c)\Rightarrow a)$ dans (19.5.1) (resp. (19.5.5)).
  \item[(iii)] Rappelons que, sans hypothèse de séparation, l'implication $c)\Rightarrow b)$ n'est plus valable même pour $n=1$ ($0$, 15.1.12, (iii)).
\end{enumerate}
\end{remarks}

\subsection{Suites régulières relativement à un module filtré quotient.}
\EGAReaderTarget{ega4.c04.subsec.19-6}{subsection:IV.19.6}

\begin{env}[19.6.1]
\label{IV.19.6.1}
Soient $A$ un anneau, $\mathbf f=(f_1,\ldots,f_n)$ une suite finie d'éléments de $A$, $\mathfrak J=f_1A+\cdots+f_nA$ l'idéal qu'elle engendre, et considérons une suite exacte de $A$-modules
\[
  0\longrightarrow R\xrightarrow{j}N\xrightarrow{p}M\longrightarrow0
\]
où l'on suppose que $N$ est muni d'une filtration décroissante $(N_k)$ formée de sous-$A$-modules, avec $N_0=N$. On pose $R_k=R\cap N_k$ (filtration induite par $(N_k)$), $M_k=p(N_k)$ (filtration quotient de $(N_k)$), et on désigne par $\operatorname{gr}_\bullet(N)$, $\operatorname{gr}_\bullet(R)$ et $\operatorname{gr}_\bullet(M)$ les $A$-modules gradués associés à ces trois filtrations. On pose d'autre part pour tout $k$,
\[
  N_k'=N_k+\mathfrak JN_{k-1}+\cdots+\mathfrak J^{k-1}N_1+\mathfrak J^kN_0
\]
de sorte que $M_k'=p(N_k')=M_k+\mathfrak JM_{k-1}+\cdots+\mathfrak J^{k-1}M_1+\mathfrak J^kM_0$; on pose d'autre part $R_k'=R\cap N_k'$ (filtration \emph{induite} par $(N_k')$), et on note $\operatorname{gr}_\bullet'(N)$, $\operatorname{gr}_\bullet'(M)$ et $\operatorname{gr}_\bullet'(R)$ les $A$-modules gradués associés à ces trois filtrations; on a donc un diagramme commutatif de suites exactes
\begin{equation}
\xymatrix{
  0 \ar[r] & \operatorname{gr}_\bullet(R) \ar[r]^-{\operatorname{gr}(j)} \ar[d] & \operatorname{gr}_\bullet(N) \ar[r]^-{\operatorname{gr}(p)} \ar[d] & \operatorname{gr}_\bullet(M) \ar[r] \ar[d] & 0 \\
  0 \ar[r] & \operatorname{gr}_\bullet'(R) \ar[r]^-{\operatorname{gr}'(j)} & \operatorname{gr}_\bullet'(N) \ar[r]^-{\operatorname{gr}'(p)} & \operatorname{gr}_\bullet'(M) \ar[r] & 0
}
\tag{19.6.1.1}\label{IV.19.6.1.1}
\end{equation}

\oldpage[IV-4]{210}

où les flèches verticales sont les homomorphismes canoniques d'un module gradué associé à une filtration dans le module gradué associé à une filtration moins fine.

On notera que si l'on pose
\[
  R_k''=R_k+\mathfrak JR_{k-1}+\cdots+\mathfrak J^{k-1}R_1+\mathfrak J^kR_0
\]
on a évidemment $R_k''\subset R_k'$, mais les filtrations $(R_k'')$ et $(R_k')$ sont en général distinctes; on notera $\operatorname{gr}_\bullet''(R)$ le $A$-module gradué associé à la filtration $(R_k'')$.
\end{env}

\begin{env}[19.6.2]
\label{IV.19.6.2}
On a défini en ($0$, 15.1.5.2) les homomorphismes surjectifs gradués de degré $0$
\begin{align*}
  \psi_N&:(\operatorname{gr}_\bullet(N)\otimes_A(A/\mathfrak J))[T_1,\ldots,T_n]\longrightarrow\operatorname{gr}_\bullet'(N),\\
  \psi_M&:(\operatorname{gr}_\bullet(M)\otimes_A(A/\mathfrak J))[T_1,\ldots,T_n]\longrightarrow\operatorname{gr}_\bullet'(M),\\
  \psi_R&:(\operatorname{gr}_\bullet(R)\otimes_A(A/\mathfrak J))[T_1,\ldots,T_n]\longrightarrow\operatorname{gr}_\bullet''(R),
\end{align*}
dont les deux premiers sont déduits des deux dernières flèches verticales de (19.6.1.1).

Comme la filtration $(R_k'')$ est plus fine que $(R_k')$ sur $R$, on a un homomorphisme canonique $\gamma:\operatorname{gr}_\bullet''(R)\to\operatorname{gr}_\bullet'(R)$; nous désignerons par $\psi_R'$ l'homomorphisme composé $\gamma\circ\psi_R$. Il résulte alors aussitôt des définitions que l'on a un diagramme commutatif
\begin{equation}
\xymatrix{
  & 0 \ar[d] \\
  (\operatorname{gr}_\bullet(R)\otimes_A(A/\mathfrak J))[T_1,\ldots,T_n] \ar[r]^-{\psi_R'} \ar[d]^{j'} & \operatorname{gr}_\bullet'(R) \ar[d]^{\operatorname{gr}'(j)} \\
  (\operatorname{gr}_\bullet(N)\otimes_A(A/\mathfrak J))[T_1,\ldots,T_n] \ar[r]^-{\psi_N} \ar[d]^{p'} & \operatorname{gr}_\bullet'(N) \ar[d]^{\operatorname{gr}'(p)} \\
  (\operatorname{gr}_\bullet(M)\otimes_A(A/\mathfrak J))[T_1,\ldots,T_n] \ar[r]^-{\psi_M} \ar[d] & \operatorname{gr}_\bullet'(M) \ar[d] \\
  0 & 0
}
\tag{19.6.2.1}\label{IV.19.6.2.1}
\end{equation}
où les colonnes sont exactes.
\end{env}

\begin{proposition}[19.6.3]
\label{IV.19.6.3}
Avec les notations précédentes, supposons que la suite $\mathbf f$ soit $\operatorname{gr}_\bullet(N)$-régulière. Considérons les propriétés suivantes:
\begin{enumerate}
  \item[a)] $\mathbf f$ est $\operatorname{gr}_\bullet(M)$-régulière.
  \item[b)] $\psi_M$ est bijectif.
  \item[c)] $\psi_R'$ est surjectif (autrement dit $\operatorname{gr}_\bullet'(R)$ est engendré comme $\operatorname{gr}_{\mathfrak J}^*(A)$-module par $\operatorname{Im}(\operatorname{gr}_\bullet(R)\to\operatorname{gr}_\bullet'(R))$).
  \item[d)] $\psi_R'$ est injectif.
  \item[d$'$)] L'homomorphisme $(\psi_R')_0:\operatorname{gr}_\bullet(R)\otimes_A(A/\mathfrak J)\to\operatorname{gr}_\bullet'(R)$ obtenu en restreignant $\psi_R'$ aux polynômes de degré $0$, est injectif.
  \item[e)] $\psi_R'$ est bijectif.
  \item[f)] $j'$ est injectif.

\oldpage[IV-4]{211}

  \item[f$'$)] L'homomorphisme $\operatorname{gr}_\bullet(R)\otimes_A(A/\mathfrak J)\to\operatorname{gr}_\bullet(N)\otimes_A(A/\mathfrak J)$ obtenu en restreignant $j'$ aux polynômes de degré $0$, est injectif.
  \item[g)] La filtration $\mathfrak J$-préadique de $\operatorname{gr}_\bullet(R)$ est induite par celle de $\operatorname{gr}_\bullet(N)$.
  \item[h)] Les filtrations $(R_k')$ et $(R_k'')$ sur $R$ sont identiques.
\end{enumerate}
On a alors les implications suivantes:
\[
\xymatrix@C=2.4em@R=1.8em{
  a) \ar@{=>}[r] \ar@{=>}[dr] & e) \ar@{=>}[r] \ar@{=>}[dr] & b) \ar@{<=>}[r] & c) \ar@{<=>}[r] & h) \\
  & g) \ar@{=>}[r] & d) \ar@{<=>}[r] & d') \ar@{<=>}[r] & f) \ar@{<=>}[r] & f')
}
\]
En outre, lorsque tout module quotient d'un sous-module d'un $\operatorname{gr}_k(M)$ est séparé pour la topologie $\mathfrak J$-préadique (ce qui sera le cas lorsque $A$ est noethérien, que les $f_i$ appartiennent au radical de $A$ et que les $\operatorname{gr}_k(M)$ sont des $A$-modules de type fini), alors les conditions $a)$ à $h)$ sont toutes équivalentes.
\end{proposition}

Notons que d'après (19.5.5), de l'hypothèse que $\mathbf f$ est $\operatorname{gr}_\bullet(N)$-régulière, il résulte que $\psi_N$ est \emph{bijectif}; comme $\operatorname{gr}'(j)$ est injectif, l'équivalence de $d)$ et $f)$ résulte du diagramme (19.6.2.1), ainsi que celle de $d')$ et $f')$; par ailleurs $f)$ et $f')$ sont trivialement équivalentes, ce qui montre l'équivalence de $d)$, $d')$, $f)$ et $f')$.

Comme on sait déjà que $\psi_M$ est surjectif et $\psi_N$ bijectif, l'équivalence de $b)$ et $c)$ découle aussi du diagramme (19.6.2.1) par un cas particulier du lemme des cinq.

La relation $\psi_R'=\gamma\circ\psi_R$ (19.6.2) et le fait que $\psi_R$ est surjectif montrent que la condition $c)$ équivaut à dire que $\gamma$ est surjectif, ce qui équivaut encore à la condition $h)$; on a donc prouvé l'équivalence de $b)$, $c)$ et $h)$.

Il est trivial que $e)$ équivaut à la conjonction de $c)$ et $d)$, donc aussi de $b)$ et $d)$. Notons maintenant que $a)$ implique $b)$ et que $b)$ implique $a)$ sous l'hypothèse de séparation (19.5.5). D'autre part, les implications $a)\Rightarrow g)\Rightarrow f')$, et l'implication $f')\Rightarrow a)$ sous l'hypothèse de séparation, résultent de (19.5.3). On a donc prouvé que $a)$ entraîne $e)$ (équivalente à la conjonction de $b)$ et $f')$), et aussi que toutes les conditions sont équivalentes sous l'hypothèse de séparation. C.Q.F.D.

\begin{corollary}[19.6.4]
\label{IV.19.6.4}
Soient $\mathfrak R$ un idéal de $A$, $\mathfrak L=\mathfrak J+\mathfrak R$. Supposons que la filtration $(N_k)$ (donc aussi $(M_k)$) soit la filtration $\mathfrak R$-préadique, de sorte que les filtrations $(M_k')$ et $(N_k')$ sont les filtrations $\mathfrak L$-préadiques. Considérons $\operatorname{gr}_\bullet(R)$ (resp. $\operatorname{gr}_\bullet'(R)$) comme un module gradué sur $\operatorname{gr}_{\mathfrak R}^*(A)$ (resp. $\operatorname{gr}_{\mathfrak L}^*(A)$).
\begin{enumerate}
  \item[(i)] Soit $S$ une partie de $\operatorname{gr}_\bullet(R)$ qui engendre $\operatorname{gr}_\bullet(R)$ comme $\operatorname{gr}_{\mathfrak R}^*(A)$-module. Alors la condition $c)$ de (19.6.3) équivaut à la condition suivante:
  \item[c$'$)] L'image de $S$ dans $\operatorname{gr}_\bullet'(R)$ par $\psi_R'$ engendre $\operatorname{gr}_\bullet'(R)$ comme $\operatorname{gr}_{\mathfrak L}^*(A)$-module.
  \item[(ii)] Supposons vérifiée la condition $e)$ de (19.6.3), et de plus que les quotients des $\operatorname{gr}_k(R)$ soient séparés pour la topologie $\mathfrak J$-préadique (ce qui sera le cas lorsque $A$ est noethérien, les $f_i$ dans le radical de $A$ et les $\operatorname{gr}_k(N)$ des $A$-modules de type fini). Alors, pour qu'une partie $S$ de $\operatorname{gr}_\bullet(R)$ formée d'éléments homogènes engendre $\operatorname{gr}_\bullet(R)$ comme $\operatorname{gr}_{\mathfrak R}^*(A)$-module, il faut et il suffit que son image par $\psi_R'$ engendre $\operatorname{gr}_\bullet'(R)$ comme $\operatorname{gr}_{\mathfrak L}^*(A)$-module.
\end{enumerate}
\end{corollary}

\oldpage[IV-4]{212}

\begin{enumerate}
  \item[(i)] Considérons l'homomorphisme de degré $0$
  \[
    \psi_A:(\operatorname{gr}_{\mathfrak R}^*(A)\otimes_A(A/\mathfrak J))[T_1,\ldots,T_n]\longrightarrow\operatorname{gr}_\bullet'(A)=\operatorname{gr}_{\mathfrak L}^*(A)
  \]
  ($0$, 15.1.5.2). Le fait que cet homomorphisme soit \emph{surjectif} entraîne que la sous-$\operatorname{gr}_{\mathfrak L}^*(A)$-algèbre de $\operatorname{gr}_\bullet'(R)$ engendrée par $S'=\psi_R'(S)$ n'est autre que $\operatorname{Im}(\psi_R')$, d'où l'équivalence de $c)$ et $c')$.
  \item[(ii)] Puisque $e)$ est vérifiée, il en est de même de $c)$, et en vertu de (i), il reste à prouver la suffisance de la condition de l'énoncé. Or, comme $\psi_R'$ est bijectif, dire que $S'$ engendre $\operatorname{gr}_\bullet'(R)$ considéré comme $\operatorname{gr}_{\mathfrak L}^*(A)$-module équivaut à dire que $S$ engendre $(\operatorname{gr}_\bullet(R)\otimes_A(A/\mathfrak J))[T_1,\ldots,T_n]$ comme $\operatorname{gr}_{\mathfrak L}^*(A)$-module, et en vertu de la surjectivité de $\psi_A$, on peut dans cet énoncé remplacer l'anneau $\operatorname{gr}_{\mathfrak L}^*(A)$ par $(\operatorname{gr}_{\mathfrak R}^*(A)\otimes_A(A/\mathfrak J))[T_1,\ldots,T_n]$. Il revient évidemment au même de dire que $S$ engendre $\operatorname{gr}_\bullet(R)\otimes_A(A/\mathfrak J)$ en tant que module sur l'anneau $\operatorname{gr}_{\mathfrak R}^*(A)\otimes_A(A/\mathfrak J)$, ou aussi en tant que $\operatorname{gr}_{\mathfrak R}^*(A)$-module. Or, si $T$ est le sous-$\operatorname{gr}_{\mathfrak R}^*(A)$-module gradué de $\operatorname{gr}_\bullet(R)$ engendré par $S$, et si l'on pose $T'=\operatorname{gr}_\bullet(R)/T$, on a par hypothèse $T'/\mathfrak JT'=0$, ou $\mathfrak JT'=T'$; mais $T'$ est un module gradué, et chaque $T_p'$ est un quotient de $\operatorname{gr}_p(R)$, donc séparé pour la topologie $\mathfrak J$-préadique; l'hypothèse $\mathfrak JT_p'=T_p'$ entraîne donc $T_p'=0$ pour tout entier $p$, donc $T'=0$, ce qui achève de prouver le corollaire.
\end{enumerate}

\subsection{Critère de platitude normale de Hironaka.}
\EGAReaderTarget{ega4.c04.subsec.19-7}{subsection:IV.19.7}
\label{IV.19.7}

\begin{theorem}[19.7.1]
\label{IV.19.7.1}
\textup{(Hironaka).} Soient $A$ un anneau, $\mathfrak R$ un idéal de $A$, $\mathbf f=(f_1,\ldots,f_n)$ une suite d'éléments de $A$ qui est $(A/\mathfrak R)$-régulière, $\mathfrak J=f_1A+\cdots+f_nA$ l'idéal engendré par $\mathbf f$, $\mathfrak L=\mathfrak J+\mathfrak R$, $M$ un $A$-module.

On a dans ce cas
\[
  (\operatorname{gr}_{\mathfrak R}^*(M))\otimes_A(A/\mathfrak J)
  =(\operatorname{gr}_{\mathfrak R}^*(M))\otimes_A(A/\mathfrak L)
  =(\operatorname{gr}_{\mathfrak R}^*(M))\otimes_{A/\mathfrak R}(A/\mathfrak L),
\]
car
\[
  (\mathfrak R^nM/\mathfrak R^{n+1}M)\otimes_A(A/\mathfrak J)
  =\mathfrak R^nM/(\mathfrak R^{n+1}+\mathfrak J\mathfrak R^n)M
  =\mathfrak R^nM/(\mathfrak R^{n+1}+\mathfrak L\mathfrak R^n)M,
\]
et la dernière égalité résulte de Bourbaki, \emph{Alg.}, chap. II, 3\textsuperscript{e} éd., § 3, n\textsuperscript{o} 7, cor. 2 de la prop. 6.

Considérons les conditions suivantes:
\begin{enumerate}
  \item[a)] $\operatorname{gr}_{\mathfrak R}^*(M)$ est un $(A/\mathfrak R)$-module plat.
  \item[b)] $\operatorname{gr}_{\mathfrak L}^*(M)$ est un $(A/\mathfrak L)$-module plat et l'homomorphisme canonique ($0$, 15.1.5.2)
  \[
    \psi_M:((\operatorname{gr}_{\mathfrak R}^*(M))\otimes_A(A/\mathfrak L))[T_1,\ldots,T_n]\longrightarrow\operatorname{gr}_{\mathfrak L}^*(M)
  \]
  est bijectif.
  \item[c)] $(\operatorname{gr}_{\mathfrak R}^*(M))\otimes_{A/\mathfrak R}(A/\mathfrak L)$ est un $(A/\mathfrak L)$-module plat et la suite $\mathbf f$ est $(\operatorname{gr}_{\mathfrak R}^*(M))$-régulière.
  \item[d)] $\operatorname{gr}_{\mathfrak L}^*(M)$ est un $(A/\mathfrak L)$-module plat, et pour tout $\mathfrak p\in\operatorname{Ass}_{A/\mathfrak R}(A/\mathfrak L)$, $(\operatorname{gr}_{\mathfrak R}^*(M))_{\mathfrak p}$ est un $(A/\mathfrak R)_{\mathfrak p}$-module plat.
\end{enumerate}
On a alors les implications
\[
\xymatrix@C=3em@R=1.6em{
  a) \ar@{=>}[r] \ar@{=>}[dr] & c) \ar@{=>}[r] & b) \\
  & d)
}
\]
Lorsque tout quotient d'un sous-module d'un $\operatorname{gr}_{\mathfrak R}^p(M)$ ($p\geq0$) est idéalement séparé pour $\mathfrak J$ (Bourbaki, \emph{Alg. comm.}, chap. III, § 5, n\textsuperscript{o} 1), les conditions $a)$, $b)$ et $c)$ sont équivalentes.

\oldpage[IV-4]{213}

Enfin, lorsque $A$ est noethérien, que les $f_i$ appartiennent au radical de $A$, que la suite $\mathbf f$ est $\operatorname{gr}_{\mathfrak R}^*(A)$-régulière et que $M$ est un $A$-module de type fini, les conditions $a)$, $b)$, $c)$, $d)$ sont équivalentes.
\end{theorem}

L'implication $a)\Rightarrow c)$ est immédiate ($0$, 15.1.13). Si $c)$ est vérifiée, il résulte de (19.5.5) que $\psi_M$ est bijectif; en outre, $\operatorname{gr}_{\mathfrak R}^*(M)\otimes_A(A/\mathfrak L)=\operatorname{gr}_{\mathfrak R}^*(M)\otimes_{A/\mathfrak R}(A/\mathfrak L)$ est un $(A/\mathfrak L)$-module plat, donc il en est de même de $(\operatorname{gr}_{\mathfrak R}^*(M)\otimes_A(A/\mathfrak L))[T_1,\ldots,T_n]$, et par suite de $\operatorname{gr}_{\mathfrak L}^*(M)$ puisque $\psi_M$ est un $(A/\mathfrak L)$-isomorphisme; donc $c)$ entraîne $b)$. Le fait que $a)$ implique $b)$ montre aussitôt que $a)$ implique aussi $d)$.

Lorsque tout quotient d'un $\operatorname{gr}_{\mathfrak R}^p(M)$ est idéalement séparé pour $\mathfrak J$, il résulte de (19.5.5) que la condition $b)$ entraîne que $\mathbf f$ est $(\operatorname{gr}_{\mathfrak R}^*(M))$-régulière; en outre, comme $(\operatorname{gr}_{\mathfrak R}^*(M)\otimes_{A/\mathfrak R}(A/\mathfrak L))[T_1,\ldots,T_n]$, isomorphe à $\operatorname{gr}_{\mathfrak L}^*(M)$, est alors un $(A/\mathfrak L)$-module plat, il en est de même de $\operatorname{gr}_{\mathfrak R}^*(M)\otimes_{A/\mathfrak R}(A/\mathfrak L)$, qui en est un facteur direct; cela prouve donc que $b)$ entraîne alors $c)$. D'autre part, sous la même hypothèse, $c)$ entraîne $a)$ en vertu de ($0$, 15.1.21) (où on peut remplacer l'hypothèse noethérienne par l'hypothèse que $M$ est idéalement séparé pour $\mathfrak J$, en vertu de ($0_{\mathrm{III}}$, 10.2.1)).

Il reste donc à prouver que lorsque $A$ est noethérien, $M$ de type fini et que les $f_i$ appartiennent au radical de $A$, $d)$ entraîne $a)$.

Il existe par hypothèse un $A$-module \emph{libre de type fini} $N$ et une suite exacte $0\to R\to N\to M\to0$. Il suffit de prouver que $d)$ entraîne $b)$, puisque tout quotient d'un sous-module d'un $\operatorname{gr}_{\mathfrak R}^p(M)$ est alors idéalement séparé pour $\mathfrak J$, en vertu du fait que $\operatorname{gr}_{\mathfrak R}^p(M)$ est un $A$-module de type fini, que $\mathfrak J$ est contenu dans le radical de $A$, et que $A$ est noethérien (Bourbaki, \emph{Alg. comm.}, chap. III, § 5, n\textsuperscript{o} 1); en d'autres termes, il s'agit de voir que $\psi_M$ est \emph{bijectif}.

Comme par hypothèse la suite $\mathbf f$ est $\operatorname{gr}_{\mathfrak R}^*(A)$-régulière, elle est aussi $\operatorname{gr}_{\mathfrak R}^*(N)$-régulière; on peut par suite appliquer (19.6.3), car tout module quotient d'un sous-module d'un $\operatorname{gr}_{\mathfrak R}^k(M)$ est alors séparé pour la topologie $\mathfrak J$-préadique puisque les $f_i$ appartiennent au radical de $A$ par hypothèse. Le diagramme commutatif (19.6.2.1) s'écrit donc ici
\begin{equation}
\xymatrix{
  & 0 \ar[d] \\
  (\operatorname{gr}_{\mathfrak R}^*(R,N)\otimes_{A/\mathfrak R}(A/\mathfrak L))[T_1,\ldots,T_n] \ar[r]^-{\psi_R'} \ar[d] & \operatorname{gr}_{\mathfrak L}^*(R,N) \ar[d] \\
  (\operatorname{gr}_{\mathfrak R}^*(N)\otimes_{A/\mathfrak R}(A/\mathfrak L))[T_1,\ldots,T_n] \ar[r]^-{\psi_N} \ar[d] & \operatorname{gr}_{\mathfrak L}^*(N) \ar[d] \\
  (\operatorname{gr}_{\mathfrak R}^*(M)\otimes_{A/\mathfrak R}(A/\mathfrak L))[T_1,\ldots,T_n] \ar[r]^-{\psi_M} \ar[d] & \operatorname{gr}_{\mathfrak L}^*(M) \ar[d] \\
  0 & 0
}
\tag{19.7.1.1}\label{IV.19.7.1.1}
\end{equation}
où $\operatorname{gr}_{\mathfrak R}^*(R,N)$ (resp. $\operatorname{gr}_{\mathfrak L}^*(R,N)$) est le gradué associé à $R$ pour la filtration des $R\cap\mathfrak R^mN$ (resp. $R\cap\mathfrak L^mN$); rappelons que dans ce diagramme les colonnes sont exactes et que $\psi_N$

\oldpage[IV-4]{214}

est \emph{bijectif}. Il s'agit, en vertu de (19.6.3), de prouver que $\psi_R'$ est \emph{surjectif}, ce qui va être fait en plusieurs étapes.

\begin{env}[19.7.1.2]
\label{IV.19.7.1.2}
Posons $B=\operatorname{gr}_{\mathfrak R}^*A\otimes_{A/\mathfrak R}(A/\mathfrak L)$, $P=\operatorname{gr}_{\mathfrak R}^*(N)\otimes_{A/\mathfrak R}(A/\mathfrak L)$, $Q=\operatorname{gr}_{\mathfrak L}^*(R,N)$; puisque $\psi_N$ est bijectif, $P[T_1,\ldots,T_n]$ s'identifie à $\operatorname{gr}_{\mathfrak L}^*(N)$, et $Q$ à un sous-$B[T_1,\ldots,T_n]$-module de $P[T_1,\ldots,T_n]$; nous allons d'abord voir que l'on a
\end{env}

\begin{env}[19.7.1.3]
\label{IV.19.7.1.3}
\[
  Q=(Q\cap P)[T_1,\ldots,T_n].
\]
\end{env}

Pour cela, posons $Z=Q/((Q\cap P)[T_1,\ldots,T_n])$; c'est un sous-$B[T_1,\ldots,T_n]$-module de
\[
  P[T_1,\ldots,T_n]/((Q\cap P)[T_1,\ldots,T_n])=(P/(Q\cap P))[T_1,\ldots,T_n].
\]
En tant que $(A/\mathfrak L)$-module, il est donc isomorphe à un sous-module d'une somme directe de $(A/\mathfrak L)$-modules isomorphes à $P/(Q\cap P)=(P+Q)/Q$. Mais $(P+Q)/Q$ est un $(A/\mathfrak L)$-module isomorphe à un sous-module de $P[T_1,\ldots,T_n]/Q$, qui n'est autre que $\operatorname{gr}_{\mathfrak L}^*(M)$ en vertu du diagramme (19.7.1.1). On a donc
\begin{equation}
  \operatorname{Ass}_{A/\mathfrak L}(Z)\subset\operatorname{Ass}_{A/\mathfrak L}(\operatorname{gr}_{\mathfrak L}^*(M)).
  \tag{19.7.1.4}\label{IV.19.7.1.4}
\end{equation}

Mais chacun des $\operatorname{gr}_{\mathfrak L}^m(M)$ est par hypothèse un $(A/\mathfrak L)$-module \emph{plat} de type fini, donc projectif puisque $A/\mathfrak L$ est noethérien, et on a par suite $\operatorname{Ass}_{A/\mathfrak L}(\operatorname{gr}_{\mathfrak L}^*(M))\subset\operatorname{Ass}_{A/\mathfrak L}(A/\mathfrak L)$. Pour voir que $Z=0$, il suffit par suite (Bourbaki, \emph{Alg. comm.}, chap. IV, § 1, n\textsuperscript{o} 1, cor. 1 de la prop. 2 et n\textsuperscript{o} 3, prop. 7) de voir que pour tout $\mathfrak q\in\operatorname{Ass}_{A/\mathfrak L}(A/\mathfrak L)$, on a $Z_{\mathfrak q}=0$. Mais les idéaux de $\operatorname{Ass}_{A/\mathfrak L}(A/\mathfrak L)$ sont les idéaux de la forme $\mathfrak p/(\mathfrak L/\mathfrak R)$, où $\mathfrak p\in\operatorname{Ass}_{A/\mathfrak R}(A/\mathfrak L)$, et il revient donc au même de voir que $Z_{\mathfrak p}=0$ pour tout $\mathfrak p\in\operatorname{Ass}_{A/\mathfrak R}(A/\mathfrak L)$. Or, si $\mathfrak p=\mathfrak r/\mathfrak R$, où $\mathfrak r$ est un idéal premier de $A$, on a $(\operatorname{gr}_{\mathfrak R}^*(M))_{\mathfrak p}=\operatorname{gr}_{\mathfrak R_{\mathfrak r}}^*(M_{\mathfrak r})$ et $(A/\mathfrak L)_{\mathfrak p}=A_{\mathfrak r}/\mathfrak L_{\mathfrak r}$, et les images des $f_i$ dans $A_{\mathfrak r}$ forment une suite $\operatorname{gr}_{\mathfrak R_{\mathfrak r}}^*(A_{\mathfrak r})$-régulière ($0$, 15.1.14); appliquant à $A_{\mathfrak r}$, $\mathfrak R_{\mathfrak r}$ et $M_{\mathfrak r}$ l'implication $a)\Rightarrow b)$ de l'énoncé, on voit que l'hypothèse $d)$ entraîne que $\psi_{M_{\mathfrak r}}$ est bijectif, donc aussi $\psi_{R_{\mathfrak r}}'$ en vertu de (19.6.3); en d'autres termes, $Q_{\mathfrak r}=Q_{\mathfrak r}'[T_1,\ldots,T_n]$, où on a posé $Q'=\operatorname{gr}_{\mathfrak R}^*(R,N)\otimes_{A/\mathfrak R}(A/\mathfrak L)$, qui est ici un sous-$B$-module de $P$; en particulier $Q_{\mathfrak r}'=Q_{\mathfrak r}\cap P_{\mathfrak r}$, donc finalement on a bien $Z_{\mathfrak r}=Z_{\mathfrak p}=0$, ce qui achève de démontrer (19.7.1.3).

\begin{env}[19.7.1.5]
\label{IV.19.7.1.5}
En vertu de (19.7.1.3), pour prouver que $\psi_R'$ est surjectif, il suffit de montrer que tout élément homogène de degré $m$ de $Q\cap P$ est l'image d'un élément de même degré de $\operatorname{gr}_{\mathfrak R}^*(R,N)\otimes_{A/\mathfrak R}(A/\mathfrak L)$. Or, les éléments de degré $m$ de $P$ s'identifient (par $\psi_N$) aux éléments de $\operatorname{gr}_{\mathfrak L}^m(N)$ qui appartiennent à $(\mathfrak R^mN+\mathfrak L^{m+1}N)/\mathfrak L^{m+1}N$; ceux qui sont images d'éléments de $\operatorname{gr}_{\mathfrak R}^m(R,N)\otimes_{A/\mathfrak R}(A/\mathfrak L)$ s'identifient aux éléments de $\operatorname{gr}_{\mathfrak L}^m(N)$ appartenant à $((R\cap\mathfrak R^mN)+\mathfrak L^{m+1}N)/\mathfrak L^{m+1}N$. Il suffira donc finalement de prouver, pour tout entier $m>0$, l'inclusion
\[
  ((R\cap\mathfrak L^mN)+\mathfrak L^{m+1}N)\cap(\mathfrak R^mN+\mathfrak L^{m+1}N)\subset(R\cap\mathfrak R^mN)+\mathfrak L^{m+1}N.
\]
Comme $\mathfrak R\subset\mathfrak L$, on vérifie aussitôt que le premier membre de cette relation est égal à $(R\cap(\mathfrak R^mN+\mathfrak L^{m+1}N))+\mathfrak L^{m+1}N$, si bien que tout se réduit à prouver
\end{env}

\begin{env}[19.7.1.6]
\label{IV.19.7.1.6}
\[
  R\cap(\mathfrak R^mN+\mathfrak L^{m+1}N)\subset(R\cap\mathfrak R^mN)+\mathfrak L^{m+1}N.
\]
\end{env}

\oldpage[IV-4]{215}

Nous procéderons par récurrence sur $m$, supposant donc (19.7.1.6) vérifiée lorsqu'on remplace $m$ par un entier $m'<m$. Nous considérerons d'autre part, pour $m$ fixé et pour un entier $d\geq0$, la relation
\begin{equation}
  R\cap(\mathfrak R^mN+\mathfrak L^{m+1}N)\subset(R\cap\mathfrak R^dN\cap\mathfrak L^mN)+\mathfrak L^{m+1}N.
  \tag{$*_d$}
\end{equation}

Il est clair qu'il suffit de démontrer $(*_d)$ pour $d=m$, et d'autre part $(*_0)$ est trivialement vraie. Nous prouverons $(*_d)$ par récurrence sur $d$; autrement dit, nous supposons, pour un $d\geq0$ fixé (avec $d<m$), que $(*_d)$ est vraie, et nous voulons prouver
\begin{equation}
  R\cap(\mathfrak R^mN+\mathfrak L^{m+1}N)\subset(R\cap\mathfrak R^{d+1}N\cap\mathfrak L^mN)+\mathfrak L^{m+1}N.
  \tag{$*_{d+1}$}
\end{equation}

Considérons pour cela, pour un $h\geq0$, la relation
\begin{equation}
  R\cap(\mathfrak R^mN+\mathfrak L^{m+1}N)\subset(R\cap(\mathfrak R^{d+1}N+\mathfrak R^d\mathfrak L^hN)\cap\mathfrak L^mN)+\mathfrak L^{m+1}N.
  \tag{$**_{d+1,h}$}
\end{equation}

L'hypothèse $(*_d)$ implique que $(**_{d+1,0})$ est vraie. Nous prouverons par récurrence sur $h$ que $(**_{d+1,h})$ est vraie pour tout $h\geq0$. Mais on a le

\begin{lemma}[19.7.1.7]
\label{IV.19.7.1.7}
Soient $A$ un anneau noethérien, $N$ un $A$-module de type fini, $E$, $F$ deux sous-modules de $N$, $\mathfrak L$ un idéal de $A$. Pour tout $k>0$ il existe $h>0$ tel que
\[
  E\cap(F+\mathfrak L^hN)\subset(E\cap F)+\mathfrak L^kN.
\]
\end{lemma}

En effet, soit $\varphi:N\to N/(E\cap F)=N_1$ l'homomorphisme canonique, et posons $E_1=\varphi(E)$, $F_1=\varphi(F)$, de sorte que $E_1\cap F_1=0$ et $\varphi(E\cap(F+\mathfrak L^hN))=E_1\cap(F_1+\mathfrak L^hN_1)$. On sait ($0_{\mathrm I}$, 7.3.2) qu'il existe $h$ tel que $(E_1+F_1)\cap\mathfrak L^hN_1\subset\mathfrak L^k(E_1+F_1)$; si $h$ est ainsi choisi et si $x_1\in E_1$ est tel qu'il existe $y_1\in F_1$ pour lequel $x_1-y_1\in\mathfrak L^hN_1$, on en déduit d'après ce qui précède qu'il existe $u_1\in\mathfrak L^kE_1$ et $v_1\in\mathfrak L^kF_1$ tels que $x_1=u_1$ et $y_1=v_1$; donc $E_1\cap(F_1+\mathfrak L^hN_1)\subset\mathfrak L^kN_1$, ce qui prouve le lemme.

Il suffit alors d'appliquer ce lemme avec $E=R\cap\mathfrak L^mN$, $F=\mathfrak R^{d+1}N$ et de prendre $k=m+1$ pour en déduire qu'avec la valeur de $h$ correspondante, $(**_{d+1,h})$ entraîne $(*_{d+1})$.

Nous supposons donc dans tout ce qui suit que $(**_{d+1,h})$ est vérifiée.

\begin{env}[19.7.1.8]
\label{IV.19.7.1.8}
Démonstration de $(**_{d+1,h+1})$: Premier cas: $h\leq m-d$.
\end{env}

Partons donc d'un élément
\begin{env}[19.7.1.9]
\label{IV.19.7.1.9}
\[
  g\in R\cap(\mathfrak R^{d+1}N+\mathfrak R^d\mathfrak L^hN)\cap\mathfrak L^mN
\]
\end{env}
congru mod. $\mathfrak L^{m+1}N$ à un élément $g_1\in R\cap(\mathfrak R^mN+\mathfrak L^{m+1}N)$; on a donc
\begin{env}[19.7.1.10]
\label{IV.19.7.1.10}
\[
  g\in\mathfrak R^mN+\mathfrak L^{m+1}N.
\]
\end{env}

Il suffira de prouver que $g$ est aussi congru mod. $\mathfrak L^{m+1}N$ à un élément $g_2\in R\cap(\mathfrak R^{d+1}N+\mathfrak R^d\mathfrak L^{h+1}N)$, car cela entraînera $g_2\in\mathfrak L^mN$, puisque l'on aura $g\in\mathfrak L^mN$ et $g_2-g\in\mathfrak L^{m+1}N$. Dans le cas $h\leq m-d$ que nous considérons, il suffit de démontrer la relation
\begin{env}[19.7.1.11]
\label{IV.19.7.1.11}
\[
  (\mathfrak R^mN+\mathfrak L^{m+1}N)\cap\mathfrak R^dN\subset\mathfrak L^{m-d+1}\mathfrak R^dN+\mathfrak R^{d+1}N.
\]
\end{env}

\oldpage[IV-4]{216}

car les relations (19.7.1.9) et (19.7.1.10) entraînent que $g$ appartient au premier membre de (19.7.1.11), et comme $h\leq m-d$, on a $\mathfrak L^{m-d+1}\subset\mathfrak L^{h+1}$, ce qui établira $(**_{d+1,h+1})$. D'ailleurs on a $\mathfrak R^mN\subset\mathfrak R^dN$, et la relation
\begin{env}[19.7.1.12]
\label{IV.19.7.1.12}
\[
  \mathfrak L^{m+1}N\cap\mathfrak R^dN\subset\mathfrak L^{m-d+1}\mathfrak R^dN
\]
\end{env}
entraînera (19.7.1.11). Comme $N$ est un $A$-module libre de type fini, on voit qu'il suffit donc de prouver
\begin{env}[19.7.1.13]
\label{IV.19.7.1.13}
\[
  \mathfrak L^{m+1}\cap\mathfrak R^d\subset\mathfrak L^{m-d+1}\mathfrak R^d.
\]
\end{env}

Il suffira de prouver que, de façon générale, pour $d<m$, on a
\begin{env}[19.7.1.14]
\label{IV.19.7.1.14}
\[
  \mathfrak L^{m+1}\cap\mathfrak R^d\subset\mathfrak L^{m-d+1}\mathfrak R^d+\mathfrak R^{d+1}.
\]
\end{env}
En effet, comme $\mathfrak R\subset\mathfrak L$, donc $\mathfrak L^{m-d+1}\mathfrak R^d\subset\mathfrak L^{m+1}$, la relation précédente entraînera
\[
  \mathfrak L^{m+1}\cap\mathfrak R^d\subset\mathfrak L^{m-d+1}\mathfrak R^d+\mathfrak L^{m+1}\cap\mathfrak R^{d+1}
\]
et par récurrence sur $d\leq m$, on en conclura
\[
  \mathfrak L^{m+1}\cap\mathfrak R^d\subset\mathfrak L^{m-d+1}\mathfrak R^d+\mathfrak R^{m+1}
\]
d'où (19.7.1.13), puisque $\mathfrak R^{m+1}\subset\mathfrak L^{m-d+1}\mathfrak R^d$.

Comme $\mathfrak L=\mathfrak J+\mathfrak R$, la relation (19.7.1.14) s'écrit aussi
\begin{env}[19.7.1.15]
\label{IV.19.7.1.15}
\[
\begin{aligned}
  &(\mathfrak J^{m+1}+\mathfrak J^m\mathfrak R+\cdots+\mathfrak J\mathfrak R^m+\mathfrak R^{m+1})\cap\mathfrak R^d\\
  &\quad\subset\mathfrak J^{m-d+1}\mathfrak R^d+\mathfrak J^{m-d}\mathfrak R^{d+1}+\cdots+\mathfrak J\mathfrak R^m+\mathfrak R^{m+1}+\mathfrak R^{d+1}\\
  &\quad=\mathfrak J^{m-d+1}\mathfrak R^d+\mathfrak R^{d+1}.
\end{aligned}
\]
\end{env}

Nous allons voir que cette inclusion est elle-même conséquence de
\begin{env}[19.7.1.16]
\label{IV.19.7.1.16}
\[
  \mathfrak R^a\mathfrak J^b\cap\mathfrak R^{a+1}\subset\mathfrak R^{a+1}\mathfrak J^{b-1}
\]
\end{env}
valable pour $a\geq0$, $b\geq1$. En effet, il suffit, pour prouver (19.7.1.15) de montrer que pour $0<q\leq d$, on a
\begin{env}[19.7.1.17]
\label{IV.19.7.1.17}
\[
\begin{aligned}
  &(\mathfrak J^{m+1}+\mathfrak J^m\mathfrak R+\cdots+\mathfrak J\mathfrak R^m+\mathfrak R^{m+1})\cap\mathfrak R^d\\
  &\quad\subset(\mathfrak J^{m+1-q}\mathfrak R^q+\mathfrak J^{m-q}\mathfrak R^{q+1}+\cdots+\mathfrak R^{m+1})\cap\mathfrak R^d
\end{aligned}
\]
\end{env}
car pour $q=d$, cela entraînera (19.7.1.15). Or, pour prouver (19.7.1.17), il suffit de procéder par récurrence sur $q$ en supposant la relation vraie pour $q<d$; un élément du premier membre s'écrit donc par hypothèse $y+z$ avec $y\in\mathfrak J^{m+1-q}\mathfrak R^q$, $z\in\mathfrak J^{m-q}\mathfrak R^{q+1}+\cdots+\mathfrak R^{m+1}$; comme $y+z\in\mathfrak R^d\subset\mathfrak R^{q+1}$, on en déduit $y\in\mathfrak R^{q+1}$, puisque $z\in\mathfrak R^{q+1}$; tenant compte de (19.7.1.16), on a $y\in\mathfrak J^{m-q}\mathfrak R^{q+1}$, donc
\[
  y+z\in\mathfrak J^{m-q}\mathfrak R^{q+1}+\cdots+\mathfrak R^{m+1},\qquad\text{et}\qquad y+z\in\mathfrak R^d
\]
par hypothèse, ce qui prouve (19.7.1.17) où $q$ a été remplacé par $q+1$.

Reste donc à établir (19.7.1.16). Un élément du premier membre s'écrit $a_1f_1+\cdots+a_nf_n$, où les $a_i$ appartiennent à $\mathfrak R^a\mathfrak J^{b-1}$; tenant compte de ce que la suite $\mathbf f$ est $(\mathfrak R^a/\mathfrak R^{a+1})$-régulière par hypothèse, donc aussi $(\mathfrak R^a/\mathfrak R^{a+1})$-quasi-régulière ($0$, 15.1.9), il résulte de la définition ($0$, 15.1.7) que la relation $a_1f_1+\cdots+a_nf_n\in\mathfrak R^{a+1}$, où les $a_i\in\mathfrak R^a$, entraîne nécessairement $a_i\in\mathfrak R^{a+1}$, donc $a_i\in\mathfrak R^a\mathfrak J^{b-1}\cap\mathfrak R^{a+1}$; il suffit de raisonner par

\oldpage[IV-4]{217}

récurrence sur $b$ pour $b\geq1$ (le cas $b=1$ étant trivial) pour achever de prouver (19.7.1.16) et par suite aussi $(**_{d+1,h+1})$ pour $h\leq m-d$.

\begin{env}[19.7.1.18]
\label{IV.19.7.1.18}
Démonstration de $(**_{d+1,h+1})$; Deuxième cas: $h>m-d$.
\end{env}

Nous allons tout d'abord démontrer que, pour tout $h\geq0$, on a
\begin{env}[19.7.1.19]
\label{IV.19.7.1.19}
\[
  R\cap(\mathfrak R^{d+1}N+\mathfrak R^d\mathfrak L^hN)\subset\mathfrak J^h(R\cap\mathfrak R^dN)+\mathfrak L^{h+1}\mathfrak R^dN+\mathfrak R^{d+1}N.
\]
\end{env}

Considérons pour cela un élément $z$ du premier membre de (19.7.1.19). Notons que l'on a $\mathfrak R^{d+1}N+\mathfrak R^d\mathfrak L^hN=\mathfrak R^{d+1}N+\mathfrak R^d\mathfrak J^hN$ puisque $\mathfrak L=\mathfrak R+\mathfrak J$; nous allons considérer la classe $\bar z$ de $z$ dans $\operatorname{gr}_{\mathfrak J}^h(\operatorname{gr}_{\mathfrak R}^d(N))$, qui, en vertu de (19.5.5) et de l'hypothèse sur $\mathbf f$, s'identifie à un sous-$(A/\mathfrak L)$-module de $\operatorname{gr}_{\mathfrak L}^{h+d}(N)$. Du fait que $z\in R$, nous allons montrer que l'on a même (avec l'identification précédente)
\begin{env}[19.7.1.20]
\label{IV.19.7.1.20}
\[
  \bar z\in\operatorname{gr}_{\mathfrak L}^{h+d}(R,N).
\]
\end{env}

En effet, on a noté, dans la démonstration de (19.7.1.3) que, pour tout idéal premier $\mathfrak r$ de $A$ tel que $\mathfrak p=\mathfrak r/\mathfrak R\in\operatorname{Ass}_{A/\mathfrak R}(A/\mathfrak L)$, l'application $\psi_{R_{\mathfrak r}}'$ est bijective, et par suite la relation (19.7.1.6), où l'on remplace tous les modules par leurs localisés en $\mathfrak r$, est vraie. Comme $z$ appartient à $R\cap(\mathfrak R^dN+\mathfrak L^{d+1}N)$, son image $z/1$ dans $R_{\mathfrak r}$ appartient à $(R_{\mathfrak r}\cap\mathfrak R_{\mathfrak r}^dN_{\mathfrak r})+\mathfrak L_{\mathfrak r}^{d+1}N_{\mathfrak r}$, donc l'image $\bar z/1$ de $\bar z$ dans $\operatorname{gr}_{\mathfrak L_{\mathfrak r}}^{h+d}(N_{\mathfrak r})$ appartient à
\[
  \operatorname{gr}_{\mathfrak L_{\mathfrak r}}^{h+d}(R_{\mathfrak r},N_{\mathfrak r})
  =(\operatorname{gr}_{\mathfrak L}^{h+d}(R,N))_{\mathfrak r}
  =(\operatorname{gr}_{\mathfrak L}^{h+d}(R,N))_{\mathfrak q},
\]
où $\mathfrak q=\mathfrak r/\mathfrak L$. En d'autres termes l'image de $\bar z$ par l'application canonique
\[
  (\operatorname{gr}_{\mathfrak L}^{h+d}(N))_{\mathfrak q}\longrightarrow(\operatorname{gr}_{\mathfrak L}^{h+d}(M))_{\mathfrak q}
\]
est nulle pour tout $\mathfrak q\in\operatorname{Ass}_{A/\mathfrak L}(A/\mathfrak L)$. Comme on a vu dans la démonstration de (19.7.1.3) que
\[
  \operatorname{Ass}_{A/\mathfrak L}(\operatorname{gr}_{\mathfrak L}^{h+d}(M))\subset\operatorname{Ass}_{A/\mathfrak L}(A/\mathfrak L),
\]
on en déduit bien (Bourbaki, \emph{Alg. comm.}, chap. IV, § 1, n\textsuperscript{o} 1, cor. 1 de la prop. 2 et n\textsuperscript{o} 3, prop. 7) que l'image de $\bar z$ par l'application canonique $\operatorname{gr}_{\mathfrak L}^{h+d}(N)\to\operatorname{gr}_{\mathfrak L}^{h+d}(M)$ est nulle, c'est-à-dire que l'on a la relation (19.7.1.20).

Cela étant, on peut écrire par définition
\begin{env}[19.7.1.21]
\label{IV.19.7.1.21}
\[
  z\equiv\sum_{|\mathbf p|=h}c_{\mathbf p}\mathbf f^{\mathbf p}\pmod{\mathfrak R^{d+1}N}
\]
\end{env}
où on pose comme d'ordinaire $\mathbf p=(p_1,\ldots,p_n)$, $\mathbf f^{\mathbf p}=f_1^{p_1}f_2^{p_2}\cdots f_n^{p_n}$, $|\mathbf p|=p_1+\cdots+p_n$ et où $c_{\mathbf p}\in\mathfrak R^dN$. En outre, comme $\psi_N$ identifie $\operatorname{gr}_{\mathfrak L}^*(N)$ à $(\operatorname{gr}_{\mathfrak R}^*(N)\otimes_{A/\mathfrak R}(A/\mathfrak L))[T_1,\ldots,T_n]$, les $c_{\mathbf p}$ sont déterminés mod. $\mathfrak R^d\mathfrak LN$, et si $\bar c_{\mathbf p}$ est la classe de $c_{\mathbf p}$ mod. $\mathfrak R^d\mathfrak LN$, on a, par l'identification précédente,
\begin{env}[19.7.1.22]
\label{IV.19.7.1.22}
\[
  \bar z=\sum_{|\mathbf p|=h}\bar c_{\mathbf p}T^{\mathbf p}.
\]
\end{env}
Utilisant (19.7.1.3), on en déduit que l'on a nécessairement, pour tout $\mathbf p$ tel que $|\mathbf p|=h$ ($\bar c_{\mathbf p}$ étant cette fois identifié par $\psi_N$ à un élément de $\operatorname{gr}_{\mathfrak R}^d(N)$)
\begin{env}[19.7.1.23]
\label{IV.19.7.1.23}
\[
  \bar c_{\mathbf p}\in\operatorname{gr}_{\mathfrak R}^d(R,N).
\]
\end{env}
Mais puisque $d<m$, l'hypothèse que $(*_d)$ est vraie implique que $\bar c_{\mathbf p}$ appartient à l'image de $\psi_R'$, donc qu'on peut supposer que $c_{\mathbf p}\in(R\cap\mathfrak R^dN)+\mathfrak R^d\mathfrak LN$. La relation (19.7.1.21) donne alors
\[
  z\in\mathfrak J^h(R\cap\mathfrak R^dN)+\mathfrak J^h\mathfrak R^d\mathfrak LN+\mathfrak R^{d+1}N
\]
et comme $\mathfrak J\subset\mathfrak L$, cela prouve (19.7.1.19).

\oldpage[IV-4]{218}

En particulier, on peut donc écrire
\[
  g=t+g_2
\]
avec $t\in\mathfrak J^h(R\cap\mathfrak R^dN)$ et $g_2\in\mathfrak R^{d+1}N+\mathfrak R^d\mathfrak L^hN$; en outre, comme $t\in R$ et $g\in R$, on a $g_2\in R$. D'autre part, utilisant l'hypothèse $h\geq m-d+1$, on voit que $t\in\mathfrak L^{m+1}N$. L'élément $g_2$ répond donc à toutes les conditions énoncées au début de (19.7.1.8), ce qui achève la démonstration de (19.7.1).

\begin{remarks}[19.7.2]
\label{IV.19.7.2}
\begin{enumerate}
  \item[(i)] Lorsque l'on suppose que $\operatorname{gr}_{\mathfrak R}^*(A)$ est un $(A/\mathfrak R)$-module \emph{plat}, l'hypothèse que la suite $\mathbf f$ est $(A/\mathfrak R)$-régulière entraîne qu'elle est aussi $\operatorname{gr}_{\mathfrak R}^*(A)$-régulière ($0$, 15.1.14). On notera que ce sera le cas lorsque $A$ et $A/\mathfrak R$ sont des anneaux réguliers ($0$, 17.3.6): en effet, $\mathfrak R$ est alors un idéal régulier de $A$ (19.1.2), donc quasi-régulier, et en localisant aux idéaux maximaux de $A$ contenant $\mathfrak R$ et utilisant ($0$, 15.1.7), on voit que $\operatorname{gr}_{\mathfrak R}^*(A)$ est un $(A/\mathfrak R)$-module \emph{projectif}.
  \item[(ii)] Il serait intéressant de pouvoir prouver la dernière conclusion de (19.7.1) en supposant seulement qu'il existe une $A$-algèbre $B$ qui est un anneau noethérien, pour lequel $\mathfrak JB$ est contenu dans le radical de $B$, et que $M$ est un $B$-module de type fini.
\end{enumerate}
\end{remarks}

\begin{corollary}[19.7.3]
\label{IV.19.7.3}
Soient $A$ un anneau local noethérien, $\mathfrak m$ son idéal maximal, $\mathfrak R$ un idéal de $A$ tel que l'anneau $A/\mathfrak R$ soit régulier et de dimension $n$, $M$ un $A$-module de type fini. Les conditions suivantes sont équivalentes:
\begin{enumerate}
  \item[a)] $\operatorname{gr}_{\mathfrak R}^*(M)$ est un $(A/\mathfrak R)$-module plat.
  \item[b)] Si $P(T)$ et $Q(T)$ sont les séries de Poincaré des $(A/\mathfrak m)$-modules gradués $\operatorname{gr}_{\mathfrak m}^*(M)$ et $\operatorname{gr}_{\mathfrak R}^*(M)\otimes_{A/\mathfrak R}(A/\mathfrak m)$, on a
\end{enumerate}
\begin{equation}
  Q(T)=P(T)(1-T)^n.
  \tag{19.7.3.1}\label{IV.19.7.3.1}
\end{equation}

Soit $\mathbf f=(f_1,\ldots,f_n)$ une suite d'éléments de $A$ tels que les images des $f_i$ dans $A/\mathfrak R$ forment un système régulier de paramètres de $A/\mathfrak R$ ($0$, 17.1.6), de sorte que $\mathbf f$ est une suite $(A/\mathfrak R)$-régulière; en outre, si $\mathfrak J$ est l'idéal engendré par $\mathbf f$, on a $\mathfrak J+\mathfrak R=\mathfrak m$, et comme $A/\mathfrak m$ est un corps, $\operatorname{gr}_{\mathfrak m}^*(M)$ est un $(A/\mathfrak m)$-module plat. Comme $A$ est noethérien et $M$ un $A$-module de type fini, on peut appliquer l'équivalence de $a)$ et $b)$ dans (19.7.1). En outre, comme $A/\mathfrak m$ est un corps, et que $\psi_M$ est surjectif, le fait que $\psi_M$ soit bijectif équivaut à dire que pour tout $h\geq0$, $\operatorname{gr}_{\mathfrak m}^h(M)$ et le sous-module des éléments de degré $h$ de $(\operatorname{gr}_{\mathfrak R}^*(M)\otimes_{A/\mathfrak R}(A/\mathfrak m))[T_1,\ldots,T_n]$ ont même rang sur $A/\mathfrak m$; mais pour le second de ces modules le rang en question est évidemment égal à
\[
  \sum_{r=0}^h\binom{r+n-1}{n-1}\operatorname{rg}(\operatorname{gr}_{\mathfrak R}^{h-r}(M)\otimes_{A/\mathfrak R}(A/\mathfrak m)).
\]
Comme
\[
  (1-T)^{-n}=\sum_{r=0}^{\infty}\binom{r+n-1}{n-1}T^r,
\]
cela prouve que la relation (19.7.3.1) est nécessaire et suffisante pour que $\psi_M$ soit bijectif.
\end{corollary}

\begin{corollary}[19.7.4]
\label{IV.19.7.4}
Soient $X$ un préschéma localement noethérien, $Y$ un sous-préschéma fermé de $X$, régulier et connexe, $\sh F$ un $\sh O_X$-Module cohérent normalement plat (6.10.1) le long de $Y$. Alors il existe une série formelle $R\in\mathbb{Q}[[T]]$ (indépendante de $x\in Y$) telle que, pour tout $x\in Y$, la série de Poincaré de $\operatorname{gr}_{\mathfrak m_x}^*(\sh F_x)$ soit donnée par la formule
\begin{equation}
  P_x(\sh F)(T)=R(T)(1-T)^{-n}\qquad\text{où }n=\dim(\sh O_{Y,x}).
  \tag{19.7.4.1}\label{IV.19.7.4.1}
\end{equation}

\oldpage[IV-4]{219}

Soit $\sh J$ l'Idéal cohérent de $\sh O_X$ définissant $Y$; l'hypothèse sur $\sh F$ signifie que $\operatorname{gr}_{\sh J}^*(\sh F)$ est un $\sh O_Y$-Module plat. On peut donc, pour tout $x\in Y$, appliquer (19.7.3) en prenant $A=\sh O_{X,x}$, $\mathfrak R=\sh J_x$, $M=\sh F_x$; l'hypothèse $a)$ étant remplie par hypothèse, la série de Poincaré $P_x(\sh F)(T)$ est donc égale à $R_x(T)(1-T)^{-n}$, où $R_x(T)$ est la série de Poincaré de $\operatorname{gr}_{\sh J_x}^*(\sh F_x)\otimes_{\sh O_{Y,x}}k(x)$. Or, chacun des $\sh O_Y$-Modules $\sh J^h\sh F/\sh J^{h+1}\sh F$ est plat et cohérent par hypothèse, donc localement libre (2.1.12), et puisque $Y$ est supposé connexe, chacun des $\sh J^h\sh F/\sh J^{h+1}\sh F$ est de rang constant dans $Y$, ce qui prouve que la série de Poincaré $R_x(T)$ est indépendante de $x\in Y$.
\end{corollary}

\begin{corollary}[19.7.5]
\label{IV.19.7.5}
Soient $X$ un préschéma localement noethérien, $Y$ un sous-préschéma fermé de $X$, $Z$ un sous-préschéma fermé de $Y$; on suppose $Y$ et $Z$ réguliers. Soit $\sh F$ un $\sh O_X$-Module cohérent.
\begin{enumerate}
  \item[(i)] Si $\sh F$ est normalement plat le long de $Y$ aux points de $Z$ (cf. (11.3.4) pour la définition de la platitude normale en un point), alors $\sh F$ est normalement plat le long de $Z$.
  \item[(ii)] Supposons en outre que $Z$ soit connexe et que $\sh O_X$ soit normalement plat le long de $Y$ aux points de $Z$ (ce qui aura lieu en particulier si $X$ est régulier aux points de $Z$). Si $\sh F$ est normalement plat le long de $Z$ et s'il existe un point $z\in Z$ tel que $\sh F$ soit normalement plat le long de $Y$ au point $z$, alors $\sh F$ est normalement plat le long de $Y$ aux points de $Z$.
\end{enumerate}

Soient $\sh K$ et $\sh L\supset\sh K$ les Idéaux cohérents de $\sh O_X$ définissant respectivement $Y$ et $Z$. Pour tout point $z\in Z$, l'immersion canonique $Z\to Y$ est régulière au point $z$, puisque $Y$ et $Z$ sont réguliers en ce point (19.1.1); autrement dit, il existe une suite $(f_i)_{1\leq i\leq n}$ d'éléments de $\sh O_{X,z}$ qui est $\sh K_z$-régulière et est telle que $\sh L_z=\sh K_z+\sum_i f_i\sh O_{X,z}$.
\begin{enumerate}
  \item[(i)] L'hypothèse signifie que $\operatorname{gr}_{\sh K_z}^*(\sh F_z)$ est un $\sh O_{Y,z}$-module plat; elle entraîne donc que $\operatorname{gr}_{\sh L_z}^*(\sh F_z)$ est un $\sh O_{Z,z}$-module plat, en vertu du fait que dans (19.7.1), $a)$ entraîne $b)$.
  \item[(ii)] L'hypothèse supplémentaire sur $X$ entraîne que la suite $(f_i)$ est $\operatorname{gr}_{\sh K_z}^*(\sh F_z)$-régulière en tout point $z\in Z$ (19.7.2). D'autre part, $Z$ étant régulier et connexe, est intègre, et si $\operatorname{gr}_{\sh K_z}^*(\sh F_z)$ est $\sh O_{Y,z}$-plat en un point $z\in Z$, alors, comme le point générique $\zeta$ de $Z$ est une générisation de $z$, on en conclut que $\operatorname{gr}_{\sh K_\zeta}^*(\sh F_\zeta)$ est $\sh O_{Y,\zeta}$-plat ($0_{\mathrm I}$, 6.3.1). On en déduit que si en tout point $z'$ de $Z$, $\operatorname{gr}_{\sh L_{z'}}^*(\sh F_{z'})$ est un $\sh O_{Z,z'}$-module plat, alors $\operatorname{gr}_{\sh K_{z'}}^*(\sh F_{z'})$ est un $\sh O_{Y,z'}$-module plat, car on peut appliquer l'équivalence de $a)$ et $d)$ dans (19.7.1).
\end{enumerate}
\end{corollary}

\subsection{Propriétés de passage à la limite projective.}
\EGAReaderTarget{ega4.c04.subsec.19-8}{subsection:IV.19.8}
\label{IV.19.8}

Dans ce numéro, les notations et conventions sur les limites projectives sont celles de (8.5.1) et de (8.8.1).

\begin{proposition}[19.8.1]
\label{IV.19.8.1}
Supposons que les morphismes de transition $S_\mu\to S_\lambda$ $(\lambda\leq\mu)$ soient plats, et en outre que l'une des deux hypothèses suivantes soit vérifiée:
\begin{enumerate}
  \item[$1^\circ$] Les préschémas $S_\lambda$ sont localement noethériens.
  \item[$2^\circ$] Les morphismes de transition $S_\mu\to S_\lambda$ sont surjectifs (donc fidèlement plats).
\end{enumerate}

\oldpage[IV-4]{220}

\emph{Dans ces conditions:}
\begin{enumerate}
  \item[(i)] Supposons $S_\alpha$ quasi-compact. Soit $\sh F_\alpha$ un $\sh O_{S_\alpha}$-Module quasi-cohérent, que l'on suppose en outre de type fini quand l'hypothèse $1^\circ$ est vérifiée. Soit $(f_{i\alpha})_{1\leq i\leq n}$ une suite finie de sections de $\sh F_\alpha$ au-dessus de $S_\alpha$. Pour que la suite des sections $f_i$ du $\sh O_S$-Module $\sh F$ au-dessus de $S$, correspondant aux $f_{i\alpha}$ $(1\leq i\leq n)$, soit $\sh F$-régulière, il faut et il suffit qu'il existe $\lambda\geq\alpha$ tel que la suite des sections $f_{i\lambda}$ de $\sh F_\lambda$ au-dessus de $S_\lambda$ soit $\sh F_\lambda$-régulière.
  \item[(ii)] Supposons $X_\alpha$ quasi-compact, et soit $j_\alpha:X_\alpha\to S_\alpha$ une immersion, que l'on suppose localement de présentation finie lorsque l'hypothèse $2^\circ$ est vérifiée. Alors, pour que l'immersion correspondante $j:X\to S$ soit régulière, il faut et il suffit qu'il existe $\lambda\geq\alpha$ tel que $j_\lambda:X_\lambda\to S_\lambda$ soit régulière.
\end{enumerate}
\end{proposition}

\begin{enumerate}
  \item[(i)] Si $p_\lambda:S\to S_\lambda$ est la projection canonique, on a
  \[
    \sh F/\Bigl(\sum_{j=1}^{i-1}f_j\sh F\Bigr)
    =p_\lambda^*\Bigl(\sh F_\lambda/\Bigl(\sum_{j=1}^{i-1}f_{j\lambda}\sh F_\lambda\Bigr)\Bigr),
  \]
  donc on est ramené au cas $n=1$, auquel cas on supprime l'indice $i$. Le fait que la condition soit suffisante résulte de ce que $p_\lambda$ est plat (8.3.8) et de ($0$, 15.1.5). Dans le cas $2^\circ$, $p_\lambda$ est fidèlement plat (8.3.8) et la nécessité de la condition résulte encore de ($0$, 15.2.5), avec $\lambda=\alpha$. Dans le cas $1^\circ$, notons $\sh N_\lambda$ (resp. $\sh N$) le noyau de l'homomorphisme $\sh F_\lambda\to\sh F_\lambda$ (resp. $\sh F\to\sh F$), multiplication par $f_\lambda$ (resp. par $f$); puisque $\sh F_\lambda$ est cohérent par hypothèse, il en est de même de $\sh N_\lambda$, donc l'hypothèse $\sh N=0$ entraîne $\sh N_\lambda=0$ pour un $\lambda\geq\alpha$ en vertu de (8.5.8, (ii)).
  \item[(ii)] Comme $p_\lambda$ est plat, la suffisance de la condition résulte de (19.1.5, (ii)). Pour prouver sa nécessité, notons que, puisque $j_\alpha(X_\alpha)$ est quasi-compact, il est contenu dans un ouvert quasi-compact de $S_\alpha$ et on peut donc se limiter au cas où $S_\alpha$ est aussi quasi-compact et $j_\alpha$ une immersion fermée, de sorte que l'image de $X_\alpha$ est définie par un Idéal quasi-cohérent $\sh J_\alpha$ de $\sh O_{S_\alpha}$, qu'on peut en outre supposer de type fini dans les cas $1^\circ$ et $2^\circ$ ($j_\alpha$ étant localement de présentation finie dans le cas $2^\circ$); l'image de $X_\lambda$ (resp. $X$) dans $S_\lambda$ (resp. $S$) est alors définie par $\sh J_\lambda=\sh J_\alpha\sh O_{S_\lambda}$ (resp. $\sh J=\sh J_\alpha\sh O_S$) qui est encore de type fini. On peut en outre, compte tenu de (8.2.11), supposer que $\sh J$ est engendré par une suite régulière $(f_i)_{1\leq i\leq n}$ de sections au-dessus de $S$, qui définissent par suite un homomorphisme surjectif $u:\sh O_S^n\to\sh J$. Compte tenu de (8.5.2, (i)) et (8.5.7), il existe un $\lambda\geq\alpha$ et un homomorphisme surjectif $u_\lambda:\sh O_{S_\lambda}^n\to\sh J_\lambda$ tels que $u=p_\lambda^*(u_\lambda)$, donc les $f_i$ sont les images canoniques de sections $f_{i\lambda}$ de $\sh J_\lambda$ au-dessus de $S_\lambda$ engendrant cet Idéal. En vertu de (i), il existe donc $\mu\geq\lambda$ tel que la suite $(f_{i\mu})$ soit $\sh O_{S_\mu}$-régulière, donc l'immersion $j_\mu$ est régulière.
\end{enumerate}

\begin{proposition}[19.8.2]
\label{IV.19.8.2}
Soient $X_\alpha$ un $S_\alpha$-préschéma quasi-compact et localement de présentation finie.
\begin{enumerate}
  \item[(i)] Soient $\sh F_\alpha$ un $\sh O_{X_\alpha}$-Module quasi-cohérent de présentation finie, $(f_{i\alpha})_{1\leq i\leq n}$ une suite de sections de $\sh O_{X_\alpha}$ au-dessus de $X_\alpha$. Pour que la suite de sections $f_i$ de $\sh O_X$ au-dessus de $X$, correspondant aux $f_{i\alpha}$, soit $\sh F$-transversalement régulière relativement à $S$ (19.2.1), il faut et il suffit qu'il existe $\lambda\geq\alpha$ tel que la suite des sections $f_{i\lambda}$ de $\sh O_{X_\lambda}$ au-dessus de $X_\lambda$ soit $\sh F_\lambda$-transversalement régulière relativement à $S_\lambda$.
  \item[(ii)] Soient $Y_\alpha$ un $S_\alpha$-préschéma quasi-compact, $j_\alpha:Y_\alpha\to X_\alpha$ une $S_\alpha$-immersion localement

\oldpage[IV-4]{221}

  de présentation finie. Pour que la $S$-immersion correspondante $j:Y\to X$ soit transversalement régulière relativement à $S$, il faut et il suffit qu'il existe $\lambda\geq\alpha$ tel que $j_\lambda:Y_\lambda\to X_\lambda$ soit transversalement régulière relativement à $S_\lambda$.
\end{enumerate}
\end{proposition}

\begin{enumerate}
  \item[(i)] La suffisance de la condition résulte de ($0$, 15.1.15). Pour prouver qu'elle est nécessaire, il suffira de montrer que tout point $x\in X$ possède un voisinage ouvert $V(x)$ pour lequel il existe un indice $\lambda=\lambda(x)\geq\alpha$ et un voisinage ouvert $V_\lambda$ de l'image $x_\lambda$ de $x$ dans $X_\lambda$ tels que $V(x)$ soit l'image réciproque de $V_\lambda$ et que la suite des restrictions à $V_\lambda$ des $f_{i\lambda}$ soit $(\sh F_\lambda|V_\lambda)$-transversalement régulière relativement à $S_\lambda$: en effet, il suffira alors d'appliquer (8.3.4) à l'ensemble ouvert $U_\lambda$ de $X_\lambda$, défini par la condition d'être le plus grand ouvert tel que les restrictions à $U_\lambda$ des $f_{i\lambda}$ forment une suite $(\sh F_\lambda|U_\lambda)$-transversalement régulière relativement à $S_\lambda$.

  Soient $s$ l'image de $x$ dans $S$, et pour tout $\lambda\geq\alpha$, soit $s_\lambda$ l'image de $s$ dans $S_\lambda$. Alors la fibre $X_s$ est limite projective des fibres $(X_\lambda)_{s_\lambda}$ et est localement noethérienne puisque le morphisme $X\to S$ est localement de présentation finie; comme en outre $(X_\alpha)_{s_\alpha}$ est quasi-compact et que les morphismes de transition $(X_\mu)_{s_\mu}\to(X_\lambda)_{s_\lambda}$ sont plats (puisqu'il en est ainsi de $\Spec(k(s_\mu))\to\Spec(k(s_\lambda))$) on peut appliquer aux sections $f_i\otimes1$ de $\sh F_s=\sh F\otimes_{\sh O_S}k(s)$ au-dessus de $X_s$ le résultat de (19.8.1, (i)); d'où l'existence d'un $\lambda\geq\alpha$ tel que les sections $f_{i\lambda}\otimes1$ de $(\sh F_\lambda)_{s_\lambda}=\sh F_\lambda\otimes_{\sh O_{S_\lambda}}k(s_\lambda)$ au-dessus de $(X_\lambda)_{s_\lambda}$ forment une suite $(\sh F_\lambda)_{s_\lambda}$-régulière. En outre (11.2.6), on peut supposer que $\sh F_\lambda$ est $S_\lambda$-plat au point $x_\lambda$. Il résulte alors de (11.3.8) qu'il existe un voisinage ouvert $V_\lambda$ de $x_\lambda$ dans $X_\lambda$ vérifiant les conditions requises.

  \item[(ii)] La suffisance de la condition résulte de (19.2.7, (ii)). Pour prouver sa nécessité, il suffit encore de montrer que tout point $x\in j(Y)$ admet un voisinage ouvert $V(x)$ pour lequel il existe un indice $\lambda=\lambda(x)\geq\alpha$ et un voisinage ouvert $V_\lambda$ de l'image $x_\lambda$ de $x$ dans $X_\lambda$ tels que $V(x)$ soit l'image réciproque de $V_\lambda$ et que la restriction $j_\lambda^{-1}(V_\lambda)\to V_\lambda$ de $j_\lambda$ soit transversalement régulière relativement à $S_\lambda$. Soient $s$ l'image de $x$ dans $S$, et pour tout $\lambda\geq\alpha$, soit $s_\lambda$ l'image de $s$ dans $S_\lambda$. Comme $X_s$ (resp. $Y_s$) est limite projective des fibres $(X_\lambda)_{s_\lambda}$ (resp. $(Y_\lambda)_{s_\lambda}$), on peut comme dans (i) utiliser (19.8.1, (ii)), et comme par hypothèse l'immersion $Y_s\to X_s$ est régulière, il existe un indice $\lambda(s)$ tel que pour $\lambda\geq\lambda(s)$ l'immersion $(Y_\lambda)_{s_\lambda}\to(X_\lambda)_{s_\lambda}$ soit régulière. En outre, en vertu de (19.2.4) il existe un voisinage ouvert quasi-compact $W(x)$ de $x$ dans $X$ tel que les morphismes structuraux $W(x)\to S$ et $Y\cap W(x)\to S$ soient plats; appliquant (11.2.6) et (8.2.11), on peut supposer qu'il existe un $\lambda\geq\lambda(s)$ et un voisinage ouvert quasi-compact $W(x_\lambda)$ de $x_\lambda$ dans $X_\lambda$, tels que $W(x)$ soit l'image réciproque de $W(x_\lambda)$ et que les restrictions à $W(x_\lambda)$ et à $Y_\lambda\cap W(x_\lambda)$ respectivement des morphismes structuraux $X_\lambda\to S_\lambda$ et $Y_\lambda\to S_\lambda$ soient plats. On conclut donc de (19.2.4) qu'il existe un voisinage ouvert $V_\lambda$ de $x_\lambda$ répondant à la question.
\end{enumerate}

\begin{remark}[19.8.3]
\label{IV.19.8.3}
Les démonstrations précédentes montrent que les énoncés de (19.8.1) et (19.8.2) subsistent lorsqu'on remplace les conditions concernant la régularité ou la régularité transversale dans tout $X$ (resp. $X_\lambda$) par ces conditions dans un voisinage d'un point $x\in X$ (resp. dans un voisinage de $x_\lambda$, projection de $x$).
\end{remark}

\subsection{\texorpdfstring{Suites $\sh F$-régulières et profondeur.}{Suites F-régulières et profondeur.}}
\EGAReaderTarget{ega4.c04.subsec.19-9}{subsection:IV.19.9}
\label{IV.19.9}

\oldpage[IV-4]{222}

\begin{env}[19.9.1]
\EGAReaderTarget{ega4.c04.statement.19-9-1}{IV.19.9.1}
Rappelons que si $X$ est un préschéma localement noethérien, $\sh F$ un $\sh O_X$-Module cohérent, $T$ une partie de $X$, on pose (5.10.1)
\begin{equation}
  \operatorname{prof}_T(\sh F)=\inf_{x\in T}\operatorname{prof}(\sh F_x).
  \tag{19.9.1.1}\label{IV.19.9.1.1}
\end{equation}

On pose d'autre part, pour tout point $t\in T$
\begin{equation}
  \operatorname{prof}_{T,t}(\sh F)=\inf_{x\in T\cap\Spec(\sh O_{X,t})}\operatorname{prof}(\sh F_x),
  \tag{19.9.1.2}\label{IV.19.9.1.2}
\end{equation}
et on dit que $\operatorname{prof}_{T,t}(\sh F)$ est la \emph{$T$-profondeur de $\sh F$ au point $t$}; il est clair que l'on a
\begin{equation}
  \operatorname{prof}_T(\sh F)=\inf_{t\in T}\operatorname{prof}_{T,t}(\sh F).
  \tag{19.9.1.3}\label{IV.19.9.1.3}
\end{equation}
\end{env}

\begin{lemma}[19.9.2]
\label{IV.19.9.2}
Soient $A$ un anneau noethérien, $M$ un $A$-module de type fini, $\mathfrak J$ un idéal de $A$. Afin que, pour tout idéal premier $\mathfrak p\supset\mathfrak J$, on ait $\operatorname{prof}_{A_{\mathfrak p}}(M_{\mathfrak p})\geq r$, il faut et il suffit qu'il existe une suite $M$-régulière de $r$ éléments appartenant à $\mathfrak J$.
\end{lemma}

La suffisance de la condition résultant aussitôt de ($0$, 16.4.5), prouvons sa nécessité. Raisonnons par récurrence sur $r$ (il n'y a rien à démontrer si $r=0$): puisque $\operatorname{prof}_{A_{\mathfrak p}}(M_{\mathfrak p})\geq r-1$ pour tout $\mathfrak p\supset\mathfrak J$, il existe par hypothèse une suite $M$-régulière $(g_i)_{1\leq i\leq r-1}$ formée d'éléments de $\mathfrak J$. Posons $N=M/(g_1M+\cdots+g_{r-1}M)$; pour tout $\mathfrak p\supset\mathfrak J$, les images des $g_i$ dans $A_{\mathfrak p}$ appartiennent donc à l'idéal maximal $\mathfrak pA_{\mathfrak p}$, et forment une suite $M_{\mathfrak p}$-régulière ($0$, 15.1.14); on conclut donc de ($0$, 16.4.6) que l'on a $\operatorname{prof}_{A_{\mathfrak p}}(N_{\mathfrak p})=\operatorname{prof}_{A_{\mathfrak p}}(M_{\mathfrak p})-(r-1)\geq1$, et on voit que tout revient à prouver le lemme pour $r=1$. Or l'hypothèse signifie alors que, pour tout $\mathfrak p\supset\mathfrak J$, $\mathfrak pA_{\mathfrak p}$ n'est pas associé à $M_{\mathfrak p}$ ($0$, 16.4.6), donc (Bourbaki, \emph{Alg. comm.}, chap. IV, § 1, n\textsuperscript{o} 2, prop. 5), $\mathfrak p$ n'est pas associé à $M$. Autrement dit, $\mathfrak J$ n'est contenu dans aucun des idéaux premiers de $\operatorname{Ass}(M)$, donc il n'est pas contenu dans leur réunion (Bourbaki, \emph{Alg. comm.}, chap. II, § 1, n\textsuperscript{o} 1, prop. 2). Mais comme cette réunion est l'ensemble des éléments qui ne sont pas $M$-réguliers (Bourbaki, \emph{Alg. comm.}, chap. IV, § 1, n\textsuperscript{o} 1, cor. 2 de la prop. 2), cela prouve le lemme.

\begin{proposition}[19.9.3]
\label{IV.19.9.3}
Soient $X$ un préschéma localement noethérien, $\sh F$ un $\sh O_X$-Module cohérent, $Y$ une partie fermée de $X$, $y$ un point de $Y$, $n$ un entier $>0$. Les conditions suivantes sont équivalentes:
\begin{enumerate}
  \item[a)] Il existe un voisinage ouvert $U$ de $y$ dans $X$ tel que $\operatorname{prof}_{Y\cap U}(\sh F|U)\geq n$.
  \item[b)] Il existe un voisinage ouvert $U$ de $y$ dans $X$ et une suite $(\sh F|U)$-régulière de $n$ sections $f_i$ de $\sh O_U$ au-dessus de $U$, telles que $f_i(x)=0$ ($0_{\mathrm I}$, 5.5.1) pour $1\leq i\leq n$ en tout point $x\in U\cap Y$.
  \item[c)] On a $\operatorname{prof}_{Y,y}(\sh F)\geq n$.
\end{enumerate}
En outre, l'ensemble des $y\in Y$ vérifiant ces conditions est ouvert dans $Y$.
\end{proposition}

La dernière assertion résulte trivialement de $a)$. L'équivalence de $a)$ et $b)$ découle de (19.9.2) appliqué à un voisinage ouvert affine de $y$ dans $X$ et de ($0$, 15.1.14). Il est clair que $a)$ entraîne $c)$, tout voisinage ouvert de $y$ dans $X$ contenant $\Spec(\sh O_{X,y})$; montrons inversement que $c)$ implique $b)$. Il résulte de $c)$ et de la définition (19.9.1.1) que l'on

\oldpage[IV-4]{223}

peut appliquer (19.9.2) à l'anneau $\sh O_{X,y}$, au $\sh O_{X,y}$-module $\sh F_y$ et à l'idéal $\sh J_y$ de $\sh O_{X,y}$, donc il existe une suite $\sh F_y$-régulière $(t_i)_{1\leq i\leq n}$ de $n$ éléments de $\sh J_y$. Il y a donc un voisinage ouvert $U$ de $y$ dans $X$ et une suite $(f_i)_{1\leq i\leq n}$ de sections de $\sh J$ au-dessus de $U$ telles que les $t_i$ soient les germes des $f_i$ au point $y$; utilisant ($0$, 15.2.4), on en déduit qu'il existe un voisinage ouvert $U'\subset U$ de $y$ dans $X$ tel que les $f_i|U'$ forment une suite $(\sh F|U')$-régulière.

\begin{corollary}[19.9.4]
\label{IV.19.9.4}
Avec les notations de (19.9.3), la fonction $y\mapsto\operatorname{prof}_{Y,y}(\sh F)$ est semi-continue inférieurement dans $Y$.
\end{corollary}

\begin{proposition}[19.9.5]
\label{IV.19.9.5}
Avec les notations de (19.9.3), soient $X'$ un préschéma localement noethérien, $g:X'\to X$ un morphisme plat, $Y'=g^{-1}(Y)$, $\sh F'=g^*(\sh F)$. Alors, pour tout point $y'\in Y'$ au-dessus de $y$, on a $\operatorname{prof}_{Y',y'}(\sh F')=\operatorname{prof}_{Y,y}(\sh F)$.
\end{proposition}

En effet, pour tout $z'\in\Spec(\sh O_{X',y'})$, il résulte de (6.3.1) que l'on a $\operatorname{prof}(\sh F'_{z'})\geq\operatorname{prof}(\sh F_{g(z')})$, et si $z''$ est un point maximal de la fibre $g^{-1}(g(z'))$, générisation de $z'$ (donc appartenant aussi à $\Spec(\sh O_{X',y'})$), on a (\emph{loc. cit.}) $\operatorname{prof}(\sh F'_{z''})=\operatorname{prof}(\sh F_{g(z')})$. La proposition résulte alors de ce que l'on a $g(\Spec(\sh O_{X',y'}))=\Spec(\sh O_{X,g(y')})$ puisque $g$ est plat (2.3.4).

\begin{proposition}[19.9.6]
\label{IV.19.9.6}
Soient $f:X\to S$ un morphisme localement de présentation finie, $Y$ une partie fermée de $X$, $\sh F$ un $\sh O_X$-Module de présentation finie et $f$-plat, $n$ un entier $>0$. Pour tout point $y\in Y$, les conditions suivantes sont équivalentes (en désignant par $X_{f(y)}$ la fibre de $f$ au point $f(y)$, par $\sh F_{f(y)}$ l'image réciproque de $\sh F$ dans $X_{f(y)}$, et en posant $Y_{f(y)}=X_{f(y)}\cap Y$):
\begin{enumerate}
  \item[a)] On a $\operatorname{prof}_{Y_{f(y)},y}(\sh F_{f(y)})\geq n$.
  \item[b)] Il existe un voisinage ouvert $U$ de $y$ dans $X$ et une suite $(g_i)_{1\leq i\leq n}$ de $n$ sections de $\sh O_X$ au-dessus de $U$, qui soit transversalement $\sh F$-régulière relativement à $S$ (cf. (19.2.1) pour la terminologie) et telle que $g_i(x)=0$ pour $1\leq i\leq n$ et pour tout $x\in U\cap Y$.
\end{enumerate}
L'ensemble $V_n$ des $y\in Y$ vérifiant ces conditions est ouvert dans $Y$; si de plus $Y$ est localement constructible dans $X$, $V_n$ est rétrocompact dans $X$ ($0_{\mathrm{III}}$, 9.1.1).
\end{proposition}

En vertu de (19.9.3), la condition $a)$ équivaut à l'existence d'un voisinage ouvert $U$ de $y$ dans $X$ et d'une suite $(\sh F_{f(y)}|(U\cap X_{f(y)}))$-régulière $(h_i)_{1\leq i\leq n}$ de sections de $\sh O_{X_{f(y)}}$ au-dessus de $U\cap X_{f(y)}$ telles que $h_i(z)=0$ pour tout $z\in U\cap Y_{f(y)}$. Si on désigne encore par $Y$ un sous-préschéma fermé de $X$ ayant $Y$ pour espace sous-jacent, par $\sh J$ l'Idéal quasi-cohérent de $\sh O_X$ définissant $Y$, on peut encore dire que l'on a $h_i\in\sh J_{f(y)}$ pour $1\leq i\leq n$. Comme on peut remplacer $U$ par un voisinage affine plus petit de $y$, on peut supposer que les $h_i$ sont les images d'une suite $(g_i)_{1\leq i\leq n}$ de sections de $\sh J$ au-dessus de $U$, de sorte que les germes $(g_i)_y$ appartiennent à l'idéal maximal $\mathfrak m_y$ de $\sh O_{X,y}$. L'équivalence de $a)$ et $b)$ résulte alors de (11.3.8), qui prouve aussi que l'ensemble $V_n$ est ouvert dans $Y$.

Reste à prouver la dernière assertion; comme il s'agit encore de propriétés locales sur $X$, on peut supposer $S=\Spec(A)$ affine et $X$ de présentation finie sur $S$; il existe alors un sous-anneau noethérien $A_0$ de $A$, un préschéma $X_0$ de type fini sur $S_0=\Spec(A_0)$ et un $\sh O_{X_0}$-Module cohérent $\sh F_0$ tels que $X=X_0\times_{S_0}S$ et $\sh F=\sh F_0\otimes_{\sh O_{X_0}}\sh O_X$ (8.9.1); on peut en outre supposer que, si $f_0:X_0\to S_0$ est le morphisme structural, $\sh F_0$ est $f_0$-plat (11.2.6). Puisque $Y$ est constructible, on peut en outre (8.3.11) supposer que $Y=p^{-1}(Y_0)$, où $p:X\to X_0$ est la projection canonique. Alors, pour tout $y\in Y$, il résulte de (19.9.5) et de la transitivité des fibres que l'on a
\[
  \operatorname{prof}_{Y_{f(y)},y}(\sh F_{f(y)})
  =\operatorname{prof}_{(Y_0)_{f_0(y_0)},y_0}((\sh F_0)_{f_0(y_0)}),
\]
en posant $y_0=p(y)$; si $V_n^0$ est l'ensemble des $y_0\in Y_0$ pour lesquels le second membre de cette relation est $\geq n$, on a donc $V_n=p^{-1}(V_n^0)$. Comme $X_0$ est noethérien, $V_n^0$ est un ouvert rétrocompact dans $X_0$, et par suite (cf. démonstration de (1.8.2)) $V_n$ est rétrocompact dans $X$.

\oldpage[IV-4]{224}

\begin{corollary}[19.9.7]
\label{IV.19.9.7}
Sous les hypothèses générales de (19.9.6), la fonction $y\mapsto\operatorname{prof}_{Y_{f(y)},y}(\sh F_{f(y)})$ est semi-continue inférieurement dans $Y$. Si de plus $Y$ est localement constructible, cette fonction est localement constructible.
\end{corollary}

\begin{proposition}[19.9.8]
\label{IV.19.9.8}
Soient $f:X\to S$ un morphisme localement de présentation finie, $Z$ une partie fermée de $X$, $\sh F$ un $\sh O_X$-Module de présentation finie et $f$-plat, $\sh G$ un $\sh O_S$-Module quasi-cohérent. Supposons que, pour tout $x\in Z$, on ait $\operatorname{prof}((\sh F_{f(x)})_x)\geq1$ (resp. $\operatorname{prof}((\sh F_{f(x)})_x)\geq2$). Alors, si $j:X-Z\to X$ est l'injection canonique, et si on pose $\sh H=\sh F\otimes_{\sh O_X}f^*(\sh G)$, l'homomorphisme canonique $\rho_{\sh H}:\sh H\to j_*(j^*(\sh H))$ relatif à $j$ est injectif (resp. bijectif).
\end{proposition}

La question est évidemment locale sur $X$ et $S$, et il suffit de démontrer la proposition dans un voisinage d'un point $x\in Z$. Autrement dit, on peut se borner au cas où $X$ et $S$ sont affines. Grâce à (19.9.6), on peut supposer qu'il existe une section transversalement $\sh F$-régulière $g_1$ de $\sh O_X$ au-dessus de $X$ relativement à $S$ (resp. une suite transversalement $\sh F$-régulière $(g_1,g_2)$ de deux sections de $\sh O_X$ au-dessus de $X$ relativement à $S$) telle que si $Z'$ est l'ensemble des $x'\in X$ pour lesquels $g_1(x')=0$ (resp. $g_1(x')=g_2(x')=0$), on ait $Z\subset Z'$.

Cela étant, revenons à la démonstration de (19.9.8); en vertu de (19.9.6) on a encore $\operatorname{prof}((\sh F_{f(x)})_x)\geq1$ (resp. $\operatorname{prof}((\sh F_{f(x)})_x)\geq2$) pour tout $x\in Z'$. Si la proposition est prouvée en remplaçant le couple $(X,Z)$ par $(X,Z')$ et $(X-Z,Z'\cap(X-Z))$, il est immédiat qu'elle le sera aussi pour $(X,Z)$; on peut donc se borner à la prouver pour $X$ et $Z'$. Autrement dit, on peut se borner à démontrer (19.9.8) lorsque $S=\Spec(A)$, $X=\Spec(B)$, où $B$ est une $A$-algèbre de présentation finie, $Z=V(\mathfrak J)$, où $\mathfrak J$ est un idéal de type fini de l'anneau $B$, $\sh F=\widetilde M$ et $\sh G=\widetilde N$, où $N$ est un $A$-module et $M$ un $B$-module de présentation finie qui est un $A$-module plat. On peut en outre se ramener au cas où $N$ est un $A$-module de présentation finie. En effet, $N$ est limite inductive d'un système inductif filtrant de $A$-modules de présentation finie $N_\alpha$ (par double limite inductive) si l'on pose $\sh G_\alpha=\widetilde{N_\alpha}$, $\sh H$ est limite inductive des $\sh H_\alpha=\sh F\otimes_{\sh O_X}f^*(\sh G_\alpha)$; comme $j_*$ et $j^*$ commutent aux limites inductives et que $\varinjlim$ est un foncteur exact dans la catégorie des Modules quasi-cohérents, si la proposition est prouvée pour chacun des $\sh H_\alpha$, elle le sera pour $\sh H$.

Il suffit en outre de prouver que l'homomorphisme canonique
\[
  \Gamma(X,\sh H)\longrightarrow\Gamma(X-Z,\sh H)
\]
est injectif (resp. bijectif). Il existe alors un sous-anneau noethérien $A_0$ de $A$, une $A_0$-algèbre de type fini $B_0$, un idéal $\mathfrak J_0$ de $B_0$ et un $B_0$-module de type fini $M_0$ qui est

\oldpage[IV-4]{225}

un $A_0$-module plat et un $A_0$-module $N_0$, tels que $B=B_0\otimes_{A_0}A$, $\mathfrak J=\mathfrak J_0B$, $M=M_0\otimes_{A_0}A$ et $N=N_0\otimes_{A_0}A$ (8.9.1, 8.5.11 et 11.2.7). De plus, soit $(A_\lambda)$ la famille des sous-$A_0$-algèbres de type fini de $A$, de sorte que $A=\varinjlim A_\lambda$; posons
\[
\begin{gathered}
  B_\lambda=B_0\otimes_{A_0}A_\lambda,\qquad
  \mathfrak J_\lambda=\mathfrak J_0B_\lambda,\qquad
  M_\lambda=M_0\otimes_{A_0}A_\lambda=M_0\otimes_{B_0}B_\lambda,\\
  N_\lambda=N_0\otimes_{A_0}A_\lambda,\qquad
  S_\lambda=\Spec(A_\lambda),\quad
  X_\lambda=\Spec(B_\lambda),\quad
  Z_\lambda=V(\mathfrak J_\lambda);
\end{gathered}
\]
si $p_\lambda:X\to X_\lambda$ est la projection canonique, on a donc $Z=p_\lambda^{-1}(Z_\lambda)$, et $X-Z=p_\lambda^{-1}(X_\lambda-Z_\lambda)$. Comme $X_\lambda-Z_\lambda$ est quasi-compact et quasi-séparé, on en conclut (8.5.2) que si l'on pose $\sh F_\lambda=\widetilde{M_\lambda}$ et $\sh H_\lambda=(M_\lambda\otimes_{A_\lambda}N_\lambda)^\sim$, les homomorphismes
\[
  \varinjlim\Gamma(X_\lambda,\sh H_\lambda)\longrightarrow\Gamma(X,\sh H)
  \qquad\text{et}\qquad
  \varinjlim\Gamma(X_\lambda-Z_\lambda,\sh H_\lambda)\longrightarrow\Gamma(X-Z,\sh H)
\]
sont bijectifs; en vertu de l'exactitude du foncteur $\varinjlim$, il suffira donc de prouver que pour tout $\lambda$ assez grand, l'homomorphisme canonique $\Gamma(X_\lambda,\sh H_\lambda)\to\Gamma(X_\lambda-Z_\lambda,\sh H_\lambda)$ est injectif (resp. bijectif). Mais pour tout $x\in Z$, si $x_\lambda$ est la projection de $x$ dans $Z_\lambda$, on a
\[
  \operatorname{prof}((\sh F_{f(x)})_x)
  =\operatorname{prof}(((\sh F_\lambda)_{f_\lambda(x_\lambda)})_{x_\lambda})
\]
en vertu de (4.2.7) et (6.7.1); pour être ramené à démontrer (19.9.8) dans le cas où $A$ est noethérien, il faut prouver que, pour $\lambda$ assez grand on a $\operatorname{prof}(((\sh F_\lambda)_{f_\lambda(x_\lambda)})_{x_\lambda})\geq1$ (resp. $\geq2$) pour tout point $x_\lambda\in Z_\lambda$. Or, notons $Z_\lambda'$ l'ensemble des points $x_\lambda\in Z_\lambda$ tels que, pour toute générisation $z_\lambda$ de $x_\lambda$ dans $Z_\lambda\cap(X_\lambda)_{f_\lambda(x_\lambda)}$, on ait $\operatorname{prof}(((\sh F_\lambda)_{f_\lambda(x_\lambda)})_{z_\lambda})\geq1$; si $p_{\lambda\mu}:Z_\mu\to Z_\lambda$ est la projection canonique pour $\lambda\leq\mu$, il résulte de l'hypothèse et de (19.9.5) que l'on a $Z_\mu'=p_{\lambda\mu}^{-1}(Z_\lambda')$ et $Z=p_\lambda^{-1}(Z_\lambda')$; d'autre part, il résulte de (19.9.6) que $Z_\lambda'$ est ouvert dans $Z_\lambda$, donc notre assertion résulte de (8.3.4).

Pour terminer la démonstration lorsque $A$ est noethérien, notons qu'en vertu de l'hypothèse de platitude et de (6.3.1), on a alors $\operatorname{prof}_Z(\sh H)\geq1$ (resp. $\operatorname{prof}_Z(\sh H)\geq2$); mais la conclusion de (19.9.8) résulte dans ce cas de (5.10.2) (resp. 5.10.4)).

\begin{remark}[19.9.9]
\label{IV.19.9.9}
La même démonstration, jointe aux résultats du chap. III, 3\textsuperscript{e} partie sur la profondeur et la cohomologie locale, permet d'établir la généralisation suivante de (19.9.8):

\emph{Sous les conditions générales de (19.9.8), supposons que, pour tout $x\in Z$, on ait $\operatorname{prof}((\sh F_{f(x)})_x)\geq k$; alors l'homomorphisme canonique}
\[
  H^i(X,\sh H)\longrightarrow H^i(X-Z,\sh H|(X-Z))
\]
\emph{est bijectif pour $i\leq k-2$ et injectif pour $i=k-1$} (ce qui, exprimé à l'aide de la notion cohomologique générale de profondeur introduite au chap. III, s'écrit aussi $\operatorname{prof}_Z(\sh H)\geq k$).
\end{remark}

% FIN SOURCE ega4/ega4-18-fr.tex

\newcommand{\shHom}{\mathcal{H}\!om}
%
% DEBUT SOURCE ega4/ega4-20-fr.tex
\section{Fonctions méromorphes et pseudo-morphismes}
\label{IV.20}

\setcounter{subsection}{-1}
\subsection{Introduction}
\EGAReaderTarget{ega4.c04.subsec.20-0}{IV.20.0}

La plupart des notions et résultats des \textsection\textsection20 et 21 se rattachent directement au chap.~I, et ne dépendent guère des chap.~II à IV, sauf en ce qui concerne l'usage occasionnel de la notion de profondeur et d'anneau local régulier (dans \sref{IV.20.6}, \sref{IV.21.11}, \sref{IV.21.13} et \sref{IV.21.15}), du « Main theorem » de Zariski dans \sref{IV.20.4} et \sref{IV.21.12}, et des propriétés des immersions transversalement régulières dans \sref{IV.20.6} et \sref{IV.21.15}.

\oldpage[IV-4]{226}
Au \textsection20, nous introduisons plusieurs variantes de la notion d'application rationnelle, déjà étudiée dans \sref[I]{section:I.7} d'un point de vue encore assez proche du point de vue classique, et pour cette raison assez mal adaptée au cas des préschémas non nécessairement réduits.
Les notions et résultats du \textsection20 sont utilisés pour développer au \textsection21 (n\textsuperscript{os} \sref{IV.21.1} à \sref{IV.21.7}) la notion générale de diviseur et ses propriétés les plus élémentaires.
Cette notion est surtout commode lorsque les anneaux locaux des préschémas considérés sont noethériens et intégralement clos, et surtout lorsqu'ils sont en outre \emph{factoriels} (\sref{IV.21.6} et \sref{IV.21.7}), en raison de son identification dans ce dernier cas avec la notion de \emph{cycle $1$-codimensionnel} (combinaison linéaire de sous-préschémas irréductibles de codimension~$1$).
Dans \sref{IV.21.9}, on détermine les diviseurs sur un préschéma noethérien de dimension~$1$ mais non nécessairement normal, ce qui est utile pour diverses applications.
Les n\textsuperscript{os} \sref{IV.21.11} et \sref{IV.21.12} donnent deux théorèmes importants, dus respectivement à Auslander--Buchsbaum et Van der Waerden, et se rattachant à la notion d'anneau factoriel (les n\textsuperscript{os} \sref{IV.21.9}, \sref{IV.21.11} et \sref{IV.21.12} sont indépendants les uns des autres).
Dans les n\textsuperscript{os} \sref{IV.21.13} et \sref{IV.21.14}, également indépendants des trois précédents, nous étudions une variante utile de la notion d'anneau local factoriel, celle d'anneau local \emph{parafactoriel}, qui s'introduit notamment~\cite{IV-41} dans le développement de théorèmes de comparaison du groupe de Picard d'un préschéma projectif $X$ sur un corps $k$ et d'une « section hyperplane ».
On verra dans \sref{IV.21.14.1} (théorème de Ramanujam--Samuel) que les anneaux locaux parafactoriels sont bien plus nombreux que l'on ne pouvait s'y attendre \emph{a priori}.

Dans \sref{IV.20.5}, \sref{IV.20.6} et \sref{IV.21.15}, nous reprenons les notions précédentes mais d'un point de vue « relatif » à un préschéma de base fixé.
Pour l'instant ces notions ne sont utilisées qu'assez rarement ; en particulier la notion de diviseur relatif n'est guère employée que lorsqu'il s'agit de diviseurs positifs, et dans ce cas elle s'explicite avantageusement sans le secours de la notion de fonction méromorphe relative, à l'aide de la notion d'immersion transversalement régulière de codimension~$1$.
On aura donc avantage à omettre ces sections en première lecture.

\subsection{Fonctions méromorphes}
\EGAReaderTarget{ega4.c04.subsec.20-1}{IV.20.1}

\begin{env}[20.1.1]
\EGAReaderTarget{ega4.c04.statement.20-1-1}{IV.20.1.1}
Soit $(X,\sh{O}_X)$ un espace annelé, et soit $\sh{S}$ un sous-faisceau \emph{d'ensembles} de $\sh{O}_X$.
Pour tout ouvert $U$ de $X$, considérons l'\emph{anneau de fractions} $\Gamma(U,\sh{O}_X)[\Gamma(U,\sh{S})^{-1}]$ (Bourbaki, \emph{Alg. comm.}, chap.~II, \textsection2, n\textsuperscript{o}1).
Il est immédiat que l'application $U\mapsto\Gamma(U,\sh{O}_X)[\Gamma(U,\sh{S})^{-1}]$ est un \emph{préfaisceau d'anneaux} (\sref[0]{0.1.5.1} et \sref[0]{0.1.5.7}).
On note $\sh{O}_X[\sh{S}^{-1}]$ le \emph{faisceau d'anneaux} associé à ce préfaisceau et on dit que c'est le \emph{faisceau d'anneaux de fractions de $\sh{O}_X$, à dénominateurs dans $\sh{S}$} ; c'est un $\sh{O}_X$-module \emph{plat}.
Il est immédiat que pour tout $x\in X$, on a un isomorphisme canonique
\[
\label{IV.20.1.1.1}
  (\sh{O}_X[\sh{S}^{-1}])_x\isoto\sh{O}_{X,x}[\sh{S}_x^{-1}],
  \tag{20.1.1.1}
\]
car le raisonnement de \sref[0]{0.1.4.5} se généralise aussitôt au cas où l'on a un système inductif $(A_\alpha,\vphi_{\beta\alpha})$ d'anneaux, et pour chaque indice $\alpha$ une partie $S_\alpha$ de $A_\alpha$ telle que
\oldpage[IV-4]{227}
$\vphi_{\beta\alpha}(S_\alpha)\subset S_\beta$ pour $\alpha\leq\beta$ ; on prend alors pour $S$ la limite inductive dans $A=\varinjlim A_\alpha$ du système inductif de parties $(S_\alpha)$.
\end{env}

\begin{env}[20.1.2]
\EGAReaderTarget{ega4.c04.statement.20-1-2}{IV.20.1.2}
Soit maintenant $\sh{F}$ un $\sh{O}_X$-module.
On pose alors
\[
\label{IV.20.1.2.1}
  \sh{F}[\sh{S}^{-1}]=\sh{F}\otimes_{\sh{O}_X}\sh{O}_X[\sh{S}^{-1}]
  \tag{20.1.2.1}
\]
et on dit que c'est le \emph{faisceau de modules de fractions de $\sh{F}$ à dénominateurs dans $\sh{S}$} ; il est immédiat qu'il est associé au préfaisceau de modules $U\mapsto\Gamma(U,\sh{F})[\Gamma(U,\sh{S})^{-1}]$, et que pour tout $x\in X$, on a un isomorphisme canonique
\[
\label{IV.20.1.2.2}
  (\sh{F}[\sh{S}^{-1}])_x\isoto\sh{F}_x[\sh{S}_x^{-1}].
  \tag{20.1.2.2}
\]
\end{env}

\begin{env}[20.1.3]
\EGAReaderTarget{ega4.c04.statement.20-1-3}{IV.20.1.3}
Nous allons nous intéresser ici au cas où $\sh{S}$ est le sous-faisceau $\sh{S}(\sh{O}_X)$ de $\sh{O}_X$ tel que pour tout ouvert $U$, $\Gamma(U,\sh{S})$ soit l'\emph{ensemble des éléments réguliers} de l'anneau $\Gamma(U,\sh{O}_X)$ ; il est immédiat qu'il s'agit d'un faisceau (et non seulement d'un préfaisceau), la régularité d'une section de $\sh{O}_X$ au-dessus de $U$ se vérifiant « fibre par fibre » (i.e. signifiant que le germe de la section en $x$ est régulier dans $\sh{O}_{X,x}$ pour tout $x\in U$), autrement dit $\sh{S}(\sh{O}_X)_x$ n'est autre que l'ensemble des éléments réguliers de $\sh{O}_{X,x}$.
Le faisceau d'anneaux correspondant
\[
  \sh{M}_X=\sh{O}_X[\sh{S}^{-1}]
\]
est appelé le \emph{faisceau des germes de fonctions méromorphes sur $X$}, et les sections de $\sh{M}_X$ au-dessus de $X$ sont appelées les \emph{fonctions méromorphes sur $X$} ; elles forment un anneau que l'on note $M(X)$.
Pour tout $\sh{O}_X$-module $\sh{F}$,
\[
  \sh{F}\otimes_{\sh{O}_X}\sh{M}_X=\sh{F}[\sh{S}^{-1}]
\]
est encore noté $\sh{M}_X(\sh{F})$ et appelé le \emph{faisceau des germes de sections méromorphes de $\sh{F}$} ; ses sections au-dessus de $X$ forment un $M(X)$-module noté $M(X,\sh{F})$, dont les éléments sont appelés \emph{sections méromorphes de $\sh{F}$ au-dessus de $X$}.
Ces définitions impliquent que pour tout ouvert $U$ de $X$, on a un isomorphisme canonique $\sh{M}_X(\sh{F})|U\isoto\sh{M}_U(\sh{F}|U)$, en particulier $\sh{M}_X|U\isoto\sh{M}_U$.
\end{env}

\begin{env}[20.1.3.1]
\label{IV.20.1.3.1}
Si $X$ est un \emph{préschéma réduit}, on notera que si un élément $s\in\Gamma(U,\sh{O}_X)$ est tel que $s_\xi\neq0$ pour tout point maximal $\xi$ de $U$, alors $s$ est \emph{régulier}.
En effet, si $st=0$ pour un $t\in\Gamma(U,\sh{O}_X)$, on a $s_\xi t_\xi=0$, donc $t_\xi=0$ puisque $\sh{O}_{X,\xi}$ est un corps, et dire que $t_\xi=0$ pour tout point maximal $\xi$ de $X$ signifie que $t=0$ : en effet, on est aussitôt ramené au cas où $U$ est affine, et un élément d'un anneau réduit qui appartient à tous les idéaux premiers minimaux est nul par définition.
La réciproque est vraie si l'ensemble des composantes irréductibles de $X$ est \emph{localement fini}.
On est en effet aussitôt ramené au cas où $X=\Spec(A)$ est affine ; si $\mathfrak p_i$ ($1\leq i\leq n$) sont les idéaux premiers minimaux de $A$ et si $s\in\mathfrak p_i$ pour un indice $i$, il existe $t\in A$ tel que $t\in\mathfrak p_j$ pour $j\neq i$ et $t\notin\mathfrak p_i$ (Bourbaki, \emph{Alg. comm.}, chap.~II, \textsection1, n\textsuperscript{o}1, prop.~1) ; on a donc $st\in\mathfrak p_i$ pour tout $i$, et par suite $st=0$ puisque $A$ est réduit ; $s$ est donc non régulier.
\end{env}

\begin{env}[20.1.4]
\EGAReaderTarget{ega4.c04.statement.20-1-4}{IV.20.1.4}
Pour tout ouvert $U$ de $X$, l'homomorphisme $t\mapsto t/1$ de $\Gamma(U,\sh{O}_X)$ dans $\Gamma(U,\sh{O}_X)[\Gamma(U,\sh{S})^{-1}]$ (qui n'est autre que l'\emph{anneau total des fractions} de
\oldpage[IV-4]{228}
$\Gamma(U,\sh{O}_X)$) est injectif ; ces homomorphismes définissent donc un \emph{homomorphisme canonique injectif}
\[
\label{IV.20.1.4.1}
  i:\sh{O}_X\to\sh{M}_X
  \tag{20.1.4.1}
\]
qui permet d'identifier $\sh{O}_X$ à un sous-faisceau de $\sh{M}_X$.
Étant donnée une fonction méromorphe $\vphi\in M(X)$, on dit que $\vphi$ est \emph{définie} dans un ouvert $U$ de $X$ si $\vphi|U$ est une \emph{section de $\sh{O}_U$} au-dessus de $U$ ; les axiomes des faisceaux montrent qu'il y a, pour une section $\vphi$ donnée, un \emph{plus grand} ouvert dans lequel $\vphi$ est définie ; on l'appelle \emph{domaine de définition} de $\vphi$ et on le note $\operatorname{dom}(\vphi)$.
\end{env}

\begin{env}[20.1.5]
\EGAReaderTarget{ega4.c04.statement.20-1-5}{IV.20.1.5}
Pour tout $\sh{O}_X$-module $\sh{F}$, on déduit de \sref{IV.20.1.4.1} un di-homomorphisme formé de $i$ et de l'homomorphisme de faisceaux de groupes additifs
\[
\label{IV.20.1.5.1}
  1_{\sh{F}}\otimes i:\sh{F}\to\sh{M}_X(\sh{F})=\sh{F}\otimes_{\sh{O}_X}\sh{M}_X.
  \tag{20.1.5.1}
\]
On notera que ce dernier n'est plus injectif en général ; lorsqu'il est injectif, on dit que $\sh{F}$ est \emph{strictement sans torsion} : cela signifie que pour tout ouvert $U$ de $X$ et toute section $s\in\Gamma(U,\sh{O}_X)$ qui est un élément régulier de cet anneau, l'homothétie $z\mapsto sz$ de $\Gamma(U,\sh{F})$ est injective ; cette condition est évidemment remplie si $\sh{F}$ est \emph{localement libre}.
\end{env}

\begin{proposition}[20.1.6]
\label{IV.20.1.6}
Soient $X$ un préschéma localement noethérien, $\sh{F}$ un $\sh{O}_X$-module quasi-cohérent.
Pour que $\sh{F}$ soit strictement sans torsion, il faut et il suffit que $\operatorname{Ass}(\sh{F})\subset\operatorname{Ass}(\sh{O}_X)$.
\end{proposition}

\begin{proof}
On est en effet aussitôt ramené au cas où $X=\Spec(A)$ est affine, $\sh{F}=\widetilde M$, et on sait que les éléments $s$ de $A$ appartenant à un idéal de $\operatorname{Ass}(M)$ sont exactement ceux pour lesquels l'homothétie $z\mapsto sz$ n'est pas injective (Bourbaki, \emph{Alg. comm.}, chap.~IV, \textsection1, n\textsuperscript{o}1, cor.~2 de la prop.~2).
\end{proof}

\begin{env}[20.1.7]
\EGAReaderTarget{ega4.c04.statement.20-1-7}{IV.20.1.7}
Si $u$ est une section de $\sh{M}_X(\sh{F})$ au-dessus de $X$, on dit que $u$ est \emph{définie} en un point $x\in X$ s'il existe un voisinage ouvert $V$ de $x$ dans $X$ tel que $u|V$ soit l'image d'une section de $\sh{F}$ au-dessus de $V$ par le di-homomorphisme \sref{IV.20.1.5.1}.
On dit que $u$ est \emph{définie} dans un ouvert $U$ de $X$ si elle est définie en tout point de $U$ ; il y a encore un plus grand ouvert dans lequel $u$ est définie, appelé \emph{domaine de définition} de $u$ et noté $\operatorname{dom}(u)$.
Lorsque $\sh{F}$ est strictement sans torsion, de sorte que $\sh{F}$ s'identifie par \sref{IV.20.1.5.1} à un sous-faisceau de $\sh{M}_X(\sh{F})$, dire que $u$ est définie dans $U$ signifie que $u|U$ est une \emph{section de $\sh{F}$ au-dessus de $U$}.
\end{env}

\begin{env}[20.1.8]
\EGAReaderTarget{ega4.c04.statement.20-1-8}{IV.20.1.8}
Conformément aux notations générales \sref[0\textsubscript{I}]{0.5.4.7}, on note $\sh{M}_X^*$ le faisceau de groupes multiplicatifs tel que $\Gamma(U,\sh{M}_X^*)$ soit (pour tout ouvert $U$ de $X$) le groupe des \emph{éléments inversibles} de $\Gamma(U,\sh{M}_X)$.
Ce faisceau n'est autre que le faisceau $\sh{S}(\sh{M}_X)$ défini en \sref{IV.20.1.3} : en effet, si $s\in\Gamma(U,\sh{S}(\sh{M}_X))$, pour tout $x\in U$, il existe un voisinage ouvert $V\subset U$ de $x$ tel que $s|V$ soit un élément régulier dans l'\emph{anneau total des fractions} de $\Gamma(V,\sh{O}_X)$ et on sait qu'un tel élément est nécessairement inversible dans cet anneau de fractions.
Nous dirons que les sections de $\sh{M}_X^*$ au-dessus de $X$ sont les \emph{fonctions méromorphes régulières} (on notera que nous nous écartons ici de la terminologie suivie par certains auteurs, qui nomment fonctions méromorphes « régulières » celles qui sont des \emph{sections de $\sh{O}_X$}, identifié à un sous-faisceau de $\sh{M}_X$).

Soit $\sh{L}$ un $\sh{O}_X$-module inversible \sref[0\textsubscript{I}]{0.5.4.1} ; alors il est clair que $\sh{M}_X(\sh{L})=\sh{L}\otimes_{\sh{O}_X}\sh{M}_X$
\oldpage[IV-4]{229}
est un $\sh{M}_X$-module inversible.
Soit $U$ un ouvert tel que $\sh{L}|U$ soit isomorphe à $\sh{O}_U$ ; comme tout automorphisme de $\sh{M}_U$ est la multiplication par un élément inversible de $\Gamma(U,\sh{M}_X)$ \sref[0\textsubscript{I}]{0.5.4.7}, il revient au même de dire qu'une section $s\in\Gamma(U,\sh{M}_X(\sh{L}))$ a une image inversible dans $\Gamma(U,\sh{M}_X)$ par \emph{un} isomorphisme ou par \emph{tout} isomorphisme sur $\Gamma(U,\sh{M}_X)$ ; on dira dans ce cas que $s$ est une \emph{section méromorphe régulière de $\sh{L}$ au-dessus de $U$} ; une section $s$ de $\sh{L}$ au-dessus de $X$ sera appelée une \emph{section méromorphe régulière de $\sh{L}$} si, pour tout ouvert $U$ tel que $\sh{L}|U$ soit isomorphe à $\sh{O}_U$, $s|U$ est une section méromorphe régulière de $\sh{L}$ au-dessus de $U$.
On notera $(\sh{M}_X(\sh{L}))^*$ le sous-faisceau de $\sh{M}_X(\sh{L})$ tel que pour tout ouvert $U$, $\Gamma(U,(\sh{M}_X(\sh{L}))^*)$ soit l'ensemble des sections méromorphes régulières de $\sh{L}$ au-dessus de $U$.
Soit $s$ une section méromorphe de $\sh{L}$ au-dessus de $X$ (i.e. une section de $\sh{M}_X(\sh{L})$) ; elle définit un homomorphisme $h_s:\sh{M}_X\to\sh{M}_X(\sh{L})$ qui, à toute section $t$ de $\sh{M}_X$ au-dessus d'un ouvert $U$, fait correspondre $(s|U)t$.
Il résulte aussitôt de ce qui précède que, pour que $s$ soit \emph{régulière}, il faut et il suffit que $h_s$ soit \emph{injectif}, et en fait $h_s$ est alors un homomorphisme \emph{bijectif} de $\sh{M}_X$ sur $\sh{M}_X(\sh{L})$, et sa restriction à $\sh{M}_X^*$ est une bijection sur $(\sh{M}_X(\sh{L}))^*$.
On en conclut que l'homothétie $t\mapsto ts$ est un isomorphisme de $M(X)$ sur $M(X,\sh{L})$.
\end{env}

\begin{env}[20.1.9]
\label{IV.20.1.9}
Soit $s$ une section méromorphe régulière du $\sh{O}_X$-module inversible $\sh{L}$ au-dessus de $X$ ; alors, pour tout $\sh{O}_X$-module $\sh{F}$, $s$ définit de même un homomorphisme $h_s\otimes1_{\sh{F}}:\sh{M}_X(\sh{F})\to\sh{M}_X(\sh{F}\otimes_{\sh{O}_X}\sh{L})$, qui est encore \emph{bijectif}.
\end{env}

\begin{env}[20.1.10]
\label{IV.20.1.10}
Soit $s$ une section méromorphe d'un $\sh{O}_X$-module inversible $\sh{L}$ au-dessus de $X$ ; pour que $s$ soit régulière, il faut et il suffit qu'il existe une section méromorphe $s'$ de $\sh{L}^{-1}$ au-dessus de $X$ telle que l'image canonique de $s\otimes s'$ dans $\sh{M}_X$ \sref[0\textsubscript{I}]{0.5.4.3} soit la section unité, et cette section $s'$ est alors unique : en effet, la nécessité de l'existence locale d'une telle section est évidente, et son unicité locale entraîne son existence (et unicité) globale ; en outre, l'existence de $s'$ est trivialement suffisante pour que $s$ soit régulière.
On posera $s'=s^{-1}$.

Enfin, si $\sh{L}'$ est un second $\sh{O}_X$-module inversible, $s$ (resp. $s'$) une section méromorphe régulière de $\sh{L}$ (resp. $\sh{L}'$) au-dessus de $X$, $s\otimes s'$ est évidemment une section méromorphe régulière de $\sh{L}\otimes\sh{L}'$ au-dessus de $X$.
\end{env}

\begin{env}[20.1.11]
\EGAReaderTarget{ega4.c04.statement.20-1-11}{IV.20.1.11}
Si $f:X'\to X$ est un morphisme d'espaces annelés, il n'y a pas en général d'application naturelle associant à une fonction méromorphe sur $X$ une fonction méromorphe sur $X'$.
Par exemple, si $X$ est le spectre d'un anneau local intègre $A$, $X'$ celui de son corps résiduel $k$, il n'y a aucun homomorphisme naturel du corps des fractions $K$ de $A$ dans $k$, et on ne peut faire correspondre à un élément de $K$ un élément de $k$ que s'il est déjà dans $A$.

De façon générale, si $f=(\psi,\theta)$, notons, pour tout ouvert $U$ de $X$, $\sh{S}_f(U)$ l'ensemble des sections \emph{régulières} $s\in\Gamma(U,\sh{O}_X)$ telles que l'image de $s$ par
\[
  \Gamma(\theta^\#):\Gamma(U,\sh{O}_X)\to\Gamma(f^{-1}(U),\sh{O}_{X'})
\]
soit une section \emph{régulière}.
Il est immédiat que $U\mapsto\sh{S}_f(U)$ est un \emph{sous-faisceau} du faisceau d'ensembles $\sh{S}(\sh{O}_X)$, que l'on note $\sh{S}_f$.
On pose $\sh{M}_f=\sh{O}_X[\sh{S}_f^{-1}]$ ; c'est un sous-faisceau
\oldpage[IV-4]{230}
d'anneaux de $\sh{M}_X$, et on déduit canoniquement de $\theta^\#:\psi^*(\sh{O}_X)\to\sh{O}_{X'}$ un homomorphisme de faisceaux d'anneaux $\theta^{\prime\#}:\psi^*(\sh{M}_f)\to\sh{M}_{X'}$ prolongeant $\theta^\#$ (Bourbaki, \emph{Alg. comm.}, chap.~II, \textsection2, n\textsuperscript{o}1, prop.~2) ; d'où, en se rappelant que $f^*(\sh{M}_f)=\psi^*(\sh{M}_f)\otimes_{\psi^*(\sh{O}_X)}\sh{O}_{X'}$, un homomorphisme canonique de $\sh{O}_{X'}$-algèbres
\[
\label{IV.20.1.11.1}
  f^*(\sh{M}_f)\to\sh{M}_{X'}.
  \tag{20.1.11.1}
\]
Pour toute fonction méromorphe $\vphi$ sur $X$, qui \emph{est une section de $\sh{M}_f$}, $\Gamma(\theta^{\prime\#})(\vphi)$ est une fonction méromorphe sur $X'$, dite \emph{image réciproque de $\vphi$ par $f$}, et notée $\vphi\circ f$ si cela n'entraîne pas confusion.

De même, si $\sh{F}$ est un $\sh{O}_X$-module, on pose $\sh{M}_f(\sh{F})=\sh{F}\otimes_{\sh{O}_X}\sh{M}_f$, et on déduit aussitôt de $\theta^{\prime\#}$ un homomorphisme canonique (qu'on écrit aussi $u\mapsto u\circ f$)
\[
  \Gamma(X,\sh{M}_f(\sh{F}))\to\Gamma(X',\sh{M}_{X'}(f^*(\sh{F}))).
\]
En outre, si $u\in\Gamma(X,\sh{M}_f(\sh{F}))$ est définie \sref{IV.20.1.7} en un point $x$, $u$ coïncide, dans un voisinage $U$ de $x$, avec une section de la forme $\sum_i h_i\otimes(t_i/s_i)$, où les $h_i$ appartiennent à $\Gamma(U,\sh{F})$, les $t_i$ à $\Gamma(U,\sh{O}_X)$ et les $s_i$ à $\Gamma(U,\sh{S}_f)$.
Comme par hypothèse les images des $s_i$ dans $\Gamma(f^{-1}(U),\sh{O}_{X'})$ sont régulières, on voit que $u\circ f$ est définie en tout point de $f^{-1}(U)$ ; autrement dit, on a
\[
\label{IV.20.1.11.2}
  f^{-1}(\operatorname{dom}(u))\subset\operatorname{dom}(u\circ f).
  \tag{20.1.11.2}
\]
Nous verrons plus loin \sref{IV.20.6.5}[(i)] des exemples (avec $\sh{F}=\sh{O}_X$) où les deux membres de \sref{IV.20.1.11.2} peuvent être distincts.

Considérons en particulier le cas où $\sh{M}_f=\sh{M}_X$ ; alors, si $\sh{L}$ est un $\sh{O}_X$-module inversible, l'image dans $\sh{M}_{X'}(f^*(\sh{L}))$, par $\Gamma(\theta^{\prime\#})$, d'une section méromorphe \emph{régulière} de $\sh{L}$ au-dessus de $X$ \sref{IV.20.1.8} est une section méromorphe \emph{régulière} de $f^*(\sh{L})$ au-dessus de $X'$, comme il résulte aussitôt de la définition de ces sections, et du fait qu'un homomorphisme d'anneaux transforme un élément inversible en un élément inversible.

Soit $f':X''\to X'$ un second morphisme d'espaces annelés, et supposons que $\sh{M}_f=\sh{M}_X$ et $\sh{M}_{f'}=\sh{M}_{X'}$ ; alors, si l'on pose $f''=f\circ f'$, on a aussi $\sh{M}_{f''}=\sh{M}_X$, et l'on voit aussitôt que pour toute section méromorphe $u$ de $\sh{F}$ au-dessus de $X$, on a $u\circ f''=(u\circ f)\circ f'$.
\end{env}

\begin{proposition}[20.1.12]
\label{IV.20.1.12}
Si le morphisme $f:X'\to X$ est plat \sref[0\textsubscript{I}]{0.6.7.1}, on a $\sh{M}_f=\sh{M}_X$, et l'homomorphisme $\vphi\mapsto\vphi\circ f$ est défini dans $M(X)$ tout entier.
En outre, si $f$ est un morphisme (plat) d'espaces annelés en anneaux locaux, on a $\operatorname{dom}(\vphi\circ f)=f^{-1}(\operatorname{dom}(\vphi))$ ; si de plus $f$ est surjectif (donc fidèlement plat), l'homomorphisme $\vphi\mapsto\vphi\circ f$ est injectif.
\end{proposition}

\begin{proof}
La première assertion résulte de ce que, si $B$ est une $A$-algèbre qui est un $A$-module plat, tout élément de $A$ non diviseur de $0$ dans $A$ est non diviseur de $0$ dans $B$ \sref[0\textsubscript{I}]{0.6.3.4}.
Pour prouver les autres assertions, notons que, pour tout $x'\in X'$, si $x=f(x')$, $\sh{O}_{X',x'}$ est un $\sh{O}_{X,x}$-module plat, et comme l'homomorphisme $\sh{O}_{X,x}\to\sh{O}_{X',x'}$ est \emph{local} par hypothèse, il est injectif (\sref[0\textsubscript{I}]{0.6.5.1} et \sref[0\textsubscript{I}]{0.6.6.2}) ; si l'on pose $A=\sh{O}_{X,x}$, $B=\sh{O}_{X',x'}$, de sorte que $A$ s'identifie à un sous-anneau de $B$, $(f^*(\sh{M}_X))_{x'}$ est égal à $S^{-1}A\otimes_A B=S^{-1}B$, où $S$ est l'ensemble des éléments réguliers de $A$, $(\sh{M}_{X'})_{x'}$ est égal à $T^{-1}B$, où $T$ est l'ensemble
\oldpage[IV-4]{231}
des éléments réguliers de $B$, et comme on a vu que $S\subset T$, l'homomorphisme $S^{-1}B\to T^{-1}B$ est injectif ; autrement dit, cela prouve que l'homomorphisme \sref{IV.20.1.11.1} $f^*(\sh{M}_X)\to\sh{M}_{X'}$ est \emph{injectif} (d'où la dernière assertion de l'énoncé).
Le quotient $f^*(\sh{M}_X)/\sh{O}_{X'}$ s'identifie à un sous-$\sh{O}_{X'}$-module de $\sh{M}_{X'}/\sh{O}_{X'}$, et $(f^*(\sh{M}_X)/\sh{O}_{X'})_{x'}$ s'identifie à $(\sh{M}_X/\sh{O}_X)_x\otimes_{\sh{O}_{X,x}}\sh{O}_{X',x'}$.
Supposons alors que $x\notin\operatorname{dom}(\vphi)$ ; l'image de $\vphi_x$ dans $(\sh{M}_X/\sh{O}_X)_x$ est donc $\neq0$ ; par fidèle platitude, on en déduit qu'il en est de même de l'image de $(\vphi\circ f)_{x'}$ dans $(\sh{M}_{X'}/\sh{O}_{X'})_{x'}$, donc $x'\notin\operatorname{dom}(\vphi\circ f)$, ce qui achève la démonstration.
\end{proof}

\begin{remark}[20.1.13]
\label{IV.20.1.13}
Soit $X$ un espace analytique complexe \emph{réduit} ; alors la notion de fonction méromorphe sur $X$ définie ci-dessus coïncide avec la notion usuelle.
Considérons d'autre part un préschéma $Y$, localement de type fini sur le corps $\bb{C}$ ; on sait alors qu'on peut associer à $Y$ un espace analytique $Y^\mathrm{an}$ ayant même espace topologique sous-jacent, et que le morphisme canonique $f:Y^\mathrm{an}\to Y$ est \emph{plat}~\cite{IV-37} ; en vertu de \sref{IV.20.1.12}, l'homomorphisme canonique $u\mapsto u\circ f$ de $M(Y)$ dans $M(Y^\mathrm{an})$ est donc partout défini et est injectif ; mais il n'est pas \emph{surjectif} en général.
Par exemple, lorsque $Y=\bb{V}_{\bb{C}}^r$ (\textbf{Err}\textsubscript{\textbf{III}},~14) est l'espace affine de dimension $r$ sur $\bb{C}$, $M(Y)$ s'identifie canoniquement au corps $R(Y)$ des fonctions rationnelles sur $Y$ \sref{IV.20.2.13}[(i)], tandis que $M(Y^\mathrm{an})$ est le corps des fonctions méromorphes usuelles sur $\bb{C}^r$.
En raison de ce fait, il est souvent préférable, en Géométrie algébrique, de s'abstenir de la terminologie introduite dans cette section, et d'utiliser la terminologie équivalente de « pseudo-fonction » qui va être définie ci-dessous.
\end{remark}

\subsection{Pseudo-morphismes et pseudo-fonctions}
\EGAReaderTarget{ega4.c04.subsec.20-2}{IV.20.2}

Les seuls espaces annelés dont il sera question dans cette section sont les \emph{préschémas}.

\begin{env}[20.2.1]
\EGAReaderTarget{ega4.c04.statement.20-2-1}{IV.20.2.1}
Rappelons \sref{IV.11.10.2} que dans un préschéma $X$, on dit qu'un ouvert $U$ est \emph{schématiquement dense} si, pour tout ouvert $V$ de $X$, l'homomorphisme canonique $\Gamma(V,\sh{O}_X)\to\Gamma(V\cap U,\sh{O}_X)$ est \emph{injectif}.

Considérons deux préschémas $X$, $Y$, deux ouverts \emph{schématiquement denses} $U$, $U'$ de $X$ ; on dit que deux morphismes $u:U\to Y$, $u':U'\to Y$ sont \emph{équivalents} s'il existe un ouvert $U''\subset U\cap U'$, \emph{schématiquement dense} dans $X$, tel que $u|U''=u'|U''$.
Comme il résulte aussitôt de la définition des ouverts schématiquement denses que l'intersection de deux tels ouverts en est un autre, il est immédiat que la relation précédente est bien une relation d'équivalence.
Une classe d'équivalence suivant cette relation s'appelle un \emph{pseudo-morphisme} de $X$ dans $Y$, ou une \emph{application rationnelle stricte} de $X$ dans $Y$.

Si $S$ est un préschéma et $X$, $Y$ deux $S$-préschémas, on appelle \emph{pseudo-$S$-morphisme} la classe d'équivalence (pour la relation précédente) d'un $S$-morphisme d'un ouvert schématiquement dense de $X$ dans $Y$.
Un pseudo-morphisme n'est donc autre chose qu'un pseudo-$S$-morphisme pour $S=\Spec(\bb{Z})$.

On note $\operatorname{Ps.hom}(X,Y)$ (resp. $\operatorname{Ps.hom}_S(X,Y)$) l'ensemble des pseudo-morphismes (resp. des pseudo-$S$-morphismes) de $X$ dans $Y$.
\end{env}

\begin{env}[20.2.2]
\EGAReaderTarget{ega4.c04.statement.20-2-2}{IV.20.2.2}
\oldpage[IV-4]{232}
Il résulte de la définition rappelée ci-dessus que si $U$ est un ouvert schématiquement dense dans $X$, alors, pour tout ouvert $V$ de $X$, $U\cap V$ est schématiquement dense dans $V$.
Si deux morphismes $u:U\to Y$, $u':U'\to Y$ d'ouverts schématiquement denses de $X$ dans $Y$ sont équivalents il s'ensuit donc que leurs restrictions $u|(V\cap U)$ et $u'|(V\cap U')$ sont aussi des morphismes équivalents (pour le préschéma induit sur $V$) ; leur classe est dite la \emph{restriction} à $V$ du pseudo-morphisme $\omega$, classe de $u$, et on note ce pseudo-morphisme $\omega|V$.
Si $W\subset V$ est un ouvert de $X$, il est clair que $(\omega|V)|W=\omega|W$.
Ceci montre que les applications de restriction définissent un \emph{préfaisceau} d'ensembles $U\mapsto\operatorname{Ps.hom}(U,Y)$ ; on notera que ce préfaisceau \emph{n'est pas} un faisceau en général (cf. \sref{IV.20.2.16}) ; le faisceau associé se note $\mathscr{P}\!s.\!\operatorname{hom}(X,Y)$.
Pour les pseudo-$S$-morphismes, on voit de même que $U\mapsto\operatorname{Ps.hom}_S(U,Y)$ est un préfaisceau d'ensembles, dont on note $\mathscr{P}\!s.\!\operatorname{hom}_S(X,Y)$ le faisceau associé.

Si $V$ est un ouvert \emph{schématiquement dense} dans $X$, alors, pour tout ouvert $U$ schématiquement dense dans $X$, $U\cap V$ est aussi schématiquement dense dans $X$, donc l'application $\omega\mapsto\omega|V$ est une \emph{bijection} de l'ensemble $\operatorname{Ps.hom}(X,Y)$ (resp. $\operatorname{Ps.hom}_S(X,Y)$) sur $\operatorname{Ps.hom}(V,Y)$ (resp. $\operatorname{Ps.hom}_S(V,Y)$).
\end{env}

\begin{env}[20.2.3]
\EGAReaderTarget{ega4.c04.statement.20-2-3}{IV.20.2.3}
Étant donné un pseudo-$S$-morphisme $\omega$ de $X$ dans $Y$, on dit qu'il est \emph{défini} en un point $x\in X$ s'il existe un voisinage ouvert $V$ de $x$ dans $X$, un ouvert $U\subset V$ \emph{contenant $x$ et schématiquement dense dans $V$}, et enfin un $S$-morphisme $u:U\to Y$ dont la classe dans $\operatorname{Ps.hom}_S(V,Y)$ soit égale à $\omega|V$ \sref{IV.20.2.2} ; on appelle \emph{domaine de définition} de $\omega$ et on note $\operatorname{dom}_S(\omega)$ (ou simplement $\operatorname{dom}(\omega)$ si $S=\Spec(\bb{Z})$) l'ensemble des $x\in X$ où $\omega$ est défini ; c'est évidemment un \emph{ouvert} de $X$.
En outre, pour tout ouvert $W$ de $X$, on a
\[
\label{IV.20.2.3.1}
  \operatorname{dom}_S(\omega|W)=(\operatorname{dom}_S(\omega))\cap W
  \tag{20.2.3.1}
\]
en vertu de la propriété des ouverts schématiquement denses rappelée dans \sref{IV.20.2.2}.
\end{env}

\begin{proposition}[20.2.4]
\label{IV.20.2.4}
Supposons que $X$, $Y$ soient des $S$-préschémas tels que le morphisme structural $Y\to S$ soit séparé ; alors, si $\omega$ est un pseudo-$S$-morphisme de $X$ dans $Y$, $\operatorname{dom}_S(\omega)$ est le plus grand des ouverts schématiquement denses $U$ de $X$ tels qu'il y ait un $S$-morphisme $u:U\to Y$ appartenant à la classe $\omega$.
\end{proposition}

\begin{proof}
Il suffit évidemment de prouver que si $U$, $U'$ sont deux ouverts schématiquement denses dans $X$ tels que deux $S$-morphismes $u:U\to Y$ et $u':U'\to Y$ soient équivalents, alors les restrictions de $u$ et $u'$ à $U\cap U'$ sont égales.
Or il existe par hypothèse un ouvert $U''\subset U\cap U'$, schématiquement dense dans $X$ et dans lequel $u$ et $u'$ coïncident ; comme $U''$ est aussi schématiquement dense dans $U\cap U'$, notre assertion résulte de \sref{IV.11.10.1}[(d)].
\end{proof}

\begin{corollary}[20.2.5]
\label{IV.20.2.5}
Soient $S$ un $S_0$-schéma, $X$, $Y$ deux $S$-préschémas tels que le morphisme composé $Y\to S\to S_0$ soit séparé (ce qui implique \sref[I]{I.5.5.1} que $Y\to S$ l'est aussi).
Pour tout pseudo-$S$-morphisme $\omega$ de $X$ dans $Y$, on a alors $\operatorname{dom}_S(\omega)=\operatorname{dom}_{S_0}(\omega)$.
En particulier, si $Y$ est un schéma, on a $\operatorname{dom}_S(\omega)=\operatorname{dom}(\omega)$.
\end{corollary}

\begin{proof}
En effet, il est clair que $\operatorname{dom}_S(\omega)\subset\operatorname{dom}_{S_0}(\omega)$ sans hypothèse de séparation ; en vertu de \sref{IV.20.2.4}, il suffit de prouver que si un $S_0$-morphisme $u_0:U_0\to Y$ d'un ouvert
\oldpage[IV-4]{233}
schématiquement dense $U_0$ de $X$ dans $Y$ est tel que le composé $v_0:U_0\xrightarrow{u_0}Y\to S$ coïncide avec la restriction du morphisme structural $w_0:U_0\to S$ dans un ouvert $U\subset U_0$ schématiquement dense dans $X$, alors on a $v_0=w_0$.
Mais en vertu de l'hypothèse que le morphisme $S\to S_0$ est séparé, cela résulte encore de \sref{IV.11.10.1}[(d)].
\end{proof}

\begin{corollary}[20.2.6]
\label{IV.20.2.6}
Sous les hypothèses de \sref{IV.20.2.4}, le préfaisceau $U\mapsto\operatorname{Ps.hom}_S(U,Y)$ est un faisceau.
\end{corollary}

\begin{proof}
En effet, soient $U$ un ouvert de $X$, $(U_\alpha)$ un recouvrement de $U$ par des ouverts de $U$ ; supposons donné pour chaque $\alpha$ un pseudo-$S$-morphisme $\omega_\alpha$ de $U_\alpha$ dans $Y$, de sorte que pour tout couple d'indices $\alpha$, $\beta$, les restrictions \sref{IV.20.2.2} des pseudo-$S$-morphismes $\omega_\alpha$ et $\omega_\beta$ à $U_\alpha\cap U_\beta$ soient égales ; en vertu de \sref{IV.20.2.3.1}, cela entraîne que $\operatorname{dom}_S(\omega_\alpha)\cap U_\beta=\operatorname{dom}_S(\omega_\beta)\cap U_\alpha$.
Soit alors $W$ la réunion des ouverts $\operatorname{dom}_S(\omega_\alpha)$, et, pour tout $\alpha$, soit $u_\alpha$ le $S$-morphisme unique $\operatorname{dom}_S(\omega_\alpha)\to Y$ appartenant à la classe $\omega_\alpha$ \sref{IV.20.2.4} ; en raison de l'hypothèse et de \sref{IV.20.2.4}, les restrictions de $u_\alpha$ et $u_\beta$ à $\operatorname{dom}_S(\omega_\alpha)\cap\operatorname{dom}_S(\omega_\beta)$ sont égales, donc il existe un morphisme $u:W\to Y$ dont la restriction à chaque ouvert $\operatorname{dom}_S(\omega_\alpha)$ est égale à $u_\alpha$ ; il est clair que $W$ est schématiquement dense dans $U$, donc définit un pseudo-$S$-morphisme $\omega$ de $U$ dans $Y$ dont les restrictions aux $U_\alpha$ sont les $\omega_\alpha$.
\end{proof}

\begin{remark}[20.2.7]
\label{IV.20.2.7}
On sait \sref{IV.11.10.4} que lorsque $X$ est \emph{réduit}, il revient au même de dire qu'un ouvert de $X$ est dense dans $X$ ou qu'il est schématiquement dense dans $X$ ; la notion de pseudo-morphisme (resp. pseudo-$S$-morphisme) de $X$ dans $Y$ coïncide alors avec celle d'\emph{application rationnelle} (resp. $S$-\emph{application rationnelle}) de $X$ dans $Y$ \sref[I]{I.7.1.2}.
Dans le cas général, la notion de pseudo-morphisme semble plus utile que celle d'application rationnelle.
\end{remark}

\begin{env}[20.2.8]
\EGAReaderTarget{ega4.c04.statement.20-2-8}{IV.20.2.8}
On appelle \emph{pseudo-fonction} sur $X$ un pseudo-morphisme de $X$ dans $\Spec(\bb{Z}[T])$ ($T$ indéterminée), ou, ce qui revient au même, un $X$-pseudo-morphisme de $X$ dans $X\otimes_{\bb{Z}}\bb{Z}[T]$ ; il revient aussi au même \sref[I]{I.3.3.15} de dire qu'une pseudo-fonction sur $X$ est une \emph{classe d'équivalence de sections de $\sh{O}_X$ au-dessus d'ouverts schématiquement denses de $X$}, deux sections $g\in\Gamma(U,\sh{O}_X)$, $g'\in\Gamma(U',\sh{O}_X)$ au-dessus de tels ouverts étant équivalentes s'il existe un ouvert $U''\subset U\cap U'$, schématiquement dense dans $X$, où $g$ et $g'$ coïncident.
On peut ici appliquer \sref{IV.20.2.4} avec $S=X$ et $Y=X\otimes_{\bb{Z}}\bb{Z}[T]$ ; donc, pour toute pseudo-fonction $\vphi$ sur $X$, il existe un \emph{plus grand} ouvert schématiquement dense $\operatorname{dom}(\vphi)$ dans $X$ tel qu'il y ait une section de $\sh{O}_X$ au-dessus de $\operatorname{dom}(\vphi)$ appartenant à la classe $\vphi$.
Il est clair que le faisceau $\mathscr{P}\!s.\!\operatorname{hom}_X(X,X\otimes_{\bb{Z}}\bb{Z}[T])$ est ici un \emph{faisceau d'anneaux}, et même une $\sh{O}_X$-\emph{Algèbre}, que nous noterons $\sh{M}'_X$ ; il résulte de \sref{IV.20.2.6} que, pour tout ouvert $U$ de $X$, $\Gamma(U,\sh{M}'_X)$ est égal à l'ensemble des pseudo-morphismes de $U$ dans $\Spec(\bb{Z}[T])$ ; on pose $M'(X)=\Gamma(X,\sh{M}'_X)$.
Dire qu'une section $\vphi$ de $\sh{M}'_X$ au-dessus de $U$ est \emph{inversible} dans l'anneau $\Gamma(U,\sh{M}'_X)$ signifie qu'il existe un ouvert $U'$ schématiquement dense dans $\operatorname{dom}(\vphi)$, donc dans $U$, tel que, si $g$ est l'unique section de $\sh{O}_X$ au-dessus de $\operatorname{dom}(\vphi)$ appartenant à $\vphi$, $g|U'$ soit \emph{inversible} dans $\Gamma(U',\sh{O}_X)$.
Il résulte de \sref[I]{I.3.3.15} que, dans la correspondance canonique entre $\Gamma(V,\sh{O}_X)$ et $\Hom(V,\bb{Z}[T])$ ($V$ ouvert de $X$), les éléments inversibles de $\Gamma(V,\sh{O}_X)$ correspondent aux morphismes qui se
\oldpage[IV-4]{234}
\emph{factorisent} en $V\to\Spec(\bb{Z}[T,T^{-1}])\to\Spec(\bb{Z}[T])$.
On en conclut que le faisceau ${\sh{M}'}^*_X$ des germes de sections inversibles de $\sh{M}'_X$ s'identifie canoniquement au faisceau $\mathscr{P}\!s.\!\operatorname{hom}_X(X,X\otimes_{\bb{Z}}\bb{Z}[T,T^{-1}])$.
\end{env}

\begin{lemma}[20.2.9]
\label{IV.20.2.9}
Soient $U$ un ouvert de $X$, $s$ un élément régulier de l'anneau $\Gamma(U,\sh{O}_X)$ ; alors l'ensemble ouvert $U_s$ des $x\in U$ tels que $s(x)\neq0$ est schématiquement dense dans $U$.
\end{lemma}

\begin{proof}
En effet, soient $V$ un ouvert de $U$, $t$ une section de $\sh{O}_X$ au-dessus de $V$ telle que $t|(V\cap U_s)=0$.
Pour tout ouvert affine $W\subset V$, il existe alors un entier $n$ tel que $s^nt|W=0$ \sref[I]{I.1.4.1} ; mais puisque $s$ est un élément \emph{régulier} de $\Gamma(U,\sh{O}_X)$, cela entraîne $t|W=0$, d'où $t=0$.
\end{proof}

\begin{env}[20.2.10]
\label{IV.20.2.10}
Considérons une fonction méromorphe $f\in M(X)$ \sref{IV.20.1.4} ; alors $\operatorname{dom}(f)$ est \emph{schématiquement dense} dans $X$ : en effet, tout point de $X$ admet un voisinage ouvert $U$ tel qu'il existe une section $s\in\Gamma(U,\sh{O}_X)$ qui est un élément régulier de cet anneau, et pour laquelle $s(f|U)\in\Gamma(U,\sh{O}_X)$ ; comme $s|U_s$ est inversible, on en conclut que $f|U_s\in\Gamma(U_s,\sh{O}_X)$, donc $U_s\subset\operatorname{dom}(f)$ par définition \sref{IV.20.1.4}, et notre assertion résulte donc du lemme \sref{IV.20.2.9}.
On peut donc associer à $f$ la pseudo-fonction égale à la classe de la section de $\sh{O}_X$ au-dessus de $\operatorname{dom}(f)$, restriction de $f$ ; opérant de même en remplaçant $X$ par un ouvert quelconque de $X$, on obtient ainsi un \emph{homomorphisme canonique de $\sh{O}_X$-Algèbres}
\[
\label{IV.20.2.10.1}
  \sh{M}_X\to\sh{M}'_X
  \tag{20.2.10.1}
\]
qui, par restriction, donne évidemment un homomorphisme de faisceaux de groupes abéliens
\[
\label{IV.20.2.10.2}
  \sh{M}_X^*\to{\sh{M}'}^*_X
  \tag{20.2.10.2}
\]
pour les faisceaux de germes inversibles de sections de $\sh{M}_X$ et $\sh{M}'_X$.
\end{env}

\begin{proposition}[20.2.11]
\EGAReaderTarget{ega4.c04.prop.20-2-11}{IV.20.2.11}
\medskip
\begin{enumerate}[label=\emph{(\roman*)}]
  \item L'homomorphisme canonique \emph{\sref{IV.20.2.10.1}} (et par suite aussi \emph{\sref{IV.20.2.10.2}}) est \emph{injectif}.
  \item Supposons, soit que $X$ soit localement noethérien, soit que $X$ soit réduit et que l'ensemble de ses composantes irréductibles soit localement fini.
  Alors l'homomorphisme canonique \emph{\sref{IV.20.2.10.1}} (et par suite aussi \emph{\sref{IV.20.2.10.2}}) est \emph{bijectif}.
\end{enumerate}
\end{proposition}

\begin{proof}
Les questions envisagées étant locales sur $X$, on peut se borner à considérer le cas où $X=\Spec(A)$ est affine, et à montrer alors que l'homomorphisme canonique $M(X)\to M'(X)$ est toujours injectif, et qu'il est bijectif dans les cas considérés dans~(ii).
Compte tenu de la définition du faisceau $\sh{M}_X$ \sref{IV.20.1.3}, on peut d'ailleurs remarquer que \sref{IV.20.2.10.1} provient en fait d'un homomorphisme de \emph{préfaisceaux}
\[
  \Gamma(U,\sh{O}_X)[\Gamma(U,\sh{S})^{-1}]\to\Gamma(U,\sh{M}'_X)
\]
et il suffit de montrer que, pour $U$ affine, ce dernier est injectif (resp. bijectif sous les hypothèses de~(ii)).
Notant $S$ l'ensemble des éléments réguliers de $A$ (de sorte que $S^{-1}A$ est \emph{l'anneau total des fractions} de $A$), on doit donc considérer l'homomorphisme canonique
\[
\label{IV.20.2.11.1}
  S^{-1}A\to M'(X)
  \tag{20.2.11.1}
\]
\oldpage[IV-4]{235}
et montrer qu'il est injectif (resp. bijectif sous les conditions de~(ii)).
Or, on peut écrire $S^{-1}A=\varinjlim A_t$, où $t$ parcourt l'ensemble des éléments \emph{réguliers} de $A$, ordonné suivant la relation « $t$ est un diviseur de $t'$ » \sref[0]{0.1.4.5} ; on peut aussi écrire $A_t=\Gamma(\mathrm{D}(t),\sh{O}_X)$.
D'autre part on a par définition $M'(X)=\varinjlim\Gamma(U,\sh{O}_X)$, où $U$ parcourt l'ensemble des ouverts schématiquement denses de $X$ (ordonné par la relation $\supset$), et l'application \sref{IV.20.2.11.1} n'est autre que l'application canonique provenant du fait que les $\mathrm{D}(t)$ constituent une partie de l'ensemble des ouverts schématiquement denses dans $X$ \sref{IV.20.2.9}.
Notons que par définition d'un ouvert schématiquement dense, l'homomorphisme $\Gamma(U,\sh{O}_X)\to\Gamma(U',\sh{O}_X)$ pour deux tels ouverts $U'\subset U$ est toujours \emph{injectif}, donc il en est de même des homomorphismes $\Gamma(U,\sh{O}_X)\to M'(X)$, et on en conclut que \sref{IV.20.2.11.1} est injectif.
Pour prouver que \sref{IV.20.2.11.1} est bijectif, il suffit de montrer que l'ensemble des $\mathrm{D}(t)$ est \emph{cofinal} à l'ensemble des ouverts schématiquement denses, autrement dit que pour un tel ouvert $U$, il existe $t$ régulier dans $A$ tel que $\mathrm{D}(t)\subset U$.
Or, lorsque $X$ est réduit et que les composantes irréductibles de $X$ forment un ensemble localement fini, cet ensemble est fini puisque $X$ a été supposé affine, autrement dit $A$ n'a qu'un nombre fini d'idéaux premiers minimaux $\mathfrak{p}_i$ ; comme $A$ est réduit, l'intersection des $\mathfrak{p}_i$ est réduite à $0$, et dire que $t$ est régulier équivaut donc à dire que $t$ n'appartient à aucun des $\mathfrak{p}_i$ ; on conclut donc par le raisonnement de \sref[I]{I.7.1.9.1}.
Lorsque $A$ est noethérien, dire que $U=X-Y$ (où $Y=V(\mathfrak{J})$ est fermé dans $X$) est schématiquement dense signifie \sref{IV.5.10.2} que $Y$ \emph{ne rencontre pas} $\operatorname{Ass}(\sh{O}_X)$, et en vertu de Bourbaki, \emph{Alg.~comm.}, chap.~IV, \textsection1, n\textsuperscript{o}~4, prop.~8, cela entraîne l'existence d'un $t\in\mathfrak{J}$ tel que $t$ soit $A$-régulier, donc $U\supset\mathrm{D}(t)$.

On a en outre prouvé au cours de cette démonstration le
\end{proof}

\begin{corollary}[20.2.12]
\label{IV.20.2.12}
Si $X=\Spec(A)$, où $A$ est noethérien, ou réduit et n'ayant qu'un nombre fini d'idéaux premiers minimaux, les ouverts schématiquement denses dans $X$ sont ceux qui contiennent un ouvert de la forme $\mathrm{D}(t)$, où $t$ est régulier dans $A$, et $M(X)=M'(X)$ est l'anneau total des fractions $S^{-1}A$, où $S$ est l'ensemble des éléments réguliers de $X$.
\end{corollary}

\begin{remarks}[20.2.13]
\EGAReaderTarget{ega4.c04.statement.20-2-13}{IV.20.2.13}
\begin{enumerate}[label=(\roman*)]
  \item Soient $\vphi$ un élément de $M(X)$, $\vphi'$ son image dans $M'(X)$ ; on a évidemment, par définition (\sref{IV.20.1.4} et \sref{IV.20.2.8}) $\operatorname{dom}(\vphi)\subset\operatorname{dom}(\vphi')$ ; mais en fait on a l'égalité $\operatorname{dom}(\vphi)=\operatorname{dom}(\vphi')$, car il existe une section de $\sh{O}_X$ au-dessus de $\operatorname{dom}(\vphi')$, appartenant à la classe $\vphi'$, et la fonction méromorphe correspondante est égale à $\vphi$ \sref{IV.20.2.11}[(i)], donc $\operatorname{dom}(\vphi')\subset\operatorname{dom}(\vphi)$.

  \item On a déjà noté que lorsque $X$ est \emph{réduit}, on a $\sh{M}'_X=\sh{R}(X)$ (faisceau des fonctions rationnelles sur $X$ \sref[I]{I.7.3.2}) ; si en outre les composantes irréductibles de $X$ forment un ensemble localement fini, on a $\sh{M}_X=\sh{M}'_X=\sh{R}(X)$.
  En général, puisque tout ensemble ouvert schématiquement dense est dense, on a un homomorphisme canonique $\sh{M}'_X\to\sh{R}(X)$, mais même lorsque $X$ est localement noethérien, cet homomorphisme n'est pas nécessairement injectif.
  Par exemple, si $X=\Spec(A)$, où $A$ est un anneau noethérien tel que $\operatorname{Ass}(A)$ contienne des idéaux premiers immergés (ce qui entraîne que $A$ n'est pas réduit), on a vu que $M(X)$ et $M'(X)$ s'identifient à l'anneau total de fractions $S^{-1}A$,
  \oldpage[IV-4]{236}
  où $S$ est l'ensemble des éléments réguliers de $A$, complémentaire de la réunion des idéaux $\mathfrak{p}\in\operatorname{Ass}(A)$ ; par contre, $R(X)$ s'identifie à $Q^{-1}A$, où $Q$ est le complémentaire de la réunion des idéaux premiers \emph{minimaux} de $A$ \sref[I]{I.7.1.9}, et l'homomorphisme canonique $A\to Q^{-1}A$ (et \emph{a fortiori} $S^{-1}A\to Q^{-1}A$) n'est donc pas injectif, puisqu'il existe dans $A-Q$ des éléments $\neq0$ de $A$ annulés par des éléments de $Q$ (Bourbaki, \emph{Alg.~comm.}, chap.~IV, \textsection1, n\textsuperscript{o}~1, cor.~2 de la prop.~1).

  \item On notera que même lorsque $X$ est localement noethérien, le $\sh{O}_X$-Module $\sh{M}'_X$ \emph{n'est pas nécessairement quasi-cohérent}.
  Considérons par exemple un anneau local noethérien $A$ de dimension $\geq2$, dont l'idéal maximal $\mathfrak{m}$ est tel que $\mathfrak{m}\in\operatorname{Ass}(A)$ et que si l'on pose $X=\Spec(A)$, le schéma induit sur l'ouvert $U$ de $X$, complémentaire de $\{\mathfrak{m}\}$, soit \emph{intègre}.
  Les seuls éléments réguliers de $A$ sont alors les éléments inversibles, donc $\Gamma(X,\sh{M}'_X)=M'(X)=A$ ; si $\sh{M}'_X$ était quasi-cohérent, il serait donc égal à $\sh{O}_X$ ; mais comme $U$ est un schéma intègre, $\sh{M}'_X|U=R(U)$ est un faisceau \emph{simple} \sref[I]{I.7.3.5}, alors que $\sh{O}_U$ n'est pas un faisceau simple puisque $\dim(A)\geq2$.

  Reste à donner un exemple d'anneau $A$ ayant les propriétés précédentes.
  Il suffit de considérer un anneau local intègre $B$ de dimension $\geq2$ et de corps résiduel $k$, et de prendre $A=B\oplus k$ avec la loi de multiplication $(b,x)(b',x')=(bb',bx'+b'x)$.

  \item Si $X$ est localement noethérien, il résulte de \sref{IV.5.10.2} que les ouverts schématiquement denses dans $X$ sont ceux qui \emph{contiennent} l'ensemble $\operatorname{Ass}(\sh{O}_X)$.
\end{enumerate}
\end{remarks}

\begin{env}[20.2.14]
\EGAReaderTarget{ega4.c04.statement.20-2-14}{IV.20.2.14}
Soient $X$ un préschéma, $\sh{F}$ un $\sh{O}_X$-Module quasi-cohérent et \emph{strictement sans torsion} \sref{IV.20.1.5}, de sorte que $\sh{F}$ s'identifie à un sous-$\sh{O}_X$-Module de $\sh{M}_X(\sh{F})$.
Pour toute section méromorphe $u$ de $\sh{F}$ au-dessus de $X$, on appelle l'\emph{Idéal des dénominateurs de $u$} l'annulateur $\sh{J}$ de la section $\overline{u}$ image de $u$ dans $\sh{M}_X(\sh{F})/\sh{F}$.
L'Idéal $\sh{J}$ est \emph{quasi-cohérent} : en effet, la question étant locale sur $X$, on peut se borner au cas où $X$ est affine, et il existe une section $s\in\Gamma(X,\sh{S}(\sh{O}_X))$ telle que $v=su\in\Gamma(X,\sh{F})$.
Dire que, pour un ouvert $U\subset X$, une section $f\in\Gamma(U,\sh{O}_X)$ appartient à $\Gamma(U,\sh{J})$ signifie que $f(u|U)\in\Gamma(U,\sh{F})$, et puisque $s|U$ est un élément régulier de $\Gamma(U,\sh{O}_X)$ et que $\sh{F}$ est strictement sans torsion, la relation précédente équivaut encore à $f((su)|U)\in\Gamma(U,s\sh{F})$ ; si $\overline{v}$ est la section de $\sh{F}/s\sh{F}$, image canonique de $v$, on voit donc que $\sh{J}$ est le noyau de l'homomorphisme $\sh{O}_X\to\sh{F}/s\sh{F}$ obtenu par multiplication par la section $\overline{v}$.
Comme $\sh{F}/s\sh{F}$ est quasi-cohérent, il en est de même de $\sh{J}$.

Il résulte aussitôt de la définition précédente que $\operatorname{dom}(u)$ est \emph{l'ouvert complémentaire du sous-préschéma fermé de $X$ défini par l'Idéal des dénominateurs de $u$}.
\end{env}

\begin{proposition}[20.2.15]
\label{IV.20.2.15}
Soient $f:X'\to X$ un morphisme, $\sh{F}$ un $\sh{O}_X$-Module quasi-cohérent, $\vphi$ une section de $\sh{M}_f(\sh{F})$ au-dessus de $X$ \sref{IV.20.1.11}.
Alors $f^{-1}(\operatorname{dom}(\vphi))$ est schématiquement dense dans $X'$.
\end{proposition}

\begin{proof}
La question étant évidemment locale sur $X$ et $X'$, on peut supposer $X=\Spec(A)$, $X'=\Spec(A')$ affines, $\sh{F}=\widetilde{M}$, et $\vphi=h\otimes(1/s)$, où $h\in M$ et $s$ est un élément régulier de $A$ dont l'image $s'$ dans $A'$ est un élément régulier.
On sait alors \sref{IV.20.2.9} que $\mathrm{D}(s')$ est un ouvert schématiquement dense dans $X'$, et c'est l'image réciproque par $f$ de $\mathrm{D}(s)$, qui est contenu dans $\operatorname{dom}(\vphi)$.
\end{proof}

\begin{remark}[20.2.16]
\label{IV.20.2.16}
Lorsque $Y$ n'est pas séparé, le préfaisceau $U\mapsto\operatorname{Ps.hom}_S(U,Y)$ sur $X$ n'est pas nécessairement un faisceau, même lorsque $X$ est noethérien, comme le montre l'exemple suivant.
Nous allons prendre pour $X$ un spectre d'un anneau semi-local noethérien $A$ de dimension $\geq2$, ayant exactement deux idéaux
\oldpage[IV-4]{237}
maximaux $\mathfrak{m}'$, $\mathfrak{m}''$ (donc tel que $X$ ait exactement deux points fermés $x'$, $x''$), de sorte que $\mathfrak{m}'$ et $\mathfrak{m}''$ appartiennent à $\operatorname{Ass}(A)$ et que l'ouvert $U=X-\{x',x''\}$ soit \emph{intègre}.
Soient $U'=X-\{x'\}$, $U''=X-\{x''\}$, de sorte que $U=U'\cap U''$.
Notons que les seuls ouverts schématiquement denses dans $X$ sont ceux qui contiennent $x'$ et $x''$ \sref{IV.20.2.13}[(iv)], donc ils sont nécessairement identiques à $X$.
Pour avoir un contre-exemple, il suffira donc de définir deux $S$-morphismes $u':U'\to Y$, $u'':U''\to Y$ (avec $S=X$) dont les restrictions à $U$ appartiennent au même pseudo-morphisme de $U$ dans $Y$, mais qui soient tels que, pour tout voisinage $V'$ de $x''$ dans $U'$ et tout voisinage $V''$ de $x'$ dans $U''$, les restrictions de $u'$ et $u''$ à $V'\cap V''$ soient distinctes.
Pour cela, considérons une partie fermée irréductible $Z$ de $X$ contenant $x'$ et $x''$, distincte de $X$ ; soit $Y$ le $X$-préschéma obtenu en recollant deux préschémas $Y'$, $Y''$ isomorphes à $X$ le long de l'ouvert partout dense $X-Z$ \sref[I]{I.2.3.1}.
On prendra alors pour $u'$ et $u''$ les restrictions respectives à $U'$ et $U''$ des isomorphismes réciproques des isomorphismes structuraux $Y'\to X$, $Y''\to X$.
Comme $V'$ et $V''$ contiennent le point générique de $Z$, les restrictions de $u'$ et $u''$ à $V'\cap V''$ sont distinctes, mais les restrictions de $u'$ et $u''$ à $U-(U\cap Z)$ sont égales, et $U-(U\cap Z)$ est un ouvert dense dans $U$, donc schématiquement dense puisque $U$ est réduit.

Reste donc seulement à définir $A$ et $Z$ de façon à avoir les propriétés précédentes.
Soient $X_0=\Spec(B)$ un schéma affine intègre (par exemple le plan affine sur un corps $k$), $Y$ une partie fermée irréductible de $X_0$ (par exemple une droite affine) contenant deux points fermés distincts $x'$ et $x''$, correspondant aux idéaux maximaux $\mathfrak{n}'$, $\mathfrak{n}''$ de $B$.
On considère l'anneau $C=B\oplus(B/\mathfrak{n}'\oplus B/\mathfrak{n}'')$ avec la multiplication $(b,z)(b',z')=(bb',bz'+b'z)$.
Si $X_1=\Spec(C)$, on a $X_0=(X_1)_\mathrm{red}$ et $X_1$ est réduit sauf aux points $x'$, $x''$.
Il suffit alors de prendre $A=R^{-1}C$, où $R$ est le complémentaire de la réunion des idéaux maximaux de $C$ aux points $x'$, $x''$, et pour $Z$ la trace de $Y$ sur $X=\Spec(A)$.
\end{remark}

\subsection{Composition des pseudo-morphismes}
\EGAReaderTarget{ega4.c04.subsec.20-3}{IV.20.3}

\begin{env}[20.3.1]
\EGAReaderTarget{ega4.c04.statement.20-3-1}{IV.20.3.1}
Soient $X$, $Y$, $Z$ trois préschémas, $\omega$ un pseudo-morphisme de $X$ dans $Y$, $f:Y\to Z$ un morphisme.
Il est clair que si $U'$, $U''$ sont deux ouverts schématiquement denses dans $X$, $u':U'\to Y$, $u'':U''\to Y$ deux morphismes de la classe $\omega$, les morphismes $f\circ u'$ et $f\circ u''$ sont équivalents (pour la relation définie dans \sref{IV.20.2.1}), donc, pour tous les morphismes $u$ de la classe $\omega$, les morphismes $f\circ u$ appartiennent à une même classe d'équivalence, que l'on note $f\circ\omega$ et que l'on appelle le pseudo-morphisme de $X$ dans $Z$ \emph{composé de $f$ et de $\omega$}.
On a $\operatorname{dom}(f\circ\omega)=\operatorname{dom}(\omega)$.
Si $g:Z\to T$ est un morphisme, il est clair que l'on a $g\circ(f\circ\omega)=(g\circ f)\circ\omega$.
\end{env}

\begin{env}[20.3.2]
\EGAReaderTarget{ega4.c04.statement.20-3-2}{IV.20.3.2}
Supposons maintenant donnés un pseudo-$S$-morphisme $\omega$ de $X$ dans $Y$, où $Y$ est \emph{séparé sur $S$}, de sorte qu'il existe un $S$-morphisme $u:\operatorname{dom}_S(\omega)\to Y$ de la classe $\omega$ \sref{IV.20.2.4}.
Soit $f:X'\to X$ un $S$-morphisme tel que l'ouvert $V'=f^{-1}(\operatorname{dom}_S(\omega))$ soit schématiquement dense dans $X'$ ; alors on dit que la classe (pour la relation d'équivalence de \sref{IV.20.2.1}) du $S$-morphisme $u\circ(f|U')$ (classe qui est définie en vertu de l'hypothèse précédente) est le pseudo-$S$-morphisme \emph{composé de $\omega$ et de $f$}, et on le note $\omega\circ f$ ; il est clair que l'on a
\[
\label{IV.20.3.2.1}
  f^{-1}(\operatorname{dom}_S(\omega))\subset\operatorname{dom}_S(\omega\circ f)
  \tag{20.3.2.1}
\]

Pour $f:X'\to X$ donné, on note $\operatorname{Ps.hom}_S(X,Y)^f$ l'ensemble des pseudo-$S$-morphismes $\omega$ de $X$ dans $Y$ qui vérifient la condition précédente.
Si $\omega$ est un tel pseudo-$S$-morphisme, il est clair que pour tout ouvert $V$ de $X$,
\[
  f^{-1}(\operatorname{dom}_S(\omega|V))=f^{-1}(V\cap\operatorname{dom}_S(\omega))=f^{-1}(V)\cap f^{-1}(\operatorname{dom}_S(\omega))
\]
est schématiquement dense dans $f^{-1}(V)$, donc, si $f^V:f^{-1}(V)\to V$ est la restriction de $f$, le composé $(\omega|V)\circ f^V$ est défini et égal à $(\omega\circ f)|f^{-1}(V)$.
On définit ainsi une
\oldpage[IV-4]{238}
application canonique de restriction $\operatorname{Ps.hom}_S(X,Y)^f\to\operatorname{Ps.hom}_S(V,Y)^{f^V}$, ce qui permet de définir un \emph{préfaisceau} d'ensembles $V\mapsto\operatorname{Ps.hom}_S(V,Y)^{f^V}$ sur $X$, sous-préfaisceau de $V\mapsto\operatorname{Ps.hom}_S(V,Y)$, d'où un faisceau d'ensembles associé que l'on note $\mathscr{P}\!s.\!\operatorname{hom}_S(X,Y)^f$.
En outre, pour tout ouvert $V$ de $X$, on a une application $\omega\mapsto\omega\circ f^V$ de $\operatorname{Ps.hom}_S(V,Y)^{f^V}$ dans $\operatorname{Ps.hom}_S(f^{-1}(V),Y)$, qui définit donc un $f$-morphisme de faisceaux d'ensembles
\[
  \mathscr{P}\!s.\!\operatorname{hom}_S(X,Y)^f\to\mathscr{P}\!s.\!\operatorname{hom}_S(X',Y).
\]
\end{env}

\begin{env}[20.3.3]
\label{IV.20.3.3}
Soit maintenant $f':X''\to X'$ un $S$-morphisme tel que l'ouvert $f'^{-1}(f^{-1}(\operatorname{dom}_S(\omega)))$ soit schématiquement dense dans $X''$ ; alors $\omega\circ(f\circ f')$ est défini et $u\circ f\circ f'$ appartient à ce pseudo-$S$-morphisme ; en outre, en vertu de \sref{IV.20.3.2.1}, $f'^{-1}(\operatorname{dom}_S(\omega\circ f))$ est \emph{a fortiori} schématiquement dense dans $X''$, donc $(\omega\circ f)\circ f'$ est aussi défini et $u\circ f\circ f'$ appartient à ce pseudo-$S$-morphisme, donc on a $(\omega\circ f)\circ f'=\omega\circ(f\circ f')$.
D'autre part, pour tout $S$-morphisme $g:Y\to Z$, on a $\operatorname{dom}_S(g\circ\omega)=\operatorname{dom}_S(\omega)$ \sref{IV.20.3.1}, donc $(g\circ\omega)\circ f$ est défini, et $g\circ u\circ f$ appartient à ce pseudo-$S$-morphisme, ce qui montre que $(g\circ\omega)\circ f=g\circ(\omega\circ f)$.
\end{env}

\begin{env}[20.3.4]
\label{IV.20.3.4}
Le cas le plus important dans la définition \sref{IV.20.3.2} est celui où l'on a $\mathscr{P}\!s.\!\operatorname{hom}_S(X,Y)^f=\mathscr{P}\!s.\!\operatorname{hom}_S(X,Y)$ : il suffit pour cela que pour tout ouvert $U$ de $X$, et tout ouvert $V$ schématiquement dense dans $U$, $f^{-1}(V)$ soit schématiquement dense dans $f^{-1}(U)$ ; lorsqu'il en est ainsi, alors, pour \emph{tout} ouvert $U$ de $X$ et \emph{tout} pseudo-$S$-morphisme $\omega:U\to Y$, on peut définir le composé $\omega\circ f^U$, \emph{même si $Y$ n'est pas séparé sur $S$}.
En effet, si $u':U'\to Y$ et $u'':U''\to Y$ sont deux morphismes de la classe $\omega$, ils coïncident dans un ouvert $U_0$ schématiquement dense dans $U$, donc les morphismes composés $f^{-1}(U')\to U'\xrightarrow{u'}Y$ et $f^{-1}(U'')\to U''\xrightarrow{u''}Y$ coïncident dans $f^{-1}(U_0)$, qui est par hypothèse schématiquement dense dans $f^{-1}(U)$ ; on peut donc définir $\omega\circ f^U$ comme la classe d'un quelconque des morphismes $f^{-1}(U')\to U'\xrightarrow{u'}Y$, où $u'$ appartient à $\omega$.
\end{env}

\begin{proposition}[20.3.5]
\EGAReaderTarget{ega4.c04.prop.20-3-5}{IV.20.3.5}
Soient $X$, $X'$ deux préschémas, $f:X'\to X$ un morphisme.
Alors, dans chacun des trois cas suivants, pour tout ouvert $U$ de $X$ et tout ouvert $V$ schématiquement dense dans $U$, $f^{-1}(V)$ est schématiquement dense dans $f^{-1}(U)$ :
\begin{enumerate}[label=\emph{(\roman*)}]
  \item $f$ est plat et localement de présentation finie.
  \item $f$ est plat et l'espace sous-jacent à $X$ est localement noethérien.
  \item $X'$ est réduit, l'ensemble des composantes irréductibles de $X$ est localement fini, et toute composante irréductible de $X'$ domine une composante irréductible de $X$.
\end{enumerate}
\end{proposition}

\begin{proof}
L'assertion~(i) résulte de \sref{IV.11.10.5}[(ii), (b)] ; l'assertion~(ii) résulte de \sref{IV.11.10.5}[(ii), (a)], car alors tout ouvert $V$ dans $U$ est rétrocompact, autrement dit l'injection canonique $j:V\to U$ est un morphisme quasi-compact.
Enfin, pour prouver~(iii), notons que puisque $X'$ est réduit, il suffit de montrer que pour tout ouvert $U$ de $X$ et tout ouvert $V$ dense dans $U$, $f^{-1}(V)$ est \emph{dense} dans $f^{-1}(U)$.
Or, pour que $f^{-1}(V)$ soit dense dans $f^{-1}(U)$, il suffit que $f^{-1}(V)$ contienne tous les points maximaux de $X'$ appartenant à $f^{-1}(U)$ ; la conclusion résulte donc de l'hypothèse que l'image par $f$ de tout point
\oldpage[IV-4]{239}
maximal de $X'$ appartenant à $f^{-1}(U)$ est un point maximal de $X$ appartenant à $U$, donc à $V$ puisque l'ensemble des composantes irréductibles de $X$ est localement fini.
\end{proof}

\begin{env}[20.3.6]
\EGAReaderTarget{ega4.c04.statement.20-3-6}{IV.20.3.6}
Soient $X$, $Y$ deux $S$-préschémas ; supposons que $X$ vérifie l'une des deux hypothèses suivantes :
\begin{enumerate}[label=\emph{\alph*)}]
  \item $X$ est localement noethérien ;
  \item $X$ est réduit et l'ensemble de ses composantes irréductibles est localement fini.
\end{enumerate}
Alors, pour tout $x\in X$, le $S$-morphisme canonique $j:\Spec(\sh{O}_{X,x})\to X$ est plat et vérifie la condition~(ii) de \sref{IV.20.3.5} dans le cas~a), la condition~(iii) de \sref{IV.20.3.5} dans le cas~b).
Pour tout pseudo-$S$-morphisme $\omega$ de $X$ dans $Y$, le composé $\omega\circ j$ est donc défini et est un pseudo-$S$-morphisme de $\Spec(\sh{O}_{X,x})$ dans $Y$, appelé la \emph{restriction} de $\omega$ à $\Spec(\sh{O}_{X,x})$.
Notons maintenant que si $X$ vérifie la condition~a) (resp. b)) de \sref{IV.20.3.6}, il en est de même de tout préschéma induit sur un ouvert $U$ de $X$ contenant $x$.
Par passage à la limite inductive, on déduit donc, des applications canoniques $\operatorname{Ps.hom}_S(U,Y)\to\operatorname{Ps.hom}_S(\Spec(\sh{O}_{X,x}),Y)$ ainsi obtenues, une application canonique
\[
\label{IV.20.3.6.1}
  (\mathscr{P}\!s.\!\operatorname{hom}_S(X,Y))_x\to\operatorname{Ps.hom}_S(\Spec(\sh{O}_{X,x}),Y)
  \tag{20.3.6.1}
\]
où le premier membre est la fibre au point $x$ du faisceau $\mathscr{P}\!s.\!\operatorname{hom}_S(X,Y)$, ensemble des germes en $x$ de pseudo-$S$-morphismes de voisinages ouverts de $x$ dans $Y$.
\end{env}

\begin{proposition}[20.3.7]
\label{IV.20.3.7}
Sous les hypothèses de \sref{IV.20.3.6}, l'application canonique \sref{IV.20.3.6.1} est injective.
Si en outre $Y$ est localement de présentation finie sur $S$, l'application \sref{IV.20.3.6.1} est bijective.
\end{proposition}

\begin{proof}
Par application de la méthode de \sref{IV.8.1.2}[(a)], cette proposition sera un cas particulier de la suivante :
\end{proof}

\begin{proposition}[20.3.8]
\EGAReaderTarget{ega4.c04.prop.20-3-8}{IV.20.3.8}
Les notations étant celles de \sref{IV.8.8.1}, supposons $X_\alpha$ quasi-compact (resp. quasi-compact et quasi-séparé) et $Y_\alpha$ localement de type fini (resp. localement de présentation finie) sur $S_\alpha$.
Supposons en outre que l'une des conditions suivantes soit vérifiée :
\begin{enumerate}[label=\emph{(\roman*)}]
  \item Les morphismes de transition $S_\mu\to S_\lambda$ (pour $\lambda\leq\mu$) sont plats, et les $X_\lambda$ et $X$ sont noethériens.
  \item Les $X_\lambda$ sont réduits, l'ensemble des composantes irréductibles de $X$ et de chacun des $X_\lambda$ est fini, et, pour $\lambda\leq\mu$, toute composante irréductible de $X_\mu$ domine une composante irréductible de $X_\lambda$.
\end{enumerate}
Alors, l'application canonique
\[
\label{IV.20.3.8.1}
  \varinjlim\operatorname{Ps.hom}_{S_\lambda}(X_\lambda,Y_\lambda)\to\operatorname{Ps.hom}_S(X,Y)
  \tag{20.3.8.1}
\]
est injective (resp. bijective).
\end{proposition}

\begin{proof}
Notons d'abord que, dans le cas~(i), les morphismes $X_\mu\to X_\lambda$ (pour $\lambda\leq\mu$) et $X\to X_\lambda$ sont plats, donc il résulte de \sref{IV.20.3.4} et \sref{IV.20.3.5} que les applications canoniques
\[
  \operatorname{Ps.hom}_{S_\lambda}(X_\lambda,Y_\lambda)\to\operatorname{Ps.hom}_{S_\mu}(X_\mu,Y_\mu)
\]
\oldpage[IV-4]{240}
pour $\lambda\leq\mu$, et $\operatorname{Ps.hom}_{S_\lambda}(X_\lambda,Y_\lambda)\to\operatorname{Ps.hom}_S(X,Y)$ sont définies (sans hypothèse de séparation sur les $Y_\lambda$ ni $Y$) ; il en est par suite de même de l'application \sref{IV.20.3.8.1}.
La proposition va résulter du lemme suivant :

\begin{lemma}[20.3.8.2]
\label{IV.20.3.8.2}
Les notations étant celles de \sref{IV.8.8.1}, supposons $X_\alpha$ quasi-compact, et soit $U_\alpha$ un ouvert dans $X_\alpha$ ; soient $U_\lambda$ et $U$ ses images réciproques dans $X_\lambda$ et $X$ pour $\lambda\geq\alpha$.
Supposons que l'une des conditions~(i),~(ii) de \sref{IV.20.3.8} soit vérifiée.
Alors, pour que $U$ soit schématiquement dense dans $X$, il faut et il suffit qu'il existe $\lambda\geq\alpha$ tel que $U_\lambda$ soit schématiquement dense dans $X_\lambda$, et dans ce cas $U_\mu$ est schématiquement dense dans $X_\mu$ pour $\mu\geq\lambda$.
\end{lemma}

Supposons d'abord vérifiée la condition~(i) ; désignons par $j_\lambda:U_\lambda\to X_\lambda$ et $j:U\to X$ les injections canoniques, par $\sh{J}_\lambda$ le noyau de l'homomorphisme canonique $\sh{O}_{X_\lambda}\to(j_\lambda)_*(\sh{O}_{U_\lambda})$.
L'immersion $j_\lambda$ étant quasi-compacte et quasi-séparée, $(j_\lambda)_*(\sh{O}_{U_\lambda})$ est un $\sh{O}_{X_\lambda}$-Module quasi-cohérent, donc $\sh{J}_\lambda$ est un Idéal quasi-cohérent de $\sh{O}_{X_\lambda}$, et puisque $X_\lambda$ est noethérien, $\sh{J}_\lambda$ est cohérent, donc de type fini.
D'autre part, le morphisme de transition $p_{\lambda\mu}:X_\mu\to X_\lambda$ (resp. $p_\lambda:X\to X_\lambda$) étant plat, il résulte de \sref{IV.2.3.1} que l'on peut identifier $(j_\mu)_*(\sh{O}_{U_\mu})$ à $p_{\lambda\mu}^*((j_\lambda)_*(\sh{O}_{U_\lambda}))$ (resp. $j_*(\sh{O}_U)$ à $p_\lambda^*((j_\lambda)_*(\sh{O}_{U_\lambda}))$).
L'assertion résulte alors de la définition d'un ouvert schématiquement dense et de \sref{IV.8.5.8}[(ii)].

Pour établir \sref{IV.20.3.8.2} lorsque la condition~(ii) est vérifiée, nous prouverons d'abord deux lemmes.

\begin{lemma}[20.3.8.3]
\label{IV.20.3.8.3}
Sous les hypothèses de \sref{IV.8.2.2}, soit $M_\lambda$ (resp. $M$) l'ensemble des points maximaux de $S_\lambda$ (resp. $S$).
Supposons que pour tout $\lambda$, l'ensemble $M_\lambda$ soit fini, et que les $M_\lambda$ forment un système projectif d'ensembles.
Alors $M$ est limite projective du système des $M_\lambda$.
\end{lemma}

Montrons d'abord qu'un point $s\in\varprojlim M_\lambda$ est maximal dans $S$ : en effet, si $s'$ est une généralisation de $s$, l'image $s'_\lambda$ de $s'$ dans $S_\lambda$ est une généralisation de l'image $s_\lambda$ de $s$, donc égale à $s_\lambda$ pour tout $\lambda$, ce qui implique $s'=s$, puisque l'ensemble sous-jacent à $S$ est limite projective des ensembles sous-jacents aux $S_\lambda$ \sref{IV.8.2.9}.
Inversement, soit $s$ un point maximal de $S$ et prouvons qu'il appartient à $\varprojlim M_\lambda$.
Soient $s_\lambda$ l'image de $s$ dans $S_\lambda$, $M'_\lambda$ l'ensemble des points maximaux de $S_\lambda$ qui sont des généralisations de $s_\lambda$ ; les $M'_\lambda$ sont des ensembles finis non vides, qui forment un système projectif, donc $M'=\varprojlim M'_\lambda$ n'est pas vide et l'application $M'\to M'_\lambda$ est surjective (Bourbaki, \emph{Ens.}, chap.~III, 2\textsuperscript{e} éd., \textsection~7, n\textsuperscript{o}~4, Exemple~I).
D'autre part, on a $\Spec(\sh{O}_{S,s})=\varprojlim\Spec(\sh{O}_{S_\lambda,s_\lambda})$ en vertu de \sref{IV.8.2.12} et \sref{IV.8.2.9}, donc les points de $\varprojlim M'_\lambda$ sont aussi points maximaux de $\Spec(\sh{O}_{S,s})$ d'après la première partie du raisonnement.
Donc $M'=\varprojlim M'_\lambda$ se réduit nécessairement au point $s$ ; on en conclut que $M'_\lambda$ se réduit au point $s_\lambda$ et par suite les $s_\lambda$ sont maximaux dans les $S_\lambda$, ce qui achève de prouver le lemme.

\begin{lemma}[20.3.8.4]
\label{IV.20.3.8.4}
Les hypothèses étant celles de \sref{IV.20.3.8.3}, supposons en outre $S_\alpha$ quasi-compact ; soit $V_\alpha$ un ensemble ouvert de $S_\alpha$, et soient $V_\lambda$ et $V$ ses images réciproques dans $S_\lambda$ pour $\lambda\geq\alpha$ et dans $S$.
Si $U_\alpha$ est dense dans $S_\alpha$, $U_\lambda$ est dense dans $S_\lambda$ pour $\lambda\geq\alpha$ et $U$ est dense dans $S$.
Inversement, si $U$ est dense dans $S$ et si de plus l'ensemble $M$ des points maximaux de $S$ est fini, il existe $\lambda\geq\alpha$ tel que $U_\lambda$ soit dense dans $S_\lambda$.
\end{lemma}

En effet, puisque $M_\alpha$ est fini, l'hypothèse que $U_\alpha$ est dense dans $S_\alpha$ entraîne que
\oldpage[IV-4]{241}
$M_\alpha\subset U_\alpha$, donc $M_\lambda\subset U_\lambda$ pour $\lambda\geq\alpha$ et $M\subset U$ en vertu de \sref{IV.20.3.8.3}, ce qui prouve la première assertion.
Inversement, supposons $M$ fini et $U$ dense dans $S$, donc $M\subset U$ ; notons que les $U_\lambda$ sont ouverts, donc ind-constructibles et les $M_\lambda$ finis, donc pro-constructibles \sref{IV.1.9.6}.
La seconde assertion résulte alors de \sref{IV.8.3.2}.

\phantomsection\label{IV.20.3.8.5}
\emph{20.3.8.5. Fin de la démonstration de \sref{IV.20.3.8.2}.}
Supposons la condition~(ii) vérifiée, et notons que $X$ est alors \emph{réduit} en vertu de \sref{IV.8.7.1} ; il revient donc au même de dire que $U_\lambda$ (resp. $U$) est schématiquement dense dans $X_\lambda$ (resp. $X$) ou qu'il est dense dans $X_\lambda$ (resp. $X$), et la conclusion résulte de \sref{IV.20.3.8.4} appliqué aux $X_\lambda$ et à $X$.

\phantomsection\label{IV.20.3.8.6}
\emph{20.3.8.6. Fin de la démonstration de \sref{IV.20.3.8}.}
Pour montrer que l'application \sref{IV.20.3.8.1} est injective, considérons deux morphismes $u'_\lambda$, $u''_\lambda$ d'un ouvert schématiquement dense $U_\lambda\subset X_\lambda$ dans $Y_\lambda$, et supposons que leurs images $u'$, $u''$, morphismes de $U$ dans $Y$, coïncident dans un ouvert schématiquement dense $V$ de $U$.
En outre, dans l'une ou l'autre des hypothèses~(i),~(ii), on peut supposer $V$ \emph{quasi-compact} ; c'est évident si $X$ est noethérien ; sinon, comme $X$ n'a qu'un nombre fini de points maximaux et est réduit, il suffit de remplacer $V$ par la réunion d'un nombre fini de voisinages ouverts affines de ces points maximaux (contenus dans $V$ par hypothèse).
Alors il existe un $\lambda\geq\alpha$ tel que $V$ soit l'image réciproque d'un ouvert quasi-compact $V_\lambda$ de $X_\lambda$ \sref{IV.8.2.11}, et il résulte de \sref{IV.20.3.8.2} qu'en prenant $\lambda$ assez grand, on peut supposer que $V_\lambda$ est schématiquement dense dans $X_\lambda$.
Il résulte alors de \sref{IV.8.8.2}[(i)] qu'en prenant $\lambda$ assez grand, $u'_\lambda$ et $u''_\lambda$ coïncident dans $V_\lambda$, donc appartiennent au même pseudo-homomorphisme.

Prouvons enfin que l'application \sref{IV.20.3.8.1} est surjective.
On considère maintenant un morphisme $u$ d'un ouvert schématiquement dense $U$ de $X$, dans $Y$ ; comme ci-dessus, on peut supposer $U$ quasi-compact et quasi-séparé ($U$ pouvant être remplacé, dans le cas~(ii), par une réunion d'un nombre fini d'ouverts affines deux à deux disjoints).
On peut donc encore supposer qu'il existe un $\lambda\geq\alpha$ tel que $U$ soit l'image réciproque d'un ouvert quasi-compact $U_\lambda$ de $X_\lambda$ \sref{IV.8.2.11} qui est automatiquement quasi-séparé, et par \sref{IV.20.3.8.2} on peut supposer en outre que $U_\lambda$ soit schématiquement dense dans $X_\lambda$.
Puisque les $Y_\lambda$ sont supposés localement de présentation finie, il résulte de \sref{IV.8.8.2}[(i)] qu'il existe un $\lambda$ tel que $u$ soit l'image d'un morphisme $u_\lambda:U_\lambda\to Y_\lambda$ ; d'où la conclusion.

\begin{remarks}[20.3.8.7]
\label{IV.20.3.8.7}
\begin{enumerate}[label=(\roman*)]
  \item Pour prouver que l'application \sref{IV.20.3.8.1} est \emph{injective}, il n'est pas nécessaire, dans l'hypothèse~(i) de \sref{IV.20.3.8}, de supposer $X$ noethérien.
  En effet, le lemme \sref{IV.20.3.8.2} n'utilise pas cette hypothèse.
  Avec les notations de \sref{IV.20.3.8.6}, soit $Z_\lambda$ le sous-préschéma des coïncidences de $u'_\lambda$ et $u''_\lambda$, et soit $Z$ le sous-préschéma des coïncidences de $u'$ et $u''$ ; il résulte de la définition \sref{IV.17.4.5} et de \sref[I]{I.3.3.10.1} que $Z$ est limite projective des $Z_\mu$ pour $\mu\geq\lambda$.
  Or, par hypothèse, $Z$ majore un ouvert schématiquement dense dans $X$ ; il en résulte que $Z$ est lui-même induit sur un ouvert de $X$ en vertu du lemme suivant :

\begin{lemma}[20.3.8.8]
\label{IV.20.3.8.8}
Soit $T$ un préschéma.
Alors tout sous-préschéma $W$ de $T$ qui majore un ouvert schématiquement dense $V$ de $T$ est induit sur un ouvert (schématiquement dense) de $T$.
\end{lemma}

  En effet, le sous-espace de $T$ sous-jacent à $W$ est localement fermé, donc ouvert
  \oldpage[IV-4]{242}
  dans son adhérence, ce qui prouve déjà que l'espace sous-jacent à $W$ est ouvert dans $T$ ; la conclusion résulte alors de \sref{IV.11.10.1}[(c)].

  Ce lemme étant établi, on en conclut que pour $\mu\geq\lambda$ assez grand, $Z_\mu$ est induit sur un ouvert de $X_\mu$ en vertu de \sref{IV.8.6.3}, car $Z_\lambda$, en tant que sous-préschéma d'un préschéma noethérien, est de présentation finie sur $X_\lambda$ \sref{IV.1.6.2}, et il en est donc de même des $Z_\mu$ pour $\mu\geq\lambda$ sur $X_\mu$ et de $Z$ sur $X$.
  On peut maintenant appliquer \sref{IV.20.3.8.2} qui montre que pour $\mu\geq\lambda$ assez grand, $Z_\mu$ est schématiquement dense dans $X_\mu$, d'où la conclusion.

  \item Si, dans l'hypothèse~(i) de \sref{IV.20.3.8}, on supprime la condition que $X$ est noethérien, on voit que le raisonnement de \sref{IV.20.3.8.6} montre encore que l'image de \sref{IV.20.3.8.1} est formée des pseudo-$S$-morphismes ayant un représentant qui est un $S$-morphisme $U\to Y$, où $U$ est schématiquement dense dans $X$ et \emph{quasi-compact et quasi-séparé}.
\end{enumerate}
\end{remarks}
\end{proof}

\begin{corollary}[20.3.9]
\label{IV.20.3.9}
Supposons vérifiées l'une ou l'autre hypothèse~a),~b) de \sref{IV.20.3.6} sur $X$, et que $Y$ soit localement de présentation finie sur $S$.
Alors, pour qu'un pseudo-$S$-morphisme $\omega$ de $X$ dans $Y$ soit défini au point $x$ \sref{IV.20.2.3}, il faut et il suffit que sa restriction à $\Spec(\sh{O}_{X,x})$ soit partout définie (autrement dit, soit un $S$-morphisme de $\Spec(\sh{O}_{X,x})$ dans $Y$).
\end{corollary}

Le résultat suivant, qui nous servira dans la démonstration de \sref{IV.20.3.11}, utilise la théorie de la descente fidèlement plate du chap.~VI.
Le lecteur pourra vérifier que les résultats de \sref{IV.20.3} ne seront pas utilisés dans cette théorie.

\begin{lemma}[20.3.10]
\label{IV.20.3.10}
Soient $f:X'\to X$ un $S$-morphisme fidèlement plat et quasi-compact, $X''=X'\times_X X'$, $p_1$ et $p_2$ les projections canoniques de $X''$ dans $X'$, $Y$ un préschéma séparé sur $S$.
Soient $U$ un ouvert de $X$, $U'=f^{-1}(U)$, $U''=p_1^{-1}(U')=p_2^{-1}(U')$, et supposons que $U''$ soit schématiquement dense dans $X''$.
Soit $g:U\to Y$ un $S$-morphisme ; alors, si $g\circ(f|U')$ se prolonge en un $S$-morphisme $X'\to Y$, $g$ se prolonge en un $S$-morphisme $X\to Y$.
\end{lemma}

On notera que les hypothèses entraînent que $U$ (resp. $U'$) est schématiquement dense dans $X$ (resp. $X'$) \sref{IV.11.10.5}[(i)] ; on peut donc encore dire que si $\omega$ désigne le pseudo-$S$-morphisme classe de $g$, l'énoncé de \sref{IV.20.3.10} signifie que \emph{si $\omega\circ f$ est partout défini, il en est de même de $\omega$}.

\begin{proof}
Pour prouver \sref{IV.20.3.10}, désignons par $g'$ un $S$-morphisme $X'\to Y$ qui prolonge $g\circ(f|U')$, et posons $g''_i=g'\circ p_i:X''\to Y$ ($i=1,2$).
Si l'on pose $f''=f\circ p_1=f\circ p_2:X''\to X$, il est clair que $g''_1$ et $g''_2$ coïncident dans $U''$ avec $g\circ(f''|U'')$.
Mais comme $Y$ est séparé sur $S$ et $U''$ schématiquement dense dans $X''$, on a $g''_1=g''_2$ \sref{IV.11.10.1}[(d)].
Comme $f$ est fidèlement plat et quasi-compact, il résulte de la théorie de la descente (chap.~VI) qu'il existe un unique $S$-morphisme $h:X\to Y$ tel que $h\circ f=g'$ ; comme la restriction $U'\to U$ de $f$ est un morphisme fidèlement plat et quasi-compact et que $U''=U'\times_U U'$, le résultat d'unicité précédent, appliqué à $U$ au lieu de $X$, montre que $h|U=g$, ce qui prouve le lemme.
\end{proof}

\begin{proposition}[20.3.11]
\label{IV.20.3.11}
Soient $Y$ un $S$-préschéma séparé sur $S$, $\omega$ un pseudo-$S$-morphisme de $X$ dans $Y$, $f:X'\to X$ un $S$-morphisme.
Supposons que $f$ soit plat, et que l'une des conditions suivantes soit réalisée :
\begin{enumerate}[label=\emph{(\roman*)}]
\oldpage[IV-4]{243}
  \item $f$ est un morphisme ouvert, ou surjectif et quasi-compact, et $\operatorname{dom}_S(\omega)$ contient un ouvert $U$ schématiquement dense dans $X$ et rétrocompact dans $X$.
  \item $f$ est localement de présentation finie.
  \item $Y$ est localement de présentation finie sur $S$, et $X$ vérifie l'une des conditions~a),~b) de \sref{IV.20.3.6}.
\end{enumerate}
Alors, $f^{-1}(\operatorname{dom}_S(\omega))$ est schématiquement dense dans $X'$, de sorte que $\omega\circ f$ est défini \sref{IV.20.3.2} et l'on a
\[
\label{IV.20.3.11.1}
  \operatorname{dom}_S(\omega\circ f)=f^{-1}(\operatorname{dom}_S(\omega)).
  \tag{20.3.11.1}
\]
\end{proposition}

\begin{proof}
Prouvons d'abord que $f^{-1}(\operatorname{dom}_S(\omega))$ est schématiquement dense dans $X'$.
La question étant locale sur $X$ et $X'$, on peut supposer $X$ et $X'$ affines et puisque $f$ est plat, il suffit de voir, en vertu de \sref{IV.11.10.5}[(ii), (a)] que $\operatorname{dom}_S(\omega)$ contient un ensemble ouvert $U$ rétrocompact et schématiquement dense dans $X$.
Cela résulte de l'hypothèse dans le cas~(i), et de \sref{IV.20.2.12} dans le cas~(iii), compte tenu du fait que dans un schéma affine, tout ouvert de la forme $\mathrm{D}(t)$ est rétrocompact ; enfin, dans le cas~(ii), on voit directement que $f^{-1}(\operatorname{dom}_S(\omega))$ est schématiquement dense dans $X'$ en appliquant \sref{IV.20.3.5}[(i)].

Prouvons maintenant \sref{IV.20.3.11.1}, autrement dit que, pour tout point $x'\in\operatorname{dom}_S(\omega\circ f)$, on a $x=f(x')\in\operatorname{dom}_S(\omega)$.
Notons d'abord qu'on peut se borner au cas où $f$ est \emph{fidèlement plat et quasi-compact}.
C'est en effet l'hypothèse dans le second cas de~(i) ; dans les autres cas, la question est locale sur $X'$, et on peut donc supposer déjà $X$ et $X'$ affines, donc $f$ quasi-compact.
Dans le premier cas de~(i) et dans le cas~(ii), $f$ est un morphisme ouvert \sref{IV.11.3.1}, donc on peut, en remplaçant $X$ par l'ouvert $f(X')$, supposer $f$ surjectif, donc fidèlement plat.
Dans le cas~(iii), en utilisant \sref{IV.20.3.9}, on peut se borner à prouver que la restriction de $\omega$ à $\Spec(\sh{O}_{X,x})$ est partout définie, et on peut donc remplacer $X$ par $\Spec(\sh{O}_{X,x})$, $X'$ par $X'\times_X\Spec(\sh{O}_{X,x})$ et $f$ par sa restriction à ce dernier préschéma, qui est un morphisme surjectif \sref{IV.2.3.4}, donc fidèlement plat.

Supposons donc $f$ fidèlement plat et quasi-compact ; avec les notations du lemme \sref{IV.20.3.10}, il suffit de voir alors que $U''$ est schématiquement dense dans $X''$, en prenant pour $U$ un ouvert schématiquement dense dans $X$, contenu dans $\operatorname{dom}_S(\omega)$ ; cela sera le cas, en vertu de \sref{IV.11.10.5}[(ii), (a)] si $U$ est pris rétrocompact dans $X$ (puisque le morphisme $f'':X''\to X$ est plat).
Or l'existence d'un tel ouvert $U$ fait partie de l'hypothèse dans le cas~(i) ; dans le cas~(iii) elle résulte de \sref{IV.20.2.12} et du fait que dans un schéma affine $\Spec(A)$, tout ouvert de la forme $\mathrm{D}(t)$ est rétrocompact.
Enfin, dans le cas~(ii), prenons $U=\operatorname{dom}_S(\omega)$ et montrons directement que $U''$ est schématiquement dense dans $X''$ : il suffit pour cela de remarquer que $f'':X''\to X$ est plat et localement de présentation finie et d'appliquer \sref{IV.11.10.5}[(ii), (b)].
\end{proof}

\begin{corollary}[20.3.12]
\label{IV.20.3.12}
Soit $\vphi$ une pseudo-fonction sur un préschéma $X$.
Alors, pour tout morphisme plat et localement de présentation finie $f:X'\to X$, la pseudo-fonction $\vphi\circ f$ est définie et l'on a $\operatorname{dom}(\vphi\circ f)=f^{-1}(\operatorname{dom}(\vphi))$.
\end{corollary}

\begin{remark}[20.3.13]
\label{IV.20.3.13}
\oldpage[IV-4]{244}
Lorsque $X$ vérifie l'une des conditions~a),~b) de \sref{IV.20.3.6}, alors on a vu \sref{IV.20.2.11} que les pseudo-fonctions sur $X$ s'identifient aux fonctions méromorphes sur $X$.
En vertu de \sref{IV.20.1.12} et de \sref{IV.20.2.13}[(i)], si l'on suppose seulement que le morphisme $f:X'\to X$ est \emph{plat}, alors, pour toute pseudo-fonction $\vphi$ sur $X$, $\vphi\circ f$ est définie et l'on a $\operatorname{dom}(\vphi\circ f)=f^{-1}(\operatorname{dom}(\vphi))$.
\end{remark}

\subsection{Propriétés des domaines de définition des fonctions rationnelles}
\EGAReaderTarget{ega4.c04.subsec.20-4}{IV.20.4}

\begin{env}[20.4.1]
\EGAReaderTarget{ega4.c04.statement.20-4-1}{IV.20.4.1}
Soient $X$, $Y$ deux $S$-préschémas, $\omega$ un pseudo-$S$-morphisme de $X$ dans $Y$.
Soit $u$ un $S$-morphisme $U\to Y$ appartenant à $\omega$, où $U$ est schématiquement dense dans $X$, et considérons le \emph{graphe} $\Gamma_u$, sous-préschéma de $U\times_S Y$ \sref[I]{I.5.3.11}, donc sous-préschéma de $X\times_S Y$.
Supposons que ce sous-préschéma admette une \emph{adhérence} \sref[I]{I.9.5.11} dans $X\times_S Y$ ; si $j:\Gamma_u\to X\times_S Y$ est l'injection canonique, ce sera le cas lorsque le $\sh{O}_{X\times_S Y}$-Module $j_*(\sh{O}_{\Gamma_u})$ est quasi-cohérent ; il résulte de la définition de la classe d'équivalence $\omega$ \sref{IV.20.2.1} que $j_*(\sh{O}_{\Gamma_u})$ ne dépend pas du représentant $u$ considéré, donc le préschéma adhérence de $\Gamma_u$ est alors bien déterminé par $\omega$ ; on dit que ce sous-préschéma fermé de $X\times_S Y$ est le \emph{graphe} du pseudo-$S$-morphisme $\omega$ et on le note $\Gamma_\omega$.
On notera que $\Gamma_\omega$ est défini s'il existe un morphisme $u:U\to Y$ de la classe $\omega$ tel que $U$ soit \emph{rétrocompact} dans $X$ et si de plus $Y$ est quasi-séparé sur $S$, car alors l'injection $j$ est un morphisme quasi-compact et quasi-séparé (\sref[I]{I.2.2} et \sref[I]{I.7.4}) ; la première hypothèse sera toujours vérifiée lorsque $X$ est \emph{localement noethérien}.
Notons aussi que lorsque $X$ est \emph{réduit}, il en est de même de $\Gamma_u$, qui est isomorphe à $U$ \sref[I]{I.5.3.11} ; alors $\Gamma_\omega$ existe et n'est autre que le sous-préschéma \emph{réduit} de $X\times_S Y$ ayant pour espace sous-jacent l'adhérence dans $X\times_S Y$ de l'espace sous-jacent à $\Gamma_u$ (\sref[I]{I.5.2.1} et \sref[I]{I.5.2.2}).
Notons enfin que si $Y$ est \emph{séparé} sur $S$, $\Gamma_u$ est \emph{fermé} dans $U\times_S Y$ \sref[I]{I.5.4.3}, donc est \emph{induit par} $\Gamma_\omega$ (lorsque ce dernier existe) sur l'ouvert $\Gamma_\omega\cap(U\times_S Y)$ \sref[I]{I.9.5.10} ; par contre, ce préschéma induit est en général différent de $\Gamma_u$ lorsque $Y$ n'est pas séparé.
En particulier, si $v:X\to Y$ est un $S$-morphisme, le graphe $\Gamma_\omega$ de la classe $\omega$ de $v$ peut être distinct du graphe $\Gamma_v$.
Aussi allons-nous dans ce qui suit, lorsqu'il sera question du graphe d'un pseudo-$S$-morphisme, nous borner au cas où \emph{$Y$ est séparé sur $S$}.
\end{env}

\begin{env}[20.4.2]
\label{IV.20.4.2}
Supposons donc $\Gamma_\omega$ défini et $Y$ séparé sur $S$ ; notons $p$ et $q$ les restrictions à $\Gamma_\omega$ des projections canoniques
\[
  \xymatrix{
  & \Gamma_\omega \ar[dl]_p \ar[dr]^q & \\
  X & & Y
  }
\]
Alors, si $U\subset\operatorname{dom}_S(\omega)$, la restriction $p^{-1}(U)\to U$ de $p$ est un isomorphisme \sref[I]{I.5.3.11} ; réciproquement, si $U$ est un ouvert de $X$ ayant cette propriété, et si $u$ est l'isomorphisme réciproque de la restriction $p^{-1}(U)\to U$ de $p$, $q\circ u$ est un $S$-morphisme de $U$ dans $Y$ qui coïncide avec un morphisme de la classe $\omega$ dans $U\cap\operatorname{dom}_S(\omega)$.
On en conclut que $\operatorname{dom}_S(\omega)$ est le \emph{plus grand ouvert} $U$ de $X$ tel que la restriction $p^{-1}(U)\to U$ de $p$ soit un
\oldpage[IV-4]{245}
isomorphisme.
Soit $v:\operatorname{dom}_S(\omega)\to\Gamma_\omega$ le $S$-morphisme réciproque de la restriction $p^{-1}(\operatorname{dom}_S(\omega))\to\operatorname{dom}_S(\omega)$ de $p$ ; on note parfois $p^{-1}$ le pseudo-$S$-morphisme de $X$ dans $\Gamma_\omega$, classe de $v$ (dont l'application rationnelle associée \sref{IV.20.2.13}[(ii)] est alors \emph{birationnelle}) ; comme $p^{-1}(\operatorname{dom}_S(\omega))$ est le graphe du $S$-morphisme $u=q\circ v:\operatorname{dom}_S(\omega)\to Y$, il est schématiquement dense dans $\Gamma_\omega$ \sref{IV.11.10.3}[(iv)], donc $\omega$ peut être considéré comme le \emph{composé} $q\circ p^{-1}$ \sref{IV.20.3.2}.
Pour toute partie $M$ de l'espace sous-jacent à $X$, on pose parfois $\omega(M)=q(p^{-1}(M))$, ce qui revient alors à considérer $\omega$ comme une application de $X$ dans $\mathfrak{P}(Y)$ (ou, comme disent certains auteurs, une « fonction multiforme »).
On notera que lorsque $x\in\operatorname{dom}_S(\omega)$, $\omega(\{x\})$ est l'ensemble $\{u(x)\}$ ; en général, pour un $x\in X$ quelconque, si l'on désigne par $s$ l'image de $x$ dans $S$, par $Y_s$ la fibre au point $s$ du morphisme structural $Y\to S$, l'ensemble $\omega(\{x\})$ est une partie du préschéma $Y_s$.
\end{env}

\begin{env}[20.4.3]
\label{IV.20.4.3}
Dans tout le reste de ce numéro, nous allons nous borner au cas où $X$ est \emph{réduit}, de sorte qu'il y a alors identité entre pseudo-$S$-morphismes et $S$-\emph{applications rationnelles} \sref{IV.20.2.7}.
En outre, le \emph{graphe} de toute $S$-application rationnelle de $X$ dans $Y$ est alors défini \sref{IV.20.4.1}.
\end{env}

\begin{proposition}[20.4.4]
\label{IV.20.4.4}
Soient $X$ un $S$-préschéma localement intègre, $Y$ un $S$-préschéma localement de type fini et séparé sur $S$, $\omega$ une $S$-application rationnelle de $X$ dans $Y$, $p:\Gamma_\omega\to X$ la projection canonique.
Alors, si $x\in X$ est un point normal tel que l'ensemble $p^{-1}(x)$ contienne un point isolé, $\omega$ est définie au point $x$.
\end{proposition}

\begin{proof}
En effet, la première projection $p_1:X\times_S Y\to X$ est alors un morphisme séparé et localement de type fini, donc il en est de même de sa restriction $p:\Gamma_\omega\to X$, qui est en outre \emph{birationnel} et comme $\Gamma_\omega$ est réduit et $X$ intègre, $\Gamma_\omega$ est \emph{intègre} ; il résulte donc de (\textbf{Err}\textsubscript{\textbf{IV}}, 30) que l'hypothèse sur $x$ entraîne l'existence d'un voisinage ouvert $U$ de $x$ tel que la restriction $p^{-1}(U)\to U$ de $p$ soit un isomorphisme ; d'où la conclusion \sref{IV.20.4.2}.

On notera que l'énoncé \sref{IV.20.4.4} est la formulation originelle donnée par Zariski à son « Main theorem » (à cela près qu'il se restreignait aux schémas algébriques sur un corps de base).
\end{proof}

\begin{proposition}[20.4.5]
\label{IV.20.4.5}
Les hypothèses et notations étant celles de \sref{IV.20.4.4}, supposons en outre $X$ normal, et soient $X'$ un préschéma réduit, $f:X'\to X$ un morphisme localement de type fini et universellement ouvert.
Alors $\omega\circ f$ est défini, et l'on a $\operatorname{dom}_S(\omega\circ f)=f^{-1}(\operatorname{dom}_S(\omega))$ (en d'autres termes, si $x'\in X'$ et $x=f(x')$, alors, pour que $\omega$ soit définie au point $x$, il faut et il suffit que $\omega\circ f$ le soit au point $x'$).
\end{proposition}

\begin{proof}
Comme $X'$ est réduit et $f^{-1}(\operatorname{dom}_S(\omega))$ partout dense dans $X'$ en vertu de \sref{IV.2.4.11}, le composé $\omega\circ f$ est défini ; pour prouver que, lorsque $\omega\circ f$ est définie au point $x'$, $\omega$ l'est au point $x$, on peut évidemment remplacer $X'$ par un voisinage ouvert de $x'$, donc supposer $\omega\circ f$ \emph{partout définie} ; en outre, comme $f(X')$ est ouvert dans $X$, on peut supposer $f$ \emph{surjectif}.
En vertu de l'hypothèse sur $f$, le morphisme $f_{(Y)}:X'\times_S Y\to X\times_S Y$ est alors ouvert, donc $\overline{f_{(Y)}^{-1}(M)}=f_{(Y)}^{-1}(\overline{M})$ pour toute partie $M$ de $X\times_S Y$ ; prenant pour $M$ l'ensemble $\Gamma_u$, où $u:\operatorname{dom}_S(\omega)\to Y$ est la restriction de $\omega$ à $\operatorname{dom}_S(\omega)$, il résulte de la relation précédente et de \sref[I]{I.3.3.12} que l'ensemble sous-jacent à $\Gamma_{\omega\circ f}$ est égal à $f_{(Y)}^{-1}(\Gamma_\omega)$ ; comme on sait
\oldpage[IV-4]{246}
que $\Gamma_{\omega\circ f}$ est un préschéma réduit \sref{IV.20.4.1}, on voit que le $X'$-préschéma $\Gamma_{\omega\circ f}$ est égal à $(\Gamma_\omega\times_X X')_\mathrm{red}$.
Mais comme par hypothèse le morphisme composé $\Gamma_{\omega\circ f}\to\Gamma_\omega\times_X X'\xrightarrow{p^{(X')}}X'$ est un isomorphisme, $p^{(X')}$ est nécessairement \emph{radiciel} ; comme $f$ est surjectif, il en est de même de $p:\Gamma_\omega\to X$ \sref[I]{I.3.4.8}, donc $p^{-1}(x)$ est un ensemble réduit à un point \sref[I]{I.3.6.4} ; il résulte alors de \sref{IV.20.4.4} que $\omega$ est définie au point $x$.
\end{proof}

La proposition suivante donne un \emph{critère valuatif} pour qu'une application rationnelle soit définie en un point :

\begin{proposition}[20.4.6]
\label{IV.20.4.6}
Soient $S$ un préschéma, $X$, $Y$ deux $S$-préschémas ; on suppose $X$ localement noethérien, $Y$ localement de type fini sur $S$.
Soient $U$ un ouvert dense dans $X$, $h:U\to Y$ un $S$-morphisme, $x\in X-U$ un point normal de $X$, $h'_x:\Spec(k(x))\to Y$ un $S$-morphisme.
Afin que $h$ puisse être prolongé en un $S$-morphisme $h':U'\to Y$, où $U'$ est un ouvert de $X$ contenant $U$ et $x$, et où le morphisme composé $\Spec(k(x))\to U'\xrightarrow{h'}Y$ est le $S$-morphisme donné $h'_x$, il faut et il suffit que la condition suivante soit vérifiée :
\begin{quote}
\emph{(P) Pour tout spectre $S_1$ d'un anneau de valuation discrète, de point fermé $a$ et de point générique $b$, et tout morphisme $u:S_1\to X$ tel que $u(a)=x$, $u(b)\in U$, le morphisme composé $h_1=h\circ(u|\{b\}):\{b\}=u^{-1}(U)\to Y$, se prolonge en un morphisme $h'_1:S_1\to Y$ tel que le diagramme}
\[
  \xymatrix{
  \Spec(k(a)) \ar[r] & S_1 \ar[d]^{h'_1} \\
  \Spec(k(x)) \ar[r]_{h'_x} & Y
  }
\]
\emph{soit commutatif.}
\end{quote}
De plus, si cette condition est vérifiée, et si $h'':U''\to Y$ est un morphisme vérifiant les mêmes conditions que $h'$, alors il existe un ensemble ouvert contenant $U\cup\{x\}$, et dans lequel $h'$ et $h''$ coïncident.
\end{proposition}

\begin{proof}
Prouvons d'abord la dernière assertion ; on peut supposer $h'$ et $h''$ définis dans le même ouvert $U_0\supset U\cup\{x\}$.
Le sous-préschéma $Z$ des coïncidences de $h'$ et $h''$ \sref{IV.17.4.5} contient $U$ et $x$, donc il y a un voisinage ouvert $V$ de $x$ dans $U_0$ tel que le sous-préschéma induit par $Z$ sur l'ouvert $Z\cap V$ soit un sous-préschéma \emph{fermé} du préschéma induit par $X$ sur $V$ ; comme ce préschéma $Z\cap V$ majore le sous-préschéma induit par $V$ sur l'ouvert $U\cap V$, et que ce dernier est schématiquement dense dans $V$, $Z\cap V$ est nécessairement égal à $V$ \sref{IV.20.3.8.8}, ce qui prouve que $h'$ et $h''$ coïncident dans $U\cup V$.

La nécessité de la condition de l'énoncé étant évidente, prouvons qu'elle est suffisante.
En vertu de la correspondance biunivoque entre $S$-morphismes de $X$ dans $Y$ et $X$-sections de $X\times_S Y$ \sref[I]{I.3.3.15}, on peut se borner au cas où $S=X$ et où $h$ est donc une $U$-section de $Y$ ; on notera qu'alors $Y$ est localement noethérien, et on peut évidemment (puisque $X$ est localement intègre et localement noethérien) supposer $X=S$ irréductible.
En outre, compte tenu de \sref{IV.20.3.7}, on peut supposer que $X=\Spec(A)$, où $A$ est un anneau \emph{local} noethérien, intègre et intégralement clos.

\oldpage[IV-4]{247}
Notons que pour tout $x'\in U$, $h(x')$ est une spécialisation de $h'_x(x)=y$.
En effet, il existe un spectre $S_1$ d'anneau de valuation discrète et un morphisme $u:S_1\to X$ tel que $u(a)=x$, $u(b)=x'$ \sref[II]{II.7.1.9}.
Appliquant la condition de l'énoncé, on obtient aussitôt notre assertion, puisque $h(x')=h_1(b)=h'_1(b)$ et $y=h'_1(a)$.
Si $Y'$ est un voisinage ouvert affine de $y$, on a donc $h(X\cap U)\subset Y'$ ; on peut par suite remplacer $Y$ par $Y'$, autrement dit supposer $Y$ \emph{affine}, donc séparé sur $X$.
Soit $\omega$ la $X$-section rationnelle de $Y$ à laquelle $h$ appartient, de sorte que son graphe a ici pour ensemble sous-jacent l'adhérence de $h(U)$ dans $Y$.
Puisque $Y$ est séparé sur $X$, on peut appliquer \sref{IV.20.4.4} à $\omega$ : il suffira de prouver que, si $p:\Gamma_\omega\to X$ est la projection canonique, $p^{-1}(x)$ est réduit à un seul point $y$ et que $h'_x(x)=y$.
En effet, par \sref{IV.20.4.4} $h$ se prolongera en une section $h'$ au-dessus d'un ouvert $U'$ contenant $U\cup\{x\}$, telle que $h'(x)=y$, et comme alors il n'existe qu'un seul $X$-morphisme $\Spec(k(x))\to Y$ transformant $x$ en $y$, cela prouvera l'identité de $h'_x$ et du composé de $h'$ et de $\Spec(k(x))\to U'$.

Puisque pour $x'\in X\cap U$, $h(x')$ est une spécialisation de $y$, on a $y\in p^{-1}(x)$.
Supposons alors que $z$ soit un point quelconque de $p^{-1}(x)$.
Comme $\Gamma_\omega$ est l'adhérence de $h(U)$ et que $h(U)$ est formé de points adhérents à $h(\xi)$, où $\xi$ est le point générique de $X$, $\Gamma_\omega$ est l'adhérence de $h(\xi)$ dans $Y$.
On prend alors un spectre $S_1$ d'anneau de valuation discrète et un morphisme $v:S_1\to Y$ tels que $v(a)=z$, $v(b)=h(\xi)$ \sref[II]{II.7.1.9}, et on pose $u=p\circ v$, de sorte que $u(a)=x$, $u(b)=\xi$.
Appliquant à $u$ la condition de l'énoncé, on voit que l'on obtient un morphisme $h'_1:S_1\to Y$ tel que $h'_1(a)=y$ et $h'_1(b)=h(\xi)$ ; mais cela entraîne $z=y$ en vertu de \sref[II]{II.7.2.3}, puisque $Y$ est séparé sur $X$ et que $v$ et $h'_1$ doivent donc coïncider, puisqu'ils sont égaux au point $b$.
C.Q.F.D.
\end{proof}

\begin{corollary}[20.4.7]
\EGAReaderTarget{ega4.c04.cor.20-4-7}{IV.20.4.7}
Les hypothèses sur $S$, $X$, $Y$, $U$ et $h$ étant celles de \sref{IV.20.4.6}, soit $E$ une partie de $X-U$ telle que $X$ soit normal en tout point de $E$, et pour chaque $x\in E$, soit $h'_x:\Spec(k(x))\to Y$ un $S$-morphisme, tel que la condition~(P) de \sref{IV.20.4.6} soit vérifiée.
On suppose en outre que la réunion $F$ des graphes des $h'_x$ (identifiés à des parties de $X\times_S Y$) pour $x\in E$, soit contenue dans la réunion d'un nombre fini d'ouverts $V_i$ de $X\times_S Y$ qui sont \emph{séparés} sur $X$ (condition automatiquement vérifiée si $Y$ est \emph{séparé} sur $S$, ou si $X$ est \emph{quasi-compact} et $Y$ de \emph{type fini} sur $S$).
Alors il existe un $S$-morphisme $h':U'\to Y$, où $U'$ est un ouvert de $X$ contenant $U\cup E$, tel que, pour tout $x\in E$, le composé
\[
  \Spec(k(x))\to U'\xrightarrow{h'}Y
\]
soit égal à $h'_x$.
\end{corollary}

\begin{proof}
Notons d'abord que, dans \sref{IV.20.4.6}, si $Y$ est supposé \emph{séparé}, il y a un plus grand ouvert $U_0\supset U$, égal au domaine de la $S$-application rationnelle correspondant à $h$, dans lequel $h$ peut être prolongé, et ce prolongement est unique \sref[I]{I.7.2.2} ; d'où la conclusion dans ce cas, grâce à \sref{IV.20.4.6}.
Dans le cas général, soit $E_i$ l'ensemble des $x\in E$ tels que $(x,h'_x(x))\in V_i$.
En vertu de \sref{IV.20.4.6}, pour tout $x\in E$, il y a un prolongement $h^{(x)}$ de $h$ à $U\cup W^{(x)}$, où $W^{(x)}$ est un voisinage de $x$ dans $X$ tel que le graphe de $h^{(x)}|W^{(x)}$ soit contenu dans \emph{tous} les $V_i$ tels que $x\in E_i$.
Puisque $V_i$ est séparé sur $X$ et $X$ réduit, pour deux points $x'$, $x''$ de $E_i$, $h^{(x')}$ et $h^{(x'')}$ coïncident dans $W^{(x')}\cap W^{(x'')}$ puisqu'ils
\oldpage[IV-4]{248}
coïncident dans l'ouvert partout dense $U\cap W^{(x')}\cap W^{(x'')}$ de cet ensemble.
Il y a donc un morphisme $h_i:U\cup U_i\to Y$ qui prolonge $h$ à un ouvert $U\cup U_i$, contenant $E_i$ ; en outre, pour tout couple d'indices $i$, $j$, les graphes des restrictions $h_i|(U_i\cap U_j)$ et $h_j|(U_i\cap U_j)$ sont contenus dans $V_i\cap V_j$ ; comme $V_i\cap V_j$ est séparé sur $X$ et que les morphismes précédents coïncident dans un ouvert partout dense $U\cap U_i\cap U_j$ de $U_i\cap U_j$, ils sont égaux.
Le morphisme $h'$ égal à $h$ dans $U$, à $h_i$ dans chacun des $U_i$, répond à la question.
\end{proof}

\begin{remark}[20.4.8]
\label{IV.20.4.8}
Lorsque $E=X-U$, il est clair que si, pour tout ouvert affine $T$ de $X$, il existe un $S$-morphisme $h'_T:T\to Y$ prolongeant $h|(U\cap T)$ et tel que le composé $\Spec(k(x))\to T\xrightarrow{h'_T}Y$ soit égal à $h'_x$ pour tout $x\in T-(U\cap T)$, alors les $h'_T$ sont les restrictions d'un $S$-morphisme $h':X\to Y$ (partout défini) en vertu de l'assertion d'unicité dans \sref{IV.20.4.6}.
Pour prouver l'existence de $h'$, on est donc ramené au cas où $X$ est \emph{affine}, et alors il suffit que l'ensemble $F$, réunion des graphes des $h'_x$, soit \emph{quasi-compact} dans $X\times_S Y$ pour que l'on puisse appliquer \sref{IV.20.4.7}.
Il en sera ainsi lorsque les $h'_x$ sont de la forme $\Spec(k(x))\to Z\xrightarrow{h_0}Y$, où $Z$ est un sous-préschéma fermé de $X$ ayant $X-U$ pour espace sous-jacent, et $h_0$ un $S$-morphisme.
\end{remark}

\begin{corollary}[20.4.9]
\label{IV.20.4.9}
Soient $S$ un préschéma localement noethérien, $X$ un préschéma localement noethérien, $f:X\to S$ un morphisme plat, $g:Y\to S$ un morphisme localement de type fini.
Soient $U$ un ouvert dense dans $S$, $h:f^{-1}(U)\to Y$ un $S$-morphisme, $T=S-U$, $Z$ le sous-préschéma fermé réduit de $X$ ayant $f^{-1}(T)$ comme espace sous-jacent, $h_0:Z\to Y$ un $S$-morphisme.
Supposons $X$ normal en tous les points de $Z$.
Afin qu'il existe un $S$-morphisme (nécessairement unique) $h':X\to Y$ prolongeant $h$ et $h_0$, il faut et il suffit que la condition suivante soit satisfaite :

\emph{Pour tout spectre $S_1$ d'un anneau de valuation discrète, de point fermé $a$ et de point générique $b$, et tout morphisme $u:S_1\to S$ tel que $u(a)\in T$ et $u(b)\in U$, il existe un $S_1$-morphisme $h'_1:X_{(S_1)}\to Y_{(S_1)}$ qui prolonge $h_{(S_1)}:f_{(S_1)}^{-1}(b)\to Y_{(S_1)}$ et est tel que le diagramme}
\[
  \xymatrix{
  \Spec(k(z_1)) \ar[r] \ar[d] & Z_{(S_1)} \ar[d]^{(h_0)_{(S_1)}} \\
  X_{(S_1)} \ar[r]_{h'_1} & Y_{(S_1)}
  }
\]
\emph{soit commutatif pour tout $z_1\in Z_{(S_1)}$.}

L'hypothèse que $f$ est plat entraîne en effet que $f^{-1}(U)$ est dense dans $X$ \sref{IV.2.3.10}, et il suffit alors d'appliquer \sref{IV.20.4.8}.
\end{corollary}

\begin{corollary}[20.4.10]
\label{IV.20.4.10}
Sous les hypothèses de \sref{IV.20.4.6}, supposons en outre $Y$ séparé et localement quasi-fini sur $S$.
Soient $U$ un ouvert dense dans $X$, $h:U\to Y$ un $S$-morphisme, $\omega$ la $S$-application rationnelle correspondante, $x\in X-U$ un point normal de $X$.
Afin que $\omega$ soit définie au point $x$, il faut et il suffit que la condition suivante soit vérifiée : il existe un spectre $S_1$ d'un anneau de valuation discrète, de point fermé $a$ et de point générique $b$, et un morphisme $u:S_1\to X$
\oldpage[IV-4]{249}
tels que $u(a)=x$, $u(b)\in U$ et que la $S$-application rationnelle $\omega\circ u$ soit partout définie.

En effet, par hypothèse toutes les fibres du morphisme projection $X\times_S Y\to X$ sont formées de points isolés, et pour appliquer \sref{IV.20.4.4} il suffit de savoir que la fibre $p^{-1}(x)$ est non vide dans $\Gamma_\omega$.
Or si $h_1$ est l'unique morphisme de la classe $\omega\circ u$, l'unique point de $X\times_S Y$ au-dessus de $x$ et de $h_1(a)$ appartient à $\Gamma_\omega$, d'où la conclusion.
\end{corollary}

\begin{proposition}[20.4.11]
\label{IV.20.4.11}
Soient $X$ un préschéma localement noethérien, $Y$ un $S$-préschéma affine sur $S$, $U$ un ouvert de $X$, $Z=X-U$.
Supposons que l'on ait $\operatorname{prof}_Z(\sh{O}_X)\geq 2$ \sref{IV.5.10.1} ; alors tout $S$-morphisme $f:U\to Y$ se prolonge de façon unique en un $S$-morphisme de $X$ dans $Y$.
\end{proposition}

\begin{proof}
On peut se borner au cas où $S$ et $X$ sont affines et (en vertu de \sref[I]{I.3.3.14}) au cas où $S=X$ ; on a donc $X=\Spec(A)$, $Y=\Spec(B)$, $B$ étant une $A$-algèbre de type fini.
Comme $B$ est quotient d'une algèbre de polynômes $A[(T_\lambda)]_{\lambda\in L}$, $Y$ est un sous-préschéma fermé de $Y'=X[(T_\lambda)]_{\lambda\in L}$.
D'autre part, l'hypothèse sur $Z$ entraîne que $U$ est schématiquement dense dans $X$ en vertu de \sref{IV.20.2.13}[(iv)] et \sref{IV.5.10.2}.
Si l'on prouve que tout $X$-morphisme $u:U\to Y$ se prolonge de façon unique en un $X$-morphisme $v':X\to Y'$, il en résultera que $v'$ se factorisera en $X\xrightarrow{v}Y\to Y'$ ; en effet, le sous-préschéma $v'^{-1}(Y)$ est fermé et majore $U$ \sref[I]{I.4.4.1}, donc est identique à $X$ \sref{IV.20.3.8.8}.
Dans ces conditions, $v$ sera l'unique $S$-morphisme de $X$ dans $Y$ prolongeant $u$.
On peut donc se borner au cas où $Y=Y'$.
Mais alors il y a correspondance biunivoque entre les $X$-morphismes d'un ouvert $U\subset X$ dans $Y'$ et les familles $(s_\lambda)_{\lambda\in L}$ de sections de $\sh{O}_X$ au-dessus de $U$ \sref[II]{II.1.7.9} ; la conclusion résulte donc de \sref{IV.5.10.5}.
\end{proof}

\begin{corollary}[20.4.12]
\label{IV.20.4.12}
Soient $X$ un $S$-préschéma localement noethérien réduit et vérifiant la condition $(S_2)$ \sref{IV.5.7.2} (par exemple un $S$-préschéma localement noethérien et normal \sref{IV.5.8.6}), $Y$ un $S$-préschéma affine sur $S$, $f$ une $S$-application rationnelle de $X$ dans $Y$ ; alors toute composante irréductible de $X-\operatorname{dom}(f)$ est de codimension $1$ dans $X$.
\end{corollary}

\begin{proof}
Il revient au même de dire que si $Z^{(2)}$ est l'ensemble des $x\in X$ tels que $\dim(\sh{O}_{X,x})\geq 2$, alors, pour toute partie fermée $Z\subset Z^{(2)}$ de $X$, tout $S$-morphisme de $X-Z$ dans $Y$ se prolonge en un $S$-morphisme de $X$ dans $Y$ ; or, cela résulte de l'hypothèse sur $X$ \sref{IV.5.7.2} et de \sref{IV.20.4.11}.
\end{proof}

\subsection{Pseudo-morphismes relatifs}
\EGAReaderTarget{ega4.c04.subsec.20-5}{IV.20.5}

\begin{env}[20.5.1]
\EGAReaderTarget{ega4.c04.statement.20-5-1}{IV.20.5.1}
Soient $X$, $Y$ deux $S$-préschémas.
Il résulte des définitions \sref{IV.11.10.8} que l'intersection de deux ouverts $U$, $U'$ de $X$, \emph{universellement schématiquement denses par rapport à $S$}, possède encore cette propriété.
On définit donc une relation d'équivalence entre $S$-morphismes $u:U\to Y$ en remplaçant dans \sref{IV.20.2.1} « schématiquement dense » par « universellement schématiquement dense par rapport à $S$ ».
Une classe d'équivalence suivant cette relation s'appelle un \emph{pseudo-morphisme de $X$ dans $Y$ relativement à $S$}, et l'ensemble de ces classes se note $\operatorname{Ps.hom}_{X/S}(X,Y)$.
\end{env}

\begin{env}[20.5.2]
\label{IV.20.5.2}
Comme tout ouvert de $X$ universellement schématiquement dense par rapport à $S$ est en particulier schématiquement dense, les éléments d'un pseudo-morphisme
\oldpage[IV-4]{250}
de $X$ dans $Y$ relativement à $S$ sont équivalents au sens de \sref{IV.20.2.1}, d'où une application canonique
\[
\label{IV.20.5.2.1}
  \operatorname{Ps.hom}_{X/S}(X,Y)\to\operatorname{Ps.hom}_S(X,Y).
  \tag{20.5.2.1}
\]
\end{env}

\begin{proposition}[20.5.3]
\label{IV.20.5.3}
Supposons $Y$ séparé sur $S$.
Alors :
\begin{enumerate}[label=\emph{(\roman*)}]
  \item L'application \sref{IV.20.5.2.1} est injective et identifie $\operatorname{Ps.hom}_{X/S}(X,Y)$ au sous-ensemble de $\operatorname{Ps.hom}_S(X,Y)$ formé des pseudo-$S$-morphismes $\omega$ tels que $\operatorname{dom}_S(\omega)$ soit universellement schématiquement dense relativement à $S$.
  \item Le préfaisceau $U\mapsto\operatorname{Ps.hom}_{U/S}(U,Y)$ sur $X$ est un faisceau.
\end{enumerate}
\end{proposition}

\begin{proof}
L'assertion~(i) est immédiate, car dire que deux $S$-morphismes $u:U\to Y$, $u':U'\to Y$ sont équivalents au sens de \sref{IV.20.2.1} signifie alors que ce sont les restrictions d'un même morphisme $\operatorname{dom}_S(\omega)\to Y$ \sref{IV.20.2.4}, et si $U$ et $U'$ sont universellement schématiquement denses par rapport à $S$, il en est de même \emph{a fortiori} de $\operatorname{dom}_S(\omega)$.
Pour prouver~(ii), notons que $U\mapsto\operatorname{Ps.hom}_S(U,Y)$ est alors un faisceau \sref{IV.20.2.6} ; d'autre part, étant donné un recouvrement ouvert $(U_\alpha)$ de $U$, et un pseudo-$S$-morphisme $\omega$ de $U$ dans $Y$, pour que $\operatorname{dom}_S(\omega)$ soit universellement schématiquement dense dans $U$ relativement à $S$, il résulte aussitôt des définitions (cf.\ \sref{IV.20.2.1}) qu'il faut et il suffit que chacun des ensembles $\operatorname{dom}_S(\omega)\cap U_\alpha=\operatorname{dom}_S(\omega|U_\alpha)$ soit universellement schématiquement dense dans $U_\alpha$ relativement à $S$ ; d'où~(ii).
\end{proof}

Lorsque $Y$ est séparé, on notera $\mathscr{P}\!s.\!\operatorname{hom}_{X/S}(X,Y)$ le sous-faisceau
\[
  U\mapsto\operatorname{Ps.hom}_{U/S}(U,Y)
\]
de $\mathscr{P}\!s.\!\operatorname{hom}_S(X,Y)$.

Lorsque $X$ est \emph{plat et localement de présentation finie} sur $S$ et $Y$ séparé sur $S$, les pseudo-morphismes de $X$ dans $Y$ relativement à $S$ s'identifient encore aux pseudo-$S$-morphismes $\omega$ tels que, pour toute fibre $X_s$ du morphisme $X\to S$, $\operatorname{dom}_S(\omega)\cap X_s$ soit \emph{schématiquement dense} dans $X_s$ \sref{IV.11.10.10}.

\begin{env}[20.5.4]
\EGAReaderTarget{ega4.c04.statement.20-5-4}{IV.20.5.4}
Un cas particulièrement important où $Y$ est séparé sur $S$ est le cas où $Y=S[T]=S\otimes_{\bb{Z}}\Spec(\bb{Z}[T])$ ($T$ indéterminée).
Il y a alors correspondance biunivoque entre les pseudo-$S$-morphismes de $X$ dans $Y$ et les \emph{pseudo-fonctions} sur $X$ \sref{IV.20.2.8} en vertu de la définition d'un produit de $\bb{Z}$-préschémas.
Les pseudo-morphismes de $X$ dans $Y$ relativement à $S$ s'identifient alors, en vertu de \sref{IV.20.5.3}, aux \emph{pseudo-fonctions $\varphi$ sur $X$ telles que $\operatorname{dom}(\varphi)$ soit universellement schématiquement dense relativement à $S$}.
Le faisceau $\mathscr{P}\!s.\!\operatorname{hom}_{X/S}(X,Y)$ est un \emph{sous-faisceau d'anneaux} de $\sh{M}'_X$, que l'on note $\sh{M}'_{X/S}$.

On pose alors $\operatorname{Ps.hom}_{X/S}(X,Y)=M'(X/S)$ et on dit que ses éléments sont les \emph{pseudo-fonctions sur $X$ relatives à $S$}.
\end{env}

\begin{env}[20.5.5]
\label{IV.20.5.5}
Soient $X$, $Y$, $Z$ trois $S$-préschémas, $f:Y\to Z$ un $S$-morphisme.
On peut, dans le raisonnement de \sref{IV.20.3.1}, remplacer partout « schématiquement dense » par « universellement schématiquement dense relativement à $S$ » ; pour tout pseudo-morphisme $\omega\in\operatorname{Ps.hom}_{X/S}(X,Y)$, les morphismes $f\circ u$, où $u$ parcourt $\omega$, appartiennent donc à une même classe d'équivalence (pour la relation définie dans \sref{IV.20.5.1}), que l'on
\oldpage[IV-4]{251}
appelle le pseudo-morphisme de $X$ dans $Z$, relativement à $S$, \emph{composé} de $f$ et de $\omega$ et que l'on note $f\circ\omega$.
Si $g:Z\to T$ est un morphisme, il est immédiat que $g\circ(f\circ\omega)=(g\circ f)\circ\omega$.
\end{env}

\begin{env}[20.5.6]
\label{IV.20.5.6}
Supposons $Y$ \emph{séparé} sur $S$, et soit $\omega$ un pseudo-morphisme de $X$ dans $Y$ relativement à $S$.
Si $f:X'\to X$ est un $S$-morphisme tel que $f^{-1}(\operatorname{dom}_S(\omega))$ soit \emph{universellement schématiquement dense dans $X'$ relativement à $S$}, il résulte de \sref{IV.20.3.2} que le pseudo-$S$-morphisme $\omega\circ f$ a un domaine universellement schématiquement dense relativement à $S$, donc \sref{IV.20.5.3} peut être considéré comme un pseudo-morphisme \emph{relativement à $S$}.
Lorsque $X'$ est plat et localement de présentation finie sur $S$, la condition que $f^{-1}(\operatorname{dom}_S(\omega))$ soit universellement schématiquement dense relativement à $S$ équivaut encore à dire que pour tout $s\in S$, $(f^{-1}(\operatorname{dom}_S(\omega)))_s$ (notation de \sref{IV.11.10.10}) est schématiquement dense dans $X'_s$, ou encore, en notant $f_s:X'_s\to X_s$ le morphisme déduit de $f$ par changement de base, que l'image réciproque par $f_s$ de $(\operatorname{dom}_S(\omega))_s$ soit schématiquement dense dans $X'_s$.
Cette dernière condition sera en particulier remplie si, pour tout $s\in S$, $X_s$, $X'_s$ et $f_s$ vérifient l'une des trois conditions~(i),~(ii),~(iii) de \sref{IV.20.3.5}.
\end{env}

\begin{env}[20.5.7]
\label{IV.20.5.7}
Supposons maintenant que $X$ et $X'$ soient tous deux des $S$-préschémas \emph{plats et localement de présentation finie sur $S$}, et que $f:X'\to X$ soit un $S$-morphisme \emph{plat} (ou, ce qui revient au même \sref{IV.11.3.10}, que pour tout $s\in S$, $f_s:X'_s\to X_s$ soit un morphisme plat).
Alors, pour tout ouvert $U$ de $X$ et tout ouvert $V\subset U$ universellement schématiquement dense dans $U$ relativement à $S$, il résulte de \sref{IV.11.10.5} et \sref{IV.11.10.9} que $f^{-1}(V)$ est universellement schématiquement dense dans $f^{-1}(U)$ relativement à $S$.
Pour tout pseudo-morphisme $\omega$ de $X$ dans $Y$ relativement à $S$, il résulte de \sref{IV.20.3.4} que le pseudo-$S$-morphisme $\omega\circ f$ est défini et est un pseudo-morphisme de $X'$ dans $Y$ relativement à $S$, \emph{même lorsqu'on ne suppose pas $Y$ séparé sur $S$}.
On en déduit que dans ce cas, pour tout $S$-morphisme $g:Y\to Z$, $(g\circ\omega)\circ f$ est encore défini et égal à $g\circ(\omega\circ f)$ (avec les définitions de \sref{IV.20.5.1}), et est donc un pseudo-morphisme \emph{relativement à $S$}.
\end{env}

\begin{env}[20.5.8]
\label{IV.20.5.8}
Soient $X$ un $S$-préschéma, $S'\to S$ un morphisme quelconque, $X'=X_{(S')}$, $g:X'\to X$ la projection canonique.
Alors, pour tout ouvert $U$ de $X$ et tout ouvert $V\subset U$ universellement schématiquement dense dans $U$ relativement à $S$, $V'=g^{-1}(V)$ est universellement schématiquement dense dans $U'=g^{-1}(U)$ relativement à $S'$ par définition \sref{IV.11.10.8}.
Soit alors $\omega$ un pseudo-morphisme de $X$ dans un $S$-préschéma $Y$ relativement à $S$ ; si $u_1$, $u_2$ sont des $S$-morphismes de la classe $\omega$, définis respectivement dans des ouverts $U_1$, $U_2$ de $X$, universellement schématiquement denses dans $X$ relativement à $S$, il résulte de ce qui précède que les morphismes $u'_1=u_1\circ(g|g^{-1}(U_1))$ et $u'_2=u_2\circ(g|g^{-1}(U_2))$ coïncident dans un ouvert $U'_3$ universellement schématiquement dense relativement à $S'$.
Or, si $Y'=Y_{(S')}$ et si $h:Y'\to Y$ est la projection canonique, $u'_1$ et $u'_2$ se factorisent canoniquement en $u'_1=h\circ v_1$, $u'_2=h\circ v_2$, et $v_1$ et $v_2$ sont deux $S'$-morphismes dans $Y'$ qui coïncident dans $U'_3$.
On voit donc que lorsque $u_1$ parcourt la classe $\omega$, les $S'$-morphismes $v_1$ correspondants appartiennent à un même pseudo-morphisme relativement à $S'$, dit \emph{image réciproque} de $\omega$ par le morphisme changement de base $S'\to S$ et noté $\omega_{(S')}$.
Il est clair que si $S''\to S'$ est un second morphisme, on a $(\omega_{(S')})_{(S'')}=\omega_{(S'')}$ (pour le morphisme changement de base composé $S''\to S'\to S$).
\end{env}

\oldpage[IV-4]{252}
\subsection{Fonctions méromorphes relatives}
\EGAReaderTarget{ega4.c04.subsec.20-6}{IV.20.6}

\begin{env}[20.6.1]
\EGAReaderTarget{ega4.c04.statement.20-6-1}{IV.20.6.1}
Soient $S$ un préschéma, $X$ un $S$-préschéma qui est \emph{plat et localement de présentation finie} sur $S$ ; pour tout $s\in S$, nous noterons $X_s$ la fibre au point $s$ du morphisme structural $X\to S$.
En général, si $\varphi$ est une fonction méromorphe sur $X$, il n'est pas possible de lui associer de façon « naturelle », pour chaque $s\in S$, une fonction méromorphe « induite » sur $X_s$ \sref{IV.20.1.11}.
Pour tout ouvert $U$ de $X$, notons $T_{X/S}(U)$ l'ensemble des sections $t\in\Gamma(U,\sh{O}_X)$ telles que, pour tout $s\in S$, l'image $t_s$ de $t$ par l'homomorphisme canonique $\Gamma(U,\sh{O}_X)\to\Gamma(U\cap X_s,\sh{O}_{X_s})$ soit une section \emph{régulière} ; cela implique d'ailleurs, par l'équivalence de a) et b) dans \sref{IV.11.3.7}, que $t$ est elle-même une section \emph{régulière}.
Il est clair que $U\mapsto T_{X/S}(U)$ est un \emph{sous-faisceau} du faisceau d'ensemble $\sh{S}=\sh{S}(\sh{O}_X)$ (notation de \sref{IV.20.1.3}), que l'on note $\sh{T}_{X/S}$ ; on pose
\[
\label{IV.20.6.1.1}
  \sh{M}_{X/S}=\sh{O}_X[\sh{T}_{X/S}^{-1}]
  \tag{20.6.1.1}
\]
(notation de \sref{IV.20.1.1}) et on dit que ce faisceau d'anneaux est le \emph{faisceau des germes de fonctions méromorphes sur $X$ relatives à $S$} ; ses sections au-dessus de $X$ sont appelées les \emph{fonctions méromorphes sur $X$ relatives à $S$} et leur ensemble est noté $M(X/S)$.
Il est clair que $\sh{M}_{X/S}$ est un \emph{sous-faisceau} de $\sh{M}_X=\sh{O}_X[\sh{S}^{-1}]$ ; pour toute fonction méromorphe $\varphi\in M(X/S)$ et tout $s\in S$, l'image réciproque de $\varphi$ par le morphisme injection canonique $j_s:X_s\to X$ est alors définie \sref{IV.20.1.11}, et notée $\varphi_s$.
\end{env}

\begin{env}[20.6.2]
\EGAReaderTarget{ega4.c04.statement.20-6-2}{IV.20.6.2}
Soit maintenant $\sh{F}$ un $\sh{O}_X$-Module quasi-cohérent ; on pose
\[
\label{IV.20.6.2.1}
  \sh{M}_{X/S}(\sh{F})=\sh{F}[\sh{T}_{X/S}^{-1}]=\sh{F}\otimes_{\sh{O}_X}\sh{M}_{X/S};
  \tag{20.6.2.1}
\]
les sections de $\sh{M}_{X/S}(\sh{F})$ sont dites \emph{sections méromorphes de $\sh{F}$ au-dessus de $X$, relatives à $S$} et leur ensemble est noté $M(X/S,F)$.
L'homomorphisme canonique $\sh{F}\to\sh{M}_{X/S}(\sh{F})$ n'est pas nécessairement injectif ; lorsqu'il l'est, on dit que $\sh{F}$ est \emph{sans torsion relativement à $S$} : cela signifie que pour tout ouvert $U$ de $X$ et toute section $t\in T_{X/S}(U)$, $t$ est $(\sh{F}|U)$-régulière ; cette condition est remplie \emph{a fortiori} lorsque $\sh{F}$ est \emph{strictement sans torsion} \sref{IV.20.1.5}.
Dans ce dernier cas, il résulte aussitôt des définitions \sref{IV.20.1.2} que l'homomorphisme canonique de $\sh{O}_X$-Modules
\[
\label{IV.20.6.2.2}
  \sh{M}_{X/S}(\sh{F})\to\sh{M}_X(\sh{F})
  \tag{20.6.2.2}
\]
est \emph{injectif}, de sorte que les sections méromorphes de $\sh{F}$ relatives à $S$ sont des sections méromorphes de $\sh{F}$ au sens de \sref{IV.20.1.3}.
\end{env}

\begin{proposition}[20.6.3]
\label{IV.20.6.3}
L'image par l'homomorphisme injectif \sref{IV.20.2.10.1} du sous-faisceau $\sh{M}_{X/S}$ de $\sh{M}_X$ est le sous-faisceau $\sh{M}'_{X/S}$ des pseudo-fonctions sur $X$ relatives à $S$ (i.e.\ des pseudo-fonctions dont le domaine de définition est universellement schématiquement dense relativement à $S$ \sref{IV.20.5.4}).
\end{proposition}

\begin{proof}
On peut évidemment se borner à prouver que l'image de $M(X/S)$ par l'homomorphisme canonique $M(X)\to M'(X)$ est égale à $M'(X/S)$ ; la proposition est alors conséquence de la proposition plus générale suivante :
\end{proof}

\oldpage[IV-4]{253}
\begin{proposition}[20.6.4]
\label{IV.20.6.4}
Soit $\sh{F}$ un $\sh{O}_X$-Module quasi-cohérent de présentation finie et strictement sans torsion.
Alors, pour qu'une section méromorphe $\varphi$ de $\sh{F}$ au-dessus de $X$ soit une section méromorphe relative à $S$, il faut et il suffit que $\operatorname{dom}(\varphi)$ soit universellement schématiquement dense relativement à $S$.
\end{proposition}

\begin{proof}
La nécessité de la condition résulte de \sref{IV.20.2.15} appliqué à chaque injection canonique $X_s\to X$ ($s\in S$), compte tenu de \sref{IV.11.10.9}.
Pour voir que la condition est suffisante, il faut prouver que pour tout $x\in X$, il existe un voisinage ouvert $U$ de $x$ dans $X$ et une section de $\sh{M}_{X/S}(\sh{F})$ au-dessus de $U$ dont la restriction à $U\cap\operatorname{dom}(\varphi)$ coïncide avec $\varphi$ dans un ouvert schématiquement dense de $U\cap\operatorname{dom}(\varphi)$.
Considérons l'Idéal des dénominateurs $\sh{J}$ de $\varphi$ \sref{IV.20.2.14}, qui est quasi-cohérent, et qui définit un sous-préschéma fermé de $X$ dont l'espace sous-jacent est $X-\operatorname{dom}(\varphi)$.
Par hypothèse, si $s$ est l'image de $x$ dans $S$, $\operatorname{dom}(\varphi)\cap X_s$ est schématiquement dense dans le préschéma localement noethérien $X_s$, donc \sref{IV.20.2.13}[(iv)] contient $\operatorname{Ass}(\sh{O}_{X_s})$ ; cela implique donc que l'idéal $\sh{J}_x$ de $\sh{O}_{X,x}$ a une image dans $\sh{O}_{X_s,x}=\sh{O}_{X,x}\otimes_{\sh{O}_{S,s}}k(s)$ qui n'est contenue dans aucun des idéaux premiers $\mathfrak{p}_i\in\operatorname{Ass}(\sh{O}_{X_s,x})$ (en nombre fini) ; donc (Bourbaki, \emph{Alg.~comm.}, chap.~II, \textsection1, n\textsuperscript{o}~1, prop.~2) il existe un élément $t_x\in\sh{J}_x$ dont l'image dans $\sh{O}_{X_s,x}$ n'appartient à aucun des $\mathfrak{p}_i$, et est par suite régulière dans cet anneau noethérien.
Soit $t$ une section de $\sh{J}$ au-dessus d'un voisinage ouvert affine $U$ de $x$, dont $t_x$ soit le germe au point $x$ ; puisque $X$ est plat et localement de présentation finie sur $S$, on peut supposer \sref{IV.11.3.8} que $t$ est une section \emph{régulière} de $\sh{O}_X$ au-dessus de $U$ et que, pour tout $s'\in S$, l'image de $t$ dans $\Gamma(U\cap X_{s'},\sh{O}_{X_{s'}})$ soit aussi \emph{régulière} ; autrement dit, on a $t\in T_{X/S}(U)$.
Mais alors, par définition de $\sh{J}$, puisque $\sh{F}$ est strictement sans torsion, $t(\varphi|(U\cap\operatorname{dom}(\varphi)))$ est une section $u$ de $\sh{F}$ au-dessus de $U\cap\operatorname{dom}(\varphi)$ ; d'autre part, $U\cap\operatorname{dom}(\varphi)$ contient l'ensemble ouvert $U_t$ des points $x'\in U$ où $t(x')\neq0$, et ce dernier contient $x$ et est schématiquement dense dans $U$ \sref{IV.20.2.9}.
On voit donc que dans $U_t$, $\varphi$ coïncide avec la restriction à $U_t$ de la section $u/t$ de $\sh{M}_{X/S}(\sh{F})$ au-dessus de $U$.
C.Q.F.D.
\end{proof}

\begin{remarks}[20.6.5]
\EGAReaderTarget{ega4.c04.statement.20-6-5}{IV.20.6.5}
\begin{enumerate}[label=(\roman*)]
  \item Soit $\varphi$ une fonction méromorphe sur $X$ relative à $S$, de sorte que pour tout $s\in S$, $\varphi_s$ est une fonction méromorphe sur $X_s$ \sref{IV.20.6.1} ; en vertu de \sref{IV.20.1.11.1}, on a
  \[
  \label{IV.20.6.5.1}
    \operatorname{dom}(\varphi)\cap X_s\subset\operatorname{dom}(\varphi_s).
    \tag{20.6.5.1}
  \]
  Mais il convient de noter que même lorsque $S$ est le spectre d'un anneau de valuation discrète $A$ et $X=S[T]$ ($T$ indéterminée), les deux membres de \sref{IV.20.6.5.1} ne sont pas nécessairement égaux : par exemple, si $\pi$ est une uniformisante de $A$, il est immédiat que $\varphi=\pi/T$ est une fonction méromorphe sur $S$ relative à $S$, car si $a$ et $b$ sont le point fermé et le point générique de $S$, $T$ est régulière dans $\Gamma(X_a,\sh{O}_{X_a})=k[T]$ et dans $\Gamma(X_b,\sh{O}_{X_b})=K[T]$, $k$ et $K$ étant le corps résiduel et le corps des fractions de $A$.
  On a $\operatorname{dom}(\varphi)=\mathrm{D}(T)$ dans $X$, mais $\operatorname{dom}(\varphi_a)=X_a$ puisque $\varphi_a=0$.

  \item Pour qu'une fonction méromorphe $\varphi$ relative à $S$ soit \emph{inversible dans l'anneau} $M(X/S)$, il faut et il suffit que pour tout $s\in S$, $\varphi_s$ soit inversible dans $M(X_s)$ (autrement dit, que $\varphi_s$ soit une fonction méromorphe régulière sur $X_s$ \sref{IV.20.1.8}).
  La condition est
  \oldpage[IV-4]{254}
  en effet trivialement nécessaire.
  Inversement, si elle est vérifiée, et si $x$ est un point quelconque de $X$, $s$ son image dans $S$, il existe par hypothèse un voisinage ouvert $U$ de $x$ dans $X$ et deux sections $u$, $t$ de $\sh{O}_X$ au-dessus de $U$ telles que $t\in T_{X/S}(U)$ et $\varphi|U=u/t$ ; l'hypothèse entraîne que si $u_s$ est l'image de $u$ dans $\Gamma(U\cap X_s,\sh{O}_{X_s})$, $u_s$ est régulière au point $x$.
  En restreignant $U$, on peut donc supposer, en vertu de \sref{IV.11.3.8}, que $u\in T_{X/S}(U)$, d'où la conclusion.

  Lorsque $\varphi$ est inversible dans $M(X/S)$, on dit encore que $\varphi$ est une \emph{fonction méromorphe régulière relative à $S$}.
  On notera que $\varphi\in M(X/S)$ peut être inversible dans $M(X)$ (autrement dit, être \emph{régulière} au sens de \sref{IV.20.1.8}) sans l'être dans $M(X/S)$, comme le montre aussitôt l'exemple de~(i).

  \item Soit $\sh{L}$ un $\sh{O}_X$-Module \emph{inversible}, et soit $\varphi$ une section méromorphe \emph{régulière} de $\sh{L}$ au-dessus de $X$ \sref{IV.20.1.8} ; on dit que $\varphi$ est \emph{régulière relativement à $S$} si, pour tout ouvert $U$ de $X$ tel que $\sh{L}|U$ soit isomorphe à $\sh{O}_U$, $\varphi|U$ correspond à un élément de $\Gamma(U,\sh{M}_X)$ qui est \emph{régulière relativement à $S$} ; en vertu de~(ii), il est immédiat qu'il faut et il suffit pour cela que, pour tout $s\in S$, $\varphi_s$ soit une section méromorphe \emph{régulière} \sref{IV.20.1.8} du $\sh{O}_{X_s}$-Module inversible $\sh{L}_s=\sh{L}\otimes_{\sh{O}_X}k(s)$.
  Si $\varphi'$ est l'inverse de $\varphi$ dans $\sh{L}^{-1}$ \sref{IV.20.1.10}, $\varphi'$ est alors aussi régulière relative à $S$.
  Si $\sh{L}_1$ est un second $\sh{O}_X$-Module inversible, $\varphi_1$ une section méromorphe de $\sh{L}_1$ au-dessus de $X$, régulière relativement à $S$, alors $\varphi\otimes\varphi_1$ est une section méromorphe de $\sh{L}\otimes\sh{L}_1$ au-dessus de $X$, régulière relativement à $S$.
\end{enumerate}
\end{remarks}

\begin{proposition}[20.6.6]
\label{IV.20.6.6}
Soient $X$ un $S$-préschéma plat et localement de présentation finie sur $S$, $\sh{F}$ un $\sh{O}_X$-Module localement libre de type fini ; pour tout $s\in S$, on note $X_s$ la fibre au point $s$ du morphisme structural $f:X\to S$.
Soit $\varphi$ une section méromorphe de $\sh{F}$ au-dessus de $X$, relativement à $S$, et supposons que $\varphi$ soit définie en tous les points $x\in X$ tels que $\operatorname{prof}(\sh{O}_{X_{f(x)},x})=1$.
Alors $\varphi$ est partout définie.
\end{proposition}

\begin{proof}
Par hypothèse, $\operatorname{dom}(\varphi)\cap X_s$ est schématiquement dense dans $X_s$ pour tout $s\in S$, donc contient les points $x$ de $X_s$ tels que $\operatorname{prof}(\sh{O}_{X_s,x})=0$ \sref{IV.5.10.2} ; l'hypothèse signifie donc que si l'on pose $Z=X-\operatorname{dom}(\varphi)$, on a $\operatorname{prof}(\sh{O}_{X_{f(x)},x})\geq2$ en tous les points de $Z$.
Il suffit donc d'appliquer \sref{IV.19.9.8}.
\end{proof}

\begin{env}[20.6.7]
\EGAReaderTarget{ega4.c04.statement.20-6-7}{IV.20.6.7}
Soient $X$, $X'$ deux $S$-préschémas plats et localement de présentation finie sur $S$, $f=(\psi,\theta):X'\to X$ un $S$-morphisme.
Pour tout ouvert $U$ de $X$, notons $T_f(U)$ l'ensemble des sections $t\in T_{X/S}(U)$ dont l'image dans $\Gamma(f^{-1}(U),\sh{O}_{X'})$ appartienne à $T_{X'/S}(f^{-1}(U))$ ; il est immédiat que $U\mapsto T_f(U)$ est un sous-faisceau du faisceau d'ensembles $\sh{T}_{X/S}$, que l'on note $\sh{T}_f$.
On pose $\sh{M}_{X/S,f}=\sh{O}_X[\sh{T}_f^{-1}]$ ; c'est un sous-faisceau d'anneaux de $\sh{M}_{X/S}$, et on en déduit canoniquement de $\theta^\#: \psi^*(\sh{O}_X)\to\sh{O}_{X'}$ un homomorphisme de faisceaux d'anneaux $\theta'^\#: \psi^*(\sh{M}_{X/S,f})\to\sh{M}_{X'/S}$ prolongeant $\theta^\#$.
Si une fonction méromorphe $\varphi$ sur $X$, relative à $S$, est \emph{une section de $\sh{M}_{X/S,f}$}, $\Gamma(\theta'^\#)(\varphi)$ est une fonction méromorphe sur $X'$, dite \emph{image réciproque de $\varphi$ par $f$}, et notée $\varphi\circ f$ si cela n'entraîne pas confusion.
On étend aussitôt de même les définitions de \sref{IV.20.1.11} relatives aux $\sh{O}_X$-Modules.
\end{env}

\begin{proposition}[20.6.8]
\label{IV.20.6.8}
Avec les notations de \sref{IV.20.6.7}, si le $S$-morphisme $f:X'\to X$ est \emph{plat}, on a $\sh{M}_{X/S,f}=\sh{M}_{X/S}$, et l'homomorphisme $\varphi\mapsto\varphi\circ f$ est défini dans $M(X/S)$ tout entier.
\end{proposition}

\begin{proof}
\oldpage[IV-4]{255}
En effet, l'hypothèse entraîne, en vertu de \sref{IV.11.3.10}, que pour tout $s\in S$, $f_s:X'_s\to X_s$ est plat ; donc, pour toute section $t\in T_{X/S}(U)$, si $t'$ est son image réciproque dans $\Gamma(f^{-1}(U),\sh{O}_{X'})$, $t'_{s}$, qui est l'image réciproque de $t_s$, est une section \emph{régulière} de $\sh{O}_{X'_s}$ au-dessus de $f^{-1}(U)\cap X'_s$, en vertu de \sref{IV.20.1.12} ; on en conclut que par définition $t'\in T_{X'/S}(f^{-1}(U))$, d'où la proposition.
\end{proof}

On déduit de là, comme dans \sref{IV.20.1.12}, un homomorphisme canonique de $\sh{O}_{X'}$-Algèbres
\[
\label{IV.20.6.8.1}
  f^*(\sh{M}_{X/S})\to\sh{M}_{X'/S}.
  \tag{20.6.8.1}
\]

\begin{env}[20.6.9]
\label{IV.20.6.9}
Considérons enfin un morphisme \emph{quelconque} $S'\to S$, et posons $X'=X\times_S S'$, qui est plat et localement de présentation finie sur $S'$ ; soit $p:X'\to X$ la projection canonique.
Soient $U$ un ouvert de $X$, $t$ une section appartenant à $T_{X/S}(U)$, $t'$ son image dans $\Gamma(p^{-1}(U),\sh{O}_{X'})$ ; pour tout $s'\in S'$, si $s\in S$ est l'image de $s'$, on a $X'_{s'}=X_s\otimes_{k(s)}k(s')$, donc le morphisme $X'_{s'}\to X_s$ est plat, et par suite \sref{IV.20.1.12} l'image réciproque $t'_{s'}$ de $t_s$ dans $\Gamma(p^{-1}(U)\cap X'_{s'},\sh{O}_{X'_{s'}})$ est régulière ; ceci prouve que l'on a $t'\in T_{X'/S'}(p^{-1}(U))$.
Ceci permet de définir canoniquement, comme dans \sref{IV.20.6.8}, un homomorphisme canonique de $\sh{O}_{X'}$-Algèbres
\[
\label{IV.20.6.9.1}
  p^*(\sh{M}_{X/S})\to\sh{M}_{X'/S'}.
  \tag{20.6.9.1}
\]

Moyennant l'identification permise par \sref{IV.20.6.3}, cette notion de changement de base pour les fonctions méromorphes relatives est un cas particulier de la notion analogue pour les pseudo-morphismes relatifs \sref{IV.20.5.8}.
\end{env}

% FIN SOURCE ega4/ega4-20-fr.tex

%
% DEBUT SOURCE ega4/ega4-21a-fr.tex
\section{Diviseurs}
\label{section:IV.21}

Sur le contenu du présent paragraphe, voir les commentaires de l'introduction du \textsection~20.
Pour les propriétés \emph{globales} des diviseurs, le lecteur se reportera au paragraphe qui leur est consacré au chap.~V.

\subsection{Diviseurs sur un espace annelé}
\EGAReaderTarget{ega4.c04.subsec.21-1}{subsection:IV.21.1}
\label{IV.21.1}

\begin{env}[21.1.1]
\label{IV.21.1.1}
Soient $(X,\sh{O}_X)$ un espace annelé, $\sh{M}_X$ le faisceau des germes de fonctions méromorphes sur $X$ \sref{IV.20.1.3}, $\sh{M}_X^*$ le faisceau (de groupes multiplicatifs) des germes de fonctions méromorphes régulières sur $X$ \sref{IV.20.1.8}.
Il est clair que le faisceau (de groupes multiplicatifs) $\sh{O}_X^*$ des germes de sections inversibles de $\sh{O}_X$ s'identifie canoniquement à un sous-faisceau (de groupes multiplicatifs) de $\sh{M}_X^*$.
\end{env}

\begin{definition}[21.1.2]
\EGAReaderTarget{ega4.c04.statement.21-1-2}{IV.21.1.2}
On appelle \emph{faisceau des diviseurs sur $X$} et on note $\sh{D}iv_X$ le faisceau quotient (de groupes commutatifs) $\sh{M}_X^*/\sh{O}_X^*$;
les sections de ce faisceau au-dessus de $X$ se nomment les \emph{diviseurs sur $X$}; ils forment un groupe commutatif noté $\operatorname{Div}(X)$.
Pour toute section $f$ de $\sh{M}_X^*$ au-dessus de $X$ (autrement dit, toute fonction méromorphe régulière sur $X$ \sref{IV.20.1.8}) on appelle \emph{diviseur de $f$} et on note $\operatorname{div}(f)$ (ou $\operatorname{div}_X(f)$) le diviseur sur $X$ image de $f$ par l'homomorphisme canonique
\[
  \Gamma(X,\sh{M}_X^*)\longrightarrow\Gamma(X,\sh{D}iv_X).
\]
\end{definition}

\oldpage[IV-4]{256}

Le \emph{support} d'un diviseur $D$ est l'ensemble fermé des $x\in X$ tels que $D_x\neq0$.
On le note $\operatorname{Supp}(D)$.

Pour tout ouvert $U$ de $X$, on a évidemment $\sh{M}_X^*|U=\sh{M}_U^*$, $\sh{O}_X^*|U=\sh{O}_U^*$, donc $\sh{D}iv_X|U=\sh{D}iv_U$, et par suite le faisceau $\sh{D}iv_X$ est égal au préfaisceau $U\mapsto\operatorname{Div}(U)$.

Lorsque $X=\operatorname{Spec}(A)$ est affine, on écrit $\operatorname{Div}(A)$ au lieu de $\operatorname{Div}(\operatorname{Spec}(A))$.

\begin{env}[21.1.3]
\label{IV.21.1.3}
Nous noterons toujours \emph{additivement} le groupe $\operatorname{Div}(X)$ des diviseurs sur $X$.
Pour deux fonctions méromorphes régulières $f$, $g$ sur $X$, on a donc
\[
\label{IV.21.1.3.1}
  \operatorname{div}(fg)=\operatorname{div}(f)+\operatorname{div}(g),
  \tag{21.1.3.1}
\]
\[
\label{IV.21.1.3.2}
  \operatorname{div}(f^{-1})=-\operatorname{div}(f).
  \tag{21.1.3.2}
\]
Par définition, pour toute fonction méromorphe régulière $f$ sur $X$, on a l'équivalence
\[
\label{IV.21.1.3.3}
  \operatorname{div}(f)=0 \iff f\in\Gamma(X,\sh{O}_X^*),
  \tag{21.1.3.3}
\]
d'où, pour deux fonctions méromorphes régulières $f$, $g$ sur $X$
\[
\label{IV.21.1.3.4}
  \operatorname{div}(f)=\operatorname{div}(g) \iff fg^{-1}\in\Gamma(X,\sh{O}_X^*).
  \tag{21.1.3.4}
\]
\end{env}

\begin{env}[21.1.4]
\EGAReaderTarget{ega4.c04.statement.21-1-4}{IV.21.1.4}
Soit maintenant $\sh{L}$ un $\sh{O}_X$-Module inversible, et soit $s$ une \emph{section méromorphe régulière} \sref{IV.20.1.8} de $\sh{L}$ au-dessus de $X$.
Tout $x\in X$ possède un voisinage ouvert $U$ tel que $\sh{L}|U$ soit isomorphe à $\sh{O}_U$, donc $\sh{M}_X(\sh{L})|U$ isomorphe à $\sh{M}_X|U$;
par un de ces isomorphismes, $s|U$ correspond à une section $f\in\Gamma(U,\sh{M}_X^*)$, et puisque deux isomorphismes de $\sh{L}|U$ sur $\sh{O}_U$ ne diffèrent que par la multiplication par un élément de $\Gamma(U,\sh{O}_X^*)$ (\hyperref[0.5.4.3]{$\mathbf{0}_{\mathrm I}$, 5.4.3}), l'élément $\operatorname{div}_U(f)$ de $\Gamma(U,\sh{D}iv_X)$ est \emph{indépendant} de l'isomorphisme choisi;
il est clair que ces éléments (pour $U$ variable) sont les restrictions d'une section de $\sh{D}iv_X$ au-dessus de $X$, que l'on appelle le \emph{diviseur de $s$} et que l'on note $\operatorname{div}(s)$ (un tel diviseur n'est pas nécessairement de la forme $\operatorname{div}(g)$ pour une fonction méromorphe régulière $g$ sur $X$; voir \sref{IV.21.2.9}).
Pour $\sh{L}=\sh{O}_X$, la définition de $\operatorname{div}(s)$ coïncide avec celle de \sref{IV.21.1.2}.
Si $\sh{L}$, $\sh{L}'$ sont deux $\sh{O}_X$-Modules inversibles, $s$ (resp. $s'$) une section méromorphe régulière de $\sh{L}$ (resp. $\sh{L}'$) au-dessus de $X$, il est immédiat que l'on a
\[
\label{IV.21.1.4.1}
  \operatorname{div}(s\otimes s')=\operatorname{div}(s)+\operatorname{div}(s')
  \tag{21.1.4.1}
\]
\[
\label{IV.21.1.4.2}
  \operatorname{div}(s^{\otimes n})=n.\operatorname{div}(s) \quad\text{pour tout}\quad n\in\mathbf{Z}
  \tag{21.1.4.2}
\]
($s^{-1}$ étant la section méromorphe régulière de $\sh{L}^{-1}$ au-dessus de $X$ définie dans \sref{IV.20.1.10})
et, pour deux sections méromorphes régulières $s$, $s'$ de $\sh{L}$ au-dessus de $X$, on a la relation
\[
\label{IV.21.1.4.3}
  \operatorname{div}(s)=\operatorname{div}(s') \iff s'=ts \quad\text{avec}\quad t\in\Gamma(X,\sh{O}_X^*).
  \tag{21.1.4.3}
\]
\end{env}

\begin{env}[21.1.5]
\label{IV.21.1.5}
Le faisceau $\sh{S}(\sh{O}_X)$ \sref{IV.20.1.3} dont les sections au-dessus d'un ouvert $U$ de $X$ sont les éléments réguliers de $\Gamma(U,\sh{O}_X)$ est un \emph{sous-faisceau de monoïdes} de $\sh{M}_X^*$; on peut écrire
\[
\label{IV.21.1.5.1}
  \sh{S}(\sh{O}_X)=\sh{O}_X\cap\sh{M}_X^*.
  \tag{21.1.5.1}
\]
\end{env}

\oldpage[IV-4]{257}

Si on note $\sh{S}(\sh{O}_X)^{-1}$ le faisceau dont les sections au-dessus de $U$ sont les inverses dans $\Gamma(U,\sh{M}_X^*)$ des éléments de $\Gamma(U,\sh{S}(\sh{O}_X))$, il est clair que l'on a $\Gamma(U,\sh{S}(\sh{O}_X))\cap\Gamma(U,\sh{S}(\sh{O}_X)^{-1})=\Gamma(U,\sh{O}_X^*)$, donc
\[
\label{IV.21.1.5.2}
  \sh{S}(\sh{O}_X)\cap\sh{S}(\sh{O}_X)^{-1}=\sh{O}_X^*.
  \tag{21.1.5.2}
\]

\begin{definition}[21.1.6]
\EGAReaderTarget{ega4.c04.def.21-1-6}{IV.21.1.6}
Le sous-faisceau d'ensembles de $\sh{D}iv_X$, image canonique du sous-faisceau $\sh{S}(\sh{O}_X)$ de $\sh{M}_X^*$, est noté $\sh{D}iv_X^+$; ses sections au-dessus de $X$ sont appelées \emph{diviseurs positifs sur $X$}, et leur ensemble est noté $\operatorname{Div}^+(X)$.
\end{definition}

Puisque $\sh{S}(\sh{O}_X)$ est un faisceau de monoïdes (multiplicatifs), on a
\[
\label{IV.21.1.6.1}
  \operatorname{Div}^+(X)+\operatorname{Div}^+(X)\subset\operatorname{Div}^+(X)
  \tag{21.1.6.1}
\]
et d'autre part, en vertu de \sref{IV.21.1.5.2} et \sref{IV.21.1.3.3}
\[
\label{IV.21.1.6.2}
  \operatorname{Div}^+(X)\cap(-\operatorname{Div}^+(X))=\{0\}.
  \tag{21.1.6.2}
\]
Ces deux relations montrent que $\operatorname{Div}^+(X)$ est l'ensemble des éléments positifs pour une \emph{structure d'ordre} sur le groupe $\operatorname{Div}(X)$, compatible avec cette structure de groupe; on note cette relation d'ordre $D\leq D'$, autrement dit, on a
\[
\label{IV.21.1.6.3}
  D\geq0 \iff D\in\operatorname{Div}^+(X).
  \tag{21.1.6.3}
\]

Nous supposerons toujours dans ce qui suit que $\operatorname{Div}(X)$ est muni de cette structure d'ordre;
il est clair que $\sh{D}iv_X^+|U=\sh{D}iv_U^+$, donc $\Gamma(U,\sh{D}iv_X^+)=\operatorname{Div}^+(U)$, et on peut donc dire que $\sh{D}iv_X^+$ définit sur $\sh{D}iv_X$ une structure de faisceau de groupes ordonnés.
La fibre $(\sh{D}iv_X^+)_x$ en un point $x$ du faisceau $\sh{D}iv_X^+$ est un sous-monoïde du groupe $(\sh{D}iv_X)_x$, ensemble des éléments $\geq0$ pour une structure d'ordre compatible avec la structure de groupe;
pour un diviseur $D$ sur $X$, il revient au même de dire que $D\geq0$ ou que $D_x\geq0$ pour tout $x\in X$.

Par définition, pour toute fonction méromorphe régulière $f$ sur $X$, on a la relation
\[
\label{IV.21.1.6.4}
  \operatorname{div}(f)\geq0 \iff f\in\Gamma(X,\sh{S}(\sh{O}_X))
  \tag{21.1.6.4}
\]
autrement dit, $\operatorname{div}(f)\geq0$ signifie que $f$ est une \emph{section régulière} de $\sh{O}_X$, ou encore une section de $\sh{O}_X$ \emph{inversible dans} $M(X)$.

Plus généralement, étant donné un diviseur $D$ sur $X$, la relation $\operatorname{div}(f)\geq D$ est équivalente à la suivante : pour tout ouvert $U$ de $X$ tel que $D|U=\operatorname{div}(g)$, où $g\in\Gamma(U,\sh{M}_X^*)$, il existe un élément régulier $t$ de $\Gamma(U,\sh{O}_X)$ tel que $f|U=tg$.

\begin{env}[21.1.7]
\label{IV.21.1.7}
Soient $\sh{L}$ un $\sh{O}_X$-Module inversible, $s$ une section méromorphe régulière de $\sh{L}$ au-dessus de $X$; on a la relation
\[
\label{IV.21.1.7.1}
  \operatorname{div}(s)\geq0 \iff s\in\Gamma(X,\sh{L})\cap\Gamma(X,(\sh{M}_X(\sh{L}))^*)
  \tag{21.1.7.1}
\]
comme il résulte aussitôt des définitions \sref{IV.21.1.4} et \sref{IV.21.1.6}.
\end{env}

\begin{proposition}[21.1.8]
\EGAReaderTarget{ega4.c04.prop.21-1-8}{IV.21.1.8}
Soient $X$ un préschéma localement noethérien, $D$ un diviseur sur $X$.
Supposons que pour tout $x\in X$ tel que $\operatorname{prof}(\sh{O}_{X,x})=1$ on ait $D_x\geq0$ (resp. $D_x=0$).
Alors on a $D\geq0$ (resp. $D=0$).
\end{proposition}

\begin{proof}
La question étant locale sur $X$, on peut supposer que $D=\operatorname{div}(f)$, $f$ étant une
\oldpage[IV-4]{258}
fonction méromorphe régulière sur $X$;
la relation $D_x\geq0$ équivaut à $x\in\operatorname{dom}(f)$, donc l'hypothèse signifie que si $T=X-\operatorname{dom}(f)$, on a $\operatorname{prof}(\sh{O}_{X,x})\geq2$ pour tout $x\in T$ (car $\operatorname{dom}(f)$ contient les points maximaux de $X$).
Par suite \sref{IV.5.10.5} l'homomorphisme de restriction $\Gamma(X,\sh{O}_X)\to\Gamma(X-T,\sh{O}_X)$ est bijectif, ce qui montre qu'il existe une section $s$ de $\sh{O}_X$ au-dessus de $X$ telle que $f=s|(X-T)$.
Mais par définition de $T$, cela entraîne $T=\varnothing$, donc $f=s$ et $D\geq0$.
L'assertion relative à la relation $D_x=0$ s'en déduit aussitôt en appliquant ce qui précède à $-D$, en vertu de \sref{IV.21.1.6.2}.
\end{proof}


\begin{corollary}[21.1.9]
\label{IV.21.1.9}
Soient $X$ un préschéma localement noethérien, $D$ un diviseur sur $X$.
Soit $S$ le support de $D$.
Alors, pour tout point maximal $x$ de $S$, on a $\operatorname{prof}(\sh{O}_{X,x})=1$.
\end{corollary}

\begin{proof}
Posons en effet $X_1=\operatorname{Spec}(\sh{O}_{X,x})$;
compte tenu de \sref{IV.20.2.11} et \sref{IV.20.3.6}, le faisceau $\sh{M}_{X_1}^*$ est induit sur $X_1$ par $\sh{M}_X^*$, donc on peut se borner au cas où $X=X_1$, auquel cas, comme $x\in S$ et que $x$ est point maximal de $S$, on a nécessairement $S=\{x\}$.
Si l'on avait alors $\operatorname{prof}(\sh{O}_{X,x})\neq1$, on en conclurait, en vertu de \sref{IV.21.1.8}, que $D=0$, ce qui contredit la définition de $S$.
\end{proof}

\begin{proposition}[21.1.10]
\label{IV.21.1.10}
Soit $A$ un anneau local noethérien;
pour que $\operatorname{Div}(A)=0$, il faut et il suffit que $\operatorname{prof}(A)=0$ (autrement dit, que l'idéal maximal $\mathfrak{m}$ de $A$ soit associé à $A$ (\hyperref[0.16.4.6]{$\mathbf{0}$, 16.4.6})).
\end{proposition}

\begin{proof}
En effet, dire que $\operatorname{Div}(A)=0$ signifie que dans $A$ tout élément régulier est inversible, ou encore que tous les éléments de $\mathfrak{m}$ sont diviseurs de zéro, ce qui signifie que $\mathfrak{m}\in\operatorname{Ass}(A)$ (Bourbaki, \emph{Alg. comm.}, chap.~IV, \textsection~1, n\textsuperscript{o}~1, cor.~3 de la prop.~2).
\end{proof}

\subsection{Diviseurs et Idéaux fractionnaires inversibles}
\EGAReaderTarget{ega4.c04.subsec.21-2}{subsection:IV.21.2}

\begin{env}[21.2.1]
\EGAReaderTarget{ega4.c04.statement.21-2-1}{IV.21.2.1}
Soit $(X,\sh{O}_X)$ un espace annelé.
On appelle \emph{Idéal fractionnaire sur $X$} un sous-$\sh{O}_X$-Module du $\sh{O}_X$-Module $\sh{M}_X$ des germes de fonctions méromorphes sur $X$.
Un Idéal fractionnaire $\sh{J}$ sur $X$ qui est un $\sh{O}_X$-Module inversible est appelé \emph{Idéal fractionnaire inversible}.
\end{env}

\begin{proposition}[21.2.2]
\label{IV.21.2.2}
Pour qu'un Idéal fractionnaire $\sh{J}$ sur $X$ soit inversible, il faut et il suffit que pour tout $x\in X$, il existe un voisinage ouvert $U$ de $x$ et une section $f\in\Gamma(U,\sh{M}_X^*)$ tels que $\sh{J}|U=\sh{O}_U.f$.
\end{proposition}

\begin{proof}
La condition est évidemment suffisante, l'application $s\mapsto s(f|V)$ de $\Gamma(V,\sh{O}_X)$ dans $\Gamma(V,\sh{J})$ étant évidemment bijective pour tout ouvert $V\subset U$.
Pour voir qu'elle est nécessaire, notons qu'il existe par hypothèse un voisinage ouvert $U$ de $x$ et un isomorphisme de $\sh{O}_X$-Modules $\sh{O}_U\xrightarrow{\sim}\sh{J}|U$.
Si $f$ est l'image de la section $1\in\Gamma(U,\sh{O}_X)$ par cet isomorphisme, on peut supposer, en restreignant $U$, que l'on a $f=u/s$, où $u\in\Gamma(U,\sh{O}_X)$ et $s\in\Gamma(U,\sh{S}(\sh{O}_X))$, et l'isomorphisme considéré fait alors correspondre, à toute section $v\in\Gamma(V,\sh{O}_X)$ (où $V$ est un ouvert contenu dans $U$), la section $v(u|V)/(s|V)$;
dire que l'application ainsi définie est bijective signifie que $u|V$ est un élément régulier de $\Gamma(V,\sh{O}_X)$, donc $f\in\Gamma(U,\sh{M}_X^*)$.
\end{proof}

On notera que pour tout ouvert $U$ de $X$ tel que $\sh{J}|U=\sh{O}_U.f$ avec $f\in\Gamma(U,\sh{M}_X^*)$, la section $f$ est déterminée \emph{à la multiplication près} par un élément de $\Gamma(U,\sh{O}_X^*)$, puisque la
\oldpage[IV-4]{259}
multiplication par ces éléments fournit tous les automorphismes du $\sh{O}_U$-Module $\sh{O}_U$ (\hyperref[0.5.4.3]{$\mathbf{0}_{\mathrm I}$, 5.4.3}).

\begin{corollary}[21.2.3]
\label{IV.21.2.3}
\begin{enumerate}[label=(\roman*)]
  \item Soit $\sh{J}$ un Idéal fractionnaire inversible;
  alors le $\sh{O}_X$-Module inversible $\sh{J}^{-1}$ s'identifie canoniquement à l'Idéal fractionnaire $\sh{J}'$ (\emph{transporteur} de $\sh{J}$ dans $\sh{O}_X$) défini de la façon suivante : pour tout ouvert $U$ de $X$ tel que $\sh{J}|U=\sh{O}_U.f$, où $f\in\Gamma(U,\sh{M}_X^*)$, on a $\sh{J}'|U=\sh{O}_U.f^{-1}$.
  \item Si $\sh{J}_1$ et $\sh{J}_2$ sont deux Idéaux fractionnaires inversibles, l'application canonique $\sh{J}_1\otimes\sh{J}_2\to\sh{J}_1\sh{J}_2$ est un isomorphisme de $\sh{O}_X$-Modules.
\end{enumerate}
\end{corollary}

\begin{proof}
L'assertion (ii) résulte aussitôt de \sref{IV.21.2.2}.
D'autre part, la remarque faite à la fin de \sref{IV.21.2.2} prouve qu'il existe bien un Idéal fractionnaire $\sh{J}'$ et un seul défini par la condition de l'énoncé;
l'isomorphisme canonique de $\sh{J}'$ sur $\sh{J}^{-1}=\shHom_{\sh{O}_X}(\sh{J},\sh{O}_X)$ s'obtient en faisant correspondre à toute section $s(f^{-1}|V)$ de $\Gamma(V,\sh{J}')$ (où $V$ est un ouvert contenu dans $U$ et $s\in\Gamma(V,\sh{O}_X)$) l'homomorphisme $t(f|V)\mapsto st$ de $\Gamma(V,\sh{J})$ dans $\Gamma(V,\sh{O}_X)$.
\end{proof}

En vertu de \sref{IV.21.2.3},~(i), on identifiera généralement les $\sh{O}_X$-Modules inversibles $\sh{J}'$ et $\sh{J}^{-1}$, considérant donc $\sh{J}^{-1}$ comme un sous-$\sh{O}_X$-Module de $\sh{M}_X$.

\begin{env}[21.2.4]
\label{IV.21.2.4}
Il résulte de \sref{IV.21.2.3} que l'ensemble $\operatorname{Id.inv}(X)$ des Idéaux fractionnaires inversibles sur $X$ est muni d'une structure de \emph{groupe commutatif} pour la loi de composition $(\sh{J}_1,\sh{J}_2)\mapsto\sh{J}_1\sh{J}_2$, l'élément neutre de ce groupe étant $\sh{O}_X$.
Il est clair que pour tout ouvert $U$ de $X$, on a $(\sh{J}_1\sh{J}_2)|U=(\sh{J}_1|U)(\sh{J}_2|U)$, donc l'application de restriction $\sh{J}\mapsto\sh{J}|U$ est un homomorphisme de groupes $\operatorname{Id.inv}(X)\to\operatorname{Id.inv}(U)$;
on définit donc ainsi un \emph{préfaisceau de groupes commutatifs} $U\mapsto\operatorname{Id.inv}(U)$;
il est immédiat qu'en fait ce préfaisceau est un \emph{faisceau de groupes commutatifs}, que l'on note $\sh{J}d.inv_X$.
\end{env}

\begin{env}[21.2.5]
\label{IV.21.2.5}
Pour toute fonction méromorphe régulière $f\in\Gamma(X,\sh{M}_X^*)$, il résulte de \sref{IV.21.2.2} que $\sh{J}(f)=\sh{O}_X.f$ est un Idéal fractionnaire inversible, et on a évidemment $\sh{J}(f_1f_2)=\sh{J}(f_1)\sh{J}(f_2)$, autrement dit l'application $f\mapsto\sh{J}(f)$ est un homomorphisme du groupe commutatif $\Gamma(X,\sh{M}_X^*)$ dans le groupe commutatif $\operatorname{Id.inv}(X)$.
Remplaçant $X$ par un ouvert quelconque $U$ de $X$ et notant que les homomorphismes obtenus sont compatibles avec les opérations de restriction, on obtient un homomorphisme canonique de faisceaux de groupes commutatifs :
\[
\label{IV.21.2.5.1}
  I_0:\sh{M}_X^*\to\sh{J}d.inv_X.
  \tag{21.2.5.1}
\]
Notons que si $f\in\Gamma(X,\sh{O}_X^*)$, on a $\sh{J}(f)=\sh{O}_X$;
on en déduit aussitôt que l'homomorphisme $I_0$ se factorise en
\[
\label{IV.21.2.5.2}
  \sh{M}_X^*\to\sh{M}_X^*/\sh{O}_X^*=\sh{D}iv_X\xrightarrow{I}\sh{J}d.inv_X
  \tag{21.2.5.2}
\]
où $I$ est un homomorphisme du faisceau de groupes additifs $\sh{D}iv_X$ dans le faisceau de groupes multiplicatifs $\sh{J}d.inv_X$;
on a donc pour tout ouvert $U$ de $X$ un homomorphisme $\sh{J}_U:\operatorname{Div}(U)\to\operatorname{Id.inv}(U)$ de groupes commutatifs, tel que, pour toute section $f\in\Gamma(U,\sh{M}_X^*)$, on ait
\[
\label{IV.21.2.5.3}
  \sh{J}_U(\operatorname{div}_U(f))=\sh{O}_U.f.
  \tag{21.2.5.3}
\]

\oldpage[IV-4]{260}

On en conclut que $\sh{J}_X(D)$, pour tout diviseur $D\in\operatorname{Div}(X)$, est l'Idéal fractionnaire inversible défini de la façon suivante : pour tout ouvert $U$ de $X$ tel que $D|U=\operatorname{div}_U(f)$, où $f\in\Gamma(U,\sh{M}_X^*)$, $\sh{J}_X(D)|U$ est l'Idéal fractionnaire inversible $\sh{O}_U.f$.
On a donc, en vertu de \sref{IV.21.1.6}, pour toute fonction méromorphe régulière $f\in\Gamma(X,\sh{M}_X^*)$, la relation
\[
\label{IV.21.2.5.4}
  f\in\Gamma(X,\sh{J}_X(D)) \iff \operatorname{div}(f)\geq D.
  \tag{21.2.5.4}
\]
\end{env}

\begin{proposition}[21.2.6]
\label{IV.21.2.6}
L'homomorphisme $I:\sh{D}iv_X\to\sh{J}d.inv_X$ est bijectif.
\end{proposition}

\begin{proof}
On définit en effet un homomorphisme $I_X'$ de $\operatorname{Id.inv}(X)$ dans $\operatorname{Div}(X)$ en faisant correspondre à tout Idéal fractionnaire inversible $\sh{J}$ sur $X$ le diviseur $I_X'(\sh{J})$ défini de la façon suivante : pour tout ouvert $U$ de $X$ tel que $\sh{J}|U=\sh{O}_U.f$, où $f\in\Gamma(U,\sh{M}_X^*)$ \sref{IV.21.2.2}, on prend $I_X'(\sh{J})|U=\operatorname{div}_U(f)$;
en vertu de la remarque suivant \sref{IV.21.2.2}, cette définition est indépendante du générateur $f$ choisi dans $\sh{J}|U$, et détermine bien un diviseur sur $X$.
En outre, cette définition montre aussitôt que les homomorphismes $\sh{J}_X$ et $I_X'$ sont réciproques l'un de l'autre.
Remplaçant $X$ par un ouvert quelconque $U$, on en déduit la définition de l'isomorphisme $I':\sh{J}d.inv_X\to\sh{D}iv_X$ réciproque de $I$, d'où la proposition.
On posera $I_X'(\sh{J})=\operatorname{div}(\sh{J})$, de sorte que l'on a, pour toute fonction méromorphe régulière $f$ sur $X$
\[
\label{IV.21.2.6.1}
  \operatorname{div}(\sh{O}_X.f)=\operatorname{div}(f).
  \tag{21.2.6.1}
\]
\end{proof}

\begin{env}[21.2.7]
\EGAReaderTarget{ega4.c04.statement.21-2-7}{IV.21.2.7}
On identifiera souvent les faisceaux $\sh{D}iv_X$ et $\sh{J}d.inv_X$ (resp. les groupes $\operatorname{Div}(X)$ et $\operatorname{Id.inv}(X)$) au moyen des isomorphismes $I$ et $I'$ (resp. $\sh{J}_X$ et $I_X'$) précédents.
On notera que l'on a la relation
\[
\label{IV.21.2.7.1}
  D\geq0 \iff \sh{J}_X(D)\subset\sh{O}_X \quad\text{pour}\quad D\in\operatorname{Div}(X)
  \tag{21.2.7.1}
\]
comme il résulte aussitôt des définitions \sref{IV.21.1.6} et de \sref{IV.21.1.5.1};
en d'autres termes, l'image $\sh{J}_X(\operatorname{Div}^+(X))$ est l'ensemble des \emph{Idéaux} de $\sh{O}_X$ (dits aussi parfois \emph{Idéaux entiers}) qui sont des $\sh{O}_X$-Modules inversibles : un tel Idéal $\sh{J}$ est encore caractérisé par le fait que pour tout $x\in X$, il y a un voisinage ouvert $U$ de $x$ dans $X$ tel que $\sh{J}|U=\sh{O}_U.f$, où $f$ est un élément régulier de $\Gamma(U,\sh{O}_X)$.
L'ensemble $\sh{J}_X(\operatorname{Div}^+(X))$ de ces Idéaux est donc un sous-monoïde de $\operatorname{Id.inv}(X)$, égal à l'ensemble des éléments positifs pour une relation d'ordre compatible avec la structure de groupe de $\operatorname{Id.inv}(X)$, et il est immédiat que cette relation n'est autre que la relation \emph{opposée} à l'inclusion;
autrement dit, on a
\[
\label{IV.21.2.7.2}
  \sh{J}_1\subset\sh{J}_2 \iff \operatorname{div}(\sh{J}_1)\geq\operatorname{div}(\sh{J}_2)
  \tag{21.2.7.2}
\]
pour $\sh{J}_1$, $\sh{J}_2$ dans $\operatorname{Id.inv}(X)$.
\end{env}

\begin{env}[21.2.8]
\label{IV.21.2.8}
Pour tout diviseur $D$ sur $X$, on pose
\[
\label{IV.21.2.8.1}
  \sh{O}_X(D)=(\sh{J}_X(D))^{-1};
  \tag{21.2.8.1}
\]
$\sh{O}_X(D)$ est donc un Idéal fractionnaire inversible, défini de la façon suivante : pour tout ouvert $U$ de $X$ tel que $D|U=\operatorname{div}_U(f)$, où $f\in\Gamma(U,\sh{M}_X^*)$, $\sh{O}_X(D)|U$ est l'Idéal
\oldpage[IV-4]{261}
fractionnaire inversible $\sh{O}_U.f^{-1}$;
en vertu de \sref{IV.21.1.6}, pour toute fonction méromorphe régulière $f$, on a la relation
\[
\label{IV.21.2.8.2}
  f\in\Gamma(X,\sh{O}_X(D)) \iff \operatorname{div}(f)\geq-D.
  \tag{21.2.8.2}
\]
En outre, il est clair que l'on a des isomorphismes canoniques \sref{IV.21.2.3}
\[
\label{IV.21.2.8.3}
  \left\{\begin{aligned}
    &\sh{O}_X(0)=\sh{O}_X,\quad \sh{O}_X(D+D')=\sh{O}_X(D)\sh{O}_X(D')\xrightarrow{\sim}\sh{O}_X(D)\otimes\sh{O}_X(D')\\
    &\sh{O}_X(nD)=(\sh{O}_X(D))^n\xrightarrow{\sim}(\sh{O}_X(D))^{\otimes n}
  \end{aligned}\right.
  \tag{21.2.8.3}
\]
pour tout entier $n\in\mathbf{Z}$, et pour deux diviseurs quelconques $D$, $D'$ sur $X$.
\end{env}

\begin{env}[21.2.9]
\label{IV.21.2.9}
Soit $\sh{J}$ un Idéal fractionnaire inversible sur $X$.
L'injection canonique $\sh{J}\hookrightarrow\sh{M}_X$ définit par tensorisation un homomorphisme de $\sh{O}_X$-Modules
\[
\label{IV.21.2.9.1}
  \sh{M}_X(\sh{J})=\sh{J}\otimes_{\sh{O}_X}\sh{M}_X\to\sh{M}_X\otimes_{\sh{O}_X}\sh{M}_X=\sh{M}_X
  \tag{21.2.9.1}
\]
qui est un \emph{isomorphisme} : en effet, si $U$ est un ouvert de $X$ tel que $\sh{J}|U=\sh{O}_U.f$, où $f\in\Gamma(U,\sh{M}_X^*)$, l'homomorphisme \sref{IV.21.2.9.1} restreint à $U$ n'est autre que l'isomorphisme qui à toute section $t$ de $\sh{M}_X(\sh{J})|U=\sh{M}_X|U$ au-dessus de $V\subset U$, fait correspondre la section $t(f|V)^{-1}$ du même faisceau.
Dans l'isomorphisme \sref{IV.21.2.9.1}, aux sections méromorphes régulières de $\sh{J}$ au-dessus de $X$ correspondent les fonctions méromorphes régulières sur $X$.

Considérons en particulier le cas où $\sh{J}=\sh{O}_X(D)$, où $D$ est un diviseur sur $X$;
on a alors un isomorphisme canonique
\[
\label{IV.21.2.9.2}
  \sh{M}_X(\sh{O}_X(D))\xrightarrow{\sim}\sh{M}_X
  \tag{21.2.9.2}
\]
et on désigne par $s_D$ la section méromorphe régulière de $\sh{O}_X(D)$ au-dessus de $X$ qui correspond par cet isomorphisme à la section $1$ de $\sh{M}_X$.
Si $U$ est un ouvert de $X$ tel que $\sh{O}_X(D)|U=\sh{O}_U.f^{-1}$, où $f\in\Gamma(U,\sh{M}_X^*)$, on a $s_D|U=1$ dans $\Gamma(U,\sh{M}_X)$;
comme alors on a $D|U=\operatorname{div}_U(f)$, on en déduit \sref{IV.21.1.4} que l'on a
\[
\label{IV.21.2.9.3}
  \operatorname{div}(s_D)=D.
  \tag{21.2.9.3}
\]
D'autre part, on déduit aussitôt des isomorphismes canoniques \sref{IV.21.2.8.3} les formules
\[
\label{IV.21.2.9.4}
  s_0=1,\quad s_{D+D'}=s_D\otimes s_{D'},\quad s_{nD}=s_D^{\otimes n}\quad(n\in\mathbf{Z}).
  \tag{21.2.9.4}
\]
\end{env}

\begin{env}[21.2.10]
\label{IV.21.2.10}
Considérons, entre deux couples $(\sh{L},s)$, $(\sh{L}',s')$, où $\sh{L}$ et $\sh{L}'$ sont deux $\sh{O}_X$-Modules inversibles, $s$ (resp. $s'$) une section méromorphe régulière de $\sh{L}$ (resp. $\sh{L}'$) au-dessus de $X$, la relation : « il existe un isomorphisme $u:\sh{L}\xrightarrow{\sim}\sh{L}'$ tel que $\overline{u}(s)=s'$ », où $\overline{u}:\Gamma(X,\sh{M}_X(\sh{L}))\xrightarrow{\sim}\Gamma(X,\sh{M}_X(\sh{L}'))$ est l'isomorphisme déduit de $u$ (on notera que l'isomorphisme $u$ vérifiant cette condition est alors déterminé de façon unique).
Il est clair que c'est une relation d'équivalence, et puisqu'il existe un ensemble de $\sh{O}_X$-Modules inversibles tel que tout $\sh{O}_X$-Module inversible soit isomorphe à un élément de cet ensemble (\hyperref[0.5.4.7]{$\mathbf{0}_{\mathrm I}$, 5.4.7}), on peut parler de l'ensemble $D(X)$ des classes d'équivalence des couples $(\sh{L},s)$ pour la relation précédente.
Pour tout $\sh{O}_X$-Module inversible $\sh{L}$ et toute section
\oldpage[IV-4]{262}
méromorphe régulière $s$ de $\sh{L}$ au-dessus de $X$, on notera $\operatorname{cl}(\sh{L},s)$ l'élément de $D(X)$ correspondant au couple $(\sh{L},s)$.
Il résulte de (\hyperref[0.5.4.3]{$\mathbf{0}_{\mathrm I}$, 5.4.3}) que si $s$, $s'$ sont deux sections méromorphes régulières de $\sh{L}$ au-dessus de $X$, la relation $\operatorname{cl}(\sh{L},s)=\operatorname{cl}(\sh{L},s')$ équivaut à l'existence d'une section $t\in\Gamma(X,\sh{O}_X^*)$ telle que $s'=ts$.

Il est immédiat que si $(\sh{L},s)$ est équivalent à $(\sh{L}_1,s_1)$ et $(\sh{L}',s')$ à $(\sh{L}_1',s_1')$, les couples $(\sh{L}\otimes\sh{L}',s\otimes s')$ et $(\sh{L}_1\otimes\sh{L}_1',s_1\otimes s_1')$ sont équivalents;
on définit donc dans $D(X)$ une loi de composition en posant
\[
\label{IV.21.2.10.1}
  \operatorname{cl}(\sh{L},s)\operatorname{cl}(\sh{L}',s')=\operatorname{cl}(\sh{L}\otimes\sh{L}',s\otimes s');
  \tag{21.2.10.1}
\]
il est immédiat que c'est une loi de groupe commutatif, dont l'élément neutre est $\operatorname{cl}(\sh{O}_X,1)$ et où l'inverse de $\operatorname{cl}(\sh{L},s)$ est $\operatorname{cl}(\sh{L}^{-1},s^{\otimes(-1)})$.
\end{env}

\begin{proposition}[21.2.11]
\EGAReaderTarget{ega4.c04.prop.21-2-11}{IV.21.2.11}
Les applications
\[
\label{IV.21.2.11.1}
  D\mapsto\operatorname{cl}(\sh{O}_X(D),s_D),\qquad \operatorname{cl}(\sh{L},s)\mapsto\operatorname{div}(s)
  \tag{21.2.11.1}
\]
sont des isomorphismes réciproques de $\operatorname{Div}(X)$ sur $D(X)$ et de $D(X)$ sur $\operatorname{Div}(X)$ respectivement.
\end{proposition}

\begin{proof}
Compte tenu de \sref{IV.21.2.8.3}, \sref{IV.21.2.9.4} et \sref{IV.21.2.10.1}, il suffit de voir que les composés de ces deux applications sont les identités dans $\operatorname{Div}(X)$ et $D(X)$ respectivement.
La première assertion n'est autre que \sref{IV.21.2.9.3}.
D'autre part, soit $D=\operatorname{div}(s)$, où $s$ est une section méromorphe régulière de $\sh{L}$ au-dessus de $X$, et soit $U$ un ouvert de $X$ tel qu'il existe un isomorphisme de $\sh{L}|U$ sur $\sh{O}_U$ transformant $s|U$ en $f\in\Gamma(U,\sh{M}_X^*)$, de sorte que $D|U=\operatorname{div}_U(f)$, $\sh{O}_X(D)|U=\sh{O}_U.f^{-1}$ et $s_D|U$ est l'élément unité de $\Gamma(U,\sh{M}_X)$.
Il y a donc un isomorphisme $v_U:\sh{L}|U\to\sh{O}_U.f^{-1}=\sh{O}_X(D)|U$ tel que $\overline{v}_U$ (notation de \sref{IV.21.2.10}) transforme $s|U$ en $f.f^{-1}=1$;
on voit aussitôt que ces isomorphismes sont compatibles avec les opérations de restriction, donc définissent un isomorphisme $v:\sh{L}\to\sh{O}_X(D)$ tel que $\overline{v}(s)=s_D$.
\end{proof}

On peut transporter par le premier des isomorphismes \sref{IV.21.2.11.1} la structure de groupe ordonné de $\operatorname{Div}(X)$ à $D(X)$;
les éléments $\geq0$ de $D(X)$ sont donc les classes $\operatorname{cl}(\sh{L},s)$ telles que $\operatorname{div}(s)\geq0$, c'est-à-dire \sref{IV.21.1.7.1} telles que
\[
  s\in\Gamma(X,\sh{L})\cap\Gamma(X,(\sh{M}_X(\sh{L}))^*).
\]

\begin{env}[21.2.12]
\label{IV.21.2.12}
Soit $D$ un diviseur \emph{positif} sur un \emph{préschéma} $X$;
l'Idéal fractionnaire $\sh{J}_X(D)$ est donc un Idéal de $\sh{O}_X$, qui est un $\sh{O}_X$-Module inversible;
soit $Y(D)$ le sous-préschéma fermé de $X$ qu'il définit.
Pour tout $x\in Y(D)$, il y a par hypothèse un voisinage ouvert $U$ de $x$ dans $X$ et une section régulière $t\in\Gamma(U,\sh{O}_X)$ telle que $\sh{J}_X(D)|U=\sh{O}_U.t$ \sref{IV.21.2.7};
en d'autres termes, l'immersion canonique $Y(D)\to X$ est \emph{régulière et de codimension $1$} en tout point de $Y(D)$ \sref{IV.19.1.4}.
Inversement, si $Y$ est un sous-préschéma fermé de $X$, régulièrement immergé dans $X$ et de codimension $1$ en tout point de $Y$, il existe un diviseur positif et un seul $D$ tel que $Y(D)=Y$, car pour tout $x\in Y$, il y a un voisinage $U$ de $x$ dans $X$ tel que $Y\cap U$ soit défini par un Idéal de $\sh{O}_U$ de la forme $\sh{O}_U.t$, où $t$ est régulier dans $\Gamma(U,\sh{O}_X)$.
On notera que l'on a alors $\operatorname{Supp}(D)=Y(D)$, car dire que $D_x\neq0$ signifie que $t_x$ (avec les notations ci-dessus) n'est pas inversible, c'est-à-dire que $x\in Y(D)$.
\end{env}

\oldpage[IV-4]{263}

\subsection{Équivalence linéaire des diviseurs}
\EGAReaderTarget{ega4.c04.subsec.21-3}{subsection:IV.21.3}

\begin{env}[21.3.1]
\EGAReaderTarget{ega4.c04.statement.21-3-1}{IV.21.3.1}
On dit qu'un diviseur $D$ sur $X$ est \emph{principal} s'il est de la forme $\operatorname{div}(f)$, où $f$ est une fonction méromorphe régulière sur $X$;
les fonctions méromorphes régulières $f'$ telles que $\operatorname{div}(f')=D$ sont alors toutes celles de la forme $tf$, où $t\in\Gamma(X,\sh{O}_X^*)$ \sref{IV.21.1.3.4}.
L'ensemble des diviseurs principaux est un sous-groupe de $\operatorname{Div}(X)$, noté $\operatorname{Div.princ}(X)$, isomorphe à $\Gamma(X,\sh{M}_X^*)/\Gamma(X,\sh{O}_X^*)$.
Deux diviseurs $D$, $D'$ sont dits \emph{linéairement équivalents} si $D-D'$ est un diviseur principal;
les diviseurs principaux sont donc les diviseurs linéairement équivalents à $0$.
\end{env}

\begin{env}[21.3.2]
\label{IV.21.3.2}
Rappelons (\hyperref[0.5.4.7]{$\mathbf{0}_{\mathrm I}$, 5.4.7}) que l'on peut parler de l'ensemble des classes d'équivalence des $\sh{O}_X$-Modules inversibles pour la relation d'isomorphie;
on désigne cet ensemble par $\operatorname{Pic}(X)$, et pour tout $\sh{O}_X$-Module inversible $\sh{L}$, on note $\operatorname{cl}(\sh{L})$ la classe d'équivalence des $\sh{O}_X$-Modules isomorphes à $\sh{L}$;
en outre, $\operatorname{Pic}(X)$ est un groupe commutatif pour la multiplication définie par $\operatorname{cl}(\sh{L})\operatorname{cl}(\sh{L}')=\operatorname{cl}(\sh{L}\otimes\sh{L}')$.
Il est clair que l'application
\[
\label{IV.21.3.2.1}
  r:\operatorname{cl}(\sh{L},s)\mapsto\operatorname{cl}(\sh{L})
  \tag{21.3.2.1}
\]
est un homomorphisme du groupe $D(X)$ \sref{IV.21.2.10} dans le groupe $\operatorname{Pic}(X)$.
Par composition, on en déduit donc un homomorphisme
\[
\label{IV.21.3.2.2}
  l:\operatorname{Div}(X)\xrightarrow{\sim}D(X)\xrightarrow{r}\operatorname{Pic}(X)
  \tag{21.3.2.2}
\]
(aussi noté $l_X$) tel que, pour tout diviseur $D$, on ait
\[
\label{IV.21.3.2.3}
  l(D)=\operatorname{cl}(\sh{O}_X(D)).
  \tag{21.3.2.3}
\]

Notons enfin que, si $u:X'\to X$ est un morphisme d'espaces annelés, $\sh{L}_1$, $\sh{L}_2$ deux $\sh{O}_X$-Modules inversibles isomorphes, les $\sh{O}_{X'}$-Modules inversibles $u^*(\sh{L}_1)$ et $u^*(\sh{L}_2)$ (\hyperref[0.5.4.5]{$\mathbf{0}_{\mathrm I}$, 5.4.5}) sont isomorphes;
comme d'autre part, pour deux $\sh{O}_X$-Modules inversibles quelconques $\sh{L}_1$, $\sh{L}_2$, on a $u^*(\sh{L}_1\otimes\sh{L}_2)=u^*(\sh{L}_1)\otimes u^*(\sh{L}_2)$ à isomorphisme canonique près (\hyperref[0.4.3.3]{$\mathbf{0}_{\mathrm I}$, 4.3.3}), on voit que le morphisme $u$ définit canoniquement un homomorphisme de groupes commutatifs
\[
\label{IV.21.3.2.4}
  \operatorname{Pic}(u):\operatorname{Pic}(X)\to\operatorname{Pic}(X').
  \tag{21.3.2.4}
\]
\end{env}

\begin{proposition}[21.3.3]
\label{IV.21.3.3}
\begin{enumerate}[label=(\roman*)]
  \item Le noyau de l'homomorphisme canonique $l:\operatorname{Div}(X)\to\operatorname{Pic}(X)$ est le sous-groupe $\operatorname{Div.princ}(X)$, autrement dit, pour que $\sh{O}_X(D)$ et $\sh{O}_X(D')$ soient isomorphes, il faut et il suffit que $D$ et $D'$ soient linéairement équivalents.
  On a donc un homomorphisme injectif canonique
  \[
  \label{IV.21.3.3.1}
    \operatorname{Div}(X)/\operatorname{Div.princ}(X)\to\operatorname{Pic}(X)
    \tag{21.3.3.1}
  \]
  déduit de $l$.
  \item Pour qu'un $\sh{O}_X$-Module inversible $\sh{L}$ soit tel que $\operatorname{cl}(\sh{L})$ soit de la forme $l(D)$, ou encore pour que $\sh{L}$ soit isomorphe à un $\sh{O}_X$-Module de la forme $\sh{O}_X(D)$, il faut et il suffit qu'il existe une section méromorphe régulière de $\sh{L}$.
\end{enumerate}
\end{proposition}

\begin{proof}
\oldpage[IV-4]{264}
La proposition résulte aussitôt des définitions et de \sref{IV.21.2.10}.
\end{proof}

\begin{proposition}[21.3.4]
\label{IV.21.3.4}
Soit $X$ un préschéma vérifiant l'une des deux hypothèses suivantes :
\begin{enumerate}[label=\emph{\alph*)}]
  \item $X$ est localement noethérien et $\operatorname{Ass}(\sh{O}_X)$ est contenu dans un ouvert affine de $X$.
  \item $X$ est réduit et l'ensemble de ses composantes irréductibles est localement fini.
\end{enumerate}
Alors l'homomorphisme canonique $l:\operatorname{Div}(X)\to\operatorname{Pic}(X)$ est surjectif, et donne par passage au quotient un isomorphisme
\[
  \operatorname{Div}(X)/\operatorname{Div.princ}(X)\xrightarrow{\sim}\operatorname{Pic}(X).
\]
\end{proposition}

\begin{proof}
Il suffit de montrer que tout $\sh{O}_X$-Module inversible $\sh{L}$ admet une section méromorphe régulière au-dessus de $X$ \sref{IV.21.3.3}.
Dans les deux cas, il suffira, grâce à \sref{IV.20.2.11},~(ii), de définir une section $s$ de $(\sh{M}_X(\sh{L}))^*$ au-dessus d'un ouvert schématiquement dense $U$ de $X$, ou encore dans le cas a), d'un ouvert $U$ contenant $\operatorname{Ass}(\sh{O}_X)$ \sref{IV.20.2.13},~(iv).
En effet, soit alors $V$ un ouvert quelconque de $X$ tel que $\sh{L}|V$ soit isomorphe à $\sh{O}_V$, de sorte qu'il existe un isomorphisme de $\sh{M}_X(\sh{L})|V$ sur $\sh{M}_X|V$, transformant $s|(U\cap V)$ en une section $f_V$ de $\sh{M}_X^*$ au-dessus de $U\cap V$.
Comme $U\cap V$ est schématiquement dense dans $V$, il résulte de \sref{IV.20.2.11} qu'il existe une fonction méromorphe régulière et une seule $g_V$ dans $V$ telle que $g_V|(U\cap V)=f_V$, et cette section correspond donc par l'isomorphisme considéré à une section méromorphe régulière $u_V$ de $\sh{L}$ au-dessus de $V$ telle que $u_V|(U\cap V)=s|(U\cap V)$.
En outre, si $V'$ est un second ouvert de $X$ tel que $\sh{L}|V'$ soit isomorphe à $\sh{O}_{V'}$, les restrictions de $u_V$ et $u_{V'}$ à $V\cap V'$ sont égales, car elles correspondent par isomorphisme à deux fonctions méromorphes qui coïncident dans un ouvert schématiquement dense $U\cap V\cap V'$, et on conclut encore par \sref{IV.20.2.11};
les $u_V$ sont donc bien les restrictions d'une section de $(\sh{M}_X(\sh{L}))^*$ au-dessus de $X$.

Cela étant, dans le cas b), on prend pour chacun des points maximaux $x_\lambda$ de $X$ un ouvert $U_\lambda$ contenant $x_\lambda$, ne rencontrant aucune composante irréductible de $X$ distincte de $\overline{\{x_\lambda\}}$ et tel que $\sh{L}|U_\lambda$ soit isomorphe à $\sh{O}_{U_\lambda}$;
on prendra pour $s$ la section telle que $s|U_\lambda$ soit la section de $\sh{L}|U_\lambda$ qui correspond par l'isomorphisme précédent à la section unité de $\sh{O}_{U_\lambda}$.

Dans le cas a), on peut prendre pour $U$ par hypothèse un ouvert affine (donc noethérien), autrement dit, supposer que $X=\operatorname{Spec}(A)$, où $A$ est noethérien, et $\sh{L}=\widetilde{P}$, où $P$ est un $A$-module projectif de rang $1$.
Si $S$ est l'ensemble des éléments réguliers de $A$, on a $\Gamma(X,\sh{M}_X)=S^{-1}A$ \sref{IV.20.2.12} et $\Gamma(X,\sh{M}_X(\sh{L}))=S^{-1}P$;
or $S$ est l'ensemble des éléments n'appartenant à aucun des idéaux associés à $A$, donc $S^{-1}A$ est un anneau semi-local dont les idéaux maximaux proviennent des éléments maximaux de $\operatorname{Ass}(A)$, et $S^{-1}P$ est un $S^{-1}A$-module projectif de rang $1$, donc ici libre de rang $1$ (Bourbaki, \emph{Alg. comm.}, chap.~II, \textsection~5, n\textsuperscript{o}~3, prop.~5);
un élément formant une base de ce $S^{-1}A$-module est par suite \sref{IV.20.1.8} une section méromorphe régulière de $\sh{L}$ au-dessus de $X$.
\end{proof}

\begin{corollary}[21.3.5]
\label{IV.21.3.5}
Si $X$ est un préschéma noethérien, tel qu'il existe un $\sh{O}_X$-Module inversible ample (\textbf{II}, 4.5.3) (par exemple un préschéma quasi-projectif sur un spectre d'anneau noethérien (\textbf{II}, 5.3.1 et 4.6.6)), l'homomorphisme canonique $l:\operatorname{Div}(X)\to\operatorname{Pic}(X)$ est surjectif.
\end{corollary}

\begin{proof}
\oldpage[IV-4]{265}
En effet (\textbf{II}, 4.5.4), il existe alors un voisinage ouvert affine de l'ensemble fini $\operatorname{Ass}(\sh{O}_X)$.
\end{proof}

\begin{remark}[21.3.6]
\label{IV.21.3.6}
Rappelons (\hyperref[0.5.4.7]{$\mathbf{0}_{\mathrm I}$, 5.4.7}) que l'on a un isomorphisme canonique $\pi:\mathrm{H}^1(X,\sh{O}_X^*)\xrightarrow{\sim}\operatorname{Pic}(X)$ défini de la façon suivante.
On part d'un $1$-cocycle $(c_{\alpha\beta})$ à valeurs dans $\sh{O}_X^*$ correspondant à un recouvrement ouvert $(U_\alpha)$ de $X$, $c_{\alpha\beta}$ étant un élément de $\Gamma(U_\alpha\cap U_\beta,\sh{O}_X^*)$, et on lui associe la classe du $\sh{O}_X$-Module inversible obtenu en recollant les $\sh{O}_{U_\alpha}$ suivant les isomorphismes $\sh{O}_{U_\alpha}|(U_\alpha\cap U_\beta)\xrightarrow{\sim}\sh{O}_{U_\beta}|(U_\alpha\cap U_\beta)$ définis par la multiplication par $c_{\alpha\beta}$.
D'autre part, on déduit de la suite exacte de faisceaux de groupes commutatifs
\[
  1\to\sh{O}_X^*\to\sh{M}_X^*\to\sh{D}iv_X\to0
\]
l'homomorphisme bord de la suite exacte de cohomologie
\[
\label{IV.21.3.6.1}
  \partial:\operatorname{Div}(X)\to\mathrm{H}^1(X,\sh{O}_X^*).
  \tag{21.3.6.1}
\]
Montrons que l'homomorphisme composé
\[
  \operatorname{Div}(X)\xrightarrow{\partial}\mathrm{H}^1(X,\sh{O}_X^*)\xrightarrow{\pi}\operatorname{Pic}(X)
\]
n'est autre que l'homomorphisme $l$ défini dans \sref{IV.21.3.2.2}.
En effet on doit partir d'un diviseur $D$ et d'un recouvrement ouvert $(U_\alpha)$ de $X$ tel que $D|U_\alpha=\operatorname{div}_{U_\alpha}(g_\alpha)$, où $g_\alpha$ est une fonction méromorphe régulière au-dessus de $U_\alpha$;
$\partial(D)$ est la classe de cohomologie du cocycle $(c_{\alpha\beta})$, où $c_{\alpha\beta}=g_{\alpha|\beta}g_{\beta|\alpha}^{-1}$, $g_{\alpha|\beta}$ désignant la restriction de $g_\alpha$ à $U_\alpha\cap U_\beta$.
Il est clair que l'image par $\pi$ de cette classe de cohomologie est la classe de l'Idéal fractionnaire inversible $\sh{L}$ tel que pour tout $\alpha$, $\sh{L}|U_\alpha=\sh{O}_{U_\alpha}.g_\alpha^{-1}$, qui n'est autre par définition que $\sh{O}_X(D)$ \sref{IV.21.2.8}.
\end{remark}

\subsection{Images réciproques de diviseurs}
\EGAReaderTarget{ega4.c04.subsec.21-4}{subsection:IV.21.4}
\label{IV.21.4}

\begin{env}[21.4.1]
\EGAReaderTarget{ega4.c04.statement.21-4-1}{IV.21.4.1}
Soit $f:X'\to X$ un morphisme d'espaces annelés;
proposons-nous de donner des conditions permettant d'associer à un diviseur $D$ sur $X$ un diviseur $D'$ sur $X'$, \emph{image réciproque} de $D$ par $f$.
Notons d'abord pour cela que pour toute section $t\in\Gamma(X,\sh{O}_X^*)$, l'image de $t$ par l'homomorphisme canonique $\Gamma(X,\sh{O}_X)\to\Gamma(X',\sh{O}_{X'})$ est encore inversible, autrement dit appartient à $\Gamma(X',\sh{O}_{X'}^*)$.
Considérons alors $D$ comme donné par la classe d'équivalence d'un couple $(\sh{L},s)$, où $\sh{L}$ est un $\sh{O}_X$-Module inversible et $s$ une section méromorphe régulière de $\sh{L}$ au-dessus de $X$ \sref{IV.21.2.11}.
Formons le $\sh{O}_{X'}$-Module inversible $f^*(\sh{L})=\sh{L}'$;
dire que l'image réciproque $s\circ f$ de $s$ par $f$ existe \sref{IV.20.1.11} et est une section méromorphe régulière de $\sh{L}'$ au-dessus de $X'$ revient à dire que les images réciproques par $f$ de $s$ et de $s^{\otimes(-1)}$ existent, autrement dit que $s\in\Gamma(X,\sh{M}_f(\sh{L}))$ et $s^{\otimes(-1)}\in\Gamma(X,\sh{M}_f(\sh{L}^{-1}))$;
la remarque faite ci-dessus montre alors que si $(\sh{L}_1,s_1)$ est un couple équivalent à $(\sh{L},s)$ au sens de \sref{IV.21.2.10}, l'image réciproque $s_1\circ f$ existe et est une section méromorphe régulière de $\sh{L}_1'=f^*(\sh{L}_1)$ au-dessus de $X'$, et les couples $(\sh{L}',s\circ f)$ et $(\sh{L}_1',s_1\circ f)$ sont équivalents.
On peut donc poser la définition suivante :
\end{env}

\begin{definition}[21.4.2]
\EGAReaderTarget{ega4.c04.def.21-4-2}{IV.21.4.2}
Étant donné un morphisme $f:X'\to X$ d'espaces annelés, on dit que l'\emph{image réciproque par $f$} d'un diviseur $D$ sur $X$ existe si l'on a $s_D\in\Gamma(X,\sh{M}_f(\sh{O}_X(D)))$ et $s_{-D}\in\Gamma(X,\sh{M}_f(\sh{O}_X(-D)))$ (cf. \sref{IV.20.1.11}).
On appelle alors \emph{image réciproque} de $D$ par $f$ et l'on note $f^*(D)$ le diviseur sur $X'$ qui correspond canoniquement \sref{IV.21.2.11} à la classe du couple $(f^*(\sh{O}_X(D)),s_D\circ f)$.
\end{definition}

Il résulte aussitôt de cette définition que si $D$ et $D'$ ont des images réciproques par $f$, il en est de même de $-D$ et de $D+D'$ (compte tenu de \sref{IV.21.2.9.4}) et que l'on a
\oldpage[IV-4]{266}
$f^*(D+D')=f^*(D)+f^*(D')$.
Autrement dit, l'ensemble $\operatorname{Div}^f(X)$ des diviseurs sur $X$ dont l'image réciproque par $f$ existe est un sous-groupe de $\operatorname{Div}(X)$, et l'application $D\mapsto f^*(D)$ est un homomorphisme croissant du sous-groupe ordonné $\operatorname{Div}^f(X)$ dans le groupe ordonné $\operatorname{Div}(X')$, rendant commutatif le diagramme
\[
\label{IV.21.4.2.1}
  \xymatrix{
    \operatorname{Div}^f(X) \ar[r]^{l_X} \ar[d]_{f^*} & \operatorname{Pic}(X) \ar[d]^{\operatorname{Pic}(f)} \\
    \operatorname{Div}(X') \ar[r]^{l_{X'}} & \operatorname{Pic}(X')
  }
  \tag{21.4.2.1}
\]

\begin{env}[21.4.3]
\label{IV.21.4.3}
La définition \sref{IV.21.4.2} montre aussitôt que, pour que $f^*(D)$ existe, il faut et il suffit que pour tout ouvert $U$ de $X$, l'image réciproque par $f^U:f^{-1}(U)\to U$ (restriction de $f$) de $D|U$ existe.
Or, si $D=\operatorname{div}(g)$, où $g$ est une fonction méromorphe régulière sur $X$, dire que l'image réciproque de $s_D$ par $f$ existe et est une section méromorphe régulière de $f^*(\sh{O}_X(D))$ signifie \sref{IV.21.2.9} que l'image réciproque de $g$ par $f$ existe et est une fonction méromorphe régulière sur $X'$.
On en déduit aussitôt une seconde description de $\operatorname{Div}^f(X)$ et de $f^*(D)$ : on considère le sous-faisceau de groupes de $\sh{M}_X^*$, noté $\sh{M}_f^*$, formé des germes de fonctions méromorphes régulières sur un ouvert de $X$ et dont l'image réciproque par $f$ existe et est régulière sur l'ouvert image réciproque \sref{IV.20.1.11}.
Alors si $f=(\psi,\theta)$, l'homomorphisme canonique \sref{IV.20.1.11} $\psi^*(\sh{M}_f)\to\sh{M}_{X'}$ donne par restriction un homomorphisme de faisceaux de groupes $\psi^*(\sh{M}_f^*)\to\sh{M}_{X'}^*$.
Posant $\sh{D}iv_X'=\sh{M}_f^*/\sh{O}_X^*$, on a $\operatorname{Div}^f(X)=\Gamma(X,\sh{D}iv_X')$, et l'application $D\mapsto f^*(D)$ correspond à l'homomorphisme de faisceaux de groupes $\psi^*(\sh{M}_f^*)/\psi^*(\sh{O}_X^*)\to\sh{M}_{X'}^*/\sh{O}_{X'}^*=\sh{D}iv_{X'}$ déduit du précédent par passage aux quotients.
\end{env}

\begin{env}[21.4.4]
\EGAReaderTarget{ega4.c04.statement.21-4-4}{IV.21.4.4}
Il résulte aussitôt des définitions précédentes que si $f':X''\to X'$ est un second morphisme d'espaces annelés, $D$ un diviseur sur $X$ tel que les images réciproques $f^*(D)$ et $f'^*(f^*(D))$ existent, alors $(f\circ f')^*(D)$ existe et est égal à $f'^*(f^*(D))$.
\end{env}

\begin{proposition}[21.4.5]
\label{IV.21.4.5}
Soit $f:X'\to X$ un morphisme d'espaces annelés.
Dans chacun des trois cas suivants, l'image réciproque par $f$ de tout diviseur sur $X$ est définie :
\begin{enumerate}[label=(\roman*)]
  \item $f$ est plat.
  \item $X$ et $X'$ sont des préschémas localement noethériens et l'on a $f(\operatorname{Ass}(\sh{O}_{X'}))\subset\operatorname{Ass}(\sh{O}_X)$.
  \item $X$ et $X'$ sont des préschémas, l'ensemble des composantes irréductibles de $X$ est localement fini, $X'$ est réduit et toute composante irréductible de $X'$ domine une composante irréductible de $X$.
\end{enumerate}
\end{proposition}

\begin{proof}
Il suffit en effet de montrer que dans les trois cas on a $\sh{M}_f=\sh{M}_X$;
dans le cas (i), cela résulte de \sref{IV.20.1.12}.
Dans le cas (iii), on peut se borner au cas où $X=\operatorname{Spec}(A)$ et $X'=\operatorname{Spec}(A')$ sont affines;
si $s\in A$ est régulier, il n'appartient à aucun idéal premier minimal de $A$ \sref{IV.20.1.3.1}, donc l'hypothèse entraîne que son image dans $A'$ n'appartient à aucun idéal premier minimal de $A'$, et est par suite un élément régulier de $A'$ \sref{IV.20.1.3.1}.
Dans le cas (ii) les fonctions méromorphes sur $X'$ s'identifient aux pseudo-fonctions sur $X'$ \sref{IV.20.2.11}, et l'hypothèse, jointe à \sref{IV.20.2.13},~(iv), assure que l'image réciproque
\oldpage[IV-4]{267}
par $f$ de tout ouvert schématiquement dense dans $X$ est un ouvert schématiquement dense dans $X'$;
on conclut donc par \sref{IV.20.3.12}.
\end{proof}

\begin{corollary}[21.4.6]
\label{IV.21.4.6}
Soit $X$ un préschéma ayant l'une des propriétés suivantes :
\begin{enumerate}[label=(\roman*)]
  \item $X$ est localement noethérien.
  \item $X$ est réduit et l'ensemble de ses composantes irréductibles est localement fini.
\end{enumerate}
Alors, pour tout $x\in X$, on a un isomorphisme canonique
\[
\label{IV.21.4.6.1}
  (\sh{D}iv_X)_x\xrightarrow{\sim}\operatorname{Div}(\sh{O}_{X,x}).
  \tag{21.4.6.1}
\]
\end{corollary}

\begin{proof}
Cela résulte en effet de \sref{IV.20.2.11}, \sref{IV.20.3.7} et de ce que $(\sh{O}_X^*)_x$ s'identifie au groupe des éléments inversibles de l'anneau $\sh{O}_{X,x}$.
\end{proof}

\begin{env}[21.4.7]
\label{IV.21.4.7}
Soient $X$, $X'$ deux préschémas, $f:X'\to X$ un morphisme.
Si $D$ est un diviseur \emph{positif} sur $X$ tel que l'image réciproque $f^*(D)$ soit définie \sref{IV.21.4.2}, alors le sous-préschéma fermé $Y(f^*(D))$ de $X'$ n'est autre que l'image réciproque $f^{-1}(Y(D))$;
cela résulte aussitôt des définitions \sref{IV.21.4.2} et \sref{IV.21.2.12}.
\end{env}

\begin{proposition}[21.4.8]
\EGAReaderTarget{ega4.c04.prop.21-4-8}{IV.21.4.8}
Soient $X$, $Y$ deux préschémas;
$f:X\to Y$ un morphisme fidèlement plat.
Alors, si un diviseur $D$ sur $Y$ est tel que $f^*(D)\geq0$ (l'existence de $f^*(D)$ résultant de \sref{IV.21.4.5}), on a $D\geq0$.
En particulier, l'application $D\mapsto f^*(D)$ de $\operatorname{Div}(Y)$ dans $\operatorname{Div}(X)$ est injective.
\end{proposition}

\begin{proof}
La question étant locale sur $Y$, on peut se borner au cas où $D=\operatorname{div}(w)$, avec $w=uv^{-1}$, $u$ et $v$ étant deux sections régulières de $\sh{O}_Y$ au-dessus de $Y$.
Par hypothèse on a $u\sh{O}_X\subset v\sh{O}_X$, donc, pour tout $x\in X$, si l'on pose $y=f(x)$, on a $u_y\sh{O}_{X,x}\subset v_y\sh{O}_{X,x}$;
on en conclut que $u_y\sh{O}_{Y,y}\subset v_y\sh{O}_{Y,y}$ en vertu de l'hypothèse que $\sh{O}_{X,x}$ est un $\sh{O}_{Y,y}$-module fidèlement plat et de Bourbaki, \emph{Alg. comm.}, chap.~I, \textsection~3, n\textsuperscript{o}~5, prop.~10,~(ii);
d'où $u\sh{O}_Y\subset v\sh{O}_Y$, puisque $f$ est surjectif, et par suite $D\geq0$.
\end{proof}

\subsection{Images directes de diviseurs}
\EGAReaderTarget{ega4.c04.subsec.21-5}{subsection:IV.21.5}
\label{IV.21.5}

\begin{env}[21.5.1]
\label{IV.21.5.1}
Soient $X$, $X'$ deux préschémas, $f:X'\to X$ un morphisme.
Nous allons, dans ce numéro, donner des conditions suffisantes pour qu'on puisse associer à tout diviseur $D'$ sur $X'$ un diviseur $D$ sur $X$, \emph{image directe} de $D'$ par $f$.
Nous nous bornerons au cas où $f$ est un morphisme \emph{fini} (pour des conditions plus générales, voir le chapitre de ce Traité consacré à la théorie des intersections).
\end{env}

\begin{lemma}[21.5.2]
\label{IV.21.5.2}
Soient $A$ un anneau, $E$ un $A$-module libre de rang fini.
Pour qu'un endomorphisme $u$ de $E$ soit injectif, il faut et il suffit que $\det(u)$ soit un élément régulier de $A$.
\end{lemma}

\begin{proof}
Cela est prouvé dans Bourbaki, \emph{Alg.}, chap.~III, 3\textsuperscript{e}~éd., \textsection~8, n\textsuperscript{o}~2, prop.~3.
\end{proof}

\begin{env}[21.5.3]
\EGAReaderTarget{ega4.c04.statement.21-5-3}{IV.21.5.3}
Supposons maintenant que le morphisme $f:X'\to X$ soit fini, et en outre que $f$ vérifie l'une des deux propriétés suivantes :
\begin{enumerate}[label=\Roman*)]
  \item $f$ est un morphisme fini \emph{localement libre}, autrement dit \sref{IV.18.2.7} le $\sh{O}_X$-Module quasi-cohérent de type fini $f_*(\sh{O}_{X'})$ est localement libre.
  \item $X$ est un préschéma localement noethérien réduit, le $\sh{R}(X)$-Module $f_*(\sh{O}_{X'})\otimes_{\sh{O}_X}\sh{R}(X)$ est localement libre, et pour toute section $s'\in\Gamma(f^{-1}(U),\sh{O}_{X'})$ ($U$ ouvert
  \oldpage[IV-4]{268}
  dans $X$), $\mathrm{N}_{f_*(\sh{O}_{X'})|\sh{O}_X}(s')$ est une section de $\sh{O}_X$ au-dessus de $U$ (cf. (\textbf{II}, 6.5.1)).
  (On rappelle que la condition (II) est vérifiée pour tout morphisme fini $f$ lorsque $X$ est un préschéma localement noethérien normal (\emph{loc. cit.}).)
\end{enumerate}

On sait alors (\textbf{II}, 6.5.5) que pour tout $\sh{O}_{X'}$-Module inversible $\sh{L}'$ on définit la norme $\sh{L}=\mathrm{N}_{X'/X}(\sh{L}')$ (que nous écrirons aussi $\mathrm{N}(\sh{L}')$), qui est un $\sh{O}_X$-Module inversible.
En outre, pour tout ouvert $U$ de $X$ et toute section régulière $s'\in\Gamma(f^{-1}(U),\sh{O}_{X'})$, la norme $\mathrm{N}_{X'/X}(s')=\mathrm{N}_{f_*(\sh{O}_{X'})|\sh{O}_X}(s')$ (que nous écrirons aussi $\mathrm{N}(s')$) est un élément régulier de $\Gamma(U,\sh{O}_X)$;
on est en effet ramené aussitôt au cas où $U=X$ est affine, et alors la conclusion résulte de \sref{IV.21.5.2} dans l'hypothèse (I);
d'autre part, dans l'hypothèse (II), le fait que $\sh{R}(X)$ soit un $\sh{O}_X$-Module plat entraîne que la section $s'\otimes1$ de $\Gamma(U,f_*(\sh{O}_{X'})\otimes_{\sh{O}_X}\sh{R}(X))$ est aussi régulière (\hyperref[0.6.1.2]{$\mathbf{0}_{\mathrm I}$, 6.1.2}), et la conclusion résulte encore de \sref{IV.21.5.2} appliqué à l'anneau $\Gamma(U,\sh{R}(X))$, compte tenu de la définition de la norme d'une section (\textbf{II}, 6.5.3).
Cela étant, soit $u'$ une section méromorphe de $\sh{L}'$ au-dessus de $X'$;
le morphisme $f$ étant affine, tout point $x\in X$ admet un voisinage ouvert $U$ tel que $u'|f^{-1}(U)$ soit de la forme $t'/s'$, où $t'\in\Gamma(f^{-1}(U),\sh{L}')$ et $s'$ est une section régulière dans $\Gamma(f^{-1}(U),\sh{O}_{X'})$;
l'élément $\mathrm{N}(t')/\mathrm{N}(s')$ (où $\mathrm{N}(t')$ est la section de $\sh{L}$ définie dans (\textbf{II}, 6.5.3)) est alors une section méromorphe de $\sh{L}$ au-dessus de $U$ en vertu de ce qui précède, et il résulte des propriétés multiplicatives de la norme (\textbf{II}, 6.5.3.1) que cette section ne dépend que de $u'|f^{-1}(U)$ et non de son écriture sous la forme $t'/s'$;
pour la même raison, lorsque $U$ varie, les sections méromorphes de $\sh{L}$ au-dessus de $U$ ainsi définies sont les restrictions d'une même section méromorphe de $\sh{L}$ au-dessus de $X$, que l'on appelle la \emph{norme} de $u'$ et que l'on note $\mathrm{N}_{X'/X}(u')$ (ou simplement $\mathrm{N}(u')$).
L'application ainsi définie
\[
\label{IV.21.5.3.1}
  \mathrm{N}_{X'/X}:\Gamma(X',\sh{M}_{X'}(\sh{L}'))\to\Gamma(X,\sh{M}_X(\mathrm{N}_{X'/X}(\sh{L}')))
  \tag{21.5.3.1}
\]
prolonge la norme définie dans (\textbf{II}, 6.5.5.3);
si $u'$ est une section méromorphe régulière de $\sh{L}'$ au-dessus de $X'$, il résulte aussitôt de ce qui précède que $\mathrm{N}(u')$ est une section méromorphe régulière de $\sh{L}$ au-dessus de $X$, car (avec les mêmes notations) $\mathrm{N}(t')$ est régulière si $t'$ l'est.
Enfin, si $\sh{L}_1'$, $\sh{L}_2'$ sont deux $\sh{O}_{X'}$-Modules inversibles, $s_1'$ (resp. $s_2'$) une section méromorphe régulière de $\sh{L}_1'$ (resp. $\sh{L}_2'$) au-dessus de $X'$, on a, en vertu de ce qui précède et de (\textbf{II}, 6.5.3.1)
\[
\label{IV.21.5.3.2}
  \mathrm{N}(s_1'\otimes s_2')=\mathrm{N}(s_1')\otimes\mathrm{N}(s_2').
  \tag{21.5.3.2}
\]
\end{env}

\begin{env}[21.5.4]
\label{IV.21.5.4}
Supposons toujours que $f$ vérifie l'une des hypothèses I), II) de \sref{IV.21.5.3}.
Si $(\sh{L}_1',s_1')$, $(\sh{L}_2',s_2')$ sont deux couples formés chacun d'un $\sh{O}_{X'}$-Module inversible et d'une section méromorphe régulière de ce Module au-dessus de $X'$, qui en outre sont équivalents au sens de \sref{IV.21.2.10}, alors les couples $(\mathrm{N}(\sh{L}_1'),\mathrm{N}(s_1'))$ et $(\mathrm{N}(\sh{L}_2'),\mathrm{N}(s_2'))$ sont aussi équivalents, car un isomorphisme de $\sh{O}_{X'}$-Modules inversibles a pour norme un isomorphisme de leurs normes (\textbf{II}, 6.5.3), et on a vu plus haut que $\mathrm{N}_{X'/X}$ transforme les sections de $\Gamma(X',\sh{O}_{X'}^*)$ en celles de $\Gamma(X,\sh{O}_X^*)$.
On peut donc poser la définition suivante :
\end{env}

\oldpage[IV-4]{269}

\begin{definition}[21.5.5]
\EGAReaderTarget{ega4.c04.def.21-5-5}{IV.21.5.5}
Étant donné un morphisme fini $f:X'\to X$ de préschémas, vérifiant l'une des conditions I), II) de \sref{IV.21.5.3}, on appelle \emph{image directe} (ou \emph{norme}) d'un diviseur $D'$ sur $X'$ par $f$, et l'on note $f_*(D')$ (ou $\mathrm{N}_{X'/X}(D')$), le diviseur sur $X$ qui correspond canoniquement \sref{IV.21.2.11} à la classe du couple $(\mathrm{N}_{X'/X}(\sh{O}_{X'}(D')),\mathrm{N}_{X'/X}(s_{D'}))$.
\end{definition}

Il résulte aussitôt de cette définition, compte tenu de \sref{IV.21.2.9.4} et de \sref{IV.21.5.3.2}, que si $D_1'$, $D_2'$, $D'$ sont des diviseurs sur $X'$, on a
\[
\label{IV.21.5.5.1}
  f_*(D_1'+D_2')=f_*(D_1')+f_*(D_2')
  \tag{21.5.5.1}
\]
et $D'\geq0$ entraîne $f_*(D')\geq0$, autrement dit, $D'\mapsto f_*(D')$ est un homomorphisme croissant du groupe ordonné $\operatorname{Div}(X')$ dans le groupe ordonné $\operatorname{Div}(X)$.
La définition \sref{IV.21.5.5} montre d'ailleurs aussitôt que pour tout ouvert $U$ de $X$, on a $f_*^U(D'|f^{-1}(U))=f_*(D')|U$ ($f^U$ étant la restriction $f^{-1}(U)\to U$ de $f$), et les homomorphismes $f_*^U$, pour $U$ variable, définissent donc un homomorphisme de faisceaux de groupes ordonnés
\[
\label{IV.21.5.5.2}
  \mathrm{N}_{X'/X}:f_*(\sh{D}iv_{X'})\to\sh{D}iv_X.
  \tag{21.5.5.2}
\]
En outre, pour tout ouvert $U$ de $X$, tout $\sh{O}_{X'}$-Module inversible $\sh{L}'$ et toute section méromorphe régulière $s'$ de $\sh{L}'$ au-dessus de $f^{-1}(U)$, on a, d'après les définitions précédentes et \sref{IV.21.1.4}
\[
\label{IV.21.5.5.3}
  \operatorname{div}_U(\mathrm{N}(s'))=f_*^U(\operatorname{div}_{f^{-1}(U)}(s')).
  \tag{21.5.5.3}
\]

\begin{proposition}[21.5.6]
\label{IV.21.5.6}
Soit $f:X'\to X$ un morphisme fini localement libre et supposons que $f_*(\sh{O}_{X'})$ soit de rang constant $n$.
Alors, pour tout diviseur $D$ sur $X$, $f_*(D)$ est défini et l'on a
\[
\label{IV.21.5.6.1}
  f_*(f^*(D))=nD.
  \tag{21.5.6.1}
\]
\end{proposition}

\begin{proof}
La première assertion résulte de ce que $f$ est plat \sref{IV.21.4.5}, et la seconde est conséquence immédiate des définitions et de (\textbf{II}, 6.5.3.2).
\end{proof}

\begin{proposition}[21.5.7]
\label{IV.21.5.7}
Soient $f:X'\to X$ un morphisme fini vérifiant l'une des hypothèses I), II) de \sref{IV.21.5.3}, $f':X''\to X'$ un morphisme fini localement libre de rang constant $n$.
Alors $f''=f\circ f':X''\to X$ vérifie la même hypothèse que $f$, et pour tout diviseur $D''$ sur $X''$, on a
\[
\label{IV.21.5.7.1}
  f_*''(D'')=f_*(f_*'(D'')).
  \tag{21.5.7.1}
\]
\end{proposition}

Compte tenu de la définition \sref{IV.21.5.5}, il suffit de prouver le résultat suivant :

\begin{lemma}[21.5.7.2]
\label{IV.21.5.7.2}
Sous les hypothèses de \sref{IV.21.5.7}, on a un isomorphisme fonctoriel
\[
\label{IV.21.5.7.3}
  \mathrm{N}_{X''/X}(\sh{L}'')\xrightarrow{\sim}\mathrm{N}_{X'/X}(\mathrm{N}_{X''/X'}(\sh{L}''))
  \tag{21.5.7.3}
\]
dans la catégorie des $\sh{O}_X$-Modules inversibles.
\end{lemma}

En effet, compte tenu de la définition de la norme d'une section de $\sh{L}''$ (\textbf{II}, 6.5.3) et de la définition \sref{IV.21.5.5}, on obtiendra aussitôt \sref{IV.21.5.7.1}.
Pour prouver \sref{IV.21.5.7.2}, il suffit, compte tenu des définitions de (\textbf{II}, 6.5.2 et 6.5.3), de prouver que pour toute section $s$ de $f_*''(\sh{O}_{X''})=f_*(f_*'(\sh{O}_{X''}))$ au-dessus d'un ouvert $U$ de $X$, on a
\[
\label{IV.21.5.7.4}
  \mathrm{N}_{f_*''(\sh{O}_{X''})|\sh{O}_X}(s)=\mathrm{N}_{f_*(\sh{O}_{X'})|\sh{O}_X}(\mathrm{N}_{f_*'(\sh{O}_{X''})|\sh{O}_{X'}}(s)).
  \tag{21.5.7.4}
\]

\oldpage[IV-4]{270}

La question est évidemment locale sur $X$, et on peut donc se borner au cas où $X=\operatorname{Spec}(A)$ est affine;
on a alors $X'=\operatorname{Spec}(A')$ et $X''=\operatorname{Spec}(A'')$, et on peut supposer que $A''$ est un $A'$-module projectif de rang $n$.
Lorsque $f$ est localement libre, on peut supposer que $A'$ est un $A$-module libre de rang $m$, et alors $A''$ est un $A$-module projectif de rang $mn$, et en restreignant $X$ à un ouvert convenable, on peut supposer que $A''$ est un $A$-module libre de rang $mn$;
la formule \sref{IV.21.5.7.4} résulte alors de la transitivité de la norme (Bourbaki, \emph{Alg.}, chap.~VIII, \textsection~12, n\textsuperscript{o}~2, prop.~7).
Lorsque $f$ vérifie l'hypothèse II), $A$ est noethérien réduit, et si $R$ est son anneau total de fractions, $A'\otimes_A R$ est un $R$-module libre de rang $m$, donc $A''\otimes_A R=A''\otimes_{A'}(A'\otimes_A R)$ est un $R$-module projectif de rang $mn$, et comme $R$ est alors un anneau semi-local, ce $R$-module est libre;
la proposition résulte alors encore de la transitivité des normes.

\begin{proposition}[21.5.8]
\label{IV.21.5.8}
Soient $f:X'\to X$ un morphisme fini, $g:Y\to X$ un morphisme;
on pose $Y'=X'\times_XY$, $f'=f_{(Y)}:Y'\to Y$, $g'=g_{(X')}:Y'\to X'$.
On suppose vérifiée l'une des conditions suivantes :
\begin{enumerate}[label=(\roman*)]
  \item $f$ est localement libre et $g$ est plat.
  \item $f$ est localement libre, $X$ et $Y$ sont localement noethériens, $g(\operatorname{Ass}(\sh{O}_Y))\subset\operatorname{Ass}(\sh{O}_X)$ et $g'(\operatorname{Ass}(\sh{O}_{Y'}))\subset\operatorname{Ass}(\sh{O}_{X'})$.
  \item $f$ vérifie l'hypothèse II) de \sref{IV.21.5.3}, $Y$ est localement noethérien, $Y$ et $Y'$ sont réduits et toute composante irréductible de $Y$ (resp. $Y'$) domine une composante irréductible de $X$ (resp. $X'$).
\end{enumerate}
Alors, pour tout diviseur $D'$ sur $X'$, $g'^*(D')$ est défini, $g^*(f_*(D'))$ est défini et l'on a
\[
\label{IV.21.5.8.1}
  g^*(f_*(D'))=f_*'(g'^*(D')).
  \tag{21.5.8.1}
\]
\end{proposition}

\begin{proof}
En effet, dans tous les cas, il résulte de (\textbf{II}, 6.5.8) que l'on a un isomorphisme fonctoriel
\[
  g^*(\mathrm{N}_{X'/X}(\sh{L}'))\xrightarrow{\sim}\mathrm{N}_{Y'/Y}(g'^*(\sh{L}'))
\]
dans la catégorie des $\sh{O}_Y$-Modules inversibles;
en outre (\textbf{II}, 6.5.4), si $s'$ est une section de $\sh{O}_{X'}$ au-dessus de $f^{-1}(U)$, $s''$ la section correspondante de $\sh{O}_{Y'}$ au-dessus de $g'^{-1}(f^{-1}(U))$ ($U$ ouvert de $X$), $\mathrm{N}_{Y'/Y}(s'')$ est la section de $\sh{O}_Y$ au-dessus de $g^{-1}(U)$ qui correspond à la section $\mathrm{N}_{X'/X}(s')$ de $\sh{O}_X$ au-dessus de $U$.
La formule \sref{IV.21.5.8.1} résultera donc des définitions si l'on prouve que $g'^*(D')$ et $g^*(D)$ sont définis, quels que soient les diviseurs $D'$ sur $X'$ et $D$ sur $X$.
En ce qui concerne $D$, cela résulte des hypothèses faites et de \sref{IV.21.4.5}.
En ce qui concerne $D'$, dans le cas (i) $g'$ est plat, donc dans tous les cas $g'^*(D')$ est défini en vertu de \sref{IV.21.4.5}.
\end{proof}

% FIN SOURCE ega4/ega4-21a-fr.tex

%
% DEBUT SOURCE ega4/ega4-21b-fr.tex
\subsection{Cycle 1-codimensionnel associé à un diviseur}
\EGAReaderTarget{ega4.c04.subsec.21-6}{subsection:IV.21.6}
\label{IV.21.6}

\begin{env}[21.6.1]
\EGAReaderTarget{ega4.c04.statement.21-6-1}{IV.21.6.1}
Soit $X$ un préschéma localement noethérien, et soit $\mathfrak{I}(X)$ l'ensemble des \emph{parties fermées irréductibles} de $X$ (qui est en correspondance biunivoque avec $X$ par l'application $x\mapsto\overline{\{x\}}$).
Dans le groupe produit $\bb{Z}^X$, considérons le sous-groupe $\mathfrak{Z}(X)$ des éléments $(n_x)_{x\in X}$ tels que l'ensemble des $\overline{\{x\}}\in\mathfrak{I}(X)$ tels que $n_x\neq0$ (ou, ce qui revient au même, l'ensemble des $x\in X$ tels que $n_x\neq0$) soit localement fini.
Il est clair que $\mathfrak{Z}(X)$ est un sous-groupe de $\bb{Z}^X$ qui contient le groupe somme directe $\bb{Z}^{(X)}$ (groupe libre ayant pour base $\mathfrak{I}(X)$), et lui est égal lorsque $X$ est noethérien.
Les éléments de $\mathfrak{Z}(X)$ sont appelés \emph{cycles} sur $X$ et ceux de $\mathfrak{I}(X)$ \emph{cycles premiers} (ils ne forment pas en général une base de $\mathfrak{Z}(X)$ lorsque $X$ n'est pas noethérien).
On considère toujours $\mathfrak{Z}(X)$ comme ordonné par l'ordre induit sur ce sous-groupe par l'ordre produit de $\bb{Z}^X$ et on notera $\mathfrak{Z}^+(X)$ l'ensemble des cycles $\geq0$.

\oldpage[IV-4]{271}

Pour tout cycle $Z\in\mathfrak{Z}(X)$, égal à $(n_x)_{x\in X}$, on écrira par abus de langage,
\[
  Z=\sum_{x\in X}n_x\cdot\overline{\{x\}}.
\]
On appelle \emph{multiplicité de $Z$ au point $x$} et on note $\operatorname{mult}_x(Z)$ l'entier $n_x$ (positif ou négatif).
Dire que $Z\geq0$ signifie que $\operatorname{mult}_x(Z)\geq0$ pour tout $x\in X$.
On appelle \emph{support de $Z$} et on note $\Supp(Z)$ la réunion des $\overline{\{x\}}$ tels que $\operatorname{mult}_x(Z)\neq0$;
par définition de $\mathfrak{Z}(X)$, c'est une partie fermée de $X$, comme réunion d'une famille localement finie de parties fermées.
On appelle \emph{dimension} (resp.\ \emph{codimension}) de $Z$ et on note $\dim(Z)$ (resp.\ $\codim(Z)$) la dimension (resp.\ codimension dans $X$) de $\Supp(Z)$.
\end{env}

\begin{env}[21.6.2]
\EGAReaderTarget{ega4.c04.statement.21-6-2}{IV.21.6.2}
On dit qu'une partie fermée $Y$ de $X$ est \emph{purement de codimension $p$} (dans $X$) si \emph{chacune} de ses composantes irréductibles est de codimension $p$ dans $X$.
On dit qu'un cycle est \emph{purement de codimension $p$}, ou (par abus de langage) est un \emph{cycle $p$-codimensionnel}, si son support est purement de codimension $p$.
On note $X^{(p)}$ l'ensemble des $x\in X$ tels que $\codim(\overline{\{x\}},X)=p$, ou, ce qui revient au même \sref{IV.5.1.2}, $\dim(\sh{O}_{X,x})=p$:
les cycles purement de codimension $p$ forment un sous-groupe $\mathfrak{Z}^p(X)$ de $\mathfrak{Z}(X)$, isomorphe au sous-groupe de $\bb{Z}^{X^{(p)}}$ formé des $(n_x)_{x\in X^{(p)}}$ tels que l'ensemble des $\overline{\{x\}}$ (ou l'ensemble des $x$) pour lesquels $n_x\neq0$ soit localement fini;
ce sous-groupe contient le groupe libre $\bb{Z}^{(X^{(p)})}$ et lui est identique lorsque $X$ est noethérien.
On note $\mathfrak{Z}^{p+}(X)$ l'ensemble des éléments $\geq0$ de $\mathfrak{Z}^p(X)$.
Il est clair que le groupe ordonné $\mathfrak{Z}(X)$ contient la somme directe des sous-groupes ordonnés $\mathfrak{Z}^p(X)$, et est identique à cette somme directe lorsque $X$ est noethérien;
dans ce dernier cas, on considère $\mathfrak{Z}(X)$ comme \emph{gradué} par les $\mathfrak{Z}^p(X)$ pour $p\in\bb{N}$.
\end{env}

\begin{env}[21.6.3]
\EGAReaderTarget{ega4.c04.statement.21-6-3}{IV.21.6.3}
Soient $Z=\sum_{x\in X}n_x\cdot\overline{\{x\}}$ un cycle sur $X$, $U$ un ouvert de $X$;
on appelle \emph{restriction} de $Z$ à $U$ et on note $Z|U$ le cycle $\sum_{x\in U}n_x\cdot(U\cap\overline{\{x\}})$ sur $U$;
on a $\Supp(Z|U)=\Supp(Z)\cap U$.
Il est clair que $Z\mapsto Z|U$ est un homomorphisme de groupes ordonnés de $\mathfrak{Z}(X)$ dans $\mathfrak{Z}(U)$ (resp.\ de $\mathfrak{Z}^p(X)$ dans $\mathfrak{Z}^p(U)$), et que si $V$ est un ouvert contenu dans $U$, on a $Z|V=(Z|U)|V$;
l'application $U\mapsto\mathfrak{Z}(U)$ (resp.\ $U\mapsto\mathfrak{Z}^p(U)$) est donc un \emph{préfaisceau} sur $X$ de groupes ordonnés commutatifs.
En fait ce préfaisceau est un \emph{faisceau}, somme directe des faisceaux $(i_x)_*(\bb{Z}_{\overline{\{x\}}})$, où $x$ parcourt $X$ (resp.\ $X^{(p)}$) et pour chaque $x\in X$, $i_x$ est l'injection canonique $\overline{\{x\}}\to X$ et $\bb{Z}_{\overline{\{x\}}}$ est le faisceau simple associé au faisceau constant $\bb{Z}$ sur l'espace $\overline{\{x\}}$:
cela résulte aussitôt de la description des sections d'une somme directe de faisceaux de groupes commutatifs (G, II, 2.7).
On note ce faisceau $\mathscr{Z}_X$ (resp.\ $\mathscr{Z}_X^p$).
On note $\mathscr{Z}_X^+$ (resp.\ $\mathscr{Z}_X^{p+}$) le sous-faisceau
\oldpage[IV-4]{272}
de monoïdes $U\mapsto\mathfrak{Z}^+(U)$ (resp.\ $U\mapsto\mathfrak{Z}^{p+}(U)$) de $\mathscr{Z}_X$ (resp.\ $\mathscr{Z}_X^p$).
Le faisceau $\mathscr{Z}_X$ est évidemment \emph{somme directe} des faisceaux $\mathscr{Z}_X^p$.

Notons enfin que les faisceaux $\mathscr{Z}_X^p$ (donc aussi $\mathscr{Z}_X$) sont \emph{flasques}.
Supposons en effet donnée une section $Z$ de $\mathscr{Z}_X^p$ au-dessus d'un ouvert $U$, de sorte que l'ensemble $S$ des $x\in U$ tels que $\operatorname{mult}_x(Z)\neq0$ est localement fini \emph{dans} $U$;
cet ensemble est aussi localement fini \emph{dans} $X$, car pour tout $z\in X$ et tout voisinage ouvert noethérien $V$ de $z$, $U\cap V$ est noethérien, donc ne contient qu'un nombre fini de points de $S$.
On définit donc une section $Z'$ de $\mathscr{Z}_X^p$ au-dessus de $X$ en posant $Z'=\sum_{x\in U}\operatorname{mult}_x(Z)\cdot\overline{\{x\}}$ et puisque l'adhérence de $x$ dans $U$ est $\overline{\{x\}}\cap U$, on a bien $Z'|U=Z$, d'où notre assertion.
\end{env}

\begin{env}[21.6.4]
\label{IV.21.6.4}
Nous nous proposons de définir un homomorphisme canonique
\[
\label{IV.21.6.4.1}
  c:\sh{D}iv_X\to\mathscr{Z}_X^1
  \tag{21.6.4.1}
\]
de faisceaux de groupes commutatifs.
Il suffit évidemment de définir un homomorphisme de $\sh{M}_X^*$ dans $\mathscr{Z}_X^1$, qui s'annule dans $\sh{O}_X^*$ \sref{IV.21.1.2};
or, par définition, $\sh{M}_X^*$ est le \emph{symétrisé} du faisceau de monoïdes $\sh{O}_X\cap\sh{M}_X^*$, et un homomorphisme $\sh{M}_X^*\to\mathscr{Z}_X^1$ est déterminé de façon unique par la donnée de sa restriction $\sh{S}(\sh{O}_X)\to\mathscr{Z}_X^1$ qui est un homomorphisme de faisceaux de \emph{monoïdes} soumis à la seule condition d'être nul dans $\sh{O}_X^*$;
il revient au même de dire que pour définir $c$, il suffit de définir un homomorphisme de \emph{faisceaux de monoïdes}
\[
\label{IV.21.6.4.2}
  c:\sh{D}iv_X^+\to\mathscr{Z}_X^{1+}.
  \tag{21.6.4.2}
\]
\end{env}

\begin{env}[21.6.5]
\label{IV.21.6.5}
Or, on a vu qu'à tout diviseur positif $D$ sur $X$ on fait correspondre le sous-préschéma fermé $Y(D)$ de $X$, défini par l'Idéal $\sh{J}_X(D)\subset\sh{O}_X$, et qui est régulièrement immergé et de codimension $1$.
En chacun des points maximaux $x$ de $Y(D)$, on a donc (\sref{IV.19.1.4} et \sref{IV.5.1.2}) $\codim(\overline{\{x\}},X)=\dim(\sh{O}_{X,x})=1$, autrement dit $x\in X^{(1)}$;
d'ailleurs l'ensemble de ces points maximaux est localement fini dans $X$ \sref{IV.3.1.6}, et $\sh{O}_{Y(D),x}$ est un anneau \emph{artinien}.
En tout point $x\in X^{(1)}$ qui n'est pas point maximal de $Y(D)$, on a nécessairement $x\notin Y(D)$, donc $\sh{O}_{Y(D),x}=0$.
Posons
\[
\label{IV.21.6.5.1}
  \operatorname{cyc}(D)=\sum_{x\in X^{(1)}}\operatorname{long}(\sh{O}_{Y(D),x})\cdot\overline{\{x\}}
  \tag{21.6.5.1}
\]
qui est donc un élément de $\mathfrak{Z}^1(X)$.
\end{env}

\begin{proposition}[21.6.6]
\label{IV.21.6.6}
L'application $D\mapsto\operatorname{cyc}(D)$ est un homomorphisme du monoïde $\operatorname{Div}^+(X)$ dans $\mathfrak{Z}^{1+}(X)$.
\end{proposition}

\begin{proof}
Tout revient à voir que pour deux diviseurs positifs $D$, $D'$, on a
\[
  \operatorname{cyc}(D+D')=\operatorname{cyc}(D)+\operatorname{cyc}(D').
\]
Or, on a \sref{IV.21.2.5} $\sh{J}_X(D+D')=\sh{J}_X(D)\sh{J}_X(D')$;
tout revient à voir que si $x\in X^{(1)}$, si l'on pose $A=\sh{O}_{X,x}$, et si $t$ et $t'$ sont deux éléments réguliers de $A$, on a $\operatorname{long}(A/tt'A)=\operatorname{long}(A/tA)+\operatorname{long}(A/t'A)$;
mais puisque $t$ est régulier, $tA/tt'A$ est isomorphe à $A/t'A$, d'où aussitôt la proposition.
\end{proof}

Il résulte aussitôt des définitions que pour tout ouvert $U\subset X$, on a
\[
  Y(D|U)=Y(D)\cap U,
\]
\oldpage[IV-4]{273}
donc $\operatorname{cyc}(D|U)=\operatorname{cyc}(D)|U$, et les homomorphismes $\operatorname{cyc}:\operatorname{Div}^+(U)\to\mathfrak{Z}^1(U)$ définissent donc un homomorphisme de faisceaux de monoïdes \sref{IV.21.6.4.2}, d'où l'homomorphisme cherché \sref{IV.21.6.4.1} de faisceaux de groupes.

On a $\Supp(\operatorname{cyc}(D))\subset\Supp(D)$ pour tout diviseur $D$ et
\begin{env}[21.6.6.2]
\label{IV.21.6.6.2}
\[
  \Supp(\operatorname{cyc}(D))=\Supp(D)\quad\text{lorsque}\quad D\geq0;
\]
\end{env}
on a déjà vu en effet la seconde relation \sref{IV.21.2.12};
lorsque $D$ est quelconque, la relation $D_x=0$ entraîne l'existence d'un voisinage ouvert $U$ de $x$ tel que $D|U=0$, d'où $\operatorname{cyc}(D)|U=0$, ce qui prouve notre assertion.

\begin{env}[21.6.7]
\EGAReaderTarget{ega4.c04.statement.21-6-7}{IV.21.6.7}
L'homomorphisme \sref{IV.21.6.4.1} donne, par passage aux groupes de sections au-dessus de $X$, un \emph{homomorphisme croissant de groupes ordonnés} $\operatorname{Div}(X)\to\mathfrak{Z}^1(X)$, que nous noterons encore $D\mapsto\operatorname{cyc}(D)$;
le nombre $\operatorname{mult}_x(\operatorname{cyc}(D))$ pour $x\in X^{(1)}$ se note encore $\operatorname{mult}_x(D)$ ou $\operatorname{mult}_{\overline{\{x\}}}(D)$ et s'appelle \emph{multiplicité du diviseur $D$ au point $x$}, ou \emph{multiplicité du cycle premier $\overline{\{x\}}$ dans $D$};
c'est un entier positif ou négatif, égal comme on l'a vu à $\operatorname{long}(\sh{O}_{Y(D),x})$ lorsque $D$ est un diviseur \emph{positif};
on a donc par définition
\[
\label{IV.21.6.7.1}
  \operatorname{cyc}(D)=\sum_{x\in X^{(1)}}\operatorname{mult}_x(D)\cdot\overline{\{x\}},
  \tag{21.6.7.1}
\]
et en vertu du fait que $D\mapsto\operatorname{cyc}(D)$ est un homomorphisme, on a
\[
\label{IV.21.6.7.2}
  \operatorname{mult}_x(-D)=-\operatorname{mult}_x(D),\qquad
  \operatorname{mult}_x(D+D')=\operatorname{mult}_x(D)+\operatorname{mult}_x(D')
  \tag{21.6.7.2}
\]
pour deux diviseurs quelconques $D$, $D'$.

Pour toute fonction méromorphe régulière $f$ sur $X$, on pose en particulier, pour $x\in X^{(1)}$
\[
\label{IV.21.6.7.3}
  o_x(f)=\operatorname{mult}_x(\operatorname{div}(f))
  \tag{21.6.7.3}
\]
et on dit que $o_x(f)$ (entier positif ou négatif) est \emph{l'ordre de $f$ au point} $x\in X^{(1)}$.
Si $f_x\in\sh{O}_{X,x}$ (élément régulier de cet anneau local par hypothèse), on a donc
\[
\label{IV.21.6.7.4}
  o_x(f)=\operatorname{long}(\sh{O}_{X,x}/f_x\sh{O}_{X,x});
  \tag{21.6.7.4}
\]
pour deux fonctions méromorphes régulières $f$, $f'$ sur $X$, on a
\[
\label{IV.21.6.7.5}
  o_x(ff')=o_x(f)+o_x(f'),\qquad o_x(f^{-1})=-o_x(f)
  \tag{21.6.7.5}
\]
pour tout $x\in X^{(1)}$.
Les cycles $1$-codimensionnels
\[
  Z^+(f)=\sum_{x\in X^{(1)}}(o_x(f))^+\cdot\overline{\{x\}},\qquad
  Z^-(f)=\sum_{x\in X^{(1)}}(o_x(f))^-\cdot\overline{\{x\}}
\]
sont appelés respectivement le \emph{cycle des zéros} et le \emph{cycle des pôles} (ou \emph{cycle polaire}) de $f$, et l'on a évidemment $\operatorname{cyc}(\operatorname{div}(f))=Z^+(f)-Z^-(f)$.
Les cycles $1$-codimensionnels de la forme $\operatorname{cyc}(\operatorname{div}(f))$ sont appelés \emph{principaux} (ou \emph{linéairement équivalents à zéro}) et forment un sous-groupe $\mathfrak{Z}^1_\mathrm{princ}(X)$ de $\mathfrak{Z}^1(X)$.
Les sections au-dessus de $X$ de l'image $c(\sh{D}iv_X)\subset\mathscr{Z}_X^1$ sont
\oldpage[IV-4]{274}
appelées \emph{cycles localement principaux};
un tel cycle $Z$ est donc caractérisé par le fait que, pour tout $y\in X$, il y a un voisinage ouvert $U$ de $y$ dans $X$ et une fonction méromorphe régulière $f$ sur $U$ telle que $Z|U=\sum_{x\in U\cap X^{(1)}}o_x(f)(\overline{\{x\}}\cap U)$.
Si $Z=\sum_{x\in X^{(1)}}n_x\cdot\overline{\{x\}}$, il revient au même de dire que, pour tout $y\in X$, si l'on pose $T_y=\Spec(\sh{O}_{X,y})$ et $Z_y=\sum_{x\in X^{(1)}\cap T_y}n_x(\overline{\{x\}}\cap T_y)$, $Z_y$ est \emph{principal}, autrement dit il existe une fonction méromorphe régulière $g$ sur $T_y$ telle que $Z_y=\operatorname{cyc}(\operatorname{div}(g))$.
Cela résulte aussitôt en effet de la définition précédente et de \sref{IV.20.3.7} qui établit une correspondance biunivoque entre les fonctions méromorphes régulières sur $T_y$ et les germes de fonctions méromorphes régulières sur les voisinages ouverts de $y$ dans $X$ lorsque $X$ est localement noethérien.
On exprime encore le fait que $Z_y$ est principal en disant que $Z$ est \emph{principal au point $y$}.
L'ensemble des points où $Z$ est principal est évidemment \emph{ouvert}, en vertu de ce qui précède.

On déduit de l'homomorphisme canonique $\operatorname{cyc}:\operatorname{Div}(X)\to\mathfrak{Z}^1(X)$ un homomorphisme $\operatorname{Div}(X)/\operatorname{Div.princ}(X)\to\mathfrak{Z}^1(X)/\mathfrak{Z}^1_\mathrm{princ}(X)$, en vertu de la définition de $\mathfrak{Z}^1_\mathrm{princ}(X)$.
On dit que le groupe $\mathfrak{Z}^1(X)/\mathfrak{Z}^1_\mathrm{princ}(X)$ est le \emph{groupe des classes de cycles $1$-codimensionnels} de $X$ et on le note $\mathfrak{Cl}(X)$.
Deux éléments de $\mathfrak{Z}^1(X)$ ayant même image dans $\mathfrak{Cl}(X)$ sont appelés \emph{linéairement équivalents}.
\end{env}

\begin{env}[21.6.8]
\label{IV.21.6.8}
Considérons en particulier le cas où $X=\Spec(A)$, où $A$ est un anneau \emph{noethérien intègre et intégralement clos}.
Alors $X^{(1)}$ est l'ensemble des \emph{idéaux premiers de hauteur $1$} de $A$, et $\mathfrak{Z}^1(X)$ s'identifie donc au \emph{groupe des diviseurs} (au sens de N.~Bourbaki) de l'anneau de Krull $A$ (Bourbaki, \emph{Alg.\ comm.}, chap.~VII, \textsection~1, n\textsuperscript{o}~3, cor.\ du th.~2 et n\textsuperscript{o}~6, th.~3).

Comme d'autre part, les fonctions méromorphes régulières sur $X$ s'identifient alors aux éléments non nuls du corps des fractions $K$ de $A$, l'application $f\mapsto\operatorname{cyc}(\operatorname{div}(f))$ de $M(X)$ dans $\mathfrak{Z}^1(X)$ s'identifie à l'application notée $f\mapsto\operatorname{div}(f)$ dans Bourbaki (\emph{loc.\ cit.}, \textsection~1, n\textsuperscript{o}~1);
$\mathfrak{Z}^1_\mathrm{princ}(X)$ s'identifie donc au \emph{groupe des diviseurs principaux} de $A$ au sens de Bourbaki, et $\mathfrak{Cl}(X)$ au \emph{groupe des classes de diviseurs} de $A$ au sens de Bourbaki (\emph{loc.\ cit.}, \textsection~1, n\textsuperscript{o}~2 et n\textsuperscript{o}~10).
\end{env}

\begin{theorem}[21.6.9]
\EGAReaderTarget{ega4.c04.thm.21-6-9}{IV.21.6.9}
Soit $X$ un préschéma localement noethérien normal.
\begin{enumerate}[label=(\roman*)]
  \item L'homomorphisme canonique $\operatorname{cyc}:\operatorname{Div}(X)\to\mathfrak{Z}^1(X)$ est injectif et son image est formée des cycles localement principaux.
  \item Les conditions suivantes sont équivalentes:
  \begin{enumerate}[label=\emph{\alph*)}]
    \item L'homomorphisme $\operatorname{cyc}:\operatorname{Div}(X)\to\mathfrak{Z}^1(X)$ est bijectif.
    \item Tout cycle $1$-codimensionnel est localement principal.
    \item Pour tout $x\in X$, l'anneau local $\sh{O}_{X,x}$ est factoriel.
  \end{enumerate}
  (On dit alors que $X$ est un préschéma \emph{localement factoriel}.)
\end{enumerate}
\end{theorem}

\begin{proof}
\begin{enumerate}[label=(\roman*)]
  \item Il suffit de démontrer que
\[
\label{IV.21.6.9.1}
  \operatorname{cyc}^{-1}(\mathfrak{Z}^{1+}(X))=\operatorname{Div}^+(X),
  \tag{21.6.9.1}
\]
puisque l'on a $\operatorname{Div}^+(X)\cap(-\operatorname{Div}^+(X))=0$ et $\mathfrak{Z}^{1+}(X)\cap(-\mathfrak{Z}^{1+}(X))=0$.
On est donc ramené à prouver que si $D$ est un diviseur sur $X$ tel que $\operatorname{mult}_x(D)\geq0$ pour tout $x\in X^{(1)}$, on a $D\geq0$.
Or, pour tout $x\notin X^{(1)}$, l'anneau local $\sh{O}_{X,x}$ est intègre et intégralement clos et
\oldpage[IV-4]{275}
de dimension $0$ ou $\geq2$, donc on a $\operatorname{prof}(\sh{O}_{X,x})=0$ ou $\operatorname{prof}(\sh{O}_{X,x})\geq2$ \sref[0]{0.16.5.1}.
D'autre part, aux points $x\in X^{(1)}$, l'anneau $\sh{O}_{X,x}$ est un anneau de valuation discrète (\sref[II]{II.7.1.6}), donc $\operatorname{prof}(\sh{O}_{X,x})=1$ \sref[0]{0.16.5.1};
les seuls points de $X$ tels que $\operatorname{prof}(\sh{O}_{X,x})=1$ sont donc les points de $X^{(1)}$, et la conclusion résulte par suite de \sref{IV.21.1.8}.

  \item L'équivalence de \emph{a)} et \emph{b)} est claire en vertu de (i).
D'après la caractérisation des cycles localement principaux donnée dans \sref{IV.21.6.7}, et la relation entre cycles $1$-codimensionnels sur $\Spec(A)$ et diviseurs (au sens de Bourbaki) de l'anneau $A$ lorsque $A$ est un anneau noethérien intégralement clos \sref{IV.21.6.8}, la condition \emph{b)} équivaut à dire que pour tout $x\in X$, tout diviseur de l'anneau $\sh{O}_{X,x}$ est principal, autrement dit que l'anneau $\sh{O}_{X,x}$ est factoriel (Bourbaki, \emph{Alg.\ comm.}, chap.~VII, \textsection~3, n\textsuperscript{o}~1), d'où l'équivalence de \emph{b)} et \emph{c)}.
\end{enumerate}
\end{proof}

\begin{env}[21.6.9.2]
\label{IV.21.6.9.2}
Lorsque $A$ est un anneau noethérien \emph{factoriel}, il est clair que $X=\Spec(A)$ est localement factoriel (Bourbaki, \emph{Alg.\ comm.}, chap.~VII, \textsection~3, n\textsuperscript{o}~4, prop.~3).
Si, dans ce cas, on écrit un élément $f\neq0$ du corps des fractions $K$ de $A$ sous la forme $r/s$, où $r$ et $s$ sont deux éléments \emph{étrangers} de $A$, dont les diviseurs sont bien déterminés (\emph{loc.\ cit.}, \textsection~3, n\textsuperscript{o}~3), ces diviseurs s'identifient respectivement au \emph{diviseur des zéros} et au \emph{diviseur des pôles} de $f$ \sref{IV.21.6.7}.
\end{env}

\begin{corollary}[21.6.10]
\label{IV.21.6.10}
Soit $X$ un préschéma localement noethérien normal.
\begin{enumerate}[label=(\roman*)]
  \item Il existe un homomorphisme canonique injectif
\[
\label{IV.21.6.10.1}
  \operatorname{Pic}(X)\to\mathfrak{Cl}(X).
  \tag{21.6.10.1}
\]

  \item Si $X$ est localement factoriel, l'homomorphisme \sref{IV.21.6.10.1} est bijectif, et réciproquement.
\end{enumerate}
\end{corollary}

\begin{proof}
On a vu \sref{IV.21.6.7} que l'image de $\operatorname{Div.princ}(X)$ par l'homomorphisme $\operatorname{cyc}:\operatorname{Div}(X)\to\mathfrak{Z}^1(X)$ est $\mathfrak{Z}^1_\mathrm{princ}(X)$;
on déduit donc de \sref{IV.21.6.9} que l'homomorphisme $\operatorname{Div}(X)/\operatorname{Div.princ}(X)\to\mathfrak{Z}^1(X)/\mathfrak{Z}^1_\mathrm{princ}(X)=\mathfrak{Cl}(X)$ déduit de $\operatorname{cyc}$ est injectif, et qu'il est bijectif si et seulement si $X$ est localement factoriel.
La conclusion résulte donc de ce que, lorsque $X$ est localement noethérien et réduit, l'homomorphisme canonique $\operatorname{Div}(X)/\operatorname{Div.princ}(X)\to\operatorname{Pic}(X)$ \sref{IV.21.3.3.1} est bijectif \sref{IV.21.3.4}, (ii).
\end{proof}

\begin{corollary}[21.6.11]
\label{IV.21.6.11}
Soit $X$ un préschéma localement noethérien et localement factoriel.
Alors le faisceau $\sh{D}iv_X$ est flasque, et pour tout ouvert $U$ de $X$, l'homomorphisme canonique $\operatorname{Pic}(X)\to\operatorname{Pic}(U)$ est surjectif.
\end{corollary}

\begin{proof}
Tenant compte de \sref{IV.21.6.9}, (ii), la première assertion équivaut à dire que le faisceau $\mathscr{Z}_X^1$ est flasque, ce qui a été vu plus haut \sref{IV.21.6.3}.
Pour tout ouvert $U$ de $X$, l'homomorphisme canonique $\mathfrak{Z}^1(X)\to\mathfrak{Z}^1(U)$ est donc surjectif;
comme en vertu de \sref{IV.21.6.10}, l'homomorphisme $\operatorname{Pic}(X)\to\operatorname{Pic}(U)$ s'identifie canoniquement à $\mathfrak{Z}^1(X)/\mathfrak{Z}^1_\mathrm{princ}(X)\to\mathfrak{Z}^1(U)/\mathfrak{Z}^1_\mathrm{princ}(U)$, il est aussi surjectif.
\end{proof}

\begin{proposition}[21.6.12]
\label{IV.21.6.12}
Soit $X$ un préschéma noethérien réduit.
Soit $(U_\lambda)_{\lambda\in L}$ une famille filtrante décroissante d'ouverts de $X$ vérifiant les conditions suivantes:
\begin{enumerate}
  \item[\rm 1\textsuperscript{o}] Si $Y_\lambda=X-U_\lambda$, on a $\codim(Y_\lambda,X)\geq2$ pour tout $\lambda\in L$.
  \item[\rm 2\textsuperscript{o}] Pour tout $x\in\bigcap_{\lambda\in L}U_\lambda$, l'anneau $\sh{O}_{X,x}$ est factoriel.
\end{enumerate}
Alors il existe des isomorphismes canoniques
\[
\label{IV.21.6.12.1}
  \varinjlim\operatorname{Div}(U_\lambda)\simto\mathfrak{Z}^1(X),\qquad
  \varinjlim\operatorname{Pic}(U_\lambda)\simto\mathfrak{Cl}(X)
  \tag{21.6.12.1}
\]
tels que, pour tout ouvert $V$ de $X$, les diagrammes
\[
\label{IV.21.6.12.2}
  \xymatrix{
    \varinjlim\operatorname{Div}(U_\lambda) \ar[r]^-\sim \ar[d] & \mathfrak{Z}^1(X) \ar[d] & & \varinjlim\operatorname{Pic}(U_\lambda) \ar[r]^-\sim \ar[d] & \mathfrak{Cl}(X) \ar[d]\\
    \varinjlim\operatorname{Div}(U_\lambda\cap V) \ar[r]^-\sim & \mathfrak{Z}^1(V) & & \varinjlim\operatorname{Pic}(U_\lambda\cap V) \ar[r]^-\sim & \mathfrak{Cl}(V)
  }
  \tag{21.6.12.2}
\]
soient commutatifs.
\end{proposition}

\begin{proof}
L'hypothèse 1\textsuperscript{o} sur $U_\lambda$ implique que pour tout $\lambda\in L$, on a $X^{(1)}\subset U_\lambda$ \sref{IV.5.1.3.1} donc l'homomorphisme de restriction $\mathfrak{Z}^1(X)\to\mathfrak{Z}^1(U_\lambda)$ est bijectif et par suite on a un isomorphisme canonique $\varinjlim\mathfrak{Z}^1(U_\lambda)\simto\mathfrak{Z}^1(X)$.
Les homomorphismes canoniques $\operatorname{cyc}:\operatorname{Div}(U_\lambda)\to\mathfrak{Z}^1(U_\lambda)$ définissent donc, par passage à la limite inductive, le premier des homomorphismes canoniques \sref{IV.21.6.12.1}, et le second s'en déduit par passage aux quotients.
En outre, il résulte de la condition 1\textsuperscript{o} que les $U_\lambda$ sont denses dans $X$, donc schématiquement denses puisque $X$ est réduit \sref{IV.11.10.4}, et par suite toute fonction méromorphe sur $X$ est entièrement déterminée par sa restriction à chaque $U_\lambda$;
on déduit aussitôt de là que dans l'isomorphisme $\mathfrak{Z}^1(X)\simto\mathfrak{Z}^1(U_\lambda)$, l'image de $\mathfrak{Z}^1_\mathrm{princ}(X)$ est $\mathfrak{Z}^1_\mathrm{princ}(U_\lambda)$, donc on a aussi un isomorphisme canonique $\varinjlim\mathfrak{Cl}(U_\lambda)\simto\mathfrak{Cl}(X)$.
Le second des homomorphismes canoniques \sref{IV.21.6.12.1} sera donc un isomorphisme si le premier l'est, et il reste donc à prouver ce dernier point, la commutativité des diagrammes \sref{IV.21.6.12.2} étant triviale.

\oldpage[IV-4]{276}
Montrons d'abord que l'homomorphisme $\varinjlim\operatorname{Div}(U_\lambda)\to\mathfrak{Z}^1(X)$ est \emph{injectif}.
Posons $T=\bigcap_\lambda U_\lambda$, et notons que les $U_\lambda$ forment un \emph{système fondamental de voisinages} de $T$.
En effet, pour tout ouvert $V\supset T$, $X-V$ est un sous-espace fermé de $X$, donc un espace noethérien dont toute partie fermée irréductible admet un point générique;
comme les ensembles $(X-V)\cap(X-U_\lambda)$ sont fermés dans $X-V$, forment une famille filtrante croissante et ont pour réunion $X-V$, notre assertion résulte de \sref[0]{0.9.2.4}.
Cela étant, il s'agit de voir que si $D\in\operatorname{Div}(U_\lambda)$ est tel que $\operatorname{cyc}(D)=0$ dans $\mathfrak{Z}^1(U_\lambda)$, alors il existe $\mu\geq\lambda$ tel que $D|U_\mu=0$.
En vertu de ce qui précède, il suffira de voir que pour tout $x\in T$, on a $D_x=0$;
en effet, il y aura alors un voisinage $W(x)$ de $x$ dans $X$ tel que $D|W(x)=0$, et la réunion des $W(x)$ pour $x\in T$ contient un $U_\mu$.
Compte tenu de \sref{IV.21.4.6}, on est donc ramené au cas où $X=\Spec(\sh{O}_{X,x})$ avec $x\in T$;
mais par hypothèse $\sh{O}_{X,x}$ est factoriel, donc intègre et intégralement clos, et $\Spec(\sh{O}_{X,x})$ est alors normal;
donc la conclusion résulte de \sref{IV.21.6.9}, (i).

\oldpage[IV-4]{277}
Pour prouver que l'homomorphisme $\varinjlim\operatorname{Div}(U_\lambda)\to\mathfrak{Z}^1(X)$ est bijectif, il suffira de prouver de même que pour tout $Z\in\mathfrak{Z}^1(X)$ et tout $x\in T$, il y a un voisinage $W(x)$ de $x$ dans $X$ et un diviseur $D^{(x)}$ sur $W(x)$ tel que $\operatorname{cyc}(D^{(x)})=Z|W(x)$.
En effet, on pourra alors recouvrir l'ensemble quasi-compact $T$ par un nombre fini de $W(x_i)$, et en vertu de la première partie de la démonstration, puisque les couples $(i,j)$ sont en nombre fini, il existera un indice $\lambda$ tel que dans chacun des $W(x_i)\cap W(x_j)\cap U_\lambda$, les restrictions de $D^{(x_i)}$ et $D^{(x_j)}$ coïncident;
comme on peut supposer en outre que $U_\lambda$ est contenu dans la réunion des $W(x_i)$, on voit que les restrictions $D^{(x_i)}|(W(x_i)\cap U_\lambda)$ sont les restrictions d'un même diviseur $D\in\operatorname{Div}(U_\lambda)$ qui sera tel que $\operatorname{cyc}(D)=Z|U_\lambda$.
On est donc ramené encore au cas où $X=\Spec(\sh{O}_{X,x})$ avec $x\in T$, mais comme $\sh{O}_{X,x}$ est factoriel, il en est de même de ses localisés (Bourbaki, \emph{Alg.\ comm.}, chap.~VII, \textsection~3, n\textsuperscript{o}~4, prop.~3), et il suffit d'appliquer \sref{IV.21.6.9}, (ii).
\end{proof}

\begin{corollary}[21.6.13]
\label{IV.21.6.13}
Soit $A$ un anneau local noethérien intègre et intégralement clos et de dimension $\geq2$.
Posons $X=\Spec(A)$, et soient $a$ le point fermé de $X$, $U=X-\{a\}$.
Pour que $A$ soit factoriel, il faut et il suffit que $U$ soit localement factoriel et que $\operatorname{Pic}(U)=0$.
\end{corollary}

\begin{proof}
En effet, dire que $A$ est factoriel équivaut à dire que $\mathfrak{Cl}(X)=0$ (Bourbaki, \emph{Alg.\ comm.}, chap.~VII, \textsection~1, n\textsuperscript{o}~4, cor.\ du th.~2 et \textsection~3, n\textsuperscript{o}~2, th.~1);
il suffit donc d'utiliser l'existence du second isomorphisme \sref{IV.21.6.12.1}, en prenant la famille $(U_\lambda)$ restreinte au seul ouvert $U$.
\end{proof}

\begin{corollary}[21.6.14]
\label{IV.21.6.14}
Soient $A$ un anneau local noethérien de dimension $\geq2$, $X=\Spec(A)$, $a$ le point fermé de $X$, $U=X-\{a\}$.
Les conditions suivantes sont équivalentes:
\begin{enumerate}[label=\emph{\alph*)}]
  \item $A$ est factoriel.
  \item $\operatorname{Pic}(U)=0$ et $\operatorname{prof}(A)\geq2$ (conditions que nous exprimerons plus loin \sref{IV.21.13} en disant que l'anneau $A$ est \emph{parafactoriel}), et $U$ est localement factoriel.
\end{enumerate}
\end{corollary}

\begin{proof}
Il est clair que si $A$ est factoriel, il est normal, donc $\operatorname{prof}(A)\geq2$ puisque $\dim(A)\geq2$ \sref[0]{0.16.5.1}, et \sref{IV.21.6.13} montre que \emph{a)} implique \emph{b)}.
Inversement, si \emph{b)} est vérifiée, il suffit de voir que $A$ est intégralement clos pour en conclure par \sref{IV.21.6.13} que \emph{b)} entraîne \emph{a)}.
En vertu du critère de Serre \sref{IV.5.8.6}, il suffit de vérifier pour $X$ les conditions $(\mathrm{R}_1)$ et $(\mathrm{S}_2)$.
Or, $U$ étant localement factoriel vérifie ces conditions, et l'hypothèse $\operatorname{prof}(A)\geq2$ entraîne que $X$ les vérifie également.
\end{proof}
\subsection{Interprétation des cycles positifs 1-codimensionnels en termes de sous-préschémas}
\EGAReaderTarget{ega4.c04.subsec.21-7}{subsection:IV.21.7}
\label{IV.21.7}

\begin{env}[21.7.1]
\label{IV.21.7.1}
Soient $X$ un préschéma localement noethérien, $C=\sum_{x\in X^{(1)}}n_x\cdot\overline{\{x\}}$ un cycle $1$-codimensionnel positif (de sorte que l'on a $n_x\geq0$ pour tout $x\in X^{(1)}$, et $n_x=0$ sauf en un ensemble localement fini de points).
Désignons par $Y(C)$ le sous-préschéma fermé de $X$, \emph{image fermée} (\sref[I]{I.9.5.3} et \sref{IV.9.5.1}) du préschéma \emph{somme} $Y'(C)=\coprod_{x\in X^{(1)}}\Spec(\sh{O}_{X,x}/\mathfrak{m}_x^{n_x})$ par le morphisme canonique, et par $\sh{J}_X(C)$ (ou $\sh{J}(C)$) l'Idéal de $\sh{O}_X$ définissant $Y(C)$.
\end{env}

\begin{proposition}[21.7.2]
\label{IV.21.7.2}
Soit $X$ un préschéma localement noethérien vérifiant la condition $(\mathrm{R}_1)$ \sref{IV.5.8.2}.
Pour qu'un sous-préschéma fermé $Y$ de $X$ soit de la forme $Y(C)$, où $C$ est un cycle $1$-codimensionnel positif, il faut et il suffit qu'il vérifie les deux conditions suivantes:
\begin{enumerate}[label=(\roman*)]
  \item $Y$ est purement de codimension $1$.
\oldpage[IV-4]{278}
  \item $Y$ vérifie la condition $(\mathrm{S}_1)$, autrement dit \sref{IV.5.7.5} n'a pas de cycle premier associé immergé.
\end{enumerate}
On a alors
\[
\label{IV.21.7.2.1}
  \operatorname{mult}_x(C)=\operatorname{long}\sh{O}_{Y,x}.
  \tag{21.7.2.1}
\]
L'application $C\mapsto Y(C)$ est une bijection de $\mathfrak{Z}^{1+}(X)$ sur l'ensemble des sous-préschémas fermés de $X$ vérifiant les conditions (i) et (ii).
\end{proposition}

\begin{proof}
Les conditions sont nécessaires (sans supposer que $X$ vérifie $(\mathrm{R}_1)$).
En effet, la question étant locale sur $X$, on peut supposer que $X=\Spec(A)$, où $A$ est un anneau noethérien;
on a alors $Y(C)=\Spec(A/\mathfrak{J})$, où $\mathfrak{J}$ est, par définition \sref[I]{I.9.5.1}, le noyau de l'homomorphisme canonique $A\to\bigoplus_i A_{\mathfrak{p}_i}/\mathfrak{p}_i^{n_i}A_{\mathfrak{p}_i}$, où les $\mathfrak{p}_i$ sont les idéaux premiers correspondant aux points $x\in X^{(1)}$ pour lesquels $n_x>0$, et où l'on a posé $n_i=n_x$.
L'image réciproque $\mathfrak{q}_i$ de $\mathfrak{p}_i^{n_i}A_{\mathfrak{p}_i}$ dans $A$ est un idéal $\mathfrak{p}_i$-primaire et on a $\mathfrak{J}=\bigcap_i\mathfrak{q}_i$.
D'ailleurs, comme les $x\in X^{(1)}$ sont tels que $\overline{\{x\}}$ est de codimension $1$, aucun point de $X^{(1)}$ ne peut être adhérent à un autre point de $X^{(1)}$.
Donc les $\mathfrak{p}_i$ sont les idéaux premiers minimaux de $A/\mathfrak{J}$ et l'ensemble des $\mathfrak{p}_i$ est égal à $\operatorname{Ass}(A/\mathfrak{J})$, ce qui prouve les conditions (i) et (ii).

Les conditions sont suffisantes.
Désignant par $\sh{J}$ l'Idéal de $\sh{O}_X$ définissant $Y$, l'hypothèse (ii) implique que $\operatorname{Ass}(\sh{O}_X/\sh{J})$ est identique à l'ensemble $F$ des points maximaux de $Y$, et l'hypothèse (i) implique que $F$ est contenu dans $X^{(1)}$;
donc \sref{IV.3.2.6} $\sh{O}_X/\sh{J}$ s'identifie à un sous-$\sh{O}_X$-Module de $\bigoplus_{x\in F}\mathscr{G}(x)$, où pour chaque $x\in F$, $\mathscr{G}(x)$ est un $\sh{O}_X$-Module irredondant tel que $\operatorname{Ass}(\mathscr{G}(x))=\{x\}$.
Or, puisque $X$ vérifie $(\mathrm{R}_1)$, pour chaque $x\in F$, $\sh{O}_{X,x}$ est un anneau de valuation discrète et par suite les idéaux primaires non nuls de cet anneau sont les puissances $\mathfrak{m}_x^{n_x}$ de l'idéal maximal;
supposant encore que $X=\Spec(A)$, on en conclut que $\sh{J}=\widetilde{\mathfrak{J}}$, où $\widetilde{\mathfrak{J}}=\bigcap_i\mathfrak{q}_i$, les $\mathfrak{q}_i$ étant les images réciproques dans $A$ d'idéaux $\mathfrak{p}_i^{n_i}A_{\mathfrak{p}_i}$, où les $\mathfrak{p}_i$ correspondent aux points de $F$.
On voit donc bien que $Y$ est de la forme $Y(C)$.
\end{proof}

\begin{corollary}[21.7.3]
\label{IV.21.7.3}
Soit $X$ un préschéma localement noethérien.
Pour tout diviseur positif $D$ sur $X$, tel que $X$ vérifie $(\mathrm{R}_1)$ aux points maximaux du sous-préschéma fermé $Y(D)$ \sref{IV.21.2.12}, $Y(D)$ majore le sous-préschéma fermé $Y(\operatorname{cyc}(D))$ \sref{IV.21.7.1}, et a même espace sous-jacent;
pour que ces deux sous-préschémas soient égaux, il faut et il suffit que $Y(D)$ vérifie la condition $(\mathrm{S}_1)$, ou encore que, pour tout $z\in Y(D)$ distinct d'un point maximal de $Y(D)$, on ait $\operatorname{prof}(\sh{O}_{X,z})\geq2$ (condition toujours vérifiée lorsque $X$ est \emph{normal}).
\end{corollary}

\begin{proof}
En effet, la question étant locale, on peut toujours supposer que $\sh{J}_X(D)=t\sh{O}_X$, $t$ étant une section régulière de $\sh{O}_X$ au-dessus de $X$;
en tout point maximal $x$ de $Y(D)$, appartenant nécessairement à $X^{(1)}$, $\sh{O}_{X,x}$ est par hypothèse un anneau de valuation discrète, donc $t_x\sh{O}_{X,x}=\mathfrak{m}_x^{n_x}$ où $n_x=\operatorname{mult}_x(D)$.
On peut supposer $X=\Spec(A)$, et alors, si $\sh{J}_X(\operatorname{cyc}(D))=\widetilde{\mathfrak{J}}$, il résulte de ce qui précède et de \sref{IV.21.7.1} que $\mathfrak{J}=tA$ et $\widetilde{\mathfrak{J}}$ ont mêmes \emph{idéaux primaires non immergés}, donc $\widetilde{\mathfrak{J}}\subset\mathfrak{J}$, puisque $\widetilde{\mathfrak{J}}$ est l'intersection
\oldpage[IV-4]{279}
de ces idéaux primaires \sref{IV.21.7.2}.
Cela prouve que $Y(D)$ majore $Y(\operatorname{cyc}(D))$ et que ces deux sous-préschémas sont égaux si et seulement si $Y(D)$ n'a pas de cycle premier associé immergé (autrement dit vérifie $(\mathrm{S}_1)$).
Comme pour tout $x\in Y(D)$, il y a par hypothèse un voisinage ouvert $U$ de $x$ dans $X$ et un élément régulier $t\in\Gamma(U,\sh{O}_X)$ tel que $\sh{O}_{Y(D)}|(U\cap Y(D))$ soit restriction à $Y(D)\cap U$ de $\sh{O}_U/t\sh{O}_U$, dire que $Y(D)$ vérifie $(\mathrm{S}_1)$ signifie aussi qu'en tout point $z\in Y(D)$ non maximal, on a $\operatorname{prof}(\sh{O}_{X,z}/t_z\sh{O}_{X,z})\geq1$, c'est-à-dire \sref[0]{0.16.4.6} $\operatorname{prof}(\sh{O}_{X,z})\geq2$.
L'assertion concernant le cas où $X$ est normal est alors immédiate puisque $X$ vérifie $(\mathrm{S}_2)$ et qu'aux points non maximaux $z$ de $Y(D)$ on a $\dim(\sh{O}_{X,z})\geq2$ \sref[0]{0.16.3.4}.
\end{proof}

\begin{env}[21.7.3.1]
\label{IV.21.7.3.1}
On voit donc que lorsque $X$ est \emph{normal}, $\operatorname{Div}^+(X)$ s'identifie canoniquement à l'ensemble des sous-préschémas fermés $Y$ de $X$ vérifiant les conditions (i) et (ii) de \sref{IV.21.7.2} et \emph{régulièrement immergés} dans $X$.
\end{env}

\begin{proposition}[21.7.4]
\label{IV.21.7.4}
Soit $X$ un préschéma localement noethérien, réduit en chacun de ses points isolés.
Les conditions suivantes sont équivalentes:
\begin{enumerate}
  \item[\emph{a)}] L'homomorphisme canonique $c:\sh{D}iv_X\to\mathscr{Z}_X^1$ \sref{IV.21.6.4} est un isomorphisme de faisceaux de groupes ordonnés.
  \item[\emph{a$'$)}] Tout cycle premier $1$-codimensionnel de $X$ est l'image par $\operatorname{cyc}$ d'un diviseur positif, et l'homomorphisme canonique $c:\sh{D}iv_X\to\mathscr{Z}_X^1$ est injectif.
  \item[\emph{a$''$)}] Pour tout sous-préschéma fermé intègre $Y$ de $X$, de codimension $1$, l'immersion canonique $Y\to X$ est régulière.
  \item[\emph{b)}] $X$ est normal et l'homomorphisme $\operatorname{cyc}:\operatorname{Div}(X)\to\mathfrak{Z}^1(X)$ est bijectif.
  \item[\emph{c)}] $X$ est localement factoriel.
\end{enumerate}
\end{proposition}

\begin{proof}
L'équivalence de \emph{b)} et \emph{c)} a été démontrée dans \sref{IV.21.6.9}, ainsi que le fait que \emph{c)} entraîne \emph{a)}.
De plus \emph{b)} entraîne \emph{a$''$)} en vertu de \sref{IV.21.7.3.1}, et \emph{a)} implique trivialement \emph{a$'$)}.
Il reste à voir que \emph{a$'$)} ou \emph{a$''$)} entraîne \emph{c)}.

Supposons donc vérifiée la condition \emph{a$'$)}, et montrons d'abord que $X$ est normal.
Notons en premier lieu que si $X$ vérifie \emph{a$'$)}, il en est de même de tout schéma local $\Spec(\sh{O}_{X,x})$.
Considérons alors $x\in X^{(1)}$, de sorte que $A=\sh{O}_{X,x}$ est un anneau local noethérien de dimension $1$.
Appliquant la condition \emph{a$'$)} à $\Spec(A)$ et au cycle premier $1$-codimensionnel formé du point fermé de $\Spec(A)$, on voit que dans $A$ l'idéal maximal est engendré par un seul élément régulier dans $A$;
donc \sref[0]{0.17.1.1} $A$ est un anneau régulier.
Comme les anneaux localisés $A_{\mathfrak p}$ sont aussi réguliers \sref[0]{0.17.3.2}, on voit que $X$ est régulier en tous ses points maximaux non isolés;
comme il est aussi réduit (donc régulier) en ses points isolés par hypothèse, on en conclut que $X$ vérifie la condition $(\mathrm{R}_1)$.
Montrons en outre que $X$ vérifie aussi $(\mathrm{S}_1)$, autrement dit que pour tout $x\in X$ tel que $\dim(\sh{O}_{X,x})\geq1$, on a $\mathfrak{m}_x\notin\operatorname{Ass}(\sh{O}_{X,x})$ \sref[0]{0.16.4.6}.
En effet, l'hypothèse \emph{a$'$)} appliquée à $X_1=\Spec(\sh{O}_{X,x})$ montre qu'il existe sur ce préschéma au moins un diviseur non nul, autrement dit que l'on a $\sh{M}_{X_1}^*\neq\sh{O}_{X_1}^*$, et il suffit d'appliquer \sref{IV.21.1.10}.
On déduit donc déjà de ces résultats que $X$ est réduit \sref{IV.5.8.5}.
Prouvons ensuite qu'il vérifie la condition $(\mathrm{S}_2)$ (ce qui établira que $X$ est normal, en vertu du critère de Serre \sref{IV.5.8.6}).
Raisonnons par l'absurde en supposant que l'ensemble $E$ des points $x\in X$ où $X$ ne vérifie pas $(\mathrm{S}_2)$ soit non vide, et soit $x\in E$ un point pour lequel $\dim(\sh{O}_{X,x})$ est la plus petite possible;
puisque $X$ vérifie $(\mathrm{S}_1)$, on a $\dim(\sh{O}_{X,x})\geq2$.
Dans $X_1=\Spec(\sh{O}_{X,x})$, l'ouvert $U=X_1-\{x\}$ vérifie $(\mathrm{S}_2)$;
montrons que $X_1$ vérifie $(\mathrm{S}_2)$, de sorte que l'on aura abouti à une contradiction.
Il suffit, en vertu de \sref{IV.5.10.5}, de montrer que toute section $f$ de $\sh{O}_{X_1}$ au-dessus de $U$ se prolonge en une section de $\sh{O}_{X_1}$ au-dessus de $X_1$.
Comme $X_1$ est réduit et $U$ dense dans $X_1$, $U$ est schématiquement dense dans $X_1$, et par suite $f$ est la restriction à $U$ d'une section méromorphe régulière $g\in M(X_1)$.
En outre, comme $\dim(\sh{O}_{X,x})\geq2$, on a $\operatorname{cyc}(\operatorname{div}(g))\geq0$ puisque $g|U=f$;
comme l'homomorphisme $\operatorname{cyc}$ est injectif, on a nécessairement $\operatorname{div}(g)\geq0$, donc \sref{IV.21.6.4} $g$ est une section de $\sh{O}_{X_1}$ au-dessus de $X_1$, ce qui établit notre assertion.

Pour montrer que \emph{a$'$)} implique \emph{c)}, remarquons alors que la condition \emph{a$'$)} implique que l'homomorphisme canonique $\operatorname{cyc}:\operatorname{Div}(X)\to\mathfrak{Z}^1(X)$ est surjectif;
puisque $X$ est normal, il suffit d'appliquer \sref{IV.21.6.9}.

Prouvons maintenant que \emph{a$''$)} entraîne \emph{c)}.
Il résulte de \sref{IV.21.6.5} que \emph{a$''$)} entraîne que tout cycle premier $1$-codimensionnel sur $X$ est l'image par $\operatorname{cyc}$ d'un diviseur positif;
la première partie du raisonnement fait ci-dessus prouve donc encore que $X$ vérifie $(\mathrm{R}_1)$ et $(\mathrm{S}_1)$.
Il reste encore à voir que $X$ est normal (la fin du raisonnement étant alors la même), c'est-à-dire que si $x\in X$ est tel que $\dim(\sh{O}_{X,x})\geq2$, on a $\operatorname{prof}(\sh{O}_{X,x})\geq2$.
Or, si $y$ est une générisation de $x$ telle que $\overline{\{y\}}$ soit de codimension $1$ dans $X$, le sous-préschéma réduit $Y$ de $X$ ayant $\overline{\{y\}}$ pour espace
\oldpage[IV-4]{280}
sous-jacent est intègre, donc régulièrement immergé dans $X$ par hypothèse.
Cela entraîne qu'il existe un élément régulier $t$ de $\sh{O}_{X,x}$ tel que l'idéal $t\sh{O}_{X,x}$ soit l'idéal premier définissant le préschéma $Y_1$ image réciproque de $Y$ dans $X_1=\Spec(\sh{O}_{X,x})$.
Comme $\sh{O}_{X,x}/t\sh{O}_{X,x}$ est intègre, cela prouve que $\operatorname{prof}(\sh{O}_{X,x})\geq2$ et termine la preuve de \sref{IV.21.7.4}.
\end{proof}

\begin{remark}[21.7.5]
\label{IV.21.7.5}
On ne peut dans \sref{IV.21.7.4} remplacer la condition \emph{a$'$)} par la condition plus faible que tout cycle premier $1$-codimensionnel de $X$ soit l'image par $\operatorname{cyc}$ d'un diviseur positif.
C'est ce que montre l'exemple où $X$ est le schéma affine défini dans \sref{IV.14.1.5} (« réunion d'un plan et d'une droite qui le coupe ») \emph{qui n'est pas normal};
avec les notations introduites dans \sref{IV.14.1.5}, les cycles premiers $1$-codimensionnels de $X$ sont ceux du plan $X_1$ et les points fermés de la droite $X_2$ distincts du point commun à $X_1$ et $X_2$;
si $t_1$, $t_2$, $t_3$ sont les images de $T_1$, $T_2$, $T_3$ dans $A/\mathfrak a$, on voit donc que les cycles premiers $1$-codimensionnels de $X$ sont définis par les idéaux premiers monogènes de générateurs $P(t_1,t_2)$ ($P$ irréductible dans $K[T_1,T_2]$) ou $t_3-a$, avec $a\neq0$ dans $K$;
ces éléments étant réguliers dans $A/\mathfrak a$, tout cycle $1$-codimensionnel est bien l'image canonique d'un diviseur positif.
\end{remark}

\begin{corollary}[21.7.6]
\label{IV.21.7.6}
Soit $X$ un préschéma localement noethérien, réduit en chacun de ses points isolés.
Les conditions suivantes sont équivalentes:
\begin{enumerate}
  \item[\emph{a)}] L'homomorphisme canonique $c:\sh{D}iv_X\to\mathscr{Z}_X^1$ est un isomorphisme de faisceaux de groupes ordonnés, et on a $\operatorname{Div}(X)=\operatorname{Div.princ}(X)$.
  \item[\emph{a$'$)}] L'homomorphisme canonique $c:\sh{D}iv_X\to\mathscr{Z}_X^1$ est injectif, et tout cycle premier $1$-codimensionnel est l'image par $\operatorname{cyc}$ d'un diviseur positif principal, autrement dit est de la forme $\operatorname{cyc}(\operatorname{div}(f))$, où $f$ est une section régulière de $\sh{O}_X$ au-dessus de $X$.
  \item[\emph{a$''$)}] Pour tout sous-préschéma fermé intègre $Y$ de $X$, de codimension $1$, il existe une section régulière $f$ de $\sh{O}_X$ au-dessus de $X$ telle que $Y$ soit défini par l'Idéal $f\sh{O}_X$.
  \item[\emph{b)}] $X$ est normal et tout cycle premier $1$-codimensionnel sur $X$ est principal.
  \item[\emph{c)}] $X$ est localement factoriel et $\operatorname{Pic}(X)=0$.
  \item[\emph{d)}] $X$ est normal, et pour tout ouvert $U$ de $X$, on a $\operatorname{Pic}(U)=0$.
\end{enumerate}
\end{corollary}

\begin{proof}
L'équivalence de \emph{a)}, \emph{a$'$)}, \emph{a$''$)}, \emph{b)} et \emph{c)} résulte aussitôt de \sref{IV.21.7.4} et de \sref{IV.21.6.11}.
En outre, il résulte aussitôt de \emph{a$''$)} que tout ouvert non vide $U$ de $X$ vérifie encore les mêmes conditions, autrement dit ces conditions impliquent \emph{d)}.
Reste à voir que \emph{d)} entraîne \emph{b)}.
Or, désignons par $U_\lambda$ les ouverts complémentaires des réunions finies d'ensembles de la forme $\overline{\{x\}}$, où $\dim(\sh{O}_{X,x})\geq2$;
il est clair que les $U_\lambda$ forment une famille filtrante décroissante et que pour tout $x\in\bigcap_\lambda U_\lambda$, on a $\dim(\sh{O}_{X,x})\leq1$, donc $\sh{O}_{X,x}$ est un anneau principal en vertu de l'hypothèse.
On peut donc appliquer à la famille $(U_\lambda)$ la proposition \sref{IV.21.6.12}, qui montre que $\mathfrak{Cl}(X)$ est isomorphe à $\varinjlim\operatorname{Pic}(U_\lambda)$, donc $\mathfrak{Cl}(X)=0$ en vertu de l'hypothèse, ce qui prouve \emph{b)} par définition.
\end{proof}

\begin{remark}[21.7.7]
\label{IV.21.7.7}
Lorsque $X=\Spec(A)$, où $A$ est un anneau noethérien intègre, les conditions équivalentes de \sref{IV.21.7.6} équivalent à dire que $A$ est un anneau \emph{factoriel}.
\end{remark}

\subsection{Diviseurs et normalisation}
\EGAReaderTarget{ega4.c04.subsec.21-8}{subsection:IV.21.8}

\begin{lemma}[21.8.1]
\label{IV.21.8.1}
Soit $f:X'\to X$ un morphisme entier de préschémas.
Pour tout $\sh{O}_{X'}$-Module localement libre $E'$ de rang constant $n$ et tout $x\in X$, il existe un voisinage ouvert $U$ de $x$ dans $X$ tel que, si l'on pose $U'=f^{-1}(U)$, $E'|U'$ soit isomorphe à $\sh{O}_{U'}^n$.
\end{lemma}

\begin{proof}
La question étant locale sur $X$, on peut se borner au cas où $X=\Spec(A)$ est affine, auquel cas on a $X'=\Spec(A')$, où $A'$ est une $A$-algèbre entière.
Alors $A'$ est limite inductive de ses sous-$A$-algèbres \emph{finies} $A'_\lambda$.
Posant $X'_\lambda=\Spec(A'_\lambda)$ et désignant par $p_\lambda:X'\to X'_\lambda$ le morphisme structural, il résulte de \sref{IV.8.5.2} et \sref{IV.8.5.5} qu'il existe un indice $\lambda$ et un $\sh{O}_{X'_\lambda}$-Module localement libre $E'_\lambda$ de rang $n$ tel que $E'=p_\lambda^*(E'_\lambda)$ et il suffira évidemment de prouver le lemme pour $X'_\lambda$ et $E'_\lambda$;
autrement dit, on peut se borner au cas où le morphisme $f$ est \emph{fini}.
Posons $B=\sh{O}_{X,x}$ et soit $B'$ l'anneau du schéma affine $X'_0=X'\times_X\Spec(\sh{O}_{X,x})$;
comme $B$ est un anneau local et que $B'$ est une $B$-algèbre finie, $B'$ est un anneau semi-local (Bourbaki, \emph{Alg.\ comm.}, chap.~V, \textsection~2, n\textsuperscript{o}~1, prop.~3);
on en conclut que le $\sh{O}_{X'_0}$-Module localement libre $E'_0=E'\otimes_{\sh{O}_{X'}}\sh{O}_{X'_0}$ est isomorphe à $\sh{O}_{X'_0}^n$ (Bourbaki, \emph{Alg.\ comm.}, chap.~II, \textsection~5, n\textsuperscript{o}~3, prop.~5).
Considérant $X'_0$ comme limite
\oldpage[IV-4]{281}
projective des préschémas induits par $X'$ sur les ouverts $f^{-1}(U)$, où $U$ parcourt l'ensemble filtrant des voisinages ouverts affines de $x$ dans $X$, suivant la méthode de \sref{IV.8.1.2}, \emph{a)}), et appliquant \sref{IV.8.5.2.5}, on obtient la conclusion cherchée.
\end{proof}

\begin{corollary}[21.8.2]
\label{IV.21.8.2}
Soit $f:X'\to X$ un morphisme entier de préschémas.
Alors:
\begin{enumerate}[label=(\roman*)]
  \item On a $\RR^1 f_*(\sh{O}^*_{X'})=0$ \sref[0]{0.12.2.1}.
  \item Le groupe $\operatorname{Pic}(X')$ est canoniquement isomorphe à $\HH^1(X,f_*(\sh{O}^*_{X'}))$.
\end{enumerate}
\end{corollary}

\begin{proof}
\begin{enumerate}[label=(\roman*)]
  \item $\RR^1 f_*(\sh{O}^*_{X'})$ est le faisceau associé au préfaisceau de groupes commutatifs $U\mapsto\HH^1(f^{-1}(U),\sh{O}^*_{X'})$ sur $X$ (\emph{loc.\ cit.});
la fibre de ce faisceau en un point $x$ peut donc \sref[0]{0.5.4.7} être identifiée au groupe commutatif $\varinjlim\operatorname{Pic}(f^{-1}(U))$, où $U$ parcourt l'ensemble des voisinages ouverts de $x$, les homomorphismes de transition $\operatorname{Pic}(f^{-1}(U'))\to\operatorname{Pic}(f^{-1}(U))$ pour deux ouverts $U\subset U'$ étant définis par \sref{IV.21.3.2.4}.
Or, pour tout $x\in X$, tout voisinage ouvert $U'$ de $x$ dans $X$, et tout élément $\zeta\in\operatorname{Pic}(f^{-1}(U'))$, il résulte de \sref{IV.21.8.1} qu'il y a un voisinage ouvert $U\subset U'$ de $x$ tel que l'image de $\zeta$ dans $\operatorname{Pic}(f^{-1}(U))$ soit nulle.
Donc la fibre de $\RR^1f_*(\sh{O}^*_{X'})$ au point $x$ est nulle.

  \item La suite spectrale de Leray pour le foncteur composé $\sh{F}'\mapsto\Gamma(X,f_*(\sh{F}'))$ de faisceaux de groupes commutatifs sur $X'$ (T, 2.4) donne la suite exacte de termes de bas degré
\[
  0\to\HH^1(X,f_*(\sh{O}^*_{X'}))\to\HH^1(X',\sh{O}^*_{X'})\to\HH^0(X,\RR^1f_*(\sh{O}^*_{X'}))
\]
et la conclusion résulte de (i) et de l'isomorphisme de $\operatorname{Pic}(X')$ et de $\HH^1(X',\sh{O}^*_{X'})$ \sref[0]{0.5.4.7}.
\end{enumerate}
\end{proof}

\begin{proposition}[21.8.3]
\label{IV.21.8.3}
Soit $f=(\psi,\theta):X'\to X$ un morphisme entier de préschémas;
supposons que, pour tout ouvert $U$ de $X$, l'homomorphisme $\Gamma(\theta):\Gamma(U,\sh{O}_X)\to\Gamma(U,f_*(\sh{O}_{X'}))$ transforme les éléments réguliers en éléments réguliers, ce qui permet de prolonger canoniquement l'homomorphisme $\theta:\sh{O}_X\to f_*(\sh{O}_{X'})$ en un homomorphisme de faisceaux d'anneaux $\theta':\sh{M}_X\to f_*(\sh{M}_{X'})$, d'où des homomorphismes de faisceaux de groupes multiplicatifs $\theta^*:\sh{O}^*_X\to f_*(\sh{O}^*_{X'})$ et $\theta'^*:\sh{M}^*_X\to f_*(\sh{M}^*_{X'})$, donnant par passage aux quotients un homomorphisme $\theta''^*:\sh{D}iv_X\to f_*(\sh{D}iv_{X'})$.
On a alors un diagramme commutatif
\[
\label{IV.21.8.3.1}
  \xymatrix{
    1 \ar[r] & \sh{O}^*_X \ar[r] \ar[d]^{\theta^*} & \sh{M}^*_X \ar[r] \ar[d]^{\theta'^*} & \sh{D}iv_X \ar[r] \ar[d]^{\theta''^*} & 0\\
    1 \ar[r] & f_*(\sh{O}^*_{X'}) \ar[r] & f_*(\sh{M}^*_{X'}) \ar[r] & f_*(\sh{D}iv_{X'}) \ar[r] & 0
  }
  \tag{21.8.3.1}
\]
dans lequel les deux lignes sont exactes.
\end{proposition}

\begin{proof}
Le seul point à vérifier est l'exactitude de la seconde ligne du diagramme, qui résulte de l'application de la suite exacte de cohomologie pour le foncteur $f_*$ à la suite exacte de faisceaux de groupes commutatifs sur $X'$
\[
  1\to\sh{O}^*_{X'}\to\sh{M}^*_{X'}\to\sh{D}iv_{X'}\to0
\]
et de \sref{IV.21.8.2}.
\end{proof}

\begin{corollary}[21.8.4]
\label{IV.21.8.4}
Si, en outre des hypothèses de \sref{IV.21.8.3}, l'homomorphisme $\theta'$ est un isomorphisme de faisceaux d'anneaux, alors $\theta^*:\sh{O}^*_X\to f_*(\sh{O}^*_{X'})$ est injectif, $\theta''^*:\sh{D}iv_X\to f_*(\sh{D}iv_{X'})$ est surjectif et $\operatorname{Ker}(\theta''^*)$ est isomorphe à $\operatorname{Coker}(\theta^*)$.
\end{corollary}

\begin{proof}
C'est une conséquence immédiate du diagramme du serpent (Bourbaki, \emph{Alg.\ comm.}, chap.~I, \textsection~2, n\textsuperscript{o}~4, prop.~2) appliqué au diagramme \sref{IV.21.8.3.1}.
\end{proof}

\oldpage[IV-4]{282}

\begin{proposition}[21.8.5]
\label{IV.21.8.5}
Soient $X$, $X'$ deux préschémas localement noethériens, $f:X'\to X$ un morphisme entier;
on suppose qu'il existe un ouvert schématiquement dense $U$ dans $X$ tel que $f^{-1}(U)$ soit schématiquement dense dans $X'$ (cf.\ \sref{IV.20.3.5}), et que le morphisme $f^{-1}(U)\to U$, restriction de $f$, soit un isomorphisme.
Alors:
\begin{enumerate}[label=(\roman*)]
  \item L'homomorphisme $\theta':\sh{M}_X\to f_*(\sh{M}_{X'})$ est défini et bijectif, l'homomorphisme
\[
  \theta^*:\sh{O}^*_X\to f_*(\sh{O}^*_{X'})
\]
est injectif, l'homomorphisme $\theta''^*:\sh{D}iv_X\to f_*(\sh{D}iv_{X'})$ est surjectif, $\operatorname{Ker}(\theta''^*)$ est isomorphe à $\sh{N}=\operatorname{Coker}(\theta^*)$ et le support du faisceau de groupes multiplicatifs $\sh{N}$ est contenu dans $S=X-U$.
  \item Supposons de plus l'ensemble $S$ discret et posons pour abréger $\sh{O}'_X=f_*(\sh{O}_{X'})$.
Alors on a un diagramme commutatif
\[
\label{IV.21.8.5.1}
  \xymatrix{
    1 \ar[r] & \prod_{s\in S}\sh{O}'^{*}_{X,s}/\sh{O}^*_{X,s} \ar[r]^-{j} \ar[d] & \operatorname{Div}(X) \ar[r] \ar[d]^{k_X} & \operatorname{Div}(X') \ar[r] \ar[d]^{k_{X'}} & 0\\
    1 \ar[r] & \bigl(\prod_{s\in S}\sh{O}'^{*}_{X,s}/\sh{O}^*_{X,s}\bigr)/\operatorname{Im}\Gamma(X',\sh{O}^*_{X'}) \ar[r]^-{j'} & \operatorname{Pic}(X) \ar[r] & \operatorname{Pic}(X') \ar[r] & 0
  }
  \tag{21.8.5.1}
\]
où les lignes sont exactes et où la flèche verticale de gauche est l'homomorphisme canonique.
\end{enumerate}
\end{proposition}

\begin{proof}
\begin{enumerate}[label=(\roman*)]
  \item L'hypothèse entraîne que l'on a un isomorphisme canonique $\sh{M}''_X\simto f_*(\sh{M}''_{X'})$ pour les faisceaux de germes de pseudo-fonctions \sref{IV.20.2.2}, d'où l'assertion relative à $\theta'$ vu l'existence des isomorphismes canoniques $\sh{M}_X\simto\sh{M}''_X$ et $\sh{M}_{X'}\simto\sh{M}''_{X'}$ \sref{IV.20.2.11}.
Le reste de l'assertion (i) découle de \sref{IV.21.8.4}, sauf ce qui concerne le support de $\sh{N}$, qui résulte directement de l'hypothèse sur $U$.

  \item Appliquant la suite exacte de cohomologie aux deux suites exactes de faisceaux de groupes commutatifs
\[
\label{IV.21.8.5.2}
\begin{gathered}
  1\to\sh{N}\to\sh{D}iv_X\to f_*(\sh{D}iv_{X'})\to0,\\
  1\to\sh{O}^*_X\to f_*(\sh{O}^*_{X'})\to\sh{N}\to1,
\end{gathered}
  \tag{21.8.5.2}
\]
il vient respectivement les suites exactes
\[
\label{IV.21.8.5.3}
\begin{gathered}
  1\to\Gamma(X,\sh{N})\xrightarrow{j}\operatorname{Div}(X)\to\operatorname{Div}(X')\to\HH^1(X,\sh{N}),\\
  \Gamma(X',\sh{O}^*_{X'})\to\Gamma(X,\sh{N})\xrightarrow{\partial}\operatorname{Pic}(X)\to\operatorname{Pic}(X')\to\HH^1(X,\sh{N}),
\end{gathered}
  \tag{21.8.5.3}
\]
\oldpage[IV-4]{283}
compte tenu de l'isomorphisme canonique $\operatorname{Pic}(X')\simto\HH^1(X,f_*(\sh{O}^*_{X'}))$ \sref{IV.21.8.2}.
Or, comme le support de $\sh{N}$ est contenu dans l'ensemble discret $S$, fermé dans $X$, on a
\[
  \HH^1(X,\sh{N})=\HH^1(S,\sh{N}|S)\quad\text{(G, II, 4.9.2)},\qquad
  \HH^1(S,\sh{N}|S)=\prod_{s\in S}\HH^1(\{s\},\sh{N}_s)=0
\]
par définition de la cohomologie (G, II, 4.4).
De même on a $\Gamma(X,\sh{N})=\prod_{s\in S}\sh{N}_s$ et $\sh{N}_s=\sh{O}'^{*}_{X,s}/\sh{O}^*_{X,s}$ en vertu de la seconde suite exacte \sref{IV.21.8.5.2}.
Reste à préciser les injections $j$ et $j'$.
Or, on peut décrire une section $t$ de $\sh{N}$ au-dessus de $X$ en prenant un recouvrement de $X$ formé de $U$ et d'ouverts $U_i$ tels que $U_i\cap S$ soit réduit à un seul point $s_i$, et que $U_i\cap U_j\subset U$ pour $i\neq j$, puis en considérant pour chaque $i$ une section $t_i$ de $\sh{N}$ au-dessus de $U_i$, nécessairement telle que $(t_i)_{s_i}\in\sh{O}'^{*}_{X,s_i}/\sh{O}^*_{X,s_i}$ et $(t_i)_x=0$ aux points $x\in U_i-\{s_i\}$.
Le germe $(t_i)_{s_i}$ provient d'une section $u_i$ de $\sh{O}^*_{X'}$ au-dessus de $f^{-1}(U_i)$, qu'on peut aussi considérer comme une section de $\sh{M}^*_{X'}$ au-dessus de $f^{-1}(U_i)$, donc une section de $\sh{M}^*_X$ au-dessus de $U_i$ (en vertu de l'hypothèse);
il correspond canoniquement à cette section une section $d_i\in\Gamma(U_i,\sh{D}iv_X)$ et comme dans $f^{-1}(U\cap U_i)$ la restriction de $u_i$ s'identifie à une section de $\sh{O}^*_X$ au-dessus de $U\cap U_i$ en vertu de l'hypothèse, les restrictions des $d_i$ à $U\cap U_i$ sont toutes nulles, donc les $d_i$ sont les restrictions d'un même diviseur de $\operatorname{Div}(X)$, qui est l'image de la section $t$ par $j$.
De même, l'image de $t$ par $\partial:\Gamma(X,\sh{N})\to\HH^1(X,\sh{O}^*_X)$ provient du cocycle égal à la restriction de $u_i$ dans $U\cap U_i$ pour chaque $i$, à $(u_i|(U_i\cap U_j))(u_j|(U_i\cap U_j))^{-1}$ dans $U_i\cap U_j$, d'où l'expression de $j'(t)$ et la commutativité du diagramme \sref{IV.21.8.5.1}.
\end{enumerate}
\end{proof}

\begin{remark}[21.8.6]
\label{IV.21.8.6}
\begin{enumerate}[label=(\roman*)]
  \item Les conditions d'application de \sref{IV.21.8.5}, (i) sont en particulier remplies lorsque $X$ est un préschéma localement noethérien réduit n'ayant qu'un nombre fini de composantes irréductibles, $X'$ son normalisé et que $X'$ est aussi localement noethérien (\sref[II]{II.6.3.8});
le $\sh{O}_X$-Module $\sh{D}iv_X$ est alors une \emph{extension} de $f_*(\sh{D}iv_{X'})$ par le conoyau $\sh{N}$ de $\sh{O}^{*}_X\to f_*(\sh{O}^*_{X'})$, que l'on peut considérer comme connu.
Si de plus $X$ est une courbe algébrique (réduite) sur un corps $k$, on est dans les conditions d'application de \sref{IV.21.8.5}, (ii).

  \item Lorsque les conditions d'application de \sref{IV.21.8.5}, (ii) sont remplies et en outre que l'homomorphisme canonique global $\Gamma(X,\sh{O}_X)\to\Gamma(X',\sh{O}_{X'})$ est \emph{bijectif}, on voit que les noyaux des homomorphismes surjectifs $\operatorname{Div}(X)\to\operatorname{Div}(X')$ et $\operatorname{Pic}(X)\to\operatorname{Pic}(X')$ sont \emph{isomorphes}.

  \item Lorsque les conditions d'application de \sref{IV.21.8.5}, (ii) sont remplies, on voit que l'homomorphisme $\operatorname{Div}(X)\to\operatorname{Div}(X')$ ne peut être injectif (auquel cas il est bijectif) que si $\sh{O}'^{*}_{X,s}=\sh{O}^*_{X,s}$ pour tout $s\in S$.
Pour un $s\in S$ tel que l'anneau $\sh{O}'_{X,s}$ soit \emph{local} (ce qui, compte tenu du fait que $X'=\Spec(\sh{O}'_X)$ \sref[II]{II.1.3}, signifie qu'il n'existe qu'un seul point $s'\in X'$ au-dessus de $s$) cela entraîne nécessairement que les corps résiduels $\boldsymbol{k}(s)$ et $\boldsymbol{k}(s')$ sont égaux et que $1+\mathfrak{m}_s=1+\mathfrak{m}_{s'}$, donc finalement équivaut à la relation $\sh{O}'_{X,s}=\sh{O}_{X,s}$ ou encore (compte tenu de l'hypothèse) au fait qu'il y a un voisinage ouvert $V$ de $s$ dans $X$ tel que le morphisme $f^{-1}(V)\to V$ soit un isomorphisme.
Dans le cas général, l'anneau $\sh{O}'_{X,s}$ est semi-local (c'est évident lorsque le morphisme $f$ est fini et dans le cas général cela résulte de \sref[0]{0.23.2.6}), l'homomorphisme canonique $\sh{O}^{*}_{X,s}\to\sh{O}'^{*}_{X,s}$ définit, par passage aux quotients, un homomorphisme du groupe multiplicatif $\boldsymbol{k}(s)^*$ dans le produit des groupes multiplicatifs $\boldsymbol{k}(s'_j)^*$, les $s'_j$ étant les points (en
\oldpage[IV-4]{284}
nombre $>1$) de $X'$ au-dessus de $s$.
Il est immédiat que cet homomorphisme ne peut être bijectif que si $\boldsymbol{k}(s)$ et tous les $\boldsymbol{k}(s'_j)$ sont égaux au corps $\bb{F}_2$ à $2$ éléments;
de plus, si cette condition est vérifiée, il faut en outre que le groupe multiplicatif $1+\mathfrak{m}_s$ ait pour image $1+\mathfrak{r}$, où $\mathfrak{r}$ est le radical de $\sh{O}'_{X,s}$, ou encore que $\mathfrak{m}_s=\mathfrak{r}$, ce qui entraîne que $\sh{O}'_{X,s}\otimes_{\sh{O}_{X,s}}\boldsymbol{k}(s)$ est composé direct de corps isomorphes à $\boldsymbol{k}(s)$ (ce qui, lorsque le morphisme $f$ est fini, entraîne qu'il est non ramifié aux points $s'_j$ \sref{IV.17.4.1}).
Si par exemple \emph{aucun} corps résiduel de $X$ n'est isomorphe à $\bb{F}_2$, l'homomorphisme canonique $\operatorname{Div}(X)\to\operatorname{Div}(X')$ ne peut être bijectif que si $f$ est un \emph{isomorphisme}.
Dans le cas où l'homomorphisme canonique $\Gamma(X,\sh{O}_X)\to\Gamma(X',\sh{O}_{X'})$ est bijectif, on conclut de ce qui précède et de (ii) que dans les considérations précédentes, on peut remplacer l'homomorphisme $\operatorname{Div}(X)\to\operatorname{Div}(X')$ par l'homomorphisme $\operatorname{Pic}(X)\to\operatorname{Pic}(X')$.
\end{enumerate}
\end{remark}

\subsection{Diviseurs sur les préschémas de dimension 1}
\EGAReaderTarget{ega4.c04.subsec.21-9}{subsection:IV.21.9}
\label{IV.21.9}

\begin{env}[21.9.1]
\label{IV.21.9.1}
Soient $X$ un espace topologique, $x$ un point de $X$, $i_x:\{x\}\to X$ l'injection canonique.
Si $A(x)$ est un groupe commutatif, on peut le considérer comme un faisceau de groupes commutatifs sur l'espace $\{x\}$ réduit à un seul point, et prendre son image directe $(i_x)_*(A(x))$ qui est un faisceau de groupes commutatifs sur $X$ \sref[0]{0.3.4.1};
il résulte aussitôt des définitions que pour tout ouvert $U$ de $X$, $\Gamma(U,(i_x)_*(A(x)))=A(x)$ si $x\in U$, et est réduit à $0$ dans le cas contraire;
donc, pour $y\in\overline{\{x\}}$ on a $((i_x)_*(A(x)))_y=A(x)$, et pour $y\notin\overline{\{x\}}$, $((i_x)_*(A(x)))_y=0$.

Si maintenant $\sh{F}$ est un faisceau de groupes commutatifs sur $X$ et $Y$ une partie de $X$, pour tout ouvert $U$ de $X$, on a un homomorphisme canonique
\[
\label{IV.21.9.1.1}
  \Gamma(U,\sh{F})\to\prod_{x\in U\cap Y}\sh{F}_x
  =\prod_{x\in Y}\Gamma(U,(i_x)_*(\sh{F}_x))
  \tag{21.9.1.1}
\]
et comme ces homomorphismes commutent aux opérateurs de restriction, ils définissent un homomorphisme canonique de \emph{faisceaux}
\[
\label{IV.21.9.1.2}
  j_Y:\sh{F}\to\prod_{x\in Y}(i_x)_*(\sh{F}_x).
  \tag{21.9.1.2}
\]
\end{env}

\begin{lemma}[21.9.2]
\label{IV.21.9.2}
Soient $X$ un espace topologique localement noethérien, $X_0$ l'ensemble de ses points fermés, $\sh{F}$ un faisceau de groupes commutatifs sur $X$.
Les conditions suivantes sont équivalentes:
\begin{enumerate}[label=\emph{\alph*)}]
  \item L'homomorphisme canonique $j_{X_0}:\sh{F}\to\prod_{x\in X_0}(i_x)_*(\sh{F}_x)$ est injectif et a pour image $\bigoplus_{x\in X_0}(i_x)_*(\sh{F}_x)$.
  \item[\rm a$'$)] Il existe une famille de groupes commutatifs $(A(x))_{x\in X_0}$ telle que $\sh{F}$ soit isomorphe à $\bigoplus_{x\in X_0}(i_x)_*(A(x))$.
  \item Toute section de $\sh{F}$ au-dessus d'un ouvert $U$ de $X$ a un support discret contenu dans $X_0$.
\end{enumerate}
Lorsqu'il en est ainsi, pour tout $x\in X_0$, le groupe $A(x)$ est déterminé à isomorphisme unique près et est isomorphe à $\sh{F}_x$.
En outre, le faisceau $\sh{F}$ est alors flasque.
\end{lemma}

\begin{proof}
\oldpage[IV-4]{285}
Comme les points de $X_0$ sont fermés dans $X$, on a $((i_x)_*(A(x)))_y=0$ pour tout $y\neq x\in X_0$;
cette remarque prouve l'assertion d'unicité relative aux groupes $A(x)$, et il est clair par ailleurs que \emph{a)} implique \emph{a$'$)}.
Pour voir que \emph{a$'$)} implique \emph{b)}, on peut se borner au cas où $U$ est noethérien;
alors (G, II, 3.10) on a
\[
  \Gamma(U,\bigoplus_{x\in X_0}(i_x)_*(A(x)))
  =\bigoplus_{x\in X_0}\Gamma(U,(i_x)_*(A(x))),
\]
donc toute section $s$ de $\sh{F}$ au-dessus de $U$ est somme directe d'un nombre \emph{fini} de sections $s_k\in\Gamma(U,(i_{x_k})_*(A(x_k)))$ ($1\leq k\leq n$) avec $x_k\in X_0$.
Mais puisque $x_k$ est fermé dans $X$, le support de $s_k$ est réduit à $\{x_k\}$, donc le support de $s$ est contenu dans l'ensemble fini des $x_k$, qui est évidemment discret puisque les points $x_k$ sont fermés dans $X$.
Montrons enfin que \emph{b)} entraîne \emph{a)};
pour tout ouvert noethérien $U$, tout support d'une section de $\sh{F}$ au-dessus de $U$ étant discret et quasi-compact, est \emph{fini}.
On déduit donc de l'hypothèse \emph{b)} que pour tout ouvert noethérien $U$, l'image de l'homomorphisme \sref{IV.21.9.1.1} (pour $Y=X_0$) est contenue dans $\bigoplus_{x\in X_0}\Gamma(U,(i_x)_*(\sh{F}_x))$ et que cet homomorphisme est injectif, ce qui établit \emph{a)}.

Enfin, pour montrer que $\sh{F}$ est flasque, considérons une section $s$ de $\sh{F}$ au-dessus d'un ouvert $U$ de $X$;
comme le support de $s$ est un sous-espace discret et fermé dans $X$, on prolonge $s$ en une section $s'$ de $\sh{F}$ au-dessus de $X$ en posant $s'_x=0$ dans $X-U$.
\end{proof}

\begin{remark}[21.9.3]
\label{IV.21.9.3}
\begin{enumerate}[label=(\roman*)]
  \item Dans la condition \emph{b)} de \sref{IV.21.9.2}, il ne suffit pas de supposer que le support de toute section de $\sh{F}$ au-dessus d'un ouvert quelconque de $X$ soit \emph{discret};
c'est ce que montre l'exemple où on prend pour $X$ un spectre d'anneau de valuation discrète, avec un point fermé $b$ et un point générique $a$, et pour $\sh{F}$ le faisceau de groupes commutatifs tel que $\sh{F}_b=0$ et $\sh{F}_a=\bb{Z}$.

  \item Les conditions de \sref{IV.21.9.2} étant supposées vérifiées, soit $E$ une partie discrète de $X_0$, et pour tout $x\in E$, soit $a(x)$ un élément de $A(x)$.
Il existe alors une section $t$ de $\bigoplus_{x\in X_0}(i_x)_*(A(x))$ et une seule au-dessus de $X$ telle que $t_x=a(x)$ pour tout $x\in E$ et que le support de $t$ soit contenu dans $E$.
En effet, pour tout ouvert noethérien $U$, $E\cap U$ est fini, et il suffit de prendre pour $t$ la section de $\bigoplus_{x\in E}(i_x)_*(A(x))$ dont la restriction à chaque ouvert noethérien $U$ est la somme des $a(x)$ pour $x\in E\cap U$.
\end{enumerate}
\end{remark}

\begin{proposition}[21.9.4]
\label{IV.21.9.4}
Soient $X$ un préschéma localement noethérien de dimension $\leq1$, $X^{(1)}$ l'ensemble des points $x\in X$ tels que $\dim(\sh{O}_{X,x})=1$.
On a alors un isomorphisme canonique
\[
\label{IV.21.9.4.1}
  \sh{D}iv_X\simto\bigoplus_{x\in X^{(1)}}(i_x)_*(\operatorname{Div}(\sh{O}_{X,x}))
  \tag{21.9.4.1}
\]
et $\sh{D}iv_X$ est flasque.
\end{proposition}

\begin{proof}
Compte tenu de l'isomorphisme \sref{IV.21.4.6.1}, l'homomorphisme \sref{IV.21.9.4.1} est défini par \sref{IV.21.9.1.2}:
prouvons que c'est une bijection.
Comme $\dim(X)\leq1$, les points de $X^{(1)}$ sont les \emph{points fermés non isolés} de $X$.
On est ramené à prouver que:
1\textsuperscript{o} $\sh{D}iv_X$ vérifie la condition \emph{b)} de \sref{IV.21.9.2};
2\textsuperscript{o} pour tout point isolé $x\in X$, on a $(\sh{D}iv_X)_x=0$.
Le second point résulte de ce qu'alors $\sh{O}_{X,x}$ est un anneau artinien et de ce que tout élément régulier de $\sh{O}_{X,x}$ est donc inversible.
\oldpage[IV-4]{286}
Pour le 1\textsuperscript{o}, il suffit de noter que pour tout ouvert $U$ de $X$ et tout diviseur $D\in\operatorname{Div}(U)$, les points maximaux du support $S$ de $D$ sont tels que $\operatorname{prof}(\sh{O}_{X,x})=1$ \sref{IV.21.1.9}, donc \emph{a fortiori} appartiennent à $X^{(1)}\cap U$;
puisque $\dim(X)\leq1$, l'ensemble de ces points est égal à $S$, et $S$ est donc discret, l'ensemble des composantes irréductibles de $S$ étant localement fini.
\end{proof}

\begin{corollary}[21.9.5]
\label{IV.21.9.5}
Soit $X$ un préschéma localement noethérien de dimension $\leq1$.
Pour toute partie discrète $E\subset X^{(1)}$ et toute famille $(D(x))_{x\in E}$ avec $D(x)\in\operatorname{Div}(\sh{O}_{X,x})$, il existe un diviseur $D$ sur $X$ et un seul tel que le support de $D$ soit contenu dans $E$ et que $D_x=D(x)$ pour tout $x\in E$.
\end{corollary}

\begin{proof}
Cela résulte de \sref{IV.21.9.4} et de \sref{IV.21.9.3}, (ii).
\end{proof}

\begin{corollary}[21.9.6]
\label{IV.21.9.6}
Soit $X$ un préschéma localement noethérien de dimension $\leq1$;
tout diviseur $D$ sur $X$ est de la forme $D'-D''$, où $D'$, $D''$ sont deux diviseurs $\geq0$ de supports contenus dans celui de $D$.
\end{corollary}

\begin{proof}
En vertu de \sref{IV.21.9.5}, on est aussitôt ramené au cas où $X=\Spec(A)$, où $A$ est un anneau local noethérien;
il suffit alors d'observer que $M(X)$ est l'anneau total des fractions de $A$, et qu'une section de $\sh{M}^*_X$ au-dessus de $X$ s'écrit par définition comme quotient $b/a$ de deux éléments réguliers de $A$.
\end{proof}

\begin{corollary}[21.9.7]
\label{IV.21.9.7}
Soient $X$ un préschéma localement noethérien de dimension $\leq1$, sans point isolé, et $U$ un ouvert dense dans $X$.
Il existe alors un diviseur $D\geq0$ sur $X$, de support contenu dans $U$ et rencontrant chacune des composantes irréductibles de $X$.
\end{corollary}

\begin{proof}
On peut supposer que $U$ est réunion d'ensembles ouverts disjoints non vides $U_\alpha$, dont chacun est contenu dans une seule composante irréductible de $X$;
il suffit alors de prendre dans chaque $U_\alpha$ un point $x_\alpha$ fermé dans $X$ (il en existe puisque l'unique point générique de $U_\alpha$ ne peut être isolé par hypothèse), et d'appliquer \sref{IV.21.9.5} à l'ensemble discret des $x_\alpha$.
\end{proof}

L'intérêt de ce corollaire est qu'il permettra de prouver qu'une courbe algébrique séparée $X$ sur un corps $k$ est quasi-projective, le diviseur $D\geq0$ défini dans \sref{IV.21.9.7} étant alors nécessairement \emph{ample} en vertu du théorème de Riemann--Roch pour les courbes (chap.~V).

\begin{env}[21.9.8]
\label{IV.21.9.8}
Pour les préschémas localement noethériens $X$ de dimension $\leq1$ la proposition \sref{IV.21.9.4} ramène la détermination de $\sh{D}iv_X$ au cas où $X=\Spec(A)$, $A$ étant un anneau local noethérien de dimension $1$.

Lorsque $A$ est un anneau local régulier de dimension $1$ (autrement dit, un anneau de valuation discrète), le groupe $\operatorname{Div}(A)$ s'identifie canoniquement à $\bb{Z}$ \sref{IV.21.6.8};
par suite, en vertu de \sref{IV.21.9.2}:
\end{env}

\begin{proposition}[21.9.9]
\label{IV.21.9.9}
Soit $X$ un préschéma localement noethérien régulier de dimension $\leq1$.
Alors le faisceau $\sh{D}iv_X$ est canoniquement isomorphe à $\bigoplus_{x\in X^{(1)}}(i_x)_*(\bb{Z}(x))$, où $\bb{Z}(x)$ est le groupe additif $\bb{Z}$ considéré comme faisceau de groupes sur l'espace $\{x\}$.
\end{proposition}

\begin{env}[21.9.10]
\label{IV.21.9.10}
Supposons seulement que l'anneau local noethérien $A$ de dimension $1$ soit \emph{réduit};
alors, si $A'$ est le \emph{normalisé} de $A$ (fermeture intégrale de $A$ dans son anneau total de fractions), $A'$ est noethérien en vertu du théorème de Krull--Akizuki (Bourbaki, \emph{Alg.\ comm.}, chap.~VII, \textsection~2, n\textsuperscript{o}~5, prop.~5), et on a vu dans \sref{IV.21.8.6} comment
\oldpage[IV-4]{287}
$\operatorname{Div}(A)$ peut s'obtenir comme \emph{extension} de $\operatorname{Div}(A')$ et ce dernier groupe est de la forme $\bb{Z}^r$.
\end{env}

\begin{proposition}[21.9.11]
\label{IV.21.9.11}
Soient $X$ un préschéma noethérien, $X_0$ un sous-préschéma fermé de $X$ ayant les propriétés suivantes:
\begin{enumerate}
  \item[\rm 1\textsuperscript{o}] $\dim(X_0)\leq1$.
  \item[\rm 2\textsuperscript{o}] Pour toute partie localement fermée $Y$ de $X$ telle que $Y\cap X_0$ soit discret, il existe une partie $Y'$ de $Y$, fermée dans $X$ et ouverte dans $Y$, contenant $Y\cap X_0$.
\end{enumerate}
Dans ces conditions:
\begin{enumerate}[label=(\roman*)]
  \item Soit $Z_0$ la réunion des ensembles $\overline{\{x\}}\cap X_0$ lorsque $x$ parcourt la partie de $\operatorname{Ass}(\sh{O}_X)$ formée des points tels que $\overline{\{x\}}\cap X_0$ soit fini.
Alors, pour tout diviseur $D_0$ sur $X_0$, dont le support ne rencontre pas $Z_0$, il existe un diviseur $D$ sur $X$ dont l'image réciproque par l'injection canonique $X_0\to X$ existe \sref{IV.21.4} et est égale à $D_0$;
si de plus $D_0\geq0$, on peut supposer $D\geq0$.

  \item Supposons en outre qu'il existe dans $X_0$ un ouvert affine $U_0$ contenant $\operatorname{Ass}(\sh{O}_{X_0})\cup Z_0$ (condition automatiquement vérifiée lorsqu'il existe un $\sh{O}_{X_0}$-Module ample \sref[II]{II.4.5.4}).
Alors l'homomorphisme canonique \sref{IV.21.3.2.4} $\operatorname{Pic}(X)\to\operatorname{Pic}(X_0)$ est surjectif.
\end{enumerate}
\end{proposition}

\begin{proof}
\begin{enumerate}[label=(\roman*)]
  \item Compte tenu de \sref{IV.21.9.6}, on peut se borner à prouver l'assertion relative aux diviseurs $D_0\geq0$;
en vertu de \sref{IV.21.9.4}, le support $T$ de $D_0$ est un ensemble discret et fermé dans $X_0$.
On peut supposer $D_0\neq0$, c'est-à-dire $T\neq\emp$, sans quoi il n'y a rien à démontrer.
Pour tout $x\in T$, il existe un élément $s_x\in\sh{O}_{X_0,x}$ qui est non diviseur de zéro dans cet anneau, appartient à son idéal maximal, et dont l'image dans $(\sh{D}iv_{X_0})_x$ est $(D_0)_x$.
Il existe un voisinage ouvert affine $U^{(x)}$ de $x$ dans $X$ et une section $g^{(x)}$ de $\sh{O}_X$ au-dessus de $U^{(x)}$ tels que $s_x$ soit l'image dans $\sh{O}_{X_0,x}$ du germe $(g^{(x)})_x\in\sh{O}_{X,x}$;
nous allons voir qu'en prenant $U^{(x)}$ assez petit, on peut faire en sorte que $g^{(x)}$ ne soit pas diviseur de zéro dans $\Gamma(U^{(x)},\sh{O}_X)$, autrement dit \sref{IV.3.1.9}, en désignant par $V^{(x)}$ la partie fermée de $U^{(x)}$ formée des $y$ tels que $g^{(x)}(y)=0$, que l'on ait $V^{(x)}\cap\operatorname{Ass}(\sh{O}_X)=\emp$.

Or, si $z\in\operatorname{Ass}(\sh{O}_X)$, l'hypothèse $x\notin Z_0$ entraîne que l'on a, ou bien $x\notin\overline{\{z\}}$, ou bien que $x$ n'est pas isolé dans $\overline{\{z\}}\cap X_0$.
En remplaçant $U^{(x)}$ par un ouvert plus petit, on peut supposer que c'est le second cas qui se produit pour un $z\in V^{(x)}\cap\operatorname{Ass}(\sh{O}_X)$.
$V^{(x)}$ contiendrait donc une composante irréductible de dimension $1$ de $X_0\cap U^{(x)}$, contenant $x$;
mais cela signifierait que l'image $\overline{g}^{(x)}$ de $g^{(x)}$ dans $\Gamma(U^{(x)}\cap X_0,\sh{O}_{X_0})$ serait dans le nilradical de cet anneau, et par suite $s_x$ appartiendrait au nilradical de $\sh{O}_{X_0,x}$, ce qui est absurde.
On a donc $x\notin\overline{\{z\}}$ pour tout point $z\in V^{(x)}\cap\operatorname{Ass}(\sh{O}_X)$, et comme cet ensemble est fini, on peut, en remplaçant $U^{(x)}$ par un voisinage ouvert plus petit de $x$, supposer que $V^{(x)}\cap\operatorname{Ass}(\sh{O}_X)=\emp$.

En vertu du choix de $U^{(x)}$, on peut définir un diviseur $D'^{(x)}$ sur $U^{(x)}$ par $D'^{(x)}=\operatorname{div}(g^{(x)})$;
d'ailleurs on a vu ci-dessus que $x$ est nécessairement isolé dans $V^{(x)}\cap X_0$, donc en remplaçant encore $U^{(x)}$ par un voisinage ouvert plus petit de $x$, on peut supposer que $V^{(x)}\cap X_0$ est réduit au point $x$.
Mais il existe alors, en vertu de la condition 2\textsuperscript{o}, une partie $W^{(x)}$ de $V^{(x)}$, fermée dans $X$ et ouverte dans $V^{(x)}$, telle que $W^{(x)}\cap X_0=\{x\}$.
Remplaçant encore $U^{(x)}$ par un voisinage ouvert plus petit de $x$, on peut donc supposer
\oldpage[IV-4]{288}
que $W^{(x)}=V^{(x)}$, autrement dit que $V^{(x)}$ est fermé dans $X$.
On peut alors définir un diviseur $D^{(x)}$ sur $X$ par la condition que $D^{(x)}|U^{(x)}=D'^{(x)}$, et $D^{(x)}|(X-V^{(x)})=0$, ce qui a un sens puisque dans $U^{(x)}\cap(X-V^{(x)})$ la restriction de $g^{(x)}$ est une section inversible, donc les restrictions à cet ouvert des deux diviseurs considérés sur $U^{(x)}$ et $X-V^{(x)}$ sont égales.
Il suffit alors, pour répondre à la question, de prendre $D=\sum_{x\in T}D^{(x)}$, ce qui a un sens puisque $T$ est fini.

  \item Compte tenu du diagramme commutatif \sref{IV.21.4.2.1}, la conclusion résultera de (i) si l'on prouve que tout $\sh{O}_{X_0}$-Module inversible $\sh{L}_0$ est isomorphe à un $\sh{O}_{X_0}$-Module de la forme $\sh{O}_{X_0}(D_0)$ \sref{IV.21.2.8}, où $D_0$ est un diviseur sur $X_0$ dont le support ne rencontre pas $Z_0$.
Or, comme $U_0$ est schématiquement dense dans $X_0$ \sref{IV.20.2.13}, (iv), il suffit pour cela de prouver qu'il existe une section $t\in\Gamma(U_0,\sh{L}_0)$ telle que $t(z_0)\neq0$ aux points de $\operatorname{Ass}(\sh{O}_{X_0})\cup Z_0$ \sref{IV.3.1.9}.
On peut donc se borner au cas où $X_0=\Spec(A_0)$ est affine;
mais alors $\sh{L}_0$ est ample \sref[II]{II.5.1.4} et comme l'ensemble $\operatorname{Ass}(\sh{O}_{X_0})\cup Z_0$ est fini, la conclusion résulte de \sref[II]{II.4.5.4}.
\end{enumerate}
\end{proof}

\begin{corollary}[21.9.12]
\label{IV.21.9.12}
Soient $A$ un anneau local hensélien \sref{IV.18.5.8}, $S=\Spec(A)$, $s_0$ le point fermé de $S$, $f:X\to S$ un morphisme séparé de présentation finie, et supposons que l'ensemble $X_0=f^{-1}(s_0)$ soit de dimension $\leq1$.
Alors, pour tout sous-schéma fermé $X'_0$ de $X$, ayant $X_0$ pour ensemble sous-jacent et de présentation finie sur $S$, l'homomorphisme canonique \sref{IV.21.3.2.4} $\operatorname{Pic}(X)\to\operatorname{Pic}(X'_0)$ est surjectif.
\end{corollary}

\begin{proof}
Montrons d'abord qu'il suffit de prouver le corollaire lorsque l'anneau local hensélien $A$ est de plus noethérien.
On sait en effet \sref{IV.18.6.15} que l'on peut écrire $A=\varinjlim A_\lambda$, où les $A_\lambda$ sont des anneaux locaux noethériens et henséliens, les homomorphismes $A_\lambda\to A$ étant locaux;
il existe en outre un indice $\alpha$ et un morphisme séparé de présentation finie $f_\alpha:X_\alpha\to S_\alpha=\Spec(A_\alpha)$ tels que $X$ et $f$ soient déduits de $X_\alpha$ et $f_\alpha$ par le changement de base $S\to S_\alpha$ \sref{IV.8.10.5}, (v);
avec les notations habituelles \sref{IV.8.8.1}, les morphismes $f_\lambda:X_\lambda\to S_\lambda$ seront séparés pour $\lambda\geq\alpha$ et on aura $X=X_\lambda\times_{S_\lambda}S$.
En outre, on peut supposer que, si $s_{0\alpha}$ est l'unique point fermé de $S_\alpha$, il y a un sous-schéma fermé $X'_{0\alpha}$ de $X_\alpha$, ayant même ensemble sous-jacent que $f_\alpha^{-1}(s_{0\alpha})$, tel que $X'_0=X'_{0\alpha}\times_{S_\alpha}S$ \sref{IV.8.6.3};
pour $\lambda\geq\alpha$, posons $X'_{0\lambda}=X'_{0\alpha}\times_{S_\alpha}S_\lambda$;
on a $\dim(X'_{0\lambda})=\dim(X'_0)\leq1$ par transitivité des fibres et \sref{IV.4.1.4}.
Il s'agit de voir que si l'on a prouvé que l'homomorphisme $\operatorname{Pic}(X_\lambda)\to\operatorname{Pic}(X'_{0\lambda})$ est surjectif pour tout $\lambda\geq\alpha$, il en est de même de $\operatorname{Pic}(X)\to\operatorname{Pic}(X'_0)$.
On a évidemment le diagramme commutatif d'homomorphismes canoniques
\[
  \xymatrix{
    \operatorname{Pic}(X_\lambda) \ar[r] \ar[d] & \operatorname{Pic}(X) \ar[d]\\
    \operatorname{Pic}(X'_{0\lambda}) \ar[r] & \operatorname{Pic}(X'_0)
  }
\]
Pour tout $\sh{O}_{X'_0}$-Module inversible $\sh{L}'_0$, il existe un $\lambda\geq\alpha$ et un $\sh{O}_{X'_{0\lambda}}$-Module inversible $\sh{L}'_{0\lambda}$ tel que $\sh{L}'_0$ s'en déduise par changement de base $S\to S_\lambda$ \sref{IV.8.5.2} et \sref{IV.8.5.5};
il suffit de considérer un $\sh{O}_{X_\lambda}$-Module inversible $\sh{L}_\lambda$ tel que $\operatorname{cl}(\sh{L}'_{0\lambda})$ soit l'image de $\operatorname{cl}(\sh{L}_\lambda)$ dans $\operatorname{Pic}(X'_{0\lambda})$, puis de prendre le $\sh{O}_X$-Module inversible $\sh{L}$ déduit de $\sh{L}_\lambda$ par changement de base;
en vertu de la commutativité du diagramme précédent, $\operatorname{cl}(\sh{L}'_0)$ sera l'image de $\operatorname{cl}(\sh{L})$.

Supposons donc $A$ noethérien, donc aussi $X$, et vérifions que $X$ et $X'_0$ satisfont aux conditions de \sref{IV.21.9.11}, (ii).
On a par hypothèse $\dim(X'_0)\leq1$;
d'autre part, pour vérifier la condition 2\textsuperscript{o} de \sref{IV.21.9.11}, considérons un sous-préschéma réduit $Y$ de $X$ ayant une partie localement fermée $Y$ pour ensemble sous-jacent;
le morphisme $g:Y\to S$, restriction de $f$, étant quasi-fini en chacun des points $x_i$ de $Y\cap X_0$ ($1\leq i\leq n$), on peut appliquer \sref{IV.18.5.11}, c), et on voit que $Y$ est somme des sous-préschémas ouverts $Y_i=\Spec(\sh{O}_{Y,x_i})$ ($1\leq i\leq n$) et d'un préschéma $Y''$ tel que $Y''\cap X_0=\emp$, et en outre que les injections canoniques $j_i:Y_i\to X$ sont telles que $f\circ j_i$ soit un morphisme fini.
Comme $f$ est séparé, $j_i$ est aussi un morphisme fini, donc $Y_i$ est fermé dans $X$;
on répond donc à la question en prenant $Y'=\bigcup_{1\leq i\leq n}Y_i$.

Il reste à vérifier l'hypothèse supplémentaire de \sref{IV.21.9.11}, (ii).
Or, $X'_0$ étant une courbe séparée sur le corps $\boldsymbol{k}(s_0)$, est un $\boldsymbol{k}(s_0)$-schéma quasi-projectif (chap.~V)\footnote{Le lecteur vérifiera que \sref{IV.21.9.12} n'est pas utilisé pour prouver cette propriété au chap.~V.}, donc il existe un $\sh{O}_{X'_0}$-Module ample \sref[II]{II.5.3.1}, et cela achève la démonstration.
\end{proof}

\begin{remark}[21.9.13]
\label{IV.21.9.13}
Sous les conditions de \sref{IV.21.9.12}, supposant en outre $f$ propre, le morphisme $f$ est même \emph{projectif}:
en effet, si $\sh{L}_0$ est un $\sh{O}_{X'_0}$-Module ample, il existe, en vertu de \sref{IV.21.9.12}, un $\sh{O}_X$-Module inversible $\sh{L}$ dont l'image réciproque dans $X'_0$ est isomorphe à $\sh{L}_0$, donc est ample.
Puisque tout voisinage de $s_0$ dans $S$ est nécessairement $S$ tout entier, on déduit alors de \sref{IV.9.6.4} que $\sh{L}$ est un $\sh{O}_X$-Module ample, d'où la conclusion (\sref[II]{II.5.3.1} et \sref[II]{II.5.5.3}).
\end{remark}

\subsection{Images réciproques et images directes de cycles 1-codimensionnels}
\EGAReaderTarget{ega4.c04.subsec.21-10}{subsection:IV.21.10}

Dans un chapitre ultérieur, consacré à la théorie des intersections, seront développées de façon systématique les notions d'image inverse et d'image directe de cycles.
Dans le présent numéro, nous nous contenterons de définir ces notions dans certains cas particuliers utiles, et pour les cycles $1$-codimensionnels, ces définitions étant choisies de telle sorte qu'elles soient compatibles avec les définitions correspondantes pour les diviseurs \sref{IV.21.4} et \sref{IV.21.5}, compte tenu de l'homomorphisme $\operatorname{cyc}$ défini en \sref{IV.21.6}.

\begin{env}[21.10.1]
\label{IV.21.10.1}
\oldpage[IV-4]{289}
Soient $X$, $X'$ deux préschémas localement noethériens, $f:X'\to X$ un morphisme, $T$ une partie de $X$;
supposons que l'image par $f$ de tout point maximal de $X'$ soit un point maximal de $X$, et que, pour tout $x'\in X'^{(1)}$ (i.e.\ tel que $\dim(\sh{O}_{X',x'})=1$), on ait l'une des trois conditions suivantes pour le point $x=f(x')$:
\begin{enumerate}[label=(\roman*)]
  \item $x\notin T$;
  \item $x\in X^{(1)}$ et $\sh{O}_{X',x'}$ est un $\sh{O}_{X,x}$-module plat;
  \item l'anneau $\sh{O}_{X,x}$ est factoriel et $\mathfrak{m}_{x'}\notin\operatorname{Ass}(\sh{O}_{X',x'})$.
\end{enumerate}
Sous ces conditions, nous nous proposons, pour tout cycle $1$-codimensionnel $Z$ sur $X$ de support contenu dans $T$, de définir un cycle $1$-codimensionnel $Z'=f^*(Z)$, de sorte que $Z\mapsto f^*(Z)$ soit un homomorphisme de groupes ordonnés du sous-groupe de $\mathfrak{Z}^1(X)$ formé des cycles de support contenu dans $T$, dans le groupe ordonné $\mathfrak{Z}^1(X')$.
Pour cela, soit
\[
\label{IV.21.10.1.1}
  Z=\sum_{x\in T\cap X^{(1)}}n_x\cdot\overline{\{x\}}
  \tag{21.10.1.1}
\]
où la famille des $x\in T\cap X^{(1)}$ tels que $n_x\neq0$ est localement finie.
Pour tout $x'\in X'^{(1)}$, définissons un entier $n_{x'}$ de la façon suivante, en posant $x=f(x')$:
\begin{enumerate}[label=\arabic*\textsuperscript{o}]
  \item si $x\notin T$, on prend $n_{x'}=0$;
  \item si $x\in X^{(1)}$ et si $\sh{O}_{X',x'}$ est un $\sh{O}_{X,x}$-module plat, on sait \sref{IV.6.1.1} que $\dim(\sh{O}_{X',x'}/\mathfrak{m}_x\sh{O}_{X',x'})=0$, autrement dit $\sh{O}_{X',x'}/\mathfrak{m}_x\sh{O}_{X',x'}$ est un $\sh{O}_{X',x'}$-module de longueur finie $\lambda_{x'}$, et on prend $n_{x'}=\lambda_{x'}n_x$;
  \item si $\sh{O}_{X,x}$ est factoriel et $\mathfrak{m}_{x'}\notin\operatorname{Ass}(\sh{O}_{X',x'})$, on sait \sref{IV.21.6.9} que l'homomorphisme canonique $\operatorname{cyc}:\operatorname{Div}(\sh{O}_{X,x})\to\mathfrak{Z}^1(\Spec(\sh{O}_{X,x}))$ est bijectif, et d'autre part comme $\dim(\sh{O}_{X',x'})=1$ et $\mathfrak{m}_{x'}\notin\operatorname{Ass}(\sh{O}_{X',x'})$, $\operatorname{Ass}(\sh{O}_{X',x'})$ se compose uniquement des points maximaux de $\Spec(\sh{O}_{X',x'})$, donc l'hypothèse sur $f$ implique que $f(\operatorname{Ass}(\sh{O}_{X',x'}))\subset\operatorname{Ass}(\sh{O}_{X,x})$ et il résulte de \sref{IV.21.4.5}, (ii) que l'homomorphisme $f^*:\operatorname{Div}(\sh{O}_{X,x})\to\operatorname{Div}(\sh{O}_{X',x'})$ est défini;
enfin, $x'$ étant l'unique point fermé de $\Spec(\sh{O}_{X',x'})$, $\mathfrak{Z}^1(\Spec(\sh{O}_{X',x'}))$ est canoniquement isomorphe à $\bb{Z}$.
\end{enumerate}
On a donc un homomorphisme canonique composé
\[
\label{IV.21.10.1.2}
  \mathfrak{Z}^1(\Spec(\sh{O}_{X,x}))
  \xrightarrow{\operatorname{cyc}^{-1}}\operatorname{Div}(\sh{O}_{X,x})
  \xrightarrow{f^*}\operatorname{Div}(\sh{O}_{X',x'})
  \xrightarrow{\operatorname{cyc}}\mathfrak{Z}^1(\Spec(\sh{O}_{X',x'}))
  \xrightarrow{\sim}\bb{Z}
  \tag{21.10.1.2}
\]
Si $Z_x$ est le cycle $\sum_{y\in\Spec(\sh{O}_{X,x})\cap X^{(1)}}n_y(\overline{\{y\}}\cap\Spec(\sh{O}_{X,x}))$ sur $\Spec(\sh{O}_{X,x})$, on prend alors pour $n_{x'}$ l'\emph{image} de $Z_x$ par l'homomorphisme \sref{IV.21.10.1.2}.
\end{env}

\oldpage[IV-4]{290}
\begin{env}[21.10.2]
\label{IV.21.10.2}
Nous nous proposons de montrer que:
\begin{enumerate}[label=\emph{\Alph*)}]
  \item Lorsque deux des conditions 1\textsuperscript{o}, 2\textsuperscript{o}, 3\textsuperscript{o} de \sref{IV.21.10.1} sont \emph{simultanément} satisfaites, les valeurs correspondantes de $n_{x'}$ \emph{coïncident}.
  \item L'ensemble des $x'\in X'^{(1)}$ tels que $n_{x'}\neq0$ est \emph{localement fini} dans $X'$.
\end{enumerate}

Pour démontrer A), supposons d'abord que $x\notin T$ et que $x$ vérifie l'une des conditions 2\textsuperscript{o} ou 3\textsuperscript{o} de \sref{IV.21.10.1};
alors $n_x=0$ et si l'on est dans le cas 2\textsuperscript{o}, on a $n_{x'}=0$;
si on est dans le cas 3\textsuperscript{o}, on a $Z_x=0$ puisque $\Supp(Z)\subset T$, donc encore $n_{x'}=0$.
Il reste à considérer le cas où l'on est à la fois dans le cas 2\textsuperscript{o} et dans le cas 3\textsuperscript{o};
alors, puisque $\dim(\sh{O}_{X,x})=1$, $\sh{O}_{X,x}$ est un anneau de valuation discrète;
si $t$ est une uniformisante de cet anneau, le diviseur correspondant à $Z_x$ est $\operatorname{div}(t^{n_x})$ dans $\operatorname{Div}(\sh{O}_{X,x})$, et son image dans $\operatorname{Div}(\sh{O}_{X',x'})$ est le diviseur de $t'^{n_x}$, où $t'$ est l'élément (régulier) de $\sh{O}_{X',x'}$ image de $t$.
On peut évidemment se borner au cas où $n_x=1$, et alors la définition \sref{IV.21.6.5.1} montre que l'image de $Z_x$ par \sref{IV.21.10.1.2} est le nombre $\lambda_{x'}$, ce qui achève de prouver A).

Démontrons maintenant B).
Posons $T_0=\Supp(Z)\subset T$, $T'_0=f^{-1}(T_0)$;
il suffit de prouver que la relation $n_{x'}\neq0$ implique que $x'$ appartient à l'ensemble des points maximaux de $T'_0$, ce dernier étant localement fini dans $X'$.
Il est immédiat que l'on a nécessairement $x'\in T'_0$;
si $x'$ n'était pas maximal dans l'ensemble fermé $T'_0$, il existerait
\oldpage[IV-4]{291}
une générisation $y'$ de $x'$ dans $T'_0$, distincte de $x'$, et puisque $x'\in X'^{(1)}$, $y'$ serait nécessairement un point maximal de $X'$;
par suite $y=f(y')$ serait point maximal de $X$ par hypothèse;
mais cela est absurde puisque $y\in T_0$ et que $T_0$ est purement de codimension $1$ dans $X$.
\end{env}

\begin{env}[21.10.3]
\EGAReaderTarget{ega4.c04.statement.21-10-3}{IV.21.10.3}
On peut maintenant poser
\[
  f^*(Z)=\sum_{x'\in X'^{(1)}}n_{x'}\overline{\{x'\}},
\]
la somme du second membre ayant un sens en vertu de ce qu'on a prouvé dans \sref{IV.21.10.2};
on dit que le cycle $1$-codimensionnel $f^*(Z)$ est l'\emph{image réciproque} de $Z$ par $f$.
Il est immédiat que l'application $f^*$ ainsi définie est un homomorphisme de groupes ordonnés.
En outre, si $U$ est un ouvert de $X$, $V$ un ouvert de $X'$ tel que $f(V)\subset U$ et $f':V\to U$ la restriction de $f$, il résulte aussitôt des définitions que l'on a
\[
\label{IV.21.10.3.1}
  f'^*(Z|U)=f^*(Z)|V.
  \tag{21.10.3.1}
\]

Désignons par $\Gamma_T(\mathscr{Z}^1_X)$ le plus grand sous-faisceau de groupes commutatifs de $\mathscr{Z}^1_X$ de support contenu dans $T$;
il résulte de la relation \sref{IV.21.10.3.1} que les applications $\Gamma(U,\Gamma_T(\mathscr{Z}^1_X))\to\Gamma(V,\mathscr{Z}^1_{X'})$ que l'on vient de définir, définissent un homomorphisme de faisceaux de groupes commutatifs ordonnés sur $X'$
\[
  \psi^*(\Gamma_T(\mathscr{Z}^1_X))\to\mathscr{Z}^1_{X'}
\]
où $\psi$ est l'application continue sous-jacente au morphisme $f$.
\end{env}

\begin{proposition}[21.10.4]
\label{IV.21.10.4}
Supposons vérifiées les conditions de \sref{IV.21.10.1}.
Alors, pour tout diviseur $D$ sur $X$ tel que $\Supp(D)\subset T$ et que $f^*(D)$ soit défini \sref{IV.21.4.2}, on a
\[
\label{IV.21.10.4.1}
  \operatorname{cyc}(f^*(D))=f^*(\operatorname{cyc}(D)).
  \tag{21.10.4.1}
\]
\end{proposition}

\begin{proof}
La question étant locale sur $X$, on peut se borner au cas où $X=\Spec(A)$ est affine, où $D=\operatorname{div}(t)$, où $t$ est un élément régulier non inversible de $A$, et où le sous-préschéma $Y(D)$ \sref{IV.21.2.12} n'a qu'un seul point maximal $y$, de sorte que $\operatorname{cyc}(D)=n_y\cdot\overline{\{y\}}$, où $n_y$ est la longueur de $\sh{O}_{X,y}/t_y\sh{O}_{X,y}$ \sref{IV.21.6.5.1}.
On a vu dans la démonstration de \sref{IV.21.10.2}, B) que les points $x'\in X'$ tels que $\operatorname{mult}_{x'}(f^*(\operatorname{cyc}(D)))\neq0$ sont des points maximaux de $f^{-1}(\overline{\{y\}})$.
Si le point $x'$ est dans le cas 3\textsuperscript{o} de \sref{IV.21.10.1}, l'égalité des multiplicités en $x'$ des deux membres de \sref{IV.21.10.4.1} résulte de la définition de $n_{x'}$ au moyen de l'homomorphisme \sref{IV.21.10.1.2}.
Supposons au contraire que $x'$ soit dans le cas 2\textsuperscript{o} de \sref{IV.21.10.1}, et soit $x=f(x')$;
puisque $x\in X^{(1)}\cap\overline{\{y\}}$, on a nécessairement $x=y$.
Remarquons maintenant que $f^*(D)=\operatorname{div}(t')$, où $t'$ est l'image de $t$ dans $\Gamma(X',\sh{O}_{X'})$, et $Y(f^*(D))=f^{-1}(Y(D))$;
la multiplicité de $f^*(D)$ au point $x'$ est donc la longueur $n'_{x'}$ du $\sh{O}_{X',x'}$-module
\[
  \sh{O}_{X',x'}/t'_{x'}\sh{O}_{X',x'}
  =(\sh{O}_{X,y}/t_y\sh{O}_{X,y})\otimes_{\sh{O}_{X,y}}\sh{O}_{X',x'};
\]
comme il résulte de \sref{IV.4.7.1} que l'on a $n'_{x'}=\lambda_{x'}n_y$, les multiplicités au point $x'$ des deux membres de \sref{IV.21.10.4.1} sont encore égales, ce qui achève la démonstration.
\end{proof}

\begin{env}[21.10.5]
\label{IV.21.10.5}
\oldpage[IV-4]{292}
Supposons maintenant que $f:X'\to X$ soit un morphisme de préschémas localement noethériens, transformant tout point maximal de $X'$ en point maximal de $X$, et supposons en outre que pour toute partie fermée rare $T$ de $X$, $f$ vérifie les conditions de \sref{IV.21.10.1};
cela signifie encore que pour tout $x'\in X'^{(1)}$, ou bien $x=f(x')$ est point maximal de $X$, ou bien $x'$ vérifie l'une des conditions (ii), (iii) de \sref{IV.21.10.1}.
Si on tient compte de ce que tout cycle $1$-codimensionnel a un support rare dans $X$, on voit que $f^*(Z)$ est défini pour tout cycle $1$-codimensionnel $Z$ sur $X$;
en vertu de \sref{IV.21.10.3.1}, on a donc défini ainsi un homomorphisme de faisceaux de groupes commutatifs ordonnés
\[
  \psi^*(\mathscr{Z}^1_X)\to\mathscr{Z}^1_{X'}.
\]
Si de plus, pour tout diviseur $D$ sur $X$, $f^*(D)$ est défini \sref{IV.21.4.5}, le fait que le support de $D$ soit rare dans $X$ \sref{IV.21.6.6} entraîne que \sref{IV.21.10.4} est applicable, et on a donc la formule \sref{IV.21.10.4.1} pour tout diviseur $D$ sur $X$.
En particulier:
\end{env}

\begin{proposition}[21.10.6]
\label{IV.21.10.6}
Soient $X$, $X'$ deux préschémas localement noethériens, $f:X'\to X$ un morphisme plat.
Alors $f^*(Z)$ est défini pour tout cycle $1$-codimensionnel $Z$ sur $X$, $f^*(D)$ est défini pour tout diviseur $D$ sur $X$ et l'on a la relation \sref{IV.21.10.4.1}.
\end{proposition}

\begin{proof}
En effet, si $x'\in X'^{(1)}$ est tel que $x=f(x')$ ne soit pas maximal, il résulte de \sref{IV.6.1.1} que l'on a nécessairement $x\in X^{(1)}$, donc on est dans le cas (ii) de \sref{IV.21.10.1}.
On peut donc appliquer \sref{IV.21.10.5}, compte tenu de \sref{IV.21.4.5} et de \sref{IV.2.3.4}.
\end{proof}

\begin{remark}[21.10.7]
\label{IV.21.10.7}
L'existence de $f^*(Z)$ pour tout cycle $1$-codimensionnel $Z$ sur $X$ résulte déjà de l'hypothèse que $f$ est \emph{plat aux points $x'$ de $X'$ de codimension $\leq1$ dans $X'$} (i.e.\ tels que $\dim\sh{O}_{X',x'}\leq1$);
en effet, pour tout point maximal $z'\in X'$, il résulte de \sref{IV.6.1.1} que $z=f(z')$ est point maximal de $X$ puisque $\dim\sh{O}_{X,z}\leq\dim\sh{O}_{X',z'}=0$.
De même, si $x'\in X'^{(1)}$, $x=f(x')$ est maximal ou appartient à $X^{(1)}$ par \sref{IV.6.1.1};
on peut donc encore appliquer \sref{IV.21.10.5}.
\end{remark}

\begin{proposition}[21.10.8]
\label{IV.21.10.8}
Soient $X$, $X'$, $X''$ trois préschémas localement noethériens, $f:X'\to X$, $g:X''\to X'$ deux morphismes;
on suppose que $f$ (resp.\ $g$) est plat en tout point de codimension $\leq1$ dans $X'$ (resp.\ $X''$).
Alors $f\circ g$ est plat en tout point de codimension $\leq1$ dans $X''$ et pour tout cycle $1$-codimensionnel $Z$ sur $X$, on a $(f\circ g)^*(Z)=g^*(f^*(Z))$.
\end{proposition}

\begin{proof}
La première assertion résulte de \sref{IV.6.1.1} et \sref{IV.2.1.6}.
La seconde résulte de ce que $f$ (resp.\ $g$, resp.\ $f\circ g$) transforme les points maximaux de $X'$ (resp.\ $X''$, resp.\ $X''$) en points maximaux de $X$ (resp.\ $X'$, resp.\ $X$) et de ce que, si $x''\in X''^{(1)}$ est tel que $x'=g(x'')\in X'^{(1)}$ et $x=f(x')\in X^{(1)}$, on a
\[
  \operatorname{long}(\sh{O}_{X'',x''}/\mathfrak{m}_x\sh{O}_{X'',x''})
  =\operatorname{long}(\sh{O}_{X'',x''}/\mathfrak{m}_{x'}\sh{O}_{X'',x''})
   \cdot\operatorname{long}(\sh{O}_{X',x'}/\mathfrak{m}_x\sh{O}_{X',x'}).
\]
En effet, si l'on pose $A=\sh{O}_{X,x}$, $A'=\sh{O}_{X',x'}$, $A''=\sh{O}_{X'',x''}$, $\mathfrak{m}=\mathfrak{m}_x$, $\mathfrak{m}'=\mathfrak{m}_{x'}$, on a $A''/\mathfrak{m}A''=(A'/\mathfrak{m}A')\otimes_{A'}A''$, et la formule
\[
  \operatorname{long}_{A''}(A''/\mathfrak{m}A'')
  =\operatorname{long}_{A'}(A'/\mathfrak{m}A')
   \cdot\operatorname{long}_{A''}(A''/\mathfrak{m}'A'')
\]
résulte de \sref{IV.4.7.1} et de l'hypothèse de platitude de $A''$ sur $A'$.
\end{proof}

\begin{env}[21.10.9]
\EGAReaderTarget{ega4.c04.statement.21-10-9}{IV.21.10.9}
Soient $X$ un préschéma localement noethérien, $\Lambda$ un groupe commutatif, noté additivement, et considéré donc comme $\bb{Z}$-module;
on notera encore par $\bb{Z}$ et $\Lambda$ les faisceaux simples sur $X$ associés aux préfaisceaux constants égaux respectivement à $\bb{Z}$ et $\Lambda$ \sref[0]{0.3.6}.
On appelle \emph{faisceau des germes de cycles $1$-codimensionnels à coefficients dans $\Lambda$} le faisceau de groupes commutatifs $\mathscr{Z}^1_X\otimes_{\bb{Z}}\Lambda$.
Si $\Lambda=\bb{Q}$, on dit que les sections
\oldpage[IV-4]{293}
de $\mathscr{Z}^1_X\otimes_{\bb{Z}}\bb{Q}$ au-dessus de $X$ sont les \emph{cycles $1$-codimensionnels à coefficients rationnels}.
Comme les fibres de $\mathscr{Z}^1_X$ sont des $\bb{Z}$-modules sans torsion \sref{IV.21.6.3}, l'homomorphisme canonique $\mathscr{Z}^1_X\to\mathscr{Z}^1_X\otimes_{\bb{Z}}\bb{Q}$ est \emph{injectif}, de sorte que les cycles $1$-codimensionnels s'identifient à des cycles $1$-codimensionnels à coefficients rationnels.
\end{env}

\begin{env}[21.10.10]
\label{IV.21.10.10}
Nous allons voir que sous certaines conditions, on peut élargir la définition de $f^*(Z)$ donnée dans \sref{IV.21.10.3} pour un cycle $1$-codimensionnel $Z$ sur $X$, mais à condition de prendre pour $f^*(Z)$ un cycle $1$-codimensionnel à coefficients rationnels sur $X'$.
Le cas plus général où nous nous plaçons est celui où $f$ transforme tout point maximal de $X'$ en un point maximal de $X$, et où, en tout point $x'\in X'^{(1)}$, on a l'une des conditions (i), (ii), (iii) de \sref{IV.21.10.1} \emph{ou une quatrième condition} (en posant $x=f(x')$):
\begin{enumerate}
  \item[(iv)] $x\in X^{(1)}$, $\mathfrak{m}_x\notin\operatorname{Ass}(\sh{O}_{X,x})$, et \emph{en outre, si l'on pose} $A=\widehat{\sh{O}}_{X,x}$, $A'=\widehat{\sh{O}}_{X',x'}$, \emph{et si $K$ désigne l'anneau total des fractions de $A$, alors $A'$ est une $A$-algèbre finie et $K'=A'\otimes_AK$ est un $K$-module} libre.
\end{enumerate}

Notons alors $r_{x'}$ le rang du $K$-module libre $K'$, $q_{x'}$ le degré de $\boldsymbol{k}(x')=\boldsymbol{k}(A')$ sur $\boldsymbol{k}(x)=\boldsymbol{k}(A)$, et posons $\mu_{x'}=r_{x'}/q_{x'}$.
Pour un cycle $1$-codimensionnel de support contenu dans $T$, donné par \sref{IV.21.10.1.1} et tout $x'\in X'^{(1)}$, on définit $c_{x'}\in\bb{Q}$, comme égal au nombre $n_{x'}\in\bb{Z}$ lorsque l'on est dans l'un des cas $1\textsuperscript{o}$, $2\textsuperscript{o}$, $3\textsuperscript{o}$ de \sref{IV.21.10.1}; mais il reste ici une quatrième possibilité:
\begin{enumerate}
  \item[$4\textsuperscript{o}$] si $x'$ vérifie la condition (iv) ci-dessus, on prend $c_{x'}=\mu_{x'}n_x\in\bb{Q}$.
\end{enumerate}
\end{env}

\begin{env}[21.10.11]
\EGAReaderTarget{ega4.c04.statement.21-10-11}{IV.21.10.11}
Il faut encore prouver que lorsque la condition $4\textsuperscript{o}$ de \sref{IV.21.10.10} est vérifiée \emph{en même temps} qu'une des conditions $1\textsuperscript{o}$, $2\textsuperscript{o}$, $3\textsuperscript{o}$ de \sref{IV.21.10.1}, on a $n_{x'}=c_{x'}$.
C'est évident si $x\notin T$ puisque alors $n_x=0$.
Pour étudier les deux autres cas, notons que $\mathfrak{m}_x\sh{O}_{X',x'}$ est fermé pour la topologie $\mathfrak{m}_{x'}$-préadique, donc le complété de $\sh{O}_{X',x'}/\mathfrak{m}_x\sh{O}_{X',x'}$ pour cette dernière topologie est $A'/\mathfrak{m}_xA'$.
Si l'on est à la fois dans le cas $2\textsuperscript{o}$ et le cas $4\textsuperscript{o}$, $\sh{O}_{X',x'}/\mathfrak{m}_x\sh{O}_{X',x'}$ est discret pour la topologie $\mathfrak{m}_{x'}$-préadique, donc isomorphe à $A'/\mathfrak{m}_xA'$.
Comme $A'$ est une $A$-algèbre finie et \emph{plate} $(\mathbf{0}_{\mathrm{III}},10.2.3)$, c'est un $A$-module \emph{libre} $(\mathbf{0}_{\mathrm{III}},10.1.3)$, et le rang de $A'/\mathfrak{m}_xA'$ sur $A/\mathfrak{m}_xA=\boldsymbol{k}(x)$ est égal au rang de $A'$ sur $A$, donc aussi à celui de $K'$ sur $K$.
D'autre part ce rang est aussi égal au produit de la longueur $\lambda_{x'}$ de $\sh{O}_{X',x'}/\mathfrak{m}_x\sh{O}_{X',x'}$ (en tant que $\sh{O}_{X',x'}$-module, ou de $\boldsymbol{k}(x')$-module) par le rang $[\boldsymbol{k}(x'): \boldsymbol{k}(x)]=q_{x'}$, ce qui prouve la relation $\mu_{x'}=r_{x'}/q_{x'}=\lambda_{x'}$ dans ce cas.

Supposons enfin que l'on soit à la fois dans le cas $3\textsuperscript{o}$ et le cas $4\textsuperscript{o}$.
Alors, puisque $\dim(\sh{O}_{X,x})=1$, $\sh{O}_{X,x}$ est un anneau de valuation discrète, donc régulier.
D'autre part, $\sh{O}_{X',x'}$ est de dimension $1$, et puisque $\mathfrak{m}_{x'}\notin\operatorname{Ass}(\sh{O}_{X',x'})$, $\sh{O}_{X',x'}$ est un anneau de Cohen--Macaulay $(\mathbf{0},16.4.6)$;
enfin, puisque $A'$ est un $A$-module de type fini, $A'/\mathfrak{m}_xA'$ est un $\boldsymbol{k}(x)$-espace vectoriel de rang fini;
puisque $\sh{O}_{X',x'}/\mathfrak{m}_x\sh{O}_{X',x'}$ est contenu dans $A'/\mathfrak{m}_xA'$, c'est aussi un $\boldsymbol{k}(x)$-espace vectoriel de rang fini, donc un anneau artinien.
L'application de \sref{IV.6.1.5} montre alors que $\sh{O}_{X',x'}$ est un $\sh{O}_{X,x}$-module \emph{plat}, donc on est aussi dans le cas $2\textsuperscript{o}$, et on conclut par ce qui a été vu plus haut.

Cela étant, dans le cas que l'on a considéré, on posera
\[
  f^*(Z)=\sum_{x'\in X'^{(1)}}c_{x'}\cdot\overline{\{x'\}}.
\]
\oldpage[IV-4]{294}
qui est donc un cycle $1$-codimensionnel sur $X'$ à coefficients \emph{rationnels}.
On a encore défini ainsi un homomorphisme $f^*$ de groupes ordonnés, vérifiant \sref{IV.21.10.3.1}, et par suite un homomorphisme de faisceaux de groupes commutatifs
\[
  \psi^*(\Gamma_T(\mathscr{Z}^1_X))\to\mathscr{Z}^1_{X'}\otimes_{\bb{Z}}\bb{Q}.
\]

Lorsque $f$ vérifie les conditions précédentes pour \emph{toute partie fermée $T$ rare} dans $X$, c'est-à-dire que pour tout $x'\in X'^{(1)}$, ou bien $x=f(x')$ est point maximal de $X$, ou bien $x'$ vérifie l'une des conditions (ii), (iii) ou (iv), $f^*(Z)$ est alors défini pour tout cycle $1$-codimensionnel $Z$ sur $X$, et on a défini un homomorphisme de faisceaux de groupes commutatifs ordonnés
\[
  \psi^*(\mathscr{Z}^1_X)\to\mathscr{Z}^1_{X'}\otimes_{\bb{Z}}\bb{Q},
\]
d'où, par tensorisation, un homomorphisme de faisceaux de $\bb{Q}$-espaces vectoriels ordonnés
\[
\label{IV.21.10.11.1}
  \psi^*(\mathscr{Z}^1_X\otimes_{\bb{Z}}\bb{Q})\to\mathscr{Z}^1_{X'}\otimes_{\bb{Z}}\bb{Q}.
  \tag{21.10.11.1}
\]
\end{env}

\begin{remark}[21.10.12]
\label{IV.21.10.12}
Lorsqu'on est dans la situation de \sref{IV.21.10.10}, on peut effectivement, pour des cycles $1$-codimensionnels $Z$ sur $X$, avoir pour $f^*(Z)$ des cycles $1$-codimensionnels à coefficients \emph{non entiers}; autrement dit les nombres $\mu_{x'}$ peuvent être non entiers.
On en a un exemple en prenant l'anneau intègre complet $A$ de $(\mathbf{0}_{\mathrm{IV}},6.15.11)$, (ii) et sa clôture intégrale $A'$: le point fermé $x'$ de $\Spec(A')$ vérifie la condition (iv) de \sref{IV.21.10.10} et l'on a $\mu_{x'}=\frac12$.
\end{remark}

\begin{lemma}[21.10.13]
\label{IV.21.10.13}
Soient $A$ un anneau local noethérien de dimension $1$, $t$ un élément régulier de $A$ appartenant à l'idéal maximal $\mathfrak{m}$ (ce qui implique que $\mathfrak{m}\notin\operatorname{Ass}(A)$).
\begin{enumerate}[label=(\roman*)]
  \item Pour tout $A$-module $M$ de type fini, le noyau $N_t(M)$ et le conoyau $P_t(M)$ de l'homothétie $t_M:M\to M$ de rapport $t$ sont de longueurs finies. On pose $d_t(M)=\operatorname{long}P_t(M)-\operatorname{long}N_t(M)$.
  \item Si $0\to M'\to M\to M''\to0$ est une suite exacte de $A$-modules de type fini, on a $d_t(M)=d_t(M')+d_t(M'')$.
  \item On a $d_t(M)\geq0$ pour tout $A$-module $M$ de type fini; pour que $d_t(M)=0$, il faut et il suffit que $M$ soit de longueur finie.
  \item Si $K$ est l'anneau total des fractions de $A$ et si $M\otimes_AK$ est un $K$-module libre de rang $n$, on a $d_t(M)=n\cdot d_t(A)=n\cdot\operatorname{long}(A/tA)$.
  \item Si $M$ vérifie les hypothèses de (iv) et si de plus $t$ est $M$-régulier, on a $\operatorname{long}(M/tM)=n\cdot\operatorname{long}(A/tA)$.
\end{enumerate}
\end{lemma}

\begin{proof}
(i) $\Spec(A)$ est formé du point $\mathfrak{m}$ et des idéaux premiers minimaux $\mathfrak{p}_i$;
comme par hypothèse $t\notin\mathfrak{p}_i$ pour tout $i$ (Bourbaki, \emph{Alg. comm.}, chap.~IV, \textsection~1, n\textsuperscript{o}~1, cor.~3 de la prop.~2), l'image de $t$ dans chacun des $A_{\mathfrak{p}_i}$ est inversible, et les supports des $A$-modules de type fini $N_t(M)$ et $P_t(M)$ sont donc vides ou réduits à $\mathfrak{m}$;
on en conclut $(\mathbf{0},16.1.10)$ que ces modules sont de longueur finie.

(ii) Puisque $t$ est régulier, on a une suite exacte
\[
  0\to A\xrightarrow{t_A}A\to A/tA\to0
\]
\oldpage[IV-4]{295}
et puisque $\operatorname{Tor}_i^A(M,A)=0$ pour $i\geq1$, la suite exacte des Tor donne d'une part la suite exacte
\[
  0\to\operatorname{Tor}_1^A(M,A/tA)\to M\xrightarrow{t_M}M\to M/tM\to0
\]
et d'autre part, pour $i\geq2$,
\[
  0=\operatorname{Tor}_i^A(M,A)\to\operatorname{Tor}_i^A(M,A/tA)\to\operatorname{Tor}_{i-1}^A(M,A)=0,
\]
donc on a $N_t(M)=\operatorname{Tor}_1^A(M,A/tA)$ et $\operatorname{Tor}_i^A(M,A/tA)=0$ pour $i\geq2$;
la suite exacte des Tor donne une suite exacte
\[
  0\to\operatorname{Tor}_1^A(M',A/tA)\to\operatorname{Tor}_1^A(M,A/tA)\to\operatorname{Tor}_1^A(M'',A/tA)\to
\]
\[
  M'\otimes_A(A/tA)\to M\otimes_A(A/tA)\to M''\otimes_A(A/tA)\to0
\]
et cette suite s'écrit, d'après ce qui précède,
\[
  0\to N_t(M')\to N_t(M)\to N_t(M'')\to P_t(M')\to P_t(M)\to P_t(M'')\to0,
\]
ce qui prouve (ii).

Pour prouver (iii), notons qu'il y a une suite de composition $(M_h)$ de $M$ dont les quotients $M_h/M_{h+1}$ sont isomorphes à $A/\mathfrak{m}$ ou à un des $A/\mathfrak{p}_i$ (Bourbaki, \emph{Alg. comm.}, chap.~IV, \textsection~1, n\textsuperscript{o}~4, th.~1), et pour que $M$ soit de longueur finie, il faut et il suffit que tous ces quotients soient isomorphes à $A/\mathfrak{m}$.
Tout revient donc (en vertu de (ii)) à prouver que $d_t(A/\mathfrak{m})=0$ et $d_t(A/\mathfrak{p}_i)>0$.
Or, l'image de $t$ dans $A/\mathfrak{m}$ étant $0$, on a $N_t(A/\mathfrak{m})=A/\mathfrak{m}$ et $P_t(A/\mathfrak{m})=A/\mathfrak{m}$, d'où la première assertion;
d'autre part, l'image de $t$ dans $A/\mathfrak{p}_i$ est régulière, donc $N_t(A/\mathfrak{p}_i)=0$ et $P_t(A/\mathfrak{p}_i)=A/(tA+\mathfrak{p}_i)$ qui n'est pas réduit à $0$, d'où la seconde assertion.

(iv) Il y a une base de $M\otimes_AK$ de la forme $(x_j/s)_{1\leq j\leq n}$, où $s$ est un élément régulier de $A$ et $x_j\in M$.
Considérons l'homomorphisme $u:A^n\to M$ qui transforme l'élément $e_j$ $(1\leq j\leq n)$ de la base canonique de $A^n$ en $x_j$ et montrons que $\operatorname{Ker}(u)$ et $\operatorname{Coker}(u)$ sont de \emph{longueur finie};
en effet, pour tout $i$, l'image de $s$ dans $A_{\mathfrak{p}_i}$ est inversible, et comme $K_{\mathfrak{p}_i}=A_{\mathfrak{p}_i}$, les images des $x_j/s$ dans $M_{\mathfrak{p}_i}$ forment une base de ce $A_{\mathfrak{p}_i}$-module;
donc $u_{\mathfrak{p}_i}:A_{\mathfrak{p}_i}^n\to M_{\mathfrak{p}_i}$ est bijectif.
Comme dans (i), on en conclut que les supports de $\operatorname{Ker}(u)$ et de $\operatorname{Coker}(u)$ sont contenus dans $\{\mathfrak{m}\}$, donc que ces modules sont de longueur finie.
Cela étant, il résulte de (ii) et (iii) que l'on a $d_t(M)=d_t(A^n)=n\cdot d_t(A)=n\cdot\operatorname{long}(A/tA)$ puisque $t$ est régulier.

Enfin, il est clair que (v) se déduit aussitôt de (iv), puisque alors $N_t(M)=0$.
\end{proof}

Ce lemme permet de généraliser \sref{IV.21.10.4}:

\begin{proposition}[21.10.13]
\label{IV.21.10.13bis}
Supposons que $f$ vérifie les conditions de \sref{IV.21.10.10}.
Alors, pour tout diviseur $D$ sur $X$ tel que $\Supp(D)\subset T$ et que $f^*(D)$ soit défini, on a
\[
\label{IV.21.10.13.1}
  \operatorname{cyc}(f^*(D))=f^*(\operatorname{cyc}(D)).
  \tag{21.10.13.1}
\]
\end{proposition}

\begin{proof}
Raisonnant comme dans \sref{IV.21.10.4}, tout revient à voir (avec les mêmes notations) que si $x'$ est dans le cas $4\textsuperscript{o}$ de \sref{IV.21.10.10} et si $n_y$ est la longueur de $\sh{O}_{X,y}/t_y\sh{O}_{X,y}$, alors la
\oldpage[IV-4]{296}
longueur $n_{x'}$ de $\sh{O}_{X',x'}/t'_{x'}\sh{O}_{X',x'}$ est égale à $\mu_{x'}n_y$, $\mu_{x'}$ étant le nombre rationnel défini dans \sref{IV.21.10.10}.
Comme $\sh{O}_{X',x'}/t'_{x'}\sh{O}_{X',x'}$ est de longueur finie, il a même longueur que son complété $\mathfrak{m}_{x'}$-préadique $A'/t'_{x'}A'$, qu'on peut aussi écrire $A'/t_yA'$;
d'ailleurs, comme $t'_{x'}$ est régulier par hypothèse dans $\sh{O}_{X',x'}$, il l'est aussi dans $A'$ par platitude $(\mathbf{0}_{\mathrm{I}},6.3.4)$, et lorsque $A'$ est considéré comme $A$-module, on peut aussi dire que $t_y$ est $A'$-régulier.
Comme $A$ est de dimension $1$ et que $A'$ est un $A$-module de type fini tel que $A'\otimes_AK$ soit un $K$-module libre de rang $r_{x'}$, on peut appliquer \sref{IV.21.10.13}, (v) à $M=A'$ et à $t_y$, et la longueur de $A'/t_yA'$ en tant que $A$-module est donc $r_{x'}n_y$.
Comme $\boldsymbol{k}(x')$ est un $\boldsymbol{k}(x)$-espace vectoriel de rang $q_{x'}$, la longueur de $A'/t_yA'$ en tant que $A'$-module est donc $r_{x'}n_y/q_{x'}=\mu_{x'}n_y$, ce qui achève la démonstration.
\end{proof}

\begin{env}[21.10.14]
\EGAReaderTarget{ega4.c04.statement.21-10-14}{IV.21.10.14}
Soient maintenant $X$, $X'$ deux préschémas localement noethériens, $f:X'\to X$ un morphisme ayant les \emph{deux propriétés} suivantes:
\begin{enumerate}[label=\emph{\alph*)}]
  \item $f$ est fini;
  \item l'image par $f$ de tout point maximal de $X'$ est un point maximal de $X$.
\end{enumerate}
Pour tout $x\in X^{(1)}$, les points $x'\in f^{-1}(x)$ appartiennent alors tous à $X'^{(1)}$, comme il résulte de l'hypothèse \emph{b)} et de l'inégalité $(\mathbf{0},16.3.9.1)$, la fibre $f^{-1}(x)$ étant discrète.
Soit alors
\[
  Z'=\sum_{x'\in X'^{(1)}}n_{x'}\cdot\overline{\{x'\}}
\]
un cycle $1$-codimensionnel sur $X'$.
Pour tout $x\in X^{(1)}$, définissons un entier $n_x$ par la formule
\[
  n_x=\sum_{x'\in f^{-1}(x)}n_{x'}\cdot[\boldsymbol{k}(x'): \boldsymbol{k}(x)],
\]
ce qui a un sens, les points de $f^{-1}(x)$ étant en nombre fini et $\boldsymbol{k}(x')$ étant un corps de degré fini sur $\boldsymbol{k}(x)$ \sref[I]{I.6.4.4}.
En outre, l'ensemble des $x\in X^{(1)}$ tels que $n_x\neq0$ est \emph{localement fini} dans $X$, car il est contenu dans l'image par $f$ de l'ensemble des $x'\in X'^{(1)}$ tels que $n_{x'}\neq0$, et la conclusion résulte de ce que le morphisme $f$ est quasi-compact.
On peut donc définir un cycle $1$-codimensionnel sur $X$ en posant
\[
\label{IV.21.10.14.1}
  f_*(Z')=\sum_{x\in X^{(1)}}n_x\cdot\overline{\{x\}}
  \tag{21.10.14.1}
\]
et on dit que $f_*(Z')$ est l'\emph{image directe} de $Z'$ par $f$.
Il est clair que l'application $f_*:\mathfrak{Z}^1(X')\to\mathfrak{Z}^1(X)$ ainsi définie est un homomorphisme de groupes ordonnés.
En outre, si $U$ est un ouvert de $X$, $f_U:f^{-1}(U)\to U$ la restriction de $f$, il résulte aussitôt des définitions que l'on a
\[
\label{IV.21.10.14.2}
  (f_U)_*(Z'|f^{-1}(U))=f_*(Z')|U,
  \tag{21.10.14.2}
\]
donc, en désignant par $\psi$ l'application continue sous-jacente au morphisme $f$, les applications $\Gamma(f^{-1}(U),\mathscr{Z}^1_{X'})\to\Gamma(U,\mathscr{Z}^1_X)$ que l'on vient de définir définissent un homomorphisme de faisceaux de groupes commutatifs ordonnés sur $X$
\[
  \psi_*(\mathscr{Z}^1_{X'})\to\mathscr{Z}^1_X.
\]
\end{env}

\begin{proposition}[21.10.15]
\label{IV.21.10.15}
\oldpage[IV-4]{297}
Soient $X$, $X'$, $X''$ trois préschémas localement noethériens, $f:X'\to X$, $g:X''\to X'$ deux morphismes vérifiant les conditions \emph{a)} et \emph{b)} de \sref{IV.21.10.14}.
Alors $f\circ g$ vérifie les mêmes conditions, et pour tout cycle $1$-codimensionnel $Z''$ sur $X''$, on a
\[
  (f\circ g)_*(Z'')=f_*(g_*(Z'')).
\]
\end{proposition}

\begin{proof}
Cela résulte aussitôt des définitions.
\end{proof}

\begin{proposition}[21.10.16]
\label{IV.21.10.16}
Soient $X$, $X'$, $X_1$ trois préschémas localement noethériens, $f:X'\to X$ un morphisme vérifiant les conditions \emph{a)} et \emph{b)} de \sref{IV.21.10.14}, $g:X_1\to X$ un morphisme plat.
On pose $X'_1=X'\times_XX_1$ (de sorte que $X'_1$ est localement noethérien) et on note $f_1:X'_1\to X_1$ et $g_1:X'_1\to X'$ les projections canoniques.
Alors $f_1$ vérifie les conditions \emph{a)} et \emph{b)} de \sref{IV.21.10.14}, et pour tout $1$-cycle codimensionnel $Z'$ sur $X'$, on a
\[
\label{IV.21.10.16.1}
  g^*(f_*(Z'))=(f_1)_*(g_1^*(Z')).
  \tag{21.10.16.1}
\]
\end{proposition}

\begin{proof}
Il est clair que $f_1$ est fini, et il vérifie la condition \emph{b)} de \sref{IV.21.10.14} en vertu de \sref{IV.2.3.7}.
Pour prouver \sref{IV.21.10.16.1}, on est aussitôt ramené au cas où $X$, $X'$ et $X_1$ sont des spectres d'anneaux locaux noethériens de dimension $1$, $A$, $A'$ et $A_1$, $A'$ étant un $A$-module fini et $A_1$ un $A$-module plat.
Désignant par $x$, $x'$ et $x_1$ les points fermés de $X$, $X'$ et $X_1$ respectivement, il s'agit de voir que l'on a
\[
\label{IV.21.10.16.2}
  \sum_{x'_1}\lambda_{x'_1}[\boldsymbol{k}(x'_1):\boldsymbol{k}(x_1)]
  =\lambda_{x_1}[\boldsymbol{k}(x'):\boldsymbol{k}(x)]
  \tag{21.10.16.2}
\]
où $x'_1$ parcourt l'ensemble des points fermés de $X'_1$ (i.e. l'ensemble des points à la fois au-dessus de $x'$ et de $x_1$) et où $\lambda_{x'_1}=\operatorname{long}(\sh{O}_{X'_1,x'_1}/\mathfrak{m}_{x'}\sh{O}_{X'_1,x'_1})$ et $\lambda_{x_1}=\operatorname{long}(A_1/\mathfrak{m}_xA_1)$.
Comme on a
\[
  [\boldsymbol{k}(x'_1):\boldsymbol{k}(x')]\cdot[\boldsymbol{k}(x'):\boldsymbol{k}(x)]
  =[\boldsymbol{k}(x'_1):\boldsymbol{k}(x_1)]\cdot[\boldsymbol{k}(x_1):\boldsymbol{k}(x)],
\]
le premier membre de \sref{IV.21.10.16.2} s'écrit aussi
\[
  \frac{[\boldsymbol{k}(x'): \boldsymbol{k}(x)]}{[\boldsymbol{k}(x_1): \boldsymbol{k}(x)]}
  \operatorname{long}_{A'}(A'_1/\mathfrak{m}_{x'}A'_1)
  =\operatorname{long}_{A_1}(A'_1/\mathfrak{m}_{x'}A'_1),
\]
où on a posé $A'_1=A'\otimes_AA_1$.
On a donc $A'_1/\mathfrak{m}_{x'}A'_1=(A'/\mathfrak{m}_{x'}A')\otimes_AA_1$, et comme $A_1$ est un $A$-module plat, on a par \sref{IV.4.7.1}
\[
  \operatorname{long}_{A_1}(A'_1/\mathfrak{m}_{x'}A'_1)
  =\operatorname{long}_{A}(A'/\mathfrak{m}_{x'}A')\operatorname{long}_{A_1}(A_1/\mathfrak{m}_xA_1)
  =[\boldsymbol{k}(x'): \boldsymbol{k}(x)]\cdot\lambda_{x_1},
\]
ce qui achève la démonstration.
\end{proof}

\begin{proposition}[21.10.17]
\label{IV.21.10.17}
Soient $X$, $X'$ deux préschémas localement noethériens, $f:X'\to X$ un morphisme fini localement libre.
Alors $f$ vérifie la condition \emph{b)} de \sref{IV.21.10.14}, et pour tout diviseur $D'$ sur $X'$, on a
\[
\label{IV.21.10.17.1}
  f_*(\operatorname{cyc}(D'))=\operatorname{cyc}(f_*(D')).
  \tag{21.10.17.1}
\]
\end{proposition}

\begin{proof}
Comme $f$ est plat et fini, la condition \emph{b)} de \sref{IV.21.10.14} et la relation $f(X'^{(1)})\subset X^{(1)}$ résultent de \sref{IV.6.1.1}.
La définition \sref{IV.21.10.14.1} montre qu'on peut se ramener au cas où $X=\Spec(A)=\Spec(\sh{O}_{X,x})$ pour un $x\in X^{(1)}$;
alors $X'=\Spec(A')$, où $A'$ est une
\oldpage[IV-4]{298}
$A$-algèbre qui est un $A$-module libre de rang fini;
en outre, on peut supposer que $D'=\operatorname{div}(t')$, où $t'$ est un élément régulier de $A'$;
on a alors $f_*(D')=\operatorname{div}(t)$, où $t=\mathrm{N}_{A'/A}(t')$ est un élément régulier de $A$ \sref{IV.21.5.2}.
On peut se borner au cas où $t'$ n'est pas inversible dans $A'$, ce qui équivaut au fait que $t$ ne l'est pas dans $A$;
l'anneau $A/tA$ est alors de longueur finie et $\neq0$, et $A'/tA'$ est composé direct d'anneaux locaux artiniens $A'_i$ $(1\leq i\leq r)$, le corps résiduel de $A'_i$ étant $\boldsymbol{k}(x'_i)$, où les $x'_i$ $(1\leq i\leq r)$ sont les points de $X'$ au-dessus de $x$.
Si $t'_i$ $(1\leq i\leq r)$ est l'image de $t'$ dans $A'_i$, $A'/t'A'$ est composé direct des $A'_i/t'_iA'_i$;
comme le produit $\operatorname{long}_{A'_i}(A'_i/t'_iA'_i)\cdot[\boldsymbol{k}(x'_i):\boldsymbol{k}(x)]$ est égal à $\operatorname{long}_A(A'_i/t'_iA'_i)$, on voit que la multiplicité au point $x$ du premier membre de \sref{IV.21.10.17.1} est $\operatorname{long}_A(A'/t'A')$, si bien que la formule à démontrer se réduit à
\[
\label{IV.21.10.17.2}
  \operatorname{long}_A(A'/t'A')=\operatorname{long}_A(A/tA).
  \tag{21.10.17.2}
\]

Cette relation découle du lemme plus général suivant.
\end{proof}

\begin{lemma}[21.10.17.3]
\label{IV.21.10.17.3}
Soient $A$ un anneau local noethérien de dimension $1$, $M$ un $A$-module libre de rang fini, $u$ un endomorphisme injectif de $M$.
Alors on a
\[
\label{IV.21.10.17.4}
  \operatorname{long}_A(\operatorname{Coker}u)=\operatorname{long}_A(A/(\det u)A).
  \tag{21.10.17.4}
\]
\end{lemma}

\begin{proof}
Distinguons plusieurs cas.

I) $A$ est un anneau de valuation discrète.
Soit en effet $\pi$ une uniformisante de $A$, et remarquons que $\operatorname{long}_A(A/\pi^kA)=k$;
si $n$ est le rang de $M$ et si $\pi^{m_i}$ $(1\leq i\leq n)$ sont les facteurs invariants de $u$, $\operatorname{Coker}(u)$ est somme directe des $A$-modules $A/\pi^{m_i}A$, donc a pour longueur $m=\sum_{i=1}^nm_i$ et $\det(u)$ est produit d'un élément inversible et de $\pi^m$, d'où la conclusion dans ce cas (Bourbaki, \emph{Alg.}, chap.~VII, \textsection~4, n\textsuperscript{o}~5, cor.~1 de la prop.~4).

II) $A$ est un anneau complet intègre (de dimension $1$).
On sait alors $(\mathbf{0},19.8.8)$, (ii) qu'il y a un sous-anneau $B$ de $A$, qui est un anneau de valuation discrète, tel que $B\to A$ soit un homomorphisme local faisant de $A$ un $B$-module de type fini;
comme ce $B$-module est évidemment sans torsion, il est libre (Bourbaki, \emph{Alg. comm.}, chap.~VI, \textsection~3, n\textsuperscript{o}~6, lemme~1).
Notons $M'$ l'ensemble $M$ muni de sa structure de $B$-module (libre), par $u'$ l'endomorphisme $u$ considéré comme $B$-endomorphisme.
Il résulte de I) que l'on a
\[
\label{IV.21.10.17.5}
  \operatorname{long}_B(\operatorname{Coker}u')=\operatorname{long}_B(B/(\det u')B).
  \tag{21.10.17.5}
\]
Mais on a (Bourbaki, \emph{Alg.}, chap.~VIII, \textsection~12, n\textsuperscript{o}~2, prop.~7)
\[
  \det u'=\mathrm{N}_{A/B}(\det u),
\]
donc, en appliquant \sref{IV.21.10.17.4} à l'homothétie $x\mapsto(\det u)x$ du $B$-module libre $A$, il vient
\[
  \operatorname{long}_B(B/(\det u')B)=\operatorname{long}_B(A/(\det u)A),
\]
d'où, en portant dans \sref{IV.21.10.17.5} et divisant par $[\boldsymbol{k}(A):\boldsymbol{k}(B)]$, longueur du $B$-module $\boldsymbol{k}(A)$, on obtient \sref{IV.21.10.17.4} dans le cas envisagé.

\oldpage[IV-4]{299}
III) $A$ est un anneau complet.
Notons qu'on peut supposer en outre que $\mathfrak{m}\notin\operatorname{Ass}(A)$;
en effet, puisque $\det(u)$ est un élément régulier de $A$, si l'on avait $\mathfrak{m}\in\operatorname{Ass}(A)$, on en déduirait $\det(u)\notin\mathfrak{m}$, donc $\det(u)$ serait inversible, $u$ un automorphisme de $M$, et la formule \sref{IV.21.10.17.4} devient alors triviale, les deux membres étant nuls.

Dans ce qui suit, pour un endomorphisme $v$ d'un module $N$ sur un anneau $R$, tel que $\operatorname{Ker}v$ et $\operatorname{Coker}v$ soient de longueur finie, on posera
\[
  \chi(N,v)=\operatorname{long}_R(\operatorname{Ker}(v))-\operatorname{long}_R(\operatorname{Coker}(v)).
\]
On notera que l'hypothèse sur $v$ revient à dire que le complexe
\[
  K^0:0\to N\xrightarrow{v}N\to0
\]
a ses modules de cohomologie de longueur finie et que $\chi(N,v)=\chi(H^0(K^0))$ $(\mathbf{0}_{\mathrm{III}},11.10)$.
On en déduit que si $N'$, $N$, $N''$ sont trois $R$-modules, si l'on a un diagramme commutatif
\[
  \xymatrix{
    0 \ar[r] & N' \ar[r]^r \ar[d]^{v'} & N \ar[r]^s \ar[d]^v & N'' \ar[r] \ar[d]^{v''} & 0\\
    0 \ar[r] & N' \ar[r]_r & N \ar[r]_s & N'' \ar[r] & 0
  }
\]
dont les lignes sont exactes, et si deux des trois nombres $\chi(N,v)$, $\chi(N',v')$ et $\chi(N'',v'')$ sont définis, alors il en est de même du troisième et on a
\[
\label{IV.21.10.17.6}
  \chi(N,v)=\chi(N',v')+\chi(N'',v'').
  \tag{21.10.17.6}
\]
Cela résulte en effet aussitôt de la suite exacte de cohomologie.

Enfin, si $N$ est un $R$-module de longueur finie, on a $\chi(N,v)=0$.

Avec ces notations, on a le lemme suivant.
\end{proof}

\begin{lemma}[21.10.17.7]
\label{IV.21.10.17.7}
Soient $A$ un anneau local noethérien de dimension $1$ dont l'idéal maximal $\mathfrak{m}$ est tel que $\mathfrak{m}\notin\operatorname{Ass}(A)$;
soient $\mathfrak{p}_i$ $(1\leq i\leq n)$ les idéaux premiers minimaux de $A$.
Soient $M$ un $A$-module libre de type fini, $u$ un endomorphisme de $M$ tel que $\chi(M,u)$ soit défini;
pour chaque $i$, posons $M_i=M\otimes_A(A/\mathfrak{p}_i)$ et soit $u_i$ l'endomorphisme $u\otimes1_{A/\mathfrak{p}_i}$ de $M_i$;
alors si $\chi(M_i,u_i)$ est défini pour chaque $i$, on a
\[
\label{IV.21.10.17.8}
  \chi(M,u)=\sum_{i=1}^n\operatorname{long}(A_{\mathfrak{p}_i})\cdot\chi(M_i,u_i).
  \tag{21.10.17.8}
\]
\end{lemma}

\begin{proof}
Puisque $\mathfrak{m}\notin\operatorname{Ass}(A)$, on a une décomposition primaire réduite unique $(0)=\bigcap_{1\leq i\leq n}\mathfrak{q}_i$, où l'idéal $\mathfrak{q}_i$ est $\mathfrak{p}_i$-primaire pour $1\leq i\leq n$.
Si l'on pose $M'_i=M\otimes_A(A/\mathfrak{q}_i)$, on a donc une suite exacte
\[
  0\to M\to\bigoplus_iM'_i\to M''\to0
\]
de $A$-modules, où $M''$ est de longueur finie:
en effet, en localisant la suite exacte précédente en chacun des $\mathfrak{p}_i$, on obtient $M''_{\mathfrak{p}_i}=0$, car $(\mathfrak{q}_i)_{\mathfrak{p}_i}=0$ et $(\mathfrak{q}_j)_{\mathfrak{p}_i}=A_{\mathfrak{p}_i}$ pour $j\neq i$ (Bourbaki, \emph{Alg. comm.}, chap.~IV, \textsection~2, n\textsuperscript{o}~4, prop.~6), donc $(M'_i)_{\mathfrak{p}_i}=M_{\mathfrak{p}_i}$ et $(M'_j)_{\mathfrak{p}_i}=0$;
\oldpage[IV-4]{300}
le support de $M''$ étant donc réduit à $\mathfrak{m}$, $M''$ est de longueur finie $(\mathbf{0},16.1.10)$.
Si l'on pose $u'_i=u\otimes1_{A/\mathfrak{q}_i}$, et si $u''$ est l'endomorphisme de $M''$ déduit de $\bigoplus u'_i$ par passage aux quotients, on déduit de \sref{IV.21.10.17.6} que $\chi(M,u)+\chi(M'',u'')=\sum_i\chi(M'_i,u'_i)$, et puisque $M''$ est de longueur finie, $\chi(M'',u'')=0$.
Pour prouver \sref{IV.21.10.17.8}, on peut donc se ramener au cas où $n=1$.
On désignera alors par $\mathfrak{p}$ l'unique idéal premier minimal, qui est le nilradical de $A$;
si $A_0=A/\mathfrak{p}$, $M_0=M\otimes_AA_0$, $u_0=u\otimes1_{A_0}$, il s'agit de prouver que si $\chi(M_0,u_0)$ est défini, on a
\[
\label{IV.21.10.17.9}
  \chi(M,u)=\operatorname{long}A_{\mathfrak{p}}\cdot\chi(M_0,u_0).
  \tag{21.10.17.9}
\]

Soit $\mathfrak{n}_j$ $(0\leq j\leq r)$ la « puissance symbolique $j$-ème » de $\mathfrak{p}$, image réciproque dans $A$ de l'idéal $(\mathfrak{p}A_{\mathfrak{p}})^j$ de $A_{\mathfrak{p}}$ $(1\leq j\leq r)$, avec $\mathfrak{n}_0=A$ et $\mathfrak{n}_r=0$;
posons
\[
  M_j=\mathfrak{n}_jM/\mathfrak{n}_{j+1}M\qquad(0\leq j\leq r-1),
\]
et notons $v_j$ l'endomorphisme de $M_j$ déduit de $u$ par restriction à $\mathfrak{n}_jM$ et passage aux quotients.
Nous allons d'abord montrer que chacun des nombres $\chi(M_j,v_j)$ est défini et que l'on a
\[
\label{IV.21.10.17.10}
  \chi(M,u)=\sum_j\chi(M_j,v_j).
  \tag{21.10.17.10}
\]
La première assertion entraînera la seconde, en appliquant \sref{IV.21.10.17.6} à chacune des suites exactes
\[
  0\to\mathfrak{n}_jM/\mathfrak{n}_{j+1}M\to M/\mathfrak{n}_{j+1}M\to M/\mathfrak{n}_jM\to0.
\]

Pour prouver la première assertion, on note que si $m$ est le rang du $A$-module libre $M$, $M_j$ est isomorphe à $(\mathfrak{n}_j/\mathfrak{n}_{j+1})^m$, ou encore (puisque $\mathfrak{p}$ annule chacun des quotients $\mathfrak{n}_j/\mathfrak{n}_{j+1}$), $M_j$ est un $A_0$-module isomorphe à $M_0\otimes_{A_0}(\mathfrak{n}_j/\mathfrak{n}_{j+1})$.
Désignons par $l_j$ le rang du $A_0$-module $\mathfrak{n}_j/\mathfrak{n}_{j+1}$;
comme le corps des fractions $K_0$ de $A_0$ est le corps résiduel de $A_{\mathfrak{p}}$, $l_j$ est aussi la longueur du $A_{\mathfrak{p}}$-module $(\mathfrak{p}A_{\mathfrak{p}})^j/(\mathfrak{p}A_{\mathfrak{p}})^{j+1}$.
Il y a un système de générateurs du $A_0$-module $M_j$ qui contient une base de $M_j\otimes_{A_0}K_0$;
donc il y a un $A_0$-homomorphisme
\[
  w_j:M_0^{l_j}\to M_j
\]
dont le localisé en l'idéal $(0)$ de $A_0$ est un isomorphisme, de sorte que le support de $\operatorname{Ker}(w_j)$ et de $\operatorname{Coker}(w_j)$ est réduit à l'idéal maximal $\mathfrak{m}/\mathfrak{p}$ de $A_0$;
$\operatorname{Ker}(w_j)$ et $\operatorname{Coker}(w_j)$ sont donc des $A_0$-modules de longueur finie $(\mathbf{0},16.1.10)$.
Comme par hypothèse $\chi(M_0,u_0)$ est défini, il en est de même de $\chi(M_0^{l_j},u_0^{l_j})=l_j\chi(M_0,u_0)$ et en vertu de \sref{IV.21.10.17.6} et du fait que $\operatorname{Ker}(w_j)$ et $\operatorname{Coker}(w_j)$ sont de longueur finie, on voit que $\chi(M_j,v_j)$ est défini et égal à $l_j\chi(M_0,u_0)$;
la relation \sref{IV.21.10.17.10} donne alors
\[
  \chi(M,u)=\left(\sum_jl_j\right)\chi(M_0,u_0),
\]
et en vertu d'une remarque précédente, $\sum_jl_j$ n'est autre que la longueur de $A_{\mathfrak{p}}$, ce qui achève de prouver le lemme \sref{IV.21.10.17.7}.
\end{proof}

\oldpage[IV-4]{301}
Pour appliquer ce lemme lorsque $A$ est un anneau complet et $\mathfrak{m}\notin\operatorname{Ass}(A)$, on observe que si $u$ est injectif, il en est de même des $u_i$ (avec les notations du lemme):
en effet, $\det u$ est alors un élément régulier de $A$, donc n'appartient à aucun des $\mathfrak{p}_i$, et son image $\det u_i$ dans $A/\mathfrak{p}_i$ est par suite un élément $\neq0$ de cet anneau intègre, ce qui prouve que $u_i$ est injectif.
Comme $\operatorname{Coker}(u_i)$ est image de $\operatorname{Coker}(u)$, il est aussi de longueur finie et $\chi(M_i,u_i)$ est donc défini pour tout $i$;
on a par suite la formule \sref{IV.21.10.17.8}.
D'autre part, puisque $\det(u)$ est un élément régulier de $A$, il n'est contenu dans aucun des $\mathfrak{p}_i$;
l'idéal $(\det u)A$ est donc $\mathfrak{m}$-primaire et le quotient $A/(\det u)A$ de longueur finie.
Appliquant le même raisonnement que ci-dessus à l'homothétie injective $t:\xi\mapsto(\det u)\xi$ de $A$ et à ses images $t_i=\det u_i$ dans les $A/\mathfrak{p}_i$, il vient
\[
  \chi(A,t)=\sum_i\operatorname{long}(A_{\mathfrak{p}_i})\cdot\chi(A/\mathfrak{p}_i,t_i).
\]
Mais les anneaux $A/\mathfrak{p}_i$ sont intègres et complets, et en appliquant le résultat de II) à chacun d'eux, on obtient encore la relation \sref{IV.21.10.17.4} pour $M$ et $u$.

IV) Cas général.
Posons $A'=\widehat{A}$, $M'=M\otimes_AA'$, $u'=u\otimes1_{A'}$;
on a $\det(u')=\det(u)$, et par platitude, $\operatorname{Coker}(u')=(\operatorname{Coker}(u))_{(A')}$ et $A'/(\det u)A'=(A/(\det u)A)_{(A')}$;
comme la formule \sref{IV.21.10.17.4} est vraie pour $A'$ et $u'$ en vertu de III), elle est aussi vraie pour $A$ et $u$ en vertu de \sref{IV.4.7.1}.

Ceci termine donc la démonstration de \sref{IV.21.10.17}.

\begin{proposition}[21.10.18]
\label{IV.21.10.18}
Soient $X$, $X'$ deux préschémas localement noethériens, $f:X'\to X$ un morphisme fini, transformant tout point maximal de $X'$ en un point maximal de $X$, et vérifiant pour tout $x'\in X'$ l'une des conditions (ii), (iii), (iv) de \sref{IV.21.10.10}.
On suppose en outre qu'il existe un entier $n$ tel que, pour tout point maximal $x$ de $X$, $(f_*(\sh{O}_{X'}))_x$ soit un $\sh{O}_{X,x}$-module libre de rang $n$.
Alors, pour tout cycle $1$-codimensionnel $Z$ sur $X$, on a
\[
\label{IV.21.10.18.1}
  f_*(f^*(Z))=n\cdot Z
  \tag{21.10.18.1}
\]
(« formule de projection »).
\end{proposition}

\begin{proof}
En vertu des définitions, on est aussitôt ramené au cas où $X$ est spectre d'un anneau local noethérien $A$ de dimension $1$, de point fermé $x$, et où $Z=\overline{\{x\}}$;
posons $X'=\Spec(A')$ où $A'$ est une $A$-algèbre finie, et, pour tout idéal premier minimal $\mathfrak{p}_i$ de $A$, $A'_{\mathfrak{p}_i}$ est un $A_{\mathfrak{p}_i}$-module libre de rang $n$.
Montrons qu'on peut se borner en outre au cas où $A$ est complet.
Faisons en effet le changement de base $h:Y\to X$, où $Y=\Spec(B)$, avec $B=\widehat{A}$, et posons
\[
  Y'=X'\times_XY=\Spec(B\otimes_AA'),\qquad g=f_{\{Y\}}:Y'\to Y;
\]
le morphisme $g$ est alors fini, et comme $h$ est plat, les points maximaux de $Y$ sont au-dessus de ceux de $X$;
au-dessus de chacun des $\mathfrak{p}_i$, il n'y a qu'un nombre fini d'idéaux minimaux $\mathfrak{p}_{ij}$ de $B$, et $(B\otimes_AA')_{\mathfrak{p}_{ij}}$ est un $B_{\mathfrak{p}_{ij}}$-module libre de rang $n$.
Enfin, si $f$ vérifie l'une des conditions (ii), (iii) ou (iv) de \sref{IV.21.10.10} en chacun des points $x'$ de $f^{-1}(x)$, $g$ vérifie la condition correspondante en l'unique point $y'$ de $g^{-1}(y)$ au-dessus de $x'$ (en désignant par $y$ le point fermé de $Y$);
c'est immédiat pour les conditions (ii) et (iv);
pour la condition (iii), elle implique que $A$ est un anneau de valuation discrète, donc il en est de même de $B$, et la condition
\oldpage[IV-4]{302}
$\mathfrak{m}_{y'}\notin\operatorname{Ass}(\sh{O}_{Y',y'})$ résulte, par platitude, de la condition $\mathfrak{m}_{x'}\notin\operatorname{Ass}(\sh{O}_{X',x'})$ $(\mathbf{0},3.3.1)$.
Le morphisme $g$ vérifie donc les mêmes conditions que $f$;
si l'on prouve la formule \sref{IV.21.10.18.1} pour $g$ et $\overline{\{y\}}$, la même formule sera valable pour $f$ et $\overline{\{x\}}$, en vertu de \sref{IV.21.10.16.1}.

On peut donc supposer que $A$ est \emph{complet};
alors il en est de même de $A'$ qui est donc composé direct d'anneaux locaux complets;
on peut par suite se borner au cas où $A'$ est un anneau \emph{local}, et il reste à vérifier \sref{IV.21.10.18.1} dans chacun des cas (ii), (iii), (iv) pour le point fermé $x'$ de $X'$.
Dans le cas (ii), $A'$ étant un $A$-module plat de type fini, est un $A$-module \emph{libre} $(\mathbf{0}_{\mathrm{III}},10.1.3)$, de rang $n$ en vertu de l'hypothèse.
Or, on a par définition \sref{IV.21.10.1} et \sref{IV.21.10.3}
\[
  f^*(Z)=\lambda_{x'}\cdot\overline{\{x'\}},
\]
où $\lambda_{x'}$ est la longueur du $A'$-module $A'/\mathfrak{m}_xA'$, puis
\[
  f_*(f^*(Z))=(\lambda_{x'}[\boldsymbol{k}(x'): \boldsymbol{k}(x)])\cdot\overline{\{x\}};
\]
mais $\lambda_{x'}[\boldsymbol{k}(x'): \boldsymbol{k}(x)]$ est la longueur de $A'/\mathfrak{m}_xA'=A'\otimes_A\boldsymbol{k}(x)$ en tant que $A$-module, ou encore son rang en tant que $\boldsymbol{k}(x)$-espace vectoriel, donc est égal à $n$.

Dans le cas (iii), $A$ est un anneau de valuation discrète, donc régulier, et l'hypothèse $\mathfrak{m}_{x'}\notin\operatorname{Ass}(\sh{O}_{X',x'})$ entraîne que $A'$ est un anneau de Cohen--Macaulay $(\mathbf{0},16.4.6)$;
comme $\dim(A')=\dim(A)=1$ et que $A'/\mathfrak{m}_xA'$ est un anneau artinien, il résulte de \sref{IV.6.1.5} que $A'$ est un $A$-module plat et on est ramené au cas (ii).

Dans le cas (iv), si $K$ est l'anneau total des fractions de $A$, $K'=A'\otimes_AK$ est par hypothèse un $K$-module libre de rang $n$ et par définition \sref{IV.21.10.10}, on a
\[
  f^*(Z)=\frac{n}{[\boldsymbol{k}(x'): \boldsymbol{k}(x)]}\cdot\overline{\{x'\}},
\]
d'où $f_*(f^*(Z))=n\cdot\overline{\{x\}}$.
\end{proof}

On notera que la formule \sref{IV.21.10.18.1} est applicable en particulier lorsque le morphisme $f$ est \emph{fini} et \emph{plat} et tel que pour tout point maximal $x$ de $X$, $(f_*(\sh{O}_{X'}))_x$ soit un $\sh{O}_{X,x}$-module libre de rang $n$.

\begin{corollary}[21.10.19]
\label{IV.21.10.19}
Sous les hypothèses de \sref{IV.21.10.18}, soit $D$ un diviseur sur $X$ tel que $f^*(D)$ soit défini \sref{IV.21.4.5}; alors on a
\[
\label{IV.21.10.19.1}
  f_*(\operatorname{cyc}(f^*(D)))=n\cdot\operatorname{cyc}(D).
  \tag{21.10.19.1}
\]
\end{corollary}

\begin{proof}
Cela résulte de \sref{IV.21.10.18} et \sref{IV.21.10.13bis}.
\end{proof}

% FIN SOURCE ega4/ega4-21b-fr.tex

%
% DEBUT SOURCE ega4/ega4-21c-fr.tex
\subsection{Factorialité des anneaux locaux réguliers}
\EGAReaderTarget{ega4.c04.subsec.21-11}{subsection:IV.21.11}
\label{IV.21.11}

\begin{theorem}[21.11.1]
\label{IV.21.11.1}
\textup{(Auslander-Buchsbaum)} ---
Un anneau local noethérien régulier est factoriel.
\end{theorem}

\begin{proof}
La démonstration qui suit est due à I.~Kaplansky.

Soit $A$ un anneau local noethérien régulier de dimension $n$; nous allons raisonner
par récurrence sur $n$. Pour $n=0$, $A$ est un corps et pour $n=1$, un anneau de
valuation discrète, donc principal (et \emph{a fortiori} factoriel). Supposons donc
$n\geq2$ et le théorème démontré pour des anneaux réguliers de dimension $<n$.
Posons $X=\operatorname{Spec}(A)$, désignons par $a$ le point fermé de $X$ et posons
$U=X-\{a\}$. En tout point $y\in U$, on a
$\dim(\mathcal O_{X,y})\leq n-1$, et puisque $A$ est régulier, les anneaux
$\mathcal O_{X,y}$ le sont aussi $(\mathbf 0, 17.3.2)$, donc l'hypothèse de
récurrence entraîne qu'ils sont factoriels. De plus on a
$\operatorname{prof}(A)=\dim(A)\geq2$ puisque $A$ est régulier, donc de
Cohen--Macaulay $(\mathbf 0, 17.1.3)$.
\oldpage[IV-4]{303}
Utilisant (21.6.14), on est ramené à prouver que $\operatorname{Pic}(U)=0$.
Considérons donc un $\mathcal O_U$-Module inversible $\mathscr L$, et prouvons qu'il
est isomorphe à $\mathcal O_U$. Il résulte de $(\mathbf I, 9.4.5)$ qu'il existe un
$\mathcal O_X$-Module \emph{cohérent} $\mathscr F$ tel que
$\mathscr F|U=\mathscr L$. Comme $A$ est régulier, donc de dimension cohomologique
finie $(\mathbf 0, 17.3.1)$, il existe une résolution gauche \emph{finie} de
$\mathscr F$ :
\[
  0\leftarrow\mathscr F\leftarrow\mathcal O_X^{n_1}
  \leftarrow\mathcal O_X^{n_2}\leftarrow\cdots
  \leftarrow\mathcal O_X^{n_h}\leftarrow0
\]
$(\mathbf 0, 17.2.8$ et $\mathbf 0, 17.2.2$, (iii)). Par restriction à $U$, on a
donc une résolution finie
\[
\label{IV.21.11.1.1}
  0\leftarrow\mathscr L\leftarrow\mathcal O_U^{n_1}
  \leftarrow\mathcal O_U^{n_2}\leftarrow\cdots
  \leftarrow\mathcal O_U^{n_h}\leftarrow0.
  \tag{21.11.1.1}
\]
\end{proof}

Le théorème résultera alors des considérations générales suivantes. Sur un espace
annelé $X$, soit $\mathscr E$ un $\mathcal O_X$-Module localement libre de rang fini;
on désignera par $\Lambda^{\max}\mathscr E$ le $\mathcal O_X$-Module inversible qui,
au voisinage de chaque point de $X$, est égal à $\Lambda^p\mathscr E$, en désignant
par $p$ le rang de $\mathscr E$ dans ce voisinage (qui peut varier avec la composante
connexe de $X$). Avec cette notation, on a le

\begin{lemma}[21.11.1.2]
\label{IV.21.11.1.2}
Soient $X$ un espace annelé en anneaux locaux, et
\[
  0\leftarrow\mathscr E_0\xleftarrow{u_0}\mathscr E_1
  \leftarrow\cdots\leftarrow\mathscr E_h\leftarrow0
\]
une suite exacte de $\mathcal O_X$-Modules localement libres de rang fini; alors le
$\mathcal O_X$-Module inversible
\[
  \bigotimes_{0\leq i\leq h}
  (\Lambda^{\max}\mathscr E_i)^{\otimes(-1)^i}
\]
est isomorphe à $\mathcal O_X$.
\end{lemma}

\begin{proof}
Montrons comment ce lemme permettra de conclure la preuve de (21.11.1). Il suffit de
noter pour cela que, pour tout entier $n$,
$\Lambda^{\max}(\mathcal O_U^n)=\Lambda^n\mathcal O_U^n$ est isomorphe à
$\mathcal O_U$, ainsi que $\mathcal O_U^{\otimes(-1)}$. Comme d'autre part
$\Lambda^{\max}\mathscr L=\mathscr L$ pour tout $\mathcal O_U$-Module inversible
$\mathscr L$, le lemme (21.11.1.2), appliqué à la suite exacte (21.11.1.1) montre
que $\mathscr L$ est isomorphe à $\mathcal O_U$.

Reste à prouver (21.11.1.2); on procède par récurrence sur $h$, le lemme étant trivial
pour $h=1$. Pour $h>1$, $\mathscr N=\operatorname{Ker}(u_0)$ est un
$\mathcal O_X$-Module localement libre de rang fini $(\mathbf 0_{\mathrm I}, 5.5.5)$,
et on a les deux suites exactes
\[
  0\leftarrow\mathscr E_0\leftarrow\mathscr E_1
  \leftarrow\mathscr N\leftarrow0
\]
\[
  0\leftarrow\mathscr N\leftarrow\mathscr E_2
  \leftarrow\cdots\leftarrow\mathscr E_h\leftarrow0.
\]
En vertu de l'hypothèse de récurrence,
$(\Lambda^{\max}\mathscr N)\otimes
(\bigotimes_{2\leq i\leq h}(\Lambda^{\max}\mathscr E_i)^{\otimes(-1)^{i-1}})$
est isomorphe à $\mathcal O_X$; il suffira donc de définir un isomorphisme canonique
\[
  (\Lambda^{\max}\mathscr N)\otimes(\Lambda^{\max}\mathscr E_0)
  \xrightarrow{\sim}\Lambda^{\max}\mathscr E_1
\]
pour achever la démonstration. Or, il existe un recouvrement ouvert $(U_\alpha)$ de
$X$ tel que dans tout $U_\alpha$, $\mathscr E_1|U_\alpha$ soit somme directe de
$\mathscr N|U_\alpha$ et d'un $\mathcal O_{U_\alpha}$-Module localement libre
$\mathscr M_\alpha$ $(\mathbf 0_{\mathrm I}, 5.5.5)$, d'où un isomorphisme canonique
$v_\alpha:\mathscr M_\alpha\xrightarrow{\sim}\mathscr E_0|U_\alpha$.
Comme on a un isomorphisme canonique
\[
  r_\alpha:(\Lambda^{\max}\mathscr N|U_\alpha)
  \otimes(\Lambda^{\max}\mathscr M_\alpha)
  \xrightarrow{\sim}(\Lambda^{\max}\mathscr E_1)|U_\alpha,
\]
on en déduit au moyen de $v_\alpha$ un isomorphisme
\[
  u_\alpha:(\Lambda^{\max}\mathscr N|U_\alpha)
  \otimes(\Lambda^{\max}\mathscr E_0|U_\alpha)
  \xrightarrow{\sim}\Lambda^{\max}\mathscr E_1|U_\alpha.
\]
\oldpage[IV-4]{304}
Et il s'agit de montrer que $u_\alpha$ et $u_\beta$ coïncident dans
$U_\alpha\cap U_\beta$ pour deux indices quelconques $\alpha$, $\beta$. Or, si
$v'_\alpha$ et $v'_\beta$ sont les restrictions à $U_\alpha\cap U_\beta$ de
$v_\alpha$ et $v_\beta$ respectivement, on a
$v'_\alpha=v'_\beta\circ w_{\beta\alpha}$, où
\[
  w_{\beta\alpha}:\mathscr M_\alpha|(U_\alpha\cap U_\beta)
  \xrightarrow{\sim}\mathscr M_\beta|(U_\alpha\cap U_\beta)
\]
est la « projection parallèle à $\mathscr N$ » telle que pour toute section
$s\in\Gamma(U_\alpha\cap U_\beta,\mathscr M_\alpha)$,
$w_{\beta\alpha}(s)=s+t$ avec
$t\in\Gamma(U_\alpha\cap U_\beta,\mathscr N)$; l'identité de $u_\alpha$ et de
$u_\beta$ résulte aussitôt de ce fait et de la définition de l'isomorphisme canonique
$r_\alpha$ (Bourbaki, \emph{Alg.}, chap.~III, 3\textsuperscript{e} éd.).
\end{proof}

\subsection{Le théorème de pureté de Van der Waerden pour l'ensemble de
ramification d'un morphisme birationnel}
\EGAReaderTarget{ega4.c04.subsec.21-12}{subsection:IV.21.12}
\label{IV.21.12}

\begin{env}[21.12.1]
\EGAReaderTarget{ega4.c04.statement.21-12-1}{IV.21.12.1}
Soient $X$ et $U$ deux préschémas, $f:U\to X$ un morphisme quasi-compact et
quasi-séparé, de sorte que $f_*(\mathcal O_U)$ est une $\mathcal O_X$-Algèbre
quasi-cohérente (1.7.4). On appelle \emph{enveloppe affine} du $X$-préschéma $U$ le
$X$-préschéma affine sur $X$
\[
  U^0=\operatorname{Aff}(U/X)=\operatorname{Spec}(f_*(\mathcal O_U))
  =\operatorname{Spec}(\mathscr A(U))\qquad(\mathbf{II}, 1.1.1).
\]
Si $f^0:U^0\to X$ est le morphisme structural, on a donc par définition
\[
  \mathscr A(U^0)=f^0_*(\mathcal O_{U^0})=f_*(\mathcal O_U)=\mathscr A(U),
\]
et à l'isomorphisme identique de $f_*(\mathcal O_U)$ il correspond par
$(\mathbf{II}, 1.2.7)$ un $X$-morphisme canonique
\[
\label{IV.21.12.1.1}
  i_U:U\to U^0.
  \tag{21.12.1.1}
\]

Pour que $i_U$ soit un \emph{isomorphisme}, il faut et il suffit évidemment que le
morphisme $f:U\to X$ soit \emph{affine}. Pour tout $X$-préschéma $V$ \emph{affine sur
$X$}, l'application $u\mapsto u\circ i_U$ est une \emph{bijection}
\[
  \operatorname{Hom}_X(U^0,V)\xrightarrow{\sim}\operatorname{Hom}_X(U,V)
\]
fonctorielle en $V$: cela résulte de l'existence des bijections canoniques
\[
  \operatorname{Hom}_X(U,V)\xrightarrow{\sim}
  \operatorname{Hom}_{\mathcal O_X}(\mathscr A(V),\mathscr A(U))
\]
et
$\operatorname{Hom}_X(U^0,V)\xrightarrow{\sim}
\operatorname{Hom}_{\mathcal O_X}(\mathscr A(V),\mathscr A(U^0))$
$(\mathbf{II}, 1.2.7)$.

On peut donc dire que $U^0$ \emph{représente} le foncteur covariant
$V\mapsto\operatorname{Hom}_X(U,V)$ dans la catégorie des $X$-préschémas affines sur
$X$ $(\mathbf 0_{\mathrm{III}}, 8.1.11)$. On en déduit
$(\mathbf 0_{\mathrm{III}}, 8.1.7)$ que pour $X$ fixé,
$U\mapsto\operatorname{Aff}(U/X)$ est un foncteur covariant de la catégorie des
$X$-préschémas quasi-compacts et quasi-séparés sur $X$, dans la catégorie des
$X$-préschémas affines sur $X$; de façon plus précise, si $U_1$, $U_2$ sont deux
$X$-préschémas quasi-compacts et quasi-séparés sur $X$, à tout $X$-morphisme
$g:U_1\to U_2$ correspond l'unique $X$-morphisme $g^0:U_1^0\to U_2^0$ rendant
commutatif le diagramme
\[
  \xymatrix{
    U_1 \ar[r]^g \ar[d]_{i_{U_1}} & U_2 \ar[d]^{i_{U_2}} \\
    U_1^0 \ar[r]_{g^0} & U_2^0.
  }
\]
\end{env}

\oldpage[IV-4]{305}
Plus généralement, considérons un diagramme commutatif
\[
  \xymatrix{
    U \ar[r]^v \ar[d]_f & U' \ar[d]^{f'} \\
    X \ar[r]_u & X'
  }
\]
où les morphismes $f$, $f'$ sont quasi-compacts et quasi-séparés et le morphisme $u$
\emph{affine}. Si l'on pose $h=u\circ f$, on a
$h_*(\mathcal O_U)=u_*(f_*(\mathcal O_U))=u_*(f^0_*(\mathcal O_{U^0}))$ et
$u\circ f^0$ est un morphisme affine; on a par suite
$U^0=\operatorname{Aff}(U/X')$ (relativement au morphisme $h$), d'où un unique
$X'$-morphisme $v^0:U^0\to U'^0$ rendant commutatif le diagramme
\[
  \xymatrix{
    U \ar[r]^v \ar[d]_{i_U} & U' \ar[d]^{i_{U'}} \\
    U^0 \ar[r]_{v^0} & U'^0.
  }
\]

\begin{proposition}[21.12.2]
\label{IV.21.12.2}
Soient $f:U\to X$ un morphisme quasi-compact et quasi-séparé,
$h:X'\to X$ un morphisme plat; on pose $U'=U\times_X X'$,
$f'=f_{\{X'\}}:U'\to X'$. Alors on a un $X'$-isomorphisme canonique
\[
\label{IV.21.12.2.1}
  \operatorname{Aff}(U'/X')\xrightarrow{\sim}
  \operatorname{Aff}(U/X)\times_X X'.
  \tag{21.12.2.1}
\]
\end{proposition}

\begin{proof}
En effet, on a
$\operatorname{Aff}(U'/X')=\operatorname{Spec}(f'_*(\mathcal O_{U'}))$, et
$\operatorname{Aff}(U/X)\times_X X'=\operatorname{Spec}(h^*(f_*(\mathcal O_U)))$
$(\mathbf{II}, 1.5.1)$; l'isomorphisme de l'énoncé provient de l'isomorphisme canonique
$h^*(f_*(\mathcal O_U))\xrightarrow{\sim}f'_*(\mathcal O_{U'})$ (2.3.1).
\end{proof}

\begin{corollary}[21.12.3]
\label{IV.21.12.3}
Pour tout morphisme quasi-compact et quasi-séparé $f:U\to X$ et tout $x\in X$, on a,
à un isomorphisme canonique près
\[
  U^0\times_X\operatorname{Spec}(\mathcal O_{X,x})
  =(U\times_X\operatorname{Spec}(\mathcal O_{X,x}))^0.
\]
\end{corollary}

Il résulte aussi de (21.12.2) que l'on a, à un isomorphisme canonique près, pour tout
ouvert $V$ de $X$
\[
\label{IV.21.12.4}
  (f^{-1}(V))^0=(f^0)^{-1}(V).
  \tag{21.12.4}
\]

\begin{env}[21.12.5]
\EGAReaderTarget{ega4.c04.statement.21-12-5}{IV.21.12.5}
Nous allons considérer en particulier le cas où $f:U\to X$ est une \emph{immersion
ouverte}, $U$ s'identifiant donc à un préschéma induit sur un ouvert de $X$. Comme le
morphisme $f^{-1}(U)\to U$ est l'identité, il résulte de (21.12.4) que le morphisme
$(f^0)^{-1}(U)\to U$ restriction de $f^0$ est un isomorphisme, $i_U:U\to U^0$ étant
donc une \emph{immersion ouverte} permettant d'identifier $U$ à un préschéma induit
sur un ouvert de $U^0$.

Plus généralement, pour un ouvert $V$ de $X$, la restriction $(f^0)^{-1}(V)\to V$ de
$f^0$ est un isomorphisme de $(f^0)^{-1}(V)$ sur le préschéma induit sur l'ouvert
$V\cap U$ si et seulement si l'immersion ouverte $U\cap V\to V$ est un morphisme
\emph{affine}. Il est clair $(\mathbf{II}, 1.2.1)$
\oldpage[IV-4]{306}
que la réunion de ces ouverts $V$ est le \emph{plus grand} d'entre eux, $U_1$, qui
contient $U$; en vertu de ce qui précède, $U_1$ est aussi le plus grand ouvert ne
rencontrant pas l'ensemble
\[
  \operatorname{Daf}(U/X)=f^0(U^0)-U
\]
(« défaut d'affinité » de l'ouvert $U$ relativement à $X$, qui est vide si et seulement
si $U$ est affine sur $X$); autrement dit, l'ensemble fermé $Z=X-U_1$ est
l'\emph{adhérence} de l'ensemble $\operatorname{Daf}(U/X)$.

On notera que pour tout morphisme \emph{plat} $h:X'\to X$, si l'on pose
$U'=h^{-1}(U)$, on a
\[
\label{IV.21.12.5.1}
  \operatorname{Daf}(U'/X')=h^{-1}(\operatorname{Daf}(U/X))
  \tag{21.12.5.1}
\]
comme il résulte aussitôt de (21.12.2.1) et de $(\mathbf I, 3.4.8)$. En particulier,
pour tout ouvert $V$ de $X$, on a
\[
\label{IV.21.12.5.2}
  \operatorname{Daf}((U\cap V)/V)=\operatorname{Daf}(U/X)\cap V
  \tag{21.12.5.2}
\]
et pour tout $x\in X$,
\[
\label{IV.21.12.5.3}
  \operatorname{Daf}((U\cap\operatorname{Spec}(\mathcal O_{X,x}))
  /\operatorname{Spec}(\mathcal O_{X,x}))
  =\operatorname{Daf}(U/X)\cap\operatorname{Spec}(\mathcal O_{X,x}).
  \tag{21.12.5.3}
\]
\end{env}

Nous allons, lorsque $X$ est localement noethérien, donner des précisions sur la nature
de l'ensemble $\operatorname{Daf}(U/X)$, qui montreront par exemple que lorsque $U$ est
partout dense dans $X$, $\operatorname{Daf}(U/X)$ n'est pas un ensemble fermé rare
quelconque :

\begin{theorem}[21.12.6]
\label{IV.21.12.6}
Soient $X$ un préschéma localement noethérien, $U$ un ouvert non vide de $X$,
$f:U\to X$ l'injection canonique. Alors :
\begin{enumerate}[label=(\roman*)]
  \item L'adhérence $Z=X-U_1$ de $\operatorname{Daf}(U/X)$ est de codimension
  $\geq2$ dans $X$.
  \item Si $T=X-U\supset Z$ est de codimension $\geq2$, le morphisme
  $f^0:U^0\to X$ est surjectif.
\end{enumerate}
\end{theorem}

\begin{proof}
(i) Montrons d'abord que pour tout point $x\in\operatorname{Daf}(U/X)$ on a
nécessairement $\dim(\mathcal O_{X,x})>1$. En effet, $x$ ne peut évidemment être point
maximal de $X$, étant contenu dans $\overline U-U$; il s'agit donc de voir que l'on ne
peut avoir $\dim(\mathcal O_{X,x})=1$. Or, cette relation entraînerait, par (21.12.5.3),
$x\in\operatorname{Daf}(U^{(x)}/\operatorname{Spec}(\mathcal O_{X,x}))$, où
$U^{(x)}=U\cap\operatorname{Spec}(\mathcal O_{X,x})$. Mais les seuls ouverts de
$\operatorname{Spec}(\mathcal O_{X,x})$ sont $\operatorname{Spec}(\mathcal O_{X,x})$
lui-même et les parties de l'ensemble (fini) des points maximaux de
$\operatorname{Spec}(\mathcal O_{X,x})$. Or, l'ensemble ouvert réduit à un point maximal
de $\operatorname{Spec}(\mathcal O_{X,x})$ est affine, par définition des préschémas; on
en conclut que \emph{tous} les ensembles ouverts dans
$\operatorname{Spec}(\mathcal O_{X,x})$ sont affines, donc
$\operatorname{Daf}(U^{(x)}/\operatorname{Spec}(\mathcal O_{X,x}))=\varnothing$
contrairement à l'hypothèse faite.

Pour prouver (i) il faut montrer davantage, savoir que si $x\in X$ est tel que
$\dim(\mathcal O_{X,x})=1$, il existe un voisinage ouvert $W$ de $x$ dans $X$ tel que
$W$ ne rencontre pas $\operatorname{Daf}(U/X)$, c'est-à-dire tel que l'injection
canonique $f_W:U\cap W\to W$ soit affine. Mais, avec les mêmes notations que ci-dessus,
on vient de voir que l'injection canonique
$f^{(x)}:U^{(x)}\to\operatorname{Spec}(\mathcal O_{X,x})$ est affine. On peut
évidemment se borner au cas où $X$ est noethérien; comme le morphisme $f$ est de
présentation finie, la conclusion résulte de (8.10.5, (viii)) appliqué suivant la
méthode de (8.1.2, a)).

\oldpage[IV-4]{307}
(ii) Il s'agit de prouver que pour tout point $x\in T$, on a $x\in f^0(U^0)$; nous
allons montrer d'abord qu'on peut se réduire au cas où $X=\operatorname{Spec}(A)$, où
$A$ est un anneau local noethérien \emph{complet}, et $x$ le point fermé de $X$. Pour
cela, il suffit de faire le changement de base
$h:X'=\operatorname{Spec}(\widehat{\mathcal O}_{X,x})\to X$; si l'on pose
$U'=h^{-1}(U)$, $f'=f_{\{X'\}}$ est l'injection canonique $U'\to X'$, et puisque le
morphisme $h$ est plat, il résulte de (21.12.2) que si l'on prouve que $x$ appartient à
$f^{\prime 0}(U^{\prime 0})$, on en déduira que $x\in f^0(U^0)$. En vertu de (6.1.1), on a
bien opéré la réduction cherchée.

Soit alors $X_1$ un sous-préschéma fermé réduit de $X$ ayant pour espace sous-jacent
une composante irréductible de $X$, de dimension \emph{maximale}, parmi celles qui
contiennent une composante irréductible de $T$, et posons
$U_1=U\cap X_1$, $T_1=T\cap X_1=X_1-U_1$; on a donc
$\operatorname{codim}(T_1,X_1)\geq2$ $(\mathbf 0, 14.2.1)$, et le couple $(U_1,X_1)$
vérifie donc les mêmes hypothèses que le couple $(U,X)$ ($X_1$ étant spectre d'un
anneau quotient de $A$, donc local noethérien et complet). On a par ailleurs un
diagramme commutatif
\[
  \xymatrix{
    U_1 \ar[r]^j \ar[d]_{f_1} & U \ar[d]^f \\
    X_1 \ar[r]_i & X
  }
\]
où $i$ et $j$ sont les injections canoniques, $i$ étant donc un morphisme affine; on en
déduit (21.12.1) l'existence d'un morphisme $j^0:U_1^0\to U^0$ rendant commutatif le
diagramme
\[
  \xymatrix{
    U_1^0 \ar[r]^{j^0} \ar[d]_{f_1^0} & U^0 \ar[d]^{f^0} \\
    X_1 \ar[r]_i & X
  }
\]
et par suite, pour prouver que $x\in f^0(U^0)$, il suffit de prouver que
$x\in f_1^0(U_1^0)$. On peut donc remplacer $X$ par $X_1$, et on peut par suite
supposer que l'anneau $A$ est en outre \emph{intègre}. Mais $A$ est quotient d'un
anneau noethérien régulier en vertu du théorème de Cohen $(\mathbf 0, 19.8.8$, (i)) et
puisqu'il est intègre, on peut appliquer (5.11.1) avec la famille $(x_\alpha)$ réduite au
seul point maximal de $X$; comme $\operatorname{codim}(T,X)\geq2$ par hypothèse, on
voit que $f_*(\mathcal O_U)$ est un $\mathcal O_X$-Module \emph{cohérent}, donc le
morphisme $f^0:U^0\to X$ est fini; comme ce morphisme est dominant ($U$ étant partout
dense), il est surjectif $(\mathbf{II}, 6.1.10)$, et par suite $x\in f^0(U^0)$. C.Q.F.D.
\end{proof}

\begin{corollary}[21.12.7]
\label{IV.21.12.7}
Soient $X$ un préschéma localement noethérien, $U$ un sous-préschéma induit sur un
ouvert de $X$, $j:U\to X$ l'immersion canonique; supposons que $j$ soit un morphisme
affine. Alors toute composante irréductible de $T=X-U$ est de codimension $\leq1$
(et par suite de codimension $1$ si $U$ est partout dense).
\end{corollary}

\begin{proof}
Supposons en effet qu'une des composantes irréductibles $T_1$ de $T$ soit de
codimension $\geq2$ dans $X$. Remplaçant au besoin $X$ par un voisinage ouvert du
point générique de $T_1$, on peut supposer $T$ irréductible et de codimension $\geq2$.
Mais alors l'hypothèse
\oldpage[IV-4]{308}
que $j:U\to X$ est affine implique que $U^0$ s'identifie à $U$ et $j^0$ à $j$,
autrement dit que $j^0(U^0)=U$; or, cela contredit la conclusion de (21.12.6) qui,
sous l'hypothèse de dimension faite implique que $j^0$ doit être surjectif.
\end{proof}

\begin{env}[21.12.8]
\label{IV.21.12.8}
Soient $A$ un anneau local noethérien, $Y=\operatorname{Spec}(A)$, $y$ l'unique point
fermé de $Y$, $Y'=Y-\{y\}$. Considérons la condition suivante :
\begin{itemize}
  \item[\textup{(W)}] \emph{Pour tout ouvert $U$ de $Y$ contenu dans $Y'$, ne
  contenant aucune composante irréductible de $Y'$ et tel que l'immersion canonique
  $U\to Y'$ soit affine, $U$ est lui-même un ouvert affine.}
\end{itemize}
Cette condition est entraînée par la suivante :
\begin{itemize}
  \item[\textup{(W')}] \emph{Pour toute partie fermée $T$ de $Y$ dont toute composante
  irréductible est de codimension $1$ dans $Y$, et tel que pour toute composante
  irréductible $Y_i$ de $Y$, on ait $Y_i\cap T\ne\{y\}$, l'ouvert $U=Y-T$ est affine.}
\end{itemize}
En effet, si (W') est vérifiée et si l'ouvert $U$ de $Y$ vérifie l'hypothèse de (W), il
résulte de (21.12.7) qu'aucune composante irréductible de $T=Y-U$ ne peut être de
codimension $\geq2$; comme par ailleurs $U$ ne contient aucune des composantes
irréductibles $Y_i-\{y\}$ de $Y'$, la condition (W') montre que $U$ est affine.

On notera que la condition (W') se simplifie lorsque $Y$ est \emph{irréductible} et est
alors équivalente à la suivante :
\begin{itemize}
  \item[\textup{(W'')}] \emph{Pour toute partie fermée irréductible $T$ de $Y$, de
  codimension $1$ dans $Y$, $Y-T$ est un ouvert affine.}
\end{itemize}
En effet, il est clair que (W') entraîne (W'') lorsque $Y$ est irréductible, et la
réciproque est vraie en considérant les composantes irréductibles de $T$ et en notant
que l'intersection d'un nombre fini d'ouverts affines est un ouvert affine
$(\mathbf I, 5.5.6)$.
\end{env}

\begin{env}[21.12.9]
\label{IV.21.12.9}
\emph{Exemples}. ---
Si $A$ est un anneau local noethérien \emph{factoriel}, il vérifie (W') (et
\emph{a fortiori} (W)), puisque tout idéal premier de hauteur $1$ est principal
$(\mathbf I, 1.3.6)$. Mais il y a des anneaux locaux noethériens non factoriels
vérifiant (W''), par exemple ceux de dimension $\leq1$: en effet, on a remarqué dans la
démonstration de (21.12.6, (i)) que \emph{tous} les ouverts de $Y$ sont alors affines.
On peut d'autre part prouver en utilisant la théorie de la dualité locale (chap.~III,
3\textsuperscript{e} partie) que tout anneau local noethérien de dimension $2$ vérifie
(W').
\end{env}

L'intérêt de la condition (W) réside dans le résultat suivant :

\begin{proposition}[21.12.10]
\label{IV.21.12.10}
Soient $X$, $Y$ deux préschémas localement noethériens, $g:X\to Y$ un morphisme,
$y$ un point fermé de $Y$, $Y'=Y-\{y\}$, $X'=g^{-1}(Y')$; supposons que le morphisme
$g':X'\to Y'$ restriction de $g$ soit une immersion ouverte, et que l'anneau local
$\mathcal O_{Y,y}$ vérifie la condition (W) (21.12.8). Alors, pour toute composante
irréductible $Z$ de $X_y=g^{-1}(y)$, ou bien $Z$ est de codimension $\leq1$ dans $X$,
ou bien son point générique est isolé dans $Z$. Si $Z$ est localement de type fini sur
$k(y)$, la seconde alternative implique que $Z$ est réduit à un seul point.
\end{proposition}

\begin{proof}
La dernière assertion de l'énoncé résulte de ce que $Z$ est alors un préschéma de
Jacobson (10.4.7), et comme l'ensemble des points fermés de $Z$ est alors dense dans
$Z$, le point générique de $Z$ ne peut être isolé dans $Z$ que si $Z$ est réduit à un
seul point.

\oldpage[IV-4]{309}
Supposons qu'il existe dans $X_y$ une composante irréductible $Z$ dont le point
générique $z$ n'est pas isolé dans $Z$, et telle que $\operatorname{codim}(Z,X)\geq2$.
La question étant locale sur $X$ et $Y$ puisque $z$ est non isolé dans $Z$, on peut, en
remplaçant $X$ par un voisinage ouvert de $z$ dans $X$, supposer $X$ et $Y$ affines,
$X_y=Z$ irréductible, non réduit à un point et telle que
$\operatorname{codim}(X_y,X)\geq2$. L'image $g(X-X_y)=U$ est un ouvert de $Y'$
isomorphe à $X-X_y$ par hypothèse; montrons qu'en remplaçant au besoin $X$ par un
voisinage ouvert de $z$ dans $X$, on peut supposer que $U$ ne contient aucune des
composantes irréductibles de $Y'$ dont l'adhérence contient $y$.

En effet, on peut tout d'abord supposer que tous les points maximaux de $X$ sont des
générisations de $z$, l'ensemble de ces points étant fini; donc l'ensemble des points
maximaux de $U$ est l'ensemble des images $y_i$ par $g$ des points maximaux $x_i$ de
$X$ (aucun point maximal de $X$ ne pouvant par hypothèse être contenu dans $X_y$,
puisque $\dim(\mathcal O_{X,z})\geq2$). Posons $X_i=\overline{\{x_i\}}$. Par
hypothèse, on a $z\in X_i$ et puisque $z$ n'est pas isolé dans $X_y$, il existe dans
$X_y$ un point $x\ne z$. Comme $X_i\ne Z$, il existe dans $X_i$ un point $t_i$ tel
que si l'on pose $T_i=\overline{\{t_i\}}$, on ait $\dim(\mathcal O_{T_i,x})=1$
(10.5.9); cela entraîne par hypothèse $t_i\notin Z$, donc $t_i$ n'est pas une
générisation de $z$. Remplaçant $X$ par $X'=X-\bigcup_iT_i$, on voit que l'image par
$g$ de $X'-X_y$ ne contient pas les images $g(t_i)$, donc ne contient aucune des
composantes irréductibles de $Y'$ dont l'adhérence contient $y$.

Cela étant, comme $X$ et $Y$ sont affines, le morphisme $g:X\to Y$ est affine, donc il
en est de même de la restriction $g':X'\to Y'$; puisque $g'$ est une immersion ouverte,
l'immersion canonique $U\to Y'$ est affine. Posons
$Y_1=\operatorname{Spec}(\mathcal O_{Y,y})$, $Y'_1=Y_1-\{y\}=Y'\cap Y_1$,
$U_1=U\cap Y_1$. Ce qui précède prouve que l'immersion canonique $U_1\to Y'_1$ est
affine, donc $U_1$ est un ouvert affine dans $Y_1$ en vertu de (W), en d'autres termes
l'immersion canonique $U_1\to Y_1$ est affine. Mais puisque cette immersion est de
présentation finie (1.6.2), il résulte de (8.10.5, (viii)) appliqué suivant la méthode de
(8.1.2, a)) qu'en restreignant au besoin $Y$ à un voisinage ouvert affine de $y$, on peut
supposer que l'immersion $U\to Y$ est affine. On en conclurait que l'ouvert $X-X_y$ de
$X$, isomorphe à $U$, serait affine, ce qui contredirait (21.12.7); la proposition est
ainsi démontrée.
\end{proof}

\begin{corollary}[21.12.11]
\label{IV.21.12.11}
Soient $X$, $Y$ deux préschémas localement noethériens, $X$ étant supposé
irréductible; soit $g:X\to Y$ un morphisme localement de type fini, et soit $V$ le plus
grand ouvert de $X$ tel que la restriction $g|V:V\to Y$ soit un isomorphisme local.
Supposons que pour tout point $y\in Y$, l'anneau local $\mathcal O_{Y,y}$ vérifie la
condition (W) (21.12.8). Alors toute composante irréductible de $T=X-V$ est, ou bien
de codimension $\leq1$, ou bien telle que son point générique $z$ soit isolé dans
$g^{-1}(g(z))$.
\end{corollary}

\begin{proof}
Posons $g(z)=y$. La question ne dépendant que de la fibre $g^{-1}(y)$ et de l'anneau
$\mathcal O_{Y,y}$, on peut par changement de base se borner au cas où
$Y=\operatorname{Spec}(\mathcal O_{Y,y})$ et où $X$ est affine
$((\mathbf I, 3.6.5)$ et $(\mathbf I, 4.5.5))$; remplaçant éventuellement $X$ par un
voisinage ouvert de $z$ dans $X$, on peut supposer que $T$ est irréductible; en outre on
peut se borner au cas où $V$ est non vide. Le morphisme $g$ peut donc être supposé
séparé et $y$ fermé dans $Y$; puisque la restriction $g|V:V\to Y$ est un isomorphisme
local et que l'ouvert non vide $V$ de $X$ est irréductible, il résulte de
$(\mathbf I, 8.2.8)$ que $g|V:V\to Y$ est une immersion ouverte. La restriction
$g|(V\cap X_y):V\cap X_y\to\{y\}=\operatorname{Spec}(k(y))$ est donc aussi une
immersion ouverte, ce qui montre que $V\cap X_y$ est, soit vide, soit réduit à un point
$x$ rationnel
\oldpage[IV-4]{310}
sur $k(y)$, donc fermé dans $X_y$ $(\mathbf I, 6.4.2)$. Remplaçant encore
éventuellement $X$ par un voisinage ouvert plus petit de $z$ dans $X$, on peut donc se
borner au cas où $V\cap X_y=\varnothing$, autrement dit $T\supset X_y$; comme d'autre
part $z\in X_y$ est le point générique de $T$, on a
$g(T)\subset\overline{\{y\}}=\{y\}$, autrement dit $T=X_y$. On est alors exactement
dans les conditions d'application de (21.12.10), d'où la conclusion.
\end{proof}

\begin{theorem}[21.12.12]
\label{IV.21.12.12}
\textup{(van der Waerden)}. ---
Soient $X$, $Y$ deux préschémas localement noethériens intègres, $g:X\to Y$ un
morphisme birationnel et localement de type fini. On suppose en outre que $Y$ est
normal et que pour tout $y\in Y$, l'anneau $\mathcal O_{Y,y}$ vérifie la condition (W)
de (21.12.8) (conditions remplies en particulier lorsque $Y$ est localement factoriel
(21.12.9)). Si $V$ est le plus grand ouvert de $X$ tel que la restriction
$g|V:V\to Y$ soit un isomorphisme local, toutes les composantes irréductibles de
$T=X-V$ sont de codimension $1$ dans $X$.
\end{theorem}

\begin{proof}
On notera que puisque $g$ est birationnel, l'ensemble ouvert $V$ n'est pas vide; il
suffit donc de prouver qu'un point maximal $z$ de $T$ ne peut être isolé dans $X_y$,
où $y=g(z)$. Mais, quitte à se restreindre à un voisinage ouvert de $z$, on peut se
borner au cas où $T=X_y$; alors toutes les fibres $g^{-1}(y')$ ($y'\in Y$) seraient
vides ou réduites à un point et il résulterait du « Main theorem » (8.12.10) que $g$
serait un isomorphisme local, contrairement à l'hypothèse $z\in T$.
\end{proof}

\begin{corollary}[21.12.13]
\label{IV.21.12.13}
Supposons vérifiées les hypothèses de (21.12.12) et en outre, supposons que $g$ soit
quasi-fini en chacun des points de $X^{(1)}$ (on rappelle que ce sont les points où
$\dim(\mathcal O_{X,x})=1$); alors $g$ est un isomorphisme local, et si de plus $g$ est
séparé, $g$ est une immersion ouverte.
\end{corollary}

\begin{proof}
Il suffit de prouver la première assertion, la seconde étant une conséquence de la
première et de $(\mathbf I, 8.2.8)$. Tout revient à prouver, avec les notations de
(21.12.12), que $T=\varnothing$; dans le cas contraire, un point générique $z$ d'une
composante irréductible de $T$ appartiendrait à $X^{(1)}$ en vertu de (21.12.12), et par
hypothèse il serait isolé dans $X_y$ avec $y=g(z)$; mais on a vu dans la démonstration
de (21.12.12) que cela n'est pas possible.
\end{proof}

\begin{remarks}[21.12.14]
\label{IV.21.12.14}
(i) La conclusion de (21.12.12) s'applique lorsque $Y$ est régulier et intègre et $X$
intègre, car en vertu du théorème d'Auslander--Buchsbaum (21.11.1), $Y$ est alors
localement factoriel. Par contre, on connaît des exemples où $X$ et $Y$ sont des schémas
algébriques normaux de dimension $3$, sur un corps algébriquement clos de
caractéristique quelconque, et où la conclusion de (21.12.12) n'est plus valable.

(ii) L'ensemble $T$ de (21.12.12) est aussi l'ensemble des points où $g$ est
\emph{ramifié}. En effet, si $g$ est non ramifié en un point $x\in X$, il est aussi non
ramifié dans un voisinage ouvert affine $U$ de $x$ (17.3.7), donc $g|U$ est un morphisme
séparé, quasi-fini et birationnel (17.4.3). Puisque $Y$ est normal, on conclut du « Main
theorem » $(\mathbf{III}, 4.4.9)$ que $g$ est un isomorphisme local au point $x$, donc
$x\in X-T$. Réciproquement, $g$ est évidemment non ramifié en tout point où il est un
isomorphisme local. Ceci justifie le titre donné à cette section.

\oldpage[IV-4]{311}
(iii) Sans supposer que $Y$ soit localement factoriel, mais en supposant par contre que
$X$ est intersection complète relativement à $Y$ (chap.~V), nous verrons au chap.~V
qu'on a un résultat plus précis que (21.12.12), en exprimant $T$ comme cycle
1-codimensionnel d'une section d'un $\mathcal O_X$-Module inversible, canoniquement
associée à $g$. Ceci s'appliquera notamment quand $X$ et $Y$ sont tous deux réguliers,
ou lorsque $g$ est un $S$-morphisme, $X$ et $Y$ étant des préschémas lisses sur $S$.

(iv) On peut se demander si (21.12.10) admet une réciproque; autrement dit, pour un
anneau local noethérien donné $A$, si l'on pose $Y=\operatorname{Spec}(A)$, $y$
désignant le point fermé de $Y$, et si, pour tout préschéma localement noethérien $X$ et
tout morphisme $g:X\to Y$ tel que $g':X'\to Y'$ soit une immersion ouverte, toute
composante irréductible de $X_y$ est de codimension $\leq1$ dans $X$ ou est réduite à
un point, alors est-il vrai que $A$ vérifie la condition (W) de (21.12.8)?

(v) Soient $Y$ un préschéma intègre \emph{régulier}, $X$ un préschéma localement
noethérien \emph{normal}, $g:X\to Y$ un morphisme localement de type fini; supposons
en outre que, si $\eta$ est le point générique de $Y$, la fibre $X_\eta$ soit \emph{étale}
sur $k(\eta)$, et soit $T'$ le complémentaire du plus grand ouvert $V'$ de $X$ tel que
$g|V':V'\to Y$ soit étale; est-il vrai alors que toutes les composantes irréductibles de
$T'$ soient de codimension $1$? Il en est bien ainsi lorsqu'on suppose en outre que $g$
est \emph{localement quasi-fini} (« théorème de pureté » de Zariski--Nagata). On peut
montrer que la conjecture précédente serait conséquence de la suivante (et lui serait
même équivalente si la conjecture de (iv) était vérifiée): pour un anneau local régulier
$A$ contenu dans un anneau local intègre et intégralement clos $B$, tel que $B$ soit un
$A$-module fini, l'ouvert $U$ des points de $\operatorname{Spec}(B)$ en lesquels
$\operatorname{Spec}(B)$ est non ramifié sur $\operatorname{Spec}(A)$ (ou, ce qui
revient au même par (18.10.1), étale sur $\operatorname{Spec}(A)$) est affine.
\end{remarks}

\begin{proposition}[21.12.15]
\label{IV.21.12.15}
Soient $S$ un préschéma, $g:X\to S$ un morphisme plat et localement de présentation
finie, dont les fibres $X_s=g^{-1}(s)$ sont géométriquement irréductibles (4.5.2),
$h:Y\to S$ un morphisme lisse; on note $Y_s=h^{-1}(s)$ les fibres de ce morphisme.
Soit $f:X\to Y$ un $S$-morphisme propre tel que, pour tout point maximal $\eta$ de
$S$, le morphisme $f_\eta:X_\eta\to Y_\eta$ soit un isomorphisme. Alors $f$ est un
isomorphisme.
\end{proposition}

\begin{proof}
Puisque $h$ est plat, tout point maximal de $Y$ est au-dessus d'un point maximal de $S$
(2.3.4), donc la réunion des $Y_\eta$, lorsque $\eta$ parcourt l'ensemble des points
maximaux de $S$, est dense dans $Y$; comme les $f_\eta$ sont des isomorphismes, $f$
est dominant et par suite \emph{surjectif} puisqu'il est propre. Il suffit donc de montrer
que $f$ est une \emph{immersion ouverte}. Compte tenu de (17.9.5), il suffit de prouver
que pour tout $s\in S$, $f_s:X_s\to Y_s$ est une immersion ouverte.

Comme tout $s\in S$ est générisation d'un point maximal $\eta$, on peut déjà, en
faisant le changement de base $\operatorname{Spec}(\mathcal O_{S',s})\to S$, où $S'$
est le sous-préschéma réduit de $S$ ayant $\overline{\{\eta\}}$ comme espace
sous-jacent, se ramener au cas où $S$ est un schéma local et intègre de point fermé $s$:
les fibres aux points $\eta$ et $s$ ne sont en effet pas changées et les propriétés des
morphismes $g$, $h$ et $f$ se conservent par changement de base. En outre, la question
est locale sur $Y$; remplaçant $Y$ par un voisinage ouvert affine $U$ dans $Y$ d'un point
de $Y_s$, et $X$ par $f^{-1}(U)$, qui est quasi-compact puisque $f$ est propre, on voit
qu'on peut supposer en outre que $X$ et $Y$ sont de présentation finie sur $S$.

Il existe alors un sous-anneau local noethérien $A_0$ de
$A=\mathcal O_{S,s}$ tel que $A_0\to A$ soit un homomorphisme local, deux morphismes
de type fini $g_0:X_0\to S_0=\operatorname{Spec}(A_0)$,
$h_0:Y_0\to S_0$ et un $S_0$-morphisme
\oldpage[IV-4]{312}
$f_0:X_0\to Y_0$ tels que $g$, $h$ et $f$ se déduisent de $g_0$, $h_0$, $f_0$ par le
changement de base $S\to S_0$ ((8.9.1) et (5.13.3)); on peut en outre supposer $g_0$
plat (11.2.7), $h_0$ lisse (17.7.9) et $f_0$ propre (8.10.5, (xii)). D'autre part, la
projection sur $S_0$ du point générique $\eta$ de $S$ est le point générique $\eta_0$ de
$S_0$ et il résulte de (2.7.1, (viii)) que le morphisme
$(f_0)_{\eta_0}:(X_0)_{\eta_0}\to(Y_0)_{\eta_0}$ est un isomorphisme; enfin le point
fermé $s$ de $S$ est au-dessus du point fermé $s_0$ de $S_0$, donc la fibre
$(X_0)_{s_0}$ de $g_0$ est géométriquement irréductible (4.5.6).

On voit donc qu'on peut se borner au cas où $S$ est le spectre d'un anneau local intègre
noethérien, de point fermé $s$, affaiblir l'hypothèse sur les fibres en supposant seulement
que $X_s$ et $X_\eta$ sont géométriquement irréductibles, et prouver que sous ces
hypothèses $f_s$ est une immersion ouverte. Il y a alors un anneau de valuation discrète
$A'$ et un morphisme $S'=\operatorname{Spec}(A')\to S$ qui transforme le point fermé
$a$ de $S'$ en $s$ et le point générique $b$ de $S'$ en $\eta$ $(\mathbf{II}, 7.1.9)$;
soient $g':X'\to S'$, $h':Y'\to S'$, $f':X'\to Y'$ les morphismes déduits de $g$,
$h$, $f$ par le changement de base $S'\to S$, qui vérifient de nouveau les hypothèses
vérifiées par $g$, $h$, et $f$ dans l'énoncé (21.12.15); si l'on prouve que
$f'_a:X'_a\to Y'_a$ est une immersion ouverte, il résultera de (2.7.1, (x)) qu'il en
sera de même de $f_s:X_s\to Y_s$. On peut donc finalement se borner au cas où $S$ est
le \emph{spectre d'un anneau de valuation discrète}.

Alors, puisque $h:Y\to S$ est lisse, $Y$ est régulier (17.5.8), et puisque la question
est locale sur $Y$, on peut se borner au cas où $Y$ est intègre. Comme $g:X\to S$ est
plat, les points maximaux de $X$ sont contenus dans $X_\eta$ (2.3.4) et puisque
$f_\eta$ est un isomorphisme, $X$ est irréductible et l'anneau local en son point
générique est un corps; de plus, par platitude (3.3.2), on voit que $X$ n'a pas de cycles
premiers immergés, donc $X$ est réduit (3.2.1) et par suite intègre et $f$ est un
morphisme birationnel et séparé. Pour prouver que $f$ est une immersion ouverte, on peut
donc appliquer le critère (21.12.13); pour voir que $f$ est quasi-fini aux points de
$X^{(1)}$, on remarque que l'assertion est évidente en ceux de ces points qui
appartiennent à $X_\eta$; le seul point de $X^{(1)}$ appartenant à $X_s$ est le point
générique $x$ de $X_s$, en vertu de (6.1.1). Or, il résulte de (2.4.6) et de (14.2.2)
que les morphismes $g$ et $h$ sont équidimensionnels; comme $X_\eta$ et $Y_\eta$ sont
isomorphes, les composantes irréductibles de $X_s$ et de $Y_s$ ont toutes même
dimension. Mais par hypothèse $X_s$ est irréductible et on a vu que $f$ est surjectif,
donc il en est de même de $f_s:X_s\to Y_s$, ce qui entraîne que $Y_s$ est aussi
irréductible; si $y$ est le point générique de $Y_s$, on a donc $f(x)=y$, et la relation
$\dim(X_s)=\dim(Y_s)$ entraîne alors, en vertu de (5.6.6), que
$\dim(f_s^{-1}(y))=0$, autrement dit $f_s$ (et par suite aussi $f$) est bien quasi-fini
au point $x$, ce qui achève la démonstration.
\end{proof}

\begin{corollary}[21.12.16]
\label{IV.21.12.16}
Soient $g:X\to S$ un morphisme propre, plat et de présentation finie, $h:Y\to S$ un
morphisme lisse, propre et de présentation finie, $f:X\to Y$ un $S$-morphisme. On
suppose que les fibres $X_s=g^{-1}(s)$ de $g$ sont géométriquement irréductibles. Soit
$U$ l'ensemble des $s\in S$ tels que $f_s:X_s\to Y_s$ soit un isomorphisme. Alors
$U$ est ouvert et fermé dans $S$, et le morphisme $g^{-1}(U)\to h^{-1}(U)$,
restriction de $f$, est un isomorphisme.
\end{corollary}

\begin{proof}
La dernière assertion résultera de la première et de (17.9.5). On sait déjà (9.6.1,
(xi)) que $U$ est localement constructible dans $S$. En vertu de (1.10.1), il suffit de
prouver que $U$ est stable par \emph{spécialisation} et par \emph{générisation}. Pour
démontrer ces assertions on
\oldpage[IV-4]{313}
peut, par un changement de base $\operatorname{Spec}(\mathcal O_{S,s})\to S$, se
ramener au cas où $S$ est un schéma local de point fermé $s$ et de point générique
$\eta$, et il faut montrer que, pour que $f_s$ soit un isomorphisme, il faut et il suffit
que $f_\eta$ le soit. Or, la suffisance de cette condition résulte de (21.12.15). Pour
montrer qu'elle est nécessaire, on raisonne comme dans la démonstration de (21.12.15)
(en remarquant, avec les mêmes notations, que la projection du point fermé de $S$ est le
point fermé de $S_0$) pour se ramener au cas où en outre $S$ est \emph{noethérien};
mais alors la conclusion résulte de $(\mathbf{III}, 4.6.7$, (ii)).
\end{proof}

\begin{remarks}[21.12.17]
\label{IV.21.12.17}
(i) La conclusion de la proposition (21.12.15) n'est plus valable si on élimine
l'hypothèse que les fibres $X_s$ sont irréductibles. Prenons en effet
$S=\operatorname{Spec}(A)$, où $A$ est un anneau de valuation discrète,
$Y=\mathbb P^1_S$, qui est propre et lisse sur $S$ (17.3.9). Notons encore $s$ et
$\eta$ le point fermé et le point générique de $S$; soit $z$ un point fermé de $Y_s$,
par exemple un point rationnel sur $k=k(s)$, et soit $X$ le $Y$-préschéma obtenu en
faisant éclater le point $z$. Comme l'anneau de polynômes $A[T]$ est régulier
$(\mathbf 0, 17.3.7)$ et de dimension $2$, on voit comme dans la démonstration de
(15.1.1.6) que, si $f:X\to Y$ est le morphisme structural, $f^{-1}(z)$ est isomorphe
à $\mathbb P^1_k$, et d'autre part, le complémentaire de $f^{-1}(z)$ dans $X_s$ est
isomorphe à $Y_s-\{z\}$, donc $Z_1=f^{-1}(z)$ et $Z_2$, adhérence dans $X_s$ du
complémentaire de $Z_1$, sont les deux composantes irréductibles de $X_s$. Par
ailleurs, $f$ est évidemment propre et $g=h\circ f$ est plat, car $X$ est intègre
$(\mathbf{II}, 8.1.4)$ et pour tout ouvert affine $U$ de $X$, l'homomorphisme
$A\to\Gamma(U,\mathcal O_X)$ est injectif $(\mathbf I, 1.2.7)$, donc
$\Gamma(U,\mathcal O_X)$ est un $A$-module sans torsion, et par suite plat puisque $A$
est un anneau de valuation discrète $(\mathbf 0_{\mathrm I}, 6.3.4)$. Il est clair que
$f_\eta:X_\eta\to Y_\eta$ est un isomorphisme, bien que
$f_s:X_s\to Y_s$ n'en soit pas un.

(ii) La conclusion de (21.12.15) tombe aussi en défaut si, dans les hypothèses, on
supprime l'hypothèse que $f$ est \emph{propre}. Il suffit pour le voir de définir $S$ et
$Y$ comme dans (i), mais de remplacer $X$ par le complémentaire $X'$ de $Z_2$ dans le
préschéma $X$ défini dans (i), $f$ par la restriction $f':X'\to Y$ de $f$; le morphisme
$g':X'\to S$, restriction de $g$, est encore plat, et cette fois $X'_s$ est
géométriquement irréductible; $X'_s$ est d'ailleurs isomorphe au complémentaire d'un
point fermé dans $\mathbb P^1_k$, donc est un schéma affine isomorphe à $\mathbb A^1_k$;
comme son image dans $Y_s$ est réduite au point fermé $z$, $f'$ n'est pas propre
$(\mathbf{III}, 4.4.2)$ et $f'_s$ n'est pas un isomorphisme, bien que $f'_\eta$ en soit
un.

(iii) Il se peut que l'énoncé de la proposition (21.12.15) reste valable quand on y
remplace le mot « isomorphisme » par « morphisme étale » (cf. (21.12.14, (v))). Il en
sera alors encore de même de (21.12.16).
\end{remarks}
\subsection{Couples parafactoriels. Anneaux locaux parafactoriels}
\EGAReaderTarget{ega4.c04.subsec.21-13}{subsection:IV.21.13}
\label{IV.21.13}

\begin{definition}[21.13.1]
\EGAReaderTarget{ega4.c04.def.21-13-1}{IV.21.13.1}
Soient $X$ un espace annelé, $Y$ une partie fermée de $X$, $U=X-Y$.
On dit que le couple $(X,Y)$ est \emph{parafactoriel} si, pour tout ouvert $V$ de $X$,
le foncteur restriction
$\mathscr L\mapsto\mathscr L|(U\cap V)$, de la catégorie des
$\mathcal O_V$-Modules inversibles dans celle des $\mathcal O_{U\cap V}$-Modules
inversibles, est une \emph{équivalence de catégories}.
\end{definition}

Dire que le couple $(X,Y)$ est parafactoriel signifie donc que, pour tout ouvert $V$
de $X$, les conditions suivantes sont vérifiées :
\begin{enumerate}[label=\arabic*\textsuperscript{o}]
  \item le foncteur $\mathscr L\mapsto\mathscr L|(U\cap V)$ est \emph{pleinement
  fidèle}, autrement dit, pour deux $\mathcal O_V$-Modules inversibles $\mathscr L$,
  $\mathscr L'$, l'application de restriction
  \[
    \operatorname{Hom}_{\mathcal O_V}(\mathscr L,\mathscr L')
    \longrightarrow
    \operatorname{Hom}_{\mathcal O_{U\cap V}}
    (\mathscr L|(U\cap V),\mathscr L'|(U\cap V))
  \]
  est \emph{bijective};
  \item le foncteur $\mathscr L\mapsto\mathscr L|(U\cap V)$ est
  \emph{essentiellement surjectif}, autrement dit, pour tout
  $\mathcal O_{U\cap V}$-Module inversible $\mathscr L_0$, il existe un
  $\mathcal O_V$-Module inversible $\mathscr L$ tel que $\mathscr L_0$ soit
  \oldpage[IV-4]{314}
  isomorphe à $\mathscr L|(U\cap V)$; cela s'exprime encore en disant que
  l'homomorphisme canonique (21.3.2.4)
  \[
    \operatorname{Pic}(V)\longrightarrow\operatorname{Pic}(U\cap V)
  \]
  est \emph{surjectif}.
\end{enumerate}

\begin{lemma}[21.13.2]
\label{IV.21.13.2}
Soit $f:X'\to X$ un morphisme d'espaces annelés; pour tout ouvert $V$ de $X$, on
note $f_V:f^{-1}(V)\to V$ la restriction de $f$. Les conditions suivantes sont
équivalentes :
\begin{enumerate}[label=\emph{\alph*)}]
  \item Pour tout ouvert $V$ de $X$, le foncteur
  $\mathscr E\mapsto f_V^*(\mathscr E)$ de la catégorie des $\mathcal O_V$-Modules
  localement libres de rang fini dans la catégorie des
  $\mathcal O_{f^{-1}(V)}$-Modules localement libres de rang fini, est fidèle
  (resp.\ pleinement fidèle).
  \item[\emph{a')}] Pour tout ouvert $V$ de $X$, le foncteur
  $\mathscr L\mapsto f_V^*(\mathscr L)$ de la catégorie des $\mathcal O_V$-Modules
  localement libres de rang $1$ dans la catégorie des
  $\mathcal O_{f^{-1}(V)}$-Modules localement libres de rang $1$ est fidèle
  (resp.\ pleinement fidèle).
  \item Pour tout ouvert $V$ de $X$, l'homomorphisme canonique
  $(\mathbf 0_{\mathrm I}, 4.4.3.2)$
  $\mathscr E\to(f_V)_*(f_V^*(\mathscr E))$ est un monomorphisme
  (resp.\ un isomorphisme) pour tout $\mathcal O_V$-Module localement libre
  $\mathscr E$ de rang fini.
  \item[\emph{b')}] L'homomorphisme canonique
  $\mathcal O_X\to f_*(\mathcal O_{X'})$ est un monomorphisme
  (resp.\ un isomorphisme).
\end{enumerate}

Supposons que l'homomorphisme canonique $\mathcal O_X\to f_*(\mathcal O_{X'})$ soit
bijectif; alors, pour qu'un $\mathcal O_{X'}$-Module localement libre de rang fini
$\mathscr E'$ soit isomorphe à un $\mathcal O_{X'}$-Module de la forme
$f^*(\mathscr E)$, où $\mathscr E$ est un $\mathcal O_X$-Module localement libre de
rang fini, il faut et il suffit que les deux conditions suivantes soient remplies :
\begin{enumerate}[label=(\roman*)]
  \item $f_*(\mathscr E')$ est un $\mathcal O_X$-Module localement libre;
  \item l'homomorphisme canonique $(\mathbf 0_{\mathrm I}, 4.4.3.3)$
  $f^*(f_*(\mathscr E'))\to\mathscr E'$ est un isomorphisme.
\end{enumerate}
Lorsque ces deux conditions sont remplies, $\mathscr E$ est isomorphe à
$f_*(\mathscr E')$.
\end{lemma}

\begin{proof}
Dire que le foncteur $\mathscr E\mapsto f_V^*(\mathscr E)$ est fidèle
(resp.\ pleinement fidèle) signifie que pour deux $\mathcal O_V$-Modules localement
libres de rang fini $\mathscr E_1$, $\mathscr E_2$, l'application
\[
  \operatorname{Hom}_{\mathcal O_V}(\mathscr E_1,\mathscr E_2)
  \longrightarrow
  \operatorname{Hom}_{\mathcal O_{f^{-1}(V)}}
  (f_V^*(\mathscr E_1),f_V^*(\mathscr E_2))
\]
est injective (resp.\ bijective); comme ceci doit avoir lieu en remplaçant $V$ par un
ouvert $W\subset V$ et $\mathscr E_1$, $\mathscr E_2$ par $\mathscr E_1|W$,
$\mathscr E_2|W$, et que
\[
  \operatorname{Hom}_{\mathcal O_W}(\mathscr E_1|W,\mathscr E_2|W)
  =\Gamma(W,\mathscr Hom_{\mathcal O_V}(\mathscr E_1,\mathscr E_2)),
\]
la condition signifie encore que l'homomorphisme canonique de faisceaux
\[
\label{IV.21.13.2.1}
  \mathscr Hom_{\mathcal O_V}(\mathscr E_1,\mathscr E_2)
  \longrightarrow
  (f_V)_*(f_V^*(\mathscr Hom_{\mathcal O_V}(\mathscr E_1,\mathscr E_2)))
  \tag{21.13.2.1}
\]
est injectif (resp.\ bijectif). Mais puisque $\mathscr E_1$ et $\mathscr E_2$ sont
localement libres de rang fini,
$\mathscr Hom_{\mathcal O_V}(\mathscr E_1,\mathscr E_2)$, isomorphe à
$\mathscr E_1^\vee\otimes_{\mathcal O_V}\mathscr E_2$
$(\mathbf 0_{\mathrm I}, 5.4.2)$, est aussi localement libre de rang fini; cela montre
déjà que $b)$ entraîne $a)$, et inversement $b)$ est un cas particulier de $a)$,
compte tenu de l'isomorphisme
$\mathscr E\xrightarrow{\sim}\mathscr Hom_{\mathcal O_V}(\mathcal O_V,\mathscr E)$.
Il est clair que $a')$ est un cas particulier de $a)$ et $b')$ un cas particulier de
$a')$. Inversement, comme $b)$ est une propriété locale, on peut, pour la vérifier,
se borner au cas où $\mathscr E=\mathcal O_V^n$, et cela montre que $b')$ entraîne
$b)$.

Si l'homomorphisme canonique $\mathcal O_X\to f_*(\mathcal O_{X'})$ est bijectif,
et si les conditions (i) et (ii) sont remplies, il est clair que
$\mathscr E'=f^*(\mathscr E)$ avec $\mathscr E=f_*(\mathscr E')$, à un isomorphisme
près.
\oldpage[IV-4]{315}
Inversement, supposons que $\mathscr E'=f^*(\mathscr E)$, avec $\mathscr E$ localement
libre; la question étant locale sur $X$, on peut supposer que
$\mathscr E=\mathcal O_X^n$, d'où $\mathscr E'=\mathcal O_{X'}^n$, et les conditions
(i) et (ii) résultent de l'hypothèse que
$\mathcal O_X\to f_*(\mathcal O_{X'})$ est un isomorphisme.
\end{proof}

Dans ce numéro, nous appliquerons le lemme précédent à une injection canonique
$j:U\to X$, où $U$ est l'espace annelé induit par $X$ sur un \emph{ouvert} de $X$.
Avec ces notations, (21.13.2) se traduit en :

\begin{corollary}[21.13.3]
\label{IV.21.13.3}
Soient $X$ un espace annelé, $Y$ une partie fermée de $X$, $U=X-Y$,
$j:U\to X$ l'injection canonique. Les conditions suivantes sont équivalentes :
\begin{enumerate}[label=\emph{\alph*)}]
  \item Pour tout ouvert $V$ de $X$, le foncteur restriction
  $\mathscr E\mapsto\mathscr E|(U\cap V)$ de la catégorie des
  $\mathcal O_V$-Modules localement libres de rang fini dans la catégorie des
  $\mathcal O_{U\cap V}$-Modules localement libres de rang fini est fidèle
  (resp.\ pleinement fidèle).
  \item[\emph{a')}] Pour tout ouvert $V$ de $X$, le foncteur restriction
  $\mathscr L\mapsto\mathscr L|(U\cap V)$ de la catégorie des
  $\mathcal O_V$-Modules localement libres de rang $1$ dans la catégorie des
  $\mathcal O_{U\cap V}$-Modules localement libres de rang $1$ est fidèle
  (resp.\ pleinement fidèle).
  \item Pour tout ouvert $V$ de $X$, l'homomorphisme canonique
  $\mathscr E\to(j_V)_*(\mathscr E|(U\cap V))$ est injectif (resp.\ bijectif) pour
  tout $\mathcal O_V$-Module localement libre de rang fini $\mathscr E$.
  \item[\emph{b')}] L'homomorphisme canonique
  $\mathcal O_X\to j_*(\mathcal O_U)$ est injectif (resp.\ bijectif).
\end{enumerate}

Supposons que l'homomorphisme canonique $\mathcal O_X\to j_*(\mathcal O_U)$ soit
bijectif. Alors, pour qu'un $\mathcal O_U$-Module localement libre de rang fini
$\mathscr F$ soit de la forme $\mathscr E|U$, où $\mathscr E$ est un
$\mathcal O_X$-Module localement libre de rang fini, il faut et il suffit que
$j_*(\mathscr F)$ soit un $\mathcal O_X$-Module localement libre, et dans ce cas, on
peut prendre $\mathscr E=j_*(\mathscr F)$.
\end{corollary}

\begin{lemma}[21.13.4]
\label{IV.21.13.4}
Soient $X$ un préschéma localement noethérien, $Y$ une partie fermée de $X$,
$U=X-Y$, $j:U\to X$ l'injection canonique. Pour que l'homomorphisme canonique
$\mathcal O_X\to j_*(\mathcal O_U)$ soit injectif (resp.\ bijectif), il faut et il
suffit que $\operatorname{prof}_Y(\mathcal O_X)\geq1$
(resp.\ $\operatorname{prof}_Y(\mathcal O_X)\geq2$) (autrement dit, que
$\operatorname{prof}(\mathcal O_{X,x})\geq1$
(resp.\ $\operatorname{prof}(\mathcal O_{X,x})\geq2$) pour tout $x\in Y$).
\end{lemma}

\begin{proof}
Ces assertions sont en effet des cas particuliers de (5.10.2) et (5.10.5).
\end{proof}

\begin{proposition}[21.13.5]
\label{IV.21.13.5}
Soient $X$ un espace annelé, $Y$ une partie fermée de $X$, $U=X-Y$,
$j:U\to X$ l'injection canonique. Pour que le couple $(X,Y)$ soit parafactoriel, il
faut et il suffit que l'homomorphisme canonique $\mathcal O_X\to j_*(\mathcal O_U)$
soit bijectif, et que pour tout ouvert $V$ de $X$ et tout
$\mathcal O_{U\cap V}$-Module inversible $\mathscr L$, $(j_V)_*(\mathscr L)$ soit un
$\mathcal O_V$-Module inversible (notation de (21.13.3)).
\end{proposition}

\begin{proof}
C'est une conséquence immédiate de la définition (21.13.1) et de (21.13.3).
\end{proof}

\begin{corollary}[21.13.6]
\label{IV.21.13.6}
Soient $X$ un espace annelé, $Y$ une partie fermée de $X$.
\begin{enumerate}[label=(\roman*)]
  \item Si le couple $(X,Y)$ est parafactoriel, il en est de même de
  $(W,Y\cap W)$ pour tout ouvert $W$ de $X$. Inversement, si $(W_\alpha)$ est un
  recouvrement ouvert de $X$ tel que chacun des couples
  $(W_\alpha,Y\cap W_\alpha)$ soit parafactoriel, alors le couple $(X,Y)$ est
  parafactoriel.
  \item Si le couple $(X,Y)$ est parafactoriel, il en est de même de $(X,Y')$ pour
  toute partie fermée $Y'$ de $Y$.
  \item Supposons que $X$ soit un préschéma, et soient $X'$ un préschéma,
  $f:X'\to X$ un morphisme fidèlement plat et quasi-compact; posons
  $Y'=f^{-1}(Y)$. Supposons que le couple $(X',Y')$ soit parafactoriel et que
  l'ouvert $U=X-Y$ soit rétrocompact dans $X$ $(\mathbf 0_{\mathrm{III}}, 9.1.1)$.
  Alors le couple $(X,Y)$ est parafactoriel.
\end{enumerate}
\end{corollary}

\begin{proof}
\oldpage[IV-4]{316}
(i) Comme le fait que l'homomorphisme canonique
$\mathcal O_X\to j_*(\mathcal O_U)$ soit bijectif est une propriété locale sur $X$,
tout revient à voir (en vertu de (21.13.5)), en désignant par
$j_\alpha:U\cap W_\alpha\to W_\alpha$ l'injection canonique, que l'on a la propriété
suivante : si, pour tout $\alpha$ et pour tout $\mathcal O_U$-Module inversible
$\mathscr L$, $(j_\alpha)_*(\mathscr L|(U\cap W_\alpha))$ est un
$\mathcal O_{W_\alpha}$-Module inversible, alors $j_*(\mathscr L)$ est un
$\mathcal O_X$-Module inversible. Mais la propriété d'être un $\mathcal O_X$-Module
inversible est locale sur $X$, et $(W_\alpha)$ est un recouvrement ouvert de $X$;
comme on a
\[
  j_*(\mathscr L)|W_\alpha
  =(j_\alpha)_*(\mathscr L|(U\cap W_\alpha)),
\]
cela prouve notre assertion.

(ii) Posons $U'=X-Y'$, et soit $j':U'\to X$ l'injection canonique. Dire que
l'homomorphisme $\mathcal O_X\to j'_*(\mathcal O_{U'})$ est bijectif revient à dire
que pour tout ouvert $V$ de $X$, l'homomorphisme
$\Gamma(V,\mathcal O_X)\to\Gamma(V\cap U',\mathcal O_X)$ est bijectif; mais
l'homomorphisme composé
\[
  \Gamma(V,\mathcal O_X)\longrightarrow\Gamma(V\cap U',\mathcal O_X)
  \longrightarrow\Gamma(V\cap U,\mathcal O_X)
\]
est bijectif par hypothèse (21.13.5), et en remplaçant $V$ par $V\cap U'$, on voit que
$\Gamma(V\cap U',\mathcal O_X)\to\Gamma(V\cap U,\mathcal O_X)$ est bijectif, donc
$\Gamma(V,\mathcal O_X)\to\Gamma(V\cap U',\mathcal O_X)$ est bijectif. Notons ensuite
que si $j'':U\to U'$ est l'injection canonique et si $\mathscr L'$ est un
$\mathcal O_{U'}$-Module inversible, $\mathscr L'|U$ est un $\mathcal O_U$-Module
inversible, donc
$j_*(\mathscr L'|U)=j'_*(j''_*(\mathscr L'|U))$ est par hypothèse un
$\mathcal O_X$-Module inversible. Comme le couple $(U',U'\cap Y)$ est parafactoriel
en vertu de (i), on a $j''_*(\mathscr L'|U)\cong\mathscr L'$, donc
$j'_*(\mathscr L')$ est un $\mathcal O_X$-Module inversible. Il suffit alors pour
conclure de remplacer dans le raisonnement précédent $X$, $U$ et $U'$ par $V$,
$V\cap U$ et $V\cap U'$, où $V$ est un ouvert de $X$.

(iii) Posons $U'=X'-Y'=f^{-1}(U)$ et notons que puisqu'il s'agit de préschémas, on
peut écrire $U'=U\times_X X'$; soit $f_U:f^{-1}(U)\to U$ la restriction de $f$ et
soit $j':U'\to X'$ l'injection canonique, qui s'écrit aussi $j_{(X')}$. Montrons
d'abord que l'homomorphisme canonique
$\rho:\mathcal O_X\to j_*(\mathcal O_U)$ est bijectif; pour cela, notons que par
hypothèse l'homomorphisme canonique
$\mathcal O_{X'}\to j'_*(\mathcal O_{U'})$ est bijectif; mais puisque le morphisme
$j$ est quasi-compact et séparé par hypothèse, et le morphisme $f$ plat,
$j'_*(\mathcal O_{U'})=j'_*(f_U^*(\mathcal O_U))$ s'identifie canoniquement à
$f^*(j_*(\mathcal O_U))$ en vertu de (2.3.1); comme
$\mathcal O_{X'}=f^*(\mathcal O_X)$, on voit que l'homomorphisme
$f^*(\rho):f^*(\mathcal O_X)\to f^*(j_*(\mathcal O_U))$ est bijectif; on en conclut
que $\rho$ est bijectif puisque $f$ est fidèlement plat (2.2.7). Considérons ensuite
un $\mathcal O_U$-Module inversible $\mathscr L$, et soit
$\mathscr L'=f_U^*(\mathscr L)$, qui est un $\mathcal O_{U'}$-Module inversible. Par
hypothèse $j'_*(f_U^*(\mathscr L))$ est un $\mathcal O_{X'}$-Module inversible, qui
est isomorphe à $f^*(j_*(\mathscr L))$ par (2.3.1) comme ci-dessus. Mais puisque $f$
est fidèlement plat et quasi-compact, cela entraîne que $j_*(\mathscr L)$ est un
$\mathcal O_X$-Module localement libre (2.5.2). Pour achever la démonstration, il
suffit de remplacer $X$ par un ouvert $V$ et $X'$ par $f^{-1}(V)$ dans ce qui
précède.
\end{proof}

\begin{definition}[21.13.7]
\EGAReaderTarget{ega4.c04.def.21-13-7}{IV.21.13.7}
On dit qu'un anneau local $A$ est \emph{parafactoriel} si, posant
$X=\operatorname{Spec}(A)$ et désignant par $a$ le point fermé de $X$, le couple
$(X,\{a\})$ est parafactoriel.
\end{definition}

\begin{proposition}[21.13.8]
\label{IV.21.13.8}
Les notations étant celles de (21.13.7), on pose $U=X-\{a\}$. Pour que $A$ soit
parafactoriel, il faut et il suffit qu'il vérifie les conditions suivantes :
\begin{enumerate}[label=(\roman*)]
  \item L'homomorphisme canonique
  $A=\Gamma(X,\mathcal O_X)\to\Gamma(U,\mathcal O_X)$ est bijectif.
  \item $\operatorname{Pic}(U)=0$.
\end{enumerate}
\oldpage[IV-4]{317}
Si en outre $A$ est noethérien, la condition (i) équivaut à
\begin{enumerate}
  \item[\textup{(i bis)}] $\operatorname{prof}(A)\geq2$.
\end{enumerate}
\end{proposition}

\begin{proof}
En effet, le seul ouvert de $X$ contenant $a$ est $X$ lui-même, et par suite tout
$\mathcal O_X$-Module inversible est isomorphe à $\mathcal O_X$, autrement dit
$\operatorname{Pic}(X)=0$; la première assertion résulte donc de la définition
(21.13.1) et de (21.13.3). La seconde assertion n'est autre qu'un cas particulier de
(21.13.4).
\end{proof}

\begin{env}[21.13.9]
\label{IV.21.13.9}
\emph{Exemples}. ---

(i) Un anneau local noethérien parafactoriel est nécessairement de dimension
$\geq2$ en vertu de (21.13.8); en d'autres termes un anneau local noethérien de
dimension $\leq1$ n'est pas parafactoriel.

(ii) Un anneau local noethérien \emph{factoriel} $A$ de dimension $\geq2$ est
parafactoriel, comme il résulte de (21.13.8) et de (21.6.14).

(iii) Si $A$ est un anneau local noethérien de dimension $\geq3$ et parafactoriel, il
n'est pas nécessairement normal ni même réduit. Considérons un anneau local régulier
$B$ de dimension $\geq3$, et soit $A=D_B(B)$ $(\mathbf 0, 18.2.3)$, isomorphe à
$B[T]/(T^2)$; on voit aussitôt que l'on a
$\operatorname{prof}(A)=\operatorname{prof}(B)=\dim(B)\geq3$. Pour montrer que $A$
est parafactoriel, il suffit donc, avec les notations de (21.13.8), de prouver que
$\operatorname{Pic}(U)=0$. Soit $\mathfrak J$ le noyau de l'augmentation $A\to B$,
qui est tel que $\mathfrak J^2=0$ et qui, en tant que $B$-module, est isomorphe à $B$.
Comme $B=A/\mathfrak J$, $X=\operatorname{Spec}(A)$ et
$X_0=\operatorname{Spec}(B)$ ont même espace sous-jacent; si
$\mathscr J=\widetilde{\mathfrak J}$, $X_0$ est le sous-schéma défini par
$\mathscr J$; nous noterons $U_0$ le sous-préschéma induit par $X_0$ sur l'ouvert
$U$. Pour tout $z\in\mathfrak J$, posons $\varphi(z)=1+z$; puisque
$(1+z)(1-z)=1$, $1+z$ est inversible dans $A$ et $\varphi(\mathfrak J)$ est le noyau
de l'homomorphisme surjectif canonique de groupes multiplicatifs $A^*\to B^*$; en
d'autres termes, on a une suite exacte de groupes commutatifs
\[
  0\longrightarrow\mathfrak J\xrightarrow{\varphi}A^*\longrightarrow B^*
  \longrightarrow1
\]
(les trois derniers groupes étant multiplicatifs, les deux premiers additifs). Pour la
même raison, pour tout $t\in A$, on a la suite exacte
\[
  0\longrightarrow\mathfrak J_t\xrightarrow{\varphi_t}(A_t)^*
  \longrightarrow(B_{t_0})^*\longrightarrow1
\]
en désignant par $t_0$ l'image de $t$ dans $B$, puisque
$B_{t_0}=A_t/\mathfrak J_t$; autrement dit, on a une suite exacte de faisceaux de
groupes commutatifs sur l'espace topologique $X$
\[
  0\longrightarrow\mathscr J\xrightarrow{u}\mathcal O_X^*
  \longrightarrow\mathcal O_{X_0}^*\longrightarrow1
\]
d'où par restriction à l'ouvert $U$, une suite exacte
\[
  0\longrightarrow\mathscr J|U\longrightarrow\mathcal O_U^*
  \longrightarrow\mathcal O_{U_0}^*\longrightarrow1.
\]
Par la suite exacte de cohomologie, on en déduit une suite exacte
\[
\label{IV.21.13.9.1}
  \mathrm H^1(U,\mathscr J)\longrightarrow\mathrm H^1(U,\mathcal O_U^*)
  \xrightarrow{u_1}\mathrm H^1(U,\mathcal O_{U_0}^*)
  \tag{21.13.9.1}
\]
Mais puisque $\mathfrak J$ est un $B$-module isomorphe à $B$, on a
$\mathrm H^1(U,\mathscr J)=\mathrm H^1(U_0,\mathcal O_{U_0})$, et il résulte du
chap.~III, 3\textsuperscript{e} partie (voir aussi [41, III, Exemple III-1]) que l'on a
\oldpage[IV-4]{318}
$\mathrm H^1(U_0,\mathcal O_{U_0})=0$ en raison de la relation
$\operatorname{prof}(B)\geq3$. Par ailleurs, comme $B$ est factoriel, on a
$\mathrm H^1(U,\mathcal O_{U_0}^*)=\operatorname{Pic}(U_0)=0$; on tire donc bien de
la suite exacte précédente que
$\operatorname{Pic}(U)=\mathrm H^1(U,\mathcal O_U^*)=0$.

(iv) Il existe aussi des anneaux locaux noethériens de dimension $3$ qui sont
intègres, intégralement clos et parafactoriels, mais non factoriels. Soit en effet $B$
un anneau local noethérien, complet, intègre et intégralement clos, de dimension
$\geq2$ et non factoriel (par exemple le complété de l'anneau localisé de l'algèbre
intègre $k[U,V,W]/(W^2-UV)$ en l'idéal image de $(U)+(V)+(W)$). Alors il résulte de
ce que l'on verra plus bas (21.14.2) que l'anneau local $A=B[[T]]$ des séries formelles
en $T$ sur $B$ est parafactoriel, mais il n'est pas factoriel, sans quoi $B$ le serait
en vertu de (21.13.12) ci-dessous.

(v) On peut montrer qu'un anneau local d'intersection complète absolue (19.3.1) de
dimension $\geq4$, est parafactoriel (cf. [41, XI, 3.13 (i)]).

(vi) On a vu (Remarque (ii)) que tout anneau local noethérien factoriel $A$ de
dimension $2$ est parafactoriel. Mais il y a des anneaux locaux noethériens de
dimension $2$ qui sont parafactoriels et non factoriels. On peut montrer que, pour
qu'un anneau local noethérien $A$ de dimension $2$ soit parafactoriel et non
factoriel, il faut et il suffit qu'il satisfasse aux trois conditions suivantes :
\begin{enumerate}[label=\arabic*\textsuperscript{o}]
  \item $A$ est un anneau de Cohen--Macaulay (autrement dit
  $\operatorname{prof}(A)=2$).
  \item $A$ est intègre et si $A'$ est sa clôture intégrale, $A'$ est factoriel et est
  une $A$-algèbre finie.
  \item Soit $\mathfrak J$ le conducteur de $A$ dans $A'$ (annulateur du $A$-module
  $A'/A$, ou encore le plus grand idéal de $A'$ contenu dans $A$); posons
  $B=A/\mathfrak J$, $B'=A'/\mathfrak J$; alors $\dim(B)=1$ (ce qui implique
  $A'\neq A$, autrement dit $A$ n'est pas intégralement clos), et l'application
  canonique $\operatorname{Div}(B)\to\operatorname{Div}(B')$ (21.4.5) est
  \emph{surjective}.
\end{enumerate}
On peut montrer en outre que ces conditions entraînent la propriété suivante :

$4^{\mathrm o}$ L'anneau $B'$ (et \emph{a fortiori} $B\subset B'$) est réduit, et le
morphisme $\operatorname{Spec}(B')\to\operatorname{Spec}(B)$ est bijectif.

Si l'on pose $X=\operatorname{Spec}(A)$, $X'=\operatorname{Spec}(A')$,
$Y=\operatorname{Spec}(B)$, $Y'=\operatorname{Spec}(B')$, de sorte que $Y$
(resp.\ $Y'$) est défini par l'idéal $\mathfrak J$ de $A$ (resp.\ $A'$), le morphisme
structural $f:X'\to X$ est un isomorphisme de $X'-Y'$ sur $X-Y$ (Bourbaki,
\emph{Alg. comm.}, chap.~V, \S1, n\textsuperscript{o}~5, cor.~5 de la prop.~16). On
voit donc en vertu de $4^{\mathrm o}$ que $f$ est un morphisme \emph{bijectif}; en
d'autres termes, $X$ est un préschéma \emph{unibranche} (6.15.1) (et en particulier
$A'$ est un anneau local noethérien); en général $X$ n'est pas géométriquement
unibranche. L'espace $Y$, de dimension $1$, est constitué par le point fermé $x$ de
$X$ et les points maximaux $y_i$ ($1\leq i\leq r$) de $Y$, et l'unique point $y_i'$
de $Y'$ au-dessus de $y_i$ est aussi un point maximal de $Y'$; comme $B$ et $B'$ sont
réduits, on en déduit que l'on a
\[
  \mathfrak m_{y_i}\mathcal O_{X',y_i'}=\mathfrak m_{y_i'}.
\]
Pour $z=y_i$ et $z'=y_i'$, cette relation est aussi vérifiée de façon évidente pour
$z\notin Y$ et $z'$ l'unique point de $f^{-1}(z)$, donc pour tout
$z\in U=X-\{x\}$ et $z'$ l'unique point de $f^{-1}(z)$. En particulier, si l'anneau
$A$ est de caractéristique $0$ (0, 21.1.1) on voit que $U'$ est \emph{non ramifié} sur
$U$ (mais non étale en général).

Nous nous bornerons ici à démontrer que les conditions $1^{\mathrm o}$,
$2^{\mathrm o}$, $3^{\mathrm o}$ ci-dessus sont \emph{suffisantes} pour que $A$ soit
parafactoriel. Or, en vertu de la condition $1^{\mathrm o}$ et de (21.13.8), il suffit
de démontrer que $\operatorname{Pic}(U)=0$. Puisque $A'$ est factoriel en vertu de
$2^{\mathrm o}$, on a $\operatorname{Pic}(U')=0$, et on déduit de (21.8.5, (ii)) une
suite exacte
\[
  1\longrightarrow
  \left(\prod_{i=1}^r
  \mathcal O_{U',y_i'}^*/\mathcal O_{U,y_i}^*\right)
  \big/\operatorname{Im}\Gamma(U',\mathcal O_{U'}^*)
  \longrightarrow\operatorname{Pic}(U)\longrightarrow\operatorname{Pic}(U')
  \longrightarrow0
\]
et il s'agit donc de montrer que le second terme de cette suite est réduit à l'élément
neutre, en prenant pour $S$ l'ensemble des $y_i$ ($1\leq i\leq r$) et en notant
qu'ici $f_*(\mathcal O_{X'})$ est la $\mathcal O_X$-Algèbre $\widetilde{A'}$. Soient
$\mathfrak p_i$, $\mathfrak p_i'$ les idéaux premiers de $A$ et $A'$ correspondant
aux points $y_i$, $y_i'$, et remarquons que $\mathfrak p_i'\cap A=\mathfrak p_i$;
donnons-nous pour chaque $i$ un élément inversible $a_i'/s_i$ de
$A'_{\mathfrak p_i'}$, avec $a_i'\in A'$, $s_i\notin\mathfrak p_i'$; comme nous
n'avons à examiner que les groupes quotients
$(A'_{\mathfrak p_i'})^*/(A_{\mathfrak p_i})^*$, on peut supposer $s_i=1$ pour tout
$i$; soit $b_i'\in B'=A'/\mathfrak J$ l'image canonique de $a_i'$, de sorte que
$b_i'/1$ est un élément non nul de $k(y_i')=\mathcal O_{Y',y_i'}$. L'ensemble
$U\cap Y$ formé par les $y_i$ est un ouvert affine de la forme $D(t)$ dans $Y$, avec
$t\in\mathfrak m$, idéal maximal de $A$; il existe donc un élément inversible
$b'\in B'_t$ dont les $b_i'/1$ sont les images canoniques, et comme $t$ est régulier
dans l'anneau intègre $A'$, on peut supposer $b'$ de la forme $b''/t^m$, où
$b''\in B'$ est inversible. Soit $a''$ un élément de $A'$ dans la classe $b''$, qui
est donc nécessairement inversible; $a'=a''/t^m$ est inversible dans $A'_t$ et pour
tout $i$, on a
\[
  u_i=a''-a_i't^m\in\mathfrak J,
  \qquad
  t^m a_i'=a''-u_i=a''(1-a''^{-1}u_i);
\]
mais $a''^{-1}u_i\in\mathfrak J\subset\mathfrak m$, donc $1-a''^{-1}u_i$ est un
élément inversible de $A$, et les classes de $a'$ et $a_i'$ dans
$(A'_{\mathfrak p_i'})^*/(A_{\mathfrak p_i})^*$ sont les mêmes, ce qui achève la
démonstration.

\oldpage[IV-4]{319}
Pour avoir un exemple explicite d'anneau parafactoriel de dimension $2$ obtenu de
cette manière et \emph{non factoriel}, considérons l'anneau
\[
  E=\mathbb R[[U,V]]/(U^2+V^2),
\]
dont la clôture intégrale $E'$ s'identifie à $\mathbb C[[U]]$ (6.15.11). Si l'on pose
$A=E[[T]]$, la clôture intégrale de $A$ est l'anneau $A'=E'[[T]]$ (Bourbaki,
\emph{Alg. comm.}, chap.~V, \S1, n\textsuperscript{o}~4, prop.~14); on vérifie
aussitôt que le conducteur de $E$ dans $E'$ est l'idéal maximal $\mathfrak n$ de
$E$, donc le conducteur $\mathfrak J$ de $A$ dans $A'$ est $\mathfrak n[[T]]$, et on
a $B=A/\mathfrak J\cong\mathbb R[[T]]$, $B'=A'/\mathfrak J\cong\mathbb C[[T]]$.
Il est alors immédiat que les conditions $1^{\mathrm o}$, $2^{\mathrm o}$ et
$3^{\mathrm o}$ énoncées ci-dessus sont bien vérifiées, mais $A$ n'est même pas
intégralement clos.

On peut varier cet exemple, et le lecteur verra sans peine que si $k$ est un corps
algébriquement clos, l'anneau $A$, localisé de l'anneau
$k[U,V,W]/(U^2-WV^2)$ en l'idéal maximal engendré par les images de $U$, $V$ et $W$,
vérifie aussi les conditions $1^{\mathrm o}$, $2^{\mathrm o}$ et $3^{\mathrm o}$
ci-dessus.

Il se pourrait que ces trois conditions impliquent même que chacun des anneaux
$B'/\mathfrak p_i'$ soit un anneau de valuation discrète, ce qui entraînerait que
$A'$ est même un anneau \emph{régulier}.
\end{env}

\begin{proposition}[21.13.10]
\label{IV.21.13.10}
Soient $X$ un préschéma localement noethérien, $Y$ une partie fermée de $X$. Pour que
le couple $(X,Y)$ soit parafactoriel il faut et il suffit que, pour tout $y\in Y$,
l'anneau local $\mathcal O_{X,y}$ soit parafactoriel.
\end{proposition}

\begin{proof}
Posons $U=X-Y$ et soit $j:U\to X$ l'injection canonique.

En vertu de (21.13.4), dire que l'homomorphisme canonique
$\mathcal O_X\to j_*(\mathcal O_U)$ soit bijectif équivaut à dire que chacun des
anneaux locaux $\mathcal O_{X,y}$ pour $y\in Y$ satisfait à la condition (i bis) de
(21.13.8).

Montrons d'abord que si le couple $(X,Y)$ est parafactoriel, chacun des anneaux
$\mathcal O_{X,y}$ ($y\in Y$) est parafactoriel. Compte tenu de la remarque
précédente, tout revient à voir que, si l'on pose
$T_y=\operatorname{Spec}(\mathcal O_{X,y})$ et $U_y=T_y-\{y\}$, tout
$\mathcal O_{U_y}$-Module inversible $\mathscr L_0$ est isomorphe à
$\mathcal O_{U_y}$. Or, lorsque $V$ parcourt l'ensemble des voisinages ouverts
affines de $y$ dans $X$, il résulte de (8.2.13) et $(\mathbf I, 2.4.2)$ que le
préschéma $U_y$ est limite projective des préschémas induits par $X$ sur les ouverts
$V-(V\cap\overline{\{y\}})$, qui sont quasi-compacts puisque $X$ est localement
noethérien. Comme les $V-(V\cap\overline{\{y\}})$ sont séparés, il résulte alors de
(8.5.2, (ii)) et (8.5.5) qu'il y a un voisinage ouvert affine $V$ de $y$ dans $X$ et
un $\mathcal O_{U_V}$-Module inversible $\mathscr L'_V$ tel que $\mathscr L_0$ soit le
faisceau induit par $\mathscr L'_V$ sur $U_y$. Mais l'hypothèse entraîne, en vertu de
(21.13.6, (i) et (iii)), que le couple $(V,V\cap\overline{\{y\}})$ est parafactoriel;
on en conclut par définition (21.13.1) qu'il existe un $\mathcal O_V$-Module inversible
$\mathscr L_V$ induisant $\mathscr L'_V$ sur $U_V$. Remplaçant au besoin $V$ par un
voisinage ouvert plus petit de $y$, on peut supposer que $\mathscr L_V$ est isomorphe
à $\mathcal O_V$, d'où on conclut que $\mathscr L_0$ est isomorphe à
$\mathcal O_{U_y}$.

Inversement, prouvons que si tous les $\mathcal O_{X,y}$ ($y\in Y$) sont
parafactoriels le couple $(X,Y)$ est parafactoriel. On est évidemment ramené, compte
tenu de la remarque du début, à montrer que pour tout $\mathcal O_U$-Module inversible
$\mathscr L$, $j_*(\mathscr L)$ est un $\mathcal O_X$-Module inversible (21.13.5).
La question étant locale sur $X$, on peut supposer $X$ noethérien. L'ensemble $V$ des
$x\in X$ tels que la restriction de $j_*(\mathscr L)$ à un voisinage ouvert de $x$
dans $X$ soit inversible est évidemment ouvert dans $X$; il s'agit de montrer que
$V=X$. Pour cela, supposons le contraire et posons $Z=X-V\neq\varnothing$. Soit $y$
un point \emph{maximal} de $Z$; puisque $Z\subset Y$ par définition,
$\mathcal O_{X,y}$ est par hypothèse un anneau parafactoriel. Remplaçant au besoin
$X$ par un voisinage ouvert de $y$ dans $X$, on peut supposer que la restriction de
$j_*(\mathscr L)$
\oldpage[IV-4]{320}
à $V-(V\cap\overline{\{y\}})$ est inversible; donc, avec les notations introduites
ci-dessus, le $\mathcal O_{U_y}$-Module $\mathscr L_0$ induit par $j_*(\mathscr L)$
sur $U_y$ est inversible. Puisque $\mathcal O_{X,y}$ est parafactoriel,
$\mathscr L_0$ est induit par un $\mathcal O_{T_y}$-Module inversible $\mathscr L'$;
mais comme $T_y$ est limite projective des voisinages ouverts $W$ de $y$ dans $X$,
on peut de nouveau appliquer (8.5.2, (ii)) et (8.5.5), prouvant l'existence d'un tel
voisinage $W$ et d'un $\mathcal O_W$-Module inversible $\mathscr L''$ induisant
$\mathscr L'$ sur $T_y$, donc induisant $\mathscr L_0$ sur $U_y$; appliquant cette
fois (8.5.2.5) et (8.5.2, (i)), on en déduit qu'en remplaçant au besoin $W$ par un
voisinage ouvert plus petit de $y$, on peut supposer que les restrictions de
$j_*(\mathscr L)$ et de $\mathscr L''$ à
$W-(W\cap\overline{\{y\}})$ soient égales. Si l'on pose alors $V_1=V\cup W$ et si
l'on désigne par $\mathscr L_1$ le $\mathcal O_{V_1}$-Module inversible égal à
$\mathscr L''$ dans $W$, à $j_*(\mathscr L)|V$ dans $V$, on a $U\subset V_1$ et
$\mathscr L_1|U=\mathscr L$. On conclut alors du fait que l'homomorphisme
$\mathcal O_X\to j_*(\mathcal O_U)$ est bijectif et de (21.13.2, $b$)) que
$j_*(\mathscr L)|V_1$ est isomorphe à $\mathscr L_1$, donc inversible. Mais cela
contredit la définition de $V$, et conclut donc la démonstration de (21.13.10).
\end{proof}

\begin{corollary}[21.13.11]
\label{IV.21.13.11}
Soient $X$ un préschéma localement noethérien, $Y$ une partie fermée de $X$,
$U=X-Y$. Supposons que le couple $(X,Y)$ soit parafactoriel et que $U$ soit
localement factoriel; en d'autres termes (21.13.10), pour tout $x\in X$, l'anneau
$\mathcal O_{X,x}$ est : $1^{\mathrm o}$ parafactoriel si $x\in Y$;
$2^{\mathrm o}$ factoriel si $x\notin Y$. Alors $X$ est localement factoriel
(autrement dit, $\mathcal O_{X,x}$ est en fait factoriel pour tout $x\in X$).
\end{corollary}

\begin{proof}
Supposons en effet le contraire, et soit $y$ un point de $Y$ tel que
$\mathcal O_{X,y}$ ne soit pas factoriel et ait une dimension \emph{minima} parmi
tous les points de $Y$ ayant cette propriété. Alors, si l'on pose
$T_y=\operatorname{Spec}(\mathcal O_{X,y})$, $U_y=T_y-\{y\}$, l'hypothèse que
$\mathcal O_{X,y}$ est parafactoriel entraîne $\operatorname{Pic}(U_y)=0$ et
$\operatorname{prof}(\mathcal O_{X,y})\geq2$ (21.13.8); de plus le choix de $y$
entraîne que $U_y$ est localement factoriel. Mais alors (21.6.14) prouve que
$\mathcal O_{X,y}$ est factoriel, contrairement à l'hypothèse.
\end{proof}

\begin{proposition}[21.13.12]
\label{IV.21.13.12}
Soient $A$, $B$ deux anneaux locaux noethériens,
$\rho:A\to B$ un homomorphisme local faisant de $B$ un $A$-module plat. Alors, si
$B$ est factoriel, il en est de même de $A$.
\end{proposition}

\begin{proof}
On sait déjà que $A$ est intègre et intégralement clos (6.5.4), donc on peut se borner
au cas où $\dim(A)\geq2$, et raisonner par récurrence sur $n=\dim(A)$. Soient
$X=\operatorname{Spec}(A)$, $a$ le point fermé de $X$,
$X'=\operatorname{Spec}(B)$, $f:X'\to X$ le morphisme structural,
$U=X-\{a\}$, $U'=f^{-1}(U)$.

Les anneaux locaux $\mathcal O_{X',x'}$ aux points $x'\in f^{-1}(a)$ sont factoriels
et de dimension $\geq2$ (6.1.2), donc parafactoriels (21.13.9, (ii)); on en conclut
(21.13.10) que le couple $(X',f^{-1}(a))$ est parafactoriel. En vertu de (21.13.6,
(iii)), cela montre que l'anneau $A$ est parafactoriel, donc
$\operatorname{prof}(A)\geq2$ et $\operatorname{Pic}(U)=0$ (21.13.8). Enfin, les
anneaux locaux de $U$ étant de dimension $<n$, l'hypothèse de récurrence prouve que
$U$ est \emph{localement factoriel}; la conclusion résulte alors de (21.6.14).
\end{proof}

\begin{env}[21.13.12.1]
\label{IV.21.13.12.1}
On notera que l'énoncé (21.13.12) où on remplace « factoriel » par « parafactoriel »
n'est plus exact, comme le montre l'exemple où $A=k$ est un corps et $B$ l'anneau
local factoriel $k[[T_1,T_2]]$. Cependant, il résulte de (21.13.6, (iii)) que, sous
les conditions de (21.13.12), si $\mathfrak m$ désigne l'idéal maximal de $A$ et si
$\mathfrak mB$ est un \emph{idéal}
\oldpage[IV-4]{321}
\emph{de définition} de $B$ (autrement dit, si $\dim(B/\mathfrak mB)=0$), alors,
lorsque $B$ est parafactoriel, il en est de même de $A$. En particulier, si le
complété $\widehat A$ d'un anneau local noethérien est parafactoriel, il en est de
même de $A$.
\end{env}

\begin{remark}[21.13.13]
\label{IV.21.13.13}
Une partie des définitions et résultats précédents se rattache à des considérations
plus générales sur la cohomologie des faisceaux de groupes commutatifs sur un espace
topologique, développées au chap.~III, 3\textsuperscript{e} partie, et nous n'en avons
donné ici un exposé indépendant que pour la commodité du lecteur. En effet, si $X$
est un espace annelé, on a une \emph{équivalence} canonique entre la catégorie des
$\mathcal O_X$-Modules inversibles et celle des \emph{faisceaux principaux homogènes
sous le faisceau de groupes} $\mathcal O_X^*$ (16.5.15). Cette équivalence se définit
en faisant correspondre à tout $\mathcal O_X$-Module inversible $\mathscr L$ le
faisceau $\mathscr L^*=\mathscr Isom_{\mathcal O_X}(\mathcal O_X,\mathscr L)$ des
germes d'isomorphismes de $\mathcal O_X$ sur $\mathscr L$; il est immédiat que
$\mathscr L^*$ est canoniquement muni d'une structure de faisceau principal homogène
sous $\mathcal O_X^*\cong\mathscr Isom_{\mathcal O_X}(\mathcal O_X,\mathcal O_X)$.
La vérification du fait que le foncteur $\mathscr L\mapsto\mathscr L^*$ est une
équivalence de catégories est immédiate.

Cela étant, considérons de façon générale un espace topologique $X$ et un faisceau de
groupes $\mathscr G$ sur $X$; soit $Y$ une partie fermée de $X$, posons $U=X-Y$, et
soit $i$ un entier tel que $0\leq i\leq2$; on dira que le couple $(X,Y)$ est
\emph{$i$-pur pour $\mathscr G$} si, pour tout ouvert $V$ de $X$, le foncteur
restriction $\mathscr P\mapsto\mathscr P|(U\cap V)$, de la catégorie de faisceaux
principaux homogènes sous le faisceau de groupes $\mathscr G|V$, dans la catégorie
analogue sous le faisceau de groupes $\mathscr G|(U\cap V)$, est \emph{fidèle}
($i=0$), resp.\ \emph{pleinement fidèle} ($i=1$), resp.\ une \emph{équivalence de
catégories} ($i=2$). Dans le cas où $X$ est un espace annelé et
$\mathscr G=\mathcal O_X^*$, dire que $(X,Y)$ est $i$-pur signifie donc pour $i=0$,
que l'homomorphisme $\mathcal O_X\to j_*(\mathcal O_U)$ est \emph{injectif}; pour
$i=1$, que cet homomorphisme est \emph{bijectif}; et enfin, pour $i=2$, que le couple
$(X,Y)$ est \emph{parafactoriel} (21.13.3).

Revenant au cas général, rappelons que les morphismes de la catégorie des faisceaux
principaux sous $\mathscr G$ sont par définition les isomorphismes. Dire que le couple
$(X,Y)$ est $0$-pur signifie donc que pour tout ouvert $V$ de $X$, l'homomorphisme
canonique $\mathrm H^0(V,\mathscr G)\to\mathrm H^0(U\cap V,\mathscr G)$ est injectif;
dire que le couple $(X,Y)$ est $1$-pur signifie que l'homomorphisme canonique
$\mathrm H^0(V,\mathscr G)\to\mathrm H^0(U\cap V,\mathscr G)$ est bijectif (et alors
l'homomorphisme canonique
$\mathrm H^1(V,\mathscr G)\to\mathrm H^1(U\cap V,\mathscr G)$ est injectif); enfin,
on peut montrer que pour que le couple $(X,Y)$ soit $2$-pur, il faut et il suffit que
les homomorphismes
$\mathrm H^i(V,\mathscr G)\to\mathrm H^i(U\cap V,\mathscr G)$ soient bijectifs pour
$i=0$ et $i=1$. Lorsque $\mathscr G$ est un faisceau de groupes commutatifs, en
introduisant les faisceaux de cohomologie $\mathscr H_Y^p(\mathscr G)$ définis au
chap.~III, 3\textsuperscript{e} partie, dire que $(X,Y)$ est $i$-pur pour $i\leq2$
signifie que $\mathscr H_Y^p(\mathscr G)=0$ pour $p\leq i$; sous cette forme, la
notion se généralise donc aussitôt pour tout entier $i$.
\end{remark}

\begin{proposition}[21.13.14]
\label{IV.21.13.14}
Soient $X$ un préschéma localement noethérien \emph{normal},
$(U_\lambda)_{\lambda\in L}$ une famille filtrante décroissante d'ouverts de $X$. Les
conditions suivantes sont équivalentes :
\begin{enumerate}[label=\emph{\alph*)}]
  \item Tout cycle $1$-codimensionnel $Z$ sur $X$ tel qu'il existe un indice
  $\lambda$ pour lequel $Z|U_\lambda$ soit localement principal (21.6.7), est
  localement principal.
  \item Pour tout $x\in X$ tel que $\dim(\mathcal O_{X,x})\geq2$ et que $x$
  n'appartienne pas à $\bigcap_{\lambda\in L}U_\lambda$, l'anneau
  $\mathcal O_{X,x}$ est parafactoriel.
\end{enumerate}
\end{proposition}

\begin{proof}
\oldpage[IV-4]{322}
La condition $b)$ peut encore s'exprimer de la façon suivante :

\emph{$b')$ Pour toute partie fermée $Y$ de $X$ telle que
$\operatorname{codim}(Y,X)\geq2$ et telle qu'il existe $\lambda$ pour lequel
$Y\subset X-U_\lambda$, le couple $(X,Y)$ est parafactoriel.}

La propriété $a)$ peut encore s'exprimer de la façon suivante : si l'ensemble fermé
$N$ des points $x\in X$ où $Z$ n'est pas principal est contenu dans l'un des
$X-U_\lambda$, alors $Z$ est localement principal. Comme $X$ est normal, la condition
$\dim(\mathcal O_{X,x})\leq1$ entraîne que $\mathcal O_{X,x}$ est un corps ou un
anneau de valuation discrète, donc $Z_x$ est principal; en d'autres termes, on a
nécessairement $\operatorname{codim}(N,X)\geq2$ (5.1.3). La propriété $a)$ est donc
équivalente à la suivante :

\emph{$a')$ S'il existe un ensemble fermé $Y$ contenu dans un des $X-U_\lambda$, tel
que $\operatorname{codim}(Y,X)\geq2$ et que $Z|(X-Y)$ soit localement principal,
alors $Z$ est localement principal.}

Notons d'autre part que $b)$ implique $b')$ en vertu de (21.13.10); inversement, si
$b')$ est vérifiée et si $x\notin U_\lambda$, alors $Y=\{x\}$ est contenue dans
$X-U_\lambda$ et on a $\operatorname{codim}(Y,X)\geq2$, donc il résulte encore de
(21.13.10) que $\mathcal O_{X,x}$ est parafactoriel, ce qui prouve l'équivalence de
$b)$ et $b')$. On est ainsi ramené à prouver l'équivalence de $a')$ et de $b')$.
Notons que puisque $X$ est normal, on a, pour toute partie fermée $Y$ de $X$ telle que
$\operatorname{codim}(Y,X)\geq2$, $\operatorname{prof}_Y(\mathcal O_X)\geq2$
(5.8.6); si l'on pose $U=X-Y$, l'homomorphisme
$\mathcal O_X\to j_*(\mathcal O_U)$ est donc bijectif (21.13.4); comme les conditions
$a')$ et $b')$ sont locales, on voit qu'il suffit de montrer l'équivalence des
conditions suivantes, lorsque $X$ est noethérien et normal et $Y$ fermé dans $X$ et
tel que $\operatorname{codim}(Y,X)\geq2$ :

\emph{$a'')$ Pour tout cycle $1$-codimensionnel $Z$ sur $X$, l'hypothèse que $Z|U$
(où $U=X-Y$) est localement principal entraîne que $Z$ lui-même est localement
principal.}

\emph{$b'')$ L'homomorphisme canonique
$\operatorname{Pic}(X)\to\operatorname{Pic}(U)$ est surjectif.}

Prouvons d'abord que $a'')$ entraîne $b'')$; un élément
$\zeta_0\in\operatorname{Pic}(U)$ a pour image dans $\mathfrak{Cl}(U)$ (21.6.11) la
classe d'un cycle $1$-codimensionnel $Z_0$ sur $U$, qui est localement principal.
L'hypothèse $\operatorname{codim}(Y,X)\geq2$ entraîne que l'homomorphisme de
restriction $\mathscr Z^1(X)\to\mathscr Z^1(U)$ est bijectif, donc $Z_0=Z|U$, où $Z$
est un cycle $1$-codimensionnel sur $X$; puisque $Z_0$ est localement principal, il en
est de même de $Z$ en vertu de $a'')$; il résulte de (21.6.11) que l'image de $Z$ dans
$\mathfrak{Cl}(X)$ est l'image d'un unique élément $\zeta\in\operatorname{Pic}(X)$,
et il est clair alors que $\zeta_0$ est l'image de $\zeta$.

Inversement, prouvons que $b'')$ entraîne $a'')$. Soit donc $Z$ un cycle
$1$-codimensionnel sur $X$ tel que $Z|U$ soit localement principal; l'image de $Z|U$
dans $\mathfrak{Cl}(U)$ est donc l'image d'un unique élément
$\zeta_0\in\operatorname{Pic}(U)$ (21.6.11). Par hypothèse il existe
$\zeta\in\operatorname{Pic}(X)$ dont l'image dans $\operatorname{Pic}(U)$ est
$\zeta_0$; l'image de $\zeta$ dans $\mathfrak{Cl}(X)$ est donc la classe d'un cycle
$1$-codimensionnel $Z'$ sur $X$ tel que $Z|U$ et $Z'|U$ soient linéairement
équivalents. Mais puisque $U$ est schématiquement dense dans $X$ l'image de
$\mathscr Z^1_{\mathrm{princ}}(X)$ par l'isomorphisme
$\mathscr Z^1(X)\to\mathscr Z^1(U)$ est $\mathscr Z^1_{\mathrm{princ}}(U)$, donc $Z$
et $Z'$ sont linéairement équivalents, et puisque $Z'$ est localement principal, il
en est de même de $Z$.
\end{proof}

\begin{corollary}[21.13.15]
\label{IV.21.13.15}
Soient $X$ un préschéma localement noethérien et normal, $S$ une partie de $X$. Les
conditions suivantes sont équivalentes :
\begin{enumerate}[label=\emph{\alph*)}]
  \item Tout cycle $1$-codimensionnel sur $X$ qui est principal aux points de $S$ est
  localement principal.
  \oldpage[IV-4]{323}
  \item Pour tout $x\in X$ qui n'est pas générisation d'un point de $S$ (autrement
  dit, tel que $\overline{\{x\}}\cap S=\varnothing$) et qui est tel que
  $\dim(\mathcal O_{X,x})\geq2$, l'anneau $\mathcal O_{X,x}$ est parafactoriel.
\end{enumerate}
\end{corollary}

\begin{proof}
Il est clair que l'ensemble $S'$ des points de $X$ générisations de points de $S$ est
l'intersection des voisinages ouverts de $S$. On peut se borner au cas où $X$ est
noethérien (les propriétés $a)$ et $b)$ étant locales sur $X$); comme tout cycle
$1$-codimensionnel sur $X$ qui est principal aux points de $S$ l'est aussi aux points
d'un voisinage ouvert de $S$, on voit que la condition $a)$ de (21.13.15) équivaut à
la condition $a)$ de (21.13.14) appliquée à la famille $(U_\lambda)$ des voisinages
ouverts de $S$. Le corollaire résulte alors de l'équivalence de $a)$ et $b)$ dans
(21.13.14).
\end{proof}

\begin{remark}[21.13.16]
\label{IV.21.13.16}
Sous les conditions générales de (21.13.14), supposons de plus que l'on ait
$\varinjlim\operatorname{Pic}(U_\lambda)=0$. Alors la condition $a)$ de (21.13.14)
est aussi équivalente à la suivante :
\begin{enumerate}
  \item[c)] \emph{Tout cycle $1$-codimensionnel sur $X$ dont le support est contenu
  dans un des ensembles $X-U_\lambda$ est localement principal.}
\end{enumerate}

On peut se borner au cas où $X$ est irréductible et tous les $U_\lambda$ sont non
vides. Il est clair que $a)$ entraîne $c)$, car si le cycle $1$-codimensionnel $Z$ sur
$X$ a son support dans $X-U_\lambda$, $Z|U_\lambda=0$, donc $Z|U_\lambda$ est
localement principal. Inversement $c)$ entraîne $a)$ : soit en effet $Z$ un cycle
$1$-codimensionnel sur $X$ tel que $Z|U_\lambda$ soit localement principal; comme $X$
est normal, $Z|U_\lambda=\operatorname{cyc}(D_\lambda)$, où $D_\lambda$ est un
diviseur sur $U_\lambda$ (21.6.10, (ii)), et l'hypothèse implique (en vertu de
(21.3.4)) qu'il y a un ensemble $U_\mu\subset U_\lambda$ tel que
$D_\mu=D_\lambda|U_\mu$ soit équivalent à $0$. Si
$D_\mu=\operatorname{div}(f_\mu)$, où $f_\mu$ est une fonction rationnelle régulière
sur $U_\mu$, $f_\mu$ peut être considérée comme une fonction rationnelle régulière
sur $X$. Si $Z'=\operatorname{cyc}(\operatorname{div}(f_\mu))$, on voit donc que
$Z''=Z-Z'$ a son support dans $X-U_\mu$; en vertu de l'hypothèse, $Z''$ est
localement principal, d'où la conclusion.

La condition $\varinjlim\operatorname{Pic}(U_\lambda)=0$ sera en particulier remplie
si $(U_\lambda)$ est la famille des voisinages ouverts d'un \emph{point} $z\in X$,
car pour tout $\mathcal O_{U_\lambda}$-Module inversible $\mathscr L_\lambda$, il
existe par définition un $U_\mu\subset U_\lambda$ tel que
$\mathscr L_\lambda|U_\mu$ soit isomorphe à $\mathcal O_{U_\mu}$. On voit donc que,
dans l'énoncé de (21.13.15), si $S=\{z\}$, les conditions $a)$ et $b)$ équivalent
aussi à la suivante :
\begin{enumerate}
  \item[c)] \emph{Tout cycle $1$-codimensionnel sur $X$ dont le support ne contient
  pas $z$ est localement principal.}
\end{enumerate}
Si de plus $\operatorname{Pic}(X)=0$, on en conclura que cette condition entraîne que
tout cycle $1$-codimensionnel dont le support ne contient pas $z$ est principal.
\end{remark}

\subsection{Le théorème de Ramanujam--Samuel}
\EGAReaderTarget{ega4.c04.subsec.21-14}{subsection:IV.21.14}
\label{IV.21.14}

\begin{theorem}[21.14.1]
\label{IV.21.14.1}
\textup{(Ramanujam--Samuel)}. ---
Soit $A$ un anneau local noethérien d'idéal maximal $\mathfrak m$, tel que son complété
$\widehat A$ soit intègre et intégralement clos. Soient $B$ un anneau local noethérien
tel que $\dim(B)\geq\dim(A)$, $\rho:A\to B$ un homomorphisme local faisant de $B$
une $A$-algèbre formellement lisse (pour les topologies préadiques
$(\mathbf 0, 19.3.1)$) et tel que le corps résiduel de $B$ soit fini sur celui de $A$.
Alors tout cycle $1$-codimensionnel sur $\operatorname{Spec}(B)$ qui est principal au
point $\mathfrak p=\mathfrak mB$, est un cycle $1$-codimensionnel principal.
\end{theorem}

\begin{proof}
Si $k=A/\mathfrak m$ est le corps résiduel de $A$, $B/\mathfrak mB$ est une
$k$-algèbre formellement lisse (pour sa topologie préadique) $(\mathbf 0, 19.3.5)$,
\oldpage[IV-4]{324}
donc régulière, et en particulier intègre, autrement dit
$\mathfrak p=\mathfrak mB$ est bien un idéal premier de $B$, ce qui justifie
l'énoncé. Tout revient évidemment à prouver que tout idéal premier $\mathfrak q$ de
$B$ non contenu dans $\mathfrak p$ est \emph{principal}.

Soient $\widehat A$, $\widehat B$ les complétés de $A$ et $B$ respectivement, de
sorte que l'idéal maximal de $\widehat A$ est $\mathfrak m\widehat A$; on sait
$(\mathbf 0, 19.3.6)$ que $\widehat B$ est une $\widehat A$-algèbre formellement
lisse pour les topologies adiques. Soit $k'$ le corps résiduel de $\widehat B$,
extension finie de $k$; il existe un homomorphisme local $\widehat A\to C$, où $C$
est un anneau local noethérien qui est un $\widehat A$-module \emph{fini} et
\emph{plat} (donc libre) et est tel que $C/\mathfrak mC$ soit isomorphe à $k'$
$(\mathbf 0_{\mathrm{III}}, 10.3.1)$; on en déduit que $C$ est \emph{complet}, et il
résulte alors de (7.5.1), (7.5.3) et (6.5.4, (ii)) que $C$ est \emph{intègre et
intégralement clos}. En outre,
\[
  D=\widehat B\widehat\otimes_{\widehat A}C
\]
est un anneau semi-local complet, composé direct d'anneaux locaux complets dont l'un,
$D_0$, a pour corps résiduel $k'$ (puisque $k'\otimes_k k'$ est composé direct
d'anneaux locaux dont l'un est isomorphe à $k'$). Comme $D$ est formellement lisse
sur $C$, il en est de même de $D_0$; par suite $D_0/\mathfrak mD_0$ est une
$k'$-algèbre formellement lisse, de corps résiduel $k'$, ce qui entraîne qu'elle est
$k'$-isomorphe à une algèbre de séries formelles
$k'[[T_1,\ldots,T_n]]$ $(\mathbf 0, 19.6.4)$; on en conclut, par
$(\mathbf 0, 19.7.1.5)$, que $D_0$ est $C$-isomorphe à
$C[[T_1,\ldots,T_n]]$, et par suite \emph{intègre et intégralement clos} (Bourbaki,
\emph{Alg. comm.}, chap.~V, \S1, n\textsuperscript{o}~4, prop.~14). Comme les
morphismes $\operatorname{Spec}(D_0)\to\operatorname{Spec}(\widehat B)$ et
$\operatorname{Spec}(D_0)\to\operatorname{Spec}(B)$ sont fidèlement plats, on en
déduit que $B$ et $\widehat B$ sont aussi \emph{intègres et intégralement clos}
(2.1.13). Cela prouve que l'idéal $\mathfrak qD_0$ est divisoriel (Bourbaki,
\emph{Alg. comm.}, chap.~VII, \S1, n\textsuperscript{o}~10, prop.~15) et non contenu
dans $\mathfrak mD_0$, sans quoi on aurait
\[
  \mathfrak q=(\mathfrak qD_0)\cap B
  \subset(\mathfrak mD_0)\cap B=\mathfrak mB=\mathfrak p
\]
par fidèle platitude $(\mathbf 0_{\mathrm I}, 6.5.1)$. On en conclut qu'aucun des
idéaux premiers $\mathfrak r_h$ de hauteur $1$ dans $D_0$ qui contiennent
$\mathfrak qD_0$ ne peut être contenu dans $\mathfrak mD_0$; si l'on prouve que ces
idéaux sont principaux, il en résultera que $\mathfrak qD_0$ est principal, les
diviseurs (au sens de Bourbaki) de $\mathfrak qD_0$ et d'un produit de puissances des
$\mathfrak r_h$ étant égaux (Bourbaki, \emph{Alg. comm.}, chap.~VII, \S1,
n\textsuperscript{o}~4, prop.~5). Comme $\mathfrak qD_0\cap B=\mathfrak q$, on en
déduit par fidèle platitude que $\mathfrak q$ est principal (2.5.2), en utilisant le
fait que $B$ est un anneau local.

On est ainsi ramené à prouver le théorème dans le cas particulier où $A$ est
\emph{complet} et $B=A[[T_1,\ldots,T_n]]$. Montrons d'abord qu'on peut se ramener au
cas où $n=1$. En effet, procédons par récurrence sur $n$, et soit
$f\in\mathfrak q$ un élément n'appartenant pas à $\mathfrak p=\mathfrak mB$; il
suffira de montrer qu'il existe un $A$-automorphisme $\sigma$ de $B$ tel que
$\sigma(f)$ n'appartienne pas à l'idéal
\[
  \mathfrak p'=\mathfrak p+BT_1+\cdots+BT_{n-1};
\]
en effet, si $B'=A[[T_1,\ldots,T_{n-1}]]$, $B'$ est complet et a pour idéal maximal
$\mathfrak mB'+B'T_1+\cdots+B'T_{n-1}=\mathfrak m'$. Comme on a
$B=B'[[T_n]]$ et $\mathfrak p'=\mathfrak m'B$, il résultera de l'hypothèse de
récurrence (compte tenu de ce que $B'$ est intègre et intégralement clos) que l'idéal
$\sigma(\mathfrak q)$ est principal dans $B$, donc aussi $\mathfrak q$. Or, si
$\overline f\in k[[T_1,\ldots,T_n]]$ est l'image de $f$ (« série réduite » de $f$),
on sait (Bourbaki, \emph{Alg. comm.}, chap.~VII, \S3, n\textsuperscript{o}~7,
lemme~3) qu'il y a un $A$-automorphisme $\sigma$ de $B$ tel que
$(\overline{\sigma(f)})(0,\ldots,0,T_n)\neq0$, ce qui implique évidemment
$\sigma(f)\notin\mathfrak p'$.

Supposons donc $n=1$ et écrivons $T$ au lieu de $T_1$, de sorte que $B=A[[T]]$. Il
suffit de montrer que dans l'anneau de polynômes $A[T]$, l'idéal
$\mathfrak q\cap A[T]$ est principal;
\oldpage[IV-4]{325}
en effet, soit $f_0$ un élément de $\mathfrak q$ n'appartenant pas à $\mathfrak mB$;
c'est une série formelle dont la série réduite $\overline f_0$ n'est pas nulle; donc,
en vertu du théorème de préparation (Bourbaki, \emph{Alg. comm.}, chap.~VII, \S3,
n\textsuperscript{o}~8, prop.~5), pour tout $f\in\mathfrak q$, il existe $g\in B$ et
un polynôme $r\in A[T]$ tels que $f=gf_0+r$ et on a donc
$r\in\mathfrak q\cap A[T]$; d'autre part (\emph{loc. cit.}, prop.~6) il existe un
polynôme distingué non constant $F_0\in A[T]$ et un élément inversible $u\in B$ tels
que $f_0=uF_0$, donc on a aussi $F_0\in\mathfrak q\cap A[T]$, ce qui prouve que
$\mathfrak q$ est engendré par $\mathfrak q\cap A[T]=\mathfrak q_1$. Puisque $B$
est plat sur $A[T]$ $(\mathbf 0_{\mathrm I}, 7.3.3)$, il résulte de Bourbaki,
\emph{Alg. comm.}, chap.~VII, \S1, n\textsuperscript{o}~10, prop.~15, que
$\mathfrak q_1$ est un idéal premier de hauteur $1$ dans $A[T]$.

En outre, on a nécessairement $\mathfrak q_1\cap A=0$; sinon,
$\mathfrak q_1\cap A$ serait nécessairement de hauteur $\geq1$, et il résulterait de
(5.5.3) que l'on aurait $\mathfrak q_1=(\mathfrak q_1\cap A)A[T]$. Mais alors,
puisque $\mathfrak q_1\cap A\subset\mathfrak m$, on aurait
$\mathfrak q_1\subset\mathfrak mA[T]$ contrairement à l'hypothèse sur
$\mathfrak q$. Si $K$ est le corps des fractions de $A$, $\mathfrak q_1K[T]$ est
donc un idéal premier distinct de $0$ et de $K[T]$ dans $K[T]$, donc de la forme
$hK[T]$, où
\[
  h(T)=T^m+a_1T^{m-1}+\cdots+a_m,
  \qquad m\geq1,
\]
avec les $a_i\in K$, $h$ étant irréductible dans $K[T]$. Mais on a vu plus haut qu'il
existe dans $\mathfrak q_1$ un polynôme distingué non constant $F_0\in A[T]$. Si $t$
est la classe de $T$ dans $K[T]/\mathfrak q_1K[T]$, $t$ est donc racine du polynôme
$F_0$ dans une extension de $K$ et par suite $h$ divise $F_0$ dans $K[T]$; mais
puisque $h$ et $F_0$ sont unitaires, cela entraîne que les coefficients $a_i$ de $h$
sont entiers sur $A$ (Bourbaki, \emph{Alg. comm.}, chap.~V, \S1,
n\textsuperscript{o}~3, prop.~11), donc appartiennent à $A$ puisque $A$ est
intégralement clos. Autrement dit, on a
\[
  h\in A[T]\cap\mathfrak q_1K[T]=\mathfrak q_1
\]
(Bourbaki, \emph{Alg. comm.}, chap.~II, \S1, n\textsuperscript{o}~5, prop.~11);
comme tout polynôme $g\in\mathfrak q_1$ est divisible par $h$ dans $K[T]$ et que $h$
est unitaire, les coefficients de $g/h$ appartiennent à $A$, donc
$\mathfrak q_1=hA[T]$. C.Q.F.D.
\end{proof}

L'énoncé (21.14.1) est \emph{équivalent} au suivant :

\begin{corollary}[21.14.2]
\label{IV.21.14.2}
Sous les hypothèses de (21.14.1) sur $A$, $B$ et $\rho$ pour tout idéal premier
$\mathfrak q$ de $B$ non contenu dans $\mathfrak p=\mathfrak mB$ et tel que
$\dim(B_\mathfrak q)\geq2$, l'anneau $B_\mathfrak q$ est parafactoriel. En
particulier, si $\dim(B\otimes_A k)>0$ (\emph{i.e.} $\dim B>\dim A$
$((\mathbf 0, 19.7.1)$ et $(\mathbf 0, 6.1.1)))$, l'anneau $B$ est parafactoriel.
\end{corollary}

\begin{proof}
En effet, on a vu dans la démonstration de (21.14.1) que les conditions de l'énoncé
entraînent que $B$ est intègre et intégralement clos; l'équivalence des énoncés
(21.14.1) et (21.14.2) résulte alors de (21.13.15) appliqué à
$X=\operatorname{Spec}(B)$ et $S=\{\mathfrak p\}$.
\end{proof}

\begin{proposition}[21.14.3]
\label{IV.21.14.3}
Soient $S$ un préschéma normal, $f:X\to S$ un morphisme lisse.
\begin{enumerate}[label=(\roman*)]
  \item Si $S$ est localement noethérien (donc aussi $X$), alors, pour tout $x\in X$
  tel que $\dim(\mathcal O_{X,x})\geq2$ et que $x$ ne soit pas point maximal de sa
  fibre $f^{-1}(f(x))$, l'anneau $\mathcal O_{X,x}$ est parafactoriel. Tout cycle
  $1$-codimensionnel $Z$ sur $X$ tel que $\operatorname{Supp}(Z)$ ne contienne aucune
  composante irréductible d'une fibre $f^{-1}(s)=X_s$, est localement principal.
  \item Soit $Y$ une partie fermée de $X$ ne contenant aucune composante irréductible
  d'une fibre $X_s$, et telle que, pour tout point maximal $\eta$ de $S$,
  $Y_\eta=Y\cap X_\eta$ soit de codimension $\geq2$ dans $X_\eta$. Alors le couple
  $(X,Y)$ est parafactoriel.
\end{enumerate}
\end{proposition}

\begin{proof}
On notera que les conditions de (ii) sont remplies si $Y$ est une partie fermée de $X$
telle que, pour tout $s\in S$, $Y_s=Y\cap X_s$ soit de codimension $\geq2$ dans
$X_s$.

Pour prouver (21.14.3), notons d'abord que sous les hypothèses de (i), si l'on pose
$Y=\overline{\{x\}}$, il existe un voisinage ouvert $V$ de $x$ dans $X$ tel que
$Y\cap V$ et $V$ vérifient les conditions de (ii) : en effet, il résulte de l'hypothèse
et de (9.5.3) que l'on peut prendre $V$ tel que $V\cap Y$ ne contienne aucune
composante irréductible d'un $X_s$; d'autre part, pour tout point maximal $\eta$ de
$S$ tel que $Y\cap X_\eta\neq\varnothing$, les points de $Y\cap X_\eta$ sont des
spécialisations de $x$, donc en un tel point $z$ on a
$\dim(\mathcal O_{X,z})\geq2$, et comme
$\mathcal O_{X,z}=\mathcal O_{X_\eta,z}$ ($S$ étant réduit au point $\eta$ en vertu
de (2.1.13)), $Y\cap X_\eta$ est bien de codimension $\geq2$ dans $X_\eta$. Comme le
couple $(V,Y\cap V)$ est alors parafactoriel en vertu de (ii), on conclut de (21.13.10)
que l'anneau $\mathcal O_{X,x}$ est parafactoriel, c'est-à-dire la première assertion
de (i). Pour prouver la seconde, on peut se borner au cas où
$Z=\overline{\{z\}}$, où $z\in X^{(1)}$; comme $\mathcal O_{X,z}$ est un anneau local
intégralement clos et de dimension $1$, c'est un anneau de valuation discrète et $Z$
est donc principal au point $z$; pour tout autre point $x$ de
$\operatorname{Supp}(Z)$, on a donc $\dim(\mathcal O_{X,x})\geq2$ et $x$ n'est pas
point maximal de sa fibre par hypothèse, donc $\mathcal O_{X,x}$ est parafactoriel.
Appliquons alors (21.13.15) en prenant pour $S$ l'ensemble formé de $z$ et des points
maximaux des fibres de $f$; il est clair que $Z$ est principal aux points de $S$
puisque $z$ est le seul point de $S$ appartenant à $\operatorname{Supp}(Z)$; comme
d'autre part la condition $b)$ de (21.13.15) est évidemment remplie, on en conclut que
$Z$ est localement principal.

On est donc ramené à prouver (ii). Posons $U=X-Y$ et soit $j:U\to X$ l'injection
canonique.
\oldpage[IV-4]{326}
Comme les hypothèses et la conclusion sont locales sur $S$ et sur $X$ en vertu de
(21.13.6, (i)), on peut se borner au cas où $S$ et $X$ sont affines, $f$ étant donc de
présentation finie. Puisque les fibres de $f$ sont régulières, l'hypothèse sur $Y$
entraîne que, pour tout point $x\in Y$, on a
\[
  \operatorname{prof}(\mathcal O_{X_{f(x)},x})
  =\dim(\mathcal O_{X_{f(x)},x})\geq1;
\]
il résulte de (19.9.8) que l'homomorphisme canonique
$\mathcal O_X\to j_*(\mathcal O_U)$ est \emph{injectif}. De plus, remplaçant $Y$ par
une partie fermée plus grande suivant la méthode de la démonstration de (19.9.8), et
utilisant (21.13.6, (ii)), on peut supposer que $Y$ est défini par un idéal \emph{de
type fini} de l'anneau de $X$, donc l'ouvert $U=X-Y$ est \emph{quasi-compact et
quasi-séparé} et l'ensemble fermé $Y$ est \emph{constructible}.

Pour prouver (ii), il suffit, pour tout $\mathcal O_U$-Module inversible $\mathscr L$
et tout point $x\in Y$ donnés, d'établir l'existence d'un voisinage ouvert $V$ de $x$
dans $X$ tel que : $1^{\mathrm o}$ l'homomorphisme canonique
$\mathcal O_V\to(j_V)_*(\mathcal O_{U\cap V})$ est surjectif;
$2^{\mathrm o}$ il existe un $\mathcal O_V$-Module inversible $\mathscr L'_V$ tel que
$\mathscr L'_V|(U\cap V)=\mathscr L|(U\cap V)$ ((21.13.5) et (21.13.3)).

Posons $s=f(x)$; nous allons voir qu'on peut se borner au cas où
$S=S_0=\operatorname{Spec}(\mathcal O_{S,s})$. Soit en effet $(S_\nu)$ un système
fondamental de voisinages ouverts affines de $s$ dans $S$, et posons
$X_\nu=f^{-1}(S_\nu)$, $Y_\nu=Y\cap X_\nu$, $U_\nu=X_\nu-Y_\nu$,
$j_\nu:U_\nu\to X_\nu$ étant l'injection canonique; $X_0=X\times_S S_0$ est alors
limite projective des $X_\nu$ (8.1.2, $a$)); supposons la proposition vraie pour le
morphisme $f_0:X_0\to S_0$ et pour $Y_0=Y\cap X_0$, et posons $U_0=X_0-Y_0$,
$j_0:U_0\to X_0$ étant l'injection canonique. La projection
$p_\nu:X_0\to X_\nu$ étant un morphisme affine, on a
\[
  (j_0)_*(\mathcal O_{U_0})
  =p_\nu^*((j_\nu)_*(\mathcal O_{U_\nu}))
  \qquad(\mathbf{II}, 1.5.2)
\]
et l'homomorphisme canonique
$\rho_0:\mathcal O_{X_0}\to(j_0)_*(\mathcal O_{U_0})$ n'est autre que
$p_\nu^*(\rho_\nu)$, où
$\rho_\nu:\mathcal O_{X_\nu}\to(j_\nu)_*(\mathcal O_{U_\nu})$ est l'homomorphisme
canonique; comme par hypothèse $\rho_0$ est surjectif, il en est de même de $\rho_\nu$
pour $\nu$ assez grand (8.5.7). Soit d'autre part $\mathscr L_0$ le
$\mathcal O_{U_0}$-Module inversible restriction de $\mathscr L$, et notons que les
$X_\nu$ sont affines, donc quasi-compacts et quasi-séparés; puisque par hypothèse il
existe un $\mathcal O_{X_0}$-Module inversible $\mathscr L'_0$ tel que
$\mathscr L'_0|U_0=\mathscr L_0$, il existe un $\nu$ assez grand et un
$\mathcal O_{X_\nu}$-Module quasi-cohérent $\mathscr L'_\nu$ de présentation finie
tels que $\mathscr L'_0$ soit isomorphe à $p_\nu^*(\mathscr L'_\nu)$ (8.5.2, (ii));
de plus (8.5.5) on peut supposer que $\mathscr L'_\nu$ est inversible. Enfin, comme
les $U_\nu$ sont quasi-compacts et quasi-séparés et que
$\mathscr L'_0|U_0=\mathscr L_0$, il existe $\nu$ assez grand tel que
$\mathscr L'_\nu|U_\nu$ et $\mathscr L|U_\nu$ soient isomorphes (8.5.2.5).

\oldpage[IV-4]{327}
On est ainsi ramené au cas où $S=\operatorname{Spec}(A)$,
$X=\operatorname{Spec}(B)$, où $A$ est un anneau \emph{local}; puisque $X$ est
normal et $f$ lisse, $S$ est normal (17.5.7), donc $A$ est \emph{intègre et
intégralement clos}. Considérons alors $A$ comme limite inductive de ses
\emph{sous-$\mathbb Z$-algèbres de type fini}; comme la clôture intégrale d'une telle
sous-$\mathbb Z$-algèbre est un sous-anneau de $A$ et est aussi une
$\mathbb Z$-algèbre de type fini (7.8.3, (ii), (iii) et (vi)), $A$ est limite inductive
filtrante des sous-$\mathbb Z$-algèbres de type fini $A_\lambda$ de $A$ qui sont des
anneaux \emph{intégralement clos}. En vertu de (1.8.4.2), il existe un indice
$\lambda$ et une $A_\lambda$-algèbre de type fini $B_\lambda$ tels que
\[
  B=B_\lambda\otimes_{A_\lambda}A
\]
à un $A$-isomorphisme près. Posons $S_\lambda=\operatorname{Spec}(A_\lambda)$,
$X_\lambda=\operatorname{Spec}(B_\lambda)$; si $p_\lambda:X\to X_\lambda$ est la
projection canonique, on peut en outre supposer (puisque $Y$ est
\emph{constructible}) que $Y=p_\lambda^{-1}(Y_\lambda)$, où $Y_\lambda$ est une
partie fermée de $X_\lambda$ (8.3.11). Soit $f_\lambda$ le morphisme
$X_\lambda\to S_\lambda$; en vertu de la transitivité des fibres et de (4.2.6),
$Y_\lambda$ ne contient aucune composante irréductible des fibres
$f_\lambda^{-1}(s_\lambda)$ pour aucun $s_\lambda\in S_\lambda$. Puisque $f$ est
lisse, on peut aussi supposer que $f_\lambda$ est lisse (17.7.8); enfin, l'image du
point générique $\eta$ de $S$ dans $S_\lambda$ est le point générique $\eta_\lambda$
de $S_\lambda$, et on a
\[
  \operatorname{codim}((Y_\lambda)_{\eta_\lambda},
  (X_\lambda)_{\eta_\lambda})
  =\operatorname{codim}(Y_\eta,X_\eta)\geq2
\]
en vertu de (6.1.4). On voit donc que $S_\lambda$, $X_\lambda$ et $Y_\lambda$
vérifient toutes les hypothèses de (ii), et en posant
$X_\mu=X_\lambda\times_{S_\lambda}S_\mu$ pour $\lambda\leq\mu$, il en est de même de
$S_\mu$, $X_\mu$ et $Y_\mu$. Montrons que si l'on prouve que le couple
$(X_\mu,Y_\mu)$ est parafactoriel pour tout $\mu\gg\lambda$, il en est de même de
$(X,Y)$. En effet, soient $U=X-Y$, $U_\mu=X_\mu-Y_\mu$, $j:U\to X$,
$j_\mu:U_\mu\to X_\mu$ les injections canoniques; la projection
$p_\mu:X\to X_\mu$ étant un morphisme affine, on a
\[
  j_*(\mathcal O_U)=p_\mu^*((j_\mu)_*(\mathcal O_{U_\mu}))
  \qquad(\mathbf{II}, 1.5.2),
\]
et par suite l'homomorphisme canonique
$\rho:\mathcal O_X\to j_*(\mathcal O_U)$ n'est autre que $p_\mu^*(\rho_\mu)$, où
$\rho_\mu:\mathcal O_{X_\mu}\to(j_\mu)_*(\mathcal O_{U_\mu})$ est l'homomorphisme
canonique; ce dernier étant par hypothèse bijectif, il en est de même de $\rho$.
D'autre part, pour tout $\mathcal O_U$-Module inversible $\mathscr L$, il existe un
$\mu\gg\lambda$ et un $\mathcal O_{U_\mu}$-Module inversible $\mathscr L_\mu$ tels
que $\mathscr L=p_\mu'^*(\mathscr L_\mu)$ (en désignant par
$p_\mu':U\to U_\mu$ la restriction de $p_\mu$) (8.5.2 et 8.5.5), $U_\lambda$ étant
quasi-compact puisque $X_\lambda$ est noethérien. Par hypothèse, il existe un
$\mathcal O_{X_\mu}$-Module inversible $\mathscr L'_\mu$ tel que
$\mathscr L'_\mu|U_\mu$ soit isomorphe à $\mathscr L_\mu$;
$\mathscr L'=p_\mu^*(\mathscr L'_\mu)$ est alors un $\mathcal O_X$-Module
inversible tel que $\mathscr L'|U$ soit isomorphe à $\mathscr L$, ce qui prouve notre
assertion.

On est ainsi ramené à prouver (ii) lorsque l'anneau $A$ est une
$\mathbb Z$-algèbre de type fini intégralement close; comme les anneaux locaux
$\mathcal O_{S,s}$ de $S$ sont alors des anneaux excellents intégralement clos, leurs
\emph{complétés} sont aussi intégralement clos (7.8.3, (ii), (iii) et (vii)). Pour
prouver que le couple $(X,Y)$ est parafactoriel, il suffit de montrer que pour tout
$x\in Y$, l'anneau $\mathcal O_{X,x}$ est parafactoriel (21.13.10). Soit $y$ un point
fermé de $Y_{f(x)}$, qui soit spécialisation de $x$ (5.1.11); si l'on pose
$s=f(x)=f(y)$, $\mathcal O_{S,s}$ a un
\oldpage[IV-4]{328}
complété intégralement clos, et $\mathcal O_{X,y}$ est une
$\mathcal O_{S,s}$-algèbre formellement lisse pour les topologies préadiques
(17.5.3); puisque $Y$ ne contient aucune composante irréductible de $X_s$, on a
$\dim(\mathcal O_{X,y})>\dim(\mathcal O_{S,s})$ (17.5.8). Si $s\neq\eta$, on a
$\dim(\mathcal O_{S,s})\geq1$; si au contraire $s=\eta$, on a par hypothèse
$\dim(\mathcal O_{X,x})\geq2$; donc dans tous les cas
$\dim(\mathcal O_{X,y})\geq2$. En outre, l'idéal premier de $\mathcal O_{X,y}$
correspondant au point $x$ n'est pas contenu dans
$\mathfrak m_s\mathcal O_{X,y}$, puisque $Y$ ne contient aucune composante
irréductible de $X_s$. Enfin, puisque $y$ est fermé dans $X_s$, $k(y)$ est extension
finie de $k(s)$ $(\mathbf I, 6.4.2)$. Dans tous les cas, on peut appliquer à
$A=\mathcal O_{S,s}$, $B=\mathcal O_{X,y}$ et
$B_\mathfrak q=\mathcal O_{X,x}$ le résultat de (21.14.2), ce qui achève la
démonstration.
\end{proof}

\begin{env}[21.14.4]
\label{IV.21.14.4}
\emph{Remarques}. ---

(i) Il se peut que les énoncés (21.14.1) et (21.14.2) restent valables \emph{sans
hypothèse} sur le corps résiduel de $B$. Énoncé avec cette généralité, le résultat
équivaudrait, en vertu de (21.13.15) et (21.13.12.1) au suivant : Soient $A$ un
anneau local noethérien complet, intègre et intégralement clos, $B$ un anneau local
noethérien qui est une $A$-algèbre formellement lisse, tel que
$\dim(B)>\dim(A)$; alors $B$ est parafactoriel.

(ii) Nous verrons au chap.~VI, en employant une technique de « descente finie »
appliquée au morphisme $Y'\to Y=\operatorname{Spec}(A)$, où $Y'$ est le normalisé de
$Y$, que la conclusion de (21.14.1) (ou de (21.14.2)) reste valable en remplaçant
l'hypothèse que $\widehat A$ est intégralement clos par l'hypothèse que
$\widehat A$ est \emph{réduit}, pourvu que $\dim B\geq\dim A+2$. Si on remplace cette
dernière condition par $\dim B\geq\dim A+3$, on peut même supprimer l'hypothèse que
$\widehat A$ est réduit. De même, la conclusion de (21.14.3) reste valable en
remplaçant l'hypothèse que $X$ est normal par l'hypothèse que $X$ est réduit, pourvu
que $\dim(\mathcal O_{X,x})\geq\dim(\mathcal O_{S,f(x)})+2$.

(iii) Au chap.~III, 3\textsuperscript{e} partie, on prouve que si $f:X\to S$ est un
morphisme lisse, $Y$ une partie fermée localement constructible de $X$ telle que,
pour tout $s\in S$, on ait (avec les notations de (21.14.3))
$\operatorname{codim}(Y_s,X_s)\geq3$, alors le couple $(X,Y)$ est parafactoriel
([41], XII, 4.8). Cette conclusion n'est plus valable lorsqu'on suppose seulement que
$\operatorname{codim}(Y_s,X_s)\geq2$ pour tout $s\in S$ et lorsqu'on ne suppose plus
$X$ réduit. Par exemple, soient $k$ un corps, $A=k[T]/(T^2)$, algèbre des nombres
duaux sur $k$, $S=\operatorname{Spec}(A)$,
$X=\operatorname{Spec}(A[T_1,T_2])$ ($T_1$, $T_2$ indéterminées), de sorte que
$X=\mathbb V_S^2$, $Y$ étant la « section nulle » de ce fibré; si $s$ est l'unique
point fermé de $S$, on a $X_s=\mathbb V_k^2$ et $Y_s$ est la « section nulle » de
$X_s$. Pour voir que le couple $(X,Y)$ n'est pas parafactoriel, il suffit de montrer
que l'anneau $B=A[T_1,T_2]_\mathfrak m$, où $\mathfrak m$ est l'idéal
$(T_1)+(T_2)$ qui définit $Y$, n'est pas parafactoriel (21.13.10). Soit
$B_0=B_{\mathrm{red}}$, et notons $U$ et $U_0$ les complémentaires du point fermé
dans $\operatorname{Spec}(B)$ et $\operatorname{Spec}(B_0)$; en raisonnant comme
dans (21.13.9), on a la suite exacte, extension de (21.13.9.1)
\[
  \Gamma(U,\mathcal O_U)^*\longrightarrow\Gamma(U_0,\mathcal O_{U_0})^*
  \longrightarrow\mathrm H^1(U_0,\mathcal O_{U_0})
  \longrightarrow\operatorname{Pic}(U)\longrightarrow\operatorname{Pic}(U_0).
\]
Or, on a $\Gamma(U,\mathcal O_U)=B$, $\Gamma(U_0,\mathcal O_{U_0})=B_0$ puisque
$\operatorname{prof}(B)=\operatorname{prof}(B_0)=2$ (5.10.5). Par ailleurs
$\operatorname{Pic}(U_0)=0$ puisque $B_0$ est factoriel, et
$\mathrm H^1(U_0,\mathcal O_{U_0})\neq0$ [41, 3, Exemple III-1]; comme
l'homomorphisme $B^*\to B_0^*$ est surjectif, on en conclut que
$\operatorname{Pic}(U)\neq0$, donc que $B$ n'est pas parafactoriel.
\oldpage[IV-4]{329}
(iv) Le résultat (21.14.3) a d'abord été démontré par C.~Seshadri [44] dans le cas
particulier où $S$ est un schéma algébrique normal sur un corps algébriquement clos
$k$ et $X=S\times_k T$, où $T$ est un $k$-préschéma algébrique lisse sur $k$. La
démonstration de Seshadri [44, p.~188--189] est de nature globale et utilise la
théorie des schémas de Picard. Elle donne en outre (\emph{loc. cit.}) des résultats
tels que le suivant (pour lequel on ne possède pas à l'heure actuelle de démonstration
par voie locale). Soient $S$, $T$ deux préschémas localement de type fini sur un
corps $k$, $X=S\times_k T$, $Z$ un cycle $1$-codimensionnel sur $X$ (considéré comme
$S$-préschéma); supposons vérifiées les conditions suivantes :
\begin{enumerate}[label=\arabic*\textsuperscript{o}]
  \item $S$ et $T$ sont géométriquement normaux sur $k$ (6.7.6);
  \item Pour tout point maximal $\eta$ de $S$, le cycle $1$-codimensionnel $Z_\eta$
  sur la fibre $X_\eta$, ayant même multiplicité que $Z$ en tout point de
  $X^{(1)}\cap X_\eta$, est localement principal (autrement dit, est l'image d'un
  diviseur de $X_\eta$, puisque $X_\eta$ est normal);
  \item Pour tout $s\in S$, $Z$ est principal aux points maximaux de la fibre $X_s$.
\end{enumerate}
Alors $Z$ est localement principal. En d'autres termes, $X$ étant normal (6.14.1),
pour tout $x\in X$ qui n'est pas maximal dans sa fibre $X_s$ et qui n'appartient à
aucune des « fibres génériques » $X_\eta$ (ce qui implique
$\dim(\mathcal O_{X,x})\geq2$ en vertu de (6.1.1)), l'anneau local
$\mathcal O_{X,x}$ est parafactoriel, en vertu de (21.12.15).
\footnote{En fait, dans l'article précité, Seshadri suppose que $k$ est algébriquement
clos, $T$ séparé et « semi-complet » (i.e.\ tel que
$\Gamma(T,\mathcal O_T)$ soit $k$-isomorphe à $k$) et remplace l'hypothèse
$3^{\mathrm o}$ par l'hypothèse plus forte que $\operatorname{Supp}(Z)$ ne contient
aucune des fibres $X_s$ pour $s\in S$. Mais comme l'énoncé est local sur $S$, on en
conclut aussitôt qu'il suffit de faire l'hypothèse $3^{\mathrm o}$, et cela prouve que
la conclusion (interprétée comme ci-dessus en termes de propriété de parafactorialité
des anneaux $\mathcal O_{X,x}$) est locale sur $S$ et sur $T$, ce qui permet
d'éliminer complètement l'hypothèse que $T$ est « semi-complet » et que $k$ est
algébriquement clos, puisque (quitte à passer d'abord à la clôture algébrique de $k$)
on peut supposer d'abord $T$ affine, ce qui permet de le plonger comme ouvert d'un
schéma projectif et normal sur $k$, auquel le résultat de Seshadri s'applique. Notons
aussi que, grâce à cette réduction, il suffit de faire la démonstration de Seshadri dans
le cas où $T$ est \emph{projectif} (et non seulement « semi-complet »), cas où la
théorie du schéma de Picard utilisée par Seshadri est contenue dans celle qui sera
développée au chap.~VI de notre Traité.}
\end{env}

\subsection{Diviseurs relatifs}
\EGAReaderTarget{ega4.c04.subsec.21-15}{subsection:IV.21.15}
\label{IV.21.15}

\begin{env}[21.15.1]
\label{IV.21.15.1}
Soient $S$ un préschéma, $f:X\to S$ un morphisme \emph{plat et localement de
présentation finie}. On a défini dans (20.6.1) le faisceau d'anneaux
$\mathscr M_{X/S}$ des \emph{germes de fonctions méromorphes sur $X$ relatives à
$S$}, sous-faisceau de $\mathscr M_X$; il est clair que l'injection canonique
$\mathcal O_X\to\mathscr M_X$ (20.1.4.1) applique $\mathcal O_X$ sur un
sous-faisceau de $\mathscr M_{X/S}$, avec lequel on l'identifie. Soit
$\mathscr M_{X/S}^*$ le faisceau (en groupes multiplicatifs) des germes de sections
inversibles de $\mathscr M_{X/S}$ qui est donc un sous-faisceau de $\mathscr M_X^*$
et contient $\mathcal O_X^*$ comme sous-faisceau.
\end{env}

\begin{definition}[21.15.2]
\EGAReaderTarget{ega4.c04.def.21-15-2}{IV.21.15.2}
On appelle \emph{faisceau des diviseurs sur $X$ relativement à $S$}, ou
\emph{faisceau des diviseurs sur $X$ transversaux à $f$}, et on note
$\mathscr{D}iv_{X/S}$ le faisceau quotient (de groupes commutatifs)
$\mathscr M_{X/S}^*/\mathcal O_X^*$; les sections de ce faisceau au-dessus de $X$ se
nomment \emph{diviseurs sur $X$ relatifs à $S$}, ou \emph{diviseurs sur $X$
transversaux à $f$}; ils forment un groupe commutatif noté $\operatorname{Div}(X/S)$.
\end{definition}

Il est clair que $\mathscr{D}iv_{X/S}$ s'identifie canoniquement à un sous-faisceau de
$\mathscr{D}iv_X$, et par suite $\operatorname{Div}(X/S)$ à un sous-groupe de
$\operatorname{Div}(X)$, que l'on note encore \emph{additivement}. Pour
\oldpage[IV-4]{330}
tout ouvert $U$ de $X$, on a
$\mathscr M_{X/S}^*|U=\mathscr M_{U/S}^*$, donc
$\mathscr{D}iv_{X/S}|U=\mathscr{D}iv_{U/S}$, et le faisceau
$\mathscr{D}iv_{X/S}$ est donc égal au préfaisceau
$U\mapsto\operatorname{Div}(U/S)$.

Puisque $\mathscr M_{X/S}^*$ est un sous-faisceau de $\mathscr M_X^*$, les
définitions, notations et formules relatives aux diviseurs des sections de
$\mathscr M_X^*$ au-dessus de $X$ (21.1.3) s'appliquent sans changement aux sections
de $\mathscr M_{X/S}^*$ au-dessus de $X$.

\begin{env}[21.15.3]
\label{IV.21.15.3}
La structure de faisceau de groupes ordonnés sur $\mathscr{D}iv_X$ (21.1.6) induit
sur le sous-faisceau $\mathscr{D}iv_{X/S}$ une structure de faisceau de groupes
ordonnés, pour laquelle le faisceau de monoïdes des germes de sections positives est
$\mathscr{D}iv_X^+\cap\mathscr{D}iv_{X/S}$, que l'on note
$\mathscr{D}iv_{X/S}^+$. On a
$\Gamma(X,\mathscr{D}iv_{X/S}^+)=\operatorname{Div}^+(X)\cap
\operatorname{Div}(X/S)$; on note ce sous-monoïde de $\operatorname{Div}(X/S)$ par
$\operatorname{Div}^+(X/S)$, et il est formé des éléments $\geq0$ pour une structure
de groupe ordonné sur $\operatorname{Div}(X/S)$. Il résulte de (21.1.5.1) que
$\mathscr{D}iv_{X/S}^+$ est l'image dans $\mathscr{D}iv_{X/S}$ du sous-faisceau de
monoïdes
\[
\label{IV.21.15.3.1}
  \mathcal O_{X/S}^{\mathrm{reg}}=\mathcal O_X\cap\mathscr M_{X/S}^*.
  \tag{21.15.3.1}
\]
\end{env}

Pour tout ouvert $U$ de $X$, $\Gamma(U,\mathcal O_{X/S}^{\mathrm{reg}})$ est
l'ensemble des sections $t$ de $\mathcal O_X$ au-dessus de $U$ telles que $t$ soit
régulière et que $1/t$ appartienne à $\Gamma(U,\mathscr M_{X/S})$, ce qui signifie,
avec les notations de (20.6.1), que $t\in\mathscr T_{X/S}(U)$, de sorte que le faisceau
$\mathcal O_{X/S}^{\mathrm{reg}}$ n'est autre que le faisceau noté $\mathscr T_{X/S}$
dans (20.6.1). On peut donc écrire
\[
\label{IV.21.15.3.2}
  \mathscr{D}iv_{X/S}^+=\mathcal O_{X/S}^{\mathrm{reg}}/\mathcal O_X^*
  \tag{21.15.3.2}
\]
à un isomorphisme canonique près.

\begin{env}[21.15.3.3]
\label{IV.21.15.3.3}
Soit $D\in\operatorname{Div}^+(X/S)$ et considérons le sous-préschéma fermé $Y(D)$
de $X$ défini par l'Idéal $\mathscr J_X(D)$ de $\mathcal O_X$ (21.6.5); en vertu de
ce qui précède, pour tout $x\in Y(D)$, il y a un voisinage ouvert $U$ de $x$ dans $X$
et une section $t\in\mathscr T_{X/S}(U)$ telle que
$\mathscr J_X(D)|U=\mathcal O_U\cdot t$; puisque $x\in Y(D)$, l'image $t_{f(x)}$ de
$t$ dans $\Gamma(U\cap X_{f(x)},\mathcal O_{X_{f(x)}})$ appartient à l'idéal maximal
de $\mathcal O_{X_{f(x)},x}$ et en outre, par définition, pour tout $s\in S$, l'image
$t_s$ de $t$ dans $\Gamma(U\cap X_s,\mathcal O_{X_s})$ est une section
\emph{régulière}. On déduit donc de (11.3.8) et (19.2.4) que l'immersion canonique
$Y(D)\to X$ est \emph{transversalement régulière relativement à $S$ et de
codimension $1$} au point $x$. La réciproque étant immédiate, on voit qu'on peut
\emph{identifier canoniquement} les diviseurs positifs sur $X$ relatifs à $S$ aux
sous-préschémas fermés $Y$ de $X$ tels que l'injection canonique $Y\to X$ soit une
immersion transversalement régulière relativement à $S$ et de codimension $1$. Nous
ferons d'ordinaire cette identification.
\end{env}

\begin{proposition}[21.15.4]
\label{IV.21.15.4}
Soient $D$ un diviseur sur $X$, $\mathscr J_X(D)$ l'Idéal fractionnaire inversible
correspondant (21.2.5). Pour que $D\in\operatorname{Div}(X/S)$, il faut et il suffit
que, pour tout $x\in X$, on ait
$(\mathscr J_X(D)\cap\mathscr M_{X/S}^*)_x\ne0$ (ou, ce qui revient au même, que
l'on ait $\mathscr J_X(D)\subset\mathscr M_{X/S}$ et
$\mathscr J_X(D)^{-1}\subset\mathscr M_{X/S}$).
\end{proposition}

\begin{proof}
En effet, dire que $D\in\operatorname{Div}(X/S)$ signifie que pour tout $x\in X$, il
existe un voisinage ouvert $U$ de $x$ et une section
$f\in\Gamma(U,\mathscr M_{X/S}^*)$ tels que $D|U=\operatorname{div}_U(f)$; comme
$\mathscr J_X(D)|U$ est l'Idéal fractionnaire inversible $\mathcal O_Uf$, la
proposition en résulte aussitôt.
\end{proof}

\oldpage[IV-4]{331}
\begin{proposition}[21.15.5]
\label{IV.21.15.5}
Soient $D$ un diviseur sur $X$, $\mathcal O_X(D)$ l'Idéal fractionnaire inversible et
$s_D$ la section méromorphe régulière de $\mathcal O_X(D)$ définis canoniquement par
$D$ (21.2.8 et 21.2.9). Pour que $D\in\operatorname{Div}(X/S)$, il faut et il suffit
que $s_D\in\Gamma(X,\mathscr M_{X/S}(\mathcal O_X(D)))$.
\end{proposition}

\begin{proof}
En effet, si $U$ est un ouvert de $X$ tel que $D|U=\operatorname{div}_U(f)$, où
$f\in\Gamma(U,\mathscr M_X^*)$, dire que
$s_D\in\Gamma(U,\mathscr M_{X/S}(\mathcal O_X(D)))$ signifie, en vertu des
définitions (20.6.2), que $f^{-1}\in\Gamma(U,\mathscr M_{X/S})$, d'où la
proposition.
\end{proof}

L'interprétation des diviseurs sur $X$ à l'aide des classes
$\operatorname{cl}(\mathscr L,s)$ (21.2.11) permet donc d'interpréter les éléments de
$\operatorname{Div}(X/S)$ comme les couples (à isomorphisme près) $(\mathscr L,s)$
tels que $\mathscr L$ soit un $\mathcal O_X$-Module inversible et que $s$ soit une
section méromorphe de $\mathscr L$ au-dessus de $X$, \emph{régulière relativement à
$S$} (20.6.5, (iii)).

\begin{proposition}[21.15.6]
\label{IV.21.15.6}
Soit $D$ un diviseur sur $X$ relatif à $S$, et supposons que pour tout $x\in X$ tel
que $\operatorname{prof}(\mathcal O_{X_{f(x)},x})=1$, on ait $D_x\geq0$
(resp. $D_x=0$). Alors on a $D\geq0$ (resp. $D=0$).
\end{proposition}

\begin{proof}
Comme dans (21.1.8), on peut se borner au cas où $D=\operatorname{div}(\varphi)$, où
$\varphi$ est une fonction méromorphe régulière relative à $S$, et tout revient à voir
que si $D_x\geq0$ en tout point $x\in X$ tel que
$\operatorname{prof}(\mathcal O_{X_{f(x)},x})=1$, $\varphi$ est partout définie dans
$X$. Mais cette hypothèse signifie que, si $T=X-\operatorname{dom}(\varphi)$, on a
$\operatorname{prof}(\mathcal O_{X_{f(x)},x})\geq2$ pour tout $x\in T$, et il suffit
pour conclure d'appliquer (20.6.6).
\end{proof}

\begin{env}[21.15.7]
\label{IV.21.15.7}
Soit $X'$ un second $S$-préschéma plat et localement de présentation finie sur $S$, et
soit $g:X'\to X$ un $S$-morphisme. Si le $S$-morphisme $g$ est \emph{plat}, on sait
(21.4.5) que l'image réciproque par $g$ de tout diviseur sur $X$ est définie; si de
plus $D\in\operatorname{Div}(X/S)$, il résulte de la définition (21.15.2) et de
(20.6.8) que l'on a alors $g^*(D)\in\operatorname{Div}(X'/S)$.
\end{env}

\begin{env}[21.15.8]
\label{IV.21.15.8}
Soit $X'$ un second $S$-préschéma plat et localement de présentation finie sur $S$, et
soit $g:X'\to X$ un $S$-morphisme \emph{fini et plat}. Notons que $g$ est alors
nécessairement de présentation finie (1.4.3, 1.4.6 et 1.6.3), donc
$g_*(\mathcal O_{X'})$ est un $\mathcal O_X$-Module plat et de présentation finie, et
par suite \emph{localement libre} (2.1.12), autrement dit $g$ est un morphisme
localement libre (18.2.7); pour tout $s\in S$, le morphisme correspondant
$g_s:X'_s\to X_s$ est donc aussi fini et localement libre. On déduit alors de (21.5.2)
et (20.6.1) que pour tout ouvert $U$ de $X$ et toute section
$t'\in\Gamma(g^{-1}(U),\mathscr T_{X'/S})$, la norme $N_{X'/X}(t')$ appartient à
$\Gamma(U,\mathscr T_{X/S})$; le raisonnement de (21.5.3) prouve ensuite que pour tout
$\mathcal O_{X'}$-Module inversible $\mathscr L'$ et toute section méromorphe $u'$ de
$\mathscr L'$ au-dessus de $X'$, \emph{régulière relativement à $S$}, la norme
$N_{X'/X}(u')$ est une section méromorphe de
$\mathscr L=N_{X'/X}(\mathscr L')$ au-dessus de $X$, régulière relativement à $S$.
L'interprétation des diviseurs relatifs à $S$ donnée dans (21.15.5) et la définition de
l'image directe d'un diviseur (21.5.5) prouvent alors que \emph{pour tout diviseur}
$D'\in\operatorname{Div}(X'/S)$, on a $g_*(D')\in\operatorname{Div}(X/S)$.
\end{env}

\begin{env}[21.15.9]
\EGAReaderTarget{ega4.c04.statement.21-15-9}{IV.21.15.9}
Considérons enfin un morphisme quelconque $S'\to S$, et (sous les hypothèses de
(21.15.1)) posons $X'=X\times_S S'$, qui est plat et localement de présentation finie
sur $S'$; si $p:X'\to X$ est la projection canonique, on a vu (20.6.9) qu'on a un
homomorphisme canonique
$p^*(\mathscr M_{X/S})\to\mathscr M_{X'/S'}$, qui transforme évidemment toute
section de $\mathscr M_{X/S}$ au-dessus d'un ouvert $U$, \emph{régulière relativement
à $S$}, en une section de $\mathscr M_{X'/S'}$ au-dessus de $p^{-1}(U)$, régulière
relativement à $S$ (20.6.5, (iii)); on conclut alors de la
\oldpage[IV-4]{332}
définition (21.15.2) et de l'exactitude à droite du foncteur $p^*$, que
l'homomorphisme précédent définit par passage aux quotients un homomorphisme canonique
\[
\label{IV.21.15.9.1}
  \operatorname{Div}(X/S)\longrightarrow\operatorname{Div}(X'/S')
  \tag{21.15.9.1}
\]
qui transforme évidemment les éléments de $\operatorname{Div}^+(X/S)$ en éléments de
$\operatorname{Div}^+(X'/S')$. On pose encore
$\operatorname{Div}(X'/S')=\operatorname{Div}_{X/S}(S')$ (resp.
$\operatorname{Div}^+(X'/S')=\operatorname{Div}_{X/S}^+(S')$), et on voit aussitôt
qu'on a ainsi défini deux foncteurs contravariants
\[
  \operatorname{Div}_{X/S}:\mathbf{Sch}_{/S}\longrightarrow\mathbf{Ab},\qquad
  \operatorname{Div}_{X/S}^+:\mathbf{Sch}_{/S}\longrightarrow\mathbf{Ens}
\]
de la catégorie des $S$-préschémas dans celle des groupes commutatifs (resp. des
ensembles). On verra plus loin (chap.~VI) des cas importants où le foncteur
$\operatorname{Div}_{X/S}^+$ est \emph{représentable} $(\mathbf 0_{\mathrm{III}},
8.1.8)$.

Pour tout diviseur $D\in\operatorname{Div}(X/S)$, l'image de $D$ par
l'homomorphisme (21.15.9.1) n'est autre d'ailleurs que l'\emph{image réciproque}
$p^*(D)$ (au sens de (21.4.2)) : l'existence de cette image réciproque et l'assertion
précédente sont en effet des conséquences immédiates de (20.6.5, (iii)) et (20.6.9).

Les éléments de $\operatorname{Div}(X'/S')$ s'appellent souvent, par abus de langage,
« \emph{familles de diviseurs sur $X$ relatifs à $S$, paramétrées par le
$S$-préschéma $S'$} »; on utilise surtout cette terminologie lorsqu'il s'agit de
diviseurs \emph{positifs}.
\end{env}

% FIN SOURCE ega4/ega4-21c-fr.tex

%
% DEBUT SOURCE ega4/ega4-backmatter-bibliography-fr.tex
% Transcription diplomatique française de EGA IV, quatrième partie, p. 333.
\clearpage
\oldpage[IV]{333}
\phantomsection
\addcontentsline{toc}{section}{Bibliographie}
\section*{BIBLIOGRAPHIE}

\begin{thebibliography}{44}

\bibitem[41]{IV-41}
A.~Grothendieck, \emph{Séminaire de Géométrie algébrique, 1962,
Cohomologie locale des faisceaux cohérents, et théorèmes de Lefschetz locaux
et globaux} (Institut des Hautes Études scientifiques, Bures-sur-Yvette).

\bibitem[42]{IV-42}
M.~Demazure et A.~Grothendieck, \emph{Séminaire de Géométrie algébrique,
1963--1964, Schémas en groupes} (Institut des Hautes Études scientifiques,
Bures-sur-Yvette).

\bibitem[43]{IV-43}
M.~Artin et A.~Grothendieck, \emph{Séminaire de Géométrie algébrique
1963--1964, Cohomologie étale des schémas} (Institut des Hautes Études
scientifiques, Bures-sur-Yvette).

\bibitem[44]{IV-44}
C.~S.~Seshadri, \emph{Quotient space by an abelian variety},
\emph{Mathematische Annalen}, t.~152 (1962), p.~185--194.

\end{thebibliography}

% FIN SOURCE ega4/ega4-backmatter-bibliography-fr.tex

%
% DEBUT SOURCE ega4/ega4-backmatter-notation-fr.tex
% Transcription diplomatique française de EGA IV, quatrième partie, pp. 334--335.
\clearpage
\oldpage[IV]{334}
\phantomsection
\addcontentsline{toc}{section}{Index des notations}
\section*{INDEX DES NOTATIONS}
\markboth{INDEX DES NOTATIONS}{INDEX DES NOTATIONS}

\begingroup
\small
\setlength{\parindent}{0pt}
\setlength{\parskip}{1.5pt}

$Y_f^{(n)}$, $Y^{(n)}$, $\mathcal{G}r_{\bullet}(f)$,
$\mathcal{G}r_n(f)$, $\mathcal{N}_{Y/X}$
($f:Y\to X$ morphisme d'espaces annelés) : 16.1.2.

$\mathcal{O}_{Y^{(\infty)}}$ : 16.1.11.

$\mathcal{P}^{n}_{f}$, $\mathcal{P}^{n}_{X/S}$,
$\mathcal{P}^{\infty}_{f}$, $\mathcal{P}^{\infty}_{X/S}$,
$\Omega^{1}_{f}$, $\Omega^{1}_{X/S}$
($f:X\to S$ morphisme de préschémas) : 16.3.1.

$d^{n}_{f}$, $d^{n}_{X/S}$, $d^{\infty}_{f}$,
$d^{\infty}_{X/S}$, $d^{n}$, $d^{\infty}$, $dt$, $d_{X/S}(t)$
($f:X\to S$ morphisme de préschémas) : 16.3.6.

$P^{n}(u)$, $\Gr_{\bullet}(u)$, $\Gr_{1}(u)$
($u:X'\to X$ morphisme de préschémas) : 16.4.3.

$\operatorname{Der}_{S}(\mathcal{O}_{X},\mathcal{F})$ : 16.5.1.

$\mathcal{D}er_{S}(\mathcal{O}_{X},\mathcal{F})$ : 16.5.4.

$\mathcal{T}_{X/S}$ : 16.5.7.

$D_i$, $\partial/\partial s_i$ : 16.5.8.

$T_{X/S}$ : 16.5.12.

$T_{X/S}(x)$, $T_Y(s)$ : 16.5.13.

$\Omega^p_{X/S}$, $\Omega^\bullet_{X/S}$, $\Omega^p_{B/A}$ : 16.6.1.

$d$, $d_{X/S}$ : 16.6.3.

$\mathcal{P}^n_{X/S}(\mathcal{F})$ : 16.7.2.

$d^n_{X/S,\mathcal{F}}$, $d^n_{X/S}$ : 16.7.5.

$\mathcal{P}^{\infty}_{X/S}(\mathcal{F})$,
$d^{\infty}_{X/S,\mathcal{F}}$ : 16.7.7.

$\operatorname{Diff}^{n}_{X/S}(\mathcal{F},\mathcal{G})$,
$\operatorname{Diff}^{n}_{X/S}$,
$\mathcal{D}iff^{n}_{X/S}(\mathcal{F},\mathcal{G})$ : 16.8.3.

$\mathcal{D}iff^{n}_{X/S}$, $\langle t,D\rangle$ : 16.8.4.

$aD$ ($D$ opérateur différentiel, $a\in\Gamma(U,\mathcal{O}_X)$) : 16.8.5.

$\operatorname{Diff}_{X/S}(\mathcal{F},\mathcal{G})$,
$\mathcal{D}iff_{X/S}(\mathcal{F},\mathcal{G})$ : 16.8.7.

$\mathcal{D}iff_{X/S}$ : 16.8.10.

$|\mathbf p|$, $\mathbf p!$, $\binom{\mathbf p}{\mathbf q}$,
$\mathbf z^{\mathbf p}$, $\boldsymbol\zeta^{\mathbf p}$
($\mathbf p,\mathbf q$ multi-indices) : 16.11.1.

$\dim_x(f)$ ($f:X\to Y$ morphisme localement de type fini, $x\in X$) : 17.10.1.

$\operatorname{Tr}_{B/A}$, $\operatorname{astr}_{B/A}$
($B$ $A$-algèbre finie et libre) : 18.2.1.

$\operatorname{Tr}_{\mathcal B/\mathcal O_X,U}$,
$\operatorname{astr}_{\mathcal B/\mathcal O_X,U}$,
$\operatorname{Tr}_{\mathcal B/\mathcal O_X}$,
$\operatorname{astr}_{\mathcal B/\mathcal O_X}$
($X$ espace annelé, $\mathcal B$ $\mathcal O_X$-Algèbre localement libre de rang fini) : 18.2.2.

$\operatorname{Tr}_f$ ($f$ morphisme fini et étale) : 18.2.6.

$d_{X/Y}$, $\mathcal D_{X/Y}$ : 18.2.7.

${}^{t}V_{\mathcal E}$ : 18.5.1.

$Al(X,\mathcal E,\mathcal J)$, $Al(\mathcal E,\mathcal J)$ : 18.5.2.

$Of(S)$, $Of(X)$, $Id(\mathcal B')$ : 18.5.3.

$\operatorname{Hom.loc}(A,B)$ : 18.6.1.

${}^{h}A$ : 18.6.5 et 18.6.7.

$\operatorname{Hom.sloc}(A,B)$ : 18.6.7.

${}^{hs}A_{(i)}$, ${}^{hs}A$ : 18.8.7 et 18.8.9.

$\operatorname{codim}^{t}_{y}(Y,X)$ : 19.1.3.

$\operatorname{prof}_{T,t}(\mathcal F)$ : 19.9.1.

$\mathcal O_X[\mathcal S^{-1}]$, $\mathcal F[\mathcal S^{-1}]$ : 20.1.1.

$\mathcal S(\mathcal O_X)$, $\mathcal M_X$, $M(X)$,
$\mathcal M_X(\mathcal F)$, $M(X,\mathcal F)$ : 20.1.3.

$\operatorname{dom}(\varphi)$ (fonction méromorphe) : 20.1.4.

$\operatorname{dom}(u)$ ($u$ section méromorphe de $\mathcal F$) : 20.1.7.

$(\mathcal M_X(\mathcal L))^{*}$
($\mathcal L$ $\mathcal O_X$-Module inversible) : 20.1.8.

$s^{-1}$ ($s$ section méromorphe d'un $\mathcal O_X$-Module inversible) : 20.1.10.

$\mathcal S_f(U)$, $\mathcal S_f$, $\mathcal M_f$
($f$ morphisme d'espaces annelés), $\varphi\circ f$ ($\varphi$ fonction méromorphe),
$u\circ f$ ($u$ section méromorphe de $\mathcal F$), $\mathcal M_f(\mathcal F)$ : 20.1.11.

\oldpage[IV]{335}
$\operatorname{Ps.hom}(X,Y)$, $\operatorname{Ps.hom}_S(X,Y)$ : 20.2.1.

$\omega|U$ ($\omega$ pseudo-morphisme),
$\mathcal P\!s.\operatorname{hom}(X,Y)$,
$\mathcal P\!s.\operatorname{hom}_S(X,Y)$ : 20.2.2.

$\operatorname{dom}_S(\omega)$, $\operatorname{dom}(\omega)$
($\omega$ pseudo-morphisme) : 20.2.3.

$\mathcal M'_X$, $M'(X)$ : 20.2.8.

$f\circ\omega$ ($f$ morphisme, $\omega$ pseudo-morphisme) : 20.3.1.

$\omega\circ f$ ($f$ morphisme, $\omega$ pseudo-morphisme) : 20.3.2.

$\operatorname{Ps.hom}_S(X,Y)^f$,
$\mathcal P\!s.\operatorname{hom}_S(X,Y)^f$ : 20.3.2.

$\Gamma_\omega$ ($\omega$ pseudo-morphisme) : 20.4.1.

$p^{-1}$ ($p:\Gamma_\omega\to X$ restriction de la première projection) : 20.4.2.

$\omega(M)$ ($\omega$ pseudo-morphisme, $M$ partie de $X$) : 20.4.2.

$\operatorname{Ps.hom}_{X/S}(X,Y)$ : 20.5.1.

$\mathcal P\!s.\operatorname{hom}_{X/S}(X,Y)$ : 20.5.3.

$\mathcal M'_{X/S}$, $M'(X/S)$ : 20.5.4.

$\mathcal S_{X/S}(U)$, $\mathcal S_{X/S}$, $\mathcal M_{X/S}$,
$M(X/S)$ : 20.6.1.

$\mathcal M_{X/S}(\mathcal F)$, $M(X/S,\mathcal F)$ : 20.6.2.

$\mathcal S_f(U)$, $\mathcal S_f$, $\mathcal M_{X/S,f}$ : 20.6.7.

$\mathcal{D}iv_X$, $\operatorname{Div}(X)$, $\operatorname{div}(f)$,
$\operatorname{div}_X(f)$ ($X$ espace annelé) : 21.1.2.

$\operatorname{Supp}(D)$ ($D$ diviseur) : 21.1.2.

$\operatorname{Div}(A)$ ($A$ anneau) : 21.2.1.

$\operatorname{div}(s)$ ($s$ section d'un Module inversible) : 21.1.4.

$\mathcal{D}iv_X^+$, $\operatorname{Div}^+(X)$, $D\leq D'$
($D,D'$ diviseurs) : 21.1.6.

$\operatorname{Id.inv}(X)$, $\sh{J}d.inv_X$ : 21.2.4.

$\mathcal J_U$, $\mathcal J_X(D)$ : 21.2.5.

$\operatorname{div}(\mathcal I)$ ($\mathcal I$ Idéal fractionnaire) : 21.2.6.

$\mathcal O_X(D)$ ($D$ diviseur) : 21.2.8.

$s_D$ ($D$ diviseur) : 21.2.9.

$D(X)$, $\operatorname{cl}(\mathcal L,s)$ : 21.2.10.

$Y(D)$ ($D$ diviseur positif) : 21.2.12.

$\operatorname{Div.princ}(X)$ : 21.3.1.

$\operatorname{Pic}(X)$, $\operatorname{cl}(\mathcal L)$, $l$, $l_X$,
$\operatorname{Pic}(u)$ : 21.3.2.

$f^*(D)$ ($f$ morphisme d'espaces annelés, $D$ diviseur),
$\operatorname{Div}^f(X)$ : 21.4.2.

$\mathcal M_f^{**}$, $\mathcal{D}iv_X^f$ : 21.4.3.

$N_{X'/X}(u')$, $N(u')$
($u'$ section méromorphe d'un $\mathcal O_{X'}$-Module inversible) : 21.5.3.

$f_*(D')$, $N_{X'/X}(D')$
($f$ morphisme fini, $D'$ diviseur sur $X'$) : 21.5.5.

$\mathfrak I(X)$, $\mathfrak Z(X)$, $\mathfrak Z^+(X)$
($X$ préschéma localement noethérien) : 21.6.1.

$\operatorname{mult}_x(Z)$, $\operatorname{Supp}(Z)$,
$\dim(Z)$, $\operatorname{codim}(Z)$ ($Z$ cycle) : 21.6.1.

$X^{(p)}$, $\mathfrak Z^p(X)$, $\mathfrak Z^{p+}(X)$ : 21.6.2.

$Z|U$ ($Z$ cycle), $\mathscr Z_X$, $\mathscr Z_X^p$,
$\mathscr Z_X^+$, $\mathscr Z_X^{p+}$ : 21.6.3.

$c$ : 21.6.4.

$\operatorname{cyc}(D)$ ($D$ diviseur) : 21.6.5 et 21.6.7.

$\operatorname{mult}_x(D)$, $\operatorname{mult}_{\{x\}}(D)$
($D$ diviseur), $o_x(f)$ ($f$ fonction méromorphe régulière) : 21.6.7.

$\mathrm Z^+(f)$, $\mathrm Z^-(f)$, $\mathfrak Z^1_{\mathrm{princ}}(X)$,
$\mathfrak{Cl}(X)$ : 21.6.7.

$Y(C)$, $Y'(C)$, $\sh{J}_X(C)$, $\sh{J}(C)$
($C$ cycle 1-codimensionnel positif) : 21.7.1.

$f^*(Z)$ ($Z$ cycle 1-codimensionnel, $f$ morphisme) : 21.10.3 et 21.10.11.

$\Gamma_T(\mathscr Z_X^1)$ : 21.10.3.

$f_*(Z')$ ($Z'$ cycle 1-codimensionnel, $f$ morphisme fini) : 21.10.14.

$\Lambda^{\max}\mathcal E$ ($\mathcal E$ module localement libre) : 21.11.1.

$U'$, $\operatorname{Aff}(U/X)$ : 21.12.1.

$\operatorname{Daf}(U/X)$ : 21.12.5.

$\mathcal{D}iv_{X/S}$, $\operatorname{Div}(X/S)$ : 21.15.2.

$\mathcal{D}iv_{X/S}^+$, $\operatorname{Div}^+(X/S)$ : 21.15.3.

$\mathcal O_{X/S}^{\mathrm{reg}}$ : 21.15.3.

$\operatorname{Div}_{X/S}(S')$, $\operatorname{Div}_{X/S}^+(S')$
($S'$ $S$-préschéma) : 21.15.9.

\endgroup

% FIN SOURCE ega4/ega4-backmatter-notation-fr.tex

%
% DEBUT SOURCE ega4/ega4-backmatter-terminology-fr.tex
% Transcription diplomatique française de EGA IV, quatrième partie, pp. 336--339.
% Les entrées et leurs renvois reproduisent l'index imprimé.
\clearpage
\oldpage[IV]{336}
\phantomsection
\addcontentsline{toc}{section}{Index terminologique}
\section*{INDEX TERMINOLOGIQUE}
\markboth{INDEX TERMINOLOGIQUE}{INDEX TERMINOLOGIQUE}

\begingroup
\small
\setlength{\parindent}{0pt}
\setlength{\parskip}{2pt}
\newcommand{\EGAIVterm}[2]{\textbf{#1}: #2.\par}
\newcommand{\EGAIVidxref}[1]{\hyperref[IV.#1]{#1}}

\EGAIVterm{Algèbre étale}{{\hyperlink{ega4.c04.subsec.17-3-2}{17.3.2}}}
\EGAIVterm{Algèbre essentiellement étale}{{\hyperlink{ega4.c04.statement.18-6-1}{18.6.1}}}
\EGAIVterm{$\sh{O}_X$-Algèbre finie et étale}{{\hyperlink{ega4.c04.prop.18-2-3}{18.2.3}}}
\EGAIVterm{Algèbre lisse}{{\hyperlink{ega4.c04.subsec.17-3-2}{17.3.2}}}
\EGAIVterm{Algèbre non ramifiée}{{\hyperlink{ega4.c04.subsec.17-3-2}{17.3.2}}}
\EGAIVterm{$K$-algèbre non ramifiée sur $Y$, sur $A$}{{\hyperlink{ega4.c04.rem.18-10-11}{18.10.11}}}
\EGAIVterm{Algèbre strictement essentiellement étale}{{\hyperlink{ega4.c04.lem.18-6-2}{18.6.2}}}
\EGAIVterm{Anneau d'intersection complète, anneau d'intersection complète absolue}{{\hyperlink{ega4.c04.def.19-3-1}{19.3.1}}}
\EGAIVterm{Anneau hensélien}{{\hyperlink{ega4.c04.def.18-5-8}{18.5.8}}}
\EGAIVterm{Anneau local parafactoriel}{{\hyperlink{ega4.c04.def.21-13-7}{21.13.7}}}
\EGAIVterm{Anneau strictement local}{{\hyperlink{ega4.c04.def.18-8-2}{18.8.2}}}
\EGAIVterm{Application linéaire tangente}{{\hyperlink{ega4.c04.statement.16-5-13}{16.5.13}}}
\EGAIVterm{Application rationnelle stricte}{{\hyperlink{ega4.c04.statement.20-2-1}{20.2.1}}}
\EGAIVterm{Codimension d'un cycle}{{\hyperlink{ega4.c04.statement.21-6-1}{21.6.1}}}
\EGAIVterm{Codimension transversale en un point d'un sous-préschéma}{{\hyperlink{ega4.c04.def.19-1-3}{19.1.3}}}
\EGAIVterm{Codimension transversale d'un morphisme d'immersion en un point}{{\hyperlink{ega4.c04.def.19-1-3}{19.1.3}}}
\EGAIVterm{Couple de morphismes transversaux en un point}{{\hyperlink{ega4.c04.cor.17-13-6}{17.13.6}}}
\EGAIVterm{Couple hensélien}{{\hyperlink{ega4.c04.def.18-5-5}{18.5.5}}}
\EGAIVterm{Couple parafactoriel}{{\hyperlink{ega4.c04.def.21-13-1}{21.13.1}}}
\EGAIVterm{Cycle}{{\hyperlink{ega4.c04.statement.21-6-1}{21.6.1}}}
\EGAIVterm{Cycle 1-codimensionnel à coefficients rationnels}{{\hyperlink{ega4.c04.statement.21-10-9}{21.10.9}}}
\EGAIVterm{Cycle $p$-codimensionnel}{{\hyperlink{ega4.c04.statement.21-6-2}{21.6.2}}}
\EGAIVterm{Cycle des pôles, cycle des zéros d'une fonction méromorphe}{{\hyperlink{ega4.c04.statement.21-6-7}{21.6.7}}}
\EGAIVterm{Cycle linéairement équivalent à 0}{{\hyperlink{ega4.c04.statement.21-6-7}{21.6.7}}}
\EGAIVterm{Cycle localement principal}{{\hyperlink{ega4.c04.statement.21-6-7}{21.6.7}}}
\EGAIVterm{Cycle polaire d'une fonction méromorphe}{{\hyperlink{ega4.c04.statement.21-6-7}{21.6.7}}}
\EGAIVterm{Cycle premier}{{\hyperlink{ega4.c04.statement.21-6-1}{21.6.1}}}
\EGAIVterm{Cycle principal}{{\hyperlink{ega4.c04.statement.21-6-7}{21.6.7}}}
\EGAIVterm{Cycle principal en un point}{{\hyperlink{ega4.c04.statement.21-6-7}{21.6.7}}}
\EGAIVterm{Cycle purement de codimension $p$}{{\hyperlink{ega4.c04.statement.21-6-2}{21.6.2}}}
\EGAIVterm{Défaut d'affinité d'un ouvert dans un préschéma}{{\hyperlink{ega4.c04.statement.21-12-5}{21.12.5}}}
\EGAIVterm{$S$-dérivation, $(X/S)$-dérivation, $f$-dérivation}{{\hyperlink{ega4.c04.statement.16-5-1}{16.5.1}}}
\EGAIVterm{Différentielle d'une section de $\sh{O}_X$}{{\hyperlink{ega4.c04.def.16-3-6}{16.3.6}}}
\EGAIVterm{Différentielle extérieure}{{\hyperlink{ega4.c04.def.16-6-3}{16.6.3}}}
\EGAIVterm{Dimension d'un cycle}{{\hyperlink{ega4.c04.statement.21-6-1}{21.6.1}}}
\EGAIVterm{Dimension relative d'un morphisme en un point, d'un $Y$-préschéma en un point}{{\hyperlink{ega4.c04.def.17-10-1}{17.10.1}}}
\EGAIVterm{Discriminant d'un $Y$-préschéma fini et localement libre}{{\hyperlink{ega4.c04.item.18-2-7}{18.2.7}}}
\EGAIVterm{Diviseur sur un espace annelé}{{\hyperlink{ega4.c04.statement.21-1-2}{21.1.2}}}
\EGAIVterm{Diviseur d'une fonction méromorphe}{{\hyperlink{ega4.c04.statement.21-1-2}{21.1.2}}}
\EGAIVterm{Diviseur d'une section méromorphe d'un Module inversible}{{\hyperlink{ega4.c04.statement.21-1-4}{21.1.4}}}
\EGAIVterm{Diviseur positif}{{\hyperlink{ega4.c04.def.21-1-6}{21.1.6}}}
\EGAIVterm{Diviseur principal}{{\hyperlink{ega4.c04.statement.21-3-1}{21.3.1}}}
\EGAIVterm{Diviseur relativement à $S$, diviseur transversal à $f$}{{\hyperlink{ega4.c04.def.21-15-2}{21.15.2}}}
\EGAIVterm{Diviseurs linéairement équivalents}{{\hyperlink{ega4.c04.statement.21-3-1}{21.3.1}}}

\oldpage[IV]{337}
\EGAIVterm{Domaine de définition d'une fonction méromorphe}{{\hyperlink{ega4.c04.statement.20-1-4}{20.1.4}}}
\EGAIVterm{Domaine de définition d'un pseudo-morphisme}{{\hyperlink{ega4.c04.statement.20-2-3}{20.2.3}}}
\EGAIVterm{Domaine de définition d'une section méromorphe d'un Module}{{\hyperlink{ega4.c04.statement.20-1-7}{20.1.7}}}
\EGAIVterm{Enveloppe affine d'un $X$-préschéma}{{\hyperlink{ega4.c04.statement.21-12-1}{21.12.1}}}
\EGAIVterm{Espace tangent à un préschéma en un point}{{\hyperlink{ega4.c04.statement.16-5-13}{16.5.13}}}
\EGAIVterm{Faisceau conormal}{{\hyperlink{ega4.c04.def.16-1-2}{16.1.2}}}
\EGAIVterm{Faisceau d'anneaux de fractions à dénominateurs dans $\sh{S}$}{{\hyperlink{ega4.c04.statement.20-1-1}{20.1.1}}}
\EGAIVterm{Faisceau d'anneaux gradués associé à un morphisme}{{\hyperlink{ega4.c04.def.16-1-2}{16.1.2}}}
\EGAIVterm{Faisceau de modules de fractions à dénominateurs dans $\sh{S}$}{{\hyperlink{ega4.c04.statement.20-1-2}{20.1.2}}}
\EGAIVterm{Faisceau des $S$-dérivations de $\sh{O}_X$}{{\hyperlink{ega4.c04.statement.16-5-7}{16.5.7}}}
\EGAIVterm{Faisceau des $S$-dérivations de $\sh{O}_X$ dans $\sh{F}$}{{\hyperlink{ega4.c04.statement.16-5-4}{16.5.4}}}
\EGAIVterm{Faisceau des $p$-différentielles de $X$ relativement à $S$}{{\hyperlink{ega4.c04.statement.16-6-1}{16.6.1}}}
\EGAIVterm{Faisceau des diviseurs sur un espace annelé}{{\hyperlink{ega4.c04.statement.21-1-2}{21.1.2}}}
\EGAIVterm{Faisceau des diviseurs relativement à $S$, ou transversaux à $f$}{{\hyperlink{ega4.c04.def.21-15-2}{21.15.2}}}
\EGAIVterm{Faisceau des germes de cycles 1-codimensionnels à coefficients dans $\Lambda$}{{\hyperlink{ega4.c04.statement.21-10-9}{21.10.9}}}
\EGAIVterm{Faisceau des germes de fonctions méromorphes sur $X$}{{\hyperlink{ega4.c04.statement.20-1-4}{20.1.4}}}
\EGAIVterm{Faisceau des germes de fonctions méromorphes relativement à $S$}{{\hyperlink{ega4.c04.statement.20-6-1}{20.6.1}}}
\EGAIVterm{Faisceau des germes de sections méromorphes d'un Module}{{\hyperlink{ega4.c04.statement.20-1-3}{20.1.3}}}
\EGAIVterm{Faisceau des parties principales d'ordre $n$ du $S$-préschéma $X$}{{\hyperlink{ega4.c04.def.16-3-1}{16.3.1}}}
\EGAIVterm{Faisceau formellement principal homogène (ou pseudo-torseur) sous un faisceau de groupes}{{\hyperlink{ega4.c04.statement.16-5-15}{16.5.15}}}
\EGAIVterm{Faisceau principal homogène (ou torseur) sous un faisceau de groupes}{{\hyperlink{ega4.c04.statement.16-5-15}{16.5.15}}}
\EGAIVterm{Faisceau principal homogène (ou torseur) trivial sous un faisceau de groupes}{{\hyperlink{ega4.c04.statement.16-5-15}{16.5.15}}}
\EGAIVterm{Faisceau tangent de $X$ relativement à $S$}{{\hyperlink{ega4.c04.statement.16-5-7}{16.5.7}}}
\EGAIVterm{Famille de diviseurs sur $X$ relativement à $S$, paramétrés par $S'$}{{\hyperlink{ega4.c04.statement.21-15-9}{21.15.9}}}
\EGAIVterm{Famille de morphismes transversaux en un point}{{\hyperlink{ega4.c04.cor.17-13-10}{17.13.10}}}
\EGAIVterm{Famille de sous-préschémas se coupant transversalement en un point}{{\hyperlink{ega4.c04.statement.17-13-11}{17.13.11}}}
\EGAIVterm{Fibré tangent de $X$ relativement à $S$}{{\hyperlink{ega4.c04.statement.16-5-12}{16.5.12}}}
\EGAIVterm{Fonction méromorphe sur un espace annelé}{{\hyperlink{ega4.c04.statement.20-1-3}{20.1.3}}}
\EGAIVterm{Fonction méromorphe définie en un point}{{\hyperlink{ega4.c04.statement.20-1-4}{20.1.4}}}
\EGAIVterm{Fonction méromorphe régulière}{{\hyperlink{ega4.c04.statement.20-1-8}{20.1.8}}}
\EGAIVterm{Fonction méromorphe régulière relativement à $S$}{{\hyperlink{ega4.c04.statement.20-6-5}{20.6.5}}}
\EGAIVterm{Fonction méromorphe relativement à $S$}{{\hyperlink{ega4.c04.statement.20-6-1}{20.6.1}}}
\EGAIVterm{Forme trace}{{\hyperlink{ega4.c04.statement.18-2-1}{18.2.1}}}
\EGAIVterm{Graphe d'un pseudo-$S$-morphisme}{{\hyperlink{ega4.c04.statement.20-4-1}{20.4.1}}}
\EGAIVterm{Groupe des classes de cycles 1-codimensionnels}{{\hyperlink{ega4.c04.statement.21-6-7}{21.6.7}}}
\EGAIVterm{Hensélisé d'un anneau local}{{\hyperlink{ega4.c04.def.18-6-5}{18.6.5}}}
\EGAIVterm{Hensélisé d'un anneau semi-local}{{\hyperlink{ega4.c04.statement.18-6-7}{18.6.7}}}
\EGAIVterm{Hensélisé strict d'un anneau local}{{\hyperlink{ega4.c04.def.18-8-7}{18.8.7}}}
\EGAIVterm{Hensélisé strict d'un anneau semi-local}{{\hyperlink{ega4.c04.statement.18-8-9}{18.8.9}}}
\EGAIVterm{Homomorphisme associé à la trace}{{\hyperlink{ega4.c04.statement.18-2-2}{18.2.2}}}
\EGAIVterm{Homomorphisme régulier de $\sh{O}_X$-Modules}{{\hyperlink{ega4.c04.prop.19-4-11}{19.4.11}}}
\EGAIVterm{Homomorphisme semi-local}{{\hyperlink{ega4.c04.statement.18-6-7}{18.6.7}}}
\EGAIVterm{Homomorphisme trace}{{\hyperlink{ega4.c04.statement.18-2-2}{18.2.2}}}
\EGAIVterm{Idéal des dénominateurs d'une section méromorphe}{{\hyperlink{ega4.c04.statement.20-2-14}{20.2.14}}}
\EGAIVterm{Idéal discriminant d'un $Y$-préschéma fini et localement libre}{{\hyperlink{ega4.c04.item.18-2-7}{18.2.7}}}
\EGAIVterm{Idéal entier}{{\hyperlink{ega4.c04.statement.21-2-7}{21.2.7}}}
\EGAIVterm{Idéal fractionnaire}{{\hyperlink{ega4.c04.statement.21-2-1}{21.2.1}}}
\EGAIVterm{Idéal fractionnaire inversible}{{\hyperlink{ega4.c04.statement.21-2-1}{21.2.1}}}
\EGAIVterm{Idéal quasi-régulier d'un faisceau d'anneaux}{{\hyperlink{ega4.c04.def.16-9-1}{16.9.1}}}
\EGAIVterm{Idéal quasi-régulier d'un anneau}{{\hyperlink{ega4.c04.statement.16-9-7}{16.9.7}}}
\EGAIVterm{Idéal régulier d'un faisceau d'anneaux}{{\hyperlink{ega4.c04.def.16-9-1}{16.9.1}}}
\EGAIVterm{Idéal régulier d'un anneau}{{\hyperlink{ega4.c04.statement.16-9-7}{16.9.7}}}
\EGAIVterm{Idéal transversalement régulier en un point d'un préschéma}{{\hyperlink{ega4.c04.def.19-2-1}{19.2.1}}}

\oldpage[IV]{338}
\EGAIVterm{Idéal transversalement régulier d'une algèbre}{{\hyperlink{ega4.c04.def.19-2-1}{19.2.1}}}
\EGAIVterm{Image directe d'un cycle 1-codimensionnel}{{\hyperlink{ega4.c04.statement.21-10-14}{21.10.14}}}
\EGAIVterm{Image directe d'un diviseur}{{\hyperlink{ega4.c04.def.21-5-5}{21.5.5}}}
\EGAIVterm{Image réciproque d'un cycle 1-codimensionnel}{{\hyperlink{ega4.c04.statement.21-10-3}{21.10.3}} et {\hyperlink{ega4.c04.statement.21-10-11}{21.10.11}}}
\EGAIVterm{Image réciproque d'un diviseur}{{\hyperlink{ega4.c04.def.21-4-2}{21.4.2}}}
\EGAIVterm{Image réciproque d'une fonction méromorphe}{{\hyperlink{ega4.c04.statement.20-1-11}{20.1.11}}}
\EGAIVterm{Image réciproque d'une fonction méromorphe relative à $S$}{{\hyperlink{ega4.c04.statement.20-6-7}{20.6.7}}}
\EGAIVterm{Immersion quasi-régulière}{{\hyperlink{ega4.c04.def.16-9-2}{16.9.2}}}
\EGAIVterm{Immersion régulière}{{\hyperlink{ega4.c04.def.16-9-2}{16.9.2}}}
\EGAIVterm{Immersion régulière en un point}{{\hyperlink{ega4.c04.prop.16-9-10}{16.9.10}}}
\EGAIVterm{Immersion transversalement régulière en un point relativement à $S$, immersion transversalement régulière relativement à $S$}{{\hyperlink{ega4.c04.def.19-2-2}{19.2.2}}}
\EGAIVterm{Intersection complète en un point relativement à $S$, intersection complète relativement à $S$ (pour un préschéma plat et localement de présentation finie)}{{\hyperlink{ega4.c04.def.19-3-6}{19.3.6}}}
\EGAIVterm{$n$-ème invariant normal d'un morphisme}{{\hyperlink{ega4.c04.def.16-1-2}{16.1.2}}}
\EGAIVterm{Invariant normal d'ordre infini d'un morphisme}{{\hyperlink{ega4.c04.rem.16-1-11}{16.1.11}}}
\EGAIVterm{$\sh{O}_X$-Module des 1-différentielles de $f$ (ou de $X$ par rapport à $S$, ou du $S$-préschéma $X$)}{{\hyperlink{ega4.c04.def.16-3-1}{16.3.1}}}
\EGAIVterm{$\sh{O}_X$-Module sans torsion relativement à $S$}{{\hyperlink{ega4.c04.statement.20-6-2}{20.6.2}}}
\EGAIVterm{$\sh{O}_X$-Module sans torsion}{{\hyperlink{ega4.c04.statement.20-1-5}{20.1.5}}}
\EGAIVterm{Morphisme différentiellement lisse}{{\hyperlink{ega4.c04.def.16-10-1}{16.10.1}}}
\EGAIVterm{Morphisme différentiellement lisse jusqu'à l'ordre $r$}{{\hyperlink{ega4.c04.statement.16-11-3}{16.11.3}}}
\EGAIVterm{Morphisme (plat) d'intersection complète}{{\hyperlink{ega4.c04.def.19-3-6}{19.3.6}}}
\EGAIVterm{Morphisme essentiellement propre}{{\hyperlink{ega4.c04.rem.18-10-20}{18.10.20}}}
\EGAIVterm{Morphisme étale}{{\hyperlink{ega4.c04.def.17-3-1}{17.3.1}}}
\EGAIVterm{Morphisme étale en un point}{{\hyperlink{ega4.c04.def.17-3-7}{17.3.7}}}
\EGAIVterm{Morphisme fini localement libre de rang $n$, morphisme fini localement libre}{{\hyperlink{ega4.c04.item.18-2-7}{18.2.7}}}
\EGAIVterm{Morphisme formellement étale, formellement lisse, formellement net, formellement non ramifié}{{\hyperlink{ega4.c04.def.17-1-1}{17.1.1}}}
\EGAIVterm{Morphisme lisse}{{\hyperlink{ega4.c04.def.17-3-1}{17.3.1}}}
\EGAIVterm{Morphisme lisse en un point}{{\hyperlink{ega4.c04.def.17-3-7}{17.3.7}}}
\EGAIVterm{Morphisme non ramifié, morphisme net}{{\hyperlink{ega4.c04.def.17-3-1}{17.3.1}}}
\EGAIVterm{Morphisme non ramifié en un point}{{\hyperlink{ega4.c04.def.17-3-7}{17.3.7}}}
\EGAIVterm{Morphisme transversal à $Y$ au point $x$ relativement à $S$}{{\hyperlink{ega4.c04.def.17-13-3}{17.13.3}}}
\EGAIVterm{Morphismes équivalents}{{\hyperlink{ega4.c04.statement.20-2-1}{20.2.1}}}
\EGAIVterm{Multiplicité d'un cycle en un point}{{\hyperlink{ega4.c04.statement.21-6-1}{21.6.1}}}
\EGAIVterm{Norme d'un diviseur}{{\hyperlink{ega4.c04.def.21-5-5}{21.5.5}}}
\EGAIVterm{Norme d'une section méromorphe}{{\hyperlink{ega4.c04.statement.21-5-3}{21.5.3}}}
\EGAIVterm{Opérateur différentiel relativement à $S$}{{\hyperlink{ega4.c04.def.16-8-1}{16.8.1}}}
\EGAIVterm{Ordre d'un opérateur différentiel}{{\hyperlink{ega4.c04.def.16-8-1}{16.8.1}}}
\EGAIVterm{Ordre d'une fonction méromorphe en un point $x\in X^{(1)}$}{{\hyperlink{ega4.c04.statement.21-6-7}{21.6.7}}}
\EGAIVterm{Partie principale d'ordre $n$, d'ordre infini, d'une section de $\sh{O}_X$}{{\hyperlink{ega4.c04.def.16-3-6}{16.3.6}}}
\EGAIVterm{Partie fermée purement de codimension $p$}{{\hyperlink{ega4.c04.statement.21-6-2}{21.6.2}}}
\EGAIVterm{Polynôme séparable sur un anneau}{{\hyperlink{ega4.c04.cor.18-4-3}{18.4.3}}}
\EGAIVterm{Préschéma de caractéristique $p$}{{\hyperlink{ega4.c04.statement.16-12-1}{16.12.1}}}
\EGAIVterm{Préschéma des coïncidences de deux morphismes}{{\hyperlink{ega4.c04.statement.17-4-5}{17.4.5}}}
\EGAIVterm{Préschéma différentiellement lisse sur $S$}{{\hyperlink{ega4.c04.def.16-10-1}{16.10.1}}}
\EGAIVterm{Préschéma étale sur $S$}{{\hyperlink{ega4.c04.def.17-3-1}{17.3.1}}}
\EGAIVterm{Préschéma étale sur $S$ en un point}{{\hyperlink{ega4.c04.def.17-3-7}{17.3.7}}}
\EGAIVterm{Préschéma formellement étale, formellement lisse, formellement non ramifié sur $S$}{{\hyperlink{ega4.c04.def.17-1-1}{17.1.1}}}
\EGAIVterm{Préschéma intersection complète absolue en un point}{{\hyperlink{ega4.c04.def.19-3-1}{19.3.1}}}
\EGAIVterm{Préschéma lisse sur $S$}{{\hyperlink{ega4.c04.def.17-3-1}{17.3.1}}}
\EGAIVterm{Préschéma lisse sur $S$ en un point}{{\hyperlink{ega4.c04.def.17-3-7}{17.3.7}}}
\EGAIVterm{Préschéma localement factoriel}{{\hyperlink{ega4.c04.thm.21-6-9}{21.6.9}}}
\EGAIVterm{Préschéma net sur $S$, non ramifié sur $S$}{{\hyperlink{ega4.c04.def.17-3-1}{17.3.1}}}

\oldpage[IV]{339}
\EGAIVterm{Préschéma net sur $S$ en un point, non ramifié sur $S$ en un point}{{\hyperlink{ega4.c04.def.17-3-7}{17.3.7}}}
\EGAIVterm{$T$-profondeur d'un $\sh{O}_X$-Module cohérent en un point}{{\hyperlink{ega4.c04.statement.19-9-1}{19.9.1}}}
\EGAIVterm{Pseudo-fonction}{{\hyperlink{ega4.c04.statement.20-2-8}{20.2.8}}}
\EGAIVterm{Pseudo-fonction relativement à $S$}{{\hyperlink{ega4.c04.statement.20-5-4}{20.5.4}}}
\EGAIVterm{Pseudo-morphisme, pseudo-$S$-morphisme}{{\hyperlink{ega4.c04.statement.20-2-1}{20.2.1}}}
\EGAIVterm{Pseudo-morphisme composé d'un morphisme et d'un pseudo-morphisme}{{\hyperlink{ega4.c04.statement.20-3-1}{20.3.1}}}
\EGAIVterm{Pseudo-morphisme composé d'un pseudo-morphisme et d'un morphisme}{{\hyperlink{ega4.c04.statement.20-3-2}{20.3.2}}}
\EGAIVterm{Pseudo-morphisme défini en un point}{{\hyperlink{ega4.c04.statement.20-2-3}{20.2.3}}}
\EGAIVterm{Pseudo-morphisme relativement à $S$}{{\hyperlink{ega4.c04.statement.20-5-1}{20.5.1}}}
\EGAIVterm{Pseudo-torseur}{{\hyperlink{ega4.c04.statement.16-5-15}{16.5.15}}}
\EGAIVterm{Restriction d'un cycle à un ouvert}{{\hyperlink{ega4.c04.statement.21-6-3}{21.6.3}}}
\EGAIVterm{Restriction d'un pseudo-morphisme à un ouvert}{{\hyperlink{ega4.c04.statement.20-2-2}{20.2.2}}}
\EGAIVterm{Restriction d'un pseudo-morphisme à un schéma local}{{\hyperlink{ega4.c04.statement.20-3-6}{20.3.6}}}
\EGAIVterm{Revêtement étale, revêtement étale localement trivial, revêtement étale trivial}{{\hyperlink{ega4.c04.item.18-2-7}{18.2.7}}}
\EGAIVterm{Schéma local hensélien}{{\hyperlink{ega4.c04.def.18-5-8}{18.5.8}}}
\EGAIVterm{Schéma strictement local}{{\hyperlink{ega4.c04.def.18-8-2}{18.8.2}}}
\EGAIVterm{Section méromorphe d'un $\sh{O}_X$-Module}{{\hyperlink{ega4.c04.statement.20-1-3}{20.1.3}}}
\EGAIVterm{Section méromorphe d'un $\sh{O}_X$-Module définie en un point}{{\hyperlink{ega4.c04.statement.20-1-7}{20.1.7}}}
\EGAIVterm{Section méromorphe régulière d'un $\sh{O}_X$-Module inversible}{{\hyperlink{ega4.c04.statement.20-1-8}{20.1.8}}}
\EGAIVterm{Section méromorphe d'un $\sh{O}_X$-Module relativement à $S$}{{\hyperlink{ega4.c04.statement.20-6-2}{20.6.2}}}
\EGAIVterm{Section méromorphe d'un $\sh{O}_X$-Module inversible régulière relativement à $S$}{{\hyperlink{ega4.c04.statement.20-6-5}{20.6.5}}}
\EGAIVterm{Sous-préschéma quasi-régulièrement immergé, régulièrement immergé}{{\hyperlink{ega4.c04.def.16-9-2}{16.9.2}}}
\EGAIVterm{Sous-préschémas se coupant transversalement en un point relativement à $S$}{{\hyperlink{ega4.c04.statement.17-13-7}{17.13.7}}}
\EGAIVterm{Support d'un cycle}{{\hyperlink{ega4.c04.statement.21-6-1}{21.6.1}}}
\EGAIVterm{Support d'un diviseur}{{\hyperlink{ega4.c04.statement.21-1-2}{21.1.2}}}
\EGAIVterm{Suite de sections transversalement régulière relativement à $S$}{{\hyperlink{ega4.c04.cor.19-1-2}{19.1.2}}}
\EGAIVterm{Symétrie canonique de $X\times_S X$}{{\hyperlink{ega4.c04.statement.16-3-3}{16.3.3}}}
\EGAIVterm{Symétrie canonique de $\sh{P}^{n}_{X/S}$, de $\sh{P}^{\infty}_{X/S}$}{{\hyperlink{ega4.c04.prop.16-3-4}{16.3.4}}}
\EGAIVterm{Système quasi-régulier (régulier) de générateurs d'un Idéal}{{\hyperlink{ega4.c04.def.16-9-1}{16.9.1}}}
\EGAIVterm{Torseur}{{\hyperlink{ega4.c04.statement.16-5-15}{16.5.15}}}
\EGAIVterm{Voisinage formel d'un sous-préschéma}{{\hyperlink{ega4.c04.rem.16-1-11}{16.1.11}}}
\EGAIVterm{$n$-ème voisinage infinitésimal}{{\hyperlink{ega4.c04.def.16-1-2}{16.1.2}}}

\endgroup

% FIN SOURCE ega4/ega4-backmatter-terminology-fr.tex

%
% DEBUT SOURCE ega4/ega4-backmatter-contents-fr.tex
% Transcription diplomatique française de EGA IV, quatrième partie, pp. 341--343.
\clearpage
\oldpage[IV]{341}
\phantomsection
\addcontentsline{toc}{section}{Table des matières}
\section*{TABLE DES MATIÈRES}
\markboth{TABLE DES MATIÈRES}{TABLE DES MATIÈRES}

\begingroup
\small
\newcommand{\EGAIVcontentsline}[3]{%
  \par\begingroup
  \rightskip=2.5em\parfillskip=-\rightskip
  \noindent\textbf{#1}\quad #2\nobreak\hspace{0.5em}%
  \dotfill\nobreak\hbox to 2.5em{\hss #3}\par
  \endgroup}

\begin{center}\textbf{CHAPITRE IV. --- Étude locale des schémas et des morphismes de schémas (fin)}\end{center}

\EGAIVcontentsline{\S16}{Invariants différentiels. Morphismes différentiellement lisses}{5}
\EGAIVcontentsline{{\texorpdfstring{\protect\hyperlink{ega4.c04.subsec.16-1}{16.1}}{16.1}}}{Invariants normaux d'une immersion}{5}
\EGAIVcontentsline{{\texorpdfstring{\protect\hyperlink{ega4.c04.subsec.16-2}{16.2}}{16.2}}}{Propriétés fonctorielles des invariants normaux d'une immersion}{9}
\EGAIVcontentsline{{\texorpdfstring{\protect\hyperlink{ega4.c04.subsec.16-3}{16.3}}{16.3}}}{Invariants différentiels fondamentaux d'un morphisme de préschémas}{14}
\EGAIVcontentsline{{\texorpdfstring{\protect\hyperlink{ega4.c04.subsec.16-4}{16.4}}{16.4}}}{Propriétés fonctorielles des invariants différentiels}{16}
\EGAIVcontentsline{{\texorpdfstring{\protect\hyperlink{ega4.c04.subsec.16-5}{16.5}}{16.5}}}{Faisceaux et fibrés tangents relatifs; dérivations}{27}
\EGAIVcontentsline{{\texorpdfstring{\protect\hyperlink{ega4.c04.subsec.16-6}{16.6}}{16.6}}}{Faisceaux de $p$-différentielles et différentielle extérieure}{34}
\EGAIVcontentsline{{\texorpdfstring{\protect\hyperlink{ega4.c04.subsec.16-7}{16.7}}{16.7}}}{Les $\sh{P}^{n}_{X/S}(\sh{F})$}{36}
\EGAIVcontentsline{{\texorpdfstring{\protect\hyperlink{ega4.c04.subsec.16-8}{16.8}}{16.8}}}{Opérateurs différentiels}{39}
\EGAIVcontentsline{{\texorpdfstring{\protect\hyperlink{ega4.c04.subsec.16-9}{16.9}}{16.9}}}{Immersions régulières et quasi-régulières}{46}
\EGAIVcontentsline{{\texorpdfstring{\protect\hyperlink{ega4.c04.subsec.16-10}{16.10}}{16.10}}}{Morphismes différentiellement lisses}{51}
\EGAIVcontentsline{{\texorpdfstring{\protect\hyperlink{ega4.c04.subsec.16-11}{16.11}}{16.11}}}{Opérateurs différentiels sur un $S$-préschéma différentiellement lisse}{53}
\EGAIVcontentsline{{\texorpdfstring{\protect\hyperlink{ega4.c04.subsec.16-12}{16.12}}{16.12}}}{Cas de la caractéristique nulle : critère jacobien pour les morphismes différentiellement lisses}{55}

\EGAIVcontentsline{\S17}{Morphismes lisses, morphismes non ramifiés, morphismes étales}{56}
\EGAIVcontentsline{{\texorpdfstring{\protect\hyperlink{ega4.c04.subsec.17-1}{17.1}}{17.1}}}{Morphismes formellement lisses, morphismes formellement non ramifiés, morphismes formellement étales}{56}
\EGAIVcontentsline{{\texorpdfstring{\protect\hyperlink{ega4.c04.subsec.17-2}{17.2}}{17.2}}}{Propriétés différentielles générales}{59}
\EGAIVcontentsline{{\texorpdfstring{\protect\hyperlink{ega4.c04.subsec.17-3}{17.3}}{17.3}}}{Morphismes lisses, morphismes non ramifiés, morphismes étales}{61}
\EGAIVcontentsline{{\texorpdfstring{\protect\hyperlink{ega4.c04.subsec.17-4}{17.4}}{17.4}}}{Caractérisations des morphismes non ramifiés}{65}
\EGAIVcontentsline{{\texorpdfstring{\protect\hyperlink{ega4.c04.subsec.17-5}{17.5}}{17.5}}}{Caractérisations des morphismes lisses}{67}
\EGAIVcontentsline{{\texorpdfstring{\protect\hyperlink{ega4.c04.subsec.17-6}{17.6}}{17.6}}}{Caractérisations des morphismes étales}{70}
\EGAIVcontentsline{{\texorpdfstring{\protect\hyperlink{ega4.c04.subsec.17-7}{17.7}}{17.7}}}{Propriétés de descente et de passage à la limite}{72}
\EGAIVcontentsline{{\texorpdfstring{\protect\hyperlink{ega4.c04.subsec.17-8}{17.8}}{17.8}}}{Critères de lissité et de non ramification par fibres}{79}
\EGAIVcontentsline{{\texorpdfstring{\protect\hyperlink{ega4.c04.subsec.17-9}{17.9}}{17.9}}}{Morphismes étales et immersions ouvertes}{79}
\EGAIVcontentsline{{\texorpdfstring{\protect\hyperlink{ega4.c04.subsec.17-10}{17.10}}{17.10}}}{Dimension relative d'un préschéma lisse sur un autre}{81}
\EGAIVcontentsline{{\texorpdfstring{\protect\hyperlink{ega4.c04.subsec.17-11}{17.11}}{17.11}}}{Morphismes lisses de préschémas lisses}{82}
\EGAIVcontentsline{{\texorpdfstring{\protect\hyperlink{ega4.c04.subsec.17-12}{17.12}}{17.12}}}{Sous-préschémas lisses d'un préschéma lisse. Morphismes lisses et morphismes différentiellement lisses}{85}

\oldpage[IV]{342}
\EGAIVcontentsline{{\texorpdfstring{\protect\hyperlink{ega4.c04.subsec.17-13}{17.13}}{17.13}}}{Morphismes transversaux}{89}
\EGAIVcontentsline{{\texorpdfstring{\protect\hyperlink{ega4.c04.subsec.17-14}{17.14}}{17.14}}}{Caractérisations locales et infinitésimales des morphismes lisses, des morphismes non ramifiés et des morphismes étales}{98}
\EGAIVcontentsline{{\texorpdfstring{\protect\hyperlink{ega4.c04.subsec.17-15}{17.15}}{17.15}}}{Cas des préschémas sur un corps de base}{99}
\EGAIVcontentsline{{\texorpdfstring{\protect\hyperlink{ega4.c04.subsec.17-16}{17.16}}{17.16}}}{Quasi-sections de morphismes plats ou lisses}{105}

\EGAIVcontentsline{\S18}{Compléments sur les morphismes étales. Anneaux locaux henséliens et anneaux strictement locaux}{109}
\EGAIVcontentsline{{\texorpdfstring{\protect\hyperlink{ega4.c04.subsec.18-1}{18.1}}{18.1}}}{Une équivalence remarquable de catégories}{109}
\EGAIVcontentsline{{\texorpdfstring{\protect\hyperlink{ega4.c04.subsec.18-2}{18.2}}{18.2}}}{Revêtements étales}{111}
\EGAIVcontentsline{{\texorpdfstring{\protect\hyperlink{ega4.c04.subsec.18-3}{18.3}}{18.3}}}{Algèbres finies et étales}{114}
\EGAIVcontentsline{{\texorpdfstring{\protect\hyperlink{ega4.c04.subsec.18-4}{18.4}}{18.4}}}{Structure locale des morphismes non ramifiés et des morphismes étales}{118}
\EGAIVcontentsline{{\texorpdfstring{\protect\hyperlink{ega4.c04.subsec.18-5}{18.5}}{18.5}}}{Anneaux locaux henséliens}{125}
\EGAIVcontentsline{{\texorpdfstring{\protect\hyperlink{ega4.c04.subsec.18-6}{18.6}}{18.6}}}{Hensélisation}{135}
\EGAIVcontentsline{{\texorpdfstring{\protect\hyperlink{ega4.c04.subsec.18-7}{18.7}}{18.7}}}{Hensélisation et anneaux excellents}{142}
\EGAIVcontentsline{{\texorpdfstring{\protect\hyperlink{ega4.c04.subsec.18-8}{18.8}}{18.8}}}{Anneaux strictement locaux et hensélisation stricte}{144}
\EGAIVcontentsline{{\texorpdfstring{\protect\hyperlink{ega4.c04.subsec.18-9}{18.9}}{18.9}}}{Fibres formelles des anneaux noethériens henséliens}{150}
\EGAIVcontentsline{{\texorpdfstring{\protect\hyperlink{ega4.c04.subsec.18-10}{18.10}}{18.10}}}{Préschémas étales sur un préschéma géométriquement unibranche ou normal}{157}
\EGAIVcontentsline{{\texorpdfstring{\protect\hyperlink{ega4.c04.subsec.18-11}{18.11}}{18.11}}}{Application aux algèbres locales noethériennes complètes sur un corps}{169}
\EGAIVcontentsline{{\texorpdfstring{\protect\hyperlink{ega4.c04.subsec.18-12}{18.12}}{18.12}}}{Applications de la localisation étale aux morphismes quasi-finis (généralisations de résultats antérieurs)}{181}

\EGAIVcontentsline{\S19}{Immersions régulières et platitude normale}{185}
\EGAIVcontentsline{{\texorpdfstring{\protect\hyperlink{ega4.c04.subsec.19-1}{19.1}}{19.1}}}{Propriétés des immersions régulières}{185}
\EGAIVcontentsline{{\texorpdfstring{\protect\hyperlink{ega4.c04.subsec.19-2}{19.2}}{19.2}}}{Immersions transversalement régulières}{190}
\EGAIVcontentsline{{\texorpdfstring{\protect\hyperlink{ega4.c04.subsec.19-3}{19.3}}{19.3}}}{Intersections complètes relatives (cas plat)}{194}
\EGAIVcontentsline{{\texorpdfstring{\protect\hyperlink{ega4.c04.subsec.19-4}{19.4}}{19.4}}}{Application : critères de régularité et de lissité pour les préschémas éclatés}{198}
\EGAIVcontentsline{{\texorpdfstring{\protect\hyperlink{ega4.c04.subsec.19-5}{19.5}}{19.5}}}{Critères de $M$-régularité}{204}
\EGAIVcontentsline{{\texorpdfstring{\protect\hyperlink{ega4.c04.subsec.19-6}{19.6}}{19.6}}}{Suites régulières relativement à un module filtré quotient}{209}
\EGAIVcontentsline{{\texorpdfstring{\protect\hyperlink{ega4.c04.subsec.19-7}{19.7}}{19.7}}}{Critère de platitude normale de Hironaka}{212}
\EGAIVcontentsline{{\texorpdfstring{\protect\hyperlink{ega4.c04.subsec.19-8}{19.8}}{19.8}}}{Propriétés de passage à la limite projective}{219}
\EGAIVcontentsline{{\texorpdfstring{\protect\hyperlink{ega4.c04.subsec.19-9}{19.9}}{19.9}}}{Suites $\sh{F}$-régulières et profondeur}{222}

\EGAIVcontentsline{\S20}{Fonctions méromorphes et pseudo-morphismes}{223}
\EGAIVcontentsline{{\texorpdfstring{\protect\hyperlink{ega4.c04.subsec.20-0}{20.0}}{20.0}}}{Introduction}{223}
\EGAIVcontentsline{{\texorpdfstring{\protect\hyperlink{ega4.c04.subsec.20-1}{20.1}}{20.1}}}{Fonctions méromorphes}{224}
\EGAIVcontentsline{{\texorpdfstring{\protect\hyperlink{ega4.c04.subsec.20-2}{20.2}}{20.2}}}{Pseudo-morphismes et pseudo-fonctions}{231}

\oldpage[IV]{343}
\EGAIVcontentsline{{\texorpdfstring{\protect\hyperlink{ega4.c04.subsec.20-3}{20.3}}{20.3}}}{Composition des pseudo-morphismes}{237}
\EGAIVcontentsline{{\texorpdfstring{\protect\hyperlink{ega4.c04.subsec.20-4}{20.4}}{20.4}}}{Propriétés des domaines de définition des fonctions rationnelles}{244}
\EGAIVcontentsline{{\texorpdfstring{\protect\hyperlink{ega4.c04.subsec.20-5}{20.5}}{20.5}}}{Pseudo-morphismes relatifs}{249}
\EGAIVcontentsline{{\texorpdfstring{\protect\hyperlink{ega4.c04.subsec.20-6}{20.6}}{20.6}}}{Fonctions méromorphes relatives}{252}

\EGAIVcontentsline{\S21}{Diviseurs}{255}
\EGAIVcontentsline{{\texorpdfstring{\protect\hyperlink{ega4.c04.subsec.21-1}{21.1}}{21.1}}}{Diviseurs sur un espace annelé}{255}
\EGAIVcontentsline{{\texorpdfstring{\protect\hyperlink{ega4.c04.subsec.21-2}{21.2}}{21.2}}}{Diviseurs et Idéaux fractionnaires inversibles}{258}
\EGAIVcontentsline{{\texorpdfstring{\protect\hyperlink{ega4.c04.subsec.21-3}{21.3}}{21.3}}}{Équivalence linéaire des diviseurs}{263}
\EGAIVcontentsline{{\texorpdfstring{\protect\hyperlink{ega4.c04.subsec.21-4}{21.4}}{21.4}}}{Images réciproques de diviseurs}{265}
\EGAIVcontentsline{{\texorpdfstring{\protect\hyperlink{ega4.c04.subsec.21-5}{21.5}}{21.5}}}{Images directes de diviseurs}{267}
\EGAIVcontentsline{{\texorpdfstring{\protect\hyperlink{ega4.c04.subsec.21-6}{21.6}}{21.6}}}{Cycle 1-codimensionnel associé à un diviseur}{270}
\EGAIVcontentsline{{\texorpdfstring{\protect\hyperlink{ega4.c04.subsec.21-7}{21.7}}{21.7}}}{Interprétation des cycles positifs 1-codimensionnels en termes de sous-préschémas}{277}
\EGAIVcontentsline{{\texorpdfstring{\protect\hyperlink{ega4.c04.subsec.21-8}{21.8}}{21.8}}}{Diviseurs et normalisation}{280}
\EGAIVcontentsline{{\texorpdfstring{\protect\hyperlink{ega4.c04.subsec.21-9}{21.9}}{21.9}}}{Diviseurs sur les préschémas de dimension 1}{284}
\EGAIVcontentsline{{\texorpdfstring{\protect\hyperlink{ega4.c04.subsec.21-10}{21.10}}{21.10}}}{Images réciproques et images directes de cycles 1-codimensionnels}{289}
\EGAIVcontentsline{{\texorpdfstring{\protect\hyperlink{ega4.c04.subsec.21-11}{21.11}}{21.11}}}{Factorialité des anneaux réguliers}{302}
\EGAIVcontentsline{{\texorpdfstring{\protect\hyperlink{ega4.c04.subsec.21-12}{21.12}}{21.12}}}{Le théorème de pureté de Van der Waerden pour l'ensemble de ramification d'un morphisme birationnel}{304}
\EGAIVcontentsline{{\texorpdfstring{\protect\hyperlink{ega4.c04.subsec.21-13}{21.13}}{21.13}}}{Couples parafactoriels. Anneaux locaux parafactoriels}{313}
\EGAIVcontentsline{{\texorpdfstring{\protect\hyperlink{ega4.c04.subsec.21-14}{21.14}}{21.14}}}{Le théorème de Ramanujam--Samuel}{323}
\EGAIVcontentsline{{\texorpdfstring{\protect\hyperlink{ega4.c04.subsec.21-15}{21.15}}{21.15}}}{Diviseurs relatifs}{329}
\EGAIVcontentsline{}{BIBLIOGRAPHIE}{333}
\EGAIVcontentsline{}{INDEX DES NOTATIONS}{334}
\EGAIVcontentsline{}{INDEX TERMINOLOGIQUE}{336}

\endgroup

% FIN SOURCE ega4/ega4-backmatter-contents-fr.tex

%
% DEBUT SOURCE ega4/ega4-errata-addenda-list3-fr.tex
% Transcription diplomatique française de EGA IV, quatrième partie, pp. 345--361.
\clearpage
\oldpage[IV]{345}
\phantomsection
\addcontentsline{toc}{section}{Errata et Addenda (Liste 3)}
\section*{ERRATA ET ADDENDA}
\markboth{ERRATA ET ADDENDA, LISTE 3}{ERRATA ET ADDENDA, LISTE 3}
\begin{center}\large (Liste 3)\end{center}

\subsection*{A) Erreurs typographiques}

\paragraph{(II, 3.3.2).}
Ligne 1 de la p.~58, remplacer $\sh{S}$ par $\sh{M}$. Ligne 6 de la
p.~58, remplacer $S(n)$ par $M(n)$.

\paragraph{(II, 4.1.1).}
Ligne 8 de la p.~71, remplacer $\sh{O}$-$P$-Module par
$\sh{O}_P$-Module.

\paragraph{(II, 5.4.5).}
Ligne 6 de la p.~102, remplacer $u$ par $f\circ u$.

\paragraph{(II, 5.4.8).}
Avant la ligne 4 du bas de la p.~102, remplacer $f\times1_{S'}$ par
$f_{(Y')}$ dans le diagramme. Ligne 4 du bas de la p.~102, remplacer
$f\times1_{S'}$ par $f_{(Y')}$.

\paragraph{(II, 6.4.4).}
Ligne 1 du haut de la p.~122, remplacer normal par intégralement clos.

\paragraph{(II, 6.4.7).}
Ligne 14 du bas de la p.~122, remplacer dimension par rang.

\paragraph{(II, 6.5.2).}
Ligne 5 de la p.~127, remplacer $(\sh{B}|U_\lambda)$ par
$(B|U_\lambda)$. Ligne 15 de la p.~127, supprimer le premier
$(N_{\sh{B}/\sh{A}}\omega_{\lambda\mu})$.

\paragraph{($0_{\mathrm{III}}$, 8.1.11).}
Ligne 8 du bas de la p.~8, remplacer $h_X(u)$ par $h'_X(u)$.

\paragraph{($0_{\mathrm{III}}$, 9.1.2).}
Ligne 10 du bas de la p.~12, remplacer $\sh{F}$ par $X$.

\paragraph{($0_{\mathrm{III}}$, 10.2.1).}
Ligne 6 du bas de la p.~18, remplacer : « en outre si » par : « si en
outre ».

\paragraph{(III, 4.4.12).}
Ligne 10 de la p.~138, remplacer : chap.~V par : chap.~IV (8.11.2).

\paragraph{(III, 5.2.5).}
Ligne 12 de la p.~153, remplacer (I, 5.5.4) par (II, 5.5.4).

\paragraph{(III, 7.4.7).}
Ligne 15 de la p.~57, remplacer : tout $A$-module par : tout
$A$-module $M$.

\paragraph{(III, 7.8.5).}
Ligne 18 du bas de la p.~74, remplacer $F$ par $\sh{F}$.

\paragraph{(III, 7.8.9).}
Ligne 9 du bas de la p.~75, remplacer $M$ par $\sh{M}$.

\paragraph{(III, 7.9.11).}
Ligne 17 du bas de la p.~79, remplacer chap.~V par chap.~VI.

\paragraph{($0_{\mathrm{IV}}$, 16.2.2.2).}
Ligne 13 de la p.~26, remplacer $M'\cap\mathfrak q^nM'$ par
$M'\cap\mathfrak q^nM$.

\paragraph{($0_{\mathrm{IV}}$, 17.3.7).}
Ligne 16 du bas de la p.~50, remplacer (17.1.10) par (17.3.3).

\paragraph{($0_{\mathrm{IV}}$, 18.1.3).}
Dans le diagramme (18.1.3.1), $p_2$ doit être à côté de la flèche
$G\to F$.

\paragraph{($0_{\mathrm{IV}}$, 19.8.7).}
Lignes 7, 8 et 9 du bas de la p.~111, remplacer (21.5.5) par (21.5.3).

\paragraph{($0_{\mathrm{IV}}$, 19.8.8).}
Ligne 2 du bas de la p.~112, remplacer $A$-module par $B$-module.

\paragraph{($0_{\mathrm{IV}}$, 20.5.4).}
Ligne 5 du bas de la p.~130, remplacer $\Omega_{B'/A'}$ par
$\Omega^1_{B'/A'}$.

\paragraph{($0_{\mathrm{IV}}$, 20.7.8).}
Ligne 9 du bas de la p.~149, remplacer (18.4.3) par (18.5.4). Ligne 4
du bas, remplacer (18.4.1) par (18.5.1).

\paragraph{($0_{\mathrm{IV}}$, 21.5.3).}
Ligne 1 du bas de la p.~166, remplacer ramifiée par ramifié.

\paragraph{($0_{\mathrm{IV}}$, 21.7.1).}
Ligne 4 de la p.~170, remplacer (20.6.15.1) par (20.6.17.1).

\paragraph{($0_{\mathrm{IV}}$, 22.1.1).}
Ligne 12 de la p.~183, remplacer (20.6.20) par (20.6.10).

\oldpage[IV]{346}
\paragraph{($0_{\mathrm{IV}}$, 22.3.3).}
Ligne 13 du bas de la p.~192, remplacer $\widehat B_{\mathfrak p}$ par
$\widehat B_{\mathfrak q}$.

\paragraph{($0_{\mathrm{IV}}$, 22.7.3).}
Ligne 21 du bas de la p.~212, remplacer $\operatorname{Ker}(i)$ par
$\operatorname{Ker}(i_0)$. Ligne 12 du bas de la p.~212, remplacer
$(\mathfrak p/\mathfrak p^2)\otimes_L L$ par
$(\mathfrak p/\mathfrak p^2)\otimes_A L$.

\paragraph{($0_{\mathrm{IV}}$, 23.1.3).}
Ligne 4 de la p.~215, remplacer $x'^q\in K'$ par $x'\in K'$. Ligne 9
de la p.~215, remplacer $x'^q\in A'$ par $x'^q\in A$.

\paragraph{({\texorpdfstring{\protect\hyperref[canon-index.ega4.c04.statement.2-1-1]{IV, 2.1.1}}{IV, 2.1.1}}).}
Ligne 4 du bas de la p.~5, remplacer (III, 1.4.15.1) par
(III, 1.4.15.5).

\paragraph{({\texorpdfstring{\protect\hyperref[canon-index.ega4.c04.prop.2-4-3]{IV, 2.4.3}}{IV, 2.4.3}}).}
Ligne 12 de la p.~20, remplacer
$(X_{\mathrm{red}}\times_{Y_{\mathrm{red}}}Y_{\mathrm{red}})_{\mathrm{red}}$
par
$(X_{\mathrm{red}}\times_{Y_{\mathrm{red}}}Y'_{\mathrm{red}})_{\mathrm{red}}$.

\paragraph{({\texorpdfstring{\protect\hyperref[canon-index.ega4.c04.prop.2-6-1]{IV, 2.6.1}}{IV, 2.6.1}}).}
Ligne 10 de la p.~27, remplacer : $X'\to X_{(S')}$ par
$X'=X_{(S')}$. 

\paragraph{({\texorpdfstring{\protect\hyperref[canon-index.ega4.c04.prop.4-2-1]{IV, 4.2.1}}{IV, 4.2.1}}).}
Ligne 20 du bas de la p.~55, remplacer
$\operatorname{deg.tr.}_kE$ par $\operatorname{deg.tr.}_LE$ et
$\operatorname{deg.tr.}_kL(t)$ par $\operatorname{deg.tr.}_LL(t)$;
ligne 19 du bas, échanger « première » et « seconde ».

\paragraph{({\texorpdfstring{\protect\hyperref[canon-index.ega4.c04.lem.4-4-2]{IV, 4.4.2}}{IV, 4.4.2}}).}
Ligne 1 du bas de la p.~59, remplacer $X$ par $X'$.

\paragraph{({\texorpdfstring{\protect\hyperref[canon-index.ega4.c04.prop.4-7-3]{IV, 4.7.3}}{IV, 4.7.3}}).}
Ligne 10 du bas de la p.~75, remplacer préschéma par $k$-préschéma.

\paragraph{({\texorpdfstring{\protect\hyperref[canon-index.ega4.c04.prop.5-6-1]{IV, 5.6.1}}{IV, 5.6.1}}).}
Ligne 16 de la p.~98, remplacer (en indice) $k(\mathfrak p($ par
$k(\mathfrak p)$.

\paragraph{({\texorpdfstring{\protect\hyperref[canon-index.ega4.c04.prop.5-6-5]{IV, 5.6.5}}{IV, 5.6.5}}).}
Ligne 20 de la p.~99, supprimer la seconde parenthèse dans $k(\xi))$.

\paragraph{({\texorpdfstring{\protect\hyperref[canon-index.ega4.c04.lem.5-6-11-1]{IV, 5.6.11.1}}{IV, 5.6.11.1}}).}
Ligne 16 de la p.~102, remplacer $K$ par $\mathfrak K$.

\paragraph{({\texorpdfstring{\protect\hyperref[canon-index.ega4.c04.cor.5-11-2]{IV, 5.11.2}}{IV, 5.11.2}}).}
Note de bas de page de la p.~123, remplacer 5.11.2) par (5.11.2).

\paragraph{({\texorpdfstring{\protect\hyperref[canon-index.ega4.c04.prop.6-1-5]{IV, 6.1.5}}{IV, 6.1.5}}).}
Ligne 2 du bas de la p.~136, remplacer (5.5.1.2) par
($0$, 16.3.9.2).

\paragraph{({\texorpdfstring{\protect\hyperref[canon-index.ega4.c04.prop.6-3-1]{IV, 6.3.1}}{IV, 6.3.1}}).}
Ligne 9 de la p.~139, remplacer $y\in\mathfrak m$ par
$y\in\mathfrak n$.

\paragraph{({\texorpdfstring{\protect\hyperref[canon-index.ega4.c04.cor.6-7-2]{IV, 6.7.2}}{IV, 6.7.2}}).}
Ligne 18 de la p.~146, remplacer $\sh{O}'_{x'}$ par $\sh{O}_{x'}$.

\paragraph{({\texorpdfstring{\protect\hyperref[canon-index.ega4.c04.statement.6-10-1]{IV, 6.10.1}}{IV, 6.10.1}}).}
Ligne 13 de la p.~155, remplacer
$\sh{F}\otimes_{\sh{O}_Y}\sh{O}_Y$ par
$\sh{F}\otimes_{\sh{O}_X}\sh{O}_Y$.

\paragraph{({\texorpdfstring{\protect\hyperref[canon-index.ega4.c04.prop.6-14-1]{IV, 6.14.1}}{IV, 6.14.1}}).}
Ligne 21 de la p.~170, remplacer : n\textsuperscript{o}~3, prop.~12
par : n\textsuperscript{o}~2, prop.~8.

\paragraph{({\texorpdfstring{\protect\hyperref[canon-index.ega4.c04.statement.6-15-2]{IV, 6.15.2}}{IV, 6.15.2}}).}
Ligne 13 de la p.~177, remplacer $t(t-v)$ par $t^2(t-v)^2$.

\paragraph{({\texorpdfstring{\protect\hyperref[canon-index.ega4.c04.prop.6-15-7]{IV, 6.15.7}}{IV, 6.15.7}}).}
Ligne 6 du bas de la p.~179, remplacer $X\otimes_Ak'$ par
$X\otimes_kk'$.

\paragraph{({\texorpdfstring{\protect\hyperref[canon-index.ega4.c04.item.7-5-4]{IV, 7.5.4}}{IV, 7.5.4}}).}
Ligne 14 du bas de la p.~206, remplacer
$\operatorname{Spec}(\widehat B\otimes k)$ par
$\operatorname{Spec}(\widehat B\otimes_{\widehat A}k)$.

\paragraph{({\texorpdfstring{\protect\hyperref[canon-index.ega4.c04.item.7-8-4]{IV, 7.8.4}}{IV, 7.8.4}}).}
Ligne 10 du bas de la p.~216, remplacer « vérifiant »
({\hyperref[canon-index.ega4.c04.def.7-8-2]{7.8.2}}, (ii) et (iii)) par
« vérifiant ({\hyperref[canon-index.ega4.c04.def.7-8-2]{7.8.2}}, (ii) et (iii)) ».

\paragraph{(IV, 2\textsuperscript{e} Partie, Bibliographie).}
Ligne 6 de la p.~224, remplacer : \emph{Amer. J. of Math.}, par
\emph{Annals of Math.}

\paragraph{({\texorpdfstring{\protect\hyperref[canon-index.ega4.c04.cor.8-5-8]{IV, 8.5.8}}{IV, 8.5.8}}).}
Ligne 6 de la p.~24, remplacer $\operatorname{Ker}g_\mu$ à la fin de
la ligne par $\operatorname{Ker}g_\lambda$.

\paragraph{({\texorpdfstring{\protect\hyperref[canon-index.ega4.c04.prop.8-10-1]{IV, 8.10.1}}{IV, 8.10.1}}).}
Ligne 13 du bas de la p.~36, remplacer (2.3.11) par
{\hyperref[canon-index.ega4.c04.cor.8-2-11]{(8.2.11)}}.

\paragraph{({\texorpdfstring{\protect\hyperref[canon-index.ega4.c04.prop.10-2-1]{IV, 10.2.1}}{IV, 10.2.1}}).}
Ligne 14 du bas de la p.~98, remplacer $y$ par $Y$.

\paragraph{({\texorpdfstring{\protect\hyperref[canon-index.ega4.c04.statement.10-4-11]{IV, 10.4.11}}{IV, 10.4.11}}).}
Ligne 10 du bas de la p.~104, remplacer (17.9.7) par
{\hyperlink{ega4.c04.prop.17-9-6}{(17.9.6)}}.

\paragraph{({\texorpdfstring{\protect\hyperref[canon-index.ega4.c04.thm.11-2-9]{IV, 11.2.9}}{IV, 11.2.9}}).}
Ligne 3 du bas de la p.~126, remplacer $A_\lambda$ par $A_\mu$.

\paragraph{({\texorpdfstring{\protect\hyperref[canon-index.ega4.c04.prop.11-3-4]{IV, 11.3.4}}{IV, 11.3.4}}).}
Ligne 16 de la p.~134, remplacer $F$ par $\sh{F}$.

\paragraph{({\texorpdfstring{\protect\hyperref[canon-index.ega4.c04.thm.11-3-10]{IV, 11.3.10}}{IV, 11.3.10}}).}
Ligne 9 du bas de la p.~139, remplacer (11.3.7) par
{\hyperref[canon-index.ega4.c04.thm.11-3-10]{(11.3.10)}}.

\paragraph{({\texorpdfstring{\protect\hyperref[canon-index.ega4.c04.statement.11-7-5]{IV, 11.7.5}}{IV, 11.7.5}}).}
Ligne 11 de la p.~159, remplacer (6.5.11) par (6.15.11).

\paragraph{({\texorpdfstring{\protect\hyperref[canon-index.ega4.c04.thm.12-1-1]{IV, 12.1.1}}{IV, 12.1.1}}).}
Ligne 13 de la p.~174, ajouter à la fin : pour $k\geq1$.

\paragraph{({\texorpdfstring{\protect\hyperref[canon-index.ega4.c04.thm.14-2-1]{IV, 14.2.1}}{IV, 14.2.1}}).}
Ligne 11 de la p.~202, remplacer (6.1.2) par (6.1.1).

\oldpage[IV]{347}
\paragraph{({\texorpdfstring{\protect\hyperref[canon-index.ega4.c04.prop.14-3-8]{IV, 14.3.8}}{IV, 14.3.8}}).}
Ligne 10 du bas de la p.~206, remplacer $(f^{-1}y)$ par $f^{-1}(y)$.

\paragraph{({\texorpdfstring{\protect\hyperref[canon-index.ega4.c04.rem.14-5-7]{IV, 14.5.7}}{IV, 14.5.7}}).}
Ligne 16 de la p.~218, remplacer (14.5.5) par (14.5.6).

\paragraph{({\texorpdfstring{\protect\hyperref[canon-index.ega4.c04.cor.14-5-11]{IV, 14.5.11}}{IV, 14.5.11}}).}
Dans le diagramme du haut de la page 223, les flèches horizontales
doivent toutes aller de droite à gauche.

\paragraph{({\texorpdfstring{\protect\hyperref[canon-index.ega4.c04.prop.15-4-2]{IV, 15.4.2}}{IV, 15.4.2}}).}
Ligne 23 de la p.~230, remplacer (6.1.4) par (6.1.5).

\paragraph{({\texorpdfstring{\protect\hyperref[canon-index.ega4.c04.cor.15-5-2]{IV, 15.5.2}}{IV, 15.5.2}}).}
Ligne 16 de la p.~233, remplacer $f^{-1}(y)$ par $f^{-1}(y')$.

\paragraph{({\texorpdfstring{\protect\hyperref[canon-index.ega4.c04.item.15-6-10]{IV, 15.6.10}}{IV, 15.6.10}}).}
Ligne 6 de la p.~242, remplacer ${}^{-1}(y')$ par $f^{-1}(y')$.
Ligne 11 de la p.~242, remplacer : $\alpha$) et $\beta$), par :
$\alpha$) ou $\beta$).

\subsection*{B) Modifications de texte}

\paragraph{(Err\textsubscript{IV}, 1).}
Dans ($0_{\mathrm{III}}$, 11.1.6), l'énoncé (reproduisant un énoncé de
(G, I, 4.4.1)) est incorrect lorsque $q(n)$ est un entier quelconque;
il n'est correct que si l'on a $q(n)=0$ ou $q(n)=n$. En effet, dans le
premier cas, le noyau de $d_r^{p,0}$ est alors $E_r^{p,0}$, puisque
l'image de $d_r^{p,0}$ est $E_r^{p+r,-r+1}=0$ par hypothèse; de même
l'image de $d_r^{p-r,r-1}$ est nulle puisque
$E_r^{p-r,r-1}=0$; on a donc bien $E_{r+1}^{p,0}=E_r^{p,0}$. On
raisonne de même lorsque $q(n)=n$ pour tout $n$.

\paragraph{(Err\textsubscript{IV}, 2).}
Dans (III, 3.4), supprimer les lignes 2 et 3 de la p.~119.

\oldpage[IV]{348}
\paragraph{(Err\textsubscript{IV}, 3).}
Dans (III, 4.4.9), ligne 15 du bas, p.~137, remplacer : finies par :
finies en tant qu'ensembles.

\paragraph{(Err\textsubscript{IV}, 4).}
Dans (III, 7.8.10), ligne 9 de la p.~76, remplacer simple par lisse.

\paragraph{(Err\textsubscript{IV}, 5).}
Dans ($0_{\mathrm{IV}}$, 16.1.4), ligne 8 du bas de la p.~23, remplacer :
Dire que $A$ est caténaire, par : Dire qu'un anneau intègre $A$ est
caténaire.

\paragraph{(Err\textsubscript{IV}, 6).}
Dans ($0_{\mathrm{IV}}$, 16.1.9), on peut supprimer l'hypothèse que
$A$ est noethérien. En effet, pour tout élément $u\in B_M$, $A[u]$
est une sous-$A$-algèbre de $C$, donc $u$ est entier sur $A$ (Bourbaki,
\emph{Alg. comm.}, chap.~V, \S1, n\textsuperscript{o}~1, cor. de la
prop.~1), et par suite aussi sur $A_M$; on conclut encore par
{\hyperlink{ega4.c04.prop.16-1-5}{(16.1.5)}}.

\paragraph{(Err\textsubscript{IV}, 7).}
Dans ($0_{\mathrm{IV}}$, 16.5.12), on notera que la relation
(16.5.11.1) équivaut à la même relation où l'on remplace $A$ par les
anneaux quotients $A/\mathfrak r$, $\mathfrak r$ parcourant l'ensemble
des idéaux premiers minimaux de $A$. On peut donc dans la preuve de
{\hyperlink{ega4.c04.statement.16-5-12}{(16.5.12)}} se borner au cas
où $A$ est intègre, et appliquer (16.1.4.2).

\paragraph{(Err\textsubscript{IV}, 8).}
Supprimer ($0_{\mathrm{IV}}$, 19.3.12) dont la démonstration est erronée
(les ensembles $W_\lambda$, $U_\lambda$ et $V_\lambda$ ne sont pas des
variétés linéaires).

\paragraph{(Err\textsubscript{IV}, 9).}
Dans ($0_{\mathrm{IV}}$, 19.8.3), lignes 6 et 8 du bas de la p.~110,
remplacer homomorphisme par homomorphisme local; de même dans
($0_{\mathrm{IV}}$, 19.8.4), ligne 1 de la p.~111.

\paragraph{(Err\textsubscript{IV}, 10).}
Dans ($0_{\mathrm{IV}}$, 20.1.2), ligne 2 du bas de la p.~117,
remplacer $a\in A$ par $b\in B$ et $aD$ par $bD$; ligne 1 du bas de la
p.~117, remplacer $A$-module par $B$-module.

\paragraph{(Err\textsubscript{IV}, 11).}
Dans ($0_{\mathrm{IV}}$, 20.2.1), ligne 8 de la p.~119, remplacer
$A$-modules par $B$-modules. Lignes 21 et 22 de la p.~119, remplacer
« di-homomorphismes de modules (relatif à $u$) » par : homomorphisme
de $B$-modules.

\paragraph{(Err\textsubscript{IV}, 12).}
Dans ($0_{\mathrm{IV}}$, 20.4.6.1), remplacer $p_2-p_1$ par $p_1-p_2$;
ligne 10 du bas de la p.~125, remplacer $x\otimes1-1\otimes x$ par
$1\otimes x-x\otimes1$.

\paragraph{(Err\textsubscript{IV}, 13).}
Dans ($0_{\mathrm{IV}}$, 20.4.8), ligne 14 de la p.~126, remplacer
$p_2=p_1-d_{B/A}$ par $p_1=p_2-d_{B/A}$.

\paragraph{(Err\textsubscript{IV}, 14).}
Dans ($0_{\mathrm{IV}}$, 20.4.13), remplacer les lignes 6 à 11 du bas
de la p.~127 par le texte suivant :

\medskip
\noindent
(iv) Soient $A$, $B$ deux anneaux intègres tels que $A\subset B$, $B$
étant entier sur $A$, et que le corps des fractions de $B$ soit une
extension séparable de celui de $A$. Alors le $B$-module
$\Omega^1_{B/A}$ est un module de torsion. En effet, pour tout $x\in B$,
soit $f$ le polynôme minimal de $x$ par rapport au corps des fractions
$K$ de $A$; il existe $a\in A$ non nul tel que $af[T]\in A[T]$; comme
$x$ est séparable sur $K$, on a $f'(x)\neq0$, et d'autre part de la
relation $af(x)=0$ on déduit $af'(x)d_{B/A}x=0$, d'où notre assertion
en vertu de {\hyperlink{ega4.c04.cor.20-4-7}{(20.4.7)}}.

\paragraph{(Err\textsubscript{IV}, 15).}
Dans ($0_{\mathrm{IV}}$, 22.2.7), ligne 5 de la p.~189, au lieu de
« formellement lisses », il faut « formellement lisses pour les
topologies préadiques ».

\paragraph{(Err\textsubscript{IV}, 16).}
Dans le Sommaire du chap.~IV, ligne 4 du bas de la p.~222, ajouter :
et anneaux strictement locaux. Lignes 1 à 3 du bas, remplacer par :
\begin{quote}
\S19. Immersions régulières et platitude normale.\\
\S20. Fonctions méromorphes et pseudo-morphismes.\\
\S21. Diviseurs.
\end{quote}
Ligne 16 de la p.~223, remplacer « de présentation finie » par
« localement de présentation finie »; ligne 19 de la p.~223, remplacer
« comme » par « comme étant ».

\paragraph{(Err\textsubscript{IV}, 17).}
Dans (IV, 3.1.10), ajouter après la ligne 16 de la p.~39 :

\begin{remark}[3.1.10.2]
\label{IV.3.1.10.2}
Soit $f:X\to Y$ un morphisme quasi-compact et quasi-séparé, et soit
$\sh{F}$ un $\sh{O}_X$-Module quasi-cohérent, de sorte que
$f_*(\sh{F})$ est un $\sh{O}_Y$-Module quasi-cohérent
{\hyperref[canon-index.ega4.c04.statement.1-7-4]{(1.7.4)}}. Alors on a
\begin{align}
 \operatorname{Ass}(f_*(\sh{F}))&\subset f(X),
 \label{IV.3.1.10.3}\tag{3.1.10.3}\\
 \operatorname{Ass}(f_*(\sh{F}))&\subset f(\operatorname{Ass}(\sh{F}))
 \quad\text{si $X$ est localement noethérien}.
 \label{IV.3.1.10.4}\tag{3.1.10.4}
\end{align}

Soit en effet $y\in\operatorname{Ass}(f_*(\sh{F}))$; en vertu de
{\hyperref[canon-index.ega4.c04.lem.2-3-1]{(2.3.1)}} et du fait que le morphisme
$\operatorname{Spec}(\sh{O}_{Y,y})\to Y$ est plat,
$(f_*(\sh{F}))_y$ est égal à $(f'_*(\sh{F}'))_y$, où
\[
 f':X\times_Y\operatorname{Spec}(\sh{O}_{Y,y})
 \longrightarrow\operatorname{Spec}(\sh{O}_{Y,y})
\]
est déduit de $f$ par changement de base et
$\sh{F}'=\sh{F}\otimes_{\sh{O}_Y}\sh{O}_{Y,y}$. On est donc ramené à
prouver (3.1.10.3) et (3.1.10.4) en y remplaçant $Y$ par
$\operatorname{Spec}(\sh{O}_{Y,y})$, $f$ par $f'$ et $\sh{F}$ par
$\sh{F}'$, compte tenu de (I, 3.6.5) et de la définition de
$\operatorname{Ass}(\sh{F}')$.

Par définition, il y a un élément
$t_y\in(f_*(\sh{F}))_y$ non nul et dont $\mathfrak m_y$ est
l'annulateur, et comme $(f_*(\sh{F}))_y$ s'identifie à
$\Gamma(Y,f_*(\sh{F}))$ puisque $Y$ est local, $t_y$ s'identifie aussi
à une section $s\in\Gamma(X,\sh{F})$. L'Idéal de $\sh{O}_X$ annulateur
de $s$ contient donc $\mathfrak m_y\sh{O}_X$, et il résulte de
($0_{\mathrm I}$, 1.7.4) que le support de $s$ est contenu dans
$f^{-1}(y)$; si ce dernier était vide, on aurait $s=0$, ce qui est
absurde. Lorsque de plus $X$ est localement noethérien, alors le
sous-$\sh{O}_X$-Module cohérent $\sh{O}_Xs$ de $\sh{F}$ n'est pas nul,
donc $\operatorname{Ass}(\sh{O}_Xs)\neq\varnothing$
{\hyperref[canon-index.ega4.c04.cor.3-1-5]{(3.1.5)}}; mais
\[
 \operatorname{Ass}(\sh{O}_Xs)\subset
 \operatorname{Supp}(\sh{O}_Xs)\subset f^{-1}(y),\qquad
 \operatorname{Ass}(\sh{O}_Xs)\subset\operatorname{Ass}(\sh{F})
\]
{\hyperref[canon-index.ega4.c04.prop.3-1-7]{(3.1.7)}}, donc
$y\in f(\operatorname{Ass}(\sh{F}))$.
\end{remark}

\oldpage[IV]{349}
\begin{remark}[3.1.10.5]
\label{IV.3.1.10.5}
On notera que même lorsque $f$ est un morphisme propre, on n'a pas
nécessairement
\[
 f(\operatorname{Ass}(\sh{F}))\subset
 \operatorname{Ass}(f_*(\sh{F})).
\]
On le voit en prenant pour $Y$ le spectre d'un corps et pour $X$ un
espace projectif sur $k$, par exemple $\mathbf P^1_k$; alors
$\sh{F}=\sh{O}_X(-1)$ est un $\sh{O}_X$-Module cohérent $\neq0$ tel
que $\Gamma(X,\sh{F})=0$ (III, 2.1.13), donc
$\operatorname{Ass}(f_*(\sh{F}))=\varnothing$ et
$\operatorname{Ass}(\sh{F})\neq\varnothing$.
\end{remark}

\paragraph{(Err\textsubscript{IV}, 18).}
Après (IV, 3.2.8), insérer après la ligne 17 de la p.~43 :

\begin{remark}[3.2.9]
\label{IV.3.2.9}
Pour tout point $s$ d'un préschéma $X$, notons
$i_s:\operatorname{Spec}(k(s))\to X$ le morphisme canonique, qui est
un monomorphisme quasi-compact. Pour tout $k(s)$-module $G(s)$,
considéré comme Module sur $\operatorname{Spec}(k(s))$,
$(i_s)_*(G(s))$ est donc un $\sh{O}_X$-Module quasi-cohérent
{\hyperref[canon-index.ega4.c04.statement.1-7-4]{(1.7.4)}}. On vérifie aussitôt
que la fibre de $(i_s)_*(G(s))$ est nulle en tout point
$x\notin\overline{\{s\}}$ et isomorphe à $G(s)$ dans le cas contraire.
Si $X$ est localement noethérien et $G(s)$ de longueur finie, il est
clair que $\operatorname{Ass}((i_s)_*(G(s)))$ est réduit au seul point
$s$, autrement dit $(i_s)_*(G(s))$ est irredondant.

Cela étant, supposons $X$ localement noethérien, et soit $\sh{F}$ un
$\sh{O}_X$-Module cohérent. Soit d'autre part $S$ une partie discrète
de $X$, et pour tout $s\in S$, posons
$\sh{F}(s)=i_s^*(\sh{F})$, qui est un $k(s)$-module; considérons pour
chaque $s\in S$ un module quotient $G(s)$ de $\sh{F}(s)$ de longueur
finie. On a donc un homomorphisme canonique
\[
 u:\sh{F}\longrightarrow
 \bigoplus_{s\in S}(i_s)_*(\sh{F}(s))\longrightarrow
 \bigoplus_{s\in S}(i_s)_*(G(s))
\]
de $\sh{O}_X$-Modules cohérents (puisque $S$, étant discret, est
localement fini dans $X$). Si $\sh{H}=\operatorname{Ker}u$,
$\sh{G}=\sh{F}/\sh{H}$ est un $\sh{O}_X$-Module cohérent et on déduit
de $u$ un homomorphisme injectif
\[
 v:\sh{G}\longrightarrow\bigoplus_{s\in S}(i_s)_*(G(s)).
\]
Si $\sh{G}^{(s)}$ est l'image de $\sh{G}$ dans $(i_s)_*(G(s))$,
$(\sh{G}^{(s)})_{s\in S}$ est une décomposition irredondante, qui est
réduite lorsque les $G(s)$ sont tous $\neq0$, et telle que
$\operatorname{Ass}(\sh{G})=S$ ($\sh{G}$ étant donc sans cycle premier
associé immergé). On a ainsi obtenu une bijection
$(G(s))_{s\in S}\leftrightarrow\sh{G}$, où $S$ parcourt l'ensemble des
parties discrètes de $X$ pour lesquelles $\sh{F}(s)\neq0$ pour tout
$s\in S$, et, pour chaque $S$, $(G(s))_{s\in S}$ parcourt l'ensemble
des quotients $\neq0$ et de longueur finie des $\sh{F}(s)$; d'autre
part $\sh{G}$ parcourt l'ensemble des $\sh{O}_X$-Modules quotients
cohérents de $\sh{F}$, sans cycle premier immergé. La bijection
réciproque associe à chaque $\sh{G}$ sa décomposition irredondante
réduite.
\end{remark}

\paragraph{(Err\textsubscript{IV}, 19).}
Dans (IV, 4.2.1), avant la ligne 17 du bas de la p.~55, insérer :

\begin{remark}[4.2.1.4]
\label{IV.4.2.1.4}
Sous les hypothèses de {\hyperref[canon-index.ega4.c04.prop.4-2-1]{(4.2.1)}}, on a
en outre
\begin{equation}
 \dim(K\otimes_kL)=
 \inf(\operatorname{deg.tr.}_kK,\operatorname{deg.tr.}_kL).
 \label{IV.4.2.1.5}\tag{4.2.1.5}
\end{equation}
où il doit être entendu que si les deux corps $K$, $L$ sont de degré de
transcendance infini (qui peuvent être des cardinaux quelconques), le
premier membre doit être pris égal à $+\infty$ (la dimension considérée
ici est celle définie dans ($0$, 16.1.1)).

Supposons en effet par exemple que
$\operatorname{deg.tr.}_kK\leq\operatorname{deg.tr.}_kL$ et en premier
lieu que $\operatorname{deg.tr.}_kK=n<+\infty$. On a vu dans la preuve
de {\hyperref[canon-index.ega4.c04.prop.4-2-1]{(4.2.1)}} que $K\otimes_kL$ est une
extension algébrique de $K'\otimes_kL$, et a donc même dimension
($0$, 16.1.5). On peut donc se borner au cas où
$K=k[T_1,\ldots,T_n]$; alors $K\otimes_kL$ est un anneau de fractions
de
\[
 k[T_1,\ldots,T_n]\otimes_kL=L[T_1,\ldots,T_n],
\]
qui est de dimension $n$ {\hyperref[canon-index.ega4.c04.cor.5-5-4]{(5.5.4)}}, et
par suite $\dim(K\otimes_kL)\leq n$ ($0$, 16.1.3.1). Pour prouver que
$\dim(K\otimes_kL)=n$, il suffit évidemment de montrer qu'il y a dans
$\operatorname{Spec}(L[T_1,\ldots,T_n])$ un point rationnel sur $L$
qui appartient à $\operatorname{Spec}(K\otimes_kL)$
{\hyperref[canon-index.ega4.c04.cor.5-2-3]{(5.2.3)}}, ou encore (I, 1.7.4) qu'il
existe un $L$-homomorphisme $u:K\otimes_kL\to L$. Mais par définition
du produit tensoriel, cela revient à prouver l'existence d'un
$k$-homomorphisme $v:K\to L$, ce qui est évident en vertu de
l'hypothèse
$\operatorname{deg.tr.}_kK\leq\operatorname{deg.tr.}_kL$.

Si $\operatorname{deg.tr.}_kK\leq\operatorname{deg.tr.}_kL$ et si ces
deux cardinaux sont infinis, on voit de la même façon qu'il y a un point
rationnel dans $\operatorname{Spec}(L[(t_\alpha)_{\alpha\in I}])$ qui
appartient à $\operatorname{Spec}(K\otimes_kL)$; mais le noyau de tout
$L$-homomorphisme de $L[(t_\alpha)_{\alpha\in I}]$ dans $L$ est un
idéal maximal engendré par une famille de la forme
$(t_\alpha-a_\alpha)_{\alpha\in I}$ avec $a_\alpha\in L$. Il est clair
qu'un tel idéal maximal contient une suite infinie décroissante
d'idéaux premiers distincts, donc l'anneau local de
$L[(t_\alpha)_{\alpha\in I}]$ correspondant est non noethérien et de
dimension infinie; comme c'est un anneau local de $K\otimes_kL$, on
voit que ce dernier est non noethérien et de dimension infinie.
\end{remark}

\oldpage[IV]{350}
\paragraph{(Err\textsubscript{IV}, 20).}
Après (IV, 4.5.19), insérer, avant la ligne 12 du bas de la p.~67 :

\begin{corollary}[4.5.19.3]
\label{IV.4.5.19.3}
Soient $k$ un corps, $X$ un $k$-préschéma de type fini sur $k$, qui est
irréductible mais non géométriquement irréductible. Alors il existe une
partie fermée $Y$ de $X$, rare dans $X$, contenant tous les points
$x\in X$ tels que $k(x)$ soit radiciel sur $k$ (et en particulier tous
les points de $X$ rationnels sur $k$).
\end{corollary}

Soient $k'$ une clôture algébrique de $k$, $X'=X\otimes_kk'$,
$p:X'\to X$ la projection canonique. Si $x\in X$ est tel que $k(x)$
soit radiciel sur $k$, il n'existe qu'un seul point $x'\in X'$ au-dessus
de $x$ (I, 3.4.9); ce point appartient nécessairement à deux composantes
irréductibles distinctes de $X'$, sans quoi $X$ serait géométriquement
irréductible en vertu de
{\hyperref[canon-index.ega4.c04.prop.4-5-19]{(4.5.19)}}, contrairement à
l'hypothèse. Soient alors $X'_i$ ($1\leq i\leq n$) les composantes
irréductibles du préschéma noethérien $X'$, et pour
$1\leq i<j\leq n$, posons $X'_{ij}=X'_i\cap X'_j$; l'ensemble
\[
 Y=p\left(\bigcup_{i<j}X'_{ij}\right)
\]
répond à la question. En effet, comme $p$ est un morphisme entier et
que les $X'_{ij}$ sont fermés dans $X'$, $Y$ est fermé dans $X$.
D'autre part, si $\xi$ est le point générique de $X$, $p^{-1}(\xi)$
est discret et formé des points génériques des $X'_i$
{\hyperref[canon-index.ega4.c04.cor.4-5-11]{(4.5.11)}}, donc ne rencontre aucun
des $X'_{ij}$; l'ensemble fermé $Y$ est par suite rare.

\paragraph{(Err\textsubscript{IV}, 21).}
Dans (IV, 5.6.11.1), ligne 7 du bas de la p.~102, après « intègre »,
ajouter : « de dimension 2 ».

\paragraph{(Err\textsubscript{IV}, 22).}
Après (IV, 6.2.3), après la ligne 18 de la p.~138, insérer :

\begin{proposition}[6.2.4]
\label{IV.6.2.4}
Sous les hypothèses de {\hyperref[canon-index.ega4.c04.prop.6-2-3]{(6.2.3)}}, soit
de plus $M$ un $A$-module de type fini et supposons $M\neq0$ et
$N\neq0$. Alors on a
\begin{equation}
 \operatorname{dim.proj}_B(M\otimes_AN)=
 \operatorname{dim.proj}_A(M)+
 \operatorname{dim.proj}_{B\otimes_Ak}(N\otimes_Ak).
 \label{IV.6.2.4.1}\tag{6.2.4.1}
\end{equation}
\end{proposition}

Considérons la résolution gauche $L_\bullet$ de $N$ introduite dans
{\hyperref[canon-index.ega4.c04.prop.6-2-3]{(6.2.3)}}, et une résolution gauche
$K_\bullet$ de $M$ par des $A$-modules libres de type fini :
\[
 \cdots\longrightarrow K_i\longrightarrow K_{i-1}\longrightarrow
 \cdots\longrightarrow K_0\longrightarrow M\longrightarrow0.
\]
On a vu que les $L_i$, $N$, les $Z_i(L_\bullet)$ sont des $A$-modules
plats; il en est de même des $B_i(L_\bullet)$ qui sont égaux aux
$Z_i(L_\bullet)$ sauf pour $i=0$, et dans ce dernier cas
$B_0(L_\bullet)$ est

\oldpage[IV]{351}
plat par ($0_{\mathrm I}$, 6.1.2); enfin, les $H_i(L_\bullet)$ sont nuls
par définition sauf $H_0(L_\bullet)=N$, qui est un $A$-module plat. On
sait que, dans ces conditions,
\[
 H_n(K_\bullet\otimes_AL_\bullet)\simeq
 \bigoplus_{p+q=n}
 H_p(K_\bullet)\otimes_AH_q(L_\bullet)
\]
(G, I, 5.5), donc $K_\bullet\otimes_AL_\bullet$ (gradué comme
d'ordinaire) est une résolution de $M\otimes_AN$ par des $B$-modules
libres. Posons
\[
 r=\operatorname{dim.proj}_A(M),\qquad
 s=\operatorname{dim.proj}_B(N)=
 \operatorname{dim.proj}_{B\otimes_Ak}(N\otimes_Ak).
\]
Il y a donc une résolution libre $K_\bullet$ de longueur $r$ et une
résolution libre $L_\bullet$ de longueur $s$, et
$K_\bullet\otimes_AL_\bullet$ est alors une résolution libre de longueur
$r+s$, ce qui prouve déjà que
$\operatorname{dim.proj}_B(M\otimes_AN)\leq r+s$. Il reste à prouver
($0$, 17.2.6) que l'on a
\[
 \operatorname{Tor}^{B}_{r+s}(M\otimes_AN,k(B))\neq0.
\]
Utilisant la résolution précédente $K_\bullet\otimes_AL_\bullet$ pour
calculer ce foncteur, on obtient
\[
 \operatorname{Tor}^{B}_{r+s}(M\otimes_AN,k(B))=
 H_{r+s}((K_\bullet\otimes_AL_\bullet)\otimes_Bk(B)).
\]
Mais on peut écrire
\[
 (K_\bullet\otimes_AL_\bullet)\otimes_Bk(B)=
 (K_\bullet\otimes_Ak)\otimes_k
 (L_\bullet\otimes_Bk(B)),
\]
qui est donc un produit tensoriel de deux complexes d'espaces
vectoriels sur $k$; tous les $k$-modules étant plats, la formule de
Künneth donne
\[
 H_{r+s}((K_\bullet\otimes_Ak)\otimes_k
 (L_\bullet\otimes_Bk(B)))\simeq
 H_r(K_\bullet\otimes_Ak)\otimes_k
 H_s(L_\bullet\otimes_Bk(B)),
\]
l'homologie du premier (resp. second) complexe étant nulle en degrés
$>r$ (resp. $>s$). On a donc
\[
 \operatorname{Tor}^{B}_{r+s}(M\otimes_AN,k(B))=
 \operatorname{Tor}^{A}_{r}(M,k)\otimes_k
 \operatorname{Tor}^{B}_{s}(N,k(B))
\]
à isomorphie près, et comme chacun des facteurs est $\neq0$ par
hypothèse, il en est de même de leur produit tensoriel sur $k$.

\paragraph{(Err\textsubscript{IV}, 23).}
Dans (IV, 6.10.1), ligne 8 de la p.~155, après : canonique, ajouter :
$\sh{F}$ un $\sh{O}_X$-Module cohérent.

\paragraph{(Err\textsubscript{IV}, 24).}
Dans ({\hyperref[canon-index.ega4.c04.prop.6-14-4]{IV, 6.14.4}}), la référence
donnée ligne 12 du bas de la p.~175 est insuffisante; il faut raisonner
de la façon suivante. Tout d'abord, si $M$ est une base de transcendance
de $K'$ sur $A$, $A(M)$ est algébriquement fermé dans $B(A(M))$
(Bourbaki, \emph{Alg.}, chap.~V, \S6, exerc.~8 c)). On peut donc se
borner au cas où $K'$ est une extension algébrique séparable de $A$.
Soit $p$ l'exposant caractéristique de $A$; il résulte de Bourbaki,
\emph{loc. cit.}, \S9, exerc.~2, que la fermeture algébrique de $K'$
dans $B(K')$ est radicielle sur $K'$; montrons que tout élément
$x\in B(K')$ tel que $x^{p^r}\in K'$ appartient déjà à $K'$. Pour cela,
considérons une base $(e_\lambda)$ de $K'$ sur $A$; puisque $K'$ est une
extension algébrique séparable de $A$, les $e_\lambda^{p^r}$ forment
aussi une base de $K'$ sur $A$ (Bourbaki, \emph{loc. cit.}, \S8,
n\textsuperscript{o}~3, cor. de la prop.~4). On a
$x=\sum_\lambda e_\lambda b_\lambda$, avec $b_\lambda\in B$ et, par
hypothèse, il existe des éléments $a_\lambda\in A$ tels que l'on ait
\[
 \sum_\lambda e_\lambda^{p^r}b_\lambda^{p^r}=
 \sum_\lambda e_\lambda^{p^r}a_\lambda;
\]
mais cela entraîne $b_\lambda^{p^r}=a_\lambda$ pour tout $\lambda$;
comme par hypothèse $A$ est algébriquement fermé dans $B$, on a donc
$b_\lambda\in A$ pour tout $\lambda$, et par suite $x\in K'$.

\paragraph{(Err\textsubscript{IV}, 25).}
Dans (IV, 7.4.8, B)), p.~203, la question a été résolue affirmativement
par R.~Kiehl, \emph{Ausgezeichnete Ringe der nichtarchimedischen
analytischen Geometrie}, à paraître dans \emph{Inventiones
Mathematicae}.

\paragraph{(Err\textsubscript{IV}, 26).}
Dans (IV, 7.8.4), ligne 9 du bas de la p.~217, remplacer :
\begin{quote}
Nous verrons plus loin (\S18) que si $k$ est un corps valué complet non
discret\ldots
\end{quote}
par :
\begin{quote}
Il a été démontré par R.~Kiehl (dans l'article cité en
(Err\textsubscript{IV}, 25)) que si $k$ est un corps valué complet non
discret\ldots
\end{quote}

\oldpage[IV]{352}
\paragraph{(Err\textsubscript{IV}, 27).}
Dans (IV, 7.8.4), ajouter, avant la ligne 6 du bas de la p.~217 :

\smallskip
\noindent
\textup{(vi)} Soit $A$ un anneau local noethérien intègre de dimension
1; il est alors universellement caténaire
{\hyperref[canon-index.ega4.c04.cor.7-2-9]{(7.2.9)}} et, pour qu'il soit excellent,
il faut et il suffit donc, en vertu de
{\hyperref[canon-index.ega4.c04.statement.7-3-19]{(7.3.19)}}, que ses fibres
formelles soient géométriquement réduites, ou encore que son complété
$\widehat A$ soit réduit et que les corps composants de son anneau total
de fractions soient des extensions séparables du corps des fractions
$K$ de $A$. En vertu de {\hyperref[canon-index.ega4.c04.thm.7-6-4]{(7.6.4)}} et
{\hyperref[canon-index.ega4.c04.thm.7-7-2]{(7.7.2)}}, il revient aussi au même de
dire que $A$ est universellement japonais. Lorsque $A$ est un anneau de
valuation discrète, il en est de même de $\widehat A$; dire que $A$ est
excellent signifie alors que le corps des fractions de $\widehat A$
est une extension séparable de $K$, condition toujours satisfaite si
$K$ est de caractéristique 0.

\paragraph{(Err\textsubscript{IV}, 28).}
Après (IV, 8.10.3), après la ligne 14 de la p.~37, insérer :

\begin{remark}[8.10.3.1]
\label{IV.8.10.3.1}
On ne peut dans {\hyperref[canon-index.ega4.c04.prop.8-10-3]{(8.10.3)}} supprimer
l'hypothèse que les $\overline{f_\alpha(X_\alpha)}$ soient
constructibles. Prenons par exemple $S=\operatorname{Spec}(A)$, où $A$
est un anneau tel que $S$ soit compact, totalement discontinu et sans
point isolé, tout idéal premier $\mathfrak p$ de $A$ étant maximal et
$A_\mathfrak p$ étant un corps (Bourbaki, \emph{Alg. comm.}, chap.~II,
\S4, exerc.~16 et 17). Soit $s\in S$, et prenons pour
$Y_\alpha=S_\alpha$ les voisinages ouverts affines de $s$ dans $S$; on
prend $X_\alpha=\operatorname{Spec}(k(s))$ pour tout $\alpha$,
$f_\alpha$ et $f$ étant les morphismes canoniques; il est clair que
$f$, qui est l'identité, est dominant, mais aucun des $f_\alpha$ ne
l'est.
\end{remark}

\paragraph{(Err\textsubscript{IV}, 29).}
Dans (IV, 8.12.8), ligne 7 de la p.~47, remplacer « tel qu'il existe un
$\sh{O}_Y$-Module ample » par : « et quasi-séparé ». Ligne 12 de la
p.~47, remplacer « L'hypothèse » par : « Il résulte de
{\hyperref[canon-index.ega4.c04.prop.8-12-3]{(8.12.3)}} que l'existence de la
factorisation de l'énoncé est une propriété locale sur $Y$; on peut
donc supposer $Y$ affine, ce qui ».

\paragraph{(Err\textsubscript{IV}, 30).}
Après (IV, 8.12.11), ajouter à la fin de la p.~48 :

\begin{corollary}[8.12.12]
\EGAReaderTarget{ega4.c04.cor.8-12-12}{IV.8.12.12}
Soient $X$, $Y$ deux préschémas intègres,
$f:X\to Y$ un morphisme séparé, birationnel et localement de type fini.
Soit $y\in Y$ un point normal de $Y$ et supposons qu'un point
$x\in f^{-1}(y)$ soit isolé dans $f^{-1}(y)$. Alors il existe un
voisinage ouvert $U$ de $y$ dans $Y$ tel que la restriction
$f^{-1}(U)\to U$ de $f$ soit un isomorphisme.
\end{corollary}

\begin{proof}
Il suffit de prouver qu'il existe un voisinage ouvert $U$ de $y$ dans
$Y$ et une $U$-section $g:U\to f^{-1}(U)$ de $f^{-1}(U)$ telle que
$g(y)=x$. En effet, comme $U$ contient le point générique $\eta$ de
$Y$ et $f^{-1}(U)$ le point générique $\xi$ de $X$, on a
nécessairement $g(\eta)=\xi$ puisque $f$ est birationnel
{\hyperref[canon-index.ega4.c04.statement.6-15-4]{(6.15.4)}}; comme $g$ est une
immersion, $g(U)$ est localement fermé dans $f^{-1}(U)$, donc ouvert
dans son adhérence dans $f^{-1}(U)$, et comme cette dernière est égale
à $f^{-1}(U)$ d'après ce qui précède, $g(U)$ est ouvert dans
$f^{-1}(U)$; comme le sous-préschéma de $f^{-1}(U)$ associé à $g$ est
isomorphe à $U$, donc réduit, et que $f^{-1}(U)$ est réduit, il résulte
de (I, 5.2.1) que $g$ est une immersion ouverte. D'autre part, puisque
$f$ est séparé, $g$ est une immersion fermée (I, 5.4.6), donc $g(U)$
est un ensemble à la fois ouvert et fermé dans $f^{-1}(U)$. Mais comme
$f^{-1}(U)$ est intègre, cela entraînera $g(U)=f^{-1}(U)$, d'où la
conclusion.

Considérons alors deux cas :

\textup{I)} Supposons d'abord $Y$ normal, et soit $V$ l'ensemble des
points de $X$ en lesquels $f$ est quasi-fini; c'est un ensemble ouvert
dans $X$ {\hyperref[canon-index.ega4.c04.cor.13-1-4]{(13.1.4)}}.\footnote{Le
lecteur vérifiera que ni {\hyperlink{ega4.c04.cor.8-12-12}{(8.12.12)}}
ni aucune de ses conséquences n'est utilisée dans la preuve de
{\hyperref[canon-index.ega4.c04.cor.13-1-4]{(13.1.4)}}.} Il résulte alors de
{\hyperref[canon-index.ega4.c04.cor.8-12-10]{(8.12.10)}}

\oldpage[IV]{353}
que la restriction $f|V:V\to Y$ est une immersion ouverte; donc il y a
un voisinage ouvert $U$ de $y\in Y$ tel que $f|V$ soit un isomorphisme
de $V$ sur le sous-préschéma induit par $Y$ sur $U$; l'isomorphisme
réciproque $g$ est alors une $U$-section de $f^{-1}(U)$ telle que
$g(y)=x$.

\textup{II)} L'obtention d'une $U$-section $g$ telle que $g(y)=x$ étant
un problème local sur $X$ et sur $Y$, on peut supposer que $X$ et $Y$
sont affines et que $f$ est de type fini, donc que $X$ est un
sous-préschéma fermé d'un $Y$-préschéma $Z$ de présentation finie sur
$Y$ (par exemple de la forme $Y[T_1,\ldots,T_n]$). Posons
\[
 Y'=\operatorname{Spec}(\sh{O}_{Y,y}),\qquad
 X'=X_{(Y')},\qquad Z'=Z_{(Y')},
\]
de sorte que $X'$ est un sous-préschéma fermé de $Z'$. En outre, si
$f'=f_{(Y')}$, on a $f'^{-1}(y)=f^{-1}(y)$ (I, 3.6.5) et $x$ est donc
isolé dans $f'^{-1}(y)$. Puisque $Y'$ est normal, on peut appliquer le
résultat de I) à $f'$, et puisque $Y'$ est le seul voisinage de $y$ dans
$Y'$, il existe une $Y'$-section $g'$ de $X'$ telle que $g'(y)=x$, que
l'on peut évidemment aussi considérer comme une $Y'$-section de $Z'$.
Puisque le morphisme structural $h:Z\to Y$ est de présentation finie,
on peut appliquer ({\hyperref[canon-index.ega4.c04.thm.8-8-2]{8.8.2}}, (i) et
(ii)) suivant la méthode de
({\hyperref[canon-index.ega4.c04.statement.8-1-2]{8.1.2}}, a)), d'où l'on conclut
à l'existence d'un voisinage ouvert $U$ de $y$ dans $Y$ et d'une
$U$-section $g$ de $h^{-1}(U)$ telle que $g(y)=x$. Il reste à voir
qu'en fait $g$ est une $U$-section de $f^{-1}(U)$. Mais puisque
$f^{-1}(U)$ est un sous-préschéma fermé de $h^{-1}(U)$ et que $U$ est
réduit, il suffit pour cela (I, 5.2.2) de voir que
$g(U)\subset f^{-1}(U)$. Or, puisque $U$ est intègre et $f^{-1}(U)$
fermé dans $h^{-1}(U)$, il suffit de prouver que si $\eta$ est le point
générique de $Y$ (et de $U$), on a $g(\eta)\in f^{-1}(U)$. Or, cela
résulte de ce que la restriction de $g$ à $Y'$ est une $Y'$-section de
$X'\subset f^{-1}(U)$ et de ce que $\eta\in Y'$.
\end{proof}

\paragraph{(Err\textsubscript{IV}, 31).}
Dans ({\hyperref[canon-index.ega4.c04.subsec.10-4]{IV, 10.4}}), ajouter, après la
proposition {\hyperref[canon-index.ega4.c04.statement.10-4-11]{(10.4.11)}} :

\begin{corollary}[10.4.12]
\EGAReaderTarget{ega4.c04.cor.10-4-12}{IV.10.4.12}
Soient $S$ un préschéma, $X$ un $S$-préschéma de présentation finie,
$f$ un $S$-endomorphisme radiciel de $X$.
\begin{enumerate}[label=\textup{(\roman*)}]
\item L'endomorphisme $f$ est bijectif et fini (donc un homéomorphisme
universel {\hyperref[canon-index.ega4.c04.prop.2-4-5]{(2.4.5)}}).
\item Supposons de plus que $X$ soit réduit, que les anneaux locaux de
$S$ soient noethériens et universellement japonais, et que $f$ soit
birationnel (condition automatiquement vérifiée si pour tout point
maximal $x$ de $X$, $k(x)$ est un corps parfait). Alors $f$ est un
$S$-automorphisme.
\end{enumerate}
\end{corollary}

\begin{proof}
\textup{(i)} La question étant locale sur $S$, on peut se borner au cas
où $S=\operatorname{Spec}(A)$ est affine. Il existe alors une
sous-$\mathbb Z$-algèbre de type fini $A_0$ de $A$ telle qu'en posant
$S_0=\operatorname{Spec}(A_0)$, on ait
$X=X_0\times_{S_0}S$, où $X_0$ est un préschéma de type fini sur $S_0$,
et que $f$ se déduise d'un $S_0$-endomorphisme $f_0$ de $X_0$ par le
changement de base $S\to S_0$
{\hyperref[canon-index.ega4.c04.prop.8-9-1]{(8.9.1)}}; on peut en outre supposer
$f_0$ radiciel ({\hyperref[canon-index.ega4.c04.thm.8-10-5]{8.10.5}}, (vii)). Il
suffira de prouver que $f_0$ est bijectif et fini, et on peut donc
supposer désormais que $A$ est une $\mathbb Z$-algèbre de type fini, ce
qui entraîne que $S$ et $X$ sont des préschémas excellents
({\hyperref[canon-index.ega4.c04.statement.7-8-3]{(7.8.3)}} et
{\hyperref[canon-index.ega4.c04.prop.7-8-6]{(7.8.6)}}) et en particulier que $X$
est un préschéma universellement japonais
({\hyperref[canon-index.ega4.c04.cor.5-11-4]{(5.11.4)}} et
({\hyperref[canon-index.ega4.c04.statement.7-8-3]{7.8.3}}, (vi))). On peut en outre
supposer $X$ réduit; en effet,
$f_{\mathrm{red}}:X_{\mathrm{red}}\to X_{\mathrm{red}}$ est un
$S_{\mathrm{red}}$-endomorphisme radiciel (I, 5.1.6); si l'on prouve que
$f_{\mathrm{red}}$ est fini, il en sera de même de $f$ : en effet,
puisque $f$ est radiciel, il est séparé
{\hyperref[canon-index.ega4.c04.lem.1-8-7-1]{(1.8.7.1)}}; comme
$f_{\mathrm{red}}$ est propre il en est de même de $f$ (II, 5.4.6);
mais puisque $f$ est évidemment quasi-fini, il est fini (III, 4.4.2).

\oldpage[IV]{354}
Soit alors $X'$ le normalisé de $X$ (II, 6.3.8); si $X_i$
($1\leq i\leq n$) sont les sous-préschémas réduits de $X$ ayant pour
espaces sous-jacents les composantes irréductibles de $X$, $X'$ est
somme des normalisés $X'_i$ des préschémas intègres $X_i$; nous
désignerons par $x_i$ (resp. $x'_i$) le point générique de $X_i$
(resp. $X'_i$). Le morphisme canonique $p:X'\to X$ est fini puisque
$X$ est universellement japonais et on a
$p^{-1}(x_i)=\{x'_i\}$, l'homomorphisme $k(x_i)\to k(x'_i)$ étant
bijectif.

Soit $f=(\psi,\theta)$; on sait déjà
{\hyperref[canon-index.ega4.c04.statement.10-4-11]{(10.4.11)}} que $\psi$ est une
bijection continue de l'espace $X$ sur lui-même; si l'on pose
$y_i=\psi(x_i)$, $y_i$ est aussi un point maximal de $X$, car si $y'_i$
était une générisation de $y_i$ distincte de $y_i$, $x_i$ serait une
générisation $\neq x_i$ de $\psi^{-1}(y'_i)$. Puisque l'ensemble $M$
des points maximaux de $X$ est fini, la restriction $\psi|M$ est une
bijection de $M$ sur lui-même, et on peut écrire
$\psi(x_i)=x_{\sigma(i)}$, où $\sigma$ est une permutation de
$[1,n]$. Alors, pour tout $i$, la restriction $f|X_i$ de $f$ se
factorise en
\[
 X_i\xrightarrow{f_i}X_{\sigma(i)}\longrightarrow X
\]
(I, 5.2.2), où la seconde flèche est l'injection canonique. Comme les
$X_i$ sont intègres, il existe (II, 6.3.9) un unique morphisme
$f'_i:X'_i\to X'_{\sigma(i)}$ rendant commutatif le diagramme
\[
\begin{array}{ccc}
X'_i&\xrightarrow{f'_i}&X'_{\sigma(i)}\\
\downarrow p_i&&\downarrow p_{\sigma(i)}\\
X_i&\xrightarrow{f_i}&X_{\sigma(i)}
\end{array}
\]
et si $f'$ coïncide avec $f'_i$ dans chacun des $X'_i$, on a donc un
diagramme commutatif
\phantomsection\EGAReaderTarget{ega4.c04.sec.10-4-12-1}{IV.10.4.12.1}
\begin{equation}
\label{IV.10.4.12.1.errata}
\begin{array}{ccc}
X'&\xrightarrow{f'}&X'\\
\downarrow p&&\downarrow p\\
X&\xrightarrow{f}&X.
\end{array}
\tag{10.4.12.1}
\end{equation}

D'autre part, comme $f$ est séparé, il en est de même de
$f\circ p=p\circ f'$, donc $f'$ est séparé (I, 5.5.1); enfin, puisque
$p$ est fini et $f$ bijectif, il résulte de
{\hyperlink{ega4.c04.sec.10-4-12-1}{(10.4.12.1)}} que $f'$ est
quasi-fini. Puisque les $X'_i$ sont intègres et normaux, il résulte de
{\hyperref[canon-index.ega4.c04.cor.8-12-11]{(8.12.11)}} que $f'_i$ se factorise en
\[
 X'_i\xrightarrow{f''_i}X''_{\sigma(i)}
 \xrightarrow{u_{\sigma(i)}}X'_{\sigma(i)},
\]
où $X''_{\sigma(i)}$ est la fermeture intégrale de
$X'_{\sigma(i)}$ relative à l'extension $k(x'_i)$ de
$k(x'_{\sigma(i)})$, et $f''_i$ une immersion ouverte. Mais comme
$k(x'_i)$ est une extension radicielle de $k(x'_{\sigma(i)})$,
$u_{\sigma(i)}$ est un morphisme radiciel (Bourbaki,
\emph{Alg. comm.}, chap.~V, \S2, n\textsuperscript{o}~3, lemme 4);
donc $f'$ est radiciel, et il résulte alors de
{\hyperref[canon-index.ega4.c04.statement.10-4-11]{(10.4.11)}} que $f'$ est
surjectif; on en conclut (puisque les $u_{\sigma(i)}$ sont injectifs)
que $f''_i$ est bijectif, donc un isomorphisme; on en déduit enfin que
le morphisme $f'$ est fini (puisqu'il en est ainsi de $u_{\sigma(i)}$,
$X'$ étant universellement japonais
{\hyperref[canon-index.ega4.c04.statement.7-8-3]{(7.8.3)}}). Cela étant, comme
$p\circ f'=f\circ p$ est alors fini et que $p$ est surjectif et $f$
séparé, $f$ est propre (II, 5.4.3) et quasi-fini, donc fini
{\hyperref[canon-index.ega4.c04.thm.8-11-1]{(8.11.1)}}.

\oldpage[IV]{355}
\textup{(ii)} Compte tenu de
({\hyperref[canon-index.ega4.c04.thm.8-10-5]{8.10.5}}, (i)) appliqué suivant la
méthode de ({\hyperref[canon-index.ega4.c04.statement.8-1-2]{8.1.2}}, a)), on peut
se borner au cas où $S$ est un schéma local. Les réductions préalables
faites dans la preuve de (i) sont donc ici inutiles, et avec les mêmes
notations on voit que l'on a un diagramme commutatif
{\hyperlink{ega4.c04.sec.10-4-12-1}{(10.4.12.1)}} dans lequel $X$ est
un $S$-préschéma de type fini, $p$ est fini, $f'$ est un isomorphisme et
l'homomorphisme canonique $\sh{O}_X\to p_*(\sh{O}_{X'})$ est injectif.
Mais on montrera au chap.~VI, dans la théorie des quotients des
préschémas par des relations d'équivalence, que sous ces conditions $f$
est lui-même un isomorphisme (le lecteur vérifiera qu'on n'utilisera pas
{\hyperlink{ega4.c04.cor.10-4-12}{(10.4.12)}} dans la théorie des
relations d'équivalence du chap.~VI).
\end{proof}

\begin{corollary}[10.4.13]
\EGAReaderTarget{ega4.c04.cor.10-4-13}{IV.10.4.13}
Soient $k$ un corps, $X$ un préschéma algébrique sur $k$. Alors tout
$k$-endomorphisme radiciel $f$ de $X$ est fini et bijectif, donc un
homéomorphisme universel. Si de plus $X$ est réduit et $f$ birationnel
(par exemple si $k$ est de caractéristique 0), $f$ est un
$k$-automorphisme.
\end{corollary}

\begin{remark}[10.4.14]
\label{IV.10.4.14}
\begin{enumerate}[label=\textup{(\roman*)}]
\item Dans {\hyperlink{ega4.c04.cor.10-4-12}{(10.4.12)}} et
{\hyperlink{ega4.c04.cor.10-4-13}{(10.4.13)}} on ne peut supprimer
l'hypothèse que $X$ est réduit pour la seconde conclusion; c'est ce que
montre l'exemple du $B$-endomorphisme correspondant à l'augmentation
$D_B(L)\to B$ de l'anneau $D_B(L)=A$ défini dans ($0$, 18.2.3), où
$B$ est identifié à un sous-anneau de $A$.

\item Soit $X$ un préschéma noethérien. Nous ignorons si tout
endomorphisme radiciel $f$ de $X$ est nécessairement surjectif et fini.
Si $f$ est supposé surjectif et si de plus $X$ est un préschéma
japonais, le raisonnement de
{\hyperlink{ega4.c04.cor.10-4-12}{(10.4.12)}} montre que $f$ est fini;
si de plus $X$ est réduit et si $f$ est birationnel, par exemple si,
pour tout point maximal $x\in X$, $k(x)$ est parfait, le même
raisonnement montre que $f$ est un automorphisme.

\item On peut donner des exemples d'endomorphismes radiciels non
surjectifs de schémas affines intègres non noethériens : si $k$ est un
anneau, $A=k[(T_\lambda)_{\lambda\in L}]$ un anneau de polynômes où $L$
est infini dénombrable, et $\mathfrak J$ l'idéal engendré dans $A$ par
une sous-famille $(T_\lambda)_{\lambda\in H}$, où $L-H$ est infini et
$H\neq\varnothing$, on remarque que $A/\mathfrak J$ est isomorphe à
$k[(T_\lambda)_{\lambda\in L-H}]$, donc à $A$; l'endomorphisme de
$\operatorname{Spec}(A)$ correspondant à l'endomorphisme d'anneaux
composé
\[
 A\xrightarrow{p}A/\mathfrak J\xrightarrow{u}A
\]
où $p$ est la projection canonique et $u$ un isomorphisme, répond à la
question, son image étant le sous-préschéma fermé de
$\operatorname{Spec}(A)$ défini par $\mathfrak J$.
\end{enumerate}
\end{remark}

\paragraph{(Err\textsubscript{IV}, 32).}
Dans (IV, 11.3.14), ligne 7 de la p.~141, avant « Soient », insérer :
(i). Après la ligne 9 de la p.~141, ajouter :

\smallskip
\noindent
\textup{(ii)} Inversement, soit $f:X\to Y$ un morphisme plat. Si $X$
est unibranche (resp. géométriquement unibranche) en un point $x$, $Y$
est unibranche (resp. géométriquement unibranche) au point $y=f(x)$.

Au début de la ligne 10 de la p.~141, insérer : (i). Après la ligne 20
de la p.~141, ajouter :

\smallskip
\noindent
\textup{(ii)} Notons que si $Y_1=Y_{\mathrm{red}}$,
$X_1=X\times_Y Y_1$ a pour espace sous-jacent celui de $X$ et le
morphisme $f_1:X_1\to Y_1$ déduit de $f$ par changement de base est
plat. On peut donc se borner au cas où $Y=\operatorname{Spec}(A)$,
$A$ étant un anneau local réduit, et $y$ le point fermé de $Y$. Posons
$B=\sh{O}_{X,x}$; par hypothèse $B_{\mathrm{red}}$ est unibranche, donc
$\operatorname{Spec}(B)$ n'a qu'un seul point maximal, et par platitude
{\hyperref[canon-index.ega4.c04.prop.2-3-4]{(2.3.4)}} on en déduit qu'il en est de
même

\oldpage[IV]{356}
de $\operatorname{Spec}(A)$; puisque $A$ est réduit, il est donc
intègre. Soient $K$ le corps des fractions de $A$, $A'$ la clôture
intégrale de $A$; par platitude, $B\otimes_AA'$ s'identifie à un
sous-anneau de $B\otimes_AK$ contenant $B$, et comme les éléments non
nuls de $A$ sont $B$-réguliers par platitude, $B\otimes_AK$ s'identifie
à un sous-anneau de l'anneau total des fractions $R$ de $B$. Pour
montrer que $A$ est unibranche, il suffit de voir que $A'$ ne peut avoir
deux idéaux maximaux distincts (nécessairement au-dessus de $y$
(Bourbaki, \emph{Alg. comm.}, chap.~V, \S2,
n\textsuperscript{o}~1, prop.~1)); s'il en était ainsi,
$B'=B\otimes_AA'$ aurait deux idéaux premiers distincts au-dessus de
$x$ (I, 3.4.7), et il suffit de voir que ce n'est pas possible.

Soit $\mathfrak N$ le nilradical de $B$; si $S$ est l'ensemble des
éléments réguliers de $B$, on a $R=S^{-1}B$, et
$S^{-1}\mathfrak N=R\mathfrak N$ est le nilradical de $R$, donc
\[
 R/R\mathfrak N=S^{-1}B/S^{-1}\mathfrak N=S^{-1}(B/\mathfrak N)
\]
est contenu dans le corps des fractions $L$ de l'anneau local intègre
$B/\mathfrak N=B_{\mathrm{red}}$; il en est donc de même de
\[
 B'_{\mathrm{red}}=B'/(B'\cap R\mathfrak N)
\]
qui est entier sur $B_{\mathrm{red}}$; par hypothèse
$B'_{\mathrm{red}}$, étant contenu dans la clôture intégrale de
$B_{\mathrm{red}}$, est un anneau local, donc il en est de même de
$B'$, ce qui établit notre assertion.

Si l'on suppose en outre $B$ géométriquement unibranche, le corps
résiduel de $B'_{\mathrm{red}}$, étant contenu dans le corps résiduel de
la clôture intégrale de $B_{\mathrm{red}}$, est une extension radicielle
du corps résiduel $k(x)$ de $B$. Notons maintenant que le corps résiduel
de $B'_{\mathrm{red}}$ est aussi celui de $B'=B\otimes_AA'$, donc est
égal à $k(x)\otimes_{k(y)}k'$, où $k'$ est le corps résiduel de $A'$;
on déduit alors de {\hyperref[canon-index.ega4.c04.prop.2-6-1]{(2.6.1)}} que $k'$
est une extension radicielle de $k(y)$, donc $A$ est aussi
géométriquement unibranche.

\paragraph{(Err\textsubscript{IV}, 33).}
Dans ({\hyperref[canon-index.ega4.c04.item.12-1-2]{IV, 12.1.2}}), l'exemple donné
en (i) n'est pas correct, le morphisme considéré étant ouvert. On
construit un exemple d'un morphisme plat $f:X\to Y$ dont la restriction
à une composante irréductible de $X$ n'est pas un morphisme ouvert de
la façon suivante. Soit $Z=\operatorname{Spec}(\mathbb C[U,V])$ le
plan affine sur le corps des nombres complexes, et considérons le
schéma $Y$ obtenu en recollant deux points fermés distincts $a$, $b$ de
$Z$; on a $Y=\operatorname{Spec}(B)$, où $B$ est le sous-anneau de
$\mathbb C[U,V]$ formé des polynômes prenant la même valeur aux points
$a$ et $b$ (l'ensemble des points fermés de $Z$ étant identifié à
$\mathbb C^2$); il est immédiat que $B$ est un sous-$\mathbb C$-espace
vectoriel de $\mathbb C[U,V]$ de codimension 1, donc $\mathbb C[U,V]$
est un $B$-module de type fini, et par suite le morphisme $Z\to Y$ est
fini, l'image réciproque du point $c$ de $Y$ obtenu par recollement de
$a$ et de $b$ étant l'ensemble $\{a,b\}$.

Considérons alors le schéma $X$ obtenu en recollant deux exemplaires de
$Z$, soient $Z_1$, $Z_2$, de sorte que si $a_i$, $b_i$ ($i=1,2$) sont
les points de $Z_i$ correspondant à $a$ et $b$, $a_1$ soit recollé à
$b_2$ et $b_1$ à $a_2$ (cf. ({\hyperref[canon-index.ega4.c04.item.6-5-5]{6.5.5}},
(ii))). On voit aisément (\emph{loc. cit.}) que le morphisme $X\to Y$
ainsi obtenu est plat et fini (et même étale) et que $Z_1$ et $Z_2$
s'identifient aux deux composantes irréductibles de $X$. Cependant, si
$D$ est une droite dans $Z_1$ passant par $a_1$ mais non par $b_1$,
l'image $V$ dans $Y$ de l'ouvert $Z_1-D$ de $Z_1$ contient l'image de
$b_1$, c'est-à-dire le point $c$, mais ne contient pas l'image du point
générique de $D$, donc $V$ n'est pas un voisinage de $c$.

\paragraph{(Err\textsubscript{IV}, 34).}
Dans le titre de (IV, 13.2), ligne 5 du bas de la p.~190, remplacer :
« équidimensionnels » par « localement équidimensionnels ». De même
dans (IV, 13.2.2), ligne 4 de la p.~191, remplacer deux fois
« équidimensionnel » par « localement équidimensionnel ».

\oldpage[IV]{357}
Ligne 5 de la p.~191, ajouter : \emph{Si de plus les fibres
$f^{-1}(f(x))$ de $f$ sont équidimensionnelles ($0_{\mathrm{IV}}$,
14.1.3), on dit que $f$ est équidimensionnel.} Ligne 6 du bas de la
p.~194, remplacer « équidimensionnelles » par « localement
équidimensionnelles ».

\paragraph{(Err\textsubscript{IV}, 35).}
Dans le titre de (IV, 13.3), ligne 4 du bas de la p.~94, remplacer
« équidimensionnels » par « localement équidimensionnels ». Dans
(IV, 13.3.2), ligne 6 du bas de la p.~196, remplacer deux fois
« équidimensionnel » par « localement équidimensionnel ». Ligne 5 du bas
de la p.~196, ajouter : \emph{Si de plus les fibres $f^{-1}(f(x))$ de
$f$ sont équidimensionnelles ($0_{\mathrm{IV}}$, 14.1.3), on dit que
$f$ est équidimensionnel.} Dans (IV, 13.3.7), ligne 14 de la p.~198,
remplacer « équidimensionnel » par « localement équidimensionnel ».
Dans ({\hyperref[canon-index.ega4.c04.statement.13-3-9]{IV, 13.3.9}}), ligne 3 du
bas de la p.~198, remplacer « équidimensionalité » par
« équidimensionalité locale ». Lignes 10, 12, 20 et 24 de la p.~199,
remplacer « équidimensionnel » par « localement équidimensionnel ».

\paragraph{(Err\textsubscript{IV}, 36).}
Dans l'introduction du \S14, ligne 6 du bas de la p.~199, remplacer
« équidimensionnel » par « localement équidimensionnel ».

\paragraph{(Err\textsubscript{IV}, 37).}
Dans (IV, 14.4.1), ligne 5 du bas de la p.~209, remplacer « localement
de type fini » par « localement de présentation finie ». L'énoncé
actuel de ({\hyperref[canon-index.ega4.c04.thm.14-4-1]{IV, 14.4.1}}) est inexact,
même lorsque $f:X\to Y$ est une immersion fermée, $X$ étant réduit à
un seul point $y$ (fermé) de $Y$, et $\sh{O}_{Y,y}$ étant un corps. Il
y a en effet des anneaux $A$ tels que $Y=\operatorname{Spec}(A)$ soit
compact, sans point isolé et totalement discontinu, tout idéal premier
$\mathfrak p$ de $A$ étant maximal et $A_\mathfrak p$ étant un corps
(Bourbaki, \emph{Alg. comm.}, chap.~II, \S4, exerc.~16 et 17); il est
clair que l'immersion fermée $\{y\}\to Y$ n'est pas un morphisme
ouvert, mais vérifie les conditions b) et b') de
({\hyperref[canon-index.ega4.c04.thm.14-4-1]{IV, 14.4.1}}).

\paragraph{(Err\textsubscript{IV}, 38).}
Dans (IV, 14.4.1.3), ligne 14 de la p.~211, remplacer « localement
quasi-fini et dominant » par « localement quasi-fini, localement de
présentation finie et dominant ».

\paragraph{(Err\textsubscript{IV}, 39).}
Dans (IV, 14.4.2), ligne 10 du bas de la p.~211, remplacer « localement
de type fini » par « localement de présentation finie ».

\paragraph{(Err\textsubscript{IV}, 40).}
Dans (IV, 14.4.3), ligne 18 de la p.~212, remplacer « localement de
type fini » par « localement de présentation finie ».

\paragraph{(Err\textsubscript{IV}, 41).}
Dans ({\hyperref[canon-index.ega4.c04.cor.14-4-4]{IV, 14.4.4}}), lignes 15 et 16 du
bas de la p.~212, remplacer « localement de type fini » par
« localement de présentation finie ». Lignes 10 et 11 du bas, remplacer
« équidimensionnel » par « localement équidimensionnel ».

\paragraph{(Err\textsubscript{IV}, 42).}
Dans (IV, 14.4.8), après la ligne 1 de la p.~214, ajouter :
\begin{quote}
\emph{Si de plus l'ensemble des composantes irréductibles de $Y$ est
localement fini, ces deux conditions sont équivalentes aux suivantes :}
\end{quote}

\paragraph{(Err\textsubscript{IV}, 43).}
Dans (IV, 14.4.10), ligne 15 de la p.~215, remplacer : quelconque, par :
dont l'ensemble des composantes irréductibles est localement fini.
Ligne 16 de la p.~215, remplacer : de type fini, par : de présentation
finie.

\paragraph{(Err\textsubscript{IV}, 44).}
Dans (IV, 15.5.6.1), ligne 15 du bas de la p.~234, après la première
phrase, ajouter :

\begin{quote}
On peut se limiter au cas où dans $X$ toute partie fermée irréductible
n'admet qu'un seul point générique. En effet, dans le cas contraire on
considère l'espace $Y=X/R$,

\oldpage[IV]{358}
quotient de $X$ par la relation d'équivalence
$\overline{\{x\}}=\overline{\{x'\}}$. Il résulte aussitôt de cette
définition que tout ouvert de $X$ est saturé pour $R$, donc la topologie
de $X$ est l'image réciproque de celle de $Y$ par l'application
canonique $\pi:X\to X/R$. Un ouvert dense dans $X$ est de la forme
$\pi^{-1}(U)$, où $U$ est un ouvert dense dans $Y$, donc un fermé rare
$Z$ dans $X$ est de la forme $\pi^{-1}(S)$, où $S=\pi(Z)$ est fermé
rare dans $Y$, et il revient au même de dire que $X-Z$ est connexe ou
que $Y-S$ est connexe. De même $\pi(T_x)$ est l'ensemble des
générisations de $\pi(x)$, et il est équivalent de dire que
$T_x-\{x\}$ est connexe ou que $\pi(T_x)-\{\pi(x)\}$ l'est. Enfin, par
construction $Y$ est un espace de Kolmogoroff localement noethérien,
dont toute partie fermée admet donc un seul point générique. On voit
alors qu'il revient au même de prouver le lemme pour $X$ ou pour $Y$,
et on peut donc supposer que toute partie fermée irréductible de $X$ a
un point générique unique.
\end{quote}

\paragraph{(Err\textsubscript{IV}, 45).}
Dans (IV, 15.6.1), ligne 18 du bas de la p.~236, supprimer « localement
noethérien » (hypothèse non utilisée dans la démonstration).

\paragraph{(Err\textsubscript{IV}, 46).}
Dans (IV, 15.6.3), ligne 11 de la p.~237, au lieu de « supposons que
pour tout $y\in Y$ », lire : « supposons que $f$ soit localement de
présentation finie et que pour tout $y\in Y$ ». Remplacer la
démonstration (lignes 14 à 24 de la p.~237) par la suivante :
\begin{quote}
En vertu de {\hyperref[canon-index.ega4.c04.prop.1-10-3]{(1.10.3)}}, il suffit de
prouver que pour tout $x\in X_y^0$ et toute générisation $y'$ de $y$,
il existe une générisation $x'$ de $x$ telle que $f(x')=y'$. Or il
suffit de prendre pour $x'$ le point générique de $X_{y'}^0$, en vertu
de {\hyperref[canon-index.ega4.c04.cor.15-6-2]{(15.6.2)}}.
\end{quote}

\paragraph{(Err\textsubscript{IV}, 47).}
Dans (IV, 15.6.7), ligne 2 de la p.~241, après « noethérien », ajouter :
« ou si la condition $\beta$) de
({\hyperref[canon-index.ega4.c04.prop.15-6-6]{15.6.6}}, (iii)) est vérifiée ».
Après la ligne 9 de la p.~241, ajouter :
\begin{quote}
Lorsque la condition $\beta$) de
({\hyperref[canon-index.ega4.c04.prop.15-6-6]{15.6.6}}, (iii)) est satisfaite, le
fait que a) entraîne b') résulte de
({\hyperref[canon-index.ega4.c04.prop.15-6-6]{15.6.6}}, (ii)), et le fait que b')
entraîne a) résulte de
({\hyperref[canon-index.ega4.c04.prop.15-6-6]{15.6.6}}, (iii)).
\end{quote}

\paragraph{(Err\textsubscript{IV}, 48).}
Dans (IV, 20.1.5), ligne 13 de la p.~228, (IV, 20.1.6), ligne 17 de la
p.~228, (IV, 20.1.7), ligne 13 du bas de la p.~228, supprimer
« strictement ».

\paragraph{(Err\textsubscript{IV}, 49).}
Dans (IV, 20.2.13), ajouter après la ligne 6 de la p.~236 : En
particulier si $X$ est réduit et n'a qu'un nombre fini de composantes
irréductibles, la notion de $\sh{O}_X$-Module sans torsion définie dans
({\hyperlink{ega4.c04.statement.20-1-5}{IV, 20.1.5}}) coïncide avec
celle de (I, 7.4.1).

\paragraph{(Err\textsubscript{IV}, 50).}
Dans (IV, 20.2.14), lignes 18 et 25 de la p.~236, supprimer
« strictement ».

\paragraph{(Err\textsubscript{IV}, 51).}
Dans (IV, 20.6.2), ligne 13 du bas de la p.~252, supprimer
« strictement ».

\paragraph{(Err\textsubscript{IV}, 52).}
Dans (IV, 20.6.4), lignes 23 et 22 de la p.~253, supprimer
« strictement ».

\paragraph{(Err\textsubscript{IV}, 53).}
Dans (IV, 21.4), après la ligne 23 de la p.~267, ajouter :

La proposition suivante et ses corollaires nous ont été signalés par
M.~Raynaud.

\begin{proposition}[21.4.9]
\EGAReaderTarget{ega4.c04.prop.21-4-9}{IV.21.4.9}
Soient $X$, $Y$ deux préschémas localement noethériens,
$f:X\to Y$ un morphisme fidèlement plat; on suppose en outre, soit que
$f$ est quasi-compact, soit que $f$ est ouvert. Pour qu'un diviseur
$D'$ sur $X$ soit de la forme $f^*(D)$, où $D$ est un diviseur sur
$Y$, il faut et il suffit que pour tout $y\in Y$ tel que
$\operatorname{prof}(\sh{O}_{Y,y})\leq1$, en posant
\[
 Y_1=\operatorname{Spec}(\sh{O}_{Y,y}),\qquad
 X_1=X\times_YY_1,\qquad
 f_1=f_{(Y_1)}:X_1\to Y_1,
\]
l'image réciproque $D'_1$ de $D'$ sur $X_1$ soit de la forme
$f_1^*(D_1)$, où $D_1\in\operatorname{Div}(Y_1)$.
\end{proposition}

\oldpage[IV]{359}
\begin{proof}
La nécessité de la condition est évidente par transitivité des images
réciproques de diviseurs, tous les morphismes considérés étant plats
{\hyperlink{ega4.c04.statement.21-4-4}{(21.4.4)}}. Pour tout ouvert
$U$ de $Y$, notons $f_U$ la restriction $f^{-1}(U)\to U$ de $f$; en
vertu de {\hyperlink{ega4.c04.prop.21-4-8}{(21.4.8)}} l'application
\[
 D_U\longmapsto f_U^*(D_U)
\]
de $\operatorname{Div}(U)$ dans $\operatorname{Div}(f^{-1}(U))$ est
injective; donc, si $U_1$, $U_2$ sont deux ouverts de $Y$, et si
\[
 D'|f^{-1}(U_1)=f_{U_1}^*(D_1),\qquad
 D'|f^{-1}(U_2)=f_{U_2}^*(D_2),
\]
on a nécessairement
$D_1|U_1\cap U_2=D_2|U_1\cap U_2$. Ceci montre aussitôt qu'il existe
un plus grand ouvert $U\subset Y$ tel que
\[
 D'|f^{-1}(U)=f_U^*(D_U),
 \qquad D_U\in\operatorname{Div}(U);
\]
il s'agit de voir que $U=Y$. Supposons le contraire, et soit $y$ un
point maximal de $Y-U$; on va montrer qu'il existe un voisinage ouvert
$V$ de $y$ dans $Y$ tel que
$D'|f^{-1}(V)=f_V^*(D_V)$ pour un $D_V\in\operatorname{Div}(V)$, ce qui
contredira le caractère maximal de $U$.

\textup{I) Cas où $f$ est quasi-compact.} Appliquant le fait que $f$
est quasi-compact et {\hyperlink{ega4.c04.prop.20-3-8}{(20.3.8)}}
suivant la méthode de ({\hyperref[canon-index.ega4.c04.statement.8-1-2]{8.1.2}},
a)), on est ramené à prouver qu'avec les notations de l'énoncé (mais
sans supposer d'abord de limitation sur
$\operatorname{prof}(\sh{O}_{Y,y})$), $D'_1$ est de la forme
$f_1^*(D_1)$. Or, si $\operatorname{prof}(\sh{O}_{Y,y})\leq1$, cette
assertion n'est autre que l'hypothèse. Supposons donc
$\operatorname{prof}(\sh{O}_{Y,y})\geq2$; puisque $y$ est point maximal
de $Y-U$, l'ouvert $W=Y_1\cap U$ de $Y_1$ est égal à $Y_1-\{y\}$, et
on peut se borner au cas où $Y=Y_1$, $U=Y-\{y\}$. Par hypothèse, il
existe un diviseur $D_U\in\operatorname{Div}(U)$ tel que
$f_U^*(D_U)=D'|f^{-1}(U)$; $D_U$ est classe d'équivalence d'un couple
$(\sh{L}_U,s)$, où $\sh{L}_U$ est un $\sh{O}_U$-Module inversible et
$s$ une section méromorphe régulière de $\sh{L}_U$ au-dessus de $U$
{\hyperlink{ega4.c04.prop.21-2-11}{(21.2.11)}}; par suite
{\hyperlink{ega4.c04.statement.21-4-1}{(21.4.1)}},
$D'|f^{-1}(U)$ est classe d'équivalence du couple
$(\sh{L}'_U,s')$, où
\[
 \sh{L}'_U=f_U^*(\sh{L}_U),\qquad s'=s\circ f_U.
\]
Notons $i:U\to Y$, $j:f^{-1}(U)\to X$ les injections canoniques, et
posons
\[
 \sh{L}=i_*(\sh{L}_U),\qquad
 \sh{L}'=j_*(\sh{L}'_U).
\]
Comme $f$ est plat, $f^*(\sh{L})$ est canoniquement isomorphe à
$\sh{L}'$ {\hyperref[canon-index.ega4.c04.lem.2-3-1]{(2.3.1)}}; d'autre part,
puisque $f$ est plat, la relation
$\operatorname{prof}(\sh{O}_{Y,y})\geq2$ entraîne
$\operatorname{prof}(\sh{O}_{X,x})\geq2$ pour tout $x\in f^{-1}(y)$
{\hyperref[canon-index.ega4.c04.prop.6-3-1]{(6.3.1)}}.

Or, $D'$ est classe d'équivalence d'un couple $(\sh{L}'',s'')$, où
$\sh{L}''$ est un $\sh{O}_X$-Module inversible et $s''$ une section
méromorphe régulière de $\sh{L}''$ au-dessus de $X$; $\sh{L}'_U$ est
par suite isomorphe à $\sh{L}''|f^{-1}(U)$, et il résulte alors de ce
qui précède et de {\hyperref[canon-index.ega4.c04.prop.5-9-8]{(5.9.8)}} et
{\hyperref[canon-index.ega4.c04.thm.5-10-5]{(5.10.5)}} que $\sh{L}'$ est un
$\sh{O}_X$-Module inversible isomorphe à $\sh{L}''$. Puisque $f$ est
fidèlement plat et quasi-compact, on en conclut que $\sh{L}$ est un
$\sh{O}_Y$-Module inversible
{\hyperref[canon-index.ega4.c04.prop.2-5-2]{(2.5.2)}} et par suite isomorphe à
$\sh{O}_Y$, puisque $Y$ est un schéma local. D'ailleurs, puisque
$\operatorname{prof}(\sh{O}_{Y,y})\geq1$, $U=Y-\{y\}$ est
schématiquement dense dans $Y$
({\hyperlink{ega4.c04.statement.20-2-13}{20.2.13}}, (iv)) et par suite
la section méromorphe régulière $s$ de $\sh{L}_U$ au-dessus de $U$ se
prolonge en une section méromorphe régulière $s_0$ de $\sh{L}$
au-dessus de $Y$ {\hyperlink{ega4.c04.prop.20-2-11}{(20.2.11)}}; la
classe d'équivalence de $(\sh{L},s_0)$ est donc un diviseur $D$ sur
$Y$; et il est clair que $f^*(D)$ et $D'$ induisent le même diviseur sur
$f^{-1}(U)=X-f^{-1}(y)$; mais comme
$\operatorname{prof}(\sh{O}_{X,x})\geq2$ pour tout $x\in f^{-1}(y)$,
on a $D'=f^*(D)$ {\hyperlink{ega4.c04.prop.21-1-8}{(21.1.8)}}, ce qui
achève la démonstration dans ce cas.

\textup{II) Cas où $f$ est ouvert.} Comme il s'agit de prouver
l'existence d'un voisinage ouvert $V$ de $y$ dans $Y$ tel que
$D'|f^{-1}(V)=f_V^*(D_V)$ pour un $D_V\in\operatorname{Div}(V)$, on
peut déjà supposer $Y$ affine. Puisque $f$ est ouvert, les images
$f(W)$, où $W$ parcourt les ouverts affines de $X$, forment un
recouvrement ouvert de $Y$, donc, puisque $Y$ est quasi-compact, il
existe un ouvert quasi-compact $T\subset X$ tel que $f(T)=Y$, et le
morphisme $f|T$ est quasi-compact (I, 1.1.1); \emph{a fortiori}, pour
tout ouvert quasi-compact $T_\alpha\supset T$ de $X$,

\oldpage[IV]{360}
$f(T_\alpha)=Y$ et $f|T_\alpha$ est quasi-compact. On peut par suite
appliquer à $f|T_\alpha$ le résultat de I), et il existe un diviseur
$D_\alpha$ sur $Y$ tel que
\[
 (f|T_\alpha)^*(D_\alpha)=D'|T_\alpha;
\]
en vertu de la transitivité des images réciproques de diviseurs par
morphismes plats {\hyperlink{ega4.c04.statement.21-4-4}{(21.4.4)}}, on
a aussi $(f|T)^*(D_\alpha)=D'|T$. On en conclut que pour deux indices
$\alpha$, $\beta$ distincts, on a
\[
 (f|T)^*(D_\alpha)=(f|T)^*(D_\beta),
\]
d'où (par {\hyperlink{ega4.c04.prop.21-4-8}{(21.4.8)}}),
$D_\alpha=D_\beta$. Si $D$ est la valeur commune des $D_\alpha$, on a
donc $f^*(D)|T_\alpha=D'|T_\alpha$ pour tout $\alpha$, et par suite
$f^*(D)=D'$.
\end{proof}

\begin{corollary}[21.4.10]
\EGAReaderTarget{ega4.c04.cor.21-4-10}{IV.21.4.10}
Supposons vérifiées les hypothèses générales de
{\hyperlink{ega4.c04.prop.21-4-9}{(21.4.9)}}, et en outre supposons que
$Y$ soit normal et que pour tout $y\in Y$ tel que
$\dim(\sh{O}_{Y,y})=1$, la fibre $X_y=f^{-1}(y)$ n'ait pas de cycle
premier associé immergé. Pour qu'un diviseur $D'$ sur $X$ soit de la
forme $f^*(D)$, où $D\in\operatorname{Div}(Y)$, il faut et il suffit
que les deux conditions suivantes soient vérifiées :
\begin{enumerate}[label=\textup{\arabic*\textdegree}]
\item pour tout point maximal $y$ de $Y$, l'image réciproque de $D'$
dans $\operatorname{Div}(X_y)$ soit nulle;
\item pour tout $y\in Y$ tel que $\dim(\sh{O}_{Y,y})=1$, il existe un
entier $n_y$ tel que, pour tout point maximal $x$ de $X_y$, l'image
réciproque de $D'$ dans
$\operatorname{Div}(\operatorname{Spec}(\sh{O}_{X,x}))$ soit égale à
$\operatorname{div}(t_y^{n_y})$, où $t_y$ est l'image par
l'homomorphisme $\sh{O}_{Y,y}\to\sh{O}_{X,x}$ d'une uniformisante de
l'anneau de valuation discrète $\sh{O}_{Y,y}$.
\end{enumerate}
\end{corollary}

\begin{proof}
Appliquons le critère {\hyperlink{ega4.c04.prop.21-4-9}{(21.4.9)}}.
Aux points maximaux $y$ de $Y$, on a $\operatorname{Div}(Y_1)=0$
puisque $\sh{O}_{Y,y}$ est artinien, et la condition 1\textdegree{} de
(21.4.10) est équivalente au critère de (21.4.9) en un tel point. Il
faut voir que les conditions 1\textdegree{} et 2\textdegree{} impliquent
de même la condition de (21.4.9) en un point $y\in Y$ où
$\operatorname{prof}(\sh{O}_{Y,y})=1$; comme $Y$ est normal, ces points
sont aussi ceux où $\dim(\sh{O}_{Y,y})=1$, ou encore où
$\sh{O}_{Y,y}$ est un anneau de valuation discrète
{\hyperref[canon-index.ega4.c04.thm.5-8-6]{(5.8.6)}}. On peut donc se borner au
cas où $Y=\operatorname{Spec}(\sh{O}_{Y,y})$. Si $t$ est une
uniformisante de $\sh{O}_{Y,y}$, l'application
$n\mapsto\operatorname{div}(t^n)$ est un isomorphisme de $\mathbb Z$
sur $\operatorname{Div}(Y)$; la condition de (21.4.9) signifie que
$D'$ est de la forme $\operatorname{div}(t^n\circ f)$ pour un
$n\in\mathbb Z$. Cela montre aussitôt que cette condition entraîne la
condition 2\textdegree{} de l'énoncé. Inversement, supposons cette
dernière vérifiée (ainsi que la condition 1\textdegree{}); posant pour
simplifier $n=n_y$, tout revient à voir que l'on a alors
$D'=\operatorname{div}(t^n\circ f)$. Si $\eta$ est le point générique
de $Y$, les germes de $D'$ et de $\operatorname{div}(t^n\circ f)$ en
un point quelconque de $X_\eta$ sont tous deux nuls en vertu de
1\textdegree{}. De même, les germes de $D'$ et
$\operatorname{div}(t^n\circ f)$ en un point maximal de $X_y$ sont
égaux par la condition 2\textdegree{}. Mais en un point $x\in X_y$ qui
n'est pas maximal, on a
$\operatorname{prof}(\sh{O}_{X_y,x})\geq1$ puisque $X_y$ n'a pas de
cycle premier associé immergé, donc aussi
$\operatorname{prof}(\sh{O}_{X,x})\geq2$ par
{\hyperref[canon-index.ega4.c04.prop.6-3-1]{(6.3.1)}}. On a donc
$D'=\operatorname{div}(t^n\circ f)$ en vertu de
{\hyperlink{ega4.c04.prop.21-1-8}{(21.1.8)}}.
\end{proof}

\begin{corollary}[21.4.11]
\EGAReaderTarget{ega4.c04.cor.21-4-11}{IV.21.4.11}
Supposons vérifiées les hypothèses générales de
{\hyperlink{ega4.c04.prop.21-4-9}{(21.4.9)}}, et supposons en outre que
$Y$ soit normal et que pour tout $y\in Y$ tel que
$\dim(\sh{O}_{Y,y})=1$, la fibre $X_y$ soit intègre. Alors, pour qu'un
diviseur $D'$ soit de la forme $f^*(D)$, où
$D\in\operatorname{Div}(Y)$, il faut et il suffit que, pour tout point
maximal $y$ de $Y$, l'image réciproque de $D'$ dans
$\operatorname{Div}(X_y)$ soit nulle.
\end{corollary}

\begin{proof}
Il suffit de vérifier que la condition 2\textdegree{} de
{\hyperlink{ega4.c04.cor.21-4-10}{(21.4.10)}} est alors
automatiquement vérifiée. Or, si $\dim(\sh{O}_{Y,y})=1$ et si $X_y$ est
intègre, on a, en l'unique point maximal $x$ de $X_y$,
$\dim(\sh{O}_{X,x})=1$ en vertu de
{\hyperref[canon-index.ega4.c04.cor.6-1-3]{(6.1.3)}}; en outre, comme
$\sh{O}_{Y,y}$ est un anneau régulier et $\sh{O}_{X_y,x}$ un corps,
$\sh{O}_{X,x}$ est régulier
{\hyperref[canon-index.ega4.c04.prop.6-5-1]{(6.5.1)}}, donc un anneau de valuation
discrète, et si $t$ est une uniformisante de $\sh{O}_{Y,y}$, son image
$t_y$ dans $\sh{O}_{X,x}$ est une uniformisante de $\sh{O}_{X,x}$
($0$, 17.3.3); l'image réciproque de $D'$ dans
$\operatorname{Div}(\operatorname{Spec}(\sh{O}_{X,x}))$ est donc bien
de la forme

\oldpage[IV]{361}
$\operatorname{div}(t_y^n)$, ce qui n'est autre que la condition
2\textdegree{} de {\hyperlink{ega4.c04.cor.21-4-10}{(21.4.10)}},
puisque $X_y$ n'a qu'un seul point maximal.
\end{proof}

\begin{remark}[21.4.12]
\label{IV.21.4.12}
Si dans {\hyperlink{ega4.c04.cor.21-4-11}{(21.4.11)}} on suppose
seulement que les $X_y$ sont réduits, il faut ajouter à la condition de
l'énoncé que les images réciproques de $D'$ dans chacun des ensembles
$\operatorname{Div}(\operatorname{Spec}(\sh{O}_{X,x}))$, où $x$
parcourt l'ensemble des points maximaux de $X_y$, sont toutes de la
forme $\operatorname{div}((t'_x)^n)$ où $t'_x$ est une uniformisante
de $\sh{O}_{X,x}$, l'entier $n$ étant le même pour tous ces points.
\end{remark}

\begin{corollary}[21.4.13]
\label{IV.21.4.13}
Supposons vérifiées les conditions de
{\hyperlink{ega4.c04.cor.21-4-11}{(21.4.11)}}. Pour tout
$\sh{O}_X$-Module inversible $\sh{L}$ tel que, pour tout point maximal
$\eta$ de $Y$, l'image réciproque
\[
 \sh{L}_\eta=\sh{L}\otimes_{\sh{O}_X}\sh{O}_{X_\eta}
\]
soit isomorphe à $\sh{O}_{X_\eta}$, il existe un
$\sh{O}_Y$-Module inversible $\sh{M}$ et un isomorphisme de
$f^*(\sh{M})$ sur $\sh{L}$.
\end{corollary}

\begin{proof}
Il suffit de montrer qu'il existe une section méromorphe régulière $s$
de $\sh{L}$ au-dessus de $X$ telle que le diviseur $D'$ correspondant
à $(\sh{L},s)$ soit tel que $D'_x=0$ en tout point $x$ d'une fibre
$X_\eta$ d'un point maximal $\eta$ de $Y$. On appliquera alors
{\hyperlink{ega4.c04.cor.21-4-11}{(21.4.11)}} à $D'$, et la conclusion
résultera de {\hyperlink{ega4.c04.def.21-4-2}{(21.4.2)}}.

Notons que les points de $\operatorname{Ass}(\sh{O}_X)$ sont au-dessus
des points maximaux de $Y$ {\hyperref[canon-index.ega4.c04.prop.3-3-1]{(3.3.1)}};
comme $Y$ est réduit et localement noethérien, il y a un ouvert dense
$V$ de $Y$, réunion d'une famille disjointe d'ouverts affines contenant
chacun un seul point maximal de $Y$, et $f^{-1}(V)$ est
schématiquement dense dans $X$
{\hyperref[canon-index.ega4.c04.prop.5-10-2]{(5.10.2)}}; on peut donc se borner au
cas où $Y$ est affine et intègre. Puisque, pour l'unique point maximal
$\eta$ de $Y$, $X_\eta$ est alors intègre,
$\operatorname{Ass}(\sh{O}_X)$ est réduit à un seul point $z$, qui est
le point générique de $X$. Soit $s_\eta$ la section de $\sh{L}_\eta$
au-dessus de $X_\eta$ qui correspond par isomorphie à la section unité
de $\sh{O}_{X_\eta}$, et soient $U$ un voisinage ouvert affine d'un
point $x\in X_\eta$ dans $X$ et $s$ une section de $\sh{L}$ au-dessus
de $U$ ayant même germe que $s_\eta$ au point $x$. Comme $U$ est
schématiquement dense dans $X$
({\hyperlink{ega4.c04.statement.20-2-13}{20.2.13}}, (iv)), $s$
s'identifie à une section méromorphe régulière bien déterminée de
$\sh{L}$ au-dessus de $X$, et il est clair
({\hyperlink{ega4.c04.prop.20-2-11}{20.2.11}} et
{\hyperlink{ega4.c04.prop.20-3-5}{20.3.5}}) que $s$ induit sur $X_\eta$
une section méromorphe régulière de $\sh{L}_\eta$ égale à $s_\eta$, et
répond donc à la question.
\end{proof}

\bigskip
\begin{center}
\small
1967. --- Imprimerie des Presses Universitaires de France. --- Vendôme
(France)\\
ÉDIT. N\textsuperscript{o} 29 470\qquad
IMPRIMÉ EN FRANCE\qquad
IMP. N\textsuperscript{o} 20 298
\end{center}

% FIN SOURCE ega4/ega4-errata-addenda-list3-fr.tex

\end{document}
% EGA-COMPONENT-END IV4

